% Options for packages loaded elsewhere
\PassOptionsToPackage{unicode}{hyperref}
\PassOptionsToPackage{hyphens}{url}
\PassOptionsToPackage{dvipsnames,svgnames,x11names}{xcolor}
%
\documentclass[
  10pt,
]{article}
\usepackage{amsmath,amssymb}
\usepackage{iftex}
\ifPDFTeX
  \usepackage[T1]{fontenc}
  \usepackage[utf8]{inputenc}
  \usepackage{textcomp} % provide euro and other symbols
\else % if luatex or xetex
  \usepackage{unicode-math} % this also loads fontspec
  \defaultfontfeatures{Scale=MatchLowercase}
  \defaultfontfeatures[\rmfamily]{Ligatures=TeX,Scale=1}
\fi
\usepackage{lmodern}
\ifPDFTeX\else
  % xetex/luatex font selection
  \setmainfont[]{DejaVu Serif}
  \setsansfont[]{DejaVu Sans}
  \setmonofont[]{DejaVu Sans Mono}
  \setmathfont[]{Latin Modern Math}
\fi
% Use upquote if available, for straight quotes in verbatim environments
\IfFileExists{upquote.sty}{\usepackage{upquote}}{}
\IfFileExists{microtype.sty}{% use microtype if available
  \usepackage[]{microtype}
  \UseMicrotypeSet[protrusion]{basicmath} % disable protrusion for tt fonts
}{}
\makeatletter
\@ifundefined{KOMAClassName}{% if non-KOMA class
  \IfFileExists{parskip.sty}{%
    \usepackage{parskip}
  }{% else
    \setlength{\parindent}{0pt}
    \setlength{\parskip}{6pt plus 2pt minus 1pt}}
}{% if KOMA class
  \KOMAoptions{parskip=half}}
\makeatother
\usepackage{xcolor}
\usepackage[margin=22mm]{geometry}
\usepackage{longtable,booktabs,array}
\usepackage{calc} % for calculating minipage widths
% Correct order of tables after \paragraph or \subparagraph
\usepackage{etoolbox}
\makeatletter
\patchcmd\longtable{\par}{\if@noskipsec\mbox{}\fi\par}{}{}
\makeatother
% Allow footnotes in longtable head/foot
\IfFileExists{footnotehyper.sty}{\usepackage{footnotehyper}}{\usepackage{footnote}}
\makesavenoteenv{longtable}
\setlength{\emergencystretch}{3em} % prevent overfull lines
\providecommand{\tightlist}{%
  \setlength{\itemsep}{0pt}\setlength{\parskip}{0pt}}
\setcounter{secnumdepth}{-\maxdimen} % remove section numbering
\usepackage{amsmath,amssymb,amsthm}
\usepackage{fvextra}
\usepackage{newunicodechar}
\providecommand{\bigboxtimes}{\mathop{\boxtimes}\displaylimits}
\providecommand{\PsiAt}{\operatorname{PsiAt}}
\newunicodechar{𝓞}{\ensuremath{\mathcal O}}
\newunicodechar{ℓ}{\ensuremath{\ell}}
\newunicodechar{∎}{\ensuremath{\blacksquare}}
\newunicodechar{⋯}{\ensuremath{\cdots}}
\newunicodechar{⋊}{\ensuremath{\rtimes}}
\newunicodechar{𝔣}{\ensuremath{\mathfrak f}}
\newunicodechar{𝜎}{\ensuremath{\sigma}}
\newunicodechar{ℒ}{\ensuremath{\mathcal L}}
\newunicodechar{〈}{\ensuremath{\langle}}
\newunicodechar{〉}{\ensuremath{\rangle}}
\DefineVerbatimEnvironment{verbatim}{Verbatim}{breaklines,breakanywhere,fontsize=\small}
\setlength{\emergencystretch}{3em}
\setcounter{tocdepth}{2}
\ifLuaTeX
  \usepackage{selnolig}  % disable illegal ligatures
\fi
\usepackage{bookmark}
\IfFileExists{xurl.sty}{\usepackage{xurl}}{} % add URL line breaks if available
\urlstyle{same}
\hypersetup{
  pdftitle={ALP\_MACHINE: Mixed CM theta sections, lift costs, and arithmetic comparison gates},
  pdfauthor={GPT-6 Astra Pro research team},
  colorlinks=true,
  linkcolor={blue},
  filecolor={Maroon},
  citecolor={Blue},
  urlcolor={blue},
  pdfcreator={LaTeX via pandoc}}

\title{ALP\_MACHINE: Mixed CM theta sections, lift costs, and arithmetic comparison gates}
\author{GPT-6 Astra Pro research team}
\date{23 November 2026 --- active research version}

\begin{document}
\maketitle

\section{Status and conventions}\label{status-and-conventions}

The requested shorter actual mixed-type lift chain, universal negative lift-cost theorem, arithmetic MainX pairing value, and Kummer Hecke RH theorem have not been established. This is the integrated active paper, containing the full proofs written so far. Checked findings and work awaiting whole-proof review are separated below. The dated register gives the exact status of each original question and each completed subproblem.

Alpöge's theorem and archive are accepted as stipulated. No Lean build or independent validity audit of that theorem is claimed. Settled prior findings retain the supplied intake statuses. A literature theorem is an input, not a newly proved finding. All new proofs have stable identifiers; the event log records who read a complete proof and distinguishes those readings from numerical diagnostics and exact computations.

The notation \nolinkurl{M} in the source arithmetic Level is distinct from the number of iterated lifts when both occur in a discussion. The actual Level requires an odd integer. Three theta germs must also be kept distinct: the value-one even-characteristic germ F1 on the D135 line, its normalized Mumford translate Sigma\_tr on the trace-one line, and the trace-one germ sigma\_q with quadratic term u1*u3-u2\^{}2. AM-R-006 computes the scalar between the latter two. Later findings refine the limitations recorded in earlier contributions; none silently replaces a selected automorphic image by an unrestricted formal function.

ALP-LIFT-01 and AM-R-001 are two independently written versions of the same elementary changed-family accounting identity. They are retained with their check histories but count as one finding in novelty summaries.

\tableofcontents
\clearpage
\part{Checked findings and their full proofs}

\clearpage

\section{ALP\_MACHINE root contribution: actual comparison and endpoint normalization}\label{alp_machine-root-contribution-actual-comparison-and-endpoint-normalization}

Author: GPT-6 Astra Pro, \nolinkurl{/root}; reasoning effort is not exposed by the runtime. Run: ALP\_MACHINE, 21 November 2026. This is an application of accepted DP\_MULT results and directly read source definitions. It is not a Lean execution or a reproof of Alpöge's theorem.

\subsection{AM-R-001. Comparison of two rebuilt families}\label{am-r-001.-comparison-of-two-rebuilt-families}

Let \(\mathscr D\) and \(\mathscr D'\) be two fully specified families of lengths \(M\) and \(M'<M\). Fix corresponding family members, but do not suppose their endpoints or coefficient fields are equal. Use DP44's scalar definitions on each family, with primes denoting the second family. Thus

\[
 Q_M=\sum_{j=1}^{M}(g_j-e_j)-T_\tau-MH+K_u,
\qquad
 Q'_{M'}=\sum_{j=1}^{M'}(g'_j-e'_j)-T'_\tau-M'H'+K'_{u'}.
\]

Then the exact difference is

\[
\begin{split}
Q'_{M'}-Q_M={}&-\sum_{j=M'+1}^{M}(g_j-e_j-H)\\
&+\sum_{j=1}^{M'}\bigl((g'_j-g_j)-(e'_j-e_j)\bigr)
-M'(H'-H)\\
&-(T'_\tau-T_\tau)+(K'_{u'}-K_u).
\end{split}\tag{R1}
\]

If the compared quantity is \(Q-B\), add \(B_M-B'_{M'}\) to the right side. No sign of \(H\), or equality between the two period normalizations, is assumed.

\textbf{Proof.} Subtract the two displayed defining expressions and split the unprimed sum into retained and deleted indices. The unprimed main-term difference is \((M-M')H-M'(H'-H)\). Combining its first term with the deleted sum gives (R1). This short proof gives the changed-family form that DP44 did not write as an equation; the frozen-family formula DP44.3 is carried, not rederived as a new finding.

\textbf{Application and limitation.} A proposed calculation of a lift saving must specify every term of (R1). The supplied curve, prime, formal-group slopes and generic rank records do not supply those numbers. \nolinkurl{SizeData₂.gain}, \nolinkurl{SizeData₂.E}, \nolinkurl{SizeData₂.tau}, the weights and level are additional source data tied to a family through the global links. \nolinkurl{LinkG_at} includes a coarse upper bound as well as the ramp inequality; that ramp inequality is only a lower certificate and cannot be reversed into a sharp upper bound on true gains. Thus the local slope upper obstructions in DP42/DP51 do not supply the missing upper bound in (R1). This is a scope statement about the available implications, not a theorem that no shorter family exists. In particular, even a negative difference alone is not equivalent to failure of the full global theorem's inequalities.

\textbf{Locators.} DP\_MULT DP44(d), equations DP44.1--DP44.4; original source \nolinkurl{Formalism/LocalRingCardinality/DataParams₂.lean}, \nolinkurl{SizeData₂} lines 76--181, \nolinkurl{sz} 203--206, \nolinkurl{gainD} 212--213, \nolinkurl{Ga} 220, \nolinkurl{Ee} 222, \nolinkurl{logOm} 224, \nolinkurl{MAIN} 232--234; \nolinkurl{Formalism/SharpRestriction/GlobalTheoremInputs.lean}, \nolinkurl{LinkDIV_at} 77--89 and \nolinkurl{LinkG_at} 122--136. Source bytes were recovered from the intake-pinned mirror and matched to the intake's original-author archive SHA256 values.

\subsection{AM-R-002. The arithmetic and archimedean cost of an endpoint detector}\label{am-r-002.-the-arithmetic-and-archimedean-cost-of-an-endpoint-detector}

This statement makes quantitative the distinction retained in DP55 between a nonzero complex pairing and an integral, controlled detector. It is an elementary application of the exact source pairing, not a claim that ideal content or dual norms are new general concepts.

Let \(K\) be a number field, \(\mathcal R\) a full projective \(\mathcal O_K\)-lattice in the row coefficient space \(W\), and \(B:V\times W\to K\) the algebraic endpoint pairing, including its stabilizer and coefficient normalization. Let \(0\ne f\in V(K)\), and suppose the functional \(\ell_f=B(f,-)\) is nonzero. Define its nonzero fractional content ideal

\[
 I_f=\ell_f(\mathcal R)\subset K.\tag{R2}
\]

At a finite place \(\mathfrak p\), write \(I_f\mathcal O_{K,\mathfrak p}=\pi_{\mathfrak p}^{s_{\mathfrak p}}\mathcal O_{K,\mathfrak p}\), with the valuation normalized by \(v_{\mathfrak p}(\pi_{\mathfrak p})=1\). Then:

\begin{enumerate}
\def\labelenumi{\arabic{enumi}.}
\tightlist
\item
  There is a detector \(z\in\mathcal R_{\mathfrak p}\) with \(B(f,z)=1\) if and only if \(s_{\mathfrak p}\le0\).
\item
  The least integer \(d\ge0\) for which such a detector exists in \(\pi_{\mathfrak p}^{-d}\mathcal R_{\mathfrak p}\) is \[d_{\min}(f,\mathfrak p)=\max\{0,s_{\mathfrak p}\}.\tag{R3}\]
\item
  If \(r_1,\ldots,r_N\) is a local lattice basis and \(b_i=B(f,r_i)\), then \[s_{\mathfrak p}=\min_i v_{\mathfrak p}(b_i).\tag{R4}\]
\end{enumerate}

For a fixed embedding \(\sigma:K\hookrightarrow\mathbb C\), put a positive Hermitian norm on \(W_\sigma\), and let \(\beta_\sigma(f)=\|\ell_{f,\sigma}\|_*\). Then

\[
 \inf_{B_\sigma(f,z)=1}\|z\|_\sigma=\beta_\sigma(f)^{-1},\tag{R5}
\]

and the infimum is attained. If the norm on rows is transported through DP55's representing map \(J\), then \(\beta_\sigma(f)=\|f\|_{\rm cl,\sigma}\), provided that transported form is a norm; alternatively apply (R5) to the quotient by \(\ker J\).

For \(a\in K^\times\),

\[
 I_{af}=aI_f,\qquad \beta_\sigma(af)=|\sigma(a)|\beta_\sigma(f).\tag{R6}
\]

Consequently the signed normalization balance

\[
 \mathcal B(f)=\frac1{[K:\mathbb Q]}
 \left(\sum_{\mathfrak p}s_{\mathfrak p}\log N\mathfrak p
       -\sum_{\sigma:K\hookrightarrow\mathbb C}\log\beta_\sigma(f)\right)\tag{R7}
\]

is invariant under \(f\mapsto af\). All complex embeddings are counted; this convention automatically counts a conjugate pair twice. It is the negative Arakelov degree of the quotient image line \(I_f\), equipped with the quotient norm \(\|1\|_\sigma=\beta_\sigma(f)^{-1}\).

\textbf{Full proof.} Since \(\mathcal R\) is a finite lattice, its image under a nonzero \(K\)-linear map is a nonzero finitely generated \(\mathcal O_K\)-submodule of \(K\), hence a fractional ideal. Localization commutes with its image. The image of \(\pi^{-d}\mathcal R_{\mathfrak p}\) is exactly \(\pi^{s_{\mathfrak p}-d}\mathcal O_{K,\mathfrak p}\); it contains 1 exactly when \(d\ge s_{\mathfrak p}\). This proves (R3), including negative \(s_{\mathfrak p}\), and assertion 1. A basis generates the image ideal, and a finitely generated ideal over a DVR has valuation equal to the least valuation of its generators, proving (R4).

By the finite-dimensional Riesz representation theorem, there is a unique nonzero vector \(y\in W_\sigma\) such that \(\ell_{f,\sigma}(z)=\langle y,z\rangle\), in the convention conjugate-linear in the first argument. Its norm is \(\beta_\sigma(f)\). Cauchy--Schwarz gives \(1\le\beta_\sigma(f)\|z\|\) for every detector. The vector \(z=y/\|y\|^2\) is a detector and attains equality. This proves (R5). DP55's identity \(B(f,z)=\langle Jz,f\rangle_{\rm cl}\), with \(J\) conjugate-linear and onto, gives the stated transported-norm specialization by the same argument; if necessary take its injective quotient first.

Bilinearity gives the ideal equality in (R6), and the definition of operator norm gives the norm equality. The change in the numerator of (R7) is \[\sum_{\mathfrak p}v_{\mathfrak p}(a)\log N\mathfrak p-\sum_\sigma\log|\sigma(a)|=0\] by the product formula. Finally, relative to the lattice \(I_f\), the rational section 1 has divisor coefficient \(-s_{\mathfrak p}\); its quotient norm is \(\beta_\sigma(f)^{-1}\). Substituting these two facts into the definition of arithmetic degree gives \(-\mathcal B(f)\). This proves every assertion.

\subsubsection{The exact source coefficient matrix}\label{the-exact-source-coefficient-matrix}

DP55's source form is the bilinear normalized Fock form

\[
 B_\lambda(x,y)=D_\lambda^{-1}\sum_\alpha x_\alpha y_\alpha\alpha!,
\quad
 D_\lambda=\prod_\tau\sum_\alpha (\operatorname{ext}_{\lambda_\tau})_\alpha^2\alpha! .\tag{R8}
\]

At class \(c\) it is multiplied by \(1/a_c\), where \(a_c\) is the actual finite stabilizer order. Choose bases of the actual invariant coefficient lattices, not a larger unconstrained polynomial module. Write \(G_c\) for their raw Fock matrix and \(f_c\) for the point coordinates. When the row lattice is expressed as the direct sum of these invariant class lattices, (R4) becomes the explicit formula

\[
 s_{\mathfrak p}
 =\min_{c,b}v_{\mathfrak p}\left(
       \frac{(f_c^{\mathsf T}G_c)_b}{a_cD_\lambda}\right).\tag{R9}
\]

If integral coordinates are tied together by an additional lattice inclusion \(C\), replace the row of all displayed coefficients by its product with \(C\). Formula (R9) then applies to generators of that actual image lattice. Thus no arbitrary basis is silently assumed unimodular.

\textbf{Proof of the application.} The source period is the finite sum of the classwise forms (R8), weighted by \(a_c^{-1}\). Evaluating it on a basis row supported on one class gives exactly the coefficient inside (R9). Formula (R4) then applies. If rows are restricted to an integral submodule, evaluating on its generators multiplies by its inclusion matrix. This proves the corrected formula for that case too.

\textbf{Locators.} DP\_MULT DP55, (ED3), (ED7), (ED11), and its explicit statement that the arbitrary detecting row has no proved denominator or norm bound. Directly read original-source declarations: \nolinkurl{Formalism/LocalToGlobal/Columns/CoefModB.lean}, \nolinkurl{fockB} line 18, \nolinkurl{extNormB} 55--56, \nolinkurl{pairB} 58--59, \nolinkurl{pairBK} 61--62, \nolinkurl{PairBDual} 113--119; \nolinkurl{Columns/Record.lean}, \nolinkurl{perCl} 315--329 and \nolinkurl{periodC} 1568--1571; \nolinkurl{Columns/TiedPer.lean}, \nolinkurl{per_jMCL_eq} 243--244 and \nolinkurl{per_jMC_onto} 250--252. Perfectness over the field does not assert that the two chosen integral lattices are mutually dual.

\textbf{Consequence for Q3e.} A unit-valued complex period cannot be inserted at zero arithmetic cost merely by rescaling. For a specified algebraic terminal lift vector, the endpoint denominator problem is exactly the content computation (R9), and the archimedean problem is the dual norm (R5). These are computable once the actual kernel/image and lattice matrices are given. This result does not construct that terminal lift vector or identify its lattice, and therefore does not decide the sign in (R1).

\subsection{AM-R-003. Powers of the genus-three Jacobian do not repair an elliptic Serre-tensor requirement}\label{am-r-003.-powers-of-the-genus-three-jacobian-do-not-repair-an-elliptic-serre-tensor-requirement}

Let \(J/\overline{\mathbb Q}\) be the Jacobian of \(y^2=x^7-1\), with its accepted CM type \(\Phi=\{1,2,3\}\) on \(E=\mathbb Q(\zeta_7)\). For every \(N>0\), no nonzero elliptic-power abelian variety is isogenous to an abelian subvariety or quotient of \(J^N\). In particular there is no nonzero isogeny, nor a nonzero homomorphism, in either direction between \(J^N\) and \(C^d\), for any elliptic curve \(C\) and \(d>0\).

This remains true after a finite extension of the base field, after twisting the displayed CM action by automorphisms, and for the geometric product \(J^P\times\bar J^Q\) used by DP65, where \(\bar J\) is the same abelian variety with conjugate \(E\)-action. Consequently an actual Serre tensor of an elliptic curve by a finite projective module cannot become such a geometric power of \(J\) through a field enlargement or an isogeny. This excludes only attempts to retain the source's elliptic \nolinkurl{IsTen} requirement. It is not a no-go theorem for a newly defined mixed-type theta correspondence.

\textbf{Proof.} The only proper CM subfield of the cyclic degree-six field \(E\) is its imaginary quadratic field \(K=\mathbb Q(\sqrt{-7})\). The two induced types from \(K\) are \(\{1,2,4\}\) and \(\{3,5,6\}\). The actual type \(\{1,2,3\}\) is neither, hence primitive. Milne, \emph{Complex Multiplication} (2020), Proposition 3.13, identifies primitive CM types with simple complex CM abelian varieties up to isogeny. Thus \(J\) is geometrically simple of dimension three.

A nonzero homomorphism \(J\to C\) would have an abelian image of positive dimension, hence image \(C\). Its connected kernel would have dimension two, contradicting simplicity. A nonzero map \(C\to J\) would give a positive-dimensional abelian subvariety of dimension at most one, again a contradiction. Every homomorphism between products is a matrix of homomorphisms between factors, so all maps \(J^N\leftrightarrow C^d\) vanish. By Poincare complete reducibility, a nonzero elliptic isogeny factor of a subvariety or quotient of \(J^N\) would give such a nonzero map, which is impossible.

The argument is over \(\overline{\mathbb Q}\), so a finite base extension changes nothing. Twisting a named endomorphism action does not change the underlying variety. For a finite projective module \(M\) over an order acting on an elliptic curve, \(M\otimes C\) is, after clearing denominators in a module basis, isogenous to an elliptic power. Hence it cannot be isogenous to a nonzero variety in the \(J\)-isotypic category. This proves the statement.

\textbf{Locators and scope.} The CM type and the earlier type-package obstruction are carried DP22, DP25, DP49, DP52(d), DP57. The exact old \nolinkurl{IsTen} input is \nolinkurl{Development/AbelianVarietyTheory/Definitions/CMLattices/IsDiffOf.lean}, lines 73--77, as identified in DP57 §7. The new assertion applies the simple-isogeny-factor invariant to arbitrary powers and field enlargements; it does not reprove the accepted local derivative results. In motives it says that the weight-one simple constituent \(h^1(J)\) cannot be replaced by an elliptic \(h^1\) constituent. Exterior and tensor powers of motives are a different operation; no vanishing of all higher-weight correspondences is claimed.

\clearpage

\section{ALP\_MACHINE: lift accounting and explicit CM metric calculations}\label{alp_machine-lift-accounting-and-explicit-cm-metric-calculations}

Author: /root/lift\_accounting, GPT-6 Astra Pro. Reasoning effort: not exposed. Started 2026-11-21 19:35 UTC. This working contribution is not a resolution of DP\_MULT Q3e. Accepted DP findings are cited and are not re-proved.

The later, complete companion findings are:

\begin{itemize}
\item
  \nolinkurl{hodge_lattice_and_cm_basis.md}: ALP-LIFT-05, the actual stable Hodge lattices at 2 and 7, full determinant metric/height, and a primitive induced CM eigenbasis with Gram matrix \(G_B I\).
\item
  \nolinkurl{directional_budget_cancellation.md}: ALP-LIFT-06, the exact directionwise prime constants and the arbitrary-order-vector cancellation against the actual diagonal Cauchy cost.
\item
  \nolinkurl{finite_theta_null_normalization.md}: ALP-LIFT-07, the exact finite theta-null normalization, explicit zero characteristic, trace-one principal line and the surviving two-torsion difference at source odd levels.
\item
  \nolinkurl{trace_one_leading_jet.md}: ALP-LIFT-08, the unique actual trace-one section with leading jet \(u_1u_3-u_2^2\), primitive at every \(p\nmid14\), and its complete finite scalar contents at 2 and 7.
\item
  \nolinkurl{actual_torus_first_lift.md}: ALP-LIFT-09, the norm-one torus class calculation, an explicit honest global splitting, actual nonzero first U(1) to U(2) theta lift, and the exact finite-index cost of a common fixed-weight first-image projector.
\item
  \nolinkurl{permitted_orbit_first_image.md}: ALP-LIFT-10, the exact DP59 half-torsion finite-input family, common actual first-lift input, exact index/norm and exceptional-prime projection costs. AM-K018 closes the ordinary-power source gauge; explicit DP59 harmonic sharp values now occur in an actual first image. Ordinary mixed-line values use a separate multiweight projector.
\item
  \nolinkurl{fixed_conductor_type_repair.md}: ALP-LIFT-11, an explicit global torus character with arbitrary integral infinity type and fixed conductor29, its exact determinant twist of the actual first lift, and removal of the specified next-edge real-type obstruction. Its fixed bilinear repair works for all degrees; the common source-character pins remain separate.
\item
  second\_definite\_lift\_nonvanishing.md: ALP-LIFT-12, an actual nonzero U(2) to U(4) lift of every nonzero repaired first-image vector, including automorphic multiplicities. The proof checks all finite local occurrence, the explicit compact-real Rallis factor, and absolute nonvanishing of the standard L-value. AM-R009 gives the concrete same-host corollary over E-prime. The second Schwartz function is unrestricted; its prescribed arithmetic row and cost data are not constructed.
\item
  fixed\_section\_sharp\_first\_images.md: ALP-LIFT-13, actual sharp first-image coefficients at3 and5 from the single fixed section in AM-K020. One fixed finite vacuum and one unit-norm source character realize every nonzero witness; the exact Fock norm and common period are retained.
\end{itemize}

\nolinkurl{cmdata_theta_constructor_audit.md} records the directly checked input and output of the generic six-identity constructor. Its principal distinction is already accepted DP52 background, not a new finding.

Each companion proof has been read in full and its actual checks are recorded in \nolinkurl{lift_accounting_CHECK_EVENTS.tsv}. None identifies the complete selected \nolinkurl{SizeData₂} or constructs the requested lift chain.

\subsection{1. Source boundary and the actual cost variables}\label{source-boundary-and-the-actual-cost-variables}

The attached DP\_MULT return contains the paper, proof snapshots and manifests, but no raw Lean modules. The present run recovered the pinned mirror modules listed in lift\_sources/FETCH\_MANIFEST.json. Their hashes are compared separately with the intake's original-author manifest; a source read is not a Lean build.

The exact primary declarations are:

\begin{itemize}
\tightlist
\item
  \nolinkurl{Formalism/LocalRingCardinality/DataParams₂.lean,} SizeData₂, lines 79--183: the fields gain, F, E, tau, supp, aT, logDerivativeAtOne and the size windows.
\item
  The same file, sz, gainD, Ga, Ee, logOm and MAIN, lines 203--234.
\item
  \nolinkurl{Formalism/LocalRingCardinality/DataHost.lean,} rdE and v, lines 140--145: \(v=\operatorname{rd}(E)/2\), not \(\sqrt d/2\).
\item
  \nolinkurl{Formalism/SharpRestriction/GlobalTheoremInputs.lean,} LinkDIV\_at, lines 77--89; flG, loG, LinkG\_at, lines 100--136.
\item
  \nolinkurl{Formalism/LocalRingCardinality/SpanPowerMembership/NZRowAsm.lean,} ExtraNZIn, lines 61--86.
\end{itemize}

In particular, a Jacobian, a prime, a CM type and an admissible Ranks record do not determine numerical gain or error fields. The two-sided LinkG\_at conditions include the coarse bound \[
 0\le {\rm gain}_{j,w}\le
 4\,|{\rm TT}|\,a/p,
\] as well as the lower ramplike certificate. The latter is not an equality. LinkDIV\_at controls each \(e_j\) within \(C B_0\) of \[
 \left(\sum_t A_t\right)\log(g(0)g(1)).
\] Its actual value still depends on the selected geometric and integral data. The exact period scalar of the old field package is \[
 \log\Omega=
 \gamma+\log(2\pi)-\log\operatorname{rd}(E)
 -\frac{L'(1,\chi_K)}{L(1,\chi_K)}.
 \tag{L0}
\] Consequently the old \(H\), before any mixed-type replacement, is \[
 H=\sum_t A_t\left[
 \log A_t-1-\gamma-\log(2\pi)+\log\operatorname{rd}(E)
 +\frac{L'(1,\chi_K)}{L(1,\chi_K)}
 \right].
 \tag{L1}
\] This is an exact substitution in source definitions, not a value of a newly constructed mixed theta family.

DP44 and DP52 already settle rank reconstruction, the generic Hermitian spaces, the conditional six-identity induction and the endpoint-period interface. DP65 supplies actual CM-labelled local rows, but explicitly does not supply one SizeData₂ package or membership in a prescribed theta image. Its final statement must not be substituted for either of those missing identifications.

\subsection{2. ALP-LIFT-01: the comparison for two rebuilt families}\label{alp-lift-01-the-comparison-for-two-rebuilt-families}

Let the old family have length \(M\), and the new family length \(M'<M\). Primes, fields and norms may differ; define their scalar totals separately: \[
 Q_M=\sum_{j=1}^{M}(g_j-e_j)-T_\tau-MH+K,\qquad
 Q'_{M'}=\sum_{j=1}^{M'}(g'_j-e'_j)-T'_\tau-M'H'+K'.
\] No identification of local place sets is needed to compare these real numbers. Then \[
\begin{split}
 Q'_{M'}-Q_M={}&-\sum_{j=M'+1}^{M}(g_j-e_j-H)\\
 &+\sum_{j=1}^{M'}[(g'_j-g_j)-(e'_j-e_j)]\\
 &+(T_\tau-T'_\tau)+(K'-K)+M'(H-H').
\end{split}
\tag{L2}
\] If an additional allowance is subtracted, the right side acquires \(B_M-B'_{M'}\).

\textbf{Proof.} Split the old gain and error sums at \(M'\), subtract the two definitions, and use \(MH-M'H'=(M-M')H+M'(H-H')\). This accounts for every displayed term and proves (L2). Adding the allowance is direct subtraction. No positivity assertion is used. \(\square\)

When the retained data and all global terms are equal, (L2) is precisely DP44.3. For a genuinely changed CM type it is the second through fourth lines, rather than the deleted tail alone, that must be computed as well. Moreover, \(Q'_{M'}<Q_M\) alone would not prove failure of a global threshold condition: both quantities could exceed that threshold. This last observation is elementary logic about inequalities; it is not an example of an admissible geometric family.

\subsection{3. ALP-LIFT-02: exact stable CM-height penalty for the proposed Jacobian}\label{alp-lift-02-exact-stable-cm-height-penalty-for-the-proposed-jacobian}

Put \(E=\mathbb Q(\zeta_7)\), \(K=\mathbb Q(\sqrt{-7})\). Let \(J\) have its actual integral CM action of type \(\Phi=\{1,2,3\}\), as in DP22 and DP49. Choose a CM elliptic curve \(B\) with full ring \(\mathcal O_K\), and define \[
 A_{\mathrm{ind}}=\mathcal O_E\otimes_{\mathcal O_K} B
\] using the embedding of \(K\) for which the induced type is \(\Psi=\{1,2,4\}\). In any common standard normalization of stable Faltings height for three-dimensional abelian varieties, \[
\boxed{\;
 h_F(J)-h_F(A_{\mathrm{ind}})
 =\log\frac{\Gamma(2/7)\Gamma(4/7)^3}
                  {\Gamma(3/7)^3\Gamma(5/7)}
 =:\delta_7>0.\;}
\tag{L3}
\] The exact rational computation accompanying this proof certifies \[
       0.05568702<\delta_7<0.05568705.
\tag{L4}
\] The higher precision decimal \(0.0556870335243028288211564027305\ldots\) is a numerical cross-check, not the rational certificate.

\subsubsection{Proof: integral comparison variety}\label{proof-integral-comparison-variety}

Let \(S=\zeta_7+\zeta_7^2+\zeta_7^4\), \(T=-1-S\). The roots \(\zeta_7,\zeta_7^2,\zeta_7^4\) give \[
 \zeta_7^3-S\zeta_7^2+T\zeta_7-1=0.
\] Since \(\mathcal O_K=\mathbb Z[S]\), \(\mathcal O_E=\mathbb Z[\zeta_7]\), and \([E:K]=3\), the elements \(1,\zeta_7,\zeta_7^2\) are an \(\mathcal O_K\)-basis of \(\mathcal O_E\). The Serre tensor construction therefore identifies \(A_{\mathrm{ind}}\) with \(B^3\) as an abelian variety (over a common field of definition), not merely up to isogeny. Thus \(h_F(A_{\mathrm{ind}})=3h_F(B)\). The \(\mathcal O_E\)-action has exactly the embeddings extending the chosen \(K\)-embedding, namely \(\Psi\). No isogeny-invariance assertion for general Faltings heights is being used.

\subsubsection{Proof: the normalized Colmez class function}\label{proof-the-normalized-colmez-class-function}

Let \(\sigma=\sigma_3\) generate \({\rm Gal}(E/\mathbb Q)\), and let \(\chi(\sigma)=\exp(2\pi i/6)\). Then \(\chi^3=\chi_{-7}\), and the exponent sets of \(\Phi\) and \(\Psi\) are respectively \(\{0,1,2\}\) and \(\{0,2,4\}\).

Use Haar probability measure in convolution: \[
 A_\Phi=\mathbf1_\Phi*\mathbf1_{\Phi^{-1}}.
\] For \(\sigma^k\), \(k=0,\ldots,5\), its values are \[
 (1/2,\,1/3,\,1/6,\,0,\,1/6,\,1/3).
\] The corresponding induced-type values are \((1/2,0,1/2,0,1/2,0)\). Discrete Fourier expansion gives \[
 A_\Phi=\frac14\mathbf1+
             \frac19(\chi+\bar\chi)+\frac1{36}\chi^3,\qquad
 A_\Psi=\frac14\mathbf1+\frac14\chi^3.
\tag{L5}
\] These expansions can also be checked at the six group elements, using \(\chi+\bar\chi=(2,1,-1,-2,-1,1)\) and \(\chi^3=(1,-1,1,-1,1,-1)\).

The abelian Colmez theorem applies to both integral CM types. Its normalization is explicit in Yang (2010), Section 3, pp.~464--466, Conjecture 3.1 and the elliptic example (3.3); the abelian case is supplied by Colmez (1993), Théorème 0.5 and Théorème III.2.9, with the possible dyadic ambiguity removed by Obus (2013), Theorem 3.9 in arXiv:1107.0684v2. Write \(Z_\eta=L'(0,\eta)/L(0,\eta)\). In the difference, the trivial character and metric constants cancel. All three nontrivial characters here have conductor \(7\), and their coefficient difference sums to zero. The conductor term therefore cancels as well. Multiplication by \([E:\mathbb
Q]=6\), which converts the type-normalized height to the actual three-dimensional height, gives \[
 h_F(J)-h_F(A_{\mathrm{ind}})
   =-\frac23(Z_\chi+Z_{\bar\chi}-2Z_{\chi^3})
   =\frac43(Z_{\chi^3}-\operatorname{Re}Z_\chi).
\tag{L6}
\]

\subsubsection{Proof: reduction to gamma values}\label{proof-reduction-to-gamma-values}

For any nontrivial Dirichlet character \(\eta\) modulo \(7\), the Hurwitz formula, together with DLMF 25.11.13 and 25.11.18, gives \[
 L(0,\eta)=-\frac17\sum_{a=1}^{6}a\eta(a),\qquad
 L'(0,\eta)=-\log7\,L(0,\eta)+
                 \sum_{a=1}^{6}\eta(a)\log\Gamma(a/7).
\tag{L7}
\] Here the constant from \(\zeta'(0,x)=\log\Gamma(x)-\tfrac12\log(2\pi)\) vanishes because the character sum is zero.

Put \(\ell_a=\log\Gamma(a/7)\). The character values give \(L(0,\chi^3)=1\) and \[
 L(0,\chi)=\frac{4+2\sqrt3\,i}{7},\qquad
 L(0,\chi)^{-1}=1-\frac{\sqrt3}{2}i.
\] Consequently \[
 Z_{\chi^3}=-\log7+\ell_1+\ell_2-\ell_3+\ell_4-\ell_5-\ell_6
\] and \[
 \operatorname{Re}Z_\chi=
 -\log7+\ell_1+\frac14\ell_2+\frac54\ell_3
                    -\frac54\ell_4-\frac14\ell_5-\ell_6.
\] Substitution in (L6) proves the equality in (L3).

\subsubsection{Proof: strict sign and exact rational certificate}\label{proof-strict-sign-and-exact-rational-certificate}

Let \(f(x)=\log\Gamma(x)\), \(h=1/7\). The expression in (L3) is \(-\Delta_h^3 f(2/7)\). By differentiating DLMF 5.15.1, \[
 f'''(x)=-2\sum_{n\ge0}(n+x)^{-3}<0 \quad(x>0).
\] Repeated integration of a finite difference therefore gives \[
 \delta_7=-\int_{[0,h]^3} f'''(2/7+s+t+u)\,ds\,dt\,du>0.
\] This proves the strict sign without numerical approximation.

For a reproducible rational bound, the Weierstrass product (or the same integral identity after integration term by term) yields \[
 \delta_7=\sum_{n=0}^{\infty}
 \log\frac{(7n+3)^3(7n+5)}{(7n+2)(7n+4)^3}
 =\sum_{n=0}^{\infty}\log(1+x_n),
\quad
 x_n=\frac{14n+7}{(7n+2)(7n+4)^3}.
\] For \(t=x/(2+x)\), the six-term sum \[
 2\sum_{k=0}^{5}\frac{t^{2k+1}}{2k+1}
\] is a lower bound on \(\log(1+x)\); the omitted positive tail is at most \(2t^{13}/(13(1-t^2))\). Bound the first \(100\) terms with exact fractions. Each remaining logarithm is a third finite difference of \(\log x\), and lies between \(2/(7n+5)^3\) and \(2/(7n+2)^3\). Integral comparison bounds the tail starting at \(N=100\) between \[
 \frac1{7(7N+5)^2},\qquad \frac1{7(7N-5)^2}.
\] The script colmez\_height\_certificate.py implements exactly these rational inequalities and verifies (L4). This completes the proof. \(\square\)

\subsubsection{Independent normalization cross-check against Wilms}\label{independent-normalization-cross-check-against-wilms}

In Yang's height convention, (L5) gives \[
 h_F^Y(J)=-\frac43\operatorname{Re}Z_\chi
               -\frac16Z_{\chi^3}-\frac34\log7
          =-2.0840659166599491401\ldots,
\] \[
 h_F^Y(A_{\mathrm{ind}})
          =-\frac32Z_{\chi^3}-\frac34\log7
          =-2.1397529501842519690\ldots.
\] Wilms (2026), Theorem 1.1, specialized to \(n=7,g=3\), gives \[
 h_F^W(J)=\log7-\frac32\log\pi+\frac3{14}\log2
     -\sum_{j=1}^{3}
       \log\frac{\Gamma((2j-1)/14)}{\Gamma((j+3)/7)}.
\tag{L8}
\] His Section 7, pp.~12--13, explicitly records the normalization \(h'_F=h_F^W+(g/2)\log\pi\). The duplication formula and \(\prod_{a=1}^6\Gamma(a/7)=(2\pi)^3/\sqrt7\) show algebraically that (L8) equals \(h_F^Y(J)-(3/2)\log\pi\). The independent numerical evaluation is \(-3.8011607454340494014\ldots\). This check catches the degree-six, dimension-three and metric factors; it is not used to prove (L3).

\subsubsection{Exact implication and limit}\label{exact-implication-and-limit}

The actual Jacobian has a strictly larger base CM stable height than the specified integral induced-type product. Thus replacing that product by J does not provide a free saving in this complete invariant of the base Hodge determinant. It does not follow that \(H'\ge H\): the source's \(H\) also incorporates a particular derivative normalization, root discriminant, unequal embedding weights, gauges and other fields. The identification of that source metric with this determinant line is a separate equation that must be proved in a mixed-type construction. There is no asserted sign for \(Q'_{M'}-Q_M\).

In Grothendieck's language, (L3) compares the arithmetic degrees of the metrized determinants of the Hodge bundles \(e^*\Omega^1_{\mathcal J}\) and \(e^*\Omega^1_{\mathcal A_{\mathrm{ind}}}\) on a common semistable integral base. The CM idempotents are those acting on the \(H^1\) motive; (L5) is the explicit Artin-character expression entering its period height. FGA's Picard-scheme theorem supplies the Jacobian, as already used in DP22; the exact FGA locators are carried from that finding. No new manuscript cote or scan is asserted to contain the Colmez formula.

\subsection{4. ALP-LIFT-03: full complex Hodge metric in the actual CM basis}\label{alp-lift-03-full-complex-hodge-metric-in-the-actual-cm-basis}

For \(C:y^2=x^7-1\), put \(\omega_i=x^{i-1}dx/y\) for \(i=1,2,3\). Equip \(H^0(C,\Omega^1)\) with \[
 (\omega,\eta)=\frac{i}{2}\int_{C(\mathbb C)}\omega\wedge\bar\eta.
\] Then \[
 (\omega_i,\omega_j)=0\quad(i\ne j),\qquad
\boxed{\;
 \|\omega_i\|^2=\frac{2\pi}{7}
 \frac{\Gamma(i/7)\Gamma(1/2-i/7)}
      {\Gamma(1-i/7)\Gamma(1/2+i/7)}.\;}
\tag{L9}
\] The norm squares are approximately \(9.47604253508638,\ 8.00728744152725,\ 15.34580357802465\).

\textbf{Proof.} The automorphism \(x\mapsto\zeta_7 x\) preserves the integral and multiplies \(\omega_i\wedge\bar\omega_j\) by \(\zeta_7^{i-j}\). For \(i\ne j\) this scalar is not one, proving orthogonality. Branch points have measure zero. Over the affine line, the map from C has degree two and the two sheets contribute equally, so \[
 \|\omega_i\|^2=
 2\int_{\mathbb C}
       \frac{|x|^{2i-2}}{|x^7-1|}\,dA(x).
\] Substitute \(z=x^7\), accounting for its degree seven and Jacobian \(49|x|^{12}\). The expression becomes \[
 \frac27\int_{\mathbb C}|z|^{2a-2}|1-z|^{-1}\,dA(z),
 \qquad a=i/7.
\tag{L10}
\] For \(a,b>0\) and \(a+b<1\), the complex beta integral is \[
 \int_{\mathbb C}|z|^{2a-2}|1-z|^{2b-2}\,dA(z)
 =\pi\frac{\Gamma(a)\Gamma(b)\Gamma(1-a-b)}
            {\Gamma(1-a)\Gamma(1-b)\Gamma(a+b)}.
\tag{L11}
\] For completeness, write each negative power as its positive gamma integral, apply Tonelli, and integrate the Gaussian in z: \[
 \frac{\pi}{\Gamma(1-a)\Gamma(1-b)}
 \int_{s,t>0}\frac{s^{-a}t^{-b}}{s+t}
                  e^{-st/(s+t)}\,ds\,dt.
\] With \(s=uv,t=u(1-v)\), this becomes \[
 \frac{\pi\Gamma(1-a-b)}{\Gamma(1-a)\Gamma(1-b)}
 \int_0^1 v^{b-1}(1-v)^{a-1}\,dv.
\] The ordinary beta identity proves (L11). All changes of variables are valid on positive integrands; the stated inequalities give convergence. Set \(b=1/2\) in (L11) and substitute into (L10). This proves (L9). \(\square\)

This supplies the full complex metric for arbitrary unequal embedding weights once the geometric differential basis is tied to the chosen theta connection and source gauge. It does not remove that tie or the integral corrections at 2 and 7. In particular the Hodge norms are not silently declared to be the source's \(\Omega\).

Geometrically, (L9) gives the Hermitian metric on the Hodge subbundle of the de Rham realization of \(h^1(J)\), split by the \(\mathcal O_E\) idempotents. Its dual metric controls invariant tangent derivations. The connection to the Poincaré biextension and its theta sections still requires the arithmetic comparison asked for in THETA3 G1.

\subsection{5. ALP-LIFT-04: the exact constant in the saturated bare-budget difference}\label{alp-lift-04-the-exact-constant-in-the-saturated-bare-budget-difference}

This result uses the accepted DP64 definition (CB1); it does not repeat the cyclic-display or global-integrality proofs behind that definition. It determines the linear term left undetermined by DP64's general \(O(N)\) comparison for this particular pair of types and equal orders.

Continue with \(\Phi=\{1,2,3\}\), \(\Psi=\{1,2,4\}\). For a rational prime \(p\ne7\), define the \textbf{formal slope totals} \[
 S_\Phi(p)=\sum_{i\in\Phi}\sum_{r\ge1}
       \frac{\mathbf1_{p^r i\notin\Phi}}{p^r},\qquad
 \Delta S(p)=S_\Phi(p)-S_\Psi(p),
\] where labels are multiplied modulo 7. Let \(S\) be a fixed finite set of primes containing 7. Set \[
 \delta_{7,S}=\delta_7-
       \sum_{p\in S\setminus\{7\}}\Delta S(p)\log p.
 \tag{L12}
\] Then the following are exact statements about these explicitly specified arithmetic certificates:

\begin{enumerate}
\def\labelenumi{\arabic{enumi}.}
\tightlist
\item
  The local difference is \[
  \boxed{\quad
  \Delta S(p)=
  \begin{cases}
  \displaystyle\frac{2p(p+1)}{p^3-1},&p\equiv2,4\pmod7,\\[4pt]
  \displaystyle-\frac{2p(p-1)}{p^3+1},&p\equiv3,5\pmod7,\\[4pt]
  0,&p\equiv1,6\pmod7.
  \end{cases}\quad}
  \tag{L13}
  \]
\item
  In ordinary increasing rational-prime order, \[
  \boxed{\quad
  \lim_{X\longrightarrow\infty}
    \sum_{\substack{p\le X\\p\notin S}}
           \Delta S(p)\log p=\delta_{7,S}.
  \quad}
  \tag{L14}
  \] This convergence is conditional. No rearrangement of the primes is authorized by (L14).
\item
  Put equal derivative order \(n\) on each of the three CM directions in both types, and use the DP64 integer certificates \[
  D_S(n)=\sum_{p\notin S}\log p
      \left(\sum_{i\in\Phi}E^\Phi_{i,p}(n)
                  -\sum_{i\in\Psi}E^\Psi_{i,p}(n)\right).
  \] The sum is finite, since a term can be nonzero only when \(p\le n\). For some constants \(c>0,C_S>0\), \[
  \boxed{\quad
      D_S(n)=n\delta_{7,S}
                 +O_S\!\left(n e^{-c\sqrt{\log n}}\right)
      \qquad(n\longrightarrow\infty).
  \quad}
  \tag{L15}
  \] In particular, with total order \(N=3n\), the exact coefficient of \(N\) is \(\delta_{7,S}/3\).
\item
  The same linear term holds before saturation. If \(D_S(n;X)\) denotes the definition of \(D_S(n)\) with \(p\le X\), then, uniformly for \(n\ge4\), \(X\ge2\), with some fixed \(c>0\), \[
  D_S(n;X)=n\delta_{7,S}
    +O_S\!\left(n e^{-c\sqrt{\log\min\{X,\sqrt n\}}}\right).
  \tag{L15a}
  \] Consequently \(D_S(n;X(n))/n\to\delta_{7,S}\) for \emph{every} \(n\to\infty\) and \(X(n)\to\infty\), including cutoffs much smaller than the derivative order. The three directions still have the same order \(n\), fixed as \(p\) varies.
\end{enumerate}

The source-supported common good-model range in DP49/DP64/DP65 excludes 2 and 7. It is therefore the specialization \[
 \boxed{\quad
 \delta_{7,\{2,7\}}=
       \delta_7-\frac{12}{7}\log2<0,
 \quad}
 \tag{L16}
\] that applies to that range. The identity with \(\delta_7\) itself uses the \emph{formal} cyclic slope at 2; it does not assert an available integral good-model certificate at 2 for the specified curve. Every additional discarded prime has the explicit correction in (L12). The strict sign in (L16) follows from (L4) and the elementary inequality \(\log2>1/2\).

\subsubsection{Proof: exact local class function}\label{proof-exact-local-class-function}

For \(a\in(\mathbb Z/7)^\times\), put \[
 b(a)=\#\{i\in\Phi:ai\notin\Phi\}
                  -\#\{i\in\Psi:ai\notin\Psi\}.
\] The Haar normalization already checked in (L5) gives \[
 b(a)=-6(A_\Phi-A_\Psi)(\sigma_a).
\] Thus \(b(1)=b(6)=0\), \(b(2)=b(4)=2\), and \(b(3)=b(5)=-2\). Write \(b(0)=0\). The function has mean zero over the six nonzero residues, and \[
 b(a)=-6\sum_{\eta\in\{\chi,\bar\chi,\chi^3\}}d_\eta\eta(a),
 \qquad(d_\chi,d_{\bar\chi},d_{\chi^3})=(1/9,1/9,-2/9).
 \tag{L17}
\] Directly from the DP64 definitions, \[
 \Delta S(p)=\sum_{r\ge1}\frac{b(p^r)}{p^r}.
 \tag{L18}
\] For residues 2 or 4 the period-three numerator is \((2,2,0)\), giving \(2(p^2+p)/(p^3-1)\). For residues 3 or 5 the period-six numerator is \((-2,2,0,2,-2,0)\). Its rational sum is \[
 \frac{-2p^5+2p^4+2p^2-2p}{p^6-1}
   =-\frac{2p(p-1)}{p^3+1}.
\] Residues 1 and 6 give zero at every power. This proves (L13), including \(\Delta S(2)=12/7\), \(\Delta S(3)=-3/7\), and \(\Delta S(11)=132/665\).

\subsubsection{Proof: ordinary prime-order convergence and the constant}\label{proof-ordinary-prime-order-convergence-and-the-constant}

The known fixed-modulus prime number theorem, as stated in Rudnick (2015), Section 2, printed p.~2 immediately before equation (1), implies for some fixed \(c_0>0\) \[
 B_1(x):=\sum_{p\le x}b(p)\log p
           =O\!\left(xe^{-c_0\sqrt{\log x}}\right).
 \tag{L19}
\] Indeed the six residue-class main terms cancel because the mean of \(b\) is zero. Replacing the displayed von Mangoldt sum in that theorem by its prime-only part changes it by \(O(\sqrt{x}\log^2x)\), which is absorbed by the stated error after possibly decreasing \(c_0\). Partial summation now shows convergence of \(\sum_p b(p)\log p/p\): its tail is at most a constant times \[
 (1+\sqrt{\log X})e^{-c_0\sqrt{\log X}}.
\] The contribution of \(r\ge2\) in (L18) is absolutely convergent, since \[
 \sum_p\log p\sum_{r\ge2}\frac{|b(p^r)|}{p^r}
 \le2\sum_p\frac{\log p}{p(p-1)}<\infty.
\] This proves the convergence assertion in ordinary prime order before any use of analytic continuation. Conversely, (L13) and the same fixed-modulus prime theorem give \[
 \sum_{p\le X}|\Delta S(p)|\log p
           =\frac43\log X+O(1),
 \tag{L20}
\] so the series is not absolutely convergent.

For real \(s>1\), the absolutely convergent Euler product, DLMF 25.15.2, and (L17) give \[
 \sum_{p\ne7}\log p\sum_{r\ge1}\frac{b(p^r)}{p^{rs}}
       =6\sum_\eta d_\eta\frac{L'(s,\eta)}{L(s,\eta)}.
 \tag{L21}
\] The left side tends as \(s\downarrow1\) to the ordinary prime sum just proved convergent. To justify the passage without a Tauberian premise, apply partial summation to (L19): the prime, first-power part is an integral with integrand bounded, uniformly for \(1\le s\le2\), by a constant times \(e^{-c_0\sqrt{\log x}}/x\), which is integrable at infinity. The higher-power part admits the absolute majorant above. Dominated convergence therefore applies. Nonvanishing of the three \(L(1,\eta)\) is the classical Dirichlet nonvanishing theorem (DLMF 25.15.9).

All three characters are odd, primitive and of conductor 7. The functional equation (DLMF 25.15.5), logarithmically differentiated at \(s=1\), yields \[
 \frac{L'(0,\eta)}{L(0,\eta)}
   +\frac{L'(1,\bar\eta)}{L(1,\bar\eta)}
       =\gamma+\log(2\pi/7).
 \tag{L22}
\] For clarity, in the DLMF normalization the sine factor has logarithmic derivative \((\pi/2)\cot(\pi s/2)=0\) at 1, while \(\Gamma'(1)/\Gamma(1)=-\gamma\); the Gauss sum is constant. The conjugate in (L22) is essential. Since the \(d_\eta\) are conjugation-invariant and sum to zero, (L21) at 1 equals \[
 6\sum_\eta d_\eta\frac{L'(1,\eta)}{L(1,\eta)}
 =-6\sum_\eta d_\eta\frac{L'(0,\eta)}{L(0,\eta)}
 =\delta_7,
\] where the last equality is the already-proved (L6). Removing the fixed finite set of prime summands proves (L12) and (L14).

\subsubsection{Proof: the actual integer floors have the same linear constant}\label{proof-the-actual-integer-floors-have-the-same-linear-constant}

First take \(S=\{7\}\). Define an arithmetic sequence \(a(d)=b(d)\Lambda(d)\), where \(\Lambda\) is the von Mangoldt function and \(b(d)\) is read modulo 7. Thus \(a(d)=0\) except at prime powers, and \[
 D_{\{7\}}(n)=\sum_{d\le n}a(d)\left\lfloor\frac nd\right\rfloor.
 \tag{L23}
\] The fixed-modulus prime theorem directly gives \[
 A(x):=\sum_{d\le x}a(d)
       =2\psi(x;7,2)+2\psi(x;7,4)
          -2\psi(x;7,3)-2\psi(x;7,5)
       =O\!\left(xe^{-c_0\sqrt{\log x}}\right).
 \tag{L24}
\] Also \(\sum_{d\le y}|a(d)|\le2\psi(y)=O(y)\), by the same known prime theorem. Partial summation proves that \(\sum_da(d)/d\) converges, with a tail \(O((1+\sqrt{\log y})e^{-c_0\sqrt{\log y}})\). Its value is \(\delta_7\): the prime-first-power sum is the same one as before, and every prime-power term of exponent at least two belongs to an absolutely convergent series. This last absolute convergence justifies the change from increasing primes to increasing prime powers.

Put \(y=\lfloor\sqrt n\rfloor\). Expanding the floor as a count of positive integers gives the exact finite identity \[
 D_{\{7\}}(n)=
 \sum_{d\le y}a(d)\left\lfloor\frac nd\right\rfloor
 +\sum_{k\le n/y}\bigl(A(n/k)-A(y)\bigr).
 \tag{L25}
\] Here the second sum enumerates exactly the pairs \((d,k)\) with \(d>y\), \(dk\le n\); if \(n/k=y\), its summand is zero. The first sum is \[
 n\sum_{d\le y}\frac{a(d)}d+O(y)
 =n\delta_7+
 O\!\left(n(1+\sqrt{\log y})e^{-c_0\sqrt{\log y}}+y\right).
\] For the second sum, (L24), \(n/k\ge y\), and the harmonic-sum bound give \[
 \sum_{k\le n/y}|A(n/k)-A(y)|
 \ll n(1+\log n)e^{-c_0\sqrt{\log y}}.
\] As \(y\asymp\sqrt n\), these errors together are \(O(n e^{-c\sqrt{\log n}})\) for some smaller \(c>0\). This proves (L15) for \(S=\{7\}\).

For one fixed omitted prime \(p\ne7\), the floor error and the geometric tail give \[
 \sum_{r\ge1}b(p^r)\left\lfloor\frac n{p^r}\right\rfloor
       =n\Delta S(p)+O_p(1+\log n).
\] Indeed there are \(O_p(1+\log n)\) nonzero floors, each rounding error has absolute value at most 2, and the omitted geometric tail after \(p^r>n\) contributes \(O_p(1)\). Removing finitely many such contributions proves (L15) and (L16). This completes the proof. \(\square\)

\subsubsection{Proof of the unsaturated-cutoff refinement (L15a)}\label{proof-of-the-unsaturated-cutoff-refinement-l15a}

It suffices first to treat \(S=\{7\}\). Split the sum into prime first powers and powers with exponent at least two. Write \[
 C_1(X)=\sum_{p\le X}b(p)\log p/p,\qquad
 C_2(X)=\sum_{p\le X}\log p\sum_{r\ge2}b(p^r)/p^r.
\] The argument above gives \[
 C_1(X)+C_2(X)=\delta_7+
             O\!\left(e^{-c_1\sqrt{\log X}}\right)
 \tag{L26}
\] for some smaller \(c_1>0\); also \(C_2(\infty)-C_2(X)=O((1+\log X)/X)\).

The second-power-and-higher rounding error is uniform in \(X\). Indeed, let \[
 H_2(t)=\sum_{\substack{p^r\le t\\r\ge2}}\log p
            =O(\sqrt t\log t).
\] This bound follows from \(\theta(u)=O(u)\), by summing \(\theta(t^{1/r})\) over at most \(\log_2t\) exponents. The floor errors for \(p^r\le n\) have total absolute value at most \(2H_2(n)\). The linear geometric terms with \(p^r>n\) have total absolute value at most \(2n\int_{(n,\infty)}t^{-1}\,dH_2(t)\), which is \(O(\sqrt n\log n)\) by integration by parts. These bounds remain valid after imposing \(p\le X\). Thus that whole contribution equals \[
                    nC_2(X)+O(\sqrt n\log n)
 \tag{L27}
\] uniformly in \(X\).

Put again \(y=\lfloor\sqrt n\rfloor\). If \(X\le y\), the first-power contribution is \(nC_1(X)+O(\theta(X))\), hence (L26) and (L27) give (L15a). If \(X\ge y\), its exact hyperbola decomposition is \[
 \sum_{p\le X}b(p)\log p\left\lfloor\frac np\right\rfloor
 =\sum_{p\le y}b(p)\log p\left\lfloor\frac np\right\rfloor
  +\sum_{k\le n/y}
       \bigl(B_1(\min\{X,n/k\})-B_1(y)\bigr).
 \tag{L28}
\] In each summand, \(y\le\min\{X,n/k\}\le n/k\). Applying (L19) and the harmonic-sum bound makes the second sum \(O(n(1+\log n)e^{-c_0\sqrt{\log y}})\). The first sum is \(nC_1(\infty)+O(n e^{-c_1\sqrt{\log y}}+y)\). Together with (L27) and \(|C_2(\infty)-C_2(X)|\ll(1+\log y)/y\), this proves (L15a) after decreasing \(c\). In both cases \(\sqrt n\log n\) is absorbed by \(n e^{-c\sqrt{\log\min\{X,\sqrt n\}}}\).

For a fixed finite exceptional set, once \(X\) contains it, the omitted-prime estimate following (L25) applies unchanged. If it does not yet contain it, the missing finite correction is bounded by a constant depending on \(S\) times \(n\), while \(X\) stays below the fixed largest exceptional prime. Enlarging the implied constant therefore proves the stated uniform assertion for all \(X\ge2\). \(\square\)

\subsubsection{Consequence, and precise limitation}\label{consequence-and-precise-limitation}

For the formal all-unramified-prime slope comparison, the finite part of the mixed-type gain equals the complete stable CM-height difference in (L3). For the actual common good-model range, the finite part is (L16), with the missing formal 2-adic amount explicitly visible. The new statement reaches the order-\(n\) term of the integer certificates both at saturation and along any growing cutoff, so it is stronger than a comparison of leading \(n\log n\) terms or an Abel-regularized identity.

It is still a theorem about the bare lower certificates, with equal orders held fixed over all primes. It does not identify those certificates with \(g_j\), does not give an upper bound on a selected theta value, and does not replace the source bad-prime, gauge, orbit, connection, row-lattice or division-error terms. In particular the height/slope equality cannot be inserted into \(Q_M\) until the coefficient with which the actual source metric occurs is proved. No conclusion of minimal chain length or impossibility of a shorter mixed-type chain follows from (L12)--(L16).

In motivic terms the same Artin class function \(A_\Phi-A_\Psi\) controls this finite part of the Frobenius-weighted bare certificate and the CM Hodge determinant height, through the functional equation. The exact identification here is a calculation for these two specified invariants; it supplies no new comparison isomorphism between the abelian theta kernel and an automorphic kernel.

\subsection{6. Approaches pursued and present Q3e boundary}\label{approaches-pursued-and-present-q3e-boundary}

\begin{enumerate}
\def\labelenumi{\arabic{enumi}.}
\tightlist
\item
  Re-read DP44's exact comparison and the actual SizeData₂ declarations. The numerical arrays are inputs and the lower certificate is not their exact value. No actual mixed family was found in the supplied return.
\item
  Re-read DP52's six-identity induction and DP65's sharp full-row witnesses. These occupy different interfaces; the latter does not instantiate the former's arithmetic theta family.
\item
  Tested whether field/type or rank restrictions alone give minimality. DP44 already rules out rank minimality; DP25 and DP52 exclude only reuse of the unchanged induced-type wrappers. They do not prohibit a new theta correspondence.
\item
  Examined the old logOm and root-discriminant normalization. This produces exact expressions (L0)--(L1), but not a mixed-family \(H\).
\item
  Calculated the actual complete base CM height, with its strict positive comparison (L3), and the complete complex Hodge Gram (L9). These remove vague dependence on the base CM period as far as these specified metrics are concerned.
\item
  No admissible rescaling counterexample is asserted. Changing gain/E fields in a bare record would not preserve the geometric ties.
\end{enumerate}

The remaining concrete problem is to identify the classical mixed-type theta kernel, its invariant differential operators and its integral lattices with one source-compatible kernel package, then evaluate its gain/error arrays and the gauge change from (L9) in that same package. The full difference is (L2). Neither its strict negativity nor a shorter valid chain has been proved in this contribution.

\subsection{7. Sources with exact locators}\label{sources-with-exact-locators}

\begin{itemize}
\tightlist
\item
  Colmez, P. (1993). Périodes des variétés abéliennes à multiplication complexe. Annals of Mathematics, 138(3), 625--683. Théorèmes 0.3, 0.5 and III.2.9. \url{https://webusers.imj-prg.fr/~pierre.colmez/artin.pdf}
\item
  Obus, A. (2013). On Colmez's product formula for periods of CM-abelian varieties. Mathematische Annalen, 356, 401--418. Theorem 3.9 in the inspected arXiv:1107.0684v2, Section 3.3. \url{https://arxiv.org/abs/1107.0684}
\item
  Yang, T. (2010). The Chowla--Selberg formula and the Colmez conjecture. Canadian Journal of Mathematics, 62(2), 456--472. Section 3, pp.~464--466, definition of \(A_\Phi\), Conjecture 3.1, and equation (3.3); p.~457 states the abelian case. \url{https://doi.org/10.4153/CJM-2010-028-x}
\item
  Wilms, R. (2026). On the Faltings height of the curve \(y^2=x^n-1\). arXiv:2601.15271v1. Theorem 1.1; Section 7, pp.~12--13; Remark 10.2 (the cyclotomic CM type). \url{https://arxiv.org/abs/2601.15271}
\item
  NIST Digital Library of Mathematical Functions. Equations 5.15.1, 25.11.13, 25.11.18. The multiplication and duplication identities used for the cross-check are 5.5.6 and 5.5.5. For ALP-LIFT-04: 25.15.2 (Euler product), 25.15.5 (primitive functional equation, including the conjugate character), and 25.15.9 (nonvanishing at 1). \url{https://dlmf.nist.gov/25.15}
\item
  Rudnick, Z. (2015, May 26). Primes in arithmetic progressions: Fixed modulus. Course notes, Section 2, printed p.~2, quantitative prime number theorem immediately before equation (1). The theorem is used as known; no proof of it is claimed here. \url{https://www.math.tau.ac.il/~rudnick/courses/sieves2015/fixed\%20progressions.pdf}
\item
  DP\_MULT checkpoint 06: DP22, DP25, DP44, DP49, DP52, DP65, each cited at statement level with the intake status retained. ALP-LIFT-04 additionally uses DP64 (CB1), with its explicit fixed exceptional-set convention and distinction from theta-selected values.
\end{itemize}

Web retrieval references for parent reuse: Yang turn507172355view2 and turn815074106view2; Wilms turn507172355view1 and turn815074106view1; Obus turn477052200view3; DLMF turn477052200view0 and turn477052200view2. The parent must itself open a source before citing its web reference. ALP-LIFT-04 primary retrievals: Rudnick turn953517703view0; DLMF functional equation/Euler product/nonvanishing turn953517703view1.

\subsection{ALP-LIFT-14: actual selected even-rank chain}\label{alp-lift-14-actual-selected-even-rank-chain}

Full proof: \nolinkurl{selected_even_chain_nonvanishing.md}, E1-E28. Starts from AM-R-012 on the common Eprime host, uses literal bilinear source convention and equal stage characters. Exact standard factors reduce to one absolutely convergent initial factor and nonzero angular Hecke factors at 1 through d; spherical local occurrence and compact real Schur tests give an actual nonzero chain with unrestricted later finite Schwartz vectors. Whole-proof author check logged 2026-11-21T22:39:18.688282+00:00. Raw Xi membership, later integral restricted-image matrices, six identities and DP44 costs remain open.

\subsection{ALP-LIFT-15: one finite unit-norm translate per edge}\label{alp-lift-15-one-finite-unit-norm-translate-per-edge}

Full proof: \nolinkurl{single_finite_heisenberg_chain.md}, H1-H10. Every later actual nonzero selected lift in LIFT14 can use one finite Heisenberg translate of the actual vacuum, norm exactly1 and coefficient1. A depth-N full window has an orthonormal basis of N\^{}(2g) such vectors and selection bound N\^{}g. Required depth and support are uncontrolled; the family is larger than DP59 and does not identify raw Xi. Author whole-proof PASS 2026-11-21T22:49:14.726742+00:00.

\subsection{ALP-LIFT-16: 29-only single translates and exact level masses}\label{alp-lift-16-29-only-single-translates-and-exact-level-masses}

Full proof: \nolinkurl{spherical_tests_and_29_only_chain.md}, S1-S27. Exact spherical-vacuum transforms, including residue2, refine LIFT15 to one norm-one translate supported only at29 on every later edge. A sufficient rank-n principal level is29\^{}(2e), with exact index I\_n=29\^{}(12 e n\^{}2) product\_(i=1)\^{}n(1-29\^{}(-i))\^{}6 and all stabilizers trivial. Its v11 is6(r+v11(r!)), r=floor(n/10), independent of unknown depth e. Every-vector and multiplicity details are retained. Whole-proof author PASS 2026-11-21T23:04:38.479472+00:00. Depth bounds, raw Xi and the restricted arithmetic image remain open. Finite arithmetic certificate: \nolinkurl{spherical_test_level_certificate.py}.

\subsection{ALP-LIFT-17: prescribed local column types exclude the spherical chain}\label{alp-lift-17-prescribed-local-column-types-exclude-the-spherical-chain}

Full proof: \nolinkurl{prescribed_column_local_type_obstruction.md}, F1-F19. For1\textless=j\textless m the actual colLev\_j has nontrivial determinant-block characters and occurs in no unramified local representation. Forj\textgreater=2, distinct block characters also exclude every tame determinant twist of an unramified representation. Thus the29-only chain cannot itself supply these explicit column modules at the prescribed G places. The variable ChIdx family and raw Xi remain separate. Any genuine source IsMu0 realization has conductor norm(g0*g1)\^{}12 on Eprime; global existence is not asserted. Author whole-proof PASS 2026-11-21T23:20:39.359309+00:00. Finite certificate: \nolinkurl{column_type_flag_certificate.py}.

\subsection{ALP-LIFT-18: exact literal-G even characters}\label{alp-lift-18-exact-literal-g-even-characters}

Full proof: \nolinkurl{literal_g_character_compatibility.md}, G1-G21. The source G half has Norm\_Eprime/Kq(zeta28)=1, so literal residue characters impose exactly sum\_(a in Phi) a e\_i,a=0 mod28; class-one norm-one torus descent proves necessity and sufficiency for IsMu0. Uniform e\_i divisible7 works. Explicit primes113/281, coefficient fieldQ(zeta560), conductor norm31753\^{}12, and numerical Hyp3 family k6,m2,t1795,C46539458987 are supplied. Odd etaT and subsequent prescribed local kernel/type nonvanishing remain open. Author whole-proof PASS 2026-11-21T23:30:15.183105+00:00. Exact certificate: \nolinkurl{literal_g_character_certificate.py}.

\subsection{ALP-LIFT-19: actual ramified theta constituents with prescribed nonterminal finite types}\label{alp-lift-19-actual-ramified-theta-constituents-with-prescribed-nonterminal-finite-types}

Full proof: \nolinkurl{ramified_column_type_chain.md}, T1-T24. At any degree-one split G place, an initial finite Steinberg type plus distinct literal tame stage characters gives irreducible actual local theta descendants and the exact inverse adelic column type at every j\textless m. Explicit finite Mackey vectors are primitive and have norm squared q\footnote{k(k-1)/2+d\emph{j}(j-1)/2+epsilon\emph{k}j} divided by the Gaussian multinomial index for(k,d\^{}j), with an exact11-valuation formula. DO033 supplies an old-host actual first input; LIFT14 extends analytically with unrestricted later finite Schwartz. Terminal HXm, raw Xi and global integral image remain open. Author whole-proof PASS 2026-11-21T23:49:36.760890+00:00. Exact certificate: \nolinkurl{ramified_column_type_certificate.py}. Independent DO033 review: \nolinkurl{review_steinberg_first_input.md}.

\subsection{ALP-LIFT-20: exact raw first G kernel and Steinberg extraction}\label{alp-lift-20-exact-raw-first-g-kernel-and-steinberg-extraction}

Full proof: \nolinkurl{raw_first_g_steinberg_bridge.md}, RG1-RG24. For m=2 and the actual allowed class y=q\^{}(-1)e\_k, both source PhiG/PhiGS give the primitive generic-unipotent residue function f\_U. Its full Steinberg component is (\textbar U\textbar\^{}2/\textbar GL\_k\textbar) f\_St, with exact compact-filter and first theta-extraction formulas. The raw finite Heisenberg coefficient absolute sum is1; raw norm squared is q\^{}(-bk)\textbar U\textbar{} and all11 averaging denominators are retained. Other raw T/Q factors, terminal PhiTop, later images and full DP44 costs remain separate. Author whole-proof PASS 2026-11-21T23:59:38.030210+00:00. Exact certificate: \nolinkurl{raw_first_g_steinberg_certificate.py}.

\subsection{ALP-LIFT-21: actual raw terminal G quotient preserves the preceding column}\label{alp-lift-21-actual-raw-terminal-g-quotient-preserves-the-preceding-column}

Full proof: \nolinkurl{raw_terminal_g_theta_quotient.md}, TG1-TG28. For m=2, the literal terminal PhiTop/ftY factor maps every nonzero vector in the actual LIFT19 preceding column into the actual small theta quotient. A compact full-rank integral and the primary Minguez boundary filtration prove the quotient identification. Exact identity value is vol(J)v (sharp vol(J)/q\^{}k); graph-plane norms are computed, and removing that volume gives a saturated injection of the natural induced11-lattices. The global automorphic scalar and exact row-lattice comparison remain open. Author whole-proof PASS 2026-11-22T00:26:12.923878+00:00. Exact finite diagnostic: \nolinkurl{raw_terminal_g_certificate.py}.

\subsection{ALP-LIFT-22: the two fixed sharp raw G edges compose on the actual middle level}\label{alp-lift-22-the-two-fixed-sharp-raw-g-edges-compose-on-the-actual-middle-level}

Full proof: \nolinkurl{composed_raw_sharp_g_chain.md}, CG1-CG21. Unwinding actual ChIdx.Ok forces the m2 sharp middle level Hsharp and its inverse character; midGSD uses this chF level, not generic ColV1. The first PhiGS forces the unique G class qy=e\_k. Its actual quotient image is the identity-parabolic section already in that middle type, and the fixed terminal sharp kernel preserves it. Their composition is injective on the initial generic type, with exact Haar scalar, composed norms, and saturated natural11-lattice image after scalar removal. Global quotient pins, other places/Y sum, and row-lattice comparison remain open. Author whole-proof PASS 2026-11-22T00:34:08.173546+00:00. Certificate: \nolinkurl{composed_raw_sharp_g_certificate.py}.

\subsection{ALP-LIFT-23: exact full coupled-Q Fourier contraction}\label{alp-lift-23-exact-full-coupled-q-fourier-contraction}

Full proof: raw\_coupled\_q\_fourier.md, Q1-Q20. At an unramified degree-one Q place the entire literal m2 Y sum collapses to the cofactor-supported Fourier function q\^{}(-d)G(chi,psi)chi(-c)\^{}(-1). The auxiliary depth cancels exactly. Joint Schwartz norm, primitive11 content, and scalar height (d-1/2)logq are explicit. This is a raw coupled-array identity; actual second theta-image survival and the global row scalar remain separate. Author whole-proof PASS 2026-11-22T00:47:32Z. Certificate: \nolinkurl{raw_coupled_q_fourier_certificate.py,10496} exact transforms.

\subsection{ALP-LIFT-24: actual fixed second T period on the first image}\label{alp-lift-24-actual-fixed-second-t-period-on-the-first-image}

Full proof: raw\_t\_projective\_period.md, MT1-MT33. The nontrivial projective-line source type is sigma\_spherical times one tame character. Whole-Schwartz GJ continuation gives the actual small theta quotient, and the raw T incidence vector has exact induced-model scalar L((d+1)/2,sigma1 eta1\^{}(-1))/(q\^{}n-1). Literal compact covariance transports the first T character into that same one-dimensional middle type. Every nonzero actual first T vector therefore survives the unchanged second T kernel. Haar denominator, L-factor valuation and global scalar comparison remain explicit. Author whole-proof PASS 2026-11-22T01:00:19Z. Certificate: \nolinkurl{raw_t_projective_period_certificate.py,58} compact vector checks and63958 Schur-complement checks.

\subsection{ALP-LIFT-25: actual selected compact highest kernels and their preceding-image contraction}\label{alp-lift-25-actual-selected-compact-highest-kernels-and-their-preceding-image-contraction}

Full proof: selected\_archimedean\_composition.md, AR1-AR26. The actual first highest polynomial outputs h1; the literal selected determinant insertion identifies the second highest source functional with h1 dual. The two compact contractions therefore preserve the actual matched line, including every multiplicity copy. Literal Haar coefficient factors are1/(d0 d1); the separate unitary quotient model gives squared factorA0²A1²/(d0 d1), with all Gaussian/CM/hook norms explicit. Exact11 Haar denominator valuations are1806 on K031 and14448 on K032. Global pins, nonhighest sharp witnesses and terminal row comparison remain separate. Author whole-proof PASS2026-11-22T01:14:48Z. \nolinkurl{Certificate:selected_archimedean_composition_certificate.py,810} exact compact projections plus polynomial and actual rank valuation checks.

\subsection{ALP-LIFT-26: full fixed sharp coupled two-edge automorphic nonvanishing}\label{alp-lift-26-full-fixed-sharp-coupled-two-edge-automorphic-nonvanishing}

Proof: \nolinkurl{fixed_global_coupled_second_lift.md}, GS1--GS31. Author whole-proof PASS 2026-11-22T01:30:35Z. On actual K031/K032 hosts and the unchanged R017 source, the complete sharp G/T/Q Y-coupled operator with LIFT25 selected highest real kernels is nonzero, retaining maximal vacua elsewhere. GQT is used in its anisotropic convergent case n=1802\textless m=3562/3436\textless2n. The exact partial standard factor is the initial degree-a Euler product at (d+1)/2, absolutely convergent with margin860/734, times Hecke L(l,chi0/chi1),1\textless=l\textless=d; the l1 boundary argument is the carried angular-type lemma. No global automorphic base change is assumed.

New assembly: local Howe multiplicity one and automorphic semisimplicity factor each actual global theta map as an injective multiplicity map tensor the product of local small quotients. Nonvanishing for every actual embedding makes that multiplicity map injective. The literal CRT Y quotient yields a tensor product of coupled local arrays. LIFT22, LIFT24, DO037 and LIFT25 give actual local left inverses on the initial types; LIFT16 supplies basic lines elsewhere. Distinct initial constituents cannot cancel after two edges by local Howe injectivity. The exact norm is a sum of positive constants c0*c1 times the squared local-tensor output norm; the constants remain the actual Rallis and quotient-mass costs.

Scope: full Y-sum nonzero, hence some admitted matched Y chain survives; the specific earlier Y0 is not asserted to survive. Nonzero G2 output does not prove its literal brC/tilQ/PtilFF restricted endpoint period, and the detecting Q graph lies outside that subgroup. Global pins, integral prior-image content, terminal branching detection, MainX mass and DP44 gain comparison remain open. The failed shortcut was inferring fixed-vector positivity directly from GQT11.6; the proof uses global factorization instead.

\subsection{ALP-LIFT-27: unchanged raw coupled operator is injective on the entire initial Mram}\label{alp-lift-27-unchanged-raw-coupled-operator-is-injective-on-the-entire-initial-mram}

Proof: full\_ramified\_raw\_coupled\_injectivity.md, MR1--MR33. Author whole-proof PASS 2026-11-22T02:01:09Z. This removes the finite Steinberg restriction at G by proving whole-Schwartz Godement--Jacquet holomorphy for every unitary initial representation and its actual first theta image. The same raw sharp G composition is injective on the entire prescribed inverse-generic space. The literal unweighted T norm kernel is the direct sum of all nontrivial projective dilation characters, each handled by the existing actual raw T theorem. The fixed full Y tensor consequently gives an injective full-G2 map on every actual initial Mram vector, including all multiplicities and rational stabilizer relations.

The exact complete-row coevaluation is d0\emph{d1 times the source-normalized highest-kernel operator, or d0}d1\emph{Lambda0}Lambda1 times its Chern version. Actual strict-covector/unipotent levels, stabilizer masses and weighted projective kernel averages are retained; coarsening does not make raw summands identical. K038's completed comparison applies after the literal terminal tilQ restriction and retains its own v0*v1 and period factors. Thus this does not prove injectivity/nonvanishing of that restricted arithmetic row or its brC/PtilFF period. It removes the need to identify a future actual bottom/Pu witness with the single R017 input. Uniform effective image costs and the actual bottom detector remain open.

\subsection{ALP-LIFT-28: fixed-level holomorphic bottom jets detect the full initial row space}\label{alp-lift-28-fixed-level-holomorphic-bottom-jets-detect-the-full-initial-row-space}

Proof: holomorphic\_bottom\_jet\_detector.md, HJ1--23. Author whole-proof PASS 2026-11-22T02:29:00+00:00. For either actual mixed CM host, the literal initial strict G/T/basic level extends to an explicit level on U(a,1), with inverse finite character for the bilinear period. Its compact U(a) packet is embedded with exactly the initial rational stabilizers, because the extra norm-one scalar is a root of unity forced to1 by the G unipotent residue level. The actual bottom HC list has selected initial K-type at the unique Pieri degree N\_tau=sum u\_tau-sum v\_tau+b0\_tau-a/2, independent of b1. Increasing all b1 by (a+1)ell changes the dual-minimal automorphic bundle by an ample canonical power and a flat determinant correction trivial on every arithmetic component. Serre vanishing on the Baily--Borel compactification therefore gives all compatible selected jets at the same finite level for sufficiently large ell in a fixed progression. PBW principal symbols transfer them to actual higher-K-type restrictions; their normalization factors are retained.

The actual cusp restriction is surjective onto the whole finite initial coefficient space, uniformly over any fixed finite word set. Every nonzero initial Mram row P thus has a bottom cusp test with exact period (v0/d\_lambda) sum\_x \textbar\textbar P\_x\textbar\textbar²/\textbar Gamma\_x\textbar. The test norm and sufficient weight are not bounded. The constructed detector has not been identified with source rJ/rChi/RealPack/Pu, and its cusp test has not been placed in the prescribed downward theta image. Consequently no literal MainX endpoint or DP44 cost comparison is claimed. The argument permits narrow-cloud interlacing profiles, but does not fit its ineffective weight increase into full SizeData2 windows.

\subsection{ALP-LIFT-29: exact source-shaped bottom level and T translation}\label{alp-lift-29-exact-source-shaped-bottom-level-and-t-translation}

Proof: bottom\_dat0\_translation\_detector.md, BT1--14. Author whole-proof PASS 2026-11-22T02:34:29+00:00. The actual rS=0 source PsharpSym(r=a+1,0,a) is exactly diag(U\_a,1), with generic U\_a character; dat0 has maximal compacts everywhere else. At T the explicit lower unipotent tau={[}{[}I,0{]},{[}varpi\^{}-1 u,1{]}{]} conjugates the maximal bottom compact so its intersection with H is exactly the original strict covector stabilizer. The translated packet x*tau has precisely the original classes/stabilizers. LIFT28 therefore gives large-weight cusp detection at this standard-frame source-shaped level, with no added T refinement and the unchanged initial mass. Translation is an ambient Petersson isometry, not an integral-lattice/height assertion.

The prescribed tS and betaY0 pins, actual two vertical kernels, and same-family downward image have not been identified with this concrete candidate. The new proof matches the finite level formula and initial-level intersection, not a completed RealPack or MainX endpoint.

\subsection{ALP-LIFT-30: primitive negative line and the finite/real bottom transport}\label{alp-lift-30-primitive-negative-line-and-the-finitereal-bottom-transport}

Proof: primitive\_negative\_line\_and\_bottom\_pin.md, NL1--14. Author whole-proof PASS 2026-11-22T02:48:20Z. The actual DO040 norm generator pi\emph{pi\_bar=T gives an explicit rational SU(a,1) matrix on a positive/negative plane, with columns (A,B)/pi and (B,A)/pi\_bar, A=(T+1)/2 and B=(T−1)/2. It sends pi\_bar times the negative coordinate to the primitive integral vector (B e1,A) of norm minus T. At T its inverse factors exactly as the LIFT29 lower unipotent for covector pi\_bar e1}, times an integral determinant-one compact element; away T it is integral and at G it is identity modulo the level. An arbitrary actual primitive T covector is obtained by an explicit integral finite H-coordinate transport, with unchanged strict class masses and Mram.

The actual identity of automorphic restrictions also transports the real vector by sigma(g\_infinity); omitting that factor would be false. Its Hilbert norm is unchanged, while the ambient matrix has singular values sqrt(T),1/sqrt(T), so algebraic height or coefficient costs are not set to zero. At 11 this matrix and inverse are integral. Exact Q(sqrt(-7)) replay for both actual DO040 primes passed24 matrix/congruence checks. Source audit distinguishes the full row compact pin tS from this Bessel embedding: rowK2 is maximal away G, and RowLevelMatchB/C further fixes the full residue action and spectral idempotent. None of these source ties, the idempotent-selected detector, or the fixed downward image is asserted complete.

\subsection{ALP-LIFT-31: the fixed vertical adjoint is injective on the full bottom space}\label{alp-lift-31-the-fixed-vertical-adjoint-is-injective-on-the-full-bottom-space}

Proof: vertical\_adjoint\_surjectivity.md, VA1--VA33, including VA22a--c, VA24a and VA31a--c.~Author whole-proof PASS 2026-11-22T03:21:43Z. For the actual mixed hosts, the unchanged conjugated DatumVec kernels give an injective upward map from the entire fixed-level cuspidal multiplicity space on U(a,1), through U(1802), to U(3561) or U(3435). This allows every initial principal-G constituent, with maximal level elsewhere; it is not restricted to a finite Steinberg or generic subspace. The two vertical increments are odd t=1759 or1633. The unequal-label whole-Schwartz quotient proves actual G and Q prior-image preservation, including exact Haar scalars and induced compact norms. The odd-rank basic test is established on the exact lattice vacuum by Gong--Grenie, Harris and the local doubling calculation; absolute ramification is retained. Both global steps apply to every actual multiplicity embedding, with Euler margins859/733 and boundary margin1.

A direct absolutely convergent Fubini/two-Schur calculation proves that the source conjugated upward maps are the exact Petersson adjoints, without assuming the source predicate or discarding measure/intertwiner factors. Hence the actual downward map is onto the full bottom space, and its postcomposition by any genuine source idempotent is onto that idempotent summand. LIFT28/29's detector therefore belongs to the downward image whenever its large-weight existence and the real-type conditions hold. Its survival under the particular e0 projection, the remaining global/source frame pins, the effective SizeData2 fit, and the norm/11-content of a preimage remain open. The formula D\emph{(DD})\^{}(-1) is a qualitative right inverse, not an effective cost bound.

Exact finite diagnostic: \nolinkurl{vertical_g_adjoint_certificate.py/.json,90,024} signed raw G-composition checks on11,400 finite target matrices, both raw and induced norms, and both actual rank margins. Independent complete review of DO041 and its certificate is saved in review\_vertical\_q\_quotient.md.

\subsection{ALP-LIFT-32: finite generic branching with exact projected inverse and central compatibility}\label{alp-lift-32-finite-generic-branching-with-exact-projected-inverse-and-central-compatibility}

Proof: finite\_generic\_branch\_restriction.md, FB1--FB21. Author whole-proof PASS 2026-11-22T03:52:33Z. For the actual completely split G groups, the generic mirabolic summand gives every generic GL\_a type inside any generic GL\_(a+1) constituent. The exact probability-normalized restriction operator has eigenvalue d\_Pi\emph{m\_sigma/d\_sigma; its canonical right inverse has norm squared sum\_sigma d\_sigma/(d\_Pi}m\_sigma)\emph{\textbar\textbar w\_sigma\textbar\textbar². The coefficient formula retains m\_sigma}\textbar H\textbar{} denominators, so no 11-integrality is inferred. The full extra torus character is necessarily omega\_Pi/omega\_sigma, and the right inverse preserves all actual finite stabilizer-equivariant fibres on that compatible subspace. Genuine squarefree Green-series projectors are generic; the exact central-character/norm-of-root dictionary fixes the Deligne--Lusztig sign. DO042/043 supply the separate actual goodTuple and aggregate mu28 compatibility construction.

The finite result is an ingredient for the parent's global projected jet packet; no global detector or effective theta inverse follows by itself. Exact certificate \nolinkurl{finite_generic_branch_certificate.py/.json} passed at03:38:32UTC: three mirabolic models, both GL2(F3)/GL3(F2) central restriction examples, 216 convolution checks, 3,408 equivariance checks and 96 incompatible-central annihilations.

\subsection{ALP-LIFT-33: actual uniform denominator improvement and non-scalar invariant pairing}\label{alp-lift-33-actual-uniform-denominator-improvement-and-non-scalar-invariant-pairing}

Proof: sharp\_stabilizer\_integral\_pairing.md, SI1--SI23. Author whole-proof PASS 2026-11-22T04:11:35Z. The explicitly certified prime18,895,241 splits both CM hosts, has order10 modulo11 and first lifting valuation1. Reduction injects every11-Sylow of an actual rank-n lattice stabilizer into GL\_n(F\_q), proving the degree-independent bound sum\_j floor(n/(10*11\^{}j)). The sharp G flags force at least a+1 fixed dimensions for each11-Sylow, proved by an integral Reynolds idempotent at G. With K042 initial neatness, the unchanged coupled map loses at most191/178 from middle masses; the terminal total losses are575 for J14 and535 for J56. This improves the post-K042 combined lower bounds by188,625 and738,503, without changing any field, rank, tuple, kernel or lattice.

For the admitted trivial tuple, a concrete standard-lattice middle class has stabilizer S\_t, t1759/1633, through its actual sharp G flags and an allowed strict-T decoration. The actual selected Weyl tensor has a primitive fixed line and a primitive regular S\_t orbit because fa\textgreater=t−1 and its first selected entries are distinct positive. Its actual invariant Weyl/dual-Weyl pairing has direct block diag(1,t!), and the class-mass pairing has exact Smith endpoint valuations−174,0 or−162,0; all its factors lie between these endpoints. The metric norms1,t! and actual mass are retained. No common scalar removes this non-scalar width. This does not place either coefficient direction in the preceding theta image or compute a global primitive minor.

Companion certificate \nolinkurl{sharp_stabilizer_integral_pairing_certificate.py/.json} PASS04:17:03UTC after the sharp-level refinement: complete full-order primality certificate plus trial division,3,561 actual order-factor valuations, all numerical bounds, both joint-weight constructions,131,712 mu28 phase checks and34,406 small regular-action entries with exact invariant Gram/Smith checks. No Lean build.

\subsection{ALP-LIFT-34: actual coupled-image cancellation of every middle class denominator}\label{alp-lift-34-actual-coupled-image-cancellation-of-every-middle-class-denominator}

Proof: coupled\_image\_sylow\_cancellation.md, MC1--MC22, including MC21a--b. Author whole-proof PASS2026-11-22T04:38:10Z. At one actual G place the raw second PhiTop support has an invertible leading b-by-b block. The actual middle11-Sylow therefore acts freely on its residue support. A0/1 finite transversal cuts the complete fixed Y-coupled tensor, retaining its integral values and identical11-vacuum. Complete integral middle subD coevaluation is invariant under the real middle group. Honest rational theta automorphy identifies the uncut class numerator with the group order times the cut-section derivative. Its direct-sum prime-to11 Heisenberg expansion separates into products, preserving the exact L0+L1 gain. Thus the actual coupled image loses no middle11 denominator. This applies to all initial MfV and its Mram lattice, without a Steinberg restriction.

The actual G primes also exceed the entire rational middle dimension. A cyclotomic-degree argument therefore makes reduction injective on the full stabilizer. The same construction represents the contribution divided by its full mass as an independent integral cut section, with exact finite norm squared reduced by \textbar Gamma1\textbar. This is an actual image-bound result, not a claim that the ambient invariant pairing of LIFT33 is unimodular or that a primitive theta minor exists. Terminal losses384/357 and the literal branch/detector cost remain. Common periods, complete-coevaluation normalization, and source creation scalars are retained.

Certificate: \nolinkurl{coupled_image_sylow_cancellation_certificate.py/.json} PASS04:37:32UTC,320458 exact assertions: order11/free sharp matrix action overF23, full-group transversals, eta phases, all529 cyclotomic Fourier inversions, integral Sym2 coevaluation and determinant character, exact finite norm and actual G dimension margins. No Lean build or global theta matrix computation.

\subsection{ALP-LIFT-35: actual terminal fixed residue prevents the same transversal argument}\label{alp-lift-35-actual-terminal-fixed-residue-prevents-the-same-transversal-argument}

Proof: terminal\_residue\_cutoff\_obstruction.md,TC1--TC5. Author whole-proof PASS04:40:57UTC. Literal RowPt has empty T-decorations. For the trivial tuple the standard normalized sharp terminal class has stabilizer (mu28\^{}t semidirect S\_t)\^{}2. An11-cycle in its second S\_t acts only on the last t target rows, which the actual raw PhiTop support forces to zero modulo every G prime. It therefore fixes every supported residue, even though the second matrix has full column rank. No0/1 first-residue cutoff can have orbit sum1. This is a precise obstruction to extending LIFT34 by the identical finite mechanism, not terminal nonintegrality, period vanishing, or a Smith-factor claim. Actual brC, tilQ and detector fixed strata remain separate.

\subsection{ALP-LIFT-36: explicit actual boundary Hecke factor and full-string cost}\label{alp-lift-36-explicit-actual-boundary-hecke-factor-and-full-string-cost}

Proof: actual\_boundary\_l\_value\_cost.md, LB1--30. Author whole-proof PASS 2026-11-22T05:00:46Z. Tate functional equation, an explicit strip/derivative bound and the positive 3+4cos+cos2 Euler polynomial prove \textbar L\_E(1,theta)\textbar{} \textgreater= {[}1024 A0\^{}3 B2 B1\^{}3{]}\^{}-1 for the actual angular stage ratio, with all conductors and gamma shifts displayed. On every K044 profile this gives \textgreater2\^{}-277051 for J14 and \textgreater2\^{}-1171147 for J56. The whole inserted string at arguments2 through t costs at most exp(n), independently of t. Including both initial-parameter factors gives the actual partial-standard-factor product \textgreater2\^{}-277221 R\_S or \textgreater2\^{}-1171821 R\_S, with the exact finite omitted-prime correction R\_S retained. No ineffective zero-free constant, global base-change assertion, or temperedness assumption is used.

This closes one quantitative component of the LIFT31 inverse problem. It does not equate the finite induced quotient norm with the globally normalized Rallis form or suppress actual real-intertwiner/Haar factors, so a full lower singular value remains separate. Exact arithmetic certificate passed1842 assertions, including588 unit-circle positivity tests and both actual field/conductor exponent calculations.

\subsection{ALP-LIFT-37: exact actual G-place Rallis factors and finite inverse cost}\label{alp-lift-37-exact-actual-g-place-rallis-factors-and-finite-inverse-cost}

Proof: vertical\_g\_rallis\_normalization.md, VG1--17. Author whole-proof PASS 2026-11-22T05:06:55Z. Actual raw G full-residue-rank support makes both orbit periods and doubled Gram integrals compact, even at the second b\textless N\textless2b edge. Exact matrix disintegration gives the missing Rallis constants C0=kappa\_b/kappa\_t and C1=kappa\_N/kappa\_t. Their composed raw Hermitian factor is q\^{}(-(b+r)\^{}2) kappa\_t\^{}2/(kappa\_b kappa\_r). Dividing by both actual raw Schwartz norms yields exactly c\_G=q\^{}(-r(b+r)) kappa\_t/(kappa\_b kappa\_r), not c\_G squared. This holds on every prescribed initial principal-G vector.

Products over all actual G places give raw lower bounds2\^{}-6433607334/2\^{}-46338383784 and unit-Schwartz lower bounds2\^{}-149943150/2\^{}-3973204872 for J14/J56. The exact products are retained and independent of large weights. Q/basic/real/global Haar factors remain to be assembled for a full inverse; no integral-preimage assertion. Certificate647 exact checks passed.

\subsection{ALP-LIFT-38: quantitative actual vertical preimages with complete local and global factors}\label{alp-lift-38-quantitative-actual-vertical-preimages-with-complete-local-and-global-factors}

Proof: vertical\_preimage\_norm\_bound.md, VI1--28. Author whole-proof PASS 2026-11-22T05:30:11.133952+00:00. The actual LIFT31 downward composition D has a least-norm complex right inverse R=U(U\emph{U)\^{}(-1), with U=D}. Its isotypic positive energy is exactly m\_b\emph{m\_N times the actual real reciprocal formal-degree/compact-dimension product, the exact raw G and Q factors, the restored standard numerator and the basic doubling denominator inverse. Actual local Haar forms make the global multiplier m\_b}m\_N; no extra a0\emph{a1²}a2 or level index is inserted. Q Euler factors are retained exactly once. The formula applies to every initial automorphic multiplicity copy and every principal-G vector.

LIFT36 supplies the positive numerator bound; LIFT37 and DO047 give all G/Q Hermitian constants; the basic denominator costs less than2. DO048 fixes the real formal degree exactly as a! dimV\_(HC-rho), including the ball Haar and compact centre. The literal Dpar1 list and all K044 weight bounds give an explicit real-dimension bound. For h\_Q=\#Q and Qnorm=product q\_v, the raw right inverse satisfies \textbar\textbar R\textbar\textbar{} \textless{} \textbar s0*s1\textbar{} 2\^{}((H+8h\_Q)/2) Qnorm\^{}(2rt), with H=71,897,687,205 and rt=75,637 for J14, H=310,338,390,993 and rt=275,977 for J56. Unit-finite kernel exponents are65,614,023,021 and267,973,212,081, with4h\_Q and Qnorm\^{}(rt/2). Exact formulas are stronger than these coarse numbers.

The actual source real-intertwiner scales s0,s1 remain explicit. Isometric analytic normalization sets them to1; arbitrary prescribed rescaling prevents a uniform bound independent of them. This closes the former qualitative complex top-preimage cost, but proves neither an11-integral preimage nor an upper non-scalar Smith index or the remaining source/terminal pins. Certificate4561 assertions passed at 2026-11-22T05:29:54.730553+00:00; proof SHA256 2f0814e8e68eaab439aa258abd441242e95afc7cddc8fcf7d613b5b2a002e392.

\subsection{ALP-LIFT-39: literal unit vacuum and explicit primitive vertical normalization}\label{alp-lift-39-literal-unit-vacuum-and-explicit-primitive-vertical-normalization}

Proof: vertical\_primitive\_fock\_normalization.md, PN1--25. Author whole-proof PASS 2026-11-22T05:56:35Z. Exact adelic Fourier orthogonality proves the Heisenberg theta realization is an isometry under the literal covolume-one normalization. IsVacO therefore requires multiplying the K035/K037 raw edge Gaussian by det(2Y\_J)\^{}(7hnm/2); the combined two-edge amplitude is det(2Y\_J)\^{}(7hb²), with the first negative-column period and creation factors conjugated explicitly. Existing raw CM formulas are unchanged.

Actual primitive joint harmonics P0=p\_A(X)y\_b\^{}B and P1=p\_Lambda(Z) have exact norms H\_A B! and H\_Lambda, and their complete vacuum characters equal literal Dpar1/Dpar2. Constructed primitive intertwiners sending these joint extremals to the normalized row extremals are rational on the abstract polynomial Harish-Chandra modules and satisfy IsItW. With the literal unit vacuum, \textbar s0*s1\textbar²=(product\_tau H\_A B! H\_Lambda)\^{}(-1). K043 PBW norms occur on both sides and leave this scale unchanged. A block-Vandermonde identity computes exact hook valuations without expanding the large polynomials.

LIFT38 now has an explicit primitive-normalization inverse bound with this hook product replacing its previously free real scalar. Raw-vacuum repinning adds the displayed CM amplitude; an explicit fixed-physical-test versus raw-test calculation prevents double counting. Rationality is not integral saturation or a net DP44 gain, and the remaining arithmetic row/CM frame pins are not asserted. Certificate1055 assertions PASS 2026-11-22T05:56:03.310848+00:00; frozen proof SHA256 2a4be598e272378ef47efe08e4a1b7b88265e40121290d8e5aefaaf1882e976b.

\subsection{ALP-LIFT-40: actual primitive vertical coefficient lattices and ambient projection poles}\label{alp-lift-40-actual-primitive-vertical-coefficient-lattices-and-ambient-projection-poles}

Proof: vertical\_primitive\_integral\_lattices.md, IL1--24. Author whole-proof PASS 2026-11-22T06:20:33.351620+00:00. The intrinsic ordinary-PBW lattice H=C(Delta\_alpha tensor Z{[}z{]}) is fixed before its oscillator embedding. K043 adjunction identifies its Shapovalov dual exactly with Nabla\_alpha tensor divided monomials. The primitive joint embeddings preserve the generated lattices, and PN15 has explicit integral retractions on them, including a two-edge section.

On the saturated actual rational Howe summands, moving integral compact divided powers and ordinary annihilators to a primitive joint weight proves normalized Fischer integrality; the large joint hook norms N0,N1 disappear. The first contraction lands in the explicit PBW dual and the second in the middle dual-Weyl lattice. At selected degree n and type lambda, n! H\_minimal H\_middle D\_lambda is a sufficient two-edge PBW denominator, with an additional H\_lambda for the intrinsic selected Weyl comparison. All factors are explicit factorial/hook/PBW products; the primitive-to-unit selected scalar sqrt(D\_lambda q\_lambda) and CM/conjugate frame transport remain distinct.

The whole ordinary Fock projection is not integral. Actual coefficient-one diagonal monomials have exact images r0=product A\_i!/H\_A and r1=product Lambda\_i!/H\_Lambda; every admissible endpoint profile has negative11 valuations. The report computes all24 actual q=0/q\_max profiles, not the larger formal q=ell envelope tested in frozen LIFT39. Exact small models distinguish an ambient11 pole from a necessary dual-Weyl enlargement inside a saturated selected summand. These are local coefficient-map results, not global theta-image saturation or an integral automorphic inverse. Certificate4694 checks passed 2026-11-22T06:19:35.232575+00:00; frozen SHA256 e7d9595772b58af4fb1e94ffd29c5ba0d97e4cd61f8a879c16e0fabac9386f66.

\subsection{ALP-LIFT-41: an actual near-origin packet point and quantitative selected one-point inverse}\label{alp-lift-41-an-actual-near-origin-packet-point-and-quantitative-selected-one-point-inverse}

Proof: near\_origin\_actual\_packet.md, NP1--20. Author whole-proof PASS 2026-11-22T07:00:23Z. For the unchanged actual primes and pi, set Q=g0*g1, pi\_k=pi−7Q²+2Q²sqrt(-7); its norm is T+77Q⁴−14Q². Explicit four-square coordinates give a primitive vector v of norm−T, principal-congruent to the negative line. The rational unitary matrix with negative column v/bar(pi) has the exact lower-T finite double coset, including its conjugate component and all other finite places. This gives a finite representative above an actual initial class in the principal-G refinement of the actual packet.

The resulting real point has sinh²s\textless2\^{}-813(J14) or2\^{}-704(J56). The actual selected rank-one coefficient has an exact finite derivative-norm polynomial. At the upper admissible gap its top-label coefficient exceeds exp(−2\^{}148) or exp(−2\^{}257); thick/thin polynomial values remain unbounded away from zero here. More substantially, conjugating the standard principal arithmetic group by this near-origin matrix preserves radius-five separation: every entry of h−I is divisible by Q but every conjugate would be \textless e\^{}11+1\textless Q. K046 then gives surjectivity onto the complete normalized selected fibre at this actual packet point for every word and admissible twist at ell\_plus, with right-inverse norm \textless=sqrt(V\_G)/(1−2f*2\^{}(-ell\_plus/8)). This does not assume a product of the rank-one coefficients is nonzero.

The genuine finite e0/Whittaker projection, simultaneous packet interpolation and permitted terminal/global period remain open. No level, prime, T covector or canonical twist changed. Exact certificate7019 assertions passed 2026-11-22T06:59:46.231244+00:00; frozen proof SHA256 1e6073f62c0e934f2c569f4dc277751702086102fccfd59ba60612c73afff298.

\subsection{ALP-LIFT-42: exact finite-character bounds for the actual near-origin packet}\label{alp-lift-42-exact-finite-character-bounds-for-the-actual-near-origin-packet}

Proof: finite\_character\_short\_orbit\_bound.md, FC1--15. Author whole-proof PASS 2026-11-22T07:12:50Z. Green's exact cuspidal formula gives the sharp nonscalar ratio1/(g\^{}a−1). After retaining arbitrary orthogonal partial-trace projectors at the two mixed G places, an algebraic-integer norm bounds the weighted number of scalar reductions: each selected scalar reduction forces vanishing at both CM partners. The complete finite origin stabilizer has size at most28\^{}(a+1)a!, and radius2/5 coarse orbit packing bounds every short correspondence.

For the actual LIFT41 packet point, radius9/2 on J14 and5 on J56 gives total normalized noncentral short character mass below2\^{}-1200 and2\^{}-646000. These are exact integer comparisons, with all stabilizer terms retained. The corresponding cardinal is below2\^{}(32af), while the accepted selected concentration tail is below2\^{}(-64af−20); J14 uses the proved refinement epsilon1/2000 and K046's moment bound969. These supply strict constants for a genuine-projector Poincare argument, whose selected-fibre positivity is separate. No Mram or terminal period is asserted here. Certificate1853 assertions passed 2026-11-22T07:12:11.939277+00:00; frozen SHA256 1fdeb38635f6a21668e64d8050adeb9ba57d205026222a3e3c47afd7cc8536ff.

\subsection{ALP-LIFT-43: exact actual T-projected packet Gram and a finer negative-column restriction}\label{alp-lift-43-exact-actual-t-projected-packet-gram-and-a-finer-negative-column-restriction}

Proof: actual\_t\_packet\_projection\_gram.md, TG1--16. Author whole-proof PASS 2026-11-22T07:33:31.235487+00:00. The actual dual detector Mram projection is the product of the nontrivial-character projections at every selected T place. Its exact Hilbert--Schmidt energy is a mass-weighted inclusion--exclusion sum of cross-centre selected kernels; the formula retains nonfree projective stabilizers and opposite coefficient conventions. One-point H-generic positivity does not determine the other Gram terms.

A principal-G rational SU transvection I+(iQ/T)vv\emph{J, v=e1+e\_n, belongs to the correct ambient relative Cartan double coset and has displacement below2\^{}-953/2\^{}-923, while its G partial trace equals1. Thus the integral short-character bound cannot extend to the ambient Hecke set. The actual LIFT41 scalar packet excludes this example: Rz must remain integral, and its exact negative residue r has norm−14Q². More generally every admitted primitive rank-one family I+(iQ a0/T)vv}J has geometric-mean displacement greater than1/(154Q³), hence at some label greater than2\^{}-143/2\^{}-221. Arbitrary correspondences and the full Mram Gram remain unresolved. Certificate866 assertions passed 2026-11-22T07:33:31.737017+00:00; frozen SHA256 324e1c4b41e075ffa6a811b980eea150069b853ceaea814ca53fac4a45f37131.

LIFT43 final precision: nontrivial tame central character implies conductor one; the converse is not used. Amended whole-proof reread 2026-11-22T07:35:48.819502+00:00,866-check replay 2026-11-22T07:35:49.346127+00:00, final SHA256 c284632159a65869656bf5a39c688593bc4f5339edba963e5a143be410b6c0ed.

\subsection{ALP-LIFT-44: actual fixed-complement pencil, residual plane and small semisimple member}\label{alp-lift-44-actual-fixed-complement-pencil-residual-plane-and-small-semisimple-member}

Proof: fixed\_scalar\_packet\_plane\_reduction.md, PR1--19. Author whole-proof PASS 2026-11-22T07:57:42.743198+00:00. The exact selected-place inverse-image lattice is (1+c\emph{pi\^{}-1}z\emph{qrow)Lambda, with c=s\^{}-1−1, qrow=−v0}J−pi\emph{m and m}z=1. Its conjugate is the exact Hermitian dual. This criterion is also sufficient for a rational relative matrix to realize the fixed scalar-packet member. All actual correspondences map the entire integral positive complement into the ambient lattice, including partial T subsets. Operator displacement below1/(10Q) forces fixation of the short(a−1)-dimensional positive part and leaves only the determinant-one plane diag(W,−1).

A general plane B\emph{S\^{}-1 has an explicit ambient U-divisibility condition and exact graph s=xprime}bar(pi\_k)/(Q\emph{x); the full−pi}m term is essential. Root supplied an allowed s=−1 semisimple member, independently verified here, with operator displacement below2\^{}-404/2\^{}-350. Thus uniform small-operator separation is false; its point is related by the integral S0=diag(−I\_a,1), whose G-deck/selected-frame action remains. On J14 the residual plane is locally anisotropic at3, so rational unipotents are absent. This is a precise local reduction and obstruction, not Mram positivity. Certificate9946 assertions PASS 2026-11-22T07:57:42.946294+00:00; frozen SHA256 0f326a7d81dcb8e92da69b990e04a5cc49b1ac25908842f781d7cd251fa26ed8.

\subsection{ALP-LIFT-45 --- single-place scalar separation and the all-T energy criterion}\label{alp-lift-45-single-place-scalar-separation-and-the-all-t-energy-criterion}

Full proof: \nolinkurl{single_place_packet_separation.md}; frozen SHA \nolinkurl{c7b0dc2cddfed44a8e735b7951dd38c5476511b85b502b9df9c594fcfda34f7f}. Author whole-proof PASS08:22:00UTC; exact certificate \nolinkurl{single_place_packet_separation_certificate.py/.json}, 21,498 assertions PASS.

The literal mass-weighted commuting projections satisfy \nolinkurl{E_ram >= E - sum_v E_v}, with every strict/projective volume and cross term displayed. A single-place scalar change at the favorable point has distance \textgreater2\^{}-25 (J14) or \textgreater2\^{}-7 (J56) in some real factor from every integral translate. The proof uses only one CM denominator pair and the actual mixed rank-one pole to exclude a compact block matrix. It does not replace the full all-base averaging energy by that favorable-base estimate, and proves no Mram nonzero or DP44 saving.

\subsection{ALP-LIFT-46 --- exact character divisors and bounded Galois CM coefficient fields}\label{alp-lift-46-exact-character-divisors-and-bounded-galois-cm-coefficient-fields}

Proof: mixed\_character\_field\_costs.md, CF1--29. Author whole-proof PASS 2026-11-22T08:55:05.710437+00:00, frozen SHA256 7ad5bff80d53629b662bf9fb04a4fa723b4baaf628f8e13e788ce9d2c8daa4ca. For the unchanged K040 characters, the actual ideal divisor is product\_tau(tau A/tau bar A)\^{}e\_tau, giving exact finite content and height at most (sum\textbar e\textbar/d)log N A, with no class-exponent loss. Normalized odd etaVal has height exactly half log N A. The induced primitive-J type reduces the odd value-field construction to three radical families. Together with the full angular orbit, one Galois CM field contains every normalized eta and every coherent stage value, degree \textless=d h\_K\^{}(f+4), root discriminant \textless=rd(K)h\_K\^{}(f+4), unramified over K outside class-number primes. The different proof handles nonmaximal monogenic orders explicitly. Exact Minkowski/zeta counting gives h\_K\textless2\^{}456/2\^{}2477 and with actual G/T roots gives degree exponents22018/427103 and root-discriminant exponents22019/427105. All bounds are weight-independent; Q roots retain their own exponent. Normalized content is invariant under coefficient extension. No literal Host completion, actual period, theta Smith factor or DP44 net gain is claimed.

Certificate: \nolinkurl{mixed_character_field_costs_certificate.py/.json,60922} exact assertions PASS 2026-11-22T08:54:21.350361+00:00. Actual CM orbit, divisors/base cancellation, source normalization, local valuation rescaling, base change, and integer bounds; no class-group or coefficient-field computation.

\subsection{ALP-LIFT-47 --- linear-rank transference bounds for the actual adapted frames}\label{alp-lift-47-linear-rank-transference-bounds-for-the-actual-adapted-frames}

Proof: transference\_adapted\_hermitian\_frames.md, TF1--21. Author whole-proof PASS 2026-11-22T09:11:07.611069+00:00, frozen SHA256 ec47e5e9683b578801fe9cf42ba9de478abf1b8748a843b3c949872d416d8673. The exact Euclidean dual is2 D\_E\^{}-1 N, and its nonzero vectors have trace norm at least2 sqrt(f)/rd(E), by primewise fractional-ideal norms including ramified2. Banaszczyk covering-radius transference gives mu(N)\textless=n sqrt(f)rd(E)/2. Reducing prescribed residues directly modulo11N constructs an actual eleven-adapted integral frame with column norm at most(11/2)n sqrt(f)rd(E), without the auxiliary short Z-basis. Independent determinant sections satisfy the rank-one version.

The unchanged actual Delta(bar L)/terminal Nabla generators now have binary degree costs27D0+31DT+13(M0+MT) on J14, or36D0+35DT+17(M0+MT) on J56, plus the exact R039 signed Lambda0 log2(T)/2 and unchanged finite projector factor. This strictly improves DO057 baseline coefficients without changing any local lattice, class, character, coefficient field or theta kernel. It proves no actual Mram period, image saturation or complete DP44 inequality.

Certificate: \nolinkurl{transference_adapted_frames_certificate.py/.json,76522} exact assertions PASS 2026-11-22T09:10:14.363446+00:00.

LIFT47 final scope clarification and complete reread PASS 2026-11-22T09:13:07.852560+00:00: the improvement is in a sufficient upper bound for integral-input cost. Formulas/certificate unchanged. Final frozen SHA256 825722c29b3674d48b971e24e544afb383ceb75cbfa4f1bfede059c02d0a0336.

\subsection{ALP-LIFT-48 --- an explicit source-normalized coupled theta operator bound}\label{alp-lift-48-an-explicit-source-normalized-coupled-theta-operator-bound}

Full proof: actual\_normalized\_theta\_operator\_bound.md, OB1--28; author whole-proof PASS 2026-11-22T09:33:20.036961+00:00, frozen SHA256 4b7b3854bdb646fd7dcdb2317ce0c3ecab061b8fb8c77ba9f466fe92f1d6cbe4. The unit-vacuum Fourier/Zak square has exact zero coefficient \textbar\textbar r\textbar\textbar² and symplectic-dual lattice Lambda, with its actual semicharacter. Normal ordering and the principal codifferent length bound 2f/rd(E) give U(D)={[}(C+1)²(D+1)+16n rd(E){]}\^{}(fn/2), n=2b², C20/28. For every K044 word/gap log2 U\textless496fn, or2\^{}37/2\^{}39. The full actual finite tensor has Heisenberg l1 costs q\^{}(n/2) at T/G and q\^{}((n+a−t)/2) at Q; the latter evaluates the entire prescribed Y sum and is depth independent.

The actual arithmetic bilinear operator is bounded by Cfin U sqrt(d1 H0 H1)\textbar bBr\textbar/\textbar M\_sigma\textbar{} times m1 sqrt(m0 mT), with the actual mass-weighted class norms. All hook/branch, CM period, creation scales and strict masses remain. This removes an unspecified theta operator norm without changing input data. Actual base genus masses and the moved Mram nonzero remain separate; no primitive Smith factor, full old Host or net DP44 saving is claimed. Exact certificate \nolinkurl{actual_theta_operator_certificate.py/.json} passed33054 assertions at2026-11-22T09:32:39.509610+00:00.

\subsection{ALP-LIFT-49 --- actual genus-mass and fixed operator costs}\label{alp-lift-49-actual-genus-mass-and-fixed-operator-costs}

Proof actual\_genus\_mass\_operator\_cost.md, GM1--21; author whole-PASS 2026-11-22T09:43:10.256759+00:00, frozen SHA256 67a2182580323db326b56ccdb42e96fef0330f58e6eeb078640ffe15d7a9d072. For every positive unimodular rank-n lattice, the TF47 frame gives at most (1+2Rn²)\^{}(fn²) Gram matrices. The finite discriminant module has size \textless=(Rn²/f)\^{}(2fn), and each O\_E-submodule has \textless=n generators by local DVR freeness plus finite CRT, giving at most its n-th power submodules. The actual standard genus is a subset of this broader count, never equated to it. This proves m2\textless2\^{}12768/2\^{}64512.

The accepted standard mass ratios give log2{[}m\_b sqrt(m\_a m\_N){]}\textless3,623,548,215/16,156,666,848. Combining the unchanged strict volumes and LIFT48 finite/theta inputs gives fixed binary operator cost284,861,333,421(\textless2\^{}39)/1,162,551,608,736(\textless2\^{}41), plus each actualQ term3+E\_Q log2q with E\_Q=6,795,239/2 or7,596,851/2. The exact sqrt(d1H0H1), branch, M\_sigma and integral input factors remain; no Mram period or net budget saving. Certificate \nolinkurl{actual_genus_mass_operator_certificate.py/.json} passed17731 assertions at2026-11-22T09:42:57.494499+00:00.

\subsection{ALP-LIFT-50 --- normalized selected hook costs and eleven certificate}\label{alp-lift-50-normalized-selected-hook-costs-and-eleven-certificate}

Proof: \nolinkurl{normalized_hook_gain_budget.md}, HG1--37. Author whole-proof PASS 2026-11-22T10:09:15.764464+00:00; frozen SHA \nolinkurl{5a5b7e36caf1ee127c1620e172bc1beff87a66ee48d42d45d5a3bebceafdc78e}. The full algebraic branch scalar sB and coordinate scalar CD are removed together with every finite valuation. Their exact modulus/phase and all remaining bad-prime contents are retained. The actual normalized mean creation factors are computed from K066. R038 at the actual terminal rank gives integral allowed Weyl field-spanning detectors, with terminal cost coefficients43/47; conditional on the same nonzero Mram row they avoid the terminal hook loss of the full dual-Weyl detector. The proved gain satisfies D/1221\textless gDelta\textless=D/1220; the existing dual-Weyl lower-bound formula is less than−D/42, a statement about that formula, not actual content. Exact hook, dimension and sparse branch products show the intrinsic complete-coevaluation/period norm factor exceeds the credited11 gain by more than220D binary units, while an explicit sufficient upper budget retains all finite tuple, genus, input and signed T costs. This is not a theta-operator lower bound, an image-content upper bound, or a DP44 no-go. Certificate5085 exact assertions passed10:06:41UTC; no Lean build.

\subsection{ALP-LIFT-51 --- concrete small-error link incompatibility}\label{alp-lift-51-concrete-small-error-link-incompatibility}

Proof: \nolinkurl{concrete_source_link_incompatibility.md}, LI1--18; author whole-proof PASS 2026-11-22T10:28:11.792218+00:00, frozen SHA \nolinkurl{efc58c4c75ac223e01a5e98dc39f43b08ae42b2c84e8dbcf13ed0cb1e035c7ab}. The literal source has KD=c Pu*Pu through full HH, dimCu=duS, and Drows=(1−delp)Dall. Exact unnormalized Hilbert--Schmidt rank deficiency gives Drows/duS \textgreater=(1−delp)(1−e0) whenever the chosen proviso activates the full LinkS diagonal estimate. For x\textgreater1 and positive eps/ez, LinkIRR and LinkS force that ratio\textgreater3/8. The accepted DO044 upper64-bit gap cohort instead has ratios\textless2\^{}-482412(J14)/2\^{}-414120(J56), all words/roles/allowed canonical parameters included. The full diagonal defect is\textgreater1−2\^{}-(L−1) per initial dimension; even its Mram compression is\textgreater1−2\^{}-(L−2). Any retained replacement input has dimension\textless3Drows.

The theorem grants an arbitrary full HH and does not identify it with projected EE. It excludes pointwise simultaneous literal links on the concrete cohort, not every asymptotic family: all relevant eventual thresholds must actually apply. It proves no period vanishing, universal DP44 obstruction, or changed-baseline saving. Exact certificate44,358 assertions passed 2026-11-22T10:26:54.541210+00:00.

\subsection{ALP-LIFT-52 --- necessary profile changes for literal link compatibility}\label{alp-lift-52-necessary-profile-changes-for-literal-link-compatibility}

Proof \nolinkurl{necessary_profile_changes_for_source_links.md}, PC1--13; author whole-PASS 2026-11-22T10:34:09.028606+00:00, frozen SHA \nolinkurl{4f4c215dfd99b0bd385ee02e2df28db6e0a5ee6afbcc8b426392b32911d5588f}. The exact non-top DO044 dimension ratio is a positive-coefficient Laurent polynomial in ell, hence convex. Its interval maximum plus the increasing top bound extends LIFT51 exclusion to ell\textgreater=ellplus/2\^{}353(J14, ratio\textless2\^{}-1218) and ell\textgreater=ellplus/2\^{}80(J56, ratio\textless2\^{}-6720), all words/roles/allowed q. Even replacing T by any split candidate above the literal Sizes bound4A cannot repair the endpoint: ratios\textless2\^{}-521598/2\^{}-1255128; uniform any-T exclusions extend to ellplus/2\^{}313 and ellplus/2\^{}46. The next-cutoff upper bounds exceed1, explicitly not an actual compatibility result.

Retaining fixed initial dimensions would require projected bottom growth by more than(3/8)2\^{}589932/(3/8)2\^{}2198952 at the endpoint; shrinking input requires dimension\textless8Drows/3 and changes Mu0. Haar, scalar and positive scale reparameterizations do not alter the rank obstruction. Source Host/Scale and effective nonvanishing at smaller gaps remain separate. Exact2139-check certificate PASS 2026-11-22T10:33:33.683897+00:00.

\subsection{ALP-LIFT-53 --- the code requirement forces a larger initial rank}\label{alp-lift-53-the-code-requirement-forces-a-larger-initial-rank}

Proof \nolinkurl{chosen_code_rank_obstruction.md}, CB1--12; author whole-PASS 2026-11-22T10:43:27.955556+00:00, frozen SHA \nolinkurl{3f4abf7b8a907b7e2514b72c0ef0ffd20f46d69ab7ed812f5dab4325e27c15f4}. Literal letter injectivity bounds Word by2\^{}(k−2), while ChCode\_at requires C0 exp(2CGam m\^{}5 fL)x\^{}CH. For m2,L343,C0\textgreater=1,CH\textgreater=0 this forces k−2\textgreater31360CGam f.~Current k=f therefore requires CGam\textless1/31360 and CH\textless140/1597 or581/1597 at every gap. The standalone SizeData2 field only says CGam\textgreater0; the every-K DesignExists requirement can legitimately request K.mid=1, forcing both CGam andCH\textgreater=1. Eventual thresholds are retained. More generally any actual ChosenCode family with fixed positive CGam has k/f\textgreater(640/7)CGam x\^{}eps along its eligible points, so bounded k/f cannot persist.

At CGam\textgreater=1 the necessary even k floors are1,317,124 and5,268,484, with N1\textgreater=2k+8,N2\textgreater=3k+16. Current J14 G primes forbid this; current J56 primes do not yet, but K.mid=13 does. This is an exact dependency and family-profile obstruction, not a universal two-edge impossibility or a repaired Host. Exact4678-check certificate PASS 2026-11-22T10:42:46.802352+00:00.

\subsection{ALP-LIFT-54 --- full source sizes exclude the actual two-edge profiles}\label{alp-lift-54-full-source-sizes-exclude-the-actual-two-edge-profiles}

Proof \nolinkurl{literal_size_chain_length_obstruction.md}, SL1--16 and SL10a--b; author whole-PASS 2026-11-22T10:52:45.058831+00:00, frozen SHA \nolinkurl{02e3d88397f8558c756412a380fdc999a4948e77a463d455f8d274bb32abd83d}. Exact \nolinkurl{SizeKRel} requires log f\textless=sqrt(m/2−1), so every full \nolinkurl{SizesAt} has even m\textgreater=2+2ceil(log²f), universally m\textgreater=20, with isolated floors30/56 for f42/168. Thus m2 is impossible at every eligible \nolinkurl{FamP} point x\textgreater x1 without waiting for code or link thresholds. \nolinkurl{SizeM} also forces m\textgreater x\^{}(10371/20800), excluding bounded-length unbounded families.

More strongly, \nolinkurl{SizeN}, \nolinkurl{SizeK} and the first \nolinkurl{SizeKRel} imply m\^{}5 f\textless k\textless(f/2)\^{}(5/2), hence f³\textgreater32m\^{}10. The exact30/56 floors contradict the upper k bounds3125/100000, so neither fixed degree admits full SizesAt at any Scale. At unchanged L343 one would need m\textgreater2\^{}231,f\textgreater2\^{}998,k\textgreater2\^{}2177. Accepted R023's retained-host lower degree gives the isolated logarithmic floor79,368; retaining o\_tie and full SizeM strengthens this to m\textgreater21\^{}(10371/4)\textgreater2\^{}11388. Actual mixed theta results remain independent, and no changed-interface DP44 comparison is decided. Exact413-check certificate PASS 2026-11-22T10:52:03.357640+00:00.

\subsection{ALP-LIFT-55 --- universal literal scale and chain floors}\label{alp-lift-55-universal-literal-scale-and-chain-floors}

Proof \nolinkurl{universal_size_scale_floor.md}, US1--10; author whole-PASS 2026-11-22T10:56:31.207358+00:00, frozen SHA \nolinkurl{371ffd8a8f3d5af04432aa28f41bcad6e57b259c05d87907cb13957331002d3f}. Combining lower SizeM/SizeN with the first SizeKRel and upper SizeK gives L\textgreater=2\^{}(400/149), retaining the factor2 in f=2n. More strongly, the simultaneous lower and logarithmic upper SizeG bounds force8(L+log k)\textless=L\^{}(101/100), so L\textgreater2\^{}300. The exact Scale epsilon\textless1/52 then implies m\textgreater2\^{}(31113/4)\textgreater2\^{}7778, f\textgreater2\^{}33598 and k\textgreater2\^{}73197. Thus rebuilding degree/ranks/primes at unchanged L343 still cannot supply full SizesAt.

These are necessary literal-interface floors, not an inconsistency at all large scales, a refutation of abstract Ranks/DP44, or an actual theta-cost sign. LIFT54's larger retained-host m\textgreater2\^{}11388 remains separate. Exact94-check certificate PASS 2026-11-22T10:56:14.878746+00:00.

\subsection{ALP-LIFT-56 --- actual longer vertical first occurrence and cuspidality}\label{alp-lift-56-actual-longer-vertical-first-occurrence-and-cuspidality}

Proof \nolinkurl{general_vertical_forward_cuspidality.md}, VF1--14; author whole-PASS 2026-11-22T11:11:00.189461+00:00, frozen SHA \nolinkurl{804c8444f8a346a35ed0ee845be972659e435d0a9dbd73be7514b3d366233def}. On the accepted R015 lists, the four real sign counts are (N\_j,q\_(j+1),0,q\_j). Every forward Witt-lower target still exceeds the source dimension but has negative index below q\_(j+1), contradicting the corrected unconditional Ichino Cor6.4. Thus all such local/global lower lifts vanish. The tower formula makes every unchanged prescribed upward finite kernel cuspidal at all stages, possibly zero. Polynomial-Fock real Rallis tests exist by Ichino Prop7.1. Every stage is in Weil convergent range, but the first positive-Witt-index edge is outside GQT11.4 and cannot use anisotropic-only Theorem7.1. Actual fixed-test nonzero propagation and positive-rank predicate3 remain separate. Exact 146091-check certificate PASS 2026-11-22T11:10:00.536935+00:00.

LIFT56 final formatting-only freeze 2026-11-22T11:11:58.322842+00:00: SHA 6a9130c36ce3d663c42711611ad8dd8c9abbda97cc1f277607b204375bc8cb86; function parentheses and inline group/rank notation corrected, with no mathematical or certificate change.

LIFT56 final source-scope freeze 2026-11-22T11:17:20.600766+00:00: SHA 4686b46a860f47275d0e163d3b28bbf96c1ef57977b464b5d094a6117df6a626. The theorem uses the independent TPar.t and its t\_ge, not Ranks.three\_m\_lt\_t; only t\textgreater=r+5\textgreater=2m+7\textgreater m is needed. Canonical ofWord matches Ranks.t. Revised certificate146135 PASS 2026-11-22T11:15:31.768279+00:00, including P.t\textless3m cases; complete author reread passed.

\subsection{ALP-LIFT-57 --- prescribed split G/Q tests through every vertical stage}\label{alp-lift-57-prescribed-split-gq-tests-through-every-vertical-stage}

Proof general\_vertical\_prescribed\_split\_tests.md, VS1--34; author whole-PASS 2026-11-22T11:34:29.542481+00:00, frozen SHA dd07aa830cb3287ad5de117c2222284e068a153f4cef67485ac4ea0eea84cb2b. For n\_j=r+jt, the exact conjugated split labels propagate the local unitary induction with all previous length-t character blocks. At G the unchanged raw map is injective on the entire actual J\_j-character type, with identity scalar vol(J\_j) and unit-input raw Rallis value exactly vol(J\_j). At Q the actual previous image has flattened support J(r,jt) and scalar C\_j. The next scalar is {[}n\_j choose r{]}\emph{q\^{}(-1) times the full product of local L(i,mu\_j/mu\_h), h\textless j,1\textless=i\textless=t. A compact tail projection precedes the absolutely convergent scalar GJ integral; no unprojected Bochner convergence is asserted after the first block. Exact compact norms and the all-stage standard-normalized Q Rallis value q\^{}(-rjt) kappa\_r kappa\_jt kappa}(j+1)t/(kappa\_t kappa\_(r+jt)) are retained. The normalized previous-image inverse factor is between q\^{}(rjt/2) and4q\^{}(rjt/2), with no global mass or arithmetic index absorbed. Other finite basics, the first positive-Witt global Rallis normalization, real pins and multiplicity propagation remain separate. Exact7775-check certificate PASS 2026-11-22T11:34:03.892281+00:00.

\subsection{ALP-LIFT-58 --- the unchanged longer vertical map and its adjoint}\label{alp-lift-58-the-unchanged-longer-vertical-map-and-its-adjoint}

Proof general\_vertical\_actual\_global\_injection.md, VI1--22; author whole-PASS 2026-11-22T11:49:40.776956+00:00, frozen SHA 88429231bf7770f8e0071339c31970020b68e13adc254925731637416429bc26. On the genuine unramified mixed hosts (and real extensions preserving the stated local data), the all-rank selfdual basics and LIFT57 fixed G/Q tests propagate every actual preceding finite image. The exact unramified Satake recurrence yields a rank-r absolutely convergent head and every pairwise length-t string of Hecke values at1,\ldots,t; literal Interlace makes each boundary ratio angular and nontrivial, so accepted LIFT14 applies without global base-change assumptions. The first positive-Witt edge is justified by Liu2011 Theorem2.1(1), probability-to-Tamagawa factor2 retained; later edges use GQT in its convergent dimension window and LIFT56 forward lower vanishing. Ichino's polynomial-Fock nonzero test plus full real Howe multiplicity one retains each actual IsItW scalar rather than choosing a highest polynomial. Every nonzero automorphic embedding survives, so the fixed upward map is injective. Cuspidal decay in both variables supplies the exact noncompact adjoint identity U=D*, hence D is onto the full fixed-level bottom cusp multiplicity space. The bottom space may be zero; actual bottom periods, positive-rank predicate3, arithmetic inverses, full sizes and DP44 costs remain open. Exact42012 diagnostic assertions PASS 2026-11-22T11:47:28.369145+00:00.

\subsection{ALP-LIFT-59 --- a genuine bottom cusp space at the unchanged literal level}\label{alp-lift-59-a-genuine-bottom-cusp-space-at-the-unchanged-literal-level}

Proof general\_good\_tuple\_bottom\_cusp.md, BC1--33; author whole-PASS 2026-11-22T12:19:04.884794+00:00, frozen SHA 0522b62257ca0350456b42625e64804b7e2a9d5a17e488653a5b35a533b5427f. The full size interface implies r\textgreater=200f(log(4r)+1) and g\textgreater=k\^{}8. An LTE count gives more than4/5 strongly good roots uniformly for all source character lists of size2m and all degrees4 throughr; Nagura supplies the actual prime-factor pattern, and two full-degree places prescribe the aggregate mu(E) central character. The normalized Flensted--Jensen coefficient has the new explicit NC-gap bound cosh(R/2)\^{}(-2delta), with the actual half-Cartan/dual-root factor and compact Blattner integrality checked. NCGaps gives delta\textgreater N³. Integral matrix packing bounds the entire noncompact rational sum by1/15; the full monomial origin stabilizer is retained, and one full-degree finite cuspidal factor bounds its nonscalars by1/4. Thus the actual unchanged-level Poincare series has \textbar P(1)-\textbar mu(E)\textbar\textbar\textless1/2, is cuspidal of the source real type, and lies in genuine e0. The generic finite-G cut is also nonzero. Combined with LIFT58 this gives a nonzero prescribed longer vertical chain. Standard actual bottom lattice and genuine host hypotheses remain explicit; no prescribed Bessel/Mram period, old Host/SizesAt construction, arithmetic inverse or DP44 sign follows. Exact60042-check diagnostic PASS 2026-11-22T12:19:05.265024+00:00.

\subsection{ALP-LIFT-60 --- actual positive-rank fractional Hida flags and finite row adapter}\label{alp-lift-60-actual-positive-rank-fractional-hida-flags-and-finite-row-adapter}

Proof positive\_rank\_hida\_finite\_adapter.md, HF1--22; author whole-PASS 2026-11-22T12:41:02.115799+00:00, frozen SHA c969d6d192bfc580e2e79ed2bc3db1c7eab7b3a0b4eaf3a8444637c944d94b56. On a finite-unramified mixed host with principal extended different, scaled hyperbolic pairs e\_a=d\^{}(rho−a)e\_a\^{}0 and f\_a=bar(d)\^{}−(rho−a)f\_a\^{}0 realize every literal aaD in the unchanged standard lattice. The primitive norm-minus-T line and all finite Hida vectors/xu conditions are explicit. Both split components give exactly PsharpSym; the negative f-chain signs and prescribed additive-character scalar are removed by an actual compact diagonal G correction. The same finite t yields full LevelMatchA/B and, with the genuine grp\_record/R026 tie, C, retaining the exact index prod\_w q\_w\^{}k prod\_(i=1)\^{}r(q\_w\^{}i−1), source field/frame and all finite Haar factors. This is not a construction of real Wu/decay, RowsJoinPinned, a Bessel/Mram period or full Host/DP44 data. Exact141764 diagnostics PASS 2026-11-22T12:37:38.178586+00:00.

\clearpage

\section{ALP-LIFT-05: stable Hodge lattices and a normalized induced CM eigenbasis}\label{alp-lift-05-stable-hodge-lattices-and-a-normalized-induced-cm-eigenbasis}

Run ALP\_MACHINE; authors /root/lift\_accounting and /root, GPT-6 Astra Pro. Reasoning effort is not exposed. The construction at 7 and the positive trace polarization were suggested independently by /root and checked here. This is an actual geometric normalization theorem. It does not identify an arithmetic theta lift with an automorphic lift, or determine a SizeData₂ gain or division-error field.

\subsection{Statement}\label{statement}

Write \(\zeta=\zeta_7\), \(E=\mathbb Q(\zeta)\), \(K=\mathbb Q(\sqrt{-7})\), and \[
 C:y^2=x^7-1,\qquad J=\operatorname{Jac}(C),\qquad
 \omega_i=x^{i-1}dx/y\quad(1\le i\le3).
 \tag{H1}
\] For a finite place above a rational prime \(p\), use the valuation normalized by \(v_p(p)=1\). Fix algebraic numbers \[
 \iota^2=-1,\quad \alpha^7=-4,\quad
 \lambda=\zeta-1,\quad \beta^2=\lambda^7,\quad \rho^4=7.
 \tag{H2}
\] After adjoining these numbers, the following assertions hold.

\textbf{(a) Actual stable Hodge lattices.} At a place above 2, the good model and its integral differential basis are \[
 Y^2+Y=X^7,\qquad
 \eta_i^{(2)}=X^{i-1}dX/(2Y+1),\qquad
 \omega_i=\alpha^i\iota^{-1}\eta_i^{(2)}.
 \tag{H3}
\] At a place above 7, they are \[
 Y^2=f_\lambda(X):=\frac{(1+\lambda X)^7-1}{\lambda^7},\qquad
 \eta_j^{(7)}=X^{j-1}dX/Y,
 \tag{H4}
\] with \[
 \omega_i=\frac{\lambda}{\beta}
       \sum_{k=0}^{i-1}\binom{i-1}{k}\lambda^k\eta_{k+1}^{(7)}.
 \tag{H5}
\] Here and below an integral differential basis means a basis of \(e^*\Omega^1_{\mathcal J/R}\), identified through the Abel map with \(H^0(\mathcal C,\Omega^1_{\mathcal C/R})\), for the good model over the local valuation ring \(R\). Consequently the signed determinant valuations of \(\omega_1\wedge\omega_2\wedge\omega_3\) are exactly \[
 \begin{array}{c|c}
 p&v_p(\omega_1\wedge\omega_2\wedge\omega_3)
       \text{ relative to an integral generator}\\ \hline
 2&12/7\\
 7&-3/4\\
 p\nmid14&0.
 \end{array}
 \tag{H6}
\] The negative entry is a valuation of a rational section, not a positive index of an inclusion of integral lattices.

\textbf{(b) Complete determinant metric and height.} Give the determinant Hodge line the normalization \[
 \|s\|_{L^2}^{\,2}
   =\frac{\mathrm i^{g^2}}{2^g}\int_{A(\mathbb C)}s\wedge\overline{s}
 \quad(g=\dim A).
 \tag{H7}
\] Let \(G=\operatorname{diag}(G_1,G_2,G_3)\) be the Gram matrix of (H1), already computed in ALP-LIFT-03: \[
 G_i=\frac{2\pi}{7}
 \frac{\Gamma(i/7)\Gamma(1/2-i/7)}
      {\Gamma(1-i/7)\Gamma(1/2+i/7)}.
 \tag{H8}
\] Then \[
 \boxed{h_{F,L^2}(J)=\frac{12}{7}\log2-\frac34\log7
                       -\frac12\log\det G.}
 \tag{H9}
\] This equals the Wilms normalization used in ALP-LIFT-02. Numerically it is \(-3.80116074543404940138079328403117913069\ldots\); the decimal is a cross-check of the exact expression.

\textbf{(c) An actual comparison variety and compatible principal polarization.} Take the elliptic curve of full CM by \(\mathcal O_K\) \[
 B:y^2+xy=x^3-x^2-2x-1,\qquad \psi=dx/(2y+x),
 \tag{H10}
\] and put \(G_B=(\mathrm i/2)\int_{B(\mathbb C)}\psi\wedge\overline\psi\), where \(\mathrm i\) is the usual positive imaginary unit in each complex metric, not the image of the algebraic symbol \(\iota\) under an arbitrary embedding. This is LMFDB 49.a4, Cremona 49a1; the full CM discriminant \(-7\) is the only catalogued fact used about it. Its Weierstrass coefficients, discriminant, local changes of variables and all following metrics are computed below.

The basis \(1,\zeta,\zeta^2\) identifies the Serre tensor variety \(A_{\mathrm{ind}}=\mathcal O_E\otimes_{\mathcal O_K}B\) with \(B^3\) as in ALP-LIFT-02. Put \[
 S=\zeta+\zeta^2+\zeta^4=(-1+\sqrt{-7})/2,\qquad
 u=-(\zeta-\zeta^{-1})^{-2}.
\] The matrix of the positive Hermitian trace form \(\operatorname{Tr}_{E/K}(u\overline x y)\) in that basis is \[
 H=\begin{pmatrix}
 2&-1&1+S\\
 -1&2&-1\\
 -S&-1&2
 \end{pmatrix},\qquad \det H=1.
 \tag{H11}
\] It defines an \(\mathcal O_E\)-compatible principal polarization on \(A_{\mathrm{ind}}\). The product polarization is not silently used as that compatible polarization.

Let \(\psi_a\), \(a=0,1,2\), denote the product invariant differentials on \(B^3\). For \(i\in\Psi=\{1,2,4\}\), set \[
 \vartheta_i=(\zeta^i-\zeta^{-i})^{-1}
                     \sum_{a=0}^{2}\zeta^{ia}\psi_a.
 \tag{H12}
\] These are CM eigenforms of type \(\Psi\). Their change-of-basis matrix \(W\), with rows in the order \(1,2,4\), has determinant \(-1\), and both \(W\) and \(W^{-1}\) have entries integral away from 7. They are therefore a primitive integral CM eigenbasis at every prime \(p\nmid7\). For the principal polarization (H11), their Hodge Gram matrix is exactly \[
                         G_B I_3.
 \tag{H13}
\] Their determinant section has normalized valuation \(-3/4\) at 7 and zero at every other prime, even though their individual lattice at 7 is not asserted to be the primitive Hodge lattice. In particular \[
 3h_{F,L^2}(B)=-\frac34\log7-\frac32\log G_B.
 \tag{H14}
\]

\textbf{(d) Exact comparison in these CM coordinates.} With \(\delta_7\) from ALP-LIFT-02 and \(\delta_{7,\{2,7\}}=\delta_7-(12/7)\log2\) from ALP-LIFT-04, \[
 \boxed{\delta_{7,\{2,7\}}
       =-\frac12\log\frac{G_1G_2G_3}{G_B^3}.}
 \tag{H15}
\] Thus the order-linear difference of the equal-direction Cauchy bounds in these actual diagonal CM coordinates is exactly the order-linear difference of the bare prime certificates over \(p\nmid14\). The precise analytic statement and its scope are given after the proof.

\subsection{Proof of (a): the good model at 2}\label{proof-of-a-the-good-model-at-2}

Make the change \(x=\alpha X\), \(y=\iota(2Y+1)\). Since \(\alpha^7=-4\), the original equation becomes exactly \(Y^2+Y=X^7\). Complete this affine chart by the chart \[
 U=X^{-1},\qquad V=YU^4,\qquad V^2+U^4V=U.
 \tag{H16}
\] On their common locus the changes are inverse to one another. Both charts are finite free of rank two over their respective affine charts of \(\mathbb P^1\), with bases \(1,Y\) and \(1,V\). They therefore form a finite flat double cover of \(\mathbb P^1\), which is proper. The algebra generators transform by \(V=U^4Y\), so the underlying rank-two pushforward is \(\mathcal O\oplus\mathcal O(-4)\); its Euler characteristic is \(1-3=-2\), giving arithmetic genus three in every fiber.

In residue characteristic 2 the finite equation has derivative 1 with respect to \(Y\). In (H16) the derivative with respect to \(U\) of \(V^2+U^4V-U\) is \(-1\). Thus the whole special fiber is smooth, including the point at infinity \((U,V)=(0,0)\). The special fiber is geometrically integral: if \(Y^2+Y-X^7\) were reducible over \(\overline k(X)\), then \(X^7=h^2+h\) for a rational function \(h\). Any finite pole of \(h\) would produce a finite pole of \(h^2+h\), and the pole at infinity of a nonconstant polynomial \(h^2+h\) has even order. Both alternatives contradict \(X^7\). The generic fiber is the original smooth curve. This proves that (H16) is a smooth proper genus-three model with an infinity section over the local valuation ring, and remains so after any finite base extension.

The residue differential on a plane curve \(F(X,Y)=0\) is \(dX/F_Y=-dY/F_X\) on the respective smooth open subsets. Thus \(\eta_i^{(2)}\) is regular on the finite chart. On the other chart it is \[
                  -U^{3-i}\frac{dU}{2V+U^4}.
 \tag{H17}
\] The same residue description shows regularity at infinity: when \(2V+U^4\) is not a unit, the alternate denominator \(4U^3V-1\) is a unit. The power \(U^{3-i}\) is nonnegative for all three indices. Reduction on the finite chart gives \(X^{i-1}dX\). These three differentials are independent over the residue field, since the Artin--Schreier function-field extension is separable and \(dX\ne0\).

For completeness, their being an integral basis follows without any unproved saturation assumption. They are a basis on the generic fiber because the genus is three. Express any global regular differential as \(\sum a_i\eta_i^{(2)}\) with coefficients in the fraction field. If some coefficient has negative valuation, multiply by the smallest positive power of a uniformizer that makes all coefficients integral and at least one a unit. This product, as a multiple of the original integral differential, reduces to zero, contradicting the independence just proved. Hence all coefficients are integral. This also proves the claim after further ramified base change.

The infinity section supplies the Abel map to the relative Jacobian. The smooth proper curve has an abelian relative Picard scheme, and the Abel map identifies its invariant differentials with those of the curve. The precise references are Milne (2021), Proposition 2.2, pp.~6--7, and Theorem 8.1 with the following paragraph, pp.~27--28; the relative assertion follows also by applying the field assertion on every fiber and Nakayama's lemma to the rank-three modules. This is the same FGA input already accepted in DP22, not a new Picard representability claim.

Finally, direct substitution gives (H3). Since \(v_2(\alpha)=2/7\), the determinant valuation is \((1+2+3)(2/7)=12/7\).

\subsection{Proof of (a): the good model at 7 and the remaining primes}\label{proof-of-a-the-good-model-at-7-and-the-remaining-primes}

The local cyclotomic uniformizer has \(v_7(\lambda)=1/6\). Write \(f_\lambda(X)=\sum_{k=1}^{7}f_kX^k\). Its top coefficient is 1, and for \(1\le k\le6\), \[
 f_k=\binom7k\lambda^{k-7},\qquad v_\lambda(f_k)=k-1.
\] Here \(v_\lambda(7)=6\) and the binomial coefficients have exactly one factor of 7. The cyclotomic identity \(\Phi_7(1+\lambda)=0\) can be rearranged as \[
 \frac7{\lambda^6}
 (1+3\lambda+5\lambda^2+5\lambda^3+3\lambda^4+\lambda^5)=-1.
\] It follows that \(f_1\equiv-1\), while \(f_2,\ldots,f_6\equiv0\). The reduction of (H4) is therefore \(Y^2=X^7-X\). Its affine chart is smooth because \(d(X^7-X)/dX=-1\) in characteristic 7. The polynomial has simple roots and is not a square, so this special fiber is geometrically integral.

The infinity chart is \[
 U=X^{-1},\quad V=YU^4,\qquad
 V^2=U+f_6U^2+f_5U^3+\cdots+f_1U^7.
 \tag{H18}
\] Away from \(U=0\) it is already covered by the finite chart. At \(U=0\) its only point is \(V=0\), and its equation has derivative \(-1\) with respect to \(U\). Thus all of this fiber is smooth. As before, the two monic charts give a proper finite flat double cover with pushforward \(\mathcal O\oplus\mathcal O(-4)\), genus three and an infinity section. Its generic fiber is the original curve by \(x=1+\lambda X\), \(y=\beta Y\).

The three differentials in (H4) are regular, since 2 is a unit and \(dX/Y\) is twice the plane residue differential. On the infinity chart they become \(-U^{3-j}dU/V\); again this is a unit multiple of the residue differential times a nonnegative power of \(U\). Their reductions are independent on the separable hyperelliptic function field. The saturation argument above proves that they form the full integral Hodge basis. Substitution of the change of variables gives (H5). Its triangular diagonal entries are \(\lambda^i/\beta\), so its determinant is \[
                      \lambda^6/\beta^3.
\] Its normalized valuation is \(6/6-3(7/12)=-3/4\), proving the second row of (H6). An explicit inverse formula, useful for the lattice rather than only its determinant, is \[
 \eta_j^{(7)}=\beta\lambda^{-j}
    \sum_{a=0}^{j-1}\binom{j-1}{a}(-1)^{j-1-a}\omega_{a+1}.
 \tag{H19}
\]

At a prime not dividing 14, the original two-chart hyperelliptic model is smooth: 2 is a unit and its branch polynomial has squarefree reduction. The same residue and independence argument gives precisely the basis (H1), proving the last row of (H6). These local models can also be made global over a common number field containing (H2). Model (H3) is smooth over \(\mathcal O_M[1/7]\); model (H4) is smooth over \(\mathcal O_M[1/2]\). On their overlap the specified changes of variables give an isomorphism, preserving infinity. Gluing yields a smooth proper curve over \(\mathcal O_M\) and hence an abelian scheme with generic fiber \(J_M\).

\subsection{Proof of (b): the determinant degree}\label{proof-of-b-the-determinant-degree}

The Hodge line is \(\det e^*\Omega^1_{\mathcal J/\mathcal O_M}\). Use its nonzero rational section \(s=\omega_1\wedge\omega_2\wedge\omega_3\). By definition of the degree of a metrized line, \[
 \frac{\widehat{\deg}(\overline\omega_{\mathcal J})}{[M:\mathbb Q]}
 =\frac1{[M:\mathbb Q]}
   \left(\sum_{v\text{ finite}}\operatorname{ord}_v(s)\log Nv
              -\sum_{\sigma:M\hookrightarrow\mathbb C}
                           \log\|s\|_\sigma\right).
 \tag{H20}
\] The local orders from (H6) give exactly \((12/7)\log2-(3/4)\log7\) after division by the degree: \(\sum_{v\mid p}e_vf_v=[M:\mathbb Q]\). The original curve and the three differentials are rational, so their complex metric is the same under every \(\sigma\). The Riemann bilinear relation identifies (H7) with \(\det G\) on this section; this is precisely the determinant convention fixed in ALP-LIFT-03. Thus the archimedean term is \(-\tfrac12\log\det G\), proving (H9). It is the stable height because the chosen model is an abelian scheme, and further base extension leaves the normalized degree unchanged. This proof does not use isogeny invariance of height.

\subsection{Proof of (c): the comparison lattice}\label{proof-of-c-the-comparison-lattice}

The Weierstrass coefficients of (H10) give \(c_4=105\), \(c_6=1323\), and discriminant \(-343\); these follow by the usual polynomial discriminant identities and can also be checked by completing the square as follows. Put \(y'=y+x/2\), \(z=x-1/4\). The equation becomes \[
                   y'^2=z^3-\frac{35}{16}z-\frac{49}{32}.
\] At 7 make the substitution \(z=\rho^2 Z\), \(y'=\rho^3 T\), where \(\rho^4=7\). The resulting equation is \[
                  T^2=Z^3-\frac5{16}Z-\frac{\rho^2}{32}.
 \tag{H21}
\] Its discriminant is exactly \(-1\), either by dividing the old discriminant by \(\rho^{12}=7^3\), or by direct substitution in \(-16(4a^3+27b^2)\). It is a smooth proper elliptic model; its special fiber is \(T^2=Z^3+Z\). The differential transforms as \[
                        \psi=\rho^{-1}\frac{dZ}{2T}.
 \tag{H22}
\] The invariant differential on the right after removing \(\rho^{-1}\) is primitive, so \(v_7(\psi)=-1/4\). At every prime different from 7 the original integral Weierstrass model has unit discriminant, and \(\psi\) is primitive. Applying (H20) in dimension one proves \(h_{F,L^2}(B)=-(1/4)\log7-(1/2)\log G_B\), and thus (H14). The value of \(G_B\) is approximately \(4.9445046002825467364981969682\); this agrees with the independent LMFDB period-area entry, but the proof of (H14) does not use that numerical entry.

\subsection{Proof of (c): the compatible polarization and eigenforms}\label{proof-of-c-the-compatible-polarization-and-eigenforms}

The integral basis \(1,\zeta,\zeta^2\) was proved in ALP-LIFT-02; it gives the specified product coordinates. Each real embedding of \(E^+\) sends \(u=-(\zeta-\zeta^{-1})^{-2}\) to \(1/(4\sin^2(2\pi i/7))>0\). Direct cyclotomic reduction gives \[
 u=\frac{6+3\zeta^2+\zeta^3+\zeta^4+3\zeta^5}{7}.
\] Taking the trace over the embeddings \(\zeta\mapsto\zeta^i\), \(i=1,2,4\), gives (H11). Its entries belong to \(\mathcal O_K\); its leading principal minors are \(2,3,1\), which proves both positive definiteness and determinant one. Multiplication by an element of \(E\) is adjoint to multiplication by its complex conjugate for this trace form. In particular, multiplication by \(\zeta\) preserves it.

Let \(\lambda_0\) be the product principal polarization on \(B^3\). The integral Hermitian matrix \(H\) is an endomorphism of that product; the homomorphism \(\lambda_H=\lambda_0\circ H\) has integral positive Riemann form, and is therefore a polarization. This standard polarization criterion is the Riemann-form criterion in Milne (2022), Section 16, p.~31. Since \(H^{-1}\) has integral entries by the adjugate formula and \(\det H=1\), \(H\) is an automorphism, so \(\lambda_H\) is principal. Its Rosati involution restricts to conjugation on \(E\) by the trace identity just proved. Thus it is the required compatible principal polarization. The polarization homomorphism is defined over a number field after the CM endomorphisms are; an ample line realizing it can be chosen after a further finite extension if necessary. No special theta trivialization is asserted at this step.

For the metric calculation, let \[
 V_{i,a}=\zeta^{ia}\quad(i=1,2,4;\ a=0,1,2),\qquad
 U=\operatorname{diag}\bigl(\sigma_i(u)\bigr).
\] Then \[
                         H=V^*UV.
 \tag{H23}
\] In product differential coordinates, the Hodge Gram for \(\lambda_H\) is \(G_B H^{-1}\). One may check this directly: the normalized product Kähler form has coefficient matrix \(G_B^{-1}H^t\); integrating \((\mathrm i/2)\psi_a\wedge\overline\psi_b\) against its second power divided by \(2!\) gives the corresponding cofactor, hence \(G_B(\det H)H^{-1}\). Here \(\det H=1\). This computation fixes the transpose convention and does not identify covectors with tangent vectors without taking the inverse metric.

The rows of \(V\) are eigen-covectors for multiplication by \(\zeta\), with eigenvalues \(\zeta,\zeta^2,\zeta^4\). Indeed evaluation at \(\zeta^i\) sends multiplication by \(\zeta\) on \(E\) to multiplication by \(\zeta^i\). It follows from (H23) that \[
 V(G_BH^{-1})V^*=G_B U^{-1}
     =G_B\operatorname{diag}(|\zeta^i-\zeta^{-i}|^2).
 \tag{H24}
\] Dividing the corresponding eigenforms by \(\zeta^i-\zeta^{-i}\) therefore proves (H13).

The Vandermonde determinant and the cyclotomic product are \[
 \det V=-\sqrt{-7},\qquad
 \prod_{i=1,2,4}(\zeta^i-\zeta^{-i})=\sqrt{-7},
\] with the same square-root convention \(2S+1=\sqrt{-7}\). Their quotient gives \(\det W=-1\). Each denominator is a unit away from 7, so all entries of \(W\) are integral there; its adjugate gives the same conclusion for its inverse. This proves primitivity at those primes. At 7 the determinant section equals \(-\psi_0\wedge\psi_1\wedge\psi_2\), exactly, so its valuation is still \(-3/4\) by (H22), regardless of the nonintegral entries of \(W\). There has been no assertion that a general change of coordinates transforms a product of derivatives by a determinant. The CM eigenbasis (H12) is chosen first, and the derivatives are then its dual coordinate directions.

Finally subtract (H14) from (H9). ALP-LIFT-02 identifies the left side with \(\delta_7\), and the finite 7-adic contributions cancel exactly. This gives \[
 \delta_7=\frac{12}{7}\log2-
                    \frac12\log\frac{\det G}{G_B^3},
\] which is (H15) and completes the proof. \(\square\)

\subsection{Exact differential rescaling and the 2-adic limitation}\label{exact-differential-rescaling-and-the-2-adic-limitation}

Let \(D_i\) be the invariant derivations dual to \(\omega_i\), and let \(E_i\) be dual to \(\eta_i^{(2)}\). Formula (H3) gives the exact operator identities \[
 D_i=\iota\alpha^{-i}E_i,\qquad
 D_1^{n_1}D_2^{n_2}D_3^{n_3}
  =\iota^{n_1+n_2+n_3}\alpha^{-(n_1+2n_2+3n_3)}
                         E_1^{n_1}E_2^{n_2}E_3^{n_3}.
 \tag{H25}
\] The associated valuation shift is \(-\sum_i(2i/7)n_i\). The same identity holds for covariant derivatives when the connection and local frame are transported along the same isomorphism: connections are linear in a vector field, the scaling factors are constant, and the individual operators are composed in the displayed order. This is a statement about the same section and connection, not a comparison of two independently normalized theta frames.

For the abstract DP64 slope certificate at 2, the three directions of \(\Phi=\{1,2,3\}\) have slopes respectively \(2/7,4/7,6/7\); the induced type \(\Psi\) has slope zero in each direction. These numerical slopes equal the valuations of the coefficients in (H3). The equality gives a concrete geometric interpretation of the omitted \((12/7)\log2\) in ALP-LIFT-04. It does not extend the accepted good-model theta theorem to 2. In particular the trace-one endomorphism \(w_0=-(\zeta+\zeta^2+\zeta^4)\) has \(w_0+w_0^\dagger=1\) but \(w_0w_0^\dagger=2\), so one of its 2-adic conjugates is a nonunit. The odd-prime unit-path hypotheses used in DP49 do not follow there. The primitive theta frame, gauge and selected gain still need their own comparison.

\subsection{Precise equal-direction Cauchy consequence}\label{precise-equal-direction-cauchy-consequence}

This paragraph states exactly what part of a theta estimate is fixed by the calculation. Let \(F_J\) and \(F_A\) be holomorphic functions in the universal-cover coordinates dual to \(\omega_i\) and \(\vartheta_i\), respectively, satisfying \[
 |F_J(z)|\le C_J\exp\!\left(\frac\pi2
                       \sum_{i=1}^3|z_i|^2/G_i\right),\qquad
 |F_A(t)|\le C_A\exp\!\left(\frac\pi{2G_B}
                       \sum_{i=1}^3|t_i|^2\right)
 \tag{H26}
\] for positive constants independent of the derivative order. These are the Chern-frame growth inequalities for bounded-norm sections of the respective principal theta lines; /root's AM-R004 constructs the specific analytic germ on \(J\) and supplies its constant separately. Define the logarithmic Cauchy upper bounds, with equal order \(n\ge1\) in the three chosen CM directions, by \[
 \mathcal C_J(n)=\log C_J+3\log(n!)+
               \frac{3n}{2}\log(\pi e/n)-\frac n2\log\det G,
\] and the same expression with \(C_A\) and \(G_B^3\) for \(\mathcal C_A(n)\). Cauchy's formula on a polydisk of radii \(R_i^2=nG_i/\pi\) proves the bound for \(J\); replacing the three \(G_i\) by \(G_B\) proves the other. Subtraction and (H15) give the exact difference of these \textbf{upper-bound expressions}: \[
 \mathcal C_J(n)-\mathcal C_A(n)
     =\log(C_J/C_A)+n\delta_{7,\{2,7\}}.
 \tag{H27}
\] Therefore ALP-LIFT-04 implies, for its bare integer certificate \(D_{\{2,7\}}(n;X)\) and every joint limit \(n\to\infty\), \(X\to\infty\), \[
 \frac{D_{\{2,7\}}(n;X)
              -\bigl(\mathcal C_J(n)-\mathcal C_A(n)\bigr)}n
                         \longrightarrow0.
 \tag{H28}
\] The same identity holds after multiplying both displayed budgets by the source's common factor 2. A tensor level \(\kappa\) in the exponent replaces \(\pi\) by \(\kappa\pi\); the added \((3n/2)\log\kappa\) cancels between the two bounds.

Neither (H27) nor (H28) is an equality between actual derivative values. The constants can change under arithmetic normalization of a theta section, and no primitive integral frame or selected automorphic row is identified by choosing an analytic sup norm. Thus this closes the order-linear \textbf{bare certificate versus specified diagonal Cauchy metric} comparison in actual CM coordinates, while leaving the source gain, division error, bad-prime theta normalization, row content, and kernel matching as separate inputs to DP44. It gives no shorter valid lift chain and no universal sign assertion for its full cost.

\subsection{Grothendieck and source scope}\label{grothendieck-and-source-scope}

The smooth models above give an actual relative Picard scheme and its Hodge bundle, not merely a characteristic-zero CM type. The Picard representability and smoothness locators are those already accepted in DP22: FGA V, Bourbaki 232, Theorem 3.1, pp.~148--149; FGA VI, Bourbaki 236, Corollary 1.7, p.~229, Corollary 2.3, p.~232, and Proposition 2.10(ii), p.~236. Those source results are used with their accepted status, not re-proved or claimed newly scan-checked. The Hodge line is the determinant of the degree-one de Rham realization of this Picard scheme. Its metric and finite lattices are evaluated explicitly in (H6)--(H15).

The principal polarization is the homomorphism to the dual Picard scheme determined by its Poincaré biextension. The trace identity in (H11) makes the \(E\)-action and its conjugate adjoint, so it is compatible with that biextension. This is the concrete content relevant to SGA 7 I, Exposés VII--VIII and the previously accepted DP49 trace-one construction; no further comparison theorem is ascribed to those references. No new claim about a lambda-ring, an unrelated K-theory class, or weights of an automorphic sheaf is needed for this Hodge determinant calculation.

\subsection{References and exact retrievals}\label{references-and-exact-retrievals}

\begin{itemize}
\tightlist
\item
  Milne, J. S. (2021, June 12). \emph{Jacobian varieties}. Proposition 2.2, printed pp.~6--7; Theorem 8.1 and following paragraph, printed pp.~27--28. \url{https://www.jmilne.org/math/xnotes/JVs.pdf} Primary text opened this run; current retrievals turn963921859view0, turn794617041view1, turn794617041view2.
\item
  Milne, J. S. (2022, January 2). \emph{Abelian varieties}. Section 16, printed p.~31, nondegenerate positive Riemann-form criterion; Section 13, printed p.~25, principal-polarization convention. \url{https://www.jmilne.org/math/xnotes/AVs.pdf} Current retrievals turn794617041view0 and turn892289938view0.
\item
  The LMFDB Collaboration. (2026). \emph{Elliptic curve 49.a4: Data}. The ec\_curvedata row identifies coefficients {[}1,−1,0,−2,−1{]}, CM discriminant −7 and Cremona label 49a1; the ec\_mwbsd row contains the independent period-area decimal. Retrieved November 21, 2026, from \url{https://www.lmfdb.org/EllipticCurve/Q/data/49.a4} Current retrieval turn963921859view1.
\item
  ALP-LIFT-02, ALP-LIFT-03 and ALP-LIFT-04 in \nolinkurl{work/agents/lift_accounting.md}, with their exact cited Colmez--Obus, Wilms, DLMF and fixed-modulus prime-theorem inputs.
\item
  DP\_MULT DP22, DP49 and DP64, with intake status unchanged. The actual theta gain is not identified with the abstract DP64 floor certificate, at 2 or at any other prime.
\end{itemize}

\clearpage

\section{ALP-LIFT-06: directional prime constants and unequal-order Cauchy cancellation}\label{alp-lift-06-directional-prime-constants-and-unequal-order-cauchy-cancellation}

ALP\_MACHINE; /root/lift\_accounting, GPT-6 Astra Pro, with the directional question proposed by /root. Reasoning effort not exposed. This refines ALP-LIFT-04 using the actual CM eigenbases and polarization constructed in ALP-LIFT-05. It concerns the bare DP64 certificate and a specified Chern metric. It does not identify actual theta-selected gain or source division error.

\subsection{Statement}\label{statement-1}

Use \(\Phi=\{1,2,3\}\), \(\Psi=\{1,2,4\}\) in \((\mathbb Z/7\mathbb Z)^\times\). For \(i=1,2,3\), define \[
 b_i(a)=\mathbf1_{ai\notin\Phi}-\mathbf1_{a\notin\Psi},
 \qquad
 d_i(p)=\sum_{r\ge1}\frac{b_i(p^r)}{p^r}\quad(p\ne7).
 \tag{D1}
\] The second indicator is the same for every selected direction of \(\Psi\), since \(\Psi\) is a subgroup. Thus it is the appropriate comparison for any pairing of the three chosen \(\Phi\)-directions with the three \(\Psi\)-directions. Let \(\ell_a=\log\Gamma(a/7)\) and \[
 Q=\ell_1+\ell_2-\ell_3+\ell_4-\ell_5-\ell_6.
\]

\textbf{(a) Directional constants in ordinary prime order.} The limits \[
 C_i=\lim_{X\longrightarrow\infty}
                  \sum_{\substack{p\le X\\p\ne7}}d_i(p)\log p
 \tag{D2}
\] exist, with the primes in their ordinary increasing order, and satisfy \[
 \boxed{C_i=\frac12\bigl(Q-2\ell_i+2\ell_{7-i}
                              -\ell_{7-2i}+\ell_{2i}\bigr).}
 \tag{D3}
\] In particular \[
\begin{split}
 C_1&=-\tfrac12\ell_1+\ell_2-\tfrac12\ell_3
                         +\tfrac12\ell_4-\ell_5+\tfrac12\ell_6,\\
 C_2&=\tfrac12\ell_1-\tfrac12\ell_2-\ell_3
                         +\ell_4+\tfrac12\ell_5-\tfrac12\ell_6,\\
 C_3&=\tfrac12\ell_2-\tfrac32\ell_3
                         +\tfrac32\ell_4-\tfrac12\ell_5
      =C_1+C_2=\tfrac12\delta_7.
\end{split}
 \tag{D4}
\] Let \(G_i\), \(G_B\) be the exact metrics in ALP-LIFT-05. Then \[
 \boxed{C_i-\frac{2i}{7}\log2
               =-\frac12\log(G_i/G_B)\qquad(i=1,2,3).}
 \tag{D5}
\] The subtraction removes exactly the formal prime-2 contribution; it leaves the actual common good-prime range \(p\nmid14\).

\textbf{(b) Integer certificates with an arbitrary growing cutoff.} For an integer \(n\ge0\), real \(X\ge2\), and fixed finite set of rational primes \(S\) containing 7, put \[
 D_{i,S}(n;X)=\sum_{\substack{p\le X\\p\notin S}}\log p
                     \sum_{r\ge1}b_i(p^r)\left\lfloor n/p^r\right\rfloor,
 \qquad
 C_{i,S}=C_i-\sum_{p\in S\setminus\{7\}}d_i(p)\log p.
 \tag{D6}
\] There is \(c>0\) such that, uniformly in \(n\ge4\), \(X\ge2\), \[
 D_{i,S}(n;X)=nC_{i,S}
    +O_S\!\left(n\exp\{-c\sqrt{\log\min(X,\sqrt n)}\}\right).
 \tag{D7}
\] The constant can be chosen for all three directions at once. Consequently, for any vector of nonnegative integer orders \(\boldsymbol n=(n_1,n_2,n_3)\), any common cutoff \(X\), and \(N=n_1+n_2+n_3\), \[
 \sum_{i=1}^3D_{i,\{2,7\}}(n_i;X)
       =-\frac12\sum_{i=1}^3n_i\log(G_i/G_B)+o(N)
 \tag{D8}
\] in every joint limit \(N\to\infty\), \(X\to\infty\). No fixed positive limiting ratio between the three orders is required.

\textbf{(c) Unequal-order Cauchy comparison.} In the actual diagonal CM coordinates of ALP-LIFT-05, let the two functions satisfy its growth inequalities (H26), with fixed constants \(C_J^\theta,C_A^\theta>0\). Pair \((\omega_1,\omega_2,\omega_3)\) with \((\vartheta_1,\vartheta_2,\vartheta_4)\). Define the logarithmic polydisk Cauchy bounds \[
 \mathcal C_J(\boldsymbol n)=\log C_J^\theta+
   \sum_{i=1}^3\left[
     \log(n_i!)+\frac{n_i}{2}\log(\pi e/n_i)
                          -\frac{n_i}{2}\log G_i\right],
 \tag{D9}
\] where a term with \(n_i=0\) is interpreted by continuity as zero, and define \(\mathcal C_A\) by replacing \(G_i\) by \(G_B\) and the leading constant by \(C_A^\theta\). Then \[
 \mathcal C_J(\boldsymbol n)-\mathcal C_A(\boldsymbol n)
 =\log(C_J^\theta/C_A^\theta)
               +\sum_{i=1}^3n_i C_{i,\{2,7\}}.
 \tag{D10}
\] It follows from (D8) that \[
 \sum_iD_{i,\{2,7\}}(n_i;X)
       -\bigl(\mathcal C_J(\boldsymbol n)-
                                 \mathcal C_A(\boldsymbol n)\bigr)
                                                        =o(N)
 \tag{D11}
\] for the same arbitrary joint limit. Multiplying both budgets by the source's factor 2 leaves this cancellation unchanged. The assertion is about the displayed bounds, not equality of actual theta values.

\subsection{Proof of the local functions and convergence}\label{proof-of-the-local-functions-and-convergence}

Direct evaluation of the six residue classes gives \[
 b_1=\mathbf1_{\{4\}}-\mathbf1_{\{3\}},\qquad
 b_2=\mathbf1_{\{2\}}-\mathbf1_{\{5\}},\qquad
 b_3=b_1+b_2.
 \tag{D12}
\] They have mean zero and are odd: \(b_i(-a)=-b_i(a)\). Geometric summation of the powers of a residue class gives the following exact local table:

\begin{longtable}[]{@{}llll@{}}
\toprule\noalign{}
\(p\bmod7\) & \(d_1(p)\) & \(d_2(p)\) & \(d_3(p)\) \\
\midrule\noalign{}
\endhead
\bottomrule\noalign{}
\endlastfoot
\(1,6\) & \(0\) & \(0\) & \(0\) \\
\(2\) & \(p/(p^3-1)\) & \(p^2/(p^3-1)\) & \(p(p+1)/(p^3-1)\) \\
\(4\) & \(p^2/(p^3-1)\) & \(p/(p^3-1)\) & \(p(p+1)/(p^3-1)\) \\
\(3\) & \(-p^2/(p^3+1)\) & \(p/(p^3+1)\) & \(-p(p-1)/(p^3+1)\) \\
\(5\) & \(p/(p^3+1)\) & \(-p^2/(p^3+1)\) & \(-p(p-1)/(p^3+1)\) \\
\end{longtable}

For example, in class 3 the nonzero terms of \(b_1(p^r)\) are \(-1\) at \(r\equiv1\pmod6\) and \(+1\) at \(r\equiv4\pmod6\); hence \((-p^{-1}+p^{-4})/(1-p^{-6})=-p^2/(p^3+1)\). In class 2 the period is three, and the nonzero \(b_1\)-term is at exponent 2, giving \(p^{-2}/(1-p^{-3})=p/(p^3-1)\). Interchanging classes 2 and 4, or 3 and 5, interchanges \(b_1\) and \(b_2\); (D12) gives the last column. These observations prove every entry. In particular \[
                  (d_1(2),d_2(2),d_3(2))=(2/7,4/7,6/7).
 \tag{D13}
\]

For clarity about convergence, the quantitative prime theorem for fixed modulus 7, already cited in ALP-LIFT-04, gives \[
 B_i(x):=\sum_{p\le x}b_i(p)\log p
                   =O(xe^{-c_0\sqrt{\log x}}).
 \tag{D14}
\] The terms at 7 are defined to be zero. The main terms cancel because \(\sum_{a=1}^6b_i(a)=0\). The prime-power version of the same theorem implies this prime version as well, since powers of exponent at least two contribute \(O(\sqrt x\log x)\). Partial summation now proves convergence of \(\sum_p b_i(p)\log p/p\), with tail \(O((1+\sqrt{\log X})e^{-c_0\sqrt{\log X}})\). The series of the remaining terms is absolutely convergent, since \[
 \sum_p\log p\sum_{r\ge2}|b_i(p^r)|p^{-r}
       \le\sum_{m\ge2}\frac{\log m}{m(m-1)}<\infty.
\] This proves (D2) in the asserted ordinary-prime order. It also justifies absorbing the harmless polynomial factor in the tail by decreasing \(c_0\).

\subsection{Proof of the exact constants}\label{proof-of-the-exact-constants}

Let \(\chi(3)=\exp(\pi\mathrm i/3)\), and write \(Z_\eta=L'(0,\eta)/L(0,\eta)\). In the ordered odd character basis \((\chi,\chi^3,\overline\chi)\), the coefficients of the three functions in (D12) are \[
 \begin{array}{c|ccc}
 &\chi&\chi^3&\overline\chi\\\hline
 b_1&(-1+\mathrm i\sqrt3)/6&1/3&(-1-\mathrm i\sqrt3)/6\\
 b_2&(-1-\mathrm i\sqrt3)/6&1/3&(-1+\mathrm i\sqrt3)/6\\
 b_3&-1/3&2/3&-1/3.
 \end{array}
 \tag{D15}
\] This follows either by substituting the six group elements into the table, or by the normalized Fourier coefficient \(6^{-1}\sum_ab_i(a)\overline\eta(a)\). The sum of each row of coefficients is zero, since \(b_i(1)=0\).

For real \(s>1\), the Euler product gives \[
 \sum_{p\ne7}\log p\sum_{r\ge1}b_i(p^r)p^{-rs}
        =-\sum_\eta\widehat b_i(\eta)
                                      \frac{L'(s,\eta)}{L(s,\eta)}.
 \tag{D16}
\] Partial summation using (D14) permits the limit \(s\downarrow1\) for the first powers by an integrable bound independent of \(s\in[1,2]\); the higher powers converge absolutely there. The functional equation for an odd primitive character of conductor 7 gives \[
 \frac{L'(1,\eta)}{L(1,\eta)}+Z_{\overline\eta}
                         =\gamma+\log(2\pi/7).
 \tag{D17}
\] The constant vanishes after summation because the coefficient row sums in (D15) are zero. Hence \[
 C_i=\sum_\eta\widehat b_i(\eta)Z_{\overline\eta}.
 \tag{D18}
\] The conjugation in this formula is essential. Write \(Z=Z_\chi\), \(Z_3=Z_{\chi^3}\). Then (D15)--(D18) give \[
 C_1=(Z_3-\operatorname{Re}Z+\sqrt3\operatorname{Im}Z)/3,
 \quad
 C_2=(Z_3-\operatorname{Re}Z-\sqrt3\operatorname{Im}Z)/3,
 \quad C_3=2(Z_3-\operatorname{Re}Z)/3.
 \tag{D19}
\]

The Hurwitz special values used in ALP-LIFT-02 yield \[
 Z_3=-\log7+Q,
\] \[
 \operatorname{Re}Z=-\log7+\ell_1+\tfrac14\ell_2
          +\tfrac54\ell_3-\tfrac54\ell_4-\tfrac14\ell_5-\ell_6,
\] and the same finite character calculation yields \[
 \sqrt3\operatorname{Im}Z=-\tfrac32\ell_1+\tfrac94\ell_2
             +\tfrac34\ell_3-\tfrac34\ell_4-\tfrac94\ell_5
             +\tfrac32\ell_6.
 \tag{D20}
\] For example, the last two formulas follow by multiplying \(\sum_a\chi(a)\ell_a\) by \(L(0,\chi)^{-1}=1-\mathrm i\sqrt3/2\), whose value was computed exactly in ALP-LIFT-02. Substituting (D20) into (D19) gives all three expansions in (D4), and these are exactly (D3).

To identify the actual metric, the elliptic case of the same normalized Colmez formula, together with (H22), gives \[
                         \log G_B=\log(\pi/7)+Q.
 \tag{D21}
\] Here is the normalization explicitly. ALP-LIFT-02's Yang convention gives \(h_F(B)=-\tfrac12Z_3-\tfrac14\log7\); passing to its L² convention subtracts \(\tfrac12\log\pi\). On the other hand (H22) gives \(h_{F,L^2}(B)=-\tfrac14\log7-
\tfrac12\log G_B\). Equating these expressions proves (D21). No isogeny invariance is used.

The gamma duplication formula, applied at \(i/7\) and \(1/2-i/7\), rewrites (H8) as \[
 \log G_i=\log(\pi/7)+\frac{4i}{7}\log2
          +2\ell_i-2\ell_{7-i}+\ell_{7-2i}-\ell_{2i}.
 \tag{D22}
\] All gamma arguments are positive for \(i=1,2,3\). Subtracting (D21) from (D22), then multiplying by \(-1/2\) and using (D3), proves (D5). The explicit prime-2 values (D13) show that its left side is exactly \(C_{i,\{2,7\}}\). This proves part (a).

\subsection{Proof of the uniform integer estimate}\label{proof-of-the-uniform-integer-estimate}

We give the estimates rather than identifying a formal slope with a theta value. First let \(S=\{7\}\), suppress the index \(i\), and put \(b=b_i\). Decompose (D6) into first powers and higher powers. Define \[
 P_1(X)=\sum_{p\le X}b(p)\log p/p,\qquad
 P_2(X)=\sum_{p\le X}\log p\sum_{r\ge2}b(p^r)/p^r.
\] The convergence proof gives \[
 P_1(X)+P_2(X)=C_i+O(e^{-c_1\sqrt{\log X}}),\quad
 P_1(X)-P_1(\infty)=O(e^{-c_1\sqrt{\log X}}),
 \tag{D23}
\] and \(P_2(X)-P_2(\infty)=O((1+\log X)/X)\).

The higher-power floor errors are uniformly \(O(\sqrt n\log n)\). Indeed \(H_2(t)=\sum_{p^r\le t,r\ge2}\log p=O(\sqrt t\log t)\) by \(\theta(t)=O(t)\). Floors for \(p^r\le n\) produce an error at most \(H_2(n)\); the tail of the linear approximation is at most \(n\int_{(n,\infty)}t^{-1}dH_2(t)
=O(\sqrt n\log n)\). Restricting to \(p\le X\) only reduces these absolute bounds. Thus the higher-power contribution is \(nP_2(X)+O(\sqrt n\log n)\).

Put \(y=\lfloor\sqrt n\rfloor\). If \(X\le y\), the first-power floor error is at most \(\theta(X)=O(X)\). Together with (D23) this proves (D7) in that range. If \(X\ge y\), expanding the floors gives the exact identity \[
 \sum_{p\le X}b(p)\log p\left\lfloor n/p\right\rfloor
 =\sum_{p\le y}b(p)\log p\left\lfloor n/p\right\rfloor
 +\sum_{k\le n/y}
       \{B_i(\min(X,n/k))-B_i(y)\}.
 \tag{D24}
\] In the second sum, \(y\le\min(X,n/k)\le n/k\). Formula (D14) and the harmonic-sum bound show that its absolute value is at most \(O(n(1+\log n)e^{-c_0\sqrt{\log y}})\). The first sum is \(nP_1(\infty)+O(ne^{-c_1\sqrt{\log y}}+y)\). Combining with the higher powers and \(P_2(X)-P_2(\infty)=O((1+\log y)/y)\) proves (D7) for \(X\ge y\), after reducing \(c>0\). In both ranges the \(O(\sqrt n\log n)\) error is absorbed by the displayed exponential bound. The replacement of \(y\) by \(\sqrt n\) only changes a fixed multiplicative constant for \(n\ge4\).

For a fixed omitted prime \(p\ne7\), the corresponding floor sum differs from \(nd_i(p)\) by \(O_p(1+\log n)\): only \(O_p(1+\log n)\) powers have a nonzero floor, and the remaining linear geometric tail is \(O_p(1)\). Its \(\log p\) factor is fixed. Once \(X\) contains the exceptional set, subtracting these finitely many terms proves (D7) for \(S\). When it does not, \(X\) stays bounded by the largest exceptional prime, and an extra \(O_S(n)\) is absorbed by the right side after enlarging its constant. This establishes the full uniform assertion.

To deduce (D8) without a restriction on relative orders, fix \(\varepsilon>0\). Choose \(K\) large enough that the error factor in (D7) is below \(\varepsilon\) whenever \(n_i\ge K\) and \(X\ge\sqrt K\). The sum of errors in these directions is at most a constant times \(\varepsilon N\). In directions with \(n_i<K\), both their floor sums and their linear approximations are bounded by a constant depending only on \(K\), since every nonzero floor has \(p\le n_i<K\), independent of \(X\). Divide by \(N\), let \(N,X\to\infty\), then let \(\varepsilon\to0\). Using (D5) now proves (D8).

\subsection{Proof of the Cauchy statement and its scope}\label{proof-of-the-cauchy-statement-and-its-scope}

Apply polydisk Cauchy to the bounds in (H26). For a positive directional order choose \(R_i^2=n_iG_i/\pi\); for order zero take \(R_i\downarrow0\). The resulting bound is (D9). Perform the same calculation in the normalized induced eigenbasis (H12), whose three metric values are all \(G_B\). The factorial terms, \(\pi e\) terms and the individual \(n_i\log n_i\) terms cancel exactly between the two expressions. The remaining metric terms give (D10) by (D5); (D8) proves (D11). A common tensor level \(\kappa\) changes \(\pi\) to \(\kappa\pi\) and hence adds \((N/2)\log\kappa\) to both expressions, so it cancels as well. This is the level dependence whose source differential normalization is treated in AM-R004. This completes the proof. \(\square\)

The input in this theorem is an actual diagonal CM eigenbasis for both principally polarized abelian varieties, primitive at every common good prime. Thus the metric term is not a hypothetical period parameter. The lower prime certificate remains precisely the accepted DP64 object. The theorem proves that assigning different orders to the three fixed CM directions does not create an order-linear advantage within this paired bare-certificate and Cauchy-bound calculation.

At 7 the normalized eigenbasis need not be integral one direction at a time, although its determinant has the exact valuation proved in ALP-LIFT-05. This theorem does not bound an unequal-order 7-adic theta error. Nor does it show that the Cauchy bounds are attained, that selected values equal their universal lower certificates, or that the same arithmetic normalization gives the constants \(C_J^\theta,C_A^\theta\). The arithmetic kernel, row content, bad-prime theta normalization and division losses are still required before applying the full DP44 comparison.

\subsection{Numerical cross-checks and citations}\label{numerical-cross-checks-and-citations}

The accompanying exact-algebra/numerical script is \nolinkurl{work/agents/hodge_lattice_certificate.py}; it checks the gamma, Fourier and metric identities to 90 decimal digits as a cross-check of the symbolic proof. Its current approximate constants are:

\begin{longtable}[]{@{}
  >{\raggedright\arraybackslash}p{(\columnwidth - 4\tabcolsep) * \real{0.2727}}
  >{\raggedleft\arraybackslash}p{(\columnwidth - 4\tabcolsep) * \real{0.3636}}
  >{\raggedleft\arraybackslash}p{(\columnwidth - 4\tabcolsep) * \real{0.3636}}@{}}
\toprule\noalign{}
\begin{minipage}[b]{\linewidth}\raggedright
direction
\end{minipage} & \begin{minipage}[b]{\linewidth}\raggedleft
all \(p\ne7\): \(C_i\)
\end{minipage} & \begin{minipage}[b]{\linewidth}\raggedleft
good \(p\nmid14\): \(C_{i,\{2,7\}}\)
\end{minipage} \\
\midrule\noalign{}
\endhead
\bottomrule\noalign{}
\endlastfoot
1 & −0.127202946857769870150102196626 & −0.325244998446325672840739945614 \\
2 & 0.155046463619921284560680397991 & −0.241037639557190320820595099985 \\
3 & 0.0278435167621514144105782013653 & −0.566282638003515993661335045599 \\
\end{longtable}

External known inputs and exact locators are unchanged from ALP-LIFT-02/04/05: the abelian Colmez theorem in Yang (2010), Section 3, pp.~464--466, with Colmez (1993) and Obus (2013), arXiv v2 Theorem 3.9; Rudnick (2015), fixed-modulus prime theorem, Section 2, printed p.~2 before equation (1); DLMF 25.15.2, 25.15.5 and 25.15.9 for the Euler product, primitive functional equation and nonvanishing at 1, 25.11.13 and 25.11.18 for Hurwitz special values, and 5.5.5 for gamma duplication. The smooth Picard/Hodge and compatible-biextension interpretation is precisely that in ALP-LIFT-05, with the accepted DP22 FGA locators; it requires no additional weight or comparison-isomorphism claim.

\clearpage

\section{AM-R-004 --- an actual CM period lattice, theta section and uniform archimedean jet bound}\label{am-r-004-an-actual-cm-period-lattice-theta-section-and-uniform-archimedean-jet-bound}

Author: GPT-6 Astra Pro, ALP\_MACHINE \nolinkurl{/root}, 21 November 2026 UTC; reasoning effort is not separately exposed. This is a geometric normalization theorem on the actual Jacobian. It is not a construction of the requested global automorphic lift chain. The period-matrix theorem itself is known: the exact source is Tadokoro (2014), Proposition 3.2, Proposition 4.1 and Theorem 4.4, arXiv:1211.6910v5, pp.~9--10 and 13--15. We cite that theorem and calculate the particular basis, polarization, sections and bounds needed here.

\subsection{Statement}\label{statement-2}

Let C be the smooth projective curve y²=x⁷−1, J its principally polarized Jacobian, ζ=exp(2πi/7), E=Q(ζ), and Φ=\{1,2,3\}. Put

\[
 \omega_i=x^{i-1}dx/y,\qquad
 c_i=-\frac{2i_{\mathbb C}}7 B(i/7,1/2)(\zeta^i-1),\quad
 D=\operatorname{diag}(c_1,c_2,c_3).
 \tag{R4.1}
\]

Here i\_C is the imaginary unit, distinct from the index i. Define the 3 by 3 matrices, with i=1,2,3 and j=0,1,2,

\[
 A_{ij}=\zeta^{2ij},\qquad
 B_{ij}=\sum_{k=j}^{2}\zeta^{i(2k+1)},\qquad P=DA,
 \qquad \tau=A^{-1}B.
 \tag{R4.2}
\]

The following normalizations are explicit.

\begin{enumerate}
\def\labelenumi{\arabic{enumi}.}
\tightlist
\item
  The actual period lattice in the integral coordinates u\_i=∫ω\_i is
\end{enumerate}

\[
 \Lambda=D\Phi(\mathcal O_E)=P(\mathbb Z^3+\tau\mathbb Z^3).
 \tag{R4.3}
\]

Its principal Riemann form in O\_E coordinates is

\[
 \mathcal E(x,y)=\operatorname{Tr}_{E/\mathbb Q}
       \left(\frac{\zeta-\zeta^{-1}}7x\bar y\right).
 \tag{R4.4}
\]

In particular this is a concrete mixed-type CM torus with its actual principal polarization, not a torus induced from the quadratic type \{1,2,4\}.

\begin{enumerate}
\def\labelenumi{\arabic{enumi}.}
\setcounter{enumi}{1}
\tightlist
\item
  The normalized matrix is the following symmetric element of the Siegel upper half-space:
\end{enumerate}

\[
\tau=\begin{pmatrix}
-\zeta^5-1&-\zeta^5-\zeta^4-\zeta^2-2&\zeta^5+\zeta^3+\zeta\\
-\zeta^5-\zeta^4-\zeta^2-2&\zeta^5-\zeta^4+2\zeta^3+\zeta&\zeta^5+\zeta^3+\zeta^2+1\\
\zeta^5+\zeta^3+\zeta&\zeta^5+\zeta^3+\zeta^2+1&\zeta^2
\end{pmatrix}.
\tag{R4.5}
\]

Write Y=Im τ, and let G=diag(G\_1,G\_2,G\_3) be the Hodge Gram matrix computed in ALP-LIFT-03. Then

\[
 G=P Y\bar P^{\mathsf T},\qquad
 \operatorname{covol}_{\mathbb R^6}(\Lambda)=\det G
 =\frac{7^{5/2}}8\prod_i|c_i|^2.
 \tag{R4.6}
\]

The Euclidean volume here refers to the displayed u coordinates. It is not a level-group volume or a \nolinkurl{MainX} covolume.

\begin{enumerate}
\def\labelenumi{\arabic{enumi}.}
\setcounter{enumi}{2}
\tightlist
\item
  For every integer κ≥1, define
\end{enumerate}

\[
 \theta_\kappa(z)=\sum_{n\in\mathbb Z^3}
     \exp(\pi i\kappa n^{\mathsf T}\tau n+2\pi i\kappa n^{\mathsf T}z),
 \quad
 F_\kappa(u)=
 \frac{\exp\!\left(\frac{\pi\kappa}{2}z^{\mathsf T}Y^{-1}z\right)
       \theta_\kappa(z)}{\theta_\kappa(0)},\quad z=P^{-1}u.
 \tag{R4.7}
\]

These are actual sections of L\^{}κ, for the symmetric principal theta line L attached to τ, written in its local Chern-normalized holomorphic frame. Their rigidified germs satisfy F\_κ(0)=1. Explicit rational bounds prove

\[
 \lambda_{\min}(Y)>6/25,\qquad
 |\theta_\kappa(0)-1|<3/8,\qquad
 |\theta_\kappa(0)|>5/8.
 \tag{R4.8}
\]

Uniformly for every u∈C³ and κ≥1,

\[
 |F_\kappa(u)|<49\exp\!\left(
       \frac{\pi\kappa}{2}\sum_{i=1}^3\frac{|u_i|^2}{G_i}\right).
 \tag{R4.9}
\]

Consequently, for every multi-index α=(α\_1,α\_2,α\_3) of nonnegative integers,

\[
 |\partial^\alpha F_\kappa(0)|
 \le49\,\alpha!\prod_{\alpha_i>0}
       \left(\frac{\pi e\kappa}{\alpha_iG_i}\right)^{\alpha_i/2}.
 \tag{R4.10}
\]

For the source-shaped level normalization κ=M² and N=\textbar α\textbar{} this gives the exact cancellation

\[
 \left|M^{-N}\partial^\alpha F_{M^2}(0)\right|
 \le49\,\alpha!\prod_{\alpha_i>0}
       \left(\frac{\pi e}{\alpha_iG_i}\right)^{\alpha_i/2}.
 \tag{R4.11}
\]

This is a uniform archimedean inequality for an explicit family of sections. It supplies neither a prescribed nonzero positive-degree jet nor its required integral finite-place lattice.

\begin{enumerate}
\def\labelenumi{\arabic{enumi}.}
\setcounter{enumi}{3}
\tightlist
\item
  With probability Haar measure dx dt on J(C), z=x+τt, the metric in the usual theta frame is exp(−2πκ t\^{}T Yt). In this metric the exact squared L² norm of the rigidified section is
\end{enumerate}

\[
 \|\theta_\kappa/\theta_\kappa(0)\|_{L^2}^2
   =\frac{\det(2\kappa Y)^{-1/2}}{|\theta_\kappa(0)|^2}.
 \tag{R4.12}
\]

Thus even in this explicit example, normalization to value 1 and normalization to L² norm 1 are different operations, with the difference recorded rather than discarded.

\subsection{Proof of the lattice and polarization}\label{proof-of-the-lattice-and-polarization}

Use Tadokoro's paths λ\_j from (0,i\_C) to (0,−i\_C) through the branch point ζ\^{}j, and write γ\_j=λ\_\{j+1\}−λ\_j, j=0,\ldots,5. Tadokoro's Proposition 3.2 states that

\[
 a_i^{\rm T}=-\gamma_{2i-1},\qquad
 b_i^{\rm T}=\sum_{k=0}^{i-1}\gamma_{2k},\quad i=1,2,3
\]

are a symplectic basis. This fixes the topological orientation; it is not inferred from a numerical period calculation. Set

\[
 a_j=\gamma_{2j},\qquad b_j=\sum_{k=j}^{2}\gamma_{2k+1},
       \quad j=0,1,2.
\]

These are integral invertible triangular combinations of the previous basis: a\_j=b\^{}T\_\{j+1\}−b\^{}T\_j with b\^{}T\_0=0, and b\_j=−Σ\_\{k=j\}²a\^{}T\_\{k+1\}. Therefore they form a Z basis. Directly their intersections are

\[
 a_i\cdot b_j=\mathbf1_{i\ge j}-\mathbf1_{i-1\ge j}=\delta_{ij},
 \qquad a_i\cdot a_j=b_i\cdot b_j=0.
\]

On the outward segment x=ζ\^{}j t, y=i\_C√(1−t⁷), the beta substitution s=t⁷ gives

\[
 \int_{\lambda_j}\omega_i
   =-\frac{2i_{\mathbb C}}7 B(i/7,1/2)\zeta^{ij}.
\]

The factor 2 includes the return segment on the opposite sheet. Subtracting adjacent paths gives ∫γ\_j ω\_i=c\_i ζ\^{}\{ij\}. Hence the periods on a\_j and b\_j are DA and DB. The six γ\_j themselves are a Z basis, so their period columns give DΦ(Z{[}ζ{]})=DΦ(O\_E), proving (R4.3). The Vandermonde matrix A is invertible because ζ²,ζ⁴,ζ⁶ are distinct.

For ξ=(ζ−ζ\^{}−1)/7, the cyclotomic trace formula is

\[
 \operatorname{Tr}(\zeta^h)=
 \begin{cases}6,&7\mid h,\\-1,&7\nmid h.\end{cases}
\]

It follows that \(\operatorname{Tr}(\xi\zeta^j\overline{\zeta^k})\) is 1 for k=j+1, −1 for j=k+1, and zero otherwise for 0≤j,k≤5. These are exactly the intersections of the γ basis determined above. Thus (R4.4) is the actual principal Riemann form. Its positivity also follows directly from Im σ\_i(ξ)\textgreater0 for i=1,2,3. The ideal ξ O\_E is the codifferent: ζ−ζ\^{}−1 generates the prime above 7, whereas 7 has valuation 6 there, so its exponent is −5, the inverse of the different of Q(ζ\_7). This is consistent with unimodularity and does not replace the explicit intersection check.

Multiplying A by the matrix in (R4.5), reducing by 1+ζ+⋯+ζ⁶=0, gives B. This proves the displayed exact matrix; symmetry is visible. Positive definiteness follows either from the Riemann bilinear relations for the symplectic basis, or from the independent positive lower bound below.

The bilinear relations give \(G=P Y\bar P^T\). Taking determinants gives det G=\textbar det P\textbar²det Y, which is the real covolume of the six columns P and Pτ. The standard complex-embedding lattice of O\_E has covolume 2\^{}−3√\textbar disc E\textbar=7\^{}\{5/2\}/8. Multiplication by D scales real volume by ∏\textbar c\_i\textbar², proving (R4.6). The discriminant formula here is also the determinant of the full cyclotomic embedding matrix; it can be checked from the polynomial discriminant of 1+X+⋯+X⁶.

\subsection{Exact nonvanishing certificate}\label{exact-nonvanishing-certificate}

Put a=sin(π/7), b=sin(2π/7), c=sin(3π/7). Formula (R4.5) gives

\[
 Y=\begin{pmatrix}
 c&a&a+b-c\\
 a&3a+b-c&a\\
 a+b-c&a&c
 \end{pmatrix}.
 \tag{R4.13}
\]

The supplied reproducible script \nolinkurl{scripts/jacobian_theta_certificate.py} uses only exact \nolinkurl{Fraction} arithmetic for its proof assertions. It brackets π by Machin's identity 16 arctan(1/5)−4 arctan(1/239), with 40-term alternating-series bounds, and brackets the three sines with 24-term alternating-series bounds. Outward rounding then gives rational intervals with denominator 10⁶. The smallest Gershgorin lower bound of this rational enclosure is strictly greater than 6/25; all signs of the off-diagonal entries are positive. This proves the first part of (R4.8).

Here is the full finite reduction behind the remaining assertions. For t≥0 define

\[
 U(t)=\left(\sum_{k=0}^{40}\frac{t^k}{k!}\right)^{-1}\ge e^{-t}.
\]

The code rounds U upward to a multiple of 10\^{}−12. For each nonzero n in the cube \textbar n\_i\textbar≤3 it uses the rational interval lower endpoint q\_n for n\^{}T Yn, choosing the lower or upper matrix endpoint according to the sign of n\_i n\_j. Thus q\_n\textgreater0 and

\[
 \sum_{\substack{n\ne0\\|n_i|\le3}}e^{-\pi n^TYn}
 \le\sum_{\substack{n\ne0\\|n_i|\le3}}
       U((314159/100000)q_n).
\]

For the complement of the cube, π λ\_min(Y)\textgreater3/4. Set a\_0=3/4, q=e\^{}−a\_0. Then

\[
 \sum_{k\in\mathbb Z}e^{-a_0 k^2}\le1+\frac{2q}{1-q},\qquad
 \sum_{|k|>3}e^{-a_0 k^2}
 \le\frac{2e^{-16a_0}}{1-e^{-9a_0}}.
\]

The second inequality uses (4+j)²≥16+9j. A union bound on which of the three coordinates lies outside the cube bounds the complement by

\[
 3\frac{2e^{-16a_0}}{1-e^{-9a_0}}
       \left(1+\frac{2e^{-a_0}}{1-e^{-a_0}}\right)^2.
\]

Replacing the three exponentials by U with upward rounding gives a rational upper bound. The final rational comparison is strictly less than 3/8 (its separately labelled decimal approximation is 0.366129006419). All numerators and denominators and all inequalities are retained in \nolinkurl{evidence/theta_certificate.json}. This proves Σ\_\{n≠0\}exp(−πn\^{}TYn)\textless3/8. Increasing κ decreases every absolute summand. Consequently \textbar θ\_κ(0)−1\textbar\textless3/8 for all κ≥1, proving (R4.8). No floating-point approximation is a premise in this certificate.

\subsection{Sections, Chern jets and the uniform bound}\label{sections-chern-jets-and-the-uniform-bound}

The series θ\_κ converges normally on compact sets. Direct reindexing gives, for a,b∈Z³,

\[
 \theta_\kappa(z+a+\tau b)
  =e^{-\pi i\kappa b^T\tau b-2\pi i\kappa b^Tz}\theta_\kappa(z).
\]

This is the factor of automorphy of L\^{}κ, so these are genuine sections on the actual abelian variety. The line L is a symmetric representative of its principal polarization. We have not asserted that this representative or its particular section is fixed by the order-seven action. That distinction matters for a prescribed CM-labelled automorphic row.

In the usual holomorphic theta frame the metric is h(z)=exp(−2πκ(Im z)\^{}T Y\^{}−1 Im z). Its Chern connection has (1,0) coefficient \(\pi\kappa Y^{-1}(z-\bar z)\). Put Q(z)=πκ z\^{}TY\^{}−1z/2. Replacing the frame e by e′=e\^{}−Q e makes the connection coefficient \(-\pi\kappa Y^{-1}\bar z\) and the metric \(\exp(-\pi\kappa\bar z^TY^{-1}z)\). Thus on the formal holomorphic directions at 0, covariant derivatives are exactly the ordinary derivatives of e\^{}Qθ\_κ. Normalization by the nonzero θ\_κ(0) proves the claimed interpretation of F\_κ.

For y=Im z, completing the square yields

\[
 |\theta_\kappa(z)|
 \le e^{\pi\kappa y^TY^{-1}y}
      \sum_{n\in\mathbb Z^3}
        e^{-\pi\kappa(n+Y^{-1}y)^TY(n+Y^{-1}y)}.
\]

For every λ\textgreater0 and real t, the elementary unimodal Gaussian bound is

\[
 \sum_{n\in\mathbb Z}e^{-\pi\lambda(n+t)^2}
      \le1+\lambda^{-1/2}.
\]

To see it, keep a nearest-to-zero sample as a term at most 1. Every sample to its left or right is at most the Gaussian integral on the adjoining unit interval toward that nearest sample; these intervals partition the real line. The remaining sum is at most the Gaussian integral λ\^{}−1/2. Applying this in three coordinates, using κ≥1 and λ\_min(Y)\textgreater6/25, bounds the displayed three-dimensional sum by (1+5/√6)³\textless(25/8)³.

Multiplication by \textbar e\^{}Q\textbar{} changes the exponential to

\[
 \exp\left(\frac{\pi\kappa}{2}\bar z^T Y^{-1}z\right).
\]

Since \(G=P Y\bar P^T\) is real diagonal, \(\bar z^TY^{-1}z=\sum_i|u_i|^2/G_i\). Dividing by \textbar θ\_κ(0)\textbar\textgreater5/8 leaves a constant smaller than (25/8)³(8/5)=3125/64\textless49. This proves (R4.9).

Cauchy's integral inequality on a polydisc of radii R\_i gives

\[
 |\partial^\alpha F_\kappa(0)|
 \le49\alpha!\prod_i
      \left(e^{\pi\kappa R_i^2/(2G_i)}R_i^{-\alpha_i}\right).
\]

For α\_i\textgreater0 minimize at R\_i²=α\_i G\_i/(πκ); for α\_i=0 take the limit R\_i→0. This proves (R4.10). With κ=M², the factor κ\^{}\{\textbar α\textbar/2\} is precisely M\^{}N, proving (R4.11). This calculation is an explicit instance of the source normalization issue already emphasized in DP50; it does not alter DP50's carried proof status.

Finally integrate the squared section norm with z=x+τt. Integration in x over {[}0,1{]}³ kills all cross terms n≠m. The remaining integrand is Σ\_n exp(−2πκ(n+t)\^{}TY(n+t)). Integration in t unfolds the unit cubes into R³, whose Gaussian integral is det(2κY)\^{}−1/2. Division by \textbar θ\_κ(0)\textbar² proves (R4.12). All integrations are justified by normal convergence or positive-term Tonelli integration. This completes the proof.

\subsection{What remains before this is a lift-count input}\label{what-remains-before-this-is-a-lift-count-input}

The statement now fixes an actual mixed-type CM lattice, polarization, complex covolume, nonzero section family, metric, Chern gauge and level-uniform archimedean jet bound. The new exact finite Hodge lattices in ALP-LIFT-05 complement it. These are concrete geometric data that can be used in a replacement arithmetic constructor.

They do not identify F\_κ with the primitive trace-one Poincaré kernel of the selected source row. Nor do they establish integral normalized coefficients, the allowable-translate expansion with controlled algebraic coefficient height, membership in the chosen automorphic theta image, all-degree \nolinkurl{ClsK.Core}, \nolinkurl{SizeData₂}, or the six identities for that tied actual family. In particular, θ\_κ(0)≠0 proves nonvanishing of the section at the origin; it is not controlled nonvanishing of the prescribed high-degree harmonic jet or terminal lift. Passing from a nonzero complex section to such a family remains a substantive arithmetic step.

In Grothendieck's language, (R4.3--6) fix the Betti lattice and polarization of the Picard variety; (R4.7--12) fix its complex metrized theta line and a Hodge-splitting jet calculation. The Poincaré biextension framework is SGA 7 I, Exposés VII--VIII; the integral crystalline realization and its filtration remain the objects of cote 54 scans 3 and 5--7 and cote 4 scan 3, (8.1.1--2). No assertion in those manuscripts is replaced by the complex analytic bounds here.

\subsection{References}\label{references}

Tadokoro, Y. (2014). The period matrix of the hyperelliptic curve w²=z\^{}(2g+1)−1. \emph{Tsukuba Journal of Mathematics, 38}(2), 137--158. The version actually consulted is arXiv:1211.6910v5, Proposition 3.2, Proposition 4.1, Theorem 4.4, pp.~9--10, 13--15. \url{https://arxiv.org/abs/1211.6910}.

Supplied programme sources: DP\_MULT DP49, DP50, DP57, DP64, DP65 and DP67; THETA3 Part II G1. Their statuses are carried without upgrade. The exact new rational certificate and the symbolic period check are documented separately.

\clearpage

\section{AM-R-005: a uniform bound for every nonsingular even theta characteristic}\label{am-r-005-a-uniform-bound-for-every-nonsingular-even-theta-characteristic}

Author: GPT-6 Astra Pro, \nolinkurl{/root}; ALP\_MACHINE, 21 November 2026 UTC. Reasoning effort is not separately exposed. This extends the one-characteristic archimedean calculation AM-R-004 to all algebraic conjugates and to normalized Mumford translates. It does not assert finite-place integrality of the reduced formal frame, membership in the selected arithmetic rows, or a lift-cost sign.

\subsection{Statement}\label{statement-3}

Retain the actual curve, differentials, period matrix and metric of AM-R-004: \[
C:y^2=x^7-1,\quad E=\mathbf Q(\zeta_7),\quad
\omega_i=x^{i-1}dx/y,\quad u=Pz,\quad Y=\operatorname{Im}\tau,
\quad P Y\overline P^{\,t}=\operatorname{diag}(G_1,G_2,G_3).
\tag{A1}
\] For a three-element subset \(S\subset\{0,\ldots,6\}\), let \(D_S=\sum_{j\in S}W_j-\infty\), with \(W_j=(\zeta_7^j,0)\). Let \(\delta_S=[a_S;b_S]\) be its nonsingular even half-characteristic in this period basis, with \(a_S,b_S\in\{0,1/2\}^3\). Put \[
F_S(u)=\exp\!\left(\frac\pi2 z^tY^{-1}z\right)
       \frac{\theta[\delta_S](z;\tau)}{\theta[\delta_S](0;\tau)}.
\tag{A2}
\] The classification and nonsingularity of these 35 characteristics are AM-DO-012. Their exact identification includes \(S_0=\{1,3,5\}\), whose characteristic is zero, in ALP-LIFT-07.

For all such \(S\), \[
\left|\frac{\theta[\delta_S](0;\tau)}{\theta[0;0](0;\tau)}\right|
   >\frac18,
\qquad
\boxed{|F_S(u)|<392\exp\!\left(\frac\pi2
                           \sum_{i=1}^3 |u_i|^2/G_i\right).}
\tag{A3}
\] In particular \(|\theta[\delta_S](0;\tau)|>5/64\). The constant is deliberately a simple certified uniform bound, not an optimum. For every multi-index \(\alpha\in\mathbf N^3\), with the usual convention for a zero component, \[
|\partial^\alpha F_S(0)|
 \le392\,\alpha!\prod_{\alpha_i>0}
             \left(\frac{\pi e}{\alpha_iG_i}\right)^{\alpha_i/2}.
\tag{A4}
\] The same bound holds for every normalized Mumford translate of \(F_S\), even when the translated section vanishes at the origin. For its pullback under \([M]\), the level-normalized derivative \(M^{-|\alpha|}\partial^\alpha\) again satisfies (A4).

Each \(F_S\) has coefficients in \(E\). If \(\sigma_a(\zeta_7)=\zeta_7^a\), then coefficientwise conjugation in the rational differential coordinates satisfies \[
                    \sigma_a(F_S)=F_{aS}.
\tag{A5}
\] Consequently (A4) bounds every conjugate of each algebraic derivative. For the Mumford translates, every embedding of their finite torsion field is bounded in the same way, after the corresponding torsion translation; the possible change of a period representative is a factor of absolute value one.

\subsection{Proof of the theta-null comparison}\label{proof-of-the-theta-null-comparison}

The only external theta-constant input is the classical first Thomae formula in the precise form stated by Eilers (2018), Theorem 2.1, equation (2.17), with the partition convention of equations (2.12)--(2.16). It applies to an odd-degree hyperelliptic equation with the branch point at infinity omitted from the Vandermonde products. Taking absolute values and a ratio removes both its common period factor and its eighth-root-of-unity ambiguity. Thus \[
\left|\frac{\theta[\delta_S](0)}{\theta[\delta_{S_0}](0)}\right|^4
 =\frac{\displaystyle
     \prod_{i<j\in S}|\zeta_7^i-\zeta_7^j|
     \prod_{i<j\notin S}|\zeta_7^i-\zeta_7^j|}
    {\displaystyle
     \prod_{i<j\in S_0}|\zeta_7^i-\zeta_7^j|
     \prod_{i<j\notin S_0}|\zeta_7^i-\zeta_7^j|}.
\tag{A6}
\] Every numerator and denominator has exactly \(\binom32+\binom42=9\) factors.

For distinct indices every factor is between \(2\sin(\pi/7)\) and 2. The lower bound is strictly greater than \(4/5\). One elementary verification uses \(\pi>3\), the monotonicity of sine on \([0,\pi/2]\), and its alternating Taylor lower bound: \[
 \sin(\pi/7)>\sin(3/7)\ge 3/7-(3/7)^3/6
                 =285/686>2/5.
\] It follows from (A6) that its right side is greater than \((2/5)^9\). The exact integer inequality \(2^{21}=2097152>1953125=5^9\) says \((2/5)^9>1/8^4\). Taking the positive fourth root gives the first inequality of (A3). AM-R-004 gives \(|\theta[0;0](0)|>5/8\), proving the asserted absolute lower bound.

\subsection{Proof of the global growth and derivative bounds}\label{proof-of-the-global-growth-and-derivative-bounds}

For real half-characteristic vectors \(a,b\) and \(z=x+iy\), completing the real quadratic form in the absolutely convergent theta series gives \[
 e^{-\pi y^tY^{-1}y}\,|\theta[a;b](z;\tau)|
 \le\sum_{n\in\mathbf Z^3}
 e^{-\pi(n+a+Y^{-1}y)^tY(n+a+Y^{-1}y)}.
\tag{A7}
\] The phase containing \(b\) has absolute value one. AM-R-004 proves \(\lambda_{\min}(Y)>6/25\) and the shifted one-dimensional Gaussian bound used there. Hence the right side of (A7) is at most \[
                         C_0=(1+5/\sqrt6)^3.
\] The same exact inequalities in AM-R-004 give \(C_0/|\theta[0;0](0)|<49\). Combining this with (A6) bounds \(C_0/|\theta[\delta_S](0)|\) by 392. Finally \[
 \operatorname{Re}(z^tY^{-1}z)/2+y^tY^{-1}y
                         =z^tY^{-1}\overline z/2,
\] and (A1) turns this into \(\tfrac12\sum_i|u_i|^2/G_i\). This proves the second inequality of (A3).

Apply polydisk Cauchy with \(R_i^2=\alpha_iG_i/\pi\) when \(\alpha_i>0\). For a zero component let its radius decrease to zero. The resulting product is exactly (A4). No equality of actual derivative size with this upper bound is claimed.

Let \(H(u,v)=\sum_i u_i\overline v_i/G_i\). A normalized translation by a period-cover representative \(t\) has the analytic form \[
 (U_tF)(u)=\varepsilon_t
       e^{-\pi H(u,t)-\pi H(t,t)/2}F(u+t),\qquad |\varepsilon_t|=1.
\tag{A8}
\] This is Bannai--Kobayashi's normalized translation convention; its Mumford version has the explicitly specified finite theta-group factor. Since \[
 H(u+t,u+t)-2\operatorname{Re}H(u,t)-H(t,t)=H(u,u),
\] the exponential cross terms cancel exactly in the absolute-value bound. Thus (A3) and (A4) hold for \(U_tF_S\) with the same constant. Substitution of \(Mu\) and the chain rule prove the level-normalized assertion. Replacing \(t\) by \(t+\lambda\), \(\lambda\in\Lambda\), changes the normalized translate only by its semicharacter and the exponential of a purely imaginary pairing, both of absolute value one.

\subsection{Algebraicity and the conjugation assertion}\label{algebraicity-and-the-conjugation-assertion}

AM-DO-011 applies Bannai--Kobayashi (2010), arXiv v4 Proposition 2.1, to these symmetric algebraic divisors, their admissible polarization, the algebraic differential coordinates, and value 1 at the origin. The same hypotheses hold for every \(D_S\), so \(F_S\in E[[u_1,u_2,u_3]]\).

We justify (A5), since algebraicity alone is not a statement about the analytic growth of conjugates. The coefficient construction in the proof of that proposition is Galois equivariant in this situation. More explicitly choose \(b\in2\mathcal O_E\) there so that every positive-degree multi-index has unequal eigenvalue products for \(b\) and \(\overline b\). The normalized periodic function \[
 R_{S,b}(u)=F_S(\iota(b)u)/F_S(\iota(\overline b)u)
\tag{A9}
\] is the unique rational function of value 1 at the origin with divisor \(\iota(b)^*\Theta_S-\iota(\overline b)^*\Theta_S\). Its Taylor coefficients are algebraic, and (A9) recursively determines every coefficient of \(F_S\) from the earlier ones. The nonzero coefficient multiplying the next unknown is exactly the difference of the two displayed eigenvalue products. These facts are the proof of the cited proposition, not a new replacement for its algebraicity theorem.

The curve and \(\omega_i\) are rational, while \(\sigma_a\iota(b)=\iota(\sigma_ab)\) because the generating endomorphism is \((x,y)\mapsto(\zeta_7x,y)\). Conjugation commutes with \(\sigma_a\), and \(\sigma_aD_S=D_{aS}\). Therefore \(\sigma_aR_{S,b}\) is exactly \(R_{aS,\sigma_ab}\), by its normalized rational divisor description. The transformed eigenvalues are the corresponding eigenvalues of \(\iota(\sigma_ab)\) on the same rational differentials. The recursive equations and initial value 1 now agree; induction proves (A5).

For torsion translations use the canonical algebraic isomorphisms in Bannai--Kobayashi Proposition 2.3 and Theorem 2.9. Formation of those isomorphisms commutes with coefficient-field embeddings. Analytically they are the translations (A8) by the embedded torsion point, up to the specified finite-order factor. The period-representative observation following (A8) removes that harmless ambiguity from absolute values. This proves the final assertion and completes the proof. \(\square\)

\subsection{Scope for the main question and Grothendieck's frameworks}\label{scope-for-the-main-question-and-grothendiecks-frameworks}

This closes the uniform archimedean normalization subgate for the 35 normalized even theta germs and their canonical algebraic torsion translations on this fixed Jacobian. It retains the exact metric \(G_i\), instead of substituting an arbitrary complex-coordinate norm. Together with ALP-LIFT-06 it shows that passing to a different such theta characteristic changes this certified bound by a fixed additive constant, not by an order-linear metric advantage. This does not bound the height of a general varying linear combination, the finite denominators of its theta-group projector, or its actual vanishing and divisibility. Those data remain in Q1 and THETA3 G1.

Geometrically the characteristic change is translation of the relative Picard theta line by a section of its finite etale two-torsion scheme over \(\mathbf Z[\zeta_7,1/14]\). The compatible translation isomorphisms belong to the symmetric-line/theta-group structure and the Poincare biextension, not to an arbitrary identification of functions. The Hodge metrics and determinant lines are the ones in ALP-LIFT-05; the precise FGA, cote54 and SGA7 dictionary is recorded separately. No new claim about number-field L-function zeros or finite-field weights follows from the uniform analytic norm bound.

\subsection{References and check provenance}\label{references-and-check-provenance}

\begin{itemize}
\tightlist
\item
  Eilers, K. (2018). Rosenhain--Thomae formulae for higher genera hyperelliptic curves. \emph{Journal of Nonlinear Mathematical Physics, 25}(1), 86--105. \url{https://doi.org/10.1080/14029251.2018.1440744}. Exact input: Theorem 2.1, equation (2.17), and the odd-degree partition convention (2.12)--(2.16). The primary article text was read at \url{https://www.atlantis-press.com/journals/jnmp/125950502/view}.
\item
  Bannai, K., \& Kobayashi, S. (2010). Algebraic theta functions and the p-adic interpolation of Eisenstein--Kronecker numbers. \emph{Duke Mathematical Journal, 153}(2), 229--295. \url{https://doi.org/10.1215/00127094-2010-024}. Exact text used: arXiv:math/0610163v4, Definition 1.8, Proposition 2.1 and its proof, Proposition 2.3, and Theorem 2.9. No ordinary-integrality result is extended to a nonordinary prime here.
\item
  AM-R-004; AM-DO-011--013; ALP-LIFT-06--07. These dependencies retain their separate whole-proof-check events. The elementary bound above is a new corollary of the cited known Thomae formula, whose proof is not repeated.
\end{itemize}

\clearpage

\section{ALP-LIFT-07: the actual theta characteristic, its finite normalization, and the source line}\label{alp-lift-07-the-actual-theta-characteristic-its-finite-normalization-and-the-source-line}

Run ALP\_MACHINE. Author: \nolinkurl{/root/lift_accounting}, GPT-6 Astra Pro; reasoning effort is not separately exposed. The good models and Hodge bases are the proved ALP-LIFT-05 inputs. The homology basis is AM-R-004. The branch half-period table also occurs independently in AM-DO-012; credit it once in a final synthesis. The new finite calculation uses the determinant of cohomology, so no unspecified Thomae scalar enters it.

\subsection{Statement and normalization}\label{statement-and-normalization}

Put \[
 C:y^2=x^7-1,\quad J=\operatorname{Jac}(C),\quad
 P_a=(\zeta_7^a,0)\ (0\leq a\leq6),\quad P_\infty=\infty.
 \tag{T1}
\] For a three-element subset \(S\subset\{0,\ldots,6\}\), set \[
 E_S=\sum_{a\in S}P_a,\qquad D_S=E_S-P_\infty.
 \tag{T2}
\] At each finite place work over a complete DVR \(R\), after a finite extension large enough for ALP-LIFT-05's good curve model and these sections. Write \(v_p(p)=1\). On the relative degree-two Picard scheme take its canonical determinant theta divisor, pull it back to \(\mathcal J\) by \(x\mapsto\mathcal O(D_S)\otimes x\), and denote the resulting line and canonical section by \((\mathcal L_S,s_S)\). The section has no vertical divisor: on each geometric fibre it is the ordinary nonzero theta section. This is the primitive integral theta model meant here. The valuation of its nonzero origin value means its valuation relative to any primitive generator of the free rank-one module \(e^*\mathcal L_S\); this is independent of the generator's integral-unit ambiguity.

\textbf{(a) Exact finite value.} All 35 choices have \(s_S(0)\ne0\) on the generic fibre, and \[
 \boxed{v_p(s_S(0))=
 \begin{cases}6/7,&p=2,\\0,&p\text{ odd}.\end{cases}}
 \tag{T3}
\] The entry at 7 is zero on the good model after extension. Normalization of the value to 1 therefore contributes the fixed scalar valuation \(-6/7\) at 2 and zero at every odd prime, in these local integral origin frames.

\textbf{(b) Characteristic of the explicit analytic germ.} In AM-R-004's symplectic basis \[
 a_j=\gamma_{2j},\qquad b_j=\sum_{k=j}^2\gamma_{2k+1}
 \quad(0\leq j\leq2),
 \tag{T4}
\] write a half-period as \((m+\tau n)/2\), with \(m,n\in\mathbb F_2^3\). The Riemann constant for infinity is \[
 K_\infty=\frac{(1,0,1)^t+\tau(1,1,1)^t}{2}
 =\operatorname{AJ}(P_1+P_3+P_5-3P_\infty).
 \tag{T5}
\] Thus \(D_{135}=P_1+P_3+P_5-P_\infty\) corresponds to the \textbf{zero analytic characteristic}. Its theta divisor is the zero locus of \(\theta(z,\tau)\), and AM-R-004's value-one germ \[
 F_1(u)=\exp\!\left(\frac\pi2z^t(\operatorname{Im}\tau)^{-1}z\right)
 \frac{\theta(z,\tau)}{\theta(0,\tau)},\qquad u=Pz,
 \tag{T6}
\] belongs to this precise principal line. AM-DO-011 identifies it as the reduced theta in the specified CM differential coordinates and applies Bannai--Kobayashi's algebraicity theorem. Those are separate carried findings; (T3) is an additional integral origin calculation.

\textbf{(c) The trace-one line is a different principal line.} Let \(a\) be induced by \(x\mapsto\zeta_7x\), let \(w_0=-(a+a^2+a^4)\), and put \(L_{w_0}=(1,w_0)^*\mathcal P_\lambda\), as in accepted DP49. Let \(L_\infty\) be the theta line obtained from \(2P_\infty\). Superscript \(\mathrm{rig}\) means tensoring a line with the inverse of its origin fibre, giving a rigidified line. On the good models as well as generically, \[
 L_{w_0}\simeq L_\infty^{\mathrm{rig}}.
 \tag{T7}
\] Its unique principal theta section vanishes to order \textbf{two} at zero. In contrast, \[
 L_{135}^{\mathrm{rig}}\simeq L_{w_0}\otimes\mathcal P_T,
 \qquad T=K_\infty\in J[2]\setminus\{0\},
 \tag{T8}
\] using the canonical polarization to identify \(J\) with its dual. For every positive even integer \(M\) there is an integral rigidified isomorphism \[
 [M]^*L_{135}^{\mathrm{rig}}\simeq L_{w_0}^{\otimes M^2}.
 \tag{T9}
\] The pulled-back primitive section retains (T3). Its normalized reduced germ is \(F_1(Mu)\), and \[
 M^{-|\alpha|}\partial_u^\alpha F_1(Mu)|_{u=0}
       =\partial_u^\alpha F_1(u)|_{u=0}.
 \tag{T10}
\] This construction is \textbf{outside the present source \nolinkurl{Level} record}, which requires \nolinkurl{Odd M}. For odd \(M\) the nonzero two-torsion twist survives pullback. An odd-order theta-group translation within \(H^0(J,L_{w_0}^{M^2})\) cannot change its underlying line into that of \([M]^*s_{135}\). This is a line distinction, not a claim that allowed odd-order translates have zero origin value.

\subsection{Proof of the determinant normalization}\label{proof-of-the-determinant-normalization}

\subsubsection{1. The determinant theta model and its dual complex}\label{the-determinant-theta-model-and-its-dual-complex}

The degree \(g-1\) theta divisor has its standard determinant-of- cohomology description: a perfect two-term complex of equal-rank vector bundles computes the universal degree \(g-1\) line's cohomology, and the determinant of its differential is the theta section. See Casalaina-Martin and Kass (2010, Definition 3.1 and the ensuing paragraphs, printed pp.~9--10). Invariance of this determinant element under a change of perfect representative is the determinant functor formalism (Stacks Project, Tag 0FJI). We use this construction over the DVR; no equal-characteristic arc theorem is asserted to cover the present mixed-characteristic family.

Write \(\omega=\Omega^1_{\mathcal C/R}\). The exact sequence \[
 0\to\omega(P_\infty-E_S)\to\omega(P_\infty)
   \to\omega(P_\infty)|_{E_S}\to0
 \tag{T11}
\] gives a simple complex. On every geometric fibre, \(H^1(\omega(P_\infty))=0\) by Serre duality and \(h^0(\omega(P_\infty))=3\) by Riemann--Roch. The inclusion \(H^0(\omega)\to H^0(\omega(P_\infty))\) is an isomorphism on each fibre: both have dimension three, or a differential cannot have a single nonzero simple-pole residue. Cohomology and base change and Nakayama therefore identify these integral rank-three modules. Since \(E_S\) is finite flat of degree three, its restricted invertible sheaf has free rank three over \(R\). Hence \[
 r_S:H^0(\omega(P_\infty))\longrightarrow
 H^0(E_S,\omega(P_\infty)|_{E_S})
 \tag{T12}
\] is a map between free rank-three modules. Its complex in degrees zero and one computes \(R\Gamma(\omega\otimes\mathcal O(-D_S))\).

Relative Serre duality identifies this with the dual shifted complex of \(R\Gamma(\mathcal O(D_S))\); see Stacks Project, Tag 0E6N. For a free complex \([A\xrightarrow d B]\) in degrees zero and one, the dual shifted complex is \([B^*\xrightarrow{\pm d^t}A^*]\), which has the same determinant valuation. Thus the theta section at \(D_S\) has valuation \(v_p(\det r_S)\). Equivalently, Smith normal form says that the determinant order of this generically acyclic complex is the sum of its elementary-divisor orders, the length of its torsion cokernel. This elementary argument works over every DVR. We have not used an unverified integral identification of \(\omega(-D_S)\) with \(\mathcal O(D_S)\).

The canonical relative theta section is primitive because its restriction to every geometric Picard fibre is a proper effective theta divisor, not that entire fibre. The translated pullback is primitive for the same reason. The determinant in (T12) therefore computes the stated model, with no arbitrary vertical scalar.

\subsubsection{2. Every odd-prime value is a unit}\label{every-odd-prime-value-is-a-unit}

For \(p\nmid14\), the original good model has seven distinct finite branch points, separate from infinity. At 7 use ALP-LIFT-05's good model \[
 Y^2=\frac{(1+\lambda X)^7-1}{\lambda^7},\qquad
 \lambda=\zeta_7-1,
 \tag{T13}
\] whose special fibre is \(Y^2=X^7-X\). The branch coordinates are \[
 X_a=(\zeta_7^a-1)/\lambda\equiv a\pmod\lambda,
 \tag{T14}
\] so they are again distinct. On each smooth odd-characteristic fibre, \(K_C\sim4P_\infty\) and \(H^0(\mathcal O(5P_\infty))=\langle1,X,X^2\rangle\). Riemann--Roch and evaluation at three distinct coordinates give \[
 h^0(D_S)=h^0(5P_\infty-E_S)=0:
 \tag{T15}
\] the quadratic-polynomial evaluation matrix is an invertible Vandermonde. This also proves the generic nonvanishing. Since the degree is \(g-1\), first cohomology vanishes as well. The matrix (T12) is therefore invertible after reduction, proving all the odd entries of (T3). No parity assertion in characteristic 2 was used.

\subsubsection{3. The dyadic collision algebra}\label{the-dyadic-collision-algebra}

Use the good model and integral Hodge basis from ALP-LIFT-05: \[
 Y^2+Y=X^7,\quad x=\alpha X,\quad y=\iota(2Y+1),\quad
 \alpha^7=-4,\quad\iota^2=-1,\qquad
 \eta_i=X^{i-1}\frac{dX}{2Y+1}\ (1\leq i\leq3).
 \tag{T16}
\] On the infinity chart \(U=1/X\), \(V=YU^4\), the equation is \(V^2+U^4V=U\). The rational function \[
 t=X^3/Y=U/V=V+U^4
 \tag{T17}
\] is a uniformizer at the infinity section. Indeed \(U\) is a formal power series in \(V\) beginning with \(V^2\), so \(t=V+O(V^8)\). The completed ring is \(R[[t]]\), and \(dt/t\) is an integral generator of \(\omega(P_\infty)\) on this disk. At \(P_a\), \[
 X_a=\zeta_7^a/\alpha,\quad Y_a=-1/2,\quad
 t_a=-2\alpha^{-3}\zeta_7^{3a}.
 \tag{T18}
\] All seven sections specialize to infinity. Since \(v_2(\alpha)=2/7\), \[
 v_2(t_a)=1/7,\qquad v_2(t_a-t_b)=1/7\quad(a\ne b).
 \tag{T19}
\] The differences of distinct seventh roots of unity are dyadic units: their reductions are distinct since \(X^7-1\) is separable in characteristic 2. The actual finite flat collision algebra is thus \[
 H^0(E_S,\mathcal O_{E_S})
   =R[[t]]\big/\!\prod_{a\in S}(t-t_a).
 \tag{T20}
\] The product is a distinguished monic cubic, so Weierstrass division gives the integral basis \(1,t,t^2\). Inverting the uniformizer splits the algebra into the three point evaluations. That split evaluation basis is not integral: its determinant relative to \(1,t,t^2\) is a Vandermonde of valuation \(3/7\).

Write \(\eta_i=f_i(t)dt/t\). Differentiating (T17) and using \(X^7=Y^2+Y\) yields the exact identity \[
 dt=-\frac{X^2(Y+4)}{Y(2Y+1)}dX,\qquad
 \frac{\eta_i}{dt}=-\frac{YX^{i-3}}{Y+4}.
 \tag{T21}
\] The ratio extends at every \(P_a\). Substitution gives \[
 f_i(t_a)=t_a\frac{\eta_i}{dt}(P_a)=-\frac27X_a^i.
 \tag{T22}
\] Up to the ordering sign, this split evaluation determinant is \[
 \det(f_i(t_a))=
 \left(-\frac27\right)^3
 \left(\prod_{a\in S}X_a\right)
 \prod_{a<b\text{ in }S}(X_b-X_a).
 \tag{T23}
\] Each \(X_a\) and each pairwise difference has valuation \(-2/7\). Thus (T23) has valuation \[
 3-6/7-6/7=9/7.
 \tag{T24}
\] The conversion matrix from the integral target basis \((1,t,t^2)dt/t\) to split evaluation has determinant \(\prod_{a<b}(t_b-t_a)\), of valuation \(3/7\). Consequently the actual integral restriction map in (T12) has \[
 v_2(\det r_S)=9/7-3/7=6/7.
 \tag{T25}
\] This proves the dyadic entry of (T3) for every triple. It also explains why using only split point values would give the incorrect answer \(9/7\). The special-fibre class \(D_S\) is \(2P_\infty\), whose space of sections has dimension two, consistent with the nonunit theta value.

\subsection{Proof of the characteristic and source-line assertions}\label{proof-of-the-characteristic-and-source-line-assertions}

\subsubsection{4. Half-periods and the Riemann constant}\label{half-periods-and-the-riemann-constant}

In this subsection only, write \(t_a=\operatorname{AJ}(P_a-P_\infty)\); this is distinct from the parameter values in (T18). Each \(t_a\) is two-torsion because \(\operatorname{div}(x-\zeta_7^a)=2P_a-2P_\infty\). The adjacent-arc relation for the paths of AM-R-004 is \[
 t_a-t_0=\frac12\sum_{j=0}^{a-1}\gamma_j
                    \pmod{H_1(C,\mathbb Z)}.
 \tag{T26}
\] This is the hyperelliptic doubled-arc calculation: the two sheets give opposite integrals between the ramification endpoints, so the closed cycle is twice the Abel integral. It also follows from AM-DO-012's independently evaluated branch half-periods. Since \(\operatorname{div}(y)=\sum_aP_a-7P_\infty\), one has \(\sum_at_a=0\). Summing (T26), using two-torsion, and reducing cycle coefficients modulo two gives \[
 t_0=(\gamma_1+\gamma_3+\gamma_5)/2=b_0/2.
 \tag{T27}
\] Substitution of \(\gamma_{2j}=a_j\), \(\gamma_1=b_0-b_1\), \(\gamma_3=b_1-b_2\), and \(\gamma_5=b_2\) now gives

\begin{longtable}[]{@{}rcc@{}}
\toprule\noalign{}
\(a\) & \(m_a\) & \(n_a\) in \(t_a=(m_a+\tau n_a)/2\) \\
\midrule\noalign{}
\endhead
\bottomrule\noalign{}
\endlastfoot
0 & 000 & 100 \\
1 & 100 & 100 \\
2 & 100 & 010 \\
3 & 110 & 010 \\
4 & 110 & 001 \\
5 & 111 & 001 \\
6 & 111 & 000 \\
\end{longtable}

The Riemann constant is determined from this table and parity, rather than assumed from the labels. The theta characteristic \(2P_\infty\) has \(h^0=2\). Its twist by any \(t_a\) is \(P_\infty+P_a\), with \(h^0=1\): Riemann--Roch identifies this dimension with the space of sections of \(3P_\infty-P_a\), a codimension-one subspace of \(\langle1,x\rangle\). The standard Riemann theta correspondence identifies algebraic and analytic characteristic parity; equivalently use Riemann singularity and analytic theta parity. These classical inputs are recalled by Mumford (1971, pp.~181--182) and Casalaina-Martin and Kass (2010, p.~1).

In the convention of DLMF 21.2.5--21.2.7, the analytic characteristic of \((m+\tau n)/2\) is \([n/2;m/2]\) and has parity \(n\cdot m\). This follows directly by replacing the theta summation index by \(-k-n\). If \((m_K,n_K)\) is the characteristic of \(2P_\infty\), its parity is even and all seven twists above are odd. Therefore \[
 n_a\cdot m_a+n_K\cdot m_a+m_K\cdot n_a=1
                         \quad\text{in }\mathbb F_2.
 \tag{T28}
\] The rows \(a=0,1,2,3,4,5\), in succession, force \[
 m_K=(1,0,1),\qquad n_K=(1,1,1).
 \tag{T29}
\] The last row is consistent and this characteristic is even. The table also gives \(t_1+t_3+t_5=K_\infty\). The divisor for \(2P_\infty\) is analytically \(\theta(z+K_\infty,\tau)=0\) up to a nowhere-vanishing factor. Translating its degree-two class by \(t_1+t_3+t_5=K_\infty\) changes it to \(D_{135}\) and adds another \(K_\infty\) to the theta argument. Since \(2K_\infty\) is a period, the divisor becomes \(\theta(z,\tau)=0\). This proves (T5)--(T6). All parity arguments here take place over \(\mathbb C\), not in residue characteristic 2.

\subsubsection{5. Identification of the invariant principal line}\label{identification-of-the-invariant-principal-line}

Use the accepted DP49 identities \(a^\dagger=a^{-1}\), \(w_0+w_0^\dagger=1\), \(w_0w_0^\dagger=2\), and its principal polarization conclusion for \(L_{w_0}\). Symmetry and multiplication isomorphisms are also in accepted DP59, proof of (TS7)--(TS8). The same Poincaré functoriality applies to each good abelian scheme in ALP-LIFT-05. In particular \[
 a^*L_{w_0}\simeq(1,a^\dagger w_0a)^*\mathcal P_\lambda
                 \simeq L_{w_0}.
 \tag{T30}
\] The theta line \(L_\infty\) is symmetric as well: relative duality sends \(2P_\infty\otimes x\) to \(2P_\infty\otimes x^{-1}\), since \(\omega\simeq\mathcal O(4P_\infty)\). This canonical- divisor identity holds on the good models, where the first integral differential has divisor exactly \(4P_\infty\). The automorphism fixes infinity and preserves effectivity, so it preserves this theta divisor.

Two rigidified principal theta lines with the same polarization differ by a rigidified line in \(\operatorname{Pic}^0\). Because both are symmetric, that difference is two-torsion; because both are \(a\)-invariant, it is fixed by the dual \(a\)-action. But \(a-1\) is invertible on two-torsion. Indeed \(\Phi_7(a)=0\), and \[
 (T-1)(T^5+2T^4+3T^3+4T^2+5T+6)=\Phi_7(T)-7
 \tag{T31}
\] shows that \(a^5+a^3+a\) is an inverse to \(a-1\) on \(J[2]\). The dual action has the same property. This is a group-scheme endomorphism identity, so it remains valid on two-torsion in residue characteristic 2. The invariant difference is therefore trivial, proving (T7) integrally with rigidifications.

The space of sections of \(2P_\infty\) is exactly \(\langle1,x\rangle\), so Riemann singularity gives theta multiplicity two at zero. Translation by \(T=K_\infty\) and the definition of the polarization homomorphism give (T8). The explicit table shows \(T\ne0\).

\subsubsection{6. Pullbacks and the exact source limitation}\label{pullbacks-and-the-exact-source-limitation}

For a rigidified symmetric line, \([M]^*L\simeq L^{M^2}\); for a rigidified degree-zero line, \([M]^*\mathcal P_T\simeq\mathcal P_{MT}\). These are the theorem-of-the-cube identities; for \(L_{w_0}\) the first is exactly the integral Poincaré identity in DP59 (TS7). Applied to (T8), they prove (T9) when \(M\) is even.

Choose a primitive local generator of \(e^*\mathcal L_{135}\) to transfer its canonical section to the rigidified line. Pullback by \([M]\) retains the origin value. It retains primitivity as well: multiplication on an abelian scheme is finite faithfully flat, so a nonzero section on the special fibre cannot pull back to zero. The isomorphism (T9) is integral and rigidified and introduces no nonunit origin scalar. This proves its claimed origin valuation. In reduced analytic frames the pullback is composition with \(u\mapsto Mu\); normalization gives \(F_1(Mu)\). The chain rule proves (T10), including the exact factor \(M^{-|\alpha|}\).

For odd \(M\), instead \(MT=T\ne0\), and hence \[
 [M]^*L_{135}^{\mathrm{rig}}
 \simeq L_{w_0}^{M^2}\otimes\mathcal P_T
       \not\simeq L_{w_0}^{M^2}.
 \tag{T32}
\] The directly recovered source file \nolinkurl{ReciprocityTheory/ReductiveGroups/ArithmeticGeometry/BirationalGeometry/ToroidalCompactification/TwoWeightEqOneAddRoot.lean}, lines 53--60, defines \nolinkurl{Level} with \nolinkurl{odd : Odd M} at line 56. Accepted DP59/DP62 likewise use odd levels with theta translations in the specified subgroup of \(J[M^2]\). A theta-group translation acts within the section space of its given line, so it cannot erase the nontrivial difference in (T32). This is the stated limitation. It does not imply that allowed odd torsion translates of the singular principal section have zero origin value; they can have nonzero value, and DP59 already supplies jet-spanning results for their prescribed span.

\subsection{What this closes, and what remains}\label{what-this-closes-and-what-remains}

This identifies the exact algebraic principal line for AM-R-004's nonvanishing analytic germ, computes its entire finite origin normalization on the stable good models, and identifies the different principal line specified by the trace-one construction. Thus this explicit analytic model has neither an unspecified characteristic nor an unspecified finite theta-null scalar.

It does not make the value-one Chern/reduced formal frame integral at nonordinary primes. A unit value at 11 does not imply integral Taylor coefficients or identify the selected gain. Conversion from an integral algebraic frame to the reduced CM frame remains a separate local problem, as AM-DO-013--014 explain. At 2 the norm \(w_0w_0^\dagger=2\) is a nonunit: neither (T3) nor (T7) verifies DP49's unit-parameter primitive-kernel hypotheses. The even-level construction (T9) is outside \nolinkurl{Level.odd}. Finally \(F_1(Mu)\) is a pullback theta, distinct from AM-R-004's other family \(\theta(M^2z,M^2\tau)\). No all-polynomial arithmetic/automorphic kernel identity, full \nolinkurl{MainX}, selected \(g_j,e_j\), source period \(H\), or shorter lift chain is asserted.

\subsection{References and exact source locators}\label{references-and-exact-source-locators}

Casalaina-Martin, S., \& Kass, J. L. (2010). \emph{A Riemann singularity theorem for integral curves} (arXiv:0907.0212v2, 1 June 2010). \url{https://arxiv.org/pdf/0907.0212} . The smooth Riemann singularity statement is recalled on printed p.~1. The determinant/Fitting construction is Definition 3.1 and the ensuing paragraphs, pp.~9--10. Proposition 3.3, p.~11, concerns an equal-characteristic test arc; its elementary determinant idea is extended here by the written DVR argument. No mixed-characteristic scope is attributed to that result.

Mumford, D. (1971). Theta characteristics of an algebraic curve. \emph{Annales scientifiques de l'École Normale Supérieure}, 4(2), 181--192. \url{https://www.numdam.org/article/ASENS_1971_4_4_2_181_0.pdf} . Parity and its quadratic law are on pp.~181--182, formula (*); the proof is in §3, pp.~188--190. Only complex/characteristic-not-two parity is used.

National Institute of Standards and Technology. (n.d.). \emph{Digital Library of Mathematical Functions}, §21.2(ii), equations 21.2.5--21.2.7. \url{https://dlmf.nist.gov/21.2} . These fix characteristic and translation conventions. The parity identity follows directly from the displayed defining series.

The Stacks Project Authors. (n.d.). \emph{Determinants of complexes}, Tag 0FJI. \url{https://stacks.math.columbia.edu/tag/0FJI} . The rank-zero two-term determinant element and its quasi-isomorphism invariance are used. \emph{The relative dualizing sheaf}, Tag 0E6N. \url{https://stacks.math.columbia.edu/tag/0E6N} . Relative duality and its base-change compatibility are used.

Accepted intake: \nolinkurl{RETURN_DP_MULT_checkpoint_06_20261121T0420Z/DP_MULT.tex}, DP49 statement and trace-one polarization proof; DP59 proof (TS7)--(TS8), symmetry and rigidified multiplication; DP62, odd-level canonical translates. These are carried at their accepted status, not newly proved. \nolinkurl{Level.odd} was read directly in the recovered pinned source module named above; no Lean execution is claimed.

\clearpage

\section{ALP-LIFT-08: the normalized actual trace-one section and its complete finite scalar content}\label{alp-lift-08-the-normalized-actual-trace-one-section-and-its-complete-finite-scalar-content}

Run ALP\_MACHINE; \nolinkurl{/root/lift_accounting}, GPT-6 Astra Pro; reasoning effort not separately exposed. This uses ALP-LIFT-05's good models and integral Hodge bases and ALP-LIFT-07's integral identification of the trace-one line with the theta characteristic \(2\infty\). It is a theorem about the actual principal section and its first nonzero jet. It does not assert integral reduced-theta coefficients from an integral first jet.

\subsection{Statement}\label{statement-4}

Let \(C:y^2=x^7-1\), \(J=\operatorname{Jac}(C)\), and \(\omega_i=x^{i-1}dx/y\), \(1\leq i\leq3\). Let \(L\) be the symmetric rigidified trace-one line \(L_{w_0}\) of DP49, identified integrally by ALP-LIFT-07 with the rigidified principal theta line for \(2\infty\). In logarithmic differential coordinates \(u_i\), with \(du_i=\omega_i\) at the origin, put \[
 q_\omega(u)=u_1u_3-u_2^2.
 \tag{Q1}
\] There is a unique algebraic section \[
 s_q\in H^0(J,L)
 \tag{Q2}
\] whose initial homogeneous jet in any origin-normalized formal frame is precisely \(q_\omega\). A formal frame is origin-normalized if its value at the identity agrees with the given rigidification. Changing such a frame cannot alter this first nonzero degree-two jet. The section is defined over \(\mathbb Q\) after using the canonical identification with the \(2\infty\) theta line, and in particular over \(E=\mathbb Q(\zeta_7)\) in the trace-one description.

At a finite place after the good-model extensions of ALP-LIFT-05, let \(s_{\mathrm{prim}}\) be any primitive generator of the rank-one module \(H^0(\mathcal J,\mathcal L)\). Define the scalar content by \[
 c_p(s_q)=v_p(a),\qquad s_q=a\,s_{\mathrm{prim}},\qquad v_p(p)=1.
 \tag{Q3}
\] This is independent of the generator. Then \[
 \boxed{c_2(s_q)=8/7,\qquad c_7(s_q)=-1/2,
             \qquad c_p(s_q)=0\quad(p\nmid14).}
 \tag{Q4}
\] The negative entry means a pole relative to the primitive integral section module. These are finite scalar contents of the entire section, not just valuations of two selected Taylor coefficients.

More explicitly, write \(v_i\) for tangent/logarithmic coordinates dual to the primitive local Hodge basis \(\eta_i\) of ALP-LIFT-05. Only their linear terms are needed. At 2 and 7 respectively, \[
 q_\omega=\frac{\alpha^4}{\iota^2}(v_1v_3-v_2^2),
 \qquad
 q_\omega=\frac{\lambda^4}{\beta^2}(v_1v_3-v_2^2)
          =\lambda^{-3}(v_1v_3-v_2^2),
 \tag{Q5}
\] where \(\alpha^7=-4\), \(\iota^2=-1\), \(\lambda=\zeta_7-1\), and \(\beta^2=\lambda^7\).

If \(\sigma_q\) denotes the reduced CM germ of \(s_q\), its initial term is again \(q_\omega\), since a normalized reduced frame has value one at the origin. The comparison of this germ with the classical hyperelliptic sigma, and integrality of its Hurwitz coefficients, are separate assertions proved by the Kummer agent. Likewise the moved germ called \(\Sigma_{\mathrm{tr}}\) in AM-DO-015 is a scalar multiple \(c\sigma_q\); it is not normalized by (Q1). Root's derivative-Thomae calculation determines that separate scalar. These notations must not be interchanged.

\subsection{Full proof}\label{full-proof}

\subsubsection{\texorpdfstring{1. The integral cohomology at the characteristic \(2\infty\)}{1. The integral cohomology at the characteristic 2\textbackslash infty}}\label{the-integral-cohomology-at-the-characteristic-2infty}

Let \(R\) be one of the complete good-model DVRs. Write \(\mathcal B=\mathcal O_{\mathcal C}(2\infty)\), to distinguish the curve line from the theta line on its Jacobian. On each good model, let \(X\) be its degree-two hyperelliptic coordinate. The modules \[
 U=H^0(\mathcal C,\mathcal B),\qquad
 W=H^0(\mathcal C,\omega_{\mathcal C/R}\otimes\mathcal B^{-1})
 \tag{Q6}
\] have integral bases \[
 U=(1,X),\qquad W=(\eta_1,\eta_2),
 \tag{Q7}
\] where \(\eta_2=X\eta_1\) and \(\eta_3=X^2\eta_1\). Indeed \(X\) has its only pole, of order two, at infinity, and the other curve generator has pole order seven. Thus the functions with pole at most two are precisely \(\langle1,X\rangle\) on every fibre. The integral Hodge forms have respective orders four, two, and zero at infinity, as verified in ALP-LIFT-05; vanishing by at least two selects precisely \(\eta_1,\eta_2\). Cohomology and base change, or reduction followed by Nakayama, proves these are the full integral modules. Relative Serre duality gives \[
 H^1(\mathcal C,\mathcal B)=W^*.
 \tag{Q8}
\] Both cohomology modules are consequently free of rank two, including at 2 and 7.

\subsubsection{2. A two-by-two determinant with integral first-order matrix}\label{a-two-by-two-determinant-with-integral-first-order-matrix}

The determinant construction of the relative theta divisor is the same one used in ALP-LIFT-07: on the completed Picard space at the class of \(\mathcal B\), a perfect two-term free complex computes the universal line's cohomology, and its determinant is a local theta equation. This standard construction is recalled in Casalaina-Martin and Kass (2010, Definition 3.1, pp.~9--10), with its determinant invariance in Stacks Project Tag 0FJI.

At the origin, the cohomology in (Q6)--(Q8) is free. The specialized complex therefore splits into these two cohomology modules, with zero differential, plus an acyclic free summand. To see that this splitting is integral, the images, kernels and cohomology are free over the DVR, so every short exact sequence with free quotient splits. Lifting the chosen decompositions to the completed base and eliminating the invertible acyclic block by row/column operations leaves a two-by-two matrix \(d\), all of whose entries belong to the origin ideal. The eliminated block has unit determinant. Thus the initial quadratic theta equation is, up to an integral unit, the determinant of the linear part \(d_1\).

We give the first-order identification explicitly. A tangent vector to the Picard scheme is a class \(\xi\in H^1(\mathcal O_C)\). If transition functions \(g_{ij}\) for \(\mathcal B\) are deformed to \(g_{ij}(1+\epsilon\xi_{ij})\), then a section \(s\) lifts to first order precisely when the Čech cocycle \(\xi_{ij}s_j\) is a coboundary (a simultaneous overall sign has no effect on the two-by-two determinant). The first differential on cohomology is therefore cup product \[
 U\longrightarrow H^1(\mathcal B),\qquad s\longmapsto\xi\smile s.
 \tag{Q9}
\] Under the Serre duality (Q8), its matrix entries are the pairings of \(\xi\) with the products \(s_a t_b\in H^0(\omega)\), for the bases in (Q7). Those products are exactly \[
 \begin{pmatrix}\eta_1&\eta_2\\\eta_2&\eta_3\end{pmatrix}.
 \tag{Q10}
\] Since \(H^1(\mathcal O_C)\) is integrally dual to \(H^0(\omega)\), in the tangent coordinates \(v_1,v_2,v_3\) dual to these forms the first-order matrix and its determinant are \[
 d_1=\begin{pmatrix}v_1&v_2\\v_2&v_3\end{pmatrix},\qquad
 \det d_1=q_\eta=v_1v_3-v_2^2.
 \tag{Q11}
\] This polynomial has an integral-unit coefficient at every place, including residue characteristic two. In particular it is nonzero there; no division by 2 or characteristic-not-two theorem was used. The argument is the usual determinant/cup-product mechanism of the Riemann--Kempf tangent cone, with the present integral case established directly rather than inferred from a complex tangent-cone statement.

The origin fibre of the determinant theta line identifies with \(\det U^{-1}\otimes\det W^{-1}\). The bases in (Q7) give it a primitive generator. Passing to its rigidification and using the integral isomorphism with \(\mathcal L_{w_0}\) from ALP-LIFT-07, one sees that the degree-two jet map on the actual principal section module has image exactly \[
 R\,q_\eta\subset\operatorname{Sym}^2 e^*\Omega^1_{\mathcal J/R}.
 \tag{Q12}
\] Here a remaining integral unit in the determinant presentation can be absorbed in the primitive section generator. The principal theta section module is free of rank one and formation commutes with base change, by the standard principal-polarization cohomology theorem used in DP59. Therefore (Q12) is a saturated rank-one image; it determines the scalar content of the \emph{entire section} from its quadratic jet.

\subsubsection{3. Algebraic normalization and finite contents}\label{algebraic-normalization-and-finite-contents}

Over the original characteristic-zero field use \(U=(1,x)\), \(W=(\omega_1,\omega_2)\). The same calculation gives the jet line spanned by (Q1). Its generator is defined over \(\mathbb Q\), as are the curve, infinity, the rigidified \(2\infty\) theta line, and this quadratic tensor. The nonzero degree-two jet map from a one-dimensional section space has a unique inverse image of \(q_\omega\). This proves the existence, uniqueness and field assertion in (Q2). A change of origin-normalized formal frame multiplies a local representative by a power series with constant term one, and cannot change its first nonzero degree-two term.

At 2, ALP-LIFT-05 gives \(\omega_i=\alpha^i\iota^{-1}\eta_i\). Thus at the tangent level \(u_i=\alpha^i\iota^{-1}v_i\), and \[
 q_\omega=\alpha^4\iota^{-2}q_\eta,
 \qquad v_2(\alpha^4\iota^{-2})=8/7.
 \tag{Q13}
\] Equation (Q12) proves \(c_2(s_q)=8/7\).

At 7, ALP-LIFT-05 gives \[
 \begin{aligned}
 u_1&=(\lambda/\beta)v_1,\\
 u_2&=(\lambda/\beta)(v_1+\lambda v_2),\\
 u_3&=(\lambda/\beta)(v_1+2\lambda v_2+\lambda^2v_3).
 \end{aligned}
 \tag{Q14}
\] Substitution cancels the \(v_1^2\) and \(v_1v_2\) terms and gives \[
 q_\omega=(\lambda^4/\beta^2)q_\eta
           =\lambda^{-3}q_\eta,
 \qquad v_7(\lambda^{-3})=-1/2.
 \tag{Q15}
\] Again (Q12) proves the stated content. At every other prime the original basis is already primitive, so (Q12) gives content zero. This proves (Q4)--(Q5).

There is also an independent determinant-frame check of the two scalings. At 2, relative to (Q7), the determinant valuations of \((1,x)\) and \((\omega_1,\omega_2)\) are \(2/7\) and \(6/7\). At 7 they are respectively \(1/6\) and \(v_7(\lambda^3/\beta^2)=-2/3\). Their sums are exactly \(8/7\) and \(-1/2\). These sums agree with (Q13) and (Q15), as required by the determinant line \(\det U^{-1}\otimes\det W^{-1}\).

\subsection{Consequence at the actual prime 11, and scope}\label{consequence-at-the-actual-prime-11-and-scope}

The section \(s_q\) is now an explicitly normalized generator of the actual integral principal section module at 11. It can be used as the fixed primitive base section in accepted DP59 and DP62. This removes an unspecified scalar from their application to this Jacobian. Their permitted subgroup, square level, theta-group lift and selected-jet conditions are unchanged. Their integral-span statements are cited as accepted results, not re-proved here.

The conclusion does not infer full reduced-theta integrality from the quadratic term. The scalar content at 7 is measured in the primitive stable Hodge/line model. It is compatible with separate Hurwitz integrality in the original \(\omega\)-coordinates, because those coordinates are not primitive at 7. The dyadic content also says nothing by itself about the norm-two, nonunit trace-one kernel parameter. Finally a principal section and its allowed translates still require an actual theta-lift row/column identification and all global costs before any integer chain improvement follows.

\subsection{Literature and status}\label{literature-and-status}

The determinant and relative-duality inputs are the same directly checked primary sources as in ALP-LIFT-07: Casalaina-Martin and Kass (2010), arXiv:0907.0212v2, Definition 3.1, pp.~9--10; Stacks Project Tags 0FJI and 0E6N. The Čech first-order argument and the integral splitting argument needed here are written in full above.

The classical tangent-cone theorem is due to Kempf, G. (1973), \emph{On the geometry of a theorem of Riemann}, Annals of Mathematics, 98(1), 178--185, \url{https://doi.org/10.2307/1970910} . The original publication metadata was checked on the Annals site; its paywalled full text was not inspected. The related cup-product matrix description was read in Hitching and Hoff, arXiv:1610.03456v2, §3.1, p.~11. Their stated Petri-injective setup is not asserted for this hyperelliptic curve; our integral proof does not use that hypothesis or their reconstruction theorem. No whole-proof check of those literature results is claimed.

\clearpage

\section{AM-R-006: the exact scalar joining the translated theta and the normalized trace-one section}\label{am-r-006-the-exact-scalar-joining-the-translated-theta-and-the-normalized-trace-one-section}

Author: \nolinkurl{/root}, GPT-6 Astra Pro; reasoning effort not externally exposed; run ALP\_MACHINE, 21 November 2026. This finding uses AM-R-004, AM-DO-015--016 and ALP-LIFT-07--08. It removes the remaining scalar between their two actual sections of the same principal theta line. It is not an identification of an iterated automorphic lift image.

\subsection{Statement}\label{statement-5}

Use the period coordinates of AM-R-004, with \(u=Pz\), \(Y=\operatorname{Im}\tau\), and \(Q(z)=\pi z^tY^{-1}z/2\). Put \[
 \theta_0=\theta[0;0](0;\tau),\qquad
 \theta_K(z)=\theta[\tfrac12(1,1,1);\tfrac12(1,0,1)](z;\tau).
 \tag{R6.1}
\] By AM-DO-015 the exact normalized Mumford translate is \[
 \Sigma_{\rm tr}(u)=-i\,e^{Q(P^{-1}u)}
                   \frac{\theta_K(P^{-1}u)}{\theta_0}.
 \tag{R6.2}
\] Let \(\sigma_q\) be the reduced germ of the actual trace-one section \(s_q\) of ALP-LIFT-08, whose first nonzero homogeneous term is \(q(u)=u_1u_3-u_2^2\). Define the nonzero scalar by the exact convergent analytic expression \[
 c=-\frac{i}{\theta_0}
     \left.\partial_{u_1}\partial_{u_3}
         \theta_K(P^{-1}u)\right|_{u=0}.
 \tag{R6.3}
\] This fixes its phase, not merely an unspecified root. Then \[
 \Sigma_{\rm tr}=c\,\sigma_q.
 \tag{R6.4}
\] For \(S=\{1,3,5\}\), \(S^c=\{0,2,4,6\}\), set \[
 R_S=\prod_{a\in S}\prod_{b\in S^c}(\zeta_7^a-\zeta_7^b)
 =49(-\zeta_7^5-\zeta_7^4-2\zeta_7^2+\zeta_7-2).
 \tag{R6.5}
\] The scalar is algebraic and satisfies the exact phase-independent identity \[
 c^8=\frac{R_S^2}{4^8},\qquad
 |c|=\frac{|R_S|^{1/4}}4.
 \tag{R6.6}
\] In normalized valuations \(v_p(p)=1\), at every place above the indicated rational prime, \[
 v_2(c)=-2,\qquad v_7(c)=\frac12,\qquad
 v_p(c)=0\quad(p\ne2,7).
 \tag{R6.7}
\] The primitive integral scalar contents of the translated section are therefore \[
 c_2(\Sigma_{\rm tr})=-\frac67,\qquad
 c_p(\Sigma_{\rm tr})=0\quad(p\text{ odd}).
 \tag{R6.8}
\] Here content is relative to the integral rigidified trace-one line of ALP-LIFT-08. Equation (R6.8) concerns the entire algebraic section, not the coefficient integrality of a Chern formal frame.

The following explicit global estimate now holds for the actual quadratic-normalized trace-one germ: \[
 |\sigma_q(u)|<\frac{6125}{16}
       \exp\!\left(\frac\pi2\sum_{j=1}^3\frac{|u_j|^2}{G_j}\right).
 \tag{R6.9}
\] Consequently, for every multiindex \(\alpha\), \[
 |\partial^\alpha\sigma_q(0)|
 <\frac{6125}{16}\,\alpha!
   \prod_{j:\alpha_j>0}
      \left(\frac{\pi e}{\alpha_jG_j}\right)^{\alpha_j/2}.
 \tag{R6.10}
\] For any positive integer \(M\), substitution of \(Mu\) multiplies the right side before level normalization by \(M^{|\alpha|}\). Thus the source factor \(M^{-|\alpha|}\) cancels this increase exactly; for the actual arithmetic \nolinkurl{Level} one retains its oddness condition. The section \(\sigma_q(Mu)\) is in the correct trace-one line for every such odd \(M\).

\subsection{Proof}\label{proof}

\subsubsection{The derivative Thomae comparison with all differential factors retained}\label{the-derivative-thomae-comparison-with-all-differential-factors-retained}

We use the derivative Thomae formula as a literature theorem, rather than reprove it. Bernatska (2020), arXiv v3, Corollary 2, equations (34)--(37), printed p.~18, and Example 4, equations (38)--(39), printed p.~19, identify the Hessian of the unique singular even theta characteristic in genus three. Her differentials are \[
 d v_1=-\frac{x^2dx}{2y},\qquad
 d v_3=-\frac{x dx}{2y},\qquad
 d v_5=-\frac{dx}{2y};
\] this follows from her equation (2), printed p.~2, with \(f=-y^2+x^7-1\). Hence \[
 (v_1,v_3,v_5)=-\tfrac12(u_3,u_2,u_1).
 \tag{R6.11}
\] The cited Hessian formula says that the degree-two term of \(\theta_K\), in these variables, is \(C_3(v_1v_5-v_3^2)\). The derivative normalization of \(C_3\) is stated immediately after (39). In the original \(u\) variables the term is therefore exactly \[
 \frac{C_3}{4}(u_1u_3-u_2^2).
 \tag{R6.12}
\] An eighth-root phase in the homology/characteristic convention is included in the actual derivative value \(C_3\); equation (R6.3) selects it in our fixed period convention. The quadratic exponential in (R6.2) has constant term one, so it cannot alter (R6.12).

Both \(\Sigma_{\rm tr}\) and \(\sigma_q\) are nonzero sections of the same rigidified principal line, represented in its same reduced frame. Its section space has dimension one. Comparing their first nonzero jets proves (R6.4) with \(c=-iC_3/(4\theta_0)\). This argument does not assume a global algebraic trivialization of that ample line.

Bernatska's Remark 3 and equation (30), printed p.~16, give \(C_3\) as the product of the square root of the period determinant factor and the fourth root of the Vandermonde in all seven finite branch points. The first Thomae formula for the nonsingular characteristic belonging to \(S\) gives the same period determinant factor times the fourth root of the two within-part Vandermondes. In our period convention that nonsingular characteristic is precisely the zero characteristic, by ALP-LIFT-07. Dividing the two formulas cancels the period determinant. The result is \[
 \frac{C_3}{\theta_0}=\epsilon R_S^{1/4},\qquad \epsilon^8=1,
 \tag{R6.13}
\] with the ordering sign absorbed in the root convention. This is also the displayed consequence of first Thomae immediately after (30). Taking absolute values and eighth powers proves (R6.6). In particular the scalar defined analytically in (R6.3) is algebraic. There is no need to select an arbitrary fourth root in place of that exact analytic definition. The cyclotomic reduction in (R6.5) is ordinary polynomial multiplication modulo \(1+T+\cdots+T^6\), and is independently reproduced by the accompanying script.

\subsubsection{Finite valuations and integral section content}\label{finite-valuations-and-integral-section-content}

For distinct seventh roots of unity, their difference is a unit away from 7 and has normalized valuation \(1/6\) at 7. For example \(\zeta^a-\zeta^b=\zeta^b(\zeta^{a-b}-1)\); the quotient of \(\zeta^{a-b}-1\) by \(\zeta-1\) is a unit and \(v_7(\zeta-1)=1/6\). There are twelve factors in (R6.5), so \(v_7(R_S)=2\), with zero valuation elsewhere. Applying valuations to (R6.6) yields \[
 8v_p(c)=2v_p(R_S)-8v_p(4),
\] which proves (R6.7) at every place after any further finite extension. Roots of unity have valuation zero, including at 2, so the phase convention causes no exception.

Scalar contents add under multiplication of a section by a scalar. ALP-LIFT-08 gives content \(8/7\) at 2, \(-1/2\) at 7, and zero elsewhere for \(s_q\). Adding (R6.7) proves (R6.8). This also checks compatibility with the finite normalization of the original value-one section in ALP-LIFT-07: the translated section has exactly the same dyadic content \(-6/7\), and is primitive at 11 and at every other odd prime on the stated good models. The equality of these contents is a computed conclusion; it was not assumed from integrality of a two-torsion translation in characteristic two.

\subsubsection{The archimedean estimate and derivatives}\label{the-archimedean-estimate-and-derivatives}

Every difference of two distinct seventh roots of unity has absolute value at least \(2\sin(\pi/7)>4/5\). The rational trigonometric certificate in AM-R-004 proves this elementary inequality. Thus \[
 |R_S|>(4/5)^{12},\qquad |c|>\frac14(4/5)^3=\frac{16}{125}.
 \tag{R6.14}
\] AM-DO-015's normalized Mumford translation preserves the AM-R-004 Gaussian estimate, including its constant 49. Divide that estimate by (R6.14), and use \(49\cdot125/16=6125/16\), to obtain (R6.9). This bound is valid on all of \(\mathbb C^3\), not only a compact set or a formal neighborhood.

Apply the multivariable Cauchy estimate on a polydisc of radii \(r_j>0\). For \(\alpha_j>0\), minimize \(e^{\pi r_j^2/(2G_j)}r_j^{-\alpha_j}\) at \(r_j^2=\alpha_jG_j/\pi\). For zero entries let \(r_j\) tend to zero after applying Cauchy. This gives (R6.10). Composition with \([M]\) acts on the universal cover by \(u\mapsto Mu\), and the chain rule gives the final level-normalization assertion.

\subsection{Consequence and limits for the main question}\label{consequence-and-limits-for-the-main-question}

The source-line section can now be used with a fixed algebraic quadratic normalization, exact finite scalar content, and an explicit archimedean Cauchy bound. Its relation to the literal finite Schwartz input of AM-DO-016 is \[
 e^{Q(z)}\Theta(\varphi_t;z,\tau)
       =\theta_0 c\,\sigma_q(Pz).
 \tag{R6.15}
\] The exact scalar is retained on the right. This closes the scalar comparison for that actual CM fibre. It neither makes the reduced formal frame integral nor computes the automorphic quotient volume, the terminal mass, a selected lift's actual gain/error, or a common all-polynomial \nolinkurl{MainX} identity. Those are separate questions.

In the Picard and biextension formulation this is an equality of sections of the same rigidified line after the two-torsion change of characteristic; the scalar valuations measure its integral section lattice. In determinant-of-cohomology terms, the chosen quadratic jet is the cup-product determinant of ALP-LIFT-08. The comparison concerns the metric and integral lattice of a line, not an identity of higher cohomological weights. See the exact FGA/SGA7/cote54 locators in \nolinkurl{GROTHENDIECK_DICTIONARY.md}.

\subsection{Primary reference and computation}\label{primary-reference-and-computation}

Bernatska, J. (2020). General derivative Thomae formula for singular half-periods. \emph{Letters in Mathematical Physics, 110}, 2983--3014. \url{https://doi.org/10.1007/s11005-020-01315-8} . The directly inspected version is \url{https://arxiv.org/pdf/1904.09333v3} , equation (2), Remark 3/equation (30), Corollary 2/equations (34)--(37), and Example 4/equations (38)--(39), printed pp.~2, 16, 18--19.

\nolinkurl{scripts/trace_one_theta_scalar.py} performs exact cyclotomic multiplication and an independent 60-digit theta-series diagnostic of the Hessian shape, scalar modulus and eighth-power identity. The latter is explicitly numerical, not an interval or formal proof. The theorem above uses the exact literature formulas and elementary algebra, independently of the numerical diagnostic.

\clearpage

\section{ALP\_MACHINE --- primitive CM reduced theta and the integral-frame boundary}\label{alp_machine-primitive-cm-reduced-theta-and-the-integral-frame-boundary}

Author: GPT-6 Astra Pro, \nolinkurl{/root/definite_oscillator}, inherited effort not separately exposed, 21 November 2026 UTC. New consequences below are labelled AM-DO-011--016. Root's AM-R-004 and the supplied THETA3/DP\_MULT results retain their own proof statuses. This note does not instantiate \nolinkurl{PinX}, \nolinkurl{ClsK.Core}, a primitive trace-one kernel, or the requested lift family.

\subsection{Literature actually read}\label{literature-actually-read}

Bannai, K., \& Kobayashi, S. (2010). Algebraic theta functions and the p-adic interpolation of Eisenstein--Kronecker numbers. \emph{Duke Mathematical Journal, 153}(2), 229--295. \url{https://doi.org/10.1215/00127094-2010-024}. The mathematical text read in this run is arXiv:math/0610163v4 (11 December 2007), \url{https://arxiv.org/pdf/math/0610163}, specifically Definition1.8 pp10--11; §2.1 and Proposition2.1 pp19--21; Theorem2.9 pp25--26; §2.3, condition(28), Lemma2.13 and Propositions2.14--2.15 pp27--29. Published bibliographic metadata was checked on Bannai's own publication page.

The used conclusions are: a reduced theta germ on a CM abelian variety has algebraic coefficients in algebraic eigen-differential coordinates, if its divisor and normalized leading value are algebraic and its line is admissible; prime-to-p torsion translations have an algebraic comparison; formal-parameter integrality requires the stated CM ordinarity condition and an arithmetic unit normalization. Those are cited theorems, not new theorems rederived in this run.

The supplied THETA3 C11 already provides an actual reduced-theta counterexample to automatic integral formal coefficients, and C19 treats integral de Rham CM splittings and derivative estimates. They are carried, not re-proved. In particular C11's y²=x³−x example at good supersingular p=7 has normalized reduced theta coefficient1411/2016 at formal degree9, of valuation−1. This is a concrete warning against promoting a generic algebraic/Chern comparison to an integral one.

\subsection{AM-DO-011. The exact R004 completion is a reduced theta, with algebraic coefficients}\label{am-do-011.-the-exact-r004-completion-is-a-reduced-theta-with-algebraic-coefficients}

Carry AM-R-004's actual Jacobian J of y²=x⁷−1, E=Q(ζ7), Φ=\{1,2,3\}, algebraic differential coordinates u\_i=∫ω\_i, u=Pz, normalized period matrix τ, and Y=Imτ. For κ≥1 write

\begin{verbatim}
θκ(z)=Σ_(n∈Z³) exp(πiκ nᵗτn+2πiκ nᵗz),
Fκ(u)=exp(πκ zᵗY⁻¹z/2) θκ(z)/θκ(0).
\end{verbatim}

The carried quantitative result gives θκ(0)≠0. Define

\begin{verbatim}
Hκ(z,w)=κ zᵗY⁻¹ conjugate(w),
ακ(a+τb)=(−1)^(κ aᵗb),  a,b∈Z³.
\end{verbatim}

Then Fκ is exactly the reduced theta function for the Appell--Humbert data(Hκ,ακ), normalized to value1 at0. It is not merely proportional to an unspecified reduced theta after an omitted quadratic exponential.

\subsubsection{Full calculation of the multiplier}\label{full-calculation-of-the-multiplier}

Direct reindexing gives

\begin{verbatim}
θκ(z+a+τb)/θκ(z)
  =exp(−πiκ bᵗτb−2πiκ bᵗz).
\end{verbatim}

For λ=a+τb, multiplication by the completion factor changes the linear coefficient in z to

\begin{verbatim}
πκ zᵗ(Y⁻¹λ−2ib)=πκ zᵗY⁻¹ conjugate(λ)=πHκ(z,λ).
\end{verbatim}

The difference between the resulting constant term and πHκ(λ,λ)/2 is

\begin{verbatim}
(πκ/2)λᵗY⁻¹(λ−conjugate(λ))−πiκ bᵗτb
   =πiκ λᵗb−πiκ bᵗτb=πiκ aᵗb.
\end{verbatim}

Therefore

\begin{verbatim}
Fκ(u+Pλ)=ακ(λ) exp(πHκ(z,λ)+πHκ(λ,λ)/2) Fκ(u),
\end{verbatim}

which is precisely Bannai--Kobayashi Definition1.8. Since ακ²=1, its semicharacter meets their finite-order hypothesis. The completion's constant is1 at the origin, so its value normalization is the stated one. This proves the assertion. ∎

\subsubsection{Admissibility and algebraicity}\label{admissibility-and-algebraicity}

In the first-argument-linear Appell--Humbert convention used here, Im Hκ is minus κ times the R004 alternating form

\begin{verbatim}
E_R4(x,y)=Tr_(E/Q)((ζ−ζ⁻¹)x conjugate(y)/7).
\end{verbatim}

The sign can be checked on the symplectic periods: Im Hκ(a\_i,b\_j)=−κδ\_ij, while E\_R4(a\_i,b\_j)=δ\_ij. The sign convention has no effect on admissibility. For every a∈E, direct substitution gives E\_R4(ax,y)=E\_R4(x,conjugate(a)y), and therefore the same equality for Im Hκ. This verifies Φ-admissibility in the exact sense of Bannai--Kobayashi §2.1. Its algebraic invariant differentials are the specified ω\_i=x\^{}(i−1)dx/y, and the CM action x↦ζx is defined over E.

For κ=1 the associated symmetric theta line and its divisor are defined over E: AM-DO-012 below identifies every possible symmetric representative with one of the E-rational hyperelliptic theta characteristics. With F1(0)=1, Proposition2.1 applies directly and gives

\begin{verbatim}
F1(u1,u2,u3) ∈ E[[u1,u2,u3]].               (DO11.1)
\end{verbatim}

Here the conclusion concerns the coefficients in the algebraic differential coordinates u, including the ordinary factorial denominators inherent in a Taylor expansion. It does not identify the theta-null period θ1(0) itself as algebraic.

For κ\textgreater1 one must also descend the selected section line, since an arbitrary complex section of L\^{}κ need not be algebraic up to scalar. In this case the descent is explicit. Let Hκ⊂J{[}κ{]} be the subgroup represented by (κ⁻¹Z³)/Z³. Translation z↦z+a/κ fixes θκ. These translations have commuting finite-order lifts to the theta group of L\^{}κ, and the corresponding joint character space in H⁰(J,L\^{}κ) has dimension1. For example the usual κ³ theta-series basis has joint eigencharacters indexed by j∈(Z/κZ)³, and θκ is its j=0 vector. Choose the matching algebraic lifts after a finite torsion-field extension. Such lifts are algebraic: starting from any algebraic isomorphism for a torsion translation, requiring its finite power to equal1 only adjoins a root of an algebraic number; the remaining exact finite commutator relations are the same finite theta-group relations. The chosen simultaneous eigenspace is the kernel of algebraic linear operators and hence is an algebraic line. Its divisor is therefore defined over a finite number field Kκ. Proposition2.1 gives Fκ∈Kκ{[}{[}u{]}{]}. This argument gives a finite descent for each κ, not a level-uniform field-degree or coefficient-height bound.

\subsection{AM-DO-012. A nonzero even theta value is a unit at every good prime for this Jacobian}\label{am-do-012.-a-nonzero-even-theta-value-is-a-unit-at-every-good-prime-for-this-jacobian}

Work over R=Z{[}ζ7,1/14{]}, with the smooth projective hyperelliptic curve C:y²=x⁷−1. Let W\_j=(ζ\^{}j,0), j=0,\ldots,6, and let ∞ be its unique point at infinity. The branch sections are pairwise disjoint over R because only7 ramifies in their differences, and2 is inverted. The standard relative hyperelliptic model is smooth and proper over R. For a subset S of the seven finite branch points with s=\textbar S\textbar≤3, put

\begin{verbatim}
D_S=Σ_(j∈S) W_j+(2−s)∞.
\end{verbatim}

Every geometric fiber has the following exact theta-characteristic classification and cohomology:

\begin{verbatim}
S=∅:             h⁰(O(D_S))=2,
|S|=1 or2:       h⁰(O(D_S))=1,
|S|=3:           h⁰(O(D_S))=0.              (DO12.1)
\end{verbatim}

There are1+7+21+35=64 such characteristics, and they exhaust all of them. In particular the35 non-effective characteristics remain non-effective in every residue characteristic other than2 and7.

\subsubsection{Full proof of the classification and cohomology calculation}\label{full-proof-of-the-classification-and-cohomology-calculation}

The differential dx/y has divisor4∞, so K\_C∼4∞. Also div(x−ζ\^{}j)=2W\_j−2∞. Hence2D\_S∼4∞ and every D\_S is a theta characteristic.

For distinct S,T with cardinalities at most3, the difference of their associated two-torsion points is represented, using2(W\_j−∞)∼0, by

\begin{verbatim}
Σ_(j∈S△T) W_j−|S△T|∞.
\end{verbatim}

Suppose this is principal for a nonempty U=S△T of size at most6. Its divisor is invariant under the hyperelliptic involution ι. If f has that divisor, then ι(f)/f is constant and its square is1. Since the characteristic differs from2, either f∈k(x) or f=y r(x) with r∈k(x). A rational function in x has even order at every branch point, contradicting the required order1 at j∈U. A function y r(x) has odd order at every one of the seven finite branch points, contradicting the required order0 at a point outside U. Thus all64 points are distinct. A genus-three Jacobian has64 geometric two-torsion points in characteristic different from2, and the theta characteristics form their torsor by addition to2∞; the list is complete.

Now compute H⁰(O(D\_S)) directly. The space is stable under ι, so every permitted function decomposes into its even and odd parts f=u(x)+y v(x), each still permitted. The even part has no finite poles, since a pole at W\_j would have even order at least2 while D\_S permits only1. Thus u is a polynomial. Its pole order2deg u at∞ must be at most2−s. This gives dimension2 when s=0, dimension1 when s=1 or2, and dimension0 when s=3.

For the odd part all finite poles must occur among S, and each possible denominator x−ζ\^{}j occurs to exponent at most1. Therefore

\begin{verbatim}
y v(x)=y Q(x)/∏_(j∈S)(x−ζ^j)
\end{verbatim}

for a polynomial Q, unless the odd part is zero. Its order at∞ is2s−7−2deg Q. To be a section of O(D\_S), this order must be at least s−2. It would follow that2deg Q≤s−5, impossible for s≤3. The odd part is therefore zero. This proves(DO12.1) over every indicated geometric fiber. ∎

\subsubsection{The actual algebraic unit rigidification}\label{the-actual-algebraic-unit-rigidification}

For \textbar S\textbar=3, Riemann--Roch gives h¹(O(D\_S))=0 too, because deg D\_S=g−1=2 and h⁰=0. There is an actual relative theta section on the Jacobian attached to D\_S, and its value at the identity is a unit over R. One can see the unit without choosing a transcendental theta null.

Let p:C×\emph{R J→J, let P be the rigidified Poincaré line on C×J, and let M=P⊗pr\_C\^{}\emph{O(D\_S). Take the effective divisor B=6∞. The line M(B) has fiber degree8, hence H¹=0 on every fiber. Thus p\_}M(B) and p}*(M(B)\textbar\_B) are locally free of the same rank6, and the determinant of their restriction map is the usual determinant-of-cohomology theta section. At the identity, its kernel and cokernel are H⁰(O(D\_S)) and H¹(O(D\_S)), both zero on every fiber. The restriction map there is an isomorphism. Its determinant is therefore a unit and supplies the canonical rigidification of the theta line at0 in which the section has value1. This constructs the integral section and its unit value; no assertion that a Chern formal trivialization is integral has been used.

The symmetric principal theta line in R004 must correspond to one of these theta characteristics. Its theta value at0 is nonzero, by R004's exact estimate, so its characteristic has h⁰=0 and therefore belongs to the35 triples. It can consequently be modeled by the corresponding D\_S over R with this unit rigidification. An explicit S is unnecessary for the conclusion, although it can be recovered from the finite level-two theta characteristic. The isomorphism to the complex R004 representative is fixed by the principal polarization and the rigidification; the resulting normalized reduced germ remains exactly F1.

\subsubsection{Explicit branch half-period data in the R004 basis}\label{explicit-branch-half-period-data-in-the-r004-basis}

This calculation makes the finite level-two comparison concrete. Use root's basis a0=γ0,a1=γ2,a2=γ4 and b0=γ1+γ3+γ5,b1=γ3+γ5,b2=γ5. Then the branch points have the following half-periods, with a point denoted(ε;δ) when its z coordinate is(ε+τδ)/2:

\begin{longtable}[]{@{}lll@{}}
\toprule\noalign{}
Point & ε & δ \\
\midrule\noalign{}
\endhead
\bottomrule\noalign{}
\endlastfoot
W0 & 000 & 100 \\
W1 & 100 & 100 \\
W2 & 100 & 010 \\
W3 & 110 & 010 \\
W4 & 110 & 001 \\
W5 & 111 & 001 \\
W6 & 111 & 000 \\
∞ & 000 & 000 \\
\end{longtable}

To prove the first row, integrate on the real path∞→W0 with y=√(x⁷−1):

\begin{verbatim}
∫∞^W0 ω_i=−B(1/2−i/7,1/2)/7.
\end{verbatim}

The gamma reflection identities give

\begin{verbatim}
B(1/2−t,1/2)/B(t,1/2)=tan(πt).
\end{verbatim}

Divide by root's c\_i=−2i\_C B(i/7,1/2)(ζ\^{}i−1)/7. The quotient is−1/{[}2(1+ζ\^{}i){]}=(ζ\^{}i+ζ\^{}\{3i\}+ζ\^{}\{5i\})/2, using(1+ζ)\^{}(−1)=−ζ−ζ³−ζ⁵. Thus AJ(W0−∞)=b0/2. The identity AJ(Wj−W0)=½Σ\_(k\textless j)γ\_k follows by taking half the difference of the two sheet-symmetric branch paths λ\_j−λ\_0. Expanding γ0=a0, γ1=b0−b1, γ2=a1, γ3=b1−b2, γ4=a2, γ5=b2 modulo2 gives every row. This proves the table. No unverified Riemann-constant convention is used to name the triple S of θ{[}0;0{]}.

\subsection{AM-DO-013. An actual pullback family with fixed normalization and ordinary integrality}\label{am-do-013.-an-actual-pullback-family-with-fixed-normalization-and-ordinary-integrality}

Let L,s be the integral symmetric theta line and unit-rigidified section from AM-DO-012 matching F1. For every M≥1, the section {[}M{]}\^{}\emph{s lives in {[}M{]}\^{}}L≃L\^{}(M²), using the rigidified symmetric theorem-of-the-cube isomorphism. Its exact normalized reduced germ is

\begin{verbatim}
G_M(u)=F1(Mu).                              (DO13.1)
\end{verbatim}

It has value1 at0, lies in E{[}{[}u{]}{]}, and satisfies

\begin{verbatim}
M^(−|α|) ∂^α G_M(0)=∂^α F1(0).             (DO13.2)
\end{verbatim}

From AM-R-004's bound at κ=1 one therefore obtains

\begin{verbatim}
|G_M(u)|<49 exp((πM²/2)Σ_i |u_i|²/G_i),
\end{verbatim}

and exactly the same level-normalized derivative bound as R4.11. This is a different actual section family from θ\_(M²)(z); the equality(DO13.2) states its chosen normalization and makes no claim that it equals a primitive trace-one kernel.

Proof. Pullback along{[}M{]} acts on the universal cover by u↦Mu. Its Appell--Humbert Hermitian form is M²H1, and its semicharacter is α1(Mλ)=α1(λ)\^{}(M²), so the reduced multiplier is that of the specified symmetric tensor power. The germ is therefore the displayed pullback, with no additional nonconstant factor; the value normalization fixes the constant. The chain rule proves(DO13.2), and substitution in the carried κ=1 estimate proves the bound. ∎

At a good prime p≠2,7 satisfying Bannai--Kobayashi condition(28), let λ(t) be the formal logarithm in an integral parameter system. The section's arithmetic divisor avoids0, and its normalized value is1. Proposition2.14, applied with local defining function f=1, gives

\begin{verbatim}
F1(λ(t)) ∈ O_v[[t1,t2,t3]]^×.              (DO13.3)
\end{verbatim}

For all M, the formal multiplication map\href{t}{M} has integral coefficients, including when p divides M. Since λ(\href{t}{M})=Mλ(t), composition gives

\begin{verbatim}
G_M(λ(t))=F1(λ([M](t)))∈O_v[[t]]^×.
\end{verbatim}

Thus this family has exact uniform formal integrality at the stated ordinary primes. Prime-to-p torsion translations have the further explicit algebraic/integral comparison of Bannai--Kobayashi Theorem2.9 and Proposition2.15; their theta-group lifts, local divisor functions and torsion field are still part of the datum.

For the primitive type Φ=\{1,2,3\} of Q(ζ7), condition(28) at an unramified rational prime is equivalent to p≡1(mod7). Indeed let D=

⊂(Z/7Z)\^{}× be its decomposition group. The ideals occurring in Φ(p) and conjugate(Φ)(p) are disjoint exactly when ΦD and conjugate(Φ)D are disjoint. Since their union is the whole six-element group, this is equivalent to Φ being a union of D-cosets, or DΦ=Φ. Its stabilizer is trivial: multiplication by2 sends3 to6; by3 sends2 to6; by4 sends1 to4; by5 sends1 to5; by6 sends1 to6. Hence D=\{1\}, equivalently p≡1(mod7). This calculation concerns the exact hypothesis of the cited theorem; it does not assert non-integrality at every other prime.

\subsection{AM-DO-014. What the CM Hodge splitting does and does not settle}\label{am-do-014.-what-the-cm-hodge-splitting-does-and-does-not-settle}

Over a number field containing E and the endomorphisms, H¹\_dR(J) is the sum of six one-dimensional E-eigenlines. Its Hodge subspace is the sum of the three Φ-eigenlines. The sum of the other three eigenlines is the unique E-stable complement. Indeed any other complement is the graph of an E-linear map from that sum to the Hodge subspace, and such a map is zero because their eigencharacters are disjoint. After base change to C the anti-holomorphic Hodge summand has precisely those complementary characters, so it is this same complement.

At a good prime not dividing7, after a finite base change splitting O\_E⊗O\_v, the same idempotents are integral, because the CM algebra is étale there. The Hodge subbundle is saturated in the de Rham lattice. Each of its eigenparts has either full rank1 or rank0, so saturation makes it respectively the entire eigenline lattice or zero. Thus the complementary sum gives an integral E-stable module splitting as well. These assertions prove the proposed comparison of the linear Hodge splittings.

They do not prove that the generic formal horizontal section of a connection, the formal exponential of a Lie splitting, or a chosen Chern/reduced theta frame is an integral formal object. The supplied C11 and C19 already establish this distinction for actual CM elliptic theta functions. It is therefore not enough to invoke uniqueness of the E-linear Hodge complement to discharge all integral-frame hypotheses of the mixed-type arithmetic constructor. The generic algebraicity needed here is independently justified by the exact reduced multiplier and Bannai--Kobayashi Proposition2.1; the ordinary integrality is the narrower result(DO13.3).

For the requested bridge to \nolinkurl{I.J.V.thMoved} and \nolinkurl{PinX}, the remaining tasks are exact and separate: identify the actual primitive/trace-one section and its fixed theta-group translates with this CM model; match the selected local coordinates, determinant roots, additive character and periods; establish the prescribed finite-place derivative or divided-power bounds for the tied family at the nonordinary primes actually used; and prove the all-polynomial kernel identity with the specified volumes and terminal rows. None is implied by a unit theta value at0 or by an algebraic normalized germ alone.

\subsection{AM-DO-015. The exact line-class and parity gate, and its torsion-translation repair}\label{am-do-015.-the-exact-line-class-and-parity-gate-and-its-torsion-translation-repair}

For this paragraph carry ALP-LIFT-07's identification of the R004 zero characteristic with D135=W1+W3+W5−∞ and of the trace-one line Lw0 with the unique ζ-invariant symmetric principal line L0 attached to2∞. Let t be their nonzero two-torsion displacement. In root's period coordinates it has

\begin{verbatim}
t=(ε+τδ)/2,  ε=(1,0,1),  δ=(1,1,1).        (DO15.1)
\end{verbatim}

Thus L1=L\_D135=t\_t\^{}\emph{L0 and t\_t\^{}}L1=L0. The source's actual arithmetic \nolinkurl{Level} requires M odd: \nolinkurl{ReciprocityTheory/ReductiveGroups/ArithmeticGeometry/BirationalGeometry/ToroidalCompactification/TwoWeightEqOneAddRoot.lean}, lines53--62, field \nolinkurl{odd}, line56. Its \nolinkurl{SAB.prim}, line82, also fixes the trace-one line by the cube identity. The actual file was fetched and read in this run; this is not inferred only from the convenient example M=29\^{}e.

Let N=L1⊗L0\^{}(−1), a nontrivial element of Pic⁰(J){[}2{]}. For every M,

\begin{verbatim}
[M]^*N=N^M,
[M]^*L1 ≃ L0^(M²) ⊗ N^M.                 (DO15.2)
\end{verbatim}

Consequently the value-one pullback family F1(Mu) lies in L0\^{}(M²) for even M, but remains in a different line for odd M. Theta-group actions on L0\^{}(M²), including those in \nolinkurl{IsVtOver}, preserve that line and cannot by themselves change its Pic⁰ class. Even M is therefore a valid geometric way to remove this particular twist, but it does not instantiate the current arithmetic \nolinkurl{Level}.

Proof. Pullback by{[}M{]} on Pic⁰ is multiplication by M, so{[}M{]}\^{}\emph{N=N\^{}M. The symmetric rigidified multiplication isomorphism gives{[}M{]}\^{}}L0=L0\^{}(M²). Their tensor product proves(DO15.2). As N has exact order2, N\^{}M is trivial exactly for even M. A theta-group action is, by definition, an automorphism of sections of its specified line via an isomorphism of a translate of that same line. It does not identify two different Pic⁰ classes. This proves the stated obstruction to that direct substitution. ∎

\subsubsection{Repair by changing the base section to the actual trace-one line}\label{repair-by-changing-the-base-section-to-the-actual-trace-one-line}

There is a precise algebraic torsion translation that fixes the line issue without changing the odd arithmetic level. Let H(z,w)=zᵗY⁻¹ conjugate(w), Q(z)=πzᵗY⁻¹z/2, and write F1 also for its expression in z. Define the normalized reduced translation

\begin{verbatim}
Σ(z)=exp(−πH(z,t)−πH(t,t)/2) F1(z+t).       (DO15.3)
\end{verbatim}

This is Bannai--Kobayashi's U\^{}M\_t convention, equations(14) and following pp14--15. Since t\_t\^{}*L1=L0, it is a nonzero section of the actual trace-one line L0. It vanishes at0, with the expected effective even theta characteristic; it is not asserted to have value1 there.

Its exact theta-series formula is

\begin{verbatim}
Σ(z)=−i · exp(Q(z))
      θ[δ/2; ε/2](z;τ) / θ[0;0](0;τ).    (DO15.4)
\end{verbatim}

To check the constant, substitute t=(ε+τδ)/2. The linear part of Q(z+t)−πH(z,t)−Q(z) is πiδᵗz. Its constant part after subtracting πH(t,t)/2 is(πi/4)(δᵗτδ+εᵗδ). By the definition of a theta characteristic,

\begin{verbatim}
θ[δ/2;ε/2](z;τ)
  =exp(πiδᵗτδ/4+πiδᵗz+πiδᵗε/2) θ(z+t;τ).
\end{verbatim}

The quotient of these two exponentials is exp(−πiεᵗδ/4)=−i because εᵗδ=2. This proves(DO15.4), including its scalar.

Bannai--Kobayashi Theorem2.9 applies to this two-torsion translation and shows that the coefficients of Σ in the algebraic u=Pz coordinates are algebraic over E(J{[}16{]}); an extra root-of-unity convention for U\_t rather than U\^{}M\_t changes only the stated finite torsion scalar. This is a generic algebraic statement. We do not infer that every coefficient, or its chosen integral-frame scalar, is a unit at a nonordinary prime.

The SAME Gaussian constant49 is preserved. Indeed R004 gives \textbar F1(z+t)\textbar\textless49 exp(πH(z+t,z+t)/2). Taking the absolute value of the exponential in(DO15.3) subtracts πReH(z,t)+πH(t,t)/2; the cross terms cancel exactly, leaving

\begin{verbatim}
|Σ(u)|<49 exp((π/2)Σ_i |u_i|²/G_i).         (DO15.5)
\end{verbatim}

For every allowed odd M=29\^{}e from carried DP59, set Σ\_M(u)=Σ(Mu). This is literally{[}M{]}\^{}*Σ in L0\^{}(M²). The standard rigidified multiplication isomorphism, with t′=t=0 and the identity theta action, verifies the generic \nolinkurl{IsVtOver} clauses, exactly as in DP59(c). Its normalized derivatives satisfy

\begin{verbatim}
M^(−|α|) ∂^αΣ_M(0)=∂^αΣ(0),
\end{verbatim}

and the R004 level-independent Cauchy bound holds with the same49. This gives an actual algebraic base section, in the correct trace-one line, and an actual generic permitted pullback family. It still does not prove membership of the selected row in an iterated automorphic theta image, any prescribed nonzero positive-degree jet, or the global \nolinkurl{PinX.mv} identity for every moved lattice.

\subsection{AM-DO-016. A literal finite Schwartz input and the retained CM period}\label{am-do-016.-a-literal-finite-schwartz-input-and-the-retained-cm-period}

The trace-one section in(DO15.4) has a concrete theta input. Fix the classical additive-character convention with ψ\_∞(x)=exp(2πix), so ψ\_f(r)=exp(−2πir) on rational r. Put α=δ/2 and β=ε/2, and define the finite Schwartz function on A\_f³ by

\begin{verbatim}
φ_t(q_f)=−i·1_(α+Zhat³)(q_f) ψ_f(−q_fᵗβ). (DO16.1)
\end{verbatim}

This is locally constant and compactly supported. Its values on rational points q=n+α are−i exp(2πi(n+α)ᵗβ), hence algebraic fourth roots of unity. Its rational theta sum is exactly

\begin{verbatim}
Σ_(q∈Q³) φ_t(q) exp(πi qᵗτq+2πi qᵗz)
   =−i θ[α;β](z;τ).                        (DO16.2)
\end{verbatim}

Multiplying by the actual Chern completion gives

\begin{verbatim}
exp(Q(z)) Θ(φ_t;z,τ)=θ[0;0](0;τ)·Σ(z).     (DO16.3)
\end{verbatim}

All derivatives of(DO16.3) follow from normally convergent theta series. Thus in this explicit one-fiber Schrödinger/Chern convention the analytic-to-normalized-algebraic comparison constant is precisely

\begin{verbatim}
c0=θ[0;0](0;τ),  |c0|>5/8.                (DO16.4)
\end{verbatim}

It has not been set to1. The phase−i is part of the finite input, so no phase is omitted from c0. The finite input is a phase and a half-lattice translation of the standard input1\_Zhat³, hence belongs to its finite Heisenberg orbit. It is supported in (1/2)Zhat³ and invariant under2Zhat³, so its finite Schrödinger support/period quotient is (1/2)Z³/2Z³. This is a modulus4 statement for the input function, not a claimed level4 stabilizer for the full metaplectic group action.

Proof. Compact support and local constancy follow from the coset of the compact open subgroup Zhat³ and the finite-conductor additive character. At a rational point the stated product formula for ψ gives its displayed value. Its support on Q³ is precisely α+Z³, so writing q=n+α identifies the sum term by term with the defining series for the theta characteristic. This proves(DO16.2), and substitution of(DO15.4) proves(DO16.3). The inequality in(DO16.4) is carried R004. ∎

This is a genuine fiberwise theta equality with an explicit finite input and an explicit retained period. It is more than an assertion of abstract nonvanishing. It is still less than the required common arithmetic \nolinkurl{ClsK.Core} for the whole family. In particular the recovered Alpöge convention is \nolinkurl{IsStdChar ψ (-1)}; a comparison with it must track the conjugate additive character and the associated polarization or coordinate convention. Nor have the other moved CM lattices, their covolumes, the fixed all-degree branching maps or the terminal mass matrix been identified by(DO16.3). A common scalar for one fixed Heisenberg module does not supply these further identities.

\clearpage

\section{AM-K012--016 --- actual reduced theta germs and derivative bounds on y²=x⁷−1}\label{am-k012016-actual-reduced-theta-germs-and-derivative-bounds-on-yuxb2xux20771}

Author: GPT-6 Astra Pro, ALP\_MACHINE agent /root/kummer\_extension, 21 November 2026 UTC. This note supplies actual germs on the specified Jacobian. The accepted Alpöge theorem is not re-examined. The distinction between the translated non-effective characteristic and the source trace-one line is maintained.

\subsection{0. Statements and conventions}\label{statements-and-conventions}

Let \(E=\mathbb Q(\zeta)\), \(\zeta=\zeta _7\), and \[
C:y^2=x^7-1,\quad J=\operatorname{Jac}(C),\quad
\omega_i=x^{i-1}dx/y,\quad u_i=\int\omega_i\quad(1\leq i\leq3).
\] Write \(P_a=(\zeta^a,0)\). The translated characteristic is \(D=P_1+P_3+P_5-\infty\). Its value-one Chern-reduced germ is AM-R-004's \[
F(u)=e^{\pi z^tY^{-1}z/2}\frac{\theta(z,\tau)}{\theta(0,\tau)},
\qquad z=P^{-1}u,\quad Y=\operatorname{Im}\tau.
\tag{K12.1}
\] ALP-LIFT-07 identifies this characteristic exactly and proves that the value-one normalization is primitive at every odd good prime. The source trace-one line is instead \(L_\infty=L_{2\infty}\). Let \(s_q\) denote its theta section normalized by the determinant-of-cohomology tangent cone \[
q(u)=u_1u_3-u_2^2.
\] Its Chern-reduced germ is denoted \(\sigma_q(u)\). ALP-LIFT-08 identifies this exact scalar normalization with an actual primitive section at every \(p\nmid14\).

A ``Hurwitz coefficient'' means \(\partial^\alpha f(0)\), so \(f=\sum_\alpha\partial^\alpha f(0)u^\alpha/\alpha!\). It does \textbf{not} mean the divided Taylor coefficient. Valuations satisfy \(v_{\mathfrak p}(p)=1\).

The results are:

\begin{enumerate}
\def\labelenumi{\arabic{enumi}.}
\tightlist
\item
  \textbf{AM-K012.} The canonical Schiffer bidifferential equals the explicit algebraic Baker bidifferential. Consequently \(F\) is precisely the Baker even sigma function for \(D\), with no unrecorded quadratic exponential. An algebraic recursion determines all its Hurwitz coefficients, which belong to \(\mathcal O_E[1/14]\). In particular \[
  \partial_1\partial_3F(0)=-\zeta^2/4.
  \tag{K12.2}
  \]
\item
  \textbf{AM-K013.} At every prime above 11, \[
  v_{\mathfrak p}(\partial_i^nF(0))
  \geq\max\left\{0,c_in-\frac{s_{11}(n)}{10}\right\},
  \quad c_1=c_3=\frac1{20},\quad c_2=\frac1{220}.
  \tag{K13.1}
  \] Here \(s_p(n)\) is the base-\(p\) digit sum. The proof uses a stable integral ring for the actual logarithmic Hessian and three finite exact coefficient checks.
\item
  \textbf{AM-K014.} The actual source-line germ is \[
  \sigma_q(u)=4\sigma(u_1/2,u_2/2,u_3/2),
  \tag{K14.1}
  \] where \(\sigma\) is the algebraically normalized \((2,7)\) sigma function in half-differential coordinates. An explicit integer determinant formula determines every coefficient. Thus \(\partial^\alpha\sigma_q(0)\in2^{2-|\alpha|}\mathbb Z\). It is even, has CM weight 4 modulo 7, and \[
  \partial_1\partial_3\sigma_q(0)=1,\quad
  \partial_2^2\sigma_q(0)=-2,\quad
  \partial_1^4\sigma_q(0)=2,\quad
  \partial_3^6\sigma_q(0)=1.
  \tag{K14.2}
  \]
\item
  \textbf{AM-K015.} The bounds (K13.1) hold for \(\sigma_q\) too, by the actual integral algebraic ratio \[
  \sigma_q(2u)=r(u)F(2u).
  \tag{K15.1}
  \] Pullback by \([2]\) trivializes the two-torsion difference of the lines. This is a comparison of germs; the final section belongs to \(L_\infty\), and no even value is inserted into the source's odd Level.M.
\item
  \textbf{AM-K016.} Each actual pure restriction \(f_i(t)=\sigma_q(te_i)\) is an entire function of order 2 and exact type \(\pi/(2G_i)\). Its nonzero coefficient orders form an infinite computable sequence \(n_{i,k}\) satisfying \(n_{i,k+1}/n_{i,k}\to1\), within the necessary congruence class in (K14.11). The coefficient limsup is \[
  \limsup_{n\to\infty}
    \left(\frac{|\partial_i^n\sigma_q(0)|}{\sqrt{n!}}\right)^{1/n}
    =\sqrt{\pi/G_i}.
  \tag{K16.1}
  \] The support and growth assertions are archimedean; they do not bound finite-place valuations above.
\end{enumerate}

These statements do not establish nonvanishing in every admissible degree, identify a prescribed automorphic harmonic row, or prove the terminal every-\(\varepsilon\) window.

\subsection{1. AM-K012: identify the actual reduced frame}\label{am-k012-identify-the-actual-reduced-frame}

Set \(f(x)=x^7-1\) and \(\mathcal F(x,z)=x^3z^3(x+z)-2\). The Baker fundamental bidifferential is \[
\Omega_B(P,Q)=
\frac{\mathcal F(x,z)+2y(P)y(Q)}
 {4(x-z)^2y(P)y(Q)}\,dx\,dz.
\tag{K12.3}
\] The \(2y(P)y(Q)\) term removes a spurious pole on the conjugate diagonal. This symmetric differential has diagonal leading term \(dt_Pdt_Q/(t_P-t_Q)^2\) and no other poles. Exact general formulas: Buchstaber, Enolskii, and Leykin (1997), §1.2, equations (1.7)--(1.9); Edelstein, Gómez-Reino, and Mariño (2001), equations (3.3)--(3.6).

Let \(B_0\) be the normalized fundamental bidifferential with zero \(a\)-periods and let \(\nu=P^{-1}\omega\) be the column of normalized holomorphic differentials. The Schiffer bidifferential is \[
\Omega_S=B_0-\pi\nu(P)^tY^{-1}\nu(Q).
\tag{K12.4}
\] It is independent of the symplectic marking. Its period relation pairs its periods with those of the intrinsic Bergman reproducing kernel and characterizes it without a marking. Shramchenko (2004), §4, printed p.1468, gives (K12.4) and these period relations. Thus every curve automorphism preserves \(\Omega_S\).

The automorphism \(a:(x,y)\mapsto(\zeta x,y)\) preserves (K12.3): \(\mathcal F\) has only terms of degree 7 and a constant, and \(dx\,dz/(x-z)^2\) is invariant. Moreover \(a^*\omega_i=\zeta^i\omega_i\). The difference \(\Omega_B-\Omega_S\) is a global holomorphic bidifferential, hence a linear combination of \(\omega_i(P)\omega_j(Q)\). Every term has weight \(i+j\in\{2,3,4,5,6\}\), none divisible by 7. Invariance forces every coefficient to vanish: \[
\boxed{\Omega_B=\Omega_S.}
\tag{K12.5}
\] This uniqueness proof makes no inference about finite-place integrality from a Hodge splitting.

The quadratic normalization includes a factor of two. If Baker's sigma theta expression has exponential \(e^{u^tKu}\), its normalized bidifferential is \[
B_0=\Omega_B+2\omega(P)^tK\omega(Q).
\] This follows by taking \(a\)-periods of the second-kind differential identity, equivalently Edelstein et al., equations (3.14)--(3.18). Comparison gives \[
2K=\pi P^{-t}Y^{-1}P^{-1}.
\tag{K12.6}
\] Hence the exponential is exactly \(e^{\pi z^tY^{-1}z/2}\). The characteristic is zero by ALP-LIFT-07, and normalization at zero divides by \(\theta(0,\tau)\). This proves (K12.1). The same calculation applies to the characteristic \(2\infty\).

\subsection{2. Actual algebraic logarithmic Hessian and recursion}\label{actual-algebraic-logarithmic-hessian-and-recursion}

Let \(a_1=\zeta,a_2=\zeta^3,a_3=\zeta^5\). Near the ordered divisor \(((a_1,0),(a_2,0),(a_3,0))\), put \[
U(X)=\prod_j(X-x_j),\qquad
\ell_j(X)=\frac{\prod_{m\ne j}(X-x_m)}{U'(x_j)}
         =\sum_{i=1}^3q_{ji}X^{i-1}.
\] The matrix \(q\) is the inverse of \(V_{ij}=x_j^{i-1}\). The invariant derivatives are \[
D_ix_j=y_jq_{ji},\qquad D_iy_j=\tfrac72x_j^6q_{ji}.
\tag{K12.7}
\] The first formula inverts \(\sum_jx_j^{k-1}dx_j/y_j=du_k\); the second differentiates \(y_j^2=x_j^7-1\). They are regular at the branch points. The \(D_i\) commute because they lift the commuting invariant derivatives on \(J\).

Let \(s_1=\sum x_j,s_2=\sum_{j<m}x_jx_m,s_3=x_1x_2x_3\). The symmetric Kleinian matrix \(W\) is determined by \[
W_{33}=s_1/4,\quad W_{23}=-s_2/4,\quad W_{13}=s_3/4
\tag{K12.8}
\] and \[
\sum_{i,k}W_{ik}x_j^{i-1}x_m^{k-1}
=\frac{\mathcal F(x_j,x_m)-2y_jy_m}{4(x_j-x_m)^2},
\qquad j<m.
\tag{K12.9}
\] These are Buchstaber et al.~(1997), Corollary 2.1.2, equation (2.9), and Theorem 2.2, equations (2.10)--(2.14). Their canonical model has leading coefficient 4. Passing from \(Y^2=4(x^7-1)\), \(du^{\rm can}=dx/Y\), to our \(u=2u^{\rm can}\) divides the logarithmic Hessian by 4 and gives exactly these factors.

After the last row is inserted, the unknowns \(W_{11},W_{12},W_{22}\) have coefficient rows \((1,x_j+x_m,x_jx_m)\). The determinant is a signed product of the three differences. Here is explicit elimination: \[
\begin{split}
R_{jm}={}&\frac{\mathcal F(x_j,x_m)-2y_jy_m}{4(x_j-x_m)^2}
-W_{13}(x_j^2+x_m^2)\\
&-W_{23}(x_j^2x_m+x_jx_m^2)-W_{33}x_j^2x_m^2,\\
t_1={R_{13}-R_{12}\over x_3-x_2},\qquad
t_2={R_{23}-R_{12}\over x_3-x_1}.
\end{split}
\] Then \[
W_{22}={t_1-t_2\over x_1-x_2},\quad W_{12}=t_1-x_1W_{22},\quad
W_{11}=R_{12}-(x_1+x_2)W_{12}-x_1x_2W_{22}.
\tag{K12.10}
\]

For \(h=\log F\), \(h(0)=0\), the characteristic translation gives \[
\partial_i\partial_jh(u)=-W_{ij}(a+u),
\quad a=\operatorname{AJ}(P_1+P_3+P_5-3\infty).
\tag{K12.11}
\] This is the Hessian of the specified actual theta germ. Translating changes an affine exponential only, whose Hessian is zero; evenness and \(F(0)=1\) fix the first and zeroth derivatives. Thus (K12.11) determines \(F\) uniquely.

\subsubsection{Stable integral ring}\label{stable-integral-ring}

Let \(R_0=\mathcal O_E[1/14]\), and \[
\mathcal R=R_0[x_1,x_2,x_3,y_1,y_2,y_3,
(x_1x_2x_3\prod_{j<m}(x_j-x_m))^{-1}]
/(y_j^2-x_j^7+1)_{j=1}^3.
\tag{K12.12}
\] Every \(D_i\) preserves this ring, by (K12.7), the quotient rule, and differentiation of the defining relations. Every \(W_{ij}\) belongs to it. Evaluation \(e:\mathcal R\to R_0\) at \((x_j,y_j)=(a_j,0)\) is defined: these roots are units, and root differences have prime divisors only above 7.

The Abel map from this ordered-divisor open set to \(J\) is étale. In a branch chart the local parameters \(y_j\) give the differential matrix \[
\left(2x_j^{i-1}/(7x_j^6)\right)_{ij},
\] with invertible determinant. At nonbranch points use \(x_j\); the same Vandermonde criterion applies. This concerns algebraic formal coordinates on \(J\); it is \textbf{not} a claim that the Abelian logarithms \(u_i\) are integral formal parameters.

For \(|\alpha|\geq2\), choose \(i,j\) with \(\alpha-e_i-e_j\geq0\). Then \[
h_\alpha:=\partial^\alpha h(0)
=-e(D^{\alpha-e_i-e_j}W_{ij})\in R_0.
\tag{K12.13}
\] The result is independent of the choice by the Kleinian Hessian identity. Also \(h_0=h_{e_i}=0\). If \(S_\alpha=\partial^\alpha F(0)\), then \[
S_0=1,\qquad
S_\alpha=\sum_{0\leq\beta\leq\alpha-e_i}
{\alpha-e_i\choose\beta}h_{\beta+e_i}S_{\alpha-e_i-\beta}
\quad(\alpha_i>0).
\tag{K12.14}
\] This is \(\partial_iF=F\partial_ih\). It involves integer binomial coefficients and no division by an order, proving \(S_\alpha\in R_0\).

\subsubsection{Initial matrix and checks}\label{initial-matrix-and-checks}

Let \(Q(X)=\prod_{a\in\{1,3,5\}}(X-\zeta^a)\), \(R(X)=f(X)/Q(X)\). Bolza's formula is \[
\sum_{i,j}W_{ij}(a)X^{i-1}Z^{j-1}
=\frac{\mathcal F(X,Z)-Q(X)R(Z)-Q(Z)R(X)}{4(X-Z)^2}.
\tag{K12.15}
\] See Edelstein et al.~(2001), equation (3.24), reversing their descending basis. Explicitly, \(4W(a)\) is \[
\begin{pmatrix}
-\zeta^5-\zeta^4-2\zeta^2+\zeta-2&\zeta^3+1&\zeta^2\\
\zeta^3+1&-2\zeta^4+\zeta^3-\zeta^2+\zeta-2&
\zeta^5+\zeta^3+\zeta^2+1\\
\zeta^2&\zeta^5+\zeta^3+\zeta^2+1&\zeta^5+\zeta^3+\zeta
\end{pmatrix}.
\tag{K12.16}
\] This proves (K12.2). Exact polynomial division modulo \(\Phi_7\) is checked in \nolinkurl{jacobian_theta_hessian_check.py.}

A separate numerical diagnostic, not a proof premise, used 65-digit independent theta summation with AM-R-004's periods. Its Hessian agrees with \(-W(a)\) to \(6.85\cdot10^{-66}\). Pure jets of orders 2,4,6,8 agree with the independent flow/Bell recurrence, largest observed discrepancy \(4.55\cdot10^{-64}\). Results and diagnostic scope are retained in jacobian\_theta\_hessian\_check.json.

\subsection{3. AM-K013: invariant-derivative nilpotence and actual logarithmic bounds}\label{am-k013-invariant-derivative-nilpotence-and-actual-logarithmic-bounds}

Fix an odd prime \(p\ne7\) and complete \(R_0\) at a prime above \(p\). The reduction is a smooth abelian variety with the order-seven automorphism. The \(D_i\) form an integral invariant basis, with distinct CM labels \(i\in\{1,2,3\}\).

In characteristic \(p\), the \(p\)-fold composition of a derivation is a derivation, because the intermediate binomial coefficients in the Leibniz formula vanish. Translation invariance is retained. The CM weight is multiplied by \(p\), so a label \(pi\bmod7\) outside \(\{1,2,3\}\) forces \(D_i^p=0\). If the label is inside this set, \(D_i^p\) is a scalar multiple of its uniquely labelled invariant derivative. Iteration gives \[
D_i^{q_i}=0\text{ on the special fibre},\qquad q_i=p^{h_i},
\tag{K13.2}
\] with these sufficient exponents:

\begin{longtable}[]{@{}lrrr@{}}
\toprule\noalign{}
\(p\bmod7\) & \(h_1\) & \(h_2\) & \(h_3\) \\
\midrule\noalign{}
\endhead
\bottomrule\noalign{}
\endlastfoot
2 & 2 & 1 & 1 \\
3 & 3 & 1 & 2 \\
4 & 1 & 2 & 1 \\
5 & 1 & 3 & 2 \\
6 & 1 & 1 & 1 \\
\end{longtable}

For \(p\equiv1\bmod7\), this argument supplies no exponent. The table is obtained by multiplying a label successively by \(p\bmod7\) until it leaves \(\{1,2,3\}\); no classification of the Newton polygon is needed.

Because the Abel map in §2 is étale, the zero derivation on the target has only the zero lift. Thus (K13.2) holds on \(\mathcal R/p\mathcal R\), and \[
D_i^{q_i}\mathcal R\subset p\mathcal R.
\tag{K13.3}
\] This is an operator assertion on a stable integral ring. Since \(D_i\mathcal R\subset\mathcal R\) and \(D_i(p)=0\), it may genuinely be iterated: \[
D_i^m\mathcal R\subset p^{\lfloor m/q_i\rfloor}\mathcal R.
\tag{K13.4}
\] For \(\ell_n=\partial_i^n\log F(0)\), applying this to \(W_{ii}\) gives \[
v_{\mathfrak p}(\ell_n)\geq
\max\{0,\lfloor(n-2)/q_i\rfloor\}\quad(n\geq2).
\tag{K13.5}
\] Also \(\ell_n=0\) for odd \(n\).

\subsubsection{Bell--Legendre lemma, with all losses}\label{belllegendre-lemma-with-all-losses}

Let \(p\) be odd, \(q=p^h\), and \[
H(t)=\sum_{j\geq2}\ell_jt^j/j!,\qquad
e^{H(t)}=\sum_{n\geq0}S_nt^n/n!.
\] Suppose \(H\) is even, all \(\ell_j\) are integral, and (K13.5) holds. Then \[
v(S_n)\geq
\max\left\{0,\frac{n}{(p-1)(q+1)}-\frac{s_p(n)}{p-1}\right\}.
\tag{K13.6}
\] If additionally \(v(\ell_{q+1})\geq1\), then \[
v(S_n)\geq
\max\left\{0,\frac{n}{2p^{h-1}(p-1)}-\frac{s_p(n)}{p-1}\right\}.
\tag{K13.7}
\]

To prove this, a complete Bell-polynomial term, indexed by \(m_j\geq0\), \(\sum jm_j=n\), is \[
\frac{n!}{\prod_jm_j!(j!)^{m_j}}\prod_j\ell_j^{m_j}.
\] Only even \(j\) occur. Legendre's formula bounds its valuation below by \[
-\frac{s_p(n)}{p-1}
+\sum_jm_j\left(v(\ell_j)+\frac{s_p(j)-1}{p-1}\right)
+\frac{\sum_js_p(m_j)}{p-1}.
\tag{K13.8}
\] The final summand is nonnegative.

For (K13.6), even \(j>0\) has \(s_p(j)\geq2\). If \(j\leq q+1\), then \((p-1)v(\ell_j)+s_p(j)-1\geq1\geq j/(q+1)\). If \(j\geq q+3\), put \(k=\lfloor(j-2)/q\rfloor\geq1\). Since \(j\leq(k+1)q+1\), \[
(p-1)v(\ell_j)+s_p(j)-1\geq(p-1)k+1\geq j/(q+1).
\] Insert this into (K13.8). Integrality supplies the maximum with zero.

For (K13.7), if \(2\leq j<q\), then \(j\leq p^{h-1}s_p(j)\), while \(s_p(j)-1\geq s_p(j)/2\). Thus the lower bound is \(j/(2p^{h-1})\). The sole even gap is \(j=q+1\): the additional assumption gives a lower bound at least \(p\), larger than \((q+1)/(2p^{h-1})\). For \(j\geq q+3\), retain \(k=\lfloor(j-2)/q\rfloor\). If \(k=1\), evenness gives \(j\leq2q\), so the lower bound \(p\) suffices. If \(k\geq2\), \[
(p-1)k+1-\frac{(k+1)q+1}{2p^{h-1}}
=\left(\frac p2-1\right)k+1-\frac p2-\frac1{2p^{h-1}}\geq0.
\] This holds for \(p\geq3,k\geq2,h\geq1\); its smallest case \(p=3,k=2,h=1\) is zero. Substitution into (K13.8) proves (K13.7).

Therefore (K13.6) is already an unconditional directional lower bound for the actual translated theta at every residue class in the table, including \(p\equiv3,5\bmod7\).

\subsubsection{Exact finite certificate for the improved 11-adic constants}\label{exact-finite-certificate-for-the-improved-11-adic-constants}

At \(p=11\), \(q_1=q_3=11\) and \(q_2=121\). The only additional checks needed are \[
\partial_1^{12}\log F(0)\equiv0,\qquad
\partial_3^{12}\log F(0)\equiv0,\qquad
\partial_2^{122}\log F(0)\equiv0\pmod{11}.
\tag{K13.9}
\] All three were verified by exact arithmetic in \[
B=\mathbb F_{11}[Z]/(1+Z+\cdots+Z^6),
\] the product of the two residue fields above 11. Zero in \(B\) proves zero at both primes.

The reproducible recurrence is \nolinkurl{jacobian_theta_direction_jets.py.} Its inputs are (K12.7)--(K12.10). To avoid power-series division, add \(z_{jm}=(x_j-x_m)^{-1}\) with \[
D_iz_{jm}=-z_{jm}^2(D_ix_j-D_ix_m).
\] This is a polynomial differential system over \(B\). The \(n\)-th Hurwitz coefficient of a product is \(\sum_{k=0}^n\binom nk a_kb_{n-k}\). The coefficient of order \(n\) in a variable's prescribed derivative gives its coefficient of order \(n+1\), with no division by \(n+1\). Initial branch-difference inverses are obtained by Euclidean polynomial inversion modulo \(\Phi_7\). Finally \(\ell_{n+2}=-D_i^nW_{ii}|_a\).

The executed command was:

\begin{verbatim}
python work/agents/jacobian_theta_direction_jets.py --prime 11 --order 160 --output jacobian_theta_direction_jets_mod11.json
\end{verbatim}

It retains all nonzero coefficient vectors. In each of (K13.9), all six coordinates in \(1,Z,\ldots,Z^5\) are zero. The companion \nolinkurl{jacobian_theta_11_certificate.py} checks the saved tables, checks the unit status of the retained coefficients, and regenerates the three exceptional coefficients using the stated recurrence.

The operator proof (K13.3)--(K13.5), the three exact checks, and (K13.7) now prove (K13.1) for every \(n\). The finite computation is used only for these three coefficients; no pattern is extrapolated.

There are exact all-order mod-11 consequences as well. In the reduced one-variable Hurwitz algebra, \(b^{11}=b(0)^{11}\): a positive divided-power monomial has zero 11-th power, as follows from \(v_{11}((11k)!/(k!)^{11})=s_{11}(k)>0\). The universal Bell/Jacobson identity is \[
\partial^{11}e^H/e^H=(H')^{11}+\partial^{11}H.
\] Iteration gives \(\partial^{11^h}e^H/e^H=\partial^{11^h}H\), since each constant \(H^{(11^j)}(0)\) is zero by evenness. Formula (K13.3) kills the logarithmic derivatives of orders at least \(q_i+2\) modulo 11; (K13.9) and oddness supply the two preceding orders. Hence \[
\partial_1^nF(0)\equiv\partial_3^nF(0)\equiv0\ (n\geq11),
\qquad\partial_2^nF(0)\equiv0\ (n\geq121).
\tag{K13.10}
\] This establishes a first congruence only. Repeated valuation gain in (K13.1) comes instead from the stable-ring estimate (K13.4) and the Bell lemma.

\subsection{4. AM-K014: an explicit determinant formula for the actual source sigma}\label{am-k014-an-explicit-determinant-formula-for-the-actual-source-sigma}

The standard theta divisor based at infinity is \(W_2\), hence its line is \(L_{2\infty}\). Buchstaber et al.~(1997), Definition 1 and equation (2.2), printed pp.7--8, normalize its Baker sigma by initial homogeneous term \(v_1v_3-v_2^2\) in the canonical half-differential coordinates. The construction below fixes the whole series and its sign, including the non-leading terms.

Use the integral arithmetic parameter \(t=-x^3/y\) at infinity. There is a unique \(A\in1+t^{14}\mathbb Z[[t^{14}]]\) satisfying \[
A=1-t^{14}A^7,\qquad x=t^{-2}A^{-1},\qquad y=-t^{-7}A^{-3}.
\tag{K14.3}
\] Substitution verifies both the curve equation and \(t=-x^3/y\). Coefficients of \(A\) are recursively determined without division, and all inverse powers \(A^{-m}\) are integral. With \(T=t^{14}\), \[
[T^0]A^{-m}=1,\qquad
[T^k]A^{-m}=\frac m k(-1)^{k-1}{7k-m-1\choose k-1}\quad(k\geq1).
\tag{K14.4}
\] For \(m=0,k>0\) this is zero. Formula (K14.4) is Lagrange inversion applied to \(1-A=T A^7\); its expression is an integer because \(A^{-m}\in\mathbb Z[[T]]\). The division in this closed expression may be avoided entirely by the recursion (K14.3).

Let the increasing nongaps be \[
n_1,n_2,\ldots=0,2,4,6,7,8,9,\ldots .
\] For even \(n=2k\), take \(\phi_n=x^k\); for odd \(n=2k+7\), take \(\phi_n=-x^ky\). Then \[
t^2\phi_n=t^{2-n}A^{-\lfloor n/2\rfloor}.
\tag{K14.5}
\] Thus the basis of \(t^2H^0(C-\{\infty\},\mathcal O_C)\) has leading exponents \[
2,0,-2,-4,-5,-6,\ldots ,
\] corresponding to the partition \(\lambda=(3,2,1)\).

Here is a finite algorithm for every coefficient. For any partition \(\mu\), choose \(L\geq\max(3,\operatorname{length}\mu)\), append zeros to \(\mu\), and define \[
\xi_\mu=
\det\left([t^{\,\mu_i-i-(2-n_j)}]\,
A(t)^{-\lfloor n_j/2\rfloor}\right)_{1\leq i,j\leq L}.
\tag{K14.6}
\] A negative power or a power not divisible by 14 has coefficient zero. Enlarging \(L\) adds an upper triangular identity tail and leaves the determinant unchanged. Each \(\xi_\mu\) is therefore an explicitly computable integer. It vanishes unless \(\mu\geq(3,2,1)\) and \(|\mu|=6+14k\) for some \(k\geq0\); \(\xi_{321}=1\).

Define polynomials \(p_m\) by \[
\sum_{m\geq0}p_m(v)z^m
=\exp(v_3z+v_2z^3+v_1z^5),\quad p_m=0\ (m<0),
\] and put \[
S_\mu(v)=\det(p_{\mu_i-i+j}(v))_{1\leq i,j\leq L}.
\] Then the actual Baker sigma is \[
\boxed{\sigma(v_1,v_2,v_3)=\sum_\mu\xi_\mu S_\mu(v).}
\tag{K14.7}
\] This is Nakayashiki's tau/sigma theorem, specialized explicitly: Nakayashiki (2010), §§4--5, equation (7), and Theorem 8; for the arithmetic parameter version, Ônishi (2015), Theorem 6.5. The affine-ring embedding is shifted by \(g-1=2\), exactly as in (K14.5), not by an unspecified convention. The monomial frame is upper unitriangularly equivalent to the normalized frame and has identical Plücker coordinates. Formula (K14.6) is its finite determinant.

We verify why there is no omitted exponential or coordinate correction in this specialization. The half-differentials \(x^{i-1}dx/(2y)\) have leading terms \(t^{6-2i}dt\); all later terms have exponents shifted by multiples of 14. Hence the gap-coordinate submatrix \(B\), at gaps \(1,3,5\), is the identity after the displayed reversal. The linear correction coefficients \(c_w\) vanish at these gaps because the square-root logarithm of the normalized highest differential contains powers only \(t^{14k}\). Likewise the regular coefficients \(q_{ij}\) of \(\Omega_B-dt_1dt_2/(t_1-t_2)^2\) vanish unless \(14\mid i+j\). Indeed the combination of \(a\) and the hyperelliptic involution acts on \(t\) by the primitive fourteenth root \(-\zeta^3\) and preserves \(\Omega_B\). For \(i,j\in\{1,3,5\}\), \(i+j\leq10\), so every retained \(q_{ij}\) vanishes. These are exactly the exponential corrections in the tau/sigma theorem.

Ônishi and Nakayashiki sometimes use \(-x^{i-1}dx/f_y\) rather than \(+x^{i-1}dx/(2y)\). Applying their convention to the identical model with \(Y=-y\) makes these differentials exactly ours. Equivalently a simultaneous sign of all three coordinates has no effect because the genus-three sigma here is even. The leading Schur polynomial fixes the remaining overall sign.

Every desired monomial has one finite weight \(5\alpha_1+3\alpha_2+\alpha_3\), so only finitely many partitions contribute to its coefficient in (K14.7). This is a checkable all-order construction of this particular function.

Each Schur polynomial in exponential times is Hurwitz integral: equivalently its coefficient of \(\prod_jT_j^{m_j}/m_j!\) is the integral symmetric-group character at cycle type \(1^{m_1}2^{m_2}\cdots\). One may instead use Ayano (2024), Lemma 7.2, or Ônishi (2015), Theorem 1.8. Since the \(\xi_\mu\) are integers, (K14.7) is Hurwitz integral over \(\mathbb Z\). Changing to \(u=2v\) and multiplying by 4 proves \[
\partial^\alpha\sigma_q(0)\in2^{2-|\alpha|}\mathbb Z.
\tag{K14.8}
\] This includes 7-adic integrality in the original, nonprimitive at 7, \(\omega\)-coordinates. It does not contradict ALP-LIFT-08's finite scalar content in primitive coordinates.

By (K12.5), the Baker gauge is the actual Chern-reduced gauge for this characteristic. The leading quadratic is \(q\); the principal theta space is one-dimensional. Thus (K14.1) agrees with the actual \(q\)-normalized section of ALP-LIFT-08, not merely with an abstractly isomorphic line.

\subsubsection{Leading coefficients and selection}\label{leading-coefficients-and-selection}

The weight-six term of (K14.7) is \[
S_{321}(v)=v_1v_3-v_2^2-\tfrac13v_2v_3^3+\tfrac1{45}v_3^6.
\tag{K14.9}
\] All subsequent weights are \(6+14k\). At weight 20, the coefficient of \(v_1^4\) is \(1/3\). This follows directly from (K14.6), or from the first-order frame shifts \(r\mapsto r+14\). In the latter calculation the only nonzero contributions after setting \(v_2=v_3=0\) are \[
(-1)(-1/6)+3(-1/2)+(-6)(-1/2)+8(-1/6)=1/3.
\] These terms come respectively from the partitions \((15,4,1)\), \((10,4,3,2,1)\), \((5,4,3,2,1^6)\), and \((3,2,2,2,1^{11})\). The script \nolinkurl{jacobian_sigma_determinant_certificate.py} checks these values without using the shift shortcut. It evaluates (K14.6) for all 627 partitions of 20, checks two extra identity-tail sizes for every nonzero coordinate, and computes the Schur determinants exactly. The retained JSON contains all 13 nonzero weight-20 Plücker coordinates and the four nonzero contributions above. It also checks (K14.9) symbolically and all four jets in (K14.2).

Consequently the beginning contains \[
\sigma_q(u)
=u_1u_3-u_2^2-\frac1{12}u_2u_3^3
+\frac1{720}u_3^6+\frac1{12}u_1^4
+\text{other terms of weights at least 20}.
\tag{K14.10}
\] In this notation \(u_1,u_2,u_3\) have weights \(5,3,1\); the displayed \(u_1^4\) is the indicated weight-20 term, and the last phrase excludes it. This proves (K14.2).

Every monomial is even in total degree, and its CM weight \(\alpha_1+2\alpha_2+3\alpha_3\) is \(4\bmod7\), since \[
\alpha_1+2\alpha_2+3\alpha_3
\equiv3(5\alpha_1+3\alpha_2+\alpha_3)\equiv18\equiv4\pmod7.
\] Thus the necessary pure-direction congruences are \[
n\equiv4\pmod{14}\ (i=1),\qquad
n\equiv2\pmod{14}\ (i=2),\qquad
n\equiv6\pmod{14}\ (i=3).
\tag{K14.11}
\] These are selection rules, not assertions that every permitted coefficient is nonzero.

\subsection{5. AM-K015: transfer to the actual source section}\label{am-k015-transfer-to-the-actual-source-section}

Fix \(p\nmid14\). On the good abelian scheme, the rigidified line difference is \[
L_D^{\rm rig}=L_\infty^{\rm rig}\otimes\mathcal P_T,\qquad 2T=0.
\] The canonical rigidified Poincaré isomorphism gives \([2]^*\mathcal P_T\simeq\mathcal P_{2T}\simeq\mathcal O\), integrally. Hence the two pulled-back theta sections can be divided to give the actual rational function \(r\) of (K15.1). The denominator has unit origin value by ALP-LIFT-07. The numerator is integral in the \(q\)-normalization by ALP-LIFT-08. Therefore \(r\) is regular and integral on an affine neighborhood \(\mathcal U\) of the origin.

The identity (K15.1) uses exactly the Chern-reduced gauges. The metrics of the two lines have the same translation-invariant curvature; their quotient is the flat metric on \(\mathcal P_T\). The normalized rigidified trivialization of \([2]^*\mathcal P_T\) is parallel and has value one at the origin. Thus it introduces neither a linear nor a quadratic exponential. Equivalently, the ratio of the two analytic reduced multipliers at \(2u\) is periodic: the difference semicharacter has order two. This confirms that the analytic ratio is the same algebraic \(r\), with no unknown scalar.

The invariant derivatives preserve \(\mathcal A=\Gamma(\mathcal U,\mathcal O)\): regular derivations preserve a localization of a regular function algebra. The CM-label proof (K13.2) applies on this actual abelian scheme, so \[
D_i^{q_i}\mathcal A\subset p\mathcal A,\qquad
v_{\mathfrak p}(D_i^kr(0))\geq\lfloor k/q_i\rfloor.
\tag{K15.2}
\] This is a separate stable-ring application. It does not assert that a Chern frame itself is an integral formal frame.

Suppose a bound already proved for \(F\) is \[
v_{\mathfrak p}(\partial_i^nF(0))
\geq c n-s_p(n)/(p-1).
\] Both constants used in §3 satisfy \[
cq_i\leq1,\qquad c\leq 1/(p^{h_i-1}(p-1)).
\tag{K15.3}
\] Differentiate (K15.1) \(n\) times along \(u_i\). Since \(2\) is a \(p\)-adic unit, the \(k\)-th Leibniz term has valuation at least \[
\begin{split}
&v_p{n\choose k}+\lfloor k/q_i\rfloor
+c(n-k)-s_p(n-k)/(p-1)\\
&=cn-s_p(n)/(p-1)
+\lfloor k/q_i\rfloor+s_p(k)/(p-1)-ck.
\end{split}
\tag{K15.4}
\] Write \(k=aq_i+b\), \(0\leq b<q_i=p^{h_i}\). Then \[
s_p(k)=s_p(a)+s_p(b),\qquad
b\leq p^{h_i-1}s_p(b).
\] Using (K15.3), \[
\lfloor k/q_i\rfloor+s_p(k)/(p-1)
\geq a+b/(p^{h_i-1}(p-1))\geq c(aq_i+b)=ck.
\] Thus every term satisfies the claimed lower bound. The sum does too. All coefficients of \(\sigma_q\) are integral at \(p\) by (K14.8), giving the maximum with zero.

This proves (K13.1) for the actual source-line germ, and also transfers the generic bound (K13.6) at every prime in the table. In particular, it gives explicit positive linear lower slopes with a logarithmic digit-sum loss at every \(p\equiv3,5\bmod7\).

The argument respects the source's odd level condition. The auxiliary map \([2]\) is used only to compare two already constructed sections and their derivatives. The section whose bounds are concluded is \(s_q\in H^0(J,L_\infty)\). A later odd-\(M\) pullback of this section has the exact normalization \[
M^{-n}\partial_i^n\sigma_q(Mu)|_{u=0}=\partial_i^n\sigma_q(0).
\] No translated \(D\)-line is substituted for \(L_\infty\).

The first congruence (K13.10) transfers too. Modulo 11, \(D_i^{q_i}\) is a derivation and kills the actual integral function \(r\). Applying it to (K15.1) and using the corresponding identity for \(F\) gives \[
\partial_1^n\sigma_q(0)\equiv\partial_3^n\sigma_q(0)\equiv0\ (n\geq11),
\qquad \partial_2^n\sigma_q(0)\equiv0\ (n\geq121).
\tag{K15.5}
\] The even order sets \(0,2,\ldots,10\) and \(0,2,\ldots,120\) found to consist of units for \(F\) are not asserted to be unit sets for \(\sigma_q\). Its different characteristic has its own selection rules.

\subsection{6. AM-K016: exact archimedean growth and controlled nonvanishing}\label{am-k016-exact-archimedean-growth-and-controlled-nonvanishing}

Fix \(i\), and put \[
f(t)=\sigma_q(te_i)=\sum_{n\geq0}a_nt^n,
\qquad \alpha=\frac{\pi}{2G_i}>0,\qquad
M(R)=\max_{|t|\leq R}|f(t)|.
\] The exact metric matrix is AM-R-004's diagonal \(G\). AM-R-006 gives the explicit bound \[
|f(t)|\leq C e^{\alpha|t|^2},\qquad C=6125/16.
\tag{K16.2}
\] The initial jets (K14.2) prove that \(f\) is nonzero for each \(i\). We now prove statements about this actual function, rather than all entire functions satisfying (K16.2).

\subsubsection{A lower growth bound from the compact theta metric}\label{a-lower-growth-bound-from-the-compact-theta-metric}

The Chern norm of the section is a continuous function on the compact torus \(J(\mathbb C)\). After a fixed constant normalization it pulls back along \(t\mapsto te_i\) as \[
g(t)=|f(t)|e^{-\alpha|t|^2}.
\tag{K16.3}
\] Restrict temporarily to real \(t\), and let \(K\) be the closure in \(J(\mathbb C)\) of this real one-parameter subgroup. It is a compact subgroup. Some real \(t_0\) has \(f(t_0)\ne0\), since otherwise the identity theorem would make \(f\) zero. Choose \(\varepsilon>0\) such that the relatively open set \[
V=\{x\in K:\|s_q(x)\|>\varepsilon\}
\] is nonempty, with the same fixed norm normalization as (K16.3).

The nonnegative real orbit is dense in \(K\). Here is an elementary justification. Its closure \(S\) is a compact subsemigroup of the torus. For any \(x\in S\), a pigeonhole argument applied to integer multiples of \(x\) in a successively finer finite cover produces positive multiples approaching zero. These return multiples may be chosen unbounded: if a bounded subsequence already returns arbitrarily close, one of its positive indices is an exact period and its multiples give unbounded exact returns. If the return indices are \(m_j\to\infty\), then \((m_j-1)x\to-x\), so \(-x\in S\). Thus \(S\) is a subgroup and must equal the closure \(K\) of the full real orbit. Positive translates of every \(x\in K\) therefore meet \(V\).

For each \(x\in K\), choose \(t_x\geq0\) with \(x+t_xe_i\in V\). The open sets \(V-t_xe_i\) cover \(K\); compactness supplies a finite subcover. Let \(L\) be the largest time used. For every \(r\geq0\), at least one \(t\in[r,r+L]\) has \(g(t)>\varepsilon\). Therefore \[
\varepsilon e^{\alpha(R-L)^2}\leq M(R)\leq Ce^{\alpha R^2}
\qquad(R\geq L).
\tag{K16.4}
\] In particular, \[
\log M(R)=\alpha R^2+O(R),
\tag{K16.5}
\] and \(f\) has order 2 and exact type \(\alpha\). It cannot be a polynomial, so it has infinitely many nonzero coefficients. The constants \(L,\varepsilon\) and the resulting thresholds are not made effective by this compactness argument.

\subsubsection{Exact coefficient limsup}\label{exact-coefficient-limsup}

Optimizing Cauchy's bound at radius \(\sqrt{n/(2\alpha)}\) gives \[
|a_n|\leq C\left(\frac{2e\alpha}{n}\right)^{n/2}\quad(n>0).
\tag{K16.6}
\] Consequently \(\limsup n|a_n|^{2/n}\leq2e\alpha\). If this inequality were strict, choose \(0<\beta<\alpha\) with \(|a_n|\leq(2e\beta/n)^{n/2}\) for all sufficiently large \(n\). With \(N=2\beta R^2\), a term at radius \(R\) is then at most \[
\exp\left(\frac N2\varphi(n/N)\right),
\qquad \varphi(x)=x(1-\log x),\qquad \varphi(1)=1.
\] The first \(\lceil e^2N\rceil\) terms sum to \(O(N+1)e^{N/2}\), and the remaining terms have a bounded geometric tail because the derivative of their logarithmic bound is at most \(-1\). The finitely many excluded terms form a fixed polynomial. Hence \(\log M(R)\leq\beta R^2+O(\log(R+2))\), contradicting (K16.5). We have proved \[
\limsup_{n\to\infty}n|a_n|^{2/n}=2e\alpha=\pi e/G_i.
\tag{K16.7}
\] Since \(\partial_i^n\sigma_q(0)=n!a_n\), Stirling's formula converts (K16.7) exactly to (K16.1). This is a limsup over all orders, with zero terms allowed. It does not assert that every nonzero derivative attains the upper bound.

\subsubsection{Every sufficiently large multiplicative window contains support}\label{every-sufficiently-large-multiplicative-window-contains-support}

Fix \(0<\delta<1\). We claim that every sufficiently large real \(N\) has a nonzero coefficient with index in \[
[(1-\delta)N,(1+\delta)N].
\tag{K16.8}
\] Suppose instead that the whole window is zero and put \(R=\sqrt{N/(2\alpha)}\). By (K16.6), each other term is bounded by \(C\exp(N\varphi(n/N)/2)\). Define \[
\kappa_\delta=1-\max\{\varphi(1-\delta),\varphi(1+\delta)\}>0.
\] There are at most \(N+1\) terms below the window, each bounded by \(C\exp((1-\kappa_\delta)N/2)\). Above the window the derivative of the logarithmic bound with respect to \(n\) is at most \(-\tfrac12\log(1+\delta)\), so those terms sum to a geometric tail with ratio at most \((1+\delta)^{-1/2}\). The constant term is zero for this source section; even if it were present it would satisfy the same estimate for large \(N\). Thus \[
M(R)\leq C\left(N+1+\frac1{1-(1+\delta)^{-1/2}}\right)
e^{(1-\kappa_\delta)N/2}.
\tag{K16.9}
\] But (K16.4) gives \(\log M(R)\geq N/2-O(\sqrt N)\), a contradiction for large \(N\). This proves (K16.8).

List the nonzero orders increasingly as \(n_1<n_2<\cdots\). If \(n_{k+1}/n_k\) stayed at least \(1+\eta\) along a subsequence, the interval between these two indices would contain a multiplicative window (K16.8) of a fixed positive relative width: take its midpoint and any \(\delta<\eta/(2+\eta)\). This contradicts (K16.8). Hence \[
\frac{n_{k+1}}{n_k}\longrightarrow1.
\tag{K16.10}
\] These orders can be enumerated exactly by (K14.6)--(K14.7): for a proposed pure order only one finite weight and finitely many partitions occur, and all coefficients are rational. The enumeration is therefore an actual algorithm with provable infinitely many outputs. Neither bounded additive gaps nor a nonzero coefficient at every allowed order is proved.

\subsection{7. What the new bounds do and do not supply to the arithmetic gate}\label{what-the-new-bounds-do-and-do-not-supply-to-the-arithmetic-gate}

The actual all-order finite lower bound at a prime in the table is \[
v_{\mathfrak p}(\partial_i^n\sigma_q(0))
\geq\max\{0,c_i(p)n-s_p(n)/(p-1)\},
\quad c_i(p)=\frac1{(p-1)(p^{h_i}+1)}.
\tag{K17.1}
\] At 11 the three finite checks replace these constants by the larger constants in (K13.1). For the two inert orbit classes specifically requested, the generic constants are

\begin{longtable}[]{@{}
  >{\raggedright\arraybackslash}p{(\columnwidth - 6\tabcolsep) * \real{0.2500}}
  >{\raggedright\arraybackslash}p{(\columnwidth - 6\tabcolsep) * \real{0.2500}}
  >{\raggedright\arraybackslash}p{(\columnwidth - 6\tabcolsep) * \real{0.2500}}
  >{\raggedright\arraybackslash}p{(\columnwidth - 6\tabcolsep) * \real{0.2500}}@{}}
\toprule\noalign{}
\begin{minipage}[b]{\linewidth}\raggedright
Prime class
\end{minipage} & \begin{minipage}[b]{\linewidth}\raggedright
\(c_1(p)\)
\end{minipage} & \begin{minipage}[b]{\linewidth}\raggedright
\(c_2(p)\)
\end{minipage} & \begin{minipage}[b]{\linewidth}\raggedright
\(c_3(p)\)
\end{minipage} \\
\midrule\noalign{}
\endhead
\bottomrule\noalign{}
\endlastfoot
\(p\equiv3\pmod7\) & \(1/((p-1)(p^3+1))\) & \(1/((p-1)(p+1))\) & \(1/((p-1)(p^2+1))\) \\
\(p\equiv5\pmod7\) & \(1/((p-1)(p+1))\) & \(1/((p-1)(p^3+1))\) & \(1/((p-1)(p^2+1))\) \\
\end{longtable}

These are uniform actual-section bounds. They are not claimed sharp. In particular, the accepted DP45/DP53 bare integral-function bounds cannot simply be assigned to this Chern-reduced germ. In the fixed global direction order, their comparison at 11 is

\begin{longtable}[]{@{}
  >{\raggedright\arraybackslash}p{(\columnwidth - 6\tabcolsep) * \real{0.2000}}
  >{\raggedleft\arraybackslash}p{(\columnwidth - 6\tabcolsep) * \real{0.2667}}
  >{\raggedleft\arraybackslash}p{(\columnwidth - 6\tabcolsep) * \real{0.2667}}
  >{\raggedleft\arraybackslash}p{(\columnwidth - 6\tabcolsep) * \real{0.2667}}@{}}
\toprule\noalign{}
\begin{minipage}[b]{\linewidth}\raggedright
Direction
\end{minipage} & \begin{minipage}[b]{\linewidth}\raggedleft
Accepted bare sharp slope from DP53
\end{minipage} & \begin{minipage}[b]{\linewidth}\raggedleft
Actual theta lower slope proved here
\end{minipage} & \begin{minipage}[b]{\linewidth}\raggedleft
Difference of the two stated constants
\end{minipage} \\
\midrule\noalign{}
\endhead
\bottomrule\noalign{}
\endlastfoot
\(\omega_1\) & \(121/1330\) & \(1/20\) & \(109/2660\) \\
\(\omega_2\) & \(11/1330\) & \(1/220\) & \(109/29260\) \\
\(\omega_3\) & \(132/1330\) & \(1/20\) & \(131/2660\) \\
\end{longtable}

The supplied exact locators are \nolinkurl{RETURN_DP_MULT_checkpoint_06_20261121T0420Z/research/DP45_integrated_fragment.tex}, the DP45 theorem, and \nolinkurl{DP53_integrated_fragment.tex}, displayed global label table and its proof paragraph beginning ``For a=3''. Their accepted \(p\equiv3,5\) triples are respectively \((p,p^3,p^2)/((p-1)(p^3+1))\) and \((p^3,p,p^2)/((p-1)(p^3+1))\). The numerical differences in the table are not proved losses of the actual theta section: its optimal finite slopes are still unknown.

There is also a uniform, explicit but much weaker upper bound on each \textbf{nonzero} actual pure jet. For \(n\geq2\), (K14.8) makes \(I_n=2^{n-2}\partial_i^n\sigma_q(0)\) a nonzero rational integer. For every odd prime \(p\), \(p^{v_p(I_n)}\leq|I_n|\). Using (K16.6) gives \[
v_p(\partial_i^n\sigma_q(0))
\leq\frac{\log(6125/16)+(n-2)\log2+\log(n!)
+\frac n2\log\!\left(\frac{\pi e}{nG_i}\right)}{\log p}.
\tag{K17.2}
\] Thus the finitely computable nonzero pure derivatives have both the actual linear lower bounds and an actual \(O(n\log n)\) upper bound. Formula (K17.2) applies at every prime above \(p\) because these derivatives are rational. It supplies no matching linear upper bound, optimal finite slope, or finite-prime analogue of (K16.1).

All mixed actual derivatives are determined by the Sato determinant and are integral away from 2. The positive linear finite bounds proved here concern the three pure directions. AM-K016 supplies nonzero pure derivatives with support ratios tending to one, and an exact archimedean limsup, but it does not ensure nonvanishing of a prescribed mixed harmonic row or charge. No transfer from one pure sequence to a source-selected harmonic image is claimed. Matching linear finite upper bounds, effective support thresholds, and the prescribed-row image comparison remain additional gates.

The full finite data for the translated \(F\) have useful but limited scope: modulo 11 its even derivatives through order 10 in directions 1 and 3, and through order 120 in direction 2, are units at both primes above 11. The exact certificate checks this by polynomial gcd with \(\Phi_7\); it does not convert those finite lists into an asymptotic finite upper bound for \(F\), much less for \(\sigma_q\).

\subsection{8. Primary sources, exact locators, and reproducibility}\label{primary-sources-exact-locators-and-reproducibility}

The literature supplies the general bidifferential, Kleinian, and tau/sigma identities cited here. The CM uniqueness argument, stable-ring estimates, Bell--Legendre constants, algebraic comparison, and growth/support arguments are proved in this note. No claim that integral Hodge eigenvectors give an integral Chern frame is used; the supplied THETA3 C11 and C19 counterexamples and actual-theta method remain consistent with every assertion above.

\begin{enumerate}
\def\labelenumi{\arabic{enumi}.}
\tightlist
\item
  Buchstaber, V. M., Enolskii, V. Z., \& Leykin, D. V. (1997). Hyperelliptic Kleinian functions and applications. \emph{American Mathematical Society Translations, Series 2, 179}, 1--33. \href{https://homepage.mi-ras.ru/~buchstab/download/advmathsci.2.179.1.pdf}{Author-hosted primary PDF}. Used: §1.2, equations (1.7)--(1.9); §2.1, Definition 1 and equations (2.1)--(2.2), pp.7--8; Corollary 2.1.2, equation (2.9), p.10; Theorem 2.2, equations (2.10)--(2.14), p.11. Local retrieved source: \nolinkurl{theta_sources/BEL1997.pdf} and \nolinkurl{.txt}.
\item
  Edelstein, J. D., Gómez-Reino, M., \& Mariño, M. (2001). \emph{Blowup formulae in Donaldson--Witten theory and integrable hierarchies} (arXiv:hep-th/0006113v3). \href{https://arxiv.org/pdf/hep-th/0006113}{Primary PDF}. Used: equations (3.3)--(3.6) for the explicit polar and bidifferential; (3.14)--(3.18) for the quadratic period normalization; (3.24), p.12, for the branch-partition Hessian. Local source: \nolinkurl{theta_sources/Edelstein2001.pdf} and \nolinkurl{.txt}.
\item
  Shramchenko, V. (2004). New Frobenius structures on Hurwitz spaces in terms of Schiffer and Bergmann kernels. \emph{Proceedings of the Institute of Mathematics of the NAS of Ukraine, 50}(3), 1466--1473. \href{https://www.slac.stanford.edu/econf/C0306234/papers/shramchenko.pdf}{Primary proceedings PDF}. Used: §4, p.1468, Schiffer formula and its intrinsic period characterization.
\item
  Nakayashiki, A. (2010). Sigma function as a tau function. \emph{International Mathematics Research Notices, 2010}(3), 373--394. \href{https://arxiv.org/pdf/0904.0846v2}{Author preprint, arXiv:0904.0846v2}. The cited v2 was revised in 2012; the journal article is from 2010. Used: §4, Plücker and Schur definitions and equation (7); §5, the shift by \(g-1\) and leading partition; §7, Theorem 4; §8, Theorem 8, normalized tau/sigma equality. All page references in this note concern the arXiv version.
\item
  Ônishi, Y. (2015). \emph{Hurwitz integrality of power series expansion of the sigma function for a plane curve} (arXiv:1510.03002v1). \href{https://arxiv.org/pdf/1510.03002}{Primary PDF}. Used: equation (2.6) for the differential sign convention; Theorem 1.8 for the general Hurwitz result; Theorem 6.5, printed p.21, for the arithmetic-parameter determinant and exponential corrections; Proposition 6.7 for weights. Local source: \nolinkurl{theta_sources/Onishi2015.pdf} and \nolinkurl{.txt}. The stronger integral statement in this specialization is proved directly from its integer frame and vanishing corrections, not imported from the general theorem without checking powers of 2.
\item
  Ayano, T. (2025). Hurwitz integrality of the power series expansion of the sigma function for a telescopic curve. \emph{Journal of the Mathematical Society of Japan, 77}(3), 833--867. \url{https://doi.org/10.2969/jmsj/90129012}. \href{https://arxiv.org/pdf/2208.10079v2}{Primary preprint, arXiv:2208.10079v2}, revised 14 April 2024. Used only as corroboration: Lemma 7.2 for Schur Hurwitz integrality and Theorem 7.5 for the tau/sigma relation. The reference in §4 to ``Ayano (2024)'' denotes this exact preprint version.
\end{enumerate}

Reproducible artifact paths, all relative to this report's directory:

\begin{longtable}[]{@{}
  >{\raggedright\arraybackslash}p{(\columnwidth - 2\tabcolsep) * \real{0.5000}}
  >{\raggedright\arraybackslash}p{(\columnwidth - 2\tabcolsep) * \real{0.5000}}@{}}
\toprule\noalign{}
\begin{minipage}[b]{\linewidth}\raggedright
Artifact
\end{minipage} & \begin{minipage}[b]{\linewidth}\raggedright
Actual check and scope
\end{minipage} \\
\midrule\noalign{}
\endhead
\bottomrule\noalign{}
\endlastfoot
\nolinkurl{jacobian_theta_direction_jets.py} / \nolinkurl{.json} & Exact algebraic ODE and Bell recurrence; rational pure \(F\) jets through order 8. \\
\nolinkurl{jacobian_theta_direction_jets_mod11.json} & Exact coefficient vectors modulo 11 through order 160 in all three directions. \\
\nolinkurl{jacobian_theta_11_certificate.py} / \nolinkurl{.json} & Regenerates the three exceptional log derivatives, checks saved coefficients and finite unit lists at both residue primes; PASS. \\
\nolinkurl{jacobian_theta_hessian_check.py} / \nolinkurl{.json} & Exact Bolza polynomial identity plus an independent finite theta-sum diagnostic for Hessians and initial \(F\) and \(\sigma_q\) jets; numerical comparisons are not proof premises. \\
\nolinkurl{jacobian_sigma_determinant_certificate.py} / \nolinkurl{.json} & Exhaustive exact weight-20 Plücker determinants, Schur polynomial, and four source-section jets; PASS. \\
\end{longtable}

The imported normalization statements are ALP-LIFT-07 in \nolinkurl{finite_theta_null_normalization.md}, ALP-LIFT-08 in \nolinkurl{trace_one_leading_jet.md}, AM-R-004 in \nolinkurl{../research/jacobian_period_theta.md}, and AM-R-006 in \nolinkurl{../research/trace_one_theta_scalar.md}. This report's own checks are recorded separately in \nolinkurl{kummer_extension_CHECK_EVENTS.tsv}; an author's whole-proof recheck is distinguished there from the parent's independent review.

\clearpage

\section{ALP\_MACHINE: definite comparison and oscillator contribution}\label{alp_machine-definite-comparison-and-oscillator-contribution}

Author: GPT-6 Astra Pro, agent \nolinkurl{/root/definite_oscillator}; effort inherited, not separately exposed. Run start 2026-11-21 19:35 UTC. These are work products for root integration, not a claim that all assigned tasks are finished.

\subsection{Source conventions and carried results}\label{source-conventions-and-carried-results}

Use the supplied run names and statement numbers. The current run carries THETA3 A7, B9, B13, B15, B17, B19, B24 and GCT\_PLETH T-052; it does not rerun their proofs. Exact local source paths follow below. The supplied source archive has reports and source manifests. In this run the needed Lean bodies were recovered centrally from the intake-pinned mirror, with hashes compared against the intake original-author archive metadata wherever that metadata contains the path. Later additional Rallis and TLiftK files have no author-hash entry in that comparison manifest and are identified as pinned-mirror reads in the events. The relevant definitions and proof bodies below were then read directly; no Lean execution or complete dependency audit is claimed.

\begin{enumerate}
\def\labelenumi{\arabic{enumi}.}
\tightlist
\item
  \nolinkurl{RETURN_THETA3_checkpoint_07_20261121T0604Z/documentation/paper_A7.tex}, equation (A7-lift), final two paragraphs: the multiplicity space P is totally definite exactly when the first-list complete CM type functions agree. In the balanced 2n-factor construction this forces all 2n type functions to agree, and the selected balanced tensor is divisor-generated.
\item
  \nolinkurl{documentation/paper_B.tex}, B9, equations (B9a) and following coefficient calculation: the canonical right coefficient on U(2,1) has compact type det(A)d\^{}\{-2\}, equivalently lowest compact weight (1,1;−2). For probability Haar and the compatible compact det test line, the LOCAL rank-zero Bessel integral is one. This is not a global period evaluation.
\item
  Same file, B13: actual MainX relation transfer comes from LawsX. Common geometry, tuple validity, finite support, common real action, double and last integral laws, kernel reduction, terminal factorization and row covariance are distinct hypotheses.
\item
  Same file, B15: literal one-active-place rank3 geometry has global Witt index zero, so abstract MainX length m=2; the remaining definite-side pattern is A2 → S → X2=T⊥Lz, with Lz of opposite sign to L. The small choice S2,T1 gives A2→S2→X2, equal ranks on the definite side. The relevant kernel identity is (B15m).
\item
  \nolinkurl{documentation/paper_B17.tex}, (B17a)--(B17g): the adjoint double lift and its positive inverse are exact analytic constructions. Its upper space is A⊥H of signature (3,1) at every real place, even if P is totally definite. It is therefore not an instantiation of the arithmetic definite upper tower.
\item
  \nolinkurl{documentation/paper_B19.tex}, (B19a)--(B19k) and final paragraph: actual CM-jet model and all period factors; the arithmetic functional is the fixed \nolinkurl{abelian₂}, not a freely chosen linear functional. The proved StdPair kernel constructor requires both ranks even and strictly increasing.
\item
  \nolinkurl{documentation/paper_B24.tex}, (B24a)--(B24e): in the repeated-type definite sector the actual integral Gram inverse is one specified positive scalar times an algebraic matrix. This is stronger than mere analytic existence but does not identify the MainX kernel or terminal masses.
\item
  \nolinkurl{RETURN_GCT_PLETH_checkpoint_09_20261121T0546Z/work/theta/oscillator_charge_report.md}, T-052, (OC.1)--(OC.12), (OC.20)--(OC.22): r=m−2, pair U(r,1)×U(m), polynomial Fock space C{[}X\_ai,Y\_i{]}, and all vacuum shifts.
\end{enumerate}

Primary literature for the last convention: Basile, T., Joung, E., Mkrtchyan, K., \& Mojaza, M. (2020), \emph{Dual pair correspondence in physics: Oscillator realizations and representations}, arXiv:2006.07102v2, §5.2, equations (5.21)--(5.30), printed pp.29--31. This source was opened anew at \url{https://arxiv.org/html/2006.07102v2} and \url{https://arxiv.org/pdf/2006.07102}. The new obstruction below is a calculation from the carried T-052 operators, not a claim of a theorem appearing in Basile et al.

\subsection{AM-DO-001. Explicit repeated-type definite geometry over Q(ζ7)}\label{am-do-001.-explicit-repeated-type-definite-geometry-over-qux3b67}

Let ζ=exp(2πi/7), E=Q(ζ), F=Q(ζ+ζ\^{}\{-1\}), and a=ζ+ζ\^{}\{-1\}. Index the three real embeddings by a↦2cos(2jπ/7), j=1,2,3. The signs of a are (+,−,−). In fact its minimal polynomial is t³+t²−2t−1, so it is an algebraic integer and a unit. Put δ=ζ−ζ\^{}\{-1\}. Its imaginary parts at the complex extensions σ1,σ2,σ3 are all positive.

Carry the CM type Φ=\{1,2,3\} of J=Jac(y²=x⁷−1) from DP\_MULT DP45/DP49. Let v be its odd type function on (Z/7Z)×. The inverse type d(σ)=v(σ\^{}\{-1\}) has signs (+,−,−) on σ1,σ2,σ3, since 1\^{}\{-1\}=1, 2\^{}\{-1\}=4, 3\^{}\{-1\}=5 modulo7. Therefore the rank-one Hermitian form h\_P=−a has exactly the required signs −d(σ)sign Imσ(δ) from THETA3 A6. Take both first-list types equal to v and

\begin{verbatim}
V=(E³,diag(1,1,−a)),  L=Ee3,  A=L⊥=(E²,I2),
P=(E²,δ(−a)I2) as a skew-Hermitian space.
\end{verbatim}

Then V has signature (2,1) at σ1 and (3,0) at the other real places; A is totally positive definite and U(P) is compact at all real places. Its Hermitian sign model is (0,2) at σ1 and (2,0) elsewhere. V and P are globally anisotropic because each has a definite real completion. These are actual global forms over E/F, not separately prescribed incoherent local signatures.

For any nonzero ξ0∈E with ξ0\^{}c=−ξ0 put q=N\_\{E/F\}(ξ0)\textgreater0 at every real place and Lz=−q\^{}\{-1\}L. Its Gram entry is a/q. Taking T=Lz gives X2=T⊥Lz with Gram (a/q)I2: positive at σ1 and negative at the other real places, hence totally definite. For the explicit ξ0=√−7/7, q=1/7, its Gram is 7a I2. Here √−7 can be taken as ζ+ζ²+ζ⁴−ζ³−ζ⁵−ζ⁶; its square is−7 by reduction modulo Φ7. The exact field identities and inverse-type signs are also recorded by \nolinkurl{definite_oscillator_geometry.py} and its JSON receipt. The identity line map satisfies G\_L=−q G\_Lz exactly, as required by the abstract last seesaw relation.

Proof: all signatures follow from the displayed cosines and the inverse type calculation; the line relation is −a=−q(a/q). The compactness criterion and the interpretation as the repeated-type A7 multiplicity space are carried A7 inputs. The polynomial identity for a follows by ζ⁷=1 and ζ≠1. No global theta nonvanishing or arithmetic MainX law is claimed merely from this geometry.

Scope: A7 supplies nonzero canonical wedges in this repeated-type family after suitable inputs and level. If a balanced last pair is requested, it must have the same type v, so this family addresses divisor-generated lines. The alternatives (adjoint T=V versus arithmetic T=Lz) remain different. The former has indefinite upper X4 and an actual analytic inverse; the latter has definite X2 but still requires a global reverse lift, actual finite classes, kernel equations, and terminal masses. This construction does not give their numerical values.

\subsection{AM-DO-002. A canonical compact-weight obstruction for the charged oscillator}\label{am-do-002.-a-canonical-compact-weight-obstruction-for-the-charged-oscillator}

\subsubsection{Precise statement}\label{precise-statement}

Let m≥3, r=m−2, and let ω\_m be the polynomial oscillator of T-052, with the fixed positive-energy convention and dual pair U(r,1)×U(m). On G=U(r,1) let K=U(r)×U(1). Let π\_can be the irreducible canonical holomorphic discrete-series Harish--Chandra module with lowest K type

\begin{verbatim}
κ_can(A,d)=det(A)d^(−r).
\end{verbatim}

For r=2 this is exactly THETA3 B9/B15's canonical ball-surface coefficient, with lowest weight (1,1;−2). There is no nonzero (g,K)-map from ω\_m or any of its algebraic subquotients to π\_can, nor from π\_can to any such subquotient. The same holds after arbitrary determinant character twists of either representation and after a simultaneous contragredient convention. In particular no charge-respecting intertwiner can identify a T-052 component with this canonical 032 component.

For odd m the bare oscillator may live on the metaplectic preimage of K. In that case every displayed map means a (g, K̃)-map on a common cover, with π\_can pulled back; equivalently the obstruction can first be read at the complexified Lie algebra and the traceless compact circle. Its generator H has the integral spectrum computed below, independently of the individual half-vacuum weights. This convention is already part of T-052 and does not assume an unprovided descent of the bare representation.

The statement concerns the canonical automorphic component, not arbitrary maps between entire oscillator representations after changing dual pairs, groups, coefficient systems, or taking additional tensor factors. When comparing precisely the ball surface U(2,1), the T-052 matching rank is m=4,r=2.

\subsubsection{Full proof}\label{full-proof-1}

T-052 (OC.7) gives

\begin{verbatim}
A_ab=Σ_i X_ai ∂_X_bi+(m/2)δ_ab,
A_00=−Σ_iY_i∂_Y_i−m/2.
\end{verbatim}

Let h=diag(1,\ldots,1,−r) in the complexified compact Lie algebra. Its imaginary multiple generates the compact circle c(z)=diag(zI\_r,z\^{}(−r))⊂SU(r,1). Its infinitesimal oscillator action is

\begin{verbatim}
H=Σ_a A_aa−r A_00=N_X+rN_Y+rm.
\end{verbatim}

All monomials in the polynomial Fock model have N\_X,N\_Y≥0. Thus the H spectrum lies in rm+Z\_\{≥0\}. More precisely it is rm+N\_X+rN\_Y. The action is diagonal on a monomial basis. Its restriction to every invariant subspace remains semisimple with this same lower bound; every algebraic quotient has only these eigenvalues as well, since it is the direct sum of its circle weight quotients.

The canonical character on c(z) is

\begin{verbatim}
det(zI_r)(z^(−r))^(−r)=z^{r(r+1)}.
\end{verbatim}

But rm=r(r+2)\textgreater r(r+1). Therefore the lowest K type of π\_can does not occur in ω\_m, or in any subquotient. A nonzero homomorphism from the irreducible π\_can is injective, contradicting absence of its lowest K type. A nonzero homomorphism to π\_can is surjective; exactness of taking an isotypic component for the compact circle (or simply its semisimple weight decomposition on K-finite vectors) would then force its lowest-circle-weight vector to have a preimage of that weight. Again impossible.

Every determinant character is trivial on c(z), since det c(z)=1. Changing the unitary splitting only adds such determinant characters to the fixed pair, so cannot change H. A unitary equivalent realization obtained by changing the polarization preserves the action and its spectrum. Simultaneously passing to contragredients negates both spectra; the distinguished weight −r(r+1) still lies outside the new spectrum, bounded above by −rm. The same absence argument applies. This completes the proof.

\subsubsection{Exact charge calculation}\label{exact-charge-calculation}

On the simultaneous charge-k subspace of T-052, with Q\_i=Σ\_aX\_ai∂\_X\_ai−Y\_i∂\_Y\_i=k, one has N\_X−N\_Y=mk. If b=N\_Y, then

\begin{verbatim}
H=m(k+r)+(r+1)b.
\end{verbatim}

The canonical equality would require m k+(r+1)b=−r, impossible for the T-052 parameters k,b≥0. For m=4,r=2 this reads 4k+3b+8=6. The contradiction is already present before charge selection.

One can also restrict the standard U(r,1) factor to a standard U(2,1) subblock for any m≥4. Its traceless compact circle generator has spectrum N\_X,selected+2N\_Y+2m≥2m≥8, whereas the canonical ball-surface coefficient has weight6. Thus this standard restriction does not evade the obstruction.

\subsubsection{Check scope and consequence}\label{check-scope-and-consequence}

This is a local archimedean obstruction. It survives global choices of finite Schwartz vectors, Haar measures, finite levels, and CM periods. It does not show that every conceivable comparison with Alpöge is impossible. It specifically excludes a route which first identifies T-052's actual charged module with the canonical 032 component while preserving the indicated group action. Q-ALPOGE and Q-ALPOGE-METHOD remain open for changed weights, other coefficient systems, auxiliary oscillator factors, or comparisons not passing through that component.

\subsection{AM-DO-003. Partial Fourier comparison with a definite dual pair, and its exact category obstruction}\label{am-do-003.-partial-fourier-comparison-with-a-definite-dual-pair-and-its-exact-category-obstruction}

\subsubsection{Algebraic statement}\label{algebraic-statement}

Over Z let A be the Weyl algebra on coordinates X\_ai,Y\_i and derivations ∂\_X\_ai,∂\_Y\_i, with {[}∂\_u,u{]}=1. Let A' be the Weyl algebra on X\_ai,Z\_i. The assignments

\begin{verbatim}
X_ai↦X_ai,  ∂_X_ai↦∂_X_ai,
Y_i↦∂_Z_i,  ∂_Y_i↦−Z_i
\end{verbatim}

define a Z-algebra isomorphism F:A→A'. (For the half-integer oscillator constants extend coefficients to Z{[}1/2{]}.) It carries the T-052 complex Lie algebra operators into the definite compact-pair oscillator operators for U(r+1)×U(m), conjugating the first matrix action by diag(I\_r,−1):

\begin{verbatim}
F(A_ab)=Σ X_ai∂_X_bi+(m/2)δ_ab,
F(A_00)=Σ Z_i∂_Z_i+m/2,
F(A_a0)=−Σ X_ai∂_Z_i,
F(A_0a)=−Σ Z_i∂_X_ai,
F(B_ij)=Σ X_ai∂_X_aj+Z_i∂_Z_j+((r+1)/2)δ_ij.
\end{verbatim}

The exact transformed charge is

\begin{verbatim}
F(Q_i)=Σ_a X_ai∂_X_ai+Z_i∂_Z_i+1.
\end{verbatim}

Thus the scalar shift1 is mandatory; charge does not become the unshifted compact column degree.

Integrally, ``Weyl algebra'' here means precisely the algebra generated by coordinate multiplication and the ordinary derivations. It is not the full Grothendieck ring of differential operators over Z. Accordingly the target below is an integral Weyl-algebra module; after tensoring with Q it is the usual algebraic delta D-module. No divided-power differential-operator action, stratification, or crystalline Frobenius is supplied. EGA IV4, Theorem16.11.2 p54 and Remarks16.11.3(ii) p55, gives the exact distinction: the integral differential-operator basis consists of divided derivatives, and ordinary derivations do not generate them in general. This primary text was read anew.

\subsubsection{Full proof of the algebra identities}\label{full-proof-of-the-algebra-identities}

The only nontrivial defining relation is {[}−Z\_i,∂\_Z\_j{]}=δ\_ij; all mixed generators commute as required. The inverse sends Z\_i↦−∂\_Y\_i and ∂\_Z\_i↦Y\_i. For A\_00 use ∂Z\_i·Z\_i=Z\_i∂Z\_i+1. For B\_ij use ∂Z\_j·Z\_i=Z\_i∂Z\_j+δ\_ij. The resulting constants are −m/2+m=m/2 and (r−1)/2+1=(r+1)/2. The off-diagonal signs are both negative and are exactly conjugation by diag(I\_r,−1). The charge formula follows from the same normal ordering. These calculations prove all displayed identities.

\subsubsection{The actual module and why it is not the definite Fock space}\label{the-actual-module-and-why-it-is-not-the-definite-fock-space}

Let Mδ be the A'-module with cyclic vector δ satisfying Z\_iδ=0 and ∂\_X\_aiδ=0. As a Z-module it has basis X\^{}u∂Z\^{}vδ. The map

\begin{verbatim}
U:Z[X,Y]→Mδ,  X^uY^v↦X^u∂Z^vδ
\end{verbatim}

is a bijection and satisfies U(a f)=F(a)U(f). Indeed, multiplication by Y becomes ∂Z, and −Z\_i∂Z\^{}vδ=v\_i∂Z\^{}\{v−e\_i\}δ reproduces ∂Y\_i. X and ∂X are unchanged. This proves exact equivariance and shows that U preserves the ORIGINAL charges, using their transformed shifted operators. In particular Z\_i∂Z\_i acts on ∂Z\^{}vδ by −(v\_i+1), so the transformed charge is deg\_X,column−v\_i as required.

Over C, Mδ is the algebraic D-module generated by the delta section on the subspace Z=0. It is not the ordinary holomorphic polynomial Fock module C{[}X,Z{]}. There is NO nonzero linear map L:C{[}X,Y{]}→C{[}X,Z{]} satisfying L(a f)=F(a)L(f) for every Weyl generator. For f=1, ∂Y\_i1=0 forces −Z\_iL(1)=0, and multiplication by Z\_i is injective on C{[}X,Z{]}, so L(1)=0. The source is cyclic under coordinate multiplications, so L=0.

It is also not a unitary oscillator intertwiner. In positive Fischer/Fock structures Y\_i†=∂Y\_i and (∂Z\_i)†=Z\_i, whereas F(∂Y\_i)=−Z\_i. Thus F fails the adjoint relation. It is a complex symplectic transformation of the algebraic Weyl algebra, not a real metaplectic equivalence of these positive-energy representations. No splitting character changes this Heisenberg adjoint obstruction.

\subsubsection{Consequence for the actual problem}\label{consequence-for-the-actual-problem}

There is an explicit algebraic comparison between the formulas underlying T-052 and a totally definite compact dual pair. Its landing object is a delta D-module and its charge has the displayed shift. It does not map the prescribed Schwartz/Fock theta inputs to the corresponding definite theta inputs, and therefore does not prove any of Alpöge's lift identities, quantitative nonvanishing, or the plethysm complexity claims. Conversely it prevents the overly strong statement that there is no comparison of any kind: there is a precise comparison after changing category.

\subsection{Current state of THETA3 G1}\label{current-state-of-theta3-g1}

The exact arithmetic equation, inherited B15m, is

\begin{verbatim}
volLev(X_j,U_{j,α}) θ_{K,j}(h_{j,α,a},h_{j+1,γ,a'};Φ_{j,P})
  = c_{K,j} abelian₂_{j,α,γ,a,a'}(mvZ(1,s_{j,τ}I)(castP P)).
\end{verbatim}

The functional \nolinkurl{abelian₂} is fixed by the specified arithmetic theta section, its character factors, and its period frame. It cannot be defined post hoc from the analytic left side. The terminal condition \nolinkurl{ClsT.Is.tail} includes top-level covolumes and reciprocal stabilizer masses. Its exact recovered statement has now been read directly in \nolinkurl{OnFramesRdApply.lean}, \nolinkurl{ClsT.Is.tail}:

\begin{verbatim}
volLev(T,U_T) Σ_{β,a} stabOrd(U_T,h_a)^(−1)
   ell_{β,a,G}(low(rd_m(ψ) evaluated on emb(h_a,1)))
  = cP·tail(ψ,G).
\end{verbatim}

The separate \nolinkurl{top} clause says that the translated terminal automorphic function at h\_a k is the functional ell\_\{β,a,G\} applied to the fixed archimedean coefficient J(v,k). Thus the terminal equation must match the actual top representation and its branching map; it is not merely the total mass of the finite class set.

Paths: \nolinkurl{Development/AbelianVarietyTheory/Definitions/LatticeFrames/OnFramesRdApply.lean}, \nolinkurl{ClsK.Core} lines108--113, class masses115--121, \nolinkurl{ClsT.Is.tail}147--150, \nolinkurl{ClsPinD}152--153; \nolinkurl{LevChiEqOne.lean}, \nolinkurl{abelian₂}73--77; \nolinkurl{SplitPart.lean}, \nolinkurl{StdPair}38--49; \nolinkurl{KernelPartLawsOfColumnPins.lean}, \nolinkurl{clsKL_is}22--52; \nolinkurl{ThetaKernelSplitPlacesEqJetClassicalTheta.lean}, \nolinkurl{thetaKernel_eq_jet_thC}46--71.

Approaches inspected:

\begin{itemize}
\tightlist
\item
  Repeated CM type supplies an actual totally definite P, and AM-DO-001 makes the forms explicit over the same cyclotomic field as the main Jacobian. It leaves the finite class matrices and global reverse lift unconstructed.
\item
  The exact adjoint double lift B17 has an indefinite upper space. Repeated types do not fix that signature; this is a carried source conclusion.
\item
  The B15 rank-one terminal alternative has the required upper signature and locally compatible archimedean representations, but neither finite local occurrence for the chosen global data nor the global reverse realization is supplied.
\item
  B24 makes the actual kernel inverse algebraic up to one retained scalar in this sector. Algebraicity/proportionality of an inverse is not equality with the fixed MainX functional, and cannot set terminal masses.
\item
  The StdPair constructor only covers strictly increasing even ranks. It cannot directly be applied to the natural rank2/3 analytic pair, nor to the small equal-rank definite side. This is a limitation of that constructor, not a proof that abstract MainX is impossible.
\item
  Arbitrarily rescaling a test vector to make one Bessel period equal1 does not solve the common kernel equations; that would lose the fixed input and uniformity.
\end{itemize}

Thus question2 is not settled by this contribution. Its current gate is simultaneous construction of the chosen global two-step theta representation and finite arithmetic class/jet matrices, with exact covolumes and terminal masses, for the concrete repeated-type data. No pairing value satisfying MainX has been computed here.

\subsection{AM-DO-004. A simultaneous sign obstruction to the common positive arithmetic pin}\label{am-do-004.-a-simultaneous-sign-obstruction-to-the-common-positive-arithmetic-pin}

\subsubsection{Statement with all hypotheses}\label{statement-with-all-hypotheses}

Let E/F be a CM extension and let X0,X1,X2 be nonzero Hermitian spaces, each definite at every real place, with signs ε\_j(τ)∈\{+1,−1\}. On the oriented matrix space W\_j=Hom\_E(X\_j,X\_\{j+1\}) use the exact source form

\begin{verbatim}
h_j(x,y)=tr(G_j^(−1) y* G_{j+1} x),
B_j(x,y)=Tr_{E/F}(ξ0 h_j(x,y)),  ξ0^c=−ξ0.
\end{verbatim}

A \textbf{label-preserving positive transport} means: identify each abstract matrix group with the corresponding positive definite one, replace all node Gram matrices by positive definite ones, use one common ξ0 and one common additive character ψ\_b(x)=ψ(bx) (b∈F×), and identify each oriented matrix oscillator space by an E-linear map, intertwining the fixed two group actions and carrying its Heisenberg form to the new form with that common additive character. No conjugation of the E action, exchange with a contragredient, or separate additive character for different edges is allowed in this definition.

A necessary condition for such a transport is

\begin{verbatim}
ε_j(τ) ε_{j+1}(τ)=sign τ(b)   for j=0,1,
\end{verbatim}

and hence ε0(τ)=ε2(τ) at every real place. More generally in a chain all even-indexed node signs must agree and all odd-indexed node signs must agree.

For the literal rank-zero MainX geometry of a one-active-place compact 032 ball surface, this necessary condition fails at every compact real place. Consequently its signed definite side cannot be transported to the common all-positive \nolinkurl{PinX}/\nolinkurl{StdPair} geometry by label-preserving positive transports, even after one common change of additive character. This does not refute the abstract signed \nolinkurl{MainX} construction.

\subsubsection{Exact source locator}\label{exact-source-locator}

\nolinkurl{Formalism/LocalToGlobal/IndefiniteReal/LocalNewformFunction.lean}: \nolinkurl{hO} lines25--26, \nolinkurl{symplO}28--29, \nolinkurl{heisO}34--35, \nolinkurl{actO}37--38, and \nolinkurl{IsWeilDatumO.om_heis}62--63. The actual formula uses the SOURCE Gram inverse; this is why a scalar edge ratio is q\_\{j+1\}/q\_j rather than its product.

\nolinkurl{Development/AbelianVarietyTheory/Definitions/LatticeFrames/SplitPart.lean}: \nolinkurl{StdPair.gS}, \nolinkurl{gT}, \nolinkurl{cc}, \nolinkurl{xi}38--49. \nolinkurl{LevChiEqOne.lean}: \nolinkurl{PinX.xs} and \nolinkurl{vc} specify Cartan involution1, orthonormal frame1 and standard lattice for every node, all with the same ψ. The abstract \nolinkurl{DefSp.Is.pm}, in \nolinkurl{JointSmooth.lean}, instead allows signs+1 or−1 separately at every real place. \nolinkurl{PlaceLaws.lean}, \nolinkurl{WeilOK.tau}165--166, imposes the opposite-line Gram equation.

\subsubsection{Full proof}\label{full-proof-2}

Fix a real place τ. Write G\_j=ε\_j P\_j with P\_j positive definite. For x≠0,

\begin{verbatim}
h_j(x,x)=ε_j ε_{j+1} tr(P_j^(−1) x* P_{j+1} x).
\end{verbatim}

The trace on the right is positive: it is the sum of squared absolute values of the matrix P\_\{j+1\}\^{}\{1/2\} x P\_j\^{}\{−1/2\}. Thus h\_j is definite of sign ε\_j ε\_\{j+1\} on the complex vector space W\_j.

Let h\_j\^{}+ be the Hermitian form for the transported positive pair and T\_j the required complex-linear map. Equality of the real symplectic forms with additive-character scaling says

\begin{verbatim}
Tr_{C/R}(ξ0 h_j(x,y))
  = b Tr_{C/R}(ξ0 h_j^+(T_jx,T_jy)).
\end{verbatim}

Because T\_j is complex linear, replacing x by ix also gives equality of the complementary real part. Since ξ0 is nonzero purely imaginary, the two real equalities imply

\begin{verbatim}
h_j(x,y)=b h_j^+(T_jx,T_jy).
\end{verbatim}

For x=y≠0 the right side has sign sign b. Therefore sign b=ε\_j ε\_\{j+1\}. Multiplying the two edge equalities gives ε0ε2=(sign b)\^{}2=1, proving the necessary condition. Repeating this argument along the chain proves the parity assertion. The same proof covers arbitrary positive Gram matrices; it is not limited to scalar forms.

For the ball surface, write V=A⊥L and X2=T⊥Lz in the literal m=2 last-step pattern. At a compact real place τ, V is definite, so A and L have the same sign. The exact line relation

\begin{verbatim}
τ_line* G_L τ_line=−N_{E/F}(ξ0) G_Lz
\end{verbatim}

has a totally positive norm factor, and therefore Lz has the opposite sign to L. Since X2 is definite and contains Lz, it has Lz's sign. Hence X0=A and X2 have opposite signs at τ. This contradicts the necessary condition. Such a compact real place exists because the 032 construction under discussion has {[}F:Q{]}\textgreater1 and only one active place.

This argument does not say that U(qI) and U(I) are different algebraic groups; they are literally the same matrix group. It says that the E-linear symplectic embeddings and their common Heisenberg character cannot simultaneously be the positive ones. Splitting-character twists multiply group actions by determinant characters but do not change the Heisenberg commutator form, so they cannot remove this obstruction. A common change ξ0↦cξ0 is already absorbed into b. Changing the middle sign multiplies both edge signs by−1 and still leaves them opposite. Rescaling the original definite completion of V changes both A and L together and again leaves the endpoint sign opposition unchanged. This completes the proof.

\subsubsection{The scalar equations after fixing a Hermitian middle-space conversion}\label{the-scalar-equations-after-fixing-a-hermitian-middle-space-conversion}

The endpoints of AM-DO-001 have q0=1 and q2=a/q, where q=Nξ0\textgreater0. Its middle object P is initially skew-Hermitian, with Gram δ(−a)I2. A comparison with the oriented Hermitian \nolinkurl{hO} model must first specify the skew-to-Hermitian conversion, the source-dual convention and the group labels. In particular one must not silently declare the middle Hermitian Gram to be −aI2 while independently choosing ξ0. Suppose that actual conversion gives a scalar Hermitian Gram q1=s∈F×. The edge ratios are then

\begin{verbatim}
r0=s,   r1=a/(qs).
\end{verbatim}

An E-linear scalar matrix map x↦c\_jx into the standard positive model would have to satisfy

\begin{verbatim}
b N_{E/F}(c_0)=s,
b N_{E/F}(c_1)=a/(qs).
\end{verbatim}

Their product gives b²N(c\_0c\_1)=a/q, independently of s. The left side is totally positive, while a/q has two negative real embeddings. Thus there is no solution. If a compatible convention separately proves s=−a, the displayed equations specialize to bN(c0)=−a and bN(c1)=−1/q; that specialization is not an additional claim here. The obstruction itself uses only the endpoints and the necessary real sign of a norm, so neither the unresolved middle conversion nor possible finite norm obstructions affect this exclusion.

\subsubsection{Meaning for the open questions}\label{meaning-for-the-open-questions}

This refutes one concrete attempt to solve THETA3 G1: transport both edges of its literal two-step definite side into the already-built positive-identity arithmetic kernel constructor while keeping the E labels, charges and a common additive character. It leaves the signed arithmetic kernel/covolume equations themselves open. They would require a new signed CM-type comparison (or a different global realization), with explicit conjugations, characters and periods. The theorem does not obstruct Q1's already-positive abstract definite tower; it is specific to importing the 032 one-active-place geometry into that same positive pin.

\subsubsection{Extension allowing whole-edge E/F conjugation}\label{extension-allowing-whole-edge-ef-conjugation}

The same route still fails if each edge identification is allowed to be either E-linear or E-conjugate-linear globally over F, with the corresponding conjugation of its group labels. Write η\_j=+1 in the former case and η\_j=−1 in the latter. Each η\_j is independent of the real place: the only nontrivial E/F automorphism is complex conjugation at all of them. The necessary condition becomes

\begin{verbatim}
sign τ(b)=η_j ε_j(τ) ε_{j+1}(τ),
ε0(τ) ε2(τ)=η0η1.
\end{verbatim}

To prove the first equality, the E-linear case is already proved. In the conjugate-linear case, put T(ix)=−iT(x) in the two symplectic equalities. They give h\_j(x,y)=−b overline(h\_j\^{}+(T x,T y)), so the sign is −sign b. Multiplying the two edge conditions proves the second equality.

For the one-active-place ball geometry, ε0ε2 is +1 at the active place because A and Lz are both positive, and −1 at every compact place because Lz has the opposite sign to A there. It is therefore not constant, contrary to η0η1. Thus whole-edge global conjugation does not repair the obstruction either. Independent conjugations at selected real places are not E/F-semilinear maps of the rational oscillator space; no claim about constructions that replace that global space and its rational theta lattice is made here.

\subsection{Additional source and proof checks}\label{additional-source-and-proof-checks}

After recovering the actual Lean files, the agent read all of \nolinkurl{OnFramesRdApply.lean} through the ClsK/ClsT definitions and conversion lemmas; all of \nolinkurl{ValLastX.lean}; \nolinkurl{StdPair} and \nolinkurl{xiN} in \nolinkurl{SplitPart.lean}; \nolinkurl{TwDat}, \nolinkurl{abelian₂}, \nolinkurl{PinX}, \nolinkurl{PinColT} in \nolinkurl{LevChiEqOne.lean}; \nolinkurl{DefSp}, \nolinkurl{KerDat}, \nolinkurl{BesAt}, \nolinkurl{JacAt} in \nolinkurl{JointSmooth.lean}; \nolinkurl{WeilOK} in \nolinkurl{PlaceLaws.lean}; and the complete proof bodies \nolinkurl{clsKL_is} and \nolinkurl{thetaKernel_eq_jet_thC}. The dependency chain has not been compiled or audited in full.

Root's AM-R-002 was independently whole-read and checked: content ideal and local detector denominator, dual norm and quotient by kernel, scaling/product-formula invariant, arithmetic-degree sign, actual normalized coefficient matrix and stabilizer masses all pass. Its formula must be applied to the actual terminal lift vector and actual invariant lattice; it does not supply them for the present signed seesaw.

\subsection{AM-DO-005. An explicit charge-preserving local bicanonical comparison}\label{am-do-005.-an-explicit-charge-preserving-local-bicanonical-comparison}

This gives a nonzero comparison after changing the target coefficient weight. It is included to make the exact scope of AM-DO-002 clear; it does not turn the canonical 032 class into an Alpöge arithmetic column.

\subsubsection{Data, splitting and polarization}\label{data-splitting-and-polarization}

Set m=4, r=2 in T-052. The real dual pair is U(2,1)×U(4) acting on the real symplectic space underlying C\^{}\{2,1\}⊗C⁴, with the T-052 convention ω=Im(h\_\{2,1\}⊗h\_4). Use its positive compatible complex structure i on the positive summand and −i on the negative summand. The polynomial Fock space is C{[}X\_ai,Y\_i{]}, a=1,2 and i=1,\ldots,4. On K\_G×U(4) the bare vacuum character has square

\begin{verbatim}
χ0²=(det A)^4 d^(−4) det g,
\end{verbatim}

so it is (det A)²d\^{}(−2)χ\_W(g), with χ\_W²=det g on the indicated metaplectic pullback. These are the exact T-052 conventions, not newly inferred global splitting data.

Twist the G=U(2,1) action by det\_G\^{}(−2), and untwist the compact U(4) factor by χ\_W\^{}(−1), just as T-052 untwists its ordinary column action. Both combined actions then descend in this selected model. The former is an honest determinant character. In the archive's \nolinkurl{splitFacO} convention, changing the character attached to the compact member by \(\kappa(z)=(z/\bar z)^2\) produces this det\_G\^{}(−2) factor: if \(x/\bar x=\det g\), then κ(x)\^{}(−1)=det(g)\^{}(−2). This states the splitting CHANGE relative to the fixed model; it does not invent a global absolute character or its finite components. The polarization is unchanged.

\subsubsection{Explicit lowest vector and intertwiner}\label{explicit-lowest-vector-and-intertwiner}

Put Δij=X\_\{1i\}X\_\{2j\}−X\_\{1j\}X\_\{2i\} and f=Δ12Δ34. Then

\begin{verbatim}
C_1 f=C_2 f=0,    Q_i f=f for every i,
(N_X,N_Y)=(4,0),  α=(2,2), k=1, b=0.
\end{verbatim}

Each column occurs once and f contains no Y, which proves the first two assertions directly. Under U(2), f transforms as det²; with the vacuum shifts, its bare lowest K\_G weight is (4,4;−2). After the determinant twist it becomes (2,2;−4), the bicanonical compact type from THETA3 B14. Its restriction to the traceless compact circle has weight12, consistently with H=8+N\_X+2N\_Y. This does not contradict AM-DO-002, which excludes weight6.

The polynomial f expands into four distinct squarefree monomials, each with coefficient±1. Under the Fischer norm its squared norm is therefore4. Let e=f/2. The G-cyclic module generated by e lies in the T-052 irreducible lowest-weight Howe summand indexed by α=(2,2),b=0. This irreducibility is a carried T-052 (OC.6),(OC.20)--(OC.22) input, not an assertion made for an arbitrary lowest vector. After the determinant twist it is the irreducible bicanonical holomorphic discrete series π\_2.

Choose a unit lowest vector v\_2 in a fixed π\_2 model. The explicit intertwiner on algebraic vectors is

\begin{verbatim}
I(u·e)=u·v_2,   u∈U(gl_3(C)),
\end{verbatim}

where the action on the source includes the det\^{}(−2) twist. This is well-defined: the carried irreducible lowest-weight identification gives the same annihilator ideal of the fixed lowest type in the two modules; equivalently both are the same irreducible lowest-weight quotient. It is a bijection. Its unit-vector normalization fixes its scalar uniquely, and invariance of the positive forms makes it unitary on completion. Because the G action commutes with every Q\_i, its whole domain has charge(1,1,1,1). Thus the comparison preserves the prescribed charge sector.

An actual nonzero coherent target may be obtained from any nonzero canonical wedge ξ supplied by A7: ξ² is a nonzero section of K\_X², and its automorphic coefficient generates a copy of π\_2. This observation specifies the archimedean recipient. No compatibility with the finite Weil action or membership in the global U(4) theta image follows from it.

\subsubsection{The multiplicity factor cannot be discarded}\label{the-multiplicity-factor-cannot-be-discarded}

The charge-one lowest space is the multilinear component of S\_(2,2)(C⁴). It has basis

\begin{verbatim}
f1=Δ12Δ34,  f2=Δ13Δ24,
\end{verbatim}

with the third pairing f3=Δ14Δ23 satisfying f1−f2+f3=0. Their Fischer Gram matrix is {[}{[}4,2{]},{[}2,4{]}{]}. The ordinary S4 action is its irreducible Specht module {[}2,2{]}, of dimension2. This is also T-052 (OC.4) at n=1,λ=(2,2,0,0), with its vacuum untwist. Consequently the full charge-selected local comparison retains

\begin{verbatim}
π_2 ⊗ [2,2],
\end{verbatim}

rather than mapping the S4 multiplicity to a scalar. For example choose the two orthonormal lowest vectors e1=f1/2 and e2=(f2−f1/2)/√3 and send u·e\_j to (u·v\_2)⊗w\_j, where w1,w2 are the corresponding orthonormal basis of {[}2,2{]}. This fixes the whole G×S4 intertwiner; S4 acts through its exact Gram-compatible matrices on the second factor.

For all n≥1 this is the T-052 parameter family λ=(n+1,n+1,n−1,n−1), since α=(2,2),k=1,b=0. The ordinary S4 comparison may include the source's sign\^{}(n+1), but {[}2,2{]} is self-conjugate and hence unchanged up to its specified Specht isomorphism. This parameter observation is carried from (OC.1) and (OC.4), not a new plethysm computation.

\subsubsection{A normalization with an actual local value}\label{a-normalization-with-an-actual-local-value}

For the compact subgroup U(2)⊂U(2,1), e1 has character det² in the twisted representation. With probability Haar measure and the compact test representation det²,

\begin{verbatim}
∫_{U(2)} <e1,π_2(diag(A,1))e1> det(A)^(−2) dA=1.
\end{verbatim}

The integrand is identically1. This is precisely the q=2 specialization of the carried THETA3 B9/B14 local calculation. It is not the global kernel/covolume or terminal-mass value requested in question2.

\subsubsection{Consequence}\label{consequence}

A blanket assertion that no charge-preserving comparison with a ball-quotient representation exists would be false: this is one, with the dual pair, covering/splitting adjustment, polarization, vector, multiplicity factor and normalization given explicitly. It uses a bicanonical coherent section. THETA3 B14 shows why that differs from the ordinary canonical cohomology class and why a diagonal cup-product argument does not return it to that class. Finite-place compatibility, actual global theta occurrence, arithmetic columns and controlled period costs remain separate. Q-ALPOGE and Q-ALPOGE-METHOD therefore remain open in their original global form.

\subsection{AM-DO-006. Full charge torus obstruction to a polarized compact tensor comparison}\label{am-do-006.-full-charge-torus-obstruction-to-a-polarized-compact-tensor-comparison}

\subsubsection{Exact scope and statement}\label{exact-scope-and-statement}

Keep the T-052 positive polarization and m≥3, r=m−2. Let T=(U(1))\^{}m be the individual column-charge torus. Its action on the creation space E\_m has character multiset

\begin{verbatim}
E_m = ⊕_{i=1}^m (C_{e_i}^{⊕r} ⊕ C_{−e_i}),
\end{verbatim}

where e\_i are the independent coordinate characters. In particular dim\_C E\_m=m(m−1). Vacuum and splitting characters do not occur in this formula: these are the commutator weights on creation operators, so a scalar twist of the whole oscillator cancels.

A \textbf{polarized full-charge tensor comparison} to a totally definite dual pair U(a)×U(b) means a real symplectic isomorphism of the full Heisenberg spaces, preserving the specified positive compatible polarizations, together with a torus homomorphism T→U(a)×U(b), such that the isomorphism intertwines the individual charge torus. The two compact factors may be dualized, and either common orientation of the symplectic form may be used. These changes only dualize one or both defining representations and do not change their tensor-product form. A comparison of this kind requires ab=m(m−1).

There is no such comparison with a,b≥2. In particular, even when dimensions happen to agree, it cannot identify the T-052 full charged oscillator with a single one of Alpöge's nonzero even-rank, totally definite tensor-pair oscillators while keeping all individual charges and the indicated positive polarizations.

This statement does not concern a map only on a chosen simultaneous-charge eigenspace, a map of Harish--Chandra modules which does not intertwine the Heisenberg action, changes of the real form or polarization not identified in the hypothesis, extra oscillator factors, or a rank-one compact factor. It is a precise obstruction to one proposed direct comparison, not a proof that Q-ALPOGE has no solution by any method.

\subsubsection{Elementary support lemma}\label{elementary-support-lemma}

Let m≥3 and S=\{±e\_1,\ldots,±e\_m\}⊂R\^{}m. If A and B are nonempty finite sets of vectors with A+B=S, then A or B is a singleton.

Proof. Suppose each set has at least two elements. Choose distinct α,α′∈A and β,β′∈B and form

\begin{verbatim}
x=α+β,  y=α′+β,  z=α+β′,  w=α′+β′.
\end{verbatim}

All lie in S, and x+w=y+z. Adjacent vertices in this displayed rectangle differ, because the selected α's and β's differ. If x=w, then x=(y+z)/2; y≠z, since otherwise x=y would follow. This contradicts the fact that every element of S is an extreme point of its convex hull. The same argument excludes y=z. Thus x,y,z,w are four distinct elements of S.

The sum of two distinct elements of S is either zero, when they are opposite, or a vector with exactly two nonzero coordinates equal to ±1. In the latter case the sum uniquely determines the unordered pair of vertices. Since the two pairs \{x,w\} and \{y,z\} are disjoint, their common sum must therefore be zero. Hence

\begin{verbatim}
α+α′+β+β′=0.
\end{verbatim}

Now fix α,α′ and β and vary β′ over B\{β\}. This equality determines β′ uniquely. It follows that B has at most two elements. By symmetry A has at most two elements. But then A+B has at most four elements, whereas S has 2m≥6. This contradiction proves the lemma. The proof uses no lattice or integrality hypothesis. ∎

\subsubsection{Full proof of the representation obstruction}\label{full-proof-of-the-representation-obstruction}

Preservation of the positive polarization identifies the complex creation spaces equivariantly for T. Restrict the defining a- and b-dimensional compact-factor representations along the claimed torus homomorphism. Compactness makes each a direct sum of characters. Write their character multisets as \{α\_1,\ldots,α\_a\} and \{β\_1,\ldots,β\_b\}, absorbing any contragredient into the signs of these characters. The creation space of the definite tensor pair then has character multiset

\begin{verbatim}
{α_p+β_q : 1≤p≤a, 1≤q≤b}.
\end{verbatim}

All multiplicities are nonnegative, so its support is exactly A+B, where A and B are the two character supports. The required isomorphism forces A+B=S. By the support lemma one factor, say A, consists of one character α. That factor representation therefore consists of a copies of α. Every character multiplicity in its tensor product is divisible by a. But every negative character −e\_i in E\_m has multiplicity1. Consequently a=1. If B is the singleton, the same argument gives b=1. Either conclusion contradicts a,b≥2.

This proof also shows why scalar splitting shifts do not help. They only multiply an entire oscillator representation by a character; the adjoint weights of the creation operators remain unchanged. If one allows an affine common shift of all creation weights in an intermediate notation, it can be absorbed into one factor's α's before applying the same argument. Passing to a common cover to define half-integral vacuum characters similarly leaves the proof intact. ∎

\subsubsection{Algebraic geometry and K-theory form}\label{algebraic-geometry-and-k-theory-form}

For the split torus T\_Z=G\_m\^{}m the creation module is the free T\_Z-equivariant vector bundle

\begin{verbatim}
E_m = ⊕_i (L_{e_i}^{⊕r} ⊕ L_{−e_i})
\end{verbatim}

on Spec Z, with character class rΣ\_i t\_i+Σ\_i t\_i\^{}(−1) in K\_0\^{}\{T\_Z\}(Spec Z)=Z{[}t\_1\^{}±1,\ldots,t\_m\^{}±1{]}. The theorem says that this positive character cannot be a product of two positive character classes of ranks at least2. The adjective positive means actual representations or vector bundles; virtual K-theory factorizations with negative coefficients are not excluded. Thus an ordinary Grothendieck-group equality is weaker than the tensor-bundle identification needed for the positive oscillator construction.

The exact Grothendieck-framework locator for the weight decomposition is SGA3, Exposé I, Proposition4.7.3; Corollary4.7.3.1 gives its graded-algebra/equivariant-affine-scheme form. Both were read anew in the French recomposed text. The relevant archive framework in the supplied materials is cote133 scan006 (pre-λ Grothendieck groups and universal identities in classifying topoi), while the exact mathematical input and new statement are T-052 (OC.3), (OC.7)--(OC.12) and AM-DO-006 here. No assertion that SGA3 or cote133 proves this factorization obstruction is made.

\subsection{AM-DO-007. Canonical rank cap for a compact middle source}\label{am-do-007.-canonical-rank-cap-for-a-compact-middle-source}

\subsubsection{Statement}\label{statement-6}

For the real dual pair U(n,1)×U(d), let the second factor be totally definite and n≥1, d≥1, q≥1. In the positive-energy orientation, if the irreducible holomorphic coefficient representation of the qth canonical power occurs as a subquotient of the oscillator, with arbitrary determinant splitting twists, then

\begin{verbatim}
d≤q(n+1).
\end{verbatim}

The same necessary bound holds in the simultaneous contragredient convention. In the opposite oscillator orientation with the original holomorphic target, no occurrence is possible. As in AM-DO-002, the comparison is made on a common metaplectic cover when necessary.

For the unchanged canonical coefficient of the 032 ball surface, n=2 and q=1, the compact source must have d≤3. Therefore the first definite edge A\_2→S\_d in a literal MainX realization cannot at the same time satisfy the existing \nolinkurl{StdPair.rk} conditions d\textgreater2 and Even d.~Enlarging the middle totally definite rank cannot put that canonical realization into this particular arithmetic constructor.

\subsubsection{Full proof}\label{full-proof-3}

Here n and d are independent; no T-052 relation between them is assumed. The general U(n,1)×U(d) oscillator formulas are the same creation-annihilation formulas carried in T-052 (OC.7), with a=1,\ldots,n and i=1,\ldots,d.~Equivalently they are Basile et al.~(2020), §5.2, equations (5.21)--(5.30), specialized to q=1 negative direction and d compact columns. The positive compact diagonal operators have vacuum shift d/2 each, and the negative diagonal operator is −N\_Y−d/2. Consequently the generator of diag(zI\_n,z\^{}(−n))⊂SU(n,1) is

\begin{verbatim}
H=N_X+nN_Y+nd,
\end{verbatim}

with every eigenvalue at least nd. The qth canonical coefficient has compact character det(A)\^{}q d\_0\^{}(−nq), where d\_0 denotes the last U(1) entry, and therefore has the weight qn(n+1) for this circle. Its lowest vector must appear in any nonzero subquotient realization. Exactness of compact weight decomposition, as proved in AM-DO-002, gives nd≤qn(n+1). Division by n\textgreater0 proves the bound.

Every determinant splitting twist is trivial on this traceless circle. Reversing the entire oscillator orientation replaces its spectrum by eigenvalues at most−nd, which excludes the unchanged positive canonical lowest weight. Simultaneously dualizing the target negates its lowest weight too and gives the same inequality. These cover the two definite signs of the auxiliary Hermitian space with fixed additive-character convention.

For the ball, d≤3 follows. The actual source \nolinkurl{SplitPart.lean}, \nolinkurl{StdPair.rk} lines43--44, requires 0\textless rank A\textless rank S and both ranks even; since rank A=2, it requires d≥4. The inequalities are incompatible. This proves the last assertion. ∎

\subsubsection{Consequence and boundary}\label{consequence-and-boundary}

This obstruction is additional to the simultaneous sign obstruction AM-DO-004. Even at the active real place alone, the unchanged canonical coefficient prevents using a larger even middle group to meet the strict kernel-constructor range. A signed equal-rank kernel theorem, a different arithmetic constructor, or a change of coherent coefficient is therefore needed for that route to G1. The bound is necessary and not sufficient: it does not prove finite local occurrence, global nonvanishing, the kernel identity, or the terminal mass equation when d≤3. At q=2, d=4 is allowed by this bound and the explicit AM-DO-005 lowest vector realizes the corresponding local bicanonical component.

\subsection{T-053, the Fourier comparison, and the precise geometric category}\label{t-053-the-fourier-comparison-and-the-precise-geometric-category}

Carry T-053, \nolinkurl{work/theta/charged_incidence_rees_report.md}, (IR.1)--(IR.5), (IR.12)--(IR.14). Its incidence ring is B=Q{[}X,Y{]}/(R\_a), R\_a=Σ\_iX\_aiY\_i, and its actual frame Rees algebra is ⊕\emph{\{k≥0\}B}(k,\ldots,k). The source gives the generator maps, multiplication and quotient scheme; these are settled results and are not rederived here. In particular its equal-charge construction is an algebraic torus-weight operation in SGA3 I4.7.3 and4.7.3.1.

The partial Fourier map AM-DO-003 has the exact additional formulas

\begin{verbatim}
F(R_a)=Σ_i X_ai∂_Zi,
F(C_a)=−Σ_i Z_i∂_Xai,
\end{verbatim}

where C\_a=Σ\_i∂\_Xai∂\_Yi. These follow by substitution into the already-proved Weyl-algebra isomorphism. They exchange the displayed multiplication constraints and differential constraints with off-diagonal compact Lie operators. They do not carry them to multiplication by polynomials in the ordinary positive definite coordinates X,Z: indeed

\begin{verbatim}
[F(R_a), Z_j]=X_aj≠0,
\end{verbatim}

whereas every such coordinate multiplication operator commutes with Z\_j. This proves that this specific Fourier map does not give a morphism to the ordinary definite polynomial incidence scheme preserving the positive coordinate polarization. It remains an exact Weyl-module comparison into the delta realization, and the original Rees algebra still has its own carried scheme interpretation. No assertion that an arbitrary algebraic comparison between selected Rees pieces is impossible follows.

Likewise AM-DO-006 concerns the full creation-space tensor factorization, not a factorization of a single equal-charge quotient B\_(k,\ldots,k). It therefore does not refute T-053, its positive row-weight counts, or its exact remaining highest-type and Specht projectors. It excludes the proposed direct identification of the whole charged oscillator with one nontrivial compact tensor-pair oscillator under the named polarization and charge requirements.

Primary framework locators read in this run: Grothendieck, EGA IV4 (1967), Theorem16.11.2 and Remarks16.11.3(ii), pp54--55, \url{https://www.iut-theory.org/divers/ega44.pdf}; Demazure--Grothendieck, SGA3, Exposé I, Proposition4.7.3 and Corollary4.7.3.1, French recomposed edition, \url{https://www.iut-theory.org/divers/SGA3I.pdf}. The former distinguishes full differential operators from the integral Weyl algebra; the latter identifies actual torus-equivariant modules and schemes with gradings. Neither identifies these realizations with an integral theta lift, a Dieudonné crystal, an ℓ-adic sheaf with a specified Frobenius, or a motive. Those require additional structures and actual comparison maps.

\subsection{AM-DO-008. The actual kernel equation as period-rescaling and lattice-membership tests}\label{am-do-008.-the-actual-kernel-equation-as-period-rescaling-and-lattice-membership-tests}

This is a reduction of the already displayed MainX equation, not a computation of its missing theta entries. It makes precise which finite calculations would certify a selected arithmetic column, and which additional identity is needed for the source's full \nolinkurl{ClsK.Core}.

\subsubsection{Polynomial and period equations}\label{polynomial-and-period-equations}

Fix one edge j, all finite charts and class representatives, the actual arithmetic \nolinkurl{abelian₂} functionals, and the corresponding analytic kernel. Index a finite list of tests by i=(α,γ,a,a′,µ), where P\_µ is a monomial polynomial. Put

\begin{verbatim}
Θ_i = volLev(X_j,U_j,α) θ_K,j(h_j,α,a,h_j+1,γ,a′; Φ_j,Pµ),
A_i = abelian₂_j,α,γ,a,a′(castP P_µ),
d_i,τ = degree of P_µ in the variables at τ.
\end{verbatim}

The exact source equation restricted to these tests is

\begin{verbatim}
Θ_i = c_j (∏_τ s_j,τ^(d_i,τ)) A_i,              (DO8.1)
\end{verbatim}

where c\_j=cK\_j≠0 and every s\_j,τ=sc\_j,τ≠0. No covolume has been omitted from Θ\_i. If the c\_j or s\_j,τ are already fixed by the arithmetic pin, (DO8.1) must use those values.

For freely allowed complex rescalings c\_j,s\_j,τ, a finite system (DO8.1) is solvable if and only if:

\begin{enumerate}
\def\labelenumi{\arabic{enumi}.}
\item
  A\_i=0 exactly when Θ\_i=0;
\item
  on the nonzero indices, with r\_i=Θ\_i/A\_i and u\_i=(1,(d\_i,τ)\_τ), one has

\begin{verbatim}
∏_i r_i^(v_i)=1
\end{verbatim}

  for every integer relation Σ\_i v\_i u\_i=0.
\end{enumerate}

In particular all ratios belonging to the same multidegree must be identical. For a single homogeneous coefficient space, every two-by-two minor Θ\_i A\_l−Θ\_l A\_i must vanish, together with the stated support condition; a nonzero analytic period alone does not impose these identities.

Proof. The scalar matrix substitution in \nolinkurl{mvZ(1,sI)} multiplies a monomial by ∏\_τs\_τ\^{}(d\_i,τ), proving (DO8.1). Nonzero scalars prove condition1; multiplying the equations according to an integer relation proves condition2. Conversely let M be the integer matrix with rows u\_i for the nonzero indices. Choose integer unimodular matrices U,V with UMV=D in Smith form, whose nonzero diagonal entries are δ1,\ldots,δh\textgreater0. Replacing target ratios by r′\_k=∏\_i r\_i\^{}(U\_ki) and the torus variables by x\_l=∏\_k y\_k\^{}(V\_lk) changes the system into

\begin{verbatim}
y_k^(δk)=r′_k  (k≤h),   1=r′_k  (k>h).
\end{verbatim}

The zero rows are integer relations and therefore their equations hold by condition2. Each nonzero complex number has a δkth root, so choose those roots and arbitrary nonzero unused variables. Undoing the unimodular substitutions solves the system. The equal-multidegree criterion is exactly the specialization where all u\_i agree. This proves every assertion. Over a number field K the same proof additionally requires the indicated r′\_k to be δkth powers in K×; it does not silently adjoin those roots or identify periods with algebraic numbers. ∎

\subsubsection{Integral finite-input equations}\label{integral-finite-input-equations}

If a finite-dimensional ansatz for the finite Schwartz input has been specified, the analytic left side of (DO8.1) is linear in its input coefficients. After choosing actual integral generators, fixing the period scalars, and making a common algebraic denominator clearing that has itself been justified, write a localized system as

\begin{verbatim}
Cx=b,   C∈Mat_(q×s)(O_v),  b∈O_v^q.          (DO8.2)
\end{verbatim}

Here b is the specified jet vector including c\_j and the period scales. Let π be a uniformizer. Suppose UCV has nonzero Smith diagonal π\^{}(a\_1),\ldots,π\^{}(a\_h), with U,V invertible over O\_v, and put b′=Ub. Then an integral solution exists exactly when

\begin{verbatim}
b′_i=0 for i>h, and v(b′_i)≥a_i for i≤h.     (DO8.3)
\end{verbatim}

Assuming the zero-row conditions hold, the least d≥0 admitting a solution x∈π\^{}(−d)O\_v\^{}s is

\begin{verbatim}
max(0, max_(i≤h)(a_i−v(b′_i))),              (DO8.4)
\end{verbatim}

with v(0)=+∞ and an empty maximum contributing no positive bound.

Proof. Smith form over a DVR follows by choosing a matrix entry of least valuation, moving it to the top left, and clearing its row and column using divisibility by that entry. Repeat on the smaller submatrix; multiplication by units turns every nonzero diagonal entry into π to its valuation. All row and column operations are invertible over O\_v. With y=V\^{}(−1)x, (DO8.2) becomes π\^{}(a\_i)y\_i=b′\_i for i≤h and 0=b′\_i afterwards. This proves (DO8.3). Permitting y\_i∈π\^{}(−d)O\_v changes the first inequalities to v(b′\_i)≥a\_i−d, giving (DO8.4). These basis changes preserve the lattice, so the criterion does not depend on the selected Smith bases. ∎

This is a conditional integral criterion for a \emph{specified} finite-input span. Allowing a larger Schwartz space may change C; choosing arbitrary analytic coordinates and treating them as integral does not instantiate (DO8.2). For the terminal detector, the analogous rank-one image-ideal and archimedean-norm costs are precisely AM-R-002, applied to the actual terminal vector and the fixed mass-weighted row lattice. Its stabilizer and Fock denominators must be retained.

\subsubsection{The remaining all-degree and family requirements}\label{the-remaining-all-degree-and-family-requirements}

The actual \nolinkurl{OnFramesRdApply.lean} definition \nolinkurl{ClsK.Core}, lines108--113, quantifies over every polynomial P. Thus a finite matrix check certifies its restriction to the chosen monomial span, not the entire predicate. A full witness needs the equality for every degree, or a proved equality of the relevant analytic theta germs from which every jet follows. The finite class sets do not make the polynomial quantifier finite.

The source's \nolinkurl{ClsK.Kval}, lines115--121, subsequently uses fixed finite weight spaces and the reciprocal stabilizer masses. The terminal \nolinkurl{ClsT.Is.top/tail}, lines143--150, must hold for all the stated terminal vectors and rows, with one fixed branching map, level and top covolume. AM-R-002 supplies the cost of normalizing one already constructed detector; it cannot substitute for those simultaneous identities. Finally \nolinkurl{PinColT.cst} in \nolinkurl{LevChiEqOne.lean} fixes cK\_j from the actual covolume, finite datum multiplier, vacuum scalar and determinant-root characters, so the free rescalings in the first criterion represent the most permissive test. Failure of that test is an obstruction; success does not override the specified normalization.

For the repeated-type geometry AM-DO-001 the actual matrices Θ,A,C,b and terminal mass matrix have not been calculated. Consequently no global value for the requested G1 pairing follows from this reduction. The concrete remaining calculation is the signed equal-rank first kernel, with its all-degree CM-jet identity, followed by the compatible second kernel and terminal row equations. AM-DO-004 and AM-DO-007 show why the existing common positive, strictly increasing constructor does not supply these entries.

\subsection{AM-DO-009. Exact tensor-to-oriented-Hom comparison, including the source dual}\label{am-do-009.-exact-tensor-to-oriented-hom-comparison-including-the-source-dual}

\subsubsection{The global symplectic and group map}\label{the-global-symplectic-and-group-map}

Use the first-argument-linear Hermitian convention of 032, §4.2, equation(4.4), printed pp14--15. Let V and P have Gram matrices G\_V and G\_P, and let δ and ξ0 be nonzero purely imaginary elements of E. The 032 space is

\begin{verbatim}
W032 = V⊗_E P,
B032(x,y)=Tr_{E/F}(δ (h_V⊗h_P)(x,y)).
\end{verbatim}

Define S to be the algebraic E-dual P\^{}∨, with Hermitian Gram

\begin{verbatim}
G_S=(ξ0/δ) G_P^(−t).                         (DO9.1)
\end{verbatim}

The ratio ξ0/δ belongs to F×. In the dual basis of S and the given basis of V, identify a tensor with its coefficient matrix:

\begin{verbatim}
I: V⊗P → Hom_E(S,V),
I(Σ_(i,a) X_ia v_i⊗p_a)=X.                 (DO9.2)
\end{verbatim}

Then I is E-linear and an exact isometry of the F-valued alternating spaces. It intertwines the tensor action of U(V)×U(P) with the oriented action of U(S)×U(V) under

\begin{verbatim}
(g,k) ↦ (k^(−t),g),
I((g⊗k)x)=g I(x) k^t.                      (DO9.3)
\end{verbatim}

This is an explicit underlying rational oscillator comparison. It is not a comparison with the all-positive \nolinkurl{StdPair} pin, and it does not alone establish any MainX kernel or covolume identity.

\subsubsection{Full proof}\label{full-proof-4}

The dual Hermitian Gram G\_P\^{}(−t) is Hermitian: this follows by transposing and conjugating the identity G\_P\emph{=G\_P. Scaling it by ξ0/δ∈F× preserves Hermitian symmetry and nondegeneracy. If k}G\_Pk=G\_P, transpose and invert this equality to obtain

\begin{verbatim}
(k^(−t))* G_P^(−t) k^(−t)=G_P^(−t).
\end{verbatim}

Thus k\^{}(−t) lies in U(S), proving that the group map is well-defined. The ordinary action on the algebraic dual is k\^{}(−t), so its inverse acts on the right of Hom matrices as k\^{}t. This proves(DO9.3).

Write x=Σ X\_ia v\_i⊗p\_a and y=Σ Y\_jb v\_j⊗p\_b. The first-argument-linear convention gives

\begin{verbatim}
(h_V⊗h_P)(x,y)
  =Σ_(i,a,j,b) conjugate(Y_jb)(G_V)_ji(G_P)_ba X_ia
  =tr(G_P^t Y* G_V X).
\end{verbatim}

By(DO9.1), G\_S\^{}(−1)=(δ/ξ0)G\_P\^{}t. The exact source definition \nolinkurl{hO} therefore gives

\begin{verbatim}
hO(S,V)(I(x),I(y))
  =tr(G_S^(−1)Y*G_VX)
  =(δ/ξ0)(h_V⊗h_P)(x,y).
\end{verbatim}

Multiplying by ξ0 and taking Tr\_\{E/F\} proves

\begin{verbatim}
symplO(S,V)(I(x),I(y))=B032(x,y).
\end{verbatim}

The same equality remains true after Tr\_\{F/Q\} and at every completion. The tensor-to-Hom map is bijective on the displayed bases, completing the proof. ∎

\subsubsection{Polarization, additive character, splitting and measure scope}\label{polarization-additive-character-splitting-and-measure-scope}

Transport 032's actual positive Lagrangian through I\_C. At the active place its two blocks are exactly the images of L⊗P\_μ and Q\emph{⊗P}\_μ from032(4.5); at the compact places they are the images of the specified embedding blocks. The positivity equation is preserved because the alternating forms are equal. This describes the chosen polarization rather than replacing it by an all-positive coordinate pin.

Choose the SAME global additive character and transport one rational Schrödinger position space and its self-dual measures through I. The associated Heisenberg groups then identify by (x,t)↦(I(x),t), with central coordinate and character unchanged. On the transported Schrödinger function spaces the intertwiner is explicitly f↦f∘I\^{}(−1) on the position spaces. The Weil actions coincide when both are pulled back from the same metaplectic representation and its same lifts. This assertion needs no change of additive character and no partial complex Fourier transform. The predicate \nolinkurl{IsStdChar D ψ (-1)} of the recovered source must separately be compared with 032's fixed theta multiplier; equality of alternating forms does not discharge that predicate. The map also does not identify separately chosen lattice measures or finite theta structures; those must be transported and compared.

The group relabelling in(DO9.3) changes determinant-root labels. Choose idele roots t,y with

\begin{verbatim}
t/c(t)=det k,  y/c(y)=det g.
\end{verbatim}

Then the source root for k\^{}(−t) is t\^{}(−1), and Alpöge's oriented character factor is

\begin{verbatim}
splitFacO(χ_S,χ_V)(t^(−1),y)
  =χ_V(t) χ_S(y).                          (DO9.4)
\end{verbatim}

This follows from \nolinkurl{LocalNewformFunction.lean}, line51. The root condition is \nolinkurl{TiedWeil.lean}, \nolinkurl{IsDetRoot}, lines64--66. At a split place \nolinkurl{IsDetAt}, lines77--78, fixes the particular root t=e·det k+(1−e); the same inverse substitution applies. The arguments of these characters are the determinant roots, not the determinants. Moreover \nolinkurl{IsWeilDatumO.om_split} retains a separate positive real factor r; (DO9.4) describes only its character factor.

Any proposed identification with already fixed 032 circle splittings must match these roots and this character factor. Transporting the existing metaplectic lifts gives one consistent comparison; matching arbitrary independently prescribed Hecke-character splittings additionally requires their determinant-character ratio. On SU(V) the ratio of two group splittings is trivial, as carried032 Lemma4.4, printed pp20--21, states. This does not set an arbitrary χ\_S(y) to1 for every possible root y. No unprovided finite components of a global Hecke character, positive normalization r or universal Weil-index scalar are assigned here.

\subsubsection{The explicit Q(ζ7) middle form}\label{the-explicit-qux3b67-middle-form}

For AM-DO-001, G\_P=−a I2. Formula(DO9.1) gives

\begin{verbatim}
G_S=s I2,  s=−ξ0/(δa).                      (DO9.5)
\end{verbatim}

With ξ0=√−7/7 and δ=ζ−ζ\^{}(−1), the signs of ξ0/δ at σ1,σ2,σ3 are(+,+,−): the quadratic Gauss sum √−7 is fixed by σ2 and negated by σ3, whereas all three imaginary parts of δ are positive. As sign(a)=(+,−,−), the signs of s are(−,+,−). In particular S is still totally definite. The scalar Gram −aI2 would give the wrong exact symplectic normalization for this independently fixed ξ0; it is not used here.

The literal definite chain has endpoint forms G\_A=I2 and G\_X2=(a/q)I2, q=Nξ0. Its actual oriented edge ratios are s and a/(qs). Their product is a/q, so AM-DO-004's simultaneous positive-pin obstruction applies unchanged. The conversion solves the rational symplectic and source-label issue; it does not solve the finite CM-jet, character, covolume, all-degree kernel or terminal-mass gates.

\subsection{AM-DO-010. The evaluation-point CM tori are explicit primitive-J powers}\label{am-do-010.-the-evaluation-point-cm-tori-are-explicit-primitive-j-powers}

Carry B19a's exact rational E-Hodge decomposition for one multiplicity line at the rational negative line x={[}Ee3{]}:

\begin{verbatim}
H^1(A_x,Q) ≃ U(d)^⊕2 ⊕ U(dflip),
A_x ∼_E B(d)^2 × B(dflip).                 (DO10.1)
\end{verbatim}

The action is the specified geometric E-action of B19a. In AM-DO-001, the positive set of d is\{1,4,5\}; changing the active pair gives the positive set of dflip as\{4,5,6\}.

For s∈(Z/7Z)×, write J\^{}(s) for the same abelian variety J with its E-action precomposed by σ\_s. Its H\^{}\{1,0\} type is sΦ, where Φ=\{1,2,3\} is the carried type of J. Therefore

\begin{verbatim}
B(d) ∼_E J^(4),   B(dflip) ∼_E J^(6),
A_x ∼_E (J^(4))^2 × J^(6).                 (DO10.2)
\end{verbatim}

Forgetting the E-action, A\_x is isogenous to J\^{}3. The two repeated multiplicity inputs of AM-DO-001 consequently use CM abelian varieties whose product is isogenous to J\^{}6.

Proof. On a type-τ eigenspace, replacing ι(a) by ι(σ\_s(a)) changes its eigencharacter to τ∘σ\_s. The Galois group is abelian and identified with multiplication modulo7, so the new type is sΦ. Direct multiplication gives

\begin{verbatim}
4Φ={4,1,5}={1,4,5},
6Φ={6,5,4}={4,5,6}.
\end{verbatim}

Each rank-one rational E-Hodge structure is determined by this type. Thus its E-Hodge structure is isomorphic to that of the indicated twist of J. The carried equivalence between rational Hodge maps in weight one and homomorphisms of complex abelian varieties up to isogeny, already used to prove B19a, supplies the two E-isogenies. Inserting them into(DO10.1) proves(DO10.2). The underlying abelian variety of every J\^{}(s) is J, giving the final two assertions. ∎

This is an isogeny and a type identification, not an identification of the actual polarization, symmetric theta line, finite theta structure, integral lattice, or algebraic tangent frame. The exact B19 comparison still retains c\_i, p\_i, N\_i, M\_i and κ\_x from(B19c)--(B19f). In particular a model by J powers makes the primitive CM source explicit but does not turn its theta-null period into1 or its arithmetic functionals into freely chosen linear maps. Root's AM-R-003 explains why these simple primitive-J powers cannot simultaneously be the old elliptic Serre-tensor model.

\clearpage

\section{ALP\_MACHINE --- question 3: the Kummer character and the precise transfer gate}\label{alp_machine-question-3-the-kummer-character-and-the-precise-transfer-gate}

Author/checker: GPT-6 Astra Pro, ALP\_MACHINE child \nolinkurl{/root/kummer_extension}; inherited reasoning effort is not separately exposed. Date: 21 November 2026. This is a mathematical contribution for integration into the main paper, not a claim that the requested RH extension has been completed.

\subsection{0. Outcome and scope}\label{outcome-and-scope}

The source has a field-independent analytic theorem that can be applied to the Kummer character. \textbf{AM-K005} below proves all of its pole-data hypotheses and identifies a single, explicit remaining arithmetic window estimate. For \(r=3\), \textbf{AM-K004} also realizes the character as a rank-one summand in a tensor construction from two explicit CM elliptic curves. These are constructive reductions. They do not prove the missing window estimate.

The literal existing character host cannot receive the Kummer splitting field or its induced representation. This is stronger than a missing notation adapter: \textbf{AM-K002} proves three exact obstructions---non-CM splitting field, failure of norm pullback from a rational Dirichlet character, and failure of a finite Artin-formalism factorization into rational Dirichlet characters. Each obstruction is explicitly scoped; none asserts that a different CM-geometric construction is impossible.

Alpöge's theorem and archive are used as accepted theorems, as directed. CX-ALG-024, CX-R02, CX-G06, DP6, DP32, DP34, DP42 and the supplied source-interface results are carried with their existing status. No check of the validity of the accepted Dirichlet GRH theorem is undertaken. Masullo's consequences remain conditional. No Lean compilation was performed in this contribution; the archive checks below are checks of actual declarations and their types.

\subsubsection{0.1 Source provenance and what was actually read}\label{source-provenance-and-what-was-actually-read}

The primary note is Alpöge, \emph{Remark}, pp.~3--4, numbered timeline steps 2--20, \url{https://alpo.ge/remark.pdf}. The note calls this a rough timeline. Its historical step numbering is not itself a one-to-one module index.

Actual Lean text was recovered at pinned mirror commit \nolinkurl{7cbf7a0c32fca1badf05931879a37c3f8a2b03b4} of \nolinkurl{CocoSgt/riemann-hypothesis-lean}, under \nolinkurl{src/}. The local copy is \nolinkurl{work/agents/kummer_sources/}; \nolinkurl{manifest.json} records each downloaded file, byte count and SHA-256. Two files checked against the supplied original-author source manifest already match: \nolinkurl{Formalism/SizeWindowBrackets/GrowthBoundGivesRiemannHypothesis.lean} and \nolinkurl{Development/AutomorphicArithmetic/ThetaCorrespondence/ThetaPairLifts/ThirdPlaceSeries/SiegelEisensteinSecondTermHypotheses.lean}. The remaining copies are pinned mirror text, corroborated where applicable by the supplied original-archive map; this contribution does not claim every mirror file has independently been compared with the original archive.

The supplied map is \nolinkurl{RETURN_THETA3_checkpoint_07_20261121T0604Z/supplied_sources/ALPOGE_LEAN_MAP_20261120.md}, especially §§3--4. Its low-confidence keyword correspondences for steps 11, 14, 17 and 19 are not promoted to a complete dependency audit here.

\subsection{\texorpdfstring{1. Setup and the \(p=r\) edge case}{1. Setup and the p=r edge case}}\label{setup-and-the-pr-edge-case}

Let \(r\) be an odd prime, \(p\) a rational prime, \(\alpha=p^{1/r}>0\), and

\[
K=\mathbf Q(\zeta_r),\qquad F=\mathbf Q(\alpha),\qquad E=K(\alpha).
\]

The user assumes \(r\nmid(p-1)\). The algebraic results below do not need this restriction. It matters for the surrounding finite-field construction. Let \(\chi=\chi_{p,r}:G_K\to\mu_r\) be \(\sigma\mapsto\sigma(\alpha)/\alpha\), viewed by class field theory as the primitive finite-order Hecke character with its actual conductor. Write \(L(s)=L_K(s,\chi)\), retaining all primitive local factors.

\subsubsection{\texorpdfstring{AM-K000. The edge \(p=r\) is permitted and is algebraically nondegenerate}{AM-K000. The edge p=r is permitted and is algebraically nondegenerate}}\label{am-k000.-the-edge-pr-is-permitted-and-is-algebraically-nondegenerate}

\textbf{Statement.} The displayed condition permits \(p=r\). For this case \([E:K]=r\), the same affine Galois group and character analysis used below apply, and the finite-field task of constructing \(\mathbf F_{p^p}\) has the elementary irreducible polynomial \(T^p-T-1\in\mathbf F_p[T]\). Its use requires no RH premise.

\textbf{Proof.} The polynomial \(T^r-r\) is Eisenstein over \(\mathbf Q\), so \([F:\mathbf Q]=r\). Because \([K:\mathbf Q]=r-1\), the intersection \(F\cap K\) is \(\mathbf Q\), and \([E:K]=r\). Equivalently, at the unique prime \(\mathfrak r\) of \(K\) above \(r\), the normalized valuation satisfies \(v_{\mathfrak r}(r)=r-1\not\equiv0\pmod r\), so \(r\) is not an \(r\)-th power in \(K\). The statement is not that \(T^r-r\) is Eisenstein over \(K\); its constant-term valuation there is \(r-1\). The automorphisms \(\alpha\mapsto\zeta_r^b\alpha\), \(\zeta_r\mapsto\zeta_r^u\) give the full affine group of order \(r(r-1)\).

For the finite-field assertion, let \(a\) be a root of \(T^p-T-1\) in an algebraic closure of \(\mathbf F_p\). Then \(a^p=a+1\), and induction gives \(a^{p^m}=a+m\) for every nonnegative integer \(m\). The least positive \(m\) with \(a^{p^m}=a\) is \(p\). Its Frobenius orbit therefore has length \(p\), so its minimal polynomial has degree \(p\) and equals the given degree-\(p\) polynomial. ∎

The previously accepted CX-ALG-024 was stated for \(p\ne r\). The finite-field observation separates the omitted wild case from that statement without changing it.

\subsection{2. The exact target representation and L-function}\label{the-exact-target-representation-and-l-function}

\subsubsection{AM-K001. A single real L-function for all nontrivial Kummer powers}\label{am-k001.-a-single-real-l-function-for-all-nontrivial-kummer-powers}

\textbf{Statement.} Set \(N=\operatorname{Gal}(E/K)\simeq C_r\), \(G=\operatorname{Gal}(E/\mathbf Q)\simeq\operatorname{AGL}_1(\mathbf F_r)\), and let \(P\) be the permutation representation of \(G\) on the \(r\) roots of \(T^r-p\). For every nontrivial character \(\vartheta\) of \(N\),

\[
\rho:=\operatorname{Ind}_N^G\vartheta\simeq P-\mathbf 1.
\tag{K1}
\]

Consequently, for \(a=1,\ldots,r-1\),

\[
L_K(s,\chi^a)=L_{\mathbf Q}(s,\rho)
=\frac{\zeta_F(s)}{\zeta_{\mathbf Q}(s)}=:L(s).
\tag{K2}
\]

The identities hold with ramified local factors as well. The resulting Dirichlet series has real coefficients, \(L(\bar s)=\overline{L(s)}\), and its representation is self-dual. Under the accepted Dirichlet GRH theorem, RH for this \(L\) is equivalent to RH for \(\zeta_F\) and for \(\zeta_E\).

\textbf{Proof of the additional representation statement.} Write an affine map as \(x\mapsto ux+b\). The permutation character has values \(r\), \(0\), \(1\), respectively, at the identity, a nonidentity translation, and an element with \(u\ne1\). Its augmentation character therefore has values \(r-1,-1,0\).

Since \(N\) is normal, the induced character vanishes off \(N\). On \(N\), its restriction is the sum of the conjugates of \(\vartheta\) under \(\mathbf F_r^\times\). This group acts transitively and freely on all nontrivial characters of \(N\). Their sum has value \(r-1\) at the identity and \(-1\) at every other translation. The two characters agree, proving (K1). This also proves that inducing any nontrivial power gives the same representation.

Artin induction compatibility for L-functions now gives (K2); the standard permutation representation identity gives \(L(P)=\zeta_F\), while \(L(\mathbf1)=\zeta_{\mathbf Q}\). These identities are local identities on inertia invariants, so they include ramified primes. Since \(P-\mathbf1\) is defined over \(\mathbf R\) and self-dual, the asserted reality and duality follow. Equivalently, conjugation of the Euler product sends \(\chi\) to \(\chi^{-1}\), which has the same L-function by (K1).

For the final equivalence, use the already checked CX-ALG-024 identity \(\zeta_E=\zeta_K L^{r-1}\); it is not rederived here. Inside \(0<\Re s<1\), both \(\zeta_{\mathbf Q}L\) and \(\zeta_K L^{r-1}\) are products of holomorphic functions, so their zero divisors add with nonnegative multiplicities. Dirichlet GRH puts every zero of \(\zeta_{\mathbf Q}\) and \(\zeta_K\) on the critical line. The equivalence follows. ∎

\textbf{Sources.} CX\_GRH/GRH\_NIGHT, CX-ALG-024, in \nolinkurl{GRH_NIGHT.tex}, subsection containing lines 14723--14826 of the supplied checkpoint. For induction, Milne (2020a), VIII §10, pp.~261--262. Milne uses ``Dirichlet'' also for ray-class characters of general number fields in that discussion; this does not make those characters rational Dirichlet characters of \((\mathbf Z/N\mathbf Z)^\times\).

\subsection{3. Three exact obstructions to literal reuse}\label{three-exact-obstructions-to-literal-reuse}

\subsubsection{AM-K002. No cyclotomic host, norm pullback, or rank-one Dirichlet factorization}\label{am-k002.-no-cyclotomic-host-norm-pullback-or-rank-one-dirichlet-factorization}

\textbf{Statement.} For odd \(r\):

\begin{enumerate}
\def\labelenumi{\arabic{enumi}.}
\tightlist
\item
  \(E\) is not a CM field and cannot embed into a CM field. In particular it cannot be the field supplied to the literal \nolinkurl{IsChiField} record.
\item
  \(\chi\) is not the pullback of a rational Dirichlet character along the idelic norm \(N_{K/\mathbf Q}\).
\item
  \(\rho\) is not a virtual finite sum of one-dimensional rational Galois characters. Thus a finite factorization of its L-function into rational Dirichlet L-functions cannot be obtained by equality of Artin Euler factors, even after enlarging the common finite Galois extension.
\end{enumerate}

These statements do not exclude CM abelian varieties defined over \(K\), tensor constructions involving their twists, or a new Hecke-character version of Alpöge's arithmetic host.

\textbf{Proof.} Let \(\sigma(\alpha)=\zeta_r\alpha\) and let \(c\) be ordinary complex conjugation. Since \(c(\alpha)=\alpha\) and \(c(\zeta_r)=\zeta_r^{-1}\),

\[
c\sigma c^{-1}=\sigma^{-1}\ne\sigma.
\tag{K3}
\]

Complex conjugation is central in the Galois group of a normal CM field: its intrinsic involution commutes with every automorphism. Thus \(E\), which is normal, is not CM. There is also a direct containment obstruction. In a CM field embedded in \(\mathbf C\), the subfield of elements fixed by ordinary complex conjugation is totally real. The positive real number \(\alpha\) cannot be in that subfield, because its minimal polynomial \(T^r-p\) has nonreal conjugates. If an embedding of the normal field \(E\) into a CM field existed, its image in \(\mathbf C\) would again be \(E\) and would contain \(\alpha\), a contradiction. This proves (1). The central-involution characterization is Milne (2020b), Proposition 1.4 and Corollary 1.5, p.~10.

For (2), conjugation acts on the Kummer character by \(\chi^c=\chi^{-1}\), from (K3). A character pulled back through \(N_{K/\mathbf Q}\) is \(c\)-invariant, since \(N(c x)=N(x)\). It would follow that \(\chi=\chi^{-1}\). An order-\(r\) character for odd \(r>1\) cannot satisfy this equality.

For (3), the commutator of a multiplier \(h_u:x\mapsto ux\) with a translation \(\sigma_b:x\mapsto x+b\) is translation by \((u-1)b\). Choosing \(u\ne1\) shows that every translation is a commutator. Since \(G/N\) is abelian, \([G,G]=N\). Every one-dimensional character is trivial on \(N\); every virtual sum of such characters therefore has the same trace at the identity and at any translation. But (K1) gives trace \(r-1\) at the identity and \(-1\) at a nonidentity translation. This proves the representation obstruction, extending the accepted cubic example CX-R02 to the displayed family.

For completeness, enlargement of the finite quotient cannot remove it. Suppose a finite collection of rational Dirichlet characters and integer exponents gave the same unramified Euler factors as \(L(s,\rho)\). Form a finite Galois extension containing \(E\) and all their cyclotomic fields. At every prime outside a finite set, equality of the coefficient of \(T\) in the Euler factors gives equality of the corresponding virtual traces. Chebotarev gives this equality on every conjugacy class of the enlarged Galois group. The induced map on commutator subgroups surjects onto \([G,G]=N\), whereas each one-dimensional character is trivial on the enlarged commutator. Lifting a nonidentity translation contradicts its trace \(-1\ne r-1\). The same argument applies to rational exponents after clearing a common denominator. ∎

\textbf{Exact source boundary.} \nolinkurl{CharacterBoxes/ExistsPageN.lean}, definition \nolinkurl{IsChiField}, line 19, explicitly requires a \(\mathbf Q\)-algebra embedding into \nolinkurl{CyclotomicField N ℚ}. It also uses a fixed field determined by rational congruences modulo \(2^{e+2}\), a primitive rational character \(\psi\), and two Jacobi symbols. \nolinkurl{Formalism/SharpRestriction/RamificationFrobenius/TotientGaloisClasses/Analytic.lean}, \nolinkurl{HostRaw.dir}, line 102, explicitly constructs a \nolinkurl{DirichletCharacter ℂ H.condN} from the rational residue map \nolinkurl{H.art}. This is the literal specialization excluded here.

\textbf{Scope.} The conclusion is not that RH for this L-function is logically independent of Dirichlet GRH. It rules out the stated finite rank-one Artin factorization and these literal host substitutions. In particular, Brauer induction permits characters over other base fields; it does not turn those bases into \(\mathbf Q\).

\subsection{4. The cubic case with its actual conductor and completed function}\label{the-cubic-case-with-its-actual-conductor-and-completed-function}

\subsubsection{AM-K003. Cubic arithmetic data and the target Euler coefficients}\label{am-k003.-cubic-arithmetic-data-and-the-target-euler-coefficients}

Let \(r=3\), \(K=\mathbf Q(\sqrt{-3})\), \(\lambda=1-\zeta_3\), and assume \(3\nmid p-1\). Thus \(p=3\) or \(p\equiv2\pmod3\). Let \(\mathfrak f\) be the primitive conductor of \(\chi_{p,3}\).

\[
\begin{array}{c|c|c|c}
\text{case}&D_F&\mathfrak f&N\mathfrak f\\ \hline
p\ne3,\ p\equiv8\pmod9&-3p^2&(p)&p^2\\
p\ne3,\ p\equiv2,5\pmod9&-27p^2&(p)\lambda^2&9p^2\\
p=3&-243&\lambda^4&81
\end{array}
\tag{K4}
\]

The analytic conductor factor is \(|D_K|N\mathfrak f=|D_F|\). With \(\Gamma_{\mathbf C}(s)=2(2\pi)^{-s}\Gamma(s)\),

\[
\Lambda(s)=|D_F|^{s/2}\Gamma_{\mathbf C}(s)L(s),
\qquad \Lambda(s)=\Lambda(1-s).
\tag{K5}
\]

In particular the root number is \(+1\); the functional equation does not force a central zero.

\textbf{Justification and specialization.} Harron (2016), §2.1, Lemma 2.1, p.~5, gives the integral bases and discriminants of pure cubic fields. Its Type II condition specializes to \(p\equiv\pm1\pmod9\), and under \(p\equiv2\pmod3\) this is \(p\equiv8\pmod9\). Otherwise the discriminant is \(-27p^2\), including \(p=3\).

The Artin conductor of an induced character is the discriminant of the base times its relative conductor norm; in the present quadratic induction case see, for example, the formula in the proof of Hida (2025), Corollary 6.18, p.~45. One can verify the normalization directly here without any local convention ambiguity. The quotient of the standard completed Dedekind functions \(\Lambda_F/\Lambda_{\mathbf Q}\) is \(|D_F|^{s/2}\Gamma_{\mathbf C}(s)L(s)\), because \(F\) has signature \((1,1)\). It satisfies the functional equation with sign \(+1\). The standard primitive Hecke completion is \((3N\mathfrak f)^{s/2}\Gamma_{\mathbf C}(s)L(s)\). Its dual has the same L-function by AM-K001. Dividing its functional equation by that of the Dedekind quotient gives \(C^{s/2}=W C^{(1-s)/2}\), where \(C=3N\mathfrak f/|D_F|>0\). On an open set where the functions are nonzero this identity forces \(C=1\) and \(W=1\), and then holds by continuation. This proves the asserted conductor normalization.

Only primes above \(3p\) can ramify, from the polynomial discriminant. If \(p\ne3\), \(p\equiv2\pmod3\) is inert in \(K\), and \(T^3-p\) is Eisenstein at \((p)\). The resulting local degree-three extension is tame, so every nontrivial local character has conductor exponent one. The remaining conductor exponent at \(\lambda\), whose norm is \(3\), is determined by the displayed norm identity. This gives (K4), including \(\lambda^4\) when \(p=3\). ∎

At a rational prime \(\ell\nmid3p\), use the factorization of \(T^3-p\) modulo \(\ell\). With \(T=\ell^{-s}\), the reciprocal Euler polynomials of \(L(s)\) are

\[
\begin{array}{c|c|c}
\text{factorization}&\det(1-\rho(\mathrm{Frob}_\ell)T)&
\operatorname{tr}(\rho(\mathrm{Frob}_\ell)^m)\\ \hline
1+1+1&(1-T)^2&2\\
1+2&1-T^2&0\ (m\text{ odd}),\ 2\ (m\text{ even})\\
3&1+T+T^2&-1\ (3\nmid m),\ 2\ (3\mid m)
\end{array}
\tag{K6}
\]

These formulas follow directly from the eigenvalues \(1,1\), \(1,-1\), and \(\zeta_3,\zeta_3^{-1}\) of the augmentation representation. Thus the rational prime-power coefficient in \(-L'/L\) is the last column times \(\log\ell\). This is the concrete coefficient system that a new geometric family must produce. It is not a periodic rational Dirichlet-character coefficient system.

The same completion argument in arbitrary odd prime degree gives

\[
\Lambda_{p,r}(s)=|D_F|^{s/2}\Gamma_{\mathbf C}(s)^{(r-1)/2}L(s),
\qquad\Lambda_{p,r}(s)=\Lambda_{p,r}(1-s),
\tag{K7}
\]

using the signature \((1,(r-1)/2)\) of \(F\) and (K2). This statement needs no explicit general formula for \(D_F\).

\subsection{5. An explicit geometric realization in the cubic case}\label{an-explicit-geometric-realization-in-the-cubic-case}

\subsubsection{AM-K004. The Kummer character is a Hom line between two CM elliptic motives}\label{am-k004.-the-kummer-character-is-a-hom-line-between-two-cm-elliptic-motives}

\textbf{Statement.} Over \(K=\mathbf Q(\zeta_3)\), take

\[
A_0:y^2=x^3+1,\qquad A_p:y^2=x^3+p^2.
\]

For an auxiliary prime \(\ell\), use a coefficient field \(C_\ell\subset\overline{\mathbf Q}_\ell\) containing \(\zeta_3\). Let \(V_0,V_p\) be the \(\zeta_3\)-eigenlines of the endomorphism \([\zeta_3]:(x,y)\mapsto(\zeta_3x,y)\) acting on \(H^1_{\mathrm{\acute et}}\). Then, with this eigenline convention,

\[
V_p\simeq V_0\otimes\chi_{p,3},\qquad
\operatorname{Hom}_{C_\ell}(V_0,V_p)\simeq\chi_{p,3}.
\tag{K8}
\]

This is realized by algebraic correspondences: the Artin line is a direct factor of \(h^1(A_p)\otimes h^1(A_0)^\vee\), with coefficients containing \(K\). It has weight zero. The realization does not provide Alpöge's required quantitative theta-lift family or the window estimate of §6.

\textbf{Proof.} Put \(\alpha=p^{1/3}\). Over \(E=K(\alpha)\) the map

\[
g:A_p\longrightarrow A_0,\qquad (x,y)\longmapsto(\alpha^{-2}x,p^{-1}y)
\tag{K9}
\]

is an isomorphism, since \(\alpha^6=p^2\). For \(\sigma\in G_K\), the ratio \(\sigma(\alpha)/\alpha=\chi(\sigma)\) lies in \(\mu_3\), and \(\chi^{-2}=\chi\). Therefore

\[
\sigma(g)=[\chi(\sigma)]\circ g.
\tag{K10}
\]

On cohomology, functoriality gives

\[
\rho_p(\sigma)g^*\rho_0(\sigma)^{-1}
=(\sigma g)^*=g^*[\chi(\sigma)]^*.
\]

The two \([\zeta_3]^*\)-eigenvalues on \(H^1\) are \(\zeta_3,\zeta_3^{-1}\). This can be seen on the differential \(dx/y\), where the eigenvalue is \(\zeta_3\), and its complex-conjugate Hodge summand; the characteristic polynomial is unchanged in étale realization. The projector

\[
e_{\zeta_3}=\frac13\sum_{j=0}^2\zeta_3^{-j}[\zeta_3]^{j*}
\tag{K11}
\]

is an algebraic-correspondence projector over the coefficient field. Restricting the displayed functoriality identity to its image yields

\[
\rho_p(\sigma)g^*=\chi(\sigma)g^*\rho_0(\sigma),
\]

which is precisely (K8). The line in the Hom motive becomes the trivial line over \(E\) with descent cocycle \(\chi\); this identifies it with the Artin motive defined by that character, not merely with an unrelated equal-dimensional representation. As a tensor of weight-one cohomology with its dual, it has weight zero. ∎

\textbf{What this does and does not change.} Both elliptic curves have CM by \(K\), are isomorphic over \(\overline{\mathbf Q}\), and have the same geometric CM type and stable Faltings height. At a good prime their formal groups become isomorphic after extension. At good rational primes inert in \(K\), these formal groups have height two, so the existing height-two derivative mechanism applies after suitable base extension; DP6 supplies its known exponent. Split good primes have height one and give no extra height-two contribution. A finite-order descent twist is actual arithmetic information, but it is not a new geometric height budget.

The invariant differentials in these equations satisfy

\[
g^*(dx/y)_{A_0}=\alpha\,(dx/y)_{A_p}.
\tag{K12}
\]

Indeed \(\alpha^{-2}/p^{-1}=p\alpha^{-2}=\alpha\). For transported complex cycles, \(\Omega_0=\alpha\Omega_p\). These algebraic factors must be retained in the integral lattices and period normalization. Removing them by consistent normalization cannot itself supply a favorable extra lift cost. No equality of \nolinkurl{MainX} kernel, covolume, mass, or terminal-period records has been inferred from this isomorphism.

This construction is also why the non-CM statement in AM-K002 must not be overstated as ``the character cannot occur in CM motives.'' It does occur in this explicit tensor category over \(K\).

\subsection{6. A completed analytic transfer and the surviving arithmetic gate}\label{a-completed-analytic-transfer-and-the-surviving-arithmetic-gate}

\subsubsection{AM-K005. Field-independent Landau reduction for the Kummer L-function}\label{am-k005.-field-independent-landau-reduction-for-the-kummer-l-function}

Define \(\Lambda_K(\mathfrak a)=\log N\mathfrak q\) when \(\mathfrak a=\mathfrak q^m\) for a prime ideal \(\mathfrak q\) and \(m\ge1\), and zero otherwise. Extend the primitive character by zero at ideals not coprime to its conductor. Put \(c_1=L'(1)/L(1)\) and, for \(x\ge1\),

\[
G_\chi(x)=\sum_{N\mathfrak a\le x}\Lambda_K(\mathfrak a)\chi(\mathfrak a)
\left(\frac{x}{N\mathfrak a}-1\right)+xc_1.
\tag{K13}
\]

\textbf{Theorem.} This is real-valued. Fix \(0<\theta<1\) and \(\delta>0\). Suppose the following single window condition holds:

\begin{quote}
\textbf{\(\mathcal W_\chi(\theta,\delta)\).} For every \(\eta>0\) and every \(Y_0\ge1\), there exist \(Y\ge Y_0\) and a measurable \(R:[Y,Y^{1+\delta}]\to\mathbf R\) such that \(0\le R(x)\), \(G_\chi(x)\le R(x)\), \(R(x)x^{-\theta-1}\) is integrable, and \[
\int_Y^{Y^{1+\delta}}R(x)x^{-\theta-1}\,dx\le Y^\eta.
\tag{K14}
\]
\end{quote}

Then \(L(s)\) has no zero with \(\Re s>\theta\). Consequently it suffices to prove \(\mathcal W_\chi(1/2+\epsilon,1/2)\) for every \(0<\epsilon<1/2\) to prove RH for \(L\), and hence the CX-ALG-024 consequences under accepted Dirichlet GRH.

\textbf{Proof.} We verify the actual \nolinkurl{Deduction.PoleData} interface and invoke the accepted source theorem \nolinkurl{Deduction.landau_window}. The analytic continuation, functional equation, Euler product for \(\Re s>1\), and nonvanishing \(L(1)\ne0\) of a nontrivial finite-order Hecke character are classical. In the Artin description (K2), the relevant holomorphy is the classical monomial case, obtained by its identification with this nontrivial Hecke L-function; it is not inferred from an unproved cancellation in a quotient of zeta functions. The standard nonvanishing at \(1\) is Milne (2020a), V Theorem 2.4(b), p.~156; its Euler product and continuation conventions are V §2 and VIII §10.

By AM-K001, the coefficients of \(L\), and therefore those of \(-L'/L\), are real. In particular \(c_1\in\mathbf R\). Grouping the ideal terms by rational norm writes

\[
-\frac{L'}L(s)=\sum_{n\ge1}\frac{b(n)}{n^s},\qquad
G_\chi(x)=\sum_{n\le x}b(n)(x/n-1)+xc_1,
\tag{K15}
\]

with \(b(n)\in\mathbf R\). At a rational prime \(\ell\), the local factor is the determinant of \(1-\rho(\mathrm{Frob}_\ell)\ell^{-s}\) on inertia invariants. All its nonzero eigenvalues are roots of unity, and there are at most \(d=r-1\) of them. Thus

\[
|b(\ell^m)|\le d\log\ell,
\qquad b(n)=0\quad\text{if }n\text{ is not a prime power}.
\tag{K16}
\]

This bound includes ramified primes and proves the absolute convergence needed below. It also gives the elementary growth estimate, for \(x\ge1\),

\[
|G_\chi(x)|\le dx\sum_{2\le n\le x}\frac{\log n}{n}+|c_1|x
\ll x(1+\log x)^2\ll x^2.
\tag{K17}
\]

At a threshold \(x=n\) the newly appearing term is zero. Hence \(G_\chi\) is continuous and piecewise affine. Set \(g(u)=G_\chi(e^u)\) for \(u\ge0\), and \(g(u)=0\) for \(u<0\). It is measurable and \(|g(u)|\le C e^{2u}\) for \(u\ge0\).

For \(\Re s>1\), direct integration gives

\[
\int_n^\infty (x/n-1)x^{-s-1}\,dx=\frac{n^{-s}}{s(s-1)}.
\]

The absolute integral of the summands is bounded by the corresponding expression with \(|b(n)|\) and \(\Re s>1\), which is finite by (K16). Termwise integration is therefore legitimate and yields

\[
\Phi(s):=\int_1^\infty G_\chi(x)x^{-s-1}\,dx
=-\frac{L'(s)}{s(s-1)L(s)}+\frac{c_1}{s-1}
=\frac{-L'/L(s)+s c_1}{s(s-1)}.
\tag{K18}
\]

At \(s=1\) the numerator vanishes. Define \(\Phi(1)\) by the removable continuation; it equals \(c_1-(L'/L)'(1)\). At zeros of \(L\), regard \(\Phi\) as meromorphic and assign arbitrary values if a total function is required by the formal interface. At points with positive real part outside those zeros it is holomorphic, including \(1\). Under \(x=e^u\), (K18) is the Laplace identity

\[
\Phi(s)=\int_0^\infty g(u)e^{-su}\,du,
\]

in particular on the source's required half-plane \(\Re s>2\).

Let \(Z=\{\rho:\Re\rho>0,\ L(\rho)=0\}\). The function \(L\) is entire, not identically zero, and \(L(1)\ne0\). Its zeros are isolated. If \(\rho\in Z\) has multiplicity \(m\), write \(L(s)=(s-\rho)^m a(s)\), with \(a\) nonvanishing near \(\rho\). Then

\[
\Phi(s)=\frac{m}{\rho(1-\rho)}\frac1{s-\rho}+h_\rho(s),
\tag{K19}
\]

where \(h_\rho\) is holomorphic near \(\rho\). The residue is nonzero since \(\rho\ne0,1\). If \(\rho\) is real, the Euler product excludes \(\rho>1\), and \(L(1)\ne0\), so \(0<\rho<1\). The residue \(m/[\rho(1-\rho)]\) is then strictly positive and real. These are exactly the positivity and simple-pole conditions of \nolinkurl{PoleData}; the pole is simple even when the zero of \(L\) has higher multiplicity.

All fields of the source record have now been supplied: measurability, exponential growth, Laplace representation, positive real parts of \(Z\), holomorphy away from \(Z\), isolation, and the required nonzero residues with positive sign at real zeros. No Hecke-family extension of its Dirichlet-specialized \nolinkurl{CharFamily} is needed for this application.

It remains to translate (K14). Given the \(U_0\in\mathbf R\) required by \nolinkurl{WindowHypU}, choose \(Y_0=\max(1,e^{U_0})\). Put \(U=\log Y\) and \(\widetilde R(u)=R(e^u)\). Then \(U\ge U_0\), \(\widetilde R\ge0\), \(g\le\widetilde R\), and the change of variables \(x=e^u\) gives

\[
\int_U^{(1+\delta)U}\widetilde R(u)e^{-\theta u}\,du
=\int_Y^{Y^{1+\delta}}R(x)x^{-\theta-1}\,dx
\le e^{\eta U}.
\]

This is precisely \nolinkurl{WindowHypU g θ δ}. The accepted theorem \nolinkurl{landau_window} gives \(\Re\rho\le\theta\) for every \(\rho\in Z\).

If the condition is known for every \(\theta=1/2+\epsilon\), then no zero has real part \(>1/2\). The self-dual functional equation (K7) maps a nontrivial zero of real part \(<1/2\) to one of real part \(>1/2\), while the gamma factor is nonzero and holomorphic in the critical strip. This proves RH under the asserted window conditions. ∎

\textbf{Actual Lean locators.} With common prefix \nolinkurl{Development/AutomorphicArithmetic/DedekindLFunctions/}:

\begin{longtable}[]{@{}
  >{\raggedright\arraybackslash}p{(\columnwidth - 4\tabcolsep) * \real{0.3333}}
  >{\raggedright\arraybackslash}p{(\columnwidth - 4\tabcolsep) * \real{0.3333}}
  >{\raggedright\arraybackslash}p{(\columnwidth - 4\tabcolsep) * \real{0.3333}}@{}}
\toprule\noalign{}
\begin{minipage}[b]{\linewidth}\raggedright
Record/theorem
\end{minipage} & \begin{minipage}[b]{\linewidth}\raggedright
File and locator
\end{minipage} & \begin{minipage}[b]{\linewidth}\raggedright
Exact transfer use
\end{minipage} \\
\midrule\noalign{}
\endhead
\bottomrule\noalign{}
\endlastfoot
\nolinkurl{Deduction.PoleData} & \nolinkurl{LandauZeroFreeCriterion.lean:15} & Arbitrary \(g:\mathbf R\to\mathbf R\), \(\Phi:\mathbf C\to\mathbf C\), \(Z\subset\mathbf C\); no character or base-field parameter. \\
\nolinkurl{Deduction.WindowHypU} & \nolinkurl{LandauZeroFreeCriterion.lean:25} & Quantifiers and weighted window integral translated in (K14). \\
\nolinkurl{Deduction.landau_window} & \nolinkurl{LandauTheorems/landauWindow.lean:83} & The accepted analytic implication used above. \\
\nolinkurl{poleData_abstract} & \nolinkurl{LandauTheorems/poleData.lean:22} & Confirms a generic analytic layer precedes the Dirichlet-family specialization at line 248. The proof above directly verifies the record instead of claiming a new Lean instantiation. \\
ramp Mellin kernel & \nolinkurl{LandauTheorems/MellinRamp.lean:10}; \nolinkurl{mellinGchiOfNeOne.lean:65,116} & Generic scalar kernel and arithmetic-sequence integration before Dirichlet specialization. \\
\end{longtable}

\textbf{A sufficient stronger estimate.} If, for every \(\epsilon>0\), \(G_\chi(x)\le x^{1/2+\epsilon}\) for all sufficiently large \(x\), take \(R(x)=x^{1/2+\epsilon}\). The integral in (K14) is \(\delta\log Y\), which is at most \(Y^\eta\) for all sufficiently large \(Y\). Thus a one-sided pointwise bound suffices; absolute-value control is unnecessary for this deduction. No such bound is proved here.

\textbf{Current gate, Q3-KUMMER-WINDOW.} Construct a relative character host and an actual arithmetic theta-lift family whose period, divisibility, covolume, integral-structure, and cost equations prove (K14), for each \(\epsilon>0\), with \(\theta=1/2+\epsilon\), already for the explicit cubic data (K4), (K8), (K13). This is an explicit analytic obligation left by the arithmetic construction. The source's generic Landau theorem eliminates a separate analytic-generalization obligation.

\subsection{7. Steps 2--20: actual interface audit and replacement obligations}\label{steps-220-actual-interface-audit-and-replacement-obligations}

The short titles in the following table identify the historical timeline. The right-hand column distinguishes a completed replacement from a required new instance. A source anchor is a declaration whose scope was actually checked, or a clearly labelled carried source interface. It is not a claim to have exhaustively inspected every proof body contributing to that step.

\begin{longtable}[]{@{}
  >{\raggedright\arraybackslash}p{(\columnwidth - 4\tabcolsep) * \real{0.3333}}
  >{\raggedright\arraybackslash}p{(\columnwidth - 4\tabcolsep) * \real{0.3333}}
  >{\raggedright\arraybackslash}p{(\columnwidth - 4\tabcolsep) * \real{0.3333}}@{}}
\toprule\noalign{}
\begin{minipage}[b]{\linewidth}\raggedright
Step
\end{minipage} & \begin{minipage}[b]{\linewidth}\raggedright
Q-sensitive interface and replacement over \(K\)
\end{minipage} & \begin{minipage}[b]{\linewidth}\raggedright
State and exact anchor
\end{minipage} \\
\midrule\noalign{}
\endhead
\bottomrule\noalign{}
\endlastfoot
2: local derivations & The local divided-power/height statement is already local arithmetic geometry. For \(r=3\), the CM elliptic curves of AM-K004 give explicit good-reduction inputs after extension. All conductor primes remain in the global bookkeeping. & DP6, DP32, DP34 and DP42 are carried; no new exponent is asserted. AM-K004 supplies the twist, not a better divisibility budget. The supplied map §4 labels its correspondence to \nolinkurl{BlockTheory/HeightDefect} and CM-good-curve files as low confidence. \\
3: quadric and biextension & Poincaré bundles and the relevant quadratic constructions can be formed over a number field. What must be newly proved is that the resulting character-weighted arithmetic functional equals the norm-indexed \(G_\chi\) in (K13), with its actual integral lattices. & The generic geometry does not manufacture this equality. Carried source families: \nolinkurl{Formalism/LocalToGlobal/GlobalTheory}; directly checked assembly record \nolinkurl{ThetaIntegralsAgree} and \nolinkurl{LocalMatching} in \nolinkurl{SiegelEisensteinSecondTermHypotheses.lean:44,69}. \\
4: determinant positivity & Positivity of a fixed Hermitian determinant is field-independent linear algebra. Identifying it with the target arithmetic sum is a separate global construction. & No new positivity theorem is needed at this level; no target determinant identity has been proved. The supplied map found no exact one-to-one source module for this historical step. \\
5: smoothed explicit formula & Replace rational integer coefficients \(\Lambda(n)\psi(n)\) by \(\Lambda_K(\mathfrak a)\chi(\mathfrak a)\), and rational conductor normalization by (K4). & Completed analytic replacement in AM-K005. Literal definitions are \nolinkurl{Analytic.lean:9,11,18,23} (\nolinkurl{abelian₂}, \nolinkurl{logDerivativeAtOne}, \nolinkurl{sumLamChi}, \nolinkurl{gChi}). \\
6: zero-free deduction & Apply a base-field-free pole theorem to the real induced L-function; do not substitute a Hecke character into \nolinkurl{PChar}. & Completed conditional deduction AM-K005. \nolinkurl{LandauZeroFreeCriterion.lean:15,25}; \nolinkurl{LandauTheorems/landauWindow.lean:83}. \\
7: periods and weights & The underlying CM/unitary setting is not restricted to \(\mathbf Q\). The new task is to choose the target ray-class character, test vectors and interlacing data so its period is exactly the coefficient functional needed in step 5. & Generic CM Galois row is visible in \nolinkurl{WeightsConjBis.lean:69} (\nolinkurl{Row0b}); it does not state a Hecke RH conclusion. Carried period anchor \nolinkurl{Formalism/PeriodIntegrals/BesselPeriodGGP.lean}. Target period identity open. \\
8: repeated definite lifts & Build an actual chain for the new ray-class family, with the six carried lift identities, controlled nonvanishing and correct rank costs. Keeping \(K\) CM is possible; using non-CM \(E\) as a CM base is excluded by AM-K002. & AM-K004 gives initial geometric descent data only. Direct assembly anchor \nolinkurl{nonempty_hyps_definite_side_of_closed_bound}, \nolinkurl{SiegelEisensteinSecondTermHypotheses.lean:44}, requires linked norm bounds and the other records. \\
9: norms and masses & Replace level/discriminant and local indices by those of the relative conductor, and include the unit and archimedean normalizations over \(K\). The cubic conductor is explicit in (K4). & Generic mass identities may be reused with hypotheses, but their evaluation for the new chain is open. Carried \nolinkurl{Formalism/PeriodIntegrals/UnitaryGroupMassFormula}; \nolinkurl{CovolumeBoundOfX} in the checked assembly theorem. \\
10: global divisibility & Local DP bounds remain available; the global sum must be normalized by \(\log N\mathfrak q\), residue degrees, bad-prime lattices and all embeddings. & No equality of the new global gain to the original gain is proved. Carried divisibility interfaces \nolinkurl{Formalism/SharpRestriction/Divisibility} and \nolinkurl{Formalism/LocalToGlobal/Divisibility}. \\
11: auxiliary tower & A relative ray-class tower over \(K\) must replace the rational cyclotomic fixed-field construction. The Kummer layer need not be CM as an absolute number field. & Literal substitution fails by AM-K002. \nolinkurl{ExistsPageN.lean:19,27,39} (\nolinkurl{IsChiField}, \nolinkurl{WinPair}, \nolinkurl{PageN}) explicitly require rational congruence and cyclotomic data. New tower with the necessary character separation and size bounds is open. \\
12: seesaw and terminal period & Seesaw/Rallis formalism over CM \(K\) is available at the generic level. The target family must satisfy the specified character independence and terminal-period normalization. & The explicit twist identity is not that terminal period identity. Actual checked bridge: \nolinkurl{ThetaIntegralsAgree}, \nolinkurl{LocalMatching}, \nolinkurl{SiegelEisensteinSecondTermHypotheses.lean:44,69}. Carried \nolinkurl{ThetaCorrespondence/RallisInnerProduct}. \\
13: slope inequality & The Arakelov slope inequality is already a number-field statement. Supply the actual metrized lattices, normalize degrees by the field degree, and include residue degrees in finite contributions. & New lattice/cost instance open. Carried source \nolinkurl{Formalism/LatticeSlopes/LatticeSlopesProofs/SlopeInequality.lean}; source assembly still needs \nolinkurl{CovolumeBoundOfX}. \\
14: cost evaluation & Compute every local and global term for the new lifts, using the actual conductor and period scaling (K12). The algebraic twist alone does not determine those terms or their cancellation. & No cost improvement or RH follows from the same geometric CM type. The supplied step-to-source map identifies \nolinkurl{UnitaryGroupMassFormula/UnipotentCost} only with low confidence; it is not treated as a complete formula for the target. \\
15: common integral structure & Replace rational primitive-character sets by finite-order ray-class characters over \(K\), and prove common coefficient-ring, saturation and congruence properties for their sections. & New host adapter required. \nolinkurl{Analytic.lean:102,106,145,149,154} (\nolinkurl{HostRaw.dir}, \nolinkurl{IsPrimOf} and analytic bounds) all quantify rational Dirichlet characters. AM-K002 excludes the naive norm-pullback route. \\
16: trace/orthogonality bound & The relative trace identity needs an actual ideal/ray-class indexing, stabilizers, unit factors and quantitative orbital count for the new host. & A formal orthogonality equality alone does not supply the needed scale-uniform bound. Carried \nolinkurl{Formalism/PeriodIntegrals/OrbitalBoundsConclusion}; no target version is proved. \\
17: cohomology control & The weighted Galois-cohomology size predicate \nolinkurl{VVle E t B} is generic in a number field \(E\). Its available supply theorem is indexed by rational \nolinkurl{PChar} and \nolinkurl{IsChiField}. & \nolinkurl{ExistsPageN.lean:12,41,44,49}; \nolinkurl{SupplyChi.lean:8}. A relative version over the new tower, with the same quantitative allowance, remains open. It is not supplied just by Dirichlet GRH. \\
18: covolumes and indices & Compute the actual determinant and finite-index lattice for the new global family; the conductor table fixes only part of its arithmetic input. & \nolinkurl{CovolumeBoundOfX} is a required hypothesis at \nolinkurl{SiegelEisensteinSecondTermHypotheses.lean:44,69}. No \nolinkurl{MainX} instance or covolume equality for \(\chi\) has been constructed. \\
19: combinatorial tower choice & Any abstract finite-set combinatorics must be realized with the new ray-class/tower parameters and the necessary cost and independence properties. & The supplied map §4 explicitly marks the use of \nolinkurl{Formalism/BehrendSets} in the final chain as uncertain. No extra Q-specific theorem is invented from that historical label. The precise checked tower restrictions remain \nolinkurl{IsChiField} and \nolinkurl{WinPair}. \\
20: final theorem & The literal endpoint is \nolinkurl{EveryCharacter}, a statement about \nolinkurl{DirichletCharacter ℂ N}. The replacement endpoint is the conclusion of AM-K005 once the target window is proved. & \nolinkurl{CharacterBoxes/EveryCharacter.lean:13}; \nolinkurl{ThirdPlaceSeries/TopPlug.lean:12,16}. The new arithmetic window remains unproved, so the requested Hecke RH conclusion is not established. \\
\end{longtable}

Two additional interface observations make the gate concrete. First, \nolinkurl{CharacterBoxes/Supply.lean:10} defines a generic request type \nolinkurl{FReq}, but its actual target request at line 42 is

\[
\texttt{GChi ψ}:
\quad\texttt{EAtH (Level.cyc ψ) H.toHostRaw TT s.x s.eps 4}
\ \land\ \texttt{VVle H.E t (sqrt s.b)}.
\]

The ambient request formalism is not the obstruction; the specific character/host supply is. Second, \nolinkurl{CharacterBoxes/ZeroFreeAllOfTopG.lean:12,26} separates the generic design/top-level family machinery from its application to \nolinkurl{PChar}. A relative version must produce a new target request and prove its supply, then connect it to (K13). Merely adding a new character type without those proofs would relocate the gap.

\subsection{8. Grothendieck-language formulation}\label{grothendieck-language-formulation}

\subsubsection{AM-K006. The finite étale and motivic dictionaries preserve the gate}\label{am-k006.-the-finite-uxe9tale-and-motivic-dictionaries-preserve-the-gate}

Let \(U=\operatorname{Spec}\mathbf Z[1/(rp)]\), \(U_K=\operatorname{Spec}\mathcal O_K[1/(rp)]\). The normalization \(U_E\to U_K\) is a finite étale \(\mu_r\)-torsor. Indeed \(p\) and \(r\) are units there; the derivative \(rT^{r-1}\) is invertible on \(T^r=p\), and \(K\) already contains \(\mu_r\). The character \(\chi\) is its Kummer local system. For coefficients containing \(\mu_r\), the idempotent

\[
e_\chi=\frac1r\sum_{g\in C_r}\chi(g)^{-1}g
\]

in the finite-cover correspondence algebra projects its rank-one Artin summand (with the chosen action convention). Every geometric Frobenius eigenvalue on that summand is a root of unity. Thus it is pointwise pure of weight zero. This is a statement about each closed fibre, not about zeros of the global L-function.

For the finite étale cover \(U_F\to U\), the permutation local system satisfies

\[
[\rho]=[h^0(U_F/U)]-[\mathbf1]
\]

in the Grothendieck group of finite-monodromy local systems, equivalently in the representation ring of \(G\). Its characteristic-polynomial identity is

\[
\det(1-gT\mid\rho)=\frac{\det(1-gT\mid P)}{1-T}.
\]

Exterior-power operations \(\lambda^j\) encode these determinants. AM-K001 and the carried CX-ALG-024 are therefore identities in the Artin-motivic \(K_0\) / representation-ring setting. AM-K002 says that this particular class is outside the subring generated by rational abelian Artin lines: every virtual object in that subring has trivial commutator action, whereas \(\rho\) does not. Adding tensor products, duals, exterior powers, symmetric powers or characteristic-zero direct summands of those abelian lines does not change that action. This is the elementary finite-group content of the carried GRH\_NIGHT AG-LAMBDA discussion, not a new global purity theorem.

For \(r=3\), AM-K004 gives the more useful relative realization over \(K\): the finite-monodromy line is the Hom of two CM elliptic eigensummands. The elliptic curves and their Picard duals extend as abelian schemes after deleting their bad primes. Their Poincaré biextensions give the usual tensor/duality pairings. The finite descent cocycle changes the Artin factor while leaving the geometric CM type and local height-two formal-group isomorphism unchanged after extension. The arithmetic integral lattice and descent factors remain essential for the cost.

\textbf{Exact framework locators.}

\begin{longtable}[]{@{}
  >{\raggedright\arraybackslash}p{(\columnwidth - 4\tabcolsep) * \real{0.3333}}
  >{\raggedright\arraybackslash}p{(\columnwidth - 4\tabcolsep) * \real{0.3333}}
  >{\raggedright\arraybackslash}p{(\columnwidth - 4\tabcolsep) * \real{0.3333}}@{}}
\toprule\noalign{}
\begin{minipage}[b]{\linewidth}\raggedright
Language
\end{minipage} & \begin{minipage}[b]{\linewidth}\raggedright
Locator and status
\end{minipage} & \begin{minipage}[b]{\linewidth}\raggedright
Meaning here
\end{minipage} \\
\midrule\noalign{}
\endhead
\bottomrule\noalign{}
\endlastfoot
Finite étale covers and finite G-sets & SGA 1, Exposé V, Théorème 4.1; primary reedition \url{https://arxiv.org/pdf/math/0206203} & The covering/local-system/representation dictionary used above. \\
Kummer torsors & SGA 1, Exposé XI §6, Proposition 6.1, Définition 6.2, Corollaires 6.3 and 6.5; reedition PDF pp.~248--250, original marginal pp.~303--305 & \(T^r=p\) on the base with \(rp\) inverted and its \(\mu_r\)-character. \\
\(\lambda\)-rings and \(K_0\) & SGA 6, Exposé V §§2 and 7, especially (7.1.1); carried GRH\_NIGHT AG-LAMBDA & Exterior-power determinant operations on Artin classes; no analytic RH follows from the class identity. \\
CM motives & Milne (2020b), §5, p.~45, for the framework; explicit new construction (K8)--(K11) here & The Artin Hom line belongs to a relative CM tensor category despite the non-CM splitting field. \\
Formal groups and crystals & Carried DP\_MULT DP10, DP6, DP32/DP34; Grothendieck cote 54 scans 5--7; cote 4 scans 5--6, 8.4.2 and Lemma 8.4.3, and scans 7--10 for torsion corrections & The local crystalline/Cartier-duality interpretation of the formal groups. No new monodromy result for cote 4 scan 12, §8.6, is asserted. The boxed zero restricted power in cote 54 concerns the vector kernel, as DP10 already established. \\
Picard schemes & FGA, Bourbaki exposé 232, Théorème 3.1; exposé 236, Corollaires 1.7 and 2.3, Proposition 2.10(ii), carried through DP\_MULT & Picard/dual abelian schemes for the elliptic curves; this gives geometric objects, not a \nolinkurl{MainX} arithmetic identity. \\
Biextensions & SGA 7 I, Exposés VII--VIII; carried DP\_MULT/THETA3 framework & Poincaré biextension and Cartier duality organize the pairings; no finer unverified proposition number is fabricated. \\
Weights versus global zeros & Carried CX\_GRH CX-G06 & Weight zero of a finite-monodromy sheaf does not imply RH over a number field. The missing input remains (K14). \\
\end{longtable}

The archive cote locators above are carried from the supplied DP\_MULT record, not newly reread scans in this child contribution. Their exact statements are not enlarged.

\subsection{9. Tried approaches and the remaining work}\label{tried-approaches-and-the-remaining-work}

\begin{longtable}[]{@{}
  >{\raggedright\arraybackslash}p{(\columnwidth - 4\tabcolsep) * \real{0.3333}}
  >{\raggedright\arraybackslash}p{(\columnwidth - 4\tabcolsep) * \real{0.3333}}
  >{\raggedright\arraybackslash}p{(\columnwidth - 4\tabcolsep) * \real{0.3333}}@{}}
\toprule\noalign{}
\begin{minipage}[b]{\linewidth}\raggedright
Approach
\end{minipage} & \begin{minipage}[b]{\linewidth}\raggedright
Result reached
\end{minipage} & \begin{minipage}[b]{\linewidth}\raggedright
Why it does not finish RH
\end{minipage} \\
\midrule\noalign{}
\endhead
\bottomrule\noalign{}
\endlastfoot
Instantiate the final theorem with \(\chi\) & Actual endpoint type read & \nolinkurl{EveryCharacter} quantifies rational Dirichlet characters. \\
Pull back a Dirichlet character by the norm & Refuted, AM-K002(2) & The target is anti-invariant under complex conjugation and has odd order. \\
Express the Artin L-function by finitely many rational Dirichlet factors & Refuted in the Euler-factor/Artin sense, AM-K002(3) & Nontrivial commutator trace. This does not refute every conceivable deduction from A\_D. \\
Insert \(E\) as the existing cyclotomic/CM host & Refuted, AM-K002(1) & \(E\) cannot embed in a CM field, whereas \nolinkurl{IsChiField} requires a rational cyclotomic embedding. \\
Realize the cubic character geometrically & Completed explicit Hom-line realization, AM-K004 & The integral lattice, period functional, nonvanishing chain, and cost/window identities are not consequences of the geometric twist. \\
Generalize the analytic Landau stage & Completed reduction, AM-K005 & All source \nolinkurl{PoleData} hypotheses verified; the arithmetic window (K14) remains. \\
Use known generic CM automorphic source rows & Type checked & \nolinkurl{Row0b} is already CM-general, but the proved outer endpoint and host supply remain rational-character-specific. \\
Use a newer cubic-Hecke zero-free source & Consulted only & OpenAI (2026), Theorem 1.1, claims \(\Re s>7/8\). Its long proof was not audited here and is not an accepted premise stipulated by the user. The downloaded \nolinkurl{ComparatorChallenges/HeckeSevenEighths.lean} is a comparator statement containing \nolinkurl{sorry}, not a proof certificate. Even the stated half-plane would not prove RH. \\
Use Masullo's claimed all-Hecke theorem & Conditional only & No Masullo premise is made unconditional. The carried GRH\_NIGHT gate Q-MAS-GEOMETRIC-NARROW remains unchanged. \\
\end{longtable}

\textbf{Masullo locators as supplied, without a new draft audit.} \nolinkurl{INTAKE_CX_GRH.md} records the all-Hecke claim in §§1.2 and 1.2.3, PDF pp.~12--16; Theorem 13.3.1.2 on PDF p.~522 is blank in that draft. The same intake records the relevant gaps at §9.3.2, Lemmas 9.3.2.4--9.3.2.5 and Theorem 9.3.2.6, pp.~425--426; §11.3, p.~495; and §§12.3--12.4. \nolinkurl{INTAKE_GRH_NIGHT_part3.md} carries the stronger narrowed gate. Those records are cited with their existing check status; no new proof of those draft claims is supplied here.

\textbf{Next concrete attack.} Build the relative host for the explicit cubic conductor in (K4), with \(\chi\) represented by (K8), and prove that its arithmetic functional is exactly (K13). Then compute the actual lift costs and common integral structure. Success must yield (K14) for every \(\theta>1/2\); an estimate at a fixed \(\theta>1/2\) is a partial result only. The previously accepted DP64 leading budget and the twist's unchanged geometric height preclude treating a mere relabelling of the same CM data as a newly computed cost saving.

\subsection{10. References and source records}\label{references-and-source-records}

\begin{itemize}
\tightlist
\item
  Alpöge, L. (2026). \emph{Remark}. \url{https://alpo.ge/remark.pdf}. Timeline steps 2--20, pp.~3--4; Lean archive as specified in the prompt. Actual declarations and pinned copies listed above and in \nolinkurl{kummer_sources/manifest.json}.
\item
  Grothendieck, A., et al.~(1971/2003). \emph{Revêtements étales et groupe fondamental (SGA 1)}. Lecture Notes in Mathematics 224; recomposed edition. Exposé V, Théorème 4.1; Exposé XI §6, Proposition 6.1, Définition 6.2, Corollaires 6.3, 6.5. \url{https://arxiv.org/pdf/math/0206203}.
\item
  Harron, R. (2016, May 24). \emph{The shapes of pure cubic fields}. §2.1, Lemma 2.1, p.~5. \url{https://math.hawaii.edu/~rharron/research/Harron-PureCubicShapes.pdf}.
\item
  Hida, H. (2025, March 27). \emph{Non-vanishing of integrals of a mod p modular form}. Proof of Corollary 6.18, p.~45, induced-conductor formula only. \url{https://www.math.ucla.edu/~hida/EllrevF.pdf}. No use of its separate nonvanishing assertions is made.
\item
  Milne, J. S. (2020a). \emph{Class field theory} (v4.03). V §2 and Theorem 2.4(b), p.~156; V Theorem 3.27 and Remark 3.28, pp.~165--166; VIII §10, pp.~261--264. \url{https://www.jmilne.org/math/CourseNotes/CFT.pdf}.
\item
  Milne, J. S. (2020b). \emph{Complex multiplication} (v0.10). Proposition 1.4 and Corollary 1.5, p.~10; §5, p.~45. \url{https://www.jmilne.org/math/CourseNotes/CM.pdf}.
\item
  OpenAI. (2026, September 30). \emph{The quasi-Riemann hypothesis: A zero-free half-plane Re(s)\textgreater7/8}. Theorem 1.1, consulted but not used as an established premise in this report. \url{https://github.com/openai/math/tree/main/preprints/The-Quasi-Riemann-Hypothesis-September-30-2026}. Saved TeX SHA-256: \nolinkurl{42a5ee0febca59fd1def55cfd6c6808c322ef4237d7726303ea52f711deac6a3}.
\item
  Supplied programme sources: \nolinkurl{INTAKE_CX_GRH.md}; \nolinkurl{INTAKE_GRH_NIGHT_part3.md}; \nolinkurl{GRH_NIGHT.tex}, CX-ALG-024 and AG-LAMBDA; \nolinkurl{DP_MULT.tex}, DP6/DP10/DP32/DP34/DP42/DP64; \nolinkurl{ALPOGE_LEAN_MAP_20261120.md}, §§3--4. These are carried under the user's stated check status.
\end{itemize}

\textbf{Status for integration.} AM-K000--AM-K006 are proved in their stated scope, with whole-proof author readings recorded separately. The requested RH theorem remains open at Q3-KUMMER-WINDOW. No unconditional deterministic-factoring conclusion is asserted from this contribution. The \(p=r\) finite-field construction in AM-K000 is independently unconditional.

\textbf{Subsequent concrete progress.} The companion kummer\_cubic\_host.md proves AM-K007--AM-K011: explicit cubic ray-class characters at every allowed prime; finite-Euler-factor invariance of the normalized window; source-compatible splitting transport and its telescoping limit; and the actual CM tower \(\mathbf Q(\zeta_{3\cdot7^n})\) with an invariant coprime half-different, positive unimodular Hermitian lattice, principally polarized products of the CM elliptic twists, and the retained Artin Hom line. Adding the split prime \(13\) supplies every mathematical field of the source's AbelianUnramifiedExtensionConditions, including \(8\mid f\). The addendum proves why a relative tower ramified at inert \(2\) cannot preserve that invariant-half-different condition, and identifies the target's nonsplit and wild ramification that its displayed IsMu0/QDat families do not directly allow. The complete HostRaw/MainX replacement and the arithmetic window remain open.

\clearpage

\section{Cubic Kummer host: finite character, splitting transport, and an explicit CM tower}\label{cubic-kummer-host-finite-character-splitting-transport-and-an-explicit-cm-tower}

\textbf{Run:} ALP\_MACHINE, question 3. \textbf{Author:} GPT-6 Astra Pro; inherited effort not separately exposed; child /root/kummer\_extension. \textbf{Date:} 2026-11-21. \textbf{Companion:} kummer\_extension.md, AM-K000--AM-K006. This addendum proves smaller arithmetic and geometric subgates. It does not assert the Kummer RH window or an instance of the archive's complete HostRaw/MainX records.

The inputs already accepted in the main contribution retain their status. In particular the conductor table AM-K003, the elliptic Hom realization AM-K004, the field-independent analytic adapter AM-K005, and CX-ALG-024 are used without reopening the stipulated Dirichlet theorem. The new field called \(B_n\) below is a CM coefficient/endomorphism field; it is not the non-CM Kummer splitting field \(K(p^{1/3})\).

\subsection{1. AM-K007: an explicit finite ray-class host for every allowed cubic prime}\label{am-k007-an-explicit-finite-ray-class-host-for-every-allowed-cubic-prime}

Put \(K=\mathbf Q(\zeta)\), \(\zeta=\zeta_3\), \(\mathcal O=\mathbf Z[\zeta]\), and \(\lambda=1-\zeta\). The ring \(\mathcal O\) is norm-Euclidean, has class number one, and has units \(\mu_6\). We use these established elementary facts; the source is Murty and Esmonde (2005), §2.3, Exercise 2.3.2 and its solution, together with §9.1. The norm is \[
N(a+b\zeta)=a^2-ab+b^2.
\]

\subsubsection{\texorpdfstring{1.1. The case \(p\ne3\), necessarily \(p\equiv2\pmod3\)}{1.1. The case p\textbackslash ne3, necessarily p\textbackslash equiv2\textbackslash pmod3}}\label{the-case-pne3-necessarily-pequiv2pmod3}

The reduction map \[
\mu_6\longrightarrow(\mathcal O/(3))^\times
\] is an isomorphism. The target has order \(9-3=6\). The only unit congruent to \(1\) modulo \((3)=(\lambda^2)\) is \(1\): the nontrivial cube roots have \(\lambda\)-adic difference from \(1\) of valuation one, and reduction modulo \(\lambda\) distinguishes the two possible signs. Thus every \(a\) prime to \(3\) has a unique unit \(\epsilon(a)\) with \(\epsilon(a)\equiv a\pmod3\). Define \(a^\circ=a/\epsilon(a)\).

For the modulus \(\mathfrak m=(3p)\), the ray-class exact sequence and class number one give \[
\operatorname{Cl}_{(3p)}(K)
 \simeq
 \bigl((\mathcal O/(p))^\times\times
             (\mathcal O/(3))^\times\bigr)/\mu_6
 \simeq \mathbf F_{p^2}^{\,\times}.
\tag{H1}
\] The last map sends the ideal class of \((a)\), with \(a\) prime to \(3p\), to \(a^\circ\bmod p\). It is independent of the generator: if \(u\in\mu_6\), then \(\epsilon(ua)=u\epsilon(a)\). Every residue class can be represented using the Chinese remainder theorem. The kernel is exactly the principal ideals with a generator congruent to \(1\) modulo \(3p\). This proves the displayed isomorphism, including its explicit formula. The exact sequence being used is Milne (2020a), V, Theorem 1.7, pp.~150--151.

Identify the three cube roots in \(\mathbf F_{p^2}\) with \(\mu_3\subset\mathbf C\) by the reduction of the fixed \(\zeta\). Then \[
\eta_p((a))=
 \left((a^\circ\bmod p)^{(p^2-1)/3}\right)^{\!\uparrow}
\tag{H2}
\] is a nontrivial order-three ray-class character. The superscript means the inverse of the chosen reduction identification, not an arbitrary numerical logarithm.

The primitive Kummer character \(\chi_{p,3}\) has conductor dividing \((3p)\) by AM-K003 and has exact order three. A cyclic group has exactly two nontrivial characters of order three, which are inverses. Therefore \[
\chi_{p,3}=\eta_p\quad\hbox{or}\quad\eta_p^{-1}
\tag{H3}
\] as ray-class characters, with the possible inverse accounting for the chosen Artin/Frobenius convention. AM-K001 proves that the two resulting L-functions are identical. Thus (H2) identifies the desired L-function without a new cubic-reciprocity input.

This also recovers the conductor distinction. The character descends to modulus \((p)\) exactly when the order-three character of \(\mathbf F_{p^2}^\times\) is trivial on the image of \(\mu_6\). Its restriction to \(\mu_3\) is trivial exactly when \(9\mid p^2-1\), or \(p\equiv8\pmod9\) under the standing congruence. The order-two unit causes no condition for an order-three character. Hence the primitive modulus is \((p)\) in that case; otherwise AM-K003 gives \((p)\lambda^2=(3p)\). The argument includes \(p=2\); no injectivity of \(\mu_6\to\mathbf F_{p^2}^\times\) was assumed.

Complex conjugation acts on (H1) by \(z\mapsto z^p\). Since \(p\equiv-1\pmod3\), it acts on \(\eta_p\) by inversion, as required.

\subsubsection{\texorpdfstring{1.2. The wild case \(p=3\)}{1.2. The wild case p=3}}\label{the-wild-case-p3}

For \(\mathfrak m=(9)=\lambda^4\), normalize again by the unique unit modulo \(3\). This gives \[
\operatorname{Cl}_{(9)}(K)
 \simeq (1+3\mathcal O)/(1+9\mathcal O)
 \simeq (\mathcal O/3\mathcal O,+)
 \simeq (\mathbf Z/3\mathbf Z)^2.
\tag{H4}
\] The middle map is \(1+3u\mapsto u\bmod3\); multiplication becomes addition because the cross term is divisible by \(9\). It is bijective, and the unit normalization proves the first isomorphism.

Write \(u=a+b\zeta\) modulo \(3\). Conjugation sends \((a,b)\) to \((a-b,-b)\). A character of this additive group has the form \(\zeta^{sa+tb}\). It is anti-invariant precisely when \[
s(a-b)+t(-b)=-sa-tb\quad\text{for every }a,b,
\] which is equivalent to \(s=0\). Thus the two nontrivial anti-invariant characters are \[
\eta_3(1+3(a+b\zeta))=\zeta^b,\qquad \eta_3^{-1}.
\tag{H5}
\] The actual Kummer character has conductor \((9)\) and is anti-invariant, by AM-K002 and AM-K003. It must be one of these two. Again AM-K001 identifies their L-functions. This explicitly includes the wild conductor; it is not a tame specialization.

\subsubsection{1.3. Character separation and the ideal-indexed Dirichlet series}\label{character-separation-and-the-ideal-indexed-dirichlet-series}

Let \(H=\operatorname{Cl}_{\mathfrak m}(K)\) with the modulus just used. Its regular module over \(\mathbf Q(\zeta)\) has the exact projector \[
e_{\eta}=\frac1{|H|}\sum_{h\in H}\eta(h)^{-1}[h].
\tag{H6}
\] Translation by an ideal class \([\mathfrak a]\) acts on its image by \(\eta(\mathfrak a)\). The elementary identity \[
1+\eta(h)+\eta(h)^{-1}
   =3\,\mathbf 1_{\eta(h)=1}
\tag{H7}
\] gives the complete finite character separation. These statements follow directly by changing the index in (H6) and summing the three roots of unity.

A smaller integral realization uses the quotient \(H/\ker\eta\simeq C_3\): the rank-two augmentation representation has trace \(2\) on the identity and \(-1\) otherwise. Over \(\mathbf Q(\zeta)\) it splits as \(\eta\oplus\eta^{-1}\). The projector onto one line uses division by \(3\), so an integral construction must retain that denominator; (H6) is not asserted to be an integral idempotent at primes dividing \(|H|\).

For \(\Re s>1\), the ray-modulus L-function is exactly \[
L_{\mathfrak m}(s,\eta)
 =\frac16\sum_{(a,b)\in\mathbf Z^2\setminus\{(0,0)\}}
 \frac{\eta((a+b\zeta))}{(a^2-ab+b^2)^s},
\tag{H8}
\] where the numerator is zero when \((a+b\zeta)\) is not coprime to \(\mathfrak m\). Each nonzero integral ideal has precisely six generators, which proves the formula by absolute convergence. This is a concrete ideal-indexed host. It does not turn the series into a rational Dirichlet L-function.

\subsubsection{1.4. The harmless extra Euler factor is harmless for the actual window}\label{the-harmless-extra-euler-factor-is-harmless-for-the-actual-window}

When \(p\equiv8\pmod9\), the convenient modulus \((3p)\) omits the unramified Euler factor at \(\lambda\) of norm \(3\). Since \(\lambda\) is fixed by conjugation and the character is odd-order and anti-invariant, \(\chi(\lambda)=1\).

More generally, remove an unramified Euler factor of norm \(q>1\) and value \(u\), \(|u|=1\): \[
L_{\rm imp}(s)=(1-uq^{-s})L_{\rm prim}(s).
\] Use the same normalized ramp as AM-K005 for each function, including its own constant \(L'/L(1)\). With \(J=\lfloor\log_q x\rfloor\), direct subtraction gives \[
G_{\rm imp}(x)-G_{\rm prim}(x)
 =\log q\left(
     \sum_{k=1}^{J}u^k+
     \frac{x(u/q)^{J+1}}{1-u/q}
            \right).
\tag{H9}
\] Indeed the removed prime powers subtract \(\log q\sum_{k\le J}u^k(x/q^k-1)\), and the change in \(L'/L(1)\) adds \(x\log q\sum_{k\ge1}(u/q)^k\). Their finite sums cancel to (H9). Consequently the difference is \(O(1+\log x)\) for fixed \(q\).

For the present real-valued ramps, their AM-K005 windows \(W(\theta,\delta)\) are equivalent for every \(\theta>0\). To prove this with the actual quantifiers, start with a majorant for one ramp at exponent \(\eta/2\), and add \(C(1+\log x)\) to majorize the other. The added integral over \([Y,Y^{1+\delta}]\) is at most \[
C\int_Y^\infty (1+\log x)x^{-\theta-1}\,dx
=O_\theta(Y^{-\theta}(1+\log Y)).
\] Increase the prescribed lower threshold so that this and \(Y^{\eta/2}\) are each at most \(Y^\eta/2\). The same argument reverses the roles. Thus the explicit nonminimal modulus in (H1) introduces no new analytic gate.

\textbf{Closed subgate.} For all allowed cubic \(p\), the finite ray-class group, the actual character up to its L-equivalent inverse, its character projector, and the ideal-indexed L-series are explicit. No Kummer extension has been inserted as a CM field.

\subsection{2. AM-K008: the Kummer character is an admissible splitting twist}\label{am-k008-the-kummer-character-is-an-admissible-splitting-twist}

This statement works for every odd \(r\), not only \(3\). Put \(K^+=\mathbf Q(\zeta_r+\zeta_r^{-1})\), let \(c\) denote CM conjugation, and let \(\chi\) be the finite-order Kummer Hecke character. AM-K002 gives \(\chi^c=\chi^{-1}\).

For \(a\in\mathbb A_{K^+}^{\times}\), conjugation fixes \(a\), so \(\chi(a)^2=1\). Its order also divides the odd number \(r\), hence \[
\chi|_{\mathbb A_{K^+}^{\times}}=1.
\tag{H10}
\] At infinity it is trivial because a continuous finite-order character of each connected group \(\mathbf C^\times\) is trivial.

Let \(T=\operatorname{Res}_{K/K^+}^{1}\mathbf G_m\) and choose \(j\) with \(2j\equiv1\pmod r\), for example \(j=(r+1)/2\). On \(T(\mathbb A_{K^+})\) define \[
\nu(z)=\chi(z)^j.
\] It is unitary, continuous, finite-order, and trivial on \(T(K^+)\subset K^\times\). Moreover \[
\nu(x/cx)=\chi(x)^{2j}=\chi(x).
\tag{H11}
\] For \(r=3\), one can take \(j=2\). For a unitary group \(U(S)\) write \(\nu_S(g)=\nu(\det g)\).

\subsubsection{2.1. Exact transport of the oriented Weil record}\label{exact-transport-of-the-oriented-weil-record}

Suppose an actual oriented Weil datum over the CM extension \(K/K^+\) is given, with splitting characters \(\chi_S,\chi_T\) and action \(\omega\). Let \(a,b\in\mathbf Z\). Set \[
\chi_S'=\chi_S\chi^a,\qquad
\chi_T'=\chi_T\chi^b,\qquad
\omega'(g,h)=\nu_S(g)^{-b}\nu_T(h)^a\,\omega(g,h).
\tag{H12}
\] Then it satisfies the same oriented Weil-datum conditions, with the new splitting characters, on the same Heisenberg space and with the same additive character.

Here is a check of every field in the actual source record, not just an appeal to the name of the representation. In LocalNewformFunction.lean, the record IsWeilDatumO is at lines 53--75.

\begin{itemize}
\tightlist
\item
  \textbf{std:} the additive character is unchanged.
\item
  \textbf{hecS, hecT, resS, resT:} products remain Hecke characters; (H10) preserves the prescribed restriction to the real subfield.
\item
  \textbf{om\_pres:} a constant scalar multiple preserves each defining quasi-periodicity, level, continuity and archimedean smoothness condition of IsThetaFunO.
\item
  \textbf{om\_heis:} the scalar commutes with every linear Heisenberg operator, leaving the same symplectic action.
\item
  \textbf{om\_aut:} on rational unitary elements the determinants lie in \(T(K^+)\), where \(\nu=1\).
\item
  \textbf{om\_cont:} the new scalar is continuous; it multiplies the old continuous matrix coefficient.
\item
  \textbf{om\_unit:} the scalar has absolute value one, so every displayed squared-norm integral is identical.
\item
  \textbf{om\_split:} the source splitting factor, line 51, is \(\chi_T(x)^{-1}\chi_S(y)\). The new factor is this times \(\chi(x)^{-b}\chi(y)^a\). For a source determinant root, \(x/cx=\det g\) and \(y/cy=\det h\), so (H11) makes these exactly the scalar in (H12). The same positive real constant works.
\end{itemize}

For the last check, the source's split-place definition IsDetAt, TiedWeil.lean:77, also implies the determinant-root identity. At the selected idempotent \(e\) it reads \(x=e\det g+(1-e)\); the supported unitary element is the identity off \(e+ce\) and satisfies \(c(\det g)=(\det g)^{-1}\). Multiplying in the two orthogonal idempotent components gives \(x/cx=\det g\).

Multiplication by the scalar character also preserves the group-homomorphism property of \(\omega\), so the ambient type of the new action is correct. This completes the transport proof.

The related actual interface TiedWeil.lean:64,83--96 defines IsDetRoot and TwRule. Its four formulas explicitly permit a splitting-character change with trivial restriction to the real base, with the sign reversed on the reversed-rank edge. TwRule requires its character to belong to the supplied family X.Chr; membership of the new character in an existing family is not deduced from that record. Formula (H12) uses the oriented convention splitFacO, and therefore keeps its signs fixed. A reversal must use the source's corresponding formula rather than silently reuse a sign.

\subsubsection{2.2. Level and integral scope}\label{level-and-integral-scope}

The archimedean representation and archimedean norms are unchanged. At a finite place where \(\chi\) is unramified, \(\nu_S\) is trivial on every compact subgroup of \(U(S)\): compact determinant images have zero valuation, and an unramified character is trivial on local units. At each remaining place, replacing a compact level \(U_v\) by \(U_v\cap\ker\nu_S\) costs group index at most \(r\).

For \(r=3\), only the rational places \(p,3\) can contribute. Thus the product of these elementary local kernel indices is at most \(9\) for \(p\ne3\), and at most \(3\) for \(p=3\). This is solely the index needed to kill this one character on a given level. It is not the complete adelic covolume, a Hecke-orbit count, or the cost of the archive's arithmetic columns. Root-of-unity multiplication preserves an integral coefficient lattice when that lattice is a module over a ring containing those roots; projectors and saturation may introduce further denominators.

\textbf{Closed subgate.} The target character has exactly the relative restriction needed for the archive's splitting-change formulas. This is an actual transport theorem for a supplied Weil datum. It does not manufacture the required D7 datum or the global arithmetic theta comparison.

\subsection{3. AM-K009: coherent twists telescope and identify the remaining comparison}\label{am-k009-coherent-twists-telescope-and-identify-the-remaining-comparison}

Fix a chain of groups \(G_0,\ldots,G_m\) and kernels \(K_i(g_i,g_{i+1})\). State the convention explicitly: \[
(\Theta_i f)(h)=\int_{[G_i]}K_i(g,h)f(g)\,dg,
\] on any function space where these integrals exist. Let \(\nu_i\) be the determinant character from AM-K008 and \(t_i\in\mathbf Z\). Define \(M_i f=\nu_i^{t_i}f\) and \[
K_i'(g,h)=\nu_i(g)^{-t_i}\nu_{i+1}(h)^{t_{i+1}}K_i(g,h).
\] Then \[
\Theta_i'=M_{i+1}\Theta_iM_i^{-1},\qquad
\Theta_{m-1}'\cdots\Theta_0'
 =M_m(\Theta_{m-1}\cdots\Theta_0)M_0^{-1}.
\tag{H13}
\] Both formulas follow by substitution and cancellation of the adjacent multiplication operators. With a convention conjugating the input function, the corresponding powers are conjugated; the same cancellation remains. Each \(M_i\) is unitary on \(L^2\), and therefore norms and pairings of simultaneously transported vectors are unchanged.

Consequently, coherent splitting transport alone supplies no new prime-weighted arithmetic identity. It gives isomorphic analytic operators and the transported endpoints. To obtain the new explicit-formula functional one must identify a character-dependent arithmetic column, Hecke action or endpoint that carries (H2)/(H5), and verify its integral structure and its actual cost.

This distinction can be seen in the checked source, rather than inferred from broad terminology:

\begin{longtable}[]{@{}
  >{\raggedright\arraybackslash}p{(\columnwidth - 4\tabcolsep) * \real{0.3333}}
  >{\raggedright\arraybackslash}p{(\columnwidth - 4\tabcolsep) * \real{0.3333}}
  >{\raggedright\arraybackslash}p{(\columnwidth - 4\tabcolsep) * \real{0.3333}}@{}}
\toprule\noalign{}
\begin{minipage}[b]{\linewidth}\raggedright
Actual source interface
\end{minipage} & \begin{minipage}[b]{\linewidth}\raggedright
What it contains
\end{minipage} & \begin{minipage}[b]{\linewidth}\raggedright
What (H12) supplies
\end{minipage} \\
\midrule\noalign{}
\endhead
\bottomrule\noalign{}
\endlastfoot
LevChiEqOne.lean:23--64, TwDat and TwDat.Is & Level charts, stabilizers, arithmetic representatives, determinant roots, and the exact character relation to the common family & The permitted determinant character; it does not construct the charts or representatives \\
Same \url{file:73–77}, abelian₂ & A prescribed product of arithmetic scalar factors times jetL of the actual abelian theta section & No new equality with a prime sum follows from scalar transport \\
Same \url{file:87–101}, VacT and PinX & The moved-vacuum theta identity, its nonzero constant, fixed positive lattice pin, and the entire family of matching characters & Once supplied, the splitting factors transform by (H12); the arithmetic identity still needs proof for the new family \\
Same \url{file:103–115}, PinColT & Finite combination datum, rational datum, and a constant combining the level volume, arithmetic weights, vacuum constant and determinant-root evaluations & Individual phase factors are explicit; the full nonzero constant and covolume are not evaluated here \\
\nolinkurl{CharacterBoxes/Supply.lean:10,42} & FReq is a predicate on Scale, Frame and HostAx; GChi adds EAtH and VVle & ``Generic request'' is generic only within this HostAx domain, not over arbitrary CM fields \\
\end{longtable}

A precise remaining task is therefore the \textbf{twisted arithmetic comparison} for the explicit ray-class character and the actual CM lattice below: construct the arithmetic section and prove that its functional is AM-K005's \(G_\chi\), with all local multiplicities, conductor factors and integral denominators. That is narrower than leaving the character, finite host and local splitting laws unspecified. It still includes a substantive theorem; this report does not rename it as a proved window.

\subsection{4. AM-K010: an explicit CM supply tower with the induced cubic type}\label{am-k010-an-explicit-cm-supply-tower-with-the-induced-cubic-type}

The previous finite host is a ray-class group, whereas the geometric construction also requires a CM field, polarization and lattice. The following explicit tower supplies several precise HostDer-shaped ingredients while keeping the Kummer character in the descent of a CM abelian variety.

\subsubsection{4.1. A criterion for the invariant half-different}\label{a-criterion-for-the-invariant-half-different}

Let \(B/K\) be a finite Galois extension of a CM field \(K\), with \(B\) CM and conjugation \(c\) restricting to conjugation of \(K\). Write \(\mathfrak D=\mathfrak D_{B/K}\). There exists a Gal\((B/K)\)-stable integral ideal \(\mathfrak b\) with \[
\mathfrak D=\mathfrak b\,c(\mathfrak b),\qquad
\mathfrak b+c(\mathfrak b)=\mathcal O_B
\tag{H14}
\] if and only if \(B/K\) is unramified at every finite prime \(\mathfrak q\) of \(K\) fixed by \(c\).

For necessity, let \(\mathfrak q=c\mathfrak q\). The Galois group is transitive on primes above \(\mathfrak q\), so stability makes the valuation of \(\mathfrak b\) the same integer at all those primes. Conjugation permutes the same set, so \(c(\mathfrak b)\) has the same valuation there. Coprimality forces that integer to be zero, hence \(\mathfrak D\) has valuation zero at all these primes. A finite number-field extension is unramified at a prime exactly when its local different is trivial.

For sufficiency, partition the non-conjugation-fixed primes in the support of the relative discriminant into pairs \(\{\mathfrak q,c\mathfrak q\}\), and choose one member of every pair. Put into \(\mathfrak b\) all prime factors of \(\mathfrak D\) above the chosen members, with their full different exponents. Since the extension is Galois, \(\mathfrak D\) and this choice are Gal\((B/K)\)-stable. Conjugation preserves the relative different and supplies precisely the complementary factors. This proves (H14). No class-number-one assumption on \(B\) is required.

For \(K=\mathbf Q(\zeta_3)\), both the prime over \(2\) and the prime over \(3\) are conjugation-fixed. Thus an induced-type invariant half-different of this form excludes relative ramification at \(2\) or \(3\). In particular a literal relative extension ramified at \(2\), as in an added nontrivial real \(2\)-power cyclotomic tower, cannot have all these properties with this distinguished \(K\). This is a specific obstruction to retaining the archive's \(2\)-adic tower, not an obstruction to all CM host fields containing \(K\).

\subsubsection{4.2. The explicit tower}\label{the-explicit-tower}

For each \(n\ge1\) set \[
B_n=\mathbf Q(\zeta_{3\cdot7^n})
   =K(\zeta_{7^n}),\quad
d_n=6\cdot7^{n-1},\quad
a_n=7^{n-1}(6n-1).
\tag{H15}
\] Then \([B_n:K]=d_n\), and \(B_n\) is a Galois CM field of degree \(2d_n\). Its Galois group over \(\mathbf Q\) is \[
(\mathbf Z/3\mathbf Z)^\times
 \times(\mathbf Z/7^n\mathbf Z)^\times.
\] Complex conjugation is \((-1,-1)\). At \(3\), inertia is contained in the first factor; at \(7\), it is contained in the second; away from \(3,7\) it is trivial. Conjugation is in neither of the two nontrivial inertia groups. The inertia group for \(B_n/B_n^+\) is their intersection with \(\{1,c\}\), and is therefore trivial at every finite prime. Hence \[
B_n/B_n^+\ \text{is everywhere unramified at finite primes}.
\tag{H16}
\]

The prime \(7\) splits in \(K\): \[
(7)=(3+\zeta)(3+\zeta^2).
\] The two factors have norm \(7\) and are distinct. Each ramifies totally in \(B_n/K\), yielding primes \(P_n,cP_n\) of residue degree one and ramification degree \(d_n\). There is no other relative ramification.

The local ring is generated by \(\zeta_{7^n}\). The derivative of \[
\Phi_{7^n}(X)=\frac{X^{7^n}-1}{X^{7^{n-1}}-1}
\] at this root has valuation \[
n\,d_n-7^{n-1}=a_n,
\] because \(v_{P_n}(7)=d_n\) and \(v_{P_n}(\zeta_{7^n}^{\,7^{n-1}}-1)=7^{n-1}\). The monogenic different formula therefore gives \[
\mathfrak D_{B_n/K}=P_n^{a_n}(cP_n)^{a_n}.
\tag{H17}
\] Thus \(\mathfrak b_n=P_n^{a_n}\) is Gal\((B_n/K)\)-stable and satisfies (H14). This is a fully specified half-different, not an existence assertion for an unspecified field.

Discriminant transitivity and (H16) yield \[
|D_{B_n}|=3^{d_n}7^{2a_n},\qquad
D_{B_n^+}=3^{d_n/2}7^{a_n},\qquad
\operatorname{rd}(B_n)=\sqrt3\,7^{\,n-1/6}.
\tag{H18}
\] The roots of unity in \(B_n\) are exactly \(\mu_{6\cdot7^n}\). One elementary check is by ramification: any root order can involve only \(2,3,7\); a factor \(4\) would ramify at \(2\), a factor \(9\) would require ramification index at least \(6\) at \(3\) instead of \(2\), and a higher \(7\)-power would require larger ramification index at \(7\). The stated roots already lie in \(B_n\).

Every one of these roots has relative norm one to \(K\). Indeed \[
N_{B_n/K}(-1)=1,\quad
N_{B_n/K}(\zeta_3)=\zeta_3^{d_n}=1,\quad
N_{B_n/K}(\zeta_{7^n})=1.
\tag{H19}
\] The last equality is the constant term of the even-degree cyclotomic polynomial. These three roots generate the full group. This verifies the actual norm-of-roots condition, not the stronger and false assertion that \(B_n\) has only \(\pm1\) as roots.

The cyclotomic facts used here are Milne (2020c), Theorem 6.4, printed p.~99 (PDF p.~101), and Remark 6.6(b), printed p.~100 (PDF p.~102), together with the elementary local cyclotomic derivative computation displayed above. The prime splitting and all discriminant exponents have been shown explicitly.

\subsubsection{4.3. A positive self-dual Hermitian lattice}\label{a-positive-self-dual-hermitian-lattice}

Set \[
L_n=\mathfrak b_n^{-1},\qquad
h_n(x,y)=\operatorname{Tr}_{B_n/K}(x\,c(y)).
\tag{H20}
\] It is a free \(\mathcal O_K\)-module of rank \(d_n\), because \(\mathcal O_K\) is a PID. Its Hermitian dual is \[
L_n^\#=\mathfrak D_{B_n/K}^{-1}\,(cL_n)^{-1}
       =(\mathfrak b_n c\mathfrak b_n)^{-1}c\mathfrak b_n
       =\mathfrak b_n^{-1}=L_n.
\tag{H21}
\] The first equality follows from the definition of the inverse different as the trace-dual of \(\mathcal O_{B_n}\), extended to fractional ideals; it can also be checked after localization at each prime. Thus \(h_n\) is integral and unimodular on \(L_n\).

Choose the CM type \(\Phi_n\) consisting of the \(d_n\) embeddings extending a fixed embedding \(K\hookrightarrow\mathbf C\). Then \[
h_n(x,x)=\sum_{\sigma\in\Phi_n}|\sigma(x)|^2
\] under that embedding, so the form is positive definite. Choose an \(\mathcal O_K\)-basis and use the matrix convention \(h_n(x,y)=\bar y^{\,t}H_nx\) for coordinate columns. Then \(H_n\) is integral Hermitian and invertible over \(\mathcal O_K\). Its determinant is a positive rational unit, hence \(1\).

This exactly supplies the mathematical shape of HostDer.principal and HostDer.unimodular for this induced type, together with the unramified CM quotient and discriminant square. It supplies them as new mathematical objects, not as an instance of the existing parameterized HostRaw record.

\subsubsection{4.4. Actual CM abelian varieties and the cubic Hom line}\label{actual-cm-abelian-varieties-and-the-cubic-hom-line}

Use the elliptic curves from AM-K004: \[
A_0:y^2=x^3+1,\qquad A_p:y^2=x^3+p^2.
\] Both have \(\mathcal O_K\)-endomorphisms over \(K\). The chosen \(\mathcal O_K\)-basis of \(L_n\) realizes multiplication by \(\mathcal O_{B_n}\) as a subring of \(M_{d_n}(\mathcal O_K)\). Consequently it acts, over \(K\), on the actual products \[
\mathcal A_{0,n}=A_0^{d_n},\qquad
\mathcal A_{p,n}=A_p^{d_n}.
\] Give each product the polarization represented relative to its product principal polarization by the same positive integral Hermitian matrix \(H_n\). This is a principal polarization: positivity is the Riemann-form criterion, and unimodularity makes the polarization morphism an isomorphism. The entries are endomorphisms defined over \(K\), so the morphism is defined over \(K\). The adjoint of multiplication by \(b\in B_n\) for \(h_n\) is multiplication by \(c(b)\), proving that Rosati induces CM conjugation.

For completeness, the complex construction can be written directly as \[
\mathbf C^{\Phi_n}/\Phi_n(L_n).
\] Choose \(\xi=-1/\sqrt{-3}\) with the fixed embedding \(\sqrt{-3}\mapsto i\sqrt3\), so \(\operatorname{Im}\sigma(\xi)>0\) for \(\sigma\in\Phi_n\). The alternating form \(\operatorname{Tr}_{B_n/\mathbf Q}(\xi x c(y))\) is integral, positive in the Riemann convention and self-dual: the inverse different of \(K/\mathbf Q\) is generated by \(\xi\), and (H21) handles the relative trace. This produces the same polarized isomorphism class over \(\mathbf C\) as the Hermitian matrix construction, after scaling the common elliptic period. The Riemann-form and CM-pair framework is Milne (2020b), §2, Example 2.9, Example 2.11, and Proposition 3.17; Proposition 3.12 describes the induced CM isogeny type.

The coordinatewise isomorphism \[
g_p^{d_n}:\mathcal A_{p,n}\longrightarrow\mathcal A_{0,n}
\] over \(K(p^{1/3})\) is \(\mathcal O_{B_n}\)-linear because the scalar elliptic twist commutes with every matrix over \(\mathcal O_K\). It respects the displayed polarizations for the same reason. Its descent cocycle is the scalar \(\chi_{p,3}\) in \(\mathcal O_K^\times\subset\mathcal O_{B_n}^\times\).

On any one-dimensional \(B_n\)-eigensummand of \(H^1\) extending the chosen \(\zeta\)-eigensummand over \(K\), the ratio of the two Galois characters is therefore exactly \(\chi_{p,3}\), by the AM-K004 intertwining computation. Thus the explicit rank-one Artin Hom line survives in a principally polarized CM model with the lattice (H20). The splitting field is used only to trivialize the descent cocycle; it is not asserted to be a CM field.

The stable Faltings heights of the underlying products are \[
h_{\rm Fal}(\mathcal A_{0,n})=d_n h_{\rm Fal}(A_0),\qquad
h_{\rm Fal}(\mathcal A_{p,n})=d_n h_{\rm Fal}(A_0).
\tag{H22}
\] This uses additivity for products and the geometric isomorphism of \(A_p\) and \(A_0\). It is independent of the chosen principal polarization. It is not an evaluated theta-lift cost: derivatives, indices, normalized kernel constants and saturation still have to be computed. In particular, on the \(d_n\)-dimensional \(\zeta\)-eigenspace the determinant twist is \(\chi_{p,3}^{d_n}=1\) because \(3\mid d_n\), while each chosen \(B_n\)-eigensummand still has nontrivial character; a determinant or height alone does not recover that summand.

\subsubsection{4.5. Precisely which source requirements have and have not been supplied}\label{precisely-which-source-requirements-have-and-have-not-been-supplied}

\begin{longtable}[]{@{}
  >{\raggedright\arraybackslash}p{(\columnwidth - 4\tabcolsep) * \real{0.3333}}
  >{\raggedright\arraybackslash}p{(\columnwidth - 4\tabcolsep) * \real{0.3333}}
  >{\raggedright\arraybackslash}p{(\columnwidth - 4\tabcolsep) * \real{0.3333}}@{}}
\toprule\noalign{}
\begin{minipage}[b]{\linewidth}\raggedright
Source requirement
\end{minipage} & \begin{minipage}[b]{\linewidth}\raggedright
Supplied here
\end{minipage} & \begin{minipage}[b]{\linewidth}\raggedright
Remaining difference
\end{minipage} \\
\midrule\noalign{}
\endhead
\bottomrule\noalign{}
\endlastfoot
ConjEMemHostFIff.lean:89,95, CM field and CM type & \(B_n\) and the type induced from \(K\) & The original record has its own Frame and distinguished quadratic-field parameters \\
Same \url{file:135}, finite unramified CM quotient & (H16) & Proven for the new tower, not transported through HostRaw \\
Same \url{file:156}, invariant coprime half-different & (H17), \(\mathfrak b_n=P_n^{a_n}\) & Its tower prime is \(7\); retaining ramification at inert \(2\) violates §4.1 \\
Same \url{file:159}, unimodular trace lattice & (H20)--(H21) & Actual later arithmetic columns still need matching to this lattice \\
Same \url{file:163}, discriminant square & (H18) & The full source root-discriminant inequalities depend on its different parameters \\
Same \url{file:170}, norm of every root of unity is one & (H19) & Its separate stronger MuPM predicate fails here; no use of that predicate may be carried silently \\
Main CM abelian variety and cubic eigensummand & Actual \(\mathcal A_{0,n},\mathcal A_{p,n}\) with principal polarizations and the Hom line & No full MainX chain, terminal vector or normalized period has been constructed \\
ExistsPageN.lean:19,39 and Supply.lean:42 & No literal instance claimed & The original cyclotomic subgroup, \(2\)-adic PageN scale, character-specific EAtH and cohomology supply require a \(7\)-adic/relative replacement \\
\end{longtable}

The phrase ``generic CM theory'' therefore conceals two separate issues: several geometric objects now exist explicitly, but the original rational congruence/tower supply and its arithmetic prime functional are not consequences of that existence.

\subsection{5. AM-K011: the actual generic field record and its finite-place boundary}\label{am-k011-the-actual-generic-field-record-and-its-finite-place-boundary}

The source contains a geometric field interface weaker than HostRaw. It can be supplied completely at the level of mathematical data. This closes a named source input; the source's arithmetic character family still needs a separate extension.

\subsubsection{5.1. A degree-divisibility adjustment using a second split prime}\label{a-degree-divisibility-adjustment-using-a-second-split-prime}

Set \[
B_n^*=\mathbf Q(\zeta_{3\cdot7^n\cdot13}),\qquad
F_n^*=(B_n^*)^+,\qquad
f_n=[B_n^*:K]=72\cdot7^{n-1}.
\tag{H23}
\] The extra prime \(13\) splits in \(K\), since \((13)=(4+\zeta)(4+\zeta^2)\). Cyclotomic disjointness gives the displayed degree, so \(8\mid f_n\). The CM quotient \(B_n^*/F_n^*\) is unramified at all finite primes: inertia at \(3\), \(7\) or \(13\) is confined to that factor of the cyclotomic Galois group, whereas conjugation is \(-1\) in all three factors. Thus it lies in none of the inertia groups.

Relative ramification over \(K\) is confined to the two split rational primes \(7,13\). The half-different criterion in AM-K010 applies. More explicitly, choose the primes \((3+\zeta)\) and \((4+\zeta)\) of \(K\), and put \[
\mathfrak b_n^*
  =\prod_{P\mid(3+\zeta)}P^{a_n}
     \prod_{Q\mid(4+\zeta)}Q^{11}.
\tag{H24}
\] Adjoining the other prime-power roots is unramified locally, so the different exponent at \(7\) remains \(a_n\) and that at \(13\) is the tame exponent \(12-1=11\). This proves \(\mathfrak D_{B_n^*/K}=\mathfrak b_n^*c(\mathfrak b_n^*)\) with coprime, Gal\((B_n^*/K)\)-stable halves. Its inverse is therefore again a free positive unimodular Hermitian \(\mathcal O_K\)-lattice, now of rank \(f_n\). All the principally polarized product and Artin Hom constructions of §4.4 apply with this lattice.

The roots of unity are \(\mu_{6\cdot7^n\cdot13}\). Their relative norms are one: the norm of \(\zeta_3\) is one because \(3\mid f_n\); the norms of primitive \(7^n\)-th and \(13\)-th roots are powers of the constant terms of their even-degree cyclotomic polynomials; and \(f_n\) is even. Coprime-conductor discriminant multiplication gives \[
|D_{B_n^*}|=3^{12d_n}7^{24a_n}13^{22d_n},
\qquad
\operatorname{rd}(B_n^*)=\sqrt3\,7^{\,n-1/6}13^{11/12}.
\tag{H25}
\] Thus satisfying the degree condition does not require introducing relative ramification at inert \(2\).

\subsubsection{5.2. Every field of AbelianUnramifiedExtensionConditions}\label{every-field-of-abelianunramifiedextensionconditions}

The exact file is \nolinkurl{Formalism/LocalToGlobal/Columns/AbelianHostFieldConditions.lean.} Its Fields structure is at lines 19--42, Fields.Ok at lines 81--108, and AbelianUnramifiedExtensionConditions at lines 110--120. To avoid its letter K being confused with the fixed quadratic field, denote its coefficient field by \(C\).

Choose the source fields \[
X.F=F_n^*,\qquad X.E=B_n^*,\qquad C=X.K=B_n^*,
\] with their natural embeddings and conjugations. Take \(X.\mathrm{Tau}=\operatorname{Gal}(B_n^*/K)\), and let its embeddings be the corresponding automorphisms of \(B_n^*\). Then every Fields.Ok condition holds:

\begin{itemize}
\tightlist
\item
  \textbf{realF, rankE, star\_F, conjE, cplxE:} these are the CM-field identities for \(B_n^*/F_n^*\).
\item
  \textbf{emb\_star:} the cyclotomic automorphisms commute with conjugation.
\item
  \textbf{card\_Tau:} the cardinal is \(f_n=[F_n^*:\mathbf Q]\).
\item
  \textbf{emb\_inj:} distinct automorphisms are distinct embeddings.
\item
  \textbf{emb\_ne:} an embedding fixing \(K\) cannot equal a conjugated one, whose restriction to \(K\) is nontrivial.
\item
  \textbf{emb\_all:} every embedding \(B_n^*\to B_n^*\) is a Galois automorphism. Its restriction to \(K\) is either identity or conjugation, so it lies in the chosen set or its conjugate.
\item
  \textbf{cplxK, normalK, conjK:} the coefficient field \(C=B_n^*\) is a normal CM field, with its displayed complex embedding and conjugation.
\item
  \textbf{quad:} the source's Kq consists of elements whose images under all chosen embeddings agree. Since the identity embedding occurs, this is the fixed field of Gal\((B_n^*/K)\), hence precisely \(K\) and of rational degree two.
\end{itemize}

The remaining AbelianUnramifiedExtensionConditions fields are \textbf{eight\_dvd}, proved by (H23); \textbf{galE} and \textbf{abelE}, from cyclotomy; and \textbf{unram}, proved above. For the last field, the source ideal inequality is the ramification-index-one condition, exactly the finite unramified CM quotient already checked.

This constructs all mathematical data and proofs for that named record. It is not a Lean term or a build certificate. A larger normal cyclotomic coefficient field \(C\supset B_n^*\) may be used when finite character values require it: all embeddings of the normal field \(B_n^*\) still have their images in \(B_n^*\), so the same verification applies.

\subsubsection{5.3. Why the target is still outside the particular tame split-prime family}\label{why-the-target-is-still-outside-the-particular-tame-split-prime-family}

In the same file, IsMu0 at lines 314--320 requires a character to be unramified away from the two auxiliary rational primes \(g_i\), and tame everywhere. The preceding Primes.Ok record requires all the selected primes over \(g_i\) to have absolute norm \(g_i\) (lines 216--218), and QDat.Ok requires its primes to have prime absolute norm (line 347). These are substantive degree-one restrictions.

For \(p\ne3\) under the present hypothesis, \(p\equiv2\pmod3\) is inert in \(K\). Therefore every prime \(P\) over \(p\) in every number field containing \(K\) has \[
f(P/p)\ge2,\qquad NP=p^{f(P/p)},
\] so no such prime has prime absolute norm. AM-K003 says that \(\chi_{p,3}\) ramifies at the prime \((p)\) of \(K\). The extension \(B_n^*/K\) is unramified at \(p\), because \(p\) is neither \(7\) nor \(13\). The pulled-back character on \(B_n^*\) remains ramified: an unramified local extension leaves the inertia subgroup unchanged. Hence it cannot be placed solely in the displayed degree-one support.

At the prime above \(3\), the cases with conductor \(\lambda^2\) or \(\lambda^4\) are wild. Indeed tame inertia at residue characteristic \(3\) has no nontrivial quotient of order three, whereas these order-three characters are ramified. In (H23) the relative extension is unramified at \(3\), so the same obstruction persists after pullback. Thus all cases except \(p\equiv8\pmod9\) require a wild local character as well. In the exceptional congruence case the prime over \(3\) is unramified, but the inert-\(p\) degree restriction remains.

Consequently the general splitting transport AM-K008 is valid, while inserting its character into these particular IsMu0/QDat families is not automatic. A relative arithmetic adapter must allow the actual degree-two prime over \(p\), the wild conductor at \(3\) when present, and the resulting norm and lattice factors. This statement concerns these named families; it does not assert that no other source construction can represent the character.

\textbf{Closed source input and surviving gate.} The geometric field predicate AbelianUnramifiedExtensionConditions has an explicit supply with \(8\mid f\). The remaining work is the target's arithmetic character/column construction and its scale-uniform cost, with these precise nonsplit and wild local conditions retained.

\subsection{6. Resulting reduced programme}\label{resulting-reduced-programme}

For the cubic question, four concrete parts are now closed:

\begin{enumerate}
\def\labelenumi{\arabic{enumi}.}
\tightlist
\item
  the target character and all its finite conductor cases have a complete ray-class model;
\item
  the character is an allowed finite-order unitary splitting change, with exact source signs;
\item
  an explicit growing CM tower has an invariant coprime half-different, a positive unimodular trace lattice, the required norm-of-roots property, and a variant satisfying every mathematical field in AbelianUnramifiedExtensionConditions;
\item
  actual principally polarized products of the cubic CM elliptic twists carry the selected Artin Hom line.
\end{enumerate}

The remaining construction must place those objects in a relative arithmetic family, prove the exact equality between its arithmetic functional and AM-K005's normalized ideal-prime ramp, and compute its lift/period/lattice costs at all scales. The source's \(2\)-adic supply cannot simply be relabelled as the \(7\)-adic tower, and its displayed tame degree-one character family must be extended at the actual nonsplit and wild places. A successful estimate must give the window for every \(\theta>1/2\); all the identities here are compatible with the present gate Q3-KUMMER-WINDOW remaining open.

\subsection{References and source provenance}\label{references-and-source-provenance}

\begin{itemize}
\tightlist
\item
  Alpöge, L. (2026). \emph{Remark} and the stipulated Lean archive. \url{https://alpo.ge/remark.pdf}. Source files actually read: \nolinkurl{Formalism/LocalToGlobal/BesselHermitianTheta/TiedWeil.lean;} \nolinkurl{Formalism/LocalToGlobal/IndefiniteReal/LocalNewformFunction.lean;} \nolinkurl{Formalism/LocalToGlobal/Columns/AbelianHostFieldConditions.lean;} \nolinkurl{Development/AbelianVarietyTheory/Definitions/LatticeFrames/LevChiEqOne.lean;} \nolinkurl{Formalism/SharpRestriction/RamificationFrobenius/TotientGaloisClasses/ConjEMemHostFIff.lean;} its \nolinkurl{InertiaGaloisSymmetry/OfRawToHostRaw.lean;} \nolinkurl{Development/AutomorphicArithmetic/CharacterBoxes/Supply.lean} and ExistsPageN.lean.
\item
  Milne, J. S. (2020a). \emph{Class field theory} (v4.03). V, Theorem 1.7, pp.~150--151. \url{https://www.jmilne.org/math/CourseNotes/CFT.pdf}.
\item
  Milne, J. S. (2020b). \emph{Complex multiplication} (v0.10). §2, Examples 2.9 and 2.11; Propositions 3.12 and 3.17; Proposition 6.1. \url{https://www.jmilne.org/math/CourseNotes/CM.pdf}.
\item
  Milne, J. S. (2020c). \emph{Algebraic number theory} (v3.08). §6, Theorem 6.4, printed p.~99 (PDF p.~101); Remark 6.6(b)--(c), printed pp.~100--101 (PDF pp.~102--103). \url{https://www.jmilne.org/math/CourseNotes/ANT.pdf}.
\item
  Murty, M. R., \& Esmonde, J. (2005). \emph{Problems in algebraic number theory} (2nd ed.). Springer. §2.3, Exercise 2.3.2 and solution; §9.1. Author-hosted copy: \url{https://www.math.toronto.edu/~ila/2005_Book_ProblemsInAlgebraicNumberTheor.pdf}.
\item
  Companion contribution: AM-K001--AM-K005 in kummer\_extension.md. The earlier conductor, functional-equation, elliptic twist and analytic adapter proofs are incorporated by reference and not rederived here.
\end{itemize}

The archive was read from the shared pinned-source recovery, with its provenance files maintained by the parent, and from the child manifest for the earlier character files. No Lean compilation or newly formalized theorem is claimed by this addendum.

\clearpage

\section{AM-DO-017--020: literal moved-theta operators, a mixed-CM family, and exact split-place phases}\label{am-do-017020-literal-moved-theta-operators-a-mixed-cm-family-and-exact-split-place-phases}

Author: GPT-6 Astra Pro, ALP\_MACHINE \nolinkurl{/root/definite_oscillator}, 21 November 2026 UTC. The inherited reasoning effort is not separately exposed. This file gives full proofs of the new finite-operator assertions. AM-R-004, AM-DO-011--016, ALP-LIFT-07 and the source definitions are retained at their recorded status. Write \textbf{Σ\_tr} for AM-DO-015's normalized Mumford translate of F₁; it is distinct from the leading-coefficient-normalized hyperelliptic sigma function σ\_q used in the parallel calculation.

\subsection{1. The recovered literal source object}\label{the-recovered-literal-source-object}

The source definition sought in \nolinkurl{I.J.V.thMoved} has now been recovered and read. Its exact file is

\nolinkurl{Development/AbelianVarietyTheory/Definitions/AutomorphicMembership/IsDiscriminantTwelfthRootOfIdealLatticeIffDelta1.lean}.

The relevant locators are:

\begin{longtable}[]{@{}
  >{\raggedright\arraybackslash}p{(\columnwidth - 2\tabcolsep) * \real{0.5000}}
  >{\raggedright\arraybackslash}p{(\columnwidth - 2\tabcolsep) * \real{0.5000}}@{}}
\toprule\noalign{}
\begin{minipage}[b]{\linewidth}\raggedright
Source lines
\end{minipage} & \begin{minipage}[b]{\linewidth}\raggedright
Exact content
\end{minipage} \\
\midrule\noalign{}
\endhead
\bottomrule\noalign{}
\endlastfoot
18--24 & \nolinkurl{lamOf}, elliptic ideal lattice; the discriminant twelfth-root pin; an un-killed elliptic two-torsion point \\
36--51 & \nolinkurl{La}, \nolinkurl{Hf}, \nolinkurl{al}, \nolinkurl{U}, \nolinkurl{Iop} \\
110--118 & \nolinkurl{VacIn}: an elliptic \nolinkurl{PeriodPair}, its prescribed lattice, half-period, \nolinkurl{IsoKG}, its complex-linear extension, and \nolinkurl{KapIn} \\
130--138 & \nolinkurl{thMM}, \nolinkurl{thZero}, \nolinkurl{thHalf}, \nolinkurl{thMoved} \\
140--141 & \nolinkurl{eNat = chiFac · thMoved} \\
173--185 & literal conversion to the earlier \nolinkurl{Formalism.L2.Ker.Base} record \\
\end{longtable}

In mathematical notation its averaging operator is

\[
 I_{21}F=q_{21}^{-1/2}
       \sum_{x\in\Lambda_2/(\Lambda_1\cap\Lambda_2)}\alpha(x)U_xF,
 \qquad q_{21}=[\Lambda_2:\Lambda_1\cap\Lambda_2],
\tag{DO17.1}
\]

where the positive real square root is used and

\[
 U_xF(z)=e^{-\pi H(z,x)-\pi H(x,x)/2}F(z+x).
\tag{DO17.2}
\]

The source's \nolinkurl{La₁.addSubgroupOf La₂} is the intersection regarded as a subgroup of Λ₂; this matters for the index. Its sign is \texttt{al(x)=(-1)\^{}\{round((√d/2)H(x,x))\}}. In the admitted integral situations it is a sign semicharacter, so α⁻¹=α.

Writing Λ₀₀, Λ₁₀, Λ₁₁ for the standard, first-lattice-moved, and both-lattices-moved period lattices, the source literally uses

\[
 \mathrm{thMoved}_{L,L'}=I_{\Lambda_{10},\Lambda_{11}}
                       I_{\Lambda_{00},\Lambda_{10}}
                       \mathrm{thZero},
 \qquad\mathrm{thZero}=\mathrm{thMM}^{\boxtimes NN'}.
\tag{DO17.3}
\]

The exact seed is not left arbitrary: \nolinkurl{Formalism/LocalToGlobal/KernelLattices/LateralTheory.lean}, lines23--29, defines \nolinkurl{IsThM} by membership in the specified theta line and the fourth-tensor equation

\[
              \theta^{\boxtimes4}=\kappa\,\mathrm{kernel4}(L,\omega_0,\Lambda_M,g_C),
\tag{DO17.4}
\]

and takes \nolinkurl{thM} by choice from that predicate. The same file, lines128--178, repeats the complete base record and moved formulas. Lines230--238 separately name \nolinkurl{SignP} (comparison with the direct one-step average, up to sign) and \nolinkurl{OrderP} (equality of the two orders of moving the lattice factors). Thus phase and order compatibility are actual source obligations; norm preservation alone cannot replace them.

These additional modules were retrieved from the pinned mirror commit \nolinkurl{7cbf7a0c32fca1badf05931879a37c3f8a2b03b4}, with individual hashes in \nolinkurl{source_recovery_20261121T204914Z.json} and \nolinkurl{source_recovery_20261121T205055Z.json}. The intake does not supply original-author hashes for these paths. Retrieval and reading do not constitute a Lean build or a transitive proof audit.

\subsection{AM-DO-017. Exact normalized averaging between principal theta lines}\label{am-do-017.-exact-normalized-averaging-between-principal-theta-lines}

Let V be a complex vector space, H a positive Hermitian form linear in its first variable, and E=Im H. Let Λ₁,Λ₂ be commensurable full lattices selfdual for E:

\[
 \Lambda_i=\{v\in V:E(v,\Lambda_i)\subset\mathbf Z\}.
\tag{DO17.5}
\]

Let αᵢ:Λᵢ→U(1) satisfy

\[
 \alpha_i(x+y)=\alpha_i(x)\alpha_i(y)e^{\pi iE(x,y)},
\tag{DO17.6}
\]

and suppose their restrictions to Λ₀=Λ₁∩Λ₂ agree. Write T(Λᵢ,αᵢ) for the holomorphic functions satisfying U\_λF=αᵢ(λ)F. On this space use the probability-Haar norm

\[
 \|F\|_i^2=\frac1{\operatorname{vol}(D_i)}
     \int_{D_i}|F(z)|^2e^{-\pi H(z,z)}\,dz,
\tag{DO17.7}
\]

where Dᵢ is any fundamental domain for Λᵢ and dz is any fixed Lebesgue measure. This is independent of Dᵢ and of rescaling dz.

Then q={[}Λ₂:Λ₀{]}={[}Λ₁:Λ₀{]} is a positive finite integer and

\[
 \boxed{I_{21}F=q^{-1/2}\sum_{x\in\Lambda_2/\Lambda_0}
             \alpha_2(x)^{-1}U_xF}
\tag{DO17.8}
\]

is a well-defined nonzero isometry T(Λ₁,α₁)→T(Λ₂,α₂). In the source's sign-valued case this is exactly (DO17.1), including its square-root index. Both spaces are one-dimensional. No unspecified positive scalar occurs in (DO17.8).

\subsubsection{Proof}\label{proof-1}

Direct multiplication of the factors in (DO17.2) gives

\[
 U_xU_y=e^{-\pi iE(x,y)}U_{x+y},\qquad
 U_xU_y=e^{-2\pi iE(x,y)}U_yU_x.
\tag{DO17.9}
\]

Also

\[
 |U_xF(z)|^2e^{-\pi H(z,z)}
       =|F(z+x)|^2e^{-\pi H(z+x,z+x)}.
\tag{DO17.10}
\]

Thus translations are isometries whenever both sides are viewed on a common lattice on which their norms are periodic.

Because Λᵢ is selfdual, the alternating integral matrix of E on any Z basis of Λᵢ is unimodular. Its Pfaffian has absolute value1. Consequently both lattices have the same symplectic covolume, hence the same covolume for dz. Since Λ₀ has finite index in each, division of covolumes gives {[}Λ₁:Λ₀{]}={[}Λ₂:Λ₀{]}=q. A fundamental domain for Λ₀ is, up to null sets, the disjoint union of q translates of a fundamental domain for Λᵢ. For a Λᵢ-periodic norm density, dividing its integral by the domain volume therefore gives precisely (DO17.7). Thus the normalized norm on the common lattice equals the norm on either original lattice; there is no extra q in changing domains.

For x∈Λ₂, the function U\_xF lies in T(Λ₀,α₀), since E(x,Λ₀) is integral. The operators V\_x=α₂(x)⁻¹U\_x form an ordinary representation of the additive group Λ₂ by (DO17.6) and (DO17.9). If x∈Λ₀, then V\_xF=F. Consequently the sum (DO17.8) is independent of coset representatives and is fixed by all V\_y, y∈Λ₂. It is therefore in T(Λ₂,α₂).

The finite pairing

\[
 (\Lambda_1/\Lambda_0)\times(\Lambda_2/\Lambda_0)
       \longrightarrow U(1),\qquad
 (\lambda,x)\longmapsto e^{2\pi iE(\lambda,x)}
\tag{DO17.11}
\]

is well-defined and nondegenerate. Indeed, if x∈Λ₂ pairs trivially with every λ∈Λ₁, then selfduality of Λ₁ implies x∈Λ₁, hence x∈Λ₀. The same argument applies to the other variable. The two groups have order q, so this is a perfect finite pairing.

For distinct x,y modulo Λ₀, choose λ∈Λ₁ for which e\^{}\{2πiE(λ,x−y)\}≠1. The operator U\_λ is unitary on T(Λ₀,α₀), by (DO17.10). By (DO17.9), U\_xF and U\_yF are its eigenvectors with eigenvalues α₁(λ)e\^{}\{−2πiE(λ,x)\} and α₁(λ)e\^{}\{−2πiE(λ,y)\}. Their eigenvalues are distinct and have absolute value1. Taking their inner product and using unitarity shows that they are orthogonal. Each has norm \textbar\textbar F\textbar\textbar₁. Therefore

\[
 \|I_{21}F\|_2^2
     =q^{-1}\sum_{x\in\Lambda_2/\Lambda_0}\|U_xF\|_0^2
     =\|F\|_1^2.
\tag{DO17.12}
\]

This proves the isometry and nonvanishing.

For completeness, dimension one is the elementary principal-theta dimension statement. Choose a symplectic Z basis of Λᵢ, normalize its first half to the real standard basis, and write the second half as the columns of τ with Im τ positive. After a linear exponential gauge, the real-period relation gives a Fourier expansion indexed by Z\^{}g. The τ-period relation successively determines every Fourier coefficient from one coefficient, with the usual Gaussian factor exp(πinᵗτn) and the semicharacter shifts. Positivity gives normal convergence of that theta series. Hence the dimension is exactly one. Equivalently this is h⁰=1 for a principal polarization. ∎

The theta line on V/Λ₀ used in this proof has degree q: its symplectic covolume is q times that of a principal lattice. The q orthogonal vectors above are consistent with its q-dimensional space of sections. The argument did not confuse this degree with q², the order of its full theta kernel.

\subsubsection{Exact transport and exact residual phase}\label{exact-transport-and-exact-residual-phase}

If C:V→V′ is complex-linear with H′(Cx,Cy)=H(x,y), CΛᵢ=Λᵢ′ and αᵢ′(Cx)=αᵢ(x), then pullback by C commutes \textbf{exactly} with (DO17.8). Indeed it commutes with each U\_x, induces a bijection of the finite quotient sets, and preserves q. Thus any specified base equality with scalar c extends to the entire two-step family with that same scalar c.

For three such lattices, choose a nonzero F∈T(Λ₁,α₁). Since all three theta lines have dimension one, there is a unique phase

\[
 \mu_{321}=\frac{\langle I_{31}F,I_{32}I_{21}F\rangle_3}{\|F\|_1^2},
 \qquad |\mu_{321}|=1,
 \qquad I_{32}I_{21}=\mu_{321}I_{31}.
\tag{DO17.13}
\]

Our convention makes the inner product linear in its second argument. This phase is independent of F. It is explicitly a finite Gauss sum: expand the three averages, use (DO17.9), and use

\[
 \frac{\langle U_zF,U_tF\rangle}{\|F\|^2}
 =\begin{cases}
 e^{\pi iE(z,t)}\alpha_1(t-z),&t-z\in\Lambda_1,\\
 0,&t-z\notin\Lambda_1.
 \end{cases}
\tag{DO17.14}
\]

The zero case follows from the same perfect-pairing argument (using a common finite-index lattice for the finitely many vectors); the nonzero case follows from U\_\{−z\}U\_t=e\^{}\{πiE(z,t)\}U\_\{t−z\}. Equations(DO17.8), (DO17.9), and (DO17.14) are an exact finite prescription for μ, with all normalizations included. Neither its value1 nor its value±1 is automatic from nonvanishing. The special \nolinkurl{SignP} and \nolinkurl{OrderP} assertions in the original constructor must therefore be checked in a new constructor.

\subsection{AM-DO-018. A trace-one mixed-CM moved family with explicit finite inputs}\label{am-do-018.-a-trace-one-mixed-cm-moved-family-with-explicit-finite-inputs}

Retain the actual principally polarized J of AM-R-004. Let V=C³ in algebraic u coordinates, Λ=DΦ(O\_E), and H(u,v)=Σu\_i conjugate(v\_i)/G\_i. In the first-argument-linear convention, E\_H=Im H is the negative of the displayed Riemann intersection form in AM-R-004; changing this sign does not change selfduality. Put

\[
 w_0=-(\zeta+\zeta^2+\zeta^4),\qquad
 w_0+\bar w_0=1,
\]

acting through the actual CM action. On every w₀-stable E\_H-integral lattice define

\[
             \alpha_{\rm tr}(x)=e^{\pi iE_H(x,w_0x)}\in\{1,-1\}.
\tag{DO18.1}
\]

This is a semicharacter, and its formula agrees on intersections of lattices. It is the Appell--Humbert semicharacter of the rigidified trace-one Poincaré line L\_\{w₀\}.

\textbf{Proof.} The Rosati adjoint of w₀ is bar(w₀). Hence

\[
 E_H(x,w_0y)+E_H(y,w_0x)
       =E_H(x,(w_0-\bar w_0)y)
       =2E_H(x,w_0y)-E_H(x,y).
\]

For integral stable lattices, its exponential times πi equals exp(πiE\_H(x,y)). Expanding E\_H(x+y,w₀(x+y)) proves (DO17.6). For the line identification, the Poincaré line pulled back by the principal polarization is m\emph{L⊗p₁}L⁻¹⊗p₂*L⁻¹. Its Hermitian form on V×V is H(z,w′)+H(w,z′) and its semicharacter on (λ,μ) is exp(πiE\_H(λ,μ)), by the semicharacter identity for any representative L. Pullback by (1,w₀) has Hermitian form H, since w₀+bar(w₀)=1, and semicharacter (DO18.1). This proves the assertion with the canonical rigidifications. ∎

As an explicit sign check in the actual period basis, E\_R004(ζʲ,w₀ζʲ)=Tr(ξ bar(w₀))=−1 for each j=0,\ldots,5, with ξ=(ζ−ζ⁻¹)/7. Thus α\_tr takes value−1 on each adjacent-cycle generator γⱼ. AM-DO-012/ALP-LIFT-07's characteristic K∞=(101\textbar111) gives exactly these six signs. This checks the trace-one line against the analytic Σ\_tr convention rather than merely matching Chern classes.

\subsubsection{Matrix-lattice family}\label{matrix-lattice-family}

Fix positive integers a,b. For full O\_E-lattices N⊂E\^{}a and N′⊂E\^{}b that are selfdual for their standard Hermitian forms, let

\[
 M(N,N')=\operatorname{Hom}_{O_E}(N,N')\subset M_{b\times a}(E),
 \qquad\Lambda(N,N')=D\Phi(M(N,N')).
\tag{DO18.2}
\]

Here DΦ is applied entry by entry, and the Hermitian form is the sum of H on all ab entries. These lattices are full, commensurable, selfdual for the product E\_H, and stable under w₀. To check selfduality, localize at a rational prime and choose arbitrary module bases g,h of N,N′. Their Hermitian Gram matrices G\_N=g* g and G\_N′=h* h are integral and unimodular because the lattices are Hermitian selfdual. Writing X=hZg⁻¹ identifies the matrix pairing with the trace pairing whose Gram operator is Z↦G\_N′ Z G\_N⁻¹; this is an integral module automorphism. The scalar trace pairing is perfect on O\_E because ξO\_E is the codifferent, as proved in AM-R-004. Hence the matrix trace pairing is perfect on every local Hom lattice. Intersecting the local duals proves the global assertion. Equivalently the argument is projective-module duality Hom(N,N′)≃N*⊗N′. This includes every lattice pair obtained from the standard lattices by finite adelic unitary moves. No local classification of ramified or dyadic unimodular Hermitian lattices, and no global rational unitary basis, is assumed.

Let

\[
 F_{00}(u)=\prod_{i=1}^{b}\prod_{j=1}^{a}\Sigma_{\rm tr}(u_{ij}),
\]

and put

\[
 F_{N,N'}=I_{\Lambda(N,O_E^b),\Lambda(N,N')}
          I_{\Lambda(O_E^a,O_E^b),\Lambda(N,O_E^b)}F_{00},
\tag{DO18.3}
\]

using (DO17.8) and the common sign(DO18.1). This is an actual, nonzero holomorphic section of the corresponding principal trace-one line on V\^{}\{ab\}/Λ(N,N′). Every averaging coefficient and every finite quotient is specified. It has exactly the same probability-Haar L² norm as F₀₀.

Write

\[
 q_1=[\Lambda(N,O_E^b):\Lambda(N,O_E^b)\cap\Lambda(O_E^a,O_E^b)],\quad
 q_2=[\Lambda(N,N'):\Lambda(N,N')\cap\Lambda(N,O_E^b)].
\tag{DO18.4}
\]

The two q's are concrete integral lattice indices, computable by Smith normal form in the integral period basis. They are not arbitrary group volumes.

\subsubsection{Explicit rational Heisenberg intertwining}\label{explicit-rational-heisenberg-intertwining}

Use normalized complex coordinates z and the convention ψ∞(r)=exp(2πir), ψ\_f(r)=exp(−2πir) on Q. For t=A+τB with A,B∈Q³ define on finite Schwartz functions

\[
 (W_t\phi)(q_f)=\psi_f(-A^tq_f+A^tB/2)\,\phi(q_f-B).
\tag{DO18.5}
\]

Let Rφ(z)=e\^{}\{Q(z)\}Σ\_\{q∈Q³\}φ(q)e\^{}\{πiqᵗτq+2πiqᵗz\}. Then

\[
                     R(W_t\phi)=U_t(R\phi).
\tag{DO18.6}
\]

\textbf{Proof.} The multiplier from Q(z+t)−Q(z)−πH(z,t)−πH(t,t)/2 equals exp(2πiBᵗz+πiBᵗτB+πiAᵗB). In the theta sum for Rφ(z+t), replace q by q′−B. The remaining coefficient is φ(q′−B)exp(2πiAᵗq′−πiAᵗB), exactly(DO18.5) on rational points. Absolute convergence on compact sets permits this reindexing and all derivatives. The same calculation entry by entry proves the matrix version. ∎

For reference this map is injective on finite Schwartz functions: compact support gives a common denominator M with rational support in M⁻¹Z³; restricting z to the real cube of side M, Fourier uniqueness recovers every φ(q) times its nonzero Gaussian coefficient. The diagonal Q is dense in A\_f, so a locally constant function vanishing on Q is zero. The same holds in all matrix dimensions.

Let φ₀ be the tensor product of AM-DO-016's specified input

\[
 -i\,1_{\delta/2+\widehat{\mathbf Z}^3}(q_f)\,
             \psi_f(-q_f^t\varepsilon/2),
 \qquad\delta=(1,1,1),\ \varepsilon=(1,0,1).
\]

Replace each U\_t in (DO18.3) by W\_t, with the same α\_tr(t) and the same positive q\_i⁻¹/², to define the finite Schwartz function φ\_\{N,N′\}. Every t lies in the rational span of the standard period lattice, so (DO18.5) applies. The output is compactly supported and locally constant, with algebraic coefficients. AM-DO-016 and (DO18.6) give the \textbf{whole moved-family equality}

\[
 \boxed{R\phi_{N,N'}=\theta[0;0](0;\tau)^{ab}\,F_{N,N'}.}
\tag{DO18.7}
\]

The period scalar is the same for every N,N′ at these fixed ranks. This is more than a nonzero period or a single-fiber identity. It is an actual family of finite inputs and all-degree holomorphic theta equalities. It still does not identify φ\_\{N,N′\} with the prescribed split Weil translate of φ₀; that further assertion is addressed below.

\subsubsection{Normalization and explicit elementary costs}\label{normalization-and-explicit-elementary-costs}

Carry the exact norm from AM-R-004(R4.12), with κ=1, and the isometry of normalized Mumford translation. Then

\[
 \|F_{N,N'}\|^2
     =\det(2Y)^{-ab/2}|\theta[0;0](0;\tau)|^{-2ab}.
\tag{DO18.8}
\]

Thus the phase-preserving scalar that changes (DO18.7)'s reduced seed to norm1 is

\[
 c_{\rm unit}=\theta[0;0](0;\tau)^{ab}\det(2Y)^{ab/4},
 \qquad \|c_{\rm unit}F_{N,N'}\|=1.
\tag{DO18.9}
\]

This is an explicit local analytic normalization, not the terminal automorphic mass pairing in G1.

AM-R-006 now computes the algebraic scalar between the two actual seeds: Σ\_tr=cσ\_q, where the phase is fixed by (R6.3), c⁸=R\_S²/4⁸, v₂(c)=−2, v₇(c)=1/2, and the other finite valuations are zero. Consequently, if the same averages are applied to σ\_q\^{}\{⊠ab\}, then (DO18.7)'s common period is exactly (θ₀c)\^{}\{ab\}. Its squared norm is (DO18.8) divided by \textbar c\textbar\^{}\{2ab\}, and the phase-preserving norm-one scalar is (θ₀c)\^{}\{ab\}det(2Y)\^{}\{ab/4\}. These equalities follow from linearity of each specified average. They do not select an arbitrary algebraic root for c or absorb it silently into an unspecified period.

If normalized valuations satisfy v\_p(p)=1, the algebraic values of the constructed Schwartz input satisfy the elementary bound

\[
 v_p(\phi_{N,N'}(q_f))\ge-\tfrac12v_p(q_1q_2)
\tag{DO18.10}
\]

where a zero value has valuation +∞. Indeed the original nonzero coefficients, α\_tr signs and finite additive-character values are roots of unity; the only displayed denominators are the two inverse square roots of positive integers. The ultrametric inequality proves the claim, and cancellation may improve it. This is a statement about \textbf{finite Schwartz coefficients}, not an assertion that the corresponding reduced CM Taylor series is integral at a nonordinary prime.

At the fixed complex embedding, AM-DO-015 and triangle inequality give the completely explicit growth bound

\[
 |F_{N,N'}(u)|<49^{ab}\sqrt{q_1q_2}\,
      \exp\!\left(\frac\pi2\sum_{i,j,k}|u_{ij,k}|^2/G_k\right).
\tag{DO18.11}
\]

Each normalized U\_t preserves the Gaussian envelope, while a q-term average costs at most √q. The corresponding polydisk Cauchy bound follows exactly as in AM-R-004. This does not claim that √(q₁q₂) is optimal. L² norm preservation does not by itself eliminate this displayed pointwise cost for varying lattices.

\subsection{AM-DO-019. Exact boundary between this mixed family and the prescribed constructor}\label{am-do-019.-exact-boundary-between-this-mixed-family-and-the-prescribed-constructor}

There are two distinct questions.

\textbf{Unchanged elliptic \nolinkurl{VacIn}.} The recovered record fixes an elliptic period lattice and a complexification of an isomorphism over the quadratic CM subfield; its base fourfold kernel is built from that elliptic data. AM-R-003 proves that the primitive Φ=\{1,2,3\} Jacobian J has no nonzero homomorphism to or from an elliptic curve and no isogeny with an elliptic power. Therefore a full, line-preserving CM/Hodge transport from that unchanged elliptic constructor to J (or the corresponding powers) does not exist. This conclusion uses the actual category demanded by the record, not the false inference that absence of a canonical map excludes all maps. Arbitrary vector-space maps, selected coefficient maps, and a new mixed-CM constructor are not excluded.

\textbf{A new mixed-CM constructor.} Equations(DO18.1)--(DO18.11) give its trace-one theta seed, every moved lattice, every finite input, the exact average/covolume factors, and one common analytic period. The remaining prescribed-Weil equality is

\[
 \omega_{\psi,\chi}(g,h)\phi_0
       \stackrel{?}=\chi(x)^{-1}\chi(y)\,\phi_{N,N'},
 \quad x/\bar x=\det g,\quad y/\bar y=\det h.
\tag{DO19.1}
\]

The roots x,y are determinant roots, as in the actual source \nolinkurl{IsDetRoot}; substituting det(g),det(h) is generally wrong. The source also fixes ψ∞ with the opposite sign from (DO18.5); its polarization/character transport must be included before testing(DO19.1). In a common realized Heisenberg representation the two sides lie on the same one-dimensional lattice vacuum line, so their quotient is a phase once both are norm1. Its value must be computed from the chosen splitting on generators, or from the finite Gauss sum(DO17.13) together with the specified determinant-root character. Merely saying that both vectors are nonzero gives no equality.

The exact recovered source makes this distinction visible. \nolinkurl{MovedWeilVacuumEqVacuumOfMovedLattices.lean}, theorem \nolinkurl{vac_moved}, proves preservation of the moved lattice vacuum, theta membership and nonvanishing. It does \textbf{not} conclude the common c₀/determinant-root formula. \nolinkurl{LevChiEqOne.lean}, lines87--101, places that formula in \nolinkurl{VacT} and \nolinkurl{PinX.mv}. \nolinkurl{ThetaKernelSplitPlacesEqJetClassicalTheta.lean}, theorem \nolinkurl{thetaKernel_eq_jet_thC}, takes the equation \nolinkurl{xiN(om(g,h)vac)=κE} as an explicit hypothesis and then transports it to every polynomial jet. Thus the missing phase is exactly where the source says it is.

A failure to prove(DO19.1) is not a proof that it is false. Conversely, (DO18.7) does not prove it by construction, because the φ\_\{N,N′\} were defined through the two fixed finite averages, while the source fixes a group representation and its splitting independently. Nor do these analytic equalities prove the full integral row-lattice identities, the nonordinary formal-gauge comparison, the \nolinkurl{MainX} cohomological/topological bounds, or the terminal-mass equation from THETA3 G1. Those remain their stated separate gates.

\subsection{AM-DO-020. The phase is exactly the splitting character for relative split-place moves}\label{am-do-020.-the-phase-is-exactly-the-splitting-character-for-relative-split-place-moves}

The previously open phase in(DO19.1) can be resolved in a substantial precise scope: moves supported on finite places of F that split in E. This includes the primes above11 and29 for E=Q(ζ₇), and also the relative split dyadic places. This assertion is conditional on a Weil representation with the displayed covariance, unitarity and split-normalization axioms; it proves the required phase from those axioms, rather than replacing the prescribed representation by a new family of unrelated operators. It does not extend the old quadratic \nolinkurl{StdPair.xi} pin to the primitive type.

\subsubsection{Local statement}\label{local-statement}

Let k be a nonarchimedean field of characteristic zero, O its integers, and ψ a nontrivial additive character. Put E\_v=k×k with conjugation interchanging factors. Choose the split Schrödinger polarization of the rational symplectic Hom space. After identifying the momentum space with the ψ-dual of the position space X, write its form and Heisenberg action as

\[
 B((x,y),(x',y'))=\langle y,x'\rangle-\langle y',x\rangle,
 \qquad
 [\rho(x,y)\phi](t)=\psi(-\langle y,t\rangle+\langle y,x\rangle/2)
                           \phi(t-x).
\tag{DO20.1}
\]

Let L⊂X be a full O-lattice and L* its annihilator under ψ. The split O\_E-lattice is Λ\_L=L⊕L*. An integral trace-one endomorphism is w₀=(w,1−w), with w,1−w∈O. Its local semicharacter is

\[
 \alpha_L(x,y)=\psi\bigl((w-\tfrac12)\langle y,x\rangle\bigr),
       \qquad (x,y)\in L\oplus L^*.
\tag{DO20.2}
\]

Here the pairing is taken in the additive-character conductor convention, so its integral values and their O-multiples are killed by ψ. Formula(DO20.2) is the local factor of α\_tr=exp(πiE\_H(x,w₀x)): globally α\_tr=ψ\_f(B(x,w₀x)/2), using E\_H=−B and the product formula for the additive character. It follows immediately that

\[
 \alpha_L(x,y)^{-1}\rho(x,y)\phi(t)
                   =\psi(-\langle y,t\rangle)\phi(t-x)
       \quad\text{for }(x,y)\in\Lambda_L.
\tag{DO20.3}
\]

Indeed the extra scalar is ψ((1−w)〈y,x〉)=1. This cancellation is valid at residue characteristic two as well. Dropping the half-pairing term without the trace-one semicharacter would not justify that dyadic assertion.

The common eigenspace

\[
 \rho(\lambda)\phi=\alpha_L(\lambda)\phi
                     \quad(\lambda\in\Lambda_L)
\tag{DO20.4}
\]

is precisely the line spanned by1\_L. For y∈L\emph{, the modulation relation forces the support into (L})*=L; for x∈L, the translation relation makes the function constant there. Conversely1\_L satisfies the relations by(DO20.3). With any fixed Haar measure dt on X, its positive norm-one generator is

\[
                         v_L=\operatorname{vol}(L)^{-1/2}1_L.
\tag{DO20.5}
\]

For two position lattices L₀,L, the averaging operator(DO17.8), realized by ρ, satisfies the exact equality

\[
                             I_{L,L_0}v_{L_0}=v_L.
\tag{DO20.6}
\]

Our index notation in(DO20.6) means that the output lattice is L; equivalently I\_\{21\}v\_\{L₁\}=v\_\{L₂\}.

\subsubsection{Proof of the exact index and value in(DO20.6)}\label{proof-of-the-exact-index-and-value-indo20.6}

Let M=L∩L₀, N=L+L₀, u={[}L:M{]}, and v={[}L₀:M{]}. Duality gives

\[
 L^*\cap L_0^*=N^*,\quad
 [L^*:N^*]=[N:L]=[L_0:M]=v.
\]

Therefore

\[
 q=[\Lambda_L:\Lambda_L\cap\Lambda_{L_0}]=uv,
 \qquad \operatorname{vol}(L)/\operatorname{vol}(L_0)=u/v.
\tag{DO20.7}
\]

Use representatives x∈L/M and y∈L\emph{/N} in the finite average. By(DO20.3), its value on1\_\{L₀\} at t is

\[
 (uv)^{-1/2}\sum_{x\in L/M}\sum_{y\in L^*/N^*}
               \psi(-\langle y,t\rangle)1_{L_0}(t-x).
\tag{DO20.8}
\]

The possible support is N. For t∈N, finite character orthogonality makes the inner sum equal v when t∈L and zero otherwise. For t∈L there is exactly one class x∈L/M with t−x∈L₀. Hence(DO20.8) equals √(v/u)1\_L(t). Multiplying by vol(L₀)⁻¹/² gives(DO20.6). In particular its value at position0 is strictly positive. This proof also gives

\[
 I_{L_3,L_2}I_{L_2,L_1}=I_{L_3,L_1}
\tag{DO20.9}
\]

on the corresponding lattice-vacuum line. Thus the triangle phase is1 and the two orders of moving the source and target Hom-lattice factors agree in this split scope.

\subsubsection{Comparison with the independently prescribed split Weil action}\label{comparison-with-the-independently-prescribed-split-weil-action}

At a split place the unitary pair is GL\_a(k)×GL\_b(k). On its position Hom space, (g,h) acts by A:Z↦hZg⁻¹. Its determinant is

\[
                    \det_k A=(\det h)^a(\det g)^{-b}.
\tag{DO20.10}
\]

The explicit unitary Schrödinger change-of-variables operator

\[
 T_A\phi(t)=|\det A|_k^{-1/2}\phi(A^{-1}t)
\tag{DO20.11}
\]

has the requisite covariance, and sends v\_\{L₀\} to v\_\{AL₀\}. Its positive value at0 is vol(AL₀)⁻¹/².

Let ω(g,h) be the prescribed unitary Weil action and let its split-normalization character be

\[
             \chi_T(x)^{-1}\chi_S(y),\qquad
               x/\bar x=\det g,\quad y/\bar y=\det h.
\tag{DO20.12}
\]

Covariance transports(DO20.4), so ω(g,h)v\_\{L₀\}=λv\_\{AL₀\} for some λ; the elementary uniqueness of that line above suffices, without a further abstract irreducibility assertion. Unitarity gives \textbar λ\textbar=1. The functional ℓ(φ)=φ(0) is invariant under all momentum Heisenberg operators. It is therefore exactly a functional to which the source \nolinkurl{om_split} condition applies. That condition gives

\[
 \ell(\omega(g,h)v_{L_0})
       =r\,\chi_T(x)^{-1}\chi_S(y)\,\ell(v_{L_0}),\qquad r>0.
\tag{DO20.13}
\]

Both ℓ(v\_\{L₀\}) and ℓ(v\_\{AL₀\}) are positive. Comparing(DO20.13) with λv\_\{AL₀\}, and using \textbar λ\textbar=1, proves

\[
 \boxed{\omega(g,h)v_{L_0}
           =\chi_T(x)^{-1}\chi_S(y)\,v_{AL_0}.}
\tag{DO20.14}
\]

The positive scalar r is forced to be

\[
                    r=\bigl(\operatorname{vol}(L_0)/
                                  \operatorname{vol}(AL_0)\bigr)^{1/2}
                      =|\det A|_k^{-1/2}.
\tag{DO20.15}
\]

This computes both the phase and the positive normalizing factor. It does not silently discard the r in the source's split axiom.

\subsubsection{Global consequence for the mixed family}\label{global-consequence-for-the-mixed-family}

Apply this place by place to finite adelic g,h supported at relative split places. The finite quotients defining the global averages decompose into their primary components. The character and Heisenberg factors decompose in the same way; on the common lattice their weighted translation operators are an ordinary additive-group representation, so the local factors commute. The positive square-root normalization is the product of the positive local square roots. Consequently the two-step input φ\_\{N,N′\} of(DO18.7) is exactly the product of the positive local lattice-vacuum vectors at the moved places, with the original factors elsewhere. Equations(DO20.6) and(DO20.14) prove(DO19.1) in this scope, with its determinant-root character and the common scalar θ₀\^{}\{ab\} retained by(DO18.7).

For the even ranks of \nolinkurl{StdPair}, the restriction of χ\_S and χ\_T to A\_F\^{}× is η\_\{E/F\}\^{}\{even\}=1. Therefore changing determinant roots by fixed-field ideles does not change(DO20.12). For arbitrary ranks the statement uses the specified supported determinant roots; one must retain the splitting's stated root dependence.

There is an explicit way to align additive-character signs with the recovered source. Use its ψ\_F=(ψ\_Q⁻¹)∘Tr\_\{F/Q\}, whose real sign is−1, and set ξ₀=−ξ in \nolinkurl{symplO}, where ξ=(ζ−ζ⁻¹)/7. Then ψ\_F(Tr\_\{E/F\}(ξ₀h))=ψ\_Q(Tr\_\{E/Q\}(ξh)), the actual Heisenberg central character used above. This is a simultaneous form/character transport, so it leaves the realized representation unchanged. It changes the old constructor's fixed quadratic ξ₀ condition, which is already part of the mixed-CM constructor gate. It does not assert that the unchanged source \nolinkurl{StdPair.xi} is satisfied.

\subsubsection{A concrete rank-two to rank-four moved lattice}\label{a-concrete-rank-two-to-rank-four-moved-lattice}

Take a=2, b=4, π=2−ζ, and r=π/bar(π). The exact norm is

\[
 N_{E/\mathbf Q}(\pi)=\Phi_7(2)=127.
\]

The primes (π) and (bar π) are distinct: ζ is2 modulo(π), while its inverse is64, and2≠64 in F₁₂₇. Put

\[
              N=rO_E\oplus O_E,\qquad N'=O_E^4.
\tag{DO20.16}
\]

The lattice N is Hermitian selfdual since r bar(r)=1. Its Hom lattice has four entries r⁻¹O\_E and four entries O\_E. For one changed entry,

\[
 O_E\cap r^{-1}O_E=\bar\pi O_E,
 \qquad [O_E:\bar\pi O_E]=[r^{-1}O_E:\bar\pi O_E]=127.
\]

Hence the exact global indices in(DO18.4) are

\[
                         q_1=127^4,\qquad q_2=1.
\tag{DO20.17}
\]

One may choose the finite adelic unitary element g with components diag(r,1) at127 and identity elsewhere, h=1. Its support is relative split because127≡1 mod7. A determinant root is the finite idele whose127-components are π and whose other components are1. For trivial splitting characters at these even ranks, the character in(DO20.14) is1, so the actual split Weil translate equals the explicitly averaged input, with no residual phase. The finite average has127⁴ terms and exact coefficient127⁻², before signs and Heisenberg characters. Its elementary coefficient denominator bound is v₁₂₇≥−2. Its fixed-embedding Gaussian bound is49⁸·127². Its exact analytic period in(DO18.7) isθ₀⁸, and its norm-one comparison scalar isθ₀⁸ det(2Y)². These are concrete values for this moved kernel family, not a terminal mass or a proof that a selected rank-four row lies in a global theta-lift image.

\subsubsection{What fails to transport at the ramified prime7}\label{what-fails-to-transport-at-the-ramified-prime7}

The recovered source phase proof is more specific than a generic metaplectic uniqueness argument. \nolinkurl{.../KernelLattices/ColumnData/ColumnDataEquations/MuFinal.lean}, theorems \nolinkurl{muOneL_of} and \nolinkurl{muOneR_of}, uses \nolinkurl{J.negOneNorm D (4*e)} to find an integral l with

\[
                         l\bar l\equiv-1\pmod{4e}.
\tag{DO20.18}
\]

Together with an even-rank row/column symmetry this makes the source's triangle factor1. The source \nolinkurl{SplitPairOfGen.lean}, theorem \nolinkurl{splitPair_of_gen}, explicitly obtains \nolinkurl{IsCMUnram} from \nolinkurl{HostDer.unramEF}.

For E=Q(ζ₇), (DO20.18) has no solution when7\textbar e. Reduction at the unique prime above7 makes conjugation the identity on F₇, so the congruence would say l²=−1 in F₇. The squares are0,1,2,4 and do not contain6. This disproves that auxiliary congruence in the mixed example and prevents directly reusing that particular proof. It does not disprove the desired phase identity, which might be obtained with a different normalization or a more precise ramified computation. The present calculation closes split-place phases and leaves a sharply located nonsplit/ramified phase check, together with the independent global integral and terminal-mass gates.

\clearpage

\section{AM-DO-021: an exact ramified compact character and a rank-two to rank-four counterexample}\label{am-do-021-an-exact-ramified-compact-character-and-a-rank-two-to-rank-four-counterexample}

Author: GPT-6 Astra Pro, ALP\_MACHINE \nolinkurl{/root/definite_oscillator}; inherited reasoning effort not separately exposed; 21 November 2026 UTC.

This result uses the standard unitary Weil splitting attached to characters χ\_S,χ\_T in Kudla's normalization, with Pan's explicit comparison to the lattice model. It computes the missing scalar at the ramified prime 7. It does not infer that every representation satisfying only some of the source's listed axioms has already been identified with this standard splitting. The split-place proof AM-DO-020 did not require this additional identification.

\subsection{Statement}\label{statement-7}

Character labels use the source convention \nolinkurl{LocalNewformFunction.splitFacO}, whose source factor is \(\chi_T(x)^{-1}\) and target factor is \(\chi_S(y)\) on determinant roots. Thus conventional Harris--Kudla--Sweet tensor parameters are \((\chi_S,\chi_T^{-1})\); corrected AM-DO-023 records the same conversion at infinity. In particular every trivial-character assertion below is independent of that labeling choice.

Let E=Q(ζ₇), F=E⁺, and let v be the unique place of F above 7. Put K=F\_v, L=E\_v, O=O\_K, O\_L=O\_\{E\_v\}, and s=ζ₇−ζ₇⁻¹. Conjugation on L is written with a bar. The residue fields of K and L are both F₇, and s is an anti-invariant uniformizer of L. Let ξ=s/7.

For positive even integers a,b, take the standard Hermitian lattices O\_L\^{}a,O\_L\^{}b and their Hom lattice Λ=M\_\{b×a\}(O\_L). Use the actual rational symplectic form of AM-DO-018 and the corresponding additive character. Let φ₀ span the trace-one Heisenberg vacuum line of Λ. Let

\begin{verbatim}
K_a=U(I_a)(O),       K_b=U(I_b)(O).
\end{verbatim}

For t∈L¹=\{t:t bar(t)=1\}, define

\[
 \epsilon(t)=\begin{cases}1,&t\equiv1\pmod{s},\\-1,&t\equiv-1\pmod{s}.
 \end{cases}
\tag{DO21.1}
\]

Every norm-one t is a unit and has residue ±1, so this is a continuous quadratic character. Let χ\_S,χ\_T be the unitary splitting characters; at these even ranks their restrictions to K× are trivial. For determinant roots x/bar(x)=det g, y/bar(y)=det h, the exact scalar is

\[
 \boxed{\omega_{\psi,\chi}(g,h)\phi_0
   =\chi_T(x)^{-1}\chi_S(y)
      \epsilon(\det g)^{b/2}\epsilon(\det h)^{a/2}\phi_0,
       \quad g\in K_a,\ h\in K_b.}
\tag{DO21.2}
\]

In particular, take a=2,b=4, g=1 and h=diag(−1,1,1,1). Both lattices are unchanged. Both prescribed finite averages in AM-DO-018 are therefore the identity. Nevertheless, with s supported at the place v and y=−2s as a determinant root at v,

\[
               \omega(1,h)\phi_0=-\chi_S(y)\phi_0.
\tag{DO21.3}
\]

Thus the unchanged two-average formula (DO19.1), whose right side is +χ\_S(y)φ₀, is false for this actual standard splitting at these ranks. This is a specified nonzero-vector counterexample, not an argument from absence of a canonical map. Twisting a splitting character changes both sides by the same determinant-root factor, so it does not remove the residual minus sign.

For this fixed standard-lattice construction, agreement with the prescribed root factor on the \emph{entire compact pair} is equivalent to

\[
                            4\mid a\quad\text{and}\quad4\mid b.
\tag{DO21.4}
\]

Alternatively it holds on K\_a⁺×K\_b⁺, where K\_n⁺ consists of elements whose determinant has residue +1. Each K\_n⁺ has index 2. This changes the level or adds an orientation to the lattice data; it is not equality on the unchanged compact group. The new index is a 7-adic unit and an 11-adic unit. No conclusion about an iterated terminal mass or the net step-8 lift count follows without the other recorded arithmetic identities.

\subsection{Proof}\label{proof-2}

\subsubsection{1. Identification of the actual trace lattice with Pan's lattice}\label{identification-of-the-actual-trace-lattice-with-pans-lattice}

Use the additive character ψ\_K=ψ\_Q∘Tr\_\{K/Q₇\}, where ψ\_Q has conductor Z₇. The different of F/Q has v-adic exponent 2 (the discriminant is 7²), so ψ\_K has conductor exponent λ\_K=−2 in the convention that ψ\_K is trivial on p\_K\^{}\{λ\_K\}. The relative extension L/K is tamely ramified quadratic. The element s has L-valuation 1, and 7 has L-valuation 6; hence v\_L(−2ξ)=−5.

Pan (2001), §3.1, equation (3.1.a), uses

\[
 \langle\!\langle\ ,\ \rangle\!\rangle
       =\tfrac12\operatorname{Tr}_{L/K}(h\otimes\overline{h'}).
\tag{DO21.5}
\]

Take V=L\^{}b with h=I\_b and V′=L\^{}a with skew-Hermitian form h′=−2ξI\_a. Conjugating the source tensor factor identifies V⊗\emph{L V′ with the Hom realization used in AM-DO-009 and AM-DO-018. Then (DO21.5) is exactly Tr}\{L/K\}(ξ h\_Hom). The finite operators (DO18.5) instead have Heisenberg cocycle ψ\_f(E\_H/2), where E\_H is the negative of that trace form. Accordingly use ψ\_K\^{}\{-1\} with Pan's displayed trace form, or equivalently negate that form and keep ψ\_K. These two descriptions realize the same central character and operators. Lemmas 7.2--7.3 below give the same displayed αβ formulas for every nontrivial ψ\_K, so the computed determinant-root factor and residual real sign are unchanged.

Set the duality exponents of Pan §1.3, equation (1.3.a), to κ=0 and κ′=−5. Both standard O\_L lattices are selfdual for these conventions: for I\_b the duality ideal is O\_L, and for −2ξI\_a it is s\^{}\{-5\}O\_L. Moreover

\[
                   \kappa+\kappa'=-5=2\lambda_K-1,
\tag{DO21.6}
\]

which is exactly Pan's compatibility condition (3.1.b) in the ramified case. His good lattice (3.2.c) is therefore

\[
 B(L_0,L_0')=L_0^*\otimes L_0'\cap L_0\otimes L_0'^*
                  =O_L^{ab}=\Lambda.
\tag{DO21.7}
\]

It is selfdual for the trace symplectic form and ψ\_K. Since 2 is a unit at 7, the trace-one local character ψ\_K(B(x,w₀x)/2) is trivial on Λ: w₀ preserves Λ, B is integral in the additive-character convention, and its conductor ideal is stable under multiplication by 1/2. Thus φ₀ is exactly the ordinary selfdual-lattice vacuum in Pan's model.

For an element preserving Λ, the lattice operator M\_B is pullback by its inverse (Pan §2.3, equation (2.3.e), or the selfdual formula in §2.4). It fixes the vacuum supported on Λ. The actual split Weil operator is β\^{}Y times the Schrödinger operator, and Pan (3.3.a) compares the latter with α M\_B. Therefore its scalar on this vacuum is exactly αβ\^{}Y. This identifies the normalization being computed; one cannot replace that scalar by the value of the spherical symplectic splitting without this comparison.

\subsubsection{2. The integral Witt decomposition of the identity forms}\label{the-integral-witt-decomposition-of-the-identity-forms}

The two-dimensional Hermitian form I₂ over L/K is anisotropic. Indeed a nonzero isotropic vector would give N\_\{L/K\}(z)=−1 after dividing by one nonzero coordinate. Such a norm must have valuation zero, so z is a unit. Conjugation is the identity on the residue field, hence the residue of its norm is a square in F₇. The residue −1=6 is not a square. This is a contradiction.

The four-dimensional symmetric form I₄ over F₇ is the sum of two hyperbolic planes. For a direct elementary check, (1,1,1,2) is a nonzero isotropic vector. A vector pairing nontrivially with it can be adjusted by a multiple of it, using that 2 is invertible, to give a hyperbolic pair. The orthogonal binary complement has determinant −1 modulo squares, because det I₄=1 and a hyperbolic plane has determinant −1. A binary nondegenerate symmetric form is isotropic precisely when the negative of its determinant is a square: this follows immediately after orthogonal diagonalization by solving A x²+B y²=0. The complement is therefore another hyperbolic plane.

This mod-7 isometry lifts to O. To see this without assuming a lifting classification, suppose an invertible matrix C\_r realizes the isometry mod p\_K\^{}r. Correct C\_r by C\_r(1+p\_K\^{}r X). The first-order correction of the Gram matrix is XᵗJ+JX. For every symmetric error matrix S, choosing X=−J\^{}\{-1\}S/2 solves this linear equation. Inducting and taking the p\_K-adic limit gives an integral invertible isometry. Consequently I₄ is integrally isometric to H⊥H over O, and hence over O\_L. By taking orthogonal sums,

\[
 I_{4k}\simeq H^{\perp 2k},\qquad
 I_{4k+2}\simeq H^{\perp 2k}\perp I_2,
\tag{DO21.8}
\]

with the standard lattice decomposed as the corresponding integral orthogonal sum. Multiplying all forms by −2ξ preserves this decomposition and the lattice admissibility. These are precisely the decompositions needed for Pan's tensor-product identity (3.6.e).

\subsubsection{3. The compact scalar on a determinant section}\label{the-compact-scalar-on-a-determinant-section}

First let the moving group act on a single nondegenerate Hermitian or skew-Hermitian line, and write its element as t=x₀+s y₀∈L¹, t≠1. Pan Lemma 7.2 (printed p.207) states that, against a hyperbolic plane,

\[
                        \alpha(t)\beta^Y(t)=\chi((t-1)s).
\tag{DO21.9}
\]

Pan Lemma 7.3 (printed p.208) gives, against an anisotropic plane,

\[
 \alpha(t)\beta^Y(t)=
 \begin{cases}
   -\chi((t-1)s),&x_0\not\equiv1\pmod{\mathfrak p_K},\\
   +\chi((t-1)s),&x_0\equiv1\pmod{\mathfrak p_K}.
 \end{cases}
\tag{DO21.10}
\]

At a ramified quadratic place, a norm-one element has residue +1 or −1. Its residue agrees with that of x₀, so the extra factor in (DO21.10) is exactly ε(t). Also

\[
       \frac{(t-1)s}{\overline{(t-1)s}}
        =\frac{(t-1)s}{(t^{-1}-1)(-s)}=t.
\tag{DO21.11}
\]

Thus the argument of χ in these formulas is the actual determinant root, not the determinant t. Since χ\textbar K×=1, its value is independent of the chosen root. At t=1 the root may be chosen to be 1 and the formula follows from the identity operator.

Use (DO21.8) for the opposite even-dimensional space. Choose splitting characters on the hyperbolic and anisotropic summands with product χ; all but one may be trivial because every summand has even dimension. Pan (3.6.e) multiplies their αβ factors. Therefore the scalar for the one-dimensional moving subgroup is

\[
                   \chi(y)\epsilon(t)^{d/2},
       \qquad y/\bar y=t,
\tag{DO21.12}
\]

where d is the opposite even rank. For d divisible by 4 the anisotropic summand is absent; for d congruent to 2 modulo 4 it occurs exactly once. This proves (DO21.12), including its phase, for diagonal elements diag(t,1,\ldots,1) of the compact unitary group.

\subsubsection{4. Extension to the full compact group}\label{extension-to-the-full-compact-group}

Pan's Proposition 3.3 (printed pp.179--180) says that the scalar αβ\^{}Y on the compact special unitary group equals the character ζ defined there. When the opposite rank is even, ζ is trivial. Thus αβ\^{}Y factors through the determinant on the compact unitary group. The determinant map is surjective onto L¹, since t↦diag(t,1,\ldots,1) is a section and every norm-one element of L is integral. Formula (DO21.12) consequently determines its value on the entire compact group.

For the target group U(I\_b), d=a and this yields χ\_S(y)ε(det h)\^{}\{a/2\}. For the source group U(I\_a), the Hom action is contragredient: its determinant on the source factor is (det g)\^{}\{-1\}. The corresponding root is x\^{}\{-1\}, the splitting factor is χ\_T(x)\^{}\{-1\}, and ε((det g)\^{}\{-1\})=ε(det g). Hence its contribution is χ\_T(x)\^{}\{-1\}ε(det g)\^{}\{b/2\}. The two commuting unitary actions multiply. This proves (DO21.2).

For a=2,b=4 and h=diag(−1,1,1,1), ε(det h)=−1, proving (DO21.3). The trace-one seed is nonzero (AM-DO-015), so its finite vacuum is nonzero by the injectivity in AM-DO-018. The equality cannot hold with the positive sign. Both lattice indices are 1 since h is integral unitary. Finally the determinant-section elements with determinant −1 exist in each compact group; thus vanishing of both residual characters is equivalent to evenness of a/2 and b/2. This proves (DO21.4). The kernels K\_n⁺ have index exactly 2 by the same determinant section.

\subsection{Scope and next gate}\label{scope-and-next-gate}

The theorem separates two questions that were previously conflated. The mixed CM trace lattice and the finite averaging construction exist, and their relative split-place phases are exactly solved in AM-DO-020. The standard ramified compact phase is now also computed. For ranks 2→4 the literal lattice-indexed phase formula fails; at ranks divisible by 4, or after the stated compact restriction, this particular obstruction disappears. Neither alternative by itself proves the noncompact ramified move formula, the rational compatibility of the fixed two-step averages, an all-polynomial arithmetic \nolinkurl{MainX} identity, or a terminal pairing value. Those remain separate equations.

An arbitrary choice of a local character twist is not a way to change (DO21.3) while retaining the same determinant-root pins: the twist appears on both sides. A new constructor could add orientation/level information to its lattice objects or alter its arithmetic column transformation law. Such a construction must be specified and checked; this theorem does not exclude it, nor does it exclude actual nonzero unitary theta lifts (including the first torus lift being constructed in ALP-LIFT).

\subsection{Primary source and check status}\label{primary-source-and-check-status}

Pan, S.-Y. (2001). Splittings of the metaplectic covers of some reductive dual pairs. \emph{Pacific Journal of Mathematics, 199}(1), 163--226. \url{https://msp.org/pjm/2001/199-1/pjm-v199-n1-p10-p.pdf} . Exact locators used: §1.1 conductor convention; §1.3 equations(1.3.a--c); §2.3 equation(2.3.e) and §2.4; §3.1 equations(3.1.a--b); §3.2 equation(3.2.c); §3.3 equation(3.3.a) and Proposition; §3.6 equations(3.6.a,e); Lemmas7.2--7.3, printed pp.207--208. These statements and the displayed proofs of Lemmas7.2--7.3 were directly read. Proposition3.3 is applied as a published theorem, not advertised as independently re-proved here.

\clearpage

\section{ALP-LIFT-09: an actual first compact unitary theta lift from the primitive CM Jacobian}\label{alp-lift-09-an-actual-first-compact-unitary-theta-lift-from-the-primitive-cm-jacobian}

Run ALP\_MACHINE; \nolinkurl{/root/lift_accounting}, GPT-6 Astra Pro; inherited reasoning effort not separately exposed. This finding uses the actual period matrix and reduced theta of AM-R-004, the trace-one section of ALP-LIFT-07--08, and AM-DO-017--018's finite Heisenberg input. Those inputs retain their own proof/check records. The arithmetic class calculation below was also obtained independently by the root agent; it is recorded once here, as agreed. This is a genuine first theta lift, not a \nolinkurl{thetaG_exists} construction with arbitrary operators.

\subsection{The result and its boundary}\label{the-result-and-its-boundary}

Put \(E=\mathbf Q(\zeta_7)\), \(F=E^+\),

\[
 T=\operatorname{U}(1)(E/F),\qquad
 K_0=\{k\in(\widehat{\mathcal O}_E)^\times:k\bar k=1\}.
 \tag{V1}
\]

The following statements hold.

\begin{enumerate}
\def\labelenumi{\arabic{enumi}.}
\tightlist
\item
  \(T(\mathbb A_{F,f})=T(F)K_0\), and \(T(F)\cap K_0=\mu_{14}\). Consequently \([T]\simeq(T(F_\infty)\times K_0)/\mu_{14}\) as a compact group. For an open subgroup \(K\subset K_0\) with \(K\cap\mu_{14}=1\), its finite component count is exactly \(h_K=[K_0:K]/14\).
\item
  For the actual trace-one finite Schwartz vacuum of the Jacobian, the whole \(K_0\) preserves its one-dimensional eigenline. For any fixed honest global unitary Weil splitting, its action is an actual finite-image character \(\lambda_f\) of \(K_0\).
\item
  Let a homogeneous Fock jet have archimedean character \(\beta_\infty\), including the vacuum shift for that same splitting. If its actual theta value is nonzero, then \(\beta_\infty\lambda_f\) descends to one automorphic character \(\chi\) of \(T\). With probability Haar measure on \([T]\), its theta lift is exactly the original theta function on the target group. There is no finite averaging denominator.
\item
  This gives an explicit nonzero example for the totally definite target \(\operatorname U(2)\): take two copies of the actual Jacobian vacuum and take the \(u_1u_3\) derivative in each copy. The lift at the identity is \((\theta[0;0](0;\tau)c)^2\ne0\), where \(\Sigma_{\mathrm{tr}}=c\sigma_q\) and \(\partial_1\partial_3\sigma_q(0)=1\).
\item
  For an arbitrary finite Schwartz span fixed by a common \(K\), one fixed finite-coset automorphic input realizes every theta function of one fixed archimedean torus weight. Its exact squared \(L^2([T])\) norm is \(h_K\). This is the appropriate generalization when Heisenberg translations do not preserve the original lattice eigencharacter.
\end{enumerate}

The splitting, Schwartz vector, automorphic character and theta integral below are actual objects. However, this result does not identify a chosen character with an independently prescribed source character, instantiate the old elliptic \nolinkurl{VacIn}, construct an entire iterated tower, or prove its six \nolinkurl{AdmS} identities and terminal mass equation. The target automorphic \(L^2\) norm and the whole integral row lattice are not computed here. In particular, the displayed nonzero first lift is not a proof that Q1's requested lift count can be shortened.

\subsection{1. The norm-one torus has one maximal-integral finite class}\label{the-norm-one-torus-has-one-maximal-integral-finite-class}

We give the arithmetic proof, rather than assume a torus class-number formula. The basic inputs are Minkowski's ideal-class bound and Dirichlet's unit theorem (Milne, 2020, Theorems 4.3 and 5.1).

The cyclotomic discriminant is \(\lvert d_E\rvert=7^5\). Its Minkowski bound satisfies

\[
 \left(\frac4\pi\right)^3\frac{6!}{6^6}\sqrt{7^5}<5.
 \tag{V2}
\]

For an entirely rational check, use \(\pi>3\) and

\(\sqrt7<8/3\); the resulting upper bound is \(31360/6561<5\). The order of \(2\pmod7\) is three and the order of \(3\pmod7\) is six. Thus every prime ideal over 2 has norm 8 and every prime ideal over 3 has norm \(3^6\). There are no nontrivial integral ideals of norms 2, 3 or 4. Each ideal class has an integral representative of norm at most 4 by (V2), so

\[
                    \operatorname{Cl}(\mathcal O_E)=1.
 \tag{V3}
\]

For later reference, \(d_F=49\) gives Minkowski bound \(14/9<2\), so \(\operatorname{Cl}(\mathcal O_F)=1\) too; the latter is not needed to prove (V5).

Let \(a=\zeta+\zeta^{-1}\) and \(b=\zeta^2+\zeta^{-2}\). They are units, since their common minimal polynomial is

\(X^3+X^2-2X-1\). At the embeddings indexed by \(1,2,3\), their signs and the signs of \(-1\) are

\begin{longtable}[]{@{}ll@{}}
\toprule\noalign{}
Unit & Signs \\
\midrule\noalign{}
\endhead
\bottomrule\noalign{}
\endlastfoot
\(a\) & \(+,-,-\) \\
\(b\) & \(-,-,+\) \\
\(-1\) & \(-,-,-\) \\
\end{longtable}

These three sign vectors generate \(\{\pm1\}^3\): their negative-sign vectors \(011,110,111\) are independent over \(\mathbf F_2\). Dirichlet gives \(\#(\mathcal O_F^\times/(\mathcal O_F^\times)^2)=8\), because \(F\) is totally real of degree three and its roots of unity are \(\{\pm1\}\). The surjective signature map between these two groups of order 8 is an isomorphism. Therefore

\[
 \mathcal O_F^{\times,+}=(\mathcal O_F^\times)^2.
 \tag{V4}
\]

Take \(t\in T(\mathbb A_{F,f})\). Its fractional \(\mathcal O_E\)-ideal \(I(t)\) satisfies \(I(t)\overline{I(t)}=\mathcal O_E\). By (V3), write

\(I(t)=a_0\mathcal O_E\). Then \(a_0\bar a_0\) is a totally positive unit of \(F\). By (V4), it is \(b_0^2\) for a unit \(b_0\in\mathcal O_F^\times\). Thus \(\gamma=a_0/b_0\in T(F)\) has the same fractional ideal as

\(t\), and \(\gamma^{-1}t\in K_0\). This proves

\[
                 T(\mathbb A_{F,f})=T(F)K_0.
 \tag{V5}
\]

A rational element in \(K_0\) is an algebraic unit whose conjugates all have absolute value one. The kernel of the logarithmic unit map is the roots of unity, so it is in \(\mu(E)=\mu_{14}\). Conversely every such root lies in \(T(F)\cap K_0\). The assertion about \(\mu(E)\) also follows directly from the ramification of the cyclotomic field: an additional prime in the root's order, a factor 4, or a factor \(7^2\) would force an excluded ramified prime or a subfield of degree greater than six.

It follows that the natural map \(T(F_\infty)\times K_0\to[T]\) is onto, with kernel the diagonally embedded \(\mu_{14}\). It is an isomorphism after taking that quotient: the source is compact and the target Hausdorff. In particular, the automorphic quotient is compact without an appeal to reduction theory. If \(K\cap\mu_{14}=1\), then

\[
 T(F)\backslash T(\mathbb A_F)/(T(F_\infty)K)
       =K_0/(\mu_{14}K),\qquad
 h_K=[K_0:K]/14.
 \tag{V6}
\]

Such levels exist. For example, intersect any given open subgroup with the principal congruence subgroup at 3. Roots of order dividing 14 inject under reduction at 3, so this kills the intersection with \(\mu_{14}\). No unknown torus class representatives enter (V6).

\subsection{2. A completely specified honest global splitting exists here}\label{a-completely-specified-honest-global-splitting-exists-here}

Let \(V=E^n\) with its positive definite Hermitian form \(\sum x_i\bar y_i\), and let \(W=E\) be the skew-Hermitian line with coefficient \(-\xi\), where \(\xi=(\zeta-\zeta^{-1})/7\). On restriction to \(\mathbf Q\), the tensor alternating form is

\[
 \mathfrak b(x,y)=-\operatorname{Tr}_{E/\mathbf Q}
       \left(\xi\sum x_i\bar y_i\right).
 \tag{V7}
\]

This is precisely the form \(E_H=\operatorname{Im}H\) in AM-DO-018. Use the global additive character with \(\psi_\infty(r)=e^{2\pi ir}\) and \(\psi_f(r)=e^{-2\pi ir}\) on rational numbers, and over \(F\) compose it with \(\operatorname{Tr}_{F/\mathbf Q}\). Thus the finite Heisenberg formula below and the chosen tensor symplectic form have the same convention. This explicit sign choice is part of this construction; it is not an identification with an independently fixed source pin.

Here is an explicit choice of the required splitting characters. Put \(K_- =\mathbf Q(\sqrt{-7})\), \(\varpi=\sqrt{-7}\) with positive imaginary part, and let \(\epsilon_7:(\mathcal O_{K_-}/\varpi)^\times\to\{\pm1\}\) be the quadratic character of \(\mathbf F_7^\times\). The Minkowski bound \((2/\pi)\sqrt7<2\) proves that \(K_-\) has class number one. For an ideal \(\mathfrak a=(a)\) prime to \(\varpi\), set

\[
             u(\mathfrak a)=\epsilon_7(a)\,a/|a|.
 \tag{V8}
\]

This is independent of the generator because the only units are \(\pm1\) and \(\epsilon_7(-1)=-1\). It is multiplicative. Define a unitary idele-class character \(\mu_-\) by the unramified local ideal values (V8), and by

\[
 \mu_{-,\infty}(z)=|z|/z,\qquad
 \mu_{-,\varpi}(v)=\epsilon_7(v)\ (v\text{ a unit}),\qquad
 \mu_{-,\varpi}(\varpi)=i.
 \tag{V9}
\]

To verify rational triviality, write any \(a\in K_-^\times\) as \(a=\varpi^m b\) with \(b\) a \(\varpi\)-adic unit. The unramified finite factors multiply to \(\epsilon_7(b)b/|b|\); the ramified factor is \(i^m\epsilon_7(b)\); the infinite factor is \((i^mb/|b|)^{-1}\). Their product is one. This verifies the idele-class character, including its ramified normalization. Its restriction to \(\mathbb A_\mathbf Q^\times\) is \(\epsilon_{K_-/\mathbf Q}\): away from 7 the uniformizer value is the quadratic residue character, at 7 its value on units is \(\epsilon_7\) and its value on 7 is one, and at infinity it is the sign character. These are the local components of the quadratic character of discriminant \(-7\).

Since \(E=FK_-\), put

\[
                   \mu_E=\mu_-\circ N_{E/K_-}.
 \tag{V10}
\]

The norm compatibility of quadratic class-field characters gives \(\mu_E|_{\mathbb A_F^\times}=\epsilon_{E/F}\). Choose \(\chi_W=\mu_E\), and choose \(\chi_V=1\) when \(n\) is even and \(\chi_V=\mu_E\) when \(n\) is odd. They satisfy exactly

\[
 \chi_W|_{\mathbb A_F^\times}=\epsilon_{E/F}^{\dim W},
 \qquad
 \chi_V|_{\mathbb A_F^\times}=\epsilon_{E/F}^{\dim V}.
 \tag{V11}
\]

Fix the global unitary Weil splitting determined by these data. The construction and its rational compatibility are the standard global splitting theorem, in the explicit form stated by Borade et al.~(2025), §§3.1 and 4.1.1. We use only that construction and the definition of the theta integral from §4.1, not their root-number or central-\(L\)-value nonvanishing theorem. The two unitary actions commute and agree with the canonical rational splitting on rational points. Thus they give an honest representation of \(G(\mathbb A_F)\times T(\mathbb A_F)\), where \(G=\operatorname U(V)\). All scalar phases below refer to this fixed representation. In particular they are not phases that may be adjusted separately on different lattices.

The Schrödinger position space used in (V12) is the rational Lagrangian supplied by the integral period basis. It need not be \(F\)-stable. This causes no change in the representation: taking the trace of the central coordinate identifies the \(F\)-Heisenberg representation with the representation of the underlying rational symplectic space having the displayed character. A rational change of polarization transports both the Schwartz representation and the rational theta distribution. Thus the usual global unitary splitting acts on this specified rational Schrödinger model, and the theta sum in that model is indexed by \(\mathbf Q^{3n}\).

\subsection{3. The actual finite vacuum and its maximal-compact eigencharacter}\label{the-actual-finite-vacuum-and-its-maximal-compact-eigencharacter}

Use the actual period coordinates (u=Pz), \(Y=\operatorname{Im}\tau\) from AM-R-004. In each copy take \(\delta=(1,1,1)^t\), \(\varepsilon=(1,0,1)^t\), and

\[
 \phi_{0,f}(q)=
   -i\,1_{\delta/2+\widehat{\mathbf Z}^{\,3}}(q)
        \psi_f(-q^t\varepsilon/2).
 \tag{V12}
\]

This is the actual input of AM-DO-016/018, not a variable in a formal theta-family interface. For \(n\) copies use its tensor power. On the rational period vector \(t=A+\tau B\), its Heisenberg action is

\[
 (W_t\phi)(q)=\psi_f(-A^tq+A^tB/2)\phi(q-B).
 \tag{V13}
\]

Let \(\widehat\Lambda\) be the completion of the product period lattice identified with \(\widehat{\mathcal O}_E^n\), and put \(w_0=-(\zeta+\zeta^2+\zeta^4)\). Its finite trace-one Heisenberg eigencharacter (a semicharacter for the operators \(W_x\)) is

\[
 \alpha_f(x)=\psi_f\bigl(\mathfrak b(x,w_0x)/2\bigr),
       \qquad x\in\widehat\Lambda.
 \tag{V14}
\]

This is an adelic additive-character expression. The notation \(e^{\pi i\mathfrak b(x,w_0x)}\) is not used for arbitrary adelic \(x\). On the rational integral lattice (V14) equals the sign semicharacter of AM-DO-018, since its value is \(\pm1\). The eigenrelations (V13) for (V12) therefore extend from that dense rational lattice to all of \(\widehat\Lambda\), by continuity.

The finite eigenspace

\[
 \mathscr V_\Lambda=
 \{\phi:W_x\phi=\alpha_f(x)\phi
                         \text{ for all }x\in\widehat\Lambda\}
 \tag{V15}
\]

inside the finite Schwartz representation has dimension one. Here is a direct justification tied to the actual theta model. AM-DO-018's reading map \(R\) is injective on finite Schwartz functions and intertwines (V13) with normalized Mumford translation. It sends (V15) into the principal theta line on the actual complex torus, which has dimension one. Formula (V12) supplies a nonzero vector in (V15), proving equality of dimensions. For tensor powers the same reading map and the product principal theta line prove the identical assertion. This is also the usual selfdual-lattice vacuum uniqueness in the Heisenberg representation.

Every \(k\in K_0\) preserves \(\widehat\Lambda\), preserves \(\mathfrak b\), and commutes with \(w_0\). Consequently it preserves (V14). Heisenberg covariance of the \emph{fixed} honest Weil representation shows that it preserves (V15). There is therefore a unique scalar

\[
 \omega_f(k)\phi_{0,f}^{\otimes n}
       =\lambda_f(k)\phi_{0,f}^{\otimes n}.
 \tag{V16}
\]

The group law makes \(\lambda_f\) a character. Unitarity makes its values have absolute value one; smoothness makes its kernel open in \(K_0\), so its image is finite and consists of roots of unity. This proves that the scalar is an actual finite-order character, without guessing a Weil index. For an explicit evaluation prescription, the value of (V12) at \(q=\delta/2\) is \(i\); therefore

\[
 \lambda_f(k)=i^{-n}
   (\omega_f(k)\phi_{0,f}^{\otimes n})
                 (\delta/2,\ldots,\delta/2).
 \tag{V17}
\]

The same line argument applies to every selfdual \(\mathcal O_E\)-lattice vacuum of AM-DO-018, because scalar norm-one integral units preserve that lattice and its trace-one character. It does not assert that an arbitrary finite Heisenberg translate has the same character (V14).

\subsection{4. The reduced algebraic jet is an actual homogeneous Fock vector}\label{the-reduced-algebraic-jet-is-an-actual-homogeneous-fock-vector}

For one copy put

\[
 K(z,x)=\exp\left(\frac\pi2z^tY^{-1}z
                    +\pi i x^t\tau x+2\pi i x^tz\right),
 \qquad \varphi_{0,\infty}(x)=K(0,x).
 \tag{V18}
\]

The first factor is exactly the Chern reduction used throughout this run. Define the commuting creation operators

\[
 r_j=-\frac i2\sum_k(Y^{-1})_{jk}
           \left(\partial_{x_k}-2\pi i(\bar\tau x)_k\right).
 \tag{V19}
\]

A direct calculation gives

\[
 r_jK(z,x)=\left(2\pi ix_j+\pi(Y^{-1}z)_j\right)K(z,x)
                  =\partial_{z_j}K(z,x).
 \tag{V20}
\]

Replace the operators by their linear combinations dual to \(u=Pz\):

\(r_{u_i}=\sum_j(P^{-1})_{ji}r_j\). Thus for every multi-index \(\alpha\),

\[
 \varphi_{\alpha,\infty}(x)
    =\left.\partial_u^\alpha K(P^{-1}u,x)\right|_{u=0}
    =r_u^\alpha\varphi_{0,\infty}(x).
 \tag{V21}
\]

It is a polynomial times a positive Gaussian, hence Schwartz. The real torus \(T(F_\infty)\simeq(S^1)^3\) preserves the positive complex polarization. Its action preserves the vacuum line and acts diagonally on the three CM creation directions. One may check preservation directly: the differentiated real Heisenberg operator with parameter \((A,B)\) acts on the vacuum as \(-2\pi i(A+\tau B)^t x\). Its complex annihilator is therefore the kernel of \(A+\tau B\), which is preserved by the CM torus. Its unique Schwartz vacuum line is consequently preserved. Consequently a monomial of fixed CM multidegree is an eigenvector. For \(n\) copies and a multi-index \(\alpha=(\alpha_{ij})\), define its actual eigencharacter by

\[
 \omega_\infty(t)\varphi_{\alpha,\infty}
      =\beta_\alpha(t)\varphi_{\alpha,\infty}.
 \tag{V22}
\]

This fixes the sign and the vacuum character for the splitting in §2. The action on creation operators determines the relative weights, while the action on the vacuum supplies its scalar shift. No replacement of (V22) by an unshifted polynomial degree is made. The eigenspace assertion can equivalently be checked from (V19): Heisenberg covariance transports the holomorphic creation space by the three CM characters, and the simultaneous monomials have the corresponding characters multiplied by the vacuum character.

For clarity, the actual Gaussian norm is also explicit. With Lebesgue measure in the fixed rational position coordinates,

\[
 \|\varphi_{\alpha,\infty}\|_2^2
  =\det(2Y)^{-n/2}\,
       \prod_{i,j}\alpha_{ij}!\left(\frac\pi{G_i}\right)^{\alpha_{ij}}.
 \tag{V23}
\]

Indeed the vacuum norm is \(\det(2Y)^{-n/2}\). From (V19), \([r_{u_i}^{*},r_{u_j}]=\pi\delta_{ij}/G_i\), and \(r_{u_i}^{*}\varphi_{0,\infty}=0\). Moving annihilation operators past creation operators proves (V23) by induction. The finite vector (V12) has norm one, because its support has volume one and every nonzero value has absolute value one. Thus (V23) is also the global Schwartz norm for this input.

The theta series at the identity is

\[
 \vartheta_\alpha(1,1)
   =\sum_{q\in\mathbf Q^{3n}}
       \phi_{0,f}^{\otimes n}(q)\varphi_{\alpha,\infty}(q)
   =\left.\partial_u^\alpha
                  R(\phi_{0,f}^{\otimes n})(u)\right|_{u=0}.
 \tag{V24}
\]

This is an equality of an actual adelic theta sum and an actual reduced CM derivative. Differentiation is legitimate because the finite support lies in a bounded-denominator rational lattice and the positive Gaussian gives locally normally convergent series.

AM-DO-018 gives

\[
 R(\phi_{0,f}^{\otimes n})(u)
   =\theta[0;0](0;\tau)^n\prod_{j=1}^n\Sigma_{\mathrm{tr}}(u_j).
 \tag{V25}
\]

ALP-LIFT-08 gives \(\sigma_q(u)=u_1u_3-u_2^2+\text{higher terms}\), and AM-DO-015 identifies \(\Sigma_{\mathrm{tr}}=c\sigma_q\) with \(c\ne0\). Taking \(\alpha_{1j}=\alpha_{3j}=1\), \(\alpha_{2j}=0\) for every \(j\), gives the exact nonzero value

\[
             \vartheta_\alpha(1,1)
                  =(\theta[0;0](0;\tau)c)^n\ne0.
 \tag{V26}
\]

For this input the norm (V23) is \(\det(2Y)^{-n/2}(\pi^2/(G_1G_3))^n\). Neither (V26) nor that norm replaces the common period by an unrecorded scalar.

\subsection{5. One actual character gives the nonzero theta lift}\label{one-actual-character-gives-the-nonzero-theta-lift}

For the fixed splitting, write

\[
 \vartheta_\alpha(g,t)=
       \sum_q(\omega(g,t)\varphi_\alpha)(q),
 \qquad
 \varphi_\alpha=\phi_{0,f}^{\otimes n}
                           \otimes\varphi_{\alpha,\infty}.
 \tag{V27}
\]

The theta function is left invariant under rational points, by the rationally compatible global Weil construction. Its two group actions commute. Equations (V16) and (V22) therefore give

\[
 \vartheta_\alpha(g,t_\infty k)
       =\beta_\alpha(t_\infty)\lambda_f(k)
                            \vartheta_\alpha(g,1).
 \tag{V28}
\]

For \(\zeta\in\mu_{14}\), the diagonal element is rational. Evaluating (V28) at \(g=1\) and using the already established nonzero value (V26) proves

\[
                  \beta_\alpha(\zeta)\lambda_f(\zeta)=1.
 \tag{V29}
\]

This is precisely the finite-root-unit compatibility which is sometimes lost if only the polynomial degree is retained. Its proof uses the actual theta series and rational splitting; it is not an imposed compatibility condition assumed for a hypothetical lift.

By (V5)--(V6) and (V29), the formula

\[
 \chi_\alpha(\gamma t_\infty k)
        =\beta_\alpha(t_\infty)\lambda_f(k),
 \quad \gamma\in T(F),\ t_\infty\in T(F_\infty),\ k\in K_0,
 \tag{V30}
\]

is well-defined and is a continuous unitary automorphic character. All automorphic forms on this anisotropic torus are cuspidal, since a torus has no proper parabolic subgroups. Equation (V28), together with rational automorphy and (V5), now holds for every adelic \(t\) in the form

\[
            \vartheta_\alpha(g,t)
                   =\chi_\alpha(t)\vartheta_\alpha(g,1).
 \tag{V31}
\]

Choose Haar measure on \([T]\) of total mass one and use the theta-lift convention with a conjugated automorphic input. Then

\[
 \boxed{\Theta(\varphi_\alpha,\chi_\alpha)(g)
   =\int_{[T]}\vartheta_\alpha(g,t)\overline{\chi_\alpha(t)}\,dt
   =\vartheta_\alpha(g,1).}
 \tag{V32}
\]

There is no limiting procedure, convergence regularization, or unproved nonvanishing theorem in this calculation. The quotient is compact; (V31) makes the integral elementary. The source character has \(L^2\) norm one. Taking \(n=2\) proves the claimed actual nonzero \(U(1)\to U(2)\) lift. With a different fixed Haar measure, the right side is multiplied by its explicitly retained total volume.

The common-period-normalized identity value is exactly one by (V26). This identifies one actual nonzero matrix value. It is not an assertion that the whole automorphic function is primitive in the source's mass-weighted integral row module, and no valuation is assigned to a transcendental period itself.

\subsection{6. General finite inputs: an exact common first-image projector}\label{general-finite-inputs-an-exact-common-first-image-projector}

This extension explains exactly what remains for translated inputs. Let \(\mathcal S\) be a finite-dimensional span of finite Schwartz functions. Let \(K\subset K_0\) fix every vector in \(\mathcal S\), and shrink it at 3 so that \(K\cap\mu_{14}=1\). Such a common open subgroup exists by smoothness and finite dimensionality. Keep any fixed homogeneous archimedean weight \(\beta\). The statement applies to any finite-dimensional subspace of that weight, not only a single monomial. Write \(h=h_K\).

The identity finite component of \([T]/K\) is the image of \(T(F_\infty)K\), and it is an injective copy of \(T(F_\infty)\), because \(K\cap\mu_{14}=1\). Define

\[
 f_{K,\beta}(t)=
 \begin{cases}
 h\,\beta(t_\infty),&[t]=[t_\infty k],\ k\in K,\\
 0,&\text{on the other finite components}.
 \end{cases}
 \tag{V33}
\]

This is a well-defined smooth automorphic form of weight \(\beta\), with

\[
             \|f_{K,\beta}\|_2^2=h,
                 \qquad \|f_{K,\beta}\|_\infty=h.
 \tag{V34}
\]

Each component has measure \(1/h\). For any \(\phi_f\in\mathcal S\) and any archimedean vector of weight \(\beta\), the theta function on the identity component is \(\beta(t_\infty)\vartheta_\phi(g,1)\). Hence

\[
                \boxed{\Theta(\phi,f_{K,\beta})(g)
                           =\vartheta_\phi(g,1).}
 \tag{V35}
\]

The same input \(f_{K,\beta}\) works for the entire specified finite span. If the source automorphic input is required to have norm one, replacing it by \(h^{-1/2}f_{K,\beta}\) multiplies the output by \(h^{-1/2}\). Thus the exact added normalization cost is \(\frac12\log h\) in logarithmic norm, not an unspecified projection constant.

There is also an explicit finite character decomposition. Set \(A=K_0/K\), with its embedded subgroup \(\mu_{14}\). Let

\[
 \mathcal X_\beta=
   \{\eta\in\widehat A:
               \eta|_{\mu_{14}}=\beta^{-1}|_{\mu_{14}}\}.
 \tag{V36}
\]

Every character of a subgroup of a finite abelian group extends: one may extend successively over cyclic generators by choosing the necessary roots of unity. The extensions form a torsor for \(\widehat{A/\mu_{14}}\). Thus \(\#\mathcal X_\beta=h\), all its finite values are roots of unity, and each

\(\chi_\eta(t_\infty k)=\beta(t_\infty)\eta(k)\) descends to \([T]\). Finite character orthogonality gives

\[
                 f_{K,\beta}=\sum_{\eta\in\mathcal X_\beta}\chi_\eta.
 \tag{V37}
\]

No coefficient \(1/h\) occurs in (V37). For one component, choosing representatives \(a\in A/\mu_{14}\) gives the exact projection

\[
 \Theta(\phi,\chi_\eta)(g)
   =\frac1h\sum_{a\in A/\mu_{14}}
                   \overline{\eta(a)}\vartheta_\phi(g,a).
 \tag{V38}
\]

The summand is independent of representative because rational-root automorphy cancels the two characters. Formula (V38) displays the only averaging denominator; its \(p\)-adic contribution is at worst \(v_p(h)\) once the specified theta values have actually been placed in a common integral lattice. If (V35) is nonzero, at least one of the \(h\) irreducible-character lifts in (V37) is nonzero. This is a consequence for the explicitly specified finite family; it is not nonvanishing for a freely prescribed Hecke character.

\subsection{7. What this closes, and what it does not close}\label{what-this-closes-and-what-it-does-not-close}

The trace-one scalar torus argument supplies an actual nonzero first compact lift with the primitive mixed CM type. It also supplies a common actual first-image map for any selected finite Schwartz span at one weight, with an exact finite-index normalization. In that limited first-lift setting, existence of an automorphic input is no longer a merely formal interface.

For the base vacuum, the finite and archimedean input normalizations are (V12), (V23), and the norm-one character in (V30). The unknown target automorphic mass norm is a different quantity. The exact finite character (V16) is fixed by the global splitting; identifying it with a source-prescribed determinant-root character requires the corresponding local calculation. AM-DO-020 performs that calculation at relative split places, including 11, for the stated unramified splitting. No omitted ramified phase is declared equal to one here.

For DP59/DP65's permitted translated sections, applying §6 requires an actual common stabilizer \(K\), the exact archimedean weight, and the source's allowed automorphic input space. An unrestricted finite Schwartz span may be larger than the prescribed one. The finite-coset input in (V33) is generally a sum of characters, and its cost (V34) must be retained if the source requires unit norm. The accepted DP59/DP65 field-span and sharp-jet results are not re-proved or promoted here to a full iterated-image statement.

Nothing in (V32) or (V35) supplies the later branching map, all-degree \nolinkurl{ClsK.Core} with its specified phases, simultaneous terminal mass identities, the global \nolinkurl{MainX} laws, or a changed DP44 cost with all \(g_j,e_j,H\) identified. Those are the remaining gates for Q1's requested shorter tower.

\subsection{Primary references actually consulted}\label{primary-references-actually-consulted}

Borade, N., Franzel, J., Girsch, J., Yao, W., Yu, Q., \& Zelingher, E. (2025). \emph{On tori periods of Weil representations of unitary groups} (arXiv:2402.16808v2). \url{https://arxiv.org/html/2402.16808v2}. Exact used locators: §3.1, tensor alternating form and Heisenberg representation; §3.1.1, splitting characters; §4.1, rational global splitting, theta sum, and conjugated-input theta integral; §4.1.1, global character restrictions and tensor-product splitting. No nonvanishing theorem from that paper is used.

Milne, J. S. (2020). \emph{Algebraic number theory} (Version 3.08). \url{https://www.jmilne.org/math/CourseNotes/ANT.pdf}. Exact used locators: Theorem 4.3, printed pp.~70--71, Minkowski bound; Theorem 5.1, printed p.~85, unit rank; the logarithmic kernel discussion in §5, printed pp.~86--87.

Prasad, D. (2000). Theta correspondence for unitary groups. \emph{Pacific Journal of Mathematics, 194}(2), 427--438. \url{https://www.math.iitb.ac.in/~dprasad/pacific.pdf}. Used as an independent convention check: §1, printed pp.~427--428, the dependence on two splitting characters and the passage from metaplectic to honest unitary representations.

\subsection{Check status}\label{check-status}

The author read the complete saved proof, including the arithmetic class and character construction, finite trace-one eigenline, Chern/Fock derivative and Gaussian norm calculation, single-character integral, and general finite-coset projector. The exact algebra certificate in actual\_torus\_certificate.py passed. Actual check events and their scope are recorded in \nolinkurl{lift_accounting_CHECK_EVENTS.tsv.} No Lean compilation or formal proof is claimed.

\clearpage

\section{AM-R-007: unitary lattice classes have representatives supported at relative split places}\label{am-r-007-unitary-lattice-classes-have-representatives-supported-at-relative-split-places}

Author: \nolinkurl{/root}, GPT-6 Astra Pro; effort not externally exposed; ALP\_MACHINE, 21 November 2026. This is a concrete application of strong approximation to the moved-lattice comparison. It is a statement about double-coset representatives, not an assertion that the two-average theta family already has the required rational equivariance.

\subsection{Statement}\label{statement-8}

Let \(E/F\) be a quadratic extension of number fields, with its nontrivial involution denoted by a bar. Let \(V\) be a nondegenerate Hermitian space of dimension \(n\geq2\), let \(G=U(V)\), \(G^{\rm der}=SU(V)\), and \(T=\ker(N_{E/F}:\operatorname{Res}_{E/F}\mathbf G_m\to\mathbf G_m)\). Let \(\mathcal S\) be the set of finite places of \(F\) that split in \(E\), and choose one \(v_0\in\mathcal S\). For every compact open subgroup \(K\subset G(\mathbb A_{F,f})\), \[
 G(\mathbb A_{F,f})=G(F)\,G_{\mathcal S}^{\rm fin}\,K,
 \tag{R7.1}
\] where \(G_{\mathcal S}^{\rm fin}\) denotes adeles equal to the identity outside a finite subset of \(\mathcal S\). The finite subset is allowed to depend on the represented class.

For the actual field \(E=\mathbb Q(\zeta_7)\), \(F=E^+\), standard form \(V=E^n\) with Gram matrix \(I_n\), and its full standard lattice stabilizer \(K_0=\prod_vU(I_n,\mathcal O_{E,v})\), the stronger statement is \[
 G(\mathbb A_{F,f})=G(F)\,G(F_{v_{11}})\,K_0,
 \tag{R7.2}
\] where \(v_{11}\) is the unique place of \(F\) above 11. It splits in \(E/F\), since the subgroup generated by \(11\bmod7=4\) has order three and does not contain complex conjugation.

Thus every standard-genus Hermitian lattice class needed at maximal level has a representative obtained by moving only at 11. At a smaller level, (R7.1) supplies representatives using finitely many relative split places. Neither assertion bounds the denominators of the representative, the index of an intersection lattice, its stabilizer order, or a terminal mass.

\subsection{Proof of the split-place reduction for arbitrary compact level}\label{proof-of-the-split-place-reduction-for-arbitrary-compact-level}

Choose a diagonal Hermitian basis over \(E\); this is possible by Hermitian Gram--Schmidt in characteristic zero. It gives an explicit section of determinant, \[
 s:T\longrightarrow G,\qquad t\longmapsto
             \operatorname{diag}(t,1,\ldots,1).
 \tag{R7.3}
\] Indeed multiplying the first diagonal coordinate by a norm-one scalar preserves the diagonal form. This also proves surjectivity of determinant on all local and global point groups.

Replace \(K\), if necessary, by a product compact open subgroup \(K'=\prod_vK'_v\subset K\), standard at almost every place. Proving (R7.1) for \(K'\) proves it for \(K\), so write this product subgroup as \(K\) from now on. Put \(K_{T,v}=\det K_v\). This is an open subgroup of \(T(F_v)\): the local analytic section (R7.3), near the identity, shows that the image of an open identity neighborhood is open. At almost all places it contains all integral points of \(T\), and in the standard unramified model it equals them.

At a nonsplit finite place every norm-one element of \(E_v^\times\) is a unit: its valuation multiplied by the residue degree is the valuation of its norm, which is zero. Thus \(T(F_v)\) is compact, and \(K_{T,v}=T(F_v)\) at all but finitely many nonsplit places.

Given \(g\in G(\mathbb A_{F,f})\), set \(t=\det g\). There exists \(t_0\in T(F)\) such that \[
 t_0^{-1}t_v\in K_{T,v}
       \quad\text{at every nonsplit finite place }v.
 \tag{R7.4}
\] Only finitely many conditions need attention. For completeness the required weak approximation on this particular torus has a direct proof. Local Hilbert 90 writes each prescribed norm-one local element as \(a_v/\bar a_v\). Weak approximation in the number field \(E\) gives a nonzero \(a\in E\) sufficiently close to these finitely many \(a_v\). The continuous map \(a\mapsto a/\bar a\) then puts \(t_0=a/\bar a\) in the required open cosets. This uses only finite local approximation, not strong approximation for a torus with a finite omitted set.

At the split places, \(t_0^{-1}t_v\) belongs to \(K_{T,v}\) for all but finitely many \(v\), since it is an adele and \(K\) is standard almost everywhere. We may therefore write \[
 t=t_0\,t_{\mathcal S}\,k_T,
 \tag{R7.5}
\] where \(t_{\mathcal S}\) is supported at finitely many split places and \(k_T\in\prod_vK_{T,v}\). Choose \(k\in K\) with \(\det k=k_T\), choosing a lift at each place, and let \(a=s(t_{\mathcal S})\), a finite-support split-place adele. Then \[
 g_2=s(t_0)^{-1}gk^{-1},\qquad h=g_2a^{-1}
                     \in SU(V)(\mathbb A_{F,f}).
 \tag{R7.6}
\]

The group \(SU(V)\) is connected, simply connected, and absolutely almost simple for \(n\geq2\); over an algebraic closure it is \(SL_n\). At the chosen split place \(v_0\), it is \(SL_n(F_{v_0})\), hence noncompact. The strong approximation theorem for simply connected absolutely almost simple groups therefore applies with the omitted set consisting of all archimedean places together with \(v_0\). The exact input is Rapinchuk (2013), Theorem 2.3, printed p.~281, recording the Kneser--Platonov theorem in characteristic zero.

Apply it to the compact open subgroup \(K''=(aKa^{-1})\cap SU(V)(\mathbb A_{F,f})\). It gives \[
 h=\gamma u_0k'',\qquad
 \gamma\in SU(V)(F),\quad u_0\in SU(V)(F_{v_0}),
 \quad k''\in K''.
 \tag{R7.7}
\] To see the factorization directly from density, approximate \(h\) outside \(v_0\) by \(\gamma\) so that \(\gamma^{-1}h\in K''_v\) at each other finite place; put these components into \(k''\) and put its \(v_0\)-component equal to one. The remaining component is \(u_0\).

Since \(k''a=a(a^{-1}k''a)\), substitution in (R7.6) yields \[
 g=s(t_0)\gamma\,(u_0a)\,(a^{-1}k''a)k.
 \tag{R7.8}
\] The first factor is rational, the middle factor is supported at finitely many relative split places, and the last two factors lie in \(K\). This proves (R7.1). The order of factors in this argument matters: no nonrational adele was commuted past a rational matrix without conjugating the compact subgroup.

\subsection{Proof of the maximal-level statement for the cyclotomic field}\label{proof-of-the-maximal-level-statement-for-the-cyclotomic-field}

The torus calculation in ALP-LIFT-09, equations (V2)--(V6), proves \[
 T(\mathbb A_{F,f})=T(F)\,T(\widehat{\mathcal O}_F)
 \tag{R7.9}
\] for \(E=\mathbb Q(\zeta_7)\). Its ingredients are the class-number one calculation for \(E\) and the explicit full signatures of \(-1,\zeta+\zeta^{-1},\zeta^2+\zeta^{-2}\), which imply that every totally positive unit of \(F\) is a square. That proof is cited once under its own finding identifier; it is not counted again here.

For the standard lattice, diagonal norm-one units in (R7.3) show \(\det K_0=T(\widehat{\mathcal O}_F)\). Thus for each \(g\), equation (R7.9) permits \(t_{\mathcal S}=1\) in (R7.5). In the preceding proof \(a=1\), and (R7.8) has its only noncompact finite factor at \(v_{11}\). This is (R7.2). The unique prime of \(F\) above 11 follows because the image of 4 in \((\mathbb Z/7\mathbb Z)^\times/\{\pm1\}\) has order three; its two lifts to \(E\) exhibit the relative splitting.

\subsection{What the representative theorem reduces, and what it does not}\label{what-the-representative-theorem-reduces-and-what-it-does-not}

There is a useful exact conditional transport principle. Suppose two assignments on \(G(\mathbb A_{F,f})\) have the same rational left-equivariance and the same right-\(K\) equivariance, with specified invertible transport operators, and agree on all finite-support split-place representatives. Factorization (R7.1) then proves equality everywhere: apply the common right transport and common rational left transport to the representative equality. This short argument proves the principle, including its explicit equivariance hypotheses.

AM-DO-017 gives exact transport when all lattices, including the base lattice, are carried by a fixed complex-linear isometry. That does not yet give the rational left-equivariance needed for the fixed-base two-average family (DO18.3). Returning the moved base to the original base introduces its finite triangle phase. Therefore the present theorem, even combined with split-place phase equality, leaves the rational triangle/order law and compact vacuum character as real proof obligations. It cannot erase them by choosing representatives.

This theorem also has no quantitative height assertion. Strong approximation supplies existence in an open adelic set; it does not bound the matrices \(\gamma\), the supports or denominators of \(a\), the Smith indices \(q_1,q_2\), or the finite stabilizer factors. Those quantities enter AM-DO-018's explicit costs and AM-R-002's terminal content formula. A shorter lift chain still requires those actual data and all six linked identities.

In geometric terms the double cosets classify the indicated Hermitian lattice genus with level structure. The determinant sequence \(1\to SU\to U\to T\to1\) separates its simply connected derived part from its norm-one torus part. This is the reductive-group and torsor side of the arithmetic comparison; it does not identify integral theta sections, crystalline filtrations, or an equality of motivic L-functions. The exact SGA3 and Picard framework locators are maintained in the geometric dictionary.

\subsection{References}\label{references-1}

Rapinchuk, A. S. (2013). Strong approximation for algebraic groups. In E. Breuillard \& H. Oh (Eds.), \emph{Thin groups and superstrong approximation} (MSRI Publications, Vol. 61, pp.~269--298). Cambridge University Press. Exact theorem: Theorem 2.3, p.~281. Directly inspected primary author exposition: \url{https://library.slmath.org/books/Book61/files/70rapi.pdf} .

Platonov, V. P. (1969). The problem of strong approximation and the Kneser--Tits conjecture for algebraic groups. \emph{Mathematics of the USSR-Izvestiya, 3}(6), 1139--1147. \url{https://doi.org/10.1070/IM1969v003n06ABEH000838} . The theorem is used in the exact form quoted above; no reproof of strong approximation is claimed.

\clearpage

\section{AM-K017 --- integral logarithmic Hessians, stronger actual theta slopes, and exact mixed-jet recurrences}\label{am-k017-integral-logarithmic-hessians-stronger-actual-theta-slopes-and-exact-mixed-jet-recurrences}

Author: \nolinkurl{/root/kummer_extension}, GPT-6 Astra Pro, ALP\_MACHINE, 21 November 2026 UTC. This supplements AM-K012--016 in \nolinkurl{jacobian_reduced_theta.md}. It concerns the actual translated germ \(F\) and the actual trace-one germ \(\sigma_q\). Every reduction below is a reduction of Hurwitz derivatives, not of divided Taylor coefficients.

\subsection{1. Statement}\label{statement-9}

Retain \[
C:y^2=x^7-1,\quad E=\mathbb Q(\zeta_7),\quad
\omega_i=x^{i-1}dx/y,\quad D_i=\partial/\partial u_i,
\quad q(u)=u_1u_3-u_2^2.
\] Let \(W_{ij}=-\partial_i\partial_j\log\sigma_q\) be the actual algebraic Kleinian Hessian. Its explicit expressions, normalization, and integral Abel chart are AM-K012. For an odd prime \(p\ne7\), let \(q_i=p^{h_i}\) be the sufficient restricted-power nilpotence exponent in the table (K13.2).

\textbf{(a) Global constant lemma.} The rational function \[
\overline{D_i^{q_i-1}W_{ij}}
\tag{K18.1}
\] on the good special fibre is a constant. It is zero unless \[
iq_i+j\equiv0\pmod7.
\tag{K18.2}
\] In particular the pure constant vanishes whenever \(q_i\not\equiv-1\pmod7\). This proves the additional log-derivative hypothesis in (K13.7) geometrically, at all such primes. The three exact 11-adic checks in AM-K013 are consistent independent finite checks, no longer necessary premises for that vanishing.

\textbf{(b) Improved actual finite bounds.} At every prime above \(p\), for either \(F\) or \(\sigma_q\), \[
v_p(\partial_i^n f(0))\geq
\max\{0,c_i(p)n-s_p(n)/(p-1)\},
\tag{K18.3}
\] where the following constants are now proved:

\begin{longtable}[]{@{}llll@{}}
\toprule\noalign{}
\(p\bmod7\) & \(c_1(p)\) & \(c_2(p)\) & \(c_3(p)\) \\
\midrule\noalign{}
\endhead
\bottomrule\noalign{}
\endlastfoot
2 & \(1/(2p(p-1))\) & \(1/(2(p-1))\) & \(1/(2(p-1))\) \\
3 & \(1/((p-1)(p^3+1))\) & \(1/(2(p-1))\) & \(1/(2p(p-1))\) \\
4 & \(1/(2(p-1))\) & \(1/(2p(p-1))\) & \(1/(2(p-1))\) \\
5 & \(1/(2(p-1))\) & \(1/((p-1)(p^3+1))\) & \(1/(2p(p-1))\) \\
6 & \(1/(p^2-1)\) & \(1/(p^2-1)\) & \(1/(p^2-1)\) \\
\end{longtable}

For \(p\equiv1\pmod7\), this argument provides no positive constant. As before, these constants are lower bounds, not proved optimal finite slopes.

\textbf{(c) Full restricted-operator identities at 11.} In the three-variable Hurwitz algebra over \(\mathbb F_{11}[\zeta_7]\), the actual germs satisfy \[
\boxed{
\partial_1^{11}f=4u_3f,\qquad
\partial_2^{11}f=-\partial_1f,\qquad
\partial_3^{11}f=-u_2f,
\qquad f=F\ \hbox{or}\ \sigma_q.}
\tag{K18.4}
\] These are identities of the complete reduced Hurwitz series. In particular \[
\partial_1^{12}\partial_3^2\sigma_q(0)\equiv8\pmod{11}.
\tag{K18.5}
\] This actual source mixed derivative is a unit at both primes above 11. It disproves the naive additive extension of the three pure-direction constants to all mixed multiindices.

\textbf{(d) A failed next strengthening.} For \(\ell_n=\partial_i^n\log F(0)\), exact arithmetic gives \[
v_{11}(\ell_{2q_i})=1\qquad(i=1,2,3)
\tag{K18.6}
\] at both primes above 11. Thus an attempted improvement based on \(p^2\mid\ell_{2q_i}\) is false. The Bell-lemma equality case at that index is a real obstruction for this particular estimate, although cancellation in the full theta derivative could still give stronger bounds by a different argument.

\subsection{2. The integral extension lemma}\label{the-integral-extension-lemma}

We first give the general local-to-global argument, including its finite-place content.

Let \(R\) be the localization of \(\mathcal O_E\) at a prime above an odd \(p\ne7\), with residue field \(k\), and let \(\mathcal J/R\) be the smooth proper abelian scheme. Completion may be taken after the argument. All rational functions and line charts can thus be taken over the number field, where the complex comparison below applies. The total space is integral and normal, and its special fibre is geometrically integral. Let \(\mathcal L\) be the primitive integral theta line of ALP-LIFT-08, and \(s_q\) its primitive section at this prime. Write \(\Theta\) for its divisor. Its restriction to the special fibre is a nonzero theta divisor; in particular \(\Theta\) has no vertical divisor component.

Take an affine open \(\mathcal U\) on which \(\mathcal L\) is trivial, and write \(f_{\mathcal U}\) for the local equation of \(s_q\). This is a regular function. At the generic point of the special fibre it is a unit, since the section does not vanish on that entire fibre. The invariant derivatives \(D_i\) preserve regular functions on \(\mathcal U\), so \[
D_iD_j\log f_{\mathcal U}
=\frac{D_iD_jf_{\mathcal U}}{f_{\mathcal U}}
 -\frac{D_if_{\mathcal U}D_jf_{\mathcal U}}{f_{\mathcal U}^2}
\tag{K18.7}
\] is a rational function with no vertical pole.

The algebraic \(W_{ij}\) has no vertical pole either. Its pullback to the ordered-divisor chart \(\mathcal R\) of (K12.12) is explicitly integral by (K12.8)--(K12.10). This chart maps étale to \(\mathcal J\), and its reduction has a nonempty open image in the irreducible special fibre. Thus it detects the generic vertical valuation. An étale pullback cannot turn a negative vertical valuation into a nonnegative one. This proves the assertion without any integral Chern frame.

On the characteristic-zero generic fibre, define \[
\beta_{ij,\mathcal U}=W_{ij}+D_iD_j\log f_{\mathcal U}.
\tag{K18.8}
\] It is regular there. To justify the cancellation explicitly, in any analytic lift of a small open set write the actual reduced sigma as \(\sigma_q=f_{\mathcal U}g\), where \(g\) is the holomorphic, nonvanishing transition from the algebraic local frame to the reduced frame. Since \(W_{ij}=-D_iD_j\log\sigma_q\), the right side of (K18.8) is \(-D_iD_j\log g\), which is holomorphic. The right side is also rational over the generic algebraic fibre, hence has no horizontal divisor pole. Equivalently one may check each horizontal valuation after embedding into \(\mathbb C\).

By (K18.7) and the vertical integrality of \(W_{ij}\), (K18.8) has no vertical divisor pole. Every codimension-one point of the smooth total space is either horizontal or generic in the special fibre. A normal affine domain is the intersection of its height-one valuation rings. Therefore \[
\boxed{\beta_{ij,\mathcal U}\in\Gamma(\mathcal U,\mathcal O).}
\tag{K18.9}
\] This is an integral statement about the logarithmic \textbf{second derivative} of a local frame. It neither asserts an integral reduced frame nor asserts integral divided Taylor coefficients. In particular it is compatible with the supplied THETA3 C11 counterexample.

Now suppose \(\overline D_i^{q}=0\) on the special fibre, where \(q=p^h\). It also vanishes on its function field: the \(q\)-th iterate is a derivation, and the quotient rule extends the zero derivation from any dense affine open. Reducing (K18.8), \[
\overline D_i^{q-1}\overline W_{ij}
=\overline D_i^{q-1}\overline\beta_{ij,\mathcal U}
 -\overline D_i^{q}
       \left(\frac{\overline D_j\overline f_{\mathcal U}}
                    {\overline f_{\mathcal U}}\right).
\tag{K18.10}
\] The last term is zero and the first term is regular on \(\mathcal U_k\). The left side is therefore regular on every open of a cover of \(\mathcal J_k\), hence is a global regular function. Properness and geometric connectedness give \(H^0(\mathcal J_k,\mathcal O)=k\). This proves that it is a constant.

There is a slightly more general version that will be used at 11. If \(\overline D_i^{p}=a\overline D_k\), with \(a\in k\), then \[
\overline D_i^{p-1}\overline W_{ij}-a\overline W_{kj}
\tag{K18.11}
\] is a constant. Lift \(a\) to the DVR arbitrarily and apply (K18.9) to both Hessians. The logarithmic term is \((\overline D_i^{p}-a\overline D_k)
 (\overline D_j\overline f/\overline f)=0\), proving regularity exactly as above.

\subsection{3. CM characters and the improved slopes}\label{cm-characters-and-the-improved-slopes}

The source section has CM character 4: \[
\sigma_q(\zeta u_1,\zeta^2u_2,\zeta^3u_3)=\zeta^4\sigma_q(u).
\tag{K18.12}
\] This was established in AM-K014, and also follows from the invariant theta line and its leading quadratic. Twice differentiating gives \[
a^*W_{ij}=\zeta^{-i-j}W_{ij},\qquad
a^*(D_i^mW_{ij})=\zeta^{-i(m+1)-j}D_i^mW_{ij}.
\tag{K18.13}
\] At \(m=q_i-1\), a nonzero constant could have this character only if \(iq_i+j\equiv0\pmod7\), because \(\zeta\) still has order 7 in every residue field under consideration. This proves (K18.1)--(K18.2).

For the pure Hessian, \(i\ne0\pmod7\), so the exceptional condition is exactly \(q_i\equiv-1\pmod7\). Otherwise \[
D_i^{q_i-1}W_{ii}\in p\mathcal R.
\tag{K18.14}
\] Together with \(D_i^{q_i}\mathcal R\subset p\mathcal R\), this even improves the logarithmic estimate to \[
v_p(D_i^mW_{ii}|_a)\geq
\left\lfloor\frac{m+1}{q_i}\right\rfloor,
\qquad
v_p(\ell_n)\geq\left\lfloor\frac{n-1}{q_i}\right\rfloor\ (n\geq2).
\tag{K18.15}
\] For clarity, if \(m=Aq_i+B\), \(0\leq B<q_i\), the usual operator estimate gives \(p^A\). When \(B=q_i-1\), apply \((D_i^{q_i})^A\) to (K18.14) to obtain \(p^{A+1}\). This proves the displayed shifted floor bound; it is an iterated stable-ring statement.

In particular \(p\mid\ell_{q_i+1}\), the hypothesis of the improved Bell--Legendre bound (K13.7). The nonexceptional columns in the table of §1 use \(1/(2p^{h_i-1}(p-1))\). The exceptional columns retain the generic constant \(1/((p-1)(p^{h_i}+1))\) of (K13.6). These are exactly the listed entries. AM-K015 transfers both bounds to \(\sigma_q\) using the actual integral \([2]\)-ratio, so the conclusion concerns the source line and respects the source's odd level.

The exceptional possibility cannot be discarded merely by symmetry. The finite certificate computes \[
\partial_1^{28}\log F(0)\equiv2\pmod3,
\qquad
\partial_2^{126}\log F(0)\equiv1\pmod5.
\tag{K18.16}
\] Here the respective exponents are \(q_1=3^3=27\) and \(q_2=5^3=125\), both congruent to \(-1\pmod7\). These are exact nonzero constants in their whole cyclotomic residue algebras.

\subsection{4. Exact restricted-operator identities at 11}\label{exact-restricted-operator-identities-at-11}

\subsubsection{The algebraic restricted powers}\label{the-algebraic-restricted-powers}

The CM-label argument gives \(D_1^{11}=D_3^{11}=0\) on \(\mathcal J_k\), and \(D_2^{11}\) is a scalar multiple of \(D_1\). The scalar is \(-1\): the exact integral ODE (K12.7), evaluated at the branch divisor, gives \[
D_2^{11}y_j=-D_1y_j\quad(j=1,2,3),
\tag{K18.17}
\] and \(D_1y_1\) is a unit in the cyclotomic residue algebra. Its coefficient vector in \(1,\zeta,\ldots,\zeta^5\) is \((5,4,6,4,5,0)\); the certificate verifies its inverse. Since the two sides are invariant vector fields in the same one-dimensional labelled space, this one unit evaluation determines the scalar globally. Thus \[
\boxed{D_1^{11}=0,\qquad D_2^{11}=-D_1,
\qquad D_3^{11}=0.}
\tag{K18.18}
\]

\subsubsection{The logarithmic constants}\label{the-logarithmic-constants}

For \(i=1\), condition (K18.2) permits only \(j=3\); for \(i=3\), it permits only \(j=2\). The exact finite calculation, using the particularly simple functions \(W_{13}=x_1x_2x_3/4\) and \(W_{23}=-(x_1x_2+x_1x_3+x_2x_3)/4\), gives \[
-D_1^{10}W_{13}=4,\qquad
-D_3^{10}W_{32}=-1\pmod{11}.
\tag{K18.19}
\] The global constant lemma makes these values valid everywhere on the special fibre, not just at the evaluation divisor.

For the middle operator, (K18.11) shows that \(D_2^{10}W_{2j}+W_{1j}\) is a constant. Its CM character is \(-(22+j)\equiv-(1+j)\pmod7\), which is nonzero for every \(j\in\{1,2,3\}\). Therefore \[
D_2^{10}W_{2j}+W_{1j}=0\pmod{11}.
\tag{K18.20}
\]

Let \(H=\log F\), whose Hurwitz coefficients are integral and whose first and all odd-total-degree derivatives at the origin vanish. Differentiating \(\partial_1^{11}H\) once in direction \(j\) gives the constants (K18.19), so in the full three-variable Hurwitz algebra \[
\partial_1^{11}H=4u_3,\qquad
\partial_3^{11}H=-u_2.
\tag{K18.21}
\] There is no integration ambiguity: a series killed by all three divided-power derivations is constant, and these constants at zero vanish by parity. Similarly (K18.20) and parity give \[
\partial_2^{11}H+\partial_1H=0.
\tag{K18.22}
\]

In characteristic 11, every positive-augmentation element of the multivariate Hurwitz algebra has 11-th power zero. For a divided-power monomial with multiindex \(\alpha\ne0\), the valuation of the scalar in its 11-th power is \(\sum_j s_{11}(\alpha_j)>0\). Frobenius is additive, and each coefficient of an infinite product involves only finitely many terms, so the same assertion holds for a complete series. The Bell-polynomial identity modulo a prime is elementary here: in the polynomial \(B_p(H',\ldots,H^{(p)})\), every term except \((H')^p\) and \(H^{(p)}\) has coefficient divisible by \(p\). Indeed its coefficient has numerator \(p!\), whereas all factorials in its denominator have arguments less than \(p\). Thus \[
\partial_i^{11}F/F=(\partial_iH)^{11}+\partial_i^{11}H
=\partial_i^{11}H.
\tag{K18.23}
\] Equations (K18.21)--(K18.22) prove (K18.4) for \(F\).

\subsubsection{Transfer to the actual trace-one section}\label{transfer-to-the-actual-trace-one-section}

Use the integral algebraic comparison \(\sigma_q(2u)=r(u)F(2u)\) of AM-K015. On reductions of its coefficient functions, the three operators \(D_1^{11}\), \(D_2^{11}+D_1\), and \(D_3^{11}\) kill \(r\), by (K18.18). Each operator is a derivation modulo 11. Apply it to the product identity. The linear terms in (K18.4) scale under \(u\mapsto2u\), while \(2^{11}=2\) in characteristic 11. Dividing by the resulting unit factor and replacing \(2u\) by \(u\) proves exactly (K18.4) for \(\sigma_q\). No integral Chern frame was introduced.

The operator identities also hold for AM-R-006's actual moved source section \(\Sigma_{\rm tr}=c\sigma_q\), since \(c\) is constant; its being an 11-adic unit preserves every unit assertion below.

\subsection{5. Mixed derivatives, finite unit families, and their scope}\label{mixed-derivatives-finite-unit-families-and-their-scope}

Write \(S_\alpha=\partial^\alpha\sigma_q(0)\), and interpret a negative index as zero. Equation (K18.4) is equivalently the all-index recurrence \[
\begin{aligned}
S_{\alpha+11e_1}&=4\alpha_3S_{\alpha-e_3},\\
S_{\alpha+11e_2}&=-S_{\alpha+e_1},\\
S_{\alpha+11e_3}&=-\alpha_2S_{\alpha-e_2}
\end{aligned}\qquad(\bmod11).
\tag{K18.24}
\] The first and third equalities are just the Hurwitz Leibniz rule for multiplying by a coordinate. The second is a constant-coefficient operator equality. Each replacement lowers the total derivative order, so these relations reduce every derivative modulo 11 to a finite base cube with all three indices less than 11. They do not assert that every element of that base cube is known without further calculation.

Taking \(\alpha=e_1+2e_3\), and using \(S_{e_1+e_3}=1\), proves (K18.5). The tempting additive extension of the pure lower constants would instead give \[
v_{11}(S_{(12,0,2)})\geq
\frac{12}{20}+\frac2{20}
-\frac{s_{11}(12)+s_{11}(2)}{10}
=\frac3{10}>0,
\] which is false. Consequently pure-direction gains cannot be assigned additively to a selected mixed harmonic row without additional analysis.

There are explicit further unit witnesses. Since \(\partial_1u_3=0\), iteration of (K18.4), as an identity of the actual formal series, gives \[
\partial_1^{11k}\sigma_q=4^ku_3^k\sigma_q\pmod{11}.
\] Differentiating the leading nonzero jets yields \[
\begin{aligned}
S_{(11k+1,0,k+1)}&=4^k(k+1)! &&(0\leq k\leq9),\\
S_{(11k,2,k)}&=-2\,4^kk! &&(0\leq k\leq10),
\end{aligned}\qquad(\bmod11).
\tag{K18.25}
\] Every displayed value is a unit. In particular \[
S_{(100,0,10)}=8,\quad S_{(110,2,10)}=2,
\quad S_{(0,1100,10)}=8\pmod{11}.
\tag{K18.26}
\] The last equality uses \(\partial_2^{1100}\sigma_q=\partial_1^{100}\sigma_q\), obtained by iterating the middle identity in (K18.4). These are finite unit families and exact all-index mod-11 recurrences; they supply no asymptotic valuation upper bound.

For completeness, the same formal algebra gives \[
\partial_1^{121}\sigma_q=0,\qquad
\partial_3^{121}\sigma_q=0,\qquad
\partial_2^{1331}\sigma_q=0\pmod{11}.
\tag{K18.27}
\] Indeed \(u_j^{11}=11!\,u_j^{[11]}=0\) in the reduced Hurwitz algebra; iterating the first or third relation eleven times proves the first two equalities, and the middle relation then proves the last. Differentiating shows vanishing for every multiindex exceeding one of these bounds in its indicated component. The stronger pure-direction cutoffs 11,121,11 from AM-K015 are also immediate by setting the other variables to zero.

Neither (K18.24) nor (K18.27) is iterated as a divisibility statement modulo powers of 11. Such an iteration would require a stable integral module and a justified operator inclusion into \(11\) times that module. No such stronger inclusion is inferred from these formal congruences.

\subsection{6. The next proposed divisibility is false}\label{the-next-proposed-divisibility-is-false}

To test whether the improved Bell constant could be increased by forcing the equality-case logarithmic derivative to acquire an additional power of 11, the exact ODE was evaluated modulo \(121\). In the coefficient basis \(1,\zeta,\ldots,\zeta^5\), it gives

\begin{longtable}[]{@{}lrl@{}}
\toprule\noalign{}
Direction & \(q_i\) & \(\ell_{2q_i}/11\pmod{11}\) \\
\midrule\noalign{}
\endhead
\bottomrule\noalign{}
\endlastfoot
1 & 11 & \(3+3\zeta+8\zeta^3+5\zeta^4+8\zeta^5\) \\
2 & 121 & \(3+3\zeta+8\zeta^3+5\zeta^4+8\zeta^5\) \\
3 & 11 & \(3+3\zeta+3\zeta^2+3\zeta^3\) \\
\end{longtable}

Each displayed polynomial has gcd 1 with \(\Phi_7\) over \(\mathbb F_{11}\), so it is a unit at both primes above 11. This proves (K18.6). At the Bell index \(j=2q_i\), the lower bound in (K13.7) has \((p-1)v(\ell_j)+s_p(j)-1=10+2-1=11\), precisely the equality value \(j/(2p^{h_i-1})\). The proposed extra divisibility is therefore unavailable. This does not prove optimality of the final theta slope, because Bell terms can cancel.

\subsection{7. Reproducibility and source scope}\label{reproducibility-and-source-scope}

The exact finite certificate is \nolinkurl{jacobian_hessian_pcurvature_certificate.py}, with complete output in \nolinkurl{jacobian_hessian_pcurvature_certificate.json}. It checks the two constants (K18.19), the restricted scalar (K18.17), both exceptions in (K18.16), and all three failures in (K18.6). For modulus \(121\), it computes the initial cyclotomic inverses over \(\mathbb Q[Z]/\Phi_7\) and reduces them only after checking their denominator units; it does not treat \(\mathbb Z/121\mathbb Z\) as a field. It verifies the reduced inverses multiplicatively and records the underlying recurrence SHA256.

The stable integral Abel chart, actual Baker/Chern identification, and the classical Kleinian Hessian identities are fully sourced in AM-K012 and its bibliography. The actual primitive line/section and good Hodge basis are ALP-LIFT-05, ALP-LIFT-07, and ALP-LIFT-08. The normality, character, logarithmic-operator, and transfer arguments needed for the new statements are written in full here. The elementary fact that a normal noetherian domain is the intersection of its height-one valuation rings is the sole normality criterion used; no assertion of reduced-frame integrality is hidden in that step.

This closes a further actual-theta coefficient subgate and exhibits a concrete failure of a naive mixed extension. It does not identify a source-selected harmonic row, prove a matching sharp finite slope, supply the terminal every-\(\varepsilon\) window, or extend the accepted Dirichlet theorem to Kummer RH.

\clearpage

\section{AM-K018 --- exact source canonical gauge for the primitive Jacobian and its ordinary powers}\label{am-k018-exact-source-canonical-gauge-for-the-primitive-jacobian-and-its-ordinary-powers}

Author: \nolinkurl{/root/kummer_extension}, GPT-6 Astra Pro, ALP\_MACHINE, 21 November 2026 UTC. This is an exact comparison of formal gauges, not a claim that a reduced theta frame is integral. The accepted DP47, DP49, DP54, DP58, and DP59 theorems and the accepted source canonical-gauge existence theorem are used without rederiving their established conclusions.

\subsection{1. Statement and scope}\label{statement-and-scope}

Let \(J=\operatorname{Jac}(y^2=x^7-1)\), with its full cyclotomic action \(\mathcal O_E=\mathbb Z[\zeta_7]\), canonical principal polarization, and holomorphic CM type \(\Phi=\{1,2,3\}\). Put \[
w_0=-(a+a^2+a^4),\qquad w_0+w_0^\dagger=1,
\qquad L_0=(1,w_0)^*\mathcal P_\lambda.
\] Use the actual rigidified identification \(L_0=L_{2\infty}^{\rm rig}\) of ALP-LIFT-07 and the normalized section \(s_q\) of ALP-LIFT-08. AM-K012--014 identify its actual Chern-reduced germ as \(\sigma_q\).

At the prime 11, in the actual integral hyperbolic formal block and invariant product frame of DP47/DP49, the following comparison is exact: \[
\boxed{\sigma_q(\log z)=f_q(z)\,U(z)\,B(z)^{-1}.}
\tag{K19.1}
\] Here \(f_q\) is the algebraic section coefficient in the stated integral frame, \(U\) is precisely the normalized integral frame factor in DP47, and \(B\) is its actual primitive arbitrary-unit Poincaré splitting. There is no additional quadratic exponential. More generally every actual section \(s\) of \(L_0^{\otimes\kappa}\), represented in the same frame, has Chern-reduced germ \[
\boxed{\Theta_s(\log z)=f_s(z)U(z)^\kappa B(z)^{-\kappa}.}
\tag{K19.2}
\] The right side is exactly the normalized formal expression in DP54, DP58, and DP59. Equality holds as a complete characteristic-zero formal series, for all \(\kappa\); it is not restricted to small jet degree or small level.

Thus their accepted selected derivative statements apply to the actual Chern/Baker germs, with the same polynomial operator in the same invariant tangent directions. This includes the allowed theta translates in DP59 once their actual section and theta-group scalar are fixed. It does not identify an unspecified external automorphic row or its direction matrix.

The argument also applies to ordinary powers \(J^m\) with diagonal \(\mathcal O_E\)-action and their product polarization, because their holomorphic characters are repetitions of \(1,2,3\). It is \textbf{not} asserted for the whole DP65 family with conjugated \(E\)-actions: such a family may have opposite holomorphic characters across factors and admit nonzero balanced bilinear forms. Extra integral self-adjoint factor idempotents could remove those cross terms, but they have not been assumed here.

\subsection{2. Exact source interfaces actually inspected}\label{exact-source-interfaces-actually-inspected}

The following source files were recovered from the intake-pinned mirror and matched against the original-author archive SHA256 manifest. No Lean compilation or rederivation of the accepted formal theorems is claimed.

\begin{longtable}[]{@{}
  >{\raggedright\arraybackslash}p{(\columnwidth - 2\tabcolsep) * \real{0.5000}}
  >{\raggedright\arraybackslash}p{(\columnwidth - 2\tabcolsep) * \real{0.5000}}@{}}
\toprule\noalign{}
\begin{minipage}[b]{\linewidth}\raggedright
Source
\end{minipage} & \begin{minipage}[b]{\linewidth}\raggedright
Exact interface used
\end{minipage} \\
\midrule\noalign{}
\endhead
\bottomrule\noalign{}
\endlastfoot
\nolinkurl{ReciprocityTheory/ReductiveGroups/ArithmeticGeometry/BirationalGeometry/ToroidalCompactification/CocycleRightZeroThd.lean} & \nolinkurl{FB} at lines 86--158; every clause of \nolinkurl{FB.IsCanonGauge} at 174--183; \nolinkurl{hL} at 185; \nolinkurl{nabVal} at 187--189. \\
\nolinkurl{Development/AbelianVarietyTheory/Proofs/FormalGroupCharts/CanonGaugeSubIsBiadd.lean} & Complete theorem \nolinkurl{canonGauge_sub_isBiadd}, lines 16--60: differences are axis-normalized, additive in both arguments, and adjoint-balanced. \\
\nolinkurl{Development/AbelianVarietyTheory/Proofs/FormalGroupCharts/CanonGaugeIndependenceAll.lean} & \nolinkurl{gaugeExp_holds}, lines 19--21; \nolinkurl{canonGauge_independence_all}, lines 23--26. We do not replace its selected-operator hypotheses by unconditional gauge independence in a different CM type. \\
\nolinkurl{ReciprocityTheory/ReductiveGroups/ArithmeticGeometry/TaylorDiscriminants/WickVanishing/CanonB.lean} & Accepted theorem \nolinkurl{canonGauge : CanonGaugeSt}, line 334, supplies a canonical gauge for a finite free reduced integral algebra with involution; its opening construction and theorem signature were inspected. \\
\nolinkurl{ReciprocityTheory/ReductiveGroups/ArithmeticGeometry/TaylorDiscriminants/ConnectionKillsSeries.lean} & \nolinkurl{CanonGaugeSt}, lines 56--59, quantifies over the full finite free reduced algebra \(O\), with no level or jet-order bound; \nolinkurl{ConnectionIntegrality}, 77--83, is a separate supplied integral-connection statement, not integral reduced-frame coefficients. \\
\nolinkurl{Formalism/LocalToGlobal/ProjectiveLattices/ThGerm.lean} & Full \nolinkurl{IsFrame} clauses, 119--149; actual primitive diagonal frame \nolinkurl{eL}, 166--170; \nolinkurl{IsGerm}, 172--181; \nolinkurl{ChartAt}, 189--206; \nolinkurl{IsNab}, 208--217. In particular the source action algebra in \nolinkurl{ChartAt} is the full \nolinkurl{𝓞 E}, not merely the quadratic field generated by \(w_0\). \\
\end{longtable}

The gauge axioms, in ordinary notation and in the \textbf{same frame}, are \[
\begin{aligned}
h(0,Y)=h(X,0)&=0,\\
\mathrm{ph}_1&=\exp(h(X,Z)+h(Y,Z)-h(X+_FY,Z)),\\
\mathrm{ph}_2&=\exp(h(X,Y)+h(X,Z)-h(X,Y+_FZ)),\\
\mathrm{sy}&=\exp(h(Y,X)-h(X,Y)),\\
\mathrm{ad}(b)&=\exp(h([b]X,Y)-h(X,[b^\dagger]Y))
\quad(b\in\mathcal O_E).
\end{aligned}
\tag{K19.3}
\] The source normalized section is \(f\exp(\kappa h_L)\), where \(h_L(X)=h(X,[w_0]X)\). The signs in (K19.3), and hence in (K19.1), are fixed by these actual source clauses.

\subsection{3. A Chern-reduced biextension frame satisfies all those axioms}\label{a-chern-reduced-biextension-frame-satisfies-all-those-axioms}

Let \(V/\Lambda\) be the complex torus \(J(\mathbb C)\). In the first-argument-linear convention of AM-DO-011 and AM-R-004, its positive Hermitian form is \(H\). A reduced theta frame for a symmetric principal line has multiplier \[
j_L(\lambda,u)=\alpha(\lambda)
 \exp\left(\pi H(u,\lambda)+\frac\pi2H(\lambda,\lambda)\right),
\tag{K19.4}
\] where \(\alpha\) is its semicharacter and \(\operatorname{Im}H\) is integral on the lattice. Formula (K19.4) is the explicit multiplier already proved in AM-DO-011 from the theta series; equivalently it is Bannai--Kobayashi's Definition 1.8.

The symmetric polarization biextension is \(m^*L\otimes\operatorname{pr}_1^*L^{-1}
\otimes\operatorname{pr}_2^*L^{-1}\), with axial rigidifications. Its induced reduced frame on \(V\times V\) has multiplier \[
j_P((\lambda,\mu),(u,v))
=\exp\left(\pi\{H(u,\mu)+H(v,\lambda)+H(\lambda,\mu)\}\right).
\tag{K19.5}
\] Indeed the quotient of the three multipliers (K19.4) gives the two cross linear terms and \(\operatorname{Re}H(\lambda,\mu)\). The semicharacter quotient is \(\exp(\pi i\operatorname{Im}H(\lambda,\mu))\), giving the last term as written.

Both partial addition maps have scalar one in this reduced frame. This follows by adding the corresponding terms in (K19.5): the multiplier for \((u+u',v)\) is exactly the product of those for \((u,v)\) and \((u',v)\), with the shared lattice shift in the second variable, and likewise with the roles reversed. Its axial values are one.

The symmetry map also has scalar one. Swapping \((u,v)\) and \((\lambda,\mu)\) changes the exponent by \(\pi(H(\lambda,\mu)-H(\mu,\lambda))
=2\pi i\operatorname{Im}H(\lambda,\mu)\), whose exponential is one. The map is normalized on the axes.

Finally the polarization adjoint relation is \[
H(bu,v)=H(u,b^\dagger v)\quad(b\in E).
\tag{K19.6}
\] Inserting it in (K19.5) shows that the multipliers of \(P(bu,v)\) and \(P(u,b^\dagger v)\) coincide. Thus their canonical normalized adjunction has scalar one in the reduced frames. These maps realize the actual rigidified Poincaré addition, symmetry, and adjunction maps; uniqueness can also be seen by dividing two such normalized geometric maps, which gives a regular unit on the proper base and hence a constant fixed to one on the axes.

Choose an algebraic formal biextension frame \(e\), normalized on the axes, and let \(e_{\rm red}\) be the reduced frame just constructed. Write \[
e=\exp(h_C)e_{\rm red}.
\tag{K19.7}
\] Its formal logarithm has zero axial values. Computing the frame ratios for the four kinds of structure map gives exactly (K19.3). For example, the source \nolinkurl{IsFrame.ph1} compares its rigidified cubic generator with \(\mathrm{ph}_1 e(X+_FY,Z)e(X,Z)^{-1}e(Y,Z)^{-1}\), giving the negative addition coboundary displayed in (K19.3). Its \nolinkurl{sy} uses \(e(Y,X)/e(X,Y)\), and its \nolinkurl{ad} uses \(e([b]X,Y)/e(X,[b^\dagger]Y)\), giving the two remaining signs. Therefore \(h_C\) is a full-\(\mathcal O_E\) source canonical gauge, not only a candidate quadratic completion of a single section.

\subsection{4. Uniqueness for this full CM type and algebraic descent}\label{uniqueness-for-this-full-cm-type-and-algebraic-descent}

Let \(h,h'\) be two source canonical gauges in the \textbf{same formal frame}, for the full cyclotomic algebra. By the accepted source difference theorem, their difference is biadditive and adjoint-balanced. In normalized group logarithms \(U,V\), biadditivity forces a bilinear form \[
L(U,V)=\sum_{i,j=1}^3c_{ij}U_iV_j.
\tag{K19.8}
\] For completeness, additivity gives \(L(nU,V)=nL(U,V)\) for every positive integer \(n\); comparing homogeneous degrees over characteristic zero leaves first-variable degree one. Applying the same argument in the second variable gives (K19.8).

Take \(b=a\), the order-seven automorphism. Since \(a^\dagger=a^{-1}\), the balance equation gives \[
c_{ij}(\zeta^i-\zeta^{-j})=0.
\tag{K19.9}
\] For \(i,j\in\{1,2,3\}\), one has \(i+j\in\{2,3,4,5,6\}\), so every scalar in parentheses is nonzero. Thus \[
\boxed{h=h'.}
\tag{K19.10}
\] This is stronger than saying merely that an invariant quadratic diagonal vanishes. It does not require the algebraic frame itself to be invariant: the same-frame adjoint cocycles cancel in the difference.

The accepted source theorem \nolinkurl{CanonB.canonGauge} supplies a gauge over the characteristic-zero algebraic coefficient field for the actual algebraic formal data. Its hypotheses apply to \(\mathbb Z[\zeta_7]\), a reduced finite free \(\mathbb Z\)-algebra with conjugation. Compare that algebraic gauge, after complex base change, with \(h_C\). Equation (K19.10) proves equality and hence descent of \(h_C\)'s formal coefficients to the same algebraic field. This avoids treating a potentially nonconvergent \(p\)-adic formal frame as an analytic frame.

Now pass to the 11-adic field. If a chosen integral formal frame \(e'\) differs from an algebraic reference frame by a normalized formal unit \(J\), set \[
h_C'=h_C+\log J.
\tag{K19.11}
\] The logarithm is a characteristic-zero formal logarithm, with no assertion of integral coefficients. Frame substitution in (K19.3) shows that \(h_C'\) is canonical for the new cocycles. Thus uniqueness applies in the precise integral product frame of DP47. There is no illicit comparison of gauges belonging to different frame cocycles.

For \(J^m\) with its ordinary diagonal action, equation (K19.9) applies to every pair of variables, whatever their factor labels, and still forces zero. If different factors have conjugated \(E\)-actions, opposite labels can occur and the coefficient in (K19.9) can vanish. No uniqueness claim for that mixed situation is made.

\subsection{5. The trace-one diagonal and the DP47 primitive kernel}\label{the-trace-one-diagonal-and-the-dp47-primitive-kernel}

Pull the reduced biextension frame back along \(z\mapsto(z,w_0z)\). Its Hermitian form is \[
\begin{aligned}
H_P((z,w_0z),(w,w_0w))
&=H(z,w_0w)+H(w_0z,w)\\
&=H(z,(w_0+w_0^\dagger)w)=H(z,w).
\end{aligned}
\tag{K19.12}
\] The line is the actual \(L_0\), and its rigidified symmetric representative is \(L_{2\infty}^{\rm rig}\) by ALP-LIFT-07. Thus this pulled-back frame is precisely the normalized Chern-reduced frame used for \(\sigma_q\). There is no remaining linear or quadratic exponential: both reduced frames have the same multiplier, and the ratio of the normalized global line identifications is a holomorphic unit on the compact torus, hence the constant one.

This also checks the source primitive diagonal frame. Its \nolinkurl{ThGerm.eL} includes the rigidified primitive cube isomorphism \nolinkurl{wP}, not just the biextension frame. Under (K19.12), that actual isomorphism has constant one in reduced frames, fixed by \nolinkurl{IsPrim} and \nolinkurl{rig0}. Consequently the source coefficient \(f\exp(\kappa h_L)\) is exactly the Chern-reduced section coordinate.

DP49 supplies the actual polarized integral block \(G\times H\) at 11, with the \textbf{full cyclotomic action}, the principal Cartier duality, and the arbitrary-unit display models. In DP47 the invariant product gauge \(h_0\) is shown to satisfy the source axis, addition, symmetry, and adjunction identities for the commutative CM action preserving this decomposition. The full \(\mathcal O_E\)-action preserves the two integral idempotents; all its scalar endomorphisms on the one-dimensional formal factor commute with the Teichmüller action, and the transpose action is its actual adjoint on the other factor. Hence this is a full-\(\mathcal O_E\) canonical gauge, not merely a gauge for the smaller quadratic order generated by \(w_0\).

Compare \(h_0\) with \(h_C'\) in that same product frame. Equation (K19.10) gives \(h_0=h_C'\). The accepted DP47 identity \[
\exp h_{0,L}=U B^{-1}
\tag{K19.13}
\] therefore proves (K19.1)--(K19.2) with its exact sign and its actual integral unit \(U\). Both sides can be nonintegral as complete ordinary power series. This comparison does not require or prove an integral reduced theta frame.

The comparison contains no level parameter. It applies to every integer tensor power, and to the actual odd levels \(M=29^e\) of DP59 with \(\kappa=M^2\). At such a level the prescribed factor \(M^{-n}\) cancels exactly the change in an \(n\)-th derivative under \(u\mapsto Mu\), as already proved in AM-K015. There is no small-degree or small-\(M\) restriction concealed in the canonical-gauge argument.

\subsection{6. Fixed actual-base sharp slopes on two explicit mixed lines}\label{fixed-actual-base-sharp-slopes-on-two-explicit-mixed-lines}

The new gauge comparison permits a direct application of the accepted DP49 theorem to the single fixed section \(s_q\), with explicit lines. In the original \(u\)-coordinates take \[
f_{23}(t)=\sigma_q(0,t,t),\qquad
f_{13}(t)=\sigma_q(t,0,t).
\tag{K19.14}
\] Their leading terms are respectively \(-t^2\) and \(t^2\), by ALP-LIFT-08 or AM-K014. Thus the algebraic prefactor has nonzero restriction on these precise lines: multiplying by any normalized gauge exponential or integral frame unit cannot remove that leading term.

DP49's unit-arrow logarithmic coordinates \((X,Y_1,Y_2)\) have CM characters \((3,2,1)\). The two invariant bases are related by diagonal units at 11, since each is an integral eigenbasis with the same one-dimensional characters. Therefore the first line has unit \(G\)-component and unit first dual component, of restricted-power length two. The second has unit \(G\)-component, zero first dual component, and unit last dual component, of length one. All hypotheses of the fixed-line clauses of DP49 are now verified, including the nonzero prefactor restriction.

Consequently \[
\begin{aligned}
v_{11}(f_{23}^{(n)}(0))
 &\geq \frac n{1220}-\frac{s_{11}(n)}{10},
&\liminf_{n\to\infty}\frac{v_{11}(f_{23}^{(n)}(0))}{n}
 &=\frac1{1220},\\
v_{11}(f_{13}^{(n)}(0))
 &\geq \frac n{120}-\frac{s_{11}(n)}{10},
&\liminf_{n\to\infty}\frac{v_{11}(f_{13}^{(n)}(0))}{n}
 &=\frac1{120}.
\end{aligned}
\tag{K19.15}
\] Zeros are assigned valuation \(+\infty\) in the liminf. These are matching lower/asymptotic upper statements for actual derivatives of a single explicitly normalized section along two fixed mixed lines. The exact derivatives are respectively \((\partial_2+\partial_3)^n\sigma_q(0)\) and \((\partial_1+\partial_3)^n\sigma_q(0)\). They are distinct from individual mixed monomial jets and from the three pure-direction bounds of AM-K013--017.

As a consistency check, AM-K017's complete mod-11 identities give \(f_{13}^{(11)}=4t f_{13}\pmod{11}\); the shortest primitive mixed logarithmic obstruction indeed occurs in degree 12. The 23 direction has algebraic restricted-power length two, consistent with the first DP49 class in the displayed table. These congruences are not used to infer the liminf: the latter is the accepted DP49 conclusion after the now-verified gauge and nonvanishing premises.

For every admitted odd level, the normalized pullbacks of the same section retain both slopes by the exact \(M^{-n}\) cancellation. DP59's prescribed translated span likewise lies in this same actual gauge; its accepted ordinary and first-row harmonic jet witnesses therefore have their stated valuations as Chern-reduced derivatives of those actual sections. Their identification with a particular later automorphic row, or with the mixed-conjugate factors of DP65, remains a separate issue.

\subsection{7. What has been used and what remains}\label{what-has-been-used-and-what-remains}

The source records use full \(\mathcal O_E\)-adjunctions and their generic canonical gauge is supplied for all formal orders. The geometric full-CM structure on this Jacobian and the actual product kernel at 11 are the accepted DP49/DP47 inputs. The new comparison proves that their canonical gauge is the specific Chern/Baker gauge already used for the explicit periods, Gaussian bounds, determinant coefficients, and K017 recurrences. Thus these calculations now concern the same actual formal normalization.

This closes the gauge comparison for ordinary primitive-J powers and the DP59 section family. It does not prove that all those sections occur in a required iterated automorphic image; it does not compare an unknown selected direction matrix or minor; it does not remove all integral lattices, projection costs, or the terminal every-\(\varepsilon\) window. In particular the no-opposite-characters uniqueness proof must not be applied unchanged to the conjugated \(E\)-actions in DP65.

Primary literature for the reduced multiplier is Bannai, K., and Kobayashi, S. (2010), \emph{Algebraic theta functions and the p-adic interpolation of Eisenstein--Kronecker numbers}, Duke Mathematical Journal, 153(2), 229--295, Definition 1.8 in the author preprint \href{https://arxiv.org/pdf/math/0610163}{arXiv:math/0610163v4}; its concrete multiplier was independently computed in AM-DO-011. Exact accepted-result locators are the supplied \nolinkurl{research/DP47_integrated_fragment.tex}, \nolinkurl{DP49_integrated_fragment.tex}, \nolinkurl{DP54_integrated_fragment.tex}, \nolinkurl{DP58_integrated_fragment.tex}, and \nolinkurl{DP59_integrated_fragment.tex} inside \nolinkurl{RETURN_DP_MULT_checkpoint_06_20261121T0420Z}. The source retrieval records are \nolinkurl{../research/source_recovery_20261121T212734Z.json}, \nolinkurl{source_recovery_20261121T213026Z.json}, and \nolinkurl{source_recovery_20261121T213717Z.json}; each compared file matched its intake author hash.

\clearpage

\section{ALP-LIFT-10: the permitted half-torsion orbit in an actual first theta image}\label{alp-lift-10-the-permitted-half-torsion-orbit-in-an-actual-first-theta-image}

Run ALP\_MACHINE; /root/lift\_accounting, GPT-6 Astra Pro; inherited effort not separately exposed. This result uses accepted DP59 and DP62 for the algebraic permitted sections and their canonical normalizations. It does not re-prove their integral span or jet separation. It uses ALP-LIFT-09 for the actual global unitary splitting, first-lift integral, and torus class calculation, and AM-DO-020--021 for the fixed split and ramified compact phases.

\subsection{Statement}\label{statement-10}

Let \(J\), \(L\), \(s_q\), \(\sigma_q\), \(\Lambda\), and \(H\) have the normalizations of ALP-LIFT-08--09. Thus \(\sigma_q\) is the actual Chern-reduced germ of \(s_q\) with leading term \(u_1u_3-u_2^2\), and \(s_q\) is primitive at every odd prime other than7. Put \[
 \delta_7^2=-7,\qquad
 \alpha=(1+2\delta_7)^e,\qquad M=29^e,\qquad \kappa=M^2,
 \qquad e\ge1.
 \tag{O1}
\] Here \(\delta_7=1+2(\zeta+\zeta^2+\zeta^4)\), so the choice is exactly DP59's actual level. For every permitted root \(t'\in J[\alpha^2]\), let \(s_{t'}\) be the canonical symmetric theta translate of \([M]^*s_q\) from DP59.

Fix an even positive integer \(n\), and use the actual compact pair \(T=U(1)(E/F)\) and \(G=U(I_n)(E/F)\), with the honest splitting \(\chi_V=1,\chi_W=\mu_E\) of ALP-LIFT-09. Let \(P\) be any polynomial differential operator in the \(3n\) algebraic differential coordinates, homogeneous of total degree \(N\) and of a fixed degree at each of the three CM embeddings. All such polynomials have one fixed archimedean torus weight, denoted \(\beta\); it includes the actual vacuum character.

There is one explicitly defined automorphic input \(f_{M,\beta}\) on \(T\), independent of the allowed tuple \(\mathbf t'=(t'_1,\ldots,t'_n)\) and of \(P\) inside this weight, such that an actual theta lift has identity value \[
 \Theta(\phi_{\mathbf t',P},f_{M,\beta})(1)
  = C^n M^{-N}
       P(\partial_u)
         \left(\prod_{j=1}^n
                    \sigma_{s_{t'_j}}(u_j)\right)_{u=0},
 \qquad C=\theta[0;0](0;\tau)c.
 \tag{O2}
\] The notation \(\sigma_{s_{t'}}\) denotes the specified Chern-reduced germ of the algebraic section \(s_{t'}\), in the line of curvature \(\kappa H\). The constant \(c\ne0\) is fixed by \(\Sigma_{\rm tr}=c\sigma_q\); it is independent of \(M\) and the tuple. The finite Schwartz vector in (O2) is an explicitly specified Heisenberg translate of the fixed vacuum, with coefficients roots of unity and finite \(L^2\) norm one.

The automorphic input is supported on one finite component and has the exact norms \[
 \|f_{M,\beta}\|_2^2=h_M
       =56^2\,29^{3(e-1)},\qquad
 \|f_{M,\beta}\|_\infty=h_M.
 \tag{O3}
\] Consequently its normalization to norm one multiplies the right side of (O2) by \(h_M^{-1/2}\) and has logarithmic cost \[
 \log\sqrt{h_M}
      =\frac32\log M+\log56-\frac32\log29.
 \tag{O4}
\] Every prime divisor of \(h_M\) lies in \(\{2,7,29\}\), so both the projection and norm normalization have zero valuation cost at11, and at every prime outside this displayed exceptional set. At the DP59 levels \(M=O(N)\), the new logarithmic norm cost is \(O(\log N)\).

This is a simultaneous realization of the specified fixed-weight jet functionals by actual first theta lifts. It is not an identification of the entire section module with an automorphic representation, an assertion about one freely prescribed irreducible torus character, or a construction of the later definite lifts.

\subsection{1. Exact multiplication and translation of the actual reduced section}\label{exact-multiplication-and-translation-of-the-actual-reduced-section}

Write \[
 (U_x^H F)(v)
       =\exp\bigl(-\pi H(v,x)-\pi H(x,x)/2\bigr)F(v+x).
 \tag{O5}
\] This is the normalized Mumford translation used in AM-DO-017--018. The rigidified multiplication isomorphism \([M]^*L\simeq L^{M^2}\) carries the reduced section to \(\sigma_q(Mu)\) exactly. Indeed the pullback Hermitian form is \(M^2H\), its semicharacter is \(\alpha_{\rm tr}(M\lambda)=\alpha_{\rm tr}(\lambda)^{M^2}\), and the two normalized Appell--Humbert presentations have the same multiplier and the same origin rigidification. Their isomorphism is therefore the identity in these presentations; no extra quadratic exponential appears.

Choose a rational period representative \(r\in\kappa^{-1}\Lambda\) of \(t'\), and put \(x=Mr\). The equality of the two translated functions is immediate from (O5): \[
 U_r^{M^2H}\bigl(\sigma_q(Mu)\bigr)
                 =(U_x^H\sigma_q)(Mu).
 \tag{O6}
\] Both linear exponential terms are \(-\pi M^2H(u,r)\), both constant exponential terms are \(-\pi M^2H(r,r)/2\), and both arguments are \(Mu+Mr\).

The canonical symmetric lift in DP59 differs from \(U_r^{M^2H}\) by an explicitly determined sign only. To see this, let \(\iota\) be normalized inversion in the line presentation. Then \[
 \iota U_r^{M^2H}\iota^{-1}=U_{-r}^{M^2H}
                       =(U_r^{M^2H})^{-1}.
 \tag{O7}
\] Since the alternating self-pairing is zero, \((U_r^{M^2H})^\kappa=U_{\kappa r}^{M^2H}\), and the latter is the scalar \(\alpha_{\rm tr}(\kappa r)^\kappa
 =\alpha_{\rm tr}(\kappa r)\in\{\pm1\}\). Because \(\kappa\) is odd, multiplying \(U_r^{M^2H}\) by \[
                        \epsilon_r=\alpha_{\rm tr}(\kappa r)
 \tag{O8}
\] gives an operator whose \(\kappa\)-th power is one and which still satisfies (O7). The uniqueness of the symmetric finite-order lift, already proved in DP59 §4, identifies it with the permitted lift. In particular, \[
                 \sigma_{s_{t'}}(u)
                    =\epsilon_r(U_x^H\sigma_q)(Mu).
 \tag{O9}
\] Changing the representative gives the same section by that uniqueness. All phases introduced in (O8) are algebraic units.

The point represented by \(x\) lies in the smaller subgroup \[
                        [M]J[\alpha^2]=J[\alpha]\subset J[M].
 \tag{O10}
\] The inclusion follows from \(\alpha M=\alpha^2\bar\alpha\) and \(\alpha\bar\alpha=M\). Equality follows from the accepted DP59 kernel calculation: the kernel of \([M]\) on \(J[\alpha^2]\) is \(J[\alpha]\), so the image has the same order as \(J[\alpha]\). This is where the actual permitted half-orbit is used; replacing it by all of \(J[M]\) would lose the sharper level below.

Taking tensor products of (O9), and using homogeneity of \(P\), gives the all-degree identity \[
 M^{-N}P(\partial_u)\prod_j\sigma_{s_{t'_j}}(u_j)\big|_0
   =\left(\prod_j\epsilon_{r_j}\right)
            P(\partial_v)\prod_j(U_{x_j}^H\sigma_q)(v_j)\big|_0.
 \tag{O11}
\] The factor \(M^{-N}\) from the source therefore cancels the multiplication pullback exactly.

\subsection{2. The exact finite Schwartz family}\label{the-exact-finite-schwartz-family}

Let \(\phi_{0,f}\) be ALP-LIFT-09 (V12). For the selected representatives define \[
 \phi_{\mathbf t',f}
       =\bigotimes_{j=1}^n\bigl(\epsilon_{r_j}W_{x_j}\phi_{0,f}\bigr).
 \tag{O12}
\] The \(W_x\) are the explicit finite operators (V13). Their translates and modulation factors preserve compact support and local constancy. Their nonzero values have absolute value one and are roots of unity, because all period coefficients are rational. They are unitary, so (O12) has norm one.

The exact reading identity of AM-DO-018 and \(R\phi_{0,f}=C\sigma_q\) give \[
 R\phi_{\mathbf t',f}(v)
       =C^n\left(\prod_j\epsilon_{r_j}\right)
                      \prod_j(U_{x_j}^H\sigma_q)(v_j).
 \tag{O13}
\] There is one common period \(C^n\) for the whole family. Use the actual archimedean Schwartz vector \[
                   \phi_{P,\infty}
                         =P(r_u)\varphi_{0,\infty}^{\otimes n},
 \tag{O14}
\] with the creation operators of ALP-LIFT-09 (V19)--(V21). Its actual torus character is \(\beta\). Formula (V20) identifies the theta sum at the identity with the derivative of (O13). Thus the identity value on the right side of (O2) is already an actual theta series. It remains to realize it by one common automorphic input with its normalization controlled.

\subsection{3. The fixed compact vacuum character for even rank}\label{the-fixed-compact-vacuum-character-for-even-rank}

For the fixed splitting \(\chi_V=1\) and even \(n\), the scalar norm-one compact character of the base vacuum is \[
                  \lambda_f(k)=\epsilon_7(k_7)^{n/2},
 \tag{O15}
\] where \(\epsilon_7(t)\) is the residue of \(t\) in \(\{1,-1\}\) at the unique ramified prime over7. All other finite local factors are trivial.

Here is the precise justification of the unramified part, supplementing AM-DO-020--021's calculations. At every odd unramified quadratic place, the norm on integral units is surjective. For example the residue norm is surjective, and successive correction in the principal units uses the surjective residue trace. Choose a unit of norm \(-1\). It gives an integral primitive isotropic vector of \(I_2\); an integral pairing partner can be adjusted, using that 2 is a unit, to make an integral hyperbolic pair. Therefore \(I_n\) is an integral orthogonal sum of hyperbolic planes. Pan (2001), Lemma 6.2, printed p.197, gives scalar \(\alpha(t)\beta^Y(t)=\chi_V((t-1)\delta)=1\) against each such plane, with the identity element treated directly. His tensor identity (3.6.e) gives one for their sum. The selfdual trace-lattice normalization is the same lattice-model normalization identified in AM-DO-021 §1.

At relative split places, including2, AM-DO-020 gives the exact split determinant-root scalar. It is one on the compact torus when \(\chi_V=1\). The only remaining ramified place is7, and AM-DO-021 (DO21.12) gives exactly (O15). The simultaneous sign transport in ALP-LIFT-09 does not change this residual real sign. Thus the fixed compact subgroup \[
                  K_7^+=\ker\epsilon_7\subset K_{0,7}
 \tag{O16}
\] has index2 and fixes the vacuum for every even \(n\). When \(4\mid n\) this restriction is redundant, but retaining it gives one uniform exact level and cost for all even ranks.

\subsection{\texorpdfstring{4. The permitted half-orbit is fixed at level \(M\), not \(M^2\)}{4. The permitted half-orbit is fixed at level M, not M\^{}2}}\label{the-permitted-half-orbit-is-fixed-at-level-m-not-m2}

The prime29 splits completely in \(E/\mathbf Q\). There are three places \(v\) of \(F\) above29, each with \(F_v=\mathbf Q_{29}\) and \(E_v=F_v\times F_v\). Since the ideals \((\alpha)\) and \((\bar\alpha)\) are coprime and their product is \((M)\), the two valuations of \(\alpha\) in each relative pair are \((e,0)\) in one order.

The actual lattice decomposes at such a place into its two relative split isotropic summands, \[
                      \Lambda_v=X_v\oplus Y_v.
 \tag{O17}
\] Both are stable under \(w_0\) and the scalar norm-one torus. A point of \(J[\alpha]\) can be represented modulo \(\Lambda_v\) entirely in the summand on which \(\alpha\) has valuation \(e\), with its coordinate in \(29^{-e}X_v\) (or in \(29^{-e}Y_v\) in the other order).

Replacing a representative by an integral vector changes \(W_x\phi_{0,f}\) by a fixed root of unity: this follows from the Heisenberg multiplication law and the lattice vacuum eigenrelation. Such a scalar does not affect its stabilizer. It is therefore legitimate to use the pure isotropic representative at each place when computing the common stabilizer of the actual vectors (O12).

Take \[
 K_v(M)=\{t=(a,a^{-1})\in K_{0,v}:
                                  a\equiv1\pmod{29^e}\}.
 \tag{O18}
\] For \(x\in29^{-e}X_v\) and \(t\in K_v(M)\), put \(d=(t-1)x\). Then \(d\in X_v\subset\Lambda_v\). Both \(x\) and \(d\) lie in the same isotropic space, so \[
                  \mathfrak b(x,d)=0,\qquad
                  \mathfrak b(d,w_0d)=0.
 \tag{O19}
\] The lattice eigencharacter of \(d\) is consequently one, and the half-pairing in the multiplication law is zero. Thus \[
                         W_{tx}\phi_{0,f}=W_x\phi_{0,f}.
 \tag{O20}
\] The base compact scalar is one by (O15). Heisenberg covariance therefore proves that \(K_v(M)\) fixes every allowed translated vacuum exactly. The identical argument applies when the denominator lies in \(Y_v\).

At every finite prime not dividing29, the vector \(x\) is integral in the trace lattice, so \(W_x\phi_{0,f}\) is already a scalar multiple of the base vacuum. Every compact element fixing the base vacuum fixes that translate. Consequently the explicit product subgroup \[
 K_M=K_7^+\times\prod_{v\mid29}K_v(M)
             \times\prod_{v\nmid7\cdot29}K_{0,v}
 \tag{O21}
\] fixes every finite vector (O12), and their entire span, for every allowed tuple.

\subsection{5. Exact finite index and automorphic input}\label{exact-finite-index-and-automorphic-input}

The three factors at29 have indices \[
 [K_{0,v}:K_v(M)]
       =\#(\mathbf Z/29^e\mathbf Z)^\times
       =28\,29^{e-1}.
 \tag{O22}
\] The index at7 is2. Reduction at29 is injective on \(\mu_{14}\), since its order is prime to29, so \(K_M\cap\mu_{14}=1\). No extra auxiliary prime is required. ALP-LIFT-09 (V6) now gives \[
 h_M=\frac{2(28\,29^{e-1})^3}{14}
       =56^2\,29^{3(e-1)}.
 \tag{O23}
\] In particular \(v_{11}(h_M)=0\).

Define \(f_{M,\beta}\) by ALP-LIFT-09 (V33), with \(K=K_M\): it equals \(h_M\beta(t_\infty)\) on the identity finite component and zero on all other finite components. The definition is unambiguous because \(K_M\) kills the rational root-unit stabilizer. That component has measure \(1/h_M\). Consequently \[
 \Theta(\phi_{\mathbf t',f}\otimes\phi_{P,\infty},
                             f_{M,\beta})(g)
       =\vartheta_{\mathbf t',P}(g,1).
 \tag{O24}
\] This is a genuine automorphic theta integral for the fixed global splitting. Equations (O11), (O13), and (O14) give (O2) on setting \(g=1\). The elementary measure calculation in (V34) gives (O3), and taking its square root gives (O4). The same input works simultaneously for every allowed tuple and every \(P\) in the specified torus weight.

The input may be decomposed into \(h_M\) automorphic characters by (V36)--(V38). Its finite coefficients are \(h_M\) times roots of unity. Individual-character projections divide by \(h_M\); the simultaneous input (O24) has no such denominator. Even after normalization to norm one, \(h_M^{-1/2}\) is an 11-adic unit after a finite coefficient-field extension. Thus the algebraic value in (O2), after removing the retained common period \(C^n\), has exactly the same 11-adic valuation with the unit-norm input as it had before that normalization. This assertion is about the displayed value; it does not assert primitivity of every target row or compute the target mass norm.

For \(n=2\) this applies directly to pairs of DP59's ordinary allowed sections, including any homogeneous two-column differential chosen on ordinary powers of \(J\). Tensoring a fixed number of ordinary factors does not change the torus level index (O23).

\subsection{6. Extracting one norm-one character without losing the local valuation}\label{extracting-one-norm-one-character-without-losing-the-local-valuation}

The simultaneous input is useful for realizing a whole specified finite span. For an existence statement one may instead extract one of its irreducible-character components with norm one. The finite torus \(K_0\) preserves \(J[\alpha]\), since its scalar integral action commutes with \(\alpha\). Therefore, for every \(k\in K_0\), Heisenberg covariance gives \[
 \omega_f(k)\phi_{\mathbf t',f}
       =\lambda_f(k)\bigotimes_j
                    \bigl(\epsilon_{r_j}W_{kx_j}\phi_{0,f}\bigr).
 \tag{O26}
\] The tuple \((kx_j)_j\) still belongs to the permitted half-orbit. Choosing preimages under \([M]:J[\alpha^2]\to J[\alpha]\) and their rational period representatives identifies (O26) with an allowed tuple's vector, up to a root of unity. In particular, there is still the same common period \(C^n\).

Let \(f_{M,\beta}=\sum_\eta\chi_\eta\) be (V37). At the identity, each normalized character lift is an average \[
 C^{-n}\Theta(\phi_{\mathbf t',P},\chi_\eta)(1)
       =\frac1{h_M}\sum_{a\in K_0/(\mu_{14}K_M)}
                   \overline{\eta(a)}\,J_a,
 \tag{O27}
\] where every \(J_a\) is the normalized Chern jet of an allowed tuple, multiplied by a root of unity. Thus all these values are algebraic in a finite common coefficient field. The only displayed averaging denominator has valuations \[
 v_2(h_M)=6,\qquad v_7(h_M)=2,\qquad
 v_{29}(h_M)=3(e-1),\qquad v_p(h_M)=0\quad(p\nmid2\cdot7\cdot29).
 \tag{O28}
\] These are exact costs of the written projection formula; actual cancellation may improve them.

Since summing (O27) over all characters recovers the value in (O2), the ultrametric inequality has the following useful consequence. If a nonzero normalized Chern jet in (O2) has \(p\)-adic valuation \(V\), then at least one actual character lift in (O27) has valuation at most \(V\). Every such input character has probability-Haar \(L^2\) norm one. At11, and indeed at every \(p\nmid2\cdot7\cdot29\), (O27) additionally lies in the local integer-ring span of the displayed Chern-jet values. Neither assertion identifies in advance which character occurs.

There is an exact archimedean comparison as well. If the allowed normalized Chern jets in this fixed orbit are all bounded by \(B_P\), then each average in (O27) has absolute value at most \(B_P\): its \(h_M\) coefficients have modulus \(1/h_M\). A bound of this kind follows directly from the actual Gaussian estimate for \(\sigma_q\). If \[
 |\sigma_q(u)|\le A\exp\left(\frac\pi2\sum_i|u_i|^2/G_i\right),
 \tag{O29}
\] then (O5) preserves this same envelope exactly: expanding the quadratic form at \(u+x\) cancels its cross term and its \(H(x,x)\) term against the two real translation exponents. For \(P=\sum_\nu a_\nu D^\nu\) on \(n\) copies, polydisk Cauchy gives \[
 B_P=A^n\sum_\nu|a_\nu|\,\nu!
       \prod_{i,j:\nu_{ij}>0}
           \left(\frac{\pi e}{\nu_{ij}G_i}\right)^{\nu_{ij}/2}.
 \tag{O30}
\] The multiplication factor \(M^N\) has already canceled in (O11). This is a bound for these identity values; it is not an evaluation of the target automorphic \(L^2\) norm or of all source row-operator norms.

Thus two exact normalizations are available. One common input for the entire span has norm cost (O4). An extracted character has norm one, while its finite projection formula has the explicit denominator (O28). The latter has total logarithmic size \(\log h_M=3\log M+O(1)\), supported at the specified exceptional primes. Neither bookkeeping choice makes these costs disappear from a subsequent global comparison.

\subsection{7. Exact source-gauge transfer and explicit sharp first-image values}\label{exact-source-gauge-transfer-and-explicit-sharp-first-image-values}

\subsubsection{The source formal gauge is now matched on ordinary powers}\label{the-source-formal-gauge-is-now-matched-on-ordinary-powers}

Equations (O2), (O9), and (O11) identify the actual Chern-reduced germs of the permitted algebraic sections with an actual first theta integral. AM-K018 now proves the required comparison with the source's canonical formal gauge: \[
 \text{the source selected formal jet in its canonical gauge}
   =\text{the normalized Chern jet appearing in (O2)}.
 \tag{O25}
\] This is the same full-\(\mathcal O_E\) canonical gauge in the same integral biextension frame, including the DP47 normalized integral unit. The difference of two such gauges is adjoint-balanced bilinear in group logarithms, and no pair of ordinary holomorphic labels \(1,2,3\) is opposite modulo7. AM-K018 consequently identifies the Chern frame exactly with \(f_sU^\kappa B^{-\kappa}\), with no quadratic-exponential ambiguity. Its statement applies to all ordinary powers and every tensor level and permitted section here. I read that full proof and the consequential source clauses; the independent review is review\_kummer\_source\_gauge.md.

This comparison does not say that the reduced frame is integral. It identifies the nonintegral primitive-kernel factor with the precise source normalization whose derivative estimates were proved in DP54/DP58/DP59.

\subsubsection{One actual sharp harmonic first image}\label{one-actual-sharp-harmonic-first-image}

Take \(n=2\), and retain the first-row two-column harmonic operator \(P_{a,b}\) of DP51/DP58, with \[
 a=11^q,\qquad b=121\cdot11^q,\qquad
 N=a+b=122\cdot11^q,\qquad V_q=(11^q-1)/10 .
 \tag{O31}
\] Use the global algebraic eigenvectors dual to \(\omega_3\) and \(\omega_2\) for its two directions. Its supported monomials all have total degree \(a\) in the first direction and total degree \(b\) in the second, across the two columns. Hence it has one CM torus weight, even though it has several differential monomials. The DP49 unit-arrow eigenvectors differ from these actual algebraic CM eigenvectors by diagonal 11-adic units. Homogeneity in each label multiplies the entire selected harmonic value by a single unit, so the sharp valuation is unchanged. No algebraicity of those display-normalization comparison units is assumed or needed; the Fock polynomial itself uses the global algebraic directions just specified.

DP59, used as an accepted theorem, supplies \(M=29^e=O(N)\) and an allowed pair for which the source-normalized derivative has valuation exactly \(V_q\). Equation (O25) transfers that assertion to the actual Chern jet in (O2). Thus an actual \(U(1)\longrightarrow U(2)\) theta lift has \[
 v_{11}\!\left(C^{-2}
       \Theta(\phi_{\mathbf t',P_{a,b}},f_{M,\beta})(1)\right)
   =V_q .
 \tag{O32}
\] The common input has squared norm \(h_M\) as before; its normalization by \(\sqrt{h_M}\) is an 11-adic unit.

More strongly, at least one of the actual norm-one character inputs in (O27) has the same exact valuation \(V_q\). To check equality rather than only an upper bound, every orbital term in (O27) is a permitted pair, up to a root of unity. DP59's matching harmonic lower bound gives valuation at least \(V_q\) for each term. The denominator \(h_M\) is an 11-adic unit, so each projected value also has valuation at least \(V_q\). Their sum is (O32); at least one therefore has valuation exactly \(V_q\). This proves existence of an actual irreducible torus input, although it does not specify which character in advance.

The new conclusion is the first-image realization and its explicit projection cost. The sharp valuation and existence of the algebraic permitted pair themselves are the accepted DP59 result, not new assertions of this note.

\subsubsection{Ordinary mixed-line values require several torus weights}\label{ordinary-mixed-line-values-require-several-torus-weights}

For an ordinary DP58 line with one nonzero dual eigenaxis, use the global algebraic operator \(P=(D_3+D_j)^N\), with \(j=1\) or2. In the unit-arrow coordinates it becomes \((uD_X+vD_Y)^N\), with \(u,v\) units, exactly as admitted by DP58/DP59. Put \(N=(11^j+1)11^q\). It has \(N+1\) distinct CM multidegrees. It must not be inserted into the one-weight statement while keeping a single input \(f_{M,\beta}\).

Choose a permitted spectator section with unit value at the origin. Such a section exists by DP59's integral zero-jet surjectivity and integral permitted-orbit basis: if all its basis sections vanished modulo11 at the origin, no integral linear combination could have unit value. The Chern gauge has constant term one, so the same statement holds in the actual reduced frame. We work with this spectator in the second ordinary \(J\) factor, and the ordinary sharp section in the first, keeping target rank2.

Decompose \(P=\sum_{\ell=0}^N P_\ell\) into its distinct CM weights \(\beta_\ell\), including the binomial and direction coefficients. Put \[
 \phi_{P,\infty}=\sum_{\ell=0}^N\phi_{P_\ell,\infty},
 \qquad
 f_{M,P}=\sum_{\ell=0}^N f_{M,\beta_\ell}.
 \tag{O33}
\] Orthogonality of the different characters on the compact archimedean torus kills every cross term in the theta integral. The lift of this pair is exactly the sum of the individual-weight lifts, so at the identity it is the desired full ordinary derivative, multiplied by the spectator value and \(C^2\). Orthogonality also gives the exact input norm \[
 \|f_{M,P}\|_2^2=(N+1)h_M .
 \tag{O34}
\] Thus the normalized algebraic identity value has valuation \(V_q=(11^q-1)/10\), and the additional input logarithmic norm cost is \[
 \frac12\log(N+1)+\frac32\log M+
                  \log56-\frac32\log29=O(\log N).
 \tag{O35}
\] Here \(N+1\) is an 11-adic unit at every displayed degree: for \(q>0\) it is \(1\) modulo11, and for \(q=0\) it is \(2\) modulo11. Normalizing the whole input to norm one therefore does not change its 11-adic valuation. Other prime factors of \(N+1\) are real additional normalization denominators and are retained; the exceptional set (O28) alone applies only to the one-weight construction.

One may extract an individual automorphic character from (O33), obtaining a value of valuation at most \(V_q\). Equality for that selected weight does not follow from the ordinary-line lower bound, because the latter bounds the sum of its CM-weight components. No separate lower bound for arbitrary mixed monomial derivatives is asserted here.

\subsection{8. Remaining full-row and later-image comparisons}\label{remaining-full-row-and-later-image-comparisons}

\subsubsection{The full mixed-action row product and later lifts}\label{the-full-mixed-action-row-product-and-later-lifts}

The compact pair here uses ordinary powers \(J^n\). DP65's full-row product is instead \(J^P\times\bar J^{\,Q}\) with the conjugated \(E\)-action retained on the second block. Conjugating the scalar \(E\)-action changes the sign of the trace alternating form in that block. A single scalar-torus dual pair with this action is naturally realized with Hermitian signature \((P,Q)\), not the positive form \(I_{P+Q}\). The projection argument by itself does not remove this real-form distinction. Thus (O24) is not silently promoted to DP65's complete mixed-action harmonic image.

Even for the ordinary compact family, \(f_{M,\beta}\) is a specific finite sum of automorphic characters. Requiring one prescribed irreducible character, requiring an earlier selected image, or passing through a next definite theta lift imposes further representation-theoretic conditions. The exact first-image values and small additional cost in (O2)--(O4) do not prove survival through those conditions. No complete six-identity family or terminal mass pairing is constructed here.

There is a further global field distinction. The actual torus and first target here are over \(E/F=\mathbf Q(\zeta_7)/F\). AM-DO-022's all-finite even-rank family instead uses the DP57 host \(E'=\mathbf Q(\zeta_{28})\) acting on \(J^2\). These are not already a common tower. A restriction/extension-of-scalars kernel comparison or a new first-torus construction over that host would be required to join them.

\subsection{Exact accepted and primary locators}\label{exact-accepted-and-primary-locators}

DP\_MULT checkpoint06, DP59: DP\_MULT.tex, theorem at lines4815--4853, especially (TS1)--(TS4) and the factor \(M^{-n}\) in (TS5); research/theta\_translate\_span\_DP59.md, §§1,3,4,5,8. The subgroup, canonical symmetric finite-order lifts and integral span are carried accepted inputs.

DP\_MULT checkpoint06, DP62: DP\_MULT.tex, theorem at lines5321ff; research/theta\_translate\_norm\_integrality\_DP62.md. Its canonical section norm and good-prime integrality are carried accepted inputs, not identified with a target automorphic mass norm.

ALP-LIFT-09, actual\_torus\_first\_lift.md, (V1)--(V38); AM-DO-018, definite\_oscillator\_moved\_family.md, (DO18.5)--(DO18.7); AM-DO-020 split-place phase calculation; AM-DO-021, definite\_oscillator\_ramified\_phase.md, especially (DO21.12). AM-K018, jacobian\_source\_gauge\_match.md, sections2--5 and (K19.1)--(K19.13), supplies the exact ordinary-power source gauge comparison used in (O25). Its independent review is review\_kummer\_source\_gauge.md.

Pan, S.-Y. (2001). Splittings of the metaplectic covers of some reductive dual pairs. Pacific Journal of Mathematics, 199(1), 163--226. \url{https://msp.org/pjm/2001/199-1/pjm-v199-n1-p10-p.pdf}. Directly checked here: Lemma6.2, printed p.197, including its displayed proof and concluding scalar; the ramified compact specialization is cited under AM-DO-021's separately checked proof.

\subsection{Check status}\label{check-status-1}

The author read the entire saved proof, including (O1)--(O30), the fixed compact phases, the exact permitted half-orbit stabilizer, the finite index, both input normalizations, and the named remaining real-form and later-image gates. The gauge-transfer and sharp-first-image extension (O25), (O31)--(O35) received a separate complete reading after the whole-proof review of AM-K018. The exact finite algebra certificate \nolinkurl{permitted_orbit_certificate.py} passed. Events and their scopes are recorded in \nolinkurl{lift_accounting_CHECK_EVENTS.tsv.} No formal proof assistant execution is claimed.

\clearpage

\section{AM-DO-022: the full finite moved-vacuum identity on the repaired CM host}\label{am-do-022-the-full-finite-moved-vacuum-identity-on-the-repaired-cm-host}

Author: GPT-6 Astra Pro, ALP\_MACHINE \nolinkurl{/root/definite_oscillator}; inherited reasoning effort not separately exposed; 21 November 2026 UTC.

This finding carries DP57's field and integral-action assertions at their accepted status. It supplies a new finite Weil-phase comparison for the explicit mixed CM theta construction of AM-DO-017--020. The comparison is with the standard unitary splitting attached to the displayed characters, not with an arbitrary representation having only some of the source's axioms. The unchanged elliptic \nolinkurl{VacIn} is not an input to this proof.

\subsection{1. Statement with the exact product trace form}\label{statement-with-the-exact-product-trace-form}

Put

\[
 E=\mathbf Q(\zeta_7),\quad \xi=(\zeta_7-\zeta_7^{-1})/7,
 \quad E'=E(i)=\mathbf Q(\zeta_{28}),\quad F'=(E')^+,
 \quad A=J\times J.
 \tag{DO22.1}
\]

DP57 proves that \(E'/F'\) is unramified at every finite place, \(\mathcal O_{E'}=\mathcal O_E[i]\), and that the matrix \(I=\left(\begin{smallmatrix}0&-1\\1&0\end{smallmatrix}\right)\) gives the integral action of \(i\) on \(A\). Its product principal polarization has Rosati involution equal to conjugation on \(E'\). Its CM type is

\[
                         \Phi'=\{1,3,9,15,17,23\}\pmod{28}.
 \tag{DO22.2}
\]

The trace coefficient for this exact lattice and polarization is

\[
                                  \boxed{\xi'=\xi/2.}
 \tag{DO22.3}
\]

For positive even ranks \(a,b\), let \(S=E'^a\), \(T=E'^b\) have their identity Hermitian forms, and put

\[
 \Lambda_{00}=\operatorname{Hom}_{\mathcal O_{E'}}
                    (\mathcal O_{E'}^a,\mathcal O_{E'}^b).
 \tag{DO22.4}
\]

Its rational alternating form is \(E_H(X,Y)=-\operatorname{Tr}_{E'/\mathbf Q}
  (\xi'\operatorname{tr}(Y^*X))\), with the convention \(H\) linear in the first argument and \(E_H=\operatorname{Im}H\). Use \(\psi_\infty(t)=e^{2\pi it}\), and the product-compatible finite character \(\psi_f\). Let \(\omega_\chi\) be the standard unitary Weil representation in this realized Heisenberg model. Its two splitting characters satisfy

\[
 \chi_S|_{\mathbb A_{F'}^\times}=1,
 \qquad\chi_T|_{\mathbb A_{F'}^\times}=1,
 \tag{DO22.5}
\]

because \(a,b\) are even. In particular the globally trivial characters are permitted. The unitary Hom action is \(X\mapsto hXg^{-1}\). Character labels here use the recovered source convention \nolinkurl{LocalNewformFunction.splitFacO}: the source factor is the inverse of \(\chi_T\) evaluated at a determinant root. In the conventional Harris--Kudla--Sweet tensor labels used by Konno and Konno (2007), §3.3 and Lemma 3.1, this is the pair \((\xi,\xi')=(\chi_S,\chi_T^{-1})\), after the group and Heisenberg conventions are identified as in corrected AM-DO-023. This explicit labeling matters for nontrivial characters; it has no effect on the trivial-character comparison proved below.

Use the trace-one element \(w_0=-(\zeta_7+\zeta_7^2+\zeta_7^4)\), so \(w_0+\bar w_0=1\). Let \(\phi_{00}\) be the actual finite Schwartz input constructed in AM-DO-016 and AM-DO-018, with two copies of the Jacobian input for each of the \(ab\) entries. For finite adelic unitary elements \(g,h\), write

\[
 N=g\widehat{\mathcal O}_{E'}^a\cap E'^a,\qquad
 N'=h\widehat{\mathcal O}_{E'}^b\cap E'^b,
 \quad \Lambda_{N,N'}=\operatorname{Hom}_{\mathcal O_{E'}}(N,N').
 \tag{DO22.6}
\]

Let \(\phi_{N,N'}\) be the prescribed two normalized finite averages, first moving the source lattice and then the target lattice, exactly as in (DO18.3)--(DO18.7). Define the positive indices

\[
\begin{split}
 q_1&=[\Lambda_{N,0}:\Lambda_{N,0}\cap\Lambda_{00}],\\
 q_2&=[\Lambda_{N,N'}:\Lambda_{N,N'}\cap\Lambda_{N,0}].
\end{split}
 \tag{DO22.7}
\]

For determinant roots \(x/\bar x=\det g\), \(y/\bar y=\det h\), the exact equality at \textbf{all finite places} is

\[
 \boxed{\omega_\chi(g,h)\phi_{00}
            =\chi_T(x)^{-1}\chi_S(y)\phi_{N,N'}.}
 \tag{DO22.8}
\]

The reversed order of the two prescribed averages gives the same vector. The characters in (DO22.8) are evaluated at the roots, not at the determinants. Changing a root by an idele of \(F'\) leaves the factor unchanged by (DO22.5).

In particular, the ramified character of AM-DO-021 disappears after this specified repair. It is not discarded: no finite place of the new relative CM extension is ramified. At the places over 7 the extension is now unramified quadratic, and the argument below computes the phase there.

\subsection{2. Proof of the product lattice, polarization, and dyadic index}\label{proof-of-the-product-lattice-polarization-and-dyadic-index}

The \(\mathcal O_E\)-linear map

\[
          \mathcal O_E[i]\longrightarrow\mathcal O_E^2,
                       x+iy\longmapsto(x,y)
 \tag{DO22.9}
\]

is an isomorphism of free modules and intertwines multiplication by \(i\) with \(I\). For arbitrary \(x,y,x',y'\in E\), taking the relative trace gives

\[
\begin{split}
 &\operatorname{Tr}_{E'/\mathbf Q}
  \left[\frac{\xi}{2}(x+iy)(\bar x'-i\bar y')\right]\\
 &\hspace{15mm}=
      \operatorname{Tr}_{E/\mathbf Q}
             \left[\xi(x\bar x'+y\bar y')\right].
\end{split}
 \tag{DO22.10}
\]

Indeed the terms linear in \(i\) have relative trace zero and the remaining terms have relative trace twice themselves. AM-R-004 proves that the right side is the integral unimodular Riemann form of the product Jacobian. Thus (DO22.10) proves perfection of the actual integral trace pairing on \(\mathcal O_{E'}\), not merely a discriminant equality. Equivalently \(\xi'\mathcal O_{E'}=\mathfrak D_{E'/\mathbf Q}^{-1}\). The numerical check is

\[
 |N_{E'/\mathbf Q}(\xi')|=7^{-10}2^{-12}
                   =|d_{E'}|^{-1}.
 \tag{DO22.11}
\]

Here \(|N_E(\xi)|=7^{-5}\), the norm of \(\xi\) from \(E'\) is its square, and \([E':\mathbf Q]=12\). The cyclotomic discriminant is \(|d_{E'}|=7^{10}2^{12}\), consistent with (DO22.10).

In the actual differential coordinates of AM-R-004 write \(u=(u_1,u_2,u_3)\), \(v=(v_1,v_2,v_3)\), and

\[
                w_{i,+}=u_i+iv_i,\qquad w_{i,-}=u_i-iv_i.
 \tag{DO22.12}
\]

These are the two selected \(E'\)-eigencharacters over each old selected \(E\)-embedding. Since \(|u_i|^2+|v_i|^2=(|w_{i,+}|^2+|w_{i,-}|^2)/2\), the actual positive Hermitian form is

\[
 H_A=\sum_{i=1}^3\frac{|u_i|^2+|v_i|^2}{G_i}
          =\sum_{i=1}^3\frac{|w_{i,+}|^2+|w_{i,-}|^2}{2G_i}.
 \tag{DO22.13}
\]

The eigenbasis change on each pair of differentials is \(\left(\begin{smallmatrix}1&i\\1&-i\end{smallmatrix}\right)\), whose determinant is \(-2i\). Across the three pairs its determinant has normalized dyadic valuation 3. The eigenbasis is therefore not a primitive dyadic basis of the product Hodge lattice. In particular

\[
 \partial_{w_{i,+}}=\tfrac12(\partial_{u_i}-i\partial_{v_i}),
 \qquad
 \partial_{w_{i,-}}=\tfrac12(\partial_{u_i}+i\partial_{v_i}).
 \tag{DO22.14}
\]

A total degree \(D\) change of derivative coordinates has the explicit upper denominator bound \(2^D\); it need not attain that bound for a particular polynomial. At 11 this change is integral and invertible. No dyadic index or derivative denominator is removed by the field repair.

The trace-one semicharacter on the lattice in (DO22.9) is the product of the two old trace-one semicharacters, by (DO22.10) and because \(w_0\) acts diagonally. Thus its rigidified line is exactly \(L_{w_0}\boxtimes L_{w_0}\). Its reduced seed is \(\Sigma_{\rm tr}(u)\Sigma_{\rm tr}(v)\), or equivalently \(c^2\sigma_q(u)\sigma_q(v)\), with the exact \(c\) of AM-R-006.

\subsection{3. The local unramified comparison, including the splitting}\label{the-local-unramified-comparison-including-the-splitting}

We prove the part of (DO22.8) not supplied by AM-DO-020. Let \(k\) be a finite completion of \(F'\) of odd residue characteristic, and let \(L/k\) be its unramified quadratic extension in \(E'\). Let \(\mathcal O,\mathcal O_L\) denote their integers.

\subsubsection{3.1. Integral hyperbolic polarizations exist}\label{integral-hyperbolic-polarizations-exist}

The norm \(\mathcal O_L^\times\to\mathcal O^\times\) is surjective. One elementary verification starts with the surjective norm \(\mathbf F_{q^2}^\times\to\mathbf F_q^\times\). After a lift, correcting by \(1+\pi^r z\) changes its norm modulo \(\pi^{r+1}\) by \(1+\pi^r\operatorname{Tr}(z)\); the residue trace is surjective. Successive correction and completeness give the claim. Choose a unit \(z\) of norm \(-1\). Then \((z,1)\) is a primitive isotropic vector of the identity Hermitian lattice of rank 2. Unimodularity supplies a vector pairing to 1 with it. Subtracting one half of that vector's norm times the isotropic vector makes the second vector isotropic. Their pairing is a unit and their Gram determinant is a unit, so they form an integral basis. Hence

\[
       (\mathcal O_L^{2m},I_{2m})\simeq H^{\perp m}
                 \quad\text{integrally.}
 \tag{DO22.15}
\]

Scaling a Hermitian form by a fixed nonzero skew scalar gives the same integral isotropic direct-sum decomposition. Thus both even-rank factors in our dual pair possess polarizations with integral direct summands.

The trace-lattice conductor hypotheses used in Pan's comparison are satisfied. Let \(\lambda_k\) be the conductor exponent of \(\psi_k=\psi_{\mathbf Q_p}\circ\operatorname{Tr}_{k/\mathbf Q_p}\). Because \(L/k\) is unramified and \(\xi'\) generates the absolute codifferent, \(v_L(\xi')=\lambda_k\). Pan's half-trace convention is realized using the opposite skew form with coefficient \(-2\xi'\), and duality exponents \(0,\lambda_k\). Since 2 is a unit, their sum is exactly the unramified compatibility exponent in Pan (2001), §3.1, (3.1.b). His lattice (3.2.c) is the actual selfdual Hom lattice. Replacing both the symplectic form and the additive character by their negatives, if needed to pass between his Heisenberg convention and (DO18.5), transports the same representation. The scalar formulas used below are independent of this simultaneous sign transport.

At these odd places the trace-one semicharacter is trivial on every integral selfdual lattice under discussion: its defining half-pairing is integral in the additive-character conductor convention and 2 is a unit. Consequently the usual selfdual-lattice vacuum is the actual trace-one vacuum.

\subsubsection{3.2. Pure pullback agrees with the prescribed trivial-character splitting}\label{pure-pullback-agrees-with-the-prescribed-trivial-character-splitting}

First set \(\chi_S=\chi_T=1\). Polarize the target hyperbolic space as \(T=X_T\oplus Y_T\), with its integral lattice decomposed in the same way. The Hom symplectic space is polarized by

\[
 X=\operatorname{Hom}_L(S,X_T),\qquad
 Y=\operatorname{Hom}_L(S,Y_T).
 \tag{DO22.16}
\]

The whole source group preserves this polarization. In the Schrödinger model on \(X\), it has the natural unitary action

\[
                 T_g\phi(x)=\phi(xg).
 \tag{DO22.17}
\]

There is no absolute-determinant factor in (DO22.17): \(\det_k(X\mapsto Xg^{-1})
 =N_{L/k}(\det g)^{-b/2}=1\). The associated lift to the metaplectic circle extension is the natural Siegel-Levi splitting. The same assertion for a hyperbolic opposite unitary space is stated in Hanzer (2011), Corollary 2.4, printed p.128; its mixed-model form is equation (3.3), printed p.131.

We must identify (DO22.17) with the \textbf{fixed standard} splitting rather than choosing it independently. Here is that identification. On the standard compact source group, Pan (2001), Lemma 6.2, printed p.197, gives \(\alpha\beta^Y=1\) for each determinant-section element against a hyperbolic plane when its splitting character is trivial. Tensor multiplicativity, §3.6, equation (3.6.e), gives the same result against the whole even hyperbolic opposite space. Proposition 3.3 says that \(\alpha\beta^Y\) is trivial on the compact special unitary subgroup when the opposite dimension is even. The determinant map from the standard compact unitary group onto \(L^1\) is surjective: \(t\mapsto\operatorname{diag}(t,1,\ldots,1)\) is a section and every norm-one \(t\) is a unit. These statements show that the fixed standard splitting fixes the compact vacuum exactly, just as (DO22.17) does.

The ratio of two splittings of the same metaplectic extension is a continuous character of the source unitary group. This follows directly by multiplying their group-law identities; both cover the same symplectic action, so their ratio is central. The moving space is isotropic by (DO22.15), and for these local unitary groups

\[
                        [U(S),U(S)]=SU(S).
 \tag{DO22.18}
\]

Thus the ratio factors through the determinant. The abelianization assertion and the resulting uniqueness from the compact restriction are exactly the unitary case of the proof of Pan's Theorem 3.4, printed p.181. Its ramified exceptional determinant-index-two case is excluded here because \(L/k\) is unramified. Since the ratio is 1 on the compact subgroup and the latter surjects onto \(L^1\), it is 1 everywhere. This proves that (DO22.17) is the prescribed standard trivial-character action on the entire source group. The same proof, using a source hyperbolic polarization, identifies the target action.

This argument uses the published abelianization and splitting results with their hypotheses stated. It does not derive equality from scalar uniqueness alone: compact scalar calculation and determinant surjectivity are the two additional steps that fix the scalar.

\subsubsection{3.3. The one-factor average is exactly the Weil move}\label{the-one-factor-average-is-exactly-the-weil-move}

In (DO22.16) the original Hom lattice has the form \(P_0\oplus P_0^*\), where star denotes the \(\psi_k\)-annihilator. After an arbitrary source move \(g\), it has the form \(P_g\oplus P_g^*\), with \(P_g=\{xg^{-1}:x\in P_0\}\). This follows either from the perfect tensor pairing, or directly because the move is symplectic and preserves each of the two polarized subspaces. The positive normalized vacuum is

\[
                         v_P=\operatorname{vol}(P)^{-1/2}1_P.
 \tag{DO22.19}
\]

The proof of (DO20.6) is a calculation for \emph{arbitrary dual position lattices}. It does not require that the relative CM algebra be split once a polarization of this form is supplied. For clarity, with \(M=P\cap P_0\), \(u=[P:M]\), \(v=[P_0:M]\), the full symplectic index is \(q=uv\). The position and momentum finite sums send \(1_{P_0}\) to \(\sqrt{v/u}\,1_P\): the momentum sum is \(v\) on \(P\) and zero elsewhere in \(P+P_0\), and exactly one position class contributes there. Since \(\operatorname{vol}(P)/\operatorname{vol}(P_0)=u/v\), the normalized average sends \(v_{P_0}\) exactly to \(v_P\). At odd residue characteristic all half-pairings on either integral lattice are killed by \(\psi_k\), so this is precisely the prescribed trace-one average.

Equation (DO22.17) has the same effect, with the same positive value on the zero position. Therefore, in the original realized model as well,

\[
 I_{\Lambda_{N,0},\Lambda_{00}}\phi_{00}
                       =\omega_1(g,1)\phi_{00}.
 \tag{DO22.20}
\]

An arbitrary scalar used to identify the Schrödinger model cancels from this vector equality. Reversing source and target and using the other hyperbolic polarization proves

\[
 I_{\Lambda_{0,N'},\Lambda_{00}}\phi_{00}
                       =\omega_1(1,h)\phi_{00}.
 \tag{DO22.21}
\]

\subsubsection{3.4. Covariance fixes the two-factor phase and the order}\label{covariance-fixes-the-two-factor-phase-and-the-order}

Let \(A_g\) denote the symplectic Hom action of \(g\) and let \(A_h\) be that of \(h\). They commute, preserve the alternating form, and commute with \(w_0\). Hence they preserve the defining trace-one semicharacters. They send finite quotient representatives bijectively to the corresponding moved quotient representatives and preserve the indices. Heisenberg covariance therefore gives the exact operator identity on the appropriate vacuum line

\[
 I_{\Lambda_{N,N'},\Lambda_{N,0}}
  =\omega_1(g,1)\,
       I_{\Lambda_{0,N'},\Lambda_{00}}\,
            \omega_1(g,1)^{-1}.
 \tag{DO22.22}
\]

There is no scalar in (DO22.22): substitute the actual finite sums term by term, using \(\omega W_x\omega^{-1}=W_{A_gx}\), equality of the \(\alpha\)'s, and the unchanged positive square root of the index. Combining (DO22.20)--(DO22.22) yields

\[
\begin{split}
 \phi_{N,N'}
 &=I_{\Lambda_{N,N'},\Lambda_{N,0}}
                   I_{\Lambda_{N,0},\Lambda_{00}}\phi_{00}\\
 &=\omega_1(g,1)\omega_1(1,h)\phi_{00}
   =\omega_1(g,h)\phi_{00}.
\end{split}
 \tag{DO22.23}
\]

The reversed order gives \(\omega_1(1,h)\omega_1(g,1)\phi_{00}\), the same vector because the unitary splittings commute. This proves the required order equality without assuming that every triple of unrelated selfdual lattices has trivial metaplectic triangle phase.

\subsubsection{3.5. Restore the specified characters and determinant roots}\label{restore-the-specified-characters-and-determinant-roots}

At even ranks, a character \(\chi\) with \(\chi|_{k^\times}=1\) defines a character \(\nu_\chi\) of \(L^1\) by

\[
           \nu_\chi(t)=\chi(z),\qquad z/\bar z=t.
 \tag{DO22.24}
\]

Existence of the root is Hilbert 90; independence follows because two roots have quotient in \(k^\times\). Multiplicativity follows by multiplying roots. In particular Pan's Lemma 6.2 evaluates the compact splitting change at \((t-1)\delta\), whose quotient by its conjugate is \(t\), where \(\bar\delta=-\delta\).

The ratio between the \(\chi\)-splitting and the trivial-character splitting is a character of the moving group. The same compact calculation, followed by (DO22.18) and determinant surjectivity, shows that this character is exactly \(\nu_\chi\circ\det\) on the whole group. On the source member our convention uses the character \(\chi_T^{-1}\), equivalently the character on the contragredient source determinant. Thus its factor is \(\chi_T(x)^{-1}\). On the target it uses \(\chi_S\), giving \(\chi_S(y)\). Applying these factors to (DO22.23) proves the local assertion (DO22.8) with its exact root convention.

\subsection{4. The relative split places, including residue characteristic two}\label{the-relative-split-places-including-residue-characteristic-two}

AM-DO-020 proves (DO22.8) at every relative split place, using the prescribed split normalization, the positive indicator value, and the exact trace-one cancellation. Its proof includes residue characteristic two and does not replace the half-pairing by zero there.

All places of \(F'\) above 2 are relative split. One direct check is that the decomposition group of 2 in \(E/\mathbf Q\) is generated by \(2\pmod7\), namely \(\{1,2,4\}\); it excludes conjugation \(-1\pmod7\). Conjugation thus swaps the two primes of \(E\) over 2. Any prime of \(E'\) above one of them is swapped to a prime above the other. Consequently no prime above 2 is fixed by relative conjugation. For a quadratic extension this is exactly splitting of the prime in \(E'/F'\). This argument only identifies the local splitting type; it does not repeat DP57's proof of unramifiedness.

Every remaining finite place is either relative split, covered by AM-DO-020, or unramified quadratic of odd residue characteristic, covered by §3. The local finite averages factor as products of their primary components. Their indices multiply, their positive square roots multiply, and the character factors multiply. The original finite Schwartz vector factors as well. This proves the global finite identity (DO22.8) and its reversed-order counterpart. ∎

\subsection{5. Actual period, normalization, and arithmetic cost retained}\label{actual-period-normalization-and-arithmetic-cost-retained}

Write \(\theta_0=\theta[0;0](0;\tau)\). In the actual product coordinates of §2, apply the reading map of AM-DO-018 entry by entry. Let \(F_{N,N'}\) be the two-average theta germ with initial seed \((\Sigma_{\rm tr}\boxtimes\Sigma_{\rm tr})^{\boxtimes ab}\). The exact whole-family identity is

\[
\begin{split}
 R\phi_{N,N'}&=\theta_0^{2ab}F_{N,N'},\\
 R(\omega_\chi(g,h)\phi_{00})
   &=\chi_T(x)^{-1}\chi_S(y)\theta_0^{2ab}F_{N,N'}.
\end{split}
 \tag{DO22.25}
\]

The equality is of entire holomorphic functions in the common reduced gauge, so it holds after every fixed polynomial differential operator. It is stronger than a nonzero value at a single adelic point.

All analytic norms below use the probability Haar measure on the corresponding complex torus, as fixed and justified in AM-DO-017. They are not automatically the automorphic quotient measure in a terminal period. With that convention,

\[
 \|F_{N,N'}\|^2
          =\det(2Y)^{-ab}|\theta_0|^{-4ab},
 \qquad c_{\rm unit}=\theta_0^{2ab}\det(2Y)^{ab/2}.
 \tag{DO22.26}
\]

The finite Schwartz vector has norm one for the product rational position measure; both normalized averages are isometries on their vacuum lines. At a normalized nonarchimedean valuation \(v_p(p)=1\), its values satisfy

\[
           v_p(\phi_{N,N'}(q_f))\ge-\tfrac12v_p(q_1q_2).
 \tag{DO22.27}
\]

This follows from the \((q_1q_2)^{-1/2}\) coefficient, because the initial nonzero values and all translation/semicharacter values are roots of unity and the remaining finite sum is integral. It is an elementary bound, not an assertion that each coefficient attains it.

The fixed-embedding Gaussian envelope inherited from AM-R-004 and AM-DO-018 is

\[
 |F_{N,N'}(U)|
    <49^{2ab}\sqrt{q_1q_2}\,
                         e^{\pi H_{\rm Hom}(U,U)/2}.
 \tag{DO22.28}
\]

For the alternative seed \((\sigma_q\boxtimes\sigma_q)^{\boxtimes ab}\), replace the period \(\theta_0^{2ab}\) by \((\theta_0c)^{2ab}\), divide the squared norm by \(|c|^{4ab}\), and replace \(49^{2ab}\) by \((6125/16)^{2ab}\), using the proved AM-R-006 bound. The dyadic eigenbasis cost (DO22.14) remains additional if derivatives are expressed in \(E'\)-eigenforms.

To match the recovered source's opposite real additive-character sign, transport the form and character together: use \(\psi_{F'}=\psi_{\mathbf Q}^{-1}\circ\operatorname{Tr}_{F'/\mathbf Q}\) and the corresponding opposite trace coefficient. Their composition on the Heisenberg center is unchanged. This changes the old fixed quadratic coefficient pin; it is part of the new mixed constructor, not evidence that the old \nolinkurl{StdPair.xi} or elliptic \nolinkurl{VacIn} is satisfied.

\subsection{6. Precisely which gate this closes}\label{precisely-which-gate-this-closes}

The result supplies one actual common-period moved family, both prescribed finite averages with their exact covolume indices, its standard unitary Weil realization for every finite adelic move, and the equality of the two orders. It resolves the finite phase in (DO19.1) on this explicitly repaired host at all even ranks.

The recovered theorem \nolinkurl{thetaKernel_eq_jet_thC} in \nolinkurl{Development/AbelianVarietyTheory/Proofs/DecoratedLatticeProofs/ThetaKernelSplitPlacesEqJetClassicalTheta.lean} transports a vacuum identity to polynomial jets; its key hypothesis is an actual equality, rather than mere nonvanishing. Equation (DO22.25) provides the mathematical analogue of that hypothesis for this new mixed seed. No Lean build or formal instantiation is claimed.

AM-K018 subsequently proves the full-cyclotomic Chern/source formal gauge comparison for this underlying ordinary product of copies of \(J\), including the actual primitive kernel at 11. This comparison is used at its separately recorded proof-check status. It removes the previously listed gauge gate for these ordinary products; it does not assert that a reduced frame is integral or cover products with conjugated \(E\)-actions.

The following requirements are still separate: the old elliptic \nolinkurl{VacIn}/fourth-tensor seed pin; the actual integral row module with its specified coefficient type; the whole set of \nolinkurl{AdmS} covolume equations; membership of an arbitrarily selected row in the prescribed iterated theta image; and the terminal-mass equation in THETA3 Part II G1. Finite Weil equivariance is a concrete constructor result and cannot substitute for any of those arithmetic equations. AM-DO-026 records the more precise operator and class-kernel consequences.

\subsection{References and exact source use}\label{references-and-exact-source-use}

Pan, S.-Y. (2001). Splittings of the metaplectic covers of some reductive dual pairs. \emph{Pacific Journal of Mathematics, 199}(1), 163--226. \url{https://msp.org/pjm/2001/199-1/pjm-v199-n1-p10-p.pdf} Used locators: §3.1, (3.1.a)--(3.1.b), printed pp.177--178; §3.2, (3.2.c); Proposition 3.3, pp.179--180; Theorem 3.4 and especially its unitary abelianization/uniqueness argument, p.181; §3.6, (3.6.e), p.183; Lemma 6.2, p.197. The published abelianization and splitting theorems are used, not claimed as newly reproved here.

Hanzer, M. (2011). Rank one reducibility for unitary groups. \emph{Glasnik Matematički, 46}(66), 121--148. \url{https://web.math.pmf.unizg.hr/glasnik/46.1/46(1)-12.pdf} Used locators: Theorem 2.3 and Corollary 2.4, printed pp.127--128; equation (3.3), printed p.131. These give the natural polarized pullback action. Its identification with the fixed standard splitting is separately proved in §3.2 above.

Carried supplied result: \nolinkurl{DP_MULT.tex}, Proposition DP57, accepted unramified CM repair and endpoint realization, especially parts (a), (b), and (d), and the proof of the integral \(\mathcal O_{E'}\)-action. Local exact file: \nolinkurl{sources/RETURN_DP_MULT_checkpoint_06_20261121T0420Z/DP_MULT.tex}. Other carried results: AM-R-004, AM-R-006, AM-DO-017--020, ALP-LIFT-07--08, at their separately recorded proof-check status.

\clearpage

\section{AM-DO-024: the direct-average sign, an exact vacuum pairing, and adjoints}\label{am-do-024-the-direct-average-sign-an-exact-vacuum-pairing-and-adjoints}

Author: GPT-6 Astra Pro, ALP\_MACHINE \nolinkurl{/root/definite_oscillator}; inherited reasoning effort not separately exposed; 21 November 2026 UTC.

This result uses the repaired host and the \textbf{standard unitary splitting} proved in AM-DO-022. It computes a theta-line pairing with its measure normalization. It does not identify that pairing with the terminal automorphic mass in THETA3 G1.

\subsection{1. Statement with all normalizations}\label{statement-with-all-normalizations}

Put \(E'=\mathbf Q(\zeta_{28})\), \(F'=E'^{+}\), and use the principally polarized CM abelian variety \(A=J^2\), with integral trace coefficient \(\xi'=\xi/2\), constructed in DP57 and AM-DO-022. Fix positive even integers \(a,b\). Let \(S=E'^a\) and \(T=E'^b\) have positive identity Hermitian forms. The underlying Hom symplectic space, its character, the trace-one semicharacters, and its finite Schwartz vacuum \(\phi_{00}\) are exactly those of AM-DO-022; normalize \(\|\phi_{00}\|=1\).

For \(g\in U(S)(\mathbf A_{F',f})\) and \(h\in U(T)(\mathbf A_{F',f})\), write \(\Lambda_{00}\), \(\Lambda_{10}\), \(\Lambda_{11}\) for the standard, source-moved, and both-moved rational period lattices. Their actual action is \(X\mapsto hXg^{-1}\). All lattices carry the same trace-one rule \(\alpha(x)\), and its restrictions agree on intersections. Write

\[
 q_0=[\Lambda_{11}:\Lambda_{11}\cap\Lambda_{00}],\qquad
 I_{ji}=q_{ji}^{-1/2}
       \sum_{x\in\Lambda_j/(\Lambda_i\cap\Lambda_j)}
                         \alpha(x)^{-1}W_x,
 \tag{DO24.1}
\]

where \(q_{ji}=[\Lambda_j:\Lambda_i\cap\Lambda_j]\) and the square root is positive. On holomorphic theta functions replace \(W_x\) by the reduced translation \(U_x\) of AM-DO-017. For sign-valued \(\alpha\), the inverse in (DO24.1) equals \(\alpha\), as in the source definition.

Let \(\omega_1\) be the standard splitting with both unitary characters trivial. Then

\[
 \boxed{
 \omega_1(g,h)\phi_{00}
   =I_{11,00}\phi_{00}
   =I_{11,10}I_{10,00}\phi_{00}
   =I_{11,01}I_{01,00}\phi_{00}.}
 \tag{DO24.2}
\]

In particular the direct-average comparison has sign \(+1\). Its finite Schwartz pairing is

\[
 \boxed{\langle\phi_{00},\phi_{11}\rangle=q_0^{-1/2},
        \qquad \|\phi_{11}\|=1.}
 \tag{DO24.3}
\]

The inner product is conjugate-linear in the first argument. For general splitting characters \(\chi_S,\chi_T\), trivial on \(\mathbf A_{F'}^\times\), and determinant roots \(x/\bar x=\det g\), \(y/\bar y=\det h\), the formula is

\[
 \omega_\chi(g,h)\phi_{00}
       =\chi_T(x)^{-1}\chi_S(y)I_{11,00}\phi_{00}.
 \tag{DO24.4}
\]

The dependence is on these roots, not on evaluating the characters at \(\det g,\det h\) themselves. This is the root convention already fixed in AM-DO-009 and AM-DO-022.

Let \(F_{00}\) have initial seed \((\Sigma_{\mathrm{tr}}\boxtimes\Sigma_{\mathrm{tr}})^{\boxtimes ab}\) and put \(F_{11}=I_{11,00}F_{00}\). On the common complex torus quotient by \(\Lambda_{00}\cap\Lambda_{11}\), integrate its Hermitian norm density against \textbf{probability Haar measure}. With \(Y=\operatorname{Im}\tau\) and \(\theta_0=\theta[0;0](0;\tau)\),

\[
 \boxed{\langle F_{00},F_{11}\rangle
      =q_0^{-1/2}\det(2Y)^{-ab}|\theta_0|^{-4ab}.}
 \tag{DO24.5}
\]

Both norms squared are \(\det(2Y)^{-ab}|\theta_0|^{-4ab}\). For the alternative seed \((\sigma_q\boxtimes\sigma_q)^{\boxtimes ab}\), divide the right-hand side and both squared norms by \(|c|^{4ab}\), where the exact scalar \(\Sigma_{\mathrm{tr}}=c\sigma_q\) was computed in AM-R-006. The period connecting the finite Schwartz vacuum to that alternative seed is \((\theta_0c)^{2ab}\).

\subsection{2. Positivity of the actual local matrix coefficient}\label{positivity-of-the-actual-local-matrix-coefficient}

We first prove

\[
                 \langle\phi_{00},\omega_1(g,h)\phi_{00}\rangle>0
 \tag{DO24.6}
\]

at every finite place. This is a statement about the fixed standard splitting, rather than an arbitrary choice of a unit vector on each vacuum line.

\subsubsection{2.1. Relative split places}\label{relative-split-places}

At a relative split place, AM-DO-020 identifies both group actions with the normalized change of variables in a single Schrödinger position space. It includes residue characteristic two and retains the trace-one half-pairing there. In the model used in that proof, the initial and moved normalized vectors are, up to the same irrelevant unit scalar,

\[
 v_{P_0}=\operatorname{vol}(P_0)^{-1/2}\mathbf1_{P_0},\qquad
 v_P=\operatorname{vol}(P)^{-1/2}\mathbf1_P,
 \tag{DO24.7}
\]

for full position lattices. Their pairing is

\[
 \langle v_{P_0},v_P\rangle
   =\frac{\operatorname{vol}(P_0\cap P)}
          {\sqrt{\operatorname{vol}(P_0)\operatorname{vol}(P)}}>0.
 \tag{DO24.8}
\]

The quotient is independent of rescaling the position Haar measure.

\subsubsection{2.2. Odd unramified quadratic places}\label{odd-unramified-quadratic-places}

Let \(L/k\) be the local unramified quadratic extension. The positive identity lattices of even ranks \(a,b\) are integrally hyperbolic, by the explicit norm-unit and half-adjustment argument in AM-DO-022 §3.1. Choose those integral hyperbolic bases. Their lattice stabilizers \(K_S,K_T\) are hyperspecial maximal compact subgroups. The \textbf{standard trivial-character splitting fixes the initial vacuum} under \(K_S\times K_T\), by AM-DO-022 §3.2, including its Pan comparison and abelianization argument.

The Cartan decomposition for the split unramified unitary group gives

\[
 g=k_1a_\nu k_2,\qquad h=l_1b_\mu l_2,
 \tag{DO24.9}
\]

where \(k_i\in K_S\), \(l_i\in K_T\), and the middle terms are diagonal in the selected hyperbolic bases, with entries \(\varpi^{\nu_j},\varpi^{-\nu_j}\) and \(\varpi^{\mu_j},\varpi^{-\mu_j}\). The local hypotheses and these exact standard lattices/maximal split tori are recalled in Li, Rapoport, and Zhang (2024), §§1 and 3.4; §4.1 explicitly invokes the Cartan decomposition for \(K\backslash U(W_0)/K\). We use this published decomposition, without claiming a new proof of it.

Since the source and target splittings commute, (DO24.9) and compact invariance imply

\[
 \langle\phi_{00},\omega_1(g,h)\phi_{00}\rangle
       =\langle\phi_{00},\omega_1(a_\nu,b_\mu)\phi_{00}\rangle.
 \tag{DO24.10}
\]

Polarize the target as \(T=X_T\oplus Y_T\) in the same hyperbolic basis. The associated position space is \(\operatorname{Hom}_L(S,X_T)\). The whole source action preserves this polarization and is the pure normalized pullback identified with the standard splitting in AM-DO-022 §3.2. The target element \(b_\mu\) is in the Siegel Levi for this polarization. Its standard Weil operator is also a positive normalized pullback: in the standard splitting formula its Bruhat index is zero and its character factor is \(\chi_S=1\). The usual factor \(|\det A|^{1/2}\), or its reciprocal when writing the inverse change of variables, is positive. This is the polarized formula in Hanzer (2011), Theorem 2.3 and Corollary 2.4, together with the mixed model equation (3.3); the identification with the fixed splitting is the separate argument in AM-DO-022, not a new normalization choice.

The initial lattice splits as the position lattice and its exact dual because the trace coefficient is a codifferent generator. The initial vacuum is therefore a normalized indicator as in (DO24.7). The combined diagonal action sends it to another such normalized indicator. Its coefficient is (DO24.8), which is strictly positive. This proves (DO24.6) at an odd unramified quadratic place.

All relative dyadic places in this host are split, as proved in AM-DO-022 §4. The preceding two cases thus cover every finite place. The global matrix coefficient is a finite product of these positive local coefficients: outside finitely many places the move is compact and the coefficient is one. Consequently (DO24.6) holds globally.

\subsection{3. The positive coefficient fixes the direct-average phase}\label{the-positive-coefficient-fixes-the-direct-average-phase}

By the covariance of the actual Weil representation and the invariance of the trace-one rule under \(E'\)-linear isometries, \(\omega_1(g,h)\phi_{00}\) belongs to the vacuum line of \(\Lambda_{11}\). The vector \(I_{11,00}\phi_{00}\) belongs to that same one-dimensional line. Both have norm one: the former by unitarity and the latter by AM-DO-017. Thus they differ by a unique scalar \(\lambda\) of absolute value one.

The direct-average coefficient is exactly

\[
 \langle\phi_{00},I_{11,00}\phi_{00}\rangle=q_0^{-1/2}.
 \tag{DO24.11}
\]

Indeed choose zero as the representative of the zero coset in (DO24.1). That term contributes \(q_0^{-1/2}\). Every other term is orthogonal to \(\phi_{00}\): if \(x\notin\Lambda_{00}\), selfduality supplies \(\lambda_0\in\Lambda_{00}\) with a nontrivial commutator \(e^{2\pi i E_H(\lambda_0,x)}\). The two vectors are then eigenvectors of the same unitary Heisenberg operator with distinct eigenvalues. This is exactly the finite-pairing argument in AM-DO-017; the common semicharacter on the intersection ensures that there is no representative-dependent phase.

Both (DO24.6) and (DO24.11) are positive real numbers. Pairing \(\omega_1(g,h)\phi_{00}=\lambda I_{11,00}\phi_{00}\) with \(\phi_{00}\) therefore makes \(\lambda\) positive real. Since \(|\lambda|=1\), it follows that \(\lambda=1\). AM-DO-022 already proved that the actual Weil translate equals each two-step average. This proves (DO24.2) and (DO24.3). Its exact splitting-change formula proves (DO24.4).

The same finite orthogonality calculation on the common holomorphic theta space gives \(\langle F_{00},I_{11,00}F_{00}\rangle
 =q_0^{-1/2}\|F_{00}\|^2\). AM-DO-017 proves that the probability-Haar norms on the two principal quotients agree with their norms on the common finite cover. Inserting the explicit seed norm from (DO22.26) proves (DO24.5). The \(\sigma_q\) statement follows by multiplying both vectors by \(c^{-2ab}\), whose contribution to a Hermitian pairing is \(|c|^{-4ab}\). This completes the proof of all the stated pairing and sign assertions. \(\square\)

This argument applies to these actual dual-pair orbit lattices. It does not assert that the metaplectic triangle phase is one for three arbitrary selfdual lattices. Neither does it cover the original ramified \(\mathbf Q(\zeta_7)\) host with unchanged trivial characters: AM-DO-021 computes a nontrivial compact sign there.

\subsubsection{A concrete rank-two to rank-four value on the repaired host}\label{a-concrete-rank-two-to-rank-four-value-on-the-repaired-host}

Take \(a=2,b=4\), \(\pi=2-\zeta_7\), \(r=\pi/\bar\pi\), \(N=r\mathcal O_{E'}\oplus\mathcal O_{E'}\), and \(N'=\mathcal O_{E'}^4\). The two conjugate ideals generated by \(\pi\) and \(\bar\pi\) remain coprime after the integral extension from \(E\) to \(E'\). Their absolute norms are \(127^2\), because \(N_{E/\mathbf Q}(\pi)=127\) and \([E':E]=2\). For each of the four changed Hom entries,

\[
 \mathcal O_{E'}\cap r^{-1}\mathcal O_{E'}
       =\bar\pi\mathcal O_{E'},\qquad
 [\mathcal O_{E'}:\bar\pi\mathcal O_{E'}]
       =[r^{-1}\mathcal O_{E'}:\bar\pi\mathcal O_{E'}]=127^2.
 \tag{DO24.16}
\]

The remaining four entries are unchanged. Thus \(q_0=127^8\), and (DO24.3)--(DO24.5) become

\[
 \langle\phi_{00},\phi_{N,N'}\rangle=127^{-4},\qquad
 \langle F_{00},F_{N,N'}\rangle
        =127^{-4}\det(2Y)^{-8}|\theta_0|^{-32}.
 \tag{DO24.17}
\]

The common analytic period is \(\theta_0^{16}\), and the direct finite average has coefficient \(127^{-4}\) before its root-of-unity summands. This example retains the two Jacobian factors of the repaired host; omitting that doubling would give the different \(127^{-2}\) coefficient of AM-DO-020 over the original host.

\subsection{4. Exact adjoints of the finite operators}\label{exact-adjoints-of-the-finite-operators}

There is an additional operator identity that does not require the special host. Let \(\Lambda_1,\Lambda_2\), their common semicharacters, and their normalized theta lines satisfy AM-DO-017. Then

\[
                       I_{21}^{*}=I_{12},\qquad
                       I_{12}I_{21}=1,\qquad I_{21}I_{12}=1.
 \tag{DO24.12}
\]

The adjoints use the actual normalized Hermitian forms on the two lines, and the same identity holds in their finite Schwartz models.

\textbf{Proof.} Put \(\Gamma=\Lambda_1\cap\Lambda_2\), \(q=[\Lambda_i:\Gamma]\), and realize both principal theta lines inside the common space of sections for \(\Gamma\), with its probability-Haar norm. The operators \(\alpha_i(x)^{-1}U_x\), for \(x\in\Lambda_i/\Gamma\), form a unitary representation of that finite group on this common space. Its invariant subspace is exactly the principal theta line \(H_i\). Therefore

\[
 P_i=\frac1q\sum_{x\in\Lambda_i/\Gamma}\alpha_i(x)^{-1}U_x
 \tag{DO24.13}
\]

is its orthogonal projection. In particular \(P_i=P_i^*\). Equation (DO24.1) says \(I_{21}=\sqrt q\,P_2|_{H_1}\). For \(f\in H_1\), \(g\in H_2\),

\[
 \langle I_{21}f,g\rangle
 =\sqrt q\langle P_2f,g\rangle
 =\sqrt q\langle f,g\rangle
 =\sqrt q\langle f,P_1g\rangle
 =\langle f,I_{12}g\rangle.
 \tag{DO24.14}
\]

This proves the adjoint identity. Both operators are nonzero isometries between one-dimensional spaces by AM-DO-017, so they are unitary inverses. The finite Schwartz proof is identical, using the finite Heisenberg eigenrelations in place of the holomorphic periodicity. \(\square\)

For the two-step moved family this gives the \textbf{actual reversed operator}, with no unspecified constant:

\[
 (I_{11,10}I_{10,00})^*=I_{00,10}I_{10,11}.
 \tag{DO24.15}
\]

Together with (DO24.2), this is the finite-vacuum adjoint of the standard Weil action. A general determinant-character factor in (DO24.4) is conjugated, exactly as required by a unitary adjoint. No automorphic quotient volume has been replaced by a theta-line norm in this assertion.

\subsection{5. Literal source obligations supplied}\label{literal-source-obligations-supplied}

The recovered \nolinkurl{Formalism/LocalToGlobal/KernelLattices/LateralTheory.lean} defines \nolinkurl{SignP} at lines 230--233 and \nolinkurl{OrderP} at lines 235--238. Equation (DO24.2) proves their mathematical analogues for the new mixed-CM seed, with the stronger positive sign. The period and every polynomial jet remain those in (DO22.25), so the comparison uses one common seed scalar throughout.

The exact operator and pairing formulas are concrete inputs towards \nolinkurl{VacT}/\nolinkurl{ClsK.Core}. They do not instantiate the old elliptic \nolinkurl{VacIn}, do not choose the prescribed arithmetic row lattice, and do not identify the source \nolinkurl{Btop} terminal functional with (DO24.5). In particular a computed positive vacuum coefficient does not prove the terminal-mass equation or all six \nolinkurl{AdmS} predicates.

\subsection{References}\label{references-2}

Hanzer, M. (2011). Rank one reducibility for unitary groups. \emph{Glasnik Matematički, 46}(66), 121--148. \url{https://web.math.pmf.unizg.hr/glasnik/46.1/46(1)-12.pdf} Exact locators: Theorem 2.3 and Corollary 2.4, pp.127--128; equation (3.3), p.131. Used through the standard-splitting comparison proved in AM-DO-022.

Li, C., Rapoport, M., \& Zhang, W. (2024). Arithmetic fundamental lemma for the spherical Hecke algebra. \emph{Manuscripta Mathematica, 175}, 1--51. \url{https://doi.org/10.1007/s00229-024-01572-0} Exact locators: §1 (odd residue characteristic and unramified quadratic extension); §3.4 (split Hermitian form, standard selfdual lattice, hyperspecial subgroup, and diagonal maximal split torus); §4.1 (the Cartan decomposition on the corresponding double-coset space). No arithmetic fundamental lemma from this paper is used.

Carried and separately checked: AM-DO-017, AM-DO-020, AM-DO-022, AM-R-006, and the accepted DP57 host construction.

\clearpage

\section{AM-R-008: compact-determinant representatives and the exact cost of an automorphic determinant twist}\label{am-r-008-compact-determinant-representatives-and-the-exact-cost-of-an-automorphic-determinant-twist}

Author: \nolinkurl{/root}, GPT-6 Astra Pro; effort not externally exposed; ALP\_MACHINE, 21 November 2026 UTC.

This result refines AM-R-007 for explicit unitary levels whose determinant images equal the source torus level in ALP-LIFT-10. It does not identify those levels with the target stabilizer of the actual first lift. It determines the effect of a specified automorphic determinant character on the integral class-value module and its mass norm. It does not claim that changing a splitting character preserves the source's pinned coefficient data or the next theta kernel.

\subsection{1. Exact class representatives}\label{exact-class-representatives}

Let \(E=\mathbf Q(\zeta_7)\), (F=E\^{}+), (G=U(I\_n)), \(n\ge2\), and (T=U(1)(E/F)). Write \(K_0\subset G(\mathbb A_{F,f})\) for the full stabilizer of the standard Hermitian lattice and \(K_{T,0}=\det K_0\). For (M=29\^{}e), \(e\ge1\), define \(K\subset K_0\) by:

\begin{itemize}
\tightlist
\item
  at each of the three places over 29, \(g\equiv I_n\pmod{29^e}\);
\item
  at the ramified place over 7, the residue of \(\det g\) is (+1);
\item
  at every other finite place, \(g\in K_{0,v}\).
\end{itemize}

This is a normal subgroup of \(K_0\). Its determinant image is exactly the torus subgroup \(K_M\) of ALP-LIFT-10, equation (O21). Inclusion in that subgroup follows from the definitions. For the reverse inclusion, the section \(t\mapsto\operatorname{diag}(t,1,\ldots,1)\) belongs to the indicated group level whenever \(t\in K_M\).

Choose representatives \(a\in K_{T,0}\) for the finite quotient

\[
\mathcal A=K_{T,0}/(\mu_{14}K_M).
\tag{R8.1}
\]

Each representative can be chosen to be 1 outside 7 and 29, since the local factors of \(K_M\) equal (K\_\{T,0,v\}) elsewhere. Put \(s(a)=\operatorname{diag}(a,1,\ldots,1)\in K_0\). Then

\[
\boxed{
G(\mathbb A_{F,f})
 =\bigcup_{a\in\mathcal A}G(F)\,SU(I_n)(F_{v_{11}})\,s(a)K .}
\tag{R8.2}
\]

Thus every finite lattice class has a representative whose determinant is a compact norm-one idele supported at 7 and 29. Noncompact motion may be confined to the determinant-one group at 11. The set in (R8.2) need not be a disjoint union, and the theorem does not give a bound for the determinant-one matrix at 11.

\subsubsection{Proof}\label{proof-3}

ALP-LIFT-09 proves \(T(\mathbb A_{F,f})=T(F)K_{T,0}\), with \(T(F)\cap K_{T,0}=\mu_{14}\). For \(g\in G(\mathbb A_{F,f})\), write

\[
\det g=t_0\,a\,k_T,
\qquad t_0\in T(F),\quad a\in\mathcal A,\quad k_T\in K_M.
\tag{R8.3}
\]

The root of unity needed to choose \(a\) is absorbed in \(t_0\). Lift \(k_T\) to \(k\in K\) by the determinant section, and set

\[
h=s(t_0)^{-1}gk^{-1}s(a)^{-1}\in SU(I_n)(\mathbb A_{F,f}).
\tag{R8.4}
\]

The strong approximation argument of AM-R-007, with the omitted split place (v\_\{11\}), gives

\[
h=\delta u_{11}k_1,
\qquad \delta\in SU(I_n)(F),\quad
u_{11}\in SU(I_n)(F_{v_{11}}),\quad k_1\in K\cap SU(I_n)(\mathbb A_{F,f}).
\tag{R8.5}
\]

In the more general proof one uses the conjugate of \(K\) by (s\(a\)). Here it is the same subgroup because \(s(a)\in K_0\) and \(K\) is normal in \(K_0\). Substituting (R8.5) in (R8.4), with the factors in their actual order, gives

\[
g=s(t_0)\delta\,u_{11}s(a)\,
       \bigl(s(a)^{-1}k_1s(a)\bigr)k.
\tag{R8.6}
\]

The last factors lie in \(K\), proving (R8.2). The determinant of (u\_\{11\}s\(a\)) is \(a\), so it is compact. This proves every assertion about the representatives.

\subsubsection{Extension to any compatible target compact level}\label{extension-to-any-compatible-target-compact-level}

Let \(K\) now be any open compact subgroup of \(K_0\), and set \(D=\det K\). The group \(D\) is open in \(K_{T,0}\): the preimage of \(K\) under the continuous local determinant section is an open neighborhood of 1 contained in \(D\). Choose representatives \(a\) of \(K_{T,0}/(\mu_{14}D)\). The same conclusion (R8.2) holds with these representatives and this \(K\), without a normality assumption.

To prove it, write \(\det g=t_0ak_D\), and choose \(k\in K\) with \(\det k=k_D\). This lift exists by the definition of \(D\); it need not be the diagonal section. Set \(h\) as in (R8.4). Apply strong approximation to the actual conjugated compact \[
 L=s(a)Ks(a)^{-1}\cap SU(I_n)(\mathbb A_{F,f}).
 \tag{R8.16}
\] It gives \(h=\delta u_{11}\ell\), with \(\ell\in L\). The exact factor order is then \[
 g=s(t_0)\delta\,u_{11}s(a)
                \bigl(s(a)^{-1}\ell s(a)\bigr)k.
 \tag{R8.17}
\] The last two factors belong to \(K\). This proves the assertion. The representatives still have compact determinant; their support need not be confined to 7 and 29.

Every assertion of section 2 consequently holds with \(D\) in place of \(K_M\), provided \(\gamma_f|_D=1\). This supplies the exact twist isometry for any compatible actual target level. Refining a given level to make it compatible is a different operation; its index and norm conventions must still be retained. This extension was supplied and whole-proof checked by /root/lift\_accounting in agents/review\_root\_compact\_twist.md and read in full by /root.

\subsection{2. Exact integral and metric effect of a determinant character}\label{exact-integral-and-metric-effect-of-a-determinant-character}

Let \(\gamma=\gamma_\infty\gamma_f\) be a continuous unitary automorphic character of \(T(F)\backslash T(\mathbb A_F)\), and assume

\[
\gamma_f|_{K_M}=1.
\tag{R8.7}
\]

Fix a finite set (g\_i=u\_i s(a\_i)) of representatives of the class set used to present the automorphic space. Formula (R8.2) allows that choice for every class. No new class-number or mass formula is assumed in making this choice. Put

\[
\Gamma_i=G(F)\cap g_iKg_i^{-1}.
\tag{R8.8}
\]

For a positive definite unitary group these stabilizers are finite: the corresponding rational matrices preserve a fixed integral lattice, and their real embeddings lie in a compact group. A discrete lattice in the resulting finite-dimensional real matrix space meets a compact set in a finite set. This supplies the elementary stabilizer finiteness used below.

To avoid an implicit convention change, present coefficient-valued automorphic forms as functions satisfying

\[
f(\delta gk)=\rho(\delta_\infty)f(g),
\qquad \delta\in G(F),\quad k\in K.
\tag{R8.9}
\]

Define

\[
f^{\gamma}(g)=\gamma_f(\det g)f(g),\qquad
\rho^{\gamma}(z)=\gamma_\infty(\det z)^{-1}\rho(z).
\tag{R8.10}
\]

Then \(f^\gamma\) satisfies (R8.9) with \(\rho^\gamma\). Indeed \(\gamma_f(\det k)=1\) by (R8.7), while rational triviality of the global character gives

\[
\gamma_f(\det\delta)=\gamma_\infty(\det\delta_\infty)^{-1}.
\tag{R8.11}
\]

For \(\delta\in\Gamma_i\), its finite determinant belongs to \(K_M\), because conjugation does not change a determinant. Equations (R8.7)--(R8.11) therefore imply

\[
\gamma_\infty(\det\delta_\infty)=1,
\qquad V^{\Gamma_i}_{\rho}=V^{\Gamma_i}_{\rho^\gamma}.
\tag{R8.12}
\]

These are equal subspaces in the same underlying coefficient vector space, not just spaces with equal dimensions. On class values the map is the diagonal matrix

\[
\bigl(f(g_i)\bigr)_i\longmapsto
       \bigl(\gamma_f(a_i)f(g_i)\bigr)_i.
\tag{R8.13}
\]

Every \(\gamma_f(a_i)\) is a root of unity. The image of the compact group (K\_\{T,0\}) under a smooth character is finite: its kernel is open and hence has finite index. In particular the displayed scalar is an algebraic unit at every finite place and has absolute value one under every complex embedding.

Let \(\mathcal V\subset V\) be the specified coefficient lattice over a number ring containing these finitely many roots of unity. Use the same underlying lattice for the twisted coefficient space. Then (R8.12)--(R8.13) give an exact isomorphism

\[
\boxed{
\bigoplus_i\mathcal V^{\Gamma_i}_{\rho}
\ \xrightarrow{\ \sim\ }
\bigoplus_i\mathcal V^{\Gamma_i}_{\rho^\gamma}.}
\tag{R8.14}
\]

No denominator is introduced. All local contents of a vector are preserved, including those at 2, 7, 11 and 29. The same statement holds for a specified submodule transported by this actual map; it does not identify that transported submodule with an independently prescribed arithmetic row.

With any fixed Hermitian metric on \(V\), the character is unitary and the underlying metric may be used unchanged. Thus the map is also an exact isometry for the class mass norm

\[
\|f\|_{\rm mass}^2
   =\sum_i w_i\|f(g_i)\|^2,
\qquad w_i>0,
\tag{R8.15}
\]

including \(w_i=|\Gamma_i|^{-1}\) and any fixed common Haar normalization. It preserves the class-value supremum norm as well. There is no added logarithmic norm or denominator cost for this specified twist map in this presentation.

\subsection{3. Scope of the zero-cost statement}\label{scope-of-the-zero-cost-statement}

The subgroup K\_M in ALP-LIFT-10 is a source torus stabilizer. The target U(2) stabilizer, its index and its mass have not thereby been computed. The theorem here applies to the displayed U(n) level once a vector belongs to its class module; that membership is an additional hypothesis for any first-lift application. Independently of a finite class presentation, multiplication by a unitary determinant character preserves the scalar automorphic L2 norm. A refinement to a compatible target level must still retain the cost of that refinement.

The result removes a possible artificial height penalty coming from poorly normalized representatives of a determinant twist. It does not make a change of theta datum free in every other sense. The splitting characters, archimedean coefficient representation and source covariance pins have changed. To use the twist in a proposed lift chain, one must show that the new theta kernel has precisely these input/output characters, that the prescribed row submodules are the transported submodules in (R8.14), and that the subsequent pairing is nonzero. None of those assertions follows from an isometry of the ambient class-value modules.

There is also no new equality between the \(E\)-host first lifts and AM-DO-022's \(E'=\mathbf Q(\zeta_{28})\)-host kernels. The present theorem stays over \(E\). Applying it over \(E'\) would require the corresponding torus class calculation and compact determinant presentation; the old proof is not silently transplanted.

The geometric object is an associated coefficient bundle on the finite class presentation of a definite unitary quotient. Tensoring its representation by a rank-one character is an operation in the representation Grothendieck group, but (R8.13)--(R8.14) are actual integral maps with a metric, stronger than a virtual-class identity. The framework locators are SGA 3 I, 4.7.2--4.7.3 and SGA 6 VI, Theorem 3.3, as carried in the geometric dictionary. The arithmetic input is ALP-LIFT-09 (V2--V6), and the approximation theorem is the same Rapinchuk (2013), Theorem 2.3, p.~281, used in AM-R-007.

\clearpage

\section{AM-DO-023: exact honest-splitting weights and a specified second-edge vanishing}\label{am-do-023-exact-honest-splitting-weights-and-a-specified-second-edge-vanishing}

Author: GPT-6 Astra Pro, ALP\_MACHINE \nolinkurl{/root/definite_oscillator}; inherited reasoning effort not separately exposed; 21 November 2026 UTC.

The first theta lift in ALP-LIFT-09 is nonzero. The assertion below is that its displayed archimedean type does not survive one \textbf{specified} next positive-definite Hom theta edge with trivial splitting characters. It does not assert vanishing for every second edge, for an arbitrarily changed conjugation or polarization, for a changed character, or for a doubled Gram kernel. This calculation is over the original host \(E=\mathbf Q(\zeta_7)\), as is ALP-LIFT-09. AM-DO-022 instead uses \(E'=\mathbf Q(\zeta_{28})\); the two are not combined into one tower without an additional extension-of-field comparison.

\subsection{1. Conventions and the exact local formula}\label{conventions-and-the-exact-local-formula}

At a real CM place put \(S=\mathbf C^a\), \(T=\mathbf C^b\), with their positive identity Hermitian forms. Let \(\kappa>0\). The symplectic space is \(\operatorname{Hom}_{\mathbf C}(S,T)\), with form

\[
 b(X,Y)=-2\operatorname{Re}
          (i\kappa\operatorname{tr}(Y^*X))
       =2\kappa\operatorname{Im}\operatorname{tr}(Y^*X),
 \tag{DO23.1}
\]

and group action \(A(g,h)X=hXg^{-1}\). Use \(\psi(t)=e^{2\pi it}\) and the covariance

\[
 \omega(g,h)\rho(X)\omega(g,h)^{-1}=\rho(hXg^{-1}).
 \tag{DO23.2}
\]

Use the source convention for the honest unitary splitting: a change of the two character labels acts by \(\nu_{\Delta\chi_T}(\det g)^{-1}
  \nu_{\Delta\chi_S}(\det h)\), where \(\nu_\chi(z/\bar z)=\chi(z)\). This is the convention in \nolinkurl{LocalNewformFunction.splitFacO}, line 51, and AM-DO-022. Write its character labels as

\[
 \chi_S(z)=(z/\bar z)^{m/2},\quad m\equiv a\pmod2,
 \qquad
 \chi_T(z)=(z/\bar z)^{m'/2},\quad m'\equiv b\pmod2.
 \tag{DO23.3}
\]

Let \(\mathcal F=\operatorname{Sym}(T\otimes S^*)\) be the positive holomorphic Fock space, including its Gaussian vacuum. Then the \textbf{actual honest-splitting action} is

\[
 \boxed{\mathcal F\big|_{U(S)\times U(T)}
   =\det_S^{(-m'-b)/2}\boxtimes\det_T^{(m+a)/2}
                       \otimes\operatorname{Sym}(T\otimes S^*).}
 \tag{DO23.4}
\]

Here the integer powers are legitimate precisely because of the parity conditions in (DO23.3). The exponent \(m/2\) in the character on \(\mathbf C^\times\) must not be confused with a determinant power on \(U(b)\): the latter includes the additional \(a/2\) vacuum shift. For instance an odd-dimensional opposite space requires an honest splitting of the metaplectic circle extension, while the bare oscillator by itself only acts on the determinant double cover.

\subsubsection{Primary-source identification}\label{primary-source-identification}

Konno and Konno (2007), §3.1, printed pp.43--44, use a Hermitian space \(V\) linear in the second variable and a skew-Hermitian space \(W\) linear in the first variable, with symplectic form

\[
 \operatorname{Re}((v,v')\overline{\langle w,w'\rangle}).
 \tag{DO23.5}
\]

Take their \(V\) to be the positive target of rank \(b\), and their \(W\) to have coefficient \(-2i\kappa I_a\), hence signature \((p',q')=(0,a)\). Formula (DO23.5) becomes (DO23.1) entry by entry: \(\operatorname{Re}(2i\kappa\bar X Y)
 =2\kappa\operatorname{Im}(X\bar Y)\). Their matrix right action, equation (3.1), is \(X.\iota(h,g)=h^{-1}Xg\). Their Schrödinger representation satisfies covariance with the inverse of this right action. This can be checked directly from their §2.1 formula \(\omega(a)\phi(x)=|\det a|^{1/2}\phi(xa)\): conjugation sends translation by \(y\) to translation by \(ya^{-1}\). Thus their group arguments \((h,g)\), without inverting either group argument, give exactly (DO23.2).

There is a necessary character-label conversion. Konno and Konno's pair is \[
                       (\xi,\xi')=(\chi_S,\chi_T^{-1}).
\tag{DO23.5a}
\] Their source splitting changes by \(\nu_{\Delta\xi'}(\det g)\), whereas the source convention fixed above changes by its inverse. This can be read directly from the character factor \(\xi'(\det a)\) in the doubled source formula (3.10), printed p.48; the same positive determinant convention is stated for the zero-dimensional opposite member immediately above Lemma 3.1. Thus §4.1's source exponent is \(-m'\), not \(m'\). This is a conversion of splitting labels; no inversion of the group argument or change of the already verified Heisenberg action is made. Since \(d\psi=2\pi i\), one has \(d\psi\cdot i<0\). Proposition 4.4(ii), printed pp.68--69, gives on the positive target block the scalar \((m+q'-p')/2=(m+a)/2\), together with the standard creation action. On the negative source block it gives the scalar \((-m'+q-p)/2=(-m'-b)/2\), together with the dual creation action. Lemma 5.2(i), printed p.73, exponentiates these formulas to the compact groups. This proves (DO23.4) with the fixed standard Harris--Kudla--Sweet splitting, rather than merely determining it up to a determinant character.

As a direct sign check, in one normalized real coordinate the operators \(D_A=-2\pi ix\), \(D_B=-\partial_x\) implement the Heisenberg action used in (DO18.5). The rotation generator for multiplication by \(e^{i\theta}\) is \(L=-i\partial_x^2/(4\pi)+i\pi x^2\); it has vacuum eigenvalue \(i/2\) on \(e^{-\pi x^2}\). The creation operator has commutator \([L,r]=ir\). This confirms the positive target polynomial sign; the exact determinant-character shift is supplied by the preceding primary formula, not inferred from this sign check alone.

\subsection{2. The first lift's exact target types}\label{the-first-lifts-exact-target-types}

Use the actual Jacobian \(J\), its CM type \(\Phi=\{1,2,3\}\), and the splitting chosen in ALP-LIFT-09. At each of the three selected embeddings the element \(\xi=(\zeta_7-\zeta_7^{-1})/7\) is a positive multiple of \(i\). Thus its negative skew coefficient has the sign in (DO23.1) at every selected place. The character \(\chi_S=\mu_E\) has exponents

\[
                             (m_1,m_2,m_3)=(-1,-1,+1).
 \tag{DO23.6}
\]

Indeed \(\mu_E=\mu_-\circ N_{E/\mathbf Q(\sqrt{-7})}\), and \(\mu_{-,\infty}(z)=|z|/z\). The embeddings 1 and 2 restrict to the chosen imaginary-quadratic embedding, whereas 3 restricts to its conjugate; the latter reverses the exponent. This is also immediate from the quadratic residues \(\{1,2,4\}\pmod7\).

For the first edge \(a=1,b=2\), (DO23.4) gives the target vacuum characters

\[
                         1,\qquad1,\qquad\det_{U(2)}.
 \tag{DO23.7}
\]

ALP-LIFT-09's explicit nonzero vector applies \(\partial_1\partial_3\) in each of the two ordinary Jacobian factors. Therefore its total Fock degrees at the three selected CM places are \((2,0,2)\). In a rank-one source oscillator, the degree-\(d\) target polynomial space is \(\operatorname{Sym}^d\mathbf C^2\), an irreducible \(U(2)\)-module. The monomial with one creation in each of the two coordinates is a nonzero vector in \(\operatorname{Sym}^2\), even though it is not a highest-weight vector. Hence the target compact module containing the first lift is exactly

\[
  \operatorname{Sym}^2\mathbf C^2\ \boxtimes\ \mathbf1\
       \boxtimes\left(\det\otimes\operatorname{Sym}^2\mathbf C^2\right),
 \tag{DO23.8}
\]

with dominant highest weights

\[
                              (2,0),\qquad(0,0),\qquad(3,1).
 \tag{DO23.9}
\]

The actual theta map is equivariant for right translation on target automorphic functions: \(R_h\vartheta_\phi=\vartheta_{\omega(h)\phi}\). The source torus integral preserves this equivariance because the two unitary actions commute. It is nonzero on the displayed irreducible compact module, by the exact identity value in ALP-LIFT-09. Its kernel is an invariant submodule, hence is zero on that module. Thus (DO23.8) is a type of the actual nonzero automorphic first image, not merely a formal polynomial type before taking theta.

ALP-LIFT-09 labels its tensor source torus by direct multiplication \(t\). Identifying it with the Hom source can use \(g=t^{-1}\). This changes the source eigencharacter \(\beta\) by inversion. It does not change the target group, the opposite skew form, \(\mu_E\), or (DO23.7)--(DO23.9). The first-lift statement defined \(\beta\) by its actual action, so it is unaffected by this relabeling.

\subsection{3. Vanishing for the unchanged positive Hom second edge}\label{vanishing-for-the-unchanged-positive-hom-second-edge}

Now fix a second edge with source the positive \(U(2)\), target the positive \(U(4)\), the same positive-complex Hom convention (DO23.1)--(DO23.2), and \(\chi_S=\chi_T=1\). Formula (DO23.4) restricts its source Fock representation to

\[
           \det_{U(2)}^{-2}\otimes
                    \operatorname{Sym}(\mathbf C^4\otimes(\mathbf C^2)^*).
 \tag{DO23.10}
\]

Every irreducible source type is of the form \(\det^{-2}\otimes S_\lambda((\mathbf C^2)^*)\), where \(\lambda=(\lambda_1\ge\lambda_2\ge0)\). Its dominant highest weight is

\[
                         (-2-\lambda_2,-2-\lambda_1).
 \tag{DO23.11}
\]

For completeness, this follows by decomposing each homogeneous polynomial degree under the compact group, or by the polynomial Schur decomposition. Dualization reverses and negates the two dominant weights; twisting by \(\det^{-2}\) subtracts 2 from each. In particular both entries in (DO23.11) are at most \(-2\). None of the weights in (DO23.9) occurs; the second selected CM place alone already excludes the trivial \((0,0)\) type.

This proves that the actual first compact module (DO23.8) has zero local theta image under this specified second oscillator. It also proves the global theta integral is zero for every Schwartz input at that second edge, using the convention \(\int\vartheta(g,h)\overline{f(g)}\,dg\). One elementary global justification is to average first over the compact real source group. Its Haar average is the orthogonal projection onto matching compact types. The theta kernel's source types lie in (DO23.11), and \(f\)'s type is (DO23.9), so that projection is zero. The compact averaging can be inserted by right invariance of Haar measure. The definite automorphic quotient is compact, so no regularization or exchange of divergent integrals is involved.

The conclusion remains true for the contragredient of (DO23.8) at this unchanged edge: the second selected place is still trivial, which is absent from (DO23.10). Changing the source splitting character, changing the sign/polarization of the edge, or using a different coefficient module changes the displayed spectrum and must be analyzed separately. In particular a Gram composition with the adjoint theta operator is a different construction; it can detect the first image and is not ruled out by this calculation.

The recovered source uses the bilinear convention \(\int\vartheta(g,h)f(g)\,dg\) for \nolinkurl{liftT}, as stated in \nolinkurl{Development/AbelianVarietyTheory/Definitions/TwistedInnerProducts/WittSmallGroupEmbeddingMemPointwiseStabilizerIso.lean}, \nolinkurl{liftT_apply}, lines 160--167. It therefore tests the contragredient of the input type against the source oscillator. The preceding trivial-place argument also proves vanishing for this exact source convention at the specified unchanged edge. It must be used when comparing a general chain's nontrivial characters and coefficients; one cannot transfer the conjugated-input matching rule unchanged.

\subsection{4. A quantitative necessary shift, without asserting a global repair}\label{a-quantitative-necessary-shift-without-asserting-a-global-repair}

If the second edge keeps its Hom geometry but has target character \(\chi_T(z)=(z/\bar z)^{m'/2}\), in the source character convention, its source vacuum exponent is \(c=(-m'-4)/2\). Its source dominant weights are \((c-\lambda_2,c-\lambda_1)\). A given dominant weight \((r,s)\), \(r\ge s\), occurs in this polynomial model if and only if \(c\ge r\): choose \(\lambda=(c-s,c-r)\). Thus the types (DO23.9) require respectively

\[
                             m'_1\le-8,\qquad
                             m'_2\le-4,\qquad
                             m'_3\le-10.
 \tag{DO23.12}
\]

These are exact local occurrence conditions for this source module. They are not a construction of a global automorphic splitting character with those infinity types, and do not prove finite occurrence or global nonvanishing for the repaired second lift. A global character would also have to satisfy rational-unit and finite-level compatibility. The unchanged even-rank trivial characters have \(m'_i=0\) and fail all three inequalities.

\textbf{Correction record.} An earlier version identified Konno's \(\xi'\) with the source label \(\chi_T\), overlooking the inverse in \nolinkurl{splitFacO}. It consequently had \((m'-b)/2\) in (DO23.4) and the reversed conditions in (DO23.12). Equations (DO23.4), (DO23.5a), and (DO23.12) now use the source convention consistently. The actual first target types and the unchanged \(\chi_T=1\) second-edge vanishing were unaffected by that error.

\subsection{Primary reference}\label{primary-reference}

Konno, T., \& Konno, K. (2007). On doubling construction for real unitary dual pairs. \emph{Kyushu Journal of Mathematics, 61}(1), 35--82. \url{https://doi.org/10.2206/kyushujm.61.35} Primary PDF: \url{https://www.jstage.jst.go.jp/article/kyushujm/61/1/61_1_35/_pdf} Exact locators: §2.1, printed pp.37--38, Heisenberg and covariance convention; §3.1, pp.43--44, forms and equation (3.1); §3.3, pp.47--48, honest splittings; §4.1, p.49, character exponents; Proposition 4.4(ii), pp.68--69, exact Fock operators; Lemma 5.2(i), p.73, compact action. The paper's displayed last page is 82; an external metadata entry listing 85 is not used for the page range.

Carried first-lift construction: ALP-LIFT-09, especially (V7)--(V11), (V18)--(V26), and (V32), with the independent whole-proof check already recorded in \nolinkurl{definite_oscillator_CHECK_EVENTS.tsv}.

\clearpage

\section{ALP-LIFT-11: a fixed-conductor determinant repair of the specified second-edge type obstruction}\label{alp-lift-11-a-fixed-conductor-determinant-repair-of-the-specified-second-edge-type-obstruction}

Author: /root/lift\_accounting, GPT-6 Astra Pro, ALP\_MACHINE. 21 November 2026 UTC. This result stays over \(E=\mathbf Q(\zeta_7)\) and \(F=E^+\). It is not a common-host identification with the \(\mathbf Q(\zeta_{28})\) family of AM-DO-022.

AM-DO-023 supplies the primary-source-checked honest-splitting Fock formula and the specified second-edge vanishing theorem. They are used here as established results. The new construction is a global splitting character of uniformly bounded finite conductor, its exact effect on the actual first theta lift, and the resulting removal of that particular real-type obstruction. It does not assert nonvanishing of the next global lift.

\subsection{1. An explicit global norm-one character of any integral infinity type}\label{an-explicit-global-norm-one-character-of-any-integral-infinity-type}

Put \(T=U(1)(E/F)\) and \(K_{T,0}=(\widehat{\mathcal O}_E)^1\). ALP-LIFT-09 proves \[
 T(\mathbb A_{F,f})=T(F)K_{T,0},\qquad
 T(F)\cap K_{T,0}=\mu_{14}.
 \tag{R1}
\] In particular \[
 T(F)\backslash T(\mathbb A_F)
       \simeq(T_\infty\times K_{T,0})/\mu_{14},
 \tag{R2}
\] where the roots of unity are embedded diagonally. These arithmetic statements are used here without re-proving their Minkowski and unit-signature arguments.

Choose the three complex coordinates indexed by \(\Phi=\{1,2,3\}\), so that the rational root \(\zeta\) has archimedean coordinates \((\zeta,\zeta^2,\zeta^3)\). For any integer vector \(\mathbf c=(c_1,c_2,c_3)\), put \[
 \beta_{\mathbf c}(z_1,z_2,z_3)=
               z_1^{c_1}z_2^{c_2}z_3^{c_3},\qquad
 C_0=c_1+c_2+c_3,\quad S_0=c_1+2c_2+3c_3 .
 \tag{R3}
\] Here each \(z_i\) lies on the unit circle.

Choose one of the three split places \(v_*\) over29 and its first \(E\) factor so that \(\zeta\) reduces to \(16\) in \(\mathbf F_{29}^{\times}\). This is possible because16 has order7 modulo29. The element2 generates \(\mathbf F_{29}^{\times}\): \(2^{14}=-1\) and \(2^4=16\) has order7. The projection of \(K_{T,0,v_*}\) onto that first factor, followed by reduction, gives a surjection \[
 K_{T,0,v_*}\longrightarrow \mathbf F_{29}^{\times},
 \qquad \zeta\longmapsto 2^4,\quad -1\longmapsto2^{14}.
 \tag{R4}
\] It is injective on the rational subgroup \(\mu_{14}\).

Choose \(r\in\{0,\ldots,13\}\) with \[
                   r\equiv 7C_0+6S_0\pmod{14}.
 \tag{R5}
\] Define the finite compact character \(\rho_{\mathbf c}\) to be trivial at every place other than \(v_*\) and to satisfy \[
            \rho_{\mathbf c}(2)=\exp(2\pi i r/28)
 \tag{R6}
\] on the residue group at \(v_*\). On \(\zeta\) its value is \(\zeta^r\), while \(\beta_{\mathbf c}(\zeta)=\zeta^{S_0}\). Equation (R5) gives \(r\equiv-S_0\pmod7\). On \(-1\) the two values are \((-1)^r\) and \((-1)^{C_0}\), and \(r\equiv C_0\pmod2\). Thus \[
       \rho_{\mathbf c}|_{\mu_{14}}
                    =\beta_{\mathbf c}^{-1}|_{\mu_{14}}.
 \tag{R7}
\] The product \(\beta_{\mathbf c}\rho_{\mathbf c}\) therefore descends through (R2) to a continuous unitary automorphic character \(\gamma_{\mathbf c}\) of \(T\).

This construction is explicit on arbitrary ideles. Write \(t_f=\delta k\) with \(\delta\in T(F)\) and \(k\in K_{T,0}\), as in (R1), and set \[
 \gamma_{\mathbf c}(t_\infty,t_f)
      =\beta_{\mathbf c}(\delta_\infty^{-1}t_\infty)
                         \rho_{\mathbf c}(k).
 \tag{R8}
\] If \(\delta\) is changed by \(\epsilon\in\mu_{14}\), the compact factor is changed by \(\epsilon^{-1}\). Equation (R7) makes the product unchanged. Multiplicativity follows because the torus is abelian, and the same formula shows triviality on rational diagonal points.

Its infinity type is exactly \(\mathbf c\). Its finite restriction to the maximal compact has conductor dividing the first power of the fixed place \(v_*\), and all those compact values are roots of unity of order dividing28. In particular \[
 \gamma_{\mathbf c,f}|_{K_M}=1
                  \qquad\text{for every }M=29^e,\ e\ge1,
 \tag{R9}
\] where \(K_M\) is ALP-LIFT-10's source torus subgroup. The finite conductor bound is independent of the size of \(\mathbf c\); the unramified values at noncompact split places are not asserted to be roots of unity.

Every finite value is algebraic. In (R8), the finite character value is a root of unity, and \(\beta_{\mathbf c}(\delta_\infty)\) is a product of integer powers of conjugates of the algebraic number \(\delta\). For a fixed \(\mathbf c\) these values belong to a fixed finite coefficient field containing the normal closure of \(E\) and \(\mu_{28}\).

At a nonsplit finite place \(v\), every norm-one element of \(E_v^\times\) is a unit, since its valuation has norm valuation zero. Hence \(T(F_v)=K_{T,0,v}\). Our character is trivial on that entire group, not merely on a smaller unit subgroup: \[
                 \gamma_{\mathbf c,v}=1
                         \quad(v\text{ nonsplit finite}).
 \tag{R10}
\] This fact removes a potentially hidden ramified splitting factor in the next step.

\subsection{2. Exact change of the actual first splitting and theta integral}\label{exact-change-of-the-actual-first-splitting-and-theta-integral}

Let \[
 j:\mathbb A_E^\times\longrightarrow T(\mathbb A_F),
                         \qquad j(x)=x/\bar x .
 \tag{R11}
\] Keep the actual first pair of ALP-LIFT-09--10: the positive group \(G=U(I_2)\) and the skew line torus \(T\), with splitting characters \(\chi_V=1\) and \(\chi_W=\mu_E\). Replace only the latter by \[
              \chi_W'=\mu_E\,(\gamma_{\mathbf c}\circ j),
              \qquad\chi_V'=1 .
 \tag{R12}
\] The restriction of the multiplying factor to \(\mathbb A_F^\times\) is1, so the required restriction \(\chi_W'|_{\mathbb A_F^\times}=\epsilon_{E/F}\) is unchanged. Rational diagonal triviality follows from that of \(\gamma_{\mathbf c}\). Thus (R12) is an actual globally admissible splitting character, not just prescribed archimedean weights.

For the standard splitting used in these results, its actual Weil operators satisfy \[
 \boxed{\quad
 \omega'(g,t)=\gamma_{\mathbf c}(\det g)\,\omega(g,t).
 \quad}
 \tag{R13}
\] Here are all local checks of this identity. At every nonsplit finite place, (R10) says that \(\gamma_{\mathbf c,v}\circ j\) is literally the trivial character of \(E_v^\times\), so the local splitting has not changed. At every split finite place, the exact determinant root formula of AM-DO-020 evaluates the ratio on an element \(y\) with \(y/\bar y=\det g\); its value is precisely \(\gamma_{\mathbf c,v}(\det g)\). This includes possible unramified values at split places outside29. At each real place, the character \(\gamma_{\mathbf c}\circ j\) changes the exponent \(m_i\) of AM-DO-023 by \(2c_i\). Its exact Fock formula therefore changes the target determinant power by \(c_i\), giving the same scalar. The source action depends on the unchanged opposite character \(\chi_V\). These checks prove (R13) at every place and hence globally.

For any of the actual finite Schwartz inputs and Fock vectors already constructed, the theta series and first theta integral consequently satisfy \[
 \vartheta'_\phi(g,t)
       =\gamma_{\mathbf c}(\det g)\vartheta_\phi(g,t),
 \qquad
 \Theta'(\phi,f)(g)
       =\gamma_{\mathbf c}(\det g)\Theta(\phi,f)(g).
 \tag{R14}
\] The source torus input \(f\) is unchanged. In particular all identity values, including their retained common period and their 11-adic normalized valuations, are unchanged exactly. The twisted first lift is nonzero whenever the original one is nonzero. Since the character has modulus one, the scalar automorphic \(L^2\) norm and supremum norm of the entire first lift are unchanged as well.

\subsection{3. Removing the specified real-type obstruction in each integral convention}\label{removing-the-specified-real-type-obstruction-in-each-integral-convention}

For a first Fock vector of fixed CM multidegree \(\mathbf d=(d_1,d_2,d_3)\), put \[
                         \boldsymbol\nu=(0,0,1).
 \tag{R15}
\] AM-DO-023 proves that the actual first target compact types are \[
              \pi_i=\det^{\nu_i}\operatorname{Sym}^{d_i}
                               (\mathbf C^2),
 \qquad\operatorname{HW}(\pi_i)=(d_i+\nu_i,\nu_i).
 \tag{R16}
\] The specified next positive Hom edge \(U(2)\longrightarrow U(4)\) with both splitting characters trivial has source types \[
 \det^{-2}S_\lambda((\mathbf C^2)^*),\qquad
 \operatorname{HW}=(-2-\lambda_2,-2-\lambda_1),
 \quad\lambda_1\ge\lambda_2\ge0.
 \tag{R17}
\] In particular the untwisted types (R16) do not occur. This is the already proved scoped vanishing of AM-DO-023, not a universal no-go for every chain.

\subsubsection{The conjugated-input integral}\label{the-conjugated-input-integral}

First use the Hermitian theta integral \(\int\vartheta(g,h)\overline{f(g)}\,dg\), as in ALP-LIFT-09. Its source compact type must match that of \(f\) itself.

Choose in the global construction of section1 \[
                         c_i=-d_i-2-\nu_i .
 \tag{R18}
\] By (R13)--(R14), the actual twisted first image now has dominant weights \[
 (d_i+\nu_i+c_i,\nu_i+c_i)=(-2,-d_i-2).
 \tag{R19}
\] These are exactly the weights in (R17) for the one-row partition \(\lambda=(d_i,0)\). Thus the twisted first compact representation occurs in the second oscillator, with the same degree \(d_i\) at each place. The repair uses a genuine global splitting of finite conductor bounded independently of \(\mathbf d\).

For the explicit nonzero first vector of ALP-LIFT-09, \(\mathbf d=(2,0,2)\) and \(\boldsymbol\nu=(0,0,1)\). One takes \[
 \mathbf c=(-4,-2,-5),\qquad C_0=-11,\quad S_0=-23,
 \qquad r=9 .
 \tag{R20}
\] The residue character at29 is \(\rho(2)=\exp(2\pi i\,9/28)\), and the repaired first weights are \[
                       (-2,-4),\quad(-2,-2),\quad(-2,-4).
 \tag{R21}
\] All these occur in the specified second source model.

A single fixed repair also works for any finite prefix of degrees. Given bounds \(D_i\) and \(0\le d_i\le D_i\), use \(c_i=-D_i-2-\nu_i\). The resulting source weights are \((d_i-D_i-2,-D_i-2)\), which occur for \[
                    \lambda_i=(D_i,D_i-d_i).
 \tag{R22}
\] The next Fock degree is then \(2D_i-d_i\). This variant changes the needed polynomial degree and any corresponding archimedean coefficient norm; it is not a zero-cost identification of the next Fock vectors.

Conversely, for this conjugated-input composition with two fixed positive Hom edges and fixed splitting characters, the second source spectrum has a fixed upper bound on its highest-weight entries. Equation (R16) then excludes sufficiently large \(d_i\). The degree-dependent repair or a finite-prefix repair changes that datum. This observation concerns the displayed orientation and character choices; it does not apply without checking the alternating signs, contragredients and pins in a different proposed source chain.

\subsubsection{The source's literal bilinear integral has a fixed repair}\label{the-sources-literal-bilinear-integral-has-a-fixed-repair}

The recovered source liftT instead uses \[
                   \int_{[G]}\vartheta(g,h)f(g)\,dg .
 \tag{R23}
\] Its local matching condition pairs the oscillator source type with the contragredient of the input type. In this convention a stronger, degree-independent repair is available. Choose \[
              \mathbf c=\mathbf 2-\boldsymbol\nu=(2,2,1).
 \tag{R24}
\] Then the actual first image has weights \((d_i+2,2)\), whose contragredient weights are \[
                        (-2,-d_i-2).
 \tag{R25}
\] They are precisely (R17) for \(\lambda=(d_i,0)\), for every nonnegative degree \(d_i\) simultaneously. The next Fock degree stays \(d_i\). Thus neither a degree-dependent splitting nor a finite-prefix enlargement of the next polynomial degree is needed for the bilinear convention.

Here \(C_0=5\), \(S_0=9\), and (R5) gives \(r=5\). The fixed finite character is \(\rho(2)=\exp(2\pi i\,5/28)\), and \[
       (m_1,m_2,m_3)+2\mathbf c=(3,3,3).
 \tag{R26}
\] Consequently one actual fixed character \(\chi_W'=\mu_E(\gamma_{(2,2,1)}\circ j)\) repairs all of these real types in the literal bilinear composition. This is an independent convention-specific conclusion; it does not turn the two different theta integral conventions into the same map.

Equivalently, (R23) is the conjugated-input lift of \(\bar f\), whose compact type is the contragredient. This reformulation is useful when applying an inner-product theorem written with the conjugated-input convention.

\subsection{4. Integral class values and the actual cost scope}\label{integral-class-values-and-the-actual-cost-scope}

The identity values and scalar automorphic norms are preserved by (R14) without choosing representatives. An integral class-value claim requires a specified finite target level and its coefficient lattice. AM-R-008 proves the relevant precise statement: with compact-determinant representatives and a character trivial on \(\det K\), multiplication by \(\gamma_f(\det g)\) is a diagonal root-of-unity map. It preserves the entire local content and the mass norm, with the coefficient action twisted by \(\gamma_\infty^{-1}\) in that note's left-equivariant convention. This inverse is consistent with the positive right-translation twist in (R14); the two statements use their respective explicit function conventions.

For the explicit group level of AM-R-008, (R9) verifies its character hypothesis. There is consequently no additional denominator, content or logarithmic norm cost for the determinant twist on that class module. All its displayed finite class scalars have order dividing28.

The source torus stabilizer \(K_M\) in ALP-LIFT-10 must not be confused with a proved target \(U(2)\) stabilizer. This note does not assert that the original first vectors are invariant under the particular target group level of AM-R-008 without checking that target action. Every finite family of their Schwartz vectors has an open target stabilizer; it can be refined inside the kernel of \(\gamma_f\circ\det\). Applying a compatible compact-determinant class presentation then transports its actual class-value module by the same root-of-unity diagonal map. Its original stabilizer index, mass and coefficient lattice remain costs to be computed if needed.

Likewise this does not identify the transported module with an independently pinned arithmetic row. The first edge has ranks1 and2 and already lies outside the source's even-rank equal-character StdPair constructor. The later PinX.mv and PinColT.om interfaces require the same specified character on both members of an even-rank edge. Changing only \(\chi_W\) in (R12) is a valid actual first theta lift, but is not a verification of those different equal-character pins.

\subsection{5. What remains for a second actual lift}\label{what-remains-for-a-second-actual-lift}

The repaired first input has exactly the real representation required by the chosen second oscillator, and it retains the actual nonzero first value and first norm. Local occurrence at infinity is a necessary condition, not a proof that its global next theta integral is nonzero. One must still select actual finite second-kernel inputs, prove a nonzero global pairing, and retain the resulting mass and norm factors. A formal Gram picture or a hyperbolic Witt-tower persistence theorem does not establish this for a sequence of positive definite groups.

The \(\mathbf Q(\zeta_{28})\) host, the full DP65 conjugated-action rows, source tied direction matrices, and the terminal six-identity family remain separate comparisons. This construction removes a specific real-type obstruction without claiming the complete requested chain.

\subsection{Exact source locators}\label{exact-source-locators}

\begin{itemize}
\tightlist
\item
  ALP-LIFT-09, actual\_torus\_first\_lift.md, (V2)--(V11), torus class number and the original global splitting, and (V32)--(V38), actual first lift and projectors.
\item
  ALP-LIFT-10, permitted\_orbit\_first\_image.md, (O21), (O23)--(O28), (O31)--(O35), the exact fixed source levels, normalized values and input costs.
\item
  AM-DO-020, exact split-place determinant-root covariance; AM-DO-023, definite\_oscillator\_real\_second\_edge.md, (DO23.4)--(DO23.11), actual honest real target weights and the trivial-source-character case. Its corrected nontrivial source-character sign is not needed here.
\item
  Konno, T., \& Konno, K. (2007). On doubling construction for real unitary dual pairs. Kyushu Journal of Mathematics, 61(1), 35--82. Proposition4.4(ii), pp.68--69, and Lemma5.2(i), p.73. \url{https://www.jstage.jst.go.jp/article/kyushujm/61/1/61_1_35/_pdf}
\item
  AM-R-008, compact\_determinant\_twist\_accounting.md, (R8.2), (R8.10)--(R8.15), the actual integral class-value and mass isometry with its stated level hypothesis.
\item
  The actual source bilinear convention is \nolinkurl{WittSmallGroupEmbeddingMemPointwiseStabilizerIso.lean,} liftT\_apply, lines160--167, as directly checked in AM-DO-023/024's source audit.
\end{itemize}

\subsection{Check status}\label{check-status-2}

The exact finite-residue and weight calculations are included in \nolinkurl{fixed_conductor_type_repair_certificate.py.} Author and independent whole-proof checks, with their actual scope and time, are recorded separately in \nolinkurl{lift_accounting_CHECK_EVENTS.tsv.} No Lean execution is claimed.

\clearpage

\section{AM-K019 --- what the full theta differential module cannot supply}\label{am-k019-what-the-full-theta-differential-module-cannot-supply}

Author: \nolinkurl{/root/kummer_extension}, GPT-6 Astra Pro, ALP\_MACHINE, 21 November 2026 UTC. This is a new, elementary obstruction to one proposed iteration argument. It does not obstruct a selected harmonic functional, an evaluated jet ideal, or a quotient with separately controlled saturation. The accepted source connection-integrality theorem is used as a theorem.

\subsection{1. Exact statement}\label{exact-statement}

Let \(J=\operatorname{Jac}(y^2=x^7-1)\), let \(u_1,u_2,u_3\) be the invariant logarithmic coordinates of AM-K012, and let \(f\) be either the actual translated, value-one reduced theta function \(F\), or the actual trace-one reduced section \(\sigma_q\). On the complement of its zero divisor put \[
z_i=\partial_i\log f.
\tag{K20.1}
\] The following assertions hold.

\begin{enumerate}
\def\labelenumi{\arabic{enumi}.}
\tightlist
\item
  The six meromorphic functions \(u_1,u_2,u_3,z_1,z_2,z_3\) are algebraically independent over the field \(\mathbb C(J)\) of abelian functions. In particular, adjoining only the three logarithmic coordinate multipliers does not put \(f\) in a finite-dimensional vector space over \(\mathbb C(J)(u_1,u_2,u_3)\) preserved by ordinary invariant differentiation.
\item
  Adjoining the three \(z_i\) does give the expected polynomial rank-one connection module. In its natural integral polynomial lattice, however, every ordinary differentiated vector \(\partial_i^N f/f\) has a monic term \(z_i^N\). No positive restricted power of that connection maps the whole lattice into \(p\) times the lattice, and no lattice commensurable with it by a bounded power of \(p\) can have that property.
\item
  The source's actual integral connection on its auxiliary-variable module has an analogous obstruction: its creation term produces a unit coefficient of auxiliary degree \(N\) after \(N\) iterations. At prime-to-\(p\) tensor power and in a primitive tangent direction, a positive restricted power cannot map that whole integral module into \(p\) times itself, even after a bounded-denominator commensurable change of lattice.
\end{enumerate}

The distinction in assertion 2 is essential. The abstract polynomial lattice injects into the characteristic-zero evaluated Hurwitz series, but its reduction modulo \(p\) need not inject. In the present translated value-one example, each \(z_i\) has zero constant term, so \(z_i^p=0\) in the reduced Hurwitz ring. Positive \(p\)-adic bounds on evaluated derivatives therefore coexist with the polynomial-lattice obstruction.

\subsection{2. The lattice shifts span all six complex directions}\label{the-lattice-shifts-span-all-six-complex-directions}

Write \(J(\mathbb C)=V/\Lambda\), where \(V=\mathbb C^3\). Its principal polarization is represented by the nondegenerate positive Hermitian form \(H\), linear in the first variable. The actual reduced theta multiplier, checked in AM-K018, is \[
f(u+\lambda)=\alpha(\lambda)
 \exp\left(\pi H(u,\lambda)+\frac\pi2H(\lambda,\lambda)\right)f(u).
\tag{K20.2}
\] Consequently \[
z_i(u+\lambda)=z_i(u)+\pi H(e_i,\lambda).
\tag{K20.3}
\] If a translated representative changes the semicharacter or a constant factor, differentiation removes that change, so (K20.3) is identical for \(F\) and \(\sigma_q\).

Let \(C\) be the nonsingular complex matrix with \((\pi H(e_i,\lambda))_i=C\overline\lambda\). The shift subgroup of the six displayed functions is \[
\Gamma=\{(\lambda,C\overline\lambda):\lambda\in\Lambda\}
 \subset\mathbb C^6.
\tag{K20.4}
\] Choose a real lattice basis \(\lambda_1,\ldots,\lambda_6\). The six vectors \((\lambda_j,C\overline{\lambda_j})\) are complex-linearly independent. Indeed, a complex-linear functional \((a,b)\) vanishing on them vanishes on their real span. Since the \(\lambda_j\) span \(\mathbb C^3\) over \(\mathbb R\), this gives \[
a\lambda+bC\overline\lambda=0\qquad(\lambda\in\mathbb C^3).
\] Replacing \(\lambda\) by \(i\lambda\) and combining the two equations gives \(a\lambda=0\) and \(bC\overline\lambda=0\) for every \(\lambda\). Thus \(a=b=0\). The six-column square matrix has full complex rank.

It follows that \(\Gamma\) is Zariski dense in \(\mathbb C^6\): after an invertible complex linear transformation this is the familiar assertion that \(\mathbb Z^6\) is Zariski dense. The latter follows by induction on the number of variables, applying the fact that a nonzero univariate polynomial has only finitely many roots.

Suppose a polynomial relation among \(u,z\) has coefficients in \(\mathbb C(J)\). Its pullback is an identity of meromorphic functions on \(V\). Choose \(u_0\) avoiding the finitely many coefficient poles and the zero divisor of \(f\). Periodicity of the coefficient functions and of the theta divisor permits substitution of every \(u_0+\lambda\). The polynomial with coefficients evaluated at \(J(u_0)\) vanishes on \((u_0,z(u_0))+\Gamma\), so Zariski density makes every one of those coefficients zero. Varying \(u_0\) on a nonempty open set makes every original rational coefficient zero. This proves assertion 1's algebraic-independence claim.

The same independence holds after any algebraic extension of the abelian-function field, for example the finite ordered-divisor cover used in AM-K012. Algebraic extension does not change the transcendence degree of the composite field. Thus the explicit Abel-coordinate ring used for integral calculations causes no exception to the argument.

\subsection{\texorpdfstring{3. Monic derivatives exclude finite rank over abelian functions and \(u\)}{3. Monic derivatives exclude finite rank over abelian functions and u}}\label{monic-derivatives-exclude-finite-rank-over-abelian-functions-and-u}

Put \(K=\mathbb C(J)\). The Hessians \[
\partial_i z_j=-W_{ij}
\tag{K20.5}
\] are abelian functions. For \(F\) they are the translated Kleinian functions of AM-K012; for \(\sigma_q\) they are the corresponding unshifted ones. In either case \(D_iK\subset K\).

Fix \(i\). Recursively set \[
P_0(Z)=1,\qquad P_{N+1}(Z)=D_iP_N(Z)+ZP_N(Z),
\quad D_iZ=-W_{ii}.
\tag{K20.6}
\] Then \[
\frac{\partial_i^N f}{f}=P_N(z_i),\qquad
P_N(Z)=Z^N+\text{terms of degree below }N,
\quad P_N\in K[Z].
\tag{K20.7}
\] Both statements follow by induction. Differentiating a coefficient preserves \(K\); differentiating \(Z^r\) lowers its \(Z\)-degree; multiplication by \(Z\) produces the next unique leading term with coefficient one.

The \(P_N(z_i)\) are linearly independent over \(K(u_1,u_2,u_3)\), since \(z_i\) is transcendental over that field and their degrees are strictly increasing. A finite-dimensional vector space over that field, consisting of actual meromorphic germs and containing \(f\), cannot be preserved by \(\partial_i\): otherwise the infinite sequence \(f,\partial_i f,\partial_i^2 f,\ldots\) would be linearly dependent. This proves the stated finite-rank obstruction. It concerns ordinary differentiation on the actual reduced functions. It does not exclude a different connection on an auxiliary bundle, or assert that no rank-one connection exists after adjoining the \(z_i\).

For algebraic constant-field variants, this conclusion persists on extension from a number field to a characteristic-zero completion. To see the needed elementary descent, clear the finitely many rational denominators in a hypothetical relation. It becomes a finite linear relation, over the extended constant field, among coefficientwise number-field formal germs. Linear independence of any finite collection of such germs is witnessed by a nonzero finite coefficient minor; that minor remains nonzero after constant-field extension. Hence no such new relation appears. This statement is about finitely many rational-function coefficients and does not assert descent of an infinite analytic relation.

\subsection{4. The natural polynomial connection is integral but has no full-module gain}\label{the-natural-polynomial-connection-is-integral-but-has-no-full-module-gain}

Let \(\mathcal O\) be a good local coefficient DVR, with uniformizer \(p\), and use the explicit \(D\)-stable integral ordered-divisor ring \(R\) of AM-K012 near the translated non-effective characteristic. Thus \(D_iR\subset R\), \(W_{ij}\in R\), and all evaluated coefficients lie in the integral Hurwitz ring. A harmless unramified extension splits the branch points; no ramified change of uniformizer is being hidden.

Define \[
A=R[Z_1,Z_2,Z_3],\qquad
\delta_i|_R=D_i,\qquad \delta_i Z_j=-W_{ij}.
\tag{K20.8}
\] The derivations commute because the mixed third derivatives of the actual theta logarithm agree. On the free rank-one \(A\)-module \(M=Ae\) define \[
\nabla_i(ae)=(\delta_i a+Z_i a)e.
\tag{K20.9}
\] Its characteristic-zero realization is \(Z_i\mapsto z_i\), \(e\mapsto F\). The algebraic independence above makes this realization faithful over the fraction field. Formula (K20.7) gives \[
\nabla_i^Ne=P_N(Z_i)e,
\qquad [Z_i^N]P_N=1.
\tag{K20.10}
\] All coefficients are integral; one of them is a unit. Therefore \[
\nabla_i^Q M\not\subset pM\qquad\text{for every }Q\ge1.
\tag{K20.11}
\] This conclusion does not depend on the nilpotence length of the derivation \(D_i\) on \(R/pR\).

Here is the precise bounded-denominator strengthening. Suppose a lattice \(M'\) in \(M[1/p]\) satisfies \[
p^a e\in M',\qquad M'\subset p^{-C}M
\tag{K20.12}
\] for fixed nonnegative integers \(a,C\). These hypotheses hold for any lattice commensurable with \(M\) by a bounded power of \(p\). If, for some \(Q\ge1\), \(\nabla_i^Q M'\subset pM'\), then iteration gives \[
p^aP_{nQ}(Z_i)e=\nabla_i^{nQ}(p^ae)
 \in p^nM'\subset p^{n-C}M.
\tag{K20.13}
\] The coefficient of \(Z_i^{nQ}\) on the left is \(p^a\). For \(n>a+C\) this contradicts (K20.13). Thus no such commensurable lattice supplies the proposed uniform restricted-power divisibility.

\subsubsection{The evaluated Hurwitz image is a different lattice problem}\label{the-evaluated-hurwitz-image-is-a-different-lattice-problem}

For the translated \(F\), \(z_i(0)=0\). Reduction of a Hurwitz series with zero constant term has \(p\)-th power zero: the multinomial coefficients in a \(p\)-fold product cancel modulo \(p\), equivalently the only surviving Frobenius term is the \(p\)-th power of the constant coefficient. Hence \[
\operatorname{ev}(Z_i^p)=z_i^p=0
 \quad\text{in the characteristic-}p\text{ Hurwitz ring},
\tag{K20.14}
\] whereas \(Z_i^p\ne0\) in \(A/pA\). The reduction of evaluation has a nontrivial kernel even though characteristic-zero evaluation is injective. It also admits additional integral elements after saturation in its characteristic-zero image; for instance the image of \(Z_i^p/p\) is integral although the abstract polynomial is not in \(A\).

Consequently (K20.11) does \textbf{not} refute the actual congruences of AM-K017, nor the actual all-order lower bounds of AM-K013--015. It refutes the attempt to infer those bounds by a uniform divisibility operator on this unsaturated polynomial lattice, or on a lattice differing from it by a fixed power of \(p\). An iterative argument in the evaluated image must specify that image, its \(p\)-saturation, and the growth of the saturation cost. The isolated mod-\(p\) theta identities supply none of those bounds by themselves.

\subsection{5. The same issue is visible in the actual source connection}\label{the-same-issue-is-visible-in-the-actual-source-connection}

The source definition is \nolinkurl{ReciprocityTheory/ReductiveGroups/ArithmeticGeometry/TaylorDiscriminants/ConnectionKillsSeries.lean}, \nolinkurl{sgOf} at 18--19, \nolinkurl{sg1} at 21--22, \nolinkurl{nabla} at 31--35, \nolinkurl{Hset} at 41--43, and \nolinkurl{ConnectionIntegrality} at 77--83. The accepted theorem realizing the latter is \nolinkurl{WickVanishing/ExpConjConnectionIntegral.lean}, lines 40--82. These exact files were inspected; no rederivation of their established integral-connection assertion is made.

Use \(X\) for its formal-group variables and \(Z=(Z_1,\ldots,Z_d)\) for its auxiliary variables. Put \(s_i(X)=-\partial_{Y_i}h(X,0)\). The source operator is \[
\nabla_{u,\phi}
 =D_u^X+\sum_i(D_us_i+\phi_i)\partial_{Z_i}
  -\kappa\sum_i(Z_i-s_i)u_i.
\tag{K20.15}
\] Its admitted integral lattice of directions has \(\phi=\phi_0-Su\) with \(u,\phi_0\) integral and \(S=ds(0)\). Conjugating by the source factor \(\exp(\kappa h_L(X))\), which depends only on \(X\), gives an operator of the form \[
T=e^{-\kappa h_L}\nabla_{u,\phi}e^{\kappa h_L}
 =D_u^X+\sum_i A_i(X)\partial_{Z_i}
  -\kappa\sum_i u_iZ_i+b(X).
\tag{K20.16}
\] The accepted source theorem states that \(T\) preserves the integral formal-power-series module \(\mathcal N=\mathcal O[[X,Z]]\). Equivalently the coefficients in this particular first-order expression are integral; this also follows directly by applying preservation to \(1,X_j,Z_i\). For the argument below, preservation itself is sufficient.

Assume \(p\nmid\kappa\) and \(u\) is primitive, so some \(u_j\) is a unit. Starting at \(1\), \(T^N1\) is polynomial in \(Z\) of degree at most \(N\), with coefficients in \(\mathcal O[[X]]\). Every term of \(T\) except multiplication by \(-\kappa u\cdot Z\) preserves or lowers \(Z\)-degree. Induction therefore gives the exact highest homogeneous part \[
\operatorname{top}_Z(T^N1)=(-\kappa u\cdot Z)^N.
\tag{K20.17}
\] Its coefficient of \(Z_j^N\) is \((-\kappa u_j)^N\), a unit. Hence \[
T^Q\mathcal N\not\subset p\mathcal N
 \qquad(Q\ge1).
\tag{K20.18}
\] The same argument as (K20.12)--(K20.13), replacing \(e\) by \(1\), rules out a bounded-denominator commensurable lattice with such a uniform gain. The source's integral connection preserves an integral lattice; it does not make a positive power of each creation operator divisible by \(p\) on the whole lattice.

This is fully compatible with the source \nolinkurl{ValueSt}, lines 71--75: after the specified constant-coefficient evaluation, the relevant connection polynomial gives an ordinary invariant derivative. That evaluation annihilates the positive auxiliary powers responsible for (K20.17). Harmonic contractions and other selected quotients can likewise remove them. Any stronger gain must be proved for the actual selected evaluation or quotient, including its saturation and normalization, rather than deduced from full-module nilpotence that is false.

\subsection{6. Consequence for the current proof-transfer attempt}\label{consequence-for-the-current-proof-transfer-attempt}

AM-K018 identifies the full-\(\mathcal O_E\) source gauge with the actual Chern/Baker gauge on ordinary powers of \(J\). The obstruction above does not undo that identification. It specifies why the full mixed mod-11 congruences of AM-K017 cannot automatically be iterated through a finite-rank ordinary-differential module, and why even the actual source integral connection does not supply the proposed uniform restricted-power inclusion on its full auxiliary module.

The viable remaining arithmetic question is therefore more precise: determine the actual evaluated or harmonic-selected theta image, bound its \(p\)-saturation in the required degree, and compare that bound with the projection and normalization costs. The accepted DP54/58/59 examples now apply in the exact Chern gauge at 11; DP45/53 at the order-six Frobenius classes remain bare integral-function inputs until their actual primitive theta kernel or selected-value interface is supplied.

Primary source for (K20.2): Bannai, K., \& Kobayashi, S. (2010). Algebraic theta functions and the \(p\)-adic interpolation of Eisenstein--Kronecker numbers. \emph{Duke Mathematical Journal, 153}(2), 229--295, Definition 1.8 in \href{https://arxiv.org/pdf/math/0610163}{arXiv:math/0610163v4}. The actual multiplier and source compatibility were proved in AM-K018; the arguments in §§2--5 above are new elementary arguments, with the accepted source integral-connection theorem explicitly separated.

\clearpage

\section{AM-K020 --- two sharp pure directions of the actual trace-one theta section}\label{am-k020-two-sharp-pure-directions-of-the-actual-trace-one-theta-section}

Author: /root/kummer\_extension, GPT-6 Astra Pro, ALP\_MACHINE, 21 November 2026 UTC. The full actual source gauge was identified in AM-K018. The supplied DP41 bounded-multiplier lemma and DP45/53 bare cyclic-display bounds are used as accepted theorems. The new result concerns the single fixed actual section, not an arbitrary integral test function.

\subsection{1. Statement}\label{statement-11}

Let \(J=\operatorname{Jac}(y^2=x^7-1)\), and use the global invariant differential basis \(\omega_i=x^{i-1}dx/y\) and its logarithmic coordinates \(u_i\). Let \(\sigma_q\) be the actual reduced theta germ of the trace-one principal line \(L_{2\infty}\), normalized by \[
\sigma_q(u)=u_1u_3-u_2^2+\text{higher terms}.
\] Its complete algebraic normalization, actual Chern/source gauge, and integral content are AM-K014, AM-K018, and ALP-LIFT-08. Write \(v_p(p)=1\), with \(v_p(0)=+\infty\).

The following are exact sharp directional statements: \[
\boxed{\begin{aligned}
v_3\!\left(\partial_1^n\sigma_q(0)\right)
 &\ge \max\left\{0,\frac n{56}-\frac{s_3(n)}2\right\},
&
\liminf_{n\to\infty}
\frac{v_3(\partial_1^n\sigma_q(0))}{n}
 &=\frac1{56},\\[2mm]
v_5\!\left(\partial_2^n\sigma_q(0)\right)
 &\ge \max\left\{0,\frac n{504}-\frac{s_5(n)}4\right\},
&
\liminf_{n\to\infty}
\frac{v_5(\partial_2^n\sigma_q(0))}{n}
 &=\frac1{504}.
\end{aligned}}
\tag{K21.1}
\] Here \(s_p(n)\) is the base-\(p\) digit sum. These are derivatives along two precisely specified global CM eigenlines, on \(J\) itself. They require no conjugated factor and no comparison of mixed-action product gauges.

There are also explicit sharp witness sequences, not just a nonconstructive liminf. Put \[
n_3(r)=28\cdot3^r+18,\qquad n_5(r)=126\cdot5^r+30.
\] Then \[
\boxed{\begin{aligned}
v_3(\partial_1^{n_3(r)}\sigma_q(0))
 &=\frac{3^r+3-s_3(n_3(r))}{2}\quad(r\ge0),\\
v_5(\partial_2^{n_5(r)}\sigma_q(0))
 &=\frac{5^r-1}{4}\quad(r\ge0).
\end{aligned}}
\tag{K21.1a}
\] For \(r\ge3\), the first expression is \((3^r-1)/2\). Its values for \(r=0,1,2\) are \(0,1,5\); the last has one extra unit of valuation from a base-three carry. Every derivative in (K21.1a) is nonzero.

The same sharp slopes hold for the translated, value-one theta \(F\) in those directions. For \(\sigma_q\) they persist under every permitted pullback with the source normalization \(M^{-n}\), since its reduced pullback is \(\sigma_q(Mu)\) and the two factors cancel exactly. In particular the previously used odd levels \(M=29^e\) are prime to both 3 and 5.

The accepted DP45/53 bare integral-function slopes in these directions are respectively \(3/56\) and \(5/504\). Thus these actual normalized theta derivatives have slope equal to \(1/p\) times that bare slope at these two primes. This comparison is an actual matching lower/asymptotic upper result for one fixed section; it is not a conclusion about the total cost of a later automorphic row.

\subsection{2. An exact weighted-initial-form lemma}\label{an-exact-weighted-initial-form-lemma}

Let \(p\) be odd, \(q=p^h\), \(s=1/(p-1)\), and \[
H(t)=\sum_{n\ge2}\ell_n\frac{t^n}{n!}.
\] Assume \[
\ell_n=0\ (n\text{ odd}),\qquad
v(\ell_n)\ge \max\{0,\lfloor(n-2)/q\rfloor\},\qquad
v(\ell_{q+1})=0.
\tag{K21.2}
\] Put \[
c=\frac{s}{q+1},\qquad w_*=s-c,\qquad
I=[s-2c,s-c].
\tag{K21.3}
\] For a rational \(w\in I\), choose a finite ramified coefficient extension and \(\gamma\) with \(v(\gamma)=w\). Write \(\delta=w_*-w\in[0,c]\) and \[
\tau=s-(q+1)\delta\in[0,s].
\tag{K21.4}
\] Then \(H(\gamma T)\) has coefficient infimum exactly \(\tau\), and its entire initial form, after division by an element of valuation \(\tau\), is a nonzero scalar times \(T^{q+1}\). In particular that initial form has nonzero formal derivative in characteristic \(p\).

\textbf{Proof.} For even \(n\), Legendre's formula gives \[
v\!\left(\frac{\ell_n\gamma^n}{n!}\right)
=v(\ell_n)+s\,s_p(n)-(c+\delta)n.
\tag{K21.5}
\] An even positive integer has \(s_p(n)\ge2\). At \(n=q+1\), the digit sum is exactly two, and the unit assumption makes (K21.5) equal \(2s-(c+\delta)(q+1)=\tau\).

If \(2\le n<q+1\), the same expression is at least \(2s-(c+\delta)n\), exceeding \(\tau\) by at least \((c+\delta)(q+1-n)>0\).

If \(n>q+1\) is even, put \(k=\lfloor(n-2)/q\rfloor\ge1\). Then \(n\le(k+1)q+1\), so (K21.5) minus \(\tau\) is at least \[
k-(c+\delta)(n-q-1)
\ge k\{1-q(c+\delta)\}
\ge k\left(1-\frac{2qs}{q+1}\right)>0.
\tag{K21.6}
\] The last inequality holds because \(p\ge3\) implies \(2s\le1\). These estimates also bound all coefficients below and show that the higher-degree excess grows at least linearly in \(k\). Odd coefficients are zero. There is therefore one and only one initial monomial. Its degree is \(q+1\equiv1\pmod p\), proving the derivative assertion. ∎

Let \(\nu_w(G)=\inf_n(v([t^n]G)+wn)\). Suppose \(V(t)\ne0\) is any formal series for which \(V(\gamma T)\) is bounded for every rational \(w\in I\), and integral for \(w=w_*\). The accepted DP41 bounded-multiplier lemma applies to the full initial form just proved, not merely to the existence of a first pole. It gives \[
\nu_w(Ve^H)=-\infty\quad(w<w_*),
\qquad
0\le\nu_w(Ve^H)<\infty\quad(w\ge w_*).
\tag{K21.7}
\] For the first assertion take \(w\in I\), \(w<w_*\), so \(0\le\tau<s\), and use the nonzero-derivative initial form. The multiplier is only required to be nonzero and bounded; a zero at the origin is allowed. Monotonicity of \(\nu_w\) extends unboundedness to every smaller nonnegative weight. At \(w=w_*\), all coefficients of \(H(\gamma T)\) have valuation at least \(s\), so its exponential and its inverse are integral, with constant term one. This gives the asserted boundedness and finite infimum; larger weights follow by monotonicity.

If \(Ve^H=\sum a_nt^n\), (K21.7) implies \[
v(a_n)\ge-w_*n,\qquad
\liminf_{n\to\infty}\frac{v(a_n)}n=-w_*.
\tag{K21.8}
\] Indeed, a strictly larger liminf would make the coefficient infimum finite at some rational weight below \(w_*\), a contradiction. Multiplication by \(n!\) and \(s_p(n)/n\to0\) give \[
v(n!a_n)\ge cn-s\,s_p(n),\qquad
\liminf_{n\to\infty}\frac{v(n!a_n)}n=c.
\tag{K21.9}
\] All assertions also hold with \(H(2t)\) in place of \(H(t)\), since \(p\) is odd and multiplication of its \(n\)-th logarithmic derivative by \(2^n\) changes no valuation or initial-form nonvanishing.

\subsection{3. Apply the lemma to the actual translated theta}\label{apply-the-lemma-to-the-actual-translated-theta}

For the actual value-one \(F\), AM-K012 proves \[
\partial_i^2\log F=-W_{ii}(a+u)
\] in the explicit integral, invariant-derivation-stable Abel-coordinate ring. AM-K013 proves the all-order operator bound \(D_i^mR\subset p^{\lfloor m/q\rfloor}R\) at the indicated nilpotence exponent, and hence exactly the logarithmic bounds (K21.2), except for its final unit condition. This is an iterated assertion on a stable integral ring, not an iteration of one coefficient congruence.

The exceptional constants were computed exactly in the primary algebraic theta recurrence of AM-K017: \[
\begin{array}{c|c|c|c}
p&i&q&\ell_{q+1}\bmod p\\ \hline
3&1&27&2\\
5&2&125&1.
\end{array}
\tag{K21.10}
\] The computations are in \(\mathbb F_p[Z]/\Phi_7(Z)\), which is a field in both cases because \(p\) has order six modulo seven. The residues shown are nonzero scalars, so the logarithmic derivatives are units. The actual trace-one section is not being substituted for \(F\) in the definition of these logarithmic derivatives: \(F(0)=1\), whereas \(\sigma_q(0)=0\).

These exact values and their full algebraic ODE inputs are saved in:

\begin{verbatim}
work/agents/jacobian_hessian_pcurvature_certificate.py
work/agents/jacobian_hessian_pcurvature_certificate.json
work/agents/jacobian_theta_direction_jets.py
\end{verbatim}

The certificate regenerates the values using the branch divisor \((\zeta,\zeta^3,\zeta^5)\), the exact Kleinian Hessian, and division-free Hurwitz recurrences. Its recorded PASS is a finite arithmetic check. The all-order deduction here uses (K21.2) and the complete initial-form lemma above.

With \(V=1\), (K21.9) proves the sharp slopes \(1/56\) and \(1/504\) for the two indicated restrictions of \(F\).

\subsection{4. Transfer to the actual vanishing source section}\label{transfer-to-the-actual-vanishing-source-section}

AM-K015 proved an equality of actual reduced germs \[
\sigma_q(2u)=r(u)F(2u),
\tag{K21.11}
\] where \(r\) is a rational function on \(J\) regular and integral on a formal neighborhood of zero at every odd prime other than 7. This comes from the canonical rigidified \([2]\)-identification of the two theta lines and their actual integral sections; it is not an arbitrary integral formal prefactor. The equality is in the Chern gauge, so no additional quadratic exponential occurs. The value-one normalization of \(F\) is a unit at 3 and 5 by ALP-LIFT-07, and the \(q\)-normalized source section is primitive there by ALP-LIFT-08.

Write \(V_i(t)=r(te_i)\). Its restriction is nonzero in the two cases: \[
\sigma_q(t,0,0)=\frac{t^4}{12}+O(t^{18}),\qquad
\sigma_q(0,t,0)=-t^2+O(t^{16}).
\tag{K21.12}
\] These initial terms are exact coefficients of the actual section in AM-K014. Since \(F(2te_i)\) is a formal unit, (K21.11) implies \(V_i\ne0\). In fact \(V_i(0)=0\) in both cases. No unit quotient is assumed.

The remaining hypothesis of (K21.7) is boundedness of this actual restricted prefactor. The accepted DP45/53 theorem gives, in these globally labelled integral eigendirections, \[
c_i^{\mathrm{bare}}(p)=\frac{p}{(p-1)(p^3+1)}=pc,
\qquad q=p^3.
\tag{K21.13}
\] For every integral formal germ \(r\), its invariant Taylor coefficients therefore satisfy \[
v([t^n]r(te_i))\ge-(s-pc)n.
\tag{K21.14}
\] This is the accepted mixed bound with only direction \(i\) retained, followed by the exact factorial formula; any zero coefficients only improve it. The global invariant basis differs from DP45's primitive display basis by diagonal units, as explicitly checked in DP53.

Every \(w\in I\) satisfies \[
w\ge s-2c>s-pc,
\tag{K21.15}
\] because \(p\ge3\). Thus \(V_i(\gamma T)\) is integral and bounded throughout the whole initial-form interval. Its zero at the origin is harmless for DP41's nonzero bounded-multiplier lemma. Apply (K21.9) to \(V_i(t)\exp(H_i(2t))=\sigma_q(2te_i)\). Finally \(2^n\) is a unit, so its derivatives have exactly the same valuations as those of \(\sigma_q(te_i)\). This proves the matching liminf and the displayed digit-loss lower bound in (K21.1). The extra lower bound zero follows from the actual odd-prime Hurwitz integrality of \(\sigma_q\) in AM-K014.

\subsection{5. Explicit sharp subsequences of the fixed source section}\label{explicit-sharp-subsequences-of-the-fixed-source-section}

\subsubsection{5.1. Exact finite Sato determinant, without partition enumeration}\label{exact-finite-sato-determinant-without-partition-enumeration}

Two additional finite residues are \[
\partial_1^{18}\sigma_q(0)\equiv1\pmod3,\qquad
\partial_2^{30}\sigma_q(0)\equiv4\pmod5.
\tag{K21.17}
\] Here is a direct finite certificate and its truncation justification. Use AM-K014's actual Sato frame \[
v_j(t)=t^{\,2-n_j}A(t^{14})^{-\lfloor n_j/2\rfloor},
\quad A=1-TA^7,\quad(n_j)=(0,2,4,6,7,8,\ldots).
\] Write \(a_{m,j}=[t^m]v_j\). For a single KP time \(t_d=v\), put \(h_{dn}=v^n/n!\) and \(h_m=0\) when \(m<0\) or \(d\nmid m\). Then \[
\tau(v)=
\det\left(\sum_m a_{m,j}h_{m+i}(v)\right)_{1\le i,j\le L}
\pmod{v^{N+1}}
\qquad(L\ge\max\{3,dN\}).
\tag{K21.18}
\] To prove this, apply Cauchy--Binet, ordering the selected exponents \(m_1>\cdots>m_L\ge-L\). Set \(\lambda_i=m_i+i\). These form a partition, and the two determinants in each summand are respectively its stable Plücker coefficient and \(\det(h_{\lambda_j-j+i})=s_\lambda\), by Jacobi--Trudi. Every term of ordinary \(v\)-degree at most \(N\) has \(|\lambda|\le dN\), hence length at most \(L\). All its Plücker coefficients already stabilize at this \(L\): the added rows and columns have upper triangular identity tail, since \(v_j\) has leading exponent \(-j\) for \(j\ge4\) and only higher corrections. Thus no omitted partition can contribute. The sum over \(m\) is finite at each truncation, which also justifies Cauchy--Binet formally.

The exact integer coefficients needed for the matrix are \[
[T^k]A^{-m}
=\frac m k(-1)^{k-1}\binom{7k-m-1}{k-1}\quad(k\ge1),
\qquad [T^0]A^{-m}=1.
\tag{K21.19}
\] The generalized integer binomial is used if the upper entry is negative. Its displayed numerator is divisible by \(k\), as also follows from the integral implicit recursion for \(A\).

The script

\begin{verbatim}
work/agents/jacobian_sigma_hurwitz_determinant.py
\end{verbatim}

evaluates (K21.18) in the truncated Hurwitz ring. Its multiplication is \((ab)_n=\sum_{k=0}^n\binom nk a_kb_{n-k}\); it never inverts \(n!\) modulo \(p\). Complete unit pivoting leaves the actual two-dimensional constant-rank defect and an exact \(2\times2\) residual determinant. Finally the factor \(4\cdot2^{-n}\) converts the \(n\)-th derivative of \(\tau(v)=\sigma(v)\) into that of \(\sigma_q(u)=4\sigma(u/2)\).

For the two computations, \((d,N)=(5,18)\) and \((3,30)\), so \(L=90\) is sufficient. Each was executed again with \(L=104\); the complete retained derivative vectors agreed. The outputs are

\begin{verbatim}
work/agents/jacobian_sigma_hurwitz_p3_order18.json
work/agents/jacobian_sigma_hurwitz_p3_order18_tail14.json
work/agents/jacobian_sigma_hurwitz_p5_order30.json
work/agents/jacobian_sigma_hurwitz_p5_order30_tail14.json
\end{verbatim}

The nonzero residues through those cutoffs are respectively \(S_4=2,S_{18}=1\) modulo 3, and \(S_2=3,S_{16}=4,S_{30}=4\) modulo 5. Separate low-order executions reproduce \(S_4=2\), \(S_2=-2\), and \(S_6=1\) in their corresponding directions. These are exact finite computations, distinct from the all-order argument that follows. The one-command replay is

\begin{verbatim}
python work/agents/jacobian_sigma_sharp_certificate.py
\end{verbatim}

It regenerates both matrix sizes, verifies the entire retained vectors, and writes jacobian\_sigma\_sharp\_certificate.json with script hashes, exact critical weights, finite possible initial degrees, and witness digit sums. Its recorded status is PASS.

\subsubsection{5.2. The critical exponential coefficients}\label{the-critical-exponential-coefficients}

At \(w=w_*\), the scaled logarithm \(H(2\gamma T)\) has exactly one coefficient of valuation \(s\), in degree \(n_0=q+1\); every other coefficient has valuation strictly above \(s\). Write \(E(T)=\exp(H(2\gamma T))\). Its constant term is one, and \[
v([T^n]E)\ge s\quad(n>0),\qquad
v([T^n]E)=s
\ \Longleftrightarrow\ n=n_0p^r\ (r\ge0).
\tag{K21.20}
\] For the equivalence, factor the exponential formally into the exponentials of its separate logarithmic monomials. A nonzero coefficient from a logarithmic coefficient of valuation \(s+\epsilon\) has valuation at least \(s\,s_p(m)+m\epsilon>s\). The unique critical monomial contributes in degree \(mn_0\) with valuation \(ms-v(m!)=s\,s_p(m)\), which equals \(s\) exactly when \(m=p^r\). All other contributions at such a degree have larger valuation, so there is no cancellation. Each coefficient involves only finitely many factors, making this a formal argument rather than an analytic product assertion.

\subsubsection{5.3. The actual algebraic prefactor has a unique critical monomial}\label{the-actual-algebraic-prefactor-has-a-unique-critical-monomial}

Set \(\widetilde V(T)=r(\gamma T e_i)\) and \(G(T)=\sigma_q(2\gamma T e_i)=\widetilde V(T)E(T)\). By (K21.14), \[
v([T^j]\widetilde V)\ge\delta_0j,\qquad
\delta_0=(p-1)c>0.
\tag{K21.21}
\] Thus its coefficient infimum \(\nu\) is attained in finitely many degrees. Since \(E(0)=1\) and every positive coefficient of \(E\) has valuation at least \(s>0\), Gauss multiplicativity and reduction show that \(G\) and \(\widetilde V\) have the same infimum and the same entire initial form. In particular their initial degrees must satisfy the actual CM selection rule of \(\sigma_q\).

For \(p=3\), the exact unit \(S_{18}\) gives \[
\nu\le -v_3(18!)+18w_*=-8+\frac{27\cdot18}{56}
=\frac{19}{28}.
\] Any initial degree of \(\widetilde V\) is at most \((19/28)/\delta_0=19\). The source CM rule leaves only degrees 4 and 18. The value at 4 is \(-1+4w_*=13/14>19/28\), so \[
\nu=\frac{19}{28},\qquad
\operatorname{in}(\widetilde V)=aT^{18}\quad(a\ne0).
\tag{K21.22}
\]

For \(p=5\), the exact unit \(S_{30}\) gives \[
\nu\le -v_5(30!)+30w_*=-7+\frac{125\cdot30}{504}
=\frac{37}{84}.
\] Initial degrees are at most \((37/84)/\delta_0=111/2\), hence the selection rule leaves only 2, 16, 30, and 44. The value at degree 2 is \(2w_*=125/252>37/84\). Hurwitz integrality alone bounds the values at 16 and 44 below by respectively \(s\,s_5(16)-16c=1-16/504\) and \(s\,s_5(44)-44c=2-44/504\), both larger than \(37/84\). Therefore \[
\nu=\frac{37}{84},\qquad
\operatorname{in}(\widetilde V)=aT^{30}\quad(a\ne0).
\tag{K21.23}
\]

\subsubsection{5.4. Unique coefficient contributions give exact valuations}\label{unique-coefficient-contributions-give-exact-valuations}

In either case put \(m=18\) or 30, respectively, and \(n=n_0p^r+m\). The contribution \([T^m]\widetilde V\,[T^{n_0p^r}]E\) has valuation exactly \(\nu+s\). Every other contribution with a positive \(E\)-degree has larger valuation, by the uniqueness in (K21.22)--(K21.23) and (K21.20). The sole constant-\(E\) contribution is \([T^n]\widetilde V\); it also has larger valuation because \[
\delta_0(n_0+m)>\nu+s:
\qquad \frac{46}{28}>\frac{33}{28},\quad
\frac{156}{126}>\frac{29}{42}.
\tag{K21.24}
\] Thus for every \(r\ge0\), \[
v([T^{n_0p^r+m}]G)=\nu+s.
\tag{K21.25}
\] Undo the weight and the unit rescaling by 2, and multiply by \(n!\). Since \(v(S_m)=0\) and \(s_p(m)=2\) in both cases, the result is \[
v(S_n)=cn+\nu+s-s\,s_p(n)
=\frac{p^r+3-s_p(n)}{p-1}.
\tag{K21.26}
\] For \(p=5\), adding \(30=5+25\) to \(126\cdot5^r\) causes no base-five carry, so \(s_5(n)=4\) for every \(r\ge0\). For \(p=3\), adding \(18=2\cdot3^2\) causes precisely one carry when \(r=2\) and none for \(r=0,1\) or \(r\ge3\). This proves every assertion of (K21.1a), including the small-order exception. It also independently yields the two liminf upper bounds in (K21.1).

\subsection{6. Source scope, exact comparison, and the remaining general-prime condition}\label{source-scope-exact-comparison-and-the-remaining-general-prime-condition}

The resulting exact comparison is

\begin{longtable}[]{@{}
  >{\raggedright\arraybackslash}p{(\columnwidth - 6\tabcolsep) * \real{0.2000}}
  >{\raggedleft\arraybackslash}p{(\columnwidth - 6\tabcolsep) * \real{0.2667}}
  >{\raggedleft\arraybackslash}p{(\columnwidth - 6\tabcolsep) * \real{0.2667}}
  >{\raggedleft\arraybackslash}p{(\columnwidth - 6\tabcolsep) * \real{0.2667}}@{}}
\toprule\noalign{}
\begin{minipage}[b]{\linewidth}\raggedright
Prime and fixed global direction
\end{minipage} & \begin{minipage}[b]{\linewidth}\raggedleft
Actual bare DP45/53 slope
\end{minipage} & \begin{minipage}[b]{\linewidth}\raggedleft
Actual \(\sigma_q\) sharp slope
\end{minipage} & \begin{minipage}[b]{\linewidth}\raggedleft
Difference
\end{minipage} \\
\midrule\noalign{}
\endhead
\bottomrule\noalign{}
\endlastfoot
\(p=3\), \(\partial_1\) & \(3/56\) & \(1/56\) & \(1/28\) \\
\(p=5\), \(\partial_2\) & \(5/504\) & \(1/504\) & \(1/126\) \\
\end{longtable}

The operators are already specified global CM directions; no prime-dependent relabelling occurs in either statement. Each pure monomial is killed by conjugate-pair contractions whenever it is placed alone in a positive CM-labelled polynomial block. That elementary harmonic observation does not assert membership in an unspecified full row with its other degree, invariance, and HC conditions. Nor does the construction identify these particular pure derivative sequences with a required later theta-lift image.

The same proof gives a precise generalization criterion. For \(p\equiv3\pmod7\), take \(i=1\); for \(p\equiv5\pmod7\), take \(i=2\). In either case \(q=p^3\), the all-order logarithmic bound and the bare bound (K21.13) already hold for the actual Jacobian. If the single actual constant \[
\partial_i^{p^3+1}\log F(0)
\tag{K21.16}
\] is a \(p\)-adic unit, then the actual fixed source germ has the exact sharp slope \(1/((p-1)(p^3+1))\) in that direction. AM-K017 makes this a global Hessian restricted-power constant; it is independent of the chosen nonzero local theta frame. Its unit status for every prime in these two residue classes is not proved here. A unit claim cannot be inferred solely from the existence of a rank-one cyclic display without identifying the relevant universal-extension or primitive-kernel pairing.

Thus the new result is unconditional for the two explicitly certified primes 3 and 5, and gives a single checkable unit condition at every other prime in the two classes. It supplies actual directional upper sharpness without a conjugated-product gauge gap. It does not settle the full prime budget or the later selected-row cost.

\subsection{7. Exact accepted-result locators and verification boundary}\label{exact-accepted-result-locators-and-verification-boundary}

The accepted source archive is RETURN\_DP\_MULT\_checkpoint\_06\_20261121T0420Z. The used statements are:

\begin{itemize}
\tightlist
\item
  research/DP41\_integrated\_fragment.tex, proof paragraph ``Exponential obstruction, including arbitrary nonzero bounded multipliers,'' lines 131--147. Its lemma allows a complete ramified coefficient DVR, a logarithm with coefficient infimum in \([0,1/(p-1))\), a nonzero-derivative entire initial form, and any nonzero bounded multiplier. Its proof and hypotheses were read to verify applicability; the accepted lemma is not rederived here.
\item
  research/DP45\_integrated\_fragment.tex, the actual integral cyclic-display theorem and its primitive directions; research/DP53\_integrated\_fragment.tex, the full globally labelled spectrum and the diagonal-unit comparison of the fixed differential basis.
\item
  AM-K012--015, the actual Hessian, stable integral ring, {[}2{]}-ratio, and source sigma normalization; AM-K017, the complete finite certificate for (K21.10); AM-K018, the exact full-\(\mathcal O_E\) Chern/source gauge.
\end{itemize}

The new weighted inequalities (K21.5)--(K21.6), the multiplier verification (K21.11)--(K21.15), the exact finite determinant specialization and residues (K21.17)--(K21.19), and the actual explicit witness sequences (K21.20)--(K21.26) are the proof contributions of this note. The underlying genus-three theta identities are cited and source-audited in AM-K012--014; this note does not introduce a new analytic theta identity.

\clearpage

\section{AM-R-009: an actual nonzero first lift on the common unramified mixed-CM host}\label{am-r-009-an-actual-nonzero-first-lift-on-the-common-unramified-mixed-cm-host}

Author: /root, GPT-6 Astra Pro, ALP\_MACHINE, 21 November 2026 UTC. This proof constructs the first image on the same field and polarized abelian variety as AM-DO-022--024. It does not identify the rank-one starting edge with the source's prescribed even-rank arithmetic chain. The finite compact-vacuum assertion in section 3 has been checked against Pan (2001), section 6.2, including the rank-one source case. The source input has norm one with the specified probability Haar measure.

\subsection{1. The precise host, period, and eigen-coordinate}\label{the-precise-host-period-and-eigen-coordinate}

Use the accepted DP57 host, as normalized in AM-DO-022: \[
 E'=\mathbf Q(\zeta_{28})=\mathbf Q(\zeta_7,i),\qquad
 F'=E'^{+},\qquad A=J^2,\qquad \xi'=\xi/2.
 \tag{R9.1}
\] The action of \(i\) on the two ordinary \(J\) factors is \((u,v)\mapsto(-v,u)\). The trace form with coefficient \(\xi'\) is exactly the product principal polarization, with no isogeny or change to the underlying four Jacobian factors used below. Put \[
 w_j^+=u_j+iv_j,\qquad w_j^-=u_j-iv_j,\qquad
 s_A(u,v)=\sigma_q(u)\sigma_q(v),\qquad C=\theta_0c.
 \tag{R9.2}
\] All symbols \(\sigma_q,\theta_0,c\), and \(G_j\) retain their normalizations in AM-R-004--006. In particular \(\sigma_q(u)=u_1u_3-u_2^2+\) terms of degree greater than two, and the actual Bargmann theta realization of the original finite vacuum is \(C\sigma_q\).

Since \(\zeta_{28}=\zeta_7^2i^{-1}\), its eigenvalue on \(w_j^+\) has label \(8j-7\bmod28\), while its label on \(w_j^-\) is \(8j+7\bmod28\). Thus \[
 \Phi'=\{1,3,9,15,17,23\},\qquad w_2^+\text{ has label }9.
 \tag{R9.3}
\] The derivative in that coordinate is \[
 \partial_{w_2^+}=\tfrac12(\partial_{u_2}-i\partial_{v_2}).
 \tag{R9.4}
\] The first nonzero homogeneous term of \(s_A\) has degree four. Consequently, in its fourth derivative only the term with two derivatives on each factor contributes: \[
 \partial_{w_2^+}^{\,4}s_A(0)
   =2^{-4}\binom42(-i)^2(-2)(-2)=-\frac32.
 \tag{R9.5}
\] This is an exact coefficient calculation, independent of analytic tail estimates.

Take the target positive space \(V=E'^2\) and the source skew line \(W\) with coefficient \(-2\xi'\), in the half-trace Hom convention of AM-DO-022. On the underlying \(A^2=J^4\), let \[
 P=\partial_{w_{2,1}^+}^{\,4}\partial_{w_{2,2}^+}^{\,4}.
 \tag{R9.6}
\] The actual finite vacuum is the tensor of the four original finite Schwartz vacua; its norm is one. Its Bargmann realization is \[
 \mathcal R\phi_{00}=C^4s_A(z_1)s_A(z_2).
 \tag{R9.7}
\] Apply the actual creation operators corresponding to \(P\), as in ALP-LIFT-09, and denote the resulting global Schwartz vector by \(\phi_P\). Equations (R9.5)--(R9.7) give \[
 \boxed{\vartheta_{\phi_P}(1,1)=\frac94\,C^4\ne0.}
 \tag{R9.8}
\] Its degree vector is eight at CM label 9 and zero at all other labels. Hence this is a single source-torus isotypic Fock vector, not a sum whose nonzero components remain unidentified.

The period \(C^4\) has not been set equal to one. In the algebraic coordinate normalization the displayed coefficient is \(9/4\); it is an \(11\)-adic unit. The coordinate change (R9.4) has a dyadic denominator. No claim of integrality at 2 or at a changed bad-prime Hodge basis is inferred from the coefficient at 11.

\subsection{2. A globally admissible splitting, with its infinity labels}\label{a-globally-admissible-splitting-with-its-infinity-labels}

Let \(\mu_E\) be the explicit global character constructed in ALP-LIFT-09, so that \(\mu_E|_{\mathbf A_F^\times}=\epsilon_{E/F}\). Define \[
 \mu'=\mu_E\circ N_{E'/E},\qquad \chi_S=\mu',\qquad \chi_T=1.
 \tag{R9.9}
\] The fields satisfy \(E'=EF'\) and \(E\cap F'=F\). Restriction of the norm in (R9.9) to \(\mathbf A_{F'}^\times\) is \(N_{F'/F}\). Compatibility of the quadratic idele-class character with base change therefore gives \[
 \mu'|_{\mathbf A_{F'}^\times}
   =\epsilon_{E/F}\circ N_{F'/F}
   =\epsilon_{E'/F'}.
 \tag{R9.10}
\] This is the required restriction for a source skew line; the even-dimensional target requires the trivial restriction and has the chosen character 1. Use the same rationally compatible standard Weil representation as in the cited constructions.

At the ordered CM labels \(1,3,9,15,17,23\), the exponents \(\mu'_\tau(z)=(z/\bar z)^{m_\tau/2}\) are \[
 \mathbf m=(-1,+1,-1,-1,+1,-1),\qquad
 \boldsymbol\nu=(\mathbf m+\mathbf1)/2=(0,1,0,0,1,0).
 \tag{R9.11}
\] Indeed restriction to \(E\) gives the three labels \(1,3,2,1,3,2\); ALP-LIFT-09 and the corrected AM-DO-023 give exponent \(-1\) on labels 1 and 2 and \(+1\) on label 3. The target compact type of the actual first image is accordingly \[
 \pi_\infty=\mathop{\boxtimes}_{\tau\in\Phi'}
       \det^{\nu_\tau}\operatorname{Sym}^{d_\tau}\mathbf C^2,
 \quad d_9=8,\quad d_\tau=0\ (\tau\ne9).
 \tag{R9.12}
\] For the Hom source action \(X\mapsto Xt^{-1}\), its actual torus eigencharacter is \[
 \beta(t)=\prod_{\tau\in\Phi'}\tau(t)^{-1}\,\tau_9(t)^{-8}.
 \tag{R9.13}
\] Here the exponent \(-1\) at each place is the rank-two target's vacuum shift with \(\chi_T=1\). The character \(\chi_S\) changes the target action, not this source action. On the rational generator \(\zeta_{28}\), (R9.13) has exponent \[
 -\sum_{\tau\in\Phi'}\tau-8\cdot9=-68-72=-140\equiv0\pmod{28}.
 \tag{R9.14}
\] Thus \(\beta\) is trivial on \(\mu_{28}\).

\subsection{3. The torus quotient and the actual input}\label{the-torus-quotient-and-the-actual-input}

Put \(T'=U(1)(E'/F')\) and \(K_0=(\widehat{\mathcal O}_{E'})^1\). Masley (1976), Main Theorem, printed pp.~179 and 184, implies \(h(E')=1\): \(\varphi(28)=12\le24\), and 28 does not occur in the displayed list with class numbers 2 through 10. This published class-number result is an input, not a new calculation or a computational class-group certificate.

It implies the exact torus decomposition \[
 T'(\mathbf A_{F',f})=T'(F')K_0,\qquad
 T'(F')\cap K_0=\mu_{28}.
 \tag{R9.15}
\] Here is a short proof of the first assertion in the required norm-one form. For a norm-one finite idele \(t\), its fractional ideal \(I(t)\) satisfies \(I(t)\overline{I(t)}=\mathcal O_{E'}\). Every nonsplit prime therefore has exponent zero, and the exponents at a split conjugate pair are opposite. Choose a fractional ideal \(B\), supported on one prime in each such pair, so that \(I(t)=B/\bar B\). Class number one gives \(B=(b)\). The rational element \(b/\bar b\) has norm one and exactly the same ideal as \(t\). Their quotient belongs to \(K_0\). This proves the decomposition without assuming a statement about signatures of real units. For the intersection, a norm-one integral unit has every complex absolute value one; Kronecker's theorem makes it a root of unity, and the roots of unity of \(E'\) are precisely \(\mu_{28}\).

The source action preserves the actual one-dimensional finite vacuum eigenspace, because norm-one integral units preserve its trace lattice and semicharacter. Therefore it has an actual smooth unitary character \(\lambda_f:K_0\to\mathbf C^\times\). This is not an unspecified phase used to identify a moved kernel: it is the character of the fixed finite representation on its named one-dimensional line. Put \(K=\ker\lambda_f\) and \(\Gamma=\mu_{28}\cap K\). Smoothness and compactness imply that \(K\) has finite index. Equation (R9.15) gives \[
 h_K=\#T'(F')\backslash T'(\mathbf A_{F',f})/K
     =\frac{[K_0:K]}{[\mu_{28}:\Gamma]}.
 \tag{R9.16}
\] Rational invariance of the nonzero theta value (R9.8) gives \(\beta(\zeta)\lambda_f(\zeta)=1\) for every \(\zeta\in\mu_{28}\). Thus \(\beta|_\Gamma=1\); in fact (R9.14) makes \(\lambda_f|_{\mu_{28}}=1\).

Normalize Haar measure on the compact torus quotient \(T'(F')\backslash T'(\mathbf A_{F'})\) to have volume one. Its \(h_K\) finite components have equal volume \(1/h_K\). Define the actual input \(f_K\) to be \(h_K\beta(t_\infty)\) on the identity finite component and zero on the other components, using rational triviality to transport that description. The condition \(\beta|_\Gamma=1\) makes this definition unambiguous. Its exact norms are \[
 \|f_K\|_2^2=h_K,\qquad \|f_K\|_\infty=h_K.
 \tag{R9.17}
\] The source action on \(\phi_P\) is \(\beta\) at infinity and is trivial on \(K\) at finite places. The target action commutes with the source action. Integrating over the one component supporting \(f_K\), with the convention \[
 \Theta_P(g)=\int_{[T']}\vartheta_{\phi_P}(t,g)
                                      \overline{f_K(t)}\,dt,
 \tag{R9.18}
\] therefore gives the equality of actual automorphic functions \[
 \boxed{\Theta_P(g)=\vartheta_{\phi_P}(1,g),\qquad
        \Theta_P(1)=\frac94 C^4.}
 \tag{R9.19}
\] For norm-one source input \(f_K/\sqrt{h_K}\), the output is \(h_K^{-1/2}\Theta_P\). This proves a nonzero first lift on the common host without any unproved claim that \(\lambda_f\) is trivial.

\textbf{The actual compact character is trivial.} At every odd unramified quadratic place the target identity plane is integrally hyperbolic, as in AM-DO-022. Pan (2001), section 6.2, Lemma and equation (6.2.a), printed p.~197, treats exactly a one-dimensional moving unitary space and a two-dimensional isotropic opposite space. Its comparison has \(\alpha=1\) and \(\beta=\chi((t-1)\delta)\). The opposite character is \(\chi_T^{-1}=1\), so both factors are one. The lattice-model action fixes the actual selfdual-lattice vacuum, as specified in Pan section 3.3 and the normalization of AM-DO-022. There is no assumption that the moving rank is even; this primary lemma is itself the rank-one case. The other splitting character \(\chi_S=\mu'\) affects the target action only. At relative split places, including all dyadic places, AM-DO-020 gives the same compact invariance with the retained trace-one convention, since its source determinant-root factor is \(\chi_T(x)^{-1}=1\). The identity element acts as one as well. These cases cover every finite place of \(E'/F'\). Consequently \[
 \boxed{\lambda_f=1,\quad K=K_0,\quad h_K=1,\quad
        \|f_K\|_2=\|f_K\|_\infty=1.}
 \tag{R9.27}
\] Thus (R9.19) is already the first lift with norm-one source input. The general-component calculation above records exactly how a nontrivial finite character would have entered, without discarding such a cost by convention.

\subsection{4. Exact Fock normalization}\label{exact-fock-normalization}

In the \(w^\pm\) coordinates the positive Hermitian form on \(A\) is \[
 H_A=\sum_{j=1}^3
       \frac{|w_j^+|^2+|w_j^-|^2}{2G_j}.
 \tag{R9.20}
\] The creation operator for \(w_2^+\) is one half of the creation operator for \(u_2\) minus \(i\) times that for \(v_2\). Orthogonality of the two modes and the normalization in ALP-LIFT-09 give its one-creation squared norm \(\pi/(2G_2)\). The actual Gaussian used in (R9.7), before applying creation operators, has squared norm \(\det(2Y)^{-2}\), because there are four ordinary Jacobian factors. Hence the exact squared norm of the actual archimedean vector in (R9.6) is \[
 \|\phi_{P,\infty}\|^2=\det(2Y)^{-2}(4!)^2
                                  \left(\frac{\pi}{2G_2}\right)^8.
 \tag{R9.21}
\] The unit archimedean vector therefore has identity value \[
 C^4\det(2Y)\frac{3}{32}\left(\frac{2G_2}{\pi}\right)^4.
 \tag{R9.22}
\] Combining unit archimedean vector and norm-one source input adds the factor \(h_K^{-1/2}\). The determinant factor is necessary: an earlier draft used a norm relative to unit Gaussian and omitted it when normalizing the actual theta Gaussian; the present formulas correct that error. These are analytic normalizations; they do not replace the algebraic coefficient \(9/4\) or the common period \(C^4\).

\subsection{5. A fixed character repair for the literal bilinear second edge}\label{a-fixed-character-repair-for-the-literal-bilinear-second-edge}

This paragraph uses the corrected source convention of AM-DO-023 and the literal bilinear theta integral in ALP-LIFT-11's amended version. Let \[
 c_\tau=2-\nu_\tau,\qquad
 \mathbf c=(2,1,2,2,1,2).
 \tag{R9.23}
\] A global norm-one character with this infinity type can be constructed explicitly from (R9.15). Choose a split place over 29 with \(\zeta_{28}\mapsto2\pmod{29}\). The element 2 generates \(\mathbf F_{29}^\times\). The infinity character \(\beta_{\mathbf c}\) has exponent \[
 \sum_{\tau\in\Phi'}c_\tau\tau=116\equiv4\pmod{28}
 \tag{R9.24}
\] on \(\mu_{28}\). Let its finite compact character be trivial elsewhere and take the value \(\exp(2\pi i\,24/28)\) on 2 at this place. Its restriction to \(\mu_{28}\) is exactly \(\beta_{\mathbf c}^{-1}\), so their product descends through (R9.15) to a unitary automorphic character \(\gamma\). Equivalently, for \(t_f=\delta k\), define \(\gamma(t_\infty,t_f)
=\beta_{\mathbf c}(\delta_\infty^{-1}t_\infty)\rho(k)\). Independence of the choice of \(\delta\) follows from the displayed equality on roots of unity. Its finite conductor divides this fixed place over 29; its compact values have order dividing 7.

Replace \(\chi_S\) in (R9.9) by \(\chi_S'=\mu'(\gamma\circ j)\), where \(j(x)=x/\bar x\); leave \(\chi_T=1\). The restriction (R9.10) is unchanged. The actual splitting ratio is \(\gamma(\det g)\) on the target. To check every place, at nonsplit finite places the local norm-one torus is its maximal compact, and \(\gamma\) is trivial there; at split finite places use the determinant-root formula of AM-DO-020; at real places the corrected Fock formula changes the target power by \(c_\tau\). Thus \[
 \Theta'_P(g)=\gamma(\det g)\Theta_P(g)
 \tag{R9.25}
\] as actual automorphic functions. All identity values and scalar automorphic norms are preserved exactly. The first target types become \[
 \pi'_\tau=\det^2\operatorname{Sym}^{d_\tau}\mathbf C^2,
 \tag{R9.26}
\] and every infinity exponent of \(\chi_S'\) is 3.

For the next positive Hom edge \(U(2)\to U(4)\) with both splitting characters 1, the source Fock type of partition \((d_\tau,0)\) is \(\det^{-2}\operatorname{Sym}^{d_\tau}(\mathbf C^2)^*\). It is exactly the contragredient of (R9.26), as required by the literal bilinear integral \(\int\vartheta\,f\). The degrees have not increased. This is local occurrence and exact first-image preservation, not a proof of the global second integral's nonvanishing.

The second edge now lies on precisely the E-prime host of AM-DO-022--024, so its finite moved-vacuum formula, positive direct-average sign, common period, and probability-Haar theta-line pairing are available without an extension-of-field comparison. A prescribed global second image, its finite input, its automorphic mass and its selected coefficient-lattice costs remain to be proved.

\subsection{6. Geometric location and limits}\label{geometric-location-and-limits}

The abelian variety is the ordinary product \(J^4\), with the extra \(E'\)-action specified in (R9.1)--(R9.3); there is no substitution of the conjugate E-action product excluded from AM-K018. The value is an actual jet of the trace-one Picard line, transferred through its reduced Poincare trivialization to a global theta integral. The integral source input is an automorphic character of the norm-one torus; the repair is tensoring an actual target representation by a rank-one character. The exact FGA/SGA locators for these constructions are those in GROTHENDIECK\_DICTIONARY.md sections 2--5 (FGA 232 Theorem 3.1; FGA 236 Corollaries 1.7 and 2.3; SGA 7 I VII 2.7 and 2.9; SGA 3 I Proposition 4.7.3).

This closes the previously separate-host obstruction for these particular first and candidate second edges. It does not instantiate the old elliptic VacIn, the even-rank starting coefficient module, the full six arithmetic identities, or the DP44 gain/cost data.

\subsection{References and exact inputs}\label{references-and-exact-inputs}

Masley, J. M. (1976). Solution of small class number problems for cyclotomic fields. \emph{Compositio Mathematica, 33}(2), 179--186. Main Theorem, pp.~179 and 184; the table was visually checked on PDF pages 2 and 7 (one-based). \url{https://numdam.org/item/CM_1976__33_2_179_0.pdf}

AM-R-004--006 (actual periods, metric, trace-one scalar); AM-K018 (ordinary-product source gauge); AM-DO-020, AM-DO-022--024 (fixed finite conventions and E-prime host); corrected AM-DO-023 (DO23.3--DO23.9); ALP-LIFT-09 (V7--V11 and V18--V38); ALP-LIFT-11 (explicit torus-character extension and bilinear repair, used with the newly proved E-prime torus decomposition).

Pan, S.-Y. (2001). Splittings of the metaplectic covers of some reductive dual pairs. Pacific Journal of Mathematics, 199(1), 163--226. Section6.2, Lemma and equation(6.2.a), p197; section3.3 for the lattice normalization. \url{https://msp.org/pjm/2001/199-1/pjm-v199-n1-p10-p.pdf}

\clearpage

\section{AM-R-010: the permitted first-image family on the common E-prime host}\label{am-r-010-the-permitted-first-image-family-on-the-common-e-prime-host}

Author: /root, GPT-6 Astra Pro, ALP\_MACHINE, 21 November 2026 UTC. This extends the actual first-image construction, not an iterated nonvanishing or a cost computation for the full source chain.

\subsection{1. Statement and exact finite index}\label{statement-and-exact-finite-index}

Use \(E'=\mathbf Q(\zeta_{28})\), \(F'=E'^+\), \(A=J^2\), and the exact polarization and standard splitting of AM-R-009. Let the positive target have an even rank \(n\), and the source be \(T'=U(1)\). Thus the underlying theta variety is \(A^n=J^{2n}\), and its fixed vacuum period is \(C^{2n}\). For the level parameter \(M=29^e\), \(e\ge1\), use the same odd norm element \[
 \alpha=(1+2\sqrt{-7})^e,\qquad \alpha\bar\alpha=M
 \tag{R10.1}
\] as ALP-LIFT-10 and DP59, now acting on \(A\) through \(\mathcal O_E\subset\mathcal O_{E'}\). This \(M\) is an arithmetic level, not the number of theta lifts in DP44.

Let \(K_M'\subset K_0=(\widehat{\mathcal O}_{E'})^1\) be principal congruence modulo \(29^e\) at all six split places of \(F'\) over 29, and full maximal compact elsewhere. The common finite stabilizer, finite class count, and its prime contents are \[
 \boxed{\quad
 h_M'=\#T'(F')\backslash T'(\mathbf A_{F',f})/K_M'
       =28^5\,29^{6(e-1)}.
 \quad}
 \tag{R10.2}
\] In particular \[
 v_2(h_M')=10,\quad v_7(h_M')=5,\quad
 v_{29}(h_M')=6(e-1),\quad v_\ell(h_M')=0
 \quad(\ell\notin\{2,7,29\}).
 \tag{R10.3}
\] There is no ramified-7 index-two correction on this host. Every permitted translated theta germ of DP59 on these ordinary Jacobian factors, and every polynomial jet of total degree at most \(D\), is an actual first theta image. For a single fixed CM weight the common input has squared source norm \(h_M'\). For a finite set \(\mathcal B\) of weights, the single input realizing their entire direct sum has squared norm \[
                         |\mathcal B|\,h_M'.
 \tag{R10.4}
\] All polynomials of degree at most \(D\) require at most \[
                         |\mathcal B|\le\binom{D+6}{6}.
 \tag{R10.5}
\] This count concerns six total CM degrees, not the number of matrix entries in the target.

The corresponding unit-input logarithmic norm cost is exactly \[
 \frac12\log(h_M'|\mathcal B|)
 =3\log M+\frac52\log28-3\log29
                        +\frac12\log|\mathcal B|.
 \tag{R10.6}
\] For \(M\) polynomially bounded in \(D\), it is \(O(\log(D+2))\). This assertion concerns the torus input and its normalization. The actual archimedean Fock norm, period, target mass, and later lift costs are separate and retained.

\subsection{2. Why the indicated subgroup fixes the actual finite family}\label{why-the-indicated-subgroup-fixes-the-actual-finite-family}

The underlying abelian variety and trace form are the ordinary product of the ones in ALP-LIFT-10. Its identity \[
 U_r^{M^2H}\bigl(s(Mu)\bigr)=(U_x^H s)(Mu),
 \qquad x=Mr,
 \tag{R10.7}
\] including the canonical finite-order lift sign, applies to each factor and hence to their tensor product. This is the actual equality of theta-group translations used in that proof, not a new formal substitution for a selected theta germ. The normalized derivative \(M^{-d}\) cancels the factor \(M^d\) from a homogeneous derivative of degree \(d\).

The allowed half-orbit has local representatives supported at 29 with denominators dividing \(29^e\) in one of the two isotropic factors at each relative split place. The same property remains true after the extension \(E'/E\): 29 splits completely in \(E'\), and the valuation of \(\alpha\) at a selected prime is repeated at the two primes above it. Conjugation exchanges the selected and unselected factors. Thus the local allowed representatives are still in a single isotropic position lattice with denominator \(29^e\).

For \(t\equiv1\pmod{29^e}\), the displacement \((t-1)x\) is integral in that same isotropic factor. Its Heisenberg pairing with \(x\) is zero. At the odd prime 29 the trace-one character on that integral factor is one. Heisenberg covariance therefore gives \[
 \omega_f(t)W_x\phi_{00}
       =W_{tx}\omega_f(t)\phi_{00}=W_x\phi_{00}.
 \tag{R10.8}
\] The initial finite vacuum is fixed by the entire source maximal compact. For target rank two this is AM-R-009, equation (R9.27). For any even target rank, decompose its integrally hyperbolic lattice at an odd unramified place into orthogonal hyperbolic planes and apply the same Pan section6.2 comparison to each plane. Every source factor uses the opposite character 1, so their product is one. At relative split places, including dyadic places, the same assertion follows from AM-DO-020. The character of the odd source affects the target action and does not enter this calculation.

Away from 29 the permitted translation is integral, and its Heisenberg line is the same original line with its specified scalar. The compact source therefore fixes it. Tensoring over entries and places proves that \(K_M'\) fixes the entire allowed finite Schwartz span, not only an individual chosen vector.

At each of the six places over 29 the unit quotient has order \(28\,29^{e-1}\), so \[
 [K_0:K_M']=(28\,29^{e-1})^6.
 \tag{R10.9}
\] The reduction of \(\mu_{28}\) at any one such place is injective, because 2 has order 28 modulo 29. Hence \(K_M'\cap\mu_{28}=1\). AM-R-009 proves \(T'(\mathbf A_f)=T'(F')K_0\). Dividing (R10.9) by 28 gives (R10.2), and the prime valuations give (R10.3).

\subsection{3. One actual source input for every retained weight}\label{one-actual-source-input-for-every-retained-weight}

Let \(\beta\) be an actual archimedean source-torus weight occurring among the chosen creation vectors. Define \(f_\beta\) to equal \(h_M'\beta\) on the identity finite component of \([T']/K_M'\), and zero on the other components. Here the component is a copy of the full compact real torus, since its rational stabilizer is \(K_M'\cap\mu_{28}=1\). There is consequently no extra rational-root restriction on \(\beta\).

The components have probability-Haar mass \(1/h_M'\). If \(\phi_\beta\) is any allowed finite Schwartz combination with this Fock weight, compact invariance, commutation of the two unitary actions, and character orthogonality give \[
 \int_{[T']}\vartheta_{\phi_\beta}(t,g)
                    \overline{f_\beta(t)}\,dt
       =\vartheta_{\phi_\beta}(1,g),\qquad
 \|f_\beta\|_2^2=h_M'.
 \tag{R10.10}
\] This is an equality of actual automorphic functions for every target \(g\), with the fixed splitting and source measure of AM-R-009.

For distinct weights \(\beta,\beta'\), the analogous cross integral is zero by compact real character orthogonality, and \(f_\beta,f_{\beta'}\) are orthogonal. Set \(f_{\mathcal B}=\sum_{\beta\in\mathcal B}f_\beta\). For \(\phi=\sum_{\beta\in\mathcal B}\phi_\beta\), \[
 \Theta(\phi,f_{\mathcal B})(g)
      =\sum_{\beta\in\mathcal B}\vartheta_{\phi_\beta}(1,g),
 \quad
 \|f_{\mathcal B}\|_2^2=|\mathcal B|h_M',
 \quad
 \|f_{\mathcal B}\|_\infty=|\mathcal B|h_M'.
 \tag{R10.11}
\] The supremum equality follows at the identity real torus point, where all characters equal one. The input depends only on the finite weight set, not on the coefficients or chosen translates of \(\phi\).

In the six eigen-coordinates of \(A\), a degree-\(d\) monomial has a total-degree vector in \(\mathbf N^6\) of sum \(d\). Its source character is this vector, with sign reversed and with the same vacuum shift. Different vectors give distinct characters of the full real torus. Counting all vectors of total degree at most \(D\) proves (R10.5). A polynomial in the original ordinary-Jacobian coordinates has such a finite decomposition because \(u_j=(w_j^++w_j^-)/2\) and \(v_j=(w_j^+-w_j^-)/(2i)\). Those transformations have only dyadic denominators; they are integral and invertible at 11. The decomposition is orthogonal in the actual Fock Hermitian metric. It introduces no additional archimedean norm beyond the norm of the original vector, although the source input cost (R10.4) must be retained.

\subsection{4. What the construction transfers, and what it does not}\label{what-the-construction-transfers-and-what-it-does-not}

Equations (R10.7) and (R10.10) transfer every specified permitted jet value on the ordinary product \(J^{2n}\) to an actual first image on the E-prime host, with the exact period \(C^{2n}\). AM-K018 identifies its full-\(\mathcal O_E\) source gauge with the actual Chern gauge on those ordinary powers. The extra \(i\)-action mixes two ordinary copies; it does not replace a factor by the conjugate E-action. Thus the comparison excluded for the latter product is not used here.

For example, an already established DP59 harmonic witness on two ordinary Jacobian factors can be tensored with two permitted spectator sections having unit zero-jet, as in ALP-LIFT-10 section7. Its value on \(A^2=J^4\) and its normalized 11-adic valuation are then unchanged. Decomposing its polynomial into the six CM weights and using (R10.11) realizes that exact value in the actual first image. For a single chosen weight, the further finite-character averaging has denominator \(h_M'\), which is an 11-adic unit. One must not assume that an individual weight summand has the same valuation unless its lower bound is also proved; the exact combined value already follows from (R10.11).

The fixed determinant repair of AM-R-009 depends only on the vacuum shifts, not on these six degrees. It therefore gives local bilinear compatibility with the next positive \(U(2)\to U(4)\) edge for every such first image when \(n=2\). That next edge is over the same E-prime host as AM-DO-022--024. No scalar, finite index, or coordinate denominator is discarded in making this common-host statement.

The present result does not choose the next finite Schwartz vector, compute its integral content or norm, identify the first images with the source's prescribed starting module, or verify the six arithmetic row identities. In particular (R10.6) is not the full \(e_j\) of DP44, and the copied derivative valuation is not automatically the determinant gain \(g_j\). These are precise remaining comparisons, not consequences of first-image existence.

\subsection{5. Framework and exact inputs}\label{framework-and-exact-inputs}

The finite part is a theta-group representation on the section line of a Picard scheme, transported along the actual odd isogeny and its isotropic kernel. Its six weights are characters of the diagonalizable CM torus, and its norm formula uses a finite compact quotient with probability measure. The corresponding exact framework locators are FGA232 Theorem3.1; FGA236 Corollaries1.7 and2.3; SGA7 I VII2.7/2.9; SGA3 I Proposition4.7.3, with the qualifications in GROTHENDIECK\_DICTIONARY.md.

The exact carried inputs are DP59 (permitted half-orbit and zero-jet spectators), AM-K018 (ordinary-product gauge), ALP-LIFT-10 (O1--O28 and section7), AM-R-009 (R9.1--R9.27), and AM-DO-020/022 (finite split and unramified normalizations). Pan (2001), section6.2, Lemma and equation(6.2.a), p197, supplies the rank-one compact comparison used in each hyperbolic plane. No Lean execution is claimed.

\clearpage

\section{ALP-LIFT-12 --- an actual second definite lift of every repaired first-image vector}\label{alp-lift-12-an-actual-second-definite-lift-of-every-repaired-first-image-vector}

Author: /root/lift\_accounting, GPT-6 Astra Pro, ALP\_MACHINE, 21 November 2026 UTC.

This result proves nonvanishing for an actual second classical theta correspondence. Its second Schwartz function is allowed to vary through the full oscillator representation. It does not identify that function with a prescribed source moved kernel, and it supplies no second-output integral content or DP44 height bound.

\subsection{1. Statement and integral conventions}\label{statement-and-integral-conventions}

Set \[
 E=\mathbf Q(\zeta_7),\qquad F=E^+,\qquad
 T=U(1),\quad G=U(I_2),\quad H=U(I_4).
 \tag{S1}
\] The positive Hermitian forms and the corresponding skew forms use the Hom convention already fixed in ALP-LIFT-09 and AM-DO-023. In particular both groups in (S1) are compact at every real place. Use the actual rationally compatible first oscillator and the globally unitary splitting characters of ALP-LIFT-09.

Let \(F_1\ne0\) be any finite sum of the actual first theta vectors constructed there, allowing finite Schwartz translations and finitely many holomorphic Fock polynomials as in ALP-LIFT-10. Equivalently, one may start with a finite sum of first theta lifts of unitary automorphic characters of \(T\) using these polynomial Fock vectors. The source torus input may be the finite projector of ALP-LIFT-10: its finite character decomposition is included in this definition.

Use the one fixed global character of ALP-LIFT-11, \[
 \gamma=\gamma_{(2,2,1)},\qquad
 F_1^+(g)=\gamma(\det g)F_1(g).
 \tag{S2}
\] This is an actual first theta vector for the repaired first splitting. It has the same identity value and the same scalar automorphic norm. At a real place with original polynomial degree \(d_i\), its compact type has highest weight \[
 (d_i+2,2),\qquad
 (\pi_i^+)^\vee=\det^{-2}\operatorname{Sym}^{d_i}(\mathbf C^2)^*.
 \tag{S3}
\] These assertions are ALP-LIFT-11, (R23)--(R26), not new character or real-weight claims in this note.

For the second edge take the actual standard oscillator for \(G\times H\), with both splitting characters equal to \(1\). The assertion is:

\textbf{Theorem.} There is a genuine global Schwartz function \(\psi\) in this second oscillator representation for which \[
 \Theta^{\mathrm{bil}}_{2,\psi}(F_1^+)(h)
   :=\int_{G(F)\backslash G(\mathbf A_F)}
                \vartheta_{2,\psi}(g,h)F_1^+(g)\,dg
 \quad\text{is not identically zero.}
 \tag{S4}
\] This holds for each specific nonzero vector \(F_1^+\), including a nonzero sum across automorphic multiplicity spaces. A target translate of \(\psi\) can make the resulting value at \(h=1\) nonzero.

Formula (S4) is bilinear, as is the literal source liftT\_apply. Gan--Qiu--Takeda use the conjugated-input convention \(\int\vartheta(g,h)\overline{f(g)}\,dg\). We therefore apply their formula to \(f=\overline{F_1^+}\), in the contragredient representation. This change is necessary; applying their compact-type criterion directly to \(F_1^+\) would reverse the relevant source type.

\subsection{2. Temperedness of the actual first constituents}\label{temperedness-of-the-actual-first-constituents}

First fix one unitary torus character contributing to \(F_1\) and one irreducible automorphic constituent of its first image. Projection to this constituent is an equivariant map from the actual first oscillator. Localizing that map shows that each local constituent occurs as a quotient in the corresponding local theta correspondence. This is the usual local equivariance of the theta integral: right translation in one local group can be moved onto the same local Schwartz factor, while integration in the source torus supplies its local character. No assertion about a formal free thetaG parameter is involved.

At a nonsplit finite place, the source is a unitary character of the compact group \(U(1)\). Its one-dimensional Langlands parameter has bounded Weil image. Gan and Ichino, Theorem 4.4(i)(a), (ii)(b), printed p.~14, give the parameter of any nonzero almost equal rank lift: \[
 \phi_{\mathrm{first}}
       =(\phi_T\otimes\chi_V^{-1}\chi_W)\oplus\chi_W .
 \tag{S5}
\] Every character in (S5) is unitary. Hence this parameter is tempered, and so is its local representation. The two cases of that theorem cover both the presence and the absence of the distinguished character in \(\phi_T\); only an already nonzero first image is being considered. Theorem 4.4(iii) also says that the nonzero big theta lift of this tempered source is irreducible. The bounded-parameter/tempered-representation equivalence is the same paper's section 2.

At a split finite place the pair is of type \(\mathrm{GL}_1\times\mathrm{GL}_2\). Mínguez (2008), Theorem 1(2), printed p.~718, says that the small lift has parameter \[
                  \phi_T^\vee\oplus 1
 \tag{S6}
\] in the standard type-II normalization. Changing the honest splitting or the source convention only introduces unitary characters and, where appropriate, a contragredient. Thus (S6) is again bounded. The uniqueness assertion in Theorem 1(1) ensures that an irreducible quotient arising from the actual global first theta integral has this small-theta parameter. This argument includes split dyadic places.

The twist in (S2) is unitary at every place. Taking a contragredient also preserves temperedness. Consequently every finite local constituent of the conjugated input for (S4) is tempered. At real places the group is compact, so its irreducible unitary representations are tempered. We have proved: \[
 \text{Every irreducible constituent of }\overline{F_1^+}
       \text{ is tempered at every place.}
 \tag{S7}
\] For a sum of torus characters, apply the same argument to each summand before decomposing into irreducible target constituents. It does not require that the original projector be one character.

\subsection{3. Every finite local constituent occurs on the second edge}\label{every-finite-local-constituent-occurs-on-the-second-edge}

\subsubsection{3.1. The four-dimensional target form is hyperbolic locally}\label{the-four-dimensional-target-form-is-hyperbolic-locally}

At a nonsplit unramified finite place, the norm map on units of the unramified quadratic extension is surjective. In particular \(-1=N(z)\) for a unit \(z\), and \((1,z)\) is isotropic in \(I_2\). Thus \(I_2\) is hyperbolic and \(I_4\) is the sum of two hyperbolic planes.

For \(E/F\) in (S1), the only relative ramified finite place is above \(7\). At that place choose \(b\in\mathbf Q_7\) with \(b^2=-5\) and \(b\equiv3\pmod7\); Hensel's lemma applies because \(3^2+5\equiv0\pmod7\) and the derivative \(6\) is a unit. In \(I_4\) the two vectors \[
                 u=(1,2,b,0),\qquad v=(0,-b,2,1)
 \tag{S8}
\] are linearly independent and have zero self-pairing and zero mutual pairing. They span a two-dimensional isotropic subspace. Since \(I_4\) is nondegenerate, it is hyperbolic. This proves the required target assertion at every nonsplit finite place without an integral classification assertion. At a split place the second pair is simply \(\mathrm{GL}_2\times\mathrm{GL}_4\).

\subsubsection{3.2. A direct quotient at nonsplit places}\label{a-direct-quotient-at-nonsplit-places}

Let \(k=F_v\), \(L=E_v\) be a nonsplit finite extension, and write the hyperbolic target as \(X\oplus X^*\), with \(\dim_LX=2\). The second Schrödinger model is \[
                    \mathcal S(X\otimes_L W)
                       \simeq\mathcal S(M_2(L)),
 \tag{S9}
\] where \(W\) is the two-dimensional source space. The source group preserves this position space. Its action is right or inverse-right change of variables, depending on the matrix identification, times a smooth unitary character. The absolute determinant multiplier is one on \(U(W)\). Either right-action convention gives the following same quotient conclusion after inversion.

The orbit of the identity matrix is \(U(W)(k)\), and it is closed in \(M_2(L)\). Indeed it is cut out by the equation \(g^*Jg=J\) for a nonsingular matrix \(J\); that equation already forces \(g\) to be invertible. Restriction therefore gives an equivariant surjection \[
 \mathcal S(M_2(L))\longrightarrow C_c^\infty(U(W)(k)),
 \tag{S10}
\] with the same character twist.

Here is the surjectivity detail. If \(C\) is a compact open subset of the closed orbit, its complement in the orbit is closed in the ambient space. Each point of \(C\) has a compact open ambient neighborhood avoiding that complement. Finitely many such neighborhoods cover \(C\); their union intersects the orbit in exactly \(C\). A compactly supported locally constant function on the orbit has a finite partition into such sets. Extending their characteristic functions, and disjointifying the finitely many ambient neighborhoods, extends that function to an element of the ambient Schwartz space.

Every irreducible smooth representation \(\pi_v\) is a quotient of the right regular representation \(C_c^\infty(U(W)(k))\). For a nonzero smooth vector \(a\in\pi_v\), the map \[
 f\longmapsto\int_{U(W)(k)} f(g)\pi_v(g^{-1})a\,dg
 \tag{S11}
\] is equivariant for right translation. Taking the normalized indicator of a compact open subgroup fixing \(a\) shows that the map is nonzero. Irreducibility makes it surjective. The group is unimodular, so the substitution in this equivariance calculation introduces no modular character. A character twist of the regular representation has the same quotient property.

Equations (S10) and (S11) prove \[
                  \operatorname{Hom}_{G(k)}
                          (\omega_{2,v},\pi_v)\ne0
 \tag{S12}
\] for every irreducible smooth \(\pi_v\), in particular every constituent in (S7). This is a quotient argument, not an inference from a subquotient of an open-orbit submodule.

\subsubsection{3.3. Split places}\label{split-places}

At a split place use Mínguez, Theorem 1(2), with \(n=2,m=4\). It states directly that \(\operatorname{Hom}_{\mathrm{GL}_2}(\omega_{2,4},\pi_v)\ne0\) for every irreducible smooth \(\pi_v\), over a nonarchimedean field of any residue characteristic. Unitary splitting characters or a group automorphism do not change this all-representations assertion.

The nonsplit closed-orbit proof must not be copied verbatim here: the full-rank \(\mathrm{GL}_2\) orbit is not closed in the whole matrix space. Its regular quotient from an open submodule alone would not establish a quotient of the full oscillator. The cited type-II theorem closes precisely that point.

\subsection{4. The real local doubling factor is explicitly nonzero}\label{the-real-local-doubling-factor-is-explicitly-nonzero}

By (S3), each real component of the conjugated input has type \[
 \pi_i=\det^{-2}\operatorname{Sym}^{d_i}(\mathbf C^2)^*,
                     \qquad \dim\pi_i=d_i+1.
 \tag{S13}
\] AM-DO-023 and ALP-LIFT-11 show that this exact type occurs in the second positive Hom oscillator with trivial splitting characters. Choose an isometric \(U(2)\)-equivariant embedding \(\iota:\pi_i\hookrightarrow\omega_{2,i}\), a unit vector \(a\), and \(\phi=\iota(a)\). Such a vector is a polynomial times a Gaussian and is an actual Schwartz vector.

For the doubled oscillator, Gan--Qiu--Takeda, section 11.3, printed p.~52, give the equivariant map \(\delta\) normalized by evaluation at zero: \[
              \delta(\phi_1\otimes\overline{\phi_2})(0)
                         =\langle\phi_1,\phi_2\rangle .
 \tag{S14}
\] Their possible determinant character is trivial here because the second target splitting character is \(1\). Consequently its Siegel--Weil section at \(s_0=1\) satisfies \[
 \Phi_\delta(g,1;1)
       =\langle\omega_{2,i}(g)\phi,\phi\rangle
       =\langle\pi_i(g)a,a\rangle .
 \tag{S15}
\] Their local zeta integral in section 11.6, printed p.~54, uses the complex conjugate of the input matrix coefficient. The bars in these two formulas have been checked on rendered PDF pages; the plain-text extraction drops them. Thus Schur orthogonality gives the exact value \[
 \begin{split}
 Z_i(1,\Phi_\delta,a,a)
    &=\int_{U(2)}
          |\langle\pi_i(g)a,a\rangle|^2\,dg\\
    &=\frac{\operatorname{vol}(U(2))}{d_i+1}>0.
 \end{split}
 \tag{S16}
\] In the GQT normalization this yields \[
 Z_i^*(1,\Phi_\delta,a,a)
   =\frac{\operatorname{vol}(U(2))}
      {(d_i+1)L_i(3/2,\pi_i)}\ne0 .
 \tag{S17}
\] The standard local factor of a tempered representation is holomorphic and nonzero in \(\operatorname{Re}s>0\); the precise doubling-factor version is Lapid--Rallis, Proposition 5, also invoked in Gan--Ichino, section 6.7, printed p.~27. For these compact real groups the factors are the usual nonvanishing gamma factors.

This direct compact calculation is needed. GQT Theorem 11.7's general converse does not cover arbitrary definite real unitary groups at \(m=2n\). We are not applying that converse without its extra hypotheses; (S16) proves its required local nonvanishing in the present actual compact types.

\subsection{\texorpdfstring{5. The global \(L\)-value and the Rallis hypotheses}{5. The global L-value and the Rallis hypotheses}}\label{the-global-l-value-and-the-rallis-hypotheses}

Let \(\pi\) be an actual irreducible cuspidal subrepresentation of the space generated by \(\overline{F_1^+}\), with its actual embedding in automorphic functions. By (S7), its unramified Satake parameters have absolute value one. Hence the partial standard Euler product at \(s=3/2\) is absolutely convergent and nonzero: the logarithms are bounded by a fixed-degree multiple of \[
       \sum_{\mathfrak p}\sum_{r\ge1}
              \frac{N\mathfrak p^{-3r/2}}r<\infty.
 \tag{S18}
\] At a split place this can be viewed as the product over the two primes of \(E\). Enlarging the finite exceptional set has no effect on this conclusion. The finitely many remaining standard local factors are holomorphic and nonzero at \(3/2\) by the same tempered doubling-factor property used in (S17). Thus the completed factor in GQT is \[
                           L(3/2,\pi)\ne0 .
 \tag{S19}
\] There is no central-value nonvanishing conjecture in (S19). In particular one must not substitute \(L(1/2,\pi)\) for this evaluation point.

For clarity, one can avoid any ramified local-LLC factor identification in the last paragraph. GQT Theorem 11.4(ii) itself gives holomorphy of the completed factor once all local theta lifts are nonzero, which has been proved in sections 3--4. Its local factors have no zeros; its partial unramified product is nonzero by (S18). These facts also give (S19). The tempered local-factor property is a direct stronger check.

The target \(I_4\) is positive definite at every real place, so it is anisotropic over \(F\). Its global Witt index is \(r=0\), even though its finite local Witt indices are two. The parameters of GQT Theorem 11.4 are exactly \[
 \epsilon_0=0,\quad n=2,\quad m=4,\quad d(n)=2,\quad
 r=0,\quad s_{m,n}=(m-n)/2=1.
 \tag{S20}
\] Thus \(d(n)<m\le2d(n)\), \(r\le n\), and \(0<s_{m,n}\le d(n)/2\). The lower-tower vanishing assumption for \(j<r\) is empty. Both automorphic quotients are compact; all their automorphic forms are cuspidal because the groups have no proper rational parabolic subgroups.

At every finite place, (S12) and the split theorem give local theta occurrence. GQT Proposition 11.6(i), printed p.~55, then gives the nonzero normalized local zeta functional. At every real place the same conclusion follows from (S17). GQT Theorem 11.4(ii), printed p.~54, and (S19) therefore imply that the second global theta lift of this actual cuspidal embedding of \(\pi\) is nonzero.

With the paper's Haar and Petersson normalizations, its identity is \[
 \langle\Theta_{2,\psi_1}(f_1),
                    \Theta_{2,\psi_2}(f_2)\rangle
       =2\,L(3/2,\pi)
          \prod_v Z_v^*(1,\psi_{1,v}\otimes
                    \overline{\psi_{2,v}},f_{1,v},f_{2,v}).
 \tag{S21}
\] Replacing these measures by probability measures introduces the corresponding positive quotient-volume factors. No such factor has been silently assigned the value one in a height or norm budget. The archimedean factors (S16) are explicit; the finite test functions supplied by local nonvanishing have not been optimized or bounded arithmetically.

\subsection{6. Why every specified first vector survives, including multiplicities}\label{why-every-specified-first-vector-survives-including-multiplicities}

It remains to pass from constituent nonvanishing to the assertion about each specified \(F_1^+\). Let \(\mathcal A\) be the closed \(G(\mathbf A_F)\)-submodule generated by the finite collection of first vectors in question, after (S2). Define the common kernel \[
 \mathcal N=\{F\in\mathcal A:
          \Theta^{\mathrm{bil}}_{2,\psi}(F)=0
                         \text{ for every }\psi\}.
 \tag{S22}
\] Equivalently, intersect the kernels of the scalar functionals \[
 F\longmapsto\int_{[G]}
                    \vartheta_{2,\psi}(g,h)F(g)\,dg
 \quad(\psi,\ h\text{ arbitrary}).
\] They are continuous on \(L^2([G])\): the quotient is compact and a fixed theta kernel is a smooth square-integrable function there. Thus \(\mathcal N\) is closed. Changing variables in the integral moves right translation of \(F\) into the source action on \(\psi\), so \(\mathcal N\) is \(G(\mathbf A_F)\)-invariant.

The compact anisotropic automorphic quotient has discrete unitary spectrum, with finite multiplicities. A nonzero closed invariant subspace of such a Hilbert direct sum contains an irreducible automorphic subrepresentation. This statement allows a subspace inside a multiplicity space: it does not assume multiplicity one. If \(\mathcal N\ne0\), choose such an irreducible copy \(\pi^+\subset\mathcal N\). Its conjugate copy \(\overline{\pi^+}\) is an actual cuspidal embedding to which sections 2--5 apply.

The GQT identity applies to this particular embedding, with its induced Petersson form. Its factorization can always be obtained by scaling one local invariant pairing; that nonzero positive scale cannot change the conclusion. It gives some \(f\in\overline{\pi^+}\) and \(\psi\) for which \(\int\vartheta_{2,\psi}(g,h)\overline{f(g)}\,dg\) is nonzero. But \(\overline f\in\pi^+\subset\mathcal N\), a contradiction. Therefore \(\mathcal N=0\), proving (S4) for each nonzero specified first vector, not only for one favorable spectral projection.

Finally, if its output is nonzero at \(h_0\), replace \(\psi\) by its target translate \(\omega_2(h_0)\psi\). The value at the target identity becomes that previous value at \(h_0\). This does not compute the arithmetic height or denominators of the translated second Schwartz function.

\subsection{7. Precisely stated transfer to another CM host}\label{precisely-stated-transfer-to-another-cm-host}

The argument is structural over a CM extension \(K/K^+\). It applies with the same ranks if the following inputs have actually been supplied:

\begin{enumerate}
\def\labelenumi{\arabic{enumi}.}
\tightlist
\item
  An actual first \(U(1)\to U(I_2)\) theta vector from unitary characters and polynomial Fock vectors, with globally unitary splitting characters.
\item
  A globally admissible fixed character twist making each real first type \(\det^2\operatorname{Sym}^{d_\tau}\), with the same Hom and bilinear integral conventions.
\item
  The positive form \(I_4\) is hyperbolic at every nonsplit finite place.
\end{enumerate}

No torus class-number calculation enters the second-step proof once input 1 exists. If \(K/K^+\) is unramified at every finite nonsplit place, input 3 follows immediately from the norm argument in section 3.1. AM-R-009 supplies the intended inputs for \(K=\mathbf Q(\zeta_{28})\): its fixed twist has \(\mathbf c=(2,1,2,2,1,2)\), and its degree-eight vector at label9 has the actual first identity value \((9/4)C^4\). These inputs have now been independently whole-checked in review\_root\_unramified\_first\_lift.md, retaining the accepted DP57 and DO022 host facts. Thus there is the concrete corollary \[
 \exists\psi:\quad
 \Theta^{\mathrm{bil}}_{U(2)\to U(4),\psi}
       (\Theta'_P)\ne0,\qquad
 \Theta'_P(1)=\frac94C^4,
 \tag{S23}
\] over this same \(E'\), with the actual first source input of norm one from (R9.27). At infinity the degrees are \(d_9=8\) and zero at the other five places, so the unnormalized compact numerator in (S16), multiplied over all six places, is \(\prod_{\tau}\operatorname{vol}(U(2)_\tau)/9\). This does not normalize the finite second test functions. The field and lattice normalizations in this corollary are the stated \(E'\) data, not an identification of \(E'\) with the original \(E\) in (S1).

\subsection{8. Remaining gate and costs}\label{remaining-gate-and-costs}

This proves a genuine nonzero two-edge composition with an unrestricted actual second Schwartz input. For the sharp first coefficients of ALP-LIFT-10, it proves that the same nonzero first vector is not annihilated by every second kernel. It does not prove that its second value has the first coefficient's valuation.

The next arithmetic gate is to choose that second kernel inside the permitted finite moved family, prove its integral lattice and content, and bound the actual mass and archimedean norm of the nonzero pairing. The full common-character pins, stage row modules, harmonic branching maps and six terminal identities are also absent from this two-edge result. In particular the rank-one starting edge is still not the source's even-rank StdPair constructor. No shorter DP44 chain or net DP44 gain is claimed from (S4) or (S21).

\subsection{Primary references and exact locators}\label{primary-references-and-exact-locators}

\begin{itemize}
\tightlist
\item
  Gan, W. T., Qiu, Y., \& Takeda, S. (2014). The regularized Siegel--Weil formula (the second term identity) and the Rallis inner product formula. arXiv:1207.4709v3. Section11.2--11.3, printed pp.51--52, the conjugated-input integral and doubled Schwartz map; section11.6--11.7, printed p.54, normalized local factors and Theorem11.4; Proposition11.6(i), printed p.55; Theorem11.7, pp.56--57, the real-place caveat deliberately bypassed. \url{https://arxiv.org/pdf/1207.4709v3}
\item
  Gan, W. T., \& Ichino, A. (2015). The Gross--Prasad conjecture and local theta correspondence. arXiv:1409.6824v2. Section2, bounded parameters and tempered representations; Theorem4.4(i)(a), (ii)(b), (iii), printed p.14, the almost equal rank parameter and irreducibility; section6.7, printed p.27, the exact tempered local doubling-factor property of Lapid--Rallis Proposition5. \url{https://arxiv.org/pdf/1409.6824}
\item
  Mínguez, A. (2008). Correspondance de Howe explicite: paires duales de type II. Annales scientifiques de l'École Normale Supérieure, 41(5), 717--741. Theorem1(1)--(2), printed p.718; Corollary6.3 is its later proved version. The theorem includes every nonarchimedean residue characteristic. \url{https://www.numdam.org/item/ASENS_2008_4_41_5_717_0.pdf}
\item
  Lapid, E. M., \& Rallis, S. (2005). On the local factors of representations of classical groups. In Automorphic representations, L-functions and applications: Progress and prospects (pp.309--360). De Gruyter. Proposition5, used in exactly the form explicitly recorded in Gan--Ichino section6.7. The primary Gan--Ichino passage was directly read; no independent full reading of the Lapid--Rallis chapter is claimed. \url{https://doi.org/10.1515/9783110892703.309}
\end{itemize}

Internal dependencies: ALP-LIFT-09, (V22)--(V38); ALP-LIFT-10, actual finite moved family and source inputs; ALP-LIFT-11, (R23)--(R26); AM-DO-023, corrected real source action and the trivial-character second case. Check events are recorded separately, only after whole proof readings. This proof has not been formalized in Lean.

\clearpage

\section{AM-DO-025: the actual common-character compact chain and its exact degrees}\label{am-do-025-the-actual-common-character-compact-chain-and-its-exact-degrees}

Author: GPT-6 Astra Pro, ALP\_MACHINE \nolinkurl{/root/definite_oscillator}; inherited reasoning effort not separately exposed; 21 November 2026 UTC.

The selected increasing even-rank chain in the recovered source has an exact local compact-type match. The failure of the unchanged trivial-character second edge in AM-DO-023 is therefore not an obstruction to this selected chain. The computation below uses the source's actual character exponents, bilinear theta convention, and coefficient weights. It also computes the polynomial degree required at each edge. It does not prove global nonvanishing or the arithmetic row identities.

\subsection{1. Source data and conventions}\label{source-data-and-conventions}

Fix one selected complex embedding and an actual \nolinkurl{HCParam} for ranks

\[
 d=t+1,\qquad N_j=k+jd,\qquad 0\le j\le m.
 \tag{DO25.1}
\]

Here \(k,m\) are even, \(t\) is odd, and the other source rank conditions are retained. The real numbers \(e_j,b_j\) are integers, the \(u_i\) are half-integers, and

\[
 u_{k-1}>-b_0,\qquad
 e_j-b_j\ge\frac{d+2}{2},\qquad
 b_{j+1}-e_j\ge\frac{d+2}{2}.
 \tag{DO25.2}
\]

These are \nolinkurl{Params.lean}, \nolinkurl{HCParam}, lines 50--69, with the decreasing \(u_i\) forced by its interlacing inequalities. The actual character on edge \(j\to j+1\) has infinity exponent \(-e_j\): \nolinkurl{ChainParSucc.lean}, \nolinkurl{Com.Is.exp}, and \nolinkurl{DataTheta.lean}, \nolinkurl{ThetaX.chrOf_inf}. This means \(\mu_{j+1}(z)=(z/\bar z)^{-e_j}\).

Use \textbf{equal source labels} \(\chi_S=\chi_T=\mu_{j+1}\) on this definite edge, exactly as required by \nolinkurl{PinColT.om} and \nolinkurl{PinX.mv} in \nolinkurl{LevChiEqOne.lean}. In conventional Konno--Konno parameters the pair is \((\mu_{j+1},\mu_{j+1}^{-1})\). The inverse is necessary because the recovered \nolinkurl{splitFacO} is \(\chi_T(x)^{-1}\chi_S(y)\) on determinant roots. This conversion is checked and corrected explicitly in AM-DO-023.

For source rank \(a=N_j\) and target rank \(b=N_{j+1}\), the oscillator's holomorphic Fock module is consequently

\[
 \det_a^{e_j-b/2}\boxtimes\det_b^{-e_j+a/2}
       \otimes\operatorname{Sym}(\mathbf C^b\otimes(\mathbf C^a)^*).
 \tag{DO25.3}
\]

The source theta operator is bilinear:

\[
                  (T_\phi f)(h)=\int_{[U(a)]}
                          \vartheta_\phi(g,h)f(g)\,dg.
 \tag{DO25.4}
\]

This is literally \nolinkurl{liftT_apply} in \nolinkurl{Development/AbelianVarietyTheory/Definitions/TwistedInnerProducts/WittSmallGroupEmbeddingMemPointwiseStabilizerIso.lean}, lines 160--167, and \nolinkurl{KerDat.lift} in \nolinkurl{JointSmooth.lean}, lines 79--82. Thus the source oscillator must contain the \textbf{contragredient} of the input coefficient type. Confusing (DO25.4) with an integral against \(\bar f\) reverses this type test.

\subsection{2. Identification with the actual selected coefficient partitions}\label{identification-with-the-actual-selected-coefficient-partitions}

Write \(w_j=(w_{j,0},\ldots,w_{j,N_j-1})\) for the source's dominant integral weight \nolinkurl{Wts.lamTrue u j τ}. Write \(\nu_j\) for \nolinkurl{Wts.nuW u (j+1) τ}. The recovered source already proves the precise identities

\[
 \nu_{j,i}=w_{j,i}+e_j-\frac{N_{j+1}}2,
 \tag{DO25.5}
\]

and

\[
 w_{j+1}=
 \left(w_j-\frac d2\mathbf1_{N_j},\,
       \left(-e_j+\frac{N_j}2\right)\mathbf1_d\right).
 \tag{DO25.6}
\]

Exact locators are \nolinkurl{Formalism/LocalToGlobal/Columns/AbelianHostFieldConditions.lean}, \nolinkurl{Wts.lamN}, \nolinkurl{lamTrue}, and \nolinkurl{nuW}, lines 462--469; and \nolinkurl{Formalism/LocalToGlobal/Columns/BaseLemmas.lean}, \nolinkurl{lamTrue_old}, lines 496--510, \nolinkurl{lamTrue_new}, lines 512--530, and \nolinkurl{nuW_eq_true}, lines 542--552. These source identities are used as settled equalities, rather than replaced by freely chosen weights.

The \(\nu_j\) are actual nonnegative partitions, as proved in §3 below. The Schur summand indexed by \(\nu_j\) in (DO25.3) has source type

\[
 \det_a^{e_j-b/2}\otimes S_{\nu_j}((\mathbf C^a)^*),
 \tag{DO25.7}
\]

whose highest weight, in decreasing order, is

\[
 \left(e_j-b/2-\nu_{j,a-1},\ldots,
                    e_j-b/2-\nu_{j,0}\right)
               =(-w_{j,a-1},\ldots,-w_{j,0}).
 \tag{DO25.8}
\]

This is precisely the contragredient required by (DO25.4). The target type of the same Schur summand is

\[
 \det_b^{-e_j+a/2}\otimes S_{\nu_j}(\mathbf C^b).
 \tag{DO25.9}
\]

Its highest weight consists of \(-e_j+a/2+\nu_{j,i}\) for the first \(a\) entries and \(-e_j+a/2\) on the final \(d\) entries. Substitution of (DO25.5) gives exactly (DO25.6). Therefore every selected adjacent pair is a matching local compact theta type with the \textbf{same prescribed character on both members of the edge}.

The Schur summand occurs with multiplicity one under the compact product group. This is the polynomial Cauchy decomposition \(\operatorname{Sym}(T\otimes S^*)=
\bigoplus_\nu S_\nu(T)\otimes S_\nu(S)^*\). The primary oscillator reference is Konno and Konno (2007), §5.1, Fact 5.1, Lemma 5.2, the highest vectors (5.4)--(5.5), and Lemma 5.3, printed pp.72--74, with the character convention fixed in AM-DO-023. Since \(b>a\), every partition of length at most \(a\) appearing above is permitted. Its total homogeneous degree is \(|\nu_j|\). No determinant twist external to the selected source data has been introduced.

This is a representation-theoretic occurrence statement. An equivariant global theta map can still vanish on an occurring local type; the computation does not assert otherwise.

\subsection{3. Exact recurrence and gap consequences}\label{exact-recurrence-and-gap-consequences}

Define

\[
       \Delta_j=e_{j+1}-e_j-d\qquad(0\le j<m-1).
 \tag{DO25.10}
\]

Then the selected polynomial partitions satisfy the exact recurrence

\[
 \boxed{\nu_{j+1}
       =(\nu_j+\Delta_j\mathbf1_{N_j},\,
                         \Delta_j\mathbf1_d).}
 \tag{DO25.11}
\]

Moreover

\[
                  \nu_{0,k-1}\ge1,\qquad \Delta_j\ge2.
 \tag{DO25.12}
\]

In particular all displayed weights \(\nu_j\) are nonnegative partitions, with strictly positive final entry.

\textbf{Proof.} Apply (DO25.5) at \(j+1\), and insert (DO25.6). For an old entry the change from \(\nu_{j,i}\) is

\[
 -\frac d2+e_{j+1}-\frac{N_{j+2}}2-e_j
                                      +\frac{N_{j+1}}2
      =e_{j+1}-e_j-d=\Delta_j.
\]

For a new entry, its value is \(-e_j+N_j/2+e_{j+1}-N_{j+2}/2=\Delta_j\). This proves (DO25.11).

The two gaps in (DO25.2) add to \(e_{j+1}-e_j\ge d+2\), so \(\Delta_j\ge2\). At stage zero the source weight is \(w_{0,i}=u_i-\rho_k(i)\), where \(\rho_k(i)=(k-1)/2-i\). This is \nolinkurl{lamN_eq} and \nolinkurl{lamTrue_zero} in \nolinkurl{BaseLemmas.lean}, together with \nolinkurl{rhoVec} in \nolinkurl{DataCoeff.lean}, line 39. Hence

\[
 \nu_{0,k-1}=u_{k-1}+e_0-\frac{d+1}{2}.
 \tag{DO25.13}
\]

Because \(u_{k-1}\) is a half-integer, \(b_0\) is an integer, and \(u_{k-1}>-b_0\), one has \(u_{k-1}+b_0\ge1/2\). Adding \(e_0-b_0\ge(d+2)/2\) proves (DO25.13) is at least one. The starting weight is dominant, so all its entries are at least one. Recurrence (DO25.11) adds a nonnegative constant to every old entry and appends entries no greater than those old entries, proving dominance and positivity inductively. All entries are integers by the parity and integrality conditions. This proves (DO25.12) and the partition assertion. \(\square\)

\subsection{4. Exact degree evolution and its sharp rank-only lower bound}\label{exact-degree-evolution-and-its-sharp-rank-only-lower-bound}

Let \(D_j=|\nu_j|\), the homogeneous degree at edge \(j\to j+1\) for this embedding. Then

\[
 \boxed{D_0=\sum_{i=0}^{k-1}u_i
                   +k e_0-\frac{k(k+d)}2,\qquad
        D_{j+1}=D_j+N_{j+1}\Delta_j.}
 \tag{DO25.14}
\]

Therefore

\[
 \boxed{D_j\ge k(2j+1)+d\,j(j+1)}
 \tag{DO25.15}
\]

and across the \(m\) selected edges,

\[
\begin{split}
 \sum_{j=0}^{m-1}D_j
   &=mD_0+\sum_{r=0}^{m-2}(m-1-r)N_{r+1}\Delta_r,\\
 \sum_{j=0}^{m-1}D_j
   &\ge k m^2+\frac d3(m^3-m).
\end{split}
 \tag{DO25.16}
\]

\textbf{Proof.} Summing (DO25.5) at \(j=0\), the entries of \(\rho_k\) sum to zero, giving the first equality in (DO25.14). Summing (DO25.11) gives the second. Since \(D_0\ge k\), \(\Delta_r\ge2\), and \(\sum_{\ell=1}^{j}N_\ell=jk+d j(j+1)/2\), induction gives (DO25.15). Summing the exact recurrence counts the increment at index \(r\) exactly \(m-1-r\) times and gives the first line of (DO25.16). Finally \(\sum_{j=0}^{m-1}(2j+1)=m^2\) and \(\sum_{j=0}^{m-1}j(j+1)=(m^3-m)/3\), proving its second line. \(\square\)

The lower bounds are attained among the displayed \nolinkurl{HCParam} conditions. Indeed choose the \(u_i\) in decreasing consecutive half-integers, take \(u_{k-1}=-b_0+1/2\), set \(e_0-b_0=(d+2)/2\), and take both successive gaps in (DO25.2) equal to \((d+2)/2\). The interlacing integer \(v_i\) is then the unique midpoint between \(u_i\) and \(u_{i+1}\). These choices have \(\nu_0=\mathbf1_k\) and \(\Delta_j=2\), and every entry is of the required parity. The partitions become

\[
 \nu_j=((2j+1)^k,(2j)^d,(2j-2)^d,\ldots,2^d),
 \tag{DO25.17}
\]

where exponents denote repeated parts and the last list is empty for \(j=0\). The source's further \nolinkurl{Hyp3} noncompact-gap conditions may force larger spacings. The sharpness statement here concerns the stated \nolinkurl{HCParam} constraints, not a claimed sharp construction for all of \nolinkurl{Hyp3}.

The degree in (DO25.14) is the degree of the selected Schur summand, not an arbitrary direction used to illustrate a formal-group slope. It is a usable exact input to subsequent derivative, norm, and denominator estimates. It does not itself measure a net arithmetic gain: the finite family, integral row lattice, normalization periods, class masses, and terminal detector remain to be inserted. On the \(E'\) realization, the dyadic eigenbasis change from AM-DO-022 still has index \(2^3\) and degree-\(D\) denominator bound \(2^D\); no type calculation removes that basis cost.

\subsection{5. Exact source locators and retained gates}\label{exact-source-locators-and-retained-gates}

All paths below are relative to the recovered source root \nolinkurl{work/research/alpoge_source}.

\begin{longtable}[]{@{}
  >{\raggedright\arraybackslash}p{(\columnwidth - 2\tabcolsep) * \real{0.5000}}
  >{\raggedright\arraybackslash}p{(\columnwidth - 2\tabcolsep) * \real{0.5000}}@{}}
\toprule\noalign{}
\begin{minipage}[b]{\linewidth}\raggedright
Statement used
\end{minipage} & \begin{minipage}[b]{\linewidth}\raggedright
Exact source
\end{minipage} \\
\midrule\noalign{}
\endhead
\bottomrule\noalign{}
\endlastfoot
Ranks and gap hypotheses & \nolinkurl{Formalism/LocalRingCardinality/Params.lean}, \nolinkurl{Ranks.N}, lines 35--36; \nolinkurl{HCParam}, lines 50--69; \nolinkurl{sExp}, \nolinkurl{sExp'}, \nolinkurl{wG}, lines 71--75 \\
Compact Harish--Chandra parameter & \nolinkurl{Formalism/LocalRingCardinality/DataTheta.lean}, \nolinkurl{TPar.Xpar}, lines 545--548 \\
Character infinity type & \nolinkurl{DataTheta.lean}, \nolinkurl{ThetaX.chrOf_inf}; \nolinkurl{Development/AbelianVarietyTheory/Definitions/LatticeFrames/ChainParSucc.lean}, \nolinkurl{Com.Is.exp}, line 40 \\
Coefficient definitions & \nolinkurl{Formalism/LocalToGlobal/Columns/AbelianHostFieldConditions.lean}, \nolinkurl{Wts.lamTrue}, \nolinkurl{Wts.nuW}, lines 464--469 \\
Exact selected recurrence & \nolinkurl{Formalism/LocalToGlobal/Columns/BaseLemmas.lean}, \nolinkurl{lamTrue_old}, \nolinkurl{lamTrue_new}, \nolinkurl{nuW_eq_true}, lines 496--552 \\
Bilinear upward kernel & \nolinkurl{Development/AbelianVarietyTheory/Definitions/TwistedInnerProducts/WittSmallGroupEmbeddingMemPointwiseStabilizerIso.lean}, \nolinkurl{liftT_apply}, lines 160--167 \\
Equal characters on the definite edge & \nolinkurl{Development/AbelianVarietyTheory/Definitions/LatticeFrames/LevChiEqOne.lean}, \nolinkurl{PinX.mv}, lines 100--101, and \nolinkurl{PinColT.om}, line 110 \\
Integral polynomial convention & \nolinkurl{Formalism/LocalRingCardinality/DataCoeff.lean}, \nolinkurl{rhoVec}, \nolinkurl{hcToWeightQ}, \nolinkurl{DomWeight.dual}, and the Weyl-module definitions \\
\end{longtable}

No Lean compilation is claimed. The source equality bodies cited in §2 were directly read; the new assertions are the complete oscillator/type identification, explicit partition recurrence, and degree formulas proved above. Membership of the desired arithmetic row in the \textbf{global} iterated image, the six actual \nolinkurl{AdmS} statements, and the \nolinkurl{MainX} terminal equation remain distinct obligations.

Primary reference: Konno, T., \& Konno, K. (2007). On doubling construction for real unitary dual pairs. \emph{Kyushu Journal of Mathematics, 61}(1), 35--82. \url{https://doi.org/10.2206/kyushujm.61.35}

\clearpage

\section{AM-DO-026: actual mixed-kernel operators and the remaining source equations}\label{am-do-026-actual-mixed-kernel-operators-and-the-remaining-source-equations}

Author: GPT-6 Astra Pro, ALP\_MACHINE \nolinkurl{/root/definite_oscillator}; inherited reasoning effort not separately exposed; 21 November 2026 UTC.

This finding combines the proved mixed-CM construction with the literal recovered source interfaces. It proves an all-polynomial geometric kernel identity, its compact covariance, a concrete finite stabilizer, and the exact adjoint and class-mass formulas for the resulting actual operators. It then identifies the additional equations needed to obtain the source \nolinkurl{PinColT}, \nolinkurl{TieIo}, and terminal identities. These last equations are not inferred from one nonzero theta value.

The carried inputs are AM-DO-017--020 and AM-DO-022--025, AM-R-004--006, AM-K018, and the accepted DP57 host. Their conclusions are used at their recorded checked status. In particular the finite direct-average phase and pairing are already proved in AM-DO-024; they are not re-proved here.

\subsection{1. A geometric kernel identity for every polynomial and finite translate}\label{a-geometric-kernel-identity-for-every-polynomial-and-finite-translate}

Use \[
 E'=\mathbf Q(\zeta_{28}),\quad F'=E'^+,
 \quad A=J^2,\quad \xi'=\xi/2,
 \quad C=\theta_0c.
 \tag{DO26.1}
\] For positive even ranks \(a,b\), let \(S=E'^a,T=E'^b\) have the positive identity Hermitian forms. All Heisenberg signs, measures, and splitting labels are those of AM-DO-022. Let \(\Lambda=\operatorname{Hom}_{\mathcal O_{E'}}
(\mathcal O_{E'}^a,\mathcal O_{E'}^b)\), with its actual perfect trace pairing, and let \(\phi_0\) be its finite trace-one vacuum.

For moved lattices \(N,N'\), define \(F_{N,N'}\) \textbf{geometrically} by the two normalized finite Mumford averages from the seed \((\sigma_q\boxtimes\sigma_q)^{\boxtimes ab}\). Thus each average uses the actual finite quotient, its trace-one theta-group multiplier, and the positive factor \(q^{-1/2}\). This definition precedes the analytic theta comparison. For finite \(g,h\) moving the standard lattices to \(N,N'\), write \[
 A_{g,h}X=hXg^{-1},\qquad
 \epsilon_\chi(g,h)=\chi_T(x)^{-1}\chi_S(y),
 \quad x/\bar x=\det g,\quad y/\bar y=\det h.
 \tag{DO26.2}
\] The independence from the determinant roots is the even-rank restriction \(\chi_S|_{\mathbb A_{F'}^\times}
=\chi_T|_{\mathbb A_{F'}^\times}=1\).

For clarity, a finite adelic translation parameter does not itself define a complex tangent vector. If \(y\) is such a parameter and \(\Lambda'\) is the rational period lattice of \(F_{N,N'}\), choose \(r\in\Lambda'\otimes\mathbf Q\) with \(\delta=y-r\in\widehat\Lambda'\). Additive approximation supplies such \(r\). Define the actual adelic translation on this theta line by \[
 \mathcal U_y^{\Lambda'}F
   =\psi_f(-b(r,\delta)/2)\alpha_{\Lambda'}(\delta)\,U_rF.
 \tag{DO26.2a}
\] The finite semicharacter is the continuous extension of the actual trace-one lattice character. The Weyl law \(W_rW_\delta=\psi_f(b(r,\delta)/2)W_y\), together with \(W_\delta\phi_{N,N'}=\alpha_{\Lambda'}(\delta)\phi_{N,N'}\), proves that (DO26.2a) is independent of the chosen rational representative and that \(\mathcal R W_y\phi_{N,N'}=C^{2ab}
\mathcal U_y^{\Lambda'}F_{N,N'}\). Equivalently, independence follows from the reduced-translation group law and the rational product formula. All displayed central factors are retained.

Choose a finite collection of rational Heisenberg parameters \(x\) and coefficients \(d_x\), and set \[
 \phi_d=\sum_x d_xW_x\phi_0,\qquad
 F_{d;g,h}=\sum_x d_x\mathcal U_{A_{g,h}x}^{\Lambda'}F_{N,N'}.
 \tag{DO26.3}
\] Here \(W_x\) and \(\mathcal U_y^{\Lambda'}\) are the actual Schrödinger and reduced Mumford translations, with their central phase retained. In particular no finite adele has been silently treated as a complex translation vector. For a polynomial \(P\) in the actual invariant tangent coordinates, let \(r_P\) be its corresponding Chern creation operator. Then \[
 \boxed{\vartheta_{r_P\phi_d}(g,h)
   =C^{2ab}\epsilon_\chi(g,h)
                  \,[P(\partial)F_{d;g,h}]_{0}.}
 \tag{DO26.4}
\] This is one common period for the entire moved family and all polynomial degrees, not a separate scalar selected at each value.

\textbf{Proof.} Exact Heisenberg covariance gives \[
 \omega(g,h)\phi_d
 =\epsilon_\chi(g,h)\sum_xd_xW_{A_{g,h}x}\phi_{N,N'}.
\] AM-DO-022 gives the vacuum identity with period \(C^{2ab}\), and AM-DO-018 identifies the actual translations under the reading map. Therefore \(\mathcal R(\omega(g,h)\phi_d)
=C^{2ab}\epsilon_\chi(g,h)F_{d;g,h}\) as entire functions. The finite action commutes with every archimedean creation operator, because they act on different factors of the restricted tensor product. The explicit identity \(\mathcal Rr_P=P(\partial)\mathcal R\), checked in ALP-LIFT-09, then gives (DO26.4) by evaluation at zero. All sums in (DO26.3) are finite. The theta sum and its fixed-order derivatives converge normally by the Gaussian estimate already proved for the reading map. Hence no interchange of divergent expressions is used. \(\square\)

AM-K018 identifies the starting Chern-reduced formal gauge with the actual source primitive gauge for the underlying ordinary product of copies of \(J\), including the integral block at 11. The extra \(E'\)-action does not replace that product by a product with conjugated \(E\)-actions. Its complete generic gauge comparison also transports through the actual CM isogenies and torsion translations in these finite averages: these preserve the same holomorphic character set, so the same no-opposite-characters uniqueness argument applies. This does not supply an integral 11-adic frame after an arbitrary isogeny of 11-power degree. The accepted original integral frame transports directly only when that isogeny is prime to 11, or when an additional integral comparison is proved. No reduced theta frame is asserted integral. The product trace remains \(\xi/2\), and the eigenbasis change has dyadic determinant valuation 3 per \(A\), as in AM-DO-022.

For specified finite class representatives \(g_i,h_j\), define the geometric jet functionals \[
 \mathcal A_{ji}(P)=\epsilon_\chi(g_i,h_j)
                    [P(\partial)F_{d;g_i,h_j}]_0.
 \tag{DO26.5}
\] They are defined from actual line-bundle translations and averages, not by dividing an analytic theta value by a desired constant. For algebraic finite characters and coefficients they have algebraic values, by AM-K018 and algebraicity of the fixed torsion translations. For any chosen source level volume \(v_S\), (DO26.4) gives the concrete all-polynomial class-kernel equation \[
 v_S\vartheta_{r_P\phi_d}(g_i,h_j)
                   =v_SC^{2ab}\mathcal A_{ji}(P).
 \tag{DO26.6}
\] This is the exact mixed-family analogue of \nolinkurl{ClsK.Core}. Identification with its independently specified \nolinkurl{abelian₂}, frame scale, and chart characters is an additional comparison, detailed in section 5.

\subsection{2. The actual compact covariance and an explicit finite stabilizer}\label{the-actual-compact-covariance-and-an-explicit-finite-stabilizer}

For the selected common-character edge of AM-DO-025, put \(a=N_j,b=N_{j+1}\) and \(\chi_S=\chi_T=\mu_{j+1}\), with infinity exponent \(-e_j\). The corrected source splitting formula gives \[
 \omega_\infty(g,h)r_P\phi_\infty
 =\prod_\tau\det(g_\tau)^{e_{j,\tau}-b/2}
                    \det(h_\tau)^{-e_{j,\tau}+a/2}
      r_{\operatorname{mvZ}(g^{-1},h)P}\phi_\infty.
 \tag{DO26.7}
\] The polynomial action is expressed in the source's creation-label convention \nolinkurl{mvZ}; it is the tensor action on \(T\otimes S^*\). Equation (DO26.7) is AM-DO-023's honest-splitting formula applied to all creation polynomials, not merely to a lowest-weight line.

The exact recovered definitions give \[
 s_S=e_j-b/2=\texttt{sSharpC},\qquad
 s_D=e_j-a/2=\texttt{sDivC}.
 \tag{DO26.8}
\] Indeed \nolinkurl{Params.sExp'} and \nolinkurl{sExp} have these values at \(j+1\); their integer floors in \nolinkurl{Columns/Record.lean}, lines 1083--1087, are unchanged because \(e_j\) is integral and both ranks are even. Thus (DO26.7) is exactly the determinant and polynomial covariance required by \nolinkurl{ClsK.Real}, after matching the actual tangent frames. AM-DO-025 proves that the source partition \nolinkurl{nuW} selects the contragredient of \nolinkurl{lamTrue} on the source and the next \nolinkurl{lamTrue} on the target. The real covariance and the local type match are therefore established for this constructed compact family.

Here is also a definite finite-level assertion. Suppose every parameter in (DO26.3) belongs to \(M^{-1}\Lambda\), where \(M\) is a positive integer. Let \(K_S,K_T\) be the standard integral unitary compacts, refined by the principal congruence condition \(1\pmod{2M^2}\). For trivial characters, both groups fix \(\phi_d\) exactly: \[
                     \omega(K_S,K_T)\phi_d=\phi_d.
 \tag{DO26.9}
\] For nontrivial even-rank characters the same assertion holds after intersecting with the kernels of the two determinant characters \(\nu_{\chi_T}^{-1}\circ\det\) and \(\nu_{\chi_S}\circ\det\). These are open subgroups of the compacts, so the refinement is finite.

\textbf{Proof.} AM-DO-022 proves that the trivial-character standard compacts fix \(\phi_0\). The induced Hom action is congruent to the identity modulo \(2M^2\), including after inversion on the source. For \(x\in M^{-1}\Lambda\), write \(\delta=A_{g,h}x-x\in2M\widehat\Lambda\). Perfection and integrality of the actual trace pairing give \[
 \psi_f(b(x,\delta)/2)=1,\qquad
 \alpha_{\rm tr}(\delta)
      =\psi_f(b(\delta,w_0\delta)/2)=1.
\] For the first equality write \(x=\lambda/M,\delta=2M\eta\); the half-pairing is \(b(\lambda,\eta)\), which is integral. For the second it is \(2M^2b(\eta,w_0\eta)\), also integral. The exact Weyl group law therefore yields \(W_{x+\delta}\phi_0=W_x\phi_0\). Covariance fixes each summand in (DO26.3), proving (DO26.9). For the specified nontrivial characters the only remaining compact scalar is the displayed determinant factor from AM-DO-022; killing it gives the assertion. \(\square\)

This bound is deliberately uniform for the named finite span. It is not asserted to be the least level; ALP-LIFT-10 obtains a smaller one for its special isotropic half-orbit. Nor does (DO26.9) prove that an independently pinned source level already lies in this stabilizer. The index of the needed refinement remains an arithmetic cost.

\subsection{3. Exact adjoint and conjugation identities for the compact kernel}\label{exact-adjoint-and-conjugation-identities-for-the-compact-kernel}

Let \(X=[U(S)]\) and \(Y=[U(T)]\), with specified compatible quotient Haar measures. Both are compact because the groups are definite. Let \(K(x,y)\) be one of the actual polynomial theta kernels from section 1. Define the source-style \textbf{bilinear} operators \[
 (T_Kf)(y)=\int_XK(x,y)f(x)\,dx,
 \qquad
 (S_Kg)(x)=\int_YK(x,y)g(y)\,dy.
 \tag{DO26.10}
\] Use the inner product conjugate-linear in its first argument, and write \(C_Xf=\bar f\), \(C_Yg=\bar g\). Then the exact identities are \[
 \boxed{T_K^*=C_XS_KC_Y,\qquad
        C_YT_KC_X=T_{\bar K}.}
 \tag{DO26.11}
\] The second equation keeps the same source-to-target orientation; in the first, \(S_K\) reverses it. In particular \nolinkurl{liftS} with the same kernel is not the Hermitian adjoint of \nolinkurl{liftT} until these pointwise conjugations are inserted.

\textbf{Proof.} At fixed level the quotient is a finite union of compact real-group pieces. The finite support is fixed on each piece, and the positive Gaussian quadratic form has a uniform positive lower bound as the real group varies over its compact set. Its polynomial theta series therefore converges normally and gives a continuous bounded kernel. It is square integrable, and Fubini applies to the displayed operators, first for continuous functions and then to all \(L^2\) functions by the bounded integral-operator estimate. Directly, \[
 \begin{aligned}
 \langle T_Kf,g\rangle_Y
  &=\int_X\bar f(x)
       \left(\int_Y\overline{K(x,y)}g(y)\,dy\right)dx\\
  &=\langle f,C_XS_KC_Yg\rangle_X.
 \end{aligned}
\] Conjugating (DO26.10) term by term gives the second equality. This proves (DO26.11). \(\square\)

The corresponding bilinear pairing \(B_X(f_1,f_2)=\int_Xf_1f_2\) satisfies \[
                  B_Y(T_Kf,g)=B_X(f,S_Kg).
 \tag{DO26.12}
\] Consequently \[
             \langle f,T_K^*T_Kf\rangle_X=\|T_Kf\|_Y^2.
 \tag{DO26.13}
\] This is strictly positive for an actual nonzero first image. It constructs an actual detecting Gram operator; it does not identify that operator with a separately prescribed positive-rank next theta edge. Complex conjugating the theta series also gives a concrete kernel in the conjugate Weil representation, with the additive character and splitting scalars conjugated. This supplies the conjugations in (DO26.11) without guessing a source-target swapping character formula.

The scalar bookkeeping is exact. If \(K=c_0\mathcal K\), then \[
 T_K=c_0T_{\mathcal K},\qquad
 T_K^*=\bar c_0T_{\mathcal K}^*,\qquad
 T_K^*T_K=|c_0|^2T_{\mathcal K}^*T_{\mathcal K}.
 \tag{DO26.14}
\] Here \(c_0=C^{2ab}\), with any explicitly chosen creation-frame scaling retained. Neither the complex-torus probability measure of AM-DO-024 nor a Gaussian unit normalization silently substitutes for the automorphic quotient measure in these formulas.

For a compact subgroup embedding \(\iota:H\to G\), the actual restriction period \(B_{\rm top}(F,G)=\int_{[H]}\overline{F(\iota h)}G(h)\,dh\) obeys \[
 B_{\rm top}(C F,G)=\int_{[H]}F(\iota h)G(h)\,dh.
 \tag{DO26.15}
\] Thus the source's \nolinkurl{conjAF} in \nolinkurl{TieBm} has a concrete role: it converts its conjugate-linear first argument into the bilinear terminal pairing. Equation (DO26.15) does not specify the arithmetic terminal row or its scalar.

\subsection{4. Actual finite class masses and the operator matrix}\label{actual-finite-class-masses-and-the-operator-matrix}

Fix an open compact \(K_f\subset G(\mathbb A_f)\), and write \(g_i\) for the finitely many double-class representatives. Put \[
 \Gamma_i=G(F)\cap g_iK_fg_i^{-1},\qquad
       v_K=\operatorname{vol}(K_f),\qquad w_i=v_K/|\Gamma_i|.
 \tag{DO26.16}
\] The real compact group has Haar volume one. The rational intersection is finite: it is discrete in the compact product of the real group and the finite compact, hence finite. Quotient integration gives \[
 \int_{[G]}F(g)\,dg
       =v_K\sum_i\frac1{|\Gamma_i|}
                    \int_{G_\infty}F(g_ik)\,dk.
 \tag{DO26.17}
\] Indeed the piece is the diagonal quotient of \(G_\infty\times K_f\) by the finite group \(\Gamma_i\). Pushing forward Haar measure divides by its order; then integrating the right \(K_f\)-invariant function contributes \(v_K\). This proves (DO26.17) with the chosen compatible quotient measure.

If \(\rho\) is an irreducible unitary real coefficient type of dimension \(d_\rho\), and a unit functional \(\ell\) is used to realize its vectors as \(k\mapsto\ell(\rho(k)v)\), then \[
 \int_{G_\infty}|\ell(\rho(k)v)|^2\,dk
                          =\frac{\|v\|^2}{d_\rho}.
 \tag{DO26.18}
\] For completeness, the averaged positive rank-one operator \(\int\rho(k)^*\ell^*\ell\rho(k)\,dk\) commutes with \(\rho\), hence is scalar by irreducibility; its trace is one, so its scalar is \(1/d_\rho\). This proves (DO26.18). Multiplication of this realization by \(\sqrt{d_\rho}\) makes it isometric. Without that normalization, the factor in (DO26.18) must remain in every class metric.

Choose orthonormal real compact coefficient bases in each class piece, using that normalization and its actual invariant subspaces. Then the metric on the class coefficient vector is \(G_X=\operatorname{diag}(w_i I_i)\). Do the same on the target, with metric \(G_Y\). Expand the relevant finite-type kernel as \[
 K((i,k),(j,l))
       =\sum_{r,s}A_{ji,sr}\overline{u_{i,r}(k)}v_{j,s}(l).
 \tag{DO26.19}
\] Such an expansion is finite for a fixed polynomial input: its compact-group orbit is finite dimensional. Orthogonality and (DO26.17) now give the \textbf{actual operator matrices} \[
 \boxed{T=A G_X,\qquad
        T^{\dagger}=G_X^{-1}\bar T^{\,t}G_Y
                         =\bar A^{\,t}G_Y.}
 \tag{DO26.20}
\] The formulas hold blockwise on the invariant coefficient subspaces; empty blocks are omitted. They follow simply by integrating \(\bar u_{i,r}u_{i,r'}\) and then evaluating the defining Hermitian-adjoint equality. Thus the stabilizer denominators, level volume, real Schur factor, and conjugation are all explicit.

This construction proves a class-mass formula for the actual automorphic operator. An independently given integral polynomial or Weyl-module basis need not be this orthonormal basis. Changing to it requires its actual Gram matrix and lattice determinant; those quantities are not removed by (DO26.20). In particular an arbitrary choice of the source's \nolinkurl{rowAF} cannot be made isometric by assertion.

\subsection{5. Exact comparison with the recovered arithmetic pins}\label{exact-comparison-with-the-recovered-arithmetic-pins}

The definitions in \nolinkurl{LevChiEqOne.lean} give the following nonoptional equations. The symbols on their right sides are independently specified source data.

First, \nolinkurl{VacT} requires (DO26.4) at degree zero with the source \nolinkurl{xiS}, \nolinkurl{I.J.V.thMoved}, and determinant-root factor \(\mu(x)^{-1}\mu(y)\), uniformly in every finite move. AM-DO-022 proves that mathematical statement for the new mixed geometric family and the standard Weil representation. \nolinkurl{PinX.mv} additionally quantifies over every \nolinkurl{om} satisfying \nolinkurl{IsWeilDatumO}; identifying all such axiomatic data with this fixed standard realization requires its uniqueness statement. The old elliptic \nolinkurl{VacIn} is not the new seed merely because both are vacua.

Second, \nolinkurl{PinColT.cst} is, literally, \[
 c_{K,j}=v_{j,\alpha}\,w_{j,\alpha\gamma}\,c_{0,j}
       \mu_{j+1}(x_{j,\alpha})
       \mu_{j+1}(x_{j+1,\gamma})^{-1}
                 \quad\text{for every }\alpha,\gamma.
 \tag{DO26.21}
\] The volume is that of the source chart level, not a theta-line norm. If every factor is nonzero, independence of the chart is equivalent to all pairwise ratio equations \[
 \frac{v_{j,\alpha}w_{j,\alpha\gamma}}
      {v_{j,\alpha'}w_{j,\alpha'\gamma'}}
 \frac{\mu(x_{j,\alpha})}{\mu(x_{j,\alpha'})}
 \frac{\mu(x_{j+1,\gamma'})}{\mu(x_{j+1,\gamma})}=1.
 \tag{DO26.22}
\] This equivalence follows by dividing two instances of (DO26.21); conversely fix one chart to define \(c_{K,j}\) and use all ratios. It is a finite concrete compatibility test. The geometric period \(c_{0,j}=C^{2N_jN_{j+1}}\) is now known; the source chart factors and their arithmetic meanings cannot be replaced by freely chosen scalars just to force (DO26.22).

Third, \nolinkurl{ClsK.Core} requires for \textbf{every} polynomial \[
 v_{j,\alpha}\Theta_{j,\alpha\gamma}(h_a,h_{a'};P)
  =c_{K,j}\,\operatorname{abelian}_2
       \bigl(\operatorname{mvZ}(1,\operatorname{sc}_j I)
                            \operatorname{castP}(P)\bigr).
 \tag{DO26.23}
\] Our equation (DO26.6) supplies its analytic-to-geometric kernel comparison. The actual source functional is \(\rho_i(c)\rho_i'(c')\operatorname{jetL}
  (I.\mathrm{th}(\operatorname{rep}c,\operatorname{rep}c'))\), as defined at lines 73--77 of \nolinkurl{LevChiEqOne.lean}. Its frame scale, finite arithmetic datum, and these two character scalars must be matched to (DO26.5). The finite-level statement in section 2 and the real statement (DO26.7) supply corresponding \nolinkurl{Lev} and \nolinkurl{Real} identities for the newly constructed datum. They do not identify an independently fixed datum without this comparison.

Fourth, \nolinkurl{ClsK.Kval} says that the prescribed arithmetic row map \(K_j\), evaluated in its specified frames, is exactly the stabilizer-weighted geometric kernel sum, including the target factor \(\det(y')^{-s_D}\). Equation (DO26.20) proves the class integral formula for the actual operator. To deduce that particular \nolinkurl{Kval}, one must exhibit the same coefficient-lattice realization and its compact projection normalization. Local type occurrence from AM-DO-025 does not choose that realization.

The abstract \nolinkurl{ThetaX} ties are equally precise. With \(\operatorname{petY}_j=m_{Y,j}\operatorname{pet}_j\) and \(\operatorname{dnY}_j=m_{Y,j+1}\Theta_j\), they require \[
 \begin{aligned}
 m_{Y,j}&=\prod_v\operatorname{parIdx}(q_v,r,jt),\\
 \operatorname{petY}_j(\operatorname{rowAF}_j x,
                      \operatorname{rowAF}_j y)&=\langle x,y\rangle,\\
 \operatorname{petY}_j(\operatorname{dnY}_j
           \operatorname{rowAF}_{j+1}y,\operatorname{rowAF}_j x)
                &=\langle\operatorname{io}_j y,x\rangle.
 \end{aligned}
 \tag{DO26.24}
\] Equations (DO26.11) and (DO26.20) give the exact analytic adjoint and its actual matrix. They establish (DO26.24) only when the preassigned \nolinkurl{rowAF}, \nolinkurl{io}, metrics, and parahoric factors are identified with those matrices. Choosing a new abstract map equal to an adjoint would prove a different statement.

\subsection{6. Terminal mass and the literal two-step MainX equation}\label{terminal-mass-and-the-literal-two-step-mainx-equation}

The source terminal equation \nolinkurl{ClsT.tail} is \[
 v_T\sum_{\beta,a}\frac{1}{|\Gamma_{T,a}|}
    \ell_{\beta,a,G}\!\left(
      \operatorname{low}\bigl(\operatorname{rd}_{m,\beta}
                    \psi\circ\iota(h_a,1)\bigr)\right)
                 =c_P\,\operatorname{tail}(\psi,G).
 \tag{DO26.25}
\] It holds for every allowed arithmetic row \(\psi\) and every terminal coefficient vector \(G\), with the same \(c_P\). The actual finite theta-line overlap \(q_0^{-1/2}\), and the holomorphic overlap with period and metric factors in AM-DO-024, are completely determined. They are not (DO26.25): the latter contains a different quotient measure, an actual restriction embedding, coefficient projections, and a prescribed arithmetic terminal functional.

For fixed finite-dimensional row and terminal spaces, (DO26.25) is a finite matrix test. Choose their actual bases, let \(B\) be the matrix of the left side computed with its masses, and let \(T\) be the matrix of \nolinkurl{tail}. The requirement is \[
                              B=c_PT.
 \tag{DO26.26}
\] If \(T\ne0\), choose one nonzero entry \(T_{rs}\). Then there is such a scalar if and only if \(B_{ij}T_{rs}=B_{rs}T_{ij}\) for all entries. This condition also forces \(B_{ij}=0\) whenever \(T_{ij}=0\). The proof is division by the chosen entry, followed by bilinearity on the two bases. Thus one nonzero period fixes a candidate scalar but does not verify its uniformity. The source also requires \(c_P\ne0\) where its later nonvanishing arguments use it.

For the literal two-step case \(m=2\), \nolinkurl{MainX} uses the implication \[
 \begin{split}
 &\sum_i\alpha_i
   c_{V,0}^{(i)}c_{K,0}^{(i)}c_{K,1}^{(i)}c_P^{(i)}
   \operatorname{tail}^{(i)}
      (K_1^{(i)}K_0^{(i)}S_0^{(i)}\psi_i,F_i)=0\\
 &\quad\Longrightarrow\quad
 \sum_i\alpha_i c_{I,0}^{(i)}c_{I,1}^{(i)}
   \operatorname{per}_0^{(i)}
      (\psi_i,\operatorname{io}_0^{(i)}
                       \operatorname{ioT}_0^{(i)}F_i;y,v)=0
                         \quad\text{for all }y,v.
 \end{split}
 \tag{DO26.27}
\] The quantifier over all finite linear relations and common frames is essential. In the general even chain the factors are exactly \[
 c_{\rm Top}=\prod_i(c_{V,2i}c_{K,2i}c_{K,2i+1})c_P,
               \qquad c_{\rm Bot}=\prod_jc_{I,j}.
 \tag{DO26.28}
\] These are \nolinkurl{Tup.cTop}, \nolinkurl{cBot}, and \nolinkurl{MainX} in \nolinkurl{ValLastX.lean}. The new compact kernels prove substantive components that enter this system. They do not replace its Bessel/Jacobi two-step unfolding, its fixed row maps, or (DO26.25).

\subsection{7. The six actual analytic predicates and exact source scope}\label{the-six-actual-analytic-predicates-and-exact-source-scope}

The source's \nolinkurl{AdmS} is a conjunction of the following six predicates, not a name for the elementary compact kernel identities alone.

\begin{longtable}[]{@{}
  >{\raggedright\arraybackslash}p{(\columnwidth - 4\tabcolsep) * \real{0.3333}}
  >{\raggedright\arraybackslash}p{(\columnwidth - 4\tabcolsep) * \real{0.3333}}
  >{\raggedright\arraybackslash}p{(\columnwidth - 4\tabcolsep) * \real{0.3333}}@{}}
\toprule\noalign{}
\begin{minipage}[b]{\linewidth}\raggedright
Actual predicate
\end{minipage} & \begin{minipage}[b]{\linewidth}\raggedright
What the new construction proves
\end{minipage} & \begin{minipage}[b]{\linewidth}\raggedright
Obligation still present in the literal source
\end{minipage} \\
\midrule\noalign{}
\endhead
\bottomrule\noalign{}
\endlastfoot
\nolinkurl{RallisSiegelWeilFormula} & Actual compact kernels, exact prescribed local coefficient types, and separate concrete nonzero first lifts & For every \nolinkurl{Hyp3} parameter, a nonzero irreducible global chain in the specified types, with every upper vector a finite sum of prescribed theta lifts \\
\nolinkurl{LiftParameterConditions} & Actual compact restriction pairings and their conjugation convention & A nonzero terminal period against the specified noncompact discrete-series family for every allowed top irreducible representation \\
\nolinkurl{ThetaUnfoldingFreeData} & Exact compact Fubini/adjoint identities (DO26.11)--(DO26.12) & The source's Bessel/Jacobi period transfer with factorizable, real-finite downward Schwartz data and its universal nonzero-period implication \\
\nolinkurl{ThetaStablePeriodVanishing} & No cuspidality problem for the constructed anisotropic compact groups themselves & The stated lower-Witt-tower vanishing implication for the noncompact \nolinkurl{Y} chain \\
\nolinkurl{RealThetaLiftDiscreteSeries} & Exact compact real formula and the selected compact Schur correspondence & The specified noncompact discrete-series preservation and lower-signature vanishing statements \\
\nolinkurl{RealConstantNonvanishing} & Translation preserves the actual compact representation and its fixed local type & Stability of the source's particular noncompact \nolinkurl{IsDv} condition under all displayed finite and real actions \\
\end{longtable}

Consequently the correct progression is: an actual mixed theta seed; its exact all-finite moved family and direct pairing; its all-degree geometric class kernel; the selected real covariance; then the independently pinned arithmetic row and terminal equations. The first four now have explicit formulas. The remaining source comparisons are concrete algebraic and automorphic obligations, not evidence that no other map or construction can exist.

\subsubsection{Exact recovered-source locators}\label{exact-recovered-source-locators}

All paths below are relative to \nolinkurl{work/research/alpoge_source}.

\begin{longtable}[]{@{}
  >{\raggedright\arraybackslash}p{(\columnwidth - 2\tabcolsep) * \real{0.5000}}
  >{\raggedright\arraybackslash}p{(\columnwidth - 2\tabcolsep) * \real{0.5000}}@{}}
\toprule\noalign{}
\begin{minipage}[b]{\linewidth}\raggedright
Interface
\end{minipage} & \begin{minipage}[b]{\linewidth}\raggedright
Source locator
\end{minipage} \\
\midrule\noalign{}
\endhead
\bottomrule\noalign{}
\endlastfoot
Bilinear \nolinkurl{liftS}/\nolinkurl{liftT} & \nolinkurl{Development/AbelianVarietyTheory/Definitions/TwistedInnerProducts/WittSmallGroupEmbeddingMemPointwiseStabilizerIso.lean}, lines 157--167 \\
Finite-combination kernel and lift & \nolinkurl{Development/AbelianVarietyTheory/Definitions/LatticeFrames/JointSmooth.lean}, \nolinkurl{KerDat.pt}, \nolinkurl{lift}, lines 74--82 \\
Actual vacuum, character roots, chart constant & \nolinkurl{Development/AbelianVarietyTheory/Definitions/LatticeFrames/LevChiEqOne.lean}, \nolinkurl{abelian₂}, lines 73--77; \nolinkurl{VacT}, 87--93; \nolinkurl{PinX}, 95--101; \nolinkurl{PinColT}, 103--119 \\
Real, level, kernel and arithmetic row identities & \nolinkurl{Development/AbelianVarietyTheory/Definitions/LatticeFrames/OnFramesRdApply.lean}, \nolinkurl{ClsK.Real}, 93--102; \nolinkurl{Lev}, 103--107; \nolinkurl{Core}, 108--113; \nolinkurl{Kval}, 115--121 \\
Terminal class mass & Same file, \nolinkurl{ClsT.Is.top} and \nolinkurl{tail}, lines 143--150 \\
Exact determinant exponents & \nolinkurl{Formalism/LocalRingCardinality/Params.lean}, \nolinkurl{sExp}, \nolinkurl{sExp'}, lines 71--73; \nolinkurl{Formalism/LocalToGlobal/Columns/Record.lean}, \nolinkurl{sSharpC}, \nolinkurl{sDivC}, lines 1083--1087 \\
Conjugation, periods and row ties & \nolinkurl{Formalism/LocalRingCardinality/DataTheta.lean}, \nolinkurl{conjAF}, 645--649; \nolinkurl{Bper} and \nolinkurl{Btop}, 657--670; \nolinkurl{petY}, \nolinkurl{dnY}, \nolinkurl{TieMY}, \nolinkurl{TieRowIsom}, \nolinkurl{TieIo}, \nolinkurl{TieBm}, 730--752 \\
Two-step and terminal relation laws & \nolinkurl{Development/AbelianVarietyTheory/Definitions/LatticeFrames/ValLastX.lean}, \nolinkurl{ValX}, \nolinkurl{ValLastX}, \nolinkurl{KerX}, \nolinkurl{PerX}, \nolinkurl{LawsX}, \nolinkurl{LastX}, \nolinkurl{MainX}; \nolinkurl{Tup.cTop} and \nolinkurl{cBot} \\
Six analytic assumptions & \nolinkurl{Formalism/SharpRestriction/CuspSpanNormBounds/CayleySignature.lean}, lines 154--202 \\
Existing vacuum-to-polynomial comparison & \nolinkurl{Development/AbelianVarietyTheory/Proofs/DecoratedLatticeProofs/ThetaKernelSplitPlacesEqJetClassicalTheta.lean}, complete theorem \nolinkurl{thetaKernel_eq_jet_thC} \\
\end{longtable}

No Lean build or formal instantiation of a source record is claimed. The cited interfaces and theorem bodies were directly read. The new operator proofs above state exactly which representation, period, measure, and coefficient normalization they use.

\clearpage

\section{AM-K021 --- the cubic relative period and the exact CM-height boundary}\label{am-k021-the-cubic-relative-period-and-the-exact-cm-height-boundary}

\textbf{Status.} New proved reductions and scoped obstructions for question 3. The stipulated Alpöge theorem, its archive, and the earlier accepted DP statements retain their accepted status. This report does not establish the relative arithmetic window, a new RH theorem, or a new automorphic cost bound.

\textbf{Dependencies.} AM-K001--AM-K005 in kummer\_extension.md and AM-K007--AM-K011 in kummer\_cubic\_host.md supply the Artin identity, conductor table, completed functional equation, elliptic Hom realization, finite ray character, norm-one splitting twist, and explicit CM supply tower. They are used with those statements' existing proofs and scope. The new result strengthens the earlier observation that coherent splitting phases telescope: even a genuine cubic twist of the actual polarized CM variety leaves the stable geometric letter degrees unchanged. A separate finite-character argument excludes the cubic standard Artin character from every finite linear combination of CM-type height classes.

\subsection{1. The exact source scalar is geometric and quadratic}\label{the-exact-source-scalar-is-geometric-and-quadratic}

All paths in this section are relative to work/research/alpoge\_source. These are inspected source interfaces, not a claim to have compiled or independently reconstructed their full upstream proofs.

\begin{longtable}[]{@{}
  >{\raggedright\arraybackslash}p{(\columnwidth - 2\tabcolsep) * \real{0.5000}}
  >{\raggedright\arraybackslash}p{(\columnwidth - 2\tabcolsep) * \real{0.5000}}@{}}
\toprule\noalign{}
\begin{minipage}[b]{\linewidth}\raggedright
Source locator
\end{minipage} & \begin{minipage}[b]{\linewidth}\raggedright
Exact role
\end{minipage} \\
\midrule\noalign{}
\endhead
\bottomrule\noalign{}
\endlastfoot
\nolinkurl{Formalism/LocalRingCardinality/DataParams₂.lean,} lines 65--71 & chiK is the quadratic Dirichlet character modulo D.host.d; dlogL is the real part of its logarithmic derivative. \\
The same file, lines 126--127 & SizeData₂.logDerivativeAtOne is constrained by logDerivative\_at\_one\_eq to equal dlogL D 1. It is not a freely chosen real parameter. \\
The same file, line 224 & logOm is Euler's constant plus log(π / host.v) minus that logarithmic derivative. \\
The same file, lines 232--234 & MAIN contains the term −R.m · aT(t) · logOm for every host type t. \\
\nolinkurl{Formalism/LocalToGlobal/ProjectiveLattices/RosatiLatticeNorms/LamObj.lean,} definitions IsVecLine and LamObjGood & A letter line must carry the actual integral tangent lattice and the norm induced by the polarized complex uniformization. \\
\nolinkurl{Formalism/LocalToGlobal/ProjectiveLattices/RosatiLatticeNorms/LamPtClauses.lean,} lines 53--56 & LamDegSt requires every actual letter to have normalized degree logOm / 2. \\
\nolinkurl{Formalism/LocalToGlobal/ProjectiveLattices/RosatiLatticeNorms/CmFrameOfEq.lean,} lines 45--58 and 78--115 & The degree identity is transported to the frame letters and is a hypothesis of the construction. \\
\nolinkurl{Formalism/LocalToGlobal/CMDeterminantIdeals/StatementS.lean,} lines 63--73 & Step5 consumes those exact letter degrees in its determinant estimate. \\
\end{longtable}

In particular the actual source chain has two separate equalities: \[
 \log\Omega=
 \gamma+\log(\pi/v)
 -\frac{L'(1,\chi_K)}{L(1,\chi_K)},
 \qquad
 \widehat{\deg}_{\rm norm}(\ell)=\frac12\log\Omega.
 \tag{R1}
\] The first equality is a prescribed auxiliary quadratic analytic constant. The second is a fact about actual integral metrized lines. Renaming a splitting character does not change either equality.

The source IsVecLine is stronger than a choice of norm on an abstract one-dimensional vector space. At every finite place its fractional ideal is characterized by membership in the tangent lattice of a model. At every complex embedding its squared norm equals π times the polarized Hermitian length of the corresponding tangent vector. Thus an arbitrary metric rescaling chosen to force a desired new value would require a new comparison with those actual geometric conditions.

\textbf{Scope.} The original Dirichlet argument can use auxiliary quadratic hosts while proving conclusions for other Dirichlet characters. Consequently (R1) does not imply that every successful relative argument must replace the quadratic logarithmic derivative by the target cubic logarithmic derivative. It proves that this replacement is not an available instantiation of the unchanged record. The task could instead require a new character-weighted arithmetic comparison using the existing auxiliary heights.

\subsection{2. Cubic descent does not create a new stable CM letter degree}\label{cubic-descent-does-not-create-a-new-stable-cm-letter-degree}

Set \[
 K=\mathbf Q(\zeta),\quad \zeta=\zeta_3,\quad
 \alpha=p^{1/3}>0,\quad M=K(\alpha),
\] and retain the standing condition p=3 or p≡2 modulo 3. Consider the actual elliptic curves \[
 A_0:\ y^2=x^3+1,\qquad A_p:\ y^2=x^3+p^2.
\] Their geometric CM action is (x,y)↦(ζx,y).

\subsubsection{Proposition R2. Stable isometry under the cubic twist}\label{proposition-r2.-stable-isometry-under-the-cubic-twist}

After a finite extension where both curves have good reduction, the polarized CM isomorphism \[
 g_p:A_p\longrightarrow A_0,\qquad
 (x,y)\longmapsto (x/\alpha^2,y/p)
 \tag{R2}
\] identifies their integral tangent eigenlines, their polarized archimedean metrics, and their normalized Arakelov degrees. The same holds for each eigenline in the polarized CM products constructed in AM-K010--AM-K011 with the same fixed positive Hermitian form.

\textbf{Proof.} The displayed substitution gives \[
 (y/p)^2=(x/\alpha^2)^3+1
\] because α³=p.~It fixes the origin, commutes with the CM action, and hence preserves the canonical principal elliptic polarization. After a further finite extension both curves have abelian-scheme models. The generic-fibre isomorphism and its inverse extend uniquely between these models, so their maps on integral Lie algebras are inverse integral isomorphisms. At a complex embedding the polarized analytic uniformizations are isomorphic, and the induced Hermitian metric is carried by the same tangent map. The resulting integral metrized lines are isometric.

Normalized Arakelov degree is unchanged by a finite base extension: both finite and infinite contributions multiply by the extension degree before normalization. It therefore gives the same degree for the two stable eigenlines over their original field of definition.

For the products in AM-K010, apply g\_p on each elliptic factor. The scalar map commutes with the given matrix action of the larger CM coefficient field and with the fixed Hermitian form. It carries the corresponding polarization and each coefficient-field eigenline to the untwisted one. The preceding argument applies to every such line, not just to the determinant of the whole tangent space. ∎

The coordinate scaling is useful for detecting an incomplete accounting. For the algebraic differential ω=dx/y, \[
 g_p^*\omega_0=\alpha\,\omega_p,
 \qquad
 \int_{\gamma_p}\omega_p
  =\alpha^{-1}\int_{g_p\gamma_p}\omega_0.
 \tag{R3}
\] Thus the displayed complex period changes by p\^{}\{-1/3\}. This is a change of algebraic basis in an isometric stable line. Its finite principal-divisor contribution cancels its archimedean contribution by the product formula. Keeping only the −(log p)/3 in the period would not compute a normalized Arakelov degree.

This assertion concerns stable lines after good reduction and their actual polarized metrics. It does not identify unstable minimal models, remove conductor terms, or set the finite level costs at 2, 3 and p to zero. Those costs remain part of any relative arithmetic adapter.

\textbf{Consequence.} In the particular induced CM supply tower already constructed, a norm-one Kummer descent twist can preserve the nontrivial Hom character while leaving every degree used in (R1) unchanged. The nontrivial Galois character is real arithmetic information, but it is not a new stable CM-type height.

\subsection{3. A representation-level obstruction for all CM-type height classes}\label{a-representation-level-obstruction-for-all-cm-type-height-classes}

Write G=Gal(M/Q). With σ(α)=ζα and complex conjugation c fixing α, \[
 G\simeq S_3,\qquad c\sigma c^{-1}=\sigma^{-1}.
 \tag{R4}
\] Let ρ be the two-dimensional augmentation representation on the three roots of X³−p.~Its character is \[
\begin{array}{c|ccc}
 \text{cycle type}&1&(12)&(123)\\ \hline
 \chi_\rho&2&0&-1.
\end{array}
\tag{R5}
\] AM-K001 gives \[
 L_K(s,\chi)=L_{\mathbf Q}(s,\rho)
            =\zeta_{\mathbf Q(\alpha)}(s)/\zeta_{\mathbf Q}(s).
 \tag{R6}
\]

Let Q\^{}\{cm\} be the composite of all CM subfields of a fixed algebraic closure. Complex conjugation is central on every finite Galois CM subfield. Such subfields are closed under finite composition and Galois closure. Equivalently Q\^{}\{cm\} is fixed by the closed subgroup generated by all commutators with complex conjugation (Milne, 2020a, I, Proposition 1.4, Corollary 1.5 and Remark 1.6, p.~10).

\subsubsection{Proposition R7. No cubic component in a finite CM-type height class}\label{proposition-r7.-no-cubic-component-in-a-finite-cm-type-height-class}

Let B/Q be a finite Galois CM field, N=BM, and H=Gal(N/B). Inflate ρ to Gal(N/Q). For every complex class function A inflated from Gal(B/Q), \[
 \langle A,\chi_\rho\rangle_{\operatorname{Gal}(N/\mathbf Q)}=0.
 \tag{R7}
\] Consequently no finite complex linear combination of class functions factoring through Q\^{}\{cm\} equals χρ. In particular no finite combination of CM-type height Artin classes gives the target cubic Artin character.

\textbf{Proof.} B∩M is Galois over Q because B and M are Galois. The normal subfields of the S₃ extension M are Q, K and M. The last possibility would imply that complex conjugation is central on Gal(M/Q), because it is central on Gal(B/Q). This contradicts (R4). Hence B∩M is Q or K, and the restriction of H to M is respectively S₃ or its normal subgroup A₃.

The restriction of ρ to A₃ is the sum of its two nontrivial one-dimensional characters. It has no invariant vectors. Thus ρ\^{}H=0. Averaging ρ(h) over h∈H is the zero projector. For each g∈Gal(N/Q), \[
 \sum_{h\in H}\chi_\rho(gh)
 =\operatorname{Tr}\!\left(\rho(g)\sum_{h\in H}\rho(h)\right)=0.
\] A is constant on each H-coset. Summing the character inner product by those cosets proves (R7).

Any finite collection of locally constant functions on Gal(Q\^{}\{cm\}/Q) factors through one finite Galois CM field after enlarging their fields of definition. Their linear span is therefore orthogonal to χρ. Since the latter is an irreducible character with inner product one with itself, it cannot belong to this span. ∎

The standard CM-type class is AΦ\^{}0, the conjugacy average of the convolution Φ*Φ∨ on Gal(Q\^{}\{cm\}/Q). Its construction is exactly on this quotient (Yang, 2010, §3, equations (3.1)--(3.12); author manuscript pp.~9--10). Proposition R7 concerns that definition and its finite Artin support. It does not require a new case of the Colmez formula.

Thus the obstruction remains even if every relevant CM-type height realization is granted. A finite identity of Artin representations, or an equality of all corresponding Artin Euler factors, cannot turn those CM-type height classes into ρ. Indeed equality of Euler factors at all unramified primes gives equality of characters on a common finite Galois group by Frobenius density, contradicting (R7).

An equivalent short detection is \[
 [c,\sigma]=c\sigma c^{-1}\sigma^{-1}=\sigma^{-2}=\sigma,
\] which acts nontrivially under ρ. Every representation factoring through Q\^{}\{cm\} kills this commutator.

\subsubsection{Limits of the obstruction}\label{limits-of-the-obstruction}

\begin{enumerate}
\def\labelenumi{\arabic{enumi}.}
\tightlist
\item
  It is an obstruction to an Artin-character identity, not a proof that an individual real number L′(1,ρ)/L(1,ρ) cannot accidentally equal some CM height expression.
\item
  It does not exclude infinite approximation schemes or an arithmetic construction with non-CM descent data.
\item
  It does not exclude a relative comparison in which CM heights remain auxiliary and the cubic character is represented elsewhere.
\item
  It does not say that the cubic Artin line is absent from all tensor categories generated by CM abelian varieties. AM-K004 explicitly realizes it as a Hom line between two CM elliptic eigensummands. The failure is in the CM-type class function used for the stable height, which forgets this finite descent twist.
\end{enumerate}

For odd prime r the same argument applies to the augmentation representation of C\_r⋊(Z/rZ)\^{}×. Complex conjugation acts by inversion on C\_r, and {[}c,σ{]}=σ\^{}\{-2\} generates C\_r. The augmentation representation has no C\_r invariants. The cubic case has the additional explicit ray-character and rank-one lattice constructions used below.

\subsection{4. A concrete relative theta period for the target logarithmic derivative}\label{a-concrete-relative-theta-period-for-the-target-logarithmic-derivative}

This section gives an actual function and an exact integral formula. It is independent of any proposed CM-height identification.

The primitive conductor from AM-K003 has the following scalar ideal generator m: \[
\begin{array}{c|c|c}
 \text{case}&\mathfrak f=m\mathcal O_K&Q=3N\mathfrak f=|D_{\mathbf Q(\alpha)}|\\ \hline
 p\ne3,\ p\equiv8\pmod9&m=p&3p^2\\
 p\ne3,\ p\equiv2,5\pmod9&m=3p&27p^2\\
 p=3&m=9&243.
\end{array}
\tag{R8}
\] Let Λ=O\_K=Z{[}ζ{]}⊂C in the fixed complex embedding. Its Euclidean covolume is V=√3/2, and \textbar a+bζ\textbar²=a²−ab+b².

Define the periodic function \[
 w(a)=
 \begin{cases}
 \chi((a)),&(a,\mathfrak f)=1,\\
 0,&(a,\mathfrak f)\ne1.
 \end{cases}
 \qquad a\in\Lambda.
 \tag{R9}
\] Here w(0)=0. Because K has class number one, the actual character is already given explicitly by AM-K007 on the finite unit quotient. The definition is independent of an ideal generator, invariant under the six units, and periodic modulo mΛ. As a nontrivial multiplicative character on (O\_K/𝔣)\^{}×, extended by zero, it has mean zero.

Set \[
 \Theta_\chi(t)=\sum_{a\in\Lambda}
                  w(a)e^{-\pi t|a|^2},
 \qquad t>0,
 \tag{R10}
\] and \[
 I_\chi(s)=\int_0^\infty\Theta_\chi(t)t^{s-1}\,dt.
 \tag{R11}
\]

\subsubsection{Proposition R12. Convergence and exact Mellin identity}\label{proposition-r12.-convergence-and-exact-mellin-identity}

The integral in (R11), and the integrals obtained by inserting any fixed power of log t, converge absolutely locally uniformly in s∈C. Therefore Iχ is entire. Moreover \[
 I_\chi(s)=6\pi^{-s}\Gamma(s)L_K(s,\chi)
 \tag{R12}
\] by analytic continuation. The function Θχ is real-valued.

\textbf{Proof.} For t bounded below away from zero, (R10) is absolutely and locally uniformly convergent by the positive-definite norm on Λ. Because w(0)=0, it decays exponentially as t→∞.

For the other endpoint, let Λ* be the Euclidean dual lattice, with pairing \(\langle z,a\rangle=\operatorname{Re}(z\bar a)\). The finite characters of Λ/mΛ are indexed by \[
 \xi\in m^{-1}\Lambda^*/\Lambda^*.
\] With \[
 \widehat w(\xi)=m^{-2}
  \sum_{b\in\Lambda/m\Lambda}
      w(b)e^{-2\pi i\langle\xi,b\rangle},
\] finite Fourier inversion gives \[
 w(a)=\sum_\xi\widehat w(\xi)e^{2\pi i\langle\xi,a\rangle}.
\] Its constant coefficient is zero. Poisson summation for the two-dimensional Gaussian yields \[
 \Theta_\chi(t)=\frac1{Vt}
   \sum_{\xi\ne0}\widehat w(\xi)
      \sum_{\lambda^*\in\Lambda^*}
       \exp\!\left(-\frac{\pi|\lambda^*-\xi|^2}{t}\right).
 \tag{R13}
\] Each nonzero coset has positive distance from Λ*. There are finitely many cosets. Hence Θχ(t)=O(t\^{}\{-1\}e\^{}\{-c/t\}) for some explicit c\textgreater0. These two endpoint bounds remain integrable after multiplication by t\^{}\{Re(s)−1\}\textbar log t\textbar\^{}j, uniformly on every compact set of s. Differentiation under the integral and entireness follow.

For Re(s)\textgreater1, absolute convergence permits termwise integration. The Gaussian integral is Γ(s)(π\textbar a\textbar²)\^{}\{-s\}. Each nonzero integral ideal of O\_K has exactly six generators, all of the same norm and character value. Thus the sum is six times the Hecke ideal series and gives (R12) in this half-plane, then everywhere by continuation.

Finally \(\chi(\bar{\mathfrak a})=\overline{\chi(\mathfrak a)}\) by AM-K002, while norms and the conductor are invariant under conjugation. The substitution \(a\mapsto\bar a\) conjugates the sum without changing it, so Θχ(t) is real. No pointwise positivity is asserted. ∎

\subsubsection{Explicit tail bounds}\label{explicit-tail-bounds}

The preceding proof is effective. Since \textbar ŵ(ξ)\textbar≤1 and Λ*=(2i/√3)Λ, define \[
 B(a)=\left(1+\frac{2e^{-a}}{1-e^{-3a}}\right)^2-1
 \qquad(a>0).
\] The elementary inequalities n²≥1+3(n−1) for n≥1 and a²−ab+b²≥(a²+b²)/2 give \[
 |\Theta_\chi(t)|\le B(\pi t/2),
 \qquad
 |\Theta_\chi(t)|\le\frac1{Vt}
                    B\!\left(\frac{2\pi}{3m^2t}\right).
 \tag{R14}
\] For the second bound, the disjoint nonzero dual cosets in (R13) form a subset of \(m^{-1}\Lambda^*\setminus\{0\}\); summing the absolute values over that larger lattice proves the formula. These inequalities give certified truncation of the integral and its logarithmic moments. They are not a useful arithmetic-window estimate by themselves.

\subsubsection{Corollary R15. An exact replacement scalar as a logarithmic heat moment}\label{corollary-r15.-an-exact-replacement-scalar-as-a-logarithmic-heat-moment}

The denominator Iχ(1) is positive. One has \[
 \frac{L'_K(1,\chi)}{L_K(1,\chi)}
 =\log\pi+\gamma+
   \frac{\displaystyle\int_0^\infty
                       \Theta_\chi(t)\log t\,dt}
        {\displaystyle\int_0^\infty
                       \Theta_\chi(t)\,dt}.
 \tag{R15}
\] Consequently, for any v\textgreater0, the scalar having the same displayed analytic form as the source scalar is \[
\begin{aligned}
 \log\Omega_\chi(v)
 &:=\gamma+\log(\pi/v)
        -\frac{L'_K(1,\chi)}{L_K(1,\chi)}\\
 &=-\log v-
   \frac{\displaystyle\int_0^\infty
                       \Theta_\chi(t)\log t\,dt}
        {\displaystyle\int_0^\infty
                       \Theta_\chi(t)\,dt}.
\end{aligned}
\tag{R16}
\]

\textbf{Proof.} The quotient in (R6) and the positive residue of the Dedekind zeta function at one give L\_K(1,χ)\textgreater0. Equation (R12) gives Iχ(1)=6L\_K(1,χ)/π\textgreater0. Its logarithmic derivative at one, using Γ′(1)/Γ(1)=−γ, is \[
 I_\chi'(1)/I_\chi(1)
  =-\log\pi-\gamma+L_K'(1,\chi)/L_K(1,\chi).
\] Rearrangement proves both assertions. ∎

This is a finite ray-character weighted lattice theta function with a completely specified Gaussian and measure. It changes actual summands according to χ; it is not the result of conjugating an unchanged oscillator chain by scalar phases. The Fourier version (R13) also makes it a finite linear combination of shifted dual-lattice heat sums.

The construction supplies a precise relative period to compare against an arithmetic model. It does not supply that comparison. In particular this report does not identify (R16) with any line satisfying the unchanged IsVecLine and LamDegSt conditions in §1. An integral or spectral construction with this logarithmic moment would need its own finite-place lattice, determinant formula, and estimate.

\subsection{5. Why a leading regulator alone does not supply the source scalar}\label{why-a-leading-regulator-alone-does-not-supply-the-source-scalar}

Let F=Q(α) and Q=\textbar D\_F\textbar. AM-K003 supplies the exact completion \[
 \Lambda(s)=Q^{s/2}\Gamma_{\mathbf C}(s)L_K(s,\chi),
 \qquad
 \Gamma_{\mathbf C}(s)=2(2\pi)^{-s}\Gamma(s),
 \qquad
 \Lambda(s)=\Lambda(1-s).
 \tag{R17}
\] The value Λ(1) is nonzero and finite. Therefore L\_K(s,χ) has a simple zero at zero. Write \[
 L_K(s,\chi)=a s+b s^2+O(s^3),\qquad a\ne0.
 \tag{R18}
\]

\subsubsection{Proposition R19. The required invariant is a next-coefficient ratio}\label{proposition-r19.-the-required-invariant-is-a-next-coefficient-ratio}

The exact relation is \[
 \frac{L'_K(1,\chi)}{L_K(1,\chi)}
 =-\log Q+2\log(2\pi)+2\gamma
       -\frac{L_K''(0,\chi)}{2L_K'(0,\chi)}.
 \tag{R19}
\] If h\_F and Reg(F) denote the ordinary class number and regulator, then \[
 L_K(1,\chi)=\frac{2\pi h_F\operatorname{Reg}(F)}{\sqrt Q},
 \qquad
 L_K'(0,\chi)=h_F\operatorname{Reg}(F).
 \tag{R20}
\] Moreover, for the Laurent expansion \[
 \zeta_F(s)=\frac{R_F}{s-1}+C_F+O(s-1),
 \qquad \gamma_F=C_F/R_F,
\] one has \[
 L_K'(1,\chi)/L_K(1,\chi)=\gamma_F-\gamma.
 \tag{R21}
\]

\textbf{Proof.} The expansion Γ(s)=s\^{}\{-1\}−γ+O(s) gives \[
 \frac{\Lambda'(0)}{\Lambda(0)}
   =\tfrac12\log Q-\log(2\pi)-\gamma+b/a.
\] At one, \[
 \frac{\Lambda'(1)}{\Lambda(1)}
   =\tfrac12\log Q-\log(2\pi)-\gamma
      +\frac{L_K'(1,\chi)}{L_K(1,\chi)}.
\] The logarithmic derivative of (R17) makes these two expressions negatives. Since b/a=L\_K″(0,χ)/(2L\_K′(0,χ)), this proves (R19).

F has signature (1,1) and only two roots of unity. The analytic class number formula therefore gives R\_F=2πh\_F Reg(F)/√Q. This is L\_K(1,χ) by (R6). This specialization uses Milne (2020b), V, Theorem 2.4(a), p.~156. At zero, Λ(0)=2a. At one, Λ(1)=√Q L\_K(1,χ)/π. Equality gives the second assertion in (R20).

Finally divide the displayed Laurent expansion of ζ\_F by ζ\_Q(s)=(s−1)\^{}\{-1\}+γ+O(s−1). The constant term of the quotient is R\_F, and its linear coefficient is C\_F−γR\_F. Their quotient is (R21). ∎

Thus an explicit unit or regulator realization of the leading coefficient L\_K′(0,χ) is valuable, but it is not the invariant needed in (R1) with χ inserted. Formula (R19) requires the next coefficient relative to the leading one. Equivalently (R21) requires the Euler--Kronecker constant of the non-Galois cubic field, relative to that of Q.

For comparison with the theta construction, \[
 I_\chi(0)=6h_F\operatorname{Reg}(F),\qquad
 I_\chi(1)=12h_F\operatorname{Reg}(F)/\sqrt Q.
 \tag{R22}
\] Its logarithmic moment Iχ′(1)/Iχ(1), rather than either displayed value alone, is precisely the new scalar in (R16).

\subsection{6. The actual Artin motive and its ordinary period matrix}\label{the-actual-artin-motive-and-its-ordinary-period-matrix}

There is no need to posit a new variety merely to realize the representation.

Over K, let σ generate Gal(M/K). In the convention where Γσ acts by σ, the projector \[
 e_\chi=\frac13\sum_{j=0}^2\chi(\sigma)^{-j}\Gamma_{\sigma^j}
 \tag{R23}
\] on h⁰(Spec M), with coefficients containing μ₃, gives the rank-one χ component. Reversing the cohomological action convention exchanges χ and χ\^{}\{-1\}, whose L-functions coincide here.

Over Q, let π:Spec F→Spec Q. The degree-three unit and trace maps give the rational projector \[
 e_\rho=1-\frac13\pi^*\pi_*,
 \qquad
 h^0(\operatorname{Spec}F)=\mathbf1\oplus\mathcal M_\rho.
 \tag{R24}
\] Indeed π\_\emph{π\^{}}=3 on the trivial summand, so eρ²=eρ. Its Betti realization is the trace-zero part of the three-element permutation representation, hence is ρ. Its Artin L-function is exactly (R6).

For an explicit de Rham/Betti comparison, the trace-zero de Rham space is \[
 \ker(\operatorname{Tr}_{F/\mathbf Q}:F\to\mathbf Q)
   =\mathbf Q\alpha\oplus\mathbf Q\alpha^2.
\] Evaluating on the three embeddings gives the matrix \[
 \begin{pmatrix}
  \alpha&\alpha^2\\
  \zeta\alpha&\zeta^2\alpha^2\\
  \zeta^2\alpha&\zeta\alpha^2
 \end{pmatrix}.
 \tag{R25}
\] Its columns sum to zero, and the first two rows have determinant p(ζ²−ζ)≠0. This is a fully explicit algebraic comparison for the weight-zero Artin motive.

An ordinary Artin comparison matrix such as (R25) has algebraic entries. By itself it does not prescribe the logarithmic heat moment (R16), its integral metric, or its derivative-of-L determinant formula. The leading regulator statement (R20) supplies a different arithmetic object, involving units; (R19) explains exactly why that leading datum alone is still insufficient.

There is also the positive-dimensional realization from AM-K004: an eigensummand of Hom(h¹(A₀),h¹(A\_p)) over K carries χ or its inverse. Proposition R2 and (R25) explain how a nontrivial Artin descent can coexist with unchanged stable CM height. These are compatible facts, not contradictory notions of ``CM''.

\subsection{7. What this changes in the steps 2--20 audit}\label{what-this-changes-in-the-steps-220-audit}

The complete original step table remains in kummer\_extension.md, §7. The present proof resolves the following narrower questions.

\begin{longtable}[]{@{}
  >{\raggedright\arraybackslash}p{(\columnwidth - 4\tabcolsep) * \real{0.3333}}
  >{\raggedright\arraybackslash}p{(\columnwidth - 4\tabcolsep) * \real{0.3333}}
  >{\raggedright\arraybackslash}p{(\columnwidth - 4\tabcolsep) * \real{0.3333}}@{}}
\toprule\noalign{}
\begin{minipage}[b]{\linewidth}\raggedright
Part of the attempted transfer
\end{minipage} & \begin{minipage}[b]{\linewidth}\raggedright
New exact conclusion
\end{minipage} & \begin{minipage}[b]{\linewidth}\raggedright
Remaining obligation
\end{minipage} \\
\midrule\noalign{}
\endhead
\bottomrule\noalign{}
\endlastfoot
CM geometric supply and character descent & The explicit polarized cubic twists and induced CM tower exist; their actual stable tangent letters are isometric to the untwisted ones. & Compute descent-sensitive finite indices and the genuinely character-weighted arithmetic quantity. \\
The height term used in MAIN & The unchanged source record fixes a quadratic dlogL and actual degrees. The cubic S₃ character cannot be obtained by a finite CM-type height Artin-class identity. & A new comparison theorem, or an auxiliary-height argument that inserts the cubic character outside that scalar. \\
Ideal/ray-character analytic side & Θχ and its Mellin transform give the exact cubic L-function, including the wild p=3 case and primitive conductor. & Relate this actual relative period to the chosen arithmetic sections and lattice estimates. \\
Motivic realization & The rank-two Q-Artin projector and rank-one K-projector are explicit; a CM elliptic Hom realization is already available. & A period/metric/determinant construction representing the logarithmic moment, with all finite terms. \\
Leading special values & L(1), L′(0), and the ordinary unit regulator are exactly related. & Control L″(0)/(2L′(0)), or equivalently the relative Euler--Kronecker invariant, in the required arithmetic comparison. \\
Final analytic implication & AM-K005 already supplies the field-independent Landau/window adapter. & Prove its full family of arithmetic window bounds. None follows merely from the formulas above. \\
\end{longtable}

The relative Mellin construction is therefore a concrete alternative target, rather than an unspecified ``Hecke analogue of Chowla--Selberg''. Its input is the already explicit ray character, its output is exactly the required logarithmic derivative, and its normalization is fixed in (R10)--(R16). The open part is an arithmetic realization and estimate of this period that meets the actual source hypotheses.

\subsection{8. Sources and verification scope}\label{sources-and-verification-scope}

\begin{itemize}
\tightlist
\item
  Alpöge, L. (2026). Author Lean archive, source locations in §1. The local mirror is the parent-recovered source set. This report inspects the listed definitions and statements; it makes no new whole-archive compilation claim.
\item
  Milne, J. S. (2020a). \emph{Complex multiplication}. I, Proposition 1.4, Corollary 1.5 and Remark 1.6, p.~10. \url{https://www.jmilne.org/math/CourseNotes/CM.pdf}. Used only for the maximal CM quotient and central conjugation facts; the obstruction in Proposition R7 is proved here.
\item
  Milne, J. S. (2020b). \emph{Class field theory}. V, §2, especially Theorem 2.4(a)--(b), pp.~155--156. \url{https://www.jmilne.org/math/CourseNotes/CFT.pdf}. Used for the Hecke ideal convention, nonvanishing at one and the analytic class number formula.
\item
  Yang, T. (2010). The Chowla--Selberg formula and the Colmez conjecture. \emph{Canadian Journal of Mathematics, 62}(2), 456--472. \url{https://doi.org/10.4153/CJM-2010-028-x}. Author manuscript: \url{https://people.math.wisc.edu/~tonghaiyang/Colmez.pdf}. §3, equations (3.1)--(3.12), author pp.~9--10, defines the CM-type class functions on the maximal CM quotient. No unproved general Colmez formula is invoked in Proposition R7.
\item
  Internal proved prerequisites: AM-K001--AM-K005 and AM-K007--AM-K011, at the companion paths identified at the beginning. All new analytic identities in §§4--6 are proved above with explicit normalizations.
\end{itemize}

\textbf{Completed author checks.} The saved report was read in full and its arguments reconstructed: stable polarized tangent-line isometry; the common-Galois-group averaging in (R7); conductor and unit factors; both finite Fourier normalizations and Gaussian tail bounds; the Mellin logarithmic derivative; both gamma-factor expansions in (R19); the class-number and regulator constants; and the Artin projectors and trace-zero comparison matrix. The listed source interfaces and the cited primary definitions were inspected directly. No numerical or finite computational certificate was run for this report. The continuation in (R12) follows from its convergence proof, not from numerical agreement. The author check and source-use events are recorded in \nolinkurl{kummer_extension_CHECK_EVENTS.tsv.}

\clearpage

\section{AM-DO-027: exact common-chart phases and an integral unit normalization}\label{am-do-027-exact-common-chart-phases-and-an-integral-unit-normalization}

Author: GPT-6 Astra Pro, ALP\_MACHINE \nolinkurl{/root/definite_oscillator}; inherited reasoning effort not separately exposed; 21 November 2026 UTC.

This result closes the scalar compatibility in \nolinkurl{PinColT.cst} once the raw arithmetic finite datum has actually been realized. The normalization is permitted by the literal \nolinkurl{IsRatDatum} definition, not by changing the vacuum or inventing a new meaning for that predicate. On equal-volume charts with compact determinants it uses only roots of unity. It therefore has no additional finite denominator, content, or norm cost. An explicit example also shows why unramifiedness of the CM extension does not make these character phases disappear automatically.

\subsection{1. The precise common-character scalar equation}\label{the-precise-common-character-scalar-equation}

Fix one even-rank edge \(a\to b\) of the common \(E'=\mathbf Q(\zeta_{28})\) host, with the same actual splitting character \(\mu\) on both members, as required by \nolinkurl{PinColT.om}. Its restriction to \(\mathbb A_{F'}^\times\) is trivial. Write \[
 \nu_\mu(t)=\mu(x),\qquad x/\bar x=t.
 \tag{DO27.1}
\] Hilbert 90 and the displayed restriction make this a well-defined character of the norm-one torus. For source and target chart base elements \(k_\alpha,k'_\gamma\), define \[
 u_\alpha=\nu_{\mu,f}(\det k_\alpha)=\mu(x_\alpha),
 \qquad
 u'_\gamma=\nu_{\mu,f}(\det k'_\gamma)=\mu(x'_\gamma).
 \tag{DO27.2}
\] The roots are exactly the source's \nolinkurl{xr}, not the determinants themselves. Its \nolinkurl{TwDat.Is.kA} requires these root equations. Let \(v_\alpha\) be the actual source chart level volume. The literal condition is \[
        c_K=v_\alpha w_{\alpha\gamma}c_0
                         u_\alpha(u'_\gamma)^{-1}
                   \quad\text{for every }\alpha,\gamma,
 \qquad c_0=C^{2ab}.
 \tag{DO27.3}
\] This is \nolinkurl{LevChiEqOne.lean}, \nolinkurl{PinColT.cst}, lines113--114.

For unchanged raw scalars \(w_{\alpha\gamma}=1\), a common nonzero \(c_K\) exists if and only if

\begin{itemize}
\tightlist
\item
  all source volumes \(v_\alpha\) are equal;
\item
  \(u_\alpha\) is constant on the source chart set;
\item
  \(u'_\gamma\) is constant on the target chart set.
\end{itemize}

Here both chart sets are assumed nonempty. To prove necessity, fix a target chart and compare two source charts in (DO27.3). The absolute values give equality of their positive volumes, because \(|u_\alpha|=1\); the phases then give equality of the \(u_\alpha\). Fixing a source chart similarly proves the target assertion. Conversely those three conditions make the right side independent of the charts. This is a necessary and sufficient criterion, not merely a sufficient normalization rule.

Equivalently, the determinant conditions are \[
 \det(k_\alpha k_{\alpha_0}^{-1})\in\ker\nu_{\mu,f},
 \qquad
 \det(k'_\gamma(k'_{\gamma_0})^{-1})\in\ker\nu_{\mu,f}
                         \quad\text{for all }\alpha,\gamma.
 \tag{DO27.4}
\] If the determinant representatives are compact and the finite character is trivial on the maximal norm-one compact, these conditions hold automatically. Relative unramifiedness of \(E'/F'\) alone imposes no such triviality on the chosen character.

\subsection{2. Linearity is present in the actual source definitions}\label{linearity-is-present-in-the-actual-source-definitions}

The relevant definitions were read directly, rather than inferred from their names. In \nolinkurl{Formalism/LocalToGlobal/IndefiniteReal/LocalNewformFunction.lean}, lines238--261, the source has:

\begin{enumerate}
\def\labelenumi{\arabic{enumi}.}
\tightlist
\item
  \nolinkurl{IsSBHalf e Λ Φ}: support in the idempotent half, a common integral period, and a common denominator support condition.
\item
  \nolinkurl{IsPlaceModel … f₀ T}: \(T\) is explicitly a complex-linear map on functions, intertwines translations and modulations, and satisfies the unchanged vacuum normalization \(T(c\,\mathrm{halfInd})=f_0\).
\item
  \nolinkurl{IsThetaAt … f₀ Φ f}: \nolinkurl{AllSome} of those models with the assertion \(f=T\Phi\).
\end{enumerate}

\nolinkurl{IsRatDatum}, at lines25--29 of its namesake file, is the existence of such \(\Phi\), with its rational restriction equal to the specified arithmetic coefficient function \(\Xi\), and with that \nolinkurl{IsThetaAt} assertion. It follows that, for any complex scalar \(z\), \[
 \boxed{\operatorname{IsRatDatum}(\Xi,f_0,f)
       \Longrightarrow
       \operatorname{IsRatDatum}(z\Xi,f_0,zf).}
 \tag{DO27.5}
\] The level idempotent, lattice, vacuum, additive character, and \nolinkurl{IsPlaceModel} normalization are unchanged.

\textbf{Proof.} Use \(z\Phi\) as the new witness. Its support is contained in that of \(\Phi\), its same integral period works, and its same denominator works. Thus it remains \nolinkurl{IsSBHalf}. Its rational restriction is \(z\Xi\). Keep exactly the existing models \(T\); each is linear, so \(T(z\Phi)=zT\Phi=zf\). This proves the new assertion for every such model as well as for its existing witness, regardless of the \nolinkurl{AllSome} packaging. The condition \(T(c\,\mathrm{halfInd})=f_0\) is untouched. This proves (DO27.5). \(\square\)

For a nonzero \(z\), the support of the finite coefficient family is unchanged. \nolinkurl{JointSmooth.lean}, lines20--24, defines \[
 \operatorname{comb}(d,f_0)=\sum_xd_x\operatorname{heis}(x)f_0,
 \tag{DO27.6}
\] and \nolinkurl{IsFinComb} says that every support point has zero infinite part. Consequently \(\operatorname{comb}(zd,f_0)=z\operatorname{comb}(d,f_0)\), and \nolinkurl{IsFinComb} is retained. In particular this is a permitted rescaling of actual finite Schwartz inputs. Any fixed compact eigencharacter or real covariance of the input is also unchanged.

\subsection{3. Actual normalized finite inputs and the common constant}\label{actual-normalized-finite-inputs-and-the-common-constant}

Assume that for every chart pair there is an \textbf{actual raw datum} \[
 \operatorname{IsRatDatum}(\Xi_{\alpha\gamma},f_0,
                         \phi^0_{\alpha\gamma}),
 \qquad
 \phi^0_{\alpha\gamma}
            =\operatorname{comb}(d^0_{\alpha\gamma},f_0),
 \tag{DO27.7}
\] with the specified finite support condition. The arithmetic \(\Xi_{\alpha\gamma}\) in this assumption is the source's \nolinkurl{A.xi}; it has not been replaced by a freely chosen function. Its existence is an independent input to this proposition.

Fix any reference positive volume \(v_*\), and set \[
 w_{\alpha\gamma}
   =\frac{v_*}{v_\alpha}\,u_\alpha^{-1}u'_\gamma,
 \quad
 d_{\alpha\gamma}=w_{\alpha\gamma}d^0_{\alpha\gamma},
 \quad
 \phi_{\alpha\gamma}=w_{\alpha\gamma}
                                  \phi^0_{\alpha\gamma}.
 \tag{DO27.8}
\] Then (DO27.5)--(DO27.6) give exactly \nolinkurl{PinColT.dat}, with its required scalar multiple \(w_{\alpha\gamma}\Xi_{\alpha\gamma}\). Substitution in (DO27.3) gives, for every chart pair, \[
                          \boxed{c_K=v_*c_0.}
 \tag{DO27.9}
\] The common period is unchanged. No new choice of the source vacuum or of its unit norm is made. This proves the simultaneous \nolinkurl{dat} and \nolinkurl{cst} assertions from (DO27.7), not only a formal consistency of the numbers in \nolinkurl{cst}.

Suppose now that the source levels are conjugate to one fixed level inside the finite unitary group. A reductive local group is unimodular, so conjugation preserves Haar measure; their volumes are all the same \(v_*\). Suppose also that all the determinants in (DO27.2) belong to the maximal compact norm-one torus. The restriction of a continuous finite-adelic unitary character to that compact group has finite image. Indeed the compact is profinite, its character image is a compact profinite subgroup of the circle, and such a subgroup is finite. Therefore \[
                 w_{\alpha\gamma}=u_\alpha^{-1}u'_\gamma
                           \quad\text{is a root of unity.}
 \tag{DO27.10}
\] It follows exactly that \[
 \|\phi_{\alpha\gamma}\|_2
       =\|\phi^0_{\alpha\gamma}\|_2,
 \qquad
 v_{\mathfrak p}(w_{\alpha\gamma})=0
                              \quad\text{at every finite prime}.
 \tag{DO27.11}
\] Multiplication by this scalar is an invertible integral map on every coefficient lattice containing the character values. It preserves its content, all elementary divisors, and its saturation index. Every complex embedding also has absolute value one. These are equalities, not upper estimates.

There is no assertion that a later selected Schur projection has unit content. The statement concerns this particular scalar normalization of the existing chart inputs. Its root-of-unity change neither cures nor worsens a previously uncomputed lattice or projection denominator.

If the chart volumes differ, the exact remaining extra scalar is \(v_*/v_\alpha\). For commensurable compact opens it is rational: with \(C_\alpha=U_*\cap U_\alpha\), \[
 \frac{v_*}{v_\alpha}
        =\frac{[U_*:C_\alpha]}{[U_\alpha:C_\alpha]}.
 \tag{DO27.12}
\] The norm changes by that positive ratio and its local content changes by its valuation. Thus the character phase and the volume cost are separated by an exact formula. Neither is discarded when equal-volume hypotheses have not been verified.

On the present host compact-determinant representatives are available. R009 proves the norm-one finite torus has class number one. If \(g\) is any finite unitary representative, write \(\det g=\delta u\), with \(\delta\in T'(F')\) and \(u\) compact. Left multiplication by the rational unitary matrix \(\operatorname{diag}(\delta^{-1},1,\ldots,1)\) makes its determinant compact without changing its finite double class. This observation does not by itself prove that an independently pinned arithmetic chart has already made this representative choice; its frames must be transported along that rational action.

\subsection{4. An explicit nonconstant phase on unchanged lattices}\label{an-explicit-nonconstant-phase-on-unchanged-lattices}

The following example uses the same character on the two members of an even-rank edge and the same unramified CM host. It shows that the root-of-unity normalization in section 3 can be necessary.

Let \(a=6,b=20\). On the six real torus coordinates indexed by \(\Phi'=\{1,3,9,15,17,23\}\), prescribe the character \[
                         \beta(t)=\prod_{\tau\in\Phi'}t_\tau^{-8}.
 \tag{DO27.13}
\] At one split place \(v_*\) above 29 choose the first component so that \(\zeta_{28}\mapsto2\bmod29\), and define its residue character by \[
                         \rho(2)=\zeta_{28}^{12}.
 \tag{DO27.14}
\] Take the other finite compact components trivial. Since \(\sum\Phi'=68\), the value of (DO27.13) on the rational \(\zeta_{28}\) is \(\zeta_{28}^{-544}=\zeta_{28}^{16}\). Thus (DO27.14) is exactly its inverse on \(\mu_{28}\). The torus class-number-one decomposition from R009 constructs an actual unitary automorphic character \(\gamma\) with these infinity and finite compact restrictions. Put \[
                  \mu=\gamma\circ(z\mapsto z/\bar z).
 \tag{DO27.15}
\] Then \(\mu|_{\mathbb A_{F'}^\times}=1\), and \(\mu_\tau(z)=(z/\bar z)^{-8}\) at every selected place. This is an admissible same-character splitting for the two even-rank members. It also has an actual source \nolinkurl{HCParam} instance with admissible numerical \nolinkurl{Ranks} data: \(k=6,t=13,m=2,\mathrm{ord}=3\), \(u=(11/2,9/2,7/2,5/2,3/2,1/2)\), \(v=(5,4,3,2,1),b=(0,16),e=(8,24)\). The required interlacings hold, the final \(u\)-entry is positive, and both successive gaps equal \((t+3)/2=8\). The numerical rank conditions are also immediate: \(m+3=5\le6\), \(k+m-1+5=12\le13\), \(3m=6<13\), \(m\le2\mathrm{ord}=6\le k\), and \(m^2=4\le34\). No claim about further \nolinkurl{Hyp3} hypotheses is needed for this phase example.

Let the source charts use the common principal congruence level modulo 29 at all six places above 29 and the standard maximal compact elsewhere. Take two source base elements \(k_0=1\) and \(k_1\), where at \(v_*\) \[
                 k_1=\operatorname{diag}((2,2^{-1}),1,1,1,1,1)
                        \in U(6)(F'_{v_*}),
 \tag{DO27.16}
\] under the relative split identification, and at all other places \(k_1=1\). Both lie in the standard integral compact and hence preserve the identical selfdual lattice. Their source levels are equal, so their volumes are equal. Choose the target identity base chart. The determinant roots give \[
                 u_0=1,\qquad u_1=\zeta_{28}^{12}
                         =\zeta_7^{3},\qquad u'_0=1.
 \tag{DO27.17}
\] The two source level classes are distinct. If \(k_1=qu\), with \(q\in U(6)(F')\) and \(u\) in the stated principal level, then \(q\) is integral at every finite place. Its determinant is therefore a norm-one unit, hence a member of \(\mu_{28}\). At any other place over 29 that determinant reduces to 1. Reduction there is injective on \(\mu_{28}\), so \(\det q=1\). At \(v_*\) the congruence would instead give \(\det q=2\bmod29\), a contradiction.

For unscaled raw inputs, (DO27.3) asks simultaneously for \[
                         c_K=vc_0,
             \qquad      c_K=\zeta_7^{3}vc_0,
 \tag{DO27.18}
\] which is impossible since \(v>0,c_0\ne0\). This failure is a real character phase on actual different level classes, despite unchanged lattices. The underlying standard trivial-character vacuum is fixed, while the prescribed common-character source Weil action is \[
                    \omega_\mu(k_1,1)\phi_0
                           =\zeta_7^{-3}\phi_0.
 \tag{DO27.19}
\] Equation (DO27.19) is the exact inverse source determinant-root factor of AM-DO-022. The theta reading is injective and its entire vacuum germ is nonzero, so this is an actual kernel phase, not just a formal determinant computation. A sufficiently high polynomial jet detects it even if the origin value of the section vanishes.

The normalization of section 3 sets \(w_{00}=1\) and \(w_{10}=\zeta_7^{-3}\). It solves (DO27.18) in the literal \nolinkurl{PinColT.cst} convention, retaining the actual determinant-root factors in the kernel. Its scalar is already in \(E'\), is a unit at every finite prime, and has absolute value one at every complex embedding. The example is therefore a counterexample to automatic phase constancy, together with an exact permitted repair. It is not a counterexample to existence of \nolinkurl{PinColT}.

\subsection{5. Uniform character phases for the selected exponents}\label{uniform-character-phases-for-the-selected-exponents}

The same one-place construction works for arbitrary integer selected exponents \(-e_{j,\tau}\). Give the real torus that infinity type, and at the chosen place over 29 put \[
 \rho_j(2)=\zeta_{28}^{r_j},\qquad
            r_j\equiv\sum_{\tau\in\Phi'}\tau e_{j,\tau}
                                               \pmod{28}.
 \tag{DO27.20}
\] The inverse root-unit condition follows by evaluating at \(\zeta_{28}\), exactly as in section 4. Extend it using the same torus class decomposition and compose with \(z/\bar z\). This produces actual admissible equal-character edges with the selected infinity exponents. For compact-determinant chart representatives every scalar in (DO27.8) then lies in \(\mu_{28}\subset E'\), uniformly in the degrees and ranks. There is no new coefficient-field extension for these phases.

If a particular source instance has already fixed different characters \(\mu_j\), the construction in (DO27.20) does not replace them. Sections 1--3 apply to those characters as they stand; their actual compact values determine the unit scalars. Likewise, normalizing each edge does not establish the separate row-lattice, Schur-projection, or terminal-mass identities. It does establish that, once the raw finite datum is available, the common-chart character phases in \nolinkurl{PinColT.cst} can be removed within the expressly permitted scalar parameter, with the exact volume cost (DO27.12) and no additional unit cost.

\subsection{Exact source and carried-result locators}\label{exact-source-and-carried-result-locators}

\begin{itemize}
\tightlist
\item
  \nolinkurl{Formalism/LocalToGlobal/IndefiniteReal/LocalNewformFunction.lean}, \nolinkurl{IsSBHalf}, lines238--241; \nolinkurl{IsPlaceModel}, lines248--255; \nolinkurl{IsThetaAt}, lines257--261; \nolinkurl{splitFacO}, line51.
\item
  \nolinkurl{Development/AbelianVarietyTheory/Definitions/LatticeFrames/IsRatDatum.lean}, definition \nolinkurl{IsRatDatum}, lines25--29.
\item
  \nolinkurl{Development/AbelianVarietyTheory/Definitions/LatticeFrames/JointSmooth.lean}, \nolinkurl{comb} and \nolinkurl{IsFinComb}, lines20--24.
\item
  \nolinkurl{Development/AbelianVarietyTheory/Definitions/LatticeFrames/LevChiEqOne.lean}, \nolinkurl{TwDat.Is.kA}, lines52--53; \nolinkurl{PinColT.dat} and \nolinkurl{cst}, lines111--114.
\item
  AM-DO-022, exact standard finite splitting and determinant roots; AM-DO-025, selected infinity exponents; AM-DO-026, actual kernel normalization and chart equation; corrected AM-R-009, the common-host norm-one torus class decomposition and character extension. The published class-number input is carried through R009, not re-proved here.
\end{itemize}

No Lean compilation or complete \nolinkurl{PinColT} instantiation is claimed. The raw arithmetic datum (DO27.7) and the remaining fields of that record are explicitly retained as separate requirements.

\clearpage

\section{AM-R-011: the selected Schur norms, their row normalization, and a finite projection denominator}\label{am-r-011-the-selected-schur-norms-their-row-normalization-and-a-finite-projection-denominator}

Author: /root, GPT-6 Astra Pro, ALP\_MACHINE, 21 November 2026 UTC.

This result inserts the \textbf{actual selected partitions} of AM-DO-025 into the Fischer metric and the coefficient normalization used by the recovered source. It computes a concrete part of the archimedean cost, and a valid integral denominator for projection onto a fixed finite polynomial component. It does not identify that denominator with the least denominator of the source row map, or compute a full DP44 gain.

\subsection{1. The two partitions that must be distinguished}\label{the-two-partitions-that-must-be-distinguished}

At one selected CM embedding use the notation of AM-DO-025:

\[
 a=N_j,\quad b=N_{j+1}=a+d,\quad
 w_j=\operatorname{lamTrue}_j,\quad
 \nu_j=w_j+e_j-\frac b2\mathbf1_a.
 \tag{R11.1}
\]

The source coefficient representation uses the partition obtained by removing the final weight:

\[
 \eta_j=w_j-w_{j,a-1}\mathbf1_a,
 \qquad q_j=\nu_{j,a-1}.
 \tag{R11.2}
\]

Thus \(\eta_j\) is literally \nolinkurl{DomWeight.natW}, and \(w_{j,a-1}\) is \nolinkurl{detShift}. AM-DO-025 proves \(q_j\ge1\) and \(\nu_j\) is an integral partition. Consequently

\[
 \boxed{\nu_j=\eta_j+q_j\mathbf1_a,
 \qquad \eta_{j+1}=(\nu_j,0^d).}
 \tag{R11.3}
\]

For the second identity, the last block of the next weight has value \(-e_j+a/2\). Every old entry exceeds that value by \(\nu_{j,i}>0\), by the exact recurrence in AM-DO-025. Subtracting the last entry therefore leaves precisely the displayed partition. This proves (R11.3) from the actual weights, without choosing a different representation. The partition \(\eta_j\) describes the row coefficient metric; \(\nu_j\) describes the homogeneous kernel polynomial. They generally have different degrees.

For a partition \(\lambda\), put

\[
 H_\lambda=\prod_{(i,l)\in\lambda}
       (\lambda_i-l+\lambda_l'-i+1),\qquad H_\varnothing=1.
 \tag{R11.4}
\]

These are ordinary hook products, not the DP44 height constant \(H\). Trailing zero rows do not change them.

\subsection{2. The normalized Fischer calculation}\label{the-normalized-fischer-calculation}

On a complex \(b\)-by-\(a\) matrix variable \(Z\), use the Fischer inner product with distinct monomials orthogonal and

\[
 \|Z^\alpha\|_{\rm F}^2=\alpha!.
 \tag{R11.5}
\]

For \(\lambda_1\ge\cdots\ge\lambda_a\ge0\), let \(\lambda_{a+1}=0\), and define the primitive polynomial

\[
 p_\lambda(Z)=\prod_{r=1}^a
       \det(Z_{[1,r],[1,r]})^{\lambda_r-\lambda_{r+1}}.
 \tag{R11.6}
\]

Its leading diagonal monomial has coefficient one. In particular it is primitive in the ordinary integral monomial lattice. The exact norm is

\[
                         \boxed{\|p_\lambda\|_{\rm F}^2=H_\lambda.}
 \tag{R11.7}
\]

Here is a short derivation fixing the normalization, rather than a new proof of Schur--Weyl duality. Let \(D=|\lambda|\), and let \(f_\lambda\) be the dimension of the corresponding symmetric-group representation. Schur--Weyl duality and the hook formula give \(D!/f_\lambda=H_\lambda\). The exact references are Fulton and Harris (1991), Theorem 6.3(2)--(4), pp.77--78; Hook Length Formula 4.12, p.50; and the rectangular Cauchy decomposition in Exercise 6.11(b), p.80. These results are cited, not re-proved.

Choose the standard unitary metrics on the Schur modules, with highest vectors of norm one, compatible under the standard isometric inclusion \(\mathbf C^a\subset\mathbf C^b\). The degree-\(D\) Fischer reproducing kernel satisfies

\[
 \begin{aligned}
 \frac{(\operatorname{tr}(ZW^*))^D}{D!}
 &=\sum_{|\lambda|=D}\frac{f_\lambda}{D!}
                  \operatorname{tr}(S_\lambda(Z)S_\lambda(W)^*)\\
 &=\sum_{|\lambda|=D}\frac1{H_\lambda}
       \sum_{r,s}p_{rs}^\lambda(Z)\overline{p_{rs}^\lambda(W)}.
 \end{aligned}
 \tag{R11.8}
\]

The first equality is the trace of \((ZW^*)^{\otimes D}\) under the cited tensor decomposition. Functoriality and the unitary metrics give \(S_\lambda(W^*)=S_\lambda(W)^*\); this also follows from singular-value decomposition and the diagonal action on weights. The second equality expands the trace in orthonormal bases. The matrix coefficients \(p_{rs}^\lambda\) are linearly independent and exhaust the Cauchy summand. Comparison with the reproducing kernel therefore gives their common squared norm \(H_\lambda\) and their orthogonality. The highest-to-highest coefficient is (R11.6): both are highest vectors and evaluate on a diagonal matrix as \(\prod_i z_i^{\lambda_i}\), with coefficient one. This proves (R11.7), including its scalar.

\subsection{3. The determinant insertion has an exact scalar}\label{the-determinant-insertion-has-an-exact-scalar}

Let \(V_\eta\) be the irreducible compact coefficient module realized as the row orbit of \(p_\eta\) in square \(a\)-by-\(a\) polynomials. Multiplication by the full determinant gives

\[
 M_q:V_\eta\longrightarrow V_{\eta+q\mathbf1_a},
                         \qquad f\longmapsto(\det Z)^q f.
 \tag{R11.9}
\]

This is an integral polynomial map. On \(U(a)\), it intertwines the two representations after the unitary determinant twist. Hence its pulled-back Fischer form is a scalar multiple of the invariant form on \(V_\eta\). Evaluating on the highest vector and using (R11.7) proves, for every \(f\in V_\eta\),

\[
 \boxed{\|M_qf\|_{\rm F}^2
      =\frac{H_{\eta+q\mathbf1_a}}{H_\eta}\|f\|_{\rm F}^2.}
 \tag{R11.10}
\]

The ratio has the explicit integral product

\[
 \frac{H_{\eta+q\mathbf1_a}}{H_\eta}
       =\prod_{i=1}^{a}\frac{(\eta_i+q+a-i)!}{(\eta_i+a-i)!}
       =\prod_{i=1}^{a}\prod_{r=1}^{q}(\eta_i+a-i+r).
 \tag{R11.11}
\]

Indeed the factorial/Vandermonde expression for the hook product has the same Vandermonde denominator before and after the uniform shift by \(q\). Cancelling it gives (R11.11). This is also the hook-product expression directly used in the recovered source \nolinkurl{Divisibility/RecBridgeNum.lean}, theorem \nolinkurl{logHook_eq}, lines149--188. The displayed ratio is strictly positive.

The source defines its row metric as Fischer divided by the squared Fischer norm of its extremal vector: \nolinkurl{DataCoeff.lean}, \nolinkurl{normFock}, lines83--84, and \nolinkurl{HermitianForms/IsLocBasis.lean}, \nolinkurl{nfAll}, lines19--20. Thus on the actual row partition it is

\[
             \|f\|_{\mathrm{row},j}^2
                         =\|f\|_{\rm F}^2/H_{\eta_j}.
 \tag{R11.12}
\]

Combining (R11.10) and (R11.12), the selected determinant insertion has squared Fischer norm \(H_{\nu_j}\|f\|_{\mathrm{row},j}^2\). In normalized coefficient metrics it is an isometry. This last assertion concerns the coefficient map (R11.9); it is not a theta-operator isometry or an identity between automorphic measures.

In particular its hook ratios obey the exact cancellation

\[
 \boxed{\prod_{j=0}^{m-1}\frac{H_{\nu_j}}{H_{\eta_j}}
              =\frac{H_{\eta_m}}{H_{\eta_0}}.}
 \tag{R11.13}
\]

This follows from (R11.3). It prevents counting the same row-metric normalization repeatedly as an independent loss. No finite projection denominator, class mass, branch map, theta value, or DP44 term is asserted to cancel by (R11.13).

\subsection{4. Actual common-host Fock norm}\label{actual-common-host-fock-norm}

On \(E'=\mathbf Q(\zeta_{28})\), use the exact period and coordinates of AM-R-009. Each of the six coordinates is \(w_i^+=u_i+iv_i\) or \(w_i^-=u_i-iv_i\), and the actual creation operator is the one read by \(\partial/\partial w_i^\pm\). The monomial creation norm relative to the unit Gaussian is \(\alpha!(\pi/(2G_i))^{|\alpha|}\). This is the explicit orthogonal coordinate change in AM-R-009, retaining the factor \(1/2\) in each derivative. Let \(i(\tau)\in\{1,2,3\}\) name the underlying Jacobian coordinate of \(\tau\).

For a selected partition \(\nu_{j,\tau}\) at every real place, set \(D_{j,\tau}=|\nu_{j,\tau}|\). The tensor product of the primitive highest polynomials (R11.6), applied to the \textbf{actual} Gaussian, has exact squared archimedean norm

\[
 \boxed{\|\phi_j^{\mathrm{hw}}\|_\infty^2
  =\det(2Y)^{-ab}
       \prod_{\tau\in\Phi'}
       H_{\nu_{j,\tau}}
       \left(\frac{\pi}{2G_{i(\tau)}}\right)^{D_{j,\tau}}.}
 \tag{R11.14}
\]

There are \(ab\) copies of \(A=J^2\), so the squared vacuum norm is \(\det(2Y)^{-ab}\). At a fixed CM place all matrix entries have the same creation scale; homogeneity supplies its degree power and (R11.7) supplies its hook product. The CM coordinate blocks are orthogonal. These three facts prove (R11.14). Its square root is the norm, and its logarithm is

\[
 -\frac{ab}{2}\log\det(2Y)
   +\frac12\sum_\tau\left[
      \log H_{\nu_{j,\tau}}+
      D_{j,\tau}\log\frac{\pi}{2G_{i(\tau)}}\right].
 \tag{R11.15}
\]

The corresponding analytic-to-algebraic kernel still has period \(C^{2ab}\), by AM-DO-026. Unit normalization divides that kernel by the square root of (R11.14), and does not delete the period. Finite-vacuum norm one was used only for the unmoved standard vector; an additional finite linear combination retains its actual finite norm. A compact coefficient isometric embedding also retains the source Schur factor \(1/\dim V_{\eta_j}\) in Haar integration, as proved in AM-DO-026.

\subsection{5. A finite-degree integral projection bound}\label{a-finite-degree-integral-projection-bound}

Let \(L_D=\operatorname{Sym}^D(\mathbf Z^b\otimes(\mathbf Z^a)^*)\) be the ordinary integral polynomial lattice. Let \(E_\lambda\) be the rational projector onto its Cauchy component of shape \(\lambda\), and let \(L_\lambda=L_D\cap\operatorname{im}(E_\lambda)\). Then

\[
 \boxed{L_\lambda\subset E_\lambda L_D
                \subset H_\lambda^{-1}L_\lambda.}
 \tag{R11.16}
\]

To prove this without assuming integral semisimplicity, start on the tensor power \((\mathbf Z^b)^{\otimes D}\otimes((\mathbf Z^a)^*)^{\otimes D}\). Act on its first factor by the central symmetric-group idempotent

\[
 e_\lambda=\frac{f_\lambda}{D!}
             \sum_{\sigma\in S_D}\chi_\lambda(\sigma^{-1})\sigma
       =\frac1{H_\lambda}
             \sum_{\sigma\in S_D}\chi_\lambda(\sigma^{-1})\sigma.
 \tag{R11.17}
\]

The characters of the symmetric group are integers, and the central numerator commutes with the simultaneous permutation of both tensor factors. It therefore descends to the integral coinvariant quotient \(L_D\). Its denominator divides \(H_\lambda\). Over \(\mathbf Q\), Schur--Weyl duality identifies the descended operator with exactly the stated Cauchy projector. Hence \(H_\lambda E_\lambda L_D\subset L_\lambda\). Conversely it fixes every vector of \(L_\lambda\), proving the other inclusion in (R11.16).

At a prime \(p\), the possible saturation quotient is consequently killed by \(p^{v_p(H_\lambda)}\); its length is at most \(\operatorname{rank}(L_\lambda)v_p(H_\lambda)\). If every hook is a \(p\)-unit, the projection is integral and split. In particular \(p>\lambda_1+a-1\) suffices. This is an explicit upper bound for one projector, not a least-denominator formula. The source's particular Weyl/dual-Weyl lattice can admit a strictly better comparison. For example multiplication by \((\det Z)^q\) is already integral, so a large uniform shift does not by itself force the denominator bound in (R11.17) to be attained on that determinant-inserted submodule.

This finite-degree statement is consistent with AM-K019. It makes no assertion about a lattice simultaneously stable under all ordinary differential powers. Nor does a bound for a projector prove that the theta jet matrix on its image has a prescribed content or that a prescribed automorphic row belongs to that image.

\subsection{6. A concrete source-rank illustration}\label{a-concrete-source-rank-illustration}

Take the smallest displayed rank choice \(k=6,m=2,t=13,d=14\), and saturate only the \nolinkurl{HCParam} gaps as in AM-DO-025. Then

\[
 N_0=6,\quad N_1=20,\quad N_2=34,\qquad
 \nu_0=(1^6),\quad \nu_1=(3^6,2^{14}),
 \quad D_0=6,\quad D_1=46.
 \tag{R11.18}
\]

The row partitions are \(\eta_0=0\), \(\eta_1=(1^6,0^{14})\), and \(\eta_2=(3^6,2^{14},0^{14})\). Thus \(H_{\nu_0}=6!=720\), and the second determinant insertion ratio is

\[
 \frac{H_{\nu_1}}{720}
 =\prod_{i=1}^{20}\prod_{r=1}^{2}
                        (\eta_{1,i}+20-i+r).
 \tag{R11.19}
\]

All integers, prime valuations and the product identity for this example are reproduced by \nolinkurl{selected_schur_metric_certificate.py}. The small-polynomial checks there also expand (R11.6) in ordinary monomials and compare its exact Fischer norm to (R11.4); they are finite checks of normalization, not the proof of the general formula. The rank choice and gaps satisfy the conditions named here. No claim is made that they satisfy all the further noncompact \nolinkurl{Hyp3} conditions or already form a nonzero global chain.

\subsection{7. Framework and retained arithmetic gate}\label{framework-and-retained-arithmetic-gate}

The polynomial lattices and determinant maps are defined over \(\mathbf Z\); the rational projector in (R11.17) specifies exactly which primes need inverting for that chosen splitting. The associated Schur operations live in the representation Grothendieck ring and the special lambda-ring formalism of SGA6, Exposé V. The actual normalized metric and (R11.16) are stronger data than an equality of virtual classes. The representation ring is related to the flag construction by SGA6, Exposé VI, Theorem3.3, with the carried-source qualification in the geometric dictionary. The CM eigenspace splitting is the diagonalizable-group grading of SGA3, Exposé I, Proposition4.7.3. The period in (R11.14) remains the line-bundle metric of the Picard/theta construction, placed in FGA232 Theorem3.1 and SGA7 I VII2.7/2.9 in that dictionary.

For DP\_MULT Q3e these formulas supply the selected homogeneous degrees, an exact archimedean creation norm and its relation to the source's row metric. To produce actual \(g_j,e_j,H\), the next required data are still the allowed finite jet matrix on these partitions, its integral row realization, the actual nonzero iteration and the endpoint detector. The separate logarithmic quantities in (R11.15), (R11.16) and the class mass cannot simply be added as independent losses without checking the source maps: (R11.13) exhibits one exact normalization cancellation already.

\subsubsection{References and source locators}\label{references-and-source-locators}

\begin{itemize}
\tightlist
\item
  Fulton, W., \& Harris, J. (1991). \emph{Representation theory: A first course}. Springer. Hook Length Formula4.12, p.50; Theorem6.3(2)--(4), pp.77--78; Exercise6.11(b), p.80. The cited theorem statements and formulas were directly read.
\item
  Alpöge source, pinned commit recorded in the source audit: \nolinkurl{Formalism/LocalRingCardinality/DataCoeff.lean}, \nolinkurl{detShift}, \nolinkurl{natW}, \nolinkurl{extremal}, \nolinkurl{normFock} and \nolinkurl{normFockQ}, lines21--23, 55--56 and83--87; \nolinkurl{Formalism/LocalToGlobal/HermitianForms/IsLocBasis.lean}, \nolinkurl{nfAll}, lines19--20; \nolinkurl{Formalism/LocalToGlobal/Divisibility/RecBridgeNum.lean}, \nolinkurl{logHook_eq}, lines149--188. These files were directly read.
\item
  AM-DO-025, (DO25.5)--(DO25.17); AM-DO-026, (DO26.7)--(DO26.8) and (DO26.17)--(DO26.20); AM-R-009, corrected actual Gaussian norm; AM-K019, the full differential-module obstruction with its exact scope.
\end{itemize}

Whole-proof check events are recorded separately. No Lean execution or automatic proof verification is asserted.

\clearpage

\section{AM-K022 --- selected Schur jets in the actual permitted half-orbit}\label{am-k022-selected-schur-jets-in-the-actual-permitted-half-orbit}

Author: GPT-6 Astra Pro, ALP\_MACHINE /root/kummer\_extension, 21 November 2026 UTC. Inherited reasoning effort is not separately exposed.

This result gives an explicit arithmetic level and a single permitted finite Schwartz vector detecting any prescribed nonzero homogeneous Schur polynomial. It applies to the ordinary product of Jacobians on the common E′ host, including the actual first even-rank edge of AM-DO-025. It does not identify the second iterated image or compute its class-mass content.

The permitted half-orbit basis is \textbf{carried DP59}, not a new representation theorem. The new conclusions are the numerical interpolation threshold, its product specialization, the factorial-exact highest-symbol detection, the norm-one individual-translate conclusion, and the equality of all finite determinantal ideals after passing from jets to the permitted section basis.

\subsection{1. Fixed geometry and the permitted lattice}\label{fixed-geometry-and-the-permitted-lattice}

Work at 11 over an unramified complete DVR R large enough to contain the CM eigenvalues and the finite torsion and symmetric lifts in DP59. Normalize v₁₁(11)=1. Set \[
 J=\operatorname{Jac}(y^2=x^7-1),\qquad
 L=L_{2\infty},\qquad s=s_q,
\] where the actual q-normalized section is primitive at 11 by ALP-LIFT-08. Its reduced germ is σq by AM-K018. For an integer N≥1 write \[
 X=J^N,\qquad \mathcal L=L^{\boxtimes N},\qquad
 s_X=s^{\boxtimes N}.
 \tag{J1}
\] For the common host E′=Q(ζ₂₈), the positive Hom variety on an edge U(a)→U(b) is X=A\^{}\{ab\}=J\^{}\{2ab\}, with A=J², so N=2ab. Its holomorphic dimension is 6ab. The actual E′ eigenbasis differs from the two ordinary J bases by a matrix invertible at 11 (AM-DO-022). It may therefore be used for the integral tangent coordinates below.

Let \[
 \delta^2=-7,\qquad
 \alpha=(1+2\delta)^e,\qquad M=29^e,\qquad
 \alpha\bar\alpha=M,\qquad \kappa=M^2.
 \tag{J2}
\] The ideals (α) and \(\bar\alpha\) are coprime. On each J factor take the exact DP59 orbit \[
 s_{t'}=U_{t'}[M]^*s,\qquad t'\in J[\alpha^2],
\] with the canonical odd-order symmetric theta lift. Choose one representative of J{[}α²{]}/J{[}α{]}. DP59(a) gives an integral basis of \[
 W_M=[\alpha]^*H^0(J,L^M)
       \subset H^0(J,L^{M^2}),
 \qquad \operatorname{rank}_R W_M=M^3.
 \tag{J3}
\] The orbit is smaller than the full level-M² section module.

On X take the tensor basis \[
 \mathcal S_{\mathbf t'}
   =\boxtimes_{\ell=1}^{N}s_{t'_\ell},
 \qquad
 \mathcal W_{M,N}=W_M^{\otimes N}
   =[\alpha]^*H^0(X,\mathcal L^M).
 \tag{J4}
\] Its rank is M\^{}\{3N\}. Its actual rescaled Heisenberg parameters are \[
 \mathbf x=M\mathbf t'\in X[\alpha]\subset X[M].
 \tag{J5}
\] Thus all vectors used here belong to the allowed half-orbit. The construction never replaces it by the full level-M² theta group.

\subsection{2. An explicit integral interpolation threshold}\label{an-explicit-integral-interpolation-threshold}

\subsubsection{Proposition J6}\label{proposition-j6}

For every integer D≥0 and every M≥3D+2, restriction is surjective: \[
 H^0(J,L^M)\longrightarrow
 H^0\!\left(J,\ L^M\otimes\mathcal O_J/I_0^{D+1}\right).
 \tag{J6}
\] The analogous assertion holds on X=J\^{}N with its product line, with the same bound M≥3D+2, independent of N.

\textbf{Known geometric input.} On an abelian variety over a field, the third power of an ample line is very ample. Milne, \emph{Abelian Varieties}, I, §6, ``Ample divisors,'' printed p.~30 (PDF page 36 of AVc.pdf), states this in the chapter's arbitrary-characteristic setting. The higher-cohomology vanishing and integral H⁰ base change used below are the classical input already read and applied in DP59 §5, equations (TS11); its primary locator is Mumford (1966), §1(II), printed p.~289. No characteristic-zero vanishing theorem is substituted for the special-fibre assertion.

\textbf{Proof.} The special-fibre line L³ is very ample. Choose a section \(\bar h_0\) nonzero at zero, and sections \(\bar h_1,\bar h_2,\bar h_3\) vanishing there whose ratios \(\bar h_i/\bar h_0\) have linearly independent differentials. These choices exist because a very ample system has a surjective evaluation map and separates tangent vectors. They can be made over the residue field: these are surjective maps of finite-dimensional vector spaces, and an invertible minor can be chosen over that field.

Integral H⁰ base change lifts these sections to sections h\_i of L³ over R. Normalize h₀(0)=1 in a rigidified fibre. Subtract h\_i(0)h₀ from h\_i for i\textgreater0. The ratios \[
 t_i=h_i/h_0,\qquad 1\le i\le3,
\] are integral formal parameters at zero because their cotangent determinant is a unit. Integral linear combinations make their first differentials any specified integral basis, including the dual of the actual CM tangent basis. Only unit inverses are used.

For each integer m≥2 there is a section of L\^{}m whose value at zero is a unit. Here is the required special-fibre fact. The square L² is generated at every geometric point: choose a translate of a nonzero divisor of L so that neither it nor its opposite translate contains that point; the theorem of the square identifies their product line with L². The forbidden translation sets are two proper closed subsets. Their complement is nonempty over an algebraic closure. Evaluation surjectivity then descends to the residue field. L³ is generated because it is very ample. Every m≥2 is a sum of copies of 2 and 3, so products give the assertion for all m. Integral H⁰ base change lifts a section nonzero at zero to one with unit value over R.

Take m=M−3D≥2 and such a section η\_m of L\^{}m. The section η\_mh₀\^{}D trivializes L\^{}M along the D-th infinitesimal neighbourhood of zero. An arbitrary target jet in this frame is \[
 Q(t)=\sum_{|\beta|\le D}q_\beta t^\beta,\qquad q_\beta\in R.
\] The global section \[
 \eta_m\sum_{|\beta|\le D}
       q_\beta h_0^{D-|\beta|}
                h_1^{\beta_1}h_2^{\beta_2}h_3^{\beta_3}
 \tag{J7}
\] has exactly this jet. This proves (J6).

For the product, tensor the single-factor restriction maps through order D. A monomial of total degree at most D in all 3N parameters has degree at most D in each three-variable block. Products of the single-factor interpolating sections therefore span every such monomial. After multiplication by the product local frame unit, they span arbitrary total D-jets. The image is contained in H⁰(X,ℒ\^{}M) by the usual product-section map. This proves the product assertion without a dimension-dependent change of M. ∎

This is the numerical specialization r=3, b=2 of the integral homogenization in accepted DP56(a), not a new claim that an arbitrary ample line on an arbitrary proper variety has that threshold.

\subsubsection{Corollary J8. The same bound for the allowed half-orbit}\label{corollary-j8.-the-same-bound-for-the-allowed-half-orbit}

For M=29\^{}e≥3D+2, \[
 \mathcal W_{M,N}\longrightarrow
 H^0\!\left(X,\mathcal L^{M^2}
                    \otimes\mathcal O_X/I_0^{D+1}\right)
 \quad\text{is surjective over }R.
 \tag{J8}
\]

\textbf{Proof.} The étale isogeny {[}α{]} has degree prime to 11, hence induces an integral formal isomorphism at zero. The rigidified equality {[}α{]}*ℒ\^{}M≃ℒ\^{}\{M²\} is part of the DP59 biextension construction. Its comparison in integral local frames is multiplication by an integral unit. Pulling (J6) back along this formal isomorphism proves (J8). The tensor orbit basis in (J4) is the carried DP59 basis on each factor, so the surjection is from exactly the permitted lattice. ∎

One explicit choice is \[
 e(D)=\max\!\left(1,\left\lceil\log_{29}(3D+2)\right\rceil\right),
 \qquad M(D)=29^{e(D)},
 \qquad
 3D+2\le M(D)<29(3D+2).
 \tag{J9}
\] In particular κ=M²\textless29²(3D+2)². The threshold is for \textbf{ordinary algebraic jets modulo I₀\^{}\{D+1\}}, not for divided powers or derivative values normalized to one.

\subsection{3. Highest-symbol detection and its exact factorial}\label{highest-symbol-detection-and-its-exact-factorial}

Use any actual integral invariant tangent basis \(D_1,\ldots,D_{3N}\), and choose formal parameters t\_i with \(D_it_j(0)=\delta_{ij}\). In the characteristic-zero logarithmic coordinates u\_i dual to these derivations, \(t_i(u)=u_i+O(u^2)\).

For a section S of ℒ\^{}\{M²\}, let F\_S(u) be its actual Chern/source reduced germ in the common normalization of AM-K018. In an integral algebraic frame normalized at zero it has the form \[
 F_S(u)=f_S(t(u))\,G_M(u),\qquad G_M(0)=1.
 \tag{J10}
\] The gauge G\_M is a characteristic-zero series; this argument does not require it to be an integral ordinary power series. The equality uses the fixed source gauge, not a section-dependent gauge.

Let \[
 P(Z)=\sum_{|\beta|=D}c_\beta Z^\beta
          \in R[Z_1,\ldots,Z_{3N}]
\] be nonzero and homogeneous, and define the normalized functional \[
 \ell_{P,M}(S)=
       M^{-D}\,[P(\partial_u)F_S(u)]_{u=0}.
 \tag{J11}
\]

\subsubsection{Proposition J12}\label{proposition-j12}

For every \(\beta\in\operatorname{supp}P\), the image ideal of ℓP,M on the permitted lattice contains \(c_\beta\beta!\): \[
 c_\beta\beta!\in\ell_{P,M}(\mathcal W_{M,N}),
 \qquad
 \beta!=\prod_i\beta_i!,
 \qquad M\ge3D+2.
 \tag{J12}
\] Consequently at least one individual allowed tuple \(\mathbf t'\) satisfies \[
 \ell_{P,M}(\mathcal S_{\mathbf t'})\ne0,\qquad
 v_{11}\!\left(\ell_{P,M}(\mathcal S_{\mathbf t'})\right)
 \le \min_{\beta\in\operatorname{supp}P}
                     v_{11}(c_\beta\beta!).
 \tag{J13}
\]

\textbf{Proof.} By (J8), choose S with algebraic jet \[
 f_S(t)\equiv M^D t^\beta\pmod{(t)^{D+1}}.
 \tag{J14}
\] The scalar M\^{}D is a unit at 11. There are no terms below degree D. Substituting \(t(u)=u+O(u^2)\) and multiplying by \(G_M(u)=1+O(u)\) therefore gives \[
 F_S(u)\equiv M^D u^\beta\pmod{(u)^{D+1}}.
 \tag{J15}
\] All higher coefficients of the possibly nonintegral gauge are irrelevant to this equality. The homogeneous differential polynomial kills every degree-D monomial except the one matching its exponent: \[
 M^{-D}P(\partial_u)M^D u^\beta\big|_0
        =c_\beta\prod_i\beta_i!.
 \tag{J16}
\] This proves (J12).

Expand S in the integral basis (J4): \[
 S=\sum_{\mathbf t'}a_{\mathbf t'}\mathcal S_{\mathbf t'},
 \qquad a_{\mathbf t'}\in R.
 \tag{J17}
\] If all basis values had valuation greater than \(v_{11}(c_\beta\beta!)\), their integral linear combination could not equal the nonzero value in (J16). Choose β minimizing that valuation. At least one basis value obeys (J13). The tuple and its canonical theta lifts are global algebraic objects by DP59(c). The conclusion does not require the coefficients in (J17) themselves to be global algebraic numbers. ∎

The factorial in (J16) is essential. Interpolating the ordinary coefficient of t\^{}β to one produces derivative β!, not derivative one. An assertion of a unit derivative at arbitrary degree would divide by β! and would not be an integral interpolation statement.

The argument neither bounds nor uses the archimedean size of the coefficients in (J17). It uses them only in the local nonarchimedean valuation inequality to select one actual basis vector.

\subsection{4. The exact selected Schur polynomial}\label{the-exact-selected-schur-polynomial}

On the even-rank edge U(a)→U(b), b\textgreater a, let \(\nu_\tau\) be the actual nonnegative source partitions identified in AM-DO-025: \[
 \nu_{\tau,i}
   =(\lambda_{\mathrm{True},\tau})_i
                        +e_\tau-b/2,\qquad 1\le i\le a.
 \tag{J18}
\] There are six holomorphic labels τ on the E′ host. Put ντ,a+1=0 and write \(Z_\tau=(Z_{\tau,i,j})\) as an a-by-b matrix. Take the standard primitive highest bideterminant \[
 P_\nu(Z)=
 \prod_\tau\prod_{k=1}^{a}
       \det(Z_{\tau,\{1,\ldots,k\},\{1,\ldots,k\}})
                     ^{\nu_{\tau,k}-\nu_{\tau,k+1}}.
 \tag{J19}
\] Its homogeneous degree is exactly \[
 D=\sum_\tau|\nu_\tau|.
 \tag{J20}
\] It is the highest vector in the actual Cauchy summand of AM-DO-025. This polynomial already lies in the selected Schur space; applying a rational Schur projector to it is unnecessary.

The diagonal monomial of Pν has coefficient exactly one: \[
 \prod_{\tau,i}Z_{\tau,i,i}^{\nu_{\tau,i}}.
 \tag{J21}
\] Indeed every determinant must contribute its diagonal term to produce a monomial containing no off-diagonal variable. The exponent of its i-th diagonal variable is \(\sum_{k\ge i}(\nu_{\tau,k}-\nu_{\tau,k+1})=\nu_{\tau,i}\). This also proves that Pν is primitive in the polynomial coefficient lattice. It does not assert that every alternative Weyl or normalized Fock lattice has that same basis; their comparison is separate.

Applying (J13) with this diagonal monomial gives a single permitted tuple with \[
 a_{\nu,\mathbf t'}
 :=M^{-D}[P_\nu(\partial)F_{\mathcal S_{\mathbf t'}}]_0
 \ne0,
\qquad
 v_{11}(a_{\nu,\mathbf t'})
 \le
 \sum_{\tau,i}v_{11}(\nu_{\tau,i}!).
 \tag{J22}
\] Equivalently, \[
 v_{11}(a_{\nu,\mathbf t'})
 \le
 \frac{D-\sum_{\tau,i}s_{11}(\nu_{\tau,i})}{10}.
 \tag{J23}
\] This is an explicit upper bound for \textbf{one selected kernel coefficient on the allowed family}. It is not a sharp derivative slope and not a bound for an arbitrary later theta-image vector.

It works for every actual ν₀ satisfying the source recurrence and gap conditions. No fixed-tail or single-row restriction is imposed. Therefore it supplies the nonzero selected first-edge kernel needed for the parent's class-supported compact extraction on U(N₀)→U(N₁).

\subsection{5. The individual Schwartz vector has no translation-combination norm cost}\label{the-individual-schwartz-vector-has-no-translation-combination-norm-cost}

Let \(\phi_{0,f}\) be the actual finite vacuum on X=A\^{}\{ab\}, normalized to finite L² norm one as in AM-DO-022. For the tuple selected in (J22), use \[
 \phi_{\mathbf x,f}=W_{\mathbf x}\phi_{0,f},
 \qquad \mathbf x=M\mathbf t'\in X[\alpha].
 \tag{J24}
\] Every local Heisenberg translation is unitary for the fixed self-dual Schrödinger measure. The exact symmetric lift differs only by its fixed finite-order unit scalar. Hence \[
                 \|\phi_{\mathbf x,f}\|_2^2=1.
 \tag{J25}
\] In its permitted-orbit expansion this chosen vector has one coefficient equal to one and all others zero. Its coefficient maximum is one at every finite and complex embedding. Thus selecting this vector introduces no coefficient height or sum-of-orbit norm cost.

The level-normalized identity of ALP-LIFT-10 and AM-R-010 is \[
 M^{-D}P(\partial_u)
 U_{\mathbf t'}^{M^2H}\bigl(\sigma_q^{\boxtimes N}(Mu)\bigr)\big|_0
 =
 P(\partial_v)
 U_{\mathbf x}^{H}\sigma_q^{\boxtimes N}(v)\big|_0,
 \tag{J26}
\] with the specified canonical finite-order scalar retained. Accordingly AM-DO-026's actual kernel has the nonzero identity coefficient \[
           \vartheta_{r_{P_\nu}\phi_{\mathbf x,f}}(1,1)
                  =C^{2ab}a_{\nu,\mathbf t'}.
 \tag{J27}
\] Here \(C=\theta[0;0](0;\tau)c\), with the fixed algebraic scalar \(c\) of AM-R-006, is the one common period for the \(\sigma_q\) seed in AM-DO-022, §5. It is not a coefficient chosen for this tuple. The archimedean norm remains precisely the norm of the creation vector \(r_{P_\nu}\phi_{0,\infty}\); (J25) removes no Fock or period factor.

This statement supplies the finite vector and its norm. The source automorphic input, its stabilizer mass, its compact Schur-orthogonality factor, and its normalization are constructed separately in the parent's AM-R-012. None of those factors is set equal to one here.

\subsection{6. Exact finite saturation reduction for every chosen polynomial family}\label{exact-finite-saturation-reduction-for-every-chosen-polynomial-family}

The same interpolation gives a useful matrix statement that does not rely on a leading-symbol approximation.

Choose finitely many homogeneous polynomials P₁,\ldots,P\_q of degree at most D in a specified integral selected lattice, and include their individual source factors M\^{}\{-deg Pj\}. Choose an ordinary algebraic monomial basis of the total D-jet module in the fixed integral theta frame. Let \[
 J_D:R^{\,M^{3N}}\longrightarrow R^{\,\binom{3N+D}{3N}}
 \tag{J28}
\] be the matrix of the actual permitted section basis's ordinary jets. Let B be the q-by-\(\binom{3N+D}{3N}\) matrix of the corresponding normalized Chern/source derivative functionals on these jets. Its entries lie in Frac(R); no integral reduced frame is assumed. Let A be the q-by-M\^{}\{3N\} matrix of their actual allowed-tuple values. Then \[
                       A=BJ_D.
 \tag{J29}
\]

\subsubsection{Proposition J30. No extra finite-translation determinantal loss}\label{proposition-j30.-no-extra-finite-translation-determinantal-loss}

For M≥3D+2, the map J\_D is split surjective over R. For every k≥1 the fractional ideals generated by the k-by-k minors satisfy \[
                     I_k(A)=I_k(B).
 \tag{J30}
\] Equivalently, after adding zero columns, A and B have the same nonzero Smith factors over the DVR.

\textbf{Proof.} Surjectivity is (J8). The jet module is finite free, so choose an R-linear right inverse C with J\_DC=1. Thus B=AC. Cauchy--Binet applied to A=BJ\_D gives \(I_k(A)\subset I_k(B)\), since J\_D is integral. Applied to B=AC it gives the reverse inclusion, since C is integral. For fractional matrices first multiply by one common denominator; both k-minor ideals scale by its k-th power, so cancellation recovers (J30). Alternatively split the free domain of J\_D as its kernel plus a copy of its target; in adapted integral bases J\_D is \((1\ 0)\), and A is \((B\ 0)\). ∎

This equality is about the actual allowed half-orbit, not a larger arbitrary Schwartz space. It quantifies the absence of additional saturation loss in passing from the finite jet module to these section generators. The integral right inverse may have large archimedean entries, but no norm estimate for it is used in (J22) or (J25).

For the second automorphic lift, its class integral and the restriction to the already produced first image impose additional linear maps on this kernel data. Those maps can change the Smith factors, can cause cancellation, and can annihilate a vector. Proposition J30 does not assert that independently chosen jets at the identity can be prescribed at every class representative. The remaining second-stage problem is precisely the content of that \textbf{restricted class matrix}, with the source's actual selected coefficient lattice and compact contraction factors included.

The full differential-module obstruction AM-K019 is consistent with (J30). Here the degree-D jet module is finite, the level grows with D, and only the specified evaluated functionals are retained. No fixed-rank D-stable lattice, bounded saturation at all degrees, or iteration of a mod-11 coefficient congruence is asserted.

\subsection{7. Sources and checks}\label{sources-and-checks}

\begin{itemize}
\tightlist
\item
  DP\_MULT DP56(a), research/quantitative\_theta\_jet\_levels\_DP56.md, §§1--3: integral homogenization and its precise ordinary-jet convention. Carried; the explicit choice r=3,b=2 is justified in Proposition J6.
\item
  DP\_MULT DP59(a)--(c), research/theta\_translate\_span\_DP59.md, §§3--8: actual rigidified pullback, exact permitted half-orbit, integral basis, unramified normalization, and global algebraicity. Carried, not rederived.
\item
  Milne, J. S. (2008). \emph{Abelian varieties} (Version 2.0), I, §6, subsection ``Ample divisors,'' printed p.~30, paragraph immediately before Example 6.5. \url{https://www.jmilne.org/math/CourseNotes/AVc.pdf}. The precise used fact is that the third power of every ample line on an abelian variety is very ample; the chapter is over an arbitrary base field.
\item
  Mumford, D. (1966). On the equations defining abelian varieties. I. \emph{Inventiones Mathematicae, 1}, 287--354. §1(II), p.~289, as already audited in DP59 for the cohomology/base-change input. \url{https://pazuki.perso.math.cnrs.fr/index_fichiers/Mumford66.pdf}.
\item
  AM-DO-022 and AM-DO-025--026: actual E′ product polarization, unitary finite Heisenberg action, selected partitions, and all-polynomial kernel identity.
\item
  AM-K018 and ALP-LIFT-08: exact source gauge on ordinary J products and primitive actual principal section at 11.
\item
  ALP-LIFT-10 and AM-R-010: exact permitted parameter \(\mathbf x=M\mathbf t'\), retained symmetric lift scalar, and level-normalized derivative identity.
\end{itemize}

The result does not assume or bound an interpolating section's global coefficient height. The chosen individual allowed translate has coefficient vector a standard basis vector, so its finite Schwartz norm and translation-coefficient bound are exact.

\clearpage

\section{AM-R-012: an actual selected even-rank first lift with a single permitted finite translate}\label{am-r-012-an-actual-selected-even-rank-first-lift-with-a-single-permitted-finite-translate}

Author: /root, GPT-6 Astra Pro, ALP\_MACHINE, 21 November 2026 UTC.

The rank-one starting correspondence of AM-R-009--010 is not the starting edge in the recovered \nolinkurl{Ranks} record. The construction here uses the actual even ranks \(N_0\to N_1\), their selected compact coefficient types, and the same splitting character on both groups. It proves a controlled nonzero \textbf{first} theta lift with an exact source input and a single permitted finite Heisenberg translate. It does not yet identify its arithmetic row with the old elliptic \nolinkurl{StdPair}, or prove the later chain and six comparison identities.

The finite interpolation input is AM-K022, which specializes the accepted DP56/DP59 allowed half-orbit theorem to the actual selected polynomial. That proof and its check status are separate. In the form used here it supplies one permitted translate with finite Schwartz norm one and the nonzero jet bound (R12.5) below.

\subsection{1. Actual fields, types, polynomial and character}\label{actual-fields-types-polynomial-and-character}

Use \[
 E'=\mathbf Q(\zeta_{28}),\quad F'=E'^+,\quad A=J^2,
 \quad G=U(I_a),\quad H=U(I_b),\quad a=k,
 \quad b=k+d=N_1.
 \tag{R12.1}
\] All six selected complex places are those of AM-R-009, and the positive Hom oscillator, trace-one vacuum, and period are exactly those of AM-DO-022/026. Let \(w_0,w_1\) and \(\nu_{0,\tau}\) be the weights and partitions of an actual \nolinkurl{HCParam}, as in AM-DO-025. Write \[
 D=\sum_{\tau\in\Phi'}|\nu_{0,\tau}|,
 \quad P(Z)=\prod_{\tau\in\Phi'}p_{\nu_{0,\tau}}(Z_\tau).
 \tag{R12.2}
\] Here each highest polynomial is the primitive bideterminant of AM-R-011. The actual source compact type of the oscillator is \(\sigma=\rho_0^\vee\), and the target compact type is \(\rho_1\), where \(\rho_j\) has highest weight \(w_j\). This is precisely the equal-character identification proved in AM-DO-025, not the trivial-character rank-two special case.

There is an explicit globally admissible common character with the required exponents \(-e_{0,\tau}\). Indeed AM-R-009 proves \(T'(\mathbf A_f)=T'(F')K_{T,0}\), with intersection \(\mu_{28}\). On the real norm-one torus set \(\beta(t)=\prod_\tau t_\tau^{-e_{0,\tau}}\). At one split place over29 choose the factor where \(\zeta_{28}\mapsto2\), and define the compact residue character \[
 \varrho(2)=\exp\!\left(\frac{2\pi i}{28}
                         \sum_{\tau\in\Phi'}e_{0,\tau}\tau\right).
 \tag{R12.3}
\] It is trivial at every other compact place. Thus \(\varrho|_{\mu_{28}}=\beta^{-1}\), so \(\beta\varrho\) descends to a unitary automorphic norm-one character \(\gamma\). Put \(\mu(x)=\gamma(x/\bar x)\). Then \(\mu|_{\mathbf A_{F'}^\times}=1\), \(\mu_\infty(z)=(z/\bar z)^{-e_{0,\tau}}\), and the finite compact conductor divides29. Take \(\chi_S=\chi_T=\mu\) on the edge. At nonsplit finite places the norm-one group is compact and \(\gamma\) is trivial there; at other split places it is unramified. All values needed for the algebraic finite kernel lie in a finite algebraic coefficient extension. This is the explicit character construction of AM-R-009 applied to the present integral exponent vector.

Choose \[
 e=\max\{1,\lceil\log_{29}(3D+2)\rceil\},\qquad M=29^e.
 \tag{R12.4}
\] By AM-K022, after a fixed unramified coefficient extension at11, one can choose a \textbf{single} permitted product-translate \(x\) such that its normalized geometric jet \(a_x\) satisfies \[
 a_x\ne0,\qquad
 v_{11}(a_x)\le
       B_\nu:=\sum_{\tau\in\Phi'}\sum_{i=1}^{a}
                              v_{11}(\nu_{0,\tau,i}!).
 \tag{R12.5}
\] The integer in the bound is the derivative of the diagonal monomial in \(P\); it is not the hook product and not an unproved lower bound for every selected jet. The translate belongs to the actual DP59 half-orbit after its specified odd-level normalization. The corresponding finite Schwartz vector is \(\phi_{f,x}=W_x\phi_0\) with its actual central multiplier, and \(\|\phi_{f,x}\|=1\) by unitarity. No interpolation coefficient norm is hidden in this selection of one orbit vector.

Let \(\phi_x=r_P\phi_{f,x}\phi_\infty\). The actual kernel identity AM-DO-026, with the normalized level factors retained, gives \[
                  \vartheta_{\phi_x}(1,1)=C^{2ab}a_x\ne0.
 \tag{R12.6}
\]

\subsection{2. A concrete source level and its exact identity-class mass}\label{a-concrete-source-level-and-its-exact-identity-class-mass}

Let \(K_{a,0}\) be the standard full integral unitary compact. Refine it to \(K_a(M)\) by principal congruence modulo29\^{}e at all six split places over29, leaving all other places unchanged. Then \(K_a(M)\) fixes \(\phi_{f,x}\) exactly in the source action. The same assertion holds for the target subgroup \(K_b(M)\).

To see this, at29 every permitted parameter has denominator dividing \(M\) and lies in the selected isotropic split position space. The linear action \(X\mapsto hXg^{-1}\) preserves that position space. For \(g,h\equiv1\pmod M\), its displacement is integral in the same space. The pairing with the original parameter is zero, and the odd-place trace-one lattice character of that displacement is one. The exact Weyl law therefore fixes \(W_x\phi_0\). The compact determinant characters also equal one on this level by (R12.3). Away from29 the parameter is integral; the line is the original trace-one vacuum line, fixed by the standard compacts and the unramified splitting. These are exactly the split/unramified normalizations of AM-DO-022 and the isotropic half-orbit argument of AM-R-010, now applied entrywise to the matrix Hom space.

At the identity finite class the rational stabilizer is trivial: \[
                   \Gamma=G(F')\cap K_a(M)=\{1\}.
 \tag{R12.7}
\] First \(\Gamma\) is finite, because it preserves the integral lattice and lies in a compact group at infinity. Its image at one split place over29 is inside \(1+29M_a(\mathbf Z_{29})\), which is torsion free. For completeness, if a nonzero matrix \(B\) has minimum entry valuation \(s\ge1\), then \((1+B)^{29}-1\) has minimum valuation exactly \(s+1\): the term \(29B\) has that valuation, and all the other terms have strictly larger valuation. Raising to an integer prime to29 leaves the minimum valuation \(s\). A finite-order element in this group would therefore remain nonidentity after raising to its order, a contradiction. The local embedding of the rational matrix group is injective, proving (R12.7).

Normalize real Haar measure to have \(\operatorname{vol}(G_\infty)=1\), and finite Haar measure by \(\operatorname{vol}(K_{a,0})=1\). Use their compatible automorphic quotient measure; its total volume is not silently set to one. The exact index and the identity-class mass are \[
 \begin{aligned}
 I_a(M)&=[K_{a,0}:K_a(M)]
       =29^{6a^2(e-1)}|\mathrm{GL}_a(\mathbf F_{29})|^6\\
       &=M^{6a^2}\prod_{i=1}^{a}(1-29^{-i})^6,
       \qquad w=I_a(M)^{-1}.
 \end{aligned}
 \tag{R12.8}
\] The split reduction maps are surjective, with successive congruence quotients of cardinality \(29^{a^2}\). Counting ordered bases gives \(|\mathrm{GL}_a(\mathbf F_{29})|=29^{a^2}\prod_i(1-29^{-i})\). There are six places, and (R12.7) removes the stabilizer divisor at this class. This proves (R12.8) with the actual chosen measure.

\subsection{3. Compact matrix-coefficient extraction}\label{compact-matrix-coefficient-extraction}

Use the standard unitary metrics on the compact representations, with highest vector of norm one. Let \(u\in\sigma\) be the unit source vector corresponding to the source factor of the highest polynomial \(P\), and put \[
 d_\rho=\dim\rho_0=\dim\sigma,
 \qquad m(k)=\langle u,\sigma(k)u\rangle.
 \tag{R12.9}
\] The Hermitian convention is conjugate-linear in the first argument. At the identity finite class, the kernel can be written \[
             \vartheta_{\phi_x}(k,1)=L(\sigma(k)u),
             \qquad L(u)=C^{2ab}a_x,
 \tag{R12.10}
\] where the scalar relating the actual highest polynomial to unit Schur vectors is included in \(L\). This is the actual source equivariance of the theta kernel; its finite translate acts on a different tensor factor and does not change the compact type.

Define an automorphic function \(f_0\) supported on the identity finite class, whose value there is \(\overline{m(k)}\), and whose value on all other finite classes is zero. Equation (R12.7) makes this definition unambiguous. It is right \(K_a(M)\)-invariant and smooth at infinity, with right compact type \(\sigma^\vee=\rho_0\). This gives an actual automorphic input, rather than a free representation parameter. It is cuspidal because the definite rational group has no proper rational parabolic subgroup.

Schur orthogonality, or the average of the rank-one projection \(v\mapsto u\langle u,v\rangle\), gives \[
 \int_{G_\infty}\sigma(k)u\overline{m(k)}\,dk
                      =u/d_\rho,
 \qquad \int_{G_\infty}|m(k)|^2\,dk=1/d_\rho.
 \tag{R12.11}
\] Insert this into the literal \textbf{bilinear} theta integral, and use the mass \(w\) of the supporting class. One obtains the exact nonzero value and norm \[
 \boxed{\Theta^{\rm bil}_{\phi_x}(f_0)(1)
        =\frac{C^{2ab}a_x}{I_a(M)d_\rho},\qquad
        \|f_0\|_2^2=\frac1{I_a(M)d_\rho},\quad
        \|f_0\|_\infty=1.}
 \tag{R12.12}
\] The supremum is attained at the real identity. In the integral coefficient presentation this is a class vector with primitive highest coefficient: the highest-to-highest matrix coefficient is the primitive bideterminant with the integral determinant character attached, as in AM-R-011. Its value at the identity is one. No claim that all scalar values on the real compact group are algebraic integers is needed or intended.

Two useful alternative normalizations are consequently \[
 \begin{array}{c|c|c}
 \text{input}&\text{squared norm}&\text{output at 1}\\\hline
 I_a(M)d_\rho f_0&I_a(M)d_\rho&C^{2ab}a_x\\
 \sqrt{I_a(M)d_\rho}\,f_0&1&C^{2ab}a_x/\sqrt{I_a(M)d_\rho}.
 \end{array}
 \tag{R12.13}
\] Thus the source normalization cost for the extracting input is exactly \(\tfrac12\log(I_a(M)d_\rho)\). It has not been removed by using a nonprimitive scalar multiple of the row.

The compact dimension is explicitly \[
 d_\rho=\prod_{\tau\in\Phi'}\prod_{1\le i<l\le a}
      \frac{\eta_{0,\tau,i}-\eta_{0,\tau,l}+l-i}{l-i}.
 \tag{R12.14}
\] With fixed ranks and \(M<29(3D+2)\), the logarithmic source cost is \(O(\log(D+2))\), with the exact finite index given in (R12.8), rather than an unspecified level constant. The full archimedean Schwartz norm is (R11.14), and the finite Schwartz norm is one. Unit normalization of both input and Schwartz vector therefore divides (R12.6) by \(\sqrt{I_a(M)d_\rho}\,\|\phi_{x,\infty}\|\).

The normalized algebraic identity coefficient in the primitive row convention has the controlled upper valuation \[
 v_{11}\!\left(C^{-2ab}\Theta^{\rm bil}_{\phi_x}(f_0)(1)\right)
          \le B_\nu-v_{11}(I_a(M)d_\rho).
 \tag{R12.15}
\] This is a selected nonzero entry, not a lower content bound for every entry of an arithmetic row map. It is therefore not itself the gain \(g_0\) in DP44.

\subsection{4. A selected first coefficient with no positive 11-adic gain}\label{a-selected-first-coefficient-with-no-positive-11-adic-gain}

For the \nolinkurl{HCParam} example in AM-R-011, take the same parameters at all six CM places. Then \(a=6,b=20\), \(\nu_{0,\tau}=(1^6)\), \(\eta_{0,\tau}=0\), \(D=36\), \(d_\rho=1\), and the explicit choice (R12.4) is \(M=29^2=841\). The polynomial is the product of six six-by-six determinants, and \(B_\nu=0\).

Furthermore \[
 I_6(841)=29^{216}|\mathrm{GL}_6(\mathbf F_{29})|^6,
                         \qquad v_{11}(I_6(841))=0.
 \tag{R12.16}
\] Indeed the order of29 modulo11 is10: \(7^5=-1\) modulo11, while \(7\) and \(7^2\) are not one. None of \(29^i-1\), \(1\le i\le6\), is divisible by11. Equations (R12.5), (R12.12) and (R12.16) therefore give an actual permitted first lift for which \[
 v_{11}\!\left(C^{-240}\Theta^{\rm bil}_{\phi_x}(f_0)(1)\right)
                                  \le0.
 \tag{R12.17}
\] Thus a strictly positive uniform 11-adic content for all these allowed selected first coefficients would be false. The statement is an upper bound; equality zero would additionally require integrality of all the relevant reduced jet coefficients. Neither that extra integrality nor the further noncompact \nolinkurl{Hyp3} conditions are assumed in this example.

Its exact archimedean Schwartz norm is \[
 \|\phi_{x,\infty}\|^2
       =\det(2Y)^{-120}(6!)^6
                    \prod_{i=1}^{3}\left(\frac{\pi}{2G_i}\right)^{12}.
 \tag{R12.18}
\] The primitive scalar input has squared norm \(I_6(841)^{-1}\). All these quantities belong to the same actual first edge. They do not bound the gains of the later restricted image or the rebuilt-family DP44 difference.

\subsection{5. What has and has not been supplied}\label{what-has-and-has-not-been-supplied}

The construction supplies one actual first edge at the source's even ranks with its selected local types, a globally admissible equal character, one permitted finite input of norm one, a specific class-supported source function, and an exact nonzero value up to the explicitly chosen algebraic jet \(a_x\). The finite selection can be made deterministic by ordering the finite permitted orbit and choosing the first tuple with the valuation bound (R12.5). No search over an infinite Schwartz space is required for this first edge.

It does not prove that every source row survives, or that the resulting vector survives the next prescribed kernel. It does not replace the old elliptic \nolinkurl{HostRaw}/\nolinkurl{StdPair} by definition, identify the independently fixed arithmetic \nolinkurl{rowAF}, or verify the noncompact analytic assumptions in \nolinkurl{AdmS}. Those are still needed for the requested shorter chain and the DP44 comparison. The exact first value, norms and index above are available for that next step, without the previous rank-one starting mismatch.

Geometrically, the permitted translates are sections in the Picard/theta-group construction, and (R12.5) is finite-jet separation on the good integral model. The compact coefficient is a Schur functor over \(\mathbf Z\), with the rational class-function realization retaining its actual metric. The relevant locators are FGA232 Theorem3.1; FGA236 Corollaries1.7 and2.3; SGA7 I VII2.7/2.9; EGA IV4 section16 on principal parts; SGA3 I Proposition4.7.3 and SGA6 V/VI as qualified in the geometric dictionary. No new crystalline or motivic L-function implication is asserted by this first lift.

Exact carried inputs: DP56--59; AM-K022; AM-DO-022, AM-DO-025 and AM-DO-026; AM-R-009--011. The source bilinear integral is \nolinkurl{WittSmallGroupEmbeddingMemPointwiseStabilizerIso.lean}, \nolinkurl{liftT_apply}, lines160--167. The selected weights are \nolinkurl{Columns/BaseLemmas.lean}, \nolinkurl{lamTrue_old}, \nolinkurl{lamTrue_new}, and \nolinkurl{nuW_eq_true}, lines496--552. The source row normalization is \nolinkurl{DataCoeff.lean}, \nolinkurl{normFock}, lines83--84. No Lean build or formal source-record instantiation is claimed.

\clearpage

\section{ALP-LIFT-13 --- fixed finite first-image witnesses at 3 and 5}\label{alp-lift-13-fixed-finite-first-image-witnesses-at-3-and-5}

Author: /root/lift\_accounting, GPT-6 Astra Pro, ALP\_MACHINE, 21 November 2026 UTC.

This specializes the already constructed actual first theta correspondence to AM-K020's sharp derivatives of one fixed theta section. It does not reprove those derivative congruences. Unlike the general permitted-orbit projector, this untranslated specialization uses one unit-norm automorphic character and one fixed finite Schwartz vector.

\subsection{1. Actual vector and its exact coefficient}\label{actual-vector-and-its-exact-coefficient}

Retain \(E=\mathbf Q(\zeta_7)\), \(F=E^+\), the actual compact pair \(T=U(1)\), \(G=U(I_2)\), and the honest splitting of ALP-LIFT-09. Write \[
 a_{i,n}=\partial_i^n\sigma_q(0),\qquad
 C=\theta[0;0](0;\tau)c .
 \tag{F1}
\] The normalization of \(\sigma_q\) is the actual algebraic Chern/source gauge established in AM-K014/K018 and ALP-LIFT-08: \(\sigma_q(u)=u_1u_3-u_2^2+\text{higher terms}\). The global period \(C\) is retained.

On two ordinary Jacobian factors use the creation polynomial \[
       P_{i,n}=\partial_{i,1}^{\,n}
                         \partial_{1,2}\partial_{3,2}.
 \tag{F2}
\] The finite Schwartz vector is the same \(\phi_{0,f}^{\otimes2}\) for every \(i,n\). Its finite \(L^2\) norm is one. Let \(\phi_{i,n}\) have this finite part and the actual Fock creation vector corresponding to (F2). ALP-LIFT-09's exact Bargmann reading gives \[
 \vartheta_{\phi_{i,n}}(1,1)
   =C^2\,P_{i,n}\bigl(\sigma_q(u_1)\sigma_q(u_2)\bigr)_{0}
   =C^2a_{i,n}.
 \tag{F3}
\] Indeed the second factor contributes exactly \(\partial_1\partial_3\sigma_q(0)=1\). It is necessary to include this spectator derivative: an undifferentiated second trace-one section vanishes at the origin and would give zero.

The polynomial has one fixed multidegree across CM embeddings: its degree vector is \((n+1,0,1)\) if \(i=1\), and \((1,n,1)\) if \(i=2\). Thus it is one torus weight, including the actual vacuum shift, despite using two different eigenlines on the spectator factor.

\subsection{2. One actual source character, without a growing finite level}\label{one-actual-source-character-without-a-growing-finite-level}

Let \(\beta_{i,n}\) be that actual infinity torus character, and let \(\lambda_f\) be the character on the fixed finite vacuum line. ALP-LIFT-09 proves \[
 T(\mathbf A_{F,f})=T(F)K_0,\qquad
 T(F)\cap K_0=\mu_{14}.
 \tag{F4}
\] For every \(a_{i,n}\ne0\), rational invariance of the actual nonzero sum (F3) gives \[
           \beta_{i,n}(\zeta)\lambda_f(\zeta)=1
                         \quad(\zeta\in\mu_{14}).
 \tag{F5}
\] This is a consequence of the actual theta sum; no compatibility is assumed from a hypothetical lift.

Consequently \[
 \chi_{i,n}(\delta\,t_\infty k)
        =\beta_{i,n}(t_\infty)\lambda_f(k),
 \quad\delta\in T(F),\quad k\in K_0,
 \tag{F6}
\] defines a unitary automorphic character. Independence of the decomposition follows from (F5). The finite compact character is fixed as \(n\) varies. In the standard splitting its only nontrivial compact factor is the orientation character \(\epsilon\) at7, by AM-DO-020--021; the proof using the actual \(\lambda_f\) does not require this optional explicit description.

Normalize the compact torus quotient to volume one. The equality of actual automorphic functions is \[
 \Theta(\phi_{i,n},\chi_{i,n})(g)
       =\int_{[T]}\vartheta_{\phi_{i,n}}(g,t)
                                 \overline{\chi_{i,n}(t)}\,dt
       =\vartheta_{\phi_{i,n}}(g,1).
 \tag{F7}
\] It follows because the theta function transforms by \(\beta_{i,n}\lambda_f\), the group actions commute, and (F4) represents every finite component. Hence \[
 \|\chi_{i,n}\|_2=\|\chi_{i,n}\|_\infty=1,\qquad
 C^{-2}\Theta(\phi_{i,n},\chi_{i,n})(1)=a_{i,n}.
 \tag{F8}
\] There is no finite Fourier projector and no factor \(h_M^{\pm1/2}\) in (F7)--(F8). The algebraic normalized value in (F8) is what will be assigned a \(p\)-adic valuation. No \(p\)-adic valuation is assigned to the unnormalized complex period \(C^2\).

As a direct root-unit check, in the fixed Hom convention \(X\mapsto Xt^{-1}\), the source vacuum shift is \(-1\) at each CM place. Therefore the exponent of \(\beta_{i,n}\) on \(\zeta_7\) is \(-in-10\). Both witness sequences below have \(in\equiv4\bmod7\). Their \(n\) are even, so \(\beta_{i,n}(-1)=-1\), matching the finite orientation value \(\lambda_f(-1)=-1\). This also verifies (F5) explicitly for these sequences using the corrected AM-DO-023 source shift. Formula (F5) remains the coordinate-independent proof.

\subsection{3. Exact sharp actual first values}\label{exact-sharp-actual-first-values}

AM-K020, jacobian\_sharp\_small\_prime\_directions.md, equations (K21.1)--(K21.1a), supplies the nonzero derivatives \[
 n_3(r)=28\cdot3^r+18,\qquad
 n_5(r)=126\cdot5^r+30
 \tag{F9}
\] and their exact valuations. Substitution into (F8) gives the following actual first-image coefficients: \[
 \begin{split}
 v_3\!\left(C^{-2}
       \Theta(\phi_{1,n_3(r)},\chi_{1,n_3(r)})(1)\right)
   &=\frac{3^r+3-s_3(n_3(r))}{2},\\
 v_5\!\left(C^{-2}
       \Theta(\phi_{2,n_5(r)},\chi_{2,n_5(r)})(1)\right)
   &=\frac{5^r-1}{4}.
 \end{split}
 \tag{F10}
\] The new assertion in (F10) is their occurrence as these specified actual theta-lift values. The numerical derivative theorem itself belongs to AM-K020.

In particular the same first-image family, with fixed finite vacuum and unit-norm source character, has \[
 \liminf_{n\to\infty}
   \frac{v_3(C^{-2}\Theta(\phi_{1,n},\chi_{1,n})(1))}{n}
       =\frac1{56},\qquad
 \liminf_{n\to\infty}
   \frac{v_5(C^{-2}\Theta(\phi_{2,n},\chi_{2,n})(1))}{n}
       =\frac1{504},
 \tag{F11}
\] where the family is indexed by the nonzero coefficients; zero coefficients may be included with valuation \(+\infty\). For a zero coefficient not satisfying (F5), no extra character is being asserted by the definition of this family. The global lower bounds and the witness subsequences in AM-K020 give both directions of (F11). Replacing the denominator \(n\) by the actual total creation degree \(n+2\) leaves these limits unchanged.

At every fixed permitted level \(M=29^e\), take the untranslated section \(s_M=[M]^*s_q\) on each factor. ALP-LIFT-10, (O6)--(O11), gives exactly \[
 M^{-(n+2)}P_{i,n}
       \bigl(\sigma_q(Mu_1)\sigma_q(Mu_2)\bigr)_{0}
                         =a_{i,n}.
 \tag{F12}
\] The point \(t'=0\) is in the permitted half-orbit and its finite phase is1, so the finite vector in that identity is the unchanged \(\phi_{0,f}^{\otimes2}\). The level can be fixed once and for all; no interpolation condition \(M\asymp n\) is required for these particular actual rows.

\subsection{4. Costs that have and have not been determined}\label{costs-that-have-and-have-not-been-determined}

The source input in (F7) has norm one and the finite Schwartz norm is one. The exact archimedean Schwartz norm is not one. ALP-LIFT-09, (V23), gives \[
 \|\phi_{i,n,\infty}\|_2^2
   =\det(2Y)^{-1}\,
           n!\left(\frac{\pi}{G_i}\right)^n
                    \frac{\pi^2}{G_1G_3}.
 \tag{F13}
\] The two spectator creation operators lie in distinct orthogonal modes and have factorial1. The first factor contributes \(n!\); the two ordinary Gaussian vacua contribute \(\det(2Y)^{-1}\). Thus \[
 \log\|\phi_{i,n,\infty}\|_2
   =\frac12\log(n!)
      +\frac n2\log\frac{\pi}{G_i}
      -\frac12\log\det(2Y)
      +\log\pi-\frac12\log(G_1G_3).
 \tag{F14}
\] Equations (F13)--(F14), the actual period \(C^2\), and the change in infinity coefficient weight must remain in any comparison of globally normalized rows. A constant finite input cost is not a zero total DP44 cost.

ALP-LIFT-11's fixed bilinear determinant repair preserves all values in (F8) at the identity and all scalar automorphic norms. ALP-LIFT-12 then shows that each repaired nonzero first vector survives some actual \(U(2)\to U(4)\) theta kernel. That theorem does not preserve the exact valuations (F10) at the second output.

The selected later harmonic row, its integral lattice, and the source DP44 gain and error fields therefore remain separate. This result supplies actual fixed-finite first rows on ordinary \(J^2\); it does not identify a DP65 conjugated-action product with those ordinary factors.

\subsection{Dependencies and checks}\label{dependencies-and-checks}

AM-K020, equations(K21.1)--(K21.1a), is the exact new derivative input; it is cited rather than reproved here. ALP-LIFT-09, (V16), (V22)--(V32), provides the actual oscillator, period and class-one character calculation. ALP-LIFT-10, (O6)--(O12), provides the exact permitted pullback normalization. AM-DO-020--021 supplies the optional explicit compact character and corrected DO023 the source vacuum shift. AM-K018 supplies the ordinary-power source-gauge identification.

This specialization is a complete proof from those inputs. Whole-proof checks are recorded separately with their actual scope; no Lean execution or second-output saturation is claimed.

\clearpage

\section{AM-DO-028: the even-rank Hom/tensor and doubling-character dictionary}\label{am-do-028-the-even-rank-homtensor-and-doubling-character-dictionary}

Author: GPT-6 Astra Pro, ALP\_MACHINE \nolinkurl{/root/definite_oscillator}; inherited reasoning effort not separately exposed; 21 November 2026 UTC.

This result fixes the character ratio in the Rallis factor for the literal bilinear source theta operator. It uses the actual Hom action and the determinant-root convention of AM-DO-022--026. It applies to even-rank edges with splitting characters trivial on the base-field ideles. It does not claim a new general local theta theorem, a full global chain, or an arithmetic estimate for the Schwartz functions obtained by nonvanishing theorems.

\subsection{1. The tensor model includes a dual source space}\label{the-tensor-model-includes-a-dual-source-space}

Write \(E/F\) for the CM extension, including the common host \(E'=\mathbf Q(\zeta_{28})\) when that host is used. Let \(S=E^a,T=E^b\) have their identity Hermitian forms, with both \(a,b\) even. Use the actual trace-zero coefficient \(\xi\) of the symplectic form; on the common host this is \(\xi'=\xi_J/2\). The underlying symplectic action and form are \[
 X\longmapsto hXg^{-1},\qquad
 B(X,Y)=-\operatorname{Tr}_{E/F}
                   (\xi\operatorname{tr}(Y^*X)).
 \tag{DO28.1}
\] Composition with \(\operatorname{Tr}_{F/\mathbf Q}\) gives the rational form used previously. The additive character and self-dual measures are transported with this form.

Put \(V=T\) and let \(U=S^\vee\) have the skew-Hermitian form \(-\xi\) times the dual identity Hermitian form, linear in the first variable. The map \[
 V\otimes_EU\longrightarrow\operatorname{Hom}_E(S,T),
                  v\otimes\ell\longmapsto v\ell
 \tag{DO28.2}
\] is a symplectic isomorphism. Indeed, on pure tensors the product of the two forms is \(-\xi(v'^*v)(\ell'^*\ell)\), whose trace is exactly (DO28.1) after expanding the matrix trace. The source group identification is \[
 \alpha:U(S)\longrightarrow U(U),\qquad
                            \alpha(g)=g^{-t}.
 \tag{DO28.3}
\] Thus the tensor action \(h\otimes\alpha(g)\) is exactly the Hom action. This is an actual representation isomorphism, with the same Heisenberg form, rather than an identification of two abstract groups followed by an unspecified oscillator twist.

For a character \(\chi\) trivial on \(\mathbb A_F^\times\), write \[
 \nu_\chi(z/\bar z)=\chi(z).
 \tag{DO28.4}
\] The standard tensor splitting changes on \(U(U)\) by \(\nu_{\Delta\chi_V}(\det g_U)\) and on \(U(V)\) by \(\nu_{\Delta\chi_U}(\det h)\). This is the usual Kudla/Harris--Kudla--Sweet character-change convention. GQT equation (2.2), printed p.13, fixes the positive source sign already in the zero-dimensional opposite-member case. Gan (2022), section2.5, printed pp.14--15, states the two character changes explicitly with the Hilbert-90 inverse.

Since \(\det\alpha(g)=\det(g)^{-1}\), pulling back this splitting through (DO28.3) gives the source convention \[
 \nu_{\Delta\chi_T}(\det g)^{-1}
                         \nu_{\Delta\chi_S}(\det h)
 \tag{DO28.5}
\] when and only when the tensor labels are \[
                     \boxed{\chi_U=\chi_S,\qquad\chi_V=\chi_T.}
 \tag{DO28.6}
\] The determinant roots, rather than the determinants themselves, are evaluated in (DO28.4). Functorial transport of the fixed standard splitting through the symplectic isomorphism fixes the baseline as well as the character changes; no free character is added. In the common-host construction the baseline is also the explicitly checked trivial-character model of DO022.

This does not contradict the corrected Konno--Konno dictionary in DO023. Their matrix source argument there was kept as \(g\), not replaced by \(g^{-t}\); its character label is therefore \(\chi_T^{-1}\). Changing that label while omitting (DO28.3) changes the representation.

\subsection{2. The input is conjugated and transported}\label{the-input-is-conjugated-and-transported}

Let \(F\) be an actual automorphic vector in a unitary representation \(\sigma\) of \(U(S)(\mathbb A_F)\). The source operator is \[
        T_\phi F(h)=\int_{[U(S)]}\theta_\phi(g,h)F(g)\,dg.
 \tag{DO28.7}
\] GQT section11.2, printed pp.51--52, uses the integral against the conjugated input. The bar was checked in the rendered PDF, since text extraction omits it. Define the input in tensor coordinates by \[
         f(u)=\overline{F(\alpha^{-1}u)},\qquad
         \pi=\overline\sigma\circ\alpha^{-1}.
 \tag{DO28.8}
\] With Haar measure pushed forward under the group isomorphism, the GQT theta integral of \(f\) is exactly (DO28.7), for the transported actual Schwartz vector. There is no unspecified volume scalar. For Tamagawa measures the same conclusion follows from invariance under the rational group isomorphism.

At the level of standard base-change parameters, contragredience from complex conjugation and contragredience from (DO28.3) cancel. More precisely, at an unramified place the automorphism \(g\mapsto g^{-t}\) sends the Satake matrix to its inverse transpose, while \(\overline\sigma=\sigma^\vee\) for the unitary representation. Hence \[
                         \phi_\pi=\phi_\sigma.
 \tag{DO28.9}
\] The same statement holds whenever the usual local Langlands parameter is available and its contragredient functoriality is used. The unramified statement suffices to identify the partial global standard Euler product. At a compact real place the two operations likewise return the original highest weight.

This is why the actual selected input in DO025 retains its old parameter block under a same-character lift. Its Hom oscillator component is the dual of the input; the tensor source component is the input after the simultaneous change (DO28.3), (DO28.8).

\subsection{3. An exact doubling-integral twist identity}\label{an-exact-doubling-integral-twist-identity}

This part can be checked directly in GQT, without guessing the meaning of a printed multiplication sign in an \(L\)-factor. Let \(\chi\) be a unitary character of \(E^\times\backslash\mathbb A_E^\times\), trivial on \(\mathbb A_F^\times\). On the doubled unitary group \(G(U\oplus U^-)\), define \[
                         d_\chi=\nu_\chi\circ\det.
 \tag{DO28.10}
\] For its Siegel Levi element \(m(A)\), the unitary determinant is \(\det A/\overline{\det A}\). Therefore \[
                       d_\chi(m(A))=\chi(\det A).
 \tag{DO28.11}
\] Multiplication by \(d_\chi\) gives an exact isomorphism between the normalized induced representations \[
 I(s,1)\longrightarrow I(s,\chi),\qquad
                        \Phi\longmapsto\Phi_\chi=d_\chi\Phi.
 \tag{DO28.12}
\] The normalization by the parabolic modulus is unchanged, and the multiplier is independent of \(s\). Standard, holomorphic, and good sections are transported accordingly. For good sections defined using the standard intertwiner, this also follows directly: \(d_\chi\) is trivial on the unipotent radical, so the intertwiner changes only by its fixed nonzero value at the chosen Weyl representative. No zero or pole in \(s\) is introduced.

On the embedded pair \((g,1)\), the multiplier is \(\nu_\chi(\det g)\). In GQT section11.6, printed p.54, the local zeta integral is the integral of this section against the \textbf{conjugated} matrix coefficient of \(\pi\). That bar was also checked in the rendered PDF. Consequently \[
 \begin{split}
 Z_\chi(s,\Phi_\chi,f_1,f_2;\pi)
 &=\int\Phi(g,1)\nu_\chi(\det g)
                    \overline{\langle\pi(g)f_1,f_2\rangle}\,dg\\
 &=Z_1(s,\Phi,f_1,f_2;
                       \pi\otimes(\nu_\chi\circ\det)^{-1}).
 \end{split}
 \tag{DO28.13}
\] This holds first in the common region of absolute convergence, and elsewhere by meromorphic continuation. It is an equality of the actual local integrals with unchanged Haar measure. The same identity holds for the global doubling integral: the automorphic character \(d_\chi\) is trivial on rational points, so it factors through the Eisenstein sum, and its two group factors cancel against the corresponding input twists. Here there is a necessary distinction between \(\chi\) and \(\nu_\chi\): for \(t\bar t=1\), \(\chi(t)=\nu_\chi(t)^2\), since \(t/\bar t=t^2\). Thus the actual factor in GQT's global integral on printed p.53 is \[
 d_\chi(g_1)d_\chi(g_2)\chi(\det g_2)^{-1}
              =d_\chi(g_1)d_\chi(g_2)^{-1}.
 \tag{DO28.13a}
\] Its product with \(\overline{f_1(g_1)}f_2(g_2)\) is \(\overline{(f_1d_\chi^{-1})(g_1)}
       (f_2d_\chi^{-1})(g_2)\), exactly as required for the representation in (DO28.14). This also checks the determinant-root normalization in the global identity.

GQT defines its local standard factor from these doubling integrals and cites the normalized GCD statement of Yamana in section11.6. Thus, in the even-character case here, its factor is the untwisted standard factor of \[
                   \pi_0=\pi\otimes(\nu_\chi\circ\det)^{-1}.
 \tag{DO28.14}
\] At nonarchimedean places the equality of good-section ideals in (DO28.13) gives the same normalized reciprocal polynomial factor; its constant-term-one normalization removes a Laurent unit ambiguity. The archimedean standard gamma-factor interpretation gives the same twist when expressed through the standard parameter. In particular, at every unramified place its standard Euler parameter is \[
                              \phi_\pi\otimes\chi^{-1}.
 \tag{DO28.15}
\] No claim that an arbitrary nonstandard \(L\)-factor convention uses the same printed symbol is necessary. Equations (DO28.12)--(DO28.15) fix the convention by the actual integrals.

\subsection{4. The factor in the source's labels}\label{the-factor-in-the-sources-labels}

Apply section3 with \(\chi=\chi_V=\chi_T\), and use (DO28.9). The Rallis factor for (DO28.7), expressed in standard base-change notation, is \[
 \boxed{L_E(s,\phi_\sigma\otimes\chi_T^{-1})
        =L_E(s,\phi_\sigma^\vee\otimes\chi_T).}
 \tag{DO28.16}
\] The equality follows from conjugation of \(E/F\): the standard parameter is conjugate self-dual, and \(\chi_T^c=\chi_T^{-1}\) because its restriction to the base field is trivial. Conjugating all places of \(E\) permutes the Euler factors. At a split place of \(F\), this permutes its two \(E\)-places, so the assertion is for the paired local standard factor, not for a single arbitrarily chosen component in isolation.

There is a useful check directly in the source model. Replacing \(\chi_T\) by \(\chi_T\delta\) multiplies the source kernel by \(\nu_\delta(\det g)^{-1}\). Replacing the input by \(F(g)\nu_\delta(\det g)\) makes the integrand unchanged. The parameter in (DO28.16) is invariant under this simultaneous change. A proposed parameter \(\phi_\sigma^\vee\chi_T^{-1}\) would fail this exact identity and is therefore incorrect for the fixed source convention.

For an already established standard parameter block \(\mu_r\otimes S_d\) in an input to edge \(j\), the corresponding character in (DO28.16) is \[
                              \mu_r\mu_j^{-1},
 \tag{DO28.17}
\] or its inverse after the simultaneous conjugation used in the right side of (DO28.16). Here \(S_d\) denotes the actual length-\(d\) segment with whichever standard/Arthur parameter description has been proved for that local lift. This calculation does not supply that local theta theorem. When the segment has the usual unramified exponents \((d-1)/2,(d-3)/2,\ldots,-(d-1)/2\), evaluating at \(s=(d+1)/2\) gives the Hecke factors at \(1,2,\ldots,d\), with the character (DO28.17).

For the selected DO025 exponents, \(\mu_j/\mu_r\) has infinity exponent \(-e_j+e_r\). If \(j>r\), the proved strict increase \(e_j-e_r\ge(j-r)(d+2)>0\) shows that this is a nontrivial angular character. This verifies the sign and nontriviality needed when applying a separate Hecke \(L(1)\)-nonvanishing theorem. It is not a proof of that theorem or a quantitative lower bound for its value.

\subsection{5. Exact scope and sources}\label{exact-scope-and-sources}

The evenness restriction is essential to the simple global character \(d_\chi\) in (DO28.10): this argument assumes \(\chi|\mathbb A_F^\times=1\). It does not silently extend to an odd opposite rank, where a fixed reference character must first be divided out. It preserves the actual host, polarization, additive character, determinant roots, Haar measures, and bilinear input convention. A Rallis nonvanishing theorem still requires its local occurrence, global range, and cuspidality hypotheses. It provides no selected arithmetic row or bound on the later finite Schwartz input.

Primary locators used here:

\begin{itemize}
\tightlist
\item
  Gan, W. T., Qiu, Y., \& Takeda, S. (2014). \emph{The regularized Siegel--Weil formula (the second term identity) and the Rallis inner product formula}. Inventiones Mathematicae, 198, 739--831. The supplied arXiv version1207.4709, section2.8 and equation(2.2), printed pp.12--13; section11.2, pp.51--52; section11.5 and section11.6, pp.53--54; Theorem11.4, pp.54--55. The PDF was inspected for the bars omitted by its text extraction.
\item
  Gan, W. T. (2022). \emph{Automorphic forms and the theta correspondence}. Arizona Winter School lecture notes, section2.5, pp.14--15, character changes; sections2.6 and3.8, pp.15--16 and26, input/quotient convention; Theorem2.6, p.18, and Theorem3.3, p.27, the rank-one inverse-character factor. \url{https://swc-math.github.io/aws/2022/2022GanNotes.pdf}
\item
  AM-DO-022, DO023 as corrected, DO025, and DO026 for the actual Hom splitting and source integral. GQT's published results are used at their stated hypotheses, not re-proved.
\end{itemize}

The full source-to-doubling twist calculation above is new work in this run; its whole-proof check is recorded separately.

\clearpage

\section{ALP-LIFT-14 --- an actual selected even-rank chain with unrestricted later finite inputs}\label{alp-lift-14-an-actual-selected-even-rank-chain-with-unrestricted-later-finite-inputs}

Author: GPT-6 Astra Pro, ALP\_MACHINE /root/lift\_accounting, 21 November 2026 UTC. The inherited reasoning effort is not separately exposed.

This theorem starts from the actual even-rank first edge AM-R-012 and continues through the literal selected ranks and compact Schur types. It uses the common mixed CM host \(E'=\mathbf Q(\zeta_{28})\), the same character on both members of each edge, and the source's bilinear theta integral. Its later finite Schwartz vectors are unrestricted. Thus it proves actual nonzero iterated theta images in these geometric groups; it does not instantiate the six comparison identities, the old elliptic arithmetic row, or their finite contents and global sizes. In particular it does not prove a shorter chain or a DP44 improvement.

The accepted rank and weight recurrences are inputs. The new argument is the character construction, exact partial-standard-factor induction and boundary nonvanishing, combined with local occurrence and a multiplicity-safe global induction. Literature theta and Rallis theorems are cited at their actual hypotheses, not re-proved or counted as new findings.

\subsection{1. Exact statement and the actual common host}\label{exact-statement-and-the-actual-common-host}

Use the host, product polarization, additive character, trace-one vacuum and Hom oscillator of AM-DO-022: \[
 E'=\mathbf Q(\zeta_{28}),\quad F'=(E')^+,\quad
 A=J^2,\quad J=\operatorname{Jac}(y^2=x^7-1),
 \qquad G_j=U(I_{N_j}).
 \tag{E1}
\] Here the rational trace coefficient is \(\xi'=(\zeta_7-\zeta_7^{-1})/14\). The identity Hermitian spaces are positive at every real place. The extension \(E'/F'\) is unramified at every finite place, and all its dyadic places split. These actual host facts, including the integral vacuum normalization, are carried from DP57 and AM-DO-022.

Take an actual source Ranks record and HCParam for which the partitions of AM-DO-025 are defined. In particular \[
 d=t+1,\quad N_j=k+jd\quad(0\le j\le m),\qquad
 k,m,d\text{ even},\quad m\ge2,\quad d\ge k+m+5.
 \tag{E2}
\] Write \(w_{j,\tau}\) for the selected highest weights, and \(\rho_j=\boxtimes_{\tau\in\Phi'}\rho_{w_{j,\tau}}\), where \[
 \Phi'=\{1,3,9,15,17,23\}.
\] Use the edge indexing \(j=0,\ldots,m-1\). The common character on edge \(j\) is denoted \(\mu_j\); thus it is the character called \(\mu_{j+1}\) in the one-based edge notation of AM-DO-025. The exact carried recurrences are \[
 \begin{split}
 \nu_{j,\tau}
   &=w_{j,\tau}+e_{j,\tau}\mathbf1
                          -N_{j+1}\mathbf1/2,\\
 w_{j+1,\tau}
   &=(w_{j,\tau}-d/2,\,
                     (-e_{j,\tau}+N_j/2)^d),\\
 e_{j+1,\tau}-e_{j,\tau}&\ge d+2.
 \end{split}
 \tag{E3}
\] Each \(\nu_{j,\tau}\) is a nonnegative partition. The exact selected polynomial component of the oscillator on edge \(j\) is \[
                    \rho_j^\vee\boxtimes\rho_{j+1}.
 \tag{E4}
\] It has polynomial degrees \(|\nu_{j,\tau}|\). Neither its degree nor its partition is replaced by an unrestricted compact type.

\textbf{Theorem.} There are globally admissible common characters \(\mu_j\) with infinity exponents \(-e_{j,\tau}\), all with finite compact conductor dividing a fixed split place above29, and there are actual nonzero automorphic forms \(F_j\) of compact type \(\rho_j\), \(1\le j\le m\), such that \(F_1\) is the first lift of AM-R-012 and \[
 F_{j+1}(h)=
       \int_{[G_j]}\vartheta_{\phi_j}(g,h)F_j(g)\,dg
                        \ne0\qquad(1\le j<m).
 \tag{E5}
\] At infinity \(\phi_j\) lies in the exact component (E4); at finite places it can be chosen to be a finite linear combination of genuine Schwartz vectors. No bound for that combination's finite level, coefficients, or norm is claimed.

One may choose the induction so that the representation generated by \(F_j\) is spherical at all nonsplit finite places and at every split place outside29. More strongly, every nonzero vector in the resulting generated automorphic module survives some next finite Schwartz input in the selected real component. This last assertion concerns actual automorphic embeddings, including their multiplicities.

\subsection{2. Global characters and strict nontrivial ratios}\label{global-characters-and-strict-nontrivial-ratios}

Let \(T'\) be the norm-one torus. AM-R-009 proves, using the class number one of \(E'\), the actual factorization \[
 T'(\mathbb A_f)=T'(F')K_{T,0},\qquad
 T'(F')\cap K_{T,0}=\mu_{28}.
 \tag{E6}
\] Its ideal proof and the cited Masley class-number statement were independently checked in review\_root\_unramified\_first\_lift.md. We use that result, not a new class-number assertion.

For edge \(j\) put \(\beta_j(t)=\prod_{\tau\in\Phi'}t_\tau^{-e_{j,\tau}}\) on \(T'(\mathbb R)\). At one split29 factor choose the residue embedding \(\zeta_{28}\mapsto2\); 2 has order28 in \(\mathbf F_{29}^{\times}\). Define a character of that residue group by \[
 \varrho_j(2)=
 \exp\!\left(\frac{2\pi i}{28}
                    \sum_{\tau\in\Phi'}e_{j,\tau}\tau\right),
 \tag{E7}
\] and take it trivial on all other compact factors. Its restriction to \(\mu_{28}\) is \(\beta_j^{-1}\). Therefore (E6) gives a unitary automorphic torus character \(\gamma_j\) with infinity component \(\beta_j\) and compact component \(\varrho_j\): to define it on an adelic point, factor its finite part as \(qk\), move the rational \(q\) to the infinite component, and evaluate \(\beta_j\varrho_j\). Changing the factorization changes the two components by one element of \(\mu_{28}\), on which their product is one.

Set \[
 \mu_j(x)=\gamma_j(x/\bar x),\qquad
             \chi_{S,j}=\chi_{T,j}=\mu_j.
 \tag{E8}
\] This character is trivial on \(\mathbb A_{F'}^\times\), has \(\mu_{j,\tau}(z)=(z/\bar z)^{-e_{j,\tau}}\), and satisfies the even-rank splitting restrictions. At a nonsplit finite place the entire local norm-one torus is compact: its elements have valuation zero. Consequently \(\gamma_{j,v}\), and hence \(\mu_{j,v}\), is identically one there. At split places other than29 it is unramified. Its compact finite values are roots of unity of order dividing28.

For \(r<j\), define \[
                           \upsilon_{rj}=\mu_r\mu_j^{-1}.
 \tag{E9}
\] Its infinity exponent at every \(\tau\) is \(e_{j,\tau}-e_{r,\tau}>0\). Thus both \(\upsilon_{rj}\) and \(\upsilon_{rj}^2\) have nontrivial angular infinity type. In particular neither is a norm character. This strict condition will exclude the pole and the possible boundary-zero obstruction in section5. It follows from the actual recurrence (E3); it is not obtained by changing source exponents.

\subsection{3. Local theta occurrence and the unramified parameter recurrence}\label{local-theta-occurrence-and-the-unramified-parameter-recurrence}

Let \(a=N_j\), \(b=a+d\). At a split finite place, the groups are \(\mathrm{GL}_a\) and \(\mathrm{GL}_b\). Mínguez (2008), Theorem1(1)--(2), printed p.718, proves nonzero small theta for every irreducible input when \(a\le b\), including dyadic residual characteristic. We use its actual type-II representation, with the Fourier transform between the two matrix polarizations; equation(1.4), printed pp.721--722, fixes the half-modulus factors.

For precision, write the local common character as \(\mu=(\eta,\eta^{-1})\) on \(E'_v=F'_v\times F'_v\). The Hom oscillator changes by \[
                     \eta_T(\det g)^{-1}\eta_S(\det h).
 \tag{E10}
\] The literal bilinear integral pairs the source oscillator with the contragredient of the actual input \(\pi\). After removing (E10), the type-II source representation is \(\pi^\vee\otimes\eta_T\). Mínguez's target is its dual together with the standard rank-\(d\) trivial-representation segment, and restoring the target twist gives \[
 \phi_{\rm out}
       =(\phi_\pi\otimes\eta_S\eta_T^{-1})\,
                       \sqcup\,\eta_S[d].
 \tag{E11}
\] Thus the old block is unchanged when \(\eta_S=\eta_T\). This checks the sign by the actual Hom action, not by an unspecified theta-convention formula.

At nonsplit finite places the residue characteristic is odd and all \(\mu_{j,v}=1\). The even identity Hermitian lattices are integrally hyperbolic. Indeed the norm map on units of an unramified quadratic extension is surjective, so the binary identity form contains a primitive isotropic vector and splits a unimodular hyperbolic plane; iteration gives the assertion. The additive character used in the actual host may have conductor inherited from the different of \(F'/\mathbf Q\). This does not change the unramified comparison: write the actual pure-imaginary trace coefficient as \(c\delta\), with \(\delta\) a pure-imaginary unit and \(c\in F_v'^\times\). The absolute-codifferent condition and the trivial relative different imply that \(\psi_{F'}(c\,\cdot)\) has integral conductor. Rescaling the skew form and additive character together preserves their composite symplectic character, the lattice model and the unitary groups. This is the integral model already fixed in AM-DO-022.

Apply Liu (2021), Definition2.8 and Theorem6.2(1)--(2), printed pp.7 and21--22, with sign \(+\), source rank \(2r=a\), target rank \(2s=b\), and trivial splitting character. Proposition3.4, printed p.9, verifies its Assumption3.3 in odd residue characteristic. For every spherical source representation, the target small theta is the spherical representation stated there, and the characteristic vector of the selfdual lattice has nonzero image. No temperedness of the source is required. The extra standard eigenvalues are \(q_F^{\pm1},q_F^{\pm3},\ldots,q_F^{\pm(d-1)}\). The same powers appear in Liu (2011), Appendix, printed p.916, but we use the stronger 2021 theorem to ensure that every specialization survives; an annihilator identity alone would not prove that.

Write \(Q_w=\#(\mathcal O_{E'}/w)\). At an unramified place let \(\mathcal A_{j,w}\) be the standard Satake multiset of an actual constituent \(\pi_j\). We use the convention that an inducing character \(\alpha|\cdot|^u\) has eigenvalue \(\alpha(\varpi_w)Q_w^{-u}\). The two local calculations give exactly \[
 \mathcal A_{j+1,w}
   =\mathcal A_{j,w}\ \sqcup\
       \left\{\mu_{j,w}(\varpi_w)
              Q_w^{(d-1)/2-\ell}:0\le\ell<d\right\}.
 \tag{E12}
\] At a nonsplit place \(Q_w=q_F^2\), giving precisely the odd \(q_F\)-powers above; inverse paired Satake parameters absorb the contragredient convention. At a split place (E11) gives each chosen component, with its conjugate component included as well. This is a formula for unramified parameters. We make no claim that the same segment describes compact real Langlands parameters.

For each actual descendant irreducible constituent choose one ancestral initial constituent \(\pi_0\) in the first input. Iteration along that ancestry gives \[
 \mathcal A_{j,w}
  =\mathcal A_{0,w}\ \sqcup\
    \coprod_{0\le r<j}
     \left\{\mu_{r,w}(\varpi_w)
                  Q_w^{(d-1)/2-\ell}:0\le\ell<d\right\}.
 \tag{E13}
\] The starting class-supported vector can contain several initial constituents. Formula(E13) is not an assertion that one common \(\pi_0\) describes the whole generated automorphic module; the unitary bound in (E18) is uniform over all such ancestors.

\subsection{4. Exact Rallis factor in the literal source convention}\label{exact-rallis-factor-in-the-literal-source-convention}

AM-DO-028 is used here with its full Hom/tensor and conjugation dictionary. In tensor coordinates the source is \(S^\vee\) and its group map is \(\alpha(g)=g^{-t}\). GQT's conjugated-input convention is obtained from the actual bilinear input by \(f=\overline{F\circ\alpha^{-1}}\). The two contragrediences cancel on standard unramified parameters. Its tensor character is \(\chi_V=\chi_T\).

There is also a direct check of the twist. For \(\chi|_{\mathbb A_{F'}^\times}=1\), define \(\nu_\chi(z/\bar z)=\chi(z)\). On the doubled unitary group multiplication by \(d_\chi=\nu_\chi\circ\det\) identifies \(I(s,1)\) with \(I(s,\chi)\), because its restriction to the Siegel Levi \(A\) is \(\chi(\det A)\). The local doubling integral uses a conjugated matrix coefficient. Hence \[
 Z_\chi(s,d_\chi\Phi;\pi)
      =Z_1(s,\Phi;\pi\otimes(\nu_\chi\circ\det)^{-1}).
 \tag{E14}
\] The good-section ideals and normalized standard factors transform accordingly. In the global integral the factor on the second group is \(\chi(\det g_2)^{-1}\), and on the norm-one torus \(\chi(t)=\nu_\chi(t)^2\); thus the two doubled determinant factors cancel with the two input twists exactly. This verifies that no square root or sign is suppressed.

Accordingly the standard partial factor in source labels is \[
 \mathcal L_j^S(s)
   =L_{E'}^S(s,\operatorname{BC}(\pi_j)\otimes\mu_j^{-1}).
 \tag{E15}
\] Here BC notation denotes the standard unramified Euler parameters. A global automorphic base-change theorem is not being assumed. Choose a finite set \(S\) containing all ramified places and omit both \(E'\)-places over every omitted split \(F'\)-place. Conjugating \(E'/F'\) gives the equivalent factor \(L^S(s,\operatorname{BC}(\pi_j)^\vee\otimes\mu_j)\). The actual formula has ratios \(\mu_r\mu_j^{-1}\), not products \(\mu_r\mu_j\).

Substituting (E13) gives an identity of Euler products in a right half-plane, and therefore of their continuations: \[
 \mathcal L_j^S(s)=
 L_{E'}^S(s,\operatorname{BC}(\pi_0)\otimes\mu_j^{-1})
 \prod_{r<j}\prod_{\ell=0}^{d-1}
 L_{E'}^S(s-(d-1)/2+\ell,\upsilon_{rj}).
 \tag{E16}
\] The Rallis point for edge \(j\) is \(s_0=d/2\) in the doubled induction and \(s_*=(d+1)/2\) in the standard factor. Therefore \[
 \mathcal L_j^S(s_*)=
 L_{E'}^S(s_*,\operatorname{BC}(\pi_0)\otimes\mu_j^{-1})
                \prod_{r<j}\prod_{\ell=1}^{d}
                      L_{E'}^S(\ell,\upsilon_{rj}).
 \tag{E17}
\]

\subsection{5. The exact nonvanishing needed in (E17)}\label{the-exact-nonvanishing-needed-in-e17}

The initial representation need not be tempered. It is unitary because it occurs in the automorphic Hilbert space of a compact quotient. Its unramified spherical matrix coefficient has absolute value at most one. Macdonald's bounded-spherical-function criterion says that the logarithmic real Satake parameter lies in the convex hull of the Weyl orbit of \(\rho\). In the standard degree-\(k\) representation this gives \[
                         |\alpha_w|\le Q_w^{(k-1)/2}.
 \tag{E18}
\] For a split place this is the familiar \(\mathrm{GL}_k\) bound. For unramified \(U(k)\), the relative \(\rho\)-coordinates are \(k-1,k-3,\ldots,1\) in \(q_F\)-units and \(Q_w=q_F^2\), giving the same bound. The unitary center contributes no extra real exponent. Exact references are Macdonald (1971), ChapterIV, section4.7, Theorem4.7.1; Macdonald (1995), ChapterV, section3, Example, printed pp.299--300 supplies the split specialization.

By (E2), \[
 s_*-(k-1)/2=(d-k+2)/2\ge(m+7)/2>1.
 \tag{E19}
\] Thus the initial Euler product in (E17) is absolutely convergent and nonzero. Indeed its logarithmic series is bounded by a constant times \(k\sum_w Q_w^{-(d-k+2)/2}\), with all higher powers dominated by a convergent geometric series. Every individual factor is finite and nonzero. This argument does not assume a Ramanujan bound for \(\pi_0\).

For \(\ell\ge2\), the Hecke factors in (E17) are absolutely convergent and nonzero because \(\upsilon_{rj}\) is unitary. The boundary \(\ell=1\) requires a separate argument.

\textbf{Boundary lemma.} Let \(\upsilon\) be a unitary Hecke character with nontrivial angular infinity type, such that \(\upsilon^2\) also has nontrivial angular infinity type. Then \(L^S(1,\upsilon)\) is finite and nonzero.

\textbf{Proof.} Hecke--Tate continuation gives holomorphy at1 for both \(L^S(s,\upsilon)\) and \(L^S(s,\upsilon^2)\). No norm-character pole occurs because of their angular types. Omitting finitely many local factors preserves holomorphy there. For real \(\sigma>1\), expand the absolutely convergent Euler logarithms: \[
 \begin{split}
 3\log\zeta_{E'}^S(\sigma)
 +4\log|L^S(\sigma,\upsilon)|
 +\log|L^S(\sigma,\upsilon^2)|
 =\sum_{w\notin S}\sum_{n\ge1}
 \frac{3+4\operatorname{Re}(\upsilon(w)^n)
              +\operatorname{Re}(\upsilon(w)^{2n})}
                     {nQ_w^{n\sigma}}\ge0.
 \end{split}
 \tag{E20}
\] The inequality follows termwise from \(3+4\cos\theta+\cos2\theta=2(1+\cos\theta)^2\). Consequently \[
 \zeta_{E'}^S(\sigma)^3
 |L^S(\sigma,\upsilon)|^4|L^S(\sigma,\upsilon^2)|\ge1.
 \tag{E21}
\] If \(L^S(s,\upsilon)\) had a zero of order \(r\ge1\) at1, the left side would be \(O((\sigma-1)^{4r-3})\), since the zeta function has a simple pole and the last factor is bounded. It would tend to zero, a contradiction. This proves the lemma. The argument asserts neither absolute convergence nor an Euler-product limit at1. ∎

The continuation input is used only for nontrivial angular characters; see Prasad (2007), Lectures on Tate's thesis, section1, Theorem2, printed p.3. The elementary positivity proof above supplies the needed boundary nonvanishing. By (E9) it applies to every \(\upsilon_{rj}\). Thus \[
                   \mathcal L_j^S(s_*)\ne0
 \quad\text{and it is holomorphic at }s_*.
 \tag{E22}
\]

\subsection{6. Rallis hypotheses, including the real local test}\label{rallis-hypotheses-including-the-real-local-test}

For every later edge \(j\ge1\), (E2) gives \[
 a=N_j\ge k+d>d,\qquad
       a<b=a+d<2a,\qquad 0<s_0=d/2<a/2.
 \tag{E23}
\] The target identity Hermitian space is positive at every real place, hence anisotropic over \(F'\): an isotropic rational vector would remain isotropic under a real embedding, contradicting positivity. Its global Witt index is therefore \(r=0\). The groups have no proper rational parabolics, their adelic quotients are compact, and their automorphic representations are cuspidal. In GQT's notation \(\epsilon_0=0\), \(d(n)=n=a\), \(m=b\), and the hypothesis that lower lifts vanish for Witt indices less than \(r\) is vacuous. Although the numerical dimensions lie in the interval treated by the second-term theorem, the actual positive target with \(r=0\) is in its convergent case. GQT Proposition11.1(ii) and Theorem11.4 explicitly include this case.

At a nonsplit finite place, the spherical local occurrence was proved in section3. At a split finite place, including the possibly ramified representation at29 and all dyadic places, Mínguez gives occurrence for every irreducible input. GQT Proposition11.6(i) therefore gives a nonzero normalized local doubling test on the actual target's Siegel--Weil image at every finite place.

At a real place we do not invoke the incomplete converse discussed in GQT Theorem11.7. The actual selected Fock summand is (E4). Choose a unit vector in the target factor and an isometric copy of the source contragredient in that summand. The doubling isomorphism satisfies \(\delta(\phi_1\otimes\overline{\phi_2})(0)=\langle\phi_1,\phi_2\rangle\), with the determinant factor retained in (E14). Choose the source vector and this Fock vector so that the oscillator coefficient is the conjugate of the input coefficient in the GQT integral. Schur orthogonality then gives \[
 Z_\tau(s_0)=
 \int_{U(a)}|\langle\rho_j(g)v,v\rangle|^2\,dg
                 =\frac{\operatorname{vol}(U(a))}
                         {\dim\rho_{j,\tau}}>0.
 \tag{E24}
\] This is the same direct compact calculation proved in ALP-LIFT-12, now with an arbitrary selected Schur representation in place of a symmetric power. The standard real gamma factors are finite and nonzero at \(s_*>0\): a compact-group parameter is tempered, the character twist is angular and unitary, and its complex gamma arguments are \(s_*\) plus nonnegative half-integer angular sizes. Hence the normalized local test also is nonzero.

The construction uses the selected Schur component, not necessarily the identical highest polynomial at every intermediate evaluation. To pair a given source highest vector, a compact Weyl translate of the polynomial factor may be needed. This preserves its partition, total degree and unitary norm. Identifying one rigid highest-vector pin throughout the source's six identities is an additional comparison, not a consequence of (E24).

All local theta lifts are now nonzero. GQT Theorem11.4(ii) implies holomorphy of the completed standard factor at \(s_*\). Equation(E22) gives a nonzero partial factor. The finitely many missing nonarchimedean standard factors are reciprocal polynomials with constant term one and have no zeros; the real gamma factors have no zeros either. Holomorphy of the completed factor therefore rules out any pole among them at this point and implies \[
                             L(s_*,\pi_j\times\chi_V)\ne0.
 \tag{E25}
\] This deduction does not assume temperedness of the later representations.

The normalized unramified test is \(1/d_v(s_0,\chi_V)\). Its infinite product converges and is nonzero because \(s_0>0\), exactly as in GQT equation(11.3), printed p.54. Combining (E24), the finite local tests, and (E25), GQT Theorem11.4 gives a nonzero global theta lift on every actual irreducible cuspidal embedding satisfying the stated unramified parameters. The use is of the theorem's actual representation and measure, through (E14); it is not an inference from a formal Euler factor alone.

\subsection{7. Actual vectors, multiplicities and retention of the spherical hypothesis}\label{actual-vectors-multiplicities-and-retention-of-the-spherical-hypothesis}

The preceding paragraph is stronger than existence of one abstract automorphic representation with a nonzero lift. We give the vector argument because actual iterated images are the issue.

Let \(\mathcal H_j\) be the closed \(G_j(\mathbb A)\)-invariant automorphic subspace generated by the actual previous theta images with the prescribed real types. Every constituent in it has the unramified parameters (E13) for an ancestral initial constituent, by uniqueness of the local small theta quotient and the local calculations in section3. The actual starting vector has fixed finite level and real type, so it has only finitely many initial constituents. No common initial constituent is required: sections5--6 apply separately to each descendant, with the uniform bound (E18). Fix the real selected component (E4), while allowing all finite Schwartz vectors and all vectors of that component.

Let \(\mathcal N_j\subset\mathcal H_j\) be the joint kernel of all these next theta maps. The automorphic quotients are compact, so the maps defined by each smooth Schwartz kernel are continuous on the relevant Hilbert spaces. The joint kernel is closed. Equivariance of the Weil representation makes it invariant under \(G_j(\mathbb A)\): translating an input is absorbed by translating the finite Schwartz vector and its real selected component. If \(\mathcal N_j\ne0\), discrete spectral decomposition of the compact automorphic quotient supplies an actual irreducible cuspidal embedding contained in \(\mathcal N_j\). The Rallis argument of section6 applies to that embedding, with its inherited positive invariant inner product, and says that some selected theta map is nonzero on it. This contradicts its definition. Therefore \[
                              \mathcal N_j=0.
 \tag{E26}
\] This handles automorphic multiplicity copies rather than implicitly identifying a representation with one preferred embedding.

It follows that every nonzero current vector survives some next Schwartz vector. The latter may be a finite sum of decomposable vectors; a nonzero sum has a nonzero summand after any needed linear choices. Its output has type \(\rho_{j+1}\).

The next induction step needs spherical input at nonsplit places. This can be retained without asserting an unsupported fixed-vacuum global nonvanishing. Local uniqueness and Liu's theorem show that every output constituent is spherical at every nonsplit place; at split places outside29 the unramified type-II theorem gives the same conclusion. For a nonzero smooth output, only finitely many local compact invariance conditions need be restored. At one such place, if all translates of the vector had zero projection to the maximal-compact invariants, its cyclic subrepresentation would have no spherical constituent, a contradiction. Thus a finite translate followed by the compact averaging projector gives a nonzero vector invariant there. Averaging may be combined at the finite set of places. The operations commute with the other local actions, and equivariance absorbs them into the finite Schwartz input. All but finitely many components were already the vacuum, so there is no infinite operation.

The resulting finite input may be a Hecke combination of the original lattice vector. This argument does not say that it remains a single prescribed move, nor that its coefficients are units. It proves exactly the spherical condition required by section3 and no finite arithmetic estimate.

\subsection{8. Initialization and the resulting chain}\label{initialization-and-the-resulting-chain}

AM-R-012 supplies the actual starting vector, using the accepted DP59 half-orbit and the new interpolation result AM-K022. It does not use a rank-one precursor. For \(a=k\), \(b=k+d\), \(D=\sum_\tau|\nu_{0,\tau}|\), choose \[
 M=29^{\max(1,\lceil\log_{29}(3D+2)\rceil)}.
\] There is one permitted translate of finite Schwartz norm one with nonzero normalized selected jet \(a_x\), and \[
 v_{11}(a_x)\le
                  \sum_{\tau,i}v_{11}(\nu_{0,\tau,i}!).
 \tag{E27}
\] The derivative factorial is retained. The exact source input \(f_0\) is supported on the identity finite class at principal level \(M\); its compact component is the primitive highest matrix coefficient. With the actual quotient measure, \[
 \begin{split}
 I_a(M)&=M^{6a^2}\prod_{i=1}^a(1-29^{-i})^6,\\
 \|f_0\|_2^2&=(I_a(M)\dim\rho_0)^{-1},\\
 F_1(1)&=\frac{C^{2ab}a_x}{I_a(M)\dim\rho_0}\ne0.
 \end{split}
 \tag{E28}
\] The period \(C=\theta[0;0](0;\tau)c\) is the actual sigma-normalized period of AM-DO-022, not a newly chosen row scalar. The first vector is invariant under the standard compact outside29. Its generated constituents are therefore spherical there and satisfy (E12) for \(j=0\). Its source representation is unitary, so the crude initial bound (E18) applies.

Starting with this actual \(F_1\), apply (E26) successively for \(j=1,\ldots,m-1\), retaining spherical invariance as in section7. This produces the actual nonzero chain (E5). Since the finite Schwartz space has a basis over an algebraic coefficient extension containing the finite splitting values, a nonzero finite linear functional cannot vanish on all such basis vectors; the finite choices may be taken with algebraic coefficients. This observation gives no bound on their height or denominator.

The concrete HC-only example in AM-R-012 has \(6\to20\), \(D=36\), \(M=841\), \(\dim\rho_0=1\), and \(v_{11}(I_6(M))=0\). Its first algebraic identity coefficient has valuation at most zero. That example does not impose the separate noncompact Hyp3 conditions and is not evidence for a later gain. It is consistent with the present existence theorem, which transports nonzero vectors but does not transport their valuations.

\subsection{9. What remains for Q1}\label{what-remains-for-q1}

This result supplies an actual common-host even-rank chain and removes the earlier unrestricted-global-nonvanishing gate. It retains the exact source selected partitions, literal bilinear integral and equal-character requirement. It does not supply the source's full D7 constructor, its six AdmS identities, or an identification with the prescribed arithmetic rowAF.

For every later edge the missing arithmetic object is the restriction of the class-integrated theta matrix to the already produced input-image lattice, with its selected coefficient normalization and compact contraction included. AM-K022's finite jet matrix has the same determinantal ideals as the corresponding derivative-function matrix before that integration; it does not imply that the restricted class matrix is saturated. The finite local tests supplied by Rallis need not lie in the prescribed DP59/DP62 moved family. Their coefficients, common denominator, finite level, archimedean norm, and content can affect every \(g_j,e_j,H\) term in DP44.

Thus this theorem establishes actual analytic existence with exact local types. It provides no implication from individual derivative slopes, or from \(L(1)\ne0\), to a smaller required number of lifts. A quantitative lower bound for those \(L\)-values would itself be insufficient without the integral image and row comparisons.

\subsection{Primary sources and exact locators}\label{primary-sources-and-exact-locators}

\begin{itemize}
\tightlist
\item
  Gan, W. T., Qiu, Y., \& Takeda, S. (2014). The regularized Siegel--Weil formula (the second term identity) and the Rallis inner product formula. Inventiones Mathematicae, 198, 739--831. arXiv:1207.4709v3, sections11.2--11.8, printed pp.51--56; Proposition11.1(ii), Theorem11.4, Proposition11.6, equation(11.3). The conjugation bars on pp.52 and54 were checked in the rendered PDF. Local copy: lift\_refs/gqt2014.pdf.
\item
  Liu, Y. (2021). Theta correspondence for almost unramified representations of unitary groups. arXiv:2105.14223v1; published Journal of Number Theory, 230 (2022), 196--224. Definition2.8, Proposition3.4, Theorem6.2, printed pp.7,9,21--22. The sign \(+\), arbitrary spherical-parameter specialization is the one used here. Local copy: lift\_refs/liu2021\_unramified\_theta.pdf.
\item
  Liu, Y. (2011). Arithmetic theta lifting and L-derivatives for unitary groups, I. Algebra \& Number Theory, 5(7), 849--921. Appendix, printed pp.912--918, especially the map p.916 and PropositionA.5. This is an independent normalization check, not the sole proof of every spherical specialization. \url{https://msp.org/ant/2011/5-7/ant-v5-n7-p01-p.pdf}.
\item
  Mínguez, A. (2008). Correspondance de Howe explicite: paires duales de type II. Annales Scientifiques de l'École Normale Supérieure, 41(5), 717--741. Theorem1, p.718; equation(1.4), pp.721--722. \url{https://www.numdam.org/item/ASENS_2008_4_41_5_717_0.pdf}.
\item
  Macdonald, I. G. (1971). Spherical functions on a group of p-adic type. Ramanujan Institute Publications2. ChapterIV, section4.7, Theorem4.7.1, bounded-spherical-function criterion. Primary chapter transcribed by Arun Ram: \url{https://math.soimeme.org/~arunram/Resources/SFOAGOPTChapterIVCalculationOfTheSphericalFunctions.html}.
\item
  Macdonald, I. G. (1995). Symmetric functions and Hall polynomials (2nd ed.). Oxford University Press. ChapterV, section3, Example, printed pp.299--300. \url{https://math.berkeley.edu/~corteel/MATH249/macdonald.pdf}.
\item
  Prasad, D. (2007). Lectures on Tate's thesis. ICTP lecture notes, section1, Theorem2, printed p.3, used only for continuation and absence of a pole for a character with nontrivial angular infinity type. \url{https://www.math.iitb.ac.in/~dprasad/ictptate.pdf}.
\item
  AM-DO-022, AM-DO-025--026 and AM-DO-028 for the actual host, compact selected summands, source integral and doubling dictionary; AM-R-009 and AM-R-012 for the class-one character construction and actual first edge; AM-K022 for the exact permitted first jet. Their claims remain separately identified and are not counted again as new results.
\end{itemize}

\clearpage

\section{ALP-LIFT-15 --- one finite Heisenberg translate at every selected edge}\label{alp-lift-15-one-finite-heisenberg-translate-at-every-selected-edge}

Author: GPT-6 Astra Pro, ALP\_MACHINE /root/lift\_accounting, 21 November 2026 UTC. Inherited reasoning effort is not separately exposed.

This refines ALP-LIFT-14's unrestricted finite Schwartz existence. Every later nonzero lift can be obtained with \textbf{one actual finite Heisenberg translate of the normalized vacuum}, with finite \(L^2\) norm exactly one and coefficient vector consisting of a single one. The required finite translation denominator is not bounded. The full orbit here is larger than the prescribed DP59 half-orbit, and the resulting rational half-model function need not equal the source's raw Xi. Those restrictions remain genuine arithmetic obligations.

The proof below gives the complete finite basis calculation, its quantitative norm bound at a specified denominator, and its application to the actual selected chain. It does not invoke Stone--von Neumann irreducibility in place of a finite calculation.

\subsection{1. The actual finite lattice model}\label{the-actual-finite-lattice-model}

On an edge \(U(a)\to U(b)\) in the common \(E'=\mathbf Q(\zeta_{28})\) host, let \(W\) be the underlying rational symplectic space and let \[
             2g=\dim_{\mathbf Q}W=12ab,\qquad g=6ab.
 \tag{H1}
\] The actual lattice is the trace lattice on \(\operatorname{Hom}_{\mathcal O_{E'}}(\mathcal O_{E'}^a,\mathcal O_{E'}^b)\), with the alternating pairing and trace-one semicharacter of AM-DO-022. It is selfdual. Write \(\widehat\Lambda\subset W(\mathbb A_f)\) for its finite completion, and \(\phi_0\) for the actual normalized finite vacuum, so \(\|\phi_0\|_2=1\).

Let \(W_x\) denote the actual finite Weyl operator, with the source's additive character and selfdual measure. Fix the commutator convention \[
                 W_xW_y=\psi_f(B(x,y))W_yW_x.
 \tag{H2}
\] Changing both signs gives the same statements below. The actual vacuum character is denoted \(\alpha\): \[
                         W_\lambda\phi_0
                            =\alpha(\lambda)\phi_0
                 \quad(\lambda\in\widehat\Lambda).
 \tag{H3}
\] Its semicharacter law is exactly the one forced by the Weyl cocycle. The trace-one construction supplies a continuous sign character in the appropriate quadratic sense; its values on the lattice are roots of unity. This is the actual source vacuum, not an arbitrary integral frame for its Chern-reduced analytic germ.

For a positive integer \(N\), define the finite-dimensional eigenspace \[
 \mathcal S_N=
  \{f\in\mathcal S(W_{\rm pos}(\mathbb A_f)):
        W_{N\lambda}f=\alpha(N\lambda)f
                    \text{ for all }\lambda\in\widehat\Lambda\}.
 \tag{H4}
\] Here a rational polarization is chosen only to write the Schrödinger space. All the definitions through (H4) are invariant under a unitary change of model.

\subsection{2. Exact orthonormal basis}\label{exact-orthonormal-basis}

\textbf{Proposition.} Choose one representative for every element of \[
                     N^{-1}\widehat\Lambda/\widehat\Lambda.
\] Then \[
                 \{W_x\phi_0:
                    x\in N^{-1}\widehat\Lambda/\widehat\Lambda\}
 \tag{H5}
\] is an orthonormal basis of \(\mathcal S_N\), after choosing the indicated representatives. Changing a representative changes its vector by a unit scalar. In particular \[
                          \dim\mathcal S_N=N^{2g}.
 \tag{H6}
\] Every finite Schwartz vector belongs to some \(\mathcal S_N\).

\subsubsection{Proof: membership and orthogonality}\label{proof-membership-and-orthogonality}

If \(x\in N^{-1}\widehat\Lambda\) and \(\lambda\in\widehat\Lambda\), selfduality gives \(\psi_f(B(N\lambda,x))=1\). Equations(H2)--(H3) therefore put \(W_x\phi_0\) in \(\mathcal S_N\).

For distinct classes \(x,y\), selfduality of \(\widehat\Lambda\) gives \(\lambda\in\widehat\Lambda\) with \(\psi_f(B(\lambda,x-y))\ne1\). The two vectors \(W_x\phi_0,W_y\phi_0\) are eigenvectors of the unitary operator \(W_\lambda\) with different eigenvalues, by (H2)--(H3). Their inner product is therefore zero. Each has norm one because Weyl operators are unitary. The quotient has order \(N^{2g}\), since the lattice has rank \(2g\). It remains to check that \(\mathcal S_N\) has exactly this dimension and that the union exhausts the Schwartz space.

\subsubsection{Proof: an explicit Schrödinger calculation}\label{proof-an-explicit-schruxf6dinger-calculation}

Choose an integral symplectic basis of the unimodular lattice. Its existence follows by the usual symplectic Gram--Schmidt argument over \(\mathbf Z\): choose a primitive vector, use unimodularity to find a partner of pairing one, split their hyperbolic plane, and iterate on its integral orthogonal complement. In the associated Schrödinger model the reference vacuum is \(\mathbf1_{\widehat{\mathbf Z}^{\,g}}\), of norm one, and its lattice semicharacter is the Weyl half-pairing character \(\alpha_{\rm std}\). The ratio \(\alpha/\alpha_{\rm std}\) is a continuous character of \(\widehat\Lambda\), since the two semicharacters have the same cocycle. The perfect finite-adelic symplectic pairing identifies the character group of this compact lattice with \(W(\mathbb A_f)/\widehat\Lambda\). Thus a finite Weyl translate \(W_t\mathbf1_{\widehat{\mathbf Z}^{\,g}}\) has precisely the character \(\alpha\).

For the trace-one sign character, the ratio has order at most two, so \(t\) may be taken in \(\tfrac12\widehat\Lambda/\widehat\Lambda\); the general argument only needs a finite denominator. The eigenspace of the full lattice for \(\alpha_{\rm std}\) is visibly one-dimensional: invariance under integral translations makes a function constant on every integral coset, while invariance under all integral modulations forces its support into \(\widehat{\mathbf Z}^{\,g}\). Consequently the actual \(\phi_0\) differs from \(W_t\mathbf1_{\widehat{\mathbf Z}^{\,g}}\) by a scalar of absolute value one. This scalar is retained and does not affect the assertions.

Conjugation by \(W_t\) reduces (H4) to the corresponding eigenspace for the standard vacuum. For the pure position and pure momentum elements of \(N\widehat\Lambda\), its conditions say exactly \[
 f(u+Nz)=f(u)\quad(z\in\widehat{\mathbf Z}^{\,g}),\qquad
       \operatorname{supp}f\subset
                  N^{-1}\widehat{\mathbf Z}^{\,g}.
 \tag{H7}
\] The mixed elements impose no additional condition: their Weyl half-pairing factor is exactly \(\alpha_{\rm std}(N\lambda)\). Thus this space is the space of functions on \(N^{-1}\widehat{\mathbf Z}^{\,g}/
                                  N\widehat{\mathbf Z}^{\,g}\). Its dimension is \(N^{2g}\), proving (H6), and hence proving completeness of the orthonormal family (H5).

Every compactly supported locally constant function is invariant under a sufficiently small open additive subgroup and has support in a sufficiently large compact open subgroup. Both can be bounded by \(N\widehat{\mathbf Z}^{\,g}\) and \(N^{-1}\widehat{\mathbf Z}^{\,g}\) for one integer \(N\). Applying the fixed inverse translate \(W_t^{-1}\) first proves the same exhaustion for the actual character \(\alpha\). This completes the proof. ∎

This proof includes dyadic places. In particular it does not set the Weyl half-pairing character equal to one at2.

\subsection{3. Exact coefficients and a finite-window bound}\label{exact-coefficients-and-a-finite-window-bound}

For \(f\in\mathcal S_N\), using the convention that the inner product is linear in its second variable, write \[
 f=\sum_x c_xW_x\phi_0,\qquad
           c_x=\langle W_x\phi_0,f\rangle,\qquad
                    \sum_x|c_x|^2=\|f\|_2^2.
 \tag{H8}
\] Thus the conversion to actual finite translates has exactly the Parseval norm, not a bound with an unspecified lattice-volume factor.

Let \(\mathcal T\) be a linear map from this finite-dimensional Schwartz window to a normed space; an evaluated theta lift is an example. The triangle inequality and Cauchy--Schwarz give \[
 \max_x\|\mathcal T(W_x\phi_0)\|
 \ \ge\
 \frac{\|\mathcal T(f)\|}{\sum_x|c_x|}
 \ \ge\
 \frac{\|\mathcal T(f)\|}{N^g\|f\|_2}
 \quad\text{if }\mathcal T(f)\ne0.
 \tag{H9}
\] The first denominator is nonzero in this case. This is a quantitative selection bound once \(N\) is known. It does not provide a bound for the \(N\) supplied by a global nonvanishing theorem.

In particular a nonzero linear functional, or a nonzero theta map on a specified input, cannot vanish on every vector in (H5). One actual translate detects it. The detecting vector itself has finite norm one; its coefficient in the full-orbit presentation is one, not a square root of the orbit cardinality. The factor \(N^g\) in (H9) measures a possible loss in a numerical lower bound for the detected output, not a norm of the selected input.

For a local factor at a fixed rational prime \(q\), the same proof works with \(N=q^e\) and only that local lattice, with dimension \(q^{2ge}\). If the input is the vacuum at every other place, the detecting translation can be chosen supported at \(q\). This last assertion concerns only an input already having that support property; it does not prove that a global Rallis test has it.

\subsection{4. A single actual translate in the iterated chain}\label{a-single-actual-translate-in-the-iterated-chain}

Apply this proposition to a later edge of ALP-LIFT-14 and its actual nonzero current vector \(F_j\). That theorem gives a selected real Schur vector and a finite Schwartz vector \(f\) with \[
                    \Theta_{f,\phi_\infty}(F_j)\ne0.
\] If the supplied Schwartz function is a finite sum of tensor products with different real vectors, at least one summand has a nonzero output. Fix that selected real vector. Choose \(N\) with \(f\in\mathcal S_N\), and expand (H8). Linearity gives \[
 \Theta_{f,\phi_\infty}(F_j)
   =\sum_x c_x\Theta_{W_x\phi_0,\phi_\infty}(F_j).
\] At least one summand is nonzero. Replace the finite input by that single \(W_x\phi_0\). Its infinite part is unchanged, and the finite part is an actual Weyl translate with norm exactly one.

The generated target representation remains spherical outside29: this is a statement about its local irreducible constituents, which are the unique local theta quotients of the spherical source constituents, not a claim that the chosen output vector itself is fixed by the maximal compact. The occurrence and Euler-factor argument of ALP-LIFT-14 depend on the constituent's sphericity, not on every current vector being its spherical vector. Its joint-kernel argument applies to every nonzero current vector. Thus it can be iterated without an additional compact averaging that would turn the single translate back into a linear combination.

We obtain the stronger form of the chain: \[
 F_{j+1}
   =\Theta_{W_{x_j}\phi_{0,j},\,\phi_{\infty,j}}(F_j)\ne0,
 \qquad
          \|W_{x_j}\phi_{0,j}\|_{2,f}=1
                \quad(1\le j<m),
 \tag{H10}
\] where every \(\phi_{\infty,j}\) belongs to the exact selected component \(\rho_j^\vee\boxtimes\rho_{j+1}\), and every \(x_j\) has zero infinite component and finite support. No coefficient combination height or finite Schwartz norm growth occurs in this presentation. The unknown denominator and support of \(x_j\), and the output norms and contents, remain unknown.

At infinity the creation-vector norm is precisely the actual norm of \(\phi_{\infty,j}\) in AM-R-011. It is not made one by the finite normalization in (H10). Likewise the source input and automorphic Haar masses are not changed or omitted.

\subsection{5. Exact source interface and remaining restrictions}\label{exact-source-interface-and-remaining-restrictions}

The recovered source definitions give the following literal facts.

\begin{itemize}
\tightlist
\item
  \nolinkurl{Development/AbelianVarietyTheory/Definitions/LatticeFrames/JointSmooth.lean,} lines20--23: comb is a finite sum of heisO translates, and IsFinComb requires only that each translation parameter have zero infinite part.
\item
  \nolinkurl{Development/AbelianVarietyTheory/Definitions/LatticeFrames/IsRatDatum.lean,} lines21--28: the alternate IsFinComb predicate has the same content, whereas IsRatDatum separately requires an actual half-Schwartz function whose rational restriction is the prescribed \(\Xi\), followed by the specified theta transform.
\item
  \nolinkurl{Development/AbelianVarietyTheory/Definitions/LatticeFrames/LevChiEqOne.lean,} lines103--118: PinColT.dat requires both the finite-combination condition and the exact IsRatDatum tie to \(w\,A.\mathrm{xi}\). They are different conditions, and proving the first does not prove the second.
\end{itemize}

Consequently the single translate in (H10) satisfies the mathematical finite-combination condition, with the delta coefficient at \(x_j\). It does not, by that fact, satisfy PinColT.dat or identify the source raw Xi. No Lean record or build is claimed.

There is also a dimension obstruction to silently identifying the basis (H5) with the accepted DP59 half-orbit. For the actual \(g\)-dimensional Hom abelian variety, a level \(M\) window of the form (H4) has dimension \(M^{2g}\). The prescribed DP59 family \([\alpha]^*H^0(\mathcal L^M)\), with \(\alpha\bar\alpha=M\), has rank \(M^g\), and its rescaled Heisenberg parameters lie in one selected isotropic half-subgroup. For \(M>1\) this is a strictly smaller family. The single translate selected by the full-window argument can have both position and momentum components and need not lie in that half-subgroup.

Similarly, the actual raw Xi of \nolinkurl{Formalism/LocalToGlobal/Columns/FiniteSumNonzeroLargeDepth.lean,} lines930--950 and1116--1120, is assembled from the specified PhiAt, omPhi, E1at and local-volume factors, together with phiT. It is not an arbitrary Schwartz function. AM-DO-029 addresses its exact finite Fourier expansion separately. The argument here has not replaced that local data, nor proved that one of its prescribed finite rows detects the previous actual image.

The precise remaining quantitative task is therefore narrower: bound or prescribe the finite translation data while preserving a nonzero action on the already produced image, and identify the resulting class-integrated matrix with the fixed raw Xi and row lattice. Its denominators, class masses and output content still enter the DP44 accounting. A norm-one full-orbit input alone does not compute those quantities.

\subsection{Inputs and check scope}\label{inputs-and-check-scope}

AM-DO-022 supplies the actual selfdual trace lattice, vacuum, additive character and unitary Heisenberg action. ALP-LIFT-14 supplies the actual nonzero selected chain before this refinement. DP59's half-orbit rank is accepted background, cited only to distinguish the two families. The finite-dimensional basis, Parseval calculation, selection bound and single-translate chain refinement above are proved directly here.

No claim is made that the later translations can all be supported at29. Such a sharpening requires a fixed-vacuum nonvanishing test at every other place, including the actual dyadic model; general local theta occurrence alone does not identify that test vector.

\clearpage

\section{AM-R-013: two literal analytic source predicates on the mixed definite host}\label{am-r-013-two-literal-analytic-source-predicates-on-the-mixed-definite-host}

Author: /root, GPT-6 Astra Pro, ALP\_MACHINE, 21 November 2026 UTC.

This result extracts the quantified analytic content of the first two members of \nolinkurl{AdmS} from the actual mixed-host construction. The first is a corollary of ALP-LIFT-14 with an algebraic, rather than merely closed, theta image. The second uses the actual terminal interlacing weights and a finite rational translate to obtain a nonzero compact restriction. Its period has an exact positive value. Neither assertion identifies that period with the prescribed arithmetic \nolinkurl{perCol} or \nolinkurl{rowAF} lattice. In particular this is not a construction of the old elliptic \nolinkurl{D7}, a proof of all six \nolinkurl{AdmS} assertions, or a solution of \nolinkurl{MainX}'s terminal mass equation.

\subsection{1. The source predicates being matched}\label{the-source-predicates-being-matched}

The exact declarations are in \nolinkurl{Formalism/SharpRestriction/CuspSpanNormBounds/CayleySignature.lean}. \nolinkurl{RallisSiegelWeilFormula}, lines154--161, requires actual irreducible cuspidal subspaces \(\pi_j\) of the selected definite types, common edge characters with infinity exponents \(-e_j\), and, for every \(\phi\in\pi_{j+1}\), a finite identity \[
 \phi=\sum_{i=1}^{r_\phi}
       \Theta^{\mathrm{bil}}_{\Phi_i}(\phi_i),
 \qquad \phi_i\in\pi_j,
 \tag{R13.1}
\] with each \(\Phi_i\) factorizable and compact-finite. The predicate does not prescribe the arithmetic row lattice or a norm bound for these kernels.

\nolinkurl{LiftParameterConditions}, lines163--166, requires that every actual irreducible automorphic subspace of the selected last definite type have a nonzero terminal period against some cusp form of the prescribed last \(Y\)-type. The latter type is imposed on the conjugate form by \nolinkurl{IsDv}, lines123--125. For the ordinary compact restriction period the literal combination of conjugations is \[
 B_{\mathrm{top}}(\overline\phi,F)
       =\int_{[H]}\phi(\iota h)F(h)\,dh,
 \quad
 B_{\mathrm{top}}(f,F)
       :=\int_{[H]}\overline{f(\iota h)}F(h)\,dh.
 \tag{R13.2}
\] Here \(G=U(I_n)\), \(n=N_m\), and \(H=U(X')\) is the unitary group of the actual positive codimension-one subspace \(X'=z^\perp\), identified with the terminal \(Y_m\). The embedding \(\iota\) acts as the identity on the orthogonal line. At each real place it is conjugate to \(h\mapsto\operatorname{diag}(h,1)\); no global isometry \(X'\simeq I_{n-1}\) is assumed. This is the actual Hermitian restriction period, with its Haar measure retained. We do not use the auxiliary identity-action Hilbert-space model \nolinkurl{thetaG} to obtain either conclusion.

\subsection{2. An actual irreducible chain and finite factorizable spanning}\label{an-actual-irreducible-chain-and-finite-factorizable-spanning}

Use the common host \(E'=\mathbf Q(\zeta_{28})\), \(F'=E'^+\), the polarized \(J^2\), the same trace coefficient, and all the characters, groups and compact types of ALP-LIFT-14. For every fixed choice of the source rank and HC data there exist actual irreducible cuspidal subspaces \[
 \pi_j\subset\mathcal A(G_j),\qquad
        (\pi_j)_\infty\simeq\rho_j,
             \qquad 0\le j\le m,
 \tag{R13.3}
\] such that (R13.1) holds for every edge. All these are specified automorphic embeddings, so no multiplicity-one assertion for the global spectrum is required.

\textbf{Proof.} The first input of AM-R-012 has fixed finite level and compact real type. The automorphic space at that level and type is finite-dimensional: it is a finite direct sum, over the definite double-coset set, of invariant subspaces of the finite-dimensional coefficient representation. Its spectral decomposition therefore involves finitely many actual irreducible automorphic embeddings. Decompose this input accordingly. Its first theta image is nonzero, so at least one summand, belonging to an actual \(\pi_0\), has a nonzero image with the same selected first real component.

Let \(\mathcal T(\pi_0)\) be the algebraic span of these theta images, allowing all finite Schwartz functions and the selected real component. It is stable under the target finite and real actions by the exact equivariance of the Weil representation. All its automorphic constituents have the unique local small-theta target at every finite place, and the unique selected target compact type at infinity. Consequently they have one global irreducible isomorphism type. The automorphic multiplicity of that type is finite. The smooth algebraic isotypic part is a finite direct sum of irreducible modules, hence is semisimple. Its submodule \(\mathcal T(\pi_0)\) is semisimple as well and contains an actual irreducible submodule \(\pi_1\).

This argument concerns algebraic spans, not just a dense image in a Hilbert-space closure. By definition, every vector of \(\pi_1\) is a finite sum of theta images of vectors of \(\pi_0\). Every finite Schwartz function is a finite sum of restricted pure tensors. Every vector in the selected real component is a finite sum of tensors over the six archimedean places, and is compact-finite. Expanding these finite sums and using bilinearity makes each kernel factorizable without changing the image. This proves (R13.1) at the first edge.

ALP-LIFT-14, sections3--7, proves nonvanishing for every actual irreducible embedding with the resulting unramified parameters, including each possible multiplicity copy. Apply it to \(\pi_1\), then repeat the preceding algebraic-image argument. Induction gives (R13.3) and (R13.1) at every edge. The input and output are cuspidal because each positive definite rational group has no proper rational parabolic. The characters are exactly the globally constructed equal characters of LIFT14, so the required \nolinkurl{IsMu} condition holds. ∎

The proof uses only even \(k,m,d=t+1\), \(d\ge k+m+5\), the HC interlacing/parity, and \(e_j-b_j,b_{j+1}-e_j\ge(d+2)/2\). Thus it applies also when the source's auxiliary \nolinkurl{TPar.t} replaces the fixed \nolinkurl{Ranks.t}: \nolinkurl{TPar.t_ge}, \nolinkurl{Hyp3.Interlace}, \nolinkurl{TParity} and \nolinkurl{Centred} imply exactly these inequalities. The other \nolinkurl{Ranks} inequalities, such as \(3m<t\), are not used by the analytic nonvanishing proof. This proves the first predicate's mathematical quantifiers for that actual interpretation, whenever the surrounding geometric data have been supplied. It is not an assertion that a Lean \nolinkurl{ThetaX D} record has been instantiated.

\subsection{3. Exact terminal branching from the source HC lists}\label{exact-terminal-branching-from-the-source-hc-lists}

Fix one of the six complex places. Since \(m\) is even, the source terminal Harish--Chandra lists have no residual \(\eta\)-shift. From \nolinkurl{DataTheta.lean}, \nolinkurl{TPar.Xpar} and \nolinkurl{TPar.Dpar}, lines538--548, their decreasing rearrangements are \[
 \begin{split}
 x={}&(u_0,\ldots,u_{k-1},\,
          B_{t+1}(-e_0),\ldots,B_{t+1}(-e_{m-1})),\\
 y={}&(v_0,\ldots,v_{k-2},-b_0,\,
       B_t(-e_0),-b_1,B_t(-e_1),\ldots,
                    -b_{m-1},B_t(-e_{m-1})),\\
 B_l(c)={}&(c+(l-1)/2,\,c+(l-3)/2,\ldots,c-(l-1)/2).
 \end{split}
 \tag{R13.4}
\] In this display the terminal line \(y\) consists of its initial \(k-1\) entries followed by the \(m\) pairs \((-b_j,B_t(-e_j))\). It has length \(n-1\), whereas \(x\) has length \(n=k+m(t+1)\). All source signature labels are positive at \(j=m\).

The actual HC assumptions imply \[
 x_i>y_i>x_{i+1}\qquad(0\le i<n-1).
 \tag{R13.5}
\] Indeed, the initial entries alternate by the original \(u,v\) interlacing. Next \(u_{k-1}>-b_0\), while \(e_0-b_0\ge(t+3)/2\) places \(-b_0\) above the top of the first \(x\)-block. Within each block the \(t\) entries of \(B_t(-e_j)\) are exactly the midpoints between the \(t+1\) adjacent entries of \(B_{t+1}(-e_j)\). Between two blocks, \(b_{j+1}-e_j\ge(t+3)/2\) and \(e_{j+1}-b_{j+1}\ge(t+3)/2\) place \(-b_{j+1}\) strictly between their endpoint entries. These observations also justify the stated decreasing rearrangements and the absence of duplicates in the source finite sets.

Put \[
 \lambda_i=x_i-(n-1)/2+i,\qquad
 \beta_i=y_i-(n-2)/2+i.
 \tag{R13.6}
\] These are the dominant integral highest weights corresponding to the two HC lists. The \(x_i\) are half-integers and the \(y_i\) integers; hence each strict gap in (R13.5) is at least \(1/2\). Substituting (R13.6) gives \[
        \lambda_i\ge\beta_i\ge\lambda_{i+1}.
 \tag{R13.7}
\] The compact branching theorem therefore gives \[
 \dim\operatorname{Hom}_{U(n-1)}
      (V_\beta,\operatorname{Res}^{U(n)}_{U(n-1)}V_\lambda)=1.
 \tag{R13.8}
\] This is the ordinary branching theorem, cited rather than re-proved: Goodman and Wallach (2009), Theorem8.1.1, p.364. Restriction from the complex algebraic groups to their compact forms preserves these finite-dimensional representation multiplicities. Taking the external tensor product over the six places gives the actual terminal compact type \(\tau\) inside \(\rho_m|_{H_\infty}\), with multiplicity one. Its HC parameter is exactly \nolinkurl{Dpar m}, not a freely selected branching constituent.

\subsection{4. A compact restriction detector for every actual terminal embedding}\label{a-compact-restriction-detector-for-every-actual-terminal-embedding}

The next elementary argument supplies the global existence needed after the local branching statement. Local occurrence alone would not justify restriction of a fixed automorphic vector at a fixed finite point.

Let \(\pi\subset\mathcal A(G)\) be any nonzero irreducible automorphic embedding with \(\pi_\infty=\rho_m\). There are \(\phi\in\pi\) and a smooth cusp form \(F\) on \(H\) such that \[
 \overline F\text{ has real type }\tau,
 \qquad
 B_{\mathrm{top}}(\overline\phi,F)
       =\|\operatorname{Res}_H\phi\|_{L^2([H])}^{2}>0.
 \tag{R13.9}
\] Here \(\phi\) is chosen so that its restriction already has type \(\tau\); no projection denominator is suppressed.

\textbf{Proof.} First, \(G(F')\) is dense in \(G_\infty\). For the identity Hermitian form, its Lie algebra over \(F'\) is the vector space of matrices satisfying \(X^*=-X\). Archimedean weak approximation for this vector space is elementary: choose an \(F'\)-basis; the diagonal image of \(F'\) is dense in the product of its real completions, since a rational basis gives an invertible real embedding matrix and \(\mathbf Q^s\) is dense in \(\mathbf R^s\). Approximate each coordinate independently. The rational Cayley map \[
             X\longmapsto(1+X)(1-X)^{-1}
 \tag{R13.10}
\] is defined on every real skew-Hermitian matrix and maps onto the open dense subset of the compact unitary group having no eigenvalue \(-1\). Its inverse there is \((g-1)(g+1)^{-1}\). The same rational identities hold over \(F'\). Approximation through this chart proves density simultaneously at all real places. No bound on denominators of the rational approximant is asserted.

Choose an actual nonzero finite vector in \(\pi_f\), and let \(v\mapsto\phi_v\) be the resulting \(G_\infty\)-equivariant map from \(V_{\rho_m}\) to the automorphic embedding. Some finite point has nonzero evaluation on this map. A right finite translation allows us to take that point to be the identity. There is then a nonzero linear functional \(L\) on \(V_{\rho_m}\) with \[
                  \phi_v(k,1)=L(\rho_m(k)v).
 \tag{R13.11}
\] Let \(V_\tau\subset V_{\rho_m}\) be the nonzero branching subspace supplied by (R13.8). The span of \(\rho_m(G_\infty)V_\tau\) is a nonzero invariant subspace of the irreducible \(\rho_m\), so it is the whole space. Consequently there exist \(k\in G_\infty\) and \(v\in V_\tau\) with \(L(\rho_m(k)v)\ne0\). This inequality is open in \(k\). By the Cayley density just proved, choose \(q\in G(F')\) with \(L(\rho_m(q_\infty)v)\ne0\).

Right-translate \(\phi_v\) by \(q_f^{-1}\), and call the result \(\phi\). For \(h\in H_\infty\), left rational automorphy gives the exact identity \[
 \begin{split}
 \phi(\iota h,1)
 &=\phi_v(\iota h,q_f^{-1})
  =\phi_v(q_\infty\iota h,1)\\
 &=L\bigl(\rho_m(q_\infty)\rho_m(\iota h)v\bigr).
 \end{split}
 \tag{R13.12}
\] Thus the restriction is nonzero, already at the identity. Finite translation commutes with the real group action. Since \(v\in V_\tau\), the full restricted automorphic function \(\operatorname{Res}_H\phi\), at all its finite points, has real type \(\tau\).

Set \(F=\overline{\operatorname{Res}_H\phi}\). Restriction is well-defined on \([H]\) by rational automorphy. It is smooth and of some finite compact level. Both groups are anisotropic over \(F'\), so their quotients are compact and all their automorphic forms are cuspidal. Hence this \(F\) is permitted in the analytic terminal predicate, and its conjugate has precisely the required real type. Substitution into (R13.2) gives (R13.9). The integral is positive because a nonzero continuous function on the compact quotient is nonzero on an open set of positive Haar measure. ∎

\subsection{5. Retained class masses and the remaining arithmetic equation}\label{retained-class-masses-and-the-remaining-arithmetic-equation}

The value in (R13.9) is exact for the chosen measure. For any compact finite level \(K_H\) fixing the restricted function, choose representatives \(h_c\) for \(H(F')\backslash H(\mathbb A_f)/K_H\), and write \(\Gamma_c=H(F')\cap h_cK_Hh_c^{-1}\). Normalize real Haar measure to volume one and retain the finite volume \(\operatorname{vol}(K_H)\). Then the same value is \[
 \sum_c\frac{\operatorname{vol}(K_H)}{|\Gamma_c|}
       \int_{H_\infty}|\phi(\iota h,\iota h_c)|^2\,dh.
 \tag{R13.13}
\] This follows by the ordinary quotient-measure disintegration over the finite classes; it is the positive mass convention already used in AM-DO-026 and AM-R-012. It neither normalizes the total automorphic volume to one nor drops the stabilizers. Scaling \(F\) by its reciprocal would force a unit period over \(\mathbf C\), but AM-R-002 shows why this would still have an arithmetic content and norm cost.

Equations(R13.1) and(R13.9) supply the first two analytic predicates with their actual quantifiers and compact types. The other four predicates involve the noncompact \(Y_j\) descent, specified Fourier--Jacobi periods, the lower-tower vanishing, and real type stability. Those are not supplied by this proof. In addition, \nolinkurl{TieRowIsom}, \nolinkurl{TieIo}, and \nolinkurl{TieBm} in \nolinkurl{DataTheta.lean}, lines739--752, still impose separate arithmetic comparisons. In particular the explicit positive restriction period (R13.13) has not been proved equal to the source's prescribed terminal coefficient matrix. The finite translation \(q_f^{-1}\), the resulting level, and the integral coefficient content remain unquantified. These are precisely the limitations relevant to the requested DP44 lift-count comparison.

\subsection{6. Geometric interpretation and sources}\label{geometric-interpretation-and-sources}

The compact branching is a restriction of algebraic representations of the reductive unitary groups, which become \(\mathrm{GL}_n\) and \(\mathrm{GL}_{n-1}\) after splitting the CM extension. The integral weights define the usual Schur functors and determinant twists; their characters live in the representation \(\lambda\)-ring. No rational projector in this proof is asserted to preserve a chosen integral lattice. This distinction is the same one recorded in the SGA6/SGA3 entries of \nolinkurl{GROTHENDIECK_DICTIONARY.md}. The theta kernels still arise from the principal polarization's biextension and the CM Picard data already identified there (SGA7I, Exp.VII, Proposition2.9.3, pp.172--174; Exp.VIII, Theorem7.1, pp.299--300; FGA232, Theorem3.1, pp.148--149). The restriction norm is an archimedean realization. A metrized integral determinant-line comparison with it is additional work, not a formal consequence of functoriality of motives or of weights for sheaves over finite fields.

Primary new literature input: Goodman, R., \& Wallach, N. R. (2009). \emph{Symmetry, representations, and invariants}. Springer. Theorem8.1.1, p.364; DOI10.1007/978-0-387-79852-3. The theorem was directly inspected in the primary book PDF at \url{https://webhomes.maths.ed.ac.uk/~v1ranick/papers/goodwallx.pdf} (PDF page383). The compact branching theorem is cited as known.

Internal inputs, with their separate check histories: ALP-LIFT-14, sections3--8; AM-DO-025--026; AM-R-002, AM-R-012; AM-K022. Exact source declarations used above are \nolinkurl{CayleySignature.lean}, lines117--166, and \nolinkurl{DataTheta.lean}, \nolinkurl{TPar}, \nolinkurl{block}, \nolinkurl{Dpar}, \nolinkurl{Xpar}, \nolinkurl{ThetaX.IsStable}, \nolinkurl{typeX}, and \nolinkurl{TieBm}. No Lean execution or old-host instantiation is claimed.

\clearpage

\section{AM-DO-029: the actual raw Xi as a finite Heisenberg vector, its exact cost, and the 29-only mismatch}\label{am-do-029-the-actual-raw-xi-as-a-finite-heisenberg-vector-its-exact-cost-and-the-29-only-mismatch}

Author: GPT-6 Astra Pro, ALP\_MACHINE \nolinkurl{/root/definite_oscillator}; inherited reasoning effort is not separately exposed; 21 November 2026 UTC.

This finding concerns the \textbf{prescribed} \nolinkurl{xiAt} and the finite clauses of \nolinkurl{IsRatDatum}. It uses the carried THETA3-A12/A13/A15 finite-model results at their checked status, and the actual mixed host and normalization of AM-DO-022/026. It does not replace a prescribed array by an arbitrary Schwartz function. The new conclusions are its explicit finite-window realization on the mixed host, exact norm and denominator accounting, and a precise obstruction to identifying it with the previously constructed input supported only at 29.

\subsection{1. Literal source interface and local windows}\label{literal-source-interface-and-local-windows}

All source paths below are relative to \nolinkurl{work/research/alpoge_source}, pinned to commit \nolinkurl{7cbf7a0c32fca1badf05931879a37c3f8a2b03b4}. Reading these recovered definitions is not a Lean compilation or a transitive dependency audit.

\begin{longtable}[]{@{}
  >{\raggedright\arraybackslash}p{(\columnwidth - 4\tabcolsep) * \real{0.3333}}
  >{\raggedright\arraybackslash}p{(\columnwidth - 4\tabcolsep) * \real{0.3333}}
  >{\raggedright\arraybackslash}p{(\columnwidth - 4\tabcolsep) * \real{0.3333}}@{}}
\toprule\noalign{}
\begin{minipage}[b]{\linewidth}\raggedright
Source
\end{minipage} & \begin{minipage}[b]{\linewidth}\raggedright
Locator
\end{minipage} & \begin{minipage}[b]{\linewidth}\raggedright
Content used
\end{minipage} \\
\midrule\noalign{}
\endhead
\bottomrule\noalign{}
\endlastfoot
\nolinkurl{Development/AbelianVarietyTheory/Definitions/LatticeFrames/IsRatDatum.lean} & 19--28 & \nolinkurl{IsFinComb}; the three clauses of \nolinkurl{IsRatDatum} \\
\nolinkurl{Development/AbelianVarietyTheory/Definitions/LatticeFrames/LevChiEqOne.lean} & 68--71 & actual \nolinkurl{TwDat.xi}, obtained from \nolinkurl{xiAt} and the specified frames \\
\nolinkurl{Development/AbelianVarietyTheory/Definitions/AutomorphicMembership/XiAtPenultimateEqXiCpl1.lean} & 146--163 & the three branches \nolinkurl{XiEntry}, \nolinkurl{XiCpl1}, \nolinkurl{xiLast}; all contain \nolinkurl{phiT} \\
\nolinkurl{Formalism/LocalToGlobal/Columns/FiniteSumNonzeroLargeDepth.lean} & 646--648; 927--948; 981--986 & \nolinkurl{phiT}; the local \nolinkurl{PhiAt}, \nolinkurl{omPhi}, \nolinkurl{E1at}, \nolinkurl{XiAt}; \nolinkurl{joinCol}, \nolinkurl{PsiAt} \\
same & 1047--1065 & \nolinkurl{ChIdx.Ok.psiSupp}, \nolinkurl{psiInv}, \nolinkurl{xiSupp}, \nolinkurl{xiInv}, \nolinkurl{psiSuppX}, \nolinkurl{psiInvX} \\
same & 1108--1138 & \nolinkurl{Sall}, \nolinkurl{cAof}, \nolinkurl{cBof}, \nolinkurl{XiEntry}, \nolinkurl{XiCpl1} \\
\nolinkurl{Formalism/LocalToGlobal/IndefiniteReal/LocalNewformFunction.lean} & 28--37; 239--264; 270--279 & \nolinkurl{symplO}, \nolinkurl{heisO}; \nolinkurl{IsSBHalf}, \nolinkurl{IsPlaceModel}, \nolinkurl{IsThetaAt}; the actual \nolinkurl{localNewformFunction} \\
\end{longtable}

Fix one nonterminal edge \(0\leq i\leq m-2\), ranks

\[
 a=N(i),\qquad b=N(i+1),\qquad n=ab,
 \tag{DO29.1}
\]

and one allowed source and target decoration, local index, and frame. Write \(\Xi\) for the resulting scalar \nolinkurl{TwDat.xi}, before or after a specified nonzero chart scalar. The scalar is retained in every formula. Let \(\mathcal S=\mathrm{Sall}\) be its actual finite half-place set. The source's selected-place and prime hypotheses are retained; in particular this is a set of relative split places with no conjugate pair selected twice.

At a \(T\)-place \(P\), the factor is supported on

\[
 \operatorname{Hom}_{\mathcal O_{E,P}}(N_P,N'_P)
 \tag{DO29.2}
\]

and is constant modulo \(P\operatorname{Hom}(N_P,N'_P)\). In residue frames it is exactly

\[
 I_{v,u}(X)=1_{\{vX=u\}},\qquad u\ne0,\quad v\ne0.
 \tag{DO29.3}
\]

This includes matrix integrality. The assertion follows directly from the localized lattice condition in \nolinkurl{phiT} and the residue-linearity in \nolinkurl{Deco.Ok}; the exact incidence description and its rank-one Fourier span are carried THETA3-A13.

At a place \(P\in\mathrm{rowSG}\), let \(f_P,f'_P\) be the specified source and target frames and put

\[
 L_P^{\rm fr}=f'_P M_{b\times a}(\mathcal O_{E,P})f_P^{-1}.
 \tag{DO29.4}
\]

For an \nolinkurl{XiEntry} branch, the literal \nolinkurl{xiSupp} and \nolinkurl{xiInv} fields give support in \(q_P^{-A}L_P^{\rm fr}\) and period lattice \(q_P^B L_P^{\rm fr}\). For an \nolinkurl{XiCpl1} branch, the literal \nolinkurl{psiSuppX} and \nolinkurl{psiInvX} fields give the same bounds in the matrix variable, with the actual fixed representative of the extra column retained. Here \(q_P=\mathrm{qOf}(P)\), and the exponents mean scalar multiplication exactly as in the source. No assumption that an arbitrary function \nolinkurl{LocCst.psi} is continuous outside its prescribed depth is needed. These are already hypotheses of the actual \nolinkurl{ChIdx.Ok} record.

The first edge of every allowed chain is one of these branches, since the source has \(m\geq2\). The present assertion does \textbf{not} silently extend these nonterminal fields to \nolinkurl{xiLast}. That branch has the additional variable additive character in \nolinkurl{ftY}; its stated depth and terminal representative conditions must be used separately.

\subsection{2. The explicit extension and finite transform}\label{the-explicit-extension-and-finite-transform}

Use the actual mixed finite symplectic space from AM-DO-022:

\[
 E'=\mathbf Q(\zeta_{28}),\quad
 X=eW_f,\quad Y=\bar eW_f,\quad
 L=e\widehat\Lambda,
 \tag{DO29.5}
\]

where \(\Lambda=M_{b\times a}(\mathcal O_{E'})\), and \(e\) is the selected finite half idempotent for \(\mathcal S\). The mixed scalar trace coefficient is \(\xi'=\xi/2\); the product lattice and the dyadic basis index remain exactly those of AM-DO-022. Write \(B\) for the effective alternating pairing with the actual fixed additive character \(\psi_f\). Thus \(W_xW_y=\psi_f(B(x,y)/2)W_{x+y}\). This notation incorporates the source's standard-character sign, rather than changing it in a Fourier formula.

Let \(L^+\supseteq L\supseteq L^-\) be compact open lattices in \(X\) such that the local bounds of §1 imply

\[
 \operatorname{supp}\Phi\subseteq L^+,
 \qquad \Phi(t+v)=\Phi(t)\quad(v\in L^-).
 \tag{DO29.6}
\]

They can be chosen without an unspecified enlargement. At each \(T\)-place take the sum of the standard lattice and (DO29.2) for \(L_P^+\), and intersect the standard lattice with \(P\) times (DO29.2) for \(L_P^-\). At each other selected place take

\[
 L_P^+=L_P+q_P^{-A}L_P^{\rm fr},\qquad
 L_P^-=L_P\cap q_P^B L_P^{\rm fr}.
 \tag{DO29.7}
\]

Products give \(L^+,L^-\). Smith normal forms of the displayed local frame matrices give their exact elementary divisors and indices. Thus no frame index is suppressed. Define

\[
 R=[L^+:L],\qquad N=[L:L^-],\qquad
 C=L^+/L,\quad U=L/L^-.
 \tag{DO29.8}
\]

Choose representatives \(a\in L^+\) of \(C\), and representatives

\[
 y\in (L^-)^\perp/L^\perp\subset Y/L^\perp
 \tag{DO29.9}
\]

of the characters of \(U\), where annihilators use \(\psi_f\circ B\). This is a perfect finite pairing; in local coordinates it is the ordinary duality between a lattice quotient and its inverse-different dual. Its order is \(N\).

There is a unique locally constant function \(\Phi\) on \(X\), extended by zero off \(X\), satisfying

\[
 \Phi(e\gamma)=\Xi(\gamma)\quad
     (\gamma\in M_{b\times a}(E')).
 \tag{DO29.10}
\]

\textbf{Proof of existence and uniqueness.} At the finitely many selected places, use the factors of §1. The support and period statements make their product constant on every coset of \(L^-\) and zero off \(L^+\). Every required finite congruence class has a rational matrix representative by additive weak approximation, or equivalently the Chinese remainder theorem after clearing denominators. Its value is independent of the representative by the period statements. This defines the function and proves (DO29.10). Conversely rational matrices meet every open coset of \(L^-\), so no other locally constant extension with these periods is possible. Choose a positive integer \(k\) with \(kL\subseteq L^-\) and \(kL^+\subseteq L\); the local elementary divisors give one explicitly. Then (DO29.6) is exactly the support and period part of \nolinkurl{IsSBHalf}, and support on \(X=eW_f\) is its remaining clause. \(\square\)

The exact Fourier coefficients are

\[
 c_{a,y}=\frac1N\sum_{u\in U}
                 \Phi(a+u)\psi_f(-B(u,y)).
 \tag{DO29.11}
\]

Finite character orthogonality gives the literal identity

\[
 \Phi(t)=\sum_{a\in C,y\in\widehat U}
 c_{a,y}\,\psi_f(B(t-a,y))1_L(t-a).
 \tag{DO29.12}
\]

Let \(T\) be any normalized \nolinkurl{IsPlaceModel}, so that

\[
 T1_L=f_0,\quad T\Phi(\cdot+x)=W_xT\Phi,
 \quad T(\psi_f(B(\cdot,y))\Phi)=W_yT\Phi.
 \tag{DO29.13}
\]

Then the required finite vector is \textbf{exactly}

\[
 \boxed{f_\Xi=\sum_{a,y}c_{a,y}W_{-a}W_yf_0
   =\sum_{a,y}c_{a,y}\psi_f(-B(a,y)/2)W_{y-a}f_0.}
 \tag{DO29.14}
\]

Indeed apply modulation and then translation to \(1_L\) in (DO29.12), and use the displayed Weyl law. Every normalized model gives the same vector. This is precisely the carried A12 construction and the literal theorem \nolinkurl{thetaAt_eq_sum_heis_translates}, now with the coefficients of the \textbf{actual} array rather than freely chosen ones. The theorem \nolinkurl{exists_isThetaAt_of_vac}, lines124--132 of \nolinkurl{Development/AbelianVarietyTheory/Proofs/DivisorProofs/ExistsIsThetaAtOfVac.lean}, supplies existence in the original source wrapper; the actual mixed Schrödinger model supplies existence on the mixed host. Section5 states their precise finite comparison.

Equations (DO29.6), (DO29.10), and (DO29.14) verify all three mathematical clauses of \nolinkurl{IsRatDatum}. They also specify the exact finite list and coefficients for \nolinkurl{IsFinComb}; every parameter has zero infinite component. They do not assert that an unchanged elliptic \nolinkurl{D7} wrapper has been instantiated for the new CM type.

\subsection{3. Exact norm, denominators, and local newform examples}\label{exact-norm-denominators-and-local-newform-examples}

Normalize the finite Haar measure by \(\operatorname{vol}(L)=1\) and the finite vacuum by \(\|f_0\|=1\). Then

\[
 \boxed{\|f_\Xi\|^2
   =\sum_{a,y}|c_{a,y}|^2
   =\frac1N\sum_{a\in C,u\in U}|\Phi(a+u)|^2.}
 \tag{DO29.15}
\]

\textbf{Proof.} The functions \(\psi_f(B(t-a,y))1_L(t-a)\) have disjoint supports for different \(a\), and character orthogonality on \(L/L^-\) gives orthogonality for different \(y\) with the same \(a\). Each has norm1. The same orthogonality holds for their images \(W_{-a}W_yf_0\): lattice modulations distinguish different position cosets, lattice translations distinguish different characters, and these operators are unitary. The vacuum norm fixes the norm of every image at1. Consequently (DO29.14) is an orthonormal expansion. Finite Parseval gives the second equality. This proves the normalization even though the bare predicate \nolinkurl{IsPlaceModel} does not separately include an isometry axiom. \(\square\)

If \(H_\Xi=\max_{a,u}|\Phi(a+u)|\), then

\[
 \|f_\Xi\|\leq\sqrt R\,H_\Xi,
 \qquad \sum_{a,y}|c_{a,y}|
       \leq\sqrt{RN}\,\|f_\Xi\|
       \leq R\sqrt N\,H_\Xi.
 \tag{DO29.16}
\]

These follow respectively from the \(RN\) summands in (DO29.15) and Cauchy--Schwarz. They are bounds, whereas (DO29.15) is exact.

Let \(K_\Xi\) contain all actual table values and all character values in (DO29.11)--(DO29.14). The latter are roots of unity. If \(\delta\Phi(a+u)\) is integral over \(\mathcal O_{K_\Xi}\), then

\[
 \delta N c_{a,y}\in\mathcal O_{K_\Xi}.
 \tag{DO29.17}
\]

Equivalently, at a finite valuation \(v\),

\[
 v(c_{a,y})\geq
       \min_{a,u}v(\Phi(a+u))-v(N).
 \tag{DO29.18}
\]

The central phase in (DO29.14) is a unit and adds no valuation. The values of the actual array still contain its \nolinkurl{volN} factors, finite sums, specified characters and chart scalar; their valuation is \textbf{not} replaced by0. If the finite level and lattice changes are prime to11, the actual theta-group isomorphisms also preserve the integral11-adic theta lattice. Thus (DO29.17) is then a bound for the geometric coefficient realization after the common period has been removed. This assertion makes no claim at an auxiliary place of residue characteristic11.

For standard frames, at a degree-one selected \(T\)-prime the local indices are \(R_T=1,N_T=q^n\). At a selected degree-one row prime they are \(R_P=q^{An},N_P=q^{Bn}\). Thus, when all these frames are integral units, one may use

\[
 R=\prod_{P\in\mathrm{rowSG}}q_P^{An},\qquad
 N=T(u)^{[F':\mathbf Q]n}
           \prod_{P\in\mathrm{rowSG}}q_P^{Bn}.
 \tag{DO29.19}
\]

Repeated rational primes in the product are intentional: each selected half place contributes its matrix quotient. For other frames or residue degrees, (DO29.7)--(DO29.8), rather than this simplified formula, give the exact indices.

\subsubsection{3.1 The actual incidence factor}\label{the-actual-incidence-factor}

At a \(T\)-place with residue field \(\kappa\) of order \(q\), choose the residue additive character induced by the actual pairing. For nonzero \(u,v\) as in (DO29.3),

\[
 1_{M_{b\times a}(\mathcal O)}(t)I_{v,u}(\bar t)
   =q^{-a}\sum_{\eta\in\kappa^a}
       \psi_\kappa(-\eta\cdot u)
       \psi_\kappa(\eta\cdot v\bar t)1_L(t).
 \tag{DO29.20}
\]

This follows by summing a character of the \(a\)-dimensional additive group \(\kappa^a\): the sum is \(q^a\) when \(v\bar t=u\), and0 otherwise. Let \(y_{v,\eta}\in P^{-1}L^\perp/L^\perp\) be the exact dual vector for the character \(t\mapsto\psi_\kappa(\eta\cdot v\bar t)\), including the coefficient of the actual trace pairing. Therefore

\[
 f_{v,u}=q^{-a}\sum_{\eta\in\kappa^a}
                 \psi_\kappa(-\eta\cdot u)W_{y_{v,\eta}}f_0,
 \quad \|f_{v,u}\|^2=q^{-a},
 \quad \sum|c_\eta|=1.
 \tag{DO29.21}
\]

The map \(\eta\mapsto y_{v,\eta}\) is injective because \(v\ne0\). Thus the norm also follows directly from \(q^a\) orthogonal terms of coefficient absolute value \(q^{-a}\). Its denominator is exactly \(q^a\) in this expansion. The frequencies are the matrices \(v^t\eta\), of rank at most1; this is the source incidence restriction, not a full matrix-delta projector.

\subsubsection{\texorpdfstring{3.2 The literal \nolinkurl{localNewformFunction}}{3.2 The literal }}\label{the-literal-localnewformfunction}

Let the local field have residue order \(q\), uniformizer \(\varpi\), and let \(P\subset\mathrm{GL}_r(\mathcal O)\) be a compact open subgroup. Assume the specified finite characters are continuous and unitary, as in the actual compact type. The source function is

\[
 F\binom{A}{B}=
 \begin{cases}
 \lambda(A)\mu(\det A)^{-1},
       &A\in P,\ B\in\varpi M_{t\times r}(\mathcal O),\\
 0,&\text{otherwise}.
 \end{cases}
 \tag{DO29.22}
\]

This is the literal definition at the locator in §1; it includes both its unit block and its lower-block divisibility. Choose \(s\geq1\) such that \(K_r(\varpi^s)\subset P\) and the product character \(\lambda\mu^{-1}\circ\det\) is trivial there. Then (DO29.11) with \(R=1,N=q^{s(r+t)r}\) is an explicit finite Heisenberg expression for this function. More sharply its norm is

\[
 \boxed{\|F\|^2
  =q^{-tr}\,
     \frac{\prod_{j=1}^{r}(1-q^{-j})}
          {[\mathrm{GL}_r(\mathcal O):P]}.}
 \tag{DO29.23}
\]

To prove the period assertion, changing \(A\) by \(\varpi^sZ\) multiplies it on the right by an element of \(K_r(\varpi^s)\) whenever \(A\in\mathrm{GL}_r(\mathcal O)\). It therefore preserves both support and value; changes in \(B\) by \(\varpi^s\) preserve its divisibility. For the norm, its absolute value on its support is1. Additive matrix measure gives \(M_r(\mathcal O)\) volume1 and \(\mathrm{GL}_r(\mathcal O)\) volume \(\prod_j(1-q^{-j})\). Multiplication by any integral unit matrix preserves this additive measure, so \(P\) has exactly the displayed index fraction of this volume. The lower block has volume \(q^{-tr}\). Multiplication gives (DO29.23). This proves every factor and does not identify a group probability measure with additive matrix measure.

The actual \nolinkurl{PhiG}, \nolinkurl{PhiGS}, \nolinkurl{PhiQ}, \nolinkurl{E1at} and their specified frame operations are retained inside the table (DO29.11). In particular a finite auxiliary sum is not replaced by just the local function (DO29.22). Equation (DO29.15) computes its exact norm after that sum, including cancellation.

\subsection{4. A nonzero actual Xi cannot be a 29-only input}\label{a-nonzero-actual-xi-cannot-be-a-29-only-input}

Let \(\mathcal V_{29}\) be the span of all finite Heisenberg translates of the actual vacuum whose parameters are supported at places over29. This includes the actual permitted half-orbit used in AM-K022 and AM-R-012, and is deliberately larger than that half-orbit. At all other places every vector in \(\mathcal V_{29}\) has its fixed vacuum factor.

For an actual valid prime and decoration datum, suppose \(\Xi\not\equiv0\) is a nonterminal \nolinkurl{xiAt}. Then

\[
 \boxed{f_\Xi\notin\mathcal V_{29}.}
 \tag{DO29.24}
\]

Here equality is tested in one common finite model, enlarging the selected half set by the29 half places if necessary. The assertion is unchanged by a nonzero scalar multiplying \(\Xi\).

\textbf{Proof.} The prime axioms give three distinct rational primes \(T(u),g_0,g_1\). This is literal \nolinkurl{Primes.Ok.g_ne} and \nolinkurl{T_ne_g}, lines200--239 of \nolinkurl{Formalism/LocalToGlobal/Columns/AbelianHostFieldConditions.lean}. There is at least one selected place above each: their stated cardinality is \([F':\mathbf Q]>0\).

If \(T(u)\ne29\), choose a \(T\)-place \(P\nmid29\). Its factor (DO29.3) vanishes at0 because \(u\ne0\), and on an open neighborhood of0 by its period lattice. If \(T(u)=29\), choose a \(G\)-place \(P\nmid29\). Carried THETA3-A13 proves locally that nonvanishing of this factor forces the underlying matrix \(x:E'^a\to E'^b\) to be injective. In its \nolinkurl{XiEntry} proof, each enlarged summand has rank at most \(\operatorname{rank}x+c_i\) but requires a unit minor of size \(a+c_i\); in the penultimate branch \nolinkurl{joinCol} gives the same inequality with \(c_i=1\). Thus its local value at \(x=0\) is0, and the actual \nolinkurl{xiInv} or \nolinkurl{psiInvX} gives an open neighborhood on which it remains0. In all cases there is a mandatory factor at some \(P\nmid29\) that vanishes near0.

In a common half model, a vector in \(\mathcal V_{29}\) has function

\[
 G(t)=G_{29}(t_{29})
              \prod_{Q\in\mathcal S,\ Q\nmid29}1_{L_Q}(t_Q).
 \tag{DO29.25}
\]

This follows by applying only29-local operators to the tensor vacuum. The29 operators need not be pure position translations; their action is all contained in \(G_{29}\). If this function equalled the nonzero \(\Phi\) of (DO29.10), choose a point where the common value is nonzero. All its coordinates away from29 are in their \(L_Q\). Replace just its \(P\)-coordinate by0. Formula (DO29.25) leaves its nonzero value unchanged, whereas the mandatory source factor makes \(\Phi\) zero there. This is a contradiction. The same argument can be made entirely on rational samples: use a sufficiently small open coset at each selected place and additive weak approximation. The finite-model map is injective on the window by (DO29.15), so different functions cannot represent the same vector. \(\square\)

This obstruction is to the \textbf{29-only span}, not to the oscillator representation, not to arbitrary finite inputs, and not to a changed valid global tuple. Equation (DO29.14) gives the required nonzero vector with the other local factors included.

\subsection{5. What the mixed finite comparison proves}\label{what-the-mixed-finite-comparison-proves}

At every selected good relative split place away from2 and7, both the source trace coefficient and the mixed coefficient \(\xi'\) are local units. Express the two effective symplectic forms using the same fixed additive character. On \(X_P\oplus Y_P\) their pairing coefficients differ by a unit \(r_P\) of the real local field. The linear map

\[
 C_P(x+y)=x+r_Py,
 \qquad B_{\rm mixed}(C_Pz,C_Pw)=B_{\rm source}(z,w)
 \tag{DO29.26}
\]

is the explicit symplectic transport; \(r_P\) is by definition the ratio of the source coefficient to the mixed coefficient. It fixes the position variable and preserves both integral half lattices. It therefore fixes \(1_{L_P}\), transports every character in (DO29.11) to its actual dual vector, and preserves the finite measure normalization. It also commutes with \(X\mapsto hXg^{-1}\), since it is scalar on the dual half. The split-place determinant-root normalization from AM-DO-020 is retained. There is no extra unidentified phase in this normalized finite Heisenberg map.

Consequently the \textbf{same} raw table \(\Xi(\gamma)\) determines the mixed vector (DO29.14), with its exact finite norm and denominators. On the actual mixed theta family, AM-DO-026 gives its all-polynomial reading:

\[
 \vartheta_{r_P f_\Xi}(g,h)
 =C^{2ab}\epsilon_\chi(g,h)
   \left[P(\partial)\sum_{a,y}
     c_{a,y}\psi_f(-B(a,y)/2)
       \mathcal U_{A_{g,h}(y-a)}F_{N,N'}\right]_{0}.
 \tag{DO29.27}
\]

Every finite phase, chart scalar, and common period \(C^{2ab}\) is displayed or defined in the checked transported operator. This comparison constructs the finite part of a \textbf{new mixed datum}. It does not identify the old elliptic \nolinkurl{VacIn} seed with the primitive mixed-CM seed, and does not assert a new \nolinkurl{D7} instance without changing that wrapper. If a chosen auxiliary place is bad for (DO29.26), its actual lattice indices and dual pairing must be used instead; the abstract finite formula still holds but the unit-lattice simplification is unavailable.

\subsection{6. Consequences for the selected first edge and MainX}\label{consequences-for-the-selected-first-edge-and-mainx}

The identity \(\Xi(0)=0\), already settled in THETA3-A13, does \textbf{not} assert \(\vartheta_{r_Pf_\Xi}(1,1)=0\). The first is one value of a finite position function. The second is the weighted theta/jet sum (DO29.27). They are different linear functionals. In particular a nonintegral position translate from R012 itself can have finite value0 at the origin and a nonzero selected theta jet.

R012 proves nonvanishing of its selected jet on some member of the actual29 half-orbit. Equation (DO29.24) prevents using that statement to identify its finite vector directly with the nonzero prescribed \(f_\Xi\). Conversely (DO29.14) removes any merely existential finite-vector obstacle: the prescribed raw array is now an explicit admissible finite Heisenberg combination on the mixed host.

The remaining first-edge question is the concrete finite functional

\[
 \sum_{a,y}c_{a,y}\psi_f(-B(a,y)/2)
       [P(\partial)\mathcal U_{y-a}F_{00}]_0
             \stackrel{?}{\ne}0,
 \tag{DO29.28}
\]

or its actual matrix version over all allowed indices and decorations. A general nonzero theta vector or a nonzero term of this sum does not prove (DO29.28). The exact criterion in carried A15 remains applicable with the selected higher jet substituted for its first derivative. Its actual local incidence and rank restrictions remain inside the coefficients.

The next result treats a fixed nonzero auxiliary operator at primes disjoint from29 and11, acting on the \textbf{actual} K022 half-orbit. It will prove a quantitative selected-jet detector using simultaneous torsion-point interpolation. Even that result does not silently make a translated29 vacuum into the fixed \nolinkurl{Q}-factor at29 of an already prescribed \(\Xi\); that is a separate image or admissible-parameter equation.

\subsubsection{Check status}\label{check-status-3}

Author whole-proof review and independent review are recorded in \nolinkurl{definite_oscillator_CHECK_EVENTS.tsv} when actually performed. All new equations are finite-model or elementary lattice/measure arguments proved above. No Lean build or full MainX constructor is claimed.

\clearpage

\section{AM-K023 --- the actual selected jet matrix modulo 11}\label{am-k023-the-actual-selected-jet-matrix-modulo-11}

Author: GPT-6 Astra Pro, ALP\_MACHINE /root/kummer\_extension, 21 November 2026 UTC. Inherited reasoning effort is not separately exposed.

AM-K022 removes the finite half-orbit saturation loss from a specified finite jet matrix. This note computes the reduction of that jet matrix by a finite polynomial normal form. It also gives actual rectangular selected rows of unexpectedly large degree whose exact 11-content is zero. These are local selected kernel statements; they do not compute a later automorphic class integral.

\subsection{1. Every integral local section, not just the vanishing sigma section}\label{every-integral-local-section-not-just-the-vanishing-sigma-section}

Work over a finite unramified complete extension R of Z₁₁ containing the CM eigenvalues and the needed prime-to-11 torsion. Let k be its residue field. The ordinary Jacobian J, its principal trace-one line L=L₂∞, its q-normalized section, and the canonical source reduced gauge are those of AM-K018 and ALP-LIFT-08. Let F be the normalized nonvanishing even theta germ at D₁₃₅ from AM-K012. Its algebraic section is an integral unit at the origin, by ALP-LIFT-07, and F(0)=1. Its Hurwitz derivatives are integral.

Let κ be any positive integer prime to 11. Write 𝓜κ for the reduced germs of all integral formal local sections of L\^{}κ at zero. The following equality of integral modules is exact: \[
 \mathcal M_\kappa
   =\{\,h(u)F(u)^\kappa:
            h\text{ is an integral formal function on }\widehat J\,\}.
 \tag{NF1}
\] Here h(u) means substitution of the characteristic-zero group logarithm into an ordinary integral formal function. It does not mean that F is an integral ordinary power series in logarithmic coordinates.

\textbf{Proof.} The difference between L and the line of F is the rigidified two-torsion line \(\mathcal P_T\). Its pullback by {[}2{]} has its canonical rigidified trivialization. The {[}2{]}-pullback of F is a nonvanishing integral section near zero, so its κ-th power is an integral local frame for {[}2{]}*L\^{}κ under this trivialization. Hence any integral formal section S satisfies \[
 F_S(2u)=r(u)F(2u)^\kappa
\] for an arbitrary integral formal function r, and every such r occurs. The comparison has no exponential factor: this is the parallel rigidified trivialization of the pulled-back flat two-torsion line, exactly as in AM-K015. Since {[}2{]} is an integral formal automorphism at 11, replacing u by u/2 identifies these r with all integral formal functions h. This proves (NF1).

Every iterated invariant derivative of h has integral value at zero, because invariant vector fields preserve the integral formal coordinate ring. Thus h has integral Hurwitz coefficients. Products and inverses of Hurwitz series with integral coefficients and unit constant term are integral, as follows from the binomial product rule and the recursive coefficient equation for an inverse. Therefore every member of 𝓜κ is Hurwitz integral. This statement does not assert ordinary Taylor integrality. ∎

In particular the finite functional matrix B in AM-K022 can be taken integral at 11. Its presentation there as a matrix over Frac(R) was deliberately more general; changing its integral formal frame multiplies its jet columns by an integral invertible matrix and does not change its determinantal ideals.

\subsection{2. Universal restricted-operator relations}\label{universal-restricted-operator-relations}

For the ordinary coordinates u₁,u₂,u₃, AM-K017 proves both the algebraic restricted powers \[
 D_1^{11}=0,\qquad D_2^{11}=-D_1,\qquad D_3^{11}=0
\tag{NF2}
\] on the special fibre, and the full Hurwitz identities for F. The same identities on the completed local ring follow by continuity. For every f in 𝓜κ they imply \[
 \boxed{
 \partial_1^{11}f=4\kappa u_3f,\qquad
 \partial_2^{11}f=-\partial_1f,\qquad
 \partial_3^{11}f=-\kappa u_2f
 }\quad(\bmod11).
\tag{NF3}
\]

\textbf{Proof.} Write H=log F. AM-K017 gives \[
 \partial_1^{11}H=4u_3,\quad
 \partial_2^{11}H=-\partial_1H,\quad
 \partial_3^{11}H=-u_2.
\] The Hurwitz augmentation ideal has zero eleventh powers. Since the first derivatives of H have zero constant term, the prime Bell identity applied to κH gives (NF3) for F\^{}κ. For f=hF\^{}κ, the eleventh derivative of a product is a derivation modulo 11. The derivative of h is prescribed by (NF2). Thus the first and third formulas multiply by h; the middle one becomes \[
 (-D_1h)F^\kappa-h\,\partial_1F^\kappa
              =-\partial_1(hF^\kappa).
\] This proves (NF3) on the full module. In particular no division by the nonunit initial quadratic of sigma is used. ∎

Apply this to J² with ordinary coordinates u,v. Choose E′=Q(ζ₂₈) coordinates \[
 z_{i,+}=u_i+i v_i,\qquad z_{i,-}=u_i-i v_i.
\tag{NF4}
\] This is an integral change of coordinates at 11. The sign is the eigenvalue of the extra CM endomorphism \(I(u,v)=(-v,u)\). Its tangent derivations are \[
 \partial_{i,+}=(\partial_{u_i}-i\partial_{v_i})/2,\quad
 \partial_{i,-}=(\partial_{u_i}+i\partial_{v_i})/2.
\] Since \(i^{11}=-i\), the exact identities on every integral local section of the product line are \[
 \boxed{
 \partial_{1,\epsilon}^{11}f=2\kappa z_{3,\epsilon}f,\qquad
 \partial_{2,\epsilon}^{11}f=-\partial_{1,-\epsilon}f,\qquad
 \partial_{3,\epsilon}^{11}f=-\frac{\kappa}{2}z_{2,\epsilon}f.
 }\tag{NF5}
\] For example, the first right side is \((4\kappa/2)(u_3+i v_3)=2\kappa z_{3,+}\). The middle formula exchanges signs because the coefficient i is conjugated by Frobenius. The same formulas hold entry by entry on the actual Hom variety A\^{}\{ab\}=J\^{}\{2ab\}. A simultaneous unit normalization of coordinates gives the conjugate formulas, with unchanged matrix ranks and content.

\subsection{3. A finite normal form computes B modulo 11}\label{a-finite-normal-form-computes-b-modulo-11}

For each matrix entry introduce six commuting derivative symbols \(Z_{i,\epsilon}\). Let \(\Lambda_f(P)=[P(\partial)f]_0\). Equation (NF5) gives the following exact evaluation rules over k: \[
\begin{aligned}
 Z_{1,\epsilon}^{11}Q
     &\longmapsto 2\kappa\,\partial_{Z_{3,\epsilon}}Q,\\
 Z_{2,\epsilon}^{11}Q
     &\longmapsto -Z_{1,-\epsilon}Q,\\
 Z_{3,\epsilon}^{11}Q
     &\longmapsto -(\kappa/2)\partial_{Z_{2,\epsilon}}Q.
\end{aligned}\tag{NF6}
\] For instance \(Q(\partial)(z_jf)|_0
=(\partial_{Z_j}Q)(\partial)f|_0\). This is the ordinary Leibniz rule, including its exact integer multiplicities.

Apply (NF6) linearly until every individual exponent is less than 11. Each nonzero replacement lowers total degree, by 12, 10, or 12 respectively. Thus the process terminates. Let Tκ(P) denote its output, in the vector space \[
 \mathcal C=\operatorname{span}_k
   \{Z^\beta:0\leq\beta_j<11\text{ for every individual variable}\}.
\tag{NF7}
\] The procedure is confluent, and the following matrix statement is exact.

\subsubsection{Proposition NF8}\label{proposition-nf8}

For any finite chosen family P₁,\ldots,Pq of degrees at most D, the rank modulo 11 of its integral source jet matrix B equals the rank of the coefficient vectors \[
                    T_\kappa(P_1),\ldots,T_\kappa(P_q).
\tag{NF8}
\] The normalization factors \(M^{-\deg P_j}\) in AM-K022 are units and do not change this rank. For one polynomial P, its derivative image ideal has valuation zero if and only if \(T_\kappa(P)\ne0\). If its normal form is zero, the image is contained in 11R. The latter assertion alone gives no exact higher content.

\textbf{Proof.} All replacements preserve the evaluation on every member of 𝓜κ, by (NF6). The evaluation functionals associated with the cube monomials of degree at most D are linearly independent over k. Indeed choose integral algebraic formal parameters with the prescribed tangent basis. The degree-ordered matrix evaluating \(Z^\gamma\), \(0\leq\gamma_j<11\), on \(t^\beta F^\kappa\) has zero entry if \(|\gamma|<|\beta|\); in equal degree its entries are zero except at \(\gamma=\beta\), where they are \(\beta!\). These diagonal entries are units because every βj\textless11. Thus this square triangular submatrix is invertible.

Any two terminal reductions of P have identical functionals on all integral prefactors. Their difference lies in the cube and is therefore zero, proving confluence. The same independence proves (NF8). All the cube degrees encountered are at most the original degree, so the argument uses precisely the finite D-jet matrix. ∎

For several selected rows this computes exactly the number of Smith factors that are 11-adic units. It does not compute positive valuations of the remaining factors, and it does not iterate a congruence modulo 11 as a divisibility statement modulo higher powers.

Combining (NF8) with AM-K022, J30, gives the identical rank and unit-content criterion for the actual permitted half-orbit tuple matrix. Thus this calculation concerns actual allowed kernels, not an enlarged arbitrary Schwartz span.

\subsection{4. Closed calculation for rectangular selected rows}\label{closed-calculation-for-rectangular-selected-rows}

Take source rank a and target rank b\textgreater a. For each of the six labels use the first a columns, and write \(X_{i,\epsilon}\) for that a-by-a matrix of derivative symbols. Define \[
 P_{\mathbf n}=
 \prod_{\epsilon=\pm}\prod_{i=1}^3
                 \det(X_{i,\epsilon})^{n_{i,\epsilon}}.
\tag{NF9}
\] It is the primitive highest bideterminant for the rectangular partitions \((n_{i,\epsilon}^a)\), with determinant twists restored as in AM-DO-025.

Put \[
 C_a(n)=n(n+1)\cdots(n+a-1).
\tag{NF10}
\] The classical ordinary Cayley identity, in its integral polynomial specialization, is \[
 \det(\partial_X)\det(X)^n=C_a(n)\det(X)^{n-1}
 \quad(n\ge1),
\tag{NF11}
\] with zero on a constant. The primary reference is Caracciolo, Sokal, and Sportiello (2013), Theorem 2.1, equation (2.1); the directly consulted author preprint arXiv:1105.6270v2 gives this on printed p.~7. The short argument below checks the normalization used here; the identity is not a new result.

For completeness, over Q the left side is a relative invariant of the left GLₐ action with character det\^{}\{n−1\}, so on the dense invertible matrices it is a constant times det\^{}\{n−1\}. At X=I, the derivative indexed by a permutation π has value \(\operatorname{sgn}(\pi)n^{c(\pi)}\): each disjoint cycle can be assigned independently to one of the n determinant factors in \(\det(I+T)^n\), and its sign is the cycle sign. The signs in det(∂) cancel these signs. The resulting sum is \(\sum_\pi n^{c(\pi)}\), equal to \(n(n+1)\cdots(n+a-1)\) by inserting the last symbol either as a new cycle or at one of the a−1 positions in the old cycles. This proves (NF11) over Q and hence as an integral polynomial identity.

Since \(\det(X)^{11}=\det(X^{[11]})\) modulo 11, rules (NF6) and (NF11) give the following exact recurrences for the evaluated polynomial:

\begin{longtable}[]{@{}
  >{\raggedright\arraybackslash}p{(\columnwidth - 4\tabcolsep) * \real{0.3333}}
  >{\raggedright\arraybackslash}p{(\columnwidth - 4\tabcolsep) * \real{0.3333}}
  >{\raggedright\arraybackslash}p{(\columnwidth - 4\tabcolsep) * \real{0.3333}}@{}}
\toprule\noalign{}
\begin{minipage}[b]{\linewidth}\raggedright
Condition
\end{minipage} & \begin{minipage}[b]{\linewidth}\raggedright
Scalar factor
\end{minipage} & \begin{minipage}[b]{\linewidth}\raggedright
Exponent change
\end{minipage} \\
\midrule\noalign{}
\endhead
\bottomrule\noalign{}
\endlastfoot
\(n_{1,\epsilon}\ge11\) & \((2\kappa)^a C_a(n_{3,\epsilon})\) & \(n_{1,\epsilon}\mapsto n_{1,\epsilon}-11,\ n_{3,\epsilon}\mapsto n_{3,\epsilon}-1\) \\
\(n_{2,\epsilon}\ge11\) & \((-1)^a\) & \(n_{2,\epsilon}\mapsto n_{2,\epsilon}-11,\ n_{1,-\epsilon}\mapsto n_{1,-\epsilon}+1\) \\
\(n_{3,\epsilon}\ge11\) & \((-\kappa/2)^a C_a(n_{2,\epsilon})\) & \(n_{3,\epsilon}\mapsto n_{3,\epsilon}-11,\ n_{2,\epsilon}\mapsto n_{2,\epsilon}-1\) \\
\end{longtable}

A factor Cₐ(0) means that the result is zero. These are finite reduction identities for functionals, not identities of the original symbol polynomials before quotienting by (NF6).

\subsubsection{Corollary NF12. Exact unit-content examples beyond the factorial threshold}\label{corollary-nf12.-exact-unit-content-examples-beyond-the-factorial-threshold}

Suppose all six exponents equal n=11q+r, where \[
 1\le q\le r,\qquad r+q\le10,\qquad a+r\le11.
\tag{NF12}
\] Then the exact 11-content of this selected row is zero. More explicitly its nonzero normal form is \[
\begin{split}
 T_\kappa(P_{\mathbf n})={}&
 \kappa^{4aq}
 \left(\prod_{\ell=0}^{q-1} C_a(r-\ell)\right)^4\\
 &{}\cdot
 \prod_{\epsilon=\pm}
    \det(X_{1,\epsilon})^{r+q}
    \det(X_{2,\epsilon})^{r-q}
    \det(X_{3,\epsilon})^{r-q}.
\end{split}\tag{NF13}
\]

\textbf{Proof.} Reduce the two first-label exponents q times. Each third-label exponent becomes \(10q+r\). Reduce these q times; each middle-label exponent becomes \(10q+r\), and the third exponents become r−q. Reduce the middle exponents q times; they become r−q and each first-label exponent becomes r+q. The scalar factors in the table give (NF13); the signs and powers of 2 cancel. Every factor in every \(C_a(r-\ell)\) has an integer representative between 1 and 10, by (NF12), so the displayed scalar is nonzero modulo 11.

Each remaining determinant power has exponent at most 10. Therefore each individual symbol exponent is less than 11, and its polynomial lies in the cube. Its diagonal monomial has coefficient one, so the result is nonzero. Proposition NF8 proves zero content. ∎

Examples include n=12 for any even a≤10. For a=6, q=r=5 gives n=60 and total degree \(6an=2160\), yet the selected row has a unit value on some actual permitted tuple once M≥6482. The least permitted power is M=29³=24389, and its finite Schwartz vector still has norm one. The simple highest-symbol factorial bound in AM-K022 would only have given \(6a\,v_{11}(60!)=180\), so the normal form improves that finite bound to the exact value zero.

These rectangles have an explicit source HC specialization at the first edge. For any already admissible rank record with k=a and d=t+1, put H=(d+2)/2 and, at all six embeddings, \[
 b_j=j(d+2),\qquad e_j=H+j(d+2),\qquad
 u_i=a+n-\tfrac32-i,\qquad v_i=a+n-2-i.
\tag{NF14}
\] Every interlacing difference is 1/2, the two e,b gaps equal H, and the last u is n−1/2\textgreater−b₀=0. The integer/half-integer conditions hold. Since \(w_{0,i}=u_i-\rho_a(i)\), equation DO25.5 gives exactly \(\nu_{0,i}=n\). Thus these are actual selected first-edge partitions, not freely relabelled Schur weights. For example the HC data a=6,d=14,n=60 give the selected edge 6→20. This checks the HC interface; any separate ambient arithmetic field or global size conditions retain their own status.

There is also a basic unbounded-rank family: if all rectangular exponents are at most 10, the highest polynomial already lies in the cube, for every a. Its row has zero content regardless of its total degree, and (NF14) realizes these first-edge HC weights whenever an admissible rank record is supplied. Neither family by itself controls every later stage of one fixed growing chain.

\subsection{5. Scope and sources}\label{scope-and-sources}

The novel step is (NF1)--(NF8): the actual K017 relations extend to every integral local section, giving a rigorous finite normal form for B rather than a congruence for only one theta function. NF12--NF13 are the new selected-row application inside that finite quotient; NF11 is the classical Cayley identity with its integral normalization explicitly checked.

Caracciolo, S., Sokal, A. D., \& Sportiello, A. (2013). Algebraic/combinatorial proofs of Cayley-type identities for derivatives of determinants and pfaffians. Advances in Applied Mathematics, 50(4), 474--594. Primary author preprint: \url{https://arxiv.org/pdf/1105.6270}, version 2 (31 December 2012), Theorem 2.1 and equation (2.1), printed p.~7. Only the ordinary determinant identity and its normalization are used here.

Carried inputs are AM-K012, AM-K015, AM-K017, AM-K018, ALP-LIFT-07--08, and AM-K022; their exact primary/source locators are in those reports. The actual E′ coordinates and selected Cauchy decomposition are AM-DO-022 and AM-DO-025, especially DO25.3 and DO25.5--9. No extra claim about reduced ordinary Taylor integrality, bounded differential-module saturation, or a later automorphic class matrix is used.

The literal HC check (NF14) uses \nolinkurl{Formalism/LocalRingCardinality/Params.lean,} lines 50--69, and the actual selected-weight identity DO25.5. The finite independent arithmetic certificate is \nolinkurl{selected_jet_normal_form_certificate.py,} with output in selected\_jet\_normal\_form\_certificate.json. It checks the integer determinant-two Cayley identity for n=1,\ldots,40; 250 one-entry full normal forms; the complete six-matrix determinant-two expansions at n=12 and n=24 (4096 and 531441 nonzero input monomials modulo 11); a vanishing example at n=23; and the rank-six n=60 level arithmetic. It does not infer any infinite or higher-content conclusion from those finite checks.

The full finite jet matrix is integral and its rank modulo 11 is now computable from finite polynomial operations. Positive Smith valuations, control as every selected weight grows, and the effect of restricting the next class matrix to a previously produced theta image remain separate problems.

\clearpage

\section{AM-K024 --- a scalable sharp upper bound in the actual ordinary E′ selected rows}\label{am-k024-a-scalable-sharp-upper-bound-in-the-actual-ordinary-e-selected-rows}

Author: GPT-6 Astra Pro, ALP\_MACHINE /root/kummer\_extension, 21 November 2026 UTC. Inherited reasoning effort is not separately exposed.

This is a fixed-rank, varying-HC-family statement for the actual ordinary host A=J² with E′=Q(ζ₂₈). It applies to the selected compact Cauchy summands identified in AM-DO-025. It uses the exact primitive-kernel coefficient from accepted DP58. It does not introduce conjugate-action J factors, and it does not identify the output of a later automorphic class integral.

The result gives both a bound for the finite jet functional matrix B and a single allowed finite Schwartz input realizing that bound. It excludes a uniform local slope larger than 1/1220 across the indicated source-selected family. There is no net global cost conclusion.

\subsection{1. Fixed selected edge and exact statement}\label{fixed-selected-edge-and-exact-statement}

Fix an admissible source rank record and a stage j with \[
 a=N_j,\qquad b=N_{j+1}=a+d,\qquad d=t+1\ge2.
\tag{SC1}
\] Fix word-independent source HC data at each of the six holomorphic embeddings of E′. Let \(\nu_\tau^*\) be the resulting actual selected partitions on this edge. They are positive antitone partitions of length a by AM-DO-025. Put \[
 D_0=\sum_\tau|\nu_\tau^*|,\qquad
 \mathcal F_0=\prod_{\tau}\prod_{i=1}^{a}\nu_{\tau,i}^*!,\qquad
 c_0=v_{11}(\mathcal F_0).
\tag{SC2}
\] All these data are fixed as r varies.

Label the six embeddings by \((i,\epsilon)\), \(i=1,2,3\), \(\epsilon=\pm\), using \(z_{i,\pm}=u_i\pm i v_i\) on A=J². These are the actual E′ eigencharacters, with the underlying J character i. Set \[
 A_r=11^r,\qquad B_r=121\,11^r,\qquad
 T_r=A_r+B_r+D_0=122\,11^r+D_0.
\tag{SC3}
\]

\subsubsection{Theorem SC4}\label{theorem-sc4}

For each r≥0 there are integers \(0\le h_r\le A_r\), \(0\le k_r\le B_r\) and actual source HC data whose selected partitions on this edge are \[
\begin{array}{ll}
 \nu_{3,+}=\nu_{3,+}^*+h_r e_1,&
 \nu_{3,-}=\nu_{3,-}^*+(A_r-h_r)e_1,\\
 \nu_{2,+}=\nu_{2,+}^*+k_r e_1,&
 \nu_{2,-}=\nu_{2,-}^*+(B_r-k_r)e_1,\\
 \nu_{1,+}=\nu_{1,+}^*,&
 \nu_{1,-}=\nu_{1,-}^*.
\end{array}\tag{SC4}
\] There is a primitive integral polynomial \(P_r\) in the exact source-and-target Cauchy summand for these selected types whose degree is Tᵣ and whose normalized source jet functional has image-ideal valuation at most \[
                         \frac{11^r-1}{10}+c_0.
\tag{SC5}
\]

Choose the least M=29\^{}e≥3Tᵣ+2 and κ=M². A single tuple of the actual DP59 half-orbit then satisfies \[
 a_r=M^{-T_r}[P_r(\partial)F_{\mathcal S_{\mathbf t'}}]_0\ne0,
 \qquad
 v_{11}(a_r)\le\frac{11^r-1}{10}+c_0.
\tag{SC6}
\] Its finite Schwartz vector has norm one; no interpolating coefficient height is used. Its actual identity theta-kernel value is \(C^{2ab}a_r\), with the one fixed period \(C=\theta_0c\).

Consequently, for any γ\textgreater1/1220 and fixed constant C₀, the lower bound \[
              v_{11}(a)\ge\gamma T-C_0
\tag{SC7}
\] cannot hold for every allowed selected local coefficient in this fixed-rank HC family. This assertion concerns the union of the explicitly displayed selected types. It does not assert a bound for every vector in one fixed type or one fixed later theta image.

\subsection{2. The source HC change is literal}\label{the-source-hc-change-is-literal}

At each embedding, increase only the initial entry u₀ of the fixed source HC record by an arbitrary nonnegative integer ητ. Leave its other u entries, every v entry, and every e,b entry fixed.

In \nolinkurl{Formalism/LocalRingCardinality/Params.lean,} lines 50--69, the only interlacing condition involving u₀ is \(u_0>v_0\). It remains true after this increase. The last-u condition is unaffected because k≥2; the half-integer condition remains true; the e,b gaps and integer conditions are unchanged. Thus this is an actual HC record. Choosing the same change for every word preserves the source's word-compatibility equations.

The initial source weight is \(w_{0,i}=u_i-\rho_k(i)\). Increasing u₀ by ητ increases its first entry by ητ and changes no other entry. The exact recurrence DO25.6 copies this change into the first entry of every subsequent source weight. Formula DO25.5 therefore gives \[
                     \nu_{j,\tau}=\nu_{j,\tau}^*+\eta_\tau e_1.
\tag{SC8}
\] Taking the four η's in (SC4) gives precisely those selected partitions. Their dominance is preserved by increasing the first entry.

All e entries remain unchanged, so the actual common source-and-target character on every edge is unchanged. This is stronger than freely choosing six unrelated partitions. It does not, however, preserve a size window that independently bounds u₀. Such a window must grow with the indicated HC changes.

\subsubsection{Full analytic Hyp3 is preserved when the fixed baseline satisfies it}\label{full-analytic-hyp3-is-preserved-when-the-fixed-baseline-satisfies-it}

Suppose additionally that the fixed baseline TPar satisfies the source's stronger Hyp3 in CayleySignature.lean, lines 110--118, including the condition \(1801\le r+t\). Then every u₀ modification used above satisfies that same Hyp3.

To check every clause, the 1801 condition and the parity and eta clauses are unchanged. Interlace was checked above. DataTheta.lean, TPar.Dpar, lines 538--544, depends only on v,b,e,t,eta, so every discrete parameter set and therefore every NCGaps inequality is literally unchanged. In TPar.Centred, lines 550--554, only the inequalities involving u₀ can change. Interlacing implies \(u_i\ge u_{\mathrm{last}}>-b_0\), and \(e_j\ge e_0>b_0\): the e,b entries alternate strictly increasingly. Thus \(u_i+e_j>0\). Increasing u₀ consequently increases its absolute value \(|u_0+e_j|\); its required lower bound is preserved. All other Centred inequalities are unchanged.

This proves full analytic Hyp3 preservation, as well as preservation of IsMu because e is fixed. Hence SC4--SC20 apply within a genuinely Hyp3-admissible analytic parameter family whenever one such baseline is supplied. They still do not preserve separately fixed SizeData₂ upper windows or identify an iterated arithmetic row.

\subsection{3. A primitive source-highest polynomial with one spare column}\label{a-primitive-source-highest-polynomial-with-one-spare-column}

Use an a-by-b symbol matrix \(Z_\tau\) for each label. Reserve column 1. On columns 2,\ldots,a+1 define the residual prefix-minor polynomial \[
 Q_0(Z)=
 \prod_\tau\prod_{\ell=1}^{a}
 \det\bigl(Z_{\tau,\{1,\ldots,\ell\},
                         \{2,\ldots,\ell+1\}}\bigr)^{
                           \nu_{\tau,\ell}^*-\nu_{\tau,\ell+1}^*},
 \qquad \nu_{\tau,a+1}^*=0.
\tag{SC9}
\] These columns exist because b≥a+1. The same columns can be reused for the six labels: they are independent sets of variables.

For \(0\le h\le A_r\), \(0\le k\le B_r\), put \[
\begin{split}
 P_{r,h,k}={}&
 Z_{3,+,1,1}^{h}
 Z_{3,-,1,1}^{A_r-h}
 Z_{2,+,1,1}^{k}
 Z_{2,-,1,1}^{B_r-k}\,Q_0(Z).
\end{split}\tag{SC10}
\] Every factor in Q₀ is invariant when an earlier row is added to a later row. A monomial using only row 1 retains this invariance. The row degrees are exactly (SC4), with h,k in place of hᵣ,kᵣ. The product of the residual diagonal terms and the displayed first-column monomial has coefficient one. Hence Pᵣ,h,k is primitive in the polynomial coefficient lattice.

\subsubsection{Exact common source-and-target Schur label}\label{exact-common-source-and-target-schur-label}

The row-invariance assertion places this polynomial in the source-highest subspace of its indicated row weight. It also determines the target representation. Indeed, for each embedding the actual Cauchy decomposition is \[
 \operatorname{Sym}(T\otimes S^*)
       =\bigoplus_\mu S_\mu(T)\otimes S_\mu(S)^*.
\tag{SC11}
\] In a source irreducible \(S_\mu(S)^*\), the specified maximal-unipotent invariant line has the source highest weight corresponding to μ, and no other such invariant weight. Since the partitions label pairwise distinct source highest weights, the invariant subspace of the row weight ν in (SC11) occurs only in its μ=ν summand. Thus \[
 P_{r,h,k}\in
 \bigotimes_\tau
      S_{\nu_\tau}(T_\tau)\otimes
              (\text{source highest line in }S_{\nu_\tau}(S_\tau)^*).
\tag{SC12}
\] Restoring the fixed determinant twists in DO25.3 gives exactly the contragredient selected source type and the selected target type DO25.7--9.

The target vector need not be its highest vector, but it lies wholly in that same irreducible target representation. It does not split among different target Schur labels. No rational projector or projection denominator is introduced: the integral polynomial itself is already in the required Cauchy summand. A demand for one separately fixed target-highest coefficient would be an additional restriction, not part of (SC12).

This argument is over characteristic zero, where the source comparison lives. Its integral assertion is the explicit polynomial with coefficients in R. Any further comparison with a differently normalized abstract Weyl pairing retains the normalization factors studied in the parent's selected-Schur metric report; none is set to one here.

\subsection{4. The actual primitive-kernel input and the exact filler factorial}\label{the-actual-primitive-kernel-input-and-the-exact-filler-factorial}

The Hom variety is \(A^{ab}=J^{2ab}\) with its ordinary J factors. On the first ordinary J factor at row 1, column 1, let X₃ and X₂ be the differential symbols for the two primitive integral directions of actual E-labels 3 and 2. They satisfy \[
             X_3=Z_{3,+,1,1}+Z_{3,-,1,1},\qquad
             X_2=Z_{2,+,1,1}+Z_{2,-,1,1}.
\tag{SC13}
\] This follows from the exact coordinate change z±=u±iv. Any unit changes to the DP49 unit-arrow display coordinates preserve the valuations below.

Accepted DP58, equations ES11--12 and the mixed monomial evaluation recorded in DP65 EH16--17, gives the actual unit-parameter primitive kernel in these ordinary directions: \[
 v_{11}\!\left(
 \partial_{X_3}^{A_r}\partial_{X_2}^{B_r}
                   F_{-\kappa}(0)\right)
                         =\frac{11^r-1}{10}.
\tag{SC14}
\] The parameter −κ is a unit for every allowed M. AM-K018 matches the full-O\_E source gauge on these ordinary J factors. DP54 and DP59 give the integral unit algebraic prefactor that cancels the local frame unit, so (SC14) is an actual permitted integral-local-section functional value. The pure primitive kernel is not substituted for the fixed sigma section without this exact prefactor comparison.

On the filler columns choose the ordinary integral algebraic monomial whose linear logarithmic term is \[
              \prod_{\tau,i}z_{\tau,i,i+1}^{\nu_{\tau,i}^*}.
\tag{SC15}
\] It has degree D₀. The residual highest polynomial Q₀ has matching monomial coefficient one; every other supported monomial has a different exponent. Since its degree is D₀, higher terms of the coordinate changes or of every reduced frame cannot contribute. Therefore the filler evaluation is exactly \[
                    Q_0(\partial)(\text{filler germ})|_0
                                      =\mathcal F_0.
\tag{SC16}
\] No factorial is inverted. The underlying change between the ordinary J coordinates and the E′ coordinates has unit determinant at 11. Thus the prefactor in (SC15) is integral in the actual formal coordinate lattice.

Take the product of the first-column primitive-kernel prefactor, these filler monomials, and normalized unit prefactors on unused entries. It is an actual integral formal local section of \(\mathcal L^\kappa\). Factorization between the first column and filler columns gives \[
 v_{11}\!\left(
 M^{-T_r}
 [X_3^{A_r}X_2^{B_r}Q_0(\partial)F_S]_0
 \right)
                         =\frac{11^r-1}{10}+c_0.
\tag{SC17}
\] The displayed \(M^{-T_r}\) is a unit. This is a construction in the finite jet module B; globalizing it is the separate integral interpolation step of AM-K022.

\subsection{5. Weight extraction without a projection denominator}\label{weight-extraction-without-a-projection-denominator}

Expanding (SC13) gives the exact identity of differential polynomials \[
 X_3^{A_r}X_2^{B_r}Q_0
    =\sum_{h=0}^{A_r}\sum_{k=0}^{B_r}
          \binom{A_r}{h}\binom{B_r}{k}P_{r,h,k}.
\tag{SC18}
\] Every coefficient is an integer. Apply this equality to the single integral formal section just constructed. If every nonzero component value had valuation greater than the right side of (SC17), their integral linear combination would have greater valuation too. Consequently at least one pair (hᵣ,kᵣ) has \[
 v_{11}\!\left(
  M^{-T_r}[P_{r,h_r,k_r}(\partial)F_S]_0
          \right)
                         \le\frac{11^r-1}{10}+c_0.
\tag{SC19}
\] Terms whose binomial coefficient is divisible by 11 cause no difficulty: the argument only uses integrality, and does not divide by a binomial coefficient. This proves the image-ideal bound (SC5) for an actual selected row. It does not assert equality or a matching lower bound for that row.

Now use AM-K022 with D=Tᵣ. Its allowed half-orbit span realizes every integral algebraic Tᵣ-jet for M≥3Tᵣ+2. Equivalently, J30 identifies its derivative image ideal with that of the finite jet module. Expanding one realizing section in its integral allowed basis and applying the nonarchimedean valuation inequality chooses an individual tuple obeying (SC6). The canonical tuple consists of global algebraic points and sections by carried DP59(c), even when the interpolating coefficients were complete-local.

The finite Schwartz vector is a single Heisenberg translate of the fixed norm-one vacuum. Hence its norm is one, and its translation coefficient vector is a standard basis vector. The actual all-polynomial identity DO26.4 gives the fixed period \(C^{2ab}\), with no degree-dependent replacement.

Finally \[
 \frac{(11^r-1)/10+c_0}{122\,11^r+D_0}
                              \longrightarrow\frac1{1220}.
\tag{SC20}
\] This proves the failure of every larger uniform slope in (SC7).

\subsection{6. What this closes and what it leaves}\label{what-this-closes-and-what-it-leaves}

The old DP65 construction used both ordinary and conjugate-action J factors in a generalized harmonic module. The present proof uses only ordinary J powers in the actual E′ Hom variety and proves membership in the precise compact source-and-target selected summand. The full-O\_E gauge comparison from AM-K018 therefore applies directly. The first-row growth is implemented inside the actual source HC record; the common edge character is unchanged.

The bound is for actual selected \textbf{kernel coefficients} in the permitted half-orbit. At the first edge, the parent's class-supported construction AM-R-012 can use such a nonzero coefficient, retaining its source mass, compact norm, selected-polynomial norm and common period. At a later edge, the input is already constrained to a previous automorphic theta image. The restricted class matrix may cancel these coefficients or avoid these witnesses. Neither (SC19) nor the exact finite saturation theorem identifies that image or preserves this value after integration.

Likewise, varying the first HC gaps grows the source size window. The archimedean Schur norm, source class mass, denominator normalizations and the global height identity remain separate costs. The local upper slope cannot by itself be promoted to a statement about the net error or gain of the full Alpöge machine.

\subsection{7. Exact input locators}\label{exact-input-locators}

\begin{itemize}
\tightlist
\item
  Accepted DP58, research/explicit\_sharp\_jets\_linear\_levels\_DP58.md, §3, ES11--12 and the mixed monomial evaluation; accepted DP65, research/embedding\_labelled\_harmonic\_witness\_DP65.md, §5, EH16--17, for the explicitly recorded evaluation of that same carried primitive kernel.
\item
  Accepted DP54, research/global\_theta\_finite\_jets\_DP54.md, and DP59, research/theta\_translate\_span\_DP59.md, §§6--8: the integral unit frame prefactor, its finite-jet realization and the actual allowed translate lattice.
\item
  AM-K018: full-O\_E source gauge match on ordinary J products; AM-K022: explicit M≥3D+2 threshold, factorial convention, permitted individual tuples, and equality of finite determinantal ideals.
\item
  AM-DO-022: ordinary E′ host and prime-to-11 coordinate change. AM-DO-025, DO25.3--9: exact determinant twists, Cauchy decomposition and selected source/target labels.
\item
  \nolinkurl{Formalism/LocalRingCardinality/Params.lean,} HCParam, lines 50--69: literal u₀ variation. \nolinkurl{Formalism/LocalToGlobal/Columns/AbelianHostFieldConditions.lean,} lines 462--469, and Columns/BaseLemmas.lean, lines 496--552: actual selected-weight identities, carried as in AM-DO-025.
\item
  \nolinkurl{Formalism/LocalRingCardinality/DataTheta.lean,} TPar.Interlace, Dpar and Centred, lines 529--554; \nolinkurl{Formalism/SharpRestriction/CuspSpanNormBounds/CayleySignature.lean,} NCGaps and the stronger Hyp3, lines 110--118: exact preservation of the full baseline analytic admissibility conditions.
\end{itemize}

All new reductions are proved above. No Lean execution, full later-chain realization, or finite computational certificate is claimed by this report.

\clearpage

\section{AM-K025 --- an adjacent-gap factorial bound for the actual highest bideterminant}\label{am-k025-an-adjacent-gap-factorial-bound-for-the-actual-highest-bideterminant}

Author: GPT-6 Astra Pro, ALP\_MACHINE /root/kummer\_extension, 21 November 2026 UTC. Inherited reasoning effort is not separately exposed.

\textbf{Scope.} The general bound below applies to every selected partition in the ordinary E′ construction. Its zero-content growing-chain example satisfies the literal HC equations only; it does not assert the additional Hyp3, NCGaps, or global SizeData₂ inequalities. Those stronger hypotheses may force much larger adjacent gaps. The result is a local kernel image-ideal bound before restriction to a previous automorphic theta image.

\subsection{1. The bound uses adjacent differences, not the row lengths}\label{the-bound-uses-adjacent-differences-not-the-row-lengths}

Fix a source rank a, target rank b≥a, and the six actual E′ holomorphic labels τ. Let \(\nu_\tau\) be nonnegative antitone partitions of length a. Put \[
 \delta_{\tau,n}=\nu_{\tau,n}-\nu_{\tau,n+1}\ge0,
 \qquad \nu_{\tau,a+1}=0,
\tag{AG1}
\] and use the actual primitive highest bideterminant \[
 P_\nu=\prod_\tau\prod_{n=1}^a
                  \det(Z_{\tau,[1,n],[1,n]})^{\delta_{\tau,n}},
 \qquad D=\sum_\tau|\nu_\tau|.
\tag{AG2}
\] Both its source and target vectors are highest vectors in the exact selected Cauchy summand. No projection onto either compact type is required.

For M=29\^{}e≥3D+2, let \(I_\nu\) be the image ideal of the normalized derivative functional \(M^{-D}P_\nu(\partial)\) on the actual integral formal section lattice, equivalently on the actual permitted DP59 half-orbit by AM-K022 J30. Then \[
 \boxed{
 0\le v_{11}(I_\nu)
       \le \sum_{\tau}\sum_{n=1}^a
                                n\,v_{11}(\delta_{\tau,n}!).
 }\tag{AG3}
\] Equivalently, the upper bound is \[
 \frac1{10}\sum_{\tau,n}
          n\bigl(\delta_{\tau,n}-s_{11}(\delta_{\tau,n})\bigr).
\tag{AG4}
\] If every adjacent difference is at most 10, then \(I_\nu=R\): its exact 11-content is zero, and an individual allowed tuple has unit value.

\textbf{Proof.} For one fixed label and one prefix size n, use the anti-diagonal determinant term \[
       (-1)^{n(n-1)/2}\prod_{i=1}^n Z_{i,n+1-i}.
\tag{AG5}
\] All these choices are simultaneously uniquely leading terms. To see this, choose a constant C\textgreater ab and assign positive variable weights \(w_{ij}=C-ij\). A term indexed by a permutation π has weight \(nC-\sum_i i\pi(i)\). Its maximum is uniquely achieved by the reversing permutation: if i\textless j but π(i)\textless π(j), swapping those two images decreases the sum by \((j-i)(\pi(j)-\pi(i))>0\). Repeated swaps give the unique reversal. Thus the product of the anti-diagonal terms is the uniquely leading term of (AG2) and has coefficient ±1.

Different prefix sizes use disjoint anti-diagonals, because their variables satisfy different equations i+j=n+1. Consequently the leading monomial has exponent exactly \(\delta_{\tau,n}\) on each of the n variables of that anti-diagonal. Its exponent factorial is therefore \[
                   \beta!=\prod_{\tau,n}
                                      (\delta_{\tau,n}!)^n.
\tag{AG6}
\] Apply the highest-symbol interpolation AM-K022 J12 to this monomial. It gives \(\pm\beta!\in I_\nu\), with no factorial inverse or reduced-frame integrality assumption. This proves the upper bound. AM-K023 NF1 proves Hurwitz integrality of the entire actual integral local section module at 11, so \(I_\nu\subset R\), giving the lower bound. If every δ≤10, (AG6) is a unit. The integral allowed basis then contains an individual tuple with unit evaluation. ∎

This can be substantially smaller than the diagonal-monomial bound \(\sum_{\tau,i}v_{11}(\nu_{\tau,i}!)\). The improvement is an exact monomial calculation inside the same highest bideterminant, with the same selected source and target types.

\subsection{2. Literal minimal HC chain: exact unit content at every stage}\label{literal-minimal-hc-chain-exact-unit-content-at-every-stage}

This section is \textbf{HC-only}. It does not claim admissibility under the stronger 1801 or NCGaps hypotheses.

Fix an already admissible rank record with \(k\) even, \(d=t+1\) even, and put H=(d+2)/2. At every holomorphic embedding use \[
 b_j=j(d+2),\qquad e_j=H+j(d+2),\qquad
 u_i=k-\tfrac12-i,\qquad v_i=k-1-i.
\tag{AG7}
\] The two interlacing gaps are 1/2; the last u is 1/2\textgreater−b₀; and both e,b gaps equal H. All required integer or half-integer conditions hold. Formula DO25.5 gives \(\nu_0=(1^k)\), while the literal recurrence DO25.11 has \(\Delta_j=e_{j+1}-e_j-d=2\). Therefore \[
 \nu_j=((2j+1)^k,(2j)^d,(2j-2)^d,\ldots,2^d)
\tag{AG8}
\] at each label. Its nonzero adjacent differences are \[
       \delta_k=1,\qquad
       \delta_{k+hd}=2\quad(1\le h\le j).
\tag{AG9}
\] Thus every selected highest bideterminant in this chain has exact local content zero at 11.

For source rank a=k+jd and target rank b=a+d, its total six-label degree and an exact interpolated value are \[
 D_j=6\{k(2j+1)+d\,j(j+1)\},
\tag{AG10}
\] \[
 M^{-D_j}P_{\nu_j}(\partial)F_S(0)
       =2^{\,6kj+3d\,j(j+1)}
\tag{AG11}
\] for a suitable integral formal section S; the six identical determinant signs multiply to +1. The level can be the least power 29 at least \(3D_j+2\). An individual permitted tuple has unit value, although its value need not be the specific power of 2 in (AG11). Its finite Schwartz norm is one.

The nonzero degrees grow quadratically with j. This rules out a positive local lower bound depending only on total selected degree across all HC records. It leaves open the more restrictive family forced by the global machine, and it does not imply that these unit coefficients survive any later class integral.

\subsection{3. The universal finite cutoff clarifies this scope}\label{the-universal-finite-cutoff-clarifies-this-scope}

AM-K023 NF5 implies, for every integral formal section on each A entry, \[
 \partial_{1,\pm}^{121}f=0,\qquad
 \partial_{3,\pm}^{121}f=0,\qquad
 \partial_{2,\pm}^{1331}f=0
                       \quad(\bmod11).
\tag{AG12}
\] Indeed iterating the first or third eleventh-power identity eleven times gives a multiple of \(z^{11}=0\) in the Hurwitz algebra; its coordinate multiplier commutes with that derivative. The middle identity then reduces its 1331st power to a 121st first-label power.

Consequently a nonzero evaluated derivative modulo 11 on A\^{}\{ab\} has, in each entry, exponents at most 120,120 on the first labels, 1330,1330 on the middle labels, and 120,120 on the third labels. Their sum is 3140 per entry. Therefore \[
             D>3140ab
 \quad\Longrightarrow\quad
 B(P)\subset11R
\tag{AG13}
\] for every homogeneous polynomial P of degree D, including every selected row. This is a uniform first divisibility statement for the full integral local section module. It cannot be iterated after dividing by 11 without a further stable-lattice argument.

There is no conflict with (AG8): its degree is far below this dimensional cutoff for the indicated ranks. More generally the adjacent-gap criterion and the cutoff give finite, independently checkable tests that any proposed admissible HC specialization must pass.

\subsection{4. Sources and status}\label{sources-and-status}

The finite interpolation and permitted span are AM-K022, using carried DP59. The integral local source gauge and Hurwitz module are AM-K018 and AM-K023. The exact selected partitions and determinant twists are AM-DO-025. The HC clauses are \nolinkurl{Formalism/LocalRingCardinality/Params.lean,} lines 50--69.

All new claims are proved above by the leading-monomial argument and the already established operator identities. No new computational certificate or Lean execution is asserted. The general bound (AG3) remains valid for the larger source gaps; only the special unit family is restricted to the stated HC scope.

\clearpage

\section{ALP-LIFT-16 --- fixed spherical tests, 29-supported single translates, and exact level masses}\label{alp-lift-16-fixed-spherical-tests-29-supported-single-translates-and-exact-level-masses}

Author: GPT-6 Astra Pro, ALP\_MACHINE /root/lift\_accounting, 21 November 2026 UTC. Inherited reasoning effort is not separately exposed.

This result sharpens ALP-LIFT-15 in the actual common \(E'=\mathbf Q(\zeta_{28})\) host. The later finite inputs can all be single norm-one Heisenberg translates supported only at29. The proof includes an exact split-place spherical-vacuum calculation valid in residue characteristic2. The translation depths remain unbounded, and the prescribed raw Xi and prior-image arithmetic lattice are not identified.

The additional quantitative output is an explicit sufficient principal level and its class mass. Its 11-adic valuation is independent of the unknown 29-depth. This does not compute a later gain: the arithmetic entries and their cancellations are separate.

\subsection{1. The split-place vacuum coefficient, with the actual measure}\label{the-split-place-vacuum-coefficient-with-the-actual-measure}

Let \(F\) be any nonarchimedean local field of characteristic zero, with residue cardinality \(q\), uniformizer \(\varpi\), and ring \(\mathcal O\). Put \[
 G=\operatorname{GL}_a(F),\quad K=\operatorname{GL}_a(\mathcal O),
                    \qquad b>a.
 \tag{S1}
\] Normalize Haar measure on \(G\) by \(\operatorname{vol}(K)=1\). At a split place of the common unitary host, AM-DO-020/022 identifies the actual normalized finite vacuum with \(\phi_0=\mathbf1_{M_{b,a}(\mathcal O)}\) after the unitary measure change determined by its lattice. Its \(L^2\) norm is one. The source action, without its determinant-root character, is \[
       \omega(g)\phi(X)=|\det g|^{b/2}\phi(Xg).
 \tag{S2}
\] For a general lattice measure, the normalized characteristic vector is \(\operatorname{vol}(L)^{-1/2}\mathbf1_L\); its matrix coefficient is unchanged by the corresponding unitary identification. Thus no additive-lattice volume is lost in the normalization used here.

Write \(a_\lambda=\operatorname{diag}(\varpi^{\lambda_1},\ldots,
\varpi^{\lambda_a})\), with \(\lambda_1\ge\cdots\ge\lambda_a\). A direct intersection-of-lattices calculation gives \[
 A_b(g):=\langle\omega(g)\phi_0,\phi_0\rangle
       =q^{-b\sum_i|\lambda_i|/2}
                  \quad(g\in Ka_\lambda K).
 \tag{S3}
\] Indeed the \(i\)-th column lies in \(\mathcal O^b\cap\varpi^{-\lambda_i}\mathcal O^b\), of volume \(q^{-b\max(-\lambda_i,0)}\), while the determinant factor is \(q^{-b\sum_i\lambda_i/2}\). The formula holds at2 as well as at odd residue characteristic. It uses the actual split Schrödinger action, not an assumed integral diagonal/antidiagonal basis for a doubled space.

\subsection{2. A matrix-disintegration proof of the exact spherical transform}\label{a-matrix-disintegration-proof-of-the-exact-spherical-transform}

Let \(\pi\) be an irreducible unramified unitary representation of \(G\), let \(m_\pi(g)\) be its normalized spherical matrix coefficient, and let \(\alpha_1,\ldots,\alpha_a\) be its standard normalized Satake parameters. Define \[
 C_{a,b}(q)=\prod_{i=0}^{a-1}(1-q^{i-b}),\qquad
                           s_*=(b-a+1)/2.
 \tag{S4}
\] We will prove, in the convergence range specified below, \[
 \boxed{\int_G A_b(g)m_\pi(g)\,dg
        =C_{a,b}(q)L(s_*,\pi)L(s_*,\pi^\vee).}
 \tag{S5}
\] For a unitary representation the two \(L\)-values at a real point are complex conjugates. Thus the right side is strictly positive whenever the factors are finite and nonzero.

\subsubsection{2.1 Exact full-rank matrix disintegration}\label{exact-full-rank-matrix-disintegration}

For an integer \(B\ge a\), put \[
             f_B(h)=\mathbf1_{M_a(\mathcal O)}(h)
                                           |\det h|^{B/2}.
\] Let \(X_0=\binom{I_a}{0}\in M_{B,a}(F)\). Every full-rank matrix can be written \(kX_0h\), with \(k\in\operatorname{GL}_B(\mathcal O)\), \(h\in G\). To see this, intersect its column space with \(\mathcal O^B\). This intersection is a saturated rank-\(a\) direct summand over the DVR, so an integral basis extends to a basis of \(\mathcal O^B\); its change-of-basis matrix is the required \(k\).

With probability Haar measure \(dk\) on \(\operatorname{GL}_B(\mathcal O)\), matrix Lebesgue measure satisfies \[
 \int_{M_{B,a}(F)}F(X)\,dX
  =C_{a,B}(q)\int_{\operatorname{GL}_B(\mathcal O)}
        \int_G F(kX_0h)|\det h|^B\,dh\,dk.
 \tag{S6}
\] The rank-deficient set has measure zero, as follows by induction and Fubini from the vanishing of a nonzero matrix minor. Both sides on the full-rank orbit are left compact-invariant and transform under right multiplication by \(h_0\in G\) with the same Jacobian \(|\det h_0|^{-B}\). Uniqueness of a relatively invariant measure on this transitive homogeneous space makes them proportional. The proportionality constant is fixed by the open set of matrices whose reduction has rank \(a\): in the right-hand coordinates this is exactly \(h\in K\), so the integral without the constant is one. Its additive measure is \[
 \frac{\prod_{i=0}^{a-1}(q^B-q^i)}{q^{Ba}}
                 =\prod_{i=0}^{a-1}(1-q^{i-B}).
\] This proves (S6), including its constant. The homogeneous-space measure assertion can equivalently be checked in a local minor chart; the same Jacobian and normalization give (S6).

Using (S6) on the product of two characteristic functions gives \[
       A_B(g)=C_{a,B}(q)\int_G f_B(hg)f_B(h)\,dh
             =C_{a,B}(q)(\check f_B*f_B)(g),
              \qquad\check f_B(h)=f_B(h^{-1}).
 \tag{S7}
\] The inverse-argument operator \(|\det g|^{-B/2}\phi(Xg^{-1})=\omega(g^{-1})\phi(X)\) has the same vacuum coefficient, since \(A_B(g^{-1})=A_B(g)\) by (S3). It is not the left action in(S2).

\subsubsection{2.2 The one-sided spherical zeta integral}\label{the-one-sided-spherical-zeta-integral}

With the normalization in (S1), the elementary unramified Godement--Jacquet calculation is \[
 \int_G\mathbf1_{M_a(\mathcal O)}(h)|\det h|^z
                     m_\pi(h)\,dh
       =\prod_{i=1}^a
            (1-\alpha_i q^{-z+(a-1)/2})^{-1}.
 \tag{S8}
\] Replacing \(m_\pi(h)\) by \(m_\pi(h^{-1})\) replaces the \(\alpha_i\) by their inverses. The product of the two formulas is therefore independent of this convention.

For an explicit normalization check, use normalized induction from unramified characters \(\chi_i\), with \(\chi_i(\varpi)=\alpha_i\), and its spherical vector \(v^\circ\) satisfying \(v^\circ(nak)=\delta_B(a)^{1/2}\prod_i\chi_i(a_i)\). Averaging \(v^\circ(kh)\) over \(K\) gives the normalized spherical function. Since the matrix indicator and determinant power are left \(K\)-invariant, that averaging disappears from the integral. In the Iwasawa integral \(h=nak\), integrality requires \(a_i\in\mathcal O\) and \(n_{ij}\in a_j^{-1}\mathcal O\). Integrating these upper off-diagonal entries contributes \(\prod_j|a_j|^{-(j-1)}\). The remaining exponent of each \(|a_i|\) is \[
          z-\frac{a+1-2i}{2}-(i-1)=z-(a-1)/2.
\] The \(a\) geometric series give (S8). This calculation is initially in a half-plane of absolute convergence. The normalized spherical function is determined by its Satake character, so it applies also at reducible principal-series parameters, by the usual spherical-function identity or its continuous specialization. It is not restricted to generic irreducible representations.

A primary independent locator for precisely this normalization is Macdonald (1995), ChapterV, section4, equations(4.4)--(4.5), printed p.301, with the convergence statement immediately following; his section3 records the inverse convention for the spherical transform. Thus (S8) is a known formula used with its constant checked here.

For \(B\) sufficiently large, \(f_B\in L^1(G)\); the spherical transform is multiplicative on \(K\)-bi-invariant convolution. Equations(S7)--(S8) give \[
 \int_G A_B(g)m_\pi(g)\,dg
  =C_{a,B}(q)
     L((B-a+1)/2,\pi)L((B-a+1)/2,\pi^\vee).
 \tag{S9}
\]

\subsubsection{2.3 Extension and absolute convergence at the actual rank}\label{extension-and-absolute-convergence-at-the-actual-rank}

Put \(r(\lambda)=\sum_i|\lambda_i|\), and introduce an independent complex variable \(T\). The Cartan sum \[
 J_\pi(T)=\sum_{\lambda_1\ge\cdots\ge\lambda_a}
   \operatorname{vol}(Ka_\lambda K)m_\pi(a_\lambda)
                                                    T^{r(\lambda)}
 \tag{S10}
\] is a power series convergent near zero. For all sufficiently large integers \(B\), (S9) identifies its value at \(T=q^{-B/2}\) with \[
 R_\pi(T)=
 \frac{\prod_{i=0}^{a-1}(1-q^iT^2)}
      {\prod_{i=1}^{a}
       (1-q^{(a-1)/2}T\alpha_i)
       (1-q^{(a-1)/2}T\alpha_i^{-1})}.
 \tag{S11}
\] These values accumulate at zero, where both functions are analytic. The identity theorem therefore proves equality of their power series near zero, and hence equality wherever the series converges and the rational function has no pole. This step does not assume that an identity at infinitely many large real dimensions alone determines a meromorphic function of the dimension.

For completeness, suppose \[
                    |\alpha_i|,\ |\alpha_i^{-1}|\le q^s
                        \qquad(1\le i\le a).
 \tag{S12}
\] Macdonald's explicit spherical formula and double-coset volume formula give constants \(C,A\), depending on this representation, such that \[
 \begin{split}
 |m_\pi(a_\lambda)|
  &\le C(1+r(\lambda))^A q^{-\langle\rho,\lambda\rangle}
                         q^{s r(\lambda)},\\
 \operatorname{vol}(Ka_\lambda K)
  &\le C q^{2\langle\rho,\lambda\rangle},
 \qquad
 \rho=((a-1)/2,(a-3)/2,\ldots,(1-a)/2).
 \end{split}
 \tag{S13}
\] These estimates follow directly from Macdonald (1995), ChapterV, (2.9) and (3.4), printed pp.298--299. For distinct Satake parameters the finite Weyl sum gives the first bound without the polynomial factor. If parameters coincide, the apparent denominators cancel; taking the finite-order limits differentiates monomials \(\alpha^{w\lambda}\), adding only a polynomial in the integer entries of \(\lambda\). There are finitely many Weyl terms and finitely many possible coincidences, so a fixed \(A\) suffices. This explains the bound also at repeated parameters. The same formula holds after a central shift of \(\lambda\); the displayed bound uses both inverse and direct parameters.

Since \(\langle\rho,\lambda\rangle\le(a-1)r(\lambda)/2\), (S13) proves absolute convergence of (S10) when \[
                       |T|q^{(a-1)/2+s}<1.
 \tag{S14}
\] Only polynomially many dominant integer tuples have a given \(r(\lambda)\), so the geometric decay suffices.

In the actual selected chain let \(a=N_j\), \(b=a+d\), \(j\ge1\). Formula(E13) of ALP-LIFT-14 and the initial bound(E18) imply \[
                |\alpha_i|,\ |\alpha_i^{-1}|
                        \le q^{(d-1)/2}.
 \tag{S15}
\] Indeed the initial rank is \(k<d\), all inserted blocks have length \(d\), and the finite character values on uniformizers have absolute value one. Unitarity gives the same bound for the inverse multiset. A further unramified unitary determinant twist preserves the bound.

At \(T=q^{-b/2}\), (S14) holds with \[
                |T|q^{(a-1)/2+(d-1)/2}=q^{-1}.
 \tag{S16}
\] Thus the actual integral in (S5) is absolutely convergent. Every denominator in (S11) has its nonconstant term of modulus at most \(q^{-1}<1\); the numerator is positive since \(b>a-1\). Substitution gives (S5), with a strictly positive value for unitary \(\pi\). This proof includes residue characteristic2.

\subsection{3. The actual split Rallis test, including the character ratio}\label{the-actual-split-rallis-test-including-the-character-ratio}

At a split place write the source common character as \(\mu=(\eta,\eta^{-1})\). Outside29 it is unramified and unitary. The actual Hom source operator contributes \(\eta(\det g)^{-1}\), by AM-DO-020/022. AM-DO-028 transports GQT's conjugated-input convention and dual-source group back to the literal bilinear integral. Its local matrix coefficient is therefore that of \[
                         \widetilde\pi=\pi\otimes\eta^{-1}.
\] The doubled evaluation identity \(\delta(\phi_0\otimes\overline{\phi_0})(0)
 =\|\phi_0\|^2=1\) shows that the section at \((g,1)\) is exactly the actual vacuum coefficient (S3) with this determinant-root factor. No choice of integral doubled coordinates is needed for this identity.

Consequently the actual local zeta value is \[
 Z_v(s_0;\phi_0,\phi_0,v^\circ,v^\circ)
  =C_{a,b}(q)
       L(s_*,\pi\eta^{-1})L(s_*,\pi^\vee\eta),
 \qquad s_0=d/2.
 \tag{S17}
\] These two factors are exactly the split local standard factor in the source convention fixed by AM-DO-028. The normalized value is \[
                            Z_v^*(s_0)=C_{a,b}(q)>0.
 \tag{S18}
\] If the chosen local group Haar measure gives \(K\) volume \(v_K\) instead of one, both (S17) and (S18) are multiplied by \(v_K\). This measure factor is retained. The formula is valid at the dyadic split places of the actual host; it closes the test-vector point explicitly left open at the end of ALP-LIFT-15.

\subsection{4. Fixed tests at all other places outside29}\label{fixed-tests-at-all-other-places-outside29}

At an odd nonsplit finite place, every stage character is trivial and the actual model is the integral unramified model, as checked in ALP-LIFT-14 section3. Liu (2021), Lemma6.1, identifies the doubled lattice characteristic vector with his unramified standard section. His Proposition5.6 with sign \(+\) gives \[
 Z_v(s_0)=\frac{L_v(s_*,\pi)}{b_a(s_0)},\qquad
 b_a(s)=\prod_{i=1}^a(1-(-1)^i q^{-2s-i})^{-1}.
 \tag{S19}
\] The odd residue characteristic makes the diagonal/antidiagonal lattice change integral, so the doubled product of the actual unit lattice vacua is this lattice characteristic vector. The normalization at the identity is one on both sides. The additive-character rescaling needed at a ramified base-field place was already carried out in LIFT14 without changing the composite symplectic character or normalized vacuum. Therefore the normalized local value is exactly \[
             b_a(s_0)^{-1}
             =\prod_{i=1}^a(1-(-1)^i q^{-d-i})>0,
 \tag{S20}
\] again multiplied by \(v_K\) for a different group Haar normalization.

Thus at every finite place outside29 the \textbf{actual fixed vacuum} and the unique spherical source line give a nonzero normalized local test. All exceptional dyadic places are split and were treated directly in section3. Almost everywhere these are the standard factors in GQT equation(11.3), so their global normalized product converges and is nonzero. At the real places the selected Schur test remains the exact positive \(\operatorname{vol}(U(a))/\dim\rho_{j,\tau}\) of LIFT14(E24). At29 arbitrary finite Schwartz functions remain available, and GQT Proposition11.6(i) supplies the local tests there.

\subsection{5. Restricting the global nonzero family to29}\label{restricting-the-global-nonzero-family-to29}

Let \(\pi\) be any actual cuspidal embedding arising in LIFT14. At each finite place outside29, \(\pi_v^{K_v}\) is one-dimensional. The Rallis argument, with the fixed tests of sections3--4, therefore gives a nonzero lift using spherical source vectors outside29 and the vacuum there, while varying the Schwartz function only at29 and in the selected real component.

This yields every-vector nonvanishing on the away-from29 spherical space, not only one representation-level occurrence. Here are the multiplicity and cancellation details.

For one global irreducible isomorphism type, its spherical space away from29 is \[
   \mathcal M_\pi\otimes\pi_{29}\otimes\rho_j
             \otimes\bigotimes_{v\nmid29,\ v<\infty}\pi_v^{K_v},
 \tag{S21}
\] where \(\mathcal M_\pi\) is its finite-dimensional actual automorphic multiplicity space. The remaining \(\pi_{29}\otimes\rho_j\) is irreducible for \(G(F'_{29})\times G_\infty\), with the product over the six29 places understood. The joint kernel of the theta family with all other finite inputs fixed is invariant under this group. In a finite sum of copies of an irreducible smooth module, a submodule is of the form \(\mathcal M'\otimes
(\pi_{29}\otimes\rho_j)\). If it were nonzero, a nonzero line in \(\mathcal M'\) would define an actual global cuspidal embedding on which all the fixed-away-from29 theta maps vanish. The Rallis argument of the preceding paragraph applies to that embedding, giving a contradiction. Hence this joint kernel is zero.

Different global source isomorphism types cannot cancel their theta images. All target constituents of one source have the unique local small-theta target. At a nonsplit place outside29 the source is spherical and its Satake multiset is recovered by subtracting the fixed inserted block in (E12). At a split place the reciprocal uniqueness in Mínguez's Howe correspondence gives the same injectivity, including at29. Thus two inequivalent global sources, whose real types are already fixed, give inequivalent global targets. Their automorphic isotypic subspaces are orthogonal. A nonzero image of one source component therefore cannot be canceled by the others.

Every smooth current vector has a finite spectral decomposition at its finite level and real type. The preceding two paragraphs show that every nonzero current vector fixed by the standard compacts outside29 survives a test varying only at29. The output is again fixed by those target compacts because the actual vacuum is fixed by them. Hence this property iterates from the first vector of AM-R-012.

Now apply the local version of ALP-LIFT-15 to the29-component. For some finite \(e_j\ge1\), it lies in the depth-\(29^{e_j}\) window. One basis vector detects the nonzero map. We obtain \[
 F_{j+1}=
  \Theta_{W_{x_j}\phi_{0,j},\,\phi_{\infty,j}}(F_j)\ne0,
 \quad
 x_j\in29^{-e_j}\Lambda_{29,j}/\Lambda_{29,j},
 \quad
 x_{j,v}=0\ (v\nmid29),
 \quad
 \|W_{x_j}\phi_{0,j}\|_{2,f}=1.
 \tag{S22}
\] Its infinite translation component is zero. The real vector stays in the exact selected Schur component. This is a full Heisenberg orbit, not the prescribed isotropic half-orbit.

\subsection{6. An explicit sufficient level and its exact index}\label{an-explicit-sufficient-level-and-its-exact-index}

Put \(M=29^e\), and let \(x\in M^{-1}\Lambda_{29}\). If \(g\) and \(h\) are principal-congruence elements modulo \(M^2\) in the source and target integral unitary groups, then \[
                 (g,h)x-x\in M\Lambda_{29}.
 \tag{S23}
\] This follows from the Hom action \(X\mapsto hXg^{-1}\). All four position/momentum matrix blocks satisfy the same integral congruence. Let the displacement be \(\delta\). The Weyl phase \(B(x,\delta)/2\) is integral because \(x\in M^{-1}\Lambda_{29}\), \(\delta\in M\Lambda_{29}\), the pairing is integral on the lattice, and29 is odd. The trace-one character of \(\delta\) is one at29. Thus \(W_{x+\delta}\phi_0=W_x\phi_0\) exactly. The common splitting character has conductor dividing29 and is one on the same principal subgroups. Covariance therefore shows that the whole translated finite vacuum is fixed by these source and target levels.

For rank \(n\), let \(K_n(M^2)\) be principal congruence modulo \(29^{2e}\) at all six split29 places and the standard compact elsewhere. Its exact index is \[
 \begin{split}
 I_n(M^2)&=[K_{n,0}:K_n(M^2)]\\
 &=29^{6n^2(2e-1)}
                  |\operatorname{GL}_n(\mathbf F_{29})|^6\\
 &=29^{12en^2}\prod_{i=1}^n(1-29^{-i})^6.
 \end{split}
 \tag{S24}
\] Therefore (S22) has an explicit sufficient target level of this form. The exponent \(e_j\) has not been bounded in terms of the selected degree.

Every rational class stabilizer at \(K_n(M^2)\) is trivial. Indeed it is finite, because it preserves the corresponding finite lattice and is compact at infinity. At one29 place it embeds into a conjugate of \(1+29M_n(\mathbf Z_{29})\), which is torsion free by the binomial valuation argument in AM-R-012. A conjugate has the same property. Consequently, with \(\operatorname{vol}(G_\infty)=1\) and \(\operatorname{vol}(K_{n,0})=1\), \textbf{every} finite class has mass exactly \(I_n(M^2)^{-1}\), not merely the identity class.

\subsection{7. The 11-adic and archimedean mass accounting}\label{the-11-adic-and-archimedean-mass-accounting}

The order of29 modulo11 is10, and \[
                           v_{11}(29^{10}-1)=1.
\] If \(10\nmid i\), \(v_{11}(29^i-1)=0\). If \(i=10\ell\), the elementary lifting-the-exponent identity gives \(v_{11}(29^i-1)=1+v_{11}(\ell)\). Hence, with \(r=\lfloor n/10\rfloor\), \[
 \boxed{v_{11}(I_n(M^2))
          =6\sum_{\ell=1}^{r}(1+v_{11}(\ell))
          =6\bigl(r+v_{11}(r!)\bigr).}
 \tag{S25}
\] It is independent of \(e\). The residue check can be made without a large integer: \(29^{10}\equiv89\pmod{121}\), so its difference from one has valuation exactly one. The standard lifting argument follows by expanding \((1+11u)^\ell-1\), first for \(\ell\) prime to11 and then iterating the 11th power. This proves the valuation formula.

For example the value in (S25) is0 at rank6,12 at rank20, 18 at rank34, and24 at rank48. Thus the unknown translation depth does not introduce an increasing 11-adic denominator through these class masses. This assertion concerns these Haar factors only.

The archimedean size of the index does depend on depth: \[
 \frac12\log I_n(M^2)
       =6en^2\log29+
                     3\sum_{i=1}^n\log(1-29^{-i}).
 \tag{S26}
\] For clarity, this is a class-coordinate normalization, not a claim that the already produced input can be replaced by an arbitrary class delta. If a scalar automorphic function of compact type \(\rho\) is fixed by this finite level, choose unitary bases and write its class restrictions as \(F_c(k)=\operatorname{tr}(A_c\rho(k))\). Since all stabilizers are trivial, Schur orthogonality gives \[
 \|F\|_2^2=
  \frac{1}{I_n(M^2)\dim\rho}
                  \sum_c\|A_c\|_{\rm HS}^2.
 \tag{S27}
\] Consequently each matrix entry of \(A_c\), in the chosen unitary bases, is at most \(\sqrt{I_n(M^2)\dim\rho}\,\|F\|_2\). Comparison with the source's primitive Schur coefficient lattice remains the separate normalization already studied in AM-R-011. Formula(S27) does not identify the arithmetic row.

\subsection{8. Remaining gate and exact scope}\label{remaining-gate-and-exact-scope}

The new construction controls the finite support and norm: all later inputs are one actual29-supported Heisenberg translate, with coefficient one and finite \(L^2\) norm one. It supplies an explicit sufficient principal level and exact class masses. The later selected real norms, the input-vector norms, and the output content are not set to one.

The depths \(e_j\) are not bounded. More fundamentally, these full-orbit inputs need not equal the source's specified raw Xi or lie in its DP59 half-orbit. AM-DO-029 records the actual finite expansion of that raw function and its separate local support conditions. Neither (S22) nor the mass formula proves that the prescribed row kernel acts nontrivially on the previous image lattice, or that its integral image is saturated.

No later \(g_j\), source \(e_j\)-error, complete \(H\), six comparison identities, or DP44 shortening is deduced from this result. The value of (S25) is useful concrete accounting data; it is not itself a row gain.

\subsection{Primary locators and carried inputs}\label{primary-locators-and-carried-inputs}

\begin{itemize}
\tightlist
\item
  Macdonald, I. G. (1995). Symmetric functions and Hall polynomials (2nd ed.). Oxford University Press. ChapterV, (2.9), p.298; (3.4), p.299; (4.4)--(4.5), p.301, including its convergence statement. \url{https://math.berkeley.edu/~corteel/MATH249/macdonald.pdf}. The formula was read directly; the matrix-disintegration normalization and the application at the actual rank are given in full above.
\item
  Liu, Y. (2021). Theta correspondence for almost unramified representations of unitary groups. arXiv:2105.14223v1. Notation5.5 and equation(5.1), p.15; Proposition5.6, p.16, sign \(+\); Lemma6.1, p.20; Theorem6.2, pp.21--22. These are the actual lattice-vector and normalized local integral statements used in section4.
\item
  Gan, W. T., Qiu, Y., \& Takeda, S. (2014). The regularized Siegel--Weil formula (the second term identity) and the Rallis inner product formula. arXiv:1207.4709v3, sections11.3--11.8, especially the doubled evaluation on p.52, equation(11.3), Theorem11.4 and Proposition11.6.
\item
  Mínguez, A. (2008). Correspondance de Howe explicite: paires duales de type II. Annales Scientifiques de l'École Normale Supérieure, 41(5), 717--741. Theorem1, p.718, including the multiplicity-one/reciprocal uniqueness and all-residue-characteristic assertion.
\item
  Actual host and phases: AM-DO-020/022; exact source-to-doubling convention: AM-DO-028. Actual nonzero chain and Satake blocks: ALP-LIFT-14. Full finite Heisenberg basis: ALP-LIFT-15. First input and torsion-free principal-level proof: AM-R-012. These retain their separate check histories.
\end{itemize}

\clearpage

\section{AM-DO-030: auxiliary finite kernels and simultaneous selected jets in the actual permitted half-orbit}\label{am-do-030-auxiliary-finite-kernels-and-simultaneous-selected-jets-in-the-actual-permitted-half-orbit}

Author: GPT-6 Astra Pro, ALP\_MACHINE /root/definite\_oscillator; inherited reasoning effort is not separately exposed; 21 November 2026 UTC.

The carried inputs are DP59/DP62, AM-K018/022, AM-DO-022/026/029, and AM-R-012. This finding keeps a \textbf{fixed nonzero auxiliary finite Heisenberg operator}, proves simultaneous integral jet interpolation on the actual DP59 half-orbit, and selects one permitted translate that remains nonzero after that operator. Its finite norm, coefficient valuation and an explicit sufficient level are retained. A prescribed local Q-factor at29 is not absorbed into an operator supported away from29.

\subsection{1. Actual variety, orbit and auxiliary vector}\label{actual-variety-orbit-and-auxiliary-vector}

For the even-rank positive edge \(U(a)\to U(b)\), use

\[
 X=A^{ab}=J^{2ab},\qquad g=\dim X=6ab,\qquad
 \mathcal L=L_{2\infty}^{\boxtimes 2ab}.
 \tag{DO30.1}
\]

Its ordinary product CM action, principal trace-one line, common vacuum period \(C^{2ab}\), actual invariant differential coordinates, and finite vacuum \(\phi_0\) are the ones already checked in the carried inputs. In particular the finite vacuum has norm1; its archimedean Gaussian norm is separate. No conjugated \(J\)-factor is substituted for an ordinary factor.

Let

\[
 \alpha=(1+2\sqrt{-7})^e,\qquad M=29^e,\qquad
 \alpha\bar\alpha=M,\qquad e\ge1.
 \tag{DO30.2}
\]

The actual allowed section lattice is

\[
 \mathcal W_M=[\alpha]^*H^0(X,\mathcal L^M)
          \subset H^0(X,\mathcal L^{M^2}).
 \tag{DO30.3}
\]

It has the \textbf{actual} integral DP59 product basis \(\mathcal S_\lambda\), indexed by the prescribed quotient of \(X[\alpha^2]\). Its rescaled parameters are

\[
 x_\lambda=M t'_\lambda\in X[\alpha].
 \tag{DO30.4}
\]

The corresponding finite vectors are \(\phi_\lambda=\epsilon_\lambda W_{x_\lambda}\phi_0\), with their specified finite-order scalars \(\epsilon_\lambda\). These scalars are units at11 and have complex absolute value1. AM-K022 identifies this basis and its exact level-normalized reading; it is not replaced by all sections of \(\mathcal L^{M^2}\).

Fix an odd integer \(d\) with \(\gcd(d,11\cdot29)=1\). Let \(\mathcal Q\) be finite places supported on the rational primes dividing \(d\). Choose finite Heisenberg parameters \(y\), supported on \(\mathcal Q\), with \(dy\in\widehat\Lambda\), where \(\Lambda\) is the actual perfect trace lattice. Consider

\[
 \mathcal O=\sum_y b_yW_y,\qquad
 \phi_{\rm aux}=\mathcal O\phi_0\ne0.
 \tag{DO30.5}
\]

The \(b_y\) are algebraic, and the operator is fixed before choosing the29 orbit. This includes the finite incidence and other auxiliary-place combinations in AM-DO-029 whenever all their parameters are supported away from29 and11. It does not assume that a general such combination is an idempotent projector.

Combine terms that differ by an integral auxiliary lattice vector, using the actual Weyl phase and lattice semicharacter. On the vacuum, and on the whole29-only span, this produces

\[
 \mathcal O=\sum_{j=1}^{s} b_jW_{y_j},\qquad b_j\ne0,
 \tag{DO30.6}
\]

where the \(y_j\) have distinct classes modulo \(\widehat\Lambda\). In (DO30.6) equality is an equality of operators \textbf{on that span}; it is not asserted on every Schwartz vector. Indeed an integral auxiliary Heisenberg operator acts on the fixed auxiliary vacuum by its actual scalar, and commutes with every29-local operator.

The classes define distinct points \(t_j\in X[d]\). The correspondence follows by rational approximation modulo the integral lattice, or directly from the finite quotient \(d^{-1}\Lambda/\Lambda\). Distinct classes remain distinct after reduction at11, since \(X[d]\) is finite étale on the good11-adic model and becomes constant after an unramified extension. We may further extend the coefficient DVR to contain the \(b_j\); this extension need not be unramified. The valuation is always normalized by \(v_{11}(11)=1\).

The finite vectors \(W_{y_j}\phi_0\) are orthonormal. This is the selfdual-lattice character argument used in AM-DO-017/029: a nonintegral difference is detected by a lattice Heisenberg eigencharacter. Consequently

\[
 S_{\rm aux}^2:=\|\phi_{\rm aux}\|^2
                    =\sum_{j=1}^{s}|b_j|^2>0.
 \tag{DO30.7}
\]

Since the auxiliary and29 factors are disjoint,

\[
 \|\mathcal O\phi_\lambda\|^2=S_{\rm aux}^2
 \quad\text{for every allowed }\lambda.
 \tag{DO30.8}
\]

More generally \(\mathcal O\) restricts to a scaled isometry from the29 span, with scale \(S_{\rm aux}\). This follows by writing that span as its variable29 factor tensored with the fixed vacuum elsewhere. There is no interpolation-coefficient norm in (DO30.8).

\subsection{2. A common torsion representative and its exact phase}\label{a-common-torsion-representative-and-its-exact-phase}

A rational representative of an auxiliary finite parameter need not act by a fixed scalar correction on every29-translated vacuum. We choose a representative that does.

Because \(M\) is invertible modulo \(d\), put

\[
 z_j=[M]^{-1}t_j\in X[d].
 \tag{DO30.9}
\]

Choose a lift \(z_j\in d^{-1}\Lambda\), and put \(r_j=M z_j\). Then \(r_j\) represents the finite class \(y_j\), and its component at29 lies in \(M\widehat\Lambda_{29}\). Thus

\[
 \delta_j=y_j-r_j\in\widehat\Lambda,\qquad
 (\delta_j)_{29}\in M\widehat\Lambda_{29}.
 \tag{DO30.10}
\]

Every allowed \(x_\lambda\) is supported at29 with denominator dividing \(M\). Perfectness of the integral trace pairing gives \(\psi_f(B(\delta_j,x_\lambda))=1\). At all other places the second parameter is0. Therefore \(W_{\delta_j}\) acts on the whole permitted orbit by the same scalar as on \(\phi_0\). Applying the Weyl law gives the explicit root of unity

\[
 \varepsilon_{j,M}
 =\psi_f(-B(r_j,\delta_j)/2)\alpha_\Lambda(\delta_j).
 \tag{DO30.11}
\]

In the convention of AM-DO-026, the reading identity is

\[
 W_{y_j}\phi_\lambda
       \ \longleftrightarrow\
 \varepsilon_{j,M}\,U_{r_j}^{H}F_\lambda.
 \tag{DO30.12}
\]

All phases in (DO30.11) are part of the actual finite reading. They are independent of \(\lambda\), and are11-adic units. Here \(F_\lambda\) is the reduced principal-level function whose pullback by \([M]\) is \(\mathcal S_\lambda\).

The reduced translation formula gives the exact identity

\[
 (U_{r_j}^{H}F_\lambda)(Mu)
       =U_{z_j}^{M^2H}\bigl(F_\lambda(Mu)\bigr).
 \tag{DO30.13}
\]

Indeed both sides have argument \(Mu+Mz_j\), linear exponential term \(-\pi M^2H(u,z_j)\), and constant term \(-\pi M^2H(z_j,z_j)/2\). Thus the auxiliary operator, after the common pullback, evaluates the section at the \textbf{distinct} points \(z_j\), with the fixed phases (DO30.11). After \([\alpha]^*\), the points on the base of (DO30.3) are

\[
 q_j=[\alpha]z_j=[\bar\alpha]^{-1}t_j.
 \tag{DO30.14}
\]

They remain distinct, including modulo11, because \(\alpha\) is invertible on \(X[d]\).

\subsubsection{The translated line and the11-adic fibre scalar}\label{the-translated-line-and-the11-adic-fibre-scalar}

Usually \(z_j\notin K(\mathcal L^{M^2})\), since its order is prime to \(M\). Thus \(U_{z_j}^{M^2H}\) is a comparison with a \textbf{translated line}, differing by the corresponding torsion line in \(\operatorname{Pic}^0(X)\). It is not an automorphism in the theta group of \(\mathcal L^{M^2}\).

This causes no loss of an integral fibre unit at11. Pull back by \([d]\). The reduced identity now uses the point \(z_j/d\), whose order divides \(d^2\), on \(\mathcal L^{d^2M^2}\). Since \(d^2\mid d^2M^2\), that point belongs to its actual theta kernel. DP62's relative theta group and odd symmetric normalization apply: the finite-order lift is integral, and its multiplier in any primitive integral fibre frames is a unit. The normalized analytic reduced lift differs by the fixed finite-order scalar, also a unit. The pullback is finite étale at11, and its map on the formal neighbourhood of the zero section is an isomorphism after choosing the zero preimage. The fibre at zero of the pulled-back line is the original fibre. Consequently the constant in the given translated-line formal comparison is this same unit. This is a local fibre assertion, not a global trivialization of the nontrivial torsion line.

The argument asserts a unit \textbf{fibre comparison}, not a unit value of the translated theta section. It does not divide by a theta-null. AM-K018 gives the same normalized canonical gauge on the ordinary \(J\)-product; its positive character CM type is unchanged by a torsion translation. The higher coefficients of this gauge need not be integral ordinary power-series coefficients, and no such claim is used below.

\subsection{3. Simultaneous integral jets with an explicit threshold}\label{simultaneous-integral-jets-with-an-explicit-threshold}

Let \(R\) be a complete coefficient DVR over11 on which the points \(q_1,\ldots,q_s\), the abelian scheme and the line are defined. Their specializations are distinct. For \(D\ge0\), set

\[
 M_0(s,D)=3(s-1)(D+1)+3D+2=3s(D+1)-1.
 \tag{DO30.15}
\]

Then for every \(M\ge M_0(s,D)\), restriction is surjective:

\[
 H^0(X,\mathcal L^M)\longrightarrow
 \bigoplus_{j=1}^{s}
 H^0\!\left(X,\mathcal L^M\otimes
                     \mathcal O_X/I_{q_j}^{D+1}\right).
 \tag{DO30.16}
\]

This is a simultaneous-point extension of AM-K022's single-point argument. The geometric inputs remain the same carried facts: \(\mathcal L^3\) is very ample on an abelian variety, every \(\mathcal L^r\) with \(r\ge2\) is generated at each point, and positive powers have the integral \(H^0\) base change used in DP59/K022. Exact primary locators are given in §7.

\textbf{Proof.} Fix \(j\). Choose a section \(h_{j0}\) of \(\mathcal L^3\) whose value at \(q_j\) is a unit, and sections \(h_{j1},\ldots,h_{jg}\) vanishing there whose ratios \(h_{j\ell}/h_{j0}\) form integral formal parameters. Very ampleness on the special fibre gives the required independent cotangent vectors. Integral base change lifts them; subtracting their values times \(h_{j0}\) makes their vanishing exact. Only unit inverses are needed.

For each \(k\ne j\), choose \(k_{jk}\) in \(H^0(X,\mathcal L^3)\) with unit value at \(q_j\) and value0 at \(q_k\). Very ampleness separates the two special-fibre points. Lift such a section, and subtract its value at \(q_k\) times a section normalized to1 there. That subtracted scalar is in the maximal ideal, so the value at \(q_j\) remains a unit. Put

\[
 H_j=\prod_{k\ne j}k_{jk}^{D+1}.
 \tag{DO30.17}
\]

It is a section of \(\mathcal L^{3(s-1)(D+1)}\), is a unit in the fibre at \(q_j\), and vanishes to order at least \(D+1\) at every other \(q_k\). Put

\[
 r=M-3(s-1)(D+1)-3D\ge2.
 \tag{DO30.18}
\]

Choose \(\eta_j\in H^0(X,\mathcal L^r)\) with unit value at \(q_j\). The section \(H_j\eta_jh_{j0}^{D}\) is a local frame of \(\mathcal L^M\) through the required jet order at \(q_j\). Express any desired jet in this frame as \(\sum_{|\beta|\le D}c_\beta t_j^\beta\), with \(c_\beta\in R\). The global section

\[
 H_j\eta_j
   \sum_{|\beta|\le D}c_\beta
       h_{j0}^{D-|\beta|}
            \prod_{\ell=1}^{g}h_{j\ell}^{\beta_\ell}
 \tag{DO30.19}
\]

has that jet at \(q_j\) and zero \(D\)-jet at every other point. Summing these sections over \(j\) realizes all desired jets. This proves (DO30.16) over \(R\), without division by factorials or interpolation determinants of nonunit valuation. \(\square\)

Pulling back by the actual étale isogeny \([\alpha]\) and its DP59 rigidified line isomorphism gives

\[
 \mathcal W_M\twoheadrightarrow
   \bigoplus_{j=1}^{s}
       \operatorname{Jet}^{D}_{z_j}(\mathcal L^{M^2}).
 \tag{DO30.20}
\]

The formal tangent maps are invertible over \(R\): \(\alpha\bar\alpha=M\) and \(M\) is a unit at11. All line comparison constants are units. Thus this is surjectivity from the actual permitted lattice, not from a larger section space.

One explicit permitted level is

\[
 e(s,D)=\max\{1,\lceil\log_{29}(3s(D+1)-1)\rceil\},
 \quad M=29^{e(s,D)},\quad
 M_0(s,D)\le M<29M_0(s,D).
 \tag{DO30.21}
\]

The threshold recovers \(3D+2\) when \(s=1\).

\subsection{4. The selected functional and exact factorial valuation}\label{the-selected-functional-and-exact-factorial-valuation}

Let \(P=\sum_{|\beta|=D}c_\beta Z^\beta\) be a nonzero homogeneous polynomial in the actual integral CM invariant tangent basis. Define

\[
 B_P=\min_{\beta\in\operatorname{supp}P}
                              v_{11}(c_\beta\beta!),
 \qquad \beta!=\prod_\ell\beta_\ell!.
 \tag{DO30.22}
\]

For the selected primitive bideterminant \(P_\nu\) of AM-K022, its diagonal coefficient is1, so

\[
 B_P\le B_\nu=\sum_{\tau,i}v_{11}(\nu_{\tau,i}!).
 \tag{DO30.23}
\]

For \(S\in\mathcal W_M\), form the exact auxiliary functional

\[
 \ell_{\mathcal O,P,M}(S)
 =M^{-D}\left[P(\partial_u)
     \sum_{j=1}^{s}b_j\varepsilon_{j,M}
             U_{z_j}^{M^2H}F_S(u)\right]_{u=0},
 \tag{DO30.24}
\]

where the translated functions are read in their actual translated-line formal charts, as in §2. All constants are the fixed ones in the reading map and the primitive fibre comparisons, not section-dependent gauge choices.

If \(M\) obeys (DO30.21), at least one \textbf{individual} permitted basis section satisfies

\[
 \boxed{\ell_{\mathcal O,P,M}(\mathcal S_\lambda)\ne0,
 \qquad
 v_{11}\!\left(\ell_{\mathcal O,P,M}(\mathcal S_\lambda)\right)
       \le \min_j v_{11}(b_j)+B_P.}
 \tag{DO30.25}
\]

\textbf{Proof.} Choose \(j\) with \(v_{11}(b_j)\) minimal and a monomial \(\beta\) attaining \(B_P\). By (DO30.20), prescribe zero \(D\)-jets at all \(z_k\), \(k\ne j\). At \(z_j\), choose integral formal parameters with first derivatives the given invariant basis. Prescribe the algebraic coefficient jet \(M^D t^\beta\) in a primitive line frame. The scalar \(M^D\) is an11-adic unit.

The translated source/Chern gauge has a unit constant \(u_{j,M}\). Its higher coefficients are immaterial: the chosen algebraic jet has no term of degree below \(D\). Consequently its degree-\(D\) reduced term is precisely \(u_{j,M}M^D u^\beta\). All other points make zero contribution. Thus for this integral section \(S\),

\[
 \ell_{\mathcal O,P,M}(S)
    =b_j\varepsilon_{j,M}u_{j,M}c_\beta\beta!.
 \tag{DO30.26}
\]

The two displayed comparison scalars are units, so this value has valuation \(\min_jv_{11}(b_j)+B_P\) and is nonzero. Expand \(S\) in the \textbf{integral} basis \(\{\mathcal S_\lambda\}\) of \(\mathcal W_M\). If every basis value had strictly larger valuation, their integral linear combination could not equal (DO30.26). This proves (DO30.25). No archimedean norm of that interpolating linear combination is used; the final selected input is one basis vector. \(\square\)

This proof also explains why the factorial is unavoidable: an ordinary coefficient prescribed to1 has derivative \(\beta!\). It does not prove a unit derivative at arbitrary degree. It also does not require integral coefficients for the full canonical reduced gauge.

The all-polynomial reading AM-DO-026 and the exact level cancellation (DO30.13) give

\[
 \vartheta_{r_P\mathcal O\phi_\lambda}(1,1)
   =C^{2ab}\,\ell_{\mathcal O,P,M}(\mathcal S_\lambda).
 \tag{DO30.27}
\]

This is a nonzero \textbf{actual} selected kernel value, with the one common period retained. The polynomial still lies in the actual selected compact type, since every auxiliary operator is finite and commutes with the archimedean creation action.

\subsection{5. Exact cost and one explicit source compact level}\label{exact-cost-and-one-explicit-source-compact-level}

The chosen finite Schwartz vector is \(\mathcal O\phi_\lambda\). Its coefficients in its reduced auxiliary orbit are the \(b_j\) times unit phases; its exact norm is \(S_{\rm aux}\) from (DO30.7). There is no dependence on \(M\) in that finite norm. The full Schwartz norm is

\[
 \|r_P\phi_{0,\infty}\|\,S_{\rm aux};
 \tag{DO30.28}
\]

the period, highest-vector hook factor and Gaussian covolume remain those of AM-R-011/012. Normalizing the finite vector to norm1 would divide its value by \(S_{\rm aux}\); that extra scalar is not silently made integral.

For the actual \(T\)-incidence product of AM-DO-029, with \(a\)-dimensional source residue space at each selected \(T\)-place, the reduced frequency set and coefficients are

\[
 s=\prod_{P\in v_T}q_P^a,\qquad
 |b_j|=s^{-1},\qquad
 S_{\rm aux}^2=s^{-1}.
 \tag{DO30.29}
\]

The distinct frequencies are the rank-one matrices \(v_P^t\eta_P\). If no \(q_P\) is11, then \(v_{11}(b_j)=0\); thus (DO30.25) retains the same factorial bound \(B_P\). The sufficient level is now \(M\ge3s(D+1)-1\). For one rational prime \(T\) split into six selected half places, \(s=T^{6a}\), not \(T^{6ab}\). Other prescribed \(G/Q\) factors are included by their exact reduced coefficients from (DO29.11), with their actual \(s\), norm and minimum valuation.

Here is an explicit finite stabilizer for the new first lift. Suppose the auxiliary places are odd relative split places, the original compact vacuum is fixed there, and each \(y_j\) has denominator at most \(\varpi_v^{c_v}\) at \(v\), \(c_v\ge1\). Take principal congruence

\[
 K_{{\rm aux},a,v}
   =1+\varpi_v^{2c_v}M_a(\mathcal O_v)
       \ \subset\mathrm{GL}_a(\mathcal O_v),
 \tag{DO30.30}
\]

and similarly on the target. Intersect with the kernel of any separately specified vacuum determinant character if it is ramified there. For the R012 character, which has compact conductor only at29, that extra intersection is unnecessary.

To check fixation, a matrix in (DO30.30) changes \(y_j\) by \(\delta\in\varpi_v^{c_v}\widehat\Lambda_v\). Both \(B(y_j,\delta)/2\) and the trace-one lattice exponent for \(\delta\) are integral at this odd place. Their character values are1. Thus the Weyl covariance formula fixes each \(W_{y_j}\phi_0\), with its actual central phase, not merely its line. At29 use the existing principal congruence of exponent \(e\); at all other places use the original vacuum compact subgroup.

In the unramified split situation its source index is explicitly

\[
 I_{{\rm aux},a}
    =\prod_{v\in\mathcal Q}
      q_v^{a^2(2c_v-1)}
                    |\mathrm{GL}_a(\kappa_v)|.
 \tag{DO30.31}
\]

The exponent comes from the \(2c_v-1\) successive additive matrix kernels. It is the index of the sufficient stabilizer, not an assertion that it is minimal. There can be11-factors in \(|\mathrm{GL}_a(\kappa_v)|\), even when the residue characteristic is not11.

Substitution into the checked class-supported construction R012 gives source identity-class mass \[
 \frac1{I_a(M)I_{{\rm aux},a}},
\] primitive compact-matrix-coefficient input norm squared \[
 \frac1{I_a(M)I_{{\rm aux},a}d_\rho},
\] and actual lifted value

\[
 \frac{C^{2ab}\,
       \ell_{\mathcal O,P,M}(\mathcal S_\lambda)}
          {I_a(M)I_{{\rm aux},a}d_\rho}.
 \tag{DO30.32}
\]

This uses the same trivial rational stabilizer supplied by the29 congruence, and the same compact Schur orthogonality. It does not change the selected real weights or the common splitting. These are source mass factors; they are distinct from the finite Schwartz norm (DO30.8).

For a fixed auxiliary operator, \(s,d\), \(I_{{\rm aux},a}\) and \(S_{\rm aux}\) are fixed and \(\log M=O(\log(D+2))\). If they vary with \(D\), the actual bound is

\[
 \log M<\log29+\log(3s(D)(D+1)-1),
 \tag{DO30.33}
\]

with the displayed auxiliary index and coefficient data retained. In particular one must not call the overhead logarithmic in \(D\) while silently allowing an exponentially growing \(s(D)\).

\subsection{6. Exact remaining raw-array equation}\label{exact-remaining-raw-array-equation}

The theorem now proves nonvanishing after every fixed auxiliary operator \(\mathcal O\) supported away from29 and11 for which \(\mathcal O\phi_0\ne0\). It does so on the actual permitted29 half-orbit, with a single selected basis vector. It is stronger than nonvanishing of some unrestricted finite Schwartz input.

It does \textbf{not} identify that new input with a fixed complete xiAt. Its finite half-function has the form

\[
 \Phi_{\rm aux}(t_{\rm aux})
                \Phi_{\lambda,29}(t_{29}),
 \tag{DO30.34}
\]

where \(\Phi_{\lambda,29}\) is the actual29-translated vacuum. If the prescribed array also contains a fixed local Q-factor at29, the required equation is still \[
 \Phi_{\lambda,29}
       \stackrel{?}{=}\Phi_{Q,29},
\] or membership of the corresponding family in the actual source's allowed local-index span, with all its frame and character rules. Neither K022's interpolation nor (DO30.25) proves that equation.

Applying a nonzero operator supported at29 cannot simply be included in this proof: the disjoint-factor tensor argument and the invertibility of \(\alpha\) on its torsion points are then unavailable. The fixed Q29 image, and the resulting prescribed row/operator identity, therefore remain precise open gates. Later restricted image saturation, Kval, terminal mass and the MainX inequalities are also not consequences of (DO30.32).

\subsubsection{6.1 Literal reference-character and prime-set audit}\label{literal-reference-character-and-prime-set-audit}

The two source symbols called mu0 have different types and must be distinguished. In Development/AbelianVarietyTheory/Definitions/AutomorphicMembership/ \nolinkurl{TrueColumnChoicesHeckeCharReferenceEqOne.lean:}

\begin{longtable}[]{@{}
  >{\raggedright\arraybackslash}p{(\columnwidth - 2\tabcolsep) * \real{0.5000}}
  >{\raggedright\arraybackslash}p{(\columnwidth - 2\tabcolsep) * \real{0.5000}}@{}}
\toprule\noalign{}
\begin{minipage}[b]{\linewidth}\raggedright
Locator
\end{minipage} & \begin{minipage}[b]{\linewidth}\raggedright
Actual object or condition
\end{minipage} \\
\midrule\noalign{}
\endhead
\bottomrule\noalign{}
\endlastfoot
175--190 & the index \(I.\mathrm{mu0}\), with \(\mathrm{tup}(I.\mathrm{mu0})=1\) \\
202--205 & \(I.\mathrm{glob.mu0}=I.C.\mathrm{mu0}\), the family of reference Hecke characters \\
211--215 & \(I.Q(a)=\mathrm{QE}(\mathrm{sel},\mathrm{tup}(a))\) \\
222--226 & the actual tying of the reference characters to Chars \\
273--280 & the proved identity \(I.Q(I.\mathrm{mu0})=\varnothing\) \\
295--299 & the reference twisting character \(I.\xi(I.\mathrm{mu0})=1\) \\
\end{longtable}

Thus the Q-set records the \textbf{finite-order twisting tuple}; it is not the conductor of the reference splitting character. In particular a Q29 factor is not forced merely by saying that a chosen reference character has conductor29.

However the reference character cannot have arbitrary ramification. In \nolinkurl{Formalism/LocalToGlobal/Columns/AbelianHostFieldConditions.lean,} IsMu0, lines314--323, requires it to be unramified away from the two G-primes, tame everywhere, and to have the specified residue character at each G-place. Chars.Tied, lines325--334, imposes these conditions on every reference character. Chars.Ok.chi\_om, lines302--305, prescribes \[
 \chi_{j,P}=\omega_G^{\,2(j+1)}
\] for one faithful residue character at each rational G-prime, and chi\_inj requires the characters to be distinct as \(j\) varies. Consequently the R012 reference character, \textbf{when its29 compact character is nontrivial}, is incompatible with keeping both G-primes away from29. Replacing it by a source-tied reference character is an additional explicit construction, with these residue conditions retained. If29 is one of the G-primes, its G-factor rather than a Q-factor remains in the raw array. If the particular29 compact character is trivial, this ramification obstruction does not apply; the other source tying conditions still do. Moreover, the later good-data clause g\_big in that same file, line488, requires \(2fk<g_i-1\). On the actual mixed host \(f=6\), and the allowed ranks have \(k\ge6\), so \(g_i>73\). Therefore29 cannot be one of these G-primes. For a reference character whose29 compact component is nontrivial, relabelling29 as G does not repair the full source compatibility. This is a further literal prime condition; the analytic characters and theta chain do not claim to satisfy it.

Even when the reference tuple is used and every G/T auxiliary prime avoids29, the new input in (DO30.27) is not automatically the same raw datum. The literal TwDat.Is.idem in \nolinkurl{Definitions/LatticeFrames/LevChiEqOne.lean,} line62, requires the predicate \[
 \operatorname{IsIdemSum}
   \bigl(A.e,\operatorname{Sall}(I.\mathrm{sel},I.\mathrm{tup}(a),u)\bigr).
\] The definition in \nolinkurl{Formalism/LocalToGlobal/IndefiniteReal/DatumVec.lean,} lines78--80, is exactly \[
 \bigl(\exists (e_P)_P,\;
   (\forall P\in S,\operatorname{IsIdemAt}(e_P,P))
   \ \land\ e=\sum_{P\in S}e_P\bigr)
       \ \land\ e\,\overline e=0.
 \tag{DO30.34a}
\] There is no function named IdemSum in this assertion. The support conclusion follows from the literal predicate, not from its name. In \nolinkurl{Formalism/LocalToGlobal/HeckeThetaElements/IsHeckeElt.lean,} lines68--70, IsIdemAt is IsSplitIdem together with \[
 x\in P\quad\Longleftrightarrow\quad
       (e_P/x)_f\notin\widehat{\mathcal O}_E
 \qquad(0\ne x\in\mathcal O_E).
 \tag{DO30.34b}
\] Here IsSplitIdem includes \(e_P^2=e_P\), \(e_P\overline{e_P}=0\), and \(e_P+\overline{e_P}\) equal to one finite base-place idempotent; see \nolinkurl{BesselHermitianTheta/TiedWeil.lean,} lines68--72. Its infinite component is zero, as also proved in EqZeroOrEqOfPlaceIdem.lean, lines134--140.

For completeness, the finite support implication has a short direct proof. Identify the finite tensor adelic algebra with the product of the completions of \(E\). Idempotence says that each component of \(e_P\) is0 or1. For any prime \(Q\ne P\), choose a nonzero \(x\in Q\setminus P\), possible because distinct maximal ideals are comaximal. If the \(Q\)-component of \(e_P\) were1, then the \(Q\)-component of \(e_P/x\) would be nonintegral. This would contradict (DO30.34b), since \(x\notin P\). Thus every such component is0. The \(P\)-component must be1: otherwise all finite components vanish, and any nonzero \(x\in P\) contradicts (DO30.34b). Consequently \(e_P\) is exactly the finite idempotent supported at \(P\). Formula (DO30.34a) therefore gives a zero29 component whenever Sall has no place above29. The normalized place model then sends its half-indicator to the \textbf{fixed} vacuum. Thus the fixed array gives \(\mathcal O\phi_0\), while (DO30.27) gives \(\mathcal O W_{x_\lambda}\phi_0\). The extra29 factor remains an actual change of finite input. One cannot add a29 component to \(e\) and still claim the same literal idem clause.

There is also a precise obstruction to absorbing a nonzero29 position shift by the available linear frame/evaluation actions. At a relative split29 place, the positive Hom polarization identifies every unitary-pair action with \(t\mapsto h^{-1}tg\), times its nonzero determinant-root and Jacobian scalar. Hence every linear translate of the vacuum has position function a nonzero scalar times the indicator of a lattice \textbf{containing0}. A pure position translate by \(x\notin L\) has function \(1_L(t+x)\), whose support does \textbf{not} contain0. They cannot be equal. This proves the claimed obstruction for all such linear source/target evaluations at29, not only integral frames. It does not rule out equality of an isolated theta value, a finite linear combination of group translates, or a genuinely enlarged datum with new affine parameters. It also does not assert that the selected tuple in (DO30.25) is nontrivial: if the zero tuple already works, it is an actual raw witness. The theorem does not force that special outcome.

Accordingly the empty-Q choice removes one possible misidentification of the ramification set, but does not eliminate the raw-array matching equation. The finite linear source model fixes the origin; its frame changes cannot be used as undeclared affine Heisenberg translations.

\subsection{7. Exact sources and check scope}\label{exact-sources-and-check-scope}

\begin{itemize}
\tightlist
\item
  AM-K022, selected\_schur\_permitted\_jets.md, J1--J5: actual permitted lattice and its integral basis; J6--J17: the carried single-point interpolation and highest-symbol factorial argument.
\item
  DP59, research/theta\_translate\_span\_DP59.md, §§3--8; DP62, research/theta\_translate\_norm\_integrality\_DP62.md, TN2--TN5 and §§2--4: actual rigidified multiplication, odd symmetric theta lifts and integral finite-order normalization. These supplied results are used as carried inputs, not reproved in full.
\item
  Milne, J. S. (2008). \emph{Abelian varieties} (Version2.0), I, §6, ``Ample divisors,'' printed p.30, paragraph before Example6.5. \url{https://www.jmilne.org/math/CourseNotes/AVc.pdf}. The arbitrary-characteristic very-ampleness input is the one already read in AM-K022.
\item
  Mumford, D. (1966). On the equations defining abelian varieties. I. \emph{Inventiones Mathematicae, 1}, 287--354, §1(II), p.289. \url{https://pazuki.perso.math.cnrs.fr/index_fichiers/Mumford66.pdf}. The positive-power cohomology/base-change input is carried from DP59/K022.
\item
  Polishchuk, A. \emph{Analogue of Weil representation for abelian schemes}, author preprint arXiv:alg-geom/9712021, §1, pp.4--5, definition preceding Lemma1.1, Proposition1.2 and Corollary1.3. \url{https://arxiv.org/pdf/alg-geom/9712021}. The precise relative theta-group input is the one already checked in DP62.
\item
  AM-K018, jacobian\_source\_gauge\_match.md, §§3--5: exact canonical reduced gauge on the ordinary product; no integral reduced power-series gauge is assumed.
\item
  AM-DO-026, (DO26.2a)--(DO26.4): actual adelic-to-Mumford phase and all-polynomial reading. AM-DO-029, (DO29.11)--(DO29.23): exact prescribed finite coefficients and norms.
\item
  AM-R-012, §§2--3: the actual compact selected first input, stabilizer mass, and Schur-orthogonality factor.
\end{itemize}

Proof-check events are recorded only when actually performed in the agent TSV. No Lean compilation or complete source MainX instance is asserted.

\clearpage

\section{AM-K026 --- a full-Hyp3 family, exact selected norm correction, and the remaining cost interval}\label{am-k026-a-full-hyp3-family-exact-selected-norm-correction-and-the-remaining-cost-interval}

Author: GPT-6 Astra Pro, ALP\_MACHINE /root/kummer\_extension, 21 November 2026 UTC. Inherited reasoning effort is not separately exposed.

This result goes beyond the HC-only unit examples. It supplies an explicit family satisfying every numerical clause of the source's stronger analytic Hyp3, computes the adjacent-gap content bound and the exact hook/mass factors, and proves that the K024 upper slope survives the selected target metric normalization. It still does not identify the second restricted class-image lattice or the full D7/SizeData₂ arithmetic adapter.

The classical/source inputs are retained as theorems. In particular the noncompact-gap spacing lemma is used at its exact recovered interface, and the hook norms are the parent's AM-R-011.

\subsection{1. A concrete full-Hyp3 baseline}\label{a-concrete-full-hyp3-baseline}

Fix an admissible source rank record with k,m even, t odd, r=k+m−1, d=t+1, and \(1801\le r+t\). Put \[
 S=r+mt,\qquad Q=S^3+2,\qquad G=2Q+d=2S^3+t+5.
\tag{HC1}
\] At every CM embedding define \[
 u_i=k-\tfrac12-i,\quad v_i=k-1-i,\quad
 b_j=jG,\quad e_j=(j+\tfrac12)G.
\tag{HC2}
\] The integer G is even. Set eta to the required \(-\operatorname{infTy}(\eta_T)\).

\subsubsection{Proposition HC3}\label{proposition-hc3}

These parameters satisfy the literal HCParam record and the stronger Hyp3 of CayleySignature.lean, lines 110--118, including the 1801 condition, Centred and NCGaps. The actual selected compact partitions are \[
 \nu_j=
 Q\,((2j+1)^k,(2j)^d,(2j-2)^d,\ldots,2^d)
\tag{HC3}
\] at every one of the six labels.

\textbf{Proof.} All u values are half-integral and v,b,e are integral. The two u,v interlacing gaps equal 1/2; the last u equals 1/2\textgreater−b₀=0. The b,e entries alternate strictly increasingly, and both HC gaps equal G/2, which exceeds \((t+3)/2\). Thus both the HC record and TPar.Interlace and TParity hold.

For Centred, u and v are positive, and every e is at least G/2. Thus the required u+e and v+e absolute values exceed \((t+3)/2\). For any i,j, \(|e_j-b_i|=|j+1/2-i|G\ge G/2\). If j′≠j and \(x\in\operatorname{block}(-e_{j'},t)\), then \[
 |x+e_j|\ge |j-j'|G-(t-1)/2
                   \ge G-(t-1)/2>(t+3)/2.
\] This checks every Centred clause in DataTheta.lean, lines 550--554.

For NCGaps apply the accepted source theorem noncompact\_gaps\_of\_spacing in \nolinkurl{ExistsTowerOfLocalSetting.lean,} lines 99--103. Its required lower spacing is \(2(r+mt)^3+t+3\), and G exceeds that by 2. Its other hypotheses are \(-b_0<v_i\), pairwise b spacing at least G, and \(e_j<b_{j+1}\); all were just checked. This gives the exact NCGaps predicate for every j≤m. The 1801 and eta clauses hold by construction.

Finally \(w_{0,i}=u_i-\rho_k(i)=k/2\). Hence DO25.5 gives \(\nu_0=e_0-d/2=Q\) in every row. The exact increment is \(\Delta_j=e_{j+1}-e_j-d=G-d=2Q\). Recurrence DO25.11 gives (HC3). ∎

This supplies the source's full analytic parameter conditions. It does not assert that these numbers satisfy every separately prescribed global upper window, or that an unchanged old induced-CM D7 record has been instantiated on the ordinary E′ host.

\subsection{2. Compulsory first divisibility and an explicit upper bound}\label{compulsory-first-divisibility-and-an-explicit-upper-bound}

Write \(a=k+jd,\ b=a+d\), and \[
 R_j=kj+d\,j(j+1)/2.
\] The total degree and the anti-diagonal upper bound from AM-K025 are \[
 D_j=6Q\{k(2j+1)+d\,j(j+1)\},
\tag{HC4}
\] \[
 U_j=6k\,v_{11}(Q!)+6R_j\,v_{11}((2Q)!).
\tag{HC5}
\] The nonzero prefix differences are Q at k and 2Q at k+d,\ldots,k+jd. In particular \[
 U_j=\frac{D_j}{10}
       -\frac6{10}\{k\,s_{11}(Q)+R_j\,s_{11}(2Q)\}.
\tag{HC6}
\]

For this family every stage satisfies the universal mod-11 cutoff: \[
                              D_j>3140ab.
\tag{HC7}
\] Indeed every part is at least Q, so \(D_j\ge6aQ\). The target rank b is at most \(N_m=k+m(t+1)=S+1\). Since S≥r+t≥1801, the elementary inequality \(6(S^3+2)>3140(S+1)\) proves (HC7).

There is also an exact minimum, rather than just this coarse comparison: \[
 \frac{D_j}{ab}
   =\frac{6Q}{d}\left(1-\frac{k^2}{ab}\right),\qquad
 \min_{0\le j<m}\frac{D_j}{ab}
                         =\frac{6Q}{k+d}.
\tag{HC7a}
\] To verify the identity, expand \(ab=k^2+kd(2j+1)+d^2j(j+1)\). The product ab strictly increases with j, so the displayed ratio strictly increases. Its minimum is at the first edge.

Consequently the exact selected highest-row image ideal satisfies \[
                      \boxed{1\le v_{11}(I_{\nu_j})\le U_j.}
\tag{HC8}
\] The lower bound is only one factor of 11. Neither the large degree nor the cutoff justifies iterating it after dividing the module by 11. The width of (HC8) is an unresolved quantitative gap, not an exact gain calculation.

Increasing the u₀ entries as in AM-K024 preserves this same full Hyp3 and only increases D at the fixed edge. Thus every member of that family remains above the cutoff, while K024 still gives selected coefficients with upper asymptotic rate at most 1/1220. These two facts are compatible: compulsory first divisibility does not determine an asymptotic slope.

\subsection{3. A fully explicit two-edge instance}\label{a-fully-explicit-two-edge-instance}

Take \[
 k=6,\quad m=2,\quad t=1795,\quad d=1796,\quad r=7,
\quad \mathrm{ord}=3.
\tag{HC9}
\] All the source rank inequalities hold, and r+t=1802≥1801. The values are \[
 S=3597,\quad Q=46\,539\,457\,175,\quad
 G=93\,078\,916\,146.
\tag{HC10}
\] Thus \[
 (b_0,b_1)=(0,93\,078\,916\,146),\qquad
 (e_0,e_1)=(46\,539\,458\,073,139\,618\,374\,219).
\] The u entries are \(11/2,9/2,7/2,5/2,3/2,1/2\); the v entries are 5,4,3,2,1.

The exact numerical consequences are:

\begin{longtable}[]{@{}lrr@{}}
\toprule\noalign{}
Quantity & Stage 0 & Stage 1 \\
\midrule\noalign{}
\endhead
\bottomrule\noalign{}
\endlastfoot
Compact edge a→b & 6→1802 & 1802→3598 \\
Selected partition at each label & \(Q^6\) & \((3Q)^6,(2Q)^{1796}\) \\
Total six-label degree D & 1,675,420,458,300 & 1,008,044,642,410,500 \\
D/(3140ab), approximately & 49,350.1104664 & 49,514.7022745 \\
Proven lower bound for \(v_{11}(I_\nu)\) & 1 & 1 \\
Adjacent-gap upper bound U & 167,542,045,704 & 100,804,464,197,676 \\
Least e with \(29^e\ge3D+2\) & 9 & 11 \\
Permitted level M & 14,507,145,975,869 & 12,200,509,765,705,829 \\
\end{longtable}

The certified factorial values are \[
 v_{11}(Q!)=4\,653\,945\,714,\qquad
 v_{11}((2Q)!)=9\,307\,891\,431.
\tag{HC11}
\] Every integer and every source rank, Interlace, Centred and NCGaps check in this example is reproduced by \nolinkurl{hyp3_selected_cost_certificate.py.} The program compares all opposite-sign discrete parameters for j=0,1,2; the minimum nonvacuous gap is G, against the required \(S^3=46\,539\,457\,173\). Common eta shifts cancel from these differences.

\subsection{4. Exact hook and class-mass factors for that instance}\label{exact-hook-and-class-mass-factors-for-that-instance}

Let Hν be the hook product of one label. The source coefficient partition is \(\eta_0=0\) and \(\eta_j=(\nu_{j-1},0^d)\) for j≥1. Thus the actual uniform insertion exponent is Q at stage 0 and 2Q thereafter. AM-R-011 gives the exact recurrence \[
 H_{\nu_j}
   =H_{\eta_j}
       \prod_{i=1}^{a}
        \frac{(\eta_{j,i}+q_j+a-i)!}{(\eta_{j,i}+a-i)!},
 \quad q_0=Q,\quad q_j=2Q\ (j\ge1).
\tag{HC12}
\] The ratio is the actual coefficient norm insertion. Its product over stages telescopes to \(H_{\nu_{m-1}}/H_{\eta_0}\); these ratios must not also be counted as independent uncancelled row-metric losses.

For the example:

\begin{longtable}[]{@{}
  >{\raggedright\arraybackslash}p{(\columnwidth - 4\tabcolsep) * \real{0.2727}}
  >{\raggedleft\arraybackslash}p{(\columnwidth - 4\tabcolsep) * \real{0.3636}}
  >{\raggedleft\arraybackslash}p{(\columnwidth - 4\tabcolsep) * \real{0.3636}}@{}}
\toprule\noalign{}
\begin{minipage}[b]{\linewidth}\raggedright
Exact 11-adic factor
\end{minipage} & \begin{minipage}[b]{\linewidth}\raggedleft
Stage 0
\end{minipage} & \begin{minipage}[b]{\linewidth}\raggedleft
Stage 1
\end{minipage} \\
\midrule\noalign{}
\endhead
\bottomrule\noalign{}
\endlastfoot
\(v_{11}(H_\nu)\), one label & 27,923,674,284 & 16,800,744,033,680 \\
\(v_{11}(H_\nu/H_\eta)\), one label & 27,923,674,284 & 16,772,820,359,396 \\
\(v_{11}(I_a(M))\), R012 class index & 0 & 1182 \\
\(v_{11}(\dim\rho)\), six-label source coefficient & 0 & 0 \\
\end{longtable}

Here the mass convention is the exact AM-R-012 convention \[
 I_a(M)=M^{6a^2}\prod_{\ell=1}^{a}(1-29^{-\ell})^6.
\tag{HC13}
\] At 11 its valuation is independent of the exponent e in M=29\^{}e, because M is an 11-unit. The source dimension is computed by the Weyl dimension formula for η; at stage 1 it is the six-label product for \((Q^6,0^{1796})\). This is an integer dimension, not a denominator inserted into a polynomial projection.

For scale, high-precision numerical evaluations give \[
 \log H_{\nu_0}\approx6.5798135173845434\cdot10^{12},\qquad
 \log H_{\nu_1}\approx4.0756480430066282\cdot10^{15}.
\tag{HC14}
\] These logarithms are numerical evaluations of the exact factorial products, not standalone interval certificates. All exact valuations are computed with integers.

The actual primitive highest creation norm remains, by AM-R-011, \[
 \log\|\phi_j^{\rm hw}\|_\infty
  =-\frac{ab}{2}\log\det(2Y)
       +3\log H_{\nu_j}
       +|\nu_j|\sum_{i=1}^3\log\frac{\pi}{2G_i}.
\tag{HC15}
\] The algebraic-to-analytic kernel period is still \(C^{2ab}\). Subtracting its logarithmic absolute value gives an additional explicit term \(-2ab\log|C|\); that period is not removed by (HC12). Source class extraction also retains (HC13) and its compact dimension exactly as in R012.

These formulas exhibit concrete archimedean and finite factors on a full-Hyp3 family. The interval (HC8) is too wide to identify the actual local gain, and no sum of the separate factors is labelled a total DP44 error.

\subsection{5. Exact metric correction for the scalable K024 selected vector}\label{exact-metric-correction-for-the-scalable-k024-selected-vector}

The target vector in K024 was shown to lie in the exact selected Cauchy summand, but it need not be its highest vector. Its normalization can be computed exactly, with no projector bound.

Fix one residual partition \(\lambda\) of length a. On the spare column let \[
       P_h=Z_{1,1}^{h}\,Q_\lambda(Z_{\text{columns }2,\ldots,a+1}).
\tag{HC16}
\] The variables in its two factors are disjoint. The Fischer monomials are orthogonal, so \[
                    \|P_h\|_{\rm F}^2=h!\,H_\lambda.
\tag{HC17}
\] Its selected source partition is \(\lambda+h e_1\). In the Cauchy realization with unit source highest vector and the target compact metric whose highest vector has norm one, AM-R-011 therefore gives the exact target squared norm \[
                  \theta_\lambda(h)
                  =\frac{h!H_\lambda}{H_{\lambda+h e_1}}.
\tag{HC18}
\] No independent hook penalty is added to (HC17): (HC18) is the comparison between the two ways of expressing that same norm.

Put \(L=\lambda_1+a-1\) and \(c_j=\lambda_1-\lambda_j+j-1\), j=2,\ldots,a. The hook factorial/Vandermonde formula from R011 gives \[
 \boxed{\theta_\lambda(h)=
  \frac1{\binom{h+L}{L}}
       \prod_{j=2}^{a}\frac{h+c_j}{c_j}.}
\tag{HC19}
\] The c\_j are distinct integers in \{1,\ldots,L\}. Cancelling them from the rising factorial gives the equivalent manifestly positive formula \[
 \theta_\lambda(h)
   =\prod_{s\in\{1,\ldots,L\}\setminus\{c_2,\ldots,c_a\}}
                         \frac{s}{h+s}.
\tag{HC20}
\] There are exactly λ₁ factors. In particular \[
        0<\theta_\lambda(h)\le1,\qquad
        |\log\theta_\lambda(h)|=O_\lambda(\log(h+1)).
\tag{HC21}
\]

The corresponding source compact dimension obeys \[
 \frac{\dim S_{\lambda+h e_1}(\mathbf C^a)}
      {\dim S_\lambda(\mathbf C^a)}
                         =\prod_{j=2}^{a}\frac{h+c_j}{c_j}.
\tag{HC22}
\] Removing a determinant twist does not change that dimension, so it is precisely the dimension ratio for the actual row coefficient η.

For the finite valuation, Kummer's carry formula, or directly Legendre's formula, gives \[
 v_{11}\binom{h+L}{L}\le\lfloor\log_{11}(h+L)\rfloor
 \quad(h+L\ge1).
\] Thus (HC19) implies \[
 v_{11}(\theta_\lambda(h))
 \ge-\lfloor\log_{11}(h+L)\rfloor-\sum_{j=2}^{a}v_{11}(c_j).
\tag{HC23}
\] The analogous upper bound and (HC22) give absolute valuation bounds of order log(h+1), with the residual rank and partition fixed.

Apply this at the four growing labels of K024, with \((h,A_r-h,k,B_r-k)\), and write Θ for the product of their θ factors. Scaling the target vector to norm one multiplies its polynomial and its kernel coefficient by Θ\^{}\{-1/2\}. This is an explicit algebraic scalar; it is not silently treated as an integral lattice unit. Combining (HC23) with K024 gives \[
 v_{11}(a_r/\sqrt{\Theta})
       \le\frac{11^r-1}{10}+c_0+O(\log T_r).
\tag{HC24}
\] The upper slope bound remains 1/1220. With the primitive polynomial kept instead, its exact actual creation norm is \[
\begin{split}
 \|\phi_{r,h,k}\|_\infty^2={}&
  \det(2Y)^{-ab}
  h!(A_r-h)!k!(B_r-k)!
       \prod_\tau H_{\nu_\tau^*}\\
 &{}\cdot
  \left(\frac{\pi}{2G_3}\right)^{A_r}
  \left(\frac{\pi}{2G_2}\right)^{B_r}
  \prod_\tau
   \left(\frac{\pi}{2G_{i(\tau)}}\right)^{|\nu_\tau^*|}.
\end{split}\tag{HC25}
\] Dividing this expression by Θ gives exactly the R011 norm with the new selected hook products. These are alternative normalizations of the same vector, not additive costs.

For fixed ranks, the AM-R-012 source class index has \(\log I_a(M)=O(\log T_r)\) at M=O(Tᵣ), and its 11-valuation is constant in e. Equation (HC22) gives \(\log\dim\rho=O(\log T_r)\) and \(v_{11}(\dim\rho)=O(\log T_r)\). Therefore the exact class-supported first extraction and the target metric comparison change this fixed-rank valuation upper bound by at most O(log Tᵣ). The creation norm (HC25), the fixed period, and all other primes remain present.

For the explicit first-edge baseline \(\lambda=(Q^6)\), this simplifies to \[
 \dim S_{(h,0^5)}(\mathbf C^6)=\binom{h+5}{5},\qquad
 \theta_\lambda(h)=
            \frac{\binom{h+5}{5}}{\binom{h+Q+5}{Q+5}}.
\tag{HC26}
\] These formulas make the full-Hyp3 normalized first family completely explicit. They do not control an arbitrary irreducible component after further global decomposition or the second restricted class-image lattice.

\subsection{6. What is still needed for actual g/e}\label{what-is-still-needed-for-actual-ge}

The actual finite local section matrix B is now bounded above on selected families, and its first forced divisibility on a fully admissible analytic family is proved. The permitted finite half-orbit adds no saturation loss by K022. The selected compact metric and first source class extraction have explicit factors.

What remains is not a choice of polynomial normalization: it is the integral image of the already produced source vector under the next selected class matrix, including its fixed branch maps and mass pairing. A larger local gain might hold on that restricted image even when it fails on the full permitted selected kernel family. Conversely no existence theorem for a later nonzero theta representation supplies its valuation content or a controlled integral projection. The current results neither prove such a gain nor rule it out.

\subsubsection{Source locators and certificate}\label{source-locators-and-certificate}

\begin{itemize}
\tightlist
\item
  \nolinkurl{Formalism/LocalRingCardinality/Params.lean,} Ranks and HCParam, lines 7--69.
\item
  \nolinkurl{Formalism/LocalRingCardinality/DataTheta.lean,} TPar.Interlace, Dpar and Centred, lines 529--554.
\item
  \nolinkurl{Formalism/SharpRestriction/CuspSpanNormBounds/CayleySignature.lean,} NCGaps and stronger Hyp3, lines 110--118.
\item
  \nolinkurl{Formalism/SharpRestriction/CuspSpanNormBounds/UnimodularSchurMeasures/ExistsTowerOfLocalSetting.lean,} noncompact\_gaps\_of\_spacing, lines 99--103; the separate integer constructor begins at line 181 and its Hyp3 theorem at line 199. The proof here uses the spacing theorem with the displayed constant-midpoint family; it does not claim that this is the different terminal midpoint choice made by that constructor.
\item
  AM-R-011, R11.7--17: exact Fischer hook norm, determinant insertion and normalized row metric. AM-R-012: exact class-supported extraction and mass. AM-K022--025: finite jets, common source gauge, normal forms, adjacent-gap upper bound and the scalable witness.
\item
  \nolinkurl{hyp3_selected_cost_certificate.py} and its JSON output: exact rank/Hyp3 comparisons, factorial and hook valuations, source dimension valuation, 29-level and index valuation; numerical logarithms are explicitly marked as such.
\end{itemize}

\clearpage

\section{AM-R-014: genuine Witt-lower cuspidality and real-type stability in the source predicates}\label{am-r-014-genuine-witt-lower-cuspidality-and-real-type-stability-in-the-source-predicates}

Author: /root, GPT-6 Astra Pro, ALP\_MACHINE, 21 November 2026 UTC.

This note verifies the mathematical contents of the fourth and sixth analytic predicates in \nolinkurl{CayleySignature.AdmS} for genuine unitary theta integrals. The fourth is an application of the established theta tower property. The sixth is a property of actual smooth isotypic subspaces. Neither is asserted as a new tower theorem. The precise application and the distinction between full smooth and compact-finite vectors are recorded because the source's quantified interfaces require them.

Together with AM-R-013 this accounts for four of the six analytic interfaces, subject to the stated genuine geometric interpretation. The Fourier--Jacobi/free-data assertion and the real discrete-series up/down assertion are not proved here. No complete arithmetic \nolinkurl{D7} or \nolinkurl{ThetaX D} record is claimed.

\subsection{1. The two literal interfaces}\label{the-two-literal-interfaces}

\nolinkurl{Formalism/SharpRestriction/CuspSpanNormBounds/CayleySignature.lean}, lines177--184, defines \nolinkurl{ThetaStablePeriodVanishing} as follows. Given a stable subspace of cusp forms on the source \(Y_{j+1}\), assume that all its theta lifts to the specified lower spaces \(Y_{j,i}^{\mathrm{low}}\), for \[
 1\le i\le q_j,
 \qquad q_j=\left\lfloor\frac{m-j}{2}\right\rfloor,
 \tag{R14.1}
\] vanish for every Schwartz input. Its lift to \(Y_j\) must then be cuspidal. This is a statement about constant terms; it does not assert that the surviving lift is nonzero.

Lines196--199 define \nolinkurl{RealConstantNonvanishing}. Despite its name, the literal conclusion here is stability of \nolinkurl{IsDv} under every finite adelic action and each actual real unitary group action. The definition of \nolinkurl{IsDv}, lines123--125, is \[
       \overline F\in\operatorname{typ}(Y_j,D_j).
 \tag{R14.2}
\] There is no lower bound on a real integral in this particular declaration. A different integral-nonvanishing assertion must not be substituted for it.

The abstract record only requires \nolinkurl{YsLow_zero} at lower index zero (\nolinkurl{DataTheta.lean}, lines639--640). To apply the tower theorem we choose the positive lower indices to be the actual Witt-lower Hermitian spaces, as described next. The auxiliary proof model \nolinkurl{SpanPowerMembership/ThetaFamilyAdmissibleThree.lean}, lines67--70, uses a zero-dimensional space at every positive lower index; that auxiliary choice is not used in the present mathematical application.

\subsection{2. Actual spaces and all rational parabolics}\label{actual-spaces-and-all-rational-parabolics}

Keep the common CM host \(E'/F'\) and even integers \(k,m\), with odd \(t\), and put \[
 n_j=k+m-1+jt,\qquad p_j=n_j-q_j.
 \tag{R14.3}
\] A concrete choice of the required rational Hermitian space is \[
 Y_j=(E')^{n_j},\qquad
       h_j=I_{p_j}\oplus(-I_{q_j}).
 \tag{R14.4}
\] Its standard integral lattice is unimodular. At every real place its signature is exactly \((p_j,q_j)\). Each positive/negative coordinate pair is a rational hyperbolic plane: if its orthogonal vectors are \(a,b\), of norms \(1,-1\), then \((a+b)/2,a-b\) are isotropic and have mutual Hermitian pairing one. Consequently the rational Witt index is at least \(q_j\). It is at most \(q_j\) by any real signature. Thus it equals \(q_j\).

For \(0\le i\le q_j\) take \[
 Y_{j,i}^{\mathrm{low}}
    =I_{p_j-i}\oplus(-I_{q_j-i}),
 \qquad
 Y_j\simeq\mathbb H^i\perp Y_{j,i}^{\mathrm{low}}.
 \tag{R14.5}
\] These are actual rational spaces in one Witt tower, of dimensions \(n_j-2i\). An equivalent construction works for an arbitrary Hermitian \(Y_j\) with the indicated real signature: take its actual Witt decomposition, whose rational index \(q\) is at most \(q_j\), and impose the lower-lift hypothesis at least for \(1\le i\le q\). No rational hyperbolic plane is inferred from signature alone in that more general case.

Every maximal proper rational parabolic of the concrete \(U(Y_j)\) is conjugate to the stabilizer of a rational isotropic subspace of dimension \(i\), for some \(1\le i\le q_j\). Witt's extension theorem gives this conjugacy for each dimension. Its Levi has the factors \(\operatorname{Res}_{E'/F'}\mathrm{GL}_i\) and \(U(Y_{j,i}^{\mathrm{low}})\). Thus the list in (R14.1) contains all proper maximal constant-term tests, not just a chosen archimedean subgroup. For \(q_j=0\) there is no such rational parabolic and cuspidality is automatic.

This choice also admits the required positive embeddings at every nonterminal stage, because \[
 p_j-N_j
   =m-1-j-\left\lfloor\frac{m-j}{2}\right\rfloor
   =\left\lceil\frac{m-j}{2}\right\rceil-1\ge0
       \quad(j<m).
 \tag{R14.6}
\] At the terminal stage \(Y_m=I_{N_m-1}\) embeds as the standard codimension-one positive subspace of \(X_m=I_{N_m}\). These statements supply this concrete analytic geometry. They do not supply every other field, auxiliary-prime, level, lattice and arithmetic condition in the full source \nolinkurl{Spaces} record.

\subsection{3. Character parity is constant along each lower tower}\label{character-parity-is-constant-along-each-lower-tower}

Use the genuine global Weil representation for each pair, with its compatible splitting characters. More precisely, take the genuine \nolinkurl{Chr} domain to be the group of unitary Hecke characters with integral angular exponents of \(E'\) whose restriction to \(\mathbb A_{F'}^\times\) belongs to \(\{1,\epsilon_{E'/F'}\}\), with the chosen odd splitting character \(\eta_T\) in that group. This restriction is part of the mathematical character domain. The literal \nolinkurl{IsMu} condition specifies infinity types alone; it would not force trivial finite base-field restriction in the group of all Hecke characters. Within the stated domain, its integral infinity types force the restriction to be one: the alternative quadratic restriction has the nontrivial sign character at every real place, whereas an integral exponent of \(z/\bar z\) is trivial on the real units. In the source notation the characters on the two members of the descending edge are \[
 \chi_{j+1}=\operatorname{car}(\mu,j+1,j+1),\qquad
 \chi_j=\operatorname{car}(\mu,j+1,j).
 \tag{R14.7}
\] Here the even-rank character \(\mu\) restricts trivially to \(\mathbb A_{F'}^\times\), and the usual odd-rank character \(\eta_T\) restricts to the quadratic character \(\epsilon_{E'/F'}\). The source \nolinkurl{car} inserts \(\eta_T^{-1}\) at an even index and makes no insertion at an odd index. Since \(k+m-1\) and \(t\) are odd, \[
 n_j\equiv j+1\pmod2,
 \qquad
 \chi_j|_{\mathbb A_{F'}^\times}
          =\epsilon_{E'/F'}^{\,n_j}.
 \tag{R14.8}
\] Removing \(i\) hyperbolic planes changes the dimension by \(2i\). The same character is therefore a valid splitting character at every lower target in (R14.5). We use these compatible restrictions, with the fixed additive character, in applying the tower formula. This argument does not identify characters of different parity or apply an even-rank doubling twist formula to an odd-rank group.

\subsection{4. Application of the unitary constant-term formula}\label{application-of-the-unitary-constant-term-formula}

The primary input is Wu's unitary tower formula: Proposition3.1, equations(3.2)--(3.4), author preprint pp.9--12. In a mixed Schwartz model, restriction to the smaller unitary Levi identifies the maximal constant term of a cuspidal theta lift with the lower theta lift of the restricted Schwartz function. The general Levi factor adds the explicitly stated character and modulus. Only vanishing is needed below. The symmetric direction is included in the proposition.

Let \(F\) be any cusp form in the source subspace of the literal predicate, and \(\Phi\) any Schwartz input. Fix a maximal parabolic indexed by \(i\) and a target point \(h\). Right equivariance replaces \(\Phi\) by \(\omega(h)\Phi\), so it suffices to test the constant term at the identity in that mixed model. The tower formula is \[
 \operatorname{CT}_{Q_i}\Theta_{Y_j}(\Phi,F)\big|_{U(Y_{j,i}^{\mathrm{low}})}
       =\Theta_{Y_{j,i}^{\mathrm{low}}}(r_i\Phi,F),
 \tag{R14.9}
\] where \(r_i\) is the Schwartz restriction at zero in the extra isotropic coordinates. The compatible Haar measures give this displayed normalization; changing them multiplies the two sides by the corresponding nonzero measure constant and has no effect on the following implication.

The lower-lift hypothesis quantifies over every Schwartz function, so the right side of (R14.9) is zero, including after the above right translation. The constant term therefore vanishes everywhere on that maximal parabolic. This holds for every \(i\), and hence for every proper rational parabolic: the constant term along a smaller parabolic is an iterated constant term through a containing maximal one. This proves cuspidality of the lift.

Although Wu states the proposition using an irreducible cuspidal representation, its pointwise identity applies to every cusp form. In the cited proof, equations(3.21)--(3.22) use only the vanishing of its proper constant terms; equations(3.23)--(3.25) give the surviving zero orbit. No projection onto an irreducible constituent is used here. In particular an arbitrary stable, possibly nonclosed source subspace in the literal predicate causes no semisimplicity gap. Absolute convergence for the cuspidal input is part of the cited formula.

Thus \nolinkurl{ThetaStablePeriodVanishing} holds in this genuine Witt-lower interpretation. Once compatible splitting characters are fixed, the constant-term argument uses no further \nolinkurl{Hyp3} inequalities and needs no nonzero lower image. Character parity was checked above using the stated parity conditions. Thus it implies the corresponding restricted source quantifiers. It supplies no bound for the norm or denominators of the descended vector.

\subsection{5. Full smooth isotypic spaces and the sixth predicate}\label{full-smooth-isotypic-spaces-and-the-sixth-predicate}

Use the usual smooth automorphic space of moderate growth for the ambient real-group action, with its cusp subspace of rapidly decreasing smooth cusp forms and its usual Schwartz seminorms. In particular the cuspidal vectors under discussion are in finite sums of smooth globalizations of actual cuspidal automorphic embeddings. This is also the cusp domain used in section4: arbitrary smooth functions satisfying only a formal zero-constant-term condition are not being substituted into an absolutely convergent theta integral. The genuine theta maps and constant-term maps are taken in these standard smooth spaces. For a parameter collection \(D\), let \(\operatorname{typ}(V,D)\) be the algebraic sum of the smooth real isotypic subspaces with the specified real representation types. If no such representation exists, this space is zero. At an indefinite real place a discrete-series type here means its full smooth globalization, rather than only its compact-finite vectors. Compact-finiteness remains a condition on the test kernels at the places where the source explicitly asks for it.

Each real-group operator preserves every such isotypic subspace. Every finite adelic right translation commutes with the real action. If \(T:V_\sigma^\infty\longrightarrow\mathcal A\) is an equivariant realization of one of the prescribed real types, then \(R(g_f)T\) is another such realization (or zero). It follows that \[
 R(g_f)\operatorname{typ}(V,D)
       \subset\operatorname{typ}(V,D),\qquad
 R(h_\infty)\operatorname{typ}(V,D)
       \subset\operatorname{typ}(V,D).
 \tag{R14.10}
\] Pointwise complex conjugation commutes with right translation of functions by a fixed real or finite group element. Consequently \[
 \overline F\in\operatorname{typ}(V,D)
 \ \Longrightarrow\ 
 \overline{R(g)F}=R(g)\overline F
                    \in\operatorname{typ}(V,D).
 \tag{R14.11}
\] This is exactly the finite and real stability in \nolinkurl{RealConstantNonvanishing}, for every parameter collection and every stage. No numerical nonzero constant is inferred from the name of the declaration.

For a noncompact real group, it would be incorrect to make this argument with only the compact-finite vectors as an ambient space carrying every real translation. Such a subspace need not be stable under that action. The smooth globalization in this interpretation resolves that actual domain issue. On the definite side of AM-R-013, the selected real representation is finite-dimensional, so that issue does not change its algebraic finite-sum identities.

\subsection{6. Framework translation and exact remaining gates}\label{framework-translation-and-exact-remaining-gates}

The parabolic and Levi groups are actual reductive group schemes over the number field. Their character and representation gradings belong to the SGA3 framework, in particular ExposeI, Proposition4.7.3 and Corollary4.7.3.1, already inspected in the geometric dictionary. The passage to a Schur representation class belongs to the carried SGA6 ExposeVI, Theorem3.3. Those geometric statements do not prove the automorphic constant-term formula; Wu's analytic theorem is the separate input used above. The polarized kernel remains an actual theta-group/biextension object in SGA7I, ExposeVII, sections2.7 and2.9, pp.169--174, with its normalizations retained as in AM-DO-026/029.

The two remaining analytic assertions are \nolinkurl{ThetaUnfoldingFreeData} (lines168--175), with its actual Fourier--Jacobi period and factorizable free-data witness, and \nolinkurl{RealThetaLiftDiscreteSeries} (lines186--194), with its precise real up/down parameter and all smaller-target vanishing. Completing these is not equivalent to the still separate arithmetic row ties, raw-Xi selected nonvanishing, terminal mass or complete cost array.

\subsection{References}\label{references-3}

Wu, C. (2013). Irreducibility of theta lifting for unitary groups. \emph{Journal of Number Theory, 133}(10), 3296--3318. Author preprint posted 2014: \url{https://arxiv.org/pdf/1409.0770} . Proposition3.1 and equations(3.2)--(3.4), printed p.9; proof, especially equations(3.21)--(3.25), pp.11--12. The primary statement and these proof passages were directly inspected in this run.

Gan, W. T., Qiu, Y., \& Takeda, S. (2014). The regularized Siegel--Weil formula (the second term identity) and the Rallis inner product formula. \emph{Inventiones Mathematicae, 198}, 739--831. \url{https://arxiv.org/pdf/1207.4709} . Section11.2, author preprint pp.51--52, records the first-occurrence cuspidality consequence for the full classical-group setting, including unitary groups. The more explicit unitary formula used here is Wu's Proposition3.1.

The exact source declarations and mathematical uses are listed above. This verification is not a new proof of the cited tower property and includes no Lean execution.

\clearpage

\section{AM-K027 --- the exact 11-adic slope in an actual full-Hyp3 selected family}\label{am-k027-the-exact-11-adic-slope-in-an-actual-full-hyp3-selected-family}

Author: GPT-6 Astra Pro, ALP\_MACHINE \nolinkurl{/root/kummer_extension}, 21 November 2026 UTC. Inherited reasoning effort is not separately exposed.

The accepted DP51 common-weight theorem supplies the missing lower bound for AM-K024. The full CM gauge comparison AM-K018 now makes that theorem applicable to the actual ordinary host \(A=J^2\), including every integral local section and the exact DP59 half-orbit. The resulting fixed-rank, varying-weight selected family has \textbf{exact asymptotic normalized derivative slope \(1/1220\)}. The statement remains a selected kernel/first-extraction result. It does not identify the later restricted automorphic image or resolve the literal raw-Xi matching equation.

\subsection{1. The common-weight bound in the actual source gauge}\label{the-common-weight-bound-in-the-actual-source-gauge}

Let R be a finite unramified complete extension of \(\mathbb Z_{11}\) containing the needed CM eigenvalues and torsion. Normalize \(v(11)=1\). Set \[
 J=\operatorname{Jac}(y^2=x^7-1),\qquad
 L=L_{2\infty},\qquad X=J^\ell,
\] with its product principal line \(\mathcal L\), where \(\ell\) is fixed. Use the actual Chern/source reduced gauge of AM-K018, and let \(\kappa\) be a positive integer prime to 11. A section is integral here if its algebraic coefficient in an integral formal frame of \(\mathcal L^\kappa\) is integral. No ordinary Taylor integrality of the reduced frame is assumed.

Put \[
                   w=\frac{121}{1220},\qquad
                   c=\frac1{1220},\qquad s=\frac1{10}.
\tag{EX1}
\]

\subsubsection{Proposition EX2. Full actual local-module estimate}\label{proposition-ex2.-full-actual-local-module-estimate}

For every integral formal local section S and its actual reduced germ \(F_S\), in the primitive ordinary J logarithmic coordinates or any integral invertible linear change of them, \[
       v([z^\alpha]F_S)+w|\alpha|\ge0.
\tag{EX2}
\] Consequently every integral homogeneous differential polynomial \(P=\sum_{|\alpha|=D}p_\alpha Z^\alpha\) satisfies \[
 v\bigl(P(\partial)F_S(0)\bigr)
 \ge \frac D{1220}
       -\frac1{10}\max_{\alpha\in\operatorname{supp}P}
                                 \sum_m s_{11}(\alpha_m).
\tag{EX3}
\] Zero values have valuation \(+\infty\). For D≥1 the weaker uniform estimate is \[
 v\bigl(P(\partial)F_S(0)\bigr)
 \ge \frac D{1220}
          -3\ell\bigl(1+\lfloor\log_{11}D\rfloor\bigr).
\tag{EX4}
\] The constants in (EX2)--(EX4) are independent of \(\kappa\) as long as \(11\nmid\kappa\), and apply to the allowed levels \(\kappa=M^2\), \(M=29^e\).

\textbf{Proof and exact accepted interface.} DP49 identifies each actual formal J with its polarized arbitrary-unit singleton/complement block of height three. AM-K018, (K19.2) and (K19.13), identifies the source/Chern reduced section with \[
              F_S=f_S\,\prod_{b=1}^{\ell}U_b^\kappa B_b^{-\kappa}
\tag{EX5}
\] in the actual integral product frame. Here \(f_S\) is an arbitrary integral formal function on the product, and the \(U_b\) are integral units. These exact identifications are carried results; in particular there is no additional rational quadratic exponential.

Use accepted DP51 §6, (HA17)--(HA21), with p=11, h=3 and primitive parameter \(-\kappa\). Its common-weight substitution makes every full primitive kernel \(B_b^{-\kappa}\) integral at weight w, and makes the inverse logarithm maps integral at that same common weight. The latter statement applies to a joint prefactor as well as a separated product of prefactors: composing the integral series \(f_S\prod_bU_b^\kappa\) with the product of those zero-constant integral inverse maps is integral. Each resulting formal coefficient uses only finitely many terms of the prefactor. Thus no separated-factor hypothesis is needed for (EX5). This gives (EX2) directly from the accepted common-weight theorem.

If \(z=Cz'\) with \(C\in\operatorname{GL}_{3\ell}(R)\), the same common scale is applied to every coordinate. Substitution by C preserves integral coefficients at that scale; its inverse gives the converse. Hence (EX2) applies to the exact \(E'=\mathbb Q(\zeta_{28})\) coordinates \(z_{i,\pm}=u_i\pm i v_i\): both this matrix and its inverse are integral at 11. It also applies after the unit diagonal comparisons with the display eigenbasis.

Finally \(\partial^\alpha F_S(0)=\alpha![z^\alpha]F_S\). The Legendre identity \[
 v(\alpha!)=\frac{|\alpha|-\sum_ms_{11}(\alpha_m)}{10}
\] and \(1/10-w=1/1220\) prove (EX3) by the nonarchimedean triangle inequality. Bounding each digit sum by \(10(1+\lfloor\log_{11}D\rfloor)\) proves (EX4). These factorial manipulations are the same normalization as accepted DP51; no new display theorem is asserted. ∎

AM-K023 separately proves Hurwitz integrality of this entire actual local module. Thus the left side of (EX3) is also nonnegative, and an integer ceiling may be taken over the present unramified coefficient ring. The bound is a genuine full weighted coefficient estimate. It is not obtained by iterating the mod-11 normal form or by assuming a finitely generated differential-stable lattice, which AM-K019 excludes in the relevant unrestricted setting.

\subsection{2. A sharper row-degree error for the selected module}\label{a-sharper-row-degree-error-for-the-selected-module}

On the actual edge \(U(a)\longrightarrow U(b)\), the Hom variety is \(J^{2ab}\), so \(3\ell=6ab\). Let the selected row degrees be the six partitions \(\nu_\tau\), and put \[
 D=\sum_\tau|\nu_\tau|,\qquad
 d_{11}(n)=\begin{cases}0&n=0,\\1+\lfloor\log_{11}n\rfloor&n>0.\end{cases}
\] Every monomial in a polynomial of these row degrees satisfies \(\sum_{j=1}^b\alpha_{\tau,i,j}=\nu_{\tau,i}\). Since every entry in that row is at most \(\nu_{\tau,i}\), (EX3) gives \[
 \boxed{
 v\bigl(P(\partial)F_S(0)\bigr)
 \ge \frac D{1220}
                 -b\sum_{\tau,i}d_{11}(\nu_{\tau,i}).}
\tag{EX6}
\] This holds for every integral polynomial in the specified row-weight space, in particular for the fixed source-highest part of the actual selected Cauchy summand used here. The entire abstract tensor product of the source and target Schur modules also has other source row weights; (EX6), with these same individual row degrees, is not an assertion about all those other weight vectors. It introduces no abstract Schur projector. For the primitive highest prefix-minor polynomial using only the first a columns, b in (EX6) may be replaced by a: \[
 v\bigl(P_{\boldsymbol\nu}(\partial)F_S(0)\bigr)
 \ge \frac D{1220}
                 -a\sum_{\tau,i}d_{11}(\nu_{\tau,i}).
\tag{EX7}
\] These are statements in the explicit polynomial coefficient lattice. Changing to a different rationally normalized Weyl lattice requires its actual scaling factor.

For every \(M\ge3D+2\), AM-K022 identifies the image ideals of the formal jet matrix and the permitted half-orbit tuple matrix. The normalizing factor \(M^{-D}\) is an 11-unit. Consequently (EX3), (EX6), and (EX7) are also bounds for their actual allowed normalized kernel coefficients. The fixed period \(C^{2ab}\) remains outside these algebraic coefficient estimates.

\subsection{3. Exact asymptotic slope for AM-K024}\label{exact-asymptotic-slope-for-am-k024}

Fix one edge and its fixed residual partitions \(\nu_\tau^*\), as in AM-K024. In particular a,b and \[
 D_0=\sum_\tau|\nu_\tau^*|,\qquad
 c_0=v\left(\prod_{\tau,i}\nu_{\tau,i}^*!\right)
\] do not vary. Let \(Q_0\) be the residual prefix-minor polynomial on columns 2,\ldots,a+1. Define the finite constant \[
 K_0=\frac1{10}\max_{\beta\in\operatorname{supp}Q_0}
                                     \sum_m s_{11}(\beta_m).
\tag{EX8}
\] It can be bounded explicitly by \(a\sum_{\tau,i}d_{11}(\nu_{\tau,i}^*)\). Set \[
 A_r=11^r,\quad B_r=121\,11^r,\quad
 T_r=122\,11^r+D_0.
\tag{EX9}
\] AM-K024 supplies \(h_r\in[0,A_r]\), \(k_r\in[0,B_r]\) and the primitive integral polynomial \[
 P_r=Z_{3,+,1,1}^{h_r}Z_{3,-,1,1}^{A_r-h_r}
     Z_{2,+,1,1}^{k_r}Z_{2,-,1,1}^{B_r-k_r}Q_0.
\tag{EX10}
\] It already lies in the \textbf{same exact source-and-target Schur summand} with the selected partitions prescribed by (SC4). It is not merely a source-highest vector awaiting an unspecified target projection.

Choose \(M_r=29^{e_r}\) least with \(M_r\ge3T_r+2\). Let \(\mathfrak I_r\) be the derivative image ideal for \(P_r\) on the integral formal \(T_r\)-jet module of \(\mathcal L^{M_r^2}\); equivalently, by AM-K022 J30, it is the ideal generated by the actual allowed half-orbit tuple values. AM-K024 proves it is nonzero and \[
          v(\mathfrak I_r)\le \frac{11^r-1}{10}+c_0.
\tag{EX11}
\]

The four growing variables in (EX10) are disjoint from all variables of \(Q_0\). Therefore the digit term in (EX3) is exactly bounded by \[
 \frac{s_{11}(h_r)+s_{11}(A_r-h_r)
       +s_{11}(k_r)+s_{11}(B_r-k_r)}{10}+K_0
 \le4d_{11}(T_r)+K_0.
\] Combining this lower bound with (EX11) proves \[
 \boxed{
 \frac{T_r}{1220}-4d_{11}(T_r)-K_0
 \ \le\ v(\mathfrak I_r)
 \ \le\ \frac{T_r}{1220}
                 +c_0-\frac1{10}-\frac{D_0}{1220}.}
\tag{EX12}
\] The same two bounds hold for an individual allowed tuple chosen to attain the image ideal. Such a tuple exists because the half-orbit basis is finite and generates that ideal; taking its least-valued nonzero entry attains the minimum. Its finite Schwartz vector is one translate of the fixed norm-one vacuum, with no interpolation coefficient height included.

It follows in particular that \[
 \boxed{
 \lim_{r\to\infty}\frac{v(\mathfrak I_r)}{T_r}
 =\lim_{r\to\infty}\frac{v(a_r)}{T_r}
 =\frac1{1220}.}
\tag{EX13}
\] This is a limit for the constructed sequence, stronger than a mere liminf upper bound. For the indices r covered by (EX13), the tuple, the selected first-row increments, and the level may vary; the ranks, residual partitions, host, and common edge character remain fixed.

By AM-K024 §2 and AM-K026 §1, choose the fixed residual partitions from the explicit full-Hyp3 baseline. Every subsequent nonnegative \(u_0\) increment preserves \textbf{all} clauses of that stronger analytic Hyp3, including Centred, NCGaps and the 1801 condition. Thus (EX13) holds inside that full analytic admissibility family. It does not preserve separately fixed SizeData₂ upper windows, nor does it by itself supply the complete old D7 arithmetic host.

\subsection{4. Exact selected-metric and first-extraction normalization}\label{exact-selected-metric-and-first-extraction-normalization}

For a fixed residual partition \(\lambda\) and first-row increment h, AM-K026 proves that the target vector represented by the primitive polynomial has squared norm \[
 \theta_\lambda(h)=\frac{h!H_\lambda}{H_{\lambda+h e_1}}
  =\prod_{s\in S_\lambda}\frac{s}{h+s},
\tag{EX14}
\] where \(S_\lambda\subset\{1,\ldots,\lambda_1+a-1\}\) is a fixed set of \(\lambda_1\) positive integers. Let \(\Theta_r\) be the product of the four such factors for the increments in (EX10). The product formula gives both valuation bounds \[
 -|S_\lambda|\log_{11}(h+\lambda_1+a)
 \le v(\theta_\lambda(h))
 \le\sum_{s\in S_\lambda}v(s).
\tag{EX15}
\] Thus \(v(\Theta_r)=O(\log(T_r+1))\), with a fixed constant depending on the residual partitions. Normalizing the target vector multiplies its algebraic value by \(\Theta_r^{-1/2}\), an explicitly retained algebraic scalar. Equations (EX12)--(EX15) imply \[
             \lim_{r\to\infty}
       \frac{v(a_r/\sqrt{\Theta_r})}{T_r}=\frac1{1220}.
\tag{EX16}
\] There is no claim that this square-root normalization is integral.

In the concrete compact class-supported first extraction of AM-R-012, the further unit-input factor is \(1/\sqrt{I_a(M_r)\dim\rho_r}\). For fixed a,b, \[
 I_a(M)=M^{6a^2}\prod_{i=1}^a(1-29^{-i})^6
\] has 11-valuation independent of the exponent e, and the dimension ratios under first-row growth are products of fixed finitely many linear factors in the growing increments. Hence \(v(\dim\rho_r)=O(\log(T_r+1))\). Retaining the common period, that specific first value has the form \[
 C^{2ab}\,
       \frac{a_r}{\sqrt{\Theta_r I_a(M_r)\dim\rho_r}},
\tag{EX17}
\] with the exact mass convention of AM-R-012. Its algebraic coefficient after removal of the fixed period again has slope \(1/1220\). This conclusion is about that explicitly constructed first extraction. The raw source-array matching and any later restriction to a previously produced automorphic image are separate gates, not consequences of the normalization calculation.

The full archimedean creation norm is still the exact expression in AM-K026 (HC25), or equivalently the corresponding product of new hooks after dividing by \(\Theta_r\). Its factorial growth has not been replaced by an \(O(\log T_r)\) assertion. Only the specified target comparison and fixed-rank class/dimension factors have logarithmic logarithm or valuation corrections. These are different quantities and must not be conflated in a net gain/error computation.

\subsection{5. Stronger compulsory bounds for the concrete full-Hyp3 baseline}\label{stronger-compulsory-bounds-for-the-concrete-full-hyp3-baseline}

The bounds above also sharpen AM-K026's first-divisibility conclusion without using any iteration of a congruence. In its exact example \[
 k=6,\quad m=2,\quad t=1795,\quad d=1796,
 \qquad Q=46\,539\,457\,175,
\] the two edge partitions are \[
 \nu_0=(Q^6),\qquad
 \nu_1=((3Q)^6,(2Q)^{1796})
\] at each of the six labels. Every positive row part has 11-adic digit length 11. Applying (EX6) and (EX7), then taking integer ceilings, gives:

\begin{longtable}[]{@{}
  >{\raggedright\arraybackslash}p{(\columnwidth - 8\tabcolsep) * \real{0.1579}}
  >{\raggedleft\arraybackslash}p{(\columnwidth - 8\tabcolsep) * \real{0.2105}}
  >{\raggedleft\arraybackslash}p{(\columnwidth - 8\tabcolsep) * \real{0.2105}}
  >{\raggedleft\arraybackslash}p{(\columnwidth - 8\tabcolsep) * \real{0.2105}}
  >{\raggedleft\arraybackslash}p{(\columnwidth - 8\tabcolsep) * \real{0.2105}}@{}}
\toprule\noalign{}
\begin{minipage}[b]{\linewidth}\raggedright
Edge
\end{minipage} & \begin{minipage}[b]{\linewidth}\raggedleft
Degree D
\end{minipage} & \begin{minipage}[b]{\linewidth}\raggedleft
Lower bound for every integral selected polynomial in the fixed source-highest row space
\end{minipage} & \begin{minipage}[b]{\linewidth}\raggedleft
Lower bound for the primitive highest prefix-minor coefficient
\end{minipage} & \begin{minipage}[b]{\linewidth}\raggedleft
Adjacent-gap upper bound for that highest coefficient
\end{minipage} \\
\midrule\noalign{}
\endhead
\bottomrule\noalign{}
\endlastfoot
6 → 1802 & 1,675,420,458,300 & 1,372,581,866 & 1,373,293,082 & 167,542,045,704 \\
1802 → 3598 & 1,008,044,642,410,500 & 825,838,183,001 & 826,051,784,873 & 100,804,464,197,676 \\
\end{longtable}

The upper bounds are AM-K026's exact adjacent-gap factorial bounds, hence apply to the image ideals and to some individual permitted tuple. The lower bounds hold for every integral section in the corresponding local module and every indicated polynomial in the stated fixed row space. The middle two columns do not assert that every element of a rationally rescaled abstract selected lattice is integral in this polynomial coefficient lattice.

These intervals are genuine local coefficient bounds, but they still leave a wide range for the exact positive Smith valuation of this particular fixed baseline. The asymptotically exact conclusion (EX13) instead concerns the explicitly varied first-row family. Neither statement computes the selected class matrix restricted to a previous theta image.

\subsection{6. Sources, new content, and check scope}\label{sources-new-content-and-check-scope}

The display and common-weight estimates are accepted results, not newly derived results:

\begin{itemize}
\tightlist
\item
  Supplied DP51, \nolinkurl{research/harmonic_asymptotic_rate_DP51.md}, §6, (HA17)--(HA21): entire polynomial module, exact common weight and factorial digit error. The full statement and proof were read for this application.
\item
  Supplied DP46, \nolinkurl{research/unit_arrow_kernel_invariance_DP46.md}, §§4--6, especially (UI16)--(UI20): arbitrary unit arrows, integral inverse-map profile and invariant integral frames.
\item
  Supplied DP59, \nolinkurl{research/theta_translate_span_DP59.md}, §§6--8, especially the lower-bound argument in §7: actual integral half-orbit basis, formal prefactor and normalized derivative convention.
\item
  AM-K018 (K19.2), (K19.13): exact actual source/Chern gauge on ordinary J powers; AM-K022 J30: equality of all finite determinantal ideals; AM-K024 (SC4)--(SC20): the exact selected witness, full-Hyp3 preservation and its fixed filler; AM-K026 (HC18)--(HC25): exact target norm and mass comparisons.
\item
  AM-DO-025: actual source-and-target Cauchy labels; AM-R-011: selected metric; AM-R-012: concrete class-supported first extraction. Their source and primary bibliographic locators are retained in those reports.
\end{itemize}

The new deductions are the transfer of the accepted full common-weight estimate to every actual ordinary-product local section, its row-support version (EX6)--(EX7), the matching selected-family sandwich and limit (EX12)--(EX13), and the retained normalized limit (EX16)--(EX17). The finite arithmetic table has a separate exact certificate. No Lean build, new display proof, full source instance, positive Smith-factor computation, or later-image valuation theorem is claimed.

\clearpage

\section{AM-DO-031: the literal terminal Xi, its finite realization, exact norms, and an explicit nonzero raw value}\label{am-do-031-the-literal-terminal-xi-its-finite-realization-exact-norms-and-an-explicit-nonzero-raw-value}

Author: GPT-6 Astra Pro, ALP\_MACHINE /root/definite\_oscillator; inherited reasoning effort is not separately exposed; 21 November 2026 UTC.

The finite realization in AM-DO-029 covered the nonterminal branches. This finding treats the actual remaining branch xiLast. It gives explicit support and period lattices, the corresponding finite Heisenberg expansion, its exact norm, and a rational matrix at which the raw array has an explicitly known nonzero value. None of these statements identifies that value with a theta value, a selected automorphic matrix coefficient, or the MainX terminal mass.

\subsection{1. Source and notation}\label{source-and-notation}

All recovered source files are relative to work/research/alpoge\_source, commit 7cbf7a0c32fca1badf05931879a37c3f8a2b03b4. The claims below are a mathematical reading of those definitions, not a claim to have compiled their Lean dependency graph.

\begin{longtable}[]{@{}
  >{\raggedright\arraybackslash}p{(\columnwidth - 4\tabcolsep) * \real{0.3333}}
  >{\raggedright\arraybackslash}p{(\columnwidth - 4\tabcolsep) * \real{0.3333}}
  >{\raggedright\arraybackslash}p{(\columnwidth - 4\tabcolsep) * \real{0.3333}}@{}}
\toprule\noalign{}
\begin{minipage}[b]{\linewidth}\raggedright
Source
\end{minipage} & \begin{minipage}[b]{\linewidth}\raggedright
Locator
\end{minipage} & \begin{minipage}[b]{\linewidth}\raggedright
Content
\end{minipage} \\
\midrule\noalign{}
\endhead
\bottomrule\noalign{}
\endlastfoot
\nolinkurl{Formalism/LocalRingCardinality/Params.lean} & 6--17, 32--40 & actual rank restrictions; r, N, rS \\
\nolinkurl{Formalism/LocalToGlobal/Columns/FiniteSumNonzeroLargeDepth.lean} & 345--347, 1103--1104 & Q-frame condition and frsF \\
same & 711--715, 780--806 & literal PhiQ, PhiG, PhiGS \\
same & 817, 838--852, 1039--1042 & qOf; limited-depth additive characters; depth hypotheses \\
same & 1123--1131, 1142--1156 & ybSet, repY, PhiTop, ftY \\
\nolinkurl{Development/AbelianVarietyTheory/Definitions/GlobalTwisting/IntegralGLSuccEmbeddingMemParahoricTimesPrincipa.lean} & 391--393, 420--441 & locTop, entB\_top, terminal rank identity \\
same & 444--509, 531--547 & G-support, congruence invariance, representative denominator \\
\nolinkurl{Development/AbelianVarietyTheory/Definitions/AutomorphicMembership/XiAtPenultimateEqXiCpl1.lean} & 124--130, 146--163 & admissible terminal representative; xiOver, xiLast, xiAt \\
\nolinkurl{Formalism/LocalToGlobal/Columns/LevGSym.lean} & 11--13, 23--28, 108--112 & symIdx and the exact symmetric unipotent support \\
\nolinkurl{Development/AbelianVarietyTheory/Definitions/AutomorphicMembership/TrueColumnChoicesHeckeCharReferenceEqOne.lean} & 242 & actual sharp=true condition \\
\nolinkurl{Development/AbelianVarietyTheory/Definitions/LatticeFrames/LevChiEqOne.lean} & 54--55, 68--71 & standard base lattices and the prescribed framed raw array \\
\end{longtable}

Put

\[
 h=m-1,\qquad r=k+m-1,\qquad
 a=N(m-1)=r+ht,\qquad
 b=N(m)=a+t+1.
 \tag{DO31.1}
\]

The actual restrictions include even \(m\ge2\), even \(k\ge m+3\), and \(t\ge r+5\). In particular every matrix block used below exists. Set \(r_S=m/2-1\). Then \(0\le r_S\) and \(k+r_S<r\).

Write \(\mathsf M=\mathrm{cst.M}\ge1\); this symbol is unrelated to the auxiliary level \(29^e\) of AM-DO-030. At a selected finite prime \(P\), write

\[
 q=q_P=|\mathcal O_E/P|,\qquad
 d_P=v_P(q).
 \tag{DO31.2}
\]

Thus \(|q|_P=q^{-d_P}\). The rational integer q appearing in the source is not replaced silently by a uniformizer. For the prescribed G-primes, Primes.Ok.vG\_card, vG\_norm, and piG\_cop give \(2[E^+:\mathbf Q]\) distinct degree-one primes after adjoining the conjugate set; the degree formula then forces ramification index1. Thus \(d_P=1\). These axioms are in \nolinkurl{AbelianHostFieldConditions.lean,} lines207--218. At a general Q-prime the formulas below retain \(d_P\). The common degree-one case is \(d_P=1\).

Fix actual allowed terminal decorations \(p,p'\) and \(Y\in\mathrm{ybSet}(p.N)\). In their prescribed local frames, write the matrix variable as

\[
 X=\begin{pmatrix}Z\\ \eta\end{pmatrix},\qquad
 Z\in M_{(b-1)\times a}(E_P),\quad
 \eta\in M_{1\times a}(E_P),
 \quad
 y'_P=f_P^{-1}\operatorname{repY}(Y).
 \tag{DO31.3}
\]

The top rank identity and entB\_top say exactly that the terminal PhiTop reads the ordinary upper matrix Z. There are no remaining auxiliary-column powers of q hidden in this assertion: \(cOf(m-1)=0\). Any piT scaling of a terminal target decoration is already in its specified frame. Applying an additional scaling would change the literal array.

With \(\epsilon_{\mathrm{Sg}}=-1\), its local factor is

\[
 F_P(X)=\Phi_{\mathrm{Top},P}(Z)\,
 \begin{cases}
  1_{\mathcal O_P^a}(\eta)\,\psi_P(-\eta y'_P),
       &P\in G,\\[2mm]
  q^{-\mathsf M a}
  1_{q^{-\mathsf M}\mathcal O_P^a}(\eta)\,
       \psi_P(-\eta y'_P),&P\in Q.
 \end{cases}
 \tag{DO31.4}
\]

Here G denotes the G branch and Q denotes the complement of G in rowSG; for the source's disjoint prime data these are its usual G and Q sets. The prescribed separation of the T-places from rowSG is retained. The complete raw xiLast is the product of these factors and the T-incidence factors.

\subsection{2. Representative bounds and finite windows}\label{representative-bounds-and-finite-windows}

\subsubsection{2.1 The actual representative has the needed denominator}\label{the-actual-representative-has-the-needed-denominator}

In the prescribed frames,

\[
 y'_P\in
 \begin{cases}
 q^{-1}\mathcal O_P^a,&P\in G,\\
 q^{-\mathsf M}\mathcal O_P^a,&P\in Q.
 \end{cases}
 \tag{DO31.5}
\]

The G assertion is precisely repY\_int at the source locator above. Here is also the Q argument, which does not need an unproved continuity assertion for psi.

By ybSet, there is a vector y with \(y_P\in q^{-\mathsf M}N_P\), and every member of Y differs from y by a vector in \(q^{\mathsf M}N_P\). Since \(\mathsf M\ge1\), \(q^{\mathsf M}N_P\subset q^{-\mathsf M}N_P\). The source's epsilon representative belongs to Y: PtAt.repY\_mem proves this for the actual terminal point; the same conclusion follows directly because the displayed coset contains y. Hence \(\operatorname{repY}(Y)_P\in q^{-\mathsf M}N_P\). DecoF's frame condition identifies \(N_P\) with \(\mathcal O_P^a\), proving the Q case.

\subsubsection{2.2 No global additive-character hypothesis is introduced}\label{no-global-additive-character-hypothesis-is-introduced}

LocCst.Ok only prescribes additivity on \(q^{-D_P}\mathcal O_P\), and value1 on \(\mathcal O_P\). These conditions suffice. At G, the arguments in (DO31.4) belong to \(q^{-1}\mathcal O_P\), and \(D_P\ge1\). At Q they belong to \(q^{-2\mathsf M}\mathcal O_P\), and ChIdx.Ok.depthQ gives \(D_P\ge2\mathsf M\).

The localized rational module modulo its integral submodule has the same finite quotient as its completion. Additivity and psi\_int make psi constant on those integral cosets, so it extends uniquely to this finite local quotient. On either such finite additive quotient, psi is nonzero: additivity applied to x and -x gives \(\psi(x)\psi(-x)=\psi(0)=1\). The quotient is finite, so every value is a root of unity. Thus it has complex absolute value1 and is a unit at every finite place of its coefficient field.

\subsubsection{2.3 Support and period lattices}\label{support-and-period-lattices}

In the chosen frames let \(L=M_{b\times a}(\mathcal O_P)\). The following windows are valid:

\begin{longtable}[]{@{}
  >{\raggedright\arraybackslash}p{(\columnwidth - 4\tabcolsep) * \real{0.3333}}
  >{\raggedright\arraybackslash}p{(\columnwidth - 4\tabcolsep) * \real{0.3333}}
  >{\raggedright\arraybackslash}p{(\columnwidth - 4\tabcolsep) * \real{0.3333}}@{}}
\toprule\noalign{}
\begin{minipage}[b]{\linewidth}\raggedright
Place
\end{minipage} & \begin{minipage}[b]{\linewidth}\raggedright
Support lattice \(L_P^+\)
\end{minipage} & \begin{minipage}[b]{\linewidth}\raggedright
Period lattice \(L_P^-\)
\end{minipage} \\
\midrule\noalign{}
\endhead
\bottomrule\noalign{}
\endlastfoot
T & \(L\) & \(P L\) \\
G & \(L\) & \(qL\) \\
Q & \(M_{(b-1)\times a}(\mathcal O_P)\oplus q^{-\mathsf M}\mathcal O_P^a\) & \(qM_{(b-1)\times a}(\mathcal O_P)\oplus q^{\mathsf M}\mathcal O_P^a\) \\
\end{longtable}

At T this is the carried exact incidence function \(1_{\{vX=u\}}\), including integral support. At G the source proves the support and congruence statements directly. They can also be checked in (DO31.4): PhiTop depends only on the reduction of Z, and changing \(\eta\) by q times an integral row changes \(\eta y'_P\) by an integral element.

At Q, PhiQ requires every entry of Z to be integral, requires the entries in its bottom-left \((b-1-r)\)-by-r block to lie in q times the local ring, and evaluates a residue character on the top r-by-r determinant. Adding q times an integral matrix preserves all these conditions and the determinant residue. Changing \(\eta\) by \(q^{\mathsf M}\) times an integral row changes its pairing with \(y'_P\) by an integral element. Both old and new character arguments remain in the permitted depth. This proves invariance.

All three support lattices contain L and all three period lattices lie in L. Therefore the indices in these frames are

\[
\begin{array}{c|c|c}
 &R_P=[L_P^+:L]&N_P=[L:L_P^-]\\ \hline
 T&1&q^{ab}\\
 G&1&q^{ab}\\
 Q&q^{d_P\mathsf M a}
   &q^{d_P((b-1)a+\mathsf M a)}.
\end{array}
\tag{DO31.6}
\]

For an arbitrary rational choice of base lattice, transport these exact lattices by \(X\mapsto f'_P Xf_P^{-1}\), and use the sum/intersection and Smith-normal-form construction of AM-DO-029. Formula (DO31.6) is not a license to discard the determinant of a nonintegral frame. For the literal standard TwDat base points, both source and target lattices are standard by TwDat.Is.base; the frames at their places are integral isomorphisms, so (DO31.6) applies without a volume factor.

\subsection{3. Exact finite Heisenberg construction}\label{exact-finite-heisenberg-construction}

Extend the local product by the vacuum half-indicator outside its selected places. The resulting finite Schwartz function \(\Phi_{\mathrm{last}}\) agrees with the literal rational xiLast on its source half-model. Its finite support and local constancy follow from the windows just proved, rather than from an assumption about psi beyond its prescribed depth.

Let \(L^+\supset L\supset L^-\) be the actual product windows, \(R=[L^+:L]\), and \(N=[L:L^-]\). For coset representatives \(a_0\in L^+/L\), and \(y\in(L^-)^\perp/L^\perp\), set

\[
 c_{a_0,y}
 =\frac1N\sum_{u\in L/L^-}
       \Phi_{\mathrm{last}}(a_0+u)\,
       \psi_f(-B(u,y)).
 \tag{DO31.7}
\]

Fourier inversion gives

\[
 \Phi_{\mathrm{last}}(t)
 =\sum_{a_0,y}c_{a_0,y}\,
      \psi_f(B(t-a_0,y))1_L(t-a_0).
 \tag{DO31.8}
\]

Consequently the literal normalized place model sends it to

\[
 f_{\mathrm{last}}
 =\sum_{a_0,y}c_{a_0,y}W_{-a_0}W_y\phi_0
 =\sum_{a_0,y}c_{a_0,y}
       \psi_f(-B(a_0,y)/2)\,W_{y-a_0}\phi_0.
 \tag{DO31.9}
\]

These are the same actual signs and central half-phase as AM-DO-029, with the terminal support and period lattices substituted. They are not a choice of an unspecified intertwining scalar.

For completeness, (DO31.8) is proved by fixing a support coset \(a_0+L\) and applying character orthogonality to the finite group \(L/L^-\). Summing \(\psi_f(B(u-u',y))\) over its full character group gives N if \(u=u'\bmod L^-\), and0 otherwise. The vectors \(W_{-a_0}W_y\phi_0\) in (DO31.9) are orthonormal: different a\_0 have disjoint position supports, and for a fixed a\_0 the same finite-character orthogonality applies. Thus

\[
 \|f_{\mathrm{last}}\|^2
 =\sum_{a_0,y}|c_{a_0,y}|^2
 =\frac1N\sum_{t\in L^+/L^-}
                  |\Phi_{\mathrm{last}}(t)|^2.
 \tag{DO31.10}
\]

If \(\delta\Phi_{\mathrm{last}}\) is integral in a coefficient DVR, then \(\delta N c_{a_0,y}\) is integral. Also

\[
 \sum_{a_0,y}|c_{a_0,y}|
 \le\sqrt{RN}\,\|f_{\mathrm{last}}\|.
 \tag{DO31.11}
\]

All finite Heisenberg parameters in (DO31.9) have zero infinite component, so this is an actual IsFinComb-type finite construction on the mixed host. It solves the finite function-equation part of a new mixed datum. It does not manufacture the old elliptic VacIn wrapper or the remaining arithmetic tying records.

\subsection{4. Exact local norms}\label{exact-local-norms}

Use additive Haar measure giving the standard position lattice volume1, equivalently the normalized finite vacuum measure already fixed in AM-DO-017/029.

\subsubsection{4.1 Incidence places}\label{incidence-places}

At T the nonzero residue row v defines a surjective map \(M_{b\times a}(\mathbf F_q)\to\mathbf F_q^a\), \(X\mapsto vX\). Each fibre has \(q^{(b-1)a}\) points. Hence

\[
 \|F_T\|^2=q^{-a}.
 \tag{DO31.12}
\]

\subsubsection{4.2 Q places, with the actual scalar q retained}\label{q-places-with-the-actual-scalar-q-retained}

The top r-by-r block must have determinant a unit; the bottom-left \((b-1-r)\)-by-r block lies in q times the local ring. All other upper entries are unrestricted integral entries. Their support has measure

\[
 q^{-d_P r(b-1-r)}
       \prod_{j=1}^{r}(1-q^{-j}).
\]

The last row has measure \(q^{d_P\mathsf M a}\), and its scalar amplitude in (DO31.4) has square \(q^{-2\mathsf M a}\). Every character value on support has absolute value1. Therefore

\[
 \boxed{\;
 \|F_Q\|^2
 =q^{(d_P-2)\mathsf M a-d_P r(b-1-r)}
       \prod_{j=1}^{r}(1-q^{-j}).
 \;}
 \tag{DO31.13}
\]

In the degree-one, unramified case this specializes to \(q^{-\mathsf M a-r(b-1-r)}
\prod_{j=1}^{r}(1-q^{-j})\). The scalar \(q^{-\mathsf M a}\) in ftY was never replaced by a Haar-normalization guess.

\subsubsection{4.3 G places in the actual sharp case}\label{g-places-in-the-actual-sharp-case}

The upper matrix Z has \(b-1=a+t\) rows. The bottom t rows vanish in the residue field. Its upper a-by-a matrix has block sizes

\[
 (r,\underbrace{t,\ldots,t}_{h\text{ times}}).
\]

The first diagonal block is upper unipotent in the order given by symIdx. In addition the k entries with row \(p<k\) and column \(k+r_S\) are zero. These are k distinct entries that would otherwise be free: symIdx(p)=p+r\_S is strictly less than \(k+r_S\), and that column is fixed by symIdx. Thus the first block contributes \(q^{r(r-1)/2-k}\) choices.

All blocks strictly below the diagonal vanish. The h diagonal t-by-t blocks must be invertible. Indeed the formula contains the inverse of the unit-character evaluation of the determinant of the whole upper a-by-a block; if it is not a unit that value, and hence the local factor, is zero. Since the first block is unipotent and the matrix is block upper triangular, unit determinant is equivalent to invertibility of each of those t-by-t blocks.

The blocks from the first block to the t-blocks contribute \(q^{rht}\), and the blocks above the diagonal among the t-blocks contribute \(q^{t^2h(h-1)/2}\). Thus the number of upper residue matrices on which PhiTop is nonzero is

\[
 q^{\,r(r-1)/2-k+rht+t^2h(h-1)/2}
       |\mathrm{GL}_t(\mathbf F_q)|^h.
 \tag{DO31.14}
\]

This counts exactly the displayed PhiGS conditions. In particular the apparently extra inequalities with a row \(p<r\) in the later block condition place no restriction on the upper rectangular blocks: the natural-number subtraction gives \((p-r)/t=0\), so its strict inequality is false.

The last row has integral support of volume1. The character values have absolute value1. Dividing (DO31.14) by \(q^{(b-1)a}\), and using \(|\mathrm{GL}_t(\mathbf F_q)|
=q^{t^2}\prod_{j=1}^{t}(1-q^{-j})\), gives

\[
 \boxed{\;
 \|F_G\|^2
 =q^{-C_G}
       \prod_{j=1}^{t}(1-q^{-j})^{m-1},\qquad
 C_G=\frac{r(r+1)}2+k+mrt+
                 \frac{m(m-1)}2t^2.
 \;}
 \tag{DO31.15}
\]

The actual source has sharp=true. If one instead uses its separate unsharp PhiG branch, there are no k extra zero entries and the same formula holds with \(C_G-k\). This is a comparison of the two literal branches, not a change to the prescribed data.

For the smallest numerical ranks previously used, \((m,k,t)=(2,6,13)\), one has \((r,a,b)=(7,20,34)\) and \(C_G=385\). Thus the actual G norm square is \(q^{-385}\prod_{j=1}^{13}(1-q^{-j})\). At a degree-one Q-place the norm square is \(q^{-20\mathsf M-182}\prod_{j=1}^{7}(1-q^{-j})\). These finite factors can be very small even though the finite function is nonzero.

\subsubsection{4.4 Global norm and chart scalar}\label{global-norm-and-chart-scalar}

For standard TwDat lattices, the exact finite norm is the product of (DO31.12), (DO31.13), and (DO31.15), over the actual selected places. For a specified raw chart scalar w it is multiplied by \(|w|^2\). If the chart uses different base lattices, its actual additive Jacobian or the equivalent Smith-normal-form formula (DO31.10) must be retained.

The archimedean selected-polynomial norm and its hook and Gaussian factors are separate, as in AM-R-011. A finite norm calculation cannot cancel them by renaming the vacuum.

\subsection{5. An exact nonzero value of the prescribed terminal array}\label{an-exact-nonzero-value-of-the-prescribed-terminal-array}

At every G-place choose the framed local matrix

\[
 X_G=\begin{pmatrix}I_a\\0\end{pmatrix}
 \in M_{b\times a}(E_P).
\]

All diagonal characters have value1, all additive arguments in PhiGS or PhiG are0, and its last row is0. Hence \(F_G(X_G)=1\). At every Q-place choose the same matrix. Its upper top-r determinant is1, its lower-left block is0, and its last row is0. Hence \(F_Q(X_Q)=q^{-\mathsf M a}\). These conclusions hold for every allowed representative \(y'_P\).

At T solve \(vX=u\). This is possible since \(v\ne0\): choose a column vector w with \(vw=1\) and take \(X=wu\). If a full-rank representative is desired, extend w and a basis of \(\ker v\) to construct one, using \(b>a\). The incidence factor is1.

Take sufficiently small local cosets around these matrices, using the periods from Section2 and the actual rational frames. Additive approximation, or a Chinese-remainder argument on a common fractional lattice, supplies a rational matrix \(\gamma\in M_{b\times a}(E)\) in all these cosets. It may simultaneously be chosen integral outside the selected finite places. Concretely, clear the finitely many denominators, impose the finitely many integral congruences by CRT, and cancel the denominator at the remaining primes by additional congruences. Distinct selected prime ideals give comaximal congruence ideals.

The exact local periodicity now proves

\[
 \boxed{\quad
 \Xi_{\mathrm{last}}(\gamma)
       =\prod_{P\in Q}q_P^{-\mathsf M a}\ne0.
 \quad}
 \tag{DO31.16}
\]

For the reference twisting tuple, the source proves \(Q=\varnothing\). Thus this exact raw value is1. With the chart scalar w it is multiplied by w.

The existence assertion is effective once the actual frames, representative and prime ideals are specified: the proof is finite CRT on explicit lattice cosets. It does not assert a height bound for the smallest such gamma; no such bound is needed for (DO31.16), and one is not hidden in the finite norm calculation.

If all these auxiliary residue characteristics are different from11, the unscaled terminal finite vector is primitive over the11-adic coefficient ring. Indeed every table value in (DO31.7) is either0 or a root of unity times \(\prod_Qq_P^{-\mathsf M a}\), hence integral; every denominator N is an11-adic unit. Therefore every Fourier coefficient is integral. At least one is a unit: otherwise Fourier inversion would make every table value nonunit, contradicting (DO31.16). This is primitivity of the finite Schwartz coefficient vector. It is not primitivity of the selected theta image.

\subsection{6. What this closes, and what it leaves open}\label{what-this-closes-and-what-it-leaves-open}

The terminal finite construction is now as explicit as the nonterminal construction in AM-DO-029. Its support/period data, Heisenberg coefficients, phase, norm, denominator and a nonzero raw value are all fixed by formulas. In particular there is no remaining existential obstacle to realizing the prescribed terminal scalar function in the mixed finite oscillator model.

After applying the actual selected archimedean creation polynomial P, its theta reading is the same checked finite sum as in AM-DO-026/029:

\[
 \vartheta_{r_P f_{\mathrm{last}}}(1,1)
 =C^{2ab}\sum_{a_0,y}c_{a_0,y}
    \psi_f(-B(a_0,y)/2)
       [P(\partial)\mathcal U_{y-a_0}F_{00}]_0.
 \tag{DO31.17}
\]

Nonvanishing of (DO31.16) does not prove nonvanishing of (DO31.17). Its terms can cancel, or the vector can lie in the kernel of the specified automorphic projection. The prescribed previous image and terminal arithmetic detector must still be paired through the actual source operator. The remaining MainX requirements include that operator/kernel identity, the integral row normalizations, and the exact terminal mass and tail inequalities. This finding does not substitute an arbitrary nonzero period for those equations.

\subsection{Review status}\label{review-status}

Author whole-proof check: PASS, 2026-11-21T23:13:49.764065+00:00. The complete saved finding and its displayed source formulas were read. The G support count, actual q-scalar at Q, CRT value, Fourier signs and finite primitivity scope were checked. Independent review is separate; no Lean build is claimed.

\clearpage

\section{AM-K028 --- the literal unramified etaT does not exist for the ordinary J² host type}\label{am-k028-the-literal-unramified-etat-does-not-exist-for-the-ordinary-juxb2-host-type}

Author: GPT-6 Astra Pro, ALP\_MACHINE \nolinkurl{/root/kummer_extension}, 21 November 2026 UTC. Inherited reasoning effort is not separately exposed.

For the actual positive host \[
 E'=\mathbb Q(\zeta_{28}),\qquad A=J^2,
 \qquad\Phi'=\{1,3,9,15,17,23\}\pmod{28},
\tag{ET1}
\] there is \textbf{no Hecke character satisfying the source's literal \nolinkurl{IsEtaT}}, when its chosen embeddings are this actual CM type. The obstruction already follows from its \nolinkurl{unr} and \nolinkurl{inf} clauses and from the defining triviality of a Hecke character on principal ideles. It is independent of all theta-lift estimates, the old \nolinkurl{IsTen} issue, and the separate raw-Xi issue.

This statement is about the specified field/type pair. It is not a claim that such a character is impossible for every CM type, every larger host, or every infinity exponent.

\subsection{1. Exact source conventions}\label{exact-source-conventions}

The definition in \nolinkurl{Formalism/LocalToGlobal/Columns/AbelianHostFieldConditions.lean}, lines 308--312, is:

\begin{verbatim}
structure IsEtaT (X : Fields) (σ₀ : X.K →+* ℂ) (χ : HChar X.E) : Prop where
  unitary : χ.IsUnitary
  unr : ∀ v, χ.UnrAt v
  res : χ.ResEta X.F
  inf : χ.HasInf (X.embC σ₀) (fun _ => 1)
\end{verbatim}

The exact imported definitions were recovered and read in \nolinkurl{Formalism/LocalToGlobal/UnitaryNorms/IsUnitaryNormOf.lean}:

\begin{longtable}[]{@{}
  >{\raggedright\arraybackslash}p{(\columnwidth - 2\tabcolsep) * \real{0.5000}}
  >{\raggedright\arraybackslash}p{(\columnwidth - 2\tabcolsep) * \real{0.5000}}@{}}
\toprule\noalign{}
\begin{minipage}[b]{\linewidth}\raggedright
Lines
\end{minipage} & \begin{minipage}[b]{\linewidth}\raggedright
Literal definition or consequence
\end{minipage} \\
\midrule\noalign{}
\endhead
\bottomrule\noalign{}
\endlastfoot
11--19 & \nolinkurl{HChar L} is a continuous homomorphism on the ideles, trivial on \nolinkurl{diagUnits L k} for every \(k\in L^\times\). \\
37 & \nolinkurl{infDiag} is the product of the infinite local embeddings of a principal element. \\
59--65 & \nolinkurl{loc}, \nolinkurl{locK}, and \nolinkurl{inf} are restrictions of that same global character. \\
71 & \nolinkurl{UnrAt v} means the local character is one on every local unit. \\
90--91 & \nolinkurl{HasInf φ n} prescribes the angular exponent n itself. \\
93--94 & \nolinkurl{HasInfQ φ a} instead requires an integer n with \(n=2a\). \\
194--204 & \nolinkurl{TrivOn} and \nolinkurl{ResEta} impose separate restrictions to the real subfield. \\
\end{longtable}

In particular the actual infinity convention is \[
 \chi_\infty(k)=
   \prod_{\tau\in\Phi'}
       \left(\frac{\tau(k)}{|\tau(k)|}\right)^{n_\tau}.
\tag{ET2}
\] The \nolinkurl{IsEtaT} exponent is exactly \(n_\tau=1\), with no factor two and no inverse. The nearby \nolinkurl{IsMu0.inf} uses \nolinkurl{HasInfQ} and has a different convention. Confusing the two would change this test.

\nolinkurl{Chars.Tied.etaT}, in \nolinkurl{AbelianHostFieldConditions.lean}, lines 325--333, actually requires this \nolinkurl{IsEtaT}; it is not an optional later specialization. The separate \nolinkurl{res} and \nolinkurl{unitary} clauses can only restrict the possibilities further.

\subsection{2. The global-unit obstruction}\label{the-global-unit-obstruction}

\subsubsection{Lemma ET3}\label{lemma-et3}

Let L be a number field, let \(\chi:\mathbb A_L^\times/L^\times\to\mathbb C^\times\) be a continuous Hecke character, and assume every finite local component is unramified. Then \[
                    \chi_\infty(\varepsilon)=1
             \qquad(\varepsilon\in\mathcal O_L^\times).
\tag{ET3}
\] Consequently, if \(\chi\) has the angular infinity prescription (ET2), then necessarily \[
       \prod_\tau
       \left(\frac{\tau(\varepsilon)}{|\tau(\varepsilon)|}\right)^{n_\tau}
                           =1
       \qquad(\varepsilon\in\mathcal O_L^\times).
\tag{ET4}
\]

\textbf{Proof.} Put \(U_f=\prod_{v<\infty}\mathcal O_{L_v}^\times\), embedded in the ideles with infinite component one. The direct sum of the local unit groups, namely the subgroup whose entries are one outside a finite set, is dense in this product. The unramified hypotheses kill every factor and hence this dense subgroup. Continuity then gives \(\chi|_{U_f}=1\). For a global unit \(\varepsilon\), its principal idele is the product of its infinite idele and a member of \(U_f\). The global character is one on that principal idele by definition. Dividing by the finite value one proves (ET3), and (ET2) gives (ET4). ∎

The density argument is included because the finite components of a principal root of unity need not be one at all but finitely many primes. One must not replace that idele by an unjustified finite product. Equivalently, the usual Hecke product formula gives (ET3), but no unproved \nolinkurl{ProductRule} field is needed here.

\subsubsection{Theorem ET5}\label{theorem-et5}

For the field and actual CM type (ET1), no character satisfies \nolinkurl{IsEtaT}.

\textbf{Proof.} Apply (ET4) to the global unit \(\varepsilon=\zeta_{28}\). Every archimedean absolute value is one. The actual six embeddings give \[
 \prod_{a\in\Phi'}\zeta_{28}^a
       =\zeta_{28}^{1+3+9+15+17+23}
       =\zeta_{28}^{68}=\zeta_{28}^{12}.
\tag{ET5}
\] This has order \(28/\gcd(28,12)=7\), so it is not one. Lemma ET3 contradicts the simultaneous \nolinkurl{unr} and \nolinkurl{inf} clauses. ∎

There is a useful check using the original primitive CM field. The type \(\Phi'\) restricts to two copies of \(\{1,2,3\}\) modulo 7. Therefore the unit \(\zeta_7\) has prescribed infinity value \(\zeta_7^{2(1+2+3)}=\zeta_7^{12}=\zeta_7^5\ne1\). This is the same order-seven obstruction. No ambiguity in the choice of a primitive 28th root affects it.

Changing the coefficient embedding or applying a Galois conjugation to the whole field/type pair multiplies the exponent 12 by a unit modulo 28, so the nontrivial value still has order seven. Complex conjugation merely inverts it. A different choice of integral frame or a norm-one theta translate does not affect a global infinity-type obstruction.

\subsection{3. What twisting or taking powers can and cannot change}\label{what-twisting-or-taking-powers-can-and-cannot-change}

Every everywhere-finite-unramified character with \textbf{uniform integer angular exponent n} on this same type must satisfy \[
                     28\mid68n,
             \quad\text{equivalently}\quad 7\mid n.
\tag{ET6}
\] Thus the literal exponent one is excluded. Taking a seventh power can kill this particular root-unit value, but changes the infinity exponent to seven and therefore does not produce the prescribed \nolinkurl{IsEtaT.inf}. Taking a root to restore exponent one would restore the same contradiction if it were still unramified everywhere finite.

Similarly, multiplying by an everywhere-unramified finite-order character cannot compensate (ET5): that character has trivial value on all finite global-unit components and trivial continuous angular infinity type. Any Hecke character with the prescribed angular type one must have some finite ramification. Allowing such ramification is a change to the literal \nolinkurl{unr} clause, with additional local and global data to control; it is not a solution of the same source record.

These statements concern a necessary unit condition. They do not assert that the exponent-seven character exists with every remaining source restriction merely because (ET6) is then satisfied. Other units and the exact restriction \nolinkurl{ResEta} must also be checked.

\subsection{4. Ordinary J powers and a type-preserving extension}\label{ordinary-j-powers-and-a-type-preserving-extension}

\subsubsection{Proposition ET7. Compatible full CM actions on ordinary J powers}\label{proposition-et7.-compatible-full-cm-actions-on-ordinary-j-powers}

Suppose a full CM field L of degree \(6s\) acts on \(J^s\), contains the specified central cyclotomic field \(E=\mathbb Q(\zeta_7)\), and its chosen CM type \(\Psi\) is the actual holomorphic type of that ordinary product. Any Hecke character of L that is unramified at every finite place and has uniform angular exponent one on \(\Psi\) necessarily satisfies \[
                                7\mid s.
\tag{ET7}
\]

\textbf{Proof.} On \(H^0(J^s,\Omega^1)\), each of the three E-characters \(1,2,3\) occurs with multiplicity s. On restricting the full CM type \(\Psi\) to E, its embeddings therefore have exactly those multiplicities. The prescribed angular value on the global unit \(\zeta_7\in E\subset L\) is \[
 \prod_{\tau\in\Psi}\tau(\zeta_7)
                    =\zeta_7^{s(1+2+3)}=\zeta_7^{6s}.
\] Lemma ET3 says this value is one. Since 6 is invertible modulo 7, one obtains \(7\mid s\). ∎

This applies to compatible maximal CM subfields in the standard matrix endomorphism algebra of an ordinary J power. It does not apply unchanged to products with conjugated E-actions, to a type unrelated to the actual Hodge representation, or to a CM field that does not contain the specified central E. For the host (ET1), s=2, recovering the contradiction above. Passing merely from J² to another ordinary multiplicity prime to seven cannot repair it.

\subsubsection{Totally real extension of the current host}\label{totally-real-extension-of-the-current-host}

Let \(F'=(E')^+\), let \(\widetilde F/F'\) be totally real of degree e, and set \(\widetilde E=E'\widetilde F\). Since every real element of E′ lies in F′, one has \([\widetilde E:E']=e\). Give \(\widetilde E\) the CM type induced from \(\Phi'\), consisting of every lift of each chosen embedding. The old root \(\zeta_{28}\) then has prescribed angular value \[
                  \zeta_{28}^{68e}.
\tag{ET8}
\] Hence any literal unramified type-one character on this induced pair necessarily has \[
                               7\mid e.
\tag{ET9}
\] For degrees prime to seven the obstruction persists, including extensions of degrees that are powers of two. This is also Proposition ET7 with the resulting ordinary multiplicity \(s=2e\). Degree seven kills this \textbf{particular} obstruction without changing the exponent one on the enlarged CM type.

There is an elementary explicit degree-seven candidate extension: let B be the unique degree-seven subfield of \(\mathbb Q(\zeta_{29})\). Its fixing subgroup in the cyclic group of order 28 has order four and contains complex conjugation, so B is totally real. Since \([F':\mathbb Q]=6\) is coprime to 7, \(F'\cap B=\mathbb Q\); therefore \(\widetilde F=F'B\) has degree seven over F′. For this candidate, (ET8) is one.

Nothing here proves that all unit angular values on \(\widetilde E\) are trivial, that an unramified character with \nolinkurl{ResEta} exists, or that the enlarged CM lattice satisfies the required source frame and polarization conditions. New units of the extension, its class group, the finite restrictions, and the enlarged arithmetic host remain exact tasks. This is a narrowly identified possible repair, not a completed one.

Nor is the obstruction universal among CM types on E′ itself. For example the type induced from the quadratic residue choice modulo seven, \[
                   \{1,9,11,15,23,25\}\pmod{28},
\] has exponent sum 84, so this root-unit test is trivial for it. That is a different type from the ordinary primitive-J² type in (ET1); passing to it does not preserve the geometric construction already checked here. No existence conclusion is inferred from its passing this one test.

\subsection{5. Consequence for the current machine adapter}\label{consequence-for-the-current-machine-adapter}

The positive ordinary host, its exact theta germ and canonical gauge, its selected compact Cauchy shapes, and its permitted finite-jet arithmetic are unaffected as independently stated results. But on this unchanged field/type pair, the literal \nolinkurl{Chars.Tied} record cannot be completed because its required \nolinkurl{etaT} does not exist. This is stronger than merely leaving its construction open. It does not question the accepted source theorem, whose hypotheses concern its own arithmetic host.

Accordingly, AM-K027's exact selected-family slope and the class-supported constructions remain results in the actual indicated positive local/global model, with the full source arithmetic adapter still absent. A valid eventual adaptation must change or enlarge the host/type or explicitly replace the unramified auxiliary-character interface and account for that change. It cannot fill this field with a conventionally normalized character on the same pair.

\subsection{6. Provenance and checks}\label{provenance-and-checks}

The source definition of \nolinkurl{IsEtaT} was already present in the intake-pinned recovery. The additional imported file \nolinkurl{IsUnitaryNormOf.lean} was retrieved in this run from the exact pinned mirror commit \nolinkurl{7cbf7a0c32fca1badf05931879a37c3f8a2b03b4}:

\href{https://raw.githubusercontent.com/CocoSgt/riemann-hypothesis-lean/7cbf7a0c32fca1badf05931879a37c3f8a2b03b4/src/Formalism/LocalToGlobal/UnitaryNorms/IsUnitaryNormOf.lean}{\nolinkurl{https://raw.githubusercontent.com/CocoSgt/riemann-hypothesis-lean/7cbf7a0c32fca1badf05931879a37c3f8a2b03b4/src/Formalism/LocalToGlobal/UnitaryNorms/IsUnitaryNormOf.lean}}

Its SHA256 is \nolinkurl{503df2de36607142871ffe2399c86c903953535658077b7e06b4f6ce593d6739}. It is \textbf{absent from the supplied selected author-file manifest}, so this run does not claim a new byte comparison of that member with the author's original archive. The precise retrieval and that null comparison status are recorded in \nolinkurl{work/research/source_recovery_20261121T231352Z.json}. The infinity convention is directly visible in the recovered source, and the mathematical obstruction is proved above from that displayed convention.

The actual host/type input is accepted DP57 and AM-DO-022 (DO22.1)--(DO22.3). The only new mathematical input is the elementary global-unit argument and its explicitly scoped consequences. The exact residue sums and orders have a separate arithmetic certificate. No Lean compilation, unramified-character existence theorem, or full enlarged-host construction is claimed.

\clearpage

\section{ALP-LIFT-17 --- prescribed column types exclude the spherical chain}\label{alp-lift-17-prescribed-column-types-exclude-the-spherical-chain}

Author: /root/lift\_accounting, GPT-6 Astra Pro, ALP\_MACHINE, 21 November2026UTC.

This result identifies an actual local representation obstruction to placing the chain of ALP-LIFT-14--16 in the source's explicit column modules. It is stronger than the observation that a generic-unipotent projector kills a spherical vector: at every noninitial, nonterminal stage, the specified determinant-block character is absent from the entire unramified local representation. From stage2 onward, even a tame determinant twist of an unramified representation cannot have the prescribed type.

The scope is the literal ColV/colLev\_j family. The later ChIdx.levG\_j is a separate variable family and is not identified with colLev\_j by the source record. No annihilation theorem for every raw Xi entry, or no-go theorem for a different ramified theta chain, is asserted.

\subsection{1. A finite flag lemma, with full proof}\label{a-finite-flag-lemma-with-full-proof}

Let \(\kappa\) be a finite field, \(G=\operatorname{GL}_n(\kappa)\), and let \(B\) be its standard upper triangular subgroup. Fix a block decomposition \[
 n=r+n_1+\cdots+n_s,\qquad n_i\ge1,\quad s\ge1.
 \tag{F1}
\] Let \(P=LN\) be the standard parabolic with Levi \[
 L=\operatorname{GL}_r(\kappa)\times
                       \prod_{i=1}^s\operatorname{GL}_{n_i}(\kappa).
\] For any subgroup \(H_0\subseteq\operatorname{GL}_r(\kappa)\), put \[
 H=(H_0\times\prod_{i=1}^s\operatorname{GL}_{n_i}(\kappa))N.
\] Let \(\lambda_0:H_0\to\mathbf C^\times\) and \(\chi_i:\kappa^\times\to\mathbf C^\times\) be characters, and define \[
 \lambda(h_0,h_1,\ldots,h_s;n)
            =\lambda_0(h_0)\prod_{i=1}^s\chi_i(\det h_i).
 \tag{F2}
\] Thus the character is trivial on the parabolic radical \(N\). If at least one \(\chi_i\) is nontrivial, then \[
 \operatorname{Hom}_H
       (\lambda,\mathbf C[G/B])=0.
 \tag{F3}
\]

Here \(\mathbf C[G/B]\) is the permutation representation on complete flags. To prove(F3), use the finite Bruhat decomposition \(G=PWB\), with permutation matrices \(W\). Since \(H\) contains \(N\) and every Levi factor except the first, every \(H\)-orbit on \(G/B\) has a representative \[
                           \ell_0 wB,
 \tag{F4}
\] where \(\ell_0\) acts only on the first \(r\)-block. The equality \(P=H\operatorname{GL}_r(\kappa)\) follows directly by writing a parabolic matrix as a block diagonal matrix times \(N\) and conjugating \(N\) past the first block.

Choose an index \(i\) with \(\chi_i\ne1\). Let \(D_i\) be the diagonal torus in the \(i\)-th determinant block, with identity in every other block. It is contained in \(H\), commutes with \(\ell_0\), and is conjugated by \(w\) into the full diagonal torus of \(B\). Consequently \[
           D_i\subset H\cap \ell_0wBw^{-1}\ell_0^{-1}.
 \tag{F5}
\] The determinant map \(D_i\to\kappa^\times\) is surjective, so \(\lambda|_{D_i}\ne1\). A \(\lambda\)-eigenfunction on the orbit of(F4) must vanish there: its value at the representative is both unchanged by the stabilizer and multiplied by a nontrivial value of \(\lambda\). It then vanishes on the entire orbit. This holds on every orbit, proving(F3).

The same conclusion holds with every character inverted. No assumption about \(\lambda_0\), genericity, or the order of the finite field is used.

For later use, write the character projector as \[
            e_\lambda=|H|^{-1}\sum_{h\in H}\lambda(h)^{-1}h.
 \tag{F6}
\] Equation(F3) says that this operator is zero on the full flag permutation module. It is therefore zero on every submodule and quotient of that module. If \(Q\) is any parabolic subgroup, the permutation module \(\mathbf C[G/Q]\) embeds in \(\mathbf C[G/B]\), after conjugating \(Q\) to contain \(B\), by pullback of functions along \(G/B\to G/Q\). Hence(F6) is also zero on every quotient of every parabolic permutation module.

This proof uses only the finite Gaussian-elimination/Bruhat decomposition and the explicit torus in(F5). It does not invoke a classification of finite unipotent representations.

\subsection{2. Every unramified local representation is excluded}\label{every-unramified-local-representation-is-excluded}

Let \(F_v\) be a nonarchimedean field, with ring \(\mathcal O\), uniformizer \(\varpi\), residue field \(\kappa\), and put \[
 K=\operatorname{GL}_n(\mathcal O),\qquad
 K_1=1+\varpi M_n(\mathcal O).
\] Let \(\pi\) be an irreducible smooth complex representation with a nonzero \(K\)-fixed vector \(v\). Its entire space is generated by the translates of \(v\). We prove directly that \(\pi^{K_1}\), as a \(K/K_1\)-module, is spanned by quotients of the parabolic permutation modules in section1.

Let \(e_{K_1}\) be averaging with probability Haar measure. Smith normal form gives \[
       \operatorname{GL}_n(F_v)=
         \bigcup_{\mathbf a}K\operatorname{diag}
                (\varpi^{a_1},\ldots,\varpi^{a_n})K.
\] Thus \(\pi^{K_1}\) is spanned by the \(K\)-translates of \[
 w_{\mathbf a}
       =e_{K_1}\pi(\operatorname{diag}(\varpi^{a_i}))v.
 \tag{F7}
\] This follows by expressing any vector as a finite linear combination of translates of \(v\), then averaging. No topological completion of the representation is being taken.

Write \(a=\operatorname{diag}(\varpi^{a_i})\). The subgroup \(K\cap aKa^{-1}\) fixes \(\pi(a)v\). Because \(K_1\) is normal in \(K\), its averaging operator commutes with \(K\), so that subgroup also fixes \(w_{\mathbf a}\). Its reduction modulo \(\varpi\) is a parabolic subgroup: entry \((i,j)\) must have valuation at least \(\max(0,a_i-a_j)\), and the resulting block zero conditions are exactly the flag stabilizer for the equal-exponent blocks. Every matrix in that parabolic has a lift satisfying these conditions, by taking the required off-block entries to be zero. Therefore the \(K/K_1\)-orbit of \(w_{\mathbf a}\) is a quotient of \[
                   \mathbf C[\operatorname{GL}_n(\kappa)/Q_{\mathbf a}]
 \tag{F8}
\] for a parabolic \(Q_{\mathbf a}\).

Inflate \(H,\lambda\) from section1 to their full inverse image \(\widetilde H\subset K\). Since \(K_1\subset\widetilde H\) and the inflated character is trivial on \(K_1\), every \(\lambda\)-eigenvector belongs to \(\pi^{K_1}\). By(F6)--(F8), its character projector vanishes on every spanning submodule. We obtain \[
 \boxed{\pi^{\widetilde H,\lambda}=0
          \quad\hbox{whenever some }\chi_i\ne1.}
 \tag{F9}
\] This excludes the type from the entire representation, including all its nonspherical vectors and all its translates. It does not assume that \(\pi\) is tempered or generic.

The proof also works for any module algebraically generated by \(K\)-fixed vectors. In particular it applies constituent by constituent, with actual automorphic multiplicities, to the away-from29 spherical module in ALP-LIFT-16.

\subsection{3. Tame determinant twists do not fix distinct blocks}\label{tame-determinant-twists-do-not-fix-distinct-blocks}

Let \(\eta:F_v^\times\to\mathbf C^\times\) be a tame character, so it is trivial on \(1+\varpi\mathcal O\). Write \(\overline\eta\) for its residue character. For \(\pi=\pi_{\mathrm{sph}}\otimes(\eta\circ\det)\), a vector of \(\widetilde H\)-type \(\lambda\) would give in \(\pi_{\mathrm{sph}}\) a vector of character \[
       \lambda\cdot(\overline\eta\circ\det)^{-1}.
 \tag{F10}
\] If \(H_0\) is upper unitriangular, its determinant is one; the determinant-block characters in(F10) are \[
                           \chi_i\overline\eta^{-1}.
\] By(F9) they must all be trivial. Thus a necessary condition is \[
                 \chi_1=\cdots=\chi_s=\overline\eta.
 \tag{F11}
\] In particular, two distinct prescribed determinant-block characters exclude every tame determinant twist of every unramified representation. This statement is deliberately restricted to tame twists; no claim about arbitrary wild twists is inferred from the finite reduction argument.

\subsection{4. Exact source identification of the finite type}\label{exact-source-identification-of-the-finite-type}

All paths are relative to the recovered source work/research/alpoge\_source, pinned to commit 7cbf7a0c32fca1badf05931879a37c3f8a2b03b4. The following were read directly.

\begin{longtable}[]{@{}
  >{\raggedright\arraybackslash}p{(\columnwidth - 4\tabcolsep) * \real{0.3333}}
  >{\raggedright\arraybackslash}p{(\columnwidth - 4\tabcolsep) * \real{0.3333}}
  >{\raggedright\arraybackslash}p{(\columnwidth - 4\tabcolsep) * \real{0.3333}}@{}}
\toprule\noalign{}
\begin{minipage}[b]{\linewidth}\raggedright
Source
\end{minipage} & \begin{minipage}[b]{\linewidth}\raggedright
Exact locator
\end{minipage} & \begin{minipage}[b]{\linewidth}\raggedright
Role
\end{minipage} \\
\midrule\noalign{}
\endhead
\bottomrule\noalign{}
\endlastfoot
\nolinkurl{Formalism/LocalToGlobal/Columns/GenericUnipotentCharacter.lean} & 155--180 & uniUT and its block upper triangular description \\
same & 232--256 & supChar, detChar, lamGen, extChi \\
same & 437--452 & HU, HX, SX and lamX \\
\nolinkurl{Formalism/LocalToGlobal/Columns/ColForms.lean} & 92--105 & colLev and ColV \\
\nolinkurl{Formalism/LocalToGlobal/Columns/ColFormsLemmas.lean} & 137--145 & colLev\_of\_lt and the terminal branch \\
\nolinkurl{Formalism/LocalToGlobal/Columns/AbelianHostFieldConditions.lean} & 273--306 & Chars, chi\_inj, chi\_om \\
same & 205--224;488 & split degree-one G-place conditions, distinct primes, and g\_big \\
same & 315--323 & IsMu0 finite restrictions and tameness \\
\nolinkurl{Formalism/LocalToGlobal/UnitaryNorms/IsUnitaryNormOf.lean} & 66--75 & LevelLe, UnrAt, TameAt \\
\nolinkurl{Formalism/LocalRingCardinality/Params.lean} & 5--20;30--40 & ranks and \(N_j=k+j(t+1)\) \\
\nolinkurl{Formalism/LocalToGlobal/Columns/FiniteSumNonzeroLargeDepth.lean} & 1017--1023 & only the initial levG is tied to colLev by zero\_G \\
\end{longtable}

For \(1\le j<m\), put \(n=N_j=k+j(t+1)\) and \(r_j=k+j\). The literal source has \[
 \begin{split}
 H_{X,j}&=\operatorname{uniUT}
                (n,\ r_j,\ [t,\ldots,t]),\\
 n&=r_j+jt,\\
 \lambda_{X,j}
       &=\psi_{S_X}\prod_{i=0}^{j-1}\chi_i(\det h_i).
 \end{split}
 \tag{F12}
\] The \(j\) determinant blocks have size \(t\), and the first block is \(U_{r_j}\). The character \(\psi_{S_X}\) is supported on the displayed selected simple-root entries of that first unitriangular block. It is trivial on the interblock radical. Thus(F12) is exactly the situation of(F1)--(F2), with \(H_0=U_{r_j}\). The omission of one simple root in \(S_X\) at some stages does not affect the argument, since \(\lambda_0\) in the finite lemma was arbitrary.

There is no additional unused trailing block: the identity \(n=r_j+jt\) is exact. The terminal declaration at \(j=m\) uses HXm and an additional fixed-column condition, so the theorem is not extended to that branch without a separate calculation.

After choosing the coefficient-field embedding into \(\mathbf C\), the characters have the same orders as in the source field. The automorphic convention may require \(\lambda_{X,j}\) or its inverse; both are excluded by the same proof.

\subsection{5. Every prescribed determinant character is nontrivial}\label{every-prescribed-determinant-character-is-nontrivial}

Let \(P\) be a selected G-place over the rational auxiliary prime \(q=g_\epsilon\), \(\epsilon\in\{0,1\}\). The source has residue field of cardinality \(q\), where \(q\) is prime. The prime-field map therefore identifies its residue units with \(\mathbf F_q^\times\).

The chi\_om condition specifies a faithful character \(\omega:\mathbf F_q^\times\to K^\times\) such that \[
                        \chi_i=\omega^{\,2(i+1)}
                         \qquad(0\le i<m).
 \tag{F13}
\] This is the exact source character, not a new choice. Its order is \[
                \operatorname{ord}(\chi_i)
                   =\frac{q-1}{\gcd(q-1,\,2(i+1))}.
 \tag{F14}
\] The source size and rank inequalities give \[
       q-1>2fk\ge2k>2m\ge2(i+1)>0.
 \tag{F15}
\] Here \(f\ge1\), and \(k\ge m+3\). Hence every \(\chi_i\) is nontrivial, and the characters for distinct \(i<m\) are distinct. The latter also appears as the separate chi\_inj field. This proves the nontriviality hypothesis of(F9) at every G-place for every \(1\le j<m\).

The conclusion for the literal column type is therefore \[
 \boxed{\pi^{\widetilde H_{X,j},\lambda_{X,j}^{\pm1}}=0
   \quad(1\le j<m)}
 \tag{F16}
\] for every unramified local representation at such a place. For \(2\le j<m\), the same vanishing holds for every tame determinant twist of an unramified representation, by the distinctness and(F11).

For the actual host \(E'=\mathbf Q(\zeta_{28})\), \(f=6\) and the admissible even rank has \(k\ge6\). Thus \(q-1>72\), so both prescribed G primes are different from29. Every local constituent of the chain in LIFT14--16 is unramified at those places. Consequently its nonzero stage-\(j\) module cannot itself lie in the explicit source ColV with this finite type, for \(1\le j<m\). Deepening the29 level, choosing a nonspherical vector in the same local representation, or taking any further finite group translate within that representation cannot remove(F16).

This is a local representation obstruction, not an inequality between gains.

\subsection{6. The initial stage and the limits of the obstruction}\label{the-initial-stage-and-the-limits-of-the-obstruction}

At \(j=0\), the source has \(H_{X,0}=U_k\) and its full nontrivial generic-unipotent character, because \(m\ge2\). There are no determinant blocks. A \(K\)-fixed vector is fixed by \(U_k\); its nontrivial character projector is therefore zero. A determinant twist does not change this calculation on \(U_k\).

This initial-vector statement does not imply that every vector of an unramified representation is killed. The finite flag permutation module can have a generic-unipotent component, and section1 gives no exclusion when there are no nontrivial determinant blocks. A source-type projection or different vector in the same representation must therefore be tested separately at \(j=0\).

The standard initial class-supported vector of AM-R-012 is spherical at the auxiliary G-places, so it is not itself a literal initial generic-type vector there. This is distinct from the stronger entire-representation obstruction at \(j\ge1\).

The source's ChIdx record requires only zero\_G: \[
                            \mathrm{levG}_0=\mathrm{colLev}_0.
 \tag{F17}
\] Its later levG\_j and conjugating indices are additional data. Therefore(F16) must not be promoted to a theorem that every permitted raw PhiG/PhiGS or Xi kernel annihilates every vector in the existing chain. That stronger assertion would require its actual right-source covariance and the identification of its later finite level. The direct raw-Xi identity and auxiliary finite interpolation in AM-DO-029/030 keep that question separate.

\subsection{7. Exact finite conductor forced by the source character pin}\label{exact-finite-conductor-forced-by-the-source-character-pin}

The existing source IsMu0 condition already requires a different finite character family from the conductor29 family in LIFT14. At each selected G-place its unit restriction is \[
                         \mu_i|_{\mathcal O_{E',P}^{\times}}
                                  =\chi_i^{-1}.
 \tag{F18}
\] It is tame there and unramified away from the two G primes. This was also audited in AM-DO-030; it is a literal source condition, not a new character-existence theorem.

The following conductor accounting follows from(F13)--(F18). The source atg condition, applied to integral representatives of each nonzero residue class, gives the displayed unit restriction; tameness makes it a character of precisely that finite quotient. At every selected G-place the character is therefore nontrivial on units but trivial on principal units, so its conductor exponent is exactly one. At the conjugate place the restriction is the inverse, because \(\mu_i\) is trivial on the diagonal base-field ideles; it is also nontrivial and has conductor exponent one.

For completeness, the degree-one prime, norm, cardinality and coprimality conditions really do account for every place above each \(g_\epsilon\). The two principal ideals generated by \(\pi_G\) and its conjugate are coprime, have product \((g_\epsilon)\), and have equal norms. Each norm is therefore \(g_\epsilon^f\). Each selected half already contains \(f\) distinct prime ideals of norm \(g_\epsilon\), so all their exponents are one and no additional prime factor is possible. Conjugation gives the other half.

Thus any genuine realization of this source character pin on the actual degree12 host has \[
 \mathfrak f(\mu_i)=(g_0g_1)\mathcal O_{E'},\qquad
                  N_{E'/\mathbf Q}\mathfrak f(\mu_i)=(g_0g_1)^{12}.
 \tag{F19}
\] The equality is conditional on the genuine source IsMu0 realization; it does not assert that the global unit compatibility needed for that realization has been solved. It records the exact finite conductor required if it is solved.

This is useful normalization data, but it is not by itself a lower bound for the full H term. A unitary determinant character can act isometrically in a suitable class presentation, and the changed kernels, levels, local tests and arithmetic content must still be compared. AM-R-008's integral isometry has its own hypotheses and does not identify the altered row matrices.

\subsection{8. What a repair would have to change}\label{what-a-repair-would-have-to-change}

A compatible new chain must admit the prescribed nontrivial local determinant blocks at the G-places. The current spherical constituents cannot acquire them merely by selecting vectors or deepening level elsewhere. From stage2 onward, one tame scalar determinant twist cannot replace the distinct blocks either.

A natural remaining route is to use globally admissible stage characters ramified as in(F18) and local theta constituents with their correspondingly ramified inserted blocks. Type-II local theta occurrence itself remains available at these split places. That observation does not prove the required local type, the exact raw-Xi action on the previous image, or global nonvanishing with the prescribed finite kernel. It also does not solve the independently audited odd-character etaT pin.

The full new result is(F16), its tame-twist strengthening, and the exact conditional conductor(F19). It neither disproves a different actual mixed theta tower nor supplies the missing \(g_j,e_j,H\) arrays, terminal row, six arithmetic identities, or a shorter DP44 chain.

\clearpage

\section{AM-K029 --- a degree-seven repair of the literal unramified etaT character}\label{am-k029-a-degree-seven-repair-of-the-literal-unramified-etat-character}

Author: GPT-6 Astra Pro, ALP\_MACHINE \nolinkurl{/root/kummer_extension}, 21 November 2026 UTC. Inherited reasoning effort is not separately exposed.

The root-unit obstruction of AM-K028 can be repaired on an explicit larger field while keeping the ordinary primitive-J Hodge type. The following construction proves the full character interface, including the base-field restriction \nolinkurl{ResEta}; it does not merely remove one necessary unit test.

Write \[
 E_0=\mathbb Q(\zeta_{28}),\quad F_0=E_0^+,
 \quad\Phi_0=\{1,3,9,15,17,23\},
\] and let B be the unique degree-seven subfield of \(\mathbb Q(\zeta_{29})\). Put \[
                         K=E_0B,\qquad F=F_0B.
\tag{ER1}
\] Give K the induced CM type \(\Psi\), containing all seven lifts of each member of \(\Phi_0\).

\subsubsection{Theorem ER2}\label{theorem-er2}

The extension K/F is CM and unramified at every finite place, \([K:\mathbb Q]=84\), and \([F:\mathbb Q]=42\). There exists a continuous unitary Hecke character \[
                   \eta:\mathbb A_K^\times/K^\times\longrightarrow S^1
\] that is unramified at every finite place, has infinity component \[
        \eta_\infty(z)=\prod_{\tau\in\Psi}\frac{z_\tau}{|z_\tau|},
\tag{ER2}
\] and restricts to the quadratic character \(\eta_{K/F}\) on \(\mathbb A_F^\times\). Thus it satisfies the literal \nolinkurl{IsEtaT} convention checked in AM-K028. Its normalized finite values \[
                 \eta(\varpi_P)\,(NP)^{-1/2}
\tag{ER3}
\] all lie in a single finite algebraic coefficient field, which may be enlarged to contain its complex conjugates. Therefore the separate \nolinkurl{etaVal} field can also be supplied after that coefficient-field enlargement.

No class-number calculation for the degree-84 or degree-42 fields is needed for this existence proof. No bound on the degree, discriminant, or heights of the auxiliary coefficient field is asserted. The full polarized host, original \nolinkurl{IsTen}, all source size predicates, and later selected-image bounds are separate conditions.

\subsection{1. The field and its roots of unity}\label{the-field-and-its-roots-of-unity}

The group \((\mathbb Z/29\mathbb Z)^\times\) is cyclic of order 28. Its subgroup of order four contains complex conjugation and fixes B; thus B is totally real, cyclic of degree seven, and ramified only at 29, with ramification index seven there. The fields E₀ and B have coprime degrees 12 and 7, so are linearly disjoint. Hence K has degree 84, and its conjugation fixes exactly F, of degree 42.

Accepted DP57 proves that E₀/F₀ is unramified at every finite place. This property is preserved by base change. More concretely, locally a finite étale extension of a DVR stays étale after base change, and its finite field components give the completions of K/F. Hence K/F is everywhere finite-unramified.

The rational ramification indices in K are \[
                     e_2=2,\qquad e_7=6,\qquad e_{29}=7,
\tag{ER4}
\] and all other finite rational primes are unramified. These follow from the standard cyclotomic indices in \(\mathbb Q(\zeta_{28})\), the totally ramified degree-seven extension at 29, and their disjoint ramification sets.

It follows that \[
                                  \mu(K)=\mu_{28}.
\tag{ER5}
\] Indeed a primitive \(\ell^a\)-th root in K would force the cyclotomic ramification index \(\varphi(\ell^a)\) to divide the corresponding index in (ER4). No odd prime other than 7 can occur: at 29 the required index 28 does not divide 7, and all remaining odd primes are unramified. A 49th root would require index 42 at 7, and an eighth root would require index four at 2. Thus the root order divides 28; E₀ already contains all these roots. The same argument gives \(\mu(E_0)=\mu_{28}\).

For a CM field L and a unit u, \(u/\bar u\) is a root of unity. Here is the elementary fact used: all conjugates of this algebraic integer unit have absolute value one. The characteristic polynomials of its positive powers have bounded integer coefficients, so only finitely many such powers can occur; two powers are equal, proving finite order. This also gives the usual CM-unit index bound below without a unit-group classification theorem.

\subsection{2. The entire unit phase image in E₀}\label{the-entire-unit-phase-image-in-eux2080}

Define \[
 a_0(x)=\prod_{\phi\in\Phi_0}\frac{\phi(x)}{|\phi(x)|}
                       \quad(x\in E_0^\times).
\tag{ER6}
\] We prove the exact identity \[
                             a_0(\mathcal O_{E_0}^\times)=\mu_7.
\tag{ER7}
\]

First every real unit \(\varepsilon\in\mathcal O_{F_0}^\times\) has phase one. The canonical quadratic Artin character of the known extension E₀/F₀ is unramified at every finite place and has local sign at every real place. On a principal real unit, its finite value is one, so global reciprocity gives \(\operatorname{sgn}N_{F_0/\mathbb Q}\varepsilon=1\). This sign is exactly \(a_0(\varepsilon)\). This uses the quadratic character on \textbf{F₀}, which exists by class field theory; it does not presume the obstructed \nolinkurl{etaT} character on E₀.

Next put \(u_0=1-\zeta_{28}\). It is a unit because \[
 \Phi_{28}(X)=X^{12}-X^{10}+X^8-X^6+X^4-X^2+1,
                      \qquad \Phi_{28}(1)=1.
\] Its conjugation quotient is \[
                       \frac{u_0}{\bar u_0}=-\zeta_{28}=\zeta_{28}^{15}.
\tag{ER8}
\] For an arbitrary unit u, write \(u/\bar u=\zeta_{28}^a\). Choose \(\epsilon\in\{0,1\}\) with \(a-15\epsilon\) even and choose j with \(2j\equiv a-15\epsilon\pmod{28}\). Then \(u/(\zeta_{28}^j u_0^\epsilon)\) is fixed by conjugation and is an integral unit. Thus \[
 \mathcal O_{E_0}^\times
   =\mu_{28}\,\mathcal O_{F_0}^\times\,\{1,u_0\}.
\tag{ER9}
\] This proves exactly the needed Hasse-index-two description; no missing generator is assumed.

On \(\mu_{28}\), the phase is \(a_0(\zeta_{28}^j)=\zeta_{28}^{68j}\), whose image has order seven. For \(0<a<28\), \[
 \frac{1-e^{2\pi ia/28}}{|1-e^{2\pi ia/28}|}
                             =-i e^{\pi ia/28}.
\] Using the six integers in \(\Phi_0\), this gives \[
                    a_0(u_0)=(-i)^6\zeta_{56}^{68}=\zeta_7^5.
\tag{ER10}
\] Equations (ER9)--(ER10), and the real-unit result, prove (ER7).

\subsection{3. All new units of K satisfy the required phase condition}\label{all-new-units-of-k-satisfy-the-required-phase-condition}

Let \[
 a(x)=\prod_{\tau\in\Psi}\frac{\tau(x)}{|\tau(x)|}
                   =a_0(N_{K/E_0}x).
\tag{ER11}
\] The equality follows by multiplying over all lifts of each embedding; the absolute values multiply in the same way.

For a unit U of K, put \(v=N_{K/E_0}U\), an E₀-unit. By (ER5), write \(U/\bar U=\zeta\in\mu_{28}\). Since this root lies in E₀, \[
                       \frac v{\bar v}=\zeta^7.
\] Consequently \[
 a(U)^2=a_0(v)^2
       =\prod_{\phi\in\Phi_0}\phi(v/\bar v)
       =a_0(\zeta^7)=1.
\tag{ER12}
\] On the other hand, (ER7) gives \(a(U)=a_0(v)\in\mu_7\). The only element common to \(\mu_2\) and \(\mu_7\) is one. Therefore \[
                            a(\mathcal O_K^\times)=\{1\}.
\tag{ER13}
\] This treats every new unit in the degree-seven extension, not only the roots inherited from E₀.

\subsection{4. Extending a unit-compatible angular type to an unramified character}\label{extending-a-unit-compatible-angular-type-to-an-unramified-character}

We use the following elementary ideal construction. Let L be a CM field, let b be a unitary angular character of \((L\otimes\mathbb R)^\times\), and suppose \(b(\mathcal O_L^\times)=1\). On the principal fractional ideals define \[
                         c((x))=b(x)^{-1}.
\tag{ER14}
\] It is well-defined by the unit condition. It extends to a homomorphism \(c:I_L\to S^1\) on all fractional ideals. To see this without an abstract extension theorem, write the finite class group as a direct sum of cyclic groups, choose representative ideals \(\mathfrak a_j\) of their generators and relations \(\mathfrak a_j^{h_j}=(x_j)\), and choose roots \(c(\mathfrak a_j)^{h_j}=b(x_j)^{-1}\) in \(S^1\). Every ideal has a principal remainder after choosing the canonical ranges of the generator exponents. The displayed root equations make its value independent of carries and make multiplication exact.

For a finite idele \(z_f\), let \(I(z_f)=\prod_PP^{v_P(z_P)}\). Then \[
                         \chi(z)=b(z_\infty)c(I(z_f))
\tag{ER15}
\] is continuous, unitary, unramified at all finite places, and trivial on L× by (ER14). The ideal-valued map is locally constant, and b is continuous, so continuity has not been omitted. A character with trivial infinity component and unramified finite components similarly factors through the ordinary finite ideal class group.

By (ER7), this constructs an unramified unitary character \(\chi_7\) on E₀ with angular type \(a_0^7\). By (ER13), it constructs a character \(\chi_0\) on K with angular type a. The latter already has all desired clauses except the exact restriction to F.

\subsection{5. Repairing the restriction without a large class-number calculation}\label{repairing-the-restriction-without-a-large-class-number-calculation}

The only numerical class-number input is \[
                                  h(F_0)=1.
\tag{ER16}
\] This is explicitly tabulated in Masley (1978), Table 3, printed p.316: conductor 28, degree 6, class number 1. I inspected that primary table image. The field is the maximal real subfield of E₀, whose Galois group is cyclic of order six. No conjectural class-number bound is used.

Let \(\eta_0=\eta_{E_0/F_0}\). The restriction of \(\chi_7\) to F₀ has sign at each real place, since seven is odd, and is unramified at all finite places. Hence \((\chi_7|_{\mathbb A_{F_0}^\times})/\eta_0\) has trivial infinity component and factors through \(\mathrm{Cl}(F_0)\). By (ER16), \[
                        \chi_7|_{\mathbb A_{F_0}^\times}=\eta_0.
\tag{ER17}
\]

Put \(\eta_1=\eta_{K/F}\) and \[
 \beta=\chi_7\circ N_{K/E_0},\qquad
 \delta=\chi_0^7\beta^{-1},\qquad
 \tau=(\chi_0|_{\mathbb A_F^\times})\eta_1^{-1}.
\tag{ER18}
\] The norm pulls local units to local units. Thus β is unramified everywhere finite and its infinity component is \(a^7\). It follows that δ is an ordinary class character of K. Likewise τ is an ordinary class character of F, because the infinity component of \(\chi_0|_F\) is the product of signs, exactly that of \(\eta_1\).

The quadratic base-change identity gives \[
                \beta|_{\mathbb A_F^\times}
             =\eta_0\circ N_{F/F_0}=\eta_1.
\tag{ER19}
\] It can also be checked locally here. At a finite prime of F₀ where E₀ splits, both sides are trivial. At an inert finite prime E₀/F₀ is unramified quadratic; after a residue extension of degree f it splits precisely when f is even. The unramified value of \(\eta_0\) on the norm of a uniformizer is \((-1)^f\), the same as the new quadratic character. Both sides are trivial on units. At infinity the norm and sign give the same product of local signs. This proves (ER19), including primes where F/F₀ itself is ramified.

Let \(j:\mathrm{Cl}(F)\to\mathrm{Cl}(K)\) be extension of ideals. From (ER18)--(ER19), \[
                                  \delta\circ j=\tau^7.
\tag{ER20}
\] The kernel of j has exponent dividing two: applying the ideal norm to an extended ideal gives its square. If x lies in this kernel, then (ER20) gives \(\tau(x)^7=1\), while \(x^2=1\) gives \(\tau(x)^2=1\). Therefore \(\tau(x)=1\). Thus \(\tau^{-1}\) is a well-defined character on \(\operatorname{im}j\).

Every character of a subgroup of a finite abelian group extends to the whole group: adjoining one generator at a time, choose a root of the specified value of its least positive power lying in the previous subgroup. All these values remain roots of unity. Apply this to obtain a class character \(\lambda\) of K with \[
                                   \lambda\circ j=\tau^{-1}.
\tag{ER21}
\] Finally set \[
                                      \eta=\chi_0\lambda.
\tag{ER22}
\] It is unitary, unramified at every finite place, has infinity component a, and its restriction to F is \(\tau\eta_1\tau^{-1}=\eta_1\). This proves all four literal \nolinkurl{IsEtaT} clauses. The argument has handled the capitulation compatibility explicitly; it has not inferred \nolinkurl{ResEta} merely from the unit condition.

\subsection{6. A single finite coefficient field for etaVal}\label{a-single-finite-coefficient-field-for-etaval}

The finite values required in the source are normalized by \((NP)^{-1/2}\). Their common algebraicity can be built into the preceding construction.

Define the algebraic type norm \[
                T(x)=\prod_{\phi\in\Phi_0}\phi(N_{K/E_0}x)\in E_0^\times.
\tag{ER23}
\] Since E₀ is Galois, every factor and the product lie in E₀. For the chosen complex embedding, \[
              a(x)=\frac{T(x)}{\sqrt{N_{K/\mathbb Q}x}},
\tag{ER24}
\] where the norm is positive because K is CM. In particular \(T(U)=1\) on every unit by (ER13).

The principal-ideal rule \(\psi((x))=T(x)^{-1}\) is therefore well-defined. Extend it to all ideals by the cyclic-class construction of §4, choosing only the finitely many roots \[
                         \psi(\mathfrak a_j)^{h_j}=T(x_j)^{-1}.
\tag{ER25}
\] All ψ-values then lie in one finite extension of E₀. Every chosen root has complex absolute value \((N\mathfrak a_j)^{-1/2}\), since \(\mathfrak a_j^{h_j}=(x_j)\); the same absolute-value identity holds on all ideals. Thus \(c(I)=\psi(I)\sqrt{NI}\) is the unitary ideal character used to construct \(\chi_0\).

The correcting class character λ has finite order. Adjoining its finitely many root-of-unity values, one obtains a fixed finite coefficient field containing \[
                        \eta(\varpi_P)(NP)^{-1/2}
                                      =\psi(P)\lambda(P)
\tag{ER26}
\] for every finite prime P. Closing that field under conjugation is still a finite extension. This supplies the existence part of the source \nolinkurl{etaVal} condition, with its exact \nolinkurl{valTw(1/2)} normalization. No bound for the necessary class orders or this coefficient field is part of the theorem.

\subsection{7. Persistence under further real extensions and remaining host work}\label{persistence-under-further-real-extensions-and-remaining-host-work}

The repaired character persists under any further totally real extension \(F_1/F\). Put \(K_1=KF_1\), use the induced CM type, and pull η back by \(N_{K_1/K}\). Its infinity exponent remains one at every induced embedding, its finite components remain unramified, and the same local quadratic base-change check proves the exact restriction \(\eta_{K_1/F_1}\). Thus the character obstruction need not be solved again when a further real extension is required for other numerical host conditions.

This construction concerns the ordinary type of a possible \(J^{14}\) host. It proves the auxiliary character and its normalized value field, with no assertion here that the full integral polarized lattice is the standard product lattice or that all old source real, rank, tensor, and size records are instantiated. Those additional comparisons must retain their own finite and archimedean costs. The degree-seven step repairs one actual previously impossible field of the source adapter; it does not complete the GRH extension or the MainX window.

\subsection{8. Primary references and precise use}\label{primary-references-and-precise-use}

\begin{itemize}
\tightlist
\item
  Masley, J. M. (1978). Class numbers of real cyclic number fields with small conductor. \emph{Compositio Mathematica, 37}(3), 297--319. Table 3, printed p.316, conductor28/degree6 row: \(h(F_0)=1\). Primary PDF: \url{https://numdam.org/item/CM_1978__37_3_297_0.pdf}. The table image was inspected directly, rather than relying on search snippets or a secondary class-number list.
\item
  Milne, J. S. (2020). \emph{Class field theory} (Version4.03, 6 August 2020). Chapter V, Proposition4.12, pp.177; Proposition5.2 and Theorem5.3, pp.178--179: the quadratic Artin character, its local restrictions, principal-idele triviality, and norms. \url{https://www.jmilne.org/math/CourseNotes/CFT.pdf}. Only these standard class-field inputs are used; the unit generation, character-extension, and capitulation correction are proved above.
\item
  Accepted DP57 and AM-DO-022: the original ordinary J² field/type and finite-unramified CM extension E₀/F₀. AM-K028: exact source \nolinkurl{HasInf}, \nolinkurl{UnrAt}, \nolinkurl{ResEta}, \nolinkurl{IsEtaT}, and \nolinkurl{etaVal} conventions and their provenance.
\end{itemize}

The exact finite arithmetic checks are separate from the proof of character existence. No Lean compilation, new analytic L-function theorem, large-field class-number computation, or effective coefficient-height estimate is claimed.

\clearpage

\section{AM-K030 --- an explicit principal product host for the degree-seven character repair}\label{am-k030-an-explicit-principal-product-host-for-the-degree-seven-character-repair}

Author: GPT-6 Astra Pro, ALP\_MACHINE \nolinkurl{/root/kummer_extension}, 21 November 2026 UTC. Inherited reasoning effort is not separately exposed.

AM-K029 supplies the literal unramified \nolinkurl{IsEtaT} character on an induced-type degree-84 CM field. This finding supplies its actual integral principally polarized ordinary-Jacobian host. The identification is an isomorphism of principally polarized complex abelian varieties, with a full integral CM action. It is stronger than an isogeny or a rational trace-form comparison.

Keep \[
 E_0=\mathbb Q(\zeta_{28}),\quad F_0=E_0^+,\quad
 \Phi_0=\{1,3,9,15,17,23\},\quad A=J^2,
\] where \(J:y^2=x^7-1\) has its actual principal polarization. Let B be the unique degree-seven subfield of \(\mathbb Q(\zeta_{29})\), let \(K=E_0B\), and let \(F=F_0B\). The induced type \(\Psi\) consists of all seven lifts of every embedding in \(\Phi_0\).

Accepted DP57 and AM-DO-022 identify A with the type-\(\Phi_0\) CM torus whose lattice is \(\mathcal O_{E_0}\) and whose exact principal trace coefficient is \[
 \xi_0=\frac{\zeta_7-\zeta_7^{-1}}{14}.
\tag{PH1}
\] All sign conventions below use that same alternating form. Nothing changes its source/target Hom convention.

\subsection{1. An explicit selfdual normal basis}\label{an-explicit-selfdual-normal-basis}

Set \(\zeta=\zeta_{29}\), \[
 H=\{1,12,17,28\}\subset(\mathbb Z/29\mathbb Z)^\times,
 \qquad C_i=2^iH\quad(i\in\mathbb Z/7\mathbb Z),
\] and define the Gaussian periods \[
             \eta_i=\sum_{h\in H}\zeta^{2^ih}.
\tag{PH2}
\] The integer 2 has order 28 modulo 29, and H is its subgroup of order four. In particular H contains \(-1\), so the periods and B are real. Let \(\sigma\) be the generator of \(\operatorname{Gal}(B/\mathbb Q)\) induced by \(\zeta\mapsto\zeta^2\). Then \(\sigma\eta_i=\eta_{i+1}\).

The following seven elements are completely explicit: \[
 \begin{split}
  a&=(-8,-4,-2,-5,-2,-4,-4),\\
  e_0&=\frac1{29}\sum_{i=0}^6a_i\eta_i,\qquad
  e_j=\sigma^j(e_0)\quad(0\le j\le6).
 \end{split}
\tag{PH3}
\]

\subsubsection{Theorem PH4}\label{theorem-ph4}

Let \(\mathfrak p\) be the unique prime of B above 29. Then \[
 \mathfrak D_{B/\mathbb Q}=\mathfrak p^6,
 \qquad \mathcal A_B:=\mathfrak p^{-3},
 \qquad
 \mathcal A_B=\bigoplus_{j=0}^6\mathbb Z e_j,
 \qquad \operatorname{Tr}_{B/\mathbb Q}(e_je_k)=\delta_{jk}.
\tag{PH4}
\] Every \(e_j\) has \(\mathfrak p\)-valuation exactly \(-3\), and \(\operatorname{Tr}_{B/\mathbb Q}e_j=1\).

\textbf{Proof of the ramification and trace identities.} The prime 29 is totally ramified in \(\mathbb Q(\zeta)\), and hence is totally ramified in its degree-seven subfield B. There are no other ramified primes. This extension is tame because \(29\nmid7\); the different exponent is \(7-1=6\). These are the standard cyclotomic and tame-different facts; precise references appear below.

The sum of the seven periods is \(-1\). Therefore \(\operatorname{Tr}\eta_i=-1\). To compute \(\operatorname{Tr}(\eta_i\eta_j)\), sum the products over the seven simultaneous translates of the two cosets. The coefficient of \(\zeta^t\) is constant for nonzero residues t, since multiplication by any nonzero residue permutes all terms. If \(i=j\), there are 28 terms with exponent zero: for each of the seven cosets, there are four ordered pairs \((x,-x)\). The remaining 84 terms distribute three to every nonzero residue. If \(i\ne j\), there is no exponent-zero term and all 112 terms distribute four to every nonzero residue. Hence \[
                  \operatorname{Tr}(\eta_i\eta_j)=29\delta_{ij}-4.
\tag{PH5}
\] In particular the periods are linearly independent, their trace determinant is \(29^6\), and, by the different calculation, they are an integral basis of \(\mathcal O_B\).

For the vector a in (PH3), direct integer multiplication gives \[
 \sum_i a_i=-29,\qquad
 \sum_i a_i a_{i+k}=\begin{cases}145&k=0,\\116&k\ne0.\end{cases}
\tag{PH6}
\] One may check the shorter equivalent computation with \(b=a+4(1,\ldots,1)=(-4,0,2,-1,2,0,0)\): its sum is \(-1\), and its cyclic correlations are 25 at zero and \(-4\) elsewhere. Combining (PH5) and (PH6) yields \[
 \operatorname{Tr}(e_0e_k)
 =\frac{29\sum_i a_i a_{i-k}-4(\sum_i a_i)^2}{29^2}
 =\delta_{0k}.
\tag{PH7}
\] The trace-one assertion follows from \(\operatorname{Tr}\eta_i=-1\). Galois invariance gives all entries in (PH4).

\textbf{Proof of membership in the square-root inverse different.} Write \[
             P(X)=\sum_i a_i\sum_{h\in H}X^{(2^ih\bmod29)},
\] using representatives in \(\{1,\ldots,28\}\). The exact polynomial identity needed for the local check is \[
\begin{split}
 P(1+T)\equiv T^{16}\bigl(&1-8T-13T^2+2T^3+7T^4-14T^5
 +14T^6-5T^7\\
 &-12T^8+6T^9-4T^{10}+4T^{11}-8T^{12}\bigr)\pmod{29}.
\end{split}
\tag{PH8}
\] This is a 29-term binomial computation; the accompanying certificate verifies every coefficient, not just the first nonzero one. In the completion of \(\mathbb Q(\zeta)\) at 29, \(t=\zeta-1\) has valuation one and 29 has valuation 28. Equation (PH8) shows that \(P(\zeta)\) has valuation exactly 16: all terms of lower degree have valuation at least 28, its degree-16 coefficient is a unit, and all higher-degree terms have strictly larger valuation. The extension \(\mathbb Q(\zeta)/B\) has ramification degree four, so \[
                    v_{\mathfrak p}(P(\zeta))=4,
                    \qquad v_{\mathfrak p}(e_0)=4-7=-3.
\tag{PH9}
\] At every other finite prime, e₀ is integral because its only rational denominator is 29. The same statements hold for all its conjugates. Thus \(L:=\sum\mathbb Z e_j\subset\mathcal A_B\).

The trace-dual of a fractional ideal I is \(I^*=\mathfrak D_{B/\mathbb Q}^{-1}I^{-1}\). Consequently \(\mathcal A_B^*=\mathcal A_B\). Equation (PH7) says \(L^*=L\). Taking duals of \(L\subset\mathcal A_B\) gives \(\mathcal A_B\subset L\), proving equality. This finishes the proof. No classification of positive unimodular lattices, approximate lattice search, or unproved integral-normal-basis assertion is needed. ∎

\subsection{2. A positive generator and an actual integral CM action}\label{a-positive-generator-and-an-actual-integral-cm-action}

Define \[
 \pi_B=N_{\mathbb Q(\zeta)/B}(1-\zeta)
       =\prod_{h\in H}(1-\zeta^h)
       =4-2\eta_0+\eta_4.
\tag{PH10}
\] The last identity is obtained by multiplication and the relation \(\sum_i\eta_i=-1\); the certificate checks it exactly in the cyclotomic algebra.

The factors with exponents \(h\) and \(-h\) multiply to \(|1-\zeta^h|^2>0\). This remains true at every real embedding of B. Thus \(\pi_B\) is totally positive. Norm transitivity and \(N_{\mathbb Q(\zeta)/\mathbb Q}(1-\zeta)=29\) give \[
              N_{B/\mathbb Q}\pi_B=29,
              \quad \mathfrak p=(\pi_B),
              \quad \mathcal A_B=\pi_B^{-3}\mathcal O_B.
\tag{PH11}
\]

Multiplication by any \(\beta\in\mathcal O_B\) on the basis e gives an integral matrix: \(\mathcal A_B\) is an \(\mathcal O_B\)-ideal. Its entries are explicitly \[
                 (M_\beta)_{jk}=\operatorname{Tr}(\beta e_k e_j),
\tag{PH12}
\] and are symmetric. The matrices give a faithful ring embedding \(\mathcal O_B\hookrightarrow M_7(\mathbb Z)\), with every B-element self-adjoint for the standard Euclidean form. As an optional directly usable generator check, multiplication by \(\eta_0\) is \[
 M_{\eta_0}=\begin{pmatrix}
 -2&0&-1&0&0&0&-1\\
 0&0&0&-1&0&0&1\\
 -1&0&0&1&0&1&-1\\
 0&-1&1&1&0&1&0\\
 0&0&0&0&0&-1&0\\
 0&0&1&1&-1&0&1\\
 -1&1&-1&0&0&1&0
 \end{pmatrix}.
\tag{PH13}
\] The conjugate-period matrices are obtained by conjugating this matrix by the cyclic permutation \(e_j\mapsto e_{j+1}\). Since the periods form an integral basis, these matrices specify the full \(\mathcal O_B\)-action. The characteristic polynomial is \[
 X^7+X^6-12X^5-7X^4+28X^3+14X^2-9X+1.
\tag{PH14}
\] Neither the characteristic polynomial nor a numerical eigenvalue computation is being used to infer the integral-lattice theorem; it was proved above.

\subsection{3. The actual product principal polarization}\label{the-actual-product-principal-polarization}

The discriminants of E₀ and B are coprime, and the degrees 12 and 7 are coprime. The integral-compositum theorem therefore gives \[
            \mathcal O_K=\mathcal O_{E_0}\otimes_{\mathbb Z}\mathcal O_B,
 \qquad |d_K|=|d_{E_0}|^7 |d_B|^{12}
                  =2^{84}7^{70}29^{72}.
\tag{PH15}
\] Consider the full lattice \[
       \Lambda=\mathcal O_{E_0}\otimes_{\mathbb Z}\mathcal A_B
               =\pi_B^{-3}\mathcal O_K\subset K.
\tag{PH16}
\] It is stable under \(\mathcal O_K\). On pure tensors define the alternating form by the product of the old principal alternating form and the real trace pairing of B. Equivalently, in K it is the trace form with coefficient \(\xi_0\). For arbitrary \(x_j,y_j\in E_0\), the exact identity is \[
 \operatorname{Tr}_{K/\mathbb Q}
  \left(\xi_0\Bigl(\sum_jx_je_j\Bigr)
                    \overline{\Bigl(\sum_ky_ke_k\Bigr)}\right)
       =\sum_{j=0}^6\operatorname{Tr}_{E_0/\mathbb Q}
                           (\xi_0x_j\bar y_j).
\tag{PH17}
\] It follows directly from (PH4), because conjugation fixes B. Thus the tensor alternating form is integral and unimodular on \(\Lambda\). Its positivity follows from the positivity of the old Hermitian form and of the real trace form. With the induced complex structure of \(\Psi\), the map \[
         (x_0,\ldots,x_6)\longmapsto\sum_jx_je_j
\tag{PH18}
\] is an isomorphism of polarized complex tori \[
      (A^7,\lambda_A^{\boxtimes7})
           \simeq (K\otimes\mathbb R/\Lambda,\lambda_\Lambda)
           \simeq (J^{14},\lambda_J^{\boxtimes14}).
\tag{PH19}
\] For the complex-linearity assertion, B acts through its regular real representation on the seven copies; E₀ supplies the same complex structure on each copy. Diagonalizing that real representation gives exactly all seven lifts of every old selected embedding. This is the stated induced type, not a change to conjugate Jacobian factors.

The commuting integral actions of \(\mathcal O_{E_0}\) and \(\mathcal O_B\), together with (PH15), supply the full integral \(\mathcal O_K\)-action on the actual product. Rosati acts as conjugation on E₀ by DP57 and as the identity on B by (PH12); it is therefore complex conjugation on K. The field degree is \(84=2\dim J^{14}\), so this is a full CM action.

If the source uses the standard lattice \(\mathcal O_K\), transport (PH16) by multiplication by \(\pi_B^{-3}\). The exact coefficient then becomes \[
             \boxed{\xi_K=\xi_0\pi_B^{-6}.}
\tag{PH20}
\] Since \(\pi_B\) is totally positive, this has the same selected positivity signs. Moreover \[
        \xi_K\mathcal O_K=\mathfrak D_{K/\mathbb Q}^{-1},
       \qquad |N_{K/\mathbb Q}\xi_K|
                      =2^{-84}7^{-70}29^{-72}.
\tag{PH21}
\] The ideal equality also follows from the established perfect trace pairing, so it is stronger than a mere equality of absolute norms. The factor \(\pi_B^{-6}\) is part of the new host data. Keeping the old \(\xi_0\) while silently replacing \(\Lambda\) by \(\mathcal O_K\) would give the wrong principal lattice.

\subsection{4. The theta line and the CM coordinate comparison}\label{the-theta-line-and-the-cm-coordinate-comparison}

Use the same trace-one element \[
                      w_0=-(\zeta_7+\zeta_7^2+\zeta_7^4)
                      \quad(w_0+\bar w_0=1).
\] It acts diagonally in (PH18). Equation (PH17) splits its actual trace-one semicharacter into the product of the old semicharacters. Thus the rigidified symmetric principal line is exactly the external product of the fourteen ordinary Jacobian trace-one lines. With the normalization of AM-R-006 and AM-LIFT-08, transport the already fixed product section through (PH18). Its actual reduced germ in product coordinates is \[
        \prod_{r=1}^{14}\Sigma_{\mathrm{tr}}(u^{(r)})
            =c^{14}\prod_{r=1}^{14}\sigma_q(u^{(r)}).
\tag{PH22}
\] Here \(\sigma_q=u_1u_3-u_2^2+\cdots\), and c is the already fixed algebraic scalar. In particular its order of vanishing at the origin is 28. No theta-null unit is asserted. This is the line at the invariant characteristic \(2\infty\), not the distinct non-effective translated characteristic.

Let \(\tau_0,\ldots,\tau_6\) be the seven real embeddings of B and put \[
                        U_{kj}=\tau_k(e_j).
\tag{PH23}
\] Equation (PH4) says \(U^tU=I_7\). Hence U is a real orthogonal matrix, \(U^{-1}=U^t\), and \(\det U=\pm1\). Write the seven A-copy coordinates as \((u_{i,j},v_{i,j})\), \(i=1,2,3\), \(j=0,\ldots,6\), and define \[
 w_{i,\pm,j}=u_{i,j}\pm i v_{i,j},\qquad
 z_{i,\pm,k}=\sum_{j=0}^6U_{kj}w_{i,\pm,j}.
\tag{PH24}
\] These are full K-eigenvectors for the actual induced CM action. The Hermitian metric is exactly \[
 H=\sum_{i,j}\frac{|u_{i,j}|^2+|v_{i,j}|^2}{G_i}
   =\sum_{i,\pm,k}\frac{|z_{i,\pm,k}|^2}{2G_i}.
\tag{PH25}
\] Thus the B-diagonalization itself has no archimedean operator-norm loss.

At every prime away from 29, all entries of U and \(U^{-1}\) are integral. At 11, (PH24) is consequently an integral invertible change of the primitive Hodge eigenbasis, including its old \(\pm i\) blocks, since 2 is a unit. This proves a precise integral comparison at 11, not merely an algebraic splitting of de Rham cohomology. AM-K018 then applies to this ordinary \(J^{14}\) product: the central \(\zeta_7\)-action has no invariant holomorphic quadratic form, and its full-CM Chern/reduced-gauge comparison is unchanged. No claim is made for the conjugate-factor products outside that lemma's scope.

At 29, every entry of U has normalized valuation \(-3/7\), by (PH9); its inverse has the same bound. A homogeneous derivative of degree D transported by this matrix therefore has coefficient valuations at least \(-3D/7\). This is an upper bound on the possible denominator, not a claim that every derivative attains it. The old \(\pm i\) coordinate change retains its dyadic denominator bound \(2^D\). The additional standard-lattice rescaling in (PH20), if used to choose logarithmic coordinates, must also be included explicitly; it is a diagonal multiplication by the displayed conjugates of \(\pi_B^{\pm3}\). All these diagonal factors are units at 11.

This section identifies the line, product section, metric, and local coordinate change. It does not assert that the old selected-Schur upper-bound witnesses retain their old individual CM weight after the seven-copy change. The common joint integral lower estimate can be transported through an integral invertible 11-coordinate change; a new selected upper witness needs its own weight calculation.

\subsection{5. Finite-level costs that change under the repair}\label{finite-level-costs-that-change-under-the-repair}

The relative CM extension K/F remains unramified at every finite place by AM-K029. At 29 it is split, because 29 splits completely in E₀. Nevertheless the absolute real field F has six primes above 29, each with residue field \(\mathbb F_{29}\) and absolute ramification index seven. Rational 29 is therefore \textbf{not} a uniformizer of these local real fields.

For example, for the standard split rank-a compact group at all six places, a principal congruence subgroup defined by the rational ideal \((M)\), with \(M=29^e\) and \(e\ge1\), has the exact index \[
     I_a^{\mathrm{new}}(M)
        =M^{42a^2}\prod_{i=1}^a(1-29^{-i})^6.
\tag{PH26}
\] Indeed the local residue length is \(7e\), so each local index is \(29^{7ea^2}\prod_{i=1}^a(1-29^{-i})\), by reduction to \(\operatorname{GL}_a(\mathbb F_{29})\) and the successive additive matrix kernels. Multiplying six places gives (PH26). This statement concerns that precisely defined principal subgroup; it is not an automatic identification of every source level symbol with rational M.

The old 29 half-orbit construction still has to be instantiated with the new actual lattice and its exact \(\alpha\bar\alpha=M\) data. The good-prime jet-surjectivity argument of AM-K022 applies at 11 once those data are specified, but the former absolute-uniformizer and degree-six mass formulas cannot simply be copied. Terminal \nolinkurl{qOf} depths must likewise retain their actual local valuations.

Combining this theorem with AM-K029 gives an actual product principal CM host and its literal auxiliary character. It does not settle the independent original \nolinkurl{IsTen}, full real/rank/size records, prescribed finite Xi equality, or the later selected-image arithmetic window. The degree-seven repair is one concrete part of that adapter, with its new finite costs exposed.

\subsection{6. Sources, novelty, and certificate scope}\label{sources-novelty-and-certificate-scope}

\begin{itemize}
\tightlist
\item
  Milne, J. S. (2020). \emph{Algebraic number theory} (Version 3.08, 19 July 2020). Proposition 6.2 and Theorem 6.4, printed pp.96--99, for prime cyclotomic ramification; Lemma 6.5, printed p.100, and Remark 6.6(c), printed p.101, for the integral compositum and its discriminant. Primary author PDF: \url{https://www.jmilne.org/math/CourseNotes/ANT.pdf}. The exact compositum statements were read in the source.
\item
  Erez, B. (1991). A survey of recent work on the square root of the inverse different. \emph{Astérisque, 198--200}, 133--152. Section 1, equations (1.1)--(1.2), printed p.134, for the tame different and trace-dual ideal; Theorem 2.9, printed p.139, records the general selfdual-normal-basis theorem in this tame abelian setting. Primary author paper: \url{https://smf.emath.fr/system/files/filepdf/AST_1991__198-199-200__133_0.pdf}. The general theorem is background only: (PH3)--(PH9) verify the particular basis directly.
\item
  Accepted DP57; AM-DO-022, §§1--2; AM-R-004/006 and AM-LIFT-08: the original Jacobian, product lattice and exact theta normalization. AM-K018: the full-CM gauge comparison for ordinary powers. AM-K029: the enlarged field, induced type and literal unramified-character construction.
\end{itemize}

The general square-root-inverse-different and selfdual-normal-basis facts are classical, not claimed as new. The concrete contributions here are the displayed basis and its exact local certificate for conductor 29, the integral product-polarization transport, the exact coefficient (PH20), and the finite coordinate/cost comparison for this host. The certificate checks rational/cyclotomic algebra, matrix identities, all coefficients in (PH8), and the displayed norm/index formulas. No Lean build, unbounded analytic derivative computation, or full source-host instance is claimed.

\clearpage

\section{AM-K032 --- an actual product principal host with degree divisible by eight}\label{am-k032-an-actual-product-principal-host-with-degree-divisible-by-eight}

Author: GPT-6 Astra Pro, ALP\_MACHINE /root/kummer\_extension, 22 November 2026 UTC; inherited reasoning effort not separately exposed.

The degree-seven construction of AM-K029/030 repairs the odd character but has real degree42. This finding supplies an explicit further real extension of degree4, together with an exact polarized product identification. The resulting host is the actual principally polarized variety \(J^{56}\), with full integral CM action, real degree168, and \(8\mid168\). Its odd character, all uniform even stage characters, and the permitted alpha half-orbit persist. The primitive-type common-field obstruction remains: the source's defined Fields.Kq is still \(\mathbb Q\). Consequently this is not a complete D7 constructor or a MainX estimate.

The new ingredient is elementary weighted trace linear algebra. No classification of positive unimodular lattices and no uncomputed class group is needed.

\subsection{1. Two exact rank-two weighted trace forms}\label{two-exact-rank-two-weighted-trace-forms}

For \(a\in\{1,3\}\), put \[
 d_a=a^2+4,\qquad R_a=\mathbb Q(\sqrt{d_a}),\qquad
 u_a=\frac{a+\sqrt{d_a}}2,
 \qquad \delta_a=\frac{u_a^{-1}}{\sqrt{d_a}}.
\tag{QH1}
\] Thus \((d_1,d_3)=(5,13)\). Each \(d_a\) is a squarefree discriminant congruent to1 modulo4. The ring of integers is \[
                 \mathcal O_{R_a}=\mathbb Z[u_a],
           \qquad u_a^2-au_a-1=0,
           \qquad u_a^{-1}=u_a-a.
\tag{QH2}
\] For a=3, the generator differs by1 from the usual \((1+\sqrt{13})/2\), so it gives the same integral ring. The discriminant of the displayed basis is \(d_a\).

\subsubsection{Proposition QH3}\label{proposition-qh3}

The element \(\delta_a\) is totally positive, generates the inverse different of \(R_a/\mathbb Q\), and satisfies \[
 \operatorname{Tr}_{R_a/\mathbb Q}(\delta_a)=1,\qquad
 \operatorname{Tr}_{R_a/\mathbb Q}(\delta_a u_a)=0,\qquad
 \operatorname{Tr}_{R_a/\mathbb Q}(\delta_a u_a^2)=1.
\tag{QH3}
\] Thus the weighted trace form \((x,y)\mapsto\operatorname{Tr}(\delta_a xy)\) has matrix \(I_2\) in the \textbf{integral} basis \((1,u_a)\). Multiplication by \(u_a\) in that basis is the symmetric integral matrix \[
                         M_a=\begin{pmatrix}0&1\\1&a\end{pmatrix}.
\tag{QH4}
\]

\textbf{Proof.} Write \(s=\sqrt{d_a}>a>0\) in the positive embedding. Then \[
 \delta_{a,+}=\frac{s-a}{2s}>0,\qquad
 \delta_{a,-}=\frac{s+a}{2s}>0.
\tag{QH5}
\] Their sum is1. Since \(\delta_a u_a=1/\sqrt{d_a}\), its trace is0. Since \(\delta_a u_a^2=u_a/\sqrt{d_a}\), its trace is1. This proves all three entries of the Gram matrix. In particular the integral lattice is selfdual for the weighted trace form, so \(\delta_a\mathcal O_{R_a}=\mathfrak D_{R_a/\mathbb Q}^{-1}\) by the definition of the inverse different. This also follows from the derivative \(2u_a-a=\sqrt{d_a}\) and the unit \(u_a^{-1}\), but the Gram computation already proves it directly. The multiplication matrix follows from \(u_a^2=1+au_a\). Finally \[
                    N_{R_a/\mathbb Q}(\delta_a)=1/d_a.
\tag{QH6}
\] All statements are exact. ∎

Let \(R=\mathbb Q(\sqrt5,\sqrt{13})\). Coprimality of the two discriminants implies \[
 \mathcal O_R=\mathcal O_{R_1}\otimes_{\mathbb Z}\mathcal O_{R_3},
 \qquad [R:\mathbb Q]=4,\qquad d_R=5^2 13^2.
\tag{QH7}
\] In the basis \[
                     1,\ u_1,\ u_3,\ u_1u_3,
\tag{QH8}
\] the form \(\operatorname{Tr}_{R/\mathbb Q}(\delta_1\delta_3xy)\) is \(I_4\). The two generating actions are commuting symmetric integral Kronecker matrices. These facts follow by multiplying the two trace identities, and the certificate checks the four-by-four matrices explicitly.

\subsection{2. Tensoring the actual CM lattice}\label{tensoring-the-actual-cm-lattice}

Retain the notation and the \textbf{actual} polarized product construction of AM-K030: \[
 \begin{split}
 E_0&=\mathbb Q(\zeta_{28}),\qquad
 B\subset\mathbb Q(\zeta_{29}),\quad[B:\mathbb Q]=7,\\
 K&=E_0B,\quad F=K^+,\quad [K:\mathbb Q]=84,\\
 \mathcal A_B&=\pi_B^{-3}\mathcal O_B,
 \qquad \pi_B=N_{\mathbb Q(\zeta_{29})/B}(1-\zeta_{29}),\\
 \Lambda&=\mathcal O_{E_0}\otimes\mathcal A_B
               =\pi_B^{-3}\mathcal O_K,\qquad
 \xi_0=\frac{\zeta_7-\zeta_7^{-1}}{14}.
 \end{split}
\tag{QH9}
\] The carried lattice \((\Lambda,\operatorname{Tr}_{K/\mathbb Q}(\xi_0x\bar y))\) is exactly the product principal CM lattice of \(J^{14}\), by the explicit trace-orthonormal basis of \(\mathcal A_B\). Its type is induced from \(\Phi_0=\{1,3,9,15,17,23\}\) and contains every real B-lift.

Set \[
                  K^*=KR,
              \qquad F^*=FR.
\tag{QH10}
\] The discriminants of K and R are coprime, and K is unramified at5 and13. Thus the fields are linearly disjoint and \[
 \begin{gathered}
 [K^*:\mathbb Q]=336,\qquad [F^*:\mathbb Q]=168,\\
 \mathcal O_{K^*}=\mathcal O_K\otimes\mathcal O_{R_1}
                               \otimes\mathcal O_{R_3}.
 \end{gathered}
\tag{QH11}
\] For completeness, any nontrivial subfield of the biquadratic R contains one of \(\mathbb Q(\sqrt5),\mathbb Q(\sqrt{13}),\mathbb Q(\sqrt{65})\), all ramified at5 or13; hence K cannot intersect R nontrivially. The integral tensor assertion is the coprime-discriminant lemma, applied twice.

Define \[
 \Lambda^*=\Lambda\otimes\mathcal O_{R_1}\otimes\mathcal O_{R_3}
                   =\pi_B^{-3}\mathcal O_{K^*},
 \qquad
 E^*(x,y)=\operatorname{Tr}_{K^*/\mathbb Q}
                       (\xi_0\delta_1\delta_3x\bar y).
\tag{QH12}
\] Use every extension of each old CM embedding to K*, so the new type has four lifts of each of the42 old embeddings.

\subsubsection{Theorem QH13}\label{theorem-qh13}

The CM lattice (QH12) is an actual product principal polarization: \[
                    (A^*,L^*)\simeq
                       (J,L_{2\infty})^{56}
\tag{QH13}
\] as principally polarized abelian varieties. It has a full action of \(\mathcal O_{K^*}\), with Rosati involution equal to complex conjugation. The CM extension \(K^*/F^*\) is unramified at every finite prime. On the standard lattice \(\mathcal O_{K^*}\), the coefficient of its alternating trace form is \[
                    \boxed{\xi^*=\xi_0\pi_B^{-6}\delta_1\delta_3.}
\tag{QH14}
\]

\textbf{Proof.} Decompose an element of \(\Lambda^*\) in the four-element real basis (QH8), with coefficients in \(\Lambda\). Taking the trace first over the two quadratic extensions and applying (QH3) kills the mixed terms and leaves the sum of four copies of the old alternating trace form. Consequently the lattice is the orthogonal direct sum of four old principal symplectic lattices. Since all four real lifts of each old holomorphic embedding are chosen, the corresponding complex structure is the product complex structure. This proves (QH13), not merely an isogeny statement. Positivity is also immediate either from this product identification or from the total positivity of both \(\delta_a\).

The two new integral actions are the matrices (QH4), tensored with identity matrices; they commute with each other and with the old \(\mathcal O_K\)-action. Their tensor ring is the full integer ring in (QH11). They are selfadjoint because their matrices are symmetric in the product principal lattice. The old Rosati involution is conjugation; together these facts give conjugation on the full K* action. The CM type is the asserted induced type.

The extension K/F is finite-unramified by AM-K029/030. Finite étale quadratic algebras remain étale under base change. Applying this to F\emph{/F proves the finite-unramified assertion. The quadratic base change is still a field: F} is totally real, whereas a generator of the quadratic CM extension K/F cannot lie in a totally real field.

Finally, write \(x=\pi_B^{-3}x'\) and \(y=\pi_B^{-3}y'\) in (QH12). The element \(\pi_B\) is real, so the coefficient becomes (QH14). This multiplication transports the entire lattice and form; it is not an assertion that the coordinate factor is a unit at29. ∎

As an independent norm/discriminant check, \[
 \begin{split}
 |d_{K^*}|&=|d_K|^4d_R^{84}
      =2^{336}7^{280}29^{288}5^{168}13^{168},\\
 N_{K^*/\mathbb Q}(\xi^*\mathcal O_{K^*})
      &=|d_{K^*}|^{-1}.
 \end{split}
\tag{QH15}
\] The latter is the norm of the inverse different ideal. Equivalently, the principal trace-dual lattice calculation already shows that \(\xi^*\mathcal O_{K^*}=\mathfrak D_{K^*/\mathbb Q}^{-1}\).

\subsection{3. Actual line, gauge, eigenbasis and finite coordinate costs}\label{actual-line-gauge-eigenbasis-and-finite-coordinate-costs}

The trace-one element \[
                 w_0=-(\zeta_7+\zeta_7^2+\zeta_7^4),
                      \qquad w_0+\bar w_0=1,
\tag{QH16}
\] acts diagonally on all56 ordinary J factors. Under the polarized product identification, its symmetric trace-one line is the external product of \(L_{2\infty}\). One way to fix the symmetric representative is to use its invariance under \(\zeta_7\): the endomorphism \(\zeta_7-1\) is invertible on the two-torsion, so the principal polarization has a unique symmetric representative fixed by this action. This is the same argument used for the actual source line in ALP-LIFT-07/08 and K030.

Transport the already fixed product section, rather than introducing a new unconstrained section. With the notation of AM-R-006 and K030 its reduced germ is \[
                    c^{56}\prod_{\ell=1}^{56}\sigma_q(u_\ell),
\tag{QH17}
\] and its origin multiplicity is112. Here \(\sigma_q\) is the actual q-normalized principal germ and c is the fixed derivative-Thomae scalar. The transported principal section is primitive at11 by the carried finite normalization; no theta-null unit at the origin is claimed, since the section vanishes there.

To compare the new CM coordinates, first retain K030's trace-orthonormal B coordinates z, whose Hermitian form has coefficient \(1/(2G_i)\) on each old E₀ label \((i,\pm)\) and B embedding. For one new quadratic extension, the two ordinary product coordinates \((z_0,z_1)\) become \[
 \begin{pmatrix}z_+\\z_-\end{pmatrix}
      =V_a\begin{pmatrix}z_0\\z_1\end{pmatrix},
 \quad
 V_a=\begin{pmatrix}1&u_{a,+}\\1&u_{a,-}\end{pmatrix},
 \quad \det V_a=-\sqrt{d_a}.
\tag{QH18}
\] This is the embedding matrix of (1,u\_a), and \[
     V_a^{\mathsf t}
           \operatorname{diag}(\delta_{a,+},\delta_{a,-})V_a=I_2.
\tag{QH19}
\] In particular it is a \textbf{weighted} identity. The unrescaled V\_a is not asserted to be orthogonal. After both extensions the Hermitian form is exactly \[
       H=\sum_{i,\pm,k,\epsilon,\epsilon'}
          \frac{\delta_{1,\epsilon}\delta_{3,\epsilon'}}{2G_i}
                      |z_{i,\pm,k,\epsilon,\epsilon'}|^2.
\tag{QH20}
\] The positive square-root rescaling by \(\sqrt{\delta_{1,\epsilon}\delta_{3,\epsilon'}}\) is an archimedean isometry to the product metric. Formula (QH20), rather than an untracked scalar bound, is the appropriate input for the selected Fock norm.

Both V\_a and V\_a\^{}\{-1\} are integral at11, since the u\_a are algebraic units and their determinants are11-units. Combined with K030's B comparison and the old \(\pm i\) comparison, the full K* eigenbasis is therefore an integral invertible change of the primitive ordinary J tangent basis at11. The central \(\zeta_7\) still has holomorphic labels1,2,3. Their pairwise sums never vanish modulo7, including repeated factors. Thus AM-K018's uniqueness proof for the full-CM canonical/source gauge applies to this ordinary product. This supplies the same gauge comparison for (QH17); no extension to a product with conjugated J factors is asserted.

The additional finite coordinate costs are explicit:

\begin{longtable}[]{@{}
  >{\raggedright\arraybackslash}p{(\columnwidth - 2\tabcolsep) * \real{0.5000}}
  >{\raggedright\arraybackslash}p{(\columnwidth - 2\tabcolsep) * \real{0.5000}}@{}}
\toprule\noalign{}
\begin{minipage}[b]{\linewidth}\raggedright
Prime
\end{minipage} & \begin{minipage}[b]{\linewidth}\raggedright
Additional effect from \(R/\mathbb Q\)
\end{minipage} \\
\midrule\noalign{}
\endhead
\bottomrule\noalign{}
\endlastfoot
11 & V\_a and its inverse are integral invertible; all \(\delta_a\) are units \\
29 & No new denominator from V\_a; the B and standard-lattice \(\pi_B^{\pm3}\) factors of K030/DO034 remain \\
5 & Entries of \(V_1^{-1}\) have valuation \(-1/2\), normalized by \(v_5(5)=1\); the other V is invertible integral \\
13 & Entries of \(V_3^{-1}\) have valuation \(-1/2\); the other V is invertible integral \\
Other primes & Both new embedding matrices and their inverses are integral invertible \\
\end{longtable}

Indeed \[
 V_a^{-1}=\frac1{u_{a,-}-u_{a,+}}
              \begin{pmatrix}u_{a,-}&-u_{a,+}\\-1&1\end{pmatrix}.
\tag{QH21}
\] At the discriminant prime every numerator entry is a unit and the denominator has valuation1/2. Thus transforming a homogeneous degree-D polynomial from new eigen-derivatives to the ordinary product derivatives has additional coefficient denominator at most D/2 at5 and at13. This is an upper bound for the coordinate change, not a claim that an actual selected theta value attains it. The old dyadic and29 costs are retained. The explicit metric (QH20) and these local changes must be included when extending DO034's selected polynomial calculation; they are not removed by the11-unit comparison.

\subsection{4. Exact ranks, conductor and explicit full-congruence G primes}\label{exact-ranks-conductor-and-explicit-full-congruence-g-primes}

The smallest real extension degree h making \(8\mid42h\) is4. The present construction realizes it. The exact source f\_le\_k therefore requires \[
       \boxed{f=168,\qquad \operatorname{ord}=84,\qquad k\ge168.}
\tag{QH22}
\] The abstract Frame can again use its parameter e=0 and distinguished element0. Its group may be \(C_{21}\times C_2\times C_2\), of order84. Its quotient to HGrp(0) is trivial and satisfies the literal Frame clauses. In fact \(\operatorname{Gal}(K^*/\mathbb Q)\simeq C_2^4\times C_{21}\), compatible with an abstract \(G_0\times A\) decomposition. This observation alone does not construct the source's full field/type record.

For m=2 and k=168, take \[
 \begin{gathered}
 t=1633,\quad r=169,\quad d=1634,\quad r+t=1802,\\
                    (N_0,N_1,N_2)=(168,1802,3436).
 \end{gathered}
\tag{QH23}
\] All source rank inequalities hold. Substituting these ranks in K031's same numerical parameter formula gives \[
 \begin{gathered}
 S=3435,\quad S^3=40\,530\,337\,875,\quad
 C=S^3+k+t+11=40\,530\,339\,687,\\
 u_i=(337-2i)C+\tfrac12,\quad
 v_i=(336-2i)C,\quad
 (b_0,b_1)=(-2C,2C),\quad(e_0,e_1)=(0,4C).
 \end{gathered}
\tag{QH24}
\] The identical proof verifies every numerical parity, interlace, centred and NCGaps clause of Hyp3 with \(\eta_{\rm num}=-1/2\). The certificate separately checks the actual new finite lists. The source selected recurrence yields \[
 \begin{split}
 \nu_{0,i}&=(337-2i)C+i-984\quad(0\le i<168),\\
 \min_i\nu_{0,i}&=3C-817>0,\qquad
 \Delta=4C-1634,\qquad
 \nu_1=(\nu_0+\Delta,\Delta^{1634}).
 \end{split}
\tag{QH25}
\] The actual new selected polynomial must use all168 CM blocks. The highest bideterminant construction in DO034 is defined separately at each CM place and therefore gives a polynomial in these168 blocks by the same formula; its diagonal coefficient remains one. DO034's saved first-extraction theorem is stated for42 places. Extending its global formula requires the present quadratic metric data and new local indices, rather than silently retaining its numerical exponents. Equations (QH20)--(QH21) supply those additional metric/coordinate data; the separate first-lift analysis retains responsibility for that extension.

The rational conductor is exactly \[
                         N^*=812\cdot5\cdot13=52\,780.
\tag{QH26}
\] Containment in \(\mathbb Q(\zeta_{52780})\) follows from the three cyclotomic containments. Conversely, containment in any cyclotomic field forces divisibility by812 by K031, and ramification of the two new quadratic subfields forces divisibility by5 and13. This is exactly the source's minimum-divisibility IsCond criterion.

The following explicit primes satisfy the full new congruence: \[
 \begin{array}{c|r|l}
 i&g_i^*&\pi_i^*\\\hline
 0&22\,285\,932\,761&52\,781+52\,780\sqrt{-7}\\
 1&64\,072\,175\,441&211\,121+52\,780\sqrt{-7}.
 \end{array}
\tag{QH27}
\] The finite certificate checks primality by direct trial division through the integer square roots149284 and253124, respectively. It also checks \[
 \pi_i^*\bar\pi_i^*=g_i^*,\qquad
 \pi_i^*\equiv1\pmod{52780\mathcal O_{\mathbb Q(\sqrt{-7})}},\qquad
 g_i^*\equiv1\pmod{52780},\qquad
                 g_i^*-1>2\cdot168^2=56\,448.
\tag{QH28}
\] The conjugate ideals are coprime because \(\gcd(g_i^*,2\operatorname{Re}\pi_i^*)=1\). The primes split completely in K*, and each chosen half contains168 degree-one primes. Thus the actual norm, full PiGCong, residue cardinality and G-size clauses are met in the same external quadratic subfield as K031. The literal piG\_Kq clause is still not met.

For faithful residue characters one may use primitive roots3 and6 at the two primes. The certificate verifies their order from the complete factorizations of \(g_i^*-1\). The common residue-root field can be chosen using \[
 L^*=\operatorname{lcm}(28,g_0^*-1,g_1^*-1)
                  =13\,526\,981\,751\,068\,240.
\tag{QH29}
\] This explicit coefficient extension is large; no claim that its degree or arithmetic height is negligible is made.

\subsection{5. Actual characters and the unchanged permitted family}\label{actual-characters-and-the-unchanged-permitted-family}

Let \(\eta\) be K029's actual odd character on K. Define \[
                    \eta^*=\eta\circ N_{K^*/K}.
\tag{QH30}
\] By the norm-compatibility proof in K029, this character is unitary, finite-unramified, has angular exponent1 at each new CM embedding, and restricts to \(\eta_{K^*/F^*}\) on \(\mathbb A_{F^*}^\times\). Its normalized etaVal values lie in one finite coefficient field: at a prime P of residue degree f over p, the normalized value is the f-th power of the old normalized value. A norm of a local uniformizer has valuation f; its unit factor is killed by the unramified character, so no choice of uniformizer affects this assertion.

For any integers e\_j choose the LIFT18 finite-order zero-infinity characters \(\nu_j\) of E₀ at the primes (QH27), and set \[
         \boxed{\mu_j^*=(\nu_j\circ N_{K^*/E_0})(\eta^*)^{-2e_j}.}
\tag{QH31}
\] The full proof is HC16--21 of K031 and applies verbatim: every new G completion is the same degree-one \(\mathbb Q_{g_i^*}\), so its local norm is the identity; the even eta power changes infinity weight while changing no compact unit character. Hence all exact uniform IsMu0 clauses hold, without a divisibility condition on e\_j and without a large-field class-number assumption. A single finite field contains all muVal and etaVal values. For j=0,1 the exact conductor is \[
     \mathfrak f(\mu_j^*)=g_0^*g_1^*\mathcal O_{K^*},
       \qquad N\mathfrak f(\mu_j^*)=(g_0^*g_1^*)^{336}.
\tag{QH32}
\] Qualitative T-prime choices use K031's ray-class construction with modulus \(52780(g_0^*g_1^*)^2\); no quantitative upper window is inferred.

The same central element \[
           \alpha=(1+2\sqrt{-7})^e,\qquad M=29^e,
                    \alpha\bar\alpha=M
\tag{QH33}
\] acts diagonally on all56 J factors and is congruent to1 modulo2. Its conjugate ideal remains coprime. Therefore on an edge of ranks a,b the permitted product DP59 basis is exactly the h=4 instance of K031: \[
 \begin{gathered}
 X=(A^*)^{ab}=J^{56ab},\quad \dim X=168ab,\\
 \mathcal W_M=[\alpha]^*H^0(X,\mathcal L^M),
       \quad\operatorname{rank}\mathcal W_M=M^{168ab},\\
 \mathcal L=(L^*)^{\boxtimes ab},\qquad
 t'\in X[\alpha^2]/X[\alpha],\qquad Mt'\in X[\alpha].
 \end{gathered}
\tag{QH34}
\] The ordinary D-jet threshold remains \(M\ge3D+2\), and each individually selected finite translate has norm1. The auxiliary multijet threshold in DO030 is unchanged under its stated coprimality assumptions. The selected highest-symbol factorial bound is K022's bound in the new integral eigenbasis. It does not transfer an old six-place sharp directional witness to one new full-CM selected weight.

At29 both real quadratics split, as witnessed by \[
                       11^2\equiv5\pmod{29},\qquad
                       10^2\equiv13\pmod{29}.
\tag{QH35}
\] Thus F* has24 primes over29, each with ramification index7 and residue field \(\mathbb F_{29}\), and all are split in K\emph{/F}. The rational principal-level index in a rank-a split unitary factor is consequently \[
        \boxed{I^*(M,a)=M^{168a^2}
                       \prod_{i=1}^a(1-29^{-i})^{24}.}
\tag{QH36}
\] This follows from \(|\operatorname{GL}_a(\mathcal O_v/29^e)|
 =29^{7ea^2}\prod_{i=1}^a(1-29^{-i})\) at each of24 places. It is the rational principal-level factor, not a substitution for every source level or the prescribed ramified G-column index. The latter uses the new primes,168 selected places at each G, and the separate LIFT19/DO034 formulas.

\subsection{6. The remaining constructor obstruction}\label{the-remaining-constructor-obstruction}

The source's eight\_dvd condition is now fulfilled by the actual host. The source's common-field requirement is not. The induced type includes every R- and B-lift of each E₀ embedding. If every selected embedding agrees on x, agreement across all lifts first forces \(x\in E_0\); the automorphisms3 and17 in \(\Phi_0\) then force \(x\in\mathbb Q\), exactly as in K031(HC28). Hence \[
                       \mathrm{Fields.Kq}=\mathbb Q.
\tag{QH37}
\] This restates the already identified primitive-type obstruction on the new host. Neither the larger field nor the odd character changes it. In particular the external \(\mathbb Q(\sqrt{-7})\) containing the explicit prime generators cannot be relabeled as the source-defined common quadratic field.

The completed pieces are therefore an actual principal CM product host with8-divisible real degree, a literal odd Hecke character, all uniform even character conditions, explicit rank/prime/congruence data, and the unchanged permitted half-orbit. The source's full field/type record, exact later selected image, prescribed raw-Xi chain and quantitative MainX window remain separate. This finding makes no conclusion about GRH or a Kummer extension of the accepted theorem.

\subsection{7. Primary and exact input locators}\label{primary-and-exact-input-locators}

\begin{itemize}
\tightlist
\item
  Milne, J. S. (2020). \emph{Algebraic number theory} (Version3.08), Example2.41, printed p.39, for quadratic integer rings and discriminants; Lemma6.5, printed p.100, and Remark6.6(c), printed p.101, for the coprime-discriminant integral tensor and discriminant formula. Primary author PDF: \url{https://www.jmilne.org/math/CourseNotes/ANT.pdf}. The weighted trace identity itself is proved in full in (QH3)--(QH6).
\item
  AM-K029/030: the actual finite-unramified degree-seven CM host, odd character, trace-orthonormal inverse-different basis and full principal product identification. AM-K031 HC1--28: source numerical records, character norm pullback, permitted isogeny and exact common-field scope.
\item
  \nolinkurl{Formalism/LocalRingCardinality/Params.lean:6–22;} \nolinkurl{Formalism/LocalToGlobal/Columns/AbelianHostFieldConditions.lean:59–76,108–125,206–238,310–331,486–494,560–579;} \nolinkurl{Formalism/LocalToGlobal/UnitaryNorms/IsUnitaryNormOf.lean:66–94:} the exact unmodified field, character, degree and rank clauses. Paths are relative to work/research/alpoge\_source.
\item
  AM-K018 and ALP-LIFT-07/08: the actual principal line, finite normalization and full-CM reduced-gauge uniqueness for ordinary J powers. AM-R-006: fixed scalar c.~Carried DP59, AM-K022 and AM-DO-030/034: actual allowed half-orbit, finite jets and new-block selected construction.
\end{itemize}

The exact certificate \nolinkurl{degree_four_principal_host_repair_certificate.py} and its JSON output check both weighted trace matrices, integral symmetric actions and tensor commutation, new degree/discriminant arithmetic, the actual new rank/gap/partition lists, direct primality and full prime congruences, primitive-root witnesses and the29 splitting/index identities. They are finite arithmetic checks. The polarized variety, character and canonical-gauge proofs are the mathematical arguments above and their explicitly stated carried inputs; no Lean build or analytic window check is claimed.

\clearpage

\section{ALP-LIFT-18 --- exact literal-G character compatibility and an explicit admissible family}\label{alp-lift-18-exact-literal-g-character-compatibility-and-an-explicit-admissible-family}

Author: /root/lift\_accounting, GPT-6 Astra Pro, ALP\_MACHINE, 21 November2026UTC.

This result solves the even stage-character compatibility problem for the actual mixed host, with the source's specified G-residue characters. It gives a necessary and sufficient congruence, constructs the characters when it holds, and supplies numerical Harish--Chandra data satisfying that congruence and the displayed source gap conditions. The next local type and prescribed-kernel nonvanishing problems are not solved by the character construction.

The norm-one torus class-number statement of AM-R-009 is used as a checked input. The exact residue characters and conductor conventions are the ones audited in ALP-LIFT-17. The independently obstructed unramified odd character etaT is not supplied here.

\subsection{1. Host, the actual G half, and the root norm}\label{host-the-actual-g-half-and-the-root-norm}

Put \[
 E'=\mathbf Q(\zeta_{28}),\qquad F'=E'^+,\qquad
 K_q=\mathbf Q(\sqrt{-7}),\qquad
 \Phi'=\{1,3,9,15,17,23\}.
 \tag{G1}
\] The actual primitive CM type on \(J^2\) is \(\Phi'\), with \(\sigma_a(\zeta_{28})=\zeta_{28}^a\). The G half in the source is defined instead by a prime generator \(\pi_G\in K_q\), with \(\pi_G\bar\pi_G=q\), where \(q\equiv1\pmod{28}\) is prime. Its selected places are the six places of \(E'\) above the chosen prime of \(K_q\), namely the support of \(\bar\pi_G\). They give one member of every conjugate pair over \(q\).

The subgroup \(\operatorname{Gal}(E'/K_q)\) consists of those units modulo28 whose reduction modulo7 is a quadratic residue: \[
                    \Psi'=\{1,9,11,15,23,25\}.
\] For an elementary check of this quadratic subfield, put \(u=\zeta_7+\zeta_7^2+\zeta_7^4\) and \(v=\zeta_7^3+\zeta_7^5+\zeta_7^6\). Then \(u+v=-1\) and direct multiplication gives \(uv=2\), so \((u-v)^2=-7\). A residue-class automorphism fixes \(u-v\) exactly when it preserves \(\{1,2,4\}\), which gives the displayed subgroup. Directly, \[
                \sum_{a\in\Psi'}a=84,\qquad
       N_{E'/K_q}(\zeta_{28})=\zeta_{28}^{84}=1.
 \tag{G2}
\] The other coset gives the same norm-one assertion. It is important that this norm uses the induced G half, rather than the primitive CM type \(\Phi'\), whose exponent sum is68.

At the selected degree-one G places, all residue fields are canonically \(\mathbf F_q\) through their prime fields. For any global root \(z\in\mu_{28}\), (G2) gives \[
               \prod_{P\mid\bar\pi_G}(z\bmod P)=1
                                      \quad\hbox{in }\mathbf F_q^\times.
 \tag{G3}
\] This is simply the reduction of its relative norm.

\subsection{2. Prescribed compact characters and the exact congruence}\label{prescribed-compact-characters-and-the-exact-congruence}

Let \(g_0,g_1\) be two such G primes, with the source's selected halves. For each choose a faithful residue character \(\omega_\epsilon:\mathbf F_{g_\epsilon}^\times
 \to\mathbf C^\times\). At stage \(i\), the literal source condition is \[
        \chi_{i,P}=\omega_\epsilon^{\,2(i+1)}
                 \qquad(P\mid\bar\pi_{G,\epsilon}),
 \tag{G4}
\] and the desired Hecke character \(\mu_i\) is required to be unitary, trivial on \(\mathbb A_{F'}^\times\), tame, unramified outside \(g_0g_1\), and to satisfy \[
         \mu_{i,P}|_{\mathcal O_{E',P}^\times}=\chi_{i,P}^{-1}.
 \tag{G5}
\] These are the IsMu0 conditions, before the independent odd-character and arithmetic-kernel pins.

Let \(T'=\ker(N_{E'/F'}:\operatorname{Res}_{E'/F'}\mathbf G_m
\to\mathbf G_m)\), and \(K_0=(\widehat{\mathcal O}_{E'})^1\). At a relative split finite place, projection to its selected E-prime identifies the compact norm-one group with the local unit group. Define a compact character \(\rho_i\) of \(K_0\) by \[
       \rho_i(t)=
          \prod_{\epsilon=0}^1
          \prod_{P\mid\bar\pi_{G,\epsilon}}
                            \chi_{i,P}(t_P\bmod P)^{-1},
 \tag{G6}
\] and take it to be one at every other finite place. Equations(G3)--(G4) imply \[
                         \rho_i|_{\mu_{28}}=1.
 \tag{G7}
\]

For any integer vector \(\mathbf e_i=(e_{i,a})_{a\in\Phi'}\), let \[
          \beta_i(t_\infty)=
                    \prod_{a\in\Phi'}\sigma_a(t_\infty)^{-e_{i,a}}.
 \tag{G8}
\] This is a unitary character of \(T'_\infty\simeq U(1)^6\). AM-R-009 proves \[
 T'(\mathbb A_{F',f})=T'(F')K_0,\qquad
                          T'(F')\cap K_0=\mu_{28}.
 \tag{G9}
\] Therefore \(\beta_i\rho_i\) descends to an automorphic character of \(T'(\mathbb A_{F'})\) if and only if its restriction to the diagonal roots is one. By(G7), the exact condition is \[
 \boxed{\sum_{a\in\Phi'}a e_{i,a}\equiv0\pmod{28}.}
 \tag{G10}
\] For a uniform vector \(e_{i,a}=e_i\), this is equivalent to \[
                                7\mid e_i,
 \tag{G11}
\] since \(\sum\Phi'=68\) and \(\gcd(68,28)=4\).

There is no further norm-one global-unit condition hidden here: the full intersection in(G9) is precisely \(\mu_{28}\). This uses the actual class-number-one torus decomposition, rather than testing only one necessary root relation in an arbitrary CM field.

\subsection{3. Explicit construction and necessity for the actual Hecke character}\label{explicit-construction-and-necessity-for-the-actual-hecke-character}

Assume(G10). Given a finite norm-one idele \(t_f\), write \(t_f=\delta k\), with \(\delta\in T'(F')\), \(k\in K_0\), using(G9), and set \[
 \gamma_i(t_\infty,t_f)
       =\beta_i(\delta_\infty^{-1}t_\infty)\rho_i(k).
 \tag{G12}
\] If another decomposition is used, the two rational factors differ by an element of \(T'(F')\cap K_0=\mu_{28}\). Equations(G7),(G10) make(G12) independent of that choice. The formula is multiplicative, continuous, unitary, and trivial on the diagonally embedded \(T'(F')\). It is therefore an actual automorphic character.

Put \(j(x)=x/\bar x\) on ideles and define \[
                              \mu_i(x)=\gamma_i(j(x)).
 \tag{G13}
\] It is trivial on the base-field ideles. At infinity its restriction is \[
           \prod_{a\in\Phi'}
               \left(\frac{\sigma_a(x)}{|\sigma_a(x)|}\right)^{-2e_{i,a}},
 \tag{G14}
\] which is exactly HasInfQ with exponent vector \(-\mathbf e_i\): the source definition sets the integral angular exponent to twice the HasInfQ entry. At a selected G-prime, an idele supported in that E-component maps under \(j\) to the norm-one pair \((x,x^{-1})\). Equations(G6),(G13) therefore give precisely(G5), with its inverse sign. At other split places it is unramified.

At a nonsplit finite place the extension \(E'/F'\) is unramified, and every norm-one element is a unit. Moreover every norm-one unit is \(a/\bar a\) for a local unit \(a\): start with local Hilbert90, then multiply its witness by a base-field uniformizer power to make its valuation zero, using relative ramification index one. Since \(\rho_i=1\) there, \(\gamma_i\) is the trivial local torus character, and \(\mu_i\) is trivial on the entire local multiplicative group. All required outside-G unramified conditions follow. The G characters factor through residue units, so tameness is exact. Their conductor is the one in LIFT17(F19).

Conversely, suppose a genuine \(\mu_i\) satisfies these exact finite restrictions and HasInfQ\((-\mathbf e_i)\). The idele map \(j:\mathbb A_{E'}^\times\to T'(\mathbb A_{F'})\) is surjective by local Hilbert90, with kernel \(\mathbb A_{F'}^\times\); the global rational version is also surjective. Hence \(\mu_i\) descends to an automorphic \(\gamma_i\) on \(T'\). At each relative split finite place its restriction to \(K_0\) is forced by the unit character. At a nonsplit place the preceding integral Hilbert90 argument shows that an unramified \(\mu_i\) descends to a character trivial on the whole compact torus. Thus \(\gamma_i|_{K_0}=\rho_i\), and its infinity character is(G8). Rational triviality on \(\mu_{28}\) forces(G10).

This proves necessity as well as sufficiency. The conclusion is about the exact IsMu0 finite and infinity conditions, with its uniformizer values defined by the constructed character. It does not allow independently prescribing incompatible muVal entries.

\subsection{4. Concrete primes and a fixed coefficient field}\label{concrete-primes-and-a-fixed-coefficient-field}

For the actual minimal numerical choice \(k=6,m=2\), take \[
 \begin{array}{c|c|c}
 q&\pi_G&\pi_G\bar\pi_G\\ \hline
 113&1+4\sqrt{-7}&113\\
 281&13+4\sqrt{-7}&281 .
 \end{array}
 \tag{G15}
\] Both rational primes are \(1\bmod28\), so they split completely in \(E'\). The elements are algebraic integers in \(K_q\). Their conjugate ideals are coprime: a common prime would divide both their norm \(q\) and their sum \(2\) or26, which is impossible. Consequently they define the required six-element selected halves, with residue fields \(\mathbf F_q\). Their lower size bound is \[
            g_\epsilon-1>2\cdot6\cdot6=72.
\] This verifies the displayed G-prime conditions used here. No additional primitive-field D7 or etaT condition is inferred.

The faithful residue characters can take values in \[
                              K=\mathbf Q(\zeta_{560}),
 \tag{G16}
\] since \(560=\operatorname{lcm}(28,112,280)\). This is a fixed degree192 cyclotomic field, independent of all the stage weights. Formula(G12) also has algebraic finite values in this field: the rational norm-one factor \(\delta\) belongs to \(E'\), all its selected complex conjugates are Galois conjugates in \(E'\), and its absolute values are one. Thus the integer-power product \(\beta_i(\delta_\infty)\) belongs to \(E'\); the remaining factor is a root of unity in the displayed cyclotomic field.

This identifies actual finite character values, including values on uniformizers, by a formula. It does not assert that all those algebraic values are units at11. On compact units their values are roots of unity and hence units at every finite place. Values on noncompact determinant representatives require the separate normalization in the class presentation.

For these two stages the exact conductor is \[
       \mathfrak f(\mu_i)=31753\,\mathcal O_{E'},\qquad
                   N_{E'/\mathbf Q}\mathfrak f(\mu_i)=31753^{12}.
 \tag{G17}
\] This conductor is independent of the size of \(\mathbf e_i\). Its norm alone is not a complete H or a theta-row error term.

\subsection{5. Nonempty source gap locus with the required character congruence}\label{nonempty-source-gap-locus-with-the-required-character-congruence}

The compatibility condition does not force the numerical Harish--Chandra gap locus to be empty. Here is a fully explicit uniform example satisfying the numerical clauses of Hyp3. Keep \[
       k=6,\quad m=2,\quad r=7,\quad t=1795,\quad d=t+1=1796.
 \tag{G18}
\] The source rank inequalities hold: \(k,m\) are even, \(t\) is odd, \(m+3\le k\), \(r+5\le t\), \(3m<t\), and \(1801\le r+t=1802\). With the actual degree-six CM type count, take its corresponding order parameter3, so \(m\le6\le k\). The remaining displayed quadratic rank bound is immediate.

Choose any positive multiple \(C\) of7 satisfying \[
                  C>(r+mt)^3+k+t+10.
 \tag{G19}
\] At every one of the six embeddings define, with zero-based indices, \[
 \begin{aligned}
 u_i&=(2(k-i)+1)C+\tfrac12
                                &&(0\le i<k),\\
 v_i&=2(k-i)C                    &&(0\le i<k-1),\\
 b_j&=(4j-2)C,\qquad e_j=4jC     &&(0\le j<m),\\
 \eta&=-\tfrac12 .
 \end{aligned}
 \tag{G20}
\] The \(u_i\) are half-integral, the other parameters are integral, and \[
 u_i>v_i>u_{i+1},\qquad
 u_{k-1}=3C+\tfrac12>-b_0=2C,\qquad
 b_0<e_0<b_1<e_1.
\] Every required centred distance is at least \(\min(2C,4C-(t-1)/2)\), except for the positive \(u+e\) and \(v+e\) distances, which are still larger. Thus each exceeds \((t+3)/2\).

For NCGaps, the only false-sign element at stage0 is \(-b_1=-2C\). The other \(b\)-element is \(2C\), and all \(v_i\) are at least \(4C\), so every opposite-sign difference is at least \(4C>(r+mt)^3\). At stages1 and2 every element in the literal Dpar list has true sign, so the cross-sign condition is empty. The common \(\eta\)-shift at odd stages does not alter any difference. These statements directly check the source definition in DataTheta.lean538--555 and \nolinkurl{CayleySignature.lean106–118.}

The splitting weights are \(e_0=0,e_1=4C\), both divisible by7. Hence(G10) holds at both stages, and sections2--3 supply actual global even characters with exactly the prescribed G residue restrictions and no29 ramification. Their infinity exponents are the required \(-e_j\). The selected first partition is nonnegative: its smallest entry is \(3C-898>0\). The degree and other selected real data are therefore genuine finite quantities, although very large.

For a concrete integer, one may take \[
 C=46539458987,\qquad
 (e_0,e_1)=(0,186157835948),\qquad
 (b_0,b_1)=(-93078917974,93078917974).
 \tag{G21}
\] Here \((r+mt)^3=46539457173\), and(G19) holds. The symbols \(e_j\) in(G20)--(G21) are the Harish--Chandra splitting weights, not the DP44 error array.

The numerical \(\eta=-1/2\) matches an odd splitting character with infinity exponent \(+1/2\). It does not supply the original all-finite-unramified etaT pin, which AM-K028 separately obstructs on this host. Thus this paragraph verifies the stated numerical gap and character conditions, not a complete D7 instance.

\subsection{6. What is now fixed and what remains}\label{what-is-now-fixed-and-what-remains}

The even global characters are now explicit on a nonempty actual parameter locus, with the literal G conductor, exact residue characters, fixed coefficient field, and exact unit compatibility. The conductor29 character choice of LIFT14 was convenient for its spherical chain, but it is not needed for existence of these different stage characters.

At a split G place, changing to these characters changes the local theta parameters: the inserted rank-d block carries the stage character with its nontrivial residue restriction. This is exactly the kind of local change that the obstruction in LIFT17 requires. It is not obtained by choosing another vector in the old spherical constituent or by one uniform tame determinant twist.

To proceed to a chain with the prescribed column types, one must identify the ramified local constituents and show that the specified generic/determinant-block type and raw Xi map survive on the previous image. Their levels, projection masses and integral content remain to be calculated. Neither the existence of the characters nor their fixed finite conductor proves that local type occurrence, global prescribed-kernel nonvanishing, an integral gain, or a shorter DP44 chain.

\subsection{Source and input locators}\label{source-and-input-locators}

\begin{itemize}
\tightlist
\item
  AM-R-009, first\_lift\_unramified\_host.md, equationR9.15 and its norm-one ideal proof: actual class-number-one torus decomposition and the complete root-unit intersection. Masley(1976) remains the checked class-number input.
\item
  ALP-LIFT-17, sections4--7: exact source residue characters, G-prime conditions, size bounds, tameness and conductor.
\item
  \nolinkurl{Formalism/LocalToGlobal/UnitaryNorms/IsUnitaryNormOf.lean,} lines66--94: LevelLe, UnrAt, TameAt, HasInf and HasInfQ.
\item
  \nolinkurl{Formalism/LocalToGlobal/Columns/AbelianHostFieldConditions.lean,} lines205--224,273--306,315--323 and488: G halves, chi\_om, IsMu0 and g\_big.
\item
  \nolinkurl{Formalism/LocalRingCardinality/DataTheta.lean,} lines501--555, and \nolinkurl{Formalism/SharpRestriction/CuspSpanNormBounds/CayleySignature.lean,} lines106--125: the exact numerical parameter and gap conditions.
\item
  AM-K028: the separate odd-character obstruction is carried, not used as a substitute for the even-character construction.
\end{itemize}

All recovered source paths are relative to work/research/alpoge\_source, pinned to commit 7cbf7a0c32fca1badf05931879a37c3f8a2b03b4. No Lean build or complete source constructor is claimed.

\clearpage

\section{AM-DO-032: the actual terminal Y-family, exact Fourier Gram, and its projection denominator}\label{am-do-032-the-actual-terminal-y-family-exact-fourier-gram-and-its-projection-denominator}

Author: GPT-6 Astra Pro, ALP\_MACHINE /root/definite\_oscillator; inherited reasoning effort is not separately exposed; 21 November 2026 UTC.

This is a further calculation on the literal xiLast of AM-DO-031. It varies only the Y-parameter already present in the recovered source. It identifies its finite Fourier family, the exact Gram matrix, and the denominator of its orthogonal projector. Its arithmetic denominator is for this specified finite family; it does not identify the global MainX terminal detector lattice.

\subsection{1. The literal finite parameter set}\label{the-literal-finite-parameter-set}

Retain the notation and all source hypotheses of AM-DO-031: \[
 a=N(m-1),\quad b=N(m),\quad r=k+m-1,\quad
 \mathsf M=\mathrm{cst.M},\quad
 q_P=|\mathcal O_E/P|,\quad d_P=v_P(q_P).
 \tag{DO32.1}
\] In particular the T-places and rowSG are disjoint, and G denotes the G branch while Q denotes its complement inside rowSG. For the actual standard TwDat base points, the local frame lattices are standard.

The source definitions are:

\begin{longtable}[]{@{}
  >{\raggedright\arraybackslash}p{(\columnwidth - 2\tabcolsep) * \real{0.5000}}
  >{\raggedright\arraybackslash}p{(\columnwidth - 2\tabcolsep) * \real{0.5000}}@{}}
\toprule\noalign{}
\begin{minipage}[b]{\linewidth}\raggedright
Source
\end{minipage} & \begin{minipage}[b]{\linewidth}\raggedright
Locator
\end{minipage} \\
\midrule\noalign{}
\endhead
\bottomrule\noalign{}
\endlastfoot
\nolinkurl{Formalism/LocalToGlobal/Columns/FiniteSumNonzeroLargeDepth.lean} & ybSet and repY,1123--1131; ftY,1150--1156 \\
\nolinkurl{Development/AbelianVarietyTheory/Definitions/AutomorphicMembership/XiAtPenultimateEqXiCpl1.lean} & actual terminal point and representative,124--130; xiOver/xiLast,146--157 \\
\nolinkurl{Development/AbelianVarietyTheory/Definitions/GlobalTwisting/IntegralGLSuccEmbeddingMemParahoricTimesPrincipa.lean} & locTop,391--393; representative bound,531--547 \\
\end{longtable}

For fixed source decoration, ybSet is the set of cosets of the product-local period module in the product-local support module. Therefore it is canonically the finite set \[
 \mathcal Y=
 \prod_{P\in G}(q_P^{-1}\mathcal O_P^a/\mathcal O_P^a)
 \ \times
 \prod_{P\in Q}
     (q_P^{-\mathsf M}\mathcal O_P^a/
                      q_P^{\mathsf M}\mathcal O_P^a).
 \tag{DO32.2}
\] This statement concerns cosets, not the particular epsilon representative chosen by repY. To prove it, send a rational coset to its local residue classes. Injectivity is exactly the definition of its period module as the intersection of the local conditions. Surjectivity follows by CRT after clearing the finitely many denominators, as in AM-DO-031 Section5. The prescribed frames identify the localized source lattice with each standard local lattice. Thus every tuple in (DO32.2) is an actual allowed Y.

Changing repY within its coset changes its local component by \(\mathcal O_P^a\) at G and by \(q_P^{\mathsf M}\mathcal O_P^a\) at Q. On the supports of ftY, the resulting character argument is integral. LocCst.Ok.psi\_int therefore proves that the local finite function is unchanged. No representative phase is suppressed.

\subsection{2. The exact local Fourier calculation}\label{the-exact-local-fourier-calculation}

At a G-place, \(d_P=1\), and the source limited-depth character induces a nontrivial additive character on \(q^{-1}\mathcal O_P/\mathcal O_P\). Since q is the rational prime and the residue degree is1, its pairing with \(\mathcal O_P/q\mathcal O_P\) is perfect. The last-row functions are exactly \[
 \varphi_y(\eta)=1_{\mathcal O_P^a}(\eta)
                       \psi_P(-\eta y),
 \qquad y\in q^{-1}\mathcal O_P^a/\mathcal O_P^a.
 \tag{DO32.3}
\] With additive measure giving \(\mathcal O_P^a\) volume1, they are an orthonormal basis of the functions supported there and constant modulo \(q\mathcal O_P^a\): \[
 \langle\varphi_y,\varphi_{y'}\rangle
          =1_{\{y=y'\}}.
 \tag{DO32.4}
\] Indeed the integral of a nontrivial character on a compact additive group is0: translation by an element at which it is not1 multiplies that integral by a nontrivial scalar. For a trivial character the integral is the group's volume1.

At a Q-place the source character is only required on \(q^{-D_P}\mathcal O_P\), where \(D_P\ge2\mathsf M\). Its annihilator of the local ring within this domain is an \(\mathcal O_P\)-module: \[
 \mathfrak a_{\psi,P}
 =\{z\in q^{-D_P}\mathcal O_P:
                  \psi_P(z\mathcal O_P)=1\}
 =P^{c_P},
 \qquad -d_P<c_P\le0.
 \tag{DO32.5}
\] Here and below P denotes its maximal ideal in the completion. To justify every part of (DO32.5), the set is closed under addition and multiplication by \(\mathcal O_P\), contains \(\mathcal O_P\), and lies in a fixed fractional ideal. Hence it is a fractional ideal \(P^{c_P}\). The source condition \(\psi_P(q^{-1})\ne1\) excludes \(q^{-1}\) from it, so \(c_P>-d_P\), while its inclusion of the ring gives \(c_P\le0\). The same conductor ideal is therefore already determined on the smaller argument range needed for ftY. This reasoning uses only the finite additive quotient and does not extend the source character arbitrarily.

The actual Q functions, including the source scalar, are \[
 \varphi_y(\eta)=q^{-\mathsf M a}
 1_{q^{-\mathsf M}\mathcal O_P^a}(\eta)
                      \psi_P(-\eta y),
 \qquad
 y\in H_P:=
 q^{-\mathsf M}\mathcal O_P^a/
                       q^{\mathsf M}\mathcal O_P^a.
 \tag{DO32.6}
\] The pairing radical is \[
 R_P=
 q^{\mathsf M}P^{c_P}\mathcal O_P^a/
                       q^{\mathsf M}\mathcal O_P^a,
 \qquad |R_P|=q^{-c_Pa}.
 \tag{DO32.7}
\] The notational product \(P^{c_P}\mathcal O_P^a\) means componentwise multiplication by that fractional ideal. The bounds in (DO32.5) imply that its numerator is inside the support module. Two allowed y's give exactly the same function if their difference is in \(R_P\). Otherwise their inner product is0. More precisely, \[
 \boxed{\quad
 \langle\varphi_y,\varphi_{y'}\rangle
 =q^{(d_P-2)\mathsf M a}
                1_{\{y-y'\in R_P\}}.
 \quad}
 \tag{DO32.8}
\]

To prove (DO32.8), the character integral is nonzero precisely when \((y-y')q^{-\mathsf M}\mathcal O_P\subset P^{c_P}\), coordinate by coordinate. This is exactly (DO32.7). Its nonzero value before the amplitude is the volume \(q^{d_P\mathsf M a}\); the squared amplitude is \(q^{-2\mathsf M a}\). Their product gives the displayed factor. If the character is trivial on the full support, the two functions agree pointwise, not just up to a scalar, since their value at \(\eta=0\) is the same nonzero number.

After removing the duplicate parameters, the number of Q functions is \[
 |H_P/R_P|=q^{(2d_P\mathsf M+c_P)a}.
 \tag{DO32.9}
\] They span exactly the functions supported on \(q^{-\mathsf M}\mathcal O_P^a\) and constant modulo \(q^{\mathsf M}P^{c_P}\mathcal O_P^a\). This follows from finite Fourier inversion for the resulting perfect quotient pairing.

In the degree-one case \(d_P=1\), (DO32.5) forces \(c_P=0\). There are no duplicates, the common norm square is \(q^{-\mathsf M a}\), and the matrix of values \[
 q^{-\mathsf M a}\psi_P(-\eta y),
 \qquad \eta,y\in H_P,
 \tag{DO32.10}
\] is unitary for counting measure: \(|H_P|=q^{2\mathsf M a}\). This explains the literal scalar in ftY without replacing it by a guessed Haar normalization.

\subsection{3. Gram matrix of the prescribed terminal family}\label{gram-matrix-of-the-prescribed-terminal-family}

Let \(f_Y\) be the actual finite Heisenberg realization of xiLast for Y, using AM-DO-031. All upper-matrix and T-incidence factors are fixed as Y varies. The local positions at distinct selected places are independent. Consequently the preceding last-row inner products tensor with the fixed nonzero upper and T factors.

Let \[
 \overline{\mathcal Y}=
 \prod_{P\in G}
        (q_P^{-1}\mathcal O_P^a/\mathcal O_P^a)
 \times\prod_{P\in Q} H_P/R_P,
 \qquad n_Y=|\overline{\mathcal Y}|.
 \tag{DO32.11}
\] Choose one representative for each reduced class. The remaining choices give identical finite vectors. Then \[
 \boxed{\quad
 \langle f_Y,f_{Y'}\rangle_f
                =\kappa\,1_{\{Y=Y'\}},
 \qquad
 \det\operatorname{Gram}(f_Y)=\kappa^{n_Y},
 \quad}
 \tag{DO32.12}
\] where the nonzero rational number \(\kappa\) is exactly the unscaled finite norm square computed in AM-DO-031: \[
\begin{split}
 \kappa={}&
 \prod_{P\in T}q_P^{-a}
 \prod_{P\in G}
    \left[
     q_P^{-C_G}
       \prod_{i=1}^{t}(1-q_P^{-i})^{m-1}
    \right]\\
 &{}\cdot
 \prod_{P\in Q}
    \left[
     q_P^{(d_P-2)\mathsf M a-d_P r(b-1-r)}
       \prod_{i=1}^{r}(1-q_P^{-i})
    \right],\\
 C_G={}&r(r+1)/2+k+mrt+m(m-1)t^2/2.
\end{split}
\tag{DO32.13}
\] The scalar is unchanged by the actual Y-representative. A prescribed common chart scalar w multiplies the Gram by \(w\overline w\); it must be retained. Nonstandard lattice frames retain their actual Jacobian as in AM-DO-031.

Thus this is a source-defined orthogonal family, not a new arbitrary choice of finite inputs. It is a basis of its actual finite span. The identity (DO32.12) is a finite oscillator pairing, not an automorphic Petersson pairing.

\subsection{4. Integral saturation versus the exact pairing denominator}\label{integral-saturation-versus-the-exact-pairing-denominator}

Assume all selected residue characteristics differ from11, and work over a coefficient valuation ring containing the finite character values. Conjugate coefficient lattices are paired when the chosen 11-place is not fixed by the coefficient involution. Every root of unity remains an11-adic unit.

The lattice generated by the distinct \(f_Y\) is saturated in the finite function module on its ambient support/period window, intersected with its characteristic-zero span. Here is a direct argument. Finite Fourier inversion uses only orders of finite p-groups with \(p\ne11\); its denominators are units. The invariant submodule under the radicals in (DO32.7) is split by the corresponding averaging idempotent, whose denominator is again a unit. Finally the fixed upper/T factor has a unit-valued coordinate: use the local identity and incidence matrices of AM-DO-031 Section5. Evaluation at that coordinate splits tensor multiplication by that fixed factor, after multiplying by its unit inverse. This proves saturation.

Nevertheless the pairing on this saturated lattice need not be unimodular. From (DO32.13), \[
 \boxed{\quad
 v_{11}(\kappa)
 =\sum_{P\in G}(m-1)
                  \sum_{i=1}^{t}v_{11}(q_P^i-1)
   +\sum_{P\in Q}\sum_{i=1}^{r}
                         v_{11}(q_P^i-1).
 \quad}
 \tag{DO32.14}
\] All powers of q are units at11, so this formula holds for the general-Q calculation as well. For \(o_P=\operatorname{ord}_{11}(q_P)\mid10\) and \(h_L=\lfloor L/o_P\rfloor\), the exact elementary lifting-the-exponent formula gives \[
 \sum_{i=1}^{L}v_{11}(q_P^i-1)
   =h_L\,v_{11}(q_P^{o_P}-1)+v_{11}(h_L!).
 \tag{DO32.15}
\] Indeed the terms vanish unless \(o_P\mid i\); for \(i=o_Pj\) their valuation is \(v_{11}(q_P^{o_P}-1)+v_{11}(j)\).

In the actual rank range, \(t\ge13\). Thus every G-place away from11 contributes at least \(m-1>0\) to (DO32.14), because \(o_P\le10\). This is compatible with the coefficient primitivity of every \(f_Y\): coefficient primitivity and a unimodular restricted pairing are different conditions.

Apply the already proved detector calculation AM-R-002 to this specified finite lattice. For an individual basis vector \(f_Y\), its pairing ideal against this same lattice is exactly \(\kappa R\). Its dual vector within this span is \(\kappa^{-1}f_Y\), with exact denominator valuation \(v_{11}(\kappa)\) and complex norm \(\kappa^{-1/2}\). With the Hermitian pairing linear in its first argument, the orthogonal projector onto the whole span is \[
 \Pi_{\mathrm{term}}
 =\kappa^{-1}\sum_{Y\in\overline{\mathcal Y}}
                  f_Y\langle\,\cdot\,,f_Y\rangle_f.
 \tag{DO32.16}
\] It has that exact denominator on the ambient finite function lattice, with its volume normalizations retained. To see that the upper bound is attained, factor the space into the fixed upper/T vector and its full last-row Fourier subspace. A unit-valued coordinate of the fixed factor gives a unit numerator in its rank-one projector. The last-row projection has unit denominators at11. Its diagonal at an allowed last-row position is a p-power unit. Hence no smaller11-power can clear the displayed projector.

This is a restriction on a detector required to remain in this terminal finite span or on its orthogonal projector. It is not a denominator lower bound for every detector in the entire finite oscillator. Since the span is saturated, an integral functional on it can extend to the ambient coefficient module; that extended functional need not be represented by an integral vector in this same orthogonal subspace. Nor is the global MainX detector lattice already identified with this finite span.

\subsection{5. Exact effect of applying a theta operator}\label{exact-effect-of-applying-a-theta-operator}

Let T denote any fixed actual selected theta operator out of this finite input space, with the other archimedean and automorphic inputs fixed and their actual measures retained. Use the source Hermitian convention, linear in the first argument. Equations (DO32.12) and (DO32.16) give the genuine adjoint identity \[
 T\Pi_{\mathrm{term}}T^\dagger
 =\kappa^{-1}
    \sum_{Y\in\overline{\mathcal Y}}
          (Tf_Y)\langle\,\cdot\,,Tf_Y\rangle.
 \tag{DO32.17}
\] Restrict the target to the finite-dimensional image of the finite input window, so every linear scalar detector has a Riesz adjoint. For any such detector \(\ell\), this says \[
 \sum_Y|\ell(Tf_Y)|^2
 =\kappa\,
   \|\Pi_{\mathrm{term}}T^\dagger\ell^\dagger\|^2.
 \tag{DO32.18}
\] This is exact positivity and finite Plancherel, including its scalar. It proves nonvanishing of one actual allowed Y precisely when the displayed projected adjoint is nonzero.

The new calculation fixes the projection, its rank, Gram and arithmetic denominator. It does not assert that a prescribed prior theta image has nonzero projection to this finite span. A zero projected adjoint is still possible; the source's local types and arithmetic row must decide it. Likewise the actual terminal mass equation uses class stabilizers, branching maps and the specified row pairing. Substituting the finite Gram scalar for that global identity would be unjustified.

The calculation therefore removes a concrete finite Fourier and normalization ambiguity, and supplies a computable factor for a future matching proof. It does not replace the remaining selected-image or terminal-mass equation by the mere fact that xiLast is a nonzero function.

\subsection{Review status}\label{review-status-1}

Author whole-proof review: PASS, 2026-11-21T23:28:15.931357+00:00. The complete saved proof was read; source-linear-first Hermitian rank-one operator notation was corrected during that review. Exact quotient/radical, Gram, saturation, projector denominator and theta-adjoint scope checked. Independent review remains separate; no Lean build claimed.

\clearpage

\section{AM-DO-033: an actual finite Steinberg first input, with exact extraction and costs}\label{am-do-033-an-actual-finite-steinberg-first-input-with-exact-extraction-and-costs}

Author: GPT-6 Astra Pro, ALP\_MACHINE /root/definite\_oscillator, 21 November2026UTC; inherited reasoning effort not separately exposed.

This is a focused extension of AM-DO-030 and AM-R-012. It supplies a nonzero selected first lift whose source input has the actual initial generic-unipotent type at the auxiliary G places. It uses a specific finite Steinberg vector, one actual permitted29 half-translate, and a compact matrix coefficient at each G place. It computes the finite Schwartz norm and the source mass exactly. The new finite vector is not identified with the prescribed raw PhiG array. No later restricted-image assertion is included.

The global construction uses the already checked common host E'=Q(zeta28), A=J², its trace-one line, and the selected even edge U(a)→U(b), with a=k and b=k+t+1. A valid common splitting character with the required infinity type is assumed fixed. It is unramified off G and29, tame at G, and has at most the residue-level compact conductor at29 used in R012. It may have the literal tame G restrictions constructed in LIFT18, in which case it is unramified at29. Replacing the splitting while keeping its infinity exponents does not change the identity kernel value or the Heisenberg reading; its finite compact action is explicitly accounted for below. The local argument works over every finite residue field; the global extension to a different host additionally requires that host's actual mixed-theta/DO030 comparison, not just an abstract extension of its number field.

\subsection{1. A primitive generic vector with a known metric}\label{a-primitive-generic-vector-with-a-known-metric}

Fix one odd relative split auxiliary place with residue field κ=F\_q, q prime to11·29. Put

\[
 \bar G=\mathrm{GL}_a(\kappa),\quad B=TU,\quad
 n_U=a(a-1)/2,\quad d_\tau=q^{n_U}.
 \tag{DO33.1}
\]

Choose the precise nondegenerate character ψ\_U of U required by the initial column convention. For the source this is the character λ\_X,0; it is nontrivial on every simple-root subgroup. Let \(\mathcal F=\mathbf C[\bar G/B]\), with its delta basis orthonormal. The standard Steinberg submodule τ is irreducible over C, has dimension d\_tau, and has a one-dimensional B-fixed space. These are the classical results of Steinberg(1957), Theorems1--2, pp.348--350, and Geck(2015), Theorem2.1 and Corollary2.3, pp.3--4. They are used as published results.

Let w0 be the longest permutation, and define

\[
 v=\sum_{u\in U}\psi_U(u)^{-1}[uw_0B],\qquad
                         u_0=q^{-n_U/2}v.
 \tag{DO33.2}
\]

Then

\[
 v\in\tau,\quad\tau(u)v=\psi_U(u)v,
 \quad\|v\|^2=q^{n_U},\quad\|u_0\|=1.
 \tag{DO33.3}
\]

Here is the explicit membership and normalization check. The Steinberg vector in the flag model is \(e=\sum_w(-1)^{\ell(w)}[wB]\). Apply \(\sum_{u\in U}\psi_U(u)^{-1}u\) to it. For w≠w0 there is a simple root subgroup in U∩wBw\^{}−1 on which ψ\_U is nontrivial: otherwise w\^{}−1 would send every simple root to a negative root, forcing w=w0. Summing over that subgroup cancels its entire contribution. For w0 its U-stabilizer is trivial, so the result is \((-1)^{\ell(w_0)}v\). This proves membership. Left translation reindexes the sum and gives the asserted character. The q\^{}\{n\_U\} flags in the longest cell are distinct, so the displayed norm is the ordinary flag norm. No orthogonality is being asserted for the different classical basis \(\{ue\}\).

Set, with Hermitian inner product conjugate-linear in its first argument,

\[
 m(g)=\langle u_0,\tau(g)u_0\rangle.
 \tag{DO33.4}
\]

Then m(1)=1, \textbar m(g)\textbar≤1, and every m(g) belongs to q\^{}−n\_U times the ring generated by the values of ψ\_U. This follows because g permutes the orthonormal flag basis in (DO33.2). In particular m(g) is integral at11, and m(1) is a unit there. Also m(gu)=ψ\_U(u)m(g), giving the exact right-generic character needed below.

\subsection{2. An actual residue-matrix Schwartz vector}\label{an-actual-residue-matrix-schwartz-vector}

Use the DO020/022 unitary split position model \(M_{b,a}(F_v)\), with measure giving its integral lattice volume1. Let η be the character on F\_v\^{}× such that the common split character is \((\eta,\eta^{-1})\) in the selected E half. It is tame at this place. The actual source action is

\[
 \omega(g)f(X)=\eta(\det g)^{-1}|\det g|^{b/2}f(Xg).
 \tag{DO33.5}
\]

Indeed DO20.11 is the inverse pullback by the Hom map \(X\mapsto Xg^{-1}\). On the integral compact the absolute-value factor is1 and η depends only on the residue determinant. The inverse-argument operator gives the same spherical vacuum coefficient but is not used as the source action here.

Put \(X_0=\begin{bmatrix}I_a\\0\end{bmatrix}\). The orbit \(X_0\bar G\subset M_{b,a}(\kappa)\) is free. For v'∈τ define a locally constant function, supported on its inverse image in \(M_{b,a}(\mathcal O_v)\), by

\[
 \mathcal E(v')(X)=
 \begin{cases}
 \eta(\det t)\langle u_0,\tau(t)v'\rangle,
                    &X\bmod\mathfrak p=X_0t,\ t\in\bar G,\\
 0,&\text{otherwise}.
 \end{cases}
 \tag{DO33.6}
\]

Freeness makes this unambiguous. The actual source action satisfies

\[
                  \omega(g)\mathcal E(v')
                        =\mathcal E(\tau(\bar g)v')
                      \qquad(g\in\mathrm{GL}_a(\mathcal O_v)).
 \tag{DO33.7}
\]

To check it, evaluate at X=X0t modulo the maximal ideal. The scalar η(det g)\^{}−1 in(DO33.5) cancels the new factor η(det g) in η(det(tg)); the matrix coefficient becomes \(\langle u_0,\tau(t)\tau(g)v'\rangle\). Off the orbit, both sides vanish. Thus the determinant correction in(DO33.6) is fixed by the prescribed splitting; it is not an arbitrary type twist. In the literal source situation η=χ\_0\^{}−1.

Each residue cell has additive measure q\^{}−ab. Finite Schur orthogonality therefore gives the exact identity

\[
 \|\mathcal E(v')\|^2=\gamma_q\|v'\|^2,
 \qquad
 \gamma_q=q^{-ab}\frac{|\mathrm{GL}_a(\kappa)|}{d_\tau}.
 \tag{DO33.8}
\]

In particular the raw choice \(f_q=\mathcal E(u_0)\) has \(f_q(X_0)=1\), norm squared γ\_q, and a pure source τ-type. Its values are algebraic and11-integral. Its unit-norm version \(f_q^\circ=\gamma_q^{-1/2}f_q\) is an isometric τ vector. This normalization is algebraic after adjoining the positive square root; its possible11-adic denominator is retained in section5.

\subsection{3. The exact auxiliary Heisenberg expansion}\label{the-exact-auxiliary-heisenberg-expansion}

The finite model is the normalized actual vacuum model of DO029, not an unrelated finite permutation representation. Here support is L, period is \(\mathfrak pL\), and the finite Fourier quotient has size q\^{}ab. Consequently

\[
 f_q=\sum_{y\in(\mathfrak pL)^\perp/L^\perp}b_{q,y}W_y\phi_{0,v},
 \qquad
 b_{q,y}=q^{-ab}\sum_{X\in M_{b,a}(\kappa)}
                              f_q(X)\psi_v(-B(X,y)).
 \tag{DO33.9}
\]

All actual Fourier and trace-one phases are included in this formula, as in DO29.11--13. Its parameters have denominator at most the rational auxiliary prime when the place is of degree one. They are supported at this auxiliary place. The nonzero coefficients number at most q\^{}ab, and

\[
 \sum_y|b_{q,y}|^2=\gamma_q,
 \qquad \min_{b_{q,y}\ne0}v_{11}(b_{q,y})=0.
 \tag{DO33.10}
\]

The first equality is Parseval with the vacuum-normalized measure. For the second, every coefficient is11-integral, since q is a unit and all character values are roots of unity. If every coefficient had positive valuation then Fourier inversion at X0 could not give f\_q(X0)=1. This also proves that the auxiliary operator acts nontrivially on the vacuum, the actual hypothesis needed in DO030.

Tensor this construction over all selected G places, and keep the vacuum at the other auxiliary places. Write

\[
 \tau_G=\bigotimes_v\tau_v,\quad
 d_G=\prod_vq_v^{a(a-1)/2},\quad
 \Gamma_G=\prod_v\gamma_{q_v},\quad
 s_G=\#\{\text{nonzero tensor Fourier coefficients}\}.
 \tag{DO33.11}
\]

The raw finite auxiliary vector has norm squared Γ\_G, its nonzero coefficients have minimum11-valuation0, and \(s_G\le\prod_vq_v^{ab}\). All these are finite, explicit quantities independent of the selected polynomial degree D. The corresponding unit auxiliary vector has norm1 exactly; its coefficient minimum is \(-\tfrac12v_{11}(\Gamma_G)\).

For the actual degree12 host, six selected relative split places occur at each of g0,g1. Thus

\[
 \begin{split}
 d_G&=(g_0g_1)^{3a(a-1)},\\
 \Gamma_G&=(g_0g_1)^{6(-ab+a(a+1)/2)}
       \prod_{\epsilon=0}^1\prod_{i=1}^a(1-g_\epsilon^{-i})^6,\\
 s_G&\le(g_0g_1)^{6ab}.
 \end{split}
 \tag{DO33.12}
\]

For a host with f selected places above each completely split G prime, the same local product gives \(d_G=(g_0g_1)^{f a(a-1)/2}\), \(\Gamma_G=(g_0g_1)^{f(-ab+a(a+1)/2)}
\prod_{\epsilon,i}(1-g_\epsilon^{-i})^f\), and \(s_G\le(g_0g_1)^{fab}\). These are local counting identities, conditional globally on that host's actual theta-family construction.

\subsection{4. A nonzero selected first kernel and a genuine source input}\label{a-nonzero-selected-first-kernel-and-a-genuine-source-input}

Fix the actual primitive selected polynomial P of degree D at infinity, as in R012, with factorial content bound B\_P from DO30.22. Apply DO030 to the raw auxiliary vector above. Choose

\[
 e\ge1,\quad M=29^e\ge3s_G(D+1)-1.
 \tag{DO33.13}
\]

Then one individual permitted DP59 half-orbit translate λ has

\[
 a_\lambda\ne0,\quad v_{11}(a_\lambda)\le B_P,\qquad
 \vartheta_{r_P\mathcal O_G\phi_\lambda}(1,1)
                                      =C^{2ab}a_\lambda.
 \tag{DO33.14}
\]

This conclusion uses simultaneous jets on the actual allowed section lattice, with the fibre units and the common period already established in DO030. It is not inferred merely from nonvanishing of the finite Schwartz vector. The29 operation commutes with every G-local compact action, so the resulting kernel still has the pure finite source type τ\_G and the selected real source type \(\sigma=\rho_0^\vee\).

At29 take the same level K\_a(M) as R012. At each G place retain the full integral compact, and let the source function vary there according to its finite matrix coefficient. At every other place retain the original integral compact. Let K* denote this product. Its volume is I\_a(M)\^{}−1, with

\[
 I_a(M)=M^{6a^2}\prod_{i=1}^a(1-29^{-i})^6.
 \tag{DO33.15}
\]

The rational stabilizer at the identity class is still trivial: it is finite by definiteness and contained in a torsion-free principal29 congruence. No extra principal congruence at G is necessary to define a matrix-coefficient function on this full compact class.

Let u\_inf be the unit real source vector used in R012, let u\_G=⊗u0, and put

\[
 m_\infty(k)=\langle u_\infty,\sigma(k)u_\infty\rangle,
 \qquad m_G(k_G)=\langle u_G,\tau_G(k_G)u_G\rangle.
 \tag{DO33.16}
\]

On the identity finite K* class define \(F(k_\infty,k_f)=\overline{m_\infty(k_\infty)m_G(k_G)}\), and set F=0 on the other classes. The trivial stabilizer makes this an unambiguous automorphic smooth function. It is cuspidal because the rational source group is definite. Its source real type is ρ0 and its finite G-type is τ\_G\^{}∨.

Two applications of compact Schur orthogonality, retaining the full compact probability measures at G, give

\[
 \boxed{
 \Theta^{\rm bil}_{r_P\mathcal O_G\phi_\lambda}(F)(1)
       =\frac{C^{2ab}a_\lambda}{I_a(M)d_\rho d_G}\ne0,
 \qquad
 \|F\|_2^2=\frac1{I_a(M)d_\rho d_G},\quad\|F\|_\infty=1.}
 \tag{DO33.17}
\]

To verify the extraction rather than just its norm, write the kernel on the supporting class as \(L(\sigma(k_\infty)u_\infty\otimes\tau_G(k_G)u_G)\), where L at the original tensor is the value in(DO33.14). Averaging each factor against its conjugated matrix coefficient returns that tensor divided by d\_rho d\_G. The support volume is I\_a(M)\^{}−1. This proves the first identity, and averaging the squares proves the second.

For every initial source unipotent u at a G place, (DO33.4) implies F(gu)=ψ\_U(u)\^{}−1F(g). This is the literal adelic initial column convention: \nolinkurl{Columns/RecLevel.lean}, lines50--52, defines \nolinkurl{levChr=lam ∘ levRed} and \nolinkurl{levTyp=levChr⁻¹}. Thus the sign agrees with the source's initial generic-unipotent type. The present conclusion does not assert that F meets every other decoration or column constraint of a full \nolinkurl{ColV} datum.

The nonzero theta image in(DO33.17) has at least one irreducible automorphic source constituent with this finite Steinberg K-type. Such a constituent has a nonzero Iwahori- fixed vector because τ\^{}B is one-dimensional. The projection to an irreducible automorphic constituent commutes with the compact actions; a nonzero compact irreducible subspace remains injective in any nonzero projected copy. Hence this is a statement about an actual source constituent, including its automorphic occurrence, not just a local abstract type.

\subsection{5. Norm-one formulas and arithmetic costs}\label{norm-one-formulas-and-arithmetic-costs}

Let \(S_G=\sqrt{\Gamma_G}>0\), and use the same positive square root at the chosen complex embedding. The raw finite Schwartz vector has norm S\_G, whereas the source input F has the norm in(DO33.17). With unit-normalized finite input and source input, the exact value is

\[
 \Theta^{\rm bil}_{r_P\,S_G^{-1}\mathcal O_G\phi_\lambda}
      (\sqrt{I_a(M)d_\rho d_G}\,F)(1)
       =\frac{C^{2ab}a_\lambda}
                    {S_G\sqrt{I_a(M)d_\rho d_G}}.
 \tag{DO33.18}
\]

The real Schwartz norm \(\|r_P\phi_{0,\infty}\|\), including its Gaussian covolume and hook factor, remains R011/R012's exact quantity. If that factor is also made unit, (DO33.18) is divided by that norm.

At11, since every q\_v is a unit,

\[
 v_{11}(d_G)=0,\qquad
 v_{11}(\Gamma_G)=
      \sum_v\sum_{i=1}^a v_{11}(q_v^i-1).
 \tag{DO33.19}
\]

Thus making the finite input unit norm has an exact local scalar cost \(\tfrac12v_{11}(\Gamma_G)\). It is not automatically a unit merely because the G primes differ from11. For the primitive raw input, however, the selected coefficient bound is

\[
 v_{11}\!\left(C^{-2ab}\Theta^{\rm bil}(F)(1)\right)
       \le B_P-v_{11}(I_a(M)d_\rho).
 \tag{DO33.20}
\]

The finite Steinberg dimension introduces no11 valuation. Its archimedean source normalization cost is the additional \(\tfrac12\log d_G\), and the finite input normalization has the explicit S\_G factor. For fixed ranks and fixed G primes all these are fixed numbers. Equation(DO33.13) gives \(\log M=O(\log(D+2))\), with the potentially large explicit coefficient s\_G retained. If the auxiliary primes or ranks vary, their displayed factors must also vary in the cost.

This proves an initial auxiliary-type construction with the chosen ramified splitting characters when those are used. It supplies neither a lower content bound for every row nor the next prescribed raw-Xi map. It also leaves the odd global etaT problem, later integral image saturation, Kval and the MainX terminal mass outside its scope.

\subsection{Sources and proof status}\label{sources-and-proof-status}

The new steps are the explicit generic flag vector metric, its determinant-corrected free-orbit embedding, its actual finite Fourier coefficients, and the finite compact extension of R012's class extraction. The mixed-family and jet comparison are the checked DO020/022/026/029/030 and R012 inputs. LIFT18 supplies the exact even character construction when its root-unit congruence holds.

Steinberg, R. (1957). Prime power representations of finite linear groups II. \emph{Canadian Journal of Mathematics,9},347--351. \url{https://doi.org/10.4153/CJM-1957-041-1} . Theorems1--2, pp.348--350, read directly in the publisher PDF.

Geck, M. (2015). On the ℓ-modular composition factors of the Steinberg representation. \url{https://arxiv.org/pdf/1504.04157} . Theorem2.1, Corollary2.3, pp.3--4, read directly; only their characteristic-zero and integral flag-model assertions are used here.

Source sign: recovered \nolinkurl{Formalism/LocalToGlobal/Columns/RecLevel.lean}, lines50--52; initial finite type \nolinkurl{GenericUnipotentCharacter.lean}437--452 and \nolinkurl{ColForms.lean}92--105, as audited in LIFT17. No Lean compilation or full literal raw-datum instantiation is claimed.

An exact finite diagnostic is saved in \nolinkurl{definite_oscillator_steinberg_certificate.py} and its JSON output. It checks GL2(F3), b=3, with a nontrivial determinant splitting: the flag vector has norm squared3, the free-orbit Schwartz vector has norm squared16/729, all source generic covariance identities hold, and Fourier inversion and Parseval agree exactly in Q(zeta3). It also checks that for a=6, v11(\textbar GL6(F113)\textbar)=1 while the corresponding valuations at281 and337 are0. The finite diagnostic is a normalization check, not a substitute for the general proof above.

\clearpage

\section{AM-R-015: the selected real discrete-series descent and its exact first occurrence}\label{am-r-015-the-selected-real-discrete-series-descent-and-its-exact-first-occurrence}

Author: /root, GPT-6 Astra Pro, ALP\_MACHINE, 21 November2026UTC.

This finding computes the real correspondence on the actual signed Harish--Chandra lists in the source. The insertion block has length t, its compact sign is positive, and its exact first-occurrence invariant is t. Every smaller target obtained by removing a positive number of hyperbolic planes is therefore zero. The analytic consequence is the fifth source predicate in the genuine oriented-kernel interpretation specified below. This is an application of established real theta theorems, with a new explicit check of the selected parameters and the source's conjugation convention. It supplies no integral row tie or arithmetic terminal value.

\subsection{1. The oriented vertical kernel}\label{the-oriented-vertical-kernel}

The direction of the integral must be distinguished from the direction of the Hom model. For the edge from the smaller space Y\_j to the larger space Y\_(j+1), write its oriented theta kernel as \(\theta_j(g,h;\Phi)\), with g on the smaller group. The actual downward operator is

\[
 \mathcal D_j(\Phi,F)(g)
       =\int_{[U(Y_{j+1})]}\theta_j(g,h;\Phi)F(h)\,dh.
 \tag{R15.1}
\]

It is bilinear and has no conjugation on the kernel. Its Hilbert-space adjoint upward operator uses the conjugated kernel. This is exactly \nolinkurl{IndefiniteReal/OrthogonalComplementCharacterisation.lean}, liftInt/IsIo, lines41--55, and IsIoUp, lines78--88; \nolinkurl{LinksChar.lean}, lines283--287, ties ioY/ioX to that convention. \nolinkurl{ChainParSucc.lean}, line132, is the adjoint upward formula. The abstract downward th has reversed group order: \nolinkurl{BesselHermitianTheta/TiedWeil.lean}, lines131--142, specifies the swapped characters and transpose transport without an additional complex conjugation.

Consequently one must not replace (R15.1) by a newly chosen standard Hom kernel from the larger space to the smaller space. That replacement would alter the representation and its character shifts.

The small-group Hom action is contragredient. Transport it by \(g\mapsto g^{-t}\) to the ordinary tensor dual pair, and transport the form and additive character together. In the standard real tensor convention of Ichino section4.1, the additive character is \(x\mapsto e^{-2\pi i x}\). The compact weight formula is his Lemma7.10, including the cross-character shifts. In Hom coordinates the corresponding joint representation is

\[
       \overline{\sigma_j}\boxtimes\sigma_{j+1}
 \tag{R15.2}
\]

when the tensor correspondence pairs \(\sigma_j\) with \(\sigma_{j+1}\). The transpose-inverse on the small factor changes its discrete-series parameter to the contragredient parameter; this can also be read on its compact Cartan, where every weight is negated, followed by the compact Weyl reordering. This statement applies to the odd and indefinite ranks as well: character changes are ratios with trivial real restriction, while the fixed odd baseline is pinned by the standard real compact-weight formula. It does not extrapolate the even-only determinant-root formula of AM-DO-028 to an unspecified odd splitting.

Explicitly, for signatures \((p,q),(r,s)\), the small Hom factor has compact vacuum weights \[
 \bigl(-(r-s)/2-c_{\rm big}\bigr)^{[p]};
 \bigl((r-s)/2-c_{\rm big}\bigr)^{[q]},
\] whereas the large factor has weights \[
 \bigl((p-q)/2+c_{\rm small}\bigr)^{[r]};
 \bigl(-(p-q)/2+c_{\rm small}\bigr)^{[s]}.
\] These are Lemma7.10's two vacuum lists after negating the small compact weights. Their integrality follows from the rank and character parities. They fix the odd baseline as well as all integer character twists. This identifies the standard genuine splitting chosen here; it does not claim that a partial list of abstract source axioms alone uniquely determines a real splitting.

Thus an input F with \(\overline F\) of type \(\sigma_{j+1}\) can have downward type only \(\overline{\sigma_j}\). The parameter calculation proving the indicated correspondence follows.

\subsection{2. Cross-character centering of the actual lists}\label{cross-character-centering-of-the-actual-lists}

Fix a selected complex embedding. Retain the source hypotheses Interlace, TParity and Centred, and write

\[
 n_j=k+m-1+jt,\qquad q_j=\lfloor(m-j)/2\rfloor,
 \qquad p_j=n_j-q_j,\qquad N_j=k+j(t+1).
 \tag{R15.3}
\]

The source odd character has half-angular weight \(h=\pm1/2\), so the parameter \(\eta=-h\) has exactly the two values permitted by TPar.eta\_half. Let \(\chi_i=\operatorname{car}(\mu,j+1,i)\) for this edge. Then

\[
 c_i=\operatorname{infTy}\chi_i
          =-e_j+1_{i\ \mathrm{even}}\eta.
 \tag{R15.4}
\]

If \(D_i^0\) denotes the unshifted signed list in Dpar i, the source definition gives \(D_i=D_i^0+1_{i\ \mathrm{odd}}\eta\). Hence the two cross-character shifts cancel exactly:

\[
       D_j-c_{j+1}=D_j^0+e_j,
       \qquad D_{j+1}-c_j=D_{j+1}^0+e_j.
 \tag{R15.5}
\]

These are the actual half-angular shifts m0/2 and n0/2 in the real theta convention. They are not arbitrary simultaneous translations. The rank parities and the odd half-integral h give the ordinary Harish--Chandra lattice \(\mathbb Z+(n_i-1)/2\) on each group.

Sort each of the following lists in decreasing order:

\[
 \begin{split}
 A_j={}&\{v_l+e_j:0\le l<k-1\}
       \cup\{e_j-b_i:0\le i\le j\}\\
      &\quad\cup\bigcup_{i<j}\bigl(B_t(-e_i)+e_j\bigr),\\
 B_j={}&\{e_j-b_i:i>j,\ i-j\ \mathrm{even}\},\\
 E_j={}&\{e_j-b_i:i>j,\ i-j\ \mathrm{odd}\},\\
 B_t(0)={}&((t-1)/2,(t-3)/2,\ldots,-(t-1)/2).
 \end{split}
 \tag{R15.6}
\]

The letter E\_j here is a list, not a number field. Interlace orders the centers e\_i and the b\_i. Every member of A\_j is positive: the v terms lie above -b\_0, the b\_i terms have i\textless=j, and each earlier block lies entirely above the current center by Centred. Every member of B\_j and E\_j is negative. More precisely Centred gives

\[
   a\ge(t+3)/2\quad(a\in A_j),\qquad
   b,e\le-(t+3)/2\quad(b\in B_j,e\in E_j).
 \tag{R15.7}
\]

All entries are integers. They are distinct: interlacing separates the v and b entries, while Centred separates each inserted block from every other block and from these entries. The centered signed lists are exactly

\[
 \begin{array}{c|c|c}
 &\text{positive signature slot}&\text{negative signature slot}\\\hline
 D_j-c_{j+1}&A_j\cup B_j&E_j\\
 D_{j+1}-c_j&A_j\cup B_t(0)\cup E_j&B_j.
 \end{array}
 \tag{R15.8}
\]

Indeed the future b-signs flip when j increases by one, the old positive entries retain their signs, and the new block is entirely in the positive signature slot. This also follows directly from \nolinkurl{DataTheta.lean}, Dpar, lines538--543.

The exact list lengths are

\[
 |A_j|=N_j,\quad |B_j|=q_{j+1},\quad |E_j|=q_j,
 \tag{R15.9}
\]

so (R15.8) has precisely the source signatures, not merely the correct total ranks:

\[
 (p_j,q_j)=(N_j+q_{j+1},q_j),\qquad
 (p_{j+1},q_{j+1})=(N_j+q_j+t,q_{j+1}).
 \tag{R15.10}
\]

\subsection{3. Nonzero upward correspondence and the exact downward type}\label{nonzero-upward-correspondence-and-the-exact-downward-type}

Apply the sufficiently separated discrete-series case of real theta lifting: Ichino (2022), Lemma8.2, printed p.35, with its explicit regular-case proof citing Li (1990). The applicable separation is \((t+1)/2\); (R15.7) is stronger. The four lists in that result are \((A_j,B_j,\varnothing,E_j)\), and the two extra signature dimensions are t and0, respectively, by (R15.10).

The resulting cohomological induction has compact Levi block U(t) and one-dimensional remaining Levi factors. In the normalized induction parameter its inserted block is zero repeated t times. The compact Levi rho adds the distinct sequence B\_t(0); all remaining entries remain the entries in (R15.8). Equivalently, refine the compact Levi to a compact Cartan and use induction in stages. The inequalities (R15.7) keep this refinement in the good range and the final parameter regular. Thus the representation is the discrete series with exactly the second signed list of (R15.8), after restoring the shift c\_j. It is \(\sigma_{j+1}\), with parameter D\_(j+1).

Therefore, in the tensor convention with these exact characters,

\[
             \theta_{p_{j+1},q_{j+1}}(\sigma_j)
                         =\sigma_{j+1}\ne0.
 \tag{R15.11}
\]

Howe duality then gives the reverse small lift

\[
             \theta_{p_j,q_j}(\sigma_{j+1})=\sigma_j.
 \tag{R15.12}
\]

Only the regular case of Lemma8.2 is used. The general theorem in Ichino's introduction invokes a global Arthur multiplicity formula; that general conditional input is not being used here. Its regular case is the earlier local Li theorem, explicitly isolated in Lemma8.2. No new proof of Li's theorem or of Howe duality is claimed.

\subsection{4. Exact first occurrence and all lower-target vanishing}\label{exact-first-occurrence-and-all-lower-target-vanishing}

For the reverse correspondence, center the larger parameter by the smaller-space character c\_j. By (R15.8) it contains the contiguous central block B\_t(0), all in the positive signature slot. This block is maximal: its next possible endpoints would be \(\pm(t+1)/2\), which are absent by (R15.7). Its length is therefore exactly t in Atobe's Definition1.6, not merely at least t.

The remaining positive entries are A\_j; the remaining negative entries in the positive slot are E\_j; the negative-slot entries are B\_j, all negative. Consequently Atobe's two base indices are

\[
     k_{\sigma_{j+1}}=t,\qquad
     r_{\sigma_{j+1}}=|A_j|+|B_j|=p_j,\qquad
     s_{\sigma_{j+1}}=|E_j|=q_j.
 \tag{R15.13}
\]

Atobe (2020), Theorem1.7(2), last sentence, determines this balanced Witt tower unconditionally for discrete series. Its permitted shifts are precisely the nonnegative integers. The smaller targets in the source have shift -i. Hence

\[
 \boxed{\quad
 \theta_{p_j-i,q_j-i}(\sigma_{j+1})=0
       \qquad(1\le i\le q_j).
 \quad}
 \tag{R15.14}
\]

The theorem concerns the big-theta nonvanishing criterion as well, so these are zero intertwining spaces, not merely zeros of a chosen Schwartz vector. The discrete-series clause is explicitly declared unconditional on printed p.5 of Atobe's published paper. The later correction to its separate k=-1 tempered case is irrelevant because here k=t\textgreater=1.

For the additive-character convention, Ichino section6.2, printed p.20, explicitly identifies Atobe's convention with his section4.1 choice, using the compact-weight comparison of Lemma7.10. His Theorem6.1(ii) gives the same balanced-tower conclusion here: its exceptional negative-shift possibility requires both next central endpoints, whereas (R15.7) excludes both. Thus a change of the additive-character sign is not being silently absorbed into the signature labels.

\subsection{5. Consequence for the genuine fifth source predicate}\label{consequence-for-the-genuine-fifth-source-predicate}

Use the genuine smooth cusp domain of AM-R-014 and the downward orientation (R15.1). If \(\overline F\) has the prescribed real type D\_(j+1), then F has the contragredient type. Equation(R15.2), Howe duality and (R15.12) imply that every nonzero cuspidal downward image has contragredient real type D\_j. Thus its complex conjugate has type D\_j, which is exactly IsDv at stage j.

One can justify the assertion constituent by constituent in the cuspidal unitary spectrum: every irreducible real constituent of the image is a quotient of the corresponding big theta module, whose unique irreducible quotient is fixed by (R15.12). The actual cusp inputs lie in finite sums of smooth automorphic embeddings, as specified in AM-R-014; no decomposition of arbitrary smooth functions is assumed. For compact-finite test data the assertion holds on the Harish--Chandra modules and then their smooth globalizations.

For a lower target, (R15.14) makes every such local intertwining space zero at each chosen real place. It therefore kills the global theta integral for every Schwartz input. Equivalently, first use the zero Harish--Chandra theta module on dense compact vectors and extend by continuity of the cuspidal theta integral in the usual Schwartz/smooth-cusp topologies. The measures and normalizing nonzero scalars do not affect this vanishing.

These are exactly the two conclusions of \nolinkurl{CayleySignature.RealThetaLiftDiscreteSeries}, lines186--194. Together with AM-R-013/014, they verify five of the six analytic predicates in the explicitly stated genuine model. The remaining analytic predicate is ThetaUnfoldingFreeData, with its precise Fourier--Jacobi period and factorizable witness. A full arithmetic ThetaX/D7 instance, the odd global character pin on a repaired host, the raw-Xi selected action and the terminal mass remain separate.

\subsection{6. Geometric framework and source status}\label{geometric-framework-and-source-status}

The signed weights, central character shifts and compact-Levi Schur classes live in SGA3 I Proposition4.7.3/Corollary4.7.3.1 and the carried SGA6 VI Theorem3.3 representation/flag framework. The oriented kernel remains an actual polarized Hom/theta-group line, in the SGA7 I VII2.7/2.9 setting, pp.169--174. Neither a virtual class nor an isocrystal slope proves the analytic real correspondence; the specific local theorems cited above supply that independent input.

The source lists, parity and centering are DataTheta.lean538--555 and \nolinkurl{CayleySignature.lean104–125.} The rank/signature formulas are Params.lean and the definitions sigYD/rankYD in the recovered source. No Lean execution is claimed.

\subsection{References}\label{references-4}

Atobe, H. (2020). On the non-vanishing of theta liftings of tempered representations of U(p,q). \emph{Advances in Mathematics, 363},106984. \url{https://doi.org/10.1016/j.aim.2020.106984} . Author repository copy: \url{https://eprints.lib.hokudai.ac.jp/repo/huscap/all/84508/Adv.\%20math\%20363_106984.pdf} . Definition1.6 and Theorem1.7, printed pp.4--5; the unconditional discrete-series scope is explicit immediately after that theorem.

Ichino, A. (2022). Theta lifting for tempered representations of real unitary groups. \emph{Advances in Mathematics, 398},108188. \url{https://doi.org/10.1016/j.aim.2022.108188} . Author preprint: \url{https://arxiv.org/pdf/2008.06174} . Section4.1 p.10; Lemma7.10 pp.30--31; regular-case Lemma8.2 p.35 and its proof. The general conditional theorem is distinguished from this local regular case. Theorem6.1(ii) and section6.2, printed pp.19--20, fix the balanced tower and the additive-character convention.

Li, J.-S. (1990). Theta lifting for unitary representations with nonzero cohomology. \emph{Duke Mathematical Journal, 61}(3),913--937. \url{https://doi.org/10.1215/S0012-7094-90-06135-6} . The precise regular case used here is read in Ichino's Lemma8.2; direct retrieval of the original publisher text was blocked by a service security page. No direct whole-paper reading of Li is claimed.

Review status: author whole-proof reading PASS, 2026-11-21T23:40:17.491120+00:00. Independent whole-proof review by /root/definite\_oscillator also PASS, recorded separately in definite\_oscillator\_review\_r015.md. The author reread the complete revised proof, including the explicit compact vacuum shifts, exact source down/adjoint distinction, regular local input, first-occurrence invariant and smooth-cusp scope. No arithmetic constructor or Lean build is certified.

\clearpage

\section{AM-R-016: the final two free-data transfers, and all six analytic predicates at lift count two}\label{am-r-016-the-final-two-free-data-transfers-and-all-six-analytic-predicates-at-lift-count-two}

Author: /root, GPT-6 Astra Pro, ALP\_MACHINE, 21 November2026UTC.

This finding proves the third analytic predicate for the final two edges of the concrete genuine tower of R014. Consequently, when m=2, it supplies the entire \nolinkurl{ThetaUnfoldingFreeData} predicate. Together with R013--015 it gives all six mathematical analytic predicates in that interpretation. The conclusion concerns freely chosen Schwartz data. It does not supply the source's arithmetic host constructor, its prescribed raw Xi maps, an integral row identification, or the cost difference of DP44.

The point is that these two edges involve only ordinary compact restriction and rank-zero Fourier--Jacobi periods. There is no positive-rank unipotent unfolding in this case. We retain the conjugations, the odd splitting character and the Haar measures.

\subsection{1. Spaces, characters and oriented kernels}\label{spaces-characters-and-oriented-kernels}

Let E/F be the CM extension used by the genuine analytic construction, with its fixed trace-zero coefficient xi and standard additive character. Let eta be its chosen unitary odd splitting character, with restriction epsilon\_(E/F) to the base-field ideles. Let mu be any even splitting character, trivial on the base-field ideles. The finite conductor of eta is allowed in this analytic assertion. On the old host its everywhere-unramified pin is obstructed by K028; K029 supplies that pin on its larger field.

Use the standard rational spaces of R014, with k,m even and t odd. Put d=t+1. At the final three indices write

\[
 A=X_{m-2}=I_a,\quad B=X_{m-1}=I_b,\quad
 C=X_m=I_c,\qquad b=a+d,\quad c=b+d.
 \tag{R16.1}
\]

Let L be the last positive coordinate line of C and put H=L\^{}perp. The corresponding vertical spaces and embeddings are exactly

\[
 Y_{m-2}=A\perp(-L),\qquad Y_{m-1}=B,\qquad Y_m=H.
 \tag{R16.2}
\]

Here L and -L have Gram matrices 1 and -1. Thus their dimensions and signatures agree with R014: the first space in(R16.2) has signature(a,1), the second is positive of rank b, and the last is positive of rank c-1. The positive embeddings are the first summand A, the identity of B, and H into C, respectively. All the quotients {[}U(A){]}, {[}U(B){]} and {[}U(H){]} are compact.

Write theta\_(S,T)\^{}(alpha,beta)(g,h;Phi) for the oriented Hom theta kernel, where alpha is the splitting character attached to S and beta the character attached to T. Its Hom symplectic form and action are

\[
 \mathcal B_{S,T}(X,Y)
   =-\operatorname{Tr}_{E/F}\bigl(\xi\operatorname{tr}(Y^*X)\bigr),
 \qquad X\longmapsto hXg^{-1}.
 \tag{R16.3}
\]

The adjoint is taken for the indicated Hermitian forms. We use the actual compatible global Weil splittings and self-dual additive measures. Orthogonal sums have product measures and the tensor product theta distribution. Every automorphic quotient carries the same fixed Haar measure at each occurrence below; its total volume is not set to one unless already so normalized.

The downward operator uses this smaller-first kernel without conjugation:

\[
 D_{S,T}(\Phi,F)(g)
       =\int_{[U(T)]}\theta_{S,T}(g,h;\Phi)F(h)\,dh.
 \tag{R16.4}
\]

This is R015's convention, and agrees with the actual downward operator after the source's rank-dependent transpose transport. We do not replace it by a Hermitian adjoint.

For edge j=m-2, which is even, \nolinkurl{car} prescribes characters (mu eta\^{}(-1),mu) on (Y\_(m-2),B). The horizontal pair(A,B) has characters(mu,mu). For edge j=m-1, which is odd, it prescribes(mu,mu eta\^{}(-1)) on(B,H), while(B,C) has(mu,mu). Here mu is the character for the particular edge; the two edge characters need not coincide.

\subsection{2. The one common rank-zero Fourier--Jacobi kernel}\label{the-one-common-rank-zero-fourierjacobi-kernel}

On the pair(B,-L) choose a Weil representation with character eta\^{}(-1) attached to -L. Its other splitting character may be any character trivial on the base field. Define

\[
 \vartheta_-(g;s)=\theta_{B,-L}(g,1;s),\qquad
 \mathcal J_B(\phi,F;s)
       =\int_{[U(B)]}\phi(g)F(g)
                         \overline{\vartheta_-(g;s)}\,dg.
 \tag{R16.5}
\]

The restriction g-\textgreater theta\_(B,-L)(g,1;s) is independent of that other splitting character after the standard character-change identification: changing the character attached to B multiplies the action of U(-L), and that argument is the identity. Therefore the same space of kernels(R16.5) can be used on both adjacent edges even when their mu characters differ.

This is precisely the rank-zero Fourier--Jacobi period: its unipotent group is trivial, its rational quotient has volume1, and its auxiliary theta factor is conjugated. The source formula is \nolinkurl{LastSeesawPeriodNonvanishing.lean}, theorem \nolinkurl{per_zero}, lines58--94, together with \nolinkurl{fjSrcFun} in \nolinkurl{WittSmallGroupEmbeddingMemPointwiseStabilizerIso.lean}, lines228--233. In the source theta-function model, evaluation of the transformed theta function at zero is the displayed theta kernel; it is not evaluation of a bare Schrödinger function.

We require two elementary compatible realizations of this same auxiliary kernel. First, conjugating the Hom Weil representation for(B,L) changes L to -L and inverts both splitting characters. For its character pair(mu,eta), it gives a conjugate-linear Schwartz isomorphism C\_L with

\[
 \theta_{B,L}^{\mu,\eta}(g,1;q)
          =\overline{\vartheta_-(g;C_Lq)}.
 \tag{R16.6}
\]

Second, Hermitian adjoint is anti-symplectic between Hom(-L,B) and Hom(B,-L). Indeed tr((Y\^{}\emph{)\^{}}X\^{}\emph{)=overline(tr(Y\^{}}X)), and bar(xi)=-xi; substitution in(R16.3) changes the sign of the form. Adjoint also interchanges the two group actions. Composing its Schwartz transport with complex conjugation therefore gives a conjugate-linear isomorphism C\_ad satisfying

\[
 \theta_{-L,B}^{\eta^{-1},\mu}(1,g;q)
             =\overline{\vartheta_-(g;C_{\rm ad}q)}.
 \tag{R16.7}
\]

Under this conjugate-adjoint exchange the splitting characters are exchanged, giving(mu,eta\^{}(-1)); they are not both inverted again. The interchange of source and target determinant actions cancels those inversions. At a split place the source scalar is chi\_T(det g)\^{}(-1) and the target scalar is chi\_S(det h), which verifies this cancellation directly. At infinity the auxiliary U(B) vacuum weight is 1/2+h when eta has half-angular weight h; the two unbarred line kernels in(R16.6)--(R16.7) both have weight -1/2-h. This also checks the odd baseline.

For completeness, the normalization of(R16.6)--(R16.7) has no free scalar. Negating one Hermitian form negates the Heisenberg form; complex conjugation changes psi to psi\^{}(-1), so the two changes cancel. Adjoint carries rational position points to rational position points when the polarizations are transported with it. Hence conjugation commutes with the rational theta sum. The self-dual measures are transported at the same time. The linear, quadratic and Fourier generators of the Weil representation respect these maps; the Fourier identity is the usual conjugation of its integral kernel. Thus the equality includes the theta distribution itself, not just an abstract representation isomorphism.

These are also the exact general theta identities recorded in \nolinkurl{IsWeilDatumConjNegHermitianSpaces.lean}, \nolinkurl{weil_conjTarget}, and \nolinkurl{ConjWeilDatumViaAdjointExchangedPairStatement.lean}. They are used as accepted source identities within their scope. The Hom computation above explains their identical normalization for the genuine E/F model without asserting that the arithmetic D7 record has been constructed.

\subsection{3. Terminal transfer: a positive line and a rank-zero period}\label{terminal-transfer-a-positive-line-and-a-rank-zero-period}

Fix a smooth cusp form phi on U(B) and a smooth cusp form F on U(H). The source's terminal bilinear period is

\[
 \mathcal P_H(f,F)=\int_{[U(H)]}f(\iota h)F(h)\,dh
                 =B_{\rm top}(\overline f,F).
 \tag{R16.8}
\]

Orthogonal factorization C=H perp L, with the character identity (mu eta\^{}(-1))*eta=mu, gives the exact theta identity

\[
 \theta_{B,C}^{\mu,\mu}(g,(h,1);\Phi\boxtimes q)
  =\theta_{B,H}^{\mu,\mu\eta^{-1}}(g,h;\Phi)
                    \overline{\vartheta_-(g;C_Lq)}.
 \tag{R16.9}
\]

To check factorization, the Hom space is the symplectic direct sum Hom(B,H) plus Hom(B,L). Its chosen position lattice over F is the corresponding direct sum; the rational theta sum of a product is the product of the two absolutely convergent sums. The tensor Weil action restricts to the displayed product. The target character multiplication above is exactly its splitting requirement. This is the standard orthogonal-sum theta tensor rule, with no index or volume factor.

Fubini on the two compact automorphic quotients now gives

\[
 \boxed{
 \mathcal P_H\bigl(\Theta_{B,C}(\phi;\Phi\boxtimes q),F\bigr)
       =\mathcal J_B\bigl(\phi,D_{B,H}(\Phi,F);C_Lq\bigr).}
 \tag{R16.10}
\]

This equality retains the two fixed Haar measures on both sides. The vertical kernel is unbarred; only the auxiliary rank-zero kernel is conjugated.

Suppose the left terminal period is nonzero for an arbitrary Schwartz input Phi'. Then it is nonzero for some input Phi boxtimes q with Phi factorizable and real compact-finite, and q likewise. Here is the needed free-data argument. The left side is a continuous linear functional of Phi' on the Schwartz space. At finite places, compactly supported locally constant functions on a direct sum are finite sums of products. At real places, finite sums of product Hermite functions are dense in the Schwartz topology. For the positive spaces B,H,L, these polynomial-Gaussian functions are finite under the relevant compact unitary groups. Finite sums of global restricted pure tensors of such functions are therefore dense. If the functional vanished on every pure product it would vanish on this dense span and hence on Phi', a contradiction. Choose one product on which it is nonzero. Equation(R16.10) then gives the required nonzero lower Fourier--Jacobi period.

The continuity assertion can be seen directly without an interchange of a merely formal series. At fixed finite support and period, the adelic quotients have finitely many charts with compact real parameters. The rational theta series and all fixed-order derivatives are uniformly bounded on those charts by a Schwartz seminorm and a convergent lattice majorant. Multiplying by the fixed smooth forms and integrating over the compact quotients is continuous. The finite-place Schwartz space is the union of these fixed support/period spaces.

\subsection{4. Penultimate transfer: ordinary Bessel restriction}\label{penultimate-transfer-ordinary-bessel-restriction}

Write V=A perp(-L), as in(R16.2), and fix a horizontal factorizable compact-finite input Phi' for(A,B). The source decomposition and the character identity mu*eta\^{}(-1)=mu eta\^{}(-1) give

\[
 \theta_{V,B}^{\mu\eta^{-1},\mu}((g,1),h;
                                      \Phi'\boxtimes q)
   =\theta_{A,B}^{\mu,\mu}(g,h;\Phi')
                         \overline{\vartheta_-(h;C_{\rm ad}q)}.
 \tag{R16.11}
\]

Let phi be a smooth cusp form on U(A), F one on U(B), and define the ordinary codimension-one Bessel period by

\[
 \mathcal B_A(\phi,G)=\int_{[U(A)]}\phi(g)G((g,1))\,dg.
 \tag{R16.12}
\]

The Bessel unipotent rank is zero. Thus there is no nontrivial unipotent character or hidden unipotent volume in(R16.12). Fubini gives the exact identity

\[
 \boxed{
 \mathcal B_A\bigl(\phi,D_{V,B}(\Phi'\boxtimes q,F)\bigr)
   =\mathcal J_B\bigl(\Theta_{A,B}(\phi;\Phi'),F;
                                                    C_{\rm ad}q\bigr).}
 \tag{R16.13}
\]

If the right side is nonzero for some auxiliary datum s, it is nonzero for some factorizable compact-finite s. The same continuity argument applies, now to the auxiliary Schwartz space alone. Polynomial-Gaussian functions for the pair(B,-L) are finite under U(B) times U(1), and their span is dense. Choose one global pure tensor with nonzero period and take q=C\_ad\^{}(-1)s. The induced rational adjoint transport preserves pure tensors and compact finiteness. Then Phi' boxtimes q is factorizable and finite under K\_V times K\_B. At infinity K\_V=U(A) times U(1), and the diagonal U(B) action on the tensor product of two finite modules is finite. Equation(R16.13) is the required lower nonzero Bessel period.

No cuspidality claim for the noncompact group U(V) has been inserted into Fubini: only its restriction to U(A) is integrated. For the selected input type, R015 gives vanishing at all smaller targets and R014 then gives cuspidality of this actual downward lift. Thus the constructed form belongs to the genuine cusp domain when that is required by the source. Independently, the literal free-data predicate quantifies over every cusp form G equal to this lift; either reading uses the same identity(R16.13).

\subsection{5. Quantifiers, scope and exact implications}\label{quantifiers-scope-and-exact-implications}

For j=m-1, a nonzero terminal \nolinkurl{PerNZ} is(R16.8) and the existential free datum is supplied by(R16.10) and section3. For j=m-2, upper \nolinkurl{PerNZ} is exactly(R16.5), and(R16.13) supplies the lower ordinary Bessel period. The \nolinkurl{pos} choices are the actual standard positive embeddings in(R16.2). The functions in the source \nolinkurl{bB} have the conjugation \nolinkurl{Bj(conjAF(phi),F)}; with the ordinary Hermitian first slot this is precisely the bilinear phi*F convention used here. More explicitly, in this genuine interpretation the rank-zero \nolinkurl{Pos} datum at the identity embedding B into B is the unique empty isotropic flag on its zero orthogonal complement. Its \nolinkurl{FJd} is the full Schwartz space for the fixed auxiliary pair (B,-L), with the cross-character eta\^{}(-1). Thus the existential upper \nolinkurl{PerNZ} ranges over exactly the kernels in(R16.5), rather than over an unexamined additional class of positions. At the lower codimension-one Bessel step the isotropic flag is again empty; the nonisotropic coordinate in the complementary line does not occur in the rank-zero unipotent character, so every such coordinate choice gives the same ordinary restriction. At j=m-1 the source uses the auxiliary conjugated theta factor, and at j=m-2 the period is independent of any auxiliary FJ datum because its Bessel unipotent rank is zero.

The conclusion is uniform in the two edge characters mu, the chosen forms and the nonzero upper period. It does not assume a special nonzero matrix coefficient or a selected automorphic multiplicity copy. It therefore proves the literal mathematical quantifiers of the third predicate at these final two indices for the genuine model specified here.

When m=2 the set of indices j\textless m is exactly\{0,1\}. Consequently R013--016 cover all six analytic predicates for that model, provided the stated genuine characters and theta spaces are supplied. This closes the analytic free-data gate for lift count2. It does not identify these freely chosen kernels with the raw arithmetic PhiY/PhiG/Xi family. In particular, the density argument provides existence and no bound on the selected finite level, coefficient denominators or norms.

For m\textgreater2 the earlier free-data edges still involve genuine positive-rank Bessel/Fourier--Jacobi unfolding. The accepted source UpBF/UpFB statements also require their local test-space package. Those additional edges are not proved in this note.

\subsection{6. Exact locators and check status}\label{exact-locators-and-check-status}

\begin{itemize}
\tightlist
\item
  Alpöge, L. (2026). \emph{Remark}, step8, with the stipulated pinned archive at commit7cbf7a0c32fca1badf05931879a37c3f8a2b03b4. \nolinkurl{Formalism/SharpRestriction/CuspSpanNormBounds/CayleySignature.lean}, \nolinkurl{car}114--115, \nolinkurl{bB}131--133, \nolinkurl{PerNZ}150--152, \nolinkurl{ThetaUnfoldingFreeData}168--175.
\item
  The same archive, \nolinkurl{Development/AbelianVarietyTheory/Proofs/TwistedOperators/LastSeesawPeriodNonvanishing.lean}, \nolinkurl{per_zero}58--94; and \nolinkurl{Definitions/TwistedInnerProducts/WittSmallGroupEmbeddingMemPointwiseStabilizerIso.lean}, \nolinkurl{fjSrcFun}228--233.
\item
  The same archive, \nolinkurl{Proofs/DecoratedLatticeProofs/ThetaTensorRuleTargetOrthogonalSum.lean}, \nolinkurl{tensorTarget}75ff.; \nolinkurl{ThetaTensorRuleOrthogonalSumSourceStatement.lean}; \nolinkurl{IsWeilDatumConjNegHermitianSpaces.lean}, \nolinkurl{weil_conjTarget}71--136; \nolinkurl{ConjWeilDatumViaAdjointExchangedPairStatement.lean}.
\item
  ALP\_MACHINE R013 section1, R014 sections2--4, R015 sections1--2: the genuine spaces, actual bilinear conventions, character parity and the already checked remaining analytic predicates.
\end{itemize}

Whole-proof review status: author complete rereading PASS, 2026-11-21T23:56:58.217024+00:00. Independent whole-proof review by /root/definite\_oscillator PASS on22November2026UTC, including the saved rank-zero Pos/FJd clarification; see agents/definite\_oscillator\_review\_r016.md and CHECK\_EVENTS.tsv. No Lean build is requested or performed. The new mathematical work is the exact final-two-edge specialization, the compatible single auxiliary kernel, the compact Fubini identities and the factorizable compact-finite witness argument.

\clearpage

\section{ALP-LIFT-19 --- ramified theta constituents with the literal nonterminal column types}\label{alp-lift-19-ramified-theta-constituents-with-the-literal-nonterminal-column-types}

Author: /root/lift\_accounting, GPT-6 Astra Pro, ALP\_MACHINE. Written 21 November2026UTC.

This theorem repairs the local representation obstruction of ALP-LIFT-17. At a degree-one split G-place, start with an irreducible representation containing the finite Steinberg type, and use the distinct tame stage characters required by IsMu0. Every later local theta representation then contains the literal adelic column type, with the correct inverse character. An explicit finite induced vector is primitive in its natural coefficient lattice and has a computable positive Hermitian norm.

The theorem is local and independent of the degree of the global CM field. It therefore applies to the old degree12 host or to an enlarged host whenever the same degree-one G-place and stage-character hypotheses hold. The global character construction of LIFT18 is used only for its stated old host. No class-number-one assertion, G-prime example, or mass formula is silently transported to the degree84 repaired host.

The finite Steinberg representation, Casselman's Iwahori theorem, Zelevinsky's irreducibility criterion, and Mínguez's local theta theorem are literature inputs. The new work here is their exact assembly with the source columns, the explicit Mackey component and its primitive vector/norm, and the consequent removal of this specific local type obstruction. The literal raw PhiG/Xi kernel and the integral lattice of the iterated global image remain separate.

\subsection{1. Local data and the source sign}\label{local-data-and-the-source-sign}

Let F be a nonarchimedean local field of characteristic zero with residue field \(\kappa=\mathbf F_q\). For the literal degree-one G places, q is the rational G prime. Put \[
 K_n=\operatorname{GL}_n(\mathcal O_F),\qquad
 K_n^1=1+\varpi M_n(\mathcal O_F).
 \tag{T1}
\] Fix integers \(k\ge1,t\ge1,m\ge2\), set \(d=t+1\), and write \[
 n_j=k+jd=(k+j)+jt,\qquad r_j=k+j.
 \tag{T2}
\] Let \(\eta_0,\ldots,\eta_{m-1}:F^\times\to\mathbf C^\times\) be unitary tame characters whose residue characters \(\bar\eta_i:\kappa^\times\to\mathbf C^\times\) are pairwise distinct and nontrivial.

At a selected source G-prime, IsMu0 and chi\_om say exactly \[
 \eta_i|_{\mathcal O_F^\times}
       =\chi_i^{-1},\qquad
 \chi_i=\omega^{\,2(i+1)},
 \tag{T3}
\] where \(\omega\) is faithful. The source inequality \(q-1>2fk>2m\) ensures the required nontriviality and distinctness. This is the character condition already audited in LIFT17; LIFT18 constructs actual global even characters with it on its specified weight locus over the degree12 host.

The adelic column convention has an inverse. In the pinned source, Columns/RecLevel.lean50--52 defines \[
 \operatorname{levChr}(\lambda)=\lambda\circ\operatorname{levRed},
 \qquad
 \operatorname{levTyp}(\lambda)=\operatorname{levChr}(\lambda)^{-1}.
 \tag{T4}
\] Consequently the desired adelic right-translation character at stage \(j<m\) is \(\lambda_{X,j}^{-1}\), whose determinant-block factors are the \(\bar\eta_i\) in(T3).

Choose a nontrivial additive character \(\psi:\kappa\to\mathbf C^\times\). For \(0\le j<m\), define \[
 \begin{split}
 H_j&=\bigl(U_{r_j}(\kappa)\times
                 \operatorname{GL}_t(\kappa)^j\bigr)N_{P_j},\\
 P_j&=P_{(r_j,t,\ldots,t)}\subseteq\operatorname{GL}_{n_j}(\kappa),\\
 \lambda_j(h_0,h_1,\ldots,h_j;n)
   &=\psi_{S_j}(h_0)^{-1}
                    \prod_{i=0}^{j-1}\bar\eta_i(\det h_{i+1}),\\
 S_j&=\{0\le a<k+j-1:\ a+1\ne k\text{ or }j+1=m\}.
 \end{split}
 \tag{T5}
\] Here \(U_r\) is upper unitriangular, and \(\psi_S\) applies \(\psi\) to the sum of the indicated simple-root entries. This is exactly HX, SX and the inverse of lamX in \nolinkurl{GenericUnipotentCharacter.lean437–452.} The character is trivial on \(N_{P_j}\). For \(j=0\), it is the full generic character of \(U_k\). For \(1\le j<m-1\), it omits only the root joining the first k coordinates to the last j coordinates in \(U_{k+j}\). For \(j=m-1\), it is fully generic in \(U_{k+j}\). The terminal HXm case \(j=m\) is not included.

\subsection{2. The exact local representation theorem}\label{the-exact-local-representation-theorem}

Let \(\pi_0\) be an irreducible admissible representation of \(\operatorname{GL}_k(F)\) such that \[
 \operatorname{St}_k(\kappa)\hookrightarrow\pi_0^{K_k^1}
 \quad\hbox{as }\operatorname{GL}_k(\kappa)\hbox{-modules}.
 \tag{T6}
\] For \(j\ge0\), form normalized parabolic induction \[
 \Pi_j=
 \pi_0\times(\eta_0\circ\det_{\operatorname{GL}_d})
       \times\cdots\times
                (\eta_{j-1}\circ\det_{\operatorname{GL}_d}).
 \tag{T7}
\]

\textbf{Theorem.} For every \(0\le j\le m\), the representation \(\Pi_j\) is irreducible. For each \(0\le j<m\), \(\Pi_{j+1}\) is the actual local theta descendant of \(\Pi_j\) under the literal bilinear Hom convention with the same splitting character \(\eta_j\) on both members of edge \(n_j\to n_{j+1}\). For \(0\le j<m\), its finite column space is nonzero: \[
                   \Pi_j^{\widetilde H_j,\lambda_j}\ne0,
 \tag{T8}
\] where \(\widetilde H_j\) is the full inverse image of \(H_j\) in \(K_{n_j}\). Thus the desired finite type occurs in the entire actual ramified theta representation. This is the opposite local conclusion from LIFT17's unramified chain.

\subsubsection{2.1 Cuspidal support and irreducibility}\label{cuspidal-support-and-irreducibility}

The finite Steinberg representation has a one-dimensional Borel-fixed subspace. For example, in its integral basis \(\{u e:u\in U_k\}\), the vector \(\sum_u u e\) is fixed by \(U_k\), and diagonal matrices permute its summands because they fix e and normalize \(U_k\). Thus(T6) gives a nonzero Iwahori-fixed vector in \(\pi_0\).

Casselman(1980), Proposition2.6, printed p.396, states that every irreducible admissible representation with a nonzero Iwahori-fixed vector embeds in an unramified principal series. Hence the cuspidal support of \(\pi_0\) consists entirely of unramified characters of \(F^\times\). This statement applies to \(\operatorname{GL}_k\), including its center; Casselman's paper is formulated for connected reductive groups.

The cuspidal support of the one-dimensional representation \(\eta_i\circ\det_{\operatorname{GL}_d}\) is \[
 \bigl\{\eta_i|\cdot|^{(1-d)/2},
             \eta_i|\cdot|^{(3-d)/2},\ldots,
                         \eta_i|\cdot|^{(d-1)/2}\bigr\}.
 \tag{T9}
\] Its residue character is uniformly \(\bar\eta_i\). The distinctness in(T3) therefore puts different factors of(T7) on disjoint cuspidal lines; none meets the unramified support of \(\pi_0\). No two segments from different factors can be linked, since an unramified absolute-value twist does not change the residue character.

Zelevinsky(1980), Proposition8.5, printed p.196, now says that the product(T7) is irreducible. Its Proposition8.6, printed p.197, records the equivalent factorization by distinct cuspidal lines. This argument does not require \(\pi_0\) to be tempered or generic. If \(\pi_0\) is unitary, so is(T7), because every inducing factor is unitary and normalized parabolic induction preserves unitarity.

\subsubsection{2.2 Identification with the actual theta descendant}\label{identification-with-the-actual-theta-descendant}

At a relative split place, write the common unitary character as \((\eta_j,\eta_j^{-1})\) on \(F^\times\times F^\times\). The actual Hom oscillator carries the scalar \[
             \eta_j(\det g)^{-1}\eta_j(\det h).
 \tag{T10}
\] The source lift is bilinear, \(\int\theta(g,h)f(g)\,dg\), so the untwisted type-II source representation is \(\Pi_j^\vee\otimes\eta_j\). Mínguez(2008), Theorem1(2), printed p.718, gives nonzero local theta for every irreducible input in the smaller rank. Its target Langlands quotient has the contragredient old parameters and the parameters of the rank-d trivial representation. Restoring(T10) therefore leaves the old \(\Pi_j\) block unchanged and inserts \(\eta_j\circ\det_{\operatorname{GL}_d}\).

This identifies the target with(T7) for \(j+1\): the induced product is irreducible by section2.1, and its Langlands quotient is precisely the one in Mínguez's theorem. Equivalently its standard parameter is the direct sum of the old parameter and(T9). This use of the Langlands quotient is not an assertion that an arbitrary reducible parabolic induction equals small theta.

The same convention was checked against the actual Hom action in LIFT14(E10--E11) and AM-DO-028. The finite right-translation character in(T4) is retained; no Hermitian-input convention is substituted.

It remains to prove(T8). We do so constructively.

\subsection{3. A finite Mackey component adapted to the source columns}\label{a-finite-mackey-component-adapted-to-the-source-columns}

For this section write \(G=\operatorname{GL}_{n_j}(\kappa)\) and \[
 Q=P_{(k,d,\ldots,d)},\qquad
 \mathcal M_j=
 \operatorname{Ind}_{Q}^{G}
       \left(\operatorname{St}_k\boxtimes
                  \mathop{\boxtimes}\displaylimits_{i=0}^{j-1}\bar\eta_i\circ\det_d\right).
 \tag{T11}
\] Finite induction is unnormalized; its function model consists of functions \(f:G\to V_{\operatorname{St}}\) with \(f(qg)=\tau(q)f(g)\), acted on by right translation.

Give the Q blocks the coordinate order \[
        A,\ L_0,T_0,\ L_1,T_1,\ \ldots,L_{j-1},T_{j-1},
 \qquad \dim A=k,\quad\dim L_i=1,\quad\dim T_i=t.
\] Let w be the permutation sending the P-coordinate order \[
            A,L_0,\ldots,L_{j-1},\ T_0,\ldots,T_{j-1}
\] to this Q-coordinate order. Put \(H=w^{-1}Qw\cap P_j\). The image of H in the Levi \(L_{P_j}=\operatorname{GL}_{r_j}\times\operatorname{GL}_t^j\) is exactly \[
              R=P_{(k,1,\ldots,1)}
                       \times\operatorname{GL}_t^j.
 \tag{T12}
\] This can be checked entry by entry. Within the first P-block, the Q-upper-triangular conditions impose the parabolic blocks \(k,1,\ldots,1\); every matrix with those conditions is allowed. The T-block diagonal matrices are unrestricted. Thus every element of R has a lift in H.

The representation \(\tau\) in(T11) is trivial on \(w(H\cap N_{P_j})w^{-1}\). Indeed off-diagonal Q blocks act trivially by inflation. Within \(L_i\oplus T_i\), the remaining intersection has matrices \(\left(\begin{smallmatrix}1&*\\0&I_t\end{smallmatrix}\right)\), on which a determinant character is trivial. On the image R, the action is \[
 \operatorname{St}_k\boxtimes
     (\bar\eta_0,\ldots,\bar\eta_{j-1})
       \boxtimes
          \mathop{\boxtimes}\displaylimits_{i=0}^{j-1}\bar\eta_i\circ\det_t.
 \tag{T13}
\]

It follows that \(\mathcal M_j^{N_{P_j}}\) contains the component \[
 \mathcal V_j\boxtimes
       \mathop{\boxtimes}\displaylimits_{i=0}^{j-1}\bar\eta_i\circ\det_t,\qquad
 \mathcal V_j=
 \operatorname{Ind}_{P_{(k,1,\ldots,1)}}^{\operatorname{GL}_{r_j}}
        (\operatorname{St}_k\boxtimes\bar\eta_0\boxtimes
                                  \cdots\boxtimes\bar\eta_{j-1}).
 \tag{T14}
\] Here is an explicit proof, so a quotient in a Mackey filtration is not mistaken for a submodule. Take functions in(T11) supported on the finite double coset QwP\_j and invariant under right \(N_{P_j}\). Their restriction to wP\_j is the induced-function model from H to P\_j. The trivial action of \(H\cap N_{P_j}\) just proved makes such a function descend to \(L_{P_j}\). Its left transformation law is exactly(T13) on R. Conversely every induced function with that law inflates to P\_j and extends uniquely to QwP\_j with the Q-equivariance of(T11), and by zero elsewhere. These constructions are inverse and \(L_{P_j}\)-equivariant. They give the actual subspace(T14).

This is the Mackey component whose intersection-size matrix has entries k in the first-first position, one from each new Q block in the first P block, t in its corresponding determinant block, and zero elsewhere. No new free trailing block is used.

\subsection{4. Generic and degenerate generic vectors in the first block}\label{generic-and-degenerate-generic-vectors-in-the-first-block}

We use the following elementary finite induction observation. Let \(R'=P_{(a_1,\ldots,a_s)}\subseteq\operatorname{GL}_r(\kappa)\), and suppose each representation \(\sigma_i\) of \(\operatorname{GL}_{a_i}(\kappa)\) has a vector for the chosen generic upper-unitriangular character. Then \(\operatorname{Ind}_{R'}^{\operatorname{GL}_r}
(\boxtimes_i\sigma_i)\) also has such a vector.

To prove it, take the permutation which reverses the order of the blocks while preserving the internal order in each block. Its conjugate parabolic meets \(U_r\) in the product of the within-block upper-unitriangular groups, with the blocks in reversed order. Choose the tensor of the given generic vectors. Define an induced function on the corresponding \(R'wU_r\) double coset by right multiplication with the desired generic character, and zero elsewhere. The transformation law on the stabilizer agrees with the tensor vector by construction, so the function is well-defined and nonzero. It is a generic eigenvector under right \(U_r\). This construction uses no averaging denominator. It works with \(\psi^{-1}\) just as with \(\psi\).

The finite Steinberg representation is regular on \(U_k\) (Steinberg1957, Theorem2(i), printed p.349), so it contains the required \(\psi^{-1}\)-eigenvector. Applying the preceding observation to(T14) proves the fully generic case \(j=m-1\).

For \(1\le j<m-1\), the source character omits the root at k. Transitivity of finite induction rewrites \(\mathcal V_j\) as \[
 \operatorname{Ind}_{P_{(k,j)}}^{\operatorname{GL}_{k+j}}
  \left(\operatorname{St}_k\boxtimes
      \operatorname{Ind}_{B_j}^{\operatorname{GL}_j}
                (\bar\eta_0\boxtimes\cdots\boxtimes\bar\eta_{j-1})\right).
 \tag{T15}
\] The subspace supported on \(P_{(k,j)}\) and fixed by its radical contains exactly the displayed Levi representation. Its first factor has the generic \(\psi^{-1}\) vector; the second has the generic \(\psi^{-1}\) vector by the preceding open-orbit construction with all block sizes one. Their tensor is therefore an eigenvector for the character of \(U_{k+j}\) which is generic inside the k and j blocks and trivial on the joining radical. That is precisely \(\psi_{S_j}^{-1}\). Together with(T14), this proves \[
                         \mathcal M_j^{H_j,\lambda_j}\ne0.
 \tag{T16}
\] For \(j=0\), the same assertion is simply the generic Steinberg vector. The proof treats both source branches exactly.

\subsection{5. Passage to the actual p-adic representation}\label{passage-to-the-actual-p-adic-representation}

By Iwasawa decomposition, the compact function model of(T7) restricts to induction from \(Q(F)\cap K_{n_j}\). Taking \(K_{n_j}^1\)-invariants gives the exact finite induced model \[
 \Pi_j^{K_{n_j}^1}\simeq
 \operatorname{Ind}_{Q(\kappa)}^{\operatorname{GL}_{n_j}(\kappa)}
       \left(\pi_0^{K_k^1}\boxtimes
                  \mathop{\boxtimes}\displaylimits_{i=0}^{j-1}\bar\eta_i\circ\det_d\right).
 \tag{T17}
\] For clarity, right \(K_{n_j}^1\)-invariance means that a function on K factors through its finite quotient. Its value is fixed by \((Q\cap K_{n_j}^1)\), whose Levi projection is \(K_k^1\times(K_d^1)^j\). These are exactly the coefficient invariants in(T17). Conversely such a finite induced function inflates to K and extends by Iwasawa; modulus characters are one on the compact intersection. This proves the isomorphism rather than merely an inclusion.

By(T6), the finite module(T11) embeds in(T17). Equation(T16) proves(T8), completing the theorem.

\subsection{6. A primitive finite vector and its exact norm}\label{a-primitive-finite-vector-and-its-exact-norm}

There is a useful quantitative statement before any global automorphic integration. Let R be a ring of integers enlarged to contain the values of \(\psi\) and the \(\bar\eta_i\). Use the Steinberg lattice \[
        \operatorname{St}_{k,R}
          =\bigoplus_{u\in U_k}R\,u e
          \ \subset R[\operatorname{GL}_k(\kappa)/B_k].
 \tag{T18}
\] The basis and its integral realization are Steinberg's Theorem1; the same statement over arbitrary coefficient rings is recorded in Geck(2015), Theorem2.1. Give the flag permutation module its Hermitian metric with flags orthonormal.

The generic vector \(v_\psi=\sum_{u\in U_k}\psi_U(u)\,u e\), with the sign chosen to give right type \(\psi_U^{-1}\), is primitive in(T18). Its coefficients in that basis are roots of unity. Moreover \[
                     \|v_\psi\|^2=|U_k|=q^{k(k-1)/2}.
 \tag{T19}
\] To see the norm directly, expand e as the alternating sum of Weyl flags. For every Weyl element except the longest, its \(U_k\)-stabilizer contains a simple-root subgroup on which \(\psi_U\) is nontrivial, so the sum cancels. Only the longest Bruhat cell remains; its \(|U_k|\) flags are distinct and have coefficients of absolute value one. This proves(T19). It does not assert that the whole basis \(\{u e\}\) is orthonormal.

Give(T11) its normalized compact-induction metric \[
                    \|f\|^2=\frac1{|G|}\sum_{g\in G}\|f(g)\|^2.
\] The constructions in sections3--4 can be made over R with no division. At one selected coset their value is \(v_\psi\); at every nonzero coset it has the same norm, because the Steinberg action is unitary and all character values have modulus1. Hence they produce a primitive vector \(v_j\) in the natural R-valued induced lattice with exact type \(\lambda_j\).

Its support size in \(Q\backslash G\) is \[
 q^{\,d j(j-1)/2+\epsilon_j k j},\qquad
 \epsilon_j=
 \begin{cases}
  1,&j=m-1,\\
  0,&0\le j<m-1.
 \end{cases}
 \tag{T20}
\] For \(j=0\) this means one coefficient. For \(j\ge1\), the fiber of \(H\backslash P_j\to R\backslash L_{P_j}\) has \[
        [N_{P_j}:H\cap N_{P_j}]=q^{t j(j-1)/2}.
\] Indeed precisely the entries from \(L_\ell\) to \(T_i\) with \(\ell>i\) are allowed in the P radical and forbidden by Q. There are \(t\,j(j-1)/2\) such entries; all others are unrestricted.

In the fully generic first-block construction the open U orbit has size \[
 |U_{k+j}|/|U_k|=q^{kj+j(j-1)/2}.
\] In the omitted-root construction(T15), the support lies in \(P_{(k,j)}\) and its flag orbit is the longest cell for \(\operatorname{GL}_j\), of size \(q^{j(j-1)/2}\). Combining these counts proves(T20).

Write the Gaussian multinomial index as \[
 \mathcal J_j(q)=[G:Q]
 =\frac{\prod_{s=1}^{k+jd}(q^s-1)}
        {\left(\prod_{s=1}^k(q^s-1)\right)
         \left(\prod_{s=1}^d(q^s-1)\right)^j}.
 \tag{T21}
\] Every left Q coset has \(|Q|\) elements, so(T19)--(T21) give \[
 \boxed{\|v_j\|^2=
       \frac{q^{\,k(k-1)/2+d j(j-1)/2+\epsilon_j k j}}
                       {\mathcal J_j(q)}.}
 \tag{T22}
\] This is a norm of an explicit primitive finite vector, before mapping into a chosen automorphic image lattice. The embedding(T6) may be scaled to be an isometry on its finite Steinberg copy, in which case the same norm holds in the compact model(T17). Such an isometric scaling does not by itself identify the arithmetic lattice inherited from a global coefficient row.

For \(11\nmid q\), let \(h=\operatorname{ord}_{11}(q)\) and \(c=v_{11}(q^h-1)\). Put \[
             D_q(n)=c\lfloor n/h\rfloor
                         +v_{11}(\lfloor n/h\rfloor!).
\] The elementary lifting-the-exponent formula gives \[
 v_{11}\mathcal J_j(q)
       =D_q(k+jd)-D_q(k)-jD_q(d),\qquad
 v_{11}(\|v_j\|^2)=-v_{11}\mathcal J_j(q).
 \tag{T23}
\] Indeed only s divisible by h contribute to \(\prod_{s=1}^n(q^s-1)\), and their valuations are \(c+v_{11}(s/h)\). This displays an actual finite normalization factor. It is not an assertion that the global theta row has that content or that this factor is a DP44 gain.

\subsection{7. Global use and remaining gates}\label{global-use-and-remaining-gates}

For this section assume the rank gaps, selected real types and angular weight inequalities of LIFT14, in addition to the local hypotheses above. These global hypotheses are stronger than the bare local inequalities \(k\ge1,t\ge1\).

The local theorem applies to any nonzero global first lift whose initial local representations at all G places contain the finite Steinberg type, and whose common stage characters have the literal unit restrictions(T3). Such a first input is constructed on the old degree12 host by the finite-type extension of AM-DO-030: use a full-rank residue-matrix orbit to place the actual Steinberg compact type in the auxiliary finite Schwartz vector, including the determinant character in(T10); simultaneous multijet interpolation supplies a permitted29 half-translate with nonzero selected kernel coefficient; compact Schur orthogonality then extracts an actual nonzero first lift. The exact extraction and normalizations belong to AM-DO-033, rather than being inferred from finite type occurrence alone.

For later edges the analytic induction of LIFT14 carries over after enlarging its omitted finite set S to include the G primes. All nonsplit local components and their fixed spherical tests are unchanged. Every G place is split, so Mínguez supplies local occurrence there. Outside S the same standard-parameter recurrence holds, and the partial Rallis factor remains \[
 L_{E}^{S}\!\left(\frac{d+1}{2},
             \operatorname{BC}(\pi_0)\mu_j^{-1}\right)
       \prod_{r<j}\prod_{\ell=1}^{d}L_E^{S}(\ell,\mu_r\mu_j^{-1}).
 \tag{T24}
\] Here BC denotes the unramified standard parameters, as in LIFT14, not a new global base-change theorem. The initial factor converges absolutely by the same crude unitary bound and source rank gap. Each ratio has nontrivial angular infinity type, so the \(\ell=1\) factors are finite and nonzero by the checked boundary argument in LIFT14; all other factors are in their absolutely convergent range.

Thus the same Rallis proof gives nonzero later lifts with unrestricted finite Schwartz vectors, with the selected real Schur component retained. Its joint-kernel proof continues to apply to each actual automorphic multiplicity copy. At the G places every constituent has the local form(T7), for its own initial ancestor, and hence contains(T8).

One can consequently choose nonzero stage vectors of the exact nonterminal finite column type. To justify this selection, take the product of the compact character projectors at the finitely many G places. If it killed every finite translate of a nonzero cyclic automorphic submodule, that submodule could have no constituent containing the corresponding local types, contrary to(T8). Therefore some finite translate followed by the projector is nonzero. These operations commute with the selected real type and are absorbed into the preceding finite Schwartz input by Weil equivariance. The output remains in the actual generated previous theta image.

This assertion does not impose a predetermined raw PhiG, PhiGS or Xi vector. A local type can occur while a particular local intertwiner vanishes on the previous image; an integral projector can also change content. Nor does it establish the terminal HXm condition, the odd etaT pin on the old host, the full arithmetic D7 constructor, or the six comparison identities. The degree84 host needs its own global even stage characters and first finite-type extraction hypotheses before this conditional transfer can be invoked there.

The new local primitive vector(T22) and the first extraction normalizations do not compute the later global image lattice. The missing quantitative object remains the class-integrated selected theta matrix restricted to that prior image, with the source's fixed row normalization. Its saturation, coefficient denominators, levels and archimedean size determine the actual \(g_j,e_j,H\) required by DP44.

\subsection{8. Finite diagnostic and scope of verification}\label{finite-diagnostic-and-scope-of-verification}

The exact certificate \nolinkurl{ramified_column_type_certificate.py} checks the first nonterminal model \(k=j=t=1\) over both \(\mathbf F_3\) and \(\mathbf F_5\), with a nontrivial quadratic determinant character. It verifies every matrix-line covariance identity for the constructed vector, its support and the norm in(T22), and the character-projector multiplicity. The generic branch has multiplicity1; the omitted-root branch has multiplicity2. The latter corrected an initial auxiliary test expectation of multiplicity1. No uniqueness claim is used or made in the theorem.

The same certificate verifies the forbidden-entry root count in 100 small block configurations and(T23) against exact integer Gaussian multinomials in1344 cases. These finite checks support the sign and normalization; the general local and global proofs are the arguments above. No Lean compilation, raw-Xi realization or global integral saturation check is claimed.

\subsection{Primary sources and exact source locators}\label{primary-sources-and-exact-source-locators}

\begin{itemize}
\tightlist
\item
  Casselman, W. (1980). The unramified principal series of p-adic groups. I. The spherical function. Compositio Mathematica, 40(3), 387--406. Propositions2.4 and2.6, printed pp.395--396. \url{https://numdam.org/item/CM_1980__40_3_387_0.pdf}.
\item
  Geck, M. (2015). On the l-modular composition factors of the Steinberg representation. arXiv:1504.04157v2. Theorem2.1, printed p.3, and Corollary2.3, printed pp.4--5. This records the integral Steinberg basis and proves the Borel-fixed line. \url{https://arxiv.org/pdf/1504.04157}.
\item
  Mínguez, A. (2008). Correspondance de Howe explicite: paires duales de type II. Annales Scientifiques de l'École Normale Supérieure, 41(5), 717--741. Theorem1, printed p.718, and equation(1.4), pp.721--722. \url{https://www.numdam.org/item/ASENS_2008_4_41_5_717_0.pdf}.
\item
  Steinberg, R. (1957). Prime power representations of finite linear groups II. Canadian Journal of Mathematics, 9, 347--351. Theorems1--2 and the integral-basis corollary, pp.348--350. \url{https://doi.org/10.4153/CJM-1957-041-1}. Original PDF: \url{https://www.cambridge.org/core/services/aop-cambridge-core/content/view/1A82E966D7A0832A4FD80689E5F67CAE/S0008414X00044709a.pdf/prime-power-representations-of-finite-linear-groups-ii.pdf}.
\item
  Zelevinsky, A. V. (1980). Induced representations of reductive p-adic groups. II. On irreducible representations of GL(n). Annales Scientifiques de l'École Normale Supérieure, 13(2), 165--210. Proposition8.5, p.196; Proposition8.6, p.197. \url{https://www.numdam.org/article/ASENS_1980_4_13_2_165_0.pdf}.
\end{itemize}

Pinned source root: work/research/alpoge\_source, commit 7cbf7a0c32fca1badf05931879a37c3f8a2b03b4.

\begin{itemize}
\tightlist
\item
  \nolinkurl{Formalism/LocalToGlobal/Columns/RecLevel.lean50–52:} exact inverse adelic type.
\item
  \nolinkurl{Formalism/LocalToGlobal/Columns/GenericUnipotentCharacter.lean437–452:} HX, SX, lamX, including omitted-root branch.
\item
  \nolinkurl{Formalism/LocalToGlobal/Columns/ColForms.lean92–105:} literal colLev and ColV.
\item
  \nolinkurl{Formalism/LocalToGlobal/Columns/AbelianHostFieldConditions.lean273–306,315–323,488:} chi\_om, distinctness, IsMu0 and g\_big.
\item
  AM-DO-028: actual Hom/bilinear splitting dictionary.
\item
  AM-DO-030/033: fixed auxiliary multijet and finite-type first extraction.
\item
  ALP-LIFT-14: later selected global analytic induction, including its every-vector and multiplicity scope.
\end{itemize}

\clearpage

\section{ALP-LIFT-20: the actual raw first G factor and its exact Steinberg component}\label{alp-lift-20-the-actual-raw-first-g-factor-and-its-exact-steinberg-component}

Author: GPT-6 Astra Pro, ALP\_MACHINE \nolinkurl{/root/lift_accounting}. Date: 21--22 November 2026 UTC. Inherited reasoning effort is not exposed.

This result concerns the actual source functions \nolinkurl{PhiG}, \nolinkurl{PhiGS}, \nolinkurl{PsiAt}, and \nolinkurl{XiCpl1} when \nolinkurl{m=2}. It identifies a particular allowed first coupled G factor exactly, rather than replacing that factor by a freely chosen Schwartz function. Its Steinberg component is nonzero, with a fully computed scalar and norm. Compact integration against the initial Steinberg input extracts this component while leaving the prescribed raw factor in the kernel.

The local theorem is independent of the global CM field and holds at every degree-one relative split G place satisfying the source coordinate conventions. The global corollary uses a separately established actual mixed theta family, a selected polynomial, and a valid common stage character. It does not instantiate the whole \nolinkurl{ChIdx.Ok}, identify other raw T/Q factors, handle \nolinkurl{PhiTop}, or prove a later restricted-image statement.

\subsection{1. Exact source specialization}\label{exact-source-specialization}

All source locations below refer to the recovered pinned commit \nolinkurl{7cbf7a0c32fca1badf05931879a37c3f8a2b03b4}. The principal file is \nolinkurl{Formalism/LocalToGlobal/Columns/FiniteSumNonzeroLargeDepth.lean}, abbreviated FS. No Lean execution is claimed.

Take \nolinkurl{m=2}, put \nolinkurl{a=k}, \nolinkurl{b=k+t+1}, and specialize the first coupled kernel to \nolinkurl{m2=0}. The actual rank definitions give

\[
 r=a+1,\quad r_S=0,\quad N_0=a,\quad N_1=b,
 \quad c_0=1,\quad c_1=0.
 \tag{RG1}
\]

Here \nolinkurl{cOf R j = R.m-1-j} (FS815), and \nolinkurl{rS=m/2-1} is in \nolinkurl{Formalism/LocalRingCardinality/Params.lean}. The formula is valid for every admissible source rank; no small-rank numerical test is being substituted for those rank constraints.

Fix a degree-one G place with uniformizer the rational prime \nolinkurl{q}, residue field \nolinkurl{kappa=F_q}, and integral coordinate lattice. Write \nolinkurl{O} for the local valuation ring. Fix the exact additive character \nolinkurl{psi_0} in the source, and let

\[
 G_q=\operatorname{GL}_a(\kappa),\qquad
 U_q=\{u\in G_q:u\text{ is upper unitriangular}\},
 \qquad
 \psi_U(u)=\psi_0\!\left(\sum_{i=0}^{a-2}u_{i,i+1}\right).
 \tag{RG2}
\]

The character is the actual initial nondegenerate column character. Its inverse is the adelic source-input character: \nolinkurl{Columns/RecLevel.lean}50--52 defines \nolinkurl{levTyp=levChr^(-1)}.

In \nolinkurl{PsiAt} the second frame is literally one (FS984--986). The first frame is also one at this stage: \nolinkurl{ChIdx.Ok.zero_h} states \nolinkurl{L.h 0 b P=1} for every initial index (FS1019--1023). Thus this specialization does not impose a new choice of initial frame beyond the source's own condition. In arbitrary global lattice frames, the following calculation is made after applying the exact \nolinkurl{inFr} and target inverse frame in \nolinkurl{XiCpl1}.

Choose the allowed target residue class

\[
                       y\equiv q^{-1}e_a\pmod{O^b}.
 \tag{RG3}
\]

The subscript is zero-based: \nolinkurl{e_a} is the first coordinate immediately after the first \nolinkurl{a} target coordinates. This class belongs to \nolinkurl{q^(-1)O^b/O^b}, exactly the G component in \nolinkurl{ybSet} (FS1125--1133). The assertion is independent of the representative chosen by \nolinkurl{repY}.

Let \nolinkurl{X} be a \nolinkurl{b by a} matrix. The definition \nolinkurl{joinCol X y} is the \nolinkurl{b by (a+1)} matrix \nolinkurl{[X,y]} here (FS981--983). Both index permutations \nolinkurl{fIdx} in \nolinkurl{entB} are the identity in (RG1). The source last-column scaling exponent is one, all its other exponents are zero, and every target exponent is zero. Hence

\[
             \operatorname{entB}(\operatorname{joinCol}(X,y))
                           =[X,qy].
 \tag{RG4}
\]

This follows directly from FS684--694; it is not a change of Schwartz normalization. Since both group arguments of \nolinkurl{omPhi} are one, all determinant, valuation, and splitting factors in \nolinkurl{Loc.om} are one (FS553--561, 940--945).

\subsubsection{The raw-factor identity}\label{the-raw-factor-identity}

Put \nolinkurl{X_0=[I_a;0]} and define

\[
 f_U(X)=
 \begin{cases}
   \psi_U(u),&X\in M_{b,a}(O),\quad\bar X=X_0u,\ u\in U_q,\\
   0,&\text{otherwise}.
 \end{cases}
 \tag{RG5}
\]

Then, for either value of the source \nolinkurl{sharp} flag,

\[
 \boxed{\ \operatorname{PsiAt}(0,P,1,X,y)=f_U(X)\ }
                  \qquad\text{whenever (RG3) holds.}
 \tag{RG6}
\]

Here is the complete predicate check. In both \nolinkurl{PhiG} and \nolinkurl{PhiGS}, the internal variable called \nolinkurl{a} is now \nolinkurl{r=a+1}, since the stage parameter is \nolinkurl{j=1}. The first integrality condition requires all entries in \nolinkurl{[X,qy]} to be integral. The condition on rows at least \nolinkurl{a+1} makes their first \nolinkurl{a+1} entries zero modulo \nolinkurl{q}. In the first \nolinkurl{a+1} rows and columns, the first \nolinkurl{a+1} diagonal entries are one and the entries below the diagonal are zero. The last column is already \nolinkurl{e_a} modulo \nolinkurl{q} by (RG3). Consequently the first \nolinkurl{a} columns have upper unitriangular top \nolinkurl{a by a} block and every subsequent row is zero. This is precisely the support in (RG5).

For \nolinkurl{PhiGS}, \nolinkurl{symIdx a 0} is the identity (\nolinkurl{Columns/LevGSym.lean}11--13). Its extra zero condition in column \nolinkurl{a} is already imposed by \nolinkurl{qy=e_a} modulo \nolinkurl{q}. Its other off-diagonal conditions are exactly the same lower zeros just used. There are no determinant-block conditions, since their column range starts at \nolinkurl{r} and ends strictly before \nolinkurl{r}.

On the support, the leading \nolinkurl{(a+1) by (a+1)} matrix is \nolinkurl{diag(u,1)} modulo \nolinkurl{q}, so its determinant is one. The inverse \nolinkurl{chi_0} determinant factor therefore equals one, even though the common splitting character itself can be ramified. In \nolinkurl{lamGvS}, the second additive sum has length \nolinkurl{r-a-1=0}. The first sum is exactly the exponent in (RG2). In \nolinkurl{lamGv}, the only additional simple-root entry is the entry in row \nolinkurl{a-1}, column \nolinkurl{a}; it is zero by (RG3). Thus both functions give \nolinkurl{psi_U(u)}. Off the stated support one of the displayed source predicates fails. This proves (RG6) using the literal definitions FS780--810 and FS927--931.

Adding an integral vector to \nolinkurl{y} changes \nolinkurl{qy} by a multiple of \nolinkurl{q}, so the same proof also establishes representative independence directly; no unproved \nolinkurl{ChIdx.Ok.psiInv} is used to obtain (RG6).

At finitely many G places the classes (RG3) can be imposed simultaneously, including after the actual target frames. Indeed the additive Chinese remainder theorem for the localized fractional lattice gives an element of the global target vector space with those finitely many prescribed classes and with the required integral classes at the other finite places in \nolinkurl{rowSG}. Its class is an element of the source \nolinkurl{ybSet}; every representative of that class has the same local factors. This is only an existence statement for that allowed class. No archimedean height bound for a chosen rational representative is asserted.

\subsection{2. Finite Fourier coefficients and exact raw norm}\label{finite-fourier-coefficients-and-exact-raw-norm}

Give \texttt{M\_\{b,a\}(O)} additive volume one, as in the actual split Schrödinger model. Each residue cell has volume \nolinkurl{q^(-ab)}. Set

\[
            n_U=a(a-1)/2,\qquad d_q=q^{n_U}=|U_q|.
 \tag{RG7}
\]

Every nonzero value of \nolinkurl{f_U} is a root of unity, its value at \nolinkurl{X_0} is one, and

\[
                         \|f_U\|^2=q^{-ab}d_q.
 \tag{RG8}
\]

This is its exact Schwartz norm, not a normalized finite group-counting norm. Its finite value lattice is primitive at every prime at which the character values are units. In particular it is primitive at11 when \nolinkurl{q} is prime to11.

There is an especially simple actual Heisenberg expansion. Let \texttt{V=M\_\{b,a\}(kappa)} and let \nolinkurl{W} be the linear subspace of matrices whose only possible nonzero entries are in the strictly upper triangular part of their top \nolinkurl{a by a} block. Then \nolinkurl{dim W=n_U}, the support of \nolinkurl{f_U} is the affine space \nolinkurl{X_0+W}, and its character there is \nolinkurl{ell(w)=sum_i w_(i,i+1)}. Use the perfect finite pairing \nolinkurl{B} and actual Fourier phase of DO29.11--13; identify \nolinkurl{ell} with the corresponding character of \nolinkurl{W}. For the actual finite vacuum \nolinkurl{phi_0} one has

\[
 f_U=\sum_y b_y W_y\phi_0,\qquad
 b_y=
 \begin{cases}
 q^{n_U-ab}\psi_0(-B(X_0,y)),&B(w,y)=\ell(w)\ \forall w\in W,\\
 0,&\text{otherwise}.
 \end{cases}
 \tag{RG9}
\]

If the finite pairing is expressed as characters rather than scalar-valued forms, the condition means equality of the two characters on \nolinkurl{W}; this is invariant under the coordinate choice. To prove the formula, insert \nolinkurl{X=X_0+w} into the finite Fourier formula \nolinkurl{q^(-ab) sum_X f_U(X) psi_0(-B(X,y))}. The sum over \nolinkurl{W} is \nolinkurl{q^(n_U)} when its character is trivial and zero otherwise. Perfectness of the pairing gives exactly \nolinkurl{q^(ab-n_U)} nonzero coefficients. Therefore

\[
 \#\{y:b_y\ne0\}=q^{ab-n_U},\qquad
 \sum_y|b_y|=1,\qquad
 \sum_y|b_y|^2=q^{n_U-ab}.
 \tag{RG10}
\]

These are the actual phase-correct Heisenberg operators, not just abstract coordinate translations. They are unitary, so the operator in (RG9) has norm at most one. Its exact value-lattice and Fourier-coefficient lattices should not be confused: the latter have explicit \nolinkurl{q}-power denominators \nolinkurl{q^(ab-n_U)}. These denominators are units at11 but must be retained at the auxiliary prime and in any global coefficient-height accounting.

\subsection{3. Exact Steinberg projection in the actual split action}\label{exact-steinberg-projection-in-the-actual-split-action}

Let \nolinkurl{eta} be the local selected-half common splitting character. In the literal source character pin, \nolinkurl{bar eta=chi_0^(-1)}. The genuine left source action is

\[
 \omega(g)f(X)=\eta(\det g)^{-1}
                    |\det g|^{b/2}f(Xg).
 \tag{RG11}
\]

This is the corrected DO33.5/LIFT16(S2) convention. On the integral compact it factors on the present functions through \nolinkurl{G_q} with its displayed determinant correction. In particular \nolinkurl{f_U} has right-source \nolinkurl{psi_U} type under \nolinkurl{U_q}; the determinant factor is trivial there.

Use the actual finite Steinberg representation \nolinkurl{St_q} and the unit generic vector \nolinkurl{u_0} constructed in DO33.2--4. Its dimension is \nolinkurl{d_q}; the \nolinkurl{psi_U} eigenspace is one dimensional. These facts are the published finite Steinberg theorems (Steinberg, 1957, Theorems1--2; Geck, 2015, Theorem2.1) and the already checked normalization DO033. With the Hermitian form conjugate-linear in its first argument put

\[
 m(g)=\langle u_0,\operatorname{St}_q(g)u_0\rangle,
 \qquad
 f_{\rm St}(X_0g)=\eta(\det g)m(g),
 \tag{RG12}
\]

and extend \nolinkurl{f_St} to residue cells, zero off the full-rank orbit \nolinkurl{X_0G_q}. This is exactly DO33.6 with \nolinkurl{v'=u_0}. It has pure source \nolinkurl{St_q} type, value one at \nolinkurl{X_0}, and

\[
                 \|f_{\rm St}\|^2=
                          q^{-ab}|G_q|/d_q.
 \tag{RG13}
\]

Let \nolinkurl{P_St} be the orthogonal projection onto the full \nolinkurl{St_q} isotypic subspace for the actual action (RG11), including all its multiplicity copies in the residue Schwartz space. Then

\[
 \boxed{\ P_{\rm St}f_U=\alpha_q f_{\rm St},\qquad
                  \alpha_q=\frac{d_q^2}{|G_q|}>0.\ }
 \tag{RG14}
\]

Here is a proof that controls the whole isotypic projection, not only its nonzero inner product. On the free orbit \nolinkurl{X_0G_q}, tensoring functions by \nolinkurl{eta(det g)} identifies the actual right-source action with the ordinary right regular representation. It leaves \nolinkurl{f_U} unchanged, since every point in its support has determinant one. The Steinberg matrix coefficients span the full Steinberg isotypic space in this regular representation. For any \nolinkurl{v,w in St_q}, their inner products with \nolinkurl{f_U} are

\[
 \begin{split}
 &q^{-ab}\sum_{u\in U_q}
       \overline{\langle w,\operatorname{St}_q(u)v\rangle}
                                         \psi_U(u)\\
 &\hspace{12mm}=q^{-ab}d_q\,
       \overline{\langle w,u_0\rangle\langle u_0,v\rangle}.
 \end{split}
 \tag{RG15}
\]

The equality is the finite \nolinkurl{psi_U} projector and the one-dimensionality of its image. Thus the projection is a multiple of the single coefficient \nolinkurl{m}, including in the multiplicity direction. By finite Schur orthogonality the squared norm of \nolinkurl{m} on \nolinkurl{G_q} is \nolinkurl{|G_q|/d_q}. Its inner product with \nolinkurl{f_U} is \nolinkurl{d_q}. The scalar is therefore exactly \nolinkurl{d_q^2/|G_q|}, proving (RG14). All other residue orbits are invariant under the source compact and have zero restriction of \nolinkurl{f_U}, so they do not add any component to this projection.

Consequently

\[
 \|P_{\rm St}f_U\|^2
       =q^{-ab}\frac{d_q^3}{|G_q|},\qquad
 \frac{\|P_{\rm St}f_U\|^2}{\|f_U\|^2}
       =\frac{d_q^2}{|G_q|}=\alpha_q.
 \tag{RG16}
\]

The same calculation gives the more precise compact filter actually used in a theta lift:

\[
 \boxed{\ \int_{\operatorname{GL}_a(O)}
       \overline{m(\bar g)}\,\omega(g)f_U\,dg
              =\frac{\alpha_q}{d_q}f_{\rm St}.\ }
 \tag{RG17}
\]

The compact has measure one. Finite Schur orthogonality makes this filter \nolinkurl{d_q^(-1)} times the projection onto \nolinkurl{u_0} on each Steinberg copy and zero on other irreducible types. Equation(RG14) then proves (RG17). This also verifies the sign: the kernel has \nolinkurl{psi_U} covariance and the source matrix coefficient is its complex conjugate, with covariance \nolinkurl{psi_U^(-1)} as required by the literal bilinear lift.

There is no claim that \nolinkurl{P_St} is an integral projector. For a fixed place above11 with \nolinkurl{q} prime to11,

\[
 v_{11}(\alpha_q)=-v_{11}(|G_q|)
             =-\sum_{i=1}^{a}v_{11}(q^i-1).
 \tag{RG18}
\]

The explicit denominators are part of the compact averaging cost. The raw factor (RG5) itself remains primitive.

\subsection{4. A first global lift retaining the raw G factors}\label{a-first-global-lift-retaining-the-raw-g-factors}

This section is a corollary of the local theorem and the already established actual-family/jet/extraction inputs, not an assertion of new global theta-family existence. Use DO033 on the actual \nolinkurl{J^2} host, or its separately verified degree-seven host version DO034 once that version and its rank/character hypotheses are in place. Keep a selected primitive polynomial \nolinkurl{P}, its actual period scalar \nolinkurl{mathcal P}, and its content upper bound \nolinkurl{B_P}. In DO033, \nolinkurl{mathcal P=C^(2ab)}; on a different host its own period formula must be used. Fix any further nonzero auxiliary finite factor for which the actual simultaneous-jet theorem has been established. It is retained in all its norms and coefficient bounds. For clarity, the displayed formulas below take that additional factor to be the normalized vacuum; its known scalar factors can be multiplied back.

At each selected G place take the exact raw factor (RG5), arising from the allowed class (RG3). Write

\[
 d_G=\prod_vd_{q_v},\quad
 A_G=\prod_v\alpha_{q_v},\quad
 \Gamma_U=\prod_vq_v^{-ab}d_{q_v},\quad
 \Gamma_{\rm St}=\prod_vq_v^{-ab}|G_{q_v}|/d_{q_v}.
 \tag{RG19}
\]

Then \nolinkurl{Gamma_U=A_G Gamma_St}. The exact raw G Schwartz norm is \nolinkurl{sqrt(Gamma_U)}, and its tensor Heisenberg expansion has coefficient absolute sum one and \nolinkurl{prod_v q_v^(ab-n_U)} nonzero terms. The Steinberg projection is \nolinkurl{A_G} times the tensor of the DO033 vectors.

Apply the actual simultaneous-jet theorem to that tensor of the nonzero \nolinkurl{f_St} vectors, not merely to an abstract finite representation. With its Fourier support size \nolinkurl{s_St}, one can take a single permitted29 half-translate at

\[
                   M=29^e\ge3s_{\rm St}(D+1)-1,
 \qquad s_{\rm St}\le\prod_vq_v^{ab}.
 \tag{RG20}
\]

It gives an actual selected coefficient \nolinkurl{a_lambda} with \nolinkurl{a_lambda != 0} and \nolinkurl{v_11(a_lambda)<=B_P}, after the same period removal and coordinate hypotheses as in DO033/034. The29 operator and the real selected differential commute with the G-compact action because they act at other places. Thus the same individual permitted translate works before the compact filtering of the raw factor; no replacement of \nolinkurl{f_U} by \nolinkurl{f_St} in the prescribed kernel is necessary.

Let \nolinkurl{I=vol(K*)^(-1)} be the exact source compact index for the selected29 level, with full integral compacts at G, and let \nolinkurl{d_rho} be the selected real source dimension. The identity class has trivial rational stabilizer at the levels in DO033/034. Use their actual class-supported source function \nolinkurl{F=conj(m_infty product_v m_v)}, zero on other classes. It has initial literal \nolinkurl{psi_U^(-1)} type at every G place. Successive application of (RG17) and real Schur orthogonality gives the exact raw-kernel value

\[
 \boxed{\ \Theta_{\rm raw}^{\rm bil}(F)(1)
       =\frac{A_G\,\mathcal P\,a_\lambda}{I d_\rho d_G}
       \ne0,\qquad
 \|F\|_2^2=\frac1{I d_\rho d_G},\quad\|F\|_\infty=1.\ }
 \tag{RG21}
\]

The proof performs the compact integral on the raw theta kernel, replacing its G factor by the right-hand side of (RG17) inside that integral. Absolute convergence of the definite theta integral and the finite compact average justify this operation. It is exactly the source's bilinear convention, not a silently conjugated-input theta map.

After making the raw finite input and source input unit norm, the value becomes

\[
 \Theta_{\rm raw,unit}^{\rm bil}(F_{\rm unit})(1)
       =\frac{A_G\,\mathcal P\,a_\lambda}
                    {\sqrt{\Gamma_U}\sqrt{I d_\rho d_G}}.
 \tag{RG22}
\]

If the real Schwartz vector is also made unit, divide by its exact norm from R011/DO034. Every factor is displayed; in particular the Steinberg averaging factor is not omitted. For \nolinkurl{q_v} prime to11,

\[
 \begin{split}
 v_{11}(\Gamma_U)&=v_{11}(d_G)=0,\\
 v_{11}(A_G)&=-\sum_v\sum_{i=1}^{a}v_{11}(q_v^i-1),\\
 v_{11}\!\left(\mathcal P^{-1}
                 \Theta_{\rm raw}^{\rm bil}(F)(1)\right)
   &\le B_P-v_{11}(I d_\rho)
                  -\sum_v\sum_{i=1}^{a}v_{11}(q_v^i-1).
 \end{split}
 \tag{RG23}
\]

This is an upper content witness for one actual first coefficient. It is neither a lower divisibility theorem for all rows nor a comparison with the full global DP44 cost. It also does not say that the original raw theta kernel has only Steinberg type.

For \nolinkurl{f} selected places above each of two completely split rational G primes \nolinkurl{g_0,g_1}, the local factors specialize to

\[
 \begin{split}
 d_G&=(g_0g_1)^{f a(a-1)/2},\\
 \Gamma_U&=(g_0g_1)^{f(-ab+a(a-1)/2)},\\
 A_G&=(g_0g_1)^{-fa}
       \prod_{\epsilon=0}^{1}\prod_{i=1}^{a}
                                     (1-g_\epsilon^{-i})^{-f}.
 \end{split}
 \tag{RG24}
\]

These identities use \nolinkurl{|GL_a(F_q)|=q^(a^2) product_i(1-q^(-i))} and are field-general local counts. The source index \nolinkurl{I} must still be the correct one for the actual global host: old \nolinkurl{J^2} and repaired \nolinkurl{J^14} have different rational29 ramification and therefore different exponents in that index.

\subsection{5. What this closes and what remains}\label{what-this-closes-and-what-remains-1}

The local equality (RG6) closes the raw initial G-factor identification for the displayed allowed coupled class when \nolinkurl{m=2}. Equations(RG14) and(RG17) prove nonzero interaction with the actual initial Steinberg input and compute its exact cost. The conditional global corollary proves that this raw G factor can be retained in the actual first lift once the cited actual-family and selected-jet hypotheses are satisfied.

The statement does not identify a complete \nolinkurl{XiCpl1} unless its other actual T/Q factors and their finite family are also supplied. It does not prove nonvanishing of \nolinkurl{XiAt} at later stages for \nolinkurl{m>2}, identify the terminal \nolinkurl{PhiTop}, construct all six \nolinkurl{AdmS} ties, or compute the full DP44 changed-data cost. Those are separate exact source gates. No isotypic occurrence is being substituted for preservation of a prescribed later arithmetic image lattice.

\subsection{6. Sources and checks}\label{sources-and-checks-1}

Accepted inputs: DO033's generic vector, actual split action, phase-correct Fourier model, simultaneous jets, and compact class extraction; DO029/030 for the actual Heisenberg and auxiliary multijet comparison; R012/R011 for the selected real polynomial and norm. Their proved scope is retained. DO034 is needed separately before the new host is used in the global corollary.

Steinberg, R. (1957). Prime power representations of finite linear groups II. \emph{Canadian Journal of Mathematics, 9}, 347--351. \url{https://doi.org/10.4153/CJM-1957-041-1} . Theorems1--2, pp.348--350: published Steinberg construction, irreducibility, and regular restriction to the upper unipotent. Read in the publisher PDF during this run.

Geck, M. (2015). \emph{On the ell-modular composition factors of the Steinberg representation}. \url{https://arxiv.org/pdf/1504.04157} . Theorem2.1 and Corollary2.3, pp.3--5: the integral flag model and characteristic-zero properties used in DO033.

Author whole-proof reading was completed on 21 November 2026, including all source predicates, the full isotypic multiplicity argument, the compact filter, the conditional global extraction, and every displayed norm/valuation factor. The exact diagnostic passed 729 residue-matrix predicate comparisons, 48 full Steinberg projections, 48 matrix-coefficient filters, and the full 729-point Fourier transform. Its raw norm is \nolinkurl{1/243}, projection scalar \nolinkurl{3/16}, and projected norm is \nolinkurl{1/1296}. The separate TSV records the check time and exact scope.

The exact diagnostic \nolinkurl{raw_first_g_steinberg_certificate.py} checks both raw predicates and the full finite Steinberg projection/filter in \nolinkurl{GL_2(F_3)}, with a nontrivial determinant splitting and all residue matrices. It also checks the explicit raw Fourier expansion and all displayed finite norms. This diagnostic is not a proof of the general rank statement or a verification of a full Lean datum.

\clearpage

\section{AM-K031 --- exact enlarged-host ranks, uniform stage characters, and the unchanged half-orbit isogeny}\label{am-k031-exact-enlarged-host-ranks-uniform-stage-characters-and-the-unchanged-half-orbit-isogeny}

Author: GPT-6 Astra Pro, ALP\_MACHINE \nolinkurl{/root/kummer_extension}, 21 November 2026 UTC; inherited reasoning effort not separately exposed.

This finding supplies concrete numerical rank/gap data, compatible split primes and actual Hecke characters on the CM host of AM-K029/030. It also proves that the permitted DP59 half-orbit survives the real-field enlargement. It is deliberately not a \nolinkurl{Dat.Ok} or D7 instance: the unmodified source still requires \(8\mid f\) and a common quadratic CM subfield for the selected type. The latter is the already audited primitive-type obstruction of AM-R-003, not a new general impossibility theorem.

Throughout the first three sections, \[
 E_0=\mathbb Q(\zeta_{28}),\quad F_0=E_0^+,\quad
 K=E_0B,\quad F=F_0B,\quad [B:\mathbb Q]=7,
\] where \(B\subset\mathbb Q(\zeta_{29})\), and the actual type contains every lift of \(\Phi_0=\{1,3,9,15,17,23\}\). Thus \([F:\mathbb Q]=f=42\), and the actual principal abelian variety is \(J^{14}\). Let \(\eta\) be the literal \nolinkurl{IsEtaT} character constructed in AM-K029.

\subsection{1. The source rank constraints and an explicit gap locus}\label{the-source-rank-constraints-and-an-explicit-gap-locus}

The exact record \nolinkurl{Formalism/LocalRingCardinality/Params.lean:6–22} imposes \[
 \begin{gathered}
 k,m\text{ even},\qquad t\text{ odd},\qquad 2\le m,
 \qquad m+3\le k,\\
 k+m+4\le t,\qquad 3m<t,\qquad
 m\le f=2\operatorname{ord}\le k,\qquad
 m^2\le k+m(t+1).
 \end{gathered}
\tag{HC1}
\] The real-place type count is tied to this f by \nolinkurl{AbelianHostFieldConditions.lean:567}. In particular the old \(k=6\) cannot remain the starting rank on the degree-84 host: \[
                         \operatorname{ord}=21,\qquad k\ge42.
\tag{HC2}
\] The inherited \nolinkurl{Frame} record itself causes no numerical obstruction to order21. Take its additive group to be \(\mathbb Z/21\mathbb Z\), its cyclotomic parameter to be zero, its projection to the trivial group \(HGrp(0)=(\mathbb Z/4\mathbb Z)^\times/\langle-1\rangle\), and its distinguished element to be zero. The projection and cyclic-kernel clauses in \nolinkurl{ZsmulGeneratorEqZeroOfProjection.lean:16–24} then hold. This is a mathematical record construction, not a claim of Lean compilation or of its compatibility with every later host record.

For the shortest allowed length \(m=2\) and smallest starting rank \(k=42\), the rank clauses alone give the least odd t as49. The numerical Hyp3 condition \(1801\le r+t\), where \(r=k+m-1\), is stronger. A convenient simultaneous choice is \[
 \boxed{f=k=42,\quad m=2,\quad r=43,\quad t=1759,
         \quad d=t+1=1760.}
\tag{HC3}
\] It gives \[
       n_V=r+t=1802,
       \qquad (N_0,N_1,N_2)=(42,1802,3562).
\tag{HC4}
\] The actual G-size clause is \nolinkurl{AbelianHostFieldConditions.lean:488}, namely \[
                 2fk<g_i-1.
\tag{HC5}
\] Here its right lower threshold is3528, not the value obtained by retaining \(k=6\).

The following uniform parameters verify all numerical parity, interlacing, centred-distance and opposite-sign gap inequalities occurring in Hyp3. Put \[
 \begin{split}
 S&=r+mt=3561,\qquad S^3=45\,156\,047\,481,\\
 C&=S^3+k+t+11=45\,156\,049\,293.
 \end{split}
\tag{HC6}
\] At every one of the 42 embeddings use zero-based indices and set \[
 \begin{aligned}
 u_i&=(2(k-i)+1)C+\tfrac12 &&(0\le i<k),\\
 v_i&=2(k-i)C &&(0\le i<k-1),\\
 b_j&=(4j-2)C,\qquad e_j=4jC &&(j=0,1),\\
 \eta_{\mathrm{num}}&=-\tfrac12.
 \end{aligned}
\tag{HC7}
\] In particular \[
 (e_0,e_1)=(0,180\,624\,197\,172),\qquad
 (b_0,b_1)=(-90\,312\,098\,586,90\,312\,098\,586).
\tag{HC8}
\]

\textbf{Verification.} The u's are half-integral, the other displayed weights are integral, and \(u_i>v_i>u_{i+1}\), while \(u_{k-1}=3C+1/2>-b_0=2C\). Both consecutive b/e gaps are \(2C\), exceeding \((t+3)/2=881\).

For the exact \nolinkurl{Centred} predicate, \(u_i+e_j\ge3C+1/2\), \(v_i+e_j\ge4C\), and \(|e_j-b_i|\ge2C\). The only cross-block term has centre difference \(4C\) and a block half-width \((t-1)/2=879\), so its absolute value is at least \(4C-879>881\). These are all four kinds of quantifiers in \nolinkurl{DataTheta.lean:550–555}.

For \nolinkurl{NCGaps}, the literal boolean in \nolinkurl{Dpar} is \(i\le j\) or \(i-j\) even. At stage0 the only false-sign entry is \(-b_1=-2C\); the other b-entry is \(2C\), and all v-entries are at least \(4C\). Thus every opposite-sign distance is at least \(4C>S^3\). At stages1 and2 every entry has true sign, so no opposite-sign comparison remains. The common odd-stage shift by \(\eta_{\mathrm{num}}\) changes no distance. This is the full finite sign analysis for \(m=2\), not a gap check only for adjacent entries.

The source's extra \nolinkurl{e_even} condition at line578 also holds: \(\sum_\tau e_j=42e_j\) is even. In the half-angular convention of \nolinkurl{HasInfQ}, the constructed odd character has infinity type \(+1/2\), so the required numerical sign is precisely \(\eta_{\mathrm{num}}=-1/2\). This checks the numerical equality; an actual \nolinkurl{ThetaX.infTy} realization remains part of the analytic adapter.

For comparison with AM-DO-025, the first selected partition at each place is \[
           \nu_{0,i}=(85-2i)C+i-921\quad(0\le i<42),
\tag{HC9}
\] whose smallest entry is \(3C-880>0\). The next partition is \[
         \nu_1=(\nu_0+\Delta,\Delta^{1760}),
         \qquad \Delta=4C-1760>0.
\tag{HC10}
\] These follow by substitution into the already checked selected-weight recurrence, not by asserting that an old six-place polynomial remains a single 42-place weight. A new selected polynomial is constructed in the new CM blocks separately.

\subsection{2. Exact conductor and split G primes, including PiGCong}\label{exact-conductor-and-split-g-primes-including-pigcong}

The abelian field K has rational conductor \[
                              N=28\cdot29=812.
\tag{HC11}
\] It lies in \(\mathbb Q(\zeta_{812})\). Conversely, a cyclotomic field containing K contains i and \(\zeta_7\), forcing respectively \(4\mid N'\) and \(7\mid N'\); ramification at29 forces \(29\mid N'\). Hence \(812\mid N'\). This verifies the minimum-divisibility formulation of \nolinkurl{IsCond} in \nolinkurl{RealCyclotomicTwoPower.lean:54}, including every conjugate embedding since K is normal over \(\mathbb Q\).

Write \(\delta=\sqrt{-7}\in E_0\), and use the following \textbf{explicit} primes and generators in the external quadratic subfield \(K_q=\mathbb Q(\delta)\): \[
 \begin{array}{c|c|c}
 i&g_i&\pi_i\\\hline
 0&4\,615\,409&1+812\delta\\
 1&5\,276\,377&813+812\delta.
 \end{array}
\tag{HC12}
\] Both rational integers are prime. The accompanying certificate checks primality by trial division through the integer square root, as well as \[
 \pi_i\bar\pi_i=g_i,\qquad
 \pi_i\equiv1\pmod{812\mathcal O_{K_q}},\qquad
                   g_i\equiv1\pmod{812}.
\tag{HC13}
\] Their conjugate ideals are coprime: a common prime would divide both the norm \(g_i\) and the sum \(2\operatorname{Re}\pi_i\), whereas the corresponding rational gcd is one. Coprimality persists after extension to K.

Thus both primes split completely in K. The support of \(\bar\pi_i\) consists of42 distinct primes, each with residue field canonically \(\mathbb F_{g_i}\) through its prime field. Equations (HC12)--(HC13) supply the norms, coprimality, cardinalities, residue norms, distinctness, full \nolinkurl{PiGCong}, and strict G-size inequalities. The value \(812\) in \nolinkurl{PiGCong} is not replaced by28. The explicitly retained exception is \nolinkurl{piG_Kq}: the source's \textbf{defined} common field \nolinkurl{Fields.Kq} is \(\mathbb Q\) for this type, rather than the external \(K_q=\mathbb Q(\sqrt{-7})\) used here. Section6 explains that literal failure.

At each i choose a faithful character \(\omega_i:\mathbb F_{g_i}^{\times}\to\mu_{g_i-1}\). For example 3 and5 are primitive roots for \(g_0\) and \(g_1\), respectively; the certificate checks their orders using the complete factorization of \(g_i-1\). Use the literal residue characters \[
             \chi_{j,P}=\omega_i^{2(j+1)}
                       \qquad(P\mid\bar\pi_i).
\tag{HC14}
\] They satisfy the equality-on-prime-field and \nolinkurl{chi_om} clauses. For \(j=0,1\) they are nontrivial and distinct. Their values lie in \[
        \mathbb Q(\zeta_L),\qquad
        L=\operatorname{lcm}(28,g_0-1,g_1-1)=14\,995\,460\,592.
\tag{HC15}
\] This is a fixed coefficient extension; its large degree is not discarded in any later height budget.

\subsection{3. Actual global stage characters for every uniform integral weight}\label{actual-global-stage-characters-for-every-uniform-integral-weight}

\subsubsection{Theorem HC16}\label{theorem-hc16}

For any integers \(e_j\), there are continuous unitary Hecke characters \(\mu_j\) of K which are trivial on \(\mathbb A_F^\times\), tame at every finite place, unramified away from \(g_0g_1\), have infinity component \[
       (\mu_j)_\infty(z)
             =\prod_{\tau\in\Psi}
                    \left(\frac{z_\tau}{|z_\tau|}\right)^{-2e_j},
\tag{HC16}
\] and satisfy \[
       \mu_{j,P}|_{\mathcal O_{K,P}^\times}=\chi_{j,P}^{-1}
                  \qquad(P\mid\bar\pi_i).
\tag{HC17}
\] These are the exact \nolinkurl{IsMu0} character conditions for a uniform splitting weight. The values required as \nolinkurl{muVal} lie in a single finite coefficient field, after adjoining the host embeddings and normal closure if required. There is no new class-number-one assumption on K or F.

\textbf{Construction and proof.} Apply the checked ALP-LIFT-18 construction on E₀ to these new split primes, with \textbf{zero} infinity weight at every embedding. Its root compatibility is then automatic. Concretely, on the compact norm-one torus its character is the product of \(\omega_i^{-2(j+1)}\) at the six selected E₀-primes. The product on global roots is one because \(N_{E_0/K_q}(\mu_{28})=1\), as proved in LIFT18(G2)--(G3). The checked class-one torus decomposition of AM-R-009 makes this an actual automorphic torus character. Composing with \(x\mapsto x/\bar x\) gives a finite-order Hecke character \(\nu_j\) of E₀, trivial on \(\mathbb A_{F_0}^\times\), with zero infinity type and the required inverse residue characters. Its values lie in the roots of unity in (HC15). This is an application of the existing character theorem; its class-number proof is not rederived here.

Pull this actual character back to K by the idele norm: \[
              \widetilde\nu_j=\nu_j\circ N_{K/E_0}.
\tag{HC18}
\] At every chosen G prime the extension is completely split, so the completion map is \(K_P=(E_0)_{P_0}=\mathbb Q_{g_i}\). On an idele supported in one such component, the local norm is the identity. Thus (HC18) has the required residue character at \textbf{each} of the seven lifts of every old selected prime; it does not raise a single local residue character to the seventh power. Norms take local units to units at every other place, so the pullback is unramified there, including the new ramification place29. The norm of an F-idele lies in \(\mathbb A_{F_0}^\times\), hence the pullback is trivial on \(\mathbb A_F^\times\).

Now use the actual odd character of AM-K029: \[
                   \boxed{\mu_j=\widetilde\nu_j\,\eta^{-2e_j}.}
\tag{HC19}
\] The restriction of η to \(\mathbb A_F^\times\) is the quadratic character \(\eta_{K/F}\). Its even power in (HC19) is therefore trivial on that group. Since η is everywhere finite-unramified, it changes no compact residue restriction or tame conductor. Its angular exponent is one, so its \(-2e_j\) power supplies exactly (HC16), which is the literal \nolinkurl{HasInfQ(-e_j)} convention. All the claimed properties follow.

For the two stages in (HC3), the finite conductor is exactly \[
          \mathfrak f(\mu_j)=g_0g_1\mathcal O_K,
          \qquad N_{K/\mathbb Q}\mathfrak f(\mu_j)=(g_0g_1)^{84}.
\tag{HC20}
\] Both halves have conductor exponent one because the selected residue characters are nontrivial and base-field triviality forces the inverse character at the conjugate place. Unramified η cannot cancel a nontrivial residue character.

To check the common algebraic field, let \(a_P=\eta(\varpi_P)(NP)^{-1/2}\), the \nolinkurl{etaVal} supplied in AM-K029. At a finite uniformizer, \[
        \mu_j(\varpi_P)
          =\widetilde\nu_j(\varpi_P)\,a_P^{-2e_j}(NP)^{-e_j}.
\tag{HC21}
\] The first factor is in the fixed finite group of roots from (HC15), the second lies in the fixed \nolinkurl{etaVal} field, and the last is rational. Thus one coefficient field works for every integer choice of the uniform weights. Taking its normal closure and adjoining K preserves finiteness. No effective small degree, discriminant or uniformizer-height bound is inferred. ∎

This is stronger than merely pulling back LIFT18's nonzero infinity weights, which would retain its sufficient \(7\mid e_j\) restriction. Here all integer uniform weights work. For example the explicit \(C\) in (HC6) is \(5\bmod7\), so \(e_1=4C\equiv6\bmod7\); it is covered by (HC19) and not by the old uniform congruence.

The construction is stable under further totally real extension: pull η back by norm, and pull \(\nu_j\) directly from E₀. At G-primes chosen completely split in the enlarged CM field, each local norm remains the identity. Every integer uniform infinity weight again works. This gives a parameterized character construction for later real-field repairs without reusing an unjustified large-field torus class-number-one assertion.

\subsection{4. The fixed alpha and the actual permitted DP59 half-orbit}\label{the-fixed-alpha-and-the-actual-permitted-dp59-half-orbit}

The alpha used in AM-K022/DO030 is in \(\mathbb Q(\sqrt{-7})\), not in \(\mathbb Q(i)\): \[
       \alpha=(1+2\delta)^e,\quad M=29^e,\quad e\ge1,
                   \qquad \alpha\bar\alpha=M.
\tag{HC22}
\] The alternative \((5+2i)^e\) has the same norm identity, but changing to it is unnecessary. We retain (HC22), which acts diagonally on every ordinary J factor.

The ideals of \(1+2\delta\) and its conjugate are coprime in every extension of E₀. A common prime would divide their sum2 and their product29, which is impossible. Thus the ideals of α and \(\bar\alpha\) remain coprime in K and in every further real extension. At rational29, their valuations in K are \(7e\) on the two disjoint six-prime halves, rather than e. This ramification does not change the element identity or the coprimality statement.

The element \(1+2\delta\) is congruent to one modulo2. It therefore acts trivially on the relevant two-torsion characteristic, and the exact odd symmetric pullback normalization carried in DP59 applies to the same invariant theta line. Under the K030 polarized product isomorphism, α is still the same diagonal endomorphism on all14 J factors: it lies in the central old E₀-action and commutes with the real trace-lattice transport.

More generally suppose a further real-field construction realizes the induced host as \(J^{14h}\) with its product line, h a positive integer. On an edge of ranks a,b put \[
       X=J^{14hab},\quad g_X=42hab,\quad
       \mathcal L=L_{2\infty}^{\boxtimes14hab}.
\tag{HC23}
\] Tensor the \textbf{carried} single-J DP59 basis, using precisely (HC22). It gives an integral basis of \[
       \mathcal W_M=[\alpha]^*H^0(X,\mathcal L^M)
                         \subset H^0(X,\mathcal L^{M^2}),
       \qquad \operatorname{rank}\mathcal W_M=M^{42hab}.
\tag{HC24}
\] Its parameters are the prescribed representatives of \(X[\alpha^2]/X[\alpha]\), with actual rescaled translation \[
                         x=Mt'\in X[\alpha].
\tag{HC25}
\] Indeed \(\alpha Mt'=\alpha^2\bar\alpha t'=0\), and coprimality makes \(\bar\alpha\) invertible on α-primary torsion. This remains the permitted half-orbit, not the full level-\(M^2\) theta group.

At11, α is étale and all new CM-coordinate changes under consideration are integral invertible by K030 or by the corresponding later host proof. Consequently the K022 ordinary-jet surjectivity, equality of all determinantal image ideals, and one-translate norm-one conclusion apply with \[
               M\ge3D+2,
       \qquad e(D)=\max\{1,\lceil\log_{29}(3D+2)\rceil\}.
\tag{HC26}
\] The threshold is independent of h,a,b. Each individually selected finite Heisenberg translate has finite Schwartz norm one; its coefficient vector is one standard basis vector. No new coefficient-combination height or norm is introduced by the real extension. The Fock norm, source representation dimension, class mass, actual period and finite29 coordinate changes are separate and remain in the cost.

The auxiliary multijet statement of DO030 also persists with its same threshold \(3s(D+1)-1\), provided its auxiliary torsion/order and Schwartz support remain prime to29 and11. Its proof uses the same half-orbit, integral selfdual trace lattice, and the isogeny's invertibility on that auxiliary torsion. Rational29 now has the new local ramification depth, so its explicit class-index factor is PH26 of K030, not the old degree-six formula. None of this identifies the DP59 translate with the separate literal source raw-Xi input.

\subsection{5. Auxiliary T-prime existence and size scope}\label{auxiliary-t-prime-existence-and-size-scope}

For any finite collection of words and any specified finite lower size threshold, the norm/congruence part of the T-prime interface can be supplied in the same external quadratic subfield. Let \[
                      L_T=812(g_0g_1)^2.
\tag{HC27}
\] Take the ray class field of \(\mathbb Q(\sqrt{-7})\) for the ideal modulus \(L_T\mathcal O_{K_q}\), and its normal closure over \(\mathbb Q\). There are infinitely many rational primes splitting completely in this finite normal extension, by the unconditional Chebotarev theorem. At such a prime q, a chosen prime ideal of \(K_q\) lies in the trivial ray class, hence has a generator \(\pi_T\equiv1\pmod{L_T}\). Since the prime has degree one and \(K_q\) is imaginary quadratic, \(\pi_T\bar\pi_T=q\). Its conjugate ideals are distinct and coprime. Moreover \(q\equiv1\pmod{812}\), so it splits completely in K and supplies the required42-element half with residue norm q. Choosing distinct such q outside a finite set gives the requested finite family.

For any fixed finite weight data one can choose a real size parameter a with \(2fN_m+1<4a\) and every source spread at most \(2a\), then choose these T-primes larger than \(4a\). This proves qualitative nonemptiness of the corresponding norm, congruence and lower-size clauses. It supplies no effective upper bound for T and no assertion that this choice meets the quantitative MainX window. The same literal \nolinkurl{piT_Kq} obstruction remains as for G, since the source's common field is not this external quadratic subfield.

\subsection{6. Literal source failures and exact scope}\label{literal-source-failures-and-exact-scope}

The following distinctions are necessary.

\begin{longtable}[]{@{}
  >{\raggedright\arraybackslash}p{(\columnwidth - 2\tabcolsep) * \real{0.5000}}
  >{\raggedright\arraybackslash}p{(\columnwidth - 2\tabcolsep) * \real{0.5000}}@{}}
\toprule\noalign{}
\begin{minipage}[b]{\linewidth}\raggedright
Interface
\end{minipage} & \begin{minipage}[b]{\linewidth}\raggedright
Result on this host
\end{minipage} \\
\midrule\noalign{}
\endhead
\bottomrule\noalign{}
\endlastfoot
Actual principal CM variety and full integral action & K030 gives \(J^{14}\), degree84, induced type, exact coefficient and lattice \\
Literal \nolinkurl{IsEtaT} and normalized \nolinkurl{etaVal} & K029 proves all character clauses and a finite common coefficient field \\
Numerical \nolinkurl{Ranks}, with f tied to degree & (HC1)--(HC4) supply \(f=k=42,m=2,t=1759\) \\
Numerical parity, interlace, centred and NCGaps inequalities & (HC6)--(HC10) verify them; this alone is not a \nolinkurl{ThetaX} instance \\
Exact uniform even stage-character conditions & (HC16)--(HC21) prove them for every integer uniform weight \\
G norms, coprimality, residue cardinalities, full PiGCong and lower bound & (HC12)--(HC14) give explicit values in an external quadratic subfield \\
DP59 half-orbit and integral 11-jets & (HC22)--(HC26) retain the actual allowed family with parameterized dimensions \\
\nolinkurl{eight_dvd : 8 ∣ f} & Fails for \(f=42\); a further real extension of degree divisible by4 is necessary numerically \\
\nolinkurl{Fields.Ok.quad}, \nolinkurl{Primes.Ok.piG_Kq}, \nolinkurl{piT_Kq} & Fail for the unchanged primitive type: its defined common field is \(\mathbb Q\) \\
Full D7, prescribed raw Xi, later selected image and quantitative cost window & Not supplied by these constructions \\
\end{longtable}

For the common-field assertion, the induced type contains all B-lifts of each E₀ embedding, so an element on which every selected embedding agrees must first lie in E₀. The selected E₀ embeddings contain the identity, 3 and17. The subgroup generated by3 is \(\{1,3,9,27,25,19\}\), and17 is outside it; together they generate all12 automorphisms of E₀. The common element is consequently rational. Thus \[
                  \mathrm{Fields.Kq}=\mathbb Q.
\tag{HC28}
\] The source defines this field in \nolinkurl{AbelianHostFieldConditions.lean:60–76}, requires degree2 at lines108 and157, and requires the prime generators to belong to it at lines210 and226. This is a precise restatement of the prior primitive-type obstruction. It cannot be repaired by relabeling the external \(\mathbb Q(\sqrt{-7})\) as \nolinkurl{Fields.Kq}.

For an induced type under any further totally real extension, the same common-field obstruction persists. Agreement across all lifts first forces descent to E₀, and the preceding group argument still gives \(\mathbb Q\). In contrast, the numerical condition \(8\mid42h\) is equivalent to \(4\mid h\). At its first possible real extension degree4, f becomes168 and the rank constraint becomes \(k\ge168\), with G-size threshold at least \(2\cdot168^2=56\,448\). A construction of that larger polarized host must be supplied separately; the formal divisibility check alone would not construct it.

\subsection{7. Sources and verification scope}\label{sources-and-verification-scope-1}

\begin{itemize}
\tightlist
\item
  Exact source: \nolinkurl{Formalism/LocalRingCardinality/Params.lean:6–22,28–40}; \nolinkurl{Formalism/LocalToGlobal/Columns/AbelianHostFieldConditions.lean:60–76,108–125,206–238,293–331,486–494,560–579}; \nolinkurl{Formalism/LocalToGlobal/UnitaryNorms/RealCyclotomicTwoPower.lean:54}; \nolinkurl{Formalism/LocalRingCardinality/DataTheta.lean:501–555}; \nolinkurl{Formalism/SharpRestriction/CuspSpanNormBounds/CayleySignature.lean:106–125}.
\item
  Inherited Frame: \nolinkurl{Formalism/SharpRestriction/RamificationFrobenius/TotientGaloisClasses/ZsmulGeneratorEqZeroOfProjection.lean:16–24} and \nolinkurl{Formalism/FrobeniusOrbits/Statement.lean:27–29}. Additional recovered members are pinned to commit \nolinkurl{7cbf7a0c32fca1badf05931879a37c3f8a2b03b4}; their retrieval manifests explicitly record when no intake-author hash exists. No original-byte match is inferred from that absence.
\item
  AM-R-009 and ALP-LIFT-18(G2)--(G14): the checked E₀ norm-one torus and its actual character construction, used with zero infinity weight. AM-K029/030: the repaired odd character, full product host and exact finite comparisons. AM-DO-025: selected partition recurrence. Carried DP59 and AM-K022/DO030: the permitted basis and ordinary/multijet interfaces.
\item
  Milne, J. S. (2020). \emph{Class field theory} (Version4.03). Chapter V, Theorem3.6, printed p.158, and Theorem3.23, printed p.164, for ray class fields and unconditional Chebotarev. Primary author PDF: \url{https://www.jmilne.org/math/CourseNotes/CFT.pdf}. These are used only for the qualitative T-prime choice, not an effective size estimate.
\end{itemize}

The finite certificate checks every numerical rank clause, all displayed gap comparisons including every actual opposite-sign pair, the prime norms/congruences/primality and primitive-root orders, the common-field finite group computation, and the alpha norm/coprimality arithmetic. It does not infer analytic nonvanishing, a selected-image bound, a large-field class number, a character theorem, or a Lean instance from those finite checks.

The independently replayable files are \nolinkurl{enlarged_host_ranks_characters_alpha_certificate.py} and its JSON output. The certificate uses integer/Fraction arithmetic and direct trial division, with no probable-prime substitution. In particular it checks the literal natural subtraction in the Dpar sign rule and reconstructs the selected partition from the source weight recurrence before comparing it with (HC9)--(HC10).

\clearpage

\section{AM-DO-034: a new selected first lift on the degree-seven principal host}\label{am-do-034-a-new-selected-first-lift-on-the-degree-seven-principal-host}

Author: GPT-6 Astra Pro, ALP\_MACHINE \nolinkurl{/root/definite_oscillator}, 21 November 2026 UTC. Inherited reasoning effort is not separately exposed.

This extends AM-DO-030/033 to the actual principally polarized \nolinkurl{J^14} host of AM-K030. The selected polynomial is constructed afresh in the forty-two CM matrix blocks. The proof does not transport the old six-block upper witness through a change of individual CM weights. It retains the standard-lattice rescaling, the ramified rational-29 level, the finite Steinberg norm, and the exact class-supported extraction.

The inputs are the proved geometry and integral action of AM-K030, the finite splitting proof of AM-DO-022, the local selected-type calculation of AM-DO-025, and the actual DP59 half-orbit theorem in the checked form AM-K022. The construction is conditional only on the displayed stage character and rank/weight data being provided. AM-K031 is supplying concrete such even characters and auxiliary primes. None of this asserts the unchanged full \nolinkurl{Dat.Ok}, the original \nolinkurl{Fields.Kq} wrapper, or every numerical source-rank condition on this new host.

\subsection{1. Actual host and two coordinate conventions}\label{actual-host-and-two-coordinate-conventions}

Keep the notation of AM-K030:

\[
 E_0=\mathbf Q(\zeta_{28}),\qquad
 B\subset\mathbf Q(\zeta_{29}),\quad[B:\mathbf Q]=7,
 \qquad K=E_0B,\quad F=K^+,\quad [F:\mathbf Q]=42.
 \tag{DO34.1}
\]

Let \nolinkurl{Psi} be the induced CM type, with forty-two members. Let \nolinkurl{pi_B} be K030's totally positive generator of the prime of B above29 and \nolinkurl{e_0,...,e_6} its explicit orthonormal trace basis of \nolinkurl{pi_B^(-3) O_B}. The actual principally polarized variety is

\[
 \mathcal A=(J^2)^7=J^{14},\qquad
 \Lambda=\pi_B^{-3}\mathcal O_K,
 \quad\xi_0=\frac{\zeta_7-\zeta_7^{-1}}{14}.
 \tag{DO34.2}
\]

In the standard \nolinkurl{O_K} lattice convention the same polarization has coefficient

\[
 \boxed{\xi_K=\xi_0\pi_B^{-6}.}
 \tag{DO34.3}
\]

This is an equality of the actual integral trace form, not only of its discriminant. It identifies the standard line with the product trace-one line. Its reduced seed in ordinary product coordinates is \nolinkurl{c^14 product sigma_q}, as in K030(PH22).

Index the CM embeddings by \nolinkurl{tau=(i,epsilon,l)}, with \nolinkurl{i=1,2,3}, \nolinkurl{epsilon=+,-}, and \nolinkurl{l=0,...,6}. For every matrix entry of a Hom space, write the fourteen ordinary-J coordinates as \nolinkurl{(u_(i,j),v_(i,j))}, and put

\[
 w_{i,\pm,j}=u_{i,j}\pm i v_{i,j},\qquad
 z_{i,\pm,l}=\sum_jU_{lj}w_{i,\pm,j},
 \quad U_{lj}=\tau_l(e_j),\quad U^tU=1.
 \tag{DO34.4}
\]

These are K030's product-eigen coordinates on the lattice \nolinkurl{Lambda}. To use the standard \nolinkurl{O_K} lattice, define instead

\[
       \zeta_\tau=\tau(\pi_B)^3z_\tau,
       \qquad
       \partial_{\zeta_\tau}
                 =\tau(\pi_B)^{-3}\partial_{z_\tau}.
 \tag{DO34.5}
\]

Indeed the map from the standard lattice to \nolinkurl{Lambda} is multiplication by \nolinkurl{pi_B^(-3)}, so its product-eigen coordinate is \nolinkurl{z_tau=tau(pi_B)^(-3) zeta_tau}. This proves (DO34.5) and explains why (DO34.3) and (DO34.5) must be used together. It is not legitimate to keep both the old coefficient \nolinkurl{xi_0} and the new standard lattice.

All the changes in (DO34.4)--(DO34.5) and their inverses are integral at11. In particular the standard \nolinkurl{zeta} derivations give an integral invariant tangent basis after a coefficient extension. The \nolinkurl{U} matrix is real orthogonal, so it introduces no real operator-norm loss. The old \nolinkurl{plus/minus i} change and the positive diagonal factors in (DO34.5) still have to be included in norms and other local valuations.

\subsection{2. A new highest polynomial in the actual forty-two blocks}\label{a-new-highest-polynomial-in-the-actual-forty-two-blocks}

Fix even ranks \nolinkurl{a=k}, \nolinkurl{b=k+t+1}, with \nolinkurl{b>a}, and actual source weight data of the form used in AM-DO-025 at all \nolinkurl{tau in Psi}. Let \nolinkurl{e_tau} be the exponent parameter and let \nolinkurl{w_(0,tau)} be the actual dominant initial coefficient weight. Define

\[
 \nu_{\tau,i}=w_{0,\tau,i}+e_\tau-b/2,
 \quad 1\le i\le a,
 \qquad d_\tau=|\nu_\tau|,
 \qquad D=\sum_{\tau\in\Psi}d_\tau.
 \tag{DO34.6}
\]

AM-DO-025 proves that these are nonnegative partitions under the displayed source HC gap conditions. That calculation is local at one real place and does not depend on the number of CM embeddings. It therefore applies to each of the forty-two places here.

For \nolinkurl{Z_tau} the \nolinkurl{a by b} matrix of formal standard-coordinate creation variables, set \nolinkurl{nu_(tau,a+1)=0} and define

\[
 P_\nu(Z)=
 \prod_{\tau\in\Psi}\prod_{r=1}^{a}
 \det\bigl(Z_{\tau,[1,r],[1,r]}\bigr)
                   ^{\nu_{\tau,r}-\nu_{\tau,r+1}}.
 \tag{DO34.7}
\]

This is a new polynomial, defined directly in the new CM eigenspaces. It has total degree D and separate degree \nolinkurl{d_tau} in each block. Its diagonal monomial is

\[
             \prod_{\tau,i}Z_{\tau,i,i}^{\nu_{\tau,i}}
 \tag{DO34.8}
\]

with coefficient exactly one. Only the diagonal term of every minor can produce a monomial with no off-diagonal variable, and summing its exponents telescopes to \nolinkurl{nu_(tau,i)}. Thus \nolinkurl{P_nu} is primitive in the standard CM polynomial coefficient lattice, in particular at11.

For the common splitting character with infinity component \nolinkurl{mu_tau(z)=(z/bar z)^(-e_tau)}, the actual compact Fock module at \nolinkurl{tau} is

\[
 \det_a^{e_\tau-b/2}\boxtimes
 \det_b^{-e_\tau+a/2}\otimes
       \operatorname{Sym}(\mathbf C^b\otimes(\mathbf C^a)^*).
 \tag{DO34.9}
\]

The bideterminant (DO34.7) belongs already to its Cauchy summand indexed by \nolinkurl{nu_tau}. Its source type is exactly \nolinkurl{rho_(0,tau)^vee} and its target type is the source's \nolinkurl{rho_(1,tau)}, by DO25.5--9. Taking the tensor product over \nolinkurl{Psi} gives source type \nolinkurl{sigma=rho_0^vee} and target type \nolinkurl{rho_1}. No Schur projector, averaging denominator, or claim that the transported old polynomial still has one CM weight is needed. This resolves the first-edge weight-projection gate for a newly constructed selected polynomial.

The standard-coordinate polynomial and the analogous product-eigen polynomial satisfy the exact differential identity

\[
 P_\nu(\partial_\zeta)
       =c_\nu P_\nu(\partial_z),\qquad
 c_\nu=\prod_{\tau\in\Psi}
                      \tau(\pi_B)^{-3d_\tau}.
 \tag{DO34.10}
\]

Separate homogeneity proves this identity. In particular it is an11-adic unit scalar, but it is not discarded at29 or at infinity. With the valuation normalized by \nolinkurl{v_29(29)=1},

\[
            v_{29}(c_\nu)=-3D/7.
 \tag{DO34.11}
\]

Each entry of \nolinkurl{U} also has valuation \nolinkurl{-3/7} at29. Expressing the full standard derivative polynomial in ordinary product derivatives therefore gives the coefficient bound \nolinkurl{-6D/7} at29. At2 it retains the old upper denominator bound \nolinkurl{2^D}. These are upper bounds on denominators, not assertions that every coefficient attains them. At11 the coordinate changes are invertible and none of these denominators appears.

\subsection{3. The actual allowed half-orbit survives the field repair}\label{the-actual-allowed-half-orbit-survives-the-field-repair}

On the edge use the actual variety and line

\[
 X=\mathcal A^{ab}=J^{14ab},\qquad
 g=\dim X=42ab,\qquad
 \mathcal L=L_{2\infty}^{\boxtimes14ab}.
 \tag{DO34.12}
\]

Choose

\[
 \alpha=(1+2\sqrt{-7})^e\in E_0\subset K,
 \quad M=29^e,\quad e\ge1,
 \qquad\alpha\bar\alpha=M.
 \tag{DO34.13}
\]

Multiplication by alpha commutes with the B-action and acts diagonally on the fourteen ordinary factors. Its actual Rosati adjoint is bar alpha. Hence its polarization norm is exactly M on the new product. The ideals \nolinkurl{(alpha)} and \nolinkurl{(bar alpha)} remain coprime after extension to K: they lie on the two different relative split halves at29. At a selected prime of K above29 their nonzero valuations are \nolinkurl{7e}, since the absolute ramification degree there is seven.

Take the exact DP59 tensor product basis on the fourteen ordinary factors, and then on all \nolinkurl{ab} entries. It is an integral basis of

\[
 \mathcal W_M=[\alpha]^*H^0(X,\mathcal L^M)
         \subset H^0(X,\mathcal L^{M^2}),
 \qquad \operatorname{rank}\mathcal W_M=M^{42ab}.
 \tag{DO34.14}
\]

The prescribed quotient representatives give parameters

\[
           t'_\lambda\in X[\alpha^2],\qquad
           x_\lambda=M t'_\lambda\in X[\alpha]\subset X[M].
 \tag{DO34.15}
\]

They retain their actual odd-order symmetric lift scalars. They are exactly the permitted product orbit, not the whole theta group of \nolinkurl{L^(M^2)}. All statements follow by the actual diagonal product identification, so they do not require a new conjectural CM theta-span theorem. At11, \nolinkurl{[alpha]} is étale and is a formal isomorphism because its degree is a power of29. The integral principal line and section are the ones in K030, so the K018 reduced/Chern gauge applies to this ordinary product.

The actual finite vacuum is the transport of the product vacuum through the exact symplectic lattice isometry in K030(PH17--20). Its finite norm is one. At every nonsplit odd place, K/F is unramified and the opposite even identity Hermitian lattice is integrally hyperbolic. The proof in DO022 §3 therefore applies with the new local field and the coefficient \nolinkurl{xi_K}, which generates the actual absolute codifferent. Its Pan conductor condition uses that exact valuation and does not require absolute unramifiedness. All dyadic places are relatively split, because splitting of E0/F0 at2 persists under base change. DO020 covers them. Consequently the standard splitting and the two normalized moved-lattice averages agree, with their actual determinant-root characters, exactly as in DO22.8. This is a transport of the proved local hypotheses and proof, not an inference from equality of vacuum lines alone.

The same Heisenberg reading therefore has the fixed product period \nolinkurl{C^(14ab)}, where C is the already fixed ordinary-J period of DO022/R006. All finite Weyl and symmetric-lift scalars remain in the reading. The equality holds as a reduced germ and hence after the actual polynomial \nolinkurl{P_nu(partial_zeta)}.

\subsection{4. Simultaneous jets and one selected nonzero coefficient}\label{simultaneous-jets-and-one-selected-nonzero-coefficient}

Fix an auxiliary finite vector

\[
 \phi_{\rm aux}=\mathcal O\phi_0\ne0,
 \qquad\mathcal O=\sum_{r=1}^{s}b_rW_{y_r},
 \tag{DO34.16}
\]

supported at odd rational primes away from11 and29. Reduce the parameters modulo the actual selfdual lattice, combine their exact phases, and retain distinct classes with nonzero coefficients as in DO030. Let a fixed odd integer \nolinkurl{d} prime to11·29 kill these classes. Then they give s distinct points of \nolinkurl{X[d]}.

Every step of DO030 is valid on (DO34.12), for the following specific reasons. First, the new lattice is perfect by K030, so the same rational representative \nolinkurl{r_j=M z_j}, where \nolinkurl{M z_j=t_j}, has the fixed phase on the whole orbit. This uses the rational inclusions \nolinkurl{delta_(j,29) in M Lambda_29} and \nolinkurl{x_lambda in M^(-1) Lambda_29}; it does not substitute29 for a uniformizer. Second, the prime-to11 torsion points remain distinct on the good model. Third, the translated-line comparison is still made by the actual \nolinkurl{[d]} pullback and its integral symmetric theta lift, so its fibre scalar is an11-unit. No translated theta-null is divided out. Finally, the integral simultaneous-jet proof depends on \nolinkurl{L^3} being very ample and on the actual product line and base change, not on the number of factors.

Thus the same explicit threshold gives the surjection

\[
 M\ge3s(D+1)-1
 \quad\Longrightarrow\quad
 \mathcal W_M\twoheadrightarrow
       \bigoplus_{j=1}^{s}
           \operatorname{Jet}^D_{z_j}(\mathcal L^{M^2}).
 \tag{DO34.17}
\]

For completeness, sections of \nolinkurl{L^3} that vanish at all but one of the s specializations supply the factor of degree \nolinkurl{3(s-1)(D+1)}. At the selected point choose g formal parameters as ratios of sections of \nolinkurl{L^3}, and homogenize an arbitrary D-jet with degree \nolinkurl{3D}. The remaining exponent is at least2 and supplies a section nonzero at that point. The construction uses only unit inverses, so it is integral. Pullback by the actual alpha isogeny gives (DO34.17).

Take the standard CM tangent basis in (DO34.5). Its diagonal monomial (DO34.8) has coefficient one, and an ordinary jet with that coefficient evaluates under \nolinkurl{P_nu} to the factorial

\[
       \prod_{\tau,i}\nu_{\tau,i}!,\qquad
       B_\nu=\sum_{\tau,i}v_{11}(\nu_{\tau,i}!).
 \tag{DO34.18}
\]

Prescribe this degree-D monomial at an auxiliary point whose \nolinkurl{b_r} has minimum valuation, and prescribe zero D-jets at all other auxiliary points. Multiply the monomial by the unit \nolinkurl{M^D} to retain the source's normalized level reading. There are no lower terms. Therefore the possibly nonintegral higher coefficients of the Chern gauge cannot affect the degree-D derivative. The resulting functional has valuation \nolinkurl{min_r v_11(b_r)+B_nu}. Expanding the interpolating section in the integral actual basis (DO34.14) implies that at least one individual basis section has valuation no larger.

Consequently one permitted lambda satisfies

\[
 a_\lambda\ne0,\qquad
 v_{11}(a_\lambda)\le\min_rv_{11}(b_r)+B_\nu,
 \qquad
 \vartheta_{r_{P_\nu}\mathcal O\phi_\lambda}(1,1)
                   =C^{14ab}a_\lambda.
 \tag{DO34.19}
\]

This proves selected first nonvanishing in the forty-two CM blocks. It does not assert the old fixed-row sharp slope, which concerned a different explicitly selected directional witness. The present factorial upper bound is its own fully proved statement.

\subsection{5. Finite Steinberg input on the new host}\label{finite-steinberg-input-on-the-new-host}

Suppose \nolinkurl{g_0,g_1} are distinct odd rational primes, different from11 and29, that split completely in K. Fix the forty-two selected relative split places above each. Let the common character mu be unramified away from these G places and29, tame at G, and trivial on the principal rational29 level used below. Its infinity exponent must be the one in (DO34.6). The concrete character supplied in K031 is unramified at29 and meets this last condition.

At each G place use exactly the determinant-corrected free residue-matrix construction of DO033. Its genuine source action is \nolinkurl{eta(det g)^(-1)|det g|^(b/2) f(Xg)}. It contains the pure finite Steinberg type, has raw value one at \nolinkurl{[I_a;0]}, and has squared norm

\[
          \gamma_q=q^{-ab}\frac{|\mathrm{GL}_a(\mathbf F_q)|}
                                         {q^{a(a-1)/2}}.
 \tag{DO34.20}
\]

Its Fourier expansion is an actual auxiliary Heisenberg combination with at most \nolinkurl{q^(ab)} terms, minimum coefficient valuation zero at11, and squared coefficient norm gamma\_q. These facts are the local theorem of DO033 and require no new field extension argument once the local fields are actually degree one.

Tensoring over the eighty-four selected G places gives

\[
 \begin{split}
 d_G&=(g_0g_1)^{21a(a-1)},\\
 \Gamma_G&=(g_0g_1)^{42(-ab+a(a+1)/2)}
       \prod_{\epsilon=0}^{1}\prod_{i=1}^{a}
                                     (1-g_\epsilon^{-i})^{42},\\
 s_G&\le(g_0g_1)^{42ab},\qquad
 \min_rv_{11}(b_r)=0,\qquad
 \|\mathcal O_G\phi_\lambda\|_f^2=\Gamma_G.
 \end{split}
 \tag{DO34.21}
\]

Apply (DO34.19) with this actual vector and choose

\[
 e=\max\{1,\lceil\log_{29}(3s_G(D+1)-1)\rceil\},
                  \qquad M=29^e.
 \tag{DO34.22}
\]

The result is one permitted half-translate, still carrying the selected real source type and the pure finite Steinberg source type, with nonzero coefficient and upper valuation \nolinkurl{B_nu}. Its finite norm is exactly as in (DO34.21); no norm of an interpolating linear combination is being used.

\subsection{6. The new level, rational stabilizer, and exact extraction}\label{the-new-level-rational-stabilizer-and-exact-extraction}

At29 use principal congruence modulo the rational ideal \nolinkurl{(M)=(29^e)} at all six real-field primes over29. At each G place keep the full integral compact, and let the source function vary there by its finite matrix coefficient. Keep the full integral compact elsewhere. Denote the resulting product by \nolinkurl{K*}.

The principal29 subgroup fixes every permitted finite translate exactly. Its parameter lies in the selected relative split isotropic position space, with denominator dividing M. Acting by \nolinkurl{g,h=1 mod M} displaces it by an integral vector in that same space. The symplectic pairing and the odd-place trace-one phase on that displacement are zero. This proves the exact Weyl fixation, even though the valuation of M at that local prime is \nolinkurl{7e}. The compact character is trivial there by hypothesis.

Normalize real Haar measure to volume one, and give the standard full finite compact volume one. K030(PH26) gives

\[
 I_a^{K}(M)=[K_{a,0}:K^*]
       =M^{42a^2}\prod_{i=1}^{a}(1-29^{-i})^6,
                  \qquad\operatorname{vol}K^*=(I_a^K(M))^{-1}.
 \tag{DO34.23}
\]

The notation here records the index caused by the29 refinement; the G compact factors have not been shrunk. There are six residue fields, each equal to \nolinkurl{F_29}, and each congruence length is \nolinkurl{7e}. Replacing the exponent6 on the residue product by42 would be incorrect.

The rational stabilizer \nolinkurl{G(F) intersect K*} is trivial. It is finite because it preserves the global integral lattice and is compact at all real places. At a place over29 it lies in \nolinkurl{1+29 M_a(O_(F_v))}. If a matrix B has minimum uniformizer valuation \nolinkurl{s>=7}, the unique lowest term of \nolinkurl{(1+B)^29-1} is \nolinkurl{29 B}, of valuation \nolinkurl{s+7}: all middle terms have valuations at least \nolinkurl{7+2s}, and \nolinkurl{B^29} has valuation \nolinkurl{29s>s+7}. Raising to an integer prime to29 preserves the minimum valuation. Thus this congruence subgroup is torsion free. The rational matrix embedding is injective, proving the assertion.

Let \nolinkurl{u_infty} be the unit source highest vector and \nolinkurl{u_G} the product of DO033's unit generic Steinberg vectors. Set \nolinkurl{d_rho=dim rho_0}. On the identity finite \nolinkurl{K*} class define

\[
 F(k_\infty,k_f)=
 \overline{\langle u_\infty,\sigma(k_\infty)u_\infty\rangle
             \langle u_G,\tau_G(k_G)u_G\rangle},
 \tag{DO34.24}
\]

and set it to zero on all other classes. Triviality of the rational stabilizer makes this a well-defined automorphic function. It is cuspidal because the rational unitary group is definite. Its real type is \nolinkurl{rho_0}, and its G compact type is the contragredient finite Steinberg type, with the literal initial inverse generic character.

On that class the actual kernel is \nolinkurl{L(sigma(k_infty)u_infty tensor tau_G(k_G)u_G)} with value \nolinkurl{C^(14ab) a_lambda} at the original tensor. Applying compact Schur orthogonality to both unit vectors returns the original tensor divided by \nolinkurl{d_rho d_G}. The class volume is the one in (DO34.23). Therefore

\[
 \boxed{
 \Theta^{\rm bil}_{r_{P_\nu}\mathcal O_G\phi_\lambda}(F)(1)
       =\frac{C^{14ab}a_\lambda}{I_a^K(M)d_\rho d_G}\ne0,
 \quad
 \|F\|_2^2=\frac1{I_a^K(M)d_\rho d_G},\quad
 \|F\|_\infty=1.}
 \tag{DO34.25}
\]

This is an actual first automorphic lift on the new host. It has an actual irreducible source constituent containing the finite Steinberg type. Finite-dimensional compact projection to a nonzero constituent is injective on that irreducible compact copy. Its Borel-fixed line therefore gives an Iwahori-fixed vector in the constituent. This supplies the initial hypothesis of the local theorem ALP-LIFT-19 on this host, whenever its stage characters and global rank-gap hypotheses are also supplied.

No class-number-one statement about K or F is needed: the input is supported on a single actual class, whose mass was computed using its trivial stabilizer.

\subsection{7. Norms and costs that remain after the repair}\label{norms-and-costs-that-remain-after-the-repair}

Let \nolinkurl{h(nu)} be the hook product of a partition. In the standard-coordinate convention (DO34.5), the exact real Schwartz norm is

\[
 \|r_{P_\nu}\phi_{0,\infty}\|^2
 =\det(2Y)^{-7ab}\,
   |c_\nu|^2
   \prod_{\tau\in\Psi}
        \left[ h(\nu_\tau)
          \left(\frac{\pi}{2G_{i(\tau)}}\right)^{d_\tau}\right].
 \tag{DO34.26}
\]

The vacuum exponent is \nolinkurl{14ab/2=7ab}. The Fischer norm of each primitive highest bideterminant is its hook product, by AM-R-011. Orthogonality of U and the metric in K030(PH25) give the creation scale \nolinkurl{pi/(2G_i)}. Equation (DO34.10) supplies the exact remaining scalar. This derives (DO34.26) in the actual product measure; it does not normalize the real Gaussian to one.

Unit-normalizing the finite input and the source input changes (DO34.25) to

\[
 \frac{C^{14ab}a_\lambda}
        {\sqrt{\Gamma_G}\sqrt{I_a^K(M)d_\rho d_G}}.
 \tag{DO34.27}
\]

If the real Schwartz vector is also made unit, divide this value by the square root of (DO34.26). Thus the standard-lattice scalar, all Gaussian/hook factors, and the auxiliary finite normalization remain visible.

At11 one has

\[
 \begin{split}
 v_{11}(d_G)&=0,\\
 v_{11}(I_a^K(M))&=6\sum_{i=1}^{a}v_{11}(29^i-1),\\
 v_{11}(\Gamma_G)&=42\sum_{\epsilon=0}^{1}
                         \sum_{i=1}^{a}v_{11}(g_\epsilon^i-1),\\
 v_{11}\!\left(C^{-14ab}\Theta^{\rm bil}(F)(1)\right)
     &\le B_\nu-v_{11}(I_a^K(M)d_\rho).
 \end{split}
 \tag{DO34.28}
\]

For the candidate primes \nolinkurl{4,615,409} and \nolinkurl{5,276,377} supplied by K031 and rank \nolinkurl{a=42}, all three integers \nolinkurl{29,g_0,g_1} have order10 modulo11, and \nolinkurl{v_11(q^10-1)=1}. Thus the sums over \nolinkurl{1<=i<=42} equal4, giving the exact fixed factors

\[
          v_{11}(I_{42}^{K}(M))=24,
          \qquad v_{11}(\Gamma_G)=336.
 \tag{DO34.29}
\]

This last calculation is independent of e. Unit finite normalization has scalar valuation cost168 in this example. It does not invalidate the raw minimum coefficient valuation zero used in (DO34.19).

For fixed ranks, fixed G primes, and variable selected weights, (DO34.22) has \nolinkurl{log M=O(log(D+2))}; the class index and compact dimension consequently give an \nolinkurl{O(log(D+2))} source extraction cost. The coordinate factor in (DO34.10), its29 denominators, and the real hook/Fock norm have their own degree dependence and are not included in that logarithmic assertion. In particular the possible \nolinkurl{6D/7} ordinary-coordinate 29 denominator is not a constant-level cost.

\subsection{8. Scope and exact inputs}\label{scope-and-exact-inputs}

This proves an actual selected first kernel, with a single permitted29 half-translate, a specified finite Steinberg auxiliary vector, and an exact class-supported source input on the new principal host. It closes the first CM-weight and actual-half-orbit transfer gap. It does not identify this auxiliary vector with the literal raw \nolinkurl{PhiG} or \nolinkurl{Xi}, prove the terminal restricted Gram nonvanishing, or compute the integral lattice of a later iterated theta image. It also does not turn the enlarged field into the unchanged original host record when unrelated field/rank wrappers fail.

The local selected Schur identity and splitting proof are used as settled earlier findings. The novelty here is their exact simultaneous realization on K030's standard principal lattice and allowed orbit, with the forty-two-block polynomial and the changed costs.

Exact locators:

\begin{itemize}
\tightlist
\item
  AM-K030(PH3--9, PH17--26): explicit trace basis, actual integral product polarization, standard coefficient, K-eigen coordinates, and rational29 congruence index.
\item
  AM-DO-022 §§3--4: standard finite splitting comparison, Pan conductor and compact phase, arbitrary odd local base field, determinant roots, and split dyadic case.
\item
  AM-DO-025(DO25.3--9): literal common-character compact Cauchy summand and source/target coefficient match; source \nolinkurl{Columns/BaseLemmas.lean} lines496--552.
\item
  AM-K022(J3--5, J6--16, J19--27): actual DP59 orbit, integral jet threshold, factorial detection and source-normalized differential reading.
\item
  AM-DO-030 §§2--4: fixed auxiliary phases, translated-line unit comparison and simultaneous integral jets.
\item
  AM-DO-033 §§1--5: primitive finite Steinberg vector, actual determinant-corrected residue model, Fourier norm and compact source extraction.
\item
  AM-R-011: exact Fischer coefficient formula, highest bideterminant hook norm, and creation/Gaussian scaling.
\item
  Primary geometric facts are the already read Milne, \emph{Abelian Varieties}, I§6, p.30, and Mumford(1966), \emph{On the equations defining abelian varieties I}, §1(II), p.289, as documented in K022/DO030.
\end{itemize}

\clearpage

\section{AM-R-017: explicit auxiliary T-primes and a first source class at maximal 29-level}\label{am-r-017-explicit-auxiliary-t-primes-and-a-first-source-class-at-maximal-29-level}

Author: /root, GPT-6 Astra Pro, ALP\_MACHINE, 22 November2026UTC. Effort is not externally exposed.

This result makes two previously qualitative choices concrete. It supplies one explicitly certified auxiliary T-prime for each of the K031 and K032 hosts. It also constructs a nonzero first automorphic input with the G generic type and nontrivial T dilation characters, without imposing a principal level at29 to remove rational stabilizers. Its exact class mass is computed. The theta lift of this input with the complete prescribed raw kernel is not proved nonzero here: averaging a translated29 kernel over the maximal29 compact remains a separate issue. The unchanged source \nolinkurl{Fields.Kq} obstruction also remains.

\subsection{1. Two explicit norm primes with the full source congruence}\label{two-explicit-norm-primes-with-the-full-source-congruence}

Use the already checked fields of K031 and K032, and write \(\delta=\sqrt{-7}\) in their common external quadratic subfield. The following data use the exact G-primes of those findings:

\begin{longtable}[]{@{}lrrrrrr@{}}
\toprule\noalign{}
Product host & f & k & Rational conductor c & g0 & g1 & n \\
\midrule\noalign{}
\endhead
\bottomrule\noalign{}
\endlastfoot
J14 & 42 & 42 & 812 & 4615409 & 5276377 & 3 \\
J56 & 168 & 168 & 52780 & 22285932761 & 64072175441 & 16 \\
\end{longtable}

In each row set \[
 L=c(g_0g_1)^2,\qquad A=1+nL,\qquad
 \pi_T=A+L\delta,\qquad T=A^2+7L^2.
 \tag{R17.1}
\] Thus every number in the construction is specified without a prime-search oracle. The exact prime integers, also recorded in the certificate, are:

\begin{itemize}
\tightlist
\item
  J14: \nolinkurl{3710360310942963702130514092646914581678165356491757211050633}.
\item
  J56: \nolinkurl{3045759889845354427724717877501093821489629367456155582288640401018208099619883018305672886530161}.
\end{itemize}

Both T are prime. Here is the complete certificate method, including the factored divisors required to make the primality conclusion deterministic. For J14 use \(F=586L\), and for J56 use \(F=44L\). Their factorizations are \[
\begin{aligned}
 F_{14}&=2^3\cdot7\cdot29\cdot293\cdot g_0^2g_1^2,\\
 F_{56}&=2^4\cdot5\cdot7\cdot11\cdot13\cdot29\cdot g_0^2g_1^2.
\end{aligned}
\tag{R17.2}
\] Each listed prime factor is verified by trial division through its integer square root. The replay verifies \(F\mid T-1\), \(F^2>T\), and, for every prime \(r\mid F\), \[
 z_r^{T-1}\equiv1\pmod T,\qquad
 \gcd(z_r^{(T-1)/r}-1,T)=1,
 \quad z_2=3,\quad z_r=2\ (r\ne2).
 \tag{R17.3}
\] These are precisely Pocklington's factored-divisor criterion. The exact modular residues and gcds are in \nolinkurl{checks/explicit_t_prime_and_stabilizer_certificate.json}. Candidate discovery used a separate primality screen, but that screen is not an input to the saved replay or to this primality proof. The replay uses only integer arithmetic, modular powers, gcds, and trial division on the small factors in(R17.2).

The norms and congruences required in the source follow exactly: \[
 \pi_T\bar\pi_T=T,\qquad \pi_T\equiv1\pmod{c(g_0g_1)^2},
 \qquad T\equiv1\pmod c.
 \tag{R17.4}
\] The certificate checks \(\gcd(T,2A)=1\). A prime dividing both \(\pi_T\) and its conjugate would divide T and their sum \(2A\), which is impossible. The two ideals therefore are coprime, also after extending to either host. Since the host is contained in the c-th cyclotomic field, T splits completely. Exactly f primes divide \(\pi_T\), all to exponent one and with residue field \(\mathbf F_T\): at each conjugate pair one must divide the product T, and coprimality forces exactly one to divide \(\pi_T\). This verifies the norm, congruence, coprimality, half-cardinality and residue-norm clauses for this explicit choice.

The first T has61 digits and the second97. This supplies exact finite values, not an asymptotic upper bound for a changing-weight family. For example, retaining the K031/K032 constants C, define \[
 A_{\rm size}=2(2k+1)C+2+2fN_2.
 \tag{R17.5}
\] The two values are respectively \nolinkurl{7676528679020} and \nolinkurl{27317450103536}. The replay verifies \(2fN_2+1<4A_{\rm size}<T\). These are the \nolinkurl{a_gt} and \nolinkurl{T_gt} inequalities of \nolinkurl{Sizes} for this size parameter. The independently specified \nolinkurl{spread_le} inequality must still be verified for whichever full collection of source words is used; it is not inferred merely from finding a large T. No quantitative MainX window is claimed by(R17.5).

The field-membership clause \nolinkurl{piT_Kq} is deliberately not asserted. K031 proved that the source's defined common field is Q for this induced primitive type. The present \(\pi_T\) lies in the actual external \(\mathbf Q(\sqrt{-7})\), as do the previously constructed G generators.

\subsection{2. The exact rational stabilizer of a standard definite lattice}\label{the-exact-rational-stabilizer-of-a-standard-definite-lattice}

Let E be either host, let F be its maximal real subfield, and let \(X=(E^a,I_a)\). Write \[
 K_0=\prod_{v<\infty}U(I_a,\mathcal O_{F_v}),\qquad
 \Gamma=U(X)(F)\cap K_0.
 \tag{R17.6}
\] At split places this notation means the standard \(\mathrm{GL}_a\) compact; equivalently K0 stabilizes \(\mathcal O_E^a\) at every finite place. Then \[
                  \Gamma=\mu(E)^a\rtimes S_a.
 \tag{R17.7}
\]

Indeed every entry x of a matrix in Gamma is integral, and at every complex embedding \(|\tau(x)|\le1\), by the unitary row norm. If x is nonzero, its algebraic norm is a nonzero integer, so the product of these absolute values is at least1. Every factor must therefore equal1. Such an integral x is a root of unity: all powers of x lie in the bounded set of the Minkowski lattice \(\mathcal O_E\) with these absolute-value bounds; that set is finite, so two powers coincide. The row norm now forces exactly one nonzero entry in each row. Orthogonality forces their column positions to be distinct. This proves containment in the displayed monomial group, and the reverse containment is immediate.

For both fields \(\mu(E)=\mu_{28}\). K031 has Galois group \(C_2^2\times C_{21}\), and K032 has \(C_2^4\times C_{21}\), each of exponent42. Their conductors are c from the table. A cyclotomic subfield arising from an additional root of unity must have conductor dividing c. The possible new factors29, and also5 or13 in the second case, would require a quotient Galois group of exponent divisible by4, which these groups do not possess. The2-part of c is4 and its7-part is7. Since i and \(\zeta_7\) already belong to both fields, exactly \(\mu_{28}\) remains.

Choose one of the T places and identify its residue field with \(\kappa=\mathbf F_T\). Reduction embeds \(\mu_{28}\) in \(\kappa^\times\), since T is prime to28. Suppose \(b\in\kappa^\times\) has order greater than a in the quotient \(\kappa^\times/\mu_{28}\), and take \[
 u=(1,b,b^2,\ldots,b^{a-1}),\quad
 P_u=\{g\in\mathrm{GL}_a(\kappa):ug=s(g)u\},\quad
 H_u=\ker(s:P_u\to\kappa^\times).
 \tag{R17.8}
\] Then the exact intersections with the rational group are \[
 \boxed{\quad\Gamma\cap\widetilde P_u=\mu_{28}I_a,
               \qquad\Gamma\cap\widetilde H_u=\{1\}.\quad}
 \tag{R17.9}
\] Tildes denote reduction preimages inside K0 at that place, with no extra restriction at the other places.

To prove(R17.9), pass to \(\kappa^\times/\mu_{28}\). A monomial g in \(P_u\) must carry the set of coordinate classes \(\{1,[b],\ldots,[b]^{a-1}\}\) to itself by multiplication by \([s(g)]\). The class \([s(g)]\) belongs to the cyclic subgroup generated by \([b]\), because it carries one member of that set to another. A proper consecutive interval of length a in a cyclic group of order greater than a has trivial translation stabilizer: its cyclic consecutive gaps have a unique gap greater than1, which any stabilizing translation must fix. The same claim is immediate for a=1. Thus \([s(g)]=1\). All coordinate classes are distinct, so the permutation of g is the identity. Every diagonal entry equals s(g), and g is scalar. Fixing u then forces that scalar to be1. Conversely every scalar root lies in \(P_u\), as required.

For both numerical rows we may take b=2. The certificate checks \[
                     2^{28j}\not\equiv1\pmod T
                         \quad(1\le j\le a),
 \tag{R17.10}
\] with a=42 or168. Thus no primitive-root search or complete factorization of T-1 is required. Formula(R17.8) is an explicit residue functional. The argument remains valid if the other chosen T places use the same functional, and if the G places are kept at their full maximal compact.

\subsection{3. A nonzero class-supported input and its exact mass}\label{a-nonzero-class-supported-input-and-its-exact-mass}

Let \(\rho_\infty\) be any of the specified irreducible compact types on \(U(a)^f\), let \(d_\rho=\dim\rho_\infty\), and choose a unit vector \(v_\rho\). At each G place use the finite Steinberg representation and its unit generic vector from DO033/LIFT20. Denote their dimensions by \(d_{G,P}\), and put \(d_G=\prod_{P\in G}d_{G,P}\). Their central characters are trivial, because they occur in the finite flag module \(\operatorname{Ind}_B^{\mathrm{GL}_a(\kappa)}1\).

Let H be K0 with its component at each of the f chosen T places replaced by \(\widetilde P_u\); use(R17.8) at every such place. With each maximal compact having volume1, \[
 d_T=[\mathrm{GL}_a(\mathbf F_T):P_u]
       =\frac{T^a-1}{T-1},\quad
 \operatorname{vol}(H)=d_T^{-f},\quad
       \Gamma_H=U(X)(F)\cap H=\mu_{28}I_a.
 \tag{R17.11}
\] The index is the number of projective nonzero row vectors: the full general linear group is transitive on them and Pu is the stabilizer. Equation(R17.9) gives the last assertion; no congruence at29 has entered.

Write \(\omega_\infty\) for the character of \(\mu_{28}\) on \(\rho_\infty\). Choose multiplicative residue characters \(\nu_P:\mathbf F_T^\times\to\mathbf C^\times\), all nontrivial, such that \[
       \omega_\infty(\zeta)\prod_{P\mid\pi_T}\nu_P(\bar\zeta_P)=1
                       \qquad(\zeta\in\mu_{28}).
 \tag{R17.12}
\] This choice always exists. Restriction from the character group of the cyclic group \(\mathbf F_T^\times\) to that of its subgroup \(\mu_{28}\) is surjective. Each restriction has \((T-1)/28>1\) extensions, so it has a nontrivial extension even if the prescribed restriction is trivial. Prescribe trivial restrictions at all but one place and the required inverse character at the remaining place. This also explains the finite coefficient extension needed for the \(\nu_P\): roots of order T-1 must be included. Their coefficient-field cost is not discarded.

For \(h\in H\) and \(k\in U(a)^f\), define \[
 m(k,h)=\langle v_\rho,\rho_\infty(k)v_\rho\rangle
   \prod_{P\in G}\langle v_P,\operatorname{St}_P(\bar h_P)v_P\rangle
   \prod_{P\mid\pi_T}\nu_P(s(\bar h_P)).
 \tag{R17.13}
\] Use the conjugate-linear-first Hermitian convention in this display. Define \(F\) on \([U(X)]\) to be \(\overline{m(k,h)}\) on \(U(X)(F)U(a)^fH\), and zero off this class. This is well-defined: two presentations differ by an element of \(\Gamma_H\), and(R17.12) says exactly that the resulting product multiplier is1. F is smooth, has the required dual real compact type and inverse G-generic character, and is invariant under \(\widetilde H_u\) at T. At29 it is invariant under the full maximal compact. It is nonzero, since F(1)=1. It is an automorphic cusp form because U(X) is anisotropic over F.

The exact scalar-dilation augmentation at each T place is zero: right multiplication by an element of Pu with s-coordinate t multiplies F by \(\nu_P(t)^{-1}\), and \(\sum_{t\in\mathbf F_T^\times}\nu_P(t)^{-1}=0\). This is the local augmentation condition used for the initial ramified column. Reading it as the full source \nolinkurl{RamT} record additionally uses its specified decoration/dilation dictionary; the scalar identity alone does not instantiate that independent record.

Normalize Haar measure on the compact real group to1, as in DO033. The quotient of the support class has measure \(d_T^{-f}/28\). Schur orthogonality at the real group and the independent G compacts therefore gives \[
 \boxed{\quad \|F\|^2
       =\frac{1}{28\,d_T^f\,d_\rho\,d_G},
       \qquad F(1)=1,\qquad\|F\|_\infty\le1.\quad}
 \tag{R17.14}
\] Every T character has modulus1 on its support. The factor28 comes from the actual rational stabilizer and is not absorbed into Haar measure. The finite G scalar denominators and the coefficient metric of \(v_\rho\) remain those of DO033/R011; this statement does not replace their primitive integral models by an arbitrarily rescaled unit vector.

For the explicit ranks the new projective-level index is an11-adic unit: \[
                         v_{11}(d_T)=0.
 \tag{R17.15}
\] For J14, T is2 modulo11 and a=42, so neither numerator nor denominator of dT is divisible by11. For J56, T is1 modulo11 and a=168, so \(v_{11}(T^a-1)-v_{11}(T-1)=v_{11}(168)=0\). The certificate checks these exact valuations. The analytic norm cost \(d_T^f\), despite its11-adic unit valuation, is still large and fully retained in(R17.14).

\subsection{4. What this removes, and what it leaves}\label{what-this-removes-and-what-it-leaves}

The old first-input construction made the rational stabilizer trivial by a principal level at29. Equations(R17.8)--(R17.14) supply a different actual input with the required first G generic behavior and T augmentation while keeping maximal level at29. Thus rational-stabilizer compatibility, by itself, is no longer an obstruction to a source input at that level. The two numerical T choices also replace the qualitative prime choice for these fixed hosts by exact certified integers.

There is still no nonvanishing assertion for the raw kernel applied to this F. If a previously chosen29 Weyl translate is inserted into the theta kernel, the source integral at maximal29 level averages that translate over the compact. Its resulting selected coefficient can vanish. Neither the explicit T-prime nor(R17.14) proves that projection nonzero. The previous-image integral lattice, the terminal row map, the global MainX mass, and the changed-family DP44 cost remain separate. Likewise(R17.4) does not repair \nolinkurl{Fields.Kq}, and(R17.5) does not prove the complete quantitative size window for an unbounded weight family.

\subsection{5. Exact sources and checks}\label{exact-sources-and-checks}

\begin{itemize}
\tightlist
\item
  K031, (HC3)--(HC16) and section5; K032, sections4--6: the two actual product hosts, G-primes, conductors, ranks and constants. These are used with their recorded whole-proof checks.
\item
  Alpöge's stipulated pinned source at commit \nolinkurl{7cbf7a0c32fca1badf05931879a37c3f8a2b03b4}: \nolinkurl{Formalism/LocalToGlobal/Columns/AbelianHostFieldConditions.lean}, \nolinkurl{Primes.Ok}206--238 and \nolinkurl{Sizes}486--494; these are the literal norm, residue, congruence, common-field and size conditions distinguished above. \nolinkurl{LevChiEqOne.lean}, \nolinkurl{RamT}85--87, is the separate global decoration/dilation condition.
\item
  Pocklington, H. C. (1914--1916). The determination of the prime or composite nature of large numbers by Fermat's theorem. \emph{Proceedings of the Cambridge Philosophical Society, 18},29--30. The precise factored-divisor form used here is also stated as Theorem1 in the FLINT official documentation, \emph{Primality proving: APRCL and ECPP}, \url{https://flintlib.org/doc/algorithms/number-theory/primality-proving.html}. The latter statement was consulted directly; this is not a claim that the inaccessible original PDF was read in full.
\item
  DO033 sections1 and4, DO034 section6, R011: matrix-coefficient normalization, actual compact masses and the separate coefficient metric. The new class calculation retains their conventions.
\end{itemize}

The exact replay is \nolinkurl{scripts/explicit_t_prime_and_stabilizer_certificate.py}. Its finite checks certify(R17.1)--(R17.5),(R17.10),(R17.15). The proofs of(R17.7)--(R17.14) are the mathematical arguments above, not consequences of a numerical experiment. Whole-proof review is recorded separately.

\clearpage

\section{ALP-LIFT-21: the raw terminal G factor survives the actual previous theta constituent}\label{alp-lift-21-the-raw-terminal-g-factor-survives-the-actual-previous-theta-constituent}

Author: GPT-6 Astra Pro, ALP\_MACHINE /root/lift\_accounting. Date: 22 November 2026 UTC. Inherited reasoning effort is not exposed. The rank-filtration reduction in section4 was suggested independently by /root while the explicit full-rank intertwiner was being derived.

This theorem concerns the actual terminal local G factor of XiCpl2 for m=2, with the same allowed coupled class used in LIFT20. It proves nonvanishing in the actual small theta quotient of each LIFT19 first constituent. It proves more than occurrence of a compatible compact type: the fixed raw terminal Schwartz vector is injective on the entire prescribed first-column type. In a specified Haar-normalized induced model its norm and its effect on a natural integral coefficient lattice are computed.

The normalization of that local induced model is explicit. Its comparison with a globally pinned automorphic coefficient model is not supplied by uniqueness of the local theta correspondence. No assertion about the final MainX row, all six AdmS identities, or the full DP44 cost follows merely from the local theorem.

\subsection{1. Ranks, actual raw factor, and its compact groups}\label{ranks-actual-raw-factor-and-its-compact-groups}

Fix a degree-one relative split G place of residue characteristic different from11. Write \(F_v\) for its local field, \(O\) for its ring of integers, and \(\kappa=\mathbf F_q\) for its residue field, with the rational prime \(q\) a uniformizer. Normalize multiplicative Haar measure by \(\operatorname{vol}(\mathrm{GL}_n(O))=1\), and additive matrix measure by \(\operatorname{vol}(M_{u,v}(O))=1\).

Take the literal source ranks with \(m=2\). Put

\[
 d=t+1,\quad r=k+1,\quad b=k+d=r+t,\quad N=b+d=k+2d.
 \tag{TG1}
\]

Thus the terminal edge is \(U(b)\to U(N)\). Every actual Ranks datum has enough coordinates for the following blocks. The argument uses \(t\ge1\), so in particular applies to the larger values required by the source.

The source locators are relative to the pinned recovered tree work/research/alpoge\_source, commit 7cbf7a0c32fca1badf05931879a37c3f8a2b03b4:

\begin{itemize}
\tightlist
\item
  \nolinkurl{Formalism/LocalRingCardinality/Params.lean,} definitions of r, N, rS, and cOf in the column file.
\item
  \nolinkurl{Formalism/LocalToGlobal/Columns/FiniteSumNonzeroLargeDepth.lean} (FS), PhiG/PhiGS780--810, PhiTop1141--1147, ftY1150--1155, and XiCpl2 1160--1167.
\item
  \nolinkurl{Development/AbelianVarietyTheory/Definitions/GlobalTwisting/IntegralGLSuccEmbeddingMemParahoricTimesPrincipa.lean,} entB\_top420--441 and the terminal rank identity.
\item
  Checked DO031, equations(DO31.1)--(DO31.5), identifies those exact raw functions and their actual terminal frames.
\end{itemize}

No Lean execution is claimed. We use the definitions and the proved finite model in DO031 with its stated scope.

Write a terminal matrix in its actual local frames as

\[
 X=\begin{pmatrix}Z\\ \ell\end{pmatrix}\in M_{N,b}(F_v),
 \qquad Z\in M_{N-1,b}(F_v),\quad \ell\in M_{1,b}(F_v).
 \tag{TG2}
\]

The terminal top function has no remaining auxiliary-column power: cOf(m-1)=0, so entB\_top makes it read \(Z\) itself. The last-row factor is the actual integral character function in DO31.4. In the convention of that equation it is

\[
             1_{O^b}(\ell)\psi_v(-\ell y),\qquad
                y\equiv q^{-1}e_k\pmod{O^b}.
 \tag{TG3}
\]

The class of \(y\) is exactly LIFT20(RG3), now in the rank-b middle space. If the last row is instead written before the source's piT unit rescaling, a fixed unit multiplies the argument in(TG3). Every proof below keeps the resulting root of unity, and all support, norm, and content assertions are unchanged. We use DO031's already-rescaled convention.

Let \(\eta_0,\eta_1\) be the actual tame selected-half characters on the first and second edges. Their residue characters are nontrivial and distinct, and the literal pin is \(\bar\eta_i=\chi_i^{-1}\). Denote the finite groups

\[
 \begin{aligned}
 H&=\left\{
   \begin{pmatrix}u&A\\0&h\end{pmatrix}:
       u\in U_r(\kappa),\ A\in M_{r,t}(\kappa),\
       h\in\mathrm{GL}_t(\kappa)\right\},\\
 H^\sharp&=\{g\in H:g e_k=e_k\}.
 \end{aligned}
 \tag{TG4}
\]

The index \(k\) is zero-based, so it is the last of the first \(r=k+1\) coordinates. Thus \(H^\sharp\) has \(u=\operatorname{diag}(u_k,1)\), with \(u_k\in U_k\), while \(A,h\) are unchanged. Write \(J,J^\sharp\) for their inverse images in \(\mathrm{GL}_b(O)\). Their cardinalities and Haar volumes are

\[
 \begin{split}
 h_q&=|H|=q^{r(r-1)/2+rt}|\mathrm{GL}_t(\kappa)|,\\
 h_q^\sharp&=q^{-k}h_q,\qquad
 v_J=\frac{h_q}{|\mathrm{GL}_b(\kappa)|},\qquad
 v_{J^\sharp}=q^{-k}v_J.
 \end{split}
 \tag{TG5}
\]

Put

\[
 \lambda(g)=
   \psi_0\!\left(-\sum_{i=0}^{r-2}u_{i,i+1}\right)
                         \bar\eta_0(\det h)\qquad(g\in H).
 \tag{TG6}
\]

This is exactly the inverse adelic first-column character from LIFT19(T1)--(T5), in the \(j=1=m-1\) fully generic branch. The source RecLevel.lean50--52 gives the inverse convention. We use the same notation for its inflation to \(J\) and its restriction to \(J^\sharp\).

\subsubsection{Exact raw support and phase}\label{exact-raw-support-and-phase}

Let \(\Phi,\Phi^\sharp\) denote the actual nonsharp and sharp terminal local factors, including(TG3). They vanish unless \(X\) is integral and, modulo \(q\), has the form

\[
 \bar X=\begin{pmatrix}B\\0_{(d-1)\times b}\\w\end{pmatrix},
 \qquad
 B=\begin{pmatrix}u&A\\0&h\end{pmatrix}\in H
       \text{ or }H^\sharp,\qquad w\in\kappa^b.
 \tag{TG7}
\]

On that support,

\[
 \begin{split}
 \Phi(X)&=
  \psi_0\!\left(\sum_{i=0}^{r-2}u_{i,i+1}-w_k\right)
                       (\bar\eta_1\bar\eta_0^{-1})(\det h),\\
 \Phi^\sharp(X)&=
  \psi_0\!\left(\sum_{i=0}^{k-2}u_{i,i+1}-w_k\right)
                       (\bar\eta_1\bar\eta_0^{-1})(\det h).
 \end{split}
 \tag{TG8}
\]

An empty sum is zero. For the source ranks the sums are nonempty; this convention also gives the small diagnostic.

We verify the raw predicates. In PhiTop the internal variable a=r+(m-1)t is \(b\), the full source width. Its integrality condition makes every entry of \(Z\) integral. The subsequent row condition makes the rows below its first \(b\) rows zero modulo \(q\). The first \(r\) columns have unitriangular leading block and zeros below. The remaining \(t\) columns give an arbitrary upper-right rectangle and a square bottom-right block. If its determinant is zero, unitV and its inverse give zero, so the function vanishes. Otherwise that block is \(h\) in(TG7).

For rS=0, symIdx k 0 is the identity. The additional PhiGS condition fixes column \(k\) to \(e_k\), precisely \(H^\sharp\). The determinant term is \(\chi_0(\det h)\chi_1(\det B)^{-1}\), with \(\det B=\det h\), giving the factor in(TG8). The additive term is the full \(U_r\) character for PhiG; in PhiGS it is the \(U_k\) character because its second root sum has length r-k-1=0. Finally(TG3) contributes \(\psi_0(-w_k)\). Changing \(y\) by an integral vector does not affect it. This proves(TG7)--(TG8) from the actual raw definitions.

Both functions are therefore supported in the full-rank open subset of \(M_{N,b}(F_v)\): the leading \(b\) by \(b\) minor is a unit. Their support is compact in that open set. They are fixed by the target principal congruence \(K_N^1\). Indeed its determinant character is trivial and it changes integral matrix entries only modulo \(q\).

Their exact additive Schwartz norms are

\[
 \Gamma_{\rm raw}=q^{-Nb+b}h_q,\qquad
 \Gamma_{\rm raw}^\sharp=q^{-Nb+b}h_q^\sharp
                              =q^{-k}\Gamma_{\rm raw}.
 \tag{TG9}
\]

The support has \(h_q q^b\) or \(h_q^\sharp q^b\) residue matrices, each of volume \(q^{-Nb}\), and every nonzero value has modulus one.

\subsection{2. The actual preceding constituent}\label{the-actual-preceding-constituent}

Let \(\pi_0\) be an irreducible actual initial local constituent with a finite Steinberg type as in LIFT19. It has Iwahori-fixed vectors and unramified cuspidal support, by the Casselman result already used there. The actual first and second local theta constituents in the literal bilinear convention are

\[
 \pi=\pi_0\times(\eta_0\circ\det_d),\qquad
 \pi^+=\pi\times(\eta_1\circ\det_d).
 \tag{TG10}
\]

The product denotes normalized parabolic induction. Both are irreducible: the added residue characters are distinct from the initial unramified line and from each other. This is LIFT19's application of Zelevinsky's disjoint-line theorem and Mínguez's actual type-II theta correspondence. These published results are used rather than claimed as new.

The source here is the actual constituent \(\pi\) of a first theta image, not an arbitrary representation chosen only by its archimedean weight. Its prescribed compact-type space

\[
                          V_\lambda=\pi^{J,\lambda}
 \tag{TG11}
\]

is nonzero by LIFT19(T16)--(T17). Every nonzero vector in this space works below. It is \(K_b^1\)-fixed. Restricting its character to \(J^\sharp\) gives the same vectors for the literal sharp kernel.

\subsection{3. The full-rank Haar-normalized intertwiner}\label{the-full-rank-haar-normalized-intertwiner}

The actual common-character split oscillator acts by

\[
 \omega(g,h)\phi(X)=
 \eta_1(\det g)^{-1}\eta_1(\det h)
 |\det g|^{N/2}|\det h|^{-b/2}\phi(h^{-1}Xg).
 \tag{TG12}
\]

This is the genuine left-action convention of DO033.5 and corrected LIFT16(S2). The source variable acts by \(X\mapsto Xg\). The source bilinear lift takes diagonal coinvariants of \(\omega\otimes\pi\).

Write \(\Omega\) for the full-rank open set and \(X_0=[I_b;0]\). For \(\phi\in S(\Omega)\), \(v\in\pi\), define

\[
 {\cal I}(\phi,v)=
   \int_{\mathrm{GL}_b(F_v)}
       \phi(X_0g)\eta_1(\det g)^{-1}
                    |\det g|^{N/2}\pi(g)v\,dg.
 \tag{TG13}
\]

This is an ordinary compact integral, with no meromorphic regularization. The image of \(g\mapsto X_0g\) is closed in \(\Omega\); the inverse image of the compact support of \(\phi\) is compact in \(\mathrm{GL}_b\). The smooth integrand takes values in a finite-dimensional span on that compact set, so its integral lies in \(\pi\).

Let \(P=P_{(b,d)}\) be the upper parabolic of \(\mathrm{GL}_N\). Set

\[
             {\cal T}(\phi,v)(h)
                      ={\cal I}(\omega(1,h)\phi,v).
 \tag{TG14}
\]

Then T factors through the source diagonal coinvariants and takes values in the actual normalized induction \(\pi^+\). The complete covariance checks are as follows.

First \({\cal I}(\omega(g_0,1)\phi,\pi(g_0)v)={\cal I}(\phi,v)\). Substitute \(g'=g g_0\) in(TG13). The two eta factors combine to \(\eta_1(\det g')^{-1}\), the absolute-value factors combine identically, and \(\pi(g)\pi(g_0)=\pi(g')\). Unimodularity preserves Haar.

Next let \(p=[A\ C;0\ D]\in P\), so \(p^{-1}X_0=X_0A^{-1}\). Substitute \(g=A g'\) in(TG13), using(TG12), to obtain exactly

\[
 {\cal I}(\omega(1,p)\phi,v)=
 |\det A|^{d/2}|\det D|^{-b/2}
          \eta_1(\det D)\pi(A){\cal I}(\phi,v).
 \tag{TG15}
\]

The unipotent C disappears because it fixes the embedded source space. The absolute-value factor is \(\delta_P(p)^{1/2}\). Therefore T(phi,v)(p h) has precisely the normalized induction transformation law. Target translation by \(h_0\) sends the section to \(h\mapsto T(\phi,v)(h h_0)\). Smoothness follows from a small compact fixing \(\phi\), and \(P\backslash\mathrm{GL}_N\) is compact.

The full-rank coinvariant identification is

\[
 (S(\Omega)\otimes\pi)_{\mathrm{GL}_b}
       \simeq
 \operatorname{Ind}_{P}^{\mathrm{GL}_N}
                  (\pi\boxtimes\eta_1\det_d),
 \tag{TG16}
\]

with map T and the displayed normalization. Here is an elementary verification in addition to the covariance calculation. Full-rank matrices form a locally trivial principal GL\_b bundle over the p-adic Grassmannian. That compact, totally disconnected space admits a finite clopen refinement of trivializing charts for each locally constant function in question. On a chart, multiply the fibre function by \(\eta_1(\det g)^{-1}|\det g|^{N/2}\). The source action becomes ordinary right translation on \(C_c^\infty(\mathrm{GL}_b)\), together with \(\pi\) on its coefficient vector. Its diagonal coinvariant map is \(f\otimes v\mapsto\int f(g)\pi(g)v\,dg\), an isomorphism onto \(\pi\).

To verify that assertion, a vector fixed by a compact open K is represented by \(\operatorname{vol}(K)^{-1}1_K\otimes v\). Changing to a smaller such K leaves its class unchanged by finite coset decomposition and the diagonal relations. Conversely, partition the compact support of f into cosets of a small compact open. The same diagonal relations and finite averaging reduce its class to that of its displayed integral. This proves both injectivity and surjectivity. Finite partitions of unity on the clopen charts glue the fibre maps, with transition law(TG15). This proves(TG16).

\subsection{4. Identification with the actual small theta quotient}\label{identification-with-the-actual-small-theta-quotient}

The open-set map must be related to the full oscillator. We use the published rank filtration in Mínguez(2008), section2, pp.722--724. The boundary quotients are the rank j strata with \(j<b\). Corollary2.3 describes their possible irreducible source quotients: after removing oscillator determinant normalization, such a quotient is a quotient of an induced representation with a nonzero trivial GL\_(b-j) factor. Modulus characters only introduce unramified twists of that factor. Its cuspidal support therefore contains an unramified line.

For(TG12) the relevant source quotient is \(\pi^\vee\). Removing \(\eta_1^{-1}\) replaces its cuspidal inertia lines by those of \(\pi^\vee\otimes\eta_1\), up to unramified powers and the inverse-transpose automorphism used in the alternate rectangular-matrix convention. Those lines are

\[
                         \bar\eta_1,\qquad
                         \bar\eta_1\bar\eta_0^{-1}.
 \tag{TG17}
\]

Both are nontrivial. Simultaneously inverting them under the alternate convention changes nothing. Uniqueness of cuspidal support therefore excludes the source quotient from every boundary stratum. No pole-free regularization is assumed.

There is an actual nonzero oscillator map to \(\pi^\vee\otimes\pi^+\) by LIFT19 and Mínguez's theta theorem. Its restriction to \(S(\Omega)\) cannot vanish: otherwise a nonzero map from a boundary stratum would follow by the finite rank filtration, contradicting the preceding argument. By(TG16) and irreducibility of \(\pi^+\), Schur's lemma makes the target-equivariant Hom space from the full-rank coinvariants to \(\pi^+\) one-dimensional. Thus the actual small-theta map restricts to a nonzero scalar multiple of T.

The coinvariant and tensor-Hom formulations agree here without a completed-tensor ambiguity. A compact open source subgroup fixing a given Schwartz vector makes the dependence on the input factor through its invariants. Admissibility of \(\pi\) makes that space finite dimensional, so the corresponding map has coefficients in the smooth contragredient tensor the target.

Every vector detected by T is therefore nonzero in the actual small theta quotient. One may normalize that local quotient map so its restriction equals(TG14). The scalar comparing this choice with a globally pinned automorphic realization is not computed by uniqueness. All norm and lattice statements retain that distinction.

\subsection{5. Exact survival of the prescribed column}\label{exact-survival-of-the-prescribed-column}

For \(v\in V_\lambda\), evaluate T at the identity. The raw factor on \(X_0g\) is supported on \(g\in J\); its last row is zero. The determinant terms give

\[
               \Phi(X_0g)\eta_1(\det g)^{-1}
                                      =\lambda(g)^{-1}.
 \tag{TG18}
\]

The determinant has absolute value one and \(\pi(g)v=\lambda(g)v\), so the integrand in(TG13) is exactly v. On \(J^\sharp\) the restricted character gives the identical calculation. Hence

\[
 \boxed{\quad
 {\cal T}(\Phi,v)(1)=v_Jv,\qquad
 {\cal T}(\Phi^\sharp,v)(1)=v_{J^\sharp}v
                   \quad(v\in V_\lambda).\quad}
 \tag{TG19}
\]

Both fixed raw kernels are injective on the entire type space. Every nonzero prior-image vector in this prescribed local column survives. This does not assert survival for arbitrary vectors outside the column. Without the type condition, identity evaluation is exactly the compact character projector: T(Phi,v)(1)=v\_J e\_(J,lambda)v, and similarly for the sharp subgroup. Thus(TG19) concerns a fixed raw function, not a Schwartz choice depending on v.

\subsection{6. Exact compact-model norms}\label{exact-compact-model-norms}

The actual automorphic constituent \(\pi\) is unitary. Give \(\pi^+\) its normalized induction metric on \(K_N=\mathrm{GL}_N(O)\), retaining the chosen unitary metric on \(\pi\). Put

\[
 {\cal J}_{b,d}(q)=[\mathrm{GL}_N(\kappa):P_{(b,d)}(\kappa)]
     =\frac{\prod_{i=1}^{N}(q^i-1)}
        {\prod_{i=1}^{b}(q^i-1)\prod_{i=1}^{d}(q^i-1)}.
 \tag{TG20}
\]

The sections are target \(K_N^1\)-fixed. Their finite parabolic quotient points have measure \(1/{\cal J}_{b,d}\). A nonzero value can occur only when the residue source b-plane is a graph

\[
 L_z=\left\{\begin{pmatrix}x\\0\\zx\end{pmatrix}:x\in\kappa^b\right\},
                    \qquad z\in\kappa^{1\times b},
 \tag{TG21}
\]

where the zero has d-1 coordinates. The raw support requires the plane to lie in the space spanned by the first b and last coordinates and project invertibly onto the first b, giving exactly these graphs. The reduction of a compact frame has full rank, so a contributing g is integral and invertible modulo q; no other finite plane can contribute. For a graph choose the integral representative \(h_z=I_N-e_Nz\), with z extended by zeros on the last d coordinates. Then \(h_z^{-1}X_0=[I_b;0;z]\).

For \(v\in V_\lambda\), all determinant and representation factors cancel as in(TG18). The finite averages are

\[
 \begin{split}
 {\cal T}(\Phi,v)(h_z)
  &=v_J\left(|H|^{-1}\sum_{g\in H}\psi_0(-zge_k)\right)v,\\
 {\cal T}(\Phi^\sharp,v)(h_z)
  &=v_{J^\sharp}\psi_0(-z_k)v.
 \end{split}
 \tag{TG22}
\]

For the sharp subgroup \(ge_k=e_k\). For H, the first k entries of \(ge_k\) vary independently over \(\kappa\), its k-th entry is one, and its subsequent entries are zero. Additive orthogonality therefore gives

\[
 {\cal T}(\Phi,v)(h_z)=
 \begin{cases}
    v_J\psi_0(-z_k)v,&z_0=\cdots=z_{k-1}=0,\\
    0,&\text{otherwise}.
 \end{cases}
 \tag{TG23}
\]

There are exactly \(q^{b-k}\) nonzero finite quotient values in the nonsharp case and \(q^b\) in the sharp case. Hence for every \(v\in V_\lambda\),

\[
 \boxed{
 \begin{split}
 \|{\cal T}(\Phi,v)\|^2
    &=\frac{q^{b-k}v_J^2}{{\cal J}_{b,d}(q)}\|v\|^2,\\
 \|{\cal T}(\Phi^\sharp,v)\|^2
    &=\frac{q^b v_{J^\sharp}^2}{{\cal J}_{b,d}(q)}\|v\|^2.
 \end{split}}
 \tag{TG24}
\]

Their squared constants have ratio \(q^{-k}\), the same as the raw Schwartz norms(TG9). After unit-normalizing the finite raw input, the two maps have the same squared scaling constant:

\[
 \boxed{
 {\cal C}_q=
 \frac{q^{Nb-k}h_q}{|\mathrm{GL}_b(\kappa)|^2{\cal J}_{b,d}(q)}
 =\frac{q^{Nb}h_q^\sharp}
        {|\mathrm{GL}_b(\kappa)|^2{\cal J}_{b,d}(q)}.}
 \tag{TG25}
\]

These are local induced-model metric identities. The real selected Schwartz norm, global mass, and scalar comparing this model to an automorphic row are separate. The finite sharp modification alone does not improve this unit-input scaling. For the primitive LIFT19(T18)--(T20) vector, its natural metric is \(\|v\|^2=q^{k(k+1)/2}/{\cal J}_{k,d}(q)\); it may be substituted into(TG24) directly.

\subsection{7. Exact content in the natural induced lattices}\label{exact-content-in-the-natural-induced-lattices}

Let R be the valuation ring at a place above11 of a coefficient extension containing the additive and tame character values. Use the natural finite models supplied by LIFT19:

\[
 \begin{split}
 L_1&=\operatorname{Ind}_{P_{(k,d)}(\kappa)}
                           ^{\mathrm{GL}_b(\kappa)}
       (\mathrm{St}_{k,R}\boxtimes\bar\eta_0\det_d),\\
 L_2&=\operatorname{Ind}_{P_{(b,d)}(\kappa)}
                           ^{\mathrm{GL}_N(\kappa)}
       (L_1\boxtimes\bar\eta_1\det_d).
 \end{split}
 \tag{TG26}
\]

Fix one finite Steinberg embedding into \(\pi_0^{K_k^1}\). These are the selected finite submodules of the actual compact models obtained from that embedding, with specified integral inducing lattices. They need not exhaust all principal-congruence invariants of either constituent. Their comparison with a global arithmetic row lattice is extra data. Put \(L_\lambda=L_1\cap V_\lambda\). It is saturated in \(L_1\): its defining equations are the kernels of R-linear maps \(\pi(h)-\lambda(h)\) into a torsion-free module.

Divide T by \(v_J\) or \(v_{J^\sharp}\). Equations(TG22)--(TG23) show that every finite quotient value is zero or a root of unity times v. The identity value is exactly v; parabolic transition maps preserve the integral lattices. Consequently

\[
 v_J^{-1}{\cal T}(\Phi,-),\quad
 v_{J^\sharp}^{-1}{\cal T}(\Phi^\sharp,-):
                         L_\lambda\longrightarrow L_2
 \tag{TG27}
\]

are primitive saturated injections. To prove saturation, suppose a uniformizer divides the image in \(L_2\). Identity evaluation shows it divides v in \(L_1\); saturation of \(L_\lambda\) puts the divided vector in the domain. Torsion-freeness then identifies its image with the divided image. This does not assume existence of an integral H-projector.

If \(c_R(v)\) denotes the content valuation in the indicated free R-module, then

\[
 \begin{split}
 c_R({\cal T}(\Phi,v))&=c_R(v)+v_{11}(v_J),\\
 c_R({\cal T}(\Phi^\sharp,v))&=c_R(v)+v_{11}(v_{J^\sharp}),\\
 v_{11}(v_J)=v_{11}(v_{J^\sharp})
              &=-\sum_{i=t+1}^{b}v_{11}(q^i-1).
 \end{split}
 \tag{TG28}
\]

Here \(v_{11}(11)=1\), and q is an11-unit. The matrix-coefficient normalization and natural induced lattices are fixed in this statement. It cannot be identified with the global DP44 gain \(g_j\) without the missing row-lattice and scalar comparison.

\subsection{8. Global implication and remaining exact gates}\label{global-implication-and-remaining-exact-gates}

The local result applies at every degree-one selected G place of a repaired host with distinct nontrivial literal tame characters. It is independent of the number of real embeddings. It therefore applies to the actual first constituents from K031/DO034, and to later repaired hosts only after their actual first-family and character hypotheses are supplied.

The fixed terminal raw G factor cannot annihilate a nonzero previous-image vector in \(V_\lambda\). Finite tensor products of these column spaces satisfy the product versions of(TG19), (TG24), and(TG28). Thus this particular G-place test is no longer an unspecified local nonvanishing condition on the m=2 route.

This does not put a particular global first vector into every required column by an integral operator of known height. LIFT19 supplies occurrence and an explicit natural vector; the comparison with the exact global previous-image lattice remains necessary. Nor does this argument treat the fixed T/Q terminal factors, the permitted29 family at the second edge, the final global row, or the normalization of all local quotient maps by one global theta period. These are separate gates before a full DP44 changed-data comparison or a shorter source-compatible chain is claimed.

\subsection{9. Sources and proof status}\label{sources-and-proof-status-1}

Mínguez, A. (2008). Correspondance de Howe explicite: Paires duales de type II. Annales scientifiques de l'École Normale Supérieure, 41(5), 717--741. \url{https://www.numdam.org/item/ASENS_2008_4_41_5_717_0/} . Theorem1, p.718; action conventions in section1.4, pp.721--722; exact rank filtration, Corollary2.3, and Theorem2.4 with its restriction argument, pp.722--724. Read directly in the primary PDF. The theorem is used as a published input.

Other published representation inputs are those already checked in LIFT19: Casselman(1980), Proposition2.6; Zelevinsky(1980), Propositions8.5--8.6; and the finite Steinberg theorem. They are not counted again as new results.

The new calculations are the raw terminal compact support and character, the Haar-normalized open-rank integral in the actual source convention, its value on the whole prescribed column, the graph-plane norm, and the saturated natural-lattice comparison. The author whole-proof reading has been completed, including the raw source predicates, all source and target covariance signs, full-rank coinvariants, the boundary restriction argument, every graph value, both norm normalizations, and saturation without an integral isotypic projector.

The exact diagnostic \nolinkurl{raw_terminal_g_certificate.py} passed two finite cases: q=5,k=1,t=1 with distinct nontrivial residue characters, and q=3,k=2,t=1 for the higher-k raw geometry. It checks1,958 raw determinant/input phase cancellations and412 full graph averages, the support counts, both Haar volumes, and equal unit-input metric scaling. The second case does not assert the distinct character hypothesis, and neither small case is an admissible full source Ranks datum. The general quotient and lattice proofs are independent of these finite checks. Exact check events are in \nolinkurl{lift_accounting_CHECK_EVENTS.tsv.}

\clearpage

\section{AM-DO-035: all first coupled raw factors, their compact filters, and the remaining 29-place gate}\label{am-do-035-all-first-coupled-raw-factors-their-compact-filters-and-the-remaining-29-place-gate}

Author: GPT-6 Astra Pro, ALP\_MACHINE \nolinkurl{/root/definite_oscillator}. 22 November 2026 UTC. Inherited reasoning effort is not separately exposed.

This completes the local G/T/Q compatibility calculation for the first coupled kernel when \nolinkurl{m=2}. It uses the actual \nolinkurl{ybSet} class, the actual common splitting character, and an initial source input satisfying the prescribed G, T and Q conditions. In particular, the Q determinant character cancels; it is not an obstruction to \nolinkurl{zero_Q}. A nontrivial T diamond character gives the initial ramification condition and an explicit nonzero compact filter.

Combined with ALP-LIFT-20 and AM-DO-030/034, this gives a nonzero first lift with \textbf{all raw G/T/Q factors retained}, after one allowed auxiliary 29-place Heisenberg translate and a source level refinement there. It does not identify that translated kernel or refined input with the literal unmodified \nolinkurl{XiCpl1}/\nolinkurl{MfV}. Section 8 states the remaining equation and proves the precise limitation of absorbing it by lattice moves.

The local identities hold for every source-admissible T prime and twisting tuple. The global corollary assumes fixed such data, with all its auxiliary primes away11 and29. This file does not certify an explicit T prime or instantiate the complete \nolinkurl{Primes.Ok} record.

\subsection{1. Source locations and conventions}\label{source-locations-and-conventions}

All source paths below are in the recovered original commit \nolinkurl{7cbf7a0c32fca1badf05931879a37c3f8a2b03b4}. Abbreviate \nolinkurl{Formalism/LocalToGlobal/Columns/FiniteSumNonzeroLargeDepth.lean} by FS. These are source readings, not a Lean execution.

\begin{longtable}[]{@{}
  >{\raggedright\arraybackslash}p{(\columnwidth - 2\tabcolsep) * \real{0.5000}}
  >{\raggedright\arraybackslash}p{(\columnwidth - 2\tabcolsep) * \real{0.5000}}@{}}
\toprule\noalign{}
\begin{minipage}[b]{\linewidth}\raggedright
Source locator
\end{minipage} & \begin{minipage}[b]{\linewidth}\raggedright
Clause used
\end{minipage} \\
\midrule\noalign{}
\endhead
\bottomrule\noalign{}
\endlastfoot
FS594--598, 600--608, 636--641 & \nolinkurl{homAt}, the two scaled modules, and the actual \nolinkurl{kerOf} sum \\
FS646--655 & \nolinkurl{phiT} and the chosen lattice/residue frames \nolinkurl{liftFr} \\
FS684--694, 714--719 & \nolinkurl{entB} and the raw Q factor \nolinkurl{PhiQ} \\
FS780--810, 981--986 & raw G/GS factors, \nolinkurl{joinCol}, and \nolinkurl{PsiAt} \\
FS999--1010 & actual uniformizer and residue tying of the common character \\
FS1019--1023 & initial index, \nolinkurl{hal=1}, prescribed G type and full trivial Q level \\
FS1038--1061 & depths, support and representative-period conditions \\
FS1103--1138 & actual frames, \nolinkurl{Sall}, \nolinkurl{cAof}, \nolinkurl{cBof}, \nolinkurl{ybSet}, \nolinkurl{repY}, and \nolinkurl{XiCpl1} \\
FS1246--1259 & inverse source and direct target determinant characters \\
\nolinkurl{Columns/Adelic.lean}139--145 & mirabolic source level at T \\
\nolinkurl{Columns/RecLevel.lean}50--52 & the adelic source character is the inverse column character \\
\nolinkurl{Columns/ScaledRamifiedLattice.lean}81--94, 160--169 & diamond operators, their sum \nolinkurl{NmG}, and \nolinkurl{MramV} \\
\nolinkurl{Development/AbelianVarietyTheory/Definitions/LatticeFrames/LevChiEqOne.lean}61, 100, 114--117 & actual common label \nolinkurl{mu0_j * xi(a)}, raw datum and scalar tie \\
\nolinkurl{Columns/AbelianHostFieldConditions.lean}314--334 & reference characters unramified away G, and their actual tying \\
\nolinkurl{Formalism/LocalToGlobal/Divisibility/Tuples.lean}, \nolinkurl{GoodPl}, \nolinkurl{tupGrp}, \nolinkurl{Qprime}, \nolinkurl{resChar} & twist ramification, tame residue characters and degree-one Q places \\
\end{longtable}

Put

\[
 m=2,\qquad a=k,\qquad b=k+t+1,\qquad m_2=0,
 \qquad r=a+1,\quad r_S=0,\quad c_0=1,\quad c_1=0.
 \tag{DO35.1}
\]

Take the base decorated lattices at the first two ranks. Their underlying lattices are the standard integral lattices. Their T functionals are nonzero residue covectors; the base choice is the last-coordinate functional. Their Q frames are the identity. At G the selected \nolinkurl{liftFr} need not literally be the identity, but preserves the integral lattice and induces the prescribed identity residue frame. Therefore all residue formulas below remain valid in those exact frames.

At a relative split degree-one place write \nolinkurl{O} for the local ring, \nolinkurl{q} for its rational uniformizer and residue cardinality, and \nolinkurl{kappa=F_q}. Additive matrix Haar gives \nolinkurl{M_(b,a)(O)} volume one. The genuine source action, with the actual common splitting character having local component \nolinkurl{eta}, is

\[
 [\omega(g)f](X)
   =\eta(\det g)^{-1}|\det g|^{b/2}f(Xg).
 \tag{DO35.2}
\]

This is the split formula checked in AM-DO-020/024 and ALP-LIFT-20. On a compact subgroup the absolute-value factor is one. The formula uses a genuine left representation; replacing \nolinkurl{Xg} by \nolinkurl{Xg^(-1)} without changing the other conventions would reverse the calculation below.

Let \nolinkurl{xi} denote the finite-order twisting tuple and \nolinkurl{mu_0} the first reference character. The actual character in (DO35.2) is

\[
                          \eta=\mu_0\xi.
 \tag{DO35.3}
\]

This is exactly \nolinkurl{TwDat.Is.chr} and \nolinkurl{PinX.mv} in the cited source. At G the tuple is unramified and the reference character is prescribed. At T both factors are unramified. At Q the reference factor is unramified, whereas \nolinkurl{xi} has its nontrivial tame residue character. The Q set belongs to the tuple, not to the reference character.

\subsection{2. One admissible global target class for all G and Q places}\label{one-admissible-global-target-class-for-all-g-and-q-places}

Write \nolinkurl{m_c=cst.M} to distinguish this fixed positive integer from the auxiliary isogeny degree \nolinkurl{M=29^e}. At every selected G place prescribe

\[
             y_P\equiv q^{-1}e_a\pmod{O^b};
 \tag{DO35.4}
\]

at every Q place prescribe

\[
             y_P\equiv e_a\pmod{q^{m_c}O^b}.
 \tag{DO35.5}
\]

Here \nolinkurl{e_a} is zero-indexed: it is the first coordinate immediately following the first \nolinkurl{a} target coordinates. These prescriptions are made in the target frames in FS1103--1106. They are compatible with the base frames described above.

\textbf{Claim.} A single rational target vector \nolinkurl{y} defines an actual class \nolinkurl{Y} in \nolinkurl{ybSet} satisfying every prescription (DO35.4)--(DO35.5).

\textbf{Proof.} The selected G and Q prime ideals are pairwise distinct; the Q set is disjoint from G and T by the tuple's good-place condition. Multiply by \nolinkurl{D_aux=g_0 g_1}. At each G prime the required congruence on \nolinkurl{D_aux y} is integral modulo that prime; at each Q prime it is an integral congruence modulo its \nolinkurl{m_c}-th power, since \nolinkurl{D_aux} is a unit there. The Chinese remainder theorem in \nolinkurl{O_K}, coordinate by coordinate, gives an integral vector satisfying all these congruences. Divide it by \nolinkurl{D_aux}. This gives a rational vector with the stated local conditions. In the actual selected G frames the same proof uses their integral coordinate maps; their reductions are the specified residue frames. In particular the congruence \nolinkurl{q y=e_a mod q} is unchanged by a frame congruent to the identity.

The vector is in the module defining \nolinkurl{ybSet}: its G valuation is at least \nolinkurl{-1}, and its Q valuation is at least zero, hence at least \nolinkurl{-m_c}. Set

\[
 Y=y+\operatorname{locLat}
       \bigl(N_1;\ e_G=0,\ e_Q=m_c\bigr).
 \tag{DO35.6}
\]

This is the literal set in FS1128--1130. Every representative of \nolinkurl{Y}, including \nolinkurl{repY Y}, has the same prescribed residues. At G ALP-LIFT-20 proves the exact representative invariance; at Q it follows from the determinant and divisibility conditions below (or from the actual \nolinkurl{psiInv} clause). No independent local \nolinkurl{y} is being substituted for the one global class. ∎

T imposes no condition on \nolinkurl{y}: its factor is \nolinkurl{phiT(X)}. Its distinct prime support makes its incidence conditions compatible with this same G/Q class. For example, one may choose a rational \nolinkurl{X} with \nolinkurl{X=X_0 mod G}, the Q block residues in Section 3, and the T last row condition in Section 4, again by the Chinese remainder theorem. Thus the combined finite raw function is nonzero before any theta or compact averaging assertion is made.

\subsection{3. The exact Q factor and cancellation of its compact character}\label{the-exact-q-factor-and-cancellation-of-its-compact-character}

Let

\[
 X=\begin{pmatrix}A\\ r\\ B\end{pmatrix},
 \quad A\in M_a(K_P),\ r\in M_{1,a}(K_P),\
 B\in M_{t,a}(K_P),
 \qquad \chi=\operatorname{resChar}(\xi,P).
 \tag{DO35.7}
\]

At Q the exponents in \nolinkurl{entB} are all zero. Both group arguments in \nolinkurl{PsiAt} are one, by \nolinkurl{zero_h} and its definition. For the class (DO35.5), the top \nolinkurl{(a+1)} square block of \nolinkurl{[X,y]} has last column \nolinkurl{e_a} modulo \nolinkurl{q}, and the lower \nolinkurl{t} rows of that column vanish modulo \nolinkurl{q}. Consequently the actual source definition gives

\[
 f_Q(X)=
 \chi(\det\bar A)\,
 1_{\{A\in\operatorname{GL}_a(O),\ r\in O^a,
                   \ B\in qM_{t,a}(O)\}}.
 \tag{DO35.8}
\]

The determinant character is extended by zero on nonunits, exactly as \nolinkurl{unitV} in FS. This proves (DO35.8) also when the top block is singular. A change \nolinkurl{y -> y+q^(m_c)z}, with integral \nolinkurl{z} and \nolinkurl{m_c>=1}, changes neither reduction nor lower-row divisibility.

For every \nolinkurl{g in GL_a(O)} the support is preserved and

\[
                     f_Q(Xg)=\chi(\det\bar g)f_Q(X).
 \tag{DO35.9}
\]

But the actual label (DO35.3) has exactly the residue character \nolinkurl{chi} at Q. Substituting (DO35.9) into (DO35.2) proves

\[
                  \boxed{\omega(g)f_Q=f_Q\quad(g\in\operatorname{GL}_a(O)).}
 \tag{DO35.10}
\]

Thus the raw Q factor is fixed by the full source compact. This is precisely the initial condition \nolinkurl{zero_Q=LevQFam.top}, whose definition is \nolinkurl{H=top, lam=1} (FS409). Using only the reference label \nolinkurl{mu_0} in (DO35.2) would incorrectly leave an extra character. The inverse source residue character \nolinkurl{rhoS=chi^(-1)} in FS1246--1248 confirms the same cancellation independently.

Its exact additive norm is

\[
 \|f_Q\|^2
  =q^{-ta}\prod_{i=1}^{a}(1-q^{-i})=:\Gamma_Q.
 \tag{DO35.11}
\]

Indeed \nolinkurl{A} has invertible reduction with the displayed probability, \nolinkurl{r} is arbitrary integral, and \nolinkurl{B} has \nolinkurl{ta} entries divisible by \nolinkurl{q}. Every nonzero value has modulus one. The value at \nolinkurl{[I_a;0;0]} is one.

The function is supported on integral matrices and periodic modulo \nolinkurl{q}. Its finite Fourier expansion in the Heisenberg vacuum basis is therefore an actual finite combination at P. All denominators divide a power of \nolinkurl{q}; at a coefficient place over11 with \nolinkurl{q!=11}, at least one Fourier coefficient is a unit. This last assertion does not require a Gauss-sum estimate: Fourier inversion over a finite group of order prime to11 is an integral automorphism, and (DO35.8) has the unit value one. The free row \nolinkurl{r} forces the corresponding dual row to vanish, so its Fourier support has size at most \nolinkurl{q^(a(b-1))}. Parseval gives exactly (DO35.11).

\subsection{4. The T factor, source level and initial ramification condition}\label{the-t-factor-source-level-and-initial-ramification-condition}

Let \nolinkurl{u in kappa^a} and \nolinkurl{v in kappa^b} be the nonzero source and target residue covectors. In the base decorations they are the last-coordinate covectors. The exact factor \nolinkurl{phiT} is

\[
 f_{T,u}(X)
  =1_{M_{b,a}(O)}(X)\,1_{\{v\bar X=u\}}.
 \tag{DO35.12}
\]

Its norm is \nolinkurl{q^(-a)} because these are \nolinkurl{a} independent residue conditions. Put

\[
 H_u=\{g\in\operatorname{GL}_a(\kappa):ug=u\},\qquad
 P_u=\{g:ug=s(g)u\text{ for some }s(g)\in\kappa^\times\}.
 \tag{DO35.13}
\]

The quotient \nolinkurl{P_u/H_u} is \nolinkurl{kappa^times}, and

\[
 d_T=[\operatorname{GL}_a(\kappa):P_u]
       =\frac{q^a-1}{q-1}.
 \tag{DO35.14}
\]

Choose any nontrivial unitary residue character \nolinkurl{nu} of \nolinkurl{kappa^times}. The rational T prime is odd in the source, so such a character exists. Work over a coefficient extension containing its values. Define

\[
 f_{T,\nu}=\sum_{s\in\kappa^\times}\nu(s)^{-1}f_{T,su}.
 \tag{DO35.15}
\]

The supports are disjoint, so this is nonzero and

\[
                    \|f_{T,\nu}\|^2=(q-1)q^{-a}.
 \tag{DO35.16}
\]

At T the character (DO35.3) is unramified. Equation (DO35.2) therefore gives, for a lift of \nolinkurl{g in P_u},

\[
 \omega(g)f_{T,u}=f_{T,s(g)^{-1}u}.
 \tag{DO35.17}
\]

With maximal-compact Haar of volume one, extend \nolinkurl{nu(s(g))} by zero from the inverse image of \nolinkurl{P_u}. Equal-size fibres of \nolinkurl{s:P_u -> kappa^times} prove the exact compact filter

\[
 \boxed{\int_{K_T}1_{P_u}(g)\nu(s(g))\,
                  \omega(g)f_{T,u}\,dg
       =\frac{f_{T,\nu}}{d_T(q-1)}
       =\frac{f_{T,\nu}}{q^a-1}.}
 \tag{DO35.18}
\]

The scalar function in this integral has absolute value one on \nolinkurl{P_u}, zero outside, norm squared \nolinkurl{1/d_T}, and is right \nolinkurl{H_u} invariant. This matches the literal T mirabolic level in \nolinkurl{Adelic.lean}139--145. Right multiplication by a lift with \nolinkurl{s(g)=c} multiplies it by \nolinkurl{nu(c)}. Since \nolinkurl{nu} is nontrivial,

\[
                    \sum_{c\in\kappa^\times}\nu(c)=0.
 \tag{DO35.19}
\]

Thus its T diamond sum vanishes. This is exactly the defining local condition in \nolinkurl{NmG} and \nolinkurl{ColRamV}, not a cuspidality claim. The source diamond operation multiplies the decorated covector by \nolinkurl{c} (\nolinkurl{StdLatUnimodular.lean}113--116); changing the identification of right translation with that operation replaces \nolinkurl{nu} by its inverse, and leaves (DO35.19) unchanged. Its scalar action on the algebraic coefficient is the identity in \nolinkurl{diamondOperatorOnForms'}.

For completeness, the finite Fourier coefficients can also be computed exactly. Fix a nontrivial residue additive character \nolinkurl{psi}. The coefficient indexed by \nolinkurl{ell in kappa^a} on the rank-one matrix frequency \nolinkurl{v^t ell} is

\[
 b_\ell=q^{-a}\sum_{s\ne0}\nu(s)^{-1}\psi(-s\,\ell u^t).
 \tag{DO35.20}
\]

It vanishes when \nolinkurl{ell u^t=0}; otherwise it is a nonzero Gauss sum divided by \nolinkurl{q^a}, of complex modulus \nolinkurl{q^(1/2-a)}. To see nonvanishing and the norm without importing a formula, put \nolinkurl{G(c)=sum_(s!=0) nu(s)^(-1) psi(cs)}. For \nolinkurl{c!=0}, substitution \nolinkurl{s -> c^(-1)s} shows all \nolinkurl{|G(c)|} agree. Orthogonality of \nolinkurl{psi} gives \nolinkurl{sum_c |G(c)|^2=q(q-1)} while \nolinkurl{G(0)=0}, hence \nolinkurl{|G(c)|^2=q}. The product of the Gauss sum and its conjugate is \nolinkurl{q}, so these coefficients are units at every place over11 when \nolinkurl{q!=11}. Their number is exactly

\[
                      s_T=(q-1)q^{a-1}.
 \tag{DO35.21}
\]

Parseval recovers (DO35.16). Their total complex absolute value is \nolinkurl{(q-1)/sqrt(q)}. No division by \nolinkurl{q-1} has been made in the auxiliary vector (DO35.15); the division belongs to the actual compact filter (DO35.18) and must be charged there.

\subsection{5. The complete auxiliary vector and exact fixed costs}\label{the-complete-auxiliary-vector-and-exact-fixed-costs}

At each G place use the Steinberg vector \nolinkurl{f_St} of AM-DO-033, with its determinant factor for the actual reference label. ALP-LIFT-20 proves, for the actual raw factor \nolinkurl{f_U},

\[
 \int_{K_G}\overline{m_G(g)}\omega(g)f_U\,dg
          =\frac{\alpha_G}{d_G}f_{St},
 \quad d_G=q^{a(a-1)/2},\quad
 \alpha_G=\frac{d_G^2}{|\operatorname{GL}_a(\kappa)|}.
 \tag{DO35.22}
\]

Here \nolinkurl{m_G} is the matrix coefficient of a unit generic vector in the finite Steinberg representation. Its inverse-unipotent source type is the literal initial \nolinkurl{colLev_0} character. The raw factor remains in the kernel on the left of (DO35.22).

Form the actual finite auxiliary vector

\[
 \phi_{\rm aux}
   =\bigotimes_{P\in G}f_{St,P}
       \otimes\bigotimes_{P\in T}f_{T,\nu_P}
       \otimes\bigotimes_{P\in Q}f_{Q,P}
       \otimes\bigotimes_{P\notin G\cup T\cup Q}\phi_{0,P}.
 \tag{DO35.23}
\]

Its finite norm squared is the exact product

\[
 \Gamma_{\rm aux}
  =\prod_{P\in G}\left(q_P^{-ab}\frac{|\operatorname{GL}_a(\kappa_P)|}{d_{G,P}}\right)
    \prod_{P\in T}(q_P-1)q_P^{-a}
    \prod_{P\in Q}\Gamma_{Q,P}.
 \tag{DO35.24}
\]

It is nonzero. Each factor is an actual finite Heisenberg combination of the local vacuum; their product is one such combination with pairwise distinct translated classes after like terms are combined. The previous arguments and DO033 show that the coefficient array is primitive at11. A valid bound for its number of terms is

\[
 s\le
  \prod_{P\in G}q_P^{ab}
  \prod_{P\in T}(q_P-1)q_P^{a-1}
  \prod_{P\in Q}q_P^{a(b-1)}.
 \tag{DO35.25}
\]

Every auxiliary rational prime here is assumed different from 11 and29. This is the exact hypothesis needed for the integral prime-to-11 translated-line comparisons in DO030. The empty Q set, supplied by the reference tuple, is explicitly allowed.

The raw finite vector with the unprojected G and T factors has norm squared

\[
 \Gamma_{\rm raw}
  =\prod_{P\in G}q_P^{-ab}d_{G,P}
    \prod_{P\in T}q_P^{-a}
    \prod_{P\in Q}\Gamma_{Q,P}.
 \tag{DO35.26}
\]

Unlike the selected G vector, the raw G norm has no GL-order denominator. If Q is empty, its finite norm has zero11 valuation. For nonempty Q the explicit factor (DO35.11) retains its possible 11 valuation. One must not equate primitive coefficients with unit Hilbert norm.

\subsection{6. Global compact extraction on the actual degree-seven host}\label{global-compact-extraction-on-the-actual-degree-seven-host}

Use the genuine host and selected highest polynomial of AM-DO-034, with the actual rank and stage characters supplied in AM-K031. In particular \nolinkurl{f=42}, \nolinkurl{a=42}, \nolinkurl{b=1802} is its \nolinkurl{m=2,t=1759} example. Keep the standard-lattice coefficient \nolinkurl{xi_K=xi_0 pi_B^(-6)} and its actual eigen-coordinate scalar. Let \nolinkurl{D} be the total selected degree and \nolinkurl{B_nu} the factorial valuation bound in DO034. Let

\[
 M=29^e\ge3s(D+1)-1.
 \tag{DO35.27}
\]

The simultaneous jet theorem DO030, in the proved host form DO034, applied to the nonzero primitive vector (DO35.23), gives one actual permitted half-orbit translate \nolinkurl{W_(x_lambda)} with

\[
 \theta_{P_\nu\otimes W_{x_\lambda}\phi_{\rm aux}}(1,1)
      =\mathcal P\,a_\lambda\ne0,\qquad
 v_{11}(a_\lambda)\le B_\nu,
 \quad \mathcal P=C^{14ab}.
 \tag{DO35.28}
\]

All gauges, the ordinary product line and the fixed root-of-unity translation phases are those already proved in DO030/034. The factor \nolinkurl{C} is the same analytic period, not a newly normalized value of the raw T or Q function.

Define \nolinkurl{K^e} to be the maximal source compact outside29 and the rational principal congruence subgroup of level \nolinkurl{29^e} at the six F places above29. Its exact index in the maximal compact is

\[
 I_a^K(M)=M^{42a^2}\prod_{i=1}^{a}(1-29^{-i})^6.
 \tag{DO35.29}
\]

The ramification index is seven and there are six residue places; these must not be replaced by forty-two unramified residue factors. DO034 proves that the finite rational stabilizer of this compact is trivial. The congruence fixes every relevant half-orbit vector; the other auxiliary places are disjoint from29.

On the identity finite class define the actual automorphic source function by

\[
 F(g_\infty,k)=
   \overline{m_\infty(g_\infty)}
   \prod_{P\in G}\overline{m_{G,P}(k_P)}
   \prod_{P\in T}1_{P_u}(k_P)\nu_P(s(k_P)),
 \quad k\in K^e,
 \tag{DO35.30}
\]

and set it to zero off that class. Here \nolinkurl{m_infty} is the unit highest source matrix coefficient used in DO034, and \nolinkurl{d_rho} is the dimension of that compact real representation. Triviality of the rational stabilizer proves this definition is well defined on the automorphic quotient, rather than only on a local chart. It is smooth at finite places, compact-finite at infinity, and square-integrable because the source unitary quotient is compact.

At G it has the initial inverse generic character; at T it has the literal mirabolic invariance and zero diamond sum from Section 4; at Q it is fixed by the full compact. These statements are about this actual global function. They do not follow merely from nonzero local factors. Its only extra source level is the displayed refinement at29. With maximal finite compact volumes one and compact real volume one, direct orthogonality gives

\[
 \|F\|^2
   =\frac1{I_a^K(M)d_\rho d_{G,\rm all}d_{T,\rm all}},
 \quad \|F\|_\infty=1,
 \tag{DO35.31}
\]

where \nolinkurl{d_(G,all)=product_G d_G} and \nolinkurl{d_(T,all)=product_T d_T}. No class-number-one assertion is used.

Now keep the \textbf{raw} G, T and Q factors in the theta kernel, and apply the same permitted29 translate to its vacuum there. The compact integrations (DO35.10), (DO35.18), (DO35.22) commute, and the real integration is the selected matrix-coefficient projection of DO034. Thus the literal bilinear lift has value

\[
 \boxed{
 \Theta_{\rm raw,G/T/Q}^{\rm bil}(F)(1)
   =\frac{A_G\,\mathcal P\,a_\lambda}
      {I_a^K(M)d_\rho d_{G,\rm all}
          \prod_{P\in T}(q_P^a-1)}\ne0,
 \qquad A_G=\prod_{P\in G}\alpha_{G,P}.}
 \tag{DO35.32}
\]

This is an equality for actual global automorphic and Schwartz inputs. It retains the raw local functions and the common character (DO35.3). The local G denominator and the T averaging denominator are not suppressed in a claim of nonzero pairing.

If the source and raw finite Schwartz vector are separately normalized to norm one, divide (DO35.32) by

\[
 \sqrt{\Gamma_{\rm raw}}\,
 \bigl(I_a^K(M)d_\rho d_{G,\rm all}d_{T,\rm all}\bigr)^{-1/2}.
 \tag{DO35.33}
\]

If the real Schwartz factor is normalized as well, divide by its exact norm in DO034, including \nolinkurl{det(2Y)^(-7ab)}, the hook factors, the eigen-coordinate metrics and \nolinkurl{|c_nu|^2}. No new real norm formula is being inferred from compact projection.

For fixed ranks and fixed G/T/Q data, (DO35.27) can be chosen with \nolinkurl{M<29(3s(D+1)-1)} (except the harmless \nolinkurl{e>=1} endpoint). Consequently the logarithm of (DO35.29) is \nolinkurl{O(log(D+1))}. All G/T/Q factors in (DO35.31)--(DO35.33) are fixed constants in that assertion. This does not remove the separate linear-degree coordinate denominator, Fock and factorial costs in DO034. If the T prime or the Q tuple varies with D, the displayed products, rather than this fixed-data simplification, apply.

\subsection{\texorpdfstring{7. What is and is not coupled to \nolinkurl{cst.M}, \nolinkurl{cAof} and \nolinkurl{cBof}}{7. What is and is not coupled to ,  and }}\label{what-is-and-is-not-coupled-to-cst.m-caof-and-cbof}

The first raw function at \nolinkurl{m2=0} contains no \nolinkurl{E1at} sum. Its G and Q local formulas are \nolinkurl{PsiAt} with both group arguments one. The integer \nolinkurl{m_c=cst.M} only controls the admitted Q target class and its period here. The choices (DO35.4)--(DO35.5) work for every \nolinkurl{m_c>=1}. The first-factor support is integral in X and periodic modulo q, so the first-factor clauses are compatible with

\[
 A\ge m_c,\quad B\ge m_c,\quad B>0,\quad
 D_G=1,\quad D_Q\ge\max(2m_c,3).
 \tag{DO35.34}
\]

For example \nolinkurl{m_c=A=B=1,D_Q=3} satisfies these displayed bounds. This is not an assertion that all group, character and later-row clauses of \nolinkurl{ChIdx.Ok} have been instantiated.

The exact integers entering the finite \nolinkurl{kerOf} expression are

\[
 c_A=\left(\prod_{P\in G\cup Q}q_P\right)^A,
 \qquad
 c_B=T(u)\left(\prod_{P\in G\cup Q}q_P\right)^B.
 \tag{DO35.35}
\]

Their exponents count the selected prime ideals; in particular a rational fully split G prime occurs f times. These integers are units at29 when all actual auxiliary places avoid29. Increasing \nolinkurl{cst.M}, A or B cannot insert a29 component into \nolinkurl{Sall} or into the actual half-position idempotent.

One must also avoid the converse false inference that every rational representative x in \nolinkurl{kerOf} is integral at29. \nolinkurl{homAt} in FS594 imposes conditions only at \nolinkurl{Sall}, so representatives may have denominators elsewhere. The finite quotient and the congruence/additivity identities are defined using that same localized module. Changing such a representative does not enlarge the actual idempotent support or change \nolinkurl{IsRatDatum}. The obstruction in the next section is about that actual finite input, not about the incidental denominator of a representative chosen in the displayed finite sum.

\subsection{8. The remaining literal datum and source-level equations}\label{the-remaining-literal-datum-and-source-level-equations}

For the degree-seven host, every29 place is absolutely ramified of index seven. A source Q place must be \nolinkurl{GoodPl}, which includes degree one and rational-uniformizer valuation one. Thus29 is not an eligible Q place. The source G primes are the large degree-one primes of K031, and the T conditions likewise require degree-one selected places and their explicit congruences. Hence the actual G/T/Q \nolinkurl{Sall} used here avoids29.

The literal idempotent predicate is \nolinkurl{IsIdemSum(A.e,Sall)} (LevChiEqOne62), and AM-DO-030 proved from \nolinkurl{IsIdemAt} that it has zero29 component. The finite input fixed by \nolinkurl{IsRatDatum} is therefore the raw G/T/Q vector with the standard vacuum at29. Equation (DO35.32) instead uses \nolinkurl{W_(x_lambda)} there, and its source function has the refined level \nolinkurl{K^e}. Neither difference is removed by Section 3's Q character cancellation.

There is a precise limitation on repair by actual lattice moves. At a relative split29 place, all source and target unitary actions in the Hom model preserve the position polarization. They act by an invertible linear change \nolinkurl{X -> h^(-1)Xg}, with a nonzero determinant-root scalar and positive Jacobian factor. A moved vacuum is consequently a nonzero scalar times the indicator of a linear lattice containing zero. The actual \nolinkurl{thMoved} family comes from those source and target lattice moves; its root variables in \nolinkurl{VacT} are determinant roots, not affine translation parameters. A nontrivial permitted pure position translate has indicator \nolinkurl{1_L(X+x)} with \nolinkurl{x notin L}; its support does not contain zero. Thus

\[
 \boxed{W_x\phi_0\ne c\,\omega(g,h)\phi_0
 \quad\text{for }x\notin L,\ c\ne0,}
 \tag{DO35.36}
\]

for every individual source/target lattice or frame move at that place. The statement includes nonintegral linear moves. Changing the selected class, the row frame, or a determinant root does not change this origin-support proof.

It does not exclude an equality of a single global theta value, a finite linear combination of group translates, nonvanishing after a compact projection, or a nonzero witness already having \nolinkurl{x_lambda=0}. These are different possible routes, and none is disproved by (DO35.36).

The remaining first-row assertion at maximal29 level can be written concretely. Let \nolinkurl{P_(29,0)} be compact averaging on the source29 factor, and let \nolinkurl{L_(G/T/Q,Pnu)} be the actual theta functional after the already computed G/T/Q and real filters. One needs, for an actual permitted global initial input at that level, a nonzero value of

\[
 L_{G/T/Q,P_\nu}\bigl(P_{29,0}\phi_{0,29}\bigr)
       =L_{G/T/Q,P_\nu}(\phi_{0,29}).
 \tag{DO35.37}
\]

The proved theorem supplies a nonzero value with \nolinkurl{W_(x_lambda) phi_(0,29)} and a smaller compact. A construction which removes the source stabilizer by a generic T decoration would allow a maximal29 class-supported input, but its integral against that translated kernel contains \nolinkurl{P_(29,0) W_(x_lambda) phi_(0,29)}. Nonzero unaveraged value does not imply nonzero averaged value. This is the exact projection gate; it is not settled by the finite nonvanishing of (DO35.23).

Likewise moving the auxiliary isogeny to another split ray prime does not alone solve (DO35.37): if that prime is outside \nolinkurl{Sall} the same issue recurs, while if it lies in Q its raw factor is the prescribed (DO35.8), not a freely selectable translate. The completed local identities provide the actual tests to use in a future local or global nonvanishing argument.

Even after such an analytic assertion is proved, identifying the source input with a specified integral \nolinkurl{MfL}/\nolinkurl{NlamC} row, or proving the literal \nolinkurl{PinCol}/\nolinkurl{MainX} mass equations, requires those actual lattice and scalar ties. The compact values and norms in this finding do not instantiate the full arithmetic records.

\subsection{9. Proof-check status}\label{proof-check-status}

Author whole-proof review and the separate finite exact diagnostic are recorded in \nolinkurl{definite_oscillator_CHECK_EVENTS.tsv} after the completed file has been read. No computational check is used as a substitute for the proofs above. No Lean build is claimed.

\clearpage

\section{AM-K033 --- explicit sharp selected content in the enlarged CM eigenbasis}\label{am-k033-explicit-sharp-selected-content-in-the-enlarged-cm-eigenbasis}

Author: GPT-6 Astra Pro, ALP\_MACHINE /root/kummer\_extension, 22 November 2026 UTC; inherited reasoning effort not separately exposed.

This finding identifies an explicit pair of CM labels on the actual \(J^{14}\) host, and on its weighted \(J^{56}\) extension, that retains the exact sharp primitive-kernel derivative. A fixed filler places this pair inside the full source-and-target selected Schur summand. The resulting image ideals in the actual permitted DP59 half-orbit have asymptotic slope \(1/1220\), with a bounded additive error before compact metric normalization.

The proof computes the full critical logarithmic form after the seven-copy eigenbasis change. It does not infer selected sharpness merely from an invertible change of the ambient tangent lattice. The accepted display and exponential theorems are used at their exact interfaces; their proofs are not rederived. The conclusion concerns the selected kernel lattice. A previously produced automorphic image or a separately imposed auxiliary finite G-vector can still impose a further restriction.

\subsection{1. Actual hosts, coordinates, and accepted primitive input}\label{actual-hosts-coordinates-and-accepted-primitive-input}

There are two cases:

\begin{longtable}[]{@{}
  >{\raggedleft\arraybackslash}p{(\columnwidth - 6\tabcolsep) * \real{0.2857}}
  >{\raggedright\arraybackslash}p{(\columnwidth - 6\tabcolsep) * \real{0.2143}}
  >{\raggedright\arraybackslash}p{(\columnwidth - 6\tabcolsep) * \real{0.2143}}
  >{\raggedleft\arraybackslash}p{(\columnwidth - 6\tabcolsep) * \real{0.2857}}@{}}
\toprule\noalign{}
\begin{minipage}[b]{\linewidth}\raggedleft
Parameter h
\end{minipage} & \begin{minipage}[b]{\linewidth}\raggedright
Host field
\end{minipage} & \begin{minipage}[b]{\linewidth}\raggedright
Principal variety
\end{minipage} & \begin{minipage}[b]{\linewidth}\raggedleft
CM type count f
\end{minipage} \\
\midrule\noalign{}
\endhead
\bottomrule\noalign{}
\endlastfoot
1 & \(K=E_0B\), \(E_0=\mathbb Q(\zeta_{28})\), \(B\subset\mathbb Q(\zeta_{29})\), \([B:\mathbb Q]=7\) & \(J^{14}\) & 42 \\
4 & \(K^*=K(\sqrt5,\sqrt{13})\) & \(J^{56}\) & 168 \\
\end{longtable}

The integral actions, principal product identifications and standard CM lattice coefficients are AM-K030 and AM-K032. Work at11 over a finite unramified complete coefficient DVR R large enough for the CM eigenvalues and torsion being used, with \(v(11)=1\). Each level can use a larger finite unramified extension; none of the bounds depends on its degree. The coefficient-field extension needed for an unrelated global etaVal is not being asserted unramified.

Use the K030 basis \(e_0,\ldots,e_6\) of \(\pi_B^{-3}\mathcal O_B\), with \[
 e_j=\sigma^j e_0,\qquad
 \sigma=(\zeta_{29}\mapsto\zeta_{29}^{\,2})|_B,\qquad
 U_{lj}=\sigma^l(e_j),\qquad U^{\mathsf t}U=UU^{\mathsf t}=I_7.
\tag{CS1}
\] Here indices l,j lie in \(\mathbb Z/7\mathbb Z\), and \(\pi_l=\sigma^l(\pi_B)\). All entries of U and its inverse and all \(\pi_l^{\pm1}\) are integral at11.

Let \(\mathcal R\) be a singleton when h=1. When h=4, let \(\mathcal R\) be the four embeddings of \(\mathbb Q(\sqrt5,\sqrt{13})\); for \(\rho=(\epsilon,\epsilon')\) put \[
 \delta_\rho=\delta_{1,\epsilon}\delta_{3,\epsilon'},
 \qquad
 u_a=\frac{a+\sqrt{a^2+4}}2,\quad
 \delta_a=\frac{u_a^{-1}}{\sqrt{a^2+4}},\quad a=1,3.
\tag{CS2}
\] For h=1 set \(\delta_\rho=1\). Every \(\delta_\rho\) is an11-unit.

The standard CM logarithmic coordinates will be denoted \[
            \zeta_{i,\varepsilon,l,\rho},
       \qquad i=1,2,3,\quad \varepsilon\in\{+,-\},
                \quad l\in\mathbb Z/7\mathbb Z,\quad\rho\in\mathcal R.
\tag{CS3}
\] For h=1, they are exactly \[
 \zeta_{i,\pm,l}=\pi_l^3\sum_{j=0}^{6}U_{lj}
                                  (u_{i,j}\pm i v_{i,j}).
\tag{CS4}
\] For h=4, apply in addition the two embedding matrices \[
 V_a=\begin{pmatrix}1&u_{a,+}\\1&u_{a,-}\end{pmatrix},
 \qquad
 V_a^{\mathsf t}\operatorname{diag}(\delta_{a,+},\delta_{a,-})V_a=I_2,
\tag{CS5}
\] on the four ordinary real tensor factors, as in K032(QH18--21). The same final \(\pi_l^3\) rescaling gives the standard lattice coordinates. All changes and their inverses are integral at11. They preserve the ordinary J label i.

For one ordinary J, the accepted DP49 decomposition and K018's actual source-gauge comparison give an integral normalized algebraic frame in which the reduced coefficient of a local section of level \(\kappa\) is \[
                   f\,U_0^\kappa\exp(-\kappa\beta).
\tag{CS6}
\] Here f is an arbitrary integral formal function, \(U_0\) is an integral unit, and \(\beta\) is the actual arbitrary-unit primitive logarithm. This formula is used on identical copies with the identical choice of ordinary tangent basis \(u_i\). Thus every copy has the same critical unit below. Choosing \(f=U_0^{-\kappa}\) is an actual integral formal-section choice. It does not assert that the fixed sigma section itself equals the pure kernel.

Put \[
                  s=\frac1{10},\qquad w=\frac{121}{1220},
                         c=s-w=\frac1{1220}.
\tag{CS7}
\] Choose a finite ramified scaling extension with \(v(\varpi)=w\), and put \[
                      \mu=\frac{\varpi^{122}}{11^{12}},
                           \qquad v(\mu)=s.
\tag{CS8}
\] Only this auxiliary Gauss-weight argument uses the ramified extension. The derivative values themselves remain in the original unramified coefficient ring.

\subsubsection{Exact accepted whole-series input}\label{exact-accepted-whole-series-input}

DP46(UI16--19), with h=3 and j=1, and the common-weight substitution of DP51 §6 give \[
 \mu^{-1}\beta(\varpi u_1,\varpi u_2,\varpi u_3)
        \text{ integral},\qquad
 \overline{\mu^{-1}\beta(\varpi u_1,\varpi u_2,\varpi u_3)}
                         =\theta\,u_3u_2^{121},
                   \quad \theta\ne0.
\tag{CS9}
\] This is the reduction of the \textbf{entire} scaled logarithm, not its first ordinary-degree term. To spell out the application: in DP46's contracted dual inverse, raising the later input \(y_2\) to the common weight w makes that input vanish in the reduction, since its profile weight is strictly smaller than w. The reduced inverse has \(Y_1=y_1\) and \(Y_2=-\bar a_1^{-1}y_1^{11}\); substitution into the complete minimum \(\bar a_2^{-1}S Y_2^{11}\) gives (CS9). The diagonal unit change from those display coordinates to the fixed ordinary CM eigenforms changes only the nonzero unit \(\theta\). Invariant integral frame changes are strictly above this critical form by DP46 §5. This is a direct use of the accepted full inverse-map and initial-form statements.

For an integral unit parameter \(\lambda\), the right side of (CS9) is multiplied by \(\bar\lambda\). In the actual tensor-level application \(\lambda=-\kappa\), where \(\kappa=M^2\), \(M=29^e\); hence it is always a unit.

\subsection{2. Frobenius computes the new critical pairs}\label{frobenius-computes-the-new-critical-pairs}

Bars denote residue-field reduction at11. For the unramified constants this is reduction modulo11; for the auxiliary ramified scaling extension it means reduction modulo its maximal ideal. On B, arithmetic Frobenius is the class of11 in \((\mathbb Z/29\mathbb Z)^\times/H\), where \(H=\{1,12,17,28\}\). One has \(11\in2^4H\) and \(121\equiv5\in2H\pmod{29}\). Therefore \[
                   \overline{U_{lj}}^{\,121}
                            =\overline{U_{l+1,j}},\qquad
                   \bar\pi_l^{\,121}=\bar\pi_{l+1}.
\tag{CS10}
\] This follows for the denominators29 as well because they are11-units. Equivalently, if P shifts the row l to l+1, then \(\bar U^{(121)}=P\bar U\).

On the two quadratic real fields, Frobenius squared is the identity, since their Galois groups have exponent2 and11 is unramified. Consequently \[
                   \bar V_a^{(121)}=\bar V_a.
\tag{CS11}
\] Also \(i^{121}=i\) modulo11.

Let \(\beta_{\rm prod}\) be the sum of the primitive logarithms on all14h ordinary J factors at one fixed matrix entry of a Hom variety. Since the principal polarization and the chosen invariant frame are products, this is the logarithm of the actual product primitive kernel in (CS6).

\subsubsection{Theorem CS12. Full critical form in standard CM coordinates}\label{theorem-cs12.-full-critical-form-in-standard-cm-coordinates}

For every unit \(\lambda\), in the coordinates (CS3), \[
 \boxed{\ 
 \overline{\mu^{-1}\lambda
       \beta_{\rm prod}(\varpi\boldsymbol\zeta)}
 =
 \frac{\bar\lambda\theta}{2}
 \sum_{\varepsilon,l,\rho}
       \bar\delta_\rho\,\bar\pi_{l+1}^{-6}\,
       \zeta_{3,\varepsilon,l+1,\rho}\,
                \zeta_{2,-\varepsilon,l,\rho}^{121}.\ }
\tag{CS12}
\] The left side denotes the same actual logarithm expressed in the new coordinates before the common scaling. Every displayed coefficient is a nonzero residue. All scaled coefficients before reduction are integral.

\textbf{Proof.} Start with (CS9) on the ordinary factors and sum them. On one pair of ordinary J copies write \(w_{i,\pm}=u_i\pm i v_i\). In characteristic11, since \(121\) is a Frobenius power and fixes i, \[
 u_3u_2^{121}+v_3v_2^{121}
  =\frac12\left(w_{3,+}w_{2,-}^{121}
                          +w_{3,-}w_{2,+}^{121}\right).
\tag{CS13}
\] For each sign the B-coordinate change is z=Uw, with inverse \(w=U^{\mathsf t}z\). Thus the coefficient pairing between the new B labels k,l is \[
 \sum_j\bar U_{kj}\bar U_{lj}^{121}
       =\sum_j\bar U_{kj}\bar U_{l+1,j}
       =\begin{cases}1&k=l+1,\\0&k\ne l+1.\end{cases}
\tag{CS14}
\] This uses the actual trace orthogonality, rather than a complex-unitary approximation or a generic invertibility assertion.

In the quadratic extension, the inverse-coordinate pairing is \[
  (\bar V_a^{-1})^{\mathsf t}
                 (\bar V_a^{-1})^{(121)}
   =\bar V_a^{-\mathsf t}\bar V_a^{-1}
   =\operatorname{diag}(\bar\delta_{a,+},\bar\delta_{a,-}),
\tag{CS15}
\] by (CS5) and (CS11). Taking the tensor product for the two quadratics introduces exactly \(\bar\delta_\rho\) and forces the same real embedding \(\rho\) on the paired labels. For h=1 this step is absent.

Finally \(z_{i,\varepsilon,l,\rho}=\pi_l^{-3}\zeta_{i,\varepsilon,l,\rho}\). On a nonzero pair the two factors contribute \[
              \bar\pi_{l+1}^{-3}
                        (\bar\pi_l^{-3})^{121}
                              =\bar\pi_{l+1}^{-6}.
\tag{CS16}
\] Equations (CS13)--(CS16) prove the displayed form.

All coordinate substitutions and inverses are integral at11, and the same scale \(\varpi\) is used on every input. Substitution therefore commutes with reduction and preserves integrality of the entire scaled logarithm in (CS9). At a fixed ordinary degree there are only finitely many coefficients involved. This verifies the whole-series assertion and rules out an unexamined logarithmic term at the same critical valuation. ∎

There are fourteen nonzero pairs when h=1 and fifty-six when h=4. In particular, fix l and \(\rho\), and set \[
                \tau_3=(3,+,l+1,\rho),\qquad
                \tau_2=(2,-,l,\rho).
\tag{CS17}
\] On the plane where only \(X=\zeta_{\tau_3}\) and \(Y=\zeta_{\tau_2}\) are nonzero, (CS12) is exactly one nonzero monomial \(\gamma XY^{121}\). Thus the critical-exponential lemma DP58(ES6--12) applies directly to this explicit new-CM plane.

For \[
                  A_r=11^r,\qquad B_r=121\,11^r,\qquad r\ge0,
\tag{CS18}
\] it gives, for the actual product primitive kernel \(F_\lambda=\exp(\lambda\beta_{\rm prod})\), \[
 \begin{split}
 v\bigl([\zeta_{\tau_3}^{A_r}\zeta_{\tau_2}^{B_r}]F_\lambda\bigr)
                  &=-\frac{121\,11^r-1}{10},\\
 \boxed{\ 
 v\!\left(
       \partial_{\tau_3}^{A_r}\partial_{\tau_2}^{B_r}
                      F_\lambda(0)\right)
                     =\frac{11^r-1}{10}.\ }
 \end{split}
\tag{CS19}
\] The first equality uses the full critical monomial, uniformly for every unit \(\lambda\). The second adds the two exact Legendre factorial valuations; both exponents are powers of11. Every derivative in (CS19) is nonzero.

This calculation locates the sharp primitive kernel in a particular pair of the enlarged CM eigenspaces. It does not enlarge the accepted singleton/complement display theorem to an arbitrary higher-height CM word.

\subsection{3. An explicit actual selected Schur family}\label{an-explicit-actual-selected-schur-family}

Fix an admissible source rank record and an edge with \[
                    a=N_j,\qquad b=N_{j+1}=a+d,
                         \qquad d=t+1\ge2.
\tag{CS20}
\] Choose fixed word-independent HC data on all f embeddings and let \(\nu_\tau^*\) be its actual selected positive partitions of length a, as in the literal identities checked in DO025. The numerical full-Hyp3 baselines of K031 and K032 are concrete choices on the two hosts. Put \[
 D_0=\sum_\tau|\nu_\tau^*|,\qquad
 \mathcal F_0=\prod_{\tau,i}\nu_{\tau,i}^*!,\qquad
                        c_0=v(\mathcal F_0).
\tag{CS21}
\] Ranks, the edge, the baseline, both fields and all stage characters remain fixed as r varies.

Reserve column1 of each a-by-b creation-symbol matrix \(Z_\tau\). Define on columns2 through a+1 \[
 Q_0(Z)=\prod_\tau\prod_{q=1}^{a}
    \det Z_{\tau,[1,q],[2,q+1]}^{
                                  \nu_{\tau,q}^*-\nu_{\tau,q+1}^*},
          \qquad\nu_{\tau,a+1}^*=0.
\tag{CS22}
\] Here \([2,q+1]\) means the consecutive q columns \(2,\ldots,q+1\). Each such column exists because b≥a+1. Set \[
 \boxed{\ 
 P_r=Z_{\tau_3,1,1}^{A_r}
             Z_{\tau_2,1,1}^{B_r}\,Q_0(Z),\qquad
 T_r=A_r+B_r+D_0=122\,11^r+D_0.\ }
\tag{CS23}
\] This polynomial is explicitly prescribed; no unknown distribution among the seven copies is selected.

Increase the first HC entry \(u_0\) by \(A_r\) at \(\tau_3\), by \(B_r\) at \(\tau_2\), and by zero at every other place. Make the same modifications for every word. The selected partitions on the fixed edge become exactly \[
 \nu_{\tau_3}=\nu_{\tau_3}^*+A_r e_1,\qquad
 \nu_{\tau_2}=\nu_{\tau_2}^*+B_r e_1,\qquad
                       \nu_\tau=\nu_\tau^*\text{ otherwise}.
\tag{CS24}
\] All e entries are unchanged. Thus the actual uniform characters supplied in K031/K032 stay the same, as do their compact G restrictions and conductors. The word row condition and the common first/last-u equations in Wts.Ok are preserved: the row data v,b,e do not change, and the first-u modification is common to every word.

Every HC inequality survives an increase of the largest u. The full numerical Hyp3 survives as well when the baseline satisfies it. The Dpar and NCGaps lists contain no u; the parity, eta and1801 clauses are unchanged. Interlacing and \(u_{\rm last}>-b_0\), together with \(e_j\ge e_0>b_0\), imply \(u_i+e_j>0\). Increasing \(u_0\) increases the absolute value in the only changing Centred inequalities. This checks all clauses, not only HC interlacing. Independently fixed upper SizeData windows are not preserved.

\subsubsection{Proposition CS25. Exact source-and-target summand}\label{proposition-cs25.-exact-source-and-target-summand}

The polynomial \(P_r\) is primitive integral and already belongs to the exact compact Cauchy summand indexed by the partitions (CS24), with the source and target determinant twists of DO025.

\textbf{Proof.} Prefix minors are invariant under adding an earlier row to a later row. Multiplication by a first-row variable retains this invariance. The row degrees telescope to (CS24). The product of the minor diagonal monomials and the two first-column powers has coefficient one, and no other determinant choices produce the same all-diagonal filler monomial. Thus \(P_r\) is primitive in the actual polynomial coefficient lattice.

Over characteristic zero, the Cauchy decomposition at each CM place is \[
       \operatorname{Sym}(T_\tau\otimes S_\tau^*)
          =\bigoplus_\lambda S_\lambda(T_\tau)\otimes S_\lambda(S_\tau)^*.
\tag{CS25}
\] In the source representation of a particular summand, its maximal-unipotent invariant line has precisely the source highest weight labelled by that \(\lambda\), and no other such invariant weight. Hence the source-invariant row weight \(\nu_\tau\) of \(P_r\) selects only \(\lambda=\nu_\tau\). The target vector can be non-highest, but it lies entirely in \(S_{\nu_\tau}(T_\tau)\), rather than being a sum of different target Schur types. Tensoring over the CM places and restoring DO025's fixed determinant twists proves the claim. No rational Schur projector or unspecified lattice denominator occurs. ∎

\subsection{4. Exact formal value and the permitted image ideal}\label{exact-formal-value-and-the-permitted-image-ideal}

On the Hom variety \[
            X=(J^{14h})^{ab}=J^{14hab},
                   \qquad \mathcal L=L_{2\infty}^{\boxtimes14hab},
\tag{CS26}
\] take the actual source/Chern reduced gauge, which is K018's full-CM gauge by K030/K032. Fix any allowed \(M=29^e\) and \(\kappa=M^2\).

At row1, column1 choose the integral formal prefactor in (CS6) that cancels the product of the normalized algebraic frame units. This produces the actual product primitive kernel with parameter \(-\kappa\). On the filler entries choose an integral algebraic monomial whose leading logarithmic term in the standard CM basis is \[
                    \prod_{\tau,i}
                         \zeta_{\tau,i,i+1}^{\,\nu_{\tau,i}^*}.
\tag{CS27}
\] Such an integral monomial exists because the new CM tangent basis is integral invertible at11: choose integral formal parameters with precisely those linear differentials. Its degree is D₀. On unused entries use normalized integral frame prefactors, giving reduced value1 at zero.

The first matrix entry and all filler entries are different product factors. The degree-D₀ filler derivative cannot see higher terms of its algebraic coordinates or of its reduced frame. The matching diagonal monomial has coefficient one in Q₀, and all other exponent vectors differ. Consequently its evaluation is \textbf{exactly} \(\mathcal F_0\); no factorial is inverted. Combining with (CS19) gives an actual integral formal section \(S_r\) such that \[
 \boxed{\ 
 v\!\left(
 M^{-T_r} P_r(\partial)F_{S_r}(0)\right)
                =\frac{11^r-1}{10}+c_0.\ }
\tag{CS28}
\] The unit \(M^{-T_r}\) is the source's normalized level reading. This is a value on an integral formal-section lattice; it is not yet a value of the fixed untranslated sigma section.

Let \(\mathfrak I_r\) be the image ideal of this normalized functional on the full integral formal \(T_r\)-jet module of \(\mathcal L^{M^2}\). Choose \[
      M_r=29^{e_r},\qquad
      e_r=\max\{1,\lceil\log_{29}(3T_r+2)\rceil\}.
\tag{CS29}
\] K031/K032 and the exact DP59/K022 half-orbit theorem give \[
 \mathcal W_{M_r}=[\alpha]^*H^0(X,\mathcal L^{M_r}),
 \quad\alpha=(1+2\sqrt{-7})^{e_r},\quad
 \operatorname{rank}\mathcal W_{M_r}=M_r^{fab},
\tag{CS30}
\] with its actual integral translated basis indexed by \(X[\alpha^2]/X[\alpha]\), and rescaled translations \(M_rt'\in X[\alpha]\). Its map to the integral \(T_r\)-jets is surjective. Therefore \(\mathfrak I_r\) is exactly the ideal generated by its actual allowed tuple values. From (CS28), \[
                           v(\mathfrak I_r)
                          \le\frac{11^r-1}{10}+c_0.
\tag{CS31}
\] This ideal is nonzero. A tuple attaining the least nonzero valuation exists because the basis is finite. It is one actual algebraic allowed tuple, although the formal interpolating coefficients may have been complete-local. Its finite Schwartz vector has norm one. The interpolation coefficient vector is not charged as an actual chosen input.

\subsection{5. A matching lower bound with constant error}\label{a-matching-lower-bound-with-constant-error}

Let \[
 K_0=\frac1{10}
       \max_{\beta\in\operatorname{supp}Q_0}
                                 \sum_m s_{11}(\beta_m).
\tag{CS32}
\] This is a finite fixed constant, independent of r and M. One computable bound is \[
                 K_0\le a\sum_{\tau,i}d_{11}(\nu_{\tau,i}^*),
 \qquad
 d_{11}(n)=
 \begin{cases}0&n=0,\\1+\lfloor\log_{11}n\rfloor&n>0.\end{cases}
\tag{CS33}
\] The accepted DP51 full common-weight bound, transferred to the actual integral local module in K027(EX2--5), applies to every section here. The new standard CM change is integral invertible, so that full bound applies in these coordinates as well.

Every monomial of \(P_r\) has the two fixed growing exponents Aᵣ and Bᵣ in columns disjoint from Q₀. Both have base11 digit sum1. The factorial digit loss is therefore exactly bounded by \(1/5+K_0\), rather than by a number growing with r. Thus \[
 \boxed{\ 
 \frac{T_r}{1220}-\frac15-K_0
 \ \le\ v(\mathfrak I_r)\
 \le\frac{T_r}{1220}
                  +c_0-\frac1{10}-\frac{D_0}{1220}.\ }
\tag{CS34}
\] The same lower bound holds for every integral section; the individual allowed tuple attaining \(\mathfrak I_r\) satisfies both bounds. All derivative values are Hurwitz-integral by K023, so the lower bound may also be replaced by its nonnegative integer ceiling.

In particular \[
             \boxed{\ 
             \lim_{r\to\infty}\frac{v(\mathfrak I_r)}{T_r}
              =\lim_{r\to\infty}\frac{v(a_r)}{T_r}
                                      =\frac1{1220}.\ }
\tag{CS35}
\] More precisely, \(v(\mathfrak I_r)-T_r/1220\) stays in the fixed interval displayed in (CS34). This proves sharpness in an \textbf{explicitly prescribed} full-CM selected family on either enlarged host. It is stronger than an ambient uniform-slope obstruction, and it has no hidden CM weight projection.

For the empty filler the same argument even gives \(v(\mathfrak I_r)=(11^r-1)/10\): the lower bound is less than this integer by only1/10. That empty-filler observation is not an assertion of the full source HC/Hyp3 conditions. The nonempty actual baseline filler is retained in (CS21)--(CS35).

For the concrete first-edge baselines already constructed in K031/K032, the constants are as follows. The last column concerns only comparison of the growing degree \(122\cdot11^r\) with D₀; it does not replace the large fixed upper error c₀.

\begin{longtable}[]{@{}
  >{\raggedright\arraybackslash}p{(\columnwidth - 8\tabcolsep) * \real{0.1579}}
  >{\raggedleft\arraybackslash}p{(\columnwidth - 8\tabcolsep) * \real{0.2105}}
  >{\raggedleft\arraybackslash}p{(\columnwidth - 8\tabcolsep) * \real{0.2105}}
  >{\raggedleft\arraybackslash}p{(\columnwidth - 8\tabcolsep) * \real{0.2105}}
  >{\raggedleft\arraybackslash}p{(\columnwidth - 8\tabcolsep) * \real{0.2105}}@{}}
\toprule\noalign{}
\begin{minipage}[b]{\linewidth}\raggedright
Actual host and fixed edge
\end{minipage} & \begin{minipage}[b]{\linewidth}\raggedleft
D₀
\end{minipage} & \begin{minipage}[b]{\linewidth}\raggedleft
c₀
\end{minipage} & \begin{minipage}[b]{\linewidth}\raggedleft
Computable bound for K₀
\end{minipage} & \begin{minipage}[b]{\linewidth}\raggedleft
First r with \(122\cdot11^r\ge D_0\)
\end{minipage} \\
\midrule\noalign{}
\endhead
\bottomrule\noalign{}
\endlastfoot
\(J^{14}\), 42→1802 & 3,504,831,920,337,006 & 350,483,192,023,230 & 899,640 & 13 \\
\(J^{56}\), 168→1802 & 194,467,812,219,985,248 & 19,446,781,221,824,712 & 60,484,032 & 15 \\
\end{longtable}

These are exact finite sums of the actual selected baseline parts and their factorial valuations, reproduced in the certificate. Thus the bounded-error asymptotic theorem retains a potentially substantial finite baseline cost. It does not itself certify a useful fixed numerical MainX window.

\subsection{6. Metric and first-image scope}\label{metric-and-first-image-scope}

The primitive target vector in the Cauchy summand has the exact metric correction already proved in K026: \[
 \Theta_r=
 \frac{A_r!\,H_{\nu_{\tau_3}^*}}
               {H_{\nu_{\tau_3}^*+A_re_1}}\,
 \frac{B_r!\,H_{\nu_{\tau_2}^*}}
               {H_{\nu_{\tau_2}^*+B_re_1}}.
\tag{CS36}
\] Here H is the hook product. Each quotient is a fixed finite product \(s/(n+s)\), so its real logarithm and its11-adic valuation have size \(O(\log(T_r+1))\), with the residual partitions fixed. Scaling the target vector to norm one multiplies the coefficient by \(\Theta_r^{-1/2}\); it is not treated as an integral unit. Hence its valuation slope is still \(1/1220\).

All real creation factors remain explicit. In standard coordinates put \[
        h_\tau=\frac{\delta_\rho\,\pi_l^{-6}}{2G_i},
           \qquad \tau=(i,\varepsilon,l,\rho).
\tag{CS37}
\] At the chosen complex realization these numbers are positive. The actual squared archimedean Schwartz norm of the primitive vector is \[
 \begin{split}
 \|r_{P_r}\phi_{0,\infty}\|^2
   ={}&\det(2Y)^{-7hab}\,
          A_r!B_r!\prod_\tau H_{\nu_\tau^*}\\
      &{}\cdot(\pi h_{\tau_3})^{A_r}
                (\pi h_{\tau_2})^{B_r}
                 \prod_\tau(\pi h_\tau)^{|\nu_\tau^*|}.
 \end{split}
\tag{CS38}
\] The two first-column powers and all filler variables are disjoint, so their Fischer norms multiply. The hook normalization is the exact parent R011/K026 normalization, and K030/K032 give the metric coefficients. This factorial archimedean cost is not an \(O(\log T_r)\) term. Dividing (CS38) by \(\Theta_r\) is an alternative target normalization of the same vector, not an additional independent loss.

For the rational29 principal class refinement, the exact index is \[
             I_a(M)=M^{fa^2}
                         \prod_{q=1}^{a}(1-29^{-q})^{6h}.
\tag{CS39}
\] Its11-adic valuation is independent of the exponent in M. If a class-supported extraction with the chosen finite compact type is supplied, its compact dimension factor varies by the two fixed-rank Weyl ratios \[
 \prod_{j=2}^{a}
 \frac{A_r+\nu_{\tau_3,1}^*-\nu_{\tau_3,j}^*+j-1}
                  {\nu_{\tau_3,1}^*-\nu_{\tau_3,j}^*+j-1}
 \quad\text{and the analogous factor with }B_r,\tau_2.
\tag{CS40}
\] These have11-valuation \(O(\log(T_r+1))\). Thus these specified metric/class corrections alone preserve the local slope. The fixed product period is \(C^{14hab}\) in the conventions of DO034, and remains separate.

There are two important restrictions on this conclusion. First, the value \(a_r\) in (CS35) comes from the unmodified allowed translated vacuum and the selected polynomial. DO034's finite Steinberg auxiliary vector, or LIFT20's literal raw G factor, defines a different functional. A separate simultaneous-germ argument is required before assigning that restricted functional the same sharp upper slope. The factorial nonvanishing theorem for an arbitrary auxiliary vector does not establish this sharper assertion. Second, the next automorphic stage acts on a previous theta image, which is a further integral sublattice. Neither (CS34) nor a nonzero theta-representation theorem determines its saturation or prevents cancellation after the class integral.

The new theorem therefore closes the enlarged-CM selected-kernel sharpness gap. It does not close the later-image arithmetic or the complete DP44 gain/error inequality. The common-type field obstruction Fields.Kq=\(\mathbb Q\) remains, including on the8-divisible host.

\subsection{7. Exact inputs and finite certificate}\label{exact-inputs-and-finite-certificate}

\begin{itemize}
\tightlist
\item
  Accepted DP46, unit\_arrow\_kernel\_invariance\_DP46.md, UI16--19 and §5: full scaled inverse maps, critical logarithmic initial form, and invariant-frame comparison. Accepted DP51, harmonic\_asymptotic\_rate\_DP51.md, §6 HA17--21: the common weight on every variable and full local-module bound. Accepted DP58, explicit\_sharp\_jets\_linear\_levels\_DP58.md, ES6--12: the critical-exponential coefficient theorem. These statements were read at their full interfaces; no new display proof is claimed.
\item
  AM-K018 K19.2/K19.13: exact source/Chern gauge on the ordinary product. K030 PH17--26 and K032 QH13--21: actual product principal lattices, full CM coordinates and integral11 comparisons. K031 HC22--26 and K032 QH33--36: the actual central alpha and permitted half-orbit.
\item
  AM-K022 J6--16/J30: ordinary jet surjectivity and equality of the permitted image ideals. AM-K024 SC9--19: filler and exact Cauchy-label argument. AM-K026 HC18--25 and parent R011: target hook comparison and Fischer normalization. AM-K027 EX2--7: accepted common-weight transfer to arbitrary joint integral local prefactors.
\item
  \nolinkurl{Formalism/LocalToGlobal/Columns/AbelianHostFieldConditions.lean:442–469,} Wts.Ok and selected weight definitions; \nolinkurl{Columns/BaseLemmas.lean:496–552,} selected recurrence; \nolinkurl{LocalRingCardinality/DataTheta.lean:529–555} and \nolinkurl{SharpRestriction/CuspSpanNormBounds/CayleySignature.lean:110–118,} exact Hyp3 clauses. Paths are relative to work/research/alpoge\_source. The old unmodified Fields.Kq condition is not changed.
\end{itemize}

The finite certificate \nolinkurl{selected_cm_frobenius_sharp_content_certificate.py} and its JSON output verify the actual seven-copy Frobenius permutation and all49 trace-pair entries over the cyclotomic residue algebra, the plus/minus pairing, both weighted quadratic pairings, the standard-lattice unit factors, the explicit support pairs, and exact Legendre and baseline arithmetic. They do not infer an infinite coefficient theorem or a selected automorphic image from finite samples. Those assertions, where made, are proved by the cited accepted whole-series theorems and the arguments above. No Lean execution is claimed.

\clearpage

\section{ALP-LIFT-22: the two prescribed sharp G kernels compose on their actual middle level}\label{alp-lift-22-the-two-prescribed-sharp-g-kernels-compose-on-their-actual-middle-level}

Author: GPT-6 Astra Pro, ALP\_MACHINE /root/lift\_accounting. Date: 22 November 2026 UTC. Inherited reasoning effort is not exposed.

This result closes a further local gate for the literal m=2 sharp coupled construction. The raw first G kernel does not require an intermediate generic projector, an arbitrary translate, or a new finite Schwartz choice before the second edge. Its actual image already has the middle level forced by ChIdx.Ok. The prescribed sharp terminal G kernel is injective on that image. The coupled Y-variable has only one possible G residue class, so no local sum over Y can cancel this component.

The result is local at G, with a fixed Haar normalization of the actual small-theta quotient models. It leaves the globally pinned quotient scalars, the other places, global Y-summation, and the exact arithmetic automorphic row lattice separate.

\subsection{1. The literal middle source level}\label{the-literal-middle-source-level}

Use LIFT21's notation \[
 m=2,\quad d=t+1,\quad r=k+1,\quad b=k+d,\quad N=b+d.
 \tag{CG1}
\] The place is degree one and relative split, with residue \(\kappa=\mathbf F_q\), \(q\ne11\), and the two actual tame characters \(\eta_0,\eta_1\) have distinct nontrivial residues. All source references below use the recovered commit 7cbf7a0c32fca1badf05931879a37c3f8a2b03b4. Let FS abbreviate \nolinkurl{Formalism/LocalToGlobal/Columns/FiniteSumNonzeroLargeDepth.lean.}

For m2=0 the coupled middle index is one. The source forces \[
             c_1=0,\qquad h_1=1.
 \tag{CG2}
\] The first is the literal cOf definition; the second is ChIdx.Ok.last\_h (FS1025--1026). Thus the following calculation does not choose a new intermediate frame.

When c=0, every Ufl in FS468 is the whole space and each extra coordinate vector eN(n,n+i) is zero. Therefore the conditions defining R1 are automatic: \[
  R1=\mathrm{GL}_b,\qquad N1=\{1\},\qquad n1Arg=0,
  \qquad \operatorname{prCat}(1,P,1)=\operatorname{Jat}(1,P).
 \tag{CG3}
\] The inclusion inX is the identity, proving the last equality from Cbe and prX. These are the exact definitions FS470--541, not a hypothesis about a new model.

For the actual sharp flag, Jat is JGsym. Since rS=0, symIdx is the identity. At c=0 the lattice flag in flagS is the usual residue prefix flag, except at index k: bd5(k,0,k)=0, so its lattice is \(qO^b\). The conditions in flagU therefore give the following exact residue subgroup: \[
 H^\sharp=
 \left\{\begin{pmatrix}u&A\\0&h\end{pmatrix}:
       u=\operatorname{diag}(u_k,1),\ u_k\in U_k(\kappa),
       A\in M_{k+1,t}(\kappa),\ h\in\mathrm{GL}_t(\kappa)
 \right\}.
 \tag{CG4}
\]

Here are the individual flag checks. For s\textless k, the condition on the s-th basis vector makes its diagonal entry one and its lower entries zero. At s=k it fixes the \emph{entire} column to e\_k because the modified lattice is qO\^{}b. Preservation of the prefix flag at r=k+1 gives the upper parabolic block shape. The next listed index is r+t=b and adds no condition; later indices also add none. Preservation of the ordinary lattice gives invertibility. Conversely every lift of(CG4) satisfies these conditions. Thus the local subgroup is the full inverse image \(J^\sharp\) of \(H^\sharp\). The flag definitions are FS738--766, with LevGSym.lean11--13 and its entrywise membership lemma108--141.

The grpG and chrG clauses of ChIdx.Ok (FS1072--1083) now force the actual middle levG group to be \(H^\sharp\). The n1Arg part of chiInv is one by(CG3); lamRow on this group is \[
        \psi_0\!\left(\sum_{i=0}^{k-2}(u_k)_{i,i+1}\right)
                                          \chi_0(\det h).
 \tag{CG5}
\] Hence the actual adelic middle type, with the inverse convention of RecLevel.lean50--52, is \[
 \lambda^\sharp=
        \psi_{U_k}^{-1}\,\bar\eta_0(\det h).
 \tag{CG6}
\] All character values are the actual residue values; the determinant inverse has not been dropped.

This middle type is not the generic ColV1 type. The source distinction is explicit: ColForms.lean105--114 defines MfV=ColV0, and FS1408--1417 imports that initial space by startC. The actual coupled kernel uses midGSD=(chF \ldots{} 1).mem (FS1220) and ChKer.cp on that space (FS1236--1237). In chF the group is FrmG of L.levG and L.levQ (FS1090--1102). It therefore uses the type(CG6), not a substituted ColV1. No assertion that the whole ChIdx.Ok record has been constructed is made.

\subsection{2. The coupled G residue class is forced}\label{the-coupled-g-residue-class-is-forced}

The allowed G classes of Y are \(q^{-1}O^b/O^b\). Before choosing any class, the first sharp raw kernel reads the matrix \([X,qy]\), by LIFT20(RG4). Its PhiGS condition fixes the last source column, which is column k in this specialization, to e\_k in every target row. Therefore \[
 \operatorname{PsiAt}_{\rm sharp}(X,y)\ne0
       \quad\Longrightarrow\quad qy=e_k\pmod q.
 \tag{CG7}
\] The diagonal condition fixes the k-th entry; the special sharp condition fixes all entries above it; the remaining lower and external-row conditions fix all entries below it. These include all b target coordinates.

Thus exactly one G residue class can contribute: \[
                         y=q^{-1}e_k\pmod{O^b}.
 \tag{CG8}
\] For that class the first kernel is exactly LIFT20's \[
 f_U(X)=
 \begin{cases}
    \psi_{U_k}(u),&X\in M_{b,k}(O),\
                       \bar X=[I_k;0]u,\ u\in U_k(\kappa),\\
    0,&\text{otherwise}.
 \end{cases}
 \tag{CG9}
\] A representative change by an integral vector does not affect either first or terminal G factor. The same allowed class is used on both coupled edges. A local coupled sum over Y therefore has only this G term. Distinct global Y classes can still differ at other places; this statement does not remove their global summation.

\subsection{3. The first raw output is an identity-parabolic section}\label{the-first-raw-output-is-an-identity-parabolic-section}

Let \(\pi_0\) be an actual initial constituent with unramified cuspidal support and a finite Steinberg type as in LIFT19/LIFT20. Put \[
 V_0=\pi_0^{J_0,\psi_{U_k}^{-1}},\qquad
 J_0=\operatorname{red}^{-1}U_k(\kappa),\qquad
 a_q=\operatorname{vol}(J_0)
                =\frac{|U_k(\kappa)|}{|\mathrm{GL}_k(\kappa)|}.
 \tag{CG10}
\] This is precisely the initial MfV local type. It contains the explicit nonzero Steinberg generic vector.

The full-rank integral of LIFT21(TG13)--(TG16), now on the edge k to b with character eta0, identifies the actual first small theta quotient with \[
                      \pi_1=
          \operatorname{Ind}_{Q_{(k,d)}}^{\mathrm{GL}_b}
                 (\pi_0\boxtimes\eta_0\det_d).
 \tag{CG11}
\] Its boundary is excluded exactly as in LIFT21: after removing eta0, every source inertia line is the nontrivial residue of eta0 or its inverse. The induction is irreducible by LIFT19. Thus the compact integral is a fixed Haar normalization of the actual quotient, with the same explicit global-normalization caveat.

Define \(s(v)\) in the compact induced model to be supported on the identity parabolic cell \(Q(O)K_b^1\), with \[
 s(v)\!\left(\begin{pmatrix}A&C\\0&D\end{pmatrix}\right)
                 =\pi_0(A)\eta_0(\det D)v
                         \qquad\text{on }Q(O).
 \tag{CG12}
\] Here v is \(K_k^1\)-fixed. The expression is well-defined on the compact cell and extends to normalized induction by Iwasawa. Its identity value is v.

The first fixed raw kernel has the exact image \[
                  {\cal T}_0(f_U,v)=a_q\,s(v)
                                      \qquad(v\in V_0).
 \tag{CG13}
\] To prove this, its residue columns span only the standard k-plane. The compact target section therefore vanishes on every other parabolic quotient point. At identity the source integral is supported on J0; the source character is trivial on its determinant, and \(\psi_{U_k}(u)\pi_0(u)v=v\). It equals \(a_qv\). Target parabolic covariance determines its values on the only remaining cell, yielding (CG13). No compact projector onto a different target type was applied.

The image s(v) has precisely(CG6). Indeed \(H^\sharp\subset Q_{(k,d)}(\kappa)\), and its Levi entries in that parabolic are u\_k and a block with determinant det h. Right translation of the identity-cell section thus acts by \(\psi_{U_k}^{-1}\bar\eta_0(\det h)\). Its support is preserved since this subgroup is contained in the parabolic. Therefore \[
                     s(V_0)\subset
                      \pi_1^{J^\sharp,\lambda^\sharp}.
 \tag{CG14}
\]

For precision, s(v) is invariant under the joining root \(1+cE_{k-1,k}\), which lies in the Q radical. A projector requiring a nontrivial generic character on that root would kill it. This is an obstruction only to imposing that different compact type on this particular output. The actual sharp coupled middle level is(CG6), so that projector is not part of the present source construction.

\subsection{4. The actual terminal raw kernel preserves this exact image}\label{the-actual-terminal-raw-kernel-preserves-this-exact-image}

Let \(\Phi^\sharp_y\) be the actual terminal PhiTop/ftY function with the same class(CG8). The sharp proof of LIFT21(TG19), (TG22), and(TG24) uses only the restricted character of \(J^\sharp\), not an extension to the larger generic group J. It consequently applies to every vector in \(\pi_1^{J^\sharp,\lambda^\sharp}\). Explicitly its identity integral is \[
          {\cal T}_1(\Phi_y^\sharp,w)(1)
                            =v_{J^\sharp}w,
     \qquad w\in\pi_1^{J^\sharp,\lambda^\sharp},
 \tag{CG15}
\] since on Jsharp the raw character times eta1's source determinant inverse is lambda-sharp inverse. On every graph-plane representative h\_z it is \(v_{J^\sharp}\psi_0(-z_k)w\), and it is zero on the other finite parabolic cells. The same boundary-rank proof identifies it with the actual second small theta quotient.

Combining(CG13)--(CG15) proves the exact composed identity \[
 \boxed{\quad
  \bigl[{\cal T}_1(\Phi_y^\sharp,
                 {\cal T}_0(f_U,v))(1)\bigr](1)
                         =a_q v_{J^\sharp}v
                          \qquad(v\in V_0).\quad}
 \tag{CG16}
\] There are two identity evaluations, in the nested induced models. Since both scalar volumes are nonzero, the two \emph{fixed raw} G kernels compose injectively on the whole initial generic type V0. No local Hecke rotation, new Schwartz input, or intermediate projection is required.

The residue class(CG8) is forced by the first raw kernel, so this is also the only possible local sharp coupled G term. This assertion does not replace the global coupled sum, its mass factors, or its other-place kernels.

\subsection{5. Exact norms and natural integral content}\label{exact-norms-and-natural-integral-content}

Put \(n_U=k(k-1)/2\) and use the Gaussian parabolic indices \({\cal J}_{u,v}(q)\) defined in LIFT21. The first identity-cell insertion satisfies \[
 \|s(v)\|^2={\cal J}_{k,d}(q)^{-1}\|v\|^2,\qquad
 \|f_U\|^2=q^{-bk}|U_k(\kappa)|.
 \tag{CG17}
\] The raw second factor has \(\|\Phi^\sharp_y\|^2=q^{-Nb+b}|H^\sharp|\). The exact composite induced-model norm is \[
 \boxed{
 \|{\cal T}_1(\Phi^\sharp_y,{\cal T}_0(f_U,v))\|^2
   =\frac{q^b(a_qv_{J^\sharp})^2}
                {{\cal J}_{k,d}(q){\cal J}_{b,d}(q)}
                                                \|v\|^2.}
 \tag{CG18}
\] The first norm is a single cell of the compact quotient; the second is the q\^{}b graph cells of LIFT21. This proves the formula for every vector in V0, not only a norm-one test vector.

If both finite raw kernels are unit-normalized, the squared scaling constant on V0 becomes \[
 \boxed{
 \frac{q^{bk}|U_k|}
          {|\mathrm{GL}_k(\kappa)|^2{\cal J}_{k,d}(q)}
 \cdot
 \frac{q^{Nb}|H^\sharp|}
          {|\mathrm{GL}_b(\kappa)|^2{\cal J}_{b,d}(q)}.}
 \tag{CG19}
\] No real selected-polynomial norm, global class mass, or automorphic quotient scalar is absorbed into this formula.

Fix an integral finite Steinberg lattice inside the actual \(\pi_0^{K_k^1}\), and use its natural finite induced lattices on both subsequent representations. Its initial generic intersection is saturated. The normalized first insertion \(a_q^{-1}{\cal T}_0(f_U,-)=s\) is integral, and identity evaluation returns the input. The normalized second map has only zero or root-of-unity multiples of that input among its finite graph values. Consequently \[
 (a_qv_{J^\sharp})^{-1}
        {\cal T}_1(\Phi_y^\sharp,{\cal T}_0(f_U,-))
 \tag{CG20}
\] is a primitive saturated injection of the initial generic lattice into the second natural induced lattice. The saturation proof is the identity-evaluation argument of LIFT21, applied twice; it uses no integral isotypic averaging.

With the valuation normalized by \(v_{11}(11)=1\), the exact scalar content change in these lattices is \[
 \begin{split}
 v_{11}(a_qv_{J^\sharp})
   &=-\sum_{i=1}^{k}v_{11}(q^i-1)
     -\sum_{i=t+1}^{b}v_{11}(q^i-1),\\
 a_qv_{J^\sharp}
   &=\frac{q^{k(k-1)+(k+1)t}
                       |\mathrm{GL}_t(\kappa)|}
           {|\mathrm{GL}_k(\kappa)|
                       |\mathrm{GL}_b(\kappa)|}.
 \end{split}
 \tag{CG21}
\] Thus this local raw two-edge map introduces exactly its displayed Haar denominator and no hidden additional content on the natural induced generic lattice. The comparison with the global arithmetic row lattice and its period normalization is still necessary before identifying this with a DP44 gain.

\subsection{6. Scope}\label{scope}

This theorem uses the actual sharp middle level forced by the source's coupled ChIdx structure, not the separate generic ColV1 level. It closes the two fixed G-edge interaction, including the local Y-variable, for the m=2 route. It applies place by place to every degree-one G prime of a repaired host once its actual first constituent and literal stage characters are supplied.

It does not construct the complete ChIdx.Ok record, prove global nonvanishing of the sum over all remaining Y data, calibrate the local quotient maps by one global pinned theta period, or establish the row-lattice comparison at11. The actual T/Q factors and the other places must be handled with their own tests. The original quadratic-type wrapper and other global constructor conditions are also outside this local calculation.

\subsection{7. Proof checks and references}\label{proof-checks-and-references}

The actual first raw factor is LIFT20. The full-rank quotient construction, published boundary theorem, and sharp terminal graph values are LIFT21. Their hypotheses and normalization limitations are retained. The new source readings are R1 at c=0, last\_h, grpG/chrG, the precise middle chF definition, and the distinction between that space and ColV1.

The finite diagnostic \nolinkurl{composed_raw_sharp_g_certificate.py} checks the forced class, the modified middle flag, first-output character, and the exact norm/content factors. The author has completed the whole-proof reading, including the exact middle source definitions, the unique G residue class, both actual quotient identifications, the identity-parabolic insertion, the larger sharp eigenspace, all composed norms, and the saturated lattice argument. The separate TSV records its time.

The finite diagnostic passed206 first-class tests and1,958 middle flag/character tests. It also checks that the different joining-root generic projector kills the particular identity section, while the actual sharp type fixes it. The displayed Haar content formula gives local composite valuations -10 at each of the two K031 G primes for k42,t1759; for the K032 prime pair and k168,t1633 it gives -73 and-184. These are local model scalars, not global DP44 gains. Small finite geometry cases are diagnostic, not full admissible Ranks examples.

\clearpage

\section{AM-K034 --- complete prime-to-11 torsion gauges and sharp selected auxiliary content}\label{am-k034-complete-prime-to-11-torsion-gauges-and-sharp-selected-auxiliary-content}

Author: GPT-6 Astra Pro, ALP\_MACHINE /root/kummer\_extension, 22 November 2026 UTC; inherited reasoning effort not separately exposed.

This finding proves the full formal-germ comparison that was deliberately absent from AM-DO-030. A prime-to-11 torsion translation has an integral horizontal difference-line frame, and this is the difference frame of the actual reduced theta translation. Consequently a fixed nonzero auxiliary finite Heisenberg operator preserves the complete selected local image ideal, up to the coefficient content of that operator. Combined with AM-K033 and simultaneous interpolation in the actual DP59 half-orbit, this transfers the sharp selected slope \(1/1220\) to the fixed finite Steinberg family of AM-DO-034. No integrality of a Chern-reduced ample-line frame, no theta-null unit, and no saturation of a later automorphic image is asserted.

\subsection{1. The two frames and their actual connections}\label{the-two-frames-and-their-actual-connections}

Let R be a complete DVR of mixed characteristic \((0,11)\), with \(v(11)=1\). The actual data below may first be defined over a number field and then completed at11. Finite coefficient extensions are allowed; no extension is claimed unramified when its extra scalar coefficients require ramification. Let \(X/R\) be an abelian scheme, \(\mathcal L\) a symmetric rigidified relatively ample line, and let its complex realization carry the translation-invariant positive Hermitian form H. In our application \(X\) is an ordinary power of J and \(\mathcal L\) is its product principal line.

Fix an integer \(\kappa\), a positive odd integer d prime to11, and a section \(t\in X[d](R)\). Denote translation by \(T_t\). The rigidified degree-zero difference line is

\[
 \mathcal Q_t=
 T_t^*(\mathcal L^\kappa)\otimes\mathcal L^{-\kappa}
       \otimes p^*(t^*\mathcal L^\kappa)^{-1},
 \qquad p:X\longrightarrow\operatorname{Spec}R.
\tag{TG1}
\]

The displayed rigidification uses the given one on \(0^*\mathcal L\). The polarization homomorphism and the rigidified theorem of the square identify \(\mathcal Q_t\) with the value at \(\kappa t\) of the Poincare family. Therefore it has a rigidification-preserving isomorphism

\[
                 \vartheta_t:\mathcal Q_t^{\otimes d}
                                     \overset{\sim}{\longrightarrow}\mathcal O_X.
\tag{TG2}
\]

This isomorphism is unique after its origin value has been fixed. Indeed two such isomorphisms differ by a global unit on the proper geometrically connected abelian scheme, hence by a unit of R, and the rigidification fixes it. We use that uniqueness, rather than an unspecified line-translation isomorphism, below.

\subsubsection{Lemma TG3. Integral torsion horizontal frame}\label{lemma-tg3.-integral-torsion-horizontal-frame}

There is a unique integrable connection on \(\mathcal Q_t\) for which (TG2) is horizontal. Its origin-normalized horizontal frame on the formal completion \(\widehat X_0\) is an integral frame. If q is an integral formal frame with origin value1 and

\[
                  a=\vartheta_t(q^{\otimes d})\in
                    1+\mathfrak m_0R[[z_1,\ldots,z_g]],
\tag{TG3}
\]

where \(\mathfrak m_0=(z_1,\ldots,z_g)\), then that frame is exactly

\[
                  q_{\mathrm{hor}}=a^{-1/d}q.
\tag{TG4}
\]

Here \(a^{-1/d}\) means the unique formal d-th root with constant1. It and its inverse belong to \(R[[z]]^\times\).

\textbf{Proof.} On a local integral frame q define \[
                        \nabla q=\frac1d\,d\log(a)\,q.
\tag{TG5}
\] This glues under a frame change \(q'=bq\), since \(a'=b^d a\). Because d is a unit and a is a local unit, this is an integral algebraic connection; it is integrable by the displayed formula. Any connection making (TG2) horizontal must have (TG5), proving uniqueness.

The recursive coefficient equation for \(b^d=a^{-1}\), \(b(0)=1\), has at each positive total degree the new coefficient multiplied by d and all other terms already determined in R. Since d is a unit, it has a unique solution in \(1+\mathfrak m_0R[[z]]\). This argument works in several variables and in ramified R. Differentiating \(b^d a=1\) gives \(\nabla(bq)=0\), proving (TG4). No division by a theta section or by one of its values occurs. ∎

\subsubsection{Lemma TG6. This is the connection in the actual reduced translation}\label{lemma-tg6.-this-is-the-connection-in-the-actual-reduced-translation}

Write the reduced translation of a coefficient function F as

\[
 U_t^{\kappa H}F(u)=
   \exp\!\left(-\pi\kappa H(u,t)
                 -\frac{\pi\kappa}{2}H(t,t)\right)F(u+t).
\tag{TG6}
\]

After normalizing the difference-line fibre at0, the difference frame between this reduced translated frame and the original reduced frame is exactly the horizontal frame of Lemma TG3. In particular its complete formal ratio to an integral difference-line frame is the unit in (TG4). This is equality of normalized connections and full germs.

\textbf{Proof.} Use the first-argument-linear convention for H. The original reduced coefficient multiplier is \[
 j_\kappa(\lambda,u)=\alpha(\lambda)^\kappa
       \exp\!\left(\pi\kappa H(u,\lambda)
                      +\frac{\pi\kappa}{2}H(\lambda,\lambda)\right).
\tag{TG7}
\] A substitution into (TG6) shows that its multiplier is \[
 j_\kappa(\lambda,u)\,\chi_t(\lambda),\qquad
 \chi_t(\lambda)=
       \exp\!\left(\pi\kappa\{H(t,\lambda)-H(\lambda,t)\}\right).
\tag{TG8}
\] The second factor is constant in u and has absolute value1. Since \(dt\in\Lambda\) and \(\operatorname{Im}H\) is integral on \(\Lambda\), one has \(\chi_t(\lambda)^d=1\).

The reduced frame of the ample line has squared norm \(\exp(-\pi\kappa H(u,u))\). In the translated frame specified by (TG6), its metric is the same function: the frame multiplier is \(\exp(\pi\kappa H(u,t)+\pi\kappa H(t,t)/2)\), whose absolute square cancels the cross and constant terms in \(\exp(-\pi\kappa H(u+t,u+t))\). The quotient frame of the rigidified line (TG1) consequently has constant norm1. It is horizontal for its flat Chern connection.

Moreover its d-th tensor power is a global holomorphic nowhere-vanishing frame of \(\mathcal Q_t^d\), since \(\chi_t^d=1\). Compare it with (TG2). Their quotient is a global holomorphic unit on the compact complex torus, hence a constant; the origin rigidification makes that constant1. Thus the actual reduced quotient frame is horizontal for exactly the connection making (TG2) horizontal. By the uniqueness in Lemma TG3 it is that connection, including its normalized formal frame.

For completeness, this comparison does not require a possibly nonconvergent integral formal frame to be evaluated over the complex numbers. Compare the algebraic torsion connection and an algebraic reference local frame after complex base change. The equality just proved descends to their finite algebraic jets. Complete at11 and then change to any integral formal frame; (TG3)--(TG5) give the resulting complete formal equality. ∎

The formal root is integral in \textbf{integral algebraic parameters}. Its expression in invariant logarithmic coordinates need not have integral divided Taylor coefficients. Similarly the original ample-line reduced frame is not claimed integral. Lemma TG6 concerns only the prime-to-11 torsion difference line and does not contradict the THETA3 C11/C19 counterexamples.

\subsubsection{Actual fibre normalization}\label{actual-fibre-normalization}

The origin normalization in Lemma TG6 must be compared with the one in the actual finite Heisenberg reading. AM-DO-030, section2, supplies this comparison: pull back by an integer prime to11 killing the torsion, use the relative symmetric theta-group lift of accepted DP62, and descend through the resulting formal etale isomorphism. Its finite-order scalar is an11-unit. One may use \(2d\) instead of d throughout this auxiliary comparison; both are11-units, and this also matches the \(2d\)-divisibility in the primary Mumford construction. It changes no level or source parity condition.

Thus the actual scalar is an11-unit, rather than an unrecorded nonzero scalar. The complete connection comparison in Lemma TG6 is the additional argument beyond DO030's fibre statement. There is no assumption that \(t\in K(\mathcal L^\kappa)\), and no claim that translation by t acts globally as a theta-group automorphism of \(\mathcal L^\kappa\).

The precise primary comparison is Bannai and Kobayashi (2010), author preprint arXiv:math/0610163v4, section1.3, formulas (14) and the definition of \(U^M_{v_0}\), PDF pp.14--15; Proposition2.3, PDF pp.21--22, with the abelian-scheme observation following Corollary2.4; Lemma2.7 and Proposition2.8, PDF pp.24--25. The latter identify this normalized reduced translation with Mumford's algebraic translation. No ordinary-reduction integrality theorem from that paper is used here.

\subsection{2. Equality of the complete translated coefficient modules}\label{equality-of-the-complete-translated-coefficient-modules}

Now take the actual ordinary J-product, with the full source gauge already identified in AM-K018. Write its accepted primitive logarithm as \(\beta_{\mathrm{prod}}\), and let \(U_0\) denote the integral normalized frame factor. In an actual integral formal frame, the complete reduced coefficient module of \(\mathcal L^\kappa\) at0 is

\[
 \mathscr C_\kappa=
 U_0^\kappa\exp(-\kappa\beta_{\mathrm{prod}})
                    R[[z_1,\ldots,z_g]].
\tag{TG9}
\]

All series are compared through the same invariant tangent coordinates. The coordinates z inside \(R[[z]]\) are algebraic formal parameters; the logarithmic variables used for differentiation are distinguished from them.

\subsubsection{Proposition TG10. Translated modules agree}\label{proposition-tg10.-translated-modules-agree}

For each prime-to-11 torsion point t, take every integral formal section near t, translate it to0 by the \textbf{actual} reduced operator (TG6), and include the finite Heisenberg fibre scalar fixed above. The resulting complete coefficient module is exactly \(\mathscr C_\kappa\).

\textbf{Proof.} Translation identifies the formal neighbourhood of t with \(\widehat X_0\) over R. Choose integral frames of \(\mathcal L^\kappa\) near0 and t, and a primitive integral vector in \(t^*\mathcal L^\kappa\). Their difference is an integral frame q of (TG1), up to an integral unit. By Lemma TG6 its reduced difference frame is \(a^{-1/d}q\); by the actual fibre comparison its remaining constant is an11-unit. Hence the translated reduced coefficient of an arbitrary integral local section is \[
 c_t\,h_t(z)\,
        U_0^\kappa\exp(-\kappa\beta_{\mathrm{prod}})\,f(z),
       \qquad c_t\in R^\times,\quad h_t\in R[[z]]^\times,
       \quad f\in R[[z]].
\tag{TG11}
\] Every f occurs. Multiplication by \(c_th_t\) is an automorphism of the complete formal function module and of every finite jet quotient. This proves equality in (TG9), not only an inclusion or equality of their leading terms. The signs of \(h_t\) depend on which integral frame is called q; either sign is an integral unit, and the normalized connection comparison has already fixed the actual reduced operator. ∎

This application uses the actual Chern/source-gauge equality of K018 on the \textbf{untranslated} ordinary product, and the explicit reduced translation (TG6) on the translated line. It does not invoke CM invariance at the translated point, which is generally false. It does not replace the integral connection in Lemma TG3 with an arbitrary holomorphic connection on a torsion line.

\subsection{3. Exact auxiliary selected image ideal}\label{exact-auxiliary-selected-image-ideal}

Use the enlarged hosts of AM-K033: h=1 or4, \(f=42h\), with actual principal variety \(A=J^{14h}\). Fix a source edge a to b, the fixed residual selected partitions, the explicit polynomials \(P_r\), and degrees

\[
                  T_r=122\,11^r+D_0
\tag{TG12}
\]

from K033(CS20--24). The full numerical Hyp3 preservation, exact source-and-target Cauchy summand, and growing first-row weights remain exactly those proved there.

Fix a nonzero auxiliary finite operator supported away from11 and29, in DO030's combined form

\[
 \mathcal O=\sum_{j=1}^{s}b_j W_{y_j},\qquad
 b_j\ne0,\qquad
 \mathcal O\phi_0\ne0,
 \qquad t_j\in X[d]\ \text{distinct},
 \quad X=A^{ab}.
\tag{TG13}
\]

Here d is odd and prime to \(11\cdot29\). The distinct classes are the classes after DO030's actual lattice phases have been combined. Set

\[
 \eta=\min_j v(b_j),\qquad
                    \mathfrak b=(b_1,\ldots,b_s)\subset\operatorname{Frac}(R).
\tag{TG14}
\]

The coefficients and operator are fixed as r varies. Put \(M=29^e\), \(\kappa=M^2\), and \(\alpha=(1+2\sqrt{-7})^e\). Choose

\[
                   M\ge 3s(T_r+1)-1 .
\tag{TG15}
\]

The exact DP59 lattice is \[
 \mathcal W_M=[\alpha]^*H^0(X,\mathcal L^M)
                 \subset H^0(X,\mathcal L^{M^2}),
\tag{TG16}
\] with its actual permitted integral basis \(\mathcal S_\lambda\). No full level-\(M^2\) orbit is substituted for it.

Let \(\ell_r=M^{-T_r}P_r(\partial)\vert_0\). Define the finite fractional ideal \[
                 \mathfrak J_{r,\kappa}=\ell_r(\mathscr C_\kappa)
\tag{TG17}
\] using full integral \(T_r\)-jets. Define \(\mathfrak I^{\rm aux}_{r,M}\) to be the ideal generated by the actual allowed auxiliary tuple values \[
 a^{\rm aux}_{r,\lambda}
 =M^{-T_r}P_r(\partial)
       \left(\sum_j b_j\varepsilon_{j,M}
                   U_{z_j}^{M^2H}F_{\mathcal S_\lambda}\right)(0),
       \qquad z_j=[M]^{-1}t_j .
\tag{TG18}
\] The \(\varepsilon_{j,M}\) are exactly the11-unit finite reading phases of DO030(11--13). Pullback notation in (TG18) uses the section-level functions, as in DO030(13); it is not a second multiplication by M.

\subsubsection{Theorem TG19. Exact content multiplication}\label{theorem-tg19.-exact-content-multiplication}

\[
 \boxed{\quad
           \mathfrak I^{\rm aux}_{r,M}
                 =\mathfrak b\,\mathfrak J_{r,\kappa},
       \qquad
           v(\mathfrak I^{\rm aux}_{r,M})
                 =\eta+v(\mathfrak J_{r,\kappa}).
       \quad}
\tag{TG19}
\]

\textbf{Proof.} By Proposition TG10 the full integral \(T_r\)-jet coefficient module at each translated point \(z_j\), read at0 with its actual reduced operator, is exactly \(\mathscr C_\kappa\). Translation keeps the invariant tangent basis; the multiplication maps and alpha are etale at11, so their maps on formal parameters are integral isomorphisms. Thus the image ideal of the j-th summand on its independent jet module is \(b_j\varepsilon_{j,M}\mathfrak J_{r,\kappa}=b_j\mathfrak J_{r,\kappa}\).

The points \(q_j=\alpha z_j=\bar\alpha^{-1}t_j\) on the base of (TG16) remain distinct modulo11. DO030's simultaneous integral jet theorem, valid on these ordinary products by DO034 and the host constructions, states that (TG15) makes the map \[
  \mathcal W_M\longrightarrow
          \bigoplus_{j=1}^{s}
             J^{T_r}_{z_j}(\mathcal L^{M^2})
\tag{TG20}
\] surjective. It is obtained from the actual \(H^0(\mathcal L^M)\) by alpha pullback and the explicit \(3s(T_r+1)-1\) interpolation bound. The map is over R; formal complete-local interpolation coefficients are sufficient for equality of image ideals.

On the independent direct sum, the image of the sum functional is the sum of the individual image ideals. Surjectivity therefore gives \[
 \mathfrak I^{\rm aux}_{r,M}
       =\sum_j b_j\mathfrak J_{r,\kappa}
       =(b_1,\ldots,b_s)\mathfrak J_{r,\kappa}.
\] This proof also shows explicitly why potential cancellation among different auxiliary terms cannot force extra content: prescribe a desired jet at one point and zero jets at all the other points. It is not an assertion that the individual global theta functions have disjoint support. Since R is a DVR, valuation of the product is the sum, proving (TG19). ∎

The finite list of basis values generates this ideal. Its minimum valuation is attained by one permitted tuple, whose actual finite vector is \(\mathcal O\phi_\lambda\). Selecting that tuple does not require charging the norm of an interpolating coefficient vector. The vector has the exact fixed norm \[
                \|\mathcal O\phi_\lambda\|_f^2
                               =\sum_j|b_j|^2,
\tag{TG21}
\] by DO030's disjoint-place scaled isometry.

\subsection{4. Sharp selected bounds after the auxiliary restriction}\label{sharp-selected-bounds-after-the-auxiliary-restriction}

K033's proof of the formal ideal bounds is uniform for every unit \(\kappa=M^2\), not just its originally chosen least interpolation level. Therefore (CS34) and (TG19) give

\[
 \boxed{\quad
 \eta+\frac{T_r}{1220}-\frac15-K_0
 \ \le\ v(\mathfrak I^{\rm aux}_{r,M})\
 \le\eta+\frac{T_r}{1220}
                  +c_0-\frac1{10}-\frac{D_0}{1220}.
 \quad}
\tag{TG22}
\]

Every basis value satisfies the lower bound, and a permitted nonzero basis value attains the ideal valuation and hence satisfies both bounds. In particular

\[
 \boxed{\qquad
 \lim_{r\longrightarrow\infty}
      \frac{v(a^{\rm aux}_{r,\lambda(r)})}{T_r}
              =\frac1{1220}.
 \qquad}
\tag{TG23}
\]

The bounded additive error includes \(\eta\) and the full fixed baseline constants from K033. The conclusion concerns the \textbf{same explicit selected family} \(P_r\) in the full enlarged CM basis. The auxiliary operator is fixed, actual, and nonzero. There is no enlargement of the harmonic subspace, no appeal to a unit in the ambient unselected space, and no restriction to the degree-\(T_r\) leading algebraic term.

The fixed auxiliary coefficients may be fractional; (TG22) retains their exact minimum valuation. If they are all integral and one is a unit, then \(\eta=0\). If a scalar normalization of the finite vector is imposed, its valuation is added explicitly. These statements use the Hurwitz derivative functional \(\ell_r\), rather than the divided Taylor coefficient. The factorial shifts are already those in K033.

\subsection{5. The actual finite Steinberg family and its costs}\label{the-actual-finite-steinberg-family-and-its-costs}

For the h=1 host, AM-DO-034 section5 supplies the finite Steinberg operator \(\mathcal O_G\), with fixed ranks a,b and the actual stage characters of K031. Its coefficients have minimum11-valuation0. It satisfies

\[
 \begin{aligned}
 d_G&=(g_0g_1)^{21a(a-1)},\\
 \Gamma_G:=\|\mathcal O_G\phi_\lambda\|_f^2
   &=(g_0g_1)^{42(-ab+a(a+1)/2)}
         \prod_{\epsilon=0}^{1}
           \prod_{i=1}^{a}(1-g_\epsilon^{-i})^{42},\\
 s_G&\le(g_0g_1)^{42ab}.
 \end{aligned}
\tag{TG24}
\]

Thus (TG22)--(TG23) apply to this actual prescribed finite compact type, with \(\eta=0\). The first-row modifications in K033 leave the e entries, finite stage characters, ranks, G-primes, and the operator itself unchanged. This fixed-operator hypothesis is essential for the asymptotic statement.

The proof of TG19 works without change on the h=4 host. When the analogous local finite Steinberg construction is made with its actual168-place record and G-primes, its parameters are the corresponding \(f=168\) local products. This local-product substitution does not silently claim that every global first-image assertion in the h=1 DO034 file has already been reconstructed for h=4.

There is a simple completely explicit sufficient level for either local product. Put \(Q=g_0g_1\), choose the smallest positive integer L with \(Q\le29^L\), and define

\[
 e_r=Lfab+\left\lceil\log_{29}\bigl(3(T_r+1)\bigr)\right\rceil,
                         \qquad M_r=29^{e_r}.
\tag{TG25}
\]

Here the ceiling means the least integer k with \(29^k\ge3(T_r+1)\); it can be computed using integer arithmetic. Then \(M_r\ge3s_G(T_r+1)\), proving (TG15). This bound intentionally keeps the auxiliary orbit size:

\begin{longtable}[]{@{}
  >{\raggedright\arraybackslash}p{(\columnwidth - 10\tabcolsep) * \real{0.1304}}
  >{\raggedleft\arraybackslash}p{(\columnwidth - 10\tabcolsep) * \real{0.1739}}
  >{\raggedleft\arraybackslash}p{(\columnwidth - 10\tabcolsep) * \real{0.1739}}
  >{\raggedleft\arraybackslash}p{(\columnwidth - 10\tabcolsep) * \real{0.1739}}
  >{\raggedleft\arraybackslash}p{(\columnwidth - 10\tabcolsep) * \real{0.1739}}
  >{\raggedleft\arraybackslash}p{(\columnwidth - 10\tabcolsep) * \real{0.1739}}@{}}
\toprule\noalign{}
\begin{minipage}[b]{\linewidth}\raggedright
Host baseline
\end{minipage} & \begin{minipage}[b]{\linewidth}\raggedleft
a
\end{minipage} & \begin{minipage}[b]{\linewidth}\raggedleft
b
\end{minipage} & \begin{minipage}[b]{\linewidth}\raggedleft
\(Q=g_0g_1\)
\end{minipage} & \begin{minipage}[b]{\linewidth}\raggedleft
L
\end{minipage} & \begin{minipage}[b]{\linewidth}\raggedleft
Fixed term \(Lfab\) in \(e_r\)
\end{minipage} \\
\midrule\noalign{}
\endhead
\bottomrule\noalign{}
\endlastfoot
K031, \(J^{14}\) & 42 & 1802 & \(4615409\cdot5276377\) & 10 & 31,787,280 \\
K032, \(J^{56}\), analogous local product & 168 & 1802 & \(22285932761\cdot64072175441\) & 15 & 762,894,720 \\
\end{longtable}

The large fixed exponent is a sufficient interpolation cost, not an assertion that the least useful level is this large. It is \(O(1)+O(\log T_r)\) for fixed source ranks and G-primes, but it must not disappear from a concrete MainX window.

For these two exact local products the certificate gives respectively \[
 (v_{11}I_a,\ v_{11}\Gamma_G,\ v_{11}d_G)
       =(24,336,0)\quad\hbox{and}\quad(408,21336,0).
\tag{TG25a}
\] The second tuple uses the24 actual real-field29 places of K032, rather than replacing their residue exponent by168. These are fixed finite normalization costs; the h=4 tuple by itself is not a new global extraction theorem.

On the h=1 host, the actual class-supported first extraction of DO034(25) now has a selected tuple with (TG22), and

\[
 \Theta_{\mathcal O_G\phi_\lambda,r_{P_r}\phi_\infty}(F_r)(1)
       =\frac{C^{14ab}a^{\rm aux}_{r,\lambda}}
                    {I_a(M)\,d_{\rho,r}\,d_G},\qquad
 \|F_r\|_2^2=\frac1{I_a(M)d_{\rho,r}d_G}.
\tag{TG26}
\]

The target vector is the primitive vector of K033(CS25), with its exact hook metric (CS36); the source highest vector used in the compact matrix coefficient is the matching source factor. The Cauchy argument already puts this possibly non-highest target vector in the single selected target summand. Compact extraction does not require it to be a target-highest vector.

Here \[
 I_a(M)=M^{42a^2}\prod_{i=1}^{a}(1-29^{-i})^6.
\tag{TG27}
\] Its11-valuation is independent of M; \(v_{11}(d_G)=0\). The changing compact dimension \(d_{\rho,r}\) has valuation \(O(\log(T_r+1))\), by the exact fixed-rank Weyl ratios in K033(CS40). Thus, after removing the displayed fixed period and retaining the stated algebraic normalization, the actual first coefficient in (TG26) has valuation slope \(1/1220\). Unit source and finite Schwartz normalization replaces the denominator by \(\sqrt{\Gamma_G}\sqrt{I_a(M)d_{\rho,r}d_G}\), exactly as in DO034(27); the fixed \(\Gamma_G\) and these compact corrections also preserve that slope.

This does not discard the archimedean Schwartz norm. It remains exactly K033(CS38), including \(A_r!B_r!\), the residual hook products, and the positive CM metric factors. Those real factorial costs are not \(O(\log T_r)\), and normalizing the full archimedean Schwartz vector changes the real cost in the complete gain/error accounting. The complete DP44 inequality is not inferred from the local slope.

The theorem applies to any other \textbf{specified nonzero fixed} auxiliary vector that has an actual finite Heisenberg expansion of the form (TG13). In particular it can be used with a literal raw G factor once its fixed expansion and the relevant compact extraction have been supplied. It does not identify a raw factor with the finite Steinberg projection by definition, and it does not assert that a later theta image contains this entire selected lattice.

\subsection{6. Scope, checks, and references}\label{scope-checks-and-references}

The classical finite-torsion flat-frame argument in Lemmas TG3/TG6 is written in full to fix the normalized connection, the integral root, and the actual reduced multiplier. The new application is the exact complete-module comparison TG10, the actual permitted auxiliary ideal identity TG19, and its sharp selected consequence TG22 on the enlarged CM family.

Carried accepted inputs are DP47/49/51/58/59/62 at the interfaces already specified in K018, K022, and K033. Exact local/source identifications are K018(K19.2/K19.13), K033(CS9/CS12/CS19/CS25/CS28/CS34), DO030(9--20), and DO034(12--29). No Lean execution is claimed. The simultaneous jet theorem and the accepted display bounds are used, not rederived.

Primary literature: Bannai, K., \& Kobayashi, S. (2010). Algebraic theta functions and the p-adic interpolation of Eisenstein--Kronecker numbers. \emph{Duke Mathematical Journal, 153}(2), 229--295. Author preprint: \href{https://arxiv.org/pdf/math/0610163}{arXiv:math/0610163v4}, exact locators given in section1. Those primary passages concern the reduced multiplier and algebraic Mumford translation; no positive-characteristic ordinary hypothesis from later sections is imported.

The companion exact certificate checks formal torsion-root and connection signs in finite rational jets, including the excluded d=11 case, and computes the displayed sufficient level constants and the actual fixed finite11-costs. These finite calculations check the indicated identities and arithmetic only. The infinite full-germ and image-ideal assertions are the proofs above, not inferences from a finite series sample.

\clearpage

\section{AM-R-018: complete numerical Sizes for the explicit fixed-prime selected families}\label{am-r-018-complete-numerical-sizes-for-the-explicit-fixed-prime-selected-families}

Author: /root, GPT-6 Astra Pro, ALP\_MACHINE, 22 November2026UTC. Reasoning effort is not externally exposed.

This result checks the remaining spread condition left explicit in R017. It gives all four clauses of the literal numerical \nolinkurl{Sizes} record on the K031/K032 baselines and on finite, explicitly bounded portions of the K033 growing selected family. It does not construct \nolinkurl{Dat.Ok}, the global \nolinkurl{SizeData₂}, or a net lift saving. In particular the common-type-field obstruction remains. The fixed T-primes cannot support an unbounded growing-weight family under this record.

\subsection{1. Literal source spread, including the rank term}\label{literal-source-spread-including-the-rank-term}

Use the actual m=2 data of K031/K032. Thus

\[
 d=t+1,\quad N_j=k+jd,\quad
 u_i=(2(k-i)+1)C+\tfrac12,\quad
 e_0=0,\quad e_1=4C.
 \tag{R18.1}
\]

Here k is42 or168, d is1760 or1634, and C is respectively 45156049293 or40530339687. The selected dominant weights are the actual \nolinkurl{lamTrue}, with their source determinant shifts. Write these weights as \(w_j\). The checked source recurrence is

\[
 \begin{aligned}
 w_{0,i}&=(2(k-i)+1)C+1-k/2+i,\\
 w_{1}&=(w_0-d/2,\ (k/2)^d),\\
 w_{2}&=(w_1-d/2,\ (-4C+(k+d)/2)^d).
 \end{aligned}
 \tag{R18.2}
\]

The notation \(x^d\) in a weight list means d repeated entries. These equations follow directly from \nolinkurl{lamTrue_old} and \nolinkurl{lamTrue_new} in \nolinkurl{Columns/BaseLemmas.lean}496--530, with \nolinkurl{hcToWeightQ} at stage zero. They are the already checked recurrence, used here to compute its extrema.

The source does \textbf{not} define spread as first minus last alone. \nolinkurl{LocalRingCardinality/DataCoeff.lean}34--35 defines it, for positive rank N, by

\[
              \operatorname{spread}(w)=w_0-w_{N-1}+N.
 \tag{R18.3}
\]

All rank terms in the following calculation are retained. The three exact baseline spreads are

\[
 \begin{aligned}
 S_0&=2(k-1)C+1,\\
 S_1&=(2k+1)C+1+d/2,\\
 S_2&=(2k+5)C+1+d/2.
 \end{aligned}
 \tag{R18.4}
\]

\textbf{Proof.} The entries of \(w_0\) decrease by \(2C-1>0\). Its last entry is \(3C+k/2\). In \(w_1\) the old last entry is \(3C+k/2-d/2>k/2\), since \(3C>d/2\), so the new constant block is the minimum. In \(w_2\) the smallest old entry is \(k/2-d/2\), which exceeds the new last entry \(-4C+(k+d)/2\) by \(4C-d>0\). At both steps the first entry remains the shifted old first entry. These observations also verify dominance at the block joins. Substituting these first and last entries and the ranks \(k,k+d,k+2d\) into (R18.3) gives(R18.4). Finally \(S_1-S_0=3C+d/2>0\) and \(S_2-S_1=4C>0\). Thus \(S_2\) is the largest baseline spread. ∎

If \(u_0\) at one embedding is increased by any integer \(n\ge0\), every weight in(R18.2) changes only in its first entry, which increases by n.~Dominance persists, and each literal spread is exactly \(S_j+n\). The same extrema hold for any interlaced changes of the middle u entries that retain the first and last entries and the same row data. Consequently the calculation applies simultaneously to any finite source word collection with those common extremes.

\subsection{2. The growing K033 family and an exact size parameter}\label{the-growing-k033-family-and-an-exact-size-parameter}

At K033's two distinguished embeddings the first u entry is increased by

\[
                   A_r=11^r,\qquad B_r=121\,11^r
                         \quad(r\ge0),
 \tag{R18.5}
\]

respectively. Other embeddings have no increase. The fixed row data, finite characters, G-primes, ranks and source interlacing remain those already checked in K033. The largest literal spread over all stages and embeddings is therefore exactly

\[
                        S(r)=S_2+B_r.
 \tag{R18.6}
\]

Choose the integral real size parameter

\[
                 a(r)=\left\lceil\frac{S_2+B_r}{2}\right\rceil.
 \tag{R18.7}
\]

This a is the \nolinkurl{Sizes} real parameter, not the source rank, the lift count, or the ordinary differential degree. The exact record \nolinkurl{AbelianHostFieldConditions.lean}486--494 is

\[
 \begin{gathered}
 2fk<g_i-1\ (i=0,1),\qquad 4a<T,\\
             2fN_2+1<4a,\qquad
 \operatorname{spread}(w_{j,\tau})\le2a\quad(0\le j\le2).
 \end{gathered}
 \tag{R18.8}
\]

The first inequalities were checked in K031/K032. The last follows exactly from(R18.6)--(R18.7). The middle lower bound holds already at r=0 by the exact values below and continues to hold because a(r) increases. Thus the only remaining condition on r is \(4a(r)<T\), with T the corresponding explicit prime of R017.

\begin{longtable}[]{@{}
  >{\raggedright\arraybackslash}p{(\columnwidth - 8\tabcolsep) * \real{0.1579}}
  >{\raggedleft\arraybackslash}p{(\columnwidth - 8\tabcolsep) * \real{0.2105}}
  >{\raggedleft\arraybackslash}p{(\columnwidth - 8\tabcolsep) * \real{0.2105}}
  >{\raggedleft\arraybackslash}p{(\columnwidth - 8\tabcolsep) * \real{0.2105}}
  >{\raggedleft\arraybackslash}p{(\columnwidth - 8\tabcolsep) * \real{0.2105}}@{}}
\toprule\noalign{}
\begin{minipage}[b]{\linewidth}\raggedright
Host
\end{minipage} & \begin{minipage}[b]{\linewidth}\raggedleft
\(S_0\)
\end{minipage} & \begin{minipage}[b]{\linewidth}\raggedleft
\(S_1\)
\end{minipage} & \begin{minipage}[b]{\linewidth}\raggedleft
\(S_2\)
\end{minipage} & \begin{minipage}[b]{\linewidth}\raggedleft
Exact permitted r interval
\end{minipage} \\
\midrule\noalign{}
\endhead
\bottomrule\noalign{}
\endlastfoot
\(J^{14}\) & 3,702,796,042,027 & 3,838,264,190,786 & 4,018,888,387,958 & \(0\le r\le55\) \\
\(J^{56}\) & 13,537,133,455,459 & 13,658,724,475,337 & 13,820,845,834,085 & \(0\le r\le90\) \\
\end{longtable}

For complete reproducibility, these endpoints are certified by the integer inequalities

\[
 \begin{split}
 121\,11^{R_*}&\le
       2\left\lfloor\frac{T-1}{4}\right\rfloor-S_2
                         <121\,11^{R_*+1},\\
 R_*&=55\text{ for }J^{14},\qquad
 R_*=90\text{ for }J^{56}.
 \end{split}
 \tag{R18.9}
\]

\textbf{Proof of the endpoints.} For positive integers S and T, \(4\lceil S/2\rceil<T\) is equivalent to \(S\le2\lfloor(T-1)/4\rfloor\). Substitute \(S=S_2+121\,11^r\), and use strict increase in r. The two exact boundary comparisons in(R18.9), all four clauses of(R18.8) throughout the stated intervals, and the failed next endpoint are verified by integer arithmetic in the saved certificate. The primality certificates for T are carried from R017; they are not replaced by a numerical primality screen here. ∎

At the first excluded r, even allowing a real size parameter cannot restore this particular fixed-T family. Indeed its spread condition forces \(a\ge S(r)/2\), and the certificate also verifies \(2S(r)\ge T\) there. The strict inequality \(4a<T\) is then impossible. This excludes only that fixed T and growing-weight choice. It does not exclude choosing a new T-prime as the weights grow.

The K033 growing degree first exceeds its fixed baseline degree at r=13 on \(J^{14}\) and r=15 on \(J^{56}\). Accordingly the nonempty ranges13--55 and15--90 also meet the entire numerical(R18.8). The large fixed valuation errors and all finite and real normalization costs in K033/K034 are unchanged; baseline domination is not a MainX gain/error inequality.

\subsection{3. Scope and exact sources}\label{scope-and-exact-sources}

R017 now has the missing \nolinkurl{spread_le} check for the stated baselines and these growing selected weights. The result supplies a complete numerical \nolinkurl{Sizes} witness, conditional only on the already specified weight/prime data. It does not make the primitive CM type satisfy \nolinkurl{Fields.Ok.quad}, identify the integral row and previous-image lattices, or compute the global DP44 budget. K033/K034's unbounded sharp-slope assertion concerns their selected section modules; its use with this fixed T is restricted to the explicit finite intervals just proved.

Exact stipulated source locators, pinned to commit \nolinkurl{7cbf7a0c32fca1badf05931879a37c3f8a2b03b4}, are \nolinkurl{Formalism/LocalRingCardinality/DataCoeff.lean}34--39, \nolinkurl{Formalism/LocalToGlobal/Columns/BaseLemmas.lean}467--530, and \nolinkurl{Formalism/LocalToGlobal/Columns/AbelianHostFieldConditions.lean} 443--494. K031 HC3--10, K032 QH rank/weight data, K033 CS20--35 and R017 R17.1--5 are carried at their recorded checked scope. The script is \nolinkurl{scripts/selected_fixed_prime_sizes_certificate.py}; its complete arithmetic output is in \nolinkurl{checks/selected_fixed_prime_sizes_certificate.json}. This is a mathematical source/parameter verification, not a Lean build. Whole-proof readings are recorded separately from the finite certificate.

\clearpage

\section{ALP-LIFT-23: exact Fourier contraction of the prescribed coupled Q family}\label{alp-lift-23-exact-fourier-contraction-of-the-prescribed-coupled-q-family}

Author: GPT-6 Astra Pro, ALP\_MACHINE /root/lift\_accounting. Date: 22 November 2026 UTC. Inherited reasoning effort is not exposed.

This finding evaluates the entire literal local Y sum in the two coupled finite kernels for m=2 at a degree-one Q place. The auxiliary depth disappears. The resulting finite function is nonzero, has an explicit Gauss-sum scalar, and has primitive content at every coefficient place over11 when q is different from11. This is an identity of the specified finite Schwartz functions. It is not a claim that the second automorphic theta integral is injective on the preceding image.

\subsection{1. Literal source data and conventions}\label{literal-source-data-and-conventions}

Use the pinned source tree work/research/alpoge\_source at commit 7cbf7a0c32fca1badf05931879a37c3f8a2b03b4. Put \[
 m=2,\quad r=k+1,\quad d=t+1,\quad b=k+d=r+t,\quad N=b+d.
 \tag{Q1}
\] The first edge has ranks k,b, and the second edge has ranks b,N. Let Q be a relative split place whose local field F is the rational q-adic field, so q is a rational prime and a uniformizer. Write O for its ring of integers. This explicit unramified degree-one hypothesis is retained; residue degree one alone would not imply it.

This hypothesis holds at every literal tuple Q place. As independently audited by /root/definite\_oscillator, Tuples.lean64--67 defines DegOne by prime absolute norm q together with q not belonging to the square of the prime ideal, and GoodPl adds relative splitting. The first condition gives residue degree one; the second gives ramification index one, since the local valuation of q is that index. The split CM completion is therefore the rational q-adic field. The tuple condition in Tuples.lean112--114 kills inertia outside GoodPl, while Qprime128 is defined by failure to kill inertia. Thus Qprime is contained in GoodPl. The larger host's ramified29 place is excluded from this Q set. These exact source clauses were read independently here in \nolinkurl{Formalism/LocalToGlobal/Divisibility/Tuples.lean.}

Take the actual standard frames and the forced initial and terminal frame choices of ChIdx.Ok.zero\_h and last\_h. The limited-depth character is trivial on O and nontrivial on q\^{}\{-1\}O. Its induced nontrivial residue additive character is \[
             \psi_0(z)=\psi(q^{-1}\widetilde z).
 \tag{Q2}
\] Only the finite argument range prescribed by LocCst.Ok is used.

Let chi be the actual tame residue character of the tuple at Q. We first assume chi is nontrivial, and extend it by zero at0. The trivial-character version is given in section5. Let \(\mathsf M=\mathrm{cst.M}\ge1\), with depth at least \(2\mathsf M\). The actual local parameter set, by the already checked DO032, is \[
       Y_Q=q^{-\mathsf M}O^b/q^{\mathsf M}O^b.
 \tag{Q3}
\] Representatives may be changed within these cosets because both literal factors are unchanged by that change.

The exact source locators are:

\begin{itemize}
\tightlist
\item
  \nolinkurl{Formalism/LocalToGlobal/Columns/FiniteSumNonzeroLargeDepth.lean} (FS): entB684--694; unitV669; PhiQ711--715; omPhi939--941; joinCol982--983 and PsiAt985--986; zero\_h1020, last\_h1025; XiCpl11134--1138; PhiTop1142--1147 and ftY1150--1156; XiCpl21160--1167.
\item
  The same file: ybSet1125--1129, and the actual coupled Y action and CplKer in1216--1237.
\item
  Accepted and checked DO031.3--DO031.5 and DO032.2--DO032.10 give the literal terminal last-row normalization and finite quotient, including the source scalar \(q^{-\mathsf M b}\).
\item
  The Q middle group is the reduction preimage of \(P_{r,t}\), by Jat909--913, last\_h, and grpQ1062--1066. It does not generally fix an individual chosen class \(y=e_k\). The present calculation retains the whole prescribed family.
\end{itemize}

The actual terminal last row in these frames is denoted by \(\eta\in F^{1\times b}\); any piT unit already absorbed by the source frame remains absorbed here. No further coordinate rescaling is introduced.

\subsection{2. The two literal functions}\label{the-two-literal-functions}

Write the first matrix as \[
                 X=\begin{pmatrix}A\\C\end{pmatrix},
 \qquad A\in M_{r,k}(F),\quad C\in M_{t,k}(F),
 \qquad y=\begin{pmatrix}z\\w\end{pmatrix}.
 \tag{Q4}
\] Because the Q exponent in lamYat is zero and the two relevant frames are the identity, entB of joinCol is precisely the ordinary matrix \([X,y]\), with the added column in position k. Thus the first raw factor is exactly \[
 F_X(y)=
 \mathbf1_{M_{r,k}(O)}(A)\mathbf1_{qM_{t,k}(O)}(C)
 \mathbf1_{O^r}(z)\mathbf1_{qO^t}(w)
       \chi\!\left(\det[\bar A,\bar z]\right).
 \tag{Q5}
\] In particular it is zero when \(\bar A\) has rank smaller than k. This includes all support conditions, including the lower auxiliary column condition; it is not just the specialization \(y=e_k\).

Write the second matrix as \(\binom Z\eta\), with \(Z\in M_{N-1,b}(F)\). Its terminal factor is \[
 H(Z)\,q^{-\mathsf M b}
 \mathbf1_{q^{-\mathsf M}O^b}(\eta)\psi(-\eta y),
 \qquad H(Z)=\Phi_{\mathrm{Top},Q}(Z).
 \tag{Q6}
\] Here H has all of the literal PhiQ conditions: Z is integral, its bottom-left \((N-1-r)\)-by-r block is divisible by q, and its leading r-by-r determinant is a unit, on which chi is evaluated.

The actual local contraction over its entire allowed Y set is therefore \[
 {\cal C}_{\mathsf M}(X,Z,\eta)
   =H(Z)q^{-\mathsf M b}
     \mathbf1_{q^{-\mathsf M}O^b}(\eta)
        \sum_{y\in Y_Q}F_X(y)\psi(-\eta y).
 \tag{Q7}
\] The source writes a finite sum, not a Haar integral over Y. The scalar in (Q7) is exactly its ftY scalar.

\subsection{3. Closed formula and independence of depth}\label{closed-formula-and-independence-of-depth}

For a matrix \(\bar A\) of rank k define its nonzero cofactor row \(\ell_A\in\kappa^{1\times r}\) by \[
                  \ell_A z=\det[\bar A,z].
 \tag{Q8}
\] Put \[
                    G(\chi,\psi_0)=
                       \sum_{s\in\kappa^\times}\chi(s)\psi_0(s).
 \tag{Q9}
\] If the support conditions on A,C in (Q5) hold and \(\bar A\) has rank k, then the exact answer is \[
 \boxed{
 {\cal C}_{\mathsf M}(X,Z,\eta)=
 \begin{cases}
 H(Z)\,q^{-d}G(\chi,\psi_0)\chi(-c)^{-1},
 &\eta\in q^{-1}O^b,\quad
   \overline{q\eta_{\mathrm{top}}}=c\ell_A,\quad c\in\kappa^\times,\\
 0,&\text{otherwise}.
 \end{cases}}
 \tag{Q10}
\] It is also zero when those A,C support or rank conditions fail. There is no condition on \(\overline{q\eta_{\mathrm{low}}}\). In particular the entire answer is independent of \(\mathsf M\ge1\).

Proof. Because of (Q5), the contributing representatives in (Q3) are exactly \[
       z\in O^r/q^{\mathsf M}O^r,\qquad
       w\in qO^t/q^{\mathsf M}O^t.
\] The determinant character depends only on \(\bar z\). The sum over all higher lifts in both blocks is zero unless \(\eta\in q^{-1}O^b\). When this condition holds, those sums contribute \(q^{(\mathsf M-1)b}\). The lower residue block is zero. Multiplication by the source scalar \(q^{-\mathsf M b}\) therefore leaves \[
       H(Z)q^{-b}\sum_{z\in\kappa^r}
            \chi(\ell_Az)\psi_0(-az),
       \qquad a=\overline{q\eta_{\mathrm{top}}}.
 \tag{Q11}
\] This argument also works for \(\mathsf M=1\): there are no higher lifts, and the support of ftY is already \(q^{-1}O^b\).

If \(a\) is not a scalar multiple of \(\ell_A\), there is a vector \(z_0\in\ker\ell_A\) with \(az_0\ne0\). Translation by a suitable multiple of \(z_0\) preserves the determinant factor and multiplies the character sum by a scalar different from1. The sum is zero. If \(a=c\ell_A\), every fibre of \(\ell_A\) has \(q^k\) elements, so the sum equals \[
           q^k\sum_{s\in\kappa^\times}\chi(s)\psi_0(-cs).
 \tag{Q12}
\] For \(c=0\), it vanishes because chi is nontrivial. For \(c\ne0\), substitute \(u=-cs\); the answer becomes \(q^k\chi(-c)^{-1}G(\chi,\psi_0)\). Since \(b-k=d\), (Q10) follows. Every manipulation takes place inside the specified finite additive quotient, so no unrestricted extension of the source psi has been assumed.

A concrete nonzero value is obtained with \(A=\binom{I_k}{0}\), C=0, \(\eta=q^{-1}e_k^\vee\), and Z having identity leading b-by-b block and all remaining rows zero. Then \(\ell_A=e_k^\vee\), H(Z)=1, and \[
      {\cal C}_{\mathsf M}(X,Z,\eta)
                       =q^{-d}G(\chi,\psi_0)\chi(-1)^{-1}\ne0.
 \tag{Q13}
\]

\subsection{4. Exact norms, content, and height of this scalar}\label{exact-norms-content-and-height-of-this-scalar}

All measures in this section give integral additive matrix lattices volume1. The standard finite Gauss identity is proved here in the needed normalization. For \(c\in\kappa\) put \(G_c=\sum_{s\ne0}\chi(s)\psi_0(cs)\). Then \(G_0=0\), and substitution shows all \(|G_c|\), \(c\ne0\), are equal. Additive orthogonality gives \(\sum_c|G_c|^2=q(q-1)\). Hence \[
                       |G(\chi,\psi_0)|^2=q.
 \tag{Q14}
\] This proof applies after every complex embedding of the finite coefficient field, since both induced characters remain nontrivial.

For fixed supported X of full residue rank, the nonzero eta set in (Q10) has measure \((q-1)q^t\): there are \(q-1\) top fractional classes and \(q^t\) arbitrary lower fractional classes, each followed by the integral lattice of volume1. Thus \[
 \int_F\cdots\int_F
   \left|{\cal C}_{\mathsf M}(X,Z,\eta)\right|^2\,d\eta
       =|H(Z)|^2(q-1)q^{-d}.
 \tag{Q15}
\] This also equals \(|H(Z)|^2\int|F_X(y)|^2dy\), since the first support has lower volume \(q^{-t}\) and the determinant is nonzero on the proportion \((q-1)/q\) of its top residue classes. Thus the depth-independent contraction is exactly the additive Fourier transform in this variable, including its normalization.

For the finite counting measure on (Q3), the number of nonzero first values is \[
                 (q-1)q^{\mathsf M b-d}.
 \tag{Q16}
\] The number of nonzero transformed values, modulo \(q^{\mathsf M}O^b\), is \((q-1)q^{t+\mathsf M b}\). Their squared amplitude \(q^{1-2d}\) gives the same number (Q16), as required by the actual unitary finite Fourier matrix.

The upper terminal factor has norm \[
 \Gamma_{\mathrm{top},Q}
  =q^{-r(N-1-r)}\prod_{i=1}^r(1-q^{-i}).
 \tag{Q17}
\] The probability that an r-by-k residue matrix has rank k is \(\prod_{i=0}^{k-1}(1-q^{i-r})\). Integrating also in X and Z gives the exact joint Schwartz norm \[
 \boxed{
 \|{\cal C}_{\mathsf M}\|_2^2
   =q^{-tk}(q-1)q^{-d}
       \prod_{i=0}^{k-1}(1-q^{i-r})\,
       \Gamma_{\mathrm{top},Q}.}
 \tag{Q18}
\] This is the norm of the contracted raw array in its product matrix variables; it is not the norm of an automorphic composite operator.

Both G and its complex conjugate are algebraic integers, and (Q14) is the algebraic identity \(G\overline G=q\). Consequently G is a unit at every coefficient place above11 if \(q\ne11\). Every nonzero value in (Q10) is then an11-adic unit. Its minimum valuation is exactly zero, by (Q13). No division by \(q-1\), a finite general linear group order, or a class mass has been introduced.

The magnitude of the scalar in (Q10) is \(q^{1/2-d}\). Its absolute logarithmic Weil height is exactly \[
                  h\!\left(q^{-d}G(\chi,\psi_0)\right)
                               =(d-\tfrac12)\log q.
 \tag{Q19}
\] Indeed its inverse is \(q^{d-1}\overline G\), an algebraic integer, and every complex conjugate of that inverse has absolute value \(q^{d-1/2}>1\). Computing the height of this algebraic integer and using \(h(x^{-1})=h(x)\) proves (Q19). This is the height of the specified finite scalar. It must not be renamed the source global H.

\subsection{5. Trivial residue character and exact scope}\label{trivial-residue-character-and-exact-scope}

If chi is trivial, still extended by zero at0, (Q11) remains valid. The value is zero unless \(a=c\ell_A\), and its two remaining values are \[
 {\cal C}_{\mathsf M}=
 \begin{cases}
 H(Z)q^{-d}(q-1),&c=0,\\
 -H(Z)q^{-d},&c\ne0.
 \end{cases}
 \tag{Q20}
\] The norm (Q15) remains the same. The main primitive-value conclusion remains true for \(q\ne11\), because the \(c\ne0\) value is a unit, although the \(c=0\) value can have positive valuation if11 divides \(q-1\). The nontrivial-character height formula(Q19) is not asserted for the trivial case.

The full allowed global Y set factors into these local finite quotients and the G quotients by the already checked CRT statement DO032.2. Therefore a finite product of the literal Q contractions is exactly the corresponding partial global Y sum. At sharp G, LIFT22 shows that the first factor leaves just its single forced local class. These statements remove cancellation in the raw finite array and retain the exact Q Fourier scalar and its norms.

They do not show that this raw function survives the next automorphic integral. At Q the two stages have the same tuple residue character, so the all-ramified boundary exclusion used in LIFT21 at G is not available. The fixed-Q-class first theta function is also not generally in the full prescribed middle parabolic type by itself: that covariance belongs to the coupled family. An argument for the second image must retain this family or supply the exact local quotient comparison. No claim about a globally pinned scalar, the iterated integral row lattice, the complete D7 data, or a DP44 gain is made.

\subsection{6. Checks}\label{checks}

The exact finite replay is \nolinkurl{raw_coupled_q_fourier_certificate.py.} It compares the literal finite Y sum with (Q10)/(Q20), including the nontrivial M=2 higher-depth cancellation, in cyclotomic integer coordinates. It also checks the support counts, Parseval factors, the cofactor orientation, and the unit normalization. Small finite examples are diagnostics of the identity, not full Ranks data. Whole-proof reading and replay events are recorded separately.

\clearpage

\section{AM-DO-036: the unchanged first raw kernel survives at maximal 29-level}\label{am-do-036-the-unchanged-first-raw-kernel-survives-at-maximal-29-level}

Author: GPT-6 Astra Pro, ALP\_MACHINE \nolinkurl{/root/definite_oscillator}. 22 November 2026 UTC. Inherited reasoning effort is not separately exposed.

The first raw T factor has a uniform nonzero local test on every unitary source vector with its prescribed mirabolic invariance when \nolinkurl{b>=4a}. The proof separates its compact term from an explicitly bounded nonunit-determinant tail. The raw G and Q tests are exact compact projectors. Together with the checked spherical tests elsewhere and AM-R-017's actual class-supported input, these facts prove nonvanishing with the \textbf{unchanged full first G/T/Q finite kernel and the maximal vacuum at29}. No auxiliary29 Heisenberg translate is inserted.

The statement uses the new mixed theta realization and its actual stage characters. It does not instantiate the unchanged elliptic \nolinkurl{VacIn}, the source's \nolinkurl{Fields.Kq} record, every \nolinkurl{ChIdx.Ok} clause, an integral prior image lattice, or the MainX terminal-mass comparison. A positive target quotient mass remains in the exact norm identity.

\subsection{1. Exact local convention and the full-rank period operator}\label{exact-local-convention-and-the-full-rank-period-operator}

Let F be a nonarchimedean characteristic-zero local field, O its ring, q its residue cardinality, and

\[
 G=\operatorname{GL}_a(F),\quad K=\operatorname{GL}_a(O),\quad
 K_b=\operatorname{GL}_b(O),\quad b\ge4a,\quad z=b/2.
 \tag{DO36.1}
\]

Normalize group Haar by \nolinkurl{vol(K)=vol(K_b)=1} and matrix Haar by \nolinkurl{vol(M_(b,a)(O))=1}. Let eta be the actual unitary local splitting component. The genuine Hom Weil action is

\[
 [\omega(g)f](X)=\eta(\det g)^{-1}|\det g|^{b/2}f(Xg).
 \tag{DO36.2}
\]

This is the source formula of AM-DO-020/024 and ALP-LIFT-20, and the common global character is exactly \nolinkurl{mu_0 xi}, as checked in DO035 from \nolinkurl{LevChiEqOne.lean}61,100. In particular it is unramified at T, has residue chi at Q, and is tame with the prescribed residue character at G.

Let pi be any unitary representation of G on a Hilbert space V. Use conjugate-linear-first inner products. Put \nolinkurl{X_0=[I_a;0]} and define the Bochner integral

\[
 B_f(v;k)=\int_G f(kX_0h)|\det h|^{b/2}
                   \eta(\det h)^{-1}\pi(h)v\,dh,
                       \qquad k\in K_b.
 \tag{DO36.3}
\]

For the raw functions below \nolinkurl{|f|<=1_(M_(b,a)(O))}. Since \nolinkurl{kX_0} has an integral left inverse, \nolinkurl{f(kX_0h)!=0} implies \nolinkurl{h in M_a(O)}. The convergence and all interchanges follow from the exact bound proved in Section 2.

The exponent is \textbf{b/2}, with no unrecorded modular-character shift. Directly, for \nolinkurl{g in G}, simultaneous replacement of \nolinkurl{f,v} by \nolinkurl{omega(g)f,pi(g)v} gives the integral with variable \nolinkurl{hg}. The determinant powers cancel after \nolinkurl{t=hg}, and

\[
 \eta(\det g)^{-1}\eta(\det(tg^{-1}))^{-1}
                         =\eta(\det t)^{-1}.
\]

Hence

\[
                 B_{\omega(g)f}(\pi(g)v;k)=B_f(v;k).
 \tag{DO36.4}
\]

This is the simultaneous coinvariant convention for the literal bilinear theta lift. It checks both the character sign and the half-determinant exponent without importing a different local induced-model normalization.

Write

\[
 C_{a,b}(q)=\prod_{i=0}^{a-1}(1-q^{i-b}).
 \tag{DO36.5}
\]

The exact full-rank disintegration proved in LIFT16(S6) is

\[
 \int_{M_{b,a}(F)}H(X)dX
   =C_{a,b}(q)\int_{K_b}\int_G
                    H(kX_0h)|\det h|^b\,dh\,dk.
 \tag{DO36.6}
\]

For clarity, every full-rank X has such a presentation because its column space intersected with \nolinkurl{O^b} is a saturated direct summand. The two measures have the same right Jacobian; their constant is the probability that a residue \nolinkurl{b by a} matrix has rank a, which is precisely (DO36.5). The rank-deficient set has additive measure zero. These facts also give a direct proof of (DO36.6) in the present convention.

Substitution into the squared norm of (DO36.3), with the second integration variable written \nolinkurl{t=hg}, now proves the exact identity

\[
 \boxed{\quad
 Z_f(v):=\int_G\langle f,\omega(g)f\rangle
                    \langle v,\pi(g)v\rangle\,dg
       =C_{a,b}(q)\int_{K_b}\|B_f(v;k)\|^2dk\ge0.
 \quad}
 \tag{DO36.7}
\]

Indeed the two eta factors on the squared-norm side become \nolinkurl{eta(det h) eta(det t)^(-1)=eta(det g)^(-1)}; the two determinant powers become \nolinkurl{|det h|^b |det g|^(b/2)}, and the vector inner product is \nolinkurl{<v,pi(g)v>}. Equation (DO36.6) gives the other matrix coefficient. With the opposite convention for writing both matrix coefficients, simultaneous \nolinkurl{g -> g^(-1)} gives the same real value. This is the actual local doubling test after the Hom/tensor and bilinear/conjugated-input transport of DO028.

\subsection{2. Exact integrability and the uniform tail inequality}\label{exact-integrability-and-the-uniform-tail-inequality}

For \nolinkurl{z>a-1}, define

\[
 D_a(z;q)=\int_{G\cap M_a(O)}|\det h|^zdh
          =\prod_{i=0}^{a-1}(1-q^{i-z})^{-1},
 \qquad E_a(z;q)=D_a(z;q)-1.
 \tag{DO36.8}
\]

Here is a self-contained normalization check. For additive matrix measure, let \nolinkurl{J_a(s)=int_(M_a(O)) |det X|^s dX}. The first column has exact common valuation n with measure \nolinkurl{q^(-an)(1-q^(-a))}. A primitive first column is carried to \nolinkurl{e_1} by an integral invertible matrix. The remaining columns then contribute a free first row and an \nolinkurl{(a-1)} square matrix. Therefore

\[
 J_a(s)=\frac{1-q^{-a}}{1-q^{-a-s}}J_{a-1}(s),\qquad J_0(s)=1.
\]

The multiplicative Haar measure is \nolinkurl{dh=prod_(i=1)^a(1-q^(-i))^(-1)|det h|^(-a)dX}, because K has the displayed additive measure. Substitute \nolinkurl{s=z-a} to obtain (DO36.8). The convergence condition is \nolinkurl{s>-1}, equivalently \nolinkurl{z>a-1}. In particular, \nolinkurl{E_a} is exactly the integral over integral matrices with nonunit determinant, not a heuristic count of Cartan double cosets.

For \nolinkurl{x_i=q^(i-z)} put

\[
 S_a(z;q)=\sum_{i=0}^{a-1}x_i
                =q^{-z}\frac{q^a-1}{q-1}.
\]

When \nolinkurl{z>=2a},

\[
 q^a S_a(z;q)\le\frac{1-q^{-a}}{q-1}<1.
\]

The elementary inequality \nolinkurl{prod(1-x_i)>=1-sum x_i} consequently gives

\[
 \boxed{\quad
 0<E_a(z;q)\le\frac{S_a(z;q)}{1-S_a(z;q)}
                 <\frac1{q^a-1}\qquad(z\ge2a).
 \quad}
 \tag{DO36.9}
\]

This includes q=2. It is strictly stronger than needed at the actual T primes of R017.

The bound \nolinkurl{||B_f(v;k)||<=D_a(z;q)||v||} proves the absolute convergence in Section 1. For the Fubini step in (DO36.7), take absolute values before integrating h,t,k; the result is bounded by \nolinkurl{C_(a,b) D_a(z;q)^2 ||v||^2}. Thus the positive identity is proved by absolutely convergent integrals, with no meromorphic continuation or local theta-occurrence theorem substituted for this fixed-vector calculation.

\subsection{3. The untouched T incidence factor has a positive test}\label{the-untouched-t-incidence-factor-has-a-positive-test}

Use the actual source and target residue covectors \nolinkurl{u!=0,v_0!=0} and the raw T function from DO035:

\[
       f_T(X)=1_{M_{b,a}(O)}(X)1_{\{v_0\bar X=u\}}.
 \tag{DO36.10}
\]

Its prescribed source compact is

\[
 J_T=\{g\in K:u\bar g=u\},\qquad
                     \operatorname{vol}(J_T)=\frac1{q^a-1}.
 \tag{DO36.11}
\]

Choose a primitive integral matrix X* of rank a with \nolinkurl{v_0 X*=u} modulo q. Such a choice exists because b\textgreater a: after an integral target-coordinate change make \nolinkurl{v_0} the last-coordinate functional, use \nolinkurl{I_a} in the first a rows, u in the last row and zero elsewhere. Extend its primitive columns to \nolinkurl{k* in K_b}, so \nolinkurl{k* X_0=X*}.

Let \nolinkurl{v in V^(J_T)}. The compact part of (DO36.3) is exactly

\[
 \int_K f_T(X*h)\eta(\det h)^{-1}\pi(h)v\,dh
                   =\frac{v}{q^a-1}.
 \tag{DO36.12}
\]

Indeed eta is unramified at T, the condition is precisely \nolinkurl{h in J_T}, and pi fixes v there. The remaining integral has norm at most \nolinkurl{E_a(b/2;q)||v||}, by unitarity of pi and eta. Define

\[
 \Delta_T(a,b;q)=\frac1{q^a-1}-E_a(b/2;q)>0.
 \tag{DO36.13}
\]

Then \nolinkurl{||B_(f_T)(v;k*)||>=Delta_T||v||}. The same equality of the integrands holds for every k congruent to k* modulo q: the raw function depends only on reduction inside integral matrices, and every such primitive matrix has an integral left inverse. This open subset has measure \nolinkurl{|GL_b(F_q)|^(-1)}. Equation (DO36.7) proves

\[
 \boxed{\quad
 Z_{f_T}(v)\ge
 \frac{C_{a,b}(q)}{|\operatorname{GL}_b(\mathbf F_q)|}
       \Delta_T(a,b;q)^2\|v\|^2>0\quad(v\ne0).
 \quad}
 \tag{DO36.14}
\]

This uses the raw incidence function, not the character-filtered auxiliary vector of DO035. A nontrivial T diamond character of the source input is fully compatible with its \nolinkurl{J_T} invariance; the lemma makes no claim that a T coefficient is spherical.

\subsection{4. The raw G and Q tests are exact compact projectors}\label{the-raw-g-and-q-tests-are-exact-compact-projectors}

At a G place let U be the upper unitriangular subgroup of \nolinkurl{GL_a(F_q)}, let psi\_U be the exact source generic character, and let \nolinkurl{J_G} be its full reduction preimage in K. ALP-LIFT-20 identified the actual first raw G/GS function, for the admitted class \nolinkurl{y=q^(-1)e_a mod O^b}, as

\[
 f_G(X)=\psi_U(\bar h)
       \quad\text{if }X\in M_{b,a}(O),\ \bar X=X_0\bar h,
                        \ \bar h\in U,
 \tag{DO36.15}
\]

and zero otherwise. If \nolinkurl{pi(j)v=psi_U(bar j)^(-1)v} for j in J\_G, then

\[
 B_{f_G}(v;1)
      =\frac{|U|}{|\operatorname{GL}_a(\mathbf F_q)|}\,v.
 \tag{DO36.16}
\]

To check this, \nolinkurl{f_G(X_0h)!=0} forces \nolinkurl{h in J_G}; its determinant has residue one, so the tame eta factor is one. The two generic characters cancel. There is no nonunit-determinant tail, since every residue matrix in U is invertible. The same equation holds on \nolinkurl{K_b(1)}, giving

\[
 Z_{f_G}(v)\ge
 \frac{C_{a,b}(q)}{|\operatorname{GL}_b(\mathbf F_q)|}
   \left(\frac{|U|}{|\operatorname{GL}_a(\mathbf F_q)|}\right)^2
                      \|v\|^2.
 \tag{DO36.17}
\]

This works on the whole actual inverse-generic source space; it is not limited to one chosen Steinberg multiplicity line.

At Q use (DO35.8), the same admissible global \nolinkurl{ybSet} class, and the actual label \nolinkurl{eta=mu_0 xi}. For every K-fixed vector,

\[
                      B_{f_Q}(v;1)=v.
 \tag{DO36.18}
\]

Here nonzero support forces \nolinkurl{h in K}, and the factor \nolinkurl{chi(det h)} in the raw function cancels \nolinkurl{eta(det h)^(-1)}. The same open \nolinkurl{K_b(1)} argument gives

\[
             Z_{f_Q}(v)\ge
              \frac{C_{a,b}(q)}{|\operatorname{GL}_b(\mathbf F_q)|}
                                       \|v\|^2.
 \tag{DO36.19}
\]

Thus all the first raw factors have a strictly positive fixed test on their actual initial source spaces. Finite nonzero Schwartz functions alone would not imply these statements; the projector and tail calculations are essential.

\subsection{5. Maximal vacua at every other finite place}\label{maximal-vacua-at-every-other-finite-place}

At a split place outside G/T/Q, eta is unramified and \nolinkurl{f=phi_0=1_(M_(b,a)(O))}. For every K-fixed v, (DO36.3) has compact term v and tail at most \nolinkurl{E_a(b/2;q)||v||}, independently of k. Hence

\[
 Z_{\phi_0}(v)\ge
       C_{a,b}(q)(1-E_a(b/2;q))^2\|v\|^2.
 \tag{DO36.20}
\]

In particular this applies at all the relative split29 places, including their absolute ramification index seven, and at the dyadic split places. The trace coefficient and additive character are the actual ones; the unitary identification with this Schrodinger model is DO022/034. Changing the absolute base-field ramification does not change the local residue-field calculation.

At an odd nonsplit place the extension is unramified, the even identity lattices are integrally hyperbolic, and the stage character is trivial: it is unramified and trivial on the base field. The actual fixed vacuum calculation LIFT16(S19--20), from Liu's unramified doubling section, gives on a unit spherical vector

\[
 Z_v= L_v(s_*,\pi_v)\prod_{i=1}^a(1-(-1)^i q_v^{-d-i}),
 \quad d=b-a,\quad s_*=(d+1)/2.
 \tag{DO36.21}
\]

The additive-character/codifferent rescaling and the actual lattice section were verified there and in DO022; no new choice of local section is made here. Macdonald's bounded spherical parameter criterion, in the form checked in LIFT14(E18), gives \nolinkurl{|alpha|<=Q^((a-1)/2)} with \nolinkurl{Q=q_v^2}. Consequently, putting \nolinkurl{x_v=q_v^(-(b-2a+2))}, a uniform positive lower bound is

\[
 Z_v\ge (1+x_v)^{-a}
                   \prod_{i=1}^a(1-q_v^{-d-i})>0.
 \tag{DO36.22}
\]

The local L-value is positive in this region. The unitary spherical parameter is stable under conjugate-inverse, and the unramified unitary-group standard parameter is stable under inverse. Thus its eigenvalue multiset is conjugation-stable. Nonreal conjugate factors have positive product, and every real factor \nolinkurl{1-alpha Q^(-s_*)} is positive because its absolute correction is less than one. The lower bound follows from \nolinkurl{|1-alpha Q^(-s_*)|<=1+x_v}. This also proves finiteness without a temperedness assumption.

The bounds (DO36.20), (DO36.22) differ from one by a summable quantity as the finite places vary. In the split case \nolinkurl{E_a=O_a(q^(-(b/2-a+1)))}, with exponent at least a+1\textgreater=2; in the nonsplit case the exponent is at least2a+2. Their infinite product is therefore strictly positive. This proves the needed fixed-vacuum assertion at29 as part of the same global family, without any later change of finite input.

\subsection{6. A direct convergent global norm identity in the correct range}\label{a-direct-convergent-global-norm-identity-in-the-correct-range}

Let E/F be either actual mixed host of K031 or K032, and \nolinkurl{G=U(I_a)}, \nolinkurl{H=U(I_b)}, with the positive real signatures, the exact common stage character, and the actual CM vacuum. The fields and trace coefficients are the ones proved in K030/K032 and DO034; this is not an extension of the old elliptic constructor. Both groups are anisotropic over F. For the two concrete first edges,

\[
 (a,b)=(42,1802)\quad\text{or}\quad(168,1802),
                              \qquad b\ge4a.
 \tag{DO36.23}
\]

Take the complete first raw finite vector from DO035 with the single admitted global Y class and the source/target T decorations of R017. At29 and every other place outside G/T/Q keep the maximal vacuum. At infinity take the actual selected highest polynomial P\_nu from DO034, with its source vector w and its true positive norm N\_infty. Choose R017's real matrix coefficient from this same w, so the input lies on the paired contragredient line at each real place. This is a choice of the matched vector in the known selected summand, not an inference that any source vector pairs with a rigid highest polynomial.

The source function F is now the \textbf{unchanged actual function} of R017, with the inverse character-name convention explained in the independent review of that finding. It is nonzero, maximal at29, has the initial inverse G character, the T mirabolic invariance and zero diamond sum, and the full trivial Q level. Its exact norm for maximal finite compact volumes one and real compact volume one is

\[
                    \|F\|^2=\frac1{28d_T^f d_\rho d_G}.
 \tag{DO36.24}
\]

\subsubsection{Correct literature interface}\label{correct-literature-interface}

GQT Theorem11.4 assumes \nolinkurl{a<b<=2a}; it must \textbf{not} be cited directly for (DO36.23). Instead use GQT Theorem7.1, printed p.33: for the anisotropic target, its Siegel--Weil formula holds in every dimension, with constant \nolinkurl{c_(b,a)=1} when \nolinkurl{s_0=(b-a)/2>0}. The doubling construction of section11.3, printed p.52, and the unfolding in sections11.5--11.6, printed pp.53--54, are the same absolutely convergent identities in this range. The character/conjugation transport is exactly DO028. No continuation through a pole is needed.

There is an especially direct unfolding here. The source Hermitian space is anisotropic. Every rational maximal isotropic subspace of its doubled space \nolinkurl{S plus (-S)} has zero intersection with either summand. Its two projections are therefore isomorphisms, so it is the graph of a rational isometry. Hence the rational double cosets in the doubling unfolding have just the open graph orbit, with stabilizer the diagonal G. Unfolding the Eisenstein sum is consequently the ordinary integral of a matrix coefficient against the right regular automorphic coefficient; there are no other rational orbits to dispose of.

Let dg be the chosen source Haar measure, let dh be the chosen target Haar measure, and put \texttt{V\_H=vol\_dh(H(F)\textbackslash{}\textbackslash{}H(A))}. Use the source inner product with dg. The resulting identity is

\[
 \boxed{\quad
 \|\Theta_\phi^{\rm bil}F\|_{dh}^{2}
       =V_H\int_{G(\mathbb A)}
            \langle\phi,\omega(g)\phi\rangle
            \langle F,R(g)F\rangle_{dg}\,dg.
 \quad}
 \tag{DO36.25}
\]

Here \nolinkurl{Theta^bil} is the integral against F, as in the literal source. For Tamagawa target measure \nolinkurl{V_H=2}; this is the factor \nolinkurl{[E:F]} in the doubled identity. Indeed GQT7.1 uses the normalized target theta integral and states \nolinkurl{tau(H)/[E:F]=1} for the present unitary group. Thus the unnormalized target integral is \nolinkurl{[E:F]} times the Siegel--Weil Eisenstein value. Scaling dh scales both sides of (DO36.25) by the same constant. Scaling dg scales the theta lift by that constant and scales the integral on the right twice, once through dg and once through the source inner product. This checks (DO36.25) for the measures used in R017 without assigning the target quotient mass the value one.

All interchanges are legitimate. The source and target quotients are compact, so their theta integrals are uniformly convergent on compact representatives. The remaining adelic matrix-coefficient integral is absolutely convergent: at each bad finite place Section 2 bounds it by \nolinkurl{C_(a,b)D_a(b/2)^2}; at almost every split place its absolute integral is exactly that same expression for the maximal vacuum, which is \nolinkurl{1+O_a(q^(-(b/2-a+1)))}. At nonsplit places the hyperspherical vacuum coefficient is nonnegative. Indeed compact invariance reduces it to a Cartan diagonal in the Levi of an integral hyperbolic polarization; the trivial splitting character gives the positive dilation action on the lattice indicator. The unitary change from that polarization preserves the selfdual lattice vacuum. Apply the same fixed-vacuum formula (DO36.21) to the trivial spherical source representation: it computes its absolute integral, and the parameter bound gives \nolinkurl{1+O_a(q^(-(b-2a+2)))}. This is a summable majorant. Infinity is compact. This is the ordinary convergent Siegel--Weil/unfolding regime, stronger than the continuation required by the second-term theorem.

\subsection{7. Strict positivity for the fixed global input, including multiplicities}\label{strict-positivity-for-the-fixed-global-input-including-multiplicities}

For a local factor define the bounded convolution operator

\[
 Q_v=\int_{G(F_v)}\langle\phi_v,\omega(g)\phi_v\rangle
                                               R(g)\,dg.
 \tag{DO36.26}
\]

At every split finite place (DO36.7) proves positivity of this operator on any unitary representation, including the actual automorphic Hilbert space. At G it preserves the inverse-generic subspace; at T it preserves the \nolinkurl{J_T}-fixed subspace; at Q and all spherical places it preserves the maximal-compact fixed subspace. These preservation assertions follow from the corresponding left/right compact covariance of its coefficient. At G/T/Q its restriction has the strictly positive uniform lower bounds (DO36.17), (DO36.14), (DO36.19). The nonsplit spherical operators are positive and have (DO36.22) on their fixed spaces. One may obtain the latter statement by decomposing the unitary spherical representation; the bound is uniform in its bounded spherical parameter.

At infinity compact Schur orthogonality gives, on the specified source contragredient line,

\[
                         Q_\infty=\frac{N_\infty^2}{d_\rho}.
 \tag{DO36.27}
\]

The actual F lies in that line in every real representation copy, by its matrix-coefficient construction in R017. It lies in the required fixed/type subspace at every finite place. Operators at different places commute. Finite products of these positive operators are bounded below on this subspace by the products of their displayed lower bounds. At the unramified places the deviations from one are uniformly summable, so the products converge on this subspace and their lower-bound product is still positive. Equivalently, decompose the finite-level, fixed-real-type automorphic space into its finitely many irreducible constituents, retaining their Hilbert multiplicity spaces; the same uniform bounds apply to every copy. The convolution operators act trivially on those multiplicity spaces. No multiplicity-one assumption and no choice of a favored irreducible projection is needed.

Let \nolinkurl{ell_v} denote (DO36.20) or (DO36.22) at the remaining spherical places and let \nolinkurl{c_G,c_T,c_Q} be the lower-bound constants in (DO36.17),(DO36.14),(DO36.19). Then

\[
 \begin{split}
 \|\Theta_\phi^{\rm bil}F\|_{dh}^{2}
 &\ge V_H\frac{N_\infty^2}{d_\rho}
       \prod_{P\in G}c_{G,P}
       \prod_{P\in T}c_{T,P}
       \prod_{P\in Q}c_{Q,P}
       \prod_{v\notin G\cup T\cup Q}\ell_v\ \|F\|^2\\
 &>0.
 \end{split}
 \tag{DO36.28}
\]

This proves that the exact nonzero R017 input survives the complete fixed first raw kernel at maximal29 level. The target value at the identity need not be nonzero; the proved assertion is positive global norm and hence a nonzero actual target automorphic form. It is stronger than a nonzero unrestricted local theta lift and addresses the particular finite input which the raw first coupled factor prescribes.

\subsection{8. Costs, source scope and range audit}\label{costs-source-scope-and-range-audit}

The finite raw Schwartz norm is exactly DO035(26). The real norm \nolinkurl{N_infty^2} is the actual DO034 Fischer/Gaussian/hook norm, including the pi\_B coordinate scalar and, on K032, its weighted quadratic metric. All finite projector and tail constants in (DO36.28) are explicit functions of the fixed ranks and residue cardinalities. They can be very small; this proof does not label them negligible. There is no auxiliary29 principal-level cost and no29 translation depth in this theorem. The T characters and their finite coefficient field are exactly those retained by R017.

The norm formula retains \nolinkurl{V_H}, the actual target quotient mass. Its positivity is enough for the present analytic nonvanishing. It is not an identification with the source's terminal mass or a solution of the MainX scalar equation. The construction is a genuine mixed-theta first image with the actual raw local G/T/Q factors and source compact types; the unchanged \nolinkurl{Fields.Kq}, the integral row comparisons and all later restricted-image contents remain separate.

Finally, the range distinction in Section 6 does not correct an earlier checked theorem. LIFT12 uses a=2,b=4, the allowed upper endpoint of GQT11.4. LIFT14(E23) explicitly applies that theorem only to later j\textgreater=1, where a=N\_j\textgreater d and \nolinkurl{a<b=a+d<2a}; its initial edge was constructed separately. LIFT16 continues those later edges, and LIFT19 has the same later-edge hypotheses after DO033's initial construction. The new first edge is the case requiring the separate GQT7.1 convergent-range derivation above.

\subsection{9. Primary references and proof checks}\label{primary-references-and-proof-checks}

\begin{itemize}
\tightlist
\item
  Gan, W. T., Qiu, Y., \& Takeda, S. (2014). The regularized Siegel--Weil formula (the second term identity) and the Rallis inner product formula. \emph{Inventiones Mathematicae, 198}, 739--831. arXiv:1207.4709v3. Theorem7.1, printed p.33; doubling section11.3, p.52; sections11.5--11.6, pp.53--54; the actual range restriction of Theorem11.4, p.54. \url{https://arxiv.org/pdf/1207.4709}. These portions were read directly in the primary PDF/text; its omitted conjugation bars use the already checked rendered-PDF convention in DO028.
\item
  Macdonald's bounded spherical criterion and Liu's actual unramified vacuum test are used in the exact checked form LIFT14(E18) and LIFT16(S19--20), with their primary locators. No new local theta classification is being asserted.
\item
  Literal raw factors, characters and target class: DO035 Sections1--4 and its source locators; LIFT20's raw first G equality; R017's exact source class, central compatibility and mass; DO034's actual new-host polynomial and metric.
\end{itemize}

The author whole-proof review and independent local review are recorded separately after the complete text is read. No Lean build is claimed.

\clearpage

\section{ALP-LIFT-24: the fixed second T kernel preserves the actual first projective-line type}\label{alp-lift-24-the-fixed-second-t-kernel-preserves-the-actual-first-projective-line-type}

Author: GPT-6 Astra Pro, ALP\_MACHINE /root/lift\_accounting. Date: 22 November 2026 UTC. Inherited reasoning effort is not exposed.

This finding replaces the first-edge determinant-tail estimate by an exact local period calculation suitable for the second edge. The raw T incidence function has a nonzero image on the actual first theta constituent. Its scalar in a specified induced model is a standard local L-factor divided by q\^{}n−1. The proof retains the nontrivial T dilation character of R017, uses the entire Schwartz space, and does not infer quotient survival from a formally nonzero full-rank integral.

The local normalization below is explicit. Its comparison with globally pinned arithmetic rows remains separate. In particular the L-factor can have nontrivial valuation at11; it is not discarded as an integral unit.

\subsection{1. Local groups and the actual T type}\label{local-groups-and-the-actual-t-type}

Let F be a characteristic-zero nonarchimedean local field with ring O, uniformizer pi, and residue field kappa of cardinality q. Assume its residue characteristic is larger than all ranks used below; this is satisfied by either explicit rational T prime of R017. Normalize vol(GL\_n(O))=1 in every rank. Write \[
 K_n=\mathrm{GL}_n(O),\quad
 J_n=\{g\in K_n:e_n^\vee\bar g=s(g)e_n^\vee,
                           s(g)\in\kappa^\times\},\quad
 H_n=\ker s.
 \tag{MT1}
\] Here e\_n denotes the last coordinate, with one-based indexing in this file. Thus J\_n is the reduction preimage of the parabolic of block sizes(n−1,1), and \[
 d_n=[K_n:J_n]=\frac{q^n-1}{q-1},\qquad
                  [K_n:H_n]=q^n-1.
 \tag{MT2}
\] Let nu be a nontrivial character of kappa\^{}times. For a smooth representation pi\_n define \[
 V_n(\nu)=\{v:\pi_n(j)v=\nu(s(j))v\quad(j\in J_n)\}.
 \tag{MT3}
\] For the R017 source, nu is the inverse of its named character nu\_P. All formulas use the actual representation character(MT3), so there is no implicit exchange of the two signs.

The literal T function for n→m is \[
 \phi_{u,v}(X)=\mathbf1_{M_{m,n}(O)}(X)
                              \mathbf1_{\{v\bar X=u\}},
 \tag{MT4}
\] where u and v are the prescribed nonzero residue covectors. This is FS646--649, phiT, as checked in DO035. The stage character eta is unramified at T. The actual Hom action is \[
 \omega_\eta(g,h)\phi(X)=
 \eta(\det g)^{-1}\eta(\det h)
 |\det g|^{m/2}|\det h|^{-n/2}\phi(h^{-1}Xg).
 \tag{MT5}
\] The source lift uses the bilinear coinvariant convention. These are the same signs as LIFT20--22 and DO036.

Integral source and target coordinate changes take u,v to the last coordinate covectors. They preserve the two Haar normalizations and the local Schwartz norm, and their eta values are1. We therefore use these coordinates locally. This does not assert an uncharged global change of arithmetic basis.

\subsection{2. Which representations contain this type}\label{which-representations-contain-this-type}

We use the following precise consequence of the principal-series type theorem and unitary classification.

\textbf{Local type lemma.} If an irreducible smooth unitary representation pi\_n has \(V_n(\nu)\ne0\), then \[
          \pi_n=\sigma\times\xi,\qquad
 \dim V_n(\nu)=1,
 \tag{MT6}
\] where sigma is an irreducible unitary spherical representation of GL\_(n−1)(F), xi is a unitary character of F\^{}times with residue nu, and the induction is normalized. Conversely this product has the displayed one-dimensional type.

Here and subsequently rank zero has its usual trivial interpretation. All actual source ranks are larger than one.

Proof. Restrict(MT3) to the upper Iwahori. Its torus character is \((1,\ldots,1,\nu)\), and it is trivial on the two relevant unipotent subgroups. Roche(1998), Theorem7.7, identifies this tame Iwahori character as the type for the principal-series inertial component. Thus pi\_n is a subquotient of a principal series with n−1 unramified characters and exactly one character with residue nu. The explicit residue-characteristic hypothesis in Roche, p.379, is satisfied here.

Separate the two inertial lines. Zelevinsky's classification and its disjoint-support irreducibility, already used in LIFT19, give \(\pi_n=\sigma\times\xi\), with sigma in the unramified principal series component and xi the unique character on the nu line. We explain why sigma is spherical and why the type has dimension1; being in that component alone would not suffice.

On the K\_n principal-congruence invariants of the full principal series, the compact representation is the finite principal series induced from the characters(1,\ldots,1,nu). Its J\_n,nu eigenspace has dimension one. Indeed the Bruhat double cosets of a Borel and P\_(n−1,1) are represented by the n possible positions of the distinguished last coordinate. If that coordinate does not match the nu coordinate, the diagonal torus in the intersection has independent unit entries at those two positions. The two characters then disagree, so the corresponding intertwiner is zero. In the matching coset they agree, and the space of scalar intertwiners is one-dimensional. This is the ordinary finite Mackey calculation; all compact invariants are exact in characteristic zero.

The unramified GL\_(n−1) principal series has its unique spherical irreducible subquotient sigma\_sph. The product sigma\_sph×xi is irreducible, by the distinct inertial supports, and has a nonzero J\_n,nu vector. Because the entire principal series has only the one such line, it is its only irreducible subquotient containing that line, with multiplicity one. Hence sigma=sigma\_sph and(MT6) follows.

Finally Tadic(1986), TheoremD, says every unitary irreducible is a product of the stated unitary Speh and complementary blocks, with unique factors, and such products are irreducible. Each block is contained in a single inertial line. Since the nu line in pi\_n has multiplicity one, its block is the unitary character xi. Removing that block leaves a unitary product, which is sigma. This proves the unitary assertion as well as the stated irreducibility.

A convenient realization of the line is the normalized induced model on P\_(n−1,1). Choose a unit spherical vector v\_sigma and let v\_0 have compact restriction \[
 v_0(k)=
 \begin{cases}
       \nu(s(k))v_\sigma,&k\in J_n,\\
       0,&k\notin J_n.
 \end{cases}
 \tag{MT7}
\] Then \(v_0(1)=v_\sigma\) and \[
                         \|v_0\|^2=d_n^{-1}.
 \tag{MT8}
\] The defining transformation under P∩K and the right J\_n character both follow directly. The lower principal unipotent congruence subgroup is in J\_n and acts trivially.

\subsection{3. The whole-Schwartz period at the actual point}\label{the-whole-schwartz-period-at-the-actual-point}

For this section pi\_n may be any irreducible smooth representation, and eta any unitary character. Put \[
 m=n+d,\qquad z_0=\frac{d+1}{2},\qquad
 X_0=\begin{pmatrix}I_n\\0\end{pmatrix}.
 \tag{MT9}
\] For Re(z) sufficiently large define \[
 I_z(\phi,w)=\int_{\mathrm{GL}_n(F)}
    \phi(X_0g)|\det g|^{z+(n-1)/2}
                      \eta(\det g)^{-1}\pi_n(g)w\,dg.
 \tag{MT10}
\] For unitary pi\_n the ordinary vector integral converges in a right half-plane, since phi is compactly supported and the determinant integral is bounded there. For general pi\_n, one can define its coordinates by smooth matrix coefficients in the usual GJ half-plane.

Suppose \[
                        L(z,\pi_n\otimes\eta^{-1})
                   \text{ is finite at }z=z_0.
 \tag{MT11}
\] Then(MT10) has a well-defined holomorphic vector value at z0 for every Schwartz phi on the entire matrix space, including functions that meet the rank-deficient locus.

Here are the details. The restriction g→phi(X0g) extends to a Schwartz function on M\_n(F). Each smooth dual matrix coefficient of(MT10) is a Godement--Jacquet integral. The published GJ theorem, in the all-irreducibles version recalled in Jiang and Luo(2023), Theorem2.4(1)--(2), pp.313--314, makes it \[
                    L(z,\pi_n\eta^{-1})\,
                              P(q^z,q^{-z})
 \tag{MT12}
\] for a Laurent polynomial P. Its normalization of Haar can change the constant polynomial, but cannot change this holomorphy claim.

The restriction of phi is invariant under left multiplication by a sufficiently small compact open subgroup J of GL\_n.~Take J also in the kernel of eta. Changing variables shows that I\_z is J-fixed wherever the vector integral converges. Admissibility makes \(V_{\pi_n}^J\) finite-dimensional. Choose smooth dual coordinates on that space; (MT12) continues each coordinate and proves holomorphy at z0. This constructs an actual smooth vector, not just an element of the algebraic bidual. It is independent of the choice of J and the coordinates by the identity theorem.

Define \[
              T_z(\phi,w)(h)=I_z(\omega_\eta(1,h)\phi,w).
 \tag{MT13}
\] The exact source covariance is \[
 I_z(\omega_\eta(g_0,1)\phi,\pi_n(g_0)w)
           =|\det g_0|^{z_0-z}I_z(\phi,w).
 \tag{MT14}
\] For \(p=\begin{pmatrix}A&B\\0&D\end{pmatrix}\in P_{n,d}(F)\), the exact target covariance is \[
 I_z(\omega_\eta(1,p)\phi,w)=
   \eta(\det D)|\det A|^{z-1/2}|\det D|^{-n/2}
                       \pi_n(A)I_z(\phi,w).
 \tag{MT15}
\] Both follow by changing g to gg0 or Ag in the convergent region; all eta factors and determinant exponents are displayed in(MT5). They are identities of rational functions in q\^{}z, so remain true at every common holomorphic value.

At z=z0, (MT14) is the literal simultaneous source coinvariant identity, and (MT15) is exactly normalized parabolic induction. Thus \[
  T_{z_0}:\ 
   \bigl(\mathcal S(M_{m,n}(F))\otimes\pi_n\bigr)_{\mathrm{GL}_n(F)}
       \longrightarrow
     \operatorname{Ind}_{P_{n,d}}^{\mathrm{GL}_m}
                    (\pi_n\boxtimes\eta\circ\det_d)
 \tag{MT16}
\] is a genuine target-equivariant map. It is smooth because a compact open subgroup fixing phi on the target fixes its induced function.

If pi\_n and eta are unitary, the target in(MT16) is irreducible by Tadic's unitary induction theorem. Any nonzero map(MT16) is therefore the actual small theta quotient map, up to a single nonzero scalar, by Mínguez(2008), Theorem1. This last statement compares nonvanishing with the actual theta quotient. It does not identify that scalar with a globally pinned arithmetic one.

This construction is useful precisely when an ordinary full-rank Bochner integral no longer converges. It avoids assuming that a map defined only on the open rank stratum automatically extends across its boundary.

\subsection{4. Exact value of the untouched incidence function}\label{exact-value-of-the-untouched-incidence-function}

Now impose(MT6), and take eta unramified. Choose a primitive matrix \(X_*\in M_{m,n}(O)\) with \[
                   vX_*=u,\qquad u=e_n^\vee.
 \tag{MT17}
\] Its columns extend to an integral basis, so \(X_*=k_*X_0\) with \(k_*\in K_m\). The corresponding value of(MT13), at \(h=k_*^{-1}\), is the period obtained by replacing X0 with X*.

Put \[
 F_u(g)=\mathbf1_{M_n(O)}(g)\mathbf1_{\{u\bar g=u\}},\qquad
 F_\nu(g)=\sum_{c\in\kappa^\times}\nu(c)^{-1}
                 \mathbf1_{M_n(O)}(g)\mathbf1_{\{u\bar g=cu\}}.
 \tag{MT18}
\] The restriction of the raw function(MT4) to X*g is exactly F\_u. For w in the line(MT3), right multiplication by an integral diagonal matrix with last entry c shows \[
                I_z(F_\nu,w)=(q-1)I_z(F_u,w).
 \tag{MT19}
\] This equation is initially in the convergent half-plane and then continues. The unit determinant has absolute value1 and its eta value is1; the character in F\_nu cancels the character of w. Moreover F\_nu is bi-J\_n equivariant with character nu\^{}\{-1\}, so its period operator has image in the one-dimensional line(MT3).

We compute its scalar by evaluating the induced vector at1. Absorb the unramified eta into the factors in(MT6), and write \(\sigma'=\sigma\eta^{-1}\), \(\xi'=\xi\eta^{-1}\). The residue of xi' is still nu. Every matrix in the support of F\_nu has a unique block form \[
 g=\begin{pmatrix}A&B\\C&e\end{pmatrix},\quad
 A\in M_{n-1}(O),\ B\in O^{n-1},
 C\in\pi O^{1\times(n-1)},\ e\in O^\times.
\] It is invertible exactly when \(A'=A-Be^{-1}C\) is invertible, and \[
 g=
 \begin{pmatrix}I&Be^{-1}\\0&1\end{pmatrix}
 \begin{pmatrix}A'&0\\0&e\end{pmatrix}
 \begin{pmatrix}I&0\\e^{-1}C&1\end{pmatrix}.
 \tag{MT20}
\] The last factor is in the principal lower congruence subgroup. In the induced model(MT7), \[
 v_0(g)=|\det A'|^{1/2}\sigma'(A')\nu(\bar e)v_\sigma.
 \tag{MT21}
\] The nu factor is canceled by F\_nu. Write \(c_j=\prod_{i=1}^j(1-q^{-i})\). Multiplicative Haar on GL\_j is \(c_j^{-1}|\det g|^{-j}\) times additive matrix Haar. The change \(A\mapsto A'\) has additive Jacobian1. Integrating B,C,e contributes \(q^{-(n-1)}(1-q^{-1})\). Converting the remaining additive measure in A' to normalized GL\_(n−1) Haar therefore contributes the exact constant \[
 \frac{c_{n-1}}{c_n}\,q^{-(n-1)}(1-q^{-1})
                       =\frac{q-1}{q^n-1}=d_n^{-1}.
 \tag{MT22}
\] The remaining determinant exponent is \[
        z+(n-1)/2+1/2-n+(n-1)=z+(n-2)/2.
\] This is exactly the rank n−1 standard GJ exponent. The spherical calculation LIFT16(S8), with vol(K\_(n−1))=1, now gives \[
                  I_z(F_\nu,v_0)=d_n^{-1}
                                   L(z,\sigma')v_0.
 \tag{MT23}
\] One can evaluate against the spherical functional to use(S8): the integral operator is bi-K\_(n−1) and hence acts on its one-dimensional spherical line. Thus no genericity or Whittaker model is required, even if sigma is a non-generic spherical quotient.

Combining(MT19) and(MT23), at every holomorphic value z0, \[
 \boxed{
   T_{z_0}(\phi_{u,v},v_0)(k_*^{-1})
      =\frac{L((m-n+1)/2,\sigma\eta^{-1})}{q^n-1}\,v_0.}
 \tag{MT24}
\] If v0 is replaced by its unit normalization, the same scalar applies. Formula(MT8) still records the norm of the displayed identity-normalized compact section.

For nonzero L and(MT11), this proves that the raw incidence function survives the actual small theta quotient on every nonzero vector of the entire line(MT3), with no change of finite Schwartz input and no supplementary compact projector.

\subsection{5. Application to the actual first image}\label{application-to-the-actual-first-image}

Use the source ranks \[
       a=k,\quad d=t+1>a,\quad n=a+d,\quad m=a+2d.
 \tag{MT25}
\] For K031 these are(42,1760,1802,3562); for K032 they are (168,1634,1802,3436). Let pi\_0 be any actual local constituent of R017's initial function. Its nontrivial T projective-line character gives \[
                         \pi_0=\sigma_0\times\xi,
 \tag{MT26}
\] where sigma0 is unitary spherical of rank a−1 and xi is unitary with the prescribed nontrivial residue. Let eta0,eta1 be the actual unramified unitary splitting characters at this T place.

The actual first local theta constituent is \[
 \pi_1=\pi_0\times(\eta_0\circ\det_d)
       =\bigl(\sigma_0\times\eta_0\circ\det_d\bigr)\times\xi.
 \tag{MT27}
\] Mínguez's parameter theorem gives this identification, and Tadic makes the displayed unitary induction irreducible. The rank n−1 factor \(\sigma_1=\sigma_0\times\eta_0\circ\det_d\) is spherical: normalized induction of its spherical vectors has a nonzero K-fixed vector, and the induction is irreducible. It is also unitary.

The fixed raw first incidence function transports its projective-line character to the actual first output. Explicitly, if \(j\in J_n\) has \(v\bar j=c v\), its target action on(MT4) is the incidence function with source covector cu. For a lift \(g_c\in J_a\) with \(u\bar g_c=cu\), that function is \(\omega(g_c^{-1},1)\phi_{u,v}\). The source coinvariant identity then gives \[
 [\omega(1,j)\phi_{u,v},w]
       =[\phi_{u,v},\pi_0(g_c)w]
       =\nu(c)[\phi_{u,v},w].
 \tag{MT28}
\] Thus every actual first output lies in \(V_{\pi_1}(\nu)\). No change of the first output by a Hecke operator is used. DO036 proves nonvanishing of the complete global first raw kernel; locally its T component is nonzero by the same first-edge test or by(MT24) at ranks a,n.~The line in(MT28) is consequently the one to which the second-edge theorem applies.

It remains to verify all holomorphy and nonzero conditions for the second edge. The unramified Satake parameters of sigma1 are the parameters alpha\_i of sigma0 and the d parameters \[
              \eta_0(\pi)\,q^{(d-1)/2-j},
                                \qquad j=0,\ldots,d-1.
 \tag{MT29}
\] Since sigma0 is unitary spherical of rank a−1, the checked bounded spherical-parameter criterion gives \(|\alpha_i|\le q^{(a-2)/2}\). At \(z_0=(d+1)/2\), with eta1 unitary, \[
 |\alpha_i\eta_1(\pi)^{-1}q^{-z_0}|
       \le q^{-(d-a+3)/2}<q^{-1},\qquad
 |\eta_0\eta_1^{-1}(\pi)q^{(d-1)/2-j-z_0}|
       =q^{-1-j}.
 \tag{MT30}
\] The ramified xi eta1\^{}\{-1\} factor has local L-function1. Therefore \(L(z,\pi_1\eta_1^{-1})=L(z,\sigma_1\eta_1^{-1})\) is finite and nonzero at z0. This proves(MT11) for all Schwartz vectors on the second edge, not only for the test being evaluated.

Consequently, for every nonzero actual first vector w in its prescribed T line, the unchanged raw second T factor has \[
 \boxed{
   T^{(1)}(\phi_T,w)(k_*^{-1})=
       \frac{L((d+1)/2,\sigma_1\eta_1^{-1})}{q^{a+d}-1}\,w
                                  \ne0.}
 \tag{MT31}
\] This is an actual preceding-image statement in the local small theta quotient. It does not merely show that the two compact types occur abstractly somewhere in a representation.

\subsection{6. Explicit size bounds and remaining normalization gates}\label{explicit-size-bounds-and-remaining-normalization-gates}

The scalar in(MT31), denoted lambda\_T, satisfies \[
 \frac{(1+q^{-1})^{-(n-1)}}{q^n-1}
       \le|\lambda_T|
       \le\frac{(1-q^{-1})^{-(n-1)}}{q^n-1}.
 \tag{MT32}
\] This follows by applying \(1-|z|\le|1-z|\le1+|z|\) to the n−1 Euler factors in(MT30). It is a uniform estimate over all initial constituents of the fixed source level and type. The large T prime does not make the basic Haar denominator \(q^n-1\) negligible.

If the parameters and the chosen local normalization are algebraic, the exact11-adic scalar valuation is \[
 v_{11}(\lambda_T)=
 v_{11}\!\left(L((d+1)/2,\sigma_1\eta_1^{-1})\right)
                              -v_{11}(q^n-1).
 \tag{MT33}
\] The first term is retained. No claim that the L-value is an11-unit is implied by its complex nonvanishing or its unitary normalization.

For the exact T primes of R017 the finite replay gives v11(q\^{}a−1)=v11(q\^{}n−1)=0 on K031, and both valuations equal1 on K032. These are only the displayed Haar denominators; the separate L-factor term in(MT33) is still retained.

The theorem settles the fixed local T test on the actual first line. Its normalization is the GJ/normalized-induction convention (MT10), Haar vol(K)=1, and compact section(MT7). Comparison with a globally pinned theta map can multiply it by a nonzero scalar. The global coupled Q and G tests have their separately recorded normalizations. Their assembly into the prescribed arithmetic row lattice, the target mass, all six tied source data, and DP44 costs require those comparisons. The analytic GJ theorem is not an integral-saturation theorem.

The scalar(MT24) is not justified by citing the generic newform test alone: sigma1 is usually non-generic because of its one-dimensional GL\_d block. The direct Schur-complement calculation and the all-irreducibles GJ theorem are essential here.

\subsection{7. Primary sources and check scope}\label{primary-sources-and-check-scope}

\begin{itemize}
\tightlist
\item
  Roche, A. (1998). Types and Hecke algebras for principal series representations of split reductive p-adic groups. Annales scientifiques de l'École Normale Supérieure,31(3),361--413. Residual-characteristic restriction p.379; tame type and Theorem7.7/Remark7.8, pp.391--392; Theorem9.4, p.404. \url{https://www.numdam.org/article/ASENS_1998_4_31_3_361_0.pdf} The theorem is used as a published input; the finite one-line calculation in section2 explains the stronger spherical conclusion.
\item
  Tadic, M. (1986). Classification of unitary representations in irreducible representations of general linear group (non-Archimedean case). Annales scientifiques de l'École Normale Supérieure,19(3),335--382. TheoremD and its block definitions, pp.338--339. The characteristic-zero appendix removes the historical dependency discussed immediately after that theorem. \url{https://www.numdam.org/article/ASENS_1986_4_19_3_335_0.pdf}
\item
  Jiang, D., \& Luo, Z. (2023). Certain Fourier operators and their associated Poisson summation formulae on GL1. Pacific Journal of Mathematics,326(2),301--372. Section2B, equation(2-10), Theorem2.4(1)--(2), pp.313--314. \url{https://msp.org/pjm/2023/326-2/pjm-v326-n2-p05-p.pdf} This is the directly read all-irreducibles restatement of Godement--Jacquet, ChapterI, Theorem3.3; the original monograph is not claimed to have been read in full. Haar constants are computed independently in(MT22) and LIFT16(S8).
\item
  Mínguez, A. (2008). Correspondance de Howe explicite: paires duales de type II. Annales scientifiques de l'École Normale Supérieure,41(5),717--741. Theorem1, printed p.718; actual oscillator normalization(1.4), pp.721--722. \url{https://www.numdam.org/item/ASENS_2008_4_41_5_717_0.pdf}
\item
  LIFT16(S8), LIFT14(E18) for the independently checked spherical Haar normalization and bounded Satake criterion; DO035--036 and R017 for the literal incidence function, actual first input, and its nontrivial T character; the source FS646--649 and actual character definition in LevChiEqOne.lean61,100.
\end{itemize}

The finite diagnostic \nolinkurl{raw_t_projective_period_certificate.py} checks the exact compact character average, block Haar constant and two-edge Satake inequalities. Such computations check the normalization and finite algebra; the GJ continuation, type theorem, and local theta identification are the proofs above and their published inputs. Whole-proof and independent reviews are recorded separately. No Lean build is claimed.

\clearpage

\section{AM-K035 --- an explicit mixed CM period and source Fock-frame adapter}\label{am-k035-an-explicit-mixed-cm-period-and-source-fock-frame-adapter}

Author: GPT-6 Astra Pro, ALP\_MACHINE /root/kummer\_extension, 22 November 2026 UTC; inherited reasoning effort is not separately exposed.

This finding supplies an actual equation downstream of the failed quadratic-field constructor. It identifies the period vector of the principally polarized hosts \(J^{14}\) and \(J^{56}\), computes their source-style embedding-lattice coordinate generators, and calculates the exact change from the source's unit-normalized creation operators to algebraic reduced-theta derivatives. The resulting all-polynomial equation has the literal shape of ClsK.Core, including its per-place RealD.sc. It does not inhabit the old Fields.Kq or elliptic VacIn records, and it does not identify an independently specified later arithmetic row map by definition.

Carried inputs are the accepted DP57 and canonical-gauge existence theorems, AM-R-004 and AM-R-006, ALP-LIFT-09, AM-K018, AM-K030/032--034, and AM-DO-017/018/022/026/027. Their proved conclusions are used at their recorded scope. The old common-quadratic obstruction is settled by AM-R-003 and is not proved again here.

\subsection{1. The exact source equations being replaced}\label{the-exact-source-equations-being-replaced}

The uniform analytic restriction occurs before the generic class-kernel equations:

\begin{itemize}
\tightlist
\item
  \nolinkurl{Development/AbelianVarietyTheory/Definitions/LatticeFrames/SplitPart.lean,} lines22--24 and38--49, uses one complex period \(\delta\), one real metric scalar sc, and the hypotheses \[
  \tau(\xi_0)=i/\sqrt d,\qquad \mathrm{sc}\,|\delta|^2=2/\sqrt d
  \quad\text{for every }\tau.
  \tag{PF1}
  \]
\item
  \nolinkurl{Development/AbelianVarietyTheory/Definitions/AutomorphicMembership/IsDiscriminantTwelfthRootOfIdealLatticeIffDelta1.lean,} lines53--67, defines the algebraic embedding lattice mmM, its coordinate-line intersection ellT, and IsDgen. This coordinate generator is logically distinct from the elliptic period in its subsequent VacIn.
\item
  \nolinkurl{Formalism/LocalToGlobal/IndefiniteReal/LocalNewformFunction.lean,} lines124--141, defines crea and the actual normalizing factor \[
  \mathrm{fockScale}(\tau)=2\sqrt{2\pi|\tau(\xi_0)|},
  \qquad a_\tau^+=\mathrm{fockScale}(\tau)^{-1}\mathrm{crea}_\tau.
  \tag{PF2}
  \] DatumVec.lean, lines130--134, uses these operators in pureTensorO. They are not the unnormalized operators \(r_\tau\) of ALP-LIFT-09.
\item
  \nolinkurl{Development/AbelianVarietyTheory/Definitions/LatticeFrames/OnFramesRdApply.lean,} lines25--39 and108--121, already allows a nonzero complex scalar RealD.sc(j,τ) separately at every place. Its Core equation is \[
  v_j\,\vartheta_{\mathrm{pureTensorO}(P)}(g_c,h_{c'})
   =c_{K,j}\,\mathrm{abelian}_2(c,c')
      \bigl(\mathrm{mvZ}(1,\operatorname{diag}(\mathrm{sc}_{j,\tau}))P\bigr).
  \tag{PF3}
  \] The following Kval clause is a separate equality for the prescribed arithmetic map \(t.K_j\).
\end{itemize}

Thus a common scalar as in (PF1) is not required by the displayed downstream equation (PF3). An exact varying period and metric profile can be tested there without asserting a new instance of the old constructor.

\subsection{2. Actual periods and metric on the two product hosts}\label{actual-periods-and-metric-on-the-two-product-hosts}

Write \(E=\mathbb Q(\zeta_7)\), \(E_0=E(i)\), and retain the degree-seven real field \(B\subset\mathbb Q(\zeta_{29})\), positive generator \(\pi_B\), and trace-orthonormal basis \(e_0,\ldots,e_6\) from AM-K030. There are two cases: \[
\begin{array}{c|c|c|c}
h&R&K_h&A_h\\ \hline
1&\mathbb Q&E_0B&J^{14}\\
4&\mathbb Q(\sqrt5,\sqrt{13})&E_0BR&J^{56}.
\end{array}
\tag{PF4}
\] The full CM type is indexed by \[
\tau=(i,\varepsilon,l,\rho),\quad
i\in\{1,2,3\},\quad \varepsilon\in\{+1,-1\},\quad
l\in\mathbb Z/7\mathbb Z,
\tag{PF5}
\] where \(\rho\) runs over the real embeddings of \(R\), omitted if \(h=1\). The restrictions are \(\zeta_7\mapsto\zeta_7^i\), \(i_{\mathbb C}\mapsto\varepsilon i_{\mathbb C}\), and the l-th B-embedding. Put \[
\begin{split}
\pi_l&=\sigma^l(\pi_B),\\
\delta_\rho&=1\quad(h=1),\qquad
\delta_\rho=\rho(\delta_1\delta_3)\quad(h=4),\\
\xi_h&=\frac{\zeta_7-\zeta_7^{-1}}{14}
                 \pi_B^{-6}\delta_1\delta_3
       \quad\text{with the last factor omitted if }h=1.
\end{split}
\tag{PF6}
\] All \(\pi_l,\delta_\rho\) are positive. Define the actual periods \[
\Omega_i=-\frac{2i_{\mathbb C}}7
          B(i/7,1/2)(\zeta_7^i-1),\qquad
\Omega_\tau=\Omega_i,
\tag{PF7}
\] and retain the actual Hodge Gram coefficients of ALP-LIFT-03 and AM-R-004: \[
G_i=\frac{2\pi}{7}
 \frac{\Gamma(i/7)\Gamma(1/2-i/7)}
      {\Gamma(1-i/7)\Gamma(1/2+i/7)}>0.
\tag{PF8}
\]

\subsubsection{Theorem PF9: an actual period-vector replacement}\label{theorem-pf9-an-actual-period-vector-replacement}

There are explicit algebraic invariant coordinates \(\zeta_\tau\) on the universal cover of \(A_h\) in which \[
\boxed{\Lambda_h=\operatorname{diag}(\Omega_\tau)
                       \Psi(\mathcal O_{K_h}),\qquad
H_h(\zeta,\zeta)=\sum_\tau h_\tau|\zeta_\tau|^2,}
\tag{PF9}
\] where \[
h_\tau=\frac{\delta_\rho\pi_l^{-6}}{2G_i},\qquad
q_\tau=-2i_{\mathbb C}\tau(\xi_h)>0,\qquad
\boxed{h_\tau|\Omega_\tau|^2=q_\tau.}
\tag{PF10}
\] The same formulas apply entry by entry on the Hom torus. Its period lattice for actual moved modules \(N,N'\) is \(\operatorname{diag}(\Omega_\tau)\Psi(\operatorname{Hom}_{\mathcal O_{K_h}}(N,N'))\).

\textbf{Proof.} On a pair of ordinary Jacobian factors use \[
w_{i,\varepsilon,j}=u_{i,j}+\varepsilon i_{\mathbb C}v_{i,j}.
\] Its period lattice is \(\Omega_i\Psi(\mathcal O_{E_0})\), directly from AM-R-004 and the integral identification \(\mathcal O_{E_0}=\mathcal O_E[i]\). For the seven B-factors put \[
z_{i,\varepsilon,l}=\sum_{j=0}^6\sigma^l(e_j)w_{i,\varepsilon,j},
\qquad
\zeta_{i,\varepsilon,l}=\pi_l^3z_{i,\varepsilon,l}.
\tag{PF11}
\] The first transformation identifies the lattice with \(\Omega_i\Psi(\pi_B^{-3}\mathcal O_{K_1})\), because the \(e_j\) form an actual integral basis of that fractional ideal. The second gives the lattice in (PF9). Orthogonality of the first matrix and the displayed rescaling give the metric when \(h=1\).

For \(h=4\), apply the embedding matrices of the integral bases \((1,u_1)\), \((1,u_3)\) in AM-K032. Their exact identities \[
V_a^{\mathsf T}\operatorname{diag}(\delta_{a,+},\delta_{a,-})V_a=I_2
\] give the factor \(\delta_\rho\), and their integral tensor product gives \(\mathcal O_{K_4}\). No isogeny or unspecified period factor occurs.

For the final normalization put \(x=i/7\). The beta-gamma and reflection formulas give \[
\begin{split}
\frac{|\Omega_i|^2}{G_i}
&=\frac87
\frac{\Gamma(x)\Gamma(1-x)}
     {\Gamma(1/2+x)\Gamma(1/2-x)}
\sin^2(\pi x)\\
&=\frac87\sin(\pi x)\cos(\pi x)
=\frac47\sin(2\pi i/7).
\end{split}
\tag{PF12}
\] Since \(\tau(\xi_h)=i_{\mathbb C}\sin(2\pi i/7)\delta_\rho\pi_l^{-6}/7\), (PF10) follows. The Hom assertion is the same scalar embedding entry by entry. The alternating-form convention is \[
\operatorname{Im}H_h(X,Y)
 =-\operatorname{Tr}_{K_h/\mathbb Q}
       \bigl(\xi_h\operatorname{tr}(Y^*X)\bigr),
\tag{PF13}
\] exactly as in AM-DO-022. ∎

\subsection{3. Exact coordinate generators for the embedding lattice}\label{exact-coordinate-generators-for-the-embedding-lattice}

Let \(M=K_h\), or any number field containing it, and \(R_M=\mathcal O_M\). The fields \(K_h\) are Galois, so all selected embeddings land in M. Define \[
\mathcal M_\Psi
   =\operatorname{span}_{R_M}
          \{(\tau(a))_{\tau\in\Psi}:a\in\mathcal O_{K_h}\}
       \subset M^{42h},
\qquad
\ell_\tau=\{c\in M:ce_\tau\in\mathcal M_\Psi\}.
\tag{PF14}
\] This is the algebraic shape of mmM and ellT for the actual standard ideal \(\mathcal O_{K_h}\); it contains no complex period. Put \[
d_i=\prod_{\substack{j\in\{1,2,3\}\\j\ne i}}
                       (\zeta_7^i-\zeta_7^j),
\] \[
\boxed{cD_\tau
   =d_i(2\varepsilon i_{\mathbb C})\pi_l^6
      \begin{cases}
      1,&h=1,\\
      \rho(\sqrt5)\rho(\sqrt{13}),&h=4.
      \end{cases}}
\tag{PF15}
\]

\subsubsection{Theorem PF16: the exact source-style generator}\label{theorem-pf16-the-exact-source-style-generator}

\[
\boxed{\ell_\tau=cD_\tau R_M.}
\tag{PF16}
\] For normalized extensions of rational valuations, \[
\begin{array}{c|ccccc}
p&2&7&29&5,13&\text{all others}\\ \hline
v_p(cD_\tau),\ h=1&1&1/3&6/7&0&0\\
v_p(cD_\tau),\ h=4&1&1/3&6/7&1/2&0.
\end{array}
\tag{PF17}
\] In particular all \(cD_\tau\) are units at11.

\textbf{Partial E-factor.} Let \(r_i=\zeta_7^i\), \(i=1,2,3\). The image after \(R_M\)-span is polynomial evaluation at these three roots. Division by the monic \(F(T)=\prod_{j=1}^3(T-r_j)\) reduces it to degree at most two. If this polynomial vanishes at the two roots other than \(r_i\), successive monic division shows that it is a scalar over \(R_M\) times \(\prod_{j\ne i}(T-r_j)\). Thus its possible values at \(r_i\), with the others zero, form precisely \(F'(r_i)R_M=d_iR_M\). This calculation uses the selected three embeddings; the full E-different would give the wrong ideal.

\textbf{Fully selected factors.} For a number field L with all embeddings in M, the dual of its full embedding lattice for the coordinate dot product is the embedding lattice of \(\mathfrak D_{L/\mathbb Q}^{-1}\), after \(R_M\)-span. This follows by writing a basis and its trace-dual basis. If \(\mathfrak D_L=(\Delta_L)\), then \[
ce_\sigma\in\operatorname{span}_{R_M}\operatorname{Emb}(\mathcal O_L)
\ \Longleftrightarrow\
c\,\sigma(\Delta_L^{-1}\mathcal O_L)\subset R_M
\ \Longleftrightarrow\
c\in\sigma(\Delta_L)R_M.
\tag{PF18}
\] The factors are therefore \(2\varepsilon i_{\mathbb C}\) for \(\mathbb Q(i)\), \(\pi_l^6\) for B, and \(\rho(\sqrt5),\rho(\sqrt{13})\) for the real quadratics. The B-different is the exact result of AM-K030; no integral power basis is assumed.

\textbf{Tensor-product step.} The coprime discriminants give the integral tensor decompositions established in AM-K030/032. Their embedding matrices are Kronecker products. For full-rank lattices \(L_1\subset M^m,L_2\subset M^n\) over a Dedekind domain, the coordinate-line ideal of \(L_1\otimes L_2\) at \((i,j)\) is the product of the two coordinate-line ideals. To prove this, localize at a prime and use basis matrices A,B over the DVR. Then \[
(A\otimes B)^{-1}(e_i\otimes e_j)
     =(A^{-1}e_i)\otimes(B^{-1}e_j).
\] The smallest entry valuation is the sum of the two minima, so the valuation threshold for multiplication by c to become integral is the sum of the two thresholds. Equality at every localization proves the ideal equality and hence (PF16).

Each nonzero difference of seventh roots has normalized 7-valuation \(1/6\) and no other finite support; \(\pi_B\) has 29-valuation \(1/7\) and no other support. The remaining different generators give the other entries in (PF17). ∎

This is an exact solution of the embedding-coordinate generator problem with \(M=K_h\). It does not assert that these are primitive Hodge bases at bad primes or that the reduced frame is integral. ALP-LIFT-08's section contents and K033's reduced-gauge denominators remain. At11 the change to jetC(cD,−) multiplies each fixed multidegree by a unit.

\subsection{4. The literal source creation factor, including its sign}\label{the-literal-source-creation-factor-including-its-sign}

The source convention in \nolinkurl{BesselHermitianTheta/TiedWeil.lean59–62} is \(\psi_\infty(t)=\exp(-2\pi it)\). In one CM coordinate \(w=x+iy\), \[
b_{\rm src}(w,v)=\operatorname{Tr}_{\mathbb C/\mathbb R}
       (\tau(\xi_h)w\bar v)
       =-q_\tau\operatorname{Im}(w\bar v).
\] Thus the literal heisO in \nolinkurl{LocalNewformFunction.lean29–35} is \[
T_v f(w)=e^{\pi iq_\tau\operatorname{Im}(w\bar v)}f(w+v).
\tag{PF19}
\] Its differentiated operators are \[
D_1=\partial_x+\pi iq_\tau y,\quad
D_i=\partial_y-\pi iq_\tau x,\quad
\mathrm{annih}=2\partial_{\bar w}+\pi q_\tau w,\quad
\mathrm{crea}=2\partial_w-\pi q_\tau\bar w.
\tag{PF20}
\] In particular a Fock-vacuum theta function is a Gaussian times a holomorphic function, with the positive Gaussian sign shown here.

For the complete Hom space define the mixed reading \[
(\mathcal R_\Omega f)(u)
 =\exp\!\left(\frac\pi2\sum_{\tau,r,s}h_\tau|u_{\tau,r,s}|^2\right)
       f\bigl((-\Omega_\tau^{-1}u_{\tau,r,s})_{\tau,r,s}\bigr).
\tag{PF21}
\] The finite component is held fixed. This is the actual varying-period counterpart of xiN.

\subsubsection{Theorem PF22: exact source polynomial normalization}\label{theorem-pf22-exact-source-polynomial-normalization}

Let \(F=\mathcal R_\Omega f\) be holomorphic and P any complex polynomial in the Hom creation labels. With identity source/target orthonormal frames, use the literal pureTensorO of (PF2). Then \[
\boxed{[\mathrm{pureTensorO}(f,P)](0)
 =\left[P(\lambda_\tau\partial_{\tau,r,s})F\right](0),\qquad
\lambda_\tau=-\frac{\Omega_\tau}{\sqrt{\pi q_\tau}}.}
\tag{PF22}
\] The square root is positive real. Every \(\lambda_\tau\ne0\), and \[
\boxed{|\lambda_\tau|^2\,\pi h_\tau=1.}
\tag{PF23}
\]

\textbf{Proof.} Conjugating crea by the Gaussian gives \[
e^{\pi q|w|^2/2}\,\mathrm{crea}\,
 e^{-\pi q|w|^2/2}=2\partial_w-2\pi q\bar w.
\tag{PF24}
\] The summands commute since \(\partial_w\bar w=0\). In its n-th power every term containing \(\bar w\) vanishes at zero. The remaining term is \(2^n\partial_w^n\), and \(u=-\Omega_\tau w\) contributes \((-\Omega_\tau)^n\). Dividing by \(\mathrm{fockScale}^n=(2\sqrt{\pi q_\tau})^n\) gives the displayed coefficient. Distinct matrix coordinates commute, proving every monomial and then every polynomial. Equation(PF23) follows from(PF10). ∎

There is no hidden parity change in(PF21). In the normalized coordinates of DO-018 let \(z=A+\tau_{\rm per}B\), \(Y=\operatorname{Im}\tau_{\rm per}\). The archimedean Schrödinger action is \[
\rho_\infty(A,B)\varphi(x)
 =e^{-2\pi iA^{\mathsf T}x+\pi iA^{\mathsf T}B}\varphi(x-B).
\] On \(\varphi_\infty(x)=e^{\pi i x^{\mathsf T}\tau_{\rm per}x}\), theta evaluation gives \[
\begin{split}
\vartheta_{\rho_\infty(A,B)\varphi_\infty\otimes\varphi_f}
 &=e^{\pi iB^{\mathsf T}z}
      \sum_q\varphi_f(q)e^{\pi iq^{\mathsf T}\tau_{\rm per}q-2\pi iq^{\mathsf T}z}\\
 &=e^{-\pi H(z,z)/2}\,R\varphi_f(-z).
\end{split}
\tag{PF25}
\] The second equality is the exact identity \[
-\frac\pi2 z^{\mathsf T}Y^{-1}\bar z
+\frac\pi2 z^{\mathsf T}Y^{-1}z=\pi iB^{\mathsf T}z.
\] Thus(PF21) returns precisely the DO-018/026 reduced function \(R\varphi_f\), with its original finite translations and phases. It applies to every finite Schwartz input. The simultaneous transport \(b_{\rm src}=-\operatorname{Im}H\), \(\psi_{\rm src}=\exp(-2\pi i\,\cdot)\) gives the same realized Heisenberg multiplier. No relabeling of finite translations is needed.

The actual algebraic-coordinate creation operator \(r_\tau\) of ALP-LIFT-09 has commutator/norm coefficient \(\pi h_\tau\); the source operator is \(\lambda_\tau r_\tau\), with coefficient1 by(PF23). Passing from K033/K034 to source pureTensorO must retain these complex period factors. No finite valuation is assigned to them.

\subsection{5. The all-polynomial arithmetic frame equation}\label{the-all-polynomial-arithmetic-frame-equation}

For a reduced germ F set \[
\mathcal J_{cD,F}(P)
 =\sum_m P_m\left(\prod_{\tau,r,s}cD_\tau^{m_{\tau,r,s}}\right)
                         \partial^m F(0).
\tag{PF26}
\] This is exactly jetL(cD,F): jetC uses ordinary derivatives, not Taylor coefficients divided by \(m!\). Let \(\mathcal D_s\) multiply every variable at τ by \(s_\tau\); the literal mvZ(1,diag(s)) does this. Monomial evaluation gives \[
\boxed{\mathcal J_{cD,F}(\mathcal D_{\lambda/cD}P)
                   =P(\lambda_\tau\partial_\tau)F(0).}
\tag{PF27}
\]

Fix actual finite class representatives \(g_c,h_{c'}\), a finite Schwartz datum \(\varphi_d\), and the independently constructed finite Mumford averages \(F_{d;c,c'}\) of DO-026 on \(A_h^{ab}\). They are defined geometrically before the analytic comparison. Let \(\epsilon_\chi(c,c')\) be the actual determinant-root splitting factor. The seed is the product of \(14hab\) copies of \(\sigma_q\), and its fixed comparison period is \(C^{14hab}\), \(C=\theta_0c\), with the phase fixed in AM-R-006. The product and finite Mumford intertwining retain this same period for every moved lattice and every polynomial. Consequently \[
\boxed{
\vartheta_{\mathrm{pureTensorO}(\varphi_d,P)}(g_c,h_{c'})
 =C^{14hab}\epsilon_\chi(c,c')\,
      \mathcal J_{cD,F_{d;c,c'}}
               (\mathcal D_{\lambda/cD}P).
}
\tag{PF28}
\] The finite Weil-phase premise is exactly the fixed unramified-CM realization of DO-022 and its product-host transport. It is not an assertion about arbitrarily chosen nonzero vacuum vectors. All finite phases and normalized-average indices of DO-026 remain in F.

The required varying source frame is explicit: \[
\boxed{\mathrm{sc}_{j,\tau}
       =-\frac{\Omega_\tau}{cD_\tau\sqrt{\pi q_\tau}}.}
\tag{PF29}
\] For a stage it is independent of degree and finite chart. Take \[
\mathrm{abelian}^{\rm geom}_2(c,c')
  =\epsilon_\chi(c,c')\mathcal J_{cD,F_{d;c,c'}}.
\] With \(c_K=v_jC^{14hab}\), equation(PF28) is(PF3) for this actual geometric family. If the chart characters have already been separated as in LevChiEqOne.abelian₂, use those characters and the common root-of-unity normalization of DO-027; its scalar criterion is unchanged. This is not a definition by division of an analytic value by a desired period: the sections and finite averages determine the functional before the comparison.

The numerical period/frame comparison left in DO-026§5 is therefore explicit for these kernels. The varying profile(PF10) replaces the uniform quadratic scale, with the source unit-Fock normalization preserved.

\subsection{6. One field for all jet orders in a fixed finite chart family}\label{one-field-for-all-jet-orders-in-a-fixed-finite-chart-family}

\subsubsection{Proposition PF30}\label{proposition-pf30}

Fix finitely many actual moved module pairs, finite translations, and character/Schwartz coefficients in(PF28). There is one number field \(M_1\), independent of derivative order, containing all values \[
\mathcal J_{cD,F_{d;c,c'}}(P),\qquad P\in M_1[Z_{\tau,r,s}].
\tag{PF30}
\]

\textbf{Proof.} The product ppav, full endomorphism action, rigidified trace-one line and normalized algebraic section are defined over a number field. Each commensurable CM lattice in the finite list is algebraically isogenous to that product: clearing denominators gives an isogeny with finite torsion kernel, whose algebraic quotient realizes the analytic torus. After a finite extension the kernels, rigidified polarizations and finitely many line isomorphisms are defined. Enlarge once for the finitely many torsion points, finite coefficients, roots of unity and square roots of averaging indices.

Choose algebraic formal frames and invariant tangent bases over this field. Their formal group laws, addition and adjunction cocycles have coefficients there. The full diagonal \(\mathcal O_E\)-action on every isogenous ordinary-Jacobian product retains holomorphic labels1,2,3 with repetitions, without opposite pairs. The accepted existence theorem and same-frame uniqueness in K018 identify the algebraic canonical gauge with the reduced Chern gauge over that same field in all orders. The rigidified translation gauges have the same algebraic descent; K034 gives the stronger integral statement at11 for prime-to11 torsion.

Hence the finitely many reduced section germs are complete series over one number field, rather than requiring a different extension for each coefficient. Invariant logarithm substitutions also have coefficients there in characteristic zero. Adjoin the finitely many explicit \(cD_\tau\) if necessary. All formal derivatives and finite sums in(PF26) remain in this field. ∎

No effective discriminant or degree bound is asserted for \(M_1\). This is not yet literal Amb.JetsRat in ThDistC.lean239--243: that statement quantifies over \textbf{every} selfdual \(N,N'\), all admitted \(c,x\), with one preassigned X.K. Passing from a finite class family to that uniform statement needs its prescribed algebraic transport and character descent. Adjoining an unrelated field for each lattice would not prove it.

\subsection{7. The next exact arithmetic equations and cost scope}\label{the-next-exact-arithmetic-equations-and-cost-scope}

The calculation separates three obligations.

\begin{enumerate}
\def\labelenumi{\arabic{enumi}.}
\tightlist
\item
  \textbf{Solved for the geometric family:} the periods, embedding-line generators, literal source creation factors, and all-polynomial equation(PF28). They do not need a common imaginary quadratic subfield.
\item
  \textbf{Still required for the prescribed arithmetic row:} identify the geometric \(\epsilon_\chi\mathcal J_{cD,F}\) with the independently fixed row functional and map. The exact next equation is ClsK.Kval, \nolinkurl{OnFramesRdApply.lean115–121:} \[
  \operatorname{val}_{j+1}(t.K_j\psi,c')
   =\sum_c|\operatorname{Stab}(c)|^{-1}
      \mathrm{abelian}_2(c,c')(\mathrm{subY}(F_c))
           \det(y')^{-s_D},
  \tag{PF31}
  \] with the actual selected polynomial, source row, target row and determinant factors. DO-026 computes this class integral for its actual operator, but local type alone does not identify an externally specified \(t.K_j\) or the integral saturation of a later image.
\item
  \textbf{Still required for a uniform arithmetic constructor:} Amb.JetsRat with its fixed descent field and the remaining constructor hypotheses or separately proved replacements. The old false Fields.Kq record remains false; PF30 is a finite-family result.
\end{enumerate}

At11 all cD factors are units, so K033/K034's selected arithmetic content bounds hold unchanged for jetC(cD,−). Their matching slope1/1220 is before multiplication by the complex factors(PF29). The Fock norm is preserved exactly by(PF23), and the analytic common period is retained. No full DP44 cost estimate, later-image saturation, terminal identity, or every-\(\varepsilon\) window follows from this adapter.

\subsection{Sources and verification}\label{sources-and-verification}

The primary period source is Tadokoro, Y. (2014), \emph{The period matrix of the hyperelliptic curve \(w^2=z^{2g+1}-1\)}, \href{https://arxiv.org/abs/1211.6910v5}{arXiv:1211.6910v5}, Propositions3.2 and4.1 and Theorem4.4, as instantiated and checked in AM-R-004. Reduced theta multipliers and normalized translation are as in Bannai, K., \& Kobayashi, S. (2010), \emph{Algebraic theta functions and p-adic interpolation of Eisenstein--Kronecker numbers}, Duke Mathematical Journal, 153(2),229--295, Definition1.8 and§1.3, \href{https://arxiv.org/pdf/math/0610163}{author preprint}. Both bibliographic titles were checked against their primary arXiv records during this review. The creation computation is directly from the recovered literal definitions above; it is not an asserted extension of a source theorem beyond its hypotheses.

The companion certificate checks exact cyclotomic interpolation/norms, finite scalar/jet normalizations, and the source-sign Gaussian identity as polynomial identities. It is a finite form check, not a substitute for the all-order proofs, a transitive audit of accepted theorems, or a Lean build. Actual hashes and check events are recorded separately.

\clearpage

\section{AM-DO-037: the prescribed coupled Q family survives the actual second local theta quotient}\label{am-do-037-the-prescribed-coupled-q-family-survives-the-actual-second-local-theta-quotient}

Author: GPT-6 Astra Pro, ALP\_MACHINE /root/definite\_oscillator. 22 November 2026 UTC. Inherited reasoning effort is not separately exposed.

The literal two-edge Q family for m=2 has a nonzero image under the actual successive local theta quotients. The initial representation is any irreducible unitary spherical representation of GL\_k at that Q place. A specified unipotent target evaluation, followed by evaluation in the first induced model, is an explicit nonzero scalar times the initial spherical vector. The proof retains the entire prescribed Y sum and both raw kernels. It does not replace either kernel by an unrestricted Schwartz input.

The actual source definition of DegOne forces the local field at every tuple Q place to be Q\_q. Thus this hypothesis is realized on both the K031 and K032 hosts, even though those hosts have absolutely ramified places elsewhere. The exact local quotient scalar below is normalized by explicit Haar integrals. Identification with the separately pinned global arithmetic row and terminal mass is not asserted.

\subsection{1. The literal Q-place hypothesis includes absolute unramifiedness}\label{the-literal-q-place-hypothesis-includes-absolute-unramifiedness}

All source locators refer to the pinned source tree at commit 7cbf7a0c32fca1badf05931879a37c3f8a2b03b4. In \nolinkurl{Formalism/LocalToGlobal/Divisibility/Tuples.lean:}

\begin{longtable}[]{@{}
  >{\raggedright\arraybackslash}p{(\columnwidth - 2\tabcolsep) * \real{0.5000}}
  >{\raggedright\arraybackslash}p{(\columnwidth - 2\tabcolsep) * \real{0.5000}}@{}}
\toprule\noalign{}
\begin{minipage}[b]{\linewidth}\raggedright
Lines
\end{minipage} & \begin{minipage}[b]{\linewidth}\raggedright
Literal condition
\end{minipage} \\
\midrule\noalign{}
\endhead
\bottomrule\noalign{}
\endlastfoot
58--60 & SplitAt: the extension of the real-field prime is a product of two distinct CM primes \\
64--65 & DegOne(v): its absolute norm is prime, and that rational prime is not in v squared \\
67 & GoodPl is SplitAt, DegOne, and avoidance of the prescribed rational set S \\
112--114 & A tuple kills local inertia at every place which is not GoodPl \\
128 & Qprime is the set where that inertia character is nontrivial \\
149--161 & IsResChar and resChar give its tame residue character \\
198--200 & IsTupleMap excludes the coefficient-place rational prime from the corresponding Q set \\
\end{longtable}

Let v belong to Qprime of a tuple. The last condition in tupGrp implies GoodPl(v): otherwise the inertia character would be trivial. Write q=Norm(v). The first DegOne clause says q is a rational prime, so the residue degree is one. In the localization O\_(F,v), the valuation of the rational prime q is the absolute ramification index e. The second clause q notin v\^{}2 says e=1. Hence

\[
                 [F_v:\mathbf Q_q]=ef=1.
 \tag{DO37.1}
\]

Relative splitting gives E\_P=F\_v for either selected CM prime P. Thus the precise local hypothesis used in LIFT23 is automatic. In particular the absolutely ramified29 places of K031/K032 cannot occur in Q; no replacement of q by a ramified uniformizer is needed for this literal tuple family.

The tuple is trivial on base-field ideles and kills wild inertia. At a split place its two residue characters are therefore inverses. Since the place belongs to Q, both are nontrivial. Consequently the selected residue character chi used below is nontrivial. If the tuple is evaluated at a coefficient place over11 through IsTupleMap, then q is different from11 as well. For a bare tuple without that additional condition the local theorem below still holds at q=11, but its Gauss sum is not claimed an11-unit.

This argument checks an actual definition. It does not use the false implication that residue degree one alone gives F\_v=Q\_q.

\subsection{2. A whole-Schwartz normalized quotient map}\label{a-whole-schwartz-normalized-quotient-map}

This section records the precise published interface needed to pass from an explicit integral to the actual small theta quotient. Let n\textless m, delta=m-n, G\_n=GL\_n(F), K\_n=GL\_n(O), and

\[
 \kappa_n(q)=\prod_{i=1}^{n}(1-q^{-i}),\qquad
 \operatorname{vol}(K_n)=1,\qquad
 \operatorname{vol}(M_n(O))=1.
 \tag{DO37.2}
\]

Thus multiplicative matrix Haar is dg=kappa\_n(q)\^{}(-1)\textbar det g\textbar\^{}(-n)d\_add g. Let eta be unitary and pi an irreducible unitary representation of G\_n.~Use the actual split Hom action

\[
 [\omega_\eta(g,h)\phi](X)=
 \eta(\det h)\eta(\det g)^{-1}
 |\det h|^{-n/2}|\det g|^{m/2}\phi(h^{-1}Xg).
 \tag{DO37.3}
\]

For E\_n={[}I\_n;0{]} and Re(z) sufficiently large, put

\[
 I_z(\phi,v)=\int_{G_n}
       \phi(E_ng)|\det g|^{z+(n-1)/2}
                 \eta(\det g)^{-1}\pi(g)v\,dg.
 \tag{DO37.4}
\]

This is defined on the whole Schwartz space M\_(m,n), not just on its full-rank open stratum. Restriction to E\_n M\_n(F) is a Schwartz function on M\_n(F). Pairing (DO37.4) with any smooth dual vector is therefore a Godement--Jacquet zeta integral. Jiang and Luo (2023), Theorem2.4(1)--(2), printed pp.313--314, states that every such integral is L(z,pi eta\^{}(-1)) times a Laurent polynomial in q\^{}z.

The vector-valued assertion follows without a weak-versus-strong holomorphy assumption. For fixed phi, the restricted Schwartz function on M\_n(F) is invariant under left multiplication by a sufficiently small open compact J contained in ker(eta det). Changing variables shows I\_z(phi,v) lies in pi\^{}J for every z where it converges. Admissibility makes pi\^{}J finite dimensional. Apply the scalar theorem to finitely many smooth dual vectors separating this space. Thus I\_z is a meromorphic vector with no pole wherever L(z,pi eta\^{}(-1)) is finite.

Direct changes of variables, first g to gg\_0 and then g to At, give

\[
 I_z(\omega(g_0,1)\phi,\pi(g_0)v)
       =|\det g_0|^{(\delta+1)/2-z}I_z(\phi,v),
 \tag{DO37.5}
\]

and, for p=block(A,B;0,D) in the upper parabolic P\_(n,delta),

\[
 I_z(\omega(1,p)\phi,v)=
 \eta(\det D)|\det A|^{z-1/2}|\det D|^{-n/2}
                       \pi(A)I_z(\phi,v).
 \tag{DO37.6}
\]

At z\_0=(delta+1)/2 these are simultaneous source coinvariance and precisely normalized parabolic covariance. Consequently

\[
 T_{\eta,\pi}(\phi\otimes v)(h)
                  = I_{z_0}(\omega(1,h)\phi,v)
 \tag{DO37.7}
\]

is, when L is finite there, an actual map from the whole oscillator tensor pi to

\[
           \pi\times(\eta\circ\det_\delta)
      =\operatorname{Ind}_{P_{n,\delta}}^{G_m}
                    (\pi\boxtimes\eta\circ\det_\delta),
 \tag{DO37.8}
\]

where induction is normalized. It intertwines the target with right translation. Holomorphy in the induced compact model is also finite dimensional for a fixed input: phi has a fixed open compact target stabilizer, so only finitely many compact cosets and the fixed finite-dimensional pi\^{}J are needed.

Tadic (1986), TheoremD(i)--(ii), printed pp.338--339, implies that normalized induction of irreducible unitary representations of general linear groups is irreducible and unitary. Indeed express the factors as products of the elementary unitary factors in part(ii), concatenate them, and apply part(i). Hence (DO37.8) is irreducible. Once (DO37.7) is nonzero, Minguez (2008), Theorem1(1), printed p.718, identifies it with the actual irreducible local theta quotient, with multiplicity one.

The source is paired bilinearly. Thus the simultaneous coinvariant uses pi, while the source irreducible quotient of the oscillator alone is pi contragredient; this is the same Hom convention as DO028. The old parameter is therefore pi in (DO37.8). No source contragredient has been dropped.

Only the holomorphic-multiple assertion of Jiang--Luo is used. Its displayed Gaussian convention is not used to fix Haar constants; all kappa factors below follow from (DO37.2). This avoids an implicit conversion between additive and maximal-compact probability Haar measures.

\subsection{3. The representations and the two raw kernels}\label{the-representations-and-the-two-raw-kernels}

At the actual Q place put

\[
 a=k,\quad r=a+1,\quad d=t+1,\quad b=a+d=r+t,\quad N=b+d.
 \tag{DO37.9}
\]

Let pi\_0 be an irreducible unitary spherical representation of GL\_a(F), and take v in pi\_0\^{}(K\_a). Write eta\_0,eta\_1 for the two actual common stage labels. Their residue characters at Q are the same chi, since each is the unramified reference label times the same tuple character. Set

\[
                  \vartheta=\eta_0\eta_1^{-1}.
 \tag{DO37.10}
\]

This is an unramified unitary character. In particular \textbar vartheta(q)\textbar=1, with \textbar q\textbar=q\^{}(-1).

For the first quotient, L(z,pi\_0 eta\_0\^{}(-1))=1: the parameter of pi\_0 is unramified, whereas twisting by the nontrivial inertia character chi\^{}(-1) leaves no inertia invariants. The map (DO37.7) is holomorphic at z=(d+1)/2 and has target

\[
                  \pi_1=\pi_0\times(\eta_0\circ\det_d).
 \tag{DO37.11}
\]

It is nonzero, for example on a Schwartz function supported near E\_a K\_a with the matching eta\_0 character. Thus it is the actual first local theta quotient by Section2.

For the second quotient,

\[
 L(z,\pi_1\eta_1^{-1})
     =\prod_{j=0}^{d-1}
           (1-\vartheta(q)q^{-z+(d-1)/2-j})^{-1}.
 \tag{DO37.12}
\]

The old pi\_0 part again has no inertia invariants. The other part is the standard parameter of the determinant character on GL\_d, with its d consecutive exponents. At z=(d+1)/2 the factors are (1-vartheta(q)q\^{}(-i))\^{}(-1), 1\textless=i\textless=d. Every factor is finite and nonzero. Section2 therefore gives an honest whole-Schwartz second map into the irreducible

\[
                  \pi_2=\pi_1\times(\eta_1\circ\det_d).
 \tag{DO37.13}
\]

We will prove that this second map is nonzero on the exact coupled first vector. The repeated tame character causes no problem for this argument; we do not invoke the distinct-inertia boundary exclusion used for G in LIFT21.

Let F\_y(X) and H\_y(Y) be the two actual local functions of LIFT23(Q5--Q7), at the actual depth mathsf M\textgreater=1. Their common Y set is q\^{}(-mathsf M)O\^{}b/q\^{}(mathsf M)O\^{}b. Write X=(A;C), with A r-by-a and C t-by-a, and write Y=(Z;eta). The checked full finite contraction is

\[
 {\cal C}(X,Z,\eta)=\sum_y F_y(X)H_y(Z,\eta)
 \tag{DO37.14}
\]

with value

\[
 {\cal C}=
 q^{-d}G(\chi,\psi_0)\chi(-c)^{-1}\chi(\det\bar Z_r)
 \tag{DO37.15}
\]

exactly when:

\begin{itemize}
\tightlist
\item
  A is integral and has full residue rank a;
\item
  C belongs to q M\_(t,a)(O);
\item
  Z is integral, its leading r-by-r block Z\_r is invertible modulo q, and its lower-left block is q-divisible;
\item
  eta belongs to q\^{}(-1)O\^{}b and overline(q eta\_top)=c ell\_A with c nonzero.
\end{itemize}

Here ell\_A z=det{[}bar A,z{]} is the oriented cofactor row. All other values are zero. The lower part of overline(q eta) is unrestricted. The Gauss sum is G(chi,psi\_0)=sum\_(s!=0) chi(s)psi\_0(s), and \textbar G\textbar\^{}2=q. This identity includes the literal q\^{}(-mathsf M b) factor and all higher-depth cancellation; it is carried from the independently checked LIFT23, not replaced by a new function.

\subsection{4. A specified target evaluation of the unchanged coupled vector}\label{a-specified-target-evaluation-of-the-unchanged-coupled-vector}

Let E={[}I\_a;0{]} be b-by-a, let e\_r\^{}* be the row with its one in position r=a+1, and put

\[
 B_*=\begin{pmatrix}I_b\\0_{(d-1)\times b}\\q^{-1}e_r^*\end{pmatrix}
                      \in M_{N,b}(F).
 \tag{DO37.16}
\]

This is a full-rank matrix. It equals h\_\emph{\^{}(-1)E\_b for the explicit lower unipotent h\_} in GL\_N(F) whose only extra last-row entry is -q\^{}(-1) in column r. In particular det h\_\emph{=1. Evaluation of the induced second output at h\_} therefore uses precisely phi(B\_*h), with no extra target character or determinant factor. This is a choice of target evaluation point, not a change of either finite Schwartz input.

These are the normalized terminal coordinates of LIFT23. In unscaled global coordinates, if the source last-row functional is piT\^{}(-1) times the coordinate projection, the last row of B\_* is piT q\^{}(-1)e\_r\^{}* and the corresponding entry of h\_* is -piT q\^{}(-1). At Q, piT is a unit and det(h\_*) remains one. Thus the literal last-row unit is retained and contributes no omitted target determinant scalar or additive Haar factor.

Form the actual coupled vector

\[
 {\cal V}(v)=
       \sum_y H_y\otimes T_{\eta_0,\pi_0}(F_y\otimes v)
                       \in\mathcal S(M_{N,b})\otimes\pi_1 .
 \tag{DO37.17}
\]

We evaluate its second quotient at h\_*, and then evaluate the resulting pi\_1-valued function at the identity in the first induced model. To compute this value, temporarily put z=(d+1)/2+u in the second meromorphic map. For Re(u) large, substitution of (DO37.4), (DO37.7) and the first target action gives exactly

\[
 \begin{split}
 {\cal E}_u(v)=\int_{\operatorname{GL}_b(F)}
       \int_{\operatorname{GL}_a(F)}
 &{\cal C}(h^{-1}Eg,Z_0h,q^{-1}e_r^*h)\,
   \vartheta(\det h)|\det h|^{d+u}\\
 &\hspace{7mm}\cdot
       \eta_0(\det g)^{-1}|\det g|^{b/2}\pi_0(g)v\,dg\,dh ,
 \end{split}
 \tag{DO37.18}
\]

where Z\_0={[}I\_b;0\_(d-1,b){]}. The middle exponent is not guessed: the second source factor has exponent z+(b-1)/2=(N/2)+u, while the first target factor has exponent -a/2. Their sum is d+u because N-a=2d. Their character ratio is exactly eta\_0 eta\_1\^{}(-1).

The following computation proves absolute convergence of this double integral already for Re(u)\textgreater-1. Hence its value at u=0 equals the holomorphic continuation defining the actual whole-Schwartz quotient, by equality in the initial common convergence half-plane and uniqueness of continuation. The argument also bounds the sum of the absolute values of the finitely many uncontracted terms. A nonzero first determinant already forces its a columns to have full residue rank; the support argument below therefore forces g into K\_a for each individual Y term before any Fourier cancellation. The same Schur-complement domain bounds all those terms.

\subsection{5. Support, cofactor cancellation and the Schur complement}\label{support-cofactor-cancellation-and-the-schur-complement}

The terminal upper factor forces h to have the block form

\[
 h=\begin{pmatrix} U&V\\qW&D\end{pmatrix},\qquad
 U\in\operatorname{GL}_r(O),\quad
 V,W,D\text{ integral}.
 \tag{DO37.19}
\]

Its Schur complement

\[
                        S=D-qWU^{-1}V
 \tag{DO37.20}
\]

belongs to M\_t(O) and is invertible over F. Write X=h\^{}(-1)Eg=(A;C). If the first raw support is nonzero, then X is integral and C is q-divisible. The equation hX=Eg shows g is integral. Modulo q its top block says

\[
                  \bar U\bar A=E_r\bar g,
          \qquad E_r=\begin{pmatrix}I_a\\0\end{pmatrix}.
 \tag{DO37.21}
\]

Since U is invertible and A has full residue rank a, g is in GL\_a(O). Conversely, once g is a unit matrix, the integrality and rank conditions on A follow from C being q-divisible and the block equations.

For such a g, the cofactor row is exactly

\[
          \ell_A=(\det\bar U)^{-1}
                         (\det\bar g)\,e_r^*\bar U.
 \tag{DO37.22}
\]

Indeed multiply the augmented matrix {[}bar A,z{]} by bar U; its determinant becomes det{[} E\_r bar g,bar U z{]}. On the other hand the upper part of q eta is e\_r\^{}*bar U. Thus the proportionality constant in (DO37.15) is c=det(bar U)/det(bar g), which is nonzero, and

\[
 \chi(\det\bar U)\chi(-c)^{-1}
                  \eta_0(\det g)^{-1}=\chi(-1)^{-1}.
 \tag{DO37.23}
\]

The unramified part of eta\_0 is one on these determinants. This is the exact cancellation of the same tame character across the two raw Q factors and the initial source action. There is no averaging which could discard the Gauss factor.

Put W'=WU\^{}(-1). Solving the block equations gives

\[
       C=-qS^{-1}W'E_r g.
 \tag{DO37.24}
\]

Consequently C is q-divisible precisely when the first a columns of W' lie in S M\_(t,a)(O). We can and must write

\[
                  W'=[S L,w],\quad
 L\in M_{t,a}(O),\quad w\in O^t .
 \tag{DO37.25}
\]

Conversely every U,V,S,L,w with these conditions produces a matrix h in the support by setting W=W'U and D=S+qW'V. Its inverse applied to E has integral top block of full residue rank and q-divisible bottom block. Thus (DO37.25) parametrizes the entire nonzero domain, not only a convenient subset of it.

Since pi\_0(g)v=v and \textbar det g\textbar=1 on the forced K\_a, the normalized g integral is simply v. Also vartheta(det h)=vartheta(det S), since det U is a unit.

\subsection{6. Exact Haar integral and nonzero quotient value}\label{exact-haar-integral-and-nonzero-quotient-value}

The block lower-left variable qW contributes the additive Jacobian q\^{}(-rt). The change D to S has Jacobian one. The change W to W'=WU\^{}(-1) has Jacobian one because U is integral unimodular. For fixed S, replacing its first a columns by S L has Jacobian \textbar det S\textbar\^{}a.

Multiplicative Haar contributes kappa\_b\^{}(-1)\textbar det S\textbar\^{}(-b). Multiplying by the weight \textbar det S\textbar\^{}(d+u) from (DO37.18), and the preceding Jacobian, leaves exactly

\[
                     |\det S|^u,
                \qquad -b+d+a=0.
 \tag{DO37.26}
\]

The U integral has additive volume kappa\_r; all V,L,w integral boxes have volume one. Therefore

\[
 {\cal E}_u(v)=q^{-d}G(\chi,\psi_0)\chi(-1)^{-1}
       q^{-rt}\frac{\kappa_r}{\kappa_b}
       \left(\int_{M_t(O)}
          \vartheta(\det S)|\det S|^u\,dS\right)v.
 \tag{DO37.27}
\]

Singular matrices have measure zero. For Re(u)\textgreater-1 the matrix integral is

\[
 J_t(u,\vartheta)
   =\prod_{i=1}^t
          \frac{1-q^{-i}}{1-\vartheta(q)q^{-u-i}}.
 \tag{DO37.28}
\]

For completeness, partition the first column by its common valuation n.~Its measure is q\^{}(-tn)(1-q\^{}(-t)), its determinant factor is vartheta(q)\^{}n q\^{}(-un), and an integral row operation takes its primitive part to the first standard vector. Such a row operation has unit determinant, on which vartheta is trivial. The remaining determinant is a (t-1)-by-(t-1) integral matrix. Summing the geometric series proves J\_t=(1-q\^{}(-t))/(1-vartheta(q)q\^{}(-t-u)) J\_(t-1), with J\_0=1, which gives (DO37.28).

Taking absolute values uses the same recurrence with vartheta=1 and Re(u). This proves the asserted absolute convergence, including at u=0. There is no noncompact tail estimated by cancellation in this computation.

We obtain the exact theorem:

\[
 \boxed{\quad
 \operatorname{ev}_1\!
   \left[T_{\eta_1,\pi_1}({\cal V}(v))(h_*)\right]
                      =\Lambda_Q\,v ,
 \quad}
 \tag{DO37.29}
\]

where

\[
 \boxed{\quad
 \Lambda_Q=
 q^{-d}G(\chi,\psi_0)\chi(-1)^{-1}
 q^{-t(a+1)}
 \frac{\kappa_{a+1}\kappa_t}{\kappa_{a+d}}
 \prod_{i=1}^t(1-\vartheta(q)q^{-i})^{-1}\ne0 .
 \quad}
 \tag{DO37.30}
\]

Every factor is specified. Nonvanishing follows from the Gauss identity and \textbar vartheta(q)\textbar=1. Thus the unchanged coupled vector has nonzero image in the actual second local theta quotient. In fact its restriction to the initial K\_a-fixed source space is injective: the two specified evaluations compose with it to the nonzero scalar Lambda\_Q times the identity.

This is a statement about the full coupled family. It does not assert that every individual Y class survives separately, or that the fixed first class y=e\_a already has the full middle parabolic covariance.

\subsection{7. Norms, arithmetic scalar, and global scope}\label{norms-arithmetic-scalar-and-global-scope}

With the given source norm,

\[
 \left\|\operatorname{ev}_1
      [T_{\eta_1,\pi_1}({\cal V}(v))(h_*)]\right\|
                            =|\Lambda_Q|\,\|v\|.
 \tag{DO37.31}
\]

This is the norm of that evaluation in pi\_0, not the norm of the full induced function or of a global automorphic operator. The raw joint Schwartz norm remains exactly LIFT23(Q18); no new raw input normalization is imposed. The local scalar has the explicit bounds

\[
 q^{-d+1/2-rt}\frac{\kappa_r\kappa_t}{\kappa_b}
       \prod_{i=1}^t(1+q^{-i})^{-1}
 \ \le |\Lambda_Q|\ \le\
 q^{-d+1/2-rt}\frac{\kappa_r}{\kappa_b}.
 \tag{DO37.32}
\]

At a coefficient place over11 with q different from11, the Gauss and character factors are units. Assuming the actual algebraic stage values are retained in the coefficient field, the exact valuation is

\[
 v_{11}(\Lambda_Q)=
 v_{11}(\kappa_r)+v_{11}(\kappa_t)-v_{11}(\kappa_b)
       -\sum_{i=1}^t v_{11}(1-\vartheta(q)q^{-i}).
 \tag{DO37.33}
\]

The q powers have valuation zero. The terms on the right need not vanish. In particular primitive values of the raw finite array do not automatically imply a primitive integrated quotient value after its Haar/Euler factors. Equation (DO37.33) records this exact local arithmetic cost rather than replacing it by a unit assertion.

For any actual source tuple on K031 or K032, Section1 shows the theorem applies at each of its Q places. The global finite Y set factors into these local quotients by the checked CRT statement DO032, and the local tensor product of (DO37.29) has scalar product\_P Lambda\_(Q,P). This proves the Q part of the prescribed two-edge local quotient comparison on the initial spherical factors.

It does not, by itself, assemble the full global coupled theta composite: the real factors, T factors, global cuspidal embeddings and their multiplicities, and the prescribed row/frame comparisons must also be included. The quotient maps used here have been normalized by (DO37.4) and our exact Haar measures. Their correspondence with an independently imposed global pin must retain that normalization; uniqueness of the abstract quotient does not set a missing global scalar to one.

\subsection{8. References and checks}\label{references-and-checks}

Carried, already checked interfaces are LIFT23(Q1--Q20), DO035's actual common labels and initial Q invariance, DO028's bilinear Hom convention, and DO032's literal Y quotient and global CRT factorization.

The raw source definitions are in \nolinkurl{Formalism/LocalToGlobal/Columns/FiniteSumNonzeroLargeDepth.lean:} entB684--694, PhiQ711--715, joinCol982--983, PsiAt985--986, zero\_h1020, last\_h1025, ybSet1125--1129, XiCpl11134--1138, PhiTop1142--1147, ftY1150--1156, and XiCpl21160--1167. The forced local-field hypothesis has the direct source proof in Section1 above.

\begin{itemize}
\tightlist
\item
  Jiang, D., \& Luo, Z. (2023). Certain Fourier operators and their associated Poisson summation formulae on GL1. Pacific Journal of Mathematics, 326(2), 301--372. Theorem2.4(1)--(2), printed pp.313--314. \url{https://doi.org/10.2140/pjm.2023.326.301}
\item
  Tadic, M. (1986). Classification of unitary representations in irreducible representations of general linear group (non-Archimedean case). Annales scientifiques de l'Ecole Normale Superieure, 19(3), 335--382. TheoremD(i)--(ii), printed pp.338--339. \url{https://www.numdam.org/item/10.24033/asens.1510.pdf}
\item
  Minguez, A. (2008). Correspondance de Howe explicite: paires duales de type II. Annales scientifiques de l'Ecole Normale Superieure, 41(5), 717--741. Theorem1(1), printed p.718, with the normalized oscillator convention in equation1.4, pp.721--722.
\end{itemize}

The author whole-proof review and exact finite Schur-complement replay passed after the complete text was read. The companion \nolinkurl{definite_oscillator_coupled_q_certificate.py} checks23,328 exact rational block matrices and cofactor phases,156 lattice Jacobians, and90 determinant-integral recurrences. These diagnostics support the calculations; the actual quotient assertion uses the full proof and the cited representation theorems. Separate check events record the author and independent reviews. No Lean execution is claimed.

\clearpage

\section{ALP-LIFT-25: the two actual selected compact kernels compose on the first highest line}\label{alp-lift-25-the-two-actual-selected-compact-kernels-compose-on-the-first-highest-line}

Author: GPT-6 Astra Pro, ALP\_MACHINE /root/lift\_accounting. 22 November 2026 UTC. Inherited reasoning effort is not exposed.

The primitive selected highest polynomial on the first definite edge produces a fixed highest right vector in the actual first compact type. The primitive selected highest polynomial on the second edge pairs nontrivially with exactly that vector. The compact contraction is computed below, including its Haar denominator and its actual Gaussian, creation and hook norms. This is a preceding-image statement for the specific first vector used in DO036, rather than merely an occurrence of the required compact types.

Two normalizations are explicitly distinguished. The literal compact coefficient integral has factor 1/dim(V\_j). An auxiliary unitary identification of the compact theta quotient has factor 1/sqrt(dim(V\_j)), because the diagonal invariant tensor has squared norm dim(V\_j). The latter identification does not set an independently pinned global theta scalar to one.

The theorem concerns the actual selected highest kernels. A fixed second highest kernel can annihilate a nonhighest vector in the same compact representation. No assertion about arbitrary first targets is made by replacing that distinction with multiplicity-one terminology.

\subsection{1. The actual selected representations and primitive vectors}\label{the-actual-selected-representations-and-primitive-vectors}

Use the genuine mixed CM host of K031 or K032, with compact real groups

\[
 K_j=\prod_{\tau\in\Psi}U(N_j),\qquad
 (N_0,N_1,N_2)=(a,b,c)=(a,a+d,a+2d).
 \tag{AR1}
\]

The actual stage characters have the same label on both members of each edge and infinity exponent -e\_(j,tau). Write V\_j for the irreducible representation with the literal weight w\_(j,tau)=lamTrue\_j at every place. It has the source row metric: Fischer divided by the Fischer norm squared of the primitive extremal. Let h\_j be that extremal, so its row norm is one. Its leading diagonal monomial has coefficient one, fixing its phase.

The following equalities are the checked DO025 and R011 identities, used here as inputs:

\[
 \begin{split}
 \nu_{j,\tau}&=w_{j,\tau}+e_{j,\tau}
                                      -\frac{N_{j+1}}2\mathbf1,\\
 \eta_{j,\tau}&=w_{j,\tau}-w_{j,\tau,N_j}\mathbf1,\\
 \nu_{j,\tau}&=\eta_{j,\tau}+q_{j,\tau}\mathbf1,
       \qquad q_{j,\tau}\ge1,\\
 \eta_{j+1,\tau}&=(\nu_{j,\tau},0^d),\\
 \nu_{1,\tau}&=(\nu_{0,\tau},0^d)
                                   +\Delta_\tau\mathbf1_b,
       \qquad \Delta_\tau=e_{1,\tau}-e_{0,\tau}-d\ge2.
 \end{split}
 \tag{AR2}
\]

Here the final entry in w is written with one-based indexing. The actual selected Cauchy summand of the edge-j oscillator is

\[
                         V_{j+1}\otimes V_j^\vee.
 \tag{AR3}
\]

The dual on the source is essential: the source theta integral is bilinear, as in DO25.4. Its compact oscillator representation is det\^{}(e\_j-N\_(j+1)/2) times the dual Schur representation on the source, and det\^{}(-e\_j+N\_j/2) times the Schur representation on the target. Equation (AR2) matches both labels exactly.

For a partition lambda on N\_j rows set

\[
 p_\lambda(Z)=\prod_{r=1}^{N_j}
      \det(Z_{[1,r],[1,r]})^{\lambda_r-\lambda_{r+1}},
                \qquad\lambda_{N_j+1}=0.
 \tag{AR4}
\]

Use these standard principal minors at every CM place, and put

\[
 P_j=\prod_{\tau}p_{\nu_{j,\tau}}(Z_\tau),\qquad
             \phi_j=r_{P_j}\phi_{0,j,\infty}.
 \tag{AR5}
\]

The vacuum and the creation operators are the actual Chern/CM Schwartz realization of DO022/026 and K030/032. Thus (AR5) specifies an actual Schwartz vector, not an abstract free oscillator vector. The first uses the fixed polynomial used in DO034 and DO036, whose source coefficient in R017 was matched to its source line. The conversion to the literal unit-normalized source creation operator is given explicitly in (AR27)--(AR30) below. The second is now specified by the same primitive formula at the actual second selected partition. Its finite factor remains the already specified second raw G/T/Q family; no finite move is introduced here.

The matrix-coefficient realization in R011 gives a unique equivariant Cauchy embedding

\[
 C_j:V_{j+1}\otimes V_j^\vee\longrightarrow\mathcal F_j,
       \qquad C_j(h_{j+1}\otimes h_j^\vee)=\phi_j,
 \tag{AR6}
\]

where h\_j\^{}vee is the unit linear functional taking h\_j to one. The source dual vector in (AR6) is a lowest vector for the usual upper Borel, as required for the dual of a highest source vector; it is not replaced by the unrelated highest vector of V\_j\^{}vee. The determinant and principal-minor convention fixes (AR6) with no sign or arbitrary phase.

Put A\_j=\textbar\textbar phi\_j\textbar\textbar\textgreater0 in the actual real Schwartz norm. By the checked Cauchy/Fischer metric calculation R011.7--8, or uniqueness of the invariant metric on the irreducible product (AR3),

\[
 \langle C_j(v\otimes\ell),C_j(v'\otimes\ell')\rangle
     = A_j^2\langle v,v'\rangle\langle\ell,\ell'\rangle.
 \tag{AR7}
\]

All row highest vectors used in this equation have norm one.

\subsection{2. The literal middle highest vector is the one paired on the next edge}\label{the-literal-middle-highest-vector-is-the-one-paired-on-the-next-edge}

The first target highest vector is the primitive row extremal of partition eta\_1=(nu\_0,0\^{}d), with its determinant shift. Adding zero rows does not change the principal-minor polynomial. On the second source, the required determinant insertion is

\[
       p_{\eta_{1,\tau}}(Z)
                 \longmapsto (\det Z)^{\Delta_\tau}
                                    p_{\eta_{1,\tau}}(Z)
                    =p_{\nu_{1,\tau}}(Z).
 \tag{AR8}
\]

This is an exact polynomial identity. Both sides have diagonal leading coefficient one. R011.10--12 shows that after the source's row metric normalization this determinant insertion is an isometry. It therefore identifies the second source highest functional with h\_1\^{}vee, and

\[
                            h_1^\vee(h_1)=1.
 \tag{AR9}
\]

This proves the needed alignment. Merely knowing that V\_1 and its dual occur in two oscillator representations would not fix (AR9). The actual common character labels and determinant shifts have already been matched in (AR2); no further character twist is made.

\subsection{3. Compact Haar projection: the exact 1/d\_j factor}\label{compact-haar-projection-the-exact-1d_j-factor}

Normalize Haar on every K\_j to have volume one, as in DO026. Let

\[
                    d_j=\dim_\mathbf C V_j.
 \tag{AR10}
\]

Give V\_j\^{}vee its dual unitary metric. For an orthonormal basis e\_1,\ldots,e\_(d\_j), write e\_r\^{}vee for its algebraic dual. The actual simultaneous source average on the selected tensor space is

\[
 \mathcal P_j(\Phi\otimes w)=
       \int_{K_j}\omega_j(g,1)\Phi\otimes\rho_j(g)w\,dg.
 \tag{AR11}
\]

Compact Schur orthogonality, with the bilinear evaluation pairing, gives the exact formula

\[
 \boxed{
 \mathcal P_j\bigl(C_j(v\otimes\ell)\otimes w\bigr)
       =\frac{\ell(w)}{d_j}
              \sum_{r=1}^{d_j}C_j(v\otimes e_r^\vee)\otimes e_r.}
 \tag{AR12}
\]

This is the same Haar normalization as DO26.18. To see the scalar without invoking a norm convention, identify V\_j\^{}vee tensor V\_j with End(V\_j). Averaging a rank-one endomorphism over conjugation gives its trace divided by d\_j times the identity. Its trace is ell(w). Thus (AR12) is basis independent and has denominator d\_j, not sqrt(d\_j).

Let

\[
 \Omega_j(v)=\sum_r C_j(v\otimes e_r^\vee)\otimes e_r.
 \tag{AR13}
\]

By (AR7),

\[
                      \|\Omega_j(v)\|^2
                             =A_j^2 d_j\|v\|^2.
 \tag{AR14}
\]

Consequently the actual highest test has invariant projection

\[
 \mathcal P_j(\phi_j\otimes h_j)
         =d_j^{-1}\Omega_j(h_{j+1})\ne0,
       \qquad
 \|\mathcal P_j(\phi_j\otimes h_j)\|^2=A_j^2/d_j.
 \tag{AR15}
\]

Because the source real group is compact, averaging identifies its coinvariants with its invariant subspace. The Cauchy decomposition is multiplicity free, and (AR2) determines a unique selected partition. Thus (AR15) is a nonzero vector in the actual local compact theta quotient. No extension from a formal type occurrence or from a proper subspace of the oscillator is being assumed.

If Haar volume is V\_(K\_j) rather than one, the integral (AR11) is V\_(K\_j) times the displayed probability projection, and its literal coefficient factor is V\_(K\_j)/d\_j. This has no bearing on the separate global finite class or terminal quotient mass.

\subsection{4. The unitary quotient model and the two-edge formula}\label{the-unitary-quotient-model-and-the-two-edge-formula}

Define the following auxiliary, explicitly normalized identification:

\[
 I_j:V_{j+1}\longrightarrow
       (\mathcal F_j^{\mathrm{sel}}\otimes V_j)^{K_j},
                 \qquad I_j(v)=\frac{\Omega_j(v)}{A_j\sqrt{d_j}}.
 \tag{AR16}
\]

Equation (AR14) says precisely that I\_j is an isometry. It is target equivariant and onto this invariant space. Therefore the corresponding unitary quotient map Q\_j=I\_j\^{}(-1) P\_j satisfies

\[
 \boxed{
 Q_j\bigl(C_j(v\otimes\ell)\otimes w\bigr)
              =\frac{A_j}{\sqrt{d_j}}\ell(w)v.}
 \tag{AR17}
\]

The square root in (AR17) arises solely from (AR16). It is not the scalar of the unnormalized compact coefficient integral.

For the two actual primitive highest vectors (AR5), (AR9) now gives

\[
 \boxed{
 Q_1\bigl(\phi_1\otimes Q_0(\phi_0\otimes w)\bigr)
     =\frac{A_0A_1}{\sqrt{d_0d_1}}h_0^\vee(w)h_2.}
 \tag{AR18}
\]

In particular the composition on the first matched line has exact squared norm multiplier

\[
                            \frac{A_0^2A_1^2}{d_0d_1}>0.
 \tag{AR19}
\]

This is an equality in the explicitly stated local unitary quotient models. If an independently pinned theta quotient map is u\_j Q\_j, the composed scalar is multiplied by u\_0 u\_1. Neither nonvanishing nor multiplicity one fixes these global comparison scalars.

There is a normalization-free version of the conclusion: (AR12) and (AR9) say that the two literal compact invariant contractions are nonzero. Thus no choice of isometric identification is required for the actual local survival assertion.

For any multiplicity space M on which the real groups act trivially, the same calculation tensored with id\_M is a nonzero scalar times the identity on M tensor C h\_0, landing in M tensor C h\_2. It does not require automorphic multiplicity one. In the actual first global kernel, every finite Y term has the same archimedean vector phi\_0. Its right compact component is therefore h\_1 in each representation copy. Finite summation preserves this line; the nonzero first image supplied by DO036 belongs to it. The second highest compact kernel cannot annihilate that line through an archimedean mismatch.

For comparison, replacing the first target h\_1 by an arbitrary v would insert the factor h\_1\^{}vee(v) in (AR18). That factor can be zero. The theorem does not assert injectivity on all of V\_1 for one fixed highest second test, and does not automatically transfer the nonhighest sharp witnesses of K033--035 to this same test.

\subsection{5. Which factor appears in the actual global compact integration}\label{which-factor-appears-in-the-actual-global-compact-integration}

Here is an explicit coefficient calculation fixing that issue. Let U be any irreducible unitary compact representation of dimension d, let h be a unit vector, and let A be a linear functional on U\^{}vee and B a linear functional on U. The source-bilinear compact integral is

\[
 \boxed{
 \int_K A(\rho^\vee(g)h^\vee)B(\rho(g)h)\,dg
                     =\frac1d\sum_r A(e_r^\vee)B(e_r).}
 \tag{AR20}
\]

It is (AR12) evaluated by A tensor B. Unitarity identifies a coefficient of rho\^{}vee(g) with the conjugate of the corresponding coefficient of rho(g), so (AR20) also gives the usual matched Schur integral. The algebraic functional formulation above is appropriate to the bilinear theta convention. Both A and B are linear; no additional conjugation of the whole input theta lift is inserted.

In a classwise theta integral, A and B contain the finite class, the theta functional, and the left coefficient or multiplicity variables. They are not assigned the value one. Applying (AR20) to the first and then the second compact group gives exactly

\[
                         \frac1{d_0d_1}
 \tag{AR21}
\]

times the contracted coefficient matrices. It retains the finite class weights vol(K\_f)/\textbar Gamma\textbar{} from DO26.17. In the unnormalized diagonal model this is exactly the successive 1/d\_j coefficient in (AR12). Passing to I\_j divides the diagonal vectors by A\_j sqrt(d\_j); that changes the output model and produces (AR17). It must not be used to replace (AR21) in the actual global integral.

One can write the same fact in the class matrix language of DO026. Using unnormalized matched compact coefficient functions, each edge matrix has the compact factor 1/d\_j and its own finite class mass. Using the isometric coefficient functions sqrt(d\_j) times those functions moves the square roots into the coefficient matrices and the class metrics. The resulting global operator is unchanged. This is a change of coordinates, not an improvement of its cost.

\subsection{6. Actual two-edge Fock norm and the visible arithmetic denominators}\label{actual-two-edge-fock-norm-and-the-visible-arithmetic-denominators}

Write h=1 for K031 and h=4 for K032. The actual variety is J\^{}(14h) and the type has 42h embeddings. In the standard CM coordinates of K033(CS3--5,CS37), set

\[
 \gamma_\tau=\frac{\delta_\rho\,\pi_l^{-6}}{2G_i}>0,
                  \qquad\tau=(i,\varepsilon,l,\rho).
 \tag{AR22}
\]

For h=1, delta\_rho=1. For h=4 it is the exact positive product of the two quadratic trace weights in K032(QH20). All coordinate factors, including the standard-lattice pi\_l\^{}(-6), are retained. Put D\_(j,tau)=\textbar nu\_(j,tau)\textbar{} and let H\_lambda be the hook product. R011 and DO34.26, with K032's metric, give

\[
 \boxed{
 A_j^2=\det(2Y)^{-7hN_jN_{j+1}}
          \prod_\tau H_{\nu_{j,\tau}}
                         (\pi\gamma_\tau)^{D_{j,\tau}}.}
 \tag{AR23}
\]

Thus the exact local unitary-quotient multiplier (AR19) is

\[
 \boxed{
 \frac{\det(2Y)^{-7h(ab+bc)}}{d_0d_1}
    \prod_\tau H_{\nu_{0,\tau}}H_{\nu_{1,\tau}}
                 (\pi\gamma_\tau)^{D_{0,\tau}+D_{1,\tau}}.}
 \tag{AR24}
\]

Write C\_theta for the fixed period C=theta0*c in DO026, to distinguish it from the numerical weight constant below. The two analytic-to-geometric kernel identities additionally retain their product period C\_theta\^{}(14h(ab+bc)); it is not deleted by a compact projection or by unit-normalizing the real vectors. This is the product of the fixed periods of the actual two kernels, not a claim about the separately pinned terminal value.

For each place the exact coefficient dimension is the Weyl product

\[
 d_{j,\tau}=\prod_{1\le i<r\le N_j}
              \frac{\eta_{j,\tau,i}-\eta_{j,\tau,r}+r-i}{r-i},
                  \qquad d_j=\prod_\tau d_{j,\tau}.
 \tag{AR25}
\]

These are the actual rational denominators in (AR21). There is no claim that they are units at11. For the uniform explicit K031/K032 weights their values can be evaluated without storing the enormous dimensions. If C\_num is the numerical constant in HC6 or QH24, then eta\_(0,i)=(a-1-i)(2C\_num-1), in zero-based indexing. Consequently

\[
 d_{0,\tau}=(2C_{\mathrm{num}})^{a(a-1)/2},\qquad
 d_{1,\tau}=d_{0,\tau}
       \prod_{i=0}^{a-1}\prod_{r=a}^{b-1}
                         \frac{\nu_{0,i}+r-i}{r-i}.
 \tag{AR26}
\]

Legendre's formula and direct Weyl-product replay give:

\begin{longtable}[]{@{}llrrr@{}}
\toprule\noalign{}
Host & a,d,f & v11(d\_(0,tau)) & v11(d\_(1,tau)) & v11(d0 d1) \\
\midrule\noalign{}
\endhead
\bottomrule\noalign{}
\endlastfoot
K031, J14 & 42,1760,42 & 0 & 43 & 1806 \\
K032, J56 & 168,1634,168 & 0 & 86 & 14448 \\
\end{longtable}

Accordingly the compact Haar factor (AR21) has valuation -1806 or -14448 in these two exact examples. This is a displayed normalization denominator, not the valuation of the final global coefficient: the finite maps and numerator content still matter.

The hook factors are also retained. Their per-place valuations are

\begin{longtable}[]{@{}lrr@{}}
\toprule\noalign{}
Host & v11(H\_nu0) & v11(H\_nu1) \\
\midrule\noalign{}
\endhead
\bottomrule\noalign{}
\endlastfoot
K031 & 8344837905338 & 40893317909328 \\
K032 & 115754650130062 & 144968918676189 \\
\end{longtable}

These are the integer hook products in the real norm formula; the real numbers A\_j themselves are not thereby assigned an 11-adic valuation. R011's determinant-insertion metric ratios telescope as H\_eta2/H\_eta0. That identity does not erase the two actual Schwartz input norms (AR23), nor the Haar dimensions, periods, finite class masses or terminal branching map.

\subsection{6a. Exact conversion to the source's unit-normalized creation operator}\label{a.-exact-conversion-to-the-sources-unit-normalized-creation-operator}

The vectors in (AR5), and therefore (AR23)--(AR24), use the Chern creation operator in algebraic logarithm coordinates. They are not literally the source's \nolinkurl{pureTensorO} vectors. K035(PF22)--(PF24) identifies the latter operator at each CM place as \[
 r^{\mathrm{src}}_\tau=\lambda_\tau r_\tau,
 \qquad
 \lambda_\tau=-\frac{\Omega_\tau}{\sqrt{\pi q^{\mathrm{met}}_\tau}},
 \qquad |\lambda_\tau|^2\pi\gamma_\tau=1.
 \tag{AR27}
\] Here \(q^{\mathrm{met}}_\tau=-2i\tau(\xi_h)>0\), the square root is positive real, and \(\Omega_\tau\) is the fixed coordinate period of K035; its metric coefficient \(h_\tau\) is exactly \(\gamma_\tau\) in (AR22). These complex scalars are fixed before choosing the polynomial degree. Put \[
 \Lambda_j=\prod_\tau\lambda_\tau^{D_{j,\tau}}\ne0.
 \qquad
 \phi_j^{\mathrm{src}}=\Lambda_j\phi_j.
 \tag{AR28}
\] The equality follows because every monomial of the selected Schur polynomial has degree \(D_{j,\tau}\) at \(\tau\). It holds for the whole selected component, not only its highest vector. Thus all preceding nonvanishing and highest-line matching assertions hold for the literal source creation vector as well, with this explicit scalar and the unchanged finite input.

In the fixed quotient coordinates of (AR16)--(AR19), the two-edge value is multiplied by \(\Lambda_0\Lambda_1\), including its phase. In particular the exact squared multiplier for the source vectors is \[
 \frac{|\Lambda_0\Lambda_1|^2 A_0^2A_1^2}{d_0d_1}
 =\frac{\det(2Y)^{-7h(ab+bc)}}{d_0d_1}
          \prod_\tau H_{\nu_{0,\tau}}H_{\nu_{1,\tau}}.
 \tag{AR29}
\] If a new unitary quotient identification is chosen using the source vector, its phase changes accordingly; this does not remove the phase from an already fixed global comparison. The literal compact Haar integration factors remain \(1/(d_0d_1)\), not their square roots. For one edge, the corresponding coevaluation vector has norm squared \(d_j\det(2Y)^{-7hN_jN_{j+1}}\prod_\tau H_{\nu_{j,\tau}}\) times the actual finite norm and target-vector norm. This is exactly K037(CR35)--(CR36).

For the geometric ordinary-derivative functional in coordinates \(cD_\tau\), K035(PF27)--(PF29) instead gives the fixed comparison \[
 C_\theta^{14hN_jN_{j+1}}\,
 \left(\prod_\tau
        (\lambda_\tau/cD_\tau)^{D_{j,\tau}}\right)
 \mathcal J_{cD,F}(P_j),
 \tag{AR30}
\] with the actual finite splitting phase and normalized finite average inside \(F\), as in that theorem. The factors \(cD_\tau\) are algebraic and are units at11; the complex \(\lambda_\tau\) are not assigned finite valuations. Formula (AR30), the Gaussian norm, the hooks, and the class/Haar factors are all retained. In particular (AR29) is a norm conversion, not permission to delete the frame factor in (AR30) from arithmetic cost accounting.

\subsection{7. Scope and sources}\label{scope-and-sources-1}

The result closes the compact prior-image gate for the specified two highest kernels and the matched R017/DO036 input. The actual finite G, T and coupled Q prior-image calculations are LIFT22, LIFT24 and DO037 respectively; their scalars have their own normalizations and are not absorbed in (AR24). Global nonvanishing still requires the global theta quotient/multiplicity comparison. The actual integral prior-image lattice, all source pins, the terminal branching and MainX mass, and DP44 costs remain separate.

Carried checked results used, not reproved:

\begin{itemize}
\tightlist
\item
  DO025, equations DO25.3--12: actual common characters, compact source dual and target types, and exact selected partitions.
\item
  R011, equations R11.7--14: Cauchy/Fischer norm, determinant insertion, row normalization and actual vacuum norm.
\item
  DO026, equations DO26.4, DO26.7--8, DO26.16--20: actual creation kernel, covariance, compact Haar and class-mass conventions.
\item
  DO034, DO34.7--11 and DO34.26; K032, QH18--21; K033, CS37: actual standard CM polynomials and complete positive metric.
\item
  DO036, section6 and DO36.27: the fixed first primitive highest polynomial and R017 coefficient matched to its source line.
\end{itemize}

The pinned Lean source identities are \nolinkurl{Formalism/LocalToGlobal/Columns/BaseLemmas.lean496–552} (lamTrue\_old, lamTrue\_new, nuW\_eq\_true), \nolinkurl{Formalism/LocalRingCardinality/DataCoeff.lean21–23,51–55,83–84} (detShift, natW, primitive extremal, normFock), and \nolinkurl{Formalism/LocalToGlobal/HermitianForms/IsLocBasis.lean19–20} (nfAll). These were directly read. No Lean build is claimed.

Primary source for the actual compact oscillator convention: Konno, T., \& Konno, K. (2007). On doubling construction for real unitary dual pairs. Kyushu Journal of Mathematics,61(1),35--82. Section5, Fact5.1, equations5.4--5.5, and Lemmas5.2--5.3, printed pp.71--74, directly read. \url{https://doi.org/10.2206/kyushujm.61.35}

The finite certificate checks exact small compact Schur projections, primitive bideterminant Fischer norms, and both actual Weyl/hook valuation tables. It is a normalization diagnostic, not a substitute for the compact representation proof or the global theta comparison.

\clearpage

\section{AM-K036 --- one algebraic jet field for every lattice on a fixed even mixed-CM edge}\label{am-k036-one-algebraic-jet-field-for-every-lattice-on-a-fixed-even-mixed-cm-edge}

Author: GPT-6 Astra Pro, ALP\_MACHINE /root/kummer\_extension, 22 November 2026 UTC; inherited reasoning effort is not separately exposed.

This finding strengthens AM-K035's finite-family rationality to all selfdual lattices at a fixed positive even-rank edge. The two additional ingredients are an exact algebraic determinant scalar under rational unitary transport and a direct expansion of the source's induction germ. In that expansion the two square-root volume factors cancel. The remaining finite coefficient denominator is explicit. The conclusion is a mixed-coordinate replacement of the displayed Amb.JetsRat statement, not an instance of the unchanged, false common-quadratic constructor.

The actual ppav, theta line, geometric finite averages, source-normalized creation operators and gauges are those fixed in AM-K018, AM-K030/032--035 and AM-DO-022/026. The standard Weil representation and its rational theta automorphy are used with the same normalization; no arbitrary representation with only a nonzero lift is substituted.

\subsection{1. Fixed data and the precise statement}\label{fixed-data-and-the-precise-statement}

Let \(K=K_h\), \(h=1\) or4, be one of AM-K035's fields, acting on the actual ordinary product \(A_h=J^{14h}\). Put \(f=[K:\mathbb Q]/2=42h\). Retain its integral trace-one element \(w_0\), exact periods \(\Omega_\tau\), trace coefficient \(\xi_h\), metric \(H\), and algebraic coordinate generators \(cD_\tau\). Fix positive even ranks \(a<b\), a finite set S of finite primes of K with \[
S\cap\bar S=\varnothing,
\] and an integer \(D>0\). These are fixed parameters throughout.

For Hermitian selfdual \(\mathcal O_K\)-lattices \(N\subset K^a,\ N'\subset K^b\), set \[
L=\operatorname{Hom}_{\mathcal O_K}(N,N'),\qquad
\Lambda=\operatorname{diag}(\Omega_\tau)\Psi(L).
\tag{UJ1}
\] It is a principal integral symplectic lattice for \(E_H=\operatorname{Im}H\). The trace-one semicharacter, defined on the entire rational span, is \[
\alpha(t)=\exp\bigl(\pi iE_H(t,w_0t)\bigr).
\tag{UJ2}
\] On \(\Lambda\) it is sign-valued. Let \(F_{N,N'}\) be the actual geometric moved reduced theta germ, starting from the q-normalized product seed and using the two specified normalized Mumford averages. This is the same family as K035(PF28), not a new choice of generators on each line.

Replace the scalar embH/HV/alV wrappers in the source by their actual versions \(\operatorname{diag}(\Omega_\tau)\Psi,\ H,\alpha\). Keep the finite induction formulas of ThDistC and XiIndCClass unchanged. Specifically, take \[
e_{N,N'}=\chi_{N,N'}F_{N,N'},\qquad
\mathrm{chiF}_{N,N'}=\chi_{N,N'}\ne0,
\tag{UJ3}
\] with the same fixed character normalization in both entries. This is the exact pattern of the source AmbIn.amb and eNat definitions, not a new normalization invented to force rationality.

\subsubsection{Theorem UJ4}\label{theorem-uj4}

There exists one number field \(M\subset\mathbb C\), containing K, such that for \textbf{every} pair \(N,N'\) above, every \(c>0\) with \(c\mid D\), every \[
x\in D^{-1}\operatorname{homAt}_S(N,N'),
\] and every derivative multi-index m, \[
\boxed{\chi_{N,N'}^{-1}
 \operatorname{jetC}_{cD}
     \bigl(\mathrm{thInd}^{\rm mix}_{N,N',c,x}\bigr)(m)\in M.}
\tag{UJ4}
\] The field is independent of N, N′, x, c and m at these fixed \(K,a,b,S,D\). A finite collection of such edges and finite parameter choices can use their one finite compositum.

The proof also gives the exact finite induction coefficient denominator \[
\boxed{n_c=\prod_{P\in S}N(P)^{ab\,v_P(c)}.}
\tag{UJ5}
\] Before derivatives, the normalized induction germ is an \(n_c^{-1}\mathbb Z[\mu_{2D^2}]\)-linear combination of the finitely many reduced \(D\)-torsion translates of \(F_{N,N'}\). In particular this finite induction step contributes no 11-adic coefficient denominator when \(S\) excludes primes above11.

Neither the degree/discriminant of M nor integral denominators in its algebraic theta germs are bounded here. Consequently this existence result is not an estimate for the full DP44 cost.

\subsection{2. The exact rational-unitary transport scalar}\label{the-exact-rational-unitary-transport-scalar}

For \(g\in U_a(K)\), \(h\in U_b(K)\), define \[
A_{g,h}X=hXg^{-1},\qquad
\Delta(g,h)=\prod_{\tau\in\Psi}
       \tau(\det g)^{b/2}\tau(\det h)^{-a/2}.
\tag{UJ6}
\] All exponents are integers, so \(\Delta(g,h)\in K^\times\) since K is Galois. Its absolute value at the fixed complex embedding is1.

\subsubsection{Proposition UJ7}\label{proposition-uj7}

For the actual normalized family, \[
\boxed{F_{gN,hN'}(A_{g,h}u)
                     =\Delta(g,h)F_{N,N'}(u).}
\tag{UJ7}
\] The equality is of complete reduced theta functions, not just one value at zero.

\textbf{Proof.} First every selfdual lattice here is a finite-adelic unitary move of the standard one. At a split finite place this is the elementary identification of a unimodular pairing between two free modules with the identity pairing. At an unramified quadratic place, reduction of a unimodular Hermitian form gives a nondegenerate Hermitian form over the quadratic residue extension. It has a vector of nonzero norm: if all norms vanished, the identities for the norms of \(x+ty\), as t ranges over the residue field, would make the nondegenerate residue trace pairing vanish on every \(h(x,y)\), a contradiction. Lift that vector to obtain unit norm. The norm map on units of an unramified quadratic local extension is surjective, including residue characteristic2: its residue norm is surjective, and corrections on successive principal-unit quotients are supplied by the surjective residue trace. Completeness gives the desired unit. Scale the lifted vector to norm1, split its orthogonal complement, and induct. Thus the local selfdual Hermitian lattice is standard. These local isometries are integral at almost every place, since N is a full lattice commensurable with the standard one; they form a finite-adelic isometry. This argument uses finite-unramifiedness and does not assert global class number1.

Use the standard unitary Weil representation with trivial splitting characters, permitted because a,b are even. AM-DO-022's finite identity gives, first from the standard lattice and then by the representation law, \[
\varphi_{gN,hN'}=\omega_f(g,h)\varphi_{N,N'}.
\tag{UJ8}
\] Its hypotheses hold on K: the relative CM extension is finite-unramified on both hosts, and the actual trace-one integral lattice and its positive averages are the ones already fixed. Thus (UJ8) includes the required phase; it is not inferred from vacuum-line dimension.

The compact archimedean action on the vacuum has the exact character from DO-026(26.7), at infinity exponent zero: \[
\omega_\infty(g,h)\varphi_\infty
   =\beta(g,h)\varphi_\infty,\qquad
\beta(g,h)=\prod_\tau
       \tau(\det g)^{-b/2}\tau(\det h)^{a/2}
              =\Delta(g,h)^{-1}.
\tag{UJ9}
\] For any archimedean Heisenberg parameter w, rational theta automorphy and Heisenberg covariance give \[
\begin{split}
\Theta\bigl(\rho_\infty(w)\omega_f(g,h)\Phi\bigr)
 &=\Theta\bigl(\omega_\infty(g,h)^{-1}\rho_\infty(w)\Phi\bigr)\\
 &=\beta(g,h)^{-1}
       \Theta\bigl(\rho_\infty(A_{g,h}^{-1}w)\Phi\bigr),
\end{split}
\tag{UJ10}
\] where \(\Phi=\varphi_\infty\otimes\varphi_{N,N'}\). In the first line one inserts \(\omega_\infty(g,h)\omega_f(g,h)\), whose rational action fixes the theta functional. K035(PF25) identifies this entire-function calculation with the actual reduced reading, including its sign.

The map \(A_{g,h}\) preserves H, and it commutes with the period matrix \(\operatorname{diag}(\Omega_\tau)\), because each \(\Omega_\tau\) is scalar on its matrix block. The Gaussian factors therefore cancel on both sides of(UJ10). The same common period \(C^{14hab}\) occurs for both finite inputs, so it cancels too. This proves(UJ7).

If one instead uses the actual common stage character with infinity exponent \(-e_\tau\), its finite determinant-root factor for rational g,h is \(\prod_\tau\tau(\det g)^{-e_\tau}\tau(\det h)^{e_\tau}\), by the global character identity. Multiplication by the archimedean character \(\prod_\tau\tau(\det g)^{e_\tau-b/2}
\tau(\det h)^{-e_\tau+a/2}\) again gives \(\Delta^{-1}\). Thus the geometric scalar does not depend on this permitted splitting-character choice. ∎

The same isometry transports reduced translations exactly: \[
(U_{A_{g,h}t}F_{gN,hN'})(A_{g,h}u)
             =\Delta(g,h)(U_tF_{N,N'})(u).
\tag{UJ11}
\] This follows by inserting(UJ7) into the explicit multiplier \(U_tF(u)=e^{-\pi H(u,t)-\pi H(t,t)/2}F(u+t)\). No extra quadratic exponential or period appears. Both its scalar and its linear coordinate matrix are algebraic over K.

\subsection{3. Finitely many unitary classes at each fixed rank}\label{finitely-many-unitary-classes-at-each-fixed-rank}

Here is a full reduction of the required finiteness to the standard Minkowski second theorem. It avoids assuming that every selfdual lattice is globally standard.

For a Hermitian selfdual \(\mathcal O_K\)-lattice \(N\subset K^n\), the underlying positive integral symmetric form is \[
Q_N(x,y)=\operatorname{Tr}_{K/\mathbb Q}
                       \sum_jx_j\bar y_j.
\tag{UJ12}
\] Its trace-dual lattice is \(\mathfrak D_K^{-1}N\). Consequently its determinant is \[
\det Q_N=|\operatorname{disc}K|^n,
\qquad \operatorname{rank}_{\mathbb Z}N=n[K:\mathbb Q].
\tag{UJ13}
\] The dual assertion follows locally from Hermitian selfduality and then globally by intersection; the index is the norm of the different to the n-th power.

For clarity, positive integral Z-lattices of fixed rank d and fixed determinant \(\mathcal D\) have finitely many isometry classes. Indeed each nonzero vector has squared norm at least1. Let \(\lambda_1\le\cdots\le\lambda_d\) be the successive Euclidean minima. Minkowski's second theorem gives \[
\prod_i\lambda_i\le
        \frac{2^d\sqrt{\mathcal D}}{\operatorname{vol}(B_d)}.
\tag{UJ14}
\] Since every \(\lambda_i\ge1\), all minima are bounded by the right side. Choose independent vectors of bounded norm attaining them. Their integral Gram entries are bounded by Cauchy--Schwarz, so only finitely many Gram matrices occur. Their generated sublattice has bounded determinant and hence bounded index in the original lattice. An overlattice of index at most I lies in \((I!)^{-1}\) times that sublattice. There are finitely many such subgroups modulo the sublattice, proving the finiteness assertion. This argument does not incorrectly assume that successive-minimum vectors form a basis.

For each of the finitely many underlying trace lattices, multiplication by every fixed generator \(\eta\) of the finitely generated \(\mathbb Z\)-algebra \(\mathcal O_K\) has operator norm at most \[
\max_{\sigma:K\hookrightarrow\mathbb C}|\sigma(\eta)|
\] for \(Q_N\). In a fixed Z-basis there are therefore only finitely many integral matrices for each generator, hence finitely many possible \(\mathcal O_K\)-actions. An isometry of trace lattices commuting with this action extends to a K-linear isometry of the original Hermitian spaces: for every \(\eta\in K\), equality of the traces \(\operatorname{Tr}(\eta h(x,y))\) determines \(h(x,y)\) by the nondegeneracy of the field trace. It is thus an element of \(U_n(K)\).

It follows that the Hermitian selfdual lattices have finitely many \(U_n(K)\)-classes. Apply this separately at ranks a and b and choose representatives \[
N_1,\ldots,N_r,\qquad N'_1,\ldots,N'_s.
\tag{UJ15}
\] Every pair is \(gN_i,hN'_j\) for rational unitary g,h. Equations(UJ7) and(UJ11) are the exact scalar control needed to use this finite list.

Minkowski's inequality is classical, not a new claim here. A primary modern proof and exact statement are Henk(2002), Theorem1.3 and§3, pp.2 and5--7, cited below. All deductions from it in this section are given explicitly.

\subsection{4. The actual induction formula and volume cancellation}\label{the-actual-induction-formula-and-volume-cancellation}

This section uses the source's actual order of quotient representatives. The definition IsReps(\(\Lambda_1,\Lambda_2,S_0\)) means that \(S_0\) lies in \(\Lambda_2\) and represents \(\Lambda_2/(\Lambda_1\cap\Lambda_2)\), not the reversed quotient.

Let \(\mathcal A_c,\mathcal B_c\) denote the images in the complex period space of halfAO and halfBO. Their local rational lattices are \[
\begin{array}{c|ccc}
&P\in S& P\in\bar S& P\notin S\cup\bar S\\ \hline
\mathcal A_c&c^{-1}L_P&cL_P&L_P\\
\mathcal B_c&cL_P&c^{-1}L_P&L_P.
\end{array}
\tag{UJ16}
\] Thus \[
D\Lambda\subset\mathcal A_c,\mathcal B_c
                  \subset D^{-1}\Lambda\qquad(c\mid D).
\tag{UJ17}
\] The lattices \(\mathcal A_c,\mathcal B_c\) are selfdual for \(E_H\); the paired exponents at conjugate primes give this immediately from the selfduality of \(\Lambda\). Their restrictions of the integral trace-one element \(w_0\) are stable, so \(\alpha\) is their actual semicharacter.

Choose the source representatives \[
S_A\subset\Lambda\cap\mathcal A_c
\quad\text{for }(\Lambda\cap\mathcal A_c)/
                        (\mathcal B_c\cap\Lambda\cap\mathcal A_c),
\] \[
S_B\subset\mathcal B_c
\quad\text{for }\mathcal B_c/(\Lambda\cap\mathcal B_c).
\tag{UJ18}
\] Local calculation from(UJ16) gives \[
|S_A|=|S_B|=n_c
       =\prod_{P\in S}N(P)^{ab\,v_P(c)}.
\tag{UJ19}
\] For \(S_B\) the product initially runs over \(\bar S\); its equality follows because c is a rational integer, so conjugate primes have equal norms and valuations.

The finite representative sets in(UJ18) exist by the displayed local indices. The lift predicate is also inhabited for every x under consideration. In fact take the full fractional lattice \(\mathcal A_D\), whose localization at \(P\in S\) is \(D^{-1}L_P\). The Chinese remainder theorem for this finitely generated projective lattice supplies \(s_0\in\mathcal A_D\) with \(s_0-x\in cL_P\) for every \(P\in S\): this is reduction modulo \((cD)\mathcal A_D\) at the finitely many relevant primes. It satisfies the lift predicate with auxiliary integer \(c'=D\). Consequently the source's choice operator does choose an element satisfying that predicate, rather than an arbitrary value of an uninhabited choice.

Let \(s_0\) be that actual source liftW, allowing any of its permitted choices. Its two conditions are: \[
s_0\in\mathcal A_{c'}\text{ for some }c'>0,\qquad
s_0-x\in c\,\operatorname{homAt}_S(N,N').
\tag{UJ20}
\] Here and below x and \(s_0\) are sent to the period space by the fixed embedding when used in U or α. If \(x\in D^{-1}\operatorname{homAt}_S\), then at every \(P\in S\), \[
s_0\in D^{-1}L_P+cL_P\subset D^{-1}L_P.
\] Outside S, the first condition in(UJ20) makes \(s_0\) integral: at conjugate primes it is in \(c'L_P\), and elsewhere in \(L_P\). Therefore \[
\boxed{s_0\in D^{-1}\Lambda.}
\tag{UJ21}
\] This bound holds for every choice permitted by liftW, without a bound on \(c'\).

The literal definitions are \[
\begin{split}
\widehat\Phi(y)&=v^{-1/2}\frac v{|S_A|}
               \sum_{t\in S_A}\Phi(t)\operatorname{brk}(y,t)^{-1},\\
\xi_S(e)&=\sum_{y\in S_B}\widehat\Phi(y)\alpha(y)U_y e,\\
\xi_{\rm Ind}(e,s_0)&=\alpha(s_0)U_{-s_0}\xi_S(e),\\
\mathrm{thInd}(e)&=v^{-1/2}\xi_{\rm Ind}(e,s_0),
\end{split}
\tag{UJ22}
\] where \(v=\mathrm{volC}>0\) and \(\Phi=1_{\mathcal B_c}\). Positivity of v follows directly from its finite product of ratios of nonzero local lattice indices; the localized full lattices have finite-index intersection in both directions. The two occurrences of \(v^{-1/2}\), together with v, cancel exactly. Furthermore only one \(t\in S_A\) belongs to \(\mathcal B_c\): it is the representative of the zero class. For this t and every \(y\in S_B\subset\mathcal B_c\), selfduality makes \(E_H(t,y)\) integral, so \(\operatorname{brk}(y,t)=1\).

Using(UJ3) and linearity in the seed gives the exact simplified formula \[
\boxed{
\chi_{N,N'}^{-1}\mathrm{thInd}^{\rm mix}_{N,N',c,x}
 =\frac{\alpha(s_0)}{n_c}
       \sum_{y\in S_B}\alpha(y)U_{-s_0}U_yF_{N,N'}.
}
\tag{UJ23}
\] In particular no square root of a varying lattice-volume remains. This cancellation does not even require the parity hypothesis; parity is used for the rational determinant powers in(UJ6).

The source character cancellation in(UJ23) is literal. \nolinkurl{TrueColumnChoicesHeckeCharReferenceEqOne.lean143–152} sets e to eNat and chiF to chiFac; \nolinkurl{IsDiscriminantTwelfthRootOfIdealLatticeIffDelta1.lean140–141} defines eNat as chiFac times thMoved. For the actual characters these factors are nonzero. If a formally chosen factor were zero, both the seed and all these induction jets would be zero; that case would give algebraicity trivially rather than justify division by zero.

\subsection{5. A bounded torsion reduction, with the coefficient denominator retained}\label{a-bounded-torsion-reduction-with-the-coefficient-denominator-retained}

For \(s,t\in D^{-1}\Lambda\), integrality of \(E_H\) and of \(w_0\) gives \[
\alpha(s)\in\mu_{2D^2},\qquad
\exp(\pi iE_H(s,t))\in\mu_{2D^2}.
\tag{UJ24}
\] The reduced Heisenberg law \[
U_sU_t=e^{-\pi iE_H(s,t)}U_{s+t}
\] therefore changes a product of these translations only by a root in that fixed finite group. For \(\lambda\in\Lambda\), \[
U_{t+\lambda}F_{N,N'}
   =e^{\pi iE_H(t,\lambda)}\alpha(\lambda)\,
                              U_tF_{N,N'}.
\tag{UJ25}
\] This scalar lies in \(\mu_{2D}\). By(UJ17),(UJ21), every \(y-s_0\) in(UJ23) lies in \(D^{-1}\Lambda\). Reduction modulo \(\Lambda\) now proves \[
\boxed{
\chi^{-1}\mathrm{thInd}^{\rm mix}_{N,N',c,x}
\in n_c^{-1}\mathbb Z[\mu_{2D^2}]
       \operatorname{-span}
       \{U_tF_{N,N'}:t\in D^{-1}\Lambda/\Lambda\}.
}
\tag{UJ26}
\] The last set has exactly \(D^{2fab}\) torsion points, with arbitrary representatives allowed by(UJ25). Neither x nor an unbounded auxiliary lift denominator enlarges it. This proves the asserted explicit coefficient denominator and the 11-unit consequence in the theorem. It does not assert integrality of the individual reduced torsion germs.

\subsection{6. Completion of the uniform field proof}\label{completion-of-the-uniform-field-proof}

For each of the finitely many representative pairs in(UJ15), take every \(D\)-torsion translate in(UJ26). K035(PF30), using the actual algebraic CM gauges of K018, supplies one number field for the complete formal germs in this finite collection. Enlarge their finite compositum once to contain K, every \(cD_\tau\), and \(\mu_{2D^2}\); call the resulting field M.

For an arbitrary pair \(N=gN_i,N'=hN'_j\), the algebraic linear map \(A_{g,h}\) identifies \[
D^{-1}\Lambda_{i,j}/\Lambda_{i,j}
   \ \longrightarrow\ D^{-1}\Lambda_{N,N'}/\Lambda_{N,N'}.
\] Equations(UJ7),(UJ11) express every translated germ on the right as an algebraic scalar \(\Delta(g,h)\in K\) times a representative translated germ composed with \(A_{g,h}^{-1}\). Its matrix in the fixed \(\zeta_\tau\) coordinates has coefficients in K, because the period vector is scalar on each τ-block. Consequently every such complete series has coefficients in M. Equation(UJ26) then places every normalized induction germ in M{[}{[}u{]}{]}. Its derivatives of every order, and the weights \(cD_\tau^m\), remain in M. This proves(UJ4). ∎

The source Amb.JetsRat parameters a,b,S,D are fixed outside its quantifiers; it does not ask for one field for all ranks simultaneously. Thus the theorem has exactly its all-lattice/all-order quantifier pattern for any fixed even edge after replacing the impossible scalar/quadratic geometry by the explicit mixed geometry. If the prescribed coefficient field is smaller than M, a separate descent to that smaller field would still be necessary. The theorem permits choosing the coefficient field once; it does not assert that every preassigned field already contains the values.

\subsection{7. What this closes, and what remains}\label{what-this-closes-and-what-remains-2}

Together K035 and K036 now give an explicit varying-period Core identity, exact algebraic coordinate generators, and a single algebraic field for every admitted induction jet on each fixed even edge. The apparently arbitrary normalization chiFac cancels from the relevant arithmetic jet, and the source induction's volume factors cancel before differentiation. These are equations for the actual geometric kernel family and the literal finite induction operation.

The following are still separate:

\begin{itemize}
\tightlist
\item
  The false Fields.Kq/elliptic IsTen wrapper is not inhabited by this result.
\item
  The selected arithmetic row and subsequent image must satisfy the actual Kval and terminal comparison equations, rather than being chosen to make them true.
\item
  The field M is not given with a useful degree, discriminant, or height bound. Rational transport matrices and representative theta sections have not been bounded uniformly in the parameters.
\item
  The exact coefficient denominator \(n_c\) controls only the finite induction combination in(UJ26). It does not control reduced-gauge denominators, selected Schur saturation, all nonarchimedean row contents, or the full e/H or DP44 cost.
\end{itemize}

For the actual S that excludes11, the finite induction coefficient ring in(UJ26) is 11-integral. This is compatible with, and does not improve by itself, the sharp selected slope1/1220 in K033/K034. That slope concerns the complete reduced-gauge module and selected derivative operator; it cannot be converted into a stronger later-row bound solely by the present algebraicity theorem.

\subsection{Source locators and verification scope}\label{source-locators-and-verification-scope}

Direct source locators:

\begin{itemize}
\tightlist
\item
  ThDistC.lean72--106: Amb, IsLift, liftW, thInd and jetK;239--243: the fixed-rank Amb.JetsRat quantifiers.
\item
  XiIndCClass.lean17--34: brk, hatS, xiS and xiInd;171--197: halfAO/halfBO;216--237: period and half-period groups and representative order.
\item
  \nolinkurl{TorsionIntegrality/QuadraticSemicharacters/IsSemiChar.lean60–61:} actual IsReps orientation.
\item
  \nolinkurl{Columns/FiniteSumNonzeroLargeDepth.lean594–605:} homAt, divLat, mulLat;617: IsSdLat.
\item
  \nolinkurl{TrueColumnChoicesHeckeCharReferenceEqOne.lean143–152} and \nolinkurl{IsDiscriminantTwelfthRootOfIdealLatticeIffDelta1.lean140–141:} exactly matching eNat/chiFac normalization.
\item
  \nolinkurl{LocalNewformFunction.lean59–78:} actual global rational automorphy and Heisenberg covariance in IsWeilDatumO; the realized character calculation is DO-026(26.7).
\end{itemize}

Primary literature for the only additional standard theorem: Henk, M. (2002). \emph{Successive minima and lattice points}, Theorem1.3 and§3. \href{https://arxiv.org/pdf/math/0204158}{Author manuscript, arXiv:math/0204158}. The paper states the Minkowski bound on p.2 and gives its proof on pp.5--7. The finiteness application and all transport/induction reductions above are spelled out, rather than cited as an unspecified arithmetic compactness argument.

The companion certificate is a finite exact normalization and local-exponent check. It does not enumerate the enormous lattice-class set, construct M explicitly, prove an asymptotic cost bound, or execute Lean. Author and independent proof checks are recorded with their actual scope in the check-event TSVs.

\clearpage

\section{AM-R-019: exact normalization of the source terminal branching polynomial}\label{am-r-019-exact-normalization-of-the-source-terminal-branching-polynomial}

Author: /root, GPT-6 Astra Pro, ALP\_MACHINE, 22 November2026UTC. Reasoning effort is not externally exposed.

This finding computes the literal kapBr in the prescribed terminal operator. The classical Gelfand--Tsetlin norm theorem is an input; the normalization bridge to the source's product of minors is proved here. The resulting factorial product gives exact real and finite-place costs without enumerating the polynomial or its spectral projector. For the concrete m=2 hosts it reduces to at most (k+1)(k+2)/2 nontrivial factors, and the integer cBr has a finite factorial formula. It does not prove that the global preceding image has a nonzero terminal restriction or identify its integral lattice.

\subsection{1. Literal source normalization}\label{literal-source-normalization}

Let N=n+1 with n\textgreater=1, and let dominant integral weights lambda=(lambda\_1,\ldots, lambda\_N), mu=(mu\_1,\ldots,mu\_n) interlace:

\[
 a_i=\mu_i-\lambda_{i+1}\ge0,\qquad
 b_i=\lambda_i-\mu_i\ge0\quad(1\le i\le n).
 \tag{R19.1}
\]

For the source's SameShift condition assume mu\_n=lambda\_N. Subtract their common determinant shift from every entry when realizing their polynomial modules. All differences below are unchanged. Thus the proof may use lambda\_N=mu\_n=0.

Let X be an N-by-N matrix. Write Delta\_i for the determinant of its first i rows and columns, and Z\_i for the determinant with rows1,\ldots,i−1,N and columns1,\ldots,i, in that order. The actual source vectors and number are

\[
 p_\lambda=\prod_{i=1}^{n}\Delta_i^{\lambda_i-\lambda_{i+1}},
 \qquad \xi_{\lambda,\mu}=\prod_{i=1}^{n}\Delta_i^{a_i}Z_i^{b_i},
 \qquad \kappa(\lambda,\mu)
    =\frac{\|\xi_{\lambda,\mu}\|_F^2}{\|p_\lambda\|_F^2}.
 \tag{R19.2}
\]

Here monomials are orthogonal and \textbar\textbar X\^{}alpha\textbar\textbar{}\emph{F\^{}2=alpha!. These are exactly xiBr and kapBr in \nolinkurl{Columns/FiniteSumNonzeroLargeDepth.lean119–133,} with aBr/bBr41--44. The accepted source interlacing theorem says that xi is the nonzero GL\_n-highest vector of weight mu in the lambda representation and coefLast(xi)=p\_mu under SameShift. For the intermediate lowering steps we also need its elementary polynomial realization without SameShift. Every Delta\_i and Z\_i is killed by E\_pq for p\textless q\textless=n: either row q is absent, or row p is already present, producing a repeated row. The row weights of xi telescope to mu\_1,\ldots,mu\_n and sum(lambda)−sum(mu). Its column raising operators also kill it, and its column weight is lambda because a\_i+b\_i=lambda\_i− lambda}(i+1). The Cauchy decomposition therefore places it in the row representation of highest weight lambda for every interlacing mu. Evaluation at the matrix X\_* of Section3 gives one, proving it is nonzero. This supplies the intermediate realization directly; it does not extend the source SameShift interface silently.

\subsection{2. Exact factorial formula}\label{exact-factorial-formula}

\subsubsection{Theorem R19.3}\label{theorem-r19.3}

\[
\boxed{
\kappa(\lambda,\mu)=
\prod_{1\le i\le j<N}
 \frac{(\lambda_i-\mu_j+j-i)!}{(\mu_i-\mu_j+j-i)!}
\prod_{1\le i<j\le N}
 \frac{(\mu_i-\lambda_j+j-i-1)!}
      {(\lambda_i-\lambda_j+j-i-1)!}.}
\tag{R19.3}
\]

Every factorial argument is a nonnegative integer. A more useful form, with x\_ij=mu\_i−mu\_j+j−i, is

\[
\boxed{
\kappa(\lambda,\mu)=
 \prod_{1\le i\le j\le n}
 \frac{(x_{ij}+b_i)!(x_{ij}+a_j)!}
      {x_{ij}!(x_{ij}+a_j+b_i)!}.}
\tag{R19.4}
\]

In particular 0\textless kappa\textless=1. A factor is one if a\_j=0 or b\_i=0. Formula(R19.4), with Legendre's factorial valuation, computes v\_p(kappa) exactly; no cancellation or integrality guess is needed.

\textbf{Proof of the last claims from(R19.3).} In the second product put j=r+1 and group it with the term (i,r) in the first. The four arguments become x+b\_i, x+a\_r, x, x+a\_r+b\_i. The ratio is \(\prod_{s=1}^{a_r}(x+s)/(x+b_i+s)\), hence is positive and at most one. Interlacing and dominance give x\textgreater=0 and all remaining arguments nonnegative. It remains to prove the precise normalization in(R19.3).

\subsection{3. Lowering-operator normalization: a complete minor calculation}\label{lowering-operator-normalization-a-complete-minor-calculation}

Use the row operators

\[
 E_{pq}=\sum_{c=1}^{N}X_{pc}\frac{\partial}{\partial X_{qc}},
 \qquad h_i=E_{ii}-i+1.
\]

They satisfy E\_pq\^{}*=E\_qp for the Fock metric. Molev's lowering operator, equation(2.10), is

\[
 z_{Ni}=\sum_{i<i_1<\cdots<i_s<N}
 E_{i_1i}E_{i_2i_1}\cdots E_{Ni_s}
       \prod_{j\notin\{i_1,\ldots,i_s\}\atop i<j<N}(h_i-h_j).
 \tag{R19.5}
\]

The empty path term is E\_Ni times all displayed diagonal factors. Molev's Lemma2.5 and multiplicity-free branching put its image of xi on the mu−epsilon\_i highest line whenever a\_i\textgreater0. We claim

\[
\boxed{
 z_{Ni}\xi_{\lambda,\mu}
   =\left(\prod_{j=i+1}^{N}
       (\mu_i-\lambda_j+j-i-1)\right)
        \xi_{\lambda,\mu-\epsilon_i}.}
\tag{R19.6}
\]

To determine the scalar, evaluate at X\_* whose first n rows are the corresponding coordinate rows and whose last row has every entry one. Each Delta\_l and Z\_l for l\textless=n is one, so both relevant xi vectors evaluate to one. Fix a path p\_0=i\textless p\_1\textless\ldots\textless p\_s\textless N. The corresponding differential product in(R19.5) is the coefficient of t\_0\ldots t\_s in xi evaluated after the simultaneous row replacements

\[
 X_{p_r,*}\mapsto e_{p_r}^*+t_r e_{p_{r+1}}^*\quad(r<s),
 \qquad
 X_{p_s,*}\mapsto e_{p_s}^*+t_s(1,\ldots,1).
 \tag{R19.7}
\]

Indeed iterated differentiation represents the operator product by left multiplication with the elementary row-update matrices in reverse order. Their pairwise cross-products in this order vanish, so the simultaneous replacements are exact. Every differentiation is first order in its own variable; coefficient and derivative agree.

Put B\_0=1 and B\_r=1−t\_(r−1)B\_(r−1). Direct triangular determinant expansion of(R19.7) gives

\[
 \Delta_l=\begin{cases}1,&l<p_s,\\1+t_sB_s,&l\ge p_s,\end{cases}
 \qquad
 Z_{p_r}=B_r\ (1\le r\le s),\qquad Z_l=1\ (l\notin\{p_1,\ldots,p_s\}).
 \tag{R19.8}
\]

For example the upper triangular path matrix has inverse-column entries 1,−t\_(s−1),t\_(s−2)t\_(s−1),\ldots, which sum to B\_s; the rank-one determinant formula gives the first identity. Replacing row p\_r by the last row gives the same inverse-column sum truncated at r, proving the second. Rows not on the path do not connect to the missing column, so their determinants remain one.

Set A\_j=sum\_(l=j)\^{}n a\_l. Substitution in xi is therefore

\[
 (1+t_sB_s)^{A_{p_s}}\prod_{r=1}^sB_r^{b_{p_r}}.
\]

First extract t\_s; it contributes A\_(p\_s) and one extra factor B\_s. Then extract t\_(s−1),\ldots,t\_0 successively using B\_r=1−t\_(r−1)B\_(r−1). The result is exactly

\[
 (-1)^s A_{p_s}\prod_{r=1}^s(b_{p_r}+1).
 \tag{R19.9}
\]

This also covers the empty path, with value A\_i. The rightmost diagonal factors in(R19.5) act first on xi, with eigenvalues M\_j=mu\_i−mu\_j+j−i. Write c\_j=b\_j+1 and M\_i=0. Summing paths according to their largest vertex, and then using A\_l=sum\_(r\textgreater=l)a\_r, gives the scalar

\[
\begin{aligned}
S_i&=A_i\prod_{j>i}M_j
 -\sum_{l>i}c_lA_l
        \prod_{i<j<l}(M_j-c_j)\prod_{j>l}M_j\\
&=\sum_{r=i}^{n}a_r
        \prod_{i<j\le r}(M_j-c_j)\prod_{j>r}M_j.
\end{aligned}
\tag{R19.10}
\]

Both equalities are finite product expansions; the second is the telescoping identity for replacing each M\_j by M\_j−c\_j in order. Now M\_j−c\_j=M\_(j−1)+a\_(j−1). Thus the last expression is

\[
 \prod_{r=i}^{n}(M_r+a_r)-\prod_{r=i}^{n}M_r
 =a_i\prod_{r=i+1}^{n}(\mu_i-\lambda_{r+1}+r-i),
\]

since M\_i=0. This is the coefficient in(R19.6). If a\_i=0, branching makes the image zero and the same scalar calculation has value zero. This completes the normalization proof.

\subsection{4. Application of the classical norm theorem}\label{application-of-the-classical-norm-theorem}

Starting at p\_lambda, apply the lowering operators in Molev's Lemma2.6/Theorem2.7 order, so the largest i acts first. Each intermediate weight interlaces; at coordinate i, decrease from lambda\_i to mu\_i. Formula(R19.6) depends on that coordinate and lambda, independently of the previously decreased coordinates. The total scalar is therefore

\[
 C_{\lambda,\mu}
 =\prod_{1\le i<j\le N}
       \frac{(\lambda_i-\lambda_j+j-i-1)!}
            {(\mu_i-\lambda_j+j-i-1)!},
 \qquad \xi_{\mu}^{\mathrm{GT}}=C_{\lambda,\mu}\xi_{\lambda,\mu}.
 \tag{R19.11}
\]

The Fock metric divided by \textbar\textbar p\_lambda\textbar\textbar\_F\^{}2 is the contravariant metric with unit highest vector used in Molev's Proposition2.4. Choose the lower rows of the GT pattern to be the initial segments of mu. The factors in that proposition from levels below N are one. Its top-level factor is

\[
 \|\xi_{\mu}^{\mathrm{GT}}\|^2
  =\left[\prod_{1\le i\le j<N}
       \frac{(\lambda_i-\mu_j+j-i)!}{(\mu_i-\mu_j+j-i)!}\right]
       C_{\lambda,\mu}.
\]

Dividing by C\_(lambda,mu)\^{}2 proves(R19.3). This invokes the published norm theorem; it does not reprove it or identify a basis by its name alone. The direct minor calculation supplies the needed normalization to the actual source vector. ∎

\subsection{5. The exact integer spectral factor}\label{the-exact-integer-spectral-factor}

Set d0=sum\_i a\_i, L\_i=lambda\_(i+1), U\_i=min(lambda\_i,lambda\_(i+1)+d0), and R\_ij=mu\_j+i−j. Then L\_i\textless=mu\_i\textless=U\_i. The literal source integer is

\[
 c_{\rm Br}=\prod_{i=1}^{n}\prod_{v=L_i\atop v\ne\mu_i}^{U_i}
                          \prod_{j=1}^{n}(v-R_{ij}).
 \tag{R19.12}
\]

For j\textless i, R\_ij\textgreater U\_i; for j\textgreater i, R\_ij\textless L\_i, by interlacing. Consequently it is a nonzero integer, and its exact signed factorial evaluation is

\[
 c_{\rm Br}=\prod_i\left[(-1)^{\mu_i-L_i}(\mu_i-L_i)!(U_i-\mu_i)!
   \prod_{j<i}\frac{(-1)^{U_i-L_i}(R_{ij}-L_i)!}
          {(R_{ij}-U_i-1)!(R_{ij}-\mu_i)}
   \prod_{j>i}\frac{(U_i-R_{ij})!}
          {(L_i-R_{ij}-1)!(\mu_i-R_{ij})}\right].
 \tag{R19.13}
\]

\textbf{Proof.} The j=i product has its single zero omitted and is the displayed two factorials with mu\_i−L\_i negative factors. For j\textless i, all factors are negative and one is omitted, leaving U\_i−L\_i of them; the full consecutive product of their absolute values is divided by the omitted R\_ij−mu\_i. For j\textgreater i all factors are positive and the same argument gives the last expression. Their product is exactly(R19.12). If L\_i=U\_i=mu\_i, every factor is one. No enormous interval needs to be enumerated. ∎

\subsection{6. The actual terminal operator norm and its scalar valuations}\label{the-actual-terminal-operator-norm-and-its-scalar-valuations}

The accepted source identities InterlaceAll.IoHw, CoefXi, BrTranspose and BeOnEq give

\[
 \operatorname{brStd}(\xi_{\lambda,\mu})
       =c_{\rm Br}\kappa(\lambda,\mu)p_\mu,
 \qquad \|\operatorname{brStd}\|=|c_{\rm Br}|\sqrt\kappa.
 \tag{R19.14}
\]

For the tensor over CM embeddings and the literal t\_tau=tau(star piT), write nBr\_tau=sum lambda\_tau−sum mu\_tau. Then

\[
\boxed{
\|\operatorname{brAt}\|
 =\prod_\tau |t_\tau|^{-nBr_\tau}
                  |c_{{\rm Br},\tau}|\sqrt{\kappa_\tau}.}
\tag{R19.15}
\]

For R017's actual norm element, every \textbar t\_tau\textbar=sqrt(T). Thus its logarithm is exactly

\[
 -\tfrac12\log T\sum_\tau nBr_\tau
   +\sum_\tau\log|c_{{\rm Br},\tau}|
   +\tfrac12\sum_\tau\log\kappa_\tau.
 \tag{R19.16}
\]

At11, t\_tau is a unit. The scalar on the primitive xi-to-p\_mu line therefore has valuation

\[
 v_{11}(c_{{\rm Br},\tau})+v_{11}(\kappa_\tau),
 \tag{R19.17}
\]

exactly calculable from(R19.4),(R19.13) and v\_p(r!)=sum\_(s\textgreater=1)floor(r/p\^{}s). This is a scalar valuation on that specified line; it is not the Smith content of the full projected integral row lattice. The latter also depends on the actual inclusion matrix and is retained as an open obligation.

\subsection{7. Concrete selected m=2 weights on J14 and J56}\label{concrete-selected-m2-weights-on-j14-and-j56}

Use K031/K032 and R018. Let k be42 or168, d be1760 or1634, C be45156049293 or40530339687, and N=k+2d. The terminal column and row weights, before removing their common determinant shift, are

\[
\begin{aligned}
 \lambda_i&=(2(k-i)+3)C+i-k/2-d &&(1\le i\le k),\\
 \lambda_{k+1},\ldots,\lambda_{k+d}&=k/2-d/2,\\
 \lambda_{k+d+1},\ldots,\lambda_N&=-4C+(k+d)/2;\\
 \mu_i&=2(k-i+1)C+i-k/2-d &&(1\le i<k),\\
 \mu_k&=2C+k/2-d,\\
 \mu_{k+1},\ldots,\mu_{k+d-1}&=k/2-d/2,\\
 \mu_{k+d}&=-2C+k/2,\\
 \mu_{k+d+1},\ldots,\mu_{N-1}&=-4C+(k+d)/2.
\end{aligned}
\tag{R19.18}
\]

These are lamTrue\_old/new and sgTrue\_eq, with RowPar.sgHC, not newly chosen interlacing weights. In particular the common shift is −4C+(k+d)/2. Their nonzero a,b entries are

\[
\begin{array}{c|cc}
 i&a_i&b_i\\ \hline
 1,\ldots,k-1&C-1&C\\
 k&2C-d/2&C\\
 k+d&2C-d/2&2C-d/2.
\end{array}
\tag{R19.19}
\]

All others vanish. Hence the nontrivial factors of(R19.4) number at most(k+1)(k+2)/2, and

\[
 d0=(k+3)C-(k-1)-d,\qquad
 nBr=(k-2)C+k/2.
 \tag{R19.20}
\]

Increasing the first u entry by an integer A\textgreater=0 changes only lambda\_1, adding A; the row mu, every a\_i and d0 stay fixed, while b\_1 and nBr increase by A. Put A\_\emph{=max(0,d0−(lambda\_1−lambda\_2)). Then cBr is independent of A for A\textgreater=A\_}: only U\_1 varies, and it has reached lambda\_2+d0 by that point. Moreover

\[
\boxed{
 \frac{\kappa(\lambda+A\epsilon_1,\mu)}{\kappa(\lambda,\mu)}
  =\prod_{j=1}^{n}
    \frac{(\lambda_1-\mu_j+j)_{a_j}}
         {(\lambda_1+A-\mu_j+j)_{a_j}},}
\tag{R19.21}
\]

where(x)\_r=x(x+1)\ldots(x+r−1), with(x)\_0=1. This follows by grouping the i=1 factors in(R19.4); every other factor is unchanged. In particular kappa is a positive rational function of A with exact decay degree d0. The real term in (R19.16) is therefore explicit even on K033's growing family: after the cBr threshold its extra logarithm is −(A/2)log T−(d0/2)log A+O(1) at the growing embedding. The constant depends only on the fixed host and fixed row. The source size restriction from R018 still limits A at a fixed T; this asymptotic formula is not an assertion that fixed-T weights remain admissible as A tends to infinity.

\subsubsection{Exact11-adic values on the actual weight family}\label{exact11-adic-values-on-the-actual-weight-family}

The saved certificate evaluates the following integers by Legendre sums and checks the actual row weights against the independent literal RowPar.sgHC formula:

\begin{longtable}[]{@{}lrr@{}}
\toprule\noalign{}
Host & d0 & A\_* \\
\midrule\noalign{}
\endhead
\bottomrule\noalign{}
\endlastfoot
J14 & 2,032,022,216,384 & 1,941,710,117,799 \\
J56 & 6,930,688,084,676 & 6,849,627,405,303 \\
\end{longtable}

\begin{longtable}[]{@{}lrrr@{}}
\toprule\noalign{}
Host & v11(kappa) & v11(cBr), baseline & v11(cBr), saturated \\
\midrule\noalign{}
\endhead
\bottomrule\noalign{}
\endlastfoot
J14 & −118 & 1,431,126,153,655,645 & 2,122,569,126,603,571 \\
J56 & −3423 & 4,747,460,542,799,389 & 7,100,307,556,522,475 \\
\end{longtable}

These are per embedding. The kappa valuation is unchanged after increasing lambda\_1 by A=11\^{}r or121\emph{11\^{}r for any r\textgreater=13. Indeed every positive integer factor in the numerator rising factorials of(R19.21) is at most S\_2−2, where S\_2 is the exact terminal spread in R018. In both hosts S\_2−2\textless11\^{}13. Adding 11\^{}r or121}11\^{}r therefore preserves the valuation of each factor. Both A choices also exceed A\_* at r\textgreater=13, so the displayed cBr valuation is the saturated one. This is an all-r proof, not an extrapolation from the two sampled endpoints in the certificate.

For the actual family with two growing CM embeddings and all other embeddings fixed, the valuation of the full highest-line branching scalar in(R19.17) is consequently

\begin{longtable}[]{@{}lr@{}}
\toprule\noalign{}
Host & v11 of the complete specified highest-line scalar, r\textgreater=13 \\
\midrule\noalign{}
\endhead
\bottomrule\noalign{}
\endlastfoot
J14 & 61,490,184,399,427,986 \\
J56 & 802,279,065,217,168,460 \\
\end{longtable}

At the two growing embeddings use the saturated cBr value; at the other f−2 embeddings use its baseline value. Add all kappa valuations. Since piT is an11-unit there is no extra T contribution at this prime. The sign of cBr is negative in both the baseline and saturated per-place cases, as checked by(R19.13). These large exact numbers are scalar normalizations of the prescribed branch on its primitive highest line. They do not give the valuation of an unknown global terminal vector or the full integral projector matrix. R018's fixed-T admissible ranges still apply when these formulas are used in that numerical family.

\subsection{8. Costs, framework, sources and check scope}\label{costs-framework-sources-and-check-scope}

This removes the unevaluated branching-polynomial norm from the terminal norm budget. It also replaces the spectral interval product by an exact finite factorial formula. It does not replace the actual global terminal vector by xi, identify a detecting row at zero cost, or assign this factor twice in DP44: the source's separate T-term already carries its own prescribed period power. R002's content/dual-norm formulas must still be applied to the actual image and row lattice.

The polynomial representation and its restriction live in the integral representation/flag setting of SGA3 I4.7.3/4.7.3.1 and SGA6 VI3.3 recorded in the dictionary. The factorial scalar and Fock norm refine this representation class to a metrized map. The polarized theta line is still the Picard/biextension object of FGA232 Theorem3.1 and SGA7 I VII2.7/2.9, VIII7.4(b); those representability results do not determine the present metric normalization. The formal/crystalline slopes of cote4 and cote54 scans5--7 are unchanged. This is a DP\_MULT Q3e/THETA3 G1 normalization result, with no consequence asserting Kummer RH.

Exact stipulated source locators, commit 7cbf7a0c32fca1badf05931879a37c3f8a2b03b4:

\begin{itemize}
\tightlist
\item
  \nolinkurl{Columns/FiniteSumNonzeroLargeDepth.lean41–55,119–151:} aBr,bBr,rSet,cBr,xiBr,kapBr,brStd,brAt.
\item
  \nolinkurl{KernelLattices/InterlaceAll.lean54–71} and125--147: BrTranspose,IoHw,CoefXi,beOnVal,BeOnEq. Their validity is accepted as instructed; no Lean build is performed.
\item
  \nolinkurl{Columns/AbelianHostFieldConditions.lean408–421,464–472:} RowPar.sgHC,sgTrue,lamTrue; \nolinkurl{Columns/BaseLemmas.lean496–530,} 575--582: the actual terminal weight formulas.
\end{itemize}

Molev, A. I. (2002). \emph{Gelfand--Tsetlin bases for classical Lie algebras} (arXiv:math/0211289v2). \url{https://arxiv.org/abs/math/0211289} . Proposition2.4, equation(2.10), Lemmas2.5--2.6 and Theorem2.7, printed pp10--11 (PDF pp11--12). These exact primary loci were read in this run. The norm theorem is cited; the source-specific minor normalization is proved in Sections3--4 above.

Finite exact diagnostics in branch\_norm\_diagnostic.py compare 250 complete Fock norms and251 lowering actions for ranks2--5; the combined path replay checks2550 paths. They support the full proof and are not substituted for it. The companion arithmetic certificate records exact actual-weight valuations without expanding huge factorials. Whole-proof reviews are logged separately after the entire finished proof is read.

\clearpage

\section{ALP-LIFT-26: the full prescribed sharp two-edge coupled lift is nonzero}\label{alp-lift-26-the-full-prescribed-sharp-two-edge-coupled-lift-is-nonzero}

Author: GPT-6 Astra Pro, ALP\_MACHINE /root/lift\_accounting. 22 November 2026 UTC. Inherited reasoning effort is not exposed.

This result assembles the checked actual local quotient maps into the fixed global m=2 coupled operator on either repaired mixed CM host. The initial automorphic function is the actual R017 function. The finite kernels are the unchanged sharp G/T/Q kernels, with maximal vacua at the remaining places. The real kernels are the two primitive selected highest polynomials of LIFT25. The full prescribed Y quotient is summed. No later Schwartz replacement or Hecke translation is made.

The distinction between this statement and unrestricted theta nonvanishing is essential. Rallis is first used to prove that the global map on every automorphic multiplicity copy is nonzero. Local Howe uniqueness then identifies the kernel of that global map with the product of the actual local small-theta quotient kernels. The explicit local calculations detect the particular coupled tensor. We do not infer fixed-vector survival from GQT Proposition11.6 alone.

The resulting nonzero analytic automorphic vector is not yet the source's arithmetic row or terminal period. Its global Rallis and class-mass normalization, integral content, and projection to the prescribed terminal branching component remain explicit obligations.

\subsection{1. Actual data and precise statement}\label{actual-data-and-precise-statement}

Let E/F be either K031's degree84 CM host with its principally polarized J14, or K032's degree336 CM host with its principally polarized J56. Write

\[
 (a,d,b,c)=(42,1760,1802,3562)
       \quad\hbox{or}\quad(168,1634,1802,3436),
 \qquad G_j=U(I_{a+jd}),\quad j=0,1,2.
 \tag{GS1}
\]

All three Hermitian spaces are positive at every real embedding. The relative CM extension is unramified at finite places; its dyadic places split. The actual selfdual trace lattices, additive character, vacuum, creation operators and splitting are those of K030/032 and DO022/034, as used in DO036. These geometric and character results are carried inputs, not reasserted constructions.

Denote the actual common splitting character on edge j by chi\_j. In the source notation it is mu\_j times the fixed tuple character xi. Both members of that edge have this same character. In particular

\[
 \upsilon=\chi_0\chi_1^{-1}=\mu_0\mu_1^{-1},\qquad
 \chi_j|_{\mathbb A_F^\times}=1.
 \tag{GS2}
\]

The fixed tuple character cancels in the ratio. At the two rational G primes and all their selected places, the two residue characters are distinct and nontrivial. At T both chi\_j are unramified. At Q they have the same specified nontrivial residue. These are the literal conditions checked in LIFT18/19, K031/032, DO035, and LIFT24. At infinity the exponents are -e\_j, with e\_1-e\_0=4C\textgreater0 for the concrete uniform data. Hence upsilon and upsilon squared have nontrivial angular infinity type.

Let F\_0 be R017's actual nonzero class-supported automorphic function on G\_0. It has the matched selected highest real line, finite Steinberg generic type at G, prescribed nontrivial projective-line type at T, and maximal compact invariance at Q and all other finite places. Its exact norm and source level are the ones already computed in R017 and DO036; no new source input is substituted.

At infinity use phi\_(j,infty) from LIFT25(AR5) on both edges.

These are the Chern-coordinate creation vectors. For the literal source \nolinkurl{pureTensorO} vectors one must multiply edge j by the fixed nonzero scalar \(\Lambda_j=\prod_\tau\lambda_\tau^{D_{j,\tau}}\) of LIFT25(AR27)--(AR28), as identified in K035(PF22)--(PF23). All statements below about the Chern-normalized operator therefore hold for the literal source operator after multiplication by \(\Lambda_0\Lambda_1\). Its squared output norm gains \(|\Lambda_0\Lambda_1|^2\); the exact source coevaluation norms are K037(CR35), or LIFT25(AR29). The coefficient comparison retains \(C_\theta^{14h(ab+bc)}\prod_\tau(\lambda_\tau/cD_\tau)
 ^{D_{0,\tau}+D_{1,\tau}}\), as in LIFT25(AR30). This conversion does not change the finite kernels or prior-image nonvanishing, and its complex period factors are not assigned finite valuations. At finite places use the literal sharp XiCpl1 and XiCpl2 factors, including ftY's exact power q\^{}(-M b), all character phases, and the actual piT in the terminal row. Let the complete coupled tensor be

\[
 \Psi=\sum_{Y\in\mathcal Y}\phi_{1,Y}\otimes\phi_{0,Y},
 \tag{GS3}
\]

where the sum is expressed in standard local frames. Section2 proves that this is exactly the literal moving-lattice Y sum, not a new choice of representatives. Define the actual bilinear coupled lift

\[
 \mathcal B F_0
   =\sum_{Y\in\mathcal Y}
       \Theta_{1,\phi_{1,Y}}^{\rm bil}
          \bigl(\Theta_{0,\phi_{0,Y}}^{\rm bil}F_0\bigr),
 \qquad
 \Theta_{j,\phi}^{\rm bil}F(h)
    =\int_{[G_j]}\theta_{j,\phi}(g,h)F(g)\,dg.
 \tag{GS4}
\]

\textbf{Theorem.} For either host and its stated actual data,

\[
                         \mathcal B F_0\ne0.
 \tag{GS5}
\]

The result holds on every nonzero vector in the same finite-level, matched-highest input subspace whose G representation lies in the finite Steinberg isotypic subspace. It includes all automorphic multiplicity copies and their finite linear combinations.

This is a theorem about the complete fixed coupled operator. In particular at least one admitted matched Y pair has a nonzero double lift. It does not say that DO036's previously selected single Y\_0 must survive the second edge: at Q the proof uses the full Y sum. The nonzero function on G\_2 does not by itself give a nonzero literal brC/tilQ row period. The source ChKer.cp2 subsequently evaluates at tilQ(d), and BtOf then applies PtilFF. The Q graph used in(GS22) is not in that embedded terminal subgroup; it cannot establish that restricted endpoint detection. This remaining terminal gate is distinct from(GS5).

\subsection{2. Exact factorization of the literal Y family}\label{exact-factorization-of-the-literal-y-family}

The source definitions are \nolinkurl{Columns/FiniteSumNonzeroLargeDepth.lean}1120--1169: locLat, ybSet, repY, XiCpl1, PhiTop, ftY, and XiCpl2. The Y support and period modules are intersections of the prescribed local modules. DO032 proves their CRT identification

\[
 \mathcal Y\simeq
 \prod_{P\in G}(q_P^{-1}O_P^b/O_P^b)
 \times\prod_{P\in Q}(q_P^{-M}O_P^b/q_P^M O_P^b).
 \tag{GS6}
\]

For completeness, localization sends a global coset injectively to its local cosets because the global period module is precisely the intersection of those local conditions. Clearing the finitely many denominators and applying CRT gives surjectivity. The prescribed middle frame identifies its localized lattice with O\_P\^{}b. This argument applies to every middle lattice, not just the identity class.

There is no archimedean chamber condition on repY. A change of repY inside its coset changes y\_P by O\_P\^{}b at G and q\_P\^{}M O\_P\^{}b at Q. The first raw factor is periodic for exactly those changes; on the terminal support, ftY's change has integral additive-character argument. LocCst.Ok.psi\_int therefore makes it one. These periodicity checks are DO031 Section2 and DO032 Section1, with the first sharp G check in LIFT22(CG7--9) and the full Q check in LIFT23/DO037. No phase depending on the global representative remains.

At a middle adelic frame h, use that frame in both raw factors. The global coset Y maps by (GS6) to the local classes of h\_P\^{}\{-1\}y. Those classes run through the entire fixed standard quotient exactly once. Therefore the finite sum of the products of the two raw local functions is exactly

\[
 \Psi_f=\bigotimes_{v<\infty}\Psi_v,\qquad
 \Psi_v=
 \begin{cases}
  \displaystyle\sum_{y\in\mathcal Y_v}
                 \phi_{1,v,y}\otimes\phi_{0,v,y},&v\in G\cup Q,\\
  \phi_{1,v}\otimes\phi_{0,v},&v\notin G\cup Q.
 \end{cases}
 \tag{GS7}
\]

All factors in the source product use the selected place over v; the opposite split component is already encoded in the honest unitary Weil representation. It is not an additional independent Y variable. At infinity there is the fixed tensor phi\_(1,infty) tensor phi\_(0,infty). Equation(GS7) is the finite distributive law after CRT, with the exact source scalars inside the factors.

At G the first sharp factor forces the sole class y=q\^{}(-1)e\_a; LIFT22(CG7--9) proves this in every target coordinate. At Q all classes in (GS6) remain. The complete local Q sum is the one in DO037, including its Fourier/Gauss scalar.

\subsubsection{Why this is one fixed Schwartz tensor under the diagonal action}\label{why-this-is-one-fixed-schwartz-tensor-under-the-diagonal-action}

The following pointwise operator identity is stronger than a bare CRT count. Every place at which Y varies is relative split. Identify its unitary group with GL over F\_v using the selected E-place, as in DO020 and LIFT24. Every prescribed local frame is then an actual GL matrix and hence an element of this split unitary group; its opposite E-component is its inverse adjoint. Write g,h,t for the initial, middle and terminal frames at this place. The genuine Hom Weil action on the two Schwartz spaces is

\[
\begin{aligned}
 [\omega_0(g,h)f](X)
   &=\kappa_0(g,h)f(h^{-1}Xg),\\
 \kappa_0(g,h)
   &=\eta_0(\det g)^{-1}\eta_0(\det h)
                         |\det g|^{b/2}|\det h|^{-a/2},\\
 [\omega_1(h,t)H](Z)
   &=\kappa_1(h,t)H(t^{-1}Zh),\\
 \kappa_1(h,t)
   &=\eta_1(\det h)^{-1}\eta_1(\det t)
                         |\det h|^{c/2}|\det t|^{-b/2}.
\end{aligned}
 \tag{GS7a}
\]

These are the honest split formulas with the prescribed character phase, not only symplectic covariance up to a scalar. For the fixed standard functions f\_(0,z),f\_(1,z) define Psi\_v by(GS7). Linearity of the two actual Weil operators gives exactly

\[
\begin{aligned}
 &[(\omega_1(h,t)\otimes\omega_0(g,h))\Psi_v](Z,X)\\
 &\quad=\kappa_1(h,t)\kappa_0(g,h)
    \sum_{z\in\mathcal Y_v}
      f_{1,z}(t^{-1}Zh)f_{0,z}(h^{-1}Xg).
\end{aligned}
 \tag{GS7b}
\]

For the literal middle lattice N\_h, its local support and period modules are hL\_+ and hL\_-, where L\_+/L\_- is the standard quotient in(GS6). The source frsF formula in FS1103--1104 uses precisely that same frame h in both XiCpl1 and XiCpl2. Thus its raw coordinate sum at (g,h,t) is the sum displayed after the two kappa factors in (GS7b): the local coordinate h\^{}(-1)repY runs through z in L\_+/L\_-, and the previously proved periodicity removes the difference between representatives. This comparison holds for arbitrary X and Z before taking theta sums. Multiplying by the actual frame factors(GS7a), as required by the genuine Weil reading of those raw arrays, gives equality with the diagonal-middle action on the \textbf{same fixed Psi\_v} at every h, not insertion of a new h-dependent projector.

In particular the middle character and modulus in that product are

\[
       (\eta_0\eta_1^{-1})(\det h)|\det h|^{(c-a)/2}
             =(\eta_0\eta_1^{-1})(\det h)|\det h|^d.
 \tag{GS7c}
\]

They are independent of z and agree with the exact middle factor in DO037's double full-rank period. They have not been replaced by one. The initial and terminal factors in(GS7a) likewise remain present.

At places outside G union Q there is no Y insertion at all: the fixed tensors transform by their genuine Weil operators. Infinity also has the fixed selected tensor. Taking the product of(GS7b) and these unchanged factors proves the global fixed-tensor identity. Finally summing over the two rational matrix variables gives the actual product of theta kernels. This proves that its global composition is evaluated through q\_1(1 tensor q\_0), as used in (GS19) and(GS27), rather than through an unproved moving-family operator. The fixed common period/creation factors in the actual mixed realization multiply both descriptions equally and are independent of Y; they remain in the chosen phi\_j and A\_j.

The first and middle quotients are compact. Theta series for these Schwartz vectors converge uniformly on compact representatives, and the Y sum is finite. Thus the double integral and finite sum in (GS4) may be interchanged without a regularization or rearrangement. The formula is the actual composition of two global theta maps evaluated on the tensor(GS3), with the actual quotient Haar measures.

\subsection{3. The standard factor for every actual first constituent}\label{the-standard-factor-for-every-actual-first-constituent}

Fix an irreducible automorphic constituent pi\_0 of the input space. It is unitary. DO036 proves nonzero first lift on every embedding of pi\_0 that occurs here, using the unchanged first raw kernel and its strict positive norm bound. Its real type is the literal selected one, and its first descendant is denoted pi\_1.

At an unramified place w of E outside a finite set S, local theta gives exactly

\[
 \mathcal A_{1,w}=\mathcal A_{0,w}\sqcup
 \{\chi_{0,w}(\varpi_w)Q_w^{(d-1)/2-l}:0\le l<d\}.
 \tag{GS8}
\]

Here Q\_w is the residue cardinality in E and A is the standard Satake multiset. At relative split places this is Mínguez's type-II formula with the actual Hom factor eta\_T(det g)\^{}(-1)eta\_S(det h). At nonsplit places it is Liu's unramified theorem in the plus tower, with the paired eigenvalues in Q\_w=q\_v\^{}2 units. These are the checked LIFT14(E10--13) formulas. Absolute ramification of the host is allowed; the assertion is only made outside S when identifying Euler parameters.

DO028 fixes the bilinear-input and Hom-to-tensor conventions. In particular the standard factor on edge1 uses chi\_1 inverse. Set

\[
 s_0=d/2,\qquad s_*=(d+1)/2.
 \tag{GS9}
\]

Substitution of (GS8), initially in a right half-plane, gives

\[
 \mathcal L_1^S(s_*)=
 L_E^S(s_*,\operatorname{BC}_{\rm unr}(\pi_0)\chi_1^{-1})
          \prod_{l=1}^{d}L_E^S(l,\upsilon).
 \tag{GS10}
\]

The notation BC\_unr means only the actual standard unramified Euler parameters. No global automorphic base-change representation or Arthur parameter is assumed. The equality at s\_* is obtained by continuing the identity of Euler products. In particular we do not infer the first constituent's real Langlands parameter from (GS8).

The crude unitary spherical bound used and checked in LIFT14 is

\[
                  |\alpha_w|\le Q_w^{(a-1)/2}.
 \tag{GS11}
\]

It applies to every initial constituent, without temperedness. The real part available in its first factor in (GS10) is

\[
 s_*-(a-1)/2=(d-a+2)/2=860\quad\hbox{or}\quad734.
 \tag{GS12}
\]

That Euler product is therefore absolutely convergent and nonzero. The factors at l\textgreater=2 are also absolutely convergent and nonzero. For l=1, upsilon and its square have nontrivial angular type by (GS2). The already proved LIFT14 boundary lemma applies over this field as well: Hecke--Tate continuation gives holomorphy at1, and

\[
 3+4\cos\theta+\cos2\theta=2(1+\cos\theta)^2\ge0
 \tag{GS13}
\]

in the Euler logarithms excludes a zero at1. Consequently (GS10) is holomorphic and nonzero at the stated point. This is not a claim of absolute Euler-product convergence for the l=1 factor, nor a quantitative lower bound for that value.

The finite G ratio is also nontrivial, which independently excludes the trivial Hecke character. The angular condition and its square are retained because they are the hypotheses of the exact boundary proof being used. The fixed stage characters are not varied.

\subsection{4. Rallis gives a nonzero global second map on every embedding}\label{rallis-gives-a-nonzero-global-second-map-on-every-embedding}

For the second edge n=b=1802 and m=c=b+d, we have

\[
                         n<m<2n,\qquad r=0.
 \tag{GS14}
\]

The target is positive at a real place and hence anisotropic over F; the same holds at all real places here. Thus its global Witt index is r=0 and the lower-tower vanishing hypothesis is empty. All automorphic representations of these anisotropic groups are cuspidal.

This is the \textbf{anisotropic convergent case within Theorem11.4's dimension window}, not the geometric second-term range. GQT Proposition11.1(ii), printed p.53, and Theorem11.4, printed pp.54--55, apply exactly in this case.

At a split finite place the actual small theta lift is nonzero by Mínguez, including the ramified G/T/Q inputs. At a nonsplit finite place the input is spherical and the relative extension is unramified; Liu's Theorem6.2 applies to the identity Hermitian lattices. The absolute-codifferent rescaling of the skew form and additive character is the one already checked in LIFT14 and DO022, so it preserves the actual lattice vector and the composite symplectic character. No classification for a split place is substituted at a nonsplit place. The actual basic vector at every remaining place is nonzero in the quotient by LIFT16; this includes places over29 with their actual absolute ramification. For K031/K032,29 is in fact relatively split, since it splits in Q(zeta28), and adjoining the real fields preserves that splitting. Dyadic places are also relatively split. The odd nonsplit places use Liu's theorem over their actual local fields.

At a real place, LIFT25 gives the exact selected summand and a nonzero compact contraction on the first highest line. More generally the matching source/target vectors give the local doubling test

\[
 \int_{U(b)}|\langle\rho_1(g)v,v\rangle|^2\,dg
               =\operatorname{vol}(U(b))/\dim\rho_1>0.
 \tag{GS15}
\]

The real standard gamma factors are finite and nonzero at s\_*\textgreater0: the compact unitary representation has bounded parameter and the character twist is angular. This supplies the real local test directly; it does not invoke the unresolved general real converse in GQT11.6.

At finite places GQT Proposition11.6(i) gives existence of a nonzero normalized test on the actual Siegel--Weil image. Therefore all local theta lifts and all required local tests are nonzero. GQT11.4(ii) gives holomorphy of the completed standard factor at s\_*. Its partial factor(GS10) is nonzero. The finitely many missing standard local factors have no zeros, and the real gamma factors have no zeros; holomorphy of the completed factor therefore excludes a pole in any missing factor and makes its value nonzero.

GQT11.4 now proves that the global second theta map is nonzero on \textbf{every actual irreducible embedding} of pi\_1. The Euler product and local tests depend on the representation, while its positive invariant inner product is the one inherited from that embedding. Replacing an embedding by a nonzero vector in its multiplicity space changes the norm by its positive squared length and does not change the argument.

At this point only existence for an irreducible embedding has been proved. The following factorization step is what makes that statement useful for the fixed raw vector.

\subsection{5. Global factorization and the automorphic multiplicity map}\label{global-factorization-and-the-automorphic-multiplicity-map}

We state the precise algebraic argument to avoid identifying a representation with a preferred automorphic copy.

Let pi be an irreducible unitary cuspidal representation in a source automorphic space, M\_pi its Hilbert multiplicity space, and let S be the algebraic Schwartz space, with K-finite vectors at infinity. Assume all local small theta lifts sigma\_v are nonzero and that the global theta map is nonzero on every embedding belonging to M\_pi. Fix nonzero actual local quotient maps in the bilinear convention

\[
 q_v:S_v\otimes\pi_v\longrightarrow\sigma_v.
 \tag{GS16}
\]

At almost every place normalize q\_v to send the basic tensor to the unit spherical vector. Let q be their restricted tensor product. Then the actual global theta map has the form

\[
 \Theta_\pi=A_\pi\otimes q:
 M_\pi\otimes(S\otimes\pi)
            \longrightarrow M_\sigma\otimes\sigma,
 \qquad \sigma=\bigotimes'_v\sigma_v,
 \tag{GS17}
\]

where A\_pi is \textbf{injective}.

Here is a proof. The target automorphic Hilbert space is a discrete unitary sum of irreducibles; fixed finite level and real type give finite-dimensional subspaces. Projecting the global theta map to an irreducible target constituent gives at each place a nonzero local balanced Hom. Local Howe duality forces that constituent to have local representation sigma\_v and makes the local Hom one-dimensional. At finite nonsplit places we use Gan--Takeda, Theorem1.2; at split places Mínguez; at compact real places this is the actual multiplicity-free Cauchy decomposition. Thus all target projections have the same global representation sigma.

The tensor product of the local one-dimensional Hom spaces is also one-dimensional. To see this without an infinite tensor ambiguity, fix the basic vectors outside a finite set. On that finite set, successive local Hom factorization makes every coefficient map a multiple of the product q\_v. Enlarging the set preserves the same multiple because the basic values were normalized to one. Every algebraic input is in one such finite tensor subspace. This gives (GS17), with all remaining coefficients contained in A\_pi.

If A\_pi killed a nonzero m in M\_pi, the embedding v -\textgreater{} m tensor v would have zero theta lift for every Schwartz vector. This contradicts the assumed nonvanishing on that actual embedding. Hence A\_pi is injective. No multiplicity-one assertion for the automorphic spectrum is used.

There is also a useful metric refinement. Give the source and target representation factors their chosen invariant unitary metrics and give the multiplicity spaces the actual quotient metrics. The global doubling identity, applied to any unit multiplicity vector, has the same representation matrix coefficient and the same local factors. Polarization in the multiplicity vectors therefore gives

\[
                    A_\pi^*A_\pi=c_\pi\,1_{M_\pi},
                           \qquad c_\pi>0.
 \tag{GS18}
\]

This also follows by applying the norm identity to every nonzero linear combination of two embeddings. The scalar is retained; it is not set to one by a choice of local quotient model.

The argument only uses irreducible quotients of the local big theta module. It does not assume that the big theta module is itself irreducible or that an arbitrary invariant form on it is positive. Automorphic unitarity provides the semisimple target projections.

For our first edge, DO036 supplies the hypothesis for every pi\_0 embedding with the initial prescribed type. For the second edge, Section4 supplies it for every embedding of its actual descendant. We obtain injective maps A\_(0,pi0) and A\_(1,pi1).

Finally, distinct initial irreducible representations cannot acquire the same second global descendant: at every place the local Howe correspondence is injective on its nonzero domain. Applying that injectivity on the second and then first edge would identify all local components of the two initial representations, hence the initial restricted tensor products themselves. Their global multiplicity copies are already handled by A\_1 A\_0.

\subsection{6. The actual local coupled tensors are injective on the input types}\label{the-actual-local-coupled-tensors-are-injective-on-the-input-types}

For a constituent pi\_0, let W\_(0,v) be the part of its local space occupied by R017's initial type. The projection of F\_0 to pi\_0 stays in the tensor product of these subspaces: spectral projections commute with all local actions and compact projectors.

At G this is the finite Steinberg isotypic subspace intersected with the prescribed inverse-generic subspace. At T it is the nontrivial projective-line character space. At Q and at all other finite places it is the maximal compact fixed line. At infinity it is the selected highest line. For every place define

\[
 C_v(w)=q_{1,v}\bigl((1\otimes q_{0,v})(\Psi_v\otimes w)\bigr).
 \tag{GS19}
\]

For a sum in Psi\_v this means the corresponding finite sum of compositions, with the literal coefficients. The checked local results give the following exact detection identities.

\subsubsection{G places}\label{g-places}

LIFT22(CG16) gives, in the nested actual induced quotient models,

\[
                   [C_v(w)(1)](1)=a_q v_{J^\sharp}\,w,
 \qquad a_q=|U_a(\mathbb F_q)|/|GL_a(\mathbb F_q)|.
 \tag{GS20}
\]

The second scalar is the exact Haar volume of the actual sharp middle group. It is nonzero. Only the forced Y class contributes. This identity holds on the whole initial generic type, so it remains injective on the finite Steinberg subspace. In particular no erroneous full generic projector on the joining root is introduced.

\subsubsection{T places}\label{t-places}

LIFT24 proves that the actual first output of the raw T factor lies on the required one-dimensional projective-line type of pi\_1. DO036 proves it is nonzero for the initial type. For every vector w\_1 on that line, the fixed second raw incidence factor satisfies

\[
 q_{1,v}(\phi_{1,v}\otimes w_1)(k_*^{-1})
  =\frac{L((d+1)/2,\sigma_{1,v}\eta_{1,v}^{-1})}{q^b-1}\,w_1.
 \tag{GS21}
\]

Here pi\_1=sigma\_1 times xi in LIFT24's exact local induction and eta\_j is the local component of chi\_j. The displayed L-factor is finite and nonzero. Thus C\_v is injective on the initial T line. The first local scalar is not replaced by one or by its compact leading term; its nonzero GJ/tail value remains part of C\_v.

\subsubsection{Q places}\label{q-places}

DO037 proves for the \textbf{complete local Y sum} that evaluation at the actual unit-transported target graph and then at the first identity has value

\[
             \ell_v(C_v(w))=\Lambda_{Q,v}w,
 \tag{GS22}
\]

where, with r=a+1, t=d-1, theta=eta\_0 eta\_1 inverse and kappa\_n=product\_(i=1)\^{}n(1-q\^{}(-i)),

\[
 \Lambda_{Q,v}=q^{-d}G(\chi,\psi_0)\chi(-1)^{-1}
       q^{-t(a+1)}\frac{\kappa_{a+1}\kappa_t}{\kappa_b}
        \prod_{i=1}^{t}(1-\theta(\varpi)q^{-i})^{-1}\ne0.
 \tag{GS23}
\]

The source input here is its full maximal compact fixed line, exactly as in R017. The graph has determinant one after the actual piT-unit transport. DO037 retains the cofactor sign, Gauss sum, all Haar constants and the initial compact restriction. Its theorem is a whole-Schwartz small-quotient identity, not merely a finite-array Fourier calculation. It proves injectivity of C\_v on that line.

\subsubsection{Remaining finite places}\label{remaining-finite-places}

Both vectors are the actual selfdual lattice vacua. The fixed-vacuum theorem LIFT16 gives nonzero image of the spherical line on each edge, with the actual additive character and local fields. Thus C\_v is nonzero on a one-dimensional space. We may normalize its quotient models so that the basic vectors have image one; the inverse normalization is then retained in A\_j. This does not change the actual global operator. In particular the absolute ramification at29 does not introduce a principal29 level or a new finite translation.

\subsubsection{Real places}\label{real-places}

LIFT25 proves that the first highest output is exactly the source line paired by the second primitive highest polynomial. In its specified auxiliary unitary quotient models,

\[
 C_\infty(h_0)
       =\frac{A_0A_1}{\sqrt{d_0d_1}}h_2\ne0,
 \tag{GS24}
\]

where A\_j is the actual selected Schwartz norm and d\_j the exact dimension of the selected source representation, taking the tensor over real places. The literal two compact coefficient integrations instead carry 1/(d\_0 d\_1), as proved there. These are different descriptions of the same map in different output metrics; neither is silently substituted for the other. The actual Gaussian, creation, hook and period factors of LIFT25 remain present.

Equations(GS20--24) give nonzero left inverses up to their displayed scalars on every W\_(0,v). Therefore their finite tensor product is injective. Outside finitely many places these are normalized basic lines. Consequently the restricted tensor map

\[
                    C_{\pi_0}=\bigotimes'_v C_v
 \tag{GS25}
\]

is injective on the actual initial type space W\_(0,pi0). There is no estimate asserting that the left inverses have norm one or are integral. Their scalar and evaluation costs are not discarded.

\subsection{7. Proof of fixed global nonvanishing}\label{proof-of-fixed-global-nonvanishing}

Decompose the actual finite-level, fixed-real-type input into its finitely many irreducible isotypic components,

\[
                  F_0=\sum_{\pi_0}F_{0,\pi_0},\qquad
 F_{0,\pi_0}\in M_{0,\pi_0}\otimes W_{0,\pi_0}.
 \tag{GS26}
\]

For each component use the two factorizations(GS17). The complete CRT tensor identity(GS7), not a choice of a nonzero summand, then gives the actual coupled output component

\[
 (\mathcal BF_0)_{\pi_2}
   =\bigl((A_{1,\pi_1}A_{0,\pi_0})\otimes C_{\pi_0}\bigr)
                                                    F_{0,\pi_0}.
 \tag{GS27}
\]

Both multiplicity maps are injective, and C\_pi0 is injective on the indicated type space. Thus a nonzero input component has a nonzero output. Distinct pi\_0 components have inequivalent pi\_2 descendants by local Howe injectivity and are orthogonal in the target automorphic space. They cannot cancel. Since F\_0 is nonzero, at least one component is nonzero, proving(GS5).

This proof permits arbitrary finite linear combinations in the initial type space and arbitrary multiplicity vectors. It does not replace a vector by a favorable irreducible projection before declaring the original sum nonzero. It also does not claim that a fixed highest second kernel preserves nonhighest real vectors; those vectors are outside the stated input space and LIFT25 explicitly records that obstruction.

\subsection{8. Exact metric identity and costs still retained}\label{exact-metric-identity-and-costs-still-retained}

With the chosen invariant unitary quotient models, write

\[
 A_{j,\pi_j}^*A_{j,\pi_j}
                =c_{j,\pi_j}\,1,\qquad c_{j,\pi_j}>0.
 \tag{GS28}
\]

Equation(GS27) and orthogonality give the exact norm identity

\[
 \boxed{\quad
 \|\mathcal BF_0\|^2=
 \sum_{\pi_0}c_{0,\pi_0}c_{1,\pi_1}
   \|(1\otimes C_{\pi_0})F_{0,\pi_0}\|^2>0.\quad}
 \tag{GS29}
\]

The constants are actual global theta norms, not arbitrary pins. The first is determined by DO036's convergent norm identity with its actual quotient mass. For the second, in Tamagawa measures, GQT gives the following exact expression. Choose decomposable local tensors x\_v=phi\_v tensor v\_v with q\_v(x\_v) nonzero, normalize q\_v/basic metrics as above, and set

\[
 b_v=\frac{Z_v^*(s_0,\delta(\phi_v\otimes\overline{\phi_v}),
                                     v_v,v_v)}
                  {\|q_v(x_v)\|^2}.
 \tag{GS30}
\]

Then

\[
                  c_{1,\pi_1}
      =2L(s_*,\pi_1\times\chi_1)\prod_v b_v>0.
 \tag{GS31}
\]

The factor on the right is expressed in the actual GQT convention through DO028, hence has the standard unramified factor(GS10). In(GS30) the local input and oscillator are transported by g -\textgreater{} g\^{}(-t) and conjugation exactly as in DO028. Thus its printed GQT matrix-coefficient bar is retained; the notation does not replace the actual bilinear input by an unrelated Hermitian lift. At almost all places b\_v=1/d\_v(s\_0,chi\_1); the product of these factors converges and is nonzero by GQT(11.3). At the finitely many remaining places b\_v is the actual normalized doubling form relative to the specified local quotient metric. Equation(GS17) and the Rallis identity imply that this ratio is independent of a nonzero detecting tensor. One can see this by fixing all other local tensors with nonzero quotient and polarizing in x\_v: the global norm factors through the local irreducible unitary quotient. This also proves that the normalized local form is a nonzero scalar multiple of its positive quotient metric in this globalizable case. Its scalar is not assumed positive separately; the complete real quantity in(GS31) is positive. This is the needed radical statement, obtained after global factorization rather than assumed from local occurrence.

For source and target measures dg\_j=h\_j dg\_j\^{}Tam with h\_j\textgreater0, the relative multiplicity metric constant on edge j changes by h\_j h\_(j+1). Indeed the theta integral changes by h\_j, its target squared norm by h\_(j+1), and its source multiplicity squared norm by h\_j. Thus the formula transports to the actual probability-real, maximal-finite-compact measures without setting any quotient mass or global comparison scalar to one.

Equations(GS20--24),(GS29--31) retain all finite Gauss/projector volumes, T local factors, real dimension/hook/Gaussian factors and global L-values. They do not evaluate the finite number of initial automorphic constituents or the constants c\_j numerically. The nonzero global L(1,upsilon) factor is not assigned a uniform lower bound. No 11-adic valuation is attached to an unspecified real norm.

Most importantly, a nonzero complex automorphic output does not identify its integral image lattice. The actual algebraic global quotient scalars, the saturation and content of the prior-image matrix, the selected terminal branching projection, the TopForm row and the MainX terminal-mass equation remain separate. The unchanged elliptic Fields.Kq/VacIn constructor is still not a constructor for the primitive mixed host. Accordingly(GS5) is not a completed DP44 comparison or a shorter-chain theorem.

\subsection{9. Sources and review status}\label{sources-and-review-status}

Published inputs used at their stated hypotheses:

\begin{itemize}
\tightlist
\item
  Gan, W. T., Qiu, Y., \& Takeda, S. (2014). The regularized Siegel--Weil formula (the second term identity) and the Rallis inner product formula. \emph{Inventiones Mathematicae, 198}, 739--831. \url{https://arxiv.org/abs/1207.4709} . Sections11.2--11.8, especially Proposition11.1(ii), Theorem11.4, Proposition11.6 and equation (11.3), printed pp.51--56. Local primary copy: \nolinkurl{work/agents/lift_refs/gqt2014.pdf} and \nolinkurl{.txt}.
\item
  Gan, W. T., \& Takeda, S. (2016). A proof of the Howe duality conjecture. \emph{Journal of the American Mathematical Society, 29}, 473--493. \url{https://arxiv.org/abs/1407.1995} . Theorem1.2 and the exact definition of the maximal semisimple quotient, pp.1--3 of arXiv v4, were read directly. This covers all residual characteristics for nonsplit unitary pairs. The split case is expressly distinguished there and supplied by Mínguez.
\item
  Mínguez, A. (2008). Correspondance de Howe explicite: paires duales de type II. \emph{Annales Scientifiques de l'École Normale Supérieure, 41}(5), 717--741. Theorem1 and equation(1.4), carried from the direct readings in LIFT14/19/21/24.
\item
  Liu, Y. (2022). Theta correspondence for almost unramified representations of unitary groups. \emph{Journal of Number Theory, 230}, 196--224. arXiv:2105.14223v1, Theorem6.2 and Proposition3.4, used in the exact spherical plus-tower scope already checked in LIFT14/16.
\item
  DO028 for the complete Hom/tensor, conjugation and standard factor dictionary; DO032 for the literal CRT/representative identity; DO036 for the first actual global map; DO037, LIFT22/24/25 for the exact actual local coupled quotients; R017 and K031/032 for the initial function and actual host.
\end{itemize}

The global multiplicity argument, the fixed coupled assembly and the exact norm decomposition are proved above. The literature theorems and carried local findings are not counted again as new results. No Lean build or automated proof of the global theorem is claimed. Whole-proof review events are recorded only after the finished text is read in full.

\clearpage

\section{AM-K037 --- an actual mixed CM arithmetic class-row adapter}\label{am-k037-an-actual-mixed-cm-arithmetic-class-row-adapter}

Author: GPT-6 Astra Pro, ALP\_MACHINE /root/kummer\_extension. 22 November 2026 UTC; inherited reasoning effort is not separately exposed.

This finding constructs the arithmetic maps of the actual mixed family from its geometric reduced-theta kernels. It then proves their exact comparison with a canonical finite sum of genuine theta lifts. The coefficient spaces, their Weyl lattices, the class stabilizers, and the geometric kernels are fixed before this comparison. No lattice is fitted to analytic values.

In particular, the separate ``identify an externally specified new-host Kval map'' obligation in K035 is unnecessary for this constructed replacement family. The map below has the literal subD/subY formula of Kval. It need not be an instance of the false old Fields.Kq constructor. The remaining quantitative question concerns the image of this actual map on its chosen lattice, its successive images, and its norms.

The accepted coefficient-module results are used as theorems. In addition to the already checked K035--036 and DO025--027, the proof uses the source's Weyl GL-stability and normalized Weyl/dual-Weyl pairing. An elementary localization argument below extends that integral stability from generic invertible matrices to all matrices. This supplies a stronger statement than mere characteristic-zero membership: the selected universal substitution itself preserves the Weyl lattice.

\subsection{1. Fixed data and the exact source interfaces}\label{fixed-data-and-the-exact-source-interfaces}

Use either actual polarized host of K030/K032: \[
 (K_h,A_h)=(\mathbb Q(\zeta_{28})B,J^{14})\quad(h=1),
 \qquad
 (\mathbb Q(\zeta_{28})B(\sqrt5,\sqrt{13}),J^{56})\quad(h=4).
 \tag{CR1}
\] Keep its actual CM type, trace pairing, source additive-character convention, theta line, finite vacuum and all chosen finite operators. For one edge let \(a=N_j<b=N_{j+1}\) be even and use the same prescribed character \(\mu_{j+1}\) on the two members of the pair. At each CM label \(\tau\), put \[
\begin{aligned}
 \lambda_\tau&=\operatorname{lamTrue}_{j,\tau},&
 \lambda'_\tau&=\operatorname{lamTrue}_{j+1,\tau},\\
 s_{S,\tau}&=e_{j,\tau}-b/2,&
 s_{D,\tau}&=e_{j,\tau}-a/2,\\
 \eta_\tau&=\lambda_\tau-\lambda_{\tau,a-1}\mathbf1_a,&
 q_\tau&=\lambda_{\tau,a-1}+s_{S,\tau},\\
 \nu_\tau&=\eta_\tau+q_\tau\mathbf1_a.&
\end{aligned}
\tag{CR2}
\] The source weight identities, as already checked in DO025 and R011, give \[
 q_\tau\ge1,\qquad
 \lambda'_{\tau,b-1}=-s_{D,\tau},\qquad
 \eta'_\tau=(\nu_\tau,0^{b-a}).
 \tag{CR3}
\] Write \(D_\tau=|\nu_\tau|\). These are the actual selected partitions; they are not substituted by more convenient highest weights.

The source locators used in the construction are as follows.

\begin{longtable}[]{@{}
  >{\raggedright\arraybackslash}p{(\columnwidth - 2\tabcolsep) * \real{0.5000}}
  >{\raggedright\arraybackslash}p{(\columnwidth - 2\tabcolsep) * \real{0.5000}}@{}}
\toprule\noalign{}
\begin{minipage}[b]{\linewidth}\raggedright
Recovered source file
\end{minipage} & \begin{minipage}[b]{\linewidth}\raggedright
Exact use
\end{minipage} \\
\midrule\noalign{}
\endhead
\bottomrule\noalign{}
\endlastfoot
\nolinkurl{Formalism/GradedWeylDivisibility/StatementA.lean,} 26--38, 68--80 & Column equations for dualWeyl, rowDiv, and the integral Weyl lattice generated from hw \\
\nolinkurl{Formalism/LocalRingCardinality/DataCoeff.lean,} 53--84, 115--143 & nablaSq, delSq, square/round actions, normalized Fischer metric and the Weyl/dual-Weyl pairing \\
\nolinkurl{Formalism/LocalRingCardinality/RankBoundMonotoneInDimension.lean,} 902--909 & Proven factSqStable, invoking the accepted WeylGLStable theorem over every commutative base ring \\
\nolinkurl{Formalism/LocalToGlobal/Divisibility/CoefMod.lean,} 60--80, 178--226, 309--330 & VV, its polynomial action, the fixed standard/transferred lattices, rhoV, rhoVr and delAtr \\
\nolinkurl{Formalism/LocalToGlobal/Columns/CoefModB.lean,} 55--90 & Literal symmetric bilinear pairBK and its rhoVr/rhoV invariance \\
\nolinkurl{Formalism/LocalToGlobal/KernelLattices/ThDistC.lean,} 19--43, 61--65, 137--148 & subD, jetLin, jetMap, retrVV, JetInVV, ordinary derivatives and the distribution map \\
\nolinkurl{Development/AbelianVarietyTheory/Definitions/LatticeFrames/EvalAtOneExtremalMul.lean,} 17--38, 90--105 & Vn, evalAt, mvZ, subY, jetL and subY\_eq\_subD \\
\nolinkurl{Development/AbelianVarietyTheory/Definitions/LatticeFrames/OnFramesRdApply.lean,} 96--121 & Real, Lev, Core and the exact Kval equation \\
\nolinkurl{Formalism/LocalToGlobal/Columns/Record.lean,} 1285--1300 & The source integral column uses delAtr, not a replaced dual-Weyl lattice \\
\end{longtable}

Take one number field \(M\), stable under complex conjugation, containing the host, the fixed finite character values, all fixed chart constants, and the all-order geometric jet field supplied by K036 for these ranks and allowed denominators. This is a finite extension chosen once. The choice has no effective degree or discriminant estimate here.

For every actual finite class pair \((c,d)\), the independently defined geometric kernel is the finite Mumford average and translation sum of DO026, with its actual character and central phases. Its normalized functional will be denoted \[
 J_{dc}(P)=\epsilon_{dc}\,
             [P(cD_\tau\partial_{\tau,rs})F_{dc}]_0\in M .
 \tag{CR4}
\] The coefficients, phases and translations in \(F_{dc}\) are fixed by the chosen finite Schwartz datum, before any analytic theta value is considered. Any already specified chart factors of DO027 are included in \(\epsilon_{dc}\). Formula (CR4) uses ordinary derivatives. No division by a multiindex factorial is inserted.

The finite input and row levels are required to have their actual covariance \[
 \omega_f(k,k')\phi_{j,f}
       =\chi_j(k)^{-1}\chi_{j+1}(k')\phi_{j,f}.
 \tag{CR4a}
\] This is not inferred from weights at infinity. DO026 supplies explicit finite stabilizers for every fixed geometric finite span, with the necessary determinant-character kernels; its finite refinement index is retained. When a more specifically prescribed raw finite family and level are used, their checked actual covariance is the input to the same theorem. Nothing here identifies an arbitrary pinned level with that stabilizer. Several chart blocks may have different fixed finite vectors, each satisfying its displayed source and target covariance.

The exact source-creation comparison proved in K035 is \[
 \vartheta_{\operatorname{pureTensorO}(\phi_f,P)}(g_c,h_d)
   =C^{14hab}J_{dc}\!\left(P(s_\tau Z_{\tau,rs})\right),
 \quad
 s_\tau=-\frac{\Omega_\tau}
                   {cD_\tau\sqrt{\pi q^{\rm met}_\tau}},
 \quad q^{\rm met}_\tau=-2i\,\tau(\xi_h)>0.
 \tag{CR5}
\] The notation \(q^{\rm met}_\tau\) distinguishes this positive metric coefficient from the integer \(q_\tau\) in (CR2). The common period \(C=\theta_0c\), the periods \(\Omega_\tau\), and the algebraic coordinate generators \(cD_\tau\) are exactly those of K035. In particular the scalar in (CR5) is not selected separately at different derivative orders.

\subsection{2. The universal substitution, including its integral Weyl image}\label{the-universal-substitution-including-its-integral-weyl-image}

For a commutative ring \(R\), let \(\nabla_\eta(R)\) and \(\Delta_\eta(R)\) be the source's dual-Weyl and Weyl polynomial modules. Their highest polynomial is the product of leading minors \[
 p_\eta(Y)=\prod_{r=1}^a
             \det(Y_{[1,r],[1,r]})^{\eta_r-\eta_{r+1}},
 \qquad \eta_{a+1}=0.
 \tag{CR6}
\] Over a characteristic-zero field these two modules have the same span. Their integral lattices are retained separately.

Use \(a\)-by-\(a\) variables \(U\), \(b\)-by-\(b\) variables \(Y\) and \(b\)-by-\(a\) variables \(Z\), at each \(\tau\). Define \[
 {\cal Q}(P)(Y,Z)=
 \left[
     \left(\prod_\tau\det(U_\tau)^{q_\tau}\right)P(U)
 \right]_{U_\tau=Z_\tau^{\mathsf T}Y_{\tau,[1,b],[1,a]}}.
 \tag{CR7}
\] This is exactly subD after the source determinant insertion. Its evaluation at \(Y=y\) is exactly subY(y) applied to that insertion. All coefficients of the substitution are integers.

\subsubsection{Theorem CR8: the actual selected substitution has integral image}\label{theorem-cr8-the-actual-selected-substitution-has-integral-image}

For every commutative ring \(R\), \[
\begin{split}
 {\cal Q}\left(\bigotimes_\tau\nabla_{\eta_\tau}(R)\right)
   &\subset
       \left(\bigotimes_\tau\nabla_{\eta'_\tau}(R)\right)[Z],\\
 {\cal Q}\left(\bigotimes_\tau\Delta_{\eta_\tau}(R)\right)
   &\subset
       \left(\bigotimes_\tau\Delta_{\eta'_\tau}(R)\right)[Z].
\end{split}
\tag{CR8}
\] Here membership means that every \(Z\)-coefficient belongs to the displayed module in the \(Y\)-variables. The second assertion retains the source's chosen Weyl lattice; it is stronger than JetInVV over a field.

\textbf{Proof of the first assertion.} It is enough to work at one label. Multiplication by \(\det(U)^q\) changes the column weight from \(\eta\) to \(\nu=\eta+q\mathbf1\), and preserves invariance under adding an earlier column to a later one. Replacing \(U\) by \(Z^{\mathsf T}Y_{[,1:a]}\) preserves these same column identities in \(Y\). The columns beyond \(a\) do not occur. Thus the polynomial has column weight \((\nu,0)\) and satisfies every defining upper-column invariance of \(\nabla_{\eta'}(R[Z])\). Comparing \(Z\)-coefficients proves the assertion, directly over \(R\).

\textbf{Integral extension of row stability.} First work over \(\mathbb Z\) and fix a target Weyl lattice \(\Delta\) in its finite-degree ambient polynomial module \(A\). Let \(G\) be a generic \(b\)-by-\(b\) matrix. The accepted factSqStable theorem, with final weight zero, implies \[
                  \operatorname{rowAct}(G)\Delta
       \subset\Delta\otimes_{\mathbb Z}\mathbb Z[G,\det(G)^{-1}].
 \tag{CR9}
\] The harmless transpose in sqAct is removed by taking its group argument to be \(G^{\mathsf T}\). The variable determinant has a leading monomial with coefficient \(1\) or \(-1\), for any monomial order. Consequently multiplication by \(\det(G)\) is injective on \((A/\Delta)[G]\), even when \(A/\Delta\) has torsion: the leading coefficient of a nonzero polynomial in this module is multiplied by a unit, and cannot disappear.

The left side of (CR9) already has coefficients in \(A[G]\). Its class in \((A/\Delta)[G]\) vanishes after inverting \(\det(G)\), so some power of \(\det(G)\) kills it. The preceding injectivity forces that class to be zero before localization. Thus \[
                 \operatorname{rowAct}(G)\Delta\subset\Delta[G].
 \tag{CR10}
\] Polynomial specialization proves stability under every row matrix, not only invertible ones. This argument uses flat base change for the polynomial and localized polynomial rings; it does not assume that the integral Weyl lattice is saturated in its ambient module.

\textbf{Proof of the Weyl assertion.} Multiplication by \(\det(U)^q\) takes \(p_\eta\) exactly to \(p_\nu\). It commutes with every off-diagonal row divided power, because an elementary row addition does not change the determinant. The diagonal row divided powers act on a row-homogeneous polynomial of row weight \(\mu\) by the integer \(\binom{\mu_i}{r}\). Every divided-power word applied to the highest polynomial remains row homogeneous (or becomes zero). Diagonal factors in such a word can therefore be removed at the cost of integer scalar factors. The defining divided-power generators show \[
                  \det(U)^q\Delta_\eta(\mathbb Z)
                                  \subset\Delta_\nu(\mathbb Z).
 \tag{CR11}
\] Embedding the first \(a\) rows and columns into the target square matrix takes \(p_\nu\) to \(p_{(\nu,0)}\). The same row divided-power words on those first \(a\) rows show that it takes the entire module in (CR11) into the target Weyl lattice.

Now apply (CR10), and specialize the first \(a\) rows of its generic matrix \(G\) to \(Z^{\mathsf T}\); the remaining rows can be zero. The embedded polynomial does not use those remaining rows of \(GY\). Its image is exactly (CR7). Therefore all its coefficients belong to the target Weyl lattice over \(\mathbb Z\). All Weyl generators are defined by integral polynomial formulas. Base change sends their integral span to the defining span over any \(R\). This proves the result over \(R\), and tensoring the separate labels proves (CR8). \(\square\)

For any \(M\)-linear polynomial functional \(J\), define \[
                    \Phi_J(P)(Y)=J_Z({\cal Q}(P)(Y,Z)).
 \tag{CR12}
\] Theorem CR8 proves the exact JetInVV premise of ThDistC for these weights. Its retrVV projection is the identity on this image, whatever complementary subspace was used to define that auxiliary projection. This is a proved membership statement, not a choice of projection which fits the desired output.

In standard integral coordinates, if all degree-profile values \(J(Z^\alpha)\) belong to a fractional ideal \(\mathfrak a\), then \[
                       \Phi_J(\Delta_\eta)
                                   \subset\mathfrak a\Delta_{\eta'}.
 \tag{CR13}
\] There is no hook-product or Weyl/dual-Weyl index denominator in (CR13). Such indices remain relevant to other projections and duality operations; the assertion concerns this particular universal substitution and this particular input Weyl lattice.

\subsection{3. The canonical Cauchy map and its exact metric normalization}\label{the-canonical-cauchy-map-and-its-exact-metric-normalization}

Let \(V_j\) denote the source polynomial space VV with action \(\rho_j^\vee=\rhoVr\), and let \(W_j\) be the same polynomial space with action \(\rho_j=\rhoV\). Use the source's nondegenerate symmetric bilinear form \[
 B_j(P,v)=\operatorname{pairBK}_{\lambda_j}(P,v)
         =\frac{\operatorname{Fischer}_{\rm bil}(P,v)}
                        {\prod_\tau H_{\eta_{j,\tau}}}.
 \tag{CR14}
\] It identifies \(W_j\) with the algebraic dual of \(V_j\), and \[
 B_j(\rho_j^\vee(g)P,\rho_j(g)v)=B_j(P,v).
 \tag{CR15}
\] This is exactly the source convention; no Hermitian conjugation is inserted into the arithmetic map. Complex conjugation gives the positive compact Hermitian form when norms are taken.

For an invertible frame \(y\), let \(v_y\in W_{j+1}\) be uniquely defined by \[
 B_{j+1}(Q,v_y)=
          \prod_\tau\det(y_\tau)^{-s_{D,\tau}}Q(y).
 \tag{CR16}
\] Such evaluation vectors span \(W_{j+1}\), because a polynomial which vanishes on all invertible matrices is zero. The target membership in CR8 therefore defines a unique bilinear map \[
 \iota_j:V_j\otimes W_{j+1}\longrightarrow M[Z],\qquad
 \iota_j(P,v_y)=
        \prod_\tau\det(y_\tau)^{-s_{D,\tau}}{\cal Q}(P)(y,Z).
 \tag{CR17}
\] Equivalently, contract the target polynomial coefficients of \({\cal Q}(P)\) with \(v\) using \(B_{j+1}\). This definition fixes its normalization without choosing a projector or an arbitrary scalar in a multiplicity-one assertion.

Under the actual source oscillator convention, \[
 \omega_\infty(g,h)\operatorname{pureTensorO}(P)
  =\prod_\tau\det(g_\tau)^{s_{S,\tau}}
                \det(h_\tau)^{-s_{D,\tau}}\,
       \operatorname{pureTensorO}
            (P(h_\tau^{\mathsf T}Z_\tau g_\tau^{-\mathsf T})).
 \tag{CR18}
\] The map \(\iota_j\), followed by this creation map, intertwines \(\rho_j^\vee\boxtimes\rho_{j+1}\) with the actual selected oscillator summand.

\textbf{Verification of the two actions.} On the source, substituting \(Zg^{-\mathsf T}\) into (CR7) changes its matrix argument to \(g^{-1}Z^{\mathsf T}Y_{[,1:a]}\). The determinant insertion contributes \(\det(g)^{-q_\tau}\); multiplication by the vacuum character \(\det(g)^{s_{S,\tau}}\) leaves \(\det(g)^{-\lambda_{\tau,a-1}}\), precisely rhoVr. On the target, substituting \(h^{\mathsf T}Z\) changes \(y\) to \(hy\). The factor \(\det(h)^{-s_D}\) combines with \(\det(y)^{-s_D}\) to give \(\det(hy)^{-s_D}\). Equation (CR15) gives \(\rho_{j+1}(h)v_y=v_{hy}\). These identities prove (CR18) on the spanning vectors \(v_y\), and hence on the whole map.

The canonical map also has a fixed norm. Normalize the compact metrics on \(V_j,W_{j+1}\) by (CR14), so the extremal polynomial has norm one. With the ordinary Fischer norm on creation-label polynomials, \[
 \boxed{\|\iota_j(P,v)\|_{\rm F}^2
       =\left(\prod_\tau H_{\nu_\tau}\right)
                              \|P\|_j^2\|v\|_{j+1}^2}
 \tag{CR19}
\] for pure tensors; the corresponding tensor-product metric identity holds on their span.

To check the scalar, take extremal source \(P=p_\eta\) and the target frame \(y=1\). Its vector \(v_1\) is exactly \(p_{\eta'}\). Indeed evaluation at \(1\) vanishes on every row-weight space except the highest row weight \(\eta'\), whose line is spanned by \(p_{\eta'}\); that polynomial evaluates to one. Row weights are orthogonal in the Fischer form, and its normalized squared norm is one. This proves \(v_1=p_{\eta'}\) with no scalar. Formula (CR17) now gives the primitive leading-minor polynomial \(p_\nu(Z)\), with coefficient one on its leading diagonal monomial. Its squared Fischer norm is \(H_\nu\), by the checked R011 calculation. The selected Cauchy component is irreducible and has multiplicity one (DO025); invariance of the metrics then proves (CR19). This is a normalization check, using the accepted Cauchy decomposition rather than re-proving it.

The integral duality is also fixed: pairBK pairs the standard dual-Weyl lattice perfectly with the standard Weyl lattice in the opposite action. Thus CR8 and (CR17) take a source integral Weyl vector and a target integral dual-Weyl test vector to an integral polynomial. This statement names the two integral lattices explicitly; their characteristic-zero spans coincide. No use of a complex orthonormal basis below changes either arithmetic lattice.

\subsection{4. Arithmetic classes, unchanged integral lattices, and Kval}\label{arithmetic-classes-unchanged-integral-lattices-and-kval}

Fix the actual finite row levels and finite characters. For clarity first use one chart at each stage, with level \(U_j\), character \(\chi_j\), finite class set \[
 {\cal C}_j=G_j(F)\backslash G_j(\mathbb A_f)/U_j,
 \qquad
 \Gamma_c=G_j(F)\cap g_cU_jg_c^{-1}.
 \tag{CR20}
\] The groups are totally definite. Their fixed-rank class sets are finite by K036; adding the specified finite decorations refines them by finite sets. Stabilizers are finite. These are the actual arithmetic and adelic classes, with their actual lattices and decorations, not a class-number-one replacement.

For \(\gamma\in\Gamma_c\), put \(\kappa_\gamma=g_c^{-1}\gamma_fg_c\in U_j\). The arithmetic class-row space is \[
 {\cal H}_j(M)=
 \left\{(P_c)\in\bigoplus_{c\in{\cal C}_j}V_j:
     \rho_j^\vee(\gamma_\infty)P_c
                =\chi_j(\kappa_\gamma)P_c
                  \text{ for every }\gamma\in\Gamma_c\right\}.
 \tag{CR21}
\] This sign is fixed by the scalar automorphic realization in (CR26) below, whose right \(U_j\)-character is \(\chi_j\).

Let \(L_{j,c}\) be the literal delAtr lattice for the lattice Lsh attached to \(c\), including its source determinant twist and any fixed finite coefficient decoration. Define \[
 {\cal L}_j={\cal H}_j(M)\cap\bigoplus_c L_{j,c}.
 \tag{CR22}
\] This is the standard integral coefficient module used here. Its definition precedes the theta map. No averaging projector is used to redefine it. It is a full lattice in its \(M\)-space: the ambient lattices are full, the invariance equations are linear over \(M\), and clearing denominators in a basis of their solution space supplies a spanning set in (CR22). Locally over a DVR it is the saturated intersection with the specified subspace.

Define the actual arithmetic map by the finite formula \[
 \boxed{({\cal K}_jP)_d
           =\sum_{c\in{\cal C}_j}\frac1{|\Gamma_c|}
                              \Phi_{J_{dc}}(P_c).}
 \tag{CR23}
\] Everything on the right was fixed geometrically or algebraically in §§1--3. It is an \(M\)-linear map, without fitting to analytic values. Each summand belongs to the prescribed target polynomial space by CR8. The target stabilizer equations follow either directly from the rational and finite-level covariance of the same geometric kernels, or from the independent genuine-theta comparison proved in §5. Thus it maps \({\cal H}_j\) to \({\cal H}_{j+1}\).

For every frame \(y'\), formula (CR23) gives \[
\boxed{\begin{split}
 &\operatorname{val}_{j+1}({\cal K}_jP,d;y')\\
 &=\sum_c|\Gamma_c|^{-1}
    J_{dc}\!\left(
        \operatorname{subY}(y')
        \left[\prod_\tau\det(U_\tau)^{q_\tau}P_c(U)\right]\right)
                       \prod_\tau\det(y'_\tau)^{-s_{D,\tau}}.
\end{split}}
\tag{CR24}
\] The source frame value is \[
 \operatorname{val}_j(P,c;y)
   =\left[\prod_\tau\det(y_\tau)^{q_\tau}P_c(y)\right]
                         \prod_\tau\det(y_\tau)^{-s_{S,\tau}},
 \tag{CR25}
\] because \(q_\tau-s_{S,\tau}=\lambda_{\tau,a-1}\). Thus (CR24) is exactly the displayed source Kval formula, with \(n=q\), \(F_c=\det^qP_c\) and abelian₂ equal to the already fixed \(J_{dc}\). The equality holds for all rows and all frames, not only for one nonzero coefficient. It uses ordinary polynomial evaluation and the actual same-field jets.

For several finite chart labels \(\alpha,\gamma\), replace the source sum by \(\sum_\alpha\sum_{c\in{\cal C}_{j,\alpha}}\) and use the predetermined block kernel \(J_{\gamma d,\alpha c}\). The same proof applies blockwise. Equal source chart volumes give the common scalar in §5. If unequal volumes are used, their exact ratios must be included by the permitted finite input normalization of DO027, or retained as explicit factors in the analytic formula. They cannot be omitted.

\subsection{5. Comparison with a canonical sum of genuine theta lifts}\label{comparison-with-a-canonical-sum-of-genuine-theta-lifts}

Let \(v_j=\operatorname{vol}(U_j)\), with infinite compact Haar of volume one. For a row \(P\) and \(v\in W_j\), define the actual scalar automorphic function by \[
 f_{P,v}(g_ck_\infty k_f)
       =\chi_j(k_f)B_j(P_c,\rho_j(k_\infty)v).
 \tag{CR26}
\] Equation (CR21) makes this well defined: replacing the representative by \(\gamma\in\Gamma_c\) multiplies the finite factor by \(\chi_j(\kappa_\gamma)\) and the bilinear coefficient by its inverse. Conversely these are precisely its stabilizer conditions. Hence this is the canonical coefficient realization of the actual row, with its prescribed right level character.

Take any basis \(v_i\) of \(W_j\) and its pairBK-dual basis \(p_i\) of \(V_j\), so \(B_j(p_i,v_l)=\delta_{il}\). For \(v\in W_{j+1}\), define \[
 {\cal T}_j(P;v)(h)
  =\sum_i\int_{[G_j]}
       f_{P,v_i}(g)\,
       \vartheta_{\operatorname{pureTensorO}
                  (\phi_{j,f},\,\iota_j(p_i,v))}(g,h)\,dg .
 \tag{CR27}
\] This is a finite sum of genuine source-style bilinear theta lifts, with the already fixed actual finite Schwartz function. It is basis independent because \(\sum_i v_i\otimes p_i\) is the canonical coevaluation tensor. It has not been defined through \({\cal K}_j\). For chart blocks use their already fixed finite Schwartz functions and sum the corresponding integrals.

\subsubsection{Theorem CR28: exact analytification of the constructed arithmetic map}\label{theorem-cr28-exact-analytification-of-the-constructed-arithmetic-map}

For every row \(P\in{\cal H}_j(M)\), every \(v\in W_{j+1}\), and every target adele \(h\), \[
 \boxed{{\cal T}_j(P;v)(h)
      =v_j C^{14hab}
              \left(\prod_\tau s_\tau^{D_\tau}\right)
                         f_{{\cal K}_jP,v}(h).}
 \tag{CR28}
\] This is the mixed arithmetic adapter. Its per-place scale is the literal RealD.sc of K035, and its class map has the exact Kval formula (CR24).

\textbf{Proof.} Write \(g=g_ck\), \(h=h_dl\) on compact real pieces. The finite source level character of the kernel is \(\chi_j^{-1}\), so it cancels (CR26); the target finite character is retained. The real covariance (CR18) sends its archimedean vector to \(\iota_j(\rho_j^\vee(k)p_i,\rho_{j+1}(l)v)\). Consequently the sum over \(i\) in the integrand is contracted by \[
 \sum_i\rho_j(k)v_i\otimes\rho_j^\vee(k)p_i
                         =\sum_i v_i\otimes p_i.
 \tag{CR29}
\] Pairing its first factor with \(P_c\) gives exactly \(P_c\). Thus the sum of integrands is independent of \(k\), and equals the theta functional at \((g_c,h_d)\) applied to \(\iota_j(P_c,\rho_{j+1}(l)v)\).

This polynomial is homogeneous of degree \(D_\tau\) at every label. Formula (CR5) therefore contributes the single common factor \(C^{14hab}\prod_\tau s_\tau^{D_\tau}\). Definitions (CR12), (CR14) and (CR17) identify the remaining value with \[
 B_{j+1}(\Phi_{J_{dc}}P_c,\rho_{j+1}(l)v).
 \tag{CR30}
\] Finally the actual quotient integration formula of DO026 is \[
 \int_{[G_j]} F(g)\,dg
  =v_j\sum_c|\Gamma_c|^{-1}
                                  \int_{G_{j,\infty}}F(g_ck)\,dk .
 \tag{CR31}
\] Its infinite integral is one after (CR29). Substitution of (CR23) gives (CR28). Normal convergence of the fixed-degree Gaussian theta kernels on the compact quotient justifies the integrals and their finite sums, as in DO026. The scalar in (CR28) is nonzero. Genuine theta automorphy also proves the target equivariance asserted after (CR23), without changing its arithmetic definition. \(\square\)

There is no unrecorded factor \(1/\dim W_j\) in (CR28). It cancels because (CR27) uses the complete coevaluation tensor. A single fixed matrix-coefficient realization instead has the Schur factor \(1/\dim W_j\), exactly as in DO026. The complete-type construction has its own cost, computed next.

\subsection{6. Exact norm and lattice bookkeeping}\label{exact-norm-and-lattice-bookkeeping}

Put \(d_j=\dim_\mathbb C W_j\). With an orthonormal compact basis \(v_i\), its pairBK-dual \(p_i\) is also orthonormal under the corresponding dual metric. For every row \(P\), \[
 \boxed{\sum_i\|f_{P,v_i}\|_{L^2([G_j])}^2
          =v_j\sum_c|\Gamma_c|^{-1}\|P_c\|_j^2.}
 \tag{CR32}
\] Indeed the sum of squared matrix coefficients at each \(k\) is \(\|P_c\|^2\), by unitarity; then use (CR31). For one unit \(v\), irreducibility instead gives \[
 \|f_{P,v}\|_{L^2}^2
       =\frac{v_j}{d_j}\sum_c|\Gamma_c|^{-1}\|P_c\|_j^2 .
 \tag{CR33}
\] Equations (CR32)--(CR33) distinguish the Hilbert-Schmidt metric of the complete coefficient realization from a single coefficient embedding.

The norm of the actual archimedean coevaluation vector \[
 {\cal E}_j(v)=
       \sum_i v_i\otimes
          \operatorname{pureTensorO}
                    (\phi_{j,f},\iota_j(p_i,v))
 \tag{CR34}
\] is, with the literal unit-normalized creation operators, \[
 \boxed{\|{\cal E}_j(v)\|^2
     =d_j\left(\prod_\tau H_{\nu_\tau}\right)
             \|\phi_{j,f}\|_f^2
             \|\phi_{\infty,0}\|_\infty^2
             \|v\|_{j+1}^2.}
 \tag{CR35}
\] This follows from CR19 and orthonormality of the \(v_i\). The actual product-host vacuum has \[
             \|\phi_{\infty,0}\|_\infty^2
                              =\det(2Y_J)^{-7hab}.
 \tag{CR36}
\] There are \(14hab\) ordinary Jacobian factors, each with squared Gaussian norm \(\det(2Y_J)^{-1/2}\). K035's change to the literal normalized creation operator preserves the Fischer normalization. The finite norm in (CR35), all hook products, the source dimension, the level volume, and the period in (CR28) remain actual factors. The arithmetic class-row lattice is still (CR22), not the orthonormal lattice used only to compute this norm.

There is also an exact local formulation of the denominator question. At a finite place \(w\) of \(M\), choose actual local bases \(B_{c,\tau}\), \(B'_{d,\tau}\) of the source and target lattices Lsh. The source definitions give \[
 L_{j,c,w}=\rho_j^\vee(B_c)\Delta_{\eta,\mathcal O_w},
 \qquad
 L_{j+1,d,w}=\rho_{j+1}^\vee(B'_d)
                                      \Delta_{\eta',\mathcal O_w}.
 \tag{CR37}
\] All determinant twists in these equalities are retained. The polynomial covariance gives exactly \[
 (\rho_{j+1}^\vee(B'))^{-1}\Phi_J\rho_j^\vee(B)
                         =\Phi_{J^{B,B'}},
 \tag{CR38}
\] where \[
 J^{B,B'}(P)=
 \prod_\tau\det(B_\tau)^{s_{S,\tau}}
                \det(B'_\tau)^{-s_{D,\tau}}\,
 J\!\left(P((B'_\tau)^{\mathsf T}
                                  Z_\tau B_\tau^{-\mathsf T})\right).
 \tag{CR39}
\] For example, (CR39) follows by writing \(\rho^\vee(B)P(U)=\det(B)^{-\lambda_{\rm last}}
P(B^{-1}U)\), absorbing \(q-\lambda_{\rm last}=s_S\) into the determinant insertion, and evaluating the target at \(B'Y\). This also verifies every transpose and sign.

Let \(\mathfrak j_{dc,w}\) be the fractional ideal generated by the finite collection \[
       J_{dc}^{B_c,B'_d}(Z^\alpha),
       \qquad |\alpha_\tau|=D_\tau\quad\text{for every }\tau .
 \tag{CR40}
\] Theorem CR8 gives the concrete bound \[
 \Phi_{J_{dc}}(L_{j,c,w})
                         \subset\mathfrak j_{dc,w}L_{j+1,d,w}.
 \tag{CR41}
\] After the actual stabilizer contraction, \[
 {\cal K}_j({\cal L}_{j,w})
     \subset \pi_w^{\,b_{j,w}}{\cal L}_{j+1,w},
 \qquad
 b_{j,w}=\min_{c,d}
       \left(v_w(\mathfrak j_{dc,w})-v_w(|\Gamma_c|)\right).
 \tag{CR42}
\] Zero ideals have valuation \(+\infty\). The saturated intersection property in (CR22) justifies passing from the ambient coefficient lattices to the equivariant row lattice in (CR42). Empty spaces are treated in the usual way.

This is an inclusion, not an equality for the sharp content. Cancellation in the class sum can increase valuations. Restricted earlier images can increase them further. The actual sharp invariant is obtained from the matrix of \({\cal K}_j:{\cal L}_{j,w}\to{\cal H}_{j+1}(M_w)\) in local integral bases and, for later stages, from its restriction to the actual preceding image. No class projection, arbitrary choice of image lattice, or ambient nonzero vector proves that invariant. Formula (CR42) isolates it using the actual map rather than leaving Kval itself unspecified.

\subsection{7. Composition, terminal compatibility, and the remaining gate}\label{composition-terminal-compatibility-and-the-remaining-gate}

Repeating the construction along the prescribed finite even-rank chain defines actual arithmetic maps \({\cal K}_j\), all over one finite compositum field, on the unchanged class coefficient lattices. At every edge their analytic complete-type theta realizations satisfy (CR28). Iteration therefore retains the explicit product \[
 \prod_j
 \left[v_j C^{14hN_jN_{j+1}}
                  \prod_\tau s_{j,\tau}^{|\nu_{j,\tau}|}\right].
 \tag{CR43}
\] The final arithmetic column is in precisely the rhoVr polynomial convention used by the literal branching operator brC and bilinear pairBK. Its ordinary terminal restriction can be paired with the rhoV terminal row as in DO038. There is no additional transpose scalar. DO038 keeps the exact branch normalization and terminal stabilizer mass; those are not redefined in this construction.

There is a further finite-assembly distinction. A single ordinary ColV row does not automatically become the source's coupled evGSD row by pulling back along tilC: the latter has a prescribed, possibly nontrivial Q character. The actual finite Y/Q expansion can be assembled from the blockwise maps (CR23), with its fixed coefficients. If its actual finite-array covariance is verified, linearity of CR24 and CR28 proves the same arithmetic/analytic comparison for that coupled row, and DO038 supplies its terminal mass identity. No such covariance is inferred from an arbitrary nonzero ordinary theta output. The literal source colC has codomain evGSD.EqFun in \nolinkurl{FiniteSumNonzeroLargeDepth.lean1410–1418,} and PerInv1451--1454 expresses its cancellation, but transferring that particular coupled assembly requires its actual finite kernel hypotheses. The basic theorem here is the ordinary class-row adapter and its blockwise extension, not an unsupported assertion that all mixed terminal arrays already satisfy those hypotheses.

This result supplies an actual replacement arithmetic family and proves its geometric Kval equation and analytification. It does not prove that the old common-quadratic Fields record exists. It also does not infer nonvanishing of every later restricted image, optimal valuation gains, the total DP44 height cost, or an every-epsilon window from the existence of these maps. Those are now questions about the concrete maps (CR23), lattices (CR22), and normalizations (CR28), (CR35), and the specified terminal pairing.

In particular K033/K034's sharp selected first-kernel content cannot automatically be substituted for the content of a later class-contracted image. Conversely an additional ``unknown externally frozen Kval map'' is not a reason to withhold this actual construction: its maps have already been fixed geometrically and their comparison proved.

\subsection{Verification and primary references}\label{verification-and-primary-references}

The companion exact certificate tests the universal substitution on small nontrivial shapes, constructs the corresponding integral Weyl lattices from divided-power generators, checks actual coefficient membership, verifies the matrix covariance (CR39), and checks coevaluation and mass normalizations. These finite tests do not replace the general proofs or assert a Lean build. Author and independent review events are recorded separately.

The representation inputs are the accepted recovered source theorems at the locators above and the primary references already checked for DO025/R011: Fulton, W., \& Harris, J. (1991). Representation theory: A first course. Springer, Theorem 6.3, pp.77--78, Hook Length Formula 4.12, p.50, and Exercise 6.11(b), p.80; and Konno and Konno (2007), section 5.1, Fact 5.1, Lemma 5.2, formulas (5.4)--(5.5), and Lemma 5.3, pp.72--74. The geometric reduced-theta input is the actual algebraic construction and all-order field comparison proved in K035--036, using Bannai and Kobayashi's primary reduced-theta conventions. No general RH extension or unstated extension of an accepted coefficient theorem is used.

\clearpage

\section{AM-DO-038 --- the literal terminal period, its top realization, and exact stabilizer cancellation}\label{am-do-038-the-literal-terminal-period-its-top-realization-and-exact-stabilizer-cancellation}

Author: GPT-6 Astra Pro, ALP\_MACHINE \nolinkurl{/root/definite_oscillator}. 22 November 2026 UTC. Inherited reasoning effort is not separately exposed.

This finding constructs the top automorphic realization from the fixed arithmetic terminal pairing and proves its terminal class-mass equation. The arithmetic tail is the source's \nolinkurl{PtilFF}; it is not defined by dividing an analytic period by a desired scalar. The final arithmetic column can be supplied by the independent geometric mixed-CM operators and class-row construction. The identity applies uniformly to every such column, including zero columns.

The branching operator remains the literal \nolinkurl{brC}, including \nolinkurl{brConst}, \nolinkurl{kapBr}, and the power of the conjugate T-uniformizer. Its closed factorial evaluation is separate. No terminal nonvanishing is inferred from nonvanishing of the last theta lift alone.

\subsection{1. Exact source data and coefficient convention}\label{exact-source-data-and-coefficient-convention}

\begin{longtable}[]{@{}
  >{\raggedright\arraybackslash}p{(\columnwidth - 2\tabcolsep) * \real{0.5000}}
  >{\raggedright\arraybackslash}p{(\columnwidth - 2\tabcolsep) * \real{0.5000}}@{}}
\toprule\noalign{}
\begin{minipage}[b]{\linewidth}\raggedright
Source
\end{minipage} & \begin{minipage}[b]{\linewidth}\raggedright
Content used
\end{minipage} \\
\midrule\noalign{}
\endhead
\bottomrule\noalign{}
\endlastfoot
\nolinkurl{Columns/FiniteSumNonzeroLargeDepth.lean}, 127--152, 201--218 & \nolinkurl{xiBr}, \nolinkurl{kapBr}, \nolinkurl{brConst}, \nolinkurl{brStd}, \nolinkurl{brAt}, \nolinkurl{brC}, \nolinkurl{tilQ}, \nolinkurl{PtilFF} \\
Same file, 220--245, 1410--1418, 1451--1454 & \nolinkurl{PtilC}, the actual \nolinkurl{evRho}/\nolinkurl{evGSD}, coupled \nolinkurl{colC}, and \nolinkurl{PerInv} \\
\nolinkurl{Columns/CoefModB.lean}, 55--66 & Bilinear Fischer \nolinkurl{pairBK}, divided by the extremal Fischer norm \\
\nolinkurl{Columns/ColForms.lean}, 22--37, 53--74; \nolinkurl{CoefCol.lean}, 59--75, 97--99; \nolinkurl{RecLevel.lean}, 50--52 & Inverse level characters, \nolinkurl{rhoVr} columns, \nolinkurl{rhoV} detectors, and the inverse character of the adelic reading \\
\nolinkurl{Columns/RowForms.lean}, 221--229, 253--269; \nolinkurl{GenericUnipotentCharacter.lean}, 407--429, 526--557, 613 & Literal terminal residue groups and their characters \\
\nolinkurl{Columns/RhoEqOfSmulEq.lean}, 134--146, 534--536 & \nolinkurl{csOf} uses stabilizer cardinal inverses \\
\nolinkurl{Columns/Record.lean}, 1558--1610, 1901--1913 & \nolinkurl{lowC_ev}, \nolinkurl{periodC}, antilinear \nolinkurl{jMCL}, \nolinkurl{recPer}, and lowering on arithmetic rows \\
\nolinkurl{KernelLattices/InterlaceAll.lean}, 74--84 & \nolinkurl{lowerK}, \nolinkurl{periodK} \\
\nolinkurl{KernelLattices/KernelLocalFactors/InterlaceDifferenceShiftInvariant.lean}, 334--345 & \nolinkurl{PtilFF_eq}, \nolinkurl{PtilC_eq} \\
\nolinkurl{HermitianForms/LastCoordinateUniformizerVectorSpansOrthogonalLi.lean}, 58--77, 84--104 & \nolinkurl{ClsPTied} and \nolinkurl{ValPTied} class, stabilizer, coefficient, and metric transports \\
\nolinkurl{HermitianForms/FormsTied.lean}, 44--46 & \nolinkurl{PerTie}: \nolinkurl{period(w,G)=ipP(jM G,w)}, with \nolinkurl{jM} antilinear \\
Same file, 23--31; \nolinkurl{Columns/StdLatUnimodular.lean}, 165--178; \nolinkurl{Columns/TiedXpU.lean}, 172--199 & Exact decorated insertion and its rational/finite frame transport \\
\nolinkurl{LatticeFrames/IsRatDatum.lean}, 139--147 & \nolinkurl{TopForm}, including its stabilizer multiplier \\
\nolinkurl{LatticeFrames/TopFormula.lean}, 27--38 & \nolinkurl{IsJ}, \nolinkurl{TopFormula} \\
\nolinkurl{LatticeFrames/StabOrd.lean}, 24--37 & \nolinkurl{volLev}, \nolinkurl{stabOrd}, \nolinkurl{IsReps} \\
\nolinkurl{LatticeFrames/OnFramesRdApply.lean}, 137--150 & \nolinkurl{ClsT.Is.reps}, \nolinkurl{inv}, \nolinkurl{top}, \nolinkurl{tail} \\
\end{longtable}

The newly retrieved HermitianForms definitions and their mirror commit, URLs, and SHA256 hashes are under \nolinkurl{definite_oscillator_source_refs/}. They are not represented as having an unavailable original-author intake hash. Their actual formulas are checked below; there is no Lean execution.

Let the final column rank be N, let n=N−1, and let V\_lambda and V\_sigma be the source polynomial coefficient spaces over C. Fix the source interlacing weights and the exact operator

\[
 B=\operatorname{brC}
 =\left(\prod_\tau
   \tau(\overline{\pi_T})^{-nBr(\lambda_\tau,\sigma_\tau)}\right)
          \operatorname{brStd}:V_\lambda\longrightarrow V_\sigma.
 \tag{DO38.1}
\]

The source's \nolinkurl{lowC_ev} extends precisely this B to frame functions. The source's \nolinkurl{PtilFF_eq} identifies its finite period with \nolinkurl{periodK} after B. These source results are cited at their existing status, not re-proved here. The same file gives \nolinkurl{cBr_ne_zero} at1364, the nonzero interlacing branch at1599--1600 and1828--1831, and \nolinkurl{chain_brC_ne} at1861--1867.

Write

\[
 [P,Q]_\sigma=\operatorname{pairBK}_\sigma(P,Q),\qquad
 \langle P,Q\rangle_\sigma=[P,\overline Q]_\sigma.
 \tag{DO38.2}
\]

The Hermitian form is linear in its first argument. Its normalization is exactly \nolinkurl{nfAll}: ordinary monomial factorials divided by the extremal Fischer norm. Positivity and the source metric comparison are the checked R011 conclusions. In a nonorthonormal algebraic basis the actual Gram matrix remains in (DO38.2). The final column uses the contragredient \nolinkurl{rhoVr} model; the terminal detector uses \nolinkurl{rhoV}, paired by the bilinear form. Complex conjugation identifies their compact Hermitian realizations.

The frame weight \nolinkurl{sgT} is the literal shifted weight \nolinkurl{sgTrue}: \nolinkurl{TiedXpV.lean}, 216--217, proves this equality, and 136--142, 163--167, 201--206 prove that its natural partition and normalized Fischer form equal those used by the unshifted row polynomial. Thus the determinant factor in the frame function is retained even though it does not alter (DO38.2). The actual \nolinkurl{brC} is equivariant under the block embedding with last coordinate one; this is \nolinkurl{branchingProjection_act} in \nolinkurl{InterlaceDifferenceShiftInvariant.lean}, 3146--3158, together with \nolinkurl{PairBInv} in \nolinkurl{CoefModB.lean}, 86--89. The fixed source weights satisfy the stated same-shift hypothesis. No unshifted replacement of the branch is made.

The bilinear invariance can also be checked without treating that predicate declaration as a theorem: the ordinary Fischer pairing is \(P(\partial)Q|_0\). Replacing P by its inverse-transpose linear substitution and Q by the matching direct substitution cancels the two matrices by the chain rule. The determinant twists are reciprocal, and the common extremal norm is a scalar. This proves the displayed \nolinkurl{rhoVr}/\nolinkurl{rhoV} invariance for the actual polynomial spaces.

\subsection{2. Actual classes and equality of their stabilizers}\label{actual-classes-and-equality-of-their-stabilizers}

Let T be the totally definite rank-n unitary group on the actual mixed CM host. Fix the finite row level U and its finite-order unitary character chi. Infinite compact Haar has volume one. Set

\[
 v_U=\operatorname{vol}(U),\quad
 \mathcal C=T(F)\backslash T(\mathbb A_f)/U,\quad
 \Gamma_c=T(F)\cap h_cUh_c^{-1}.
 \tag{DO38.3}
\]

The fixed-rank lattice-class proof of K036 section3 applies at this odd rank as well; only its later determinant-power argument needs even ranks. Finite decorations refine each lattice class by a finite set. The local standardness argument there identifies the actual selfdual lattice genus with adelic unitary moves. A local frame supplies h\_c; its residue frame supplies the stipulated split-place decoration. Changing the frame by U or the rational basis gives exactly the two equivalences in (DO38.3). No global class-number-one assertion is used.

The literal \nolinkurl{ClsPTied.y_bij} expresses this class bijection, and its \nolinkurl{stabI_inj} and \nolinkurl{stabI_range} identify the arithmetic stabilizer with Gamma\_c.~To verify the adelic side explicitly, a rational gamma is counted by \nolinkurl{stabOrd(U,h_c)} exactly when

\[
                       h_c^{-1}\gamma_fh_c\in U.
 \tag{DO38.4}
\]

In \nolinkurl{gamma h_c=h_c k kappa}, the infinite element is forced to be gamma\_infinity and the finite element to be (DO38.4); conversely these choices give the equality. The source range condition is precisely (DO38.4). Thus

\[
 \boxed{\operatorname{stabOrd}(U,h_c)=|\Gamma_c|
       =\operatorname{Nat.card}(\operatorname{stabilizer}_2(c)).}
 \tag{DO38.5}
\]

Also \nolinkurl{volLev=v_U}, by its definition as the measure of the product of the infinite compact group and U. DO026(26.17) then gives the actual class integration coefficient v\_U/\textbar Gamma\_c\textbar. This is equality of the specific stabilizers, not an equality inferred from two total masses.

For the constructed decorated module this identification is direct: an arithmetic stabilizer has a rational isometry preserving the lattice and a residue-frame element in the chosen subgroup. The residue-frame map is surjective, so its residue element is uniquely determined by that rational isometry. Conversely (DO38.4) supplies that element by reduction. These are inverse maps and prove (DO38.5) for the constructed family, rather than assuming an uninhabited \nolinkurl{ClsPTied} record on the new host.

\subsection{3. The arithmetic pairing before its automorphic realization}\label{the-arithmetic-pairing-before-its-automorphic-realization}

Let M be the terminal coefficient-row module and retain its fixed arithmetic pairing

\[
 \mathscr P(W,G)=\sum_{c\in\mathcal C}|\Gamma_c|^{-1}
                                      [W_c,G_c]_\sigma.
 \tag{DO38.6}
\]

This is \nolinkurl{periodK}. The source \nolinkurl{periodC} is this pairing in the chosen row basis; its matrix is not replaced by an identity matrix. Let H\_G,c be the Hermitian Riesz vector of \nolinkurl{W_c -> [W_c,G_c]_sigma}. In the identity frame and real polynomial model,

\[
                         H_{G,c}=\overline{G_c}.
 \tag{DO38.7}
\]

This holds termwise in the Fischer formula, even in a nonorthonormal primitive basis. The polynomial of the source \nolinkurl{eW(jM G)_c} is exactly this Riesz vector. In another orthonormal frame the actual Gram transport gives the equivalent Riesz equation, without a new scalar.

Indeed \nolinkurl{ValPTied.ipP(z,z')} is the sum of \nolinkurl{nfAll(Q,P)}, where P belongs to z and Q to z'. Together with \nolinkurl{PerTie} it gives

\[
 \mathscr P(W,G)=\sum_c|\Gamma_c|^{-1}
                               \langle W_c,H_{G,c}\rangle_\sigma.
 \tag{DO38.8}
\]

The order and the antilinearity of jM make this linear in both W and G. The fixed arithmetic tail is therefore

\[
 \boxed{\mathscr T(\psi,G)=\operatorname{PtilFF}(\psi,G)
   =\sum_c|\Gamma_c|^{-1}
                       [B\psi(d_c),G_c]_\sigma.}
 \tag{DO38.9}
\]

Here, exactly as in the source, \nolinkurl{PtilFF} takes an already restricted function \(\psi:\operatorname{RowPt}\to V_\lambda\). If f is a full final column, put \(f^\downarrow(d)=f(\operatorname{tilC}d)\). Then \nolinkurl{PtilC(f,G)=PtilFF(f^downarrow,G)}. The source \nolinkurl{tilC} is \nolinkurl{tilQ} followed by forgetting the additional Q decorations. A final coupled column \nolinkurl{colC} instead lands directly in \nolinkurl{evGSD.EqFun}; it is already such a restricted function. These two domains are not identified by notation. In either case B and the arithmetic period are fixed before the automorphic realization.

Write \(\lambda_U=\operatorname{rowTy}.\mathrm{chr}\). The actual source actions are

\[
 \psi(gd)=\lambda_U(g_f)\rho^\vee_\lambda(\operatorname{ext}g_\infty)
                  \psi(d),\qquad
 G(gd)=\lambda_U(g_f)^{-1}\rho_\sigma(g_\infty)G(d).
 \tag{DO38.9a}
\]

The first is \nolinkurl{evRho}; the second is \nolinkurl{formRho} for \nolinkurl{rowTy.inv}. Equivariance of B and invariance of \nolinkurl{pairBK} prove that every summand in (DO38.9) is independent of the representative. They also prove that \(W=B\psi\) lies in the opposite coefficient space. This is the literal \nolinkurl{PerInv} equation, with the two finite characters canceled, not a new character convention. In the source adelic reading \nolinkurl{levTyp=(levChr)^-1}; hence both the lowered column and its Riesz family have reading character \(\chi=\lambda_U^{-1}\), and the scalar top automorphic form will have character \(\overline\chi\).

For a full column f, (DO38.9a) requires the actual block level inclusion and character restriction (\nolinkurl{LevExt}, \nolinkurl{Record.lean}, 1501--1519). For the source m=2 terminal groups it holds when \(P.H\le U_r\) and the displayed characters restrict as prescribed; in particular \nolinkurl{rPsym} has this inclusion by \nolinkurl{GenericUnipotentCharacter.lean}, 526--557. For the coupled column it is part of its actual \nolinkurl{evGSD.EqFun} codomain. Mere nonvanishing of a finite Schwartz factor does not establish this equivariance.

\subsubsection{3.1 The decorated insertion is the actual group restriction}\label{the-decorated-insertion-is-the-actual-group-restriction}

This transport can be constructed on the new mixed host. Choose the base final decoration to be the literal insertion of the base row decoration. The source formula \nolinkurl{Deco.til} has underlying lattice \(N\oplus O_Ee_N\), T-functional \(x\mapsto\overline{\pi_T}x_N\), and G/Q residue frame with last coordinate \(x_N/\pi_T\). Thus \(\operatorname{diag}(h_c,1)\) sends the base decoration to \(\operatorname{tilQ}(d_c)\) at every place, preserving these exact last-coordinate normalizations. Here \(\pi_T\) is a unit at the G/Q places, as in the literal source disjointness assumptions.

Let a be the chosen final decorated class containing this insertion. There are a rational unitary \(\gamma_c\) and a finite final-level element \(\kappa_c\) such that

\[
 \gamma_{c,f} h_a=\operatorname{diag}(h_c,1)\kappa_c
                         \quad\text{in }T_N(\mathbb A_f).
 \tag{DO38.9b}
\]

This is the defining double-coset equality, so it does not require class number one or a new existence theorem. The rational decorated move and its residue component are exactly the arithmetic move used to pass from the chosen final representative to \nolinkurl{tilQ(d_c)}. Substitution in \nolinkurl{rdS} gives the actual restricted frame function

\[
 \Phi_{f,c}(y)=\chi_N(\kappa_c)^{-1}
                            f_a(\gamma_{c,\infty}^{-1}y).
 \tag{DO38.9c}
\]

The arithmetic equivariance gives precisely the same expression for the reading of \(f^\downarrow(d_c)\). This is \nolinkurl{EvzTie}, including its inverse finite character, derived from (DO38.9b). The embeddings fix the line generated by \(z_0=\pi_Te_N\), since fixing that line vector is equivalent to fixing \(e_N\). Changing representatives modifies both sides by the same rational and finite actions, already canceled in (DO38.9a). The factor \(\overline\pi_T^{-nBr}\) in B is still present: neither (DO38.9b) nor (DO38.9c) discards it.

Thus \nolinkurl{evzC}, \nolinkurl{lowC_ev}, and \nolinkurl{recPer_lower_emb} apply to this actual canonical geometric reading of an admitted full column. For a coupled final row, the same terminal mass calculation applies once its fixed decorated kernel expansion has this reading and the equivariance (DO38.9a). The source \nolinkurl{colC} has exactly that equivariant codomain, but transplanting a standard local Y-sum to moving middle classes needs its actual covariance and transport proof. Neither a nonzero quotient of the standard local sum nor pointwise CRT alone supplies that proof. An unrelated class map or an arbitrary finite rephasing would not satisfy (DO38.9c).

\subsection{4. A canonical actual top automorphic realization}\label{a-canonical-actual-top-automorphic-realization}

Let V be the compact coefficient representation in the source frame convention, with \texttt{rho(k)Phi(y)=Phi(k\^{}\{-1\}y)}. Use (DO38.2) and let J identify the top compact representation with V. One may take the actual polynomial representation and J the identity; an equivalent compact DSModel is transported by that same intertwiner. This fixes \nolinkurl{IsJ}, rather than freely rescaling each row or each class.

The admitted coefficient families satisfy the literal \nolinkurl{FiC} condition

\[
 \rho(\gamma_\infty)H_c
   =\chi(h_c^{-1}\gamma_fh_c)H_c
                  \quad(\gamma\in\Gamma_c).
 \tag{DO38.10}
\]

For g=gamma h\_c k kappa define

\[
 H(g)=\chi(\kappa)\rho(k^{-1})H_c,
 \qquad e_T(G,v)(g)=\langle Jv,H_G(g)\rangle_\sigma.
 \tag{DO38.11}
\]

The first equation is exactly the \nolinkurl{rdP}/\nolinkurl{rdS} reading, whose frame argument is \texttt{gamma\_infinity\^{}\{-1\}\ g\_infinity\ y}. Any two decompositions of g differ by an element of Gamma\_c.~Substituting (DO38.10) shows that they give the same value. The resulting function is left T(F)-invariant, smooth, and finite under the infinite compact group. It transforms by conjugate(chi) on the right under U. Since T is anisotropic, its quotient is compact and has no proper rational parabolic constant terms; the function is cuspidal.

The formula is linear in G because H\_G is antilinear and occupies the second Hermitian argument. It is linear in v. Unitarity gives

\[
                       e_T(G,v)(gk)=e_T(G,\pi(k)v)(g).
 \tag{DO38.12}
\]

At a finite representative it gives

\[
 \boxed{e_T(G,v)(h_c)
   =\langle Jv,H_{G,c}\rangle_\sigma
   =|\Gamma_c|\mathscr P(\operatorname{single}_c(Jv),G).}
 \tag{DO38.13}
\]

This is the literal \nolinkurl{TopForm}/\nolinkurl{TopFormula}, with its forced stabilizer multiplier. The single-class family here belongs to the ambient space of coefficient families; it need not itself satisfy the stabilizer constraint. The source period is defined on that ambient space. Equivalently one may first average Jv over the finite stabilizer with its prescribed character: pairing with H\_G,c gives the same value, since H\_G,c has exactly that character. For the full coefficient-row module, (DO38.11) gives the full top multiplicity space at this level: evaluating any intertwining automorphic embedding at each h\_c gives a functional on V, whose Riesz vector satisfies (DO38.10). Riesz duality recovers the unique coefficient family. A separately restricted arithmetic subspace need not be onto that full multiplicity space.

The actual norm, for the displayed J, is

\[
 \boxed{\|e_T(G,v)\|^2
   =\frac{v_U}{\dim V}\,\|Jv\|^2
                       \sum_c|\Gamma_c|^{-1}\|H_{G,c}\|^2.}
 \tag{DO38.14}
\]

This is the checked class and compact-Schur formula of DO026(26.17)--(26.18). Thus this realization is not declared isometric to an independently normalized source row. Its level volume and Schur dimension remain explicit; they are not class-dependent constants in (DO38.13).

\subsubsection{4.1 Which side the terminal branch acts on}\label{which-side-the-terminal-branch-acts-on}

For clarity, let \(P_c\) be the finite-class coefficient of a full column, and write its complete frame function as \(\Phi_{P_c}(y)=\operatorname{fnOf}(P_c)(y)\), including the source determinant shift. At a finite class representative and an infinite compact element k, the literal reading is

\[
       \operatorname{rd}(h_c k,y)=\Phi_{P_c}(k y),\qquad
       \operatorname{lowC}(\Phi_{P_c})
                         =\operatorname{ev}_\sigma(BP_c).
 \tag{DO38.14a}
\]

The second equality acts on the whole frame function before a value of y is chosen. It therefore branches the class coefficient P\_c. A single highest-weight Schwartz kernel instead fixes one right compact seed in the scalar automorphic output. Choosing the identity frame can make that particular scalar function restrict only to the top H-type. This does not imply that P\_c is a highest polynomial: in a matrix coefficient, the class coefficient and the fixed right seed occupy the two different factors. Right translations vary the seed while retaining the class coefficient.

For example, in the elementary compact representation Sym\^{}2(C\^{}2), let the right seed be e\_1\^{}2 and the class functional be the coefficient of e\_2\^{}2. The scalar matrix coefficient on diag(U(1),1) is identically zero, although that class functional has a nonzero restriction to the e\_2\^{}2 branch; a right rotation of the seed detects it. This example only illustrates the distinction. No such new rotation is inserted in the fixed arithmetic operator. Its actual test is BP\_c after the literal decorated transport (DO38.9c).

In the source \nolinkurl{PerX}, the proportionality factor a(y,v) is permitted to vanish for particular frames or compact vectors. Consequently a vanishing period at a single highest seed neither proves nor disproves the arithmetic tail. A complete-type Kval realization is what recovers P\_c and makes the test in (DO38.14a) meaningful.

\subsection{5. The exact terminal class-mass equation}\label{the-exact-terminal-class-mass-equation}

Choose the canonical terminal finite trace \nolinkurl{delta_1} and one terminal chart. Define the class functional from the fixed period by

\[
 \ell_{c,G}(\Phi)=|\Gamma_c|
              \mathscr P(\operatorname{single}_c\Phi,G).
 \tag{DO38.15}
\]

For an arbitrary frame function, (DO38.15) uses the literal \nolinkurl{periodC} extension through \nolinkurl{Ev.extend}; it agrees with the fixed bilinear polynomial pairing on V by \nolinkurl{periodC_eq_ipCl}. Thus it is linear on the full domain required by \nolinkurl{ClsT.ell}; on V it equals \nolinkurl{<Phi,H_G,c>_sigma}. Equations (DO38.12)--(DO38.13) give \nolinkurl{ClsT.Is.top}. The opposite final-column and top finite characters cancel in their product, giving \nolinkurl{ClsT.Is.inv}; section2 gives \nolinkurl{reps}.

Let Phi\_psi,c be the final column restriction at h\_c, read by the stipulated decorated insertion and frame map. The literal \nolinkurl{lowC} transport says, for every coefficient row,

\[
 \operatorname{lowC}(\Phi_{\psi,c})
      =\operatorname{ev}_\sigma(B\psi(d_c)).
 \tag{DO38.16}
\]

This is the source \nolinkurl{lowC_ev}/\nolinkurl{recPer_lower_emb} identity on the entire polynomial space. In the new family it is implemented by those same decorated insertions and readings. It would not follow for an unrelated class map merely from a geometric theta identity.

Substituting (DO38.5), (DO38.15), and (DO38.16) into the literal terminal class-mass sum gives

\[
 \begin{split}
 v_U\sum_c|\Gamma_c|^{-1}
             \ell_{c,G}(\operatorname{lowC}\Phi_{\psi,c})
 &=v_U\sum_c|\Gamma_c|^{-1}
                     [B\psi(d_c),G_c]_\sigma\\
 &=v_U\operatorname{PtilFF}(\psi,G).
 \end{split}
 \tag{DO38.17}
\]

Therefore, uniformly in both rows and in every class,

\[
                         \boxed{c_P=v_U=\operatorname{volLev}(U)>0.}
 \tag{DO38.18}
\]

The coefficient Gram, hooks, \nolinkurl{kapBr}, and T-uniformizer powers remain inside the fixed bilinear pairing and B on both sides. None has been normalized away. For maximal finite compact volumes one,

\[
                 c_P=\prod_{v\text{ carrying a row level}}
                                         [K_v:U_v]^{-1}.
 \tag{DO38.19}
\]

At a split residue level H\_v in GL\_n(k\_v), its factor is \nolinkurl{|H_v|/|GL_n(k_v)|}, since reduction is surjective and its principal kernel lies in U\_v. The value is thus an explicit positive rational index product, without enumerating individual global class stabilizers. A different Haar choice rescales both sides by that same factor.

\subsubsection{5.1 Literal G/Q volumes and their denominators}\label{literal-gq-volumes-and-their-denominators}

Put \(r=k+m-1\) and \(n=r+mt=N_m-1\). The top row has T-parameter one and no T-decoration. Its G-level is the full preimage of \(H_G=\operatorname{HPar}(n;r,t,\ldots,t;P.H)\), and its Q-level is the full preimage of \(H_Q=P_{r,mt}\). These are the groups in \nolinkurl{rowTyG} and \nolinkurl{rowTyQ}, not the kernels of their characters. Their orders are

\[
 \begin{split}
 |H_G(\mathbf F_q)|
  &=|P.H(\mathbf F_q)|\,|\operatorname{GL}_t(\mathbf F_q)|^m
                  q^{rmt+\binom m2t^2},\\
 |H_Q(\mathbf F_q)|
  &=|\operatorname{GL}_r(\mathbf F_q)|
    |\operatorname{GL}_{mt}(\mathbf F_q)|q^{rmt}.
 \end{split}
 \tag{DO38.19a}
\]

To prove the first formula, choose the leading diagonal block in P.H, each of the m remaining diagonal blocks in GL\_t, and all upper block entries freely. Their dimensions are rmt and binom(m,2)t\^{}2. All lower blocks are zero, and block triangularity is exactly invertibility of the displayed diagonal blocks. The second formula is the same argument with two unrestricted diagonal blocks. Consequently

\[
 \boxed{c_P=
 \prod_{v\in G}
   \frac{|P.H(\mathbf F_{q_v})|\,|GL_t(\mathbf F_{q_v})|^m
             q_v^{rmt+\binom m2t^2}}{|GL_n(\mathbf F_{q_v})|}
 \prod_{v\in Q}
   \frac{|GL_r(\mathbf F_{q_v})|\,|GL_{mt}(\mathbf F_{q_v})|
             q_v^{rmt}}{|GL_n(\mathbf F_{q_v})|}.}
 \tag{DO38.19b}
\]

Here every actual real-field place is counted once; the selection of one CM prime above it is only the split identification. If an additional finite refinement is imposed, multiply by its index inverse. The displayed canonical choice has maximal compact at 29; no principal 29 refinement is introduced to make the class formulas work. The exact logarithmic denominator cost of c\_P is \(\sum_v\log[K_v:U_v]\), with possible cancellation when combined with other independently specified constants left visible.

For q prime to 11, let e=ord\_11(q) and c=v\_11(q\^{}e−1). Directly from \(|GL_j(\mathbf F_q)|=q^{j(j-1)/2}\prod_{i=1}^j(q^i-1)\) and the odd-prime lifting-of-exponent identity,

\[
 v_{11}|GL_j(\mathbf F_q)|
   =\lfloor j/e\rfloor c+v_{11}(\lfloor j/e\rfloor!).
 \tag{DO38.19c}
\]

The identity follows because 11 divides q\^{}i−1 exactly when e divides i, and then v\_11(q\^{}(e s)−1)=c+v\_11(s), obtained by expanding successive 11th powers and then a power prime to 11. For the source \nolinkurl{rPsym}, P.H is a subgroup of the finite upper unipotent group, so its order is a power of q. For m=2, r=k+1 and rS=0 it consists of U\_k with the last column fixed, and its exact order is q\^{}(k(k−1)/2).

\begin{longtable}[]{@{}
  >{\raggedright\arraybackslash}p{(\columnwidth - 10\tabcolsep) * \real{0.1364}}
  >{\raggedright\arraybackslash}p{(\columnwidth - 10\tabcolsep) * \real{0.1364}}
  >{\raggedleft\arraybackslash}p{(\columnwidth - 10\tabcolsep) * \real{0.1818}}
  >{\raggedleft\arraybackslash}p{(\columnwidth - 10\tabcolsep) * \real{0.1818}}
  >{\raggedleft\arraybackslash}p{(\columnwidth - 10\tabcolsep) * \real{0.1818}}
  >{\raggedleft\arraybackslash}p{(\columnwidth - 10\tabcolsep) * \real{0.1818}}@{}}
\toprule\noalign{}
\begin{minipage}[b]{\linewidth}\raggedright
Host
\end{minipage} & \begin{minipage}[b]{\linewidth}\raggedright
k, t, n, number f of selected places per G prime
\end{minipage} & \begin{minipage}[b]{\linewidth}\raggedleft
G prime q
\end{minipage} & \begin{minipage}[b]{\linewidth}\raggedleft
v11(
\end{minipage} & \begin{minipage}[b]{\linewidth}\raggedleft
GL\_n
\end{minipage} & \begin{minipage}[b]{\linewidth}\raggedleft
)
\end{minipage} \\
\midrule\noalign{}
\endhead
\bottomrule\noalign{}
\endlastfoot
K031 & 42, 1759, 3561, 42 & 4615409 & 390 & 191 & −8 \\
K031 & same & 5276377 & 390 & 191 & −8 \\
K032 & 168, 1633, 3435, 168 & 22285932761 & 754 & 357 & −40 \\
K032 & same & 64072175441 & 1888 & 896 & −96 \\
\end{longtable}

Thus the exact G contribution is v11(c\_P,G)=−672 on K031 and −22848 on K032 for this row group. The Q factors are the inverse Gaussian binomial indices in (DO38.19b) and must still be included at their actual primes. These numbers concern the terminal volume only; they are not the cost of the branch, the theta kernel, or an optimal integral terminal detector.

\subsubsection{5.2 The full terminal period, including its real scalar}\label{the-full-terminal-period-including-its-real-scalar}

There is also an exact real normalization, useful for distinguishing the arithmetic branch from a single highest seed. Let B* be the Hermitian adjoint of the fixed B for the two source normalized Fischer metrics, and put \(d_\sigma=\dim V_\sigma\). The compact branch has multiplicity one. Since B is a nonzero intertwiner, Schur's lemma and positivity give

\[
 BB^*=\beta_B I,\qquad
 \Pi_\sigma=\beta_B^{-1}B^*B,\qquad
 \beta_B=\|B\|^2>0.
 \tag{DO38.19d}
\]

The scalar is explicit in root R019(19.15):

\[
 \boxed{\beta_B=\prod_\tau
  |\tau(\overline\pi_T)|^{-2nBr_\tau}
                    |c_{{\rm Br},\tau}|^2\kappa_\tau.}
 \tag{DO38.19e}
\]

For a source frame y let \(E_y(P)=\Phi_P(y)\). This is a fixed linear evaluation functional on the complete coefficient space, with its determinant factor included. Then the actual period for an admitted full-column reading is

\[
 \boxed{\int_{[T]} e_T(G,v)(h)\,
                  \operatorname{rd}(f,y)(\iota h)\,dh
       =v_U\,\frac{E_y(B^*Jv)}{\beta_Bd_\sigma}
                         \operatorname{PtilC}(f,G).}
 \tag{DO38.19f}
\]

For an actually assembled coupled reading, replace PtilC(f,G) by PtilFF(psi,G) under precisely the reading condition of section3.1. Thus the source PerX scalar is \(a(y,v)=E_y(B^*Jv)/(\beta_Bd_\sigma)\); c\_P remains v\_U.

\textbf{Proof.} Work on one finite class after its fixed decorated transport, and write P for that class coefficient and \(\ell_G(W)=[W,G_c]_\sigma\). By (DO38.11)--(DO38.12), the two compact factors are

\[
   \ell_G(\rho_\sigma(k)Jv),\qquad
              E_y(\rho_\lambda(\iota k)^{-1}P).
 \tag{DO38.19g}
\]

The integral kills the H-types other than sigma. On its single sigma summand substitute \(\Pi_\sigma P=B^*BP/\beta_B\). B* intertwines, so the second factor becomes \(\beta_B^{-1}(E_y\circ B^*)(\rho_\sigma(k)^{-1}BP)\). The bilinear compact matrix-coefficient identity is

\[
 \int_{T_\infty}\ell_1(\rho(k)v)\ell_2(\rho(k)^{-1}w)\,dk
             =d_\sigma^{-1}\ell_1(w)\ell_2(v).
 \tag{DO38.19h}
\]

It follows by averaging the rank-one endomorphism \(v\otimes\ell_2\): its average commutes with the irreducible representation, and its trace is \(\ell_2(v)\). Applying the average to w and then \(\ell_1\) proves (DO38.19h). Thus the one-class integral is \([BP,G_c]_\sigma E_y(B^*Jv)/(\beta_Bd_\sigma)\). Multiplication by the exact class measure v\_U/\textbar Gamma\_c\textbar{} and summation gives (DO38.19f). No complex conjugation of the bilinear period is inserted. The conjugation is already in the adjoint and Riesz vectors.

If v is nonzero then B\emph{Jv is a nonzero polynomial frame function, so some frame y has nonzero E\_y(B}Jv). Even unitary frames suffice: their Zariski density follows by complexifying the Lie algebra, or equivalently by the fact that the compact group's matrix coefficients separate its irreducible vectors. This proves that the proportionality factor is not identically zero, while allowing it to vanish at an individual highest frame. It does not prove that the particular arithmetic coefficient family has BP\_c nonzero.

There is a quantitative archimedean normalization. Let \(d_\lambda=\dim V_\lambda\), take \(\|Jv\|=1\), and use the standard unitary frame orbit. Evaluation at the identity frame is the unit highest-coefficient functional: all other weight spaces vanish there, because the column weights are fixed and only diagonal monomials survive; the normalized highest polynomial has value and norm one. Hence every E\_y on this compact orbit has norm one. Compact Schur orthogonality on V\_lambda gives

\[
 \int_{U(N)_\infty}|E_y(B^*Jv)|^2\,dy
        =\frac{\beta_B}{d_\lambda}.
 \tag{DO38.19i}
\]

Some unitary frame therefore satisfies \(|a(y,v)|\ge1/(d_\sigma\sqrt{\beta_Bd_\lambda})\). Rescaling v by a(y,v)\^{}−1 makes the factor a exactly one, with norm at most \(d_\sigma\sqrt{\beta_Bd_\lambda}\) for the canonical isometric J. The resulting actual period is exactly v\_U PtilC(f,G). This is an archimedean bound and a normalization of a freely chosen compact test vector. It neither changes the fixed arithmetic branch nor supplies a height bound for an algebraic frame or an integral detector.

\subsection{6. Exact nonzero values and retained finite-projector scope}\label{exact-nonzero-values-and-retained-finite-projector-scope}

Suppose the lowered family

\[
                         W_c=B\psi(d_c)
 \tag{DO38.20}
\]

is nonzero and belongs to the admitted opposite coefficient space. The full canonical terminal module contains its Riesz dual, given in the identity frame by G\_c=conjugate(W\_c). For that actual row,

\[
 \boxed{\operatorname{PtilFF}(\psi,G)
           =\sum_c|\Gamma_c|^{-1}\|W_c\|_\sigma^2>0.}
 \tag{DO38.21}
\]

The terminal class-mass value is c\_P times this number. If W is algebraic over a conjugation-stable number field, G is algebraic over that same field. Its integral denominator still depends on its actual coefficient lattice. A smaller prescribed detector space need not contain G. Nonzero final theta output does not imply (DO38.20) nonzero: the fixed branch and fixed decorated insertion can kill it. That exact coefficient test remains; generic branching or arbitrary rational rotation is not substituted for it.

For a nonidentity prescribed finite trace, retain its arithmetic coefficient operator. If e preserves the admitted row module and its exact action is E, then \nolinkurl{trc(e)e_T(G,v)=e_T(EG,v)}; the same proof replaces G by EG in (DO38.17). This follows by applying the actual class reading to every finite translation, including its character and rational frame factors. Several terminal charts give the corresponding sum of coefficient actions. Equality with an independently fixed, unmodified \nolinkurl{PtilFF} requires that additional operator equality. No such nonidentity-projector equality is asserted here.

\subsection{7. Consequence for the canonical two-step family}\label{consequence-for-the-canonical-two-step-family}

For identity terminal trace the literal m=2 top factor becomes

\[
                  c_{\mathrm{Top}}=c_{V,0}c_{K,0}c_{K,1}\,v_U.
 \tag{DO38.22}
\]

The kernel periods and creation frames come from the independent DO026/K035/K036 family and its canonical class-row operator. Equation (DO38.17) supplies the terminal finite mass tie for its admitted ordinary columns read by the literal decorated insertion and brC, and for any coupled final row with the actual reading and covariance specified in section3.1. The top realization has been constructed from the fixed arithmetic pairing, not left as an arbitrary external normalization. An ordinary output whose pullback fails the Q character condition is not silently treated as such a coupled row.

To invoke MainX one must verify its actual \nolinkurl{LawsX} identities for this same family, including the geometric arithmetic column maps and any coupled finite assembly. MainX itself is a vanishing implication and does not require a separately proved nonzero lowered family. Its contrapositive would give a nonzero prescribed tail from a nonzero bottom period of the stipulated descended top input. Alternatively, (DO38.20)--(DO38.21) give a direct terminal detector if that lowered family is known to be nonzero. The integral-image and cost estimates remain additional arithmetic questions in either route. The old common-quadratic field wrapper is not inhabited. This is a mathematical adapter for the specified mixed geometry and coefficient rows, with its finite trace explicitly fixed.

\subsection{Verification}\label{verification}

The companion finite certificate checks unequal stabilizer masses, nonidentity Gram matrices, the bilinear/Hermitian Riesz convention, a nontrivial small branching polynomial, the exact real Schur scalar with a nonidentity source metric, and finite-index normalization. It passes 144 unequal-mass identities, 432 Riesz identities, 144 positive full-detector tests, 64 real Schur tests, six enumerated finite-group cases, and 814 direct-versus-LTE valuation checks, as well as the literal branching and zero-seed examples. It does not enumerate the actual large-rank class set or establish nonvanishing of (DO38.20) for the selected theta image. Author and independent whole-proof events are recorded separately. No Lean is run.

\clearpage

\section{AM-R-020 --- the whole integral terminal branch and a uniform content bound}\label{am-r-020-the-whole-integral-terminal-branch-and-a-uniform-content-bound}

Author: GPT-6 Astra Pro, ALP\_MACHINE /root. 22 November 2026 UTC. Reasoning effort is not externally exposed.

R019 evaluated the literal terminal branch on its distinguished polynomial highest line and computed its exact complex norm. That polynomial line is not, in general, the complete integral image of the branch. This finding bounds the branch on the actual standard Weyl lattice, without identifying it with that line or assuming that its image is a scalar multiple of the target Weyl lattice. The bound is uniform in the growing first weight of the two concrete hosts. It concerns the whole branch operator; its restriction to an earlier theta image is explicitly excluded.

\subsection{1. Fixed conventions and accepted inputs}\label{fixed-conventions-and-accepted-inputs}

First remove the common final weight from the interlacing pair, so that we have partitions lambda of length N=n+1 and mu of length n with lambda\_N=mu\_n=0. The general same-shift pair is obtained by the common determinant twist; it has exactly these polynomial lattices. The literal T-uniformizer exponent is retained separately below and is not changed by this notational reduction.

Write Delta\_lambda and Nabla\_lambda for the integral Weyl and dual-Weyl polynomial lattices in V\_lambda over Q, and similarly at mu. Use the source symmetric bilinear normalized Fischer forms

\[
 [P,Q]_\lambda=\frac{\operatorname{Fischer}(P,Q)}{H_\lambda},
 \qquad H_\lambda=\|p_\lambda\|_{\rm Fischer}^2,
 \tag{R20.1}
\]

where p\_lambda is the primitive extremal, of normalized norm one. These are the exact \nolinkurl{pairBK} conventions of K037 and DO038. The following accepted results are used, not proved anew:

\begin{itemize}
\tightlist
\item
  \nolinkurl{Formalism/LocalRingCardinality/FactDelDualNabla.lean}, theorem \nolinkurl{factDelDualNabla}, line75 in the recovered copy; its statement is \nolinkurl{DataCoeff.lean}126--129. Delta is the full dual lattice of Nabla under (R20.1). Biduality therefore makes Nabla the full dual of Delta.
\item
  \nolinkurl{Formalism/GradedWeylDivisibility/StatementA.lean}34--38 and68--80: Delta is generated from p\_lambda by integral row divided-power words; Nabla is the column-equation lattice, and these row operations preserve it. Its row-weight decomposition is an integral direct sum.
\item
  The multiplicity-one characteristic-zero branching and the literal highest embedding are those used in R019. In particular there is a rational map i:V\_mu to V\_lambda with i(p\_mu)=xiBr. The polynomial xiBr is integral, belongs to Nabla\_lambda, and its last-row coefficient at bBr is exactly p\_mu. R019 proves the required normalization, including its positive rational normalized squared norm kappa.
\item
  \nolinkurl{Columns/FiniteSumNonzeroLargeDepth.lean}118--141 defines B=brStd=cBr\emph{kappa}projVV\emph{coefLast. The accepted branch equivariance is \nolinkurl{InterlaceDifferenceShiftInvariant.lean}3146--3158; its normalized adjoint identity is \nolinkurl{brTranspose}, 304--314, or \nolinkurl{inD_adj},3870--3883. R019 gives B(xiBr)=cBr}kappa\emph{p\_mu and \textbar\textbar B\textbar\textbar\^{}2=\textbar cBr\textbar\^{}2}kappa. None of these identifies an integral image.
\end{itemize}

All integer lattices in this note are the displayed lattices. We do not replace Delta by Nabla when they have different index.

\subsection{2. Exact lattice inclusions for the whole branch}\label{exact-lattice-inclusions-for-the-whole-branch}

Let c=cBr, which is a nonzero integer. Then

\[
 \boxed{\quad
 c\Delta_\mu\ \subset\ B(\Delta_\lambda)
       \ \subset\ c\nabla_\mu
       \ \subset\ \frac{c}{H_\mu}\Delta_\mu .\quad}
 \tag{R20.2}
\]

The inclusions are inclusions of fractional lattices in V\_mu; the sign of c has no effect on them. In particular the whole branch image need not equal c times Delta\_mu.

\textbf{Proof.} Let i\^{}dagger denote the adjoint of i for (R20.1). All coefficients are rational, so this is also the Hermitian adjoint for the corresponding compact metrics. Branching multiplicity one and the normalization on p\_mu give

\[
 i^\dagger i=\kappa I,\qquad B=c\,i^\dagger .
 \tag{R20.3}
\]

Indeed the first equality follows from invariance and the squared norm of i(p\_mu). Both sides of the second equality are the same equivariant projection up to a scalar; applying them to xiBr determines that scalar, since both values are c\emph{kappa}p\_mu. Equivalently (R20.3) is the literal normalized transpose identity cited above with the highest embedding divided by c.

Every integral divided-power word in the rank-n row operators commutes with i and preserves Nabla\_lambda. Applying the word to p\_mu and to xiBr proves

\[
                         i(\Delta_\mu)\subset\nabla_\lambda.
 \tag{R20.4}
\]

Consequently, for v in Delta\_lambda and w in Delta\_mu, {[}i\^{}dagger v,w{]}\_mu={[}v,iw{]}\_lambda is integral. Perfect duality makes i\^{}dagger v an element of Nabla\_mu. This proves the middle inclusion in (R20.2).

For the first inclusion, xiBr is a primitive element of Nabla\_lambda. In fact coefLast(xiBr)=p\_mu, and p\_mu has an integer monomial coefficient one. Thus xiBr also has an integer coefficient one; that coefficient functional on Nabla\_lambda takes value one at xiBr. Perfect duality supplies v in Delta\_lambda satisfying

\[
                              [v,\xi_{\rm Br}]_\lambda=1.
 \tag{R20.5}
\]

Replace v by its integral row-weight component having the weight of xiBr. This component remains in Delta: its divided-power generators are row-weight homogeneous, so the generator span is the direct sum of its row-weight spans. Distinct row weights are Fischer orthogonal, so (R20.5) is preserved. The first n entries of that weight are mu. The image i\^{}dagger v is therefore in the highest weight line Q*p\_mu of V\_mu. Pairing it with p\_mu and using (R20.5) shows that it equals p\_mu. Thus p\_mu belongs to i\^{}dagger(Delta\_lambda).

The latter image is stable under all integral rank-n row divided powers. One way to check this in the polynomial convention is to transpose an intertwining row divided-power identity for i: the adjoint of the (r,s) divided power is the (s,r) divided power for the normalized Fischer form. Both source Weyl lattices are stable under these operators. The divided-power generation of Delta\_mu now gives Delta\_mu contained in i\^{}dagger(Delta\_lambda), proving the first inclusion after multiplication by c.

Finally any two elements of Nabla\_mu are integer polynomials. Their ordinary Fischer pairing is an integer. Equation (R20.1) gives {[}H\_mu*u,w{]}\_mu in Z for all u,w in Nabla\_mu. Perfect duality then gives

\[
                          H_\mu\nabla_\mu\subset\Delta_\mu,
 \tag{R20.6}
\]

which proves the last inclusion. This proof uses no saturation of Delta in the ambient polynomial module. The integral inclusions remain valid after tensoring with the integers of a number field and localizing at any finite place. QED.

\subsection{3. Exact finite-place formulation, tensor labels and frames}\label{exact-finite-place-formulation-tensor-labels-and-frames}

Let O\_w be a DVR, with valuation v\_w, and put I\_w(b)=\{x in Frac(O\_w):v\_w(x) at least b\} for b in its valuation group. Write content\_w(B) for the largest such b for which B(Delta\_lambda,O\_w) is contained in I\_w(b)*Delta\_mu,O\_w. This convention also permits the later normalization v\_w(11)=1 without writing a fractional power of a uniformizer. The image is nonzero, and (R20.2) implies

\[
 v_w(c)-v_w(H_\mu)
                \ \le\ \operatorname{content}_w(B)\ \le\ v_w(c).
 \tag{R20.7}
\]

The left inequality is the image inclusion in the largest displayed lattice. The right inequality follows because c*p\_mu is an actual image and p\_mu is primitive in Delta\_mu. Notice that the inequalities have opposite uses; a lower divisibility certificate is not reversed into the upper inequality.

For several CM labels the full coefficient spaces and maps are tensor products. Put

\[
 C=\prod_\tau c_\tau,\quad H=\prod_\tau H_{\mu_\tau},\quad
 u_T=\prod_\tau\tau(\overline\pi_T)^{-nBr_\tau}.
 \tag{R20.8}
\]

Tensoring (R20.2), and then applying the literal scalar u\_T, gives

\[
 u_TC\Delta_{\boldsymbol\mu}
 \subset \operatorname{brAt}(\Delta_{\boldsymbol\lambda})
 \subset u_TC\nabla_{\boldsymbol\mu}
 \subset \frac{u_TC}{H}\Delta_{\boldsymbol\mu}.
 \tag{R20.9}
\]

Thus (R20.7) has v\_w(u\_T C) and v\_w(H) in this setting. In a finite frame whose source and target changes are integral GL matrices with integral inverses, the same bounds hold, since the Weyl lattices are stable under those changes. A nonintegral frame has its exact source and target lattice factors, as in K037(CR37)--(CR39); they are not silently suppressed. This statement is about standard coefficient lattices and their indicated transports, before stabilizer invariants or restriction to a prior image.

\subsection{4. Evaluation on the actual growing J14 and J56 weights}\label{evaluation-on-the-actual-growing-j14-and-j56-weights}

Use precisely R019's lambda and mu lists, and increase only lambda\_1 at the two designated CM labels by A=11\^{}r and A=121*11\^{}r. The target mu, and hence its natural partition and hook product, stay fixed. R019 proves that both cBr factors at the growing labels have reached their constant saturated values for every r at least13. All other labels have their unchanged baseline values. The T primes of R017 are distinct from11, so every factor in u\_T is an11-unit.

For eta\_i=mu\_i-mu\_n and n=N-1, the exact hook valuation used here is

\[
 v_{11}(H_\eta)=
 \sum_{i=1}^{n}v_{11}((\eta_i+n-i)!)
       -\sum_{1\le i<j\le n}
                        v_{11}(\eta_i-\eta_j+j-i).
 \tag{R20.10}
\]

This is the standard hook-product factorial identity, used with the same normalization as R011. Legendre's floor sum evaluates its first sum; the second sum has finitely many explicit integer factors. The script retains the actual literal RowPar check of the mu list from R019 before evaluating it. The results are:

\begin{longtable}[]{@{}
  >{\raggedright\arraybackslash}p{(\columnwidth - 8\tabcolsep) * \real{0.1579}}
  >{\raggedleft\arraybackslash}p{(\columnwidth - 8\tabcolsep) * \real{0.2105}}
  >{\raggedleft\arraybackslash}p{(\columnwidth - 8\tabcolsep) * \real{0.2105}}
  >{\raggedleft\arraybackslash}p{(\columnwidth - 8\tabcolsep) * \real{0.2105}}
  >{\raggedleft\arraybackslash}p{(\columnwidth - 8\tabcolsep) * \real{0.2105}}@{}}
\toprule\noalign{}
\begin{minipage}[b]{\linewidth}\raggedright
Host
\end{minipage} & \begin{minipage}[b]{\linewidth}\raggedleft
n
\end{minipage} & \begin{minipage}[b]{\linewidth}\raggedleft
Number of CM labels
\end{minipage} & \begin{minipage}[b]{\linewidth}\raggedleft
v11(H\_mu), one label
\end{minipage} & \begin{minipage}[b]{\linewidth}\raggedleft
v11(C), full tensor
\end{minipage} \\
\midrule\noalign{}
\endhead
\bottomrule\noalign{}
\endlastfoot
J14 & 3561 & 42 & 40694631292531 & 61490184399432942 \\
J56 & 3435 & 168 & 144279902901612 & 802279065217743524 \\
\end{longtable}

Therefore the content of the \textbf{whole} literal branch on the standard Weyl lattice satisfies

\[
\begin{array}{ll}
\text{J14:}&
 59781009885146640\ \le\operatorname{content}_{11}(\operatorname{brAt})
                                  \le61490184399432942,\\[2mm]
\text{J56:}&
 778040041530272708\ \le\operatorname{content}_{11}(\operatorname{brAt})
                                  \le802279065217743524.
\end{array}
 \tag{R20.11}
\]

Valuations here are normalized by v(11)=1. The bounds hold at each coefficient-field place above11 with the standard integral tensor model. They are independent of r. Their intersection with the already certified actual fixed-T families is r=13 through55 for J14 and r=15 through90 for the K033 J56 growing witnesses; the branch assertion itself holds for all r at least13. No enlargement of the R018 size intervals follows.

This proves that the normalization/content uncertainty of the whole terminal branching map is bounded independently of the growing first weight. It does not prove that the branch restricted to the actual preceding theta image has bounded content: that image can enter a deeper congruence sublattice or even the branch kernel.

\subsection{5. A complete small example with unequal Smith factors}\label{a-complete-small-example-with-unequal-smith-factors}

Take lambda=(2,1,0), mu=(2,0), and write x\_i and y\_i for the first two entries of row i in a three-row matrix. Then

\[
 p_\lambda=x_1(x_1y_2-x_2y_1),\quad
 \xi=x_1(x_1y_3-x_3y_1),\quad c=4,\quad\kappa=1.
 \tag{R20.12}
\]

The target is Sym\^{}2 in x\_1,x\_2, so projVV is the identity after extracting the coefficient of y\_3 with x\_3-degree zero. Consequently B is four times that extraction. Every source Weyl polynomial is integral, hence

\[
 B(\Delta_{(2,1,0)})\subset
       4\mathbb Zx_1^2+4\mathbb Zx_1x_2+4\mathbb Zx_2^2.
 \tag{R20.13}
\]

Let D\_rs=sum\_j X\_rj partial/partial X\_sj, with divided powers as in the source. The following three source vectors belong to Delta\_lambda:

\[
 \xi=D_{32}p_\lambda,\qquad
 \xi'=D_{21}^{(2)}\xi=x_2(x_2y_3-x_3y_2),\qquad
 v=D_{31}p_\lambda
       =2x_1x_3y_2-x_2x_3y_1-x_1x_2y_3.
 \tag{R20.14}
\]

Their images are 4x\_1\^{}2,4x\_2\^{}2,-4x\_1x\_2, respectively. Thus (R20.13) is equality: the image is 4*Nabla\_(2,0). The integral target Weyl basis is (x\_1\^{}2,2x\_1x\_2,x\_2\^{}2), directly from the divided powers of x\_1\^{}2. The image matrix is diagonal (4,2,4), up to ordering and a sign. Its Smith factors are exactly

\[
                                      2,\quad4,\quad4.
 \tag{R20.15}
\]

The highest-line scalar is4, while the full content is2. This is an explicit counterexample to identifying the whole integral branch with its primitive highest-line scalar, or with a homothetic copy of the target Weyl lattice. It is not a counterexample to R019, whose scope was expressly that highest line and the complex norm.

\subsection{6. Check scope and implications}\label{check-scope-and-implications}

The companion diagnostic constructs the exact divided-power Weyl lattices, applies the literal coefficient extraction and normalized Fischer projection, and computes their full small Smith matrices. It covers70 small branching pairs and finds nine nonhomothetic images. The three-variable proof in section5 is complete without relying on those computations. The actual-host certificate computes the hook and cBr valuation bounds in section4; it does not enumerate the enormous actual matrices.

For DP\_MULT Q3e this isolates another actual terminal factor: the whole branch has a fixed finite content interval in the growing family. It leaves the prior-image restriction, class contraction, terminal detector lattice and full g\_j,e\_j,H comparison. For THETA3 G1, the same B is the branch in DO038's exact mass and real period equation; its nonzero image on a full Weyl module does not imply a nonzero restricted theta period. The polynomial-lattice statement lives over Z and its localizations, and its rational Schur functors have the lambda/K-theory interpretation already recorded at SGA6 VI3.3. No new motive, weight, or RH comparison is asserted.

Author and independent whole-proof readings are logged separately. This file does not claim a Lean build or a new proof of the accepted Weyl duality and characteristic-zero branching theorems.

\clearpage

\section{AM-K038 --- the actual sharp coupled G/Q row and its integral class matrix}\label{am-k038-the-actual-sharp-coupled-gq-row-and-its-integral-class-matrix}

Author: GPT-6 Astra Pro, ALP\_MACHINE /root/kummer\_extension. 22 November 2026 UTC; inherited reasoning effort is not separately exposed.

This finding constructs the two-edge coupled arithmetic row from the actual mixed-CM geometric kernels of K035--037 and the unchanged raw finite functions. It proves their finite level covariance directly. The resulting row belongs to the prescribed evGSD module and is the actual argument of PtilFF. An ordinary ColV row is not substituted for this coupled row.

The statement concerns the sharp m=2 family, with the actual ranks, characters, G/T primes, maximal level at29, and selected partitions already specified in K031/K032, R017, and LIFT26. Every admitted fixed tuple is allowed; its full finite Q set and its entire prescribed Y quotient are retained. This is a constructed arithmetic replacement family on the mixed host. It does not assert an instance of the incompatible old Fields.Kq record, a uniform coefficient field for an unbounded collection of new tuples, or a numerical DP44 gain.

\subsection{1. Data and exact source interfaces}\label{data-and-exact-source-interfaces}

Use the polarized host K\_h and A\_h of K037(CR1), with h=1 or4. Put \[
 a=k,\quad d=t+1,\quad r=a+1,\quad
 b=a+d=r+t,\quad c=a+2d,\quad n_T=c-1=r+2t.
 \tag{CA1}
\] The two concrete rank quadruples (a,d,b,c) are (42,1760,1802,3562) and (168,1634,1802,3436). All three column ranks a,b,c are even. The terminal row has rank n\_T. Keep the actual sharp parameter rS=0. On the coupled chain m2=0 the source forces cOf(0)=1, cOf(1)=0 and h0=h1=1; there are no preceding uncoupled matrix edges.

At every G or Q place the local field is the rational q-adic field, with q a uniformizer. For G this is part of the chosen degree-one split prime data. For Q it follows from literal GoodPl: DegOne requires prime residue norm and q outside the square of the prime. The tuple is tame and unramified outside GoodPl. These are the definitions in \nolinkurl{Tuples.lean64–67,112–114,128,} not an inference from residue degree alone.

Write O for this local valuation ring and kappa for its residue field. Let psi have conductor O and put psi0(z)=psi(q\^{}\{-1\}z). Its finite argument ranges are those of the prescribed LocCst. At G let chi0,chi1 be the two actual residue characters in Ch.chi; the honest splitting characters eta\_j restrict to chi\_j\^{}\{-1\}. At Q let chi be the residue character of the fixed tuple; both eta\_j restrict to chi. These exact conventions are LIFT26(GS7a), LIFT23, and FS1004--1013. No character is conjugated implicitly.

Here and below FS denotes the recovered primary file \nolinkurl{Formalism/LocalToGlobal/Columns/FiniteSumNonzeroLargeDepth.lean.} The other paths in this table have the same recovered source root.

\begin{longtable}[]{@{}
  >{\raggedright\arraybackslash}p{(\columnwidth - 2\tabcolsep) * \real{0.5000}}
  >{\raggedright\arraybackslash}p{(\columnwidth - 2\tabcolsep) * \real{0.5000}}@{}}
\toprule\noalign{}
\begin{minipage}[b]{\linewidth}\raggedright
Source locator
\end{minipage} & \begin{minipage}[b]{\linewidth}\raggedright
Exact interface used
\end{minipage} \\
\midrule\noalign{}
\endhead
\bottomrule\noalign{}
\endlastfoot
FS194--245 & brC, RowPt, tilQ, PtilFF, evRho and evGSD \\
FS270--323 & Cpl, CplKer, its two covariance equations, class-weighted comp and colCpl \\
FS350--445 & DecoF full Q frames, their action, frmRho and frmGSD \\
FS610--639 & ThDist, determinant/rational covariance, and kerOf as an unaveraged finite latSum \\
FS640--665 & inFr, phiT, and liftFr \\
FS684--715 & entB, Q parabolic lattice, lamQB and PhiQ \\
FS738--810 & sharp G flags, lamGBS and PhiGS \\
FS909--941,982--986 & Jat, lamRow, muAt, omPhi and PsiAt \\
FS1019--1026,1062--1083 & zero/last h, initial levels, prescribed middle groups and characters \\
FS1103--1167 & actual frames, XiCpl1, PhiTop, ftY and XiCpl2 \\
FS1180--1210,1220--1275 & Y transport, midGSD, ybCpl, ChKer and its exact Shape \\
FS1408--1418,1451--1454 & initial startC, coupled colC with evGSD codomain, and PerInv \\
\nolinkurl{Columns/RowForms.lean221–231,253–269} & actual rowTyG, rowTyQ, gluing and inverse detector convention \\
\nolinkurl{Columns/RowFormsSym.lean16–20;} Columns/LevGSym.lean11--13 & rPsymF and symIdx, the latter the identity for rS=0 \\
\nolinkurl{Columns/GenericUnipotentCharacter.lean526–557,603–620} & HPsym, chiSym, HQ and the bottom determinant character \\
\nolinkurl{KernelLattices/ThDistC.lean349–455} & exact Cp1Equiv/Cp2Equiv, cp1K/cp2K, and the chKer constructor \\
Columns/RecLevel.lean50--52 & arithmetic finite character and inverse adelic right character \\
\end{longtable}

The companion certificate records SHA256 hashes of the recovered source bytes actually checked here; their mirror/author provenance is in the existing source-recovery records. This proof invokes the already accepted coefficient-module results at K037's stated interfaces. There is no Lean execution.

Fix one number field M containing the actual finite character values and the all-order geometric jet field of K036 for these two edges and this fixed tuple. All equations below are algebraic over M after the common period and source-frame scale of K035 are removed. No effective degree, discriminant or height bound for M is asserted.

\subsection{2. Coordinate and honest frame rules}\label{coordinate-and-honest-frame-rules}

At a local place a decoration frame F changes to F kappa\^{}\{-1\} under a pure finite action. For a simultaneous initial, middle and terminal change kappa0,kappa1,kappaT, the first coordinates and the terminal coordinates therefore change as follows: \[
\begin{split}
 (X,y)&\longmapsto
       (\kappa_1X\kappa_0^{-1},\kappa_1y),\\
 (Z,\eta,y)&\longmapsto
       (\kappa_TZ\kappa_1^{-1},\eta\kappa_1^{-1},\kappa_1y).
\end{split}
\tag{CA2}
\] Here X has size b by a, Z has size n\_T by b, and eta is the final row, including the literal piT\^{}\{-1\} in FS1166. The terminal finite action on the full rank-c target is diag(kappaT,1), so that row is fixed. In particular \[
                 (\eta\kappa_1^{-1})(\kappa_1y)=\eta y.
                 \tag{CA3}
\] This is the exact phase cancellation, with its original minus sign inside ftY. At G/Q piT is a unit; it has not been removed or replaced.

The honest split Weil action is the one checked in LIFT26(GS7a): \[
 [\omega_j(g,h)f](X)
 =\eta_j(\det g)^{-1}\eta_j(\det h)
       |\det g|^{N_{j+1}/2}|\det h|^{-N_j/2}
                    f(h^{-1}Xg).
 \tag{CA4}
\] For integral frame changes as in(CA2), the extra factor is \[
              \eta_j(\det\kappa_{\rm source})
                        \eta_j(\det\kappa_{\rm target})^{-1}.
                        \tag{CA5}
\] All absolute values in(CA5) are one. At Q this is chi(det source)/chi(det target). At G it is chi\_j(det source)\^{}\{-1\}chi\_j(det target). These are also exactly the inverse-frame determinant factors of ThDist.Ok.det with FS1246--1253's rhoS/rhoT. They are necessary: raw polynomial covariance alone leaves the wrong determinant character.

At Q the middle frames in DecoF are full local integral frames, so (CA2) is exact. At G a chosen liftFr is determined only modulo the principal congruence subgroup. Two choices differ by I+qA. The G raw matrix factors below depend only on the residue matrix. The G Y coordinate changes by an integral vector, and eta is integral on ftY's support; hence its additive phase is unchanged. At the terminal Q place the arbitrary liftFr affects only Z, and PhiQ depends only on its residue and integral support. Thus the computations do not depend on a preferred liftFr.

\subsection{3. Individual-Y covariance at Q}\label{individual-y-covariance-at-q}

Let P\_\{r,s\} be the upper block parabolic of sizes(r,s). The actual initial group is GL\_a(O), the actual middle group is the inverse image of P\_\{r,t\}(kappa), and the terminal group is the inverse image of P\_\{r,2t\}(kappa). Write \[
 \bar\kappa_1=
 \begin{pmatrix}A_1&B_1\\0&D_1\end{pmatrix},\qquad
 \bar\kappa_T=
 \begin{pmatrix}A_T&B_T\\0&D_T\end{pmatrix}.
 \tag{CA6}
\] The arithmetic characters are \[
 \lambda_{0,Q}=1,\quad
 \lambda_{1,Q}(\kappa_1)=\chi(\det D_1)^{-1},\quad
 \lambda_{T,Q}(\kappa_T)=\chi(\det D_T)^{-1}.
 \tag{CA7}
\] For the middle group, this follows directly from lamQB and cOf(1)=0. For the terminal it is the inverse lamQP in RowForms230--231.

Write X=(A\_X;C\_X) and y=(z;w), with top block size r. The exact first raw factor, LIFT23(Q5), is \[
 F_Q(X,y)=
 {\bf1}_{X\in M_{b,a}(O),\,C_X\in qM_{t,a}(O),\,
                  z\in O^r,\,w\in qO^t}
                         \chi(\det[\bar A_X,\bar z]).
 \tag{CA8}
\] The character is extended by zero at singular determinants. The r columns are precisely the a columns of X and the last y column. The support is preserved by(CA2) and the stated groups. Modulo q its top square matrix changes to \[
       A_1[\bar A_X,\bar z]\operatorname{diag}
                                      (\bar\kappa_0^{-1},1).
\] Consequently, including zero values, \[
 F_Q(\kappa_1X\kappa_0^{-1},\kappa_1y)
   =\chi(\det A_1)\chi(\det\bar\kappa_0)^{-1}F_Q(X,y).
 \tag{CA9}
\] Multiplication by(CA5) for the first edge gives exactly lambda\_\{1,Q\}(kappa1)lambda\_\{0,Q\}(kappa0)\^{}\{-1\}.

Let H\_Q(Z) be the literal terminal PhiQ. Its support consists of integral Z with lower-left 2t by r block divisible by q and invertible leading r by r block; its value is chi of that leading determinant. The terminal factor is \[
 H_Q(Z)\,q^{-\mathsf M b}
              {\bf1}_{\eta\in q^{-\mathsf M}O^b}
                       \psi(-\eta y),\qquad \mathsf M\ge1.
 \tag{CA10}
\] Its Y parameter lies in q\^{}\{-M\}O\^{}b/q\^{}MO\^{}b. The raw support is preserved under(CA2), and its leading block changes to A\_T Z\_\{11\} A\_1\^{}\{-1\}. The eta support and phase are invariant by(CA3). Thus the raw factor changes by chi(det A\_T)chi(det A\_1)\^{}\{-1\}. Combining with(CA5) for the second edge gives \[
       \chi(\det D_1)\chi(\det D_T)^{-1}
           =\lambda_{1,Q}(\kappa_1)^{-1}
                              \lambda_{T,Q}(\kappa_T).
           \tag{CA11}
\] No Fourier summation or division by the size of the Y set was used. In particular(CA9)--(CA11) hold for every individual Y term.

The statement is independent of the chosen Y representative. In(CA8), adding q\^{}M O\^{}b to y changes neither integral support nor residue determinant; in(CA10), its pairing with q\^{}\{-M\}O\^{}b is integral. This also proves that arbitrary representatives in FS.repY are allowed.

\subsection{4. Individual-Y covariance at G}\label{individual-y-covariance-at-g}

For an upper unipotent a by a matrix u set \[
             \psi_U(u)=\psi_0\left(\sum_{i=1}^{a-1}u_{i,i+1}\right).
             \tag{CA12}
\] It is a group character because the indicated superdiagonal entries add under multiplication. The exact middle residue group is \[
 H^\sharp_{r,t}=
 \left\{\begin{pmatrix}\operatorname{diag}(u,1)&A\\0&D\end{pmatrix}
       :u\in U_a(\kappa),\ D\in GL_t(\kappa)\right\}.
 \tag{CA13}
\] This is the direct flag calculation LIFT22(CG4), not the generic ordinary ColV1 level. The initial residue group is U\_a. The terminal group has upper block form of sizes(r,t,t), with first diagonal block diag(u\_T,1) and other diagonal blocks D\_\{T,0\},D\_\{T,1\}. Its exact characters, from lamGBS and rowTyG(rPsymF), are \[
\begin{split}
 \lambda_{0,G}(u_0)&=\psi_U(u_0),\\
 \lambda_{1,G}(\kappa_1)&=\psi_U(u_1)\chi_0(\det D_1),\\
 \lambda_{T,G}(\kappa_T)&=
          \psi_U(u_T)\chi_0(\det D_{T,0})\chi_1(\det D_{T,1}).
\end{split}
\tag{CA14}
\] For rS=0, symIdx is the identity, the distinguished r-th column in the first r block is fixed, and chiSym omits its joining root. This gives(CA12), rather than a generic character on U\_r.

The first raw factor is zero unless \[
                  y=q^{-1}e_a\pmod{O^b}.
                  \tag{CA15}
\] On that class it is exactly \[
 F_G(X,y)=
 \begin{cases}
  \psi_U(u_X),&X\in M_{b,a}(O),\
                    \bar X=\binom{I_a}{0}u_X,\quad u_X\in U_a(\kappa),\\
  0,&\text{otherwise}.
 \end{cases}
 \tag{CA16}
\] Here e\_a uses zero-based indexing, so is the last vector of the first r block. Both statements are the explicit raw PhiGS/PsiAt calculation LIFT22(CG7--9). They are retained with no projection.

The group(CA13) fixes e\_a, so(CA15) is stable. Under(CA2) the matrix u\_X changes to u\_1 u\_X u\_0\^{}\{-1\}. The raw factor therefore acquires psi\_U(u\_1)psi\_U(u\_0)\^{}\{-1\}. The honest first factor(CA5) is chi0(det kappa0)\^{}\{-1\}chi0(det kappa1)=chi0(det D1). The result is precisely lambda\_\{1,G\}lambda\_\{0,G\}\^{}\{-1\}. This proves the first covariance also off its support.

For the terminal raw factor, write the first b rows of Z as B and the last t rows as Z\_2. Its support is \[
         Z\in M_{n_T,b}(O),\quad \bar Z_2=0,\quad
         \bar B\in H^\sharp_{r,t}.
 \tag{CA17}
\] If the blocks of B are u\_B,D\_B as in(CA13), its value is \[
 H_G(Z)=\psi_U(u_B)\chi_0(\det D_B)\chi_1(\det D_B)^{-1}.
 \tag{CA18}
\] The full terminal local factor multiplies this by 1\_\{eta in O\^{}b\}psi(-eta y). These are the literal PhiGS and ftY: PhiGS has a=b at j=2, its last inverse unitV is chi1(det B)\^{}\{-1\}, and det B=det D\_B on the support.

Right multiplication by kappa1\^{}\{-1\} preserves(CA17) and gives \[
 H_G(Z\kappa_1^{-1})
  =\psi_U(u_1)^{-1}\chi_0(\det D_1)^{-1}
                                      \chi_1(\det D_1)H_G(Z).
 \tag{CA19}
\] For the terminal left action let C\_T be its leading b by b block. It belongs to Hsharp\_\{r,t\}. The zero lower block in(CA17) shows that the new leading block is C\_T B, while the last t rows remain zero. Thus \[
 H_G(\kappa_TZ)
 =\psi_U(u_T)\chi_0(\det D_{T,0})
                           \chi_1(\det D_{T,0})^{-1}H_G(Z).
 \tag{CA20}
\] These are identities of functions, including all zero cases, because the inverse group actions preserve the same supports. The honest second frame factor is \[
             \chi_1(\det D_1)^{-1}
                  \chi_1(\det D_{T,0}\det D_{T,1}).
                  \tag{CA21}
\] Equations(CA19)--(CA21), together with(CA3), leave exactly lambda\_\{1,G\}\^{}\{-1\}lambda\_\{T,G\}. This is the full compact-group covariance, not merely its restriction to a diagonal torus or the identity cell of a quotient model.

Adding an integral vector to y leaves(CA15)--(CA16) unchanged and pairs integrally with eta in the terminal factor. Thus it is an identity on the actual G quotient q\^{}\{-1\}O\^{}b/O\^{}b.

\subsection{5. Rational transport, T, and the actual polynomial kernels}\label{rational-transport-t-and-the-actual-polynomial-kernels}

At T the raw factor is exactly phiT: it says that the linear map sends the first local lattice into the second and pulls its designated covector back to the first covector. Under rational isometries gamma0,gamma1 and x -\textgreater{} gamma1 x gamma0\^{}\{-1\}, this incidence condition is unchanged. The same statement holds on the second edge with gammaT extended to diag(gammaT,1). The literal tilQ/tilH extension and its piT coordinate give exactly this covector transport. The finite identity does not require erasing a character at T. Its prescribed initial type is the actual one of R017/DO035/LIFT24. Away from G,Q,T retain the fixed maximal vacuum and the actual unramified splitting; its compact covariance is already that of K035. In particular no principal29 refinement is introduced.

The Y set is the quotient of the two localized middle lattices specified in FS1125--1129. The CRT description is \[
 {\cal Y}(p)\simeq
 \prod_{G}q^{-1}O^b/O^b\ \times\
 \prod_Qq^{-\mathsf M}O^b/q^{\mathsf M}O^b .
 \tag{CA22}
\] It is finite. A rational gamma1 sends it bijectively to the quotient for gamma1 N; a pure finite decoration action leaves the global coset Y unchanged and changes only its coordinates as in(CA2). This is exactly the action in ybCpl. The source proof FS1180--1210 is the elementary identity gamma1(z-y)=gamma1z-gamma1y.

Define the geometric functionals J\_\{1;p0,p,Y\}, J\_\{2;p,Y,d\} by the finite Mumford translation and induction sums of K035--037, using the literal XiCpl1 and XiCpl2 coefficients just displayed. For an inadmissible Y set the two functionals equal zero. They are defined from algebraic sections, character phases, and finite sums before their analytic theta comparison. No new Schwartz vector is selected.

Let Phi\_J be K037's universal subD operator on the selected coefficient spaces. Define \[
 k_1(p_0,p,Y)=\Phi_{J_{1;p_0,p,Y}},\qquad
 k_2(p,Y,d)=\Phi_{J_{2;p,Y,d}}.
 \tag{CA23}
\] The first map goes from V\_0 to V\_1; the second goes from V\_1 to V\_2. Every V\_j is the source rhoVr polynomial model with its actual selected weight. Its integral lattice is the literal delAtr Weyl lattice, not a replacement by the dual Weyl lattice.

K036 gives exact rational transport of the geometric kernel, including its determinant half-rank scalar. K037(CR18) gives the matching selected polynomial action, with source rhoVr and target rhoVr after the subD contraction. Equivalently, substitute X -\textgreater{} gamma1 X gamma0\^{}\{-1\} in the geometric germ, use K036's actual determinant scalar, and apply the ordinary-derivative chain rule inside K037(CR12). The character exponents in(CR2)--(CR3) cancel that scalar and yield the two displayed polynomial actions. This is the exact rational covariance proved in K036--037; no quadratic elliptic period is inserted.

Combining that rational identity with the finite identities (CA9)--(CA11),(CA14)--(CA21) proves \[
 k_1(g_0p_0,g_1p,g_1Y)
   =\rho_1(g_1)\,k_1(p_0,p,Y)\,\rho_0(g_0^{-1}),
\tag{CA24}
\] and \[
 \boxed{\ k_2(g_1p,g_1Y,g_Td)
       =\rho_{\rm ev}(g_T)k_2(p,Y,d)\rho_1(g_1^{-1}).\ }
 \tag{CA25}
\] Equations(CA24) and(CA25) are exactly CplKer.equiv1/equiv2. Here rho1 is frmRho, including its arithmetic middle character, and \[
 \rho_{\rm ev}(\gamma_T,\kappa_T)
       =\lambda_T(\kappa_T)\,
                      \rho_2^\vee(\operatorname{diag}(\gamma_T,1)).
 \tag{CA26}
\] Indeed evRho takes the inverse of rowTy.inv's character, which is rowTy's original character. Equations(CA7),(CA14) are exactly that character. The zero extension outside the allowed Y set is equivariant because(CA22) is equivariant.

Only the middle group acts on Y in these two equations.

\subsection{6. The actual coupled class map and its terminal domain}\label{the-actual-coupled-class-map-and-its-terminal-domain}

Take the finite decorated class sets C0,C1 for the initial and sharp middle groups, and their fixed representatives p\_\{0,e\},p\_\{1,f\}. Let Gamma0e,Gamma1f be their actual finite stabilizers. Finiteness comes from the definite selfdual lattice-class argument of K036 and finite residue decorations. The odd terminal rank uses the same finite-class argument, whose proof does not require even rank. Even rank was needed only for the subsequent all-order period field transport on the two column edges.

The source module consists of its prescribed equivariant rows P\_e, with its integral lattice defined by intersection with delAtr as in K037(CR21)--(CR22). Define, for every terminal decoration d, \[
 \boxed{\ 
 P_{\rm cpl}(d)=
  \sum_{e\in{\cal C}_0}\frac1{|\Gamma_{0,e}|}
  \sum_{f\in{\cal C}_1}\frac1{|\Gamma_{1,f}|}
  \sum_{Y\in{\cal Y}(p_{1,f})}
       k_2(p_{1,f},Y,d)\,
       k_1(p_{0,e},p_{1,f},Y)P_e .
 \ }
 \tag{CA27}
\] There is no factor 1/\textbar Y\textbar. The scalar q\^{}\{-Mb\} of ftY is already inside k2. This is the actual CplKer.comp followed by the initial class operator in FS298--323. Since m2=0, no preceding kk edge occurs. The singleton initial chart is the source's startC inclusion.

This initial module is the actual MfV=ColV0, not only its finite Steinberg subspace. ColForms.lean93--115 gives its initial G residue group U\_a and arithmetic character psi\_U; RecLevel50--52 gives the inverse adelic right character. At T its finite level fixes the designated nonzero covector, with trivial character on that strict stabilizer. It is maximal at the other places, including 29 and the initial Q places. The ramified source subspace is \[
 {\cal M}_{\rm ram}
       =\bigcap_{P\mid T}\ker\left(\sum_{z\in\kappa_P^\times}
                                      \operatorname{diamond}_P(z)\right)
                                      \subseteq{\cal M}_{f}.
 \tag{CA27a}
\] \nolinkurl{ScaledRamifiedLattice.lean84–91,155–175} is exactly this definition: the norm is an unweighted sum, without an auxiliary character. Record.lean1811--1848 proves its literal transport to the adelic scalar-dilation action. Over C the finite cyclic quotient acts semisimply, so(CA27a) retains precisely the nontrivial dilation characters at each T place. Formula(CA27) is defined on all MfV and restricts to this actual Mram and its inherited integral lattice. There is no extra initial projector in this restriction.

\textbf{Theorem CA28.} Formula(CA27) is an M-linear map from the actual initial row module to the literal evGSD.EqFun module with target polynomial space V\_2 and its prescribed terminal finite character. It is independent of the chosen class representatives, Y representatives, and liftFr choices, with the natural row transport.

\textbf{Proof.} Equations(CA24) and(CA25) cancel the two rho1 factors in each composition. Changing a middle representative therefore only bijects its Y set, and conjugates its stabilizer, whose cardinal remains unchanged. Changing an initial representative has the same effect after applying the equivariance equation for P\_e. This proves representative independence without choosing trivial stabilizers. The earlier periodicity proofs and liftFr argument prove the other two independence assertions.

For a terminal action g\_T, (CA25) left-multiplies every summand of (CA27) by rho\_ev(g\_T). Hence \[
                   P_{\rm cpl}(g_Td)=\rho_{\rm ev}(g_T)P_{\rm cpl}(d).
                   \tag{CA29}
\] This is precisely the defining evGSD equation. The target polynomial space is already correct by K037's integral subD theorem, so no post hoc isotypic projection is used. All sums are finite and all entries lie in M, proving linearity and algebraicity. QED.

This theorem supplies the missing raw covariance input in K037§7 and the admitted coupled-row input of DO038. In particular the exact arithmetic terminal functional is \[
 \operatorname{PtilFF}(P_{\rm cpl},G)
   =\sum_{z\in{\cal C}_T}\frac1{|\Gamma_{T,z}|}
       \operatorname{pairBK}
                 \bigl(\operatorname{brC}P_{\rm cpl}(d_z),G(d_z)\bigr).
 \tag{CA30}
\] Its brC is unchanged, including brConst/kapBr and the specified conjugate piT power. This is FS215--218, not a pairing defined by division of an analytic period. The final column uses rhoVr, the detector rhoV; no transpose or conjugation scalar has been inserted. DO038's actual terminal class transport and top realization therefore apply to this row after their stated hypotheses are retained.

Throughout(CA27)--(CA30), retain R017's maximal29 level. Distinguish the literal MfV level from R017's coarser support realization. At T the former uses the strict covector stabilizer H\_T; R017 uses the larger projective-line stabilizer P\_T with its chosen nontrivial character. At G, R017 keeps the full maximal compact and realizes the finite Steinberg type by its matrix coefficient; literal MfV uses the inverse image of U\_a with right character psi\_U\^{}\{-1\}. R017(R17.9) gives stabilizer1 after the strict generic T decoration even with G still maximal, and stabilizer mu28 for the coarser projective-T/full-G support. No principal29 refinement is made in any of these realizations.

Here is the exact refinement rule. For compact levels H contained in P of index I, the strict classes above a projective class with stabilizer Gamma\_P are the orbits of Gamma\_P on P/H. If their stabilizers are Gamma\_s, orbit-stabilizer gives \[
 \sum_s\frac1{|\Gamma_s|}=\frac I{|\Gamma_P|},\qquad
 v_H=\frac{v_P}{I},\qquad
 \sum_s\frac{v_H}{|\Gamma_s|}
                      =\frac{v_P}{|\Gamma_P|}.
 \tag{CA30a}
\] For the T-only refinement of R017, with G still maximal, the map from mu28 into P\_T/H\_T is injective. Its action is free. Put I\_T=product\_\{P\textbar T\}(q\_P-1). There are I\_T/28 intermediate classes, each with trivial stabilizer. The additional literal G refinement has index \[
 I_G=\prod_{P\in G}
       [GL_a(\kappa_P):U_a(\kappa_P)].
 \tag{CA30c}
\] Because each intermediate stabilizer is trivial, every such class refines into I\_G classes. Thus the full initial strict MfV support has I\_T I\_G/28 classes, each with stabilizer1, and volume v\_P/(I\_T I\_G) per class, where P here is R017's projective-T/full-G support level. Their total mass is v\_P/28. For norms the T character has modulus one, while the actual finite G matrix coefficients have average square1/d\_G by finite Schur orthogonality. Together with the real factor 1/d\_rho this recovers exactly R017(R17.14); it does not assign constant magnitude to the G coefficients.

It does not make the raw theta summands identical. The raw incidence kernel is H\_T-invariant and need not be nu\^{}\{-1\}-eigen under P\_T. Denote its actual right input character by nu (the inverse of R017's nu\_P if that notation is used). On this component coarsening the T integration is the exact identity below, with P=P\_T, H=H\_T and I={[}P\_T:H\_T{]}: \[
 \int_P F(gp)\theta_\phi(gp,h)\,dp
   =v_P F(g)\theta_{\Pi_{\nu^{-1}}\phi}(g,h),\qquad
 \Pi_{\nu^{-1}}\phi=
      \frac1I\sum_{z\in P/H}\nu(z)\omega(z)\phi .
 \tag{CA30b}
\] The character is trivial on H, and the chosen H-fixed finite vector makes this formula well-defined. Change variables in the sum to check that the average is nu\^{}\{-1\}-eigen; its factor 1/I is retained. Formula(CA30b) is an identity obtained from the original integration, not an extra Schwartz replacement. Thus(CA27) always uses the actual strict stabilizers and kernels. Using R017's projective scalar requires its actual transported nu-weighted kernel sum as well. Scalar masses alone do not justify replacing that sum by one representative kernel. At G the analogous coset decomposition retains the finite Steinberg matrix-coefficient weights: the input transforms by psi\_U\^{}\{-1\} and the raw source kernel by psi\_U, so their product descends to K\_G/J\_\{0,G\}. Summing that product over the cosets, with measure vol(J\_\{0,G\}), is exactly its integration over K\_G. It does not replace the raw kernel by an unweighted average. This is the finite G normalization already retained in R017/DO036/LIFT20.

\subsection{7. Exact integral image containment and the finite matrix gate}\label{exact-integral-image-containment-and-the-finite-matrix-gate}

Fix a finite place w of M and write R=O\_\{M,w\}. Use the actual transferred delAtr lattices L0e,L1f,L2d. All are free over R after localization. For each matched triple let j1(e,f,Y;w) and j2(f,Y,d;w) be the fractional ideals generated by the finitely many ordinary derivatives in the exact local frames of K037(CR39) for the selected multidegree on that edge. Include their geometric phase/average coefficients, their cD factors, and the determinant factors of those frames. These are fixed ideals, not the content of an ambient arbitrary integral function.

The integral Weyl-to-Weyl substitution theorem K037(CR8) gives \[
 \begin{split}
 k_1(L_{0,e})&\subseteq j_1(e,f,Y;w)L_{1,f},\\
 k_2(L_{1,f})&\subseteq j_2(f,Y,d;w)L_{2,d}.
 \end{split}
 \tag{CA31}
\] Consequently \[
 \boxed{\ 
 P_{\rm cpl}({\cal L}_0)_d\subseteq
   \left(\sum_{e,f,Y}
       \frac{j_2(f,Y,d;w)\,j_1(e,f,Y;w)}
            {|\Gamma_{0,e}|\,|\Gamma_{1,f}|}\right)L_{2,d}.
 \ }
 \tag{CA32}
\] This statement follows by composition and the triangle inequality for fractional ideals. It does not require an inverse Gram matrix. In valuation form, its lower bound is the minimum of v\_w(j1)+v\_w(j2)-v\_w(\textbar Gamma0\textbar)-v\_w(\textbar Gamma1\textbar) over the indicated terms. Since(CA29) already holds, intersecting the right side with the equivariant target space supplies the actual target row lattice.

When all G and Q residue characteristics differ from11, their raw coefficients are roots of unity, zero, or q\^{}\{-Mb\} times roots of unity. They are integral, and their nonzero scalar factors are 11-units. Thus this raw finite assembly introduces no new 11-denominator beyond the displayed geometric jet/frame ideals and actual stabilizer factors. This conditional statement excludes any tuple having a Q place over11; the proof of covariance itself allows such a tuple. In R017's generic initial strict classes the stabilizer factor is1; in its projective realization28 is an11-unit, but the weighted average(CA30b) has the displayed factor1/I. No claim is made that every intermediate or terminal stabilizer has order prime to11.

For complete precision, choose R-bases of the initial invariant row lattice and target invariant row lattice. Their existence follows from the saturated-intersection description and the DVR property. Formula(CA27), with its literal subD matrices, gives a finite matrix A\_cpl over Frac(R). Formula(CA30) adds the fixed branch and terminal pairing matrix. The least entry valuation of A\_cpl is the exact content of its image in the chosen target lattice; its minors give its other determinantal ideals. On an already produced input-image lattice, use that lattice's basis instead. The preliminary containment(CA32) need not be sharp: Y summation, class summation, and terminal branching can increase valuations or annihilate an image. No saturation assertion follows from nonzero analytic theta lift alone.

The gate is therefore a specific finite algebraic image matrix on fixed arithmetic lattices, with all its entries defined by the actual geometric kernels. It is no longer a missing unspecified coupled covariance or terminal-domain identification. Enumeration and a useful uniform bound for this matrix are not supplied here.

\subsection{8. Analytic comparison, complete type, and fixed highest vectors}\label{analytic-comparison-complete-type-and-fixed-highest-vectors}

Let v0,v1 be the actual strict-covector finite level volumes used for the classes in(CA27), with compact archimedean Haar of mass one. Projective nu-isotypic realizations instead use the precise coarsening(CA30a)--(CA30b). Let D\_\{j,tau\} be the two prescribed selected degrees, and let s\_tau=lambda\_tau/cD\_tau be K035's literal source frame scale. Set \[
 {\cal P}_{01}
   =v_0v_1 C_\theta^{14h(ab+bc)}
            \prod_\tau s_\tau^{D_{0,\tau}+D_{1,\tau}}.
 \tag{CA33}
\] Keep the same actual finite tensor as LIFT26: sum\_Y phi1Y tensor phi0Y. Its pointwise identity(GS7a--c) proves that the moving-lattice Y sum is the action of this one fixed tensor; no h-dependent projection is inserted.

Apply K037's complete coevaluation construction at both real types to this tensor. The coevaluation tensors are the canonical sum v\_i tensor p\_i with pairBK-dual bases. Class integration and pointwise coevaluation give, on the actual terminal restriction, \[
          {\cal T}_{\rm cpl}(P;v)
                      ={\cal P}_{01}\,f_{P_{\rm cpl},v}.
                      \tag{CA34}
\] To check the coefficient directly, first decompose both compact adelic quotients into their actual classes. Their masses are v0/\textbar Gamma0\textbar{} and v1/\textbar Gamma1\textbar. The complete source coevaluation contracts to P\_e; the complete middle coevaluation contracts the resulting V1 vector with the second subD kernel. The remaining sum is precisely(CA27). Formula K035(PF28) supplies the two common period/degree factors in(CA33), independent of Y and of the class. This proves(CA34) without rescaling the arithmetic row.

Complete coevaluation does not create a factor 1/dim V\_j. For a single source coefficient line, compact Schur integration does give that dimension factor, as in K037. In particular the fixed highest Chern creation vectors of LIFT25/LIFT26 are not silently identified with the source's unit-normalized pureTensorO: \[
 \operatorname{pureTensorO}(P_j)
          =\Lambda_j\,r_{P_j},\qquad
 \Lambda_j=\prod_\tau\lambda_\tau^{D_{j,\tau}}\ne0.
 \tag{CA35}
\] Replacing the two Chern vectors by the literal source vectors multiplies the two-edge operator by Lambda0 Lambda1 and its squared norm by \textbar Lambda0 Lambda1\textbar\^{}2. The further single-line Schur and R017 input normalizations remain the explicit factors of LIFT25/LIFT26; they are not absorbed into(CA27) or(CA33). These complex period factors are not assigned finite valuations.

For the actual normalized highest vectors, one can make the relation to the fixed detector exact. Let h\_j be the norm-one highest vector in the rhoV coefficient model and let p\_j be its corresponding dual functional with p\_j(h\_j)=1. The latter is a lowest vector for the dual upper Borel, as in LIFT25(AR6), not the unrelated highest vector of that dual representation. The real polynomial of LIFT25 is iota\_j(p\_j,h\_\{j+1\}), with its explicit Chern/unit-creation conversion(CA35). Compact Schur averaging of h\_j tensor p\_j equals the canonical coevaluation divided by d\_j=dim V\_j: its average is an invariant endomorphism, hence scalar, and its trace is1. Thus, for the same fixed coupled finite tensor and every initial arithmetic row P, \[
 {\cal T}_{\rm cpl}(P;h_2)
    =d_0d_1\,{\cal B}_{\rm source,highest}(f_{P,h_0})
    =d_0d_1\Lambda_0\Lambda_1\,
                   {\cal B}_{\rm Chern,highest}(f_{P,h_0}).
 \tag{CA36}
\] This argument is independent of the finite datum on each summand and so commutes with its prescribed full Y sum. It retains all factors in(CA33) independently. It does not assert that a fixed highest tensor itself is the full coevaluation tensor.

Formula(CA34) identifies the canonical complete-type analytic map with the newly constructed arithmetic row. It does not prove that the nonzero unrestricted full G2 output in LIFT26 has nonzero restriction to this terminal decorated subgroup, or survives brC. That is the actual image/terminal detection question in(CA30).

\subsection{9. The empty-Q reference member and family scope}\label{the-empty-q-reference-member-and-family-scope}

The reference tuple is the identity in tupGrp: Tuples.lean118--120. Its Qprime is empty by definition, and the actual source realization proves this at \nolinkurl{TrueColumnChoicesHeckeCharReferenceEqOne.lean273–279.} Thus Q=empty is an admitted reference branch; it retains the stage G and T characters. Equations(CA22)--(CA34) then simply omit the Q factors. The sharp middle group(CA13) is still the actual coupled middle group and need not equal ordinary ColV1.

This reference branch does not remove the nonreference tuples from the source comparison. DataParams₂.lean108's untwisted condition sets Aj and F to zero at every place whose muAt=mu0. LinksChar.lean386--402 requires every nonzero-index place to lie among the twist places and its tuple there to differ from mu0. Therefore a reference-only example supplies one actual operator and its costs; it does not prove the DP44 family inequalities or the gains carried by the nonreference rows.

The proof here works for every fixed admitted tame tuple with its full Q set. It also works for every fixed word for which the stated two-edge ranks, characters, selected weights and finite T data hold. It does not assert that the explicitly checked finite word/size ranges of R018 are unbounded.

\subsection{10. Checks and residual scope}\label{checks-and-residual-scope}

The companion finite certificate tests the two block covariance identities, including nonzero and zero-support cases, over small prime fields with exact formal character phases. It separately tests the complete ftY pairing modulo q\^{}\{2M\}, the sharp finite group character law, the double class/stabilizer contraction, and the source anchors. These are finite exact form checks; the proofs above establish the arbitrary-rank identities.

New result: the unchanged sharp raw G/Q family supplies actual Cp1Equiv/Cp2Equiv, hence a geometric arithmetic coupled map on the actual Weyl lattices with exact evGSD and PtilFF target. The coefficient matrix, stabilizer weights, common periods and creation scales are specified.

Still required for a numerical improvement: the restricted successive-image content of that matrix, the terminal branch image of the actual theta row, effective field/global height costs, and the complete source family comparison. None is deduced by replacing the image with an ambient integral module or by normalizing away a period, class mass or hook factor.

\clearpage

\section{ALP-LIFT27: injectivity of the unchanged sharp coupled operator on the entire initial ramified type}\label{alp-lift27-injectivity-of-the-unchanged-sharp-coupled-operator-on-the-entire-initial-ramified-type}

Author: GPT-6 Astra Pro, /root/lift\_accounting, ALP\_MACHINE. Date: 22 November 2026 UTC. Inherited reasoning effort is not exposed.

This strengthens LIFT26 from its explicitly chosen R017 input to every vector in the actual initial ramified automorphic row space. The new local point is that finite Steinberg type was unnecessary at G. A direct unitary Godement--Jacquet bound makes the whole-Schwartz quotient map holomorphic on every permitted initial constituent and its actual first theta image. At T, the literal unweighted norm operator is the trivial character projector, so all remaining dilation characters are handled by LIFT24. The source's complete coevaluation operator is also compared exactly with the previously checked highest-kernel operator.

The prescribed finite G/T/Q functions and full coupled Y sum are unchanged. No local Whittaker-genericity, local cuspidality, temperedness, initial finite Steinberg projection, extra finite translation, or new real rotation is assumed. The conclusion is before the literal terminal brC, tilQ, and PtilFF restriction. Injectivity of the unrestricted terminal row does not imply injectivity of that restriction.

\subsection{1. Fixed spaces, initial conditions, and theorem}\label{fixed-spaces-initial-conditions-and-theorem}

Use either actual mixed CM host K031 or K032, the common edge characters and selected real representations of LIFT26, and ranks \[
 (a,d,b,c)=(42,1760,1802,3562)
       \quad\hbox{or}\quad(168,1634,1802,3436),
 \qquad b=a+d,\quad c=b+d.
 \tag{MR1}
\] Write \(G_0=U(I_a),G_1=U(I_b),G_2=U(I_c)\). These groups are positive at every real place. Fix exactly the raw sharp coupled finite tensor \(\Psi_f\) of LIFT26(GS3),(GS6)--(GS7c), the source-style bilinear theta convention, and the unit-normalized source creation vectors of LIFT25(AR27)--(AR30). Their Chern-coordinate versions differ by the fixed nonzero scalar \(\Lambda_0\Lambda_1\).

At a selected split degree-one G place put \[
 K_a=\mathrm{GL}_a(O),\quad K_a^1=\ker(K_a\to\mathrm{GL}_a(\kappa)),
 \quad J_G=\operatorname{red}^{-1}U_a(\kappa).
 \tag{MR2}
\] The initial right condition is \(\psi_U^{-1}\) on \(J_G\), for the literal source generic upper-unipotent character. There is no restriction to the finite Steinberg representation.

At a selected T place let \(u\ne0\) be the actual residue covector, and \[
 H_a=\{g\in K_a:u\bar g=u\},\qquad
 J_a=\{g\in K_a:u\bar g=s(g)u,\ s(g)\in\kappa^\times\}.
 \tag{MR3}
\] The input is \(H_a\)-fixed and killed by the unweighted sum of the \(\kappa^\times\) dilation operators. At Q and all other finite places, including29, it has the actual maximal compact condition of LIFT26. At infinity its right vector is the fixed unit highest vector \(h_0\) of the selected compact representation \(W_0\). Let \(\mathcal H_{\mathrm{ram},h_0}\) be this finite-dimensional automorphic subspace. All rational stabilizer relations, including central roots of unity, are retained. An incompatible tuple of local characters gives a zero summand; it is not artificially admitted.

Let \(\mathcal B\) be the same full finite Y-sum of the two actual highest-kernel theta lifts as in LIFT26. Then \[
 \boxed{\ker(\mathcal B|_{\mathcal H_{\mathrm{ram},h_0}})=0.}
 \tag{MR4}
\] The canonical complete-coevaluation operator from actual coefficient rows to full G2 automorphic functions is injective as well, with the exact factor in section7. Section2 identifies these initial conditions with the source's actual \(M_{\mathrm{ram}}\). The final arithmetic comparison still retains its independently determined periods and class volumes.

\subsection{2. Identification with the literal initial ramified space}\label{identification-with-the-literal-initial-ramified-space}

All Lean locators are relative to work/research/alpoge\_source, commit 7cbf7a0c32fca1badf05931879a37c3f8a2b03b4.

\nolinkurl{Formalism/LocalToGlobal/Columns/ColForms.lean93–115} defines colLev0, MfV, and equivariant forms. At the exact first index \(n=k\), \nolinkurl{GenericUnipotentCharacter.lean437–451} gives \(HX(k,k,t,0)=U_k\) and \(\operatorname{lamX}=\psi_U\). RecLevel.lean50--52 takes the inverse level character for the adelic reading. Thus(MR2) is the entire prescribed G condition.

\nolinkurl{Columns/ScaledRamifiedLattice.lean84–91} defines NmG as the unweighted sum of diamondOperatorOnForms′ over residue units. Lines122--127 define diaOp and NmOp; lines155--169 give \[
 M_{\mathrm{ram}}=\bigcap_{P\in v_T}
       \ker\!\left(\sum_{s\in\kappa_P^\times}\operatorname{dia}_P(s)\right).
 \tag{MR5}
\] The decoration fixes the nonzero covector, not merely its line, giving \(H_a\). Scaling it gives \(J_a/H_a=\kappa^\times\). \nolinkurl{Columns/Record.lean1811–1848} proves transport of precisely that sum under the faithful adelic reading and the exact Mram equivalence. There is no determinant character in that sum. The stage splitting characters are unramified at T and hence trivial on compact determinants. Changing the integral frame conjugates these conditions.

For a unitary representation \(\pi\), the quotient acts unitarily on \(\pi^{H_a}\), so \[
 \pi^{H_a}=\bigoplus_{\nu\in\widehat{\kappa^\times}}\pi^{J_a,\nu},
 \qquad
 \ker\sum_s\pi(\delta(s))
             =\bigoplus_{\nu\ne1}\pi^{J_a,\nu}.
 \tag{MR6}
\] Here \(\delta(s)\) is any lift of s; its action on \(H_a\)-fixed vectors is independent of the lift. Indeed the character projectors are \((q-1)^{-1}\sum_s\nu(s)^{-1}\pi(\delta(s))\); the unweighted sum is \((q-1)\) times the trivial projector. Conditions at different T places commute. Thus the simultaneous kernel is the direct sum over tuples of nontrivial characters. This decomposition is used in the proof; it is not an extra projector in the theta construction.

K037(CR26),(CR32)--(CR33) gives the actual coefficient realization \(P\mapsto f_{P,h_0}\), including arithmetic stabilizers. It is injective: at each finite class a zero scalar function means its coefficient vanishes on the orbit of \(h_0\); that orbit spans the irreducible \(W_0\). With the actual class-row metric, \[
 \|f_{P,h_0}\|^2
   =\frac{v_0}{d_0}\sum_x|\Gamma_x|^{-1}\|P_x\|^2,
 \quad d_0=\dim W_0,\quad v_0=\operatorname{vol}(U_0)>0.
 \tag{MR7}
\] The level volume, stabilizers, source dimension, and arithmetic coefficient lattice are not replaced by unitary-coordinate data. Every actual \(M_{\mathrm{ram}}\) row therefore gives an element of \(\mathcal H_{\mathrm{ram},h_0}\).

Here \(U_0\) is the literal strict covector/unipotent level, not R017's coarser projective-T/full-G support. The exact refinement is important even though it does not affect injectivity. If \(H\subset P\) has finite index I and a P-class has rational stabilizer \(\Gamma_P\), then its H-classes are the \(\Gamma_P\)-orbits on \(P/H\). For their stabilizers \(\Gamma_s\), \[
 \sum_s|\Gamma_s|^{-1}=I/|\Gamma_P|,\qquad
 v_H=v_P/I,\qquad
 \sum_s v_H/|\Gamma_s|=v_P/|\Gamma_P|.
 \tag{MR7a}
\] This is orbit-stabilizer and Haar refinement, with no freeness assumption. In R017's particular generic support, the T refinement has \(I_T/28\) free classes and the further G refinement gives \(I_TI_G/28\) classes of stabilizer1, exactly K038(CA30a),(CA30c). The finite G coefficient squares retain their Schur average \(1/d_G\); they are not all1.

Coarsening raw theta integration also retains a vector average. If the input has actual right character \(\nu\) on P, the H-fixed kernel satisfies \[
 \int_P F(gp)\theta_\phi(gp,h)\,dp
   =v_P F(g)\theta_{\Pi_{\nu^{-1}}\phi}(g,h),\qquad
 \Pi_{\nu^{-1}}\phi=I^{-1}\sum_{z\in P/H}\nu(z)\omega(z)\phi .
 \tag{MR7b}
\] This follows directly by decomposing P into H-cosets; changing variables proves the displayed average is \(\nu^{-1}\)-eigen. It is the average supplied by the original integral, not a replacement of the prescribed Schwartz vector. Thus equal total class masses do not make raw kernel summands identical. All arguments below use the genuine integrals at their actual strict levels. If expressed in R017's projective realization, they use (MR7b) with its factor \(1/I\), as in corrected K038(CA30b). The analogous G refinement retains its finite matrix-coefficient weights. No principal29 level is inserted.

\subsection{3. Uniform local holomorphy for an arbitrary unitary input}\label{uniform-local-holomorphy-for-an-arbitrary-unitary-input}

Let F be nonarchimedean of characteristic zero and residue cardinality q. Normalize additive Haar by \(\operatorname{vol}M_n(O)=1\) and multiplicative Haar by \(\operatorname{vol}K_n=1\). Put \(\kappa_n=\prod_{i=1}^n(1-q^{-i})\), so \(dg=\kappa_n^{-1}|\det g|^{-n}dX\). For an irreducible smooth unitary \(\pi\) of \(\mathrm{GL}_n(F)\), a unitary character \(\eta\), a Schwartz function f on \(M_n(F)\), and a smooth matrix coefficient m, use \[
 Z(z,f,m;\eta)=\int_{\mathrm{GL}_n(F)}
       f(g)m(g)\eta(\det g)^{-1}|\det g|^{z+(n-1)/2}\,dg.
 \tag{MR8}
\] Every such integral is holomorphic in \[
                          \Re z>(n-1)/2.
 \tag{MR9}
\]

\textbf{Proof.} Smooth coefficients of a unitary representation are bounded. A smooth dual functional is represented by a Hilbert vector: take a compact subgroup fixing it, use Riesz on the finite-dimensional invariant space, and compose with the orthogonal compact average. Cauchy--Schwarz then bounds the full coefficient.

The checked determinant integral LIFT16(S8), with these measures, is \[
 \int_{M_n(O)}|\det X|^u\,dX
       =\prod_{i=1}^n\frac{1-q^{-i}}{1-q^{-u-i}},
                       \qquad\Re u>-1.
 \tag{MR10}
\] This accepted matrix-Iwasawa identity and its convergence range are used, not claimed as a new general integration theorem. Substituting \(u=\Re z-(n+1)/2\) gives precisely(MR9). Any Schwartz support is contained in \(\varpi^{-N}M_n(O)\); its determinant integral differs by a finite scaling factor. On every compact subset of(MR9), the same expression gives a locally uniform integrable majorant. Thus the zeta integral converges absolutely and holomorphically. ∎

The published all-irreducibles Godement--Jacquet ideal theorem, as restated by Jiang and Luo(2023), Theorem2.4(1)--(2), pp.313--314, says these integrals generate \(L(z,\pi\eta^{-1})\mathbf C[q^z,q^{-z}]\). The generator is therefore a finite Laurent-polynomial linear combination of such integrals. Consequently \[
 \boxed{L(z,\pi\eta^{-1})\text{ is holomorphic and nonzero for }
                                      \Re z>(n-1)/2.}
 \tag{MR11}
\] Holomorphy follows from the finite combination, and nonvanishing from the standard local expression \(1/P(q^{-z})\), \(P(0)=1\). No temperedness, local cuspidality, or local genericity is assumed.

Now take any unitary irreducible \(\pi_0\) of \(\mathrm{GL}_a(F)\) and unitary characters \(\eta_0,\eta_1\). Set \[
 z_*=(d+1)/2,\quad
 \pi_1=\pi_0\times(\eta_0\circ\det_d),\quad
 \pi_2=\pi_1\times(\eta_1\circ\det_d),\qquad d>a-2.
 \tag{MR12}
\] Both normalized inductions are irreducible and unitary by Tadic's TheoremD: combine the lists of unitary blocks of its factors and use TheoremD(i). The characteristic-zero appendix removes the historical unpublished dependency mentioned on p.339. No disjoint cuspidal support assumption is needed.

The standard parameter of irreducible normalized induction is the direct sum of its factors. The same multiplicativity used in LIFT24(MT27)--(MT30), now without a spherical assumption on \(\pi_0\), therefore gives \[
 L(z_*,\pi_1\eta_1^{-1})
   =L(z_*,\pi_0\eta_1^{-1})
                     \prod_{r=1}^{d}L(r,\eta_0\eta_1^{-1}).
 \tag{MR13}
\] The first factor is holomorphic and nonzero by(MR11). Each rank-one unitary factor at \(r\ge1\) is1 if ramified and otherwise \((1-\eta_0\eta_1^{-1}(\varpi)q^{-r})^{-1}\), so is finite and nonzero. In the actual G application the ratio is ramified, making that product equal1. The first-edge factor \(L(z_*,\pi_0\eta_0^{-1})\) is likewise holomorphic and nonzero. The actual rank margins are \[
                  z_*-(a-1)/2=(d-a+2)/2=860\ \hbox{or}\ 734.
 \tag{MR14}
\] These statements are uniform over all unitary initial constituents, including arbitrary depth-zero cuspidal support at G.

\subsection{4. The actual G quotient and raw norms on the entire generic type}\label{the-actual-g-quotient-and-raw-norms-on-the-entire-generic-type}

The whole-Schwartz construction LIFT24(MT10)--(MT16), at the holomorphic values just proved, gives target-equivariant balanced maps \[
 q_0:(\mathcal S(M_{b,a})\otimes\pi_0)_{G_a}\longrightarrow\pi_1,
 \quad
 q_1:(\mathcal S(M_{c,b})\otimes\pi_1)_{G_b}\longrightarrow\pi_2.
 \tag{MR15}
\] Its exact common-character action is \[
 \omega_\eta(g,h)f(X)=
  \eta(\det g)^{-1}\eta(\det h)
      |\det g|^{m/2}|\det h|^{-n/2}f(h^{-1}Xg).
 \tag{MR16}
\] The exponent at the continued point is \(z_*+(n-1)/2=m/2\), so these are the literal bilinear source coinvariants. MT12--MT15 continues finite-dimensional compact-fixed coordinates, not merely bidual functionals, and proves the full parabolic covariance. A nonzero map(MR15) is the actual small theta quotient by Mínguez(2008), Theorem1 and Tadic irreducibility. No map on an open-rank stratum is simply presumed to extend.

Let \[
 V_0=\pi_0^{J_G,\psi_U^{-1}},\qquad
                 a_q=\frac{|U_a(\kappa)|}{|\mathrm{GL}_a(\kappa)|}.
 \tag{MR17}
\] The unchanged first sharp raw factor is LIFT20(RG6): it is supported on integral matrices with residue \([I_a;0]u\), with value \(\psi_U(u)\). The unique allowed local coupled class is \(y=q^{-1}e_a\bmod O^b\), as in LIFT22(CG7)--(CG8). For every \(v\in V_0\), \[
                             q_0(f_U,v)=a_qs(v),
 \tag{MR18}
\] where s(v) is the identity parabolic-cell section of CG12, with s(v)(1)=v. To check this without Steinberg: the support forces the source variable into \(J_G\), with determinant residue1; the two generic characters cancel, leaving Haar volume times v. The target support is the identity cell and parabolic covariance gives all remaining values. Nothing here uses the cuspidal support of \(\pi_0\).

The section s(v) has the exact sharp middle character \(\lambda^\sharp\) of LIFT22(CG4)--(CG6): on the displayed \(\operatorname{diag}(U_a,1)\times\mathrm{GL}_{d-1}\) blocks it is \(\psi_U^{-1}\eta_0(\det h)\), and is trivial on the specified radical. This is the actual coupled middle level, not the different full generic ColV1 condition.

The terminal raw \(\Phi_y^\sharp\) is exactly LIFT21(TG7)--(TG8). Its top b-by-b block is invertible on its support. Thus its restriction to a fixed full-rank orbit has compact support, and the continued map agrees there with the ordinary compact integral. The calculation of LIFT22 gives \[
 q_1(\Phi_y^\sharp,w)(1)=v_{J^\sharp}w,\quad
 w\in\pi_1^{J^\sharp,\lambda^\sharp},\qquad
 v_{J^\sharp}=\frac{|H^\sharp|}{|\mathrm{GL}_b(\kappa)|}.
 \tag{MR19}
\] It uses only the specified compact character. Consequently \[
 \boxed{[q_1(\Phi_y^\sharp,q_0(f_U,v))(1)](1)
                                    =a_qv_{J^\sharp}v.}
 \tag{MR20}
\] This proves actual small-quotient occurrence and injectivity on all of \(V_0\), without a finite Steinberg restriction.

The exact support-cell norm calculation CG17--CG19 also only uses unitarity. Hence \[
 \|q_1(\Phi_y^\sharp,q_0(f_U,v))\|^2
  =\frac{q^b(a_qv_{J^\sharp})^2}
        {\mathcal J_{a,d}(q)\mathcal J_{b,d}(q)}\|v\|^2,
 \quad
 \mathcal J_{n,d}=|\mathrm{GL}_{n+d}(\kappa)/P_{n,d}(\kappa)|.
 \tag{MR21}
\] The raw squared norms remain \(q^{-ab}|U_a|\) and \(q^{-bc+b}|H^\sharp|\). If both raw vectors are unit-normalized, the squared multiplier is \[
 \frac{q^{ab}|U_a|}{|\mathrm{GL}_a|^2\mathcal J_{a,d}}\,
 \frac{q^{bc}|H^\sharp|}{|\mathrm{GL}_b|^2\mathcal J_{b,d}}.
 \tag{MR22}
\] These are exact local induced-model norms. The globally pinned row normalization is not thereby1, nor is an integral image lattice for arbitrary \(\pi_0\) determined.

\subsection{5. All nontrivial T characters and the other places}\label{all-nontrivial-t-characters-and-the-other-places}

For every nontrivial \(\nu\) in(MR6), checked LIFT24(MT6)--(MT8) applies to every unitary irreducible constituent with that type: \[
 \pi_0=\sigma_0\times\xi,\qquad
 \sigma_0\text{ unitary spherical of rank }a-1,\qquad
                       \xi|_{O^\times}=\nu.
 \tag{MR23}
\] The space \(\pi_0^{J_a,\nu}\) is one-dimensional. The proof uses Roche's tame type theorem, the finite parabolic invariant calculation, and Tadic's unitary blocks; the actual T primes satisfy the residual characteristic condition. It does not assume local cuspidality. An irreducible \(\pi_0\) cannot have two different nontrivial \(\nu\): its inertial support has exactly one nontrivial rank-one character, whose identity is unique by cuspidal-support uniqueness.

The unchanged first raw incidence kernel carries this line to the same \(\nu\)-line in \(\pi_1=(\sigma_0\times\eta_0\det_d)\times\xi\), by MT28. DO036's positive full-rank test applies uniformly to all \(H_a\)-fixed vectors, so it is nonzero on each of these lines. The second raw incidence kernel has exactly the graph value \[
 \lambda_T(\pi_0,\nu)w_1,\qquad
 \lambda_T(\pi_0,\nu)=
  \frac{L((d+1)/2,\sigma_1\eta_1^{-1})}{q^b-1}\ne0,\quad
                    \sigma_1=\sigma_0\times(\eta_0\det_d).
 \tag{MR24}
\] All Euler and graph/frame factors are MT29--MT33. Alternatively (MR11)--(MR13), at initial rank a−1, proves uniform holomorphy; the inserted factors are at1,\ldots,d and are unitary. The prescribed raw function is not replaced by its \(\nu\)-projector. Its actual target covariance preserves the character, so distinct characters in(MR6) cannot cancel. We retain each exact \(\lambda_T\), input norm, and compact coefficient normalization; complex nonvanishing is not a p-adic valuation statement.

DO037's Q calculation already proves injectivity of the full coupled Q quotient on every \(K_a\)-fixed vector, with its exact nonzero \(\Lambda_Q\). All other fixed finite basic vectors are nonzero on the corresponding spherical line by LIFT16. That result distinguishes relative splitting from absolute host ramification, including29. At infinity LIFT25 proves the actual consecutive-highest alignment, with source dual, dimensions, hooks, Gaussian norms and \(\Lambda_j\) conversion. These are carried checked findings used with their existing scope.

\subsection{6. Global injectivity and all automorphic multiplicities}\label{global-injectivity-and-all-automorphic-multiplicities}

The finite-level, fixed-real-type initial space has finitely many irreducible constituents and finite-dimensional multiplicity spaces. All are globally cuspidal and unitary since the positive groups are anisotropic; their finite local representations need not be cuspidal. For each constituent, DO036's first norm proof applies to every embedding and every nonzero vector of its actual initial type: DO36.15--17 already treats the entire inverse-generic G space, its T test treats all \(H_a\)-fixed vectors, and Q and basic tests are unchanged. This first edge has b\textgreater2a and uses GQT7.1 with convergent unfolding, not GQT11.4 outside its dimension window.

The first unramified theta parameter is still the original parameter plus the length-d \(\chi_0\) string. Thus the second global standard factor remains the one in LIFT26: \[
 L_E^S((d+1)/2,\operatorname{BC}_{\rm unr}(\pi_0)\chi_1^{-1})
            \prod_{r=1}^d L_E^S(r,\chi_0\chi_1^{-1}).
 \tag{MR25}
\] Only unramified parameters are used; no automorphic global base-change representation is assumed. The initial factor has the same absolute convergence margin(MR14). The boundary Hecke factor at1 is nonzero by the checked angular-type argument of LIFT14/LIFT26; absolute convergence at1 is not claimed. All fixed stage characters remain unchanged.

Local occurrence is now given at G by(MR15)--(MR20), at T by(MR23)-- (MR24), and elsewhere by the carried tests. The real converse is the explicit compact Schur integral. GQT11.4 applies on the second edge because b\textless c\textless2b and target Witt index is0: the anisotropic convergent case inside that theorem's dimension window. Therefore the second lift is nonzero on every actual embedding of every first constituent.

Apply exactly the multiplicity factorization of LIFT26(GS19)--(GS27): \[
       \Theta_\pi=A_\pi\otimes\bigotimes'_v q_{\pi,v},
                           \qquad\ker A_\pi=0.
 \tag{MR26}
\] Local Howe uniqueness gives the quotient factors; every-embedding nonvanishing makes the actual multiplicity map injective. No automorphic multiplicity-one hypothesis is used. All hypotheses have now been checked on the larger initial space.

Every local fixed coupled map \(q_{1,v}(1\otimes q_{0,v})(\Psi_v\otimes\cdot)\) is injective on the whole allowed initial type, by sections4--5. At basic places it is nonzero on a line, and at infinity it is nonzero on the fixed highest seed. Their restricted tensor product is injective, also on any subspace imposed by rational stabilizer compatibility. LIFT26(GS7a)-- (GS7c) identifies precisely this SAME fixed tensor with the literal moving-frame Y family, including \((\eta_0\eta_1^{-1})(\det h)|\det h|^d\). No middle-dependent projection has replaced it. On each initial constituent the full operator is consequently \[
                       (A_1A_0)\otimes C_{\pi_0},
                              \qquad\ker C_{\pi_0}=0.
 \tag{MR27}
\] It is injective. Distinct initial irreducible representations cannot coalesce after two edges, by the injectivity part of local Howe correspondence. Their target irreducible subspaces are orthogonal, so no cross-constituent cancellation occurs. This proves(MR4) for the whole ramified space, not merely for pure tensors or one class.

\subsection{7. Exact highest-kernel/complete-row normalization}\label{exact-highest-kernelcomplete-row-normalization}

This step is independent of the GJ extension. Use K037's compact metrics and selected Cauchy map \(\iota_j:W_j^\vee\otimes W_{j+1}\to\mathcal F_j\), normalized by the leading coefficient1 of its highest polynomial. For a coefficient row P and target seed v, K037(CR27) defines the actual complete coevaluation lift \[
 \mathcal T_j(P;v)=
   \sum_{i=1}^{d_j}
    \Theta_{j,\operatorname{pureTensorO}
                      (\phi_{j,f},\iota_j(p_i,v))}^{\rm bil}
                                         (f_{P,v_i}),
 \tag{MR28}
\] where \(v_i\) is an orthonormal basis and \(p_i\) its bilinear dual.

For any unit h and its unit dual \(h^\vee\), compact Schur averaging with Haar volume1 gives \[
 \int_{G_{j,\infty}}\rho_j(k)h\otimes\rho_j^\vee(k)h^\vee\,dk
                    =\frac1{d_j}\sum_i v_i\otimes p_i.
 \tag{MR29}
\] Indeed the integral is invariant and has trace1 under the endomorphism identification. The invariant line is the identity, whose trace is \(d_j\). This fixes coefficient and sign.

Apply(MR29) on every finite class of the genuine theta integral. The finite vector, level character, and class mass remain unchanged. Multilinearity gives \[
 \boxed{\mathcal T_j(P;v)=d_j\,
   \Theta_{j,\operatorname{pureTensorO}
                     (\phi_{j,f},\iota_j(h^\vee,v))}^{\rm bil}
                                             (f_{P,h}).}
 \tag{MR30}
\] For consecutive highest vectors h and v, the polynomial \(\iota_j(h^\vee,v)\) is exactly LIFT25's primitive principal-minor polynomial: K037's leading-minor substitution fixes its scalar. No source volume, creation norm, or finite mass is absorbed here.

Apply(MR30) twice to the same coupled finite tensor. The first complete-coevaluation output is a genuine \(W_1\)-coefficient family, so the second use has the correct middle representation and dual. The finite Y-sum commutes with both compact averages. Hence \[
 \boxed{\mathcal T_{10}^{\rm coev}(P;h_2)
                         =d_0d_1\,\mathcal B(f_{P,h_0}).}
 \tag{MR31}
\] Here \(\mathcal B\) uses source-normalized creation. For the Chern operator the right side is instead \(d_0d_1\Lambda_0\Lambda_1\mathcal B^{\rm Chern}(f_{P,h_0})\). These factors are all nonzero. Equations(MR7),(MR4),(MR31) prove injectivity of the complete raw coevaluation operator, valued in full G2 automorphic functions, on every initial \(M_{\mathrm{ram}}\) row.

For the actual restricted arithmetic class map constructed in K038(CA27)--(CA34), the precise comparison scalar is \[
 c_{10}=v_0v_1 C_\theta^{14h(ab+bc)}
     \prod_\tau(\lambda_\tau/cD_\tau)^{D_{0,\tau}+D_{1,\tau}}\ne0.
 \tag{MR32}
\] This is corrected K038(CA33), where \(v_0,v_1\) are the actual strict-covector level volumes at the two stages. They are not the coarser projective-T support volumes. Changing to a projective character realization requires both(MR7a) and the actual weighted kernel average(MR7b). Unequal chart volumes must retain their stipulated ratios, as in K037§4. K038(CA34) identifies the complete coevaluation family with its constructed arithmetic row \textbf{after the literal terminal tilQ restriction}; CA36 agrees with(MR31) after that restriction. Thus that actual row map is fully specified, but its injectivity does NOT follow from(MR4). The latter concerns the full G2-valued operator. Passing to a terminal subgroup or applying brC can still create a kernel. If an unrestricted arithmetic class realization is used instead, its own faithful all-class comparison would transfer (MR4) to characteristic zero. Such a comparison must not be silently substituted for K038's restricted one. No integral elementary-divisor or terminal branch-injectivity assertion is made here.

\subsection{8. Retained norms and the endpoint still to be detected}\label{retained-norms-and-the-endpoint-still-to-be-detected}

For every vector in this larger finite-dimensional ramified space the actual norm identity of LIFT26 remains \[
 \|\mathcal Bf\|^2
   =\sum_{\pi_0}c_{0,\pi_0}c_{1,\pi_1}
              \|(1\otimes C_{\pi_0})f_{\pi_0}\|^2,
                  \qquad c_{0,\pi_0},c_{1,\pi_1}>0.
 \tag{MR33}
\] The c's are actual Rallis and class/Haar factors, not arbitrarily chosen isometries. G retains(MR21)/(MR22), T retains(MR24), Q retains DO037's scalar, and source infinity retains AR27--30. The complete-row value(MR31) has squared norm \((d_0d_1)^2\) times(MR33); comparison with the row input also retains(MR7).

A positive least singular value exists on this fixed finite-level space by injectivity and finite dimensionality. No effective uniform global lower bound is claimed. This does not determine a DP44 gain, a saturated integral prior-image lattice, or finite valuations of complex periods.

For the bottom problem, the precise improvement is that a nonzero initial Mram vector obtained from an actual bottom detector need not be R017's vector: the same unrestricted raw coupled operator survives EVERY such input. Its prescribed terminal restriction must still be detected. The exact source interfaces clarify the remaining issue:

\begin{itemize}
\tightlist
\item
  ValLastX.lean244--251 (MainX) transports top-tail vanishing to bottom vanishing against the SAME downward image dn(ioT F). At m2 StepX is vacuous.
\item
  JointSmooth.lean264--292 (Bot and LastVan) separates all-bottom period detection from testing only that downward image.
\item
  PairForm.lean143--145 (PuNZ) asks for a nonzero Poincare image of some ramified row; it is not an assertion about R017's chosen row.
\item
  \nolinkurl{ExistsCentralBottomDatum.lean48–49} (one\_central\_pin) explicitly assumes \(1\le r_S\). At m2, \(r_S=m/2-1=0\), so this constructor cannot supply the common central pin. BottomRelationOfPin.lean101 instead applies to an already pinned datum. This is a constructor limitation, not an impossibility theorem for m2 theta lifts.
\item
  \nolinkurl{BottomTriangleNonzeroPeriod.lean36–47} is a normalized kernel identity for independently supplied Pair data. \nolinkurl{BottomNonvanishingFromInputs.lean28–31} separately assumes LNe, BotSmall, NormBig and TriEx. Faithfulness and unrestricted G2 nonvanishing do not supply them. Its uniform scalar coordinate condition retains the old host normalization; mixed application needs K035's actual varying comparison.
\end{itemize}

No implication from unrestricted G2 nonvanishing to the actual brC/tilQ/PtilFF period is used here.

\subsection{9. Sources and proof/check scope}\label{sources-and-proofcheck-scope}

New steps proved: uniform unitary holomorphy(MR9)--(MR14); the unchanged G quotient/norm on the entire generic type(MR15)--(MR22); the full-Mram global application(MR25)--(MR27); and the exact highest/coeval comparison(MR29)--(MR31). The finite cyclic decomposition is written out to match the source's unweighted norm.

Primary references:

\begin{itemize}
\tightlist
\item
  Jiang, D., \& Luo, Z. (2023). Certain Fourier operators and their associated Poisson summation formulae on GL1. Pacific Journal of Mathematics,326(2),301--372. §2B,(2-10),Theorem2.4(1)--(2), pp.313--314. \url{https://msp.org/pjm/2023/326-2/pjm-v326-n2-p05-p.pdf} Directly checked copy: work/agents/lift\_refs/jiang\_luo\_2023.pdf. This restates Godement--Jacquet, ChapterI, Theorem3.3; the original monograph is not claimed to have been read in full. Haar constants here are fixed independently.
\item
  Tadic, M. (1986). Classification of unitary representations in irreducible representations of general linear group (non-Archimedean case). Annales scientifiques de l'École Normale Supérieure,19(3),335--382. TheoremD,pp.338--339 and the characteristic-zero appendix. \url{https://www.numdam.org/article/ASENS_1986_4_19_3_335_0.pdf}
\item
  Mínguez, A. (2008). Correspondance de Howe explicite: paires duales de type II. Annales scientifiques de l'École Normale Supérieure, 41(5),717--741. Theorem1,p.718 and normalization(1.4),pp.721--722. \url{https://www.numdam.org/item/ASENS_2008_4_41_5_717_0.pdf}
\item
  Roche, A. (1998). Types and Hecke algebras for principal series representations of split reductive p-adic groups. Annales scientifiques de l'École Normale Supérieure,31(3),361--413. Theorem7.7/Remark7.8,pp.391--392;Theorem9.4,p.404;residue characteristic restriction p.379. Used through checked LIFT24. \url{https://www.numdam.org/article/ASENS_1998_4_31_3_361_0.pdf}
\item
  Gan, W. T., Qiu, Y., \& Takeda, S. The regularized Siegel--Weil formula and the Rallis inner product formula. Theorem7.1 and §§11.3--11.6, with the exact dimension ranges and conjugation dictionary checked in DO028,DO036,LIFT26.
\end{itemize}

Checked run inputs, not reproved: LIFT16(S8 and actual spherical test),LIFT20(RG6),LIFT21(TG7--TG9),LIFT22(CG4--CG21), LIFT24(MT6--MT33),LIFT25(AR5--AR30),LIFT26(GS7a--GS31), DO036/037,K035(PF22--PF29),K037(CR19,CR26--CR36). K038(CA27a,CA30a--CA30c,CA33--CA36) supplies the completed initial-level and restricted arithmetic comparison used in section7; that completed restriction is explicitly distinguished from the full G2 map.

No Lean execution or new finite computation certificate is claimed. Author whole-proof and independent reviews are recorded only after the corresponding complete text has been read.

\clearpage

\section{AM-R-021 --- the current equally spaced family fails the full budget weight windows}\label{am-r-021-the-current-equally-spaced-family-fails-the-full-budget-weight-windows}

Author: GPT-6 Astra Pro, ALP\_MACHINE /root. 22 November 2026 UTC. Reasoning effort is not externally exposed.

The concrete K031/K032 weights satisfy the analytic Hyp3 inequalities, and R018 verifies the four numerical \nolinkurl{Sizes} clauses on explicit finite growing intervals. Neither result claimed the much stronger global \nolinkurl{SizeData₂} record. This finding checks two of that record's additional conditions and proves that the current equally spaced family cannot satisfy them, even after arbitrary common determinant translations at individual CM labels. A different weight profile is required for the actual DP44 budget. This is an obstruction to this specified family, not to all mixed theta chains or all shorter lift counts.

\subsection{1. Exact conditions and a shift-invariant necessary inequality}\label{exact-conditions-and-a-shift-invariant-necessary-inequality}

The source \nolinkurl{Formalism/LocalRingCardinality/DataParams₂.lean}157--164 requires the following, writing S\_c for its positive real parameter Sc. For every CM label tau and every i=0,\ldots,k−2,

\[
                  0\le v_{\tau,i}\le S_c.
                  \tag{R21.1}
\]

For every label outside \nolinkurl{host.topPlaces}, it also requires

\[
                  S_c\le u_{\tau,0}-v_{\tau,0}
                         \le C_c S_c.
                  \tag{R21.2}
\]

Only the first inequality in (R21.2) will be used. There is no dependence on a hypothetical sign of the height term, the chosen local gain, the coefficient field, or an endpoint detector in this argument.

Allow an arbitrary real number s\_tau at each label, and replace

\[
 (u_{\tau,i},v_{\tau,i},b_{\tau,j},e_{\tau,j})
 \longmapsto
 (u_{\tau,i}-s_\tau,\ v_{\tau,i}-s_\tau,
                         b_{\tau,j}+s_\tau,\ e_{\tau,j}+s_\tau).
 \tag{R21.3}
\]

This is the common weight translation induced by a determinant twist, when such a twist exists with the required characters. We permit all real s\_tau here, a larger set than the admissible twists; therefore impossibility in this larger set also excludes every actual integral or character-compatible determinant repair.

\textbf{Necessary inequality.} At every non-top label, any translated profile satisfying (R21.1)--(R21.2) must satisfy

\[
 \boxed{\quad
 \max_i v_{\tau,i}-\min_i v_{\tau,i}
                  \le u_{\tau,0}-v_{\tau,0}.\quad}
                  \tag{R21.4}
\]

\textbf{Proof.} Subtracting the smallest translated v entry from the largest in (R21.1) gives max(v)−min(v) at most S\_c.~The common translation cancels from this difference and from u\_0−v\_0. Apply (R21.2). This proves (R21.4) with no assumption that the translated profile satisfies any of the other budget conditions. QED.

\subsection{2. Actual weights and enough unchanged non-top labels}\label{actual-weights-and-enough-unchanged-non-top-labels}

K031(HC7) and K032's corresponding baseline use

\[
 u_{\tau,i}=(2(k-i)+1)C+\tfrac12,\qquad
 v_{\tau,i}=2(k-i)C.
 \tag{R21.5}
\]

The growing K033 family adds A=11\^{}r to u\_0 at one label and 121A to u\_0 at one other label. It leaves every v entry unchanged and leaves u\_0 unchanged at all other labels. Thus at each unchanged label

\[
 \max v-\min v=2(k-2)C,\qquad u_0-v_0=C+\tfrac12.
 \tag{R21.6}
\]

There is necessarily an unchanged non-top label under any source frame indexing. In fact the literal CM label set is the disjoint union of the two images tau1(A) and tau2(A), each of size f/2. The top set is TT mapped by tau1 for a subset TT of A. Hence there are at least f/2 non-top labels. Removing the at most two growing labels leaves at least f/2−2 unchanged non-top labels. For the actual hosts these numbers are19 and82, respectively.

The exact accepted locators for that counting statement are \nolinkurl{DataHost.lean}100 (TT),129 (topPlaces),190--191 (FactPlaces), and the proven bijection \nolinkurl{tauPl_bij} in \nolinkurl{FactDelDualNabla.lean}26--45. The proof of the last theorem uses the disjoint tau1/tau2 parametrization and the frame's iota bijection. We use these accepted frame statements, not a new audit of the archive. Relabeling the two selected growing places cannot evade the cardinality argument.

Apply (R21.4) at any one of these unchanged non-top labels. It would give

\[
                         2(k-2)C\le C+\tfrac12,
                         \tag{R21.7}
\]

which is false for both concrete hosts. The full exact numbers are:

\begin{longtable}[]{@{}
  >{\raggedright\arraybackslash}p{(\columnwidth - 8\tabcolsep) * \real{0.1579}}
  >{\raggedleft\arraybackslash}p{(\columnwidth - 8\tabcolsep) * \real{0.2105}}
  >{\raggedleft\arraybackslash}p{(\columnwidth - 8\tabcolsep) * \real{0.2105}}
  >{\raggedleft\arraybackslash}p{(\columnwidth - 8\tabcolsep) * \real{0.2105}}
  >{\raggedleft\arraybackslash}p{(\columnwidth - 8\tabcolsep) * \real{0.2105}}@{}}
\toprule\noalign{}
\begin{minipage}[b]{\linewidth}\raggedright
Host
\end{minipage} & \begin{minipage}[b]{\linewidth}\raggedleft
k=f
\end{minipage} & \begin{minipage}[b]{\linewidth}\raggedleft
C
\end{minipage} & \begin{minipage}[b]{\linewidth}\raggedleft
Required lower bound for S\_c
\end{minipage} & \begin{minipage}[b]{\linewidth}\raggedleft
Required upper bound for S\_c
\end{minipage} \\
\midrule\noalign{}
\endhead
\bottomrule\noalign{}
\endlastfoot
J14 & 42 & 45156049293 & 3612483943440 & 90312098587/2 \\
J56 & 168 & 40530339687 & 13456072776084 & 81060679375/2 \\
\end{longtable}

The lower bound exceeds the upper by7134655788293/2 in the first case and26831084872793/2 in the second. This proves:

\textbf{Theorem R21.8.} No member of the specified K031/K032--K033 growing family, including every member of the R018 certified fixed-T ranges, satisfies the full literal \nolinkurl{SizeData₂} weight conditions. This remains true after every relabeling of CM places and every common per-label translation (R21.3) compatible with the indicated source frame. In particular determinant twists alone cannot turn this family into a full DP44 budget instance.

\textbf{Proof.} The counting argument supplies an unchanged non-top label. Equations (R21.6) and the necessary (R21.4) give the contradiction (R21.7), quantified exactly in the table. QED.

\subsection{3. Scope, repair direction, and framework translation}\label{scope-repair-direction-and-framework-translation}

This obstruction is independent of the earlier incompatible old quadratic \nolinkurl{Fields.Kq}/elliptic \nolinkurl{IsTen} wrapper. Even a replacement mixed host that retained the source's two-half CM index and these weight-window inequalities would encounter (R21.7). Conversely, a replacement theorem changing those inequalities would require its own quantitative proof; it could not cite \nolinkurl{SizeData₂} unchanged.

The counterexample to full admissibility does not invalidate the proved Hyp3, numerical Sizes, local selected content, actual raw theta injectivity, geometric class maps, or terminal normalization on this family. Their scopes did not include (R21.1)--(R21.2). In particular R020's fixed whole-branch content interval remains a correct calculation on the stated weights. It is not a budget certificate for a family that fails the additional weight windows.

For a new non-top profile, (R21.4) gives a concrete repair criterion: the v cloud must fit below the first u\_0−v\_0 gap. Changing only the common determinant shift never changes that ratio. One can instead alter the spacings among the v entries and the relative position of u\_0. Top labels also require their distinct \nolinkurl{top_u}, \nolinkurl{top_b}, \nolinkurl{aT_win} and \nolinkurl{win_top} conditions; this note does not assert that the repair criterion alone supplies a full profile. Every such change must be carried through the selected kernels, character compatibility, arithmetic costs and terminal detector.

In the source's representation language the translation (R21.3) is a determinant-character operation. In algebraic geometry it tensors the coefficient representation by a rank-one character, organized by SGA3 I4.7.3/4.7.3.1 and SGA6 VI3.3 as in R008. The two differences in (R21.4) are unchanged by that operation; this is a numerical weight constraint, not an equality of integral images or a new Picard/biextension theorem. The existing formal Cartier--Dieudonne, Picard, motive and ell-adic interpretations of the kernels remain those in the dictionary. No RH implication or universal negative DP44 difference follows.

Primary source citation: Alpöge, L. (2026). \emph{Remark}, step8, with its stipulated archive at commit 7cbf7a0c32fca1badf05931879a37c3f8a2b03b4, exact declarations above. Prior run input: DP\_MULT DP44, which rebuilds every required datum and does not equate the smaller numerical Sizes record with \nolinkurl{SizeData₂}. Whole-proof readings and the arithmetic certificate are recorded separately; no Lean build is claimed.

\clearpage

\section{AM-DO-039: literal split-line tests and the two local Q restriction orbits}\label{am-do-039-literal-split-line-tests-and-the-two-local-q-restriction-orbits}

Author: GPT-6 Astra Pro, ALP\_MACHINE \nolinkurl{/root/definite_oscillator}. 22 November 2026 UTC. Inherited reasoning effort is not separately exposed.

This finding isolates two exact facts about terminal detection. First, varying all integral orthogonal line insertions into the standard terminal lattice does not separate its entire coefficient representation for the actual selected weights. Second, at a Q place the unchanged coupled theta vector is zero on one induced-model H orbit and nonzero on another. The latter has an explicit H-equivariant realization with the correct terminal compact character and a nonzero normalized local test.

Neither statement decides the global terminal period. The first concerns a fixed ambient coefficient space, before its rational-stabilizer and previous-image restrictions. The second concerns functions in the actual local induced representation, not a chosen global automorphic embedding. Their exact scope prevents either an unjustified universal-spanning claim or an unjustified global vanishing claim.

\subsection{1. What the literal inserted line is}\label{what-the-literal-inserted-line-is}

All source locators below refer to commit \nolinkurl{7cbf7a0c32fca1badf05931879a37c3f8a2b03b4} in the recovered source tree. Write E for either K031/K032 host and use the standard Hermitian form \(h(x,y)=\sum_i\overline{x_i}y_i\) on \(E^N\).

\nolinkurl{Columns/StdLatUnimodular.lean}165--176 defines

\[
 \operatorname{tilN}(L)=L\oplus\mathcal O_E e_N,
 \qquad h(e_N,e_N)=1.                         \tag{DO39.1}
\]

Its T covector is defined using \(z_0=\pi_T e_N\), and its last G/Q frame coordinate is \(x_N/\pi_T\). The lattice in (DO39.1) contains the norm-one vector \(e_N\), independently of the size of T. \nolinkurl{FiniteSumNonzeroLargeDepth.lean}204--205 defines the actual \nolinkurl{tilQ} using this \nolinkurl{DecoSet.tilH} construction. Consequently an isometric realization of one of these lattices inside the standard terminal lattice supplies an integral norm-one orthogonal summand there.

This must be distinguished from the original nonstandard-coordinate \nolinkurl{TilTied} record: \nolinkurl{LastCoordinateUniformizerVectorSpansOrthogonalLi.lean} 109--115 supplies a vector of norm T in its fixed Gram matrix and requires an exact orthogonal-sum lattice equation. A collection of arbitrary primitive norm-T vectors in the standard lattice is not automatically a collection of literal (DO39.1) insertions. To give that interpretation in the standard chart, the vector must be \(\pi_T\) times an integral norm-one summand, with the stipulated complement and decorations.

Root R017(R17.7) proves

\[
 U(I_N)(F)\cap\operatorname{GL}_N(\mathcal O_E)
                   =\mu(E)^N\rtimes S_N.       \tag{DO39.2}
\]

We use that settled theorem. Its entrywise argument also proves directly that every integral norm-one vector is a root of unity times one coordinate vector. Indeed every conjugate absolute value of every entry is at most one. A nonzero integral entry has nonzero integral algebraic norm, so all its conjugate absolute values equal one; its powers lie in a bounded subset of the Minkowski lattice, making it a root of unity. The sum of squared absolute values then permits exactly one nonzero entry. Thus there are precisely N such lines in the standard lattice.

\subsection{2. The full coefficient space is not separated by these N lines}\label{the-full-coefficient-space-is-not-separated-by-these-n-lines}

Let dominant integral weights \(\lambda\) of length N and \(\mu\) of length n=N−1 interlace and satisfy the source SameShift condition \(\mu_n=\lambda_N\). Common determinant shifts do not affect the argument. In (DO39.3--9) write the shifted partitions again as lambda and mu. Let \(V_\lambda,V_\mu\) be the characteristic-zero polynomial models after removing that shift, and let

\[
 B:V_\lambda\longrightarrow V_\mu              \tag{DO39.3}
\]

be the nonzero literal branching map, with the scalar computed by R019 and its integral lattice interpretation in R020. Multiplicity-one branching, the source \nolinkurl{BrTranspose/IoHw/CoefXi} interface, and its normalization are accepted carried results. No normalization of B to one is made here; only its kernel is being considered.

For each i choose the coordinate permutation \(g_i\) interchanging \(e_i\) and \(e_N\), and put \(B_i=B\rho(g_i)\). Changing the permutation by an integral isometry preserving the last line does not change its kernel: that isometry has a block \(A\in U(I_n)\) and a root of unity z on the last line. The branch has fixed last-coordinate weight \(|\lambda|-|\mu|\), so

\[
 B\rho(\operatorname{diag}(A,z))
     =z^{|\lambda|-|\mu|}\rho_\mu(A)B.          \tag{DO39.4}
\]

Equation (DO39.4) uses the positive polynomial row action. The literal source \nolinkurl{rhoVr} precomposes that action with inverse transpose (\nolinkurl{Divisibility/CoefMod.lean}320--321); it accordingly replaces A by A\^{}\{-t\} and z by z\^{}\{-1\}. The same kernel identity follows, and coordinate transpositions are unchanged. All factors remain invertible. This shows that changing the root-of-unity generator of a coordinate line or the integral complement basis does not add another independent kernel condition. The \(\overline\pi_T^{-nBr}\) factor of literal \nolinkurl{brC} is nonzero and does not alter this assertion.

The combined map for all integral norm-one line splittings at this fixed standard coefficient value is therefore

\[
 \mathcal B_{\mathrm{split}}:V_\lambda\longrightarrow
       \bigoplus_{i=1}^{N}V_\mu,\qquad
       v\longmapsto(B_i v)_i,\quad
 \operatorname{rank}\mathcal B_{\mathrm{split}}\le N\dim V_\mu.
                                                        \tag{DO39.5}
\]

There is an elementary lower bound for its domain dimension. For every integer s with \(\mu_1\le s\le\lambda_1\), the weight \(\mu(s)=(s,\mu_2,\ldots,\mu_n)\) still interlaces \(\lambda\). These are distinct summands in the multiplicity-free restriction. The Weyl dimension formula gives

\[
 \frac{\dim V_{\mu(s)}}{\dim V_\mu}
       =\prod_{j=2}^{n}
          \frac{s-\mu_j+j-1}{\mu_1-\mu_j+j-1}\ \ge1.
                                                        \tag{DO39.6}
\]

Every denominator is positive, including when the weight has repeated entries. Summing these particular branching summands yields

\[
 \dim V_\lambda\ge(\lambda_1-\mu_1+1)\dim V_\mu.
                                                        \tag{DO39.7}
\]

Thus if \(\lambda_1-\mu_1+1>N\), the common kernel in (DO39.5) is nonzero, with dimension at least

\[
       (\lambda_1-\mu_1+1-N)\dim V_\mu.        \tag{DO39.8}
\]

For the CM tensor coefficient space use bold weights and tensor B at all f CM labels. A global coordinate line has the same coordinate index at every embedding, so there are N tests, not \(N^f\). Root-of-unity entries merely change the nonzero factors in (DO39.4). Therefore

\[
 \dim\ker\mathcal B_{\mathrm{split}}
 \ \ge\left[\prod_\tau
    (\lambda_{\tau,1}-\mu_{\tau,1}+1)-N\right]
                              \dim V_{\boldsymbol\mu}.
                                                        \tag{DO39.9}
\]

This bound concerns the full tensor coefficient module.

\subsubsection{Actual selected weights}\label{actual-selected-weights}

R019(R19.18) gives the actual weights, including their common shift. At every unchanged label,

\[
 \lambda_1-\mu_1=C,\qquad
 \begin{array}{c|r|r|r}
 \text{host}&f&N&C\\ \hline
 J^{14}&42&3562&45\,156\,049\,293\\
 J^{56}&168&3436&40\,530\,339\,687
 \end{array}                                           \tag{DO39.10}
\]

At a label with selected increase A, the difference is \(C+A\). It already exceeds N at A=0. For K033's two growing labels, the exact lower bound in (DO39.9) is

\[
 \left[(C+1)^{f-2}(C+11^r+1)(C+121\cdot11^r+1)-N\right]
                         \dim V_{\boldsymbol\mu}>0.     \tag{DO39.11}
\]

These are the already constructed selected representations. R018 gives the limited numerical \nolinkurl{Sizes} ranges, whereas the independently checked R021 proves that this equally spaced family does not satisfy the full \nolinkurl{SizeData₂} windows. The present algebraic statement neither extends those ranges nor supplies a full budget instance.

\subsubsection{Rational and integral qualifications}\label{rational-and-integral-qualifications}

Ignoring the common nonzero T scalars, the branching and permutation matrices are rational in the fixed polynomial model. Thus their common kernel is defined over the coefficient field. Its intersection with any full source Weyl lattice contains nonzero vectors, by clearing finitely many denominators. At a fixed coefficient DVR this intersection is saturated: if a nonzero scalar times a lattice vector lies in the kernel, the vector already does. Hence a nonzero intersection contains a vector primitive in the ambient local Weyl lattice. No 11-denominator is forced merely by the existence of this ambient kernel.

This does \textbf{not} show that such a vector satisfies all rational-class stabilizer relations, terminal finite characters, or lies in the actual preceding theta image. Varying decorations can also change the value of a coefficient family; (DO39.5) does not replace all such values by one vector. Other terminal lattice classes can have different orthogonal line configurations. Accordingly the result rules out the proposed universal separation of the entire standard coefficient space by its literal split lines. It is neither a counterexample to MainX nor a vanishing theorem for the constructed arithmetic row.

\subsection{3. The actual Q vector and its closed H orbit}\label{the-actual-q-vector-and-its-closed-h-orbit}

Now work at a single actual Q place. By DO037.1 the local field is \(F=\mathbf Q_q\); its nontrivial tame character is \(\chi\). Write

\[
 a=k,\quad d=t+1,\quad b=a+d,\quad N=b+d,
 \quad H=\operatorname{diag}(\operatorname{GL}_{N-1}(F),1).
                                                        \tag{DO39.12}
\]

Use the exact raw functions \(F_y\) and \(H_y\) of LIFT23(Q5)--(Q7), and DO037's normalized whole-Schwartz quotient maps \(q_0,q_1\). Their targets are \(\pi_1=\pi_0\times\eta_0\det_d\) and \(\pi_2=\pi_1\times\eta_1\det_d\), for the actual spherical unitary \(\pi_0\), with \(\eta_0|_{O^\times}
=\eta_1|_{O^\times}=\chi\). Let

\[
 C(v)=q_1\left(\sum_y H_y\otimes q_0(F_y\otimes v)\right)
                      \in\pi_2,\qquad v\in\pi_0^{K_a}.
                                                        \tag{DO39.13}
\]

This is precisely the unchanged full coupled family, not one chosen Y class or a replacement finite vector.

For every first matrix X,

\[
                         \sum_yF_y(X)=0.                \tag{DO39.14}
\]

Indeed on its support the top r=a+1 rows A have full residue rank and \(F_y(X)\) is \(\chi(\ell_A\bar z)\), with uniform higher-lift multiplicity and the stated integral lower coordinates. The nonzero linear form \(\ell_A\) has equally sized fibres, so summing reduces to a nonzero integer times \(\sum_{s\in\mathbf F_q^\times}\chi(s)=0\). Off that support every summand is zero. This is also the zero-frequency case of the checked LIFT23 Fourier formula.

Regard \(\pi_2\) in the normalized parabolic model \(\operatorname{Ind}_{P_{b,d}}^{\operatorname{GL}_N}
(\pi_1\boxtimes\eta_1\det_d)\). Then

\[
                              C(v)(h)=0\quad(h\in H).
                                                        \tag{DO39.15}
\]

To prove this, initially put the second quotient parameter in its convergent half-plane. The matrix \(h^{-1}E_b g\) has last row zero. The actual terminal factor \(H_y\) is therefore independent of y: its Fourier factor \(\psi(-\eta y)\) equals one, and every other factor is independent of y. Its sum with the first output is the operator on \(\sum_y F_y=0\), by linearity. This proves the vector integral is zero in that half-plane. The checked finite-dimensional meromorphic continuation and holomorphy at the actual quotient point then prove (DO39.15). There is no interchange of a divergent integral at the actual point in this argument.

\subsection{4. A nonzero H-equivariant open-orbit restriction}\label{a-nonzero-h-equivariant-open-orbit-restriction}

In normalized coordinates set \(c=q^{-1}e_{a+1}^*\); in the literal unscaled terminal chart set \(c=\pi_Tq^{-1}e_{a+1}^*\). Here \(\pi_T\) is a unit. Relative to block sizes \(b,d-1,1\), put

\[
 g_* =\begin{pmatrix}I_b&0&0\\0&I_{d-1}&0\\-c&0&1\end{pmatrix}.
                                                        \tag{DO39.16}
\]

Its determinant is one. The actual DO037 evaluation, retaining this unit convention, is

\[
 \operatorname{ev}_1(C(v)(g_*))=\Lambda_Q v\ne0
                    \quad\text{if }v\ne0,              \tag{DO39.17}
\]

where \(\operatorname{ev}_1:\pi_1\to\pi_0\) is the first induced model's identity evaluation, and the carried exact scalar is

\[
 \Lambda_Q=q^{-d}G(\chi,\psi_0)\chi(-1)^{-1}
 q^{-t(a+1)}\frac{\kappa_{a+1}\kappa_t}{\kappa_b}
       \prod_{i=1}^{t}(1-(\eta_0\eta_1^{-1})(q)q^{-i})^{-1}.
                                                        \tag{DO39.18}
\]

Here \(\kappa_j=\prod_{i=1}^j(1-q^{-i})\). The Gauss phase and every Haar/Euler factor are retained. Formula (DO39.17) is a carried proved value, not a new normalization.

Define the restriction along this orbit by

\[
 (R_{g_*}F)(h)=F(g_*h),\qquad h\in H.                     \tag{DO39.19}
\]

It is an H-intertwiner for the actual right-translation action. Its left covariance is explicit. Write

\[
 J=H\cap g_*^{-1}P_{b,d}g_*.
\]

A direct matrix calculation gives

\[
 J=\left\{\begin{pmatrix}A&B\\0&D\end{pmatrix}
 \in\operatorname{GL}_{N-1}(F):\ cA=c\right\},
 \quad A\in\operatorname{GL}_b,
 \quad D\in\operatorname{GL}_{d-1}.                     \tag{DO39.20}
\]

To see necessity as well as sufficiency, the b-plane \(g_*^{-1}E_bF^b\) consists of \((x,0,cx)\). An element of H preserves it precisely when its lower-left block is zero and its upper-left block satisfies \(cA=c\). For such j,

\[
 g_*jg_*^{-1}=
 \begin{pmatrix}A&B&0\\0&D&0\\0&-cB&1\end{pmatrix}.
                                                        \tag{DO39.21}
\]

Its lower d-by-d block has determinant \(\det D\). Normalized parabolic covariance therefore gives

\[
 (R_{g_*}F)(jh)=\tau(j)(R_{g_*}F)(h),\qquad
 \tau(j)=|\det A|^{d/2}|\det D|^{-b/2}
                         \eta_1(\det D)\pi_1(A).
                                                        \tag{DO39.22}
\]

Thus (DO39.19) is a specified map into the smooth function space on H with this left covariance. No compact-support condition or unmentioned normalization by the modular function of J is imposed. It is nonzero on the actual vector C(v) by (DO39.17).

This is the open orbit: a b-plane is in it exactly when it does not contain the last coordinate line and is not contained in its complementary hyperplane. Projecting to that hyperplane writes it as the graph of a nonzero functional on a b-dimensional subspace. GL\_(N−1) is transitive on these pairs. These two conditions are open and nonempty. Equivalently, \(\dim H-\dim J=bd\), the dimension of the Grassmannian. The closed orbit used in (DO39.15) consists of b-planes contained in the hyperplane.

\subsection{5. The actual terminal compact character and local mass}\label{the-actual-terminal-compact-character-and-local-mass}

Let \(U_H\subset H\) be the inverse image of \(P_{r,2t}(\mathbf F_q)\), where r=a+1. This is the literal terminal Q compact group in the standard chart. Put

\[
 \chi_H(u)=\chi(\det\overline D_u),\qquad
 v_{U_H}=\frac{|P_{r,2t}(\mathbf F_q)|}
                    {|\operatorname{GL}_{N-1}(\mathbf F_q)|}.
                                                        \tag{DO39.23}
\]

The maximal compact in H has Haar mass one. Then

\[
                 \pi_2(u)C(v)=\chi_H(u)C(v)
                              \quad(u\in U_H).          \tag{DO39.24}
\]

This can be checked directly, without identifying a row by convention. The target Weil action is \(\eta_1(\det u)\phi(u^{-1}X)\) on the compact group. The raw terminal top determinant contributes \(\chi(\det\overline A_u)^{-1}\), the last row and its Fourier factor remain fixed, and the honest Weil factor contributes \(\chi(\det\overline A_u)
\chi(\det\overline D_u)\). Their product is (DO39.24). It is the inverse of the arithmetic \nolinkurl{rowTyQ} in K038, exactly as the source inverse level-character convention requires.

Combining (DO39.17), (DO39.19) and (DO39.24) proves the actual local test

\[
 \boxed{\quad
 \int_{U_H}\chi_H(u)^{-1}
       \operatorname{ev}_1\bigl((R_{g_*}C(v))(u)\bigr)\,du
                         =v_{U_H}\Lambda_Q v.\quad}      \tag{DO39.25}
\]

It has no unknown local scalar. The normalized average divides by \(v_{U_H}\), which need not be an 11-unit. At a coefficient place over11 with q different from11, DO037.33 gives the exact 11-valuation of \(\Lambda_Q\); the terminal parabolic index adds \(v_{11}(v_{U_H})\). These are already the source local volume factors, not a declaration of a primitive integral image.

\subsection{6. Consequence and remaining global gate}\label{consequence-and-remaining-global-gate}

Equations (DO39.15) and (DO39.25) coexist: the same nonzero actual local vector has zero restriction on the closed induced-model orbit and a nonzero open-orbit H-model with the prescribed compact character. An abstract nonzero theta quotient therefore does not identify the terminal functional, but a closed-orbit evaluation zero does not prove that every H-equivariant test is zero either.

A global automorphic realization of \(\pi_2\) is not its local parabolic function model evaluated at one place. Consequently (DO39.15) is not a vanishing assertion for the global function on \(H(F)\backslash H(\mathbf A)\). Conversely (DO39.25) does not show that the actual global brC/tilQ/PtilFF period uses this open-orbit functional or that its global scalar is nonzero. That requires the prescribed global period or a suitable proved seesaw identity, with all actual class insertions and the selected archimedean mu component.

The full actual coupled arithmetic map is now constructed by K038, and DO038 gives its exact terminal mass/branch transport. LIFT27 gives injectivity of the unrestricted full-G2 map on all initial Mram vectors. The remaining issue is survival under this particular global terminal restriction and branch, not the existence of an unspecified arithmetic row map or a free choice of finite Schwartz data.

\subsection{7. Checks and carried references}\label{checks-and-carried-references}

The companion exact certificate checks a nonzero integral common branch kernel in a small Weyl module, the dimension inequalities on the actual weight formulas, the open-orbit matrix conjugation and finite orbit cardinality, and the zero Fourier frequency in small residue fields. The proofs above establish arbitrary-rank statements; the certificate is not a replacement for them or a Lean check.

Carried checked findings used without rederiving their hard inputs: R017(R17.7), R019(R19.18--21), R020(R20.2--8), DO037.1--33, DO038.9--22, K038(CA15--30,CA33--36), LIFT23(Q5--12), LIFT27(MR4,MR32). The classical branching and dimension formulas are the same polynomial representation theory used in the checked R019; see Molev, A. I. (2002), \emph{Gelfand--Tsetlin bases for classical Lie algebras}, arXiv:math/0211289v2, §2 and the interlacing basis. The new results here are the exact scope of (DO39.5--11), the local zero statement (DO39.15), the explicit stabilizer and covariance (DO39.20--22), and the unchanged-character test (DO39.25).

\clearpage

\section{AM-K039 --- an actual coupled-row divisibility bound on the fixed raw family}\label{am-k039-an-actual-coupled-row-divisibility-bound-on-the-fixed-raw-family}

Author: GPT-6 Astra Pro, ALP\_MACHINE \nolinkurl{/root/kummer_extension}, 22 November 2026 UTC; inherited reasoning effort is not separately exposed.

This finding bounds the \textbf{actual arithmetic coupled map} constructed in AM-K038, on its unchanged integral Weyl row lattices and with its actual finite G/T/Q data. It proves the integral-section premise needed to apply the already checked full-module bound AM-K027(EX2--EX4). It then subtracts explicit uniform stabilizer denominators and applies the actual terminal branch AM-R-020 and the source's integral dual-Weyl detector lattice.

The conclusion is a lower bound on image divisibility. It supplies no sharp upper bound on a prior image, no nonvanishing after terminal restriction, and no DP44 gain/error or GRH conclusion. In particular the uniformly spaced HC profiles used below fail the separately required full SizeData2 cloud inequalities, as established in the parent's current profile audit. The numerical Sizes intervals and analytic Hyp3/kernel construction remain useful checked statements at their recorded scope; they are not called full Scale data here.

\subsection{1. Actual objects, source interfaces, and the prime 11}\label{actual-objects-source-interfaces-and-the-prime-11}

There are two principally polarized hosts:

\begin{longtable}[]{@{}
  >{\raggedleft\arraybackslash}p{(\columnwidth - 14\tabcolsep) * \real{0.1333}}
  >{\raggedright\arraybackslash}p{(\columnwidth - 14\tabcolsep) * \real{0.1000}}
  >{\raggedright\arraybackslash}p{(\columnwidth - 14\tabcolsep) * \real{0.1000}}
  >{\raggedleft\arraybackslash}p{(\columnwidth - 14\tabcolsep) * \real{0.1333}}
  >{\raggedleft\arraybackslash}p{(\columnwidth - 14\tabcolsep) * \real{0.1333}}
  >{\raggedleft\arraybackslash}p{(\columnwidth - 14\tabcolsep) * \real{0.1333}}
  >{\raggedleft\arraybackslash}p{(\columnwidth - 14\tabcolsep) * \real{0.1333}}
  >{\raggedleft\arraybackslash}p{(\columnwidth - 14\tabcolsep) * \real{0.1333}}@{}}
\toprule\noalign{}
\begin{minipage}[b]{\linewidth}\raggedleft
h
\end{minipage} & \begin{minipage}[b]{\linewidth}\raggedright
CM field K\_h
\end{minipage} & \begin{minipage}[b]{\linewidth}\raggedright
A\_h
\end{minipage} & \begin{minipage}[b]{\linewidth}\raggedleft
f={[}K\_h:Q{]}/2
\end{minipage} & \begin{minipage}[b]{\linewidth}\raggedleft
a
\end{minipage} & \begin{minipage}[b]{\linewidth}\raggedleft
d
\end{minipage} & \begin{minipage}[b]{\linewidth}\raggedleft
b=a+d
\end{minipage} & \begin{minipage}[b]{\linewidth}\raggedleft
c=a+2d
\end{minipage} \\
\midrule\noalign{}
\endhead
\bottomrule\noalign{}
\endlastfoot
1 & Q(zeta28) B, B the degree-seven subfield of Q(zeta29) & J\^{}14 & 42 & 42 & 1760 & 1802 & 3562 \\
4 & K\_1(sqrt5,sqrt13) & J\^{}56 & 168 & 168 & 1634 & 1802 & 3436 \\
\end{longtable}

Here J is the Jacobian of y\^{}2=x\^{}7−1. The integral product-polarization identifications and actual CM eigenbases are K030/K032; K033 checks that the change from primitive ordinary J coordinates and its inverse are integral at11. Fix a coefficient-field place w above11 and normalize v(11)=1. We may make a finite extension of its complete DVR R to define the finitely many torsion points, phases, and frames. Such an extension is not asserted unramified. All bounds below retain this normalization.

Use the actual two edges a→b→c, the complete selected coefficient spaces of K037, and the fixed raw coupled finite datum of K038. Initial and middle classes use the \textbf{literal strict} levels of K038(CA27--CA33). At T this is the strict covector stabilizer, and at G the actual unipotent-character level. The level at29 stays maximal. Projective-T or full-G realizations, if used to describe an input, retain the actual weighted sums and matching volumes of K038(CA30a--c).

The final row is the actual evGSD row, with rhoVr coefficients. The source terminal detector lattice is exactly

\[
 \mathcal M_T=\mathrm{MtildeL}
       =(\mathrm{mtildeGSD}) .\mathrm{intNab},
 \tag{IG1}
\]

defined in \nolinkurl{Formalism/LocalToGlobal/Columns/RowForms.lean}, lines281--283. Its intDel variant is a separately named object at286. Thus the detector uses the transported dual-Weyl lattice in the opposite rhoV action. The pairing with the branch column is the literal perfect Delta/nabla pairing pairBK, as in K037(CR17--CR19), DO038 section1 and \nolinkurl{Formalism/LocalRingCardinality/FactDelDualNabla.lean}, theorem at75. No orthonormal lattice replaces either integral lattice.

\subsubsection{Lemma IG2. Every admitted Q place has residue characteristic other than11}\label{lemma-ig2.-every-admitted-q-place-has-residue-characteristic-other-than11}

The real field F\_h contains F\_0=Q(zeta28)\^{}+. The powers of11 modulo28 are

\[
                 11,9,15,25,23,1.
 \tag{IG2}
\]

None of the first five is 1 or −1. Hence the Frobenius of11 in Gal(F\_0/Q)=(Z/28)\^{}times/\{±1\} has order6. Every F\_0 place above11 has residue degree6. By the residue-field tower, every F\_h place above11 has residue degree divisible by6. Its absolute norm therefore is not a rational prime. Literal \nolinkurl{DegOne} requires that norm to be prime: \nolinkurl{Formalism/LocalToGlobal/Divisibility/Tuples.lean}, lines64--67. Consequently no such place is GoodPl, and the tuple definition at lines112--114 kills its inertAt subgroup. The exact Qprime definition at128 is the set where that subgroup is not killed. Hence this place cannot belong to the admitted Q ramification set. QED.

For a second exact check, the real cyclotomic polynomial is X\^{}6−7X\^{}4+14X\^{}2−7. Its reduction modulo11 is irreducible. The proof uses the explicit Frobenius order, not an inference from a sample of Q tuples. Both hosts are unramified at11; their ramified rational primes belong to \{2,7,29\} or \{2,5,7,13,29\}. The already fixed G and T primes also differ from11. Thus the conditional Q-over11 caveat in K038 section7 is absent for these two exact hosts and every admitted tuple.

\subsection{2. Representatives and unchanged integral frames}\label{representatives-and-unchanged-integral-frames}

\subsubsection{Lemma IG3. Standard local representatives at11}\label{lemma-ig3.-standard-local-representatives-at11}

For each of the actual initial, middle, and terminal unitary class sets, choose representatives whose finite components are the identity at every place above11. This changes neither the class problem nor its level, scalar masses, or integral coefficient lattices.

\textbf{Proof.} Let U(H) be one of the actual unitary groups over F\_h, and let g\_v be the finitely many components of a representative at v\textbar11. It suffices to find gamma in U(H)(F\_h) with gamma g\_v in K\_v, the unchanged maximal compact subgroup. Then multiply on the right by the element of the allowed compact level whose v-component is (gamma g\_v)\^{}−1 and whose other components are the identity.

Here is a direct weak-approximation argument for the required gamma. On the F\_h-vector space of H-skew-adjoint endomorphisms use

\[
             C(X)=(1-X)(1+X)^{-1},\qquad X^*=-X.
 \tag{IG3}
\]

Matrix multiplication gives C(X)\^{}*C(X)=1. Its inverse on the open set det(1+g)≠0 is X=(1−g)(1+g)\^{}−1. Every nonempty local open subset of U(H)(F\_\{h,v\}) meets this set: multiply any element by a central norm-one scalar z sufficiently close to1. The local norm-one torus has infinitely many elements in each identity neighbourhood, for example from its rational Cayley parameter. The polynomial det(1+zg) has only finitely many zeros, so one can avoid them while staying in the open subset.

Apply this to K\_v g\_v\^{}−1. Choose a Cayley preimage in each such open set. Weak approximation in F\_h, applied to a fixed F\_h-basis of the skew-adjoint vector space, supplies one rational X in all the specified inverse-image neighbourhoods. Its C(X) is the required gamma. The argument does not use a determinant-one constraint, and therefore applies also to the odd terminal rank c−1.

Rational left multiplication and allowed compact right multiplication are precisely the existing class equivalences. Their covariance and lattice transport were proved in K037(CR37--CR39) and K038(CA24--CA29). The \textbf{new representative's} local lattice at11 is standard, so its own transported delAtr or intNab lattice is the standard local Weyl or dual-Weyl lattice. We do not assert that the rational change from an earlier display of this class is integral. Its exact transport is used instead. The final tilQ decoration only adds the specified T element, an11-unit. Thus it also gives integral local frames. QED.

There are finitely many classes for these fixed ranks and levels by K036. The preceding choices can be made once, independently of the growing archimedean weight. No bound for their global heights, no new principal29 level, and no additional arithmetic-field complexity bound is asserted.

\subsection{3. Actual geometric class germs are integral local sections}\label{actual-geometric-class-germs-are-integral-local-sections}

This section proves the hypothesis that cannot be replaced by generic algebraicity or an integral Hodge splitting. In the statement, an integral local section means that its coefficient in an integral algebraic formal frame is integral. The reduced Chern frame itself need not be integral.

\subsubsection{Lemma IG4. Finite character scalars away from11 are11-units}\label{lemma-ig4.-finite-character-scalars-away-from11-are11-units}

Let eta be one of the algebraic angular Hecke characters in the actual stage splitting, including a finite-order tuple twist. If a finite idele x has component1 at all primes above11, then eta\_f(x) is an 11-unit at every chosen algebraic coefficient-field place w.

\textbf{Proof.} Principalize the ideal of x: for some positive integer m, it has m-th power (a). Since that ideal has zero valuation at all primes above11, a is a unit at all those primes. We can write x\^{}m=a u with u a finite unit idele. The restriction of the actual Hecke character to finite units has open conductor and finite image, so eta\_f(u) is a root of unity. The global relation gives

\[
               \eta_f(x)^m=\eta_\infty(a)^{-1}\eta_f(u).
 \tag{IG4}
\]

For the actual angular infinity types, eta\_infty(a) is an algebraic product of embedded a's and roots of their a bar a norms. Every embedded a is an11-unit, since a is a unit at \textbf{all} primes over11. All factors on the right of (IG4) have valuation zero. Hence so does eta\_f(x). Finite-order twists already have root-of-unity values. Determinant roots at a place where g\_v=1 may be chosen to be1; their choice-independence is the actual even-rank restriction retained in DO022/DO026. This proves the assertion for those splitting factors as well. QED.

This lemma does not claim that the finite value of an infinite-order algebraic Hecke character is a root of unity. Only its11-valuation is needed, and (IG4) proves exactly that assertion.

\subsubsection{Lemma IG5. Integral prime-to11 Heisenberg expansion}\label{lemma-ig5.-integral-prime-to11-heisenberg-expansion}

Let phi be a finite Schwartz function in the fixed rational Schrodinger realization. Suppose its component at11 is the standard vacuum and its values are algebraic and integral at w. Then it is a finite linear combination

\[
               \phi=\sum_t b_t W_t\phi_0,
               \qquad b_t\in R,
 \tag{IG5}
\]

where phi\_0 is the actual trace-one vacuum, each rational Heisenberg parameter t represents prime-to11 torsion of the base torus, and all central phases are retained. A finite extension of R is allowed.

\textbf{Proof.} In the standard position lattice Lambda, choose a positive integer D prime to11 such that phi has support in D\^{}−1 Lambda\_hat and is invariant under D Lambda\_hat. This is possible because its 11-component is the standard lattice indicator. On each coset B+Lambda\_hat, B in D\^{}−1 Lambda/Lambda, expand the function in the characters of Lambda/D Lambda. With A running over its dual group, the coefficient is

\[
 b_{A,B}=\frac1{[\Lambda:D\Lambda]}
  \sum_{x\in\Lambda/D\Lambda}
        \phi(B+x)\,\psi_f(-\langle A,x\rangle),
 \tag{IG6}
\]

up to the explicitly prescribed central phase of W\_(A,B). All character values and that phase are roots of unity. The denominator is a power of D and is an11-unit. Thus every coefficient is in R. Finite character orthogonality proves the expansion exactly on each support coset and outside the support.

The actual phi\_0 is a unit phase times a half-lattice Heisenberg translate of the standard indicator: DO016(1), with coordinates alpha=delta/2 and beta=epsilon/2, reads −i 1\_(alpha+Zhat\^{}3)(x) psi\_f(−x\^{}t beta). Its product realization on the actual hosts has the same property. Undoing this fixed translate only multiplies (IG6) by roots of unity and enlarges the denominator by a power of2. Those factors are11-units. Thus (IG5) holds relative to the \textbf{specified} trace-one vacuum, with rational parameters whose denominators divide a power of2D. They represent prime-to11 torsion. No averaging over a group of order q−1 occurs. QED.

\subsubsection{Proposition IG7. The full actual class-kernel module premise}\label{proposition-ig7.-the-full-actual-class-kernel-module-premise}

After the representative choices of Lemma IG3, every normalized geometric germ appearing in either edge of K038(CA27), at every actual class and every allowed Y, is the reduced germ of an integral formal local section of the ordinary-product principal theta line. Its splitting scalar, cD coordinates, and actual local frames add no negative11-valuation.

\textbf{Proof.} At11 the actual raw G/T/Q data are absent by Lemma IG2, and both class representatives are the identity by Lemma IG3. The finite input is therefore precisely the local vacuum. At the other places the raw values in K038 sections3--5 are zero, roots of unity, or powers of the residue cardinality times roots of unity, including the exact ftY factor q\^{}−Mb. All are integral at w. The finite Weil action preserves this coefficient integrality at places of residue characteristic different from11: its Schrodinger generators are linear changes of variable with absolute-determinant half powers, root-of-unity phases, and normalized Fourier transforms. On each compact support/period quotient the Fourier denominator and the measure are powers or square roots of that residue characteristic. These are11-units. Determinant splitting characters have valuation zero by Lemma IG4. This argument applies also at the primes2 and7; it does not assert integrality at their own residue valuations.

Equivalently the prescribed two moved-vacuum averages of DO017(8)/DO026(3) have indices prime to11, because all moved lattices agree with the standard lattice locally at11. Their exact positive square-root normalizers are therefore11-units. This is the same finite Weil realization checked in DO022 and K035(PF28), with no new vacuum normalization.

Lemma IG5 now writes each actual moved input as a sum with integral coefficients of prime-to11 Heisenberg translates of phi\_0. Under the already proved reading K035(PF25--PF28), remove the common complex period C\_theta\^{}(14h n\_j n\_(j+1)). The resulting function is exactly the same linear combination of the corresponding \textbf{actual reduced Mumford translates} of the seed product of sigma\_q's. K035 retains their central phases; it is an all-function equality before taking derivatives, so it identifies the germ already used to define the arithmetic class kernel.

LIFT08 proves that sigma\_q, normalized by leading coefficient u1u3−u2\^{}2, is primitive in the actual integral theta line at11. K034(TG3--TG10) proves that every prime-to11 torsion translate of an integral section, with precisely this reduced translation and its finite Heisenberg fibre scalar, belongs to the same complete coefficient module as an integral local section at0. It compares the normalized torsion connection, not merely the translated line. For even torsion order the same proof uses an exponent2d prime to11; the formal root and the rigidification argument are unchanged. Taking the integral linear combination proves the asserted premise. In the explicit notation of K034(TG9), the resulting coefficient belongs to

\[
            U_0\exp(-\beta_{\rm prod})\,R[[t_1,\ldots,t_g]],
\]

where the t\_i are integral algebraic formal parameters. The exponential is retained when passing to logarithmic derivative coordinates. It is not replaced by an integral Chern Taylor unit.

Finally K033(CS1--CS5) gives integral invertible actual CM coordinates, K035(PF15--PF16) gives cD\_tau an11-unit, and Lemma IG3 gives standard local class frames. The tilQ T-scaling is also an11-unit. Thus the framed functional in K037(CR39--CR40), including determinant twists, has exactly the asserted integral-module bound. QED.

The common complex period, the source unit-Fock creation factors lambda\_tau, and the literal real frame sc\_tau=lambda\_tau/cD\_tau have not been assigned finite valuations. They stay outside the arithmetic functional exactly as in K038(CA33--CA36). This proof does not claim ordinary Taylor integrality of an ample-line reduced frame. Its integral object is the algebraic local-section module, transferred using the checked torsion difference connection.

\subsection{4. A bound for every entry in the complete selected module}\label{a-bound-for-every-entry-in-the-complete-selected-module}

Write D\_(j,tau) for the selected degree at CM label tau on edge j, j=0,1, and D\_j=sum\_tau D\_(j,tau). Put n\_0=a,n\_1=b,n\_2=c and

\[
 d_{11}(n)=\begin{cases}0&n=0,\\1+\lfloor\log_{11}n\rfloor&n>0,\end{cases}
 \qquad
 E_j=n_jn_{j+1}\sum_\tau d_{11}(D_{j,\tau}),
 \qquad
 L_j=\max\{0,\lfloor D_j/1220\rfloor-E_j\}.
 \tag{IG8}
\]

\subsubsection{Theorem IG9. Actual local selected matrix-entry bound}\label{theorem-ig9.-actual-local-selected-matrix-entry-bound}

For every matched actual class pair and Y in K038,

\[
       k_j(\mathcal L_{j,c})
                       \subset 11^{L_j}\mathcal L_{j+1,c'}.
 \tag{IG9}
\]

\textbf{Proof.} By Proposition IG7 its actual germ F satisfies the checked full-module estimate K027(EX2). That estimate remains valid over the present R: its proof uses integral inverse logarithm maps and arbitrary integral formal prefactors, and every coefficientwise inequality is preserved by an integral extension. It makes no claim that this extension is unramified.

For an integral monomial differential operator of the selected multidegree, K027(EX3), obtained there from the exact Legendre formula, gives

\[
 v(\partial^\alpha F(0))
 \ge D_j/1220-\frac1{10}\sum_{\tau,r,s}s_{11}(\alpha_{\tau,r,s}).
 \tag{IG10}
\]

Every exponent at label tau is at most D\_(j,tau), and there are n\_j n\_(j+1) entries at that label. Thus its digit sum is at most 10 d\_11(D\_(j,tau)); (IG10) is at least D\_j/1220−E\_j. The full Hurwitz integrality carried in K023/K027 gives the additional lower bound0. Taking the conservative integer floor yields (IG8). No integer-valued valuation assumption on a ramified R is needed.

All polynomials produced by the actual universal subD substitution have this \textbf{label multidegree} and integral coefficients, on the whole source Weyl lattice. Their source row weights need not be the highest weight. The integral Weyl-to-Weyl theorem K037(CR8--CR14), not an abstract rational Schur projector, therefore applies to this complete family of monomial functionals. It gives (IG9), with the actual transported frames fixed in Proposition IG7. QED.

This argument uses K027(EX2--EX4), whose full proof was checked, and does not invoke the narrower fixed-source-row estimate EX6 on all weights. It introduces neither a hook-product inverse nor a Gram matrix inverse. All hook and compact metric factors remain where the analytic normalization puts them.

\subsection{5. Uniform actual class-stabilizer denominators}\label{uniform-actual-class-stabilizer-denominators}

\subsubsection{Lemma IG11. An elementary uniform11-bound}\label{lemma-ig11.-an-elementary-uniform11-bound}

For every finite stabilizer Gamma of an actual rank-n Hermitian O\_(K\_h)-lattice, including the additional level or row decorations,

\[
 v_{11}(|\Gamma|)\le S_f(n):=2m+v_{11}(m!),
                    \qquad m=\lfloor2fn/5\rfloor.
 \tag{IG11}
\]

\textbf{Proof.} The underlying integral lattice has Z-rank D=2fn, and Gamma embeds faithfully in GL\_D(Z). It is finite because the Hermitian form is positive definite at every real place: preserving the summed positive real trace form bounds all matrix entries of an integral automorphism.

The kernel of reduction GL\_D(Z)→GL\_D(F\_3) is torsion-free. Here is the needed direct proof. A nonidentity matrix in that kernel has the form A=1+3\^{}s B, s≥1, with B not identically zero modulo3. If A has prime order ell different from3, then (A\^{}ell−1)/3\^{}s is ell B modulo3, a contradiction. If ell=3, then (A\^{}3−1)/3\^{}(s+1) is B modulo3, again a contradiction; the remaining two binomial terms are divisible by3 since s≥1. Any nontrivial finite-order element has a power of prime order in the same kernel. Thus reduction is injective on Gamma, and its order divides \textbar GL\_D(F\_3)\textbar.

The exact powers 3,9,5,4,1 modulo11 give ord\_11(3)=5, and 3\^{}5−1=242=2·11\^{}2. Therefore 11 divides 3\^{}i−1 exactly when5 divides i. The elementary lifting formula for odd prime11 gives v\_11(3\^{}(5j)−1)=2+v\_11(j). It follows from the standard product \textbar GL\_D(F\_3)\textbar=3\^{}(D(D−1)/2) product\_(i=1)\^{}D(3\^{}i−1) that

\[
 v_{11}(|GL_D(F_3)|)
       =\sum_{j=1}^{\lfloor D/5\rfloor}(2+v_{11}(j))
       =2m+v_{11}(m!).
 \tag{IG12}
\]

For completeness, the used lifting formula follows by writing A=1+11\^{}e u with e≥1 and 11 not dividing u: raising A to a power prime to11 preserves e, while the binomial expansion of A\^{}11 gives valuation e+1. Factor j into its11-power and prime-to11 parts. This proves (IG12) and the lemma without enumerating classes. QED.

The rank D is the actual \textbf{2fn}, not n or fn. Taking level stabilizers or invariant-row stabilizers only makes the finite group smaller. This estimate does not assert that all actual stabilizers are prime to11, and it makes no principal congruence refinement.

\subsubsection{Theorem IG13. Actual coupled row and actual terminal scalar}\label{theorem-ig13.-actual-coupled-row-and-actual-terminal-scalar}

Let P\_cpl be K038(CA27) on the literal integral MfV lattice, or on its actual Mram sublattice. Then

\[
 \boxed{\quad
 P_{\rm cpl}(\mathcal L_0)
       \subset 11^{G_{\rm cpl}}\mathcal L_{\rm ev},
 \qquad G_{\rm cpl}=L_0+L_1-S_f(a)-S_f(b).
 \quad}
 \tag{IG13}
\]

For the growing family below and r≥13, let B\_h be the actual whole-Weyl branch lower bound of R020(R20.11):

\[
 B_1=59\,781\,009\,885\,146\,640,\qquad
 B_4=778\,040\,041\,530\,272\,708.
 \tag{IG14}
\]

For every actual terminal detector G in MtildeL of (IG1),

\[
 \boxed{\quad
 v\bigl(\operatorname{PtilFF}(P_{\rm cpl}(P),G)\bigr)
       \ge G_{\rm term}:=
          G_{\rm cpl}+B_h-S_f(c-1)
 \quad (P\in\mathcal L_0).
 \quad}
 \tag{IG15}
\]

Zero values have valuation+infinity.

\textbf{Proof.} Apply (IG9) to each of the two actual kernels in every summand of K038(CA27). Their composition maps the initial lattice into 11\^{}(L\_0+L\_1) times the target lattice. Lemma IG11 bounds the two actual denominators \textbar Gamma\_0\textbar{} and \textbar Gamma\_1\textbar, term by term. The Y sum is an actual unnormalized sum: it contributes no inverse cardinality. The nonarchimedean triangle inequality gives the same bound for all class sums. K038(CA29) puts the result in the actual evGSD equivariant subspace. Intersecting the full coefficient lattice with that subspace is saturated over the DVR, so the ambient containment gives exactly (IG13). Restricting to Mram is legitimate because it is a sublattice; no projection onto it is introduced.

R020 proves that the actual brC, including brConst/kapBr, maps the whole final Weyl lattice into 11\^{}B\_h times the terminal Weyl lattice for every r≥13. Its retained conjugate piT power is an 11-unit. Lemma IG3 supplies compatible integral local frames, so this inclusion applies to the preceding image as an inclusion. It does not require that the preceding image fill the whole Weyl lattice. PairBK takes that terminal Weyl lattice and the actual intNab detector lattice (IG1) into R. Finally K038(CA30) divides each terminal term by the actual rank-(c−1) stabilizer. Apply Lemma IG11 one more time and sum. This proves (IG15). QED.

The upper endpoint of the \textbf{whole-Weyl} branch interval in R020 has not been applied to a restricted image. Such an application would be invalid: a previous image can enter a deeper congruence lattice or the branch kernel. No nonzero vector or finite-index surjectivity is obtained from (IG13)--(IG15).

\subsection{6. Exact values on the already constructed growing weights}\label{exact-values-on-the-already-constructed-growing-weights}

Use the numerical HC/NCGaps baseline of K031/K032. At all f CM labels, let

\[
 C=\begin{cases}45\,156\,049\,293&h=1,\\40\,530\,339\,687&h=4,\end{cases}
 \qquad
 \nu_{0,i}=(2(a-i)+1)C+1-a+i-d/2
                    \quad(0\le i<a),
 \tag{IG16}
\]

and Delta=4C−d.~The exact selected recurrence is nu\_1=(nu\_0,0\^{}d)+Delta·1\_b. Define

\[
 \delta_0=\sum_i\nu_{0,i},\qquad
 \delta_1=\delta_0+b\Delta.
 \tag{IG17}
\]

The actual values, independently reconstructed from the checked K031 parameter function, are:

\begin{longtable}[]{@{}
  >{\raggedright\arraybackslash}p{(\columnwidth - 8\tabcolsep) * \real{0.1579}}
  >{\raggedleft\arraybackslash}p{(\columnwidth - 8\tabcolsep) * \real{0.2105}}
  >{\raggedleft\arraybackslash}p{(\columnwidth - 8\tabcolsep) * \real{0.2105}}
  >{\raggedleft\arraybackslash}p{(\columnwidth - 8\tabcolsep) * \real{0.2105}}
  >{\raggedleft\arraybackslash}p{(\columnwidth - 8\tabcolsep) * \real{0.2105}}@{}}
\toprule\noalign{}
\begin{minipage}[b]{\linewidth}\raggedright
Host
\end{minipage} & \begin{minipage}[b]{\linewidth}\raggedleft
delta\_0
\end{minipage} & \begin{minipage}[b]{\linewidth}\raggedleft
f delta\_0
\end{minipage} & \begin{minipage}[b]{\linewidth}\raggedleft
delta\_1
\end{minipage} & \begin{minipage}[b]{\linewidth}\raggedleft
f delta\_1
\end{minipage} \\
\midrule\noalign{}
\endhead
\bottomrule\noalign{}
\endlastfoot
J14 & 83,448,379,055,643 & 3,504,831,920,337,006 & 408,933,179,188,067 & 17,175,193,525,898,814 \\
J56 & 1,157,546,501,309,436 & 194,467,812,219,985,248 & 1,449,689,186,828,864 & 243,547,783,387,249,152 \\
\end{longtable}

The K033 growth increases the first weight at two designated labels by A=11\^{}r and B=121·11\^{}r, respectively. On \textbf{each} edge the selected partition increases in its first row by that same respective amount. Therefore

\[
 D_j(r)=f\delta_j+122\,11^r,\qquad
 \sum_\tau d_{11}(D_{j,\tau}(r))
 =(f-2)d_{11}(\delta_j)
      +d_{11}(\delta_j+11^r)+d_{11}(\delta_j+121\,11^r).
 \tag{IG18}
\]

Substitution in (IG8) gives a completely explicit integer lower bound for every selected coefficient of every actual class kernel. The stabilizer bounds do not depend on r:

\begin{longtable}[]{@{}lrrr@{}}
\toprule\noalign{}
Host & S\_f(a) & S\_f(b) & S\_f(c−1) \\
\midrule\noalign{}
\endhead
\bottomrule\noalign{}
\endlastfoot
J14 & 1,479 & 63,572 & 125,628 \\
J56 & 23,705 & 254,294 & 484,744 \\
\end{longtable}

For example:

\begin{longtable}[]{@{}
  >{\raggedright\arraybackslash}p{(\columnwidth - 10\tabcolsep) * \real{0.1304}}
  >{\raggedleft\arraybackslash}p{(\columnwidth - 10\tabcolsep) * \real{0.1739}}
  >{\raggedleft\arraybackslash}p{(\columnwidth - 10\tabcolsep) * \real{0.1739}}
  >{\raggedleft\arraybackslash}p{(\columnwidth - 10\tabcolsep) * \real{0.1739}}
  >{\raggedleft\arraybackslash}p{(\columnwidth - 10\tabcolsep) * \real{0.1739}}
  >{\raggedleft\arraybackslash}p{(\columnwidth - 10\tabcolsep) * \real{0.1739}}@{}}
\toprule\noalign{}
\begin{minipage}[b]{\linewidth}\raggedright
Host
\end{minipage} & \begin{minipage}[b]{\linewidth}\raggedleft
r
\end{minipage} & \begin{minipage}[b]{\linewidth}\raggedleft
L\_0(r)
\end{minipage} & \begin{minipage}[b]{\linewidth}\raggedleft
L\_1(r)
\end{minipage} & \begin{minipage}[b]{\linewidth}\raggedleft
G\_cpl(r)
\end{minipage} & \begin{minipage}[b]{\linewidth}\raggedleft
G\_term(r)
\end{minipage} \\
\midrule\noalign{}
\endhead
\bottomrule\noalign{}
\endlastfoot
J14 & 13 & 6,325,039,610,289 & 17,526,248,479,794 & 23,851,288,025,032 & 59,804,861,173,046,044 \\
J14 & 15 & 420,597,585,034,725 & 431,798,774,950,794 & 852,396,359,920,468 & 60,633,406,244,941,480 \\
J56 & 13 & 162,851,354,098,892 & 203,065,992,654,567 & 365,917,346,475,460 & 778,405,958,876,263,424 \\
J56 & 15 & 577,123,898,917,856 & 617,338,519,806,723 & 1,194,462,418,446,580 & 779,234,503,948,234,544 \\
\end{longtable}

These numbers are bounds on valuations normalized by v(11)=1, not logarithms of real norms. The exact certificate also records every r in the checked fixed-T numerical Sizes ranges,13--55 for J14 and13--90 for J56. The sharp selected \textbf{upper-witness} argument K033 begins at13 or15 respectively; (IG13) does not need that witness, so its J56 statement at13 is valid as a lower bound.

For the abstract growing kernel family, the digit error is O(r) at fixed ranks, whereas D\_j(r)=122·11\^{}r+O(1). Hence

\[
 \lim_{r\to\infty}
    \frac{G_{\rm cpl}(r)}{D_0(r)+D_1(r)}=\frac1{1220},\qquad
 \lim_{r\to\infty}
    \frac{G_{\rm term}(r)}{D_0(r)+D_1(r)}=\frac1{1220}.
 \tag{IG19}
\]

These are limits of the \textbf{proved lower-bound formulas}, not equalities for actual image content. Actual content can be larger, and the map or terminal scalar can vanish. Moreover fixed T only allows the finite numerical intervals already stated; no unbounded fixed-SizeData2 family follows. The new cloud-profile obstruction means that even those numerical intervals are not full DP44 Scale instances for this baseline.

\subsection{7. Exact remaining cost and image questions}\label{exact-remaining-cost-and-image-questions}

The common analytic comparison remains exactly K038(CA33--CA36):

\[
 v_0v_1 C_\theta^{14h(ab+bc)}
       \prod_\tau s_\tau^{D_{0,\tau}+D_{1,\tau}},
 \tag{IG20}
\]

with literal strict Haar volumes, actual source unit-Fock scales, and complete-coevaluation versus fixed-highest factors d\_0d\_1. Terminal mass is the further actual volume from DO038, not part of the arithmetic PtilFF valuation. These scalar, metric, and period factors have not been replaced by free normalizations.

To deduce a main gain/cost statement one still needs the corrected full parameter profile, a nonzero actual terminal value (or a quantitative statement on the restricted preceding image), and the source's complete norm/height/error accounting in the same normalization. In particular K034's sharp auxiliary first-extraction upper witness uses its explicitly chosen growing29 half-orbit. It does not furnish an upper bound for the fixed raw class map here, whose maximal29 level and complete fixed G/T/Q tensor are retained.

What is gained here is the integral-section comparison for every actual framed class kernel and a uniform bound for the actual composite and terminal maps on their original integral lattices. The previous purely formal or ambient coefficient lower bound is therefore no longer the unresolved premise for this fixed family.

\subsection{Sources and check scope}\label{sources-and-check-scope}

The new arguments are IG2--IG7, the full-label matrix application IG8--IG10, the elementary uniform stabilizer bound, and their actual coupled/terminal consequence IG13--IG19. The accepted primitive common-weight theorem is used through the already fully checked K027(EX2--EX4), not rederived and not promoted from a declaration. The standard Cayley formula, finite character orthogonality, congruence-kernel torsion argument and lifting valuation formula have their full proofs above.

The geometric translation comparison uses the proved K034(TG3--TG10), whose primary source conventions were directly checked against:

Bannai, K., \& Kobayashi, S. (2010). Algebraic theta functions and the p-adic interpolation of Eisenstein--Kronecker numbers. \emph{Duke Mathematical Journal, 153}(2), 229--295. \url{https://doi.org/10.1215/00127094-2010-024}. Author preprint \url{https://arxiv.org/pdf/math/0610163}, version4, section1.3, formula(14), Proposition2.3, Lemma2.7 and Proposition2.8. No ordinary-reduction integrality theorem from that paper is used at11 here.

All source-specific levels, scalar conventions and integral lattices are located above and in the carried K035--K038 reports. The exact certificate verifies the finite residue-degree and Fourier identities, rational Cayley form identities, degree formulas, GL-order/LTE arithmetic, and the numerical bounds. It does not enumerate unitary classes, evaluate theta coefficients or a terminal image matrix, prove nonvanishing, or run Lean. Whole-proof check events and hashes are recorded separately after completion.

\clearpage

\section{AM-R-022: a narrow-cloud replacement of the two-edge weight profile}\label{am-r-022-a-narrow-cloud-replacement-of-the-two-edge-weight-profile}

Author: /root, GPT-6 Astra Pro, ALP\_MACHINE, 22 November2026UTC. Reasoning effort is not externally exposed.

This finding replaces the equally spaced profile excluded by R021. It constructs explicit rational Harish--Chandra data, including every binary interior word, and verifies the \textbf{weight and Scale clauses} listed below. It is not a full \nolinkurl{Host}, \nolinkurl{D7}, \nolinkurl{SizeData₂}, \nolinkurl{SizesAt}, or \nolinkurl{Dat.Ok} instance. In particular the primitive-type/common-field obstruction, global gain and height records, and the terminal nonzero comparison are not removed by a numerically feasible weight profile. The numerical example also needs a new T prime; R017's old fixed primes are too small for its spread.

\subsection{1. Source clauses and parameters}\label{source-clauses-and-parameters}

All source locators use commit \nolinkurl{7cbf7a0c32fca1badf05931879a37c3f8a2b03b4}. \nolinkurl{LocalRingCardinality/Params.lean}6--80 gives the numerical ranks and \nolinkurl{HCParam}. \nolinkurl{DataParams₂.lean}9--49 defines Scale; 79--90 gives \nolinkurl{aT_win}, \nolinkurl{bT_ge}, \nolinkurl{top_u}, \nolinkurl{top_b}; 106 gives \nolinkurl{b_le}; 148--164 gives \nolinkurl{Sc_eq}, the three role windows, \nolinkurl{cloud_gap}, \nolinkurl{cloud_v}, \nolinkurl{cloud_u}, and \nolinkurl{be_close}. The binary interior formula is \nolinkurl{LocalToGlobal/Divisibility/Numbers.lean}, definition \nolinkurl{uW}. The source top/thick sets are disjoint and thin is their complement within the CM labels: \nolinkurl{DataHost.lean}129--136 and193--196.

Use the two already constructed CM hosts and ranks of K031/K032:

\begin{longtable}[]{@{}lrrrrrr@{}}
\toprule\noalign{}
Host & f=k & n=f/2 & t & d=t+1 & S=k+1+2t & N₂=k+2d \\
\midrule\noalign{}
\endhead
\bottomrule\noalign{}
\endlastfoot
J¹⁴ & 42 & 21 & 1759 & 1760 & 3561 & 3562 \\
J⁵⁶ & 168 & 84 & 1633 & 1634 & 3435 & 3436 \\
\end{longtable}

Put

\[
 \ell=28\left\lceil\frac{S^3+k+t+11}{28}\right\rceil,
 \qquad \delta=\ell,\quad
 S_c=4k\delta,\quad C_c=2,\quad h_0=2\ell,
 \quad\eta_{\rm sz}=4\ell=2h_0.                 \tag{R22.1}
\]

In both cases \(28\mid\ell\), \(S^3<\ell<2^{36}\), \(S_c<2^{46}\), and \((t+3)/2<\ell\). These are elementary integer inequalities, recorded exactly in the companion certificate. The source symbol \nolinkurl{Del} is to be assigned \(\delta\); this verifies the equality \nolinkurl{Sc_eq} if and when the remaining source record is built. It does not construct that record.

Choose any disjoint top and thick subsets of the f CM labels and let thin be the remaining labels. This includes every partition arising from the source sets just cited. There can be any number of top labels; all top labels below have the same prescribed extremes.

\subsection{2. An exact Scale whose main weight is an integer}\label{an-exact-scale-whose-main-weight-is-an-integer}

Define

\[
 A=2^{960},\quad B=2^{955}=A/32,\quad
 x=Ae^{-343},\quad
 \epsilon=\frac{\log343}{960\log2-343},\quad J=1.
                                                        \tag{R22.2}
\]

Then this is a Scale and its derived values are exactly

\[
 L=343,\quad b=e^{343},\quad a=A,\quad X'=e^7,
 \quad b'=Ae^{-7},\quad b'/x=e^{336},\quad
 w''=e^{336}/n.                                        \tag{R22.3}
\]

\textbf{Proof.} The elementary bounds \(0.69<\log2<0.7\), \(\log343<6\), proved for example by the convergent exponential series with its positive remainder bound, imply \(\log x>960(0.69)-343=319.4>52\cdot6>52\log343\). Thus \(x>1\) and \(0<\epsilon<1/52\). The definition gives \(x^\epsilon=343\); since \(343=7^3\), all the displayed Scale identities follow from their definitions. Also \(J\ge1\). Taking \(a_T=A\) satisfies both ends of the arbitrarily short \nolinkurl{aT_win} because its lower endpoint equals A exactly. Taking \(b_T=B\) satisfies \nolinkurl{bT_ge}, since \(e^7>2^7>32\) and hence \(Ae^{-7}<A/32\). No assertion about rounding a generic irrational a into a short interval is needed. ∎

For each label let the exact real role centre be

\[
 w_\tau=\begin{cases}
 B+e^{336}&\tau\text{ top},\\
 e^{343}&\tau\text{ thick},\\
 e^{336}/n&\tau\text{ thin}.
 \end{cases}                                           \tag{R22.4}
\]

Put

\[
 E_\tau=28\lceil w_\tau/28\rceil,\quad
 (e_{\tau,0},e_{\tau,1})=(E_\tau,E_\tau+2\ell),\quad
 (b_{\tau,0},b_{\tau,1})=(E_\tau-\ell,E_\tau+\ell).
                                                        \tag{R22.5}
\]

These are actual integers, specified uniquely by the displayed formula. The certificate determines their complete decimal values using rational upper and lower exponential-series bounds; it does not use a floating point rounding decision. Both e entries are multiples28 and their difference is the same \(2\ell\) at every CM label. Existence and finite conductors of compatible nonuniform stage characters are a separate character construction, not inferred from the integer formula.

\subsection{3. The full binary interior cloud}\label{the-full-binary-interior-cloud}

For \(0\le i\le k-2\) set

\[
 v_{\tau,i}=4\delta(k-1-i).                             \tag{R22.6}
\]

For every binary word \(\omega\in\{0,1\}^{k-2}\), define the interior entries, with \(1\le i\le k-2\), by

\[
 u^\omega_{\tau,i}=v_{\tau,i-1}-\tfrac12-
                             \delta\omega_{i-1}.        \tag{R22.7}
\]

At the endpoints use

\[
 (u^\omega_{\tau,0},u^\omega_{\tau,k-1})=
 \begin{cases}
 (A+(k-1)/2,-B-(k-1)/2)&\tau\text{ top},\\
 (v_{\tau,0}+S_c+1/2,2\delta+1/2)&\tau\text{ not top}.
 \end{cases}                                           \tag{R22.8}
\]

These formulas are compatible with the literal \nolinkurl{uW} rule, simultaneously at all labels. They include all words, not just one choice of each interior coefficient. If a later application uses a subset of the words, all the assertions below continue to hold on that subset.

\subsubsection{Proposition R22.9: HC, centredness, and noncompact gaps}\label{proposition-r22.9-hc-centredness-and-noncompact-gaps}

For every word and every label, (R22.5--8) is an \nolinkurl{HCParam} with the same numerical ranks as K031/K032. It satisfies the centredness and \nolinkurl{NCGaps} inequalities used in K031(HC5--HC8), including the actual natural-subtraction sign rule of \nolinkurl{Dpar}.

\textbf{Proof.} The interior interlacing gaps are exactly \(v_{i-1}-u_i=1/2+\delta\omega_{i-1}>0\) and \(u_i-v_i=4\delta-1/2-\delta\omega_{i-1}>0\). At the non-top endpoints the first gap is \(S_c+1/2>0\) and the last gap is \(2\delta-1/2>0\). At the top endpoints the first gap is positive since \(A>2^{46}>v_0\), and the last is positive since \(v_{k-2}>0>u_{k-1}\). Parity follows because k is even, all v,b,e are integers, and all u are strict half-integers. The be and eb gaps are all \(\ell\), larger than \((t+3)/2\).

Every role centre is at least \(e^{336}/84>2^{329}\), while all rank and cloud constants used here are less than \(2^{47}\). At a top label, \(u_{k-1}+b_0\ge e^{336}-\ell-(k-1)/2>0\). At other labels \(u_{k-1}>0\) and \(b_0>0\). This is the remaining HC inequality.

For centredness the distances from \(-e_j\) to u and v are larger than \((t+3)/2\); the only cancellation requiring an explicit check is the top last u entry, where \(u_{k-1}+e_0\ge e^{336}-(k-1)/2\). The distances to the two b entries are \(\ell\) or \(3\ell\). The block for the other e entry has distances at least \(2\ell-(t-1)/2>(t+3)/2\). These are precisely the complete centredness comparisons recorded in K031, and they hold for every interior word because those entries remain positive.

At stage0 the sole negative-sign entry of \nolinkurl{Dpar} is \(-b_1\). The positive-sign entries are the v entries and \(-b_0\). Their distances from \(-b_1\) are at least \(\min(2\ell,4\delta+b_1)>S^3\). At stages1 and2 the literal sign rule makes all entries positive; there are no opposite-sign comparisons. The odd-stage common shift \(-1/2\) changes no distance. This verifies every \nolinkurl{NCGaps} clause. No analytic nonvanishing is deduced just from this numerical check. ∎

\subsubsection{Proposition R22.10: all listed weight-window clauses}\label{proposition-r22.10-all-listed-weight-window-clauses}

The Scale (R22.2), top integers A,B, and (R22.1,5--8) satisfy \nolinkurl{aT_win}, \nolinkurl{bT_ge}, \nolinkurl{top_u}, \nolinkurl{top_b}, \nolinkurl{b_le}, \nolinkurl{Sc_eq}, \nolinkurl{win_thick}, \nolinkurl{win_top}, \nolinkurl{win_thin}, \nolinkurl{cloud_gap}, \nolinkurl{cloud_v}, \nolinkurl{cloud_u}, and \nolinkurl{be_close}, with positive Sc and etaSz.

\textbf{Proof.} The first three clauses were checked above. For \nolinkurl{top_b}, \(b_0-B\ge e^{336}-\ell>(k-1)/2+3\). Every e value lies between its role centre and that centre plus \(2\ell+28\le4\ell=\eta_{\rm sz}\); this gives each role window, including the allowed additional1 in the thick upper end. Every b/e difference has absolute value at most \(3\ell\), giving \nolinkurl{be_close}. The v entries lie between \(4\delta\) and \(4\delta(k-1)<S_c\). Every interior u lies strictly between0 and that same upper bound, for both choices of its word bit. At every non-top label \(u_0-v_0=S_c+1/2\in[S_c,2S_c]\). This proves every cloud clause and removes exactly R021's contradiction.

For the uniform b upper bound, \(e<4\) gives \(e^{343}<2^{686}\) and \(e^{336}<2^{672}\). At top labels \(b_1<B+2^{672}+28+2^{36}<2^{956}\); at other labels \(b_1<2^{687}\). But \(\log2<0.7\) gives \(e^{686}>2^{980}\). Thus every \(b_j<e^{686}=b^2\). The remaining equalities and positivity assertions follow directly from (R22.1). ∎

\subsection{4. Selected partitions and the numerical size of a new prime}\label{selected-partitions-and-the-numerical-size-of-a-new-prime}

Use the already checked source recurrence (R018.2), now with the actual label-dependent e entries. Write

\[
 w_{0,i}=u_i-(k-1)/2+i,\quad
 \nu_{0,i}=w_{0,i}+e_0-(k+d)/2,\quad
 \Delta=2\ell-d>0.                                    \tag{R22.11}
\]

Every \(\nu_0\) is an integral decreasing partition with positive last entry. Indeed at top labels its last entry is \(E_\tau-B-(k+d)/2>0\); at other labels it is \(E_\tau+2\delta-d/2>0\). Interlacing gives the required dominance and the exact adjacent differences are integers. The next selected partition is exactly

\[
                 \nu_1=(\nu_0,0^d)+\Delta\mathbf1_{k+d}.
                                                        \tag{R22.12}
\]

This is the same two-edge recurrence as before with its new e values; it retains every determinant shift.

The literal spread is first minus last plus rank. At top labels its three values are exactly

\[
 A+B+k,\qquad A+E_\tau+(k+d)/2,\qquad
 A+E_\tau+2\ell+(k+d)/2.                               \tag{R22.13}
\]

At a non-top label put \(W_0=v_0+S_c+1-k/2\). Then they are

\[
 u_0-u_{k-1}+1,\qquad W_0+E_\tau+(k+d)/2,\qquad
 W_0+E_\tau+2\ell+(k+d)/2.                            \tag{R22.14}
\]

The displayed selected last entries verify that each new constant block is the minimum at its stage; the first entry remains the shifted old first entry. Therefore the recurrence proves (R22.13--14) directly. They do not depend on the binary interior word. Each is strictly less than \(2A\), by the powers-of-two bounds used above, and \(2fN_2+1<4A\). The old G-size inequalities are unchanged. Consequently a norm/congruence T-prime with \(T>4A\) suffices for all four clauses of the numerical \nolinkurl{Sizes} record R018.8. This statement is conditional on choosing that new prime; it does not recycle the old fixed T values or prove the stronger \nolinkurl{AdmT}/height window.

\subsection{5. Remaining source and cost interfaces}\label{remaining-source-and-cost-interfaces}

This replacement proves the explicitly listed HC, centredness, gap, cloud and role-window conditions. The globally unramified character extension for its nonuniform e entries, a suitably large T prime, and the actual analytic/arithmetic transport must be supplied separately. Even after those steps, no full \nolinkurl{SizeData₂} is asserted here. Among its other fields are \nolinkurl{o_tie}, \nolinkurl{QS_tie}, \nolinkurl{split_tie}, the gain/error/support data, the logarithmic height equality, and the prescribed common host. The source's additional \nolinkurl{Host.o_pow}, \nolinkurl{qd_big}, and common quadratic CM-type requirements are not consequences of this weight calculation.

For example at the explicit Scale above, \(x^{\epsilon/100}=343^{1/100}<2\). Thus this particular Scale cannot satisfy \nolinkurl{o_tie} for a frame with o at least2. It is deliberately a concrete \textbf{weight-window repair}, not an assertion that all the source's global size records follow. Increasing the scale and rebuilding the host is a separate question. The terminal bottom detector and all actual DP44 gains, errors, heights and cross-family endpoint costs remain to be established in their common normalization.

The geometric interpretation is a replacement of coefficient weights in the same integral Schur functors; it changes neither the CM abelian scheme nor its Picard/biextension line. Exact programme locators remain SGA3I4.7.3/4.7.3.1, SGA6VI3.3, FGA232 Theorem3.1 pp148--149, FGA236 Corollary1.7 p229/2.3 p232, SGA7I VII2.7/2.9 pp169--174 andVIII7.4(b) pp310--311, EGAIV4 16.11.2--3(ii) pp54--55. No integral image, motivic correspondence, ell-adic weight theorem, or Kummer RH result is implied by this numerical repair.

\clearpage

\section{ALP-LIFT-28: a fixed-level holomorphic bottom detector with the exact selected jet degree}\label{alp-lift-28-a-fixed-level-holomorphic-bottom-detector-with-the-exact-selected-jet-degree}

Author: /root/lift\_accounting, GPT-6 Astra Pro, ALP\_MACHINE. 22 November 2026 UTC. Reasoning effort is not separately exposed.

This is a qualitative bottom-period construction. For either actual mixed CM host K031/K032, the literal initial compact level and its finite character extend to an explicit level on \(U(a,1)\). Its compact \(U(a)\) packet has distinct points and precisely the original rational stabilizers. The selected initial compact representation occurs in a single explicitly computed holomorphic jet degree. After increasing the last bottom parameter through a fixed arithmetic progression, cusp forms at this same finite level realize all these selected jets. Consequently every nonzero initial coefficient row has a nonzero ordinary compact bottom period against such a cusp form.

The theorem does not identify this constructed level with the literal rJ/finite frame in a completed source RealPack. It does not identify the new cusp forms with the prescribed downward theta image. There is no effective weight threshold, cusp-form norm bound, or DP44 saving here. The source's central-pin constructor requires positive rS, whereas the current m=2 problem has rS=0; that constructor is not supplied by this argument. These distinctions remain even though the finite packet and the selected real type below are actual, rather than parameterized existence interfaces.

\subsection{1. Fixed groups, initial rows, and conventions}\label{fixed-groups-initial-rows-and-conventions}

Let \(E/F\) be either K031's degree-84 CM host or K032's degree-336 CM host. Its chosen CM type has size f=42 or f=168, respectively. Let \(a=k=42\) or 168 and put

\[
 H=U(I_a),\qquad G^+=U(I_a\oplus(-1)),\qquad
 H\hookrightarrow G^+,\quad h\longmapsto\operatorname{diag}(h,1).
                                                        \tag{HJ1}
\]

At all real places H is compact and \(G^+\) has signature \((a,1)\). The Hermitian domain D is the product of the a-dimensional balls of negative lines, with base point \(o=[e_{a+1}]\).

Retain the actual initial strict compact level \(J_0\subset H(\mathbf
A_f)\) and actual right character \(\chi_0\) of LIFT27(MR4--7) and K038(CA27a,CA30a--c). At G its residue group is the full upper unipotent \(U_a(\kappa)\), with inverse generic character. At T it is the strict covector stabilizer, with trivial character. At Q, 29, and the other finite places it is the stipulated maximal basic compact group, with trivial character. This is the strict level, not the projective T level used to describe an individual nontrivial dilation eigenspace.

The positive compact highest weight at a CM label is denoted by \(\lambda_\tau\), and \(W_\lambda=\bigotimes_\tau W_{\lambda_\tau}\). Equip it with a fixed invariant positive Hermitian norm. Let \(\mathcal A_\lambda(J_0,\chi_0)\) be its actual coefficient-row space, equivalently the multiplicity space for \(W_\lambda\) in functions on \([H]=H(F)\backslash H(\mathbf A)\) with right character \(\chi_0\). All rational-stabilizer relations are imposed. Initial Mram is its specified subspace, including the unweighted T norm-kernel conditions.

For a unit vector w in \(W_\lambda\), write \(f_{P,w}\) for the scalar automorphic realization of a row P. Use real compact Haar mass one and the same finite Haar as K038/LIFT27. If \(x\) runs over the actual finite classes and \(\Gamma_x=H(F)\cap xJ_0x^{-1}\), then

\[
 \|f_{P,w}\|_{[H]}^2
   ={v_0\over d_\lambda}\sum_x{ \|P_x\|^2\over|\Gamma_x|},
 \quad v_0=\operatorname{vol}(J_0),\quad
 d_\lambda=\dim W_\lambda.                         \tag{HJ2}
\]

This is the carried compact Schur/mass identity, not a normalization of the future bottom cusp form. In the bilinear period the bottom finite character must restrict to \(\chi_0^{-1}\).

\subsection{2. The exact bottom Harish--Chandra list and jet order}\label{the-exact-bottom-harishchandra-list-and-jet-order}

Use zero-based indices. At a fixed label suppress tau and assume the source interlacing and parity conditions. In particular,

\[
 u_i>v_i>u_{i+1}\quad(0\le i<a-1),\qquad
 u_{a-1}>-b_0,\qquad b_0<e_0<b_1.
                                                        \tag{HJ3}
\]

The relevant half-integer/integer differences are at least \(1/2\). The exact definitions are \nolinkurl{Formalism/LocalRingCardinality/DataTheta.lean,} TPar.Interlace 529--533, Dpar 538--543, and Xpar 545--548, in the recovered commit 7cbf7a0c32fca1badf05931879a37c3f8a2b03b4. For m=2, j=0, the signed bottom parameter is

\[
       (p_0,\ldots,p_{a-1};q)
       =(v_0,\ldots,v_{a-2},-b_0;-b_1).             \tag{HJ4}
\]

Every positive-slot entry exceeds the negative-slot entry. There is no eta shift at j=0. Let \(\sigma\) be the discrete series with this parameter. For \(K=U(a)\times U(1)\), its minimal K type is

\[
 \tau_{\min}=W_\alpha\boxtimes z^\beta,\quad
 \alpha_i=p_i-{a-2\over2}+i,\qquad
 \beta=-b_1-{a\over2}.                             \tag{HJ5}
\]

Indeed the compact half-sum is \(\rho_c=((a-1)/2-i;0)\), while the noncompact half-sum for this chamber is \(\rho_n=((1/2)^{[a]};-a/2)\). Thus the Blattner parameter is the HC parameter plus \(\rho_n-\rho_c\), giving (HJ5). This is the usual minimal-type formula; see Harris, Kobayashi, and Speh (2025), Definition 2.7 and (2.24), printed p.16.

The initial compact type has highest weight

\[
                \lambda_i=u_i-{a-1\over2}+i.       \tag{HJ6}
\]

This is also the literal initial lamTrue: Columns/BaseLemmas.lean, lamN\_eq and lamTrue\_zero, 475--494, together with \nolinkurl{LocalRingCardinality/DataCoeff.lean41,} hcToWeightQ.

At o the holomorphic tangent representation is \(T_oD=\mathbf C^a\boxtimes z^{-1}\). The Harish--Chandra module of \(\sigma\), in the polarization with this tangent representation, has K-module structure

\[
       \tau_{\min}\otimes
       \operatorname{Sym}^{\bullet}
          (\mathbf C^a\boxtimes z^{-1}).            \tag{HJ7}
\]

This is the standard free generalized-Verma description for a holomorphic or antiholomorphic discrete series, with the name depending on the complex-structure convention. We use the explicit module in (HJ7). Its dual is realized by holomorphic sections on the negative-line ball with fibre \(\tau_{\min}^{\vee}\). Howe and Kraft (1997), opening discussion and equation (3), give this general description before their additional symplectic-group theorem. We do not invoke that latter theorem for unitary groups.

The required Pieri inequalities are exact:

\[
 \lambda_i-\alpha_i=u_i-v_i-\tfrac12\ge0
       \quad(i<a-1),\qquad
 \alpha_i-\lambda_{i+1}=v_i-u_{i+1}-\tfrac12\ge0,
\] \[
 \lambda_{a-1}-\alpha_{a-1}
                  =u_{a-1}+b_0-\tfrac12\ge0.      \tag{HJ8}
\]

Consequently \(\lambda/\alpha\) is a horizontal strip after any common determinant shift making these weights partitions. Pieri's formula therefore gives one copy of \(W_\lambda\), in precisely degree

\[
 \boxed{\quad N=|\lambda|-|\alpha|
       =\sum_{i=0}^{a-1}u_i-\sum_{i=0}^{a-2}v_i
                          +b_0-{a\over2}\in\mathbf Z_{\ge0}.\quad}
                                                        \tag{HJ9}
\]

Its extra \(U(1)\) weight is \(\beta-N\). The degree is unique because the scalar centre of \(U(a)\) counts polynomial degree. The multiplicity is one by the horizontal-strip Pieri rule (Macdonald, 1995, Chapter I, equation (5.16), printed p.73). Apply this at every label; the selected multidegree is \(\mathbf N=(N_\tau)_\tau\).

In particular, \(N_\tau\) does not involve \(b_{1,\tau}\). It is also unchanged by a common HC determinant shift \((u,v,-b_0)\mapsto(u+s,v+s,-b_0+s)\). No equal-spacing assumption on v occurs anywhere in (HJ3--9). These formulas therefore apply to a future narrow-cloud weight profile satisfying the same interlacing.

\subsection{3. An explicit bottom finite level extending the initial level}\label{an-explicit-bottom-finite-level-extending-the-initial-level}

At a relative split G place use the compatible identifications \(H(F_v)=\operatorname{GL}_a(F_v)\) and \(G^+(F_v)=\operatorname{GL}_{a+1}(F_v)\). Set

\[
 J_v^+=\operatorname{red}^{-1}U_{a+1}(\kappa_v).
                                                        \tag{HJ10}
\]

Extend the initial inverse generic character by the same additive character on all successive upper-diagonal entries, including the new last one. Denote the resulting character by \(\chi_v^+\). Then \(\chi_v^+|_{J_{0,v}}=\chi_{0,v}\).

At T, if the actual initial covector is u, take the strict stabilizer of \((u,0)\) modulo the same prime, inside \(\operatorname{GL}_{a+1}(O_v)\), with trivial character. At every other place take the stabilizer of the standard selfdual lattice in \(I_a\oplus(-1)\), again with trivial character. Denote the product by \((J^+,\chi^+)\). The form is integral unimodular at all these places; the relative CM extension in K031/K032 is unramified. At a split place this just means the displayed integral general linear compact. We have the literal equality

\[
             J^+\cap H(\mathbf A_f)=J_0,\qquad
                         \chi^+|_{J_0}=\chi_0.     \tag{HJ11}
\]

The bottom forms constructed below have right character \((\chi^+)^{-1}\), so the bilinear H period is invariant at the initial level. This inverse is essential.

At least one selected G prime has residue characteristic prime to \(|\mu(E)|\). The actual primes of K031/K032 satisfy this, and root reduction is injective there. For every finite component g, put \(\Gamma_g=G^+(F)\cap gJ^+g^{-1}\). Then

\[
                         \det\Gamma_g=\{1\}.       \tag{HJ12}
\]

To prove it, \(\det\gamma\) is a norm-one E-unit: every finite valuation is zero because every component of J+ has integral-unit determinant. At all complex embeddings it has absolute value one. Its powers are bounded in the Minkowski lattice of \(O_E\), so it is a root of unity. At the selected G place its residue is one, because the determinant of an upper unipotent residue is one; determinant is unchanged by the conjugating component g. Injective reduction of roots of unity makes the root equal to one. This proof works on every adelic component.

The construction is an explicit genuine level, not a claim that an uninstantiated source rJ record has these values. Section 8 states the remaining comparison precisely.

\subsection{4. The compact packet is embedded, with the correct stabilizers}\label{the-compact-packet-is-embedded-with-the-correct-stabilizers}

Map the finite class set

\[
        Z=H(F)\backslash H(\mathbf A_f)/J_0
      \longrightarrow
        G^+(F)\backslash(D\times G^+(\mathbf A_f))/J^+,
              \qquad [x]\longmapsto[(o,x)].         \tag{HJ13}
\]

It is finite because H is definite. It is injective. Suppose two displayed points agree. The relating rational element \(\gamma\) fixes o at every real embedding; hence it preserves the E-line \(Ee_{a+1}\), and has the form \(\operatorname{diag}(h,t)\), with \(h\in H(F)\) and \(t\in E^1(F)\). The finite level condition makes t an integral unit at every finite place. It is a root of unity by the same bounded-powers argument. At G the last diagonal residue of an element of (HJ10) is one, so t=1. Thus h gives exactly the original H double-coset equivalence, proving injectivity.

The same argument with the two representatives equal identifies the orbifold stabilizer at \([(o,x)]\) with

\[
                   \Gamma_x=H(F)\cap xJ_0x^{-1},   \tag{HJ14}
\]

including its action on the fibre and the finite character. There is no extra central U(1) stabilizer and no quotient by an additional root-of-unity group. In particular all compatible coefficient rows are legitimate equivariant jet data at these packet points. This establishes the actual finite input fit for the detector, rather than assuming that independently chosen packet values descend.

\subsection{5. The weight direction is an ample canonical twist}\label{the-weight-direction-is-an-ample-canonical-twist}

Hold u, v and b0 fixed at every label. Set, for an integer \(\ell\ge0\),

\[
 b_{1,\tau}(\ell)=b_{1,\tau}(0)+(a+1)\ell,\qquad
 \beta_\tau(\ell)=\beta_\tau(0)-(a+1)\ell.          \tag{HJ15}
\]

The compact type \(\lambda\), minimal positive weight \(\alpha\), and all jet degrees \(N_\tau\) stay fixed. The separation in (HJ4) only improves. If one also wants the elementary interlacing \(b_1<e_1\), increase e1 by the same amount. This statement alone is not a full SizeData2 assertion.

At a real place write a maximal compact element as \((A,z)\). Its determinant on the tangent space is \(\det A\,z^{-a}\); the canonical line has character

\[
       \kappa(A,z)=\det A^{-1}z^a,\qquad
       z^{a+1}=\det(\operatorname{diag}(A,z))\,\kappa(A,z).
                                                        \tag{HJ16}
\]

The bundle with fibre \((W_\alpha\boxtimes z^{\beta(\ell)})^\vee\) consequently changes by \(\kappa^\ell\) and the flat determinant character \(\det_{G^+}^{\,\ell}\). By (HJ12) the latter character is trivial on every arithmetic component group. Hence, as a holomorphic bundle on that component, the change is its canonical automorphic line to the \(\ell\)-th power. This is a statement about the quotient bundle, not an equality of maximal-compact representation weights: the flat determinant factor in (HJ16) must be retained when reading the real type. It is precisely why \(\alpha\) stays fixed.

On a neat finite cover choose a smooth projective toroidal compactification \(\overline X\) with reduced boundary \(D_\infty\) and morphism \(\pi:\overline X\to X^*\) to Baily--Borel. The canonical automorphic line extends to \(K_{\overline X}(D_\infty)\); it is the pullback of an ample line bundle on \(X^*\), after a fixed positive power if necessary for descent to a coarse quotient. These are Mumford's Main Theorem 3.1 and Proposition 3.4(a),(b), printed pp.256--257 and 264--265. Take such a fixed power q and henceforth allow \(\ell\) in one progression modulo q. Nothing depends on altering J+ or J0. The neat cover is used to prove descent statements; the constructed forms will have the original level.

\subsection{6. Cusp sections separate the selected finite jets}\label{cusp-sections-separate-the-selected-finite-jets}

Here is the finite-jet argument, including the cuspidal condition. On the neat cover let \(\overline{\mathcal E}_0\) be the canonical extension of the initial automorphic vector bundle, including the finite character after passage to its kernel. Its subcanonical extension is

\[
       \mathcal E_0^{\rm sub}
                  =\overline{\mathcal E}_0(-D_\infty).
                                                        \tag{HJ17}
\]

Its global sections are vector-valued holomorphic cusp forms. Canonical extension and the boundary condition are the established automorphic-bundle results of Mumford (1977), Proposition 3.3 and the proof of Corollary 3.5, and Lan and Stroh (2014), introduction and Section 1. In the present F-rank-one case this condition also has a direct Fourier interpretation: vanishing on the toroidal boundary removes the constant term at each rational cusp. We need no higher direct-image vanishing theorem.

Let L be the ample line on \(X^*\) furnished in Section 5 and put \(\mathcal F=\pi_*\mathcal E_0^{\rm sub}\). It is coherent by properness. Over the interior, \(\pi\) is an isomorphism, so \(\mathcal F\) is the desired automorphic vector bundle there. Let \(Z_{\mathbf N}\) be the finite infinitesimal scheme at all packet preimages imposing all Taylor coefficients through degree \(N_\tau\) in each ball variable. One may instead use all total degrees at most \(\sum_\tau N_\tau\), which contains these data. It is disjoint from the boundary.

Let \(\mathcal K\) be the kernel of the coherent restriction \(\mathcal F\to\mathcal F|_{Z_{\mathbf N}}\). For all sufficiently large n, Serre vanishing gives

\[
 H^1(X^*,\mathcal K\otimes L^n)=0.
                                                        \tag{HJ18}
\]

The long exact sequence gives the surjection

\[
 H^0(X^*,\mathcal F\otimes L^n)
       \longrightarrow
 H^0(Z_{\mathbf N},\mathcal F\otimes L^n|_{Z_{\mathbf N}}).
                                                        \tag{HJ19}
\]

The projection formula identifies the source with the cusp sections of the appropriate weight on the toroidal compactification. Serre vanishing is used only in this elementary coherent-sheaf form; see Grothendieck (1961), EGA III.1, Theorem (2.2.1)(ii), printed p.100. The same formulation is The Stacks Project, Lemma 30.17.1, equivalence (1)--(2).

Take the neat cover normal and deep enough to kill the finite character. Its finite deck group acts on (HJ19), on the line linearization, and on all the packet preimages. Averaging in characteristic zero makes passage to its specified character and invariant subspace exact. Thus (HJ19) descends to the original level and gives every jet satisfying precisely the stabilizer conditions (HJ14). This does not refine the actual bottom level.

At a packet point the selected homogeneous quotient of the jet space is

\[
 \bigotimes_\tau
 \left[(W_{\alpha_\tau}\boxtimes z^{\beta_\tau})
       \otimes\operatorname{Sym}^{N_\tau}
                 (\mathbf C^a\boxtimes z^{-1})\right]^\vee
    \supset W_\lambda^\vee\boxtimes
                       z^{-\boldsymbol\beta+\mathbf N}.
                                                        \tag{HJ20}
\]

The direct summand has multiplicity one by (HJ8--9). Its restriction to the actual stabilizers is exactly the dual initial coefficient module. Projecting (HJ19) to this summand remains surjective. For all sufficiently large n, therefore, all its compatible packet values arise from genuine cusp forms of the required discrete series \(\sigma_\ell^\vee\).

For completeness, the last type assertion uses the holomorphic discrete-series realization, not a dimensional guess. The minimal K type in the section construction is \(\tau_{\min,\ell}^\vee\), and the antiholomorphic lowering algebra annihilates its section vectors. Their generated Harish--Chandra modules are the corresponding generalized-Verma quotients. In the regular discrete-series chamber of (HJ4) this is the irreducible holomorphic discrete-series module, whose free polynomial K structure is (HJ7) dual. The cusp forms are in \(L^2\) and decompose into copies of that representation. Large ell remains in this chamber. We impose no Hecke eigenvector or irreducible finite-adelic constituent restriction.

\subsection{7. From jets to the actual compact restriction and period}\label{from-jets-to-the-actual-compact-restriction-and-period}

One must distinguish a jet of the minimal-type section from evaluation of its selected higher K-type vector. This distinction does not spoil surjectivity, but carries a real normalization factor.

Fix at each label a nonzero Pieri inclusion \(j_\tau:W_{\lambda_\tau}\to
W_{\alpha_\tau}\otimes\operatorname{Sym}^{N_\tau}T_oD\). The Poincaré--Birkhoff--Witt map into \(\sigma_\ell\) is injective in this discrete-series range by (HJ7). Its restriction to \(j_\tau\) is therefore nonzero. After choosing invariant Hilbert norms it is a nonzero scalar times the isometric selected K-type inclusion. Denote the scalar by \(b_{\ell,\tau}\ne0\); its absolute value is the norm of that fixed PBW image of a unit vector. No value one is assigned.

The transpose operation on holomorphic cusp sections is a differential operator of order \(N_\tau\). Its leading symbol is the transpose of \(j_\tau\), times the inverse of the selected normalization scalar (and the sign dictated by the chosen Lie-action convention). It is nonzero. This can be checked in Harish--Chandra coordinates: each noncompact Lie generator acts by a first-order operator whose principal part is the corresponding constant directional derivative at o; a product of \(N_\tau\) generators has the symmetric product as principal symbol. Lower-order terms come only from differentiating the frame multiplier. The prescribed Pieri projection gives precisely the indicated selected symbol. This argument also proves that no different derivative order can substitute for (HJ9).

Prescribe all lower jets to be zero, and prescribe the selected multidegree jet with the inverse of this nonzero symbol factor. Equation (HJ19) then realizes any desired selected higher-K-type coefficient value at every packet point. More generally its triangular lower-jet correction can be inverted because the top symbol is surjective onto this summand. Thus the actual restriction map

\[
 \operatorname{Hom}_{G^+_\infty}
       \bigl(\sigma_\ell^\vee,
         \mathcal A_{\rm cusp}(G^+,J^+,(\chi^+)^{-1})\bigr)
  \longrightarrow
        \mathcal A_{\lambda^\vee}(J_0,\chi_0^{-1})
                                                        \tag{HJ21}
\]

is surjective for all ell sufficiently large in the progression of Section 5. The map uses the fixed isometric inclusion of the selected \(W_\lambda^\vee\) in \(\sigma_\ell^\vee\), then the literal restriction along (HJ1). It does not insert a rational or real translation to force a nonzero value.

In particular let \(P\ne0\) be any initial row, including any nonzero Mram row, and w any unit seed. Choose the target row in (HJ21) to be the Hermitian conjugate of P. The resulting actual cusp form F in the selected \(W_\lambda^\vee\) has

\[
                F|_{[H]}=\overline{f_{P,w}}.       \tag{HJ22}
\]

Its bilinear bottom period is exactly

\[
 \boxed{\quad
   \int_{[H]}f_{P,w}(h)F(\operatorname{diag}(h,1))\,dh
       ={v_0\over d_\lambda}
             \sum_x{ \|P_x\|^2\over|\Gamma_x|}
       >0.\quad}                                  \tag{HJ23}
\]

This equality is an interpolation prescription on a freely chosen bottom cusp form; it is not a change of normalization of either prescribed theta kernel. The factors \(b_{\ell,\tau}\), the norm of a right inverse of (HJ19), and the Petersson norm of F are not estimated. Normalizing F to norm one divides (HJ23) by its actual positive norm; there is no uniform lower bound for that normalized period asserted here.

For one sufficiently large ell, (HJ21) is surjective on the entire finite initial coefficient space, so the same ell works for all initial rows, not merely one chosen vector. Taking Hilbert adjoints, the corresponding period/Poincaré operator is injective on that finite coefficient space. This last statement refers to the explicitly constructed period operator (HJ21), not automatically to the source operator D.s.Pu.

The same conclusion is uniform over any fixed finite collection of admissible word profiles: use the union of their finite jet conditions and the maximum of their finitely many Serre thresholds. When v and b0 are fixed across the words, the minimal-type bundle is already common and only the selected jet degrees/components vary. This is a qualitative uniform statement over the source's finite word set; it does not estimate that maximum.

\subsection{8. Exact remaining gate to the source endpoint}\label{exact-remaining-gate-to-the-source-endpoint}

The current result supplies an actual bottom representation, finite level, finite character, compact packet, selected type, and nonzero bottom period. It removes the need to guess that some initial vector is compatible with this particular geometric detector: every compatible row is detected.

The following identifications have not been proved here.

\begin{enumerate}
\def\labelenumi{\arabic{enumi}.}
\tightlist
\item
  The displayed \((J^+,(\chi^+)^{-1})\), lattice frame and coordinate embedding have not been identified with rJ/rChi and the prescribed finite translate in a completed actual RealPack. PairForm.lean42--45 defines rJ via dat0, betaY and levR; these pins contain more data than (HJ11).
\item
  A cusp form furnished by (HJ21) has not been shown to lie in the actual prescribed downward theta image from the terminal input. R016's free-data seesaw and LIFT27's full-G2 injectivity do not supply this.
\item
  Even an abstract identification of the full bottom cusp space does not construct BotDat.Pin, ValLastX, or the same-family ChainOnto/LastX transport needed by MainX. The existing \nolinkurl{ExistsCentralBottomDatum.lean,} \nolinkurl{MatchingTranslatesBottom.lean} and \nolinkurl{PlacedBesselCompatibility.lean} constructors used in this route explicitly assume \(1\le rS\); m=2 has rS=0.
\item
  The threshold for ell is ineffective in this proof. It has not been fitted into the b/e and SizeData2 windows. Root R021 separately rules out the previously equally spaced family as a full SizeData2 instance; no consequence here repairs that obstruction. No derivative gain, saturation, global row content, or complete H cost is inferred from (HJ23).
\end{enumerate}

Thus a source PuNZ conclusion would require the actual detector identification in item 1 together with the fixed real/Poincaré normalizations; a terminal nonzero conclusion additionally requires items 2--3. These are concrete comparisons, not an assumption that all bottom cusp forms are automatically in the stipulated image.

\subsection{9. Sources, checks, and status}\label{sources-checks-and-status}

Primary mathematical inputs used with the indicated scope:

\begin{itemize}
\tightlist
\item
  Harris, M., Kobayashi, T., \& Speh, B. (2025, September 21). Translation functors, branching problems, and applications to the restriction of coherent cohomology of Shimura varieties. Definition 2.7, equation (2.24), printed p.16; the opposite holomorphic convention and contragredient rule are explicit in Sections 10.3.1 and 11.2. \url{https://www.ms.u-tokyo.ac.jp/~toshi/texpdf/tk2025l-HKS_250921.pdf}
\item
  Howe, R., \& Kraft, H. (1997). Principal covariants, multiplicity-free actions, and the K-types of holomorphic discrete series. In Geometry and representation theory of real and p-adic groups (pp.147--161). Birkhäuser. Opening general discussion, equation (3), preprint pp.1--2. \url{https://doi.org/10.1007/978-1-4612-4162-1_8} The original paper's opening text was directly inspected in the transcription hosted by Academia.
\item
  Macdonald, I. G. (1995). Symmetric functions and Hall polynomials (2nd ed.). Oxford University Press. Chapter I, (5.16), p.73. \url{https://math.berkeley.edu/~corteel/MATH249/macdonald.pdf}
\item
  Mumford, D. (1977). Hirzebruch's proportionality theorem in the non-compact case. Inventiones Mathematicae, 42, 239--272. Main Theorem 3.1; Propositions 3.3--3.4; proof of Corollary 3.5. \url{https://www.dam.brown.edu/people/mumford/alg_geom/papers/1977b--HirzProp-GT.pdf} The relevant original text was inspected through its paperzz transcription because the author and repository PDF hosts returned errors in this session; no PDF inspection is claimed.
\item
  Lan, K.-W., \& Stroh, B. (2014). Relative cohomology of cuspidal forms on PEL-type Shimura varieties. Algebra \& Number Theory, 8(8), 1787--1799. Introduction and Section 1 for the subcanonical cusp-form interpretation; its stronger direct-image theorem is not needed. \url{https://msp.org/ant/2014/8-8/ant-v8-n8-p01-s.pdf}
\item
  The Stacks Project Authors. Ample invertible sheaves and cohomology, Lemma 30.17.1 (Tag 0B5U; contained in Section 01XO). \url{https://stacks.math.columbia.edu/tag/0B5U}
\item
  Grothendieck, A. (1961). Éléments de géométrie algébrique: III. Étude cohomologique des faisceaux cohérents, première partie. Publications Mathématiques de l'IHÉS, 11, 5--167. Theorem (2.2.1)(ii), p.100. \url{https://www.numdam.org/item/PMIHES_1961__11__5_0.pdf} The original PDF was recovered and the theorem directly inspected in its extracted text, work/agents/lift\_refs/ega\_iii\_1.txt.
\end{itemize}

Carried accepted run inputs: K031/K032 actual hosts and finite primes; R015's source signed HC convention; K038's strict initial row level and character; LIFT27's realization of full Mram and exact compact mass; R021's separate failure of the previous weight cloud.

Author whole-proof review completed 2026-11-22T02:29:00+00:00: PASS in the stated qualitative fixed-level bottom-detector scope. Every section and HJ1--23 was read; the HC signs, exact Pieri degree, root-of-unity packet/stabilizer argument, determinant-flat correction, coherent jet exact sequence, finite-level descent, and higher-K-type principal symbol were checked. The actual source initial lamTrue formula was also directly checked. No numerical Serre threshold, finite norm bound, source rJ identification, or Lean verification is claimed.

\clearpage

\section{AM-K040 --- nonuniform CM stage characters with fixed finite data}\label{am-k040-nonuniform-cm-stage-characters-with-fixed-finite-data}

Author: GPT-6 Astra Pro, ALP\_MACHINE \nolinkurl{/root/kummer_extension}, 22 November 2026 UTC; inherited reasoning effort is not separately exposed.

This finding removes the character-existence restriction to uniform archimedean weights in K031/K032. On either actual host, every vector e in (28 Z)\^{}Psi has an everywhere finite-unramified unitary Hecke character alpha\_e, trivial on the full base-field idele group, with angular infinity exponents −2e\_tau. The construction is coherent in e and has the exact uniform-direction identity

\[
                  \alpha_{c\mathbf1}=\eta^{-2c}
                     \qquad(c\in28\mathbb Z).
 \tag{NC1}
\]

Multiplying by the already constructed zero-infinity G character therefore supplies the actual stage characters for arbitrary such nonuniform profiles, without adding finite ramification or changing the prescribed tame G characters. One finite Galois coefficient field works for all profiles. The proof explicitly handles compatibility with principal ideals extended from the real field; checking units alone would not establish the required base-field restriction.

This is an arithmetic character theorem. It does not establish the new profile's cloud inequalities, full SizeData2, selected content, terminal nonvanishing, or the remaining DP44 cost window.

\subsection{1. Hosts and literal infinity convention}\label{hosts-and-literal-infinity-convention}

Use the fields and actual CM types of K029--K032:

\[
 E_0=\mathbb Q(\zeta_{28}),\quad
 K_1=E_0B,\quad [B:\mathbb Q]=7,\quad B\subset\mathbb Q(\zeta_{29}),
 \qquad K_4=K_1(\sqrt5,\sqrt{13}).
 \tag{NC2}
\]

Let F\_h=K\_h\^{}+ and let Psi\_h be the actual induced CM type of the principal variety J\^{}(14h), with f=42h labels. Both CM extensions are unramified at every finite place. Retain the actual odd character eta\_h from K029 and its norm pullback in K032:

\[
 (\eta_h)_\infty(z)=\prod_{\tau\in\Psi_h}z_\tau/|z_\tau|,
 \qquad \eta_h|_{\mathbb A_{F_h}^\times}=\eta_{K_h/F_h}.
 \tag{NC3}
\]

The right side is quadratic, and eta\_h is everywhere finite-unramified. Its normalized etaVal values belong to the already constructed fixed field. This accepted result is not rederived here.

The exact source definitions are \nolinkurl{AbelianHostFieldConditions.lean}, lines314--331: IsMu0 requires unitary, TrivOn F, HasInfQ(−e), unramified away from the two G primes, tame everywhere, and the prescribed inverse residue character at each selected G prime. Chars.Tied additionally requires all unramified finite values in one coefficient field and retains the same etaT/etaVal. In \nolinkurl{Formalism/LocalToGlobal/UnitaryNorms/IsUnitaryNormOf.lean}, lines90--94, HasInfQ(a) means the \textbf{angular integer exponent 2a}. Thus the needed infinity component for e is

\[
 b_e(z)=\prod_{\tau\in\Psi_h}
              (z_\tau/|z_\tau|)^{-2e_\tau}.
 \tag{NC4}
\]

All e\_tau below may be positive, negative, or zero. They are multiples of28. No uniformity over labels is imposed.

\subsubsection{Root-of-unity fact}\label{root-of-unity-fact}

\[
                         \mu(K_h)=\mu_{28}\quad(h=1,4).
 \tag{NC5}
\]

For K\_1 this is K029(ER4--ER5). For K\_4, the added ramification indices at5 and13 are2, whereas the indices at2,7,29 remain2,6,7. A primitive ell\^{}a-th root would force phi(ell\^{}a) to divide the rational ramification index at ell. Thus neither fifth nor thirteenth roots can occur; an eighth, forty-ninth, or twenty-ninth root is excluded by the same test. All other odd primes are unramified. The order therefore divides28, and E\_0 already contains mu28. This proves (NC5).

\subsection{2. Exact principal/base-ideal compatibility}\label{exact-principalbase-ideal-compatibility}

Write I\_K for the fractional ideal group, P\_K for its principal subgroup, and j:I\_F→I\_K for extension of ideals. Define

\[
 H=P_Kj(I_F),\qquad
 C_{K/F}=I_K/H=\operatorname{Cl}(K)/j\operatorname{Cl}(F).
 \tag{NC6}
\]

This quotient is finite; no value of either class number is assumed. For a in K\^{}times put

\[
 T_e(a)=\prod_{\tau\in\Psi}
                    \tau(a/\bar a)^{e_\tau}\in K^\times.
 \tag{NC7}
\]

Here K is Galois over Q, so all its embeddings have image K after the fixed complex embedding. Each complex value of T\_e(a) has absolute value1. Equation (NC4) gives T\_e(a)=b\_e(a)\^{}−1.

\subsubsection{Lemma NC8}\label{lemma-nc8}

If (a) belongs to j(I\_F), then T\_e(a)=1. Consequently

\[
           c_e((a)j(\mathfrak b))=T_e(a)
                    \qquad((a)j(\mathfrak b)\in H)
 \tag{NC8}
\]

is a well-defined homomorphism H→K\^{}times, trivial on j(I\_F), with values of complex absolute value1.

\textbf{Proof.} If (a)=j(b), that ideal is fixed by complex conjugation. Thus (a/bar a)=1, and a/bar a is a global unit. Every complex embedding of this unit has absolute value1. The elementary CM unit argument, proved in K029 section1, makes it a root of unity: the monic polynomials of its positive powers have bounded integer coefficients, so two powers coincide. By (NC5) it belongs to mu28. Each exponent e\_tau is divisible by28; every factor of T\_e(a) is therefore1, proving the first assertion.

If (a)j(b)=(a')j(b'), then (a/a')=j(b'/b). Apply the first assertion to a/a' to obtain T\_e(a)=T\_e(a'). This proves well-definedness of (NC8). Multiplicativity follows from (NC7), and a=1 proves its triviality on j(I\_F). QED.

In particular this handles any capitulating real ideal class. It does not first extend a character on P\_K and then presume that the resulting restriction to F is trivial. The compatibility is checked on the full intersection P\_K intersect j(I\_F) before extension.

\subsection{3. Explicit extension across the finite relative class quotient}\label{explicit-extension-across-the-finite-relative-class-quotient}

Choose an invariant-factor decomposition

\[
 C_{K/F}=\bigoplus_{s=1}^q\mathbb Z/d_s\mathbb Z,
 \qquad d_s>1,
 \qquad H_{\rm rel}=\prod_s d_s,
 \quad d_{\rm rel}=\operatorname{lcm}_s d_s.
 \tag{NC9}
\]

For the trivial group use H\_rel=d\_rel=1. Choose fractional ideals A\_s representing the generators, and fixed relations

\[
                  A_s^{d_s}=(a_s)j(\mathfrak b_s).
 \tag{NC10}
\]

Choose z\_(s,e) with z\_(s,e)\^{}(d\_s)=T\_e(a\_s). Its chosen complex value has absolute value1, as do all d\_s-th roots. Every ideal can be written uniquely in quotient-exponent ranges as

\[
 I=h\prod_s A_s^{r_s},\qquad h\in H,\quad0\le r_s<d_s.
\]

Define

\[
                 \widetilde c_e(I)
                    =c_e(h)\prod_s z_{s,e}^{r_s}.
 \tag{NC11}
\]

This is a homomorphism: when multiplying two ideals, each carry of d\_s moves A\_s\^{}(d\_s) into H and contributes exactly the same factor T\_e(a\_s)=z\_(s,e)\^{}(d\_s). It extends c\_e and is trivial on j(I\_F). Every value belongs to the field generated by K and the finitely many chosen roots.

For a finite idele x\_f set I(x\_f)=product\_P P\^{}(v\_P(x\_P)). Then

\[
             \alpha_e(x)=b_e(x_\infty)\widetilde c_e(I(x_f))
 \tag{NC12}
\]

is continuous, unitary, everywhere finite-unramified, and trivial on K\^{}times. The last assertion follows from c\_e((a))=b\_e(a)\^{}−1. Continuity follows because the finite ideal map is locally constant and b\_e is continuous. For an F-idele, its finite ideal extends from F and so has value1. At infinity all its components are real, and the even angular exponents in (NC4) make b\_e identically1 there. Thus

\[
                   \alpha_e|_{\mathbb A_F^\times}=1.
 \tag{NC13}
\]

Equations (NC8)--(NC13) prove existence with the exact base restriction. They use only the finite relative ideal quotient, not class number one of the enlarged host or a character extension with unverified capitulation compatibility.

\subsection{4. A coherent family with the prescribed odd character}\label{a-coherent-family-with-the-prescribed-odd-character}

Independent choices in (NC11) would not automatically respect addition of weight vectors. We now make coherent choices and enforce the exact uniform relation (NC1).

Enumerate Psi=\{tau\_1,\ldots,tau\_f\}. For i\textless f, construct once a character xi\_i using (NC12) for the basis vector e=28 e\_i. Put

\[
             \xi_f=\eta_h^{-56}\prod_{i<f}\xi_i^{-1},
 \qquad
             \alpha_e=\prod_{i=1}^f\xi_i^{e_{\tau_i}/28}.
 \tag{NC14}
\]

Since eta\_h has the quadratic base restriction (NC3), its even power eta\_h\^{}−56 is trivial on the full F-idele group. Therefore xi\_f is also trivial there. All xi\_i are finite-unramified and unitary. The infinity exponent of eta\_h\^{}−56 is−56 at every label; dividing by the first f−1 basis characters leaves exponent−56 exactly at tau\_f.~Thus xi\_f has the required last basis type, and alpha\_e has (NC4) at every label. All finite conditions remain as proved above.

The construction is a homomorphism from the free weight lattice:

\[
       \alpha_{e+e'}=\alpha_e\alpha_{e'},\qquad
       \prod_i\xi_i=\eta_h^{-56},\qquad
       \alpha_{c\mathbf1}=\eta_h^{-2c}\quad(c\in28\mathbb Z).
 \tag{NC15}
\]

In particular, for two nonuniform profiles with

\[
       e_{1,\tau}-e_{0,\tau}=2\ell
                     \quad\text{for every tau},\qquad2\ell\in28\mathbb Z,
\]

one has the \textbf{identity of full Hecke characters}

\[
                         \alpha_{e_1}/\alpha_{e_0}=\eta_h^{-4\ell}.
 \tag{NC16}
\]

It is not merely an equality of their infinity components up to an unspecified finite class character.

\subsection{5. The actual fixed-G stage characters}\label{the-actual-fixed-g-stage-characters}

Let nu\_j be the already constructed zero-infinity finite-order character of E\_0 with the chosen G data, and let

\[
       \widetilde\nu_j=\nu_j\circ N_{K_h/E_0},\qquad
       \boxed{\ \mu_j=\widetilde\nu_j\alpha_{e_j}.\ }
 \tag{NC17}
\]

Use exactly the existing K031 or K032 G primes and local characters. Their completely split degree-one completions make each local norm the identity. Thus at every selected G place P,

\[
                 \mu_{j,P}|_{\mathcal O_{K_h,P}^\times}
                                      =\chi_{j,P}^{-1}.
 \tag{NC18}
\]

The conjugate place has the inverse character, as required by triviality on A\_F. Nothing in alpha\_e changes a compact finite restriction, because alpha\_e is unramified at every finite place. Consequently the exact tame conductor remains

\[
       \mathfrak f(\mu_j)=g_0g_1\mathcal O_{K_h},\qquad
       N\mathfrak f(\mu_j)=(g_0g_1)^{2f}
                       \quad(j=0,1),
 \tag{NC19}
\]

for the original nontrivial distinct G characters. In particular no ramification is added at29,11,T,2,7,5, or13. Unitarity, TrivOn F, HasInfQ(−e\_j), the exact atg clause and every finite unramified/tame clause in literal IsMu0 all follow. The same eta\_h and etaVal are retained in Chars.Tied. The numerical sum e\_even is automatic since each e\_(j,tau) is a multiple of28.

For the constant stage difference above, (NC16) gives

\[
          \mu_1\mu_0^{-1}
          =(\widetilde\nu_1\widetilde\nu_0^{-1})\eta_h^{-4\ell}.
 \tag{NC20}
\]

Thus the exact common relative character used in the two-stage analytic comparison is preserved. A common finite tuple twist, if subsequently applied as in the original construction, cancels from this ratio and retains its own separately prescribed Q data.

\subsection{6. One coefficient field for every profile, with an explicit bound}\label{one-coefficient-field-for-every-profile-with-an-explicit-bound}

Let M\_base be a finite Galois field containing K\_h, the existing normalized etaVal field, and the finite G roots. It is fixed before the profiles e\_j are chosen. Write L\_h=lcm(28,g\_0−1,g\_1−1) for the latter root exponent.

For each relative class relation (NC10) and every selected tau, choose a root, using the root already chosen for the corresponding basis character when that character is among the first f−1:

\[
          z_{s,\tau}^{d_s}=\tau(a_s/\bar a_s)^{28}.
 \tag{NC21}
\]

All finite values of the first f−1 basis characters lie in K\_h(z\_(s,tau):s,tau), by (NC11). The last basis character in (NC14) also takes values in the compositum with M\_base: on an unramified prime P, eta\_h(P)\^{}−56 equals

\[
       \bigl(\eta_h(P)(NP)^{-1/2}\bigr)^{-56}(NP)^{-28},
 \tag{NC22}
\]

which is in the old etaVal field. Multiplication extends this to every ideal. Hence

\[
 M=M_{\rm base}\bigl(\mu_{d_{\rm rel}},
                             z_{s,\tau}:s,\tau\in\Psi\bigr)
 \tag{NC23}
\]

contains every finite value of every alpha\_e and mu\_j, for \textbf{all} e in (28 Z)\^{}Psi. It therefore supplies the literal muVal field, not merely a different algebraic extension for each prime or weight.

This M is Galois over Q. Indeed K\_h and M\_base are Galois, and complex conjugation is central on K\_h. A Q-automorphism sends a radicand in (NC21) to the corresponding radicand for another selected embedding or its inverse. All its d\_s-th roots are already present because mu\_(d\_rel) is present. Every conjugate of each generator therefore lies in M. Character coherence is not changed by adjoining the unused last-label roots.

The field degree satisfies

\[
       [M:M_{\rm base}]
       \le\varphi(d_{\rm rel})\prod_s d_s^f
       =\varphi(d_{\rm rel})H_{\rm rel}^{f}
       \le h(K_h)^{f+1}.
 \tag{NC24}
\]

This gives a profile-independent bound without computing the large class group. A wholly numerical bound follows from the known discriminants and the standard Minkowski theorem. The exact absolute discriminants are

\[
 |\Delta_{K_1}|=2^{84}7^{70}29^{72},\qquad
 |\Delta_{K_4}|=2^{336}7^{280}29^{288}5^{168}13^{168}.
 \tag{NC25}
\]

The first follows from \textbar Delta\_(E0)\textbar=2\^{}12·7\^{}10, \textbar Delta\_B\textbar=29\^{}6 and the coprime-discriminant compositum. The second uses \textbar Delta\_(Q(sqrt5,sqrt13))\textbar=5\^{}2·13\^{}2 and the same formula. These are the integral composita already proved in K030/K032.

Minkowski's ideal-class representative bound is (4/pi)\^{}f (2f)!/(2f)\^{}(2f) sqrt(\textbar Delta\_K\textbar); see Milne (2020), Theorem4.3, printed p.70. In particular every class has an integral ideal of norm at most

\[
 B_1=2^{84}7^{35}29^{36},\qquad
 B_4=2^{336}7^{140}29^{144}65^{84},
 \tag{NC26}
\]

because these are exactly 2\^{}f sqrt(\textbar Delta\_(K\_h)\textbar). The number of integral ideals of norm m in a degree-n field is at most d\_n(m), the number of ordered n-tuples of positive integers with product m. To prove the bound, at a rational prime p the ideal-counting Euler factor is product\_(P\textbar p)(1−T\^{}(f\_P))\^{}−1. Coefficientwise it is bounded by (1−T)\^{}−n, since there are at most n primes and all f\_P≥1. Multiplication over p proves the assertion. Also d\_n(m)≤m\^{}(n−1), by choosing its first n−1 entries. Therefore

\[
                  h(K_h)\le H_h:=B_h^{2f},\qquad
                  [M:M_{\rm base}]\le H_h^{f+1}.
 \tag{NC27}
\]

These crude bounds are finite, explicit, and independent of the new weight profile. They are not proposed as economical arithmetic costs.

One can also bound the previously existing M\_base rather than treating its degree as a fixed input. In K029(ER25--ER26), the etaVal construction adjoins, for each cyclic generator of Cl(K\_1), one root of degree its class order, and the correcting class character has order dividing h(K\_1). For each root relation, adjoin one root of every conjugate radicand under the84 automorphisms of K\_1, and also mu\_(h(K\_1)); this contains every conjugate of each chosen root. The resulting normal field has degree at most h(K\_1)\^{}85 over K\_1. The normalized etaVal field for K\_4 is contained in the same field by its norm-pullback formula. Hence the deliberately loose completely numerical bound

\[
             [M:\mathbb Q]
                  \le 2f\,L_h\,H_1^{85}\,H_h^{f+1}
 \tag{NC28}
\]

is available. In this formula H\_1=B\_1\^{}84 refers to the original degree84 field in (NC26), even when h=4. No class number, regulator, or discriminant of the coefficient field has been computed here.

The bounded \textbf{field degree} does not bound all finite values by a constant height. Profile-dependent values are powers of finitely many fixed basis characters. At each fixed finite place P their logarithmic heights are bounded above by the sum of \textbar e\_tau\textbar/28 times the heights of those basis values, plus the fixed finite G factor. An effective global height/error statement would still need those arithmetic constants. No such cost is suppressed in (NC24)--(NC28).

\subsection{7. What this repairs in the new profile}\label{what-this-repairs-in-the-new-profile}

Arbitrary top/thick/thin assignments of e\_tau in28Z are covered, including negative assignments, independently at all42 or168 CM labels. If the second stage is obtained by the same increment2ell at every label, with2ell in28Z, the exact stage-ratio formula(NC20) is available. The G primes, their generic tame characters and residue fields, the odd etaT, the maximal29 level, and unramified 11 components are unchanged.

Thus nonuniform stage-character existence and a common algebraic value field impose no additional obstruction on that proposed profile. The field/type constructor's old common-quadratic failure is not altered. The profile still needs its own numerical proof and the full norm/height/terminal-image analysis; this theorem does not certify those independent requirements.

\subsection{References and check scope}\label{references-and-check-scope}

Milne, J. S. (2020). \emph{Algebraic number theory} (Version3.08). Theorem4.3, printed p.70; Theorem4.4, printed pp.70--71; Lemma6.5 and Remark6.6(c), printed pp.100--101. \url{https://www.jmilne.org/math/CourseNotes/ANT.pdf}. Theorem4.3 was directly read from the primary author PDF in this run. The two latter compositum statements are retained from the checked K030/K032 application. The finite relative-ideal extension, base restriction, coherent eta relation and all field-degree deductions are proved above.

Source conventions: \nolinkurl{AbelianHostFieldConditions.lean},314--331; \nolinkurl{IsUnitaryNormOf.lean},90--94. Carried inputs: K029(ER2,ER5,ER25--26), K031(HC16--21), K032(QH27--31), and LIFT18's zero-infinity finite character. No new class-number-one assertion or analytic L-function theorem is used. The companion exact certificate checks root-order exclusions, finite quotient carry identities, coherent type exponents, conductor and discriminant arithmetic, the exact fixed G root exponents, and the stated field bounds as formulas. It does not compute the large class groups, their ideals, Hecke values or coefficient heights, and it does not run Lean.

\clearpage

\section{AM-DO-040: two certified replacement T-primes for the large weight profile}\label{am-do-040-two-certified-replacement-t-primes-for-the-large-weight-profile}

Author: GPT-6 Astra Pro, ALP\_MACHINE \nolinkurl{/root/definite_oscillator}. 22 November 2026 UTC. Inherited reasoning effort is not separately exposed.

This finding replaces the two auxiliary T-primes of R017 by explicit larger primes, retaining the same CM hosts and G-primes. They satisfy T\textgreater16 A for A=2\^{}960, the exact norm and full source congruence, and an explicit projective residue-vector condition. The latter uses powers of the integer 2, with a certified order bound; no primitive generator of the whole residue-field multiplicative group or complete factorization of T−1 is assumed.

A standalone replay proves primality using only integer arithmetic, trial division of the specified factor primes, modular powers, and gcds. The probable-prime discovery was a separate calculation and is not an input to that proof. The result does not claim that these T choices alone fill a complete SizeData₂ record or repair the original quadratic host-field wrapper.

\subsection{1. Exact parameters}\label{exact-parameters}

Retain the checked hosts, conductors and G-primes of K031/K032:

\begin{longtable}[]{@{}lrrrrrr@{}}
\toprule\noalign{}
Host & f=a & c & g0 & g1 & B & m \\
\midrule\noalign{}
\endhead
\bottomrule\noalign{}
\endlastfoot
J\^{}14 & 42 & 812 & 4615409 & 5276377 & 392 & 263 \\
J\^{}56 & 168 & 52780 & 22285932761 & 64072175441 & 334 & 54 \\
\end{longtable}

Let delta=sqrt(−7) in their common external quadratic subfield and set

\[
 D=c(g_0g_1)^2\,2^B m,\qquad
 \pi_T=1+2D+D\delta,\qquad
 T=11D^2+4D+1,\qquad F=4D.                    \tag{DO40.1}
\]

The quantities in this formula are integers except for the explicitly specified quadratic generator in piT. Thus it completely specifies the prime and generator without a search oracle. The exact decimal primes are recorded below and in the replay JSON.

\textbf{J14} (1001 bits, 302 decimal digits):

\begin{verbatim}
17952167255185635051321154351674895613987764293769673456729204696357336089011562929338665498585310931065658058642583455560352833168574896704942121915304013574614983189345694485846080474422936718464345734525547995223578746517703781220872222478207325026191727208735314247959847795300609580930248088748033
\end{verbatim}

\textbf{J56} (996 bits, 300 decimal digits):

\begin{verbatim}
454943247013778122512187780897314598574665539127335259605458539964582491987337758099655723658178473308877448900224099993773653967422266630524027386542260738612598969923877926849959335045299704408598932706573786754076850608752614606900351396549840892191161193206683853263893019307714081695789188382721
\end{verbatim}

The replay verifies D\textgreater=2\^{}490 and 4 divides D in each case. In fact

\[
 \left\lfloor\frac{T}{4\cdot2^{960}}\right\rfloor
 =\begin{cases}
 460\,533\,319\,762&J^{14},\\
 11\,670\,820\,624&J^{56}.
 \end{cases}                                           \tag{DO40.2}
\]

Hence both T\textgreater16 A, more than the requested T\textgreater4 A. This inequality is an exact integer comparison. No floating-point exponential estimate enters the certificate.

\subsection{2. Deterministic primality proof}\label{deterministic-primality-proof}

The exact known factorizations are

\[
\begin{aligned}
 F_{14}&=2^{396}\,7\,29\,263\,4615409^2\,5276377^2,\\
 F_{56}&=2^{339}\,3^3\,5\,7\,13\,29\,
                      22285932761^2\,64072175441^2.
\end{aligned}                                          \tag{DO40.3}
\]

The replay checks each equality with F=4D. It proves every listed factor prime by testing all odd divisors through its integer square root (and handles 2 directly). The largest trial-division bound is only the integer square root of 64072175441. Primality of the large T is proved by the following complete factored-divisor argument.

First,

\[
 T-1=4D\left(1+\frac{11D}{4}\right),\qquad
 F^2-T=(5D+1)(D-1)>0.                                  \tag{DO40.4}
\]

Thus F divides T−1 and F\textgreater sqrt(T). For every distinct prime q dividing F, use the specified witness z in the following table:

\begin{longtable}[]{@{}lll@{}}
\toprule\noalign{}
Host & q & z \\
\midrule\noalign{}
\endhead
\bottomrule\noalign{}
\endlastfoot
J14 & 2 & 5 \\
J14 & 7,29,263,4615409,5276377 & 2 \\
J56 & 2 & 11 \\
J56 & 3 & 3 \\
J56 & 5,7,13,29,22285932761,64072175441 & 2 \\
\end{longtable}

Every one of the exact congruences and gcd identities

\[
 z^{T-1}\equiv1\pmod T,\qquad
       \gcd\left(z^{(T-1)/q}-1,T\right)=1               \tag{DO40.5}
\]

is replayed. The complete quotient-power residues, not only PASS flags, are saved in the JSON.

Here is why these checks prove primality. If p is any prime divisor of T, the multiplicative order of the q-specific z modulo p divides T−1. The second condition in (DO40.5) says it does not divide (T−1)/q. Its q-adic order therefore contains the full q-primary part of T−1, in particular the q-primary part of F. Since that multiplicative order divides p−1, this holds for each q, and F divides p−1. Consequently every prime divisor p of T satisfies p\textgreater=F+1\textgreater sqrt(T). A composite positive integer has a prime divisor at most its square root. Hence T is prime. This also excludes prime powers of exponent at least two.

This is the elementary Pocklington factored-divisor criterion used in R017, with its proof included. No unverified property of the complementary factor 1+11D/4 is needed; it is not asserted prime or fully factored.

\subsection{3. Exact source norm, congruence and residue conditions}\label{exact-source-norm-congruence-and-residue-conditions}

Direct multiplication in Q(delta) gives

\[
 \pi_T\overline\pi_T=(1+2D)^2+7D^2=T,\qquad
 \pi_T\equiv1\pmod{c(g_0g_1)^2}.                       \tag{DO40.6}
\]

Both generators are algebraic integers. The replay also checks

\[
 \gcd(T,2(1+2D))=1,\quad
 \gcd(T,11cg_0g_1)=1,\quad T\equiv1\pmod c.            \tag{DO40.7}
\]

A prime ideal dividing both piT and its conjugate would divide their product T and sum 2(1+2D), contradicting the first equality. Their ideals are therefore coprime, also after extending to either host.

Each checked host is contained in the c-th cyclotomic field. The prime T is coprime to c and congruent to one modulo c, so it splits completely in that cyclotomic field and therefore in the host. At all host primes over T its valuation is one. Equation (DO40.6) and coprimality show that precisely one prime of each complex-conjugate pair divides piT. There are exactly f such primes, each with residue field F\_T and piT-valuation one. These are the norm, congruence, coprimality, residue-norm and half-cardinality conditions used in R017. The G-primes and their previously checked congruences are unchanged.

The literal old \nolinkurl{piT_Kq} clause is not claimed: as K031 recorded, its original source-defined common field is Q on this induced mixed host, whereas piT belongs to the actual external Q(sqrt(−7)). The replacement geometric family retains the latter explicit quadratic subfield. Constructing the remaining generalized source data remains separate from the elementary prime certificate.

\subsection{4. A completely specified residue covector}\label{a-completely-specified-residue-covector}

Both hosts have roots of unity mu28, by the carried R017 host check. Reduction at T is injective on this group. Since T is congruent to one modulo28, its image is the unique subgroup of order28 in F\_T\^{}times. The replay checks, for every integer 1\textless=j\textless=a, that

\[
                         2^{28j}\not\equiv1\pmod T.     \tag{DO40.8}
\]

It follows that the image of 2 in the cyclic quotient F\_T\^{}times/mu28 has order h\textgreater a. In particular the residue covector

\[
                         u=(1,2,2^2,\ldots,2^{a-1})    \tag{DO40.9}
\]

has nonzero coordinates in distinct mu28 orbits. This condition has been proved without factoring all of T−1 or claiming that 2 is a primitive residue-field generator.

The projective stabilizer in the integral monomial group is exactly central mu28. For completeness, suppose a monomial matrix with entries in mu28 maps u to c u. Passing to F\_T\^{}times/mu28 says that translation by the class of c preserves the set I=\{0,1,\ldots,a−1\} in the cyclic group generated by the class of 2. That class of c lies in this cyclic subgroup because one member of I is zero. Around its cyclic ordering, I has a−1 successive gaps of length one and one gap of length h−a+1\textgreater1. The unique larger gap is preserved by a translation stabilizing I, so that translation is zero. Each coordinate is therefore fixed, and every monomial entry equals the same c in mu28. Conversely these scalar matrices clearly preserve the projective line. If u itself is fixed rather than its line, c=1. Thus

\[
 \Gamma\cap P_u=\mu_{28}I_a,\qquad
 \Gamma\cap H_u=1,\quad
 \Gamma=U(I_a)(F)\cap\operatorname{GL}_a(\mathcal O_K).
                                                        \tag{DO40.10}
\]

Here the exact monomial description of Gamma is R017(R17.7). A single T-place already forces these conclusions; imposing all f T decorations does not enlarge either group. This is the same class support needed for R017's matrix-coefficient input, now with the explicit covector (DO40.9). No principal29 refinement is introduced.

The restriction map from residue characters of F\_T\^{}times to characters of mu28 is surjective. Every fibre has (T−1)/28 members, which is larger than one. Even the trivial root-character therefore admits a nontrivial residue-character extension. Choose such nontrivial T characters at f−1 places; at the last choose a nontrivial extension of the root-character needed to cancel the actual archimedean and G central character. This gives precisely R017's central compatibility for any fixed selected archimedean coefficient representation. If these particular character-valued inputs are used as algebraic rows, their finite character values must be included in the coefficient field; adjoining mu\_(T−1) suffices. No cost-free field enlargement is claimed.

\subsection{5. Exact inherited class mass and its T-part at11}\label{exact-inherited-class-mass-and-its-t-part-at11}

Write

\[
 d_T=\frac{T^a-1}{T-1},\qquad
 d_G=\prod_{P\in G}|U_a(\mathbf F_{g_P})|,
 \qquad d_\rho=\dim W_0.
                                                        \tag{DO40.11}
\]

For the actual R017 projective-T/full-G matrix-coefficient support, with its chosen unit real and finite Steinberg vectors, the squared norm is still exactly

\[
                        \frac1{28\,d_T^f d_\rho d_G}.
                                                        \tag{DO40.12}
\]

This is an application of R017's checked mass and Schur-normalization formula to the newly verified stabilizer and level; it does not declare the raw kernel constant on this support. The literal strict refinement remains K038(CA30a--c): with

\[
 I_T=(T-1)^f,\qquad I_G=\prod_{P\in G}
                    [\operatorname{GL}_a(\mathbf F_{g_P}):U_a],
                                                        \tag{DO40.13}
\]

there are I\_T I\_G/28 strict classes of stabilizer one, each with volume v\_P/(I\_T I\_G). Both actual weighted kernel sums are retained. The dimensions and G factors have not been dropped from the norm.

The new T contributions have particularly simple exact11-valuations:

\begin{longtable}[]{@{}lrrrrr@{}}
\toprule\noalign{}
Host & T mod11 & order of T mod11 & v11(T−1) & v11(d\_T) & v11(d\_T\^{}f) \\
\midrule\noalign{}
\endhead
\bottomrule\noalign{}
\endlastfoot
J14 & 9 & 5 & 0 & 0 & 0 \\
J56 & 10 & 2 & 0 & 2 & 336 \\
\end{longtable}

For J14, 5 does not divide a=42, so T\^{}a−1 is nonzero modulo11. For J56, the replay gives v11(T+1)=2; since a=168, the odd-prime lifting formula gives v11(T\^{}a−1)=2+v11(84)=2. Both denominators T−1 are units at11. Thus the T-character coarsening factor 1/I\_T is an11-unit on both hosts, whereas the J56 projective-line mass retains its valuation −336. The unchanged G refinement and variable real dimension factors must still be accounted separately.

\subsection{6. Sources, checks and scope}\label{sources-checks-and-scope}

The full deterministic replay is \nolinkurl{definite_oscillator_large_t_certificate.py}; its JSON contains the exact primes, generators, factorization, every witness residue and gcd, all210 projective-order tests, and the exact11 mass valuations. The separate discovery script is explicitly labelled as such. The proof above explains why the replay establishes primality and the geometric level conditions; numerical checks are not substituted for those implications.

Carried findings: K031/K032's fields, conductors and G-primes; R017(R17.7--14) for the monomial group and matrix-coefficient mass; K038(CA30a--c) for exact strict refinement and weighted averages. The literal source conditions are \nolinkurl{Columns/AbelianHostFieldConditions.lean}206--238 (\nolinkurl{Primes.Ok}) and 486--494 (\nolinkurl{Sizes}), with the common-field clause separated as above.

The prime criterion is the factored-divisor argument of Pocklington, H. C. (1914--1916), ``The determination of the prime or composite nature of large numbers by Fermat's theorem,'' \emph{Proceedings of the Cambridge Philosophical Society, 18}, 29--30, already recorded in R017. Its complete required implication is proved in section2, so no additional inaccessible literature assertion is used here.

These data supply the requested large auxiliary prime and actual class support for the new weight-profile construction. They do not establish all its weight windows, a terminal nonzero period, or the complete arithmetic gain/cost inequality. No Lean execution is claimed.

\clearpage

\section{ALP-LIFT-29: the rank-zero sharp bottom level and an exact T translation detector}\label{alp-lift-29-the-rank-zero-sharp-bottom-level-and-an-exact-t-translation-detector}

Author: /root/lift\_accounting, GPT-6 Astra Pro, ALP\_MACHINE. 22 November 2026 UTC. Reasoning effort is not separately exposed.

This refines LIFT28's convenient bottom level. In the actual m=2, rS=0 ranks, the source bottom sharp residue group is exactly diag(U\_a,1), with its generic character on U\_a. The standard-frame dat0 level is maximal at T. An explicit lower unipotent at each T place makes its intersection with the compact subgroup H exactly the original strict covector level. The translated compact packet has exactly the original classes and stabilizers. LIFT28's holomorphic interpolation argument therefore applies at this source-shaped bottom level, without a new T level restriction.

This gives a concrete candidate for the finite bottom level and translate. The fixed tS, betaY, real model, and downward theta family of a completed source RealPack have not thereby been identified. No assertion about their terminal endpoint is made.

\subsection{1. Exact specialization of the source sharp subgroup}\label{exact-specialization-of-the-source-sharp-subgroup}

Use the recovered source commit 7cbf7a0c32fca1badf05931879a37c3f8a2b03b4. Let m=2, a=k, r=k+m-1=a+1, and rS=0. All finite residue statements below concern the actual selected relative split G places.

\nolinkurl{Divisibility/DefGroup.lean186–189} defines unipCol(r,A,C) as the upper unipotent subgroup with all entries (i,C) equal to zero for A\textless=i\textless C. Lines258--261 set

\[
 \operatorname{unipSym}(r,\rho,k)
       =\operatorname{unipCol}(r,\rho,\rho+k).
                                                        \tag{BT1}
\]

Thus, at r=a+1, rho=0, k=a, all entries above the last diagonal entry in the last column are zero. Upper triangularity already forces the last row below the diagonal to vanish. Consequently

\[
 \operatorname{unipSym}(a+1,0,a)
                =\{\operatorname{diag}(u,1):u\in U_a\}.
                                                        \tag{BT2}
\]

The symEdge definition at lines218--220 removes the adjacent edge ending at column a; its extra bridge edge would require i+1=0, which is impossible. All and only the adjacent edges within the first a coordinates remain. Its additive character is therefore exactly the usual generic character on U\_a. PsharpSym is defined from these local groups and characters at lines608--609.

The permutation in LinksChar.lean109--110 satisfies IsSymOrd (defined at lines75--76). At rS=0 every one of its cases gives pi(i)=i. There is no residual reordering in this specialization.

Further, LinksDefs.lean430--431 and \nolinkurl{SubgroupCharExtensionOneApply.lean32–43} show that dat0 has n=r=a+1, j=s=0, Q empty and Frm.B. At G the block function blkOf(r,t) is constant, so parLev adds no condition beyond integrality. The character product over inserted blocks is empty, as is the Q product; see DefGroup.lean635--637 and704--707. With the standard lattice frame beta=1, the exact finite data are

\[
 J_{B,v}=
 \begin{cases}
   \operatorname{red}^{-1}\operatorname{diag}(U_a(\kappa_v),1),
                                                 &v\in G,\\
   K_{a+1,v},                                   &v\notin G,
 \end{cases}
\quad
 \chi_{B,v}=
 \begin{cases}\psi_{U_a},&v\in G,\\1,&v\notin G.\end{cases}
                                                        \tag{BT3}
\]

Here \(K_{a+1,v}\) is the unitary stabilizer of the standard selfdual lattice; at a relative split place it is the ordinary general linear maximal compact. The character at G is the inverse of the actual initial right character. This statement specializes the source level formula; it does not assert that every other RealPack tie has been constructed.

\subsection{2. The exact T conjugation}\label{the-exact-t-conjugation}

Work at one actual T place. It is relative split and unramified of residue characteristic different from11 in the current constructions. Let \(O\), \(\varpi\), and \(\kappa\) be its local integer ring, uniformizer, and residue field. Let \(u\in O^{1\times a}\) be a lift of the actual nonzero covector defining the initial strict T level. Set

\[
 \tau_u=
 \begin{pmatrix}I_a&0\\ \varpi^{-1}u&1\end{pmatrix}
           \in\operatorname{GL}_{a+1}(F_v)
                     \simeq G^+(F_v).             \tag{BT4}
\]

The relative split identification is the one compatible with H embedded as diag(h,1). It sends every displayed matrix to an actual unitary point by its paired component \(J^{-1}g^{-t}J\), where \(J=\operatorname{diag}(I_a,-1)\); no rational global unitary matrix is being postulated.

Direct multiplication gives, for h in GL\_a(F\_v),

\[
 \tau_u^{-1}\operatorname{diag}(h,1)\tau_u
   =\begin{pmatrix}
      h&0\\ \varpi^{-1}u(1-h)&1
     \end{pmatrix}.                              \tag{BT5}
\]

This matrix and its inverse are integral if and only if \(h\in\operatorname{GL}_a(O)\) and \(\overline u\,\overline h=\overline u\). Hence

\[
 H(F_v)\cap\tau_u K_{a+1,v}\tau_u^{-1}
   =J_{0,v}
   =\{h\in\operatorname{GL}_a(O):
                         \overline u\,\overline h=\overline u\}.
                                                        \tag{BT6}
\]

This is the strict covector stabilizer, not the projective stabilizer and not an additional principal congruence group. There is no dilation average in (BT6). Both relevant characters are trivial on this group.

For later use, the identical calculation for an extra scalar t is

\[
 \tau_u^{-1}\operatorname{diag}(A,t)\tau_u
       =\begin{pmatrix}
            A&0\\
            \varpi^{-1}(t u-uA)&t
         \end{pmatrix}.                          \tag{BT7}
\]

Membership in the maximal compact implies both \(A\in\operatorname{GL}_a(O)\) and \(t\in O^\times\); this follows from the upper-left and lower-right blocks of the matrix and its inverse. Once t=1 it is exactly the original condition (BT6).

Take the finite adele \(\tau\) with components (BT4) at all T places, identity elsewhere, and identity at infinity. It is a genuine finite unitary adele of determinant one. From (BT2--6),

\[
 H(\mathbf A_f)\cap\tau J_B\tau^{-1}=J_0,\qquad
 \chi_B(\tau^{-1}h\tau)=\chi_0(h)^{-1}
                                         \quad(h\in J_0).
                                                        \tag{BT8}
\]

In particular, if F has bottom right level (J\_B,chi\_B), then h maps to F(diag(h,1)tau) has precisely the required inverse initial finite character for the bilinear period.

\subsection{3. The translated packet and all rational stabilizers}\label{the-translated-packet-and-all-rational-stabilizers}

Map

\[
 Z=H(F)\backslash H(\mathbf A_f)/J_0
   \longrightarrow
 G^+(F)\backslash(D\times G^+(\mathbf A_f))/J_B,
 \qquad [x]\longmapsto[(o,x\tau)].                 \tag{BT9}
\]

It is well defined by (BT8). Suppose two points agree. The relating rational element fixes the negative line at every real place and is diag(h,t), with \(h\in H(F)\), \(t\in E^1(F)\). At a place away T the finite level condition makes t an integral unit; at T this follows from (BT7). Thus t is an integral norm-one E-unit, hence a root of unity. At a selected G place, (BT2) forces its last diagonal residue to be one. The residue characteristic is prime to the root order, so t=1. Now (BT8) says exactly that the original two H classes agree. Thus (BT9) is injective.

Applying the same argument to a single packet point identifies its entire orbifold stabilizer, including its finite character, with

\[
                       \Gamma_x=H(F)\cap xJ_0x^{-1}.
                                                        \tag{BT10}
\]

This retains the actual strict-level stabilizers. It introduces no extra factor from the complementary U(1).

For an arbitrary arithmetic component \(\Gamma_g=G^+(F)\cap gJ_Bg^{-1}\), its determinant is a global norm-one unit. The G determinant in (BT2) is one modulo the selected prime, so the same root argument gives

\[
                              \det\Gamma_g=1.     \tag{BT11}
\]

Therefore the determinant-flat correction in LIFT28(HJ16) is still trivial on every component. The proof of jet separation does not change when the packet points move from x to x tau.

\subsection{4. The actual translated bottom detector}\label{the-actual-translated-bottom-detector}

Keep the actual selected initial HC u, bottom positive list \((v,-b_0)\), and the selected degrees

\[
 N_\tau=\sum_i u_{\tau,i}-\sum_i v_{\tau,i}
                              +b_{0,\tau}-a/2.    \tag{BT12}
\]

Vary b1 as in LIFT28(HJ15). For all sufficiently large ell in its fixed progression, the cusp restriction map at the actual finite level (BT3), using the fixed selected K-type embedding, is onto:

\[
 \operatorname{Hom}_{G^+_\infty}
    \bigl(\sigma_\ell^\vee,
       \mathcal A_{\rm cusp}(G^+,J_B,\chi_B)\bigr)
 \longrightarrow
       \mathcal A_{\lambda^\vee}(J_0,\chi_0^{-1}),
       \qquad F\longmapsto[h\mapsto F(\operatorname{diag}(h,1)\tau)].
                                                        \tag{BT13}
\]

Proof: repeat LIFT28 Sections5--7 with the finite level (BT3) and packet (BT9). Equation(BT10) supplies exactly the same compatible jet target, and(BT11) supplies the same ample canonical twist. The PBW-to-isometric factors are unchanged: tau is finite and the differential operators are real. No new finite Schwartz input is used in this argument.

For every nonzero initial row P, including every nonzero Mram row, and every unit seed w, there is consequently an actual bottom cusp form F such that

\[
 F(\operatorname{diag}(h,1)\tau)=\overline{f_{P,w}(h)}
                                             \quad(h\in[H]),
\] \[
 \int_{[H]} f_{P,w}(h)F(\operatorname{diag}(h,1)\tau)\,dh
       ={v_0\over d_\lambda}
           \sum_x{\|P_x\|^2\over|\Gamma_x|}>0.     \tag{BT14}
\]

All Haar and stabilizer factors are those of the original strict initial level. Right translation by tau is unitary on the ambient automorphic L2 space, so it does not itself change the Petersson norm of F. This is not an assertion that tau preserves an integral coefficient lattice or has zero arithmetic height cost. The norm of the interpolating cusp form remains uncontrolled, as does the sufficient ell.

\subsection{5. What has and has not been matched}\label{what-has-and-has-not-been-matched}

The new conclusions are (i) the exact rS=0 source residue specialization (BT2--3), and (ii) the explicit maximal-T translation and packet comparison (BT4--11). They refine the bottom detector to the standard-frame dat0 level formula, while preserving the literal initial strict level.

The source RealPins.tSpin condition is stronger. It identifies the entire conjugated bottom level and character with the fixed rowJ/rowChi in a particular real/arithmetic model; see \nolinkurl{BottomHypothesisMatchingCentral.lean22–44} and PairForm.lean66--73. It also uses the prescribed tS, not a translate freely redefined after the kernel construction. The equality of this tau with that fixed tS, up to the allowed embedded/rational/central/level corrections, has not yet been proved. Nor has betaY0 been tied to the standard frame while preserving the actual two vertical raw kernels and their terminal lattice family.

Thus (BT13--14) does not assert ValLastX, BotDat.Pin, or that the interpolated cusp form belongs to the prescribed downward theta image. Those remain the endpoint comparison. The conclusion is nevertheless stronger than choosing an unrelated bottom finite level: its exact residue subgroup, maximal T condition, character, and initial-level intersection are now fixed explicitly.

\subsection{6. Verification}\label{verification-1}

The source specialization was read directly in DefGroup.lean 138--275, 601--609, 620--707, LinksDefs.lean430--431, and \nolinkurl{SubgroupCharExtensionOneApply.lean32–43;} the reordering was checked in LinksChar.lean75--76 and109--110. These files were recovered from the pinned source; retrieval manifests are source\_recovery\_20261122T022749Z.json and source\_recovery\_20261122T022927Z.json. Recovery is not a Lean check.

LIFT28's whole checked interpolation proof is a carried input. The displayed block identities, compact intersection in both directions, root-unit packet argument, character inverse and Haar scope are the new proof here. Author whole-proof review completed 2026-11-22T02:34:29+00:00: PASS after reading every section and BT1--14, with direct source specialization and both directions of the block intersection checked. No effective image-cost or terminal-nonzero conclusion is claimed.

\clearpage

\section{AM-R-023: literal Host rank and Scale obstructions for the repaired CM fields}\label{am-r-023-literal-host-rank-and-scale-obstructions-for-the-repaired-cm-fields}

Author: GPT-6 Astra Pro, ALP\_MACHINE \nolinkurl{/root}, 22 November 2026 UTC. Reasoning effort is not separately exposed.

This finding fixes the scope of the remaining global parameter adapter. The actual mixed theta families on J14 and J56 exist in the senses proved in K029--K040 and LIFT26--29. Neither underlying CM field can, however, be the field of a literal source \nolinkurl{HostRaw}, for any choice of its abstract Frame. Changing only the weights, Scale, or T prime cannot repair this. Even if one discards the already proved CM-type obstruction and enlarges the field, the unchanged source Frame/Host inequalities impose the exact rank and Scale lower bounds below. These are necessary bounds, not a construction of the enlarged host or a sign computation for DP44.

\subsection{1. Source data and exact scope}\label{source-data-and-exact-scope}

Write n for Frame.ord, A for its additive group of order n, e for its cyclotomic parameter, a0 for its distinguished generator, and

\[
 n_0=\operatorname{ord}(a_0),\qquad o=n/n_0.
                                                        \tag{R23.1}
\]

The literal Frame contains a surjective homomorphism

\[
 h:A\longrightarrow H_e,
 \qquad H_e=(\mathbb Z/2^{e+2}\mathbb Z)^\times/\{\pm1\},
                                                        \tag{R23.2}
\]

and requires that h be injective on the cyclic subgroup generated by a0. The target is cyclic of order 2\^{}e. This last standard cyclotomic group fact is the source interface in \nolinkurl{FrobeniusOrbits/Statement.lean99;} it is used here, not presented as a new cyclotomic theorem.

A \nolinkurl{HostRaw} further requires e\textgreater=2 and an isomorphism

\[
 \operatorname{Gal}(E/\mathbb Q)\simeq C_2^2\times A.
                                                        \tag{R23.3}
\]

The full \nolinkurl{Host} imposes o\^{}64\textless=2\^{}e. The source rank and Scale conditions include

\[
 f=2n\le k,\quad k+m+4\le t,\qquad
 o\le x^{\epsilon/100},\quad x>1,\quad0<\epsilon<1/52,
                                                        \tag{R23.4}
\] \[
 L=x^\epsilon,\qquad \mathfrak a=x\exp L.
                                                        \tag{R23.5}
\]

The different symbol mathfrak a is used for the source Scale.a, to avoid confusing it with a unitary-group rank. The conclusion below uses only these displayed clauses. It does not use, weaken, or audit the source's accepted theorem.

\subsection{2. No literal HostRaw on J14 or J56}\label{no-literal-hostraw-on-j14-or-j56}

The actual fields already constructed are

\[
 K_{14}=\mathbb Q(\zeta_{28})B,\qquad
 K_{56}=K_{14}(\sqrt5,\sqrt{13}),
                                                        \tag{R23.6}
\]

where B is the cyclic degree-seven subfield of Q(zeta29). Their checked Galois groups are

\[
 \operatorname{Gal}(K_{14}/\mathbb Q)=C_2^2\times C_{21},
 \qquad
 \operatorname{Gal}(K_{56}/\mathbb Q)=C_2^4\times C_{21}.
                                                        \tag{R23.7}
\]

Both have exponent 42. Therefore neither field supports a literal HostRaw, under any relabeling or choice of Frame.

\textbf{Proof.} Equations (R23.2--3) would give a surjection from its Galois group to H\_e, hence to C4 because e\textgreater=2. Every quotient of a group of exponent 42 has exponent dividing 42. Since 4 does not divide 42, C4 cannot be such a quotient. This is a contradiction. The argument does not depend on the previously chosen e=0 Frames; it excludes every possible Frame on either fixed field. Their earlier numerical \nolinkurl{Ranks} records did not claim to instantiate HostRaw and remain valid. ∎

The same proof excludes every compositum of either field with finitely many quadratic extensions of Q: its Galois group embeds in the product of the original group with a group of exponent two and still has exponent dividing 42. This assertion concerns quadratic extensions of Q, not arbitrary iterated quadratic towers, which can have cyclic quartic quotients.

\subsection{3. Necessary rank of any enlargement retaining the literal inequalities}\label{necessary-rank-of-any-enlargement-retaining-the-literal-inequalities}

Let q0 be the odd part of n.~Every Frame satisfies

\[
 2^e\mid n,\qquad n_0\mid2^e,
 \qquad q_0\mid o.
                                                        \tag{R23.8}
\]

If it satisfies the full Host inequality, set

\[
 e_{\min}(q_0)=\min\{j\ge0:q_0^{64}\le2^j\}.
\] Then

\[
 e\ge e_{\min}(q_0),\qquad
 n\ge q_0\,2^{e_{\min}(q_0)},\qquad
 k\ge2q_0\,2^{e_{\min}(q_0)}.
                                                        \tag{R23.9}
\]

\textbf{Proof.} Surjectivity in (R23.2) gives 2\^{}e\textbar n by the first isomorphism theorem. Injectivity on the cyclic subgroup gives n0\textbar2\^{}e by Lagrange's theorem. In particular n0 is a power of two, so the integer o=n/n0 is divisible by the odd part q0 of n.~Thus 2\^{}e\textgreater=o\^{}64\textgreater=q0\^{}64. The integer 2\^{}e q0 divides n because its two factors are coprime and both divide n. These observations prove (R23.9), including k\textgreater=2n. ∎

In particular, if E contains K14 or K56 and has the literal degree {[}E:Q{]}=4n, the tower formula gives 21\textbar n.~One may apply the same argument directly with the odd divisor 21, even if the enlargement introduces additional odd factors. Exact integer arithmetic gives

\[
 2^{281}<21^{64}<2^{282},\qquad
 e\ge282,\quad n\ge21\,2^{282},\quad
 f,k\ge42\,2^{282},\quad t\ge42\,2^{282}+m+4.
                                                        \tag{R23.10}
\]

The numerical lower bound for the Frame plus o\_pow is sharp: take A=C21 x C\_(2\^{}282), let h be projection to the second factor, and choose a0=(0,1). Then n0=2\^{}282, o=21 and the displayed Frame conditions and o\_pow all hold. This verifies sharpness only for these group inequalities; it does not construct a number field, CM type, polarization, full Host, or admissible theta family.

For comparison, suppose a principally polarized host is isogenous to an ordinary power J\^{}N and has a full CM field of degree 2dim(J\^{}N)=6N. If its degree also equals 4n, then N is even, n=3N/2 is divisible by3, and

\[
 2^{101}<3^{64}<2^{102},\qquad
 n\ge3\,2^{102},\quad k\ge6\,2^{102},\quad N\ge2^{103}.
                                                        \tag{R23.11}
\]

This is a conditional numerical bound obtained before imposing the CM type. AM-R-003 already excludes the actual quadratic-induced Host type on every nonzero J power. The two statements must not be confused: (R23.11) is not a construction that avoids R003.

\subsection{4. The Scale cost cannot be repaired at fixed a=2\^{}960}\label{the-scale-cost-cannot-be-repaired-at-fixed-a2960}

For any o\textgreater=1 satisfying (R23.4--5),

\[
 L\ge o^{100},\qquad
 \log\mathfrak a>L+52\log L
                    \ge o^{100}+5200\log o.
                                                        \tag{R23.12}
\]

\textbf{Proof.} The o condition is o\textless=L\^{}(1/100), and taking its hundredth power gives the first inequality. Since x\textgreater1 and epsilon\textgreater0, L\textgreater1 and log L=epsilon log x. The strict upper bound epsilon\textless1/52 gives log x\textgreater52log L. Use log mathfrak a=log x+L. The function y+52log y is increasing for y\textgreater0, proving the final inequality. ∎

Consequently every enlarged host still containing K14 or K56 requires

\[
 L\ge21^{100},\qquad
 \mathfrak a>\exp(21^{100})\,21^{5200}.
                                                        \tag{R23.13}
\]

For any o\textgreater=2, there is no Scale with mathfrak a=2\^{}960 satisfying the literal o\_tie clause, regardless of how x, epsilon or J are changed. Indeed (R23.12) would give log mathfrak a\textgreater2\^{}100, whereas log(2\^{}960)=960log2\textless960\textless2\^{}100. Thus R022's specified weight profile and DO040's corresponding primes solve their stated weight/numerical Sizes problems; reparametrizing that same fixed Scale.a cannot complete the literal global record. A successful replacement must change the scale and coefficient profile, or prove an appropriately modified global budget statement for the mixed host.

\subsection{5. Consequences and geometric language}\label{consequences-and-geometric-language}

For DP\_MULT Q3e, this converts the fixed-host realization issue into an explicit negative statement: neither K14 nor K56 can be inserted in the literal HostRaw interface. The unchanged global Host inequalities also make the rank and Scale costs of a field-enlargement attempt explicit. They do not evaluate the actual g\_j, e\_j, H or endpoint terms in R001, and do not prove that every adapted mixed construction has negative DP44 difference. The mixed Picard/section and actual two-edge kernel constructions remain available independently of this literal interface.

Over Spec(Q) the quotient obstruction is a statement about finite etale Galois realizations and their constant finite quotient groups. The CM-type distinction belongs to the simple weight-one isogeny constituent h1(J), exactly as in R003; it is not a statement that every higher tensor or exterior-power correspondence vanishes. The previously fixed schemes, Picard functors and biextension lines retain the locators FGA232 Theorem3.1, pp148--149; FGA236 Corollary1.7, p229, and2.3, p232; SGA7I VII2.7/2.9, pp169--174, andVIII7.4(b), pp310--311. Representation lambda-rings/K-theory (SGA6 VI3.3) do not alter these scalar inequalities. The cote4/cote54 Cartier--Dieudonne realization is unchanged; no new ell-adic purity, motivic comparison, or Kummer RH assertion follows.

\subsection{6. Exact locators and verification}\label{exact-locators-and-verification}

\begin{itemize}
\tightlist
\item
  Alpöge, L. (n.d.). \emph{Remark}, step8 and the associated Lean archive, pinned commit7cbf7a0c32fca1badf05931879a37c3f8a2b03b4: \nolinkurl{Formalism/SharpRestriction/RamificationFrobenius/TotientGaloisClasses/ZsmulGeneratorEqZeroOfProjection.lean}, Frame16--24, n0/o68--70; \nolinkurl{Formalism/FrobeniusOrbits/Statement.lean}, HGrp27--28 and cyclotomic group interface97--101; \nolinkurl{.../TotientGaloisClasses/ConjEMemHostFIff.lean}, HostRaw12--43, cmType66, HostDer degree91--92; \nolinkurl{Formalism/LocalRingCardinality/DataHost.lean}, o\_pow50; \nolinkurl{Params.lean}, Ranks6--22 and f/r28--32; \nolinkurl{DataParams₂.lean}, Scale9--38 and o\_tie131.
\item
  ALP\_MACHINE AM-K029/030 and AM-K032, especially K032(QH22) and its preceding Galois decomposition: actual fields and ordinary J powers. AM-R-003: the independent isogeny/CM-type obstruction, used as a settled result without rederiving its proof.
\item
  AM-R-022 and AM-DO-040: the exact scope of the fixed-a replacement weights and prime certificates; none claimed a full Host.
\end{itemize}

The accompanying certificate uses only integer powers, divisibility, group exponents, and the displayed rank arithmetic. It does not certify a Host, a number-field extension, or a transcendental equality. The Scale proof above is symbolic and independent of floating-point logarithm approximations. Whole-proof checks are recorded separately in CHECK\_EVENTS.tsv.

\clearpage

\section{AM-K041 --- actual coupled divisibility and full-theta survival for the narrow cloud}\label{am-k041-actual-coupled-divisibility-and-full-theta-survival-for-the-narrow-cloud}

Author: GPT-6 Astra Pro, ALP\_MACHINE /root/kummer\_extension, 22 November 2026 UTC. Inherited reasoning effort is not separately exposed.

This finding specializes the \textbf{actual} class-kernel estimate K039 to the repaired profile R022, K040's actual nonuniform stage characters, and DO040's actual new T-prime. It covers every binary interior word and every top/thick/thin assignment. It also checks the changed hypotheses needed to transport LIFT25--27's full-theta injectivity. No old uniform numerical Schwartz norm or coefficient dimension is carried over.

The results are an arithmetic coupled-image lower bound and, separately, analytic injectivity before terminal restriction. They do not assert nonvanishing of the restricted row, PtilFF, or MainX. They do not fill Fields.Kq, Host, D7, SizeData₂, a downward-image condition, or a gain/error and height record. R022's stated o\_tie limitation remains. The old R020 numerical branch bounds are not used for the new partitions.

\subsection{1. Actual family and integral modules}\label{actual-family-and-integral-modules}

\begin{longtable}[]{@{}lrrrrrr@{}}
\toprule\noalign{}
Host & h & f=a=k & d & b=a+d & c=b+d & ell \\
\midrule\noalign{}
\endhead
\bottomrule\noalign{}
\endlastfoot
J\^{}14, K031 & 1 & 42 & 1760 & 1802 & 3562 & 45,156,049,316 \\
J\^{}56, K032 & 4 & 168 & 1634 & 1802 & 3436 & 40,530,339,696 \\
\end{longtable}

Here J is the Jacobian of y²=x⁷−1, and the CM type has f labels. Keep the product-polarized integral models, actual standard CM eigenbases, source creation operators, and period of K030/K032/K033/K035. Choose any partition of the labels into top, thick, thin, with counts \(N_T,N_H,N_L\) summing to f.~Set

\[
 A=2^{960},\quad B=2^{955},\quad n=f/2,\qquad
 E_T=28\left\lceil\frac{B+e^{336}}{28}\right\rceil,\quad
 E_H=28\left\lceil\frac{e^{343}}{28}\right\rceil,\quad
 E_L=28\left\lceil\frac{e^{336}/n}{28}\right\rceil.
 \tag{NW1}
\]

These are R022's exact integers. The replay independently encloses the exponentials by rational Taylor bounds and checks each ceiling. At a label of role beta use \(e_{0,\tau}=E_\beta,\ e_{1,\tau}=E_\beta+2\ell\). K040 constructs actual unitary characters with these infinity types, the prescribed real-adele restriction, the unchanged tame G-unit characters, and unramified components at29,11,T. Its exact relation is

\[
 \mu_1\mu_0^{-1}
    =(\widetilde\nu_1\widetilde\nu_0^{-1})\eta_T^{-4\ell}.
 \tag{NW2}
\]

Use DO040's new T and piT, the unchanged G-primes, and any fixed admitted common tuple character xi and its Q set. Both members of edge j have splitting character \(\chi_j=\mu_j\xi\). Keep K038's exact raw G/T/Q tensor, unnormalized Y sum, literal strict initial and middle levels, actual class stabilizers, and terminal tilQ decoration. The level at29 stays maximal. The changed numerical T and residue covector are exactly DO040's; no auxiliary Heisenberg translation or principal29 refinement is introduced.

For a word \(\omega\in\{0,1\}^{k-2}\), write \(\mathcal L_0^\omega,\mathcal L_{\rm ev}^\omega\) for K037/K038's actual integral Weyl coefficient lattices at these weights, with their stabilizer relations and transported class frames. The target is the coupled evGSD row with rhoVr action, not an arbitrary ordinary ColV row. The ramified source is its actual Mram sublattice. Allow a finite extension of the coefficient DVR R at a place over11; always normalize v(11)=1. This extension includes the actual root-of-unity fields of the new T input characters, with no claim that their field cost is trivial.

\subsection{2. Exact all-word selected degrees}\label{exact-all-word-selected-degrees}

Put \(W=\sum_i\omega_i\). The literal source weights and selected partitions of R022 are, in zero-based indices,

\[
 w_{0,i}=u_i-\frac{k-1}{2}+i,\qquad
 \nu_{0,i}=w_{0,i}+E_\beta-\frac{k+d}{2},\qquad
 \Delta=2\ell-d>2,\qquad
 \nu_1=(\nu_0,0^d)+\Delta\mathbf1_b.
 \tag{NW3}
\]

Let \(\delta_{j,\beta}(W)=|\nu_{j,\tau}|\) for role beta. Then

\[
\begin{split}
 \delta_{0,T}(W)
  &=A-B+kE_T+\ell\{2k(k-1)-4-W\}
       -\frac{k-2}{2}-\frac{k(k+d)}2,\\
 \delta_{0,\beta}(W)
  &=kE_\beta+\ell\{2k^2+6k-6-W\}
       +\frac{4-k}{2}-\frac{k(k+d)}2
                      \quad(\beta=H,L),\\
 \delta_{1,\beta}(W)&=\delta_{0,\beta}(W)+b\Delta.
\end{split}
 \tag{NW4}
\]

\textbf{Proof.} The rho shifts in w0 sum to zero. The R022 interior sum is \[
 \sum_{i=1}^{k-2}u_i
   =2\ell k(k-1)-4\ell-\frac{k-2}{2}-\ell W.       \tag{NW5}
\] The endpoints sum to A−B at a top label, and \(\ell(8k-2)+1\) otherwise. Add these sums and then \(kE_\beta-k(k+d)/2\). This proves the first two equations; the last follows from (NW3). All terms are integers since k,d are even. The arrangement of the bits can change the partition and its dimension, but its degree depends only on W. QED.

For \(d_{11}(m)=1+\lfloor\log_{11}m\rfloor\), the exact digit lengths on \textbf{both} edges, for every \(0\le W\le k-2\), are:

\begin{longtable}[]{@{}lrrr@{}}
\toprule\noalign{}
Host & top & thick & thin \\
\midrule\noalign{}
\endhead
\bottomrule\noalign{}
\endlastfoot
J14 & 278 & 145 & 141 \\
J56 & 279 & 146 & 141 \\
\end{longtable}

Indeed (NW4) decreases by ell at each increment of W. The exact endpoint checks are \[
 11^{z_\beta-1}\le\delta_{j,\beta}(k-2)
        \le\delta_{j,\beta}(0)<11^{z_\beta}.        \tag{NW6}
\] No floating point logarithm enters. Put \[
 Z=N_Tz_T+N_Hz_H+N_Lz_L,\qquad
 D_j(W)=\sum_\beta N_\beta\delta_{j,\beta}(W)
                         =D_j(0)-f\ell W.         \tag{NW7}
\]

\subsection{3. The actual local-section premise transports}\label{the-actual-local-section-premise-transports}

\subsubsection{Proposition NW8}\label{proposition-nw8}

Every actual normalized framed class germ in the two changed-weight K038 kernels satisfies K039(IG7)'s integral algebraic local-section premise in the same normalized Chern/CM gauge. Thus K039(IG9) applies with (NW7).

\textbf{Proof.} Only the infinity types, selected polynomials, and numerical T/residue covector have changed. The CM varieties, theta line, primitive sigma section, and source/logarithm comparison are unchanged. K039(IG2) excludes11 from every admitted Q set by the literal DegOne condition and residue degree6 in Q(zeta28)\^{}+. DO040's T differs from11 and splits completely; G is unchanged; K040's stages are unramified at11.

Apply K039(IG3)'s Cayley weak approximation to the new finite-level class sets. Representatives have component one at11 and their \textbf{own transported} standard integral Weyl frames. A rational change from an earlier display is not presumed integral. These choices depend on the fixed ranks and levels, not the archimedean word. The piT in tilQ is an11-unit.

K039(IG4)'s principal-ideal proof applies to the new algebraic angular characters: their finite splitting values on ideles trivial over11 are11-units. That proof never assumed uniform infinity type or finite-order global stages. The raw G/T/Q values and finite Weil transforms remain integral at11, since their Fourier/measure denominators involve only residue characteristics different from11. K039(IG5) expresses them as integral linear combinations of prime-to11 Heisenberg translates of the specified vacuum. There is no division by T−1, by a residue character group, or by the cardinality of Y.

K035's all-polynomial identity then gives the actual reduced torsion translates of the sigma product, after removing the fixed period. LIFT08's primitive normalization and K034's normalized torsion difference connection place them in the complete module \(U_0\exp(-\beta_{\rm prod})R[[t]]\). The exponential is retained: reduced Chern Taylor integrality is not claimed. K033's CM coordinate change and inverse, and K035's algebraic cD factors, are still integral and invertible at11.

Finally K037's integral Weyl-to-Weyl subD theorem applies to every actual dominant partition (NW3). Every resulting monomial has the selected degree at each label. Therefore the full-module bound K027(EX2--EX4), for arbitrary integral algebraic prefactors in the retained gauge, applies to the \textbf{whole} selected module. No rational isotypic or Gram projector is introduced. These are exactly the hypotheses of K039(IG9). QED.

\subsection{4. Explicit actual coupled-row divisibility}\label{explicit-actual-coupled-row-divisibility}

Define \[
 S_f(r)=2m+v_{11}(m!),\qquad m=\lfloor2fr/5\rfloor,
 \quad E_0=abZ,\quad E_1=bcZ,
\] \[
 L_j(W)=\left\lfloor D_j(W)/1220\right\rfloor-E_j,\qquad
 G(W;N_T,N_H,N_L)=L_0(W)+L_1(W)-S_f(a)-S_f(b).       \tag{NW9}
\] All L\_j here are positive, as proved below, so K039's maximum with zero is unnecessary. The stabilizer bounds are:

\begin{longtable}[]{@{}lrr@{}}
\toprule\noalign{}
Host & S\_f(a) & S\_f(b) \\
\midrule\noalign{}
\endhead
\bottomrule\noalign{}
\endlastfoot
J14 & 1,479 & 63,572 \\
J56 & 23,705 & 254,294 \\
\end{longtable}

\subsubsection{Theorem NW10}\label{theorem-nw10}

For every word and admitted fixed tuple, the actual K038 arithmetic map at the new profile satisfies \[
 \boxed{\ 
 P_{\rm cpl}^\omega(\mathcal L_0^\omega)
   \subset11^{G(W;N_T,N_H,N_L)}\mathcal L_{\rm ev}^\omega.
 \ }                                                     \tag{NW10}
\] The same holds on Mram. Uniformly in all words at fixed roles, use \(G(k-2;N_T,N_H,N_L)\).

\textbf{Proof.} NW8 and K039(IG8--IG10) give the two entrywise bounds, with digit errors abZ and bcZ. K039(IG11)'s reduction-mod3 bound uses the actual underlying integral rank 2fr and applies to every stabilizer of the new strict classes without a level refinement. Divide by the two actual stabilizers in the class sum. The Y sum is unnormalized. The nonarchimedean triangle inequality gives (NW10), exactly K039(IG13). K038 covariance applies because its finite-unit characters are unchanged and its integral subD identities are valid at the new weights. Intersection with the actual evGSD equivariant subspace is saturated; Mram is a sublattice, with no added projector. Finally Z is independent of W and D\_j decreases with W, proving the uniform-word version. QED.

A bound independent even of the role assignment is \[
\begin{split}
 G_{\min,h}={}&
 \left\lfloor f\delta_{0,L}(k-2)/1220\right\rfloor
 +\left\lfloor f\delta_{1,L}(k-2)/1220\right\rfloor\\
 &-141f(ab+bc)-S_f(a)-S_f(b).
\end{split}                                              \tag{NW11}
\] It is the minimum of the \textbf{lower-bound formula}, not a claim about the actual image's minimum content. Exact integers are in the replay; useful rigorous sizes are \[
                  G_{\min,1}\ge2^{481},\qquad
                  G_{\min,4}\ge2^{483}.                 \tag{NW12}
\]

\textbf{Proof of minimization and positivity.} For beta=H,T the difference \(R_\beta=\delta_{0,\beta}(W)-\delta_{0,L}(W)\) is independent of W and identical on edge1. The exact certificate verifies \[
 \lfloor R_\beta/1220\rfloor
         >\max(ab,bc)(z_\beta-141)>0.                    \tag{NW13}
\] Replacing a thin label by beta increases either floor in (NW9) by at least floor(R\_beta/1220), and its digit error by only ab(z\_beta−141) or bc(z\_beta−141). Thus each L\_j increases. Both L\_j and G are smallest for all thin and W=k−2. Direct substitution proves (NW12) and positivity of each L\_j there. This minimizes over a possibly larger role set than any completed source Host, which is valid for a uniform lower bound; no all-thin Host is asserted. QED.

There is also an exact finite degree comparison. For \(D=D_0(W)+D_1(W)\), \[
 \boxed{\qquad
     G(W;N_T,N_H,N_L)/D>1/1220-2^{-455}.
 \qquad}                                                 \tag{NW14}
\] The two floors lose less than2. With zmax=278 or279, the whole error is at most f(ab+bc)zmax+2+S\_f(a)+S\_f(b), respectively 75,828,772,861 or304,406,169,777. These are less than2\^{}37 and2\^{}39. The minimum D is at least2\^{}492 or2\^{}494. The error/D is therefore strictly less than2\^{}-455. This is a ratio of the proved arithmetic lower bound to raw derivative degree, not exact image content or a normalized DP44 gain/cost ratio.

\subsection{5. Actual selected highest nonvanishing transports}\label{actual-selected-highest-nonvanishing-transports}

\subsubsection{Proposition NW15: both compact weight ties and the middle line}\label{proposition-nw15-both-compact-weight-ties-and-the-middle-line}

The new compact oscillator on each edge has the actual summand \(W_{j+1}\otimes W_j^\vee\), with the source dual convention of LIFT25. The two primitive highest kernels pair the same middle line with value one.

\textbf{Proof.} At each label the literal recurrence is \[
 w_1=(w_0-\tfrac d2\mathbf1_a,
                   (-E_\beta+\tfrac a2)^{[d]}),\qquad
 w_2=(w_1-\tfrac d2\mathbf1_b,
                   (-E_\beta-2\ell+\tfrac b2)^{[d]}).
 \tag{NW16}
\] This is the accepted source lamTrue\_old, lamTrue\_new and nuW\_eq\_true, \nolinkurl{Columns/BaseLemmas.lean496–552,} at the new e entries. Thus, for j=0,1, \[
 \nu_j=w_j+e_j-\frac{N_{j+1}}2\mathbf1,\qquad
 w_{j+1}=(\nu_j,0^d)-e_j+\frac{N_j}2\mathbf1.             \tag{NW17}
\] These identities match \textbf{both} group actions coordinate by coordinate, for either value of every interior bit. The last nu0 entry is \(E_T-B-b/2>0\) at top, and \(E_\beta+2\ell-d/2>0\) otherwise. The last nu1 entry is Delta\textgreater2. Writing eta\_j=w\_j−min(w\_j), one obtains \(\eta_1=(\nu_0,0^d)\), \(\eta_2=(\nu_1,0^d)\), and \(\nu_1=\eta_1+\Delta\mathbf1_b\).

The principal-minor polynomials consequently satisfy \(p_{\nu_1}=\det^\Delta p_{\eta_1}\), with leading diagonal coefficient one. The determinant-insertion isometry in the fixed row metrics, LIFT25(AR8--AR9), gives \(h_1^\vee(h_1)=1\). Its proof is labelwise; neither the partitions nor e0 need be uniform across labels. Compact Schur averaging therefore gives the same nonzero consecutive highest contraction, with its unchanged dimension factors. QED.

\subsubsection{Theorem NW18: full analytic coupled injectivity}\label{theorem-nw18-full-analytic-coupled-injectivity}

For each word, let \(\mathcal H_{\rm ram,h_0}^\omega\) be LIFT27's actual initial ramified automorphic space, at the new T, K040 stages, and present compact weight. The raw coupled operator with these two primitive highest polynomials and source-normalized creation is injective on that space. The complete-coevaluation operator from its coefficient rows to full G2 functions is also injective.

\textbf{Proof by checking all changed hypotheses of LIFT27.} The local GJ holomorphy and rank margins use a,d and unitarity: \((d-a+2)/2=860\) or734 is unchanged. At G, K040 preserves the distinct nontrivial tame unit characters. The exact generic and sharp-middle quotient calculations use these characters and ranks, so LIFT27(MR15--MR22) applies without a finite Steinberg assumption.

At the new T the stages are unramified, the source has the same strict covector condition, and its norm-kernel is the sum of nontrivial residue dilation types. DO040 gives degree-one splitting and a residue characteristic exceeding the type bounds. The first tail test uses only b\textgreater=4a and q\textgreater=2; the second whole-Schwartz quotient uses unitary unramified stages and the nontrivial dilation type. Hence LIFT27(MR23--MR24) applies to every permitted T component, without recycling an old T mass.

At Q the two stage unit characters agree. Their ratio is unramified and unitary, so DO037's Gauss/graph quotient remains nonzero for the new values. At other finite places the relative field, selfdual lattices, and basic vectors are unchanged. Their tested spherical quotients allow arbitrary unitary unramified stage values. In particular29 remains maximal. NW15 proves the exact real selected occurrence and consecutive-highest alignment. The compact local parameters and angular twists give the same finite nonzero real gamma factors.

Globally put \(\upsilon=\chi_0\chi_1^{-1}=\mu_0\mu_1^{-1}\); the tuple cancels. HasInfQ(−e\_j) means angular exponent −2e\_j. Thus upsilon has angular exponent4ell at \textbf{every} label, and its square has exponent8ell. Both are nontrivial. LIFT26's exact boundary lemma at1, carried in LIFT27(MR25), therefore applies. Higher Hecke factors are at integers at least2, and the initial Euler product retains absolute-convergence margin860 or734. Neither e0 uniformity nor the old numerical equation e1−e0=4C is used once these hypotheses are checked.

The first global edge still has b\textgreater2a and the second b\textless c\textless2b with anisotropic target. Thus the accepted norm identities and every-embedding nonvanishing have exactly the same valid scope. LIFT27's multiplicity factorization, fixed moving-frame tensor identity, and local injectivity prove the assertion. The complete coevaluation comparison is its exact factor d0 d1. All arguments use actual characters and their stated local properties, not an old uniform numerical specialization. QED.

The initial space is nonzero for each word. DO040's explicit covector gives projective class stabilizer mu28 and strict stabilizer one on the identity support. Every character of mu28 extends to a \textbf{nontrivial} residue character at T, since each restriction fibre has (T−1)/28\textgreater1 elements. Choose the T characters with the last cancelling the actual compact/G central character, as in R017/DO040. The class-supported highest coefficient with the finite G Steinberg generic vector is a nonzero Mram input. This provides one input; NW18 itself covers the full Mram space. Its norm uses the new T and actual word-dependent compact dimension.

\subsection{6. Normalizations and the unresolved terminal step}\label{normalizations-and-the-unresolved-terminal-step}

Write d\_j=dim W\_j, D\_(j,tau)=\textbar nu\_(j,tau)\textbar, and H\_nu for the hook product. With LIFT25(AR22)'s actual positive gamma\_tau, the Chern highest-vector squared Schwartz norm is \[
 A_j^2=\det(2Y)^{-7hN_jN_{j+1}}
       \prod_\tau H_{\nu_{j,\tau}}
                         (\pi\gamma_\tau)^{D_{j,\tau}}. \tag{NW19}
\] Source creation multiplies by \(\Lambda_j=\prod_\tau\lambda_\tau^{D_{j,\tau}}\); its squared norm is \(\det(2Y)^{-7hN_jN_{j+1}}\prod_\tau H_{\nu_{j,\tau}}\). The literal compact coefficient integrals have 1/(d0 d1). The auxiliary isometric quotient has the corresponding square-root dimension factors. Hooks and dimensions use the actual partitions; LIFT25(AR26)'s old uniform numerical tables are inapplicable.

The arithmetic-to-analytic comparison, only after literal tilQ, is \[
 v_0v_1 C_\theta^{14h(ab+bc)}
       \prod_\tau(\lambda_\tau/cD_\tau)^{
                         D_{0,\tau}+D_{1,\tau}}.        \tag{NW20}
\] Here v0,v1 are the actual strict level volumes at the new T. Projective-T/full-G descriptions retain K038's weighted averages and matching volume refinements. The cD factors are11-units. The complex periods, source-creation factors, and real metrics are not assigned finite valuations in NW10. The finite input character field remains an arithmetic cost.

NW18 concerns full G2 functions. Restriction by tilQ, terminal brC, and PtilFF can still create a kernel. NW10 is valid even if the restricted map is zero. An upper content witness on the preceding image, a terminal nonzero scalar, and full e/H and gain/error accounting remain separate. This is actual new-profile image divisibility, not a source gainD value or GRH conclusion.

\subsection{Sources and check scope}\label{sources-and-check-scope-1}

Carried checked inputs are K027(EX2--EX4), K033(CS1--CS5), K034(TG3--TG10), K035(PF22--PF29), K037(CR8--CR14,CR35--CR40), corrected K038(CA27--CA36), K039(IG2--IG13), K040(NC17--NC28), R022(R22.1--R22.14), DO040, LIFT25(AR2--AR30), and LIFT27(MR8--MR32). R017 supplies the class-supported construction with DO040's new prime/covector and central-character extension. No accepted global theta theorem is rederived.

Direct pinned source anchors: \nolinkurl{Formalism/LocalToGlobal/Columns/BaseLemmas.lean496–552} (lamTrue\_old, lamTrue\_new, nuW\_eq\_true); \nolinkurl{Formalism/LocalRingCardinality/DataCoeff.lean21–23,39–41,51–55,83–84} (determinant shift, HC conversion, Weyl extremal, row norm); \nolinkurl{Formalism/LocalToGlobal/UnitaryNorms/IsUnitaryNormOf.lean90–94} (HasInf, HasInfQ), all in recovered commit 7cbf7a0c32fca1badf05931879a37c3f8a2b03b4.

The new proof is the all-word/role degree specialization, the actual geometric and character premise check, and the coupled-image bound. The replay independently encloses the defining exponentials, checks each Hamming weight by an explicit word, checks all local bit choices needed for arbitrary words, verifies both oscillator weight ties and digit endpoints, evaluates the bounds, and enumerates every role-count triple as a redundant arithmetic check. It does not enumerate 2\^{}(k−2) words, calculate a theta coefficient or image Smith form, prove a Rallis theorem by arithmetic tests, or build Lean.

\clearpage

\section{AM-R-024: actual terminal divisibility for the repaired profile}\label{am-r-024-actual-terminal-divisibility-for-the-repaired-profile}

Author: GPT-6 Astra Pro, ALP\_MACHINE \nolinkurl{/root}, 22 November 2026 UTC. Reasoning effort is not separately exposed.

This finding transports the terminal arithmetic lower bound to the R022/K040/DO040 family whose actual coupled kernel is K041. It uses the already proved whole-Weyl inclusions R020, rather than identifying the full branch with its highest line. There are two conclusions. For the literal Weyl detector sublattice MtildeLdel, every role assignment has a large terminal divisibility bound without a hook denominator. For the full literal MtildeL=intNab detector lattice, every assignment with at least one top label also has a positive explicit bound, now obtained from a long interval in the actual integer spectral factor. Neither conclusion asserts that a terminal value is nonzero or that the normalization is cheap in the DP44 budget.

\subsection{1. Actual terminal weights, with the word retained}\label{actual-terminal-weights-with-the-word-retained}

Keep the notation of K041: k=f is42 or168, d is1760 or1634, c=k+2d, n=c−1, ell is45,156,049,316 or40,530,339,696, A=2\^{}960, B=2\^{}955, and the role integer is E=E\_T,E\_H,E\_L. Write W for the Hamming weight of the binary word. At one CM label,

\[
 e_0=E,\quad e_1=E+2\ell,\qquad b_0=E-\ell,\quad b_1=E+\ell,
 \qquad v_i=4\ell(k-1-i)\quad(0\le i\le k-2).
                                                        \tag{R24.1}
\]

The literal column weight lambda is K041's w2, not its selected polynomial partition nu1. The literal row weight mu is RowPar.sgTrue. In zero-based initial indices their complete lists are

\[
\begin{split}
 \lambda={}&
 (u_i-(k-1)/2+i-d)_{i=0}^{k-1},\quad
 (-E+(k-d)/2)^{[d]},\quad
 (-E-2\ell+(k+d)/2)^{[d]},\\
 \mu={}&
 (v_i-(c-2)/2+i)_{i=0}^{k-2},\quad
 -b_0+k/2-d,\quad
 (-E+(k-d)/2)^{[d-1]},\quad
 -b_1+k/2,\quad
 (-E-2\ell+(k+d)/2)^{[d-1]}.
\end{split}                                             \tag{R24.2}
\]

Commas between displayed groups mean concatenation. These formulas follow by substituting (R24.1) in the source lamTrue\_old/new and RowPar.sgHC/sgTrue\_eq formulas already used in R019 and K041. All entries are integers. They interlace and have the common last entry s=−E−2ell+(k+d)/2. Set eta=mu−s1\_n and D\_mu(E)=sum eta\_i. Direct summation gives the word-independent integer

\[
 D_\mu(E)=(k-1)
       \left(E+2(k+1)\ell-\frac{k+3d}{2}\right)
                  +(d+1)(2\ell-d).
                                                        \tag{R24.3}
\]

For clarity, put a\_i=mu\_i−lambda\_(i+1), b\_i\^{}Br=lambda\_i−mu\_i, using zero-based indices 0\textless=i\textless n.~The nonzero a entries are

\[
 a_i=\ell\omega_i\ (0\le i\le k-3),\qquad
 a_{k-2}=
 \begin{cases}
 B+4\ell+(k-2)/2,&\text{top},\\
 2\ell-1,&\text{thick or thin},
 \end{cases}
\] \[
 a_{k-1}=a_{k+d-1}=\ell-d/2,
                                                        \tag{R24.4}
\] and all others vanish. The initial b\^{}Br entries are u\_i−v\_i−1/2 for i\textless k−1, the next is u\_(k−1)+b0−1/2, and b\^{}Br\_(k+d−1)=ell−d/2; all remaining entries are zero. R022 makes all these quantities nonnegative for every arrangement of the bits. This also checks interlacing without replacing arbitrary words by a Hamming representative.

Consequently R019's d0=sum a\_i and nBr=sum lambda−sum mu are

\[
\begin{array}{c|cc}
 &d_0&n_{\rm Br}\\ \hline
\text{top}&B+(W+6)\ell+(k-2)/2-d
          &A-B-(W+6)\ell+1-k/2\\
\text{thick/thin}&(W+4)\ell-d-1
          &(8k-8-W)\ell+2-k/2.
\end{array}                                             \tag{R24.5}
\]

These are exact changed-weight formulas. In particular the old uniform numerical R019/R020 tables do not apply.

\subsection{2. The two actual detector lattices}\label{the-two-actual-detector-lattices}

At each coefficient-field place above11 use v(11)=1, and choose the actual class representatives and transported integral frames of K039(IG3)/K041(NW8). The source explicitly distinguishes

\[
 \mathcal M_\nabla=\mathrm{MtildeL}
             =\mathrm{mtildeGSD.intNab},\qquad
 \mathcal M_\Delta=\mathrm{MtildeLdel}
             =\mathrm{mtildeGSD.intDel}.
                                                        \tag{R24.6}
\]

They have the same EqFun character, residue-action and stabilizer conditions. RowForms.lean281--286 gives these exact definitions. The coefficient inclusion Delta\_mu subset Nabla\_mu implies M\_Delta subset M\_Nabla; thus M\_Delta is an allowed sublattice of the literal detector lattice, not an invented unconstrained row space.

Let C=product\_tau cBr\_tau, computed with the actual weights (R24.2), and H=product\_tau H\_(eta\_tau). R019 proves every cBr\_tau is a nonzero integer. R020(R20.9) gives, with uT the literal piT power,

\[
 B_{\rm term}(\Delta_{\boldsymbol\lambda})
       \subset u_TC\nabla_{\boldsymbol\mu}
       \subset (u_TC/H)\Delta_{\boldsymbol\mu}.
                                                        \tag{R24.7}
\]

The new piT is an11-unit, so v(uT)=0. The normalized Fischer pairing is integral on Nabla\_mu x Delta\_mu. If G\_cpl is K041(NW9), the exact class formula DO038(DO38.9)/K038(CA30) therefore gives

\[
 v\bigl(\operatorname{PtilFF}(P_{\rm cpl}(P),Z)\bigr)
       \ge G_{\rm cpl}+v(C)-S_f(n)
                         \qquad(Z\in\mathcal M_\Delta),
                                                        \tag{R24.8}
\] \[
 v\bigl(\operatorname{PtilFF}(P_{\rm cpl}(P),Z)\bigr)
       \ge G_{\rm cpl}+v(C)-v(H)-S_f(n)
                         \qquad(Z\in\mathcal M_\nabla).
                                                        \tag{R24.9}
\]

Here P is in the actual initial integral lattice or Mram, and S\_f(n)=2floor(2fn/5)+v11(floor(2fn/5)!) is K039's proved stabilizer bound. It is125,628 on J14 and484,744 on J56.

\textbf{Proof.} K041 places every coupled coefficient in 11\^{}G\_cpl times the actual final Weyl lattice, including all initial and middle stabilizer denominators. Apply the first inclusion of (R24.7) and pair with the actual Weyl detector coefficient. Perfect Weyl/dual-Weyl integrality gives (R24.8), after dividing by each terminal class stabilizer and applying K039's bound. For the full dual-Weyl detector instead use the second inclusion of (R24.7). The same pairing and class sum prove (R24.9). K041's local frames and piT-unit transport preserve all these inclusions. No projection onto an invariant row has been applied with an unrecorded denominator; the rows already satisfy their exact EqFun constraints. ∎

Since v(C)\textgreater=0, the exact K041 lower bound minus S\_f(n) is still at least2\^{}481 on J14 and2\^{}483 on J56. Thus (R24.8) gives those uniform terminal divisibility exponents for every role assignment and every binary word. This result does not require a top label.

\subsection{3. A large spectral interval controls the full detector lattice}\label{a-large-spectral-interval-controls-the-full-detector-lattice}

Assume there is at least one top label. Put

\[
 a_* = B+4\ell+(k-2)/2,\qquad
 C_*=n\,v_{11}(a_*!),\qquad
 D_T=D_\mu(E_T),\quad D_H=D_\mu(E_H),\quad D_L=D_\mu(E_L).
                                                        \tag{R24.10}
\]

At every top label, the \textbf{actual integer} spectral factor satisfies

\[
                         v_{11}(c_{\rm Br})\ge C_*.
                                                        \tag{R24.11}
\]

\textbf{Proof.} Use the spectral product R019(R19.12) at index i=k−2. By (R24.4), its lower endpoint is L\_i=mu\_i−a\_* and its upper endpoint U\_i is at least mu\_i. Its product therefore contains, for each of the n column indices j, every factor

\[
     \prod_{v=\mu_i-a_*}^{\mu_i-1}
                         (v-(\mu_j+i-j)).
                                                        \tag{R24.12}
\]

The omitted zero in the full product is v=mu\_i,j=i; it is absent from (R24.12). For j=i the absolute product is a\_\emph{!. For j\textless i all terms are negative, and for j\textgreater i all terms are positive: these are the strict endpoint inequalities proved in R019 for interlacing weights. In either case the absolute values are a\_} consecutive positive integers. Their product is a\_\emph{! times an integer binomial coefficient. Each of the n products is thus divisible by a\_}!, while all remaining spectral factors are nonzero integers. This proves (R24.11), without enumerating the long interval. ∎

For every partition eta of D, the carried hook-length identity D!/H\_eta=dim Specht\_eta is a positive integer (R011, with Fulton and Harris's exact locator below). Therefore

\[
                    v_{11}(H_\eta)\le v_{11}(D!)\le D/10.
                                                        \tag{R24.13}
\]

For the exact two profiles, D\_T\textgreater=D\_H\textgreater=D\_L and 10C\_*\textgreater D\_T−D\_H. Define the explicit integer

\[
 B_* = C_*-\left\lfloor\frac{D_T+(f-1)D_H}{10}\right\rfloor.
                                                        \tag{R24.14}
\]

Then, for every assignment with N\_T\textgreater=1 top labels and every word,

\[
 \boxed{\quad v(C)-v(H)\ge B_*\ge2^{963}.\quad}
                                                        \tag{R24.15}
\]

\textbf{Proof.} All cBr factors are integral, and each top factor contributes at least C\_\emph{. The hook loss is at most floor((N\_T D\_T+N\_H D\_H+N\_L D\_L)/10), by (R24.13). Replacing each non-top label by a thick label only increases this upper bound. Increasing N\_T by one increases the difference between the proposed lower and upper bounds by at least C\_}−ceil((D\_T−D\_H)/10), which is nonnegative because 10C\_\emph{\textgreater D\_T−D\_H. Thus its minimum for N\_T\textgreater=1 is at N\_T=1, with all remaining labels thick. This gives (R24.14). The exact integers E\_T,E\_H from R022/K041, formula (R24.3), and Legendre's floor sum give B\_}\textgreater=2\^{}963 on both hosts. These last inequalities are exact integer comparisons in the saved certificate; there is no asymptotic approximation to a factorial or exponential. The one-top/all-thick comparison is a numerical worst-case bound, not an assertion that this role assignment supplies a full Host. ∎

Combining (R24.9), (R24.15), and the positive K041 lower bound minus the terminal stabilizer loss gives

\[
 \boxed{\quad
 v\bigl(\operatorname{PtilFF}(P_{\rm cpl}(P),Z)\bigr)
       \ge G_{\rm cpl}+B_*-S_f(n)\ge2^{963}
       \quad(Z\in\mathcal M_\nabla,\ N_T\ge1).
 \quad}                                                  \tag{R24.16}
\]

The coefficient inclusion B\_term(Delta\_lambda) subset 11\^{}B\_* Delta\_mu is itself valid on the whole standard tensor Weyl module in these local frames. It is an inclusion before restricting to the preceding image, and remains an inclusion after restriction. No upper content estimate on that restricted image follows.

\subsection{4. Nonzero values and costs remain distinct}\label{nonzero-values-and-costs-remain-distinct}

All displayed valuation bounds include zero, whose valuation is infinite. K041's injectivity into full G2 functions does not prove that the literal lowered class family is nonzero. If that lowered family is nonzero and the full compatible terminal row space of DO038 is admitted, its nondegenerate class pairing has a detecting row in M\_Delta after clearing denominators: M\_Delta is a full lattice in that compatible field-vector space. This is the usual field-spanning argument for a lattice; it is not a unit-valued or bounded-norm detector construction.

The large number in (R24.16) includes the mandatory integer cBr. R019(R19.15--16) retains the corresponding real norm

\[
 \log\|B_{\rm term}\|
 =-\tfrac12\log T\sum_\tau n_{{\rm Br},\tau}
    +\sum_\tau\log|c_{{\rm Br},\tau}|
    +\tfrac12\sum_\tau\log\kappa_\tau.
                                                        \tag{R24.17}
\]

Here nBr is the new formula (R24.5), and cBr/kappa are the exact factorial products R019(R19.4,R19.13) at (R24.2). Thus the integer gain is not inserted without its archimedean normalization, nor counted twice as a new local slope. The global mixed Host/budget replacement isolated by R023, terminal image nonvanishing, a detector of adequate norm, and the actual complete DP44 comparison remain open.

There is an exact way to separate this scalar, rather than only a warning about it. For the fixed word and role assignment, C and uT are independent of the terminal class. Put

\[
 s_B=u_TC,\qquad
 \widehat B=s_B^{-1}B_{\rm term}
                      =\bigotimes_\tau i_\tau^\dagger.
                                                        \tag{R24.18}
\]

The last equality is the already proved R020(R20.3). Its full integral inclusions are

\[
 \Delta_{\boldsymbol\mu}
       \subset\widehat B(\Delta_{\boldsymbol\lambda})
       \subset\nabla_{\boldsymbol\mu},\qquad
 -v(H)\le\operatorname{content}_v(\widehat B)\le0.
                                                        \tag{R24.19}
\]

Thus the normalized whole branch has no positive compulsory divisibility on its full Weyl source; its image contains the primitive target Weyl lattice. This does not bound its content on the preceding theta image above, because that image need not contain the vectors realizing the left inclusion.

Choose a number field M containing the actual algebraic scalars. For every complex embedding sigma of M, all conjugates of piT have absolute value sqrt(T), because piT belongs to Q(sqrt(−7)) and has rational norm T. Hence, with D\_Br=sum\_tau nBr\_tau,

\[
 \log|\sigma(s_B)|=\log|C|-\frac{D_{\rm Br}}2\log T,
 \qquad \log\|\widehat B\|=\frac12\sum_\tau\log\kappa_\tau.
                                                        \tag{R24.20}
\]

With ord\_w normalized by a local uniformizer, the product formula and R002's scalar-invariant normalization balance give the exact scalar cancellation

\[
 \frac1{[M:\mathbb Q]}
   \sum_w\operatorname{ord}_w(s_B)\log Nw
 =\frac1{[M:\mathbb Q]}
   \sum_\sigma\log|\sigma(s_B)|
 =\log|C|-\frac{D_{\rm Br}}2\log T.
                                                        \tag{R24.21}
\]

This is an application of R002 and R020, not a new product-formula theorem. It computes the signed contribution of this actual scalar as zero after both finite and real terms are included. It does not declare the entire terminal restriction cost zero, nor assign the remaining coupled-image, detector, field or error terms of DP44. In particular the large cBr contribution in (R24.16) cannot by itself be a net lift saving.

\subsection{5. Framework, sources and verification}\label{framework-sources-and-verification}

This is an inclusion of actual coefficient-image lattices over the DVR above11, paired with one of the source's two specified integral row lattices. It refines the integral representation/flag picture of SGA3I4.7.3/4.7.3.1 and SGA6VI3.3, rather than replacing lattice indices by equality of rational representation classes. The actual theta section retains the Picard and biextension locators FGA232 Theorem3.1, pp148--149; FGA236 Corollary1.7, p229, and2.3, p232; SGA7I VII2.7/2.9, pp169--174, andVIII7.4(b), pp310--311. Its formal jet realization retains EGAIV4 16.11.2--3(ii), pp54--55, and the Cartier--Dieudonne context at cote4/cote54 scans5--7 already fixed in the dictionary. There is no new motivic, ell-adic-weight, cycle nonvanishing, or Kummer RH conclusion.

Exact inputs and locators:

\begin{itemize}
\tightlist
\item
  AM-K041(NW1--20): actual all-word coupled kernels, frames and lower bounds, with R022, K040 and DO040's specified data.
\item
  AM-R-020(R20.2,R20.7,R20.9): whole Weyl/dual-Weyl inclusions and normalized pairing; AM-R-019(R19.1,R19.4,R19.12--16): exact spectral product, real norm and same-shift normalization.
\item
  AM-K039(IG3,IG11--13), AM-K038(CA27--30) and AM-DO-038(DO38.9): actual local frames, class stabilizers and literal PtilFF formula.
\item
  Alpöge, L. (n.d.). \emph{Remark}, step8, associated archive commit 7cbf7a0c32fca1badf05931879a37c3f8a2b03b4: \nolinkurl{Formalism/LocalToGlobal/Columns/RowForms.lean},281--286; \nolinkurl{AbelianHostFieldConditions.lean},408--421; \nolinkurl{BaseLemmas.lean},496--552,575--582; \nolinkurl{FiniteSumNonzeroLargeDepth.lean},41--55,118--151; \nolinkurl{Formalism/LocalRingCardinality/FactDelDualNabla.lean},75.
\item
  Fulton, W., \& Harris, J. (1991). \emph{Representation theory: A first course}. Springer. Theorem6.3(2)--(4), pp77--78, and Exercise6.11(b), p80, as already checked and carried in AM-R-011. The hook identity is an input; it is not rederived here.
\end{itemize}

The companion certificate checks the literal row formula independently, the complete terminal interlacing lists, exact d0/nBr and hook degrees, 624 Hamming representatives with the bitwise identities (R24.4), the long top spectral interval, every displayed finite integer bound, and the terminal stabilizer loss after the carried coupled-map bound. It does not compute a theta coefficient or decide its nonvanishing. Whole-proof readings are recorded separately in CHECK\_EVENTS.tsv.

\clearpage

\section{ALP-LIFT-30: the primitive norm-minus-T line behind the bottom translation}\label{alp-lift-30-the-primitive-norm-minus-t-line-behind-the-bottom-translation}

Author: /root/lift\_accounting, GPT-6 Astra Pro, ALP\_MACHINE. 22 November 2026 UTC.

This identifies LIFT29's finite translation with an explicit rational primitive negative line. It also identifies the necessary real-vector transport. The finite translation is not the source's tS merely because its intersection with the compact subgroup has the desired T level: the source tS conjugates the entire bottom level to the row level. The norm-minus-T Bessel embedding is a separate part of that source.

\subsection{1. Actual arithmetic inputs}\label{actual-arithmetic-inputs}

Let E/F be either checked mixed host K031 or K032, and let a be its chosen initial rank. Work on the actual rational Hermitian space

\[
 V=E^{a+1},\qquad J=\operatorname{diag}(I_a,-1),\qquad
 H=U(I_a)\subset G=U(V).
 \tag{NL1}
\]

Let T be one of the actual primes of DO040, or any odd rational prime with the same norm and congruence properties. Thus

\[
 \pi\bar\pi=T,\qquad (\pi,\bar\pi)=1,\qquad
 \pi\equiv1\pmod{(g_0g_1)^2},
 \tag{NL2}
\]

with pi integral and T completely split in E. These properties for the specified large primes are carried from DO040.1 and DO040.6--7 in definite\_oscillator\_large\_t\_primes.md; their primality is not reproved. At a selected T place, pi is a uniformizer and its conjugate is a unit. The other component of the relative split unitary group is always the paired inverse-transpose component; a displayed GL matrix below refers to the selected component.

Set

\[
 A=(T+1)/2,\qquad B=(T-1)/2.
 \tag{NL3}
\]

Both are rational integers, A-B=1, and A squared minus B squared is T. In the plane spanned by the first positive coordinate and the negative coordinate, define

\[
 g=
 \begin{pmatrix}
  A/\pi&B/\bar\pi\\
  B/\pi&A/\bar\pi
 \end{pmatrix}.
 \tag{NL4}
\]

On every other coordinate g is the identity.

\textbf{Proposition.} The element g lies in SU(V)(F). The integral vector

\[
 z_T=B e_1+A e_{a+1}
       =g(\bar\pi e_{a+1})
 \tag{NL5}
\]

is primitive, has norm minus T, and its stabilizer is

\[
 H_T=gHg^{-1}.
 \tag{NL6}
\]

Its Hermitian orthogonal complement is therefore rationally isometric to I\_a, with the explicit isometry obtained by restricting g.

\textbf{Proof.} The first column of the displayed two-plane matrix has norm (A squared minus B squared)/(pi times conjugate pi)=1. The second has norm minus one. Their Hermitian cross product is AB/(conjugate pi squared) minus BA/(conjugate pi squared), hence zero. The determinant is (A squared minus B squared)/(pi times conjugate pi)=1. This proves the special-unitary assertion and (NL5). The norm of z\_T is B squared minus A squared, or minus T. Because A-B=1, the coordinates generate the unit ideal of O\_E; equivalently the map O\_E to O\_E\^{}(a+1) with value z\_T has an integral left inverse. Its image is a saturated rank-one summand. Stabilizing the nonzero scalar multiple bar(pi)e\_last is the same as stabilizing e\_last. Conjugation by g proves (NL6). The restriction of g to the orthogonal complement of e\_last gives the last assertion. QED.

This supplies the actual primitive negative vector and positive rational complement. It does not by itself assert every extra field of the source's FlagS, FixedHyp, SBes, or transport records.

\subsection{2. Exact finite factorization}\label{exact-finite-factorization}

At a selected T place put c=bar(pi) and

\[
 \tau_c=
 \begin{pmatrix}I_a&0\\\pi^{-1}c e_1^*&1\end{pmatrix}.
 \tag{NL7}
\]

On the distinguished two-plane, direct multiplication gives

\[
 \begin{split}
 g^{-1}&=\tau_c k,\\
 k&=
 \begin{pmatrix}
 A/\bar\pi&-B/\bar\pi\\
 -\bar\pi&\bar\pi
 \end{pmatrix}\in\operatorname{SL}_2(\mathcal O_{F_v}).
 \end{split}
 \tag{NL8}
\]

Indeed the lower-left entry of tau\_c inverse times g inverse is minus (A+B)/pi=minus bar(pi), and its lower-right entry is bar(pi). The determinant of k is A-B=1. All entries are integral because bar(pi) is a unit. This proves membership, including integrality of the inverse, not merely a determinant valuation.

Away from T both pi and bar(pi) are units, so g and its inverse are integral. At each selected G place, (NL2) gives T congruent to one, A congruent to one, and B congruent to zero modulo g\_0g\_1. The prime 2 is distinct from both G primes. Consequently g is identity modulo the relevant G prime.

Let J\_B be LIFT29's standard-frame dat0-shaped bottom compact: at G its reduction is diag(U\_a,1), and elsewhere it is maximal. Let chi\_B be its generic G character and be trivial elsewhere. Let tau\_f be the product of the elements (NL7) at T and identity elsewhere. Then (NL8) and the preceding integrality assertions prove

\[
       g_f^{-1}=\tau_f k_f,\qquad
       k_f\in J_B,\qquad \chi_B(k_f)=1.
 \tag{NL9}
\]

There is no finite scalar in this identity for automorphic functions of right type chi\_B inverse.

\subsection{3. An arbitrary actual primitive T covector}\label{an-arbitrary-actual-primitive-t-covector}

The chosen covector of DO040 is generally not e\_1. Let u\_v be any primitive integral row at the selected T place, including its actual row (1,2,\ldots,2\^{}(a-1)). Choose

\[
 a_v\in\operatorname{GL}_a(\mathcal O_{F_v}),
       \qquad c e_1^*a_v=u_v.
 \tag{NL10}
\]

Such a matrix exists: a primitive row extends to an integral basis, and c is a unit. Complete its relative unitary pair in H(F\_v). Let a\_f be these matrices at T and identity elsewhere, and put hat(a\_f)=diag(a\_f,1). There is the exact identity

\[
 \tau_{u_v}
   =\widehat a_v^{-1}\tau_c\widehat a_v.
 \tag{NL11}
\]

The rightmost hat(a\_v) is in J\_B and has trivial chi\_B. Therefore, for any right-chi\_B-inverse automorphic function F and any initial automorphic row function P,

\[
 \begin{split}
 \int_{[H]}P(h)F(h\tau_u)\,dh
 &=\int_{[H]}P(h a_f)F(h\tau_f)\,dh.
 \end{split}
 \tag{NL12}
\]

Here a\_f is the H element, without its ambient block embedding. The equality is right invariance of Haar measure after the substitution h\_old=h\_new a\_f, followed by (NL11). The strict groups are related by J\_u=a\_f inverse J\_c a\_f.~Hence P\_c(h)=P(h a\_f) is exactly the corresponding initial row at the coordinate-covector level. This is an isometric change of finite class presentation. Rational stabilizers, class masses and the unweighted T diamond-norm kernels are transported by conjugation, so the whole Mram subspace is carried to the corresponding Mram subspace. No projective-level coarsening or character average is involved.

The chosen a\_f is an actual finite coordinate transport, not a claim that the unchanged source transport tr already has this value.

\subsection{4. The real component must also be transported}\label{the-real-component-must-also-be-transported}

A rational element is not a finite adelic element. Let sigma be the actual unitary bottom discrete-series representation and let

\[
 {\cal W}:\sigma^\infty\longrightarrow{\cal A}_{\rm cusp}(G)
\]

be an equivariant automorphic intertwiner, with right action R(x)F(y)=F(yx). For any smooth vector v put

\[
                  v_T=\sigma(g_\infty)v.
 \tag{NL13}
\]

Then, with h in H(A\_F),

\[
 \begin{split}
 {\cal W}(v_T)(g h g^{-1})
 &= {\cal W}(v)(g h g^{-1}(g_\infty,1))\\
 &= {\cal W}(v)(h(1,g_f^{-1}))\\
 &= {\cal W}(v)(h\tau_f).
 \end{split}
 \tag{NL14}
\]

The first equality is equivariance, the second is left G(F) invariance and the separation of finite and infinite components, and the last is (NL9). This proves the exact norm-minus-T Bessel interpretation of the translated detector.

If v belongs to the selected compact H type at the standard negative line, v\_T belongs to that same type under H\_T at the conjugate negative line. Unitarity gives norm(v\_T)=norm(v). Equivalently, the selected compact-type embedding is conjugated together with the negative line. It is incorrect to omit (NL13) while keeping the original fixed compact vector.

At every real place the two nontrivial singular values of g are sqrt(T) and 1/sqrt(T). Indeed its columns differ by unit phases from the real matrix with entries A/sqrt(T),B/sqrt(T), whose eigenvalues are (A+B)/sqrt(T)=sqrt(T) and (A-B)/sqrt(T)=1/sqrt(T). Thus the equality of Hilbert norms is a unitary-representation fact; it is not a bound of one for algebraic matrix entries, coordinate heights, finite polynomial coefficients, or selected jet frames. Those distinct costs have not been discarded.

At 11 the entire matrix g and its inverse are integral because T is different from 11. In particular the ambient algebraic lattice change is invertible at 11. This assertion does not identify an integral theta trivialization or an integral discrete-series model.

\subsection{5. Exact comparison with the source pins}\label{exact-comparison-with-the-source-pins}

The following source statements were read in this run at recovered commit 7cbf7a0c32fca1badf05931879a37c3f8a2b03b4.

\begin{itemize}
\tightlist
\item
  \nolinkurl{HaarTransversals/HidaSetting.lean117–130} defines K2 by congruence modulo g=g0g1 inside the full lattice compact. Hence its localization away G is maximal. Lines210--226 specify the norm-minus-T vector, its lattice and flag conditions. Lines267--276 put xu in that lattice compact and require OneOff away G.
\item
  \nolinkurl{Definitions/LatticeFrames/ThetaIdentityOfMapsAll.lean68–73} defines rowJ, rowChi and Matches. Matches conjugates the entire compact and its character, not only their intersection with H. Lines150--153 define the additional finite shift using xu.
\item
  \nolinkurl{IndefiniteReal/RealPins.lean47–52} requires a finite tS and the full conjugacy together with RowLevelMatchAt. Its betaStd field at line56 constrains betaV and betaX; it does not set betaY to one.
\item
  \nolinkurl{FlagStabilizerTensors/RowLevelMatch.lean22–47} divides the further pin into LevelMatchA, LevelMatchB and LevelMatchC. A is the compact character match. B also requires the actual full residue-group action on the automorphic multiplicity module, and C equates the selected idempotent coefficients. These latter clauses are not consequences of a compact intersection calculation.
\end{itemize}

For our standard-frame J\_B, tS therefore cannot be identified with the nonintegral T element tau merely from LIFT29: the row compact is maximal at T, whereas tau conjugates that maximal compact. Equations (NL5)--(NL14) place tau in the Bessel embedding and its finite/real coordinate transport instead.

They provide an explicit primitive negative line, its rational positive complement, and the exact transported detector. They do not identify the completed source's chosen betaY, tS, tr, SBes, RowLevelMatchB/C, or its idempotent-selected bottom space. The holomorphic detector still has to be placed in the actual prescribed downward theta image. No terminal period, DP44 gain, or full D7 construction is claimed here.

\clearpage

\section{AM-K042 --- the existing strict levels are neat, and the bottom twist has an exact finite window}\label{am-k042-the-existing-strict-levels-are-neat-and-the-bottom-twist-has-an-exact-finite-window}

Author: GPT-6 Astra Pro, ALP\_MACHINE /root/kummer\_extension, 22 November 2026 UTC. Inherited reasoning effort is not separately exposed.

This finding proves that the \textbf{existing} strict initial level and the dat0 bottom level have neat arithmetic groups. Consequently every initial definite class stabilizer is trivial, improving the actual K041 bound without a level refinement. It also computes the exact selected bottom jet degree and the entire canonical-twist interval that retains R022's windows and K040's character construction.

The comparison with the jet degree is not an interpolation-impossibility theorem. An effective coherent-jet or selected-period estimate remains required. No terminal nonzero value, full Host/SizeData₂, or DP44 cost sign is asserted.

\subsection{1. A direct local neatness criterion}\label{a-direct-local-neatness-criterion}

Call a characteristic-zero matrix neat if the subgroup of the algebraic multiplicative group generated by all its eigenvalues contains no root of unity other than one.

\subsubsection{Lemma NB1}\label{lemma-nb1}

Let q be a rational prime and \(M\in{\rm GL}_D(\mathbb Q_q)\). Suppose its characteristic polynomial lies in \(\mathbb Z_q[X]\) and reduces to \((X-1)^D\). If q\textgreater D+1, then M is neat.

\textbf{Proof.} If lambda is an eigenvalue, y=lambda−1 satisfies \[
 y^D+c_{D-1}y^{D-1}+\cdots+c_0=0,\qquad c_i\in q\mathbb Z_q.
 \tag{NB1}
\] Normalize v(q)=1. If v(y)\textless1/D, the leading term has strictly smaller valuation than every other term. For v(y)\textless=0 this follows from \(Dv(y)<1+iv(y)\); for \(0<v(y)<1/D\) it follows from \(Dv(y)<1\le1+iv(y)\). Both contradict (NB1). Therefore \[
                       v(\lambda-1)\ge1/D.              \tag{NB2}
\] This includes lambda=1. Each integer power of lambda, including negative powers, has the same lower bound on its difference from one. Multiplying finitely many such units preserves the lower bound. Every element of the eigenvalue-generated group is consequently in \(1+\{z:v(z)\ge1/D\}\).

For completeness, a nontrivial root of unity zeta with v(zeta−1)\textgreater0 has q-power order. Indeed a root of order m prime to q cannot be nontrivially near one: dividing (1+y)\^{}m−1 by y gives a unit when v(y)\textgreater0. The prime-to-q component of an arbitrary root must therefore be one. If zeta has exact order q\^{}s, then \[
 v(\zeta-1)=\frac1{q^{s-1}(q-1)}\le\frac1{q-1}.           \tag{NB3}
\] To prove the equality, the cyclotomic identity \(\prod_{(j,q)=1}(1-\zeta^j)=q\) has q\^{}(s−1)(q−1) factors. Their valuations agree because \((1-\zeta^j)/(1-\zeta)\) reduces to the q-unit j. Taking valuations proves (NB3). Roots with a nontrivial prime-to-q component instead have valuation zero. Since 1/(q−1)\textless1/D, no nontrivial root can satisfy (NB2). QED.

This proof controls \textbf{products} of eigenvalues, not merely the absence of finite-order matrices. It is stronger than the usual torsion-free conclusion needed for finite stabilizers.

\subsection{2. Application to the actual initial and bottom levels}\label{application-to-the-actual-initial-and-bottom-levels}

Let K be K031 or K032, with {[}K:Q{]}=2f and f=42 or168. For rank n, use the faithful rational representation on the underlying vector space of K\^{}n, of dimension \[
                              D=2fn.                    \tag{NB4}
\] The actual G data contain one chosen K prime from every conjugate pair above each rational g0,g1. Those rational primes split completely. Thus they contain every place of the real field above each g\_i. K031(HC12--HC13) gives42 chosen primes per rational prime; K032(QH27--28) gives168. This all-place condition matters when controlling products of eigenvalues from different embeddings.

At the literal initial level, the reduction at every selected G place is \(U_a(\mathbb F_g)\). The exact source is \nolinkurl{Columns/ColForms.lean93–115} together with \nolinkurl{Columns/GenericUnipotentCharacter.lean437–451:} HX(k,k,t,0)=HU(k), the full upper unipotent subgroup. At the dat0 bottom level of LIFT29(BT2--BT3), reduction is \[
                  {\rm diag}(U_a(\mathbb F_g),1).         \tag{NB5}
\] This is the exact m=2,rS=0 specialization of \nolinkurl{Divisibility/DefGroup.lean186–189,258–261,608–609.} It is not a projective T level or a full maximal G compact.

At a split place the second K component of a unitary matrix is the conjugate inverse transpose of the first. Unipotence is preserved by that operation. Therefore, in the rational representation (NB4), every block of every level element is unipotent modulo g. Conjugation by a finite class representative preserves the characteristic polynomial. Hence each arithmetic group \(G(\mathbb Q)\cap xJx^{-1}\) satisfies NB1 in dimension D.

The sufficient numerical bounds hold with a large margin:

\begin{longtable}[]{@{}lrrr@{}}
\toprule\noalign{}
Host & g0 & initial D=2fa & bottom D=2f(a+1) \\
\midrule\noalign{}
\endhead
\bottomrule\noalign{}
\endlastfoot
J14 & 4,615,409 & 3,528 & 3,612 \\
J56 & 22,285,932,761 & 56,448 & 56,784 \\
\end{longtable}

In particular g0\textgreater D+1 in both columns.

\subsubsection{Theorem NB6}\label{theorem-nb6}

Every arithmetic component group for the existing strict initial level and the existing dat0 bottom level is neat. Every initial definite class stabilizer is therefore trivial: \[
                         \Gamma_{0,x}=\{1\}.             \tag{NB6}
\] The statement remains true for every word, role assignment, new T-prime of DO040, and admitted tuple used in K041.

\textbf{Proof.} Apply NB1 to the characteristic polynomial just described. This works for every element and every class representative, in the faithful rational representation, so it is neatness of the actual restriction-of-scalars group. A finite subgroup of a neat group is trivial: a finite-order matrix has root-of-unity eigenvalues, hence all eigenvalues one; it is diagonalizable in characteristic zero, so it is the identity. Initial stabilizers are finite because the unitary group is definite. This proves (NB6). Neither T nor the archimedean weight changes the G-level premise. QED.

No principal29 level or new level at G was introduced. The statement does not say that R017's \textbf{coarser} projective-T/full-G support has trivial stabilizer: that support still has mu28, and its norm formula retains the factor28 and its G Schur denominator. Its refinement to the literal strict level still uses K038's weighted finite averages.

Nor does NB6 apply to the actual sharp middle level, whose reduction contains a GL\_(d−1) block. The terminal level has its own blocks. Their previous stabilizer estimates are not removed.

\subsubsection{Corollary NB7: exact improvement of K041}\label{corollary-nb7-exact-improvement-of-k041}

With the notation of K041(NW7--NW10), the actual arithmetic coupled map satisfies \[
 P_{\rm cpl}^\omega(\mathcal L_0^\omega)
       \subset11^{\,G(W;N_T,N_H,N_L)+S_f(a)}
                                  \mathcal L_{\rm ev}^\omega.
 \tag{NB7}
\] Thus one removes the loss1479 for J14 and23705 for J56.

\textbf{Proof.} In the literal K038 class sum the initial denominator is \(|\Gamma_{0,x}|=1\) in every summand. The middle denominator, the unnormalized Y sum, and every framed jet bound stay unchanged. Repeat K041(NW10)'s proof with no initial stabilizer subtraction. This also applies to its Mram sublattice. QED.

For LIFT29's translated compact packet the point stabilizers equal the initial ones, so they are all trivial as well. The bottom arithmetic quotient is already neat; a further cover is not needed to remove orbifold stabilizers. This does not supply an effective ample power, a jet interpolation bound, or a norm of a right inverse.

\subsection{3. Exact selected bottom jet order on the repaired profile}\label{exact-selected-bottom-jet-order-on-the-repaired-profile}

Retain R022's a=k, ell, A,B, E\_beta and binary word of Hamming weight W. LIFT28(HJ9) gives the exact selected degree \[
 N_\tau=\sum_i u_{\tau,i}-\sum_i v_{\tau,i}
                                      +b_{\tau,0}-a/2.
 \tag{NB8}
\] Since b0=E\_beta−ell and \(\sum_i v_i=2\ell k(k-1)\), substitution of the already proved sums in K041(NW4--NW5) gives \[
 \begin{split}
 N_T(W)&=A-B+E_T-\ell(W+5)-k+1,\\
 N_\beta(W)&=E_\beta+\ell(8k-7-W)+2-k
                                      \quad(\beta=H,L).
 \end{split}                                             \tag{NB9}
\] This is an exact formula for all words, not the degree of an unrelated ambient polynomial. The selected Pieri component has multiplicity one and this degree is unique, by the carried LIFT28 argument. A common HC determinant shift leaves NB8 unchanged.

The minimum occurs at W=k−2. Exact certified lower bounds are:

\begin{longtable}[]{@{}lrrr@{}}
\toprule\noalign{}
Host & top N\_min & thick N\_min & thin N\_min \\
\midrule\noalign{}
\endhead
\bottomrule\noalign{}
\endlastfoot
J14 & \textgreater2\^{}960 & \textgreater2\^{}494 & \textgreater2\^{}480 \\
J56 & \textgreater2\^{}960 & \textgreater2\^{}494 & \textgreater2\^{}478 \\
\end{longtable}

Full integer values and direct checks against NB8 are in the companion certificate. Large N by itself is not an obstruction: the minimal automorphic vector bundle also depends on the profile.

\subsection{4. The canonical-twist interval preserving these windows}\label{the-canonical-twist-interval-preserving-these-windows}

Use a new variable s to distinguish the canonical twist from the fixed gap ell. Hold u,v,b0,e0 fixed and put, at every label, \[
 q_s=(a+1)s,\qquad
 b_1(s)=E_\beta+\ell+q_s,\qquad
 e_1(s)=E_\beta+2\ell+q_s,\qquad s\in\mathbb Z_{\ge0}.
 \tag{NB10}
\] This is precisely LIFT28's ample canonical direction, with the retained flat determinant factor. It leaves the selected degree NB8 and the initial compact type unchanged. K040's specified character construction still applies exactly when e1(s) is a multiple28. Since gcd(a+1,28)=1 for a=42,168, this is 28\textbar s.

\subsubsection{Theorem NB11}\label{theorem-nb11}

Within the direction (NB10), retaining R022's fixed Scale, role centres, etaSz=4ell, and other listed window parameters, all the named R022 HC/centred/gap/window clauses hold exactly for \[
 s\in28\mathbb Z_{\ge0},\qquad
        s\le s_{\max}=28\left\lfloor
                           \frac{\ell}{28(a+1)}\right\rfloor.
 \tag{NB11}
\] Here ``exactly'' also includes the required multiple28 character condition; no assertion of a full SizeData₂ record is made.

\begin{longtable}[]{@{}lrrrr@{}}
\toprule\noalign{}
Host & s\_max & q\_smax & ell−q\_smax & number of candidate s \\
\midrule\noalign{}
\endhead
\bottomrule\noalign{}
\endlastfoot
J14 & 1,050,140,672 & 45,156,048,896 & 420 & 37,505,025 \\
J56 & 239,824,480 & 40,530,337,120 & 2,576 & 8,565,161 \\
\end{longtable}

\textbf{Proof.} The literal be\_close clause quantifies \textbf{all} pairs b\_i,e\_j: DataParams₂.lean164. The pair b0,e1(s) has difference \[
                          e_1(s)-b_0=3\ell+q_s.
 \tag{NB12}
\] Thus be\_close forces q\_s\textless=ell. Conversely, for 0\textless=q\_s\textless=ell, the four differences are ell,3ell+q\_s,ell+q\_s,ell, all at most4ell. This proves the exact restrictive inequality.

Write E\_beta=w\_beta+r\_beta with 0\textless=r\_beta\textless28. Then e1(s)\textless=w\_beta+3ell+28\textless w\_beta+4ell, so every role window still holds; the thick window allows an additional1 anyway. The earlier e0 is unchanged. The cloud clauses and top\_u/top\_b are unchanged, and the small increment of b1 preserves R022's proved b² bound.

All three HC interlacing gaps b0\textless e0\textless b1(s)\textless e1(s) remain positive, and the last positive/negative HC comparison still uses b0. The original centred comparisons either stay at ell, increase, or stay at least their previous lower bound; in particular the old e0 block is separated from e1(s) by 2ell+q\_s. At stage0 the opposite-sign gap between −b0 and −b1 increases. At stages1,2 the literal Dpar rule still gives no negative-sign entries. Thus the checked NCGaps condition is preserved.

The stage0 and1 compact spreads stay unchanged, and the stage2 spread increases by q\_s\textless=ell\textless2\^{}36. R022's elementary bounds still give every spread\textless2A, so DO040's same T\textgreater16A supplies the numerical Sizes inequalities. Positivity and the fixed eta/cloud equalities are unchanged. Finally gcd(a+1,28)=1 gives the displayed progression. Its next member violates NB12; the exact excesses are784 and2156. QED.

The two selected theta partitions become \[
 \nu_0(s)=\nu_0,\qquad
 \nu_1(s)=(\nu_0,0^d)+(\Delta+q_s)\mathbf1_b.             \tag{NB13}
\] Thus D0 is fixed, and D1 increases by fbq\_s. The digit lengths in K041 stay unchanged throughout the interval, as checked at its two endpoints. Every actual algebraic/gauge argument of K041 is unchanged, so NB7 holds with this new D1. The actual compact highest ties remain NB13, and the global stage ratio now has angular type4ell+2q\_s, which is nonzero. Hence K041's full-theta injectivity proof transports to the entire interval, with the new hook, dimension, and creation exponents retained.

If LIFT28's construction uses an additional descent power q\_can, one must intersect NB11 with its stipulated progression. Neatness alone is not used to set q\_can=1. The finite interval above is an upper allowance, not a proof of interpolation for any member.

\subsection{5. The remaining effective interpolation problem}\label{the-remaining-effective-interpolation-problem}

At the actual translated packet Z, NB6 removes every point stabilizer. Let \(\mathcal F_s\) be LIFT28/LIFT29's subcanonical automorphic sheaf, pushed to the minimal compactification, with its actual finite character. Let Z\_N impose the full jets through degree N\_tau at each ball factor at every packet point. The restriction kernel \(\mathcal K_s\) fits into \[
 0\longrightarrow\mathcal K_s\longrightarrow\mathcal F_s
  \longrightarrow\mathcal F_s|_{Z_{\mathbf N}}\longrightarrow0.
 \tag{NB14}
\] A sufficient explicit target is H¹(K\_s)=0 for one s in NB11's allowed progression. A sharper target is surjectivity only after projection to the selected Pieri quotient. Neither vanishing nor surjectivity has been proved in this interval.

For clarity, with packet cardinality z and minimal fibre rank r0, the unconstrained full-jet target has exact dimension \[
 z\,r_0\prod_\tau {N_\tau+a\choose a}.                    \tag{NB15}
\] The selected homogeneous target instead has dimension z dim W\_lambda. These give necessary dimension comparisons with the actual cusp section space, not estimates of that space. The neatness proof justifies the absence of an orbifold-invariant correction in these dimensions.

The generic Serre argument establishes a threshold for unbounded s; NB11 shows precisely the bounded interval in which that threshold must be controlled here. Neither N\_tau≫s\_max nor the size of the full-jet target alone proves failure: the sheaf's rank and weight also vary with the original profile. A proof of an upper section- dimension bound or a quantitative jet-separation theorem is still required.

I checked Lan's primary automorphic-cohomology vanishing theorem: Lan (2016), Theorem4.1(2), printed p.18, in the corrected author PDF \url{https://www.kwlan.org/articles/van-coh-aut.pdf}. It treats subcanonical automorphic bundles under weight conditions. The point-jet kernel in NB14 is an additional ideal modification and is not itself one of those homogeneous bundles. Thus that theorem alone does not give H¹(K\_s)=0. No such implication is used.

\subsection{6. Geometric interpretation and scope}\label{geometric-interpretation-and-scope}

The neatness result concerns the arithmetic quotient of the fixed unitary group; it changes neither the CM abelian scheme nor its Poincare biextension. NB7 is an improved coherent-module inclusion over Spec(R) for the same actual Schur coefficient map. NB8--NB15 concern principal-parts sheaves and coherent interpolation on the actual bottom Shimura variety, with its fixed finite character.

The inherited framework locators are EGA IV.4,16.11.2--3(ii) for jets/principal parts and EGA III.1,2.2.1(ii), printed p.100, for the qualitative coherent vanishing used by LIFT28; SGA6 VI3.3 for the lambda-ring description of Schur functors. Picard and biextension objects retain FGA232 Theorem3.1, pp148--149; FGA236 Corollary1.7,p229 and2.3,p232; SGA7I VII2.7/2.9, pp169--174 andVIII7.4(b),pp310--311. The carried formal-group inputs remain DP\_MULT's cote4/cote54,scans5--7 dictionary. No new ell-adic weight theorem, motivic purity statement, or Kummer RH assertion follows from these level and jet calculations.

\subsection{Check scope}\label{check-scope}

The companion script checks the exact source level anchors, field degrees and G-prime inequalities, small characteristic-polynomial and cyclotomic identities used in NB1, all displayed jet formulas, the complete interval endpoints, and preservation of all changed R022 comparisons. Endpoint and local-bit checks support the affine/monotone all-word and all-twist proofs; they do not enumerate all words or every integer in the twist interval.

The mathematical proof of neatness and its actual level application are checked separately from those finite diagnostics. No arithmetic class enumeration, cusp dimension calculation, successful jet interpolation, or Lean build is claimed.

\clearpage

\section{AM-DO-041: the actual vertical Q kernels survive both local theta quotients}\label{am-do-041-the-actual-vertical-q-kernels-survive-both-local-theta-quotients}

Author: GPT-6 Astra Pro, ALP\_MACHINE /root/definite\_oscillator. 22 November 2026 UTC. Inherited reasoning effort is not separately exposed.

For the two upward steps Y0→Y1→X′ at m=2, the unchanged prescribed Q functions have a nonzero composite image on every initial spherical unitary constituent. In the explicit whole-Schwartz quotient maps below, evaluation at the identity in both induced models is \[
 \boxed{\quad
 \Lambda_Q v=\binom{r+t}{r}_q^{-1}
       \prod_{i=1}^{t}(1-\theta(\varpi)q^{-i})^{-1}v,
 \qquad \theta=\mu_1/\mu_0.\quad}                 \tag{DO41.1}
\] The input v is the actual K\_r-fixed vector. No Y summation, auxiliary translation, Gauss projector, or replacement Schwartz function occurs. The two functions are the conjugates of the literal IsStepVec functions. All maximal compact groups have Haar volume1. Comparison with a separately pinned global theta map remains a nonzero scalar comparison; this local theorem does not compute a terminal period or prove a global vertical image assertion.

\subsection{1. Literal source convention and local hypotheses}\label{literal-source-convention-and-local-hypotheses}

All locators refer to the recovered source at commit 7cbf7a0c32fca1badf05931879a37c3f8a2b03b4.

\begin{longtable}[]{@{}
  >{\raggedright\arraybackslash}p{(\columnwidth - 4\tabcolsep) * \real{0.3000}}
  >{\raggedleft\arraybackslash}p{(\columnwidth - 4\tabcolsep) * \real{0.4000}}
  >{\raggedright\arraybackslash}p{(\columnwidth - 4\tabcolsep) * \real{0.3000}}@{}}
\toprule\noalign{}
\begin{minipage}[b]{\linewidth}\raggedright
Source
\end{minipage} & \begin{minipage}[b]{\linewidth}\raggedleft
Lines
\end{minipage} & \begin{minipage}[b]{\linewidth}\raggedright
Exact content used
\end{minipage} \\
\midrule\noalign{}
\endhead
\bottomrule\noalign{}
\endlastfoot
IndefiniteReal/DatumVec.lean & 22--25,32--46 & frame matrix, integrality, top square block, and LowZero \\
IndefiniteReal/DatumVec.lean & 64--70,83--98 & omegaQ on LowZero support; IsStepVec has exponent +1 \\
\nolinkurl{IndefiniteReal/LinksChar.lean} & 255--269 & actual character pairs and the two vertical IsStepVec data \\
\nolinkurl{IndefiniteReal/OrthogonalComplementCharacterisation.lean} & 13--18 & tuple mu·xi and odd/even chrStep \\
\nolinkurl{Development/AbelianVarietyTheory/Definitions/LatticeFrames/ChainParSucc.lean} & 132--133 & Step.up integrates the conjugate theta kernel \\
Divisibility/Tuples.lean & 58--67,112--129 & split absolutely degree-one Q places and tame tuple ramification \\
\end{longtable}

The four abbreviated IndefiniteReal paths and Divisibility path start with Formalism/LocalToGlobal. Conjugation in Step.up inverts both unitary splitting labels and the tame raw character. It does not replace IsStepVec by the separately defined adjoint IsThPVec.

Let F be a characteristic-zero nonarchimedean field, O its valuation ring, varpi a uniformizer and q its residue cardinality. In the actual Q set, F=Q\_q: DegOne requires prime absolute norm and requires that rational prime not belong to the square of the prime ideal. Thus both absolute residue degree and ramification index equal1, as proved from the literal clauses in DO037 §1. Absolutely ramified places of the enlarged hosts are not silently admitted. The local calculation itself works over general F with these uniformizer and measure conventions.

Set \[
 r=a+1,\quad t=d-1,\quad b=r+t,\quad N=r+2t,\quad z_*=(t+1)/2.
                                                        \tag{DO41.2}
\] Assume t\textgreater=1 and z\_\emph{\textgreater(r−1)/2. The actual ranks are \[
\begin{array}{c|rrrrrr}
 &a&r&t&b&N&z_*\\ \hline
 J^{14}&42&43&1759&1802&3561&880\\
 J^{56}&168&169&1633&1802&3435&817 .
\end{array}                                             \tag{DO41.3}
\] The holomorphy margins z\_}−(r−1)/2 are859 and733.

Let mu0,mu1,eta be unitary unramified characters of F\^{}times, and xi a unitary tame character with the prescribed residue restriction. These are the actual vertical data at Q; eta and mu\_j are unramified there. An arbitrary value of xi on varpi is allowed. Conjugating the source odd/even pairs gives, with target label alpha and source label beta, \[
 (\alpha_0,\beta_0)=((\mu_0\xi)^{-1}\eta,(\mu_0\xi)^{-1}),\qquad
 (\alpha_1,\beta_1)=((\mu_1\xi)^{-1},(\mu_1\xi)^{-1}\eta).
                                                        \tag{DO41.4}
\] The labels mu0,mu1 here name the successive two stage characters; they are not an assertion that the first source chrStep index is zero. The honest split Hom action on M\_(m,n) is \[
 \omega_{\alpha,\beta}(g,h)\phi(X)=
 \alpha(\det h)\beta(\det g)^{-1}
 |\det g|^{m/2}|\det h|^{-n/2}\phi(h^{-1}Xg).       \tag{DO41.5}
\] This is the determinant-root convention checked in DO020/DO028: for a split norm-one determinant(c,c\^{}\{-1\}), choose root(c,1). Evaluating the splitting character at the full norm-one determinant would generally square it and would be incorrect. The fixed split vacuum and common honest realization are those carried in DO020/DO022. No Fourier transform or unnamed Weil phase enters this linear action.

Normalize additive measure on each matrix O-lattice to1 and vol(K\_n)=1, where K\_n=GL\_n(O). Write \[
 \kappa_n=\prod_{i=1}^n(1-q^{-i}),\qquad
 dg=\kappa_n^{-1}|\det g|^{-n}\,dg_{\rm add},\qquad
 \binom mr_q=q^{r(m-r)}\frac{\kappa_m}{\kappa_r\kappa_{m-r}}.
                                                        \tag{DO41.6}
\] The Gaussian binomial is the index in K\_m of the reduction parabolic with blocks(r,m−r).

\subsection{2. Actual whole-Schwartz quotients with unequal labels}\label{actual-whole-schwartz-quotients-with-unequal-labels}

The all-irreducibles Godement--Jacquet input and finite-dimensional continuation are the checked LIFT24(MT10--16)/LIFT27(MR8--14) inputs. The required unequal-label version is as follows. For n→m=n+s, put X\_n=(I\_n,0)\^{}transpose and initially in a right half-plane set \[
 I_{\beta,z}(\phi,v)=\int_{GL_n(F)}
 \phi(X_ng)|\det g|^{z+(n-1)/2}\beta(\det g)^{-1}\pi(g)v\,dg.
                                                        \tag{DO41.7}
\] If L(z,pi beta\^{}\{-1\}) is finite at z\_s=(s+1)/2, this has a holomorphic smooth-vector value there for every Schwartz function on the entire matrix space, including functions meeting the rank-deficient locus.

Indeed g↦phi(X\_ng) extends to a Schwartz function on M\_n.~Every smooth dual coordinate is L(z,pi beta\^{}\{-1\}) times a Laurent polynomial in q\^{}z by Jiang and Luo(2023), Theorem2.4(1)--(2), pp.313--314. Choose a small compact open J killing beta and fixing this matrix restriction under left multiplication. Change of variables makes the integral J-fixed. Admissibility makes pi\^{}J finite-dimensional; continuing its smooth dual coordinates gives an actual vector, independently of J by the identity theorem. Thus extension from the open-rank stratum has been proved, not assumed.

Define \[
 T_z(\phi,v)(h)=I_{\beta,z}(\omega_{\alpha,\beta}(1,h)\phi,v).
                                                        \tag{DO41.8}
\] Changes g↦gg0 and g↦Ag give respectively \[
 I_{\beta,z}(\omega(g_0,1)\phi,\pi(g_0)v)
       =|\det g_0|^{z_s-z}I_{\beta,z}(\phi,v),       \tag{DO41.9}
\] and, for p=(A,B;0,D) in P\_(n,s), \[
 T_z(\phi,v)(p)=
 (\alpha\beta^{-1})(\det A)\alpha(\det D)
 |\det A|^{z-1/2}|\det D|^{-n/2}\pi(A)I_{\beta,z}(\phi,v).
                                                        \tag{DO41.10}
\] They hold first by convergent integration, then by continuation. At z\_s they give a balanced target-equivariant map \[
 q_{\alpha,\beta}:
 (\mathcal S(M_{m,n})\otimes\pi)_{GL_n}
 \longrightarrow(\pi\otimes\alpha\beta^{-1})
                      \times(\alpha\circ\det_s),          \tag{DO41.11}
\] with normalized induction. For unitary factors the target is irreducible by Tadić(1986), TheoremD, pp.338--339. A nonzero map is consequently the actual small theta quotient up to one nonzero scalar, by Mínguez(2008), Theorem1, p.718. That comparison scalar is not set equal to a global arithmetic pin.

Let pi0 be any irreducible unitary spherical representation of GL\_r and v in pi0\^{}K\_r. The two actual targets are \[
 \pi_1=(\pi_0\eta)\times((\mu_0\xi)^{-1}\eta)\circ\det_t,
 \qquad
 \pi_2=\pi_0\times(\mu_0\xi)^{-1}\circ\det_t
                    \times(\mu_1\xi)^{-1}\circ\det_t.      \tag{DO41.12}
\] Both whole-Schwartz maps are holomorphic at z\_\emph{. The first uses the carried unitary bound Re(z)\textgreater(r−1)/2 from LIFT27(MR11). The second uses normalized-induction multiplicativity: \[
 L(z,\pi_1\beta_1^{-1})=
 L(z,\pi_0\mu_1\xi)
 \prod_{j=0}^{t-1}L(z+(t-1)/2-j,\theta),\quad
                         \theta=\mu_1/\mu_0.      \tag{DO41.13}
\] At z\_} the rank-one arguments are1,\ldots,t, and their unitary factors are finite and nonzero. The initial factor satisfies the same large initial-rank bound. No genericity of pi1 or open-rank boundary vanishing is required.

\subsection{3. Prescribed functions and their first compact image}\label{prescribed-functions-and-their-first-compact-image}

For an integral rectangular matrix write A\_r for its top-left r-square block. LowZero\_r means that columns1,\ldots,r below row r vanish modulo varpi. The actual conjugate raw functions are \[
 f_0(X)=\mathbf1_{M_{b,r}(O)}(X)
 \mathbf1_{\mathrm{LowZero}_r(X)}
 \mathbf1_{\det\bar A_r\ne0}\,\xi(\det A_r)^{-1},    \tag{DO41.14}
\] \[
 f_1(Z)=\mathbf1_{M_{N,b}(O)}(Z)
 \mathbf1_{\mathrm{LowZero}_r(Z)}
 \mathbf1_{\det\bar A_r\ne0}\,\xi(\det A_r)^{-1}.    \tag{DO41.15}
\] In f1, columns after r are arbitrary integral columns. Its LowZero rank is the unchanged initial Levi rank in datYt, not b. At Q the omegaG factor is absent by disjointness. The IsStepVec half-position realization relative to the fixed vacuum is exactly these functions in the standard split chart.

Put F0=q\_(alpha0,beta0)(f0,v). At the identity, f0(X\_rg) forces g in K\_r. Its xi(det g)\^{}\{-1\} cancels beta0(det g)\^{}\{-1\}, and pi0(g)v=v. Thus \[
                              F_0(1)=v.          \tag{DO41.16}
\] If J\_(r,t) is the reduction preimage of P\_(r,t)(F\_q) in K\_b, its full compact restriction and normalized-induction norm are \[
 F_0(k)=
 \begin{cases}
 \xi(\det\bar D)^{-1}v,&k\in J_{(r,t)},\\
 0,&k\notin J_{(r,t)},
 \end{cases}
 \qquad \|F_0\|^2=\binom br_q^{-1}\|v\|^2.         \tag{DO41.17}
\] Here D is the lower-right t-square residue block. To check this, the first r columns of k\^{}\{-1\}X\_rg can have the required residue span only when k preserves the first r-space and g belongs to K\_r. Its top block A contributes xi(det A) from the raw factor. The target splitting factor xi(det k)\^{}\{-1\} leaves xi(det D)\^{}\{-1\}, while the g factors cancel as before. The support has the reciprocal parabolic index as Haar measure. In particular the first actual quotient is nonzero.

\subsection{4. Exact two-edge integral and all character cancellations}\label{exact-two-edge-integral-and-all-character-cancellations}

Put F2=q\_(alpha1,beta1)(f1,F0). Evaluate at1 in the outer induced model, then at1 in pi1's induced model; denote the latter map by ev1. Deform only the second quotient by z, and write w=z−z\_*. Initially on a right half-plane the nested integral is \[
\begin{split}
 \operatorname{ev}_1T_z(f_1,F_0)(1)=
 \int_{GL_b}\int_{GL_r}&f_1(X_bh)f_0(h^{-1}X_rg)
 |\det h|^{t+w}\theta(\det h)\\
 &\cdot|\det g|^{b/2}\beta_0(\det g)^{-1}
                         \pi_0(g)v\,dg\,dh .
\end{split}                                             \tag{DO41.18}
\] The ratio is alpha0 beta1\^{}\{-1\}=mu1/mu0, and the exponent is z+(b−1)/2−r/2=t+w. These identities fix both the character direction and the quotient point.

Evaluation ev1 is not assumed to be a bounded Hilbert functional. Choose a compact open J fixing the left matrix restriction of f1 and killing beta1. The outer integral lies in finite-dimensional pi1\^{}J. Apply ev1 there, equivalently first applying the compact projector e\_J. In the convergent region h↦j\^{}\{-1\}h removes this projector, since the matrix restriction, determinant measure and beta1 are invariant. The calculation below proves absolute convergence of the resulting evaluated double integral for Re(w)\textgreater−1, including w=0. Its continuation therefore agrees with ev1 of the whole-Schwartz vector at z\_*.

In the support of f1(X\_bh), write \[
 h=\begin{pmatrix}U&V\\\varpi W&D\end{pmatrix},\quad
 U\in GL_r(O),\quad V,W,D\text{ integral},\quad
 S=D-\varpi WU^{-1}V\in M_t(O),\quad\det S\ne0.    \tag{DO41.19}
\] The additional f0 support condition is exactly \[
                  g\in K_r,\qquad
                  C=S^{-1}WU^{-1}\in M_{t,r}(O). \tag{DO41.20}
\] For the forward implication, Y=h\^{}\{-1\}X\_rg is integral with invertible top block and lower block in varpi M\_(t,r). Then hY=X\_rg shows g integral, and its top reduction is U times an invertible block. Thus g lies in K\_r. The lower block of Y is −varpi S\^{}\{-1\}WU\^{}\{-1\}g, so C is integral. Conversely, these two conditions give \[
 Y_{\rm low}=-\varpi Cg,\qquad
 Y_{\rm top}=U^{-1}g+\varpi U^{-1}VCg,              \tag{DO41.21}
\] which satisfy precisely the required residue conditions.

On this support the tame characters cancel: \[
 \xi(\det U)^{-1}
 \xi(\det(U^{-1}g))^{-1}\beta_0(\det g)^{-1}=1,
 \qquad\pi_0(g)v=v.                               \tag{DO41.22}
\] Also det h=det U det S and theta(det U)=1. The change(W,D)→(C,S), with W=SCU and D=S+varpi SCV, has additive Jacobian \textbar det S\textbar\^{}r: first D↦S has Jacobian1, then W=SCU scales the r columns by S; det U is a unit. The lower block varpi W contributes q\^{}\{-rt\}. Haar measure contributes kappa\_b\^{}\{-1\}\textbar det S\textbar\^{}\{-b\}; the weight gives \textbar det S\textbar\^{}\{t+w\}. As b=r+t all exponents except w cancel. Integrals over g,U,V,C give1,kappa\_r,1,1. Therefore \[
 \operatorname{ev}_1T_z(f_1,F_0)(1)=
 \frac{q^{-rt}\kappa_r}{\kappa_b}
 \int_{M_t(O)}\theta(\det S)|\det S|^w\,dS\ v.     \tag{DO41.23}
\] The singular matrices have additive measure zero.

\subsection{5. Determinant recursion and exact compact norms}\label{determinant-recursion-and-exact-compact-norms}

In the stated normalization, \[
 \int_{M_t(O)}\theta(\det S)|\det S|^w\,dS
 =\prod_{i=1}^t\frac{1-q^{-i}}{1-\theta(\varpi)q^{-i-w}},
 \qquad\Re(w)>-1.                                \tag{DO41.24}
\] Here is an elementary proof of the required determinant integral. Partition the first column by the common valuation j of its entries; that set has measure(1−q\^{}\{-t\})q\^{}\{-tj\}. GL\_t(O) is transitive on primitive columns. Compact averaging can take the column to varpi\^{}j e1 and preserves the distribution of all other columns. Expansion along it leaves a free top row, an independent(t−1)-square matrix, and factor theta(varpi)\^{}j q\^{}\{-jw\}. Compact determinants are units, on which theta is trivial. Summing j gives the recursion with factor (1−q\^{}\{-t\})/(1−theta(varpi)q\^{}\{-t-w\}). Induction from t=0 proves the identity. The same argument with absolute values proves convergence and locally uniform majorants for Re(w)\textgreater−1.

Equations(DO41.23--24) at w=0 give(DO41.1), with no missing additive phase. Every Euler factor is nonzero since \textbar theta(varpi)\textbar=1 and q\^{}\{-i\}\textless1. Thus the second quotient and the composite are nonzero.

There is an exact compact norm refinement. Set F20(k)=ev1 F2(k) for k in K\_N. Then \[
 F_{20}(k)=
 \begin{cases}
 \Lambda_Q\xi(\det\bar D_{2t})^{-1}v,&k\in J_{(r,2t)},\\
 0,&k\notin J_{(r,2t)}.
 \end{cases}                                      \tag{DO41.25}
\] For support, if both raw functions in the analogue of(DO41.18) at k are nonzero, Y1=k\^{}\{-1\}X\_bh integral implies h integral. Its first r columns have rank r modulo varpi. If Y0=h\^{}\{-1\}X\_rg lies in f0's support, hY0=X\_rg first makes g integral and then shows g has rank r modulo varpi: Y0's top block is invertible and its lower block vanishes. It also shows that the first r columns of h span the first r-space in F\_q\^{}b. Thus those of X\_bh span the first r-space in F\_q\^{}N. But X\_bh=kY1 says the same space is k times that first r-space. Hence k preserves it. Conversely, for k in J\_(r,2t), the raw f1 transforms by xi(det A\_r), while alpha1(det k)=xi(det k)\^{}\{-1\}. Their product is xi(det D\_(2t))\^{}\{-1\}. This proves the formula first under convergent integration and then at the holomorphic value.

Only the nested scalar F20 is asserted to have this support; the pi1-valued F2 need not vanish outside J\_(r,2t). With unitary normalized induction norms, one obtains \[
       \boxed{\ \|F_2\|^2=
       |\Lambda_Q|^2\binom Nr_q^{-1}\|v\|^2.\ }    \tag{DO41.26}
\] Indeed its norm is the integral over K\_N×K\_b of \textbar\textbar F2(k)(h)\textbar\textbar². By(DO41.10), left multiplication by diag(h,I\_t) turns this, up to a unit-modulus character, into \textbar\textbar F20(diag(h,I\_t)k)\textbar\textbar². Left Haar invariance on K\_N removes the h integral. Equation(DO41.25) then gives exactly the reciprocal parabolic index. No compact Schur projection factor is suppressed. The raw additive-Schwartz norms, which are different quantities, are \[
             \|f_0\|^2=\kappa_rq^{-rt},\qquad
             \|f_1\|^2=\kappa_rq^{-2rt}.          \tag{DO41.27}
\] They follow by counting the prescribed residue blocks.

\subsection{6. Retained size, valuation and global scope}\label{retained-size-valuation-and-global-scope}

The exact normalization is(DO41.1). For a coarse uniform bound, kappa\_n\textgreater1/4 for every q\textgreater=2: at q=2 its first two factors give3/8, while the remaining product is at least1−sum\_(i\textgreater=3)2\^{}\{-i\}=3/4. Also product\_(i\textgreater=1)(1+q\^{}\{-i\})\textless exp(1/(q−1))\textless=e\textless3. Consequently \[
                    q^{-rt}/48<|\Lambda_Q|<4q^{-rt}.
                                                        \tag{DO41.28}
\] The large Haar cost is therefore retained.

At a coefficient-field place above11, suppose theta(varpi) is an algebraic11-unit, as it is for the actual angular characters away11 by the principal-ideal unit calculation of K039/K040. The actual Q set excludes11, so q is an11-unit. Then \[
 v_{11}(\Lambda_Q)=
 -v_{11}\binom br_q
 -\sum_{i=1}^t v_{11}(1-\theta(\varpi)q^{-i})\le0. \tag{DO41.29}
\] Both terms are retained. The integrated quotient maps are not asserted to be saturated integral lattice maps. Changing their normalization would change precisely the displayed scalar and quotient norms.

This applies at every actual tuple Q place on either enlarged host, with the unchanged vertical raw functions and the initial spherical line. It is injective on that line and supplies a specified nonzero evaluation in the final actual local theta representation. An empty Q set contributes the empty product1. It settles the Q component in the adjoint vertical tower; the other finite factors, real orientation, global Rallis map, class identifications and tSpin/RowLevelMatch remain separate. Neither this result nor full global theta nonvanishing alone asserts a nonzero literal brC/tilQ/PtilFF value.

\subsection{7. Sources and proof-check scope}\label{sources-and-proof-check-scope}

The source clauses in §1 were read directly. These published inputs are used at their already checked scope in LIFT24/LIFT27:

\begin{itemize}
\tightlist
\item
  Jiang, D., \& Luo, Z. (2023). Certain Fourier operators and their associated Poisson summation formulae on GL1. Pacific Journal of Mathematics,326(2),301--372. §2B, equation(2-10), Theorem2.4(1)--(2), pp.313--314. \url{https://msp.org/pjm/2023/326-2/pjm-v326-n2-p05-p.pdf}
\item
  Tadić, M. (1986). Classification of unitary representations in irreducible representations of general linear group (non-Archimedean case). Annales scientifiques de l'École Normale Supérieure,19(3),335--382. TheoremD and definitions, pp.338--339. \url{https://www.numdam.org/article/ASENS_1986_4_19_3_335_0.pdf}
\item
  Mínguez, A. (2008). Correspondance de Howe explicite: paires duales de type II. Annales scientifiques de l'École Normale Supérieure, 41(5),717--741. Theorem1,p.718; normalization(1.4), pp.721--722. \url{https://www.numdam.org/item/ASENS_2008_4_41_5_717_0.pdf}
\item
  LIFT24(MT10--16), LIFT27(MR8--14): whole-Schwartz continuation, unitary holomorphy, multiplicativity and normalized induction.
\item
  DO020/022/028: actual determinant roots, split vacuum and honest Hom action. DO037 §1: literal Q local field. K039/K040: angular11-units.
\end{itemize}

The new proofs are the unequal-label covariance, unchanged raw compact sections, exact two-edge support and Schur-complement integral, determinant recursion, nested compact support and quotient norms. The companion exact certificate tests these finite identities in small ranks and the actual rank arithmetic. It is not a global theta or terminal-value certificate. Whole-proof readings are recorded separately in the agent check log. No Lean build is claimed.

The exact replay passed at03:09:45UTC:100 Gaussian/Haar identities; 10,716 finite-ring matrices for the determinant recursion;96 rational Schur-complement/support cases, including nonunit determinants and failed reverse-support cases;2,304 tame phase cancellations; all168 GL3(F2) and11,232 GL3(F3) compact targets in the nested-section test; four exhaustive raw-support norm counts; and both actual rank margins. The full script and output are \nolinkurl{definite_oscillator_vertical_q_certificate.py} and its matching JSON.

\clearpage

\section{AM-R-025: the complete standard bottom finite level and its residue action}\label{am-r-025-the-complete-standard-bottom-finite-level-and-its-residue-action}

Author: GPT-6 Astra Pro, ALP\_MACHINE \nolinkurl{/root}, 22 November2026 UTC. Reasoning effort is not separately exposed.

This result supplies a concrete finite realization of the source's rank-zero bottom compact, its character, its complete residue quotient, and the action identity in RowLevelMatchA/B. It uses the actual hosts, G places and fixed generic characters of K031/K032/K033/K040 and the unchanged dat0 level of LIFT29. It does not identify an arbitrary finite translation with tS. In these explicitly chosen standard frames the level-matching translation is1; LIFT29/30's T translation remains in the Bessel embedding.

The source requires a further central-idempotent coefficient identity, RowLevelMatchC. That condition is stated precisely below and is not inferred from a level equality. R023's obstruction to a literal full Host on the current fields is also retained: the result realizes these finite mathematical equations on the mixed hosts, not a completed D7 or RealPack record.

\subsection{1. Groups, lattice and actual residue fields}\label{groups-lattice-and-actual-residue-fields}

Let E/F be either actual K031 or K032 host, let a=42 or168, and set r=a+1. Put

\[
 V=E^r,\qquad J=\operatorname{diag}(I_a,-1),\qquad
 \Lambda=O_E^r,\qquad K=\operatorname{Stab}_{U(J)(\mathbf A_{F,f})}
                                      (\widehat\Lambda).
                                                        \tag{R25.1}
\]

Use the actual selected prime set P\_G, containing one E prime from every conjugate pair over each g0,g1. Both rational primes split completely, so every selected residue field has degree one and its identification with F\_g is the unique unital prime-field isomorphism. Every F place over g0g1 occurs. Write

\[
 \mathcal G=\prod_{w\in P_G}\mathrm{GL}_r(\mathbf F_{g(w)}),
 \quad \operatorname{red}:K\longrightarrow\mathcal G,
 \quad K^1=\ker\operatorname{red}.
                                                        \tag{R25.2}
\]

At a split place the chosen E component identifies U(J)(F\_v) with GL\_r(F\_v); its other component is determined by the inverse transpose and J. In particular reduction at all chosen components determines reduction at the conjugate components. At places away from G there is no congruence restriction in K\^{}1.

The map red is a surjective homomorphism, and

\[
 K^1=\{k\in K:(k-1)\widehat\Lambda
                           \subset(g_0g_1)\widehat\Lambda\}.
                                                        \tag{R25.3}
\]

\textbf{Proof.} Matrix reduction is a homomorphism. At each selected place, any matrix in GL\_r(F\_g) lifts to an integral invertible matrix: lift its entries arbitrarily; its determinant remains a unit. Supply the other split component by the exact unitary relation. Choose identity at all other places. These finitely many independent lifts give an adelic element of K with the prescribed full residue tuple. The kernel is congruent to identity at all selected components. The unitary relation implies the same at their conjugates. Since the rational primes split completely, this is exactly congruence modulo g0g1 at every prime where that ideal is nonunit; away from those primes the right side of (R25.3) is automatic. QED.

This also verifies the literal source GoodAt condition. In RowDict.lean49--69 use the standard global basis e1,\ldots,er. Its reduction is a basis at each degree-one prime, because every O\_E residue has an integer lift. The difference between k times a basis vector and its residue-column lift lies in the completed w-lattice: at w this is precisely the reduction identity, and at every other prime it is an integral lattice vector. Thus red is the explicit witness in RowLevelMatch.IsRedHom, including its surjectivity and exact kernel. For the literal type of that source predicate, extend this map from K to the whole finite adelic group by the identity element of the residue group outside K. Every homomorphism, reduction, surjectivity and kernel clause in IsRedHom concerns the restriction to K, so this extension changes none of the preceding assertions. No choice of an unrecorded conjugacy class replaces this quotient.

\subsection{2. A Hida residue frame with the actual generic character}\label{a-hida-residue-frame-with-the-actual-generic-character}

Let U\_a(F\_g) be the upper unitriangular group. For the fixed generic character at w write

\[
 \Psi_w(u)=\psi_w\left(\sum_{i=1}^{a-1}u_{i,i+1}\right),
                                                        \tag{R25.4}
\]

where psi\_w is exactly the nontrivial additive character in LIFT29's PsharpSym specialization. Set

\[
 K_B=\operatorname{red}^{-1}
          \prod_{w\in P_G}\operatorname{diag}(U_a(F_{g(w)}),1),
 \qquad\chi_B(k)=\prod_{w\in P_G}\Psi_w(\operatorname{red}_w k).
                                                        \tag{R25.5}
\]

This is the complete dat0 bottom group and right character in standard betaY0=I frames, by LIFT29(BT1--4), not only its intersection with H. Its subgroup at every non-G place, including T and Q, is the complete standard lattice stabilizer. The initial horizontal character is chi\_B\^{}\{-1\} on H, as in LIFT28/29.

Here is an actual HidaSetting finite residue frame producing precisely this group and character. Since rS=0 there are no e0 or f0 vectors. Set z0=e\_r. At every G split place let ew be the integral idempotent of the chosen E component; use zero away from G. Its conjugate is the opposite idempotent, and ew+bar(ew)=1 at G. Choose units c\_i at the selected components, 1\textless=i\textless=a, and put

\[
 b_i=ew\,c_i e_i,\qquad b_i^d=\overline{ew}\,\overline{c_i}^{\,-1}e_i.
                                                        \tag{R25.6}
\]

Here c\_i denotes an adelic unit at the chosen component, with its opposite-component conjugate used in b\_i\^{}d; only these components enter the displayed vectors. Equivalently, at a split pair the two vectors are (c\_i e\_i,0) and (0,c\_i\^{}\{-1\}e\_i). Hence h(b\_i\^{}d,b\_j)=delta\_ij ew, and all pairings with z0 vanish. The vectors z0,b\_i,b\_i\^{}d span the lattice modulo g0g1. They are integral. They satisfy exactly the rank-zero finite orthogonality, duality, spanning and idempotent clauses in HidaSetting.lean233--266.

The units c\_i are chosen to match the existing psi\_w, rather than change it. The Hida additive character evaluated at ew*x/(g0g1) induces a nontrivial character psi\^{}H\_w of F\_g. Nontriviality follows from complete splitting, the standard finite additive character of conductor O at g, and the fact that the other rational G prime is a unit there. On a prime field every additive character is a unique scalar translate of any fixed nontrivial one. Choose the unique h\_w in F\_g\^{}times such that

\[
                  \psi^H_w(h_wx)=\overline{\psi_w(x)}.
                                                        \tag{R25.7}
\]

Choose c1=1 and c\_(i+1)/c\_i congruent to h\_w at this component. Unit lifts exist. The choices are independent at the finitely many G places. They affect neither the residue flag nor K\_B.

For x in K, the rank-zero Hida congruences say

\[
 (x-1)e_r\in(g_0g_1)\widehat\Lambda,\qquad
 (x-1)b_i\in\langle b_1,\ldots,b_{i-1}\rangle
                                     +(g_0g_1)\widehat\Lambda.
                                                        \tag{R25.8}
\]

At the chosen split component they say that column i has diagonal1 and only earlier coordinates above it, that its last coordinate is0, and that the last column is e\_r. Therefore they are exactly the block-diagonal reduction in (R25.5). At the conjugate component the conditions on b\_i have zero input and the condition on e\_r follows from the inverse-transpose unitary relation. At non-G places the modulus is a unit, so these conditions impose nothing beyond K. Consequently the full source K2set is K\_B.

Finally the root coordinate h(b\_i\^{}d,x b\_(i+1)) at w equals (c\_(i+1)/c\_i)x\_(i,i+1) in the chosen component. By (R25.7), the literal Hida chi2f is

\[
                    \chi_2(k)=\overline{\chi_B(k)}
                                  \qquad(k\in K_B).
                                                        \tag{R25.9}
\]

This includes the g0g1 denominator and every additive-character phase. There is no assumption that the two rational-prime characters have identical unscaled coordinates.

Equations(R25.5),(R25.8),(R25.9) prove complete compact equality and character equality with level translation t=1. They provide both ConjLevel(1) and the rank-zero finite content of LevelMatchA. In particular the negative-line translation tau of LIFT29 is not the level-matching element in this construction.

\subsection{3. The complete residue-group action}\label{the-complete-residue-group-action}

Let \(\mathcal V\) be the space of scalar automorphic functions right invariant under K\^{}1, with the given real type if desired. Every function with right (K\_B,chi\_B) type belongs to this space, because chi\_B is trivial on K\^{}1. Define for gamma in Gcal

\[
             (R_\gamma F)(g)=F(g\widetilde\gamma),
                \quad\operatorname{red}(\widetilde\gamma)=\gamma.
                                                        \tag{R25.10}
\]

This is independent of the chosen lift: two lifts differ by an actual element of K\^{}1, which acts trivially. It gives the usual right-translation representation of the full residue group. In particular, for every k in K, every F in V, and every adelic g,

\[
                       F(gk)=R_{\operatorname{red}k}F(g).
                                                        \tag{R25.11}
\]

The source actL uses dat0(triv), not dat0(Psharp), for this full residue action. In the present rank-r bottom, dat0 has n=r,s=0,Q=empty. Its full flag parahoric is K and its flag character is1: the extra block-character product in DefGroup.lean703--707 is empty, as is its Q product. LevelP.triv gives precisely K\^{}1. Thus every lift in (R25.10) satisfies the literal IsLevelLift, and its levChi(flagL) normalizing factor is exactly1. \nolinkurl{SubgroupCharExtensionOneApply.lean95–110} therefore identifies actG0, and hence LinksDefs.lean183--193 actL, with (R25.10).

Use the canonical prime-field identifications to identify the source G0g with Gcal, and its indexing group G' with Gcal. Then sigma is that canonical isomorphism, and (R25.11) is exactly the LevelMatchB action identity, with t=1 and the red map of section1. It holds for every function in the relevant multiplicity space and every real vector, not merely on the initially selected bottom cusp form.

The quotient has the exact finite index

\[
 [K:K^1]=|\mathcal G|,
 \qquad
 [K:K_B]=\prod_{w\in P_G}
     { |\mathrm{GL}_r(F_{g(w)})|\over
                               g(w)^{a(a-1)/2}}.
                                                        \tag{R25.12}
\]

For any chosen finite Haar these imply vol(K\^{}1)=vol(K)/\textbar Gcal\textbar{} and vol(K\_B)=vol(K)/{[}K:K\_B{]}. The action identity does not omit either of these normalization factors when passed to a convolution operator.

At11 these two indices have the same valuation, because every g(w) is an11-unit. Using \textbar GL\_r(F\_q)\textbar=q\^{}\{r(r-1)/2\} product\_(i=1)\^{}r(q\^{}i−1) gives the exact values

\begin{longtable}[]{@{}lrrr@{}}
\toprule\noalign{}
Host & v11( & GL\_r(F\_g0) & ) \\
\midrule\noalign{}
\endhead
\bottomrule\noalign{}
\endlastfoot
J14 & 4 & 4 & 336 \\
J56 & 36 & 91 & 21,336 \\
\end{longtable}

There are f=42 or168 copies of each rational-prime factor. These are actual finite group-order valuations; a specific central idempotent may have smaller denominators. The table is not a claim that every projector incurs this entire loss or that it has no additional norm cost.

\subsection{4. The remaining spectral identity is a separate condition}\label{the-remaining-spectral-identity-is-a-separate-condition}

Let e0 be the actual finite central idempotent selected by the source spectral construction, and let e\_A be the actual idempotent of the mixed row family. In the coordinates of section3, the literal remaining condition is

\[
               e_A(\sigma\gamma)=\overline{e_0(\gamma)}
                                      \quad(\gamma\in\mathcal G).
                                                        \tag{R25.13}
\]

This is precisely RowLevelMatchC. If the separately required grp\_record identity e\_A(sigma gamma)=e0(gamma) is imposed, (R25.13) is equivalent to coefficientwise reality of e0. RowDict.lean82 names this latter condition IdempotentRealFixed. Centrality and self-adjointness of a convolution idempotent do not by themselves prove coefficientwise reality; complex conjugation also interacts with inversion. We do not silently add such a hypothesis or replace e0 by a different projector.

Even when (R25.13) has been established, projection of LIFT29's detector may erase its compact period. The sufficient actual condition is a nonzero restriction of that period functional to e\_A V; it is not merely nonvanishing of e\_A as a group-algebra element. LIFT31's proposed surjectivity addresses downward-image membership in the full cusp space and, after projection, onto e\_A V. It does not imply this additional period condition.

For reference, the exact e0 is the sum of tensor products of the finite series idempotents eSer over goodTuple: \nolinkurl{LevelLetterCocycles/DataN.lean960–966,996–1007.} Its goodTuple invokes both the prescribed tame characters and their inverses. The complete IsGood condition is \nolinkurl{CyclotomicNorms/StatementT.lean19–29,} including squarefreeness, prime-to-part exclusions, and factor-degree clauses. No choice of a single Steinberg block or an arbitrary generic finite representation identifies this actual projector.

\subsection{5. Compatibility with the explicit negative line}\label{compatibility-with-the-explicit-negative-line}

LIFT30's vector z\_T=B e1+A e\_r satisfies norm(z\_T)=−T and is primitive integral. DO040's congruence gives T congruent1 modulo(g0g1)\^{}2, so z\_T is congruent e\_r modulo g0g1. In the Hida finite frame above we can therefore take x\_u=1 for this congruence: its xu\_z relation holds, and xu\_K/xu\_one hold literally. This is only the stated finite congruence portion of that source record.

The actual compact embedding remains H\_T=gHg\^{}\{-1\}, and its period uses the finite/real identity of LIFT30 with the mandatory real vector transport. The t=1 level witness does not make g or tau equal to1, and does not set the algebraic transport cost to zero. The full HidaSetting also contains the actual flag, global lattice, real representation and spectral data. Those are not constructed solely by the finite congruence calculation here.

\subsection{6. Framework, citations and check scope}\label{framework-citations-and-check-scope}

The result identifies the complete compact subgroups through their integral standard lattice models, followed by their finite etale residue group at the split good primes. It supplies an actual action on principal-level functions, rather than a virtual equality of representations. Its group-scheme framework retains SGA3I4.7.3 and 4.7.3.1; its coefficient Schur functors retain SGA6VI3.3. The fixed CM abelian scheme, Picard and biextension comparison retains FGA232 Theorem3.1 pp148--149, FGA236 Corollary1.7 p229/2.3 p232, SGA7I VII2.7/2.9 pp169--174 andVIII7.4(b) pp310--311. The formal-jet context is EGAIV.4 16.11.2--3(ii) pp54--55 and the established cote4/cote54 crystal dictionary. No new motivic, ell-adic, cycle-nonvanishing or Kummer RH statement follows.

Exact source locators at the pinned Alpöge archive commit above:

\begin{itemize}
\tightlist
\item
  \nolinkurl{HidaSetting.lean117–136,233–276:} full congruence subgroup, generic character, integral residue frame and xu conditions.
\item
  RowDict.lean49--82: GoodAt/RedMat, exact residue interpretation and separately named coefficientwise-reality condition.
\item
  RowLevelMatch.lean18--47: all three finite matching equations.
\item
  \nolinkurl{DefGroup.lean581,608–609,620–707:} principal residue level, rank-zero sharp group, full parahoric and exact character.
\item
  \nolinkurl{SubgroupCharExtensionOneApply.lean32–43,59–73,95–110:} actual dat0/flag levels, normalizing character and residue action.
\item
  \nolinkurl{LinksDefs.lean183–193,427–431,628–637;} LinksChar.lean106--119: actL, dat0, the central idempotent and grp\_record.
\item
  RealPins.lean35--52: the independent group/idempotent and tS pins.
\end{itemize}

Settled mathematical inputs are K031/K032 actual complete-split G places; K033/K040 fixed generic/stage-character choices; LIFT29 rank-zero dat0; LIFT30 primitive negative line and finite/real transport; DO040 prime congruence; and R023 fixed-host limitation. They are used at that stated scope. No Lean execution or archive soundness claim is involved. The proof is the explicit construction and matrix/residue-action calculation above; whole-proof check events are recorded separately.

The companion exact certificate enumerates all11,232 matrices of GL3(F3) and all20,160 matrices of GL4(F2), checks both split components of the complete Hida congruence against the block-diagonal unipotent group, and checks20 exact additive-character phase identities. It also computes the displayed actual group-order valuations by finite integer powers. It does not enumerate the actual high-rank groups, test the spectral projector, or compute an automorphic period.

\clearpage

\section{ALP-LIFT-31: the prescribed vertical adjoint is injective and its downward map is onto}\label{alp-lift-31-the-prescribed-vertical-adjoint-is-injective-and-its-downward-map-is-onto}

Author: /root/lift\_accounting, GPT-6 Astra Pro, ALP\_MACHINE. 22 November 2026 UTC. Inherited reasoning effort is not separately exposed.

\textbf{Author whole-proof review completed: 2026-11-22T03:21:43Z.} This statement concerns the actual two vertical finite Schwartz functions in DatumVec, with their conjugation in IsIoUp. It is not an assertion about freely chosen later Schwartz functions. The starting group is the indefinite U(a,1), the two targets are positive definite, and both rank increments are t. The conclusion is injectivity on the entire fixed-level bottom cuspidal multiplicity space, and hence surjectivity of the actual adjoint downward map onto that space. Projection afterwards gives surjectivity onto its prescribed idempotent summand as well.

The theorem does not identify a particular holomorphic jet detector with that idempotent summand, does not complete the source tS, RowLevelMatch, or arithmetic row pins, and does not give an effective lower singular value or an integral saturation estimate. All global theta, real intertwiner, and Haar normalizations remain actual factors.

\subsection{1. Actual ranks, spaces, characters, and statement}\label{actual-ranks-spaces-characters-and-statement}

Use either actual mixed CM host K031 or K032 and the corresponding literal stage characters. Use the R022/K040 profiles, or the earlier profiles when only their real interlacing and centering conditions are needed. Put

\[
 \begin{array}{c|rrrrrr}
 &a&d&r=a+1&t=d-1&b=r+t&N=b+t\\ \hline
 K031&42&1760&43&1759&1802&3561\\
 K032&168&1634&169&1633&1802&3435 .
 \end{array}                                                     \tag{VA1}
\]

Thus r and t are odd, b is even, N is odd, b\textgreater4r, and b\textless N\textless2b. Let

\[
 G_0=U(I_a\oplus(-1)),\qquad G_1=U(I_b),\qquad G_2=U(I_N).
                                                               \tag{VA2}
\]

The first group has relative F-rank one; the last two are anisotropic over F because they are positive definite at every real embedding. Both vertical increments are t. In particular the standard Rallis argument is evaluated at

\[
                 s_0=t/2,\qquad z_0=s_0+\tfrac12=d/2.           \tag{VA3}
\]

Write the two horizontal stage characters as \(\mu_0,\mu_1\), the tuple character as \(\xi\), and the fixed odd character as \(\eta\). The unbarred vertical characters are, in order,

\[
 ((\mu_0\xi)\eta^{-1},\,\mu_0\xi),\qquad
 (\mu_1\xi,\,(\mu_1\xi)\eta^{-1}).                              \tag{VA4}
\]

The upward adjoint uses the complex-conjugated kernels. In a relative split Hom model its unitary character pairs are consequently

\[
 (\alpha_0,\beta_0)=((\mu_0\xi)^{-1}\eta,\,(\mu_0\xi)^{-1}),
 \quad
 (\alpha_1,\beta_1)=((\mu_1\xi)^{-1},\,(\mu_1\xi)^{-1}\eta).
                                                               \tag{VA5}
\]

All characters in this display are unitary analytic characters. This does not remove their algebraic normalization in an arithmetic model. The relevant relative character is

\[
                            \vartheta=\mu_1/\mu_0.              \tag{VA6}
\]

It is nontrivial. In the actual narrow-cloud profiles its angular weights are \(-4\ell\), up to the fixed simultaneous inverse convention, and its square is nontrivial. It also has the nontrivial literal successive character ratio at G. K040/K041 and LIFT14's checked boundary-value lemma give \(L_E(1,\vartheta)\ne0\), finite.

Use the genuine real types identified in AM-R-015. To avoid confusing the signed parameter with the actual scalar automorphic function, write \(\tau_j=\overline{\sigma_j}\), where \(\sigma_j\) is R015's representation. The upward conjugated oscillator contains \(\tau_j^\vee\boxtimes\tau_{j+1}\). In particular \(\tau_0\) is a holomorphic or antiholomorphic discrete series on U(a,1), with the orientation fixed by R015, and \(\tau_1,\tau_2\) are finite-dimensional unitary representations of compact real groups.

Let \(\mathcal C_j\) be the full cuspidal automorphic multiplicity space of this real type and the finite level specified in section 2. At \(j=1,2\) cuspidality is automatic. At \(j=0\) it is imposed. These spaces are finite dimensional. Let

\[
 D=D_0D_1:\mathcal C_2\longrightarrow\mathcal C_0              \tag{VA7}
\]

be the actual two downward maps, defined by the unbarred source kernels and nonzero real intertwiners. Let U be the composition of the two upward maps using the conjugated kernels and adjoint real intertwiners. Section 10 proves directly that \(U=D^*\) for the actual Petersson multiplicity metrics; this is not assumed from a source predicate.

The conclusion is

\[
 \boxed{\quad\ker U=0,\qquad D(\mathcal C_2)=\mathcal C_0.\quad}  \tag{VA8}
\]

For any idempotent \(e\) on \(\mathcal C_0\), in particular the genuine source e0 after its action has been identified,

\[
                         eD(\mathcal C_2)=e\mathcal C_0.        \tag{VA9}
\]

The proof below uses the specified finite functions at every place. No finite translate, auxiliary principal level, or tuple change is introduced.

\subsection{2. Literal source finite data and upward orientation}\label{literal-source-finite-data-and-upward-orientation}

All source locators are relative to work/research/alpoge\_source, commit 7cbf7a0c32fca1badf05931879a37c3f8a2b03b4.

Formalism/LocalToGlobal/IndefiniteReal/ \nolinkurl{OrthogonalComplementCharacterisation.lean41–55} defines the downward integral as theta times the larger-group input, with no conjugation. Lines78--88 define the upward integral with the conjugated theta kernel. Lines74--76 define the corresponding adjoint real intertwiner. LinksChar.lean253--269 ties the two actual Weil data and step vectors to chrStep and to datYt with LevelP.triv. LevelP.triv is the bottom residue subgroup with trivial character, not the full residue group: \nolinkurl{Divisibility/DefGroup.lean581.}

Let \(K_n=\mathrm{GL}_n(O)\), \(K_n^1=\ker(K_n\to\mathrm{GL}_n(k))\) at a selected degree-one relative split place. At a G place the actual bottom compact is \(K_r^1\). The middle compact and character are

\[
 J_{G,1}=\operatorname{red}^{-1}
   \left\{\begin{pmatrix}I_r&B\\0&D\end{pmatrix}:D\in\mathrm{GL}_t(k)\right\},
 \qquad \lambda_{G,1}=\chi_1(\det D).                           \tag{VA10}
\]

The top compact is the block upper triangular subgroup for \((r,t,t)\), with first diagonal block \(I_r\), and character \(\chi_1(\det D_1)\chi_2(\det D_2)\). At Q, the bottom compact is \(K_r\), the middle compact is the parahoric for the flag \((r,t)\) with character \(\xi(\det D)^{-1}\), and the top compact is the parahoric for \((r,2t)\) with character \(\xi(\det D)^{-1}\). At T and every other finite place all three levels are the actual maximal lattice stabilizers. \textbf{There is no vertical T incidence factor.}

These are DefGroup.lean620--707 and \nolinkurl{IndefiniteReal/SubgroupCharExtensionOneApply.lean32–43,71–73.} The latter uses right transformation \(F(h\kappa)=\lambda(\kappa)F(h)\). It is not the inverse convention of the initial horizontal row. DatumVec.lean49--98 gives the exact finite functions below, including the integral support and the half-idempotent at G and Q.

For \(E_n=\binom{I_n}{0}\), the first conjugated G function is

\[
 f_{G,0}(X)=
  \mathbf1_{\{X\in M_{b,r}(O),\ \overline X=E_r\}}.             \tag{VA11}
\]

The second conjugated G function is supported on integral matrices whose reduction is \[
 X=\begin{pmatrix}I_r&B\\0&D\\0&0\end{pmatrix},
                      \qquad D\in\mathrm{GL}_t(k),
\] and its value is \[
                    \chi_1(\det D)^{-1}\chi_2(\det D).         \tag{VA12}
\]

The determinant in the second factor is equivalently the full invertible top b-block determinant. Keeping that form is useful when checking source and target covariance.

At Q the first function is the integral indicator of a top r-block invertible modulo the uniformizer and all lower rows zero modulo it, times \(\xi(\det X_{\rm top})^{-1}\). The second imposes the same condition on the first r columns of \(X\in M_{N,b}(O)\) and the same top-r determinant character; the last t columns are merely integral. The first and second functions at all other places are the prescribed selfdual lattice vacua.

For clarity, these statements concern the scalar half-Schwartz functions after the actual relative split Heisenberg identification. The prescribed half-idempotents and their measure normalization are already included. Their inverse images are the actual DatumVec vectors; they are not replacement test functions.

The target covariance follows directly from

\[
 [\omega_{\alpha,\beta}(g,h)f](X)
 =\alpha(\det h)\beta(\det g)^{-1}
          |\det g|^{m/2}|\det h|^{-n/2}f(h^{-1}Xg).             \tag{VA13}
\]

For example, at the second G edge, left multiplication of X by the inverse of a target upper block matrix contributes \(\chi_1(\det D_1)\chi_2(\det D_1)^{-1}\) from f, whereas \(\alpha_1(\det h)=\chi_2(\det D_1D_2)\). Their product is the top character stated above. At Q the top-r character contributes \(\xi(\det A)\), and \(\alpha_1(\det h)=\xi(\det A\det D)^{-1}\); the result is \(\xi(\det D)^{-1}\). The first edge is the same calculation with no old t-block. Thus the actual upward maps land in \(\mathcal C_1,\mathcal C_2\), rather than only in a larger unrestricted finite-level space.

LinksChar.lean180--182 identifies the top model HX with the complete AdForms space at precisely datXt, with its actual class metric. LinksDefs.lean649--655 defines the CC spaces and the complete downward chain. No restriction to a selected subspace of the top coefficient space is used in (VA8).

\subsection{3. A cross-character local quotient with no open-orbit extension ambiguity}\label{a-cross-character-local-quotient-with-no-open-orbit-extension-ambiguity}

We record the exact split quotient needed for both edges. Let \(m=n+t>n\), let \(\pi\) be an irreducible unitary representation of \(\mathrm{GL}_n(F_v)\), and let \(\alpha,\beta\) be unitary characters. For a Schwartz function on \(M_{m,n}(F_v)\) put

\[
 I_z(f,v)=\int_{\mathrm{GL}_n(F_v)}
 f(E_ng)|\det g|^{z+(n-1)/2}\beta(\det g)^{-1}\pi(g)v\,dg.
                                                               \tag{VA14}
\]

The whole-Schwartz continuation argument of LIFT24(MT10--16) and LIFT27(MR8--13) applies without change to different target character alpha: alpha does not enter this source integral. For completeness, the standard Godement--Jacquet ideal theorem implies holomorphy whenever \(L(z,\pi\beta^{-1})\) is holomorphic. Pair with every smooth dual functional, use the fractional-ideal statement, and then use the finite-dimensional compact-invariant space containing the vector integral. This produces a vector-valued holomorphic map on the entire Schwartz space. It is not an arbitrarily extended functional on the full-rank open stratum.

At \(z=(t+1)/2\), simultaneous application of \(\omega(g)\) to f and \(\pi(g)\) to v leaves \(I_z\) invariant, by the substitution \(h\mapsto hg\). For \(p=\begin{pmatrix}A&*\\0&D\end{pmatrix}\) in the target parabolic, the function \(T(f,v)(h)=I_z(\omega(1,h)f,v)\) satisfies

\[
 T(f,v)(ph)=\delta_P(p)^{1/2}
    \alpha(\det D)(\alpha/\beta)(\det A)\pi(A)T(f,v)(h).
                                                               \tag{VA15}
\]

Indeed \(p^{-1}E_n=E_nA^{-1}\); changing variable \(g=Au\) gives \(|\det A|^{m/2}|\det p|^{-n/2}=\delta_P(p)^{1/2}\) and the two character factors in (VA15). The target is therefore the normalized induction

\[
             (\pi\otimes(\alpha/\beta)\det_n)
                              \times(\alpha\det_t).            \tag{VA16}
\]

It is irreducible and unitary by the unitary induction theorem of Tadic, in exactly the use checked in LIFT27. A compact full-rank test makes this map nonzero. Minguez's type-II local correspondence, with the same source coinvariant convention, identifies it with the actual small theta quotient, up to a nonzero scalar. The scalar is retained and is not declared to be the source global pin.

For the present two edges this gives

\[
 \pi_1=(\pi_0\otimes\eta\det_r)
                   \times((\mu_0\xi)^{-1}\eta)\det_t,
 \qquad
 \pi_2=\pi_0\times(\mu_0\xi)^{-1}\det_t
                      \times(\mu_1\xi)^{-1}\det_t.              \tag{VA17}
\]

The needed holomorphy holds for every unitary \(\pi_0\). The first point \(z_0=d/2\) exceeds \((r-1)/2\), so the bounded-coefficient determinant integral of LIFT27 proves holomorphy and nonvanishing. For the second edge, multiplicativity gives

\[
 L(z_0,\pi_1\beta_1^{-1})
   =L(z_0,\pi_0\mu_1\xi)
                       \prod_{i=1}^{t}L(i,\vartheta).          \tag{VA18}
\]

Every rank-one factor is finite and nonzero because its character is unitary and \(i\ge1\); a ramified factor is 1. Thus the whole-Schwartz map is holomorphic at the actual point for all initial constituents, without local genericity or temperedness.

\subsection{4. The two actual G tests are injective on the full prescribed type}\label{the-two-actual-g-tests-are-injective-on-the-full-prescribed-type}

Normalize group Haar by \(\operatorname{vol}K_n=1\). For \(v\in\pi_0^{K_r^1}\), (VA11) forces the first orbit variable g to lie in \(K_r^1\), where every splitting character is trivial. Consequently

\[
                    I_0(f_{G,0},v)=|{\rm GL}_r(k)|^{-1}v.       \tag{VA19}
\]

Its entire compact induced-model section is supported on the identity \((r,t)\)-parabolic cell. If \(k=\begin{pmatrix}A&B\\0&D\end{pmatrix}\pmod{\varpi}\), its value is the fixed factor in (VA19), times \(\chi_1(\det D)\pi_0(A)v\). Here the choice of a lift of A is harmless because v is \(K_r^1\)-fixed. Its squared compact norm is exactly

\[
 |{\rm GL}_r(k)|^{-2}
             {b\brack r}_q^{-1}\|v\|^2.                       \tag{VA20}
\]

This uses probability Haar on \(K_b\); the identity parabolic cell has measure the inverse Gaussian binomial coefficient. In particular the first output has precisely (VA10).

For any \(w\in\pi_1^{J_{G,1},\lambda_{G,1}}\), the second raw function forces the orbit variable h to lie in \(J_{G,1}\). On this set \(\beta_1^{-1}=\chi_2^{-1}\) on unit determinants. The \(\chi_2\) factor in (VA12) cancels it; the remaining \(\chi_1^{-1}\) cancels the source type of w. Hence

\[
 I_1(f_{G,1},w)=\operatorname{vol}(J_{G,1})w,\qquad
 \operatorname{vol}(J_{G,1})
       ={q^{rt}|{\rm GL}_t(k)|\over|{\rm GL}_b(k)|}.            \tag{VA21}
\]

Both integrals have compact full-rank support, so no meromorphic regularization is needed for their evaluation. Combining (VA19) and (VA21), and evaluating the first induced section at its identity, gives the exact local left inverse

\[
  {q^{rt}|{\rm GL}_t(k)|\over
                  |{\rm GL}_b(k)|\,|{\rm GL}_r(k)|}\,v.        \tag{VA22}
\]

It detects every nonzero \(K_r^1\)-fixed initial vector. There is no finite Steinberg, generic-unipotent, or cuspidal restriction. All displayed finite denominators belong to the actual local calculation. They are not absorbed into a claim of integral saturation.

There is also an exact second induced norm. Put \[
 c_G={q^{rt}|{\rm GL}_t(k)|\over
                  |{\rm GL}_b(k)|\,|{\rm GL}_r(k)|}.
\] The second section is supported on the identity \((b,t)\)-parabolic cell of \(K_N\). To verify support, integrality of \(k^{-1}E_bh\), followed by its prescribed full residue rank, forces \(h\in K_b\); the raw lower-zero condition then forces k to preserve the first b-plane. On that parabolic cell, covariance(VA15) and(VA21) give a unitary translate of \(\operatorname{vol}(J_{G,1})F_1\), where \(F_1\) is the actual first section. Therefore \[
 \|F_2\|^2
   =|c_G|^2{N\brack b}_q^{-1}{b\brack r}_q^{-1}\|v\|^2.
                                                               \tag{VA22a}
\] The raw additive-Schwartz norms are different quantities: \[
 \|f_{G,0}\|^2=q^{-br},\qquad
 \|f_{G,1}\|^2=q^{-Nb+rt}|{\rm GL}_t(k)|.
                                                               \tag{VA22b}
\] The first is one residue matrix; the second has exactly \(q^{rt}|{\rm GL}_t(k)|\) admitted residue matrices. At a coefficient prime above11, with \(q\ne11\), the exact scalar has valuation \[
 v_{11}(c_G)=v_{11}|{\rm GL}_t(k)|
             -v_{11}|{\rm GL}_b(k)|-v_{11}|{\rm GL}_r(k)|\le0.
                                                               \tag{VA22c}
\] This exhibits an actual averaging cost, not an integral unit.

\subsection{5. The fixed Q pair preserves the same previous image}\label{the-fixed-q-pair-preserves-the-same-previous-image}

Use AM-DO-041's independent exact quotient computation, whose character convention is (VA5). Its main formula is

\[
 \Lambda_Q={b\brack r}_q^{-1}
                \prod_{i=1}^{t}(1-\vartheta(\varpi)q^{-i})^{-1}
       =q^{-tr}{\kappa_r\kappa_t\over\kappa_b}
                \prod_{i=1}^{t}(1-\vartheta(\varpi)q^{-i})^{-1},
 \quad \kappa_n=\prod_{i=1}^{n}(1-q^{-i}).                     \tag{VA23}
\]

The composed local map has an evaluation equal to \(\Lambda_Qv\) for every \(v\in\pi_0^{K_r}\). The first image is the identity-parabolic section with value \(\xi(\det D)^{-1}v\), and squared compact norm \({b\brack r}_q^{-1}\|v\|^2\). Thus this is a calculation on the actual preceding image.

Here is the mechanism that fixes the determinant exponents. Write the middle variable \[
 h=\begin{pmatrix}U&V\\\varpi W&D\end{pmatrix},\quad
 U\in{\rm GL}_r(O),\quad S=D-\varpi WU^{-1}V.
 \tag{VA24}
\] The first support forces the initial variable \(g\in K_r\). Its lower support is equivalent to \(WU^{-1}\in S M_{t,r}(O)\). Write \(W=SCU\). The additive Jacobian is \(|\det S|^r\). The two oscillator determinant factors give \(|\det h|^{(N-r)/2}=|\det h|^t\); multiplicative Haar contributes \(|\det h|^{-b}\). Thus the remaining exponent is \(r+t-b=0\). All \(\xi\) factors cancel. The only remaining character on S is \(\alpha_0\beta_1^{-1}=\vartheta\). Integration over \(M_t(O)\) gives (VA23), with the normalization \(\operatorname{vol}K_n=1\).

In particular \(\Lambda_Q\ne0\). This is an absolutely convergent composed integral; the separate second quotient is the whole-Schwartz continuation of section 3. There is no unproved interchange at a singular determinant locus and no replacement of the vertical function by the horizontal Y-coupled Gauss transform.

The stronger compact calculation in DO041(25)--(27) gives \[
 \|F_{2,Q}\|^2=|\Lambda_Q|^2{N\brack r}_q^{-1}\|v\|^2,
 \qquad
 \|f_{Q,0}\|^2=\kappa_rq^{-rt},\quad
 \|f_{Q,1}\|^2=\kappa_rq^{-2rt}.                               \tag{VA24a}
\] It uses the nested evaluated section supported on the \((r,2t)\)-parabolic cell, then integrates the middle compact argument by left Haar invariance. Evaluation in an induced model is first restricted to a finite-dimensional compact-invariant space before continuation; it is not assumed to be a bounded functional on the whole Hilbert representation. The exact Euler denominators and their 11-valuations remain as in DO041(29).

\subsection{6. Exact basic vectors at odd-rank and absolutely ramified places}\label{exact-basic-vectors-at-odd-rank-and-absolutely-ramified-places}

The vertical kernels are basic at T and at every place outside G and Q. The first source rank is odd, so the even-rank theorem used in LIFT14/16 is not, by itself, a justification here. The following argument supplies both occurrence and the actual test-image assertion.

At a relative split place the actual normalized vacuum is \(\mathbf1_{M_{m,n}(O)}\). The source character is beta. LIFT16(S3--S14) applies in every rank and residue characteristic. Its direct quotient value on a spherical vector is

\[
 I(f^\circ,v^\circ)
       =L((m-n+1)/2,\pi\beta^{-1})v^\circ,                    \tag{VA25}
\]

with \(\operatorname{vol}K_n=1\). It is finite and nonzero at both present steps, by (VA18) and the initial unitary bound. Thus the same basic vector survives at every split place, including all dyadic places of K031/K032. It also covers every place above 29 in these hosts: those places are relatively CM-split, although the real field is absolutely ramified there.

Now let E/F be unramified nonsplit and the residue characteristic odd. The actual lattice is selfdual for the actual additive pairing. Write its trace-zero coefficient as a trace-zero unramified unit times an element of \(F^\times\), and transport that scalar into the additive character. Selfduality makes the resulting character have conductor O. This changes neither the oscillator representation nor its normalized lattice vacuum. It applies even when F itself is absolutely ramified, including the relatively nonsplit places above 7 in the actual hosts. The places above 29 are handled by the split calculation.

Let \(v_A\) be the normalized lattice vacuum. In the doubled oscillator, the canonical partial Fourier transform of \(v_A\otimes\overline{v_A}\) is the characteristic function of the doubled selfdual lattice. Indeed, with 2 a unit the diagonal/antidiagonal change preserves that lattice, and the Fourier transform of a selfdual lattice indicator is itself. Its value at zero is \(\|v_A\|^2=1\). This fixes the scalar and distinguishes the actual vector from an arbitrary spherical Siegel section.

Gong and Grenie (2011), Lemma A.1, printed pp.199--200, computes the Siegel--Weil image of that indicator in all ranks: it is the spherical section whose value on the hyperspecial compact is 1. Their proof checks the partial Fourier operators, the selfdual measure, and the Weil index. The hypotheses are relative unramifiedness, an unramified splitting character, and odd residue characteristic, all satisfied here after the preceding simultaneous transport.

For the unit spherical source vector, the standard local doubling calculation is consequently

\[
 Z_v(s_0,\delta(v_A\otimes\overline{v_A}),v^\circ,v^\circ)
       ={L_v(z_0,\pi\times\beta)\over d_v(s_0,\beta)}\ne0.
                                                               \tag{VA26}
\]

The displayed multiplication sign is GQT's convention. On Satake parameters its numerator in the present Hom orientation is \(\phi_\pi\beta^{-1}\). At a nonsplit unramified place \(\beta(\varpi)=\pm1\), because the uniformizer belongs to F and the character restriction is the appropriate quadratic parity. Thus inverse versus direct beta gives the same unramified factor. At split places (VA14--25) fixes the inverse sign directly. This avoids extending DO028's even-character determinant twist argument to an unspecified odd baseline.

Here the initial unramified unitary bound on the standard Satake roots is \(q_E^{(r-1)/2}\). This is the rank-independent bounded-spherical-function argument used in LIFT14(E18): a unit spherical matrix coefficient is bounded by 1, so its real logarithmic Satake parameter lies in the convex hull of the Weyl orbit of rho. The standard representation of the dual unitary group then gives the displayed bound; the odd central Satake root has modulus1. Thus \(z_0\) exceeds their exponents by 859 or733. For the actual first theta image the newly inserted roots have exponents at most \((t-1)/2\); the second value exceeds those by 1. Every local Euler denominator is therefore nonzero. The normalizing \(d_v(s_0,\beta)\) is a product of finite nonzero rank-one Euler factors. Formula (VA26) is GQT section11.6, not a deduction from a parameter match alone.

The invariant pairing defined by nonzero Z implies that the big local theta module is nonzero. This implication uses the defining coinvariant pairing and does not require the converse assertion of GQT11.6 or its dimension window. Local Howe duality gives a nonzero small quotient. Harris (2007), Theorem2.A.1(b), then identifies the image of the \textbf{same lattice vacuum} with the tensor of the two nonzero spherical lines. Therefore the basic vector itself is nonzero in that quotient.

Finally, Harris Corollary2.2.7 and the Hecke-scheme argument in Corollary2.A.11 identify its Satake parameter on every point of the actual theta domain. The last sentence of the proof explicitly extends the generic parameter formula to all such points. This is used only after occurrence and vacuum-image nonvanishing have been established.

In the first step the semisimple unramified parameter is

\[
 \phi_{\pi_1}
 =(\phi_{\pi_0}\eta)\ \oplus\
  ((\mu_0\xi)^{-1}\eta)
       \bigoplus_{j=0}^{t-1}|\cdot|_E^{(t-1)/2-j}.             \tag{VA27}
\]

The displayed sum denotes Satake roots. It is not a claim that the local representation is a Steinberg representation, nor a global automorphic base-change assertion. The split formula (VA17) and the nonsplit formula agree. The second standard factor is consequently exactly (VA18) at every unramified place.

\subsection{7. Positive real tests for the actual selected types}\label{positive-real-tests-for-the-actual-selected-types}

AM-R-015 already verifies the signed Harish--Chandra lists, the opposing character shifts, and the real theta correspondences \(\tau_0\leftrightarrow\tau_1\leftrightarrow\tau_2\). We use that result and do not reprove the real classification. A further test is needed because real occurrence alone does not automatically prove nonvanishing of a particular normalized local Rallis pairing.

For a real dual pair with one member compact, the compact isotypic projection in the oscillator Hilbert space is the actual tensor product of that compact representation with its unitary Howe partner. This is the compact-partner form of Howe duality, explicitly realized by the orthogonal projector in McKee, Pasquale, and Przebinda (2025), equations(1),(4)--(5), printed pp.259--260.

Apply it to the first edge. Its selected Hilbert summand is \(\tau_0^\vee\widehat\otimes\tau_1\). Choose unit vectors \(v\in\tau_0\), \(h\in\tau_1\), and the corresponding unit oscillator tensor \(\overline v\otimes h\). The source oscillator matrix coefficient is the conjugate of the matrix coefficient of \(\tau_0\). Therefore the actual source coinvariant local pairing is

\[
 \int_{G_0(F_v)}
          |\langle v,\tau_0(g)v\rangle|^2\,dg
                       =\operatorname{fdeg}(\tau_0)^{-1}>0.   \tag{VA28}
\]

The equality is discrete-series Schur orthogonality for the stated Haar measure. The center of U(a,1) is compact; its volume is included in the formal degree. The integral converges absolutely. For compact-finite v and h, the tensor lies in the polynomial Fock module used by the actual creation maps. Indeed the central weights of the compact groups U(a) and U(1) fix the two Fock polynomial bidegrees in the positive and negative variables. Their joint fixed-type subspace is finite dimensional, so a compact-finite tensor in this Hilbert summand is already a polynomial Fock vector, not only a Hilbert completion limit. The source permits all such real polynomials. No new finite Schwartz vector is introduced.

At the second edge both real groups are compact. The same selected summand gives

\[
 \int_{G_1(F_v)}
          |\langle w,\tau_1(g)w\rangle|^2\,dg
                 ={\operatorname{vol}G_1(F_v)\over\dim\tau_1}>0
                                                               \tag{VA29}
\]

for a unit vector w. These are tests in an explicitly chosen isometric local model. A prescribed nonzero real intertwiner differs by its actual nonzero scalar; that scalar is not set to 1. The equations establish positivity and distinguish the formal-degree or compact dimension factor from the global multiplicity metric.

\subsection{8. Global nonvanishing for every automorphic embedding}\label{global-nonvanishing-for-every-automorphic-embedding}

Let \(\pi_0\) be any irreducible cuspidal automorphic constituent occurring in \(\mathcal C_0\), with any chosen nonzero automorphic multiplicity embedding. Its finite components are unitary. At G they have the actual \(K_r^1\)-fixed vectors; at Q and elsewhere they have the specified spherical vectors.

For the first edge, b\textgreater2r, so one must not cite the dimension window of GQT Theorem11.4. The target is anisotropic. The convergent Siegel--Weil identity, GQT Theorem7.1, and the doubling/unfolding argument of sections11.3,11.5--11.6 give the convergent Rallis formula, as in the checked application DO036. At G and Q its fixed finite local test is strictly positive: DO036(3)--(7), the full-rank disintegration identity, applies in every rank to (VA13). For the first G function its orbit period is (VA19), and for the first Q function its identity period is v. These values persist on an open compact set of target frames, so the integral of their squared norms is positive. The hypothesis b\textgreater4r ensures all its absolute-convergence bounds.

At basic places the actual tests are nonzero by section6. At real places (VA28) is strictly positive. The partial first standard Euler product has degree r and is absolutely convergent and nonzero at \(z_0\): its distance from the crude unitary Satake bound is 859 or733, both greater than1. Thus the actual first global lift is nonzero for this embedding. Changing the nonzero multiplicity embedding only multiplies the Petersson form by its positive norm; the conclusion holds for every such embedding.

Its output \(\pi_1\) is cuspidal because G1 is anisotropic. For the second edge \(b<N<2b\), and the target is again anisotropic. This is the anisotropic convergent case within GQT11.4's dimension window, not a claim that the global target has positive Witt index. The lower-tower premise is vacuous.

For any first descendant choose its actual ancestral \(\pi_0\). No common ancestor is asserted for a whole sum of constituents. Formula(VA27), or its split counterpart, gives the partial factorization

\[
 \mathcal L_1^S(z_0)
   =\mathcal L^S(z_0,\phi_{\pi_0}\mu_1\xi)
                         \prod_{i=1}^{t}L_E^S(i,\vartheta).
                                                               \tag{VA30}
\]

The first factor is the standard unramified Euler product, not an assumed global base change. It is absolutely convergent and nonzero by the same bound. The factors with \(i\ge2\) converge absolutely and are nonzero. The \(i=1\) factor is finite and nonzero by the carried Hecke boundary-value argument for the actual angular character \(\vartheta\).

There is nonzero local theta occurrence at every place with the required target: sections4--6 give the finite quotient tests; (VA29) gives the real normalized Rallis test. At finite places GQT11.6 identifies local occurrence with nonzero normalized Rallis pairing in this second dimension window. Theorem11.4 consequently gives a nonzero global second theta map for every actual multiplicity embedding of \(\pi_1\).

This last use is an existence assertion for the global multiplicity map. It does not say that every finite Schwartz vector has nonzero local norm. The preservation of the prescribed finite vectors is the separate quotient calculation in sections4--6.

\subsection{9. From local preservation to the actual fixed upward injection}\label{from-local-preservation-to-the-actual-fixed-upward-injection}

Fix unitary realizations of each local small theta quotient. By local Howe multiplicity one, the full global theta map on one irreducible automorphic constituent factors as

\[
             A_j\otimes\bigotimes_v q_{j,v}.                  \tag{VA31}
\]

Here \(A_j\) acts on the actual automorphic multiplicity space and each \(q_{j,v}\) is the local small quotient. The archimedean factor is the selected real intertwiner. This is the same algebraic factorization, with its multiplicity issue made explicit, proved in LIFT26(GS21--29): pair against target irreducible constituents and use the one-dimensional local intertwining spaces; the actual theta span is semisimple in the cuspidal unitary spectrum.

If \(A_j\) had a nonzero kernel vector, that vector would define a nonzero automorphic embedding whose theta lift vanished for all Schwartz data. Section8 excludes this for every embedding, so both multiplicity maps are injective. No multiplicity-one theorem for the automorphic spectrum is assumed.

Now insert the fixed finite tensor of DatumVec. At G, the composed local map has the nonzero left inverse(VA22). At Q it has(VA23). At every other finite place the actual first spherical line and then its actual second image survive by section6. At infinity the complete selected intertwiners are nonzero; their adjoints identify the multiplicity maps in IsIoUp. Tensoring these finite local injections gives an injection on every permitted finite initial type.

The source upward map is therefore injective on each initial isotypic constituent, including its entire multiplicity space. Distinct initial constituents cannot cancel after the two steps: if two terminal irreducibles agreed, local Howe injectivity at every place would identify their middle representations, and then their initial representations. Orthogonal decomposition of the cuspidal spectrum completes the proof \(\ker U=0\).

Only finite-dimensional fixed-weight, fixed-level multiplicity spaces are used in this last step. One does not claim that a set of K-finite scalar vectors is stable under the full real group. The actual smooth real model and its full intertwiners are those of R014/R015.

\subsection{10. Adjoint surjectivity, costs, and exact remaining gates}\label{adjoint-surjectivity-costs-and-exact-remaining-gates}

The source downward maps take the actual cuspidal domains into each other. At the noncompact bottom this follows from R014/R015: the only lower Witt target is killed by the exact real first-occurrence block. Their integrals and adjoints are therefore the genuine maps on \(\mathcal C_j\).

We prove the exact adjoint identity rather than invoking the source predicate AdjointLiftPeterssonIdentity. For one edge let H be the larger group and G the smaller group; let F be a larger-group multiplicity vector and u a smaller-group multiplicity vector. Write \(T_\phi:D_H\to D_G\) for the prescribed real intertwiner and \(T_\phi^+:D_G\to D_H\) for its actual adjoint. All inner products in this paragraph are linear in the first variable, as in the source. Because T is nonzero, choose compact-finite \(\phi,v',v\) with \[
 \beta=\langle T_\phi v',v\rangle_{D_G}
       =\langle v',T_\phi^+v\rangle_{D_H}\ne0.                \tag{VA31a}
\] Consider the double integral \[
 \mathcal I=\int_{[G]}\int_{[H]}
   \theta(g,h;\phi)F(v';h)\overline{u(v;g)}\,dh\,dg.
                                                                  \tag{VA31b}
\] It is absolutely convergent. For the first edge, {[}H{]} is compact, the theta function has uniform moderate growth in g, and the cuspidal coefficient \(u(v;g)\) is rapidly decreasing on Siegel sets. Their product is integrable; the bounded H-coefficient introduces no extra growth. For the second edge both automorphic quotients are compact. This proves the absolute Fubini premise for the fixed actual functions and measures.

First integrate over H, then use the actual smaller-group Schur identity for the multiplicity metric. The result is \[
                  \mathcal I=\langle DF,u\rangle\,\beta.
\] Interchanging the two integrals instead gives \[
 \mathcal I=\int_{[H]}F(v';h)
   \overline{\int_{[G]}\overline{\theta(g,h;\phi)}u(v;g)\,dg}
     \,dh
  =\langle F,\mathrm{Up}\,u\rangle
        \langle v',T_\phi^+v\rangle_{D_H}
  =\langle F,\mathrm{Up}\,u\rangle\,\beta.                    \tag{VA31c}
\] Here the inner integral is exactly IsIoUp, including its conjugated kernel; the last equality is IsTPlus. Cancelling the nonzero beta proves that Up is the Petersson adjoint of D on this edge, without altering a Haar, real-intertwiner, or Schur scalar. Apply this to both edges to obtain the claimed \(U=(D_0D_1)^*\). The definitions being verified are \nolinkurl{OrthogonalComplementCharacterisation.lean} 74--88 and 278--298. The latter is a predicate, so its required identity has been proved here in the genuine analytic setting.

For finite-dimensional Hilbert spaces, \[
            (\operatorname{im}D)^\perp=\ker D^*=\ker U=0.
\] Hence \(\operatorname{im}D=\mathcal C_0\), proving(VA8). Applying any fixed idempotent proves(VA9). This is onto the full finite-level bottom cusp multiplicity space, not merely onto an unspecified nonzero theta image.

No normalization is lost in this implication. In chosen unitary local quotient models the norm on each initial constituent can be written \[
 \|U x\|^2
   =\big\|(A_1A_0)\otimes
         \bigotimes_{v<\infty}
          q_{1,v}(f_{1,v},q_{0,v}(f_{0,v},-))\,x\big\|^2,
                                                               \tag{VA32}
\] with the nonzero real-intertwiner factors retained inside \(A_0,A_1\). If the actual global Haar or multiplicity metrics differ from probability or Tamagawa choices, their positive conversion factors remain in this expression. Local models may be rescaled only with the inverse change in the multiplicity maps. Thus(VA19),(VA21),(VA23), (VA28), and(VA29) are explicit local factors, not a declaration that the entire global scalar is their product.

At a fixed complete weight and level, injectivity gives some positive smallest singular value, and the least-norm right inverse of D is \[
                         D^*(DD^*)^{-1}.                      \tag{VA33}
\] Its size is not bounded here. No integral saturation, 11-adic denominator bound for this inverse, or DP44 improvement follows from qualitative surjectivity.

In particular:

\begin{itemize}
\tightlist
\item
  LIFT28/29's detector belongs to the full bottom cusp space and thus, at weights where that detector exists and the R015 conditions hold, to this actual two-step downward image.
\item
  The source map Th also includes e0. A particular detector has not thereby been proved to lie in e0's image or to retain its period after that projection. Root's separate finite idempotent audit is needed for that assertion.
\item
  The finite frames here are the standard genuine realization of the explicit source formulas. A completed RealPack additionally requires tS, RowLevelMatchA/B/C, betaY, top coefficient and arithmetic period identifications. This theorem does not assert them.
\item
  The ineffective holomorphic interpolation threshold must still be fitted to full SizeData2. The actual terminal brC/tilQ/PtilFF identity needs the same-family source transport.
\end{itemize}

\subsection{11. References and accepted inputs}\label{references-and-accepted-inputs}

Gan, W. T., Qiu, Y., \& Takeda, S. (2014). The regularized Siegel--Weil formula (the second term identity) and the Rallis inner product formula. Inventiones Mathematicae, 198, 739--831. Theorem7.1; sections11.3,11.5--11.8, especially Theorem11.4 and Proposition11.6. Local copy: work/agents/lift\_refs/gqt2014.pdf and .txt.

Gong, Z., \& Grenie, L. (2011). An inequality for local unitary Theta correspondence. Annales de la Faculte des Sciences de Toulouse, 20(1), 167--202. \url{https://doi.org/10.5802/afst.1289} . LemmaA.1 and its full proof, printed pp.199--200. \url{https://numdam.org/item/10.5802/afst.1289.pdf} . Local copy: work/agents/lift\_refs/gong\_grenie2011.pdf and .txt.

Harris, M. (2007). Cohomological automorphic forms on unitary groups, II: Period relations and values of L-functions. In Harmonic Analysis, Group Representations, Automorphic Forms and Invariant Theory. Corollary2.2.7, printed p.21; Appendix2.A, Theorems2.A.1--2, Corollaries2.A.3 and2.A.11, printed pp.26--28. \url{https://webusers.imj-prg.fr/~michael.harris/periodrelations.pdf} . Local copy: work/agents/lift\_refs/harris2007\_periodrelations.pdf and .txt.

Harris, M., Kudla, S. S., \& Sweet, W. J. (1996). Theta dichotomy for unitary groups. Journal of the American Mathematical Society, 9(4), 941--1004. Proposition2.2, especially the character factor under doubling; section6 for the spherical standard-factor convention. Primary published PDF retrieved at \url{https://scispace.com/pdf/theta-dichotomy-for-unitary-groups-1x5vgrzgv7.pdf} .

McKee, M., Pasquale, A., \& Przebinda, T. (2025). Symmetry breaking operators for dual pairs with one member compact. Journal of Mathematical Sciences, the University of Tokyo, 32(3), 257--400. Equations(1),(4)--(5), printed pp.259--260. \url{https://www.ms.u-tokyo.ac.jp/journal/1e7cc0870888f546c6a1c8893158d493cf517a8d.pdf} . Local copy: work/agents/lift\_refs/mpp2025\_symmetry.pdf and .txt.

Accepted inputs are used with their checked scope: R014/R015 real types and lower-tower cuspidality; K031/K032 actual hosts; K040/K041 current characters and weight transport; LIFT14's nonzero boundary Hecke value; LIFT16's split spherical calculation; LIFT24/27's whole-Schwartz quotient and unitary Godement--Jacquet bound; DO036's positive first full-rank local pairing; DO041's vertical Q calculation; and LIFT26's explicit automorphic multiplicity factorization. None is re-labelled as a newly proved general theorem here.

\subsection{12. Verification record}\label{verification-record}

The author read the complete final proof, including the unequal-label quotient, the exact G and Q preceding-image norms, the all-rank basic test, both real local pairings, the every-embedding multiplicity argument, and the direct Fubini/Schur adjoint proof. The finite Q calculation was also independently reviewed in work/agents/review\_vertical\_q\_quotient.md, with its complete certificate replayed.

The exact finite diagnostic \nolinkurl{work/agents/vertical_g_adjoint_certificate.py} passed at 2026-11-22T03:15:47.945903+00:00. It exhausts the raw two-kernel G composition for ranks 1 to 2 to 3 over residue fields of orders 2 and 3, including all admitted tame character triples. It checks 90,024 signed composition values on 11,400 target matrices, their flag support, both raw norms, the exact composed scalar and induced compact norm, together with the two actual rank and Euler margins. The first script run exposed an ancillary enumerator-count assertion that omitted the independently varying initial diagonal unit; that assertion was corrected from q(q-1) to q(q-1)\^{}2 before the successful full replay. No mathematical formula in this proof changed.

These exact computations support the finite formulas. They are not formal certificates for the global analytic theorem, a Lean build, or a bound for the global right inverse in (VA33).

\clearpage

\section{AM-K043 --- the exact selected bottom PBW normalization}\label{am-k043-the-exact-selected-bottom-pbw-normalization}

Author: GPT-6 Astra Pro, ALP\_MACHINE /root/kummer\_extension, 22 November 2026 UTC. Inherited reasoning effort is not separately exposed.

This is a source-aligned application of Nakahama's known norm formula, not a claim of a new holomorphic discrete-series norm theorem. It replaces the unspecified nonzero PBW scalar in LIFT28 Section 7 by an explicit positive square root of an integer, in the actual matrix-unit and Fischer conventions. It gives its exact finite valuations and elementary real bounds for all R022 words and every allowed K042 canonical twist.

The local scalar does not bound a global interpolation right inverse. Section 6 records the precise quantitative source-counting hypothesis which the current bottom profile does not supply. No cusp multiplicity estimate, terminal nonzero value, full Host/SizeData₂, or DP44 cost sign is asserted.

\subsection{1. The primary norm theorem and its conventions}\label{the-primary-norm-theorem-and-its-conventions}

Nakahama (2016), Theorem 5.1, printed pp.960--961, treats SU(q,s), q\textless s, with an arbitrary polynomial irreducible U(s) fibre. For q=1, the norm ratio on the Pieri component of highest weight n in V\_k tensor Sym\^{}N(C\^{}s), with constants normalized to their given fibre norm, is \[
 \frac{\|P\|_{\mathcal H}^{2}}{\|P\|_{\rm F}^{2}}
   =\prod_{j=1}^{s}
       (\Lambda-j+1+k_j)_{n_j-k_j}^{-1}.
 \tag{PB1}
\] Here (r)\_m=r(r+1)\ldots(r+m−1), with (r)\_0=1. The Fischer norm is the Gaussian norm in his equation (3.2), printed p.940; in coordinates it has \(\|z^\gamma v\|_{\rm F}^{2}=\gamma!\|v\|^{2}\). His action is equation (5.5), printed p.959.

Primary PDF: \url{https://jolt.centre-mersenne.org/item/10.5802/jolt.919.pdf}. The author version \url{https://arxiv.org/pdf/1506.05919} has the same Theorem 5.1 on printed p.32. Both PDFs and extracted text have been retrieved to work/research/nakahama2016 and nakahama\_arxiv1506. The proof below uses the published statement, not an inferred scalar-case extension.

\subsection{2. Exact translation to the bottom parameters}\label{exact-translation-to-the-bottom-parameters}

At one CM label use LIFT28's notation with zero-based indices: \[
 p_i=v_i\ (i<a-1),\quad p_{a-1}=-b_0,\quad q=-b_1,
\] \[
 \alpha_i=p_i-\frac{a-2}{2}+i,\qquad
 \beta=q-\frac a2,\qquad
 \lambda_i=u_i-\frac{a-1}{2}+i.
 \tag{PB2}
\] Thus the minimal compact type is V\_alpha tensor z\^{}beta and the selected type is V\_lambda tensor z\^{}(beta−N), where \[
 m_i=\lambda_i-\alpha_i=u_i-p_i-\tfrac12,\qquad
 N=\sum_i m_i.
 \tag{PB3}
\] The actual interlacing makes every m\_i a nonnegative integer and lambda/alpha a horizontal strip. Its Pieri multiplicity is one.

Swap the two Hermitian blocks to identify the source U(a,1) with U(1,a); multiplying the form by −1 changes no unitary group. Write a compact element as diag(z,A). Restricted to SU(1,a), z det A=1. Choose an integer c with alpha+c(1,\ldots,1) nonnegative, and apply PB1 with \[
 k=\alpha+c\mathbf1,\quad
 n=\lambda+c\mathbf1,\quad
 \Lambda=-\beta-c.
 \tag{PB4}
\] For a compact element, Nakahama's action on constants is det(A)\^{}Lambda V\_k(A), which is det(A)\^{}(-beta)V\_alpha(A)=z\^{}beta V\_alpha(A). His polynomial variables transform by z\^{}(-1) times the standard A action, exactly as in LIFT28(HJ7). This verifies both the minimal type and the selected type, including the U(1) character. The auxiliary determinant shift c cancels in every factor of PB1. The theorem's convergence condition is Lambda+k\_a\textgreater a, equivalently b1−b0\textgreater0, and holds here. The compact central extension from SU to U adds a unitary scalar character and changes neither the norm ratio nor the noncompact matrix-unit operators.

Consequently define the positive integer \[
 \begin{split}
 D_{\lambda,\alpha,\beta}
 &=\prod_{i=0}^{a-1}(\alpha_i-\beta-i)_{m_i}\\
 &=\left[\prod_{i=0}^{a-2}
       (v_i+b_1+1)_{\,u_i-v_i-1/2}\right]
       (b_1-b_0+1)_{\,u_{a-1}+b_0-1/2}.
 \end{split} \tag{PB5}
\] Every starting factor is positive, because alpha\_i−beta−i=p\_i−q+1\textgreater1. On this actual selected polynomial summand the Hilbert norm is the Fischer norm divided by D.

\subsection{3. From the polynomial norm to the PBW map}\label{from-the-polynomial-norm-to-the-pbw-map}

Let C be the PBW creation map from V\_alpha tensor Sym\^{}N(C\^{}a) to the discrete series, using the noncompact matrix units. Choose the sign of the lowering operators so they act as the constant derivatives partial\_i in the holomorphic coordinates. The adjoints of these lowering operators are the creation operators; the opposite matrix-unit sign changes only the common phase (−1)\^{}N, never any following norm or valuation.

For homogeneous vector-valued polynomials P,Q of degree N, moving the N creation operators by adjunction gives \[
                    \langle C(P),Q\rangle_{\mathcal H}
                        =\langle P,Q\rangle_{\rm F}.
 \tag{PB6}
\] Indeed, for P=z\^{}gamma v and Q=sum\_delta z\^{}delta q\_delta, the resulting constant term is gamma! times the fibre pairing with q\_gamma. Constants have the given fibre norm by PB1's normalization. Linearity proves PB6.

Both C and the two metrics commute with the compact group. On the multiplicity-one selected summand PB1 and PB6 therefore give the exact operator and norm identities \[
 C|_{V_\lambda}=D\,{\rm id},\qquad
 \|C(P)\|_{\mathcal H}^{2}=D\|P\|_{\rm F}^{2}.
 \tag{PB7}
\] The first identity uses the polynomial realization and the chosen adjoint signs. It is not a statement that C is the identity map in an arbitrary realization.

\subsubsection{Theorem PB8}\label{theorem-pb8}

Choose a Fischer-isometric Pieri inclusion j:V\_lambda→V\_alpha tensor Sym\^{}N(C\^{}a), and normalize the minimal-type inclusion into the unitary discrete series isometrically. Then the PBW scalar in LIFT28 Section 7 has absolute value \[
                  \rho=\sqrt{D_{\lambda,\alpha,\beta}}.
 \tag{PB8}
\] The selected isometry is J=rho\^{}(-1) Cj, up to its fixed unit phase. For a non-isometric inclusion j with j* j=c\_j id, the formula is rho²=c\_j D. The invariant scalar c\_j must then be retained.

\textbf{Proof.} Apply PB7 to a unit vector in the isometric image of j. This yields PB8, and Schur's lemma gives the assertion on the whole selected component. Rescaling j gives the last formula. The dual/contragredient realization uses the dual isometry and has the same positive rho. QED.

The lift agent confirmed that LIFT28 uses these matrix-unit normal coordinates and the Fischer convention gamma!, so PB8 supplies its local scalar. Any scalar in a chosen global automorphic realization remains separate.

\subsubsection{The factorial distinction and a cancellation}\label{the-factorial-distinction-and-a-cancellation}

The raw jet functional associated to P=sum\_gamma z\^{}gamma p\_gamma is the Fischer contraction \[
 {\cal D}_P F(0)
   =\sum_\gamma\gamma!\langle p_\gamma,F_\gamma\rangle.
 \tag{PB9}
\] It uses derivatives, not divided Taylor coefficients. The transpose of the unit selected inclusion is rho\^{}(-1) times this raw contraction, with the chosen sign and frame terms. In the universal Hilbert model at the origin this normalized functional has norm one on the selected component: \[
       \|{\cal D}_P\|=\rho\|P\|_{\rm F},\qquad
       \|\rho^{-1}{\cal D}_{jv}\|=\|v\|.
 \tag{PB10}
\] These follow from PB6--PB8 by Riesz duality. Thus a large rho is not by itself an unavoidable global interpolation loss. It cancels with the norm of the corresponding raw local derivative. The norm of the actual automorphic packet evaluation map and of its right inverse is a different quantity. No equality between that global map and the universal local functional norm is asserted.

\subsection{4. Evaluation on every R022 word and allowed K042 twist}\label{evaluation-on-every-r022-word-and-allowed-k042-twist}

Write ell for the fixed R022 gap, s for the K042 canonical twist, and q\_s=(a+1)s. Throughout the exact interval in K042, 0≤q\_s≤ell and s is a multiple of28. At a label of role beta, \[
 b_0=E_\beta-\ell,\qquad b_1=E_\beta+\ell+q_s.
 \tag{PB11}
\] For an interior bit omega\_i, 1≤i≤a−2, the exponents in PB5 are \[
 m_i=4\ell-1-\ell\omega_i.
 \tag{PB12}
\] The endpoint exponents are exactly \[
\begin{array}{c|cc}
 &m_0&m_{a-1}\\ \hline
 {\rm top}
 &A+\frac{a-2}{2}-4\ell(a-1)
 &E_T-B-\ell-\frac a2\\
 {\rm thick\ or\ thin}
 &4a\ell&E_\beta+\ell .
\end{array} \tag{PB13}
\] Their sum is K042(NB9). In particular the endpoint exponents do not depend on the word or on s; only the Pochhammer starting points vary with s. They are \[
 r_i=4\ell(a-1-i)+E_\beta+\ell+q_s+1\ (i<a-1),\qquad
 r_{a-1}=2\ell+q_s+1.
 \tag{PB14}
\] Equations PB11--14 are obtained by direct substitution in PB2--5, so they apply to the actual selected component, not an ambient monomial of the same total degree.

For all words and twists, the maximum of r\_i+m\_i−1 equals u\_0+b1−1/2. Exact endpoint checks give these digit bounds:

\begin{longtable}[]{@{}lr@{}}
\toprule\noalign{}
Role & base-11 digits of every r\_i+m\_i−1, both hosts \\
\midrule\noalign{}
\endhead
\bottomrule\noalign{}
\endlastfoot
top & at most278 \\
thick & at most144 \\
thin & at most139 \\
\end{longtable}

The comparisons use the largest twist and all bit choices; no monotonicity of finite valuations in s is assumed.

\subsection{5. Exact finite valuations and real bounds}\label{exact-finite-valuations-and-real-bounds}

For any rational prime p and positive integers r,m, \[
 v_p((r)_m)
 =\sum_{h\ge1}\left(
       \left\lfloor\frac{r+m-1}{p^h}\right\rfloor
       -\left\lfloor\frac{r-1}{p^h}\right\rfloor\right)
 =\frac{m-s_p(r+m-1)+s_p(r-1)}{p-1}.
 \tag{PB15}
\] This is the factorial identity (r)\_m=(r+m−1)!/(r−1)! and Legendre's formula; the latter follows by counting multiples of each p\^{}h in the factorial. Hence PB5 specifies v\_p(D) exactly without constructing its enormous integer value.

If L\_i is the number of base-p digits of r\_i+m\_i−1, then both digit sums in the difference lie between0 and(p−1)L\_i. Therefore \[
 \left|v_p(D)-\frac{N}{p-1}\right|\le\sum_iL_i.
 \tag{PB16}
\] For p=11 at a label of role beta, this gives the uniform statement \[
 \left|v_{11}(\rho)-\frac{N}{20}\right|
       \le \frac{a L_\beta}{2},
 \qquad (L_T,L_H,L_L)=(278,144,139).
 \tag{PB17}
\] Here rho is the positive algebraic square root of the integer D; the normalized extension of v11 gives half its valuation. For a product of labels, the exact valuations and errors add. This is a local normalization calculation, not an image-content claim for a bottom cusp lattice.

The companion certificate also computes the exact minimum and maximum v11(D) over \textbf{all} words at s=0 and s=s\_max. At a fixed s each interior bit contributes independently to the sum PB15, so these extrema are obtained by choosing the smaller or larger of its two contributions at every coordinate. They are not estimated by treating word order as a Hamming-only invariant. No extremum over all intermediate twists is claimed.

\subsubsection{One explicit field for all word-dependent PBW scalars}\label{one-explicit-field-for-all-word-dependent-pbw-scalars}

Fix s and the CM role assignment. Let D\_0 be the product of PB5 over all labels for the zero word. For 1≤i≤a−2 put \[
 B_i=\prod_\tau (r_{\tau,i}+3\ell-1)_\ell\in\mathbb Z_{>0}.
\] Changing bit i from zero to one removes exactly this block from the Pochhammer product, and therefore \[
 D_{\rm all}(\omega)=D_0\prod_i B_i^{-\omega_i},\qquad
 \rho_{\rm all}(\omega)=
       \frac{\sqrt{D_0}}{\prod_i(\sqrt{B_i})^{\omega_i}}.
 \tag{PB17a}
\] Positive real square roots fix the signs. Consequently every global scalar for all 2\^{}(a−2) words lies in the one Galois field \[
 \mathbb Q(\sqrt{D_0},\sqrt{B_1},\ldots,\sqrt{B_{a-2}}),
 \qquad [\,\cdot:\mathbb Q]\le2^{a-1}. \tag{PB17b}
\] If the individual label scalars, rather than only their product, are needed, R022 has at most three distinct role profiles. Using these radicals separately for each occurring role gives degree at most2\^{}(3(a−1)), independently of f and of the number of words. These bounds are for a fixed s and for the PBW scalars only. They do not bound the full coefficient field, nor all intermediate twists simultaneously.

Since each D\_all(omega) is a positive integer, its positive square root has absolute logarithmic Weil height \(h(\rho_{\rm all}(\omega))=\tfrac12\log D_{\rm all}(\omega)\). This follows directly from its conjugates, which are the two signs of that root, or from its rational integer value if it is a square. In particular its height is at most(log D\_0)/2. This is not an assertion that adjoining the scalar is necessary: one may instead retain the unnormalized PBW map and its exact norm PB8.

For an elementary real estimate put R=max\_i(r\_i+m\_i−1). Because (r)\_m≥m!, and the multinomial coefficient is at most a\^{}N, \[
 D\ge\prod_i m_i!\ge N!/a^N,\qquad D\le R^N.
\] Consequently, for N≥1, \[
 \frac12\{N\log(N/a)-N+1\}\le\log\rho
                   \le\frac N2\log R.
 \tag{PB18}
\] The lower estimate uses log(N!)≥integral\_1\^{}N log(t)dt. These bounds and the exact product are compatible with PB10; neither supplies a global detector norm.

\subsection{6. The quantitative multiplicity route has a separate gap}\label{the-quantitative-multiplicity-route-has-a-separate-gap}

The source has a number-field-generic counting interface. AutInst in \nolinkurl{SharpRestriction/CuspSpanNormBounds/GlobAut.lean32–74} uses a CM extension, a selfdual Hermitian lattice and real isometries. It does not itself impose the failed Fields.Kq elliptic tensor. PlaceDegreeData.lean55--59 defines HidaSetting; CountData is at183--194. However its PointsLemma, lines360--377, requires, in addition to the actual idempotent and representation data, \[
      C n^{19}f^5\log(b_{\rm cnt}\,{\rm rd})\le h,\qquad
      |\operatorname{HC}_i|\le b_{\rm cnt}^2 .
 \tag{PB19}
\] The uniform exported thresholds n0,C0 are existential: \nolinkurl{PlaceDegreeData.lean411–424.} This is distinct from LIFT28's coherent-jet Serre threshold.

In the present bottom group n=a+1. The smallest positive/negative HC gap is b1−b0=2ell+q\_s≤3ell. Thus any admissible h is at most 3ell, by \nolinkurl{FormalDegreeBessel/StatementDS.lean16–17’s} literal opposite-sign gap definition. Also b\_cnt²≥b1≥E\_L, and rd≥1. Exact R022 arithmetic gives E\_L\textgreater2\^{}P with P=480 for J14 and P=478 for J56. Since log2\textgreater2/3, \[
 \log(b_{\rm cnt}\,{\rm rd})>P/3,\qquad
 C<\frac{9\ell}{P(a+1)^{19}f^5}.
 \tag{PB20}
\] The last right side has the following exact values/bounds:

\begin{longtable}[]{@{}
  >{\raggedright\arraybackslash}p{(\columnwidth - 4\tabcolsep) * \real{0.3000}}
  >{\raggedright\arraybackslash}p{(\columnwidth - 4\tabcolsep) * \real{0.3000}}
  >{\raggedleft\arraybackslash}p{(\columnwidth - 4\tabcolsep) * \real{0.4000}}@{}}
\toprule\noalign{}
\begin{minipage}[b]{\linewidth}\raggedright
Host
\end{minipage} & \begin{minipage}[b]{\linewidth}\raggedright
rational upper allowance for C
\end{minipage} & \begin{minipage}[b]{\linewidth}\raggedleft
simpler bound
\end{minipage} \\
\midrule\noalign{}
\endhead
\bottomrule\noalign{}
\endlastfoot
J14 & 1612716047/2703882435438619499115494475276645461760 & \textless2\^{}(-100) \\
J56 & 120626011/45210562042634360301430306501976426261415827190902784 & \textless2\^{}(-148) \\
\end{longtable}

In particular PB19 fails throughout K042's interval for every C≥1. This does not prove the optimal source constant C0 is at least1, and does not negate the accepted counting theorem. Rather, no effective exported constant meeting PB20 has been provided. Choosing the harmless weaker export max(C0,1) cannot prove the desired estimate on this profile.

Even an unrestricted common determinant shift of the entire bottom HC list cannot evade the simpler bounds in the last column: its spread is at least E\_L, so some shifted HC coordinate has absolute value at least E\_L/2. Replacing P by P−1 in PB20 gives the same certified bounds C\textless2\^{}(-100) and C\textless2\^{}(-148).

I also inspected Stadler (2025), Theorem3.4 and Theorem5.2, arXiv:2504.01183v2, printed pp.11 and18. Those formulas concern one real group SU(a,1) over an imaginary quadratic field. Our restriction-of-scalars arithmetic lattice has f real factors and is not thereby covered. They are not used as a multiplicity bound here.

Even a rough dimension heuristic must retain the G levels. At a split place their exact finite indices obey \[
 \frac{[{\rm GL}_{a+1}(\mathbb F_q):
                      {\rm diag}(U_a,1)]}
      {[{\rm GL}_{a}(\mathbb F_q):U_a]}
       =q^a(q^{a+1}-1).
 \tag{PB21}
\] This follows from the standard GL order formula, since the two unipotent denominators are identical. These factors occur at every selected G place. Comparing only T\^{}a with E² and dropping PB21 would not be a valid selected-interpolation dimension test.

The remaining effective target is still K042(NB14): a selected jet surjection, or an adequate quantitative period estimate, for some s in its exact finite progression, with the actual global norm and arithmetic lattice controlled. PB8 and PB15 remove an unspecified local scalar from that target.

\subsection{References and checks}\label{references-and-checks-1}

Nakahama, R. (2016). Norm computation and analytic continuation of vector valued holomorphic discrete series representations. Journal of Lie Theory, 26(4),927--990. \url{https://doi.org/10.5802/jolt.919}.

Stadler, A. (2025). Explicit multiplicities in the cuspidal spectrum of SU(n,1) (arXiv:2504.01183v2). \url{https://arxiv.org/abs/2504.01183}.

The certificate independently computes low-rank PBW creation operators in the symmetric-tensor fibre, checks the Pieri eigenvalue and factorial normalization symbolically, checks PB11--17 for the actual profile, obtains exact all-word endpoint valuation extrema, and verifies PB20--21 by integer arithmetic. These are finite diagnostics and exact profile calculations. The known general norm theorem is cited rather than rederived; no cusp dimension, global right inverse, or Lean build is claimed.

\clearpage

\section{AM-R-026: reciprocal good tuples and the genuine finite projector}\label{am-r-026-reciprocal-good-tuples-and-the-genuine-finite-projector}

Author: GPT-6 Astra Pro, ALP\_MACHINE \nolinkurl{/root}, 22 November 2026 UTC. Reasoning effort is not separately exposed.

This result verifies coefficientwise reality of the source's prescribed finite projector when its series idempotents are realized by the genuine finite general-linear Lusztig series. It uses the entire \nolinkurl{goodTuple} predicate, including both tame-character orientations. It does not infer this realization from the abstract idempotent axioms alone. It also computes a universal denominator bound for this actual projector.

Nonvanishing of a finite projector and survival of an automorphic period are different assertions. The latter is not proved here. In particular, an extra negative-line torsion stabilizer at a coarse CM point can impose a central-character condition not detected by restriction to GL\_a alone.

\subsection{1. Literal predicates and the reciprocal involution}\label{literal-predicates-and-the-reciprocal-involution}

The pinned Alpöge archive commit is 7cbf7a0c32fca1badf05931879a37c3f8a2b03b4. The exact definitions are \nolinkurl{Formalism/CyclotomicNorms/StatementT.lean}19--29 and \nolinkurl{Formalism/SharpRestriction/LevelLetterCocycles/DataN.lean}996--1007. For a rational prime g, let f\_v be monic of degree n with nonzero constant term. Define

\[
 f_v^{\vee}(X)=f_v(0)^{-1}X^n f_v(X^{-1}).                 \tag{R26.1}
\]

This is monic of degree n, its constant term is f\_v(0)\^{}\{-1\}, its roots are the inverses of the roots of f\_v, and (f\_v\^{}\{\vee\})\^{}\{\vee\}=f\_v. The superscript here denotes the normalized reciprocal polynomial, not complex conjugation.

For a nonzero finite-field element x, the literal \nolinkurl{primeToPart} in StatementT.lean9--12 is its prime-to-l component in the finite cyclic group generated by x. Indeed the defining exponent is divisible by the full l-part of its order and congruent to1 modulo the prime-to-l part. In particular

\[
 \operatorname{pt}_{l}(x^{-1})
       =\operatorname{pt}_{l}(x)^{-1}.                     \tag{R26.2}
\]

\subsubsection{Proposition R26.3}\label{proposition-r26.3}

For any finite indexing sets of places and tame characters,

\[
 \operatorname{IsGood}(n,g,(f_v),\chi)
 \quad\Longleftrightarrow\quad
 \operatorname{IsGood}(n,g,(f_v^{\vee}),\chi^{-1}).         \tag{R26.3}
\]

\textbf{Proof.} We check all six clauses of the literal definition. Equation(R26.1) preserves monicity, degree and nonzero constant term. Inversion is a bijection on the nonzero roots, so it preserves squarefreeness. At every prime l outside the same BB(n,g), equality of the prime-to-l components of two reciprocal roots is, by(R26.2), equivalent to equality of the corresponding original components; the original injectivity clause then identifies the roots. At every prime l different from g, equality between the projected reciprocal root and a projected inverse tame character would invert to the forbidden equality for f\_v and chi. This checks the two prime-to-part clauses, with their distinct prime quantifiers unchanged.

If q divides f\_v, then q has nonzero constant term after multiplication by a scalar to make it monic, and q\^{}vee divides f\_v\^{}vee with the same degree. The reciprocal construction is involutive, so it bijects divisors of every specified degree. Hence the clause excluding a degree-a divisor at some place is preserved for each1\textless=a\textless=n/2. An irreducible nonconstant factor remains irreducible under this bijection: a factorization of its reciprocal would reciprocate to a factorization of the original. Its degree is unchanged, so the prime-degree factor with n\textless2deg(q) and deg(q)+3\textless=n is preserved. Applying the same argument twice gives both implications. QED.

The source \nolinkurl{goodTuple} requires IsGood for chi and also for chi inverse, at both rational G primes. PropositionR26.3 therefore gives

\[
       \operatorname{goodTuple}(f)
          \Longleftrightarrow\operatorname{goodTuple}(f^{\vee}).
                                                               \tag{R26.4}
\]

No disjointness claim about BB(n,g0) and BB(n,g1) is used. No existence of a good tuple is needed for this involutive symmetry.

\subsection{2. Genuine finite-series idempotents}\label{genuine-finite-series-idempotents}

For G\_g=GL\_n(F\_g), identify its dual finite group with G\_g and use the fixed finite-field character duality in the actual spectral realization. A monic degree-n polynomial f with f(0) nonzero determines the unique semisimple conjugacy class s\_f having characteristic polynomial f. Let E(G\_g,s\_f) be its genuine rational Lusztig series, and set

\[
 e_f=\sum_{\pi\in E(G_g,s_f)}e_\pi,\qquad
 e_\pi={\dim\pi\over|G_g|}
            \sum_{x\in G_g}\chi_\pi(x^{-1})[x].             \tag{R26.5}
\]

The primitive central idempotent formula is the usual finite-group character projector. All occurrences below use this genuine choice for the source eSer. The listed abstract eSer axioms alone would not identify the series or force each allowed e\_f to be nonzero.

Complex conjugation carries the series at s to the series at s\^{}\{-1\} (Srinivasan \& Vinroot, 2019, Lemma3.2, author PDF pp.7--8). Since the characteristic polynomial of s\_f\^{}\{-1\} is f\^{}vee, the character-projector formula gives

\[
                         \overline{e_f}=e_{f^{\vee}}.       \tag{R26.6}
\]

The bar acts on coefficients, without inverting group elements. This distinction is essential for the source RowLevelMatchC.

Each E(G\_g,s\_f) is nonempty: the Jordan decomposition bijection of Srinivasan and Vinroot(2019), Theorem4.1, author PDF pp.12--13, identifies it with the unipotent characters of the connected centralizer, which include the trivial character. When f is squarefree that centralizer is a torus, so the series is a singleton. In the finite-GL Green parametrization its nonempty partitions all equal(1). Gow and Vinroot(2007), Theorems3.2 and4.3, author PDF pp.9 and12, then identify this constituent as generic, with multiplicity1 in the Gelfand--Graev representation. These are applications of the known finite-group results, not new classification theorems.

\subsection{3. The exact source e0 is real and orthogonal}\label{the-exact-source-e0-is-real-and-orthogonal}

Let

\[
 \mathcal G=\prod_{k=0}^{1}\prod_{v\in P_k}GL_n(F_{g_k}),\qquad
 e_0=\sum_{f\;\mathrm{goodTuple}}
                        \bigotimes_{k,v}e_{f_{k,v}}.        \tag{R26.7}
\]

This is the exact finite coefficient sum of DataN.lean1001--1007: its coefficient-vector indexing parametrizes every monic degree-n polynomial once, and goodTuple discards zero constant terms and all other forbidden tuples. Distinct tuples index orthogonal central idempotents. Consequently(R26.7) is itself a central idempotent, and it is self-adjoint for the unitary group-algebra involution.

Applying coefficient conjugation to(R26.7), using(R26.6), and reindexing by the bijection(R26.4) proves

\[
                   \overline{e_0(\gamma)}=e_0(\gamma)
                         \quad(\gamma\in\mathcal G).        \tag{R26.8}
\]

This is coefficientwise reality of the whole prescribed sum. It does not require any single admitted irreducible representation to be self-dual. On every unitary representation of Gcal the operator e0 is an orthogonal projection, hence has norm at most1. Its group-algebra value is nonzero if and only if the set of good tuples is nonempty, because every genuine series in(R26.7) is nonempty and different summands are orthogonal.

Under the canonical finite group identification of R025, the source grp\_record equality e\_A(sigma gamma)=e0(gamma) now implies

\[
              e_A(\sigma\gamma)=\overline{e_0(\gamma)},      \tag{R26.9}
\]

which is exactly RowLevelMatchC, with R025's same t=1 and the same residue action. Thus coefficientwise reality is no longer an additional condition once the genuine same-projector row identification is made. A separate arithmetic family cannot be silently assigned a new projector; its grp\_record equality still must be the actual equality in that family's construction.

\subsection{4. A retained arithmetic denominator bound}\label{a-retained-arithmetic-denominator-bound}

Let M be the exponent of Gcal and L=Q(zeta\_M). Every irreducible character value is a sum of roots of unity of order dividing M, since each finite-order matrix is diagonalizable. Thus its values are algebraic integers in L. From(R26.5)--(R26.7),

\[
                |\mathcal G|\,e_0\in O_L[\mathcal G].       \tag{R26.10}
\]

The coefficients themselves lie in the real subfield of L by (R26.8). At every extension of the normalized11-adic valuation, R025's exact group-order computation therefore gives

\[
  v_{11}(e_0(\gamma))\ge
  \begin{cases}
      -336,&J14,\quad n=43,\ f=42,\\
      -21336,&J56,\quad n=169,\ f=168.
  \end{cases}                                               \tag{R26.11}
\]

For zero coefficients the valuation is infinite. Equivalently, \textbar Gcal\textbar e0 preserves every Gcal-stable lattice over O\_L. This is an upper bound on a sufficient clearing denominator, not a sharp content calculation or a claim that the entire group order is always necessary. It supplies no denominator bound for an inverse restriction or an automorphic interpolation map. The possibly large cyclotomic field L is explicit here; it is not declared to be the minimal or previously fixed coefficient field.

\subsection{5. Scope and frameworks}\label{scope-and-frameworks}

This closes the reciprocal/reality subproblem of R025 for the genuine finite-series realization and gives its elementary operator-norm and denominator controls. It does not prove that the actual bottom cusp space contains the required e0 type, that the LIFT29 period survives projection, or that a particular jet interpolates in K042's finite canonical-twist interval. A coarse CM point has a stabilizer in H times the negative-line unitary torus, so compatibility with that torus must be included in any finite-orbit interpolation argument. The full mixed Host/RealPack and DP44 costs remain subject to their previously stated gates.

Geometrically this is an actual idempotent of the characteristic-zero group algebra of a finite etale split residue group, with an explicit integral multiple. It is not merely a K-theory class. The unchanged framework locators are SGA3I4.7.3/4.7.3.1 and SGA6VI3.3, FGA232 Theorem3.1 pp.148--149, FGA236 Corollary1.7 p.229/2.3 p.232, SGA7I VII2.7/2.9 pp.169--174 andVIII7.4(b) pp.310--311, and EGAIV.4 16.11.2--3(ii) pp.54--55. No motivic projector, algebraic cycle nonvanishing, ell-adic weight or Kummer RH consequence is inferred from this finite representation statement.

\subsection{References and exact source locators}\label{references-and-exact-source-locators-1}

Gow, R., \& Vinroot, C. R. (2007). \emph{Extending real-valued characters of finite general linear and unitary groups on elements related to regular unipotents} (arXiv:0704.2022v1). Theorem3.2 p.9, Lemma3.1(i) p.10, Theorem4.3 p.12. \url{https://arxiv.org/pdf/0704.2022}.

Srinivasan, B., \& Vinroot, C. R. (2019). \emph{Galois group action and Jordan decomposition of characters of finite reductive groups with connected center}. \url{https://doi.org/10.1016/j.jalgebra.2019.04.025}. Lemma3.2, author PDF pp.7--8; Theorem4.1 pp.12--13. \url{https://www.math.wm.edu/~vinroot/JordanGaloisRefereed.pdf}.

The local PDFs, extracted texts and retrieval hashes are in \nolinkurl{research/finite_gl_refs/}. The exact archive equations are StatementT.lean9--29, \nolinkurl{DataN.lean960–966/996–1007,} LocalToGlobal/SiegelLatticeTheory/FlagStabilizerTensors/ RowLevelMatch.lean39--47 and RowDict.lean82. R025 supplies the canonical residue action and exact group orders. No Lean execution or archive soundness check is involved.

\clearpage

\section{AM-DO-042: full literal goodTuple existence at the actual auxiliary primes}\label{am-do-042-full-literal-goodtuple-existence-at-the-actual-auxiliary-primes}

Author: GPT-6 Astra Pro, ALP\_MACHINE /root/definite\_oscillator. 22 November2026 UTC. Inherited reasoning effort is not separately exposed.

The complete finite predicate goodTuple has solutions at all four current G primes, for ranks43 and169 and the actual42 or168 selected places. This includes both prescribed character values and their inverses, all prime-to-part conditions, and both factor-degree conditions. No factorization of g\^{}d−1 is needed.

The construction is a rigorous counting construction, with an explicit choice rule from finite sets. It proves the actual finite predicate, not merely squarefreeness or regular semisimplicity. At two distinguished places use the irreducible-factor patterns \[
 \begin{array}{c|c|c}
 N&v_0&v_1\\ \hline
 43&(43)&(37,6)\\
 169&(169)&(163,6).
 \end{array}                                             \tag{DO42.1}
\] Use the first pattern at all remaining places. A conditional corollary for the literal e0 formula is stated separately: when eSer is the genuine finite-GL Lusztig-series projector, the resulting e0 is nonzero. The listed abstract eSer axioms alone do not assert that every valid polynomial has a nonzero series projector.

\subsection{1. Exact source predicate and actual parameters}\label{exact-source-predicate-and-actual-parameters}

The source is the recovered archive at commit 7cbf7a0c32fca1badf05931879a37c3f8a2b03b4.

\begin{itemize}
\tightlist
\item
  \nolinkurl{Formalism/CyclotomicNorms/StatementI.lean10–12} defines Exc(g,d,l) by l prime and g\^{}d\textless l\^{}(4 v\_l(g\^{}d−1)), and defines BB(N,g) to contain g and all such primes for1\textless=d\textless=N.
\item
  \nolinkurl{Formalism/CyclotomicNorms/StatementT.lean9–16} defines primeToPart and rootsBar; lines19--29 define IsGood.
\item
  \nolinkurl{Formalism/SharpRestriction/LevelLetterCocycles/DataN.lean996–999} requires IsGood for chiT and also for chiT inverse. Lines1001--1007 define the exact coefficientwise group-algebra sum e0 over good tuples.
\end{itemize}

In ordinary notation, IsGood(N,g,(f\_v),(chi\_vj)) requires all of:

\begin{enumerate}
\def\labelenumi{\arabic{enumi}.}
\tightlist
\item
  Every f\_v is monic of degree N with nonzero constant term.
\item
  Every f\_v is squarefree.
\item
  For each prime l outside BB(N,g), primeToPart\_l is injective on the complete root set of each f\_v.
\item
  For every prime l different from g, every root x of f\_v, and every j, primeToPart\_l(x) differs from primeToPart\_l(chi\_vj).
\item
  For every1\textless=a\textless=N/2, some f\_v has no factor of degree a.
\item
  At some v, f\_v has an irreducible factor of prime degree p with N\textless2p and p+3\textless=N.
\end{enumerate}

For goodTuple, condition4 must hold also for chi\_vj inverse. Since m=2, at each place it is enough to avoid the set \[
 C_v=\{\chi_{v,0},\chi_{v,1},
                 \chi_{v,0}^{-1},\chi_{v,1}^{-1}\}\subset\mathbf F_g^\times,
 \qquad |C_v|\le4.                               \tag{DO42.2}
\] No assumption that these four values are distinct is needed.

The actual primes and ranks are \[
 \begin{array}{c|r|r|r|r}
 \text{host}&f&N&g_0&g_1\\ \hline
 J^{14}&42&43&4\,615\,409&5\,276\,377\\
 J^{56}&168&169&22\,285\,932\,761&64\,072\,175\,441 .
 \end{array}                                             \tag{DO42.3}
\] They are the proved K031/K032 primes, not new auxiliary choices. They all satisfy \[
                  2^{22}<g<2^{36},\qquad812\mid g-1,
                  \qquad812=2^2\cdot7\cdot29.    \tag{DO42.4}
\] The companion certificate rechecks their primality by exact trial division and these elementary inequalities. The actual selected place sets have size f and hence at least two elements.

\subsection{2. What the source primeToPart does on a finite cyclic group}\label{what-the-source-primetopart-does-on-a-finite-cyclic-group}

Let x have finite order m, and let l be prime. Write m=l\^{}a m′ with l not dividing m′. The exponent in the source primeToPart definition is congruent to0 modulo l\^{}a and to1 modulo m′. Hence \[
 x_{\ell'}:=\operatorname{primeToPart}_\ell(x)
\] is precisely the prime-to-l component of x. The case m′=1 gives1, as it should. If x belongs to a fixed cyclic group C\_n, this operation agrees with the single cyclic-group projection C\_n→C\_(n/l\^{}v\_l(n)): the two exponents have the same residues modulo the order of x. It is therefore a homomorphism, commutes with Frobenius, and has kernel of order \[
                             L_\ell(n)=\ell^{v_\ell(n)}.
                                                        \tag{DO42.5}
\] For l not dividing n it is the identity. This verifies that the counting below concerns the literal source operation, not the different operation x↦x\^{}(l\^{}a) without its inverse exponent.

For n=g\^{}d−1 and s dividing d, the number of x in F\_(g\^{}d)\^{}times whose prime-to-l component lies in F\_(g\^{}s)\^{}times is \[
 L_\ell(n)\gcd\!\left(n/L_\ell(n),g^s-1\right)
                         \le L_\ell(n)(g^s-1).    \tag{DO42.6}
\] This follows by intersecting the image subgroup with the unique subgroup of order g\^{}s−1 and multiplying by the kernel size. For any fixed c in F\_g\^{}times, the number of solutions of x\_(l′)=c\_(l′) is exactly L\_l(n).

\subsection{3. A uniform bound for all character collisions}\label{a-uniform-bound-for-all-character-collisions}

For n\textgreater=2 define \[
                   S(n)=\frac{\sum_{\ell\mid n,\ \ell\ {\rm prime}}
                                      \ell^{v_\ell(n)}}n.
                                                        \tag{DO42.7}
\] This quantity is nonincreasing under divisibility. To prove it, increase one existing prime-power factor A to A u: its own contribution to S stays fixed and every other contribution is divided by u. If a new prime-power factor B is added, then \[
              S(nB)=S(n)/B+1/n\le S(n),
                                                        \tag{DO42.8}
\] because the sum of the existing prime-power factors is at least2 and B/(B−1)\textless=2. Iterating proves the claim.

Since812 divides g−1 and hence n=g\^{}d−1, \[
                     S(n)\le S(812)=\frac{4+7+29}{812}
                                         =\frac{10}{203}. \tag{DO42.9}
\] Fix any set C in F\_g\^{}times of cardinality at most4. By(DO42.5), the union of all character-collision sets for l dividing n has size at most \[
       |C|\sum_{\ell\mid n}L_\ell(n)\le\frac{40}{203}\,n.
                                                        \tag{DO42.10}
\] This bound includes all prime divisors, including exceptional primes in BB; it does not merely bound the nonexceptional collisions. For l not dividing n, a collision is x=c and is excluded by the full-degree condition below when d\textgreater=2. Thus the infinitely quantified character condition has been reduced exactly to this finite union.

\subsection{4. Full degree survives every nonexceptional projection}\label{full-degree-survives-every-nonexceptional-projection}

Fix N\textless=169 and6\textless=d\textless=N. Call x in F\_(g\^{}d)\^{}times strongly good for (N,g,C,d) when:

\begin{itemize}
\tightlist
\item
  x has degree d over F\_g;
\item
  for every prime l outside BB(N,g), x\_(l′) has degree d over F\_g;
\item
  for every prime l different from g and every c in C, x\_(l′) differs from c\_(l′).
\end{itemize}

The set is Frobenius stable. In particular all conjugates of one strongly good root satisfy the same conditions.

The number of elements lacking full degree is at most \[
              \sum_{\substack{s\mid d\\s<d}}(g^s-1)
                           \le d\,g^{d/2}.       \tag{DO42.11}
\] For a prime l outside BB(N,g), the literal definition of BB and d\textless=N imply \[
              L_\ell(g^d-1)^4\le g^d,\qquad
              L_\ell(g^d-1)\le g^{d/4}.          \tag{DO42.12}
\] If l does not divide g\^{}d−1, the projection is the identity, so its degree loss is already covered by(DO42.11). For l dividing g\^{}d−1, equation(DO42.6) bounds the degree-loss set by g\^{}(d/4) times the sum in(DO42.11). The number of distinct prime divisors of g\^{}d−1 is at most log\_2(g\^{}d−1)\textless36d. Thus all additional nonexceptional degree-loss sets have total size at most \[
                              36d^2 g^{3d/4}.    \tag{DO42.13}
\] Double counting only enlarges this upper bound and is harmless.

Divide(DO42.10--13) by n=g\^{}d−1. Since g\textgreater2\^{}22 and6\textless=d\textless=169, the fraction of elements which are not strongly good is at most \[
 B_{\rm bad}:=\frac{40}{203}
   +\frac{169\,2^{-66}+36\cdot169^2\,2^{-33}}{1-2^{-132}}
                         <\frac15.              \tag{DO42.14}
\] All inequalities here are strict in the available margin. The exact rational number is \[
 B_{\rm bad}=
 \frac{2394648453263414000006332163567196866312}
      {12145462942408880542077370603718496162735}
 \simeq0.19716403274361063.
                                                        \tag{DO42.15}
\] The displayed decimal is diagnostic only; the strict comparison with1/5 is checked by integer cross multiplication. It follows that \[
 \boxed{\quad
  \#\{\text{strongly good }x\in\mathbf F_{g^d}^\times\}
                          >\frac45(g^d-1).\quad} \tag{DO42.16}
\] This proves existence without finding a factor of g\^{}d−1 or an exceptional prime. The estimates use only the exact definition of BB.

\subsection{5. Construction of the actual good tuple}\label{construction-of-the-actual-good-tuple}

For a strongly good x of degree d let \[
                      P_x(X)=\prod_{i=0}^{d-1}(X-x^{g^i}).
                                                        \tag{DO42.17}
\] This is a monic irreducible polynomial in F\_g{[}X{]} of degree d with nonzero constant term. Its roots satisfy all character exclusions. For l outside BB(N,g), their prime-to-l parts are distinct: the projection x\_(l′) has full Frobenius degree d, so its d conjugates are distinct.

Choose distinct v0,v1 in each actual place set. At every place other than v1 choose a strongly good degree-N root for its own C\_v and put f\_v=P\_x. At v1 choose strongly good roots of degrees p and6 for the same C\_v1, where p=37 if N=43 and p=163 if N=169; put \[
                             f_{v_1}=P_xP_y.     \tag{DO42.18}
\] The two factors have different degrees, so they are distinct and the product is squarefree. Their roots have no cross-factor collision after a nonexceptional prime-to-l projection: those projections have degrees p and6, respectively, and two equal algebraic elements cannot have different degrees over F\_g. This is why preserving full degree in §4 is stronger than just checking distinct roots within one orbit.

We now verify all source clauses.

\begin{enumerate}
\def\labelenumi{\arabic{enumi}.}
\tightlist
\item
  All f\_v are monic of degree N with nonzero constant term.
\item
  They are squarefree: irreducibles over a finite field are separable, and the two factors at v1 have different degrees.
\item
  Within each factor, full projection degree gives distinct projected conjugates. Between the two factors, their different preserved degrees exclude equality. Thus projection is injective on the complete root set for every l outside BB(N,g).
\item
  Each root avoids all projections of C\_v for every prime l different from g. By(DO42.2), this proves the clauses for both chi and inverse.
\item
  The polynomial at v0 is irreducible of degree N. Its nonzero factors have degree0 or N, so it has no factor of any degree1\textless=a\textless=N/2. The same v0 witnesses every such a.
\item
  At v1 the factor P\_x has prime degree p.~The exact inequalities are 43\textless2·37 and37+3\textless=43, or169\textless2·163 and163+3\textless=169.
\end{enumerate}

The construction is independent at the two rational G primes and therefore gives the literal DataN.goodTuple. It works for every choice of the actual two chiT values at each place; their equality or inverse conventions require no guessed generator.

To make the choice rule definite, fix any concrete polynomial model of each finite extension, order its coefficient vectors using 0,\ldots,g−1, and select the first strongly good root. These are nonempty finite sets by(DO42.16). This rule specifies an actual tuple. The result is a counting/existence construction, not a claim that its very large coefficient vectors were enumerated.

There is also an explicit lower count of available monic irreducibles. Frobenius stability and full degree give at least \[
                 I_d(g):=\left\lfloor\frac{4(g^d-1)}{5d}\right\rfloor+1
                                                        \tag{DO42.19}
\] choices of P\_x at each place. Keeping v0,v1 fixed, the construction therefore produces at least \[
                       I_N(g)^{f-1}I_p(g)I_6(g)  \tag{DO42.20}
\] different good polynomial tuples at one rational G prime. Products at v1 are distinct by unique factorization and the different degrees.

\subsection{6. Exact relation to the prescribed e0}\label{exact-relation-to-the-prescribed-e0}

Let G\_fin be the actual product of GL\_N(F\_(g\_k)) over the two G place sets. Suppose eSer\_k(f) is the genuine orthogonal central projector for the rational Lusztig series with semisimple characteristic polynomial f.~A squarefree f has regular semisimple parameter. The standard finite-GL realization has a nonzero series for this parameter; in fact its centralizer is a torus and the series has one irreducible character. This last representation-theoretic input is separate from the elementary counting theorem above and is to be used with its independently verified primary anchor in the finite-model adapter.

For the constructed tuple f, set \[
                        E_f=\bigotimes_{k,v}eSer_k(f_{k,v}).
                                                        \tag{DO42.21}
\] Every factor is nonzero in that genuine realization, so E\_f is nonzero. The source map a↦monicOf(a) is a bijection between coefficient vectors of length N and monic degree-N polynomials. Consequently the literal coefficientwise e0 sum in DataN is exactly the finite sum of E\_h over all good tuples h, with no extra multiplicity or coefficient. Orthogonality gives \[
                              E_f e_0=E_f,\qquad e_0\ne0. \tag{DO42.22}
\] This implication is purely algebraic once the genuine nonzero series projectors have been specified. For their unitary realizations each projector is self-adjoint, so e0 is also an orthogonal projector. No cancellation or unspecified sign in a virtual character can make it zero.

The abstract eSer\_idem/central/orth/zero/sum axioms in DataN960--966 alone do not force nonzero eSer(f) for every valid f.~One could assign zero to a selected good label while still satisfying those axioms. Thus(DO42.22) is explicitly a corollary for the genuine finite-GL series realization; it is not asserted solely from that abstract list. Likewise e0 nonzero on the finite regular representation does not by itself prove nonzero action on a given automorphic cusp space. The actual finite branching and equivariant jet detector are separate steps, with their own source levels and character signs.

\subsection{7. Scope, normalization and verification}\label{scope-normalization-and-verification}

This theorem proves the complete actual finite goodTuple predicate on the current hosts and gives its exact relation to the literal e0 formula. It changes neither the G primes nor the two stage characters. It introduces no auxiliary projector and no relaxed prime-to-part condition. It does not assert BB(N,g0) intersect BB(N,g1) is empty: that is a distinct NTypes pin and an additional hypothesis of the source irreducibility criterion. Nor does it instantiate the full NTypes/Host/SizeData2 record.

The count itself has no Fourier or Haar normalization. Realizing E\_f in a group algebra retains the usual finite representation dimensions and group-order denominators; this argument does not make them11-units or bound an arithmetic right inverse. Nonzero finite e0, its occurrence in the actual bottom automorphic space, effective jet separation in K042's bounded window, and a nonzero terminal PtilFF value are different assertions.

All new counting and finite-field claims are proved in §§2--5. The exact diagnostic checks the four actual primes and ranges, the rational bad-density comparison, all relevant factor-degree patterns, source CRT prime-to projection on small cyclic examples, the divisibility monotonicity and collision fibers, and finite-field Frobenius degree counts in small examples. Such examples check the algebra; the uniform proof uses the displayed inequalities rather than extrapolating from them. The genuine-series primary theorem in §6 is an explicitly separate input. No Lean execution, enormous polynomial enumeration or global automorphic nonvanishing is claimed.

\subsection{Author verification record}\label{author-verification-record}

Whole proof and complete diagnostic script read on 2026-11-22T03:31:58.687640+00:00. All six literal IsGood clauses, both character signs, ambient/source CRT projection equality, divisibility-monotone collision sum, complete nonexceptional degree losses, cross-factor collisions, and the exact e0 scope were checked. The certificate passed 85,932 literal projection/Frobenius element checks, 376 exact fibers, 4,536 monotonicity checks, 18 field-degree counts and all four actual prime/range checks. Author whole-proof PASS. This records an actual read and computation; no external independent review is asserted.

\clearpage

\section{AM-DO-043: prescribed CM central character inside the literal goodTuple sum}\label{am-do-043-prescribed-cm-central-character-inside-the-literal-goodtuple-sum}

Author: GPT-6 Astra Pro, ALP\_MACHINE /root/definite\_oscillator. 22 November 2026 UTC. Inherited reasoning effort is not separately exposed.

For either current host and any prescribed character of the diagonal CM root group mu\_28, the literal goodTuple sum contains a nonzero genuine finite-GL series constituent with that central character. The G primes, stage characters and all six IsGood conditions remain fixed. The repair varies two previously free irreducible degree-N factors; it does not scalar-twist an already chosen polynomial and assume goodness survives.

This solves a finite central-character compatibility condition. It does not by itself construct a global cusp form, prove that a prescribed automorphic constituent occurs, or give a nonzero bottom or terminal period.

\subsection{1. Settled counting input and the free places}\label{settled-counting-input-and-the-free-places}

Use AM-DO-042 with N=43, f=42 or N=169, f=168 and the actual four G primes. For each place v, let S\_v be its strongly good subset of F\_(g\^{}N)\^{}times. All four prescribed values chi\_vj and their inverses are included in its exclusion set. The complete proved count is

\[
 |S_v|>\frac45(g^N-1).                         \tag{DO43.1}
\]

Every x in S\_v and every nonexceptional prime-to-part has degree N; the set is Frobenius stable. It supplies monic irreducibles of degree N with all literal projection and character conditions. At one different place retain the factors (37,6) or (163,6), respectively, as in DO042.

Choose two full-degree places v and w in one G half, disjoint from that mixed-factor place. The actual half has at least42 places, so this is possible. Fix all parameters at every other place and both rational G primes. The factor-degree witnesses in DO042 remain valid for every choice at v and w: there are still irreducible degree-N places, and the prime-degree mixed place is unchanged.

\subsection{2. Exact finite central-character dictionary}\label{exact-finite-central-character-dictionary}

Fix a generator t of F\_(g\^{}d)\^{}times, put n\_d=g\^{}d−1 and

\[
 t_F=t^{n_d/(g-1)}.
\]

Choose compatible complex primitive roots of unity. For a semisimple parameter alpha=t\^{}b, use the dual torus character theta\_alpha defined by

\[
 \theta_\alpha(t)=\zeta_{n_d}^{\,b},\qquad
 N_{\mathbf F_{g^d}/\mathbf F_g}(\alpha)=t_F^b.
                                                        \tag{DO43.2}
\]

The generator choices can be made compatible under norm maps: take L to be a common multiple of the finitely many degrees in use, choose a generator of F\_(g\^{}L)\^{}times, and take its norm to each subfield. Those norms are generators, and norm transitivity gives compatibility. These choices specify the finite torus duality used for the genuine eSer realization. They are not an arbitrary orbit bijection unrelated to the source polynomial labels. Frobenius on both sides is exponent b↦gb.

For an irreducible factor of degree d the associated torus is F\_(g\^{}d)\^{}times, and central c in F\_g\^{}times acts through theta\_alpha(c). Here is the exact sign check. Deligne and Lusztig (1976), Theorem4.2, printed p.124, gives the character at a central semisimple element z by substituting s=z and u=1. Since its connected centralizer is G, all torus conjugates contain z and the theta factor is theta(z). The remaining sum is precisely the same formula at1. Thus

\[
 R_T^\theta(z)=\theta(z)R_T^\theta(1).          \tag{DO43.3}
\]

The possible overall sign of the virtual Deligne--Lusztig character cancels. For the regular parameter it gives the actual irreducible central character. For a product of distinct factors, the torus is their product and the central character is the product of the characters in (DO43.2). Its base-field parameter is the product of the norms of the factor parameters, namely the semisimple determinant, with no unrecorded inverse. The genuine nonzero singleton and genericity are the standard finite-GL realization proved in AM-R026; they are also visible in Gow--Vinroot's Green parametrization, Theorem3.2, and Gelfand--Graev decomposition, Theorem4.3.

Because 28 divides g−1, for the element c=t\_F\^{}((g−1)j/28) one gets

\[
 \theta_\alpha(c)=\zeta_{28}^{bj}.              \tag{DO43.4}
\]

Thus the map from roots to restricted central characters is the quotient

\[
 \mathbf Z/n_d\mathbf Z\longrightarrow\mathbf Z/28\mathbf Z,
 \qquad b\longmapsto b\pmod {28}.               \tag{DO43.5}
\]

It has exactly n\_d/28 elements in each fibre. It is Frobenius invariant, since g is1 modulo28. Equivalently, take the norm into F\_g\^{}times and then its quotient by 28th powers. Each actual embedding of the fixed CM mu\_28 into the residue field changes (DO43.5) by a unit automorphism of Z/28, because it preserves the order of a primitive28th root. The opposite duality or contragredient convention would multiply every exponent by−1, also a unit; it would change the named target character, but not the existence or counting assertion below. In this report the direct theta sign is fixed by (DO43.3).

\subsection{3. Two places realize every aggregate character}\label{two-places-realize-every-aggregate-character}

Let A\_v be the set of mu\_28 characters realized by roots in S\_v, using the actual residue embedding. Since all28 fibres have equal size, (DO43.1) implies

\[
 |A_v|>\frac45\,28=22.4,
 \qquad |A_v|\ge23.                             \tag{DO43.6}
\]

The same holds at w, even if their prescribed chi values differ. For any c in Z/28,

\[
 |A_v\cap(c-A_w)|\ge|A_v|+|A_w|-28\ge18.        \tag{DO43.7}
\]

Hence A\_v+A\_w is all of Z/28. Choose c to be the desired global central exponent minus the already fixed contribution of every other place. Select a in the nonempty intersection in (DO43.7), then a good root at v with exponent a and one at w with exponent c−a. The resulting monic polynomials have the required aggregate central character. Every original IsGood condition remains true by DO042, including both signs of chi, all exceptional-prime character exclusions, and the cross-factor projection conditions at the mixed place.

The argument also gives a uniform lower count of choices. Put n=g\^{}N−1, M=n/28, and let a\_r,b\_r be the numbers of strongly good roots in the corresponding character fibres at v,w. Thus 0≤a\_r,b\_r≤M. Write a\_r=M(1-u\_r), b\_r=M(1-v\_r), with u\_r,v\_r in{[}0,1{]}. For every c,

\[
 \begin{split}
 \sum_r a_r b_{c-r}
 &=M^2\sum_r(1-u_r)(1-v_{c-r})\\
 &\ge M^2\left(28-\sum_r u_r-\sum_rv_r\right)\\
 &=\frac{n^2}{28}
       \left(\frac{|S_v|}{n}+\frac{|S_w|}{n}-1\right)
 >\frac{3n^2}{140}.
 \end{split}                                             \tag{DO43.8}
\]

Each polynomial corresponds to exactly N good roots, and the character is constant on their Frobenius orbit. Consequently, with all other parameters fixed, every target central character occurs for more than

\[
                     \frac{3(g^N-1)^2}{140N^2}            \tag{DO43.9}
\]

ordered pairs of irreducible polynomial factors at v and w. This is a finite counting existence result, not a claim that these enormous polynomial lists were enumerated.

\subsection{4. Literal e0 and the extra line character}\label{literal-e0-and-the-extra-line-character}

Let Z be the diagonal CM mu\_28 in the actual product of finite GL\_N groups. For any Omega in the character group of Z, choose the good tuple just constructed, and let E\_f be its genuine nonzero product series projector. Let

\[
 p_\Omega=\frac1{28}\sum_{z\in Z}\Omega(z)^{-1}z.
                                                        \tag{DO43.10}
\]

This is the orthogonal projector to the Omega-isotypic part. The exact source e0 includes E\_f with coefficient1 and all distinct tuple projectors are orthogonal. Therefore

\[
 p_\Omega E_f=E_f,
 \qquad p_\Omega e_0E_f=E_f\ne0.                \tag{DO43.11}
\]

Thus the actual goodTuple projector has a nonzero block for every prescribed CM central character. The auxiliary central average has operator norm1 and denominator28, an11-unit. This does not remove the group-order denominators already present in E\_f or e0; AM-R026 retains them. Nor does it say that these blocks act nontrivially on a particular global automorphic representation.

The negative-line compatibility in a finite branching construction has an exact elementary form. Let Pi be an irreducible GL\_(a+1)(F\_g) representation, and let a GL\_a constituent sigma occur on the upper block H. For A in GL\_a and c in F\_g\^{}times,

\[
 \operatorname{diag}(A,c)
 =cI_{a+1}\operatorname{diag}(c^{-1}A,1).
\]

Hence the action on this sigma-isotypic space extends to H times F\_g\^{}times as

\[
 \sigma(A)\,\frac{\omega_\Pi(c)}{\omega_\sigma(c)}.
                                                        \tag{DO43.12}
\]

The same formula holds on every multiplicity copy. If all other finite and archimedean factors prescribe an aggregate character for the extra line on CM mu\_28, the choice in §3 can meet it exactly, after multiplying by the already fixed initial omega\_sigma factors. No claim is made that arbitrary noncentral finite stabilizers are controlled by this central calculation. The separate finite mirabolic branching theorem and actual packet stabilizer identification must still be used.

\subsection{5. Scope and references}\label{scope-and-references}

The new proof is the two-place prescription, with exact torus-duality and central-character formulas. It uses the fully checked root count DO042 and the genuine finite-series realization R026. It leaves BB-set disjointness and the full source NTypes record separate. It also leaves the global equivariant jet-surjectivity threshold and its norm and arithmetic right inverse separate. In particular no argument from finite e0 nonzero to a nonzero prescribed terminal value is made.

Deligne, P., \& Lusztig, G. (1976). Representations of reductive groups over finite fields. \emph{Annals of Mathematics, 103}(1),103--161. Theorem4.2, printed p.124; §5 specifies the compatible finite torus duality choices. \url{https://publications.ias.edu/sites/default/files/Number27.pdf}.

Gow, R., \& Vinroot, C. R. (2007). \emph{Extending real-valued characters of finite general linear and unitary groups on elements related to regular unipotents} (arXiv:0704.2022v1). Theorem3.2, p.9; Theorem4.3, p.12. \url{https://arxiv.org/abs/0704.2022}.

The companion certificate checks the exact quotient-fibre arithmetic, Frobenius invariance, central exponent and norm formula, the finite two-set and weighted-pair inequalities, all four actual prime congruences, and the exact positive lower counts. These are diagnostics for the displayed proof; no global finite-series or cusp computation is claimed.

\subsection{Author whole-proof check}\label{author-whole-proof-check}

Whole proof DO43.1--12 and the complete certificate read on 2026-11-22T03:37:06.722456+00:00. Directly read Deligne--Lusztig Theorem4.2 at printed p.124 and Gow--Vinroot Theorems3.2/4.3. Checked compatible finite torus generators, exact central sign, equal fibres and Frobenius invariance, the two-place and weighted counts, all retained goodTuple conditions, literal e0 block and extra-line character. Corrected a typesetting comma and clarified that the determinant is the product of factor norms. Author PASS; no global occurrence or terminal value inferred.

\clearpage

\section{AM-K044 --- an exact larger gap interval inside the same weight windows}\label{am-k044-an-exact-larger-gap-interval-inside-the-same-weight-windows}

Author: GPT-6 Astra Pro, ALP\_MACHINE /root/kummer\_extension, 22 November 2026 UTC. Inherited reasoning effort is not separately exposed.

The small value of ell selected in R022 is not forced by its weight windows. This finding determines an exact interval of gap parameters for which its entire binary-word construction, the K042 canonical-twist interval, and the named HC/Scale/weight conditions remain valid. It gives an explicit larger example and its quantitative source-counting margins.

Changing ell changes the coefficient weights and the bottom sheaf. It does not prove an effective interpolation theorem for the original sheaf, or turn an existential cusp-counting constant into a known number. The finite levels, CM hosts and Scale remain fixed. The source Host/Fields obstructions and all unproved global norm or image assertions remain.

\subsection{1. Fixed data and the interval}\label{fixed-data-and-the-interval}

Retain the exact R022 values

\[
 A=2^{960},\quad B=2^{955},\quad
 E_T=28\left\lceil(B+e^{336})/28\right\rceil,\quad
 E_H=28\left\lceil e^{343}/28\right\rceil,\quad
 E_L=28\left\lceil e^{336}/(28n_F)\right\rceil,
 \tag{EG1}
\]

where n\_F=f/2, and the two rank choices are

\begin{longtable}[]{@{}lrrrrr@{}}
\toprule\noalign{}
Host & a=f & t & d=t+1 & bottom rank n=a+1 & S=a+1+2t \\
\midrule\noalign{}
\endhead
\bottomrule\noalign{}
\endlastfoot
J14 & 42 & 1759 & 1760 & 43 & 3561 \\
J56 & 168 & 1633 & 1634 & 169 & 3435 \\
\end{longtable}

The symbol n in this report denotes the bottom matrix rank, while n\_F denotes the old half-degree. Let R=E\_T−B and put

\[
 \ell_- =28\left(\left\lfloor S^3/56\right\rfloor+1\right),
 \qquad
 \ell_+=28\left\lfloor
               \frac{R-a/2-3}{28}\right\rfloor .
 \tag{EG2}
\]

The allowed parameters are

\[
 \ell\in28\mathbb Z,\quad \ell_-\le\ell\le\ell_+,
 \qquad s\in28\mathbb Z_{\ge0},\quad q=(a+1)s\le\ell.
 \tag{EG3}
\]

For each ell set Del=ell, Sc=4a ell, etaSz=4ell, Cc=2 and use exactly R022's word formulas with this ell. Replace the two e and b entries by

\[
 (e_0,e_1)=(E,E+2\ell+q),\qquad
 (b_0,b_1)=(E-\ell,E+\ell+q).
 \tag{EG4}
\]

The remaining u and v formulas are unchanged:

\[
 v_i=4\ell(a-1-i),\qquad
 u_i=v_{i-1}-\tfrac12-\ell\omega_{i-1}\quad(1\le i\le a-2),
 \tag{EG5}
\]

with top endpoints A+(a−1)/2 and −B−(a−1)/2, and other endpoints v\_0+4a ell+1/2 and 2ell+1/2.

The exact lower endpoints are22578023748 and20265168952. At the upper endpoints R−ell\_+ is40 for J14 and96 for J56. Their full integers are in the reproducible certificate; in both cases ell\_+\textless2\^{}485.

\subsubsection{Theorem EG6}\label{theorem-eg6}

Every parameter in EG3, for every word and every role assignment, satisfies the named conditions proved in R022 and K042: HCParam, Centred, strict NCGaps, the unchanged numerical Hyp3 rank bound, aT\_win, bT\_ge, top\_u, top\_b, b\_le, Sc\_eq, all three role windows, cloud\_gap, cloud\_v, cloud\_u and be\_close. Both selected partitions are dominant and positive, and the same DO040 primes satisfy the four numerical Sizes clauses.

If at least one top label is required and q=0 is included, the endpoints in EG2 are sharp for this conjunction of conditions: the next lower 28-multiple fails NCGaps, and the next upper28-multiple fails top\_b. If there is no top label, EG3 is still a valid interval, but no optimal upper endpoint is claimed.

\subsection{2. Proof of the weight and gap assertions}\label{proof-of-the-weight-and-gap-assertions}

The strict source condition is

\[
 S^3<|x-y|\quad\hbox{for every opposite-sign pair in Dpar}.
 \tag{EG7}
\]

Its exact locator is \nolinkurl{SharpRestriction/Divisibility/FunctionModelRallisFromInputs.lean20–22;} \nolinkurl{CuspSpanNormBounds/CayleySignature.lean110–120} additionally displays the 1801 rank premise. At stage zero, the unique negative-sign entry is −b\_1. The positive-sign entries are v\_i and −b\_0. Their closest distance is exactly

\[
 \min\{b_1-b_0,\min_i(v_i+b_1)\}
     =\min\{2\ell+q,E+5\ell+q\}=2\ell+q.
 \tag{EG8}
\]

At the other two stages all signs are positive, exactly as in R022; the odd-stage common shift changes no distance. Hence every q in EG3 satisfies EG7 exactly when 2ell\textgreater S\^{}3 at its smallest value q=0. This gives ell\_- and its sharp preceding-step failure. The rank a+1+t equals1802 for both hosts, independently of ell.

Every interior interlacing gap in EG5 is either1/2 or ell+1/2 on one side, and either4ell−1/2 or3ell−1/2 on the other. Non-top endpoint gaps are Sc+1/2 and2ell−1/2. At a top label, u\_0\textgreater v\_0 because ell≤ell\_+\textless2\^{}485 and A=2\^{}960. Its last HC inequality is

\[
            u_{a-1}+b_0=R-\ell-(a-1)/2>3.
 \tag{EG9}
\]

This is implied by the strict source top\_b condition, which is itself exactly R−ell\textgreater(a−1)/2+3. Since ell and R are integers and a is even, that condition is equivalent to ell≤R−a/2−3. Taking the largest 28-multiple proves EG2's upper endpoint and the next-step failure. At a non-top label,

\[
                  u_{a-1}+b_0=E+\ell+1/2>0.
\]

In particular b\_0 is permitted to become negative for some thin roles; the source HCParam does not require b\_0≥0. No positivity condition on b\_0 has been silently added. All b/e gaps are at least ell, which exceeds (t+3)/2 at ell\_-. Parity and integrality are immediate.

For centredness, all u\_i+e\_j and v\_i+e\_j are positive. The only top cancellation is u\_last+e\_0=R−(a−1)/2, independent of ell. The b/e distances are ell, ell+q or3ell+q. The nearest point in an other-e block has distance at least2ell+q−(t−1)/2. All exceed the required (t+3)/2. These are the complete clauses in \nolinkurl{LocalRingCardinality/DataTheta.lean550–554,} including both e indices.

The Scale has not changed, so aT\_win, bT\_ge and top\_u are carried literally from R022. Each rounded E is at most its centre plus28. Thus the largest e value lies at most3ell+28 above that centre, which is at most4ell since ell≥ell\_-\textgreater28. Every b/e difference is at most3ell+q≤4ell. This proves the role windows and be\_close. All v and interior u entries lie in(0,Sc), and the non-top first gap is Sc+1/2∈{[}Sc,2Sc{]}. This proves the remaining cloud clauses.

The bound b\_le holds with ample margin. The inequalities ell\_+\textless2\^{}485, E\_H\textless2\^{}495 and E\_T\textless B+2\^{}485 imply b\_1\textless A for all roles, whereas R022 proves e\^{}686\textgreater2\^{}980\textgreater A. The same bounds give \textbar v\_i\textbar\textless A and \textbar b\_j\textbar\textless A; these will also be used below.

The selected first partition is nu\_0,i=u\_i−(a−1)/2+i+E−(a+d)/2. Its last entry is R−(a+d)/2 at top labels and E+2ell−d/2 elsewhere, both positive. Dominance follows directly from the adjacent interlacing gaps. The second partition is exactly

\[
       \nu_1=(\nu_0,0^d)+(2\ell+q-d)\mathbf1_{a+d}.
 \tag{EG10}
\]

Its new constant is positive throughout EG3. The largest top spread is A+E\_T+2ell+q+(a+d)/2\textless2A. At a non-top label it is (8a−4)ell+1−a/2+E+2ell+q+(a+d)/2\textless2A. Earlier spreads are smaller. The unchanged rank/G requirements and DO040's actual T\textgreater16A therefore retain all four numerical Sizes clauses. This is not a claim of full SizeData2. QED.

The companion certificate independently reconstructs the exponential roundings, and checks all coordinate-bit possibilities and all endpoints of these affine inequalities. The two extreme Hamming weights suffice for their degree extrema because the total degree is affine in W, with coefficient−ell. Checking ell endpoints and q=0,ell covers the entire triangle in EG3; q=ell is a harmless larger set than the actual integer progression. Absolute values above have already had their signs resolved in the proof, so an untested interior zero cannot invalidate this argument.

\subsection{3. A concrete choice and the exact canonical-twist interval}\label{a-concrete-choice-and-the-exact-canonical-twist-interval}

Choose, for either host,

\[
                    \ell_* =28\lceil2^{400}/28\rceil.
 \tag{EG11}
\]

It lies strictly between ell\_- and ell\_+. For any fixed ell in EG3, the exact allowed canonical-twist endpoint is still

\[
       s_{\max}(\ell)=28\left\lfloor
                           \frac{\ell}{28(a+1)}\right\rfloor.
 \tag{EG12}
\]

The next step fails be\_close because e\_1−b\_0=3ell+(a+1)s. The original K042 values are the special case of its much smaller ell. This enlargement changes both the initial selected type and the minimal bottom fibre. Consequently its much longer interval is not an effective Serre threshold for the old fixed packet or old sheaf.

K040 supplies the actual new stage characters: all E and2ell+q are multiples28, their relative angular type is constant over the labels, and its construction allows these arbitrary exponents in the same finite value field. Ramification at G and the absence of29/11 ramification persist. All CM varieties, finite level subgroups, sigma gauge and prime-to11 local integral-section premises of K039 are unchanged.

The analytic transport argument of K041 uses the labelwise two compact ties, EG10, positive2ell+q−d, unitarity and the same G/T/Q conditions; it does not use the original numerical value of ell. It therefore retains its stated full-theta nonzero scope. The arithmetic K039 bound also applies to these actual changed partitions, with freshly computed degrees and base11 digit counts. Old hook norms, dimensions, creation factors, PBW products and numerical content exponents are not reused.

The certificate supplies the endpoint degrees and safe digit counts for the entire new interval. K043's exact PBW product still applies, now with the new ell and q. Its per-word degree and all-word squareclass identities continue to hold; the old numerical digit bounds need not.

\subsection{4. Explicit quantitative counting margins}\label{explicit-quantitative-counting-margins}

Let b\_cnt=2\^{}480. By the proof of EG6 every bottom HC coordinate has absolute value less than b\_cnt²=A. The exact host discriminants of K030/K032 give

\[
 \operatorname{rd}(K_{14})=2\,7^{5/6}29^{6/7},\qquad
 \operatorname{rd}(K_{56})=2\,7^{5/6}29^{6/7}\sqrt{65}.
 \tag{EG13}
\]

Both are less than2·7·29·9\textless2\^{}12. Since log2\textless0.7,

\[
       \log(b_{\rm cnt}\operatorname{rd})
                  <492\log2<345.                 \tag{EG14}
\]

The actual bottom noncompact gap is h=2ell+q. Consequently the numerical gap hypothesis of source PointsLemma is satisfied whenever

\[
         0<C\le\frac{2\ell+q}{345(a+1)^{19}f^5}.
 \tag{EG15}
\]

This is a sufficient bound, in contrast to K043's necessary small-gap cap. Exact integer arithmetic gives the following convenient powers, already at q=0:

\begin{longtable}[]{@{}lrr@{}}
\toprule\noalign{}
Host & ell=ell\_* & ell=ell\_+ \\
\midrule\noalign{}
\endhead
\bottomrule\noalign{}
\endlastfoot
J14 & every0\textless C≤2\^{}262 & every0\textless C≤2\^{}347 \\
J56 & every0\textless C≤2\^{}214 & every0\textless C≤2\^{}299 \\
\end{longtable}

The certificate records the exact rational allowances and also the slightly larger values at s\_max. These are finite explicit ranges, not claims about the unknown constant exported by the source theorem.

Other displayed numerical conditions can also be separated exactly. One has log f\textless sqrt(a+1) for both hosts. Moreover log log b\_cnt\textless6, log43\textgreater15/4 and log169\textgreater5. Therefore

\[
 \log\log b_{\rm cnt}<c_b\,(a+1)f\log(a+1)
 \tag{EG16}
\]

is ensured by c\_b≥2\^{}(-10) for J14 and by c\_b≥2\^{}(-14) for J56. Also log log rd\textless a+1. These numerical assertions are certified with positive exponential-series bounds, not floating-point logarithms.

\subsection{5. Exact remaining interfaces}\label{exact-remaining-interfaces}

The source locators are PlaceDegreeData.lean55--59 for HidaSetting, 364--378 for PointsLemma, and414--423 for its existential exports. EG15--16 do not establish the following separate premises:

\begin{enumerate}
\def\labelenumi{\arabic{enumi}.}
\tightlist
\item
  The exported threshold n0 must be at most43 or169 respectively. Neither increasing ell nor changing a common determinant twist changes that rank condition. The accepted export supplies existence of n0, not a numerical value for the present application.
\item
  Its C0 and positive c\_b must meet the finite ranges above. A larger gap removes the severe small-gap mismatch of K043, but does not prove C0≤2\^{}262/2\^{}214 or c\_b above the displayed sufficient thresholds.
\item
  The actual finite IsAdmIdem, its congruence level g≤b\_cnt\^{}A\_cnt, CountData, and the relevant analytic projector/theorem hypotheses must be supplied. They are not created by changing infinity weights. In particular the existing uniform theorem export still has its own D7/admitted hypotheses; a generic AutInst definition alone does not instantiate the failed old host wrapper.
\item
  A positive cusp count would not itself give the selected packet-jet surjection or the norm and arithmetic lattice of a right inverse. K042(NB14) remains the precise interpolation target for each fixed ell. The norm cancellation of K043 is local and does not settle it.
\end{enumerate}

If all premises of PointsLemma are supplied with constants in EG15--16, its conclusion applies literally: the count differs from Dstar by at most exp(−c n f log n)Dstar. For a nonzero positive idempotent with Dstar\textgreater0 and c\textgreater0 this implies a positive count. This last implication is just the inequality exp(−c n f log n)\textless1. It does not assert the stronger selected nonvanishing or period lower bound needed by the MainX cost.

Thus this finding is an actual repair of one numerical counting gate within the constructed mixed family, with an exact bounded range of available constants. It is not an assertion of full DP44 scale data, a Kummer RH extension, or a favorable total gain/error/height balance.

\subsection{Source and verification record}\label{source-and-verification-record}

All Lean locators refer to the recovered pinned author source. The strict top condition is \nolinkurl{LocalRingCardinality/DataParams₂.lean88–89,} and its all-pairs be\_close condition is163. HCParam is Params.lean50--69; Centred and Dpar are DataTheta.lean538--554; strict NCGaps is cited above. These definitions are used exactly in their old scope.

The report uses the already checked R022 Scale proof, K030/K032 actual host discriminants, K040 characters, K039/K041 local-section and selected module theorems, K042 twist/level comparisons, and K043 PBW formula. No new global multiplicity theorem is cited or inferred. The certificate performs rational exponential rounding, affine endpoint/bit comparisons, integer partitions, exact base11 digit bounds and the displayed counting margins. It does not compute a cusp dimension, a global theta coefficient, a selected-jet determinant, an interpolation norm, or a Lean build.

\clearpage

\section{AM-K045 --- the actual auxiliary-prime exceptional sets are disjoint}\label{am-k045-the-actual-auxiliary-prime-exceptional-sets-are-disjoint}

Author: GPT-6 Astra Pro, ALP\_MACHINE /root/kummer\_extension, 22 November 2026 UTC. Inherited reasoning effort is not separately exposed.

For both current mixed CM hosts, the actual G-primes satisfy the exact BB\_disj condition required by the source NTypes record. The primes and ranks are unchanged. The verification uses30410 exact integer gcds and complete trial-division factorizations of only54 small residual gcds. It does not require factoring any full number g\^{}d−1.

\subsection{1. Literal statement and the two instances}\label{literal-statement-and-the-two-instances}

The accepted source definitions are \nolinkurl{CyclotomicNorms/StatementI.lean10–12:}

\[
 \operatorname{Exc}(g,d,p)
   \Longleftrightarrow p\text{ is prime and }
                   g^d<p^{4v_p(g^d-1)},
\] \[
 \operatorname{BB}(N,g)
       =\{g\}\cup\{p:\exists d\in\{1,\ldots,N\},
                                     \operatorname{Exc}(g,d,p)\}.
 \tag{BD1}
\]

The literal NTypes field is \nolinkurl{SharpRestriction/LevelLetterCocycles/DataN.lean917,} with N=R.r. Params.lean32 defines R.r=k+m−1, and m=2 here. Thus the two instances are

\begin{longtable}[]{@{}lrrrr@{}}
\toprule\noalign{}
Host & k & N & g0 & g1 \\
\midrule\noalign{}
\endhead
\bottomrule\noalign{}
\endlastfoot
J14 & 42 & 43 & 4615409 & 5276377 \\
J56 & 168 & 169 & 22285932761 & 64072175441 \\
\end{longtable}

These are exactly the K031/K032 auxiliary primes. Their primality is rechecked by division by every odd integer through the square root; the largest such limit is253124. No probable-prime test is used.

\subsubsection{Theorem BD2}\label{theorem-bd2}

For each of the two rows above,

\[
                    \operatorname{BB}(N,g_0)
                       \cap\operatorname{BB}(N,g_1)=\varnothing.
 \tag{BD2}
\]

The conclusion is the entire BB intersection, including the distinguished primes g0 and g1 themselves. It does not list either entire BB set.

\subsection{2. An elementary gcd exclusion}\label{an-elementary-gcd-exclusion}

\subsubsection{Lemma BD3}\label{lemma-bd3}

Let A,B\textgreater1 be integers and let C=gcd(A−1,B−1). If a prime p satisfies

\[
                 A<p^{4v_p(A-1)},\qquad
                 B<p^{4v_p(B-1)},
\]

then C\^{}4\textgreater min(A,B). Consequently C\^{}4≤min(A,B) excludes every such p.

\textbf{Proof.} Write u=v\_p(A−1), v=v\_p(B−1). Both are positive, since otherwise the corresponding right side is1. If u≤v, then p\^{}u divides C and A\textless p\^{}(4u)≤C\^{}4. If v≤u, then B\textless p\^{}(4v)≤C\^{}4. In either case C\^{}4 is strictly greater than min(A,B). QED.

Apply the lemma with A=g0\^{}d and B=g1\^{}e for every1≤d,e≤N. If a prime is exceptional for both bases, at one of its two witness exponents it must occur in a pair failing the gcd exclusion, and it must divide that pair's gcd. Hence it belongs to the union of the prime divisors of the remaining gcds. Fully factoring only those gcds therefore gives a rigorous finite superset of every possible common exceptional prime.

The exact screen counts are

\begin{longtable}[]{@{}
  >{\raggedright\arraybackslash}p{(\columnwidth - 10\tabcolsep) * \real{0.1304}}
  >{\raggedleft\arraybackslash}p{(\columnwidth - 10\tabcolsep) * \real{0.1739}}
  >{\raggedleft\arraybackslash}p{(\columnwidth - 10\tabcolsep) * \real{0.1739}}
  >{\raggedleft\arraybackslash}p{(\columnwidth - 10\tabcolsep) * \real{0.1739}}
  >{\raggedleft\arraybackslash}p{(\columnwidth - 10\tabcolsep) * \real{0.1739}}
  >{\raggedleft\arraybackslash}p{(\columnwidth - 10\tabcolsep) * \real{0.1739}}@{}}
\toprule\noalign{}
\begin{minipage}[b]{\linewidth}\raggedright
Host
\end{minipage} & \begin{minipage}[b]{\linewidth}\raggedleft
all exponent pairs
\end{minipage} & \begin{minipage}[b]{\linewidth}\raggedleft
excluded pairs
\end{minipage} & \begin{minipage}[b]{\linewidth}\raggedleft
remaining pairs
\end{minipage} & \begin{minipage}[b]{\linewidth}\raggedleft
distinct remaining gcds
\end{minipage} & \begin{minipage}[b]{\linewidth}\raggedleft
largest prime factor of a remaining gcd
\end{minipage} \\
\midrule\noalign{}
\endhead
\bottomrule\noalign{}
\endlastfoot
J14 & 1849 & 1672 & 177 & 27 & 61 \\
J56 & 28561 & 27943 & 618 & 27 & 991 \\
\end{longtable}

Every remaining gcd is factored completely by trial division through the proved primes below1000. The quotient is explicitly required to be1, and the product of the displayed prime powers is checked against the original gcd. This is a certificate of complete factorization, not an assertion that no larger prime was found in a bounded search.

\subsection{3. Complete candidate checks}\label{complete-candidate-checks}

For J14 the candidate set is exactly

\[
             \{2,3,5,7,19,23,29,31,61\}.
 \tag{BD4}
\]

For every candidate p and every1≤d≤43, the certificate evaluates the exact integer v\_p(g\^{}d−1) and the strict inequality in BD1. The complete nonempty candidate exceptions are

\begin{longtable}[]{@{}lrl@{}}
\toprule\noalign{}
Base & candidate p & all exceptional witness pairs (d,v\_p(g\^{}d−1)) \\
\midrule\noalign{}
\endhead
\bottomrule\noalign{}
\endlastfoot
g0=4615409 & 7 & (1,3) \\
g0=4615409 & 29 & (1,2) \\
g1=5276377 & 19 & (1,2) \\
g1=5276377 & 61 & (2,2) \\
\end{longtable}

All other candidate/base combinations have no exceptional exponent in the indicated range. Thus the two candidate exception sets are disjoint. The word candidate is essential: this table makes no claim that these are all exceptional primes of either individual base.

For J56 the candidate set is exactly

\[
 \{2,3,5,7,11,13,17,19,29,31,61,67,79,83,331,811,991\}.
 \tag{BD5}
\]

None of these primes is exceptional for either base at any exponent 1≤d≤169. This is checked by the same strict integer comparison, without logarithmic rounding. As a second consistency check, every prime factor of every unscreened pair is tested again at that pair's actual two exponents; no common exception occurs.

Finally the distinguished-prime cases in BD1 are checked separately. For each p=g0 and base g1, and each p=g1 and base g0, every exponent through N is tested. No cross-prime exception occurs. Since g0≠g1, the two distinguished singleton sets do not intersect. Together with Lemma BD3 and the complete candidate tables, this proves BD2. QED.

\subsection{4. Reproducible certificate and exact source consequence}\label{reproducible-certificate-and-exact-source-consequence}

The script \nolinkurl{actual_bb_disjointness_certificate.py} is self-contained apart from reading the three stated source files for anchors and hashes. It uses Python's exact integers, Euclidean gcd, integer square root, trial division, repeated division for valuations and direct comparison of positive integers. It has no symbolic-factoring or primality-oracle dependency. Its output stores every remaining exponent pair and gcd, all27 factorizations per host, all candidate exception lists, both cross-prime checks and a digest of the complete ordered gcd scan. Replaying the script verifies the excluded pairs directly as well.

The certificate records17 grouped assertions and30410 complete gcd pairs. Those group counts are not presented as30410 independent mathematical theorems. The proof is the elementary reduction above; the finite arithmetic certificate supplies its two concrete instances.

This closes the displayed NTypes.BB\_disj pin for the current prime/rank data. In particular it supplies the disjointness premise in the accepted \nolinkurl{CyclotomicNorms/StatementT.lean36–42} IrreducibilityCriterion. DO042 separately supplies the full IsGood/goodTuple predicate and DO043 its prescribed aggregate CM central character. Neither assertion is inferred from the gcd calculation; they have their own proofs.

No complete NTypes record, global Galois representation, projected automorphic period, coefficient-field comparison, MainX bound or Kummer RH extension follows from this one arithmetic pin. All other hypotheses of any use of the accepted irreducibility criterion remain literal. There is no Lean build or source-soundness claim.

\clearpage

\section{ALP-LIFT-32: generic finite branching, exact projected restriction, and the extra torus character}\label{alp-lift-32-generic-finite-branching-exact-projected-restriction-and-the-extra-torus-character}

Author: /root/lift\_accounting, GPT-6 Astra Pro, ALP\_MACHINE. 22 November 2026 UTC. Inherited reasoning effort is not separately exposed.

\textbf{Author whole-proof review passed, 22 November 2026, 03:52 UTC.} The finite theorem below applies to the actual completely split G residue groups. It supplies a quantitative restriction map from any admitted generic GL\_(a+1) constituent to every generic initial GL\_a type. It also states the exact extra scalar-torus compatibility. That compatibility cannot be omitted at coarse CM points, whose stabilizers need not be trivial.

The classical finite branching and Green parametrization are cited as published results. The contribution here is their application to the literal generic character, the exact restriction/extension normalization, the stabilizer-equivariant version, and the central-character dictionary needed to use a genuine goodTuple projector. No global jet interpolation, arithmetic row isomorphism, or DP44 comparison is inferred from the finite theorem.

\subsection{1. Finite groups and the exact generic character}\label{finite-groups-and-the-exact-generic-character}

Let k be a finite field with q elements, n=a+1, and put \[
 G=\mathrm{GL}_n(k),\qquad
 H=\{\operatorname{diag}(A,1):A\in\mathrm{GL}_a(k)\}.
 \tag{FB1}
\] Let U\_j be the standard upper unitriangular subgroup in GL\_j(k). Fix the actual nondegenerate character psi\_a of U\_a. At the G places in the source this is the prescribed simple-root additive character, with its sign fixed by the positive bottom character of R025; the initial bilinear row uses its inverse.

Choose a nondegenerate character psi\_n of U\_n whose restriction to U\_a is psi\_a. This means keeping all a-1 prescribed simple-root coefficients and choosing any nonzero coefficient on the last simple root. Define a representation to be psi\_j-generic when its psi\_j eigenspace under U\_j is nonzero.

All nondegenerate characters of U\_j are conjugate by the diagonal torus: for specified nonzero simple-root coefficients, choose consecutive diagonal ratios to transport one coefficient list to the other. Thus a generic irreducible is generic for the exact psi\_j, including its inverse. Conjugating the Whittaker vector does not alter a central projector.

Write \[
 P_n=\left\{\begin{pmatrix}A&v\\0&1\end{pmatrix}:
                   A\in\mathrm{GL}_a(k),\ v\in k^a\right\}
\] for the mirabolic subgroup and \[
 I_n=\operatorname{Ind}_{U_n}^{P_n}\psi_n,\qquad
 \Gamma_a=\operatorname{Ind}_{U_a}^{H}\psi_a.
 \tag{FB2}
\] Induction here uses functions f with f(up)=psi\_n(u)f(p), acted on by right translations. Equip all finite groups with probability measure, and all induced spaces with the corresponding invariant Hilbert norm.

Gan, Gross, and Prasad (2012), Proposition 4.2 and Corollary 4.3, identify the highest-derivative summand of the mirabolic restriction and its Gelfand--Graev restriction to H. The following elementary formulation also fixes the normalizations needed here.

\subsection{2. An isometric generic branching summand}\label{an-isometric-generic-branching-summand}

\textbf{Lemma.} I\_n is irreducible. Its restriction to H is isometrically Gamma\_a. Every irreducible generic representation Pi of G contains an isometric H-copy of Gamma\_a.

To verify the first assertion, write P\_n=H semidirect N, where N=k\^{}a is its last-column unipotent radical. The restriction of psi\_n to N is a nonzero additive linear functional. H acts transitively on the nonzero characters of N, and its stabilizer is the rank-(n-1) mirabolic group P\_(n-1). Induction in stages writes I\_n as induction from this stabilizer of I\_(n-1) tensored with that N-character. The finite little-group argument proves irreducibility inductively: distinct H-translates have distinct N-characters, so an invariant subspace is determined by its component at one character, and that component is irreducible by induction. The base I\_1 is one-dimensional.

Every p in P\_n lies in U\_n H, and U\_n intersect H=U\_a. Restriction of functions therefore identifies I\_n with Gamma\_a as an H-module. The norm is unchanged, since the N-fibres all have equal cardinality and the inducing character has modulus1. In particular \[
 d_I:=\dim I_n
       =\frac{|H|}{|U_a|}
       =\prod_{i=1}^{a}(q^i-1).
 \tag{FB3}
\]

Now choose a unit vector v in the psi\_n eigenspace of Pi. Define \[
 T_v f=\sum_{p\in U_n\backslash P_n}f(p)\Pi(p^{-1})v.
 \tag{FB4}
\] The summand is independent of a coset representative, because its two psi\_n factors cancel. Changing the summation variable proves T\_v R\_x=Pi(x)T\_v. It is nonzero: the function delta\_U supported on U\_n with value psi\_n maps to v. Its probability norm is d\_I\^{}(-1/2). Since I\_n is irreducible and T\_v is an intertwiner between unitary representations, T\_v\^{}*T\_v is a scalar. Evaluating it on delta\_U gives \[
 T_v^*T_v=d_I\,1,\qquad
                 \iota_v=d_I^{-1/2}T_v
                 \text{ is an isometric embedding.}
 \tag{FB5}
\] This proves the lemma, with the exact square-root normalization retained.

Consequently, if sigma is any psi\_a-generic irreducible of H and \[
 m_\sigma=\dim\operatorname{Hom}_H(\sigma,\Pi|_H),
\] then m\_sigma is positive. There is no requirement that sigma be Steinberg or cuspidal. The finite Gelfand--Graev representation for GL\_a is multiplicity-free; this follows, with its exact Green labels, from Gow and Vinroot (2007), Theorem 4.3. Positivity alone would suffice for the restriction result below.

\subsection{3. The complete projected restriction and its exact norm}\label{the-complete-projected-restriction-and-its-exact-norm}

Use probability Hilbert spaces L²(G) and L²(H). Let R denote restriction to H. Its adjoint is the actual extension by zero with the index factor \[
 (R^*w)(x)=[G:H]\,\mathbf1_H(x)w(x).
 \tag{FB6}
\] Let e\_Pi be the central orthogonal projector \[
 e_\Pi=\frac{d_\Pi}{|G|}
             \sum_{g\in G}\overline{\chi_\Pi(g)}\,R_g,
 \qquad d_\Pi=\dim\Pi,\quad (R_gf)(x)=f(xg).
 \tag{FB7}
\] Let P\_psi be the orthogonal projector onto the right psi\_a eigenspace: \[
 P_\psi=\frac1{|U_a|}
                  \sum_{u\in U_a}\psi_a(u)^{-1}R_u.
\] The target of interest is W\_psi=P\_psi L²(H). It is stable under left H; it is not asserted to be stable under all right H translations. Its right-H Fourier components are exactly the generic sigma components.

Because e\_Pi is central, it commutes with P\_psi. The positive operator \[
                         A_\Pi=R e_\Pi R^*
\] is central for the H regular action. Directly, for x in H, \[
 (A_\Pi w)(x)=\frac{d_\Pi}{|H|}
                \sum_{h\in H}\overline{\chi_\Pi(h)}\,w(xh).
 \tag{FB8}
\] Character orthogonality or Schur's lemma now gives, on the sigma isotypic component, \[
                  A_\Pi|_\sigma
                    =\frac{d_\Pi m_\sigma}{d_\sigma}\,1.
 \tag{FB9}
\] Indeed the trace of the central operator in sigma is d\_Pi times the inner product of chi\_Pi restricted to H with chi\_sigma. This is d\_Pi m\_sigma, and division by d\_sigma proves the formula.

Every eigenvalue in (FB9) on W\_psi is positive by section 2. Therefore \[
 R:e_\Pi P_\psi L^2(G)\longrightarrow W_\psi
                       \text{ is onto.}
 \tag{FB10}
\] An explicit right inverse is \[
                \mathcal E_\Pi=e_\Pi R^* A_\Pi^{-1}
                   \quad\text{on }W_\psi.
 \tag{FB11}
\] There is no unspecified choice of a finite test vector in this formula. If w=sum\_sigma w\_sigma is its orthogonal right-H Fourier decomposition, then \[
 \|\mathcal E_\Pi w\|_G^2
       =\sum_\sigma\frac{d_\sigma}{d_\Pi m_\sigma}
                         \|w_\sigma\|_H^2
       \le \|w\|_H^2.
 \tag{FB12}
\] The equality follows from e\_Pi=e\_Pi\^{}*=e\_Pi² and (FB8). The inequality follows from m\_sigma\textgreater=1 and d\_Pi\textgreater=d\_sigma whenever sigma occurs in Pi restricted to H. It is a probability-measure statement; the index in (FB6) is essential.

For a single sigma component the actual coefficient formula is \[
 (\mathcal E_\Pi w_\sigma)(x)
       =\frac{d_\sigma}{m_\sigma |H|}
             \sum_{h\in H}\chi_\Pi(h^{-1}x)\,w_\sigma(h).
 \tag{FB13}
\] To check it, substitute g=x\^{}(-1)h in (FB7) applied to R\^{}*w and then use (FB9). Thus the values lie in the field generated by the finite character and the input coefficients. Character values are algebraic integers, but the displayed m\_sigma \textbar H\textbar{} denominator is retained. If the input must first be projected onto sigma, that central H projection contributes its own exact factor d\_sigma/\textbar H\textbar. No integral-unit or 11-saturation assertion follows from (FB12).

The same formulas apply to a central orthogonal projector e which is a sum of irreducible projectors. Its eigenvalue is \[
               A_e|_\sigma
                 =\frac1{d_\sigma}
                    \sum_{\Pi\in\operatorname{supp}e}
                               d_\Pi m_\sigma(\Pi).
 \tag{FB14}
\] One generic constituent in the support makes this positive on all of W\_psi. The extension eR\^{}\emph{(ReR\^{}})\^{}(-1) has the corresponding exact norm.

\subsection{4. The extra scalar torus: a necessary condition}\label{the-extra-scalar-torus-a-necessary-condition}

Let C be any subgroup of k\^{}times and embed \[
 K=H\times C\hookrightarrow G,\qquad (A,c)\mapsto\operatorname{diag}(A,c).
 \tag{FB15}
\] Let omega\_Pi and omega\_sigma denote the central characters. The elementary matrix identity \[
 \operatorname{diag}(A,c)
       =c I_n\,\operatorname{diag}(c^{-1}A,1)
\] shows that the sigma isotypic summand of Pi restricted to K has exactly the additional character \[
                 \nu_\sigma(c)
                  =\omega_\Pi(c)\omega_\sigma(c)^{-1}.
 \tag{FB16}
\] The multiplicity space carries no additional C action: the displayed central matrix acts by omega\_Pi, and c\^{}(-1)I\_a acts by omega\_sigma(c)\^{}(-1) on every copy of sigma. Hence \[
 \operatorname{Hom}_K(\sigma\boxtimes\nu,\Pi)\ne0
 \quad\Longleftrightarrow\quad
 m_\sigma>0\ \text{ and }\
 \nu=\omega_\Pi|_C\,\omega_\sigma|_C^{-1}.
 \tag{FB17}
\] This is not an arbitrary line-character choice.

Apply the restriction proof of section 3 with K in place of H. On each compatible sigma tensor nu block the eigenvalue remains d\_Pi m\_sigma/d\_sigma; on an incompatible block it is zero. Therefore the projected restriction is onto the entire generic, central-compatible subspace of L²(K), with exactly (FB11)--(FB13), replacing H by K.

All maps in this assertion intertwine both left and right K translations on their full compatible isotypic domain. Restriction and extension by zero do so, the G central projector commutes with both actions, and A\_Pi is central on K. After applying the right-U\_a projector, left K equivariance remains. Consequently the construction commutes with every actual finite stabilizer Gamma contained in K, including a stabilizer acting simultaneously on a jet fibre. Tensoring with the fibre and taking Gamma invariants preserves this explicit right inverse.

This last statement still requires the input to have the compatible central character. It does not turn an incompatible fibre into a compatible one. In the mixed CM application C contains the residue images of mu\_28. Strict-level neatness does not imply that the coarser point stabilizer is trivial.

For products over the actual G places, take tensor products of these groups and operators. The compatibility condition for the diagonal global mu\_28 is the product of (FB16), with the actual embedding of each root of unity into each residue field. It is the aggregate central character, not an arbitrarily identified identical generator at every place.

\subsection{5. Genuine squarefree series are generic}\label{genuine-squarefree-series-are-generic}

The source DataN.lean960--966 lists central, orthogonal and exhaustive eSer axioms; those axioms alone do not identify a generic finite representation. The application here uses the genuine finite-GL realization.

Choose the usual compatible Green duality between Frobenius orbits of nonzero finite-field elements and orbits of multiplicative characters. If \[
                 f=\prod_j f_j
\] is monic squarefree, of degree n, and has nonzero constant term, its factors correspond to distinct cuspidal orbit labels. In the Green parametrization the associated irreducible Pi\_f has partition (1) at each of those labels and the empty partition elsewhere. Gow--Vinroot Theorem3.2 gives this irreducible parametrization; their Theorem4.3 shows it is generic because every nonempty partition has one row. Equivalently, Pi\_f is the irreducible parabolic product of the mutually distinct cuspidal factors; GGP Proposition4.1 also gives its one-dimensional highest derivative.

More generally the genuine eSer(f) is the sum of the central irreducible projectors whose partition sizes at the orbit labels equal the factor multiplicities of f, and is zero for an invalid f.~This supplies precisely the finite eSer axioms. In the squarefree case it is the single projector e\_(Pi\_f). Thus any actual goodTuple of squarefree polynomials supplies a tensor product of generic representations in the support of the genuine e0.

This identifies the finite series used here. It does not manufacture the remaining source ties to local automorphic parameters or an existing global row-family record. Those ties retain their own statements.

The inverse convention causes no hidden issue: taking duals inverts the orbit labels and the generic character. The parent's coefficient-reality audit is separate. For a bilinear period pairing a positive bottom character with an inverse-generic initial row, complex conjugation gives the corresponding opposite-character Hilbert-space statement; it cannot be dropped from the source dictionary.

\subsection{6. The central character and the norm of a root}\label{the-central-character-and-the-norm-of-a-root}

We give the exact character dictionary required for the mu\_28 condition, retaining the choice of duality. For a factor of degree d choose a cyclic generator t of F\_(q\^{}d)\^{}times and a compatible primitive complex root zeta\_(q\^{}d-1). Write alpha=t\^{}b and define \[
 \theta_\alpha(t)=\zeta_{q^d-1}^{\,b},\qquad
 t_1=t^{(q^d-1)/(q-1)}.
 \tag{FB18}
\] For c=t\_1\^{}j in k\^{}times, \[
 \theta_\alpha(c)=\zeta_{q-1}^{\,bj},\qquad
             \operatorname{Norm}_{F_{q^d}/k}(\alpha)=t_1^b.
 \tag{FB19}
\]

The corresponding Deligne--Lusztig constituent has central character theta\_alpha restricted to k\^{}times. This sign follows directly from Deligne--Lusztig (1976), Theorem4.2: at a central semisimple element z, the character formula has u=1 and connected centralizer G, so \[
                   R_T^\theta(z)=\theta(z)R_T^\theta(1).
 \tag{FB20}
\] Passing from the virtual character to its irreducible signed representative does not change the quotient by its dimension. Parabolic products multiply central characters. Therefore for Pi\_f the character is labelled by the product of the factor norms, equivalently by the semisimple determinant (-1)\^{}n f(0), in the chosen Green duality.

If 28 divides q-1 and c is a chosen generator of the residue mu\_28, write \[
                  c=t_1^{(q-1)j/28},\quad (j,28)=1.
\] Then \[
                        \omega_{\Pi_\alpha}(c)
                             =\zeta_{28}^{\,b j}.
 \tag{FB21}
\] Frobenius changes b to qb, which leaves this residue unchanged because q=1 modulo28. Reversing the global duality convention simultaneously inverts these characters; it does not alter the cardinalities or the compatibility test.

In particular, in a full-degree root space, the 28 exponent classes b modulo28 have equal size. If an admitted root set has density greater than4/5, it meets at least23 of these classes. Two independently variable such places realize every aggregate mu\_28 character, since any two subsets of Z/28 of cardinality at least23 have sum all of Z/28. This final combinatorial implication uses only the stated density and independence. The actual goodTuple count and factor-degree clauses are DO042's separate construction; DO043 proves the two-place central-character refinement with the actual residue embeddings. They are not silently assumed invariant under scalar twisting.

\subsection{7. Consequence and remaining global interface}\label{consequence-and-remaining-global-interface}

For an admitted generic Pi, every inverse-generic initial finite GL\_a row can be paired with a positive-character bottom finite coefficient in its genuine projector support, with the exact finite restriction norm (FB12). If the coefficient space includes the additional scalar torus, the same assertion holds precisely on the central-compatible subspace (FB17). The canonical extension respects any nontrivial finite stabilizer and its fibre action.

This gives the finite representation-theoretic ingredient for an equivariant global jet construction. It does not itself prove that those jets are globally realizable in a fixed automorphic bundle, satisfy every actual stabilizer character, or survive the source's particular period functional. Parent's separate construction must use the actual coarse CM orbit and its full stabilizers. After that passage, LIFT31 supplies membership in the actual downward image; its global inverse norm remains unbounded here.

All probabilities, group indices, dimension ratios and finite coefficient denominators are visible above. An orthogonal projection has Hilbert norm1, but this says nothing by itself about its integral lattice or about the arithmetic height of the resulting global cusp form.

\subsection{8. Primary references and exact source locators}\label{primary-references-and-exact-source-locators}

Gan, W. T., Gross, B. H., \& Prasad, D. (2012). Restrictions of representations of classical groups: Examples. In Sur les Conjectures de Gross et Prasad, I (Astérisque346, pp.111--170). Société mathématique de France. Proposition4.2 and Corollary4.3, author-preprint printedpp10--11; Proposition4.1 for the highest derivative of a parabolic product. Primary author copy: \url{https://abel.math.harvard.edu/~gross/preprints/ggp-evidence4.pdf} . Published copy: \url{https://www.numdam.org/item/AST_2012__346__111_0.pdf} .

Gow, R., \& Vinroot, C. R. (2007). Extending real-valued characters of finite general linear and unitary groups on elements related to regular unipotents. arXiv:0704.2022v1. Theorem3.2, PDFp9; Lemma3.1, PDFp10; Theorem4.3, PDFp12. \url{https://arxiv.org/pdf/0704.2022} . Local primary copy: work/agents/lift\_refs/gow\_vinroot2007.pdf and .txt.

Deligne, P., \& Lusztig, G. (1976). Representations of reductive groups over finite fields. Annals of Mathematics,103(1),103--161. Theorem4.2, printedp124, with the central specialization (FB20); section5 fixes compatible torus character dualities. \url{https://publications.ias.edu/sites/default/files/Number27.pdf} .

The literal source locations are DataN.lean960--966 and996--1007 for eSer/e0; \nolinkurl{CyclotomicNorms/StatementT.lean19–29} for the complete IsGood predicate; RecLevel.lean50--52 and ColForms.lean93--115 for the inverse-generic initial row; and R025's direct Hida/RowLevelMatch audit for the positive bottom character and full residue action. The actual G places and field embeddings are the carried K031/K032/K033/K040 inputs. No new proof of those carried results or a Lean build is claimed.

\subsection{9. Verification record}\label{verification-record-1}

The author read the complete proof FB1--FB21, including all finite measure and character conventions, the full compatible K-block extension, and its stabilizer-equivariance. The primary GGP §4, Gow--Vinroot Theorems3.2/4.3 and Deligne--Lusztig Theorem4.2 were read directly. No statement is promoted from an unread source axiom to an actual representation.

The exact companion certificate \nolinkurl{finite_generic_branch_certificate.py} completed successfully at 2026-11-22T03:38:32.779906+00:00; its full output is \nolinkurl{finite_generic_branch_certificate.json}. It independently constructs the mirabolic characters for (n,q)=(2,3),(3,2),(3,3), checks irreducibility and all 56 H-character values, and constructs actual finite Steinberg characters and central restriction matrices for GL2(F3) and GL3(F2). There are 216 exact central-idempotent convolution checks, 3,408 left-K equivariance checks and 96 wrong-central-character annihilation checks. It verifies the displayed right inverse, all exact norm Gram matrices and positive eigenvalues. These checks are finite diagnostics, not a formal proof of the general theorem, a global jet construction, an integral-saturation computation, or a Lean build.

\clearpage

\section{AM-R-027: the genuine e0 projector preserves a qualitative bottom detector}\label{am-r-027-the-genuine-e0-projector-preserves-a-qualitative-bottom-detector}

Author: GPT-6 Astra Pro, ALP\_MACHINE \nolinkurl{/root}, 22 November 2026 UTC. Reasoning effort is not separately exposed.

This theorem concerns the actual mixed hosts, finite levels and T translation of LIFT29. For sufficiently large canonical twists it constructs the LIFT29 detector inside the genuine source e0 summand. The finite projector is not replaced by a convenient generic projector: a genuine irreducible block contained in that exact e0 supplies the finite extension, after which the full e0 acts as the identity.

The statement is qualitative in the twist. It does not assert that its threshold lies in K042's bounded progression, and it gives no norm or integral denominator for the global interpolation inverse. Consequently it does not yet give a controlled terminal value at an admissible current weight or a DP44 saving. The entire same-family arithmetic/real adapter and R023's global Host/scale obstruction remain separate.

\subsection{1. Principal level, the actual packet, and the central character}\label{principal-level-the-actual-packet-and-the-central-character}

Write E/F for K031 or K032, H=U(I\_a), G=U(I\_a,-1), and n=a+1. Let Jmax be the full standard finite lattice stabilizer in G(A\_f). Let J1 be principal modulo g0g1 at all selected G places and equal to Jmax at every other finite place. The two actual residue groups are

\[
 \mathcal G=\prod_{k,v}GL_n(F_{g_k}),\qquad
 \mathcal H=\prod_{k,v}GL_a(F_{g_k})\subset\mathcal G.
                                                               \tag{R27.1}
\]

R025 supplies Jmax/J1=Gcal with the exact full action. Put Ucal equal to the product of the upper unipotent subgroups of Hcal, and let psi be the actual generic character, inverse to the initial character. Then the source-shaped bottom dat0 group is J\_B=red\^{}\{-1\}Ucal, with right character psi. Its cusp space is the psi part of the principal-level cusp space. The genuine e0 commutes with that part.

Retain the actual finite lower-unipotent translation tau at T from LIFT29(BT4--8), identity at G and at infinity. Set

\[
 J_H^1=H(A_f)\cap\tau J^1\tau^{-1},\qquad
 J_H^{\max}=H(A_f)\cap\tau J^{\max}\tau^{-1}.               \tag{R27.2}
\]

At T both are exactly the original strict covector stabilizer. At G they are respectively principal and maximal. The original initial J0 is the inverse image of Ucal between them, with inverse generic character. Thus its opposite coefficient rows are exactly the psi part of the principal H-level coefficient space.

Use LIFT28's discrete series sigma\_s, minimal compact type W\_alpha tensor z\^{}beta(s), and actual bottom type tau\_s=sigma\_s\^{}vee. Here b1 increases by ns, beta decreases by ns, while u,v,b0 and the selected degrees N\_tau are fixed. At one label,

\[
       N=|\lambda|-|\alpha|,\qquad
       |\lambda|+\beta-N=|\alpha|+\beta.                    \tag{R27.3}
\]

Let Z be the image of the global CM roots mu(E)=mu\_28 as scalar matrices in Gcal, and let Omega\_s be the actual character by which Z acts on the principal-level bottom automorphic space. With right translation on scalar automorphic functions its value is

\[
 \Omega_s(z)=\omega_{\tau_s,\infty}(z)^{-1}
       =\prod_{\tau} \tau(z)^{|\alpha_\tau|+\beta_\tau(s)}.
                                                               \tag{R27.4}
\]

Indeed the finite adele z\_G, supported only at G, equals a global central z times z\_infty\^{}\{-1\} and central integral components away from G. The latter act trivially at J1. Left global invariance and the real central character give(R27.4). This determines the sign; the opposite automorphic-action convention would invert both sides. For s a multiple of28, Omega\_s is independent of s. It is also independent of the interior word when v,b0 are fixed, as in R022.

\subsection{2. The complete coarse packet and its deck stabilizers}\label{the-complete-coarse-packet-and-its-deck-stabilizers}

Let o be the negative-line point and form the principal-level quotient X1=G(F)(D times G(A\_f))/J1 and the coarse-level quotient Xmax with Jmax in place of J1. The finite deck action is the right Gcal action. Every arithmetic group at J1 is neat, since J1 is contained in the already neat level of K042. No trivial-stabilizer assertion is made for the coarser level.

\subsubsection{Lemma R27.5}\label{lemma-r27.5}

The map

\[
 H(F)\backslash H(A_f)/J_H^{\max}
   \longrightarrow X^{\max},\qquad [x]\mapsto[(o,x\tau)]
                                                               \tag{R27.5}
\]

is injective. For the principal-level point p\_x={[}(o,x tau){]}, its stabilizer S\_x in Gcal satisfies

\[
 S_x=(S_x\cap\mathcal H)\times Z,\qquad
 S_x\cap\mathcal H\simeq
           H(F)\cap xJ_H^{\max}x^{-1}.                     \tag{R27.6}
\]

The product uses the scalar central Z inside Gcal. Its intersection with Hcal=diag(GL\_a,1) is trivial. The Hcal orbit above x is the actual initial principal-level packet over this coarse H class.

\textbf{Proof.} A rational element relating two displayed points fixes o at every real place. It therefore preserves the rational negative line and is diag(h,t), with h in H(F) and t in E\^{}1(F). At places away T the block level condition makes t and its inverse integral. At T the exact multiplication LIFT29(BT7) has bottom-right entry t and inverse bottom-right entry t\^{}\{-1\}; hence the same conclusion holds. Thus t is a global norm-one integral unit. Its conjugates all have absolute value1, so the established root-unit argument gives t in mu\_28.

The scalar global t preserves every standard finite lattice. It also commutes with tau. Multiplying diag(h,t) by that scalar's inverse replaces it by diag(t\^{}\{-1\}h,1) and preserves every finite level condition. The resulting equivalence is precisely an H equivalence modulo J\_Hmax, proving injectivity in(R27.5).

The same argument for a stabilizer shows that its residue lies in Hcal times Z. Conversely every z in Z fixes each geometric point p\_x: the global scalar z has trivial action on the negative-line domain, and its components outside G are central elements of J1. This assertion concerns geometric points; its nontrivial action on the jet fibre is exactly Omega\_s from(R27.4). Since Z is contained in S\_x, dividing any s in S\_x by its Z factor puts it in S\_x cap Hcal. This proves the product decomposition.

For the H factor, the rational element after scalar removal lies in H(F) cap xJ\_Hmax x\^{}\{-1\}; its residue gives every element of S\_x cap Hcal. The residue map has trivial kernel here. Such a kernel element is an arithmetic stabilizer at the principal H level, a subgroup of the neat principal bottom level. It is finite because H is definite; a finite neat group is trivial. This proves the last identification. The same double-coset argument identifies the Hcal orbit with the initial principal-level packet, including its real fibre action. QED.

In particular different coarse H classes give disjoint full Gcal orbits in X1. Their packet data can be treated independently; no unrecorded relation between different initial classes is imposed.

\subsection{3. The selected jet module and finite extension}\label{the-selected-jet-module-and-finite-extension}

Fix one word and one coarse class x. Let V\_x be the selected homogeneous jet fibre from LIFT28(HJ20), with its actual unitary action of S\_x cap Hcal. That action is the original opposite initial coefficient action. The scalar Z acts by Omega\_s. With consistent deck transport, the full orbit jet module and its initial packet submodule are

\[
 \operatorname{Ind}_{S_x}^{\mathcal G}
       (V_x\boxtimes\Omega_s),\qquad
 \operatorname{Ind}_{S_x\cap\mathcal H}^{\mathcal H} V_x.
                                                               \tag{R27.7}
\]

These are genuine induced finite modules with their fibre relations, not free assignments at a coarse point. Concretely they are V\_x-valued functions satisfying F(sg)=rho\_x(s)F(g). Restriction to Hcal gives the displayed initial module. The central part of its extension to Hcal times Z is fixed by Omega\_s. Each point change fixing o is in U(a) times U(1) at infinity; in Harish--Chandra coordinates it acts linearly on the ball variables and preserves the selected homogeneous Pieri summand. Therefore(R27.7) includes the full stabilizer action on these jets, without lower-degree mixing.

By DO043 choose a genuine irreducible product representation Pi of Gcal, contained in the literal goodTuple/e0 sum, whose character on Z is Omega\_s. Every factor is generic by R026. LIFT32's finite mirabolic restriction theorem shows that Pi restricted to Hcal contains every generic irreducible of Hcal. Its normalized restriction and canonical right inverse commute with left Hcal and with Z. They also preserve the right psi type. Consequently that right inverse, applied componentwise in V\_x, preserves the left S\_x relations of(R27.7).

Explicitly, first extend an initial function w on Hcal to Hcal times Z by w(hz)=Omega\_s(z)w(h), and then apply

\[
 \mathcal E_{\Pi}=e_{\Pi}R^*(R e_{\Pi}R^*)^{-1}           \tag{R27.8}
\]

on the generic compatible part. All constituents of the right psi space are generic. LIFT32 gives invertibility there, including its exact eigenvalues; it is not an inference from e0 nonzero alone. Because the operator commutes with the left group action, the S\_x cap Hcal relation on w is retained; because Pi has character Omega\_s, the central Z relation is retained. Therefore

\[
 R:e_{\Pi}\bigl(\operatorname{Ind}_{S_x}^{\mathcal G}
                   (V_x\boxtimes\Omega_s)\bigr)^{U,\psi}
       \longrightarrow
       \bigl(\operatorname{Ind}_{S_x\cap\mathcal H}^{\mathcal H}
                   V_x\bigr)^{U,\psi}
 \quad\hbox{is onto}.                                     \tag{R27.9}
\]

In(R27.9), e\_Pi denotes the Pi-isotypic projection; the explicit right inverse is(R27.8). Since e0 e\_Pi=e\_Pi, the same onto assertion holds with the full e0 summand in the source. Summing over the disjoint coarse classes proves finite selected-jet restriction surjectivity onto the entire opposite initial coefficient space.

The finite probability-Haar norm estimates and group-algebra denominators are precisely those of LIFT32 and R026. Passing to unweighted packet sums retains their indices and stabilizer orders. No such finite norm estimate is identified with a global Petersson interpolation norm.

\subsection{4. Equivariant coherent interpolation at large twist}\label{equivariant-coherent-interpolation-at-large-twist}

Take the union of the full Gcal orbits of the principal-level packet, and thicken each point to include all required lower and selected jets. This is a finite Gcal-stable subscheme in the interior of X1. Work on the same toroidal/minimal compactification and subcanonical cusp bundle used in LIFT28, now at the principal level. The carried Mumford/Serre argument applies to this finite invariant thickening: for all sufficiently large twists in a fixed positive progression, the global cusp-section map onto its complete jet space is surjective. The map is equivariant for the actual finite deck action on the automorphic bundle.

Choose the progression to be a multiple of28 as well as the fixed canonical-line descent period. Then Omega\_s is fixed by(R27.4). The chosen Pi can be fixed on this progression. In characteristic zero, applying e\_Pi and the right psi projector to an equivariant surjection preserves surjectivity onto the corresponding target summand. Thus every extended datum from §3 is realized by an actual cusp section in e0 and of the required bottom finite character.

The passage from a selected jet to the actual selected K-type value is LIFT28 §7, with the now explicit PBW scalar of K043. Prescribe lower jets to be zero and the selected jet to be the inverse of the nonzero leading-symbol scalar times the desired coefficient. The compact stabilizer action preserves this prescription. Equivalently one can invert the finite triangular lower-jet correction. This gives the following actual restriction surjectivity, for all sufficiently large s in the fixed progression:

\[
 \begin{split}
 \mathcal R_s:
 e_0\operatorname{Hom}_{G_\infty}
    (\tau_s,\mathcal A_{\rm cusp}(G,J_B,\psi))
       &\longrightarrow \mathcal A_{\lambda^\vee}
                              (H,J_0,\chi_0^{-1}),\\
 F&\longmapsto[h\mapsto F(\operatorname{diag}(h,1)\tau)] .
 \end{split}                                                \tag{R27.10}
\]

The selected K-type embedding is the fixed isometric one of LIFT28/K043. The finite tau is unchanged. Neither finite theta kernel is altered by this interpolation construction.

For a fixed finite word family, use all the corresponding selected jet conditions in one finite thickening and take the largest of their finitely many sufficient twists. In R022 the minimal bundle and central character are common across words; the selected components and degrees may differ. The statement is uniform over that entire finite family, with no effective bound on this maximum.

\subsection{5. Exact positive period and the remaining admissible-weight gate}\label{exact-positive-period-and-the-remaining-admissible-weight-gate}

Let P be any nonzero actual initial coefficient row, in particular any nonzero Mram row, and choose a unit seed w in its compact type. Surjectivity(R27.10) supplies F in the genuine e0 bottom summand with restriction equal to the Hermitian conjugate of f\_(P,w). The actual compact period is then

\[
 \boxed{\quad
  \int_{[H]} f_{P,w}(h)
       F(\operatorname{diag}(h,1)\tau)\,dh
       ={v_0\over d_\lambda}
             \sum_x {\|P_x\|^2\over|\Gamma_x|}>0.\quad}     \tag{R27.11}
\]

This is LIFT29's original strict-level measure and original real isometry normalization. The root-unit stabilizers at the coarse level have already been accounted for in §2--3; no additional28 is inserted into(R27.11). At the actual strict level K042 gives Gamma\_x=1, but the displayed formula retains the full general stabilizer form for comparison with the preceding records.

If, at such a weight, the same actual vertical family has the real types and all hypotheses of LIFT31, its unchanged downward map D is onto the full bottom space. Thus e0 D is onto its e0 summand, and F has an actual top preimage. This is the correct conditional use of LIFT31. The present proof does not assert that a sufficiently large s occurs within K042's finite admissible interval or satisfies every other source/global weight constraint simultaneously.

The remaining effective problem is now a bound putting the equivariant selected-jet threshold inside that permitted interval, together with the actual Petersson and integral size of its right inverse and of the LIFT31 preimage. Normalizing F to unit norm divides(R27.11) by its actual norm; no uniform lower estimate for that quotient is proved here. The same-projector arithmetic grp\_record, full real/frame transport, literal brC/tilQ/PtilFF identity at the chosen admissible family, and the specific032 cycle comparison remain required. Analytic large-weight surjectivity does not by itself solve these arithmetic equations or determine DP44.

\subsection{6. Sources, frameworks and check scope}\label{sources-frameworks-and-check-scope}

New arguments are the complete coarse packet/stabilizer factorization, the actual induced jet module with its scalar character, and the equivariant finite-extension passage to the genuine e0 cusp detector. The finite representation theorem is LIFT32(FB1--21); goodTuple and its scalar prescription are DO042/043 and R026; the exact source finite action is R025. LIFT28(HJ15--23) and LIFT29(BT3--14) supply the checked automorphic-bundle, selected-jet and period framework. K042 supplies neatness and the finite permitted window. K043 supplies the selected PBW norm. LIFT31(VA1--33) is used only conditionally at weights satisfying its real-family hypotheses.

The known coherent inputs retain their exact locators: Mumford(1977), Main Theorem3.1 and Propositions3.3--3.4, pp.256--257/264--265; Grothendieck(1961), EGAIII.1 Theorem2.2.1(ii), p.100; the complete references and checked scope are in LIFT28. Principal-parts schemes are EGAIV.4 16.11.2--3(ii), pp.54--55. The finite residue action and Schur coefficient frameworks remain SGA3I4.7.3/4.7.3.1 and SGA6VI3.3; the fixed CM Picard/biextension comparison remains FGA232 Theorem3.1 pp.148--149, FGA236 Corollary1.7 p.229/2.3 p.232, and SGA7I VII2.7/2.9 pp.169--174 andVIII7.4(b) pp.310--311. The analytic positive period is not an algebraic-cycle nonvanishing theorem, a motivic projector, an ell-adic weight claim or a Kummer RH theorem.

No Lean build, global numerical period evaluation or effective interpolation computation is claimed. Whole-proof review events are recorded separately; the theorem is not promoted before the finite representation input and this entire argument are checked.

\clearpage

\section{AM-R-028: an explicit finite restriction inverse integral at11}\label{am-r-028-an-explicit-finite-restriction-inverse-integral-at11}

Author: GPT-6 Astra Pro, ALP\_MACHINE \nolinkurl{/root}, 22 November 2026 UTC. Reasoning effort is not separately exposed.

This gives a different inverse from the minimum-norm inverse of LIFT32. It has an exact probability-Hilbert norm and is integral at every prime different from the residue characteristic. In particular the actual finite extension needed by R027 has no11-adic denominator. It does not bound a global cusp-section interpolation inverse or the integral inverse of the vertical theta map.

\subsection{1. The normalized finite Bessel coefficient}\label{the-normalized-finite-bessel-coefficient}

Use LIFT32's exact groups

\[
 G=GL_n(F_q),\quad H=\operatorname{diag}(GL_{n-1}(F_q),1),
 \quad U=U_n,\quad U_H=U_{n-1},\quad n=a+1.
 \tag{R28.1}
\]

Let psi be a nondegenerate character of U extending the actual positive bottom character psi\_H on U\_H. Let Pi be any irreducible generic complex representation of G, of dimension d\_Pi and character chi\_Pi. Equip Pi with an invariant positive Hermitian form, conjugate linear in the first argument. Finite GL generic multiplicity one is Gow--Vinroot(2007), Theorem4.3: the Gelfand--Graev character is the multiplicity-free sum of the height-one Green characters. Hence

\[
 P_\psi={1\over|U|}\sum_{u\in U}\psi(u)^{-1}\Pi(u)
                  \quad\hbox{has rank one}.
 \tag{R28.2}
\]

For a unit vector v spanning its image define

\[
 \begin{split}
 B_\Pi(g)&=\operatorname{tr}(\Pi(g)P_\psi)
               =\langle v,\Pi(g)v\rangle\\
 &= {1\over|U|}\sum_{u\in U}
                         \psi(u)^{-1}\chi_\Pi(gu).
 \end{split} \tag{R28.3}
\]

This formula fixes both the sign and normalization. In particular B\_Pi(1)=1, its absolute value is at most1, and

\[
 B_\Pi(u_1gu_2)=\psi(u_1)\psi(u_2)B_\Pi(g),\qquad
 B_\Pi(zg)=\omega_\Pi(z)B_\Pi(g)
 \tag{R28.4}
\]

for u1,u2 in U and central z. These identities follow immediately from the rank-one coefficient and Pi(u)v=psi(u)v. The trace formula also shows that its values lie in the field generated by chi\_Pi and psi, and that \textbar U\textbar B\_Pi(g) is an algebraic integer.

\subsubsection{Lemma R28.5}\label{lemma-r28.5}

On the embedded H the coefficient is exactly

\[
             B_\Pi(h)=\mathbf1_{U_H}(h)\psi_H(h).
 \tag{R28.5}
\]

\textbf{Proof.} The irreducible mirabolic representation I\_n=Ind\_U\^{}\{P\_n\}psi of LIFT32(FB2--5) occurs in Pi restricted to P\_n. Frobenius reciprocity and (R28.2) show that its multiplicity is one. The entire psi line therefore lies in that summand. In the probability induced model, the function delta\_U, equal to psi on U and zero elsewhere, has norm d\_I\^{}(-1/2), where

\[
       d_I=[H:U_H]=\prod_{i=1}^{n-1}(q^i-1).
 \tag{R28.6}
\]

The vector d\_I\^{}(1/2)delta\_U is a unit psi vector. Its right-translation matrix coefficient on P\_n is psi(p) if p is in U and zero otherwise: the two supporting left U cosets are disjoint in the latter case. The isometric inclusion from LIFT32 identifies this vector with v up to a unit scalar, which does not change its coefficient. Restrict to H and use U intersect H=U\_H. QED.

\subsection{2. Explicit extension and exact norm}\label{explicit-extension-and-exact-norm}

Let W\_psi be the functions w on H satisfying w(hu)=psi\_H(u)w(h), with the probability norm on H. Define

\[
 (\mathcal E_Bw)(g)=
       \sum_{h\in H/U_H}w(h)B_\Pi(h^{-1}g).
 \tag{R28.7}
\]

The notation H/U\_H means right cosets. A representative change h to hu multiplies w(h) by psi\_H(u) and the Bessel coefficient by its inverse; thus the summand is independent of the representative.

\subsubsection{Theorem R28.8}\label{theorem-r28.8}

The map E\_B is a linear right inverse of restriction to H. Its image is in the Pi-isotypic component of L²(G), and has the full right (U,psi) character. It intertwines left H and the required central character. Its exact probability-Hilbert norm is

\[
       \|\mathcal E_Bw\|_G^2
          ={d_I\over d_\Pi}\|w\|_H^2
          \le\|w\|_H^2. \tag{R28.8}
\]

For any number field L containing all values of chi\_Pi and psi, and every prime ideal l of O\_L over a rational prime different from the characteristic of F\_q, it maps the value lattice

\[
 W_\psi\cap O_{L,l}^{\,H}
       \longrightarrow e_\Pi L^G\cap O_{L,l}^{\,G}.
 \tag{R28.9}
\]

Thus restriction is a split surjection between these actual function lattices at l; no claim that the central projector e\_Pi itself is integral at l is required.

\textbf{Proof.} For x in H, the only nonzero summand in (R28.7) is the coset xU\_H, by (R28.5). Write x=hu. That summand equals w(h)psi\_H(u)=w(x), proving the restriction identity. The right-U and central transformations follow from (R28.4). Every summand is a matrix coefficient of Pi, so e\_Pi acts as the identity. For left translation L\_aw(h)=w(a\^{}(-1)h), changing the right-coset variable shows E\_B L\_a=L\_a E\_B for every a in H.

The vectors Pi(h)v, for h in H/U\_H, are orthonormal. Their mutual inner product is B\_Pi(h\^{}(-1)k), and (R28.5) gives orthogonality and unit diagonal. In the chosen Hermitian convention,

\[
 (\mathcal E_Bw)(g)=
   \left\langle\sum_{h\in H/U_H}\overline{w(h)}\Pi(h)v,
                                                   \Pi(g)v\right\rangle.
\]

Finite Schur orthogonality therefore gives norm squared (sum\_h\textbar w(h)\textbar²)/d\_Pi. The numerator is d\_I\textbar\textbar w\textbar\textbar\_H² by the probability convention. This proves the equality in (R28.8); LIFT32's inclusion I\_n into Pi gives d\_Pi\textgreater=d\_I and the inequality. This inverse need not minimize the norm among all right-U\_H extensions: its extra right-U character is a genuine additional constraint.

Finally chi\_Pi(gu) and psi(u) are algebraic integers, and psi(u) is a unit. Formula(R28.3) has only \textbar U\textbar, a power of the residue characteristic, in its denominator. Formula(R28.7) is a finite sum with no averaging or additional divisor. This proves (R28.9). QED.

\subsection{3. Fibres, stabilizers and the actual goodTuple block}\label{fibres-stabilizers-and-the-actual-goodtuple-block}

The construction applies componentwise to any finite-dimensional coefficient space V with a fixed value lattice. Its left-H equivariance preserves every relation w(sh)=rho(s)w(h) for a subgroup S of H, with the same fibre action rho. Its central character similarly preserves the relation for a scalar subgroup Z acting on the fibre by omega\_Pi. No average over S or Z is used, so no stabilizer-order denominator is introduced by this operation.

Tensor products give the same statements on the actual product groups Gcal and Hcal. Choose the genuine Pi inside e0 with the aggregate scalar character Omega supplied by DO042/043. R027 proves that each actual coarse-packet stabilizer is SH times the diagonal global mu28. Consequently (R28.7), with the product Bessel function, is an explicit integral-at11 substitute for R27.8 on that packet. It preserves all its fibre relations and its exact restricted value. The full e0 acts as the identity on its image. No coefficientwise integrality of e0 is asserted or needed for this finite extension.

This is a statement about the finite orbit value/jet module. Identifying its pointwise integral lattice with the image of global integral cusp sections is a separate problem. It is not solved by an integral finite extension, even though the characteristic-zero interpolation of R027 can use that extension as its prescribed target datum.

\subsection{4. Exact finite cost for the two current hosts}\label{exact-finite-cost-for-the-two-current-hosts}

Each of the two residue primes occurs at f places, with (n,f)=(43,42) or(169,168). A sufficient global scalar clearing all coefficients of the product Bessel kernel is

\[
 D_B=(g_0g_1)^{\,f n(n-1)/2}. \tag{R28.10}
\]

Its normalized11-adic valuation is exactly zero. The other-prime denominator costs are retained. The following strict binary upper bounds follow from the actual prime bit lengths:

\begin{longtable}[]{@{}lrr@{}}
\toprule\noalign{}
Host & exponent of each g\_i in D\_B & log2(D\_B), strict upper bound \\
\midrule\noalign{}
\endhead
\bottomrule\noalign{}
\endlastfoot
J14 & 37,926 & 1,744,596 \\
J56 & 2,384,928 & 169,329,888 \\
\end{longtable}

For a common value field one can use L=Q(zeta\_M), where

\[
 M=\operatorname{lcm}_{i=0,1}
       \left(g_i\operatorname{lcm}_{1\le d\le n}(g_i^d-1)\right).
 \tag{R28.11}
\]

Indeed every actual g\_i is larger than n.~The unipotent part of a matrix therefore has order dividing g\_i, while the semisimple eigenvalues have orders dividing some g\_i\^{}d−1 with d\textless=n.~Thus M is divisible by all group element orders and by the orders of psi. It follows that all required finite character values lie in L. In particular the explicit, nonminimal bounds are

\[
 [L:Q]\le M<2^{43562}\quad(J14),\qquad
 [L:Q]\le M<2^{1019986}\quad(J56). \tag{R28.12}
\]

For example, replace each lcm by its product and use 1+sum\_(d=1)\^{}n d, multiplied by the sum of the two prime bit lengths. The field is independent of the word and canonical twist. It is not identified with the previously fixed arithmetic value field without the corresponding compositum and degree cost.

Every complex conjugate of B\_Pi(g) has absolute value at most1: a Galois conjugate replaces Pi and psi by another generic irreducible and a nondegenerate character, so the rank-one unitary coefficient argument applies again. Accordingly its absolute logarithmic height is at most log\textbar U\textbar, and a product-kernel coefficient has height at most log D\_B. This uses the displayed global clearing scalar and all-embedding bound; it does not declare that denominator optimal.

The exact finite Hilbert factor can also be evaluated for DO042's specific factor-degree pattern. At every full-degree n place Pi is cuspidal and d\_Pi=d\_I. At the one p+6 place in each G half, where p=n−6, its degree is

\[
 d_\Pi={\prod_{j=1}^n(g_i^j-1)
                  \over(g_i^p-1)(g_i^6-1)}.
 \tag{R28.13}
\]

Here the cuspidal degree formula is the specialization of Deligne--Lusztig(1976), Theorem7.1 and Proposition7.4, pp.140--141, to a regular character of the degree-p or degree-6 torus. Ordinary parabolic induction multiplies the two degrees by its Gaussian index, giving(R28.13). Distinct cuspidal factors make that induced representation irreducible, as in LIFT32 and GGP §4. There are exactly two such mixed places, hence

\[
 \|\mathcal E_Bw\|_{\mathcal G}^2
 =\kappa\|w\|_{\mathcal H}^2,\qquad
 \kappa=\prod_{i=0,1}
        {(g_i^{n-6}-1)(g_i^6-1)\over g_i^n-1}<1.
 \tag{R28.14}
\]

The scalar is positive. Its strict upper bound follows by expanding the numerator as g\_i\^{}n−g\_i\^{}(n−6)−g\_i\^{}6+1\textless g\_i\^{}n−1. Changing the two free full-degree labels to prescribe Omega does not change their dimensions. Thus(R28.14) is uniform for that entire DO043 selection.

\subsection{5. Sources, geometric meaning and remaining costs}\label{sources-geometric-meaning-and-remaining-costs}

The generic multiplicity-one input and Green labels are Gow, R., \& Vinroot, C. R.(2007), \emph{Extending real-valued characters of finite general linear and unitary groups on elements related to regular unipotents}, arXiv:0704.2022v1, Theorems3.2/4.3, pp.9/12. LIFT32 already checked the mirabolic induction with its exact isometry; GGP(2012), Proposition4.2/Corollary4.3, author-preprint pp.10--11, are the published branching inputs. Deligne, P., \& Lusztig, G.(1976), \emph{Representations of reductive groups over finite fields}, Annals of Mathematics,103(1),103--161, Theorem7.1 and Proposition7.4, pp.140--141, supply the cited degree. Primary URLs are retained in LIFT32/R026; the DL original is \url{https://publications.ias.edu/sites/default/files/Number27.pdf} .

The explicit integral splitting concerns finite etale residue actions and their characteristic-zero coefficient modules. It strengthens an equality of representation classes to a concrete split restriction on localized value lattices. The inherited frameworks are SGA3I4.7.3/4.7.3.1 and SGA6VI3.3; its use on principal parts retains EGAIV.4 16.11.2--3(ii), pp.54--55. The fixed Picard/biextension context remains FGA232 Theorem3.1 pp.148--149, FGA236 Corollary1.7 p.229/2.3 p.232, SGA7I VII2.7/2.9 pp.169--174 andVIII7.4(b) pp.310--311.

No denominator at11 is lost in the finite extension itself. This does not remove PBW scalars, global integral jet evaluation, arithmetic fibre-frame changes, integral saturation of the vertical theta image, or any global Petersson inverse. The G-prime denominators and explicit coefficient-field enlargement above remain real costs. No favorable total DP44 balance, algebraic-cycle nonvanishing, ell-adic weight or Kummer RH conclusion follows from this finite splitting.

\clearpage

\section{AM-R-029: a genuine goodTuple block whose projector is integral at11}\label{am-r-029-a-genuine-goodtuple-block-whose-projector-is-integral-at11}

Author: GPT-6 Astra Pro, ALP\_MACHINE /root, 22 November2026UTC. Reasoning effort is not separately exposed.

The full source e0 is unchanged. This constructs a block inside that same projector with any prescribed aggregate mu28 character and with an integral central projector at11. It refines the finite choice in DO042 by replacing its distinguished p+6 polynomial by p+3+3. No global cusp occurrence, integral jet surjectivity or theta-image saturation is inferred from this finite construction.

\subsection{1. Two separated strongly good cubic orbits}\label{two-separated-strongly-good-cubic-orbits}

Retain the exact four G primes, N=43 or169, and all character values and their inverses in DO042(DO42.1--4). At each place let C have at most four forbidden tame values, as in that statement. A strongly good root of degree d has full Frobenius degree d, retains that degree under every prime-to-l projection with l outside BB(N,g), and avoids every projection of every value of C, for every prime l other than g. These are the exact three requirements in DO042 §4.

\subsubsection{Lemma R29.1}\label{lemma-r29.1}

For each actual g and every such C, more than four fifths of F\_(g\^{}3)\^{}times are strongly good cubic roots. One can choose two distinct Frobenius orbits of these roots such that their prime-to-l images remain disjoint for every l outside BB(N,g).

\textbf{Proof.} Put b=bitlength(g), M=g\^{}3−1. The cyclic projection and character-collision argument DO42.5--10 is unchanged in degree3: the forbidden tame values exclude at most(40/203)M elements. The full-degree failure count is at most3g\^{}(3/2). For each prime l outside BB(N,g) dividing M, the projection kernel has order at most g\^{}(3/4). The inverse image of all proper subfields therefore has at most3g\^{}(9/4) elements. There are fewer than3b distinct prime factors of M. Consequently the excluded fraction is at most

\[
 {40\over203}
   +{3g^{-3/2}+9b\,g^{-3/4}\over1-g^{-3}}.       \tag{R29.1}
\]

This is the same elementary count as DO42.11--13 with its actual prime-factor bound3b in place of108; the degree restriction d\textgreater=6 there was used in its subsequent numerical estimate, not in these cyclic or subfield counts.

For the two smaller primes, b=23 and g\textgreater2\^{}22. The exact integer inequality90000\^{}4\textless2\^{}66 gives g\^{}(-3/4)\textless1/90000. For the two larger primes, b\textless=36 and g\textgreater2\^{}34, and 36·2\^{}(-25)\textless23/90000. In all four cases(R29.1) is strictly less than

\[
 {40\over203}+
 {3\cdot2^{-33}+207/90000\over1-2^{-66}}< {1\over5}.
                                                        \tag{R29.2}
\]

All comparisons in(R29.2) are rational. This proves the density.

Choose one strongly good root alpha. For a second root beta, exclude the three conjugates of alpha and every inverse image of their prime-to-l projections, for all nonexceptional l dividing M. For each l this excludes at most3g\^{}(3/4) values, and there are fewer than3b such primes. Thus at most

\[
                    3+9b\,g^{3/4}                        \tag{R29.3}
\]

values are additionally excluded. For l not dividing M the projection is the identity; those collisions are already excluded by the three conjugates themselves. Frobenius commutes with every prime-to-l projection, so the displayed exclusion covers every pair of conjugates from the two orbits, not only their chosen representatives.

The fraction in(R29.3) is at most

\[
 {3\cdot2^{-66}+324\cdot2^{-49}\over1-2^{-66}}
                  <{1\over1000}.                         \tag{R29.4}
\]

Indeed b\textless=36 and g\textgreater2\^{}22 give g\^{}(-9/4)\textless2\^{}(-49). The strongly good set has fraction greater than4/5, so it contains a beta outside this additional set. Its two Frobenius orbits have the claimed full prime-to-part separation. QED.

\subsection{2. The complete goodTuple and its central character}\label{the-complete-goodtuple-and-its-central-character}

At the distinguished place in each G half, take one strongly good degree-p orbit, p=N−6=37 or163, and the two cubic orbits of Lemma R29.1, all for that place's actual C. Set

\[
                f_v=P_\alpha P_\beta P_\gamma,
             \qquad(\deg P_\alpha,\deg P_\beta,\deg P_\gamma)
                            =(p,3,3).                    \tag{R29.5}
\]

At every other place keep a strongly good irreducible degree-N polynomial as in DO042. This is a literal goodTuple for both character signs.

\textbf{Proof.} All polynomials are monic of degree N with nonzero constant term. Their factors are distinct: the two cubic orbits were separated, and the p orbit has a different degree. At a nonexceptional prime, each orbit retains its full degree, the two cubic projected orbits stay disjoint, and a degree-p projected orbit cannot equal a degree-3 projected orbit. This proves the whole injectivity clause. The tame character-avoidance clause holds for every root by strong goodness, including exceptional primes.

There is a degree-N irreducible place in each half, which rules out every factor degree a with1\textless=2a\textless=N. The distinguished place has its prime-degree factor p\textgreater N/2 with p+3=N−3\textless=N. These are exactly the two remaining factor-degree clauses. Thus every conjunct of StatementT.lean19--29 holds, for the characters and their inverses. No clause is relaxed. QED.

Keep the distinguished place fixed. In one G half at least two other degree-N places remain available. DO043's norm-compatible central-character argument applies without any change to those two free places: each realizes at least23 of the28 scalar characters, hence their products realize every prescribed aggregate mu28 character. It follows that for every Omega the unchanged e0 contains a genuine generic product irreducible Pi of the pattern (R29.5) with central character Omega.

As in DO042, a finite lexicographic choice in the nonempty root sets specifies the polynomials. This is a proved finite construction, not a claim that all these enormous fields were enumerated.

\subsection{3. Integral central projection at11}\label{integral-central-projection-at11}

The exact orders of the four primes in F\_11\^{}times are

\begin{longtable}[]{@{}lrrr@{}}
\toprule\noalign{}
Host & g & g mod11 & multiplicative order mod11 \\
\midrule\noalign{}
\endhead
\bottomrule\noalign{}
\endlastfoot
J14 & 4615409 & 7 & 10 \\
J14 & 5276377 & 7 & 10 \\
J56 & 22285932761 & 4 & 5 \\
J56 & 64072175441 & 10 & 2 \\
\end{longtable}

None of these orders divides its corresponding N, p=N−6, or3. Therefore every torus order in the following formula is an11-unit:

\[
 T_v=\begin{cases}
       g^N-1,&\text{at a full-degree place},\\
       (g^p-1)(g^3-1)^2,&\text{at the distinguished place}.
     \end{cases}                                        \tag{R29.6}
\]

The genuine squarefree-series irreducible Pi\_v has degree

\[
 d_{\Pi_v}={\prod_{j=1}^N(g^j-1)\over T_v},\qquad
 {d_{\Pi_v}\over|GL_N(F_g)|}
          ={1\over g^{N(N-1)/2}T_v}.                    \tag{R29.7}
\]

This is the regular torus degree of Deligne--Lusztig Theorem7.1/ Proposition7.4, or equivalently the product of its distinct cuspidal degrees and the exact parabolic index. Repeating degree3 is allowed: the two cubic Green labels are distinct, so this parabolic induction is irreducible. LIFT32 identifies it with the genuine eSer(f\_v).

Let L contain all character values, for example R28.11's single fixed cyclotomic field. At every prime ideal l above11,

\[
 e_{\Pi_v}={d_{\Pi_v}\over|GL_N(F_g)|}
       \sum_{x\in GL_N(F_g)}\overline{\chi_{\Pi_v}(x)}\,x
              \ \in O_{L,l}[GL_N(F_g)].                 \tag{R29.8}
\]

Indeed its character values are algebraic integers and its displayed denominator is prime to11. Taking products proves

\[
 e_\Pi\in O_{L,l}[\mathcal G],\qquad
 e_\Pi^2=e_\Pi,\quad e_0e_\Pi=e_\Pi,\quad
                  p_\Omega e_\Pi=e_\Pi.                 \tag{R29.9}
\]

These identities hold in the genuine characteristic-zero group algebra and its stated localization. Thus e\_Pi preserves every Gcal-stable O\_(L,l) lattice and its image is a direct summand. Its Hilbert operator norm is1 wherever the summand occurs. No claim that the entire e0 has integral coefficients is made. Neither is this an assertion that a particular global cusp lattice is already identified with its geometric integral model.

Combining with R028 gives both an integral finite extension and an integral projection onto its prescribed Pi block at11. If a Gcal-equivariant global integral jet map is surjective, then applying this actual integral idempotent preserves surjectivity onto its Pi summand, without a new11 denominator. The hypothesis about that global integral map remains unproved.

\subsection{4. The retained other-prime and metric costs}\label{the-retained-other-prime-and-metric-costs}

A sufficient clearing scalar for this product projector is

\[
 D_\Pi=(g_0g_1)^{\,fN(N-1)/2}
     \prod_{i=0,1}
        (g_i^N-1)^{f-1}(g_i^{N-6}-1)(g_i^3-1)^2.
                                                        \tag{R29.10}
\]

Its normalized11 valuation is zero. The other prime divisors have not disappeared. The exact bit-length estimates are

\[
 \log_2 D_\Pi<
 \begin{cases}
 1827672,&J14,\\
 171345720,&J56.
 \end{cases}                                             \tag{R29.11}
\]

For each g the sum of the torus degrees is fN, so its upper exponent is fN(N+1)/2; multiply by the two prime bit lengths as in R028. The same fixed field bounds R28.11--12 apply. Every conjugate coefficient of e\_Pi has modulus at most1: \textbar chi\_Pi(x)\textbar\textless=d\_Pi and d\_Pi²\textless=\textbar Gcal\textbar{} by the finite character degree sum. Consequently each coefficient has absolute logarithmic height at most log D\_Pi. No minimal field or optimal denominator is claimed.

R028's Bessel inverse for this different admitted block has the exact squared probability-norm factor

\[
 \kappa_3=\prod_{i=0,1}
   {(g_i^{N-6}-1)(g_i^3-1)^2\over g_i^N-1},\qquad
                          0<\kappa_3<1.                 \tag{R29.12}
\]

The two mixed places are the only noncuspidal factors; their degrees in(R29.7) give this identity. It holds for all scalar characters Omega selected using the two free full-degree places. All comparisons are at the finite probability measures; global Petersson interpolation and vertical inverse norms remain separate.

\subsection{5. Sources, scope and geometric framework}\label{sources-scope-and-geometric-framework}

The exact source clauses are \nolinkurl{CyclotomicNorms/StatementT.lean19–29} and \nolinkurl{DataN.lean960–966/996–1007.} K045 retains BB-disjointness for the unchanged primes and ranks. DO042(DO42.5--13) supplies the cyclic-projection and forbidden-character counts; only their degree3 estimate and the two-orbit separation are new. DO043 supplies the two-free-place central character argument. R026/LIFT32 supply the genuine finite series, and R028 supplies the explicit Bessel inverse.

The literature degree input is Deligne, P., \& Lusztig, G.(1976). Representations of reductive groups over finite fields. Annals of Mathematics,103(1),103--161, Theorem7.1 and Proposition7.4, pp.140--141. The exact Green and generic parametrizations remain Gow--Vinroot(2007), Theorems3.2/4.3, pp.9/12, and the parabolic branching GGP(2012), Proposition4.2/Corollary4.3, pp.10--11 of the author version. These inputs have already been checked in R028 and LIFT32; no new re-proof of their general theorems is claimed.

This is an actual integral finite residue-group correspondence over O\_L localized at11. It is stronger than its representation-ring class, and it preserves equivariant principal-parts modules when those modules are supplied. The inherited exact framework locators are SGA3I4.7.3/4.7.3.1, SGA6VI3.3 and EGAIV.4 16.11.2--3(ii), pp.54--55; the fixed FGA232/236 and SGA7I VII--VIII Picard/biextension comparisons remain unchanged.

The new block lies inside the original e0 and uses the original finite hosts, primes, generic character and every goodTuple clause. It does not finish the source same-projector arithmetic record, an admitted global detector, the integral theta inverse, the exact DP44 difference or the specific032 arithmetic comparison. No motivic nonvanishing, ell-adic purity or Kummer RH result is inferred.

\clearpage

\section{AM-K046 --- exact selected Bergman concentration and one-point interpolation}\label{am-k046-exact-selected-bergman-concentration-and-one-point-interpolation}

Author: GPT-6 Astra Pro, ALP\_MACHINE /root/kummer\_extension, 22 November 2026 UTC. Inherited reasoning effort is not separately exposed.

At the upper gap endpoint of K044, the actual selected bottom component has an explicit uniform L1 concentration estimate. The estimate keeps the exact PBW normalization of K043. It is uniform in every binary word, every allowed canonical twist, every CM role, and every vector in the selected component, including vectors which are not pure tensors.

For the standard principal-level component of R025, an elementary algebraic-integer argument verifies the separation condition needed by the Poincare criterion. This gives an actual nonzero cusp Poincare series and a quantitative right inverse for the selected fibre at that one standard point. This component and this point are explicitly specified. The actual detector packet in R027 is at (o,x tau), with the nonintegral T translation tau. The separation proof here does not cover that packet, the finite e0 projection, or simultaneous interpolation on the packet. It gives no cusp-dimension upper bound or DP44 cost sign.

\subsection{1. The actual component and its normalization}\label{the-actual-component-and-its-normalization}

Fix one of K044's two hosts, with a=f equal to42 or168, and use its upper endpoint ell=ell\_+. Write s for the allowed canonical twist and q=(a+1)s, so that 0\textless=q\textless=ell. At a CM label tau of role T,H or L, put E=E\_T,E\_H orE\_L as in K044(EG1). Retain all its word formulas.

Use the zero-based K043 notation p\_i=v\_i for i\textless a-1, p\_(a-1)=-b0, and p\_-=-b1. The minimal compact weights and selected weights are

\[
 \alpha_i=p_i-(a-2)/2+i,\qquad
 \beta=p_--a/2,\qquad
 \lambda_i=u_i-(a-1)/2+i.
 \tag{BL1}
\]

The Pieri exponents and norm factors are

\[
 m_i=\lambda_i-\alpha_i=u_i-p_i-1/2,\qquad
 r_i=\alpha_i-\beta-i=p_i-p_-+1,
 \qquad D_\tau=\prod_{i=0}^{a-1}(r_i)_{m_i}.
 \tag{BL2}
\]

All m\_i are nonnegative integers and the actual selected Pieri component has multiplicity one. Its Fischer metric has monomial norm gamma!, and its invariant Hilbert metric is the Fischer metric divided by D\_tau. These are K043(PB1--PB8), obtained from Nakahama(2016), Theorem5.1, printed pp.960--961. In particular, for a Fischer-isometric selected inclusion j\_tau, the unit selected inclusion in the polynomial Hilbert model is

\[
             J_\tau=\sqrt{D_\tau}\,j_\tau,
             \qquad \rho_\tau=\sqrt{D_\tau}.
 \tag{BL3}
\]

This uses the matrix-unit sign convention of K043. Its unit phase has no effect on any norm below. A nonisometric choice of j\_tau must retain the scalar c\_j of K043(PB8); it is not used here.

Let V\_sel be the tensor product of the selected compact modules at all f labels. For a unit vector v in V\_sel, set

\[
 P=\rho_{\rm all}j(v),\qquad
 \rho_{\rm all}=\prod_\tau\sqrt{D_\tau},\qquad
 j=\bigotimes_\tau j_\tau.
 \tag{BL4}
\]

Then P has Hilbert norm1. The formula applies to every v, not only to decomposable vectors. The raw derivative associated with j(v), divided by rho\_all, is the normalized selected functional of K043(PB9--10). No large PBW factor is omitted or estimated as a unit.

\subsection{2. Haar measure, fibre metric and exact radial moments}\label{haar-measure-fibre-metric-and-exact-radial-moments}

On the ball B\^{}a put Y(z)=1-\textbar z\textbar\^{}2 and use invariant measure

\[
                  d\mu(z)=\pi^{-a}Y(z)^{-a-1}\,dV(z).
 \tag{BL5}
\]

Normalize compact Haar measure to1. In Nakahama's conventions the normalizing constant for the q=1 case of Theorem5.1 is

\[
              c_\tau=\prod_{i=0}^{a-1}(r_i-1),
              \qquad C=\prod_\tau c_\tau.
 \tag{BL6}
\]

Indeed his constant is product\_j(Lambda-(j-1)+k\_j-1), with Lambda=-beta-c and k=alpha+c1, so its factors are exactly r\_i-1. His displayed metric(5.7), printed p.960, uses pi\^{}(-a)dV and the power Y\^{}(Lambda-a-1). Thus (BL5--6) retain the Haar and metric constants, rather than an unspecified positive multiple of the norm.

Write F\_P(g) for the vector-valued canonical lift of P to the product of the real unitary groups. In the holomorphic model this is evaluation at0 after the unitary translate of P. Equivalently it is the canonical fibre metric times P(g.0), as in Zunar(2025), equations(5.9--5.13), printed pp.9--10. Put b(g)=\textbar\textbar F\_P(g)\textbar\textbar. Then

\[
                0\le b(g)\le1,\qquad
                I:=\int b(g)^2\,dg=1/C.
 \tag{BL7}
\]

For the first inequality, evaluation at0 is the orthogonal projection onto the constant compact type, whose Hilbert metric is the given fibre metric. It is a contraction. Unitarity and \textbar\textbar P\textbar\textbar=1 prove the bound. The second identity is exactly (BL5--6) and the norm formula.

Give the product of balls the probability measure b\^{}2 dg/I. For an integer j\textgreater=0 smaller than every r\_i-1 at a fixed label, its negative radial moment is

\[
 M_\tau(j):={\bf E}(Y_\tau^{-j})
      =\prod_{i=0}^{a-1}
          \frac{(r_i+m_i-j)_j}{(r_i-1-j)_j}.
 \tag{BL8}
\]

More generally, all nonnegative integer moments are

\[
 {\mathbb E}(Y_\tau^h)
      =\prod_i\frac{(r_i-1)_h}{(r_i+m_i)_h}
                      \quad(h\ge0),
 \tag{BL9}
\]

and joint moments at distinct labels are products of these expressions.

\textbf{Proof.} Holding the polynomial, its fibre metric and its compact Pieri component fixed, replacing Lambda by Lambda+h multiplies the integrand in Nakahama's metric by Y\^{}h. The unnormalized integral is the Fischer norm divided by c\_(Lambda+h)D\_(Lambda+h). Dividing by its value at h=0 gives

\[
 \frac{\prod_i(r_i-1)\,(r_i)_{m_i}}
      {\prod_i(r_i+h-1)\,(r_i+h)_{m_i}}.
 \tag{BL10}
\]

Cancelling factorials gives(BL9), and using h=-j gives(BL8). The shifted metric remains in the convergent chamber because j\textless min\_i(r\_i-1) is equivalent to the last shifted starting factor being positive. On an irreducible selected component each shifted metric is the stated scalar times the Fischer metric. On its tensor product the scalars multiply. This proves the joint-moment assertion also for an entangled v in(BL4). QED.

For another exact description, the law of Y\_tau is the law of the product of independent Beta(r\_i-1,m\_i+1) variables. Their moments are(BL9), and equality of all polynomial moments on{[}0,1{]} determines the probability measure. The joint law factors for the same reason. This description is not needed for the estimate. It is a useful normalization check: in the scalar case r\_i=Lambda-i and m=(N,0,\ldots,0), the product telescopes to

\[
           {\mathbb E}(Y^h)=
             \frac{(\Lambda-a)_h}{(\Lambda+N)_h},
 \tag{BL11}
\]

the Beta(Lambda-a,N+a) radial law. In particular the dimension shift a is retained.

\subsection{3. Uniform exact moment bounds for the actual profiles}\label{uniform-exact-moment-bounds-for-the-actual-profiles}

Set

\[
                   t=\ell/2,\qquad s_0=a+1.
 \tag{BL12}
\]

Both are integers, and t\textgreater s0. Direct substitution of the K044 words in (BL2) gives, for 0\textless=i\textless a-1,

\[
 r_i=E+(4a-3-4i)\ell+q+1,
 \qquad r_{a-1}=2\ell+q+1,
 \tag{BL13}
\]

and, for 1\textless=i\textless a-1, writing omega\_i for the bit in interior row i (the bit denoted omega\_(i-1) in K044),

\[
 m_i=4\ell-1-\ell\omega_i.
 \tag{BL14}
\]

The endpoint exponents are

\begin{longtable}[]{@{}lll@{}}
\toprule\noalign{}
Role & m0 & m\_(a-1) \\
\midrule\noalign{}
\endhead
\bottomrule\noalign{}
\endlastfoot
top & A+(a-2)/2-4ell(a-1) & E\_T-B-ell-a/2 \\
thick or thin & 4a ell & E+ell \\
\end{longtable}

At the exact upper endpoint, the last top exponent is19 for J14 and12 for J56. In particular this exponent does not silently grow with ell. All r\_i-1 are at least2ell, so(BL8) is valid through j=t.

Define a rational number, using q=0 and the zero word, by

\[
        A_\tau=
        \prod_{i=0}^{a-1}\frac{r_i+m_i}{r_i-1-t}.
 \tag{BL15}
\]

For every actual word, every allowed q and every integer0\textless=j\textless=t,

\[
                       M_\tau(j)\le A_\tau^j.
 \tag{BL16}
\]

To see this, each numerator factor in(BL8) is at most r\_i+m\_i, and each denominator factor is at least r\_i-1-t. Replacing an interior zero bit by1 only decreases m\_i. Increasing q adds the same nonnegative number to numerator and denominator of every ratio in(BL15); because the numerator is larger, the ratio decreases. Thus the zero word and q=0 give a simultaneous upper bound.

Exact integer product comparisons give the following strict bounds. They use the exact rational exponential rounding already certified in K044 for E\_T,E\_H,E\_L; they do not use floating point values.

\begin{longtable}[]{@{}lrrr@{}}
\toprule\noalign{}
Host & top A\_tau \textless{} & thick A\_tau \textless{} & thin A\_tau \textless{} \\
\midrule\noalign{}
\endhead
\bottomrule\noalign{}
\endlastfoot
J14 & 45 & 969 & 222 \\
J56 & 45 & 1722 & 1059 \\
\end{longtable}

In particular

\[
                    M_\tau(j)\le 2^{11j}
                    \quad(0\le j\le t).
 \tag{BL17}
\]

For an independent check of the indices, the zero-word first negative moment telescopes. Its denominator is (E+5ell+q-1)(2ell+q-1). For non-top roles its numerator is (E+(8a-3)ell+q)(E+3ell+q). For the top role its numerator is (A+E+ell+q+(a-2)/2)(E\_T-B+ell+q-a/2). This follows because each interior numerator r\_i+m\_i-1 equals r\_(i-1)-2. The certificate verifies both forms exactly.

\subsection{4. L1 concentration, not only L2 concentration}\label{l1-concentration-not-only-l2-concentration}

The moment calculation directly concerns b\^{}2. The following argument converts it to the L1 estimate required by the Poincare criterion.

Let W=product\_tau Y\_tau\^{}s0. The elementary radial integral is

\[
 \int_{B^a}Y^{s_0}\,d\mu
      =\frac{\Gamma(s_0-a)}{\Gamma(s_0)}=\frac1{a!},
 \qquad V_W:=\int W\,dg=(a!)^{-f}\le1.
 \tag{BL18}
\]

This follows by polar coordinates with x=\textbar z\textbar\^{}2: pi\^{}(-a)dV gives x\^{}(a-1)dx/(a-1)!, and the remaining integral is the ordinary beta integral. All r\_i\textless2\^{}960 for the actual profiles and all q\textless=ell, so

\[
                            C<2^{960af}.
 \tag{BL19}
\]

First, Cauchy--Schwarz and(BL17--18) show that b is in L1 and

\[
 \frac1C\le\int b\,dg
 \le C^{-1/2}(a!)^{-f/2}2^{11f(a+1)/2}.
 \tag{BL20}
\]

The left inequality uses b\textgreater=b\^{}2. For the right one use integral b \textless=(integral b\^{}2 W\^{}-1)\^{}(1/2)(integral W)\^{}(1/2), and the joint moment formula.

Fix epsilon=2\^{}-12 and a label tau. On the event Y\_tau\textless epsilon,

\[
 1_{Y_\tau<\epsilon}Y_\tau^{-s_0}
       \le\epsilon^{t-s_0}Y_\tau^{-t}.
\]

Apply Cauchy--Schwarz again, then(BL8),(BL17--20). The fraction of L1 mass on this event is at most

\[
 \begin{split}
 \frac{\int_{Y_\tau<\epsilon}b}{\int b}
 &\le\left(
       C V_W\epsilon^{t-s_0}
       M_\tau(t)\prod_{\tau'\ne\tau}M_{\tau'}(s_0)
                    \right)^{1/2}\\
 &\le 2^{(B_0-t)/2},\qquad
 B_0=960af+(11f+1)(a+1).
 \end{split}
 \tag{BL21}
\]

The exact upper-endpoint integers satisfy4B0\textless ell. Since t=ell/2, the last bound is strictly smaller than2\^{}(-ell/8). Summing over the f labels proves

\[
 \frac{\displaystyle
       \int_{\{\exists\tau:Y_\tau<2^{-12}\}}b(g)\,dg}
      {\displaystyle\int b(g)\,dg}
          <\delta:=f\,2^{-\ell/8}<2^{-100}.
 \tag{BL22}
\]

There is no exchange of an L2 probability estimate for an L1 estimate without proof: the weight W and both Cauchy--Schwarz factors in(BL21) are essential. The finite certificate checks the exact domination 4B0\textless ell and the final union bound for both hosts.

For the Cartan coordinate u\textgreater=0 of SU(a,1), \textbar z\textbar=tanh u. Therefore Y\textgreater=2\^{}-12 implies

\[
 u\le\operatorname{arctanh}\sqrt{1-2^{-12}}
      <\tfrac12\log(4\cdot2^{12})=7\log2<5.
 \tag{BL23}
\]

For the last elementary estimate, e\^{}(7/10)\textgreater2 follows already from the first four positive terms of its exponential series. Consequently log2\textless7/10. Let C5 be the product of the Cartan sets K exp({[}0,5))K. Then(BL22) holds with the complement of C5 in the numerator. This proves uniform concentration on the complete selected component.

\subsection{5. A quantitative Poincare lemma}\label{a-quantitative-poincare-lemma}

The known general nonvanishing criterion used here is Zunar(2025), Proposition8.1, printed pp.15--16, a finite-dimensional vector version of Muic's criterion. Its hypotheses are a unimodular second-countable group, a discrete subgroup Gamma, an L1 vector function F, and a Borel set C with CC\^{}-1 intersect Gamma=\{1\}, containing more than half of the L1 mass. No Gaussian-integer special case is being transferred to the present number fields.

The following quantitative form follows from its elementary proof. If the L1 mass outside C is at most delta\textbar\textbar F\textbar\textbar1, delta\textless1/2, then

\[
          \|P_\Gamma F\|_{L^1(\Gamma\backslash G)}
               \ge(1-2\delta)\|F\|_{L^1(G)}.
 \tag{BL24}
\]

Indeed C injects into Gamma\G. On C subtract all nonidentity translates from the identity summand using the triangle inequality. The sets gamma C, gamma in Gamma, are disjoint, and the nonidentity ones lie outside C. Their total error integral is at most the L1 mass outside C. The integral of the identity term on C is at least(1-delta)\textbar\textbar F\textbar\textbar1. This proves(BL24). Absolute almost-everywhere convergence and the L1 periodization identity follow by Tonelli.

For our seeds, whenever C5 C5\^{}-1 intersect Gamma=\{1\},(BL20--24) give

\[
          \|P_\Gamma F_P\|_1\ge(1-2\delta)/C>0.
 \tag{BL25}
\]

To read the series as cusp sections, use the integrable discrete-series Poincare theorem of Muic(2019), Theorem6-4 and Theorem7-6(i), author version printed pp.21--22 and26. Its integrability hypothesis holds here. In type A\_a, for a noncompact root alpha, sum\_(roots eta)\textbar eta(H\_alpha)\textbar=4a. Zunar's criterion(5.4), printed p.8, therefore says that the positive HC gap p\_i-p\_- must exceed a. Its minimum is b1-b0=2ell+q\textgreater a. The sign is the highest-weight dual of(BL1), so the criterion's HC value is p\_--p\_i, not its negative. This also agrees with the scalar ball threshold Lambda\textgreater2a.

The K-finite vector F\_P is a finite sum of matrix coefficients with constant-type test vectors; Zunar(2025), Section6 and(7.4), supplies this exact holomorphic realization. Muic's theorem makes its Poincare sum L2, smooth and cuspidal for an arithmetic congruence lattice. The unitary groups have compact real centres; one may pass to their derived groups and then use the prescribed compact central character. This changes neither the radial metric nor the normalized compact Haar factor. In the standard congruence group below the rational determinant is in fact1, so the arithmetic group already lies in the derived group.

No effective global upper norm is inferred from Muic's theorem: his proof chooses an unspecified small-support reproducing function beta. It is used here for convergence, boundedness and cuspidality only.

\subsection{6. Separation at the actual standard principal-level component}\label{separation-at-the-actual-standard-principal-level-component}

Use the field E/F and standard lattice of R025(R25.1--3), and put

\[
 \Gamma_1=U(I_a,-1)(F)\cap J^1,
 \qquad J^1=\ker(J^{\max}\longrightarrow\mathcal G).
 \tag{BL26}
\]

All selected components over both rational primes occur in R025. Its exact split-unitary relation implies full rational congruence: every entry of gamma-1 for gamma in Gamma1 belongs to (g0g1)O\_E. Every entry of gamma itself is integral at all finite places, hence is an algebraic integer. These statements concern the standard global lattice O\_E\^{}(a+1), not an arbitrary translated lattice.

\subsubsection{Lemma BL27}\label{lemma-bl27}

For this Gamma1, C5 C5\^{}-1 intersect Gamma1=\{1\}.

\textbf{Proof.} In the standard Hermitian coordinates, a Cartan element of parameter u has Euclidean operator norm e\^{}u. Elements of K are Euclidean unitary. Thus, if gamma lies in C5 C5\^{}-1, every archimedean operator norm of gamma is less than e\^{}10. Every conjugate of every entry x of gamma-1 consequently has modulus less than e\^{}10+1. The chosen f complex embeddings and their conjugates exhaust all embeddings of E. Since e\textless3,

\[
 e^{10}+1<3^{10}+1=59050<g_0
 \tag{BL27}
\]

for each actual host. If x were nonzero, x/g0 would be a nonzero algebraic integer all of whose conjugates have modulus less than1. Its nonzero integer norm would then have absolute value less than1, a contradiction. Every entry is therefore zero. QED.

Likewise det(gamma) is an integral unit, all its conjugates have modulus1, and det(gamma)-1 is divisible by g0. Applying the same norm argument with the bound2\textless g0 proves det(gamma)=1. Thus the semisimple congruence-group hypothesis in §5 is available without an unproved determinant restriction.

Combining the lemma with(BL25) gives actual nonzero cusp Poincare series at the unchanged principal level J1, for every word and every canonical twist in the K044 interval at ell=ell\_+. No extra prime, principal29 refinement or changed finite lattice was used.

\subsection{7. One-point selected-jet Gram bound and right inverse}\label{one-point-selected-jet-gram-bound-and-right-inverse}

There is a useful consequence beyond individual nonvanishing. Let H\_cusp be the Hilbert space of the corresponding holomorphic cusp sections on the standard component Gamma1\G, with the invariant fibre metric and Haar convention of§2. Its contragredient is the source bottom type sigma\_s\^{}vee. Let

\[
             T:V_{\rm sel}\longrightarrow H_{\rm cusp},
             \qquad Tv=P_{\Gamma_1}P_v,
             \quad P_v=\rho_{\rm all}j(v).
 \tag{BL28}
\]

Here periodization is read in the canonical holomorphic frame; its lift is the vector Poincare sum of§5. Let H0 be the finite-dimensional image of T, and let J0 be the normalized selected jet at the standard origin restricted to H0. With inner products linear in the first slot, its defining pairing is

\[
       \langle J_0S,v\rangle
         =\rho_{\rm all}^{-1}
             \langle S_{\rm sel},j(v)\rangle_{\rm F},
 \tag{BL29}
\]

where S\_sel is the selected homogeneous Taylor summand in these Harish--Chandra coordinates. This is the K043 normalized derivative, not an unnormalized divided Taylor coefficient.

Unfolding gives the exact adjoint identity, with T* taken from H0 to V\_sel,

\[
                     J_0=C\,T^*.
 \tag{BL30}
\]

For clarity, pair S in H0 with Tv and unfold the defining series. The resulting integral over G is the canonical weighted pairing of S with P\_v. Averaging over the compact group first selects exactly S\_sel: its Pieri summand has multiplicity one, and its central compact weight fixes the homogeneous degree. Nakahama's metric on that summand then gives

\[
 \langle S,Tv\rangle_{\rm Pet}
       = C^{-1}\langle S_{\rm sel},P_v\rangle_{\mathcal H}
       = C^{-1}\rho_{\rm all}^{-1}
                           \langle S_{\rm sel},j(v)\rangle_{\rm F}.
 \tag{BL31}
\]

This proves(BL30). The unfolding is absolutely convergent: sections in H0 are finite linear combinations of the bounded Poincare sections of Muic's theorem, and the seed is L1. Compact averaging may first be done on compact radial subdomains; the absolute integral permits passage to the limit. No boundedness assertion about a completion of the entire cusp space is required for this finite Gram identity. The scalar C is forced by(BL6), not chosen to fit the periodized value.

Put V\_G=vol(Gamma1\G) in this Haar measure. By(BL24--25) and Cauchy--Schwarz on the finite-volume quotient, for every v,

\[
             \|Tv\|_2\ge
       \frac{1-2\delta}{C\sqrt{V_G}}\,\|v\|.
 \tag{BL32}
\]

Thus

\[
 J_0T=C T^*T\succeq
       \frac{(1-2\delta)^2}{CV_G}I,
 \qquad
 J_0J_0^*\succeq\frac{(1-2\delta)^2}{V_G}I.
 \tag{BL33}
\]

In particular J0 from H0 is onto the complete selected fibre; therefore the actual full cusp-section jet map is onto it as well. The right inverse with image in H0 satisfies

\[
       R_0=J_0^*(J_0J_0^*)^{-1}
           =T(C T^*T)^{-1},\qquad
       \|R_0\|\le\frac{\sqrt{V_G}}{1-2\delta}.
 \tag{BL34}
\]

The second inequality in(BL33) follows from J0=C T\emph{, so J0 J0}=C\^{}2 T*T. No inference that nonvanishing alone implies surjectivity is used: the uniform lower bound holds for every v, and the explicit Gram matrix is positive definite.

The PBW factors cancel in this normalized local adjoint calculation. For a raw derivative of unit Fischer vector j(v), the rho\_all of (BL29) must be restored. A different global automorphic realization or a source measure rescales the displayed maps by its prescribed scalar; it must be applied on both sides. Under a common positive Haar rescaling, the final bound(BL34) has the same form with the rescaled quotient volume. No period or finite class mass is hidden in C.

\subsection{8. Exact remaining arithmetic obstruction}\label{exact-remaining-arithmetic-obstruction}

R027's actual packet consists of points (o,x tau), where tau is the LIFT29 lower-unipotent T translation. Its arithmetic group is

\[
       \Gamma_{x\tau}=G(F)\cap x\tau J^1\tau^{-1}x^{-1}.
 \tag{BL35}
\]

Its standard-coordinate entries need not be integral at T or at the places in the ideal class of x. Consequently the norm argument in BL27 cannot be applied to gamma-1 in those old coordinates. Changing to an integral lattice frame changes the archimedean comparison; no bound for that change is supplied by the group-level statement that the lattice is selfdual. This is a concrete remaining separation condition, not a permission to replace tau by1.

Moreover the full finite deck packet includes distinct points which need not be separated by the radius5 neighbourhoods. The actual genuine e0 projection can mix their series. The one-point bound does not imply that this projection preserves its chosen jet, nor that the whole prescribed initial coefficient module can be interpolated in the allowed twist window. R027's finite equivariant right inverse is a different map from the unprojected one-point R0 in(BL34).

Nothing here bounds the multiplicity of the holomorphic discrete series in the bottom cusp space by a small multiple of volume times formal degree. Such an upper bound needs a trace/orbital estimate, a vanishing/Riemann--Roch bound, or another actual global argument. The local mass estimate alone does not resolve the possible full-packet dimension deficit being examined separately.

Thus the result closes the universal local concentration calculation and the standard one-point inverse, while reducing the needed actual detector estimate to the translated/projected packet Gram matrix. It does not count this partial inverse as the source's terminal norm or as a saving in its total gain/cost expression.

\subsection{9. Sources, certificate and check scope}\label{sources-certificate-and-check-scope}

Primary sources used directly:

\begin{itemize}
\tightlist
\item
  Nakahama, R. (2016). Norm computation and analytic continuation of vector valued holomorphic discrete series representations. \emph{Journal of Lie Theory, 26}, 927--990. Theorem5.1, pp.960--961; equation(5.7), p.960; action(5.5), p.959; Fischer metric(3.2), p.940. \url{https://jolt.centre-mersenne.org/item/10.5802/jolt.919.pdf}
\item
  Zunar, S. (2025). \emph{On the non-vanishing of Poincare series on irreducible bounded symmetric domains}. Author version arXiv2501.06876, Lemma5.1(vii), p.8; equations(5.9--5.16), pp.9--10; Sections6--7; Proposition8.1, pp.15--16; Theorem8.3, pp.16--17. The SU(p,q;Z{[}i{]}) specialization of Theorem9.3 is not applied to our number fields. \url{https://arxiv.org/pdf/2501.06876}
\item
  Muic, G. (2019). Smooth cuspidal automorphic forms and integrable discrete series. \emph{Mathematische Zeitschrift, 292}, 895--922. Author version arXiv1610.05483v2, Theorem6-4, pp.21--22, and Theorem7-6(i), p.26. \url{https://arxiv.org/pdf/1610.05483}
\end{itemize}

Internal proved inputs are K043(PB1--PB14), K044(EG1--EG10), and the standard lattice and full selected-component congruence of R025(R25.1--R25.3). The scope comparison uses LIFT29(BT4--BT8) and R027(R27.2,R27.5--R27.10). No accepted theorem is widened to an unproved new field or new period map.

The standalone standard-library script \nolinkurl{selected_bergman_localization_certificate.py} writes its exact checks and input hashes to the matching JSON file. It checks the scalar beta-law dimension shift, the positive and negative vector moment index identities, all six exact rational profile bounds, the first-moment telescoping, endpoint monotonicity form checks, the actual normalizer and gap inequalities, the L1 exponent and Cartan radius constants, and the actual principal-level numerical separation margin. The continuous monotonicity and probability/ Poincare arguments are proved in this report, not tested by endpoint sampling. The certificate computes no automorphic Gram matrix and does not certify the translated packet, e0 projection, cusp dimension, full source records or a Lean build.

\clearpage

\section{ALP-LIFT-33: sharper actual class denominators and an exact non-scalar middle pairing}\label{alp-lift-33-sharper-actual-class-denominators-and-an-exact-non-scalar-middle-pairing}

Author: GPT-6 Astra Pro, ALP\_MACHINE \nolinkurl{/root/lift_accounting}. 22 November 2026 UTC. Inherited reasoning effort is not separately exposed.

\textbf{Author whole-proof review passed at04:11 UTC; sharp-level refinement checked at04:18 UTC.} This finding gives a quantitative improvement on the actual coupled arithmetic map, and an actual example showing why its invariant lattices cannot all be treated as unimodular. It does not infer a primitive global theta minor from local quotient nonvanishing.

For both constructed CM hosts, every rank-n lattice stabilizer has 11-valuation at most

\[
 M_{11}(n)=\sum_{j\ge0}\left\lfloor\frac{n}{10\,11^j}\right\rfloor.
 \tag{SI1}
\]

The bound is independent of the field degree. It replaces K039's reduction in the much larger underlying rational dimension. At the actual middle rank1802 it equals197. The actual sharp G levels further reduce the faithful11-Sylow dimension, giving middle bounds191 and178 and terminal bounds384 and357. The initial strict stabilizers are already trivial by K042. These are actual uniform improvements of K039/K041/R024, for every admitted tuple and every coefficient word.

Separately, for the admitted trivial tuple there is an actual middle class with stabilizer S\_t, t=1759 or1633. Its literal Weyl/dual-Weyl invariant coefficient pairing, divided by its actual class order, has 11-adic elementary factors with both endpoint valuations

\[
           (-174,0)\quad\hbox{or}\quad(-162,0).
 \tag{SI2}
\]

All its elementary valuations lie between those endpoints. Thus this is an exact non-scalar integral effect in the actual middle class contraction. The assertion concerns the invariant coefficient lattices, not their as-yet uncomputed intersections with a preceding theta image.

\subsection{1. A single explicit split prime}\label{a-single-explicit-split-prime}

Use the actual fields of K031/K032. They lie in Q(zeta\_52780), and the first lies already in Q(zeta\_812). This is the carried explicit cyclotomic construction, not the incompatible old Fields.Kq package. Set

\[
 q_* =18\,895\,241=1+358\cdot52\,780,
 \qquad q_*\equiv2\pmod {121}.
 \tag{SI3}
\]

This integer is prime. Here is an elementary certificate with no probable-prime assumption. Its predecessor factors completely as

\[
 q_*-1=2^3\cdot5\cdot7\cdot13\cdot29\cdot179.
 \tag{SI4}
\]

All listed factors are prime, as trial division up to their square roots verifies. The base3 satisfies 3\^{}(q\_\emph{-1)=1 modulo q\_}, and for r=2,5,7,13,29,179 the residues of 3\^{}((q\_*-1)/r) are respectively

\[
 18\,895\,240,\ 6\,545\,031,\ 5\,879\,497,\
 6\,338\,700,\ 17\,742\,864,\ 359\,875.
 \tag{SI5}
\]

Each residue minus one is coprime to q\_\emph{. Consequently the order of3 modulo every prime divisor l of q\_} is divisible by every prime-power factor of q\_\emph{-1: otherwise the corresponding quotient exponent would annihilate3 modulo l, contrary to that gcd. Its order is therefore q\_}-1. Since it divides l-1, l\textgreater=q\_*, proving primality. The companion certificate checks every residue and gcd using exact integers.

The congruence in(SI3) makes q\_* split completely in both actual fields. At any chosen prime above q\_* the residue field is F\_(q\_*). It is unrelated to an auxiliary-prime choice in the theta family: we do not change the G, T, Q,29 or11 local data by using it to bound an automorphism group.

Also ord\_11(q\_\emph{)=10: its residue2 has order10 in F\_11\^{}times. Moreover q\_}\^{}10-1 is55 modulo121, so

\[
                   v_{11}(q_*^{10}-1)=1.
 \tag{SI6}
\]

\subsection{2. Uniform bound for the actual stabilizers}\label{uniform-bound-for-the-actual-stabilizers}

Let L be any rank-n projective O\_E-lattice in any of the actual definite Hermitian spaces, and let Gamma be a finite subgroup of its lattice automorphisms. A level or decoration only takes a subgroup, so the same statement applies to every actual initial, middle or terminal class stabilizer. Localize L at one prime above q\_*; it is free of rank n over this DVR. Reduction gives a homomorphism

\[
                \Gamma\longrightarrow GL_n(F_{q_*}).
 \tag{SI7}
\]

The kernel has no element of order11. Indeed if A=1+varpi\^{}s B with s\textgreater=1 and B nonzero modulo varpi had order11, division of A\^{}11-1 by varpi\^{}s would give11B modulo varpi, which is nonzero because the residue characteristic is q\_\emph{≠11. A finite nontrivial11-group has an element of order11. Thus reduction is injective on a Sylow11 subgroup of Gamma, even if the whole reduction map is not injective. It follows that v\_11\textbar Gamma\textbar\textless=v\_11\textbar GL\_n(F\_(q\_}))\textbar.

The latter order is q\_\emph{\^{}(n(n-1)/2) times product\_(i=1)\^{}n(q\_}\^{}i-1). By(SI6), its factors are divisible by11 exactly for i=10j, and then

\[
                  v_{11}(q_*^{10j}-1)=1+v_{11}(j).
 \tag{SI8}
\]

To verify the last formula, write q\_*\^{}10=1+11u with11 not dividing u. A prime-to11 power preserves the valuation of its difference from1; an11th power raises it by one, directly by the binomial expansion. Factoring j into its11-power and prime-to11 part proves(SI8). Summing gives exactly(SI1). No rational-dimension enlargement occurs.

The numerical comparison is:

\begin{longtable}[]{@{}lrrr@{}}
\toprule\noalign{}
Actual host & rank n & old K039 S\_f(n) & new M\_11(n) \\
\midrule\noalign{}
\endhead
\bottomrule\noalign{}
\endlastfoot
J14 & initial42 & 1,479 & 4 (actual0 by K042) \\
J14 & middle1802 & 63,572 & 197 \\
J14 & terminal3561 & 125,628 & 390 \\
J56 & initial168 & 23,705 & 17 (actual0 by K042) \\
J56 & middle1802 & 254,294 & 197 \\
J56 & terminal3435 & 484,744 & 376 \\
\end{longtable}

The proof uses the actual E-rank n, the complete splitting of q\_* and only the11-Sylow subgroup. It remains valid when the coefficient field is further enlarged or ramified over11; valuations are still normalized by v(11)=1.

\subsubsection{Sharp-level refinement}\label{sharp-level-refinement}

For the actual middle and terminal groups there is a stronger bound. Set r=a+1. The G reductions have upper-left block \nolinkurl{diag(U_a,1)} and preserve its r-dimensional coordinate subspace: the middle shape is K038(CA13), and the terminal shape is CA14's(r,t,t) flag. Let P be a Sylow11 subgroup of an actual class stabilizer. At any one actual G place, its homomorphism to that upper-left unipotent block is trivial, since a group of11-power order has no nontrivial image in a group of g-power order. Thus the reduction of its standard representation has at least r fixed dimensions.

This lower bound lifts to characteristic zero. In a local integral class frame the Reynolds matrix

\[
                 e_P=|P|^{-1}\sum_{\gamma\in P}\gamma
 \tag{SI8a}
\]

is integral at this G place, because g is different from11. It is an idempotent, whose image and kernel are free direct summands over the local DVR. Its rank is consequently unchanged by reduction. The reduced image is exactly the invariant space in the residue representation. Extension of scalars also preserves the dimension of the simultaneous kernel of the matrices gamma-1. Therefore the E-rational representation has at least r fixed dimensions.

Its E-rational complement \nolinkurl{ker(e_P)} is P-stable and faithful: P already acts trivially on the complementary invariant space. The complement has rank at most n-r, and it admits a P-stable full O\_E-lattice by summing the P-translates of any initial lattice. Apply(SI1) to that lattice and use monotonicity of M\_11. This gives

\[
 \boxed{
 v_{11}|\Gamma_1|\le M_{11}(t),\qquad
 v_{11}|\Gamma_T|\le M_{11}(2t).}
 \tag{SI8b}
\]

The exact actual values are

\[
 \begin{array}{c|ccc}
       &M_{11}(t)&M_{11}(2t)&\text{sum}\\ \hline
 J14&191&384&575\\
 J56&178&357&535.
 \end{array}
 \tag{SI8c}
\]

This uses the specified sharp middle and terminal flags. It is not an assertion about an arbitrary maximal-compact rank-n class. The integral average(SI8a) is taken at G to prove a characteristic-zero fixed-dimension statement; no integrality of that average at11 is claimed.

\subsection{3. Consequence for the unchanged arithmetic coupled map}\label{consequence-for-the-unchanged-arithmetic-coupled-map}

Keep the notation of K041(NW7--NW10), K042(NB6--NB7) and K038(CA27). Let L\_0,L\_1 be their proved whole-module derivative lower bounds for the actual two selected partitions, with the changed words, role assignments and canonical gap, if any, retained. Then

\[
 P_{\rm cpl}(\mathcal L_0)
       \subset 11^{L_0+L_1-M_{11}(t)}\mathcal L_{\rm ev}.
 \tag{SI9}
\]

The same containment holds on the actual Mram sublattice. The proof is the literal class sum: K042 removes every initial denominator, (SI8b) bounds every middle denominator by191 or178, each kernel has the carried integral bound, and the Y sum is unnormalized. No new projection, row lattice or finite Schwartz input is introduced.

If an actual terminal branch has a coefficient-lattice lower bound B, its terminal scalar against the corresponding integral dual lattice consequently has lower bound

\[
 v\bigl(PtilFF(P_{\rm cpl}(P),Z)\bigr)
       \ge L_0+L_1+B-
 \begin{cases}575,&J14,\\535,&J56.\end{cases}
 \tag{SI10}
\]

Here B is precisely the previously established branch bound for the chosen detector lattice, not a new independent local slope. For R024's actual branch, its two valid choices are B=v(s\_B) for the Weyl detector and B=v(s\_B)-v(H) for the full dual-Weyl detector, using R24.7--R24.9. For the scalar-normalized branch, subtract v(s\_B) from both choices. The common scalar s\_B still cancels against its real cost as in R24.18--R24.21; (SI10) does not reverse that cancellation.

Compared with the already sharpened K042 bounds, the coupled gain increases by63,381 or254,116, respectively. Including the terminal class denominator improves the lower bound by188,625 or738,503. These are fixed-rank arithmetic improvements. They do not prove a sharp upper bound, a nonzero restricted terminal value, or a net DP44 saving, and do not replace the image-content problem by a scalar.

\subsection{4. An actual middle class with symmetric-group stabilizer}\label{an-actual-middle-class-with-symmetric-group-stabilizer}

For this part take the admitted trivial tuple mu0; it is literally \nolinkurl{Tuples.lean120}, and has empty Qprime by its definition at128. All actual G/T data, ranks, selected weights and maximal29 level are unchanged. This is a member of the constructed K038 family.

Put b=a+1+t, r=a+1, and take the standard positive Hermitian lattice O\_E\^{}b. R017(R17.7), whose proof applies in every rank, gives its full rational integral automorphism group

\[
                       \mu_{28}^{b}\rtimes S_b.
 \tag{SI11}
\]

At every G place choose the standard sharp flag. The actual middle residue group is LIFT22(CG4)/K038(CA13): its upper-left r-block is \nolinkurl{diag(u_a,1)} and its remaining diagonal block is GL\_t. The intersection of(SI11) with that group is exactly the monomial group on the last t coordinates, acting as the identity on the first r coordinates. Indeed a monomial upper-unipotent matrix is the identity; once those r rows and columns are occupied, no upper-right monomial entry remains. Reduction on mu28 is injective at these prime-to28 places.

At each T place let u\_T be the fixed nonzero residue covector defining the middle strict level. Choose an integral invertible frame g\_T such that

\[
              u_T\overline g_T^{-1}
                   =(0,\ldots,0,1,\ldots,1),
 \tag{SI12}
\]

with r zeros and t ones. Such a frame exists because GL\_b acts transitively on nonzero residue covectors, and lifts to GL\_b(O). The split unitary identification turns it into an actual maximal-compact unitary element. This changes a decorated class representative, not the fixed level or its covector. It leaves the underlying lattice standard. Use these frames at T and identity outside G union T. At G the original source \nolinkurl{lamY} lattice is transported to the standard local lattice by the fixed split diagonal isometry cancelling its \nolinkurl{entB} coordinate powers. Its two split components are that diagonal matrix and its inverse transpose, so this is an actual unitary isometry. The matrix need not be integral relative to the ambient standard coordinates: its purpose is to send the source's weighted base lattice to the standard lattice. After this transport, the normalized K038(CA13) flags are precisely the displayed standard flags. Every localization of the resulting adelic lattice is now standard, so its global intersection is exactly O\_E\^{}b. This realizes the asserted global lattice from the source base genus by an allowed finite class representative; it changes neither the prescribed level nor the coefficient normalization.

For a rational monomial automorphism remaining after the G conditions, (SI12)'s strict equality forces every nonzero monomial entry to be1. All permutations of the last t coordinates remain. Thus the exact stabilizer of this actual middle class is

\[
                              \Gamma_1=S_t.
 \tag{SI13}
\]

Its finite character is trivial. The unipotent part at G is the identity, and the actual chi0 is an even power of the residue Teichmuller character, so chi0(-1)=1. The strict T character is trivial and there are no Q factors. The integral T frames add no finite character here. In particular this is a genuine nonzero coefficient-invariant class, not a group with a deliberately incompatible fibre character.

Consequently the class denominator really has

\[
 v_{11}(|\Gamma_1|)=v_{11}(t!)=
 \begin{cases}159+14+1=174,&t=1759,\\148+13+1=162,&t=1633.
 \end{cases}
 \tag{SI14}
\]

This proves that the new uniform191/178 bounds cannot be replaced by zero for the actual sharp middle class space. It does not assert that either bound is the optimal uniform maximum.

\subsection{5. Primitive fixed and regular orbits in the actual Weyl lattice}\label{primitive-fixed-and-regular-orbits-in-the-actual-weyl-lattice}

Let R be a complete coefficient DVR above11, enlarged to contain the actual CM labels, with v(11)=1. At this class its11-component is standard. The actual middle Weyl lattice is

\[
 L=\bigotimes_{\tau=1}^{f}
      \det^{s_\tau}\Delta_{\eta_\tau}(R),\qquad
 \eta_\tau=(\nu_{0,\tau},0^d),\quad
 s_\tau=a/2-e_{0,\tau}.
 \tag{SI15}
\]

This is R011(R11.1--R11.3) and the exact source lamTrue\_new recurrence, not an unrestricted polynomial lattice. In R022/K040 and K044 every entry of nu\_(0,tau) is positive and the a entries are pairwise distinct. The interior differences follow from the actual gaps, at least3ell up to the displayed half-integer shifts; the two endpoint inequalities are stricter. Positivity is their checked last-entry inequality. Also each e0\_tau is a multiple of28. Since f=a is even,

\[
            \sum_\tau s_\tau=a^2/2-\sum_\tau e_{0,\tau}
                              \quad\hbox{is even}.
 \tag{SI16}
\]

The combined determinant twist is therefore trivial on S\_t.

Let p\_fix be the tensor of the primitive highest polynomials, using the first a rows at every label. S\_t fixes it. Its weight space is rank one and primitive in L. More generally every Weyl-group translate of this polynomial is a primitive generator of its extremal weight space. These are the actual integral Weyl lattices of K037; a permutation matrix and its inverse preserve them, so no factorial or hook denominator is introduced by such a translate.

There is also a primitive regular S\_t orbit in L. The elementary inequality needed for it is

\[
                  fa\ge t-1:
              \quad1764\ge1758,\qquad28224\ge1632.
 \tag{SI17}
\]

Assign each of the first t-1 spare coordinates to a different pair (tau,i), 1\textless=tau\textless=f,1\textless=i\textless=a. At label tau, move the positive extremal weight nu\_(0,tau),i to its assigned spare coordinate, and keep any unassigned positive weights on their original first-a coordinates. There is a permutation matrix doing this, using disjoint swaps of those assigned original and spare coordinates. Apply it to the primitive highest vector at that label and tensor the results; call the result p\_reg.

The t spare coordinates now have pairwise distinct joint weight profiles: each of the first t-1 has a positive entry at its assigned label, with distinct positive values if labels coincide, while the last has all entries zero. No nonidentity element of S\_t preserves this joint weight. Hence the vectors

\[
                     p_\gamma=\gamma p_{\rm reg}
                                      \quad(\gamma\in S_t)
 \tag{SI18}
\]

have distinct extremal joint weights. They are linearly independent, primitive, and form a copy of the regular permutation R{[}S\_t{]} lattice. Their span is a saturated direct summand of L: integral weight spaces give a direct sum decomposition, and these extremal weight spaces each have rank one. Its complement is S\_t-stable. The fixed weight line Rp\_fix is disjoint from this orbit and is another such stable direct summand.

Let L\^{}vee be the actual opposite-action dual-Weyl lattice, identified with Hom\_R(L,R) by K037(CR17--CR19)'s perfect pairBK. Let p\_fix\^{}vee and p\_gamma\^{}vee be the corresponding primitive dual weight vectors. Set

\[
              u=\sum_{\gamma\in S_t}p_\gamma,
              \qquad u^\vee=\sum_{\gamma\in S_t}p_\gamma^\vee.
 \tag{SI19}
\]

Both invariant lattices have the indicated rank-two orthogonal direct summands, and the pairing on them is exactly

\[
 \left.\operatorname{pairBK}\right|_{
 (Rp_{\rm fix}\oplus Ru)\times(Rp_{\rm fix}^\vee\oplus Ru^\vee)}
                     =\begin{pmatrix}1&0\\0&t!\end{pmatrix}.
 \tag{SI20}
\]

This is an actual invariant coefficient-lattice calculation. It is not a small-rank surrogate or an inference from dimensions.

\subsection{6. Exact range of the class-pairing elementary factors}\label{exact-range-of-the-class-pairing-elementary-factors}

For any finite group Gamma acting on a finite free R-lattice L, the ordinary perfect duality restricts to an injective full-rank map

\[
       L^\Gamma\longrightarrow
                    \operatorname{Hom}_R((L^\vee)^\Gamma,R).
 \tag{SI21}
\]

Its cokernel is killed by \textbar Gamma\textbar. To see this, \nolinkurl{(L^vee)^Gamma} is saturated in L\^{}vee, being the intersection of the kernels of gamma-1. Over a DVR it is a direct summand. Extend any functional on it to L\^{}vee, obtaining x in L. The norm sum N\_Gamma x is invariant, and its pairing with an invariant dual vector is \textbar Gamma\textbar{} times that functional. This proves the asserted annihilation. Over the fraction field the invariant pairing is nondegenerate by averaging, proving full rank. Thus all ordinary elementary-factor valuations lie in {[}0,v(\textbar Gamma\textbar){]}.

For the actual class(SI13), the two direct summands(SI20) attain both endpoints. The invariant coefficient contraction used at this class in the literal K038(CA27) divides this pairing by \textbar Gamma\textbar; this is not an assertion that either fixed kernel realizes every coefficient direction. Its elementary valuations therefore lie in

\[
                 [-v_{11}(t!),0],
 \tag{SI22}
\]

and both endpoints occur. The content of the full fractional class pairing is exactly -v\_11(t!). After removal of any common scalar the width174 or162 remains unchanged. In particular no scalar makes this whole invariant pairing both integral and perfect.

The metric contribution is also explicit. In the normalized compact coefficient metric of R011, primitive extremal vectors have norm one; distinct weights are orthogonal. Thus

\[
 \|p_{\rm fix}\|^2=1,\qquad \|u\|^2=t!,
 \quad\hbox{and identically for the dual vectors}.
 \tag{SI23}
\]

The middle class metric multiplies these squared norms by its actual mass v1/t!, giving v1/t! and v1. Its normalized scalar pairing has matrix \nolinkurl{diag(1/t!,1)} before that common v1. This retains the metric which accompanies the integral anisotropy. The apparent integer t! on the regular invariant is not an uncharged local saving.

\subsection{7. What is and is not decided}\label{what-is-and-is-not-decided}

Equations(SI9--SI10) are improved uniform bounds on the actual coupled arithmetic map and terminal pairing, using their unchanged lattices. Equations(SI13--SI23) compute a sharp non-scalar class-pairing effect inside that same family for its trivial admitted tuple. They disprove the proposed shortcut that strict initial neatness, local primitive quotient tests, or a common period normalization make every middle invariant contraction unimodular at11.

They do not show that both displayed invariant directions occur in the preceding theta image, or that the next raw kernel detects them independently. The fixed kernels can impose additional congruences or cancellation across classes. Determining that restricted image, its primitive minors and its least integral inverse remains an actual arithmetic calculation. Complex injectivity in LIFT26/LIFT27 does not supply it. No universal no-go for a shorter lift chain, complete DP44 comparison, global endpoint nonzero claim, or Lean build follows.

\subsection{8. Source locators and verification scope}\label{source-locators-and-verification-scope-1}

Carried actual inputs: K031/K032 cyclotomic fields and integral hosts; K033 integral CM coordinates at11; R011(R11.1--R11.3) actual middle weights and normalized coefficient metrics; R017(R17.6--R17.9) standard lattice automorphisms; K037(CR17--CR19) actual perfect Weyl/dual-Weyl pairing; K038(CA13--CA15,CA27,CA33) actual sharp group, character, class contraction and mass; K039(IG7--IG13) whole-module integrality; K041(NW7--NW10), K042(NB6--NB7), K044(EG4--EG10) changed weights and strict neatness; R024(R24.7--R24.21) terminal lattice and scalar balance.

Direct primary code locators, at the recovered pinned commit: \nolinkurl{Formalism/LocalToGlobal/Divisibility/Tuples.lean}104--120,128 for the actual trivial tuple and ramification set; \nolinkurl{Formalism/LocalToGlobal/Columns/FiniteSumNonzeroLargeDepth.lean} 738--805 for sharp flags and characters, and298--323 for class-weighted composition; \nolinkurl{Columns/BaseLemmas.lean}496--552 for the exact lamTrue recurrence. These are source-interface checks, not a Lean execution.

The new proofs are(SI3--SI10)'s split-prime Sylow estimate and its actual quantitative application, and(SI12--SI23)'s constructed middle class and exact invariant lattice effect. The companion certificate is \nolinkurl{sharp_stabilizer_integral_pairing_certificate.py/.json}. It passed at 2026-11-22T04:17:03.956396+00:00 after adding the sharp-level refinement. It checked the complete primality certificate and an independent trial division, all3,561 exact local order factors, both actual rank/valuation tables, both explicit joint-weight constructions,131,712 full mu28 phase tests,34,406 small regular-action entries and their exact invariant Gram/Smith matrices. There were10,968 successful assertions. These are exact diagnostics; the general group/lattice proofs above were separately read in full by the author. No global theta image was numerically computed.

\clearpage

\section{AM-K047 --- a short rational unipotent at an actual translated packet point}\label{am-k047-a-short-rational-unipotent-at-an-actual-translated-packet-point}

Author: GPT-6 Astra Pro, ALP\_MACHINE /root/kummer\_extension, 22 November 2026 UTC. Inherited reasoning effort is not separately exposed.

The radius-five separation condition in K046 fails at an actual point of R027's prescribed translated packet, on both current hosts. The counterexample is a completely explicit rational special-unitary unipotent. Its membership is checked with both components of the CM split unitary group, through the global rational unitary matrix of LIFT30. No selected-component matrix is mistaken for a global unitary element.

This refutes the separation hypothesis of that sufficient Poincare criterion. It does not assert that a Poincare sum vanishes, that a particular projected period vanishes, or that every possible analytic method fails. The standard unshifted result of K046 is unchanged.

\subsection{1. The actual field, level and prime}\label{the-actual-field-level-and-prime}

Use either K031 or K032, with E/F the actual CM extension, a=42 or168, n=a+1, and J=diag(I\_a,-1). Both fields contain Q(zeta\_28), hence contain i and delta=sqrt(-7), with complex conjugation sending each of these two imaginary generators to its negative.

Retain R025's standard lattice O\_E\^{}n and its principal compact J1: it is the identity modulo both rational G primes at every split component over them, and the full standard lattice compact elsewhere. Write

\begin{verbatim}
Q = g0*g1,
Gamma1 = U(J)(F) intersect J1.
\end{verbatim}

The exact full congruence is R025(R25.1--R25.3). The actual prime pairs are (4615409,5276377) and (22285932761,64072175441).

Use the unchanged DO040 prime, with its completely specified integers D,T and norm generator

\begin{verbatim}
pi = 1+2D+D*delta,       T=pi*bar(pi)=11D^2+4D+1.
\end{verbatim}

In particular pi and bar(pi) are coprime, T splits completely in E, and pi is1 modulo Q\^{}2. At the selected T components pi is a local uniformizer and c=bar(pi) is a unit. These are DO040.1 and DO040.6--7, retained in LIFT30(NL2). We use those certified primes, not a new choice of T.

Put A=(T+1)/2 and B=(T-1)/2. On the first positive coordinate and the negative coordinate use the rational matrix

\begin{verbatim}
g = [[A/pi, B/bar(pi)],
     [B/pi, A/bar(pi)]],
\end{verbatim}

and the identity on the other a-1 coordinates. LIFT30(NL3--NL6) proves g is in SU(J)(F). Directly, its columns have Hermitian norms 1,-1 and are orthogonal, since A\^{}2-B\^{}2=T. Its determinant is1.

\subsection{2. The finite equality holds at principal G level}\label{the-finite-equality-holds-at-principal-g-level}

Let tau\_c be the finite adele which is identity away from T and has selected T matrix

\begin{verbatim}
tau_c = [[I_a,0], [pi^(-1)*c*e1_star,1]].
\end{verbatim}

At the conjugate component use the actual unitary partner determined by J and inverse transpose. LIFT30(NL7--NL9) gives

\begin{verbatim}
g_f^(-1) = tau_c * k_f.                       (SP1)
\end{verbatim}

For completeness, in the displayed two-plane its selected T factor is

\begin{verbatim}
k = [[A/bar(pi), -B/bar(pi)],
     [-bar(pi),   bar(pi)]].
\end{verbatim}

Its determinant is A-B=1; all its entries and its inverse are integral at that selected component. Its unitary partner is likewise integral. Away from T the two norm generators are units, so g and its inverse are integral. At G, pi is1 modulo Q\^{}2, T is1 modulo Q\^{}2, A is1 and B is0 modulo Q\^{}2, where2 is a unit. Thus g and its inverse are identity modulo every G prime. Consequently

\begin{verbatim}
                     k_f is in J1.             (SP2)
\end{verbatim}

This strengthens the stated J\_B membership of LIFT30(NL9) to the principal compact actually used in R027 and K046. It uses the already retained norm-generator congruence; no level is refined. From(SP1--2),

\begin{verbatim}
tau_c J1 tau_c^(-1) = g_f^(-1) J1 g_f,
Gamma_tau_c = g^(-1) Gamma1 g.                 (SP3)
\end{verbatim}

The second equality is between rational subgroups, because g itself is a rational unitary point. The equality checks both E halves automatically.

\subsection{3. The explicit unipotent and its conjugate}\label{the-explicit-unipotent-and-its-conjugate}

Let v=e1+e\_n.~Then v\_star J v=0. Define

\begin{verbatim}
t=i*Q,
u0=I+t*v*v_star*J.
\end{verbatim}

Here star means the CM-conjugate transpose. The rank-one matrix N=v\emph{v\_star}J has N\^{}2=0. Since t+bar(t)=0,

\begin{verbatim}
u0_star J u0=J,     det(u0)=1,
u0^(-1)=I-t*N.                                 (SP4)
\end{verbatim}

These identities follow by multiplying the two rank-one factors; the quadratic term contains v\_star J v and vanishes. All entries are integral, and u0-I belongs to Q M\_n(O\_E). Hence

\begin{verbatim}
                     u0 is in Gamma1.          (SP5)
\end{verbatim}

It is not the identity, because t is nonzero and N has a nonzero entry. Put gamma=g\^{}(-1)u0g. The rational vector w=g\^{}(-1)v is

\begin{verbatim}
w = bar(pi)^(-1)*e1 + pi^(-1)*e_n.
\end{verbatim}

Indeed the two entries use A-B=1. The unitary identity for g yields

\begin{verbatim}
gamma = I+t*w*w_star*J,
gamma is in Gamma_tau_c,
gamma is not I.                               (SP6)
\end{verbatim}

Its selected two-plane is explicitly

\begin{verbatim}
gamma-I = t * [[1/T,          -1/bar(pi)^2],
               [1/pi^2,       -1/T       ]].    (SP7)
\end{verbatim}

Thus the proposed small denominators are justified by an actual global conjugation. A bare matrix with entries Q/T at one split component would not suffice to establish(SP6).

\subsection{4. Archimedean displacement at every CM label}\label{archimedean-displacement-at-every-cm-label}

At every complex embedding tau of E, \textbar tau(pi)\textbar\^{}2=T and \textbar tau(i)\textbar=1. Consequently \textbar\textbar tau(w)\textbar\textbar\^{}2=2/T. For a rank-one matrix xy\_star its Euclidean operator norm is \textbar\textbar x\textbar\textbar*\textbar\textbar y\textbar\textbar. As J is Euclidean unitary,

\begin{verbatim}
||tau(gamma)-I||_op = 2Q/T.                    (SP8)
\end{verbatim}

In particular \textbar\textbar tau(gamma)\textbar\textbar\_op \textless=1+2Q/T. In U(a,1), a Cartan decomposition k1 a(u) k2 has singular values e\^{}u,1,\ldots,1,e\^{}(-u), with u\textgreater=0. Therefore its Cartan parameter is exactly the logarithm of its operator norm. The actual parameter u\_tau of gamma satisfies

\begin{verbatim}
0 < u_tau <= log(1+2Q/T) < 2Q/T.              (SP9)
\end{verbatim}

It is strictly positive: a nonidentity unipotent cannot lie in the compact group of Euclidean unitary matrices, whose elements are diagonalizable over C. The strict upper bound log(1+x)\textless x follows by integrating 1/(1+x)\textless1 on a positive interval.

The exact prime sizes give the following convenient, deliberately nonoptimal rational bounds:

\begin{longtable}[]{@{}llrr@{}}
\toprule\noalign{}
Host & bit lengths of g0,g1 & bit length of T & 2Q/T \textless{} \\
\midrule\noalign{}
\endhead
\bottomrule\noalign{}
\endlastfoot
J14 & 23,23 & 1001 & 2\^{}(-953) \\
J56 & 35,36 & 996 & 2\^{}(-923) \\
\end{longtable}

For example Q\textless2\^{}46 and T\textgreater=2\^{}1000 in the first case, so 2Q/T\textless2\^{}(-953). The other case is identical. The companion exact certificate checks these bounds with the actual integers.

Let C5 be K046's product of K exp({[}0,5)) K over all labels. Then

\begin{verbatim}
gamma is in C5,
gamma is in Gamma_tau_c and gamma is not I.   (SP10)
\end{verbatim}

Since the identity belongs to C5, it follows in particular that C5 C5\^{}(-1) intersect Gamma\_tau\_c is not \{I\}.

\subsection{5. This is a point of the actual prescribed packet}\label{this-is-a-point-of-the-actual-prescribed-packet}

The actual primitive T covector is u=(1,2,\ldots,2\^{}(a-1)), as in DO040.9. At each selected T place choose an integral invertible a\_v in GL\_a(O\_v) satisfying

\begin{verbatim}
                     c*e1_star*a_v=u.         (SP11)
\end{verbatim}

One explicit choice takes first row u/c and all later rows equal to the corresponding standard rows. Its determinant is1/c, a unit. Supply the conjugate split component by its unitary partner. Take a\_f in H(A\_f)=U(I\_a)(A\_f) to equal these matrices at T and identity elsewhere. Its embedded matrix hat(a\_f)=diag(a\_f,1) lies in J1, because it is integral at T and identity at G.

The exact multiplication of LIFT30(NL10--NL11) gives

\begin{verbatim}
tau_u = hat(a_f)^(-1) tau_c hat(a_f),
hat(a_f) tau_u J1 = tau_c J1.                 (SP12)
\end{verbatim}

If a different integral lift of the same residue covector had been chosen, the quotient of the two lower-unipotent matrices would be integral at T, so the right J1 coset would be unchanged.

R027's actual packet includes the point represented by (o,x tau\_u) for every x in H(A\_f). Choose x=a\_f.~Then(SP12) identifies that point with (o,tau\_c) at the same principal level. Its arithmetic group is Gamma\_tau\_c, which contains the nontrivial short gamma of(SP6--10). This is not a replacement of the prescribed tau\_u by a different detector; it is one actual finite class of its packet.

The same statement holds after grouping the initial classes at their literal strict T covector level. No projective stabilizer, dilation average or equal-summand mass substitution was used. It is also not necessary to assert that a\_f is rational: it is explicitly a finite unitary adele, and x is allowed to range over precisely those classes.

\subsection{6. What is and is not excluded}\label{what-is-and-is-not-excluded}

The standard unshifted Gamma1 satisfies K046(BL27), whereas this particular translated packet group fails that radius-five separation condition by an enormous explicit margin. Therefore the uniform K046 concentration plus its one-point Poincare separation argument cannot simply be applied to the whole R027 packet.

The group Gamma\_tau\_c is still discrete and neat; a short unipotent has infinite order and causes no contradiction. Failure of the sufficient nonvanishing criterion says nothing by itself about cancellation in the Poincare series, the projected global period, or the smallest singular value of the actual packet evaluation map. No cusp-dimension conclusion follows from this example. Those remain separate analytic or cohomological questions.

The proof uses only the exact retained arithmetic data and elementary rank-one matrix identities. The referenced general analytic criterion is Zunar(2025), Proposition8.1, used with all conventions in K046; this report asserts that one of its hypotheses fails at this point.

\subsection{7. Exact certificate and source locators}\label{exact-certificate-and-source-locators}

The standalone standard-library certificate \nolinkurl{translated_packet_short_unipotent_certificate.py} works exactly in Q(i,sqrt(-7)), using four rational coordinates and explicit CM conjugation. For each actual prime it checks the norm generator, g\_star J g=J, determinant1, the inverse and finite two-plane factor, N\^{}2=0, u0\_star J u0=J, the rational conjugation formula(SP7), the norm of w, the G congruences and the exact displacement bounds. It also checks the triangular local covector transport for both actual ranks. This is a finite exact matrix/form check, not an enumeration of the arithmetic lattice or a global cusp computation.

Direct internal source inputs are R025(R25.1--R25.3), LIFT30(NL1--NL11), LIFT29(BT4--BT8), and DO040.1,6--9. Their original source level predicates and explicit prime certificates are retained; no unsupported common-quadratic Host constructor is claimed. In particular the old Fields.Kq obstruction is not changed.

Primary analytic reference: Zunar, S. (2025). On the non-vanishing of Poincare series on irreducible bounded symmetric domains. Author version arXiv2501.06876, Proposition8.1, printed pp.15--16; the explicit type A III Cartan coordinates are in Section9, pp.18--19. \url{https://arxiv.org/pdf/2501.06876} . The actual number-field membership argument above is elementary and does not apply that paper's Gaussian-integer principal-level example to another field.

\clearpage

\section{AM-DO-044 --- an actual dimension obstruction to the full projected bottom packet}\label{am-do-044-an-actual-dimension-obstruction-to-the-full-projected-bottom-packet}

Author: GPT-6 Astra Pro, ALP\_MACHINE \nolinkurl{/root/definite_oscillator}. 22 November 2026 UTC. Inherited reasoning effort is not separately exposed.

This finding concerns the genuine e0 bottom space and the entire initial strict-T coefficient space of R025/R027, on the actual K031/K032 hosts. At the largest gap in K044, every permitted canonical twist and every binary word have far too few bottom cusp forms to detect every initial row. The same is true after restricting the initial space to the literal Mram norm kernels. The result is a dimension statement about these fixed spaces. It does not say that every initial row is undetected, that the particular theta-image row is undetected, or that its terminal period is zero. It does not contradict R027's assertion for sufficiently large, unrestricted twists.

The global upper bound below is obtained from cohomology, including the actual finite projector and its boundary vanishing. It is not inferred from a local Bergman concentration estimate. The compact-dual Weyl dimension is kept distinct from a Haar-normalized discrete-series formal degree.

\subsection{1. Actual spaces and the conclusion}\label{actual-spaces-and-the-conclusion}

Retain the exact fields, primes, and standard lattices of K031, K032, DO040, R025, and R027. Write E/F for either CM host, f={[}F:Q{]}, and

\[
 (a,f,n)=(42,42,43)\quad\hbox{or}\quad(168,168,169),
 \qquad H=U(I_a),\qquad G=U(I_a,-1),\qquad n=a+1.
                                                        \tag{DO44.1}
\]

The relative extension E/F is unramified at every finite place. Both G primes g0,g1 split completely. The DO040 prime T also splits completely. There are f places above each of these rational primes.

Let J0 be the actual initial strict level: at G its reduction is U\_a with the stipulated inverse generic character; at each T place it is the strict stabilizer of the prescribed nonzero covector; elsewhere it is the standard maximal compact, with trivial character. Let A\_lambda be its coefficient-row space, and let A\_ram be the literal Mram intersection of unweighted T norm kernels. All stabilizers of these definite classes are trivial, by K042's all-place neatness argument.

The bottom level J\_B and character are exactly R025(R25.5): at G the reduction is diag(U\_a,1) with the opposite generic character psi; at every non-G place, including T,Q,7,29, the level is maximal and the character is trivial. Thus no non-G determinant character has been discarded. Put

\[
 \mathcal G=\prod_{j=0,1}\prod_{v\mid g_j}GL_n(\mathbf F_{g_j}),
 \qquad \mathcal U=\prod_{j,v}\operatorname{diag}(U_a,1).
                                                        \tag{DO44.2}
\]

Let B\_(e0,psi)(Lambda) be the actual bottom cuspidal multiplicity space with this finite type, the genuine source goodTuple projector e0, and the chosen holomorphic discrete series (or its contragredient in the literal period convention). It is the space of vector-valued cusp forms, not that space multiplied by the dimension of its lowest K type. The translation at T used by R027 is fixed and changes neither this dimension nor the following upper bound.

Fix A=2\^{}960 and B=2\^{}955. Use K044 at ell=ell\_+, and allow every word, every role assignment, and every canonical twist q=(a+1)s with s in 28 Z\_(\textgreater=0), q\textless=ell. Then

\[
\begin{array}{c|c|c}
 &\dim B_{e0,\psi}/\dim A_\lambda
 &\dim B_{e0,\psi}/\dim A_{\rm ram}\\ \hline
 J14&<2^{-589932}&<2^{-589931}\\
 J56&<2^{-2198952}&<2^{-2198951}.
\end{array}                                               \tag{DO44.3}
\]

In particular the R027 selected-jet restriction map cannot be onto A\_lambda\^{}dual or A\_ram\^{}dual at these weights. The bilinear period map from A\_ram to B\_(e0,psi)\^{}dual has a nonzero kernel of dimension at least dim A\_ram minus dim B\_(e0,psi). This is an actual space obstruction, independent of the norm chosen for the interpolation inverse.

\subsection{2. Exact initial dimensions, including Mram}\label{exact-initial-dimensions-including-mram}

Let m\_r be the standard definite genus mass of (O\_E\^{}r,I\_r), with each class weighted by 1/\textbar Aut\textbar. Let d\_lambda be the product of the dimensions of the actual U(a) coefficient representations. The finite index of the G level in the maximal compact is

\[
 I_{G,a}=\prod_{j=0,1}
   \left({|GL_a(\mathbf F_{g_j})|\over|U_a(\mathbf F_{g_j})|}\right)^f.
                                                        \tag{DO44.4}
\]

The strict covector stabilizer has index T\^{}a-1 in GL\_a(F\_T). Since all strict-level stabilizers are1, mass refinement and coefficient evaluation on the classes give the exact identity

\[
       \dim A_\lambda=m_a I_{G,a}(T^a-1)^f d_\lambda.
                                                        \tag{DO44.5}
\]

Here arbitrary right characters at G cause no loss: once the rational stabilizer is trivial, a coefficient value at a class is unrestricted. This is the literal argument underlying \nolinkurl{SiegelMassRatios/ColsDim.lean}, \nolinkurl{ColsDecomp} and \nolinkurl{ColsDim}, lines63--85.

For Mram, enlarge the T level from the strict covector stabilizer H\_u to its projective stabilizer P\_u, retaining the strict U\_a level at all G places. This enlarged level is still neat by K042: its proof used the G congruences, not the T restrictions. Thus every projective-T class also has trivial stabilizer. The strict class cover is free under

\[
             \mathcal D_T=(\mathbf F_T^\times)^f.
\]

Each coarse class contributes its regular D\_T representation tensored with its coefficient fibre. The source norm at a T place is exactly the unweighted sum over that factor, so its kernel is the augmentation ideal in the corresponding regular representation. Intersecting these kernels gives

\[
 \dim A_{\rm ram}=
       \left({T-2\over T-1}\right)^f\dim A_\lambda
       >\tfrac12\dim A_\lambda.                         \tag{DO44.6}
\]

The last inequality follows from T-1\textgreater2f and (1-1/(T-1))\^{}f\textgreater=1-f/(T-1)\textgreater1/2. There is no central28 divisor and no character-weighted projector in this calculation. Source locators are \nolinkurl{ScaledRamifiedLattice.lean}84--91 (\nolinkurl{NmG}),126--127 (\nolinkurl{NmOp}),160--169 (\nolinkurl{ColRamV},\nolinkurl{MramV}), and \nolinkurl{Columns/Record.lean}1811--1848 (diamond and norm transport). Inversion of the diamond action leaves the regular and augmentation dimensions unchanged.

\subsection{3. A free cover and an equivariant Euler bound}\label{a-free-cover-and-an-equivariant-euler-bound}

\subsubsection{3.1. Auxiliary level3 is used only to enlarge the bottom space}\label{auxiliary-level3-is-used-only-to-enlarge-the-bottom-space}

Put principal rational level3 at every place above3, and leave the other non-G bottom levels unchanged. Let X3 have maximal level at G, and let Y3 have principal level at G. Then

\[
                 Y3\longrightarrow X3
\]

is a free finite cover with full deck group Gcal. This auxiliary level does not change the claimed detector family: its cusp space contains the original one, and is used only for an upper bound.

For completeness, principal3 is neat. Eigenvalues of an element congruent to1 modulo3, in the restriction-of-scalars lattice, lie in 1+3 O\_(Q3bar). So do all finite products of such eigenvalues and their inverses. A nontrivial root of unity cannot lie in this group: a root of order prime to3 has nontrivial reduction, while a primitive root of 3-power order has v3(zeta-1)\textless=1/2. Thus the multiplicative group of eigenvalues has no torsion. This argument applies to every rational conjugate of the level. In particular the center mu28 does not create fixed points for the deck cover.

Let I3 be the exact principal3 index in the standard compact of rank n. The fields are unramified at3. Smooth reduction and the formulas for GL\_n or U\_n over a residue field give

\[
                       I3\le 2^f3^{fn^2}.              \tag{DO44.7}
\]

Indeed \textbar U\_n(F\_Q)\textbar=Q\^{}(n\^{}2) product\_(i=1)\^{}n(1-(-1)\^{}i Q\^{}(-i)), and the product is at most product\_(i\textgreater=1)(1+Q\^{}(-i)) \textless=exp(1/(Q-1))\textless2 for Q\textgreater=3. The split formula is smaller than Q\^{}(n\^{}2). There are at most f residue places, and their residue degrees sum to f.

\subsubsection{3.2. The actual good projector annihilates the boundary}\label{the-actual-good-projector-annihilates-the-boundary}

The F-Witt index of (I\_a,-1) is1: it has an isotropic vector and at each real place its negative index is1. Consequently every proper rational parabolic is an isotropic-line parabolic. At a selected split G place its residue Levi is GL1 x GL\_(n-2) x GL1.

For each genuine series block admitted by e0, literal IsGood clause5 with degree1 supplies a selected G place whose polynomial has no linear factor. The exact source is \nolinkurl{CyclotomicNorms/StatementT.lean}19--29. At that place the finite representation has zero Harish--Chandra restriction to this Levi. Here is the finite reason. Under Green's characteristic map, ordinary parabolic induction is multiplication in the symmetric-function algebra indexed by Frobenius orbits. A GL1 factor contributes a degree-one orbit. Multiplication preserves the degree in each orbit variable; hence an irreducible with no degree-one orbit cannot occur in such an induction. Frobenius reciprocity proves the asserted zero restriction. This uses the genuine Green series, as in R026, not merely a formal idempotent. Primary locator: Gow--Vinroot(2007), §3, the definition of the product before Theorem3.2 and Theorem3.2 itself.

We explain why the algebraic local system does not invalidate this finite argument. A rational boundary face is a nilmanifold fibration over its Levi arithmetic quotient. With algebraic coefficients F\_Lambda, its fibre cohomology is H*(n,F\_Lambda), naturally for parabolic conjugation and congruence-level changes. The unipotent radical acts trivially on this cohomology: on the Chevalley--Eilenberg complex its infinitesimal action satisfies L\_X=d i\_X+i\_X d; the connected unipotent group therefore acts trivially on cohomology. The Nomizu--van Est identification with the nilmanifold cohomology retains these algebraic coefficients. Exact primary descriptions are Goresky--Harder--MacPherson--Nair(2002), §5.1 equations(5.1.3--4), Lemma5.2, and Remark5.3, pp254--255; they refer explicitly to Goresky--Harder--MacPherson(1994), §§11.4--11.5 and12.6.

At the finite principal G level, assemble the boundary faces over a fixed rational cusp class. Iwasawa decomposition organizes the deck translates by the corresponding finite parabolic. An individual cusp stabilizer may be a smaller subgroup S of that parabolic; it contains the full finite unipotent radical, by unipotent strong approximation. Inducing first from S to the parabolic and then to Gcal gives the same description, with the remaining data in the Levi module. The Leray spectral sequence for the nilmanifold fibration has Levi cohomology with the preceding fibre coefficients. From its E2 page onward it is a finite parabolic induction whose finite unipotent radical is trivial. One can see the latter directly: unipotent strong approximation realizes a finite residue unipotent deck translation by a rational unipotent translation, whose action on the fibre cohomology is the just-proved trivial action. The remaining arithmetic and level data sit in the Levi module. Nonstandard cusps only conjugate the finite parabolic. The boundary-stratification spectral sequence has these modules as its graded pieces; rank1 has no additional proper-parabolic intersections.

Thus each admitted tensor block kills the boundary at its supplied place. Finite group representations over C are semisimple, so vanishing on the graded pieces is vanishing on the boundary cohomology itself:

\[
                 e0\,H^*(\partial Y3,F_\Lambda)=0.      \tag{DO44.8}
\]

This statement holds for the full source sum of good blocks, regardless of which place supplies degree-one exclusion in each block. It is not the assertion that a finite level cover alone has trivial unipotent boundary action.

\subsubsection{3.3. Actual regular highest weight and the middle degree}\label{actual-regular-highest-weight-and-the-middle-degree}

At one CM label the signed bottom HC parameter is exactly LIFT28(HJ4)

\[
 \Lambda=(v_0,\ldots,v_{a-2},-b_0,-b_1),\quad
 v_i=4\ell(a-1-i),\quad b_0=E-\ell,\quad b_1=E+\ell+q.
                                                        \tag{DO44.9}
\]

The consecutive differences are 4ell, E+3ell, and 2ell+q. All exceed1, including thin labels where b0 can be negative. Therefore the algebraic highest weight mu=Lambda-rho\_n has every simple positive-coroot pairing at least1. The dual weight has the same property. This is regularity of the highest weight itself; regular infinitesimal character alone would not be enough.

Lan(2016), Theorem4.10(1--3), equation(4.11), and Remark4.15 apply: all noncompact factors are of type A, so this regularity implies ordinary de Rham vanishing below the complex dimension D=af, compact support vanishing above D, and interior cohomology only in degree D. These are statements for the actual product domain; no Kunneth decomposition of the arithmetic quotient is assumed. Together with the boundary exact sequence and (DO44.8), they imply

\[
           e0 H^i(Y3,F_\Lambda)=0\quad(i\ne D).          \tag{DO44.10}
\]

The holomorphic-discrete-series cuspidal multiplicity space injects into the cohomology with the matching coefficient F\_Lambda (or its dual). Each such discrete series has a nonzero middle-degree relative Lie algebra cohomology space, of dimension1 in this holomorphic case. The cuspidal cohomology injection is the direct-summand statement of Borel--Labesse--Schwermer(1996), §5.5, p15; its representation decomposition is §6.5. Applying it componentwise to the derived groups, with the fixed compact central character retained, gives the assertion for U. We only need the injection, not equality with all middle cohomology.

\subsubsection{3.4. Equivariant Euler characteristic}\label{equivariant-euler-characteristic}

The arithmetic quotient X3 has a finite CW homotopy model, for example the Borel--Serre compactification. Lift a CW model through the free cover. For any descended local system of rank dim F\_Lambda, its cellular cochains are sums of dim F\_Lambda copies of the regular Gcal module, one sum for every cell of the base. Consequently, in the representation ring,

\[
 \sum_i(-1)^i[H^i(Y3,F_\Lambda)]
       =\chi(X3)\dim(F_\Lambda)[\mathbf C[\mathcal G]].   \tag{DO44.11}
\]

This proof does not assume that deck transport is trivial on a fibre: choosing one lift of a cell and transporting its fibre basis gives the regular module because the cell stabilizer is trivial.

Put e\_psi=\textbar Ucal\textbar\^{}(-1) sum\_(u in Ucal) psi(u)\^{}(-1)u. It commutes with the central e0, and its rank on the regular module is \textbar Gcal\textbar/\textbar Ucal\textbar. Equations(DO44.10--11) and the cusp injection give

\[
 \dim B_{e0,\psi}(\Lambda)
 \le |\chi(X3)|\,\dim F_\Lambda\,
                      { |\mathcal G|\over|\mathcal U|}.
                                                        \tag{DO44.12}
\]

\subsection{4. Exact Euler-to-definite-mass normalization}\label{exact-euler-to-definite-mass-normalization}

We now prove the equality needed for (DO44.12), rather than identifying an indefinite covolume with a definite mass by convention:

\[
                     |\chi(X3)|=n^f m_n I3.             \tag{DO44.13}
\]

At every finite place, -1 is the norm of a unit of E\_v: this is immediate at split places and follows from norm surjectivity on units in an unramified quadratic extension otherwise. Multiplying the last basis vector by such a unit identifies the integral forms I\_n and (I\_a,-1). Thus their standard compact groups and their principal3 subgroups are integrally isomorphic at every finite place.

Here is a full reduction to the semisimple Euler formula, including all determinant components. Let S=SU(I\_a,-1), S\_c=SU(I\_n), and T\^{}1=ker(N\_(E/F):Res\_(E/F) G\_m -\textgreater{} G\_m). Let K3 and K\_c3 be the maximal finite levels with principal3 imposed. The same algebraic section

\[
                 s(t)=\operatorname{diag}(t,1,\ldots,1)
\]

splits det:U-\textgreater T\^{}1 for both forms, and maps the standard compact of T\^{}1 into both standard compacts, and principal3 into principal3. Consequently both determinant images are exactly

\[
 K_{T,3}=\{t\in T^1(\widehat O_F):t\equiv1\pmod3\}.
\]

Crucially,

\[
             T^1(F)\cap K_{T,3}=\{1\}.                 \tag{DO44.14a}
\]

An element of this intersection is an integral global norm-one unit. All its complex conjugates have absolute value1, so it is a root of unity. Its congruence1 modulo3 excludes every nontrivial root of unity by the proof of neatness in Section3.1. This also proves that the determinant of every arithmetic component group at K3 or K\_c3 is1: conjugation of the finite level does not change its determinant. Thus discrete central arithmetic determinants have been removed by the actual auxiliary level, not silently divided out.

Let C\_T=T\^{}1(F)\T\^{}1(A\_f)/K\_(T,3), a finite torus class set. For each representative t in T\^{}1(A\_f), set

\[
 K_{S,t}=S(A_f)\cap s(t)K3s(t)^{-1},\qquad
 K_{S_c,t}=S_c(A_f)\cap s(t)K_{c3}s(t)^{-1}.
                                                        \tag{DO44.14b}
\]

The determinant fiber over {[}t{]} of X3 is \textbf{exactly} S(F)(D x S(A\_f))/K\_(S,t); the corresponding compact-form fiber is S\_c(F)\S\emph{c(A\_f)/K}(S\_c,t). To verify this, use the rational and compact sections s to move the determinant of a representative to t, then write its finite part as x s(t) with x in S(A\_f). Two such representatives are equivalent only if their rational multiplier gamma and compact multiplier k satisfy det(gamma)det(k)=1. This determinant lies in T\^{}1(F) cap K\_(T,3), hence equals1 by(DO44.14a). Therefore both multipliers are in the respective special unitary groups. The remaining compact conjugation is precisely(DO44.14b). Conversely those SU equivalences are plainly U equivalences. At infinity the section lies in the maximal compact, so D=S\_infty/(S\_infty cap K\_infty). This proves the fiber identifications without any residual torus quotient.

Choose the trace bilinear densities on the two special unitary Lie algebras and their associated compatible Tamagawa measures. The local inner-form isometries preserve the trace pairing. Their finite absolute densities, discriminant normalization, and convergence factors therefore agree. Denote the common finite compact volume by v\_S, and the volume of the compact dual S\_c,infty by V\_S. Conjugation by s(t) preserves these Haar measures: it preserves the trace pairing, and its action on the semisimple Lie algebra has determinant1. Thus every K\_(S,t) and K\_(S\_c,t) has the same finite volume v\_S, independent of t.

Both simply connected special unitary groups have Tamagawa number1. This is the Weil--Langlands--Lai--Kottwitz theorem (type A here); an exact accessible statement is the Automorphic Project, version ce6f683, Theorem13.8.17,p176, with the original Kottwitz(1988) reference and the Hasse-principle qualification. No exceptional-type issue occurs here. Every compact principal3 stabilizer is trivial. Hence every compact fiber has exactly

\[
                    {1\over V_S v_S}
\]

classes. Summing over C\_T and using ordinary mass refinement gives

\[
             m_n I3={|C_T|\over V_S v_S}.               \tag{DO44.14}
\]

For the noncompact SU fiber, Harder's arithmetic Gauss--Bonnet theorem and the compact-dual Euler density give

\[
 \chi\bigl(S(F)\backslash(D\times S(A_f))/K_{S,t}\bigr)
   ={(-1)^{af}n^f\over V_S v_S}.                       \tag{DO44.15}
\]

Indeed the Weyl ratio is \textbar W(SL\_n)\^{}f\textbar/\textbar W(S(GL\_a x GL\_1))\^{}f\textbar=n\^{}f.~The exact density formula is Kionke(2014), Theorem3.1,p383, and its adelic integration is Theorem3.3, pp386--387, derived there from Harder(1971) and Serre. These theorems are applied only to the semisimple groups S and S\_c.~The quotient is neat; Harder's formula holds also when it is noncompact. Summing (DO44.15) over the \textbf{same} finite set C\_T and comparing with (DO44.14) proves(DO44.13). There is no unexplained central volume, unitary Tamagawa multiplicativity, or determinant-class correction.

For clarity, dim F\_Lambda is the compact-dual Weyl dimension

\[
 \dim F_\Lambda=\prod_\tau
 {\prod_{i<j}(\Lambda_{\tau,i}-\Lambda_{\tau,j})
       \over\prod_{i<j}(j-i)}.                         \tag{DO44.16}
\]

This is the literal \nolinkurl{fdegHC} numerical expression in \nolinkurl{FormalDegreeBessel/StatementDS.lean}19--22 and \nolinkurl{FdegWeyl}60--64, but here it is used as a Weyl dimension. For the explicit ball Haar of K046, the actual formal degree differs from this quantity by a! at each label. Neither that Haar factor nor a lowest-K-type dimension is inserted in (DO44.12--13).

\subsection{5. The exact ratio and uniform archimedean bounds}\label{the-exact-ratio-and-uniform-archimedean-bounds}

Combine (DO44.5),(DO44.12),(DO44.13). The identity \textbar GL\_(a+1)(F\_g)\textbar/\textbar GL\_a(F\_g)\textbar=g\^{}a(g\^{}(a+1)-1) gives

\[
 {\dim B_{e0,\psi}\over\dim A_\lambda}
 \le n^f I3\,{m_n\over m_a}
  \prod_\tau{\dim F_{\Lambda_\tau}\over\dim W_{\lambda_\tau}}
  \left[
   {g_0^a(g_0^n-1)g_1^a(g_1^n-1)\over T^a-1}
  \right]^f.                                          \tag{DO44.17}
\]

The prime T appears only on the initial side because the bottom has maximal T level. This is the actual class-count difference that a selected packet interpolation theorem would have to overcome.

Let u be the initial HC list, with the exact K044 coordinates

\[
 u_i=v_{i-1}-\tfrac12-\ell\omega_{i-1}\ (1\le i\le a-2),
\]

top endpoints A+(a-1)/2, -B-(a-1)/2, and non-top endpoints (8a-4)ell+1/2, 2ell+1/2. For a decreasing list x write V(x) for its Vandermonde. Put p=(v\_0,\ldots,v\_(a-2),-b0), and

\[
 R_a=\prod_{r=1}^{a-3}\left({4r\over4r-1}\right)^{a-2-r}.
\]

The exact one-label ratio is

\[
 W_\tau:={\dim F_{\Lambda_\tau}\over\dim W_{\lambda_\tau}}
  ={V(p)\over V(u)}{\prod_{i=0}^{a-1}(p_i+b_1)\over a!}.
                                                        \tag{DO44.18}
\]

For interior u pairs with j-i=r, their difference is at least (4r-1)ell. Hence cancellation of the interior Vandermondes gives

\[
 {V(v_0,\ldots,v_{a-2})\over V(u_1,\ldots,u_{a-2})}
       \le R_a(4\ell)^{a-2}(a-2)!.                     \tag{DO44.19}
\]

At a top label and ell=ell\_+, every v\_i+b\_j is \textless2B, u0-u\_i\textgreater A/2, u\_i-u\_last\textgreater B, and u0-u\_last\textgreater A. Also b1-b0\textless=3ell. Inserting these factors into(DO44.18), with A=32B, yields

\[
 W_{\rm top}\le
        U_{\rm top}:=R_a{3\over8a(a-1)}B\ell^{a-1}.
                                                        \tag{DO44.20}
\]

To check the constant, the endpoint denominator is at least A(A B/2)\^{}(a-2), the added numerator is at most (2B)\^{}(2a-2)3ell, and (4ell)\^{}(a-2)(a-2)!/a! supplies the remaining factors.

At a non-top label, u0-u\_last=(8a-6)ell, u0-u\_i\textgreater=4aell, and u\_i-u\_last\textgreater=4(a-i-1)ell. Each v\_i+b\_j is at most E+4aell. Therefore

\[
 W_{\rm non-top}\le
 U_E:=R_a{3(E+4a\ell)^{2a-2}
      \over a!(8a-6)(4a)^{a-2}\ell^{a-2}}.              \tag{DO44.21}
\]

These bounds are uniform over every binary word and q\textless=ell. Since E\_thin\textless E\_thick, set U=max(U\_top,U\_(E\_thick)); it bounds every role. The proof is by the displayed factorwise inequalities, not by sampling words or estimating logarithms asymptotically.

\subsection{6. Definite mass ratio and exact arithmetic}\label{definite-mass-ratio-and-exact-arithmetic}

Use the stipulated and source-exact Siegel mass ratio, whose formula was recovered in \nolinkurl{HermitianLattices/StatementMass.lean}27--51,78--89. For n=a+1 and rd=\textbar disc(E)\textbar\^{}(1/{[}E:Q{]}), it says

\[
 {m_n\over m_a}=
 \left[{a!\,\operatorname{rd}^{a+1/2}
              \over(2\pi)^{a+1}}\right]^f L_F(n,\chi).
                                                        \tag{DO44.22}
\]

Here L\_F is exactly the positive Euler product encoded by \nolinkurl{lValEta}: the paired split E primes supply (1-Nv\^{}(-n))\^{}(-1), and an inert prime supplies (1-(-1)\^{}n Nv\^{}(-n))\^{}(-1). Thus L\_F(n,chi)\textless=zeta\_F(n)\textless=zeta(n)\^{}f\textless2\^{}f.

The actual root discriminants satisfy rd\textless2\^{}8 on J14 and rd\textless2\^{}11 on J56. These inequalities can be checked without real approximation by raising to the42nd power:

\[
 2^{42}7^{35}29^{36}<2^{336},\qquad
 2^{42}7^{35}29^{36}65^{21}<2^{462}.
\]

Since 2pi\textgreater6, putting r=8 or11 gives the rational upper bound

\[
 {m_n\over m_a}<\mu_a^f,\qquad
 \mu_a={a!\,2^{ra+\lceil r/2\rceil+1}\over6^{a+1}}.
                                                        \tag{DO44.23}
\]

Using(DO44.7), the right side of(DO44.17) is strictly less than R\^{}f, where the following is an exact positive rational number:

\[
 R=n(2\,3^{n^2})\mu_a U
       {g_0^a(g_0^n-1)g_1^a(g_1^n-1)\over T^a-1}.
                                                        \tag{DO44.24}
\]

The companion Python certificate reads the full actual primes and weight integers from the already checked DO040/K044 certificates. It performs integer and Fraction arithmetic only. Its certified ceiling base-two logarithms are

\begin{longtable}[]{@{}
  >{\raggedright\arraybackslash}p{(\columnwidth - 14\tabcolsep) * \real{0.0968}}
  >{\raggedleft\arraybackslash}p{(\columnwidth - 14\tabcolsep) * \real{0.1290}}
  >{\raggedleft\arraybackslash}p{(\columnwidth - 14\tabcolsep) * \real{0.1290}}
  >{\raggedleft\arraybackslash}p{(\columnwidth - 14\tabcolsep) * \real{0.1290}}
  >{\raggedleft\arraybackslash}p{(\columnwidth - 14\tabcolsep) * \real{0.1290}}
  >{\raggedleft\arraybackslash}p{(\columnwidth - 14\tabcolsep) * \real{0.1290}}
  >{\raggedleft\arraybackslash}p{(\columnwidth - 14\tabcolsep) * \real{0.1290}}
  >{\raggedleft\arraybackslash}p{(\columnwidth - 14\tabcolsep) * \real{0.1290}}@{}}
\toprule\noalign{}
\begin{minipage}[b]{\linewidth}\raggedright
host
\end{minipage} & \begin{minipage}[b]{\linewidth}\raggedleft
R\_a
\end{minipage} & \begin{minipage}[b]{\linewidth}\raggedleft
U\_top
\end{minipage} & \begin{minipage}[b]{\linewidth}\raggedleft
U\_thick
\end{minipage} & \begin{minipage}[b]{\linewidth}\raggedleft
mu\_a
\end{minipage} & \begin{minipage}[b]{\linewidth}\raggedleft
finite G/T ratio
\end{minipage} & \begin{minipage}[b]{\linewidth}\raggedleft
auxiliary3
\end{minipage} & \begin{minipage}[b]{\linewidth}\raggedleft
R
\end{minipage} \\
\midrule\noalign{}
\endhead
\bottomrule\noalign{}
\endlastfoot
J14 & 51 & 20868 & 20783 & 400 & -38251 & 2932 & -14046 \\
J56 & 295 & 82187 & 82763 & 2423 & -143551 & 45270 & -13089 \\
\end{longtable}

Each last entry is obtained by exact comparison of the entire rational R with the power of2, not by adding rounded logarithms. Raising to f=42 or168 gives(DO44.3), and(DO44.6) gives its Mram column. The certificate additionally checks1512 exact small-rank/word instances of(DO44.18--21), all actual endpoint inequalities, discriminant bounds, and the Mram fraction. These diagnostics supplement the proof; they do not establish the cohomological statements.

There is also a broad explicit interval consequence. If ell is any allowed28-multiple with ell\_+/2\^{}64\textless=ell\textless=ell\_+, top bounds only decrease, and(DO44.21) is at most2\^{}(64(a-2)) times its endpoint bound (use the endpoint numerator and the lower bound for ell in its denominator). Consequently the full-row ratios are respectively \textless2\^{}-482412 and \textless2\^{}-414120 throughout this upper part of K044's interval. The Mram ratios remain less than twice these numbers. No claim is made here for the remainder of that interval.

\subsection{7. What this settles, and what remains}\label{what-this-settles-and-what-remains}

For the specified weights, a projected bottom interpolation map onto every initial row is impossible. This conclusion retains the literal strict T index, both G-prime groups, genuine e0, generic character, coefficient rank, and every finite class. It also retains a deliberately generous auxiliary principal3 enlargement on the bottom side.

The dimension argument produces undetected nonzero Mram rows, but it does not locate a particular such row, identify the kernel of the fixed coupled theta map followed by tilQ/brC, or show that the existing selected witness belongs to that kernel. The earlier complete-target theta injectivity statements remain compatible with it: their target is a different space, and the endpoint restriction is an additional map. A single chosen bottom detector, or a detector on a smaller subspace, is not ruled out. The source Host/Scale typing obstructions are also independent of this genuine mixed-family dimension statement.

\subsection{References and exact primary anchors}\label{references-and-exact-primary-anchors}

Borel, A., Labesse, J.-P., \& Schwermer, J. (1996). On the cuspidal cohomology of S-arithmetic subgroups of reductive groups over number fields. \emph{Compositio Mathematica, 102}(1),1--40. §5.5,p15; §6.5,pp18--20. \url{https://www.numdam.org/item/CM_1996__102_1_1_0.pdf}

Gan, W. T., Hanke, J. P., \& Yu, J.-K. (2001). On an exact mass formula of Shimura. \emph{Duke Mathematical Journal, 107}(1),101--131. §2; §4, equations(4.2--3), Propositions4.4--4.5. \url{https://web.math.princeton.edu/~jonhanke/Web-02/Mass-Formula-of-Shimura/shimura-mass.pdf}

Goresky, M., Harder, G., \& MacPherson, R. (1994). Weighted cohomology. \emph{Inventiones Mathematicae,116},139--213. §§11.4--11.5,12.6. \url{https://doi.org/10.1007/BF01231560}

Goresky, M., Harder, G., MacPherson, R., \& Nair, A. (2002). Local intersection cohomology of Baily--Borel compactifications. \emph{Compositio Mathematica,134},243--268. §5.1,Lemma5.2,Remark5.3,pp254--255. \url{https://www.math.ias.edu/~goresky/pdf/localIH.jour.pdf}

Gow, R., \& Vinroot, C. R. (2007). \emph{Extending real-valued characters of finite general linear and unitary groups on elements related to regular unipotents} {[}Preprint{]}. §3,Theorem3.2. \url{https://arxiv.org/abs/0704.2022}

Harder, G. (1971). A Gauss-Bonnet formula for discrete arithmetically defined groups. \emph{Annales Scientifiques de l'Ecole Normale Superieure, 4}(3),409--455. §§1.3,2.1. \url{https://numdam.org/item/10.24033/asens.1217.pdf}

Kionke, S. (2014). Lefschetz numbers of symplectic involutions. \emph{Pacific Journal of Mathematics,271}(2),369--414. Theorems3.1,3.3, pp383,386--387. \url{https://msp.org/pjm/2014/271-2/pjm-v271-n2-p.pdf}

Kottwitz, R. E. (1988). Tamagawa numbers. \emph{Annals of Mathematics, 127}(3),629--646. Weil's theorem, with its type-A Hasse principle; accessible statement in the Automorphic Project(2020),Theorem13.8.17,p176. \url{https://doi.org/10.2307/2007007} \url{https://math.jhu.edu/~sakellar/automorphic-files/book.pdf}

Lan, K.-W. (2016). Vanishing theorems for coherent automorphic cohomology. \emph{Research in the Mathematical Sciences,3},Article39. Theorem4.10,equation(4.11),Remark4.15,pp20--21. The one-page author errata changes only type-D and exceptional-root material; type A and Theorem4.10 are unchanged. \url{https://www.kwlan.org/articles/van-coh-aut-err.pdf} \url{https://doi.org/10.1186/s40687-016-0085-6}

The locally retrieved primary PDFs and extracted text are retained under \nolinkurl{definite_oscillator_dimension_refs/}; the GH-MN and BLS passages were also read directly through their public PDF pages. Source formulas are read from the recovered pinned archive; this finding does not claim a Lean build.

\subsection{Verification record}\label{verification-record-2}

Author whole-proof check completed22 November2026,04:20:40 UTC, by GPT-6 Astra Pro, ALP\_MACHINE \nolinkurl{/root/definite_oscillator}; inherited effort not separately exposed. The complete proof and certificate were read. The check included the literal goodTuple degree1 clause and finite Green support, the algebraic-coefficient boundary radical action, Lan's actual highest-weight hypothesis, the free-cover Euler representation, the compact-central-torus/inner-form measure cancellation, all exact Weyl and mass factors, the uniform word bounds, and the full-row versus particular-row scope. PASS in the stated cohort. The final exact certificate replay at04:19:17 UTC passed1512 diagnostics and all actual integer inequalities. An independent source check of the exact Mram fraction was supplied by \nolinkurl{/root/lift_accounting} at04:14:36 UTC.

The strengthened Section4 was rechecked in full at04:27 UTC. It now identifies every determinant fiber explicitly, proves that principal3 removes all residual rational torus determinants, and applies the semisimple Euler formula only after that reduction. The accessible numbered Tamagawa theorem is retained with its exact review scope; no inaccessible original theorem locator has been invented. The current author PDF of Lan and its entire one-page errata were also read; the two errata concern only typeD and an exceptional root, and do not change the typeA vanishing used here. No arithmetic certificate or stated numerical bound changed in this refinement.

\clearpage

\section{AM-R-030: the actual ball Haar volume in terms of the definite mass}\label{am-r-030-the-actual-ball-haar-volume-in-terms-of-the-definite-mass}

Author: GPT-6 Astra Pro, ALP\_MACHINE \nolinkurl{/root}, 22 November 2026 UTC. Reasoning effort is not separately exposed.

This is an exact normalization corollary of DO044 and K046. It evaluates the total adelic volume that bounds the standard-component volume in K046. It does not extend that theorem to the translated packet, remove the genuine projector, or give a global integral inverse.

\subsection{1. Statement and fixed measures}\label{statement-and-fixed-measures}

Use either actual host E/F of K031/K032, with a=f=42 or168, n=a+1, G=U(I\_a,-1), and m\_n the definite genus mass of (O\_E\^{}n,I\_n). Let K\_max be the standard maximal finite compact. At infinity give each maximal compact volume1 and use precisely the K046 quotient Haar

\[
 d\mu_{\mathrm{ball}}(z)
   =\pi^{-a}(1-\|z\|^2)^{-a-1}\,dV(z).
                                                        \tag{R30.1}
\]

For a finite compact K contained in K\_max, define the total volume by summing the real quotient volumes over all finite adelic components. If K is not neat, retain the actual finite stabilizer weights; in particular an ineffective central stabilizer is not discarded. Then

\[
 \boxed{\quad V(K)=\frac{m_n}{(a!)^f}[K_{\max}:K].\quad}  \tag{R30.2}
\]

For the actual principal G level J1 and actual bottom level J\_B of R025, respectively, this becomes

\[
 V(J1)=\frac{m_n}{(a!)^f}
             \prod_{j=0,1}|GL_n(\mathbf F_{g_j})|^f,
                                                        \tag{R30.3}
\] \[
 V(J_B)=\frac{m_n}{(a!)^f}
       \prod_{j=0,1}
       \left(\frac{|GL_n(\mathbf F_{g_j})|}
                         {|U_a(\mathbf F_{g_j})|}\right)^f.
                                                        \tag{R30.4}
\]

These are volumes of the underlying levels. Prescribed characters and central spectral projectors change spaces of sections; they do not change these quotient measures. The bottom level is maximal at T, so no T index is inserted in (R30.4).

\subsection{2. Compact-dual density and its exact volume}\label{compact-dual-density-and-its-exact-volume}

The compact dual of the ball is CP\^{}a. In its standard affine chart, the compact-dual density matching (R30.1) at the origin is

\[
 d\mu_c(z)=\pi^{-a}(1+\|z\|^2)^{-a-1}\,dV(z).           \tag{R30.5}
\]

Here the normalization can be checked directly. The Hermitian matrix of the potential log(1+\textbar\textbar z\textbar\textbar\^{}2) has determinant (1+\textbar\textbar z\textbar\textbar\^{}2)\^{}(-a-1), by the rank-one determinant identity. The potential -log(1-\textbar\textbar z\textbar\textbar\^{}2) similarly has determinant (1-\textbar\textbar z\textbar\textbar\^{}2)\^{}(-a-1). Both matrices are the identity at the origin. Multiplying the resulting volume density by pi\^{}(-a) therefore gives the same tangent density on the noncompact and compact duals. This is the absolute trace-density correspondence used in the Euler measure formula of Kionke (2014), Theorem3.1, with the common maximal compact given the same volume1.

The complementary hyperplane has measure zero. Polar coordinates and the beta integral give

\[
 \begin{split}
 \operatorname{vol}(\mathbf {CP}^a)
 &=\frac1{(a-1)!}\int_0^\infty
                     t^{a-1}(1+t)^{-a-1}\,dt\\
 &=\frac1{(a-1)!}\int_0^1 u^{a-1}\,du
  =\frac1{a!},
 \end{split}                                           \tag{R30.6}
\]

where u=t/(1+t). The compact-dual Weyl ratio is n=a+1. Consequently, for the product of f factors, the Euler--Poincare measure relative to the stated ball Haar is

\[
              d\mu_{\mathrm{EP}}
                 =(-1)^{af}(n!)^f d\mu_{\mathrm{ball}}.
                                                        \tag{R30.7}
\]

This is the compact-dual formula applied to SU on each determinant fibre. The determinant fibres and their common compact central measure are exactly those checked in DO044.14a--15. No application of a semisimple theorem directly to reductive U is being assumed.

\subsection{3. Removing the auxiliary level and restoring the actual levels}\label{removing-the-auxiliary-level-and-restoring-the-actual-levels}

Let K3 be the auxiliary principal3 subgroup used only for DO044's normalization proof, and I3={[}K\_max:K3{]}. DO044.13 proves

\[
                 |\chi(X3)|=n^f m_n I3.
                                                        \tag{R30.8}
\]

The quotient is neat. Harder's Gauss--Bonnet formula, in the precise form Kionke (2014), Theorems3.1 and3.3, therefore combines with (R30.7--8) to give

\[
             V(K3)=\frac{m_n I3}{(a!)^f}.               \tag{R30.9}
\]

For Haar quotient volume, passage to a finite-index compact subgroup multiplies the sum of component volumes by that index. This remains true when the coarser space has finite stabilizers: the group quotient measure, equivalently the stabilizer-weighted class decomposition, counts precisely those multiplicities. Dividing (R30.9) by I3 gives V(K\_max)=m\_n/(a!)\^{}f.~Multiplying by {[}K\_max:K{]} proves (R30.2).

At each of the f places above g\_j the principal residue quotient is GL\_n(F\_(g\_j)). The actual J\_B quotient replaces the identity residue subgroup by diag(U\_a,1). These are R025's exact levels, so their indices give (R30.3--4). This proves all asserted volume identities.

\subsection{4. The standard selected inverse with no unevaluated component volume}\label{the-standard-selected-inverse-with-no-unevaluated-component-volume}

Write Gcal for the principal G residue deck group, and let delta=f*2\^{}(-ell\_max/8)\textless2\^{}-100, as in K046. Its selected-jet right inverse R0 is defined at the standard principal component. That component's volume is at most the sum over all components. Thus K046's actual normalized Hilbert estimate implies

\[
 \boxed{\quad
  \|R_0\|\le
  \frac{\sqrt{m_n|\mathcal G|/(a!)^f}}{1-2\delta}.
 \quad}                                                \tag{R30.10}
\]

The PBW/Fischer scalar rho is retained exactly as in K046: (R30.10) bounds its normalized selected evaluation, not an unnormalized derivative. Formula (R30.10) is expressed in the actual definite mass, which is determined by the stipulated mass theorem and the actual field; it does not claim a newly computed decimal mass or class number. No transitivity of the finite adelic components was required.

The translated-point obstruction K047 and full-packet dimension obstruction DO044 are unchanged. This corollary supplies no lower bound for the prescribed projected terminal period, no integral global interpolation, and no inverse bound for the vertical theta map.

\subsection{5. Sources and geometric interpretation}\label{sources-and-geometric-interpretation}

K046, BL1--35, fixes (R30.1), the selected normalization and its standard-point inverse. DO044.13--15 supplies the exact determinant-fibre Euler/mass comparison. R025.1--5 supplies both actual finite levels. These are cited checked findings, not proved again here.

Kionke, S. (2014). Lefschetz numbers of symplectic involutions. \emph{Pacific Journal of Mathematics,271}(2),369--414. Theorems3.1 and3.3, pp383,386--387. \url{https://msp.org/pjm/2014/271-2/pjm-v271-n2-p.pdf}

The new content is the explicit ball-to-compact-dual normalization and its substitution into the already checked measure and inverse formulas. It concerns metrized analytic principal parts and arithmetic class volumes. The coherent/jet comparison retains EGAIII.1 Theorem2.2.1(ii), p100, and EGAIV.4 16.11.2--3(ii),pp54--55, with the Picard/biextension locators in the run dictionary. No algebraic-cycle, motivic projector, ell-adic purity or Hecke RH claim follows.

\clearpage

\section{ALP-LIFT-34 --- the actual coupled image cancels every middle 11-primary class denominator}\label{alp-lift-34-the-actual-coupled-image-cancels-every-middle-11-primary-class-denominator}

Author: GPT-6 Astra Pro, ALP\_MACHINE \nolinkurl{/root/lift_accounting}. 22 November 2026 UTC. Inherited reasoning effort is not separately exposed.

\textbf{Author whole-proof reading passed at04:38:10 UTC.} The separate finite certificate passed at04:37:32 UTC; its scope is recorded below.

This result concerns the unchanged arithmetic two-edge map of K038, on its actual initial Weyl row lattice, with its prescribed sharp G/T/Q functions and its full Y quotient. It improves the actual image bound of LIFT33. The denominator improvement comes from a free action on the finite support of the actual second kernel, followed by the exact middle coefficient contraction. An infinite theta-orbit count alone would not establish the required integral assertion. The proof instead constructs an auxiliary finite Schwartz cutoff, with integral values at11 and the same11-vacuum, that represents the divided contribution.

The cutoff is a proof device. It does not replace the fixed raw family or impose a new finite level on its source. It need not have the original level covariance by itself. No assertion of a primitive theta minor, exact positive Smith factors, or terminal detection is made.

\subsection{1. Actual data and the stronger image bound}\label{actual-data-and-the-stronger-image-bound}

Use either actual host and the literal ranks

\[
\begin{array}{c|rrrrr}
 &a&r=a+1&t&b=r+t&c=b+t+1\\ \hline
 J^{14}&42&43&1759&1802&3562\\
 J^{56}&168&169&1633&1802&3436.
\end{array}
\tag{MC1}
\]

Keep K038's actual sharp coupled map, its strict T levels, its genuine G characters, every admitted common tuple and Q set, and maximal level at29. This includes the repaired coefficient profiles of K041 and K044, with their exact changed selected partitions and source creation scales. Choose class representatives standard at11 as in K039(IG3). All coefficient valuations are normalized by v(11)=1, also after a finite extension of the coefficient DVR R.

For an initial class e, a middle class f and a terminal decoration d, write

\[
 N_{e,f,d}(P)=\sum_{Y\in\mathcal Y_f}
                 k_2(f,Y,d)k_1(e,f,Y)P,
 \qquad \Gamma_f=\Gamma_{1,f}.
\tag{MC2}
\]

The k\_j are the actual geometric ordinary-derivative/subD maps, not abstract local quotient models. Let L\_0,L\_1 be the whole-module bounds of K039(IG8--IG10), or their exact changed-profile versions in K041. Denote the source and output transported Weyl lattices by L\_\{0,e\} and L\_\{2,d\}. Then

\[
\boxed{\quad
 N_{e,f,d}(L_{0,e})
      \subseteq 11^{v_{11}(|\Gamma_f|)+L_0+L_1}L_{2,d}.
 \quad}
\tag{MC3}
\]

Consequently the actual class-weighted contribution satisfies

\[
             |\Gamma_f|^{-1}N_{e,f,d}(L_{0,e})
                       \subseteq11^{L_0+L_1}L_{2,d}.
\tag{MC4}
\]

The initial stabilizers are trivial by K042(NB6). The unchanged full coupled row therefore obeys

\[
\boxed{\quad
 P_{\rm cpl}(\mathcal L_0)
          \subseteq11^{L_0+L_1}\mathcal L_{\rm ev}.
 \quad}
\tag{MC5}
\]

The same containment holds on its actual Mram sublattice. No assumption that an initial local constituent is Steinberg is used here. Compared with LIFT33(SI9), this removes its remaining middle losses191 and178. Compared with K039's post-K042 bound, it removes the entire middle loss63,572 or254,294. These are lower-bound improvements, not positive Smith-factor equalities or net DP44 savings.

\subsection{2. The actual second G support is free under a middle Sylow subgroup}\label{the-actual-second-g-support-is-free-under-a-middle-sylow-subgroup}

Fix e,f,d and a Sylow11 subgroup P\_f of the finite group Gamma\_f. Write h\_f for the middle finite class frame. At any one actual selected G place, put

\[
                 A_\gamma=h_f^{-1}\gamma_fh_f
                    \quad(\gamma\in P_f).
\tag{MC6}
\]

These are integral invertible matrices in the actual sharp compact group. Its exact residue shape is K038(CA13). Reduction is injective on P\_f.~Indeed an element in its kernel has the form 1+varpi\^{}s B, s\textgreater=1, B nonzero modulo varpi. If it had order11, the binomial expansion of its11th power, divided by varpi\^{}s and reduced modulo varpi, would be11B, which is nonzero. A nontrivial finite 11-group has an element of order11. The residue characteristic is the actual G prime, different from11. This argument uses the actual local class lattice and requires no auxiliary split prime.

The complete second matrix has c rows and b columns. In the notation of K038(CA17), its first n\_T=c-1 rows are Z and its last row is eta, including the specified piT factor. Whenever the actual second G function is nonzero,

\[
 \bar Z=\begin{pmatrix}B\\0\end{pmatrix},
 \qquad B\in H^\sharp_{r,t}\subseteq GL_b(\mathbf F_g).
\tag{MC7}
\]

Thus Z itself has full column rank modulo g. This is the literal PhiGS/PhiTop support proved in LIFT21(TG7--TG8) and K038(CA17--CA18), not an inferred generic open-orbit condition. The eta row and its Fourier phase are retained. The right action

\[
                       \bar Z\longmapsto\bar Z\bar A_\gamma
\tag{MC8}
\]

is free on this support: a left inverse of bar Z makes a fixed-point equation imply bar A\_gamma=1, hence gamma=1 by injectivity. The support is stable by K038(CA19). Choose one representative from each P\_f orbit in this finite set of residues, and let c\_f(bar Z) be its0/1 indicator. Then, on the full raw support,

\[
                  \sum_{\gamma\in P_f}
                           c_f(\bar Z\bar A_\gamma)=1.
\tag{MC9}
\]

No division by \textbar P\_f\textbar{} is used in defining this indicator. The choice does not depend on a selected polynomial or derivative order.

\subsection{3. Cut the fixed full-Y tensor, not an individual moving summand}\label{cut-the-fixed-full-y-tensor-not-an-individual-moving-summand}

Let Psi\_raw be the actual finite tensor on the direct sum of the two Hom Schrödinger spaces. At G and Q it is the complete local Y sum; at the other places it is the given product. LIFT26(GS6--GS7c) proves that the literal moving-lattice Y construction is precisely the diagonal-middle Weil action on this one fixed tensor. Its proof keeps the common middle frame in both factors, the exact ftY scalar, and all periodic phases. In particular we may cut this tensor before taking either theta sum.

Let C\_f multiply a finite Schwartz function by c\_f(bar Z) at the one G place just chosen, and by1 elsewhere. Set

\[
                        \Psi'_{\rm raw}=C_f\Psi_{\rm raw}.
\tag{MC10}
\]

This is a genuine Schwartz function: the multiplier depends only on a finite residue, and the original support is compact. It is a finite sum of decomposable two-edge tensors, as every locally constant compactly supported function on a product is. The full Y sum remains inside it. The proof does not choose a surviving Y or divide by the size of that quotient.

The multiplier is defined in the exact split Hom Schrödinger model of LIFT26(GS7a). If the fixed global rational position model uses a different local polarization, transport this function through the same local Weil intertwiner used in K039(IG7). Its linear changes, phases and finite Fourier measures at G have only G-primary denominators and their square roots, hence preserve11-integrality. This local model choice neither acts at11 nor changes the tensor.

The exact middle finite covariance of the two original factors is inverse on the two edges. At G this is K038(CA14--CA21); at Q it is CA9--CA11; at T it is the same strict-covector incidence transport; elsewhere it is the actual basic-vacuum compact covariance. Therefore their full tensor is fixed by the diagonal middle compact action. The middle character cancels before a coefficient is evaluated. These are the actual honest Weil operators, with the splitting factors of LIFT26(GS7a), not merely symplectic actions modulo phase.

More explicitly, at this split G place the diagonal action pulls the position variables back by

\[
               (X,Z,\eta)\longmapsto
                  (A_\gamma^{-1}X,ZA_\gamma,\eta A_\gamma).
\tag{MC11}
\]

Any determinant scalar in that operator commutes with multiplication by c\_f.~In fact det gamma=1 here: its11-power order makes its determinant an11-power root of unity in the actual CM field, whose roots of unity are mu28. This observation is convenient but not needed to cancel the two prescribed finite coefficient characters. The operators at the other places commute with C\_f.~Since the full uncut tensor is invariant, equations(MC9)--(MC11) give the exact identity

\[
            \Psi_{\rm raw}
                =\sum_{\gamma\in P_f}
                   \omega_{\rm mid,f}(A_\gamma)\Psi'_{\rm raw}.
\tag{MC12}
\]

Now apply the actual initial, middle and terminal finite class frames on both Hom spaces. Denote the resulting inputs by Psi\_\{e,f,d\} and Psi'\_\{e,f,d\}. Conjugation of the diagonal middle action gives

\[
             \Psi_{e,f,d}
                  =\sum_{\gamma\in P_f}
                      \omega_{\rm mid,f}(\gamma)\Psi'_{e,f,d}.
\tag{MC13}
\]

This is why the rational stabilizer and its local class matrices may not be silently identified. Their conjugating frame is included in both cut and uncut inputs.

The cut tensor has exactly the same product vacuum at11: the cut is at G, all three class frames are identity at11 by K039(IG3), and no raw G/T/Q datum occurs at11 by IG2. Its values are integral at w. Before transport this follows from the raw roots of unity and the prime-to11 powers in ftY. After transport it follows from the same finite Weil generator calculation and algebraic splitting-unit lemma used in K039(IG4--IG7). The0/1 multiplier introduces no new scalar denominator, character field, or choice of square root.

For later norm accounting there is also an exact finite identity. The summands in(MC12) have disjoint support at the chosen G residue and are unitary translates. Consequently

\[
                     \|\Psi'_{e,f,d}\|_f^2
                         =|P_f|^{-1}\|\Psi_{e,f,d}\|_f^2.
\tag{MC14}
\]

It applies also when the original Y sum has cancellations, because it is an identity for the full function. This finite Hilbert norm is not identified with a global inverse norm or an endpoint height.

\subsection{4. The complete middle coefficient contraction is integral and invariant}\label{the-complete-middle-coefficient-contraction-is-integral-and-invariant}

Let Delta\_j be the actual Weyl lattice of the source rhoVr model and Nabla\_j its pairBK-perfect dual in the opposite action. This is the literal integral duality of K037(CR17--CR19). For P\_0 in Delta\_0 and a target test v\_2 in Nabla\_2 choose an R-basis p\_i of Delta\_1 and its pairBK-dual basis v\_i of Nabla\_1. Define the two-variable polynomial

\[
 Q(P_0,v_2)(Z_0,Z_1)
        =\sum_i\iota_0(P_0,v_i)(Z_0)\,
                     \iota_1(p_i,v_2)(Z_1).
\tag{MC15}
\]

Both iota maps are the canonically normalized subD maps of K037(CR7--CR18). The first sends a Weyl input and a dual-Weyl test to an integral polynomial, as does the second. Thus(MC15) is integral. Equivalently it is the actual integral universal substitution Q\_1 composed with Q\_0:

\[
   \Delta_0\longrightarrow\Delta_1[Z_0]
                \longrightarrow\Delta_2[Z_0,Z_1],
\tag{MC16}
\]

followed by the perfect target test. This basis-free formulation is K037(CR8); it uses no inverse Fischer matrix, inverse hook product, or rational Schur projector. Every monomial of(MC15) has exactly the two label multidegrees D\_\{0,tau\} and D\_\{1,tau\}.

Let q\_infty(P\_0,v\_2) be its actual two-edge creation vector. The middle action on the first iota target is rho\_1 and on the second iota source is its dual. Since

\[
                \sum_i\rho_1(h)v_i\otimes\rho_1^\vee(h)p_i
                          =\sum_i v_i\otimes p_i,
\tag{MC17}
\]

the actual creation vector q\_infty is fixed by the diagonal real middle unitary group. This includes every gamma\_infty for gamma in P\_f.~The assertion concerns the complete coevaluation, not just the fixed source-and-target highest tensor. Its normalization is exactly K037(CR27--CR30) and K038(CA36), so there is no1/dim factor in(MC15).

The source unit-Fock creation scalars Lambda\_0 Lambda\_1 and the periods and cD comparisons are common to every term of(MC15), since both multidegrees are fixed. They stay outside the algebraic ordinary-derivative functional just as in K038(CA33--CA36). In particular this argument does not assign finite valuations to a complex period or conflate the Chern and literal source norms.

\subsection{5. Exact rational theta identity for the divided contribution}\label{exact-rational-theta-identity-for-the-divided-contribution}

Let Theta be the ordinary rational theta functional on the direct sum of the two actual Hom spaces. The two-edge theta product is this theta functional, because the rational summation lattice is a direct product. A finite sum of separable tensors gives the same identity by linearity. For rational gamma in the middle unitary group, its honest global Weil action preserves Theta. This is the actual rational theta automorphy used in K035/K036 and LIFT26, with the global splitting character trivial on rational points. In a Schrödinger model this identity includes Poisson summation when the rational isometry does not preserve the chosen position polarization; no permutation formula for such a position space is assumed here.

The fixed-degree Gaussian series is normally convergent, so the following finite sum and applications of rational theta automorphy are valid. Apply(MC13) and use the real invariance(MC17). Then

\[
\begin{split}
 \Theta(q_\infty\otimes\Psi_{e,f,d})
  &=\sum_{\gamma\in P_f}
      \Theta(q_\infty\otimes
                 \omega_{\rm mid,f}(\gamma)\Psi'_{e,f,d})\\
  &=\sum_{\gamma\in P_f}
      \Theta(\omega_{\rm mid}(\gamma)
                           (q_\infty\otimes\Psi'_{e,f,d}))\\
  &=|P_f|\Theta(q_\infty\otimes\Psi'_{e,f,d}).
\end{split}
\tag{MC18}
\]

Remove the same nonzero common period and source-creation scalar on the two sides, using the actual all-function identity of K035 and its K038 coupled comparison. The left side is precisely pairBK(N\_\{e,f,d\}(P\_0),v\_2). The right side is \textbar P\_f\textbar{} times the ordinary derivative of the geometric germ attached to the cut finite input, tested by the same polynomial(MC15). Hence

\[
 |P_f|^{-1}\operatorname{pairBK}(N_{e,f,d}(P_0),v_2)
                  =Q(P_0,v_2)(cD\partial)F'_{e,f,d}(0).
\tag{MC19}
\]

This is an identity in the algebraic coefficient field after a common finite extension. Its right side has an independent geometric construction from the integral finite cutoff. It has not been defined by division of the left side. The equality thus supplies the missing integral content information that an analytic orbit count alone would not give.

\subsection{6. The cut germ retains the separated derivative gain L\_0+L\_1}\label{the-cut-germ-retains-the-separated-derivative-gain-l_0l_1}

Apply K039(IG5--IG7) to the direct sum of the two actual Schrödinger spaces. The conditions were checked after(MC13): integral finite values at w and exactly the product vacuum at11. Its prime-to11 finite Fourier expansion therefore has the form

\[
 \Psi'_{e,f,d}=
    \sum_{t_0,t_1}b_{t_0,t_1}
           (W_{t_0}\phi_{0,\rm vac})\otimes
           (W_{t_1}\phi_{1,\rm vac}),
                  \qquad b_{t_0,t_1}\in R.
\tag{MC20}
\]

Here both Heisenberg parameters have denominators prime to11, and their prescribed central phases are retained. To see the separated form directly, the Heisenberg group of a symplectic direct sum acts as the tensor product on the two factors; its finite Fourier characters and cosets separate accordingly. The only Fourier denominators are powers of a common prime-to11 integer D. There is no division by \textbar P\_f\textbar{} in this construction, even though \textbar P\_f\textbar{} is an 11-power.

After removal of the fixed common period, each summand in(MC20) is the product of two actual reduced Mumford-translate germs, each belonging to the complete integral local-section module on its respective edge. This is the torsion-connection comparison K034(TG3--TG10), used in K039(IG7). The reduced Chern frame itself need not be integral; the retained formal factor is exactly the one in that result.

For each monomial of(MC15), its derivative of a product in(MC20) is the product of one ordinary derivative in each disjoint variable block. K039(IG9--IG10) bounds these by L\_0 and L\_1, respectively. All polynomial and expansion coefficients are integral. Therefore

\[
             v\bigl(Q(P_0,v_2)(cD\partial)F'_{e,f,d}(0)\bigr)
                                     \ge L_0+L_1.
\tag{MC21}
\]

This uses the separated finite Heisenberg expansion, not an unjustified addition of separate bounds for an arbitrary mixed coefficient. It also does not use the fixed-source-row version of the derivative theorem on other Weyl weights: the invoked bounds are the whole-module label-degree estimates of K039.

Since pairBK is a perfect pairing of the actual target Weyl and dual-Weyl lattices, equations(MC19)--(MC21), for all v\_2 in Nabla\_2, prove(MC3). Write \textbar Gamma\_f\textbar=\textbar P\_f\textbar m\_f, where m\_f is an11-unit; this proves(MC4). The original finite initial and middle class sums, without changing any representative or Y multiplicity, now prove (MC5). Equivariance is inherited from the unchanged original map, not asserted for an individual cut function.

\subsubsection{The actual large G primes also permit cancellation of the full middle mass}\label{the-actual-large-g-primes-also-permit-cancellation-of-the-full-middle-mass}

For these two exact hosts one can replace P\_f everywhere in the cutoff construction by the entire Gamma\_f.~This is a separate strengthening; it does not follow just from injectivity on an11-Sylow. The actual prime and dimension comparison is

\[
\begin{array}{c|r|rr}
 &D_{\mathbb Q}=2fb&g_0&g_1\\ \hline
 J^{14}&151368&4615409&5276377\\
 J^{56}&605472&22285932761&64072175441.
\end{array}
\tag{MC21a}
\]

These are the unchanged, proved prime G data of K031/K032. At either chosen place, g-1\textgreater D\_Q. The finite reduction kernel has no element of prime order l different from g, by exactly the local binomial argument in section2 with l in place of11. If that kernel were nontrivial, Cauchy's theorem would therefore give an element of order g. But Gamma\_f acts faithfully on the underlying D\_Q-dimensional rational space. A nonidentity rational matrix of prime order g has a primitive g-th root as an eigenvalue, so its rational characteristic polynomial is divisible by Phi\_g, of degree g-1. This contradicts(MC21a). Consequently reduction is injective on the full Gamma\_f.

The action(MC8) is now free for the full group. The covariance and real coevaluation arguments used no Sylow assumption. In particular the determinant scalar in(MC11) commutes with the cutoff even when det gamma is nontrivial; the original finite character still cancels between the two edges, and(MC17) includes the honest real vacuum characters. Choose a full Gamma\_f transversal and denote its framed input by Psi\^{}full\_\{e,f,d\} and its geometric germ by F\^{}full\_\{e,f,d\}. Equations(MC19) and(MC14) become

\[
\begin{split}
 |\Gamma_f|^{-1}\operatorname{pairBK}
                  (N_{e,f,d}(P_0),v_2)
     &=Q(P_0,v_2)(cD\partial)F^{\rm full}_{e,f,d}(0),\\
 \|\Psi^{\rm full}_{e,f,d}\|_f^2
     &=|\Gamma_f|^{-1}\|\Psi_{e,f,d}\|_f^2.
\end{split}
\tag{MC21b}
\]

The same prime-to11 Heisenberg expansion proves the same L\_0+L\_1 integral bound. Thus the entire middle class mass has an independent cut-section representation, with its exact finite Hilbert norm. This does not establish integrality at the G residue characteristic, an inverse norm of the global map, or cancellation of the terminal mass.

\subsection{7. What has and has not been cancelled}\label{what-has-and-has-not-been-cancelled}

The improvement is an actual image statement. At the middle class constructed in LIFT33, the ambient invariant coefficient pairing still has Smith endpoints(-174,0) or(-162,0). Equation(MC19) proves that the actual paired theta numerators have the extra11-primary divisibility required to cancel this class denominator. It does not make that ambient pairing unimodular, or show that the image reaches either of its extremal elementary directions.

For the terminal period, applying only LIFT33's remaining terminal bound and the same actual branch/detector bound B gives

\[
 v\bigl(\operatorname{PtilFF}(P_{\rm cpl},G)\bigr)
     \ge L_0+L_1+B-
       \begin{cases}384,&J^{14},\\357,&J^{56}.
       \end{cases}
\tag{MC22}
\]

The branch normalization is exactly R024's: B is v(s\_B) on its Weyl test, or v(s\_B)-v(H\_mu) on the literal dual-Weyl detector as stated there, and normalizing the branch subtracts v(s\_B) on both sides. No terminal denominator cancellation is asserted in(MC22).

The reason the present free-action proof is genuinely a middle argument is that the second raw map is injective on its source columns. A terminal group acts on a larger target, and may fix that image while acting on a complementary subspace. The literal terminal branch, fixed extra line and detector must all be retained in any further argument. They are not supplied by(MC8).

Nothing here bounds an inverse of the fixed global theta matrix, identifies its exact nonzero invariant factors, proves the restricted endpoint nonzero, or supplies all source SizeData2/Host predicates. The finite cutoff norm(MC14) and the unchanged common analytic scales are retained, so the divisibility improvement is not called a net gain/error improvement without the remaining actual costs.

\subsection{8. Source and verification locators}\label{source-and-verification-locators}

The carried results used here were already proved and checked in this run; their statements are applied, not re-proved as new display or theta-correspondence results:

\begin{itemize}
\tightlist
\item
  K038(CA13--CA21): exact sharp G support, phases and full compact covariance; CA22--CA27: full Y quotient and actual coupled class map.
\item
  LIFT26(GS6--GS7c): one fixed finite Schwartz tensor under the same diagonal middle frame, before both theta sums.
\item
  K037(CR7--CR19): integral subD composition, exact selected oscillator intertwining, and perfect Weyl/dual-Weyl duality; CR27--CR30: complete coevaluation and its normalization.
\item
  K039(IG2--IG7): no admitted Q-over11, standard local frames, finite splitting units, prime-to11 Heisenberg expansion and integral-section reading; IG8--IG10 and K041: complete-module derivative bounds.
\item
  K042(NB6): actual initial strict stabilizers are trivial.
\item
  LIFT33(SI8a--SI10): remaining terminal bound and exact class scope.
\item
  Primary recovered source \nolinkurl{Formalism/LocalToGlobal/Columns/FiniteSumNonzeroLargeDepth.lean}: 298--323 actual class masses; 780--810 PhiGS; 1103--1167 the common middle frame and actual XiCpl1/PhiTop/ftY/XiCpl2; 1180--1275 exact Y transport and CplKer covariance. No Lean build is made.
\end{itemize}

The companion finite certificate tests the free residue action, the actual support and phase invariance in a small sharp model, the cutoff partition and norm identity, and an integral coevaluation diagnostic. It is a diagnostic of these finite mechanisms, not a computation of the enormous actual global theta matrix.

\subsubsection{Completed author verification}\label{completed-author-verification}

The author read all of sections1--8 and equations(MC1)--(MC22), including the full-mass corollary(MC21a--b), after the final changes. The review checked exact support and inverse finite covariance, the common-frame conjugation, the distinction between local and global position models, the real coevaluation with perfect integral duality, honest rational theta automorphy, the separated Heisenberg expansion, the whole-module derivative estimates and target-lattice descent. It also re-read the literal recovered FS1103--1185 definitions and the cited K037--039 identities. No terminal restriction or positive Smith assertion was imported.

The exact certificate \nolinkurl{coupled_image_sylow_cancellation_certificate.py} and its JSON output completed320,458 assertions. They include all12,144 matrices in SL2(F23), an order11 subgroup and its free sharp-support action, 12,144 full-group inverse/transversal checks,2,277 nontrivial row/phase transports, all529 exact Fourier inversions over Z{[}zeta23{]}, the actual Weyl/dual-Weyl Sym2 contraction with its factor2,729 integral evaluation checks, the honest determinant-character cancellation, and both actual G-versus-rational-dimension margins. No probabilistic numerical test, global theta matrix computation, or Lean build is involved.

\clearpage

\section{AM-DO-045 --- two obstructions to the direct near-equal-rank theta shortcut}\label{am-do-045-two-obstructions-to-the-direct-near-equal-rank-theta-shortcut}

Author: GPT-6 Astra Pro, ALP\_MACHINE \nolinkurl{/root/definite_oscillator}. 22 November2026 UTC. Inherited reasoning effort is not separately exposed.

The direct theta lift from the actual compact H=U(I\_a) to the actual bottom G=U(I\_a,-1), where a=42 or168, cannot supply the prescribed bottom detector. There are two independent reasons. At any top CM label, the infinitesimal characters are incompatible throughout the entire K044 gap interval, even after arbitrary determinant twists. At the actual split G primes, the genuine goodTuple projector excludes every cuspidal constituent of such a direct lift, whether or not the role assignment has a top label.

These are obstructions to this one-edge construction. They are not vanishing theorems for the H-branching period, for the previously constructed two-step vertical tower, or for the restricted terminal value of the actual coupled theta map.

\subsection{1. An unconditional infinitesimal-character necessity}\label{an-unconditional-infinitesimal-character-necessity}

For a real type-I unitary dual pair of ranks a and a+1, use the metaplectic preimages and the bare oscillator representation. In standard complex Cartan coordinates, the infinitesimal-character map from the smaller member to the larger is

\[
        (x_1,\ldots,x_a)\longmapsto(x_1,\ldots,x_a,0),
                                                        \tag{DO45.1}
\]

up to the two type-A Weyl groups. This is Przebinda(1996), Theorems1.13 and1.19, equations(1.16)--(1.18), printed p97. For GL\_a and GL\_(a+1), his added sequence has its single entry (a+1+(a+1))/2−(a+1)=0. His theorem applies to all irreducible oscillator quotients, independently of the real form. It does not require an Arthur multiplicity formula or a stable-range hypothesis.

An honest unitary splitting multiplies the bare actions by characters. On either connected unitary group, the differential of any character vanishes on the derived special unitary Lie algebra. It therefore adds a constant to every Cartan coordinate. A determinant twist has the same effect. Contragradients negate the coordinates, and Weyl permutations or changing the ordered signature permute them. A change of sign of the real additive character replaces the oscillator by its contragredient; a positive rescaling is an isomorphic oscillator after symplectic dilation. All these operations preserve the width

\[
                 w(x)=\max_i x_i-\min_i x_i.            \tag{DO45.2}
\]

Consequently every occurrence of a prescribed pair of real representations in any of these oscillator conventions satisfies

\[
               w(\mathrm{HC}_H)\le w(\mathrm{HC}_G).
                                                        \tag{DO45.3}
\]

Indeed, after undoing the constant shifts and possible negations, the larger multiset contains the smaller multiset by(DO45.1). Containment cannot decrease width. This proof also applies on a common cover when one bare compact action is genuine rather than an action of the linear unitary group.

\subsection{2. The actual top labels fail the necessity}\label{the-actual-top-labels-fail-the-necessity}

Retain K044(EG1--EG5), including every permitted word, role assignment, canonical twist and gap. Put

\[
 A=2^{960},\quad B=2^{955},\quad
 \ell_-\le\ell\le\ell_+,\quad0\le q\le\ell,
 \qquad a=42\text{ or }168.
                                                        \tag{DO45.4}
\]

At a top label, the actual initial compact Harish-Chandra endpoints are

\[
 u_0=A+(a-1)/2,\qquad u_{a-1}=-B-(a-1)/2.
                                                        \tag{DO45.5}
\]

Their difference is independent of all word bits and of ell and q:

\[
                         w(u)=A+B+a-1.                 \tag{DO45.6}
\]

The prescribed bottom signed Harish-Chandra list, before taking the contragredient required by the literal period convention, is

\[
 \Lambda=(v_0,\ldots,v_{a-2},-b_0,-b_1),\quad
 v_i=4\ell(a-1-i),\quad b_0=E_T-\ell,\quad b_1=E_T+\ell+q.
                                                        \tag{DO45.7}
\]

These entries are strictly ordered in the displayed order. Its width is

\[
 w(\Lambda)=4\ell(a-1)+E_T+\ell+q
       \le E_T+(4a-2)\ell_+<2B.                        \tag{DO45.8}
\]

The last inequality is exact on both actual hosts. One can either use the sharper K044 certificate or the sufficient bounds E\_T\textless B+2\^{}485, ell\_+\textless2\^{}485 and4a−1\textless=671, together with B=2\^{}955. The companion certificate verifies the actual integer inequality. Since A=32B, (DO45.6) is strictly larger than(DO45.8).

It follows that no direct rank-a to rank-(a+1) unitary theta lift of any actual initial automorphic constituent can have this prescribed bottom infinite type at a top label. This includes every choice of Schwartz function, every determinant twist of either member, all honest unitary splitting characters, both contragredient conventions and both signs of the real additive character. A global occurrence would give the excluded local occurrence at that label.

This argument needs at least one top label. It does not assume that all role assignments have one. The following finite argument has no such qualification.

\subsection{3. The genuine finite e0 excludes the entire direct cusp image}\label{the-genuine-finite-e0-excludes-the-entire-direct-cusp-image}

\subsubsection{3.1. A GL1 block in every local near-equal theta quotient}\label{a-gl1-block-in-every-local-near-equal-theta-quotient}

Let F\_v be a nonarchimedean local field, K\_j=GL\_j(O\_v), K\_j\^{}1=ker(K\_j→GL\_j(k\_v)), and let pi be an arbitrary irreducible smooth representation of GL\_a(F\_v). At a split unitary place the oscillator is the type-II GL\_a×GL\_(a+1) oscillator, up to the two actual determinant characters and the source contragredient from the Hom convention. Mínguez(2008), Theorem1(2), printed p718, determines its unique irreducible theta quotient: if pi is the Langlands quotient of tau\_1×\ldots×tau\_s, the quotient on the larger group is the Langlands quotient of

\[
                  1\times\widetilde\tau_1
                         \times\cdots\times\widetilde\tau_s.
                                                        \tag{DO45.9}
\]

Here each tau\_i is essentially square-integrable and the first factor is a GL1 character. The exact splitting/double-dual dictionary of DO028 may twist these factors and invert pi, but it never removes the GL1 factor. No assertion about whole-Schwartz holomorphy at a special zeta point is required for this quotient theorem.

More generally, let Pi be any irreducible subquotient of normalized parabolic induction with one GL1 character factor eta. Then Pi\^{}(K\_(a+1)\^{}1) is a subquotient of finite ordinary parabolic induction with a GL1 character factor, or is zero. Here is a direct verification. By local Iwasawa decomposition, restricting the smooth parabolic induction to K\_(a+1) is induction from P∩K\_(a+1). Taking principal congruence invariants first reduces each Levi factor to its own principal-congruence invariants; the finite unipotent radical acts trivially. Thus the result is ordinary induction from the corresponding finite parabolic. If eta is nontrivial on1+varpi O\_v, its invariants are zero and so are these invariants of the induction. The modulus character is one on the compact. Compact invariants are exact over C, so passage to any subquotient preserves the assertion.

Under the Green characteristic map, ordinary parabolic induction is multiplication of symmetric functions in the orbit alphabets. A GL1 character contributes a Schur factor of positive degree in an orbit alphabet of degree1. Multiplication cannot remove that orbit support. Hence every finite irreducible constituent of the preceding module has a degree-one factor in its semisimple parameter polynomial. This is the precise graded-product consequence of Gow--Vinroot(2007), §3, printed pp8--9, Theorem3.2, already checked in DO044§3.2. In particular every genuine semisimple-series projector whose polynomial has no linear factor annihilates it.

\subsubsection{3.2. Applying the literal goodTuple condition}\label{applying-the-literal-goodtuple-condition}

The actual source condition is \nolinkurl{Formalism/CyclotomicNorms/StatementT.lean}19--29. Its fifth conjunct is

\[
 \forall d\ge1,\ 2d\le n\Longrightarrow
 \exists v:\ \text{no divisor of }f_v\text{ has degree }d.
                                                        \tag{DO45.10}
\]

Take d=1 and n=a+1≥43. For every admitted tensor block of the genuine e0, this gives a selected G place at which its polynomial has no linear factor. The \nolinkurl{goodTuple} definition requires this at each of the two rational G primes, although only one such place is needed. The literal definition is \nolinkurl{Formalism/SharpRestriction/LevelLetterCocycles/DataN.lean}996--1007: e0 is the finite sum of tensor products of the eSer(f\_v) for goodTuple. R026/R029 supply the genuine Green-series realization. No abstract idempotent with arbitrarily chosen action is substituted here.

At that supplied place, §3.1 kills the corresponding block on every local rank-a to rank-(a+1) theta quotient. Therefore the whole tensor block is zero, and summing the actual blocks proves:

\[
 \boxed{\text{No cuspidal automorphic constituent of a direct theta
 lift }U(I_a)\to U(I_a,-1)\text{ survives the actual }e0.}
                                                        \tag{DO45.11}
\]

To make the global implication precise, pair any such theta output with a cuspidal irreducible bottom representation. The source unitary group is anisotropic, so its automorphic quotient is compact. Cusp decay on the target justifies the usual local equivariant functionals obtained from this pairing. A nonzero global pairing gives a nonzero local oscillator intertwiner at every place; local Howe duality makes its target the small theta quotient described above. The relevant principal-congruence fixed subspace is killed by the supplied eSer, so the global e0 pairing is zero. We claim this for the cuspidal projection; no semisimplicity of the entire big theta module is assumed.

The argument permits every initial automorphic constituent, every Schwartz input and every splitting character. A wild added character only makes the principal-congruence fixed space zero. A tame twist can change which linear factor occurs, but cannot change its degree. Inversion of the parameter has the same property. No local unramified assumption on the initial representation is used.

\subsection{4. Why this does not decide the actual branching period}\label{why-this-does-not-decide-the-actual-branching-period}

The period in R027 is restriction from G to its embedded compact H, with the actual finite T translation and the selected derivative. The groups H and G in that restriction are not being used as a reductive dual pair. Thus (DO45.1) is not a branching rule.

In fact LIFT28/K043 prove that the actual compact H type occurs in the prescribed bottom discrete series, with the stated Pieri multiplicity and PBW normalization. R028/LIFT32 likewise give the actual finite generic restriction modules and central compatibility. Those local branching occurrences are consistent with both obstructions above. They do not assert nonvanishing of an automorphic period, whose global spectral occurrence remains an independent gate.

The previously proved vertical adjoint tower has ranks a+1→a+1+t→ a+1+2t and prescribed different oscillator data; it is not the rank-a→a+1 shortcut excluded here. The full two-step coupled map, its arithmetic row assembly and its restricted terminal value are not changed by this finding. In particular (DO45.11) is not a vanishing theorem for every bottom cusp form or every endpoint period.

\subsection{References and verification}\label{references-and-verification}

Mínguez, A. (2008). Correspondance de Howe explicite: paires duales de type II. \emph{Annales scientifiques de l'École Normale Supérieure, 41}(5),717--741. Theorem1(2),p718; normalization(1.4),pp721--722. \url{https://www.numdam.org/item/ASENS_2008_4_41_5_717_0.pdf}

Przebinda, T. (1996). The duality correspondence of infinitesimal characters. \emph{Colloquium Mathematicum,70}(1),93--102. Theorems1.13,1.19, equations(1.16)--(1.18),p97. \url{https://doi.org/10.4064/cm-70-1-93-102} Author PDF: \url{https://tomasz.przebinda.com/correspondence_infinites_char_1996.pdf}

Gow, R., \& Vinroot, C. R. (2007). \emph{Extending real-valued characters of finite general linear and unitary groups on elements related to regular unipotents} {[}Preprint{]}. §3,pp8--9,Theorem3.2. \url{https://arxiv.org/abs/0704.2022}

The primary Przebinda PDF/text is retained in \nolinkurl{definite_oscillator_near_equal_refs/}; the other two primary PDFs/texts are in the shared \nolinkurl{lift_refs/}. All exact quoted code predicates were read from the recovered pinned source. No Lean build is claimed.

Author whole-proof check completed 2026-11-22T04:36:34.808873+00:00. The complete proof and certificate were read, including the primary infinitesimal-character formula, every splitting/duality width change, the actual endpoint inequalities, the type-II Langlands quotient, exact principal-congruence invariants of parabolic induction, genuine Green orbit support, the literal goodTuple quantifiers and the cuspidal global implication. PASS in the stated direct-theta scope. The exact certificate passed at04:35:29 UTC for both actual hosts and28 parabolic-index diagnostics. The finite Green and global occurrence proofs are mathematical arguments, not inferred from those elementary diagnostics.

\clearpage

\section{ALP-LIFT-35 --- an actual terminal stabilizer fixes the raw G residue support}\label{alp-lift-35-an-actual-terminal-stabilizer-fixes-the-raw-g-residue-support}

Author: GPT-6 Astra Pro, ALP\_MACHINE \nolinkurl{/root/lift_accounting}. 22 November2026 UTC. Inherited reasoning effort is not separately exposed.

This is a precise obstruction to extending LIFT34's first-residue transversal proof to the terminal class mass. It is not a terminal period vanishing theorem, an assertion that a terminal denominator survives in the actual theta image, or a positive Smith-factor result.

\subsection{1. The literal terminal class has no T decoration}\label{the-literal-terminal-class-has-no-t-decoration}

In the recovered primary source \nolinkurl{Formalism/LocalToGlobal/Columns/FiniteSumNonzeroLargeDepth.lean}, lines201--204, the actual terminal row points are

\[
 \operatorname{RowPt}
   =\operatorname{DecoSet}(E,n_T,\varnothing,G\cup Q),
 \qquad n_T=c-1=r+2t,
\tag{TC1}
\]

and tilQ subsequently adds the distinguished norm line and T data. Thus a terminal strict T-covector condition must not be imposed on the row class itself. This differs from the actual middle class in LIFT33. Keep the admitted trivial tuple, so Q is empty. Use the standard global Hermitian lattice O\_E\^{}\{n\_T\} and its standard sharp G flags of block sizes(r,t,t).

In literal weighted lamY coordinates choose the fixed split diagonal isometry sending that local source lattice to the standard lattice, with inverse-transpose action on the opposite split component. The normalized G flag is then the standard K038(CA14) flag. Outside G take the identity frame. This gives an actual member of the source base genus: all resulting local lattices are standard, hence their global intersection is O\_E\^{}\{n\_T\}. No new level has been chosen.

The full integral unitary automorphism group of the standard lattice is mu28\^{}\{n\_T\} semidirect S\_\{n\_T\}, by the rank-independent proof of R017(R17.7). Intersect it with the terminal sharp flag group. A monomial matrix in the first block diag(U\_a,1) is the identity on those r coordinates. The two later GL\_t diagonal blocks allow exactly the monomial groups on their own coordinates. An upper-block parabolic monomial matrix cannot exchange these blocks: its occupied columns and rows force every off-diagonal monomial entry to vanish. Reduction is injective on mu28 at the actual G primes. There is no Q or T row condition. Therefore this actual terminal row stabilizer is

\[
 \Gamma_T=
       (\mu_{28}^{t}\rtimes S_t)
                 \times(\mu_{28}^{t}\rtimes S_t).
\tag{TC2}
\]

The two factors act on rows r through b-1 and b through n\_T-1, respectively, where b=r+t. This group statement does not presuppose that a particular branching detector has a nonzero invariant vector.

\subsection{2. A nontrivial11-subgroup fixes every supported residue}\label{a-nontrivial11-subgroup-fixes-every-supported-residue}

Both t=1759 and t=1633 exceed11. In the second S\_t factor take the 11-cycle on rows b,b+1,\ldots,b+10, and let P be the group it generates. It is an actual rational unitary subgroup of Gamma\_T of order11, with determinant1. Its finite coefficient character is1, because the unipotent blocks and both determinant characters in K038(CA14) see the identity. The tilQ extension is diag(gamma,1), so it also fixes the distinguished last line and its literal piT scaling.

On the second raw G support, write the first n\_T rows of the full Hom matrix as Z and its last row as eta. K038(CA17) is exactly

\[
               \bar Z=\begin{pmatrix}B\\0_{t\times b}\end{pmatrix},
                     \qquad B\in H^\sharp_{r,t}.
\tag{TC3}
\]

Every element of P acts only on the zero rows of(TC3). Hence

\[
        \overline{\gamma Z}=\bar Z,
             \qquad \eta\longmapsto\eta
                    \quad(\gamma\in P)
\tag{TC4}
\]

on \textbf{every} supported residue, although Z has full column rank b. The first edge, the middle Y coordinate and its full sum are unchanged by a terminal action. Thus adding those variables does not restore freeness. The assertion holds simultaneously at all G places, since the same rational row permutation lies in every standard sharp flag.

Consequently no0/1 function C of the existing first G residues can satisfy the orbit partition identity required in LIFT34:

\[
       \sum_{\gamma\in P}C(\gamma\cdot\bar z)
                  =11C(\bar z)\ne1
          \quad\text{on the nonempty raw support}.
\tag{TC5}
\]

This is an exact fixed-point obstruction to that finite transversal, not a failure of reduction injectivity. P itself reduces faithfully: its permutation matrix is visibly nonidentity. The different issue is that a larger target group can fix a full-column-rank embedding pointwise while acting on its complementary rows.

\subsection{3. Consequences and remaining scope}\label{consequences-and-remaining-scope}

The middle cutoff theorem LIFT34 remains valid. Its source-side action on a full-column-rank matrix is free, whereas the terminal action(TC4) is a target-side action. The latter cannot be substituted into the former proof.

A deeper G residue may distinguish matrices whose last t rows are nonzero multiples of g. Exact fixed subspaces still exist, and their theta contributions cannot be discarded merely because they have additive measure zero. A terminal argument would have to retain the actual brC, the literal tilQ extra line, the detector's stabilizer equations and the arithmetic normalization on any such fixed stratum. No result about those theta values is proved here.

In particular(TC5) does not imply that the terminal class period is nonintegral, that its11-mass fails to cancel by some other mechanism, that the terminal branching map kills the prior image, or that no detector exists. It isolates precisely why the already proved finite-transversal argument stops at this actual terminal example.

\subsection{4. Verification}\label{verification-2}

Author whole-proof reading completed and passed at04:40:57 UTC on 22 November2026.

The author checked the literal empty-ST RowPt and tilQ definitions, the actual terminal block/character formulas K038(CA14,CA17--CA21), and the complete standard-lattice monomial proof R017(R17.7). The group intersection in(TC2), determinant-one11-cycle, and fixed-support identity(TC4) are proved above for both actual ranks. These are algebraic identities; no global theta numerical test or Lean build is claimed.

\clearpage

\section{AM-K048 --- actual packet frequencies and a normalized local spectral tail}\label{am-k048-actual-packet-frequencies-and-a-normalized-local-spectral-tail}

Author: GPT-6 Astra Pro, ALP\_MACHINE \nolinkurl{/root/kummer_extension}. 22 November 2026 UTC; inherited reasoning effort is not separately exposed.

This is a quantitative \textbf{local representation} result at the actual translated packet point of K047. It keeps the actual selected compact module, every allowed word and canonical twist at the K044 upper gap, and the exact PBW isometry. With at most half of the labels top, the unitary spectral transform of that module is extremely small on the actual arithmetic central-frequency lattice. The result includes a summed sampling bound, rather than only a single-frequency estimate.

It is not yet a bound for the automorphic evaluation distribution or for a global selected-jet matrix. The latter distinction matters: a small value of a normalized local spectral function does not bound the normalization of a global Fourier coefficient. Section 8 gives the precise outstanding comparison and a possible invariant-cutoff route. The current mixed numerical profile permits arbitrary role partitions; the at-most-half-top hypothesis is explicit and is not inferred from an unavailable literal source Frame.

\subsection{1. The actual point and its central periods}\label{the-actual-point-and-its-central-periods}

Retain K047, with Q=g0 g1, actual norm generator pi, prime T=pi bar(pi), and the rational two-plane matrix

\[
 g=\begin{pmatrix}(T+1)/(2\pi)&(T-1)/(2\bar\pi)\\
                  (T-1)/(2\pi)&(T+1)/(2\bar\pi)\end{pmatrix}.
 \tag{ST1}
\]

The other a−1 coordinates are fixed. Here pi denotes the algebraic norm generator, not the circular constant appearing below. At every complex label the image of the ball origin is

\[
       g\cdot0=\frac{T-1}{T+1}e_1.
 \tag{ST2}
\]

Indeed it is the quotient of the last-column entries B/bar(pi) and A/bar(pi). In the right-half-plane coordinate w=(1+z1)/(1−z1), at zero transverse coordinates, this point has w=T. Thus the actual cusp height is exactly T, with no period or norm-generator phase left. The two-plane matrix is a positive dilation taking w=1 to w=T, followed by a compact SU(1,1) element. That compact element preserves the selected compact module and its norm.

Put v=e1+e\_n and N=v v* J, where J=diag(I\_a,−1). Then N²=0. For every b in O\_F,

\[
       u_b=1+iQbN\ \in\Gamma_1.
 \tag{ST3}
\]

All entries are integral, the determinant is1, and the matrix is the identity at both principal G congruences. This is K047(SP4--SP5) with Q replaced by Qb. Its action on w is w↦w−2iQb at an embedding where i has value i. At a general chosen CM label tau write

\[
                \epsilon_\tau=\tau(i)/i\in\{1,-1\}.
\]

The real central-period lattice is therefore the signed image of Q O\_F. The signs must be retained: the CM labels do not necessarily all restrict to the same complex embedding of Q(i).

Use the self-adjoint complexified sl2 generator X=−N. With H=diag(1,−1), B=E12 and C=E21 it is X=−H+B−C. The group element in (ST3) is exp(−iQbX). The positive discrete-series spectral variable for −X is x\textgreater0. Its possible nonzero arithmetic frequencies before dilation are

\[
 x_\tau=\frac{2\pi}{Q}\epsilon_\tau\sigma_\tau(\xi)>0,
 \qquad \xi\in\mathfrak d_F^{-1},\quad\xi\ne0.
 \tag{ST4}
\]

Here and henceforth the unitalicized mathematical constant in 2π is the circular constant. Formula (ST4) is exactly trace duality for O\_F: Q times the frequency divided by 2π must pair integrally with all b. Depending on the contragredient convention all the signs can be reversed. Every norm and bound below is unchanged.

Dilation to height T replaces the spectral argument by

\[
             x_\tau(\xi)=\frac{2\pi T}{Q}|\sigma_\tau(\xi)|.
 \tag{ST5}
\]

One can verify the scale without a Fourier convention: g\^{}−1 N g is T\^{}−1 times a compact conjugate of N, by K047(SP6--SP8). The unitary spectral measure will be dx/x, which adds no dilation scalar.

\subsection{2. Exact compact energy of the selected component}\label{exact-compact-energy-of-the-selected-component}

Use the actual K043/K044 parameters at ell=ell\_+. At one label the minimal compact type is V\_alpha tensor z\^{}beta and the selected type is V\_lambda tensor z\^{}(beta−N\_tau), where

\[
 \alpha_i=p_i-(a-2)/2+i,\quad \beta=-b_1-a/2,\quad
 \lambda_i=u_i-(a-1)/2+i,\quad
 N_\tau=\sum_i(u_i-p_i-1/2).
 \tag{ST6}
\]

Consider the SU(1,1) subgroup on the first positive and the negative coordinate. Its compact circle is diag(e\^{}(it),1,\ldots,1,e\^{}(−it)). On a weight vector of V\_lambda with first weight coordinate mu1, its compact energy is mu1−beta+N\_tau. Every such mu1 lies between the last and first entries of lambda. In particular the maximum is

\[
 K_\tau=\lambda_0-\beta+N_\tau
       =u_0+b_1+N_\tau+\tfrac12.
 \tag{ST7}
\]

This is an integer, and bounds the whole selected compact module, not just its highest vector. In the full holomorphic model the compact energies are positive and have finite-dimensional eigenspaces. To see this directly, the minimal fibre has lowest energy alpha\_(a−1)−beta=b1−b0+a=2ell+q+a, and a polynomial variable contributes one or two positive units of energy. Thus bounded energy permits only bounded polynomial degree. In particular

\[
                     K_\tau\ge K_*:=2\ell+a.
 \tag{ST8}
\]

The restriction to this SU(1,1) is an orthogonal sum of positive lowest-weight discrete series. An elementary proof suffices here. Take the lowest remaining circle eigenspace. Its lowering operator annihilates it. The sl2 adjunction and commutator identities generate orthogonal standard lowest-weight ladders from an orthonormal basis of that space. Remove their closed invariant sum and repeat. The positive integral energies and finite-dimensional energy spaces ensure exhaustion of the dense algebraic vectors. Every ladder has an integral lowest weight kappa\textgreater0. On the part reached by the selected module its ladder index n satisfies

\[
                         \kappa+2n\le K_\tau.
 \tag{ST9}
\]

No scalar lowest-weight replacement of the vector-valued discrete series is being made; all these ladders and their multiplicity spaces are retained orthogonally.

For later uniform bounds put

\[
 K_T=3A,\qquad K_H=2E_H+(16a-9)\ell+4,
 \qquad A=2^{960}.
 \tag{ST10}
\]

Then every top label has K\_tau\textless K\_T, and every other label has K\_tau\textless K\_H. For completeness, with w the number of interior ones,

\[
 N_T=A+E_T-B-(5+w)\ell-a+1,
 \qquad
 N_{H,L}=E+(8a-7-w)\ell-a+2.
 \tag{ST11}
\]

Substitution in (ST7), and 0≤q≤ell, give (ST10). At a top label use E\_T=B+ell+40 or B+ell+96 and A=32B. The non-top bound is obtained by putting w=0,q=ell,E=E\_H. The exact endpoint integers give

\[
 K_T<2^{962},\qquad K_H<2^{497},\qquad K_*\ge2^{485}.
 \tag{ST12}
\]

The actual selected inclusion is J\_tau=rho\_tau j\_tau, with rho\_tau²=product\_i(r\_i)\_(m\_i), as in K043(PB5--PB10). It is isometric for the given invariant Hilbert norm. All local vectors in the next sections are measured using this isometry. Thus the large raw PBW factor has already been divided out in exactly the prescribed manner.

\subsection{3. A normalized Laguerre model and two elementary estimates}\label{a-normalized-laguerre-model-and-two-elementary-estimates}

For integers kappa≥1,n≥0 define

\[
 f_{\kappa,n}(x)=
 \sqrt{\frac{n!}{\Gamma(n+\kappa)}}\,
 x^{\kappa/2}e^{-x/2}L_n^{\kappa-1}(x),\qquad x>0.
 \tag{ST13}
\]

These form the unitary spectral transform of the circle basis of lowest weight kappa, using measure dx/x. The exact primary anchor is Groenevelt (2003), Section 3.1, equation (3.2) and its orthogonality, Section 3.2, equation (3.7), and Proposition 3.2 in Section 3.3, author PDF pages9--12. His lowest-weight parameter is kappa/2; multiplying his normalized polynomial by the square root of its weight times x changes his measure to dx/x and gives (ST13). This also verifies the factor in (ST4), because his X has eigenvalue −x in the positive series.

Write D=x d/dx. The two elementary Laguerre recurrences give

\[
 D f_{\kappa,n}
  =\tfrac12\sqrt{(n+1)(n+\kappa)}f_{\kappa,n+1}
   -\tfrac12\sqrt{n(n+\kappa-1)}f_{\kappa,n-1}.
 \tag{ST14}
\]

The term with index−1 is absent. Orthogonality therefore gives

\[
 \|Df_{\kappa,n}\|_2^2
  =\tfrac14\{(n+1)(n+\kappa)+n(n+\kappa-1)\}.
 \tag{ST15}
\]

In the coordinate y=log x, integrate the derivative of \textbar f\textbar² from −infinity to y and apply Cauchy--Schwarz. Since \textbar\textbar f\textbar\textbar\_2=1,

\[
 |f_{\kappa,n}(x)|^2\le2\|Df_{\kappa,n}\|_2
          \le2(K+1)\quad\hbox{if }\kappa+2n\le K.
 \tag{ST16}
\]

This uniform polynomial bound is important. Replacing it by an exponential bound in the largest selected degree would lose the useful joint-label estimate.

There is also a particularly simple large-x estimate at the present scales. Suppose 1≤K\textless2\^{}962, kappa+2n≤K and x≥2\^{}180 K. The explicit coefficient formula

\[
 L_n^{\kappa-1}(x)=
       \sum_{j=0}^n(-1)^j\binom{n+\kappa-1}{n-j}\frac{x^j}{j!}
\]

and n!/Gamma(n+kappa)≤1 show, since x≥1,

\[
 |f_{\kappa,n}(x)|
  \le(K+1)2^K x^{K/2}e^{-x/2}\le e^{-x/4}.
 \tag{ST17}
\]

Here is an explicit check of the last inequality. At x0=2\^{}180 K, log x0\textless1142, log(K+1)\textless963 and log2\textless1, so the logarithm of the prefactor is less than1535 K. This is less than x0/4. The derivative of x/4−(K/2)log x is positive for x\textgreater2K, so the inequality persists for every x≥x0. The constant1535 harmlessly covers K=1 as well. The certificate checks1535\textless2\^{}178 exactly.

\subsection{4. Bounds for the whole selected vector space}\label{bounds-for-the-whole-selected-vector-space}

Let T\_tau(x) be evaluation of the preceding unitary spectral transform on the selected compact module, with target the orthogonal sum of all its ladder multiplicity spaces. It is a well-defined finite-dimensional operator at each x\textgreater0. By Cauchy--Schwarz over the at most K+1 eligible indices n in each fixed ladder, (ST16) gives

\[
                   \|T_\tau(x)\|_{op}^2\le8K_\tau^2.
 \tag{ST18}
\]

The orthogonal sum over different lowest weights costs no extra multiplicity factor. Similarly (ST17) gives, when x≥2\^{}180 K\_tau,

\[
                  \|T_\tau(x)\|_{op}^2
                       \le2K_\tau e^{-x/2}.
 \tag{ST19}
\]

For a vector in the tensor product of selected modules, including an entangled vector, the product spectral evaluation is the tensor product of these operators. Hence its operator norm is bounded by the product of their operator norms. No assumption of a decomposable vector, of one scalar fibre, or of one selected weight vector enters these bounds. The compact phase in (ST1) is unitary on the selected module and leaves them unchanged.

\subsection{5. The inverse different forces a far-tail label}\label{the-inverse-different-forces-a-far-tail-label}

For a nonzero xi in the inverse different,

\[
 |N_{F/\mathbf Q}(\xi)|\ge |\operatorname{disc}F|^{-1}.
 \tag{ST20}
\]

Indeed the integral ideal (xi) d\_F has positive integer norm. The actual extension E/F is unramified at finite places, so rd(F)=rd(E). The already checked discriminant inequalities in DO044 give rd(F)\textless2\^{}8 for J14 and rd(F)\textless2\^{}11 for J56. Write r=8 or11. From (ST5), (ST20), and 2π\textgreater6,

\[
       \prod_\tau x_\tau(\xi)
                    >\left(\frac{6T}{Q2^r}\right)^f.
 \tag{ST21}
\]

Assume now that at most f/2 labels are top. Since K\_T\textgreater K\_H,

\[
              \prod_\tau K_\tau\le(K_TK_H)^{f/2}.
 \tag{ST22}
\]

The exact actual integers satisfy

\[
 \frac{6T}{Q2^r\sqrt{K_TK_H}}
       >\begin{cases}2^{221}&J14,\\2^{187}&J56.\end{cases}
 \tag{ST23}
\]

This is certified by squaring and comparing integers, without computing a square root or a logarithm. In particular for every admissible nonzero frequency xi some label obeys

\[
                         x_\tau(\xi)>2^{180}K_\tau.
 \tag{ST24}
\]

The sign pattern in (ST4) has no effect on the norm argument. It is not legitimate to call xi totally positive without first checking that sign pattern; the proof uses its absolute norm.

\subsection{6. A summed arithmetic sampling bound}\label{a-summed-arithmetic-sampling-bound}

Define the sampling map on the unitary selected module by

\[
 S_T(v)=\bigl((\mathop{\otimes}_\tau T_\tau(x_\tau(\xi)))v
       \bigr)_{\xi},
 \tag{ST25}
\]

where xi ranges over the nonzero inverse-different lattice points in the sign orthant of (ST4). Give the target the ordinary Hilbert sum norm, with counting measure on this lattice. This defines a specific local sampling operator; no global Fourier coefficient is used in its definition.

\subsubsection{Theorem ST26}\label{theorem-st26}

At ell=ell\_+, for every word and allowed q, and with at most f/2 top labels,

\[
             \|S_T\|_{op}^2<\exp(-2^{663}),
             \qquad \|S_T\|_{op}<\exp(-2^{662}).
 \tag{ST26}
\]

\textbf{Proof.} Put y\_tau=x\_tau/K\_tau and K\_* as in (ST8). Distinct inverse-different lattice points differ by a nonzero element of that same ideal. Thus the scaled y-lattice has infinity-norm separation at least

\[
 \delta=\frac{6T}{Q2^r K_T}>2^{-47}.
 \tag{ST27}
\]

The exact integer inequality is checked for both hosts. The points all lie in the positive orthant. Cubes of side2\^{}−47 show that the number with max y\_tau≤R is at most (1+2R·2\^{}47)\^{}f.~This remains an upper bound if some coordinates or boundary points are absent.

Every sample has max y\_tau\textgreater R0=2\^{}180 by (ST24). In the shell m≤max y\_tau\textless m+1, select a label attaining the maximum. Its spectral argument is at least K\_*m, so (ST18--19) give the uniform squared operator bound

\[
          P e^{-K_*m/2},\qquad
                P=\prod_\tau8K_\tau^2<2^{1927f}.
 \tag{ST28}
\]

Summing over the shells m≥R0 bounds the squared norm of (ST25) by

\[
 P\sum_{m\ge R0}(1+2(m+1)2^{47})^f e^{-K_*m/2}.
 \tag{ST29}
\]

The ratio of consecutive summands is at most exp(f/(m+1)−K\_\emph{/2)\textless1/2, since f≤168 and K\_}≥2\^{}485. Hence the sum is less than twice its first term. Also 1+2(R0+1)2\^{}47\textless2\^{}230. Its entire logarithmic prefactor is therefore less than

\[
       (1+168(1927+230))\log2<2^{19},
\]

whereas K\_*R0/2≥2\^{}664. Consequently (ST29) is less than exp(2\^{}19−2\^{}664)\textless exp(−2\^{}663). Taking the square root proves the second inequality. This proof applies to the operator norm, not just to a sampled collection of vectors. QED.

\subsection{7. What is retained, and what is not inferred}\label{what-is-retained-and-what-is-not-inferred}

The theorem retains the actual T, the actual principal-G periods, the ordinary J14/J56 local unitary representations, the full selected compact modules, their large derivative orders, and the exact PBW normalization. A raw selected derivative reinstates rho\_all, exactly as in K043; it must not be compared directly with (ST26) without it. The theorem concerns the normalized selected vector/jet convention.

The genuine finite projector e0 is orthogonal in a unitary finite realization. Tensoring (ST25) with a finite Hilbert space and applying such a projector cannot enlarge this local bound. This observation does not identify the actual global Fourier coefficient map, replace its level volumes, or establish an automorphic evaluation bound.

The role hypothesis is substantive. R022/K044 allow all-top role partitions. The product comparison using only K\_T on every label has 6T/(Q2\^{}r K\_T)\textless1 on both current hosts, so the present argument cannot force (ST24) in that case. This is a failure of this sufficient frequency-window comparison. It is not a proof that an admissible frequency lies near all the local peaks or that the actual global norm is large in the all-top case.

\subsection{8. Exact remaining automorphic comparison}\label{exact-remaining-automorphic-comparison}

To turn (ST26) into a detector-norm lower bound one must control the actual automorphic evaluation distribution at the standard component on this local spectral realization. Global Fourier coefficients can have nontrivial normalizations and their unweighted square sum need not converge. Therefore no unweighted bounded coefficient operator is assumed here. A weighted coefficient norm or a direct invariant spectral-cutoff/Sobolev comparison would suffice, but must be proved with the actual Haar measures and finite level volumes.

There is a concrete alternative route. Because the evaluation distribution is invariant under the central lattice (ST3), a smooth spectral multiplier Phi equal1 on its dual lattice and given by an absolutely convergent Fourier series of those lattice translations leaves that distribution unchanged. A sufficiently narrow periodic cutoff can retain only small neighbourhoods of the dual lattice. Away from zero, the estimates above have large margins. Near zero, one must additionally bound the normalized Laguerre functions using their exact vanishing power kappa/2. Commuting a Lie derivative past Phi involves at most two multiplier derivatives, because ad\_N\^{}3=0 on this rank-one parabolic grading. The standard-origin separation of K046 could then provide a local Sobolev evaluation bound.

This route still needs an explicit global Sobolev constant, control of all the commuted vector-valued derivatives, the zero-frequency window, convergence in the necessary smooth topology, and the exact finite-volume conversion. None of these is declared complete by (ST26). In particular no actual terminal e\_j, endpoint interpolation cost, or DP44 sign is asserted in this finding.

\subsection{Primary reference and verification}\label{primary-reference-and-verification}

Groenevelt, W. (2003). Laguerre functions and representations of su(1,1). \emph{Indagationes Mathematicae, 14}(3--4),329--352. \url{https://doi.org/10.1016/S0019-3577(03)90050-4} . Author PDF \url{https://arxiv.org/pdf/math/0302342} , Section 3.1 equation (3.2), Section 3.2 equation (3.7), Proposition 3.2, PDF pages9--12. The retrieved primary PDF and extracted text are \nolinkurl{work/research/groenevelt2003_laguerre.pdf} and \nolinkurl{.txt}.

The unitary Laguerre model is a known theorem. The elementary bounds, actual compact-energy/CM-frequency comparison and summed sampling application are the new calculations here. The companion certificate checks the Laguerre orthogonality and dilation recurrence exactly at small integer parameters, all actual profile inequalities and integer frequency margins, and the displayed exponential comparison constants. It does not certify the analytic proof by finite sampling and does not compute an automorphic Fourier coefficient or global jet matrix.

\subsection{Author verification record}\label{author-verification-record-1}

Author whole saved report and complete certificate read and reconstructed PASS at 2026-11-22T04:43:05.639533+00:00. The review includes the exact cusp-height/central-frequency scale, CM signs, positive-energy restriction and whole-selected-module energy bound, unitary Laguerre/PBW normalization, polynomial global bound and far tail, inverse-different arithmetic and summed sampling estimate. The local/global distinction and conditional top-count scope are retained. The final certificate replay passed775 exact checks at2026-11-22T04:42:51.083157+00:00. No automorphic coefficient or global Sobolev estimate is claimed.

\clearpage

\section{AM-R-031 --- an explicit global evaluation bound at the actual standard component}\label{am-r-031-an-explicit-global-evaluation-bound-at-the-actual-standard-component}

Author: GPT-6 Astra Pro, ALP\_MACHINE \nolinkurl{/root}, 22 November 2026 UTC. Reasoning effort is not separately exposed.

This proves the standard-point Sobolev comparison needed for the translated-point spectral-cutoff approach. It is a bound for actual automorphic functions, with the actual ball Haar and its finite-volume conversion. It does not identify a local sampled spectral function with an automorphic Fourier coefficient, and does not complete that cutoff argument by itself.

\subsection{1. Statement, derivatives and measures}\label{statement-derivatives-and-measures}

Let either actual host have a=f=42 or168, and put n=a+1. Set

\[
 G_\infty=\prod_{\tau=1}^f U(a,1),\qquad
 K_\infty=\prod_{\tau=1}^f(U(a)\times U(1)),\qquad
 d=fn^2<2^{23}.
 \tag{SB1}
\]

The exact dimensions are77,658 and4,798,248. Let Gamma1 be the standard arithmetic component at the unchanged principal-G level J1. Use the Haar measure with K\_infinity volume1 and, on each ball,

\[
 d\mu_{\mathrm{ball}}(z)
       =\pi^{-a}(1-\|z\|^2)^{-a-1}dV(z).
 \tag{SB2}
\]

For the whole finite level, the \textbf{raw} norm is the sum of these quotient L2 norms over all finite adelic components. A function on one component can instead use just that component's raw norm.

Use the following real basis X\_1,\ldots,X\_d of the product Lie algebra. Within each factor, use iE\_ii for all n diagonal coordinates; E\_ij-E\_ji and i(E\_ij+E\_ji) for1\textless=i\textless j\textless=a; and E\_i,n+E\_n,i and i(E\_i,n-E\_n,i) for1\textless=i\textless=a. Each has matrix operator norm1 and Frobenius norm at most sqrt2. They are mutually orthogonal for Re Tr(X*Y), with squared norms1 or2. Define

\[
 (R_XF)(g)=\left.\frac d{dt}\right|_{t=0}F(g\exp(tX)).
 \tag{SB3}
\]

For an ordered word I of length m, R\_I denotes the corresponding composition. The empty word is the identity. Let

\[
 \mathcal S_d(F)^2
   =\sum_{m=0}^d\ \sum_{I\in\{1,\ldots,d\}^m}
                      \|R_I F\|_{\mathrm{raw}}^2.
 \tag{SB4}
\]

The theorem holds for scalar functions and for functions with values in any Hilbert space, with no factor depending on that space's dimension. It suffices that F be smooth and that the displayed norms be finite. Then

\[
 \boxed{\quad
   \|F(1)\|\le C_d\mathcal S_d(F),\qquad
   C_d=2^{11d/2}r^{-d/2}(5d/r)^d,\quad
   r=\frac1{100d^2},\quad \log C_d<2^{30}.
 \quad}
 \tag{SB5}
\]

In particular, with the maximum over all displayed ordered words,

\[
 \boxed{\quad
   \|F(1)\|\le\exp(2^{32})
           \max_{|I|\le d}\|R_I F\|_{\mathrm{raw}}.
 \quad}
 \tag{SB6}
\]

No cusp-form or multiplicity-dimension estimate is used in this bound.

\subsection{2. An embedded product-exponential chart}\label{an-embedded-product-exponential-chart}

On the cube(-r,r)\^{}d define

\[
      \Phi(t)=\exp(t_1X_1)\cdots\exp(t_dX_d).
 \tag{SB7}
\]

All matrices can be viewed in their common block-diagonal product representation. Put tau=sum\textbar t\_j\textbar\textless=dr\textless=1/100. Then

\[
 \|\Phi(t)\|_{op},\ \|\Phi(t)^{-1}\|_{op}\le e^\tau,
 \qquad \|\Phi(t)-1\|_{op}\le e^\tau-1\le2\tau.
 \tag{SB8}
\]

In particular the principal matrix logarithm is defined. Let B(t)=sum t\_jX\_j, viewed as a real linear map from Euclidean R\^{}d to the product Frobenius space. Its least singular value is at least1. Write h=e\^{}tau-1. For the derivative in any coordinate, product differentiation gives

\[
 \|\partial_i\Phi-X_i\|_F\le\sqrt2h,
 \qquad \|\partial_i\Phi\|_F\le\sqrt2 e^\tau.
 \tag{SB9}
\]

The convergent series for log(1+H), differentiated term by term for \textbar\textbar H\textbar\textbar\_op\textless1, gives

\[
 \|D\log_{1+H}(V)-V\|_F
                \le\frac{\|H\|_{op}}{1-\|H\|_{op}}\|V\|_F.
 \tag{SB10}
\]

Combining(SB9--10), h\textless=2tau, e\^{}tau\textless2 and1-h\textgreater1/2 bounds each column of D(log Phi)-B by10sqrt2 tau\textless15tau. Hence

\[
 \|D(\log\Phi)-B\|_{op}
       <15\sqrt d\,dr=\frac{15}{100\sqrt d}<\frac12.
 \tag{SB11}
\]

Integrate this derivative estimate on the straight segment between two points of the convex cube. The norm of their logarithmic image difference is at least one half the Euclidean distance. Thus Phi is injective on this cube, as well as an immersion.

The Cartan displacement in each U(a,1) factor is at most tau\textless1, by(SB8). The image therefore lies in the set C1 of K046. The checked standard-component separation K046 gives Gamma1 intersect C5 C5\^{}(-1)=\{1\}. Since C1 is contained in C5, this chart maps injectively to Gamma1\G\_infinity. Its use does not claim separation at any translated point of K047.

\subsection{3. A lower bound for the actual Haar density}\label{a-lower-bound-for-the-actual-haar-density}

Give G\_infinity the left-invariant metric whose value at identity is Re Tr(X*Y), summed over the factors. This metric is right-K\_infinity invariant. Its restriction to K\_infinity is the ordinary bi-invariant Frobenius metric. Let Vol(K\_infinity) denote its Riemannian volume.

An elementary volume upper bound suffices:

\[
             \operatorname{Vol}(K_\infty)<2^{4d}.
 \tag{SB12}
\]

Here are explicit details to fix its normalization. Elimination of the entries below the diagonal of a unitary matrix by complex Givens rotations expresses every U(a) matrix as a product of a(a-1)/2 factors

\[
 \begin{pmatrix}\cos\theta&-e^{i\phi}\sin\theta\\
 e^{-i\phi}\sin\theta&\cos\theta\end{pmatrix}
 \tag{SB13}
\]

in specified coordinate planes, followed by a diagonal unitary matrix. Choose0\textless=theta\textless=pi/2,0\textless=phi\textless=2pi, and a diagonal phases in{[}0,2pi{]}. This is a surjective smooth parametrization from a compact rectangle with a\^{}2 parameters. Each parameter derivative has Frobenius norm at most sqrt2, even after the other unitary factors are multiplied in. The volume Jacobian is consequently at most (sqrt2)\^{}(a\^{}2). The area/change-of-variables formula, with multiplicity at least one on regular values, gives Vol U(a)\textless=(2pi sqrt2)\^{}(a\^{}2)\textless2\^{}(4a\^{}2). The remaining U(1) has volume2pi\textless2\^{}4. Multiplication over the f factors proves(SB12), since f(a\^{}2+1)\textless=d.~The exceptional critical values have measure zero and do not alter this volume upper bound.

The associated quotient metric on the ball has value2 times its Euclidean real metric at the origin: each noncompact basis vector in section1 has squared length2 and moves the origin at unit speed in its indicated real coordinate. Its invariant volume density is therefore2\^{}a(1-\textbar\textbar z\textbar\textbar\^{}2)\^{}(-a-1)dV. Riemannian submersion and compact fibre integration now identify the Haar(SB2) exactly as

\[
 d\mu=\frac{d\operatorname{vol}_{G_\infty}}
                 {(2\pi)^{af}\operatorname{Vol}(K_\infty)}.
 \tag{SB14}
\]

For completeness, the left-trivialized i-th tangent vector of Phi is Ad(exp(-t\_dX\_d)\ldots exp(-t\_(i+1)X\_(i+1)))X\_i. Its difference from X\_i has Frobenius norm at most sqrt2(e\^{}(2tau)-1)\textless6tau. Thus the least singular value of this column map is at least1-6sqrt(d)tau\textgreater1/2. The pullback Riemannian volume density is at least2\^{}(-d).

Since2pi\textless2\^{}3, af\textless=d and(SB12) holds, (SB14) implies the convenient conservative bound

\[
                    \Phi^*d\mu\ge2^{-9d}\,dt.
 \tag{SB15}
\]

The quotient chart is embedded, so for any measurable function H,

\[
  \int_{(-r,r)^d}\|H(\Phi(t))\|^2dt
       \le2^{9d}\|H\|_{L^2(\Gamma_1\backslash G_\infty,d\mu)}^2.
 \tag{SB16}
\]

The full raw sum of component norms can only increase the right side.

\subsection{4. Coordinate derivatives in ordered right Lie derivatives}\label{coordinate-derivatives-in-ordered-right-lie-derivatives}

Let f=F composed with Phi. Product differentiation gives

\[
 \partial_i f(t)
     =\sum_{j=1}^d c_{ij}(t)(R_{X_j}F)(\Phi(t)),
 \tag{SB17}
\]

where the c\_ij are the coefficients of the left-trivialized tangent vector just computed. These coefficients are entire functions of the complex variables t. On the complex polydisc \textbar t\_j\textbar\textless2r the operator norm estimate for the adjoint, and orthogonality of the complexified matrix-unit basis, give

\[
               \sup\sum_j|c_{ij}(t)|<4d.
 \tag{SB18}
\]

Indeed the sum of the moduli of the coefficients is at most sqrt(d) times the Frobenius norm; the latter is at most sqrt2 exp(4dr)\textless2sqrt2. This also applies to the block-diagonal product basis.

After any ordered m coordinate differentiations, the chain rule expresses the result as

\[
        \sum_{|I|\le m} A_I(t)(R_I F)(\Phi(t)).
 \tag{SB19}
\]

We claim that, on \textbar t\_j\textbar\textless2r-mr/d,

\[
                    \sum_I|A_I(t)|\le(5d/r)^m.
 \tag{SB20}
\]

For m=0 the only coefficient is1. For the induction step, Cauchy's inequality across the decrease r/d of polydisc radius bounds the sum of the differentiated coefficient moduli by d/r times the previous bound. When differentiation falls on the group function, (SB17--18) bounds the new coefficient sum by4d times the previous bound. Since r\textless=1, their sum is at most5d/r times that bound. This proves(SB20). Cauchy's inequality here is applied only to the explicit analytic coefficients, not to the arbitrary smooth F.

In particular every mixed derivative partial\^{}epsilon f, with each epsilon\_j equal0 or1, satisfies on the real cube

\[
 \|\partial^\epsilon f(t)\|^2
       \le(5d/r)^{2d}
                   \sum_{|I|\le d}\|(R_I F)(\Phi(t))\|^2.
 \tag{SB21}
\]

This follows from the coefficient l1 bound and Cauchy--Schwarz; it holds for Hilbert-valued F as well.

\subsection{5. Elementary Sobolev inequality and the numerical constant}\label{elementary-sobolev-inequality-and-the-numerical-constant}

For a continuously differentiable Hilbert-valued function h on{[}-r,r{]}, the fundamental theorem of calculus, Cauchy--Schwarz and averaging give

\[
 \|h(0)\|^2\le r^{-1}\int_{-r}^r\|h(t)\|^2dt
                       +2r\int_{-r}^r\|h'(t)\|^2dt.
 \tag{SB22}
\]

Apply this in all d coordinates. The result is

\[
 \|f(0)\|^2
    \le r^{-d}\sum_{\epsilon\in\{0,1\}^d}
                (2r^2)^{|\epsilon|}\|\partial^\epsilon f\|_2^2
    \le2^dr^{-d}\sum_\epsilon\|\partial^\epsilon f\|_2^2.
 \tag{SB23}
\]

There are2\^{}d displayed epsilon. Combining(SB16),(SB21),(SB23) therefore gives

\[
       \|F(1)\|^2
          \le2^{11d}r^{-d}(5d/r)^{2d}\mathcal S_d(F)^2,
 \tag{SB24}
\]

which proves the first assertion of(SB5).

The logarithm of its constant is

\[
 \log C_d
     =d\{(11/2)\log2+(1/2)\log100+\log500+4\log d\}
     <108d<2^{30}.
 \tag{SB25}
\]

For the strict bound use log2\textless1, log100\textless5, log500\textless7 and log d\textless23; the two actual d are less than2\^{}23. The number of ordered words in(SB4) is at most(d+1)\^{}(d+1). Its square-root logarithm is less than12(d+1)\textless=24d. Thus the logarithm of the maximum-norm constant is less than132d\textless2\^{}31\textless2\^{}32, proving(SB6).

\subsection{6. Actual automorphic norm conversions and selected vectors}\label{actual-automorphic-norm-conversions-and-selected-vectors}

Let V(J1) be the total raw quotient volume. By the already checked R030 and the actual principal residue group Gcal,

\[
       V(J1)=\frac{m_n|\mathcal G|}{(a!)^f}.
 \tag{SB26}
\]

For the normalized-average norm, raw norm equals sqrt(V(J1)) times average norm. Consequently(SB6) becomes

\[
 \|F(1)\|\le
   \exp(2^{32})\sqrt{\frac{m_n|\mathcal G|}{(a!)^f}}
                \max_{|I|\le d}\|R_I F\|_{\mathrm{av}}.
 \tag{SB27}
\]

Alternatively, if the finite maximal compact has volume1 at all places, the actual finite J1 volume is1/\textbar Gcal\textbar. The unnormalized adelic norm squared is that finite volume times the raw component norm squared. In that convention the multiplier in(SB27) is instead exp(2\^{}32)sqrt(\textbar Gcal\textbar). Any different overall Haar normalization must be inserted explicitly through this same identity.

Suppose U is an isometric unitary automorphic realization of a representation, including its multiplicity Hilbert space, in the raw L2 sum. It intertwines d pi(X) with R\_X. Therefore the same bounds apply to evaluation of U(v), replacing the derivative norms by \textbar\textbar d pi(I)v\textbar\textbar. There is no factor from the number of automorphic copies. A bounded realization of norm c contributes exactly c. This is also valid on a chosen isotypic space; no simplicity or temperedness assertion is needed for this analytic inequality.

K043's normalized selected inclusion J=rho j is isometric. It may therefore be used in this application with its prescribed rho already retained. An unnormalized selected derivative, a separately rescaled real intertwiner, or a vector-valued automorphic convention with a different scalar-to-multiplicity metric must retain its actual scalar. This lemma does not identify such a scalar with1.

\subsection{7. Scope, sources and geometric interpretation}\label{scope-sources-and-geometric-interpretation}

K046 supplies only the standard-component C5 separation, used once to embed the small chart. K047 explains why it cannot be assumed at the translated point. R030 supplies the exact raw total volume and the matched ball Haar. K043 supplies the selected PBW isometry. Those are carried checked inputs, not new theorems proved here.

The new content is the fully explicit standard-origin automorphic Sobolev constant and the scalar/multiplicity/volume bookkeeping. The proof is given above using a product-exponential chart, a direct compact-volume parametrization, coordinate derivative bounds and the one-dimensional fundamental theorem of calculus. Its arbitrary Hilbert-valued form avoids a hidden coefficient-rank multiplier.

In the run's geometric dictionary this is an analytic estimate on metrized automorphic sections and their differential evaluations. It supplies no algebraic principal-parts surjectivity; the latter retains the already located EGA III.1 Theorem2.2.1(ii),p100 and EGA IV.4 16.11.2--3(ii),pp54--55. No new motivic, ell-adic purity, cycle comparison or Hecke RH claim is made.

The pending translated-point argument still has to prove that its periodic spectral multiplier preserves the actual evaluation distribution, controls every word in(SB4), and handles its window at zero frequency. This lemma does not infer any of those statements from K048's local sampling bound.

\subsection{Author whole-proof verification}\label{author-whole-proof-verification}

The author read every section and equation SB1--27 in the saved proof, including the complex-coefficient Cauchy estimate, exact Riemannian-to-ball density, compact Givens parametrization, Hilbert-valued tensor Sobolev inequality, all ordered derivative counts and both volume conversions. The exact standard-component separation statement K046 BL26--27 was reread. Completed PASS at 2026-11-22T04:57:57.935086+00:00. No translated-point cutoff theorem, global Fourier-coefficient bound or raw-to-normalized selected scalar was assumed.

\clearpage

\section{AM-DO-046 --- actual generic parameters, temperedness, and the remaining branching-period criterion}\label{am-do-046-actual-generic-parameters-temperedness-and-the-remaining-branching-period-criterion}

Author: GPT-6 Astra Pro, ALP\_MACHINE \nolinkurl{/root/definite_oscillator}. 22 November 2026 UTC. Inherited reasoning effort is not separately exposed.

This finding removes a global representation-theoretic hypothesis from the actual mixed family. It does not invoke the failed literal HostRaw wrapper. Every irreducible constituent of the actual initial Mram space has a generic elliptic Arthur parameter and is tempered at every place. Every cuspidal constituent in the actual genuine e0 bottom space has, in addition, a \textbf{simple} parameter: its base change is a conjugate-selfdual cuspidal representation of GL\_(a+1).

The statements apply throughout K044's allowed gap and canonical-twist interval. They assert properties of every existing bottom constituent, not the existence of such a constituent at an allowed bounded weight. The actual initial Mram space is nonzero, so the initial assertion is not vacuous. The exact global GGP formula consequently applies to any actual initial--bottom pair. It leaves a specific central L-value and the prescribed local test vectors to be checked; it does not establish the terminal restricted period or its lower bound.

\subsection{1. Actual groups, characters, and coefficient scope}\label{actual-groups-characters-and-coefficient-scope}

Use either K031 or K032 host E/F and retain its actual induced CM type Psi. Put

\[
 H=U(I_a),\qquad G=U(I_a,-1),\qquad n=a+1,
 \qquad (a,[F:\mathbb Q])=(42,42)\ \text{or}\ (168,168).
 \tag{AP1}
\]

Both groups are actual algebraic groups over F. H is anisotropic and G has F-rank one: its displayed form contains an isotropic vector e\_1+e\_n, and at every real place its Witt index is one. These facts do not depend on the old archive Host records.

Retain the unitary character eta constructed in K029/K032 and recorded in K040(NC3):

\[
 \eta_\infty(z)=\prod_{\tau\in\Psi}z_\tau/|z_\tau|,
 \qquad \eta|_{\mathbb A_F^\times}=\omega_{E/F}.
 \tag{AP2}
\]

It is everywhere unramified at finite places. Also eta\^{}c=eta\^{}(-1): for an E idele x, eta(x)eta(c(x))=eta(N\_(E/F)x)=omega\_(E/F)(N\_(E/F)x)=1. The last equality is the defining norm-kernel property of the quadratic class-field character. Thus eta is an actual conjugate-selfdual unitary Hecke character, with angular exponent one, and not a formal half-character chosen only at infinity.

Fix any role assignment, word, ell and q allowed by K044(EG1--EG5). Let pi\_H be any irreducible automorphic constituent of the actual initial Mram space with its prescribed compact infinity type. Let pi\_G be an irreducible cuspidal constituent of the genuine source e0 bottom space with its prescribed holomorphic or contragredient holomorphic discrete series at infinity. The right-action convention of R027 uses the latter contragredient; the argument below treats both because negation and reversal preserve all gap and integrality conditions.

The bottom projector is the actual finite operator supplied by R025--R029, formed from the literal \nolinkurl{goodTuple} and its genuine Green-series idempotents. Its action is on principal G-level vectors before taking the actual U\_a character part. It is not replaced by an arbitrary generic finite projector.

\subsection{2. The actual Harish--Chandra parameters have no difference one}\label{the-actual-harishchandra-parameters-have-no-difference-one}

For a group of rank N, write its Harish--Chandra parameter as an unordered multiset. The standard unitary normalization is the one in Mok(2015), §7.1, equation(7.1.3), printed190: the character (z/bar z)\^{}x contributes coordinate x. Integral algebraic coefficient weights give the grid Z+(N−1)/2.

At an initial label the ordered coordinates are the actual u\_i. At a bottom label they are

\[
       v_0>\cdots>v_{a-2}>-b_0>-b_1.
 \tag{AP3}
\]

K044 gives

\[
 v_i=4\ell(a-1-i),\qquad
 b_0=E-\ell,\qquad b_1=E+\ell+q.
 \tag{AP4}
\]

The bottom adjacent gaps are therefore 4ell, E+3ell and 2ell+q. All exceed one. The initial interior adjacent gaps are 4ell+ell(omega\_i−omega\_(i−1)), hence are at least3ell. At a non-top label its first and last adjacent gaps are at least4a ell+1 and5ell−1. At a top label the first endpoint A+(a−1)/2 is larger than every interior coordinate by much more than one, and its last endpoint −B−(a−1)/2 is smaller than every interior coordinate by much more than one, using ell\textless2\^{}485, A=2\^{}960 and B=2\^{}955. Thus

\[
       |x-y|>1\quad\text{for any two distinct coordinates}
       \quad x,y\in\operatorname{IC}(\pi_{H,\tau})
       \text{ or }\operatorname{IC}(\pi_{G,\tau}).
 \tag{AP5}
\]

The initial coordinates lie in Z+1/2, the bottom coordinates in Z, exactly the respective grids Z+(a−1)/2 and Z+a/2. Negation and a common determinant translation preserve(AP5); the actual profiles already have the displayed grids.

The certificate checks all adjacent coordinate-bit choices and the affine ell,q endpoints. It gives minimum tested gaps 45,156,047,496 and40,530,337,904; the proof of affine endpoint reduction covers the complete stated interval. No exponential list of words is sampled in place of that reduction.

\subsection{3. An exact passage from the actual infinity type to the global Arthur parameter}\label{an-exact-passage-from-the-actual-infinity-type-to-the-global-arthur-parameter}

The following lemma is useful independently of the finite projector.

\subsubsection{Lemma AP6: wide gaps eliminate the Arthur SL2}\label{lemma-ap6-wide-gaps-eliminate-the-arthur-sl2}

Let U be an actual unitary group over a number field, and let pi be an irreducible square-integrable automorphic representation of U. Suppose at one real place pi\_v is tempered and its Harish--Chandra parameter contains no pair of coordinates differing by one. Then the global elliptic Arthur parameter of pi is generic.

\textbf{Proof.} Chen--Zou(2025), Theorem2.1 and Proposition4.1, attach an elliptic parameter

\[
                 \Psi=\boxplus_i\Pi_i\boxtimes S_{d_i}
 \tag{AP6a}
\]

to the actual near-equivalence class. Here each Pi\_i is a conjugate-selfdual cuspidal GL representation, and S\_d is the d-dimensional irreducible representation of SL2(C). This theorem applies to actual unitary inner forms; no archive Host or NTypes statement is an input.

Weak transfer alone would not identify the infinity parameter. We make that extra step explicitly. Choose the auxiliary quasi-split unitary group U\_aux of sufficiently large rank M in Chen--Zou §3.1, with the opposite parity and Witt index exceeding dim U. Use honest global splitting characters chi\_V,chi\_W, available for example as suitable powers of(AP2). The auxiliary lift is used here only in the proof of a classification statement; it is not an additional edge of the proposed arithmetic tower and contributes no asserted lift-cost saving.

By Chen--Zou Theorem3.1, the abstract large theta lift Sigma=theta\^{}abs(pi) occurs in the discrete automorphic spectrum of U\_aux. Their Proposition4.1, proof and equation(4.2), identify its global Arthur parameter as

\[
 \Theta(\Psi)=\Psi\chi_W^{-1}\chi_V
                      \boxplus\chi_V\boxtimes S_{M-N},
 \qquad N=\dim U.
 \tag{AP6b}
\]

At the chosen real place, let phi\_v be the actual tempered local Langlands parameter of pi\_v. Apply the \textbf{local} Chen--Zou Theorem6.3, with their displayed equation(6.1), to phi\_v. It says that the same nonzero local theta lift Sigma\_v belongs to the quasi-split Arthur packet associated with

\[
 \Theta(\phi_v)=\phi_v\chi_W^{-1}\chi_V
                      \oplus\chi_V\boxtimes S_{M-N}.
 \tag{AP6c}
\]

The equality of the two local lifts uses Howe duality and the bijection in Theorem6.3, not an assumed equality between a global weak transfer and an arbitrary local parameter.

All members of a real quasi-split Arthur packet have the infinitesimal character of its associated Langlands parameter. The exact primary construction is Mok Proposition8.2.1, printed231--232: every standard character used there has that same infinitesimal character. Arancibia Robert--Mezo(2022), equations(114) and Theorems8.2/9.3, pp42/44, also give this statement directly by placing the entire packet in the fixed infinitesimal-character category, including singular parameters. For the possibly unbounded localizations of global parameters, Chen--Zou §3.4, printed4784--4785, uses normalized parabolic induction; its constituent infinitesimal character is still unchanged.

Consequently(AP6b) and(AP6c) have the same infinitesimal-character multiset. Both apply the \textbf{same} scalar character shift chi\_W\^{}(-1)chi\_V and append the \textbf{same} multiset from chi\_V times S\_(M−N). Cancel that common finite multiset, with multiplicities, and undo the common scalar shift. This proves

\[
                     \operatorname{IC}(\Psi_v)
                         =\operatorname{IC}(\pi_v).
 \tag{AP6d}
\]

There is no guessed vacuum constant, contragredient, or missing half-determinant shift in this comparison: the identical two occurrences of the published formula(6.1) fix them all. The comparison can also be obtained from the type-A center map of Przebinda, but that separate normalization is unnecessary here.

Now restrict any Pi\_i,v to the complex Weil group, where it is a sum of characters. If one of its holomorphic infinitesimal exponents is z, then the contribution of Pi\_i,v times S\_d to the associated parameter of Psi\_v contains

\[
                  z+(d-1)/2,\ z+(d-3)/2,\ldots,
                                  z-(d-1)/2.
 \tag{AP6e}
\]

This is Chen--Zou equation(2.4), with \textbar z\textbar\_C=z bar z. For d\textgreater1 the displayed multiset contains two entries differing by one. No cancellation is possible in a direct-sum parameter. Equation(AP6d) contradicts the hypothesis. Therefore every d\_i=1. QED.

Both actual groups in(AP1) satisfy the hypotheses: the compact real representations of H are tempered, as are the prescribed genuine discrete series of G. Equation(AP5) therefore gives

\[
        \Psi_H=\boxplus_{i=1}^{r_H}\Pi_{H,i},\qquad
        \Psi_G=\boxplus_{i=1}^{r_G}\Pi_{G,i}.
 \tag{AP7}
\]

Each sum is elliptic, so its cuspidal summands are pairwise distinct. Chen--Zou Theorem2.6 now applies: the actual representation belongs to the corresponding \textbf{local Langlands packet at every place}, and its automorphic multiplicity is one. This all-place conclusion is taken only after genericity has been proved. The full rank-one Theorem7.7 is consistent with the argument but is not needed to substitute for this step.

\subsection{4. Finite temperedness of the actual initial and bottom constituents}\label{finite-temperedness-of-the-actual-initial-and-bottom-constituents}

Generic global parameters alone are not a proof of the Ramanujan conjecture. In this particular family the actual algebraic infinity grids allow the known conjugate-selfdual theorem to be applied, with its rank parity checked.

Let Pi be one cuspidal block of rank d in a generic parameter for a group of rank N. By(AP7) and the all-place local-packet statement, its complex infinity characters are a subset of the bounded characters of the actual pi\_tau. Write their exponents x\_1\textgreater\ldots\textgreater x\_d.~They lie in Z+(N−1)/2, are distinct, and have adjacent gaps greater than one.

Choose epsilon in\{0,1\} with epsilon congruent to d−N modulo2, and put

\[
                         \widetilde\Pi=\Pi\otimes\eta^\epsilon.
 \tag{AP8}
\]

At a selected complex embedding its exponents become x\_j+epsilon/2, which lie in Z+(d−1)/2. Hence

\[
        \mu_j=x_j+\epsilon/2-\bigl((d-1)/2-(j-1)\bigr)
                           \in\mathbb Z
 \tag{AP9}
\]

are dominant integral highest-weight coordinates. At the conjugate embedding the parameters are the corresponding negatives, with their order reversed; they give the conjugate coefficient weight. Thus tilde Pi is regular C-algebraic. It remains cuspidal and unitary, and is conjugate-selfdual by(AP2).

For cuspidal GL representations, regular C-algebraicity implies cohomological infinity type: Shin--Templier, Lemma2.12 and Remark2.13, pp7--8 of the author version, state the exact implication and distinguish it from the false assertion for arbitrary local representations. They cite Clozel's Lemma3.14. We use the published implication on the cuspidal representation tilde Pi, not merely integrality of a local infinitesimal character.

Caraiani(2012), Theorem1.2, printed2312, therefore implies that tilde Pi is tempered at every finite place. For d=1 the same conclusion follows directly because it is a unitary Hecke character. Since eta is unitary everywhere, untwisting(AP8) preserves temperedness. Thus every block Pi is tempered at every finite place. Its infinity parameters were already bounded. Their direct sum is bounded at every place, so the actual local packets in(AP7) are tempered. We have proved

\[
       \boxed{\pi_H\text{ and every existing }\pi_G
                       \text{ are tempered at every place.}}
 \tag{AP10}
\]

The initial representation is not assumed Steinberg, globally simple, or already known tempered from a first theta lift. This proves temperedness for \textbf{all} actual initial Mram constituents. It does not assert that the later large compact theta descendants are tempered: their inserted strings are different and are not covered by(AP5).

The assertion has concrete initial content. H has compact automorphic quotient, the actual strict initial stabilizers are trivial by K042, and the exact diamond action gives

\[
 \dim M^{\rm ram}/\dim M
                   =((T-2)/(T-1))^{[F:\mathbb Q]}>0.
 \tag{AP11}
\]

This is the literal unweighted-norm calculation in DO044 and the independent LIFT review of that calculation. The coefficient dimension and the finite class mass are positive. Hence the actual initial Mram space contains an irreducible constituent to which(AP10) applies. No bottom cusp-counting assertion is inferred from(AP11).

\subsection{5. The literal goodTuple degree condition forces simple bottom base change}\label{the-literal-goodtuple-degree-condition-forces-simple-bottom-base-change}

We need the following elementary local observation. It avoids an unproved transfer from a finite regular series to a depth-zero inertial parameter.

\subsubsection{Lemma AP12: a global block degree survives finite parabolic restriction}\label{lemma-ap12-a-global-block-degree-survives-finite-parabolic-restriction}

Let k be a nonarchimedean field, K\_N=GL\_N(O\_k), K\_N\^{}1 its first principal subgroup, and let tau\_i be smooth representations of GL\_(d\_i)(k), with sum d\_i=N. For normalized parabolic induction I=Ind\_P\^{}(GL\_N)(tensor tau\_i), there is an exact finite-group identification

\[
 I^{K_N^1}\simeq
 \operatorname{Ind}_{\bar P}^{GL_N(\kappa)}
                     \bigotimes_i\tau_i^{K_{d_i}^1}.
 \tag{AP12a}
\]

If an irreducible finite Green-series constituent with semisimple polynomial f occurs on the left, then f has a monic divisor of degree d\_i for every i.

\textbf{Proof.} Iwasawa decomposition realizes the induced functions on K\_N, with covariance under P intersect K\_N. Right K\_N\^{}1 invariance makes their values invariant under each K\_(d\_i)\^{}1. The unipotent radical acts trivially on the inducing representation, and the modulus is one on the compact group. Reduction gives(AP12a). Compact invariants are exact in characteristic zero by normalized averaging, so the same constituent implication holds for any irreducible subquotient of I.

The finite group modules on the right are semisimple. Choose irreducible constituents rho\_i of the factors. The Green characteristic ring is the tensor product of symmetric-function rings indexed by Frobenius character orbits; finite Harish--Chandra induction is multiplication in this ring. Use the same compatible character-orbit/root-orbit identification as the genuine eSer realization in R026. Its grading records, for each orbit, the total partition size. Therefore the semisimple polynomial of every constituent of the induction is the product of the semisimple polynomials of the rho\_i, each of degree d\_i. This is the product defined in Gow--Vinroot(2007), §3, pp8--9, and its character-ring isomorphism Theorem3.2. It proves the divisor assertion, regardless of the number of finite constituents in each tau\_i\^{}(K\_i\^{}1). QED.

Now let pi\_G have a nonzero vector in the actual e0 bottom space. There is an admitted tuple f whose tensor of genuine finite Green-series projectors acts nontrivially. This is the literal finite sum defining e0, realized as in R025/R026. At either actual rational G prime g, write f\_v for its degree-n polynomial at each selected place v over g.

Suppose the generic parameter Psi\_G from(AP7) has more than one global cuspidal block. Some block has rank d with1\textless=d and2d\textless=n.~At every selected split G place v, the actual local representation pi\_G,v has Langlands parameter equal to the direct sum of those block parameters. Choosing the conjugate split embedding replaces the blocks by their duals and does not change any block rank or divisor-degree conclusion. It is the corresponding irreducible tempered parabolic induction, or equivalently an irreducible subquotient of the induced product; the latter weaker statement suffices. Lemma AP12 shows that every finite polynomial occurring in its K\_n\^{}1 invariants has a divisor of degree d.

But the fifth displayed conjunct of the \textbf{actual} predicate \nolinkurl{IsGood} is

\[
 \forall d\ge1,\ 2d\le n\Longrightarrow
      \exists v:\ \text{no divisor of }f_v\text{ has degree }d.
 \tag{AP12b}
\]

Its exact locator is \nolinkurl{Formalism/CyclotomicNorms/StatementT.lean}19--29, and \nolinkurl{Formalism/SharpRestriction/LevelLetterCocycles/DataN.lean}996--1007 supplies the actual goodTuple/e0 connection. The place supplied by(AP12b) contradicts the nonzero local projector. Hence r\_G=1, and

\[
 \boxed{\operatorname{BC}(\pi_G)=\Pi_G
       \text{ is conjugate-selfdual cuspidal on }GL_{a+1}(\mathbb A_E).}
 \tag{AP13}
\]

This proof applies to the entire literal e0, not only our explicitly chosen irreducible-polynomial subblock. If the right-action convention replaces the finite series by its contragredient, R026's proved reciprocal closure of full IsGood gives the same conclusion.

The squarefree conjunct is present and was read literally. It makes each admitted semisimple finite series regular; R026 supplies its genuine unique generic constituent. \textbf{Squarefreeness by itself is not used as an uncited theorem about a p-adic inertial parameter.} The Arthur SL2 was eliminated at infinity in Lemma AP6. Once genericity is known, the exact finite parabolic identity and the literal degree-divisor clause already prove simplicity. Thus no unexplained depth-zero type dictionary, or unproved NTypes consequence, remains in this argument.

\subsection{6. Exact GGP normalization for an actual pair}\label{exact-ggp-normalization-for-an-actual-pair}

The groups of the GGP theorem are U(h) inside U(h plus the norm-one positive line). Our displayed pair is put in that form by taking h=−I\_a: U(−I\_a)=U(I\_a), and U(−I\_a,1)=U(I\_a,−1), with the same block embedding. Simultaneously multiplying the forms by−1 changes neither algebraic group nor this embedding. There is no Weil representation or additive-character rescaling in this identification.

Let pi\_H and pi\_G be actual constituents as above, and set

\[
 \Pi_H=\boxplus_{i=1}^{r_H}\Pi_{H,i},\qquad
 \Pi=\Pi_H\boxtimes\Pi_G.
 \tag{AP14}
\]

These are Hermitian Arthur parameters in the exact sense of Beuzart-Plessis--Chaudouard--Zydor(2022), §1.1.3: generic elliptic, with the prescribed conjugate-selfdual Asai parity. The component-group order in their convention is

\[
                         |S_\Pi|=2^{r_H+1}.
 \tag{AP15}
\]

It is four only if the initial base change is also simple; that extra assertion has not been made.

Let tau be the actual finite T translation of LIFT29/R027 and let f tensor F be a nonzero factorizable vector of pi\_H times pi\_G. Put

\[
 \mathcal P_\tau(f,F)=\int_{[H]} f(h)F(\iota(h)\tau)\,dh.
 \tag{AP16}
\]

This is the source-style bilinear period. The representation of F is the actual contragredient bottom representation when that is the source convention. No additional dual is to be inserted into(AP14). Right translation by tau preserves its Petersson norm and is retained in every local test vector.

For Tamagawa measures, define

\[
 \mathcal L(s,\pi_H\boxtimes\pi_G)=
 \left(\prod_{j=1}^{a+1}L(s+j-\tfrac12,\omega_{E/F}^{\,j})\right)
 \frac{L(s,\Pi_H\times\Pi_G)}
      {L(s+\tfrac12,\pi_H\boxtimes\pi_G,\operatorname{Ad})}.
 \tag{AP17}
\]

At each place, with phi\_v=f\_v tensor pi\_G,v(tau\_v)F\_v, put

\[
 \mathcal P_v^\natural(\phi_v,\phi_v)=
 \mathcal L_v(\tfrac12,\pi_H\boxtimes\pi_G)^{-1}
 \int_{H(F_v)}
    \langle(\pi_H(h)\otimes\pi_G(\iota(h)))\phi_v,
                                          \phi_v\rangle\,dh_v.
 \tag{AP18}
\]

Temperedness(AP10) makes these integrals absolutely convergent. The local normalizing factor is finite and nonzero at1/2. Then BPCZ Theorem1.1.6.1, printed187, gives the exact equality

\[
 \boxed{
 \frac{|\mathcal P_\tau(f,F)|^2}
      {\|f\|_{\rm Tam}^2\|F\|_{\rm Tam}^2}
   =2^{-(r_H+1)}\mathcal L(\tfrac12,\pi_H\boxtimes\pi_G)
     \prod_v\frac{\mathcal P_v^\natural(\phi_v,\phi_v)}
                         {\|\phi_v\|_v^2}.}
 \tag{AP19}
\]

The product is the convergent, measure-compatible product of that theorem; at almost all places its factor is vol(H(O\_v)). Their Tamagawa convention includes global normalizing L-values, explicitly warned in the theorem's footnote. The identity uses complete L-functions, including infinity factors. It is not a claim that an arbitrarily truncated L-value controls the fixed source period.

For other fixed Haar measures mu\_H=c\_H times Tamagawa and mu\_G=c\_G times Tamagawa, the corresponding normalized left side is c\_H/c\_G times(AP19). Indeed the period is multiplied by c\_H, while the two Petersson norms together are multiplied by c\_H c\_G. This records the measure conversion without hiding it in an unspecified period scalar. If the initial and bottom Haar choices are the canonical maximal-finite/compact-real and ball measures used in R030, their already computed quotient volumes can be substituted for c\_H and c\_G after the common Tamagawa normalization is fixed. No new numerical Tamagawa constant is assumed here.

Equation(AP19) extends to finite sums of pure tensors by polarization. That extension does not remove cancellations in a prescribed nonfactorizable arithmetic row.

\subsection{7. What the criterion proves, and the precise remaining gate}\label{what-the-criterion-proves-and-the-precise-remaining-gate}

The new content is the verification of the actual hypotheses for(AP19): generic global parameters, all-place LLC compatibility, finite temperedness with the half-parity repair, and simplicity of the actual e0 bottom parameter. It is stronger than merely restating GGP under unidentified parameter assumptions.

It does \textbf{not} prove any of the following separate statements:

\begin{enumerate}
\def\labelenumi{\arabic{enumi}.}
\tightlist
\item
  That the genuine e0 bottom cusp space is nonzero at one of the bounded current weights. R027 supplies a qualitative detector after sufficiently large canonical twist; DO044 excludes full-packet surjectivity over an upper portion of the allowed interval. Neither statement supplies this bounded existence assertion.
\item
  That an existing actual pair satisfies L(1/2,Pi\_H times Pi\_G) nonzero. The genericity and local-temperedness theorem is not a central-value theorem. BPCZ Theorem1.1.5.3 constructs some simple rank-a parameter for a given bottom parameter; it does not preserve our initial weight and G/T levels.
\item
  That every prescribed local test in(AP18) is nonzero. The actual real compact branching occurs and its PBW metric is known (K043); the finite generic restriction in R028/LIFT32 is also known. A finite residue branching statement alone is not a proof about the complete p-adic GGP functional of an arbitrary actual pair. In particular the actual maximal-T bottom vector and its fixed nonintegral tau translation remain in(AP18).
\item
  That these periods detect the prescribed final brC/tilQ image or give its value in the fixed integral terminal detector lattice. DO038 supplies exact class/branch transport, not the missing nonzero restricted value.
\end{enumerate}

For a \textbf{chosen actual factorizable pair}, the exact remaining criterion is nonvanishing of the right side of(AP19), with every actual local vector and normalizing factor retained. The global existence theorem BPCZ1.1.5.1 states nonvanishing for some member of the appropriate pure-inner-form packet when the complete central L-value is nonzero. It cannot replace the fixed-inner-form, fixed-level and fixed-vector checks in the source task.

Thus the global L-value obstruction is now a legitimate, precisely normalized remaining obstruction for the actual family. It is not being used to declare the period zero, or to turn the failure of the direct theta shortcut in DO045 into a branching-period no-go.

\subsection{8. Primary source and source-code record}\label{primary-source-and-source-code-record}

The following primary passages were read directly. Local copies are retained under \nolinkurl{definite_oscillator_ggp_refs}.

\begin{itemize}
\tightlist
\item
  Chen, R., \& Zou, J. (2025). \emph{Arthur's multiplicity formula for even orthogonal and unitary groups}. Journal of the European Mathematical Society, 27, 4769--4843. \url{https://doi.org/10.4171/JEMS/1564}. Theorem2.1 p4777; Theorem2.6 p4780; Theorem3.1 p4783; Proposition4.1 and its proof pp4788--4790; §3.4 pp4784--4785; equation(6.1) and Theorem6.3 p4796; Theorem7.7 and Remark7.8 p4823. Official PDF: \url{https://ems.press/content/serial-article-files/51452}.
\item
  Mok, C. P. (2015). \emph{Endoscopic classification of representations of quasi-split unitary groups}. Memoirs of the American Mathematical Society, 235(1108). \url{https://doi.org/10.1090/memo/1108}. Directly read the author version arXiv:1206.0882: §7.1, printed190, and Proposition8.2.1, printed231--232. These establish the exponent convention and constant infinitesimal-character construction used above.
\item
  Arancibia Robert, N., \& Mezo, P. (2022). \emph{Equivalent definitions of Arthur packets for real unitary groups}. arXiv:2204.10715v2. \url{https://arxiv.org/abs/2204.10715}. Exact read locators are equation(114) and Theorem8.2 p42; Theorem9.3 p44. The later journal publication is Representation Theory 29 (2025), 968--1027, DOI10.1090/ert/707; the locators here refer explicitly to the retained preprint, not to unverified journal pagination.
\item
  Caraiani, A. (2012). \emph{Local-global compatibility and the action of monodromy on nearby cycles}. Duke Mathematical Journal, 161(12), 2311--2413. \url{https://doi.org/10.1215/00127094-1723706}. Theorem1.2, p2312, is the exact finite-temperedness result. Author PDF: \url{https://www.ma.imperial.ac.uk/~acaraian/papers/lgc1.pdf}.
\item
  Shin, S. W., \& Templier, N. (2014 author version). \emph{On fields of rationality for automorphic representations}. Lemma2.12 and Remark2.13, pp7--8. \url{https://math.berkeley.edu/~swshin/FieldRat.pdf}. These state the cuspidal regular-C-algebraic to cohomological implication, with its local-temperedness qualification, citing Clozel's Lemma3.14. The original Clozel volume was located but not successfully downloaded in this check; no claim to have read that original lemma is made.
\item
  Gow, R., \& Vinroot, C. R. (2007). \emph{Extending real-valued characters of finite general linear and unitary groups on elements related to regular unipotents}. arXiv:0704.2022. The characteristic-ring multiplication and Theorem3.2, pp8--9. \url{https://arxiv.org/abs/0704.2022}. These are the exact finite Green-product facts used in Lemma AP12.
\item
  Beuzart-Plessis, R., Chaudouard, P.-H., \& Zydor, M. (2022). \emph{The global Gan--Gross--Prasad conjecture for unitary groups: The endoscopic case}. Publications Mathématiques de l'IHÉS, 135, 183--336. \url{https://doi.org/10.1007/s10240-021-00129-1}. Definitions1.1.2--1.1.4 pp184--185; Theorems1.1.5.1/1.1.5.3 pp185--186; all normalizations and Theorem1.1.6.1 pp186--187. Primary PDF: \url{https://www.numdam.org/item/10.1007/s10240-021-00129-1.pdf}.
\end{itemize}

Literal code locators: \nolinkurl{Formalism/CyclotomicNorms/StatementT.lean}19--29 and \nolinkurl{Formalism/SharpRestriction/LevelLetterCocycles/DataN.lean}996--1007. Actual finite action/projector realization, reciprocal closure and concrete admitted blocks are the already checked R025--R029 and DO042--043. The actual weight formulas are K044(EG1--EG5), and the genuine odd global character is K029/K032 and K040(NC3). No Lean build is made.

\subsection{9. Verification scope}\label{verification-scope}

The companion script and JSON are \nolinkurl{definite_oscillator_actual_arthur_parameters_certificate.py} and \nolinkurl{.json}. The exact replay completed46,455 assertions at2026-11-22T04:52:18.683592+00:00. It checks every adjacent coordinate-bit possibility at the affine parameter endpoints, all actual rank parity repairs, integral dominant coefficient weights in every possible block rank, the unit-step Arthur strings, exact common-multiset cancellation and finite polynomial-degree bookkeeping.

The certificate does not compute an automorphic parameter, a bottom cusp dimension, a global L-value, a local GGP integral, or a terminal period. Those logical steps are the written proofs above with the primary inputs specified. The whole-proof author verification is recorded separately after the final complete reading.

\subsubsection{Completed author whole-proof check}\label{completed-author-whole-proof-check}

At 2026-11-22T04:57:16.284698+00:00, the author reread the complete report(AP1--AP19), including the full proof of Lemma AP6, the actual odd-character parity repair, the exact compact-invariant parabolic formula, the complete genuine e0 argument and the Tamagawa GGP normalization. The reading checked the actual source predicates and the primary theorem passages listed above. The auxiliary large theta lift is used only to identify infinitesimal characters; genericity precedes the all-place packet application. All existing-initial versus conditional-existing-bottom, whole-space versus fixed-vector, and central-L versus local-test distinctions remain explicit. \textbf{PASS in this stated scope.}

\clearpage

\section{ALP-LIFT-36: an explicit boundary L-value cost for the actual vertical tower}\label{alp-lift-36-an-explicit-boundary-l-value-cost-for-the-actual-vertical-tower}

Author: /root/lift\_accounting, GPT-6 Astra Pro, ALP\_MACHINE. 22 November 2026 UTC. Inherited reasoning effort is not separately exposed.

\textbf{Author whole-proof review completed: 2026-11-22T05:00:46Z.}

This result supplies an effective lower bound for the actual finite Hecke factor at the boundary of the second vertical Rallis formula. It also bounds the complete inserted Hecke string and the absolutely convergent initial-parameter factors. It does not identify the local quotient norms of LIFT31 with globally normalized Rallis forms, or remove the source's real-intertwiner and Haar scalars. Thus it is one explicit component of the top-preimage cost, not yet a lower singular value for the entire map.

\subsection{1. Analytic conventions and the explicit general bound}\label{analytic-conventions-and-the-explicit-general-bound}

Let E be totally imaginary of degree n=2f, with positive absolute discriminant D. Let chi be a unitary Hecke character whose local archimedean components have no radial part:

\[
 \chi_\tau(z)=(z/|z|)^{m_\tau},\qquad
 a_\tau=|m_\tau|/2\ge0.                              \tag{LB1}
\]

Assume that at least one m\_tau is nonzero. Both chi and chi squared are then different from every pure norm power. Their finite L-functions are entire. Throughout this report L\_E(s,chi) means the finite Euler product, with local factor1 at a ramified prime, analytically continued. Put F=N(f\_chi), A=DF, and

\[
 \begin{split}
 A_0&=2^nD^{3/2}5^{3f+1},\\
 B_1&=2^nA^{3/2}\prod_\tau(a_\tau+6)^3,\\
 B_2&=2^nA^{3/2}\prod_\tau(2a_\tau+6)^3.
 \end{split}                                           \tag{LB2}
\]

The conductor of chi squared divides the conductor of chi, so the same A is a valid upper bound for its finite conductor-discriminant factor.

\textbf{Theorem.} With these conventions,

\[
 \boxed{\quad
 |L_E(1,\chi)|\ge
       \frac1{1024 A_0^3B_2B_1^3}>0.
 \quad}                                                \tag{LB3}
\]

All constants are displayed. In particular there is no ineffective constant or ineffective exceptional-character alternative.

The analytic inputs are Tate's continuation, functional equation and split Mellin-integral argument, not a quantitative zero-free-region theorem. In Tate (1950), section2.5, printed pages2.14--2.16 (PDF pages 21--23), the angular character exp(im theta) has gamma shift \textbar m\textbar/2. Theorem4.4.1 and its proof give the two convergent Mellin tails; section4.5 identifies the finite Euler product and the conductor factors. The original source is saved as lift\_refs/tate1950.pdf and .txt; the gamma formula was read directly in lift\_refs/tate1950\_complex\_gamma2.png. The local-global normalization is also recorded in Kudla (2001), Proposition3.5 and Theorem4.3/Corollary4.4, saved as lift\_refs/kudla\_tate.pdf and .txt. The remainder of the argument is proved below.

\subsection{2. A uniform bound in a vertical strip}\label{a-uniform-bound-in-a-vertical-strip}

Up to an s-independent nonzero constant, the completed function is

\[
 \Lambda(s,\chi)=A^{s/2}(2\pi)^{-fs}
             \prod_\tau\Gamma(s+a_\tau)L_E(s,\chi).
 \tag{LB4}
\]

It satisfies Lambda(s,chi)=epsilon Lambda(1-s,chi\_bar), with \textbar epsilon\textbar=1. Conjugation changes m\_tau to minus m\_tau and leaves all a\_tau unchanged. The omitted constants such as (2pi) to powers of a\_tau occur on both sides and cancel.

For Re(s)\textgreater=2, the Euler product and unitarity give

\[
 |L_E(s,\chi)|\le\zeta_E(2)
                 \le\zeta_{\mathbf Q}(2)^n<2^n.
 \tag{LB5}
\]

The second inequality follows prime by prime: a rational prime has at most n prime ideals above it and each residue norm is at least that prime. On s=-1+iy, (LB4) and the gamma recurrence therefore give

\[
 \begin{split}
 |L_E(-1+iy,\chi)|
 &\le2^nA^{3/2}(2\pi)^{-3f}
  \prod_\tau|(a_\tau-1-iy)(a_\tau-iy)(a_\tau+1-iy)|\\
 &\le2^nA^{3/2}\prod_\tau|a_\tau+2+iy|^3.
 \end{split}                                           \tag{LB6}
\]

At a gamma pole the identity is understood by continuity; its right side is still a polynomial. For a\textgreater=0 each of \textbar a-1\textbar, a, a+1 is at most a+2, proving the last inequality.

Consider the holomorphic function on the closed strip -1\textless=Re(s)\textless=3

\[
       F_\chi(s)=\frac{L_E(s,\chi)}
                         {\prod_\tau(s+a_\tau+3)^3}.
 \tag{LB7}
\]

The zeros of its denominator lie to the left of the strip. Equations (LB5) and (LB6) bound it on both vertical edges by 2\^{}n A\^{}(3/2). For completeness the growth needed for the strip maximum principle does not require an unspecified convexity estimate. Choose Tate's polynomial Gaussians at infinity and the usual conductor tests at the finite places. After splitting the global zeta integral at idele norm1, the two Mellin tails are bounded on each closed vertical strip: the imaginary power has modulus1, and rapid decay gives a common integrable majorant. The finite local factors used here are constants times positive real numbers to the power s. Dividing the resulting gamma product and using Stirling gives, for this fixed character, \textbar F\_chi(s)\textbar=O(exp(C\textbar Im(s)\textbar{} log(2+\textbar Im(s)\textbar))) on the strip. No value of C will enter a bound below.

Apply the maximum principle on rectangles to F\_chi(s)exp(epsilon s\^{}2). The horizontal-edge bound tends to zero as their height tends to infinity. On the two vertical edges the extra modulus is at most exp(9epsilon). Let the height tend to infinity and then epsilon to0. We obtain the explicit bound

\[
 |L_E(s,\chi)|\le2^nA^{3/2}
                  \prod_\tau|s+a_\tau+3|^3
             \quad(-1\le\Re(s)\le3).                 \tag{LB8}
\]

For each real x in {[}1,2{]}, the closed disc \textbar s-x\textbar\textless=1/4 is inside this strip and \textbar s+a+3\textbar\textless a+6. Cauchy's derivative bound yields

\[
 \sup_{1\le x\le2}|L_E'(x,\chi)|\le4B_1=:M.
 \tag{LB9}
\]

Applying (LB8) to chi squared, and using its conductor divisibility, also gives

\[
          |L_E(x,\chi^2)|\le B_2\qquad(1\le x\le2).
 \tag{LB10}
\]

\subsection{3. An explicit simple-pole bound for the Dedekind zeta function}\label{an-explicit-simple-pole-bound-for-the-dedekind-zeta-function}

It is important to keep the simple pole, rather than to use zeta\_Q(x)\^{}n near1. Apply the same argument to

\[
             F_\zeta(s)=
             \frac{(s-1)\zeta_E(s)}{(s+3)^{3f+1}}.
 \tag{LB11}
\]

This is holomorphic on the strip. On Re(s)=-1 the gamma ratio has the three factors (-1-iy),(-iy),(1-iy) at each complex place. Each has modulus at most \textbar2+iy\textbar. The additional \textbar s-1\textbar{} is exactly \textbar2+iy\textbar{} and is cancelled by the last denominator factor. Hence \textbar F\_zeta(-1+iy)\textbar\textless=2\^{}nD\^{}(3/2). On Re(s)=3, the Euler bound applies, \textbar s-1\textbar\textless=\textbar s+3\textbar, and the remaining denominator has modulus at least1. The same damped-rectangle argument gives

\[
       (x-1)\zeta_E(x)\le
       2^nD^{3/2}(x+3)^{3f+1}\le A_0
                      \qquad(1<x\le2).               \tag{LB12}
\]

All zeros and removable gamma singularities are harmless in this argument; the only pole of zeta\_E was removed in (LB11).

\subsection{4. Positivity and optimization at the boundary}\label{positivity-and-optimization-at-the-boundary}

For real x\textgreater1, expansion of the absolutely convergent Euler logarithms proves

\[
       \zeta_E(x)^3|L_E(x,\chi)|^4|L_E(x,\chi^2)|\ge1.
 \tag{LB13}
\]

Indeed, at a prime unramified for chi and each positive Euler exponent, the coefficient is

\[
           3+4\cos\theta+\cos(2\theta)
                    =2(1+\cos\theta)^2\ge0.           \tag{LB14}
\]

If chi is ramified, its Euler factor is1. If chi squared is also ramified the remaining coefficient is3. If chi squared becomes unramified, the remaining coefficient has the form3+cos(theta), which is at least2. Thus (LB13) applies to the actual primitive L-functions with their possibly different conductors.

For 0\textless epsilon\textless=1, (LB10), (LB12) and (LB13) imply

\[
 |L_E(1+\epsilon,\chi)|\ge c\epsilon^{3/4},
       \qquad c=A_0^{-3/4}B_2^{-1/4}.                 \tag{LB15}
\]

All constants in (LB2) are at least1, so c\textless=1 and M\textgreater=1. Set

\[
                      \epsilon=(c/(2M))^4\le1.
 \tag{LB16}
\]

By integrating (LB9) along the real interval,

\[
 \begin{split}
 |L_E(1,\chi)|
 &\ge c\epsilon^{3/4}-M\epsilon
   =\frac{c^4}{16M^3}\\
 &=\frac1{1024A_0^3B_2B_1^3}.
 \end{split}                                           \tag{LB17}
\]

This proves (LB3). The argument also proves nonvanishing; it has not assumed the qualitative boundary-value statement of LIFT14 or LIFT31.

\subsection{5. Application to the actual coherent stage characters}\label{application-to-the-actual-coherent-stage-characters}

Use either K031's J14 host or K032's J56 host, with the coherent nonuniform characters of K040. Their uniform stage difference is

\[
 \delta=e_{1,\tau}-e_{0,\tau}=2\ell+q,
       \qquad0\le q\le\ell,\quad\ell\le\ell_+<2^{485}.
 \tag{LB18}
\]

Here q is the canonical weight increment in K044; it is not a residue cardinality. The homomorphism identity K040(NC15), specialized as in NC16/NC20, gives the exact global character

\[
 \vartheta=\mu_1/\mu_0
   =(\widetilde\nu_1\widetilde\nu_0^{-1})\eta^{-2\delta}.
 \tag{LB19}
\]

The odd character eta is unitary and finite-unramified, with angular exponent1 at every selected complex embedding. Thus all a\_tau in (LB1) are exactly delta. The tuple character cancels from the ratio. The finite ratio in (LB19) is tame at the two completely split G primes and unramified elsewhere. Therefore

\[
     N\mathfrak f_\vartheta\le Q^n,
         \qquad Q=g_0g_1.                             \tag{LB20}
\]

The existing nontrivial successive residue characters make equality hold, but the upper bound is all that is used. These facts follow from K031(HC18--20), K032(QH30--32), and K040(NC15--20). They do not require a bound for the coefficient field or for a relative class group. Nonzero delta ensures that both vartheta and its square are outside the pure-norm exception in Tate's theorem.

Substitution into (LB3) gives the explicit expression

\[
 |L_E(1,\vartheta)|\ge
 \left[
 2^{7n+10}D^{21/2}(N\mathfrak f_\vartheta)^6
 5^{9f+3}(\delta+6)^{9f}(2\delta+6)^{3f}
 \right]^{-1}.                                        \tag{LB21}
\]

All archimedean gamma shifts are included here, with their actual product over f complex places. In particular the bound is not the one-dimensional imaginary-quadratic conductor bound reused f times.

\subsection{6. Fully numerical bounds for the two actual hosts}\label{fully-numerical-bounds-for-the-two-actual-hosts}

The already verified discriminant formulas are

\[
 \operatorname{rd}(E_{14})=2\,7^{5/6}29^{6/7}<2^8,
 \quad
 \operatorname{rd}(E_{56})=2\,7^{5/6}29^{6/7}\sqrt{65}<2^{11}.
 \tag{LB22}
\]

These are K044(EG13), with the sharper integer inequalities used in DO044 and K048. They can also be checked by taking the42nd or84th power, so no floating-point logarithm is needed. The actual G data are

\begin{longtable}[]{@{}lrrrlr@{}}
\toprule\noalign{}
Host & n & f & r with D\textless2\^{}(rn) & (g0,g1) & qG with Q\textless2\^{}qG \\
\midrule\noalign{}
\endhead
\bottomrule\noalign{}
\endlastfoot
J14 & 84 & 42 & 8 & (4615409,5276377) & 45 \\
J56 & 336 & 168 & 11 & (22285932761,64072175441) & 71 \\
\end{longtable}

Since delta\textless3 times2\^{}485, we have delta+6\textless2\^{}487 and 2delta+6\textless2\^{}488. Also5\textless2\^{}3. The base2 logarithm of the denominator in (LB21) is strictly less than

\[
       (5888+21r+12q_G)f+19.
 \tag{LB23}
\]

Indeed the contributions are respectively14f+10,21rf,12qG f, 27f+9,4383f,1464f. Hence, uniformly over all actual K044 profiles, their binary words, their role-dependent centres, and their permitted canonical increments,

\[
 \boxed{
 \begin{array}{c|c}
 \text{host}& |L_E(1,\vartheta)|\\\hline
 J14&>2^{-277051}\\
 J56&>2^{-1171147}.
 \end{array}}
 \tag{LB24}
\]

No restriction on the number of top labels is used in this result. The simultaneous inverse convention replaces vartheta by its complex conjugate and leaves every displayed absolute bound unchanged.

\subsection{7. The whole inserted string and removed Euler factors}\label{the-whole-inserted-string-and-removed-euler-factors}

Let S be any finite set of prime ideals of E and write L\_E\^{}S for the Euler product with those local factors omitted. Define the explicit positive number

\[
 R_S(\vartheta)=
   \prod_{\substack{\mathfrak p\in S\\
                  \chi\text{ unramified at }\mathfrak p}}
                (1-(N\mathfrak p)^{-1}),
       \qquad\chi=\vartheta.
 \tag{LB25}
\]

Then \textbar L\_E\^{}S(1,vartheta)\textbar\textgreater=R\_S(vartheta)\textbar L\_E(1,vartheta)\textbar, simply by multiplying the omitted factors. This retains every exceptional finite place. A ramified prime contributes1, as it must.

For every positive integer t\textgreater=2,

\[
             \prod_{i=2}^{t}|L_E^S(i,\vartheta)|>e^{-n}.
 \tag{LB26}
\]

To prove it, each reciprocal Euler factor has modulus at least 1/(1+(Np)\^{}(-i)). Since log(1+x)\textless=x, the negative logarithm of the asserted lower bound is at most

\[
 \sum_{\mathfrak p}\sum_{i=2}^\infty(N\mathfrak p)^{-i}
 \le n\sum_{p\text{ rational prime}}\frac1{p(p-1)}
 < n\sum_{m=2}^\infty\frac1{m(m-1)}=n.
 \tag{LB27}
\]

Omitting primes only removes nonnegative terms from this majorant. In particular there is no loss proportional to t from multiplying the entire long Hecke string. If C is the exponent in (LB24), then

\[
       \prod_{i=1}^{t}|L_E^S(i,\vartheta)|
                  >2^{-C}e^{-n}R_S(\vartheta).
 \tag{LB28}
\]

\subsection{8. The initial-parameter factors in the actual vertical formula}\label{the-initial-parameter-factors-in-the-actual-vertical-formula}

LIFT31(VA18,VA30) identifies the second standard factor as an initial rank-r Euler product times the string in (LB28). At the first and second evaluations the initial Euler products both have the verified exponent margin h=859 for J14 or h=733 for J56. This uses the already proved unramified unitary Satake bound, not a new global base-change or temperedness assertion.

For any such rank-r Euler product P\^{}S, its eigenvalues at an unramified E-prime satisfy \textbar alpha\_p\textbar(Np)\^{}(-z0)\textless=(Np)\^{}(-h). Thus

\[
 \begin{split}
 |P^S|&\ge\exp\left(-rn\left(2^{-h}
                          +\frac{2^{1-h}}{h-1}\right)\right)\\
       &>\frac12.
 \end{split}                                           \tag{LB29}
\]

The first line follows by the same reciprocal-factor inequality and by summing all integers\textgreater=2, with the integral test for the tail from3. The last line holds for r=43,n=84,h=859 and r=169,n=336,h=733; for example the exponent is smaller than2\^{}(-716) in both cases. The estimates are uniform over all actual automorphic constituents and multiplicity copies with these ranks and prescribed finite level.

Consequently the product of the two finite partial standard factors appearing in the vertical Rallis formulas satisfies

\[
 \boxed{
 |\mathcal L_0^S(z_0)\mathcal L_1^S(z_0)|>
 \begin{cases}
 2^{-277221}R_S(\vartheta),&J14,\\
 2^{-1171821}R_S(\vartheta),&J56.
 \end{cases}}
 \tag{LB30}
\]

Here e\^{}(-n)\textgreater2\^{}(-2n), and the two factors1/2 in (LB29) have been included. All exceptional local zeta forms remain outside this partial-Euler estimate. In particular no absolute-ramified basic vector or G/Q quotient scalar is silently absorbed into (LB30).

\subsection{9. Exact scope for a top-preimage bound}\label{exact-scope-for-a-top-preimage-bound}

LIFT31 proves U=D* and ker(U)=0. The minimal-norm right inverse is D\emph{(DD})\^{}(-1). Its norm is the reciprocal of the least singular value of U. Equations (LB24) and (LB30) remove one former ineffective ingredient of that singular-value problem: the actual boundary Hecke value and its complete long string now have explicit positive lower bounds uniform in the initial automorphic constituent.

To turn these into a bound for U one must still compare the exact finite quotient models in LIFT31(VA22b,VA24a) with the local Rallis Hermitian forms, retain the basic-vector normalization and the explicit doubling denominators, and divide by the actual nonzero real-intertwiner norms used to define the source multiplicity maps. The real first discrete-series formal degree and second compact dimension in VA28--29 are not set to1. Changing a source real intertwiner by c changes its associated multiplicity map by c\^{}(-1) and correspondingly changes the inverse cost. Therefore (LB30) alone is not asserted to be the lower singular value of U.

\subsubsection{Primary references}\label{primary-references}

Tate, J. T. (1950). \emph{Fourier analysis in number fields and Hecke's zeta-functions} {[}Doctoral dissertation, Princeton University{]}. Sections2.5, 4.4,4.5. Original scan: \url{https://sites.math.rutgers.edu/~alexk/2023S572/Tate1950.pdf} .

Kudla, S. S. (2001). \emph{Tate's thesis} {[}Lecture notes, School on Automorphic Forms, L-functions and Number Theory, Hebrew University, March2001{]}. Proposition3.5, Theorem4.3 and Corollary4.4. Lecture manuscript: \url{https://u.cs.biu.ac.il/~reznikov/courses/kudla-1.pdf} .

Gan, W. T., Qiu, Y., \& Takeda, S. (2014). The regularized Siegel--Weil formula (the second term identity) and the Rallis inner product formula. \emph{Inventiones Mathematicae, 198},739--831. Sections11.3--11.7, in particular the normalization (11.3) and Theorem11.4. Local copy: lift\_refs/gqt2014.pdf.

The new quantitative inequalities (LB3), (LB21), (LB24), (LB26), and (LB30) are proved here. The host, character, and actual tower identifications are cited from their checked findings as indicated.

\clearpage

\section{AM-K049 --- an invariant spectral cutoff at the actual translated point}\label{am-k049-an-invariant-spectral-cutoff-at-the-actual-translated-point}

Author: GPT-6 Astra Pro, ALP\_MACHINE \nolinkurl{/root/kummer_extension}. 22 November 2026 UTC; inherited reasoning effort is not separately exposed.

Status: author whole-proof check PASS; exact finite certificate PASS. Independent and root reviews are separate events.

This finding addresses the global comparison left open in K048. The main operator theorem below is an explicit estimate in the actual unitary representation, uniform on its full selected compact module and on every multiplicity space. It uses a periodic spectral multiplier built from the actual rational central lattice. Automorphic evaluation is unchanged by this multiplier. Consequently an ordinary standard-origin Sobolev bound can be applied after the multiplier has made the translated vector small.

The role hypothesis remains \textbf{at most f/2 top labels}, with ell=ell\_+ and every permitted word and canonical twist in K044. No condition on the zero central Fourier coefficient is used. In particular this proof does not infer centre-zero vanishing from cuspidality of the full Heisenberg constant term.

\subsection{1. Hilbert and differential conventions}\label{hilbert-and-differential-conventions}

Write G\_infty=product\_tau U(a,1), n=a+1 and

\[
 d=f n^2<2^{23},\qquad a=f=42\ \hbox{or}\ 168.
 \tag{IC1}
\]

Let H be the actual holomorphic discrete series of K043/K044, with its given unitary norm, or a tensor product with an arbitrary Hilbert multiplicity space on which G\_infty acts trivially. Its selected compact subspace is included by the exact PBW isometry J=rho\_all\^{}(-1) Cj=rho\_all j of K043, where the last equality uses its polynomial realization Cj=D\_all j and rho\_all²=D\_all. A vector v in that subspace below has norm1. Raw derivatives retain the factor rho\_all separately.

At each label use the following real matrix-unit basis of u(a,1):

\begin{itemize}
\tightlist
\item
  i E\_ii for each diagonal entry;
\item
  E\_ij−E\_ji and i(E\_ij+E\_ji) in the positive compact block;
\item
  E\_i,n+E\_n,i and i(E\_i,n−E\_n,i) for the noncompact pairs.
\end{itemize}

There is also the negative diagonal basis element already included in the first item. Each basis matrix has operator norm at most1. Across labels use their disjoint union, denoted X\_1,\ldots,X\_d.~For an ordered word I=(i1,\ldots,im) write d pi(X\_I)=d pi(X\_i1)\ldots d pi(X\_im). The estimate below controls every word of length at most d; it does not replace these noncommuting derivatives by commuting scalar derivatives.

Let g be exactly K047(ST1), at all real labels. The actual packet point (o,tau\_c) equals (g\_infty o,1\_f) in the principal-J1 quotient, because K047 proves g\_f\^{}−1=tau\_c k\_f with k\_f in J1. Thus, for the right regular representation and evaluation E at the standard identity component,

\[
                    E_g(v)=E(\pi(g)v).
 \tag{IC2}
\]

This fixes which of g and g\^{}−1 occurs. The compact factor in g is in the displayed SU(1,1) circle; its noncompact dilation has singular values sqrt(T),1,\ldots,1,1/sqrt(T).

\subsection{2. A bounded integral basis for the actual real field}\label{a-bounded-integral-basis-for-the-actual-real-field}

Let F0=Q(zeta28)\^{}+, let B be the degree-seven real subfield of Q(zeta29), and set F=F0 B or F0 B(sqrt5,sqrt13), as in K030/K032. Put t=zeta28+zeta28\^{}−1. A Gaussian period eta=sum\_(h in H) zeta29\^{}h, where H has order4, generates B: its seven conjugates cannot all be rational, since their sum is−1 and a rational algebraic integer common to them would equal−1/7. The prime degree7 then proves generation. The fields in this compositum are linearly disjoint, as already checked from their disjoint discriminant supports in K030/K032.

Consequently the tensor monomials

\[
       t^i\eta^j\,(\sqrt5)^e(\sqrt{13})^h,
       \quad 0\le i<6,\quad0\le j<7,\quad e,h\in\{0,1\},
 \tag{IC3}
\]

with the last two factors omitted for J14, form a Q-basis of integral elements. Every conjugate of each element has absolute value less than 2\^{}21, using \textbar t\textbar≤2, \textbar eta\textbar≤4 and sqrt65\textless16.

There exists an integral basis b1,\ldots,bf of O\_F whose coordinate matrix relative to (IC3) is upper triangular with every entry of absolute value at most1. Here is a direct lattice proof. For any rational full lattice Lambda containing Z\^{}f, its intersection with the first coordinate axis is generated by c e1 with 0\textless c≤1. Quotient by that axis, apply induction to the projected lattice, lift the resulting basis, and reduce the first coordinate of every lift modulo c into {[}0,c). Repeating gives the assertion. Apply it to the coordinates of O\_F in the tensor basis.

Let B\_F=(sigma\_tau(b\_j)) be the real embedding matrix. Then

\[
 \max_{\tau,j}|(B_F)_{\tau j}|<2^{29},\qquad
 \|B_F\|_{op}<2^{37},\qquad
 \|B_F^{-1}\|_{\infty}<2^{7000}.
 \tag{IC4}
\]

The first two inequalities use f≤168\textless2\^{}8. For the last, \textbar det B\_F\textbar²=disc(F)≥1. The singular-value inequality gives \textbar\textbar B\_F\^{}−1\textbar\textbar\_op≤\textbar\textbar B\_F\textbar\textbar\_op\^{}(f−1), and the infinity norm adds at most sqrt(f). Thus the exponent7000 is a conservative integer upper bound. The same bound applies to the transpose inverse. An actual basis can be fixed by this fractional Hermite procedure; no class-number, unit, regulator or unproved basis-height estimate is being used.

\subsection{3. The periodic multiplier and exact evaluation invariance}\label{the-periodic-multiplier-and-exact-evaluation-invariance}

For each label put N\_tau=(e1+e\_n)(e1+e\_n)*J in the selected two-plane, Y\_tau=iN\_tau, and

\[
                      A_\tau=-i\,d\pi(Y_\tau).
 \tag{IC5}
\]

These are commuting positive self-adjoint operators in the holomorphic representation. Their unitary spectral realization is the Laguerre model of K048. Write epsilon\_tau=tau(i)/i. The actual rational elements

\[
             u_b=1+iQbN\in\Gamma_1,\qquad b\in O_F,
 \tag{IC6}
\]

act as exp(iQ sum\_tau epsilon\_tau sigma\_tau(b) A\_tau). In particular any standard-component automorphic evaluation distribution E satisfies E pi(u\_b)=E. The finite companion lies in J1 at every place, exactly as in K047. This invariance is not imposed on an arbitrary nonarithmetic real translation.

Choose

\[
      \delta=2^{-12000},\qquad L=8d+8f+100<2^{27}.
 \tag{IC7}
\]

Let h\_delta be the 1-periodic function equal to (1−(t/delta)²)\^{}L on \textbar t\textbar\textless delta modulo Z and equal to0 elsewhere. It is C\^{}(L−1), takes values in {[}0,1{]}, and h\_delta(0)=1. Set

\[
 \begin{split}
 h(\theta)&=\prod_{j=1}^f h_\delta(\theta_j),\\
 \Phi(x)&=h\left(\frac{Q}{2\pi}B_F^{\mathsf T}
                       \operatorname{diag}(\epsilon_\tau)x\right),
 \qquad \Phi(A)=\Phi(A_1,\ldots,A_f).
 \end{split}\tag{IC8}
\]

The circular constant is denoted by π in (IC8); Q is the actual product of the two G primes. This bounded multiplier is defined by the joint spectral theorem. It is not a finite averaging projector.

\subsubsection{Lemma IC9: exact invariance}\label{lemma-ic9-exact-invariance}

For every smooth automorphic vector w on which the indicated evaluation and derivatives are defined,

\[
                              E(\Phi(A)w)=E(w).
 \tag{IC9}
\]

\textbf{Proof.} Expand h in its Fourier series sum\_(m in Z\^{}f) c\_m exp(2πi m·theta). The corresponding operator is sum\_m c\_m pi(u\_(sum m\_j b\_j)). Because h is C\^{}(L−1) and L\textgreater2d+f+3, its Fourier coefficients are summable after multiplication by any fixed polynomial of degree2d. Conjugation by a central unipotent acts on the real Lie algebra by a polynomial of degree at most2 in its parameter. It follows that this Fourier series converges absolutely in every Sobolev norm through order d.~In that topology the standard local evaluation functional is continuous. Therefore its value is sum\_m c\_m E(w)=h(0)E(w)=E(w). This proof retains the zero-frequency window; it uses no vanishing assertion at frequency zero. QED.

The smooth vectors needed below have all finite Sobolev norms. The same Fourier argument justifies each subsequent commutation through order d.

\subsection{4. Support geometry, including the zero window}\label{support-geometry-including-the-zero-window}

The support condition in (IC8) says that theta lies within delta in each coordinate of an integer vector m. Its centre in the A spectral variables is

\[
       x^0(m)=\frac{2\pi}{Q}\operatorname{diag}(\epsilon_\tau)
                                     B_F^{-\mathsf T}m.
 \tag{IC10}
\]

After the dilation g, the spectral variable of the original unit vector is x=T x\^{}0. By (IC4), 2π\textless8, Q≥1 and T\textless2\^{}1001, each dilated window has infinity radius less than

\[
                         2^{-3000}.
 \tag{IC11}
\]

In fact these elementary bounds give radius less than2\^{}−3996. The integer vector m corresponds exactly to an inverse-different element xi; the sign-correct centre is (2πT/Q)(epsilon\_tau sigma\_tau(xi)). Thus K048's norm comparison applies to every nonzero centre, whether or not every coordinate of that centre is positive.

We will need the slightly larger finite-energy spaces with local compact energies at most

\[
                   \widetilde K_\tau=K_\tau+2d\le2K_\tau<2^{963}.
 \tag{IC12}
\]

Their lowest SU(1,1) weights are still at least K\_*=2ell+a≥2\^{}485, because the ambient representation has not changed. K048(ST23), divided by at most2 for these enlarged energies, shows that each nonzero window centre has some coordinate with absolute value greater than2\^{}186 times its local tilde-K. A window meeting the positive spectral orthant cannot have that coordinate large and negative. Subtracting the radius (IC11) therefore gives, at every point in its positive intersection,

\[
                      x_\tau>2^{180}\widetilde K_\tau
                  \quad\hbox{for some label tau}.
 \tag{IC13}
\]

The window with m=0 instead has 0\textless x\_tau\textless2\^{}−3000 at every label. This window is treated directly; it is not discarded as a cuspidal constant.

\subsection{5. Filtered Hilbert norm on the whole enlarged energy space}\label{filtered-hilbert-norm-on-the-whole-enlarged-energy-space}

Let E\_tilde be the tensor product of the subspaces whose compact energies are bounded by (IC12), with any multiplicity Hilbert space attached. For a unit vector w in E\_tilde, including an entangled vector,

\[
               \|1_{\operatorname{supp}\Phi}(A)\pi(g)w\|
                                      <\exp(-2^{492}).
 \tag{IC14}
\]

\textbf{Proof of the nonzero-window bound.} The argument of K048(ST17) remains valid with K\textless2\^{}963, changing its harmless numerical constant1535 to1536. On a single positive discrete-series ladder, for x≥2\^{}180 K, \textbar f\_(kappa,n)(x)\textbar≤exp(−x/4). At most K+1≤2K eligible indices occur in one ladder. Cauchy--Schwarz, orthogonality of different ladders, and integration with dx/x show that the spectral projection onto x≥2\^{}180 K has squared operator norm at most

\[
  2K\int_{2^{180}K}^{\infty} e^{-x/2}\frac{dx}{x}
       \le 2^{-178}\exp(-2^{179}K).
 \tag{IC15}
\]

The union over all f labels, using K≥K\_*≥2\^{}485, has squared norm less than exp(−2\^{}663). The other factors are unitary spectral transforms, so this union bound is valid for entangled vectors and multiplicities. By (IC13), it covers every nonzero cutoff window.

\textbf{Proof of the zero-window bound.} The explicit Laguerre coefficient formula also gives, for positive integers kappa, integers n≥0 and n+kappa≤K,

\[
            |f_{\kappa,n}(x)|\le(Kx)^{\kappa/2}e^{Kx}.
 \tag{IC16}
\]

To verify it without an asymptotic formula, divide each polynomial coefficient by its constant coefficient. The absolute ratio at degree j is (n)\_falling,j/((kappa)\_j j!)≤(n/kappa)\^{}j/j!. Summing gives exp(nx/kappa). The normalized constant coefficient equals sqrt(Gamma(n+kappa)/n!)/Gamma(kappa), which is at most K\^{}(kappa/2): the numerator squared is the product of kappa−1 integers at most K, and Gamma(kappa)=(kappa−1)!≥1. All ladder weights here are integral, as proved in K048(ST9) from the actual compact-circle action. Dropping exp(−x/2) proves (IC16).

On 0\textless x\textless h0=2\^{}−3000 and K\textless2\^{}963, Kx\textless2\^{}−2037. Integrating (IC16) and using kappa≥K\_* yields, on the whole truncated energy space at one label,

\[
 \|1_{(0,h0)}(A)\|^2
      \le \frac{2e^2K}{K_*}(K h0)^{K_*}
      <\exp(-2^{494}).
 \tag{IC17}
\]

For the last inequality use log2\textgreater1/2, K\_*≥2\^{}485 and log K\textless963. The zero-window product projection is contained in this one-coordinate projection. Combining (IC15--17) gives squared norm less than exp(−2\^{}493), hence (IC14). QED.

The projection notation in (IC17) means the norm of the spectral projection restricted to the indicated finite-energy subspace. Its norm on the entire infinite-dimensional representation is of course1.

\subsection{6. All multiplier and noncommuting derivative constants}\label{all-multiplier-and-noncommuting-derivative-constants}

\subsubsection{6.1. Derivatives of the multiplier}\label{derivatives-of-the-multiplier}

For 0≤m≤2d, the explicit polynomial defining the bump gives

\[
 \|h_\delta^{(m)}\|_\infty
                      \le2^L(2L/\delta)^m.
 \tag{IC18}
\]

This follows by summing the absolute derivatives of its binomial expansion on its support. Its derivatives through this order extend continuously by zero because 2d\textless L. For the product bump and the linear change in (IC8), the operator norms of B\_F and the actual Q\textless2\^{}71 give

\[
 \max_{|\alpha|\le2d}\|\partial^\alpha\Phi\|_\infty
       \le2^{fL+(12000+28+108)2d}<\exp(2^{42}).
 \tag{IC19}
\]

The factor108 includes the coordinate sum in the linear chain rule, not just one matrix entry. These derivatives have the same support as Phi.

\subsubsection{6.2. Moving a Lie derivative past the cutoff}\label{moving-a-lie-derivative-past-the-cutoff}

Since N\_tau²=0, ad\_(N\_tau)\^{}3 is zero on the entire matrix Lie algebra. For any basis matrix X at that label, commuting d pi(X) past the Fourier series of Phi therefore produces only multiplier derivatives of orders 0,1,2, followed by d pi(X), d pi({[}Y\_tau,X{]}) or one half d pi({[}Y\_tau,{[}Y\_tau,X{]}{]}), with the appropriate unit-modulus powers of i. This identity follows directly by expanding Ad(exp(−tY\_tau))X in the Fourier integral; it is not a commuting- derivative approximation.

In the displayed real matrix basis, every bracket of two basis matrices has coefficient l1 norm at most16, by direct entry multiplication. Y\_tau=iE11−iEnn−i(E1n−En1) has coefficient l1 norm3. Thus the sum of the coefficient norms in one such commutation is at most 1+48+48²/2\textless2\^{}12. Iterating an ordered word of length m gives terms with multiplier derivative order at most2m, a Lie word of length m, and total coefficient l1 norm at most2\^{}(12m).

\subsubsection{6.3. Lie-word norms before and after the actual g}\label{lie-word-norms-before-and-after-the-actual-g}

For every word of length m≤d and every unit selected vector v,

\[
                       \|d\pi(X_I)v\|\le2^{1000m}.
 \tag{IC20}
\]

Here are the algebraic bounds behind this deliberately loose constant. Every compact highest-weight coordinate initially has absolute value less than2A. Acting by m matrix-unit derivatives changes compact weights by at most m, because the adjoint compact weights are differences of standard coordinate weights. Thus all compact types reached through order d have coordinate size at most M=2A+d.

Let C\_n=sum\_(i,j)Eij Eji in the complexified enveloping algebra. Its infinitesimal-character scalar is sum\_i(Lambda\_i²−rho\_i²), whose absolute value is at most n((2A)²+n²). In the real matrix basis,

\[
 2C_n=-2\sum_i(iE_{ii})^2-
       \sum_{X\in\mathcal K_{\mathrm{off}}}X^2+
       \sum_{X\in\mathcal P}X^2.
 \tag{IC21}
\]

Here K\_off and P are respectively the compact off-diagonal and noncompact basis sets listed in Section1. The identity is checked pair by pair from Eij Eji+Eji Eij. Taking its unitary quadratic form expresses the sum of the noncompact derivative norms in terms of the compact Casimir and the displayed scalar. The compact Casimir on weights bounded by M is at most2n(M+n)². Consequently each basis derivative has norm at most8n(M+n) on this finite sum of compact types. Since n≤169, d\textless2\^{}23 and A=2\^{}960, this is less than2\^{}1000. Induction proves (IC20). The same argument applies with multiplicities or tensor products, since it is an operator-norm bound.

The actual g has \textbar\textbar g\textbar\textbar\_op=\textbar\textbar g\^{}−1\textbar\textbar\_op=sqrt(T). Expanding Ad(g\^{}−1)X in the basis therefore has coefficient l1 norm at most 2n²T\textless2\^{}1018. Hence

\[
        \|d\pi(\operatorname{Ad}(g^{-1})X_I)v\|
                              \le2^{2100m}.
 \tag{IC22}
\]

Every vector on the left has compact energy at label tau at most K\_tau+2m, because the chosen SU(1,1) circle has matrix-unit weights between−2 and2. It consequently belongs to E\_tilde from Section5.

\subsubsection{Theorem IC23: an actual Sobolev bound}\label{theorem-ic23-an-actual-sobolev-bound}

For every ordered word I of length at most d,

\[
 \boxed{\quad
       \|d\pi(X_I)\Phi(A)\pi(g)v\|<\exp(-2^{491}).
       \quad}
 \tag{IC23}
\]

\textbf{Proof.} Use Section6.2 to move the word past Phi. Move each remaining Lie word past pi(g), which produces (IC22). Each resulting vector belongs to the enlarged energy space, so (IC14) applies. Each multiplier derivative is bounded by (IC19) on the same support. Thus the full norm is at most

\[
        2^{2112d}\exp(2^{42})\exp(-2^{492})
                                    <\exp(-2^{491}).
 \tag{IC24}
\]

All exponent comparisons use d\textless2\^{}23. This proof includes the zero window and all mixed noncommuting derivatives. QED.

Summing over all ordered words through length d adds at most (d+1)\^{}(d+1) terms. Its logarithm is less than2\^{}29. Thus

\[
        \sum_{|I|\le d}\|d\pi(X_I)\Phi(A)\pi(g)v\|
                                     <\exp(-2^{490}).
 \tag{IC25}
\]

\subsection{7. An actual global evaluation and selected-jet bound}\label{an-actual-global-evaluation-and-selected-jet-bound}

R031(SB1--SB27), independently whole-checked in this run, proves for the actual raw R030 Haar and scalar automorphic L2 norm

\[
 |F(1)|\le C_{\rm ev}\sum_{|I|\le d}\|R_{X_I}F\|_2,
                  \qquad \log C_{\rm ev}<2^{32}<2^{200}.
 \tag{IC26}
\]

Its proof uses a small embedded product-exponential chart at the standard origin. The exact matrix-unit basis is the one in Section1. No separation at the translated point is needed. The sum of all principal-J1 component norms only enlarges its right side. The lemma extends by Sobolev continuity to the finitely differentiable cutoff vectors used here: the Fourier partial sums converge in every required derivative norm.

Let U be the isometric realization of an actual holomorphic-discrete-series isotypic subspace in this scalar automorphic Hilbert space, including its multiplicity space. Apply (IC26) to U(Phi(A)pi(g)v), and use (IC9),(IC25). It follows that

\[
 \boxed{\quad
        |E(U(\pi(g)v))|\le\exp(-2^{489})\,\|v\|
        \quad(v\hbox{ in the whole selected isotypic subspace}).
        \quad}
 \tag{IC27}
\]

For nonzero v the displayed estimate is strict by (IC25--26); the weak inequality also includes v=0. The same convention applies to (IC31) and (IC33), which are strict for nonzero forms.

The displayed notation E(U(\ldots)) means evaluation of the automorphic function at the standard identity component. By (IC2) it is evaluation at the actual translated packet point. Equivalently one can suppress U and write E\_g on the automorphic isotypic space itself. The estimate is valid whether or not the particular projected cuspidal space is nonzero. It holds on the genuine finite e0/generic-character subspace because that is a subspace of this actual principal-level space; no substitute finite projector is used.

\subsubsection{Canonical vector-valued and PBW normalization}\label{canonical-vector-valued-and-pbw-normalization}

Here is the exact scalar-to-vector-valued conversion, so that an arbitrary global realization scalar is not silently set to1. Let V\_alpha be the minimal compact type, of dimension d\_alpha, with its given metric, and let j\_alpha:V\_alpha→H be its isometric inclusion. Write the scalar isotypic space as M tensor H, with U its unitary identification. For u in M define the V\_alpha-dual-valued automorphic form by its components

\[
       \langle F_u(g),e_i\rangle
                  =U(u\otimes j_\alpha e_i)(g)
 \tag{IC28}
\]

in an orthonormal basis e\_i. Orthogonality of these scalar components in the raw automorphic Hilbert space gives exactly

\[
                       \|F_u\|_{\rm raw}^2
                                   =d_\alpha\|u\|^2.
 \tag{IC29}
\]

Conversely the actual minimal-type-valued holomorphic forms in this isotypic component have this description. This is the elementary compact coefficient realization: their scalar components intertwine the minimal compact type, and its multiplicity-one inclusion extends to the given holomorphic discrete series. It introduces no choice of eigenfunction normalization beyond the unitary isotypic identification.

For the selected isometry J\_lambda=rho\_all\^{}(-1)Cj=rho\_all j of K043, in that fixed polynomial realization, define the normalized selected jet by

\[
   \langle\mathcal J F_u(g),v\rangle
                   =U(u\otimes J_\lambda v)(g).
 \tag{IC30}
\]

Moving the PBW creation operators through U shows that this is exactly rho\_all\^{}−1 times the source's raw Fischer/ordinary-derivative contraction, with the fixed K043 matrix-unit sign convention. The contragredient reading is its conjugate and has the same norm. No other period or intertwiner factor occurs in this canonical coefficient realization. If a separately rescaled realization is used, its actual scalar remains.

Apply (IC27) to u tensor v for all unit v. The norm of the resulting functional in v is the norm of the selected-jet vector. Thus

\[
 \boxed{\quad
 \|\mathcal JF_u(o,\tau_c)\|
       \le\frac{\exp(-2^{489})}{\sqrt{d_\alpha}}
                                 \|F_u\|_{\rm raw,J1}.
 \quad}
 \tag{IC31}
\]

There is no factor sqrt(dim V\_lambda): (IC27) is an operator estimate on the full selected isotypic Hilbert tensor, rather than a separate bound on individual coordinates. The finite component is the literal scalar principal-level realization. An additional abstract finite vector-valued realization must keep its own coefficient norm map; it is not silently substituted for this scalar finite-level model.

\subsubsection{Exact volume conventions and necessary detector norm}\label{exact-volume-conventions-and-necessary-detector-norm}

Refinement from the actual bottom level J\_B to J1 multiplies both raw norm squared and total volume by the same finite index. Consequently the normalized-average norm is unchanged. R030 proves

\[
               V(J1)=\frac{m_n}{(a!)^f}
                         \prod_{j=0,1}|GL_n(\mathbf F_{g_j})|^f.
 \tag{IC32}
\]

Therefore (IC31) becomes, in the actual bottom normalized-average norm,

\[
 \boxed{\quad
 \|\mathcal JF(o,\tau_c)\|
       \le\sqrt{\frac{V(J1)}{d_\alpha}}\,
                    \exp(-2^{489})\|F\|_{\rm av,J_B}.
 \quad}
 \tag{IC33}
\]

For the unnormalized adelic convention in which every finite maximal compact has volume1, replace V(J1) in (IC33) by \textbar Gcal\textbar, exactly as in R031(SB27). These two conventions must not be mixed.

In particular, if an actual detector in this space has a prescribed normalized selected jet of norm1 at this packet point, then necessarily

\[
              \|F\|_{\rm av,J_B}
                      >\sqrt{\frac{d_\alpha}{V(J1)}}
                                         \exp(2^{489}).
 \tag{IC34}
\]

This is a necessary norm bound, not an existence statement. A raw jet reinstates rho\_all: the right side of (IC33) is multiplied by rho\_all for the raw contraction, and a prescribed raw value1 has the accordingly different reciprocal bound.

The conclusion concerns this one actual packet point, the exact upper gap, and the role cohort top-count≤f/2. It does not locate the earlier coupled image within the packet dual, prove a terminal period nonzero, or show that a particular actual prior-image row requires this coordinate normalization. It therefore supplies an actual endpoint-evaluation cost for this specified normalization, not the completed DP44 net cost or all-source Host/Scale constructors.

\subsection{8. Sources and verification scope}\label{sources-and-verification-scope-2}

Groenevelt, W. (2003). Laguerre functions and representations of su(1,1). \emph{Indagationes Mathematicae, 14}(3--4), 329--352.

The known unitary Laguerre transform is Groenevelt (2003), Section3.1, equation(3.2), Section3.2, equation(3.7), and Proposition3.2; its exact normalization was directly checked in K048. The source is \url{https://arxiv.org/pdf/math/0302342} and \url{https://doi.org/10.1016/S0019-3577(03)90050-4} . The polynomial estimates, lattice cutoff, matrix commutators and Casimir bounds are proved above. No unproved global Fourier-coefficient bound or cusp-neighbourhood interpolation theorem is substituted.

K030/K032 supply the actual fields and product polarizations. K043 supplies the PBW isometry; K044 the exact all-word parameters; K047 the actual packet point, rational g, and unchanged principal finite level; K048 the signed central lattice and unitary spectral estimates. R030 supplies the exact Haar volumes. R031 supplies the complete actual standard-component Sobolev theorem with this same differential basis. These are applied within their stated scopes.

The companion finite certificate checks the free degree-two Casimir expression, exact small matrix commutators and their coefficient bounds, small bump/Laguerre coefficient diagnostics, and every actual arithmetic exponent margin. It does not compute an automorphic Fourier coefficient, period, class basis, or global jet matrix. Those analytic statements are proved through the invariant cutoff and the independently checked Sobolev comparison, rather than inferred from finite samples.

\subsection{Author verification record}\label{author-verification-record-2}

Author whole saved report and complete certificate read and reconstructed PASS at 2026-11-22T05:06:50.812247+00:00. The proof check includes the actual rational central invariance and packet direction, fractional integral-basis bound, finite-Sobolev Fourier convergence, nonzero and zero spectral windows, all noncommuting derivative/Casimir constants, the exact R031 Haar comparison, and scalar-to-vector-valued PBW normalization. Final saved notation clarifies J=rho\_all\^{}(-1)Cj=rho\_all j in the polynomial realization. The certificate replay passed3342 exact checks at2026-11-22T05:06:31.730644+00:00. The top-count, exact upper-gap, normalized-jet, one-packet-point and actual-prior-image limitations remain explicit. No global automorphic basis or Fourier coefficient was computed.

Final statement alignment at 2026-11-22T05:08:35.269032+00:00: root review noticed that homogeneous norm bounds must include the zero vector. IC27, IC31 and IC33 now use weak inequalities, with an explicit strictness statement for nonzero vectors/forms. IC34 retains its strict necessary norm bound. These statements and their unchanged derivation were targeted reread; no estimate or certificate calculation changes.

\clearpage

\section{ALP-LIFT-37: exact G-place Rallis factors for the actual vertical composition}\label{alp-lift-37-exact-g-place-rallis-factors-for-the-actual-vertical-composition}

Author: /root/lift\_accounting, GPT-6 Astra Pro, ALP\_MACHINE. 22 November 2026 UTC. Inherited reasoning effort is not separately exposed.

\textbf{Author whole-proof review completed: 2026-11-22T05:06:55Z.}

This computes the local positive Hermitian form, not only a nonzero theta quotient. The finite functions are exactly the two conjugated DatumVec G functions of LIFT31(VA11--12). Their normalization remains the raw additive-Schwartz normalization. An additional table gives the effect of normalizing both of these vectors to Hilbert norm1.

\subsection{1. Measures, representations and the actual functions}\label{measures-representations-and-the-actual-functions}

Let F\_v be an actual degree-one relative split G completion, with ring O, residue cardinality q and uniformizer varpi. Put

\[
 r=a+1,\quad b=r+t,\quad N=r+2t,\qquad
 \kappa_j(q)=\prod_{i=1}^{j}(1-q^{-i}).                \tag{VG1}
\]

The two actual profiles have (r,t)=(43,1759) or(169,1633). Group Haar is normalized by vol(GL\_j(O))=1, and additive matrix Haar by vol(M\_(m,n)(O))=1. At these G places this is the actual unramified additive normalization in the split Hom model. Any conversion to a different global group-Haar convention remains a separate factor.

Let pi\_0 be any unitary representation of GL\_r(F\_v), and take v in pi\_0\^{}(K\_r\^{}1). In the actual upward convention the first function f0 is the indicator of integral matrices reducing to E\_r={[}I\_r;0{]}. The second function f1 is supported on integral matrices reducing to

\[
       \begin{pmatrix}I_r&B\\0&D\\0&0\end{pmatrix},
             \quad B\in M_{r,t}(k),\quad D\in\mathrm{GL}_t(k),
 \tag{VG2}
\]

and has the unit value chi1(det D)\^{}(-1)chi2(det D). These are the literal functions in LIFT31(VA11--12), read there from DatumVec.lean49--98. The source Hom action is

\[
    [\omega(g)f](X)=\beta(\det g)^{-1}
                              |\det g|^{m/2}f(Xg).     \tag{VG3}
\]

Use conjugate-linear-first inner products, the same convention as DO036 and LIFT31. The compact-model period operator is

\[
 B_f(v;k)=\int_{\mathrm{GL}_n(F_v)}
 f(kE_nh)|\det h|^{m/2}\beta(\det h)^{-1}\pi(h)v\,dh,
                  \quad k\in K_m.                    \tag{VG4}
\]

Its target representation has the unitary compact induced-model metric with probability Haar on K\_m. This metric is not declared equal to the oscillator quotient metric; their exact comparison is the main point below.

\subsection{2. Why the Gram identity is absolutely convergent at both edges}\label{why-the-gram-identity-is-absolutely-convergent-at-both-edges}

Suppose f is supported on integral matrices of full residue rank n. For k in K\_m, the condition f(kE\_nh)!=0 implies h is integral, since kE\_n has an integral left inverse. Its reduction has rank n only if h is invertible modulo varpi. Hence h belongs to K\_n. Both actual functions (VG2), including the first edge, have this property. Therefore every orbit integral (VG4) has compact support. The double integral obtained by expanding its squared norm has both orbit variables in K\_n, uniformly in k. All integrands are bounded, so absolute Fubini is valid without a rank inequality such as m\textgreater=2n. In particular it is valid at the second actual edge, where b\textless N\textless2b.

The full-rank matrix disintegration of LIFT16(S6), or DO036(6), is

\[
 \int_{M_{m,n}(F_v)}H(X)dX
 =C_{n,m}(q)\int_{K_m}\int_{\mathrm{GL}_n(F_v)}
                   H(kE_nh)|\det h|^m\,dh\,dk,
 \quad C_{n,m}(q)=\frac{\kappa_m(q)}{\kappa_{m-n}(q)}.
 \tag{VG5}
\]

This constant is the additive probability that an m by n residue matrix has full rank. Expand the square of (VG4), write its second orbit variable as hg, and apply (VG5). The factors of beta become beta(det g)\^{}(-1); the determinant factors become \textbar det h\textbar\^{}m\textbar det g\textbar\^{}(m/2). This proves the exact positive identity

\[
 \boxed{
 Z_f(v):=\int_{\mathrm{GL}_n(F_v)}
       \langle f,\omega(g)f\rangle\langle v,\pi(g)v\rangle dg
       =C_{n,m}(q)\|B_f(v)\|_{K_m}^{2}.}
 \tag{VG6}
\]

This is DO036(7) with its convergence argument specialized to the actual full-residue-rank support. The same identity applies to sesquilinear pairings by polarization. The doubling/Weil and bilinear-versus-conjugated transport is the one fixed in DO028 and LIFT31; no unspecified local Schur proportionality remains in (VG6).

\subsection{3. The first and second actual compact sections}\label{the-first-and-second-actual-compact-sections}

The exact first period and its complete compact norm are already computed in LIFT31(VA19--20):

\[
 B_{f0}(v;1)=|\mathrm{GL}_r(k)|^{-1}v,
 \qquad
 \|B_{f0}(v)\|^2=
       |\mathrm{GL}_r(k)|^{-2}{b\brack r}_q^{-1}\|v\|^2.
 \tag{VG7}
\]

The section is supported on the identity(r,t)-parabolic cell. Its actual second-source sharp character is the required chi1 character; it is not an arbitrary vector in the first theta module.

Write w=B\_f0(v). The actual second compact projector has volume

\[
 v_J=\frac{q^{rt}|\mathrm{GL}_t(k)|}{|\mathrm{GL}_b(k)|},
 \quad
 c_G=\frac{v_J}{|\mathrm{GL}_r(k)|}.
 \tag{VG8}
\]

LIFT31(VA21--22a) proves that the second section is supported on the identity(b,t)-parabolic cell and has unitary translates of v\_J w as its values there. Hence

\[
 \|B_{f1}(B_{f0}(v))\|^2
       =c_G^2{N\brack b}_q^{-1}{b\brack r}_q^{-1}\|v\|^2.
 \tag{VG9}
\]

All character values have modulus1. The arguments establishing support and covariance in VA19--22a apply to every v in the full initial K\_r\^{}1-fixed subspace, so these are exact scalar identities there, not bounds only on one chosen generic or Steinberg line.

\subsection{4. The actual composed Hermitian factor}\label{the-actual-composed-hermitian-factor}

For clarity about the first metric, (VG6) says that its Rallis quotient metric is C\_(r,b) times the compact induced-model metric. Using this metric as the source inner product on the next step multiplies its Gram form by C\_(r,b). Therefore the composed local positive form, in the unchanged compact quotient models, is

\[
 \begin{split}
 \mathcal Z_{G,10}(v)
 &=C_{r,b}(q)C_{b,N}(q)
          \|B_{f1}(B_{f0}(v))\|^2\\
 &=A_G^{\rm raw}(q)\|v\|^2,
 \end{split}                                           \tag{VG10}
\]

where the exact factors are

\[
 C_{r,b}=\kappa_b/\kappa_t,
 \qquad C_{b,N}=\kappa_N/\kappa_t,
 \tag{VG11}
\]

and

\[
 \boxed{\quad
 A_G^{\rm raw}(q)=q^{-(b+r)^2}
                        \frac{\kappa_t(q)^2}
                             {\kappa_b(q)\kappa_r(q)}.
 \quad}                                                \tag{VG12}
\]

To verify the simplification, use

\[
 |\mathrm{GL}_j(k)|=q^{j^2}\kappa_j,
 \quad {m\brack n}_q=q^{n(m-n)}
                         \frac{\kappa_m}{\kappa_n\kappa_{m-n}},
 \quad
 c_G=q^{-r(b+r)}\frac{\kappa_t}{\kappa_b\kappa_r}.
 \tag{VG13}
\]

Substitution into (VG9)--(VG11) cancels kappa\_N, leaves kappa\_t squared/(kappa\_b kappa\_r), and gives exponent -(2r+t)(b+r)=-(b+r)\^{}2. Thus this calculation fixes the Hermitian normalization which a nonzero local-quotient theorem alone would leave unspecified.

\subsection{5. Raw versus unit-Schwartz normalization}\label{raw-versus-unit-schwartz-normalization}

The raw additive norms from LIFT31(VA22b) are

\[
 \|f0\|^2=q^{-br},\qquad
 \|f1\|^2=q^{-Nb+rt}|\mathrm{GL}_t(k)|,
 \quad
 \|f0\|^2\|f1\|^2=q^{-b(b+r)}\kappa_t.
 \tag{VG14}
\]

Divide (VG12) by this product when, and only when, both finite vectors have been normalized to Hilbert norm1. The resulting composed factor is

\[
 \boxed{\quad
 A_G^{\rm unit}(q)=q^{-r(b+r)}
                     \frac{\kappa_t}{\kappa_b\kappa_r}=c_G.
 \quad}                                                \tag{VG15}
\]

In particular the normalized two-edge factor is c\_G, not c\_G squared and not1. No square root of a finite group order was silently included in either raw vector.

\subsection{6. Explicit products over every actual G place}\label{explicit-products-over-every-actual-g-place}

For each of the two rational G primes g\_i there are exactly f places of F above it, all with residue cardinality g\_i. Thus

\[
 \mathcal A_G^{\rm raw}
     =\prod_{i=0}^1 A_G^{\rm raw}(g_i)^f,
 \qquad
 \mathcal A_G^{\rm unit}
     =\prod_{i=0}^1 A_G^{\rm unit}(g_i)^f.
 \tag{VG16}
\]

Both are explicit positive numbers, independent of the automorphic constituent and the large weight/word parameter. The field, ranks, raw kernels, and finite G characters are the actual K031/K032 data. For q\textgreater=3 and t\textgreater=r we have

\[
 \frac{\kappa_t}{\kappa_b\kappa_r}>1,
 \qquad
 \frac{\kappa_t^2}{\kappa_b\kappa_r}
       >\frac{\kappa_t}{\kappa_r}
       =\prod_{j=r+1}^t(1-q^{-j})>\frac12.
 \tag{VG17}
\]

The first inequality uses kappa\_t\textgreater kappa\_b and kappa\_r\textless1. For the second, the elementary product lower bound1-sum(q\^{}-j) is at least1-1/(q-1), with strictness for the finite product. Using the exact G-product bounds Q\textless2\^{}45 or2\^{}71 from LIFT36 gives

\begin{longtable}[]{@{}
  >{\raggedright\arraybackslash}p{(\columnwidth - 6\tabcolsep) * \real{0.2500}}
  >{\raggedright\arraybackslash}p{(\columnwidth - 6\tabcolsep) * \real{0.2500}}
  >{\raggedright\arraybackslash}p{(\columnwidth - 6\tabcolsep) * \real{0.2500}}
  >{\raggedright\arraybackslash}p{(\columnwidth - 6\tabcolsep) * \real{0.2500}}@{}}
\toprule\noalign{}
\begin{minipage}[b]{\linewidth}\raggedright
Host
\end{minipage} & \begin{minipage}[b]{\linewidth}\raggedright
(r,b,t,f)
\end{minipage} & \begin{minipage}[b]{\linewidth}\raggedright
Raw squared-norm factor
\end{minipage} & \begin{minipage}[b]{\linewidth}\raggedright
Unit-Schwartz squared-norm factor
\end{minipage} \\
\midrule\noalign{}
\endhead
\bottomrule\noalign{}
\endlastfoot
J14 & (43,1802,1759,42) & \textgreater2\^{}(-6433607334) & \textgreater2\^{}(-149943150) \\
J56 & (169,1802,1633,168) & \textgreater2\^{}(-46338383784) & \textgreater2\^{}(-3973204872) \\
\end{longtable}

The raw exponents are qG\emph{f}(b+r)\^{}2+2f; the unit exponents are qG\emph{f}r*(b+r). The factor2f in the raw table bounds exactly the two sets of f kappa ratios by1/2. These deliberately coarse binary bounds can be replaced by the exact products (VG16).

\subsection{7. What this settles in the inverse calculation}\label{what-this-settles-in-the-inverse-calculation}

Equations(VG10)--(VG16) are the actual finite G contribution to the local Rallis quotient norms. They hold on every prescribed initial vector and hence on every automorphic multiplicity copy; there is no new selection of a better vector or a new finite test. Together with LIFT36 they make the G factors and the boundary global Euler factors effective.

To bound the full vertical adjoint one must additionally retain the Q factors, the basic vectors and doubling denominators away G/Q, the first real formal degree and second compact dimension, and the actual source real-intertwiner and global Haar factors. Those remaining factors are not set to1 here. This result is a local Hermitian normalization and a strict reduction of the cost problem; it is not by itself a global lower singular value or an integral preimage assertion.

\subsubsection{Checked input locators}\label{checked-input-locators}

\begin{itemize}
\tightlist
\item
  LIFT31, vertical\_adjoint\_surjectivity.md, VA11--13 and VA19--22c: exact DatumVec functions, characters, compact sections and raw norms.
\item
  LIFT16, spherical\_tests\_and\_29\_only\_chain.md, S6: full-rank matrix measure and its probability normalization.
\item
  DO036, definite\_oscillator\_raw\_first\_nonvanishing.md, DO36.2--7: actual Hom operator and doubled Gram identity.
\item
  DO028, definite\_oscillator\_doubling\_dictionary.md: literal bilinear/source convention and its doubling transport.
\item
  K031(HC18--20), K032(QH27,QH30--32): actual fully split G places.
\end{itemize}

The extension of the absolutely convergent Gram identity to these second-edge raw G functions follows from the compact support proof in section2. The exact simplifications (VG12), (VG15), and the global products (VG16) are proved here.

\clearpage

\section{AM-DO-047: exact Rallis normalization of the two prescribed vertical Q factors}\label{am-do-047-exact-rallis-normalization-of-the-two-prescribed-vertical-q-factors}

Author: GPT-6 Astra Pro, ALP\_MACHINE /root/definite\_oscillator. 22 November 2026 UTC. Inherited reasoning effort is not separately exposed.

The actual two vertical Q functions have an exact positive local Hermitian normalization on the first identity-parabolic image. This determines the scalar omitted by an abstract small-theta identification. With the normalized group and additive measures of DO041, put

\[
 r=a+1,\quad b=r+t,\quad N=r+2t,\qquad
 \kappa_j=\prod_{i=1}^j(1-q^{-i}).                         \tag{QN1}
\]

At a genuine Q place, after dividing by both actual raw Schwartz norms and by the two split standard L-factors, the two-step local Hermitian factor is exactly

\[
 \boxed{\ c_Q=q^{-rt}\frac{\kappa_r\kappa_{2t}}{\kappa_b},
 \qquad \frac{q^{-rt}}{16}<c_Q<q^{-rt}.\ }                \tag{QN2}
\]

The corresponding local inverse norm on the actual initial spherical line, with the normalized local Rallis quotient metric specified in (QN24)--(QN25), is

\[
 c_Q^{-1/2}=q^{rt/2}
       \left(\frac{\kappa_b}{\kappa_r\kappa_{2t}}\right)^{1/2},
 \qquad q^{rt/2}<c_Q^{-1/2}<4q^{rt/2}.                   \tag{QN3}
\]

These are local factors. Global Rallis constants, archimedean metrics, global L-values and the arithmetic class/period comparison remain separate. This is not a global lower singular value or a nonzero terminal brC/tilQ value.

\subsection{1. Actual conventions and the carried quotient calculation}\label{actual-conventions-and-the-carried-quotient-calculation}

Use exactly DO041 sections 1--5, not a different occurrence map. The actual ranks are (r,t,b,N)=(43,1759,1802,3561) or (169,1633,1802,3435); in particular t\textgreater=r. Let F be the actual local field, O its valuation ring, varpi its uniformizer and q its residue cardinality. At the literal Q places DegOne gives F=Q\_q. The local argument also works over a general nonarchimedean F with the stated measures. Every matrix O-lattice has additive measure one and every K\_j=GL\_j(O) has group Haar volume one. All representation inner products are conjugate-linear in the first argument.

The honest Hom action and unequal splitting labels are

\[
 [\omega_{\alpha,\beta}(g,h)f](X)
 =\alpha(\det h)\beta(\det g)^{-1}
  |\det g|^{m/2}|\det h|^{-n/2}f(h^{-1}Xg),              \tag{QN4}
\]

\[
 (\alpha_0,\beta_0)=((\mu_0\xi)^{-1}\eta,(\mu_0\xi)^{-1}),
 \qquad
 (\alpha_1,\beta_1)=((\mu_1\xi)^{-1},(\mu_1\xi)^{-1}\eta).
                                                               \tag{QN5}
\]

Here mu0,mu1,eta are unitary and unramified and xi is the actual tame tuple character. The determinant-root convention and the source's conjugate upward kernel have already been checked in DO020/028/041. No Fourier or additive-character phase is introduced in (QN4).

Let pi0 be an irreducible unitary spherical representation of GL\_r and v in its K\_r-fixed line. Nothing below requires pi0 to be tempered. The two target representations, with their probability compact normalized-induction norms, are

\[
 \pi_1=(\pi_0\eta)\times((\mu_0\xi)^{-1}\eta)\det_t,
 \qquad
 \pi_2=\pi_0\times(\mu_0\xi)^{-1}\det_t
                       \times(\mu_1\xi)^{-1}\det_t.       \tag{QN6}
\]

The notation chi det\_t means chi composed with determinant. For X\_n=(I\_n,0)\^{}transpose, the whole-Schwartz quotient is the holomorphic value at z=(m-n+1)/2 of

\[
 Q_z(f,w)(k)=\alpha(\det k)\int_{GL_n(F)}
 f(k^{-1}X_nh)|\det h|^{z+(n-1)/2}
                \beta(\det h)^{-1}\pi(h)w\,dh.            \tag{QN7}
\]

DO041 proves the continuation into the actual irreducible unitary induction using Jiang--Luo's all-irreducibles local zeta ideal, not an unproved extension from the full-rank locus. Retain this exact map.

The prescribed conjugate Q functions are

\[
 f_0(X)=\mathbf1_{M_{b,r}(O)}(X)
  \mathbf1_{\bar X_{>r,*}=0}\mathbf1_{\det\bar X_{[r],[r]}\ne0}
                     \xi(\det X_{[r],[r]})^{-1},           \tag{QN8}
\]

\[
 f_1(Z)=\mathbf1_{M_{N,b}(O)}(Z)
  \mathbf1_{\bar Z_{>r,[r]}=0}\mathbf1_{\det\bar Z_{[r],[r]}\ne0}
                     \xi(\det Z_{[r],[r]})^{-1}.           \tag{QN9}
\]

The columns after r in f1 are unrestricted integral columns. Let F0=Q0(f0,v) and F2=Q1(f1,F0), where Qi denotes (QN7) at the actual point. Write

\[
 A={b\brack r}_q=q^{rt}\frac{\kappa_b}{\kappa_r\kappa_t},
 \quad B={N\brack r}_q=q^{2rt}\frac{\kappa_N}{\kappa_r\kappa_{2t}},
 \quad L_t(\theta)=\prod_{i=1}^t(1-\theta(\varpi)q^{-i})^{-1},
 \quad\theta=\mu_1/\mu_0.                                \tag{QN10}
\]

The complete support, Schur-complement and Jacobian proof is already DO041(16)--(27), independently checked. Its settled conclusions are

\[
 F_0(k)=\mathbf1_{J_{r,t}}(k)\xi(\det\bar D_t)^{-1}v,
 \quad \|F_0\|^2=A^{-1}\|v\|^2,\quad
 \|F_2\|^2=A^{-2}|L_t(\theta)|^2B^{-1}\|v\|^2,            \tag{QN11}
\]

\[
 \|f_0\|^2=\kappa_rq^{-rt},\qquad
 \|f_1\|^2=\kappa_rq^{-2rt}.                             \tag{QN12}
\]

Here J\_(r,t) is the reduction parabolic in K\_b and D\_t its lower diagonal block. The linear quotient scalar remains Lambda\_Q=A\^{}\{-1\}L\_t(theta), including its phase. This report does not replace it by its absolute value in a linear identity.

\subsection{2. A positive majorant for the actual first-image vector}\label{a-positive-majorant-for-the-actual-first-image-vector}

The usual large-rank absolute Bochner estimate does not apply to the separate second integral at the actual point: N/2 may be at most b-1. We prove the needed Gram convergence directly.

Let rho be normalized unitary induction from the trivial representations of GL\_r and GL\_t to GL\_b, and let e be its compact spherical vector of constant value one. Its compact norm is one. The representation rho preserves the cone of nonnegative compact functions: Iwasawa reexpression multiplies by the positive factor delta\_(P\_(r,t))\^{}(1/2). Consequently its spherical coefficient m\_rho(g)=\textless e,rho(g)e\textgreater{} is nonnegative.

In the same induced realization of pi1, the pointwise norm of F0 on K\_b is at most \textbar\textbar v\textbar\textbar. All inducing characters and pi0 are unitary. Taking absolute values inside the compact coefficient integral gives

\[
 |\langle F_0,\pi_1(g)F_0\rangle|
                  \le \|v\|^2m_\rho(g)\quad(g\in GL_b(F)).
                                                               \tag{QN13}
\]

More explicitly, write kg=p(g,k)k'. The integrand on the left is bounded by delta(p)\^{}(1/2)\textbar\textbar v\textbar\textbar\^{}2 times the two support indicators 1\_J(k)1\_J(k'), which can both be discarded. The remaining integral is the right side. This proves the bound for the actual identity-cell vector, not merely for a spherical vector in pi1.

The Satake parameters of rho are the union of the two trivial-block strings q\^{}((r-1)/2),\ldots,q\^{}(-(r-1)/2) and q\^{}((t-1)/2),\ldots,q\^{}(-(t-1)/2). Since t\textgreater=r, all parameters and their inverses have modulus at most q\^{}((t-1)/2). The positive vacuum coefficient in S(M\_(N,b)) is

\[
 A_N(g)=q^{-N\sum_i|\lambda_i|/2}
 \quad\text{for }g\in K_b\operatorname{diag}(\varpi^{\lambda_i})K_b.
                                                               \tag{QN14}
\]

The exact Macdonald Cartan estimate used and proved from its finite Weyl formula in LIFT16(S12)--(S16) gives

\[
 \int_{GL_b}A_N(g)m_\rho(g)\,dg<\infty,\qquad
 q^{-N/2}q^{(b-1)/2}q^{(t-1)/2}=q^{-1}.                  \tag{QN15}
\]

This is a strict geometric convergence margin, including coincident Satake parameters, whose limiting Weyl formula contributes only a polynomial in the Cartan exponents. It remains valid if rho is realized before taking its spherical irreducible quotient: its spherical coefficient is the same. No temperedness or genericity of pi1 has been assumed.

Since \textbar f1\textbar\textless=1\_(M\_(N,b)(O)), (QN13)--(QN15) prove absolute convergence of the actual local doubled oscillator Gram integral

\[
 Z_1=\int_{GL_b}
  \langle f_1,\omega_{\alpha_1,\beta_1}(g,1)f_1\rangle
  \langle F_0,\pi_1(g)F_0\rangle\,dg.                    \tag{QN16}
\]

This is the local Rallis value at s0=t/2, under the carried Hom/tensor and bilinear/conjugated-input dictionary of DO028. The standard L-function point is z*=s0+1/2=(t+1)/2.

\subsection{3. Exact disintegration and the continuation scalar}\label{exact-disintegration-and-the-continuation-scalar}

For n\textless m, matrix additive measure and probability compact Haar satisfy

\[
 \int_{M_{m,n}(F)}H(X)dX
  =C_{n,m}\int_{K_m}\int_{GL_n}
       H(k^{-1}X_nh)|\det h|^m\,dh\,dk,\qquad
 C_{n,m}=\prod_{i=0}^{n-1}(1-q^{i-m})
              =\frac{\kappa_m}{\kappa_{m-n}}.             \tag{QN17}
\]

The full-rank complement has measure zero. On its open orbit the Jacobian determines the measure up to scalar; the matrices with primitive reduction have measure C\_(n,m), whereas their h-coordinate lies in K\_n of volume one. This is the exact normalization already proved in LIFT16(S6)/DO036(6), recalled to fix its use here.

We claim that the actual whole-Schwartz second quotient obeys

\[
 \boxed{\quad Z_1=C_{b,N}\|F_2\|^2,\qquad
                  C_{b,N}=\kappa_N/\kappa_t.\quad}        \tag{QN18}
\]

Here is a complete justification despite possible nonabsolute Bochner convergence of the single orbit integral. In (QN7) write z=z*+w and first truncate h to a compact K\_b-biinvariant union of Cartan double cosets H\_R, increasing to GL\_b. Denote the compact function so obtained by Q\_(w,R). Its squared norm is an absolutely convergent double integral in h and u. The unweighted absolute Gram majorant, after u=hg and (QN17), is at most

\[
 C_{b,N}^{-1}\|v\|^2
             \int_{GL_b} A_N(g)m_\rho(g)\,dg<\infty.      \tag{QN19}
\]

Indeed the integrand before this bound is \textbar f1(k\^{}(-1)X\_bh)f1(k\^{}(-1)X\_bu)\textbar{} \textbar det h det u\textbar\^{}(N/2) times the absolute coefficient of F0 at h\^{}(-1)u. Equations (QN13), (QN17) and \textbar f1\textbar\textless=vacuum give precisely (QN19). Thus this bound controls absolute values before the two orbit integrals are interchanged, not just their integrated complex value.

Nonzero f1(k\^{}(-1)X\_bh) forces h integral, since k is integrally invertible and X\_b has an integral left inverse. Therefore the extra factors \textbar det h\textbar\^{}w and \textbar det u\textbar\^{}conjugate(w) have modulus at most one for Re(w)\textgreater=0. Dominated convergence using (QN19) proves that Q\_(w,R) is Cauchy in compact L2, uniformly on compact subsets of Re(w)\textgreater=0. The limit is holomorphic for Re(w)\textgreater0 and continuous at w=0.

There is no hidden weak-vector ambiguity. Each Q\_(w,R) is right K\_N(1)-invariant and has values in pi1\^{}(K\_b(1)). For the latter, left replacement h→j\^{}(-1)h for j in K\_b(1) preserves H\_R, beta1 and the raw matrix function: on integral h it preserves its reduction; on nonintegral h both raw functions vanish. Admissibility makes pi1\^{}(K\_b(1)) finite-dimensional, and K\_N/K\_N(1) is finite. Hence all limits lie in one finite-dimensional compact-function space. The truncated functions are entire in w there. Their limit agrees with the Bochner integral for Re(w) sufficiently large. The identity theorem identifies it with the whole-Schwartz holomorphic quotient (QN7) for Re(w)\textgreater0, and continuity identifies its w=0 limit with F2. DO041(13) ensures this quotient is holomorphic at and to the right of z*: the initial unitary rank-r factor is within its convergent half-plane, and the added unitary Hecke factors have arguments 1+w,\ldots,t+w.

Finally expand the w=0 Gram integral, put u=hg and apply (QN17). The determinant powers become \textbar det h\textbar\^{}N\textbar det g\textbar\^{}(N/2), the character ratio beta1(h)beta1(hg)\^{}(-1) is beta1(g)\^{}(-1), and the representation coefficient is \textless F0,pi1(g)F0\textgreater. The result is C\_(b,N)\^{}(-1) times (QN16), proving (QN18). The target alpha1 on the compact variable has absolute value one and cancels. This derivation fixes the scalar; it does not use Schur uniqueness to choose it.

\subsection{4. Exact values, actual ramification and normalized local costs}\label{exact-values-actual-ramification-and-normalized-local-costs}

Equations (QN10)--(QN12) and (QN18) give

\[
 Z_1=q^{-4rt}\frac{\kappa_r^3\kappa_t\kappa_{2t}}{\kappa_b^2}
                 |L_t(\theta)|^2\|v\|^2,                \tag{QN20}
\]

\[
 \frac{Z_1}{\|f_1\|^2\|F_0\|^2}
       =q^{-rt}\frac{\kappa_r\kappa_{2t}}{\kappa_b}
                    |L_t(\theta)|^2.                    \tag{QN21}
\]

For the first edge the corresponding exact disintegration gives

\[
 Z_0=C_{r,b}\|F_0\|^2
       =\kappa_rq^{-rt}\|v\|^2
       =\|f_0\|^2\|v\|^2,\qquad C_{r,b}=\kappa_b/\kappa_t.
                                                               \tag{QN22}
\]

Here a direct compact orbit calculation suffices: the full rank of the raw reduction forces the orbit variable into K\_r. Alternatively the large first gap gives the absolute Bochner proof of DO036(7).

The literal Q set is the ramification set of xi. Precisely, Tuples.lean defines Qprime(xi)=\{v:not inertAt(E,v)\textless=ker(xi)\} at lines 128--129, while tupGrp at lines 112--114 kills wildAt and the diagonal F-ideles. At a split place, the two local components of xi are inverse characters. Their restriction to O\^{}times is therefore nontrivial when v belongs to Qprime(xi). Thus xi in (QN5) is actually tamely ramified at each place now under consideration. This is not an arbitrary additional nontriviality hypothesis. The tuple of trivial character simply has empty Q set.

Since pi0 is spherical and mu\_j unramified, its standard local L-factor after twisting by mu\_j xi is one. Multiplicativity and the literal source label beta1 give

\[
 L(z*,\pi_0\beta_0^{-1})=1,\qquad
 L(z*,\pi_1\beta_1^{-1})=L_t(\theta).                    \tag{QN23}
\]

The split standard factor in the local Rallis formula is this L-value times its contragredient/unitary-conjugate value. Hence normalizing Z0,Z1 by their actual split standard factors gives

\[
 Z_0^*=\|f_0\|^2\|v\|^2,\qquad
 \frac{Z_1^*}{\|f_1\|^2\|F_0\|^2}=c_Q.                  \tag{QN24}
\]

For clarity, composing the two local Hermitian quotients retains both disintegration factors, not just the second one:

\[
 \frac{C_{r,b}C_{b,N}\|F_2\|^2}
      {\|f_0\|^2\|f_1\|^2\|v\|^2|L_t(\theta)|^2}
                      =c_Q.                            \tag{QN25}
\]

Thus the result is on the actual preceding image with its actual first Rallis metric. It is not a comparison using an independently renormalized middle representation.

For all q\textgreater=2 and all j, kappa\_j\textgreater1/4: at q=2 the first two factors are 3/8 and the remaining product is at least 1−sum\_(i\textgreater=3)2\^{}(-i)=3/4. Omitting factors handles j=0,1,2. Since 2t\textgreater=b, kappa\_(2t)\textless=kappa\_b; also kappa\_r\textless1. These prove the two bounds in (QN2) and hence (QN3). The exact logarithmic inverse cost is

\[
 \tfrac12rt\log q+
       \tfrac12(\log\kappa_b-\log\kappa_r-\log\kappa_{2t}).
                                                               \tag{QN26}
\]

If the two actual local group measures give K\_r,K\_b volumes nu\_r,nu\_b instead of one, each orbit quotient integral scales by its source volume. In (QN17) its disintegration coefficient becomes C\_(n,m)/nu\_n.~Consequently the composed Hermitian factor in (QN25) is multiplied by nu\_r nu\_b. Compact induction norms still use probability compact Haar. This states the complete measure change needed before inserting the result into a global Rallis identity.

\subsection{5. Sources, scope and proof verification}\label{sources-scope-and-proof-verification}

The established inputs are DO041's entire literal quotient proof, DO028's actual Hom/Rallis dictionary, DO036(5)--(7)'s disintegration, and LIFT16(S3)--(S18)'s exact vacuum coefficient and Cartan estimate. No later global scalar or arithmetic row normalization is presumed.

Direct source locators are under the recovered source commit 7cbf7a0c32fca1badf05931879a37c3f8a2b03b4:

\begin{itemize}
\tightlist
\item
  \nolinkurl{Formalism/LocalToGlobal/IndefiniteReal/DatumVec.lean,} 64--70,83--98: omegaQ and the unchanged IsStepVec exponent.
\item
  \nolinkurl{Formalism/LocalToGlobal/IndefiniteReal/LinksChar.lean,} 255--269: the actual odd/even vertical character pairs.
\item
  \nolinkurl{Formalism/LocalToGlobal/Divisibility/Tuples.lean,} 58--67,112--129: split degree-one good places, tame diagonal-trivial tuple characters and the actual Qprime ramification set.
\item
  Development/AbelianVarietyTheory/Definitions/LatticeFrames/ ChainParSucc.lean, 132--133: the conjugate kernel in Step.up.
\end{itemize}

Primary mathematical sources used through the already checked results are:

Macdonald, I. G. (1995). \emph{Symmetric functions and Hall polynomials} (2nd ed.). Oxford University Press. Chapter V, (2.9), (3.4), pp.298--299; section 4, (4.4)--(4.5), p.301. Exact spherical coefficient, Cartan bound and normalization.

Jiang, D., \& Luo, Z. (2023). Certain Fourier operators and their associated Poisson summation formulae on GL1. \emph{Pacific Journal of Mathematics, 326}(2), 301--372. Theorem 2.4(1)--(2), pp.313--314. The entire-Schwartz local zeta ideal; the continuation itself is proved with its finite-dimensional-vector justification in DO041.

Gan, W. T., Qiu, Y., \& Takeda, S. (2014). The regularized Siegel--Weil formula (the second term identity) and the Rallis inner product formula. \emph{Inventiones Mathematicae, 198}, 739--831. Sections 11.3, 11.5--11.6, with the actual source-to-tensor dictionary in DO028. Here the local form is computed directly, so no qualitative nonvanishing theorem is used to choose its normalization.

This result proves exact positive local Hermitian factors and their inverse cost. It does not identify an arithmetic inverse, a Smith minor, a global multiplicity-space scalar or a terminal period.

Verification: author whole-proof read of sections 1--5 and (QN1)--(QN26) completed after the exact certificate. The finite-dimensional Abel limit, absolute pre-Fubini majorant, actual tuple ramification and the two successive Hermitian metrics were checked separately. The companion \nolinkurl{definite_oscillator_vertical_q_rallis_certificate.py} passed 7,447 exact assertions at 2026-11-22T05:08:05.594588+00:00. It independently counts finite full-rank and raw Q residue supports, checks both disintegration factors and all scalar simplifications for rational unitary phases, both actual rank margins, Haar scaling and explicit bounds. This is a diagnostic for the proved identities, not an automorphic L-value or local representation computation.

\clearpage

\section{AM-DO-048 --- the actual real formal degree in the fixed ball Haar}\label{am-do-048-the-actual-real-formal-degree-in-the-fixed-ball-haar}

Author: GPT-6 Astra Pro, ALP\_MACHINE /root/definite\_oscillator. Inherited reasoning effort is not separately exposed. 22 November 2026 UTC.

This is an exact normalization of a known formal-degree formula for the actual bottom holomorphic discrete series. It supplies the real factor left explicit in LIFT31 and the subsequent inverse-cost assembly. It is valid throughout the K044 interval, for every CM role and binary word. It does not identify a freely normalized real intertwiner with a prescribed source intertwiner, or prove a new global nonzero value.

\subsection{1. Measure, parameters, and result}\label{measure-parameters-and-result}

Put G=U(a,1), K=U(a)×U(1), and normalize compact Haar to probability. The quotient measure is exactly R030/K046's measure

\[
 d\mu_B(z)=\pi^{-a}(1-|z|^2)^{-a-1}\,dV(z),\qquad z\in B^a.
 \tag{FD1}
\]

Thus integration on G is integration on B\^{}a followed by probability integration on K. For an irreducible unitary square-integrable representation τ, its formal degree d(τ) in this measure is defined by

\[
 \int_G |\langle\tau(g)v,w\rangle|^2\,dg
     =\frac{\|v\|^2\|w\|^2}{d(\tau)}.
 \tag{FD2}
\]

The centre of G is compact. FD2 is therefore an ordinary group integral, with the centre already included in the probability normalization of K.

At one actual CM label use the settled K043 dictionary, with zero-based indices:

\[
 p_i=v_i\ (0\le i<a-1),\qquad p_{a-1}=-b_0,
 \qquad p_-=-b_1,
\] \[
 \alpha_i=p_i-\frac{a-2}{2}+i,
 \qquad\beta=p_--\frac a2.
 \tag{FD3}
\]

The actual holomorphic representation has minimal compact type Wα⊗z\^{}β, and its Harish-Chandra multiset, in decreasing order, is

\[
 h=(p_0,\ldots,p_{a-1},p_-).
 \tag{FD4}
\]

All differences p\_i−p\_- are positive. For these actual parameters the minimal compact weight α is dominant integral. Since a=42 or168, h−ρ is also dominant integral for GL\_(a+1), where ρ=(a/2,a/2−1,\ldots,−a/2).

\subsubsection{Theorem FD5}\label{theorem-fd5}

In precisely the measure FD1,

\[
 \boxed{
 d(\tau)=\dim(W_\alpha)\prod_{i=0}^{a-1}(p_i-p_-)
    =a!\,\dim V_{h-\rho}.
 }
 \tag{FD5}
\]

The same positive number is the formal degree of τ∨ or its complex conjugate. A unitary determinant twist has the same formal degree. There is no additional compact-centre, 2π, or finite-cover factor.

\subsection{2. Exact Bergman calibration and proof}\label{exact-bergman-calibration-and-proof}

Nakahama (2016), equation (5.7), printed p.960, uses π\^{}(−a)dV in the rank-one ball case of SU(1,a). His Theorem5.1, printed pp.960--961, normalizes constants to the given fibre metric, with normalizing factor

\[
 c_\Lambda=\prod_{j=1}^a(\Lambda-(j-1)+k_j-1).
 \tag{FD6}
\]

The exact source-to-Nakahama translation has already been proved in K043(PB2--PB4). Namely, choose an integer c making k=α+c1 polynomial and put Λ=−β−c.~For j=i+1, each factor in FD6 becomes

\[
 \Lambda-i+k_{i+1}-1
     =\alpha_i-\beta-i-1=p_i-p_-.
 \tag{FD7}
\]

The auxiliary determinant shift c cancels. The convergence condition Λ+k\_a\textgreater a becomes p\_(a−1)−p\_-\textgreater0, which holds in every actual case. Consequently, if Hτ is the invariant Hilbert realization in which constant functions have exactly the fibre metric, then

\[
 \|P\|_{H_\tau}^2
   =c_\tau\int_G\|F_P(g)\|_{W_\alpha}^2\,dg,
 \qquad c_\tau=\prod_i(p_i-p_-).
 \tag{FD8}
\]

Here F\_P is the canonical vector-valued lift, equivalently evaluation at0 after translating P (using g or g\^{}−1 changes no integral). Formula FD8 is exactly Nakahama's displayed metric after its fibre-metric factor is included; it is also the already checked K046(BL5--BL7). The compact central extension from SU(1,a) to U(1,a), and swapping the two Hermitian blocks, leave that metric unchanged.

Let dα=dim(Wα) and choose an orthonormal basis e\_1,\ldots,e\_dα of the constant compact type. Evaluation at0 is the orthogonal projection onto that type: positive-degree polynomials are orthogonal to constants by the scalar compact circle, and constants have their prescribed fibre metric. Therefore

\[
 \|F_P(g)\|^2
       =\sum_{j=1}^{d_\alpha}
          |\langle\tau(g)P,e_j\rangle_{H_\tau}|^2.
 \tag{FD9}
\]

Apply the correct Schur orthogonality FD2 to each summand. Since the sum is finite, no exchange of an infinite series is involved. Combining FD8 and FD9 gives

\[
 \frac{\|P\|^2}{c_\tau}
       =\frac{d_\alpha\|P\|^2}{d(\tau)}.
 \tag{FD10}
\]

Choose P nonzero. This proves the first identity in FD5, including the factor dα. It is precisely the distinction between the squared norm of a vector-valued evaluation and one scalar matrix coefficient.

For the second identity, Weyl's dimension formula and FD3 give

\[
 d_\alpha
  =\prod_{0\le i<j<a}\frac{\alpha_i-\alpha_j+j-i}{j-i}
  =\frac{\prod_{i<j<a}(p_i-p_j)}{\prod_{j=1}^{a-1}j!}.
 \tag{FD11}
\]

Multiplying by all a differences to p\_- yields

\[
 d(\tau)=\frac{\prod_{0\le i<j\le a}(h_i-h_j)}
                        {\prod_{j=1}^{a-1}j!}
       =a!\frac{\prod_{i<j\le a}(h_i-h_j)}
                        {\prod_{j=1}^{a}j!}.
 \tag{FD12}
\]

The last quotient is the GL\_(a+1) dimension of V\_(h−ρ). This proves FD5. Complex conjugation or contragredience preserves the absolute matrix coefficient in FD2 after exchanging its two vectors and possibly g with g\^{}−1. A unitary determinant twist multiplies that coefficient by a scalar of modulus1. These observations prove the remaining claims.

To verify the central normalization directly, the finite central cover Z×SU(a,1)→U(a,1) maps probability Haar on the compact factors to probability Haar on K. The quotient ball measure is FD1 on both sides. The absolute squared coefficient is invariant under Z. Thus integrating on the cover produces the same number, rather than a factor equal to the order of the finite kernel. QED.

\subsection{3. Scalar calibration and independent root-product anchor}\label{scalar-calibration-and-independent-root-product-anchor}

There is an elementary scalar check of the entire measure convention. For the scalar holomorphic discrete series with parameter λ\textgreater a, its normalized constant matrix coefficient has squared absolute value Y\^{}λ, where Y=1−\textbar z\textbar². Radial integration gives

\[
 \pi^{-a}\int_{B^a}Y^{\lambda-a-1}\,dV
       =\frac{\Gamma(\lambda-a)}{\Gamma(\lambda)}
       =\prod_{j=1}^a(\lambda-j)^{-1}.
 \tag{FD13}
\]

Thus its formal degree is ∏\_(j=1)\^{}a(λ−j). In rank1 this is λ−1. This also calibrates the Harish-Chandra root-product formula to FD1.

The primary paper of van Haastrecht and Zhang (2026), Section4.1, equations(11)--(13), explicitly states Schur orthogonality and the Harish-Chandra root-product formula up to its measure constant. Its Lemma4.1 and Proposition4.3 give, on a rank-one ball with its unnormalized Lebesgue quotient density, c\_G=π\^{}(−a)a!. Multiplying that Haar measure by π\^{}(−a) multiplies the formal degree by π\^{}a, leaving exactly a! times the compact Weyl dimension, as in FD5. The proof above independently fixes this factor through Nakahama's vector norm and FD10; no arbitrary proportionality constant remains.

\textbf{Source-audit precision.} This argument does not use Theorem4.4 of van Haastrecht--Zhang. Its displayed equation(16) replaces the product of squared norms in Schur orthogonality by \textbar\textless v,w\textgreater\textbar². For orthogonal nonzero fibre vectors those quantities differ. The correct sum in FD9--FD10 has dα scalar matrix coefficients, and gives Bergman kernel constant d(τ)/dα when the raw fibre integral is used. The scalar calibration and Section4.1 root-product statement are unaffected. No reliance on the vector-valued assertion of that theorem is made.

The original Harish-Chandra (1956) bibliographic attribution is given below. Its original text was not retrieved in this check; the direct source readings were the published Nakahama theorem and the expressly identified portions of the 2026 primary article.

\subsection{4. Closed formula on the actual K044 family}\label{closed-formula-on-the-actual-k044-family}

For one CM role write E=E\_T,E\_H orE\_L. Throughout K044's exact interval,

\[
 p_i=4\ell(a-1-i)\quad(i<a-1),\quad
 p_{a-1}=-E+\ell,\quad p_-=-E-\ell-q,
 \qquad0\le q\le\ell.
 \tag{FD14}
\]

These bottom parameters do not depend on the binary word. The selected higher compact type does depend on that word; FD5 is the formal degree of the ambient bottom discrete series and is unchanged by that choice. The two factors of FD5 simplify to

\[
 d_\alpha=
  \frac{(4\ell)^{(a-1)(a-2)/2}}{(a-1)!}
       \prod_{j=1}^{a-1}\bigl(E+(4j-1)\ell\bigr),
 \tag{FD15}
\] \[
 c_\tau=(2\ell+q)
       \prod_{j=1}^{a-1}\bigl(E+(4j+1)\ell+q\bigr).
 \tag{FD16}
\]

Formula FD15 follows by separating the first a−1 entries in the Vandermonde product FD11. Their differences are 4ell(j−i), whose Vandermonde denominator cancels all but (a−1)!. FD16 lists the a differences to the negative entry. Their product is an exact positive rational expression, in fact a positive integer by FD5. It is directly computable without constructing the representation.

The crude but explicit estimate

\[
 1\le\prod_\tau d(\tau_\tau)
       <2^{960 f a(a+1)/2}
 \tag{FD17}
\]

holds on both actual hosts. To see this, the K044 inequalities give ell\textless2\^{}485, E\_T\textless2\^{}955+2\^{}485, E\_H\textless2\^{}495, and E\_L\textless2\^{}485. For a≤168 and q≤ell the largest HC difference in FD14 is E+(4a−3)ell+q\textless2\^{}956\textless2\^{}960. Every positive-root difference is therefore less than2\^{}960. The denominator in FD12 is at least1, and there are a(a+1)/2 differences at each of the f labels. The left inequality follows from FD5 and integrality. The exact formula, rather than FD17, should be used when a sharper cost matters.

\begin{longtable}[]{@{}lrr@{}}
\toprule\noalign{}
Host & a=f & Exponent in FD17 \\
\midrule\noalign{}
\endhead
\bottomrule\noalign{}
\endlastfoot
J14 & 42 & 36,408,960 \\
J56 & 168 & 2,289,530,880 \\
\end{longtable}

\subsection{5. The real Rallis factors and their precise scope}\label{the-real-rallis-factors-and-their-precise-scope}

For an isometrically normalized real Howe tensor on the first vertical edge, LIFT31's positive real local test is

\[
 Z_{0,\infty,\tau}=d(\tau_{0,\tau})^{-1}.
 \tag{FD18}
\]

If the oscillator carries the contragredient, FD5 gives the same number. At the second vertical edge its source is compact; with probability Haar the corresponding unit tensor has

\[
 Z_{1,\infty,\tau}=\dim(\tau_{1,\tau})^{-1}.
 \tag{FD19}
\]

Thus the exact product of these particular real local factors is

\[
 \prod_\tau
 \frac1{a!\,\dim V_{h_\tau-\rho}\,\dim(\tau_{1,\tau})}.
 \tag{FD20}
\]

Two different scalar operations must be distinguished. Multiplying the actual oscillator test vector on edge j by s\_(j,τ) multiplies its raw Gram test by \textbar s\_(j,τ)\textbar². In contrast, suppose the fixed theta operator is factored as A⊗I and the prescribed source intertwiner is I\_(j,τ)=s\_(j,τ)I\_(j,τ)\^{}iso, while the oscillator test is held fixed. The same operator then has multiplicity map A\_(j,τ)=s\_(j,τ)\^{}−1 A\_(j,τ)\^{}iso. Its squared norm is divided by \textbar s\_(j,τ)\textbar². Under the latter convention of LIFT38, the two-edge multiplicity-map energy is therefore divided by ∏\emph{τ\textbar s}(0,τ)s\_(1,τ)\textbar². These are distinct operations; computing the formal degree does not set either prescribed scalar to1.

Likewise, multiplying the Haar on the noncompact real group by ν changes its formal degree to d(τ)/ν, and its local test to ν/d(τ). Global multiplicity norms, mass prefactors, finite local tests, and completed or partial L-factors must retain their separately fixed conventions. Nothing in FD20 provides the literal terminal restricted period or its integral image.

\subsection{6. Sources and verification}\label{sources-and-verification-1}

\begin{itemize}
\tightlist
\item
  Nakahama, R. (2016). Norm computation and analytic continuation of vector valued holomorphic discrete series representations. \emph{Journal of Lie Theory, 26}, 927--990. Equation(5.7), p.960; Theorem5.1, pp.960--961; the normalization of constants is specified with the invariant norm in Section3. \href{https://jolt.centre-mersenne.org/item/10.5802/jolt.919.pdf}{Published primary PDF}.
\item
  van Haastrecht, R., \& Zhang, G. (2026). Wehrl inequalities for matrix coefficients of holomorphic discrete series. \emph{The Journal of Geometric Analysis, 36}, Article324. Section4.1, equations(11)--(13), Lemma4.1, and Proposition4.3, PDFpp.10--16. \url{https://doi.org/10.1007/s12220-026-02573-z}. Theorem4.4 is explicitly not used for the reason above.
\item
  Harish-Chandra. (1956). Representations of semisimple Lie groups. VI. Integrable and square-integrable representations. \emph{American Journal of Mathematics, 78}(3), 564--628. \url{https://doi.org/10.2307/2372674}. Bibliographic attribution only; original theorem text not independently retrieved in this check.
\item
  Settled local inputs: K043(PB1--PB4), K044(EG1--EG5), K046(BL1--BL7), R030(R30.1--R30.10), and LIFT31's unit-tensor real test. Their proofs are not reclassified as new here.
\end{itemize}

Exact diagnostic certificate: \nolinkurl{definite_oscillator_actual_real_formal_degree_certificate.py} and its JSON output. It checks the two Weyl products, the closed actual-profile factorization, scalar beta calibration, determinant/dual invariance, and actual-rank gap bounds using exact rational arithmetic. These are normalization diagnostics, not a numerical automorphic evaluation.

\clearpage

\section{AM-K050 --- bounded nonvanishing of the genuine e0 bottom space}\label{am-k050-bounded-nonvanishing-of-the-genuine-e0-bottom-space}

Author: GPT-6 Astra Pro, ALP\_MACHINE /root/kummer\_extension. 22 November 2026 UTC. Inherited reasoning effort is not separately exposed.

\textbf{Status: author whole-proof check PASS; exact finite certificate PASS.} Root and independent reviews are separate events.

This finding concerns either actual K031/K032 host, every role assignment and word in K044, every permitted canonical twist, and every allowed gap ell in {[}ell\_dagger,ell\_+{]}, where ell\_dagger=22,578,023,748 for J14 and 587,937,122,584 for J56. The first value is the complete J14 lower endpoint. It constructs a nonzero cusp form in the genuine source e0 bottom space at these bounded parameters. The finite level is exactly R027's bottom level: maximal at T and29, principal then generic at the two G primes. No new refinement is introduced.

The construction uses the minimal compact type at the standard origin. It does not interpolate the selected jets at the translated packet, prove a central GGP value, or show that a terminal period is nonzero on the earlier coupled image. The small translated-point bound K049 and the full-packet dimension obstruction DO044 are compatible with this existence statement.

\subsection{1. Fixed groups and the standard compact stabilizer}\label{fixed-groups-and-the-standard-compact-stabilizer}

Write E/F for the actual host, G=U(I\_a,-1), n=a+1 and a=f=42 or168. Let Jmax be the stabilizer of O\_E\^{}n at every finite place. Let J1 be the unchanged principal congruence level at all G places of R025/R027. Thus

\[
 \mathcal G=J^{\max}/J^1
   =\prod_{k,v}GL_n(\mathbf F_{g_k}),\qquad
 \Gamma=G(F)\cap J^{\max},\quad
 \Gamma_1=G(F)\cap J^1.
 \tag{BE1}
\]

All entries of Gamma are algebraic integers. Let K\_infty be the stabilizer of the standard negative line at every real label. Then

\[
 S=\Gamma\cap K_\infty
       =(\mu_{28}^{\,a}\rtimes S_a)\times\mu_{28},
 \qquad |S|=28^{a+1}a!.
 \tag{BE2}
\]

The first factor acts by monomial matrices on the positive coordinates; the last factor acts on the negative line. Its diagonal scalar subgroup Z is the actual global mu28.

\textbf{Proof.} Fixing the negative line and preserving its orthogonal complement makes an element block diagonal diag(A,t). Every embedding of each entry of A has absolute value at most1, because A is unitary at every real place. If such an entry is nonzero, its absolute norm is a positive integer at most1. Thus every conjugate has absolute value1, and the established root-unit argument gives an element of mu(E)=mu28. The column norm identity then permits exactly one nonzero entry in each column and row. The same root-unit argument gives t in mu28. Conversely all displayed matrices preserve the lattice and form. QED.

The reduction of S into Gcal is injective. At any G place the roots of unity of order28 retain their exact order, since g is prime to28. The monomial permutation is visible in the zero pattern. In particular an element of S which reduces to a scalar at one G place was already a global scalar. Also g\textgreater a and g is prime to \textbar S\textbar, so every such residue element is semisimple.

For r\textgreater0 let C\_r be the product of the Cartan sets K exp({[}0,r))K. Then

\[
                 \Gamma\cap C_{1/4}C_{1/4}^{-1}=S.
 \tag{BE3}
\]

\textbf{Proof.} Any element on the left has Cartan displacement less than1/2 at every label. The positive part of its last column has Euclidean norm sinh(u)\textless sinh(1/2)\textless1. Consequently every entry above the last diagonal has all conjugates of absolute value less than1. An algebraic integer with that property is zero. Preservation of the Hermitian form then makes the element fix the negative line, so it belongs to S. The converse follows from S being contained in K. The elementary inequality sinh(1/2)\textless1 follows from exp(1/2)\textless2. QED.

This geometric statement is for the standard maximal lattice. It makes no assertion of separation at the nonintegral T-translated packet.

\subsection{2. A genuine finite constituent contains the required stabilizer type}\label{a-genuine-finite-constituent-contains-the-required-stabilizer-type}

Let pi\_infty be the actual prescribed bottom discrete series, including the contragredient convention of R027. Let V be its minimal compact type and d\_V=dim V. Its central character on the global roots is omega\_infty. Put Omega=omega\_infty\^{}\{-1\}.

DO043 supplies an actual irreducible product representation Pi of Gcal inside the literal goodTuple/e0 sum, with central character Omega on Z. Its construction retains at least one place at which the semisimple polynomial is irreducible of degree n.~The finite constituent there is the cuspidal representation attached to the regular character of the elliptic torus F\_(g\^{}n)\^{}times. All other places retain their literal goodTuple factors and restrictions.

For every s in S,

\[
 \chi_\Pi(s)=
 \begin{cases}
       (\dim\Pi)\Omega(s),&s\in Z,\\
       0,&s\notin Z.
 \end{cases}
 \tag{BE4}
\]

\textbf{Proof.} At a full-degree place the residue of s is semisimple and has the negative-line eigenvalue t in F\_g\^{}times. A semisimple matrix conjugate into the elliptic torus can have an eigenvalue in F\_g only when it is scalar: in the multiplication representation of F\_(g\^{}n)\^{}times its eigenvalues form one Frobenius orbit, with the appropriate repetition. A Frobenius orbit containing an F\_g element has length1. By the faithful root and monomial reduction above, a nonscalar global s stays nonscalar at this place.

The Deligne--Lusztig semisimple character formula therefore makes the cuspidal character zero at this residue. The relevant formula is Deligne--Lusztig(1976), Theorem4.2, printed124: if the semisimple element has no conjugate in the defining torus, its character sum is empty. The overall sign of the elliptic virtual character has no effect on vanishing. This is the same finite-torus realization fixed in DO043. One zero factor makes the product character vanish. At a scalar the central character gives the first line. QED.

Set R=Pi\^{}dual. Character orthogonality and the exact central compatibility give

\[
 \dim\operatorname{Hom}_S(V,R)
    =\frac{|Z|}{|S|}\,d_V\dim\Pi>0.
 \tag{BE5}
\]

Indeed both V and R have central character omega\_infty; the character sum has no contribution outside Z. There is no need to embed the full V into R. Choose any nonzero intertwiner

\[
                T:V\longrightarrow R,\qquad \|T\|_{\rm HS}=1.
 \tag{BE6}
\]

Then its operator norm is at most1. This normalizes a finite linear map which exists by (BE5); it does not assume that every minimal-type vector survives a fixed finite projection.

\subsection{3. The minimal-type seed and its exact radial moments}\label{the-minimal-type-seed-and-its-exact-radial-moments}

Normalize the unitary model of pi\_infty as in K043/K046, with minimal type inclusion and orthogonal projection pr\_V. Choose a unit vector u in V. Define the R-valued function

\[
            f_u(g)=T\,\operatorname{pr}_V\pi_\infty(g)u.
 \tag{BE7}
\]

It has the exact covariance f\_u(sg)=R(s)f\_u(g) for s in S, and b(g)=\textbar\textbar f\_u(g)\textbar\textbar{} is at most1. All formulas below also hold in the contragredient model by conjugation.

Use the actual R030 Haar, with compact volume1 and ball measure pi\^{}(-a)Y\^{}(-a-1)dV at each label, where Y=1-\textbar z\textbar\^{}2. Put

\[
 r_{\tau,i}=p_{\tau,i}-p_{\tau,-}+1,\quad
 C=\prod_{\tau,i}(r_{\tau,i}-1),\quad
 s_0=a+1,\quad t=\ell/2.
 \tag{BE8}
\]

These are the same actual r\_i as K046(BL13), and r\_i-1\textgreater=2ell. The minimal-type case has all Pieri exponents m\_i=0. Consequently

\[
 I:=\int b(g)^2\,dg=\frac1{C d_V},
 \qquad
 \mathbf E(Y_\tau^{-j})
       =\prod_{i=0}^{a-1}
            \frac{r_{\tau,i}-1}{r_{\tau,i}-1-j}
 \quad(0\le j<\min_i(r_{\tau,i}-1)).
 \tag{BE9}
\]

Here the expectation uses the probability measure b²dg/I. Joint moments at distinct labels multiply, and in particular

\[
          \mathbf E(Y_\tau^{-j})\le(4/3)^a
                         \qquad(0\le j\le t).
 \tag{BE10}
\]

\textbf{Proof.} Before applying T, these are K046(BL8--10) with m=0, whose factors telescope to the ratio in (BE9). The convention g rather than g\^{}\{-1\} does not alter the integrals, since Haar and all Y\_tau are invariant under inversion.

It remains to check the finite map rather than assume it is an isometry. For any function of the radial coordinates, left-K averaging replaces T\^{}*T by its compact Schur average (\textbar\textbar T\textbar\textbar\_HS²/d\_V)Id\_V. The tensor-product minimal compact type is irreducible for the product K. Therefore every squared radial integral is exactly 1/d\_V times its value before T. Its probability ratios are unchanged. Since r\_i-1\textgreater=2ell and j\textless=ell/2, each ratio is at most4/3. This proves all assertions, including entangled vectors and the actual finite intertwiner. QED.

\subsection{4. An effective L1 concentration bound at radius1/4}\label{an-effective-l1-concentration-bound-at-radius14}

The actual profiles satisfy C\textless2\^{}(960af). The Weyl dimension formula gives the crude but explicit bound

\[
 d_V<2^{963 f a(a-1)/2}.
 \tag{BE11}
\]

Indeed each compact highest-weight coordinate has absolute value less than2A, with A=2\^{}960. Each positive-root numerator is smaller than 8A=2\^{}963; its denominator is at least1. The negative U(1) factor has dimension1. This argument applies separately at every label.

Put

\[
 B=960af+963f a(a-1)/2+af+s_0.
 \tag{BE12}
\]

Define ell\_dagger to be the larger of the actual K044 lower endpoint and the least multiple of28 strictly greater than256B. The exact values are:

\begin{longtable}[]{@{}lrrr@{}}
\toprule\noalign{}
Host & B & 256B & ell\_dagger \\
\midrule\noalign{}
\endhead
\bottomrule\noalign{}
\endlastfoot
J14 & 36,519,253 & 9,348,928,768 & 22,578,023,748 \\
J56 & 2,296,629,385 & 587,937,122,560 & 587,937,122,584 \\
\end{longtable}

Both values lie in the actual K044 interval. Thus every allowed ell\textgreater=ell\_dagger satisfies256B\textless ell. All the chamber, dimension and C bounds above hold throughout that interval: each individual compact weight, Weyl numerator and C factor is affine in ell,q, so checking its factorwise bound at the K044 endpoint pairs suffices. Also ell is even and t\textgreater s0. Let W=product\_tau Y\_tau\^{}s0. K046(BL18) gives integral W=(a!)\^{}\{-f\}\textless=1. Cauchy--Schwarz and (BE9--10) imply b is in L1, while b\textgreater=b² gives integral b\textgreater=I.

For one label, and y0=16/17,

\[
 \frac{\int_{\{Y_\tau<y_0\}}b}{\int b}
 \le
 \left(
 C d_V(a!)^{-f} y_0^{t-s_0}
 M_\tau(t)\prod_{\tau'\ne\tau}M_{\tau'}(s_0)
 \right)^{1/2}
 <\exp(-\ell/128).
 \tag{BE13}
\]

The first inequality is exactly the weighted Cauchy--Schwarz argument of K046(BL21), now using I=1/(Cd\_V). To verify the last inequality, use log(17/16)\textgreater=1/17, log2\textless1, log(4/3)\textless1 and (BE11). Its logarithm is at most (960af+963fa(a-1)/2+af-(ell/2-s0)/17)/2. The inequality 256B\textless ell bounds the positive part by ell/512, and 1/68-1/512\textgreater1/128.

Since cosh(1/4)\textgreater=1+1/32, its square is at least1089/1024\textgreater17/16. Thus u\_tau\textgreater=1/4 implies Y\_tau\textless16/17. Summing (BE13) gives

\[
   \int_{G_\infty\setminus C_{1/4}}b
      <\delta\int b,\qquad
      \delta=f\exp(-\ell/128)<1/4.
 \tag{BE14}
\]

All estimates are uniform in the role assignment, word and permitted canonical twist in the displayed interval. No at-most-half-top restriction is needed for this minimal-type statement.

\subsection{5. Poincare nonvanishing with the actual finite stabilizer}\label{poincare-nonvanishing-with-the-actual-finite-stabilizer}

Extend R from Gcal to Gamma by reduction. Consider

\[
       P_u(g)=\frac1{|S|}
             \sum_{\gamma\in\Gamma}
                     R(\gamma)^{-1}f_u(\gamma g).
 \tag{BE15}
\]

Its norm is automorphic under Gamma, and P\_u(delta g)=R(delta)P\_u(g). The finite-dimensional representation R is unitary. The seed is S-covariant, so each summand indexed by S equals f\_u(g); the normalizing factor makes their total exactly f\_u(g).

Here is the precise nonvanishing argument, with finite stabilizers retained. Set C=C\_(1/4). Each translate gamma C with gamma outside S is disjoint from C by (BE3). At any point of G\_infty the family of all translates gamma C has multiplicity at most \textbar S\textbar: if gamma1 c1=gamma2 c2, then gamma2\^{}\{-1\}gamma1 lies in CC\^{}\{-1\} intersect Gamma=S. Therefore triangle inequality and a change of variables give

\[
 \begin{split}
 \int_C\|P_u(g)\|\,dg
 &\ge \int_C b-
    \frac1{|S|}\sum_{\gamma\notin S}\int_C b(\gamma g)\,dg\\
 &\ge (1-2\delta)\int b>0.
 \end{split}
 \tag{BE16}
\]

Absolute almost-everywhere convergence follows by the usual L1 periodization/Tonelli argument, with the finite matrix coefficients bounded by1. The image of C in Gamma\G has multiplicity \textbar S\textbar; if desired (BE16) yields \textbar\textbar P\_u\textbar\textbar\_(L1(Gamma\G)) \textgreater=(1-2delta)\textbar\textbar f\_u\textbar\textbar\_1/\textbar S\textbar. In particular P\_u is nonzero.

The functions in (BE7) are finite sums of matrix coefficients of the prescribed integrable discrete series. Its actual noncompact HC gaps are at least2ell+q\textgreater a, so the same integrability criterion and integrable-discrete-series Poincare theorem used in K046 apply. To avoid any new torsion convention, decompose the sum in (BE15) over Gamma/Gamma1. It is a finite linear combination of Gamma1 Poincare sums, left translated by rational normalizers of Gamma1. Those sums are smooth, L2 and cuspidal by the primary Muic theorem already audited in K046; rational normalizer translation preserves all these properties. Thus P\_u is a nonzero cusp section. Right real translation intertwines the original pi\_infty action. The nonzero image therefore has exactly the desired real discrete-series type.

\subsection{6. Actual finite Pi type and the genuine bottom space}\label{actual-finite-pi-type-and-the-genuine-bottom-space}

Pair the R=Pi\^{}dual-valued function with Pi. On the standard coarse finite component define, for v in Pi and k in Jmax,

\[
          F_{u,v}(g_\infty,k)
                    =P_u(g_\infty)\bigl(\Pi(k)v\bigr).
 \tag{BE17}
\]

The Gamma covariance from (BE15) makes this well defined under global left invariance: R(gamma) is the dual of Pi(gamma), so its action cancels Pi(gamma) on the second argument. Right J1 acts trivially. Extend by zero to the other coarse Jmax components. This is the usual finite coefficient realization of the standard coarse component, including all principal-level components lying over it; it does not assume Gamma/Gamma1=Gcal.

Some v gives a nonzero scalar form because P\_u is nonzero. The map Pi→the actual cusp space obtained by varying v intertwines the full right Gcal action. Irreducibility makes it injective. Consequently the genuine central block e\_Pi acts as the identity, and so does the literal full e0 containing it.

Every factor of Pi is the genuine generic constituent in its regular finite series. Extend the actual nondegenerate character of the embedded U\_a to a nondegenerate character of U\_n.~Its finite Whittaker vector is nonzero and restricts to the required U\_a character. Therefore Pi has a nonzero Ucal,psi subspace. Choosing v there gives an actual nonzero form with the bottom J\_B character and the genuine e0 condition of R027. The orientation is fixed by using exactly its psi in this extension, not an unspecified inverse.

We have proved:

\subsubsection{Theorem BE18}\label{theorem-be18}

For either actual K031/K032 host, for every K044 role assignment, word and permitted canonical twist at an allowed gap ell\_dagger\textless=ell\textless=ell\_+, the prescribed genuine-e0 bottom cuspidal space is nonzero at the unchanged finite levels. Its real representation is the prescribed holomorphic discrete series or its source contragredient, as appropriate.

The theorem removes the bounded-bottom-existence qualification for this explicit bounded interval in DO046. It covers every K044 gap for J14 and the stated upper subinterval for J56. It does not select a pair with nonzero central GGP value, establish the prescribed local T functional, or give a nonzero terminal period on the earlier image. It supplies neither a norm for a packet-interpolation inverse nor an integral detector denominator. K049 continues to give a large necessary normalized selected-jet cost at its actual translated point under its separate role-count hypothesis.

\subsection{7. Source and verification record}\label{source-and-verification-record-1}

The actual field/root group and all-place lattice are K030--K032 and R025. The exact finite projector and its root/character convention are R026 and DO043; the source-shaped bottom group and central character are R027. The unitary model, actual matrix-unit metric and minimal-type moments are K043/K046, ultimately Nakahama(2016), Theorem5.1. The same integrable-discrete-series Poincare theorem is Muic(2019), Theorems6-4 and7-6(i), with the primary passages directly audited in K046. The finite-stabilizer variant (BE16) is proved here, rather than silently applying a criterion which assumes trivial stabilizer. The finite elliptic semisimple character statement is Deligne and Lusztig(1976), Theorem4.2, printed124.

The accompanying exact certificate checks the actual minimal weight bounds at both affine interval endpoints, radial inequalities, stabilizer orders and prime conditions, dimension exponent and exact concentration thresholds. It cannot compute a global cusp form or prove Poincare nonvanishing by finite sampling; those statements follow from the written representation and integral argument.

\subsubsection{Primary finite-character reference}\label{primary-finite-character-reference}

Deligne, P., \& Lusztig, G. (1976). Representations of reductive groups over finite fields. \emph{Annals of Mathematics, 103}(1),103--161. \url{https://doi.org/10.2307/1971021} . Primary author archive: \url{https://publications.ias.edu/sites/default/files/Number27.pdf} . Theorem4.2, printed124, was directly read in this run; Proposition7.4, printed141, and Theorem8.3, printed147, give the irreducible/cuspidal elliptic constituent used here. The full-degree root condition makes the torus character regular. Genericity of the chosen genuine finite series remains the already audited R026 theorem.

The remaining primary analytic inputs, with the already checked K046 normalizations, are:

\begin{itemize}
\tightlist
\item
  Nakahama, R. (2016). Norm computation and analytic continuation of vector valued holomorphic discrete series representations. \emph{Journal of Lie Theory, 26},927--990. Theorem5.1, pp.960--961, and metric(5.7). \url{https://jolt.centre-mersenne.org/item/10.5802/jolt.919.pdf} .
\item
  Muic, G. (2019). Smooth cuspidal automorphic forms and integrable discrete series. \emph{Mathematische Zeitschrift, 292},895--922. Author version arXiv:1610.05483v2, Theorem6-4, pp.21--22, and Theorem7-6(i), p.26. \url{https://arxiv.org/pdf/1610.05483} .
\end{itemize}

\subsection{Author verification record}\label{author-verification-record-3}

Author whole saved report BE1--18 and complete certificate read and reconstructed PASS at 2026-11-22T05:26:06.149742+00:00. The review covers the exact maximal-lattice stabilizer, reduction and elliptic finite-character vanishing, all dual/central signs, Hilbert--Schmidt seed normalization and left-compact Schur averaging, degree-zero moments and exact full-interval thresholds, finite-stabilizer translate multiplicities, principal-level cuspidal Poincare reduction, and the genuine finite Pi/e0/generic bottom realization. The final certificate replay passed1177 exact checks at2026-11-22T05:24:24.282126+00:00. Direct Deligne--Lusztig primary passages were checked, and previously audited K043/K046 analytic theorems are used with their exact conventions. No actual global cusp Fourier expansion, selected packet interpolation, central L-value, detector denominator or terminal period was computed or asserted.

\clearpage

\section{AM-K051 --- genuine e0 bottom nonvanishing throughout the full K044 interval}\label{am-k051-genuine-e0-bottom-nonvanishing-throughout-the-full-k044-interval}

Author: GPT-6 Astra Pro, ALP\_MACHINE /root/kummer\_extension. 22 November 2026 UTC. Inherited reasoning effort is not separately exposed.

Status: author whole-proof check PASS; exact finite certificate PASS. Root and independent reviews are separate events.

This is a short extension of frozen K050, whose report hash is c12ad0489beff63187ef1f0ae598ee238f277851b8cbb99a0ec25b2eb11a43b5. It changes only the radius and negative-moment parameter in its nonvanishing argument. In particular it preserves the actual finite constituent, Hilbert--Schmidt normalization, group levels and contragredient conventions.

The resulting theorem covers \textbf{every allowed K044 gap for both hosts}, including the previous uncovered lower part of the J56 interval. It is still an existence theorem for the genuine-e0 bottom cusp space, not a packet-interpolation or terminal-period theorem.

\subsection{1. The larger standard neighborhood has exactly the same stabilizer}\label{the-larger-standard-neighborhood-has-exactly-the-same-stabilizer}

Retain K050's Gamma, S and C\_r. Replace its radius1/4 by r=2/5. Then

\[
             \Gamma\cap C_{2/5}C_{2/5}^{-1}=S.
 \tag{FI1}
\]

The same integral-entry proof applies once sinh(4/5)\textless1 is verified. Here is an exact elementary check. For j\textgreater=2, j!\textgreater=2·3\^{}(j−2). Consequently

\[
 e^{2/5}\le1+\frac25+
       \frac{(2/5)^2}{2(1-(2/5)/3)}
       =\frac{97}{65}<\frac32.
 \tag{FI2}
\]

It follows that e\^{}(4/5)\textless9/4 and e\^{}(−4/5)\textgreater4/9, whence

\[
              \sinh(4/5)<\frac12\left(\frac94-\frac49\right)
                         =\frac{65}{72}<1.
 \tag{FI3}
\]

An element of Gamma in the indicated product neighborhood therefore has every positive last-column entry of absolute value less than1 at every embedding. Those integral entries vanish, giving exactly S as in K050(BE2--BE3).

Moreover cosh(2/5)\textgreater=1+(2/5)\^{}2/2=27/25, and

\[
                  (27/25)^2=\frac{729}{625}>\frac87.
 \tag{FI4}
\]

Thus the complement of C\_(2/5) is contained in the union of the events Y\_tau\textless7/8.

\subsection{2. The stronger moment parameter and complete lower endpoints}\label{the-stronger-moment-parameter-and-complete-lower-endpoints}

Keep K050's HS-unit T, unit u, seed b=\textbar\textbar f\_u\textbar\textbar, and I=1/(C d\_V). Set

\[
                 t=\ell,\qquad y_0=7/8,\qquad s_0=a+1.
 \tag{FI5}
\]

For every actual allowed gap, t\textgreater s0 and r\_i−1\textgreater=2ell, so t is strictly inside the negative-moment convergence range. K050(BE9) gives

\[
        M_\tau(j)=\prod_i\frac{r_i-1}{r_i-1-j}\le2^a
                        \qquad(0\le j\le\ell).
 \tag{FI6}
\]

Let B be exactly K050(BE12), so

\[
 B=960af+963f a(a-1)/2+af+s_0.
\]

Its two values and the actual lower-gap comparisons are

\begin{longtable}[]{@{}lrrr@{}}
\toprule\noalign{}
Host & B & ell\_min & ell\_min−8B \\
\midrule\noalign{}
\endhead
\bottomrule\noalign{}
\endlastfoot
J14 & 36,519,253 & 22,578,023,748 & 22,285,869,724 \\
J56 & 2,296,629,385 & 20,265,168,952 & 1,892,133,872 \\
\end{longtable}

For a single label, the same weighted Cauchy--Schwarz proof as K050(BE13), now with (FI5--FI6), yields

\[
 \frac{\int_{\{Y_\tau<7/8\}}b}{\int b}
 \le
 \left(Cd_V(a!)^{-f}(7/8)^{\ell-s_0}2^{af}\right)^{1/2}
 <\exp\left(-\frac{\ell-8B}{16}\right).
 \tag{FI7}
\]

Indeed the logarithm of the middle expression is at most \[
 \frac12\left((B-s_0)\log2-(\ell-s_0)\log(8/7)\right)
 <\frac12\left(B-s_0-\frac{\ell-s_0}{8}\right)
 =-\frac{\ell-8B+7s_0}{16}
 <-\frac{\ell-8B}{16}.
\] The elementary bound log(8/7)\textgreater=1/8 follows by integrating 1/(1+x) on {[}0,1/7{]}. All previous C and d\_V bounds are factorwise uniform throughout K044, as checked in K050.

Summing over labels and using (FI4) gives

\[
 \int_{G_\infty\setminus C_{2/5}}b
    <\delta_{\rm full}\int b,\qquad
 \delta_{\rm full}
    = f\exp\left(-\frac{\ell-8B}{16}\right)
    <\exp(-2^{26})<\frac14.
 \tag{FI8}
\]

For the explicit last bound it suffices to use log f\textless f and the two table entries: (ell\_min−8B)/16−f exceeds2\^{}26 for both hosts. There is no asymptotic or numerical floating-point approximation.

\subsection{3. Full-interval theorem and exact scope}\label{full-interval-theorem-and-exact-scope}

Insert (FI1) and (FI8) into the finite-stabilizer Poincare argument K050(BE15--BE17). The identity stabilizer still contributes exactly the seed, and the total nonidentity tail has integral at most delta\_full times its full L1 mass. The periodized section is nonzero, smooth and cuspidal by the same principal-level decomposition.

The finite Pi type and its actual Ucal,psi subspace are unchanged. Thus the literal full e0 acts as the identity on the resulting nonzero bottom vector.

\subsubsection{Theorem FI9}\label{theorem-fi9}

For either actual K031/K032 host, every role assignment and word, every allowed gap in the \textbf{complete K044 interval}, and every allowed canonical twist, the genuine source e0 bottom cuspidal space of R027 is nonzero at its unchanged G/T/29 and other finite levels.

This removes the bounded-bottom-existence qualification in DO046 throughout this entire numerical family. It does not imply a nonzero central GGP value for any prescribed initial--bottom pair. It does not show that the translated selected-jet map detects the earlier coupled image, nor give a norm or denominator for an actual terminal detector. The K049 upper-gap translated-point estimate and DO044 full-packet dimension obstruction have their separate stated scopes and are not contradicted by nonvanishing of the bottom space.

\subsection{4. Verification scope}\label{verification-scope-1}

The only changes to K050 are the elementary exact radius, the choice t=ell, and the resulting sharper concentration comparison. The companion certificate independently recalculates the two B values, complete lower endpoints, all actual moment-factor inequalities at the affine endpoint pairs, and rational exponential/cosh margins. No global Poincare series, cusp-form basis, finite character table, or period is numerically computed.

All representation-theoretic and arithmetic inputs, including the primary source locators, are precisely those recorded in K050.

\subsection{Author verification record}\label{author-verification-record-4}

Author whole saved FI1--FI9 and complete certificate read and reconstructed PASS at 2026-11-22T05:30:36.535208+00:00. The exact radius2/5 stabilizer criterion, elementary exponential and cosh bounds, moment convergence at t=ell, uniform dimension/measure factors, complete lower-endpoint margins and reuse of the unchanged genuine finite-type Poincare construction were checked. The certificate passed84 exact assertions at2026-11-22T05:30:19.074209+00:00. No period, selected packet inverse or new finite family is asserted.

\clearpage

\section{ALP-LIFT-38: a quantitative preimage bound for the actual vertical map}\label{alp-lift-38-a-quantitative-preimage-bound-for-the-actual-vertical-map}

Author: /root/lift\_accounting, GPT-6 Astra Pro, ALP\_MACHINE. 22 November 2026 UTC. Inherited reasoning effort is not separately exposed.

\textbf{Author whole-proof check completed; final verification details are recorded below.}

This report assembles the actual fixed G/Q/basic kernels of LIFT31, the explicit boundary-value estimate LIFT36, the exact G-place Hermitian comparison LIFT37, and DO047's exact Q-place comparison. Its result concerns the genuine analytic downward map and its least-norm right inverse. It does not assert an integral preimage or identify the remaining arithmetic source pins.

The estimate retains the actual source real-intertwiner scales. There cannot be a numerical bound uniform in arbitrary nonzero rescalings of those intertwiners. A separate specialization uses isometric analytic intertwiners, for which their scales are fixed to one.

\subsection{1. Quantitative inverse theorem}\label{quantitative-inverse-theorem}

Use LIFT31's actual ranks r=a+1, b=r+t, N=r+2t and its full fixed-level bottom cusp multiplicity space C0. Let D:C2-\textgreater C0 be its actual downward composition and U=D* the prescribed conjugated upward composition. Use finite group Haar with maximal compact volume1, the R030 real ball Haar at G0, and probability real Haar at the compact groups G1,G2. Let m\_b,m\_N be the full definite masses for these same measures.

Write d\_tau for the first real discrete-series formal degree and d1\_tau=dim(tau1\_tau). If s0,s1 are the actual real intertwiners relative to isometric selected oscillator contraction, put

\[
 \mathscr R_\infty=
       \frac{1}{|s_0s_1|^2}\prod_\tau(d_\tau d1_\tau)^{-1}.
 \tag{VI1}
\]

For the literal finite tuple Q define the exact raw products

\[
 \begin{split}
 A_G&=\prod_{v\in G}q_v^{-(b+r)^2}
            \frac{\kappa_t(q_v)^2}{\kappa_b(q_v)\kappa_r(q_v)},\\
 A_Q&=\prod_{v\in Q}q_v^{-4rt}
            \frac{\kappa_r(q_v)^3\kappa_{2t}(q_v)}{\kappa_b(q_v)}.
 \end{split}                                            \tag{VI2}
\]

The empty Q product is1. The uniform bound, proved in the following sections, is

\[
 \|Ux\|^2>B_{\rm raw}\|x\|^2\quad(x\ne0),\qquad
 B_{\rm raw}=m_bm_N\mathscr R_\infty A_GA_Q\,2^{-(C+1)},
 \tag{VI3}
\]

where C=277221 for J14 and C=1171821 for J56. Consequently the least-norm right inverse R=U(U*U)\^{}(-1) has

\[
                       \|R\|\le B_{\rm raw}^{-1/2}.
 \tag{VI4}
\]

The strict inequality in(VI3) is stated only for a nonzero vector. If C0=0 the right inverse is zero and the asserted upper bound is automatic. All symbols in(VI1)--(VI4) will retain their stated measures and literal finite kernels throughout the assembly.

\subsection{2. Real intertwiners and the two archimedean Gram factors}\label{real-intertwiners-and-the-two-archimedean-gram-factors}

Use the genuine unitary real types tau0,tau1,tau2 of R015 and LIFT31. Their oscillator Hilbert summands are tau\_j dual tensor tau\_(j+1). Choose an isometric identification of each selected summand. Its canonical contraction sends the unit tensor v\_bar tensor h, evaluated on v, to h. The prescribed source real upward intertwiner is a nonzero scalar s\_j times this contraction. This defines \textbar s\_j\textbar{} intrinsically relative to the actual Hilbert metrics. If the real intertwiner is factored over the labels, s\_j is the product of those scalars; no factorization is needed to define it globally.

These scalars do not depend on the finite automorphic constituent. The archimedean types and intertwiners have already been fixed. The selected contraction has multiplicity one by the compact-partner unitary Howe correspondence. LIFT31 section7 also proves that the compact-finite unit tensors used to test it lie in the actual polynomial Fock module, so the following tests are permitted tests of the same source multiplicity map.

At the first edge choose unit v in tau0 and unit h in tau1 and use v\_bar tensor h. Discrete-series Schur orthogonality for the actual real Haar gives

\[
 Z_{0,\tau}=\int_{U(a,1)}
       |\langle v,\tau_0(g)v\rangle|^2dg=d_\tau^{-1}.
 \tag{VI5}
\]

At the second edge use h\_bar tensor u with unit u in tau2. Compact Schur orthogonality, with probability Haar, gives

\[
                         Z_{1,\tau}=d1_\tau^{-1}.
 \tag{VI6}
\]

Equations(VI5)--(VI6) are LIFT31(VA28--29) with their measures unchanged. The scalar theta lift of these tests equals the multiplicity lift evaluated on s\_j times the target unit vector. Dividing its squared norm by \textbar s\_j\textbar\^{}2 therefore gives the actual multiplicity-map norm. This proves exactly the factor(VI1). In particular there is no additional coevaluation or dimension factor besides the compact Schur factor already in(VI6).

\subsection{3. Global measure conversion without double counting}\label{global-measure-conversion-without-double-counting}

Let mu\_j be the stated actual global group Haar and let mu\_j=a\_j mu\_j\^{}Tam. GQT section2.5 gives Tamagawa number2 for the unitary groups. The two targets are compact at infinity and their real Haar is probability, so

\[
            2a_1=m_b,\qquad 2a_2=m_N.
 \tag{VI7}
\]

At the bottom, R030 gives 2a0=m\_r/(a!)\^{}f, but that value will cancel from the formula expressed in actual local measures.

Here is the conversion on one edge. Let a globally unit vector for the actual source Petersson metric be represented as a\_j\^{}(-1/2) times a Tamagawa-unit vector. Integrating the theta kernel with mu\_j multiplies the Tamagawa theta integral by a\_j. The actual target Petersson squared norm then contributes a\_(j+1). Thus the scalar norm ratio changes from Tamagawa by a\_j a\_(j+1).

The product of all actual-source local Gram integrals, however, is a\_j times their Tamagawa product: every Gram integral has exactly one source-group integration. Consequently the factor 2 in the Tamagawa Rallis formula becomes 2a\_(j+1) when the local forms are expressed in the actual measures used in LIFT37 and DO047. For our two edges their product is

\[
                    (2a_1)(2a_2)=m_bm_N.
 \tag{VI8}
\]

Equivalently, if every local form were kept in Tamagawa measures, the two-edge map and Petersson metrics would contribute a0 a1\^{}2 a2. With finite maximal volumes1, the real Tamagawa forms instead contribute the compensating factor1/(a0 a1). Both descriptions give(VI8). One must not place a0 a1\^{}2 a2 on top of already rescaled local Gram forms.

These identities concern full adelic Petersson metrics. The finite G and Q level indices already occur in the actual raw Gram factors. They are not multiplied into(VI7)--(VI8) a second time.

\subsection{4. The exact finite contributions}\label{the-exact-finite-contributions}

At each G place, LIFT37(VG10--12) proves that the two successive raw Gram forms, with their actual first quotient metric, give the scalar in the first line of(VI2). This applies on the entire initial K\_r\^{}1-fixed space. The proof uses actual full-rank residue support, so both orbit integrals and doubled Gram integrals are compact, also at the second edge b\textless N\textless2b.

At each Q place, DO047(QN18,QN22,QN25) gives the composed raw factor

\[
 q^{-4rt}\frac{\kappa_r^3\kappa_{2t}}{\kappa_b}
       \left|\prod_{i=1}^t(1-\vartheta(\varpi)q^{-i})^{-1}\right|^2.
 \tag{VI9}
\]

The first source is spherical there. Literal Tuples.Qprime makes xi genuinely ramified, while mu0,mu1,eta are unramified. Thus its first standard local factor is1; the second split standard factor is exactly the Euler absolute square in(VI9). Dividing by those actual standard factors leaves the second line of(VI2). The proof of this assertion is the source audit and exact continuation in DO047, not a supposition that every Q place has an arbitrary test character. For trivial xi, Q is empty.

At every remaining finite place both functions are the actual normalized lattice vacua. LIFT31 section6 proves on these exact vectors the unramified formula

\[
             Z_{j,v}=L_v(z_0,\pi_j\times\beta_j)
                                  d_{j,v}(s_0)^{-1}.
 \tag{VI10}
\]

It includes dyadic split places and the absolutely ramified places above29 and7: the former are relatively split, the latter have the relative-unramified selfdual lattice and character transport checked there. There is no exceptional finite test omitted from(VI2),(VI9), and(VI10). In particular the vertical T functions are basic; the horizontal T incidence estimates do not enter this calculation.

\subsection{5. Exact Euler assembly and a uniform positive numerator}\label{exact-euler-assembly-and-a-uniform-positive-numerator}

Let G in an Euler superscript mean the removal of the E-primes over the actual finite G places. Define the finite standard factors in the actual Hom/source convention of LIFT31 by

\[
                  \mathcal N_{\pi_0}
                  =\mathcal L_0^G(z_0)\mathcal L_1^G(z_0).
 \tag{VI11}
\]

At G, no standard factor is inserted: we already use the raw Gram form. At Q, the standard Euler factor has been retained in(VI11) and removed exactly once from(VI9). At basic places(VI10) retains the standard factor in the same way. This explains precisely which partial product occurs; no Q Euler factor is discarded twice.

The unramified parameter identity LIFT31(VA18,VA27--30) and the ramified Q computation give

\[
 \mathcal N_{\pi_0}=
   P_{0,\pi_0}^{GQ}(z_0)P_{1,\pi_0}^{GQ}(z_0)
                       \prod_{i=1}^t L_E(i,\vartheta).
 \tag{VI12}
\]

Each P is the initial rank-r parameter Euler product with its stated unitary twist. It is absolutely convergent with margin859 or733. At Q the initial twisted factors are1 because xi is ramified. The inserted character vartheta is ramified at G and unramified elsewhere, so its complete string already has local factor1 at G. This is why there is no removed-prime factor R\_S left in(VI12).

LIFT36(LB24--30) consequently gives, uniformly in every initial automorphic constituent and every actual K044 word/profile,

\[
           |\mathcal N_{\pi_0}|>2^{-C},\qquad
 C=\begin{cases}277221,&J14,\\1171821,&J56.\end{cases}
 \tag{VI13}
\]

This uses only the verified unramified unitary Satake bound for the initial Euler products. The stronger all-place temperedness theorem DO046 is compatible with it but is not needed here. Nor is a global base-change identification of the large middle theta representation assumed.

\subsection{6. The basic doubling denominator has a uniform explicit bound}\label{the-basic-doubling-denominator-has-a-uniform-explicit-bound}

Write epsilon\_v for the unramified quadratic character of E\_v/F\_v. Both vertical rank increments t are odd. The actual beta restriction and the unramified unitary doubling formula give

\[
 d_{j,v}(s_0)=\prod_{k=t+1}^{n_j+t}
          (1-\epsilon_v(\varpi_v)^k q_v^{-k})^{-1},
             \qquad n_0=r,\quad n_1=b.
 \tag{VI14}
\]

One exact primary locator is Harris(2007), equations(1.3.4) and (2.1.8.3): its argument is2s+n\_j-l=t+n\_j-l and its quadratic exponent n\_j-1+l has the same parity as that argument because t is odd. This is also the d\_v of GQT section11.6. At split places(VI14) equals C\_(n\_j,n\_j+t)\^{}(-1), independently checking the same normalization against LIFT16's matrix calculation.

Set

\[
              \mathcal D=\prod_{v\notin G\cup Q}
                                      d_{0,v}(s_0)d_{1,v}(s_0).
 \tag{VI15}
\]

This is a positive absolutely convergent Euler product. For q\textgreater=2, k\textgreater=t+1 and t\textgreater=2, -log(1-q\^{}(-k))\textless=2q\^{}(-k). There are at most f primes of F above a rational prime. Enlarging both finite k-ranges to all k\textgreater=t+1 and applying the integral test gives

\[
 \begin{split}
 \log\mathcal D
 &\le4\sum_{v}\sum_{k=t+1}^{\infty}q_v^{-k}
 \le8\sum_vq_v^{-(t+1)}\\
 &\le8f\left(2^{-(t+1)}+\frac{2^{-t}}t\right)
 \le8f\,2^{-t}<\log2.
 \end{split}                                            \tag{VI16}
\]

If a quadratic factor is less than1 it only improves this upper bound. Hence

\[
                         \mathcal D^{-1}>\tfrac12.
 \tag{VI17}
\]

The exact product(VI15), rather than this deliberately weak bound, can be used throughout the result if desired.

\subsection{7. Polarization, multiplicities and the exact scalar identity}\label{polarization-multiplicities-and-the-exact-scalar-identity}

Fix an initial irreducible cuspidal constituent pi0. The first theta image lies in the actual local small-theta tensor product, and the second lies in its next small-theta tensor product. This is LIFT31's algebraic factorization. In particular it does not replace the incoming middle vector by an unrelated test vector.

The first global theta norm formula is the anisotropic convergent Siegel--Weil formula GQT Theorem7.1, followed by the doubling and unfolding in sections11.3,11.5--11.6. The first rank b\textgreater2r is outside the printed dimension window of Theorem11.4, so that theorem is not used for the first edge. The second has b\textless N\textless2b and anisotropic target. It is the convergent case of Proposition11.1(ii) and Theorem11.4. Both have positive s0=t/2 and factor2 with Tamagawa measures, as used in section3.

Rallis is a sesquilinear identity. On two different automorphic multiplicity embeddings of the same pi0, the unfolded matrix coefficient includes their Petersson multiplicity pairing. Thus orthogonal embeddings have zero cross term. On the finite principal-G type, LIFT37 polarizes to the scalar(VI2) times its original inner product. At Q the starting spherical line is one-dimensional and DO047 supplies its scalar; at all other finite places(VI10) supplies the scalar on the spherical line. The real tests(VI5)--(VI6) have already divided by the actual source scales.

Consequently on the entire initial pi0-isotypic fixed-level space, including its full multiplicity space and every prescribed principal-G vector, the actual composition satisfies

\[
 \boxed{\quad
 \langle Ux,Uy\rangle=\lambda_{\pi_0}\langle x,y\rangle,
 \quad
 \lambda_{\pi_0}=m_bm_N\mathscr R_\infty A_GA_Q
                       \mathcal N_{\pi_0}\mathcal D^{-1}.
 \quad}                                                \tag{VI18}
\]

Every factor other than the displayed standard numerator is positive. Positivity of the actual Rallis form makes that numerator positive in this identity; one may equivalently use its absolute value. Linear Q Gauss/Euler phases have not been replaced in a linear theta identity: they occur here only in the proved Hermitian forms, where they pair with their conjugates.

Distinct initial irreducibles have orthogonal terminal images. Indeed equality of terminal irreducibles would imply equality of the middle local representations at every place by local Howe injectivity, and then equality of the initial representations. Orthogonal decomposition of the cuspidal spectrum therefore allows (VI13) and(VI17) to be applied simultaneously to every initial isotypic summand. This proves(VI3), uniformly on the full C0.

LIFT31 proves the exact Petersson adjoint identity U=D\emph{. Thus U}U=DD* is positive definite on C0. Spectral calculus gives the right inverse R=U(U*U)\^{}(-1) and its norm bound(VI4). No dimension factor from summing multiplicity copies is required or inserted.

\subsection{8. Finite products, masses and an explicit coarse cost}\label{finite-products-masses-and-an-explicit-coarse-cost}

Let hQ be the number of F-places in Q and Qnorm=product\_(v in Q)q\_v. For every q\textgreater=2 the elementary estimate kappa\_j(q)\textgreater1/4 holds. For completeness, at q=2 the product of its first three factors is 21/64, and the remaining product is at least 1-sum\_(i\textgreater=4)2\^{}(-i)=7/8. Their product147/512 exceeds1/4; shorter products and larger q only increase it. Consequently

\[
            A_Q>2^{-8h_Q}\operatorname{Qnorm}^{-4rt}.
 \tag{VI19}
\]

Together with LIFT37's raw G table and(VI13),(VI17), this gives

\[
 B_{\rm raw}>
   \frac{m_bm_N}{|s_0s_1|^2\prod_\tau d_\tau d1_\tau}
       2^{-H_{\rm raw}-8h_Q}\operatorname{Qnorm}^{-4rt},
 \quad
 H_{\rm raw}=\begin{cases}
     6433884556,&J14,\\46339555606,&J56.
 \end{cases}
 \tag{VI20}
\]

All finite G constants in this bound are fixed independently of the large weights. The actual tuple contribution is explicit, including each of its Q places, rather than an unspecified bad-prime constant.

The masses can also be bounded without evaluating the full genus formula. The identity maximal-lattice class contributes 1/(28\^{}n n!) to m\_n.~Indeed the standard lattice stabilizer is the monomial group mu28\^{}n semidirect S\_n: each integral matrix entry has all complex conjugates of modulus at most1 by the unitary column identities; a nonzero such algebraic integer is a root of unity, and exactly one entry in each row and column can be nonzero. The actual roots of unity are mu28 by K040. This is the standard stabilizer already used in R017 and LIFT33, now used only to bound the full maximal-level mass. It follows that

\[
 (m_bm_N)^{-1}\le28^{b+N}b!N!
       <2^{16b+17N}
       =\begin{cases}2^{89369},&J14,\\2^{87227},&J56.
         \end{cases}
 \tag{VI21}
\]

The first inequality uses just one class, not a false identification of the entire genus with that class. The second uses b\textless2\^{}11, N\textless2\^{}12 and28\textless2\^{}5. No principal-G or T index is added to this mass.

Combining(VI4),(VI20),(VI21) gives a concrete upper cost for any prescribed real intertwiners:

\[
 \boxed{\quad
 \|R\|< |s_0s_1|
   \left(\prod_\tau d_\tau d1_\tau\right)^{1/2}
   2^{(H_{\rm raw}+H_{\rm mass}+8h_Q)/2}
                         \operatorname{Qnorm}^{2rt},
 \quad}                                                \tag{VI22}
\]

where H\_mass=89369 or87227. The exact bound(VI4) is stronger because it retains the true masses, kappa products and standard Euler values.

\subsection{9. Unit finite kernels and the unavoidable real-scale condition}\label{unit-finite-kernels-and-the-unavoidable-real-scale-condition}

For comparison, normalize both entire finite Schwartz tensors to norm1. This changes the upward composition by division by the product of the two raw finite norms. LIFT37 and DO047 give the corresponding exact finite factors

\[
 \begin{split}
 A_G^{\rm unit}&=\prod_{v\in G}q_v^{-r(b+r)}
                     \frac{\kappa_t}{\kappa_b\kappa_r},\\
 A_Q^{\rm unit}&=\prod_{v\in Q}q_v^{-rt}
                     \frac{\kappa_r\kappa_{2t}}{\kappa_b}.
 \end{split}                                            \tag{VI23}
\]

The analogue of(VI22) is

\[
 \|R_{\rm unit}\|< |s_0s_1|
   \left(\prod_\tau d_\tau d1_\tau\right)^{1/2}
   2^{(H_{\rm unit}+H_{\rm mass}+4h_Q)/2}
                         \operatorname{Qnorm}^{rt/2},
 \quad
 H_{\rm unit}=\begin{cases}150220372,&J14,\\3974376694,&J56.
          \end{cases}
 \tag{VI24}
\]

This comparison does not silently replace the literal raw source kernel; its normalization has been stated explicitly.

The factors s0,s1 cannot be omitted. Multiplying a prescribed downward real intertwiner by c multiplies its adjoint by conjugate(c). The defining source IsIo/IsIoUp identities then divide the corresponding multiplicity maps by c and conjugate(c), respectively. Its minimal right-inverse norm changes by \textbar c\textbar. Thus no uniform upper bound independent of arbitrary nonzero real-intertwiner scales can hold. Equations(VI4),(VI22),(VI24) are quantitative and scaling-covariant statements for those actual scales.

If the actual analytic construction chooses the isometric selected contractions as its real intertwiners, then s0=s1=1 by definition. The bounds specialize accordingly, with no remaining real scalar. This is an explicit analytic normalization, not an assertion that an as-yet-unmatched arithmetic source real pin has that value.

\subsection{10. Exact real degrees and a fully numerical specialization}\label{exact-real-degrees-and-a-fully-numerical-specialization}

DO048(FD5,FD8--FD12) computes the formal degree in exactly the real Haar used here. If h\_tau is the decreasing bottom Harish-Chandra list and rho\_r=((r-1)/2,\ldots,-(r-1)/2), then

\[
 d_\tau=a!\,\dim V_{h_\tau-\rho_r}
       =\dim W_{\alpha_\tau}
                    \prod_{i=0}^{a-1}(p_i-p_-).
 \tag{VI25}
\]

There is no undetermined measure constant in this identity. Nakahama(2016),(5.7) and Theorem5.1 normalize the constants in the ball Hilbert space by c\_tau=product\_i(p\_i-p\_-). Summing the scalar Schur relation over an orthonormal basis of the minimal compact type gives d\_alpha/d\_tau=1/c\_tau. Weyl's formula then gives (VI25). DO048 proves this calibration, including the compact centre and the contragredient convention, in full. The compact degree d1\_tau is the ordinary Weyl dimension for the actual middle Harish-Chandra list; unitary determinant shifts change neither degree.

Both degrees are explicitly bounded on the entire actual K044 interval, not only its largest-gap endpoint. At a CM role write E=E\_T,E\_H or E\_L, and let q\_can be the canonical twist parameter. The bottom list is

\[
 (4\ell(a-1),4\ell(a-2),\ldots,4\ell,
                    -E+\ell,-E-\ell-q_{\rm can}).
 \tag{VI26}
\]

The middle list adds the block B\_t(-E) and a common shift eta. This is literally DataTheta.Dpar at j=1: its v and b entries are unchanged, the only added block is centred at -e0=-E, and the common half-integral shift cancels from every root difference. In particular the horizontal top value A is not a new entry in this vertical compact list.

The checked bounds ell\textless2\^{}485, 0\textless=q\_can\textless=ell, E\_T\textless2\^{}955+2\^{}485, E\_H\textless2\^{}495, E\_L\textless2\^{}485, a\textless=168 and t\textless=1759 show that every positive difference in either list is less than 2\^{}960. For example their full range is bounded by E+(4a+3)ell+t+2\textless2\^{}956. Applying the root-product formula (VI25), whose denominator is product\_(j=1)\^{}(a-1) j!\textgreater=1, and the ordinary compact Weyl formula therefore gives

\[
 \prod_\tau d_\tau d1_\tau<2^{H_\infty},\qquad
 H_\infty=\frac{960f}{2}
                      \{a(a+1)+b(b-1)\}.
 \tag{VI27}
\]

The factor a! has already cancelled against the last compact Vandermonde denominator in (VI25); no second a! is inserted in (VI27). The exact closed formal-degree products are also supplied by DO048(FD15--FD16) if a sharper estimate is wanted.

\begin{longtable}[]{@{}lrrrr@{}}
\toprule\noalign{}
Host & H\_infty & H\_raw+H\_mass+H\_infty & H\_unit+H\_mass+H\_infty & rt \\
\midrule\noalign{}
\endhead
\bottomrule\noalign{}
\endlastfoot
J14 & 65,463,713,280 & 71,897,687,205 & 65,614,023,021 & 75,637 \\
J56 & 263,998,748,160 & 310,338,390,993 & 267,973,212,081 & 275,977 \\
\end{longtable}

Combining (VI22),(VI24),(VI27) gives, with the table's H\_tot and H\_tot\_unit,

\[
 \boxed{\quad
 \|R\|<|s_0s_1|\,2^{(H_{\rm tot}+8h_Q)/2}
                       \operatorname{Qnorm}^{2rt},\qquad
 \|R_{\rm unit}\|<|s_0s_1|\,2^{(H_{\rm tot,unit}+4h_Q)/2}
                       \operatorname{Qnorm}^{rt/2}.
 \quad}                                               \tag{VI28}
\]

For the explicit isometric analytic normalization s0=s1=1 these are numerical upper bounds once the actual finite tuple Q is given. If Q is empty, h\_Q=0 and Qnorm=1. Their very large sizes are retained; neither the local or global estimates nor the intertwiner scalars have been replaced by an unspecified constant.

\subsection{11. Exact remaining arithmetic distinction}\label{exact-remaining-arithmetic-distinction}

The result bounds a complex Petersson-norm preimage for every vector in the fixed bottom cusp space. It proves no 11-integral preimage, no primitive Smith minor, and no upper bound for the non-scalar arithmetic image index. LIFT34's middle-mass divisibility and the terminal arithmetic branch remain separate statements. Applying the genuine e0 projection gives a bounded complex preimage of every vector in its bottom image, but does not create a detector when that projected bottom vector is zero.

The estimate uses the actual mixed CM host and fixed vertical kernels of the checked analytic adapter. It does not repair the remaining literal Host/Fields pins or claim a complete DP44 shorter-chain comparison. What it removes is the former wholly qualitative top-preimage norm in LIFT31: the bound is now an explicit product, with its only real-normalization dependence displayed and with every finite G/Q and global boundary factor controlled.

\subsubsection{References and exact inputs}\label{references-and-exact-inputs-1}

Gan, W. T., Qiu, Y., \& Takeda, S. (2014). The regularized Siegel--Weil formula (the second term identity) and the Rallis inner product formula. \emph{Inventiones Mathematicae, 198},739--831. Sections2.5,7.1, 11.3--11.7; Theorem7.1, Proposition11.1(ii), Theorem11.4 and(11.3). Local primary files: lift\_refs/gqt2014.pdf and .txt.

Harris, M. (2007). \emph{Cohomological automorphic forms on unitary groups, II: Period relations and values of L-functions}. Equations(1.3.4),(2.1.8.3), together with their unitary local normalization. Local primary files: lift\_refs/harris2007\_periodrelations.pdf and .txt. The exact basic test is LIFT31 section6, with its Gong--Grenie/Harris primary verification.

The other carried proofs are LIFT31(VA11--33), LIFT36(LB1--30), LIFT37(VG1--17), DO047(QN1--26), R015's actual real types and R030's actual ball/definite Haar comparison. None of those accepted construction or classification results is being re-proved here.

Nakahama, R. (2016). Norm computation and analytic continuation of vector valued holomorphic discrete series representations. \emph{Journal of Lie Theory,26},927--990. Equation(5.7),p.960; Theorem5.1,pp.960--961. The full exact formal-degree application and compact-centre calibration are DO048(FD1--FD20), definite\_oscillator\_actual\_real\_formal\_degree.md. The actual middle list is DataTheta.lean, Dpar,lines538--543; its real representation identification is R015. The uniform actual weight bounds are K044(EG1--EG5).

\subsubsection{Verification record}\label{verification-record-3}

The author read the complete proof, including the two different global-measure descriptions, the restored Q factors, all multiplicity pairings, the exact source-scale transformation, and the literal Dpar middle list. The companion \nolinkurl{vertical_preimage_norm_certificate.py} checks the Haar identities, exact Q normalizations, a finite-dimensional inverse identity, all actual parity ranges and the final numerical costs. Its4561 assertions are arithmetic diagnostics; they do not compute an automorphic matrix or an integral Smith form. Independent reviews of the carried LIFT31/36/37, DO047 and DO048 inputs are separately recorded in the project check logs.

\clearpage

\section{AM-R-032 --- the genuine finite projector on the actual vertical preimage and terminal rows}\label{am-r-032-the-genuine-finite-projector-on-the-actual-vertical-preimage-and-terminal-rows}

Author: GPT-6 Astra Pro, ALP\_MACHINE /root. 22 November 2026 UTC. Reasoning effort is not separately exposed.

\textbf{Status: author whole-proof check PASS; independent review pending.}

This result fixes a previously unchosen spectral part of the constructed mixed arithmetic family. It uses the genuine full goodTuple idempotent, the unchanged vertical finite kernels, and the canonical terminal class module. It does not choose a projector to fit a theta image. It proves that the actual vertical map and its least-norm inverse preserve this same finite projector, and identifies its arithmetic and analytic coefficient actions. The integral theta inverse, the fixed full seesaw comparison, and terminal nonvanishing remain separate.

\subsection{1. Exact groups and the leading residue action}\label{exact-groups-and-the-leading-residue-action}

Use either actual K031/K032 host and all K044 parameters. Set r=a+1, b=r+t, N=r+2t, with the ranks of LIFT31. At a G-place v let K\_m=GL\_m(O\_v), K\_m\^{}1=ker(K\_m-\textgreater GL\_m(k\_v)). The actual finite levels of the three vertical spaces are the ones in LIFT31(VA10--12):

\[
 J_{0,v}=K_r^1,
\] \[
 J_{1,v}=\operatorname{red}^{-1}
 \left\{\begin{pmatrix}I_r&B\\0&D_1\end{pmatrix}\right\},
\] \[
 J_{2,v}=\operatorname{red}^{-1}
 \left\{\begin{pmatrix}I_r&B_{01}&B_{02}\\0&D_1&B_{12}\\0&0&D_2
                      \end{pmatrix}\right\},
 \qquad D_i\in GL_t(k_v).
 \tag{R32.1}
\]

Their characters are respectively 1, chi1(det D1), and chi1(det D1)chi2(det D2). All other finite levels and characters, including the Q parahorics and the maximal T levels, remain fixed. Let C\_j be the corresponding full automorphic multiplicity spaces, with the real types and Petersson metrics of LIFT31/38. Initially the leading residue subgroup is trivial; imposing the desired generic subgroup is a further character subspace inside these same C\_j.

For gamma in GL\_r(k\_v), choose a lift gamma\_tilde in K\_r and put

\[
 k_{0,v}(\gamma)=\widetilde\gamma,\quad
 k_{1,v}(\gamma)=\operatorname{diag}(\widetilde\gamma,I_t),\quad
 k_{2,v}(\gamma)=\operatorname{diag}(\widetilde\gamma,I_t,I_t).
 \tag{R32.2}
\]

These elements normalize J\_j and its character: conjugation merely changes the upper off-diagonal blocks and leaves D1,D2 unchanged. Different lifts have quotient in J\_j with character one. Hence

\[
 \rho_j(\gamma)F(g)=F(g k_j(\gamma))
 \tag{R32.3}
\]

is a well-defined unitary action of the same finite group

\[
 \mathcal G=\prod_{v\in G}GL_r(k_v)
            =\prod_{i=0,1}\prod_{v\mid g_i}GL_r(\mathbf F_{g_i})
 \tag{R32.4}
\]

on C\_j. Multiplication holds because the discrepancy between the chosen lift of gamma delta and the product of chosen lifts lies in the stated character-one subgroup. Right translation preserves the Petersson inner product. The finite action commutes with the real representation, so it acts on multiplicity spaces as well as scalar automorphic functions.

These are the literal leading-block residue actions: DefGroup.lean 620--707 defines parLev, rTup, rHom, lev and levFn; rTup is the upper left r-block. LevelP.triv at581 means the trivial residue subgroup, not the full residue group. R025 supplies the exact residue-field isomorphisms, complete G places and actual bottom reduction.

\subsection{2. Invariance of both unchanged vertical G kernels}\label{invariance-of-both-unchanged-vertical-g-kernels}

Write f0,f1 for the actual conjugated upward half-Schwartz functions. By LIFT31(VA11--13), f0 is the integral indicator with residue E\_r={[}I\_r;0{]}, whereas f1 is supported on residues

\[
 X=\begin{pmatrix}I_r&B\\0&D\\0&0\end{pmatrix},\quad
 D\in GL_t(k_v),
 \qquad f_1(X)=\chi_1(\det D)^{-1}\chi_2(\det D).
 \tag{R32.5}
\]

The full split Hom action is

\[
 [\omega_{\alpha,\beta}(g,h)f](X)=
 \alpha(\det h)\beta(\det g)^{-1}
 |\det g|^{m/2}|\det h|^{-n/2} f(h^{-1}Xg).
 \tag{R32.6}
\]

On the first upward edge alpha0/beta0=eta; on the second alpha1/beta1=eta\^{}\{-1\}, as fixed in LIFT31(VA4--5). The actual eta is unramified at all finite places by K029/K040, so these ratios are one on det(gamma\_tilde). The tuple xi cancels from each ratio. All determinants in (R32.2) are units, so the absolute-value factors are also one.

For the first function,

\[
 \operatorname{diag}(\widetilde\gamma,I_t)^{-1}
       E_r\widetilde\gamma=E_r.
 \tag{R32.7}
\]

Integral changes of variables preserve the integral lattice, and (R32.7) identifies its residue support in both directions. For the second function, simultaneous action by k1,k2 sends the residue in (R32.5) to

\[
 \begin{pmatrix}I_r&\gamma^{-1}B\\0&D\\0&0\end{pmatrix}.
 \tag{R32.8}
\]

Thus its support and its exact character value are unchanged. The same transformations are invertible, so points outside the support remain outside it. We have proved the identities of actual Schwartz vectors

\[
 \boxed{\quad
 \omega_{j,v}(k_{j,v}(\gamma),k_{j+1,v}(\gamma))f_{j,v}=f_{j,v},
                 \qquad j=0,1.
 \quad}                                                    \tag{R32.9}
\]

The half-idempotent and split Heisenberg transport have already been included in VA11--13; (R32.9) is not a replacement of the original DatumVec vector by an invariant surrogate. The action is identity at all places outside G and at infinity. Consequently the complete fixed global Schwartz tensor satisfies the same simultaneous invariance. This uses the actual unramified eta; a ramified odd splitting character would require retaining its determinant twist.

\subsection{3. Exact equivariance and the least-norm inverse}\label{exact-equivariance-and-the-least-norm-inverse}

Let U\_j:C\_j-\textgreater C\_(j+1) be the actual upward map. In the scalar theta integral, simultaneous invariance gives Theta(g k\_j,h k\_(j+1))=Theta(g,h). Changing the source integration variable by right multiplication therefore gives

\[
 U_j\rho_j(\gamma)=\rho_{j+1}(\gamma)U_j.
 \tag{R32.10}
\]

The same computation holds for the conjugated upward kernel. The fixed real intertwiners commute with these finite actions, so (R32.10) is an equality of the actual multiplicity maps, with their source scales retained. No Haar index occurs: this is a change of variable in the full quotient integral, not an unnormalized finite average.

LIFT31 proves D\_j=U\_j\^{}*. Taking adjoints in (R32.10), and replacing gamma by gamma\^{}\{-1\}, gives

\[
 D_j\rho_{j+1}(\gamma)=\rho_j(\gamma)D_j.
 \tag{R32.11}
\]

Put U=U1 U0 and D=D0 D1=U\emph{. LIFT38 proves that U}U is positive invertible on C0 and supplies its uniform lower bound. Equations (R32.10--11) show that U*U and its inverse commute with rho0. Thus the exact least-norm right inverse

\[
 R=U(U^*U)^{-1}:C_0\longrightarrow C_2
 \tag{R32.12}
\]

obeys R rho0(gamma)=rho2(gamma)R. This holds for every group-algebra element, not only for a central idempotent.

Let e be the genuine full goodTuple idempotent R026(R26.7), and set E\_j=rho\_j(e). Then

\[
 \boxed{\quad DE_2=E_0D,\quad UE_0=E_2U,\quad RE_0=E_2R.\quad}
 \tag{R32.13}
\]

These are orthogonal projections. Every x in E0 C0 therefore has the actual same-projector top preimage Rx in E2 C2, with

\[
 DRx=x,\qquad \|Rx\|\le B_{\rm raw}^{-1/2}\|x\|.
 \tag{R32.14}
\]

B\_raw is exactly LIFT38(VI3); its raw numerical bounds VI28 apply without an additional projector or dimension factor. The same is true with the explicitly normalized finite tensors. The zero-space case is immediate.

The action also intertwines the character projector of any subgroup P of Gcal. Hence if x transforms by the actual generic character psi under the embedded product diag(U\_a,1), Rx has that same leading character. At the terminal level this is precisely rowTyG with its P.H and character, rather than a refinement of the final level. The inverse/conjugate row convention is the one already fixed in DO038 and \nolinkurl{RowForms.lean221–229,253–269.} Centrality makes E\_j commute with this character projector. The group-equivariance assertion is made on C\_j before the nonnormal subgroup cut, so no action on a nonnormal character subspace is being silently assumed.

Every individual genuine block, including the 11-integral block of R029, satisfies the same intertwining identities. Such a block need not be coefficientwise real; the full e is used below for the literal same-projector/reality equations.

\subsection{4. An arithmetic character list with exactly this idempotent}\label{an-arithmetic-character-list-with-exactly-this-idempotent}

Let M be the already fixed conjugation-stable geometric value field of K036--038 for the chosen finite tuple. Let L=Q(zeta\_mG), where mG is the group exponent used in R028(R28.11), and set K=ML inside the fixed complex embedding. This is one explicitly specified finite extension. R028 proves that every finite character value lies in L and gives the retained relative degree bounds

\[
 [K:M]\le[L:Q]<
 \begin{cases}2^{43562},&J14,\\2^{1019986},&J56.\end{cases}
 \tag{R32.15}
\]

This is only the additional finite-character field cost. It is not a bound for the old geometric field M or a free change in DP44's field and height terms.

For every literal goodTuple f, take the genuine finite product representation Pi\_f of R026. Because every f\_v is squarefree, each series is a singleton; different tuples give distinct irreducible characters. Let X\_good be the finite set of their K-valued characters, with the canonical residue identification eR of R025. This character list is fixed by the finite classification and the literal IsGood predicate, independently of every global theta image or value.

The source's exact formula is

\[
 \operatorname{chiIdem}(\chi)
  ={\chi(1)\over|\mathcal G|}\sum_{g\in\mathcal G}\chi(g^{-1})[g],
 \qquad
 e_A=\operatorname{idemOf}(X_{\rm good})
      =\sum_{\chi\in X_{\rm good}}\operatorname{chiIdem}(\chi).
 \tag{R32.16}
\]

It is \nolinkurl{GoodPartHermitianSpec.lean17–24.} The characters are absolutely irreducible in the sense of that file: they are the genuine irreducible characteristic-zero finite-group characters already used in R026, whose algebraic character values can be viewed in an algebraic closure of K. This does not require assigning an arbitrary central idempotent to an uninhabited representation.

For a product group, its order, the character degree and the character value in (R32.16) are products of their individual factors. Thus chiIdem of a product character is the tensor product of its primitive idempotents. Summing over exactly the literal good tuples is exactly \nolinkurl{DataN.lean1001–1007/R026(R26.7).} Consequently

\[
 \boxed{\quad \iota_K(e_A(g))=e(eR(g))\quad(g\in\mathcal G).\quad}
 \tag{R32.17}
\]

There is no omitted character multiplicity, normalization, or changed finite series. The place bijection in grp\_record sends (i,v) to the same chosen degree-one E-prime over g\_i. The required residue isomorphism is its unique prime-field isomorphism. The divisibility clause holds because that prime lies over g\_i. The final residue matrix clause is the equality of upper-left reductions in R32.2. These verify the mathematical coefficient/residue equations of LinksChar.lean112--117 for this constructed finite part. They do not assert that the incompatible full Host/Fields record exists.

R026 proves coefficientwise reality of e. Combining it with (R32.17) gives precisely the conjugated coefficient identity RowLevelMatchC/R25.13, in the unchanged standard frames and with the same level-matching translation one. The independent T translation in the Bessel point is unchanged.

\subsection{5. The actual terminal arithmetic action and its analytification}\label{the-actual-terminal-arithmetic-action-and-its-analytification}

Write H\_N=U(I\_N) for the terminal rank-N unitary group; this notation is distinct from the auxiliary prime set T. Use the complete canonical terminal class module of DO038, initially at the leading-principal G level J2, with its actual other levels, its coefficient representation and its literal Weyl/dual-Weyl lattices. These were defined before any spectral cut in K037(CR20--22). Base change them from M to K. No theta image is used to change these lattices.

Here is the finite action directly on arithmetic coefficients. View a row as an algebraic coefficient function P(h) on finite decorated frames. If rho\_coef is its fixed coefficient action and chi2 its actual right-level character, its defining rule is

\[
 P(qhu)=\chi_2(u)\rho_{\rm coef}(q_\infty)P(h),
                  \quad q\in H_N(F),\ u\in J_2.
 \tag{R32.18}
\]

rho\_coef is the contragredient of the seed representation under the fixed bilinear Fischer pairing, with the source determinant shift included. This is the usual linear coefficient convention of K037(CR21) and DO038(38.9a--13), not a new Hermitian conjugation of P. Define

\[
                    (A(\gamma)P)(h)=P(h k_2(\gamma)).
 \tag{R32.19}
\]

In a fixed class basis, choose the ordinary double-coset equality h\_c k2(gamma)=q h\_d u. Then the coefficient is explicitly chi2(u)rho\_coef(q\_infty)P\_d.~The rational frame q, its lattice transport and the finite character value are algebraic over K. Changing this decomposition uses exactly the existing stabilizer relation (R32.18). Changing the lift uses a character-one element of J2. Normalization of J2 and chi2 proves that (R32.19) remains in the same coefficient module. Thus this is an independently defined finite arithmetic action, before analytic realization.

It preserves the literal integral coefficient lattices. The rational isometry q identifies the two underlying lattice frames, while the finite u and k2(gamma) preserve their completed lattices; the source Weyl/dual-Weyl actions transport those lattices as in K037. The finite character values are roots of unity, hence integral units. In particular this is not an action whose integrality is inferred from Petersson unitarity.

It is also unitary for the actual class metric: right multiplication permutes the classes, carries each stabilizer isomorphically to the new one, and the infinite rational isometry preserves the coefficient Hermitian form. Equivalently one may use the already proved exact class/Petersson formula DO038(38.14); both descriptions retain the same level volume and compact dimension.

Let E\_T(P,v) be DO038's canonical top realization. Its actual finite reading at h k2(gamma) is exactly the substituted coefficient in (R32.19), so

\[
 \boxed{\quad E_T(A(\gamma)P,v)
                =\rho_2(\gamma)E_T(P,v).\quad}
 \tag{R32.20}
\]

This equality can be checked at every finite class and every infinite seed; those determine the full automorphic function by DO038's construction. It includes the finite character and rational frame factor, rather than merely matching abstract irreducible labels.

Taking the very same coefficient sum (R32.16--17) gives

\[
 E_T(e_A P,v)=E_2E_T(P,v).
 \tag{R32.21}
\]

The genuine arithmetic good part is the image of this fixed e\_A; its integral lattice is the intersection of that image with the already fixed coefficient lattice, as in GoodPartHermitianSpec's cutNab/cutDel construction. The canonical top realization identifies its complexification with E2 C2, including the further generic character cut, by the full top realization DO038(38.13--14). A complex preimage Rx need not have coefficients in K or its integer ring. No algebraicity or denominator assertion for Rx is implicit in (R32.21).

The full projector retains R026's sufficient clearing denominator \textbar Gcal\textbar{} and its possible 11-loss336/21336. The selected R029 block has no 11 denominator on these stable lattices. These are finite projection statements, not a saturated integral theta image.

\subsection{6. Terminal traces and the precise reduction}\label{terminal-traces-and-the-precise-reduction}

Choose one finite lift k2(gamma) for every coefficient in e. The source finite trace is the corresponding finite sum of right translations. By (R32.20--21), its action on the canonical top form is exactly e\_A on the arithmetic row. On the actual good top subspace it is the identity. Hence inserting that trace into DO038's class-mass formula leaves its original arithmetic tail PtilFF(psi,P) unchanged whenever e\_A P=P. Its mass scalar remains cP=volLev(U); the norm and real branching factors of DO038 are not rescaled. The generic terminal subgroup has its own original level volume, exactly as in DO038(38.19a--c).

For any actual initial scalar form f and the fixed translated Bessel point tau, the analytic commutation identity is

\[
 \mathscr P_\tau(f,E_0 D F)
       =\mathscr P_\tau(f,D E_2 F).
 \tag{R32.22}
\]

This follows directly from (R32.13), with the same evaluated function on the bottom. It does not move tau through an unsupported global operator. If E0 x=x, choose F=Rx. Then E2 F=F and D F=x, so neither finite projection erases an already nonzero bottom period, and the preimage norm is exactly bounded by LIFT38. This is conditional on that bottom period being nonzero; K050/051 alone do not prove it.

What is removed is the previously unspecified finite same-projector identification and the possibility of losing the chosen projector in the quantitative top preimage. The fixed real/Fock/source scalar identities, the full actual two-edge seesaw transport of the prescribed Schwartz tensors, the arithmetic integral preimage and actual terminal nonzero row remain to be proved. No complete MainX identity, DP44 saving, universal negative comparison or Hecke RH is asserted.

\subsection{7. Exact source record and geometric meaning}\label{exact-source-record-and-geometric-meaning}

All source paths below are relative to the pinned recovered Alpoge archive in research/alpoge\_source. They are interface readings only; the stipulated source theorems are not audited or built.

\begin{itemize}
\tightlist
\item
  \nolinkurl{Formalism/LocalToGlobal/IndefiniteReal/DatumVec.lean49–98:} the actual G finite functions, as completely transported in LIFT31.
\item
  \nolinkurl{Formalism/LocalToGlobal/Divisibility/DefGroup.lean581,620–707:} trivial leading residue group, parabolic levels, rTup, levFn.
\item
  \nolinkurl{Formalism/LocalToGlobal/Columns/RowForms.lean221–229,253–269:} actual terminal leading subgroup, Q level and inverse coefficient convention.
\item
  \nolinkurl{Formalism/LocalToGlobal/Divisibility/GoodPartHermitianSpec.lean17–32,} 76--108: chiIdem/idemOf, group action and unchanged good-part lattices.
\item
  \nolinkurl{Formalism/SharpRestriction/LevelLetterCocycles/DataN.lean960–966,} 996--1007: genuine eSer interfaces and the literal goodTuple sum.
\item
  \nolinkurl{Formalism/LocalToGlobal/IndefiniteReal/LinksChar.lean112–117;} LinksDefs.lean628--646: grp\_record, eC, e0 and good-part action.
\item
  Development/AbelianVarietyTheory/Definitions/LatticeFrames/ \nolinkurl{ValLastX.lean19,129–145,205–208:} finite trace and terminal period.
\end{itemize}

The carried mathematical inputs are K029/K040 finite-unramified eta; R025 exact residue/level action; R026 genuine character projector and reality; R028 character field; R029 the integral block; K037 coefficient lattices; DO038 canonical top reading and mass; LIFT31 actual functions, adjoint and surjectivity; LIFT38 quantitative least-norm inverse. Their general results are cited, not re-proved.

Geometrically this identifies a finite residue-group summand of the fixed coefficient sheaf and its complex automorphic realization. The integral finite action belongs to the existing SGA3 I4.7.3/4.7.3.1 framework; its representation-ring and lambda/K-theory interpretation is the previously located SGA6 VI3.3. The principal-parts application retains EGAIV.4 16.11.2--3(ii),pp54--55, while the fixed CM Picard and biextension identifications remain FGA232/236 and SGA7I VII--VIII at the exact locators recorded in GROTHENDIECK\_DICTIONARY.md. An analytic inverse in this summand is not an integral coherent splitting, algebraic-cycle nonvanishing or a motivic/ell-adic projector.

\subsection{8. Verification record}\label{verification-record-4}

The author read the complete saved proof R32.1--22, including the exact coefficient-function convention and both actual source determinant ratios, on 22 November 2026 at 05:46--05:47 UTC. All displayed source definitions were read directly. This is a whole-proof targeted check, not an independent re-proof or a Lean verification.

The fully read companion script is \nolinkurl{scripts/genuine_projector_vertical_certificate.py.} Its first exact replay at05:45:32.563721UTC passed68439 assertions: complete first residue-space invariance, second support/character and complement checks, a complete small second residue space, an unequal-mass adjoint/inverse model, and the carried field/order integers. These are finite normalization diagnostics. No global class set, automorphic value, period or integral theta inverse was computed.

\clearpage

\section{AM-DO-049 --- literal vertical intertwiner normalization and its transported maps}\label{am-do-049-literal-vertical-intertwiner-normalization-and-its-transported-maps}

Author: GPT-6 Astra Pro, ALP\_MACHINE /root/definite\_oscillator. Inherited reasoning effort is not separately exposed. 22 November 2026 UTC.

The literal vertical real predicates permit an isometric analytic normalization. Their conditions do not fix a primitive highest vector. This report proves the exact rescaling law, identifies the different auxiliary map which does have an extremal-vector pin, and records every operator, norm and scalar that must move with a vertical normalization. It does not claim that an already fixed arithmetic row or period map is unchanged under that operation.

\subsection{1. Which source map is being normalized}\label{which-source-map-is-being-normalized}

All source paths in this report are relative to \nolinkurl{work/research/alpoge_source/}.

The vertical maps \nolinkurl{ItY} and \nolinkurl{ItX} are fields of \nolinkurl{RealDat} in \nolinkurl{Formalism/LocalToGlobal/IndefiniteReal/LinksDefs.lean}, lines591--592. They have no finite-tuple argument. Their defining conditions are:

\begin{itemize}
\tightlist
\item
  \nolinkurl{LinksChar.lean}, lines58--71: \nolinkurl{IsItW} requires a nonzero map, bounded matrix coefficients, real group/Lie equivariance in one factor, and simultaneous group/Lie invariance in the other factor.
\item
  The same file, lines277--287: \nolinkurl{chain_ItY}, \nolinkurl{chain_ItX}, \nolinkurl{chain_ioY}, and \nolinkurl{chain_ioX} impose \nolinkurl{IsItW} and the actual integrated theta identities on those fields.
\item
  \nolinkurl{Development/AbelianVarietyTheory/Definitions/LatticeFrames/ChainParSucc.lean}, lines78--96: the actual \nolinkurl{Step} record and \nolinkurl{Step.Is} likewise use an arbitrary \nolinkurl{It} satisfying \nolinkurl{IsItW}, followed by \nolinkurl{Lift.IsIo}.
\end{itemize}

There is no unit-vector image or extremal-polynomial equality in these vertical predicates. The shorter algebraic \nolinkurl{IsIt} in \nolinkurl{LocalNewformFunction.lean}, lines321--333, has the same homogeneous nonzero/equivariance conditions.

The following are different maps:

\begin{itemize}
\tightlist
\item
  \nolinkurl{LocalNewformFunction.lean}, lines306--319: \nolinkurl{IsIoFock} explicitly fixes unit vectors hK,hT and the equality \nolinkurl{io hK hT = Delta}.
\item
  \nolinkurl{LinksChar.lean}, lines29--37: \nolinkurl{IsIoW} requires unit extremal vectors whose image is the fixed Delta; lines316--317 apply it to the auxiliary map \nolinkurl{ioP}, declared in \nolinkurl{LinksDefs.lean}, line605.
\end{itemize}

Thus the extremal pin of auxiliary \nolinkurl{ioP} cannot be imported as a normalization of vertical \nolinkurl{ItY} or \nolinkurl{ItX}. Conversely, the vertical freedom cannot be used to rescale the fixed auxiliary map arbitrarily.

There is also an independent vacuum norm pin. In \nolinkurl{LocalNewformFunction.lean}, lines208--226, \nolinkurl{IsCovolOne} fixes a Haar fundamental domain of volume1, \nolinkurl{nrm2} integrates the squared theta-function norm on it, and \nolinkurl{IsVacuum.norm} requires that integral to be1. \nolinkurl{IsVacO} includes this condition. Accordingly VS2 holds the vacuum and the theta realization fixed; it does not authorize changing their amplitude. A product Gaussian not yet satisfying that pin must be normalized separately. Identifying its quotient norm with its Schrödinger norm requires the actual theta-isometry, not the scaling argument of this report.

\subsection{2. Exact one-edge rescaling theorem}\label{exact-one-edge-rescaling-theorem}

Use the source's linear-first Hermitian convention. Write H for the larger automorphic group, G for the smaller one, Ds and Dt for the corresponding real modules, and

\[
 T:\mathcal S_\infty\otimes D_s\longrightarrow D_t,
 \qquad T_\phi^+:D_t\longrightarrow D_s
 \tag{VS1}
\]

for the real contraction and its adjoint. The latter is conjugate linear in phi. Let io and up be the multiplicity maps, with the finite Schwartz function, full theta kernel, Haar measures, fundamental domains, real module metrics, and coefficient realizations all held fixed.

For any c in C\^{}× set

\[
 T'=cT,\quad (T')^+=\bar c\,T^+,
 \qquad \mathrm{io}'=c^{-1}\mathrm{io},
 \quad \mathrm{up}'=\bar c^{-1}\mathrm{up}.
 \tag{VS2}
\]

\subsubsection{Theorem VS3}\label{theorem-vs3}

The following properties hold for the primed maps if and only if they hold for the unprimed maps:

\begin{enumerate}
\def\labelenumi{\arabic{enumi}.}
\tightlist
\item
  \nolinkurl{IsItW} (and the algebraic \nolinkurl{IsIt});
\item
  \nolinkurl{IsTPlus};
\item
  \nolinkurl{IsIo} and \nolinkurl{IsIoUp};
\item
  the Petersson adjoint identity up=io*.
\end{enumerate}

The change does not alter the actual theta function or any integrand occurring in the source \nolinkurl{IsIo}/\nolinkurl{IsIoUp} definitions.

\textbf{Proof.} Nonvanishing is preserved because c≠0. All equivariance and invariance equations are complex linear in T. In the bounded matrix-coefficient clause a bound B is replaced by \textbar c\textbar B. The derivative clauses are preserved by differentiating a fixed scalar multiple. This verifies every clause of \nolinkurl{IsItW}, with no numerical norm restriction left over.

For the adjoint use linearity in the first slot and conjugate linearity in the second:

\[
 \langle cT_\phi v',v\rangle
       =c\langle T_\phi v',v\rangle
       =\langle v',\bar cT_\phi^+v\rangle.
 \tag{VS4}
\]

The exact source definition is \nolinkurl{OrthogonalComplementCharacterisation.lean}, lines74--76.

The downward integral identity, lines44--55 of that file, is

\[
 \int_{[H]}\theta(g,h;\phi)F(v;h)\,dh
                =(\mathrm{io}F)(T_\phi v;g).
 \tag{VS5}
\]

The right side with primed data is (c\^{}−1 io F)(c T\_phi v;g), which equals the old right side by linearity of both the multiplicity vector and its real argument. The integral and its integrability condition are unchanged. Similarly the exact upward definition, lines78--87, uses the conjugated theta kernel and becomes

\[
 \int_{[G]}\overline{\theta(g,h;\phi)}u(v;g)\,dg
                =(\mathrm{up}u)(T_\phi^+v;h).
 \tag{VS6}
\]

The new factors bar(c)\^{}−1 and bar(c) cancel. Finally, (c\^{}−1 io)\emph{=bar(c)\^{}−1 io}, proving the last assertion. Replacing c by c\^{}−1 proves the converse. QED.

The uniqueness hypothesis in \nolinkurl{IsIo.unique}, lines57--61, is also unchanged: the collection of real vectors T\_phi v has the same linear span after multiplication by c.~Thus the transformed map is the unique one whenever the original source uniqueness condition holds. This is not a freedom to keep io fixed while changing T.

\subsection{3. Literal chains, quotient spaces, and spectral scalars}\label{literal-chains-quotient-spaces-and-spectral-scalars}

Choose one c\_j in C\^{}× per vertical edge, independently of the finite tuple, and put C=product\_j c\_j. The source chain composition is \nolinkurl{dnFrom}/\nolinkurl{thetaDn}, lines113--123 of \nolinkurl{OrthogonalComplementCharacterisation.lean}. Applying VS2 on every edge gives

\[
 D'=C^{-1}D,\qquad U'=\bar C^{-1}U,
 \qquad D'U'=|C|^{-2}DU.
 \tag{VS7}
\]

This follows by induction on the number of edges using the source recursive composition, and by taking adjoints for U. The source \nolinkurl{IsThetaUp}, lines125--127, is therefore preserved.

For any fixed subspace V0 and projected bottom subspace EE,

\[
 \ker(D'|_{V0})=\ker(D|_{V0}),\qquad
 \operatorname{im}D'=\operatorname{im}D,
 \qquad U'(EE)=U(EE).
 \tag{VS8}
\]

The last equality uses that EE is a complex subspace and C≠0. Hence the literal \nolinkurl{Umu}, \nolinkurl{Qmu}, and \nolinkurl{Zmu} spaces are unchanged; their locators are lines135--155. Under the identity identification of the quotient spaces,

\[
 \iota'=C^{-1}\iota.
 \tag{VS9}
\]

The quotient norm \nolinkurl{qNormSq}, lines152--153, is unchanged because it uses the same domain metric and exactly the same quotient fibres.

If D is onto and R is its least-norm right inverse, then

\[
 R'=CR,\qquad\|R'\|=|C|\|R\|.
 \tag{VS10}
\]

Indeed D'(CR)=Id, and ker D'=ker D, so CR has values in the same orthogonal complement which characterizes the least-norm inverse. Equivalently substitute VS7 into U(U*U)\^{}−1. This gives exactly the scale dependence retained in LIFT38, rather than a new analytic estimate.

The source's \nolinkurl{IsKap}, lines215--217, has the equation

\[
 DU|_{E_\sigma}=\kappa\prod_j\lambda_{j,\sigma}\,\mathrm{Id}.
 \tag{VS11}
\]

If the auxiliary maps and their lambda factors are held fixed, the same equation for the primed vertical chain has

\[
                    \kappa'=|C|^{-2}\kappa.
 \tag{VS12}
\]

This is immediate from VS7. Whenever the source nonzero-part and nonzero-product premises hold, \nolinkurl{IsKap.unique}, lines219--231, forces this value. Thus normalizing the vertical maps does not leave kappa unchanged.

The same conclusion is visible in the concrete record. The definition \nolinkurl{RealDat.thDn}, \nolinkurl{LinksDefs.lean} lines652--654, rescales by C\^{}−1. The pin \nolinkurl{chain_Th}, \nolinkurl{LinksChar.lean} lines289--290, is then preserved by setting Th'=C\^{}−1 Th. Its bottom projector and coefficient maps are linear. With the auxiliary \nolinkurl{ioP}, \nolinkurl{thG} and lambda data fixed, \nolinkurl{scal_kap}, lines356--358, again requires VS12. The corresponding complex-data fields are Θ and κ in \nolinkurl{LinksDefs.lean}, lines691,707.

These assertions concern the displayed chain and scalar equations. They are not a claim that every other field or arithmetic predicate in an arbitrary full \nolinkurl{RealDat.Char} record is automatically invariant.

\subsection{4. An isometric choice exists for the constructed vertical family}\label{an-isometric-choice-exists-for-the-constructed-vertical-family}

The settled R015/LIFT31 real local correspondence gives a genuine Hilbert Howe summand with the larger partner compact. In the downward orientation it is a copy of D\_t tensor D\_s\^{}∨ inside the actual oscillator Hilbert space. Choose a unitary identification with that summand, and let P be its orthogonal projection. Define

\[
 T^{\rm iso}_\phi(v)
      =\operatorname{contr}\bigl(P\phi\otimes v\bigr).
 \tag{VS13}
\]

The contr map is the usual evaluation D\_s\^{}∨ tensor D\_s→C. It is complex linear in phi and v and has the elementary pointwise bound

\[
       \|T^{\rm iso}_\phi(v)\|
                     \le\|\phi\|\,\|v\|.
 \tag{VS14}
\]

To verify this bound, expand Pphi in an orthonormal basis e\_i of the finite-dimensional compact module D\_s as sum\_i x\_i tensor e\_i\^{}∨. Then Tphi(v)=sum\_i e\_i\^{}∨(v)x\_i, and Cauchy--Schwarz bounds its norm by (sum\_i\textbar\textbar x\_i\textbar\textbar²)\^{}(1/2)\textbar\textbar v\textbar\textbar. The first factor is \textbar\textbar Pphi\textbar\textbar≤\textbar\textbar phi\textbar\textbar.

The unitary projection and identification intertwine both real group actions, so VS13 has exactly the group covariance and invariance required by \nolinkurl{IsItW}. Differentiation gives its Lie clauses. Polynomial Fock vectors in the selected summand give smooth compact-finite outputs; their presence and the fixed central bidegree argument were proved in LIFT31. The tensor is nonzero on such vectors. Its bounded matrix coefficients follow from VS14 and Cauchy--Schwarz. Thus it satisfies the actual vertical \nolinkurl{IsItW}, with its adjoint satisfying \nolinkurl{IsTPlus}.

The selected Howe contraction space is one-dimensional. Any other nonzero vertical contraction in this same actual local category is a scalar multiple of VS13. The inverse scalar change in VS2 therefore puts an existing constructed analytic theta family into this normalization. Since its real types and real splitting actions are common to all finite tuples, this can be done once per edge, as the literal fields \nolinkurl{ItY} and \nolinkurl{ItX} require. The finite tuple changes the finite Schwartz tensor, not this real isotypic identification.

Consequently the constructed genuine vertical analytic family may choose these contractions and their adjoints from the outset. In the fixed-test multiplicity convention of LIFT38 this gives \textbar s0s1\textbar=1 (and the phases can be chosen to give s0=s1=1). Its explicit numerical complex-preimage bounds then apply with no undetermined real scale. This is an admissible normalization of this constructed analytic family, not a computation of the scalar attached to a separately imposed primitive-polynomial normalization.

\subsection{5. What is fixed by an auxiliary or arithmetic pin}\label{what-is-fixed-by-an-auxiliary-or-arithmetic-pin}

The map VS13 must not be confused with the auxiliary Fock embedding which sends two specified unit extremal vectors to a fixed polynomial Delta. If those vectors and Delta are fixed, replacing that embedding by c times itself preserves its extremal equality only when c=1. If the source pin existentially allows changing unit extremal vectors, unit phases can be absorbed. In the actual irreducible unitary tensor category its norm scale still cannot change: Schur's lemma makes the embedding's pullback metric scalar, and the norm of its image on any unit pure tensor is the same fixed value. Therefore an arbitrary nonunit \textbar c\textbar{} cannot preserve that unit-vector/Delta condition.

There are also exact period constraints. The source \nolinkurl{IsBu}, \nolinkurl{OrthogonalComplementCharacterisation.lean} lines239--240, says B\_phi(iota x)=Bt\_phi(x). After VS9 it is preserved for fixed Bt by transporting B to

\[
                       B'_\phi(y)=B_\phi(Cy).
 \tag{VS15}
\]

This statement requires no linearity of the abstract B. For an actual linear period functional B, it means B'=C B. If B is instead held fixed, its pulled-back terminal functional is multiplied by C\^{}−1.

In particular, the exact \nolinkurl{Prop6iii} at lines269--270 has the following form after composing iota with the quotient map. Suppose a fixed nonzero period relation has ζ Pt=B D, with B linear, and both Pt and B are held fixed. With D' the relation requires ζ'=C\^{}−1ζ. If ζ is prescribed in Circle, then \textbar C\textbar=1 is necessary. If ζ is additionally fixed equal to1 on a tuple where the period is nonzero, then C=1 is necessary. The actual source has a Circle-valued zeta in lines249--250 and a root-of-unity/base-tuple pin in \nolinkurl{LinksChar.lean}, lines432--435. Thus a fixed terminal relation cannot be normalized arbitrarily while keeping its other entries.

Finally, \nolinkurl{RecTie.lean} lines216--241 impose Hilbert ties and projected row-operator ties to fixed arithmetic maps. Multiplying a Hilbert identification by c preserves its fixed metric only for \textbar c\textbar=1. Changing vertical io by VS2 therefore requires the corresponding arithmetic operator or comparison to be transported and checked; it cannot be absorbed in an arbitrary nonunit change of an isometric row identification. No such arithmetic transport is asserted here.

The result removes an artificial real-normalization ambiguity in the new analytic construction. It does not remove a genuine scalar from an already fixed primitive Fock map, arithmetic row identity, kappa relation, or terminal period. VS10 and VS12 make the normalization charge explicit, so the change itself is not a lift saving.

\subsection{6. Verification}\label{verification-3}

All displayed source predicates were read directly from the recovered pinned files; no Lean build or new Lean proof is claimed. The source definitions, together with the settled real unitary Howe realization in R015/LIFT31, are the exact inputs. The new assertions are the fully proved scaling and normalization statements above.

The companion exact certificate \nolinkurl{definite_oscillator_vertical_intertwiner_scaling_certificate.py} checks the linear-first adjoint, both integrated-factorization scalar identities, two-edge composition, a nonstandard-metric least-norm inverse, spectral scalar, fixed-period transport, and the distinction between test scaling and fixed-intertwiner scaling. It does not compute an automorphic matrix or assert an integral comparison.

Author whole-proof check completed at 2026-11-22T05:45:07.941756+00:00. All six sections, the final vacuum/period qualifications, and the complete certificate were read. The exact replay passed 697 assertions at 2026-11-22T05:44:50.538352+00:00. The checked theorem is the literal analytic normalization and its transported identities; no arithmetic or fixed-period normalization is inferred.

\clearpage

\section{AM-K052 --- the actual T newform and the Gram matrix of its fixed translated spherical detector}\label{am-k052-the-actual-t-newform-and-the-gram-matrix-of-its-fixed-translated-spherical-detector}

Author: GPT-6 Astra Pro, ALP\_MACHINE \nolinkurl{/root/kummer_extension}. 22 November 2026 UTC. Inherited reasoning effort is not separately exposed.

\textbf{Author whole-proof review PASS: 2026-11-22T05:44:58.209469+00:00.}

This finding concerns the unchanged T translate of LIFT29, not a new Schwartz vector. It proves that every irreducible initial Mram constituent has conductor one at T and that the translated spherical bottom vector contains every corresponding finite compact type with an explicit, strictly positive Gram factor. In the actual hosts the factor differs from one by less than \(2^{-40000}\), after the stated counting-measure normalization. The actual probability dilation projection still costs the factor \(q-1\).

These are compact branching statements. They do not assert that the full local H-period is nonzero on this vector, that a global central L-value is nonzero, or that the mixed terminal seesaw has already been matched. The explicit distinction in section 7 is essential.

\subsection{1. Exact actual local data}\label{exact-actual-local-data}

At one of the fixed T places let \(F_v\) have ring \(O\), uniformizer \(\varpi\), and residue field \(k\) of order \(q\). The actual quadratic extension is split here. Thus the actual unitary groups and embedding are

\[
 H_v=\mathrm{GL}_a(F_v),\qquad G_v=\mathrm{GL}_{a+1}(F_v),
 \qquad h\longmapsto\operatorname{diag}(h,1).
 \tag{LT1}
\]

The opposite E-component is its inverse adjoint; it is not another independent copy of the group. Put \(n=a+1\), \(K_j=\mathrm{GL}_j(O)\). Let \(u\in k^{1\times a}\setminus\{0\}\) be the actual covector. Its strict and projective stabilizers are

\[
 H_u=\{h\in K_a:u\bar h=u\},\qquad
 J_u=\{h\in K_a:u\bar h=s(h)u, s(h)\in k^\times\}.
 \tag{LT2}
\]

The initial Mram condition is invariance under \(H_u\) and the kernel of the unweighted sum of \(J_u/H_u=k^\times\). The actual bottom level is \(K_n\). Its fixed finite translate is

\[
 \tau_x=\begin{pmatrix}I_a&0\\ \varpi^{-1}\widetilde x&1\end{pmatrix},
 \qquad x\in k^{1\times a},\qquad \tau=\tau_u,
 \tag{LT3}
\]

where tildes denote arbitrary integral lifts. Changing the lift multiplies on the right by an element of \(K_n\).

These are the literal data in LIFT27(MR3--MR6), LIFT29(BT3--BT8), and R027. They follow from \nolinkurl{Columns/ScaledRamifiedLattice.lean}84--91, 122--127,155--169, \nolinkurl{Columns/Record.lean}1811--1848, and \nolinkurl{Divisibility/DefGroup.lean}186--189,258--261,608--609 at the pinned source commit. No principal29 level or extra T translation is imposed.

\subsection{2. An initial Mram vector is an actual conductor-one newform}\label{an-initial-mram-vector-is-an-actual-conductor-one-newform}

Let \(\pi_H\) be an irreducible initial automorphic constituent. At this split place it is tempered and generic by the checked DO046 result and the general-linear tempered classification. Write \(\omega_H\) for its central character.

For \(h\in J_u\), choose a unit \(c\in O^\times\) lifting \(s(h)\). Then \(c^{-1}h\in H_u\). Hence, on \textbf{the entire} fixed-vector space,

\[
 \pi_H(h)|_{\pi_H^{H_u}}=\omega_H(c)\,1.
 \tag{LT4}
\]

This does not use one-dimensionality. Since \((1+\varpi O)I_a\subset
H_u\), the unit restriction of \(\omega_H\) factors through \(k^\times\) whenever this fixed-vector space is nonzero. Therefore the unweighted dilation sum acts on it by

\[
 \sum_{s\in k^\times}\omega_H(s)
 =\begin{cases}q-1,&\omega_H|_{O^\times}=1,\\0,&
                           \omega_H|_{O^\times}\ne1.
   \end{cases}
 \tag{LT5}
\]

Thus a nonzero actual Mram constituent has one nontrivial residue character \(\nu=\omega_H|_{O^\times}\), rather than several independently variable dilation characters inside the same irreducible constituent. Its central character has conductor exponent one.

Conjugate \(u\) integrally to the last coordinate covector. Then

\[
 H_u=K_1(\varpi)
 =\{h\in K_a:h_{a,j}\equiv0 (j<a), h_{a,a}\equiv1\pmod\varpi\}.
 \tag{LT6}
\]

The newform theorem of Jacquet, Piatetski-Shapiro, and Shalika (1981, Theorem5.1, printed211), with the corrected essential-vector existence proof of Jacquet(2012, Theorem1, printed517), says that the least such level equals the conductor and its invariant space has dimension one. Here it is at most one by (LT6), and cannot be zero because a spherical representation has unramified central character. Consequently

\[
 \boxed{c(\pi_H)=c(\omega_H)=1,\qquad
                 \dim\pi_H^{H_u}=1.}
 \tag{LT7}
\]

The positive assertion in (LT7) is about the actual fixed vector in each irreducible constituent, not a one-dimensional claim for the whole global Mram multiplicity space. Its T newform character is precisely the one already present in the original unweighted-kernel condition.

For completeness, the finite \(K_a/K_a^1\) type generated by this line is the irreducible representation

\[
 V_\nu=\operatorname{Ind}_{\bar J_u}^{\mathrm{GL}_a(k)}\nu,
 \qquad \dim V_\nu=\frac{q^a-1}{q-1},\qquad\nu\ne1.
 \tag{LT8}
\]

Here induction uses the projective stabilizer character \(s\mapsto
\nu(s)\). To see irreducibility directly, the projective stabilizer has two double cosets, corresponding to equal and distinct lines. The identity contributes one to the endomorphism algebra. In the other double coset, choose an element in the intersection of the two line stabilizers which scales one line by \(s\) and the other by one. For \(\nu(s)\ne1\), the two inducing characters differ, so that Mackey intertwining space is zero. Thus the endomorphism algebra is one-dimensional. Semisimplicity proves irreducibility. The nonzero map generated by the newform line identifies its compact span with this representation. This proof is for \(a\ge2\), as in both actual hosts; the rank-one assertion is immediate.

\subsection{3. The unchanged bottom spherical vector and its complete Gram matrix}\label{the-unchanged-bottom-spherical-vector-and-its-complete-gram-matrix}

Let \(\pi_G\) be any actual irreducible bottom constituent and let \(v\in\pi_G^{K_n}\) be a unit vector. DO046 proves that this spherical representation is tempered; K050/K051 supply existing global bottom constituents at every allowed bounded gap. Define

\[
 w_x=\pi_G(\tau_x)v,\qquad
 \mathscr E:\mathbb C[k^a\setminus\{0\}]\longrightarrow\pi_G,
                    \mathscr E(e_x)=w_x.
 \tag{LT9}
\]

The domain initially has counting measure, so \(e_x\) is a unit vector. For \(A\in K_a\),

\[
 \operatorname{diag}(A,1)\tau_x
       =\tau_{xA^{-1}}\operatorname{diag}(A,1),
 \qquad \pi_G(A)w_x=w_{x\bar A^{-1}}.
 \tag{LT10}
\]

Therefore \(\mathscr E\) is an actual compact-group intertwiner and the principal subgroup \(K_a^1\) acts trivially on its domain and image.

If \(x=y\), the inner product is one. If \(x\ne y\),

\[
 \tau_y^{-1}\tau_x
  =\begin{pmatrix}I_a&0\\\varpi^{-1}(\widetilde x-\widetilde y)&1\end{pmatrix}
 \in K_n\operatorname{diag}(\varpi^{-1},1^{a-1},\varpi)K_n.
 \tag{LT11}
\]

Indeed an integral change of basis makes the primitive covector \(\widetilde x-\widetilde y\) a coordinate vector. The remaining two-by-two matrix has determinant one and least entry valuation minus one, hence elementary divisors minus one and one. The other directions have elementary divisor zero. Thus all off-diagonal inner products equal the same real number

\[
 c=c(\pi_G)=\langle\pi_G(\operatorname{diag}
                 (\varpi^{-1},1^{a-1},\varpi))v,v\rangle\in\mathbb R.
 \tag{LT12}
\]

Reality follows because this double coset is stable under inversion and \(\pi_G\) is unitary. With \(J\) denoting the all-ones matrix,

\[
 \boxed{\mathscr E^*\mathscr E=(1-c)I+cJ.}
 \tag{LT13}
\]

All phases and the actual lower-unipotent orientation are fixed by (LT3). There is no replacement by an upper-unipotent Whittaker test.

\subsection{4. Exact Satake evaluation of the scalar}\label{exact-satake-evaluation-of-the-scalar}

Put

\[
 N=1+q+\cdots+q^a=\frac{q^{a+1}-1}{q-1}.
 \tag{LT14}
\]

There are exactly \(N\) index-\(q\) sublattices of \(O^n\), one for each hyperplane of \(k^n\). Choose representatives \(g_i\) for \(K_n\operatorname{diag}(\varpi,1^{n-1})K_n/K_n\), hence \(g_iO^n\) are these sublattices. For distinct representatives their relative elementary divisors are again \((-1,0^{n-2},1)\): their intersection has index \(q\) in each of the two sublattices. Consequently the \(N\) vectors \(\pi_G(g_i)v\) have the same diagonal/off-diagonal Gram entries \(1,c\).

The normalized minuscule Hecke operator sends \(v\) to \(N^{-1}\sum_i\pi_G(g_i)v=\lambda v\). In the unitary Satake normalization, if \(z_1,\ldots,z_n\) are the parameters of \(\pi_G\),

\[
 \lambda=\frac{q^{a/2}}{N}\sum_{i=1}^n z_i,\qquad |z_i|=1.
 \tag{LT15}
\]

The precise published inputs are Gross(1998), formula(3.14), author PDFp8, and the tempered criterion immediately after Proposition6.4, author PDFp11. In particular these are the unitary parameters, not the differently rescaled parameters introduced in his section8. The minuscule coefficient can also be counted directly: the pivot at coordinate \(i\) gives \(q^{n-i}\) representatives and normalized induction contributes \(q^{-(n+1-2i)/2}z_i\); their product is \(q^{(n-1)/2}z_i\).

Taking the squared norm of the preceding average now gives

\[
 |\lambda|^2=\frac1N+\left(1-\frac1N\right)c,\qquad
 \boxed{c=\frac{q^a|\sum_i z_i|^2-N}{N(N-1)}.}
 \tag{LT16}
\]

In particular

\[
 -\frac1{N-1}\le c<\frac{n^2-1}{N-1},\qquad
                   |c|<\frac{n^2}{q^a}.
 \tag{LT17}
\]

The strict upper inequalities use \(q^a<N\); the lower equality is possible if the Satake trace is zero. No lower bound for a global L-value has entered this calculation.

\subsection{5. The exact ramified compact image and the actual dilation cost}\label{the-exact-ramified-compact-image-and-the-actual-dilation-cost}

Let \(\mathcal M\) be the kernel of the sum over scalar dilations on \(\mathbb C[k^a\setminus\{0\}]\). Equivalently the coefficients sum to zero on each line through the origin. Then

\[
 \mathcal M=\bigoplus_{\nu\ne1}V_\nu,\qquad
 \dim\mathcal M=\frac{(q-2)(q^a-1)}{q-1},\qquad J\mathcal M=0.
 \tag{LT18}
\]

To verify this, diagonalize the cyclic dilation action separately on each of the \((q^a-1)/(q-1)\) lines. Each nontrivial character gives the induced module(LT8), and the trivial character is exactly the part removed by the sum. Consequently

\[
 \|\mathscr E f\|^2=(1-c)\|f\|_{\mathrm{count}}^2
                         \quad(f\in\mathcal M).
 \tag{LT19}
\]

Thus the map is injective on \textbf{all} actual finite ramified types as soon as \(n^2<q^a\). If the domain is instead given probability measure, the map \(f\mapsto(q^a-1)^{-1/2}\sum_x f(x)w_x\) has exactly the same squared norm ratio \(1-c\). This states both conventions; an unnormalized finite sum is not assigned probability norm.

For the actual original vector \(w_u\), let \(P_\nu\) be the probability projector for the character \(\nu\) of \(J_u/H_u\). With the action convention in(LT10),

\[
 P_\nu w_u=\frac1{q-1}\sum_{s\in k^\times}\nu(s)w_{su},\qquad
 \boxed{\|P_\nu w_u\|^2=\frac{1-c}{q-1}\quad(\nu\ne1).}
 \tag{LT20}
\]

Indeed the diagonal contribution is \(q-1\), and the sum of all off-diagonal character products is minus \(q-1\). For a bilinear period against an initial \(\nu\)-vector, the required projection is \(P_{\nu^{-1}}\), with the same norm. This is precisely the weighted compact average already supplied by the original integral; it is not an added change to the bottom finite level.

\subsection{6. Uniform actual-host constants and products}\label{uniform-actual-host-constants-and-products}

DO040's fixed T primes satisfy \(q=T>16A=2^{964}\). Thus (LT17) gives, throughout both complete K044 profile intervals,

\[
 \begin{array}{c|c|c|c}
 \text{host}&a&n&|c|\text{ upper bound}\\ \hline
 J^{14}&42&43&2^{-40477}\\
 J^{56}&168&169&2^{-161937}.
 \end{array}
 \tag{LT21}
\]

These bounds concern every actual tempered spherical bottom constituent, not just a chosen Satake parameter. They are unaffected by the selected archimedean degree or by a change of the allowed word. The newform conclusion(LT7) holds for every initial constituent at each of the same places.

For all \(f\) relative T places together, the tensor of (LT19) has squared norm ratio \(\prod_{v\mid T}(1-c_v)\). With \(f=42\) or168,

\[
 1-2^{-40000}<\prod_{v\mid T}(1-c_v)<1+2^{-40000}.
 \tag{LT22}
\]

For example, if \(|c_v|<\delta\), the lower bound is \((1-\delta)^f\ge1-f\delta\). The upper bound is \((1+\delta)^f\le(1-f\delta)^{-1}\le1+2f\delta\) when \(f\delta\le1/2\). Insert(LT21). For the literal tensor of probability dilation projections the squared norm is instead

\[
             \prod_{v\mid T}\frac{1-c_v}{q_v-1}.
 \tag{LT23}
\]

Thus the local projection cost is explicit and nonzero. It cannot be discarded merely because the Gram correction is small.

\subsection{7. What these local statements settle, and the remaining exact test}\label{what-these-local-statements-settle-and-the-remaining-exact-test}

At G, LIFT32(FB9--FB17) already gives a positive finite restriction eigenvalue \(d_\Pi m_\sigma/d_\sigma\) on every permitted generic initial type, subject to the actual global root-of-unity compatibility supplied by DO043. At T, (LT7) identifies the actual one-dimensional newform line, and (LT19--LT23) show that its finite compact type survives in the fixed translated spherical vector. In particular no local annihilation follows from the \textbf{compact} G/T type or from the unweighted T norm operator.

There is also no abstract representation-level split-place branching obstruction. For two generic tempered representations of \(\mathrm{GL}_{a+1}(F_v)\) and \(\mathrm{GL}_a(F_v)\), the Rankin--Selberg functional at \(s=1/2\) is invariant under the diagonal \(\mathrm{GL}_a\). Booker, Krishnamurthy, and Lee(2020), Theorem1.1 and section3, give a tensor on which it is a nonzero formal Euler factor. That factor is finite and nonzero at \(1/2\) for tempered parameters. The fractional-ideal continuation makes the functional well defined at this point on the whole smooth tensor product. This proves existence of a nonzero local invariant functional. It does not evaluate it on the prescribed vectors.

For the actual T data the outstanding local scalar is precisely

\[
 \mathcal L_v\bigl(v_{H,\mathrm{new}},
                         \pi_G(\tau_u)v_{G,\mathrm{sph}}\bigr).
 \tag{LT24}
\]

Its relevant compact projection has the nonzero norm(LT20), but a nonzero invariant functional can still vanish on this specific copy of a compact type. The restriction of an infinite-dimensional representation to \(K_a\) can contain further copies of that type; Schur's lemma on one finite copy does not fix the functional's scalar.

The explicit test in Booker--Krishnamurthy--Lee(3.1), in the present conductor-one case, uses the \textbf{upper} last-column average

\[
 q^{-(a-1)}\sum_{y\in k^{a-1}}
 \pi_G\!\begin{pmatrix}I_a&\varpi^{-1}(y,1)^t\\0&1\end{pmatrix}
                                      v_{G,\mathrm{sph}}.
 \tag{LT25}
\]

Their(3.2) and Lemma3.1 evaluate its pairing with the newform as \(L(s,\pi_G\times\pi_H)\) times a nonzero Gauss sum divided by \([K_a:J_u]\). That is a genuine test but has not been identified with(LT3) in the fixed local period. Transposing a matrix alone also transports the embedding, Whittaker convention and newform line; none of those changes is silently made here.

Finally, even nonvanishing of(LT24) at every T place would still leave the global central L-value in DO046(AP13), the prescribed other local vectors, and the exact fixed-Schwartz mixed seesaw comparison. The bounded bottom existence of K051 and vertical surjectivity of LIFT31 do not imply that any of these scalars is nonzero. No finite-dimensional rank statement about the earlier theta image is inferred from(LT19).

\subsection{8. Primary references and verification scope}\label{primary-references-and-verification-scope}

Gross, B. H. (1998). On the Satake isomorphism. In A. J. Scholl \& R. L. Taylor(Eds.), \emph{Galois representations in arithmetic algebraic geometry} (London Mathematical Society Lecture Note Series254, pp.223--237). Cambridge University Press. \url{https://people.math.harvard.edu/~gross/preprints/sat.pdf} . Exact used author locators: (2.5) for maximal volume one; (3.14), p8, for the minuscule transform; Proposition6.4 and the next paragraph, p11, for unitary tempered parameters.

Jacquet, H., Piatetski-Shapiro, I. I., \& Shalika, J. (1981). Conducteur des représentations du groupe linéaire. \emph{Mathematische Annalen,256},199--214. \url{https://doi.org/10.1007/BF01450798} . Exact used locator: Theorem5.1, printed211. Primary article copy: \url{https://lesesvre.perso.math.cnrs.fr/newforms-references/jpss.pdf} .

Jacquet, H. (2012). A correction to Conducteur des représentations du groupe linéaire. \emph{Pacific Journal of Mathematics,260}(2),515--525. \url{https://doi.org/10.2140/pjm.2012.260.515} . Exact used locator: Theorem1, printed517, correcting the essential vector existence used in the 1981 theorem. \url{https://msp.org/pjm/2012/260-2/pjm-v260-n2-p09-s.pdf} .

Booker, A. R., Krishnamurthy, M., \& Lee, M. (2020). Test vectors for Rankin--Selberg L-functions. \emph{Journal of Number Theory,209},37--48. \url{https://doi.org/10.1016/j.jnt.2019.08.010} . The directly read primary preprint is arXiv:1903.03458v1, Theorem1.1, (3.1), (3.2), Lemma3.1, PDFpp1,7--8: \url{https://arxiv.org/pdf/1903.03458} .

Local PDFs and their retrieval hashes are in \nolinkurl{local_t_branch_refs/manifest.json}. The elementary Gram computation, the actual Mram central-character reduction and the stated complete scope are proved above. The companion certificate is an exact finite diagnostic of matrix elementary divisors, compact actions, character sums, Hecke counts, and actual numerical inequalities. It does not compute an automorphic form, an actual Satake parameter, a local Rankin--Selberg integral, or a global period.

\subsection{9. Completed author check}\label{completed-author-check}

The complete saved proof LT1--LT25 and the full certificate were read and reconstructed. The central-character reduction, conductor-one application, Mackey irreducibility, lower-translation orientation, minuscule neighbor Gram identity, probability/counting normalization and actual product bounds all passed. The four primary PDFs were read at the exact locators above; their local copies and hashes are retained. The BKL upper average is explicitly kept distinct from the actual lower translate.

The certificate passed 15,618 exact assertions at 2026-11-22T05:43:21.475759+00:00, including 719 lower-translate pairs, 13,098 compact-action tests and 243 minuscule-neighbor pairs. It does not claim a conductor computation for an actual unknown automorphic representation, a local-period evaluation or a global nonzero pairing.

\clearpage

\section{ALP-LIFT-39 --- literal vacuum normalization and primitive vertical intertwiners}\label{alp-lift-39-literal-vacuum-normalization-and-primitive-vertical-intertwiners}

Author: /root/lift\_accounting, GPT-6 Astra Pro, ALP\_MACHINE. 22 November 2026 UTC. Inherited reasoning effort is not separately exposed.

\textbf{Author whole-proof check completed; the exact verification record is below.}

This report fixes two previously separate real normalizations for the actual vertical chain of LIFT31. First, the literal source IsVacO asks for norm one on a rational Heisenberg quotient. The raw CM Gaussian in K035/K037 must be multiplied by its explicit reciprocal norm to satisfy that clause. Second, explicit primitive joint-harmonic vectors fix the otherwise free vertical ItY/ItX scalars. Their exact hook norms compute the scales s0,s1 in LIFT38. The higher selected bottom type introduces the known PBW factor of K043 on both sides, so it does not change these intertwiner scales.

These are constructed choices of the real models and intertwiners. They do not identify an arbitrary pre-existing IsItW map with that choice. A polynomial rational structure, an integral image, and the CM arithmetic period are explicitly distinguished below.

\subsection{1. Literal source clauses and conventions}\label{literal-source-clauses-and-conventions}

The recovered primary source is the pinned Alpoge archive under work/research/alpoge\_source. Exact locators are:

\begin{itemize}
\tightlist
\item
  \nolinkurl{Formalism/LocalToGlobal/IndefiniteReal/LocalNewformFunction.lean,} lines205--228: nrm2, IsVacuum, IsVacO. In particular IsVacuum.norm asserts nrm2=1 for the rational-lattice quotient with covolume one.
\item
  The same file, lines116--174: Cartan complex structures, annihilation and creation operators, fockScale, aPlus, orthonormal frames and polynomial Fock vectors.
\item
  LinksChar.lean58--71: IsItW consists of nonzero, boundedness and group/Lie intertwining conditions. It imposes no extremal scalar.
\item
  LinksChar.lean90--94 and277--287: the actual vertical ptY/ptX are pureTensorO, and their ItY/ItX enter the literal IsIo identities.
\item
  \nolinkurl{OrthogonalComplementCharacterisation.lean41–88:} IsIo and IsIoUp, with conjugated-kernel adjoint and the actual IsTPlus convention.
\end{itemize}

The extremal equality in IsIoW, LinksChar29--37, belongs to the auxiliary ioP; it is not silently transferred to the vertical ItW. The choices below therefore provide an additional explicit vertical normalization, consistent with the homogeneous IsItW axioms.

We use the actual ranks

\begin{longtable}[]{@{}lrrrrrr@{}}
\toprule\noalign{}
Host & a=f & r=a+1 & t & b=r+t & N=r+2t & h, with host J\^{}(14h) \\
\midrule\noalign{}
\endhead
\bottomrule\noalign{}
\endlastfoot
J14 & 42 & 43 & 1759 & 1802 & 3561 & 1 \\
J56 & 168 & 169 & 1633 & 1802 & 3435 & 4 \\
\end{longtable}

At a CM role use K044's E, ell and canonical parameter q, with 0\textless=q\textless=ell. Write d=t+1 and delta=2ell+q. The unshifted bottom Harish-Chandra entries are

\[
 p_i=4\ell(a-1-i)\quad(0\le i<a-1),\qquad
 p_{a-1}=-E+\ell,\qquad p_-=-E-\ell-q.
 \tag{PN1}
\]

The minimal compact type is (alpha;beta), where alpha\_i=p\_i-(a-2)/2+i and beta=p\_--a/2. Its primitive highest vector has norm one in the normalized compact metric. The compact middle and final models likewise give their primitive extremals norm one. These are the normalized Fischer/compact conventions already used by R011, K043 and DO048.

For identifying types we use the standard tensor convention of Ichino(2022), Lemma7.10. R015's transpose-inverse and actual downward/ adjoint conjugation transport it to the literal Hom model. Thus a vector called highest in the tensor model may be the corresponding dual extremal in the source model. The transport preserves its Fischer norm, integrality of its polynomial coefficients, and all positive scalars computed below. We do not call both dual vectors highest in the same Borel.

\subsection{2. The actual theta realization is an isometry}\label{the-actual-theta-realization-is-an-isometry}

Let W=X+Y be the rational symplectic space of one vertical Hom kernel, with a rational polarization. Take global additive Haar measures dual for the fixed standard additive character, normalized so that each of X(F)\textbackslash X(A) and Y(F)\textbackslash Y(A) has volume one. Their product is the covolume-one measure on W(F)\textbackslash W(A). Specifying both quotient volumes matters: the product measure alone would not fix the separate Schrödinger L2 normalization. In Schrödinger coordinates, up to the fixed modulus-one trace-one quadratic phase, theta reads

\[
 \Theta_\varphi(x,y)=\psi(\tfrac12\langle x,y\rangle)
       \sum_{\xi\in X(F)}\varphi(\xi+x)
                                \psi(\langle\xi,y\rangle).
 \tag{PN2}
\]

The sign or the trace-one replacement of the half-pairing changes no squared norm. For every Schwartz function used here,

\[
 \boxed{\quad
 \int_{W(F)\backslash W(\mathbb A)}
                |\Theta_\varphi(w)|^2dw
       =\int_{X(\mathbb A)}|\varphi(x)|^2dx.
 \quad}                                                \tag{PN3}
\]

\textbf{Proof.} Choose product fundamental domains for X(F) and Y(F). For fixed x, the functions psi(\textless xi,y\textgreater) are the orthonormal characters on Y(F)\textbackslash Y(A). The dual lattice is precisely X(F) by rational symplectic duality and the standard additive character. Fourier orthogonality therefore makes the y-integral equal to sum\_xi \textbar varphi(xi+x)\textbar². Integrating x unfolds its rational translates to X(A), proving(PN3). For our finite Schwartz and real polynomial Gaussians the lattice series converges normally on the compact adelic quotient; alternatively Parseval followed by Tonelli on the nonnegative diagonal terms gives the same identity. The modulus-one initial phase cancels. A rational symplectic change of coordinates is unitary in the Weil realization and preserves the covolume-one measure, so the computation is independent of the chosen polarization. QED.

The source's IsCovolOne fixes the Haar uniquely, since two additive Haar measures differ by a positive scalar and the fundamental domain must have measure one. Thus(PN3) is the source nrm2, not an additional arbitrary quotient-volume convention.

At every finite place start with the same selfdual lattice vacuum as LIFT31, of norm one. Consequently the literal IsVacO.norm is satisfied by a real vacuum of norm one. Multiplying a raw real Gaussian by its reciprocal norm changes neither its annihilator, its finite lattice-vacuum equations nor its rational theta transformation law. This constructs the normalization required by the literal source; it is not an inference that the raw vacuum was already unit.

For the product-CM Gaussian of K035/K037, let Y\_J=Im(tau\_J) be R004's exact genus-three matrix. At an edge of ranks n,m the raw real vacuum has squared norm

\[
 \nu(n,m)=\det(2Y_J)^{-7hnm}.
 \tag{PN4}
\]

There are14hnm Jacobian factors, each with squared norm integral\_R³ exp(-2pi x\^{}tY\_Jx)dx=det(2Y\_J)\^{}(-1/2). The first vertical source has one negative Hermitian coordinate. Its corresponding factors use the conjugated Gaussian exp(-pi i x\^{}t bar(tau\_J)x), with the same positive matrix Y\_J. Hence the count and norm in(PN4) are unchanged. The actual orthogonal CM-frame transports of K030/K032 and their Weil realizations are unitary. No extra determinant amplitude is inserted into this norm calculation.

The literal unit-vacuum scalar is therefore

\[
 a_{\rm vac}(n,m)=\nu(n,m)^{-1/2}
                    =\det(2Y_J)^{7hnm/2}>0.
 \tag{PN5}
\]

Its phase is chosen positive in the fixed raw product realization; another permitted vacuum phase is retained as a unit scalar.

\subsection{3. The first primitive joint harmonic and its exact types}\label{the-first-primitive-joint-harmonic-and-its-exact-types}

Put

\[
 A_i=p_i+E-a-\frac{t-1}{2}+i\quad(0\le i<a),\qquad
 B=\ell+q-\frac d2,\qquad L_*=\ell-\frac d2.
 \tag{PN6}
\]

Then A is a positive decreasing partition with A\_(a-1)=L\_\emph{\textgreater0, B=L\_}+q\textgreater0. These inequalities are implied by the actual K044 lower gap. All entries are integers. Use the polynomial Fock variables X\_(c,i), 1\textless=c\textless=b,1\textless=i\textless=a, for the positive part, and y\_c,1\textless=c\textless=b, for the negative source coordinate. Define

\[
 P_0(X,y)=p_A(X)y_b^B,\qquad
 p_A(X)=\prod_{i=1}^a
        \det X_{[1,i],[1,i]}^{A_{i-1}-A_i},\quad A_a=0.
 \tag{PN7}
\]

The leading diagonal monomial has coefficient one. Thus P0 is primitive in the ordinary integer polynomial lattice.

It is joint harmonic: each contraction sum\_c partial\_(X\_(c,i)) partial\_(y\_c) kills it, since p\_A uses only the first a target rows and y\_b uses the last row b\textgreater a. It is a highest vector for U(a) and U(b): the usual leading-minor raising operators kill p\_A, and the dual variable y\_b is killed by every target raising operator. On the target its unshifted weight is (A,0\footnote{t},-B), because b=a+t+1.

The first tensor-model source vacuum weights are b/2-E on U(a) and -b/2-E on U(1). Adding the weights (A;-B) gives exactly (alpha;beta). The target vacuum shift is

\[
                     v_0=(a-1)/2-E+\eta,
                         \qquad\eta=\pm1/2.
 \tag{PN8}
\]

Adding rho\_b to (A,0\footnote{t},-B)+v0 therefore gives the decreasing Harish-Chandra list

\[
     (p_0,\ldots,p_{a-1},B_t(-E),p_-)+\eta.
 \tag{PN9}
\]

This is literally DataTheta.Dpar1; the separation inequalities order its three displayed pieces as written. This verifies the actual characters and the odd vacuum shift, rather than merely matching a polynomial degree. Ichino's joint-harmonic correspondence and the accepted real lift R015 identify P0 with the lowest compact extremal of the actual bottom discrete series tensored with the actual middle compact extremal. Multiplicity one fixes that pure joint-extremal line.

Let H\_lambda denote the ordinary hook product. R011's exact Fischer convention gives

\[
                 \|P_0\|_{\rm F}^2=H_A\,B!.
 \tag{PN10}
\]

Indeed the two variable sets are disjoint, the primitive bideterminant norm is H\_A, and the monomial norm is B!.

\subsection{4. The second primitive joint harmonic}\label{the-second-primitive-joint-harmonic}

Write Delta=delta-d=2ell+q-d=L\_*+B and form the positive partition of length b

\[
 \Lambda=(A_0+\Delta,\ldots,A_{a-1}+\Delta,
                              \Delta^{[t]},L_*).
 \tag{PN11}
\]

The source is now compact of rank b and the target compact of rank N=b+t. In the standard matrix Fock model take

\[
                         P_1(Z)=p_\Lambda(Z).
 \tag{PN12}
\]

It is primitive and joint highest, and

\[
                         \|P_1\|_{\rm F}^2=H_\Lambda.
 \tag{PN13}
\]

Here is a direct character check. If h1\^{}0 denotes the list in (PN9) before its common eta shift, then

\[
 \Lambda_i=(h1^0)_i+E+2\ell+q-b-\frac{t-1}{2}+i.
 \tag{PN14}
\]

For its first a entries this is A\_i+Delta. On the next t entries it is Delta, and on the final entry it is L\_*. The second source vacuum is N/2-E-2ell-q+eta, so subtracting it from the middle highest weight gives(PN14). The target vacuum is b/2-E-2ell-q. Adding target rho\_N now gives h1\^{}0 followed by B\_t(-E-2ell-q), which is precisely Dpar2. The eta shifts have cancelled in the source difference, exactly as R015 requires. This supplies the second actual joint-highest line and its full partition, not an unrestricted compact type with an adjustable character.

\subsection{5. Explicit primitive ItY/ItX and the exact scalar}\label{explicit-primitive-ityitx-and-the-exact-scalar}

For each label realize the bottom representation on its K043 polynomial Harish-Chandra module, normalized so the primitive minimal compact extremal v0 has norm one. Give the two compact models the normalized Fischer metric with primitive extremals v1,v2 of norm one. In the standard tensor convention identify the corresponding joint lines with P0 and P1. Transport to the actual Hom/downward convention of R015 and call the resulting corresponding extremals by the same labels.

By compact-partner Howe multiplicity one there is a unique vertical intertwiner on the selected isotypic component satisfying

\[
                 T_0^{\rm prim}(P_0,v_1)=v_0,\qquad
                 T_1^{\rm prim}(P_1,v_2)=v_1.
 \tag{PN15}
\]

Extend it by zero on the other compact-partner isotypic components. These maps satisfy IsItW: nonzero follows from(PN15), the group and Lie identities from the intertwining construction, and boundedness for each fixed polynomial from its finite-rank tensor realization. The adjoint TPlus is then the actual partial Hilbert adjoint in IsTPlus. No independent scalar is selected for that adjoint.

The construction is algebraic on polynomial Harish-Chandra modules. The first module is generated by its rational minimal compact type under the ordinary matrix-unit PBW operators. Sending its joint extremal tensor to the integer polynomial P0 defines the rational intertwining embedding. The second is the rational Cauchy embedding normalized by P1. On each polynomial degree the compact-partner isotypic projector is rational: its image is a rational subspace, and the Fischer Gram matrix and its orthogonal projector have rational coefficients. Inverting the embeddings on their images and contracting gives(PN15) over the corresponding Chevalley/Fock field (adjoining i if the selected real Lie-basis convention uses it). This assertion is rationality, not an assertion that every such inverse has integral coefficients. Normalizing a primitive compact extremal to norm one divides its invariant Fischer form by a rational hook norm; the Gram matrices and their inverses remain rational. This rational structure is the abstract polynomial/Chevalley one. It does not declare the CM Schrödinger realization or its positive vacuum amplitude algebraic in the same coefficient field.

Use the literal unit vacuum from Section2, and define the exact positive integers

\[
 N_0=\prod_\tau H_{A_\tau}B_\tau!,\qquad
 N_1=\prod_\tau H_{\Lambda_\tau},\qquad \mathcal H=N_0N_1.
 \tag{PN16}
\]

Let I\_j\^{}iso be the isometric selected oscillator contraction used in LIFT38, with the positive phase fixed by the joint extremal. The image of the unit tensor under the inverse isometry is P\_j/sqrt(N\_j), since both row extremals have norm one. Therefore I\_j\^{}iso(P\_j,v\_(j+1))=sqrt(N\_j)v\_j. Equation(PN15) gives

\[
 \boxed{\quad
 T_j^{\rm prim}=N_j^{-1/2}I_j^{\rm iso},\qquad
 |s_0s_1|^2=\mathcal H^{-1}.
 \quad}                                                \tag{PN17}
\]

If the raw product vacuum is used instead, while the polynomial condition(PN15) is held fixed, its physical joint-extremal norms are nu(r,b)N0 and nu(b,N)N1. Thus the exact raw formula is

\[
 |s_0s_1|^2=
       [\nu(r,b)\nu(b,N)\mathcal H]^{-1}.
 \tag{PN18}
\]

These are different explicitly stated normalizations of the actual kernel map. The unit-vacuum case is the one selected by IsVacO.norm. A prescribed different ItW map is a scalar multiple of(PN15), and that scalar must still be retained; this theorem does not alter it.

Here is the fixed-test check required when using(PN18) in LIFT38. Write F\_raw(P)=sqrt(nu) F\_unit(P) in the same physical oscillator Hilbert space. Repinning(PN15) gives T\_raw=nu\^{}(-1/2)T\_unit. For the physical unit joint tensor the Rallis Gram factor is still 1/d\_real, while its contraction scalar is1/sqrt(nu N\_j); the multiplicity-map energy is therefore nu N\_j/d\_real. Alternatively, the raw primitive polynomial has Gram nu N\_j/d\_real and its T\_raw contraction is exactly1, giving the same answer. Thus the raw Gram is not multiplied into LIFT38 a second time. Merely rewriting a fixed physical T in raw polynomial coordinates, without repinning(PN15), would not change the multiplicity map at all.

\subsection{6. The higher selected bottom type and K043}\label{the-higher-selected-bottom-type-and-k043}

For an actual word let lambda be its selected bottom compact type, and let j\_lambda be a Fischer-isometric Pieri inclusion into W\_alpha tensor Sym\^{}D(C\^{}a). K043 supplies the exact positive number

\[
 D_{\rm PBW}=
  \prod_{i=0}^{a-2}(v_i+b_1+1)_{u_i-v_i-1/2}
                 (b_1-b_0+1)_{u_{a-1}+b_0-1/2}.
 \tag{PN19}
\]

Its PBW creation map satisfies \textbar\textbar C\_DS(j\_lambda w)\textbar\textbar²= D\_PBW\textbar\textbar w\textbar\textbar² at one label. For all labels put D\_PBW,all=product\_tau D\_PBW,tau; in(PN20) and its subsequent normalization D\_PBW denotes this product, consistently with N0 from(PN16). Apply the same matrix-unit operators to the first Fock joint-extremal embedding, on the source factor only. Denote this polynomial by Phi\_lambda(w). Equivariance and(PN15) give

\[
 T_0^{\rm prim}(\Phi_\lambda(w),v_1)
                  =C_{\rm DS}(j_\lambda w),\qquad
 \|\Phi_\lambda(w)\|^2=N_0D_{\rm PBW}\|w\|^2.
 \tag{PN20}
\]

The norm identity follows from the invariant oscillator embedding, whose scalar is N0 on the entire irreducible unitary tensor product, as fixed on its extremal line by(PN10). The source representation and target middle factor commute, so no additional target compact projection is inserted here.

Dividing both vectors in(PN20) by sqrt(D\_PBW) gives the isometric selected compact type and the same scalar s0 from(PN17). Thus the PBW norm does not become a new multiplicative scale of T0. The partitions A,B,Lambda use v,b,e only and are independent of the binary word u. The word affects the selected compact type and(PN19), not the primitive vertical normalization. A non-isometric Pieri inclusion retains its factor j* j as in K043(PB8).

\subsection{7. Exact hook arithmetic without expanding the polynomials}\label{exact-hook-arithmetic-without-expanding-the-polynomials}

The scalar mathcal H is explicitly computable from finite products. For a partition lambda of length n, including trailing zeros,

\[
 H_\lambda=
 \frac{\prod_{i=0}^{n-1}(\lambda_i+n-1-i)!}
      {\prod_{0\le i<j<n}(\lambda_i-\lambda_j+j-i)}.
 \tag{PN21}
\]

This is the standard hook formula already used by R011. Its valuation at11 is the sum of the corresponding Legendre floor sums minus the valuations of the displayed linear differences. No factorial of size E needs to be materialized.

There is a useful shorter exact expression for this family. Put C\_i=A\_i+a-1-i=p\_i+E-d/2. The middle block in(PN11) gives

\[
 \frac{H_\Lambda}{H_A}=
 \left[\prod_{i=0}^{a-1}
     \frac{(C_i+\delta)!}
          {(C_i+t)!(C_i+B+d)}\right]
 \frac{L_* !\,B!}{(B+t)!}
 \frac{\prod_{u=1}^{t}(\Delta+u)!}
      {\prod_{u=1}^{t-1}u!}.
 \tag{PN22}
\]

To prove this, separate the Vandermonde in(PN21) into its first-a, constant-t, and final-one blocks. The first Vandermonde cancels that of H\_A. The first-versus-middle product is product\_i (C\_i+t)!/C\_i!, the first-versus-last product is product\_i(C\_i+B+d), and the middle-versus-last product is (B+t)!/B!. The middle Vandermonde is product\_(u=1)\^{}(t-1)u!. The numerator first entries are(C\_i+delta)!, the middle entries (Delta+u)!, and the last entry L\_*!. This gives(PN22) exactly.

These hook valuations describe the normalization scalar, not the content of a theta image. In particular a primitive polynomial can have a highly divisible Fischer norm. No integral preimage is inferred from(PN17) or from rationality of the intertwiner.

\subsection{8. Exact relation to K037's CM period and source frame}\label{exact-relation-to-k037s-cm-period-and-source-frame}

Keep every raw K035/K037 identity as written. On an ordinary positive product at an edge n,m its raw common period is C\_theta\^{}(14hnm), with the same algebraic coordinate factors cD and source creation factors lambda\_tau. For that existing reading, multiplying only the vacuum by(PN5) gives exactly

\[
 \Theta_{{\rm unit},P}=
   a_{\rm vac}(n,m)\,C_\theta^{14hnm}
       \left[\prod_\tau
          (\lambda_\tau/cD_\tau)^{\deg_\tau P}\right]
                        \mathcal J_{cD,F}(P).
 \tag{PN23}
\]

All finite phases, finite averages and raw finite Schwartz functions inside F are unchanged. This is simply linearity in the vacuum, applied to the exact K035(PF27--PF29)/K037(CR28) identity; it is not a new algebraic normalization of C\_theta.

For the two vertical edges the additional positive period factor is

\[
 a_{\rm vac}(r,b)a_{\rm vac}(b,N)
          =\det(2Y_J)^{7h b^2},\qquad r+N=2b.
 \tag{PN24}
\]

For the first vertical edge there are a positive source coordinates and one negative coordinate. The factorwise conjugate CM model therefore has raw period C\_theta\^{}(14hab) times bar(C\_theta)\^{}(14hb), rather than replacing the latter factor by C\_theta without its phase. Its negative-column creation/frame factors are the conjugates of the corresponding positive ones. This dictionary is obtained by complex-conjugating the entire K035 identity on that factor, including its Heisenberg phase and finite input transport, and tensoring the identities. Complex conjugation sends exp(pi i x\^{}t tau\_J x) to exp(-pi i x\^{}t bar(tau\_J)x), and commutes with polynomial differentiation after replacing the coordinate by its conjugate. This proves the asserted factorwise relation, with the unchanged positive scalar(PN5). On the second vertical edge all coordinates are positive and(PN23) applies directly.

In particular the source norm correction is the same(PN24) for both readings. A subsequent identification of the chosen algebraic conjugate CM frames with every literal arithmetic row pin is still required. Neither a positive-copy derivative formula nor its 11-integral gauge properties are transferred to a negative copy without this conjugate-frame transport. The scalar normalization proved here does not assert those additional pins.

The creation factors lambda\_tau are unchanged by making the vacuum unit; their identity \textbar lambda\_tau\textbar² pi gamma\_tau=1 still normalizes the Fischer creation operators. The new positive period scalar in(PN23) must be retained in any analytic/arithmetic comparison. It is not assigned an11-adic valuation.

\subsection{9. Consequence for the actual inverse cost and scope}\label{consequence-for-the-actual-inverse-cost-and-scope}

For the literal unit vacuum and the constructed primitive maps, substitute(PN17) into LIFT38. The exact positive lower singular value squared is multiplied by mathcal H relative to its s=1 isometric-intertwiner normalization, and

\[
 \boxed{\quad
 \|R_{\rm prim}\|
 <\mathcal H^{-1/2}\,
    2^{(H_{\rm tot}+8h_Q)/2}\operatorname{Qnorm}^{2rt}.
 \quad}                                                \tag{PN25}
\]

Here H\_tot and every finite symbol are LIFT38's explicit table; the unit-finite analogue uses its corresponding table and Q powers. The raw-vacuum version has the extra multiplier(PN24) on the right-hand side, exactly as(PN18) requires: its squared singular value has the factor nu(r,b)nu(b,N)mathcal H. The fixed-test calculation following(PN18) already accounts for the raw test Gram, so no further vacuum norm is inserted. There is no remaining arbitrary positive real intertwiner scalar after the choice(PN15).

This is a norm statement for those constructed multiplicity maps. It is not a claim of a net DP44 saving from a large hook product: changing a scalar normalization also changes its arithmetic lattice comparison. The literal IsVacO normalization, the primitive rational Fock normalization, and the CM frame/period conversion are now explicitly related. Full arithmetic pin compatibility, integral saturation of the vertical image, and the terminal detector remain separate obligations.

\subsection{10. Sources and verification}\label{sources-and-verification-2}

Ichino, A. (2022). Theta lifting for tempered representations of real unitary groups. Advances in Mathematics,398,108188. Author primary PDF \url{https://arxiv.org/pdf/2008.06174}, section4.1 and Lemma7.10, printedpp.30--31; direct web reading turn459255364view0 and turn895096777view0. The precise contragredient convention and actual selected real correspondence are already proved in R015.

Nakahama, R. (2016). Norm computation and analytic continuation of vector valued holomorphic discrete series representations. Journal of Lie Theory,26,927--990. Equation(5.7), Theorem5.1, pp.960--961. The exact selected PBW application is the accepted K043(PB1--PB10).

The hook normalization is the accepted R011(R11.5--R11.7), with its Fulton--Harris primary anchors. The actual unit creation/CM comparison is K035(PF22--PF29) and K037(CR28,CR35--CR36). Exact source file locators are in Section1. The proof of(PN3) is given here explicitly.

\subsubsection{Verification record}\label{verification-record-5}

The author read the complete proof, checked the literal IsVacuum, IsItW, IsIo and IsTPlus clauses, and rechecked both joint partitions and their vacuum characters against the accepted R015 transport. The review includes the separate additive quotient measures, the rational normalized compact metrics, and the fixed-physical-test versus raw-primitive-test calculation following(PN18). In particular the raw-vacuum Gram is counted exactly once. The actual two-edge inverse comparison retains the conjugate-column CM phase and does not identify a complex norm estimate with an integral image index.

The companion \nolinkurl{vertical_primitive_fock_certificate.py} checks finite Fourier orthogonality, exact small Fock polynomials and PBW norms, the block hook identity, the actual large partition lists and 11-valuations, and the independent fixed-test scaling identity. These diagnostics do not evaluate an automorphic theta period, prove an integral Smith minor, or instantiate the full source D7.

The companion exact certificate checks finite Fourier orthogonality, small joint-harmonic/highest polynomials, hook norms and(PN22), all actual type lists and character shifts, and the normalization/PBW identities. It does not evaluate an automorphic period or certify an integral theta-image minor. Whole-proof check events are separate.

\clearpage

\section{AM-K053 --- the prescribed lower T translate is a nonzero local newform test}\label{am-k053-the-prescribed-lower-t-translate-is-a-nonzero-local-newform-test}

Author: GPT-6 Astra Pro, ALP\_MACHINE \nolinkurl{/root/kummer_extension}. 22 November 2026 UTC. Inherited reasoning effort is not separately exposed.

\textbf{Author whole-proof review PASS: 2026-11-22T06:03:20.060168+00:00.}

This finding closes the specific local test left open in K052(LT24). For every actual initial irreducible Mram constituent and every actual spherical bottom constituent, the unchanged lower-unipotent T translate pairs nontrivially with the actual initial newform under the local general-linear Rankin--Selberg functional at \(s=1/2\).

The qualitative statement is proved first. A separate normalization argument gives the exact modulus in the \(W(1)=1\) Whittaker convention. Neither assertion supplies a nonzero global central L-value, an actual terminal arithmetic pairing, or an integral vertical preimage.

\subsection{1. Fixed local data and the theorem}\label{fixed-local-data-and-the-theorem}

Retain exactly K052(LT1--LT7). Write \(F\) for the local field, \(O\) for its integer ring, \(k\) for its residue field of cardinality \(q\), and \(\varpi\) for the retained uniformizer. Let

\[
 H=\mathrm{GL}_a(F),\qquad G=\mathrm{GL}_{a+1}(F),\qquad
 h\longmapsto\operatorname{diag}(h,1).
 \tag{NT1}
\]

The actual ranks are \(a=42,168\). Let \(\psi\) have conductor \(O\). Let \(\Pi\) be an irreducible tempered spherical representation of \(G\), and let \(\pi\) be an irreducible tempered generic representation of \(H\) satisfying

\[
 c(\pi)=c(\omega_\pi)=1.
 \tag{NT2}
\]

For the actual family these hypotheses have already been proved: DO046 supplies temperedness and genericity, K052(LT7) supplies (NT2), and K050/K051 supply actual existing bottom constituents throughout the allowed bounded gap interval. We make no new global existence assumption.

Conjugate the retained covector integrally to \(e_a\), transporting the newform and the embedding together. Put

\[
\begin{split}
 K&=\mathrm{GL}_a(O),\\
 K_r&=\{h\in K:e_a\bar h=e_a\},\\
 K_c&=\{h\in K:\bar h e_a^t=e_a^t\},\\
 t&=\operatorname{diag}(\varpi I_{a-1},1),\\
 \tau&=I_{a+1}+\varpi^{-1}E_{a+1,a},\qquad
 u_\pm=I_{a+1}\pm\varpi^{-1}E_{a,a+1}.
\end{split}
 \tag{NT3}
\]

Here \(K_r\) is the original strict row-covector level, not the larger projective stabilizer. Direct block multiplication gives

\[
 K_c=tK_rt^{-1},\qquad
 \operatorname{Stab}_K\!\left(\Pi(u_\pm)v_{\mathrm{sph}}\right)
 \supset K_c.
 \tag{NT4}
\]

The second inclusion follows because \(h e_a^t-e_a^t\in\varpi O^a\) for \(h\in K_c\); the corresponding change in the upper column is integral and fixes the spherical vector.

Let \(W_\Pi^0\in\mathcal W(\Pi,\psi)\) and \(W_\pi^0\in\mathcal W(\pi,\psi^{-1})\) be the spherical and essential Whittaker functions normalized by value one at the identity. Define

\[
 Z(\Pi,\pi)=
 \left.
 \int_{N_a(F)\backslash H}
 W_\Pi^0(\operatorname{diag}(h,1)\tau)\,
 W_\pi^0(h)|\det h|^{s-1/2}\,dh
 \right|_{s=1/2}.
 \tag{NT5}
\]

Haar measures are those of Booker, Krishnamurthy, and Lee (2020, section2): \(K\) and \(N_a(O)\) each have volume one, and the quotient measure is the resulting right-invariant measure. Evaluation in (NT5) can be read through the holomorphic continuation of the local zeta integral. All local L-factors below are regular and nonzero at \(1/2\).

\textbf{Theorem.} For these exact vectors,

\[
 Z(\Pi,\pi)\ne0,\qquad
 \boxed{\left|\frac{Z(\Pi,\pi)}{L(1/2,\Pi\times\pi)}\right|
          =\frac{q^{a/2}}{q^a-1}.}
 \tag{NT6}
\]

The nonzero assertion is invariant under any nonzero rescaling of the Whittaker functionals or of the actual newform vectors. The displayed modulus uses specifically the conventions in (NT5).

\subsection{2. Conductor one with ramified determinant leaves exactly one tame character}\label{conductor-one-with-ramified-determinant-leaves-exactly-one-tame-character}

We record the precise classification needed for the proof; it also excludes the possible Steinberg conductor-one alternative.

\textbf{Lemma.} The representation in (NT2) is the normalized induction

\[
 \pi=\chi\times\chi_1\times\cdots\times\chi_{a-1},
 \tag{NT7}
\]

where \(\chi\) is a unitary character of conductor one and the \(\chi_j\) are unramified unitary characters. In particular, with \(\beta_j=\chi_j(\varpi)\),

\[
 L(s,\pi)=\prod_{j=1}^{a-1}(1-\beta_jq^{-s})^{-1},
 \qquad |\beta_j|=1.
 \tag{NT8}
\]

\textbf{Proof.} Use the usual Weil--Deligne parameter \((r,N)\) of \(\pi\) on \(V\), including monodromy. The compatibility of central characters and epsilon factors with local Langlands gives determinant conductor one and total Artin conductor one. Deligne's precise formula is

\[
 1=a(r,N)=\operatorname{Sw}(r)
              +a-\dim(\ker N)^{I_F}.
 \tag{NT9}
\]

This is Deligne (1973), (8.12.1), printed572. The ramified determinant implies that inertia is not trivial, hence

\[
 \dim(\ker N)^{I_F}\le\dim V^{I_F}\le a-1.
\]

Every term in (NT9) is nonnegative. Therefore Swan conductor is zero and both displayed dimensions equal \(a-1\). In particular \(N\) vanishes on \(V^{I_F}\). The finite inertia action is semisimple; its complement is a one-dimensional nontrivial inertia character. This complement is stable under the Weil group, because \(I_F\) is normal and its invariant averaging projector commutes with the Weil action. The nilpotent \(N\), which commutes with inertia, also preserves that line and hence vanishes on it. Thus \(N=0\).

Frobenius semisimplicity now splits \(V^{I_F}\) into \(a-1\) unramified Weil characters; the remaining line is a tame ramified character. Temperedness makes all these characters unitary. The local Langlands classification then gives (NT7), and its standard local L-factor is (NT8). The induced representation is irreducible: ratios of unitary characters cannot equal \(|\cdot|^{\pm1}\), the reducibility values for general-linear rank-one induction.

The local Langlands inputs are the established correspondence and its standard extension from supercuspidals by the discrete-series classification, not a claim that a semisimplified Weil parameter alone detects monodromy. Scholze (2013), Theorem14.1(i),(iii)--(v), provides the character, dual, central-character, and factor compatibilities on the supercuspidal generators; the extension is the usual classification described explicitly in BKL section2.1. Deligne (8.12.1) is applied to the full \((r,N)\) above before proving that \(N=0\). This distinction avoids using Scholze's deliberately monodromy-forgetting notation as a conductor formula. \(\square\)

\subsection{3. The one-dimensional row/column projection}\label{the-one-dimensional-rowcolumn-projection}

For any compact subgroup \(J\) write \(e_J\) for probability averaging in the given unitary representation. The newform theorem and (NT2) give

\[
 \dim\pi^{K_r}=\dim(\pi^\vee)^{K_r}=1,\qquad
 \dim(\pi^\vee)^{K_c}=1.
 \tag{NT10}
\]

Let \(v^\vee\) be the \(W(1)=1\) newform in \(\pi^\vee\). Define scalars by

\[
 e_{K_r}\pi^\vee(t)v^\vee=\lambda_+v^\vee,\qquad
 e_{K_r}\pi^\vee(t^{-1})v^\vee=\lambda_-v^\vee.
 \tag{NT11}
\]

These use probability averages. They are not the unnormalized double-coset Hecke operators in Miyauchi's notation. Unitarity gives \(\lambda_-=\overline{\lambda_+}\).

The quotient \(K_r/(K_r\cap K_c)\) has the \(q^{a-1}\) representatives

\[
 h_y=\begin{pmatrix}I_{a-1}&y\\0&1\end{pmatrix},
 \qquad y\in k^{a-1}.
 \tag{NT12}
\]

The vector \(\pi^\vee(t)v^\vee\) is \(K_c\)-fixed, so (NT12) computes the first probability average in (NT11). Every \(h_y\) is an integral upper unipotent. Its Whittaker character equals one. Evaluating at the identity therefore yields

\[
 \lambda_+=W_{\pi^\vee}^0(t)
 =q^{-(a-1)/2}\prod_{j=1}^{a-1}\beta_j^{-1},
 \qquad
 \lambda_-=q^{-(a-1)/2}\prod_{j=1}^{a-1}\beta_j.
 \tag{NT13}
\]

The diagonal value is exactly Miyauchi (2014), Theorem4.1, printed21--22: the Schur polynomial for \((1^{a-1},0)\) at \((\beta_1^{-1},\ldots,\beta_{a-1}^{-1},0)\) is the product of the nonzero parameters. Also \(\delta_{B_a}(t)=q^{-(a-1)}\). Hence both eigenvalues are nonzero and have modulus \(q^{-(a-1)/2}\).

Conjugating the probability projector in (NT11) gives the key exact identity

\[
 \boxed{e_{K_c}v^\vee
      =\pi^\vee(t)e_{K_r}\pi^\vee(t^{-1})v^\vee
      =\lambda_-\,\pi^\vee(t)v^\vee.}
 \tag{NT14}
\]

Consequently, if \(\ell\) is any \(H\)-invariant bilinear functional and \(w\) in its \(G\)-factor is \(K_c\)-fixed, then

\[
 \ell(w,v^\vee)=
 \lambda_-\,\ell(w,\pi^\vee(t)v^\vee).
 \tag{NT15}
\]

This is a one-dimensional projection identity. It makes no multiplicity-one assertion about all copies of a finite \(K\)-type inside the restriction of the \(G\)-representation.

\subsection{4. Correct transpose-dual convention}\label{correct-transpose-dual-convention}

Let \(w_j\) be the anti-diagonal permutation matrix of size \(j\) and put

\[
 \widetilde W(g)=W(w_j\,{}^tg^{-1}).
 \tag{NT16}
\]

Right translation is \((\rho(g)W)(x)=W(xg)\); thus the stabilizer of \(\rho(g)W\) is \(g\operatorname{Stab}(W)g^{-1}\). Since transpose-inverse exchanges \(K_r\) and \(K_c\), \(\widetilde W_\pi^0\) is \(K_c\)-fixed. Therefore \(\rho(t^{-1})\widetilde W_\pi^0\) is a nonzero \(K_r\)-fixed newform for \(\pi^\vee\), and

\[
 \widetilde W_\pi^0=C_\pi\,\rho(t)W_{\pi^\vee}^0,\qquad
 C_\pi=\widetilde W_\pi^0(t^{-1})=W_\pi^0(w_at)\ne0.
 \tag{NT17}
\]

This is the direction \(t^{-1}\) for the dual newform and \(t\) on the right side of (NT17). Reversing these directions would give an incorrect upper/lower comparison.

Similarly, direct multiplication gives

\[
 \widetilde{\rho(\tau)W_\Pi^0}
       =\rho({}^t\tau^{-1})\widetilde W_\Pi^0
       =\rho(u_-)\,W_{\Pi^\vee}^0.
 \tag{NT18}
\]

The last equality uses \(w_{a+1}\in\mathrm{GL}_{a+1}(O)\), so the transpose-dual spherical Whittaker function still has value one at the identity. The additive characters in the dual pair are \(\psi^{-1}\) on \(G\) and \(\psi\) on \(H\).

\subsection{5. The qualitative nonzero test}\label{the-qualitative-nonzero-test}

Write \(\ell^\vee\) for the central Rankin--Selberg functional of the dual pair. This functional is \(H\)-invariant at \(s=1/2\). The local functional equation in the anti-diagonal convention (NT16) says

\[
 \ell^\vee\bigl(\rho(u_-)W_{\Pi^\vee}^0,
                        \widetilde W_\pi^0\bigr)
   =\omega_\pi(-1)^a\,
      \gamma(1/2,\Pi\times\pi,\psi)\,Z(\Pi,\pi).
 \tag{NT19}
\]

The exact source is JPSS (1983), section2.7, also stated with the same right-translation and anti-diagonal matrices in Jacquet and Liu (2018), Theorem2.3 and (2.2), author PDFpp6--7. Their extra matrix \(w_{a+1,a}=\operatorname{diag}(I_a,w_1)\) is the identity here. The factor \(\gamma(1/2,\Pi\times\pi,\psi)\) is nonzero.

Let \(w_-=\rho(u_-)W_{\Pi^\vee}^0\). It is \(K_c\)-fixed by (NT4). Combining (NT17), (NT15), and (NT19) gives

\[
 Z(\Pi,\pi)=
 \frac{C_\pi}
 {\omega_\pi(-1)^a\gamma(1/2,\Pi\times\pi,\psi)\lambda_-}
           \ell^\vee(w_-,W_{\pi^\vee}^0).
 \tag{NT20}
\]

It remains to evaluate exactly the last factor. Averaging \(w_+\) over \(K_r\) gives

\[
 e_{K_r}w_+
 =q^{-(a-1)}
     \sum_{y\in k^{a-1}}
       \Pi^\vee\!\left(
       \begin{pmatrix}I_a&\varpi^{-1}(y,1)^t\\0&1\end{pmatrix}
                 \right)W_{\Pi^\vee}^0.
 \tag{NT21}
\]

Indeed \(h\in K_r\) sends \(e_a^t\) modulo \(\varpi\) to \((y,1)^t\); all \(y\) occur equally often. Integral changes of this vector fix the spherical Whittaker function. Equation (NT21) is precisely the probability average in BKL(3.1) for conductors one, without a changed embedding or a changed newform.

Because the other factor is \(K_r\)-fixed, the average does not change the pairing. BKL(3.2) and Lemma3.1 give

\[
 \ell^\vee(w_+,W_{\pi^\vee}^0)=
 \frac{\mathfrak G(\omega_\pi^{-1},\psi^{-1},\varpi^{-1})}{d_a}
          L(1/2,\Pi^\vee\times\pi^\vee),
 \quad
 d_a=\frac{q^a-1}{q-1},
 \tag{NT22}
\]

where

\[
 \mathfrak G(\omega,\psi_0,\varpi^{-1})
       =\int_{O^\times}\omega(z)\psi_0(\varpi^{-1}z)\,d^\times z,
 \qquad \operatorname{vol}(O^\times)=1.
 \tag{NT23}
\]

The naive L-factor in BKL equals the actual Rankin--Selberg factor because \(\Pi^\vee\) is unramified. The Gauss sum is nonzero since \(\omega_\pi\) has conductor one.

The integral element \(h_-=\operatorname{diag}(I_{a-1},-1)\) sends \(w_+\) to \(w_-\). Its action on the dual newform is \(\omega_\pi^{-1}(-1)\). Invariance therefore gives

\[
 \ell^\vee(w_-,W_{\pi^\vee}^0)
    =\omega_\pi(-1)\ell^\vee(w_+,W_{\pi^\vee}^0).
 \tag{NT24}
\]

Every factor on the right side of (NT20), after (NT22)--(NT24), is nonzero. This proves the qualitative assertion in (NT6).

In particular the actual local test is now proved, rather than inferred from the nonzero compact projection of K052. At the split place the local invariant Hom-space is one-dimensional; thus any nonzero normalization of the local invariant period has the same nonvanishing on this specific tensor. This does not assert a nonzero global period.

\subsection{6. Exact modulus in the stated Whittaker normalization}\label{exact-modulus-in-the-stated-whittaker-normalization}

We verify the scalar \(C_\pi\) separately. Let \(\sigma\) be the unramified tempered principal series of \(\mathrm{GL}_{a-1}(F)\) with every Satake parameter equal to one. Its central character is one. This is an irreducible generic representation; no ratio of its unitary inducing characters is \(|\cdot|^{\pm1}\).

The essential-vector test, proved in Matringe (2013), Corollary3.3, gives the actual L-factor for \(W_\pi^0\) against the normalized spherical Whittaker function of \(\sigma\). The same statement holds for the dual pair. To fix the measure in that corollary's phrase ``appropriate normalisation,'' use its proof in the case \(m=r=a-1\), printed1212, or use Miyauchi's diagonal formula directly. With the probability compact and unit-diagonal measures of (NT5), the Iwasawa sum is \[
 \sum_{\lambda_1\ge\cdots\ge\lambda_{a-1}\ge0}
       s_\lambda(\beta)\,s_\lambda(1,\ldots,1)
                   q^{-s|\lambda|}
   =\prod_{i,j=1}^{a-1}(1-\beta_iq^{-s})^{-1}.
\] The diagonal unit integrals equal one and the constant partition contributes one, so there is no additional Haar scalar. The dual calculation has exactly the same normalization.

In the dual integral insert (NT17). The occurrence of \(t\) amounts to replacing the embedded \((a-1)\)-by-\((a-1)\) variable \(h\) by \(\varpi h\). At \(s=1/2\) there is no determinant weight; the quotient Haar measure is right invariant; and the spherical \(\sigma^\vee\) factor has trivial central character. The functional equation therefore gives exactly

\[
 C_\pi\,L(1/2,\pi^\vee\times\sigma^\vee)
  =\varepsilon(1/2,\pi\times\sigma,\psi^{-1})
            L(1/2,\pi^\vee\times\sigma^\vee).
 \tag{NT25}
\]

No zero is divided out: this local L-value is a finite product of nonzero factors at \(1/2\). Thus

\[
 C_\pi=\varepsilon(1/2,\pi\times\sigma,\psi^{-1}),\qquad
 |C_\pi|=1.
 \tag{NT26}
\]

The modulus follows either from the unitary epsilon-factor identity or directly from (NT7): the only nontrivial factors are the usual normalized tame character Gauss sums, each of modulus one.

For the main pair, use \(\gamma=\varepsilon L^\vee/L\) in (NT20). The full exact identity is

\[
 \boxed{
 \frac{Z(\Pi,\pi)}{L(1/2,\Pi\times\pi)}
   =
 \frac{\omega_\pi(-1)^{1-a} C_\pi
           \mathfrak G(\omega_\pi^{-1},\psi^{-1},\varpi^{-1})}
      {\lambda_-\,d_a\,
           \varepsilon(1/2,\Pi\times\pi,\psi)} .}
 \tag{NT27}
\]

This retains the sign, the normalized newform Hecke factor, the Gauss sum, the projective index, and the epsilon factor. The unitary epsilon factor of the main pair also has modulus one. Finally,

\[
 |\mathfrak G|=\frac{\sqrt q}{q-1},\qquad
 |\lambda_-|=q^{-(a-1)/2},\qquad
 d_a=\frac{q^a-1}{q-1}.
 \tag{NT28}
\]

Their quotient is \(q^{a/2}/(q^a-1)\), proving the second part of (NT6). The Gauss modulus can be checked without any analytic estimate: for a nontrivial multiplicative character \(\nu\) of \(k^\times\) and a nontrivial additive character \(\bar\psi\) of \(k\), the squared modulus of \(\sum_{z\ne0}\nu(z)\bar\psi(z)\) is \[
 \sum_{r\ne0}\nu(r)\sum_{y\ne0}\bar\psi((r-1)y)
       =(q-1)-\sum_{r\ne1}\nu(r)=q.
\] Division by \(q-1\) gives the probability normalization in (NT28).

\subsubsection{Hilbert normalization and the actual family}\label{hilbert-normalization-and-the-actual-family}

Equation (NT27) belongs to the \(W(1)=1\) Whittaker models. If a specified unitary realization gives these vectors Hilbert norms \(h_\Pi,h_\pi\), the same Rankin--Selberg functional on their Hilbert-unit versions has value \(Z/(h_\Pi h_\pi)\). If instead the Whittaker functionals are chosen to send Hilbert-unit newforms to value one, the normalization of the Rankin--Selberg functional has changed accordingly. Neither convention may be substituted silently for a separately fixed normalized local GGP matrix-coefficient integral.

The nonvanishing established here is unaffected. A numerical comparison with that canonical GGP integral must retain its Whittaker/Hilbert and L-factor normalizations. No such comparison constant is silently set equal to one.

\subsection{7. Uniform actual-prime consequence}\label{uniform-actual-prime-consequence}

The local factor has degree \((a+1)(a-1)=a^2-1\), and every nonzero Euler parameter is unitary. Consequently

\[
 (1+q^{-1/2})^{-(a^2-1)}
 \le |L(1/2,\Pi\times\pi)|
 \le (1-q^{-1/2})^{-(a^2-1)}.
 \tag{NT29}
\]

The actual fixed T primes in DO040 satisfy \(q=T>2^{964}\), while \(a\le168\). Hence

\[
 \left|\log |L(1/2,\Pi\times\pi)|\right|
       <2^{-466}.
 \tag{NT30}
\]

Indeed \(a^2-1<2^{15}\), \(q^{-1/2}<2^{-482}\), and \(-\log(1-x)\le 2x\) for \(0\le x\le1/2\). The same estimate bounds the other side of (NT29).

Across the \(f=42\) or \(168\) actual T places, all of residue cardinality \(T\), the product of the \(W(1)=1\) local test moduli therefore satisfies

\[
 \prod_{v\mid T}|Z_v|
 =T^{-af/2}(1-T^{-a})^{-f}
          \prod_{v\mid T}|L_v(1/2,\Pi_v\times\pi_v)|,
 \tag{NT31}
\]

and

\[
 \left|\log\!\left(T^{af/2}\prod_{v\mid T}|Z_v|\right)\right|
          <2^{-450}.
 \tag{NT32}
\]

The exponents \(af/2\) are \(882\) and \(14112\), respectively. These are local Whittaker-test estimates, not a global central L-value estimate or a total source cost.

\subsection{8. Scope, exact sources, and verification}\label{scope-exact-sources-and-verification}

The actual covector, embedding, retained uniformizer, strict initial T level, and spherical bottom level are those already pinned in K052 section1 and LIFT29(BT3--BT8). Conjugating the covector to \(e_a\) uses an integral \(H\)-element and transports all objects together. The opposite split E-component remains the inverse adjoint. No extra translate, changed conductor, principal29 refinement, or replacement initial vector has been introduced.

The earlier proposed direct absorption of the upper average after duality had the direction of \(t\) reversed. It was rejected before this finding was written. Equations (NT16)--(NT20) use the corrected direction, and (NT14) supplies the required column-space projection. This is part of the proof audit rather than a promoted earlier result.

Primary references used:

\begin{itemize}
\tightlist
\item
  Booker, A. R., Krishnamurthy, M., \& Lee, M. (2020). Test vectors for Rankin--Selberg L-functions. \emph{Journal of Number Theory, 209}, 37--48. \url{https://doi.org/10.1016/j.jnt.2019.08.010}. Original author manuscript \url{https://arxiv.org/pdf/1903.03458}: section2, (3.1), (3.2), and Lemma3.1, manuscriptpp3--4,7--9.
\item
  Deligne, P. (1973). Les constantes des équations fonctionnelles des fonctions L. In \emph{Modular Functions of One Variable II}, Lecture Notes in Mathematics349, pp501--597. Springer. \url{https://publications.ias.edu/sites/default/files/Number20.pdf}. Formula(8.12.1), printed572, is the full Weil--Deligne conductor.
\item
  Jacquet, H., Piatetski-Shapiro, I. I., \& Shalika, J. A. (1983). Rankin--Selberg convolutions. \emph{American Journal of Mathematics, 105}(2), 367--464. Section2.7; author-hosted original at \url{https://www.math.columbia.edu/~hj/Rankin\%20Selberg\%20convolutions.pdf}.
\item
  Jacquet, H., \& Liu, B. (2018). On the local converse theorem for p-adic GL(n). \emph{American Journal of Mathematics, 140}(5),1399--1422. \url{https://doi.org/10.1353/ajm.2018.0035}. Exact restatement Theorem2.3 and(2.2), manuscriptpp6--7, with transpose-dual convention onp5: \url{https://archive.ymsc.tsinghua.edu.cn/pacm_download/318/8960-JL-LCT-final.pdf}.
\item
  Matringe, N. (2013). Essential Whittaker functions for GL(n). \emph{Documenta Mathematica, 18},1191--1214. \url{https://doi.org/10.4171/dm/425}. Corollaries3.2--3.3, printed1210--1212; original at \url{https://lesesvre.perso.math.cnrs.fr/newforms-references/matringe.pdf}.
\item
  Miyauchi, M. (2014). Whittaker functions associated to newforms for GL(n) over p-adic fields. \emph{Journal of the Mathematical Society of Japan, 66}(1),17--24. \url{https://doi.org/10.2969/jmsj/06610017}. Lemma3.1 and Theorem4.1, printed20--22: \url{https://www.jstage.jst.go.jp/article/jmath/66/1/66_17/_pdf}.
\item
  Scholze, P. (2013). The local Langlands correspondence for GL(n) over p-adic fields. \emph{Inventiones Mathematicae, 192},663--715. \url{https://doi.org/10.1007/s00222-012-0420-5}. Theorem14.1, author PDFp35: \url{https://people.math.uni-bonn.de/scholze/LocalLanglands.pdf}. Its introductory footnote explicitly suppresses monodromy; our use of the full Weil--Deligne conductor was separated above.
\end{itemize}

The accompanying certificate is a finite diagnostic of the row/column subgroups, the transpose directions, the one-dimensional projector identity, elementary symmetric factor, exact Gauss modulus, and the actual T-prime inequalities. It is not a numerical computation of an actual automorphic constituent or global period. The mathematical nonvanishing proof is (NT7)--(NT24).

Author verification: the complete saved proof NT1--NT32 and the complete certificate were read and reconstructed after the final measure-normalization paragraph was inserted. The corrected transpose direction, the column projector, all probability factors, the W(1)=1 dual scalar, and the scope of the actual T test were checked. The certificate completed 2026-11-22T06:01:34.229211+00:00 with 52,454 exact assertions. No global period was numerically computed.

\clearpage

\section{AM-R-033 --- the full local period at every cuspidal G place}\label{am-r-033-the-full-local-period-at-every-cuspidal-g-place}

Author: GPT-6 Astra Pro, ALP\_MACHINE \nolinkurl{/root}. 22 November 2026 UTC. Reasoning effort is not separately exposed.

\textbf{Author whole-proof review PASS: 2026-11-22T06:14:41.950825+00:00. Independent review pending.}

This finding evaluates the complete p-adic matrix-coefficient integral, with the actual finite generic character and unchanged levels, at all but two of the G places in the genuine family constructed by DO042/043 and K050/051. The two mixed-factor places are explicitly retained as exceptions. Compact branching alone was insufficient; compact induction now makes the whole integral a finite probability average at the other places.

\subsection{1. Actual data and the statement}\label{actual-data-and-the-statement}

Fix one of the full-degree places in the DO042/043 goodTuple. Let F be its local field, O its integer ring, k its residue field and q=\textbar k\textbar. Set \[
 n=a+1,\quad G=\operatorname{GL}_n(F),\quad
 H=\operatorname{GL}_a(F),\quad
 \iota(h)=\operatorname{diag}(h,1),\quad
 K_j=\operatorname{GL}_j(O),\quad K_j^1=1+\varpi M_j(O).
 \tag{GP1}
\] The notation K\_j\^{}1 means the kernel of reduction inside K\_j. The actual ranks are a=42 or168. The actual G primes are completely split, so these are precisely the local general-linear groups, not an inner form substituted for them.

Let pi\_G be an existing actual unitary irreducible bottom constituent containing the inflation of the finite cuspidal representation Pi of GL\_n(k) selected at this place. The existence and the prescribed finite type are K050/051; temperedness is DO046. Replacing all finite labels by their contragredients when required by the source convention preserves cuspidality and every assertion below.

Let pi\_H be the corresponding actual unitary initial constituent and let y be any Hilbert-unit vector in \[
 V^{\psi^{-1}},\qquad
 V=\pi_H^{K_a^1},\qquad
 \pi_H(u)y=\psi(u)^{-1}y\quad(u\in U_a(k)).
 \tag{GP2}
\] Here psi is the exact nondegenerate character in R025/LIFT32: the bottom uses psi and the initial bilinear row its inverse. The initial level includes K\_a\^{}1, so (GP2) is a property of the actual local vectors, not an extra principal-level refinement.

For unit x in the selected finite Pi-subspace of pi\_G, put \[
 \alpha_v(x,y)=\int_H
    \langle\pi_G(\iota(h))x,x\rangle
    \langle\pi_H(h)y,y\rangle\,dh,
 \qquad \operatorname{vol}(K_a,dh)=1.
 \tag{GP3}
\] At these G places the source translate tau is the identity; its nonintegral lower component is at T and its nonzero Whittaker test is K053. All inner products below are invariant positive Hermitian forms.

\textbf{Theorem.} There is a unit vector x with the exact positive bottom character psi and with the same finite type Pi, such that \[
 \boxed{\alpha_v(x,y)\ge\frac1{D_a(q)}>q^{-a(a+1)/2}>0,\qquad
        D_a(q)=\prod_{i=1}^a(q^i-1).}
 \tag{GP4}
\] More precisely, write the orthogonal finite H(k)-decomposition as \[
 V=\bigoplus_\sigma\sigma\otimes M_\sigma,\qquad
 y=\sum_{\sigma\ \mathrm{generic}}e_\sigma\otimes m_\sigma,
 \quad\|e_\sigma\|=1,\quad
 e_\sigma\in\sigma^{\psi^{-1}},\quad
 \sum_\sigma\|m_\sigma\|^2=1.
 \tag{GP5}
\] For each nonzero m\_sigma one can choose a unit x\_sigma in the sigma-dual copy inside Pi, with positive character psi, for which \[
                 \alpha_v(x_\sigma,y)
                        =\frac{\|m_\sigma\|^2}{\dim\sigma}.
 \tag{GP6}
\] In particular the value is exactly 1/dim sigma for a unit initial vector in a single irreducible generic finite type. These are full noncompact-group integrals, not merely compact projections.

\subsection{2. Compact induction identifies the support exactly}\label{compact-induction-identifies-the-support-exactly}

Morris (1999), Proposition4.1, printed151, says that an irreducible smooth representation containing a cuspidal representation of the reductive quotient of a maximal parahoric is compactly induced from its normalizer. For GL\_n(F), that normalizer is ZK\_n, where Z is the scalar center. Thus \[
        \pi_G\simeq\operatorname{c-Ind}_{ZK_n}^{G}\rho.
 \tag{GP7}
\] The restriction of rho to K\_n is precisely the inflation of Pi, and rho(z k)=omega\_pi\_G(z)Pi(bar k). To justify the last statement, the center acts by the scalar central character of pi\_G. Any copy of Pi in rho is therefore stable under ZK\_n, and irreducibility of rho makes it the whole representation. Compatibility on Z intersect K\_n is supplied by the actual existing pi\_G, not freely imposed.

For clarity, the normalizer statement has an elementary lattice verification. If g normalizes K\_n, the lattice gO\^{}n is K\_n-stable. Scale it so it is contained in O\^{}n but not in varpi O\^{}n.~Its nonzero reduction is GL\_n(k)-stable, hence all of k\^{}n.~Nakayama's lemma then gives the full lattice O\^{}n.~Consequently g lies in ZK\_n; the reverse inclusion is immediate.

Realize (GP7) as functions f:G to Pi satisfying f(jg)=rho(j)f(g), compactly supported modulo ZK\_n, with right translations. The Hilbert norm is the sum of squared values over the discrete left cosets ZK\_n\G. The unitarity of rho follows from the unitary central character and finite Pi. This fixes an isometric finite-type copy \[
 x\longmapsto f_x,\qquad
 f_x(j)=\rho(j)x\ (j\in ZK_n),\qquad f_x(g)=0\ (g\notin ZK_n).
 \tag{GP8}
\] All these vectors are K\_n\^{}1-fixed. Their matrix coefficients satisfy \[
 \langle\pi_G(g)f_x,f_{x'}\rangle=
 \begin{cases}
   \langle\rho(g)x,x'\rangle,&g\in ZK_n,\\
   0,&g\notin ZK_n.
 \end{cases}
 \tag{GP9}
\] Indeed the two supports ZK\_n g\^{}(-1) and ZK\_n intersect exactly when g belongs to ZK\_n, and in that case the only occupied left coset in the second vector is ZK\_n itself. No formal-degree scalar enters this Hilbert normalization. Transport to the actual unitary realization by an isometry preserves (GP9).

Only this canonical copy is needed. We do not assert, without a separate intertwining calculation, that every occurrence of the finite type in the entire restricted representation is this copy. The actual bottom finite vector is allowed to be chosen inside its existing automorphic constituent. The copy (GP8) is at the same K\_n\^{}1 level and has the same genuine central projector.

The decisive support intersection is \[
             \iota(H)\cap ZK_n=\iota(K_a).
 \tag{GP10}
\] If diag(h,1)=z k with k and k\^{}(-1) integral, its last diagonal entry and that of its inverse show v(z)\textless=0 and v(z)\textgreater=0 respectively. Thus z is a unit, h and h\^{}(-1) are integral, and h is in K\_a. The converse is immediate. In particular (GP3) is compactly supported and absolutely convergent directly.

\subsection{3. The full period is a finite invariant projection}\label{the-full-period-is-a-finite-invariant-projection}

Substituting (GP9) and (GP10) in (GP3) gives \[
 \alpha_v(x,y)=\frac1{|H(k)|}\sum_{h\in H(k)}
       \langle\Pi(\iota(h))x,x\rangle
       \langle\pi_H(h)y,y\rangle
       =\|P_{H(k)}(x\otimes y)\|^2,
 \tag{GP11}
\] where P\_H(k) is probability averaging in Pi restricted to H(k) tensor V. Reduction is legitimate because both vectors are fixed by the corresponding first principal subgroup. The last equality uses that probability averaging is the orthogonal projection onto finite-group invariants.

LIFT32(FB1--FB5) supplies an isometric copy of the finite Gelfand--Graev representation Gamma\_a=Ind\_(U\_a)\^{}(H(k)) psi inside Pi restricted to H(k). The finite cuspidal Pi is generic, as already verified in DO042/043 and K050. Every generic sigma-dual therefore occurs in Pi restricted to H(k). LIFT32, citing Gow--Vinroot Theorem4.3, also supplies \[
 \dim\sigma^{\psi^{-1}}=1\quad(\sigma\text{ generic}),\qquad
 \sum_{\sigma\text{ generic}}\dim\sigma
       =\dim\Gamma_a=|H(k)|/|U_a(k)|=D_a(q).
 \tag{GP12}
\] These settled finite branching results are used with their exact positive/inverse character conventions, not reproved here.

Choose an isometric inclusion j\_sigma:sigma-dual to Pi restricted to H(k), and set x\_sigma=j\_sigma(bar e\_sigma), using the Hermitian identification of the dual with the conjugate representation. Its character is psi, as required. On sigma-dual tensor sigma the unit invariant is \[
              \Omega_\sigma=(\dim\sigma)^{-1/2}
                      \sum_{i=1}^{\dim\sigma}\bar e_i\otimes e_i.
 \tag{GP13}
\] The projection of bar e\_sigma tensor e\_sigma onto this line has squared norm 1/dim sigma. The multiplicity vector m\_sigma is unchanged; all other irreducible summands of y are orthogonal and have zero invariant projection with x\_sigma. Equations (GP11) and (GP13) prove the exact formula (GP6).

Set w\_sigma=\textbar\textbar m\_sigma\textbar\textbar\^{}2. Since sum w\_sigma=1, there is a sigma with w\_sigma/dim sigma at least 1 divided by the sum of the dimensions of the occurring generic sigmas. That sum is at most D\_a(q) by (GP12). This proves the first inequality in (GP4). Finally each q\^{}i-1 is strictly less than q\^{}i, so D\_a(q) is strictly less than q\^{}(a(a+1)/2). This proves the strict second inequality.

The argument allows arbitrary multiplicities M\_sigma in the actual initial finite module; they were not assumed one-dimensional. Nor is a new initial finite vector substituted for the prescribed y. Only the bottom vector x is chosen, inside the unchanged finite Pi and with the prescribed U\_a character.

\subsection{4. Which actual places this closes}\label{which-actual-places-this-closes}

The DO042 construction, preserved in the central-character repair DO043, uses one mixed-factor place in each rational G half and full-degree irreducible polynomials at all the other places. The mixed degrees are (37,6) for n=43 and (163,6) for n=169. The full degree finite representations are cuspidal. Each half has f places, where f=42 or168. Consequently the theorem applies at exactly \[
           2(f-1)=82\text{ or }334\text{ places},
 \tag{GP14}
\] with two mixed-factor exceptions in either host.

The exceptions cannot simply be replaced by full-degree factors: the literal goodTuple/IsGood clause6 requires a prime factor-degree witness p satisfying n\textless2p and p+3\textless=n.~DO042 supplies that witness at the retained mixed place. Altering it would require a new admissible construction; no such alteration is made here.

For the unchanged full-degree places, independent choices in the local tensor factors give the exact product lower bound \[
 \prod_{v\in G_{\mathrm{cusp}}}\alpha_v(x_v,y_v)
 \ge [D_a(g_0)D_a(g_1)]^{-(f-1)}
 >2^{-C_G},\quad
 C_G=\frac{a(a+1)}2(f-1)
       (\operatorname{bits}(g_0)+\operatorname{bits}(g_1)).
 \tag{GP15}
\] The actual data are

\begin{longtable}[]{@{}lrrrr@{}}
\toprule\noalign{}
Host & a=f & g\_0 & g\_1 & C\_G \\
\midrule\noalign{}
\endhead
\bottomrule\noalign{}
\endlastfoot
J\^{}14 & 42 & 4,615,409 & 5,276,377 & 1,703,058 \\
J\^{}56 & 168 & 22,285,932,761 & 64,072,175,441 & 168,321,972 \\
\end{longtable}

This bound is on a product of squared-period integrals with maximal-compact volume1. It is not a bound on a linear detector, an integral denominator, or the DP44 net gain.

\subsection{5. Canonical GGP factor and arithmetic scope}\label{canonical-ggp-factor-and-arithmetic-scope}

Retain DO046(AP17--AP18) exactly. In the present maximal-volume1 measure its normalized local factor is \[
 \mathcal P_v^\natural(x\otimes y,x\otimes y)
       =\mathcal L_v(1/2,\pi_H\boxtimes\pi_G)^{-1}
                              \alpha_v(x,y).
 \tag{GP16}
\] At a split unitary place the factor is explicitly \[
 \mathcal L_v(1/2)
    =\left(\prod_{j=1}^{a+1}\zeta_q(j)\right)
       \frac{|L(1/2,\pi_G\times\pi_H)|^2}
       {L(1,\pi_G\times\pi_G^\vee)
        L(1,\pi_H\times\pi_H^\vee)}.
 \tag{GP17}
\] The adjoint convention here includes the central direction of the unitary adjoint representation, as in DO046. All these factors are finite and nonzero for the actual tempered data. Thus the exact canonical factor is (GP6) divided by (GP17), and in particular is nonzero. This proof leaves the explicit factor (GP17) in the bound; it does not set it equal to1 or remove an adjoint L-factor. If a fixed measure is c\_v times maximal-volume1, (GP3), (GP6) and (GP16) are each multiplied by c\_v.

At a pure initial finite type the exact factor is 1/(dim sigma times mathcal L\_v). For arbitrary initial y it is w\_sigma/(dim sigma times mathcal L\_v). The positivity of the raw finite projection is explicit in (GP11); only nonzero finite normalization is needed for the canonical nonvanishing conclusion.

Global occurrence of pi\_G allows its local vectors to be changed within the same restricted tensor product. The new vector has the same level, genuine projector, finite central character and archimedean type. Such a local choice does not by itself compute a global period. The global central L-value, the two exceptional mixed-factor local tests, the fixed full-Schwartz arithmetic pins, a nonzero permitted terminal row, and its integral preimage and cost remain separate. No unrestricted interpolation or integral central projector is used in selecting x\_sigma.

In particular (GP4) cannot be interpreted as an11-integral right inverse. The isometric finite embeddings and projections are complex Hilbert-space maps. Their arithmetic denominators must be assessed by the separate R028/R029 finite lattice results and the pending LIFT40 analysis and the global theta-image problem.

\subsection{6. Sources, geometric dictionary, and verification scope}\label{sources-geometric-dictionary-and-verification-scope}

\begin{itemize}
\tightlist
\item
  Morris, L. (1999). Level zero G-types. \emph{Compositio Mathematica, 118},135--157. \url{https://doi.org/10.1023/A:1001019027614}. Proposition4.1, printed151, gives precisely the compact-induction theorem used in (GP7). Original publisher PDF: \url{https://www.cambridge.org/core/services/aop-cambridge-core/content/view/0881CA03D211FCCD82DE58C054D7AB46/S0010437X99000512a.pdf/level-zero-g-types.pdf}. The complete statement and its hypotheses were read directly; the theorem is cited rather than independently reproved.
\item
  ALP-LIFT-32(FB1--FB5,FB12), finite generic branching and its exact positive/inverse generic character. Its primary sources are Gan--Gross--Prasad(2012), Proposition4.2 and Corollary4.3, and Gow--Vinroot(2007), Theorem4.3; exact locators are retained there.
\item
  AM-DO-042/043, full-degree and mixed goodTuple factors; AM-K050/051, genuine existing bottom constituents at the unchanged levels; AM-R-025/032, actual residue/character and finite-projector pins; AM-DO-046(AP10,AP17--18), actual temperedness and GGP normalization. These are used in their checked scopes and not reproved.
\end{itemize}

The finite representations here are representations of the smooth GL group schemes over O and their reductive special fibres. Their character idempotents act on the same coefficient Schur functors located in SGA3I4.7.3/4.7.3.1 and SGA6VI3.3 by the existing dictionary. Compact induction and the analytic period belong to a local representation realization, not to a newly constructed integral motive or a change of the CM Picard/biextension object. The existing FGA232/236, SGA7I VII--VIII and cote4/cote54 locators continue to describe that underlying object; no new Cartier--Dieudonne slope, ell-adic weight or Hecke RH consequence follows from (GP11).

The companion certificate checks the support-intersection exponent criterion, exact finite Schur projection and multiplicity examples, the actual place counts and the rational product bounds. It is a finite diagnostic, not a computation of an actual p-adic integral, global L-value, terminal row or Smith matrix. The author read the complete proof GP1--17, reconstructed its support, finite Schur and arbitrary-multiplicity arguments, and checked the actual place counts, generic signs and retained GGP factor. The exact certificate passed28,461 assertions at06:14:06UTC. A syntax-only script failure before that successful run is retained in the tool history; no mathematical assertion failed. Independent review remains pending before integration.

\clearpage

\section{AM-K054 --- exact canonical T-period factor on the actual unit vectors}\label{am-k054-exact-canonical-t-period-factor-on-the-actual-unit-vectors}

Author: GPT-6 Astra Pro, ALP\_MACHINE \nolinkurl{/root/kummer_extension}. 22 November 2026 UTC. Inherited reasoning effort is not separately exposed.

\textbf{Author whole-proof PASS --- 2026-11-22T06:13:33.362104+00:00.}

This is a normalization companion to K053. It converts its nonzero Whittaker test to the actual unit-vector local matrix-coefficient period. The measure on \(\mathrm{GL}_a(F)\) gives its maximal compact subgroup volume one, exactly as at the retained T places.

\subsection{1. Canonical quantity and conclusion}\label{canonical-quantity-and-conclusion}

Use K053(NT1--NT3): \(\Pi\) is tempered spherical on \(G=\mathrm{GL}_{a+1}(F)\), \(\pi\) is tempered generic on \(H=\mathrm{GL}_a(F)\), \(c(\pi)=c(\omega_\pi)=1\), and the actual translation is \(\tau=I+\varpi^{-1}E_{a+1,a}\). Let \(v_\Pi,v_\pi\) be Hilbert-unit spherical and newform vectors, respectively. Put

\[
 A_T=\int_H
  \langle\Pi(\operatorname{diag}(h,1))\Pi(\tau)v_\Pi,
                              \Pi(\tau)v_\Pi\rangle
  \langle\pi(h)v_\pi,v_\pi\rangle\,dh_{\mathrm{prob}}.
 \tag{CP1}
\]

The invariant Hermitian forms can be any unitary realizations: normalizing the vectors to unit length makes (CP1) independent of their scalar choices. The tensor product in (CP1) has the same bilinear period convention as K053; if the global reading uses the contragredient initial representation, replace \(\pi\) throughout by that representation. All hypotheses and the final value remain the same.

Define

\[
\begin{split}
 L_G&=L(1,\Pi\times\Pi^\vee),\qquad
 L_H=L(1,\pi\times\pi^\vee),\\
 L_0&=L(1/2,\Pi\times\pi),\\
 \zeta_q(s)&=(1-q^{-s})^{-1},\qquad
 \Delta_{a+1}(q)=\prod_{j=1}^{a+1}\zeta_q(j),\\
 \mathcal L_T&=\Delta_{a+1}(q)\frac{|L_0|^2}{L_GL_H},\qquad
 A_T^\natural=A_T/\mathcal L_T .
\end{split}
 \tag{CP2}
\]

At a split unitary place, these are the literal local factors in the unitary GGP normalization: base change has the two general-linear factors, so its central Rankin--Selberg factor is \(|L_0|^2\); the two adjoint factors are \(L_G,L_H\); and the split quadratic character makes \(\Delta\) the displayed zeta product. This includes the central direction of the unitary adjoint representation. It is not the adjoint factor with the general-linear scalar line removed.

\textbf{Theorem.} On the unchanged vectors and embedding,

\[
 \boxed{A_T^\natural=\frac{q}{(q-1)(q^a-1)}>0,}
 \qquad
 A_T=\Delta_{a+1}(q)\frac{|L_0|^2}{L_GL_H}
                       \frac{q}{(q-1)(q^a-1)}.
 \tag{CP3}
\]

In particular, (CP3) closes the Hilbert/GGP normalization comparison explicitly left open in K053 section6. It remains a local theorem. It implies no global central L-value lower bound or terminal arithmetic-image nonvanishing.

\subsection{2. All Haar changes}\label{all-haar-changes}

Let \(dx\) be additive Haar measure with \(O\) volume one. For \(\mathrm{GL}_j(F)\), write

\[
 dg_{\mathrm{raw}}=
       \frac{\prod_{r,s}dx_{rs}}{|\det g|^j},\qquad
 \kappa_j=\prod_{i=1}^{j}(1-q^{-i}),\qquad\kappa_0=1.
 \tag{CP4}
\]

Counting invertible residue matrices gives

\[
 \operatorname{vol}_{\mathrm{raw}}\mathrm{GL}_j(O)
 =q^{-j^2}|\mathrm{GL}_j(k)|=\kappa_j,\qquad
 dg_{\mathrm{raw}}=\kappa_j\,dg_{\mathrm{prob}}.
 \tag{CP5}
\]

Both measures on \(N_j\) are the additive product measure with \(N_j(O)\) volume one. Therefore the corresponding quotient measures also differ by \(\kappa_j\).

The precise primary normalization is Zhang (2014), section2.1, printed550--551, together with the \textbf{change of measures} paragraph immediately before Lemma4.7, printed569. His later local section4.5 uses \(dg_{\mathrm{raw}}\), with the earlier convergence factor \(\zeta_q(1)\) removed. Ignoring that paragraph would change (CP3).

For any generic unitary representation \(\rho\) of \(\mathrm{GL}_j(F)\), Zhang's invariant Whittaker inner product is

\[
 \vartheta_{\rho,\mathrm{raw}}(W,W)=
 \int_{N_{j-1}\backslash\mathrm{GL}_{j-1}(F)}
   |W(\operatorname{diag}(h,1))|^2\,dh_{\mathrm{raw}} .
 \tag{CP6}
\]

Absolute convergence and invariance are the standard unitary Whittaker pairing, stated in Zhang(3.2), printed553--554. Consequently

\[
 \vartheta_{\rho,\mathrm{raw}}
      =\kappa_{j-1}\vartheta_{\rho,\mathrm{prob}}.
 \tag{CP7}
\]

The Rankin--Selberg functional uses \(N_a\backslash H\), so its raw value on the normalized functions in K053 is

\[
 \lambda_{\mathrm{raw}}=\kappa_a Z(\Pi,\pi).
 \tag{CP8}
\]

This is a measure change, not a redefinition of the actual vectors.

\subsection{3. Exact two Whittaker norms}\label{exact-two-whittaker-norms}

First let \(W_\Pi^0(1)=1\) be spherical. The diagonal Casselman--Shalika/Miyauchi formula and probability Iwasawa measures give

\[
 \vartheta_{\Pi,\mathrm{prob}}(W_\Pi^0,W_\Pi^0)
  =\sum_{\lambda_1\ge\cdots\ge\lambda_a\ge0}
     q^{-|\lambda|}\,
       s_{(\lambda,0)}(z)\,
       s_{(\lambda,0)}(\bar z)
  =\frac{L_G}{\zeta_q(a+1)} .
 \tag{CP9}
\]

Here \(z\) is the unitary Satake tuple of length \(a+1\). For completeness, the last equality has an exact normalization check. The complete Cauchy sum over all partitions of length at most \(a+1\) is \(L_G\). Every such partition is uniquely \(\mu+m(1^{a+1})\), where \(\mu_{a+1}=0\) and \(m\ge0\). The products of the determinant characters cancel, since \(|\prod z_i|=1\). The extra sum is \(\sum_{m\ge0}q^{-(a+1)m}=\zeta_q(a+1)\). This proves (CP9), without inserting an arbitrary Whittaker scalar. It also agrees with Zhang(3.3): in raw measures the first norm is \(\kappa_{a+1}L_G\), since \(\kappa_a/\zeta_q(a+1)=\kappa_{a+1}\).

Next use K053(NT7--NT8): the standard L-factor of the actual initial representation has the \(a-1\) unramified unitary parameters \(\beta_1,\ldots,\beta_{a-1}\) and one zero parameter. For its normalized newform, Miyauchi's diagonal formula gives

\[
 \vartheta_{\pi,\mathrm{prob}}(W_\pi^0,W_\pi^0)
   =\sum_{\lambda_1\ge\cdots\ge\lambda_{a-1}\ge0}
        q^{-|\lambda|}s_\lambda(\beta)s_\lambda(\bar\beta)
   =\prod_{i,j=1}^{a-1}(1-\beta_i\bar\beta_jq^{-1})^{-1}
   =\frac{L_H}{\zeta_q(1)}.
 \tag{CP10}
\]

The last equality uses the exact conductor-one classification. The tensor product of the tame character with its dual contributes one unramified trivial character, hence \(\zeta_q(1)\). Its cross terms with the unramified characters remain ramified and have local L-factor one. The unramified/unramified terms give the displayed product.

The Haar factors in these Iwasawa sums are fixed as follows. The compact and diagonal-unit integrals each have volume one. The quotient of the two Borel moduli in the squared diagonal value contributes \(|\det h|\), giving exactly \(q^{-|\lambda|}\). For the ramified newform its compact invariance is the embedded \(\mathrm{GL}_{a-1}(O)\), which is contained in the actual \(K_r\); full \(\mathrm{GL}_a(O)\)-invariance is not assumed.

Combining (CP7), (CP9), and (CP10) yields

\[
 \vartheta_{\Pi,\mathrm{raw}}(W_\Pi^0,W_\Pi^0)
        =\kappa_a\,\frac{L_G}{\zeta_q(a+1)},\qquad
 \vartheta_{\pi,\mathrm{raw}}(W_\pi^0,W_\pi^0)
        =\kappa_{a-1}\,\frac{L_H}{\zeta_q(1)}.
 \tag{CP11}
\]

Translation by \(\tau\) preserves the first norm, by invariance of (CP6).

\subsection{4. The published matrix-coefficient comparison and cancellation}\label{the-published-matrix-coefficient-comparison-and-cancellation}

Zhang (2014), Proposition4.10, printed571, with its completed proof on printed574, proves for tempered representations and the raw Whittaker norms (CP6) that

\[
 \int_H
     \langle(\Pi\otimes\pi)(h)W,W\rangle_{\vartheta}\,
                  dh_{\mathrm{raw}}
        =|\lambda_{\mathrm{raw}}(W)|^2.
 \tag{CP12}
\]

The statement is exactly for the standard block embedding of the general-linear groups at a split place. We use it as a theorem; we do not replace it by a formal multiplicity-one argument with an unspecified scalar. Its \(\lambda\) is the Rankin--Selberg functional of section3.3, with exponent zero at the central point. The opposite Whittaker characters are the ones in K053.

Apply (CP12) to \(W=\rho(\tau)W_\Pi^0\otimes W_\pi^0\). Divide by the two norms in (CP11) to pass to Hilbert-unit vectors, and divide the outer Haar measure by \(\kappa_a\) to obtain (CP1). By (CP8),

\[
\begin{split}
 A_T
 &=\frac{1}{\kappa_a}
       \frac{\kappa_a^2|Z|^2}
        {[\kappa_a L_G/\zeta_q(a+1)]
         [\kappa_{a-1}L_H/\zeta_q(1)]}\\
 &=\frac{\zeta_q(a+1)\zeta_q(1)}{\kappa_{a-1}}\,
                   \frac{|Z|^2}{L_GL_H}.
\end{split}
 \tag{CP13}
\]

On the other hand, \(\Delta_{a+1}(q)=\zeta_q(a+1)\zeta_q(a)/\kappa_{a-1}\). Thus

\[
 A_T^\natural
   =\frac{\zeta_q(1)}{\zeta_q(a)}
                       \left|\frac{Z}{L_0}\right|^2
   =\frac{1-q^{-a}}{1-q^{-1}}\,
                       \frac{q^a}{(q^a-1)^2}
   =\frac{q}{(q-1)(q^a-1)}.
 \tag{CP14}
\]

The middle equality is precisely K053(NT6). This proves (CP3). All representation-dependent adjoint and central L-factors have been retained until the definition of \(A_T^\natural\), where the displayed cancellation is exact.

There is also an independent normalization check on the measure convention. With both representations spherical and \(W(1)=1\), the two raw mirabolic norms are \(\kappa_{a+1}L_G\) and \(\kappa_aL_H\), and the raw period is \(\kappa_aL_0\). The same formula (CP12), divided by the outer \(\kappa_a\), gives \(A_T=\kappa_{a+1}^{-1}|L_0|^2/(L_GL_H)=\mathcal L_T\). Thus the canonical maximal-compact unit-vector factor is one in the unramified case, as required by Zhang(1.3)--(1.4).

In particular this is stronger than the compact Gram statement K052: integration over the full noncompact \(H\) is now evaluated on the fixed tensor. It does not rely on a claim that all finite compact-type occurrences have the same invariant functional.

\subsection{5. Both actual hosts and the eleven-adic value}\label{both-actual-hosts-and-the-eleven-adic-value}

Across the \(f\) actual T places, with their retained common residue cardinality \(T\), (CP3) gives the rational number

\[
 \boxed{\prod_{v\mid T} A_{T,v}^\natural
   =\left(\frac{T}{(T-1)(T^a-1)}\right)^f
   =T^{-af}(1-T^{-1})^{-f}(1-T^{-a})^{-f}.}
 \tag{CP15}
\]

It is independent of the particular actual tempered constituents and of the allowed role/word/canonical-twist profile. The only assumptions used were those already proved at T for the actual vectors.

The exact prime certificate of DO040 gives:

\begin{longtable}[]{@{}
  >{\raggedright\arraybackslash}p{(\columnwidth - 10\tabcolsep) * \real{0.1304}}
  >{\raggedleft\arraybackslash}p{(\columnwidth - 10\tabcolsep) * \real{0.1739}}
  >{\raggedleft\arraybackslash}p{(\columnwidth - 10\tabcolsep) * \real{0.1739}}
  >{\raggedleft\arraybackslash}p{(\columnwidth - 10\tabcolsep) * \real{0.1739}}
  >{\raggedleft\arraybackslash}p{(\columnwidth - 10\tabcolsep) * \real{0.1739}}
  >{\raggedleft\arraybackslash}p{(\columnwidth - 10\tabcolsep) * \real{0.1739}}@{}}
\toprule\noalign{}
\begin{minipage}[b]{\linewidth}\raggedright
Host
\end{minipage} & \begin{minipage}[b]{\linewidth}\raggedleft
\(a=f\)
\end{minipage} & \begin{minipage}[b]{\linewidth}\raggedleft
\(T\bmod11\)
\end{minipage} & \begin{minipage}[b]{\linewidth}\raggedleft
\(v_{11}(T-1)\)
\end{minipage} & \begin{minipage}[b]{\linewidth}\raggedleft
\(v_{11}(T^a-1)\)
\end{minipage} & \begin{minipage}[b]{\linewidth}\raggedleft
\(v_{11}(\prod A_{T,v}^\natural)\)
\end{minipage} \\
\midrule\noalign{}
\endhead
\bottomrule\noalign{}
\endlastfoot
\(J^{14}\) & 42 & 9 & 0 & 0 & 0 \\
\(J^{56}\) & 168 & 10 & 0 & 2 & \(-336\) \\
\end{longtable}

These valuations are values of the \textbf{rational local normalization factor}, not valuations of the global period or of an integral arithmetic row. No statement about the saturation of that row follows from this table.

Since \(T>2^{964}\) and \(f\le168\),

\[
 0<
 \log\!\left(T^{af}\prod_{v\mid T}A_{T,v}^\natural\right)
   =-f\log(1-T^{-1})-f\log(1-T^{-a})
   <2^{-954}.
 \tag{CP16}
\]

For example, use \(-\log(1-x)\le2x\) and bound the right side by \(2f(2^{-964}+2^{-964a})<2^{-954}\). Thus the exact squared-period normalization has the expected principal factor \(T^{-af}\), and every remaining local correction is explicitly fixed. Taking a square root when estimating a linear detector cost must retain that this is a squared-period factor.

\subsection{6. Primary sources and checks}\label{primary-sources-and-checks}

\begin{itemize}
\tightlist
\item
  Zhang, W. (2014). Automorphic period and the central value of Rankin--Selberg L-function. \emph{Journal of the American Mathematical Society, 27}(2),541--612. \url{https://doi.org/10.1090/S0894-0347-2014-00784-0}. Original publisher offprint: \url{https://archive.ymsc.tsinghua.edu.cn/pacm_download/271/8715-jams784_off_print.pdf}. Section2.1 printed550--551; (3.2) printed553; the local raw-measure change preceding Lemma4.7 printed569; Proposition4.10 printed571 and its proof through574; local split base change on569--570; definition of the unitary \(\mathcal L\) on542--543 and the adjoint identification with the appropriate Asai factor in Remark7, printed560. At a split place both Asai signs restrict to the tensor-product parameter, explaining the full factors in (CP2).
\item
  Miyauchi (2014), Theorem4.1, and BKL (2020), the exact sources already recorded and checked in K053, supply the newform diagonal values and probability quotient measures.
\item
  The exact lower-translate test is K053(NT6), with the full phase identity in NT27. No new local vector is selected here.
\end{itemize}

The complete primary passages above were directly read from the downloaded original offprint. The attached exact certificate checks the rational measure cancellation, finite determinant-shift Cauchy normalizations, actual prime valuations and positive rational logarithmic majorants. It does not compute a global period, automorphic representation, or central global L-value.

\subsection{Author check record}\label{author-check-record}

The complete saved CP1--16 proof and complete certificate were read and reconstructed. The raw/probability Haar conversion, both actual Whittaker norms, split-unitary adjoint convention, published matrix-coefficient identity, exact scalar cancellation and both-host valuation bounds passed. The primary original offprint was directly read through the full proof of Proposition4.10. The finite certificate passed 942 exact assertions at 2026-11-22T06:12:37.470718+00:00; it enumerated 20,922 finite matrices and checked 128 finite partitions. This is an author proof check plus exact finite/formal diagnostics, not a Lean build or a computation of an automorphic representation.

\clearpage

\section{AM-DO-050 --- the selected final-two-step real recoupling}\label{am-do-050-the-selected-final-two-step-real-recoupling}

Author: GPT-6 Astra Pro, ALP\_MACHINE /root/definite\_oscillator. Inherited reasoning effort is not separately exposed. 22 November 2026 UTC.

\textbf{Author whole-proof check completed; exact verification record below.}

This report identifies the selected real component of the literal lastBi tensor. It uses the complete compact-middle coevaluation, then states separately the result for a pair of matched extremal vectors. The recoupling coefficient is nonzero and its absolute value is an exact ratio of two finite compact dimensions. The primitive-polynomial version retains all four hook norms, the selected bottom PBW norm, and the actual terminal branching norm. No dimension of an infinite discrete series occurs.

This is a local real identity. It does not replace the fixed finite G/Q/T tensor laws, identify an arbitrary highest tensor with lastBi, or assert a nonzero global terminal period. Section 8 states exactly how it enters K038 and DO038.

\subsection{1. Spaces, actions, and settled inputs}\label{spaces-actions-and-settled-inputs}

At one real label write

\[
 A=U(a),\quad V=U(a,1),\quad B=U(b),\quad
 H=U(N),\quad C=U(c),
 \qquad b=a+1+t,\quad N=b+t,\quad c=N+1.
 \tag{RC1}
\]

The horizontal edges have ranks a,b,c.~The vertical edges have ranks r=a+1,b,N. All groups other than V are compact. Let L be the negative line in V and Z the positive line in C, so that V=A perpendicular L and C=H perpendicular Z. The notation also denotes the corresponding Hermitian spaces when an oscillator space is written.

The actual downward Hom convention of R015 is retained. In its notation a selected oscillator summand is bar(sigma\_source) external-tensor sigma\_target. Write

\begin{itemize}
\tightlist
\item
  sigma0, sigma1, sigma2 for the three actual horizontal compact representations of DO025/K037;
\item
  tau0, tau1, tau2 for the actual vertical discrete-series/compact representations of R015 and LIFT39.
\end{itemize}

Thus tau0 is infinite-dimensional, while tau1 and tau2 are finite compact representations. The source arithmetic coefficient action rhoVr is the conjugate/dual of the corresponding sigma action; K037(CR14--18) fixes this convention through pairBK. Every dual below means that actual contragredient, with the conjugate Hilbert metric.

The following results are used as checked inputs:

\begin{enumerate}
\def\labelenumi{\arabic{enumi}.}
\tightlist
\item
  R015 and LIFT39: the actual unitary Howe embeddings on the vertical edges, with the stated odd-character and Hom transport. The first contains bar(tau0) tensor tau1 and the second contains bar(tau1) tensor tau2.
\item
  DO025 and K037(CR17--19): the horizontal Cauchy summands are bar(sigma0) tensor sigma1 and bar(sigma1) tensor sigma2, each with multiplicity one, with the specified common characters.
\item
  K043(PB5--8): there is an isometric selected compact inclusion J: sigma0 to tau0. Its construction is the actual Pieri/PBW map divided by the square root of D\_PBW. Its multiplicity is one, including the compact central weight.
\item
  R019/R020: the actual codimension-one terminal branching has a multiplicity-one inclusion i: tau2 to sigma2, normalized by i(p\_mu)=xiBr, and i* i=kappa I. Hence J\_mu=i/sqrt(kappa) is an isometric inclusion. The literal branch is B\_T=u\_T cBr i*, with u\_T=bar(pi\_T)\^{}(-nBr).
\item
  LIFT39(PN3--5): the literal theta vacuum is unit, and the product-CM-to-unit factors are known explicitly.
\end{enumerate}

The source definitions of lastBi and kap are Development/AbelianVarietyTheory/Definitions/LatticeFrames/ \nolinkurl{CombAtSingleZero.lean29–35.} The exact skew adjoint is Development/AbelianVarietyTheory/BlockTheory/TwistedDiscreteSeries/ \nolinkurl{SkewGramHalfModule.lean21–28.} The latter file has now been recovered under the own-prefix seesaw-source directory, with its intake hash matched. The finite tuple does not occur in the real spaces in RC1.

\subsection{2. The literal line shuffle is a unitary real oscillator map}\label{the-literal-line-shuffle-is-a-unitary-real-oscillator-map}

For an orthogonal decomposition o of V and o' of C, the source says

\[
 \begin{split}
 \operatorname{kap}_\tau(w)
   &=-\xi_0^{-1}G_B^{-1}(\tau w)^*G_L,\\
 \operatorname{lastBi}_\tau(\Phi_0,\Phi_1)(x,z)
   &=\Phi_1(z_H)\,
        \Phi_0\bigl(x,\operatorname{kap}_\tau(z_Z)\bigr).
 \end{split}                                                    \tag{RC2}
\]

At rank zero the flatBF and the top extension in Comm.lst are the identity under the fixed orthogonal coordinates. This follows literally from the zero-dimensional Fourier/unipotent integrations; \nolinkurl{SeeSawLastTwoStepsValueFixedLevelPlaces.lean39–62} proves its split coordinate counterpart flCol\_one without a measure factor.

The required line identity is

\[
       \tau^*G_L\tau=-\xi_0\bar\xi_0G_Z.
 \tag{RC3}
\]

It is WeilOK.tau in PlaceLaws.lean165--167. The corresponding characters are those of R016: with the finite tuple character xi included, the vertical labels are (mu0 xi eta\^{}-1,mu0 xi) and (mu1 xi,mu1 xi eta\^{}-1), whereas the horizontal labels are (mu0 xi,mu0 xi) and (mu1 xi,mu1 xi). The negative-line cross character is eta\^{}-1. Therefore the orthogonal-sum and adjoint transports give exactly the same A times diagonal-B times H action on the two sides of RC2. No inverse character is omitted. At infinity xi is trivial. These are the explicit equations of PlaceLaws.WeilOK.last, lines156--164, and the established R016 conjugate-target/adjoint transport.

The two edge characters mu0 and mu1 need not agree. Their apparent discrepancy on the exchanged line does not change the asserted three-group action: the character attached to B controls the action of the line group, whose argument here is the identity. The B action on the exchanged summand is controlled by the line character eta\^{}-1 and its adjoint inverse. This is the exact character-change observation in R016 Section 2. We assert equivariance for A, the diagonal B, and H; we do not assert an unproved equality of actions of the two discarded line groups.

Here is also a direct metric and complex-structure check. In orthonormal B coordinates write G\_L=-ell and G\_Z=z, with ell,z\textgreater0. Equation RC3 is ell\textbar tau\textbar²=\textbar xi0\textbar²z. Thus

\[
 \operatorname{kap}_\tau(w)
        ={\ell\bar\tau\over\xi_0}w^*,\qquad
 {1\over\ell}\|\operatorname{kap}_\tau(w)\|^2
        =z\|w\|^2.                                           \tag{RC4}
\]

The negative-line Hom summand has opposite complex structure to the positive-line target summand; the adjoint in RC4 reverses i and intertwines these compatible complex structures. Its symplectic form is preserved as well: the negative Hermitian sign and complex conjugation each reverse the imaginary part, and their two reversals cancel. Equivalently substitute RC3 into the source symplO formula. Consequently RC2 rearranges an orthogonal direct sum of polarized real symplectic spaces by an isometry. Its oscillator realization

\[
 \mathcal S:\Omega(V,B)\widehat\otimes\Omega(B,H)
          \longrightarrow
          \Omega(A,B)\widehat\otimes\Omega(B,C)                 \tag{RC5}
\]

is unitary and complex linear. It takes polynomial Fock vectors to polynomial Fock vectors. The scalar phase of its realization is the one fixed by the literal theta-function pullback RC2 and the chosen vacua; it is not reset in the argument below. A rational change of line coordinates preserves the covolume-one rational Heisenberg measure, so the global theta isometry of LIFT39 gives the same unitarity. No Fourier covolume is to be inserted a second time.

For normalized lines G\_L=-1,G\_Z=1, tau=-xi0 gives kap(w)=w\emph{; tau=xi0 gives kap(w)=-w}. This sign is part of the map. In the source finite-line condition valuation(lambda)=1 means a unit for its multiplicative valuation, not additive valuation one.

\subsection{3. The horizontal Hom space is exactly one-dimensional}\label{the-horizontal-hom-space-is-exactly-one-dimensional}

Let j0 and j1 be the unitary vertical Howe embeddings. Let k0,k1 be the unitary horizontal Cauchy embeddings. At the polynomial level their restrictions have the respective source-target types listed in Section 1.

Consider the space of A times H intertwiners from bar(sigma0) tensor tau2 to the diagonal-B-invariant part of Omega(A,B) tensor Omega(B,C). We claim it has dimension exactly one.

\textbf{Proof.} The full polynomial Cauchy decomposition on the first horizontal edge is a direct sum over partitions nu of bar(sigma\_nu,A) tensor sigma\_nu,B, with the fixed vacuum determinant characters. Fixing the full U(a) type bar(sigma0), including its central weight, fixes every part of nu, not only its differences. It therefore fixes nu0 and its unique partner sigma1. There is no additional determinant-twist family with the same U(a) type.

On the second horizontal edge, taking diagonal U(b) invariants with that sigma1 forces its dual in the second oscillator. Cauchy multiplicity one then fixes nu1 and sigma2 uniquely. The remaining space is Hom\_H(tau2,sigma2). It has dimension one by codimension-one unitary branching and the actual interlacing in R019. The inclusion J\_mu is nonzero, so the Hom space is not zero. The argument applies in the Hilbert completion too: the compact A-central character already fixes a finite polynomial degree on the first edge, and the selected B type fixes a finite degree on the second. There is no unaccounted infinite direct sum. QED.

This proof is stronger than fixing only initial and terminal weights: it uses the entire first Cauchy correspondence and the diagonal middle invariants. It also explains why omitting the middle compact projection would leave extra components.

\subsection{4. Exact complete-coevaluation recoupling}\label{exact-complete-coevaluation-recoupling}

Put d\_v=dim(tau1) and d\_h=dim(sigma1), both finite positive integers. Choose orthonormal bases e\_i of tau1 and f\_j of sigma1. For v in bar(sigma0) and w in tau2 set

\[
 \begin{split}
 E(v,w)&=\sum_{i=1}^{d_v}
    j_0(\bar Jv\otimes e_i)\otimes j_1(\bar e_i\otimes w),\\
 F(v,w)&=\sum_{j=1}^{d_h}
    k_0(v\otimes f_j)\otimes k_1(\bar f_j\otimes J_\mu w).
 \end{split}                                                   \tag{RC6}
\]

Here bar J denotes the conjugate inclusion, and the displayed conjugates implement the canonical dual coevaluations. The formulas are independent of the orthonormal bases. The arithmetic version uses a basis and its pairBK-dual and is the same tensor.

\subsubsection{Theorem RC7}\label{theorem-rc7}

There is one nonzero scalar z, fixed by the chosen intertwiners and the literal line shuffle, such that

\[
 \boxed{\quad \mathcal S E(v,w)=zF(v,w),\qquad
                  |z|^2={d_v\over d_h}.\quad}                 \tag{RC7}
\]

The equality holds for every v,w. Its phase is not an arbitrary Howe scalar: in the specified models it is given exactly by

\[
 z={\langle\mathcal S E(v,w),F(v,w)\rangle\over d_h}
       \quad\text{for unit }v,w.                              \tag{RC8}
\]

The Hermitian convention is linear in the first entry. Formula RC8 is independent of the chosen unit vectors. It is a finite Fischer contraction of the explicitly specified polynomials and line map.

\textbf{Proof.} E is a nonzero A times H intertwiner into the diagonal-B invariants, since its middle tensor is the canonical coevaluation. The same holds for F. Equation RC5 preserves those actions and invariants. Section 3 therefore gives a scalar z with the first equality. Orthogonality and the four isometric embeddings give

\[
 \|E(v,w)\|^2=d_v\|v\|^2\|w\|^2,\qquad
 \|F(v,w)\|^2=d_h\|v\|^2\|w\|^2.                              \tag{RC9}
\]

Unitarity of S proves the second equality in RC7, and in particular z is nonzero. Pairing with F proves RC8. No infinite-dimensional trace has been taken: the only trace is over compact tau1. QED.

A change of the physical vacuum or line-intertwiner phase changes RC8 by that same phase. A change of a vertical IsIt normalization without changing the physical test instead changes its coefficient identification inversely, as proved in DO049. These two operations must not be confused.

\subsection{5. Matched simple tensors and the exact compact average}\label{matched-simple-tensors-and-the-exact-compact-average}

For middle vectors e,f in tau1, probability Haar on B gives

\[
 \int_B (\tau_1(k)e)\otimes\overline{\tau_1(k)f}\,dk
      ={\langle e,f\rangle\over d_v}
                       \sum_i e_i\otimes\bar e_i.             \tag{RC10}
\]

This follows from Schur orthogonality, or by identifying the tensor with a rank-one operator and averaging its conjugates. In particular its trace is exactly \textless e,f\textgreater, fixing the coefficient.

Consequently the middle-compact projection of the simple vertical pair is

\[
 \mathcal S\operatorname{Av}_B
 \{j_0(\bar Jv\otimes e)\otimes j_1(\bar f\otimes w)\}
       ={\langle e,f\rangle z\over d_v}F(v,w).                \tag{RC11}
\]

If e=f is a unit extremal, this is nonzero. If e and f are orthogonal, it is zero. Thus RC7 is not a claim about every simple tensor. LIFT39(PN20) keeps the same middle extremal while applying PBW on the bottom factor, and PN12/PN15 uses its matching middle dual at the second edge. Those prescribed matched extremals have the nonzero coefficient in RC11.

The global middle integral is unchanged by this probability compact average: substitute its variable h by hk and use Haar invariance before integrating k. The other two automorphic factors do not depend on k. In the compact real realization the integrals are absolutely convergent; this is precisely the compact Fubini operation already used in R016 and K038. No replacement of a finite Schwartz function is involved. The source op\_T is supported at finite T places, so it commutes with this real compact averaging. This observation does not identify its finite scalar or its prescribed finite kernel.

\subsection{6. Primitive hooks, selected PBW, and the actual terminal branch}\label{primitive-hooks-selected-pbw-and-the-actual-terminal-branch}

For a single label let

\[
 N_0=H_A B!,\qquad N_1=H_\Lambda,\qquad
 H_0=H_{\nu_0},\quad H_1=H_{\nu_1},\qquad D=D_{\rm PBW}.
 \tag{RC12}
\]

A and B, Lambda are exactly LIFT39(PN6--14); nu0,nu1 are the actual horizontal partitions of DO025, and D is K043(PB5). All compact extremals have norm one in the source normalized Fischer metrics. Thus the primitive embeddings are sqrt(N0)j0, sqrt(N1)j1 and sqrt(H0)k0,sqrt(H1)k1. Applying the unnormalized selected PBW map on the bottom multiplies the first embedding by sqrt(D), relative to the isometry J. This is PN20; using P0 alone would miss that selected type and this factor.

For the matched primitive vertical pair, its averaged horizontal coefficient tensor, expressed through the primitive horizontal Cauchy embeddings, is therefore

\[
 {z\over d_v}\sqrt{{N_0N_1D\over H_0H_1}}
       \sum_j (v\otimes f_j)\otimes
                            (\bar f_j\otimes J_\mu w).
 \tag{RC13}
\]

This is an equality after applying the stated primitive embeddings; it does not introduce a projection with an unspecified scalar. For the complete vertical coevaluation, replace z/d\_v by z. For an isometrically normalized selected PBW vector, omit D.

The literal terminal branch acts by

\[
 B_TJ_\mu=u_Tc_{\rm Br}\sqrt\kappa\,I,
             \qquad u_T=\bar\pi_T^{-nBr}.                     \tag{RC14}
\]

Indeed B\_T=u\_TcBr i\emph{, J\_mu=i/sqrt(kappa), and i}i=kappa I. This is the exact R019/R020 normalization, not its scalar highest-line value cBr*kappa substituted into an isometric line. Combining RC13 and RC14 gives the nonzero scalar

\[
 \boxed{\quad
 r_{\rm prim}={z\over d_v}
       \sqrt{{N_0N_1D\over H_0H_1}}
                       u_Tc_{\rm Br}\sqrt\kappa,
 \qquad
 |r_{\rm prim}|^2=
 {N_0N_1D\over d_vd_hH_0H_1}
                 |u_Tc_{\rm Br}|^2\kappa.\quad}              \tag{RC15}
\]

For complete vertical coevaluation the squared norm factor is instead d\_v/d\_h times N0N1D/(H0H1) times \textbar u\_TcBr\textbar² kappa. For all CM labels, tensor the maps and multiply every displayed factor, including z and u\_T. Each dimension is the finite Weyl product of its actual compact highest weight; all factorials and valuations can be computed by the checked R019/R011 formulas. This is an exact finite expression even at the enormous permitted weights. No factorial expansion or monomial enumeration is needed to compute its absolute value.

The phase in RC15 consists of the fixed Fischer contraction RC8, the actual sign of cBr, and the actual algebraic u\_T. This report has not evaluated RC8 to a universal sign for the source's already fixed CM/line frames. It supplies an exact polynomial formula and a nonzero absolute value; declaring z positive without transporting the chosen frame phase would be an unjustified extra pin.

\subsection{7. Vacuum and CM frame balance}\label{vacuum-and-cm-frame-balance}

Let a\_vac(n,m)=det(2Y\_J)\^{}(7hnm/2), as in LIFT39(PN5). Then

\[
 a b+b c=r b+b N=2b^2,
\] so \[
 a_{\rm vac}(a,b)a_{\rm vac}(b,c)
 =a_{\rm vac}(r,b)a_{\rm vac}(b,N)
 =\det(2Y_J)^{7hb^2}.                                        \tag{RC16}
\]

The positive raw-to-unit vacuum factor therefore cancels between the two complete chains. Individual edges still have their own factors. This cancellation does not remove their primitive It normalizations N0,N1 or the horizontal hook norms.

The raw CM common factors are also explicit. Horizontal is C\_theta\^{}(14h(ab+bc)). On the first vertical edge its a positive source columns contribute C\_theta\^{}(14hab), while its single negative column contributes bar(C\_theta)\^{}(14hb); the second vertical edge contributes C\_theta\^{}(14hbN). Their ratio is

\[
 {\text{horizontal raw CM factor}\over
          \text{vertical raw CM factor}}
             =\left({C_\theta\over\bar C_\theta}\right)^{14hb}.
 \tag{RC17}
\]

It has absolute value one. The creation-coordinate factors lambda\_tau/cD\_tau and their negative-line conjugates remain those of K035/LIFT39. RC17 is a CM reading ratio, not a replacement for the fixed skew-adjoint phase RC8. The literal lastBi formula RC2 and the chosen field embeddings determine how they combine.

\subsection{8. Exact use in the arithmetic class and terminal adapters}\label{exact-use-in-the-arithmetic-class-and-terminal-adapters}

K037(CR17--19) is the primitive horizontal Cauchy embedding used in RC13. Its complete pairBK coevaluation is canonical. K038 (CA33--36) therefore applies to the selected horizontal component of the full tensor S E, not just to one primitive highest product. Section 3 proves that the indicated compact projection is exactly that component; no arbitrary retrVV projection is needed.

DO038 realizes the literal terminal mass pairing, with the actual branch B\_T, actual decorated tilQ/tilC classes and their stabilizer weights. RC14 is precisely the coefficient-side branch in that adapter. It is not a statement about the restriction of one fixed right-highest scalar automorphic vector. In particular this result resolves the former right-highest/top-H-type ambiguity locally: the complete coefficient tensor contains the prescribed mu branch, and its norm is the positive expression RC15.

The global equation still needs the same finite tensor on both sides. In literal ValLastX it is

C=(op\_T).app(lastBi(Phi0,Phi1)),

with its full finite G/Q/T modifications and the actual finite trace/frame insertions. The source constructs this tensor before varying the finite tuple: \nolinkurl{SeeSawLastTwoStepsDoubleLiftPeriodValue.lean} 321--345, with tuple independence supplied by flat\_indep. It does not assert that C equals a product of two highest polynomials.

The remaining finite task is therefore concrete: prove those actual local G/Q/T identities in the new mixed family, with the fixed scalar cV=(sc0 sc1)cQ cG cT. No middle quotient-volume kappa occurs at m=2; see PlaceLaws.CVEq266--268 and ValLastX194--201. RC7/RC15 do not supply that finite equality or a global bottom restriction functional. They do show that the correctly selected real tensor and terminal branch are nonzero, with an exact local norm and an explicitly defined phase.

\subsection{9. Verification scope}\label{verification-scope-2}

The companion certificate checks the finite compact coevaluation and projection identities, exact hook/branch ratios in a concrete nontrivial Fock family, all vacuum-rank balances for the two hosts, and primitive versus isometric scaling. These checks are diagnostics for the full proof above, not a computation of an automorphic period or an assertion that every simple tensor has the stated coefficient.

The concrete Fock family uses a=1,b=2,H=U(2),C=U(3), with R=X1\emph{y1+X2}y2 and det Z of a two-by-two old-row block. Its shuffled polynomial is (X1\emph{z31+X2}z32)\^{}m det(Z)\^{}L. The vertical compact type is one-dimensional, the horizontal middle type has dimension m+1, and the terminal branch of (m+L,L,0) to (L,L) has kappa=(m+1)!(L+1)!/(m+L+1)!. Direct polynomial Fischer calculations verify RC7 with \textbar z\textbar²=1/(m+1), including nonzero non-top branching. This small family tests normalization; it is not substituted for the actual rank/character argument in Sections 1--6.

The certificate \nolinkurl{definite_oscillator_fixed_real_recoupling_certificate.py} completed at 2026-11-22T06:15:43.743831+00:00 with 6,336 exact assertions. It uses rational Gaussian-number Schur averaging, ordinary integer polynomial dictionaries and exact factorial/Fraction arithmetic; no floating-point equality is used. It covers 49 nontrivial and degenerate rank-two Fock examples, both actual host rank balances, and the primitive/complete/isometric scale distinctions.

The author then whole-read Sections 1--9, directly rechecked CombAtSingleZero.lastBi/kap, SkewGramHalfModule.ofAdj, PlaceLaws.WeilOK.last/tau/CVEq and ValLastX, and reread the settled K037(CR14--19), K043(PB2--8), LIFT39(PN15--20), R016 character transport and R020 adjoint branch formulas. The proof passes in the exact local selected-real and compact-average scope stated. The finite tensor laws, the fixed global restricted period and the unevaluated source-frame phase RC8 are expressly not promoted by this check.

\clearpage

\section{ALP-LIFT-40 --- integral generated vertical modules and the exact ambient projection obstruction}\label{alp-lift-40-integral-generated-vertical-modules-and-the-exact-ambient-projection-obstruction}

Author: GPT-6 Astra Pro, ALP\_MACHINE \nolinkurl{/root/lift_accounting}. 22 November 2026 UTC. Inherited reasoning effort is not separately exposed.

\textbf{Author whole-proof review completed; exact verification is recorded below.}

LIFT39 fixed the real scalar of the primitive vertical contraction. This report determines its arithmetic behavior on three different, independently defined polynomial lattices. On the module generated from the primitive joint vector it has an integral retraction. On the saturated selected oscillator summand it takes values in an explicit dual lattice, without the large joint hook norm. On the whole ordinary Fock lattice it need not be integral: explicit coefficient-one monomials have negative11-adic valuations in every actual endpoint profile checked below. The last phenomenon is an isotypic-projection obstruction and is not confused with the first.

The construction is local to the actual real/polynomial models. It does not put the indicated polynomial vectors in the fixed global theta image, compute a global Smith matrix, or suppress the CM period and conjugate-frame comparison of LIFT39.

\subsection{1. Fixed modules, pairings and lattices}\label{fixed-modules-pairings-and-lattices}

Work first over Z and Q; all assertions then localize at11 and extend to the actual coefficient DVR. Harmless matrix-unit phases may require adjoining i, which is an11-unit. Use exactly LIFT39's actual first and second partitions A,B,Lambda and joint vectors

\[
 P_0=p_A(X)y_b^B,\qquad P_1=p_\Lambda(Z),\qquad
 N_0=H_A B!,\quad N_1=H_\Lambda
 \tag{IL1}
\]

at one CM label. Tensor products over labels will be stated below. The primitive vertical maps satisfy T0(P0,v1)=v0 and T1(P1,v2)=v1. The actual unit vacuum is held fixed throughout.

Let V0,min be the minimal compact type of the bottom discrete series, and V1,V2 the middle and top compact representations. Remove their common last weights only for writing polynomial partitions. Their natural partitions are

\[
 \eta_0=A-L_*\mathbf1_a,\qquad
 \eta_1=(A+B\mathbf1_a,B^{[t]},0),\qquad
 \eta_2=(\Lambda,0^{[t]}).
 \tag{IL2}
\]

Indeed Lambda-L\_*1\_b=eta1. Determinant twists do not change the normalized compact pairings or the divided-power lattices. Write Delta\_i and Nabla\_i for the actual Weyl and dual-Weyl lattices of these modules. Their primitive extremals vi have normalized contravariant square norm1. The contragredient partner uses the same underlying lattice with the contragredient action; a superscript c will indicate that action when needed.

The following accepted source results are used with their exact normalizations, rather than reproved:

\begin{itemize}
\tightlist
\item
  \nolinkurl{GradedWeylDivisibility/StatementA.lean}34--38,68--80: Delta is generated from the primitive highest polynomial by integral row divided powers; polynomial row substitutions preserve the relevant integral modules.
\item
  \nolinkurl{LocalRingCardinality/DataCoeff.lean}126--129 and the proved \nolinkurl{FactDelDualNabla.lean} theorem: Delta and Nabla are perfect duals for the normalized Fischer form. Hence
\end{itemize}

\[
 \Delta_i\subset\nabla_i,\qquad
 H_{\eta_i}\nabla_i\subset\Delta_i.
 \tag{IL3}
\]

The second inclusion also follows directly: the ordinary Fischer pairing of integer polynomials is integral, and its normalized form divides by H\_eta\_i. This is the same argument as R020(R20.6).

Let F\_Z be the ordinary polynomial Fock lattice: coefficients in Z in the unit creation variables, with Fischer pairing {[}z\^{}gamma,z\^{}delta{]}\emph{F=delta}(gamma,delta) gamma!. It is distinct from its Fischer-dual divided-monomial lattice. In particular no division by a factorial is hidden in the phrase integral Fock polynomial.

\subsection{2. The intrinsic ordinary-PBW lattice and its explicit dual}\label{the-intrinsic-ordinary-pbw-lattice-and-its-explicit-dual}

Let H0 be the bottom holomorphic Harish--Chandra module in K043's polynomial realization. Put

\[
 \mathcal P_n=\Delta_0\otimes\mathbb Z[z_1,\ldots,z_a]_n,
 \qquad
 \mathcal H_n=C(\mathcal P_n),\qquad
 \mathcal H=\bigoplus_{n\ge0}\mathcal H_n,
 \tag{IL4}
\]

where C is the ordinary matrix-unit PBW creation map of K043. This defines the arithmetic lattice in the discrete series before the oscillator embedding is considered. It is not fitted to theta values. K043(PB6--PB7) implies that C is a rational isomorphism in each degree and, on each actual Pieri type lambda, C=D\_lambda id, with the explicit positive integer D\_lambda in K043(PB5). Thus H\_n is a free lattice of the full expected rank.

The compact divided-power action preserves P\_n and H\_n.~To see this without a nonintegral exponential argument, the compact group acts by the algebraic tensor action on Delta0 and the variables z; the coefficients of an elementary unipotent parameter are exactly the divided powers. The PBW map is compact equivariant. The ordinary noncompact creators map H\_n into H\_(n+1).

Give H0 its actual invariant form B\_H, with the minimal compact form normalized as above. K043(PB6), in bilinear contravariant form, is the exact identity

\[
 B_H(C(P),Q)=B_F(P,Q).
 \tag{IL5}
\]

It holds on the whole homogeneous polynomial space, not only on one Pieri component. Therefore the B\_H-dual of the intrinsic PBW lattice is the following explicit lattice in the polynomial model:

\[
 \boxed{\quad
 \mathcal H_n^{\# B_H}
  =\mathcal P_n^{\# B_F}
  =\nabla_0\otimes
      \bigoplus_{|\gamma|=n}\mathbb Z\frac{z^\gamma}{\gamma!}.
 \quad}
 \tag{IL6}
\]

\textbf{Proof.} Pairing against every C(P), P in P\_n, is equivalent by (IL5) to the Fischer pairing against every P in P\_n.~Perfect compact Delta/Nabla duality and the diagonal monomial factorials give the last equality. This proves equality of lattices, not just one containment. It also makes all dual denominators explicit. QED.

Write P\_n\^{}sharp for the lattice in(IL6). In particular

\[
             n!H_{\eta_0}\,\mathcal P_n^\sharp
                           \subset\mathcal P_n,
 \tag{IL7}
\]

because gamma! divides n! for every multiindex of total degree n.

\subsection{3. Integral joint embeddings and generated-lattice retractions}\label{integral-joint-embeddings-and-generated-lattice-retractions}

Let j0 and j1 be the actual rational Howe embeddings normalized by

\[
 j_0(v_0\otimes v_1^c)=P_0,\qquad
 j_1(v_1\otimes v_2^c)=P_1.
 \tag{IL8}
\]

The positive norm scalars are exactly N0,N1 from LIFT39. Thus the normalized Fischer form on each selected summand, divided by N\_j, is the tensor product of the displayed row forms. Denote their rational images by S0 and S1.

Define the generated polynomial lattices

\[
 \mathcal G_0=j_0(\mathcal H\otimes\Delta_1^c),\qquad
 \mathcal G_1=j_1(\Delta_1\otimes\Delta_2^c).
 \tag{IL9}
\]

These are integral sublattices of F\_Z. Here is a direct verification. At the minimal compact degree, the two independent compact groups generate Delta0 tensor Delta1\^{}c from the joint primitive extremal. Their row divided powers act integrally on the polynomial variables: they are coefficients of the integer unipotent substitution. The same assertion proves integrality of j1. On the first edge the ordinary noncompact matrix-unit creator is, up to the fixed sign,

\[
                    E_i=\sum_{c=1}^b X_{ci}y_c.
 \tag{IL10}
\]

It multiplies by an integer polynomial. Applying all ordinary PBW words proves integrality of j0 on H tensor Delta1\^{}c.~These exact operators and their derivative adjoints are also the literal \nolinkurl{RealPlaceAnalysis/FockPairGl.lean}15--31; no half or factorial is inserted in a noncompact PBW creator. The actual Hom/conjugate transport of LIFT39 preserves these assertions and their11-units.

The primitive contractions consequently obey

\[
 \begin{split}
 T_0(\mathcal G_0,\Delta_1)&=\mathcal H,\\
 T_1(\mathcal G_1,\Delta_2)&=\Delta_1.
 \end{split}
 \tag{IL11}
\]

\textbf{Proof.} On a generated tensor, T0(j0(x tensor z),w)=B1(z,w)x. The compact normalized form is integral on Delta1\^{}c by Delta1, so the image is in H. Taking z=v1\^{}c,w=v1 gives x, since their pairing is1. The section x maps to j0(x tensor v1\^{}c) is integral. The same argument proves the second equality and its section. QED.

For the two edges together there is an explicit integral section of their coefficient contraction on H:

\[
 x\longmapsto
 j_0(x\otimes v_1^c)\otimes P_1\otimes v_2.
 \tag{IL12}
\]

Successive contraction gives x exactly. This is a local polynomial section with the finite theta data untouched. It is not a theorem that this polynomial tensor is the fixed global Schwartz/theta image or that its automorphic preimage is integral.

\subsection{4. The saturated selected summand needs no joint hook denominator}\label{the-saturated-selected-summand-needs-no-joint-hook-denominator}

The saturated lattices in the independently fixed rational summands are

\[
                \mathcal S_j=F_{\mathbb Z}\cap S_j.
 \tag{IL13}
\]

They need not equal the generated lattices G\_j. Nevertheless

\[
 \boxed{\quad
 \begin{split}
 \mathcal H_n&\subset
 T_0(\mathcal S_{0,n},\Delta_1)
                    \subset\mathcal P_n^\sharp,\\
 \Delta_1&\subset
 T_1(\mathcal S_1,\Delta_2)
                    \subset\nabla_1.
 \end{split}\quad}
 \tag{IL14}
\]

Here S0,n has positive/negative bidegrees \textbar A\textbar+n and B+n.~The notation T(L,L') denotes the Z-span of its bilinear values.

\textbf{Proof of the upper inclusions.} The normalized summand form is B=Fischer/N\_j. Every generator of G\_j is an integral word in compact divided powers and, on the first edge, ordinary noncompact creators, applied to P\_j. The adjoint word preserves F\_Z: compact divided powers transpose to the opposite row divided powers; the noncompact adjoints are integer derivative contractions. All these operators preserve the rational selected summand.

More explicitly, a compact root operator acts on X and y through an integer elementary-unipotent substitution. The coefficient of z\^{}r in that substitution is its divided power E\_(ij)\^{}r/r!, hence preserves the ordinary polynomial lattice. Its Fischer adjoint is the opposite compact divided power, with the same property. For a noncompact pair the annihilator is D\_i=sum\_c partial\_(X\_ci) partial\_(y\_c). Even D\_i\^{}r/r! is integral on ordinary monomials: its multiindex term has coefficient product\_c beta\_c! binom(m\_(X\_ci),beta\_c) binom(m\_(y\_c),beta\_c). Our PBW lattice uses the ordinary creators E\_i\^{}r, whose adjoints D\_i\^{}r are consequently integral as well. The divided creator E\_i\^{}r/r! need not preserve this lattice and is not used in(IL4), (IL9), or this proof.

For v in the integral lattice mathcal S\_j and such a word A, move A across the Fischer pairing. The joint-extremal weight component of A* v lies in S\_j and F\_Z. That weight space of S\_j is one-dimensional, generated by P\_j. The polynomial P\_j has an ordinary monomial coefficient1, so this component is an integer multiple of P\_j. Since B(P\_j,P\_j)=1,

\[
                         B(v,A P_j)\in\mathbb Z.
 \tag{IL15}
\]

The weight projection just used preserves F\_Z: the compact torus weights are integer monomial weights, so projection selects a sum of monomials. On the first edge the negative-source U(1) weight selects the minimal bottom compact degree, and the remaining source and partner highest weights then give the stated one-dimensional joint weight. There is no additional weight multiplicity at that endpoint of the selected summand.

Equation(IL15) says that the inverse j0-image of mathcal S0 lies in the dual of H tensor Delta1\^{}c.~Contracting against Delta1 gives the dual H\^{}\#=P\^{}sharp of(IL6). The same argument with j1 gives Nabla1. The lower inclusions follow from G\_j subset S\_j and(IL11). QED.

This is the useful improvement over a direct unnormalized Fischer bound: neither N0 nor N1 occurs in(IL14). Its proof uses membership in the actual rational Howe summand, and cannot be applied to an arbitrary ordinary integral Fock polynomial.

\subsection{5. Explicit selected-PBW denominator bounds}\label{explicit-selected-pbw-denominator-bounds}

For an actual selected bottom type lambda of degree n, set

\[
 \mathcal H_{n,\lambda}=\mathcal H_n\cap V_\lambda,
 \quad
 \mathcal S_{0,n,\lambda}
       =F_{\mathbb Z}\cap j_0(V_\lambda\otimes V_1^c).
 \tag{IL16}
\]

These intersections are saturated in their stated ambient free lattices. They are defined from the fixed rational representations, not from the contraction image. Combining(IL7), K043's C=D\_lambda id on this Pieri component, and(IL14) gives

\[
 \boxed{\quad
 \mathcal H_{n,\lambda}\subset
 T_0(\mathcal S_{0,n,\lambda},\Delta_1)
 \subset
 \frac{1}{n!H_{\eta_0}D_\lambda}\,
                         \mathcal H_{n,\lambda}.
 \quad}
 \tag{IL17}
\]

Indeed n!H\_eta0 times an output belongs to P\_n and to V\_lambda. Applying C multiplies it by D\_lambda and puts it in H\_n.~The lower inclusion uses the integral section of(IL11) on the chosen intersection. This supplies an explicit denominator at11 for every actual word: v11(n!)+v11(H\_eta0)+v11(D\_lambda), with D\_lambda given by the actual K043 rising factorials. There is no uncomputed inverse Gram matrix in this bound.

The second selected edge satisfies the purely compact bound

\[
 \Delta_1\subset T_1(\mathcal S_1,\Delta_2)
             \subset\nabla_1\subset H_{\eta_1}^{-1}\Delta_1.
 \tag{IL18}
\]

For successive selected-summand tensors, multiplying the first bound by H\_eta1 clears the middle dual-Weyl comparison. Thus n!H\_eta0 H\_eta1 D\_lambda is an explicit denominator for the two-edge coefficient contraction into the intrinsic selected PBW lattice. This is an upper bound, not a claim that the factor H\_eta1 is necessary or that no cancellation occurs in a complete source coevaluation.

\subsubsection{Comparison with a single selected Weyl lattice}\label{comparison-with-a-single-selected-weyl-lattice}

Let xi\_lambda be a primitive generator of the highest weight line in H\_(n,lambda), with fixed positive leading coefficient. The Weyl lattice Delta\_lambda\^{}PBW generated from xi\_lambda is intrinsic. The usual maximal/minimal admissible-lattice argument gives

\[
 \Delta_\lambda^{\rm PBW}\subset\mathcal H_{n,\lambda}
 \subset\nabla_\lambda^{\rm PBW}
 \subset H_{\eta_\lambda}^{-1}\Delta_\lambda^{\rm PBW}.
 \tag{IL19}
\]

Here eta\_lambda is lambda minus its last weight. For completeness, the first inclusion is divided-power stability. Weight projections preserve H\_(n,lambda): P\_n has its integral direct sum of compact weight lattices, and the compact-equivariant isomorphism C transports that direct sum to H\_n.~To prove the second, pair an element of H\_(n,lambda) with any divided-power word in xi\_lambda using the contravariant form normalized to give xi\_lambda norm1. Move that word across the pairing. The resulting highest weight component is an integer multiple of xi\_lambda because the highest line is primitive in H\_(n,lambda). Perfect Delta/Nabla duality gives membership in Nabla. The last inclusion is(IL3).

Consequently an entirely explicit extra factor H\_eta\_lambda converts (IL17) to this selected Weyl lattice. If one instead uses the unit selected compact isometry, its comparison scalar is retained: write xi\_lambda=C(P\_lambda), where P\_lambda is primitive in the corresponding highest line of P\_n, and put q\_lambda=B\_F(P\_lambda,P\_lambda)\textgreater0. Then q\_lambda is an integer and

\[
 \|\xi_\lambda\|^2=D_\lambda q_\lambda,\qquad
 J_\lambda(p_\lambda)
       =(D_\lambda q_\lambda)^{-1/2}\xi_\lambda.
 \tag{IL20}
\]

This is exactly K043's non-isometric Pieri convention. It is not permissible to silently replace the primitive PBW lattice by the unit-normalized selected lattice. The integer q\_lambda is the finite Fischer norm of that primitive Pieri highest vector; no new global period or theta-image index is hidden in it. The bounds above hold on the intrinsic PBW lattice without evaluating q\_lambda.

\subsection{6. Whole ambient Fock lattices: bounds and exact nonintegral values}\label{whole-ambient-fock-lattices-bounds-and-exact-nonintegral-values}

For arbitrary integer Fock polynomials the weight argument(IL15) is unavailable, because their projections to S\_j can be nonintegral. The unnormalized Fischer pairing still gives the honest bounds

\[
 \begin{split}
 T_0(F_{\mathbb Z,n},\Delta_1)&\subset N_0^{-1}\mathcal P_n^\sharp,\\
 T_1(F_{\mathbb Z,|\Lambda|},\Delta_2)&\subset N_1^{-1}\nabla_1.
 \end{split}
 \tag{IL21}
\]

To prove them, pair the output against H or Delta1, respectively. Adjunction gives N\_j times that row pairing equal to the ordinary Fischer pairing against the integral joint embedding(IL9), which is an integer. Perfect duality and(IL6) give exactly(IL21). These bounds also tensor over the CM labels.

The loss is real, even though P\_j itself is primitive. Define the ordinary coefficient-one monomials

\[
 m_0=\prod_{i=1}^{a}X_{ii}^{A_{i-1}}\,y_b^B,\qquad
 m_1=\prod_{i=1}^{b}Z_{ii}^{\Lambda_{i-1}}.
 \tag{IL22}
\]

Their images under the primitive contractions are exactly

\[
 \boxed{\quad
 T_0(m_0,v_1)=r_0v_0,\quad
 r_0=\frac{\prod_i A_i!}{H_A};\qquad
 T_1(m_1,v_2)=r_1v_1,\quad
 r_1=\frac{\prod_i\Lambda_i!}{H_\Lambda}.
 \quad}
 \tag{IL23}
\]

\textbf{Proof.} The row and partner weights are the corresponding joint-extremal weights. After contraction the output is therefore in the one-dimensional extremal line; on the first edge its U(1) weight also forces the minimal compact degree. Pairing with the normalized extremal, adjunction gives the Fischer coefficient of m\_j in P\_j divided by N\_j. That coefficient is1. The ordinary Fischer monomial norm is product\_i A\_i! times B! on the first edge, and product\_i Lambda\_i! on the second. The B! cancels in the first ratio, proving(IL23). QED.

The actual valuations in the endpoint profiles of K044 are listed below. Here q\_max=28(a+1) floor(ell/{[}28(a+1){]}) is the largest allowed canonical twist, so every listed endpoint is admissible. The larger formal envelope q=ell is not substituted for q\_max. LIFT39's frozen certificate also tests q=ell as a formal-envelope endpoint of its formulas, which hold throughout0\textless=q\textless=ell. That diagnostic is not an additional admissible canonical twist. No change to LIFT39's theorem or prior-check hash is made here.

\begin{longtable}[]{@{}llllrr@{}}
\toprule\noalign{}
Host & Role & ell & q\_can & v11(r0) & v11(r1) \\
\midrule\noalign{}
\endhead
\bottomrule\noalign{}
\endlastfoot
J14 & top & min & 0 & -27 & -3390 \\
J14 & top & min & q\_max & -27 & -1799 \\
J14 & top & max & 0 & -30 & -3422 \\
J14 & top & max & q\_max & -30 & -3981 \\
J14 & thick & min & 0 & -30 & -3395 \\
J14 & thick & min & q\_max & -30 & -1800 \\
J14 & thick & max & 0 & -29 & -3423 \\
J14 & thick & max & q\_max & -29 & -3988 \\
J14 & thin & min & 0 & -27 & -3386 \\
J14 & thin & min & q\_max & -27 & -1800 \\
J14 & thin & max & 0 & -20 & -3388 \\
J14 & thin & max & q\_max & -20 & -3968 \\
J56 & top & min & 0 & -151 & -3001 \\
J56 & top & min & q\_max & -151 & -2794 \\
J56 & top & max & 0 & -153 & -2266 \\
J56 & top & max & q\_max & -153 & -2746 \\
J56 & thick & min & 0 & -149 & -2993 \\
J56 & thick & min & q\_max & -149 & -2794 \\
J56 & thick & max & 0 & -154 & -2264 \\
J56 & thick & max & q\_max & -154 & -2747 \\
J56 & thin & min & 0 & -155 & -3001 \\
J56 & thin & min & q\_max & -155 & -2788 \\
J56 & thin & max & 0 & -155 & -2270 \\
J56 & thin & max & q\_max & -155 & -2749 \\
\end{longtable}

Every entry is an exact Legendre-floor and finite Vandermonde calculation. The certificate retains the complete actual integers and also computes H\_eta0 and H\_eta1. No factorial of the enormous weights is materialized. The displayed table concerns precisely these endpoint profiles; negativity at every intermediate ell or q\_can is not inferred from it.

Thus no integral map on the entire ordinary Fock lattice can satisfy the primitive normalization(PN15) in these profiles. The obstruction is already present before any automorphic class sum. It is not a statement that a selected theta row has these poles.

The image on the whole first Fock lattice is stable under the integral compact divided powers and ordinary noncompact creators, since T0 intertwines them. It contains r0 v0, hence contains r0 H. In each degree and in each saturated selected PBW type the image therefore contains r0 times that lattice. Together with the generated-lattice section it contains H+r0 H. This is a statement about the available ambient polynomial inputs; it does not place the monomials or their projections in the fixed arithmetic kernel.

Similarly the second ambient image contains Delta1+r1 Delta1. The composed ambient contraction sends m0 tensor m1 tensor v2 to r0 r1 v0, with no omitted coefficient. This gives an explicit necessary denominator for an ambient integral normalization.

\subsubsection{Two exact small diagnostics of the distinction}\label{two-exact-small-diagnostics-of-the-distinction}

These examples illustrate the mechanism; their ranks are not substituted for the actual host ranks. On2-by-2 matrices, for shape (10,1) take P=x00\^{}9 det(X). Its Fischer norm is11\emph{9!, whereas the diagonal monomial x00\^{}10}x11 pairs with it as10!. The primitive Cauchy contraction therefore takes that monomial to10/11 times the normalized extremal. This is an ordinary ambient projection pole at11.

For shape(12,1), put P=x00\^{}11 det(X). Its first row-lowering divided-power generator gives11\emph{x00\^{}10}x10\emph{det(X). Hence the integer polynomial Q=x00\^{}10}x10*det(X) belongs to the saturated selected summand and contracts to(1/11) times the corresponding primitive Weyl lowering vector. That vector has normalized square norm11, so its quotient by11 belongs to the dual-Weyl lattice. Thus the dual-Weyl upper bound in(IL14) cannot in general be strengthened to the Weyl lattice. The generated-lattice retraction and the saturated-summand bound have genuinely different scope.

\subsection{7. Source frames, CM conjugation and remaining global scope}\label{source-frames-cm-conjugation-and-remaining-global-scope}

For all CM labels, tensor the displayed lattices, maps and integral sections. In each finite multidegree the saturated sublattices are direct summands as Z-modules; their tensor product is therefore the intersection of the tensor rational summand with the ambient tensor polynomial lattice. The selected two-edge PBW denominator is exactly the displayed sufficient product

\[
 \mathfrak d_{\boldsymbol\lambda}
  =\prod_\tau n_\tau!\,H_{\eta_{0,\tau}}
                     H_{\eta_{1,\tau}}D_{\lambda_\tau}.
 \tag{IL24}
\]

For the selected Weyl comparison multiply additionally by product\_tau H\_eta\_lambda,tau. For full ambient Fock inputs the sufficient products additionally contain product\_tau N0,tau N1,tau. The monomial values are the products of r0,tau and r1,tau. Flat base change and localization preserve every inclusion. These statements give finite explicit11-adic denominator bounds, not assertions that those bounds are sharp.

All inclusions are proved for the exact unit-creation polynomial coordinates and the primitive row conventions of LIFT39. An integral matrix change with integral inverse preserves the Weyl lattices and transports the statements unchanged. For a nonintegral source frame its actual matrix and inverse must be applied; there is no invariance claim for an arbitrary real change of basis.

In particular LIFT39(PN23--PN24) is retained verbatim as a separate analytic/arithmetic conversion. The raw CM vacuum contributes the explicit positive determinant amplitude, the positive and negative first-edge columns have conjugate C\_theta/creation/frame factors, and the algebraic coordinates cD remain fixed. A transcendental unit-creation scale is not assigned an11-adic valuation. The lattice bounds here do not identify the remaining negative CM row frame with every literal arithmetic pin.

The generated joint module has an explicit integral local retraction. The saturated selected module has the explicit dual lattice and selected-PBW denominator bounds. The full ambient projection is demonstrably nonintegral. These settle different pieces of the primitive contraction normalization. None proves that a fixed global theta-image lattice is saturated, that the arithmetic vertical inverse has the complex norm of LIFT38, or that a net DP44 cost improvement follows from a common scalar.

\subsection{8. Sources and verification scope}\label{sources-and-verification-scope-3}

The carried exact inputs are LIFT39(PN1--PN25), K043(PB5--PB8), R020(R20.1--R20.6), and the source Weyl/dual-Weyl theorem identified in Section1. The integer creation/adjoint operators are read directly in \nolinkurl{Formalism/RealPlaceAnalysis/FockPairGl.lean}9--31,48--53. The PBW duality, primitive-weight integrality argument, ambient monomial calculation, and actual valuation table are proved here.

The author read the complete IL1--IL24 proof, including the intrinsic PBW dual identity, both integral joint embeddings, the adjoint divided-power and weight-projection argument, the distinct ambient projection, and the actual canonical-twist restriction. All source operator and Weyl-duality locators above were checked directly. The selected PBW/Weyl comparison retains its unit-isometry scalar instead of silently treating the two lattices as equal.

The companion \nolinkurl{vertical_primitive_integral_lattices_certificate.py} tests ordinary versus divided monomials,891 divided-annihilator coefficients,649 exact small Cauchy generated/saturated joint weights, both explicit denominator11 examples, and all24 actual admissible endpoint profiles. It also checks48 actual selected PBW degrees and their exact factorial/rising-factorial denominator bounds. These are arithmetic diagnostics, not evaluations of a global theta matrix or a global integral preimage.

\clearpage

\section{AM-K055 --- the actual finite mirabolic test at the two mixed G places}\label{am-k055-the-actual-finite-mirabolic-test-at-the-two-mixed-g-places}

Author: GPT-6 Astra Pro, ALP\_MACHINE /root/kummer\_extension. 22 November 2026 UTC. Inherited reasoning effort is not separately exposed.

\textbf{Author whole-proof PASS --- 2026-11-22T06:31:04.653070+00:00.}

This supplies a full local test in the actual mixed finite type. The goodTuple factors remain (37,6) or (163,6); no finite type, level, global constituent, or initial vector is replaced. The bottom vector is chosen inside its admitted finite representation. The argument also applies at the full-degree places of R033. Its compact support is support modulo the Whittaker unipotent, rather than the matrix-coefficient support used in R033.

\subsection{1. Exact local statement}\label{exact-local-statement}

Let F be a nonarchimedean local field, O its ring of integers, p=(varpi), k=O/p and q=\textbar k\textbar. For j\textgreater=1 put

\[
 G_j=\operatorname{GL}_j(F),\quad K_j=\operatorname{GL}_j(O),
 \quad K_j^1=\ker(K_j\to\operatorname{GL}_j(k)),
 \quad P_j=\{g:e_jg=e_j\}.
 \tag{MK1}
\]

Let N\_j be the upper unitriangular subgroup. Thus P\_1=\{1\}. Use G\_0=N\_0=K\_0=\{1\} for the zero-rank integrals below. Write P\_j(O)=P\_j intersect K\_j and P\_j\^{}1=P\_j intersect K\_j\^{}1. Let V\_j be the last-column unipotent radical of P\_j, so P\_j=G\_(j-1) semidirect V\_j.

Fix a nontrivial additive character chi of F with conductor p: it is trivial on p and nontrivial on O. Choose the compatible nondegenerate characters

\[
 \psi_j(u)=\chi\!\left(\sum_{i=1}^{j-1}u_{i,i+1}\right).
 \tag{MK2}
\]

The actual simple-root coefficients may be arbitrary units. A diagonal element of K\_a with last entry1 transports them to (MK2). Extend it to diag(d\_a,1) in K\_(a+1). This transports both representations, both characters, and the embedding simultaneously; it preserves the standard lattices and all Haar normalizations. The last new simple-root coefficient can be chosen to be1. Consequently the proof applies to the literal positive/inverse characters of R025/LIFT32, rather than an unspecified generic character. The choice of Whittaker character is a model choice; it does not change the source additive character or finite action.

Set n=a+1. Let pi\_G be a generic tempered representation of G\_n whose K\_n\^{}1-fixed vectors contain the inflation of the admitted irreducible generic finite representation Pi of GL\_n(k). At the two exceptional places Pi is the actual irreducible parabolic product with the retained mixed factors. Let pi be the actual generic tempered initial representation of G\_a and let y be any Hilbert-unit vector satisfying

\[
 y\in\pi^{K_a^1},\qquad
 \pi(u)y=\psi_a(u)^{-1}y\quad(u\in N_a(O)).
 \tag{MK3}
\]

These are exactly the unchanged initial G-level conditions. The source translate at G is the identity. Give K\_a volume1.

\textbf{Theorem.} There is a Hilbert-unit vector x inside the same finite Pi-subspace, with the prescribed positive N\_a(O) character, for which

\[
\begin{split}
 \alpha_G(x,y)
 &=\int_{G_a}
    \langle\pi_G(\operatorname{diag}(h,1))x,x\rangle
    \langle\pi(h)y,y\rangle\,dh_{\mathrm{prob}}\\
 &=\boxed{\frac1{\dim\sigma}}>0.
\end{split}
 \tag{MK4}
\]

Here sigma is the unique irreducible generic finite GL\_a(k) representation generated by y. Its existence and uniqueness are proved below; they are not assumptions on the initial finite type. In particular

\[
 \alpha_G(x,y)\ge\frac1{D_a(q)}>q^{-a(a+1)/2},
 \qquad D_j(q)=\prod_{i=1}^j(q^i-1),\quad D_0=1.
 \tag{MK5}
\]

This is an evaluation over the entire noncompact G\_a. It proves neither a global central value nonzero nor an integral-lattice bound for the chosen unit vector.

\subsection{2. An elementary finite-mirabolic support lemma}\label{an-elementary-finite-mirabolic-support-lemma}

Inflate to P\_j(O) the irreducible finite mirabolic representation

\[
 I_j=\operatorname{Ind}_{N_j(k)}^{P_j(k)}\psi_j,
 \qquad \dim I_j=D_{j-1}(q).
 \tag{MK6}
\]

Finite induction uses functions f(up)=psi\_j(u)f(p) and right translation. LIFT32(FB2--FB5) proves its irreducibility by the finite little-group argument. In particular V\_j(k) has no invariant vectors in I\_j. Its nonzero character spaces form one orbit under GL\_(j-1)(k). The character

\[
       \xi_0(v)=\chi(v_{j-1})\qquad(v\in V_j(O))
 \tag{MK7}
\]

has stabilizer P\_(j-1)(k), and its eigenspace carries I\_(j-1). The base I\_1 is the one-dimensional representation of \{1\}. These assertions also follow immediately by induction in stages from (MK6), and fix the actual row/column convention.

Let F\_j be the space of smooth functions on P\_j(F) satisfying f(up)=psi\_j(u)f(p) for u in N\_j(F), with right translation. No support condition is imposed on F\_j. Define its subspace

\[
 C_j=\{f\in F_j:
       \operatorname{supp}f\subset N_j(F)P_j(O),
       \ f(pp_1)=f(p)\ (p_1\in P_j^1)\}.
 \tag{MK8}
\]

Restriction to P\_j(O), followed by reduction, identifies C\_j with I\_j. Indeed N\_j(F) intersect P\_j(O)=N\_j(O), and the character is trivial on N\_j(O) intersect P\_j\^{}1. Extension from the finite induced function by left equivariance and zero outside N\_j(F)P\_j(O) is well defined and smooth. The image of P\_j(O) in N\_j(F)\P\_j(F) is compact and open, so this extension has compact support modulo N\_j.

\textbf{Support lemma.} Every P\_j(O)-submodule of F\_j isomorphic to I\_j equals C\_j. In particular the I\_j-isotypic part of F\_j has multiplicity one.

\textbf{Proof.} Proceed by induction on j. The case j=1 is immediate. Let S be a copy of I\_j. Since I\_j is inflated from the finite group, S is fixed pointwise by P\_j\^{}1. It suffices first to treat f in the xi\_0 eigenspace S\_0 of V\_j(O).

For h in GL\_(j-1)(F), the Whittaker identity gives

\[
 f\!\left(\begin{pmatrix}h&0\\0&1\end{pmatrix}
                  \begin{pmatrix}1&v\\0&1\end{pmatrix}\right)
       =\chi(e_{j-1}h v)\,
                  f\!\left(\begin{pmatrix}h&0\\0&1\end{pmatrix}\right).
 \tag{MK9}
\]

If the displayed value of f is nonzero, its right xi\_0 character forces

\[
             e_{j-1}h\equiv e_{j-1}\pmod{pO^{j-1}}.
 \tag{MK10}
\]

This follows from the exact annihilator of O for chi: it is p. It simultaneously proves integrality of that row and its required nonzero reduction. There is k in K\_(j-1)\^{}1 with last row e\_(j-1)h. For example, keep its first j-2 rows equal to those of the identity and take that last row; its last entry is a unit congruent1, so it is invertible and congruent to the identity. Then

\[
       p_0=h k^{-1}\in P_{j-1}(F),\qquad
       f(\operatorname{diag}(h,1))
              =f(\operatorname{diag}(p_0,1)).
 \tag{MK11}
\]

The equality uses the right P\_j\^{}1 invariance, not a nonintegral change of frame. Restriction of S\_0 to the embedded P\_(j-1)(F) is a P\_(j-1)(O)-intertwining map into F\_(j-1). Since S\_0 is I\_(j-1), its image is zero or a copy of that representation. By induction, any nonzero restricted value requires p\_0 in N\_(j-1)(F)P\_(j-1)(O). Hence h is in N\_(j-1)(F)K\_(j-1).

Every element of P\_j(F) is in N\_j(F)diag(GL\_(j-1)(F),1), so this proves support in N\_j(F)P\_j(O) for f in S\_0. The finitely many P\_j(O)-translates of S\_0 span S by the little-group description. The asserted support set is right P\_j(O)-stable; therefore the support assertion holds for all S. It follows that S is contained in C\_j. Both are I\_j, so they are equal. The same reasoning applies to every copy, proving multiplicity one. This completes the induction. \(\square\)

The lemma does not assert compact support of a matrix coefficient on G\_j. Its support is exactly the support needed by the Rankin--Selberg functional.

\subsection{3. Applying the lemma to the actual unitary representations}\label{applying-the-lemma-to-the-actual-unitary-representations}

Use the Whittaker models W(pi\_G,psi\_n) and W(pi,psi\_a\^{}\{-1\}). In either sign, Zhang (2014), (3.2) and the following paragraph, printed553--554, gives the invariant positive Whittaker inner product

\[
 \theta_{j,\mathrm{prob}}(W,W)
   =\int_{N_{j-1}(F)\backslash G_{j-1}}
        |W(\operatorname{diag}(h,1))|^2\,dh_{\mathrm{prob}}.
 \tag{MK12}
\]

Here K\_(j-1) and N\_(j-1)(O) have volume1; changing these positive Haar constants only rescales the inner product. It is finite and positive definite for generic unitary representations. Consequently restriction of W to P\_j(F) is injective: a zero restriction has zero positive norm in (MK12).

The actual finite Pi is generic. LIFT32 supplies its I\_n submodule under P\_n(O). By injectivity and the support lemma, its restriction is exactly C\_n.~In particular, restricting further to diag(G\_a,1), this is the whole space of functions

\[
 F(h)\text{ supported on }N_a(F)K_a,
 \quad F(uh)=\psi_a(u)F(h),\quad F(hk_1)=F(h)
                     \ (k_1\in K_a^1).
 \tag{MK13}
\]

Its restriction to K\_a is the finite Gelfand--Graev representation Ind\_(N\_a(k))\^{}(GL\_a(k)) psi\_a. Every such finite function is obtained, because the restriction I\_n\textbar GL\_a(k) is exactly that representation, with the same dimension D\_a(q).

We must also check the prescribed initial y. Its P\_a(O)-span is a copy of I\_a for the inverse character. Indeed finite Frobenius reciprocity maps I\_a to that span by sending the function supported on N\_a(k) to y; it is nonzero, and I\_a is irreducible. The support lemma and the injective Whittaker restriction show that W\_y\textbar P\_a is a nonzero function in C\_a. The right N\_a(O) inverse-character eigenspace of C\_a is one-dimensional by finite Frobenius reciprocity and irreducibility. Its generator is the function supported on N\_a(F)P\_a\^{}1 with value1 at1. Thus

\[
 W_y(1)=c\ne0,\qquad
 W_y|_{P_a}=c\,\delta_{N_a}
 \quad\text{in this normalized finite induced model}.
 \tag{MK14}
\]

In particular W\_y\textbar K\_a is nonzero. This is proved for the given y, rather than obtained by choosing a different initial vector.

The same argument shows that the whole finite space (pi\^{}\{K\_a\^{}1\})\^{}\{N\_a(O),psi\_a\^{}\{-1\}\} has dimension at most one: all its I\_a copies inject into the single copy C\_a. It is nonzero because it contains y. Finite semisimplicity and the one-dimensional generic eigenspace in each irreducible finite GL\_a representation therefore imply that exactly one generic irreducible sigma occurs, with multiplicity one, and y is in that sigma. The K\_a-span of y is exactly sigma.

\subsection{4. The fixed generic test and its exact two norms}\label{the-fixed-generic-test-and-its-exact-two-norms}

On K\_a define

\[
                  F_x(k)=\overline{W_y(k)}.
 \tag{MK15}
\]

This is a finite left-psi\_a function and extends to (MK13). By that statement there is a unique x in the chosen I\_n inside the unchanged finite Pi with this restriction. Because W\_y has right inverse character under N\_a(O), F\_x has right positive character. Injectivity of restriction proves the same exact positive character for x. Thus x has the literal bottom G-level and genuine central projector. It does not require a translated or deeper finite level.

Write d\_sigma=dim(sigma), and put

\[
 B=\int_{K_a}|W_y(k)|^2\,dk_{\mathrm{prob}}.
 \tag{MK16}
\]

The evaluation functional W-\textgreater W(1) on sigma has the inverse generic character. By one-dimensionality, its Hermitian Riesz vector is a scalar multiple of y. Finite Schur orthogonality therefore gives, for any scaling of y,

\[
                         B=\frac{|c|^2}{d_\sigma}.
 \tag{MK17}
\]

For clarity, if \textbar\textbar y\textbar\textbar\^{}2 is not1 in this Whittaker realization, the Riesz vector has squared norm \textbar c\textbar\^{}2/\textbar\textbar y\textbar\textbar\^{}2; Schur's formula multiplies it by \textbar\textbar y\textbar\textbar\^{}2/d\_sigma, giving exactly (MK17).

The support in (MK13) and the fact that N\_a(O) has volume1 give

\[
 \theta_{n,\mathrm{prob}}(W_x,W_x)=B,\qquad
 \lambda_{\mathrm{prob}}(W_x,W_y)
       =\int_{N_a\backslash G_a}F_x(h)W_y(h)\,dh_{\mathrm{prob}}
       =B>0.
 \tag{MK18}
\]

The product is independent of the representative in N\_a\G\_a because the two Whittaker characters are opposite. Its support modulo N\_a is the compact image of K\_a, so (MK18) is an actual convergent local Rankin--Selberg integral at its central point.

On the initial side (MK14) restricts to a function of constant modulus \textbar c\textbar{} supported on N\_(a-1)(F)K\_(a-1)\^{}1. Its quotient volume is exactly

\[
 \frac{\operatorname{vol}(K_{a-1}^1)}
      {\operatorname{vol}(N_{a-1}(O)\cap K_{a-1}^1)}
   =\frac{|N_{a-1}(k)|}{|\operatorname{GL}_{a-1}(k)|}
   =\frac1{D_{a-1}(q)}.
 \tag{MK19}
\]

Hence

\[
      \theta_{a,\mathrm{prob}}(W_y,W_y)
           =\frac{|c|^2}{D_{a-1}(q)},\qquad
 \frac{|\lambda_{\mathrm{prob}}(W_x,W_y)|^2}
      {\theta_{n,\mathrm{prob}}(W_x,W_x)
       \theta_{a,\mathrm{prob}}(W_y,W_y)}
           =\frac{D_{a-1}(q)}{d_\sigma}.
 \tag{MK20}
\]

These are exact norms of the actual functions, not replacements of their norms by dimensions of an ambient representation.

\subsection{5. Conductor-p Haar normalization and the complete period}\label{conductor-p-haar-normalization-and-the-complete-period}

The character in this proof has conductor p, which matters for the selfdual measures in Zhang's Proposition4.10. Let dx\_0 give O volume1. The selfdual measure for chi is q\^{}(1/2)dx\_0. Indeed the Fourier transform of the characteristic function of O is q\^{}(1/2) times the characteristic function of p; applying it twice returns that function. Write

\[
 \kappa_j=\prod_{i=1}^j(1-q^{-i}),\qquad
 g_j=q^{j^2/2}\kappa_j,\qquad
 r_j=q^{j(j+1)/4}\kappa_j.
 \tag{MK21}
\]

The raw differential-form measure on GL\_j is g\_j times its maximal-compact-probability measure. The selfdual product measure on N\_j is q\^{}\{j(j-1)/4\} times the measure giving N\_j(O) volume1. Consequently the raw quotient measure on N\_j\GL\_j is r\_j times the quotient used in (MK12) and (MK18). The exponents here count j\^{}2 matrix entries and j(j-1)/2 unipotent entries separately.

Zhang(2014), Proposition4.10, proves in precisely these selfdual raw measures that the complete matrix-coefficient period is the squared Rankin--Selberg functional. Apply it to W\_x tensor W\_y, divide by their raw mirabolic norms, and then divide the outer GL\_a measure by g\_a. The result for Hilbert-unit vectors is

\[
\begin{split}
 \alpha_G(x_{\mathrm{unit}},y_{\mathrm{unit}})
 &=\frac1{g_a}
    \frac{r_a^2|\lambda_{\mathrm{prob}}(W_x,W_y)|^2}
         {r_a\theta_{n,\mathrm{prob}}(W_x,W_x)
          r_{a-1}\theta_{a,\mathrm{prob}}(W_y,W_y)}\\
 &=\frac1{D_{a-1}(q)}
    \frac{|\lambda_{\mathrm{prob}}(W_x,W_y)|^2}
         {\theta_{n,\mathrm{prob}}(W_x,W_x)
          \theta_{a,\mathrm{prob}}(W_y,W_y)}
   =\frac1{d_\sigma}.
\end{split}
 \tag{MK22}
\]

The middle equality uses the exact identity

\[
 \frac{r_a}{g_a r_{a-1}}
   =\frac{q^{-a(a-1)/2}}{\kappa_{a-1}}
   =\frac1{D_{a-1}(q)}.
 \tag{MK23}
\]

This proves (MK4). Arbitrary invariant Hilbert realizations give the same answer on their unit vectors. In particular no scalar depending on an unspecified Whittaker functional or its value at1 has been omitted.

Finally d\_sigma\textless=D\_a(q), because finite GL Gelfand--Graev is multiplicity free and the sum of dimensions of its generic irreducibles is D\_a(q). This proves (MK5).

\subsection{6. Actual mixed places, all G factors, and remaining global scope}\label{actual-mixed-places-all-g-factors-and-remaining-global-scope}

The proof uses only genericity of the admitted finite Pi, not cuspidality or a full-degree polynomial. LIFT32 section5 proves this genericity for the actual squarefree DO042/043 finite series. It therefore applies at the unchanged two mixed places with factor degrees (37,6) or (163,6). No use of the optional (p,3,3) eleven-integral alternative is needed.

It also applies at all other G places. For independent prescribed local initial vectors y\_v, choosing the bottom vectors as above gives

\[
 \prod_{v\mid G}\alpha_v(x_v,y_v)
   =\prod_{v\mid G}\frac1{\dim\sigma_v}
   \ge[D_a(g_0)D_a(g_1)]^{-f}>2^{-C_G^{\mathrm{all}}},
 \tag{MK24}
\]

where

\[
 C_G^{\mathrm{all}}
   =\frac{a(a+1)}2f
       (\operatorname{bits}(g_0)+\operatorname{bits}(g_1)).
 \tag{MK25}
\]

\begin{longtable}[]{@{}lrrr@{}}
\toprule\noalign{}
Host & a=f & Number of G places & C\_G\^{}all \\
\midrule\noalign{}
\endhead
\bottomrule\noalign{}
\endlastfoot
J\^{}14 & 42 & 84 & 1,744,596 \\
J\^{}56 & 168 & 336 & 169,329,888 \\
\end{longtable}

The bound concerns squared-period integrals with K\_a volume1. The canonical unitary GGP factors remain

\[
 \alpha_v^\natural
 =\frac1{\dim\sigma_v\,\mathcal L_v(1/2)},\qquad
 \mathcal L_v(1/2)
 =\left(\prod_{j=1}^{a+1}\zeta_q(j)\right)
  \frac{|L(1/2,\pi_G\times\pi)|^2}
       {L(1,\pi_G\times\pi_G^\vee)
        L(1,\pi\times\pi^\vee)}.
 \tag{MK26}
\]

They are finite and positive for the actual tempered data. These factors are retained; no adjoint or central local L-factor is silently set to1. A measure c\_v times the displayed measure multiplies the period and its canonical factor by c\_v.

Each choice x\_v is inside the already existing local constituent, with the same admitted finite Pi, positive generic level character, and genuine e0. The global restricted tensor product therefore permits these simultaneous local choices without changing the archimedean representation or the other finite vectors. This closes the two exceptional local-test qualifications in R033. It does not select a global pair with nonzero central L-value, prove the real fixed-tensor seesaw identity, or identify a nonzero permitted terminal arithmetic row or its integral preimage. The maps in this proof are Hilbert-space and Whittaker maps; their coefficient-field and integral-denominator costs are not claimed to be units or bounded by (MK24).

\subsection{7. Primary anchors and verification scope}\label{primary-anchors-and-verification-scope}

\begin{itemize}
\tightlist
\item
  Zhang, W. (2014). Automorphic period and the central value of Rankin--Selberg L-function. \emph{Journal of the American Mathematical Society, 27}(2),541--612. \url{https://doi.org/10.1090/S0894-0347-2014-00784-0}. Original publisher offprint: \url{https://archive.ymsc.tsinghua.edu.cn/pacm_download/271/8715-jams784_off_print.pdf}. Measures and Fourier selfduality in section2.1, printed550--551; invariant positive mirabolic norm (3.2) and its unitary convergence, printed553--554; central Rankin--Selberg functional (3.23)--(3.26), printed559; raw-measure change before Lemma4.7, printed569; complete Proposition4.10 and proof, printed571--574. These primary passages were directly read in the K054 source audit.
\item
  LIFT32(FB2--FB5, section5), with its primary finite mirabolic and finite GL Gelfand--Graev sources, supplies I\_j and the generic finite Pi. The support lemma here explicitly checks the passage from that finite module to the actual p-adic Whittaker model. No unproved transfer is supplied by the name of a finite type.
\item
  K050/051 supplies the existing genuine bottom constituents; DO046 supplies actual temperedness and the unitary GGP normalization. K038/LIFT27 and LIFT29 fix the initial and bottom levels and their opposite G characters. R033 is used only for comparison and its already checked actual G primes.
\end{itemize}

The new argument is the explicit finite-mirabolic support and given-vector evaluation lemma, its actual mixed-place test, and the exact conductor-p Haar cancellation. The published unitary Whittaker pairing and matrix-coefficient identity are used as theorems. A finite certificate checks representative matrix reductions, finite induced models, Haar exponents, Schur norms and actual prime exponents; it does not compute an automorphic representation or a global period.

\subsection{Author check record}\label{author-check-record-1}

The complete saved MK1--26 proof and full certificate were read and reconstructed. The finite mirabolic character/stabilizer convention, last-row reduction, support induction and uniqueness, prescribed initial-vector evaluation, finite Schur norm, conductor-p selfdual Haar factors, full Zhang period comparison and actual mixed-place application all passed. The exact certificate passed 3,465 assertions at 2026-11-22T06:29:48.570003+00:00, including seven complete finite mirabolic models (1,810 group elements and 35,224 induced-character actions), 240 rational row reductions, 159 annihilator cases and110 Haar rank/prime cases. The general support lemma is proved in the text; it is not inferred from these finite diagnostics. No Lean build, actual global central value or arithmetic-image computation occurred.

\clearpage

\section{AM-K056 --- all remaining finite GGP tests are the actual unramified tests}\label{am-k056-all-remaining-finite-ggp-tests-are-the-actual-unramified-tests}

Author: GPT-6 Astra Pro, ALP\_MACHINE /root/kummer\_extension. 22 November 2026 UTC. Inherited reasoning effort is not separately exposed.

\textbf{Author whole-proof and direct-source audit complete; see the check record below.}

This is a source and local-model audit for the initial--bottom period of LIFT29. It proves that at every finite place outside G and T its actual local vectors are spherical for compatible hyperspecial subgroups. In particular, a ramified auxiliary tuple at Q does not ramify these endpoint vectors. Its nontrivial higher vertical kernels and their normalization remain separate.

\subsection{1. The Q endpoint data, read literally}\label{the-q-endpoint-data-read-literally}

All source locators are relative to work/research/alpoge\_source/Formalism/LocalToGlobal, commit 7cbf7a0c32fca1badf05931879a37c3f8a2b03b4.

\begin{longtable}[]{@{}
  >{\raggedright\arraybackslash}p{(\columnwidth - 4\tabcolsep) * \real{0.3333}}
  >{\raggedright\arraybackslash}p{(\columnwidth - 4\tabcolsep) * \real{0.3333}}
  >{\raggedright\arraybackslash}p{(\columnwidth - 4\tabcolsep) * \real{0.3333}}@{}}
\toprule\noalign{}
\begin{minipage}[b]{\linewidth}\raggedright
Source locator
\end{minipage} & \begin{minipage}[b]{\linewidth}\raggedright
Literal datum
\end{minipage} & \begin{minipage}[b]{\linewidth}\raggedright
Consequence at an auxiliary Q place
\end{minipage} \\
\midrule\noalign{}
\endhead
\bottomrule\noalign{}
\endlastfoot
\nolinkurl{Columns/ColForms.lean92–114} & colGSD uses only Pr.vT(u) and Pr.vGs; MfV=ColV(u,0) & The initial coefficient module has no Q decoration or Q character. \\
\nolinkurl{IndefiniteReal/LinksDefs.lean430–431} & dat0=datY(\ldots,j=0,\ldots,Q=empty,\ldots,xi=1,Frm.B) & The bottom endpoint has empty Q data even when the tuple is ramified at Q. \\
\nolinkurl{IndefiniteReal/SubgroupCharExtensionOneApply.lean32–43} & datY stores s=j, its Q set and xq; Frm.B inverts the stored xi & At dat0, s=0 and Q is empty, so the inversion creates no Q character. \\
\nolinkurl{Divisibility/DefGroup.lean635–637} & parLev is unitary lattice stabilizer intersected with G and Q parahorics & Away from the displayed nonempty sets, only the lattice stabilizer remains. \\
\nolinkurl{Divisibility/DefGroup.lean639–640,689–703} & rTup and the residual level/character involve Pg; levFn's last product is over D.Q & At Q for dat0 there is no remaining residual restriction or determinant character. \\
\nolinkurl{IndefiniteReal/SubgroupCharExtensionOneApply.lean71–73} & right transformation is the literal levChi & The bottom endpoint character at Q is exactly1. \\
\end{longtable}

The initial interpretation is the checked K038(CA27a): its G level has the inverse generic character, its T level is the strict covector stabilizer, and it is the maximal lattice stabilizer with trivial character at every other place, including Q and29. LIFT29(BT3) gives the same explicit endpoint statement on the bottom group. LIFT31 section2 also identifies its bottom Q compact as K\_r, while its middle and top Q compacts are character-bearing parahorics.

Thus, with r=a+1, at every Q place the actual pair has

\[
 (J_{H,v},\chi_{H,v})=(K_{a,v},1),\qquad
 (J_{B,v},\chi_{B,v})=(K_{a+1,v},1),\qquad \tau_v=1 .
 \tag{QF1}
\]

The last identity is LIFT29(BT4--BT8): the prescribed nonintegral translation is supported at T only. There is no substitution of a different Q vector in (QF1).

The tuple character xi does enter datYt and datXt in LinksDefs427--434. It changes the higher vertical parahorics and the actual raw Q functions in DatumVec. That does not change dat0. Consequently (QF1) is compatible with the nontrivial scalar and Gauss/Euler factors in LIFT31(VA23--VA24a), DO041 and the vertical Rallis computation; none of those factors is asserted to be1 here.

\subsection{2. Why every remaining maximal stabilizer is hyperspecial}\label{why-every-remaining-maximal-stabilizer-is-hyperspecial}

Use the actual J\^{}14 or J\^{}56 host. The relative CM extension E/F is unramified at every finite place: the original E0/F0 assertion is the accepted DP57, carried by DO022; K029/030 preserves it in the real degree-seven extension, and K032(QH13) proves it after the two further real quadratic extensions by finite étale base change. Therefore every finite completion is either split or an unramified quadratic field extension. There is no relatively ramified finite place omitted from this list.

The two endpoint Hermitian spaces are the actual standard ones,

\[
             H:\ h_H=I_a,\qquad
             B:\ h_B=I_a\oplus(-1).
 \tag{QF2}
\]

Their standard lattices are O\_E\^{}a and O\_E\^{}(a+1). Both matrices are unimodular, so the lattices are selfdual for these E/F Hermitian forms, and the inclusion is the standard orthogonal direct summand.

Here is a direct integral-model justification of hyperspecialness. Let A=O\_(F\_v), and let B be the finite étale quadratic A-algebra obtained from E. For either invertible Hermitian matrix H in (QF2) the group functor

\[
 \mathscr U_H(R)
 =\{g\in\operatorname{GL}_j(B\otimes_A R):
                         \bar g^{\,t}H g=H\}
 \tag{QF3}
\]

is an A-form of GL\_j. After the finite étale faithfully flat base change A-\textgreater B, the quadratic algebra splits and the equation solves one of the two matrix components uniquely in terms of the other. Thus its base change is GL\_j over B. Smoothness, affineness and geometrically connected reductive fibres descend under this faithfully flat base change. Hence (QF3) is a reductive A-model. Its A-points are exactly the unitary stabilizer of the displayed lattice. By the definition of a hyperspecial subgroup these stabilizers are hyperspecial, with compatible integral embedding.

This argument concerns relative E/F forms and their actual lattices. It does not assume that F\_v is unramified over Q\_p. In particular it covers primes of the absolute discriminant, including those over7,29 and, for J\^{}56,5 and13. The trace-form coefficient used to realize the CM polarization can be nonunit at these rational primes; it does not change the unimodular E/F Hermitian matrices (QF2) or their lattice stabilizers.

All dyadic places are relatively split on the original DO022 host, and remain split after the real extensions. Thus they also admit the direct GL\_j integral description. The preceding étale-model argument in any case avoids an unproved odd-residue assumption about the group scheme. The relatively nonsplit places have odd residue characteristic in this actual family.

The unchanged R017 family keeps maximal level at29. No principal29 refinement from the separate finite-jet interpolation arguments is introduced here.

\subsection{3. The actual local factors}\label{the-actual-local-factors}

Let pi\_H,v and pi\_B,v be any actual irreducible constituents in the initial and genuine bottom spaces. At v outside G and T, (QF1) and its identical outside-Q counterpart give nonzero fixed vectors under the compatible hyperspecial compacts. Their spherical lines are one-dimensional. Therefore the actual local vectors, when made Hilbert-unit, are precisely the spherical unit vectors up to phases. The phases disappear from the matrix-coefficient period.

DO046 proves temperedness of these actual constituents. The unramified local period formula in Zhang(2014), introduction, printed543, (1.3), property(2), with footnote3, gives

\[
 \alpha_v(x_v,y_v)
       =\mathcal L_v(1/2,\pi_H\boxtimes\pi_B)
                      \operatorname{vol}(K_{H,v}),\qquad
 \alpha_v^\natural(x_v,y_v)=\operatorname{vol}(K_{H,v}).
 \tag{QF4}
\]

For the actual local convention vol(K\_H,v)=1, this is

\[
                  \boxed{\alpha_v^\natural=1
                    \quad\text{for every finite }v\notin G\cup T.}
 \tag{QF5}
\]

The hypotheses are now checked: relative unramifiedness, compatible hyperspecial groups, actual spherical vectors and temperedness. This is not inferred merely from absence of a character in a displayed theta kernel.

At a relative split place (QF4) can also be verified by the exact raw Whittaker normalization in K054 section4: spherical norms kappa\_(a+1)L\_G and kappa\_a L\_H, raw central integral kappa\_a L0, and outer measure kappa\_a produce alpha=Delta\_(a+1)\textbar L0\textbar\^{}2/(L\_GL\_H). This applies in particular to Q and the dyadic places. At a nonsplit unramified place use the unitary unramified formula itself; the split zeta product is not substituted for its actual unitary local normalization.

The complete finite inventory is consequently:

\begin{longtable}[]{@{}
  >{\raggedright\arraybackslash}p{(\columnwidth - 6\tabcolsep) * \real{0.2500}}
  >{\raggedright\arraybackslash}p{(\columnwidth - 6\tabcolsep) * \real{0.2500}}
  >{\raggedright\arraybackslash}p{(\columnwidth - 6\tabcolsep) * \real{0.2500}}
  >{\raggedright\arraybackslash}p{(\columnwidth - 6\tabcolsep) * \real{0.2500}}@{}}
\toprule\noalign{}
\begin{minipage}[b]{\linewidth}\raggedright
Finite place
\end{minipage} & \begin{minipage}[b]{\linewidth}\raggedright
Actual initial vector
\end{minipage} & \begin{minipage}[b]{\linewidth}\raggedright
Actual bottom vector and translate
\end{minipage} & \begin{minipage}[b]{\linewidth}\raggedright
Canonical local result
\end{minipage} \\
\midrule\noalign{}
\endhead
\bottomrule\noalign{}
\endlastfoot
G, including the two mixed places & K\_a\^{}1-fixed inverse generic vector & The K055 vector in the same finite Pi, positive generic; tau=1 & 1/(dim(sigma\_v) mathcal L\_v), K055(MK26) \\
T & Actual strict-level newform in its nontrivial dilation component & Spherical unit vector translated by the prescribed lower tau & q/((q-1)(q\^{}a-1)), K054(CP3) \\
Auxiliary Q & Maximal-lattice fixed, trivial character & Maximal-lattice fixed, trivial character; tau=1 & 1 \\
Other relatively split places & Hyperspecial fixed & Compatible hyperspecial fixed; tau=1 & 1 \\
Other relatively nonsplit places & Hyperspecial fixed & Compatible hyperspecial fixed; tau=1 & 1 \\
Relatively ramified places & None on either actual host & None & No missing factor \\
\end{longtable}

For decomposable unit vectors chosen at G as in K055, the full canonical finite product is exactly

\[
 \prod_{v<\infty}\alpha_v^\natural
   =
   \left(\frac{T}{(T-1)(T^a-1)}\right)^f
       \prod_{v\mid G}
          \frac1{\dim\sigma_v\,\mathcal L_v(1/2)}.
 \tag{QF6}
\]

Here the actual T primes are degree one and completely split in the real host, as checked in DO040, so there are f of them and their common residue cardinality is T. The factors use the full unitary adjoint convention in K054/K055. The separate conversion from these local probability Haar measures to the global measure is not absorbed into (QF6).

There is no auxiliary-Q factor in this GGP product. There are still the actual auxiliary-Q costs in the horizontal and vertical theta maps; (QF6) does not cancel them from a source MainX cost. It also does not replace the central global L-value, the archimedean selected-vector calculation, or the actual arithmetic image and integral-preimage questions.

\subsection{4. Sources and check scope}\label{sources-and-check-scope-2}

\begin{itemize}
\tightlist
\item
  Zhang, W. (2014). Automorphic period and the central value of Rankin--Selberg L-function. Journal of the American Mathematical Society,27(2),541--612. \url{https://doi.org/10.1090/S0894-0347-2014-00784-0}. Direct original publisher offprint, printed542--543, (1.1), (1.3), (1.4) and footnote3; the measure normalization audit is K054.
\item
  Recovered Lean source at the exact locators in section1. The relevant definitions were directly reread for this finding. No Lean build was performed and no literal Fields/Host constructor obstruction was bypassed.
\item
  Accepted DP57 and checked DO022/K029/K030/K032(QH13): the actual relative finite-unramified CM hosts and lattices. The short étale-model proof in section2 verifies the needed endpoint hyperspecial property, rather than presuming it from absolute discriminant data.
\item
  K038(CA27a), LIFT29(BT3--BT8), LIFT31 section2: the actual endpoint levels, characters and T-only translation. DO046 supplies actual all-place temperedness.
\item
  K054 and K055 supply the two nontrivial rows of the finite factor inventory. This source audit proves the remaining rows.
\end{itemize}

The asserted result is the stated endpoint GGP test. No raw vertical Q function is replaced by a spherical vector, and no unrestricted local test vector is substituted into an earlier theta image.

\subsection{5. Author check record}\label{author-check-record-2}

The complete saved proof QF1--QF6 was read and reconstructed. The displayed Lean definitions were directly reread, including the actual levFn locator700--703, together with K038(CA27a), LIFT29(BT3--BT8), LIFT31 sections2 and6, K030's finite-place statement, K032(QH13), and DO022 section4. The original Zhang offprint printed542--543, including footnote3, was directly read. The integral-model descent proves hyperspecialness independently of absolute ramification. This is a whole mathematical/source audit, not a numerical certificate or a Lean build. It found no omitted relatively ramified finite place and no endpoint Q character. The exact UTC event and frozen SHA256 are recorded in the author's check ledger.

\clearpage

\section{AM-K057 --- the exact selected archimedean GGP factor}\label{am-k057-the-exact-selected-archimedean-ggp-factor}

Author: GPT-6 Astra Pro, ALP\_MACHINE /root/kummer\_extension. 22 November 2026 UTC. Inherited reasoning effort is not separately exposed.

\textbf{Author whole-proof check PASS; exact certificate PASS.} Root and independent review events are recorded separately.

This computes the local GGP normalization for the actual compact initial type and the actual contragredient bottom discrete series. The selected compact inclusion is the isometry of K043. The raw compact integral is exactly the reciprocal compact dimension for matched unit vectors. The canonical factor is that value divided by an explicit positive gamma product, with a rational factorial expression in the actual source parameters. There is no unspecified real-period scalar in this local calculation. An arbitrary global automorphic realization, the genuine same-Schwartz seesaw, and a nonzero central global value remain separate.

\subsection{1. Actual Harish--Chandra parameters and the dual convention}\label{actual-harishchandra-parameters-and-the-dual-convention}

Fix one selected CM embedding of either actual host, and write a=42 or168. Put n=a+1 and k=a/2. The real pair is

\[
 H=U(a),\qquad G=U(a,1),\qquad
 h\longmapsto\operatorname{diag}(h,1),\qquad \operatorname{vol}(H)=1.
 \tag{AF1}
\]

Use zero-based indices and the actual source lists

\[
 u_0>p_0>u_1>p_1>\cdots>u_{a-1}>p_{a-1}>p_a,
 \qquad p=(v_0,\ldots,v_{a-2},-b_0,-b_1).
 \tag{AF2}
\]

The u\_i are in Z+1/2 and the p\_j are in Z. These are precisely DataTheta.lean529--548, TPar.Interlace, Dpar(0), and Xpar(0): for m=2 the positive bottom slots are v and −b0, the negative slot is −b1, and the eta shift at j=0 is zero. The files are under work/research/alpoge\_source/Formalism/LocalRingCardinality, recovered commit7cbf7a0c32fca1badf05931879a37c3f8a2b03b4.

The initial compact representation W\_lambda has highest weight

\[
 \lambda_i=u_i-\frac{a-1}{2}+i,
 \qquad
 d_\lambda=\dim W_\lambda
       =\prod_{0\le i<j<a}\frac{u_i-u_j}{j-i}.
 \tag{AF3}
\]

The highest-weight formula is literal hcToWeightQ in DataCoeff.lean41 and lamTrue\_zero in \nolinkurl{LocalToGlobal/Columns/BaseLemmas.lean491–494.} The dimension formula follows from the Weyl dimension formula, since lambda\_i−lambda\_j+j−i=u\_i−u\_j.

Let sigma be LIFT28's discrete series with signed parameter (p\_0,\ldots,p\_\{a-1\};p\_a). Its minimal compact type is

\[
 W_\alpha\boxtimes z^\beta,
 \quad \alpha_i=p_i-\frac{a-2}{2}+i,
 \quad \beta=p_a-\frac a2.
 \tag{AF4}
\]

The actual right-action bottom representation is sigma\^{}vee, exactly as R027 section1 and LIFT28(HJ20) prescribe. Thus the actual GGP pair is W\_lambda external-tensor sigma\^{}vee. Its standard base-change parameters at this complex embedding are

\[
 \phi_H|_{W_{\mathbb C}}=\bigoplus_{i<a}\chi_{u_i},\qquad
 \phi_B|_{W_{\mathbb C}}=\bigoplus_{j\le a}\chi_{-p_j},
 \qquad \chi_x(z)=(z/|z|)^{2x}.
 \tag{AF5}
\]

The expression (z/bar z)\^{}x is only shorthand for this single-valued character; 2x is an integer. In particular chi\_x(−1)=(−1)\^{}(2x). This is the standard stable base change with no auxiliary character shift: in Mok(2015), section7.1, (7.1.2)--(7.1.3), printed190, take kappa=+1 and the base-change character chi=1. The character exponent is then the actual Harish--Chandra coordinate. The construction's eta characters do not add a further half shift to the standard parameters of the resulting representations. The dual negates the p coordinates; sorting them in reverse order does not affect any product below.

The two characters in each Rankin factor therefore have exponent u\_i−p\_j. Replacing this difference by a sum would correspond to the wrong bottom representation. Simultaneously conjugating both members of the actual bilinear pair negates every coordinate and leaves all following factors unchanged.

\subsection{2. The actual selected vector and the raw compact period}\label{the-actual-selected-vector-and-the-raw-compact-period}

The horizontal-strip calculation of LIFT28(HJ8--HJ9) gives the unique selected degree

\[
 m_i=u_i-p_i-\tfrac12\in\mathbb Z_{\ge0},\qquad
 N=\sum_{i<a}m_i=\sum_{i<a}u_i-\sum_{j<a}p_j-\frac a2.
 \tag{AF6}
\]

K043 supplies the H-isometry J:W\_lambda→sigma on this selected type. Its dual conjugate isometry is denoted by bar J:overline(W\_lambda)→sigma\^{}vee. For unit vectors v,w in W\_lambda, put x=bar J(bar w). Compact Schur orthogonality gives the full local matrix-coefficient integral

\[
 \begin{split}
 \alpha_\infty(v,x)
 &=\int_H\langle\sigma_H(h)v,v\rangle
                   \langle\sigma^\vee(h)x,x\rangle\,dh\\
 &=\frac{|\langle v,w\rangle|^2}{d_\lambda}.
 \end{split}
 \tag{AF7}
\]

Indeed the second coefficient is the complex conjugate of the coefficient of w in W\_lambda. In column-vector notation, averaging the rank-one operator v w* under the irreducible compact representation gives (w* v/d\_lambda)I: it is scalar by Schur's lemma and has trace w* v. Taking the matrix element v* (averaged operator) w proves (AF7), with no infinite-dimensional formal degree: the integral is over the compact group H and the selected finite type. For the actual matched choice w=v,

\[
                 \boxed{\alpha_\infty=1/d_\lambda.}
 \tag{AF8}
\]

The statement is not that any two unit vectors have this value. Orthogonal v,w give zero. It also does not claim that an arbitrary unselected bottom vector has a nonzero selected component.

For clarity, the PBW integer of K043 is

\[
 D_{\rm PBW}=\prod_{i<a}(p_i-p_a+1)_{m_i}.
 \tag{AF9}
\]

If j is the Fischer-isometric Pieri inclusion, then J=D\_PBW\^{}(-1/2) Cj, where C is the matrix-unit creation map. Thus (AF8) has already included the PBW normalization. For an unnormalized selected polynomial P whose Fischer norm squared is c\_P, the vector C(P) has norm squared D\_PBW c\_P. Matched against an initial unit vector, its raw value is D\_PBW c\_P/d\_lambda; dividing by its Hilbert norm squared recovers (AF8). Source hook norms, creation-frame constants, and scalars of a non-isometric global realization are retained when such unnormalized vectors are used. They cannot be inserted a second time into the unit-vector value (AF8).

\subsection{3. The full unitary adjoint, including its centre}\label{the-full-unitary-adjoint-including-its-centre}

Set

\[
 \Gamma_{\mathbb R}(s)=\pi^{-s/2}\Gamma(s/2),\qquad
 \Gamma_{\mathbb C}(s)=2(2\pi)^{-s}\Gamma(s).
 \tag{AF10}
\]

Knapp(1994), printed404, (3.6), and printed406, (4.6), give L(s,sgn)=Gamma\_R(s+1), L(s,chi\_x)=Gamma\_C(s+\textbar x\textbar), and L(s,Ind\_(W\_C)\^{}(W\_R) chi\_x)=Gamma\_C(s+\textbar x\textbar) for x nonzero. These are the usual archimedean local factors, not gamma quotients defined only up to a scalar.

Here is a direct computation of the adjoint in the needed unitary convention. For a rank-m discrete parameter let x\_1,\ldots,x\_m be its distinct coordinates in Z+(m−1)/2. Mok section2.1, (2.1.1), printed6--7, gives

\[
 {}^LU(m)=GL_m(\mathbb C)\rtimes W_{\mathbb R},\qquad
 j:g\longmapsto\Phi_m\,{}^tg^{-1}\Phi_m^{-1},\qquad
 \Phi_m^t=(-1)^{m-1}\Phi_m,\quad\Phi_m^2=(-1)^{m-1}I.
 \tag{AF11}
\]

One representative of the parameter has phi(z)=diag(chi\_(x\_i)(z)) on W\_C and phi(j)=(Phi\_m\^{}(-1),j). It satisfies jzj\^{}(-1)=bar z because the two Phi matrices cancel, and the matrix-component equation is Phi\_m\^{}(-1) theta(Phi\_m\^{}(-1))=(-1)\^{}(m−1)I=diag(chi\_(x\_i)(−1)), using the coordinate grid. Thus the full parameter satisfies phi(j)\^{}2=phi(−1). On the full adjoint Lie algebra its j action is consequently

\[
                             X\longmapsto-{}^tX.
 \tag{AF12}
\]

Each diagonal line has the sign character. Each pair of off-diagonal lines E\_ij,E\_ji is the induced two-dimensional representation with W\_C characters chi\_(x\_i−x\_j) and its inverse. Twisting that two-dimensional representation by sign gives an isomorphic representation, so the displayed minus signs do not alter its L-factor. We obtain

\[
 L_\infty(s,\pi_x,\operatorname{Ad}_U)
  =\Gamma_{\mathbb R}(s+1)^m
        \prod_{i<j}\Gamma_{\mathbb C}(s+|x_i-x_j|).
 \tag{AF13}
\]

This is the full unitary adjoint convention of Zhang(2014), Remark7, printed560, and DO046(AP17). The one-dimensional centre has the sign action; it has not been removed as for a trace-zero GL adjoint. Since that centre has no W\_R invariants, quotienting only the invariant dual-centre Lie algebra makes no change here.

\subsection{4. Exact gamma and factorial expressions}\label{exact-gamma-and-factorial-expressions}

At a real place the quadratic character of E/F is sign. Define

\[
 \Delta_{a+1,\infty}
    =\prod_{r=1}^{a+1}\Gamma_{\mathbb R}
                         (r+(r\bmod2)).
 \tag{AF14}
\]

Combining (AF5), (AF10), (AF13), and Zhang(1.1), or equivalently DO046(AP17) at s=1/2, yields the completely specified positive number

\[
 \boxed{\mathcal L_\infty(\tfrac12)
  =\frac{\displaystyle
       \Delta_{a+1,\infty}
       \prod_{i<a,\ j\le a}\Gamma_{\mathbb C}
                              (\tfrac12+|u_i-p_j|)}
     {\displaystyle
       \Gamma_{\mathbb R}(2)^{2a+1}
       \prod_{i<j<a}\Gamma_{\mathbb C}(1+u_i-u_j)
       \prod_{i<j\le a}\Gamma_{\mathbb C}(1+p_i-p_j)}}.
 \tag{AF15}
\]

All arguments are positive integers. There is no pole, zero, phase, or unspecified choice of square root in this expression. The canonical local factor for matched unit vectors is therefore

\[
          \boxed{\alpha_\infty^\natural
                    =\frac1{d_\lambda\mathcal L_\infty(1/2)}.}
 \tag{AF16}
\]

For an exact rational form put

\[
 h_{ij}=|u_i-p_j|-\tfrac12\in\mathbb Z_{\ge0},\qquad
 A_k=k!\prod_{r=0}^{k-1}(r!)^2,
\] \[
 R(u,p)=
   \frac{\displaystyle\prod_{i<a,\ j\le a}h_{ij}!}
        {\displaystyle\prod_{i<j<a}(u_i-u_j)!
                           \prod_{i<j\le a}(p_i-p_j)!}.
 \tag{AF17}
\]

Then (AF15) is exactly

\[
 \boxed{\mathcal L_\infty(\tfrac12)
     =A_k R(u,p)\,
           2^{a(a-1)/2-2N}\,
           \pi^{k(k-1)-2N}.}
 \tag{AF18}
\]

\textbf{Proof.} Pair the factors r=2s−1 and r=2s in (AF14), leaving the last r=2k+1. Since Gamma\_R(2s)=(s−1)! pi\^{}(−s), this gives Delta\_(a+1,infty)=A\_k pi\^{}(−(k+1)\^{}2). The numerator in (AF15) has a(a+1) complex gamma factors; its two adjoint products have a\^{}2. The factorials remaining after Gamma\_C(m)=2(2pi)\^{}(−m)(m−1)! are exactly (AF17).

To check the powers, the ordered interlacing (AF2) gives

\[
 \sum_{i,j}|u_i-p_j|
   -\sum_{i<j}(u_i-u_j)-\sum_{i<j}(p_i-p_j)
       =2\left(\sum_{i<a}u_i-\sum_{j<a}p_j\right)=2N+a.
 \tag{AF19}
\]

For each u\_i the coefficient on the left is (a+1−2i)−(a−1−2i)=2. For p\_j with j\textless a it is (a−2j−2)−(a−2j)=−2; for p\_a it is −a−(−a)=0. This proves (AF19) without a numerical approximation. After including the constant arguments, the difference between the sums of the complex gamma arguments is 2N+a(3−a)/2. Finally Gamma\_R(2)=pi\^{}(−1). Substitution gives the exponents in (AF18). QED.

These formulae are invariant under a common translation of all u and all p, as they must be for mutually cancelling determinant twists of the actual pair. The selected N is also unchanged.

\subsection{5. Full local product and the exact remaining scope}\label{full-local-product-and-the-exact-remaining-scope}

At all f labels take the actual matched isometric selected vectors. K056 then gives the complete canonical local product

\[
 \prod_v\alpha_v^\natural
 =\left(\prod_{\tau\mid\infty}
          \frac1{d_{\lambda_\tau}\mathcal L_\tau(1/2)}\right)
   \left(\frac{T}{(T-1)(T^a-1)}\right)^f
       \prod_{v\mid G}\frac1{\dim\sigma_v\,\mathcal L_v(1/2)}.
 \tag{AF20}
\]

This uses probability Haar on each compact H\_infty and volume one for the finite standard maximal H compacts. The global Haar conversion in R030/DO046 remains explicit. There is no missing finite Q factor, but the separate Q theta/Rallis factors remain.

In the completed GGP expression the infinity normalizer cancels label by label:

\[
 \left(\prod_\tau\mathcal L_\tau(1/2)\right)
       \left(\prod_\tau\alpha_\tau^\natural\right)
                    =\prod_\tau d_{\lambda_\tau}^{-1}.
 \tag{AF21}
\]

Thus neither the explicit gamma factor nor the PBW square root may be charged twice when moving between a complete global L-value and the Hilbert-unit local tests. This cancellation does not give a lower bound for the finite central L-value. It also gives no nonzero assertion for a prescribed nonfactorizable arithmetic row: polarization can still have cancellations, and the source integral preimage/detector lattice is a different input.

There is nevertheless an exact nonvanishing reduction for chosen factorizable vectors. Fix an actual initial Mram constituent and an existing genuine bottom constituent carrying an admitted generic finite Pi to which K055 applies. At G choose K055's vector in that same Pi, at T use the actual spherical vector and prescribed lower translation of K054, elsewhere use the actual spherical vectors of K056, and at infinity use the matched selected isometry above. Every canonical local factor is then positive. DO046(AP19), the proved global GGP formula with its actual parameter hypotheses, therefore gives

\[
 \boxed{\mathcal P_\tau(f,F)\ne0
       \quad\Longleftrightarrow\quad
       L(\tfrac12,\Pi_H\times\Pi_B)\ne0.}
 \tag{AF22}
\]

The global adjoint denominator and the other normalizing factors are the nonzero factors in that theorem. The chosen local vectors remain in the stated finite levels and genuine projector; no new ramification or change of inner form occurs. This is not a central value nonvanishing theorem, and it does not place the chosen global pair inside a prescribed earlier arithmetic image. Those two questions are not consequences of the positivity of local tests.

The new content here is the exact actual-dual parameter audit, the full unitary-adjoint computation, and their selected-vector normalization and reduction (AF15--AF22). The compact Schur identity, Weyl dimension formula and archimedean local-factor definitions are standard primary inputs, not newly claimed theorems.

\subsection{6. Sources and verification}\label{sources-and-verification-3}

\begin{itemize}
\tightlist
\item
  Knapp, A. W. (1994). Local Langlands correspondence: The archimedean case. In Motives, Proceedings of Symposia in Pure Mathematics55, Part2 (pp.393--410). American Mathematical Society. \url{https://doi.org/10.1090/pspum/055.2/1265560}. Original author reprint: \url{https://www.math.stonybrook.edu/~aknapp/pdf-files/motives.pdf}. Printed404 (3.6), printed406 (4.6), read from the original page images. Local PDF: archimedean\_period\_refs/knapp1994.pdf, SHA256684de4bcfc50e448fe52fddc863392012581b43097f40f8bcc4c83dbc5a2dfbb.
\item
  Mok, C. P. (2015). Endoscopic classification of representations of quasi-split unitary groups. Memoirs of the American Mathematical Society235(1108). \url{https://doi.org/10.1090/memo/1108}. Retained author version arXiv:1206.0882: section2.1,(2.1.1), printed6--7; section7.1,(7.1.2)--(7.1.3),printed190.
\item
  Zhang, W. (2014). Automorphic period and the central value of Rankin--Selberg L-function. Journal of the American Mathematical Society27(2),541--612. \url{https://doi.org/10.1090/S0894-0347-2014-00784-0}. Printed542--543,(1.1),(1.3),(1.4); Remark7,printed560.
\item
  Literal DataTheta/DataCoeff/Columns locators in section1; LIFT28(HJ4--HJ9,HJ20), R027 section1, K043(PB5--PB8), and DO046(AP17--AP19). K056 supplies the finite inventory.
\end{itemize}

The companion certificate checks finite exact matrix identities, the coordinate/grid/interlacing algebra, and the rational gamma factorial reduction in small ranks and at actual-profile argument counts. Its final replay at2026-11-22T06:49:50.500426+00:00 passed 1,247 named exact checks, including16 compact-Schur vector pairs, and separately verified285 matrix units,240 small gamma reductions and108 actual-profile argument cases. The all-coordinate proof (AF19) covers every word; the108 actual cases are explicitly samples of the parameter values, not an enumeration of all words. It does not numerically compute enormous factorials in the actual profiles, an automorphic L-value, or a global period.

The complete saved proof AF1--AF22 and full certificate were read and reconstructed by the author. Direct primary checks include Mok's L-group and archimedean exponent definitions, the original Knapp page images for both real and complex factors, and Zhang's full unitary adjoint and canonical normalizer. The literal source HC lists, R027's bottom contragredient, LIFT28's selected dual type and K043's isometric PBW convention were checked together. The single-valued character and full semidirect-product j-square conventions are explicit. \textbf{PASS in the stated matched-vector and chosen-factorizable-pair scope.} The author ledger records the exact UTC completion and frozen report SHA256. No Lean build or global central-value computation was performed.

\clearpage

\section{AM-DO-051 --- the actual T frames, finite Fourier operator and fixed last-two tensor}\label{am-do-051-the-actual-t-frames-finite-fourier-operator-and-fixed-last-two-tensor}

Author: GPT-6 Astra Pro, ALP\_MACHINE /root/definite\_oscillator. Inherited reasoning effort is not separately exposed. 22 November 2026 UTC.

\textbf{Author whole-proof PASS --- 2026-11-22T06:54:10.814862+00:00.}

This report computes the T part of the literal final-two-step seesaw in the constructed mixed family. It keeps the actual primitive negative line of LIFT30, the positive terminal line of the source, the original basic vertical vacua, and the prescribed two horizontal incidence functions. The T input operator and period scalar are explicit. A finite Heisenberg operator produces one common tensor before any finite tuple is varied. The calculation does not replace the remaining fixed G/Q laws or assert a global restricted period.

The result is independent of the elliptic vacuum constructor: the proof uses the actual mixed Weil representation, its unit lattice vacua, and the standard split Schrödinger model already constructed in DO022/K035/LIFT39. All T places on both actual hosts are absolutely split, odd, and disjoint from the different and the fixed character conductors. Those arithmetic facts are the checked DO040 inputs.

\subsection{1. Actual rational line and the exact glued lattice}\label{actual-rational-line-and-the-exact-glued-lattice}

Use the actual ranks a,b,N,c=N+1 and the actual prime T of DO040. Let pi=pi\_T, pi bar(pi)=T, and let xi0 be the chosen trace-zero coefficient of the genuine mixed oscillator. It is a unit at T. Write c\_v=bar(pi) at a selected T place; pi is a uniformizer and c\_v is a unit. The two-plane special-unitary matrix of LIFT30 is

\[
 g=\begin{pmatrix}A/\pi&B/\bar\pi\\B/\pi&A/\bar\pi\end{pmatrix},
 \qquad A=(T+1)/2,\quad B=(T-1)/2.
 \tag{FT1}
\]

It is the identity on the other a−1 positive coordinates. Retain the primitive negative vector z\_T=(B,A)=g(bar(pi)e\_last) and the positive complement isometry g\textbar E\^{}a. Thus the intrinsic negative line has Gram −1, its distinguished vector has coordinate zeta=bar(pi), and the physical embedding is g. The terminal line has Gram +1 and distinguished vector zeta0=pi e\_last, exactly StdLatUnimodular.z0/Deco.til.

Define the rational line map

\[
          \tau=-\xi_0\,{\bar\pi\over\pi}.
 \tag{FT2}
\]

Then tau\emph{zeta0=−xi0}zeta, so the source lambda is −xi0, a T-unit, and tau* (−1) tau=−xi0 bar(xi0). These are both literal conditions of CombAtSingleZero.PinOff.tau, lines161--162. The multiplicative valuation equals one; no additive-valuation-one condition is used. At infinity kap(w)=(pi/bar(pi))w*, with the opposite complex structure on the negative line. Its unit phase remains in DO050.

Fix any common local uniformizer varpi, put rho=varpi/pi, and take the sum frame

\[
                 \alpha_0=g\,\operatorname{diag}(I_a,\bar\pi).
 \tag{FT3}
\]

The actual standard bottom lattice has the following coordinates:

\[
 \alpha_0^{-1}O^{a+1}
 =\{(x,z):x\in O^a,\ \pi z\in O,\ x_1\equiv\pi z\pmod\pi\}
 =\operatorname{gluedS}(\varpi,1,\ker P,v_P),
 \tag{FT4}
\]

where P has lattice O\^{}a, functional f\_P(x)=red(rho)x\_1 and vector v\_P=rho\^{}-1 e\_1. Its kernel is x\_1=0 modulo pi.

\textbf{Proof.} LIFT30(NL8) gives g\^{}-1={[}{[}I,0{]},{[}pi\^{}-1 bar(pi)e\_1*,1{]}{]}k with k integral invertible. Multiplying by diag(I,bar(pi)\^{}-1) proves the first equality. The source definition of gluedS is the O-span of (ker P,0) and (v\_P,varpi\^{}-1). A vector in this span has varpi z=s in O and x\_1=rho\^{}-1 s modulo pi, equivalently x\_1=pi z modulo pi. This is exactly the displayed lattice. Conversely those conditions give the same s and a vector of ker P. QED.

Thus the bottom vacuum in FT4 is the unchanged standard physical vacuum, expressed in its actual Bessel sum frame. It is not a new freely translated Schwartz function. The arbitrary actual primitive covector of DO040 is reached by LIFT30(NL10--12)'s integral finite coordinate transport, simultaneously on the initial row, the sum frame, and the finite character. No principal T refinement or projective averaging is introduced.

\subsection{2. Exact positive terminal pair and diamond phase}\label{exact-positive-terminal-pair-and-diamond-phase}

The source terminal decorated lattice is the old terminal lattice plus O e\_last; its functional is red(bar(pi)x\_last). In the line frame gamma0=diag(I\_N,pi), its kernel is O\^{}N plus O e\_last. The vector varpi\^{}-1 e\_last has value

\[
                      s=\operatorname{red}(T/\varpi).
 \tag{FT5}
\]

Indeed its physical coordinate is pi/varpi and multiplication by bar(pi) gives T/varpi. The source lastPair has value one on this vector. Therefore the exact LastKernelDatumInLastLineFrame comparison uses u=s\^{}-1, not an arbitrary residue scalar.

The source line Fourier theorem KapStep.kapStep explicitly chooses epsilon'=red(lambda)\^{}-1. Combined with FT2 and FT5, the actual initial diamond in the last-two law is

\[
 \boxed{\quad
 \epsilon=(\epsilon')^{-1}u
       =\operatorname{red}(-\xi_0\varpi/T).
 \quad}                                                    \tag{FT6}
\]

In the retained rational uniformizer varpi=T, epsilon=−red(xi0) and the terminal comparison s is exactly one. All formulas below allow any varpi; its unit change must also change the residue additive character and FT6. Whenever this residue unit occurs in a linear expression over F, use the exact fixed unit lift epsilon\_tilde=−xi0*varpi/T. Fix O-linear lifts of f\_P and f\_Q as well. Thus equalities involving pi, epsilon\_tilde and varpi below are equalities in F, not an identification of a residue class with an unspecified lift.

\subsection{3. The literal finite Fourier operator and its metric}\label{the-literal-finite-fourier-operator-and-its-metric}

Let chi have conductor O and let chi\_bar(x)=chi(varpi\^{}-1 x) on the residue field k, of cardinality q=T. For a lattice Lambda of rank a, the source FFStar block is exactly

\[
 (\mathcal F^-f)(\ell)
       =\sum_{x\in\Lambda/\varpi\Lambda,\ x\ne0}
                      \bar\chi(-\ell(x))f(x),\qquad \ell\ne0.
 \tag{FT7}
\]

This is HasConductorMulShift.FFStar\_ofDual, lines136--138. The domain labels are the pairs up(Lambda,varpi,x), whose kernel is Lambda; the output labels are ofDual(Lambda,ell), whose lattice is Lambda. The source ramified condition is the unweighted sum zero on every scalar orbit, not a character-weighted average.

With counting metrics on these two finite fibres,

\[
 (\mathcal F^-)^*\mathcal F^-=q^a I-J,
 \qquad \|\mathcal F^-f\|^2=q^a\|f\|^2\quad(f\in M_{\rm ram}),
 \tag{FT8}
\]

where J has every entry one. The inverse on M\_ram is q\^{}-a times the positive-sign Fourier block. Multiplication by the source diamond epsilon permutes the residue labels, so it changes no norm.

\textbf{Proof.} The (x,y) Gram entry is the sum over nonzero ell of chi\_bar(ell(y−x)). It is q\^{}a−1 when x=y and −1 otherwise, by the complete character sum including ell=0. On M\_ram the sum over all nonzero x is zero, since the nonzero vectors partition into scalar orbits. This kills J and proves the norm and inverse. Moreover the scalar-orbit sum of the output at ell is sum\_x f(x) sum\_\{c!=0\}chi\_bar(−c ell(x)). The inner sum is q−1 or −1 according as ell(x)=0 or not. Each of those x sets is a union of scalar orbits, so that output sum is zero too. QED.

Equation FT8 is an exact finite-fibre metric assertion. It does not silently replace the actual global class/stabilizer measure by counting measure. The adelic input operator is literally PlacesToAdeles.sOp=eval\_P(dia(epsilon) FFStar formAt\_P); it is the operator used in the identity below. The source's equivariant form realization transports the finite Fourier identity, while any numerical global class norm still uses its own stated masses.

\subsection{4. A proof of the exact source constant}\label{a-proof-of-the-exact-source-constant}

For clarity write Lambda\_B=O\^{}b, P=(Lambda\_A,f\_P), and choose v\_P with f\_P(v\_P)=1. In the sum coordinates (X,y), the actual glued first vacuum has the indicator

\[
 G_P(X,y)=1_{X\ker P\subset\Lambda_B}
                  1_{y+\varpi Xv_P\in\varpi\Lambda_B}.
 \tag{FT9}
\]

Let theta(v,Phi) be the first theta integration over the compact initial unitary group, with all other variables held fixed. Its exact source transformation law for g in GL\_a(F) is

\[
 \theta(\pi_A(g)v,\Phi(Xg,y))
      =l_S(\det g)^{-1}\theta(v,\Phi),
 \qquad
 l_S(\varpi)=q^{-b/2}\chi_W(\operatorname{idele}_v\varpi)^{-1}.
 \tag{FT10}
\]

This is the actual split Weil formula: the determinant modulus is b/2 and the source cross character is the inverse of chi\_W. It is also the explicitly proved exists\_theta\_exact in the recovered GluedPlaceSourceConstantProof, lines147--200. The linear model is based on the same vacuum throughout. No separate renormalization of G\_P and its replacement is permitted.

Finite Fourier expansion of the second indicator in FT9 gives q\^{}-b times a sum over all residue covectors f on Lambda\_B. The zero covector term is 1\_\{X ker P subset Lambda\_B\}1\_\{Lambda\_B\}(y). For f!=0, its y factor is charInd(Q\_f)(y)=chi\_bar(−f(y)); its X factor is

\[
 1_{X\ker P\subset\Lambda_B}\bar\chi(-f(\varpi Xv_P))
 =1_{X\ker P\subset\ker Q_f}
   +\sum_{\ell\ne0}
       \bar\chi(-\ell(\varpi v_P))
              \operatorname{phiOne}(P_\ell,Q_f)(X),
 \tag{FT11}
\]

where P\_ell has lattice ker P and nonzero functional ell. The identity follows by classifying the single functional f composed with X on ker P: it is either zero or exactly one nonzero ell. The signs agree with FT7 and with the source charInd definition.

If the initial input belongs to M\_ram, the zero terms vanish after theta integration. For example the X-lattice indicator is fixed by the stabilizer of ker P. Averaging the input over this compact group averages its nonzero residue-position labels, whose sum is zero by the ramified condition. The same argument applies to the X indicator into ker Q\_f.~All relevant unramified determinant characters are one on these compact stabilizers.

The remaining finite sum transposes to FFStar on the input. Here is a direct justification, so no normalization is hidden in that transposition. The compact GL(ker P/varpi ker P) acts transitively on nonzero covectors ell. Transport each ell to a fixed ell0 and use FT10 with determinant a unit. The residual stabilizer of ell0 makes the resulting theta value depend only on ell0(x), where x is the transported nonzero input position. For each c in k, the number of nonzero ell with ell(x0)=c equals the number of nonzero x with ell0(x)=c: both are q\^{}(a−1) if c!=0 and q\^{}(a−1)−1 if c=0. Thus the two weighted sums agree term by term in c.~Their common weight is chi\_bar(−c), giving precisely FT7 and no finite index.

The resulting pair initially has lattice ker P. Choose g taking it to P. Because {[}P.lat:ker P{]}=q, det(g)*varpi is a unit. Unramifiedness and FT10 therefore give

\[
 l_S(\det g)=q^{b/2}\chi_W(\operatorname{idele}_v\varpi).
 \tag{FT12}
\]

Combining the Fourier coefficient q\^{}-b with FT12 proves the source first-lift coefficient q\^{}-b/2 chi\_W(varpi). There is no change to the actual initial Haar measure.

The second vertical test is the unchanged basic vacuum and is fixed by the full middle compact at T. Moving its middle variable permutes all q\^{}b−1 nonzero covectors Q\_f of Lambda\_B, without a character factor. Haar invariance and compact Fubini therefore turn their sum into q\^{}b−1 copies of any fixed Q. Applying the diamond epsilon to that fixed choice gives the exact local last-two period constant

\[
 \boxed{\quad
 c_{T,v}=q^{-b/2}(q^b-1)
                \chi_W(\operatorname{idele}_v\varpi),
 \qquad \chi_W=\mu_0\xi .
 \quad}                                                    \tag{FT13}
\]

This reproduces the literal source: sourceExact chooses r0=q\^{}-b q\^{}(b/2), e=1; SmallPairPairingAtOneGluedPlace.smallOne multiplies it by the number of nonzero covectors. The proof just given explains every factor. In particular q\^{}b−1 is not replaced by a projective index, and the character is not inverted.

\subsection{5. The exact transformed raw T tensor}\label{the-exact-transformed-raw-t-tensor}

We now calculate the same tensor, not merely its period. Work in the fixed Schrödinger models with additive measure O\^{}b of volume one. In the actual frames FT3 and gamma0=diag(I\_N,pi), the unitary Fourier-and-dilation model of the line exchange is

\[
 (\mathcal S_v^0\Phi)(X,Z,t)
   =q^{b/2}\int_{F^b}
      \Phi_0(X,y)\,
      \chi(-\xi_0\,(t/\pi)\cdot y)\,dy\ \Phi_1(Z).
 \tag{FT14}
\]

The Fourier character has conductor O because xi0 is a unit. The factor q\^{}b/2 is the positive unitary dilation from the row coordinate t/pi to the physical row t. The negative-line frame has unit determinant at this place. The formula agrees with the source \nolinkurl{KapStep.fourier}, which uses \nolinkurl{hat psi_k nu (-1)} with psi\_k(x)=chi(xi0*x), and \nolinkurl{KapStep.kapStep}, where s=−lambda/xi0=1.

The actual source \nolinkurl{lastBi} line exchange, read in the fixed unitary local models for the original vacua, is denoted S\_v. It is a Heisenberg intertwiner for the same symplectic line map as S\_v\^{}0; this follows directly from \nolinkurl{kap} and the Gram equality in FT2. The irreducible Schrödinger models and unitarity therefore give a unique fixed scalar

\[
       \mathcal S_v=\gamma_v\mathcal S_v^0,
                       \qquad |\gamma_v|=1.
       \tag{FT14a}
\]

We retain this phase. It is not inferred to be1 from intertwiner uniqueness or from an abstract Fourier convention. It can equivalently be computed from the actual transformed unmodified local tensor by

\[
 \gamma_v=q^{b/2}
   [\mathcal S_v(G_P\otimes\Phi_1)](0,0,0),
                       \qquad \Phi_1=1_{\operatorname{Mat}_{N,b}(O)}.
 \tag{FT14b}
\]

The equality follows from the explicit calculation FT17 below. Thus gamma\_v is fixed by the actual source maps and fixed vacuum models, not chosen from the desired period. The phase is independent of all finite tuple characters. A convention using S\_v\^{}0 as the local exchange has gamma\_v=1, but its comparison to an independently fixed global CM normalization would have to transport the accompanying phases. No such extra identification is asserted here. Every operator identity below is independent of this phase, and keeps the actual S\_v.

Take a fixed nonzero middle covector Q with functional f\_Q and replace the first test, in the same local model, by

\[
 \Phi'_0(X,y)=\operatorname{phiOne}(P,Q)(X)
               1_{O^b}(y)\bar\chi(-\epsilon f_Q(y)),
 \qquad \Phi_1(Z)=1_{\operatorname{Mat}_{N,b}(O)}(Z).
 \tag{FT15}
\]

The Fourier integral in FT14 is one exactly when xi0 t/pi+varpi\^{}-1 epsilon f\_Q lies in the dual lattice, and is zero otherwise. Using FT6, this condition is bar(pi)t=f\_Q modulo pi. That is exactly the terminal raw phiOne(Q,P\_terminal) condition. Thus

\[
 \boxed{\quad
 \mathcal S_v(\Phi'_0\otimes\Phi_1)
     =\gamma_v q^{b/2}\,\operatorname{phiOne}(P,Q)
                   \otimes\operatorname{phiOne}(Q,P_{\rm terminal}).
 \quad}                                                    \tag{FT16}
\]

The top N rows are integral on both sides; the last physical row has the indicated affine residue. This is an equality of the entire fixed Schwartz tensor, before theta summation or class integration. It is independent of the finite tuple character xi, which is unramified at T and does not change the Heisenberg model.

There is an independent norm check. The first incidence function imposes a independent residue equations on b*a matrix entries, so its squared norm is q\^{}-a. The terminal incidence function imposes b equations on the last row, so its squared norm is q\^{}-b. The left side of FT16 has squared norm q\^{}-a by unitarity; the right side has q\^{}b q\^{}(-a-b)=q\^{}-a. Omitting the positive factor q\^{}b/2 would violate this identity.

The original transformed glued vacuum is obtained from FT14 with FT9, and in the above coordinates equals

\[
 \gamma_v q^{-b/2}
 1_{X\ker P\subset O^b}\,1_{t\in O^b}\,
 \chi\bigl(\xi_0(\varpi/\pi)t\cdot Xv_P\bigr)\,
                         1_{Z\in\operatorname{Mat}_{N,b}(O)}.
 \tag{FT17}
\]

Its squared norm is one: the X domain has volume q\^{}b and its amplitude squared is q\^{}-b. This is the actual initial finite lastBi tensor; it is not asserted to equal a product vacuum.

\subsection{6. One fixed finite Heisenberg operator and common C}\label{one-fixed-finite-heisenberg-operator-and-common-c}

Let v\_bas be the product horizontal basic vacuum, of norm one, and let R\_T be the product of the two raw incidence functions in FT16. We give actual finite Heisenberg sums, without an abstract endomorphism-algebra choice. Write T\_u f(w)=f(w−u) for position translation and M\_l f(w)=chi(l(w))f(w) for modulation. Both are pure Heisenberg operators in the fixed Schrödinger reading. The central phase when a product is written as one Weyl operator is retained in its coefficient.

Choose residue representatives in O. In the following operator displays, an expression chi\_bar(L(w)) means chi(varpi\^{}-1 L\_tilde(w)) with these fixed O-linear lifts, even when w is outside the integral lattice. This defines a genuine modulation on the whole Schwartz space. On integral arguments it is exactly the residue character already used above. On the first vertical space define

\[
\begin{aligned}
 P_{\rm int}&=q^{-b}\sum_{z\in k^b}M_{z\cdot Xv_P},
       &S_y&=\sum_{r\in k^b}T_{(0,r)},\\
 P_0f(X,y)&=q^{-a}\sum_{c\in k^a}
       \bar\chi\bigl(c\cdot(f_QX-f_P)\bigr)f(X,y),
       &M_\epsilon f(X,y)&=
          \bar\chi(-\epsilon f_Q(y))f(X,y),\\
 E_0&=M_\epsilon P_0 S_y P_{\rm int},
       &O_0&=\mathcal S_v(E_0\otimes1)\mathcal S_v^{-1}.
\end{aligned}
\tag{FT18a}
\]

The expression for P\_0 means the corresponding sum of pure modulations, with constant coefficient q\^{}-a chi\_bar(−c·f\_P). Its value on an integral X is the incidence indicator f\_QX=f\_P. On the support of G\_P, Xv\_P belongs to varpi\^{}-1 O\^{}b. Character orthogonality therefore proves that P\_int removes precisely its nonintegral possibilities. Formula FT9 then becomes 1\_(Mat\_(b,a)(O))(X) 1\_(varpi O\^{}b)(y). The q\^{}b y-translations fill all residue cosets once, giving the product integral indicator. Finally P\_0 and M\_epsilon give FT15. Thus E\_0 G\_P=Phi'\_0 exactly. This proof uses no replacement of the original glued vacuum or its amplitude.

There is also an explicit formula for O\_0 itself. Under FT14, translation of y by r becomes modulation by chi(−xi0(t/pi)·r), and modulation by chi(l·y) becomes translation of t by pi\emph{l/xi0. Since −pi}epsilon/(xi0*varpi)=bar(pi)\^{}-1, put

\[
 B_t=\sum_{r\in k^b}M_{-\xi_0(t/\pi)\cdot r},\qquad
 u_Q=\bar\pi^{-1}f_Q.
 \tag{FT18d}
\]

Then O\_0=(P\_0 P\_int) tensor (T\_(0,u\_Q) B\_t), with the latter operator acting only on the last row of the second matrix. This follows directly by substituting y+r in the Fourier integral, so its signs and central phases are fixed. Independently, apply this formula to FT17: P\_int restricts X to the integral lattice, where its displayed character is1 and gamma\_v is unchanged; B\_t then contributes q\^{}b on t in pi O\^{}b and zero elsewhere. The final translation gives bar(pi)t=f\_Q modulo pi. This recovers FT16 with its exact gamma\_v*q\^{}b/2 factor. In particular the same O\_0 is valid for the actual exchange; its scalar cancels under conjugation.

The source horizontal dT operators have the following explicit realizations on the standard basic vacua:

\[
\begin{split}
 d_{T,0}&=P_0,\\
 d_{T,1}f(Z,t)&=q^{-b}\sum_{c\in k^b}
       \bar\chi\bigl(c\cdot(\bar\pi t-f_Q)\bigr)f(Z,t).
\end{split}
\tag{FT18b}
\]

They give exactly the two raw phiOne functions, with coefficient1. This is the literal \nolinkurl{Reps}/\nolinkurl{Over} normalization, not a separate normalization of the raw functions. Specifically, \nolinkurl{IsSumFrame.Over} uses \nolinkurl{IsThetaAt ... 1 f0 ... f}, and \nolinkurl{PlacesToAdeles.OverBasic} fixes the basic lattice indicator to f0 itself. The actual Heisenberg action intertwines this reading, so the displayed character sums give the required raw \nolinkurl{Over} images for every vector with this local vacuum factor. Any other already prescribed finite combination having the same \nolinkurl{Reps} property has the identical action on that vacuum line. Nothing requires those combinations to be equal on the whole Schwartz space. In particular we retain the literal dT operators when they have already been specified.

For an explicit inverse on the needed image, choose w\_Q in O\^{}b with f\_Q(w\_Q)=1 and put

\[
 D'_0=\sum_{c\in k^a}T_{w_Qc^t},\qquad
 D'_1=\sum_{c\in k^b}T_{(0,\bar\pi^{-1}c)},\qquad
 D'_T=D'_0\otimes D'_1,\qquad O_T=D'_T O_0.
\tag{FT18c}
\]

In D'\_0 the position variable is the first b-by-a matrix. In D'\_1 it is the last row of the second matrix; the N top rows are unchanged. The first translates run once through all values of f\_QX, and the second through all values of bar(pi)t. Hence their affine supports are disjoint and partition the product integral lattice. Therefore D'\_T R\_T=v\_bas, with no averaging denominator and no undetermined scalar. FT16 now gives, exactly,

\[
        C_{T,v}:=O_T\mathcal S_v(G_P\otimes\Phi_1)
                         =\gamma_v q^{b/2}v_{\rm bas}.
 \tag{FT19}
\]

By FT18b, the actual raw operator dT0 tensor dT1 sends v\_bas to R\_T under the same source Over normalization. Therefore (dT0 tensor dT1) C\_T equals FT16, exactly as required in the source LastTAll/LevelPrimeLaw construction. The operator O\_T and the tensor C\_T are chosen once, before any tuple character varies. They depend only on the fixed line, finite frames, additive character and incidence functions. This supplies a concrete realization of the source finite op, rather than dropping it from C.

These are operators of the exact source type. \nolinkurl{JointSmooth.comb} is a finite sum of \nolinkurl{heisO} operators; \nolinkurl{biComb} is their product action, and \nolinkurl{PlaceLaws.BiOp} is a finite sum of such products. The line shuffle is a symplectic Heisenberg intertwiner. It sends each term of E\_0 to one Weyl operator on the orthogonal direct sum of the two horizontal symplectic spaces, hence to a product of two Weyl operators. Position/modulation products are reduced using their actual central character. The split local section embeds every parameter as an adele supported at this T place, so both factors satisfy the literal \nolinkurl{CombAt} and therefore \nolinkurl{IsFinComb}; their infinite component is zero. This proves \nolinkurl{OpAt} for O\_T. The same construction tensored with the identity on all other local and real factors works for every vector in the relevant local vacuum line. It therefore has the uniform quantification of the source last-two operator law.

More explicitly, FT18c--d give O\_T=(D'\_0 P\_0 P\_int) tensor (D'\emph{1 T}(0,u\_Q) B\_t). Thus the source \nolinkurl{BiOp.n} can be1. Each of its two \nolinkurl{Finsupp} entries is obtained by expanding the indicated finite product and collecting Weyl operators with their central phases. This is stronger than merely proving that some finite sum of product operators agrees on a quotient.

For definiteness all these functions lie in a finite Schrödinger space supported in varpi\^{}-M O\^{}d and invariant under varpi\^{}M O\^{}d for some fixed M. The displayed sums act on the entire Schwartz space, not just on this quotient, and the proved identities hold on the actual stated vectors. Thus no extension of an abstract matrix unit is being assumed.

Here q=T is a rational prime. All displayed Heisenberg coefficients lie in Z{[}zeta\_T,1/T{]}; one may retain the explicit finite coefficient field enlargement K0(zeta\_T), of relative degree at most T−1. Their common denominator divides T\^{}(a+b). Indeed P\_int and P\_0 contribute T\^{}-b and T\^{}-a; all other sums have integral root-of-unity coefficients. After collecting products, the Weyl central phases are T-th roots of unity, since every occurring pairing is in T\^{}-1 O and 2 is a unit. The parameter transport through the shuffle preserves the symplectic pairing. The explicit estimates are

\[
 \|E_0\|_{\ell^1}\le q^b,\qquad
 \|O_T\|_{\ell^1}\le q^{a+2b},\qquad
 \|C_{T,v}\|=q^{b/2}.
 \tag{FT19a}
\]

The first two are bounds for the displayed finite Weyl expansions, and also bound their operator norms because every Weyl operator is unitary. Their number of terms before collecting is at most q\^{}(2a+3b). The positive q\^{}b/2 is an integer because b is even. Thus these finite operator coefficients are integral at11. This coefficient-field assertion does not assign a field or an integral valuation to the retained source-model phase Gamma\_T. These bounds do not imply an integral arithmetic row, a stabilizer cancellation, a real primitive normalization, or a primitive image minor.

For all f selected T places, tensoring FT13 and FT19 yields

\[
 c_T=q^{-fb/2}(q^b-1)^f
       \prod_{v\mid T}\!(\mu_0\xi)(\operatorname{idele}_v\varpi),
 \qquad C_T=\Gamma_T q^{fb/2}v_{\rm bas,T},
 \qquad \Gamma_T=\prod_{v\mid T}\gamma_v,\quad |\Gamma_T|=1.
 \tag{FT20}
\]

The positive amplitudes cancel if c\_T is multiplied into the physical tensor C\_T, leaving Gamma\_T*(q\^{}b−1)\^{}f times the exact Hecke phase. The source stores c\_T and C separately, so both must be retained until that multiplication. This is not a new lift saving.

\subsection{7. Relation to DO050 and the remaining common-tensor gates}\label{relation-to-do050-and-the-remaining-common-tensor-gates}

DO050 remains frozen. Its word ``primitive'' in RC13 means the PN15 primitive embedding normalization with a Fischer-isometric selected Pieri input. It does not assert that the higher selected PBW polynomial has primitive integer coefficients. If an integral primitive Pieri map has j* j=q\_lambda I, its squared RC15 factor also contains q\_lambda, exactly as K043(PB8) and LIFT40 specify.

The real compact averaging in DO050 is valid inside the middle period integral. It is not an equality between every simple lastBi tensor and its invariant projection. Here FT16--FT19 are genuine equalities of entire finite tensors, so the same common C can be used before and after taking that real compact average. The exact real selected coevaluation factor, its fixed phase, the PBW factor and the terminal branch remain DO050's factors.

At G the recovered source SeeSawLastTwoStepsValueFixedLevelPlaces.last\_value explicitly chooses cG=1 after its actual GLoc conditions are supplied. At Q, LastStepsKernelEqSumCouplingCosets.kernel\_eq\_sum is an exact pointwise full-Y tensor identity. Neither statement is promoted to an already verified new-host datum solely because it exists in the source: the actual finite G/Q frames, trace insertions and character conventions still have to be matched. The current report closes their T counterpart with unchanged tests.

Even after all three finite laws hold, the literal global equation is ValLastX with C=op\_T.app(lastBi(Phi0,Phi1)) and cV=(sc0 sc1)cQ cG cT. It is not the assertion that every fixed horizontal highest product has that period or that the actual global bottom restricted period is nonzero. The terminal mass and coefficient-side branch are the checked DO038 adapter.

\subsection{8. Exact source locators and verification record}\label{exact-source-locators-and-verification-record}

All recovered files match their intake hashes at pinned commit 7cbf7a0c32fca1badf05931879a37c3f8a2b03b4. Newly recovered source is under work/agents/definite\_oscillator\_seesaw\_source.

\begin{itemize}
\tightlist
\item
  \nolinkurl{Definitions/TwistedInnerProducts/HasConductorMulShift.lean:} dia/Nm/Mram29--46, FFStar79--80, exact block136--138, phiOne/maps336--359, gluedS405--417 and its membership formula.
\item
  \nolinkurl{BlockTheory/TwistedDiscreteSeries/DualLiftInvariance.lean:} datum/charInd/lastPair, TransformFormula163--189 and LastStep213--237.
\item
  \nolinkurl{Proofs/DecoratedLatticeProofs/GluedPlaceSourceConstantProof.lean:} exists\_theta\_exact147--200 and sourceExact207 onward, with explicitly chosen r0=q\^{}-b q\^{}(b/2), exponent1.
\item
  \nolinkurl{SmallPairPairingAtOneGluedPlace.lean:} integral\_sum\_translates and smallOne, including the exact q\^{}b−1 compact orbit count and no source/target Haar reset.
\item
  \nolinkurl{LatticePairCharacterIndicatorTransformEqCosetInd.lean:} fourier and kapStep, epsilon'=red(lambda)\^{}-1, full negative-sign Fourier phase and volume constant retained in the local model.
\item
  \nolinkurl{LastKernelDatumInLastLineFrameAtLevelPrime.lean:} hasVal\_transport/maps\_transport and lastKernel\_reps, u=s\^{}-1.
\item
  \nolinkurl{LevelPrimeLawLastTwoStepsAllTuples.lean47–79:} epsilon=(epsilon')\^{}-1 u and fixed op constructed before tuples.
\item
  \nolinkurl{Definitions/LatticeFrames/JointSmooth.lean20–28:} actual comb, IsFinComb and biComb; PlaceLaws.lean20--27: actual finite-product BiOp. CombAtSingleZero.lean21 and SmallPair.lean26 give the exact support predicates.
\item
  \nolinkurl{Definitions/LatticeFrames/IsSumFrame.lean187–190} and \nolinkurl{PlacesToAdeles.lean169–179:} literal Over/IsThetaAt coefficient1, its uniqueness, and the basic lattice indicator mapping to the original vacuum. SmallPair.lean28--32 is the uniform Reps property used in FT18b; this is an equality of actions on the local vacuum line, not equality of arbitrary finite operators.
\item
  \nolinkurl{Formalism/LocalToGlobal/Columns/StdLatUnimodular.lean44,} 165--179: z0=pi e\_last and actual terminal functional.
\end{itemize}

I reconstructed all eight sections of the saved proof, including FT1--FT20, the explicit FT18a--d finite sums, their literal \nolinkurl{BiOp.n=1} realization, and the retained source-model phase. I read the complete certificate and the cited exact source clauses. The source constant, Fourier signs, actual terminal functional, positive dilation, \nolinkurl{OverBasic} normalization, and ramified input operator were checked separately. The unit phase in FT14a--b is retained and is not evaluated by a model-uniqueness argument.

The companion \nolinkurl{definite_oscillator_fixed_t_seesaw_certificate.py} passed \textbf{362,363 exact assertions at 2026-11-22T06:53:14.134251+00:00}. These include bidirectional integral lattice changes; 120,400 terminal phase cases; the nonzero Fourier Gram and all tested scalar-orbit augmentations; the finite glued expansion; 161,296 complete-quotient E0 checks; 32,432 modified-shuffle and 32,432 explicit-O0 checks; the exact affine image inverse; norms, coefficient counts and both actual host parameters. The finite calculations support the written proof and do not replace its general source or Heisenberg arguments. No Lean execution, global period evaluation, or integral theta-image computation was performed.

\clearpage

\section{ALP-LIFT-41 --- a principal-G point of the actual T packet near the standard origin}\label{alp-lift-41-a-principal-g-point-of-the-actual-t-packet-near-the-standard-origin}

Author: GPT-6 Astra Pro, ALP\_MACHINE \nolinkurl{/root/lift_accounting}. 22 November 2026 UTC; inherited reasoning effort is not separately exposed.

\textbf{Author whole-proof review PASS, 2026-11-22 07:00:23 UTC.} Independent whole-proof review is pending. Exact diagnostic certificate: 7,019 assertions PASS, 2026-11-22 06:59:46 UTC.

This finding constructs a different point of the unchanged T-translated packet. Unlike K047's particular cuspidal representative, this point is very close to the standard origin. Its membership uses the actual lower T translate, including the conjugate component, and the principal congruence at both G primes. The exact selected compact coefficient along the corresponding real rank-one subgroup is retained below. At the upper admissible gap, the complete normalized selected jet at this actual packet point has a principal-level right inverse with the quantitative bound in NP20. No nonzero genuine-e0 automorphic jet, terminal row or global period is claimed.

\subsection{1. Fixed integers and an explicit norm identity}\label{fixed-integers-and-an-explicit-norm-identity}

Use the unchanged data of DO040 on either K031/K032 host. Write

\[
 Q=g_0g_1,\quad \delta=\sqrt{-7},\quad
 \pi=1+2D+D\delta,\quad T=\pi\bar\pi=11D^2+4D+1.
 \tag{NP1}
\]

The actual congruence is \(\pi\equiv1\pmod {Q^2}\). Put

\[
 k=Q^2,\qquad \pi_k=\pi-7k+2k\delta,\qquad
 R=7k(11k-2)=77Q^4-14Q^2.
 \tag{NP2}
\]

Direct multiplication gives

\[
 \pi_k\bar\pi_k=T+R.
 \tag{NP3}
\]

Indeed the coefficient of D in the cross term cancels, and the remaining difference is \(-14k+77k^2\). The following exact four-square representations fix the positive coordinates; they are data, not an ineffective appeal to a representation theorem.

\begin{longtable}[]{@{}
  >{\raggedright\arraybackslash}p{(\columnwidth - 4\tabcolsep) * \real{0.3000}}
  >{\raggedleft\arraybackslash}p{(\columnwidth - 4\tabcolsep) * \real{0.4000}}
  >{\raggedright\arraybackslash}p{(\columnwidth - 4\tabcolsep) * \real{0.3000}}@{}}
\toprule\noalign{}
\begin{minipage}[b]{\linewidth}\raggedright
Host
\end{minipage} & \begin{minipage}[b]{\linewidth}\raggedleft
Q
\end{minipage} & \begin{minipage}[b]{\linewidth}\raggedright
u1,u2,u3,u4
\end{minipage} \\
\midrule\noalign{}
\endhead
\bottomrule\noalign{}
\endlastfoot
J14 & 24352637893193 & 2,40538469,57165115,213693530251813 \\
J56 & 1427908193729121522601 & 2,95192390527,383092671043,12529843548438452388609 \\
\end{longtable}

In each row,

\[
 \sum_{j=1}^4u_j^2=7(11Q^2-2).
 \tag{NP4}
\]

For \(a=42\) or168 set

\[
 w=Q(u_1,u_2,u_3,u_4,0,\ldots,0)^t,\quad
 v=(w,\pi_k)^t,\quad z=v/\bar\pi.
 \tag{NP5}
\]

With \(J=\operatorname{diag}(I_a,-1)\),

\[
 v^*Jv=-T,\qquad z^*Jz=-1,\qquad
 z\equiv e_{a+1}\pmod Q.
 \tag{NP6}
\]

The last congruence is at every selected and conjugate G component, where both \(\pi\) and \(\bar\pi\) are1 modulo \(Q^2\). The division by \(\bar\pi\), rather than \(\pi\), is essential for the lower-unipotent T convention.

\subsection{2. Rational unitary extension and the finite orbit}\label{rational-unitary-extension-and-the-finite-orbit}

There is an explicit rational unitary matrix sending \(e=e_{a+1}\) to z. Let \(c=z_{a+1}\), \(b=z-e\), and set

\[
 \gamma=I+\frac{bb^*J}{1-\bar c}.
 \tag{NP7}
\]

Here \(c\ne1\): the real part of \(\pi_k-\bar\pi\) is \(-7k\). Since \(b^*Je=1-\bar c\) and \(b^*Jb=c+\bar c-2\), multiplying the rank-one matrices proves

\[
 \gamma^*J\gamma=J,\qquad \gamma e=z.
 \tag{NP8}
\]

The determinant is a norm-one rational scalar. It can be corrected by multiplying on the right by an element of H=U(I\_a), without changing either the negative vector or any double coset below. Special-unitary membership is not needed for the packet assertion.

We prove

\[
 \gamma_f^{-1}\in H(\mathbb A_f)\,\tau\,J^1,
 \tag{NP9}
\]

where \(\tau\) is the actual lower-unipotent T translation and \(J^1\) is principal at G, maximal elsewhere, as in R025/R027.

Away from T, z is integral and has unit norm. For the everywhere finite-unramified CM extension, a norm-unit vector in a unimodular Hermitian lattice extends to an integral orthogonal decomposition. Its orthogonal complement is unimodular: projection onto the norm-unit line is integral, and the determinant of the Gram matrix splits as its unit norm times the complement determinant. Here is the local classification argument, including the dyadic case. A nonsingular Hermitian form over the quadratic residue-field extension has a vector of nonzero norm; otherwise polarization, including multiplication of one argument by an element with nonzero residue trace, would make the whole form zero. Lifting such a vector and splitting its norm-unit line diagonalizes the lattice by induction. Every diagonal unit is a unit norm. The residue norm is onto; successive corrections use

\[
 N(1+\varpi^m x)\equiv1+\varpi^m\operatorname{Tr}(x)
                       \pmod {\varpi^{m+1}},\qquad m\ge1,
\]

and the residue trace is onto even in characteristic2. Completeness therefore lifts any desired unit norm. Rescaling the diagonal basis shows that two such unimodular forms of the same rank are integrally isometric. Apply this to the complement of z and the standard complement of e to extend the isometry between the two norm−1 lines. At split places the same assertion is just extension of a primitive vector and its dual covector to a basis. Therefore a local integral unitary matrix carries e to z, and \(\gamma\in J^{\max}H\).

At G the matrix can be chosen congruent to I. An explicit split description suffices: if the desired column is q and its dual row is l, with \(lq=1\) and both congruent to e modulo the residue prime, take the first a columns to be \(e_j-(l_j/l_{a+1})e_{a+1}\) and the last column to be q. The resulting matrix has determinant \(1/l_{a+1}\), hence it and its inverse are integral; it is I modulo the prime, and the inverse last row is l. Its other CM component is the inverse adjoint. Thus the actual principal condition is retained.

At a selected T component \(\bar\pi\) is a unit, so z is integral and primitive. Its dual row is

\[
 -z^*J=(-w^t/\pi,\bar\pi_k/\pi).
 \tag{NP10}
\]

It has valuation exactly−1: the first nonzero positive coordinate is \(2Q\), prime to T. Its pairing with z is1. These are precisely the integral-column/denominator-one-row invariants of \(\tau^{-1}\). To see completeness, first carry the primitive column to e by an integral GL matrix. The dual row is then \((\pi^{-1}u',1)\), where u' is primitive modulo \(\pi\). An integral change in the first a coordinates carries u' to the actual retained covector modulo \(\pi\); integral row changes remove the remaining integral part. This gives \(\gamma\in J^{\max}\tau^{-1}H\) at T. Supplying the inverse-adjoint component proves the same unitary assertion at its CM partner. Combining the local decompositions gives NP9.

Consequently for a finite representative x lying above an actual initial class,

\[
 [(o,x\tau)]=[(\gamma_\infty o,1_f)]
 \quad\hbox{in the principal-level bottom quotient.}
 \tag{NP11}
\]

Precisely, the principal-level packet is indexed by \(H(F)\backslash H(\mathbb A_f)/(H\cap\tau J^1\tau^{-1})\). It is the principal-G refinement of the original initial character level. An original initial double coset alone need not specify a unique point of this refinement. NP9 proves the stronger principal membership needed for NP11.

No new prime, changed T covector, principal29 refinement or new canonical twist is introduced. A common prime divisor of all entries of v would divide T by NP6; it cannot divide the positive coordinate 2Q. Thus v is primitive at every finite place as well.

\subsection{3. Exact real distance}\label{exact-real-distance}

At every CM label the Cartan distance s from o to \(\gamma o\) has the same value, with

\[
 \sinh^2s=R/T,\qquad \cosh^2s=1+R/T.
 \tag{NP12}
\]

Indeed the negative-line vector z has positive Euclidean squared length R/T and negative-coordinate squared modulus1+R/T. The compact phase and positive direction can be absorbed into two compact factors. The real selected compact module is unchanged by choosing the corresponding left and right compact vectors.

\subsection{4. Exact selected rank-one coefficient}\label{exact-selected-rank-one-coefficient}

Retain the unit compact-highest vector in the K043 selected component. Use one-based indices in this section. Let r=\(\lambda_a-\alpha_a\), let \(d_a=\alpha_a-\beta-(a-1)=2\ell+q+1\), and let \(\kappa=\lambda_a-\beta+N\) be its weight for the last positive/negative rank-one subgroup. The lowering operator is \(D_a=\partial_{z_a}\). It commutes with every compact positive root, so its j-th nonzero image is in the unique Pieri component \(\lambda-j\epsilon_a\).

The polynomial degree in \(z_a\) is at most r: every weight of the minimal coefficient representation has its last coordinate at least \(\alpha_a\). Hence \(D_a^{r+1}\) kills this vector. Using K043's Hilbert/Fischer ratio separately on each image component gives

\[
 \|D_a^jv_{\lambda}\|^2
 \le {r!\over(r-j)!}\,(d_a+r-j)_j,
 \qquad0\le j\le r.
 \tag{NP13}
\]

Here is the complete norm argument. Choose the highest vector P in the polynomial Pieri summand with Fischer norm1. Its Hilbert-unit version is \(v_\lambda=\sqrt{D_\lambda}P\), where \(D_\lambda\) is K043(PB5). The compact weight equation for a monomial term in P is coefficient weight plus monomial exponent equal to \(\lambda\). Since a weight of \(V_\alpha\) has every coordinate at least \(\alpha_a\), the last exponent is at most r. Therefore

\[
 \|\partial_a^jP\|_{\mathrm F}^2
 \le r(r-1)\cdots(r-j+1)\|P\|_{\mathrm F}^2.
\]

This follows term by term from the Fischer norm: differentiating a last-coordinate monomial of degree h multiplies its squared norm by \(h(h-1)\cdots(h-j+1)\). The derivatives of distinct monomials remain orthogonal in their coefficient-weight spaces.

For every compact positive matrix unit \(E_{i,i+1}\), \(i<a\), its commutator with \(D_a\) is zero. Thus the nonzero derivative is a compact highest vector of weight \(\lambda-j\epsilon_a\). This weight still interlaces \(\alpha\) for \(j\le r\); its Pieri multiplicity is one. On that whole image component the Hilbert norm is the Fischer norm divided by \(D_{\lambda-j\epsilon_a}\). All factors in their ratio cancel except the last one, giving

\[
 {D_\lambda\over D_{\lambda-j\epsilon_a}}
       ={(d_a)_r\over(d_a)_{r-j}}=(d_a+r-j)_j.
\]

This proves NP13, including a zero derivative. It uses the actual selected vector; it is not an estimate on an arbitrary ambient monomial of the same total degree.

The rank-one Gauss decomposition gives the exact coefficient

\[
 \langle v_{\lambda},\pi(a_s)v_{\lambda}\rangle
 =(\cosh s)^{-\kappa}
 \sum_{j=0}^r {(-\sinh^2s)^j\over(j!)^2}
                         \|D_a^jv_{\lambda}\|^2.
 \tag{NP14}
\]

To check its normalization, choose the matrix-unit signs so that the raising operator E is the adjoint of D and \([D,E]=H\), \([H,E]=2E\). The real rank-one Cartan element is \(\exp(s(E-D))\). Its two-by-two Gauss factorization is

\[
 \exp(\tanh(s)E)\,
       \exp(-\log(\cosh s)H)\,
       \exp(-\tanh(s)D).
\]

Expand the last exponential on v, where it terminates after r. The j-th term has H-weight \(\kappa-2j\). Pairing the first exponential with v uses only its j-th term, and adjunction gives \(\|D^jv\|^2\). The scalar is \((\cosh s)^{-\kappa}(-\sinh^2s)^j/(j!)^2\), as in NP14. The identity holds first in the local analytic action and then for all real s by the holomorphic fractional-linear realization. Reversing both matrix-unit signs changes no displayed coefficient.

The inner product and the normalized compact vector are the same as K043; primitive PBW vectors carry its already computed norm factor. In particular, the difference of the polynomial in NP14 from1 is at most

\[
 \exp\{r(d_a+r)R/T\}-1.
 \tag{NP15}
\]

Indeed \(r!/(r-j)!\le r^j\), \((d_a+r-j)_j\le(d_a+r)^j\), and \((j!)^{-2}\le(j!)^{-1}\). Summing the absolute values gives NP15. The inequality is deliberately uniform in the chosen compact highest vector and the actual word.

\subsection{5. Actual numerical top-label estimates}\label{actual-numerical-top-label-estimates}

The fixed integers give the following rigorous bounds. They are verified by integer cross-multiplication in the companion certificate.

\begin{longtable}[]{@{}
  >{\raggedright\arraybackslash}p{(\columnwidth - 8\tabcolsep) * \real{0.1579}}
  >{\raggedleft\arraybackslash}p{(\columnwidth - 8\tabcolsep) * \real{0.2105}}
  >{\raggedleft\arraybackslash}p{(\columnwidth - 8\tabcolsep) * \real{0.2105}}
  >{\raggedleft\arraybackslash}p{(\columnwidth - 8\tabcolsep) * \real{0.2105}}
  >{\raggedleft\arraybackslash}p{(\columnwidth - 8\tabcolsep) * \real{0.2105}}@{}}
\toprule\noalign{}
\begin{minipage}[b]{\linewidth}\raggedright
Host
\end{minipage} & \begin{minipage}[b]{\linewidth}\raggedleft
\(\sinh^2s=R/T\)
\end{minipage} & \begin{minipage}[b]{\linewidth}\raggedleft
top r at \(\ell_+\)
\end{minipage} & \begin{minipage}[b]{\linewidth}\raggedleft
polynomial error in NP14
\end{minipage} & \begin{minipage}[b]{\linewidth}\raggedleft
absolute coefficient lower bound
\end{minipage} \\
\midrule\noalign{}
\endhead
\bottomrule\noalign{}
\endlastfoot
J14 & \(<2^{-813}\) & 19 & \(<2^{-320}\) & \(>\exp(-2^{148})\) \\
J56 & \(<2^{-704}\) & 12 & \(<2^{-211}\) & \(>\exp(-2^{257})\) \\
\end{longtable}

For completeness, \(Q<2^{45}\), \(T\ge2^{1000}\) in J14, and \(Q<2^{71}\), \(T\ge2^{995}\) in J56. Since \(R<77Q^4<2^7Q^4\), these already prove the displayed radius bounds. K044 has \(\ell_+<2^{485}\) and \(0\le q\le\ell_+\). Thus \(d_a+r<2^{487}\) and \(r<2^5\). The exponent in NP15 is less than \(2^{-321}\) or \(2^{-212}\), respectively. For \(0\le x<1/2\), \(e^x-1\le2x\), proving the error bounds.

The actual top weights give \(0<\kappa<2A=2^{961}\): use K043(PB13--14), sum the Pieri exponents, and note that every remaining term is bounded by \(4a^2\ell+10\ell+2a+100<A\). Therefore

\[
 -\log\left|\langle v_\lambda,\pi(a_s)v_\lambda\rangle\right|
 \le {\kappa\over2}\log(1+R/T)
       -\log(1-\varepsilon)
 < A R/T+2\varepsilon,
 \tag{NP16}
\]

where \(\varepsilon\) is the polynomial-error bound. This is smaller than the two stated powers of2. The compact factors carrying the actual positive direction to the last coordinate are absorbed into the chosen left and right vectors inside the same selected compact module. Their norms stay one; any associated unit phases remain part of that choice.

At thick or thin labels, r is \(E+\ell\), so NP15 can exceed one. No nonzero lower bound for these labels is asserted here. The exact finite polynomial in NP14 remains valid, with all its derivative norms and the envelope retained. In particular, the favorable top estimate is not multiplied into a claimed all-label nonzero value.

\subsection{6. A quantitative selected one-point inverse at this actual packet point}\label{a-quantitative-selected-one-point-inverse-at-this-actual-packet-point}

The near-origin point permits a direct application of the already proved K046 selected-jet theorem, including all thick and thin labels. This consequence does not use a nonzero product of the rank-one coefficients in NP14.

Let \(\Gamma_1=G(F)\cap J^1\), in the standard global coordinates, and let \(C_5\) be the product of the real Cartan sets of parameter less than5. Define

\[
 \Gamma_z=\gamma_\infty^{-1}\Gamma_1\gamma_\infty.
 \tag{NP17}
\]

Then

\[
                 \Gamma_z\cap C_5C_5^{-1}=\{1\}.
 \tag{NP18}
\]

Here is the complete separation proof. Every entry of \(h-I\), for \(h\in\Gamma_1\), is in \(Q\mathcal O_E\), by the actual standard split frames and both selected and conjugate G components in R025/K046. Its entries are integral at every other finite place. By NP12 and the radius bounds, \(s<1/2\): indeed \(\sinh s> s\) for \(s>0\), and \(R/T<1/4\). The Euclidean operator norms of \(\gamma_\infty\) and its inverse are \(e^s\) at each CM label, since their Cartan parameter is s. Thus, if \(\gamma_\infty^{-1}h\gamma_\infty\in C_5C_5^{-1}\), every conjugate of every entry of \(h-I\) has modulus less than

\[
             e^{10+2s}+1<e^{11}+1<3^{11}+1=177148<Q.
 \tag{NP19}
\]

If such an entry were nonzero, dividing it by Q would give a nonzero algebraic integer all of whose conjugates have modulus less than1. Its nonzero integer norm cannot have absolute value less than1. Every entry is zero, proving NP18. The embeddings used here and their conjugates exhaust all embeddings of E; no bound at only the selected half of the CM places is being substituted for a norm bound.

Now take \(\ell=\ell_+\), every actual binary word, and every admissible canonical twist \(0\le q\le\ell_+\) in K044. Write \(V_{\rm sel}\) for the complete selected compact fibre, with the K043 normalized PBW metric. Under the ball isometry carrying the origin to \(\gamma_\infty o\), K046(BL20--25) applies to \(\Gamma_z\) with the same uniform concentration error \(\delta_0=f2^{-\ell_+/8}<2^{-100}\). Its exact adjoint and Gram argument BL28--34 consequently gives a finite-dimensional space of cuspidal Poincare sections on the unchanged principal component and a right inverse to the normalized selected jet at the actual point NP11 satisfying

\[
 \mathcal J_z\mathcal R_z=I_{V_{\rm sel}},\qquad
 \|\mathcal R_z\|\le
       {\sqrt{V_G}\over1-2\delta_0},\qquad
 V_G=\operatorname{vol}(\Gamma_1\backslash G(\mathbb R)).
 \tag{NP20}
\]

For clarity, the transfer is an isometry between sections on \(\Gamma_z\backslash G(\mathbb R)\) and on \(\Gamma_1\backslash G(\mathbb R)\), using left translation by \(\gamma_\infty\) and the invariant coefficient-bundle metric. It identifies the origin jet with the transported jet at \(\gamma_\infty o\). Inner conjugation preserves Haar measure, so the quotient volume is precisely the same \(V_G\). The compact fibre is transported in full. Thus NP20 concerns every vector of the selected component, not only its highest line. The raw derivative requires K043's already retained PBW factor; a different global Haar normalization changes \(V_G\) as in K046(BL34).

This removes the translated-point separation condition of K046(BL35) for this explicit representative of the actual packet. It does not resolve the full finite-deck packet or the genuine-e0 projection of the resulting sections.

\subsection{7. Remaining global issue}\label{remaining-global-issue}

NP11 is a favorable point of the prescribed packet, not a replacement packet. It does not contradict K047/K049, which treat another actual point. Neither the explicit local coefficient nor principal-level Poincaré nonvanishing proves that the genuine finite e0 projection retains a selected coefficient at this point. That same-projector automorphic issue, and the actual terminal pairing, remain open.

NP20 extends K046's standard-point criterion to this actual packet point. It does not prove that the genuine-e0 projection retains that selected coefficient, nor prescribe all values of an initial Mram row. Distinct points of the full finite-deck packet can remain close. The representation-theoretic and automorphic condition therefore remains separate from the explicit one-point inverse.

\subsection{8. Source and contribution ledger}\label{source-and-contribution-ledger}

The exact CM extension and its finite-unramified property are carried from accepted DP57 and K029/K032. R025/R027 identify the actual standard finite frames, principal G quotient and packet. DO040 supplies the unchanged norm generator, both G primes and T, including primality, complete splitting and \(\pi\equiv1\pmod{Q^2}\). LIFT29/LIFT30 fix the lower-unipotent T convention. K047 supplies the earlier cuspidal packet point; it is not re-proved or contradicted.

The geometric construction is a joint root/lift contribution. The initial observation was \(N(\pi-7+2\delta)=T+63\); root supplied the principal-G repair \(k=Q^2\). The present proof supplies the actual four-square coordinates, rational matrix, full finite orbit, primitive integral vector and selected rank-one coefficient. The actual one-point inverse NP20 uses the accepted uniform selected concentration and Gram theorem K046(BL20--34), after the new complete separation proof NP17--19. It does not re-prove that concentration theorem or its cuspidal Poincare construction.

For NP13 the representation norm theorem applied is the accepted K043 translation of Nakahama (2016), Theorem5.1, printed pp.960--961: \url{https://jolt.centre-mersenne.org/item/10.5802/jolt.919.pdf}. NP13 is deduced from that norm ratio and the ordinary Fischer derivative identity. NP14 is proved by the displayed two-by-two Gauss factorization; no additional theta normalization is assumed.

\subsection{9. Exact diagnostic certificate}\label{exact-diagnostic-certificate}

\nolinkurl{near_origin_actual_packet_certificate.py} and its JSON output verify the explicit four-square identities, the actual norm/G/T relations, primitivity, rational unitary matrix and its negative column, the principal split-frame construction, the actual numerical bounds, and the exact integer margin in the new separation proof. They also compare the rank-one Gauss formula with an independent power-series action in702 small scalar discrete-series cases and check the Fischer/Hilbert lowering inequality on explicit GL2 Pieri highest polynomials. No actual automorphic function, Poincaré series, e0 projection or global period is evaluated.

\clearpage

\section{AM-R-034 --- the actual Q tensor, passive core frame and source scalar}\label{am-r-034-the-actual-q-tensor-passive-core-frame-and-source-scalar}

Author: GPT-6 Astra Pro, ALP\_MACHINE /root. 22 November 2026 UTC; reasoning effort is not separately exposed.

\textbf{Whole-proof PASS, 22 November 2026 UTC.} Root read the entire proof and its 80,222-assertion exact certificate. Independent whole-proof review by \nolinkurl{/root/definite_oscillator} passed at07:07:23 UTC on the same mathematical text; only this status paragraph and the final check-status sentence were subsequently updated. Check events retain both hashes.

This identifies the unchanged raw Q tests with the full horizontal Y tensor for the actual LIFT30 negative line. It fixes the phase of the passive source-frame comparison, rather than declaring an integral frame change to have scalar one. The identity is an equality of entire Schwartz tensors in the same models. It does not identify an arbitrary initial trace with the source trace, remove a class weight, or assert the complete global \nolinkurl{QLast} record without its remaining hypotheses. The G vector adaptation is separate and is not inferred here.

\subsection{1. Actual data and normalized passive model}\label{actual-data-and-normalized-passive-model}

Use either actual K031/K032 host, a=42 or168, r=a+1, b=r+t=1802, N=r+2t, and the actual split auxiliary Q place of LIFT23 section1. Its local field is Q\_q, with ring O and residue field k of size q; q is disjoint from T, G and the fixed bad set. The actual character psi\_k is the one in \nolinkurl{Place.IsChar}:

\[
 \psi_k(x)=\psi_E(\xi_0\iota_v(x)),\qquad
 \operatorname{cond}(\psi_k)=O,\qquad \operatorname{vol}(O)=1.
 \tag{R34.1}
\]

This psi\_k is not to be replaced by the bare split character. Let pi bar(pi)=T, c=bar(pi), and g be LIFT30(NL4). At this Q place g and g inverse are integral, det(g)=1, and c and pi are units. The actual negative and positive line frames and map are

\[
 \alpha=g\operatorname{diag}(I_a,c),\quad
 \gamma=\operatorname{diag}(I_N,\pi),\quad
 \tau=-\xi_0c/\pi.
 \tag{R34.2}
\]

The positive horizontal initial frame is the one transported by the actual isometry g on the complement; the middle and terminal positive frames are unchanged. Hence alpha is exactly the source \nolinkurl{IsCoreFrame} with s=1, and tau\emph{pi=-xi0}c is the required line equation. The final physical row t has coordinate eta=t/pi.

Let T\_std be the fixed local Heisenberg model of the first vertical space, normalized by sending the integral matrix indicator to the original basic vacuum. Define a new reading of that same model by

\[
 T_\alpha(\Phi)=T_{\rm std}\bigl(W\longmapsto\Phi(W\alpha)\bigr).
 \tag{R34.3}
\]

This is a passive coordinate substitution. It is a model for source frame alpha and target frame I, with the same basic vacuum.

\textbf{Proof.} The source \nolinkurl{Place.up/down} definitions are up(x)=beta\_target\emph{x}beta\_source inverse and down(W)=beta\_target inverse\emph{W}beta\_source. Thus translation by x in alpha coordinates is translation by x alpha inverse in physical coordinates. Substitution in R34.3 gives exactly the \nolinkurl{IsModel.trans} identity. The same substitution in the symplectic character gives \nolinkurl{IsModel.modul}; the conjugate half is transported by the dual map. The theta-function condition is preserved by the original model. Since alpha is in GL\_r(O), the integral matrix indicator is fixed by W-\textgreater W alpha. Therefore its image is the same prescribed vacuum, not a new unit multiple. These are the complete local model clauses used in the source \nolinkurl{LastLoc} identity. QED.

Here and below the plus/minus CM components are paired by the actual Hermitian adjoint. No independent change of the other half of a unitary element is permitted.

\subsection{2. The unchanged raw first function and its exact scalar}\label{the-unchanged-raw-first-function-and-its-exact-scalar}

For W in M\_(b,r)(F), write W\_top for its first r rows. The unbarred vertical Q function of DatumVec and LIFT31 is

\[
 \Phi_Q(W)=\mathbf1_{W\in M_{b,r}(O),\ W_{\rm low}\in qM_{t,r}(O)}
                   \chi(\det\bar W_{\rm top}),
 \tag{R34.4}
\]

where chi is the actual tuple residue character, extended by zero on singular determinants. This is the unbarred function used by the downward correspondence. Its conjugate is used by the adjoint; the distinction is retained. For every W, including outside the support,

\[
 \Phi_Q(W\alpha^{-1})=\chi(\bar c)^{-1}\Phi_Q(W).
 \tag{R34.5}
\]

\textbf{Proof.} Right multiplication by an integral invertible r-by-r matrix preserves both displayed lattice conditions and their complements. The top determinant is multiplied by its determinant. Here alpha inverse=diag(I\_a,c inverse) g inverse and det(g)=1. The resulting determinant multiplier is c inverse. The same identity remains true for singular top matrices because they stay singular. QED.

Consequently the actual, unchanged vertical vector is

\[
 T_{\rm std}(\Phi_Q)
       =\chi(\bar c)^{-1}T_\alpha(\Phi_Q).
 \tag{R34.6}
\]

The bar over c in chi(bar c) denotes reduction modulo q in R34.5--R34.6, not a second CM conjugation. Equivalently write the scalar as chi\_v(c)\^{}(-1), where chi\_v is inflated to O\^{}times. Since c=bar(pi) as a CM element, the actual scalar is

\[
                         s_{Q,0}=\chi_v(\bar\pi)^{-1}.
 \tag{R34.7}
\]

There is no corresponding first-frame change on the second vertical edge, so s\_(Q,1)=1 in these retained positive frames. The scalar is an explicitly specified root of unity, a unit at11. It can depend on the tuple through chi. It is not asserted to be1.

For comparison, the honest unbarred split Weil action is

\[
 \omega(h,1)\Phi(W)
   =\beta_0(\det h)^{-1}|\det h|^{b/2}\Phi(Wh),
                    \qquad\beta_0=\mu_0\xi.
 \tag{R34.8}
\]

At Q, beta\_0 on units is chi\_v, because the fixed mu\_0 and eta are unramified there. Thus omega(alpha,1)Phi\_Q=Phi\_Q: its honest character chi\_v(c)\^{}(-1) cancels the raw coordinate factor chi\_v(c). This does not remove R34.7. The canonical vector in the passive alpha model is T\_alpha(Phi\_Q)=chi\_v(c)T\_std(Phi\_Q), so the same actual moved vector still has coefficient chi\_v(c)\^{}(-1) in that model. In the literal \nolinkurl{VForm.vf} comparison it is this coefficient that belongs to w.sc(0). Applying an additional honest-Weil frame factor to R34.6 would count the same comparison twice.

\subsection{3. Equality with the full horizontal tensor}\label{equality-with-the-full-horizontal-tensor}

Write X=(A;C) in M\_(b,a)(F), y=(z;w) in F\^{}b, and let H\_Q(Z) be the unchanged second vertical function on M\_(N,b)(F). In the core model, Phi\_Q({[}X,y{]}) is exactly LIFT23(Q5):

\[
 F_X(y)=\mathbf1_{A\in M_{r,a}(O),\ C\in qM_{t,a}(O)}
        \mathbf1_{z\in O^r,\ w\in qO^t}
                       \chi(\det[\bar A,\bar z]).
 \tag{R34.9}
\]

For the prescribed depth M\textgreater=1 put Y=q\^{}(-M)O\^{}b/q\^{}M O\^{}b. The accepted source \nolinkurl{LastLoc} evaluates the line exchange in compatible Heisenberg models by

\[
 H_Q(Z)\int_{F^b}F_X(y)\psi_k(-\eta y)\,dy,
                         \qquad\eta=t/\pi.
 \tag{R34.10}
\]

Both the entire raw sum and every support condition have already been proved in LIFT23(Q5--Q10). In precisely its normalization, R34.10 equals

\[
 H_Q(Z)q^{-Mb}\mathbf1_{\eta\in q^{-M}O^b}
              \sum_{y\in Y}F_X(y)\psi_k(-\eta y).
 \tag{R34.11}
\]

The q\^{}(-Mb) is the volume of each coset, not the reciprocal of the number of Y classes. The fixed conductor in R34.1 determines the support of the transform. The source terminal row pi inverse has not been discarded just because pi is a unit. No new proof of LIFT23's Gauss-sum evaluation is claimed here.

Let S\_v denote the actual source \nolinkurl{lastBi} line exchange, let v\_0,v\_1 be its original vertical basic vacua, and set C\_(Q,0)=S\_v(v\_0 tensor v\_1). For each allowed y let d\_0(y),d\_1(y) be the actual finite Heisenberg operators representing the two horizontal raw functions in R34.11 relative to their basic vacua. Then the equality of complete tensors is

\[
 \boxed{
 S_v\bigl(T_{\rm std}(\Phi_Q)\otimes T_1(H_Q)\bigr)
  =\chi_v(\bar\pi)^{-1}
       \sum_{y\in Y}(d_0(y)\otimes d_1(y))C_{Q,0}.}
 \tag{R34.12}
\]

\textbf{Proof.} Use R34.6 in the first slot. Apply \nolinkurl{LastLoc} to the unchanged models and frames R34.2--R34.3. Its linear two-variable model sends the transform of the two integral indicators to the actual C\_(Q,0). With conductor O and volume one, that transform is the product horizontal integral indicator, since pi is a unit. The source \nolinkurl{IsCombAt}/\nolinkurl{IsBiModel} intertwining therefore sends each product summand of R34.11 to precisely the displayed (d\_0 tensor d\_1)C\_(Q,0). Linearity gives the finite sum. This also proves equality for zero values outside the supports. If the independently fixed horizontal and vertical local realizations differ from a Fourier convention by a fixed scalar, that scalar is already in C\_(Q,0) on both sides and cancels from this argument. No evaluation or removal of that model scalar is needed. QED.

This equality is before any theta summation, terminal restriction or global period. K038(CA2--CA11) supplies its unchanged honest character covariance under the actual source, middle and terminal class-frame transports. It retains the entire Y sum; choosing one representative without its orbit weight would be a different tensor.

\subsection{4. Exact orbit indices and the limit of the assertion}\label{exact-orbit-indices-and-the-limit-of-the-assertion}

The actual middle and terminal Q groups are the inverse images of P\_(r,t) and P\_(r,2t). Put kappa\_j(q)=product\_(i=1)\^{}j(1-q\^{}(-i)). Their indices in the respective maximal compacts are

\[
 I_1={b\brack r}_q
       =q^{rt}\frac{\kappa_b}{\kappa_r\kappa_t},\qquad
 I_2={N\brack r}_q
       =q^{2rt}\frac{\kappa_N}{\kappa_r\kappa_{2t}}.
 \tag{R34.13}
\]

This follows by counting full-rank r-frames and dividing by GL\_r(k): the standard parabolic is the stabilizer of an r-plane. The source \nolinkurl{QAt.wt_pen/wt_top} and \nolinkurl{QFam.Pins.cQ} consequently give

\[
 w_1(\mathcal O)=I_1|\mathcal O|,\quad w_2=I_2,\quad
 c_{Q,v}=(I_1I_2)^{-1}
       =q^{-3rt}\frac{\kappa_r^2\kappa_t\kappa_{2t}}
                        {\kappa_b\kappa_N}.
 \tag{R34.14}
\]

If an integrated summand is constant on a middle orbit, the product c\_(Q,v) w\_2 w\_1 is exactly its orbit cardinality. The unchanged honest covariance K038 proves the required invariance whenever the initial and terminal functions have their stipulated characters; the paired middle characters cancel, and eta*y is unchanged. This explains precisely how representatives reproduce the full sum in R34.12.

The initial weight w\_0 and initial trace epsilon\_0 are retained. In the complete source \nolinkurl{QLast} equation they occur on both sides, and \nolinkurl{LastStepsKernelEqSumCouplingCosets.kernel_eq_sum} inserts the specified epsilon\_0 before replacing the kernel by the full Y sum. This report does not replace that trace by an identity operator or assert its remaining global row/type hypotheses. Thus R34.12 and R34.14 close the actual local Q tensor and frame-scalar comparison; they do not alone instantiate every field of \nolinkurl{QAt}, \nolinkurl{QFam.Pins}, or the mixed arithmetic D7 host.

The Q endpoint GGP factor1 of K056 is a different quantity. It neither cancels R34.14 nor removes the Q factors in the vertical Rallis metric. The controlled global restricted row, its detector, integral preimage and complete DP44 comparison remain open.

\subsection{5. Exact sources and check scope}\label{exact-sources-and-check-scope-1}

All source references use pinned commit 7cbf7a0c32fca1badf05931879a37c3f8a2b03b4; the archive is accepted as stipulated and no Lean build or validity audit is performed.

\begin{itemize}
\tightlist
\item
  \nolinkurl{Definitions/LatticeFrames/ExistsUnitarySupportedAtPlaceOfMulEqIdem.lean}, Place.IsChar43--44, up/down49--54, IsModel60--70, IsCoreFrame75--76.
\item
  \nolinkurl{Definitions/LatticeFrames/IsCombAt.lean}, IsCombAt24--27, LastLoc49--66, including its quantification over the actual models.
\item
  \nolinkurl{Formalism/LocalToGlobal/IndefiniteReal/DatumVec.lean}, frMapO23--25, omegaQ66--72, IsStepVecOf84--96.
\item
  \nolinkurl{Definitions/LatticeFrames/PlaceLaws.lean}, VForm.vf133--135, QLast240--252 and CVEq267--268.
\item
  \nolinkurl{Definitions/LatticeFrames/CosInj.lean}, QAt23--118, especially wt\_mid/wt\_pen/wt\_top113--118, and QFam.Pins129--150.
\item
  \nolinkurl{LastStepsKernelEqSumCouplingCosets.lean}, full kernel\_eq\_sum statement and its full-Y, same-basic-tensor conclusion.
\item
  LIFT23(Q5--Q10), LIFT30(NL2--NL8), LIFT31(VA4--VA5, section2), and K038(CA2--CA11) are used in their checked scopes.
\end{itemize}

The new proof is the actual passive-model normalization and its compatibility with the unchanged vertical Q vector. Its arithmetic scalar R34.7 is distinct from an independently fixed local Fourier phase. Numerical diagnostics, whole-proof readings and the independent review are recorded separately; their scopes do not include an automorphic value or a complete QLast host.

\clearpage

\section{AM-DO-052 --- the actual generic G projection and scaled source frame}\label{am-do-052-the-actual-generic-g-projection-and-scaled-source-frame}

Author: GPT-6 Astra Pro, ALP\_MACHINE /root/definite\_oscillator. Reasoning effort is inherited and is not separately exposed. 22 November 2026 UTC. No Lean execution.

\textbf{Status: proof complete; author whole-proof check recorded below.}

This proves a precise comparison between the uncut vertical kernels of LIFT31/R032 and the canonical generic functions required by the source \nolinkurl{GLoc.eg_is} and \nolinkurl{egT_is}. It also checks the passive Q-frame scalar. The comparison uses the actual generic character subspace, the original diagonal G scaling, and the actual finite Haar measures. It does not identify the uncut identity-cell function with a generic function.

\subsection{1. Fixed conventions and the source distinction}\label{fixed-conventions-and-the-source-distinction}

At a selected G place write \(F=E_P=F_v\), \(O=O_F\), \(k=\mathbf F_q\), and use the original rational uniformizer \(q\). The accepted primes are absolutely split, so this is genuinely \(\mathbf Q_q\). Put

\[
 r=a+1,\qquad b=r+t,\qquad N=r+2t,
 \qquad P=\operatorname{diag}(U_a(k),1)\subset GL_r(k),
 \qquad n_P=|P|=q^{a(a-1)/2}.
 \tag{GP1}
\]

Here \(U_a\) is upper unitriangular and \(\psi_P(u)=\psi_q(q^{-1}\sum_{i<a-1}u_{i,i+1})\), with precisely the fixed residue additive character. The unbarred downward oscillator is written

\[
 [\omega_j(g,h)f](W)=\alpha_j(\det h)\beta_j(\det g)^{-1}
 |\det g|^{m_j/2}|\det h|^{-n_j/2}f(h^{-1}Wg).
 \tag{GP2}
\]

Thus on the first edge \(n_0=r,m_0=b,\ \beta_0=\mu_0\xi\). These are the conjugates of LIFT31(VA5,VA13), which uses the upward conjugated kernel. There is no ambiguity from determinant characters on \(P\), since every member has determinant one.

Let \(\delta_j\) denote the two actual \emph{unbarred} identity-cell kernels. Their conjugates are LIFT31(VA11--12). Explicitly,

\[
 \delta_0(W)=\mathbf1_{W\in M_{b,r}(O),\ \bar W=[I_r;0]},
 \tag{GP3}
\]

and \(\delta_1\) is supported on integral matrices with reduction \(\left(\begin{smallmatrix}I_r&B\\0&D\\0&0\end{smallmatrix}\right)\), where \(D\in GL_t(k)\), with the unchanged unbarred determinant factor \(\chi_1(\det D)\chi_2(\det D)^{-1}\).

The source requires a different scalar function before frame transport: replace \(I_r\) in these descriptions by an element \(u\in P\), multiply the value by \(\psi_P(u)\), and call the resulting functions \(\phi_{P,0},\phi_{P,1}\). These are exactly \nolinkurl{omegaG} for the generic leading subgroup: \nolinkurl{DatumVec.lean49–63} includes the leading group and its character, the old diagonal-block characters and the final inverse determinant character. No averaging denominator occurs in \nolinkurl{omegaG}.

\nolinkurl{HzUnit.lean66–70} asks \nolinkurl{eg_is} and \nolinkurl{egT_is} to realize the corresponding \nolinkurl{SymData.Phi} and \nolinkurl{PhiTop}. The uncut functions (GP3) are not equal to those functions. All subsequent comparisons retain this distinction.

\subsection{2. Exact averaging on the actual character subspace}\label{exact-averaging-on-the-actual-character-subspace}

Choose arbitrary integral lifts of the finite groups. The kernels are fixed by the principal kernel, so every expression below is independent of those choices. For \(u\in P\), put \(k_0(u)=u\) on the first edge and \(k_1(u)=\operatorname{diag}(u,I_t)\) on the second. Then

\[
 \boxed{\quad
 \phi_{P,j}=\sum_{u\in P}\psi_P(u)^{-1}
                 \omega_j(k_j(u),1)\delta_j.
 \quad}                                                     \tag{GP4}
\]

\textbf{Proof.} For the first edge, if \(\bar W=[v;0]\) with \(v\in P\), then \(\delta_0(Wk_0(u))\ne0\) precisely for \(u=v^{-1}\). The unique summand has coefficient \(\psi_P(v^{-1})^{-1}=\psi_P(v)\). If no such \(v\) exists every summand is zero. Integral support is preserved. On the second edge the first \(r\) columns are multiplied by \(u\), whereas the old last \(t\) columns, their diagonal block \(D\), and its determinant character are unchanged. Again the unique possible summand is \(u=v^{-1}\). Formula (GP2) contributes no character or modulus for these determinant-one integral matrices. This proves (GP4) on the entire Schwartz space, including its zero set. ∎

Let \(\rho_j\) be the actual right finite-group actions of R032 and define

\[
 E_{P,j}=\frac1{n_P}\sum_{u\in P}\psi_P(u)^{-1}\rho_j(u).
 \tag{GP5}
\]

It is the orthogonal projector onto the right \(\psi_P\)-character subspace. For completeness, replacing \(u\) by \(h^{-1}u\) in the sum gives \(\rho_j(h)E_{P,j}=\psi_P(h)E_{P,j}\). It is the identity on that character subspace. The involution in the unitary group algebra sends each coefficient at \(u\) to the same coefficient at \(u^{-1}\), so it is self-adjoint. These statements imply \(E_{P,j}^2=E_{P,j}\).

Let \(D_j\) denote the actual downward map with kernel \(\delta_j\), including the fixed real intertwiner, Haar measures and all other places. Right covariance of the theta sum and (GP4) give

\[
 D_{\phi_{P,j}}=n_P E_{P,j}D_j.
 \tag{GP6}
\]

This is an identity of maps, not a proportionality inferred from multiplicity one. R032 proves, for these actual uncut kernels, \(E_{P,j}D_j=D_jE_{P,j+1}\). Therefore the normalized physical tests

\[
                \delta_j^P=\phi_{P,j}/n_P                  \tag{GP7}
\]

satisfy

\[
 D_{\delta_j^P}=E_{P,j}D_j=D_jE_{P,j+1},\qquad
 D_{\delta_j^P}|_{E_{P,j+1}\mathcal C_{j+1}}
                    =D_j|_{E_{P,j+1}\mathcal C_{j+1}}.
 \tag{GP8}
\]

In particular, the two-edge composition is

\[
 D_{\delta_0^P}D_{\delta_1^P}
   =E_{P,0}D_0E_{P,1}D_1
   =D_0D_1E_{P,2}.                                         \tag{GP9}
\]

The two normalized averages collapse to one projector under middle integration. They do not create a factor \(n_P^{-2}\) in the actual map on the prescribed character subspace. If the \textbf{unnormalized} canonical functions \(\phi_{P,j}\) are used, the two-edge operator on that same subspace is instead \(n_P^2D_0D_1\).

The same formulas hold simultaneously at all G places by taking products of commuting projectors. The genuine central idempotent commutes with these projectors. Hence R032's same-projector inverse restricts to the same generic cut, and LIFT38's least-preimage norm bound is unchanged when the normalized tests (GP7) are used. This is an equality of actual maps and metrics; it does not assert an integral right inverse. The new averaging denominators are powers of the G residue primes and consequently are 11-units.

\subsection{3. The required diagonal basis and its honest Weil factor}\label{the-required-diagonal-basis-and-its-honest-weil-factor}

The source \nolinkurl{StandardOrderExtendedBasisApplyLowBlock.lean617–618} defines

\[
 \operatorname{typeFrame}(q).\operatorname{bas}
      =\operatorname{diag}(I_n,q,q^2,\ldots,q^c).
 \tag{GP10}
\]

\nolinkurl{LatFrame.coords/transport}, lines496--500, reads a matrix by \(\operatorname{bas}_{\rm target}^{-1}W
\operatorname{bas}_{\rm source}\). \nolinkurl{typeChain}, lines990--992, uses these frames. With \(m=2\), its first source has \(c=1\), while the middle and terminal frames have \(c=0\). Thus, in standard ordered coordinates and with \(d=\operatorname{diag}(I_a,q)\),

\[
 \Phi_0(X,y)=\phi_{P,0}([X,qy]),\qquad
 \Phi_1(Z)=\phi_{P,1}(Z).                                  \tag{GP11}
\]

The first function is therefore exactly the LIFT22 sharp function

\[
 f_{U_a}(X)\,\mathbf1_{q^{-1}e_a+O^b}(y),                  \tag{GP12}
\]

with the extra leading coordinate indexed by \(a\) in zero-based notation. The second is its prescribed canonical \nolinkurl{PhiTop}. There is no Fourier or Jacobian factor in (GP11): it is a statement about the actual scalar functions.

Let \(\alpha=g\operatorname{diag}(I_a,c)\), \(c=\bar\pi_T\), be the actual core frame of LIFT30/DO051. At G the fixed congruences give \(\alpha\equiv I\pmod{q^2}\). In the physical standard source frame \(\beta_{Y,0}=I\), the source \nolinkurl{MovedTo} clause requires

\[
 b_0=\alpha d,\qquad b_1=b_T=1.                            \tag{GP13}
\]

More generally it is the exact equation \(b_0\beta_{Y,0}=\alpha d\), with the already fixed integral frame included. This is \nolinkurl{IsCombAt.lean150–153,196–204}. A unit core frame alone cannot provide the factor \(d\).

Fix the passive model \(T_\alpha(\Phi)=T_{\rm std}(W\mapsto\Phi(W\alpha))\). This is the actual source model: substitution in \nolinkurl{Place.IsModel.trans} and \nolinkurl{.modul} verifies it, and since \(\alpha\in GL_r(O)\), it sends the standard characteristic function to the same basic vector. By (GP2), (GP7), (GP11), and (GP13),

\[
 \boxed{\quad
 \omega_0(b_0,1)\delta_0^P
  =n_P^{-1}q^{-b/2}\beta_0(q)^{-1}T_\alpha(\Phi_0),\qquad
 \delta_1^P=n_P^{-1}T_{\rm std}(\Phi_1).
 \quad}                                                     \tag{GP14}
\]

Here equality includes the physical model maps on the left. The downward source character is \(\beta_0=\mu_0\xi\). The factor from \(c=\det\alpha\) is one because it is congruent to one modulo \(q^2\) and these characters are tame. The forward comparison has \(q^{-b/2}\); \(q^{b/2}\) is its inverse, not an additional factor.

Thus the G contribution to the literal \nolinkurl{VForm.vf/vfT} scalars for these normalized physical tests is

\[
 s_{G,0}=n_P^{-1}q^{-b/2}(\mu_0\xi)(q)^{-1},\qquad
 s_{G,1}=n_P^{-1}.                                        \tag{GP15}
\]

The \nolinkurl{eg}/\nolinkurl{egT} Heisenberg combinations continue to realize the unnormalized canonical functions \(\Phi_0,\Phi_1\). The denominator in (GP15) belongs to \nolinkurl{w.sc}, and is not inserted into their raw \nolinkurl{IsCombAt} functions. This gives the exact finite-function part of \nolinkurl{VForm} for the constructed normalized realization. A statement that the normalized physical functions themselves satisfy the unscaled \nolinkurl{IsStepVecOf} clause for the generic subgroup would be false: that clause prescribes the unnormalized \nolinkurl{omegaG}. One can instead use that literal raw datum and retain the corresponding \(n_P\) in each map. The two descriptions are linked by (GP6)--(GP9); they cannot be mixed.

As a measure check, the disjoint residue cells give \(\|\delta_j^P\|^2=\|\delta_j\|^2/n_P\). For the first edge, \(\|\delta_0\|^2=q^{-br}\) and \(\|\Phi_0\|^2=n_Pq^{-ab}\). Multiplying the latter by the squared absolute factor in (GP14) gives \(q^{-br}/n_P\), as required. There is no hidden square root of a group order.

For the two actual hosts \(a=42,168\), \(b=1802\), the combined absolute scalar per G place is respectively \(q^{-2623}\) and \(q^{-28957}\). These are comparison scalars, not a loss in the unchanged maps (GP9).

\subsection{4. The scaled evaluation retains the old local branching period}\label{the-scaled-evaluation-retains-the-old-local-branching-period}

The diagonal matrix in (GP10) has the factorization

\[
 d=qI_r\cdot\operatorname{diag}(q^{-1}I_a,1)\in Z_{GL_r}H,
 \qquad H=GL_a.
 \tag{GP16}
\]

Also \(d^{-1}\alpha d\in K_r^1\): for a matrix in \(q^2M_r(O)\), conjugation by \(d\) changes the valuation of an entry by at most one. Consequently any old \(K_r^1\)-fixed bottom vector \(F\) satisfies

\[
                         F(h\alpha d)=F(hd).               \tag{GP17}
\]

Let \(\omega_B,\omega_H\) be the actual unitary central characters of a bottom vector \(F\) and an initial vector \(f\), respectively. For the bilinear local period, whenever the integral or its fixed canonical continuation is defined,

\[
 \int_H f(h)F(\operatorname{diag}(h,1)\alpha d)\,dh
 =\omega_B(q)\omega_H(q)
                  \int_H f(h)F(\operatorname{diag}(h,1))\,dh.
 \tag{GP18}
\]

\textbf{Proof.} Use (GP17), then (GP16), and substitute \(h'=hq^{-1}I_a\). Both left and right Haar measures on the reductive local group agree; right translation has no modulus. Central equivariance gives the two displayed characters. The proof first holds in any absolutely convergent realization, and the exact covariance continues with the period when that is its stipulated definition. For the actual tempered local tests of K055 their canonical period has this equivariance directly. ∎

The scalar has absolute value one. It does not destroy a nonzero K055 G-period or alter its Hilbert-normalized absolute square. Equation (GP18) is a statement about the period and not a replacement of the full fixed Schwartz tensor by an unscaled tensor.

\subsection{5. Independent verification of the passive Q scalar}\label{independent-verification-of-the-passive-q-scalar}

At Q keep the original physical vertical vectors and standard Schrödinger model. Here \(\alpha=g\operatorname{diag}(I_a,c)\in GL_r(O)\) with \(\det g=1\) and \(c=\bar\pi_T\in O^\times\). The raw \textbf{unbarred} first Q function is the integral top-r determinant function with residue character \(\chi_Q\), as in \nolinkurl{DatumVec.lean65–71,84–98}. Its support is preserved by right \(\alpha\) and

\[
 \Phi_Q(W\alpha)=\chi_Q(c)\Phi_Q(W).
 \tag{GP19}
\]

Indeed, the lower-zero condition is invariant under \(GL_r(O)\), the top-r determinant is multiplied by \(\det\alpha=c\), and no other columns occur on this first edge. Therefore, in the passive model fixed before (GP14),

\[
 T_{\rm std}(\Phi_Q)=\chi_Q(\bar\pi_T)^{-1}T_\alpha(\Phi_Q),
 \qquad T_\alpha(\mathbf1)=T_{\rm std}(\mathbf1).
 \tag{GP20}
\]

This is the exact scalar contribution to \nolinkurl{w.sc 0} when the physical vector is retained and only its coordinate model is changed. There is no honest Weil action in this argument. Consequently the honest splitting determinant factor in K038(CA5) must \textbf{not} be multiplied into (GP20) as an additional cancellation. If an honest Weil operator is used as a different method of transporting the model, its action on the basic vector must also be included; after fixing the same basic vector the result is again (GP20).

The terminal core coordinate is the prescribed \(\eta=t/\pi_T\). \nolinkurl{LastLoc} fixes \(\psi_k(x)=\psi_E(\xi_0\iota_v(x))\), so its Fourier character is exactly \(\psi_k(-\eta y)=\psi_E(-\xi_0\iota_v(\eta y))\). These facts are unaffected by the passive first-source substitution. The entire Q coset sum, its \(q^{-Mb}\) coefficient, and its orbit weights remain the ones in LIFT23/DO037 and the source \nolinkurl{CosInj}/\nolinkurl{QLast}.

\subsection{6. Precisely what is established}\label{precisely-what-is-established}

The canonical G functions, their generic projection factors, the required nonintegral source basis, the honest Weil modulus/character, and the passive Q scalar are all explicit. On the actual generic subspace, the normalized finite realization preserves the original vertical map and its least-preimage norm. The scaled G evaluation preserves the old nonzero local branching test up to an explicit unit.

These facts supply the finite-function comparison needed for the G/Q part of a constructed \nolinkurl{VForm}. They do not themselves instantiate the remaining global source record, identify the normalized functions with an unscaled raw datum, evaluate DO051's fixed source-model phase, or prove a nonzero global central Rankin value. The common finite tensor must still be the actual lastBi tensor with all these fixed comparisons; equality of a local branching period alone is insufficient.

\subsection{7. Verification record}\label{verification-record-6}

The accompanying certificate checks the finite character projector, both generic-kernel formulas, the actual type-frame substitution, the scaled-lattice congruence, the passive Q determinant character, and exact norm/exponent identities in finite and rational models. It does not evaluate an automorphic integral or run Lean. The author whole-proof check and replay are recorded in the agent check ledger.

\subsubsection{Exact retained source locations}\label{exact-retained-source-locations}

All Lean bytes are from the supplied archive/pinned commit \nolinkurl{7cbf7a0c32fca1badf05931879a37c3f8a2b03b4}, with intake SHA matches.

\begin{itemize}
\tightlist
\item
  \nolinkurl{Formalism/LocalToGlobal/IndefiniteReal/DatumVec.lean49–98}: raw \nolinkurl{omegaG}, \nolinkurl{omegaQ}, and the unscaled \nolinkurl{IsStepVecOf} prescription.
\item
  \nolinkurl{Definitions/LatticeFrames/HzUnit.lean32–82,93–111}: G frames, \nolinkurl{eg_is}, \nolinkurl{egT_is}, and actual local splitting characters.
\item
  \nolinkurl{Definitions/LatticeFrames/IsCombAt.lean24–27,49–66,150–153,196–204,   236–241}: fixed-model combinations, lastBi Fourier formula, \nolinkurl{MovedTo}, the required bas transport and actual moved vector.
\item
  \nolinkurl{Definitions/LatticeFrames/ExistsUnitarySupportedAtPlaceOfMulEqIdem.lean   43–44,49–70,75–76}: source additive character, exact up/down and passive Heisenberg-model laws, core last column.
\item
  \nolinkurl{Definitions/TwistedInnerProducts/StandardOrderExtendedBasisApplyLowBlock.lean   179–184,281–292,496–500,617–618,954–962,990–992}: character, raw type function, coordinate transport and scaled chain.
\item
  \nolinkurl{Proofs/DecoratedLatticeProofs/DescentDataStepExponents.lean75–142}: the exact moved scalar-function comparison; its scalar is one only after the stated moved-frame and generic-level hypotheses.
\item
  \nolinkurl{Definitions/LatticeFrames/PlaceLaws.lean128–147}: literal \nolinkurl{VForm}.
\end{itemize}

These paths after \nolinkurl{Definitions/} or \nolinkurl{Proofs/} are relative to \nolinkurl{Development/AbelianVarietyTheory/} and are recovered under this agent's \nolinkurl{definite_oscillator_seesaw_source} prefix where absent centrally.

\clearpage

\section{ALP-LIFT-42 --- exact finite-character suppression of short packet correspondences}\label{alp-lift-42-exact-finite-character-suppression-of-short-packet-correspondences}

Author: GPT-6 Astra Pro, ALP\_MACHINE \nolinkurl{/root/lift_accounting}. 22 November 2026 UTC; inherited reasoning effort is not separately exposed.

\textbf{Author whole-proof review PASS, 2026-11-22 07:12:50 UTC.} Independent whole-proof review is pending. Exact diagnostic certificate: 1,853 assertions PASS, 2026-11-22 07:12:11 UTC.

This result supplies the finite-character and arithmetic packing constants for a projected Poincare argument at the actual LIFT41 packet point. It does not itself prove positivity of that projected reproducing kernel, an actual initial Mram period, or terminal-image nonvanishing. The selected fibre and finite-stabilizer positivity argument is being treated separately by \nolinkurl{/root/kummer_extension}.

\subsection{1. Actual groups, finite representations, and the quantity bounded}\label{actual-groups-finite-representations-and-the-quantity-bounded}

Use either actual K031 or K032 host, with

\[
 (a,f,n)=(42,42,43)\quad\hbox{or}\quad(168,168,169),
 \qquad Q=g_0g_1.
 \tag{FC1}
\]

The two rational G primes are completely split in the CM field E of degree2f. There are f selected relative places over each prime. The genuine goodTuple representation \(\Pi\) of their product of \(\mathrm{GL}_n\) groups is cuspidal at every place except one over each rational prime; at those two places it is the actual irreducible parabolic product of degrees n−6 and6. These are precisely the fixed DO042/DO043 choices, including the allowed aggregate central character. Their central-character adjustment does not change any bound below.

Let \(P_0,P_1\) be arbitrary nonzero orthogonal projectors on the two mixed factors, and let

\[
 \mathcal P=I_{\rm cusp}\otimes P_0\otimes P_1,\qquad
 b(h)=\frac{\operatorname{Tr}(\mathcal P\Pi(\operatorname{red}h))}
                  {\operatorname{Tr}\mathcal P}.
 \tag{FC2}
\]

In particular the projectors may be the two canonical highest-mirabolic summand projectors of LIFT32. No character formula for a mixed factor is assumed. The elementary inequality \(\left|\operatorname{Tr}(PU)\right|\le\operatorname{rank}P\) for a unitary U and orthogonal P bounds each of its normalized partial traces by1. At a scalar, its partial trace equals the same central character as the full irreducible representation.

Let

\[
 \Gamma=U(I_a,-1)(F)\cap J^{\max},\qquad
 Z=\{zI:z\in\mu_{28}\},\qquad
 S=\operatorname{Stab}_{\Gamma}(o).
 \tag{FC3}
\]

All coordinates here are the actual standard integral frames of R025. Thus every matrix entry of an element of Gamma is in \(\mathcal O_E\), and the selected and conjugate split reductions are tied by the actual unitary relation. The group Z is contained in S. Write \(u_\tau(h)\) for its nonnegative Cartan parameter at label tau.

For B\textgreater0 put

\[
 N(B)=\#\{h\in\Gamma:\ u_\tau(h)\le B\text{ for all }\tau\},
 \qquad
 \epsilon(B)=\sup_{h\in\Gamma\setminus Z,\ u_\tau(h)\le B}|b(h)|.
 \tag{FC4}
\]

Counting modulo Z only makes N smaller. The estimate includes every finite origin-stabilizer term outside Z; none is deleted by a neatness claim or by assuming that the packet stabilizer is scalar.

\subsection{2. A sharp uniform cuspidal character ratio}\label{a-sharp-uniform-cuspidal-character-ratio}

For n\textgreater=3 and a prime power q\textgreater=n+1, every irreducible cuspidal character chi of \(\mathrm{GL}_n(\mathbb F_q)\) satisfies

\[
 \frac{|\chi(g)|}{\chi(1)}\le\frac1{q^{n-1}-1}
       \qquad(g\text{ not scalar}).
 \tag{FC5}
\]

We cite Green's character formula and prove this elementary consequence. Green(1955), Section5, Example(ii), printed pp.430--431, together with the regular-label irreducibility theorem, gives the formula. Its explicit Jordan form is also stated as Theorem2.1, equation(2.1), in Gorodetsky and Hazan(2017), p.5. If the semisimple part does not come from the degree-n torus, the character is zero. Otherwise let d divide n be the degree of its eigenvalue, and let t be the number of Jordan blocks for one eigenvalue. Then

\[
 |\chi(g)|\le d\prod_{j=1}^{t-1}(q^{dj}-1),\qquad
 \chi(1)=\prod_{j=1}^{n-1}(q^j-1),\qquad t\le n/d.
 \tag{FC6}
\]

For d=1 and a nonscalar element, t\textless=n−1, so cancellation immediately gives FC5. For d\textgreater=2, cancel all factors whose indices are multiples of d, up to n−d.~The denominator left over contains both \(q^{n-1}-1\) and \(q-1\): the indices n−1 and1 are distinct and are not divisible by d.~Since q−1\textgreater=n\textgreater=d, the latter factor absorbs the numerator d.~Every remaining factor is at least1. This proves FC5, including t\textless n/d.~For a scalar times a transvection the bound is an equality, so the exponent n−1 is not an estimate obtained from a weaker general character theorem.

\subsection{3. An algebraic norm bounds the scalar-reduction places}\label{an-algebraic-norm-bounds-the-scalar-reduction-places}

Let h be in \(\Gamma\setminus Z\), with all Cartan parameters at most B. There is a nonzero algebraic integer x that is either an off-diagonal entry of h or the difference of two diagonal entries. Otherwise h is a scalar integral unitary matrix, hence is in Z by the actual norm-one unit calculation carried in R027/K032.

The Euclidean operator norm of h at each CM embedding is at most \(e^B\). Therefore

\[
          |\sigma(x)|\le2e^B\quad\text{at all embeddings }\sigma.
 \tag{FC7}
\]

Suppose h has scalar reduction at a selected relative place v over one of the G primes. Its conjugate reduction is scalar too, because the two components are inverse adjoints. Thus x vanishes at both degree-one CM prime ideals over v. They are distinct and unramified, and each has norm q\_v. Taking the absolute field norm of x gives

\[
 \prod_{v:\ \operatorname{red}_v(h)\text{ scalar}}q_v^2
       \le|N_{E/\mathbb Q}x|\le(2e^B)^{2f},
 \quad
 \sum_{v:\ \operatorname{red}_v(h)\text{ scalar}}\log q_v
       \le f(B+\log2).
 \tag{FC8}
\]

The two mixed places together have log-weight log Q. Hence the total log-weight of nonscalar reductions at cuspidal places is at least \((f-1)\log Q-f(B+\log2)\). Apply FC5 at those places and the partial-trace bound at the mixed places. Since \(2f/g_{\min}^a<1/2\) for the actual integers, the elementary product bound \(\prod(1-x_i)\ge1-\sum x_i\) proves

\[
 \epsilon(B)<
 2\exp\{-a((f-1)\log Q-f(B+\log2))\}.
 \tag{FC9}
\]

This bound is uniform in all regular cuspidal parameters and in both nonzero mixed projectors. It retains the loss from dropping exactly two places. It also applies if one normalizes the actual trace by a smaller-rank mixed mirabolic projector, as in FC2.

\subsection{4. Coarse orbit packing with the complete finite stabilizer}\label{coarse-orbit-packing-with-the-complete-finite-stabilizer}

For the standard origin,

\[
                   |S|\le28^{a+1}a!.
 \tag{FC10}
\]

Indeed an integral matrix fixing o is block diagonal with a positive unitary a-by-a block and one norm-one unit. Every conjugate of every entry of the positive block has modulus at most1. A nonzero such entry has nonzero integral norm of absolute value at most1; equality forces every conjugate modulus to be1. The norm-one condition on a column then leaves exactly one nonzero entry. These entries are roots of unity: all powers have conjugates bounded by1, so the integer coefficients of their characteristic polynomials are uniformly bounded; only finitely many such algebraic integers of the fixed degree exist, and two powers coincide. The actual roots of unity are mu28. Thus the positive block is monomial, proving FC10. A determinant or component restriction can reduce S, which only strengthens the bound.

Take r=2/5. Distinct points of \(\Gamma o\) have disjoint open product balls of radius r. To prove this, if all Cartan parameters of \(h\in\Gamma\) were at most2r=4/5, every positive entry of its last column would have every conjugate bounded by \(\sinh(4/5)<1\). Such an algebraic integer must be zero, and h fixes o. The strict elementary inequality follows from

\[
 e<11/4,\qquad 4\cdot11^4<9^5,
 \quad e^{4/5}<9/4,
 \quad\sinh(4/5)<\tfrac12(9/4-4/9)=65/72<1.
\]

The invariant ball measure, in the R030 convention, has exact volume \(\sinh^{2a}(u)/a!\) in one factor. This follows by integrating \(x^{a-1}(1-x)^{-a-1}/(a-1)!\) up to x=tanh²u. Comparing the disjoint balls with the containing product ball of radius B+r proves

\[
 N(B)\le |S|
       \left({\sinh(B+2/5)\over\sinh(2/5)}\right)^{2af}.
 \tag{FC11}
\]

This is an orbit-packing estimate, followed by the complete finite stabilizer factor. It is not a count of arbitrary integral matrix entries, and it does not remove S from the arithmetic group.

\subsection{5. Actual radii and exact numerical margins}\label{actual-radii-and-exact-numerical-margins}

At the actual LIFT41 point the real conjugating matrix has Cartan parameter s at every label, with \(2s<1/100\). Therefore every element of the near-centred \(C_R C_R^{-1}\) has standard Cartan parameters at most \(B=2R+1/100\). Use

\[
 \begin{array}{c|ccc}
  \text{host}&R&B&m\\ \hline
  J14&9/2&901/100&28\\
  J56&5&1001/100&31
 \end{array}.
 \tag{FC12}
\]

Here m is a convenient integer upper bound for 3B+4/5. The elementary inequalities \(\sinh(2/5)>2/5\), \(\sinh(B+2/5)<e^{B+2/5}/2\), and e\textless11/4, together with FC9--11, give the completely rational bound

\[
 N(B)\epsilon(B)<
 {2\,28^{a+1}a!\,(25/8)^{af}(11/4)^{maf}
       \over Q^{a(f-1)}}
 <\begin{cases}
       2^{-1200}&(J14),\\
       2^{-646000}&(J56).
   \end{cases}
 \tag{FC13}
\]

Both final comparisons are integer cross-multiplications in the companion certificate, not floating point logarithms. In particular the sum of the absolute normalized finite partial traces over all noncentral short correspondences is below the indicated bounds.

The radius9/2 on J14 is compatible with the accepted K046 concentration calculation. Its exact simultaneous moment bound is A\_tau\textless969. Replacing its cutoff epsilon=2\^{}-12 by epsilon=1/2000 leaves epsilon*A\_tau\textless1/2. K046(BL18--22), with the same t=ell/2 and s0=a+1, consequently gives a tail less than \(f2^{-\ell/8}\) at ell=ell\_+, since its retained inequality 4B0\textless ell dominates the unchanged constant terms. Finally e³\textgreater20, proved already by the first nine positive terms of its power series, implies \(\cosh(9/2)>\sqrt{2000}\). Thus Y\textgreater=1/2000 lies inside radius9/2. On J56 the original radius5 and original concentration bound are retained. This is a numerical refinement of the existing moment proof, not a new assertion about the genuine projector.

\subsection{6. A separate cardinal bound and the retained long-tail scale}\label{a-separate-cardinal-bound-and-the-retained-long-tail-scale}

The same packing inequality, without any character factor, also gives

\[
          N(B)<2^{32af}<2^{64af}.
 \tag{FC14}
\]

Indeed bound it by \(28^{a+1}a!(5/4)^{2af}(11/4)^{m_Naf}\), with m\_N=19 on J14 and21 on J56. These integers exceed \(2(B+2/5)\). Comparison with \(2^{32af}\) is again an exact integer inequality in the certificate. On both hosts, the actual ell\_+ also satisfies

\[
 \delta_0=f2^{-\ell_+/8}<2^{-64af-20}
                  <\frac1{8\sqrt{2N(B)}}.
 \tag{FC15}
\]

For the first inequality use f\textless=2\^{}f and the certified \(\ell_+>8(64af+20+f)\). The second follows already from the weaker cardinal bound \(N(B)<2^{64af}\). These inequalities retain the full stabilizer in N and supply numerical inputs for an L1-tail argument; they do not assert that a tail can be discarded from a projected Poincare series.

\subsection{7. Scope and sources}\label{scope-and-sources-2}

FC13 is a short-correspondence estimate with the complete arithmetic origin stabilizer, the two actual mixed finite factors, and the actual near-origin packet point all retained. To deduce a projected selected-jet detector one must still prove the relevant twisted Poincare or reproducing-kernel positivity statement, including its long-correspondence tail, vector-valued fibre and central character. No such assertion is folded into this finite estimate.

Primary character sources:

\begin{itemize}
\tightlist
\item
  Green, J. A. (1955). The characters of the finite general linear groups. \emph{Transactions of the American Mathematical Society, 80}(2), 402--447. Section5, Example(ii), pp.430--431, with Section7's regular-label irreducibility and Section8's character theorem. \url{https://www.ams.org/tran/1955-080-02/S0002-9947-1955-0072878-2/S0002-9947-1955-0072878-2.pdf}
\item
  Gorodetsky, O., \& Hazan, Z. (2017). \emph{On certain degenerate Whittaker models for cuspidal representations of GL\_(kn)(F\_q)}. arXiv:1707.07308v2. Theorem2.1 and equation(2.1), p.5, state the exact Jordan form of Green's formula used in FC6. \url{https://arxiv.org/pdf/1707.07308}
\end{itemize}

The actual field, finite root group, goodTuple degrees and frames are accepted K031/K032, DO042/DO043 and R025/R027. The selected moment and metric statements are accepted K043/K046. Their conclusions are used with the displayed restrictions, rather than proved again here.

\clearpage

\section{AM-DO-053 --- the literal initial Q trace, its unit mass, and the complete two-step Q law}\label{am-do-053-the-literal-initial-q-trace-its-unit-mass-and-the-complete-two-step-q-law}

Author: GPT-6 Astra Pro; ALP\_MACHINE /root/definite\_oscillator. 22 November 2026 UTC. Inherited reasoning effort is not separately exposed.

Status: author whole-proof check PASS, 22 November 2026 UTC. The complete proof, literal source clauses, and full companion certificate have been read. The exact finite certificate passes 5,325,427 assertions. Independent review is pending.

This finding resolves the initial Q trace left open in R034. In the actual two-step family the local trace is exactly the identity and its weight is exactly one, by the literal definitions. The terminal trace is a nontrivial index scalar and is retained. The allowed middle coupling classes form one orbit, but replacing its complete Schwartz sum by a weighted representative is justified only after the middle theta integration. No equality of those two unintegrated tensors is asserted.

The result concerns the actual mixed geometric family with the already fixed frames, Heisenberg models and honest splitting characters. It does not assert that the obstructed original HostRaw record exists on these fields, prove a nonzero restricted period, or remove any arithmetic frame or terminal pairing cost.

\subsection{1. Data, conventions, and the initial representative}\label{data-conventions-and-the-initial-representative}

Use either K031/K032 host, with \[
 a=42\ \text{or}\ 168,\qquad r=a+1,\quad b=r+t=1802,\quad N=r+2t.
 \tag{QT1}
\] Fix an actual auxiliary Q place of the tuple. Its local field is \(F=\mathbf Q_q\), its ring is \(O\), and the actual additive character has conductor \(O\) with \(\operatorname{vol}(O)=1\), as proved in R034. The conclusions below also hold over any nonarchimedean local field with uniformizer \(\varpi\), residue cardinality \(q\), and those same conductor/volume conventions after replacing every occurrence of the scalar uniformizer by \(\varpi\). For the literal source comparison here the actual degree-one tuple places and rational uniformizer \(q\) are retained.

Let \(\chi\) be the tuple residue character inflated to \(O^\times\). The unbarred raw first function and second function are the ones of R034 and DatumVec, not new test functions. Write their common tame parameter in the source twistChain as \(\delta|_{O^\times}=\chi^{-1}\), choosing \(\delta(q)=1\); parPhi uses \(\delta^{-1}\). The remaining unramified stage factors retain their actual values on \(q\). The local levels and their right characters are \[
\begin{array}{c|c|c}
 \text{stage}&J_j&\lambda_j\\ \hline
 0&K_r=\mathrm{GL}_r(O)&1\\
 1&\operatorname{red}^{-1}P_{(r,t)}&\chi(\det D)^{-1}\\
 2&\operatorname{red}^{-1}P_{(r,2t)}&\chi(\det D)^{-1}.
\end{array}                                                   \tag{QT2}
\] These are precisely the vertical spaces retained by LIFT31(VA10--13) and R032, with the same actual right-character convention.

For \(m=2\), QAt.hs forces \(s=0\), and its only epsilon clause has \(j=\rho=d=0\). The source Chain.IsRepSystem.zero states \[
                         A_0=\{1\}.                           \tag{QT3}
\] Since QAt.rep\_bij is a bijection from its finite index set onto \(A_0\), that index set is a singleton and its representative is exactly 1. This is stronger than merely choosing an equivalent double-coset representative.

The source initial twistChain level is the parahoric with blocks \((r,0)\). Its lower-left congruence condition is empty, and the determinant of its empty lower block is 1. Thus its group is \(K_r\) and its character is 1, proving the first line of QT2 directly from parahoric and parahoric\_lam. All permutations in the plain frame are integral and preserve this maximal compact and character.

\subsection{2. The exact IsOneComb witness and the weight}\label{the-exact-isonecomb-witness-and-the-weight}

The source defines toOne by \(P-1\) iterations of exDown and proves toOne\_one by reflexivity. Consequently, for every smooth representation \(\sigma\), every linear functional \(l\), and every vector \(v\), \[
 \operatorname{toOne}(\nu,\sigma,1,l)(v)=l(v)
      =\operatorname{comb}(\mathbf1_{\{1\}})
                           \bigl[g\longmapsto l(\sigma(g)v)\bigr].
 \tag{QT4}
\] In particular the finitely supported function \[
                         e_{Q,0}=\delta_1                    \tag{QT5}
\] is a literal IsOneComb witness for \(P=1,\alpha=1\). It works universally, before imposing smoothness, the level character, or N-equivariance; those hypotheses of IsOneComb cause no further restriction. Mapping it through the local adelic group homomorphism still gives \(\delta_1\), since a homomorphism sends 1 to 1.

Here is the accompanying mass calculation, without invoking a generic unfolding assertion. In SimpleSum the level function on \(a\oplus\operatorname{Fin}(1)\) has constant value 1 when \(p=1\). Every pair of indices satisfies the defining inequality for InN. Thus its only generator is the identity, and \[
 N(F,a,1,1)=\{1\},\qquad
 R(F,a,1,1)=\{\operatorname{diag}(h,1):h\in\mathrm{GL}_a(F)\}.
 \tag{QT6}
\] The source Cgrp and prX therefore give \[
\begin{split}
 C(K_r,1)&=\{\operatorname{diag}(h,1):h\in K_a\},\\
 \operatorname{prX}C(K_r,1)&=K_a .
\end{split}                                                    \tag{QT7}
\] Indeed a block diagonal matrix and its inverse are integral if and only if the \(a\)-block and its inverse are integral. No unipotent quotient remains to supply an extra volume.

QAt.haar normalizes the Haar measure of \(K_a\) to 1. Hence its literal wt\_mid formula gives \[
                 w_{Q,0}=1,\qquad \epsilon_{Q,0}=\delta_1.
 \tag{QT8}
\] This proves the exact initial scalar and operator, not just their action on one nonzero theta vector.

For several Q places, QFam.Pins.epsilon takes the product of these weights and the convolution of these finitely supported functions. Both products remain identities: \[
                    w_0=1,\qquad \epsilon_0=\delta_1 .
 \tag{QT9}
\] The genuine goodTuple idempotent of R032 is a different operator, supported at G. It is not being identified with epsilon. The identity QT9 commutes with that projector on every vector and every class module. The complete G residue actions, their generic cuts in DO052, and the same-projector least right inverse therefore retain their previous statements and norms.

There is a necessary global distinction. VForm.e maps \(g\) to \(g\,b_0\). It follows from QT9 that the actual moved initial trace in that record is \[
                             t.e_0=\delta_{b_0},               \tag{QT10}
\] not \(\delta_1\) unless \(b_0=1\). The G scaling and central period character of DO052 and the T translation of DO051 are not removed. This finding resolves the Q trace; it does not replace the complete moved trace by an identity.

\subsection{3. The exact coupling orbit}\label{the-exact-coupling-orbit}

Fix the actual depth \(M\ge1\). Put \[
 Y=q^{-M}O^b/q^MO^b,\qquad
 \mathcal O_Q=
 \{(z,w)\bmod q^MO^b:z\in O^r\setminus qO^r,\ w\in qO^t\}.
 \tag{QT11}
\] For \(X\in M_{b,a}(F)\), the first raw summand is \[
 f_y(X)=
 \mathbf1_{\,[X,y]\text{ integral},\ [X,y]_{\rm low}\in qM_{t,r}(O)}
                  \chi\bigl(\det\overline{[X,y]}_{\rm top}\bigr),
 \tag{QT12}
\] with the residue character extended by zero on singular matrices. This is T.xiPen with the unique representative 1. It is the unchanged raw first Q factor.

The source coupIdx is exactly \(\mathcal O_Q\). If \(f_y(X)\ne0\), the lower coordinates of \(y\) are in \(qO^t\), and the last column of an invertible top \(r\)-matrix is primitive. Conversely, every primitive \(z\) extends to an \(O\)-basis of \(O^r\). Use its other \(r-1=a\) columns for the top block of \(X\), and set the lower block of \(X\) to zero. Then \(f_y(X)\ne0\). This proves both inclusions, including the absence of any additional character restriction on \(y\).

The group \(J_1=\operatorname{red}^{-1}P_{(r,t)}\) acts transitively on \(\mathcal O_Q\). To prove this, choose a representative \(y=(z,w)\) with primitive \(z\). Choose \(A\in\mathrm{GL}_r(O)\) whose last column is \(z\), and a matrix \(C\in qM_{t,r}(O)\) whose last column is \(w\). Then \[
 \kappa=\begin{pmatrix}A&0\\C&I_t\end{pmatrix}\in J_1,\qquad
                         \kappa e_r=y.                        \tag{QT13}
\] All choices descend to the indicated finite quotient. Conversely \(J_1\) preserves the integral lattice, its \(q^M\)-sublattice, and the primitive-top/lower-zero residue conditions. Thus there is exactly one relevant orbit and \(e_r\) is a valid representative. Its exact cardinality is \[
                    h_M=(q^r-1)q^{b(M-1)}.                    \tag{QT14}
\] The top primitive choices number \(q^{Mr}-q^{(M-1)r}\), and the lower choices number \(q^{(M-1)t}\), which proves QT14.

One may consequently instantiate the QAt local index sets at stages 0, 1, 2 by singletons, with initial representative 1 and middle coset \(e_r+q^MO^b\). Choose \(c_1=q^{-M}\) and \(c_2=q^M\). Equation QT12 is supported in the first box and is periodic under the second, proving adapt. Transitivity proves cos\_cov; cos\_inj and the terminal one clause are immediate. This is an orbit-indexed realization of the law. It does not delete the full global ybSet sum in K038.

\subsection{4. Honest covariance and why the orbit can be compressed}\label{honest-covariance-and-why-the-orbit-can-be-compressed}

Let \(g_y(Z,\eta)\) be the second raw horizontal function. It is the unchanged second vertical Q function \(H_Q(Z)\) multiplied by the Fourier transform of the coset \(y+q^MO^b\): \[
 g_y(Z,\eta)=H_Q(Z)\,q^{-Mb}
             \mathbf1_{\eta\in q^{-M}O^b}\psi_k(-\eta y).
 \tag{QT15}
\] Here \(\eta\) is the actual terminal coordinate \(t_{\rm row}/\pi_T\) of R034. Neither its unit \(\pi_T^{-1}\) nor the xi0 in \(\psi_k\) is discarded.

For \(\kappa\in J_1\), write its reduction as \(\left(\begin{smallmatrix}A&B\\0&D\end{smallmatrix}\right)\). On \(J_1\), both horizontal honest splitting characters have the same unit restriction \(\chi\). Their unramified stage factors are 1. The actual Hom formulas give \[
\begin{split}
 \omega_0(1,\kappa)f_y&=\chi(\det D)\,f_{\kappa y},\\
 \omega_1(\kappa,1)g_y&=\chi(\det D)^{-1}\,g_{\kappa y}.
\end{split}                                                     \tag{QT16}
\]

For the first equality, the honest target factor is \(\chi(\det\kappa)\), whereas applying \(\kappa^{-1}\) to the top \(r\)-block of \([X,\kappa y]\) supplies \(\chi(\det A)^{-1}\). Their product is \(\chi(\det D)\). For the second equality, the honest source factor is \(\chi(\det\kappa)^{-1}\), the raw right multiplication of \(Z\) supplies \(\chi(\det A)\), and \(\psi_k(-\eta\kappa y)\) is precisely the Fourier phase for \(\kappa y\). The matrix \(\kappa\) preserves the dual support lattice, because it is integral with integral inverse. All absolute determinant factors are 1. These proofs apply outside the supports as well: each displayed support condition is preserved in both directions.

Therefore, as an equality of entire tensor vectors, \[
 f_{\kappa y}\otimes g_{\kappa y}
   =\bigl(\omega_0(1,\kappa)\otimes\omega_1(\kappa,1)\bigr)
                              (f_y\otimes g_y).                \tag{QT17}
\] The two characters have canceled; neither has individually been set to 1.

Let \(\mathcal B_\phi\) denote the actual two consecutive horizontal theta integrations with fixed initial automorphic input \(\phi\), including the complete coefficient coevaluation of K037/K038. For any \(\kappa\), simultaneous middle translation of the two kernels is removed by the change of variable \(h\mapsto h\kappa\) in the middle integral. Its quotient Haar measure is right invariant. The horizontal groups are definite, their adelic quotients are compact, and these smooth theta integrals converge absolutely. Thus \[
 \mathcal B_\phi(f_{\kappa y}\otimes g_{\kappa y})
                  =\mathcal B_\phi(f_y\otimes g_y).             \tag{QT18}
\] This is equality of the entire resulting automorphic function, before terminal restriction or pairing. It does not require that each raw tensor be identical to its translate.

The fixed-tensor identity R034/LIFT26 is used before QT18: the same local y variable occurs in both factors, in the same middle frame, and the identity is before theta summation. Consequently there is no substitution of an unrelated source kernel or an independent second y sum.

The arithmetic finite-class expression has the same covariance. Its frame changes are \((X,y)\mapsto(\kappa_1X\kappa_0^{-1},\kappa_1y)\) and \((Z,\eta,y)\mapsto
(\kappa_TZ\kappa_1^{-1},\eta\kappa_1^{-1},\kappa_1y)\). The Fourier pairing is invariant. K038(CA24--CA29) proves the remaining source and terminal characters with these same honest factors. Under a rational change of representative the stabilizer is conjugated and has the same order; refining a finite level uses its exact orbit-stabilizer masses K038(CA30a--CA30c). Hence QT18 transports to the actual class-weighted coupled map with all of its \(|\Gamma_0|^{-1}|\Gamma_1|^{-1}\) factors. No freeness, constant stabilizer, or new principal level is assumed.

Combining QT14 and QT18 now gives the permitted compression: \[
 \mathcal B_\phi\!\left(\sum_{y\in Y}f_y\otimes g_y\right)
                  =h_M\,\mathcal B_\phi(f_{e_r}\otimes g_{e_r}).
 \tag{QT19}
\] Terms outside \(\mathcal O_Q\) are zero by QT12. Equation QT19 is an integrated identity. The unintegrated sum in its left side remains the one fixed in R034 and K038.

\subsection{5. All three weights and the exact QLast equation}\label{all-three-weights-and-the-exact-qlast-equation}

Retain the exact indices of R034: \[
 I_1={b\brack r}_q,\qquad I_2={N\brack r}_q,\qquad
 I_1=q^{rt}\frac{\kappa_b(q)}{\kappa_r(q)\kappa_t(q)},\quad
 I_2=q^{2rt}\frac{\kappa_N(q)}{\kappa_r(q)\kappa_{2t}(q)},
 \tag{QT20}
\] where \(\kappa_j(q)=\prod_{i=1}^j(1-q^{-i})\). QAt.wt\_pen, wt\_top and QFam.Pins now read exactly \[
 w_0=1,\quad w_1=I_1h_M,\quad w_2=I_2,\quad
 dQ_0=dL_0,\quad dQ_1=I_1h_M\,dL_1,\quad
 \epsilon_0=\delta_1,\quad \epsilon_T=I_2\delta_1,\quad
 c_Q=(I_1I_2)^{-1}.                                           \tag{QT21}
\] Here the dL's are the actual finite Heisenberg combinations representing \(f_{e_r}\) and \(g_{e_r}\) in the fixed basic models. Their existence and allowed finite support follow from the Heisenberg basis result LIFT15/DO017, as used already in R034. They do not act on a newly normalized basic tensor.

For finitely many Q places take products of the displayed weights and indices and convolutions of the actual supported operators, precisely as QFam.Pins specifies. Let \[
 I_j^{\rm all}=\prod_{v\in Q}I_{j,v},\qquad
 h^{\rm all}=\prod_{v\in Q}h_{M,v}.
\] The scalar product on the right of the source QLast law is \[
 (I_1^{\rm all}I_2^{\rm all})^{-1}
           (I_2^{\rm all})(I_1^{\rm all}h^{\rm all})
                          =h^{\rm all}.                       \tag{QT22}
\] The initial trace contributes exactly one by QT9. The terminal trace contributes \(I_2^{\rm all}\), not one. The first dQ contributes no further weight. This identifies every scalar in that equation.

Let \(C_{Q,0}\) be the image of the two original basic vacua under the actual lastBi operator. It is fixed before the tuple character is chosen, with the actual word, negative line and common Heisenberg conventions fixed. This independence is literal: CombAtSingleZero.lastBi32--35 uses only the two orthogonal decompositions, the fixed line map and the two Heisenberg functions. It contains no splitting character. At rank zero, PlaceLaws.lst44--47 composes it with the fixed coordinate restriction flatBF; IsThetaFam.flatBF107--108 likewise introduces no splitting character. Its zero-dimensional measures are the fixed probability measures. Thus the basic tensor does not secretly acquire a tuple-dependent scalar.

First express the physical first input as \(s_{Q,0}\) times its canonical function in the passive model, as in R034.6--7. Divide R034.12 by this nonzero scalar to obtain the canonical eQ version used by QLast. The scalar stays in VForm.sc; it is not part of QLast or its basic tensor. Apply that entire Schwartz identity with this \(C_{Q,0}\); apply QT19 place by place; then restrict the resulting equality of automorphic functions to the terminal subgroup and pair with any admissible terminal automorphic function. The equality obtained is precisely PlaceLaws.QLast with the data QT21 and product convention QT22. Initial and terminal index sets are singletons; the middle ones are the orbit-indexed singletons just proved. The R032 projector may be applied to the original vertical input throughout, since QT9 is universal and the other changes occur at the disjoint Q places.

This proves the actual m=2 QLast equality in the genuine model. The geometric frame and basic-vacuum clauses are those already proved in R034/R016, and the finite Q type clauses are the actual ones in QT2, not arbitrarily enlarged function spaces. For clarity, the remaining local QAt clauses have the following explicit realization.

\begin{longtable}[]{@{}
  >{\raggedright\arraybackslash}p{(\columnwidth - 2\tabcolsep) * \real{0.5000}}
  >{\raggedright\arraybackslash}p{(\columnwidth - 2\tabcolsep) * \real{0.5000}}@{}}
\toprule\noalign{}
\begin{minipage}[b]{\linewidth}\raggedright
QAt clauses
\end{minipage} & \begin{minipage}[b]{\linewidth}\raggedright
Realization and justification
\end{minipage} \\
\midrule\noalign{}
\endhead
\bottomrule\noalign{}
\endlastfoot
T, reg, laws, bas, le, leT, openT & The actual plain-frame twistChain with \(\delta\) above; parahoric groups are compact open, the two functions are supported in integral matrix lattices and constant on \(q\)-cosets. The honest source/target calculation gives first characters \(1,\lambda_1^{-1}\) and second characters \(\lambda_1,\lambda_2^{-1}\), exactly HasTypes. \\
mu & The actual source IsComp identification; split the local character into its displayed tame unit part and its unchanged unramified part. Open kernels at a finite place are retained. \\
FX, FY, FT, FM, bes, jac, tau, latY, latT & R034's actual core and terminal frames, fixed negative line and rank-zero decompositions; its integral core change preserves the specified selfdual lattices. \\
gl, gl\_loc & Embed \(\mathrm{GL}\) at the chosen split factor and its inverse Hermitian adjoint at the conjugate factor, conjugating by the displayed frame; use identity at every other place. This is a group homomorphism and gives exactly that local frame action. \\
A, reps, Ch, rep, rep\_bij, adapt, cos, one & QT3 and QT11--14, with singleton orbit indices. No other recursive representative clause occurs at \(s=0\). \\
eL, eLT, dL\_pen, dL\_top & The fixed local Heisenberg expansions of the actual raw functions, QT12 and QT15; LIFT15/DO017 gives expansions supported at this same place and their model-independent IsCombAt property. There is no dL\_mid clause at \(s=0\). \\
epsilonL, epsilonL\_at, idx, m, wt & QT4--8 and QT20--21, with maximal compact Haar volume1. \\
\end{longtable}

The QFam cross-character clauses are the existing R016 split labels: the intermediate positive-rank mid clause is empty for \(m=2\), and lastc is the same fixed source character used in R034. The FY spaces can be the actual finite-level spaces QT2; LIFT31 proves that the unchanged upward maps land in them. R032 acts at disjoint G places and preserves these Q conditions.

The orbit-indexed Tup is an explicit realization of this law, and QT19 is its comparison with the unchanged full-Y arithmetic map. If one instead freezes a redundant all-Y index set, its operators must be related by that sum/orbit map; the two unintegrated finite arrays are not identified by a claimed bijection of index sets. The existential ambient D7 packaging is not claimed on a field where HostRaw is separately obstructed.

\subsection{6. Scalar, integral, and fixed-family scope}\label{scalar-integral-and-fixed-family-scope}

In the actual passive core coordinates the unchanged first vertical vector has coefficient \[
                s_{Q,0}=\chi_v(\bar\pi_T)^{-1},\qquad
                s_{Q,1}=1.                                   \tag{QT23}
\] This is R034's root of unity. It belongs to VForm.sc, and can depend on the tuple. It is not absorbed into the fixed \(C_{Q,0}\), and it is not canceled by an additional honest Weil-frame factor. Thus the initial trace simplification does not alter the before-tuples choice of C.

Globally the moved terminal trace is \[
                        t.e_T=I_2^{\rm all}\delta_{b_T}.        \tag{QT24}
\] It is distinct from both the initial trace QT10 and the genuine projector of R032. Equations QT10, QT23 and QT24 are all needed when comparing with a terminal mass or an unchanged detector.

There is no new coefficient denominator or norm charge from the initial Q trace: its support, coefficient and operator norm are exactly \(1,1,1\). The middle and terminal indices remain as in QT21. Although their product cancels in QT22, one may not erase them from separately normalized maps or from a prescribed integral lattice. In particular an index can have nonzero 11-valuation.

The local raw Fourier coefficients have the already retained \(q^{-Mb}\) and roots of unity, and the frame coefficient QT23 is a root of unity. At actual Q places \(q\ne11\), these are 11-units or 11-integral combinations, as previously proved for the unchanged raw kernel. Nothing here computes a least integral right inverse or a restricted terminal nonzero value.

The exact orbit weight has \[
 \log h_M=b(M-1)\log q+\log(q^r-1).
 \tag{QT25}
\] It is a genuine counting weight, not an additional free normalization parameter. All global class covolumes and the separate archimedean recoupling, G projections and T operator of DO050--DO052 remain in their stated normalizations.

\subsection{7. Source locators and verification scope}\label{source-locators-and-verification-scope-2}

All archive references use the accepted pinned source commit 7cbf7a0c32fca1badf05931879a37c3f8a2b03b4. No Lean build or validity audit is performed.

\begin{itemize}
\tightlist
\item
  \nolinkurl{Definitions/LatticeFrames/ToOneSucc.lean35–37} and97: toOne and toOne\_one.
\item
  \nolinkurl{Definitions/LatticeFrames/IsOneComb.lean22–27:} the complete universal IsOneComb predicate.
\item
  \nolinkurl{Definitions/Foundations/SimpleSum.lean13–34:} lev, InN, N, embX and R; lines73--74 give IsNEquiv.
\item
  Definitions/TwistedInnerProducts/ \nolinkurl{StandardOrderExtendedBasisApplyLowBlock.lean384–407:} parahoric, its group and character, and parPhi; 675--703 Cgrp, prX and IsColChar; 880--921 Psi, Adapted, coupIdx, xiPen and xiTop; 932--937 IsRepSystem, including initial singleton; 964--975 and996--998 actual plain-frame twistChain.
\item
  \nolinkurl{Definitions/LatticeFrames/CosInj.lean23–118:} QAt, especially reps/rep\_bij78--83, epsilon104--106, weights113--118;129--150 QFam/Pins with all weight placements.
\item
  \nolinkurl{Definitions/LatticeFrames/PlaceLaws.lean137–138:} moved trace pins;240--252 the literal QLast equation.
\item
  \nolinkurl{Definitions/LatticeFrames/CombAtSingleZero.lean32–35,} PlaceLaws.lean44--47 and IsThetaFam.lean107--108: the tuple-independent lastBi/flatBF function, with the fixed zero-dimensional measures at rank zero.
\item
  Proofs/DecoratedLatticeProofs/ \nolinkurl{ExchangeCombinationMovesLastStepDatumToChart.lean15–77:} the source also proves N(a,1,1) trivial in its last exchange.
\item
  \nolinkurl{LastStepsKernelEqSumCouplingCosets.lean30–125:} the full same-basic-tensor Fourier sum, with initial epsilon.
\item
  Formalism/LocalToGlobal/Columns/ \nolinkurl{FiniteSumNonzeroLargeDepth.lean1019–1026:} zero\_uniq/zero\_h/zero\_Q and last\_uniq/last\_h; 1123--1167 the full ybSet, XiCpl1/2 and ftY.
\item
  R034, K038, LIFT26, LIFT31 and R032 are invoked only in their already checked genuine-model, covariance and normalization scopes. The original mixed HostRaw obstruction is unchanged.
\end{itemize}

The companion certificate checks finite rank-one traces, probability-volume weights, the exact coupling orbit, its transitivity and cardinality, both honest character identities, Fourier phase transport and the complete product of the source weights in finite examples. It is not a computation of an automorphic period or of a D7 record.

\clearpage

\section{AM-R-035 --- the common finite last-two tensor with the actual G and Q comparisons}\label{am-r-035-the-common-finite-last-two-tensor-with-the-actual-g-and-q-comparisons}

Author: GPT-6 Astra Pro, ALP\_MACHINE /root. 22 November 2026 UTC; reasoning effort is not separately exposed.

\textbf{Whole-proof PASS, 22 November 2026 UTC.} Root read the entire proof; the independent DO whole-proof/source review passed at07:24:25 UTC. Only this status paragraph and the final check-status paragraph were updated after that review. All event timestamps and hashes are retained.

The local G comparison is now an equality of the whole Schwartz tensor, using the normalized generic tests of DO052 and the actual negative line. Together with R034 and DO051 it gives one fixed finite tensor before the auxiliary Q tuple is varied. The T statement uses the explicit DO051 replacement operator; the corresponding law for the original glued vacuum is a period identity on M\_ram. These two statements are kept distinct. No unproved initial trace, global arithmetic phase, or nonzero arithmetic row is absorbed into this finite identity.

\subsection{1. Entire G tensor in the fixed models}\label{entire-g-tensor-in-the-fixed-models}

Keep the ranks a=42 or168, r=a+1, b=1802, N=r+2t and all original G primes. At a G place v, let q be its rational uniformizer and residue cardinality, n\_P=q\^{}(a(a-1)/2), alpha=g diag(I\_a,bar(pi\_T)), and d=diag(I\_a,q). The actual moved first frame is b\_0=alpha d; the other two moved frames are1. Write psi\_k(x)=psi\_E(xi0 iota\_v(x)), with conductor O and additive volume vol(O)=1. All of these are fixed before the auxiliary tuple changes.

Let v\_0,v\_1 be the same physical integral basic vacua as DO052. Use its passive source model T\_alpha, the unchanged middle model T\_1, and the actual source line exchange S\_v. Put

\[
 C_{G,v}=S_v(v_0\otimes v_1).
 \tag{R35.1}
\]

The canonical first and second scalar functions are exactly the checked LIFT22/DO052 functions

\[
 \Phi_0(X,y)=f_{U_a}(X)\mathbf1_{q^{-1}e_a+O^b}(y),
 \qquad \Phi_1(Z)=H_G(Z).
 \tag{R35.2}
\]

Here e\_a denotes the extra source coordinate with zero-based index a; f\_U and H\_G retain their entire integral support and their actual generic and old diagonal-block characters. They are not the uncut identity-cell functions. The horizontal final row coordinate is eta=t\_last/pi\_T.

With the literal source character and frames, LastLoc gives the transform

\[
 \Xi_G(X,Z,\eta)=f_{U_a}(X)H_G(Z)
       \mathbf1_{O^b}(\eta)\psi_k(-q^{-1}\eta e_a).
 \tag{R35.3}
\]

Indeed substitute y=q\^{}\{-1\}e\_a+u in the Fourier integral in LastLoc. Translation has no measure factor, vol(O\^{}b)=1, and its integral in u is exactly the characteristic function of the dual lattice O\^{}b. This proves the entire formula, including its zero set. The last row pi\_T inverse is retained, even though pi\_T is a unit at G.

Let d\_(G,0),d\_(G,1) be the actual source Heisenberg combinations realizing the two factors of R35.3. LIFT22(CG7--CG9) and K038(CA15--CA18) identify these functions with the fixed penultimate and terminal G functions. Then

\[
 S_v\bigl(T_\alpha(\Phi_0)\otimes T_1(\Phi_1)\bigr)
       =(d_{G,0}\otimes d_{G,1})C_{G,v}.
 \tag{R35.4}
\]

To justify the same basic tensor, apply LastLoc to the two basic integral indicators first. Its transform is the product horizontal integral indicator, because pi\_T is a unit here. The two-variable model in that source statement is the same linear model for the basic and actual functions. Its IsBiModel law takes each Heisenberg operator on the product indicator to the corresponding operator on C\_(G,v). Apply this to R35.3. Any fixed scalar relating that model to a chosen Fourier convention is already in C\_(G,v) on both sides; none is set to1. Thus R35.4 is an identity of actual vectors, not an inference from a nonzero local period.

For the unchanged-map normalized physical tests delta\_j\^{}P=phi\_(P,j)/n\_P of DO052, the exact statement is

\[
 \boxed{\quad
 S_v\bigl(\omega_0(b_0,1)\delta_0^P\otimes\delta_1^P\bigr)
   =s_{G,v}(d_{G,0}\otimes d_{G,1})C_{G,v},\quad
 s_{G,v}=q^{-a(a-1)-b/2}(\mu_0\xi)(q)^{-1}.
 \quad} \tag{R35.5}
\]

The physical model maps on the left are understood as in DO052(GP14). Equation R35.5 follows by inserting its two individual scalars into R35.4. The modulus exponent is2623 on J14 and28957 on J56. Both factors n\_P inverse are present. They are not an additional loss in the original two-edge map on its prescribed P-character subspace: that map equality and its unchanged least-preimage norm are DO052(GP8--GP9). If the literal raw generic datum is used instead, multiply R35.5 by n\_P squared; the corresponding map is also n\_P squared times the old map. These two normalizations may not be mixed.

\subsection{2. Restricted tensor product and one fixed finite C}\label{restricted-tensor-product-and-one-fixed-finite-c}

At every finite place take the same original vertical basic vacua v\_(0,v),v\_(1,v), and let C\_(0,v)=S\_v(v\_(0,v) tensor v\_(1,v)). At T the first is the actual glued-coordinate reading of the original standard bottom vacuum, as in DO051(FT4,FT9). At G and Q it is the same physical integral vacuum used above and in R034. At all other places the actual vacua and exchange are simply retained.

Let O\_T be the product of DO051's fixed finite Heisenberg BiOp at the places above T, and set

\[
 C_f=O_T\bigotimes_{v<\infty}'C_{0,v}.
 \tag{R35.6}
\]

The restricted product is taken in the already fixed global oscillator realization; no independent local phases are selected. Only finitely many factors are modified. O\_T, all basic vacua and S are fixed before the auxiliary tuple changes. In particular C\_f is independent of that tuple. Its T component is the retained Gamma\_T T\^{}(fb/2) times the horizontal product basic vacuum, with the phase Gamma\_T of DO051(FT14a--b) unchanged.

For a given auxiliary tuple, form the two physical finite tests as follows: at G use the left-hand tests of R35.5; at Q use the unchanged raw Q tests of R034; at T use (E\_(0,v)v\_(0,v),v\_(1,v)) from DO051(FT18a); elsewhere use the original basic tests. Denote their tensor product by F\_f\^{}mod. The superscript records the explicit T replacement. For each tuple of Y labels at Q, let D\_y be the product of the actual horizontal G, Q(y) and T Heisenberg operators, identity elsewhere. Operators on different local factors commute. Then

\[
 \boxed{\quad S_f F_f^{\rm mod}
  =s_G(\xi)s_Q(\xi)\sum_{y\in Y_Q}D_y C_f,\quad}
 \tag{R35.7}
\]

where

\[
 s_G(\xi)=\prod_{v\in G}
       q_v^{-a(a-1)-b/2}(\mu_0\xi)(q_v)^{-1},\qquad
 s_Q(\xi)=\prod_{v\in Q}\chi_v(\bar\pi_T)^{-1}.
 \tag{R35.8}
\]

Here Y\_Q is the full product of the prescribed Q quotients, and the Q operators include the coefficients q\_v\^{}(-M\_v b) of R034. These coefficients are not replaced by the reciprocals of quotient cardinalities. A placewise proof of R35.7 is immediate but precise: R35.5 gives its G factor, R034(R34.12) gives each full Q sum, and DO051(FT18) gives (d\_T0 tensor d\_T1)O\_T C\_0=S\_v(E\_0v\_0 tensor v\_1). The other factors are their defining C\_0. Multiply these identities in the same restricted tensor product and distribute the finite Q sums. This proves R35.7 before theta summation or class integration.

Because each rational G prime has f selected places,

\[
 |s_G|=(g_0g_1)^{-f(a(a-1)+b/2)},\qquad |s_Q|=1.
 \tag{R35.9}
\]

The tame/unramified Hecke values in s\_G and the residue-character values in s\_Q retain their actual phases. No product-formula cancellation is asserted. The individual scalar list for VForm is still s\_(G,0),s\_(G,1),s\_(Q,0),s\_(Q,1) of DO052/R034, whose product is R35.8. Equations R35.7--R35.9 supply no additional Petersson normalization or arithmetic valuation of Gamma\_T.

\subsection{3. Original T period, trace weights and the remaining global gate}\label{original-t-period-trace-weights-and-the-remaining-global-gate}

R35.7 concerns the explicit T-modified tests. For the original glued first vacuum and unchanged second vacuum, DO051(FT10--FT13) supplies the separate source period identity on an initial M\_ram input. Its coefficient, in the stated original Haar convention, is

\[
 c_T=T^{-fb/2}(T^b-1)^f
       \prod_{v\mid T}(\mu_0\xi)(\operatorname{idele}_v\varpi_v).
 \tag{R35.10}
\]

That identity uses the literal diamond epsilon\_v and FFStar on the initial input. It does not say that applying E\_0 to a vacuum leaves the entire vertical map unchanged. The positive T powers in c\_T and the T component of C\_f cancel only in the product in the period law, leaving the fixed Gamma\_T, the displayed finite factors and the prescribed input transform. Any input outside M\_ram requires its extra terms; no such terms have been discarded here.

At Q the middle and terminal weights and c\_Q remain exactly R034.14. Their product changes an invariant orbit representative to its full cardinality. The initial epsilon\_0 and w\_0 are not supplied by R35.7 and must be established from their actual source definitions. Likewise, the real tensor is DO050's full recoupling; its selected compact average and Fischer/PBW phase cannot be substituted pointwise for a product of highest vectors.

Thus the finite common-tensor assembly, including the previously missing G identity, is explicit. To obtain the actual global arithmetic MainX value one must still establish the remaining initial/global row and trace pins, keep the real arithmetic comparison and fixed model phase, and produce a permitted nonzero M\_ram restricted row with a controlled integral preimage. No central Rankin value, universal DP44 sign or shorter lift chain is proved by this tensor identity.

\subsection{4. Exact inputs and verification scope}\label{exact-inputs-and-verification-scope}

The new claim is the compatible actual G full-tensor identity R35.5 and its restricted-product assembly R35.7 with one unchanged C\_f. The Fourier/coset integral is proved above. The normalized physical function comparison, Q tensor, and T operator/period are cited results, not re-proved here.

\begin{itemize}
\tightlist
\item
  AM-DO-052, GP1--GP20, especially GP11--GP18.
\item
  AM-R-034, R34.1--R34.14; ALP-LIFT-22, CG7--CG9; AM-K038, CA15--CA18.
\item
  AM-DO-051, FT4, FT9--FT18 and its fixed BiOp/common-C construction.
\item
  AM-DO-050, the selected real recoupling in its stated primitive Fischer-isometric convention; ALP-LIFT-40 retains the alternative integral primitive Pieri norm.
\item
  Pinned stipulated source commit 7cbf7a0c32fca1badf05931879a37c3f8a2b03b4: \nolinkurl{Definitions/LatticeFrames/IsCombAt.lean24–27,49–66} for IsCombAt, LastLoc and the same two-variable model; \nolinkurl{HzUnit.lean66–82} for eg/egT and horizontal cg; \nolinkurl{PlaceLaws.lean128–147} for the scalar, row and operator requirements of VForm. No Lean execution or validity audit is performed.
\end{itemize}

Whole-proof and independent source-review events are recorded separately. The existing exact DO051/DO052 and R034 diagnostics retain their own scopes; no new numerical automorphic or Fourier evaluation is claimed for this assembly. They cannot certify the unproved global arithmetic period or its integral image.

\clearpage

\section{AM-K058 --- a genuine finite-projector selected jet in the actual packet}\label{am-k058-a-genuine-finite-projector-selected-jet-in-the-actual-packet}

Author: GPT-6 Astra Pro, ALP\_MACHINE \nolinkurl{/root/kummer_extension}. 22 November 2026 UTC. Inherited reasoning effort is not separately exposed.

\textbf{Author whole-proof check PASS --- 2026-11-22T07:23:45.485095+00:00.} Exact diagnostic certificate: 237 assertions PASS; independent reviews are separate events.

At the upper admissible gap, this finding constructs a nonzero selected jet of an actual genuine-e0 bottom cusp form on an H-deck translate of LIFT41's favorable point of the prescribed T packet. Every role assignment, binary word and permitted canonical twist is retained. The selected component is the actual K043 component, with its PBW normalization.

The initial restricted row is generic at every G place. It is not asserted to be in the strict-T Mram summand, a prescribed initial automorphic constituent, or the earlier coupled arithmetic image. The result therefore does not assert a terminal period on that image or a DP44 cost sign.

\subsection{1. Fixed data, dual conventions and the good finite subspace}\label{fixed-data-dual-conventions-and-the-good-finite-subspace}

Use the actual E/F, a=f=42 or168, n=a+1, Q=g0g1 and T of K031/K032, DO040 and K044. Let

\[
 G=U(I_a,-1)(\mathbb R\otimes F),\quad
 \Gamma=U(I_a,-1)(F)\cap J^{\max},\quad
 Z=\mu_{28}I_n,\quad z_0=|Z|=28.
 \tag{PG1}
\]

All finite levels, including maximal29 and maximal bottom level at T, are unchanged. J1 is the existing principal level at the two G primes. Write Gcal=Jmax/J1 and Hcal for its product of leading GL\_a subgroups. The bottom group J\_B is the inverse image of the actual Ucal, with its prescribed positive generic character psi, as in R027.1--4.

Let sigma be the holomorphic Hilbert representation of K043/K046, V its minimal compact type and W its complete selected compact component. Put d=dim W, and let pi denote its unitary action. The actual scalar bottom right representation is sigma-dual, as in R027 and K057. The canonical holomorphic lift of w is

\[
                 F_w(g)=\operatorname{pr}_V\pi(g^{-1})w.
 \tag{PG2}
\]

This fixes the inverse convention explicitly. It is the same canonical holomorphic lift as K046, expressed by orthogonal projection. In particular F\_w(gk)=V(k)\^{}\{-1\}F\_w(g). If omega is the central character of sigma, then F\_w(zg)=omega(z)\^{}\{-1\}F\_w(g).

Choose the fixed genuine goodTuple constituent Pi of DO042/DO043, with central character Omega=omega on Z. It is cuspidal at all but the two retained mixed places, one over each G prime. The mixed degrees remain (37,6) or(163,6). Put R=Pi-dual, with its invariant unitary metric. Its central character is omega-inverse, so R and the lift in PG2 have the same left central covariance. No central twist is added.

At each mixed place take the orthogonal projector onto the highest mirabolic I\_n submodule of R. It is the dual of LIFT32(FB2--5)'s canonical copy; the nondegenerate character is the inverse of the prescribed bottom character. At a full-degree cuspidal place use the identity projector. The I\_n submodule is then the whole representation, since both dimensions are the cuspidal degree. Denote the tensor product projector by Qgood and its image by Rgood. Thus

\[
 R_{\rm good}|_{\mathcal H}\simeq
       \bigotimes_{k,v}\operatorname{Ind}_{U_a(\mathbb F_{g_k})}
                           ^{GL_a(\mathbb F_{g_k})}\psi^{-1},
 \qquad
 D:=\dim R_{\rm good}
       =\prod_{k=0}^1\prod_{i=1}^a(g_k^i-1)^f.
 \tag{PG3}
\]

The induction and embedding carry the probability metrics of LIFT32; Qgood is an orthogonal projection in the given finite representation, not a new finite idempotent replacing e0. The product in PG3 is multiplicity-free and consists exactly of the generic Hcal irreducibles. The full R can have other Hcal constituents. We do not discard them in the automorphic target; only the domain of the selected test map is restricted to Rgood.

\subsection{2. The actual packet point and its exact compact stabilizer}\label{the-actual-packet-point-and-its-exact-compact-stabilizer}

Let g\emph{=gamma\_infty be the rational unitary matrix of LIFT41 and let p=g}o. LIFT41(NP11) identifies (p,1) with a principal-G point {[}(o,x tau){]} of the unchanged T packet. Its four nonzero positive coordinates are proportional to Q(u1,u2,u3,u4), where the four positive integers u\_i in NP4 are pairwise distinct. The same proportionality is used at every label. The Cartan parameter satisfies 2s\textless1/100.

The standard origin stabilizer and the actual point stabilizer are

\[
 \begin{split}
 S_0&=\Gamma\cap K
       =(\mu_{28}^{a}\rtimes S_a)\times\mu_{28},\\
 S_p&=\Gamma\cap g_*Kg_*^{-1}
       =\{\operatorname{diag}(tI_4,A_0,t):
                  t\in\mu_{28},\ A_0\in\mu_{28}^{a-4}\rtimes S_{a-4}\}.
 \end{split}
 \tag{PG4}
\]

The second display uses the four supported positive coordinates first. Consequently \textbar S\_p/Z\textbar=28\^{}(a-4)(a-4)!.

\textbf{Proof.} The description of S0 is K050(BE2): every entry of an integral positive unitary matrix has all conjugates at most1 in absolute value; the algebraic norm and column norm identities force a monomial root-unit matrix. If h fixes p, its Cartan parameters are at most2s\textless1/2. The positive entries in its last column have all conjugates bounded by sinh(1/2)\textless1. They are integral, hence zero. Thus h already belongs to S0. Writing h=diag(A,t), preservation of the line through the explicit negative vector forces A w=t w, because its last coordinate is nonzero. A monomial matrix can permute two supported positive coordinates only if their absolute values agree. They do not; the first four diagonal entries must be t. The zero coordinates admit exactly the arbitrary monomial block A0 displayed above. Every such matrix fixes p.~This proves both inclusions and the order formula. All embeddings, including the conjugate CM half, enter the norm argument. \(\square\)

The compact action k\_s=g\emph{\^{}-1 s g} for s in S\_p preserves W. The subspace Rgood is S0-stable: S0 is contained in Z times the leading Hcal, and PG3 is Hcal-stable. Hence on

\[
          E=R_{\rm good}\otimes W
 \tag{PG5}
\]

the operators R(s) tensor pi(k\_s) give a unitary action of S\_p/Z. Their central factors cancel. Its exact orthogonal invariant projection is

\[
 A_p={1\over |S_p/Z|}
      \sum_{s\in S_p/Z}R(s)\otimes\pi(k_s)|_W.
 \tag{PG6}
\]

The choice of a coset representative does not affect a summand.

For every s in S0 outside Z,

\[
              \operatorname{Tr}(Q_{\rm good}R(s))=0.
 \tag{PG7}
\]

Indeed at any full-degree cuspidal place its reduction is semisimple, nonscalar, and has the negative-line eigenvalue in the residue field. It cannot belong to the elliptic torus defining that cuspidal character. The Deligne--Lusztig semisimple formula therefore makes that character zero, exactly as in K050(BE4). Root and monomial reduction is faithful at these primes. The two mixed partial traces cannot remove the zero. At a scalar, the trace is D times the actual central character. Thus

\[
                  \operatorname{Tr}_E A_p={Dd\over |S_p/Z|}>0.
 \tag{PG8}
\]

This is an exact finite-stabilizer contribution, not the assertion that the stabilizer is scalar. It also identifies the precise allowed selected fibre E\^{}(S\_p/Z). Below the global selected-jet Gram is proved nonzero on that fibre; it is not claimed positive definite on the entire fibre.

\subsection{3. A matrix-valued selected seed and its moments}\label{a-matrix-valued-selected-seed-and-its-moments}

Define a map W→V at every real point by

\[
       B(g)w=F_w(g_*^{-1}g)
              =\operatorname{pr}_V\pi(g^{-1}g_*)w.
 \tag{PG9}
\]

Set b(g)=\textbar\textbar B(g)\textbar\textbar\_HS. Summing the exact K046 moments over an orthonormal basis of W gives

\[
 b\le\sqrt d,\qquad \int_G b^2={d\over C},\qquad
              \int_G b\ge{\sqrt d\over C},
 \quad C=\prod_{\tau,i}(r_{\tau,i}-1).
 \tag{PG10}
\]

The probability moments of b² are the same as those for a single unit selected vector. This follows because the shifted metric is scalar on the complete irreducible selected component; it is not an independence assumption on a chosen tensor basis. The last inequality uses b²\textless=sqrt(d)b. Translation by g* preserves all integrals.

Take R0=9/2 for J14 and R0=5 for J56, and write

\[
        C_* = g_* C_{R_0},\qquad
        \delta_0=f2^{-\ell_+/8}.
 \tag{PG11}
\]

K046 and LIFT42(FC12--15), applied to b/sqrt(d), give

\[
       \int_{G\setminus C_*}b<\delta_0\int_G b.
 \tag{PG12}
\]

The smaller J14 radius uses the exact moment bound A\_tau\textless969 and cutoff Y\textgreater=1/2000. No smaller-radius assertion for the old 2\^{}-12 cutoff is used. These statements hold for every word, role assignment and allowed twist at ell=ell\_+. The PBW factor rho\_all remains in the definition of W's unitary inclusion, as in K043/K046.

For every g define a linear map from E to R tensor V by

\[
        L(g)(r\otimes w)=r\otimes B(g)w.
 \tag{PG13}
\]

Its Hilbert--Schmidt norm is sqrt(D)b(g). It has the exact central covariance L(zg)=(R(z) tensor1)L(g), so the following sum is independent of representatives:

\[
       (Tt)(g)=\sum_{\gamma\in Z\backslash\Gamma}
                   (R(\gamma)^{-1}\otimes1)L(\gamma g)t,
             \qquad t\in E.
 \tag{PG14}
\]

The resulting sections satisfy Tt(delta g)=R(delta)Tt(g) and the actual right minimal-compact covariance. The finite covariance is proved by the change of variable gamma'=gamma delta. Absolute L1 convergence follows from PG10--12 and Tonelli. The full sum is smooth, holomorphic, L2 and cuspidal: decompose it over the existing finite quotient Gamma/Gamma1, and apply the accepted K046/LIFT41 cusp Poincare theorem to each principal-level translated seed. Translation by the rational g* is the conjugated arithmetic-group application in LIFT41(NP17--20). No new finite-level refinement or uniform norm assertion is inserted into that convergence theorem.

\subsection{4. Exact intertwining at the finite stabilizer}\label{exact-intertwining-at-the-finite-stabilizer}

Let H\_sigma be the Hilbert representation space of sigma, and embed E isometrically into R tensor H\_sigma by

\[
               i_p(r\otimes w)=r\otimes\pi(g_*)w.
 \tag{PG15}
\]

On this tensor space the diagonal S0 action is R(s) tensor pi(s). Its central Z action is trivial. Let A0 be its average over S0/Z. The seed corresponding to A0 i\_p(t) has exact left S0 covariance R. To check this, invariance under R(s) tensor pi(s) gives (1 tensor pi(s\^{}-1))A0 i\_p(t)=(R(s) tensor1)A0 i\_p(t); applying 1 tensor pr\_V pi(g\^{}-1) proves the covariance. Grouping PG14 by left S0 cosets gives exactly

\[
 Tt=|S_0/Z|\sum_{\gamma\in S_0\backslash\Gamma}
       (R(\gamma)^{-1}\otimes1)
          F_{A_0 i_p(t)}(\gamma\,\cdot).
 \tag{PG16}
\]

Thus the finite-intertwining seed is explicit. We have not picked an unrelated minimal-type vector. Unlike Ap, A0 need not preserve i\_p(E), because S0 need not fix the displaced point. Its compressed trace is nevertheless exactly Dd/\textbar S0/Z\textbar{} by PG7. For the actual point stabilizer, A0's analogue preserves i\_p(E), identifies with Ap, and gives

\[
                         T=T A_p.
 \tag{PG17}
\]

This follows either by grouping over Sp or by replacing a summation index by its left multiple in Sp. The tensor action, central factors and inverse convention in PG14--16 are all retained.

\subsection{5. A quantitative partial-trace Poincare lemma}\label{a-quantitative-partial-trace-poincare-lemma}

Put B0=2R0+1/100 and let N=N(B0), epsilon=epsilon(B0) be the actual integer count and normalized partial-trace supremum in LIFT42(FC4), applied to R and Qgood. Complex conjugation of Pi changes no bound. Its exact results give

\[
 N\epsilon<2^{-1000}<\tfrac12,\quad
 N<2^{32af},\quad
 \delta_0<2^{-64af-20}<{1\over8\sqrt{2N}}.
 \tag{PG18}
\]

The count includes the entire finite S0; no neatness is assumed here. The original sharper short-sum margins are 2\^{}-1200 and2\^{}-646000.

Truncate L by the indicator of C\emph{, and denote its periodization by T\_C. At a fixed g there are at most N active cosets gamma in Z\Gamma. Indeed, after fixing one active delta, all gamma delta\^{}-1 lie in Gamma intersect g} C\_R0 C\_R0\^{}-1 g*\^{}-1, whose Cartan parameters are at most2R0+2s\textless B0. Counting modulo Z can only reduce the bound N.

For two active cosets the exact Hilbert--Schmidt cross term is

\[
 \operatorname{Tr}\!\left(Q_{\rm good}R(\delta\gamma^{-1})\right)
       \langle B(\gamma g),B(\delta g)\rangle_{\rm HS},
 \tag{PG19}
\]

for the inner product linear in the first slot. Taking its simultaneous complex conjugate gives the opposite ordered term in the complete real sum. This follows by multiplying the two operators in PG13--14 and cyclically moving the domain projection. At equal cosets its finite factor is D. At distinct cosets the difference is noncentral, so its absolute finite factor is at most epsilon D. For a difference in S0 outside Z the factor is exactly zero by PG7. The arbitrary two mixed I\_n projectors are retained in the trace.

Writing b\_gamma=b(gamma g) for the active cosets, Cauchy--Schwarz gives

\[
 \begin{split}
 \|T_C(g)\|_{\rm HS}^2
   &\ge D\{\sum b_\gamma^2-
                       \epsilon\sum_{\gamma\ne\delta}b_\gamma b_\delta\}\\
   &\ge D(1-N\epsilon)\sum b_\gamma^2,\\
 \|T_C(g)\|_{\rm HS}
   &\ge\sqrt D\,\eta\sum b_\gamma,
          \qquad \eta=\sqrt{(1-N\epsilon)/N}.
 \end{split}
 \tag{PG20}
\]

There is no assertion that the individual summands have one sign. The finite partial traces control the complete Hilbert--Schmidt square.

The raw quotient measure is the R030 real Haar with discrete Gamma counting measure. Because PG14 sums over Gamma/Z rather than Gamma, unfolding a nonnegative scalar function gives

\[
 \int_{\Gamma\backslash G}
       \sum_{\gamma\in Z\backslash\Gamma}b(\gamma g)
           ={1\over z_0}\int_G b.
 \tag{PG21}
\]

The factor z0 is essential: the full Gamma sum contains z0 identical terms. The same identity applies to either truncated region. The tail is bounded above by its triangle inequality and Tonelli. Thus PG20 and PG12 imply

\[
 \|T\|_{L^1({\rm HS})}\ge
 {\sqrt D\over z_0}
       [\eta(1-\delta_0)-\delta_0]\int_G b
 \ge {\sqrt{Dd}\over2\sqrt{2N}\,z_0 C}>0.
 \tag{PG22}
\]

For the last inequality, eta\textgreater=1/sqrt(2N), delta0\textless1/4 and delta0\textless1/(8sqrt(2N)); hence the bracket is greater than 1/(2sqrt(2N)). PG10 supplies the remaining integral. This is an L1 long-tail proof, not a deletion of the non-short correspondences from a reproducing-kernel sum.

In particular T is nonzero on the exact invariant selected fibre in PG8/PG17. Its Hilbert--Schmidt norm as a map to the raw cusp Hilbert space satisfies, with V\_Gamma=vol(Gamma\G),

\[
          \|T\|_{\rm HS}\ge
        {\sqrt{Dd}\over2\sqrt{2N}\,z_0 C\sqrt{V_\Gamma}}.
 \tag{PG23}
\]

This is Cauchy--Schwarz on the finite-volume quotient applied to the matrix-valued section. It does not require an upper bound for a global Poincare norm or a cuspidal multiplicity estimate.

\subsection{6. The exact selected adjoint and positive Gram}\label{the-exact-selected-adjoint-and-positive-gram}

Let H0 be the finite-dimensional span of all the sections Tt. Let J\_p be the actual K043-normalized selected jet at p, in the transported LIFT41 chart. It takes values in R tensor W. Let

\[
                 J_{\rm good}=(Q_{\rm good}\otimes1)J_p.
 \tag{PG24}
\]

The precise adjoint identity on H0 is

\[
                 J_{\rm good}=z_0 C\,T^*,\qquad
                 J_{\rm good}T=z_0 C\,T^*T\succeq0.
 \tag{PG25}
\]

\textbf{Proof.} Pair a section S in H0 with Tt and unfold PG14. Its Gamma covariance cancels R(gamma), and the quotient-by-Z convention leaves the factor1/z0 in PG21. The remaining integral is the weighted holomorphic pairing with the seed at p.~Translate p to the origin; the invariant bundle metric and Haar measure are unchanged. Compact averaging selects exactly the W summand: its central compact weight fixes the degree, and the actual Pieri multiplicity is one. The K043 metric is Fischer divided by D\_PBW, so K046(BL29--31) gives

\[
       \langle S,Tt\rangle_{\rm Pet}
           ={1\over z_0 C}\langle J_{\rm good}S,t\rangle.
 \tag{PG26}
\]

All constants are therefore fixed by the original geometric metric, not fitted to the periodized value. Absolute unfolding follows from the L1 seed and bounded cuspidal Poincare sections in H0; one may first average compact radial truncations and then pass to the limit. This is the same finite-dimensional unfolding justification as K046(BL31), now with the explicit finite covariance. Equation PG26 is precisely PG25. \(\square\)

In particular the selected jet is nonzero. More quantitatively, rank T\textless=Dd and PG23 imply

\[
          \|J_{\rm good}|_{H_0}\|_{\rm op}
                   \ge {1\over2\sqrt{2N V_\Gamma}}.
 \tag{PG27}
\]

Indeed \textbar\textbar T\textbar\textbar op\textgreater=\textbar\textbar T\textbar\textbar HS/sqrt(Dd) and PG25 cancels z0 C exactly. Choose a maximizing singular vector in the finite-dimensional H0. There is an actual R-valued cusp section S with raw norm1 whose good selected jet has norm at least the right side of PG27. By PG17 the nonzero positive Gram lies in E\^{}(S\_p/Z). We have proved positive trace and a positive eigenvalue there, rather than surjectivity onto every allowed fibre vector.

The unitary selected jet in PG24--27 is the raw derivative divided by rho\_all=product\_tau sqrt(D\_tau), with the K043 matrix-unit convention. That factor must be restored for an unnormalized derivative. It is not an arithmetic unit, and no integral denominator is inferred from PG27.

\subsection{7. Genuine e0 and the actual H-generic packet}\label{genuine-e0-and-the-actual-h-generic-packet}

The R tensor V-valued section produces V-valued holomorphic automorphic sections by

\[
           \Phi_v(g,k)=S(g)(\Pi(k)v),\qquad v\in\Pi.
 \tag{PG28}
\]

The finite covariance makes this well defined under Gamma. Extend by zero to other coarse Jmax components, just as in K050(BE17). The map v→Phi\_v is a nonzero Gcal intertwiner whenever S is nonzero; Pi is irreducible, so it is injective. The actual e\_Pi and full genuine e0 act as the identity. This construction does not assume that Gamma/Gamma1 surjects onto Gcal.

It remains to obtain the prescribed positive psi type on an actual H-deck translate of p, rather than an arbitrary Gcal-deck point. Let Vgood=Rgood-dual inside Pi, let Ppsi be its actual psi projector, and put qpsi=dim Ppsi Vgood. By PG3, Vgood is the multiplicity-free Gelfand--Graev product. If dmax is the largest dimension of one of its irreducible Hcal constituents, Schur averaging gives

\[
 {1\over|\mathcal H|}\sum_{h\in\mathcal H}
              \Pi(h)P_\psi\Pi(h)^{-1}
        =\bigoplus_{\sigma\subset V_{\rm good}}{1\over\dim\sigma}I_\sigma
        \succeq {1\over d_{\max}}I_{V_{\rm good}}.
 \tag{PG29}
\]

Each generic constituent has a one-dimensional psi line. Its orthogonal complement outside Vgood is immaterial because all test vectors in PG29 belong to Vgood. Apply PG29 to the squared norm of the good selected jet from PG27, and sum over an orthonormal basis of Ppsi Vgood. There exist h in Hcal and a unit psi vector v in Vgood such that

\[
       \|J_{p h}\Phi_v\|
          \ge {\|J_{\rm good}S\|\over\sqrt{d_{\max}q_\psi}}>0.
 \tag{PG30}
\]

A basis vector suffices, so there is no unexplained finite coefficient normalization. The equality of the finite frame with the standard frame follows from LIFT41(NP11)'s principal-G representative; the residual adjustment belongs to J1. Since h belongs to Hcal, p h is in the initial principal-level H packet above the same maximal-H coarse class, exactly as in R027.5--7. A Gcal-deck translate outside this H orbit is not used. The original generic-character-level initial class can change; membership in the full unchanged initial packet is the assertion.

With finite Jmax probability normalization, compact Schur orthogonality gives the exact section norm

\[
                \|\Phi_v\|_{\rm sec}^2={\|S\|_{\rm sec}^2\over\dim\Pi}
                       \quad(\|v\|=1).
 \tag{PG31}
\]

Indeed average \textbar S(g)(Pi(k)v)\textbar² over the finite compact fibre first. That average is \textbar\textbar S(g)\textbar\textbar²/dim Pi; integrating the raw real quotient retains its stabilizers and gives PG31. Any alternate adelic or normalized-average Haar convention must apply the R030 conversion to both sides.

Here the distinction between a V-valued section and a scalar automorphic vector matters. Put d\_V=dim V. For a unit minimal-type test u in V, the scalar function attached to Phi is I\_Phi(bar(u))(g,k)=\textless Phi(g,k),u\textgreater. Compact averaging gives

\[
       \|\mathcal I_\Phi(\bar u)\|_{\rm aut}^2
               ={\|\Phi\|_{\rm sec}^2\over d_V}.
 \tag{PG31a}
\]

The actual holomorphic Poincare construction extends this minimal-type map to a linear intertwiner from sigma-dual: each seed is a matrix coefficient of that one discrete series, and periodization commutes with right translation. Its norm ratio is consequently the same on every unit vector, by Schur's lemma. In particular, for a unit vector w in the selected component, K043's exact PBW normalization gives

\[
       \mathcal I_\Phi(\bar w)(p h)
                  =\langle J_{p h}\Phi,w\rangle,\qquad
       \|\mathcal I_\Phi(\bar w)\|_{\rm aut}
                  ={\|\Phi\|_{\rm sec}\over\sqrt{d_V}}.
 \tag{PG31b}
\]

One may check the first identity directly on each Poincare seed: replace pr\_V in PG9 by the matrix coefficient against w and evaluate at p; this is the projection onto pi(g*)W. The normalized coefficient is exactly K046(BL29), so there is no additional selected-dimension or Pieri scalar. The extension can also be read in the source's dual convention of LIFT28(HJ20--21). Formula PG31a is the separate compact scalarization factor, not a new choice of global normalization.

Combining PG27--31b, there is an actual Hilbert-unit scalar bottom vector F in the prescribed selected compact type, genuine e0 block and positive psi type, for which

\[
       |F(p h)|\ge
        {1\over2\sqrt{2N V_\Gamma}}
              \sqrt{{d_V\dim\Pi\over d_{\max}q_\psi}}
       \ge {1\over2\sqrt{2N V_\Gamma d_{\max}}}.
 \tag{PG32}
\]

Choose w in PG31b in the direction of the nonzero jet in PG30. The second inequality only uses qpsi\textless=dim Pi and d\_V\textgreater=1. Equivalently, before scalarization the unit section has selected-jet norm at least the first bound with d\_V omitted. The finite bound dmax\textless=D from PG3 is available if an explicit product is preferable. No claim that the finite projections are integral accompanies this unitary norm estimate.

\subsection{8. Exact conclusion and remaining image condition}\label{exact-conclusion-and-remaining-image-condition}

At ell=ell\_+, every current role assignment, binary word and permitted canonical twist has an actual genuine-e0 bottom cusp form whose normalized selected restriction is nonzero in the prescribed T packet and has the actual positive generic type at all G places. Its initial H restriction is therefore a nonzero row of the opposite initial coefficient type and original strict T level. Equations PG27 and PG32 give norm bounds in their explicitly stated metrics.

This is a strict improvement over unprojected one-point interpolation and over minimal-type bottom existence. It still does not prescribe a particular initial row. In particular, nonzero restriction to the strict T level does not imply that its nontrivial-dilation projection Mram is nonzero. The construction does not identify the row with the earlier coupled theta image, supply an integral detector or vertical preimage, or prove the fixed terminal arithmetic pairing nonzero. The packet dimension obstruction and the small value at K047/K049's different packet point are unchanged.

\subsection{9. Sources and verification scope}\label{sources-and-verification-scope-4}

The accepted exact representation and geometry inputs are K043(PB1--10), K046(BL5--22,BL28--34), K050(BE1--7,BE15--17), R027.1--10, LIFT32(FB2--5), LIFT41(NP1--20), and LIFT42(FC1--15). The two latter reports are used at their frozen stated scopes. The field, root units, actual prime levels and goodTuple are K031/K032, DO040/042/043 and R025.

Primary anchors carried with those results are:

\begin{itemize}
\tightlist
\item
  Nakahama, R. (2016). Norm computation and analytic continuation of vector valued holomorphic discrete series representations. \emph{Journal of Lie Theory, 26}, 927--990. Theorem5.1 and(5.7), pp.960--961. \url{https://jolt.centre-mersenne.org/item/10.5802/jolt.919.pdf} .
\item
  Muic, G. (2019). Smooth cuspidal automorphic forms and integrable discrete series. \emph{Mathematische Zeitschrift, 292}, 895--922. Author version Theorems6-4 and7-6(i), pp.21--22 and26. \url{https://arxiv.org/pdf/1610.05483} .
\item
  Gan, W. T., Gross, B. H., \& Prasad, D. (2012). Restrictions of representations of classical groups: Examples. \emph{Astérisque, 346}, 111--170. Proposition4.2 and Corollary4.3. \url{https://www.numdam.org/item/AST_2012__346__111_0.pdf} .
\item
  Deligne, P., \& Lusztig, G. (1976). Representations of reductive groups over finite fields. \emph{Annals of Mathematics, 103}(1), 103--161. Theorem4.2, p.124. \url{https://publications.ias.edu/sites/default/files/Number27.pdf} .
\item
  Gorodetsky, O., \& Hazan, Z. (2017). \emph{On certain degenerate Whittaker models for cuspidal representations of GL\_(kn)(F\_q)}. arXiv:1707.07308v2. Theorem2.1 and(2.1), p.5. \url{https://arxiv.org/pdf/1707.07308} .
\end{itemize}

The new argument is the partial-trace matrix Poincare estimate with its retained long tail, the exact selected adjoint and Gram, the actual point-stabilizer calculation, and the finite H-generic extraction from the two mixed mirabolic subspaces. The finite diagnostics do not evaluate a Poincare series or an automorphic period; the global nonvanishing is the proved L1 and adjoint argument above.

\subsection{Author verification record}\label{author-verification-record-5}

Author whole saved PG1--32 proof, including PG31a--b, and complete final certificate read and reconstructed PASS at 2026-11-22T07:23:45.485095+00:00. The review checked the actual point and origin stabilizers, fixed mixed mirabolic domains, central tensor covariance, exact ordered noncommuting partial trace, retained long L1 tail, quotient-by-mu28 Haar factor, selected adjoint and positive Gram, H-generic extraction, and both finite and compact scalarization factors. The final diagnostic replay completed at 2026-11-22T07:23:12.518768+00:00 with237 named exact assertions and100608 auxiliary monomial candidates. It does not evaluate an actual automorphic value. LIFT42 FC1--15 was independently whole-read and used at its frozen scope; its root1853-check replay is retained, without a duplicate replay claimed here. Primary Gorodetsky--Hazan Theorem2.1 was directly consulted. Mram, a prescribed initial constituent, the earlier arithmetic image and an integral preimage remain unproved by this finding.

\clearpage

\section{ALP-LIFT-43 --- the actual T packet projection and its additional Gram}\label{alp-lift-43-the-actual-t-packet-projection-and-its-additional-gram}

Author: GPT-6 Astra Pro, ALP\_MACHINE \nolinkurl{/root/lift_accounting}. 22 November 2026 UTC.

Author whole-proof review PASS: 2026-11-22T07:33:31.235487+00:00. Whole-proof reread after the conductor-direction precision: 2026-11-22T07:35:48.819502+00:00, PASS. The exact certificate is recorded separately; no actual cross Gram is computed.

This report fixes the precise additional operator needed to pass from the H-generic strict-level detector of K058 to the literal initial Mram space. It also proves two concrete geometric statements: the ambient relative Hecke set contains principal-G correspondences explicitly very close on the scale of the actual T, whereas the fixed LIFT41 packet excludes the most obvious such example and forces a quantitative displacement on an entire specified rank-one unipotent family. No nonzero Mram projection or global restricted period is inferred.

\subsection{1. The literal projection, including every T place}\label{the-literal-projection-including-every-t-place}

Let V be the complex row realization in which the selected detector takes values, with the raw adelic mass-weighted Hermitian metric of K038/R030. Its coefficient type is opposite to the initial row, as in K058; its ramified summand is the unitary dual of the source Mram. Replacing every diamond by its inverse gives the same complete averages and kernel condition. Thus all formulas below apply in this actual detector convention. The strict T level and coefficient type are unchanged. If the class metric omits the common finite compact volume, divide both sides of each displayed norm identity by that same known volume. For each selected place v above T, write D\_v=k\_v\^{}times, q\_v=T, and let R\_v(s) be the actual diamond action. The source norm is unweighted: \nolinkurl{ScaledRamifiedLattice.lean84–91} defines the sum, and155--175 defines Mram as the intersection of its kernels. Record.lean1811--1848 identifies this with the adelic dilation action; K038(CA27a) records the exact scope.

The operators are unitary: they permute decorated classes, transport the coefficient metric by isometries, and preserve automorphism-group orders. They commute at different places. Therefore

\[
 A_v={1\over q_v-1}\sum_{s\in D_v}R_v(s),\qquad
 P_{\rm ram}=\prod_{v\mid T}(1-A_v)
 \tag{TG1}
\]

are commuting orthogonal projections, and the range of P\_ram is exactly the appropriate dual realization of the literal complex Mram subspace. The denominator in TG1 normalizes an operator; the source kernel is the same for its unnormalized norm sum. It does not assert an integral projector.

Let B be any finite-dimensional actual bottom cusp subspace of the fixed genuine-e0 block, prescribed G-generic type, maximal T level and selected real type, and let J:B→V be its actual normalized selected restriction. One may take the span produced in K058, if that report is used at its checked scope. All statements below also hold for the full finite-level bottom space. A nonzero strict restriction does not determine any of its T characters.

For an irreducible initial automorphic constituent with strict fixed vectors, the local identity delta(s)=sI\_a h\_s, h\_s in the strict covector stabilizer, makes R\_v(s) equal to its central character on that whole fixed-vector space. K052(LT4--7) then shows that a nontrivial such character forces conductor one; Mram is exactly the nontrivial-character condition. The converse from conductor one alone is not used. The H-generic cut at G does not impose this condition at T. In particular, the actual joint condition in TG1 is nontriviality at every one of the f selected T places, not merely at one.

\subsection{2. Exact orbit coordinates and all mass factors}\label{exact-orbit-coordinates-and-all-mass-factors}

At one T place put u for the actual nonzero residue covector and choose h\_s in GL\_a(O\_v) with u h\_s=s u modulo the residue prime. Write

\[
 \delta_s=\operatorname{diag}(h_s,1),\qquad
 \tau_u=\begin{pmatrix}I_a&0\\\varpi^{-1}u&1\end{pmatrix}.
 \tag{TG2}
\]

Integral lifts of u and s only change a right maximal-compact factor. Direct multiplication gives

\[
 \delta_s\tau_u=\tau_{s^{-1}u}\delta_s,
 \qquad
 \tau_u^{-1}\delta_s\tau_u
  =\begin{pmatrix}h_s&0\\\varpi^{-1}(1-s)u&1\end{pmatrix}
       \pmod{\text{right }K_{a+1}}.
 \tag{TG3}
\]

The first equality is literal when u h\_s=s u for the chosen lifts; the right-coset statement is enough in general. Thus the fixed projective-T orbit over an H class x is represented by the same x and all the lower translates tau\_(s\^{}-1 u). It is not a family of newly chosen Schwartz functions. The finite T action commutes with the actual G character and finite e0 projector. The selected real fibre is unchanged.

Put D=product\_v D\_v. For a projective-T class x, let Gamma\_x be its actual rational stabilizer, acting on D and on the coefficient fibres. Let v\_P be its projective-level compact volume, and let v\_H=v\_P/\textbar D\textbar{} be the strict-level volume. The strict classes over x are the Gamma\_x orbits on D, with their actual coefficient transports. For a descended vector F=(F\_s), orbit-stabilizer counting gives the exact identity

\[
 \sum_{[s]\in\Gamma_x\backslash D}
       {v_H\over|\Gamma_{x,s}|}\|F_s\|^2
  ={v_P\over|\Gamma_x|}\,{1\over|D|}\sum_{s\in D}\|F_s\|^2.
 \tag{TG4}
\]

This requires neither a free action nor a scalar stabilizer. Norms are constant on an orbit because the coefficient transport is unitary. It follows by \textbar orbit(s)\textbar=\textbar Gamma\_x\textbar/\textbar Gamma\_(x,s)\textbar. Formula TG4 is the strict/projective mass comparison of K038, now retained for the entire joint dilation orbit. In particular one cannot replace its strict sum by q-1 identical kernel values.

Write J\_(x,s) for evaluation of J on the lifted point x delta\_s tau, including the actual selected fibre transport. For every subset S of the selected T places, put D\_S=product\_(v in S)D\_v. The empty product is the identity. The exact projected Hilbert--Schmidt energy is

\[
 \begin{split}
 \|P_{\rm ram}J\|_{\rm HS}^2
  =\sum_x {v_P\over|\Gamma_x|}
   \sum_{S\subseteq\{v\mid T\}}{(-1)^{|S|}\over|D|\,|D_S|}
    \sum_{t\in D}\sum_{s\in D_S}
       \operatorname{Tr}\bigl(J_{(x,t)}^*J_{(x,ts)}\bigr).
 \end{split}
 \tag{TG5}
\]

The inverse convention for the right diamond can replace s by s\^{}-1; the complete group sum is identical. The simultaneous conjugate terms make TG5 real. To prove it, expand TG1, use self-adjointness and idempotence to write \textbar\textbar P J\textbar\textbar²=Tr(J*P J), and apply TG4 to each scalar product. This is a finite identity, so no interchange with an infinite theta or Poincare sum is involved.

Equivalently, writing P\_chi for the probability character projector on D, the same nonnegative number is

\[
 \|P_{\rm ram}J\|_{\rm HS}^2
   =\sum_{\chi:\ \chi_v\ne1\ (\forall v\mid T)}
                    \|P_\chi J\|_{\rm HS}^2.
 \tag{TG6}
\]

TG5 is the exact extra Gram that a finite-character projection of the K058 construction has to control. Its empty-S term is the full strict restriction energy; its other terms are actual cross-point kernels. The positive selected Gram at one point in K058 does not evaluate them. The ambient local spherical Gram of K052 is a different Hilbert pairing. A restricted global H-period form is only H-invariant; replacing it by the ambient G Hilbert form would require another proved identity.

\subsection{3. The actual local relative double coset}\label{the-actual-local-relative-double-coset}

For s not equal to1 in the residue field, the second matrix in TG3 has Cartan invariants

\[
                       (-1,0^{a-1},1).
 \tag{TG7}
\]

Indeed conjugate u integrally to the last coordinate and take h\_s=diag(I\_(a-1),s). The only nontrivial block is {[}{[}s,0{]},{[}varpi\^{}-1(1-s),1{]}{]}. Its entry ideal has valuation-1 and its determinant is a unit; hence its two elementary divisors have valuations -1 and1. The other a-1 coordinates are integral identity factors. Different choices of h\_s lie in the same relevant strict-level coset and do not change this maximal-compact double coset.

Thus cross terms in TG5 belong, after transporting the fixed frame, to rational correspondences with denominator T. LIFT42 bounds integral Gamma elements. Its field-norm proof cannot be used on this larger set without an additional argument.

\subsection{4. Explicit failure of the ambient short-correspondence bound}\label{explicit-failure-of-the-ambient-short-correspondence-bound}

Use the actual Q=g0g1 and rational prime T. Both hosts contain i with bar(i)=-i. Put J=diag(I\_a,-1), e=e\_(a+1), v=e\_1+e and

\[
                 U=I+{iQ\over T}\,v v^*J.
 \tag{TG8}
\]

Since v\emph{Jv=0, the rank-one operator N=v v}J satisfies N²=0. Moreover N\emph{J=JN. With alpha=iQ/T and bar(alpha)=-alpha, direct expansion gives U}JU=J. The determinant lemma gives det U=1. Therefore U is an actual rational special-unitary matrix.

It is integral away T and congruent to I at both G primes. At every selected T place the nontrivial2 by2 block has an entry of valuation-1 and determinant1. Its inverse has the same bound; its Cartan invariants are exactly TG7. Yet at every real label

\[
        \|U-I\|_{\rm op}={2Q\over T}
          <\begin{cases}2^{-953}&J14,\\2^{-923}&J56.\end{cases}
 \tag{TG9}
\]

The equality uses \textbar\textbar v\textbar\textbar²=2 and \textbar\textbar v*J\textbar\textbar=sqrt2. The actual prime bounds T\textgreater=2\^{}1000,Q\textless2\^{}45 and T\textgreater=2\^{}995,Q\textless2\^{}71 are certified in LIFT41. At G the reduction of U is the identity, so the normalized partial trace used in LIFT42 is exactly1, for the same genuine goodTuple and both mixed projectors. Consequently there is no analogue of LIFT42's uniform small noncentral partial trace on this ambient T-Hecke set. This is an explicit obstruction to that proposed proof, not a theorem that U occurs in the specific packet cross term in TG5.

\subsection{5. The literal near packet excludes this example}\label{the-literal-near-packet-excludes-this-example}

There is a useful finer restriction, absent from the ambient double coset. Retain the exact LIFT41 notation k=Q²,

\[
 d=k(-7+2\sqrt{-7}),\quad r=(w,d),\quad
 v_0=\pi e+r=(w,\pi_k),\quad z=v_0/\bar\pi.
 \tag{TG10}
\]

Then at every real label

\[
 r^*Jr=-14Q^2<0,\qquad
              \|r\|^2=154Q^4-14Q^2<154Q^4.
 \tag{TG11}
\]

Here sum\textbar w\_i\textbar²=77Q\^{}4-14Q² and \textbar d\textbar²=77Q\^{}4, exactly as in LIFT41. The vector r is a genuine negative vector, not an isotropic residue approximation.

If a rational relative matrix R connects the fixed LIFT41 point to another point of its T scalar orbit in the same maximal finite class, then at every selected T place it necessarily satisfies

\[
                            Rz\in O_{E,v}^{a+1}.
 \tag{TG12}
\]

To see this with the actual frames, write gamma\^{}-1=x tau k with k in the maximal compact, as in LIFT41(NP9). Then k\^{}-1 e=gamma e=z. If h\_s x delta\_s tau lies in the maximal compact and R=h\_s gamma\^{}-1, then R lies in K tau\^{}-1 delta\_s\^{}-1 tau k. This implies TG12 because tau\^{}-1 delta\_s\^{}-1 tau fixes the last column e. Choices of representatives only change the displayed integral compact factor. This is a necessary condition, not a sufficient characterization of the orbit.

The matrix U in TG8 fails TG12. Its added last-column term is (iQ/T)v(v*Jz), and

\[
              v^*Jv_0=2Q-\pi-d\equiv2Q-d\pmod\pi.
 \tag{TG13}
\]

The nonzero quadratic integer 2Q-d has norm 77Q\^{}4+28Q\^{}3+4Q², strictly between0 and T for both actual hosts. It cannot be divisible by the prime pi in Q(sqrt(-7)), hence by any selected degree-one extension prime. Thus U z is nonintegral there. The explicit tiny ambient correspondence is therefore excluded from the fixed scalar packet, rather than being an example of a large uncontrolled term in TG5.

\subsection{6. Quantitative exclusion for a whole unipotent family}\label{quantitative-exclusion-for-a-whole-unipotent-family}

The same negative residue yields a uniform necessary displacement. Let a0 be a nonzero integer of F that is a unit at all selected T places, and let v be a nonzero isotropic vector in O\_E\^{}(a+1), primitive at each selected T place. Suppose

\[
       R_v=I+{iQ a_0\over T}v v^*J
 \tag{TG14}
\]

satisfies TG12. The vector z is integral at those places, while the coefficient has exact valuation-1. Primitivity of v implies pi divides v\emph{Jv\_0 in O\_E. Reducing v\_0=pi e+r gives pi divides b:=v}Jr.~The integer b is nonzero: at any real label the orthogonal complement of the negative vector r is positive definite, so it contains no nonzero isotropic v.

Since Norm\_(E/Q)(pi)=T\^{}f and pi divides b,

\[
                 \prod_{\tau\in\Phi}|b_\tau|^2\ge T^f.
 \tag{TG15}
\]

Cauchy--Schwarz and TG11 give product\_tau \textbar\textbar v\_tau\textbar\textbar² \textgreater= T\^{}f/(154Q\^{}4)\^{}f.~Also product\_tau \textbar a0\_tau\textbar=\textbar Norm\_(F/Q)(a0)\textbar\textgreater=1. The exact rank-one operator norm is \textbar\textbar R\_(v,tau)-I\textbar\textbar=Q\textbar a0\_tau\textbar\textbar\textbar v\_tau\textbar\textbar²/T. Multiplying over all labels proves

\[
 \left(\prod_\tau\|R_{v,\tau}-I\|_{\rm op}\right)^{1/f}
       >{1\over154Q^3},\qquad
 \max_\tau\|R_{v,\tau}-I\|_{\rm op}
       >\begin{cases}2^{-143}&J14,\\2^{-221}&J56.\end{cases}
 \tag{TG16}
\]

Strictness follows from the strict Euclidean bound in TG11. This keeps the actual coefficient field, all selected T places and every real embedding. It does not assume one small archimedean value controls the other embeddings. It is a necessary bound for this specified primitive rank-one family. Arbitrary relative matrices, partial T subsets and the complete Poincare cross Gram are not reduced to this family here.

\subsection{7. What remains to be evaluated}\label{what-remains-to-be-evaluated}

The precise sufficient and necessary condition for the proposed route is positivity of TG5 on the actual finite bottom subspace. A quantitative result would also bound that energy, or one positive singular value, in the original mass and PBW metrics. Existing results establish the empty-subset contribution and local branching tests; they do not yet bound the other subset terms. TG9 proves that maximal-compact relative position alone loses the finite-character estimate. TG12--16 give a specific finer packet restriction that removes the obvious obstruction and a whole very small primitive unipotent family.

The K049 small selected value at its cusp packet point cannot presently be compared with LIFT41's favorable point by a scalar-orbit argument: LIFT41 does not identify its positive complement with the standard positive lattice underlying that cusp point. No such global integral isometry or projective-T orbit identification has been proved.

All statements are about the unchanged finite family. No new prime, principal-T level, canonical twist, changed goodTuple or unconstrained Schwartz vector has been introduced. The full Mram detector and the prescribed terminal theta pairing remain open after this result.

\subsection{8. Source and verification scope}\label{source-and-verification-scope}

The source clauses directly read are \nolinkurl{ScaledRamifiedLattice.lean84–91,} 155--175 and the carried exact transport in Record.lean1811--1848. Accepted inputs are K038(CA27a and strict/projective mass correction), K052(LT1--7), LIFT41(NP1--11), LIFT42(FC1--15), and the explicitly scoped K058 selected Gram when available. All new matrix, orbit-counting, projection and field-norm arguments are proved above. The certificate checks exact finite projection identities, strict orbit masses, local Cartan minors, rational unipotent equations and both actual numerical bounds. It does not evaluate any automorphic period or the unknown cross terms in TG5.

\clearpage

\section{AM-DO-054 --- the literal T law and the common tensor for the original two vertical tests}\label{am-do-054-the-literal-t-law-and-the-common-tensor-for-the-original-two-vertical-tests}

Author: GPT-6 Astra Pro, ALP\_MACHINE /root/definite\_oscillator. Inherited reasoning effort is not separately exposed. 22 November 2026 UTC.

Status: author whole-proof check PASS, 22 November 2026 UTC. The complete proof, its literal source clauses, and the complete companion certificate have been read. Independent review is pending.

This finding proves the remaining T interface in the actual two-step mixed construction. The source TLaw requires a period identity involving the original vertical tests. It expressly allows a fixed finite Heisenberg operator on the horizontal common tensor. DO051's operator and original-vacuum period are exactly the two parts of that clause. Consequently no equality between the original vertical map and its E0-modified Schwartz map is required.

The result includes all local vacua, arbitrary off-T Heisenberg insertions, all places above T, the same tensor before the auxiliary tuple varies, and the exact last-step CVEq scalar. It supplies an analytic ValLastX identity for the constructed orbit-indexed finite family. It does not construct the old HostRaw, identify every arithmetic row with an externally frozen record, or prove that an actual restricted row survives Mram.

\subsection{1. Fixed data and the exact source clause}\label{fixed-data-and-the-exact-source-clause}

Keep either actual host K031 or K032, the actual T prime and pi of DO040, and the fixed rank profile

\[
 a\in\{42,168\},\quad r=a+1,\quad b=1802,\quad
 c=2b-a,\quad N=c-1,\quad m=2.
 \tag{TC1}
\]

The horizontal spaces are A=I\_a, B=I\_b, X=I\_c, and the vertical spaces are V=A perpendicular to a negative line, B, and H=I\_N. All occurrences of A,B,H use the same compact-quotient Haar measures as the accepted theta and period constructions. The rank-zero middle identification V\_1=B=X\_1 is the identity, with the same Haar measure. Thus the literal Comm.VolLast equation holds with no additional middle volume factor (HzUnit.lean23--24).

Write Tcal for the f places of F above T. The DO051 frames and line are fixed:

\[
 \zeta=\bar\pi,\quad \zeta_0=\pi,\quad
 \tau=-\xi_0\bar\pi/\pi,\quad
 \epsilon_v=\operatorname{red}(-\xi_0\varpi_v/T).
 \tag{TC2}
\]

Here varpi\_v=T is available at every place. The formulas also retain any consistently chosen unit multiple of this uniformizer. The negative embedding is the actual primitive norm-minus-T line of LIFT30. Its finite covector transport to the prescribed primitive T covector is kept in both the frame and input action.

For a finite tuple xi put m\_j=mu\_j xi. The actual unbarred vertical labels are

\[
       (m_0\eta^{-1},m_0),\qquad(m_1,m_1\eta^{-1}),
 \tag{TC3}
\]

and the horizontal labels are (m\_0,m\_0), (m\_1,m\_1). In the source WeilOK.last choose chi\_A=m\_0, chi\_z=eta\^{}\{-1\}, chi\_S'=m\_1, chi\_T=m\_1 eta\^{}\{-1\}. Then chi\_A chi\_z=m\_0 eta\^{}\{-1\} and chi\_T chi\_z\^{}\{-1\}=m\_1, exactly. The auxiliary exchanged line is fixed at group argument1, so its other character does not enter the common tensor. This is the checked R016/DO050 convention, including conjugation of the negative line. Every one of these characters is unramified at T: the stage characters are K040's actual characters, eta is the accepted finite-unramified odd character, and the allowed tuple support Q is disjoint from T.

For any two vertical theta functions Psi\_0,Psi\_1 let

\[
        L(\Psi_0,\Psi_1)
        =\operatorname{lastBi}_{o,o_T,\tau}
                           (\Psi_0,\Psi_1).
 \tag{TC4}
\]

The rank-zero flatBF and ext0 transports are understood; they are the identity maps in these chosen rank-zero coordinates. For a horizontal tensor C define the actual bilinear terminal functional

\[
 \mathcal P_{t,\phi,F}(C)
   =\int_{[U(H)]}F(h)\,
        \operatorname{biLift}_{t}(\phi,C)(\iota_Hh)\,dh.
 \tag{TC5}
\]

The source TLaw last clause, PlaceLaws.lean216--226, asks precisely for

\[
\begin{split}
 &\operatorname{BiVacAt}_{\mathcal T}
                       \bigl(O_{\mathcal T}L(\Psi_0,\Psi_1)\bigr),\\
 &\mathcal P_{t,\phi,F}\bigl(D_{\rm off}L(\Psi_0,\Psi_1)\bigr)
   =c_T(t)\,
     \mathcal P_{t,S_{\mathcal T}\phi,F}
        \bigl(D_{\rm off}D_{\mathcal T}
                         O_{\mathcal T}L(\Psi_0,\Psi_1)\bigr).
\end{split} \tag{TC6}
\]

It quantifies over all Psi\_j with the indicated T-vacuum condition, all two-sided finite Heisenberg insertions D\_off supported outside Tcal, all phi in RamAd\_0, and all continuous automorphic F satisfying the stated terminal compact condition. Here D\_T=comb(dT\_0) tensor comb(dT\_1). The other TLaw clause has j+2(sigma+1)+2=2 and is empty. Equation TC6 never requires E0 Psi\_0=Psi\_0 or equality of the two vertical operators before taking the period.

\subsection{2. Actual finite operator and all-vacuum quantification}\label{actual-finite-operator-and-all-vacuum-quantification}

Use the explicit E0\_v, O0\_v, D'\_v and O\_v=D'\_v O0\_v of DO051(FT18a--FT19). They are actual finite sums of Heisenberg operators, supported at v. In particular this is an operator of the literal BiOp/OpAt type, not an arbitrary endomorphism of a finite quotient. Its two Finsupp entries have zero infinite part and may be retained with BiOp.n=1, after collecting the actual central phases.

Let V0\_v be the original first glued vacuum and V1\_v the original second integral vacuum. DO051 proves the entire-vector identities

\[
\begin{split}
 O_v L_v(V0_v,V1_v)&=\gamma_v q_v^{b/2}V_{{\rm bas},v},\\
 D_vO_v L_v(V0_v,V1_v)&=L_v(E0_vV0_v,V1_v),\\
 D_vV_{{\rm bas},v}&=R_{T,v},\qquad |\gamma_v|=1.
\end{split} \tag{TC7}
\]

R\_T,v is the product of the two unchanged raw incidence functions, with coefficient1. The phase gamma\_v is defined by the actual lastBi and the fixed local vacuum models; it is retained. Neither it nor O\_v depends on the varying tuple character.

These identities hold with arbitrary off-v factors and for every vector satisfying the same local vacuum condition. Here is the representation-theoretic justification of the quantifier. In a Schroedinger realization the self-dual lattice vacuum space is one dimensional. Equivalently, impose the lattice eigenrelations in a sufficiently deep finite Heisenberg quotient: the nondegenerate central-character representation is its unique Schroedinger module, and the vacuum eigenspace in that module is one dimensional. The multiplicity space consists of the other local and real factors. Both sides of TC7 are the same finite Heisenberg operator on that one-dimensional factor, tensor the identity on this multiplicity space. Thus equality on the fixed vacuum proves equality on the whole required vacuum subspace. Enlarging the finite depth gives the same operators and identities, so this argument is independent of the chosen quotient.

This also follows from the actual source Over/Reps quantification: SmallPair.lean28--33 defines Reps by its action on every vacuum; IsSumFrame.Over and PlacesToAdeles.OverBasic fix the lattice indicator to the original vacuum with coefficient1. DO051 constructed its operators under that convention. No new scalar may be introduced when the off-place vector changes.

Form the products over the distinct T places:

\[
 O_{\mathcal T}=\prod_{v\in\mathcal T}O_v,\quad
 D_{\mathcal T}=\prod_{v\in\mathcal T}D_v,\quad
 E_{\mathcal T}=\prod_{v\in\mathcal T}E0_v,\quad
 \Gamma_T=\prod_{v\in\mathcal T}\gamma_v.
 \tag{TC8}
\]

Disjoint-place symplectic pairings are zero, so these finite Heisenberg operators commute exactly, without a residual central character. TC7 therefore gives, for every allowed pair,

\[
 D_{\mathcal T}O_{\mathcal T}L(\Psi_0,\Psi_1)
                =L(E_{\mathcal T}\Psi_0,\Psi_1),
 \qquad
 \operatorname{BiVacAt}_{\mathcal T}
                (O_{\mathcal T}L(\Psi_0,\Psi_1)).
 \tag{TC9}
\]

The tensor on the right of the first equation is modified; the vertical input on the left of the period in TC6 remains original.

\subsection{3. The original period law, with all off-T insertions}\label{the-original-period-law-with-all-off-t-insertions}

Take the actual ramified space

\[
 \mathcal R_{\mathcal T}
 =\{\phi:\ \phi\text{ is continuous automorphic and }
       \operatorname{RamAt}(\beta_{0,v},P_v,\phi)
                   \text{ for every }v\in\mathcal T\}.
 \tag{TC10}
\]

This is literally TDat.Ram, VacOfShape.lean43. Put

\[
 S_{\mathcal T}\phi
 =\prod_{v\in\mathcal T}
        \operatorname{eval}_{P_v}
           \bigl(\operatorname{dia}(\epsilon_v)
                  \operatorname{FFStar}_v
                        \operatorname{formAt}_{P_v}(\phi)\bigr).
 \tag{TC11}
\]

The product means the fixed source fold order, VacOfShape.lean45--46; different-place actions commute. It is precisely the source Sad. DO051(FT7--FT8) proves that each factor preserves the unweighted ramified condition. Applying a different-place factor preserves the earlier conditions as well. Consequently S\_T preserves R\_T. This is the ram field of VCom.Is, not an unproved nonzero projection onto R\_T.

The original-vacuum law DO051(FT10--FT13) is the actual SmallOne law, with

\[
 r_v=q_v^{-b/2}(q_v^b-1)>0,\qquad e_v=1,\qquad
 c_{T,v}(t)=r_v\,m_0(\operatorname{idele}_v\varpi_v).
 \tag{TC12}
\]

Its proof uses the original glued first indicator, the original second vacuum and the unweighted ramified condition. The two zero Fourier terms vanish on that condition, the remaining finite sum is exactly FFStar, and the middle compact moves the q\_v\^{}b-1 nonzero covectors transitively. No new Haar index appears. The accepted result is for every continuous automorphic terminal function, which is at least as strong as the terminal compact condition in TLaw. SmallPair.lean134--164 records precisely these one-place and all-place quantifiers.

Apply this law successively at every v. At each stage the original vacuum at any not-yet-used place is unchanged, since all already used operators are supported elsewhere; TC11 retains its required ramified condition. The constants multiply, giving

\[
 \mathcal P_{t,\phi,F}\bigl(L(\Psi_0,\Psi_1)\bigr)
 =c_T(t)\,
   \mathcal P_{t,S_{\mathcal T}\phi,F}
                    \bigl(L(E_{\mathcal T}\Psi_0,\Psi_1)\bigr),
\qquad
 c_T(t)=T^{-fb/2}(T^b-1)^f
               \prod_{v\mid T}m_0(\operatorname{idele}_v\varpi_v).
 \tag{TC13}
\]

This is the source monomial with the explicit r\_v,e\_v above. No field norm or arbitrary inverse of m\_0 has replaced that product.

Now let D\_off=comb(d) tensor comb(d') be arbitrary and supported outside Tcal. The line exchange is an actual Heisenberg intertwiner. Pulling back one such product through it gives a product Weyl operator on the two vertical factors, with its actual central phase. Pulling back the finite sum gives

\[
 D_{\rm off}L(\Psi_0,\Psi_1)
       =\sum_{\ell=1}^M
           L(A_\ell\Psi_0,B_\ell\Psi_1),
 \tag{TC14}
\]

where all A\_ell,B\_ell are supported outside Tcal. This is exactly the accepted source LastOp second clause, SmallPair.lean98--100; the support is preserved because the line map is local. These operators preserve the T vacuum conditions. They commute with E\_T and with all T operators, including their central phases. Apply TC13 to each summand of TC14, then use linearity and TC9. This proves both equations in TC6 with the displayed c\_T. The finite sums and the compact A,B,H integrals justify the interchanges; no sum over a noncompact bottom quotient has been introduced. Thus the full literal last-step TLaw is proved.

The proof is the concrete version of \nolinkurl{LevelPrimeLawLastTwoStepsAllTuples.lean,} TLastAll.main (statement at line35; its hdec and final SmallAll application). The settled finite laws are being instantiated, not re-audited.

\subsection{4. Vacuum away from T and the actual common real tensor}\label{vacuum-away-from-t-and-the-actual-common-real-tensor}

For completeness the current standard lattices also give the remaining off-T vacuum assertion used to construct the common C. The primitive matrix g of DO051 has entries integral at every finite place outside T: A=(T+1)/2 and B=(T-1)/2 are rational integers and both pi and bar(pi) are units there. It is unitary with determinant1, so its inverse is integral there as well. The negative sum decomposition therefore identifies the original standard lattice with O\_E\^{}a perpendicular to O\_E; the positive terminal decomposition does the same.

There is no uncanceled different in the exchanged line. In the standard Gram matrices (-1,I\_b), the literal ofAdj formula SkewGramHalfModule.lean24 gives

\[
 \operatorname{kap}_{\tau}(w)
       =\xi_0^{-1}\overline{\tau}\,w^*
       =(\pi/\bar\pi)\,w^* .
 \tag{TC15}
\]

The actual mixed principal lattice is O\_E, preserved by conjugation. The scalar pi/bar(pi) is an O\_E unit away from T. Hence TC15 preserves the exchanged integral lattice at every such place, including the absolute different and dyadic places. In the literal IsLatVac convention, conjugation sends a split idempotent e to bar(e), and multiplication by this unit preserves its half-space; the doubled-lattice and half-idempotent clauses are both preserved \nolinkurl{(LocalNewformFunction.lean199–202).} Thus the actual line intertwining carries both lattice vacuum conditions. This proves the last part of VacLaw, PlaceLaws.lean203--205, on these standard mixed lattices. No claim about arbitrary lattices or a new compact-splitting comparison is used.

Fix any real polynomial inputs P0,P1 and their original unit-vacuum theta functions Psi\_0(P0),Psi\_1(P1). Define

\[
 C(P0,P1)
       =O_{\mathcal T}L\bigl(\Psi_0(P0),\Psi_1(P1)\bigr).
 \tag{TC16}
\]

It is the actual tensor, chosen before xi. The T part is

\[
                  \Gamma_T T^{fb/2}
                         V_{{\rm bas},\mathcal T}.
 \tag{TC17}
\]

At all other finite places it has the just-proved horizontal vacuum conditions. The T-supported operator does not change them. The line exchange and finite Heisenberg operators preserve the two-variable theta laws and real smoothness. Combining the on-T and off-T lattice conditions proves the literal Comm.IsReal=C is BiVac and BiSmooth, ValLastX.lean67--68. Equivalently, this is the explicit last case of the accepted VReal.real\_of/v\_real construction in \nolinkurl{SeeSawCommonRealDatumBiThetaVacuumContraction.lean.}

No real simple tensor is substituted for C. In the selected coefficient calculation its full real factor is exactly DO050's recoupling. If PN15 primitive embeddings are used with a Fischer-isometric Pieri input, DO050's scalar applies directly. Using a higher primitive integral Pieri vector requires LIFT40's additional q\_lambda. None of these choices changes TC6.

The common tensor is independent of xi because lastBi is defined only by the fixed orthogonal maps, tau and the theta-function arguments, \nolinkurl{CombAtSingleZero.lean30–34.} At rank zero the flatBF map adds no character factor. The real tuple character is trivial, and all finite changes of splitting occur in the Weil actions, not in TC16. The source valLast proof makes exactly this choice at one reference tuple and proves it independent by flat-independence \nolinkurl{(SeeSawLastTwoStepsDoubleLiftPeriodValue.lean326–339).}

\subsection{5. Exact last-step scalar and analytic ValLastX}\label{exact-last-step-scalar-and-analytic-vallastx}

Keep the normalized G tests which agree with the original vertical maps on their actual generic P-character spaces, as in DO052. Keep all raw horizontal G,Q,T incidence operators. Let I1\_v={[}GL\_b(k\_v):P\_(r,t)(k\_v){]} and I2\_v={[}GL\_N(k\_v):P\_(r,2t)(k\_v){]} at Q. The DO053 orbit-indexed family has

\[
 \epsilon_0=\delta_1,\quad w_0=1,\quad
 \epsilon_T=\Bigl(\prod_{v\in Q}I2_v\Bigr)\delta_1,\quad
 c_Q(t)=\prod_{v\in Q}(I1_v I2_v)^{-1}.
 \tag{TC18}
\]

Its actual moved traces remain delta\_b0 and (product I2\_v) delta\_bT. Its weighted second horizontal Q operator retains I1\_v times the actual orbit size h\_Mv. Thus neither terminal averaging nor full-Y counting is erased.

The full G and passive Q comparisons R035/DO052/R034 give

\[
\begin{split}
 s_0(t)s_1(t)
  &=\prod_{v\in G}q_v^{-a(a-1)-b/2}m_0(q_v)^{-1}
                  \prod_{v\in Q}\chi_v(\bar\pi)^{-1},\\
 c_G&=1.
\end{split} \tag{TC19}
\]

The individual two factors 1/\textbar U\_a(k\_v)\textbar{} occur in s\_0,s\_1 as already fixed. Taking the unnormalized generic tests instead changes both these equations and the vertical maps; it is not the convention here.

Define the actual last-step value scalar by the computed factors

\[
 \boxed{\quad
 c_V(t)=
   \left[\prod_{v\in G}q_v^{-a(a-1)-b/2}m_0(q_v)^{-1}\right]
   \left[\prod_{v\in Q}
             {\chi_v(\bar\pi)^{-1}\over I1_v I2_v}\right]
   T^{-fb/2}(T^b-1)^f
       \prod_{v\mid T}m_0(\operatorname{idele}_v\varpi_v).
 \quad} \tag{TC20}
\]

Every factor is nonzero. This is precisely CVEq.last, PlaceLaws.lean268. There is no kap factor in this rank-zero last clause. The middle quotient Haar equality was fixed in TC1, so none is being hidden in c\_V.

Here is the resulting analytic identity, stated without assuming an unspecified arithmetic K map. Let the actual initial family phi\_alpha have the DO053 finite covariance and belong to R\_T. Let F be an actual terminal automorphic vector at the prescribed finite type, maximal at T, and retain the moved b0,bT. Let t.two use the actual horizontal finite G,Q,T operators and the DO053 orbit weights. Then

\[
 \operatorname{perF}_0
       \bigl(\phi,\operatorname{lift}_0(P0,
                                \operatorname{lift}_1(P1,F))\bigr)
 =c_V(t)\,
       \operatorname{perT}_2
             \bigl(t.\operatorname{two}_0(S_{\mathcal T}\phi,C(P0,P1)),
                    F\bigr).
 \tag{TC21}
\]

\textbf{Proof.} The actual rank-zero compact seesaw (R016, with the fixed line map) rewrites the left period as the biLift terminal integral with tensor L, using the same Haar at each occurrence. Insert the two actual VForm scalars s\_0,s\_1. Apply DO053's QLast to the unchanged Q kernels, with initial identity trace and retained terminal trace. Apply R035's entire G identity, whose constant is c\_G=1 in the present raw eG convention. The resulting GQ operators have support outside T. TC6 then applies to each finite index term. It replaces the initial input by the actual Sad and inserts D\_T O\_T. Disjoint-place commutation puts D\_Q D\_G D\_T in the actual dd order. The terminal Q trace commutes with the T compact, so its presence does not invalidate the terminal condition. The product of the constants is exactly TC20; the remaining tensor is TC16. All sums are finite and all three period integrals are over the fixed compact A,B,H quotients. This proves TC21. QED.

This is the literal mathematical ValLastX equation for the constructed orbit-indexed analytic family, not a definition of tail by that period. It is also exactly the sequence of left\_side, QLast, right\_side and CVEq in \nolinkurl{SeeSawLastTwoStepsDoubleLiftPeriodValue.lean204–240,249–319,} 321--360. At m=2 the nonlast OffLaw, GSupp, QLaw, QFin, ValX and LocFin quantifiers with j+2(sigma+1)+2=2 are empty. Their disappearance does not remove VForm's actual row, operator and level requirements. Those used in TC21 are the stated genuine finite kernels and input/detector spaces.

\subsection{6. Costs, distinctions, and remaining arithmetic work}\label{costs-distinctions-and-remaining-arithmetic-work}

The operator TC8 has exactly the DO051 finite coefficient bound

\[
 O_{\mathcal T}\ \text{has coefficients in }\
       \mathbb Z[\zeta_T,1/T],\qquad
 \operatorname{den}(O_{\mathcal T})\mid T^{f(a+b)},\qquad
 \|O_{\mathcal T}\|_{\ell^1}\le T^{f(a+2b)}.
 \tag{TC22}
\]

These are upper bounds after expanding and collecting the actual Heisenberg operators. They do not assert optimal denominators. Since T is prime to11, the displayed coefficients add no 11-adic denominator. This says nothing further about the independently fixed scalar phase Gamma\_T or the integral image of an inverse automorphic map.

The positive T-vacuum factor in TC17 and the T\^{}\{-fb/2\} in TC20 cancel in the complete period product:

\[
        c_T(t)\,\Gamma_T T^{fb/2}
          =\Gamma_T (T^b-1)^f
                    \prod_{v\mid T}m_0(\operatorname{idele}_v\varpi_v).
 \tag{TC23}
\]

The source FFStar/diamond still acts on the input. TC23 is not an assertion of norm preservation or of an arithmetic gain. All Q indices and moved trace coefficients remain as in TC18.

What is now proved is the actual finite T constructor and its fixed common-tensor role for the original vertical period, together with the explicit analytic last-two identity TC21. There is no remaining requirement to prove that E0 leaves the entire vertical map unchanged; the source never imposed it. The old source HostRaw is still not instantiated. To invoke its MainX conclusion for an arithmetic family one must use the independently fixed arithmetic rows and prove their KerX/PerX, stability and level interfaces in the same finite indexing. DO053's integrated orbit reduction cannot be silently called a bijection with a previously frozen all-Y index record. Most importantly, TC21 applies to inputs already in Mram; it does not prove that the nonzero K058 restricted row has a nonzero Mram component or belongs to the prescribed theta image.

\subsection{7. Source and verification record}\label{source-and-verification-record-2}

All source locators are at the stipulated pinned commit 7cbf7a0c32fca1badf05931879a37c3f8a2b03b4. Source equations are accepted as instructed; no Lean execution or validity audit has been performed.

\begin{itemize}
\tightlist
\item
  \nolinkurl{PlaceLaws.lean21–27,43–47,112–117,128–147,199–226,} 243--251 and266--284: BiOp, lst, finite support/ramification, VForm, VacLaw/TLaw, QLast and the exact last-step scalar.
\item
  \nolinkurl{VacOfShape.lean22–52,69–104,114–120:} TDat, fixed fold and character monomial, frame/line shape, unramified conditions.
\item
  \nolinkurl{SmallPair.lean26–44,92–100,102–164:} OpAt/Reps, support-preserving pullback, compact seesaw, LastTAll and SmallOne/SmallAll.
\item
  \nolinkurl{LevelPrimeLawLastTwoStepsAllTuples.lean,} TLastAll.main, line35 and its explicit lastOp/SmallAll assembly.
\item
  \nolinkurl{SeeSawCommonRealDatumBiThetaVacuumContraction.lean,} real\_of and v\_real: assembling on/off finite vacuum and joint real smoothness of the same common tensor.
\item
  \nolinkurl{SeeSawLastTwoStepsDoubleLiftPeriodValue.lean204–360:} actual left period, right period and reference-tuple tensor.
\item
  \nolinkurl{LocalNewformFunction.lean199–202} and \nolinkurl{SkewGramHalfModule.lean21–28:} literal lattice vacuum and the xi0 cancellation in kap.
\end{itemize}

The substantive predecessors are DO051, DO052, DO053, R016, R034--035, DO050 and the fixed primitive normalization LIFT39.

The author reconstructed all TC1--TC23, including the all-vacuum multiplicity argument, the pullback of arbitrary off-T insertions, the actual xi0 cancellation in kap, the reference-tuple choice, the complete original-map period identity and every Q/G/T factor. The companion \nolinkurl{definite_oscillator_literal_t_law_certificate.py} passes 2,128 exact assertions: finite character Gram and augmentation inverse, diamond compatibility, disjoint-place and multiplicity tensor identities, parabolic/orbit weights, complete-chain vacuum exponents, the actual two G exponents, and rational complex line-metric identities. It does not evaluate an automorphic period or certify a nonzero arithmetic row. No Lean is run. Their established calculations are cited rather than repeated. The new result is their fully quantified source interface, the off-T insertion extension, actual standard off-T lattice check, and the original-test ValLastX assembly with all constants.

\clearpage

\section{AM-K059 --- bounded integral terminal spanning families for the actual coupled map}\label{am-k059-bounded-integral-terminal-spanning-families-for-the-actual-coupled-map}

Author: GPT-6 Astra Pro, ALP\_MACHINE \nolinkurl{/root/kummer_extension}. 22 November 2026 UTC. Inherited reasoning effort is not separately exposed.

\textbf{Author whole-proof check PASS --- 2026-11-22T07:56:35.227350+00:00.} The exact companion certificate passes 445 assertions. Independent reviews are separate events.

This finding removes an unnecessary integral-inverse requirement from an existential terminal nonzero argument. The full admitted terminal coefficient space has an explicitly described finite spanning family of integral vectors. A bilinear functional that is nonzero on the prescribed coupled image is nonzero on one of these vectors. No integral lift of a previously chosen analytic vector is required.

The result does not replace the prescribed image by the full final column space. Its exact nonzero premise, and the smaller premise available after the original-map seesaw is matched, are stated in Sections 6--7. The current K058 result gives a G-generic selected restriction, but its Mram projection remains unproved. Global class frames and their explicit denominator bounds remain costs.

\subsection{1. Actual modules, scalar field and class metric}\label{actual-modules-scalar-field-and-class-metric}

Fix one of the two existing mixed hosts, one allowed word, one allowed tuple, the fixed real profile, and the actual maximal29 levels. Let K be the conjugation-stable coefficient field of the actual arithmetic data, including CharT.Ok and R032's finite character list. The construction through Section 5 uses that field and makes no further extension. For any additional arithmetic kernel, K must contain its coefficients.

The T-th roots needed for DO054's O\_T are already forced by CharT.Ok. Indeed \nolinkurl{ScaledRamifiedLattice.lean141–148} gives a nontrivial additive character psi:O\_E/P→K. Choose x with psi(x)≠1. The residue characteristic is the prime T, so psi(x)\^{}T=psi(Tx)=psi(0)=1. Its order therefore is exactly T. Its powers are all T distinct roots of X\^{}T−1 in characteristic zero. Thus mu\_T is contained in K, even without using residue degree one; at a degree-one place psi(1) is primitive. The ring Z{[}zeta\_T,1/T{]} of O\_T is already contained in K{[}1/T{]}. No new extension is required after the actual Sad/FFStar/ CharT.Ok data have been supplied.

If a provisional earlier field omitted mu\_T, it did not yet carry that full arithmetic character family. Completing it at that step requires an extension of degree at most T−1. This changes the coefficient field, not the CM host E/F, f, the ranks or the CM type. It is unramified at11 because T≠11: X\^{}T−1 has unit derivative at each root over Z\_11, so its root algebra is finite etale, also after base change to any local coefficient valuation ring. No global field-degree, discriminant or DP44 normalization cost is inferred from this local observation.

The terminal definite group has rank N=3561 for J14 and N=3435 for J56. Its actual compact coefficient representation V\_T is the detector representation \nolinkurl{rowVal}, with the source rhoV convention. The coupled column uses rhoVr. Their pairing is the literal bilinear \nolinkurl{pairBK}; its Hermitian realization is linear in its first argument. The terminal lattice is

\[
 \Lambda_T=(\mathrm{mtildeGSD}) .\mathrm{intNab},\qquad
 \Lambda_T^{\Delta}=(\mathrm{mtildeGSD}) .\mathrm{intDel}.
 \tag{IS1}
\]

These are separate lattices. RowForms.lean281--286 defines them with these names. The finite level character is the inverse in \nolinkurl{mtildeGSD}, at268--269, as in DO038(38.9a). No Hermitian orthonormal lattice replaces either of them.

Let U be the actual terminal compact level, with its prescribed character, and let C be its finite arithmetic double-class set. Use the representatives h\_c and actual finite stabilizers Gamma\_c.~The coefficient family space has the form

\[
 \mathcal V_T=\bigoplus_{c\in C}V_T^{A_c(\Gamma_c)},\qquad
 \|Z\|_{\rm cls}^2
     =\sum_{c\in C}{\|Z_c\|_{\rm nfAll}^2\over|\Gamma_c|}.
 \tag{IS2}
\]

Here A\_c is the actual twisted stabilizer action of the GSD. More explicitly, if its coefficient relation is rho(gamma)v=lambda\_c(gamma)v, then A\_c(gamma)=lambda\_c(gamma)\^{}\{-1\}rho(gamma). Thus fixed vectors are exactly admitted class values. This definition retains the finite character rather than guessing its inverse from an adelic convention. Every A\_c is unitary at the chosen complex realization and preserves the literal local coefficient lattice L\_c.~Character values on the finite compact groups are roots of unity. The same assertions apply at any other complex embedding with the corresponding compact form.

Let L\_c denote the full coefficient lattice before its stabilizer condition. Then

\[
 \Lambda_T=\bigoplus_{c\in C}
                 (L_c\cap V_T^{A_c(\Gamma_c)}).
 \tag{IS3}
\]

Equations IS2--3 are the actual finite class presentation, with stabilizer inverse masses. They follow from the GSD covariance and the transported coefficient lattices, as in K037(CR20--22) and DO038 Sections 2--4. They do not require class number one or neatness.

\subsection{2. Explicit local denominator clearing}\label{explicit-local-denominator-clearing}

At each CM label tau choose the accepted integral semistandard dual-Weyl basis in the source polynomial model, and tensor these bases. Write the resulting finite K-basis as p\_1,\ldots,p\_d.~The determinant shift in the representation is retained. It has absolute value one on the compact real group but does affect finite lattice membership.

The accepted statements \nolinkurl{SSBasis}, \nolinkurl{NablaLBasis}, \nolinkurl{NablaLKBasis}, \nolinkurl{NablaLTwMem} and \nolinkurl{TensorBasis} are in GradedWeylDivisibility/ \nolinkurl{StatementD.lean24–29,43–56,} and its later twist/tensor clauses, with theorems in FinalD.lean. They supply the bases and the local lattice test. No new proof of those accepted theorems is attempted here.

For each finite place w of K choose an O\_(K,w)-basis of L\_c tensor O\_(K,w), and let T\_(c,w) be the matrix whose columns are the coordinates of the fixed p\_i in this local basis. Equivalently, the columns include the exact frame action and determinant power appearing in \nolinkurl{NablaLTwMem}. Thus a coefficient column x belongs to this local lattice if and only if T\_(c,w)x is integral. This convention fixes the direction of every inverse in the denominator formula below.

Normalize ord\_w by ord\_w(uniformizer)=1, and put e\_w=ord\_w(l) for the rational prime l below w. Define the nonnegative integers

\[
 m_{c,l}=\max_{w\mid l,\,i,j}
       \left\lceil{-\operatorname{ord}_w(T_{c,w;ji})\over e_w}
       \right\rceil_+,
 \qquad d_c=\prod_l l^{m_{c,l}}.
 \tag{IS4}
\]

Zero entries contribute no negative valuation. Only finitely many factors are nontrivial because L\_c and the standard lattice are full lattices in the same finite dimensional K-space. The local bases need not glue to an O\_K-basis; no freeness of the global lattice is assumed.

\textbf{Lemma IS5.} Each d\_c p\_i belongs to L\_c.~The family d\_c p\_i spans V\_T over K, and d\_c is a positive rational integer.

\textbf{Proof.} At w above l, the valuation of the (j,i) entry of d\_c T\_(c,w) is at least e\_w m\_(c,l)+ord\_w(T\_(c,w;ji)), which is nonnegative by IS4. Other rational prime factors of d\_c are units at w. Hence d\_c p\_i is integral at every finite place, which is membership in L\_c.~Multiplication of a K-basis by the nonzero scalar d\_c preserves its K-span. QED.

Choose the actual class representatives in the standard11 frames of K039(IG3). The local lattice of each new representative at11 is the standard dual-Weyl lattice. Consequently m\_(c,11)=0 and

\[
                         11\nmid d_c.                         \tag{IS6}
\]

This uses the new representative and its transported lattice. It does not assert that a rational change from an earlier display is 11-integral. K039's Cayley weak approximation applies also to this odd terminal rank. The tilQ decoration uses an11-unit, so it does not change this conclusion. These choices give no bound for d\_c at other primes and no bound for the global heights of the frames.

There is also an explicit bound in terms of the underlying local frame matrices, without forming the large coefficient matrices T. Suppose the literal membership test at w is the tensor product of det(g\_(c,tau,w))\^{}(m\_tau) rowAct(g\_(c,tau,w)), where m\_tau is the retained determinant exponent and D\_tau the polynomial degree. Put a\_(c,tau,w)=max(0,-min\_(i,j) ord\_w(g\_(c,tau,w;ij))). Then IS4 may be replaced by the sufficient, possibly larger clearing exponent

\[
 m_{c,l}'=\max_{w\mid l}
   \left\lceil{\sum_\tau\bigl(D_\tau a_{c,\tau,w}
            -m_\tau\operatorname{ord}_w\det g_{c,\tau,w}\bigr)
                        \over e_w}\right\rceil_+ .           \tag{IS6a}
\]

Indeed a local uniformizer to the power a\_(c,tau,w) times g is an integral matrix. Its polynomial row action preserves the integral dual-Weyl lattice and is homogeneous of degree D\_tau in its matrix entries. Thus the original row action has denominators at most D\_tau a\_(c,tau,w). Multiplying by the displayed determinant power gives exactly the numerator in IS6a, and tensoring adds the valuations. This proves the sufficient bound. The fact that an integral matrix, even with nonunit determinant, acts integrally follows directly from the polynomial Schur construction; equivalently Cauchy--Binet acts on each integral minor and the accepted integral straightening basis. This argument uses no inverse of that integral matrix. Formula IS6a makes the weight dependence explicit once the actual class frame valuations are known; it does not bound those frame valuations uniformly for free.

For later norm estimates define, at a fixed complex embedding,

\[
 B_i^2=\|p_i\|_{\rm nfAll}^2
    =\sum_{\alpha}|c_{i,\alpha}|^2
                   {\alpha!\over H},\qquad B_T=\max_i B_i.
 \tag{IS7}
\]

The sum is the literal normalized Fischer norm, with the tensor extremal norm H; determinant shifts are represented in the source frame but have compact norm one. This is an exact finite formula. If D\_tau is the polynomial degree at tau, the accepted basis consists of products of minors. A minor of size s has coefficient l1 norm at most s!; since s!\textless=N\^{}s, its product has coefficient l1 norm at most N\^{}(D\_tau). A degree-D\_tau monomial has alpha!\textless=D\_tau!, and H\_tau\textgreater=1. Therefore the optional coarse bound is

\[
 B_T\le\prod_\tau N^{D_\tau}
                    \sqrt{D_\tau!/H_\tau}
       \le\prod_\tau N^{D_\tau}\sqrt{D_\tau!}.
 \tag{IS8}
\]

This estimate is deliberately coarse. It is not a claim that the terminal coefficient cost is small enough for DP44.

\subsection{3. Integral class vectors and the genuine finite projector}\label{integral-class-vectors-and-the-genuine-finite-projector}

For every c,i form the actual stabilizer sum

\[
 s_{c,i}=\operatorname{single}_c
          \left(\sum_{\gamma\in\Gamma_c}
                         A_c(\gamma)d_c p_i\right).
 \tag{IS9}
\]

The single-class expression is extended by the original equivariance. It is well defined because its value is fixed by A\_c(Gamma\_c).

\textbf{Lemma IS10.} The vectors s\_(c,i) belong to Lambda\_T and span V\_T's full admitted class-family space calV\_T over K. Moreover

\[
          \|s_{c,i}\|_{\rm cls}
               \le\sqrt{|\Gamma_c|}\,d_c B_i.                 \tag{IS10}
\]

\textbf{Proof.} Each summand is in L\_c, which the stabilizer preserves; the sum is invariant by reindexing Gamma\_c.~On the fixed subspace the unnormalized sum of A\_c(gamma) is \textbar Gamma\_c\textbar{} times the identity. Since the d\_c p\_i span all V\_T and K has characteristic zero, their stabilizer sums span precisely the fixed subspace. Taking the direct sum over c proves the spanning assertion. The triangle inequality and unitarity bound the value at c by \textbar Gamma\_c\textbar d\_c B\_i. Its class norm is divided by sqrt(\textbar Gamma\_c\textbar), which gives IS10. QED.

No stabilizer denominator was inserted into these integral vectors. The stabilizer denominator in the metric remains exactly the one in IS2. This distinction is essential at a prime dividing \textbar Gamma\_c\textbar.

Let Gcal be the fixed finite group of R032, let X\_good be the actual full genuine goodTuple character list, and put

\[
 d_G=|\mathcal G|,
 \qquad e_A=\sum_{\chi\in X_{\rm good}}
       {\chi(1)\over d_G}\sum_{g\in\mathcal G}\chi(g^{-1})[g].
 \tag{IS11}
\]

For the two fixed primes this is d\_G=\textbar GL\_(a+1)(g0)\textbar\^{}f \textbar GL\_(a+1)(g1)\textbar\^{}f.~The exact finite group and character convention are those already matched to the analytic e0 in R032; no selected constituent replaces the full good projector.

Every d\_G e\_A coefficient is an algebraic integer: irreducible character values are sums of roots of unity and chi(1) is an integer. The actual finite Gcal action preserves the full coefficient lattice and class metric before the generic-character cut. Centrality makes e\_A commute with that cut. Hence on the admitted terminal generic space e\_A is orthogonal and d\_G e\_A is integral. This is exactly the source \nolinkurl{LvlOpIntegral} convention in \nolinkurl{GoodPartHermitianSpec.lean248–258,} not an assertion that every individual g preserves a nonnormal character subspace.

Define the finite family

\[
                        z_{c,i}=d_G e_A s_{c,i}.              \tag{IS12}
\]

\textbf{Theorem IS13.} These vectors lie in the unchanged genuine-good terminal lattice Lambda\_T intersect im(e\_A), span its K-vector space, and satisfy

\[
 \boxed{\quad
   \|z_{c,i}\|_{\rm cls}
       \le d_G\sqrt{|\Gamma_c|}\,d_c B_i.
 \quad}                                                       \tag{IS13}
\]

\textbf{Proof.} Integrality follows from IS10 and the preceding coefficient argument. They span the image because IS10 spans the domain and d\_G is nonzero. Orthogonality gives \textbar\textbar e\_A s\textbar\textbar\textless=\textbar\textbar s\textbar\textbar; combine with IS10. QED.

The family may include zero or dependent vectors. It need not generate the whole integral lattice over O\_K and is not asserted saturated. Only its integral membership and K-span are needed. The source \nolinkurl{cutNab} is the pullback of the unchanged lattice through \nolinkurl{cutE}; no integrality of an arbitrary image-coordinate basis or orthogonal matrix inverse is assumed.

For the actual fixed primes, the sufficient d\_G clearing factor has the following exact11 valuations:

\begin{longtable}[]{@{}lrlr@{}}
\toprule\noalign{}
Host & f & g0,g1 & ord\_11(d\_G) \\
\midrule\noalign{}
\endhead
\bottomrule\noalign{}
\endlastfoot
J14 & 42 & 4615409,5276377 & 336 \\
J56 & 168 & 22285932761,64072175441 & 21336 \\
\end{longtable}

Indeed \textbar GL\_r(q)\textbar=q\^{}(r(r−1)/2) product\_(i=1)\^{}r(q\^{}i−1), q is an11-unit, and r=a+1. Summing those elementary valuations and multiplying by f gives the table. No optimality of d\_G is claimed. Its archimedean size also has the elementary bound d\_G\textless Q\^{}(f(a+1)\^{}2), with Q=g0g1. These are fixed finite-projector costs, not new divisibility of the coupled image.

If one instead starts at principal G level and then imposes the actual subgroup character psi of P, put q\_P=\textbar P\textbar. The literal projector is e\_psi=q\_P\^{}(-1) sum\_(h in P) psi(h\^{}(-1)){[}h{]}. The family d\_G q\_P e\_A e\_psi s is integral and spans the same image. The metric on this level-character realization is q\_P\^{}(-1) times the ambient class metric, by \nolinkurl{lvlH} and \nolinkurl{LvlClassSum} in the same source file. Consequently its norm bound is

\[
      d_G\sqrt{q_P}\sqrt{|\Gamma_c|}\,d_c B_i.                \tag{IS14}
\]

The direct construction IS9--13 at the actual generic level avoids this extra clearing factor. IS14 records it when a principal-level description is used; it must not be silently set to one.

\subsection{4. Exact automorphic norm of an integral terminal vector}\label{exact-automorphic-norm-of-an-integral-terminal-vector}

Let v\_U be the actual Haar volume of U and let J be the fixed compact intertwiner used in DO038. For the canonical top map and a compact test vector v, DO038(38.11--14) gives

\[
 \|e_T(Z,v)\|_{\rm aut}^2
    ={v_U\over\dim V_T}\|Jv\|^2\|Z\|_{\rm cls}^2.           \tag{IS15}
\]

In particular each generator in IS12 has

\[
 \|e_T(z_{c,i},v)\|_{\rm aut}
 \le \sqrt{v_U/\dim V_T}\,\|Jv\|\,
                   d_G\sqrt{|\Gamma_c|}\,d_c B_i.             \tag{IS16}
\]

The same realization identifies the complexification of the full good terminal module with the full genuine-good top multiplicity space at this finite level. Evaluation on each class and finite dimensional Riesz duality prove this statement; DO038 Section 4 gives its exact inverse. It does not identify the final column image with this detector space.

For the identity terminal trace the source mass scalar is precisely c\_P=v\_U. For a nonidentity finite trace E\_T, its actual arithmetic coefficient action must be retained. A scalar trace t\_T delta\_1 multiplies this mass by t\_T. In the orbit-indexed Q construction DO053/DO054 the standard terminal trace is

\[
            t_T\delta_1,\qquad t_T=\prod_{v\in Q} I2_v,       \tag{IS17}
\]

and moved charts retain t\_T delta\_(bT). Thus its transported action is not replaced by an unmodified identity on every chart. All branch constants brConst/kapBr and the conjugate pi\_T power remain inside the fixed brC. IS15--17 introduce no redefinition of that operator.

\subsection{5. Optional integral spanning family of the actual initial Mram}\label{optional-integral-spanning-family-of-the-actual-initial-mram}

This section is useful when both inputs of the terminal pairing are to be chosen integrally. It is unnecessary if a fixed integral initial vector has already been supplied.

Use the actual initial strict-T/generic-G class space, its Weyl coefficient lattice delAtr, its actual stabilizers Gamma\_(0,c), and its unchanged maximal29 level. Choose a standard integral Weyl K-basis. It can be obtained as the pairBK-dual of the accepted standard dual-Weyl basis, using the accepted FactDelDualNabla.lean75 and K037(CR17--19). This is a Weyl basis, not a second use of the dual-Weyl lattice. Let its exact norm bound be B\_(0,i), and let d\_(0,c) be its local clearing integer defined as in IS4. Stabilizer sums give integral vectors u\_(c,i) spanning the entire initial coefficient space, with

\[
 \|u_{c,i}\|_{0,\rm cls}
     \le\sqrt{|\Gamma_{0,c}|}\,d_{0,c}B_{0,i}.                \tag{IS18}
\]

At each actual selected place v above T let Delta\_v(t) be the literal diamond action and q\_v=\textbar k\_v\^{}times\textbar=T−1. It changes the covector decoration and has the identity coefficient map: \nolinkurl{ScaledRamifiedLattice.lean79–91.} It preserves the integral lattice and the actual class metric, since it is an invertible decoration permutation commuting with rational transport. Different places commute. Thus

\[
 N_v=\sum_{t\in k_v^\times}\Delta_v(t),\qquad
 p_{\rm ram}=\prod_{v\mid T}(1-q_v^{-1}N_v),\qquad
 C_T=\prod_{v\mid T}q_v=(T-1)^f                              \tag{IS19}
\]

is the orthogonal projector onto the literal Mram=∩ker(N\_v). This is exactly ColRamV/MramV at157--173 of that source, with unweighted norms. No prescribed nontrivial central character is substituted for this intersection.

The finite family

\[
 x_{c,i}=C_Tp_{\rm ram}u_{c,i}
          =\prod_{v\mid T}(q_v I-N_v)u_{c,i}                 \tag{IS20}
\]

belongs to the actual ramified integral lattice, spans Mram over K, and obeys

\[
 \|x_{c,i}\|_{0,\rm cls}
     \le C_T\sqrt{|\Gamma_{0,c}|}\,d_{0,c}B_{0,i}.            \tag{IS21}
\]

\textbf{Proof.} The expression on the right of IS20 is integral. Each average q\_v\^{}(-1)N\_v is an orthogonal projection, because it averages a finite unitary group action. Their commuting complements project onto the intersection of their kernels. Therefore p\_ram has norm at most one and its image is Mram. Applying it to a spanning family proves the K-span and IS21. QED.

The source Sad/FFStar operator S\_0 is a different operator. It is retained in the actual coupled input whenever the original-map law requires it. IS20 is not a replacement for S\_0, and IS21 does not bound its norm or denominator.

\subsection{6. Finite arithmetic matrix and the exact nonzero premise}\label{finite-arithmetic-matrix-and-the-exact-nonzero-premise}

Let A\_cpl be the actual K038 full-Y coupled arithmetic map with its literal subD coefficient contractions. Schematically its entries are

\[
 (A_{\rm cpl}P)(\delta)=
     \sum_{e,f,Y}{k_2(f,Y,\delta)k_1(e,f,Y)P_e
                         \over|\Gamma_{0,e}|\,|\Gamma_{1,f}|}.
 \tag{IS22}
\]

The notation abbreviates the fixed contraction matrices, not their norms or a scalar surrogate. The whole actual Y sum and every raw finite weight are retained. Its codomain is evGSD.EqFun, as proved in K038(CA27--CA30); an ordinary ColV row is not substituted.

For that fixed map define the bilinear K-valued functional

\[
 \mathscr B(P,Z)=\operatorname{PtilFF}(A_{\rm cpl}S_0P,Z)
   =\sum_{c\in C}{
       \operatorname{pairBK}(\operatorname{brC}
                   (A_{\rm cpl}S_0P)(\delta_c),Z_c)
                         \over|\Gamma_c|}.                  \tag{IS23}
\]

Here delta\_c is the actual row decoration, whereas d\_c in IS4 is the denominator-clearing integer. If the application has no S\_0 on the input, replace it by the identity in IS23. If the original source requires another fixed trace or coefficient operator, it must be inserted in IS23 at its actual place, before using this theorem.

\textbf{Theorem IS24.} For the unchanged map in IS23, the following are equivalent:

\begin{enumerate}
\def\labelenumi{\arabic{enumi}.}
\tightlist
\item
  B is not identically zero on Mram times the full genuine-good admitted terminal coefficient space.
\item
  At least one of the finite matrix entries B(x\_(c,i),z\_(d,j)) is nonzero.
\item
  The actual branched image brC A\_cpl S\_0(Mram) is not contained in the bilinear annihilator of that terminal detector space.
\end{enumerate}

\textbf{Proof.} Equivalence of 1 and 3 is the definition of the annihilator with the actual stabilizer masses. By IS13 and IS20, the two displayed families span the two domains over K. If all their pairings vanish, bilinearity makes B zero on the product of the spans. The converse is immediate. Scalar extension to C preserves the zero/nonzero assertion for this K-defined matrix. QED.

For a fixed integral P for which B(P,-) is nonzero, IS13 alone gives a suitable bounded integral terminal z\_(d,j); no assumption on an integral preimage of an analytic vector is used. If only a smaller prescribed initial subimage I is allowed, replace Mram in IS24 by I and use the images of a spanning family under that prescribed prefix map. Replacing I by Mram without such a reason would invalidate the conclusion. The two actual horizontal edges in A\_cpl are never removed from IS23.

The estimates IS13, IS16 and IS21 bound the chosen inputs. They do not bound the size of a nonzero matrix entry from below and do not bound the operator norm of A\_cpl or S\_0. Their denominators, class frame heights, scalar field costs and source creation/period factors remain explicit parts of the eventual total cost.

\subsection{7. What the original-map seesaw removes, and what it does not}\label{what-the-original-map-seesaw-removes-and-what-it-does-not}

The literal source offers a vanishing implication, which is enough. For m=2 its \nolinkurl{LastX} specialization is

\[
\begin{split}
 &c_Vc_{K,0}c_{K,1}c_P\,
      \operatorname{tail}(K_1K_0S_0P,F)=0\quad\Longrightarrow\\
 &c_{I,0}c_{I,1}\,
       \operatorname{per}(P,D_0D_1F;y,v)=0
                   \quad\text{for every }y,v.
\end{split}                                                  \tag{IS25}
\]

This is ValLastX.lean237--243, with its actual original maps io and ioT, not a newly fitted top map. It is also MainX at n=0. All c\_I are required nonzero; the actual displayed c\_V and the source kernel/terminal scales are nonzero in the checked construction. The source proof first obtains identities for all real test polynomials and then passes through the common kernel/frame laws. IS25 is not an equality between one chosen compact frame value and one prescribed coefficient pairing.

Assume the following exact interfaces for the constructed mixed family:

\begin{itemize}
\tightlist
\item
  its original downward map is D=D\_0D\_1, and R032/LIFT38 supplies D onto the actual genuine-good, generic bottom multiplicity space;
\item
  its arithmetic row K\_1K\_0S\_0 is precisely the map in IS23, including the complete Y/orbit comparison, finite traces, and the fixed common tensor;
\item
  its terminal reading is DO038, so the tail for e\_T(Z) is the fixed arithmetic pairing in IS23, with the trace action and c\_P retained;
\item
  its initial reading identifies the selected bottom period with \nolinkurl{per(P,-;y,v)} on the literal Mram space.
\end{itemize}

\textbf{Corollary IS26.} Under these interfaces, if there exist P in Mram, a bottom vector g in the stated genuine-good space, and y,v such that per(P,g;y,v) is nonzero, then B is nonzero and IS24 supplies bounded integral inputs with a nonzero pairing. There is no additional independent prescribed-image premise.

\textbf{Proof.} Onto D gives a top F with DF=g. DO038 and R032 identify that full top multiplicity space with the complexification of the genuine-good admitted coefficient module. If all actual tail pairings of K\_1K\_0S\_0P with this module were zero, the first line of IS25 would hold for F. Its second line contradicts the chosen nonzero period. Thus at least one tail pairing is nonzero. This is B after the stipulated arithmetic/terminal interfaces. Apply IS24; if P was not initially K-rational, K-spanning of Mram and linearity first extract a K-defined input. No norm bound for a chosen analytic inverse F is needed for this existential argument. QED.

The arithmetic/analytic interface in this corollary is substantive. The current DO054 draft(TC21) states the original-test analytic last-step equation for its orbit-indexed family. Its initial Q trace is the identity (DO053), its terminal Q trace is IS17, and its Sad is the literal diamond/FFStar fold. Its Section 6 retains the need to match that orbit-indexed arithmetic KerX/PerX family to any earlier frozen all-Y record. Until that matching is checked, IS26 remains conditional at this named interface. It is not licensed merely by nonzero unrestricted theta images or D's onto property.

Even after that interface is proved, K058 alone does not supply IS26's Mram premise. It supplies a genuine-good selected restriction whose initial row is generic at G. The actual T projection is the commuting intersection in IS19. LIFT43's cross-centre Gram retains that unresolved distinction. Thus this finding removes the need for an integral inverse, but does not assert that the presently known generic row has the required nonzero Mram component.

\subsection{8. A small exact cost corollary of K058}\label{a-small-exact-cost-corollary-of-k058}

For completeness, the already checked K058(PG31b) unit scalar bottom vector satisfies

\[
 |F(ph)|\ge {1\over2\sqrt{2NV_\Gamma d_{\max}}},
 \quad N<2^{32af},\quad
 d_{\max}\le D< Q^{fa(a+1)/2}.                               \tag{IS27}
\]

Here N is K058's short-orbit count, not the terminal rank of Section 1. With a=f and the actual Q, Q\textless2\^{}45 for J14 and Q\textless2\^{}71 for J56. Consequently

\[
 |F(ph)|>
 \begin{cases}
   2^{-881561}/\sqrt{V_\Gamma},&J14,\\
   2^{-85116530}/\sqrt{V_\Gamma},&J56.
 \end{cases}                                                \tag{IS28}
\]

Indeed the respective exact exponents before rounding are 1763121/2 and170233059/2. This is a cost for the already constructed G-generic selected value only. It is not a lower bound after Mram, not an arithmetic nonzero-value height, and not a bound for the terminal generators' remaining d\_c costs.

\subsection{9. Source and check record}\label{source-and-check-record}

All Lean source locators refer to the stipulated pinned mirror commit 7cbf7a0c32fca1badf05931879a37c3f8a2b03b4. Accepted source theorems are used at their stipulated status. No Lean execution is claimed.

\begin{itemize}
\tightlist
\item
  \nolinkurl{Formalism/LocalToGlobal/Columns/RowForms.lean268–286:} inverse row character, MtildeV, intNab and separate intDel.
\item
  \nolinkurl{Formalism/GradedWeylDivisibility/StatementD.lean24–29,43–56} and the twist/tensor clauses; FinalD.lean: accepted actual semistandard bases and transported lattice tests.
\item
  \nolinkurl{Formalism/LocalRingCardinality/FactDelDualNabla.lean75:} accepted perfect source Weyl/dual-Weyl pairing.
\item
  \nolinkurl{Formalism/LocalToGlobal/Divisibility/GoodPartHermitianSpec.lean} 17--31,78--132,235--258: genuine/level idempotents, stable lattices, cutNab, the q\_P metric factor and the d\_G integral projection.
\item
  \nolinkurl{Formalism/LocalToGlobal/Columns/ScaledRamifiedLattice.lean79–91,} 141--148,157--173: actual diamond sum, the nontrivial CharT character valued in K, and the literal Mram lattice. Columns/BaseInst.lean 340--362 supplies the matching cyclotomic existence construction.
\item
  Development/AbelianVarietyTheory/Definitions/LatticeFrames/ \nolinkurl{ValLastX.lean19,108–145,153–168,194–208,237–251:} fixed input/terminal traces, original downward maps, ValLastX, LastX and MainX with S\_0 retained.
\item
  K037(CR17--22,37--39), K038(CA27--36), K039(IG1--3), DO038(38.9--17,21--22), R032(R32.10--21), LIFT38 and K058 supply the previously checked actual mixed-model interfaces.
\item
  DO053 and the current DO054 draft are used only with the explicit status and conditional orbit-to-arithmetic scope stated in Section 7.
\end{itemize}

The companion certificate tests exact denominator clearing, unnormalized invariant spanning, genuine/character projection, metric scaling, ramified projector spanning, bilinear nonzero extraction and the two exact IS28 exponent roundings on finite models. Such models check the algebraic normalizations; they do not enumerate the actual global class set or prove its nonzero period. The full mathematical argument above is separate from those finite checks.

\subsubsection{Author whole-proof check}\label{author-whole-proof-check-1}

The complete saved IS1--28, including IS6a, and the entire companion script were read and reconstructed. The local frame direction, ramification-index rounding, Weyl/dual-Weyl distinction, literal stabilizer masses, genuine and generic projector normalizations, canonical top norm, Mram spanning, restricted-image equivalence, and conditional LastX argument were checked. The numerical IS28 roundings and fixed d\_G eleven-valuations were replayed exactly. DO054 is retained as a draft interface, not promoted by this check. No global class enumeration, coupled image primitivity, nonzero Mram period, or integral inverse is asserted.

Revision check: the entire final report was reread after replacing the provisional O\_T field caveat by the exact CharT.Ok proof. The coefficient-field/CM-host distinction and finite-etale local11 claim were checked directly. The original-map matching remains explicitly conditional at the historical DO054/DO055 interface. The unchanged 445-assertion certificate was replayed at07:56:17.114728 UTC.

\clearpage

\section{ALP-LIFT-44 --- the exact fixed-complement pencil and a residual unitary plane}\label{alp-lift-44-the-exact-fixed-complement-pencil-and-a-residual-unitary-plane}

Author: GPT-6 Astra Pro, ALP\_MACHINE \nolinkurl{/root/lift_accounting}. 22 November2026 UTC.

Author whole-proof review PASS: 2026-11-22T07:57:42.743198+00:00. Exact diagnostics are recorded separately; no automorphic Gram is evaluated.

For the unchanged favorable LIFT41 point, this proves an actual local lattice description of every scalar-T packet correspondence. It gives integrality on the full fixed positive complement. Consequently a sufficiently small relative operator fixes that complement's short (a-1)-dimensional part, and acts only on an explicit unitary plane. Every subset of changed T places is allowed. The remaining plane arithmetic, point-versus-operator comparison and the projected global Gram are retained as separate problems.

\subsection{1. Fixed data and what a relative correspondence means}\label{fixed-data-and-what-a-relative-correspondence-means}

Use the exact actual Q=g0g1, pi, T and LIFT41 vector

\[
 w=Q(u_1,u_2,u_3,u_4,0,\ldots,0),\quad u_1=2,\quad
 W=w^*w=77Q^4-14Q^2,\quad
 v_0=(w,\pi_k),\quad z=v_0/\bar\pi.
 \tag{PR1}
\]

The Hermitian matrix is J=diag(I\_a,-1), and v0*Jv0=-T. The explicit u\_i are the fixed positive integers of LIFT41(NP4), not new choices. In both hosts each u\_i\textless9Q, and sqrt(u\_i²+4)\textless10Q. The exact certificate below verifies these integer inequalities.

Fix the actual finite decomposition

\[
                 \gamma_f^{-1}=x\tau k,\qquad
             x\in H(\mathbb A_f),\quad k\in J^1,
 \tag{PR2}
\]

from LIFT41(NP9), where gamma e=z. For a tuple s of residue units at any subset of the selected T places, let delta\_s be its projective-T representative in H, equal to the identity away that subset. Suppose h is a rational unitary matrix bringing the packet point x delta\_s tau to the same standard maximal finite component:

\[
             h_f x\delta_s\tau\in J^{\max},\qquad R=h\gamma^{-1}.
 \tag{PR3}
\]

This is the relative matrix considered here. No assertion that all packet points lie in this coarse component is needed. A cross term between different coarse components is zero in the corresponding componentwise Poincare construction. The theorem applies whenever a relative rational R as in PR3 occurs. At unchanged places R is integral. At G it need not be congruent to I; no such hypothesis is used below.

At a selected split T component, let Lambda=O\_v\^{}(a+1), and denote by u the actual primitive covector in the lower translation tau. Its conjugate component is always the inverse-adjoint one. All equalities in the next section are at the selected component, not simultaneous matrix equalities at the opposite CM component.

\subsection{2. The full selected-place lattice pencil}\label{the-full-selected-place-lattice-pencil}

Choose delta\_s=diag(s I\_a,1). This is a valid representative of the same strict/projective quotient as any h\_s with u h\_s=s u. Different representatives differ by the actual strict stabilizer, and their conjugates by tau are integral. Direct multiplication gives

\[
 \kappa_s=\tau^{-1}\delta_s\tau
  =\begin{pmatrix}sI_a&0\\(1-s)\pi^{-1}u&1\end{pmatrix},
 \qquad
 \kappa_s\Lambda=
       (1+c\pi^{-1}e(u,0))\Lambda,\quad c=s^{-1}-1.
 \tag{PR4}
\]

The right diagonal factor diag(s\^{}-1 I\_a,1) is an integral unit. Thus c ranges over the residue classes other than-1; c=0 is the unchanged lattice. Integral lifts of c give the same lattice because the rank-one operator e(u,0) squares to zero.

At this component PR2 gives k\^{}-1e=z, since x and the selected lower tau both fix e. Define the integral row

\[
 m=e^t k,\qquad q=(u,0)k.
 \tag{PR5}
\]

Both z and q are primitive, qz=0 and mz=1. The last row of gamma\^{}-1 is -z*J; its selected component also equals(pi\^{}-1u,1)k. Therefore

\[
             q=-v_0^*J-\pi m.
 \tag{PR6}
\]

Here z\emph{=v0}/pi. Formula PR6 is exact, not merely a residue congruence. Using PR3, R belongs to K kappa\_s\^{}-1 k. Equivalently its inverse-image lattice is exactly

\[
 \boxed{\quad R^{-1}\Lambda
       =k^{-1}\kappa_s\Lambda
       =(1+c\pi^{-1}zq)\Lambda.\quad}
 \tag{PR7}
\]

No unspecified left-compact factor remains in this lattice identity.

For clarity, put Lambda\_-=\{y in Lambda:qy in pi O\_v\} and Lambda\_+=Lambda+pi\^{}-1 O\_v z. Primitivity and qz=0 imply that Lambda\_+/Lambda\_- is a two-dimensional residue vector space. The lattices in PR7 are the q\_v-1 lines with parameter c in k\_v except-1; there are two omitted projective lines, infinity and c=-1. Indeed, choose y with qy=1. Modulo Lambda\_-, the lattice in PR7 is spanned by y+c pi\^{}-1z. The line Lambda itself is c=0. This description keeps the actual c normalization, so that the excluded value is not lost under a projective reparameterization.

The conjugate-place lattice is the Hermitian dual of PR7 under the split pairing. This follows directly from R*JR=J: the inverse image of a selfdual pair of standard lattices is again a selfdual pair. Thus PR7 and its dual give the complete local inverse-image lattice condition, including the conjugate place. They do not determine a left compact representative or a Weil-model phase. Conversely, if a rational unitary R has these selected-place lattices and their duals, and is integral invertible elsewhere, then h=R gamma satisfies PR3: at T the matrix R k\^{}-1 kappa\_s preserves Lambda, and the other places are immediate. If R is principal at G, the resulting G frame lies in J1 because k does. Thus the criterion also constructs the stated packet member, without a new rational approximation argument.

\subsection{3. Integral transport of the entire fixed positive complement}\label{integral-transport-of-the-entire-fixed-positive-complement}

Define the actual global integral complement

\[
                  L=\{y\in O_E^{a+1}:v_0^*Jy=0\}.
 \tag{PR8}
\]

It is positive definite at every real label. We claim

\[
                              R L\subset O_E^{a+1}.
 \tag{PR9}
\]

At a selected T place, PR6 gives qy=-pi my for y in L, so y belongs to Lambda\_-. The latter is contained in every lattice PR7: if qy=pi b, then (1-c pi\^{}-1 zq)y=y-cbz is integral. Since (zq)²=0, this explicitly gives the inverse of the lattice generator. Hence R y is integral there.

At the conjugate place use the exact duality just proved. Every selected-place vector in PR7 belongs to Lambda+pi\^{}-1 O\_v z. An integral y orthogonal to v0 pairs integrally with Lambda and pairs as zero with z, because y*Jz=0. Consequently its conjugate component belongs to the dual lattice, so R y is integral there as well. This proves both halves without dropping the inverse-adjoint action. At every other finite place R is integral by PR3. Integrality at all finite places proves PR9. The proof is unchanged if only some T places are varied: at the remaining places the lattice is Lambda.

\subsection{4. An exact rank-two reduction for small operators}\label{an-exact-rank-two-reduction-for-small-operators}

Let

\[
                   \epsilon=\max_{\tau\in\Phi}
                                  \|R_\tau-I\|_{\rm op}
                         <{1\over10Q},
 \tag{PR10}
\]

where this is the Euclidean operator norm on the standard coordinates. This hypothesis is stronger than a small displacement in the symmetric space; no converse is assumed.

The following a-1 independent vectors lie in L:

\[
 e_5,\ldots,e_a,\qquad
       l_i=u_i e_1-2e_i\quad(i=2,3,4).
 \tag{PR11}
\]

Their negative coordinates are zero. The equations w*l\_i=0 use u1=2, and their norms are1 or sqrt(u\_i²+4)\textless10Q. By PR9, R l\_i-l\_i and R e\_j-e\_j are integral vectors. At every complex embedding each coordinate has absolute value strictly below1, by PR10. A nonzero algebraic integer cannot have all conjugates below1, since its absolute norm is a positive integer. Thus all these differences vanish. They span the positive hyperplane w-perp. Hence

\[
 R|_{w^\perp}=1,\qquad
 R\bigl(\operatorname{span}_E(w,e)\bigr)
                  =\operatorname{span}_E(w,e),\qquad
                         \operatorname{rank}_E(R-I)\le2.
 \tag{PR12}
\]

The invariance of the orthogonal plane follows from the unitary identity and the pointwise fixation of its nondegenerate orthogonal complement. There is no Jordan-form or unipotence assumption. The exact displayed threshold is larger than2\^{}-49 for J14 and larger than2\^{}-75 for J56, by Q\textless2\^{}45 or2\^{}71. In particular it covers the proposed2\^{}-230 small-operator range with a large margin.

\subsection{5. Determinant and the exact residual form}\label{determinant-and-the-exact-residual-form}

The determinant of R is a finite unit: the local lattices in PR7 have relative determinant valuation zero, including their conjugates, and R is integral invertible elsewhere. It has norm one at every real embedding. Therefore it is a root of unity by the elementary Kronecker argument already used in K050, and the actual hosts have root group mu28 (K031/K032).

On the orthonormal plane spanned by w/sqrt(W),e, the residual matrix U satisfies \textbar\textbar U-I\textbar\textbar\textless=epsilon. Since det R=det U, \textbar det R-1\textbar\textless=2epsilon+epsilon². For PR10 this is less than1/7. Every nonidentity28th root has distance at least2 sin(pi/28)\textgreater=1/7 from1, using sin x\textgreater=2x/pi for0\textless=x\textless=pi/2. Consequently det R=1.

In the E-basis(w,e), the plane has Gram diag(W,-1). Inverting a 2 by2 determinant-one matrix and comparing it with its Hermitian inverse gives the exact shape

\[
 U=\begin{pmatrix}\alpha&\beta\\W\bar\beta&\bar\alpha\end{pmatrix},
 \qquad \alpha\bar\alpha-W\beta\bar\beta=1,
 \qquad \alpha,\beta\in E.
 \tag{PR13}
\]

Its ambient action is

\[
 R=\begin{pmatrix}
  I_a+{\alpha-1\over W}ww^*&\beta w\\
  \bar\beta w^*&\bar\alpha
 \end{pmatrix}.
 \tag{PR14}
\]

These coefficients must still satisfy the actual local pencil PR7 at every selected T place, its conjugate dual, and integrality at all other places. PR13 alone is not a sufficient packet condition.

\subsection{6. A reusable exact graph test in the residual plane}\label{a-reusable-exact-graph-test-in-the-residual-plane}

This section is joint with root's residual-plane calculation. Put U=W/Q²=77Q²-14 and use the primitive positive vector u=w/Q. Its first coordinate is2 and its second is odd, so its integral coordinate ideal is the unit ideal. In the E-basis(u,e), define

\[
 S=\begin{pmatrix}\bar\pi_k&Q\\QU&\pi_k\end{pmatrix},\qquad
 B=\begin{pmatrix}x&x'\\U\bar x'&\bar x\end{pmatrix},\qquad
 x\bar x-Ux'\bar x'=T.
 \tag{PR15}
\]

Here x,x' are integral and assume U divides x\emph{pi\_k-T in O\_E. Both S and B are Hermitian similitudes of multiplier T. Extend B S\^{}-1 as the identity on u-perp. Its plane coefficients are alpha=(x}pi\_k-x'QU)/T and beta=(-xQ+x'bar(pi\_k))/T. The ambient top block is I+(alpha-1)uu*/U. The displayed divisibility makes this block integral outside T; the other entries already have denominator at most T. The ambient matrix and its inverse are therefore integral outside T, and determinant1 and unitarity hold globally. This ambient divisibility is essential; plane-integral B alone would not suffice.

Suppose also that x and x' are units at a selected T place. Then its exact parameter in PR7 is

\[
                 s={x'\bar\pi_k\over Qx}\pmod\pi,
             \qquad c=s^{-1}-1.
 \tag{PR16}
\]

In particular it is an allowed scalar because s is a residue unit. To prove this with the full original row q, put t=-v0\emph{J and t'=(-(bar x')u},x). Since Rv0=(x'u,bar x), unitarity gives the exact identity t R\^{}-1=t'. Reduction of S adj(B)/T, including the identity extension on u-perp, gives pi R\^{}-1=z lambda t' modulo pi, with lambda=1-x'bar(pi\_k)/(Qx). Consequently q R\^{}-1=t'-pi m R\^{}-1 is integral and reduces to (1-lambda)t'=s t'. We used the exact mz=1; discarding the -pi*m term would give the wrong parameter. The pole coefficient in the Lambda\_+/Lambda\_- pencil is lambda t', while its ordinary q coordinate is s t'. Their ratio is lambda/s, which is s\^{}-1-1. The adjugate of S also shows that the pole row of R annihilates Lambda\_-, so that Lambda\_- is contained in R\^{}-1 Lambda. Since the residue row t' is primitive and both lattices have relative determinant valuation zero, this proves equality with the line in PR7. If lambda=0 the same argument gives the unchanged lattice. The conjugate component is the dual already proved in section2.

A useful sufficient principal-G condition is x congruent bar(pi\_k) modulo Q² and x' divisible by Q. It makes alpha-1 divisible by Q², beta divisible by Q, and U is a unit at G, so the ambient R is congruent to I at both G primes.

\subsection{7. A small actual semisimple member: separation is false}\label{a-small-actual-semisimple-member-separation-is-false}

Root supplied the following matrix and the author independently verified the full graph test, including the -pi*m term in PR6. In PR15 take x=bar(pi\_k), x'=-Q. Then PR16 gives s=-1 at every selected T place, and the resulting ambient matrix is

\[
 R_-=
 \begin{pmatrix}
 I_a+2ww^*/T&-2\bar\pi_k w/T\\
 -2\pi_k w^*/T&1+2W/T
 \end{pmatrix}.
 \tag{PR17}
\]

All ambient integrality and principal-G conditions of section6 hold. Thus R\_- is an actual allowed relative scalar-packet matrix, not only an element of the ambient Hecke double coset. It has determinant1, is semisimple and is not the identity. With x0=W/T, its two residual real eigenvalues are (sqrt(1+x0) plus or minus sqrt(x0))². Therefore

\[
 \|R_--I\|_{\rm op}
   =2x_0+2\sqrt{x_0(1+x_0)}<4\sqrt{x_0}
   <\begin{cases}2^{-404}&J14,\\2^{-350}&J56.\end{cases}
 \tag{PR18}
\]

The middle bound holds for x0\textless1/16, which follows from LIFT41. The last inequalities follow by squaring and using its certified exact W/T values. Hence no lower operator-separation statement on the whole actual scalar orbit at the scale2\^{}-242 can follow from PR12 alone. The primitive rank-one bound of LIFT43 is unaffected: R\_- has rank2 and is semisimple.

The precise point and decoration distinction is

\[
 R_-v_0=(-w,\pi_k),\qquad R_-p=S_0p,
                         \quad S_0=\operatorname{diag}(-I_a,1).
 \tag{PR19}
\]

The two representatives in the symmetric space are distinct, while S0 is an integral rational unitary matrix. Thus they are equivalent in the coarse maximal-level quotient. S0 is not in the principal-G subgroup, and its finite G action and selected real frame action must remain in the genuine-e0 character computation. Nor does a diagonal s=-1 comparison decide nontriviality at every separate T place. This actual example is a limit on a proposed separation argument, not a claim that a projected Poincare or Mram detector vanishes.

\subsection{8. One exact anisotropy check on J14}\label{one-exact-anisotropy-check-on-j14}

On J14 the plane diag(W,-1) is anisotropic over F. Indeed the exact actual Q has v3(11Q²-2)=1 and3 does not divide7Q; hence v3(W)=1. The actual field is Q(zeta28) times the real degree7 subfield of Q(zeta29). At3, the zeta28 Frobenius has order6 and its cube is complex conjugation, since3³=-1 modulo28. The real degree7 factor has Frobenius order1 or7. In either case the midpoint of the full Frobenius cyclic group is complex conjugation: its exponent is3 or21. Thus E/F is an unramified inert quadratic extension at3. All these fields are absolutely unramified at3.

Norms in an unramified quadratic local extension have even valuation, whereas W has valuation1. So W is not a local norm and the plane has no nonzero isotropic vector. In particular its special unitary group has no nonidentity rational unipotent: a nontrivial unipotent unitary 2 by2 matrix has an isotropic fixed line, as follows from its rank-one nilpotent part and the unitary equation, or from pairing a fixed vector with a generalized fixed vector.

This statement is only for J14. The real sqrt5 factor in J56 changes the Frobenius calculation at3, so the same obstruction cannot be transferred. On J14 every nonidentity residual PR13 is therefore semisimple; the remaining near-identity semisimple problem is not resolved here.

\subsection{9. Exact retained scope}\label{exact-retained-scope}

PR7--9 are valid for every subset of changed T places in the fixed projective-T orbit of the LIFT41 class. PR12--14 require small operator norm at every real label. A small symmetric-space displacement permits an additional compact rotation and does not supply PR10 by itself. A transformation fixing the symmetric-space point may still act on its coefficient fibre or decorations; it cannot be deleted without its actual stabilizer character.

The result reduces a specified small-operator part of the extra Gram in LIFT43 to the exact plane PR13 plus the pencil PR7. PR17 proves that actual small semisimple members remain. The result does not bound all short correspondences, prove separation of the full scalar orbit, or prove a nonzero Mram detector. No integral comparison of LIFT41's positive complement with the standard K049 cusp complement is assumed. No finite level, tuple, period or theta input is changed.

\subsection{10. Verification and sources}\label{verification-and-sources}

The actual frame and packet inputs are LIFT41(NP1--11), LIFT43(TG2--3, TG10--12), R025/R027 and the original strict/projective T definitions. The field and root-unit data are K031/K032. All new lattice-pencil, complement, rank and determinant arguments are given above. The certificate checks the actual small-complement norms, exact finite pencil counts, direct rational inverse-image conditions and both CM components, residual determinant/unitary identities, and the J14 valuation/Frobenius calculation. It does not classify the residual semisimple solutions or evaluate an automorphic Gram or period.

\clearpage

\section{AM-DO-055 --- the unchanged coupled arithmetic image in the actual last-two seesaw}\label{am-do-055-the-unchanged-coupled-arithmetic-image-in-the-actual-last-two-seesaw}

Author: GPT-6 Astra Pro, ALP\_MACHINE /root/definite\_oscillator. 22 November 2026 UTC; inherited reasoning effort is not separately exposed.

Status: author whole-proof check PASS, 22 November 2026 UTC. The complete proof and companion certificate have been read, together with the literal source and all normalization inputs cited below. Independent review is pending.

This result identifies the orbit-indexed Q calculation with the fixed full-Y arithmetic operator of K038. It retains the actual class masses, the intermediate level volume, the terminal trace, the original T period law, the real coevaluation, and the literal terminal branch. It then proves the actual LastX vanishing implication for this constructed mixed family. Consequently a nonzero \textbf{actual Mram bottom period} and the checked vertical surjectivity imply nonvanishing on the prescribed coupled arithmetic image. No additional assumption that an unrelated row already lies in that image is required.

The assertion is conditional on that Mram bottom period. It neither proves the missing Mram projection of the K058 row nor supplies an integral vertical right inverse. The old HostRaw is not instantiated.

\subsection{1. Fixed maps, levels and conventions}\label{fixed-maps-levels-and-conventions}

Retain either K031/K032 host, the K044 weights, the actual fixed finite data and the ranks

\[
 a\in\{42,168\},\quad r=a+1,\quad b=1802=r+t,
 \quad c=2b-a,\quad N=c-1=r+2t.
 \tag{CI1}
\]

Let \(\mathcal M\) be the exact initial arithmetic row module, with its strict G-unipotent and T-covector levels, and let \(\mathcal M_{\rm ram}\) be its literal intersection of the kernels of the unweighted T diamond sums. These are K038(CA27a) and the actual Record transport. Let \(\mathcal Z^{\rm good}\) be the canonical terminal coefficient module cut out by the genuine arithmetic idempotent of R032, including its prescribed generic subgroup character. The coefficient lattices are the already chosen Weyl/dual-Weyl lattices; they are not defined from an analytic theta image.

Write

\[
 \mathcal A= A_{\rm cpl}:\mathcal M\longrightarrow
                     \operatorname{evGSD},\qquad
 B(P,Z)=\operatorname{PtilFF}(\mathcal A P,Z).
 \tag{CI2}
\]

Here \(\mathcal A\) is exactly K038(CA27): the full prescribed Y sum, with coefficients \(|\Gamma_0|^{-1}|\Gamma_1|^{-1}\), the unchanged geometric jet kernels, and the fixed frame transports. The arithmetic pairing is exactly CA30/DO038, including brC and the terminal class weights. Nothing in this report defines that pairing from a period.

Let \(S\) denote the actual DO054(TC11) Sad on initial rows. Its analytic and arithmetic readings coincide because it is the same finite T-supported diamond/Fourier sum on those readings. It preserves \(\mathcal M_{\rm ram}\). Its coefficients lie in the already fixed character field: a nontrivial source additive character of the residue field of characteristic T takes a value of exact order T, so the literal CharT data already contain \(\zeta_T\). Thus this finite operator introduces no new coefficient-field extension. In the nonvanishing argument only preservation by S is needed, not an integral inverse of S.

Let \(E_T(Z,v)\) be the canonical top realization of DO038, and let \(D=D_0D_1\) be the original two-step downward multiplicity map of LIFT31/38, with the chosen real intertwiners and their nonzero IsIo constants retained. R032 proves that D is onto the same genuine good bottom multiplicity space and that \(E_T\) identifies the complexified good terminal coefficient module with its top multiplicity space. DO052's normalized generic tests give exactly these original maps on the prescribed generic character subspace.

The actual initial period will always include its moved trace:

\[
 \mathscr P_{\rm act}(P,Q;y,w)
  =\operatorname{perF}_0\bigl(\operatorname{rd}_0(P,y),Q(w)\bigr).
 \tag{CI3}
\]

In DO053 this trace is \(\delta_{b_0}\), even though the \textbf{Q component} of the trace is the identity. Thus CI3 evaluates the actual translated bottom function. It is not an assertion about an arbitrary unshifted restriction. At G the explicit central-period unit of DO052 remains if one rewrites CI3 in standard frames. The T translation/primitive negative line remains the one of LIFT30/DO051.

\subsection{2. Full-Y compression as an equality of entire top functions}\label{full-y-compression-as-an-equality-of-entire-top-functions}

At each Q place use DO053's orbit

\[
 \mathcal O_v=\{(z,w)\bmod q_v^{M_v}:
       z\in\mathcal O_{F_v}^r\setminus q_v\mathcal O_{F_v}^r,
                  \ w\in q_v\mathcal O_{F_v}^t\},\qquad
 h_v=(q_v^r-1)q_v^{b(M_v-1)}.
 \tag{CI4}
\]

The middle compact is \(J_v=\operatorname{red}^{-1}P_{r,t}\). Let \(y_v=e_r\) and let \(J_{v,y}\) be its stabilizer in this finite orbit. Then \([J_v:J_{v,y}]=h_v\). Put \(h_Q=\prod_{v\in Q}h_v\) and \(I_j=\prod_{v\in Q}I_{j,v}\), with the two parabolic indices of DO053(QT20).

For the actual first and second functions, that result gives

\[
 f_{\kappa y}\otimes g_{\kappa y}
   =(\omega_0(1,\kappa)\otimes\omega_1(\kappa,1))
                                        (f_y\otimes g_y).
 \tag{CI5}
\]

The two honest determinant characters in CI5 are reciprocal and have been included. They are not individually trivial. The Fourier phase uses the same y, the same middle frame and the actual terminal coordinate \(t/\pi_T\). Equation CI5 is an equality of entire Schwartz tensors.

Let \(\mathcal H_P\) be the two actual horizontal theta integrations, with any fixed initial input and with the top variable g still free. For each such g, right translation of the middle integration variable by \(\kappa\) removes the simultaneous action in CI5. All horizontal groups are definite; their adelic quotients are compact and the smooth theta integrals converge absolutely. The substitution is therefore valid pointwise for every top g, and uniformly for the finite sums in question. Consequently

\[
 \mathcal H_P\left(\sum_{y\in Y_Q}f_y\otimes g_y\right)(g)
   =h_Q\,\mathcal H_P(f_{e_r}\otimes g_{e_r})(g)
                                      \qquad\text{for every }g.
 \tag{CI6}
\]

All other fixed finite factors, the real polynomial, and the initial input are arbitrary here. Operators at other places commute with the middle Q action. Terms outside the displayed primitive orbit vanish by the exact first raw support. This proves CI6 before restriction to the terminal subgroup or multiplication by a terminal form. It does not assert equality of the unintegrated full sum and a single tensor.

\subsection{3. The same equality with every actual class mass}\label{the-same-equality-with-every-actual-class-mass}

There is an elementary exact class calculation which makes the coefficient content of CI6 explicit. Fix a middle decorated class c at the original level J and let \(\Gamma_c\) be its actual finite stabilizer. Under the class frame, \(\Gamma_c\) acts on the finite orbit \(O=\prod_v\mathcal O_v\). The honest characters in the two kernels cancel as in CI5, and their coefficient actions are contracted by the same pairBK. The resulting paired value \(a_y\) is therefore constant on each \(\Gamma_c\)-orbit, with the fixed rational frame transport included. Orbit-stabilizer gives

\[
 {1\over|\Gamma_c|}\sum_{y\in O}a_y
   =\sum_{[y]\in\Gamma_c\backslash O}
                              {a_y\over|\Gamma_{c,y}|}.
 \tag{CI7}
\]

This is valid with nontrivial and unequal stabilizers. The right side is precisely the canonical decorated class calculation at the auxiliary middle level \(J_y\). Indeed the refined double-coset classes over c are \(\Gamma_c\backslash J/J_y\), and their stabilizers are \(\Gamma_{c,y}\). A change of representative conjugates these groups, preserves their orders, and changes both coefficient factors by inverse frame/character actions. Those are exactly K038(CA24--25), not an assumed freeness of the class action.

If \(v_1=\operatorname{vol}(J)\), then

\[
                \operatorname{vol}(J_y)=v_1/h_Q.
 \tag{CI8}
\]

Multiplying CI7 by \(v_1\) shows that the full-Y coarse class integral equals \(h_Q\) times the one-representative refined class integral. This is the class-by-class form of CI6. In particular, the initial factor \(|\Gamma_0|^{-1}\), every middle factor \(|\Gamma_1|^{-1}\), and every finite frame phase in CA27 are retained. No class number or common stabilizer order is used.

This auxiliary refinement is a calculation, not a replacement of the prescribed final operator or a newly imposed final level. Equivalently one may use the associated finite bundle over the original middle classes with fibre indexed by O: equivariant families on O are identified with their value at \(e_r\), carrying the stabilizer condition for \(J_y\). Equation CI7 proves the identity of its contracted coefficient sum with the original full-Y sum. Neither individual first-edge output nor an arbitrary redundant array is identified by a bijection with an old scalar middle row.

For emphasis, in the raw arithmetic normalization (class weights but no level volumes), the refined one-representative composition equals the original full-Y coefficient composition by CI7. Its analytic comparison has the volume \(v_1/h_Q\), instead of \(v_1\). The QAt operator on the second edge has the literal factor \(I_1h_Q\). Together they give

\[
 (I_1h_Q)\,{v_1\over h_Q}=I_1v_1.
 \tag{CI9}
\]

Thus the orbit-indexed \textbf{analytic} two-edge operator is exactly \(I_1\) times the unchanged full-Y operator. One must not remove \(h_Q\) both from the raw coefficient map and from the middle volume.

The terminal Q trace is \(I_2\delta_{b_T}\), while \(c_Q=(I_1I_2)^{-1}\). Hence the complete scalar in the last-two period satisfies

\[
                   c_Q I_1 I_2=1.
 \tag{CI10}
\]

In the current frames the terminal move is identity (or a terminal maximal-T element acting trivially); if one changes those frames, its exact coefficient translation must be applied to Z on both sides. The initial move \(b_0\) in CI3 has not been canceled by CI10.

\subsection{4. Faithful coefficient reading and the real selected component}\label{faithful-coefficient-reading-and-the-real-selected-component}

K037's subD substitution and K038's coupled assembly fix the class kernel independently of the theta integral. Their comparison holds for every coefficient row, every actual finite class and every real frame, with complete coevaluation. Evaluations at those real frames span the dual of each finite-dimensional coefficient space. Therefore their complete reading is faithful: if its value is zero at all such frames, the coefficient polynomial is zero. A finite class imposes only its actual stabilizer relation, which does not change this fact.

It follows that CI7 is also an equality in the actual coupled coefficient module, with the volume distinction in CI8--9. In particular it supplies the arithmetic comparison used by the orbit-indexed last-two law. There is no unspecified new K map: its full-Y version is exactly \(\mathcal A\) in CI2. For this two-step claim one can prove the resulting TwoX comparison directly; there is no need to assert that two incorrectly identified individual edge maps satisfy KerX.

Here is the real coefficient reduction in a form that also records all relevant normalizations. Let \(V_0,V_1,V_2\) be the three horizontal compact coefficient representations, in the actual source dual conventions. Let \(W\) be the terminal compact representation, and put \(d_0=\dim V_0\), \(d_W=\dim W\). Use the primitive Cauchy embeddings \(Q_0,Q_1\) fixed by the subD model, not independently rescaled isometric embeddings. For a real polynomial tensor C, project to the initial selected type and then average its diagonal middle-compact action with probability Haar. DO050 Section3 proves that these operations select exactly the two prescribed Cauchy summands; no second intermediate partition survives.

In orthonormal coefficient/dual bases the resulting tensor is uniquely of the form

\[
 \Pi C=\sum_{i,k}c_{ik}
      \sum_j Q_0(e_i\otimes f_j)
                       \otimes Q_1(f^j\otimes w_k).
 \tag{CI11}
\]

The conjugates/duals in CI11 are the source pairBK duals, as in DO050(RC6). They are suppressed only in the basis letters; the source and target representations have not been identified. In a nonorthonormal algebraic basis the inverse Gram coefficients in the dual bases are retained. The coefficients \(c_{ik}\) are obtained from the actual finite polynomial C by these fixed projections. They depend on C, not on P, Z, a class representative or the tuple.

For the initial reading at frame y define the final evaluation functional

\[
 E_{C,y}={1\over d_0}\sum_{i,k}c_{ik}E_y(e_i)E_{w_k}.
 \tag{CI12}
\]

The factor \(1/d_0\) is the first compact Schur integration. The middle coevaluation in CI11 already cancels its middle Schur denominator. To check this normalization, a single matched middle tensor averages to \(1/\dim V_1\) times that coevaluation. For unit highest seeds CI12 therefore gives the two factors \(1/(\dim V_0\dim V_1)\), exactly K038(CA36). This also proves CI12 directly by applying the compact rank-one averaging identity DO038(19h) successively. It does not assign the coevaluation scalar to each simple tensor before averaging.

Write \(B_T=\operatorname{brC}\), \(J:W\to V_\sigma\) for the canonical top coefficient identification, and \(B_TB_T^*=\beta_B I\). The exact terminal factor is

\[
       a_C(y,v)={E_{C,y}(B_T^*Jv)\over \beta_Bd_W}.
 \tag{CI13}
\]

This is DO038(19f) with the functional CI12; it includes the complete branch and its actual normalized Fischer metrics. It may vanish at particular y,v,C. Its dependence on C is a fixed finite linear calculation and does not depend on an arithmetic row or an analytic period value.

For comparison with the primitive real formulas, write \(b_{\rm Br}=\prod_\tau u_{T,\tau}c_{{\rm Br},\tau}\sqrt{\kappa_\tau}\). Then \(B_TJ_\mu=b_{\rm Br}I\) and \(\beta_B=|b_{\rm Br}|^2\). Let \(H_j=\prod_\tau H_{\nu_{j,\tau}}\) be the two horizontal primitive Fischer norms, and let z, \(d_v\), \(N_0,N_1\) and \(D_{\rm PBW}\) have exactly the DO050(RC7,RC12) meanings, with products over all real labels. With matched unit source and terminal seeds, normalizing \(E_y\) of the source seed to one, a unitary complete vertical coevaluation has

\[
 a_C={z\over d_0\sqrt{H_0H_1}}\,
                          {1\over b_{\rm Br}d_W}.
 \tag{CI14}
\]

For a matched simple vertical pair replace z by \(z/d_v\); for PN15 primitive embeddings and the unnormalized selected PBW vector multiply by \(\sqrt{N_0N_1D_{\rm PBW}}\), exactly DO050(RC13). A primitive integral higher Pieri input carries its additional LIFT40 \(q_\lambda\). The branch factor is inverted in CI14 because this equation compares a restricted \textbf{period} with PtilFF, which already contains the branch. Applying the branch to a vector instead gives the forward factor in DO050(RC15). Those are different operations.

\subsection{5. Exact last-two scalar identity on the prescribed image}\label{exact-last-two-scalar-identity-on-the-prescribed-image}

Let \(v_0,v_1,v_T\) be the original strict finite volumes, and let the K038 raw source-creation comparison be

\[
 \mathfrak p_{01}=v_0v_1
   C_\theta^{14h(ab+bc)}
      \prod_\tau(\lambda_\tau/cD_\tau)^{D_{0,\tau}+D_{1,\tau}}.
 \tag{CI15}
\]

This is CA33. The literal source unit-vacuum comparison is

\[
 \mathfrak p_{01}^{\rm unit}
       =a_{\rm vac}\mathfrak p_{01},\qquad
 a_{\rm vac}=\det(2Y_J)^{7hb^2}>0.
 \tag{CI16}
\]

LIFT39 supplies CI16; using the raw product Gaussian without this factor would mix the two vacuum normalizations. The real tensor coefficients in CI11 use the same unit source vacuum. The negative column conjugated CM period and the actual line-shuffle phase are therefore those of DO050/LIFT39, not silently replaced by their absolute values.

The common finite tensor is a scalar multiple of the horizontal product basic tensor. At T its scalar is the computed \(\Gamma_TT^{fb/2}\). At all other places both tensors belong to the same one-dimensional lattice-vacuum space. Retain their fixed unit phase \(\Gamma_{\rm bas}\), defined by their Hilbert inner product in the already chosen local Over models. It is one at places where those models identify the two basic indicators with coefficient one, and is one at almost all places. Thus

\[
                \Omega_f=\Gamma_{\rm bas}\Gamma_TT^{fb/2}\ne0.
 \tag{CI17}
\]

This is a definition from the two independently fixed basic vectors, not a scalar fitted to a desired theta period. The vacuum-line argument proves that it is independent of P, Z, the tuple and all real polynomials. Its phase is not asserted to be one or assigned an unproved arithmetic valuation. One may instead retain the entire basic model comparison; CI17 records exactly that remaining fixed phase in scalar notation.

Let \(s_Gs_Q\) be DO054(TC19), and let \(c_T\) be TC13. For the real inputs P0,P1 of TC21, use their common real tensor in CI11--13. Then the preceding calculations prove, for every actual \(P\in\mathcal M_{\rm ram}\), every good terminal row Z, every initial frame y and terminal seed v,

\[
 \boxed{\begin{split}
 &\operatorname{perF}_0\left(\operatorname{rd}_0(P,y),
       \operatorname{lift}_0(P0,
             \operatorname{lift}_1(P1,E_T(Z,v)))\right)\\
 &\qquad=s_Gs_Qc_T\,\Omega_f\,
       \mathfrak p_{01}^{\rm unit}v_T\,
       a_C(y,v)\,
       \operatorname{PtilFF}(\mathcal A(SP),Z).
 \end{split}} \tag{CI18}
\]

The actual moves in CI3 and the same passive Q comparison are understood on both sides. There is no extra initial Q projector. If a nonidentity terminal move is retained, its explicitly transported row replaces Z, as in DO038 Section6; in the current construction it acts trivially as noted after CI10. The genuine finite good projector is the R032 operator on Z and on D, not the initial Q identity.

\textbf{Proof.} Apply DO054(TC21). Its scalar is \(c_V=s_Gs_Qc_T/(I_1I_2)\). CI6--9 turn its orbit-indexed two-edge function into \(I_1\) times the unchanged full-Y function. Its terminal trace contributes \(I_2\). The Q scalars cancel by CI10. K038's canonical full-type comparison, with CI16--17, identifies the remaining function on the actual terminal class reading with \(\Omega_f\mathfrak p_{01}^{\rm unit}\) times the reading of \(\mathcal A(SP)\) at the evaluation functional CI12. This equality holds on every actual class by CI7 and K038's two covariance laws. Finally DO038(19f) gives exactly \(v_Ta_C(y,v)\) times its literal PtilFF. This proves CI18. Every scalar comes from a fixed kernel, level volume, line map or coefficient metric. QED.

In IsIo notation the left side is also

\[
 c_{I,0}c_{I,1}\,
   \mathscr P_{\rm act}\bigl(P,DE_T(Z);
                   y,\operatorname{It}_0(P0,
                                    \operatorname{It}_1(P1,v))\bigr).
 \tag{CI19}
\]

The nonzero \(c_{I,j}\) depend on the constructed real normalization; they are not dropped. Formula CI19 is exactly the source IsIo/IsIoT equation applied twice, as in ValLastX.lean148--150. LIFT39 gives the primitive normalization if that choice is desired.

Only the fixed real factor \(a_C\) can vanish in CI18. In particular the complete T scalar is

\[
 c_T\Omega_f=\Gamma_{\rm bas}\Gamma_T(T^b-1)^f
               \prod_{v\mid T}(\mu_0\xi)(\varpi_v)\ne0.
 \tag{CI20}
\]

All Q indices have canceled in this complete equality, but the Q denominators and weights remain in its separate arithmetic and analytic constituents. CI20 is not a claim that their integral costs or the complete DP44 costs disappear.

\subsection{6. The actual-image consequence and its precise hypothesis}\label{the-actual-image-consequence-and-its-precise-hypothesis}

Suppose

\[
 \operatorname{PtilFF}(\mathcal A P,Z)=0
       \quad\text{for every }P\in\mathcal M_{\rm ram},
                           \ Z\in\mathcal Z^{\rm good}.
 \tag{CI21}
\]

Since S preserves Mram, CI18 makes the left side of that equation zero for every P0,P1,y,v and every such P,Z. By CI19, all the corresponding bottom periods of \(DE_T(Z)\) vanish. These real vectors span the needed selected bottom compact type: PN15 gives \(\operatorname{It}_1(P_1,v_2)=v_1\), and PN20 applies the fixed Pieri/PBW operators to the first joint vector, giving every vector in that selected compact type after division by its nonzero PBW scalar. More generally the source ItSpans clauses give the full span. Thus

\[
 \mathscr P_{\rm act}(P,DE_T(Z);y,w)=0
 \quad\text{for every selected }w,y,P,Z.
 \tag{CI22}
\]

R032/LIFT31 supplies \(D(E_T(\mathcal Z^{\rm good}_\mathbb C))
=\mathcal C_0^{\rm good}\), with the actual generic character cut and the unchanged other finite levels. Therefore CI22 holds for every actual good bottom multiplicity vector. Taking the contrapositive proves the following theorem.

\subsubsection{Actual-image theorem}\label{actual-image-theorem}

If there exist an actual initial Mram row P, a genuine-good bottom vector Q and selected real/frame vectors y,w such that \(\mathscr P_{\rm act}(P,Q;y,w)\ne0\), then there exist \(P'\in\mathcal M_{\rm ram}\) and \(Z\in\mathcal Z^{\rm good}\) with

\[
                  \operatorname{PtilFF}(\mathcal A P',Z)\ne0.
 \tag{CI23}
\]

One may take \(P'=SP\) for a nonzero witness arising directly from CI18 after choosing a terminal preimage and spanning real inputs. No separate assertion that a prescribed nonzero ambient coefficient row lies in the image of \(\mathcal A\) is used. This is the mathematical m=2 LastX/MainX contrapositive, now with the actual coupled operator and its trace weights identified.

Only a complex top preimage is needed for this existence argument. The arithmetic coefficient modules and their pairings are defined over the fixed value field. Their complexifications are the actual finite-level coefficient/multiplicity spaces by DO038 and R032. A bilinear form nonzero after complexification is nonzero over that field, since its matrix is not the zero matrix. Multiplying selected rows by nonzero field denominators yields integral rows whenever one has chosen full integral lattices in those spaces. This elementary descent gives no effective denominator estimate. K059's independently constructed bounded integral spanning families can instead detect the nonzero bilinear form without constructing an integral right inverse of D.

The hypothesis in CI23 is specifically the actual moved period CI3 with an input in literal Mram. A nonzero K058 selected restriction before projection to Mram does not by itself satisfy it. A local nonzero test, a nonzero unrestricted double theta output, or a nonzero period at another level/point does not replace that hypothesis. Nor does the theorem establish the obstructed old HostRaw, a DP44 net saving, or an effective arithmetic preimage bound.

\subsection{7. Exact source anchors and verification scope}\label{exact-source-anchors-and-verification-scope}

The source statements are read at the stipulated pinned commit 7cbf7a0c32fca1badf05931879a37c3f8a2b03b4 and used at their accepted status. No Lean is run.

\begin{itemize}
\tightlist
\item
  \nolinkurl{LatticeFrames/ValLastX.lean148–150,194–210,228–254:} IsIo, ValLastX, KerX/PerX, TwoX and the exact LastX/MainX implication.
\item
  \nolinkurl{PlaceLaws.lean216–226,243–251,266–268:} original-map TLaw, QLast and the last scalar.
\item
  \nolinkurl{Columns/FiniteSumNonzeroLargeDepth.lean298–323,1080–1178,} 1410--1418,1451--1454: the coupled composition, the actual moving Y kernel and the invariant terminal coefficient pairing.
\item
  Columns/CoefModB.lean55--90 and the K037 source table: the fixed bilinear duality, Cauchy/subD comparison and complete contraction.
\item
  DO053(QT11--QT22), K038(CA24--CA36): exact orbit and paired covariance, original class weights and complete-type reading.
\item
  DO038(19d--19h), DO050(RC6--RC17), LIFT39(PN15--PN20): actual terminal branch, finite-dimensional recoupling and real span.
\item
  DO054(TC1--TC23), R032(R32.13--R32.22): common tensor, original-map period and the same-projector top realization.
\end{itemize}

The author reconstructed all CI1--CI23, including the entire-function middle substitution, the class-orbit formula with unequal stabilizers, the refined level volume, the exact compact coevaluation and inverse branch-period factor, the raw-to-unit vacuum conversion, every Q weight, the moved initial trace, and the actual-image contrapositive. The complete K037(CR17--CR29), K038(CA24--CA36), DO038 terminal proof, DO050 real projection and the literal ValLastX/LastX equations were re-read for this check.

The companion \nolinkurl{definite_oscillator_coupled_image_certificate.py} passes 6,660 exact diagnostics. They include genuine unequal finite stabilizers and level-volume refinement; an entire finite top-function orbit computation with both determinant characters; nonidentity-Gram compact Schur averages; complex branch phases and the literal terminal period; faithful polynomial frame evaluations; and the distinct complete/matched coevaluation factors. Its first development replay required simplifying an exact symbolic complex expression before comparison; no mathematical formula was changed by that diagnostic correction. No numerical automorphic period or vanishing/nonvanishing of the current Mram projection is claimed. No Lean is run.

\clearpage

\section{AM-R-036 --- an explicit fifth-power point in the same scalar T packet}\label{am-r-036-an-explicit-fifth-power-point-in-the-same-scalar-t-packet}

Author: GPT-6 Astra Pro, ALP\_MACHINE /root. 22 November 2026 UTC; effort is not separately exposed.

Status: author whole-proof and independent whole-proof checks PASS, 22 November 2026 UTC. Exact check events are recorded separately.

This construction supplies a second point in the same projective-T class as LIFT41's favorable point, with the same original principal G level. On J14 the new point is still very close to the standard origin, but it is separated from every integral arithmetic translate of the old point by an explicit positive distance. The proof uses the actual scalar-orbit lattice pencil, not merely the ambient Hecke double coset. No automorphic value, all-place ramified projection or terminal period is inferred from the geometric separation.

\subsection{1. The fixed data and a fifth-power identity}\label{the-fixed-data-and-a-fifth-power-identity}

Retain LIFT41's actual integers and vectors. Put

\[
 Q=g_0g_1,\quad k=Q^2,\quad U=7(11k-2),\quad W=kU,
 \quad M=11k-1,\quad \delta=\sqrt{-7},
 \tag{R36.1}
\] \[
 \pi=1+2D+D\delta,\quad T=\pi\bar\pi,
 \quad \pi_j=\pi-7j+2j\delta,
 \quad v_0=(Qu,\pi_k),\quad z=v_0/\bar\pi.
 \tag{R36.2}
\]

Here u is the exact four-supported primitive integral vector of LIFT41(NP4), extended by zeros, with u* u=U and u\_1=2. In both actual hosts Q is congruent to1 modulo7, D is divisible by Q squared, and T is the already fixed prime. No new prime is chosen.

Define the following explicit positive integers:

\[
 X_5=16M^5-20M^3+5M,\qquad
 t_5=16M^4-12M^2+1,\qquad j_5=(X_5+1)/11.
 \tag{R36.3}
\]

They satisfy

\[
 X_5^2-(M^2-1)t_5^2=1,\qquad
 77j_5^2-14j_5=W t_5^2,
 \tag{R36.4}
\] \[
 j_5-k=4M(4M^2-1)k(11k-2),\quad
 k\mid j_5,\quad U\mid(j_5-k).
 \tag{R36.5}
\]

For a direct proof, expand (M+sqrt(M squared-1)) to the fifth power. Its two coefficients are X\_5 and t\_5, and its norm is1. Since M is -1 modulo11, X\_5 is -1 modulo11, so j\_5 is integral. Replacing X\_5 by11j\_5-1 in the norm equation, and using M squared-1=11k(11k-2), gives the second equation of R36.4. Subtracting M from X\_5 and factoring gives R36.5's first equation. It proves k divides j\_5. Finally M is3 modulo7, so4M squared-1 is divisible by7; the same equation proves U divides j\_5-k.

For every integer j, direct multiplication in Q(delta) gives

\[
             \pi_j\bar\pi_j=T+77j^2-14j.
 \tag{R36.6}
\]

In this identity the terms involving D cancel. Thus

\[
       v_5=(Qt_5u,\pi_{j_5}),\qquad z_5=v_5/\bar\pi,
       \qquad v_5^*Jv_5=-T,\quad z_5^*Jz_5=-1.
 \tag{R36.7}
\]

The subscripts in this report refer to a fifth algebraic power, not a fifth theta lift. The construction remains a two-edge family.

\subsection{2. The ambient integral unitary matrix}\label{the-ambient-integral-unitary-matrix}

The rational plane spanned by u and the negative basis vector e has Gram diag(U,-1). Set

\[
 S=\begin{pmatrix}\bar\pi_k&Q\\QU&\pi_k\end{pmatrix},\qquad
 B_5=\begin{pmatrix}\bar\pi_{j_5}&Qt_5\\QUt_5&\pi_{j_5}\end{pmatrix}.
 \tag{R36.8}
\]

Both have determinant T and carry diag(U,-1) to T diag(U,-1). Therefore R\_5=B\_5 S inverse is special unitary on this plane. Extend it by the identity on u-perp. If x=bar(pi\_j5), x'=Qt\_5, its entries are

\[
 \alpha={x\pi_k-QUx'\over T},\qquad
 \beta={x'\bar\pi_k-Qx\over T},\qquad
 \alpha\bar\alpha-U\beta\bar\beta=1,
 \tag{R36.9}
\] \[
 R_5=\begin{pmatrix}
 I_a+{\alpha-1\over U}uu^*&\beta u\\
 \bar\beta u^*&\bar\alpha
 \end{pmatrix},\qquad R_5v_0=v_5.
 \tag{R36.10}
\]

The ambient factor U inverse is essential. An arbitrary plane-integral B would not by itself prove integrality of this ambient matrix. Here R36.5 supplies it. Indeed

\[
 x\pi_k-T=kU+(j_5-k)(-7-2\delta)\pi_k.
 \tag{R36.11}
\]

The right side is divisible by U, and in fact by kU. The latter assertion uses the first factored formula in R36.5 and its divisibility by7. Subtracting QUx'=kUt\_5 leaves the same divisibilities. Thus (alpha-1)/U belongs to (k/T)O\_E, while beta belongs to(Q/T)O\_E. Consequently R\_5 is integral away T and congruent to the identity at both actual G primes, with their conjugate components as well. Its unitary inverse is integral there too. There are no additional primes dividing U in its denominator.

The exact actual integer computations give gcd(Qt\_5,T)=1 on both hosts. Hence v\_5 is primitive at every finite place: a common prime divisor would divide its norm T, but cannot divide its first positive coordinate 2Qt\_5. These are the existing fields and lattices.

\subsection{3. Membership in the actual scalar packet}\label{membership-in-the-actual-scalar-packet}

We use the exact pencil of LIFT44(PR4--7,PR15--16), whose proof keeps the selected place and its inverse-adjoint CM partner. The relevant calculation is included here to fix the scalar unambiguously.

At a selected place above pi write the original finite decomposition gamma inverse=x\_H tau k\_v, where gamma e=z and x\_H belongs to the initial definite group. Let Lambda be the standard integral lattice, and put m=e-transpose k\_v. The actual primitive row q satisfies

\[
 q=-v_0^*J-\pi m,\qquad qz=0,\quad mz=1.
 \tag{R36.12}
\]

For the scalar representative delta\_s=diag(sI\_a,1), the actual relative inverse-image lattice is

\[
 k_v^{-1}\tau^{-1}\delta_s\tau\Lambda
   =(1+c\pi^{-1}zq)\Lambda,\qquad c=s^{-1}-1.
 \tag{R36.13}
\]

The permitted parameters are exactly s nonzero in the residue field. The projective pencil also has two omitted lines; they cannot be restored merely from the maximal-compact elementary divisors.

For the general matrix in R36.9--10 put trow=-v\_0\emph{J and trow'=-v\_5}J=(-bar(x')u*,x). Unitarity gives the exact relation trow R\_5 inverse=trow'. Direct multiplication modulo pi gives

\[
 \pi R_5^{-1}=z\lambda\,\mathrm{trow}'\pmod\pi,\qquad
 \lambda=1-{x'\bar\pi_k\over Qx},
 \tag{R36.14}
\]

where the equality describes the residue of the integral matrix pi R\_5 inverse. Keeping the full -pi m in R36.12 then gives q R\_5 inverse=(1-lambda)trow' modulo pi. Thus the pole coefficient of a vector in R\_5 inverse Lambda is lambda/(1-lambda) times its q-coordinate. This is exactly the graph in R36.13, with

\[
 \boxed{\qquad s_5={t_5\bar\pi_k\over\bar\pi_{j_5}}\pmod\pi.
 \qquad} \tag{R36.15}
\]

To justify equality of lattices, the residue calculation puts the entire inverse-image lattice inside that graph. Both lattices have determinant valuation zero, so the containment is equality. The denominator in R36.15 is a unit: its norm modulo T is W t\_5 squared, and W,t\_5 are T-units. Therefore s\_5 is permitted. The conjugate place has exactly the dual lattice by the unitary equation.

All these quantities lie in Q(sqrt(-7)), so at the selected degree-one places the scalar has the same residue in F\_T. An exact way to evaluate it is to replace delta by -(1+2D)/D modulo T. The certificate records its integer residue for both hosts and verifies s\_5\^{}28 is not1. In particular this change is neither identity nor the image of a global root of unity.

Set gamma\_5=R\_5 gamma. The lattice equality proves gamma\_5 x\_H delta\_(s5) tau is compact at T; away T compactness follows from R36.10--11. At G it remains in the original principal subgroup. Therefore

\[
 [(o,x_H\delta_{s_5}\tau)]=[(\gamma_{5,\infty}o,1_f)],
 \qquad \gamma_5 e=z_5.
 \tag{R36.16}
\]

The projective-T H class x\_H has not changed. Thus this is a point in the same literal scalar orbit, not a new choice of positive complement, Schwartz function, prime, or auxiliary finite level.

\subsection{4. Exact radii and separation on J14}\label{exact-radii-and-separation-on-j14}

At every real label let r\_0 and r\_5 be the Cartan distances of the two negative lines from the standard origin. Their exact values are

\[
       \sinh^2 r_0=W/T,\qquad
       \sinh^2 r_5=Wt_5^2/T.
 \tag{R36.17}
\]

On J14 the positive integer Wt\_5 squared has932 bits, whereas T has 1001 bits. Direct integer cross-multiplication gives the strict bounds

\[
 2^{-70}<{Wt_5^2\over T}<2^{-68},\qquad
 {W\over T}<2^{-813}.
 \tag{R36.18}
\]

Consequently r\_0\textless2\^{}-406 and r\_5\textless2\^{}-34. Also sinh(r\_5)\textgreater2\^{}-35. For 0\textless x\textless1/2 the integral formula for asinh gives asinh(x)\textgreater3x/4. Since r\_0\textless2\^{}-37, these estimates imply

\[
                  r_5-r_0>2^{-36}.\tag{R36.19}
\]

Let Gamma be the actual integral maximal group of K058, and let S\_0=Gamma intersect K be its full monomial compact stabilizer. Write p=gamma-infinity o, p\_5=gamma\_5-infinity o, and d\_tau for the Cartan distance in each ball. Then

\[
 \gamma_0\in S_0\quad\Longrightarrow\quad
       d_\tau(p_5,\gamma_0p)>2^{-36}\quad\hbox{for every tau},
 \tag{R36.20}
\] \[
 \gamma_0\in\Gamma\setminus S_0\quad\Longrightarrow\quad
       d_\tau(p_5,\gamma_0p)>1/4\quad\hbox{for some tau}.
 \tag{R36.21}
\]

For R36.20, compact elements preserve distance from the origin, so the triangle inequality bounds the displayed distance below by r\_5-r\_0. For R36.21, suppose every displayed distance were at most 1/4. Another triangle inequality gives d\_tau(o,gamma\_0 o)\textless=r\_5+1/4+r\_0\textless1/2 at every label. Each positive entry of the last column of gamma\_0 would then have all conjugates of modulus at most sinh(1/2)\textless1. They are algebraic integers, so all vanish. Unitarity forces gamma\_0 into the compact block, contradicting gamma\_0 not in S\_0. This is the same integral-entry argument used in K058, with its exact original maximal group. It proves both assertions without assuming that every arithmetic translate stays in the plane.

In particular the two points are distinct in the actual maximal arithmetic quotient. This argument rules out arithmetic folding for these two points inside the specified small origin ball.

On J56 the same fifth-power construction is an actual scalar packet point, but Wt\_5 squared has1448 bits while T has996 bits. Its radial quotient exceeds2\^{}450. The small-ball argument and the J14 separation conclusion are not asserted on J56.

\subsection{5. Scope, failed shortcuts, and verification}\label{scope-failed-shortcuts-and-verification}

The established result is an explicit same-projective-T point and its J14 geometric separation, retaining the principal G and conjugate T conditions. It does not estimate an automorphic value there. In particular K046's selected W-to-minimal-type seed has only its proved radius4.5/5 L1 concentration. Its interior degrees grow with ell; an ell\^{}-1/2 concentration assertion for that full seed would be unsupported. Likewise a small lowering depth19/12 at a highest vector's final coordinate does not bound the depth of every vector in the full selected compact module. A separate W-to-W coefficient or period argument is required.

Even proving that one bottom restriction changes under the displayed diagonal scalar would only force some nontrivial joint dilation character. It would not force a character nontrivial at every one of the selected T places. The literal Mram condition is the joint intersection in LIFT43(TG1--6). No MainX nonzero arithmetic row or DP44 cost sign follows here without that additional argument.

Inputs are LIFT41(NP1--12), the exact pencil and ambient matrix test of LIFT44(PR4--7,PR15--16), and K058(PG1--4)'s actual integral group and compact stabilizer. Their source locators and check histories remain in those reports. The new fifth-power arithmetic, actual ambient integrality, permitted scalar, radii and separation are proved above. The companion root certificate checks the polynomial identities on finite diagnostics and both actual integer instances, the exact quadratic-field unitary matrices, denominator and G congruences, permitted scalar residues, and all asserted bit bounds. It does not compute a Poincare series, selected Gram or theta value.

\clearpage

\section{AM-K060 --- a uniform full-selected-component rank-one compression bound}\label{am-k060-a-uniform-full-selected-component-rank-one-compression-bound}

Author: GPT-6 Astra Pro, ALP\_MACHINE \nolinkurl{/root/kummer_extension}. 22 November 2026 UTC. Inherited reasoning effort is not separately exposed.

\textbf{Author whole-proof check PASS --- 2026-11-22T08:08:19.089080+00:00.} The exact companion certificate passes975 assertions; independent reviews are separate events.

This is an application of Haagerup and Schlichtkrull's uniform Jacobi inequality to the actual vector-valued holomorphic discrete series. It controls the full selected compact component, with the K043 PBW isometry, independently of its large polynomial degree. The proof uses the negative compact-circle grading to diagonalize the compression on each rank-one ladder. It does not claim that the selected-to-minimal seed is localized at radius ell\^{}(-1/2).

The numerical application is to R036's second J14 point in the same actual scalar T packet as LIFT41. It gives small local selected-to- selected coefficients against every integral translate, with two different separation bounds. It is not an absolute estimate of the infinite automorphic sum and does not prove a nonzero Mram projection.

\subsection{1. The actual Hilbert model and two commuting gradings}\label{the-actual-hilbert-model-and-two-commuting-gradings}

At a CM label use exactly K043(PB2--PB8). Let H be the unitary holomorphic discrete series of U(a,1), whose minimal compact type is V\_alpha tensor z\^{}beta. Its polynomial realization is the Hilbert completion of

\[
       \bigoplus_{N\ge0}\mathcal H_N,
 \qquad \mathcal H_N=V_\alpha\otimes\operatorname{Sym}^N(\mathbb C^a).
 \tag{RC1}
\]

The inner product on different Pieri components is the actual Nakahama/K043 inner product, not one unmodified Fischer metric on the entire direct sum. The negative U(1) factor acts on H\_N by z\^{}(beta−N). Since it acts unitarily, the finite dimensional spaces H\_N are mutually orthogonal. Write L for the nonnegative degree operator, with L\textbar H\_N=N. The full compact group preserves this grading.

Fix a unit positive coordinate and the negative coordinate, giving a subgroup SU(1,1). Let B be its creation operator, C=B* its lowering operator, and H\_1={[}C,B{]} the positive compact energy operator. Choose the matrix-unit signs as in K043/K048. On the algebraic vectors,

\[
 [H_1,B]=2B,\quad [H_1,C]=-2C,\quad
 [L,B]=B,\quad [L,C]=-C,\quad [L,H_1]=0.
 \tag{RC2}
\]

Thus L−H\_1/2 commutes with this rank-one Lie algebra. These are operator identities, not approximations on a highest vector.

For a monomial z\^{}gamma v of degree N, if v has first fibre weight nu\_1, its H\_1 energy is

\[
                 \nu_1-\beta+N+\gamma_1.
 \tag{RC3}
\]

Every fibre weight coordinate is at least alpha\_(a−1): after subtracting that common determinant weight, V\_alpha is a polynomial Schur representation and all its weight coordinates are nonnegative. Hence all energies obey

\[
  H_1\ge\kappa_0:=\alpha_{a-1}-\beta
       =b_1-b_0+a=2\ell+q_s+a>1.                            \tag{RC4}
\]

At bounded H\_1 energy, RC3 bounds N, so the energy eigenspaces are finite dimensional. The fixed actual selected module W is one compact Pieri summand of H\_(N\_tau), with N\_tau and the weights specified in K043. Its isometric inclusion is rho\_tau\^{}(-1) C\_PBW j, where rho\_tau\^{}2 is the exact product of Pochhammer factors PB5.

\subsection{2. Orthogonal ladders compatible with the degree}\label{orthogonal-ladders-compatible-with-the-degree}

\textbf{Lemma RC5.} There is an orthogonal Hilbert sum decomposition into positive SU(1,1) lowest-weight ladders, with each lowest vector also homogeneous for L. If a ladder begins in degree N0 and energy kappa, then its normalized basis e\_n has

\[
 \begin{aligned}
 L e_n&=(N_0+n)e_n,& H_1e_n&=(\kappa+2n)e_n,\\
 B e_n&=\sqrt{(n+1)(n+\kappa)}\,e_{n+1},&
 C e_n&=\sqrt{n(n+\kappa-1)}\,e_{n-1}.
 \end{aligned}                                                \tag{RC5}
\]

Here kappa\textgreater=kappa0. Multiplicity spaces are retained, and the decomposition is in the actual invariant Hilbert norm.

\textbf{Proof.} At each positive energy choose the orthogonal complement to the ladders already generated from lower energies. It is annihilated by C: by adjunction, C of such a vector pairs to zero against every lower-energy vector; the latter have already been exhausted. This kernel is L-stable by RC2. Choose an orthonormal basis of simultaneous H\_1,L eigenvectors in it. For a vector of energy kappa, the commutator relation and C e\_0=0 give \textbar\textbar B\^{}n e\_0\textbar\textbar\^{}2=n!(kappa)\_n by induction, proving RC5.

Distinct starting vectors produce orthogonal ladders. To check this even when their energies differ, move the smaller number of creation operators across the inner product. Unless the residual degrees in B are zero, the resulting inner product is zero because one lowest vector is annihilated by C. In the remaining case it is the original orthogonality of the starting vectors. Repeating the construction over the discrete energies exhausts the dense algebraic subspace: a nonzero orthogonal remainder would have a lowest energy, which the construction would have included. Finally RC2 raises L by one at each step, and RC4 gives kappa\textgreater=kappa0. QED.

The closed ladders in this argument are invariant under the actual rank-one group, not just its formal Lie algebra. Indeed RC5 bounds the norm of (B−C)\^{}m e\_n by 2\^{}m(n+kappa+2)\_m. After division by m!, its power series converges for \textbar s\textbar\textless1/2. Thus each basis vector is analytic for the Cartan generator, its small-time exponential stays in the closed ladder, and iteration gives invariance for every real s. The same applies to the compact circle. This also identifies each ladder with the disk model below from their identical normalized generator actions.

This proof can equivalently be made degree by degree, taking the orthogonal complement of B H\_(N−1) in H\_N and diagonalizing H\_1 there. The positive lower bound in RC4 excludes finite ladders.

Let P\_N be the orthogonal projection to H\_N. A fixed H\_N meets each ladder in at most one weight line: its index is exactly n=N−N0. Consequently, for a positive Cartan element a\_s, the operator P\_N pi(a\_s) P\_N is the orthogonal direct sum of the diagonal coefficients on those lines. In particular

\[
 \|P_N\pi(a_s)P_N\|\le
    \sup_{\kappa\ge\kappa_0,\ n\ge0}
       |\langle\pi_\kappa(a_s)e_n,e_n\rangle|.               \tag{RC6}
\]

The assertion would fail without retaining this homogeneous grading: an unrestricted unitary ladder operator has norm one. No summation of the bounds over n or over ladder multiplicities occurs in RC6.

\subsection{3. Exact diagonal coefficient and the primary uniform bound}\label{exact-diagonal-coefficient-and-the-primary-uniform-bound}

Write c=cosh(s), h=sinh(s), t=tanh(s), with s\textgreater0. The positive lowest-weight representation of energy kappa\textgreater1 has the unit disk model with monomial norm

\[
              \|z^n\|^2={n!\over(\kappa)_n}.
 \tag{RC7}
\]

For example, this follows by integrating against (kappa−1)(1−\textbar z\textbar\^{}2)\^{}(kappa−2) dA/pi and the elementary beta integral. The Cartan action may be written

\[
       f(z)\longmapsto(c+hz)^{-\kappa}
                           f\left({h+cz\over c+hz}\right).
 \tag{RC8}
\]

Changing the inverse convention replaces s by −s and does not change the following diagonal coefficient. Normalizing z\^{}n as in RC7 cancels the same scalar in its input and output. The coefficient of z\^{}n in (h+cz)\^{}n(c+hz)\^{}(-kappa−n) is therefore

\[
 \begin{split}
 d_{\kappa,n}(s)
  &=c^{-\kappa}\sum_{j=0}^n
           (-1)^j {n\choose j}{(n+\kappa)_j\over j!}t^{2j}\\
  &=c^{-\kappa}\,{}_2F_1(-n,n+\kappa;1;t^2)\\
  &=c^{-\kappa}P_n^{(0,\kappa-1)}(1-2t^2).
 \end{split}                                                \tag{RC9}
\]

The first equality is a finite binomial expansion, so it also checks the matrix-unit and factor-of-two convention independently of an asymptotic formula. For n=0 it is c\^{}(-kappa), and for n=1 it is c\^{}(-kappa)(1−(kappa+1)t\^{}2).

Haagerup and Schlichtkrull (2014, Theorem 1.1, p.228) define

\[
 g_n^{(A,B)}(x)=
  \left({\Gamma(n+1)\Gamma(n+A+B+1)
          \over\Gamma(n+A+1)\Gamma(n+B+1)}\right)^{1/2}
  \left({1-x\over2}\right)^{A/2}
  \left({1+x\over2}\right)^{B/2}P_n^{(A,B)}(x)
 \tag{RC10}
\]

and prove a universal inequality with a constant strictly below12:

\[
       (1-x^2)^{1/4}|g_n^{(A,B)}(x)|
           \le C(2n+A+B+1)^{-1/4},\qquad C<12,
 \tag{RC11}
\]

for n\textgreater=0, A,B\textgreater=0 and −1\textless=x\textless=1. This is a scalar polynomial inequality; its use here does not confuse a compact SU(2) representation with our noncompact SU(1,1) representation.

At A=0, B=kappa−1, x=1−2t\^{}2, the gamma factor in RC10 is exactly one and g=c\^{}(-(kappa−1))P. Thus d=c\^{}(-1)g. Moreover

\[
        c^4(1-x^2)=4h^2.
 \tag{RC12}
\]

Combining RC9--12 gives the uniform, degree-independent bound

\[
 \boxed{\quad
 |d_{\kappa,n}(s)|
   <{12\over\sqrt{2\sinh s}\,(2n+\kappa)^{1/4}}.
 \quad}                                                       \tag{RC13}
\]

Unitarity also bounds its absolute value by one. This bound retains the full polynomial rather than estimating its alternating sum by the sum of absolute coefficients.

\subsection{4. The full selected-module theorem}\label{the-full-selected-module-theorem}

Let r(g)\textgreater=0 be the Cartan parameter of g in U(a,1), normalized so its noncompact two-plane block has diagonal cosh(r) and off-diagonal sinh(r). The compact central scalar has modulus one.

\textbf{Theorem RC14.} For any fixed homogeneous degree N, and any subspace W of H\_N with its inherited Hilbert norm, the compression satisfies, when r(g)\textgreater0,

\[
 \boxed{\quad
 \|P_W\pi(g)|_W\|
  \le \min\left\{1,
    {12\over\sqrt{2\sinh r(g)}\,\kappa_0^{1/4}}\right\}.
 \quad}                                                       \tag{RC14}
\]

The norm is strictly less than the second displayed quantity; the minimum also records the contraction bound. In particular this applies to the complete actual selected compact type W, for every R022/K044 word and role, with kappa0=2ell+q\_s+a. It does not require a bounded lowering depth.

\textbf{Proof.} RC6 and RC13 bound the homogeneous compression at a\_s. In a Cartan decomposition g=k1 a\_s k2, the compact factors preserve H\_N. The full homogeneous compression is therefore unitarily equivalent on its two sides to P\_N pi(a\_s)P\_N. Compressing further to W cannot increase the norm. This argument does not even require W to be compact-stable, although the actual selected W is. The bound one is the contraction norm of a projection of a unitary operator. QED.

For the product of CM labels, write W=the tensor product of the actual W\_tau and g=(g\_tau). The exact tensor norm is the product, so

\[
 \|P_W\pi(g)|_W\|
 \le\prod_\tau
   \min\left\{1,{12\over\sqrt{2\sinh r_\tau(g)}
                          (2\ell+q_s+a)^{1/4}}\right\}.
 \tag{RC15}
\]

A factor with r\_tau=0 is interpreted as one. No selected dimension factor occurs in this operator norm. A Hilbert--Schmidt or trace bound derived from it must retain the usual dimension factor; for example \textbar Tr(A P\_W pi(g)\textbar W)\textbar\textless=rank(W)\textbar\textbar A\textbar\textbar{} times the right side.

\subsection{5. Application to the actual fifth-power packet point}\label{application-to-the-actual-fifth-power-packet-point}

Use the J14 upper gap ell\_+, with every allowed word, role assignment and canonical twist. K048(ST12) gives

\[
                    \kappa_0\ge2\ell_++a>2^{485}.
 \tag{RC16}
\]

Let p be the actual near point of LIFT41 and p5 the fifth-power point of R036. That construction retains the original principal G level and the same projective-T class, with a scalar whose28-th power is not one. Write Gamma for the original maximal integral arithmetic group and S0=Gamma intersect K for its full compact stabilizer. R036(R36.20--21) gives the two geometric statements

\[
 \begin{split}
 \gamma\in S_0&\quad\Longrightarrow\quad
       r_\tau(p_5,\gamma p)>2^{-36}\ \text{at every label},\\
 \gamma\notin S_0&\quad\Longrightarrow\quad
       r_\tau(p_5,\gamma p)>1/4\ \text{at some label}.
 \end{split}                                                \tag{RC17}
\]

For the local operator below choose any g\_p,g\_5 carrying the origin to these points, with their fixed coefficient-frame identifications. Put g\_gamma=g\_5\^{}(-1) gamma g\_p.~Its Cartan parameter is the distance in RC17; compact frame choices only multiply the compression by unitaries because W is an actual compact type.

\textbf{Corollary RC18.} Subject only to the geometric statement RC17, the complete normalized selected-to-selected operator has

\[
 \boxed{\quad
 \|P_W\pi(g_5^{-1}\gamma g_p)|_W\|
    <\begin{cases}
       2^{-4200},&\gamma\in S_0,\\
       2^{-117},&\gamma\in\Gamma\setminus S_0.
      \end{cases}
 \quad}                                                       \tag{RC18}
\]

\textbf{Proof.} If r\textgreater2\^{}(-36), then sinh(r)\textgreater2\^{}(-36). The fourth power of RC13's bound is smaller than 12\^{}4/(4 times2\^{}(-72) times2\^{}485). Since12\^{}4=20736\textless2\^{}15, this is less than2\^{}(-400). Thus every one of the42 factors in the first case is smaller than2\^{}(-100). If r\textgreater1/4, the corresponding fourth power is smaller than12\^{}4/(4 times(1/4)\^{}2 times2\^{}485), which is less than2\^{}(-468). That factor is smaller than2\^{}(-117), and all remaining factors are at most one. QED.

This corollary uses only the displayed distance thresholds, so its operator estimate can be checked independently of the Pell arithmetic. At the time of this draft R036 is a separately authored geometric finding; its status is retained in the verification record below. No assertion is made here about the J56 fifth-power point, whose distance regime is different.

\subsection{6. Exact normalization and limitations of the conclusion}\label{exact-normalization-and-limitations-of-the-conclusion}

The local normalized selected evaluation at gp is P\_W pi(gp\^{}(-1)). Its adjoint is pi(gp) followed by the isometric inclusion of W. The two-point local Gram is therefore exactly

\[
  J_{p_5}J_p^*=P_W\pi(g_5^{-1}g_p)|_W.
 \tag{RC19}
\]

The same equation with gamma inserted is the local term associated to that arithmetic translate. It uses K043's selected isometry. Replacing both normalized selected derivatives by raw Fischer derivatives multiplies this Gram at a label by rho\_tau\^{}2, with the exact Pochhammer product PB5. A non-isometric compact Pieri inclusion also retains its scalar as in K043(PB8). Those factors are not suppressed by RC18.

In particular, RC18 is not an assertion that the actual global automorphic cross Gram equals its identity-translate term. Such a Gram contains finite-projector matrices, rational translates, class masses and, for Mram, all actual scalar orbit terms. Its infinite sum cannot be bounded by summing the uniform2\^{}(-117) in RC18. Its large-distance tail and its finite-character correlations must still be treated with their actual measures and indices. Nor does a small local coefficient by itself show that the G-generic row constructed in K058 has nonzero all-place Mram component.

The full-W route avoids a false bounded-depth shortcut. At a top label the selected highest weight has lambda\_0=A and lambda\_(a−1)=−B, so its spread is A+B. Only the highest-vector last skew degree lambda\_(a−1)−alpha\_(a−1) equals19 or12. A Weyl-conjugate selected weight does not retain that small bound in a fixed positive coordinate. None of RC5--18 assumes it. One could restrict to the appropriate U(a−1) branch with that last weight, but this additional restriction is unnecessary for RC14.

\subsection{7. Primary sources and verification record}\label{primary-sources-and-verification-record}

Haagerup, U., \& Schlichtkrull, H. (2014). Inequalities for Jacobi polynomials. \emph{The Ramanujan Journal, 33}, 227--246. \url{https://doi.org/10.1007/s11139-013-9472-4} The exact primary anchor is Theorem1.1 and the definition of g on printed p.228, with C\textless12 stated immediately after the theorem. Published PDF: \url{https://www.isibang.ac.in/~soumyashant/misc/collected-works-of-haagerup/2010s/2014_Inequalities_for_Jacobi_polynomials.pdf} Author version: \url{https://arxiv.org/pdf/1201.0495v2}, pp.1--2. Both PDFs and extracted text are retained in \nolinkurl{work/agents/selected_rank_one_refs/}.

Nakahama, R. (2016). Norm computation and analytic continuation of vector-valued holomorphic discrete series representations. \emph{Journal of Lie Theory, 26}, 927--990. \url{https://jolt.centre-mersenne.org/item/10.5802/jolt.919.pdf} The actual action and norm conventions are equation(5.5), p.959, and Theorem5.1, pp.960--961, as already directly checked in K043.

K043(PB2--PB10) fixes the actual selected isometry and raw PBW factor. K044(EG1--5) gives the real profile and canonical twist. K048(ST6--12) gives the compact energy and its exact upper-gap lower bound. R036 supplies only RC17's arithmetic geometry.

The companion certificate checks the finite binomial/hypergeometric/ Jacobi identities, the fourth-power change of variables, graded ladder normalizations in finite ranges, both exact dyadic constants, and the two actual degree/weight conventions. These are exact normalization checks. They do not numerically sum an automorphic series or establish its Mram nonvanishing.

\subsubsection{Author whole-proof check}\label{author-whole-proof-check-2}

The entire saved RC1--19 and companion certificate were read and reconstructed. Checked the actual negative-circle grading, the positive energy floor, orthogonal degree-compatible ladders and analytic-vector group invariance, disk coefficient normalization, exact Jacobi conversion, absence of a ladder multiplicity factor, compact-frame transport, tensor operator norm, raw PBW scale and both dyadic numerical margins. The published primary Theorem1.1 and its C\textless12 statement were read directly in the original PDF text and via the retrieved primary page. R036 arithmetic geometry is a separate input; no automorphic summation or Mram nonzero claim has been introduced.

\clearpage

\section{AM-R-037 --- sharp normalized Fischer bounds for the actual standard dual-Weyl generators}\label{am-r-037-sharp-normalized-fischer-bounds-for-the-actual-standard-dual-weyl-generators}

Author: GPT-6 Astra Pro, ALP\_MACHINE /root. 22 November 2026 UTC; effort is not separately exposed.

Status: author and independent whole-proof reviews PASS; exact events recorded separately.

This result replaces the coarse coefficient factor B\_T in K059 by one for its literal standard semistandard dual-Weyl basis. It concerns that fixed standard frame. It neither bounds the clearing integers d\_c for arbitrary class frames nor asserts the corresponding bound for a pairBK-dual Weyl basis.

\subsection{1. Literal polynomials and their normalization}\label{literal-polynomials-and-their-normalization}

Let V=C\^{}N have its standard Hermitian metric. Choose a finite list of column sizes s\_1,\ldots,s\_m, with 0\textless=s\_i\textless=N. Write {[}s{]}=\{1,\ldots,s\}. For every increasing subset I\_i of \{1,\ldots,N\} of size s\_i define

\[
 F_I(Z)=\prod_{i=1}^m\det Z_{I_i,[s_i]},\qquad
 F_{\rm hi}(Z)=\prod_{i=1}^m\det Z_{[s_i],[s_i]}.
 \tag{R37.1}
\]

Empty determinants and empty tensor factors are one. For a polynomial F=sum\_alpha c\_alpha Z\^{}alpha use the source Fischer norm

\[
 \|F\|_{\rm F}^2=\sum_\alpha |c_\alpha|^2\alpha!,\qquad
 H=\|F_{\rm hi}\|_{\rm F}^2>0.
 \tag{R37.2}
\]

These are exactly StatementA.initMinor/bideterminant and hw, followed by DataCoeff.extremal and CoefModB.extNormB. If nu\_j is the number of s\_i at least j, the highest polynomial in R37.1 is product\_j det Z\_{[}j{]},{[}j{]}\^{}(nu\_j-nu\_(j+1)); it therefore equals the literal hw for this shape, independent of the ordering of the list.

\subsubsection{Theorem R37.3}\label{theorem-r37.3}

Every product in R37.1 satisfies

\[
                         \|F_I\|_{\rm F}^2\le H.
 \tag{R37.3}
\]

In particular this is true of every semistandard standard-basis polynomial of the accepted SSBasis theorem. The normalized bound is sharp, since F\_hi attains equality.

\subsection{2. Proof by orthogonal projection in an exterior tensor space}\label{proof-by-orthogonal-projection-in-an-exterior-tensor-space}

Give each exterior power of V the metric in which its increasing wedge basis is orthonormal. Consider the finite unitary tensor space

\[
 \mathcal T=\bigotimes_{i=1}^m\bigwedge^{s_i}V,
 \quad h=\bigotimes_i e_{[s_i]},
 \quad e_I=\bigotimes_i e_{I_i}.
 \tag{R37.4}
\]

Both h and every e\_I have norm one. Let rho\_T be its polynomial representation of all N by N matrices. The GL\_N orbit of h spans the irreducible highest-weight summand U of weight nu, often called the Cartan summand. Only its usual finite-dimensional compact representation meaning is used here. To fix precisely which copy is meant, h is a highest vector, and the weight-nu space of T is one dimensional: equality in every initial-coordinate weight sum forces each wedge subset to be {[}s\_i{]}. Finite-dimensional unitary complete reducibility and the highest-weight decomposition then put h in exactly one irreducible summand. Thus U is unambiguous and has an orthogonal projection P\_U with norm one.

Let the Hermitian product be linear in its first argument. Define the conjugate-linear polynomial map

\[
             \Phi(v)(Z)=\langle\rho_T(Z)h,v\rangle,
                    \qquad v\in U.
 \tag{R37.5}
\]

The image rho\_T(Z)h belongs to U for every matrix Z. It does so for invertible Z by the definition of U, and for arbitrary Z by polynomial continuity and density of invertible matrices. The determinant coefficient formula for an exterior power therefore gives the exact identity

\[
                    F_I=\Phi(P_Ue_I),\qquad
                    F_{\rm hi}=\Phi(h).
 \tag{R37.6}
\]

This identity has no factorial or antisymmetrizer scale: the wedge bases in R37.4 are unit bases, and the coefficient of e\_I in wedge\^{}s(Z)e\_{[}s{]} is precisely det Z\_I,{[}s{]}.

The pulled-back Fischer quadratic form v-\textgreater\textbar\textbar Phi(v)\textbar\textbar\_F\^{}2 is invariant under U(N). Indeed, for a unitary matrix k,

\[
          \Phi(\rho_T(k)v)(Z)=\Phi(v)(k^{-1}Z).
 \tag{R37.7}
\]

Left multiplication of the matrix variables by k inverse is unitary on their Euclidean space, so it preserves the Fischer form. This last assertion is also immediate from its normalized Gaussian integral: the integral of Z\^{}alpha conjugate(Z)\^{}beta is zero unless alpha=beta, in which case it is alpha!. A unitary change of variables preserves both the Gaussian and Lebesgue measure.

Since U is irreducible, the finite-dimensional unitary Schur lemma makes this invariant quadratic form a scalar multiple of its given Hermitian form. Evaluating at the unit vector h fixes that scalar:

\[
                        \|\Phi(v)\|_{\rm F}^2=H\|v\|^2.
 \tag{R37.8}
\]

This also proves that Phi is injective; alternatively its kernel is orthogonal to the spanning GL\_N orbit of h. Finally R37.6 and the contractivity of an orthogonal projection give

\[
       \|F_I\|_{\rm F}^2
          =H\|P_Ue_I\|^2\le H\|e_I\|^2=H.
 \tag{R37.9}
\]

This proves the theorem. No estimate of an alternating determinant expansion, coefficient l1 norm or polynomial degree is needed.

\subsection{3. Actual source basis and tensor products of CM labels}\label{actual-source-basis-and-tensor-products-of-cm-labels}

The accepted SSBasis statement chooses a basis of the standard dual-Weyl lattice, and explicitly says that each basis vector is one bideterminant with shapeDeg=nu and strictly increasing row sets. Every such vector is covered by R37.3, even if its bideterminant description is not unique. DataCoeff uses nu=lam.natW. The possibly negative determinant exponent is carried separately by sqAct; it does not insert a determinant denominator into these polynomials. On the compact real group this extra determinant character has absolute value one.

Consequently, at each CM label tau and for every one of those basis vectors p\_(tau,i), the literal source metric is

\[
             \|p_{\tau,i}\|_{\rm nf}^2
                  =\|p_{\tau,i}\|_{\rm F}^2/H_\tau\le1.
 \tag{R37.10}
\]

For the tensor basis over all labels used in K059, polynomial variables at different labels are disjoint. Both the Fischer norm and H are products. Therefore

\[
                    B_i=\|p_i\|_{\rm nfAll}\le1,
                    \qquad B_T=1.
 \tag{R37.11}
\]

The highest-weight space is one dimensional. Its SSBasis vector has the same leading monomial and leading coefficient one as hw, so it is exactly hw. This proves equality B\_T=1. The inequality B\_T\textless=1 alone is sufficient for all following bounds. A zero-degree factor has H=1 and causes no exception.

Thus K059's actual terminal generators obey the sharper bounds

\[
 \|s_{c,i}\|_{\rm cls}\le\sqrt{|\Gamma_c|}\,d_c,\qquad
 \|z_{c,i}\|_{\rm cls}\le d_G\sqrt{|\Gamma_c|}\,d_c,
 \tag{R37.12}
\] \[
 \|E_T(z_{c,i},v)\|_{\rm aut}
 \le \sqrt{v_T/\dim V_T}\,\|Jv\|\,
                       d_G\sqrt{|\Gamma_c|}\,d_c.
 \tag{R37.13}
\]

The factors in these equations are unchanged K059 factors; only B\_i has been bounded by one. If starting at principal G level, retain the additional sqrt(q\_P) from IS14. This is not a claim that the individual standard generators are orthonormal. Their Gram matrix may have small eigenvalues, and the dual Weyl basis may have large norm. In particular no corresponding B\_(0,i)\textless=1 is asserted for K059's pairBK-dual initial Weyl basis. The representation determinant shift still affects finite frame membership and d\_c.

\subsection{4. Scope and exact anchors}\label{scope-and-exact-anchors}

The theorem is about the fixed standard polynomial lattice. Arbitrary transported class frames can have large d\_c; this result does not bound their heights, their reduction cost, the actual image content, or a nonzero arithmetic matrix entry. It supplies no new integral vertical inverse and proves no Mram period. It does not change the DP44 host/Scale obstruction or the specific032 cycle identification.

Direct source anchors, at the stipulated pinned archive commit:

\begin{itemize}
\tightlist
\item
  \nolinkurl{Formalism/GradedWeylDivisibility/StatementA.lean42–67:} initMinor, bideterminant, shapeDeg and hw.
\item
  \nolinkurl{Formalism/GradedWeylDivisibility/StatementD.lean24–29:} accepted SSBasis and its literal product-of-minors basis vectors.
\item
  \nolinkurl{Formalism/GradedWeylDivisibility/FinalD.lean18:} accepted sSBasis theorem. Its theorem status is stipulated; this argument does not audit or build Lean.
\item
  \nolinkurl{Formalism/LocalRingCardinality/DataCoeff.lean49–62,84–85:} nablaSq, extremal, determinant twist and normalized Fischer form.
\item
  \nolinkurl{Formalism/LocalToGlobal/Columns/CoefModB.lean18,55–69:} literal fockB, extNormB and pairBK.
\item
  AM-K059 IS2,IS4,IS7--IS16: class masses, clearing integers, actual tensor basis, genuine projector and top norm.
\end{itemize}

The compact highest-weight decomposition and unitary Schur lemma are the same standard finite-dimensional representation inputs used in the previously checked coefficient and compact averaging results. The new assertion is the exact norm comparison R37.3 and its consequence for the prescribed integral generators.

\clearpage

\section{ALP-LIFT-45 --- single-place scalar separation and the exact all-T projection criterion}\label{alp-lift-45-single-place-scalar-separation-and-the-exact-all-t-projection-criterion}

Author: GPT-6 Astra Pro, ALP\_MACHINE \nolinkurl{/root/lift_accounting}. 22 November 2026 UTC. Inherited reasoning effort is not separately exposed.

Author whole-proof check PASS: 2026-11-22T08:22:00.655654+00:00. The full proof and complete exact certificate have been read; independent review is separate. Whole-proof reread after the SP11 norm-display repair: 2026-11-22T08:23:52.187485+00:00, PASS; no mathematical change.

This result makes two precise reductions. First, the intersection of all the required nontrivial T-character subspaces can be tested by the sum of the \textbf{single-place} averaging energies, with the original class masses. Second, at LIFT41's fixed favorable point a scalar change at one T-place has a uniform geometric separation from every integral translate of that point. The latter uses the actual lattice pencil and the small number of denominator primes, rather than an operator-norm estimate mistaken for point separation.

These statements do not yet prove a nonzero Mram projection. The projection criterion averages over all scalar-orbit base points, whereas the proved geometric bound has the favorable point as its base. That remaining distinction is explicit below.

\subsection{1. The exact projection and mass convention}\label{the-exact-projection-and-mass-convention}

Retain the actual detector row realization V, opposite to the initial arithmetic row, and its raw adelic mass-weighted metric from LIFT43. Let \(\mathcal T\) be the f selected places above the fixed rational prime T, and put

\[
D=\prod_{v\in\mathcal T}D_v,\qquad D_v=\mathbb F_T^\times,
\quad A_v=\frac1{T-1}\sum_{s\in D_v}R_v(s),
\quad P_{\rm ram}=\prod_v(1-A_v).
\tag{SP1}
\]

The \(R_v(s)\) are the actual unitary diamond operators. They commute between places. Each \(A_v\) is an orthogonal projection, and

\[
\boxed{\qquad P_{\rm ram}\ \ge\ I-\sum_{v\in\mathcal T}A_v.
\qquad}\tag{SP2}
\]

Here the inequality is of Hermitian operators. Simultaneously diagonalize the commuting projections. On a joint eigenspace let k be the number of \(A_v\)'s whose eigenvalue is one. The two sides have eigenvalues \(\mathbf1_{k=0}\) and \(1-k\), respectively. Their difference is zero for \(k=0,1\) and is \(k-1\ge0\) otherwise. This proves SP2 without a dimension or semisimplicity assumption beyond the finite-dimensional row space.

Let \(J:B\to V\) be the actual selected restriction of any fixed finite-dimensional genuine-good bottom space. It may be the K058 constructed space or the full fixed-level space. Define

\[
E=\|J\|_{\rm HS}^2,\qquad E_v=\|A_vJ\|_{\rm HS}^2,
\qquad E_{\rm ram}=\|P_{\rm ram}J\|_{\rm HS}^2.
\]

Taking the trace of \(J^*(\cdot)J\) in SP2 gives

\[
\boxed{\qquad E_{\rm ram}\ge E-\sum_vE_v.\qquad}
\tag{SP3}
\]

In particular, \(\sum_vE_v<E\) is sufficient for actual Mram nonvanishing. This is a sufficient criterion, not an equivalence.

Use the exact projective-T classes x and their stabilizers \(\Gamma_x\) from LIFT43(TG4). If \(v_P\) is the projective compact volume, let

\[
w_x=\frac{v_P}{|\Gamma_x|},\qquad
J_{x,t}:B\longrightarrow W_{x,t},\quad t\in D,
\tag{SP4}
\]

where every evaluation includes the actual coefficient transport. The strict volume is \(v_P/|D|\). Hence the complete formulas are

\[
\begin{split}
E&=\sum_x\frac{w_x}{|D|}\sum_{t\in D}
                         \operatorname{Tr}(J_{x,t}^*J_{x,t}),\\
E_v&=\sum_x\frac{w_x}{|D|(T-1)}
       \sum_{t\in D}\sum_{s\in D_v}
                         \operatorname{Tr}(J_{x,t}^*J_{x,ts}).
\end{split}\tag{SP5}
\]

The mixed scalar products use the diamond's unitary fibre identification. Replacing the diamond convention by its inverse merely reindexes the complete s-sum. The total in the second line is real and nonnegative because it is \(\operatorname{Tr}(J^*A_vJ)\).

These formulas do not assume trivial stabilizers. The identity

\[
\sum_{[t]\in\Gamma_x\backslash D}
 \frac{v_P/|D|}{|\Gamma_{x,t}|}\|F_t\|^2
 =\frac{w_x}{|D|}\sum_{t\in D}\|F_t\|^2
\]

is exactly the orbit-stabilizer calculation of TG4; the same argument applies to the transported scalar products. Thus no mass or factor \(T-1\) is suppressed in SP5.

Separating the term \(s=1\), put

\[
C_v=\sum_x\frac{w_x}{|D|(T-1)}
       \sum_t\sum_{s\ne1}
              \operatorname{Tr}(J_{x,t}^*J_{x,ts}).
\]

Then

\[
E_v=\frac{E}{T-1}+C_v,\qquad
E_{\rm ram}\ge\left(1-\frac f{T-1}\right)E-\sum_vC_v.
\tag{SP6}
\]

For example, if a proved uniform estimate on these \textbf{actual} cross kernels were

\[
|\operatorname{Tr}(J_{x,t}^*J_{x,ts})|
 \le\eta_v\|J_{x,t}\|_{\rm HS}\|J_{x,ts}\|_{\rm HS}
 \quad(s\in D_v\setminus\{1\}),
\tag{SP7}
\]

Cauchy--Schwarz and the permutation \(t\mapsto ts\) would give

\[
E_{\rm ram}\ge
 \left[1-\frac f{T-1}-\frac{T-2}{T-1}\sum_v\eta_v\right]E.
\tag{SP8}
\]

SP7 is a sufficient hypothesis, not a bound proved in this finding. The ambient local spherical Hilbert pairing of K052 cannot replace the restricted automorphic kernel in SP7 without an additional theorem.

\subsection{2. The actual denominator ideal for a scalar change}\label{the-actual-denominator-ideal-for-a-scalar-change}

Now fix the favorable point p of LIFT41. Use LIFT44(PR1--7)'s notation

\[
v_0=(w,\pi_k),\quad z=v_0/\bar\pi,\quad p=\gamma_\infty o,
\quad\gamma e=z,
\qquad w_1=2Q.
\tag{SP9}
\]

Let \(S\subseteq\mathcal T\) be nonempty, and change the scalar parameters precisely at the places in S, with each \(s_v\ne1\). Suppose the resulting point is in the same standard maximal finite component, so that a rational relative matrix R as in LIFT44(PR3) exists. If it is in a different component, the corresponding componentwise Poincare cross term is zero; no matrix is needed.

At a changed selected CM prime \(P_v\), R and its inverse have entries of valuation at least \(-1\). This follows directly from

\[
R^{-1}\Lambda_v
       =(1+c_v\pi^{-1}zq_v)\Lambda_v,\qquad c_v=s_v^{-1}-1,
\tag{SP10}
\]

because the displayed matrix and its inverse are integral after multiplication by \(\pi\). The inverse-adjoint component has the same bound at \(\bar P_v\). At unchanged places R and its inverse are integral. In particular, for every entry x of R,

\[
\mathfrak d(x)\mid\prod_{v\in S}P_v\bar P_v,
\qquad \mathrm N(P_v)=\mathrm N(\bar P_v)=T.
\tag{SP11}
\]

The norm equalities use the actual degree-one, unramified split T places. Each absolute residue degree is one. This is a denominator statement about entries of R, and does not assert that z is globally integral.

If \(x\ne0\), the fractional ideal norm, or equivalently the norm of its numerator ideal, gives

\[
|N_{E/\mathbb Q}x|\ge T^{-2|S|},\qquad
\prod_{\tau\in\Phi}|\tau(x)|\ge T^{-|S|}.
\tag{SP12}
\]

Indeed the denominator ideal has norm at most \(T^{2|S|}\), and a nonzero integral numerator ideal has positive integer norm at least one. The second assertion is its square root under the CM pairing of embeddings. Consequently some label satisfies

\[
|\tau(x)|\ge T^{-|S|/f}.
\tag{SP13}
\]

For a single changed place this is \(T^{-1/f}\), substantially stronger than a bound using a denominator T at every absolute prime.

\subsection{3. The actual pole rules out a compact block matrix}\label{the-actual-pole-rules-out-a-compact-block-matrix}

At every selected T place, both the first positive coordinate and the negative coordinate of z are units. The former is

\(2Q/\bar\pi\), and T divides neither 2 nor Q. For the latter,

\[
\pi_k\bar\pi_k=T+W,\qquad W=77Q^4-14Q^2,
\]

and the exact actual integers satisfy \(0<W<T\) and

\(\gcd(W,T)=1\). Thus \(\pi_k/\bar\pi\) is a unit at each selected prime as well.

At a changed place \(c_v\ne0\) modulo the residue prime. Since q and z are primitive and \(qz=0\), SP10 implies

\[
\overline{\pi R^{-1}}\text{ has rank one and image }
                              \mathbb F_T\bar z.
\tag{SP14}
\]

To verify this without choosing a lattice basis implicitly, write

\(R^{-1}=(1+c_v\pi^{-1}zq)C\) with

\(C\in\mathrm{GL}_{a+1}(O_v)\), as follows from equality of the inverse-image lattices. Reduction of its product with \(\pi\) is

\(\bar c_v\bar z\bar q\bar C\). It has exactly the rank and image in SP14.

Suppose R were block diagonal in the original decomposition

\(E^a\perp Ee\). Then \(\overline{\pi R^{-1}}\) would also be block diagonal. A block diagonal matrix of rank one has image entirely in one block: if both blocks were nonzero, their ranks would add to at least two. This contradicts SP14, since \(\bar z\) has nonzero coordinates in both blocks. Therefore R is not block diagonal.

For a unitary matrix, vanishing of all positive entries in its last column would make it block diagonal: its last entry has norm one, and orthogonality to that column then kills the remaining entries of the last row. We have proved that at least one of the entries

\[
                         R_{i,a+1},\qquad 1\le i\le a,
\tag{SP15}
\]

is nonzero. This is the step converting an entry-denominator bound into a point-displacement bound. No compact factor is discarded and no rational compact alignment is asserted.

\subsection{4. Geometric separation, including all integral translates}\label{geometric-separation-including-all-integral-translates}

Let \(u_\tau(R)=d_\tau(o,R_\tau o)\) be the usual ball Cartan distance. The negative unit column identity gives

\[
\sinh^2u_\tau(R)=\sum_{i=1}^{a}|\tau(R_{i,a+1})|^2.
\]

Apply SP13 to one nonzero entry in SP15. At some real label,

\[
u_\tau(R)\ge\operatorname{asinh}(T^{-|S|/f}).
\tag{SP16}
\]

The old point has the same radius \(r_0\) at every label, with

\(\sinh^2r_0=W/T\). The triangle inequality and isometry by R give

\[
u_\tau(R)\le d_\tau(o,p)+d_\tau(p,Rp)+d_\tau(Rp,Ro)
             =2r_0+d_\tau(p,Rp).
\]

Consequently

\[
\boxed{\quad
\max_\tau d_\tau(p,Rp)
 \ge\operatorname{asinh}(T^{-|S|/f})-2r_0.
\quad}\tag{SP17}
\]

If the right side is nonpositive, the inequality is merely vacuous. For the simultaneous all-place change this is consistent with the small semisimple example in LIFT44. The present proof does not rule that example out.

Let \(\Gamma=U(I_a,-1)(F)\cap J^{\max}\), exactly as in K058. For any \(\gamma_0\in\Gamma\), the matrix \(\gamma_0R\) has the same inverse-image lattices as R, because \(\gamma_0\) is integral with integral inverse. The denominator bound and mixed pole image remain identical. Thus the same lower bound holds for

\[
\max_\tau d_\tau(p,\gamma_0Rp)
                      \quad\hbox{for every }\gamma_0\in\Gamma.
\tag{SP18}
\]

In particular SP17 measures separation in the actual maximal arithmetic quotient, with the maximum product distance. It is not merely separation of two convenient representatives. The G reduction of \(\gamma_0R\) may act nontrivially on the finite coefficient space; the geometric estimate neither removes nor evaluates that factor.

\subsection{5. Exact single-place numerical bounds}\label{exact-single-place-numerical-bounds}

For J14, the accepted actual integers satisfy

\[
f=42,\quad 2^{1000}\le T<2^{1001}<2^{24f},
\qquad W/T<2^{-813}.
\]

Hence \(2^{-24}<T^{-1/f}<1/2\), and \(r_0<2^{-406}\). For \(0<x<1/2\), the integral formula for asinh gives

\(\operatorname{asinh}x>3x/4\). Therefore a single changed place has

\[
\operatorname{asinh}(T^{-1/f})-2r_0
 >3\cdot2^{-26}-2^{-405}>2^{-25}.
\tag{SP19}
\]

For J56 the corresponding exact bounds are

\[
f=168,\quad 2^{995}\le T<2^{996}<2^{6f},
\qquad W/T<2^{-704}.
\]

Thus \(2^{-6}<T^{-1/f}<1/2\), \(r_0<2^{-352}\), and

\[
\operatorname{asinh}(T^{-1/f})-2r_0
 >3\cdot2^{-8}-2^{-351}>2^{-7}.
\tag{SP20}
\]

Both bounds are for at least one real label, uniformly over the nonidentity residue scalar at the chosen place and over every integral translate. They make no assertion that the same label works for all changes.

\subsection{6. What this does and does not supply to the selected Gram}\label{what-this-does-and-does-not-supply-to-the-selected-gram}

SP19--20 can be used with a valid bound for the complete selected W-to-W matrix coefficient, with its finite deck and fibre factors retained. In particular the large degree of W cannot be replaced by the depth of one highest vector. This finding itself evaluates no such coefficient or rational Poincare sum.

The exact projection criterion SP3 requires \(E_v\) from SP5. Those sums range over every \(t\in D\), whereas SP17--20 use the fixed favorable point corresponding to \(t=1\). The other base points need not be near the standard origin, and their transported negative columns need not have the same simple unit-coordinate description. Therefore a one-point separation or a one-point selected Gram does not establish SP7 or the inequality

\(\sum_vE_v<E\). The uniform cross-kernel estimate or an exact averaged substitute remains the task.

If that inequality is established, the implication to the prescribed coupled arithmetic image is now the checked DO055 actual-image theorem, with its exact Sad, period, branch and class constants. There is no additional requirement to place an unrelated ambient row in the theta image. This report does not claim the missing inequality, an effective final DP44 sign, or a result on Question 2.

\subsection{7. Inputs and verification}\label{inputs-and-verification}

The source norm and mass conventions are exactly LIFT43(TG1--6), including its source locators in \nolinkurl{ScaledRamifiedLattice.lean} and \nolinkurl{Record.lean}. The packet and conjugate-place lattice identities are LIFT44(PR1--7); the actual numbers are LIFT41(NP1--12) and DO040. K058(PG1--4) fixes the maximal arithmetic group. These accepted findings are cited at their established scopes, not reproved here.

The new claims are SP2--8 and SP11--20, all proved above. The exact certificate checks the commuting-projection inequality and its normalization on finite Fourier models, the mixed rank-one pole obstruction, and the actual integer comparisons used in SP19--20. It does not compute an automorphic cross Gram or run Lean.

\clearpage

\section{AM-R-038 --- bounded integral Weyl spanning by the existing cyclotomic grid}\label{am-r-038-bounded-integral-weyl-spanning-by-the-existing-cyclotomic-grid}

Author: GPT-6 Astra Pro, ALP\_MACHINE /root. 22 November 2026 UTC; effort is not separately exposed.

Status: author whole-proof check PASS, 2026-11-22T08:30:00 UTC. The complete proof, source conventions and complete finite diagnostic script have been read. Independent review is recorded separately.

This supplies an explicit bounded integral K-spanning family for the standard initial Weyl module. It avoids an estimate for the inverse Gram matrix of a prescribed dual basis. The actual character field already contains the needed roots of unity, and the actual initial polynomial degrees are smaller than T. Class-frame clearing and stabilizer sums remain necessary. No nonzero Mram period is inferred.

\subsection{1. The exact standard Weyl lattice and polynomial action}\label{the-exact-standard-weyl-lattice-and-polynomial-action}

Let lambda be a dominant integral weight of GL\_a. Write

\[
 c=\lambda_{a-1},\qquad \nu_i=\lambda_i-c\ge0,
 \qquad D=\sum_{i=0}^{a-1}\nu_i.
 \tag{R38.1}
\]

These are precisely DataCoeff.detShift and natW. Let Delta\_std be the actual standard Weyl lattice in the source polynomial space, generated from the highest polynomial h=hw(nu) by the accepted divided-root construction. Its normalized Fischer metric has \textbar\textbar h\textbar\textbar=1. Let K be the existing characteristic-zero coefficient field with mu\_T contained in it, as already proved from literal CharT.Ok in K059 Section1. Assume D\textless T.

For g in SL\_a(K), use the literal polynomial action

\[
                 \rho(g)F(Z)=F(g^{\mathsf T}Z).
 \tag{R38.2}
\]

It is sqAct on determinant-one matrices; the separate detShift is therefore one. The standard Weyl lattice is preserved by SL\_a(O\_K). This is also a consequence of the accepted perfect standard Weyl/dual-Weyl pairing: the integral polynomial action and its inverse preserve the standard dual-Weyl lattice, and pairBK transports the two inverse-transpose actions. No integrality of an orthogonal projection is required.

Using rho in R38.2 constructs vectors in the fixed polynomial lattice. It does not change the subsequent rhoV/rhoVr choice in the actual adelic covariance. The eventual stabilizer sums use exactly that original covariance, as in K059(IS9,IS18).

\subsection{2. A finite family using only existing coefficients}\label{a-finite-family-using-only-existing-coefficients}

Put m=a(a-1)/2 and choose an ordering of the positions (i,j) with i\textgreater j. For z=(z\_1,\ldots,z\_m) define U(z) to be the lower unitriangular matrix with these entries z\_k. This is the matrix with those direct entries, not a product whose entries are left uncontrolled. Set

\[
            v_\zeta=\rho(U(\zeta))h,
                   \qquad \zeta\in\mu_T^m.
 \tag{R38.3}
\]

Each U(zeta) has determinant one, integral entries and integral inverse. Its entries are roots of unity, zero or one; thus

\[
                    v_\zeta\in\Delta_{\rm std}.
 \tag{R38.4}
\]

There is no division by T in R38.3 or R38.4 and no coefficient-field extension. In particular these input vectors are integral at11 and at T. The family can be redundant and very large.

\subsubsection{Theorem R38.5}\label{theorem-r38.5}

The vectors v\_zeta span the whole Weyl representation over K. At every complex embedding of the coefficient field,

\[
                     \|v_\zeta\|_{\rm nf}\le a^D.
 \tag{R38.5}
\]

For a=1 the grid has one vector and the assertion is immediate. The two actual ranks satisfy a\textgreater1.

\subsection{3. Proof of the span, including the finite grid denominator}\label{proof-of-the-span-including-the-finite-grid-denominator}

The vector-valued polynomial

\[
                         V(z)=\rho(U(z))h
 \tag{R38.6}
\]

has degree at most D in each parameter: every substitution in a degree-D polynomial is linear in the matrix entries. Write V(z)=sum\_alpha v\_alpha z\^{}alpha, with 0\textless=alpha\_k\textless=D\textless T. For a primitive root omega of order T, the finite Fourier identity gives

\[
 v_\alpha=\frac1{T^m}\sum_{j\in\{0,\ldots,T-1\}^m}
            \omega^{-\sum_k j_k\alpha_k}
                         V(\omega^{j_1},\ldots,\omega^{j_m}).
 \tag{R38.7}
\]

Indeed summation in each coordinate is T if the two exponents are equal modulo T and zero otherwise. The degree bound excludes aliasing. Thus the K-span of the finite grid contains all the coefficient vectors v\_alpha and hence every V(z). Conversely every grid value is a combination of the coefficient vectors. This proves equality of the two K-spans. R38.7 is used only to prove a K-span; its T-power denominator is not inserted in the integral vectors. No claim of integral saturation follows from this argument.

The lower-unipotent orbit of h spans the full irreducible characteristic-zero highest-weight module. For precision with the transpose convention in R38.2, upper unitriangular matrices fix h: their transposes perform lower row additions inside every initial minor. Diagonal matrices multiply h by its highest-weight character. The product N\^{}- T\_diag N\^{}+ is the open Gauss cell of GL\_a. A linear functional vanishing on rho(N\^{}-)h therefore vanishes on rho(g)h throughout that open cell, then on all GL\_a by polynomial density. Since h generates the irreducible highest-weight module, the functional is zero. This proves the orbit span and, by R38.7, the finite-grid span. The determinant twist makes no change to this argument on N\^{}- and only a nonzero scalar on the diagonal.

\subsection{4. Proof of the norm bound in the literal Fischer metric}\label{proof-of-the-norm-bound-in-the-literal-fischer-metric}

At every complex embedding, all entries zeta\_k have modulus one. Therefore

\[
 \|U(\zeta)\|_{\rm op}\le\|U(\zeta)\|_{\rm Frob}
       =\sqrt{a(a+1)/2}\le a.
 \tag{R38.8}
\]

The degree-D polynomial Schur representation with the normalized invariant Hermitian metric has operator norm at most \textbar\textbar g\textbar\textbar\_op\^{}D on a matrix g. One exact way to see the needed statement here is the Cartan tensor realization in R037: it is a subrepresentation of tensor\_i wedge\^{}(s\_i)(C\^{}a), with sum\_i s\_i=D. The tensor action has norm at most product\_i \textbar\textbar g\textbar\textbar\_op\^{}(s\_i)=\textbar\textbar g\textbar\textbar\_op\^{}D. Restricting to its invariant Cartan summand does not increase the norm, and R037(R37.8) identifies the normalized source metric with that unit tensor metric. The transpose/conjugate-linear coefficient realization uses the adjoint action, which has the same norm.

Applying this to the unit highest vector and to U(zeta), whose determinant twist equals one, proves R38.5. The same proof applies after every complex embedding of K because a root of unity still has modulus one. No assertion that U(zeta) is a unitary matrix is made; R38.8 is its stated operator bound.

\subsection{5. Verification of the actual degree hypothesis}\label{verification-of-the-actual-degree-hypothesis}

Use the original initial coefficient rank a=42 or168, and the complete K044 allowed interval and word profiles. At a top label, the highest and lowest selected weights are exactly A and -B, with A=2\^{}960 and B=2\^{}955. At a non-top label their spread is (8a-6)ell+1-a, as follows from the K044 u profile and the rho shift. In particular it is less than8a ell, and ell\textless=ell\_+\textless2\^{}485. All intermediate word bits leave these extremes unchanged. Hence every actual initial label satisfies

\[
 0\le D_\tau\le a(\lambda_{0,\tau}-\lambda_{a-1,\tau})
       \le a(A+B)<2^{969}.
 \tag{R38.9}
\]

The fixed actual T has1001 bits on J14 and996 bits on J56, so T\textgreater=2\^{}995 in both cases. Thus D\_tau\textless T, with a large margin. The canonical allowed twist changes the negative-circle parameter, not these selected positive weights. No extra prime is chosen. This is a statement about the initial coefficient degree; it does not assert D\textless T for every other representation in the tower.

Tensor the families R38.3 over every CM label in the actual split coefficient-field realization of the standard tensor lattice. These are independent factors of the scalar-extended product representation; one does not assert that independent label choices are the embeddings of one matrix over the original CM field. Each factor action is integral on its standard split Weyl lattice, so their tensor action is integral on the actual standard coefficient lattice. The resulting finite family \{v\_j\} is integral in the actual standard initial Weyl tensor lattice, spans its K-space and satisfies

\[
               \|v_j\|_{\rm nfAll}\le a^{D_\Sigma},
                  \qquad D_\Sigma=\sum_\tau D_\tau.
 \tag{R38.10}
\]

It uses at most T\^{}(f a(a-1)/2) vectors. This cardinality bound is not proposed as an efficient enumeration algorithm.

\subsection{6. Transport to the actual integral initial Mram family}\label{transport-to-the-actual-integral-initial-mram-family}

Let d\_(0,c) be K059's actual clearing integer with d\_(0,c) Delta\_std subset L\_(0,c). Its construction from any integral basis remains valid for every vector of Delta\_std. Define, using the original twisted stabilizer action A\_(0,c),

\[
 u_{c,j}=\operatorname{single}_c
          \sum_{\gamma\in\Gamma_{0,c}}
                  A_{0,c}(\gamma)d_{0,c}v_j.
 \tag{R38.11}
\]

These vectors are integral and span the full admitted initial class space over K by the same unnormalized averaging argument as IS10. Their exact class metric gives

\[
 \|u_{c,j}\|_{0,\rm cls}
     \le\sqrt{|\Gamma_{0,c}|}\,d_{0,c}a^{D_\Sigma}.
 \tag{R38.12}
\]

Apply the literal integral ramified operator C\_T P\_ram=product\_v((T-1)I-N\_v), C\_T=(T-1)\^{}f.~Then

\[
 x_{c,j}=C_TP_{\rm ram}u_{c,j}\in\mathcal M_{\rm ram,L},
 \qquad
 \|x_{c,j}\|_{0,\rm cls}
 \le C_T\sqrt{|\Gamma_{0,c}|}\,d_{0,c}a^{D_\Sigma}.
 \tag{R38.13}
\]

The family spans actual Mram over K. Combining it with R037's terminal family in the DO055/K059 nonzero extraction retains the same prescribed arithmetic coupled map. The conclusion remains conditional on a nonzero actual moved Mram bottom period. The norm bound is explicit but is not asserted to fit a net DP44 saving.

\subsection{7. Scope and exact source anchors}\label{scope-and-exact-source-anchors}

Class-frame clearing integers and heights, stabilizer/projector costs, actual arithmetic pairing values and the mixed host budget remain. The grid construction does not bound an integral inverse, change the coefficient field, prove Mram nonzero, or make the arithmetic map saturated. Its Fourier extraction has a denominator only in the proof of K-spanning, not in its integral input vectors.

Direct source/interface anchors at the stipulated archive commit:

\begin{itemize}
\tightlist
\item
  \nolinkurl{DataCoeff.lean20–23,49–62:} detShift, natW, delSq, extremal, and the transpose convention in sqAct.
\item
  \nolinkurl{GradedWeylDivisibility/StatementA.lean59–76:} highest polynomial and divided-root Weyl lattice.
\item
  \nolinkurl{CoefModB.lean55–69,113–119} and the accepted perfect standard Weyl/dual-Weyl pairing in FactDelDualNabla.lean75.
\item
  K044(EG1--5), K059 Section1 and IS18--21, R037(R37.4--10), DO055(CI21--23): actual weights, existing mu\_T, class and Mram normalization, source metric and prescribed-image consequence.
\end{itemize}

Accepted source theorems are cited at their stipulated status. No Lean audit or build is part of this argument.

\clearpage

\section{AM-K061 --- the actual two-point projected selected Gram}\label{am-k061-the-actual-two-point-projected-selected-gram}

Author: GPT-6 Astra Pro, ALP\_MACHINE \nolinkurl{/root/kummer_extension}. 22 November 2026 UTC. Inherited reasoning effort is not separately exposed.

\textbf{Author whole-proof check PASS --- 2026-11-22T08:26:27.789895+00:00.}

This finding sums the actual rational selected-kernel series for the two J14 points of R036. It proves a quantitative almost-orthogonality statement for their genuine finite-projector, G-good selected evaluations. In particular an actual G-generic bottom restriction changes under the specified simultaneous T scalar. This only detects a nontrivial joint T-dilation character somewhere. It does not prove that every T component is nontrivial, an Mram period, or a terminal arithmetic-image value.

The large compact representation dimension is divided out before the infinite sum is estimated. The long terms are bounded by a complete L1 sampling argument. A constant pointwise bound is not summed over an infinite arithmetic group.

\subsection{1. Fixed data and exact kernel convention}\label{fixed-data-and-exact-kernel-convention}

Retain all data and metrics of K058(PG1--15). Thus a=f=42 for the final two-point conclusion, G is the product of U(a,1), Gamma is its actual integral maximal arithmetic group, Z=mu28, and z0=28. Let pi be the holomorphic Hilbert representation, V its minimal compact type, and W its entire selected compact component. Put dW=dim W. Its inclusion is the unitary K043 inclusion, with the actual PBW factors retained.

The genuine finite representation is R=Pi-dual. Let Qgood and Rgood be the actual highest-mirabolic partial projectors at the two mixed places and the identity at every full-degree cuspidal place, as in K058(PG3). Put D=dim Rgood and E=Rgood tensor W. This uses the unchanged goodTuple and finite levels. In particular it is not a new finite projector.

For a real matrix g\_p representing a packet point p, set

\[
 B_p(x)=\operatorname{pr}_V\pi(x^{-1}g_p)|_W,\qquad
 b_p(x)=\frac{\|B_p(x)\|_{\rm HS}}{\sqrt{dW}}.
 \tag{CG1}
\]

The periodization T\_p is exactly K058(PG14), with B replaced by B\_p. Let J\_p be its actual good selected evaluation. On the finite-dimensional span of all Poincare sections under consideration,

\[
                       J_p=z_0 C T_p^*,\qquad
 C=\prod_{\tau,i}(r_{\tau i}-1).
 \tag{CG2}
\]

This is K058(PG25--26), with raw R030 Haar, not normalized total volume. The same unfolding applies simultaneously to finitely many translated seeds. Both R036 points have real Cartan parameters below1/200. The principal-G separation proof in LIFT41(NP17--20) therefore applies to each: its bound e\^{}11+1\textless Q is unchanged. The accepted cusp Poincare theorem consequently puts both periodizations in the same actual cuspidal Hilbert space. This argument does not refine the finite level.

The exact cross-kernel is

\[
 J_qJ_p^*=z_0 C\sum_{\gamma\in Z\backslash\Gamma}
 P_E\bigl(R(\gamma)\otimes
       P_W\pi(g_q^{-1}\gamma g_p)|_W\bigr)P_E.
 \tag{CG3}
\]

Here P\_E restricts the finite factor to Rgood; it does not commute an arbitrary gamma through Qgood. Absolute convergence is established in sections2--5 below. Unfold the two periodizations using K058(PG19--21), and replace the remaining gamma by its inverse. The full Gamma sum gives the factor1/z0, while \(\int\langle F_u,F_v\rangle=C^{-1}\langle u,v\rangle\). Combining with CG2 gives precisely CG3. The central characters of R and pi cancel, so each summand is independent of its Z representative.

Define the normalized scalar trace

\[
 \mathcal K(q,p)=\frac{\operatorname{Tr}_E(J_qJ_p^*)}{z_0 C DdW}
   =\sum_{\gamma\in Z\backslash\Gamma}a_\gamma c_\gamma(q,p),
 \tag{CG4}
\] \[
 a_\gamma=\frac{\operatorname{Tr}(Q_{\rm good}R(\gamma))}{D},\qquad
 c_\gamma(q,p)=\frac{\operatorname{Tr}_W
            (P_W\pi(g_q^{-1}\gamma g_p)|_W)}{dW}.
 \tag{CG5}
\]

In particular \textbar a\_gamma\textbar\textless=1 and \textbar c\_gamma\textbar\textless=1. The finite and compact dimensions have both cancelled. They must not be reinserted into the tail estimate below. These are normalized traces, not operator norms of the full automorphic evaluation map.

\subsection{2. Radial moments after all required Lie derivatives}\label{radial-moments-after-all-required-lie-derivatives}

For this section either host a=f=42 or168 is allowed. Write

\[
 d_L=f(a+1)^2<2^{23},\quad \ell=\ell_+,\quad
 t=\ell/2,\quad s_0=a+1,\quad
 \delta=f2^{-\ell/8}.
 \tag{CG6}
\]

Use the actual real matrix basis of R031(SB1), each operator norm at most1. For an ordered word I of length m\textless=d\_L, put

\[
 B_I(x)=\operatorname{pr}_V\pi(x^{-1})d\pi(X_I)|_W,\qquad
 b_I(x)=\|B_I(x)\|_{\rm HS}/\sqrt{dW}.
 \tag{CG7}
\]

Every compact type reached by this word has Pieri exponents m'\_i with

\[
             0\le m'_i\le m_i+2d_L.
 \tag{CG8}
\]

To check the coordinate bound, the complexified Lie algebra as a compact module is contained in the sum of the trivial representation, the standard module, its dual, and standard tensor dual. Tensoring a dominant U(a) module with the standard or dual changes one highest-weight coordinate by one, by the elementary one-box Pieri rule. Thus each Lie derivative changes a coordinate by at most two. Apply this through the word at its actual labels. The fixed holomorphic representation still has the same minimal alpha,beta and the same starting factors r\_i. Its compact decomposition is the multiplicity-one Pieri decomposition used in K043. The shift bound does not assert that the result is a single compact type.

For every compact type in CG8 the exact shifted-metric moment proof of K046(BL8--10) still applies. The shifted inner product is compact invariant, hence orthogonal between different compact types. The moment of a vector in their sum is therefore bounded by the largest component moment, with no factor counting the types or their dimensions.

The actual zero-word constants of K046 satisfy A\_tau\textless969 on J14 and A\_tau\textless1722 on J56. Since r\_i+m\_i\textgreater=2ell and the exact integer inequality 10000 a d\_L\textless ell holds, CG8 gives

\[
 A'_\tau\le A_\tau(1+d_L/\ell)^a
       \le\frac{A_\tau}{1-a d_L/\ell}
       <\begin{cases}970&J14,\\1723&J56.\end{cases}
 \tag{CG9}
\]

The middle bound follows by comparing the binomial coefficients with a\^{}j and the resulting geometric series. Increasing the actual canonical twist or changing an interior zero bit to one only improves these inequalities, as in K046(BL16). Thus all roles, words and allowed twists remain covered.

K049(IC20--21)'s noncommuting Casimir estimate gives \(\|d\pi(X_I)|_W\|\le2^{1000m}\). Its argument also holds with the conservative coordinate allowance 2d\_L in CG8: the bound 8(a+1)(2A+2d\_L+a+1)\textless2\^{}1000 is unchanged. Consequently

\[
 \int b_I^2\le\frac{2^{2000m}}C,
 \qquad
 \int b_I\le2^{1000m}G_0,
 \qquad
 G_0=C^{-1/2}(a!)^{-f/2}2^{11f(a+1)/2}.
 \tag{CG10}
\]

For the second estimate use Cauchy--Schwarz with weight \(\prod_\tau Y_\tau^{a+1}\), whose integral is(a!)\^{}-f, and the joint negative moments through s0. CG9 bounds each moment by2\^{}(11s0). Summing over an orthonormal basis of W and dividing by dW proves exactly the same statements for the displayed Hilbert--Schmidt average.

Take R=9/2 on J14 and R=5 on J56, and put R'=R-1/1000. Then

\[
 \int_{G\setminus C_{R'}}b_I
                      <\frac{2^{1000m}\delta}{C}.
 \tag{CG11}
\]

Here are the cutoff constants. On J14, cosh\^{}2(R')\textgreater1950 and 970/1950\textless1/2. On J56, cosh\^{}2(R')\textgreater4096 and1723/4096\textless1/2. For a rational verification, e\^{}3\textgreater20 and e\^{}(-1/500)\textgreater499/500 give cosh\^{}2(4499/1000)\textgreater8000(499/500)/4=1996. Also e\textgreater5/2 gives cosh\^{}2(4999/1000)\textgreater4990\textgreater4096. Thus outside C\_R' one Y\_tau is less than epsilon, with epsilon=1/1950 or1/4096 respectively. In either case epsilon A'\_tau\textless1/2 and epsilon\^{}-1\textless=2\^{}12.

The weighted Cauchy--Schwarz calculation in K046(BL18--22) now gives, for one such label,

\[
 \int_{Y_\tau<\epsilon}b_I
  \le\frac{2^{1000m}}C\,
     2^{(B_0-t)/2},\qquad
 B_0=960af+(11f+1)(a+1).
 \tag{CG12}
\]

It uses only C\textless2\^{}(960af), the Haar-weight integral just stated, and the joint moment of order t at this label and s0 at the others. The actual inequality 4B0\textless ell makes the last power less than2\^{}(-ell/8). Sum the f possible exceptional labels to obtain CG11. This is an absolute L1 tail estimate, not an iteration of a coefficientwise congruence.

\subsection{3. An explicit L1 sampling bound}\label{an-explicit-l1-sampling-bound}

Put rho=1/(100d\_L\^{}2) and \(\Phi(x)=\prod_{j=1}^{d_L}\exp(x_jX_j)\) on(-rho,rho)\^{}d\_L. The chart, injectivity and Haar computations of R031(SB8--20) give

\[
 \|\Phi\|,\|\Phi^{-1}\|\le e^{d_L\rho},\quad
 \Phi^*d\mu\ge2^{-9d_L}\,dx,
 \quad\sum_I|A_I(x)|\le(5d_L/\rho)^m
 \tag{CG13}
\]

for the expression of any ordered m\textless=d\_L coordinate derivatives in ordered left Lie derivatives. R031 states right derivatives; here one uses the right-trivialized tangent instead of the left-trivialized tangent. Prefix and suffix products have the same operator norm and complex-polydisc bounds, so its proof, including Cauchy's coefficient estimate, gives exactly CG13.

The one-variable fundamental theorem and averaging imply the Banach-valued inequality

\[
 \|h(0)\|\le(2\rho)^{-1}\int_{-\rho}^{\rho}\|h(u)\|du
                      +\int_{-\rho}^{\rho}\|h'(u)\|du.
 \tag{CG14}
\]

Apply it successively in all coordinates. Since2rho\textless1, then use CG13. For any smooth matrix-valued F and any x0 this proves

\[
 \|F(x_0)\|\le A_L\sum_{|I|\le d_L}
          \int_{Vx_0}\|L_I F(x)\|d\mu(x),\qquad
 A_L=2^{9d_L}\rho^{-d_L}(5d_L/\rho)^{d_L},
 \quad V=\Phi(({-\rho},\rho)^{d_L}).
 \tag{CG15}
\]

For clarity, the 2\^{}d\_L mixed derivatives from CG14 cancel the factor 2\^{}-d\_L in(2rho)\^{}-d\_L. The coefficient l1 bound is then bounded above by the sum over all ordered Lie words; no scalar or matrix dimension appears in CG15.

The groups U(a,1), and hence their product, are unimodular. Thus right translation from V to Vx0 preserves the same Haar measure. This is why no point-dependent Jacobian enters CG15 or the sampling sum.

For any g\_q with Cartan parameters below1/200, set Gamma\_q=g\_q\^{}-1 Gamma g\_q. The sets V gamma x, gamma in Gamma\_q, are disjoint for every x. Indeed an overlap would put a nonidentity gamma in V\^{}-1 V. Conjugate it back to Gamma. By CG13 its difference from the identity has operator norm less than16d\_L rho\textless1 at every CM embedding. Each matrix entry of that difference is an algebraic integer. A nonzero one would have nonzero integral norm of absolute value at least1, whereas all its conjugates have modulus below1. This is impossible. The same argument includes the finite central and compact stabilizers; it does not assume that Gamma is neat.

Apply CG15 to B\_emptyset divided by sqrt(dW), sum over these disjoint charts, and use CG10--11 for its left derivatives. Left differentiation of \(\operatorname{pr}_V\pi(x^{-1})|_W\) acts on the domain W by the corresponding signed, possibly reversed Lie word; all such words were included above. There are at most(d\_L+1)\^{}(d\_L+1) words. The explicit inequalities

\[
 A_L(d_L+1)^{d_L+1}2^{1000d_L}<K_L,
                 \qquad K_L=2^{1200d_L}
 \tag{CG16}
\]

follow from d\_L\textless2\^{}23, log2 A\_L\textless140d\_L and log2((d\_L+1)\^{}(d\_L+1))\textless25d\_L. For the actual two d\_L these strict integer comparisons are replayed by the certificate.

Translation back from Gamma\_q gives the following uniform estimates, where full Gamma rather than Gamma/Z can be used to strengthen the upper bounds:

\[
 \sup_x\sum_{\gamma\in Z\backslash\Gamma}b_q(\gamma x)
                   \le K_LG_0,
 \tag{CG17}
\] \[
 \sup_x\sum_{\substack{\gamma\in Z\backslash\Gamma\\
                       \gamma x\notin g_qC_R}}
                 b_q(\gamma x)<\frac{K_L\delta}{C}.
 \tag{CG18}
\]

For CG18 each chart centred outside C\_R lies outside C\_R', because its Cartan displacement is at most d\_L rho\textless1/1000. Thus it is legitimate to use the derivative tail CG11 throughout that chart. No derivative of a discontinuous cutoff is taken. The full-Gamma sampling bound costs no omitted stabilizer factor: all of its disjoint charts are included; passing to Gamma/Z only reduces a nonnegative sum.

\subsection{4. The complete long rational sum}\label{the-complete-long-rational-sum}

Orthogonality for the actual canonical lift gives the convolution bound

\[
 |c_\gamma(q,p)|
       \le C\int_G b_q(x)b_p(\gamma^{-1}x)d\mu(x).
 \tag{CG19}
\]

To prove it, write the normalized trace in CG5 as the sum of dW diagonal matrix coefficients, use the exact integral realization of each such coefficient with constant C, and use Hilbert--Schmidt Cauchy--Schwarz. The division by dW turns the two Hilbert--Schmidt norms into b\_q and b\_p. This is the point at which a spurious factor dim W would invalidate the useful estimate.

Let both centres have Cartan parameters below1/200 and put B=2R+1/100. For a gamma with max\_tau u\_tau(gamma)\textgreater B, the two inclusions x in g\_q C\_R and gamma\^{}-1 x in g\_p C\_R cannot both hold. This follows from subadditivity of the Cartan parameter, equivalently from the operator norm. In the integral in CG19, first put the part with x outside g\_q C\_R against the full sampled sum of b\_p.~On its complement only the sampled tail of b\_p can contribute. More explicitly the sum is at most

\[
 C\left\{
   \left(\sup_x\sum_\gamma b_p(\gamma^{-1}x)\right)
                 \int_{G\setminus g_qC_R}b_q
 + \left(\sup_x\sum_{\gamma:\,\gamma^{-1}x\notin g_pC_R}
                  b_p(\gamma^{-1}x)\right)\int_Gb_q\right\}.
 \tag{CG19a}
\]

Use CG10--11 and CG17--18 in this expression. Tonelli gives

\[
 \sum_{\substack{\gamma\in Z\backslash\Gamma\\
                  \max_\tau u_\tau(\gamma)>B}}
           |c_\gamma(q,p)|
       <2K_LG_0\delta
       \le f\,2^{1+1200d_L+11f(a+1)/2-\ell/8}
       <2^{-2000}.
 \tag{CG20}
\]

The same bound applies after multiplication by \textbar a\_gamma\textbar\textless=1. Its numerical margin is extremely large: the certificate uses the actual integer ell and checks the last exponent inequality directly. Short sets are finite by LIFT42's arithmetic packing estimate. Thus the whole normalized trace series is absolutely convergent. The same convolution majorant, with finite dimension factors restored if needed, proves absolute convergence of every matrix entry and justifies the unfolding in CG3--4. No finite numerical truncation of the automorphic series is being substituted for this proof.

\subsection{5. The actual J14 pair and its Gram}\label{the-actual-j14-pair-and-its-gram}

Now take p and q=p5 from R036, with its exact rational representatives g\_p=gamma\_infinity and g\_q=(R5 gamma)\_infinity. Both have Cartan parameters below1/200. The finite matrix R5 is principal at both G primes, so the finite factors in CG3 are the original R(gamma). There is no extra G-deck factor. R036(R36.12--16) proves that q is in the very same original projective-T scalar orbit as p; it does not use an ambient Hecke double coset in place of the actual pencil.

At R=9/2 and B=901/100, LIFT42(FC13) supplies

\[
 \sum_{\substack{\gamma\in Z\backslash\Gamma,\ \gamma\notin Z\\
                       \max u_\tau(\gamma)\le B}}
                         |a_\gamma|<2^{-1200}.
 \tag{CG21}
\]

This estimate includes the full compact origin stabilizer. In fact its noncentral terms have exactly zero finite partial trace by K058(PG7). Both mixed I\_n projectors are retained; there is no replacement by a full character at those places.

For the identity coset in the cross-kernel, R036(R36.20) and K060(RC14--18) give

\[
                |c_1(q,p)|<2^{-4200}.
 \tag{CG22}
\]

This uses all42 labels, each separated by more than2\^{}-36, and the full selected W compression, not only its highest vector. Compact frame factors in g\_p and g\_q preserve W and do not change that bound. For either diagonal kernel, a1=c1(p,p)=c1(q,q)=1.

Combining CG20--22 proves the actual normalized Gram estimates

\[
 |\mathcal K(p,p)-1|<\eta,\qquad
 |\mathcal K(q,q)-1|<\eta,\qquad
 |\mathcal K(q,p)|<\eta,\qquad \eta=2^{-1199}.
 \tag{CG23}
\]

Equivalently, in the Hilbert space of Hilbert--Schmidt operators from the common cuspidal section space to E, the two actual selected evaluations have diagonal squared norms within eta of z0 C DdW and cross-inner-product of absolute value below eta times that number. Hence for every complex number lambda of modulus1,

\[
 \|J_p-\lambda J_q\|_{\rm HS}^2
                 >2z_0 C DdW(1-2\eta)>0.
 \tag{CG24}
\]

More generally the2-by-2 Hilbert--Schmidt Gram is positive definite, with smallest eigenvalue greater than z0 C DdW(1-2eta). This asserts linear independence of these two evaluation operators. It does not assert surjectivity onto either full selected fibre.

The same proof bounds the cross trace after any unitary change in the common W frame: in CG19 one inserts that unitary into one matrix factor, which does not change its Hilbert--Schmidt norm. Thus CG24 is insensitive to the actual compact frame comparison. It does not permit an arbitrary finite G-deck factor, which could change CG21.

\subsection{6. An actual generic restricted row, and the remaining T condition}\label{an-actual-generic-restricted-row-and-the-remaining-t-condition}

The rank of J\_p-J\_q is at most DdW. CG24 with lambda=1 therefore gives an actual unit cusp section S in the genuine Pi block with

\[
        \|(J_p-J_q)S\|>
                      \sqrt{2z_0 C(1-2\eta)}.
 \tag{CG25}
\]

The two evaluations use the same finite principal-G frame and the same right selected type. Apply the exact Hcal average K058(PG29) to this difference. It supplies one common h in Hcal and one actual positive psi vector in the good finite subspace whose two selected values differ. The finite h is away T, so it commutes with the scalar T action. Consequently the two points remain in the same unchanged scalar orbit. The compact and finite scalarization formulas K058(PG31--31b) apply to this linear combination just as to a single evaluation. In particular there is a Hilbert-unit scalar bottom vector F in the genuine e0 block, the prescribed G-generic type and the actual selected real type with

\[
 |F(ph)-F(qh)|>
             \sqrt{\frac{2z_0 C(1-2\eta)}{d_{\max}}}>0.
 \tag{CG26}
\]

Here dmax is the largest generic Hcal constituent dimension; the sharper factor is multiplied by sqrt(dim V dim Pi/qpsi), exactly as in PG31--32. The raw derivative version restores rho\_all, or the explicitly chosen nonisometric Pieri scalar. No arithmetic unit or integral normalization is inferred from this Hilbert-space estimate.

R036 identifies q with the simultaneous scalar s5 at every T place. Thus CG26 proves that this actual restriction is not invariant under that one joint scalar. On the finite product dilation group its character expansion has at least one nontrivial character component. It does not imply that the component is nontrivial at every place. The all-place operator is product\_v(1-A\_v); nonconstancy only controls the complement of the product-invariant space. An exact all-place projected Gram, or a valid union-bound estimate for all A\_v, remains necessary before the Mram and actual terminal-image implications of K059/DO055 can be invoked.

\subsection{7. Verification and source locators}\label{verification-and-source-locators}

The companion certificate checks the actual derivative-moment margins, both chart/sampling exponents, tail exponent and rational cutoff bounds, the finite trace normalization in noncommuting finite models, and the two-point Gram inequalities. It does not numerically evaluate an automorphic series or a period.

Primary analytic inputs are exactly those already source-audited in K043/K046/K060: Nakahama(2016), Theorem5.1 and action(5.5), printed pp.959--961, and Haagerup and Schlichtkrull(2014), Theorem1.1, printed p.228, with its stated universal constant below12. Finite character input is LIFT42(FC5--13), from Green(1955), Section5 Example(ii), pp.430--431, or its explicit statement in Gorodetsky and Hazan(2017), Theorem2.1, p.5. The accepted Poincare theorem is used through the precisely matched K046/LIFT41 interface. The sampling argument in CG13--20 is elementary and proved here; R031 supplies its already checked actual Haar density and coordinate coefficients.

Reports used in full: K058(PG1--32), K060(RC1--19), R036(R36.1--21), LIFT42(FC1--15). The derivative and chart statements are rechecked at K046(BL1--22), K049(IC20--21) and R031(SB8--24).

\subsubsection{Primary references}\label{primary-references-1}

Nakahama, R. (2016). Norm computation and analytic continuation of vector-valued holomorphic discrete series representations. \emph{Journal of Lie Theory, 26}, 927--990. \url{https://jolt.centre-mersenne.org/item/10.5802/jolt.919.pdf}

Haagerup, U., \& Schlichtkrull, H. (2014). Inequalities for Jacobi polynomials. \emph{The Ramanujan Journal, 33}, 227--246. \url{https://doi.org/10.1007/s11139-013-9472-4}

Gorodetsky, O., \& Hazan, Z. (2017). On certain degenerate Whittaker models for cuspidal representations of GL(kn,Fq). Author primary version, Theorem2.1. \url{https://arxiv.org/abs/1707.07308}

\subsubsection{Author whole-proof check}\label{author-whole-proof-check-3}

The complete saved CG1--26, including CG19a, and the complete certificate were read and reconstructed. The check covered the finite/real inverse order, central quotient Haar, actual PBW norm, compact-type perturbation, dimension-free Hilbert--Schmidt moments, left-chart disjointness, complete L1 sampling and rational tail, actual principal-G Pell frame, and the generic/scalar extraction. The finite certificate passed1,629 assertions at2026-11-22T08:25:14.856735+00:00. Its finite noncommuting group model checks normalization only; the infinite sum is proved in CG13--20. All-place Mram and the actual terminal-image implication remain separate.

\clearpage

\section{AM-DO-056 --- actual MainX normalization, integral generators and the remaining frame cost}\label{am-do-056-actual-mainx-normalization-integral-generators-and-the-remaining-frame-cost}

Author: GPT-6 Astra Pro, ALP\_MACHINE \nolinkurl{/root/definite_oscillator}. 22 November 2026 UTC. Inherited reasoning effort is not separately exposed.

This finding applies the proved mixed-CM kernel, coupled-image and terminal-mass identities to their actual integral lattices. It gives an exact cancellation of the large T powers in the Hilbert-normalized last-two relation, replaces an unspecified Sad norm by its exact value, and gives an explicit bounded integral nonzero-witness consequence of K059 and DO055. It also proves that determinant/covolume information alone cannot bound K059's particular class-frame clearing integers.

The nonzero hypothesis is precisely an actual moved Mram bottom period. It is not replaced by the presently known nonzero selected restriction before all T projections. Once that hypothesis holds, no additional assertion that an ambient nonzero row belongs to the prescribed coupled image is required. These statements concern the constructed J14/J56 mixed family; they do not identify the signed 032 seesaw with the old positive-identity constructor or supply the obstructed old HostRaw.

\subsection{1. Fixed actual spaces, metrics and kernel scalar}\label{fixed-actual-spaces-metrics-and-kernel-scalar}

Use one allowed word and gap in K044. The ranks and fields are those already fixed in K031/K032 and DO040:

\begin{longtable}[]{@{}lrrrrrr@{}}
\toprule\noalign{}
Host & f=a & r=a+1 & t & b=r+t & c=a+2(t+1) & N=c-1 \\
\midrule\noalign{}
\endhead
\bottomrule\noalign{}
\endlastfoot
J14 & 42 & 43 & 1759 & 1802 & 3562 & 3561 \\
J56 & 168 & 169 & 1633 & 1802 & 3436 & 3435 \\
\end{longtable}

Put h=1 or4 respectively. Let A be the unchanged full-Y arithmetic coupled operator of K038(CA27). Let B\_T be the literal terminal brC, and write

\[
 \mathcal B(P,Z)=\operatorname{PtilFF}(A P,Z),\qquad
 P\in\mathcal M_{\rm ram},\quad Z\in\mathcal Z^{\rm good}.
 \tag{AC1}
\]

The source and detector lattices are respectively the literal initial Weyl lattice MramL and the good part of terminal intNab. The latter is not replaced by intDel. Their class norms are

\[
 \|P\|_0^2=\sum_e|\Gamma_{0,e}|^{-1}\|P_e\|^2,
 \qquad
 \|Z\|_T^2=\sum_d|\Gamma_{T,d}|^{-1}\|Z_d\|^2.
 \tag{AC2}
\]

All coefficient norms here are the source normalized Fischer norms, with their determinant characters. Let v0,v1,vT be the original strict finite level volumes, with compact real Haar probability. Let d0 be the initial compact coefficient dimension and dW the terminal detector coefficient dimension. Unit compact test vectors give the exact readings

\[
 \|\operatorname{rd}_0(P,y)\|_{\rm aut}^2
       ={v_0\over d_0}\|P\|_0^2,
 \qquad
 \|E_T(Z,v)\|_{\rm aut}^2
       ={v_T\over d_W}\|Z\|_T^2.
 \tag{AC3}
\]

Here y is a unit coefficient functional and the canonical J identifies v isometrically with its terminal coefficient vector. If J or the chosen vectors are scaled, their displayed squared norms multiply AC3. The first formula is the same compact Schur integration used in K037(CR26--29); the second is DO038(38.14)/K059(IS15).

Separate the actual source unit-vacuum/CM factor from the two finite volumes:

\[
 \mathfrak M_\sigma
  =\det(2Y_J)^{7hb^2}
     C_\theta^{14h(ab+bc)}
       \prod_\tau(\lambda_\tau/cD_\tau)^{D_{0,\tau}+D_{1,\tau}}.
 \tag{AC4}
\]

The subscript records the chosen complex realization; its CM periods and frames are those of that realization. Thus CI15--16 say p01\_unit=v0 v1 M\_sigma. Formula AC4 retains the actual period, its phase, the source creation scale, and the unit-vacuum factor. It assigns none of these complex quantities an unproved finite valuation.

Let L\_C(P,Z;y,v) denote the left side of DO055(CI18), using the same actual real tensor C, original vertical lifts and moved b0. Then

\[
 L_C(P,Z;y,v)
   =s_Gs_Qc_T\Omega_f\,\mathfrak M_\sigma
                v_0v_1v_T a_C(y,v)\,\mathcal B(SP,Z).
 \tag{AC5}
\]

S is the literal Sad. The fixed real factor is

\[
 a_C(y,v)={E_{C,y}(B_T^*Jv)\over\beta_Bd_W},
 \qquad B_TB_T^*=\beta_B I.
 \tag{AC6}
\]

In the exact matched complete-unitary-coevaluation normalization of CI14,

\[
 a_C={z\over d_0\sqrt{H_0H_1}\,b_{\rm Br}d_W},\qquad
 b_{\rm Br}=\prod_\tau u_{T,\tau}c_{{\rm Br},\tau}\sqrt{\kappa_\tau}.
 \tag{AC7}
\]

The recoupling scalar z has its fixed phase and \textbar z\textbar\^{}2=d\_v/d\_h. For one matched simple pair replace z by z/d\_v. PN15 primitive embeddings with the unnormalized selected PBW input additionally multiply AC7 by sqrt(N0 N1 D\_PBW); a primitive integral higher Pieri input has its further q\_lambda. These are the already proved CI14/DO050 conventions, not new choices of a scalar. In particular b\_Br is inverted in this period comparison although the arithmetic branch itself contains the forward factor.

The all-polynomial kernel/covolume identities have not been reduced to one nonzero period: they are K035--K038's actual jet/subD comparisons and the faithful complete class readings used in CI5--CI10. Their exact middle refinement identity is

\[
             (I_1h_Q)(v_1/h_Q)=I_1v_1,
             \qquad (I_1I_2)^{-1}I_1I_2=1.
 \tag{AC8}
\]

The first equality uses the full function after middle integration and the actual stabilizers of the refined classes. The second includes the terminal trace. Neither is an unintegrated tensor identity, a bijection of the redundant Y array, or a license to omit a Q coefficient in the raw kernel.

\subsection{2. Exact norm and away-T lattice behavior of the prescribed Sad}\label{exact-norm-and-away-t-lattice-behavior-of-the-prescribed-sad}

At one selected split T place put q=T and let the source rank be a. The literal local operator in PlacesToAdeles.lean22--28 is dia(epsilon) FFStar. Its two Fourier kernels are defined in \nolinkurl{HasConductorMulShift.lean76–79.} On a lattice L they are the punctured finite Fourier matrices

\[
 F(x,f)=\psi(f(x)),\qquad F^*(f,x)=\psi(-f(x)),
 \quad x\in L/qL\setminus\{0\},\ f\in(L/qL)^*\setminus\{0\}.
 \tag{AC9}
\]

The source pairing conventions identify the two sides with \nolinkurl{withKer L} and \nolinkurl{withLat L}. No probability factor is included. Changing the local uniformizer or the nontrivial residue character only changes the indicated finite Fourier matrix by a permutation or a unitary character dilation; the same calculation applies with the literal choice. Character orthogonality gives

\[
                 F^*F=q^a I-J,
 \tag{AC10}
\]

where J sends a function to its total sum on the punctured fibre. Indeed the (f,g) coefficient is sum\_(x nonzero) psi((g-f)(x)), which is q\^{}a-1 when f=g and -1 otherwise. This is exactly the accepted \nolinkurl{FFStar_FF} in \nolinkurl{ExistsFFNeZeroThree.lean78–101;} its \nolinkurl{isAdj_FF}47--49 and \nolinkurl{tr_eq_zero}114--116 use the same signs.

Literal Mram imposes the unweighted sum zero on each nonzero scalar orbit. Each fibre is a disjoint union of such orbits, so J vanishes on Mram. The Fourier operators preserve these conditions: one may sum their formulas over a scalar orbit and use the same line-sum-zero identity, or use the source commuting diamond law. Thus AC10 is q\^{}a I on that subspace.

There is no omitted arithmetic stabilizer in passing to the global class metric. To see this explicitly, expand a class-pairing sum over all incident decorated pairs (P,Q). For a finite group Gamma acting on a finite incident fibre S, the orbit identity is

\[
 {1\over|\Gamma|}\sum_{s\in S}a_s
       =\sum_{[s]\in\Gamma\backslash S}{a_s\over|\Gamma_s|}
 \tag{AC11}
\]

for invariant paired values a\_s. Apply this first at the P end and then at the Q end of the same correspondence. Reversing an incidence leaves its stabilizer unchanged. Its additive-character coefficient is conjugated, and its actual coefficient transport is unitary because the rational group is compact at every real place. Therefore the reverse Fourier correspondence is the Hilbert adjoint for AC2, even when the vertex stabilizers differ. This is the same finite groupoid mass argument already used in CI7; it does not give equal masses to unequal stabilizers.

The quotient is a finite-dimensional automorphic class space. On its ramified subspace F\^{}\emph{F=q\^{}a I makes F injective and hence surjective, so FF\^{}}=q\^{}a I there as well. The diamond is unitary. Distinct T places commute and each preserves the other ramified conditions. Consequently

\[
 \boxed{\quad S^*S=SS^*=T^{fa}I\quad\hbox{on }\mathcal M_{\rm ram},
     \quad\|SP\|_0=T^{fa/2}\|P\|_0,
     \quad U_S=T^{-fa/2}S\text{ is unitary}.\quad}
 \tag{AC12}
\]

This includes the zero space without an assertion that it is nonzero. The arithmetic reverse S\^{}\# is the explicitly reversed correspondence; at this complex realization it is S*. Its inverse is T\^{}(-fa) S\^{}\#.

At a coefficient-field place w whose residue characteristic is not T, every one of these neighbor lattices agrees locally with the original one, and the finite character coefficients are roots of unity. The original G decorations and their exact coefficient transports are retained. Thus S and S\^{}\# preserve the actual transported Weyl lattice at w. Since T is a unit there, AC12 implies

\[
             S\bigl(\mathrm{MramL}_w\bigr)
                         =\mathrm{MramL}_w\qquad(w\nmid T).
 \tag{AC13}
\]

In particular the actual Sad contributes no denominator at11. The local lattice statement follows from the source neighbor equalities away P and its conjugate \nolinkurl{(DecoFin.lean653–656,736–742),} followed by the functorial coefficient lattice test. It does not require the coordinates of a rational change of representative to be integral in an earlier, untransported frame. No claim that S preserves the T-adic lattice is made: the lattice Lsh in CoefCol.lean96--99 changes at that place. The unitary U\_S in AC12 is used for complex norms, not asserted to be an integral endomorphism.

\subsection{3. The exact T cancellation and the remaining finite volumes}\label{the-exact-t-cancellation-and-the-remaining-finite-volumes}

At T the initial and middle levels are the strict nonzero-covector stabilizers in GL\_a(O) and GL\_b(O). The finite groups act transitively on nonzero covectors, with q\^{}n-1 of them; at every selected T place the residue degree is one. Therefore

\[
 v_{0,T}=(T^a-1)^{-f},\qquad
 v_{1,T}=(T^b-1)^{-f},\qquad v_{T,T}=1.
 \tag{AC14}
\]

The source T scalar and the actual finite basic comparison give CI20:

\[
 c_T\Omega_f=\Gamma_{\rm bas}\Gamma_T(T^b-1)^f
                 \prod_{v\mid T}(\mu_0\xi)(\varpi_v).
 \tag{AC15}
\]

In particular the entire factor (T\^{}b-1)\^{}f cancels v1,T in AC5. This cancellation is a consequence of the literal strict middle level, not of replacing it by a maximal compact. Q refinements have already been accounted for in AC8.

Write v0\^{}o,v1\^{}o for the volumes with their T factors removed; vT has no T factor. Let

\[
 \Xi_\sigma=s_Gs_Q\Gamma_{\rm bas}\Gamma_T
                           \prod_{v\mid T}(\mu_0\xi)(\varpi_v).
 \tag{AC16}
\]

The finite splitting characters and fixed basic phases in this definition are not suppressed. At the given unitary complex realization, only the explicit positive modulus in sG contributes to \textbar Xi\_sigma\textbar:

\[
 |\Xi_\sigma|=\prod_{v\in G}q_v^{-a(a-1)-b/2}.
 \tag{AC17}
\]

No finite valuation of Gamma\_bas or Gamma\_T is inferred from AC17. The splitting values may be algebraic nonunits at their own finite primes even when their displayed complex modulus is one.

For unit compact y,v define the bilinear expression on the coefficient spaces by normalizing the actual scalar automorphic input and detector:

\[
 \mathscr L_C(P',Z)=
 {L_C(U_S^{-1}P',Z;y,v)\over
                  \sqrt{v_0/d_0}\sqrt{v_T/d_W}}.
 \tag{AC18}
\]

Its argument U\_S\^{}(-1)P' has the same class norm as P'. Substitution of AC12--AC16 in AC5 gives the exact identity

\[
 \boxed{\quad
 \mathscr L_C(P',Z)=
  \Xi_\sigma\mathfrak M_\sigma\,
       v_1^o\sqrt{v_0^o v_T}\,
       (1-T^{-a})^{-f/2}\sqrt{d_0d_W}\,a_C(y,v)
                         \mathcal B(P',Z).
 \quad}                                                     \tag{AC19}
\]

For example the T factor is computed without an estimate:

\[
 T^{fa/2}(T^b-1)^f v_{1,T}\sqrt{v_{0,T}}
               =(1-T^{-a})^{-f/2}.
 \tag{AC20}
\]

Thus all large positive or negative powers of T that arise solely from Sad, the two strict-T volumes and the finite basic-vacuum scalar cancel in this Hilbert normalization. Their remaining logarithmic contribution is

\[
 0<\log(1-T^{-a})^{-f/2}\le fT^{-a},
 \qquad (1-T^{-a})^{-f/2}<1+2^{-40000}
 \tag{AC21}
\]

for the actual DO040 primes. To prove the first bound, put x=T\^{}(-a)\textless1/2 and use -log(1-x)\textless=2x. The inequality exp(y)\textless=1+2y for 0\textless=y\textless=1, T\textgreater2\^{}962, a\textgreater=42 and f\textless=168 then proves the second by an exact integer comparison. This bound concerns AC19, not a T-adic integral normalization of the source row or the branch.

For the matched complete coevaluation AC7, the coefficient in AC19 becomes

\[
 \boxed{\quad
 \Xi_\sigma\mathfrak M_\sigma\,
  v_1^o\sqrt{v_0^o v_T}\,(1-T^{-a})^{-f/2}
       {z\over\sqrt{d_0d_W H_0H_1}\,b_{\rm Br}}.
 \quad}                                                     \tag{AC22}
\]

All additional simple-pair, primitive-PBW and higher integral Pieri factors remain exactly those following AC7. The genuine bottom multiplicity period is obtained by dividing AC19 or AC22 by the actual nonzero cI0 cI1, as CI19 requires. Rescaling an IsIo intertwiner changes that multiplicity normalization inversely; it is not a further cancellation of a fixed test.

The other finite factors are explicit. Put u=a(a-1)/2 and G\_n(q)=\textbar GL\_n(F\_q)\textbar. At a G place the actual strict initial, sharp middle and terminal levels give

\[
 v_{0,G}={q^u\over G_a(q)},\quad
 v_{1,G}={q^{u+rt}G_t(q)\over G_b(q)},\quad
 v_{T,G}={q^{u+2rt+t^2}G_t(q)^2\over G_N(q)}.
 \tag{AC23}
\]

The middle group is LIFT22(CG4), and the terminal group is DO038(38.19a). Multiplying the positive sG modulus by v1,G sqrt(v0,G vT,G) gives exactly

\[
 \mathcal V_G(q)=
 q^{2rt+(t^2-b)/2}
       {G_t(q)^2\over G_b(q)\sqrt{G_a(q)G_N(q)}}.
 \tag{AC24}
\]

At a Q place the initial group is maximal, and the two actual groups are the (r,t) and (r,2t) parabolics. If I1 and I2 are those residue indices, their norm factor is I1\^{}(-1) I2\^{}(-1/2). Thus cancellation of the orbit-trace indices in AC8 does not remove these independently prescribed group-volume contributions to AC19. All finite class frames, passive chi\_Q(bar pi)\^{}(-1), exact subgroup characters and the moved b0 remain in that equation.

\subsection{4. Bounded integral witnesses without an integral vertical inverse}\label{bounded-integral-witnesses-without-an-integral-vertical-inverse}

Assume there is a nonzero actual moved Mram bottom period of DO055(CI3). The proved CI23 implies that AC1 is not the zero bilinear form on Mram times the genuine-good terminal coefficient space. This statement already concerns the image of the fixed A.

Use K059's integral spanning families. Let d0,e and dT,d be its actual local clearing integers for initial and terminal classes. Let B0,i be the exact norms of its initial standard Weyl basis. Root R037, independently whole-checked in this run, proves the sharp bound B\_T,i\textless=1 for every actual standard terminal dual-Weyl generator. Put dG=\textbar Gcal\textbar{} and C\_T=(T-1)\^{}f.~Then there are integral vectors x\_ei and z\_dj spanning the two spaces over the fixed value field and satisfying

\[
 \|x_{e,i}\|_0\le C_T\sqrt{|\Gamma_{0,e}|}\,d_{0,e}B_{0,i},
 \qquad
 \|z_{d,j}\|_T\le d_G\sqrt{|\Gamma_{T,d}|}\,d_{T,d}.
 \tag{AC25}
\]

The x are the unnormalized product of the unweighted augmentation projectors, and the z are dG times the genuine good projector applied to stabilizer sums. Neither family is asserted to be an integral basis or saturated.

There exist indices with

\[
            \mathcal B(x_{e,i},z_{d,j})\ne0.                 \tag{AC26}
\]

Indeed, if every such value were zero, bilinearity and the two field-spanning properties would make AC1 zero, contrary to CI23. Both arguments of AC26 belong to the actual integral lattices. This proves bounded integral existence without clearing an analytic preimage, without making an integral inverse of D, and without paying to apply S to an integral generator. One applies the spanning argument directly to the nonzero form B\_A, not unnecessarily to B\_A composed with S. The actual S and its scalar remain essential to the proof of CI23 and to AC19; they are not erased from the analytic identity.

For a quantitative value comparison the generator sizes in AC25 give the actual Petersson sizes after multiplication by sqrt(v0/d0) and sqrt(vT/dW), respectively. Therefore the mass vT does not disappear from detector costs: it contributes a square root to the detector norm but a full factor to the unnormalized period in AC5. Equivalently the coefficient of a detector normalized by its Petersson norm is sqrt(vT dW) aC before the other factors, exactly as AC19 records.

There is an explicit weight-independent bound for the finite factors in AC25. Initial stabilizers are trivial by K042. For a terminal stabilizer Gamma choose the first actual G prime q=g0. Reduction embeds Gamma into the finite subgroup in AC23. Here is the complete argument for injectivity. A finite element in the principal reduction kernel can have only q-power order, by expanding (1+varpi\^{}s B)\^{}l for a prime l!=q. If the finite kernel were nontrivial, Cauchy's theorem would give an element of order q. On the underlying rational space of dimension 2fN, such a nonidentity rational matrix has a primitive q-th root eigenvalue and hence characteristic-polynomial degree at least q-1. The actual q-1\textgreater2fN excludes this. Thus

\[
 |\Gamma_{T,d}|\le q^{u+2rt+t^2}G_t(q)^2
                       <q^{u+2rt+3t^2}.                    \tag{AC27}
\]

This uses reduction solely as a bound; it does not add a level or change a class. K059 supplies dG\textless(g0 g1)\^{}(f(a+1)\^{}2). The certified bit lengths of the unchanged primes now give

\begin{longtable}[]{@{}
  >{\raggedright\arraybackslash}p{(\columnwidth - 8\tabcolsep) * \real{0.1579}}
  >{\raggedleft\arraybackslash}p{(\columnwidth - 8\tabcolsep) * \real{0.2105}}
  >{\raggedleft\arraybackslash}p{(\columnwidth - 8\tabcolsep) * \real{0.2105}}
  >{\raggedleft\arraybackslash}p{(\columnwidth - 8\tabcolsep) * \real{0.2105}}
  >{\raggedleft\arraybackslash}p{(\columnwidth - 8\tabcolsep) * \real{0.2105}}@{}}
\toprule\noalign{}
\begin{minipage}[b]{\linewidth}\raggedright
Host
\end{minipage} & \begin{minipage}[b]{\linewidth}\raggedleft
log2 C\_T upper bound
\end{minipage} & \begin{minipage}[b]{\linewidth}\raggedleft
log2 dG upper bound
\end{minipage} & \begin{minipage}[b]{\linewidth}\raggedleft
log2 sqrt(Gamma\_T) upper bound
\end{minipage} & \begin{minipage}[b]{\linewidth}\raggedleft
sum upper bound
\end{minipage} \\
\midrule\noalign{}
\endhead
\bottomrule\noalign{}
\endlastfoot
J14 & 42042 & 3494610 & 108495347 & 112031999 \\
J56 & 167328 & 340675608 & 149905858 & 490748794 \\
\end{longtable}

The inequalities are strict where bit lengths were used. Consequently, with maxima over actual classes and initial basis indices,

\[
 \|x_{e,i}\|_0\,\|z_{d,j}\|_T
    <2^{K_{\rm fin}}\,d_{0,\max}d_{T,\max}B_{0,\max},
 \quad K_{\rm fin}=112031999\text{ or }490748794<2^{29}.
 \tag{AC28}
\]

The table bounds the construction's finite projection/stabilizer costs; it does not bound the remaining three factors. In particular R037 does not give B0\textless=1 for the dual initial Weyl basis. These are not bounds on the value field's degree or discriminant. The literal nontrivial CharT already forces the field to contain the T-th roots of unity, as checked in K059; no further field extension is needed for these particular projector vectors, but the cost of that actual field is not erased by this observation.

At11, d0,e and dT,d are units by the actual standard11 representatives. Also 11 does not divide T-1: DO40.4 writes T-1=4D(1+11D/4) and its certified factorization has 11 not dividing D. Thus C\_T is an11-unit. The sufficient dG factor has exactly the K059 valuations336 and21336; no optimality is claimed. These are costs of one spanning construction, not new compulsory divisibility of the original operator.

\subsection{5. Terminal covolumes and scalar normalization}\label{terminal-covolumes-and-scalar-normalization}

Let nT denote the dimension of the full admitted genuine-good terminal class coefficient space. This is not the compact coefficient dimension dW. On any basis of that space, let G\_cls be its Hermitian Gram matrix for AC2. For a fixed unit compact vector v, AC3 gives

\[
 G_{\rm aut}={v_T\over d_W}G_{\rm cls},\qquad
 \sqrt{\det G_{\rm aut}}
       =\left({v_T\over d_W}\right)^{n_T/2}
                                      \sqrt{\det G_{\rm cls}}.
 \tag{AC29}
\]

The square root here is the determinant norm at one complex embedding; a real Euclidean covolume in the associated complex vector space is its square. Global Arakelov conventions sum the logarithms over all embeddings. Formula AC29 is proved by scaling the Gram matrix and makes the dimension exponent explicit. K059's rectangular, possibly dependent generator list is not such a basis. Its norms give Hadamard bounds after an independent subset is chosen, not an equality with the covolume of the saturated lattice.

For the algebraic terminal matrix the exact detector ideal is already R002(R2--R9), applied to B\_T A P and the unchanged good lattice. Thus a complex value normalized to one requires the actual image ideal, not the norm bound AC25. At a finite place w the least denominator is the positive part of that ideal's valuation. At infinity it is the inverse dual norm. These are simultaneous costs on the actual restricted image.

The same scalar cancellation applies to kernel determinant lines. Apply R002 to the top exterior power of a fixed rank-s algebraic operator, or simply take determinants of its nonzero Smith block. Multiplying that operator by c in K\^{}times shifts every nonzero elementary valuation by ord\_w(c), hence shifts its determinant valuation by s ord\_w(c). It shifts the logarithm of the corresponding archimedean determinant norm by s log\textbar sigma(c)\textbar. The product formula cancels their normalized global sums exactly. This applies in particular to the rational finite level volumes and to the known algebraic branch scalar.

Root R024 already computes the latter without a bound:

\[
 B_T=s_B\widehat B_T,\quad s_B=u_T\prod_\tau c_{{\rm Br},\tau},
 \quad
 {1\over[K:\mathbb Q]}\sum_w\operatorname{ord}_w(s_B)\log Nw
   ={1\over[K:\mathbb Q]}\sum_\sigma\log|\sigma(s_B)|.
 \tag{AC30}
\]

Its common value is log\textbar product cBr\textbar-(sum nBr/2)log T. The scalar branch gain therefore cannot be counted once as divisibility and then omitted from the archimedean normalization. AC7/AC22 place its inverse in exactly the needed period comparison. The non-scalar hook/saturation effects and the actual theta-image restriction remain the separate operators already described in R020 and LIFT34.

None of AC29--30 uses a product formula for an unspecified CM period or an unnamed complex phase. The algebraic coefficient map and the analytic normalization are compared by AC5; the finite valuation of that algebraic map is computed in its actual coefficient field.

\subsection{6. The precise global frame-height issue}\label{the-precise-global-frame-height-issue}

K059(IS4,IS6a) gives a finite formula for each clearing integer. If g\_(c,tau,w) is the actual local frame matrix, D\_tau is its polynomial degree and m\_tau its determinant exponent, define

\[
 A_{c,\tau,l}=\max_{w\mid l}
 {\max(0,-\min_{ij}\operatorname{ord}_w g_{ij})\over e_w},
 \quad
 B_{c,\tau,l}=\max_{w\mid l}
 { |\operatorname{ord}_w\det g|\over e_w}.
 \tag{AC31}
\]

Use the actual inverse/dual frame required by the relevant Weyl or dual-Weyl membership test, rather than silently substituting g for its inverse. Let S\_c be the finite rational-prime support of these frames. The ceiling formula proves the sufficient bound

\[
 \log d_c\le
   \sum_\tau D_\tau\sum_{l\in S_c}A_{c,\tau,l}\log l
 + \sum_\tau |m_\tau|\sum_{l\in S_c}B_{c,\tau,l}\log l
 + \sum_{l\in S_c}\log l.
 \tag{AC32}
\]

Indeed max of a sum is at most the sum of the maxima, the signed determinant term is at most its absolute value, and rounding a nonnegative number upward costs less than one at each prime. This proves AC32 directly from IS6a. For fixed underlying classes and levels the constants in the prime sums are finite and independent of polynomial degree. It is therefore a linear-in-weight bound once those actual frame heights are bounded. It is not yet a numerical bound by the source host budget.

There is a concrete reason that the determinant and the standard11 pin do not supply such a bound. Let m be a positive even multiple of11, set d=m\^{}2+1 and define the rational unitary matrix

\[
 U_m={1\over d}
 \begin{pmatrix}m^2-1&-2m\\2m&m^2-1\end{pmatrix}.
 \tag{AC33}
\]

Then U\_m\^{}*U\_m=I, det U\_m=1, and U\_m is integral invertible at11. These facts hold over either actual CM field because all its entries are rational. The lattice L\_m=U\_m O\_E\^{}2 has exactly the same covolume and every determinant valuation as the standard lattice. It is rationally unitarily isometric to it.

For the degree-D standard symmetric-power polynomial lattice transported from L\_m, the membership matrix on the fixed standard monomial basis is Sym\^{}D(U\_m\^{}(-1)). The sufficient clearing integer of IS4 is exactly

\[
                              d_c=d^D.                    \tag{AC34}
\]

Every entry has denominator dividing d\^{}D. Conversely its pure-power diagonal coefficient is (m\^{}2-1)\^{}D/d\^{}D. Since m is even, d is odd and gcd(m\^{}2-1,d)=gcd(2,d)=1. No rational-prime divisor of d cancels in that entry. Thus IS4 needs exactly its full D-th power at every prime dividing d, proving AC34. Also d is prime to11. As m grows, log d\_c=D log(m\^{}2+1) is unbounded with field, rank, all determinant/covolume data and the standard11 condition fixed.

This is a counterexample to an inference about the particular fixed-standard-basis clearing method, not a lower bound on optimal integral generators. In this example an adapted U\_m-transported standard polynomial basis has the same norms as the original basis (at most one in the normalized Fischer metric, with the highest vector of norm one) and no such clearing cost. It proves precisely that a quantitative reduction or adapted-basis construction is additional data. Bounding underlying Hermitian rank-N frames before applying the Schur functor is a possible positive route; it is not supplied by the determinant of the enormous coefficient representation.

\subsection{7. What is available for the actual gain/error comparison}\label{what-is-available-for-the-actual-gainerror-comparison}

Let X\_sigma and Z\_sigma be the proved bounds in AC25 or AC28 at every complex realization, and let

\[
 L_\sigma=\|P_{\rm good}B_TA\|_{0\to T,\sigma}
 \tag{AC35}
\]

with the actual class metrics. The Riesz realization of the class pairing gives, for every pair of generators,

\[
                   |\sigma(\mathcal B(x,z))|
                       \le L_\sigma X_\sigma Z_\sigma.
 \tag{AC36}
\]

The projection in AC35 is the genuine-good orthogonal projection; the bilinear/contragredient convention is the one in DO038, not a change of the arithmetic pairing. If AC26 supplies a nonzero algebraic value b, the product formula gives the exact necessary inequality

\[
 {1\over[K:\mathbb Q]}\sum_w\operatorname{ord}_w(b)\log Nw
       \le {1\over[K:\mathbb Q]}
                  \sum_\sigma\log(L_\sigma X_\sigma Z_\sigma).
 \tag{AC37}
\]

Any proved local lower bound, including K041/LIFT34/R024 at11, may be inserted on the left only together with lower bounds at the other primes. Equivalently subtract a proved total negative-valuation allowance there. One cannot infer a global saving from the11 valuation alone. On the right AC19--22 give the exact conversion from an analytic period or inverse-lift norm to the coefficient normalization, including all CM and class factors. The finite volume/Sad T cancellation is established, while the frame heights AC32, initial Weyl norms and actual operator height are separate quantitative inputs.

The original source budget asks for more than finite bounded generators: \nolinkurl{GlobalTheoremInputs.LinkDIV_at}77--89 requires the actual map height and compares it with MAIN, the top-weight logarithmic allowance and the B0/Nc terms. \nolinkurl{DataParams₂.lean}203--234 fixes the local weights, gains, errors and MAIN; its definitions do not assign the free frame constants in AC32 to one of those allowances. The Host/Scale wrapper obstruction already isolated by R023 is unchanged. Thus the present exact cancellations and bounded integral witness statement remove specific normalization and inverse-image charges without proving the final DP44 inequality or replacing its full family hypotheses.

\subsection{8. Source locators and verification record}\label{source-locators-and-verification-record}

The pinned source commit is 7cbf7a0c32fca1badf05931879a37c3f8a2b03b4. Accepted definitions and theorems are used at their stipulated status; no Lean build or audit is undertaken.

\begin{itemize}
\tightlist
\item
  \nolinkurl{OnFramesRdApply.lean108–121,143–150:} all-polynomial Core, class-weighted Kval, and terminal top/tail equation.
\item
  LevChiEqOne.lean73--77 and PinColT.cst: the independently specified arithmetic jet and its scalar pin.
\item
  \nolinkurl{ValLastX.lean148–150,228–254:} exact cI, LastX and MainX conventions.
\item
  \nolinkurl{HasConductorMulShift.lean76–79,125–138:} actual FF/FFStar and punctured fibre formulas.
\item
  \nolinkurl{ExistsFFNeZeroThree.lean47–49,61–65,78–116:} adjoint, diamond covariance, Fourier composition and ramified trace zero.
\item
  PlacesToAdeles.lean22--28: actual formAt, sOp and diaAt.
\item
  \nolinkurl{ScaledRamifiedLattice.lean79–115,141–173:} integral neighbor coefficients, CharT.Ok and unweighted Mram.
\item
  \nolinkurl{DecoFin.lean653–656,736–742,1143–1190:} neighbor lattices and the actual Lsh convention.
\item
  CoefCol.lean59--65,96--99; StatementD.lean131--141: actual coefficient lattices and determinant-twisted frame membership.
\item
  Settled checked inputs: DO038, DO040, DO050, DO054--055; K035--K039, K042/K044/K059; LIFT22/LIFT33/LIFT34/LIFT39; root R002/R020/R024/R032/R037. Their scopes are kept as stated above.
\end{itemize}

The companion exact certificate checks finite Fourier and line-sum-zero matrices, nontrivial stabilizer-weighted Fourier quotients, every scalar and volume identity used here, fixed-host numerical constants, the rational unitary/frame counterexample, and detector determinant-norm scaling. It does not evaluate the actual automorphic periods or the enormous global theta matrix. The author whole-proof check is recorded only after the complete final text and certificate have been read.

\subsubsection{Completed author verification}\label{completed-author-verification-1}

Author whole-proof check: \textbf{PASS}, 2026-11-22T08:30:23.400106+00:00. The complete AC1--AC37 proof, all eight sections and the entire companion certificate were read. The check included the literal FF/FFStar signs and diamond convention, finite-dimensional ramified adjoint, inverse-stabilizer groupoid masses, actual away-T coefficient lattices, all source class volumes, CM/IsIo/branch factors, conditional integral witness reduction, exact terminal finite subgroup bound, product-formula scope and the determinant-free rational frame example. No actual Mram nonzero value, effective class-frame reduction, saturated generator basis or DP44 budget completion is asserted.

The exact certificate \textbf{passed 142,222 assertions} at 2026-11-22T08:28:06.343439+00:00. Its finite Fourier checks use exact cyclotomic arithmetic; the scalar, volume, frame and host bounds use exact rational/integer arithmetic. These are diagnostics supporting the separately written general proofs. They do not compute an automorphic period or execute Lean.

Final author whole-reread: \textbf{PASS}, 2026-11-22T08:32:51.497903+00:00, after the root review corrected the description of the adapted symmetric-power basis: its normalized Fischer norms are at most one, not all equal to one. The entire report was reread. No equation or certificate changed.

\clearpage

\section{AM-R-039 --- the exact highest-weight cost of the actual T-covector kernel}\label{am-r-039-the-exact-highest-weight-cost-of-the-actual-t-covector-kernel}

Author: GPT-6 Astra Pro, ALP\_MACHINE /root. 22 November 2026 UTC; effort is not separately exposed.

Status: author whole-proof check PASS, 2026-11-22T08:49:48 UTC. The whole proof, actual source conventions and complete diagnostic script have been read. Independent review is recorded separately.

The initial coefficient lattice is formed from Lsh, not merely from the underlying unimodular Hermitian lattice N. Its polynomial frame lattice is the conjugate of the covector kernel L. This finding gives the exact scalar inclusion for that modification. For a single cyclic elementary divisor the exponent is the largest actual highest weight, including its signed determinant shift. It is not the total polynomial degree. The statement concerns the full coefficient lattice before the Mram projector and before the theta image; no optimal image-content claim or net lift saving follows.

\subsection{1. Local covariant lattices and the precise action}\label{local-covariant-lattices-and-the-precise-action}

Let R be a characteristic-zero discrete valuation ring with fraction field K, uniformizer varpi and normalized valuation v. Fix rank a\textgreater0 and a dominant integral weight lambda=(lambda\_0,\ldots,lambda\_(a-1)). Put

\[
 c=\lambda_{a-1},\qquad \nu_i=\lambda_i-c\ge0,
 \quad D=\sum_i\nu_i,\quad \nu_{a-1}=0.
 \tag{R39.1}
\]

Write W\_lambda for either the source standard Weyl lattice Delta or the standard dual-Weyl lattice Nabla, with its determinant twist c. The actual action on its fixed polynomial realization is

\[
 \rho_\lambda(g)F(Z)=\det(g)^c F(g^{\mathsf T}Z).
 \tag{R39.2}
\]

For a free lattice M with basis matrix B, its covariant coefficient lattice is rho\_lambda(B)W\_lambda(R\^{}a). This agrees with the source transLat membership test rho\_lambda(B\^{}(-1))F integral. An integral change of basis preserves both standard lattices, so the definition is independent of B. These statements use the accepted source NabDelEquiv and NabDelStable interfaces, not an integral orthogonal Cartan projection.

Suppose M' is a sublattice of M and, in some R-basis of M,

\[
 M'=\varpi^e R e_1\oplus\cdots\oplus\varpi^e R e_r
          \oplus R e_{r+1}\oplus\cdots\oplus R e_a,
 \quad e\ge0,\quad 0\le r\le a.
 \tag{R39.3}
\]

Define the two signed partial weight sums

\[
 H_r(\lambda)=\sum_{i=0}^{r-1}\lambda_i,
 \qquad L_r(\lambda)=\sum_{i=a-r}^{a-1}\lambda_i.
 \tag{R39.4}
\]

\subsubsection{Theorem R39.5}\label{theorem-r39.5}

For either lattice functor,

\[
 \boxed{\quad
 \varpi^{eH_r(\lambda)}W_\lambda(M)
       \subseteq W_\lambda(M')
       \subseteq\varpi^{eL_r(\lambda)}W_\lambda(M).
 \quad} \tag{R39.5}
\]

For e\textgreater0 both exponents are sharp as scalar inclusions of the full coefficient lattices. In particular, when M/M' is cyclic of length e,

\[
 \varpi^{e\lambda_0}W_\lambda(M)
       \subseteq W_\lambda(M')
       \subseteq\varpi^{e\lambda_{a-1}}W_\lambda(M).
 \tag{R39.6}
\]

Negative powers in these formulas denote fractional lattices. They must not be replaced by their positive parts unless a separate rational-integer clearing convention is being imposed.

\subsection{2. Proof using the literal integral weight decomposition}\label{proof-using-the-literal-integral-weight-decomposition}

The accepted semistandard basis of Nabla consists of products of initial-column minors. Their column lengths s\_j have shape nu, so the number of columns of length at least i+1 is nu\_i. Every such minor has strictly increasing row indices. For any fixed set I of r rows, each column contains at most min(r,s\_j) indices in I. Consequently every row weight mu occurring in the polynomial module satisfies

\[
 \sum_{i\in I}\mu_i\le\sum_j\min(r,s_j)
              =\sum_{i=0}^{r-1}\nu_i.
 \tag{R39.7}
\]

Applying this upper bound to the complementary rows and using total degree D gives the lower bound

\[
 \sum_{i=a-r}^{a-1}\nu_i
       \le\sum_{i\in I}\mu_i.
 \tag{R39.8}
\]

For Nabla the source IsWeightBasis gives an integral direct sum of these row-weight submodules. For Delta, each divided-root word applied to the highest polynomial is either zero or row-homogeneous. Its integral defining span is therefore also the direct sum of its row-weight submodules. The accepted Delta subset Nabla inclusion places their weights in the same bounds R39.7--8. No rational idempotent is used to extract the integral weight summands.

On the summand of weight mu, the diagonal matrix in R39.3 acts by

\[
 \varpi^{e(rc+\sum_{i=0}^{r-1}\mu_i)}.
 \tag{R39.9}
\]

By R39.7--8 its exponent lies between e L\_r(lambda) and e H\_r(lambda). Applying this separately to every integral weight summand proves both inclusions R39.5. The arbitrary basis B of M then transports them to the original lattices.

The highest polynomial has row weight nu and a coefficient equal to one. Its R-line is primitive in both lattices: it belongs to Delta by definition and to Nabla by the accepted highest-polynomial identity; division by varpi would destroy its integral coefficient. It realizes the upper exponent on the first r rows. A permutation of the rows realizes the lower exponent, and the corresponding polynomial remains primitive because that permutation and its inverse preserve the integral lattice. Thus neither scalar exponent can be improved for the whole lattice. This also covers r=0,a, where the two exponents coincide. The case e=0 is equality.

This is a short derivation of the explicit lattice form from the accepted weight basis and divided-root definitions; it does not reprove those representation-theoretic source theorems.

\subsection{3. Global cyclic kernels and the split coefficient tensor}\label{global-cyclic-kernels-and-the-split-coefficient-tensor}

Let E be the actual CM field and let K be the existing coefficient field, with its fixed embeddings tau:E into K. Let N be an integral unimodular Hermitian O\_E lattice in the standard positive space. Let pi generate the product of the selected primes over the rational prime T. The actual choices satisfy

\[
 \pi\bar\pi=T,
 \qquad (\pi)=\prod_{P\in\mathcal T}P,
 \qquad \mathrm N(P)=T.
 \tag{R39.10}
\]

Each selected covector is nonzero. Its simultaneous kernel L therefore satisfies pi N subset L subset N, with N/L locally a free rank-one module over O\_E/P at each selected prime, and zero at all other primes. Equivalently each changed elementary-divisor list is (1,0,\ldots,0). The source mem\_locP\_L and the covector surjectivity give exactly this local description; no freeness of N over O\_E globally is assumed.

Apply R39.6 after each embedding and localize at every finite place of K. Ramification in K simply multiplies v(pi) and does not change the formula. For lambda=(lambda\_tau), put

\[
 \beta_\lambda=\prod_\tau\tau(\bar\pi)^{\lambda_{0,\tau}},
 \qquad \Lambda_0=\sum_\tau\lambda_{0,\tau}.
 \tag{R39.11}
\]

For the covariant split tensor lattice on the conjugate kernel,

\[
 \boxed{\quad
 \beta_\lambda W_\lambda(\bar N)
                \subseteq W_\lambda(\bar L).
 \quad} \tag{R39.12}
\]

This follows factor by factor from R39.6 and then by tensor product. The local-to-global equality test for finitely generated lattices gives the displayed global inclusion. The same argument proves that the fractional ideal of scalars b satisfying b W\_lambda(bar N) subset W\_lambda(bar L) is exactly (beta\_lambda): at each coefficient place the independent extremal weight in every split factor attains the sum of the maximal local exponents. Independent choices concern the scalar-extended tensor representation and do not assert a common matrix over E.

\subsection{4. The literal initial coefficient lattice}\label{the-literal-initial-coefficient-lattice}

The source colVal is roundPk on Lsh. In standard Hermitian coordinates Lsh is the Hermitian dual L\^{}\#, while the coefficient construction delAtr uses the bilinear frame dual. Since

\[
             \operatorname{dualFr}(L^\#)=\bar L,
 \tag{R39.13}
\]

the polynomial lattice in actual colVal.Del is precisely the covariant Weyl lattice on bar L, including its determinant shifts. To verify R39.13, bilinear dual is conjugate Hermitian dual for the identity Gram matrix, and the double dual of a full lattice over O\_E is itself. The source RoundLat expresses the same identity after tau base change. Thus R39.12 applies to the actual initial lattice; it is not being applied to N in place of Lsh. The terminal rowVal is squarePk on N and has no analogous initial T-kernel factor.

All conjugates of pi have absolute value sqrt(T). Hence at every complex embedding sigma of K,

\[
 |\sigma(\beta_\lambda)|=T^{\Lambda_0/2},\qquad
 \|\beta_\lambda x\|_{\rm nfAll}
           =T^{\Lambda_0/2}\|x\|_{\rm nfAll}.
 \tag{R39.14}
\]

Moreover beta\_lambda is a unit at every coefficient place over11. For the actual initial weights all lambda\_0,tau are positive, so it is integral as well. If a previously proved integral K-spanning family for W\_lambda(bar N) has norm bound B\_N, multiplication by this one common scalar gives a K-spanning family in the actual coefficient lattice with norm bound T\^{}(Lambda\_0/2) B\_N. Integral stabilizer sums and (T-1)\^{}f P\_ram can then be applied exactly as in K059/R038; the input spanning assertion does not require saturation.

The scalar is the same for all decorated classes at fixed weight, so its use does not introduce a new class-dependent height. It is also common to all K044 interior words at fixed ell, by the next explicit formula. All original covariance and character data remain.

\subsection{5. Actual weight exponent and quantitative scope}\label{actual-weight-exponent-and-quantitative-scope}

Let n\_top be the number of top CM labels. The K044/R022 endpoints and the exact rho shift give

\[
 \lambda_{0,\tau}=\begin{cases}
 A,&\tau\text{ top},\\
 (8a-4)\ell+1-a/2,&\tau\text{ not top},
 \end{cases}
 \quad A=2^{960},
\] \[
 \boxed{\quad
 \Lambda_0=n_{\rm top} A+
    (f-n_{\rm top})\bigl((8a-4)\ell+1-a/2\bigr).
 \quad} \tag{R39.15}
\]

The formula is independent of the interior word and of the canonical allowed twist. Both non-top roles have the same u endpoints. Every allowed ell gives 0\textless lambda\_0,tau\textless=A and consequently 0\textless Lambda\_0\textless=fA. These are exact integer comparisons with ell\_+\textless2\^{}485 and a=42 or168.

Using only pi N subset L and total polynomial degree would give the larger exponent D\_tau+c\_tau at one factor. R39.12 replaces it by nu\_0,tau+c\_tau=lambda\_0,tau. The reduction in that sufficient archimedean cost is exactly

\[
 {\log T\over2}\sum_\tau(D_\tau-\nu_{0,\tau})\ge0.
 \tag{R39.16}
\]

This concerns the initial coefficient construction only. The true finite valuations of beta\_lambda must be retained as well. In fact the product formula gives their normalized total exactly Lambda\_0 log(T)/2, matching R39.14. Thus the scalar cannot be counted as a free arithmetic gain after its norm has been removed. R39.5's sharpness is on the entire coefficient lattice, not on its Mram subspace or the prescribed theta image. Actual all-place period nonvanishing, remaining unimodular-frame/character/operator heights, the mixed Host/Scale adapter and the full DP44 comparison remain.

\subsection{6. Source locators and verification}\label{source-locators-and-verification}

All source interfaces are at the stipulated archive commit 7cbf7a0c32fca1badf05931879a37c3f8a2b03b4; no Lean build or audit.

\begin{itemize}
\tightlist
\item
  \nolinkurl{GradedWeylDivisibility/StatementA.lean59–76:} highest polynomial and homogeneous divided-root construction.
\item
  \nolinkurl{GradedWeylDivisibility/StatementD.lean8–38:} row weights, IsWeightBasis, literal semistandard minor basis and diagonal action.
\item
  \nolinkurl{LocalRingCardinality/DataCoeff.lean20–23,49–67:} signed determinant shift, polynomial degree convention and transpose action.
\item
  \nolinkurl{LocalToGlobal/Divisibility/CoefMod.lean184–217,314–330:} standard and transported coefficient lattices and bilinear dual.
\item
  \nolinkurl{CoefModLemmas.lean129–145,199–205:} Delta subset Nabla, covariant frame transport, integral stability and RoundLat.
\item
  \nolinkurl{LocalToGlobal/Columns/CoefCol.lean59–65,96–108:} actual initial roundPk on Lsh and terminal squarePk on N.
\item
  DecoFin.lean1131--1205: covector kernel, inclusions and localization; K044 EG2--5 and R022 R22.6--8: exact endpoint weights.
\item
  DO040 and R017: the unchanged completely split T prime, balanced norm generator pi and its selected-prime factorization.
\end{itemize}

The finite diagnostic script checks literal minor row-weight ranges, signed determinant shifts, sharp scalar exponents, split tensor valuation addition and all actual endpoint comparisons. It does not enumerate the actual coefficient module or compute a period.

\clearpage

\section{AM-DO-057 --- bounded adapted generators for the actual Hermitian coefficient lattices}\label{am-do-057-bounded-adapted-generators-for-the-actual-hermitian-coefficient-lattices}

Author: GPT-6 Astra Pro, ALP\_MACHINE \nolinkurl{/root/definite_oscillator}. 22 November 2026 UTC. Inherited reasoning effort is not separately exposed.

This finding replaces the uncomputed class-frame factors in DO056(AC25--28) by explicit, uniform bounds for a new integral field-spanning family. It first reduces the underlying Hermitian lattice, whose rank is a or N, and only then applies the polynomial representations. It preserves the actual eleven-adic lattice and adds no coefficient field. The construction is not an integral basis or a saturated inverse.

The initial source lattice is \textbf{delAtr(Lsh)}. Its polynomial lattice is the ordinary dual of Lsh, which is \textbf{bar(L)} for the standard Hermitian form. It is not the unimodular lattice N. The covector-kernel condition at T therefore remains in the bounds. The independent determinant-line construction below retains its signed determinant contribution exactly. A sharper highest-weight inclusion, if proved separately, can replace the elementary T bound without changing this frame construction.

\subsection{1. Literal source conventions and the lattices to be bounded}\label{literal-source-conventions-and-the-lattices-to-be-bounded}

Write E for the CM field, F=E\^{}+, {[}E:Q{]}=2f, and O=O\_E. The field of coefficient values is K and is not denoted by E. All Hermitian lattices in this finding use the standard definite form h=I, positive at every complex embedding. They are not the alternating CM lattices with coefficient xi\_K.

For a dominant coefficient weight lambda\_tau of rank n set

\[
 m_\tau=\lambda_{n-1,\tau},\qquad
 \nu_{i,\tau}=\lambda_{i,\tau}-m_\tau\ge0,\qquad
 D_\tau=\sum_i\nu_{i,\tau}.
 \tag{BF1}
\]

Use D\_Sigma=sum D\_tau, M\_Sigma=sum \textbar m\_tau\textbar, and S\_Sigma=sum m\_tau. The literal polynomial action is det(g)\^{}m times substitution by g\^{}t. Thus the coefficient functor on an underlying lattice M is

\[
             (\det M)^{m_\tau}\otimes\Delta_{\nu_\tau}(M)
       \quad\hbox{or}\quad
             (\det M)^{m_\tau}\otimes\nabla_{\nu_\tau}(M).
 \tag{BF2}
\]

This is an equality of the actual transported models, not an identification of Delta with Nabla. The defining source checks are:

\begin{itemize}
\tightlist
\item
  \nolinkurl{CoefMod.lean}79--86 defines detTwK/trAct/actP; 193--210 defines transLat as all finite local-basis membership tests.
\item
  \nolinkurl{CoefMod.lean}213--222 defines dualFr by the ordinary dot product; 309--329 distinguishes rhoV/nabAt/delAt from rhoVr/nabAtr/delAtr.
\item
  \nolinkurl{CoefCol.lean}59--72 defines roundPk and squarePk; 96--99 uses roundPk on d.Lsh in colVal; 104--107 uses squarePk on d.N in rowVal.
\item
  \nolinkurl{DataCoeff.lean}20--23 defines the determinant shift and nonnegative polynomial weight. \nolinkurl{StatementD.lean}24--49,131--141 records the standard weight basis and its exact determinant twist.
\end{itemize}

Let N be the unimodular lattice at an initial decorated class and L the kernel of its primitive residue covector at every selected T place. The source Lsh is L's Hermitian dual. Since h=I and L is reflexive,

\[
        \operatorname{dualFr}(L^\sharp)=\overline L,
        \qquad
        \mathrm{colVal.Del}=\Delta_\lambda(\overline L).
 \tag{BF3}
\]

Indeed ordinary dot-product duality is conjugation of Hermitian duality, and taking the Hermitian dual twice returns L. This remains true after the indicated coefficient embeddings and local scalar extension. The final detector lattice is instead Nabla\_lambda(N) for the terminal unimodular N.

Let pi=pi\_T be the fixed element with pi bar(pi)=T. The selected T primes are exactly those dividing pi, each to exponent one. Nonzero residue covectors are surjective, so their kernels have quotient length one. Localizing the kernel at all finite primes gives

\[
 \pi N\subset L\subset N,\qquad
 \det L=\pi\det N,\qquad
 \det\overline L=\overline\pi\det\overline N.
 \tag{BF4}
\]

The first inclusion is also the literal \nolinkurl{Deco.Ok.smul_N_mem_L}, \nolinkurl{DecoFin.lean}1156--1160. Its \nolinkurl{mem_L}1125--1144 identifies the actual kernel; \nolinkurl{L_le_N}, \nolinkurl{N_le_Lsh} and lattice statements occur at1146--1190. For the determinant equality choose a local basis in which the surjective covector is the first coordinate; the kernel basis is diag(pi\_P,1,\ldots,1). At every other prime the kernel is N. The resulting determinant ideal is the product of the selected primes, namely (pi). In particular L=N locally at11. No choice of a different T decoration is made.

\subsection{2. Trace covolume, minimum, and an E-spanning short ball}\label{trace-covolume-minimum-and-an-e-spanning-short-ball}

For any integral unimodular positive Hermitian O-lattice N of rank n, give its real space the Euclidean norm

\[
 \|x\|_{\mathrm{tr}}^2
     =\operatorname{Tr}_{F/\mathbf Q}h(x,x)
     =\sum_{\tau\in\Phi}h_\tau(x,x),\qquad d=2fn.
 \tag{BF5}
\]

Every nonzero x has h(x,x) a totally positive algebraic integer. Its norm to Q is a positive integer. Arithmetic--geometric mean therefore gives \textbar\textbar x\textbar\textbar\_tr\^{}2\textgreater=f.~Thus every successive minimum of the trace lattice is at least sqrt(f).

Its exact covolume is

\[
              c_n=2^{-fn}|d_E|^{n/2}.
 \tag{BF6}
\]

Here is a proof that also allows nonfree O-lattices. The real bilinear pairing is (1/2)Tr\_E h. Because N is selfdual for h, its Z-dual for this pairing is 2 D\_E\^{}\{-1\} N. This follows locally from the definition of the inverse different: testing against all O-multiples of each local basis vector is precisely the condition h(x,N) subset 2 D\_E\^{}\{-1\}. Multiplication by a fractional ideal A changes trace covolume by \textbar Norm(A)\textbar\^{}n; this can be checked on every local elementary divisor and hence on the global finite index after clearing a denominator. Consequently

\[
 c_n^{-1}=\operatorname{covol}(N^*)
      =2^{2fn}|d_E|^{-n}c_n,
\]

which proves BF6. This argument does not assume that N has a global O-basis or that its determinant ideal is principal.

Let V\_d be the volume of the Euclidean unit ball. We use the established Minkowski second theorem, in exactly the form of Henk (2002, Theorem1.3, PDFp2): product\_i lambda\_i V\_d \textless=2\^{}d c\_n.~Its Definition1.1 uses linearly independent lattice vectors, not a lattice basis. Put k=2f(n-1)+1. Applying the lower bound sqrt(f) to the first 2f(n-1) minima and monotonicity to the remaining 2f gives

\[
 \lambda_k\le R_E(n):=
 \left({2^{fn}|d_E|^{n/2}\over
                    V_{2fn} f^{f(n-1)}}\right)^{1/(2f)}
       \le\sqrt f\,(n\,\mathrm{rd}(E))^{n/2}.
 \tag{BF7}
\]

For the last inequality use V\_(2fn)=pi\^{}(fn)/(fn)!, pi\textgreater2 and (fn)!\textless=(fn)\^{}(fn). The real span of an E-subspace of dimension at most n-1 has dimension at most2f(n-1). The k independent short vectors supplied by this successive minimum therefore E-span the entire space. Selecting an E-independent subset gives n vectors v\_i in N with trace norm at most R\_E(n). This is the reason for using lambda\_k rather than lambda\_d.~No claim that the selected vectors form an integral basis is made.

\subsection{3. A bounded frame with the prescribed residue basis at11}\label{a-bounded-frame-with-the-prescribed-residue-basis-at11}

Fix an integral Z-basis omega\_1,\ldots,omega\_(2f) of O with

\[
                  |\sigma(\omega_j)|\le M_E
                  \quad\hbox{for every complex embedding}.
\]

The d vectors omega\_j v\_i are Z-independent, lie in N, and have trace norms at most M\_E R\_E(n). Their Z-span N\_0 is a full sublattice. Its centered fundamental parallelepiped has covering radius at most d M\_E R\_E(n)/2, by rounding each real coordinate in this basis to the nearest integer.

Choose any prescribed O/11O-basis of N/11N, for example the reduction of the already fixed standard11 frame. The module is free of rank n over that finite semilocal ring, since N is locally free of rank n.~Lift its n vectors arbitrarily to N. Subtract elements of11N\_0 using the preceding rounding. The resulting vectors y\_i have the same prescribed residues and satisfy

\[
 \|y_i\|_{\mathrm{tr}}\le R_n:={11dM_E\over2}R_E(n)
       \le 11fnM_E\sqrt f\,(n\,\mathrm{rd}(E))^{n/2}.
 \tag{BF8}
\]

They are E-independent, since their reductions are bases at every place over11. Write A=(y\_1,\ldots,y\_n). Then AO\^{}n subset N, and A is a local O-basis at all those places. The original lattice, class representative, level, covector and standard11 identification have not been altered. The new frame agrees with that identification modulo11; it need not equal the original frame to every order.

For completeness its index and inverse admit the effective bounds

\[
 \begin{split}
 I_A&=[N:AO^n]=\prod_{\tau\in\Phi}
       \det(h_\tau(y_i,y_j))\le R_n^{2fn},\qquad 11\nmid I_A,\\
 \|A_\tau\|_{\mathrm{op}}&\le\sqrt n R_n,\qquad
 \|A_\tau^{-1}\|_{\mathrm{op}}
       \le n^{(n-1)/2} R_n^{fn-1}.
 \end{split}
 \tag{BF9}
\]

The index equality follows by comparing BF6 for N with the covolume of the free sublattice in the complex frame A; Hadamard gives its bound. As I\_A\textgreater=1 and each other determinant has modulus at most R\_n\^{}n, \textbar det A\_tau\textbar\textgreater=R\_n\^{}(-n(f-1)). The cofactor/singular-value bound \textbar\textbar A\_tau\^{}-1\textbar\textbar\textless=\textbar\textbar A\_tau\textbar\textbar\^{}(n-1)/\textbar det A\_tau\textbar{} proves the last formula. These inverse estimates are not raised to a determinant power in the construction below.

The proof is effective in the following finite sense. A presentation of N and its positive forms allows enumeration in the BF7 ball until an E-basis is found, followed by integral linear algebra and the stated rounding to the prescribed residue lifts. The bound is uniform in the class; no efficiency or actual enumeration of the large class set is claimed.

\subsection{4. Independent determinant sections and integral polynomial transport}\label{independent-determinant-sections-and-integral-polynomial-transport}

The determinant lattice det(N), with its induced Hermitian metric, and its linear-dual lattice, with the dual metric, are rank-one integral unimodular positive Hermitian lattices. Apply BF8 with n=1 to each, prescribing a nonzero residue basis at11. This gives

\[
 \delta_+\in\det N,\qquad
 \delta_-\in(\det N)^\vee,\qquad
 0\ne\delta_\pm,\qquad |\sigma\delta_\pm|\le R_1.
 \tag{BF10}
\]

Both are local generators at11. They need not be global generators, and their product is not required to be1. Each spans its determinant line over E. For an integer m put s\_m=delta\_+\^{}m if m\textgreater=0 and s\_m=delta\_-\^{}(-m) if m\textless0. It is an integral, locally eleven-primitive section of (det N)\^{}m with archimedean norm at most R\_1\^{}\textbar m\textbar. This is why the factor I\_A\^{}\textbar m\textbar{} is unnecessary.

We use the elementary polynomial-lattice fact

\[
 AO^n\subset M\ \Longrightarrow\
 \rho_\nu(A)\Delta_{\nu,\mathrm{std}}\subset\Delta_\nu(M),
 \qquad
 \rho_\nu(A)\nabla_{\nu,\mathrm{std}}\subset\nabla_\nu(M).
 \tag{BF11}
\]

Here rho\_nu is only the degree-D polynomial action, without the determinant twist. To verify BF11 using the literal local definition, fix a local basis B of M. The matrix B\^{}-1 A is integral. Over the coefficient DVR factor it into two invertible integral matrices and a diagonal matrix with nonnegative valuations. The standard Weyl and dual-Weyl lattices are preserved by the invertible factors. Both have integral decompositions by polynomial row weight, with nonnegative weights; the diagonal factor multiplies each such summand by an integral scalar. For Nabla this is also direct Cauchy--Binet on the source bideterminants. For Delta the weight decomposition follows from its divided-root homogeneous generators, equivalently from the perfect Weyl/dual-Weyl pairing and the integral weight basis. This proves every \nolinkurl{transLat} membership test. It does not identify the two lattices or use an integral orthogonal projection.

If delta is an integral section of (det M)\^{}m, then delta rho\_nu(A)p belongs to the twisted lattice BF2. In a local basis B the extra scalar is delta/(det B)\^{}m, which is integral by precisely that line membership. Thus the polynomial frame and determinant section may be chosen independently. At11 both transformations are invertible integral; no new eleven-adic clearing scalar occurs.

The normalized Fischer metric is the invariant metric used in R037/R038. The degree-D polynomial action satisfies

\[
              \|\rho_\nu(A_\tau)\|\le(\sqrt n R_n)^D.
 \tag{BF12}
\]

This is the Cartan tensor realization in products of exterior powers, whose total exterior degree is D; the transpose realization has the same operator norm. The determinant section contributes its own R\_1\^{}\textbar m\textbar. No inverse Gram matrix is estimated.

\subsection{5. The actual initial and terminal adapted families}\label{the-actual-initial-and-terminal-adapted-families}

\subsubsection{Initial delAtr(Lsh)}\label{initial-delatrlsh}

Apply BF8 and BF10 to the unimodular conjugate lattice bar(N). Its norms are the same. By BF4 the columns of bar(pi) A lie in bar(L). Suitable determinant sections of bar(L) and its dual are bar(pi) delta\_+ and bar(pi)\^{}(-1) delta\_-. Since every embedding of pi has modulus sqrt(T), the resulting polynomial/determinant transport has the exact T factor displayed below and operator norm at most

\[
        T^{(D_\tau+m_\tau)/2}
            (\sqrt a R_a)^{D_\tau}R_1^{|m_\tau|}.
 \tag{BF13}
\]

The exponent m\_tau is signed. Replacing it by a\textbar m\_tau\textbar{} would be a different and needlessly larger estimate. For m\_tau\textless0, the reciprocal pi is integral as a section of the dual determinant lattice; it is not asserted to be an algebraic integer in the untwisted scalar line.

Use the existing standard Weyl spanning family of R038. It is integral, field-spanning, has norm at most a\^{}D\_tau at every complex embedding, and uses only the already present T-th roots of unity. Its degree condition D\_tau\textless T holds for every actual initial word/profile in K044, as proved in R38.9. Transport that family by the preceding polynomial frame and determinant sections, independently in the split coefficient-field factors. The resulting finite family p\_(c,i) lies in the actual delAtr(Lsh) at each decorated class c, field-spans that coefficient space, and satisfies uniformly

\[
 \boxed{\ \|p_{c,i}\|\le
 B_0:=a^{3D_{0,\Sigma}/2}R_a^{D_{0,\Sigma}}
           R_1^{M_{0,\Sigma}}
           T^{(D_{0,\Sigma}+S_{0,\Sigma})/2}.\ }
 \tag{BF14}
\]

All choices are in E and the existing coefficient field K. Independent CM-label grid choices refer to the scalar-extended product representation, not to an assertion that one E-matrix has arbitrary prescribed embeddings. The actual action on class values is still rhoVr and its actual character; BF14 constructs vectors in the fixed polynomial module and does not change that action.

\subsubsection{Terminal intNab}\label{terminal-intnab}

For each final unimodular rank-N lattice use BF8/BF10 directly. The literal standard semistandard dual-Weyl basis has normalized Fischer norms at most1 by R037. Transporting it gives an integral field-spanning family z\_(d,j) in the actual terminal intNab with

\[
 \boxed{\ \|z_{d,j}\|\le
 B_T:=(\sqrt N R_N)^{D_{T,\Sigma}}R_1^{M_{T,\Sigma}}.\ }
 \tag{BF15}
\]

This does not replace intNab by intDel. The family need not be an integral basis globally, since the underlying short frame can have a nontrivial index. Its eleven-adic frame change is integral invertible. Initial and terminal polynomial local tests are the original tests from BF3, not tests in an auxiliary unimodular lattice at T.

\subsection{6. Numerical constants on the unchanged two hosts}\label{numerical-constants-on-the-unchanged-two-hosts}

The existing constructions K030(PH2,PH5) and K032(QH2,QH7--11,QH15) give integral bases with explicit archimedean bounds. For J14 take powers of zeta28 times the seven Gaussian periods eta\_j. A period is a sum of four roots of unity, so M\_E=4. The periods form an integral basis by the exact trace determinant29\^{}6, and the two discriminants are coprime. For J56 tensor this basis with1,u\_1,u\_3,u\_1u\_3. At all embeddings \textbar u\_1\textbar\textless2, \textbar u\_3\textbar\textless4, so M\_E=32. These are bases of O\_E; they are not the trace-orthonormal fractional basis of the abelian variety's polarization.

The exact discriminants and conservative root-discriminant bounds are

\[
 \begin{array}{c|c|c|c}
 &f&M_E&|d_E|\ \ \ (\mathrm{rd}(E)\text{ bound})\\\hline
 J14&42&4&2^{84}7^{70}29^{72}\quad(<2^8)\\
 J56&168&32&2^{336}7^{280}29^{288}5^{168}13^{168}\quad(<2^{11}).
 \end{array}
 \tag{BF16}
\]

The root-discriminant comparisons mean \textbar d\_E\textbar\textless2\^{}(2f b), with b=8 or11; they are checked by exact integer inequalities. Substituting into BF8 gives

\begin{longtable}[]{@{}lrrrrr@{}}
\toprule\noalign{}
Host & initial a & terminal N & R\_1 & R\_a & R\_N \\
\midrule\noalign{}
\endhead
\bottomrule\noalign{}
\endlastfoot
J14 & 42 & 3561 & \textless2\^{}19 & \textless2\^{}315 & \textless2\^{}35637 \\
J56 & 168 & 3435 & \textless2\^{}27 & \textless2\^{}1625 & \textless2\^{}39536 \\
\end{longtable}

In particular the bounds BF14--15 imply

\[
 \begin{array}{c|c|c}
 &\log_2 B_0\text{ upper bound}&\log_2 B_T\text{ upper bound}\\\hline
 J14&324D_{0,\Sigma}+19M_{0,\Sigma}
       +\tfrac12(D_{0,\Sigma}+S_{0,\Sigma})\log_2T
       &35643D_{T,\Sigma}+19M_{T,\Sigma}\\
 J56&1637D_{0,\Sigma}+27M_{0,\Sigma}
       +\tfrac12(D_{0,\Sigma}+S_{0,\Sigma})\log_2T
       &39542D_{T,\Sigma}+27M_{T,\Sigma}.
 \end{array}
 \tag{BF17}
\]

If the signed coefficient of log T is negative, it stays negative; an upper estimate for T cannot then be substituted without reversing the inequality. BF17 leaves it exact. The numerical constants are deliberately elementary and are not asserted to be optimal. The coefficient normalization is the normalized source Fischer metric; CM derivative periods and creation scales still belong to the operator comparison DO056(AC4--7), not to BF17.

\subsection{7. Actual integral witnesses and exactly what this closes}\label{actual-integral-witnesses-and-exactly-what-this-closes}

Use the unchanged stabilizer actions, the literal ramified projector and the genuine terminal projector exactly as in K059. Initial stabilizers are1 by K042. Apply C\_T P\_ram, C\_T=(T-1)\^{}f, to the class-supported BF14 generators. On terminal classes take their actual stabilizer sums and then dG e\_A, dG=\textbar Gcal\textbar. These operations preserve the actual integral modules and their images field-span the required two subspaces. Hence the generator norm product is bounded by

\[
         2^{K_{\rm fin}}B_0B_T,
       \qquad K_{\rm fin}=112031999\text{ or }490748794,
 \tag{BF18}
\]

using the exact finite factors already checked in DO056(AC25--28). No uncomputed d\_(0,c), d\_(T,d) or norm of a dual standard Weyl basis remains in BF18. The new adapted family replaces the old clearing method; it does not bound that method's particular integers, whose unboundedness was proved in AC33--34.

If the actual moved Mram bottom period in DO055(CI3) is nonzero, its established actual-image implication makes B(P,Z)=PtilFF(A\_cpl P,Z) nonzero on these spaces. Bilinearity then gives at least one pair from the preceding finite integral families with

\[
                    B(x,z)\ne0,
                    \qquad\|x\|_0\|z\|_T<2^{K_{\rm fin}}B_0B_T.
 \tag{BF19}
\]

The bound is uniform at every complex embedding of K. No integral inverse of the vertical lift or Sad is used. The actual terminal mass is retained: the two scalar automorphic norm bounds are respectively sqrt(v0/d0)\textbar\textbar x\textbar\textbar{} and sqrt(vT/dW)\textbar\textbar z\textbar\textbar, exactly DO056(AC3). The unit-vacuum/CM factor, fixed phase, real recoupling, branch and class masses remain in AC19; none is absorbed into a redefined coefficient lattice.

This is a positive effective replacement for the class-frame-height gate in that spanning argument. It does not prove the all-T Mram nonzero period, does not evaluate a nonzero algebraic value, and does not assert that the resulting bounds satisfy the requested DP44 budget. The height/norm of the fixed arithmetic operator, the actual value field and the remaining normalized nonzero comparison must still be combined with BF18. A large sufficient bound is not a lower bound on optimal cost. A separate sharper inclusion for the one-covector T lattice can reduce BF13--14 without changing any of BF5--12 or BF15.

\subsection{Sources and proof-check scope}\label{sources-and-proof-check-scope-1}

Henk, M. (2002). \emph{Successive minima and lattice points}. arXiv:math/0204158. \url{https://arxiv.org/abs/math/0204158} . Definition1.1, PDFp1, and Theorem1.3, PDFp2, are the only external geometry-of-numbers input. The actual primary PDF and text are retained in \nolinkurl{definite_oscillator_frame_refs/henk2002.pdf} and \nolinkurl{.txt}. The established theorem is cited, not reproved here.

All other mathematical inputs are either proved above or the already checked source polynomial-lattice conventions, K030/K032 integral field bases, R037/R038 coefficient families, K042/K059 stabilizer/projector construction and DO055--056 actual-image/metric identities. The accompanying exact diagnostic certificate tests the discriminant/radius bounds, residue rounding, determinant-twist transport and polynomial integral examples. It does not enumerate actual high-rank classes or prove an automorphic nonzero assertion. Author whole-proof review and independent reviews are recorded separately in the event ledgers.

\subsection{Author verification record}\label{author-verification-record-6}

Author whole-proof PASS at 2026-11-22T08:51:58.274924+00:00. Read all seven proof sections/BF1--19 and the complete certificate, the actual local coefficient/source kernel definitions, the integral field-basis inputs and primary Henk Definition1.1/Theorem1.3. Reconstructed the half-trace covolume, E-spanning successive-minimum bound, prescribed-residue rounding, independent determinant sections, both actual coefficient directions and the signed T factor. The exact certificate passed15,738 diagnostics at2026-11-22T08:50:43.094384+00:00. The claim remains uniform adapted integral field-spanning bounds, conditional nonzero witnesses; no saturated inverse, actual class enumeration, all-T period or DP44 success is asserted.

\clearpage

\section{ALP-LIFT-46 --- exact character divisors and bounded coefficient fields for the actual mixed CM hosts}\label{alp-lift-46-exact-character-divisors-and-bounded-coefficient-fields-for-the-actual-mixed-cm-hosts}

Author: GPT-6 Astra Pro, ALP\_MACHINE /root/lift\_accounting. 22 November 2026 UTC. Inherited reasoning effort is not separately exposed.

Status: author whole-proof PASS after reading all eight sections and the complete certificate. The separate exact arithmetic replay passed 60,922 assertions at 08:54:21 UTC; the author-check timestamp is recorded in \nolinkurl{lift_accounting_CHECK_EVENTS.tsv.} No Lean compilation is claimed.

\subsection{1. Result and its precise scope}\label{result-and-its-precise-scope}

Retain the actual hosts, CM type, fixed odd character, and coherent stage characters of AM-K029--032 and AM-K040. Put \[
 E=\mathbb Q(\zeta _7),\qquad
 K_1=\mathbb Q(\zeta _{28})B,\qquad
 K_4=K_1(\sqrt5,\sqrt{13}),\qquad
 f=42h,\quad d=[K_h:\mathbb Q]=2f,\quad h\in\{1,4\},
 \tag{CF1}
\] where B is the real degree-seven subfield of \(\mathbb Q(\zeta _{29})\). Write K=K\_h, F=K\^{}+, \(\Psi=\Psi_h\), and \(H=h(K)\). Let \[
 H_{\rm rel}=\bigl|\operatorname{Cl}(K)/j\operatorname{Cl}(F)\bigr|.
\] The accepted construction supplies an everywhere finite-unramified unitary odd character eta and characters alpha\_e for \(e\in(28\mathbb Z)^\Psi\), with the exact principal-ideal rules used below. These characters themselves are unchanged in this finding.

For every integral ideal A of K, define the fractional ideal \[
 \mathcal D_e(A)=
       \prod_{\tau\in\Psi}
           \left(\frac{\tau A}{\tau\bar A}\right)^{e_\tau}.
 \tag{CF2}
\] The new conclusions are:

\begin{enumerate}
\def\labelenumi{\arabic{enumi}.}
\item
  In any coefficient extension containing the values, \[
   (\alpha_e(A))=\mathcal D_e(A),\qquad
   h_{\rm W}(\alpha_e(A))
      =\frac{\log N_K(\mathcal D_e(A)^-)}d
      \le \frac{\sum_\tau|e_\tau|}{d}\log N_K A.
   \tag{CF3}
  \] Ideals on the right are extended to the coefficient field in the first equality. The minus superscript means the integral denominator ideal after cancellation. Every complex conjugate of alpha\_e(A) has modulus one. No ideal-class exponent multiplies the height bound.
\item
  For the exact source normalization \[
         \beta(A)=\eta(A)(N_K A)^{-1/2},
  \] one has \[
   (\beta(A))=\left(\prod_{\tau\in\Psi}\tau A\right)^{-1},
      \qquad h_{\rm W}(\beta(A))=\tfrac12\log N_K A.
   \tag{CF4}
  \] This concerns normalized etaVal, not all unnormalized odd eta values in a single finite field.
\item
  One Galois CM field L contains K, every beta(A), and every alpha\_e(A), simultaneously for all ideals and all weight profiles, with \[
   [L:\mathbb Q]\le d\,\varphi(H)H^3H_{\rm rel}^{\,f}
                    \le d H^{f+4},\qquad
   \operatorname{rd}(L)
       \le \operatorname{rd}(K)H^4H_{\rm rel}^{\,f}
       \le \operatorname{rd}(K)H^{f+4}.
   \tag{CF5}
  \] Here \(\varphi(1)=1\). The extension L/K is unramified outside primes dividing H. Its construction uses the existing fixed characters and does not change a class-character choice.
\item
  Adjoining any prescribed finite cyclotomic exponent M multiplies the degree bound by at most \(\varphi(M)\le M\), and the root discriminant bound by at most M. In particular M may include the fixed G characters and the actual additive T character. The separately prescribed finite Q roots can be included by replacing M with their common multiple; no bound for an unspecified Q tuple is asserted.
\end{enumerate}

For the actual G primes, write \[
 M_G=\operatorname{lcm}(28,g_0-1,g_1-1).
\] The exact arithmetic certificate proves the following deliberately coarse, unconditional bounds. Every exponent is a strict base-two upper bound; logarithmic heights below use natural logarithms.

\begin{longtable}[]{@{}
  >{\raggedright\arraybackslash}p{(\columnwidth - 8\tabcolsep) * \real{0.1579}}
  >{\raggedleft\arraybackslash}p{(\columnwidth - 8\tabcolsep) * \real{0.2105}}
  >{\raggedleft\arraybackslash}p{(\columnwidth - 8\tabcolsep) * \real{0.2105}}
  >{\raggedleft\arraybackslash}p{(\columnwidth - 8\tabcolsep) * \real{0.2105}}
  >{\raggedleft\arraybackslash}p{(\columnwidth - 8\tabcolsep) * \real{0.2105}}@{}}
\toprule\noalign{}
\begin{minipage}[b]{\linewidth}\raggedright
Actual host
\end{minipage} & \begin{minipage}[b]{\linewidth}\raggedleft
\(H<2^{\,u}\)
\end{minipage} & \begin{minipage}[b]{\linewidth}\raggedleft
\(M_G\)
\end{minipage} & \begin{minipage}[b]{\linewidth}\raggedleft
\([L(\mu_{M_G}):\mathbb Q]<2^{\,v}\)
\end{minipage} & \begin{minipage}[b]{\linewidth}\raggedleft
\(\operatorname{rd}(L(\mu_{M_G}))<2^{\,w}\)
\end{minipage} \\
\midrule\noalign{}
\endhead
\bottomrule\noalign{}
\endlastfoot
\(J^{14}\) & \(u=456\) & 14995460592 & \(v=21017\) & \(w=21018\) \\
\(J^{56}\) & \(u=2477\) & 13526981751068240 & \(v=426107\) & \(w=426109\) \\
\end{longtable}

Including the actual DO040 additive T roots gives respectively degree/root-discriminant exponents 22018/22019 and 427103/427105. The field bounds are independent of all admissible e, ell, and canonical determinant twists. They do not identify a minimal field.

These results close one explicit coefficient-field and determinant-character cost. They do not complete the original quadratic-type Host, give an actual terminal Mram period, compute theta-image Smith factors, or establish a positive DP44 net budget.

\subsection{2. Carried inputs and exact character conventions}\label{carried-inputs-and-exact-character-conventions}

The actual type on K is induced from \(\Phi=\{1,2,3\}\) on E. Indeed the six E\_0 labels in K029/K040 are \(\{1,3,9,15,17,23\}\); reduction modulo seven gives each element of Phi twice. Each subsequent real extension retains all lifts. Consequently, for a in K, \[
 T_\Psi(a):=\prod_{\tau\in\Psi}\tau(a)
      =\prod_{i=1}^3\sigma_i(N_{K/E}a)\in E.
 \tag{CF6}
\] K029 ER23--26 constructs beta on ideals with the principal rule \[
                     \beta((a))=T_\Psi(a)^{-1}.
 \tag{CF7}
\] The correcting unramified class character is already part of beta. It is trivial on principal ideals. For all sigma in \(\operatorname{Gal}(K/\mathbb Q)\), \[
 |\sigma T_\Psi(a)|=|N_{K/\mathbb Q}a|^{1/2},\qquad
 T_\Psi(a)\overline{T_\Psi(a)}=N_{K/\mathbb Q}a.
 \tag{CF8}
\] The first equality holds because a Galois translate of a CM type still chooses one embedding from each conjugate pair.

K040 NC7--16 gives \[
 \alpha_e((a))
      =T_e(a):=\prod_{\tau\in\Psi}\tau(a/\bar a)^{e_\tau},
 \quad
 \alpha_e|_{jI_F}=1,\quad
 \alpha_{e+e'}=\alpha_e\alpha_{e'},\quad
 \alpha_{c\mathbf1}=\eta^{-2c}\quad(c\in28\mathbb Z).
 \tag{CF9}
\] Every Galois conjugate of T\_e(a) has modulus one. For i\textless f, K040 chooses a fixed basis character \(\xi_i=\alpha_{28e_i}\), and defines \[
 \xi_f=\eta^{-56}\prod_{i<f}\xi_i^{-1},\qquad
 \alpha_e=\prod_i\xi_i^{e_i/28}.
 \tag{CF10}
\] The finite-order factors \(\widetilde\nu_j\) in the actual \(\mu_j=\widetilde\nu_j\alpha_{e_j}\) take their values in \(\mu_{M_G}\). This is the zero-infinity construction of LIFT18 G12, followed by K031/K032 norm pullback; its rational factor has infinity character one. Multiplying by these roots changes neither the divisor nor the Weil height in CF3.

No existence theorem from K029 or K040 is reproved or counted as a new result. In particular their capitulation compatibility is retained.

\subsection{3. Exact ideal divisors and heights}\label{exact-ideal-divisors-and-heights}

Let A be integral. Choose a positive integer m annihilating its ordinary ideal class and write \(A^m=(a)\), with a integral. From CF9, \[
 \alpha_e(A)^m=T_e(a).
 \tag{CF11}
\] Taking the principal ideal of the right side gives \[
 (T_e(a))
       =\prod_\tau(\tau A/\tau\bar A)^{m e_\tau}
       =\mathcal D_e(A)^m.
\] Unique factorization of fractional ideals in the coefficient field allows cancellation of m from every valuation. This proves the first equality of CF3, including negative e\_tau and nonprincipal A. It also shows directly that alpha\_e(A) is a unit above every rational prime away from the support of A and its conjugates. In particular A prime to 11 implies an 11-unit value, even if the coefficient field itself is ramified at 11.

All roots in CF11 have complex modulus one, in every embedding: the conjugate polynomial has right side of modulus one. The absolute logarithmic Weil height therefore consists entirely of finite denominator terms. Extension of an ideal from K to L multiplies its logarithmic norm by {[}L:K{]}, so \[
 h_{\rm W}(\alpha_e(A))
       =d^{-1}\log N_K\mathcal D_e(A)^-.
 \tag{CF12}
\] For e\_tau\textgreater=0, the possible denominator contributed by that factor is \((\tau\bar A)^{e_\tau}\); for e\_tau\textless0 it is \((\tau A)^{-e_\tau}\). Cancellation only decreases the denominator. Every conjugate ideal has norm N\_K A. This proves the last bound in CF3.

Similarly CF7 gives \[
 \beta(A)^m=T_\Psi(a)^{-1},\qquad
 (\beta(A))=\left(\prod_{\tau\in\Psi}\tau A\right)^{-1}.
 \tag{CF13}
\] Every complex value has modulus \((N_K A)^{-1/2}\le1\), and the denominator ideal has norm \((N_K A)^f\). Division by d=2f proves CF4. Equivalently beta(A)\^{}-1 is an algebraic integer all of whose conjugates have modulus \(\sqrt{N_K A}\).

These proofs apply to the fixed characters, not only to some alternative extension of their infinity types. A different allowed unramified class-character factor is a root of unity and leaves the formulas unchanged.

For reference, the uniform-direction convention in CF10 is consistent with the same divisor formula: \[
  (\eta(A)^{-56})
      =(\beta(A)^{-56}(N_K A)^{-28})
      =\prod_{\tau\in\Psi}(\tau A/\tau\bar A)^{28}.
 \tag{CF14}
\] Here the ideal of the rational number N\_K A is the product of all d Galois conjugates of A. Thus no square-root value has been silently adjoined in using the final basis character.

\subsection{4. A finite Galois CM field containing the unchanged family}\label{a-finite-galois-cm-field-containing-the-unchanged-family}

Choose invariant-factor generators A\_j of Cl(K), of orders h\_j, and generators B\_s of Cl(K)/jCl(F), of orders d\_s: \[
 A_j^{h_j}=(a_j),\qquad
 B_s^{d_s}=(b_s)j(\mathfrak c_s),\qquad
 \prod_jh_j=H,\quad\prod_sd_s=H_{\rm rel}.
 \tag{CF15}
\] All representatives may be chosen integral. Replacing a representative for this analysis does not change any character: its value is already determined and satisfies the new relation.

Start with K(mu\_H). For each j adjoin roots of \[
                         X^{h_j}-\sigma(T_\Psi(a_j)^{-1})
 \tag{CF16}
\] for three representatives sigma of \(\operatorname{Gal}(E/\mathbb Q)/\langle\text{conjugation}\rangle\). Choose the root for sigma=1 to be the actual beta(A\_j), and arbitrary roots for the other representatives. All roots of each polynomial are then present because h\_j divides H. For a complementary conjugate, its roots are \[
                         (N_K A_j)^{-1}\zeta/\beta_\sigma,
                 \qquad \zeta\in\mu_{h_j},
 \tag{CF17}
\] so no fourth, fifth, or sixth root family is required. The degree added by the odd-character families is at most H\^{}3. The principal rule CF7 and the class decomposition show that the resulting field contains every beta(I), including the fixed class correction.

For each relative generator B\_s and every selected tau adjoin one root \[
                  z_{s,\tau}^{d_s}
                       =\tau(b_s/\bar b_s)^{28}.
 \tag{CF18}
\] For tau\_i with i\textless f, use the actual value xi\_i(B\_s). Its equation follows from CF9 and triviality on jI\_F. The unused last-label root may be chosen arbitrarily. Every Q-conjugate of a radicand in CF18 is another selected radicand or its inverse, because K is Galois CM. With mu\_H present, all its conjugate roots are therefore already included. The degree added is at most H\_rel\^{}f.~Principal/base-ideal decomposition gives all xi\_i values for i\textless f, and CF10/CF14 gives the last basis values and all alpha\_e values.

The resulting field L is Galois over Q. The entire set of odd root polynomials is stable under Galois conjugation using CF17, and the same assertion for the angular roots follows from CF18. All their roots are included; K itself is Galois. Its degree has the first bound in CF5.

L is also a CM field. Every conjugate of an angular root has modulus one; hence conjugation sends it to its inverse. For a root beta\_sigma from CF16, every conjugate has modulus \((N_K A_j)^{-1/2}\), so \[
                  \overline{\beta_\sigma}
                         =(N_K A_j)^{-1}\beta_\sigma^{-1}.
 \tag{CF19}
\] The constants on the right are rational and fixed by all automorphisms. Consequently complex conjugation commutes with every Q-automorphism on each generator, and hence on L. Its fixed field is totally real: for any embedding, the image of a conjugation-fixed element is real. Since L contains the totally imaginary field K, L is a totally imaginary quadratic extension of that fixed field. The assertion remains true after adjoining any roots of unity.

\subsection{5. Root-discriminant and ramification control}\label{root-discriminant-and-ramification-control}

We use the following elementary local lemma, stated with the integral-closure issue included.

\textbf{Lemma CF20.} Let A be the valuation ring of a characteristic-zero nonarchimedean local field k. Suppose c in k\^{}times has valuation divisible by a positive integer m. For any m-th root alpha of c, the different of k(alpha)/k divides the ideal generated by m. Thus a global extension obtained by adjoining an m-th root of a number with all finite valuations divisible by m satisfies \[
       \operatorname{rd}(k'_{\rm global})
             \le m\,\operatorname{rd}(k_{\rm global}),
 \tag{CF20}
\] and is unramified at primes not dividing m.

\textbf{Proof.} Locally rescale alpha by a power of a base uniformizer. This does not change its field and makes alpha a unit satisfying \(f(X)=X^m-u\), with u a unit of A. Let g be its monic minimal polynomial. Both g and f/g belong to A{[}X{]}, since A is integrally closed. Set O=A{[}alpha{]} and let B be the full ring of integers. The trace-dual formula for the monogenic order is \[
              O^*=g'(\alpha)^{-1}O.
 \tag{CF21}
\] This follows from the monogenic different formula, or directly from the dual basis obtained by polynomial division. The definition of the trace dual gives B* contained in O*. Hence \[
          g'(\alpha)B^*\subset O\subset B,
\] so g'(alpha) belongs to the different of B/A. Multiplying by (f/g)(alpha) shows that \(f'(\alpha)=m\alpha^{m-1}\) belongs to that different. Alpha is a unit, so m belongs to it. Thus the different ideal contains (m), equivalently its valuation is at most that of m. The norm-of-different and discriminant tower formulas give CF20. This argument does not assume that O=B. QED.

The trace-dual definition and its identification with the dualizing module are Stacks Project 49.8, especially Lemmas 49.8.2--3. The monogenic derivative formula is Example49.12.5; the identification in Lemma49.12.4 gives CF21.

For CF16, the principal ideal of its radicand is \[
 \left(\sigma T_\Psi(a_j)^{-1}\right)
        =\left(\prod_{\tau\in\Psi}\sigma\tau A_j\right)^{-h_j}.
 \tag{CF22}
\] All valuations are divisible by h\_j. For CF18, extension of the real ideal cancels on taking b\_s/bar b\_s, so \[
  \left(\tau(b_s/\bar b_s)^{28}\right)
         =(\tau B_s/\tau\bar B_s)^{28d_s}.
 \tag{CF23}
\] Again all valuations are divisible by the root order d\_s. Divisibility is preserved after every intermediate extension. Applying CF20 to the three odd families and f angular families gives factors H\^{}3 and H\_rel\^{}f.~Adjoining mu\_H gives an additional factor at most H, by the same lemma applied to X\^{}H−1. This proves the root-discriminant bound in CF5 and the ramification assertion. A subsequent mu\_M contributes at most M. No unknown size of a class-relation generator enters.

\subsection{6. Explicit class and generator bounds}\label{explicit-class-and-generator-bounds}

The discriminants, already proved in K030/K032 and recorded in K040 NC25, are \[
 |\Delta_{K_1}|=2^{84}7^{70}29^{72},\qquad
 |\Delta_{K_4}|=2^{336}7^{280}29^{288}5^{168}13^{168}.
 \tag{CF24}
\] Let \[
 B_h=\left\lceil
       \frac{4^f d!\sqrt{|\Delta_{K_h}|}}{3^f d^d}
                          \right\rceil.
 \tag{CF25}
\] Minkowski's ideal-class theorem and pi\textgreater3 imply that every class has an integral representative of norm at most B\_h. The square root in this formula is the exact integer given by CF24. Relative-class generators may be lifted to ordinary classes and represented with the same bound.

The coefficientwise ideal-count comparison \(\zeta_K(s)\le\zeta(s)^d\) for real s\textgreater1 is the elementary Euler-factor argument already supplied in K040 NC27. For any integer k\textgreater=1, sum the ideals of norm at most B\_h and use \(\zeta(1+1/k)<k+1\). It follows that \[
 H < B_h^{\,1+1/k}(k+1)^d,\qquad
 H^k < B_h^{\,k+1}(k+1)^{dk}.
 \tag{CF26}
\] The zeta inequality follows by bounding the tail by the integral of x\^{}(-1-1/k) from one to infinity; no zero-free estimate or class-number computation is used.

The certificate verifies with integers only:

\begin{longtable}[]{@{}
  >{\raggedright\arraybackslash}p{(\columnwidth - 6\tabcolsep) * \real{0.2000}}
  >{\raggedleft\arraybackslash}p{(\columnwidth - 6\tabcolsep) * \real{0.2667}}
  >{\raggedleft\arraybackslash}p{(\columnwidth - 6\tabcolsep) * \real{0.2667}}
  >{\raggedleft\arraybackslash}p{(\columnwidth - 6\tabcolsep) * \real{0.2667}}@{}}
\toprule\noalign{}
\begin{minipage}[b]{\linewidth}\raggedright
Host
\end{minipage} & \begin{minipage}[b]{\linewidth}\raggedleft
\(\log_2 B_h\) upper bound
\end{minipage} & \begin{minipage}[b]{\linewidth}\raggedleft
k used in CF26
\end{minipage} & \begin{minipage}[b]{\linewidth}\raggedleft
\(\log_2 H\) upper bound
\end{minipage} \\
\midrule\noalign{}
\endhead
\bottomrule\noalign{}
\endlastfoot
\(J^{14}\) & 216 & 3 & 456 \\
\(J^{56}\) & 1357 & 4 & 2477 \\
\end{longtable}

For example the last column is checked by raising both sides of CF26 to the integer k and comparing with \(2^{ku}\); it is not a rounded floating-point computation. Together with CF5, M\_G, and the exact DO040 T values this proves the table in section1.

There are also small explicit heights for the radical generators themselves. The odd generators satisfy \[
 h_{\rm W}(\beta_\sigma(A_j))\le\tfrac12\log B_h.
 \tag{CF27}
\] For each angular root in CF18, CF23 and modulus one give \[
 h_{\rm W}(z_{s,\tau})\le(28/d)\log B_h.
 \tag{CF28}
\] Thus these heights are respectively less than 108 log2 and 72 log2 on J14, and (1357/2) log2 and (1357/12) log2 on J56. The actual K040 last basis character is the product in CF10; it need not equal an arbitrarily chosen unused last-label radical. The field bound uses the latter only to close the Galois orbit.

\subsection{7. Consequences for the actual arithmetic bookkeeping}\label{consequences-for-the-actual-arithmetic-bookkeeping}

For a finite idele whose ideal A is integral and prime to11, the actual even-stage character values are11-units. Their complex conjugates all have modulus one. Exact finite valuations are given by CF2; denominator height is CF3. Consequently a determinant-ideal transition of known norm can be charged explicitly without multiplying by a class order or adding an archimedean norm penalty for alpha\_e. The finite-order nu factors do not alter this statement. For a fractional ideal the same divisor equality holds; one applies the positive/negative ideal decomposition instead of silently using the integral-ideal norm bound.

The odd normalization has its different, explicitly retained factor in CF4. Unnormalized eta(A) is not substituted for beta(A) in the coefficient field. Formula CF14 is why the even stage family nevertheless lies in the fixed field.

Normalized global arithmetic sums are invariant under coefficient extension. Indeed for an ideal I of L0 and an extension L/L0, \[
 \log N_L(I\mathcal O_L)=[L:L_0]\log N_{L_0}I,
 \qquad
 \frac{\log N_L(I\mathcal O_L)}{[L:\mathbb Q]}
       =\frac{\log N_{L_0}I}{[L_0:\mathbb Q]}.
 \tag{CF29}
\] Locally this is the identity \(\sum_{W\mid v}e(W/v)f(W/v)=[L:L_0]\). Archimedean averages repeat each old embedding the same number of times. Thus enlarging the coefficient field does not multiply a normalized theta-content gain or a normalized height by its absolute degree. If a particular arithmetic argument charges the coefficient field's discriminant, CF5 supplies an explicit upper bound. It does not turn that discriminant into an additional term in a source formula which does not contain it.

The exact source distinction is retained: DataParams2 uses its own local weights and ramification indices in the global sz/gain/error accounting. The field and type requirements of the original Host remain separately obstructed as in AM-R023. This finding does not replace those requirements, identify the analytic Omega with a new height, or infer an actual terminal row.

\subsection{8. Verification plan and references}\label{verification-plan-and-references}

The companion \nolinkurl{mixed_character_field_costs_certificate.py/.json} checks the actual CM-type orbit/complement identities, finite ideal-divisor and cancellation identities on deterministic finite Galois-orbit models, source odd/even normalization, relative-base cancellation, normalized base-change identities, and every actual integer field/height bound. It does not compute Cl(K), construct a minimal coefficient field, factor its discriminant, or evaluate a theta period.

Primary references:

\begin{itemize}
\tightlist
\item
  Milne, J. S. (2020). Algebraic number theory (Version3.08). Theorem4.3, printed p.70, for the exact Minkowski representative bound; Theorem8.8, printed p.138, for the normalized product formula. \url{https://www.jmilne.org/math/CourseNotes/ANT.pdf}
\item
  The Stacks Project Authors. (2026). The Stacks Project, Section49.8, Lemmas49.8.2--3, and Section49.12, Lemma49.12.4/Example49.12.5. Trace complements, monogenic duality, and the derivative different formula. \url{https://stacks.math.columbia.edu/tag/0BW0} \url{https://stacks.math.columbia.edu/tag/0BWB} \url{https://stacks.math.columbia.edu/tag/0BWI}
\end{itemize}

Carried exact source/results:

\begin{itemize}
\tightlist
\item
  AM-K029 ER23--26; K032 section8; AM-K040 NC6--27: the fixed actual odd/even characters, principal and relative-class rules, exact host discriminants and the elementary ideal-counting comparison.
\item
  LIFT18 G12--16 and K031/K032: the fixed zero-infinity G factors and their cyclotomic value field.
\item
  \nolinkurl{Formalism/LocalToGlobal/Columns/AbelianHostFieldConditions.lean} 314--331 and Formalism/LocalToGlobal/UnitaryNorms/ IsUnitaryNormOf.lean 90--94: exact IsMu0 and angular HasInfQ convention; K028/K029 record normalized etaVal.
\item
  \nolinkurl{Formalism/LocalRingCardinality/DataParams₂.lean,} definitions of sz, Ga, Ee, and FactEtaMean: normalized source accounting is a separate interface.
\end{itemize}

No statement of actual theta-image saturation or a successful DP44 comparison is included.

\clearpage

\section{AM-K062 --- averaging the actual T lattice pencils before the rational sum}\label{am-k062-averaging-the-actual-t-lattice-pencils-before-the-rational-sum}

Author: GPT-6 Astra Pro, ALP\_MACHINE \nolinkurl{/root/kummer_extension}. 22 November 2026 UTC. Inherited reasoning effort is not separately exposed.

\textbf{Author whole-proof check PASS --- 2026-11-22T08:56:58.377809+00:00.}

This report gives an exact averaged substitute for the all-base kernel hypothesis in LIFT45. The actual all-T projection is performed first, at the finite lattice pencils. Its local coefficient is an explicit fixed-point count. After a single rational transport of the whole favorable orbit, the all-T energy is the sum of a controlled integral contribution, a controlled infinite tail, and one finite short sum of nonintegral rational matrices. The sign of that last sum is not asserted. Consequently this is a concrete reduction of the Mram gate, not an Mram nonvanishing theorem or a DP44 cost conclusion.

\subsection{1. The actual orbit has a complete projective lattice pencil}\label{the-actual-orbit-has-a-complete-projective-lattice-pencil}

At a selected split degree-one place above T, let O be the local DVR, varpi its uniformizer, k its residue field, and q=\textbar k\textbar=T. Let Lambda=O\^{}n, where n=a+1. In the notation of LIFT43(TG2--3),

\[
 \tau=\begin{pmatrix}I_a&0\\\varpi^{-1}u&1\end{pmatrix},\qquad
 \delta_s=\operatorname{diag}(sI_a,1),\qquad s\in k^\times.
 \tag{AT1}
\]

The row u is primitive. The scalar representative sI\_a is permitted because u(sI\_a)=s u. A different lift changes only the existing strict level or a right maximal-compact factor. No new finite vector is chosen.

Let gamma be LIFT41's fixed rational matrix and write gamma\_f\^{}-1=x tau k0, with k0 in the actual maximal compact and principal at G. Use this same gamma for the entire scalar orbit. Left rational translation, followed by the same right k0 adjustment, represents it by

\[
              (p,h_s),\qquad p=\gamma_\infty o,\qquad
              h_s=\gamma_f\delta_s\gamma_f^{-1}.
 \tag{AT2}
\]

Indeed delta\_s is central in the embedded initial H, so it commutes with x. The base h1 is the identity and every h\_s is the identity away T. At G the fixed correction is principal, hence acts trivially on the given finite coefficient representation. The right T correction is in the maximal bottom level. Thus AT2 retains the exact original orbit and its coefficient frames. Its real centre p is fixed, even when a scalar point cannot individually be represented by a real point near the origin.

At the selected place, direct multiplication of AT1 gives

\[
 h_s\Lambda=(1+c\varpi^{-1}zv)\Lambda,
 \qquad c=s^{-1}-1,\qquad z=k_0^{-1}e_n,\quad v=(u,0)k_0,
 \quad vz=0.
 \tag{AT3}
\]

Both z and v are primitive. This is the same actual graph identity as LIFT44(PR4--7), with its two omitted points retained. The conjugate CM place is determined by the inverse-adjoint lattice and contributes no independent scalar variable.

Put

\[
 M=\{w\in\Lambda:vw\in\varpi O\},\qquad
 N=\Lambda+\varpi^{-1}Oz.
 \tag{AT4}
\]

Then varpi N is contained in M, and N/M is a two-dimensional k-space. Choose w with vw=1. Its classes represented by varpi\^{}-1 z and w form a basis of N/M. The lattice in AT3 is the inverse image of the line spanned by (c,1). Adding c=infinity gives all q+1 intermediate lattices

\[
     \mathscr L=\{L:M\subset L\subset N,\ \dim_k(L/M)=1\}
                                  \simeq\mathbb P^1(k).
 \tag{AT5}
\]

The actual orbit consists of the q-1 values

\[
       D=k\setminus\{-1\}\subset\mathbb P^1(k),\qquad
       c=\infty,-1\text{ omitted}.
 \tag{AT6}
\]

In particular c=0 is included and its lattice is the standard Lambda. The full pencil is a device for calculating intersections; its two omitted lattices are not added to the actual family.

\subsection{2. A partial permutation and its exact projector trace}\label{a-partial-permutation-and-its-exact-projector-trace}

For B in GL\_n of the local field, define the matrix on C{[}D{]}

\[
 (M_B)_{d,c}=\mathbf1_{B L_c=L_d},\qquad
 F(B)=\#\{c\in D:B L_c=L_c\},\quad
 C(B)=\#\{(d,c)\in D^2:B L_c=L_d\}.
 \tag{AT7}
\]

It is a partial permutation: every row and every column has at most one nonzero entry. Let m=q-1, let J be the all-ones matrix, and put P=I-J/m. This is exactly the probability projection removing the trivial scalar character, expressed in the orbit coordinates. Direct matrix multiplication gives

\[
 \operatorname{Tr}M_B=F(B),\qquad
 \operatorname{Tr}((J/m)M_B)=C(B)/m,\qquad
 r(B):=\operatorname{Tr}(PM_B)=F(B)-C(B)/m.
 \tag{AT8}
\]

The traces are unchanged on replacing B by B\^{}-1. All T local types of the bottom are maximal-compact fixed. Thus the nonzero entries in AT7 are exactly1; there is no omitted finite character or coefficient phase at T. Rational class transports are handled by the complete adelic kernel below, before passage to strict-class representatives.

The following elementary incidence alternatives give sharp useful bounds.

\textbf{Different full pencils.} If B(mathscr L) differs from mathscr L, then C(B)\textless=1. Indeed two distinct lattices in a pencil have intersection M and sum N. If two lattices belonged to both pencils, their intersection and sum would identify both defining pairs, and the pencils would be equal. Thus an unequal pencil meets this one at most once. It follows that \textbar r(B)\textbar\textless=1.

\textbf{Equal pencil, nontrivial projective action.} If B preserves the full pencil, it preserves M and N and induces an invertible k-linear map on N/M. If its projective action is not the identity, it has at most two fixed projective lines. Since D omits exactly two points,

\[
                  F(B)\le2,\qquad m-2\le C(B)\le m,
                  |r(B)|\le1+2/m.
 \tag{AT9}
\]

The fixed-line assertion follows by solving the eigenline equation of a non-scalar2-by-2 matrix; it also covers a repeated eigenvalue or no k-rational eigenvalues. The intersection bound follows because two sets obtained by removing two points from the same q+1 point set have at least q-3=m-2 points in common.

\textbf{Identity projective action.} In this case M\_B=I, so

\[
                 F(B)=C(B)=m,\qquad r(B)=m-1=q-2.
 \tag{AT10}
\]

It fixes every pencil lattice, in particular the included Lambda. Thus B and B\^{}-1 are integral. Define rho(B)=r(B)/(q-2). Since M\_B is a contraction and P has rank q-2, \textbar rho(B)\textbar\textless=1 in every case. If the projective action is not the identity, including the unequal-pencil case, the stronger bound is

\[
                     |\rho(B)|\le\frac2{T-2}.
 \tag{AT11}
\]

There is no termwise positivity statement. For example the projective translation c maps to c+b, b nonzero, fixes only infinity. For the actual D in AT6, F=0 and C=m-1, so

\[
                  r(B)=-(m-1)/m,\qquad\rho(B)=-1/m.
 \tag{AT12}
\]

Including b=0, the sum of r(B) over the whole additive translation subgroup is zero. This is an explicit finite-field orthogonality cancellation. It explains why replacing all projected terms by positive counts would be incorrect.

\subsection{3. Exact all-base energy, with the original class masses}\label{exact-all-base-energy-with-the-original-class-masses}

First consider any of the actual projective-T class orbits in the literal LIFT43/LIFT45 row. Its weight is w\_x=v\_P/\textbar Gamma\_x\textbar{} and its strict-volume factor is w\_x/m\^{}f.~No assumption on the stabilizer is made. In the scalar-orbit coordinates, the all-place projector is tensor\_v P\_v. Therefore, for the actual kernel K\_x(t,u)=Tr(J\_(x,t) J\_(x,u)\^{}*), the exact energy is

\[
 E_{x,\mathrm{ram}}=\frac{w_x}{m^f}
          \sum_{t,u\in D^f}(P_{\mathrm{ram}})_{u,t}K_x(t,u).
 \tag{AT13}
\]

This is LIFT43(TG4--5) without expanding the product of projections. It is also valid for the associated finite G-good coefficient bundle used in K058, at its corresponding coarse G level. In that model the finite fibre is Rgood tensor W, with the invariant probability metric; the class weight is the weight of that very bundle. Extracting a scalar G-generic component later retains K058's finite Schur factors. The bundle metric is not silently identified with the scalar generic-level metric.

Now use the favorable orbit and the single transport AT2. Let d\_R=dim Rgood and d\_W=dim W, with exactly K058's fixed genuine Pi and its two mixed I\_n projectors. Work in the actual cuspidal section space, or in the finite-dimensional span of the representers for these finitely many evaluations. For a rational unitary matrix g define

\[
 a(g)=\frac{\operatorname{Tr}(Q_{\rm good}R(g_G))}{d_R},\qquad
 c(g)=\frac{\operatorname{Tr}_W
        (P_W\pi(g_p^{-1}g_\infty g_p)|_W)}{d_W}.
 \tag{AT14}
\]

The expression a(g) is set to zero unless g is in the actual maximal compact at every finite place away T; at G its reduction is then read in the original fixed frame. All other away-T local factors are the existing maximal-level characteristic functions. No hyperspecial assertion at a ramified CM prime is needed. At T write B\_v=g\_v in the standard transported pencil AT3--6.

Here is a direct adelic derivation of the kernel, to make its extension from the standard component explicit. Realize the R-valued sections by right finite covariance F(gk)=R(k)\^{}-1F(g), with R trivial outside G. For the evaluation indexed by u, use the finite seed supported on h\_u K, whose value at h\_u k is R(k)\^{}-1, tensored with the actual real seed pr\_V pi(g\_infinity\^{}-1 g\_p)\textbar W. Periodize this seed over Z backslash G(F). It has exactly the unchanged right maximal-T invariance. The finite compact Haar is probability Haar, so unfolding contributes no additional finite volume. The central quotient contributes1/z0 and real orthogonality contributes1/C\_infinity, just as in K061(CG2--3).

Evaluating against the seed indexed by t leaves the condition h\_t\^{}-1 g h\_u in K after replacing the rational summation index by its inverse. Its finite G factor is R(g\_G), because both h\_t and h\_u are the identity at G. Its T factor is the characteristic function of that condition. The real factor is the compressed operator at g\_p in AT14. Thus the exact Poincare reproducing-kernel expansion is

\[
 K_x(t,u)=z_0 C_\infty d_Rd_W
      \sum_{g\in Z\backslash G(F)}a(g)c(g)
                       \prod_{v\mid T}(M_{B_v})_{t_v,u_v},
 \quad z_0=28,
 \quad C_\infty=\prod_{\tau,i}(r_{\tau i}-1).
 \tag{AT15}
\]

The local condition is h\_t\^{}-1 g h\_u in the bottom maximal compact, which is exactly g L\_u=L\_t. The real centre is the same g\_p at every t,u, so c(g) is independent of these two indices. This is the reason for averaging before introducing separate real representatives. At G there is no t-dependent finite deck factor, because h\_t=1 there. The central factors cancel as in K061; hence the same z0 C\_infinity normalization and raw Haar apply. Summing t,u in AT15 first, by AT8, gives the exact identity

\[
 \boxed{\quad
 E_{x,\mathrm{ram}}
   =\frac{w_x z_0 C_\infty d_Rd_W}{m^f}
     \sum_{g\in Z\backslash G(F)}a(g)c(g)
                         \prod_{v\mid T}r(B_v).
 \quad}\tag{AT16}
\]

This formula keeps every scalar-orbit base. It neither substitutes the favorable-base kernel for the other bases nor assumes that the bottom maximal-T space is invariant under conjugated diamonds. The analogous unprojected and single-average identities replace the local product in AT16 by, respectively,

\[
       \prod_v F(B_v),\qquad
       \frac{C(B_{v_0})}{m}\prod_{v\ne v_0}F(B_v).
 \tag{AT17}
\]

Thus it is also an exact averaged substitute for LIFT45(SP5--8). Absolute convergence, sufficient to exchange the finite sums, follows quantitatively in section4. Orbit-stabilizer descent preserves exactly w\_x/m\^{}f, so no scalar projector denominator or class mass is removed.

\subsection{4. All contributing rational matrices have bounded denominators}\label{all-contributing-rational-matrices-have-bounded-denominators}

Let mathscr C be the rational matrices satisfying the away-T compact condition in AT14 and C(B\_v)\textgreater0 at every T place. This contains the support of AT16. At each selected place a matrix in this support maps one lattice between M\_v and N\_v to another such lattice. Since

\[
           \varpi\Lambda\subseteq L_t\subseteq\varpi^{-1}\Lambda,
 \]

both B and B\^{}-1 have entries of valuation at least-2. The conjugate place obeys the same bound, by the inverse-adjoint unitary relation. Globally this proves

\[
          g,g^{-1}\in T^{-2}M_n(O_E)\qquad(g\in\mathscr C).
 \tag{AT18}
\]

In particular a ratio g1 g2\^{}-1 has entries in T\^{}-4 M\_n(O\_E). This uniform denominator bound holds for the full all-base sum; it does not assert integrality of each correspondence.

The power two is necessary for this all-base statement. Already on a two-dimensional local summand, take L={[}{[}1,varpi\^{}-1{]},{[}0,1{]}{]} and let S exchange its two coordinates. The matrix L S L\^{}-1 maps L Lambda to itself and has an entry 1-varpi\^{}-2. Both this example and its inverse are legitimate split local matrices; the conjugate CM component is their inverse adjoint. Thus the one-place, favorable-base denominator-one statement of LIFT45 must not be substituted in AT18.

Apply the K061 L1 sampling proof with

\[
       \rho_T=\frac1{100d_L^2T^4},\qquad
       K_T=2^{1200d_L}T^{8d_L},\qquad d_L=f(a+1)^2.
 \tag{AT19}
\]

The translated charts for any distinct g1,g2 in mathscr C are disjoint. If they overlapped, their ratio would differ from the identity by less than16d\_L rho\_T at every archimedean embedding after the bounded g\_p conjugation. Multiplication by T\^{}4 produces an algebraic-integer matrix whose entries have all conjugates below1, hence zero. This contradicts distinctness. No group structure on mathscr C is required for this disjoint-chart argument.

All derivative moments in K061(CG8--12) are unchanged. Its chart constant contains rho\^{}-2d\_L, so shrinking rho by T\^{}4 multiplies the complete sampling bound by T\^{}8d\_L, precisely as displayed. Its convolution proof therefore gives, for R=9/2 or5 respectively and B=2R+1/100,

\[
 \sum_{\substack{g\in Z\backslash\mathscr C\\
                   \max_\tau\log\|g_\tau\|_{op}>B}}
                       |c(g)|
  <2K_TG_0\delta<2^{-2000},
 \quad \delta=f2^{-\ell_+/8}.
 \tag{AT20}
\]

Here G0 is exactly K061(CG10). The actual T\textless2\^{}1001 and d\_L\textless2\^{}23 make the additional exponent at most8008d\_L, still vastly below ell\_+/8. Since \textbar a(g)\textbar\textless=1 and \textbar rho(B\_v)\textbar\textless=1, the same estimate bounds the normalized projected series. The support with bounded real operator norm is finite: T\^{}2g has algebraic-integer entries in a bounded region of every archimedean embedding. This proves absolute convergence of AT15--16 and of their matrix entries. It also supplies a uniform tail estimate for all scalar-orbit bases simultaneously.

\subsection{5. An explicit finite short-sum gate}\label{an-explicit-finite-short-sum-gate}

Use the exact algebraic operator-norm cutoffs

\[
 H=2^{13}\quad(J14),\qquad H=2^{15}\quad(J56).
 \tag{AT21}
\]

The rational lower series for log2 proves log H\textgreater B in both cases. Thus AT20 applies outside this cutoff. For the integral subgroup Gamma of K058, the LIFT42 finite-character proof and packing give the exact bound

\[
 \sum_{\substack{g\in Z\backslash\Gamma,\ g\notin Z\\
                          \|g_\tau\|_{op}\le H\ (\forall\tau)}}
        |a(g)|
 <\frac{2\,28^{a+1}a!\,(225/32)^{af}H^{3af}}
              {Q^{a(f-1)}}<2^{-2000}.
 \tag{AT22}
\]

For the first inequality use LIFT42(FC9--11), e\^{}(4/5)\textless9/4, and sinh(2/5)\textgreater2/5. Unlike an exponential cutoff, H allows an exact algebraic comparison with every candidate matrix. Integer comparisons give margins exceeding2,438 bits on J14 and620,304 bits on J56. The certificate checks the stated weaker2,000-bit bound directly.

If a local projective action is the identity, it fixes the included standard Lambda by AT10. Thus a correspondence scalar on every pencil belongs to the original integral Gamma. Conversely integral matrices need not be scalar on the pencils, but their local normalized factors are still bounded by1. Hence every noncentral integral contribution inside AT21 is bounded by AT22.

Define the genuinely finite sum

\[
 S_T=\sum_{\substack{g\in Z\backslash\mathscr C,\ g\notin\Gamma\\
                    \|g_\tau\|_{op}\le H\ (\forall\tau)}}
           a(g)c(g)\prod_{v\mid T}\rho(B_v).
 \tag{AT23}
\]

For a nonintegral term, at least one local pencil action is not the identity. If b(g) counts such places, its local multiplier has absolute value at most(2/(T-2))\^{}b(g), with b(g)\textgreater=1. This factor is exact in scope; it is not asserted to dominate the number of correspondences. Terms may have either sign or phase, as the example AT12 shows.

The identity coset contributes1. Pairing inverse matrices shows the full sums and their integral/nonintegral cutoff parts are real: both F(B) and C(B) are invariant under inverse, while a(g)c(g) is conjugated. Combining AT16 and AT20--23 yields

\[
 \boxed{\quad
 \frac{E_{x,\mathrm{ram}}}
      {w_x z_0 C_\infty d_Rd_W\bigl((T-2)/(T-1)\bigr)^f}
         =1+S_T+\varepsilon_T,
                  \qquad |\varepsilon_T|<2^{-1999}.
 \quad}\tag{AT24}
\]

In particular the explicit sufficient finite condition

\[
                         S_T>-1+2^{-1999}
 \tag{AT25}
\]

would prove a nonzero actual all-T projected selected restriction in the genuine G-good carrier. The Hcal averaging and finite/compact Schur argument of K058 commute with every T projector. They then extract an actual G-generic scalar bottom vector with nonzero Mram restriction, retaining their dimension and metric factors. The checked DO055/K059 source bridge would turn that into a nonzero prescribed coupled-image terminal pairing. None of these implications supplies AT25 itself.

The set in AT23 is explicitly finite and checkable. Enumerate the algebraic-integer matrices A=T\^{}2g with every conjugate entry bounded by T\^{}2H, impose g*Jg=J and the exact finite lattice conditions, and test the operator norms against the algebraic integer H. The coefficient c(g) is the actual selected Hilbert matrix coefficient, with its PBW normalization; it must not be replaced by a highest-vector depth or by an ambient finite Hilbert pairing. This report does not enumerate that immense set or assert a sharp cost for doing so.

The trace c(g) is independent of the choice of unitary extension of the fixed negative column z: two extensions differ on the right by the compact positive U(a), which preserves W and conjugates the compressed operator. The local pencils are equally explicit from LIFT44's relation v=-v0*J-varpi m modulo varpi: their M, N and the two omitted lines depend on the given v0 and its actual uniformizer frame. Thus AT23 has not introduced an arbitrary real extension or an unspecified replacement pencil.

\subsection{6. Scope and checks}\label{scope-and-checks}

The actual all-base projection has been reduced to AT23, with all original T factors and orbit masses. The favorable point has only been used to choose one fixed rational transport and to control the common real centre in the infinite tail. No uniform near-origin assertion for the individual scalar points has been made.

The remaining finite nonintegral sum is a material gate. Its local weights have cancellation, and an O(1/T) factor at one place is not by itself a bound on the whole sum. The original integral finite character argument cannot be applied directly to matrices with T denominators. Full Mram nonvanishing, a controlled integral detector, and the final mixed DP44 budget therefore remain unproved here.

Source locators used: LIFT43(TG1--6) and LIFT45(SP1--8) for the literal norm/projector and mass; \nolinkurl{ScaledRamifiedLattice.lean84–91,155–175} and Record.lean1811--1848 for its source meaning; LIFT41(NP11), LIFT44(PR4--7) and LIFT43(TG2--3) for the actual scalar pencil; K058(PG1--32), K061(CG1--26) and R030 for the selected kernel and Haar; LIFT42(FC5--13) for the exact finite-character/packing input.

The lattice incidence, projector traces and averaging identity are proved here. The companion certificate checks all projective linear actions over several finite fields, the omitted endpoints, the actual tensor projector and mass factors in finite models, the denominator sampling margin and the two actual algebraic-cutoff bounds. It computes neither AT23 nor an automorphic period.

\subsubsection{Primary references for retained analytic inputs}\label{primary-references-for-retained-analytic-inputs}

Green, J. A. (1955). The characters of the finite general linear groups. \emph{Transactions of the American Mathematical Society, 80}(2), 402--447. Section5, Example(ii), pp.430--431. \url{https://www.ams.org/tran/1955-080-02/S0002-9947-1955-0072878-2/S0002-9947-1955-0072878-2.pdf}

Nakahama, R. (2016). Norm computation and analytic continuation of vector-valued holomorphic discrete series representations. \emph{Journal of Lie Theory, 26}, 927--990. Theorem5.1, pp.960--961. \url{https://jolt.centre-mersenne.org/item/10.5802/jolt.919.pdf}

\subsubsection{Author whole-proof check}\label{author-whole-proof-check-4}

The entire saved AT1--25 and complete companion certificate were read and reconstructed. This includes the fixed rational transport, both omitted lattice points, the actual finite spherical factors, the associated-bundle versus scalar generic metric, the central/finite Haar normalization, strict/projective class masses, denominator-two support and denominator-four pair ratios, the complete infinite tail, and the remaining signed finite sum. The certificate passed20,639 exact assertions at2026-11-22T08:54:23.068459+00:00. The rational short sum S\_T was not computed and AT25 was not asserted.

\clearpage

\section{ALP-LIFT-47 --- transference bounds for the actual adapted Hermitian frames}\label{alp-lift-47-transference-bounds-for-the-actual-adapted-hermitian-frames}

Author: GPT-6 Astra Pro, ALP\_MACHINE /root/lift\_accounting. 22 November 2026 UTC. Inherited reasoning effort is not separately exposed.

Status: author whole-proof check PASS, 2026-11-22T09:13:07.852560+00:00. All TF1--21/eight sections and the complete companion script were read. Independent review is recorded separately.

This finding improves the actual integral field-spanning construction of AM-DO-057. It preserves its literal initial and terminal coefficient lattices, its prescribed residue frames at eleven, and its independent determinant-line sections. It changes neither the class representatives nor any theta kernel. The improvement comes from bounding the minimum of the Euclidean dual by the codifferent and applying the covering-radius transference theorem directly to each prescribed residue class. No auxiliary short integral basis of the CM field is needed.

The resulting polynomial-degree cost coefficients are 27 and 31 on the J14 host, and 36 and 35 on the J56 host, for the initial and terminal families respectively. Their corresponding determinant-degree coefficients are 13 and 17. The exact signed cyclic-kernel factor from AM-R-039 remains separate. This is a bound on genuine integral field-spanning generators, not an integral basis, a saturated theta image, or a proof that the actual Mram period is nonzero.

\subsection{1. Fixed conventions and accepted inputs}\label{fixed-conventions-and-accepted-inputs-1}

Let E be the actual CM field, F its maximal real subfield, and \([E:\mathbb Q]=2f\). Write \(\Delta_E=|\operatorname{disc}(E)|\), \(\operatorname{rd}(E)=\Delta_E^{1/(2f)}\), and \(\mathfrak D_E\) for the different of E over Q. The coefficient-value field is separate. Let \(\mathcal N\) be an integral unimodular Hermitian \(O_E\)-lattice of rank n in the standard positive definite Hermitian space. It may be projective and need not be free.

Use exactly the trace norm and underlying Euclidean pairing of DO057(BF5--6):

\[
 \|x\|_{\mathrm{tr}}^2=\operatorname{Tr}_{F/\mathbb Q}h(x,x),
 \qquad
 \langle x,y\rangle=\tfrac12\operatorname{Tr}_{E/\mathbb Q}h(x,y),
 \qquad d=2fn.
 \tag{TF1}
\]

The latter expression is real because the trace is unchanged by conjugation; equivalently it is the real part of the sum over a CM type. The Euclidean dual lattice, in this normalization, is

\[
                  \mathcal N^*=2\mathfrak D_E^{-1}\mathcal N.
 \tag{TF2}
\]

DO057(BF6) proves TF2 locally, including projective lattices. For clarity, testing the pairing against every \(O_E\)-multiple of a local lattice basis vector says precisely that the Hermitian pairing belongs to \(2\mathfrak D_E^{-1}\). Hermitian unimodularity then gives TF2. There is no replacement of the Euclidean dual by the Hermitian dual.

The other carried inputs are:

\begin{enumerate}
\def\labelenumi{\arabic{enumi}.}
\tightlist
\item
  DO057(BF3--4, BF10--15): the literal initial coefficient lattice is \(\Delta_\lambda(\overline{\mathcal L})\), where \(\mathcal L\) is the primitive T-covector kernel in \(\mathcal N\); the terminal lattice is \(\nabla_\lambda(\mathcal N)\). Integral polynomial transport and separate determinant sections preserve these actual transported lattices.
\item
  AM-R-038: the standard initial Weyl module has an integral field-spanning cyclotomic family of norm at most \(a^{D_\tau}\), using roots of unity already present in the coefficient field. The actual degree condition \(D_\tau<T\) is proved there.
\item
  AM-R-037: the actual standard semistandard dual-Weyl basis has normalized Fischer norms at most one.
\item
  AM-R-039(R39.5--13): the exact cyclic covector-kernel inclusion multiplies the initial coefficient lattice by \(\beta_\lambda=\prod_\tau\tau(\overline\pi_T)^{\lambda_{0,\tau}}\), whose archimedean modulus is \(T^{\Lambda_0/2}\), where \(\Lambda_0=\sum_\tau\lambda_{0,\tau}\). This is a signed exponent.
\item
  K059 and DO056: the fixed stabilizer/projector operations on these generators have the already established finite factor \(2^{K_{\mathrm{fin}}}\), with \(K_{\mathrm{fin}}=112031999\) or \(490748794\). The actual ramified Fourier transform, Haar-volume cancellation and common scalar product formula remain as established there.
\end{enumerate}

These accepted results are cited, not reproved.

\subsection{2. A primewise lower bound for the Euclidean dual minimum}\label{a-primewise-lower-bound-for-the-euclidean-dual-minimum}

\subsubsection{Proposition TF3}\label{proposition-tf3}

Every nonzero \(y\in\mathcal N^*\) satisfies

\[
 \|y\|_{\mathrm{tr}}^2\ge
                  4f\,\operatorname{rd}(E)^{-2}.
 \quad\text{Thus}\quad
 \lambda_1(\mathcal N^*)\ge
                  {2\sqrt f\over\operatorname{rd}(E)}.
 \tag{TF3}
\]

\textbf{Proof.} Complex conjugation preserves the different: the defining trace duality is invariant under every Q-automorphism of E, in particular its complex conjugation. Put \(I=2\mathfrak D_E^{-1}\). By TF2, \(y\in I\mathcal N\). At a finite prime P, express y in a local \(O_{E,P}\)-basis of the projective lattice. Integrality of h and conjugation invariance of the different give

\[
       \alpha:=h(y,y)\in I\overline I
             =4\mathfrak D_E^{-2},\qquad
       v_P(\alpha)\ge 2v_P(2)-2v_P(\mathfrak D_E).
 \tag{TF4}
\]

The calculation uses the conjugate prime in the second coefficient of h: its different exponent is the same after transporting by conjugation. Thus TF4 holds also at ramified primes, including those above two. It is not an assertion that \(\alpha\) is integral.

Since h is positive definite at every embedding, \(\alpha\in F\) is nonzero and totally positive. Taking the norm of the inclusion of fractional ideals \((\alpha)\subseteq4\mathfrak D_E^{-2}\) gives

\[
 |N_{E/\mathbb Q}\alpha|
   \ge {4^{2f}\over\Delta_E^2}.
 \tag{TF5}
\]

Equivalently, multiplying TF4 by \(\log N(P)\) and summing over all finite primes proves TF5. The direction is correct: inclusion of fractional ideals reverses the inequality of their positive norms. Because \(\alpha\in F\), its norm from E is the square of its positive norm from F. Consequently

\[
 N_{F/\mathbb Q}\alpha\ge{4^f\over\Delta_E},
 \qquad
 \operatorname{Tr}_{F/\mathbb Q}\alpha
       \ge f\left({4^f\over\Delta_E}\right)^{1/f}
       =4f\,\operatorname{rd}(E)^{-2}.
 \tag{TF6}
\]

The second inequality is arithmetic--geometric mean applied to the f positive real conjugates. This proves TF3. No discriminant exponent n occurs in TF3: this is a lower bound for one nonzero vector, not for a lattice covolume. ∎

\subsection{3. Direct covering-radius adaptation at eleven}\label{direct-covering-radius-adaptation-at-eleven}

The only new external theorem used is Banaszczyk (1993, Theorems2.1 and2.2, printed pp631--632). For a Euclidean lattice L of real dimension d,

\[
 \lambda_d(L)\lambda_1(L^*)\le d,\qquad
 \mu(L)\lambda_1(L^*)\le d/2,
 \tag{TF7}
\]

where \(\mu\) is its Euclidean covering radius. The paper states strict upper inequalities in dimensions greater than one; the weak forms above are sufficient and include the one-dimensional endpoint. The definitions of the dual, successive minima and covering radius appear on printed pp625--626. Here our dimension \(d=2fn\) is at least two.

Combining TF3 with both parts of TF7 gives

\[
 \lambda_d(\mathcal N)\le n\sqrt f\,\operatorname{rd}(E),
 \qquad
 \mu(\mathcal N)\le
               {n\sqrt f\over2}\operatorname{rd}(E).
 \tag{TF8}
\]

The second inequality will be used directly; going through a short Z-sublattice would unnecessarily enlarge the constant.

\subsubsection{Theorem TF9 --- every prescribed residue frame}\label{theorem-tf9-every-prescribed-residue-frame}

Fix any \(O_E/11O_E\)-basis of \(\mathcal N/11\mathcal N\), in particular the reductions of the already fixed local eleven-adic frames. There exist \(n\) vectors \(y_i\in\mathcal N\) with precisely those residues such that

\[
 \boxed{\quad
  \|y_i\|_{\mathrm{tr}}\le
  \widehat R_n:={11\over2}\,n\sqrt f\,\operatorname{rd}(E).
  \quad}
 \tag{TF9}
\]

They are E-independent; their O-span is an integral sublattice of \(\mathcal N\), with index prime to eleven. The new frame agrees with the specified one modulo eleven, and is an integral invertible frame at every place above eleven. It is not asserted to be a global O-basis.

\textbf{Proof.} Lift the prescribed residue vectors arbitrarily to \(r_i\in\mathcal N\). By the definition of the covering radius, there is \(v_i\in\mathcal N\) such that \(\|r_i/11-v_i\|_{\mathrm{tr}}\le\mu(\mathcal N)\). Let \(y_i=r_i-11v_i\). Then \(y_i\equiv r_i\pmod{11\mathcal N}\) and TF8 gives TF9.

At each prime P above eleven the reductions of the \(y_i\) generate \(\mathcal N_P/11\mathcal N_P\), and hence \(\mathcal N_P/P\mathcal N_P\). Nakayama's lemma makes them a local basis. There are exactly n vectors, so they are E-independent. Their global O-span therefore has finite index in \(\mathcal N\), and the local basis property shows that this index has no factor eleven. ∎

This proof works for a rational prime p with 11 replaced by p.~In the application the value eleven is fixed. It preserves the existing local lattice and level; it does not create a congruence subgroup or change any decoration.

The construction is effective in a finite, possibly expensive sense: given the integral lattice presentation and the positive algebraic metrics, enumerate its vectors in the compact TF9 ball and test the prescribed residues. The theorem guarantees a suitable vector in each residue class. One may instead solve the corresponding closest-vector problems. No polynomial-time algorithm or actual high-rank class enumeration is claimed.

For later bookkeeping, the same Hadamard argument as DO057(BF9) gives, with \(A=(y_1,\ldots,y_n)\),

\[
 [\mathcal N:AO_E^n]\le\widehat R_n^{\,2fn},\quad
 \|A_\tau\|_{\mathrm{op}}\le\sqrt n\,\widehat R_n,\quad
 \|A_\tau^{-1}\|_{\mathrm{op}}
       \le n^{(n-1)/2}\widehat R_n^{\,fn-1}.
 \tag{TF10}
\]

The inverse bound is recorded but is not raised to the signed determinant weight below.

\subsection{4. Determinant lines and the unchanged integral coefficient models}\label{determinant-lines-and-the-unchanged-integral-coefficient-models}

The determinant lattice \(\det\mathcal N\), with its induced Hermitian metric, and its linear dual with the dual metric are positive integral unimodular rank-one Hermitian lattices. Apply TF9 to each, prescribing their fixed nonzero residue bases. This gives nonzero sections

\[
 \delta_+\in\det\mathcal N,\qquad
 \delta_-\in(\det\mathcal N)^\vee,\qquad
 |\tau(\delta_\pm)|\le\widehat R_1
              ={11\over2}\sqrt f\,\operatorname{rd}(E).
 \tag{TF11}
\]

Their sizes are the induced Hermitian sizes of line sections, without choosing a globally principal determinant ideal. Each is primitive at eleven. Their product need not be one. For signed m, use \(\delta_+^m\) if \(m\ge0\), and \(\delta_-^{-m}\) otherwise. The integral norm bound is \(\widehat R_1^{|m|}\).

For each actual dominant weight write

\[
 m_\tau=\lambda_{n-1,\tau},\quad
 \nu_{i,\tau}=\lambda_{i,\tau}-m_\tau,\quad
 D_\tau=\sum_i\nu_{i,\tau}.
 \tag{TF12}
\]

Transport the accepted standard polynomial Weyl or dual-Weyl lattice by the frame A, and its signed determinant factor by the separate section TF11. DO057(BF11--12) proves that this lands in the original transported integral coefficient lattice and has archimedean operator bound

\[
         (\sqrt n\,\widehat R_n)^{D_\tau}
                       \widehat R_1^{|m_\tau|}.
 \tag{TF13}
\]

At eleven both the polynomial frame and determinant section are integral invertible. The statement holds for Delta and Nabla separately; no identification of these lattices or integral orthogonal projection is being made.

For the actual initial class first apply this construction to \(\overline{\mathcal N}\). Then multiply the full weight tensor by the accepted cyclic scalar of R039:

\[
 \beta_\lambda=
      \prod_\tau\tau(\overline\pi_T)^{\lambda_{0,\tau}},
 \qquad |\sigma\beta_\lambda|=T^{\Lambda_0/2},
 \quad\Lambda_0=\sum_\tau\lambda_{0,\tau}.
 \tag{TF14}
\]

R039 places the result in the literal \(\Delta_\lambda(\overline{\mathcal L})=\mathrm{delAtr}(\mathcal L^\sharp)\). The multiplier is a unit at eleven. Its exponent is signed; when it is negative, its meaning is membership in the relevant fractional coefficient lattice, not integrality as an isolated scalar in an untwisted line.

\subsection{5. Actual initial and terminal generator costs}\label{actual-initial-and-terminal-generator-costs}

Let \(a\) be the initial rank and \(N\) the terminal rank. Put \(D_{0,\Sigma}=\sum_\tau D_{0,\tau}\), \(D_{T,\Sigma}=\sum_\tau D_{T,\tau}\), and similarly let \(M_{0,\Sigma}\) and \(M_{T,\Sigma}\) sum the absolute signed determinant weights.

Use the R038 standard initial Weyl family, of norm at most \(a^{D_{0,\tau}}\). On every actual initial class the construction above gives an integral field-spanning family in the original delAtr lattice with bound

\[
 \boxed{\quad
 \widehat B_0
   =a^{3D_{0,\Sigma}/2}\widehat R_a^{D_{0,\Sigma}}
       \widehat R_1^{M_{0,\Sigma}}T^{\Lambda_0/2}.
 \quad}
 \tag{TF15}
\]

Use the R037 standard terminal semistandard dual-Weyl family, of norms at most one. On every terminal class it gives an integral field-spanning family in the original intNab lattice with bound

\[
 \boxed{\quad
 \widehat B_T
     =(\sqrt N\,\widehat R_N)^{D_{T,\Sigma}}
                       \widehat R_1^{M_{T,\Sigma}}.
 \quad}
 \tag{TF16}
\]

These families use no new coefficient field. Their mod-eleven frame changes and determinant factors are units, so this improvement adds no eleven-adic denominator or clearing factor. They need not span the full integral lattice globally. Tensoring their field spans, and then applying the unchanged finite projectors and stabilizer sums, gives exactly the same field-spanning subspaces as in DO057(BF18).

\subsection{6. Exact bounds on the two actual hosts}\label{exact-bounds-on-the-two-actual-hosts}

The accepted K030/K032 discriminants and ranks are

\[
 \begin{array}{c|c|c|c|l}
  \text{host}&f&a&N&\Delta_E\\ \hline
  J14&42&42&3561&2^{84}7^{70}29^{72}\\
  J56&168&168&3435&2^{336}7^{280}29^{288}5^{168}13^{168}.
 \end{array}
 \tag{TF17}
\]

The following comparisons use these exact integers, not floating-point logarithms:

\begin{longtable}[]{@{}
  >{\raggedright\arraybackslash}p{(\columnwidth - 10\tabcolsep) * \real{0.1304}}
  >{\raggedleft\arraybackslash}p{(\columnwidth - 10\tabcolsep) * \real{0.1739}}
  >{\raggedleft\arraybackslash}p{(\columnwidth - 10\tabcolsep) * \real{0.1739}}
  >{\raggedleft\arraybackslash}p{(\columnwidth - 10\tabcolsep) * \real{0.1739}}
  >{\raggedleft\arraybackslash}p{(\columnwidth - 10\tabcolsep) * \real{0.1739}}
  >{\raggedleft\arraybackslash}p{(\columnwidth - 10\tabcolsep) * \real{0.1739}}@{}}
\toprule\noalign{}
\begin{minipage}[b]{\linewidth}\raggedright
Host
\end{minipage} & \begin{minipage}[b]{\linewidth}\raggedleft
\(\widehat R_1\)
\end{minipage} & \begin{minipage}[b]{\linewidth}\raggedleft
\(\widehat R_a\)
\end{minipage} & \begin{minipage}[b]{\linewidth}\raggedleft
\(\widehat R_N\)
\end{minipage} & \begin{minipage}[b]{\linewidth}\raggedleft
\(a^{3/2}\widehat R_a\)
\end{minipage} & \begin{minipage}[b]{\linewidth}\raggedleft
\(\sqrt N\widehat R_N\)
\end{minipage} \\
\midrule\noalign{}
\endhead
\bottomrule\noalign{}
\endlastfoot
J14 & \(<2^{13}\) & \(<2^{19}\) & \(<2^{25}\) & \(<2^{27}\) & \(<2^{31}\) \\
J56 & \(<2^{17}\) & \(<2^{25}\) & \(<2^{29}\) & \(<2^{36}\) & \(<2^{35}\) \\
\end{longtable}

For example, the assertion \(n^{j/2}\widehat R_n<2^C\) is equivalent to the exact positive-integer comparison

\[
              (121 f n^{j+2})^f\Delta_E
                        <2^{\,2f(C+1)}.
 \tag{TF18}
\]

Apply this with \(j=0\) for a frame radius, \(j=3\) for the initial polynomial coefficient, and \(j=1\) for the terminal one. The companion certificate also checks that the preceding power-of-two bound fails in each column, so the displayed integers are the least such integer exponents for these particular formulas. No optimality of the underlying frame construction is claimed.

Thus the exact generator bounds imply

\[
 \begin{array}{c|l|l}
 &\log_2\widehat B_0\ \text{upper bound}
 &\log_2\widehat B_T\ \text{upper bound}\\ \hline
 J14&
 27D_{0,\Sigma}+13M_{0,\Sigma}
                  +\tfrac12\Lambda_0\log_2T&
 31D_{T,\Sigma}+13M_{T,\Sigma}\\
 J56&
 36D_{0,\Sigma}+17M_{0,\Sigma}
                  +\tfrac12\Lambda_0\log_2T&
 35D_{T,\Sigma}+17M_{T,\Sigma}.
 \end{array}
 \tag{TF19}
\]

The signed T term remains exact in TF19. It is not replaced by an upper bound for T when its coefficient is negative.

Compared with the already checked DO057 baseline, the polynomial coefficients decrease from \((324,35643)\) to \((27,31)\) on J14, and from \((1637,39542)\) to \((36,35)\) on J56. The determinant coefficients decrease from 19 to 13 and from 27 to 17. Both comparisons use the same R039 cyclic kernel factor; the improvement is entirely in the actual frame costs.

\subsection{7. The resulting integral-witness bound and remaining gate}\label{the-resulting-integral-witness-bound-and-remaining-gate}

Apply the fixed K059 ramified projector and stabilizer/projector sums exactly as in DO057(BF18). If the actual moved Mram pairing is nonzero, some pair of these genuine integral field-spanning generators detects it. Its product of Hilbert norms has the bound

\[
        \|x\|\,\|z\|
             \le 2^{K_{\mathrm{fin}}}\widehat B_0\widehat B_T.
 \tag{TF20}
\]

In particular its binary logarithm is bounded above by

\[
 \begin{cases}
 112031999+27D_{0,\Sigma}+31D_{T,\Sigma}
       +13(M_{0,\Sigma}+M_{T,\Sigma})
       +\tfrac12\Lambda_0\log_2T,&J14,\\
 490748794+36D_{0,\Sigma}+35D_{T,\Sigma}
       +17(M_{0,\Sigma}+M_{T,\Sigma})
       +\tfrac12\Lambda_0\log_2T,&J56.
 \end{cases}
 \tag{TF21}
\]

The input hypothesis is actual nonvanishing. A field-spanning family detects a nonzero bilinear form by elementary linear algebra, without requiring an integral inverse or a bound for the number of generators. This does not establish that hypothesis.

The common scalar product-formula cancellation, normalized vacuum/period factor, actual class masses, and theta operator norm in DO056 remain unchanged. TF20 improves the sufficient upper bound for the integral-input cost. It is not a complete DP44 inequality, a literal Host instance, a shorter chain, an image-saturation assertion, or a claim that the source and target integral lattices have index one.

\subsection{8. Sources and verification scope}\label{sources-and-verification-scope-5}

\textbf{Primary new source.} Banaszczyk, W. (1993). New bounds in some transference theorems in the geometry of numbers. \emph{Mathematische Annalen, 296}, 625--635. \url{https://doi.org/10.1007/BF01445125}. The Euclidean definitions are on pp625--626, Theorem2.1 on p631, and Theorem2.2 with its proof on p632. The original-paper transcription was directly read at \url{https://paperzz.com/doc/7332762/new-bounds-in-some-transference-theorems-in-the-geometry-of}, retrieval turn355565772view0 and focused retrievals turn875992504view0--1. Its OCR renders the fraction one-half as ``89''; the displayed balancing inequalities in the original proof are both \(2^{-1/2}\sqrt d\), and the dimension-one case is one-half. We use the precise bound \(d/2\). The publisher metadata was independently checked at turn737650521view0.

\textbf{Carried source-aligned results.} AM-DO-057(BF1--18), AM-R-037, AM-R-038(R38.9), AM-R-039(R39.5--13), K059 and AM-DO-056. The complete source lattice audit is in DO057 and its independent review; this finding changes only its quantitative frame construction.

The companion script tests the exact integer inequalities TF18 and the resulting cost coefficients, primewise fractional-ideal norm bookkeeping, and finite Gaussian-integer residue-frame diagnostics. These diagnostics do not prove Banaszczyk's theorem, enumerate the actual high-rank classes, compute an actual automorphic period, or certify theta-image saturation. The proof of the new general assertions is TF3--TF16 above.

\clearpage

\section{AM-DO-058 --- the smallest denominator stratum in the actual all-T sum}\label{am-do-058-the-smallest-denominator-stratum-in-the-actual-all-t-sum}

Author: GPT-6 Astra Pro, ALP\_MACHINE \nolinkurl{/root/definite_oscillator}. 22 November 2026 UTC. Inherited reasoning effort is not separately exposed.

\textbf{Author whole-proof check PASS --- 2026-11-22T09:20:36.728836+00:00.}

This report bounds an actual part of K062's finite signed sum, without changing its class orbit, finite representation, scalar projection, or real selected vector. On the J56 host the entire stratum whose entry-denominator ideal has norm T squared contributes less than \nolinkurl{2^-101768} in absolute value. The remaining higher-denominator terms are not bounded by that conclusion. On J14 the same argument does not give a useful bound even for this smallest stratum.

The proof also gives an exact finite-stabilizer bound from the two actual G primes and a general denominator-stratum estimate. These estimates count all possible source and target pencil lattices. No favorable-base substitution, positivity of individual projector traces, or class-mass cancellation is used.

\subsection{1. The fixed normalized sum and denominator invariant}\label{the-fixed-normalized-sum-and-denominator-invariant}

Keep all hypotheses, genuine finite representations, selected real types, weights, and the fixed scalar orbit of K062. In particular its archimedean cutoff is

\[
 (a,f,n,H)=
 \begin{cases}
 (42,42,43,2^{13}),&J14,\\
 (168,168,169,2^{15}),&J56,
 \end{cases}
 \qquad n=a+1.
 \tag{DS1}
\]

The group is the actual \nolinkurl{U(I_a,-1)} over F, with CM field E of degree 2f. The prime T is completely split in E. At each of its f selected relative places, K062(AT3--6) gives the same fixed family of m=T−1 lattices \nolinkurl{L_c}, with \texttt{c\ in\ D=F\_T\ minus\ \{-1\}}. Each is obtained from the standard lattice by the displayed determinant-one unipotent, so every \nolinkurl{L_c} has the same local determinant ideal as the standard lattice. The conjugate CM component is always the inverse-adjoint component; it is not a second independent scalar choice.

K062's normalized nonintegral short sum is

\[
 S_T=\sum_{g\in Z\backslash\mathscr C_H,\ g\notin\Gamma}
       a(g)c(g)\prod_{v\mid T}\rho(g_v).
 \tag{DS2}
\]

Here \nolinkurl{mathscr C_H} is its actual support, including maximal compactness away T, the existence of a source/target pencil pair at every T place, and \nolinkurl{||g_tau||_op <= H} at every real label. The factors are exactly K062(AT14,AT23): a is the normalized genuine finite G partial trace, c is the normalized selected real trace, and rho is the signed \nolinkurl{(F-C/m)/(T-2)} pencil trace. In particular \nolinkurl{|c(g)|<=1}. The quotient Z is the existing scalar group of order 28.

For such a g define its entry-denominator ideal in the fixed standard coordinates by

\[
 \mathfrak D(g)=\prod_{\mathfrak P}
   \mathfrak P^{d_{\mathfrak P}(g)},\qquad
 d_{\mathfrak P}(g)=
      \max\{0,-\min_{i,j}v_{\mathfrak P}(g_{ij})\}.
 \tag{DS3}
\]

Only primes above T occur. K062(AT18) proves that every exponent is at most two, both for g and its inverse. Let s(g) be the number of relative T places at which g is not integral in the standard coordinates. Then

\[
 N\mathfrak D(g)=T^{e(g)},\qquad
                  2s(g)\le e(g)\le4s(g).
 \tag{DS4}
\]

\textbf{Proof of the paired lower bound.} A supported local g maps one L\_c to another. Their determinant ideals are equal, hence \nolinkurl{v(det g)=0}. An integral square matrix with unit determinant has integral inverse. The conjugate CM component of g is its inverse adjoint, with the standard unit entries of the Hermitian matrix. Consequently

\[
 d_{\mathfrak P}=0\iff d_{\overline{\mathfrak P}}=0,
 \quad
 (d_{\mathfrak P},d_{\overline{\mathfrak P}})
     \in\{(1,1),(1,2),(2,1),(2,2)\}
 \tag{DS5}
\]

at a nonintegral relative place. All these prime ideals have norm T, which proves DS4. In particular the smallest nonzero e is two. If e=2, there is exactly one nonintegral relative place and \nolinkurl{mathfrak D(g)=P bar(P)} there. This conclusion uses the actual determinant condition, not just the entrywise denominator-two bound. ∎

Write \texttt{S\_\{s,e\}} for the part of DS2 with these two invariants. Since inversion exchanges conjugate denominator exponents and preserves the cutoff, each such stratum is inverse closed. Its sum is real by the same conjugation identity used in K062; the estimates below nevertheless bound the absolute values before any cancellation.

\subsection{2. Finite G suppression with the actual fractional denominator}\label{finite-g-suppression-with-the-actual-fractional-denominator}

For every nonintegral supported g,

\[
 \boxed{\quad
 |a(g)|\le \min\left\{1,
  {2(2H)^{af}T^{ae(g)/2}\over Q^{a(f-1)}}\right\},
                       \qquad Q=g_0g_1.
 \quad}
 \tag{DS6}
\]

\textbf{Proof.} Such a g is not a scalar matrix. Indeed a scalar supported matrix has valuation zero at all T primes by \nolinkurl{v(det g)=0}, and is already integral away T; it would belong to Gamma. Thus choose a nonzero entry x off the diagonal, or a nonzero difference of two diagonal entries. Every complex conjugate of x has absolute value at most 2H and \nolinkurl{x D(g)} is an integral ideal.

At any selected G relative place where g reduces to a scalar, its conjugate reduction is scalar too by the inverse-adjoint relation. Therefore x vanishes at both corresponding degree-one CM primes. These are prime to the denominator ideal. Taking norms gives

\[
 \prod_{v\in G:\,\operatorname{red}_v(g)\ \mathrm{scalar}}q_v^2
 \le |N_{E/\mathbb Q}x|\,N\mathfrak D(g)
 \le (2H)^{2f}T^{e(g)}.
 \tag{DS7}
\]

This is the fractional-ideal version of LIFT42(FC8), with its direction and its full denominator retained. Thus the total logarithmic weight of scalar G places is at most \nolinkurl{f log(2H)+(e/2)log T}.

At all but the two mixed places, the fixed genuine finite factor is cuspidal of degree n.~Use the already checked Green-character ratio in LIFT42(FC5): at a nonscalar residue the ratio is at most \nolinkurl{1/(q^a-1)}. At the two mixed factors the normalized partial trace has absolute value at most one. Removing those two places costs log Q, and the same correction product as LIFT42(FC9) is less than two. Substitution of DS7 gives DS6. No new finite-character formula or change of e0 is used. ∎

For orientation, exact integer comparisons give these pointwise bounds for the expression on the right before taking its minimum with one:

\begin{longtable}[]{@{}lrr@{}}
\toprule\noalign{}
Host & e & Strict upper bound \\
\midrule\noalign{}
\endhead
\bottomrule\noalign{}
\endlastfoot
J14 & 2 & \nolinkurl{2^-9847} \\
J14 & 3 & \nolinkurl{2^11169} \\
J56 & 2 & \nolinkurl{2^-1352799} \\
J56 & 4 & \nolinkurl{2^-1185564} \\
J56 & 18 & \nolinkurl{2^-14924} \\
J56 & 19 & \nolinkurl{2^68693} \\
\end{longtable}

These are individual-term bounds, not bounds for their multiplicities. The companion certificate retains the full e=0 through 20 table.

\subsection{3. An exact common bound for the compact isometry fibres}\label{an-exact-common-bound-for-the-compact-isometry-fibres}

Let S be any finite subgroup of the actual rational unitary group whose entries are integral at the two G primes. Then

\[
 |S|\mid
 G_n:=\gcd\bigl(|\operatorname{GL}_n(\mathbb F_{g_0})|,
               |\operatorname{GL}_n(\mathbb F_{g_1})|\bigr).
 \tag{DS8}
\]

The two actual primes are completely split in E. Reduction at one fixed CM place over each is sufficient. The exact results are

\begin{longtable}[]{@{}lrrrr@{}}
\toprule\noalign{}
Host & n & g0 & g1 & Binary length of G\_n \\
\midrule\noalign{}
\endhead
\bottomrule\noalign{}
\endlastfoot
J14 & 43 & 4,615,409 & 5,276,377 & 814 \\
J56 & 169 & 22,285,932,761 & 64,072,175,441 & 5,022 \\
\end{longtable}

\textbf{Proof of injectivity.} At either residue prime q, the kernel of reduction on integral matrices is a pro-q group. One can see this directly: its successive congruence quotients are subgroups of the additive matrix space over the residue field, and the intersection of the congruence kernels is trivial. A finite subgroup in that kernel is therefore a q-group. If nontrivial it contains an element of order q.

Over characteristic zero, a nonidentity matrix of order q has a primitive qth root of unity among its eigenvalues. Its minimal polynomial over E has degree at least

\[
                 [E(\zeta_q):E]\ge {q-1\over2f}>n.
 \tag{DS9}
\]

The strict final inequality holds for both actual primes on both hosts, as checked by exact integers. Such a polynomial cannot divide the characteristic polynomial of an n-by-n matrix. Thus reduction is injective. Lagrange's theorem at both primes proves DS8.

The exact order formula used to compute the gcd is

\[
 |\operatorname{GL}_n(\mathbb F_q)|
          =q^{n(n-1)/2}\prod_{i=1}^n(q^i-1).
 \tag{DS10}
\]

The certificate stores the exact gcd integers, their hashes, the order bit lengths, and both divisibilities. No factorization of these large group orders is required. The bound concerns the actual finite subgroup; it does not replace a class mass by its reciprocal or remove stabilizers. ∎

\subsection{4. Counting last columns instead of whole matrices}\label{counting-last-columns-instead-of-whole-matrices}

Fix a denominator ideal D supported above T with \nolinkurl{ND=T^e}. Every nonzero x in \nolinkurl{D^-1} satisfies

\[
 |N_{E/\mathbb Q}x|\ge T^{-e},\qquad
           \sum_{\tau}|\tau x|^2\ge fT^{-e/f}.
 \tag{DS11}
\]

The first statement follows because \nolinkurl{(x)D} is an integral nonzero ideal; the second is arithmetic--geometric mean. It follows that any two distinct vectors in \nolinkurl{(D^-1)^n} have Euclidean trace distance at least \nolinkurl{sqrt(f) T^(-e/(2f))}: apply DS11 to one nonzero coordinate difference.

For a matrix in the real cutoff, its last column v satisfies \nolinkurl{sum_tau ||v_tau||^2 <= fH^2}. Disjoint Euclidean balls of half the above separation radius around all such columns lie inside the ball enlarged by that radius. Comparing volumes in real dimension 2fn gives

\[
 \#\{\hbox{possible last columns}}
 \le\left(1+2HT^{e/(2f)}\right)^{2fn}
 \le (2H+1)^{2fn}T^{ne}.
 \tag{DS12}
\]

This is a direct packing proof. It does not assume a reduced lattice basis, a principal denominator ideal, or a global integral frame for a modified class. It overcounts by ignoring the unit negative Hermitian norm of the last column.

Now fix a set of s nonintegral relative places and a choice of denominator exponents satisfying DS5. For each possible ordered pair \nolinkurl{(c_v,d_v)} of pencil points at those places, form the two global O\_E lattices with these local lattices and the standard lattice elsewhere. They exist as the intersection of their adelic lattices with E\^{}n.~Any g in this portion of the support is an isometry from one of these source lattices to its target. There are at most \nolinkurl{m^(2s)} choices of the ordered pencil pairs. At every other T place g is integral and hence preserves the standard lattice; therefore no further pencil choice is needed for this upper bound.

Within one such isometry torsor, fix one possible last column. If g and g* have that same column, then \nolinkurl{gamma=(g*)^-1g} preserves the source lattice and fixes the standard negative vector e\_n.~All such gammas lie in the group of source-lattice automorphisms fixing e\_n, which is finite. Indeed fixing e\_n makes its real action compact on the positive orthogonal complement at every label, while preservation of the fixed finite lattices bounds all finite denominators; their entries are therefore a finite set of algebraic numbers in a bounded Minkowski region. They are integral at G because all the source lattices there are the unchanged standard lattice. DS8 bounds the group order by G\_n.

Consequently the number of matrices for this fixed denominator pattern is at most

\[
               m^{2s}G_n(2H+1)^{2fn}T^{ne}.
 \tag{DS13}
\]

Counting matrices instead of Z-cosets only increases the bound. A matrix may realize more than one ordered pencil pair, which also only increases this union bound. Every actual class/frame compatibility is retained in the set being bounded; none is inferred from an arbitrary local matrix.

There are \nolinkurl{binomial(f,s)} choices of the nonintegral relative places. For a fixed such set, the exact number of possible denominator patterns of total exponent e is

\[
 [X^e](X^2+2X^3+X^4)^s
                  =\binom{2s}{e-2s}.
 \tag{DS14}
\]

The two distinct orientations at exponent three are both counted. A binomial coefficient outside its natural range means zero.

\subsection{5. The general stratum inequality and its actual specialization}\label{the-general-stratum-inequality-and-its-actual-specialization}

At each nonintegral relative place the pencil action cannot be the identity: that would fix the included standard lattice. K062(AT11) therefore gives \nolinkurl{|rho|<=2/(T-2)} there. At the remaining places use \nolinkurl{|rho|<=1}. Combining DS6 and DS13--14 proves

\[
\begin{split}
 |S_{s,e}|\le{}&
 \binom fs\binom{2s}{e-2s}
       m^{2s}G_n(2H+1)^{2fn}T^{ne}
       \left({2\over T-2}\right)^s\\
 &\quad\times
       \min\left\{1,
       {2(2H)^{af}T^{ae/2}\over Q^{a(f-1)}}\right\}.
\end{split}
 \tag{DS15}
\]

This is an absolute estimate for the actual signed sum. In particular it does not replace the trace \nolinkurl{F-C/m} by F, assume positivity of a(g)c(g), or omit the two noncuspidal finite factors.

For the smallest nonzero denominator exponent, \nolinkurl{e=2} forces \nolinkurl{s=1}. The untruncated character bound in DS6 is less than one on both hosts. Thus the entire smallest stratum satisfies

\[
 \boxed{\quad
 |S_{1,2}|\le
 {4f(T-1)^2G_n(2H)^{af}(2H+1)^{2fn}T^{a+2n}
          \over (T-2)Q^{a(f-1)}}.
 \quad}
 \tag{DS16}
\]

All quantities in DS16 are actual fixed integers. The certificate computes its numerator A and denominator B and verifies the exact comparisons \nolinkurl{2^k A < B <= 2^(k+1) A}, with the inequalities interpreted by integer shifts when k is negative. The results are

\[
 \begin{array}{c|r|r|r}
 &\operatorname{bits}(A)&\operatorname{bits}(B)&
  \max\{k\in\mathbb Z:A/B<2^{-k}\}\\ \hline
 J14&206183&77577&-128606\\
 J56&1870846&1972614&101768.
 \end{array}
 \tag{DS17}
\]

Therefore

\[
                \boxed{|S_{1,2}|<2^{-101768}\quad(J56).}
 \tag{DS18}
\]

On J14 the bound in DS16 is greater than one. DS17 makes that failure explicit; it is not an obstruction to positivity or to nonzero Mram. The estimate DS15 likewise is not asserted to control all higher strata.

\subsection{6. Exact consequence for the remaining finite gate}\label{exact-consequence-for-the-remaining-finite-gate}

On J56 decompose the actual finite sum as

\[
              S_T=S_{1,2}+S_{\ge3},
 \tag{DS19}
\]

where the second sum retains every actual supported nonintegral matrix of denominator exponent at least three, with its original factor \nolinkurl{a(g)c(g) product rho(g_v)}. By DS4 no other nonintegral case exists. K062(AT24) becomes

\[
 {E_{x,\mathrm{ram}}\over
       w_x\,28C_\infty d_Rd_W((T-2)/(T-1))^f}
       =1+S_{\ge3}+\varepsilon'_T,
 \qquad
 |\varepsilon'_T|<2^{-1999}+2^{-101768}.
 \tag{DS20}
\]

The original class weight, strict scalar-orbit denominator, genuine finite projector, and selected real normalization are visibly unchanged. For example \texttt{S\_\{\textgreater{}=3\}\textgreater{}-1+2\^{}-1998} is a sufficient condition for positivity. This report does not establish it. On J14 the original K062 gate remains.

The main scope distinction is between three true assertions: individual finite-character suppression DS6, an exact multiplicity upper bound DS15, and the actual small total DS18 for one host and one complete stratum. None of them alone establishes the sign of the entire finite short sum. The local spherical Gram positivity carried in K052 also cannot be substituted for the actual automorphic selected restriction in DS20.

\subsection{7. Carried sources and verification}\label{carried-sources-and-verification}

\begin{itemize}
\tightlist
\item
  \textbf{AM-K062}, \nolinkurl{averaged_scalar_pencil_trace.md}, AT3--8, AT11, AT14--18, AT21--24: actual pencil, signed trace, finite support, cutoff, class weights and normalized sum. Independently whole-reviewed by this agent in \nolinkurl{definite_oscillator_review_k062.md}.
\item
  \textbf{ALP-LIFT-42}, \nolinkurl{finite_character_short_orbit_bound.md}, FC1--2, FC5--9: actual G factors, checked Green-character ratio and product correction. DS7 proves the new fractional-denominator refinement.
\item
  \textbf{AM-DO-040}, \nolinkurl{definite_oscillator_large_t_certificate.json}: actual certified T, g0, g1 and host data. The complete splitting and actual field degrees are the accepted K031/K032 host facts.
\end{itemize}

The new finite-group injection, Euclidean column packing, torsor-fibre bound and denominator-stratum estimates are proved in full above. They invoke no new unproved automorphic input.

The companion \nolinkurl{definite_oscillator_minimal_denominator_certificate.py} checks the two exact group-order gcds, the actual cyclotomic-degree inequalities, all binary bounds, the denominator-pattern coefficients, and rational elementary-matrix tests of the actual local isometry-torsor and paired-pole constraints. It does not enumerate the actual S\_T or compute its real matrix coefficients, and is not a Lean proof.

\subsection{Author verification record}\label{author-verification-record-7}

At 2026-11-22T09:20:36.728836+00:00, I read the complete final proof DS1--20 and all seven sections, plus the complete certificate, and reconstructed the fractional-ideal inequality, paired pole invariant, finite reduction injection, Euclidean last-column count, compact isometry-fibre bound, and actual normalized-energy consequence. The finite-fibre wording explicitly bounds a subset by the full source-lattice stabilizer; it does not call a cutoff subset a group. No unresolved step is used in DS18. Higher strata, all J14 nonintegral terms, and the full Mram sign remain open.

The exact certificate passed 24,774 assertions at 2026-11-22T09:17:55.190051+00:00. Its numerical comparisons use only integers and rational arithmetic. No actual global matrix enumeration, real period evaluation, or Lean execution was performed.

\clearpage

\section{AM-K063 --- a quantitative bound for the complete small-field plane part of S\_T}\label{am-k063-a-quantitative-bound-for-the-complete-small-field-plane-part-of-s_t}

Author: GPT-6 Astra Pro, ALP\_MACHINE \nolinkurl{/root/kummer_extension}. 22 November 2026 UTC. Inherited reasoning effort is not separately exposed.

\textbf{Author whole-proof check PASS --- 2026-11-22T09:18:27.220653+00:00.} The complete saved PF1--PF25 proof and full certificate were read and reconstructed. The separate exact replay passed20,901 assertions. No rational short sum or automorphic value was enumerated.

This finding bounds a complete, explicitly defined part of the signed finite remainder S\_T in K062. It is not limited to the Pell sequence of R036. On the J14 host it covers every short actual matrix in the fixed unitary plane with coefficients in Q(zeta\_7). On the J56 host it covers the larger coefficient field Q(zeta\_28). The remaining matrices, including matrices whose coefficients use the new real extensions, are not bounded here. No all-place Mram nonvanishing or DP44 budget conclusion is asserted.

\subsection{1. The actual fixed plane and the asserted bound}\label{the-actual-fixed-plane-and-the-asserted-bound}

Let calE be either actual CM host of K031/K032, with real degree f=42 or168. Retain exactly the finite primes, favorable orbit, cutoff, finite representation and selected compact type in K062. Write

\[
 k_-=\mathbb Q(\sqrt{-7}),\quad
 E_4=\mathbb Q(i,\sqrt{-7}),\quad E_6=\mathbb Q(\zeta_7),
 \quad E_{12}=\mathbb Q(\zeta_{28}),
 \tag{PF1}
\]

and let L be any of these four subfields, of degree 2r. Thus r=1,2,3,6 and m=f/r is an integer. Every L is stable under complex conjugation. The actual rational prime T splits completely in calE, and hence in L. Every selected T prime lies over the same prime (pi) of k\_-, where

\[
 \pi=1+2D+D\sqrt{-7},\qquad \pi\bar\pi=T.
 \tag{PF2}
\]

There are r selected primes of L above pi, each with exactly m lifts to calE. No prime or local level is changed.

Use LIFT41/R036's actual rational vector

\[
 u=(2,u_2,u_3,u_4,0,\ldots,0)\in\mathbb Z^a,
 \qquad u^*u=U=7(11Q^2-2),\quad Q=g_0g_1.
 \tag{PF3}
\]

The ambient Hermitian form is diag(I\_a,-1). Let P\_L consist of all unitary matrices with coefficients in L which fix u-perp pointwise. They preserve the plane spanned by u and the last, negative basis vector. We restrict to the matrices in the actual support calC of K062 and its short region

\[
             \|g_\tau\|_{\rm op}\le H\quad\hbox{for every tau},
 \qquad H=8192\ (J14),\quad H=32768\ (J56).
 \tag{PF4}
\]

Let S\_L denote precisely the sub-sum of K062(AT24)'s S\_T over these matrices. A central class has at most one representative in P\_L: both representatives fix the nonzero space u-perp, so a scalar relating them must be one. The following estimate therefore counts the actual central quotient without adding a multiplicity.

\textbf{Theorem PF5.} Put

\[
 x_L=4^mT^{8-m}.
 \]

Then

\[
 \boxed{\quad
 |S_L|\le28\,15^{2r}H^{4r}
                  \bigl((1+x_L)^r-1\bigr).
 \quad}\tag{PF5}
\]

In particular the exact actual arithmetic gives the following strict binary bounds. An omitted entry means only that PF5 is not useful there.

\begin{longtable}[]{@{}
  >{\raggedright\arraybackslash}p{(\columnwidth - 4\tabcolsep) * \real{0.2727}}
  >{\raggedleft\arraybackslash}p{(\columnwidth - 4\tabcolsep) * \real{0.3636}}
  >{\raggedleft\arraybackslash}p{(\columnwidth - 4\tabcolsep) * \real{0.3636}}@{}}
\toprule\noalign{}
\begin{minipage}[b]{\linewidth}\raggedright
Coefficient field L
\end{minipage} & \begin{minipage}[b]{\linewidth}\raggedleft
J14: \(\lvert S_L\rvert<\)
\end{minipage} & \begin{minipage}[b]{\linewidth}\raggedleft
J56: \(\lvert S_L\rvert<\)
\end{minipage} \\
\midrule\noalign{}
\endhead
\bottomrule\noalign{}
\endlastfoot
\(\mathbb Q(\sqrt{-7})\) & \(2^{-33851}\) & \(2^{-158791}\) \\
\(\mathbb Q(i,\sqrt{-7})\) & \(2^{-12831}\) & \(2^{-75309}\) \\
\(\mathbb Q(\zeta_7)\) & \(2^{-5784}\) & \(2^{-47436}\) \\
\(\mathbb Q(\zeta_{28})\) & --- & \(2^{-19425}\) \\
\end{longtable}

The bound is uniform in the actual finite goodTuple, its two fixed mixed projectors, the selected compact module, and all admitted profiles. It uses only the unitary normalized inequalities \textbar a(g)\textbar\textless=1 and \textbar c(g)\textbar\textless=1. It neither changes a projector nor assumes positivity of individual summands. In fact the proof bounds the sum of the absolute values of all the indicated summands by the same right side.

\subsection{2. Exact local endpoint algebra}\label{exact-local-endpoint-algebra}

This section refines, but does not replace, K062's signed local trace. Let X=P\^{}1(k), q=\textbar k\textbar, E=\{infinity,-1\}, D=X minus E, and m0=q-1. For a matrix B preserving the full pencil let

\[
 k_B=|E\cap BE|,\qquad F_E(B)=|\operatorname{Fix}(B)\cap E|,
 \qquad \chi_{\rm aug}(B)=|\operatorname{Fix}_X(B)|-1.
 \tag{PF6}
\]

Since D has m0 elements, its transported intersection has size C(B)=q-3+k\_B. Substitution in K062(AT8) gives the exact formula

\[
 \boxed{\quad r(B)=\chi_{\rm aug}(B)-F_E(B)
                         +\frac{2-k_B}{q-1}.\quad}
 \tag{PF7}
\]

There is a literal rank-two projector behind this identity. In C{[}X{]} let P\_aug=I-J/(q+1), and put v\_e=delta\_e-1/(q+1) for the two endpoints. Their Gram matrix and inverse are

\[
       G=I_2-\frac{J_2}{q+1},\qquad
       G^{-1}=I_2+\frac{J_2}{q-1}.
 \tag{PF8}
\]

Writing V for the matrix with these two columns, Q\_E=VG\^{}\{-1\}V\^{}* is the orthogonal projector onto their span. The actual projector, extended by zero outside D, is

\[
                    P_D^0=P_{\rm aug}-Q_E.
 \tag{PF9}
\]

Indeed its image is the functions supported in D whose coordinate sum is zero; this is exactly the perpendicular complement of both endpoint vectors inside the augmentation space. Taking its trace against the full permutation action gives PF7. The augmentation representation is the usual rank-one finite Steinberg representation: the action on X is two-transitive, so its permutation commutant has dimension two; after removing constants the remaining complex summand is irreducible of dimension q. No decomposition theorem beyond this elementary argument is required for PF7--9.

For unequal full pencils there is at most one surviving pair. Therefore the exact normalized values, not merely their bounds, are

\[
 \rho(B)=
 \begin{cases}
 0,&C(B)=0,\\
 1/(q-1),&C(B)=1\text{ and the sole pair is fixed},\\
 -1/((q-1)(q-2)),&C(B)=1\text{ and it is not fixed}.
 \end{cases}\tag{PF10}
\]

For the additive projective translations c maps to c+b, one has rho(0)=1 and rho(b)=-1/(q-1) for b nonzero. Hence for any complex weights h(b),

\[
 \sum_{b\in k}\rho(b)h(b)
 =\frac q{q-1}\left(h(0)-\frac1q\sum_{b\in k}h(b)\right).
 \tag{PF11}
\]

This is an exact weighted cancellation, including the two omitted points. It cannot be inserted into the rational sum unless one first identifies a complete actual rational family and its unchanged other local and real factors. PF5 does not require that additional assertion.

\subsection{3. Descent of the actual pencils, including the omitted points}\label{descent-of-the-actual-pencils-including-the-omitted-points}

The replication of a nonintegral place is the arithmetic reason for PF5. It concerns the actual pencil, not a newly chosen model.

Let v0=(Qu,pi\_k) and z=v0/bar(pi), exactly as in R036. Both lie in k\_-\^{}(a+1). At a selected T place, LIFT44(PR5--7) gives a primitive integral row q\_v with

\[
       q_v=-v_0^*J-\pi m_v,\qquad m_vz=1,\quad q_vz=0.
 \tag{PF12}
\]

The first coordinate z1=2Q/bar(pi) is a unit at every selected T place. Choose the single row m\_\emph{=e1\^{}t/z1 over k\_- and put q\_}=-v0\^{}\emph{J-pi m\_}. It is integral at all these places, satisfies q\_*z=0, and has the same reduction modulo pi as every q\_v.

The actual lattice (1+c pi\^{}-1 z q\_v)Lambda does not change on replacing q\_v by q\_\emph{. To check this directly, write q\_v-q\_}=pi ell. Then ell z=0 and the two rank-one operators commute with zero product; their quotient is the integral invertible matrix 1+c z ell. It commutes with the other lattice generator and therefore acts as a right integral change of basis. This proves equality of lattices, not just equality of their intersection and sum. The permitted parameter set remains exactly c in k minus \{-1\}; infinity is still omitted.

Consequently the complete actual pencil and its two endpoints are defined over the completion of k\_- at pi. If g has coefficients in L, then its local partial-permutation matrix, F, C, r and rho are identical at all m={[}calE:L{]} selected calE-primes above a given L-prime. The completion maps are the identity because T splits completely.

For g in the actual support calC, every local determinant is a unit: at T it maps one determinant-zero pencil lattice to another; away T it is integral invertible. Thus det g is a global unit of L and has absolute value one at every embedding. The elementary CM unit argument makes it a root of unity. It lies in the already checked host root group mu28; at most28 possibilities occur.

Let S be the set of selected L-primes above pi at which g is not integral. It is nonempty for the sub-sum S\_L. Unitarity and the unit determinant show that integrality at a selected prime is equivalent to integrality at its conjugate: an integral matrix with unit determinant has integral inverse, and the other component is its inverse-adjoint. There are therefore exactly paired denominator places, and no missing conjugate component.

At every prime in S the action cannot be the identity projectively on the full pencil, since that action fixes the included standard lattice. K062(AT11), and the replication just proved, yield

\[
       \left|\prod_{v\mid T}\rho(g_v)\right|
       \le\left(\frac2{T-2}\right)^{m|S|}
       \le\left(\frac4T\right)^{m|S|}.
 \tag{PF13}
\]

At the other places only \textbar rho\textbar\textless=1 is used. Nonintegrality, rather than a chosen scalar or a chosen local representative, defines S.

\subsection{4. Plane coordinates and a denominator-sensitive lattice count}\label{plane-coordinates-and-a-denominator-sensitive-lattice-count}

On the plane with basis (u,e), the form is diag(U,-1). Let epsilon be the determinant of the plane action; it is also det g. Comparing the ordinary inverse of a two-by-two matrix with its Hermitian inverse shows that every g under consideration has the unique form

\[
 \begin{pmatrix}\alpha&\beta\\
      \epsilon U\bar\beta&\epsilon\bar\alpha\end{pmatrix},
 \quad \alpha\bar\alpha-U\beta\bar\beta=1,
 \quad \alpha,\beta\in L,
 \quad \epsilon\in\mu_{28}\cap L.
 \tag{PF14}
\]

The ambient matrix is consequently

\[
 g=\begin{pmatrix}
 I_a+(\alpha-1)uu^*/U&\beta u\\
 \epsilon\bar\beta u^*&\epsilon\bar\alpha
 \end{pmatrix}.
 \tag{PF15}
\]

The lower-left entry in the plane basis is epsilon U bar(beta), whereas the ambient lower-left row is epsilon bar(beta)u*. Keeping that distinction prevents an extra or missing factor U. Determinants other than one are included throughout.

The original short-region condition gives \textbar sigma(alpha)\textbar\textless=H and sqrt(U)\textbar sigma(beta)\textbar\textless=H at every embedding of L, by using the orthonormal plane basis (u/sqrt(U),e). In particular both absolute values are at most H.

Put

\[
       \mathfrak d_S=\prod_{P\in S}P\bar P,
       \qquad N_L\mathfrak d_S=T^{2|S|}.
 \tag{PF16}
\]

K062's actual pairwise lattice support puts every entry of g in d\_S\^{}-2. This includes the possible order-two poles; an order-one bound would be false in general. From the last diagonal entry and the first positive entry in the last column of PF15, using u1=2, we obtain

\[
           \alpha\in\mathfrak d_S^{-2},\qquad
           \beta\in\tfrac12\mathfrak d_S^{-2}.
 \tag{PF17}
\]

No divisibility assumption on U is used in this count. Outside T alpha is integral and 2beta is integral, exactly as PF17 states.

Here is the needed elementary lattice estimate. If a fractional ideal I has norm T\^{}-4b, every nonzero x in I satisfies \textbar N\_L/Q x\textbar\textgreater=T\^{}-4b. Hence in the r complex embeddings,

\[
                    \max_\sigma|\sigma(x)|\ge T^{-2b/r}=:\lambda.
 \tag{PF18}
\]

The same lower bound holds for differences of distinct elements of I. In C\^{}r with the maximum complex absolute-value norm, the open balls of radius lambda/2 about such elements are disjoint. For centers in the radius-H ball they lie in the ball of radius H+lambda/2. Ordinary Euclidean volume, in real dimension2r, therefore gives

\[
 \#\{x\in I: \max|\sigma x|\le H\}
       \le(1+2H/\lambda)^{2r}
       \le(3H/\lambda)^{2r}.
 \tag{PF19}
\]

We used H\textgreater=1 and lambda\textless=1. Boundary touching is harmless because the balls are open. For (1/2)I the separation is lambda/2 and the corresponding bound is (5H/lambda)\^{}(2r). Thus for a fixed set S of b primes the number of possible pairs (alpha,beta) is at most

\[
                 15^{2r}H^{4r}T^{8b}.
 \tag{PF20}
\]

This deliberately ignores the norm equation, finite support constraints and G conditions. It is still an upper bound. Multiplication by28 accounts for every actual determinant possibility. The proof uses only the product formula and an elementary packing estimate; it does not assume that the fractional ideal is principal or has a small chosen basis.

\subsection{5. Proof of the quantitative bound and the actual numbers}\label{proof-of-the-quantitative-bound-and-the-actual-numbers}

For every g in K062's summation, the finite normalized coefficient and the selected real normalized coefficient satisfy \textbar a(g)\textbar\textless=1 and \textbar c(g)\textbar\textless=1. The first inequality follows from the trace of a unitary operator against an orthogonal projection divided by its rank. The second follows in the same way from compression to W. They apply to the actual two mixed finite projectors and to every actual selected W, regardless of its dimension or highest weight.

For a fixed support S of size b, multiply PF13 by the count PF20. There are binomial(r,b) possible supports. Summing the absolute values over nonempty supports gives

\[
 |S_L|\le28\,15^{2r}H^{4r}
 \sum_{b=1}^r\binom rb
                  \bigl(4^mT^{8-m}\bigr)^b,
 \tag{PF21}
\]

which is PF5. The central quotient has no repeated plane representative, as checked in section1. No sign is discarded in a lower bound: PF21 is an upper bound on the absolute value of the entire indicated sub-sum.

For the actual instances T\textgreater2\^{}1000 on J14 and T\textgreater2\^{}995 on J56, while H=2\^{}13 and2\^{}15 respectively. If m\textgreater8, x\_L\textless1, and

\[
       (1+x_L)^r-1\le(2^r-1)x_L<2^r x_L.
 \tag{PF22}
\]

Also28\textless2\^{}5 and15\^{}(2r)\textless2\^{}(8r). Therefore the following strict and entirely integral sufficient exponent is available, where b\_T is 1000 or995 and h\_H is13 or15:

\[
       |S_L|<2^{-B_L},\qquad
       B_L=(m-8)b_T-2m-5-9r-4r h_H.
 \tag{PF23}
\]

For r=1 the exact identity (1+x)\^{}1-1=x saves one extra bit relative to PF23. The table in section1 uses this improvement for its first row. The certificate also compares the exact rational expression PF5 directly with every displayed power of two, using the unchanged DO040 T integers. Thus none of the entries depends on floating-point logarithms.

For E12 on J14 one has m=7. In that case x\_L=4\^{}7 T grows, so the same estimate does not give a small bound. This failure is explicitly retained; no extrapolation from its proper subfields is made.

\subsection{6. Consequence for the exact remaining gate}\label{consequence-for-the-exact-remaining-gate}

On J14, split K062's S\_T into its complete E6-plane part and all remaining matrices. On J56 use the complete E12-plane part. Then

\[
 \begin{array}{ll}
 J14:&S_T=S_{E_6}+S_{\rm rest},\quad |S_{E_6}|<2^{-5784},\\
 J56:&S_T=S_{E_{12}}+S_{\rm rest},\quad |S_{E_{12}}|<2^{-19425}.
 \end{array}\tag{PF24}
\]

These are exact partitions of the original finite signed sum. The integral part and infinite tail were already bounded in K062, and are not counted a second time here. In particular neither this complete central-CM plane nor its explicit Pell sequence can by itself cancel the normalized identity contribution in K062(AT24).

The unbounded S\_rest includes non-plane matrices and plane matrices whose coefficients use a larger field than the indicated L. No rank-two classification of all short correspondences is asserted. A proof that S\_rest is greater than the remaining negative threshold is still required for all-T Mram nonvanishing. The local spherical Hilbert Gram is not substituted for the actual global evaluation kernel.

More explicitly, retain the positive original normalization

\[
 \mathcal N=w_x\,28C_\infty d_Rd_W
                       \left(\frac{T-2}{T-1}\right)^f.
 \]

With B=5784 or19425 respectively, K062 and PF24 give

\[
 \frac{E_{x,\mathrm{ram}}}{\mathcal N}
       =1+S_{\rm rest}+\varepsilon',\qquad
       |\varepsilon'|<2^{-1999}+2^{-B}.
 \tag{PF25}
\]

The two parts are real because each is closed under inversion and the normalized coefficients conjugate under inversion, as in K062. Thus the explicit sufficient remaining inequality is S\_rest \textgreater{} -1+2\^{}(-1999)+2\^{}(-B). It is not proved here. The known plane contribution has been charged in this threshold, rather than deleted from the original sum without an error allowance.

\subsection{7. Inputs and verification scope}\label{inputs-and-verification-scope}

The carried inputs are K062(AT1--11,AT14--24), LIFT41(NP1--12), LIFT44(PR4--7), R036(1--16), and the actual fields, roots and splitting data K031/K032/K040. Their proof/source records are retained. The exact endpoint projector algebra, subfield descent of the actual pencil, complete plane count, and quantitative sub-sum bound are proved here. No new analytic or class-field theorem is invoked.

The companion certificate tests the endpoint projector and partial permutation formulas, weighted additive cancellation, the plane coordinate/ambient normalization, repeated prime counts, and all actual integer bounds. It does not enumerate S\_L or S\_rest, compute a theta value, or run Lean.

\clearpage

\section{AM-K064 --- the plane norm equation controls the complete Q(zeta\_28) plane}\label{am-k064-the-plane-norm-equation-controls-the-complete-qzeta_28-plane}

Author: GPT-6 Astra Pro, ALP\_MACHINE \nolinkurl{/root/kummer_extension}. 22 November 2026 UTC. Inherited reasoning effort is not separately exposed.

\textbf{Author whole-proof check PASS --- 2026-11-22T09:26:40.611436+00:00.} The complete saved NR1--NR19 proof and full certificate were read and reconstructed. The separate exact replay passed1,705 assertions. No actual rational sum or real theta coefficient was evaluated.

This is a separate refinement of frozen K063. It retains its actual sub-sum, all determinant components, finite projectors and original class normalization. The plane norm equation cuts the necessary residue count from T\^{}(8b) to (T\^{}6+T\^{}5-T\^{}3)\^{}b. As a consequence the \textbf{entire fixed Q(zeta\_28)-defined unitary plane} contributes less than2\^{}-614 on J14 and less than2\^{}-21414 on J56. The rest of K062's finite S\_T remains open.

\subsection{1. Exact statement with the unchanged plane and support}\label{exact-statement-with-the-unchanged-plane-and-support}

Use all notation of K063(PF1--PF5,PF12--PF17). In particular L is one of the four listed CM subfields, {[}L:Q{]}=2r, m=f/r, H is the original 8192/32768 cutoff, and P\_L fixes the actual u-perp pointwise. Let S\_L be the same nonintegral short sub-sum of K062's S\_T. Put

\[
       N_2(T)=T^6+T^5-T^3,\qquad
       y_L=4^mN_2(T)T^{-m}.
 \tag{NR1}
\]

\textbf{Theorem NR2.} The sum of absolute values of the indicated summands, and therefore \textbar S\_L\textbar, is at most

\[
 \boxed{\quad
 28\,15^{2r}H^{4r}\bigl((1+y_L)^r-1\bigr).
 \quad}\tag{NR2}
\]

For L=Q(zeta\_28) this gives, for the actual fixed T,

\[
              |S_L|<2^{-614}\quad(J14),\qquad
              |S_L|<2^{-21414}\quad(J56).
 \tag{NR3}
\]

No term of the original sum is reweighted. The improved estimate counts only a necessary congruence of the actual unitary equation; it does not replace its global solutions by arbitrary local ones.

\subsection{2. Paired fractional-ideal residues without a global generator}\label{paired-fractional-ideal-residues-without-a-global-generator}

Fix the nonempty set S of selected L-primes above pi at which a matrix is nonintegral, and put b=\textbar S\textbar{} and

\[
 \mathfrak d_S=\prod_{P\in S}P\bar P,
 \qquad I=\mathfrak d_S^{-2}.
 \tag{NR4}
\]

K063 supplies the unique plane coordinates alpha,beta and determinant epsilon in mu28 intersect L, with

\[
 \alpha\bar\alpha-U\beta\bar\beta=1,\quad
 \alpha\in I,\quad b_0:=2\beta\in I,\quad
 |\sigma(\alpha)|\le H,\quad |\sigma(b_0)|\le2H.
 \tag{NR5}
\]

The symbol b0 here is only the coordinate 2beta, not a weight parameter. The actual integers satisfy gcd(2U,T)=1, as follows also from the exact certificate. Thus U/4 is a unit in every T residue ring.

For each selected P, set R\_P=O\_(L,P)/pi\^{}2 O\_(L,P). Its residue field has T elements and R\_P has T squared elements. Since pi has valuation one at P, multiplication by pi squared identifies the local quotient P\^{}-2/O with R\_P. At the conjugate prime use bar(pi) squared and then conjugation back to P. Fractional-ideal CRT therefore gives an additive isomorphism

\[
 I/O_L\simeq\prod_{P\in S}(R_P\oplus R_P),
 \qquad
 x\longmapsto
       (\pi^2x\bmod P^2,\ \pi^2\bar x\bmod P^2)_{P\in S}.
 \tag{NR6}
\]

This is a local quotient construction. It does not assert that P or I has a global generator. The conjugate coordinate in NR6 is precisely the residue at bar(P), transported by the ring isomorphism induced by complex conjugation. Hence the two coordinates are independent before the norm equation is imposed.

Let a,a' be the two residues of alpha and b,b' those of b0 at P. Multiplying the actual equation NR5 by pi to the fourth power gives the necessary relation

\[
           aa'-(U/4)bb'=0\quad\hbox{in }R_P.
 \tag{NR7}
\]

The right side before reduction is pi to the fourth power, and hence vanishes modulo pi squared. Because U/4 is a unit, the four residues are in bijection with the entries of the matrix

\[
                  \begin{pmatrix}a&b\\(U/4)b'&a'\end{pmatrix}
 \tag{NR8}
\]

subject to its determinant being zero. All determinant phases epsilon are still allowed: the plane norm equation NR5 is independent of epsilon.

\subsection{3. Exact singular-matrix count over the length-two residue ring}\label{exact-singular-matrix-count-over-the-length-two-residue-ring}

For any length-two unramified local ring R with residue cardinality q, the number of two-by-two matrices of determinant zero in R is

\[
                         q^6+q^5-q^3.
 \tag{NR9}
\]

Here is a direct count that also applies to R\_P. If the first row is primitive, there are q\^{}4-q\^{}2 choices. Its determinant functional on the second row is surjective R\^{}2 to R and has q\^{}2 solutions with value zero. This gives q\^{}6-q\^{}4 matrices.

If the first row is a nonzero multiple of the uniformizer, there are q\^{}2-1 choices. The determinant condition is one nonzero linear equation over the residue field on the second row. It has q\^{}3 lifts to R\^{}2, giving q\^{}5-q\^{}3 matrices. If the first row is zero, there are q\^{}4 choices for the second row. Adding the three disjoint cases proves NR9.

Consequently at most N2(T)\^{}b pairs of residue classes (alpha,b0) modulo O\_L squared can satisfy all the necessary congruences NR7. Some of these classes may contain no actual global norm-one pair; including them only enlarges the upper bound.

\subsection{4. Integral lifts and the new support sum}\label{integral-lifts-and-the-new-support-sum}

Two distinct elements of one residue class modulo O\_L differ by a nonzero algebraic integer. Its field norm is a nonzero integer, so the maximum complex absolute value of its conjugates is at least one. Packing the open maximum-complex-norm balls of radius1/2 gives, in real dimension2r,

\[
 \#\{\alpha\hbox{ in a fixed class}:\max|\sigma\alpha|\le H\}
       \le(1+2H)^{2r}\le(3H)^{2r},
 \tag{NR10}
\] \[
 \#\{b_0\hbox{ in a fixed class}:\max|\sigma b_0|\le2H\}
       \le(1+4H)^{2r}\le(5H)^{2r}.
 \tag{NR11}
\]

These are counts within arbitrary additive cosets; no globally small coset representative or principal-ideal assumption is needed. The product gives at most15\^{}(2r)H\^{}(4r) lifts per residue pair. Including all at most28 determinant choices, the number of short matrices for a fixed support S is therefore bounded by

\[
                         28\,15^{2r}H^{4r}N_2(T)^b.
 \tag{NR12}
\]

The actual nonintegral-place replication and its signed-projector bound remain exactly K063(PF13): the absolute product of rho factors is at most(4/T)\^{}(mb). The actual normalized finite and real traces have absolute value at most one. Multiplying by NR12 and summing over the binomial(r,b) nonempty supports proves NR2.

The new bound improves the count, not the local projector exponent. In particular it does not treat the unequal-pencil sole-fixed-pair case as though its rho had order T\^{}-2. That case retains its actual order T\^{}-1, already covered by the same bound.

\subsection{5. Exact numerical specialization}\label{exact-numerical-specialization}

For q\textgreater1, N2(q)\textless2q\^{}6. Thus

\[
                  y_L<2^{1+2m}T^{6-m}.
 \tag{NR13}
\]

All four actual subfields have m\textgreater=7, so the actual values satisfy y\_L\textless1. As in K063, (1+y\_L)\^{}r-1\textless2\^{}r y\_L. With T\textgreater2\^{}(b\_T), H=2\^{}(h\_H), one obtains

\[
 |S_L|<2^{-B'_L},\qquad
 B'_L=(m-6)b_T-2m-6-9r-4rh_H.
 \tag{NR14}
\]

For r=1 the exact one-term binomial saves one bit. For the whole E12 plane, r=6. On J14, m=7, b\_T=1000 and h\_H=13, giving B'\_L=614. On J56, m=28, b\_T=995 and h\_H=15, giving21414. The companion certificate also compares the exact rational NR2 directly to these dyadic bounds. No rounded logarithm is used.

\subsection{6. An explicit actual nonintegral member of the removed family}\label{an-explicit-actual-nonintegral-member-of-the-removed-family}

The subfamily just bounded is nonempty in the actual short support on both hosts. Retain W=Q\^{}2U and the actual pi\_k. Set

\[
 \alpha_-=1+\frac{2W}{T},\qquad
 \beta_-=-\frac{2Q\bar\pi_k}{T},\qquad\epsilon=1,
 \tag{NR15}
\]

and form the ambient matrix R\_minus by K063(PF15). Since pi\_k bar(pi\_k)=T+W, one computes

\[
 \alpha_-\bar\alpha_--U\beta_-\bar\beta_-=1,
 \qquad R_- (Qu,\pi_k)=(-Qu,\pi_k).
 \tag{NR16}
\]

Its entries lie in k\_-, it fixes u-perp, and it is integral invertible away T. Its positive diagonal update is (2Q\^{}2/T)uu*, and its off-diagonal entries are divisible by Q away T. Hence it is principal at both actual G primes, exactly in the original frame. In particular its actual finite normalized G coefficient is one.

Substitute j=k and t=-1 in the already proved graph calculation R036(12--15). The permitted scalar is s=-1 at every selected T place, and the inverse-image lattice is the actual L\_(c=-2). Its conjugate component is the required dual lattice. Thus this is an actual member of calC, not a newly selected Hecke or Schwartz translate. Since T is odd and c=-2 is neither0 nor the omitted-1, it is nonintegral.

The exact inequalities W\textless T on both hosts give alpha\_-\textless3. Its Cartan operator norm is alpha\_-+sqrt(alpha\_- squared-1)\textless2alpha\_-\textless6\textless H. This proves membership in the original short region at every real label.

One can also determine its local trace category. Its ambient trace is n+4W/T, which is nonintegral at each selected T prime because T does not divide4W. A full-pencil stabilizer preserves the full O-lattice M and therefore has integral trace. Hence R\_minus does not preserve the full pencil. There is exactly one surviving pair, and it is the nonfixed pair L\_(-2) to Lambda. Consequently

\[
                   \rho((R_-)_v)
                       =-\frac1{(T-1)(T-2)}
              \quad\hbox{at every selected T place}.
 \tag{NR17}
\]

This explicitly locates an allowed nonintegral rational matrix in the removed subfamily. It does not assert a nonzero real compressed trace for that particular matrix; strictness here concerns the actual index set, not an uncomputed automorphic coefficient. R036's fifth-power matrix is another member of this family on J14, where it is short.

\subsection{7. The exact remaining inequality}\label{the-exact-remaining-inequality}

On either host split the original signed sum as

\[
                   S_T=S_{E_{12}}+S_{\rm rest}.
 \tag{NR18}
\]

The removed part is now the complete E12-defined fixed plane on both hosts. With the original positive normalization calN of K063(PF25), K062 and NR3 give

\[
 \frac{E_{x,\mathrm{ram}}}{\mathcal N}
    =1+S_{\rm rest}+\varepsilon'',\qquad
 |\varepsilon''|<2^{-1999}+2^{-B},
 \quad B=614\ (J14),\ 21414\ (J56).
 \tag{NR19}
\]

Thus the remaining sufficient inequality is S\_rest \textgreater{} -1+2\^{}(-1999)+2\^{}(-B). It has not been proved. The inverse-closed partition retains reality of both sub-sums. Non-plane matrices and matrices with coefficients using the additional real extensions remain in S\_rest with their exact original factors. No full rank-two reduction, all-place Mram theorem, or completed mixed DP44 cost is inferred.

\subsection{Verification scope}\label{verification-scope-3}

The proof uses the checked actual K063 field/pencil/plane interfaces, K062's exact normalized sum, and R036's fixed packet calculation. The paired fractional-ideal CRT argument and the length-two determinant count are proved above. The certificate independently enumerates finite length-two matrix models, checks the conjugate-pair exponent and lift normalizations, the actual R\_minus identities, and every dyadic bound. It does not enumerate the actual S\_T or evaluate a real theta coefficient.

\clearpage

\section{AM-DO-059 --- a sharper complete tail and the full single-place short sum}\label{am-do-059-a-sharper-complete-tail-and-the-full-single-place-short-sum}

Author: GPT-6 Astra Pro, ALP\_MACHINE /root/definite\_oscillator. 22 November 2026 UTC. Inherited reasoning effort is not separately exposed.

\textbf{Author whole-proof check PASS --- 2026-11-22T09:30:40.151685+00:00.}

On the actual J56 host at the upper K044 gap endpoint, all supported nonintegral short correspondences that are nonintegral at exactly one relative T place have total absolute contribution less than \(2^{-104546}\). This includes every permitted pole depth at that place, every source and target pencil pair, and both conjugate CM components. It retains the exact signed projector and class normalization of K062.

There are two new ingredients. A smaller negative-moment order permits the algebraic cutoff \(H=8192\) on J56 while preserving the complete infinite-tail bound. Then the last column is counted in the actual target-lattice over-lattice, whose finite quotient is smaller than the ambient denominator window. The moment-order proposal came from /root; the derivative, sampling, and exact arithmetic checks are given here. Neither ingredient proves the sign of the remaining sum with two or more nonintegral places.

\subsection{1. The unchanged actual kernel and a different moment order}\label{the-unchanged-actual-kernel-and-a-different-moment-order}

Keep all data of K062 and DO058, now on J56: \[
 a=f=168,\qquad n=a+1=169,\qquad
 d=f(a+1)^2,\qquad \ell=\ell_+,\qquad s_0=a+1.
 \tag{SC1}
\] No automorphic representation, weight, finite projector, class frame, level, or scalar-orbit point is changed. Only the proof of the analytic tail uses a different moment order: \[
                         t=\lfloor\ell/16\rfloor.
 \tag{SC2}
\]

Use the exact K043/K046 Pieri notation. At a label of role top, thick, or thin, the shifted-metric starting factors and polynomial exponents are \(r_i,m_i\), with the actual word and canonical twist. K046(BL8--10) gives the negative radial moments \[
 M_\tau(j)=
 \prod_{i=0}^{a-1}\frac{(r_i+m_i-j)_j}{(r_i-1-j)_j}.
 \tag{SC3}
\] This identity is for the actual invariant Hilbert metric with its PBW normalization. It applies to arbitrary vectors in each selected compact component, not only a highest vector.

All Lie words needed by the full L1 sampling argument have length at most d.~K061(CG8) proves that every compact type reached by such a word has Pieri exponents between zero and \(m_i+2d\), with the same \(r_i\). Compact invariance makes the shifted metrics orthogonal on distinct types. Thus the same upper bound applies to their sum without multiplying by the number or dimension of the types.

At the zero word and zero canonical twist, put \[
 A'_\tau(t)=
       \prod_{i=0}^{a-1}\frac{r_i+m_i+2d}{r_i-1-t}.
 \tag{SC4}
\] The actual integers satisfy \(t>s_0\) and \(t<\min_i(r_i-1)\). Increasing the canonical twist adds the same nonnegative number to numerator and denominator in each factor, and an interior one-bit only decreases \(m_i\). Therefore the zero-word, zero-twist expression bounds every actual word and allowed twist, exactly as in K046(BL15--16). Every negative moment through order t is at most \(A'_\tau(t)^j\).

Substitution of the accepted K044 integer profiles gives the strict, exact rational bounds \[
 \begin{array}{c|ccc}
 \text{role}&\text{top}&\text{thick}&\text{thin}\\ \hline
 A'_\tau(t)&<35&<1266&<453 .
 \end{array}
 \tag{SC5}
\] The numerator in SC4 includes the full \(2d\) allowance in every coordinate. Thus SC5 is already a derivative bound; it is not just a bound on the original selected vector. The certificate compares the integer products directly and verifies all endpoint monotonicities. In particular all these constants are less than 1280 and less than \(2^{11}\).

\subsection{\texorpdfstring{2. The full derivative L1 tail at radius \(9/2\)}{2. The full derivative L1 tail at radius 9/2}}\label{the-full-derivative-l1-tail-at-radius-92}

Set \[
 \epsilon=1/1536,\qquad R=9/2,\qquad R'=4499/1000,\qquad
 \delta_*=f\,2^{-\ell/1024}.
 \tag{SC6}
\] The exact inequalities \[
       \epsilon A'_\tau<5/6<2^{-1/4},\qquad
       \epsilon^{-1}<2^{11}
 \tag{SC7}
\] follow from SC5 and \(2\cdot5^4<6^4\).

For a Lie word I of length m, use precisely K061(CG7): \[
 B_I(x)=\operatorname{pr}_V\pi(x^{-1})d\pi(X_I)|_W,\qquad
 b_I(x)=\|B_I(x)\|_{\mathrm{HS}}/\sqrt{\dim W}.
 \tag{SC8}
\] The actual Casimir bound and Haar normalization remain \[
 \|d\pi(X_I)|_W\|\le2^{1000m},\qquad
 C=\prod_{\tau,i}(r_{\tau i}-1)<2^{960af}.
 \tag{SC9}
\] These are K061(CG9--10) and K046(BL19); changing t changes neither inequality.

For completeness, the weighted Cauchy--Schwarz calculation gives the absolute estimate needed for sampling. Let \(\mathcal W=\prod_\tau Y_\tau^{s_0}\). Under the actual ball Haar, \[
 \int\mathcal W\,dg=(a!)^{-f}\le1.
 \tag{SC10}
\] On \(Y_\tau<\epsilon\), \[
 1_{\{Y_\tau<\epsilon\}}Y_\tau^{-s_0}
          \le\epsilon^{t-s_0}Y_\tau^{-t}.
 \tag{SC11}
\] Apply Cauchy--Schwarz with this weight, then SC3--5 and the bound in SC9. The joint shifted metrics factor over labels, including for entangled vectors and the Hilbert--Schmidt average. One obtains \[
 \int_{\{Y_\tau<\epsilon\}}b_I
 \le {2^{1000m}\over C}
       2^{B_*/2-t/8},
 \qquad B_*=960af+11f(a+1).
 \tag{SC12}
\] Indeed the expression inside the square root is bounded by \[
 C\,\epsilon^{-s_0}
       (\epsilon A'_\tau)^t\,2^{11(f-1)s_0}
             <2^{B_*-t/4}.
 \tag{SC13}
\] This calculation retains the factor C in converting the squared norm to an absolute L1 tail. It does not identify an L2 probability tail with an L1 tail.

The actual integer inequality \[
                     128t-512B_*>\ell
 \tag{SC14}
\] implies \(B_*/2-t/8<-\ell/1024\). Furthermore \[
 \cosh^2 R'>e^{\,2R'}/4>1996>1536.
 \tag{SC15}
\] An elementary rational certificate for the middle inequality uses \(e^3>20\) and \(e^{-1/500}>499/500\), giving \(e^{8.998}/4>8000(499/500)/4=1996\). Thus outside the product Cartan set \(C_{R'}\), at least one label has \(Y_\tau<\epsilon\). Summing SC12 over those labels proves \[
             \int_{G\setminus C_{R'}}b_I
                    <{2^{1000m}\delta_*\over C}.
 \tag{SC16}
\] The full L1 bound remains the one already used in K061: \[
 \int_Gb_I\le 2^{1000m}G_0,\qquad
 G_0=C^{-1/2}(a!)^{-f/2}2^{11f(a+1)/2}.
 \tag{SC17}
\] It follows from the same weighted calculation through order \(s_0\).

\subsection{3. Compatibility with the complete rational sampling argument}\label{compatibility-with-the-complete-rational-sampling-argument}

Keep the actual entire finite support \(\mathscr C\) of K062. No group structure on this support is required. K062(AT18--19) proves that a ratio of two supported matrices has entry denominator at most \(T^4\), and constructs disjoint left charts using \[
 \rho_T={1\over100d^2T^4},\qquad
                         K_T=2^{1200d}T^{8d}.
 \tag{SC18}
\] The radial displacement of a chart is less than \(1/1000\). Therefore a chart centred outside \(C_R\) lies outside \(C_{R'}\), exactly the margin in SC6. All derivative orders, Casimir factors, Haar Jacobians and ordered-word counts in K061(CG13--18) are unchanged. Substituting SC16 for its previous derivative tail is legitimate in every chart.

The same normalized convolution identity K061(CG19--20), already extended to the whole support by K062(AT19--20), now yields \[
 \sum_{\substack{g\in Z\backslash\mathscr C\\
          \max_\tau\log\|g_\tau\|_{\mathrm{op}}>9.01}}
                     |c(g)|
                <2K_TG_0\delta_*<2^{-2000}.
 \tag{SC19}
\] Here the common packet centre still has Cartan parameters below \(1/200\), and \(9.01=2R+1/100\). The use of normalized Hilbert--Schmidt traces in SC8 means that no factor \(\dim W\) is inserted in this estimate.

For the last inequality, \(T<2^{1001}\) and \(C,(a!)^f\ge1\) give the sufficient exponent comparison \[
 {\ell\over1024}>
 2000+1+9208d+\lceil\log_2f\rceil
                         +{11f(a+1)\over2}.
 \tag{SC20}
\] It is checked exactly. Thus the changed moment order remains sufficient even after the full \(T^{8d}\) denominator-dependent sampling cost.

We may consequently use the smaller algebraic cutoff \[
                                H=2^{13}=8192
 \tag{SC21}
\] on J56. The exact lower series \(\log2>2\sum_{j=0}^{14}3^{-2j-1}/(2j+1)\) verifies \(\log H>9.01\), so every term outside this cutoff is included in SC19. All integral noncentral terms inside it have absolute sum less than \(2^{-2000}\), by K062(AT22) at the older larger cutoff.

Therefore, if \(S_T^{\mathrm{new}}\) denotes the actual nonintegral sum inside SC21, \[
 {E_{x,\mathrm{ram}}\over
 w_x\,28C_\infty d_Rd_W((T-2)/(T-1))^f}
       =1+S_T^{\mathrm{new}}+\varepsilon_*,\qquad
                           |\varepsilon_*|<2^{-1999}.
 \tag{SC22}
\] This is a fresh decomposition of the same complete kernel. The annulus between the two cutoffs is covered by SC19; it is not dropped or counted twice.

As a first consequence, DO058(DS15) at the new H gives \[
             |S_{1,2}^{\mathrm{new}}|<2^{-271780},
             \qquad
             |S_{1,3}^{\mathrm{new}}|<2^{-19933}.
 \tag{SC23}
\] For e=3 there is exactly one nonintegral relative place, with either oriented pole pair \((1,2)\) or \((2,1)\). Both are counted. These inequalities are verified by actual integer products, without asymptotic logarithmic estimates.

\subsection{4. The target-lattice over-lattice improves the whole support count}\label{the-target-lattice-over-lattice-improves-the-whole-support-count}

Let s be the number of nonintegral relative T places, as in DO058. Fix a choice of ordered source/target pencil points at those places. At either CM component over one of them, let L and L' be the chosen source and target lattices and \(\varpi\) the local uniformizer. K062's actual pencil satisfies \[
            \varpi\Lambda\subset L,L'\subset\varpi^{-1}\Lambda,
            \qquad \operatorname{vol}(L)=
                   \operatorname{vol}(L')=\operatorname{vol}(\Lambda).
 \tag{SC24}
\] For an isometry g from L to L', \[
          \varpi e_n\in L
          \quad\Longrightarrow\quad
                       g e_n\in\varpi^{-1}L'.
 \tag{SC25}
\] The lattice on the right contains the standard lattice and has exactly \[
                [\,\varpi^{-1}L':\Lambda\,]=T^n
 \tag{SC26}
\] cosets. The conjugate component obeys the same statement; its pencil lattice is the paired dual, still of determinant valuation zero.

At all remaining finite places the last column is integral. Thus the last column lies in a global O\_E over-lattice \(\mathcal A\) containing the standard \(O_E^n\), with \[
                       [\mathcal A:O_E^n]=T^{2ns}.
 \tag{SC27}
\] This construction depends on the actual target pencil points. It uses their exact determinant ideals; it is not the replacement of the target by an arbitrary denominator window.

Within each coset of \(O_E^n\), distinct columns have trace distance at least \(\sqrt f\), since a nonzero integral coordinate has absolute field norm at least one. Their trace norms are at most \(\sqrt f\,H\). The same Euclidean packing proof as DO058(DS12), now in each coset, therefore counts at most \((2H+1)^{2fn}\) columns per coset. Hence \[
       \#\{\text{possible last columns for this pencil pair}\}
                 \le(2H+1)^{2fn}T^{2ns}.
 \tag{SC28}
\]

For a fixed column, DO058(DS8--13) bounds the isometry fibre by \[
 G_n=\gcd\bigl(|\operatorname{GL}_n(\mathbb F_{g_0})|,
             |\operatorname{GL}_n(\mathbb F_{g_1})|\bigr).
 \tag{SC29}
\] This is the same exact 5,022-bit integer on J56. The fibre is a subset of the finite source-lattice group fixing \(e_n\); no cutoff subset is claimed to be a group. There are \((T-1)^{2s}\) ordered pencil pairs and \(\binom fs\) possible support sets. Consequently \[
 \#\{\text{all supported matrices with exactly s nonintegral places}\}
 \le \binom fs(T-1)^{2s}G_n(2H+1)^{2fn}T^{2ns}.
 \tag{SC30}
\] The bound is independent of the entry-denominator exponent e. It already counts all its possible values, so no extra pole-pattern multiplicity is introduced.

\subsection{5. The complete single-place contribution}\label{the-complete-single-place-contribution}

DO058(DS4--7) gives \(e\le4s\) and hence \[
 |a(g)|\le\min\left\{1,\,
             {2(2H)^{af}T^{2as}\over Q^{a(f-1)}}\right\}.
 \tag{SC31}
\] At each nonintegral relative place, K062's normalized signed trace is at most \(2/(T-2)\) in absolute value. Combining SC30--31 yields the complete support-size estimate \[
 \begin{split}
 |S_s^{\mathrm{new}}|\le{}&
 \binom fs(T-1)^{2s}G_n(2H+1)^{2fn}T^{2ns}
          \left({2\over T-2}\right)^s\\
 &\times
 \min\left\{1,\,
             {2(2H)^{af}T^{2as}\over Q^{a(f-1)}}\right\}.
 \end{split}
 \tag{SC32}
\] This bounds the sum of absolute values of the actual terms, including all permissible e and every pencil pair.

At s=1 the character bound in SC31 is less than one. The exact numerator and denominator are therefore \[
 A_1=4f(T-1)^2G_n(2H)^{af}(2H+1)^{2fn}T^{2n+2a},
 \qquad B_1=(T-2)Q^{a(f-1)}.
 \tag{SC33}
\] The actual integer comparison is \[
   2^{104546}A_1<B_1\le2^{104547}A_1,
 \quad
 \operatorname{bits}(A_1)=1868068,\quad
 \operatorname{bits}(B_1)=1972614.
 \tag{SC34}
\] It proves the announced result \[
                  \boxed{|S_1^{\mathrm{new}}|<2^{-104546}.}
 \tag{SC35}
\] All three possible total pole exponents e=2,3,4 at this one place are included. SC23 records the earlier denominator-specific estimates: its e=2 bound remains sharper, whereas SC35 also improves the e=3 bound. These estimates are not added to SC35.

The same crude worst-pole bound at s=2 is greater than one: its largest strict inverse binary exponent is \(-567385\). Thus SC32 does not establish the desired sign of the complete remaining sum. There is no claim that every nonintegral correspondence has s=1.

\subsection{6. Exact remaining normalized gate}\label{exact-remaining-normalized-gate}

Let \(S_{\ge2}^{\mathrm{new}}\) retain every original supported term inside H whose matrix is nonintegral at two or more relative T places. Then \[
 {E_{x,\mathrm{ram}}\over
 w_x\,28C_\infty d_Rd_W((T-2)/(T-1))^f}
       =1+S_{\ge2}^{\mathrm{new}}+\varepsilon_{**},
 \quad
 |\varepsilon_{**}|<2^{-1999}+2^{-104546}.
 \tag{SC36}
\] No class weight, finite projector, selected real factor, or scalar-orbit denominator has changed.

One may also remove the already checked complete \(\mathbb Q(\zeta_{28})\)-plane part from this remainder. K063 bounds the sum of absolute values in that plane by \(2^{-19425}\), even at the older larger cutoff. Any nonintegral matrix in that plane has nonintegrality replicated at least 28 times on J56, so it is disjoint from the s=1 stratum. For the actual remainder outside that plane, \[
 {E_{x,\mathrm{ram}}\over
 w_x\,28C_\infty d_Rd_W((T-2)/(T-1))^f}
       =1+S_{\mathrm{rest}}^{\mathrm{new}}+\varepsilon_{\mathrm{rest}},
 \quad
 |\varepsilon_{\mathrm{rest}}|
   <2^{-1999}+2^{-104546}+2^{-19425}<2^{-1998}.
 \tag{SC37}
\] The indicated remainder is still uncomputed and may contain positive, negative, and phased summands. Its lower bound, rather than local spherical module positivity alone, is the remaining all-T Mram gate. No conclusion on J14 is asserted by this refinement.

\subsection{7. Verification and carried inputs}\label{verification-and-carried-inputs}

The quantitative inputs are the exact actual K044 profiles and K046(BL8--21), K061(CG8--20), K062(AT18--24), DO058(DS4--15), and K063(PF13--25). Their proof/source records are retained. The new moment constants, absolute derivative tail, over-lattice coset count, and complete single-place bound are proved above.

The companion \nolinkurl{definite_oscillator_sharper_cutoff_certificate.py} checks the exact derivative moment products, word/twist endpoint monotonicity, all floor and sampling exponents, rational exponential/logarithm comparisons, the exact group-order gcd, target-lattice index diagnostics, and actual stratum inequalities. It does not enumerate the rational short sum, evaluate an automorphic period, or prove a sign for the two-place remainder. No Lean is run.

\subsection{Author verification record}\label{author-verification-record-8}

At 2026-11-22T09:30:40.151685+00:00, I read all final SC1--37 and the seven sections in full, together with the complete certificate. I reconstructed the exact shifted-metric bounds after all required Lie derivatives, weighted absolute L1 tail with the floor moment order, disjoint-chart and convolution transfer, target over-lattice index at both conjugate components, compact isometry-fibre count, and the fresh normalized-kernel decomposition. The subset and complete-support bounds are distinguished. No unproved sign or global occurrence assertion is used.

The exact certificate passed 2,180 assertions at 2026-11-22T09:28:47.248783+00:00. It includes every actual host inequality used here and explicit small moment/index checks. The two-or-more-place remainder remains uncomputed.

\clearpage

\textbf{Normalization correction.} Read the scalar-removal paragraph below together with the checked Branch scalar normalization erratum (BE1--5), included in this part. Dividing only the rational integer C retains the modulus of u\_T; dividing the full algebraic scalar u\_T C removes it. The displayed main operator bound is unchanged.

\section{ALP-LIFT-48: an explicit upper bound for the actual normalized coupled theta operator}\label{alp-lift-48-an-explicit-upper-bound-for-the-actual-normalized-coupled-theta-operator}

Author: GPT-6 Astra Pro, ALP\_MACHINE \nolinkurl{/root/lift_accounting}. Date: 22 November 2026 UTC. Status: author whole-proof review PASS; the companion exact diagnostic passed 33,054 assertions. Independent whole review remains separate.

This result concerns the actual two-edge family of K038, with its full prescribed Y sum, original finite input, source creation operators and terminal branch. It supplies an upper bound for its norm. It does not assert that its terminal restriction is nonzero, determine its integral Smith factors, or instantiate the old quadratic-field Host. The actual Mram period remains the nonvanishing gate isolated by DO055.

\subsection{1. Conventions and carried inputs}\label{conventions-and-carried-inputs}

Use either principal CM host of K031/K032, with CM field E, real field F, degree {[}E:Q{]}=2f, root discriminant R\_E, and actual positive trace coefficient xi. Thus xi O\_E is the absolute codifferent, bar(xi)=-xi, and q\_tau=-2i tau(xi)\textgreater0 at each chosen CM label. The actual ranks are

\[
 (f,a,b,c,t)=(42,42,1802,3562,1759)
 \quad\hbox{or}\quad(168,168,1802,3436,1633).
 \tag{OB1}
\]

Put d=t+1=b-a, r=a+1, N=c-1, and

\[
 n=ab+bc=2b^2=6,494,408,\qquad m=fn.
 \tag{OB2}
\]

Here n is the Hermitian rank of the direct sum of the two Hom spaces, and m is its complex dimension. The terminal insertion is the literal norm-one line: its lattice is L\_T direct-sum O\_E e\_c.~The pi\_T factor belongs to its covector and to the indicated G/Q frame, where it is a unit. It does not replace this norm-one Hermitian lattice by a lattice of norm T. See DO039 section1 and the source \nolinkurl{Columns/StdLatUnimodular.lean}165--176.

The following are carried, not re-proved:

\begin{itemize}
\tightlist
\item
  K030(PH20--21), K032(QH14--16): the actual principal trace polarization and codifferent identity, including positivity at every CM label.
\item
  LIFT39(PN2--5): the literal unit-vacuum theta transform is the Zak transform with its covolume-one Heisenberg metric; its norm is the Schrödinger norm. The raw CM Gaussian has its separate determinant factor. The source pin \nolinkurl{IsVacuum.norm} is \nolinkurl{IndefiniteReal/LocalNewformFunction.lean}208--226.
\item
  K037(CR17--19, CR27--35), K038(CA27--36), and DO056(AC4): the actual creation-polynomial, complete-coevaluation, arithmetic-class and CM period comparison, with the prescribed Weyl/dual-Weyl lattices.
\item
  DO029(DO29.7--19): the exact finite Heisenberg expansion in a principal lattice model, including its phases, Parseval identity and l1 bound.
\item
  K038(CA8--21), LIFT23(Q4--10): the complete raw local tensors. These statements apply to the actual frames, not just one selected value.
\end{itemize}

All norms below are the source normalized compact/Fischer norms. An orthonormal basis used to estimate a norm changes no arithmetic lattice. The complex quantities in the period comparison receive no unproved finite valuation.

\subsection{2. The unit-vacuum Fourier-square identity}\label{the-unit-vacuum-fourier-square-identity}

Let (V,B,Lambda) be a real principally polarized symplectic torus, with positive compatible complex metric H and a chosen lattice semicharacter epsilon. Let r be a homogeneous degree-D vector in its unitary Fock model. Write Theta\_r(z) for the corresponding theta section, and use its principal Hermitian line metric. Then

\[
 \|\Theta_r(z)\|^2
 =\sum_{\lambda\in\Lambda}
   \epsilon_\lambda\,
   \langle r,W(\lambda)r\rangle
   \exp(2\pi i B(\lambda,z)),\qquad |\epsilon_\lambda|=1.
 \tag{OB3}
\]

The indexing lattice is the symplectic dual of Lambda; it is Lambda itself because the polarization is principal. In particular its zero Fourier coefficient is exactly \textbar\textbar r\textbar\textbar\^{}2. There is no theta-null divisor, unrecorded covolume, or arbitrary vacuum amplitude in (OB3).

Here is a direct verification fixing the normalization. Choose a real symplectic basis of Lambda, identify its Schrödinger Lagrangian with R\^{}m, and write the unitary Zak expression as

\[
 Z_f(q,p)=e^{i\vartheta(q,p)}
       \sum_{u\in\mathbf Z^m}f(q+u)e^{2\pi i u\cdot p}.
 \tag{OB4}
\]

The trace-one semicharacter specifies the displayed unit phase; its absolute square is one. For the (s,t) Fourier coefficient of \textbar Z\_f\textbar\^{}2, integration over p first forces the difference of the two summation indices to be t. Unfolding q over all translates of {[}0,1{]}\^{}m leaves

\[
 \int_{\mathbf R^m}f(x)\overline{f(x-t)}
                             e^{-2\pi i s\cdot x}\,dx.
 \tag{OB5}
\]

The standard half-pairing in W(s,t) changes (OB5) only by its exact unit phase, which combines with the lattice semicharacter to give epsilon\_lambda. For s=t=0 the integral is \textbar\textbar f\textbar\textbar\^{}2. The metaplectic change from this symplectic basis to the actual Fock model is unitary, so (OB3) follows with its actual metric H. Polynomial Gaussian Schrödinger vectors are Schwartz. Their Zak functions and all derivatives are smooth, and the Fourier series is absolutely convergent. The following bound also proves absolute convergence directly. This argument proves (OB3) for every homogeneous r, including linear combinations within its degree, not only monomial vectors.

Set u=pi H(lambda,lambda). Normal ordering of the unitary displacement operator, in creation/annihilation coordinates with {[}a\_i,a\_j\^{}*{]}=delta\_ij, gives the following estimate. Its pi normalization can also be fixed without convention: in orthonormal Schrödinger coordinates the unit vacuum is 2\^{}(m/4) exp(-pi\textbar\textbar x\textbar\textbar\^{}2), and direct Gaussian integration gives its displacement coefficient exp(-pi H(lambda,lambda)/2). Thus the annihilator's squared length is precisely u, and

\[
 |\langle r,W(\lambda)r\rangle|
 \le \|r\|^2 e^{-u/2}
       \sum_{j=0}^D {D\choose j}\frac{u^j}{j!}
 \le\|r\|^2 e^{-u/2+2\sqrt{Du}}.
 \tag{OB6}
\]

Indeed in exp(a\^{}*) exp(-a) only equal creation and annihilation powers pair back with degree D. For an annihilator of squared length u, \textbar\textbar a\^{}j r\textbar\textbar\^{}2 \textless= (D)\_j u\^{}j \textbar\textbar r\textbar\textbar\^{}2. The coefficient is (j!)\^{}-2, giving the first sum. The inequality binom(D,j)\textless=D\^{}j/j! and binom(2j,j)\textless=4\^{}j bound this sum by cosh(2 sqrt(Du))\textless=exp(2 sqrt(Du)). Unitarity independently bounds the same matrix coefficient by \textbar\textbar r\textbar\textbar\^{}2. If u\textgreater=64D, then 2 sqrt(Du)\textless=u/4. Consequently

\[
 \|\Theta_r\|_\infty^2\le\|r\|^2
 \left(\#\{\lambda:H(\lambda,\lambda)\le64D/\pi\}
             +\sum_{\lambda\in\Lambda}
                         e^{-\pi H(\lambda,\lambda)/4}\right).
 \tag{OB7}
\]

Counting the transition shell in both terms only enlarges the bound.

\subsection{3. A uniform lower lattice length for every actual class}\label{a-uniform-lower-lattice-length-for-every-actual-class}

For any positive selfdual Hermitian O\_E lattice L, including any actual Hom lattice, alpha=h(x,x) is a totally positive algebraic integer when x is nonzero. The Hermitian Hom form is trace(X\^{}\emph{X): integrality follows because X\^{}} is integral on selfdual lattices and the trace of an integral projective-module endomorphism lies in O\_E. It is fixed by conjugation and therefore lies in O\_F. Its field norm is a positive integer. The codifferent identity gives

\[
 \prod_{\tau\in\Phi}q_\tau
       =2^f\,|\operatorname{disc}E|^{-1/2}.
 \tag{OB8}
\]

Applying arithmetic-geometric mean to the positive numbers q\_tau tau(alpha) proves, for every nonzero vector,

\[
 H(x,x)=\sum_\tau q_\tau\tau(\alpha)
 \ge f\left(2^f|\operatorname{disc}E|^{-1/2}
                         N_{F/\mathbf Q}(\alpha)\right)^{1/f}
 \ge {2f\over R_E}=:s^2.
 \tag{OB9}
\]

This bound has no factor depending on the Hermitian rank. Both Hom lattices are selfdual because their source and target lattices are; their direct sum has this same property and the same bound. Its principal trace lattice has real dimension 2m. No shortest-vector estimate for an unrelated standard lattice is being substituted.

Disjoint balls of radius s/2 centered on lattice points yield

\[
 \#\{\lambda:\|\lambda\|\le R\}\le(1+2R/s)^{2m}.
 \tag{OB10}
\]

For the Gaussian sum, integrate this count by layers. The inequality (1+x)\^{}(2m)\textless=2\^{}(2m-1)(1+x\^{}(2m)) and a gamma integral give

\[
 \sum_\lambda e^{-\pi\|\lambda\|^2/4}
 \le2^{2m-1}\left[1+
             \left({8R_E\over\pi f}\right)^m m!\right]
 \le(16nR_E)^m.
 \tag{OB11}
\]

For the last inequality use pi\textgreater2, m!\textless=m\^{}m, m=fn, and nR\_E\textgreater=1: the expression in brackets is at most 2(4nR\_E)\^{}m.

On J14 take C=20, and on J56 take C=28. The exact root-discriminant bounds R\_E\textless2\^{}8 and R\_E\textless2\^{}11 give 128R\_E/(pi f)\textless C\^{}2. Equations (OB7--11) therefore prove the explicit uniform estimate

\[
 \boxed{\quad
 \|\Theta_r\|_\infty\le U_E(D)\|r\|,\qquad
 U_E(D)=\bigl[(C+1)^2(D+1)+16nR_E\bigr]^{m/2}.
 \quad}\tag{OB12}
\]

For the first counting term, (1+C sqrt(D))\^{}2 \textless= (C+1)\^{}2(D+1). Then combine its m-th power with (OB11) using A\^{}m+B\^{}m\textless=(A+B)\^{}m. This proves the inequality also at D=0.

For every word and every allowed gap in K044, its EG10 entry bounds give D=sum\_tau(\textbar nu\_0,tau\textbar+\textbar nu\_1,tau\textbar)\textless2A f(a+b)\textless2\^{}980, where A=2\^{}960. Since D is an integer, D+1\textless=2\^{}980. Also C+1\textless=32 and 16nR\_E\textless2\^{}38. Thus

\[
 \log_2 U_E(D)<496m,
 \quad m=272,765,136\text{ or }1,091,060,544.
 \tag{OB13}
\]

The respective bounds 496m are less than 2\^{}37 and 2\^{}39. These are deliberately coarse but uniform explicit bounds. The dependence on the creation degree is logarithmic for fixed lattice dimension. This does not bound the separate hook, CM-period, or finite-input factors.

\subsection{4. Exact finite inputs and their Heisenberg coefficients}\label{exact-finite-inputs-and-their-heisenberg-coefficients}

Work at each finite place in the principal lattice model attached to the actual class lattices. If a function Phi is supported on L\_plus, periodic by L\_minus, with L\_minus contained in L contained in L\_plus, put R={[}L\_plus:L{]}, N0={[}L:L\_minus{]}, and H0=max\textbar Phi\textbar. DO029's exact finite Fourier expansion in the unit-vacuum translates has

\[
 \Phi=\sum_{a,y}c_{a,y}W_{-a}W_y\phi_L,
 \qquad \sum_{a,y}|c_{a,y}|\le R\sqrt{N_0}\,H_0.
 \tag{OB14}
\]

Its coefficients include the half-pairing and trace-one phases. The bound follows from Parseval on each L\_plus/L coset and Cauchy--Schwarz for its N0 Fourier coefficients. The measure and normalization are those in which phi\_L is a unit vector. Equivalently, divide by the value of the normalized lattice vacuum before describing the displayed finite function. This is the exact principal-model normalization of DO029, not an assertion that an arbitrary raw characteristic function has norm one in an unrelated Haar measure.

At G use the fixed split isometry from the literal lamY/entB lattice to the CA13 normalized lattice, as in LIFT33's final frame clarification. This realizes the underlying selfdual class lattice; it does not replace it by a different genus. Subsequent integral frame changes contribute only the honest unit characters CA4--5. At T and Q use their actual integral adapted frames. Thus the raw functions cited below are the coefficients relative to these very vacua. The global unit-vacuum conversion remains the separate factor in (OB20), and is not applied again here. An adapted symplectic basis can have a different raw period determinant from the fixed host presentation. Its metaplectic map is unitary and carries the fixed unit vacuum; no such raw determinant is separately identified with det(2Y\_J) or inserted a second time in (OB20).

For the direct sum of the two horizontal Schrödinger spaces there are n=ab+bc coordinates at one split F-place. The following is an upper bound for the full local tensor after its prescribed Y sum:

\[
\begin{array}{c|c|c|c|c}
\text{place}&R&N_0&H_0&\text{bound in (OB14)}\\\hline
T&q^0&q^n&1&q^{n/2}\\
G&q^0&q^n&1&q^{n/2}\\
Q&q^b&q^{n-b}&q^{-d+1/2}&q^{(n+a-t)/2}.
\end{array}\tag{OB15}
\]

Here q=T at T. For T both literal incidence functions are supported on integral matrices, are periodic modulo q, and have absolute value at most one. There is no T member of the Y set. For G, CA15 forces the unique class y=q\^{}-1 e\_a modulo O\^{}b; CA16--18 give integral support and q-periodicity, including the remaining additive phase in the integral last row. Its absolute value is at most one. It would be incorrect to multiply this answer by the cardinality of the redundant Y parameter set.

For Q, use the complete Fourier contraction LIFT23(Q10) before estimating. In its notation the first matrix X and the first c-1 rows Z of the second matrix are integral and periodic modulo q. The last row eta lies in q\^{}-1 O\^{}b and is periodic modulo O\^{}b. Hence R=q\^{}b and N0=q\^{}(ab+(c-1)b)=q\^{}(n-b). The nontrivial tame character has Gauss sum of modulus sqrt(q) under every complex embedding. Its nontriviality is literal: Qprime is the ramification set of the tuple, by DO047(QN22--23) and Tuples.lean128--129; the trivial tuple has empty Q set. Its exact scalar is q\^{}-d G(chi,psi)chi(-c0)\^{}-1, so H0=q\^{}-d+1/2. Substituting into (OB14) gives b+(n-b)/2-d+1/2=(n+a-t)/2. This retains the full Y family and is independent of its auxiliary depth M\textgreater=1. The actual Q places are degree-one unramified F-places by DO037 section1's direct DegOne audit, and pi\_T is a unit there. The raw Gauss and character phases remain in the exact coefficients.

Outside G,T,Q the actual finite vector is the principal vacuum. Taking products of (OB15), over the actual selected half-places, therefore gives

\[
 C_{\rm fin}=
 T^{fn/2}(g_0g_1)^{fn/2}
             \prod_{v\in Q}q_v^{(n+a-t)/2}.
 \tag{OB16}
\]

There are f selected T places and f places over each g\_j. The Q product is over the actual tuple, with its actual multiplicity. No unknown cardinality has been replaced by f.~If the common datum has no Q places, that product is one. Finite splitting characters have unit absolute value; their algebraic coefficients and denominators remain in the exact expansion. Since all these places are different from11, the displayed finite Fourier denominators do not introduce an11-denominator. This observation does not remove the pre-existing algebraic CM/frame factors.

The tensor expansion (OB14) is an expansion on the direct sum of both Schrödinger spaces, not a separated estimate for an arbitrarily chosen Y term. A finite Heisenberg coset may be represented by a rational translation by the lattice Chinese remainder theorem. Global rational Heisenberg invariance identifies its theta value with Theta\_r at the translated real torus point, with the exact unit central phase. Thus (OB12) and (OB16) give for every homogeneous r, uniformly over every actual class triple,

\[
 |\Theta_{\Psi,r}(0)|\le C_{\rm fin} U_E(D)\|r\|.
 \tag{OB17}
\]

The passage to translated points uses the principal theta line metric in (OB3). It is not a bound on an uncorrected holomorphic trivialization whose translation multipliers have varying modulus.

\subsection{5. The complete selected polynomial and exact norm factors}\label{the-complete-selected-polynomial-and-exact-norm-factors}

Let V\_j and W\_j be the source compact coefficient space and its paired opposite realization. Put

\[
 H_j=\prod_{\tau}H_{\nu_{j,\tau}},\qquad
 d_1=\dim V_1.
 \tag{OB18}
\]

By K037(CR19), iota\_j is sqrt(H\_j) times an isometry on V\_j tensor W\_(j+1), with the actual source creation labels. Choose orthonormal pairBK-dual middle bases (v\_i,w\_i). The complete polynomial associated to source P and terminal test Z is

\[
 Q(P,Z)=\sum_{i=1}^{d_1}
       \iota_0(P,w_i)\otimes\iota_1(v_i,B_T^*Z),
 \qquad
 \|Q(P,Z)\|^2=d_1H_0H_1\beta_B\|P\|^2\|Z\|^2.
 \tag{OB19}
\]

Indeed both tensor factors are orthogonal as i varies, with respective squared norms H0\textbar\textbar P\textbar\textbar\^{}2 and H1\textbar\textbar B\_T\^{}\emph{Z\textbar\textbar\^{}2. The actual branch obeys B\_T B\_T\^{}}=beta\_B I. Consequently sqrt(beta\_B)=\textbar b\_Br\textbar, where b\_Br=product\_tau u\_T,tau cBr\_tau sqrt(kappa\_tau), with its actual unit phase retained. This is a complete coevaluation. It has no 1/d1 factor and invokes no inverse Fischer Gram matrix. Its total homogeneous degree is D=D0\_sum+D1\_sum.

The exact source unit-vacuum/CM scalar is DO056(AC4):

\[
 \mathfrak M_\sigma=
 \det(2Y_J)^{7hb^2}
 C_\theta^{14h(ab+bc)}
 \prod_\tau(\lambda_\tau/cD_\tau)^{D_{0,\tau}+D_{1,\tau}}.
 \tag{OB20}
\]

Here h is the actual host multiplicity used there, and sigma indicates the chosen complex realization. The lambda\_tau are the literal source normalized creation scales; cD\_tau are the actual algebraic coordinate scales. K037/K038 and DO056 identify the unit-vacuum theta coefficient for the complete finite tensor with M\_sigma times the arithmetic complete coefficient. No finite level volume occurs in this pointwise class identity. Those volumes arise only on passing to adelic integrals. In particular the det(2Y\_J) factor in (OB20) is applied exactly once.

Combining (OB17--20) proves the classwise estimate

\[
 |K_{efd}(P_e,Z_d)|\le
 {C_{\rm fin}U_E(D)\sqrt{d_1H_0H_1}\,|b_{\rm Br}|
                  \over|\mathfrak M_\sigma|}
                         \|P_e\|\|Z_d\|.
 \tag{OB21}
\]

K\_efd is the entire literal Y-summed two-edge arithmetic coefficient at the indicated initial, middle, and terminal classes, after the literal terminal tilQ insertion and brC branch. This inequality does not require its value to be nonzero.

\subsection{6. The actual mass-weighted operator bound}\label{the-actual-mass-weighted-operator-bound}

Let Gamma\_0,e, Gamma\_1,f and Gamma\_T,d be the actual finite stabilizers. Set

\[
 \mathsf m_j=\sum_{c\in\mathcal C_j}|\Gamma_{j,c}|^{-1},
 \qquad
 \|P\|_0^2=\sum_e{\|P_e\|^2\over|\Gamma_{0,e}|},
 \quad
 \|Z\|_T^2=\sum_d{\|Z_d\|^2\over|\Gamma_{T,d}|}.
 \tag{OB22}
\]

The actual bilinear arithmetic form is K038(CA27--30)

\[
 \mathcal B(P,Z)=
 \sum_{e,f,d}{K_{efd}(P_e,Z_d)\over
             |\Gamma_{0,e}||\Gamma_{1,f}||\Gamma_{T,d}|}.
 \tag{OB23}
\]

Apply (OB21), sum the middle mass, and apply Cauchy--Schwarz separately to the initial and terminal weighted sums. The result is

\[
 \boxed{\quad
 |\mathcal B(P,Z)|\le \mathcal A_\sigma\|P\|_0\|Z\|_T,
 \quad
 \mathcal A_\sigma=
 {C_{\rm fin}U_E(D)\sqrt{d_1H_0H_1}\,|b_{\rm Br}|
                    \over|\mathfrak M_\sigma|}
       \mathsf m_1\sqrt{\mathsf m_0\mathsf m_T}.
 \quad}\tag{OB24}
\]

This is valid on the entire initial coefficient space, hence on literal Mram and on any actual subspace or image. Character conditions restrict the space but do not replace its groupoid measure by a free orbit count.

Equivalently, if m\_l denotes the standard definite genus mass of (O\_E\^{}l,I\_l), and v0,v1,vT are the actual strict finite level volumes with maximal compacts of volume one, mass refinement gives

\[
 \mathsf m_0=m_a/v_0,\quad
 \mathsf m_1=m_b/v_1,\quad
 \mathsf m_T=m_N/v_T,
 \qquad
 \mathsf m_1\sqrt{\mathsf m_0\mathsf m_T}
 ={m_b\sqrt{m_am_N}\over v_1\sqrt{v_0v_T}}.
 \tag{OB25}
\]

This is an exact alternative expression, not another level index to multiply into (OB24). The volumes are the explicit strict T/G/Q volumes in DO056(AC14, AC23--24). The genus masses remain the actual ones; their standard rank ratios are explicitly

\[
 {m_l\over m_{l-1}}=
 \left[{(l-1)!R_E^{l-1/2}\over(2\pi)^l}\right]^f
       L_F(l,\eta^l),\qquad l\ge3,
 \tag{OB26}
\]

as in DO044(DO44.22) and the stipulated source \nolinkurl{HermitianLattices/StatementMass.lean}27--51,78--89. Formula (OB24) can be used with its actual finite mass sums without making a numerical claim about the remaining base mass m2. It removes an unknown theta-operator norm; it does not claim to have evaluated every class mass or every DP44 datum.

\subsection{7. Arithmetic normalization and the resulting cost component}\label{arithmetic-normalization-and-the-resulting-cost-component}

For every actual nonzero algebraic row supplied conditionally by K059/DO055, LIFT47 supplies adapted integral field-spanning input and terminal generators with norm product at most

\[
 2^{K_{\rm fin}}\widehat B_0\widehat B_T.
 \tag{OB27}
\]

The same actual vector may therefore be chosen with the bound

\[
 |\sigma\mathcal B(P,Z)|\le
 \mathcal A_\sigma\,
                       2^{K_{\rm fin}}\widehat B_0\widehat B_T.
 \tag{OB28}
\]

All quantities are those of that complex realization. Taking logarithms produces a sufficient upper bound for this archimedean coefficient cost, with log U bounded by (OB13). Actual finite valuations of the algebraic row must still be combined with the exact source character field/discriminant costs and gains. The CM period in (OB20) has not been silently replaced by a unit scalar, nor is its logarithm assigned an independent algebraic height.

If the common rational branch scalar product\_tau cBr\_tau is divided out of the arithmetic map, its archimedean contribution \textbar b\_Br\textbar{} in (OB24) becomes product\_tau sqrt(kappa\_tau). The finite image content shifts by the valuations of that same rational scalar. This is the normalization invariance already proved in R024 and the endpoint content calculation; it is not a new gain. Likewise the signed T\^{}(Lambda0/2) in LIFT47 comes from the actual initial lattice inclusion. It cannot simply be dropped while retaining the old finite image bound. The separate DO056(AC19) cancellation of the Sad norm with strict T volumes remains valid when comparing the actual global period; the T factors in (OB16) bound a finite theta input and are not those already canceled volumes.

\subsection{8. What has and has not been computed}\label{what-has-and-has-not-been-computed}

The new assertions are the all-polynomial unit-vacuum estimate (OB3--13), the complete fixed-input bound (OB15--17), and their source-normalized arithmetic operator consequence (OB24). They apply uniformly to every actual class and every allowed K044 word/gap. They neither enlarge the finite Schwartz family nor replace a selected source polynomial by an arbitrary real vector.

No new nonvanishing assertion is made. The actual all-T Mram projection and its moved bottom period remain separate. No theta image saturation, exact primitive minor, full old Host, or numerical net DP44 improvement is inferred. The CM period, hook factors, finite Fourier costs, class masses and branch factor are displayed so a later comparison can retain their cancellations honestly.

The companion \nolinkurl{actual_theta_operator_certificate.py/.json} passed 33,054 exact checks: Fock annihilation contractions, number-state Laguerre majorants, finite Parseval, full coupled Q Fourier sums, coevaluation norms, nonfree masses, support exponents, and both host inequalities. These finite diagnostics do not evaluate the actual high-rank theta series or its automorphic periods.

\clearpage

\section{AM-K065 --- one transverse lift per residue class and a larger coefficient field}\label{am-k065-one-transverse-lift-per-residue-class-and-a-larger-coefficient-field}

Author: GPT-6 Astra Pro, ALP\_MACHINE \nolinkurl{/root/kummer_extension}. 22 November2026 UTC. Inherited reasoning effort is not separately exposed.

\textbf{Author whole-proof check PASS --- 2026-11-22T09:30:43.184134+00:00.} The complete saved TL1--TL7 proof and full certificate were read and reconstructed; the independent finite arithmetic component passed30 exact assertions. No actual rational or automorphic sum was evaluated.

This short refinement keeps the original K062 cutoff and normalization. It bounds the full fixed plane over Q(zeta\_28) on J14 by2\^{}-794, and the full fixed plane over Q(zeta\_28,sqrt(5),sqrt(13)) on J56 by2\^{}-135. The latter field is the whole degree48 CM subfield complementary to the new degree-seven real field. Matrices using that degree-seven extension, and non-plane matrices, remain outside this bound.

\subsection{1. The small transverse coordinate}\label{the-small-transverse-coordinate}

Use the actual u,U,Q,H and plane coordinates of K063/K064. The orthonormal plane basis gives the stronger bound retained in K063:

\[
       |\sigma(\beta)|\le H/\sqrt U,
       \qquad |\sigma(b_0)|\le2H/\sqrt U,
       \qquad b_0=2\beta.
 \tag{TL1}
\]

For both actual hosts, exact integer arithmetic gives

\[
                             U>16H^2.
 \tag{TL2}
\]

Two possible b0 values in the same residue class modulo O\_L differ by an algebraic integer. If they are distinct, its absolute field norm is a positive integer. But TL1--TL2 would give every conjugate of their difference absolute value strictly below4H/sqrt(U)\textless1, contradicting that norm. Therefore \textbf{there is at most one admissible b0 lift in each residue class}. This argument is valid in any of the CM subfields under consideration and does not require a small coset representative.

The alpha coordinate still has at most(3H)\^{}(2r) lifts, by K064(NR10). The exact paired-prime norm congruence is unchanged and has N2(T)=T\^{}6+T\^{}5-T\^{}3 possibilities at each pair of primes. Thus K064's count improves to

\[
             28(3H)^{2r}N_2(T)^b
 \tag{TL3}
\]

for a fixed nonintegral support of b selected L-primes.

\subsection{2. Extension to the actual degree48 subfield}\label{extension-to-the-actual-degree48-subfield}

On J56 take

\[
          L_{48}=\mathbb Q(\zeta_{28},\sqrt5,\sqrt{13}).
 \tag{TL4}
\]

K032's actual disjoint compositum has degree48. Its CM involution fixes sqrt5 and sqrt13; it contains k\_-=Q(sqrt(-7)). The full actual host is L48 B with {[}B:Q{]}=7 and intersection L48 intersect B=Q. Consequently r=24, f=168 and m={[}calE:L48{]}=7.

For completeness, all hypotheses of the preceding counting argument continue to hold for this larger subfield. The actual T splits completely in the full host and hence in L48. Each selected L48-prime above pi has exactly seven lifts. The actual z and the row q\_* in K063(PF12) are defined over k\_-, so the complete labelled pencil, including both omitted points, descends at every such lift. Matrices with L48 coefficients therefore have identical local rho values at the seven lifts. Actual support gives the same paired order-two pole bound, and the full host root group still bounds all determinant possibilities by28. The unitary plane shape and the inequalities on alpha and beta are unchanged because u and U are rational.

The paired fractional-ideal CRT and singular-matrix count of K064 therefore apply without a global generator or any change of local level. On J14 use the already covered L12=Q(zeta\_28), with r=6 and m=7.

\subsection{3. Complete sub-sum bound and exact numbers}\label{complete-sub-sum-bound-and-exact-numbers}

For these fields, and for every previously admitted smaller coefficient field, put y=4\^{}m N2(T) T\^{}(-m). Multiplying TL3 by the unchanged local factor (4/T)\^{}(mb) and summing over nonempty supports gives

\[
 \boxed{\quad
 \sum_{g\text{ in the indicated plane subfamily}}
       |a(g)c(g)\prod_v\rho(g_v)|
 \le28(3H)^{2r}\bigl((1+y)^r-1\bigr).
 \quad}\tag{TL5}
\]

All determinant components, finite partial projectors, compact types and the central quotient are exactly as in K063/K064. In particular no positivity of a rational summand is assumed.

Use T\textgreater2\^{}(b\_T), H=2\^{}(h\_H), N2(T)\textless2T\^{}6, 3\textless4 and28\textless32. With y\textless1, the exact same binomial estimate gives

\[
 |S_L|<2^{-B''_L},\qquad
 B''_L=(m-6)b_T-2m-6-5r-2rh_H.
 \tag{TL6}
\]

For J14, (r,m,b\_T,h\_H)=(6,7,1000,13), so B'`=794. For J56 and L48 these data are(24,7,995,15), so B'`=135. For comparison, keeping only L12 on J56 gives B''=21618. The certificate also compares the exact rational TL5 to all these bounds, without rounded logarithms.

The actual R\_minus of K064(NR15--17) lies in each bounded family. Thus the removed index set is nonempty in the unchanged short support. No individual real matrix coefficient is asserted to be nonzero.

\subsection{4. Exact remaining gate}\label{exact-remaining-gate}

Split the original S\_T into its L12-plane part on J14 or its L48-plane part on J56, and call the complementary sum S\_rest. Keeping the original positive class/central/selected normalization calN,

\[
 \frac{E_{x,\mathrm{ram}}}{\mathcal N}
       =1+S_{\rm rest}+\varepsilon,\qquad
 |\varepsilon|<2^{-1999}+2^{-B},\quad
 B=794\ (J14),\quad135\ (J56).
 \tag{TL7}
\]

This enlarged partition is inverse closed, so the two sub-sums remain real. The explicit sufficient remaining inequality is S\_rest \textgreater{} -1+2\^{}(-1999)+2\^{}(-B). It is not proved here.

The J56 bound removes more matrices than K064 at the price of a larger error allowance. Both statements remain available, with their different explicit partitions. This result does not classify arbitrary short matrices as plane matrices, eliminate dependence on B in the full family, prove all-place Mram nonvanishing, or establish the DP44 budget.

\subsection{Verification scope}\label{verification-scope-4}

The carried source/proof inputs are K032's actual field construction, K063's actual labelled-pencil and all-determinant plane formulas, K064's paired fractional-ideal norm congruence and count, and K062's original normalized sum. The new uniqueness of the transverse lift is proved in section1 and the larger field hypotheses are checked in section2. The certificate verifies exact actual inequalities and dyadic bounds only; it does not enumerate a rational or automorphic sum.

\clearpage

\section{AM-R-040 --- the combined larger-field plane and single-place remainder}\label{am-r-040-the-combined-larger-field-plane-and-single-place-remainder}

Author: GPT-6 Astra Pro, ALP\_MACHINE /root. 22 November 2026 UTC. Reasoning effort is not externally exposed.

This is an explicit corollary of the checked K065 and DO059 bounds. It changes no kernel or source datum. It records the stronger constants available when K065's enlarged coefficient field is used at DO059's smaller J56 cutoff. The remaining finite sum is not evaluated.

\subsection{1. The same plane estimate at the smaller cutoff}\label{the-same-plane-estimate-at-the-smaller-cutoff}

On J56 let \(L=\mathbf Q(\zeta_{28},\sqrt5,\sqrt{13})\), so \([L:\mathbf Q]=48\), \(r=24\), and \([E:L]=m=7\). Put \(H=8192=2^{13}\). The actual \(T\) still satisfies \(T>2^{995}\). K065's proof uses only \(H\ge1\), \(U>16H^2\), and the exact short-region bound at each embedding. Its original \(U>16(32768)^2\) implies the required inequality at this smaller H. Every pencil, field-descent, paired CRT and transverse uniqueness assertion remains unchanged. Therefore that proof, with the smaller height ball, gives

\[
 \sum_{g\in\mathcal P_{48}(H)}
       |a(g)c(g)\prod_v\rho(g_v)|
 \le 28(3H)^{48}\bigl((1+y)^{24}-1\bigr),\qquad
 y=4^7(T^6+T^5-T^3)T^{-7}.
 \tag{R40.1}
\]

Here \(\mathcal P_{48}(H)\) is precisely the original nonintegral, supported fixed-plane index set, in the central quotient, with matrix coefficients in L and every archimedean operator norm at most H.

As in K065(TL6), use \(T^6+T^5-T^3<2T^6\), \(3<4\), \(28<32\), and \((1+y)^{24}-1<2^{24}y\) for \(0<y<1\). The exponent is now

\[
 (7-6)995-2\cdot7-6-5\cdot24-2\cdot24\cdot13=231.
 \tag{R40.2}
\]

Thus the left side of R40.1 is strictly less than \(2^{-231}\). The companion certificate also verifies the exact rational R40.1 at the actual T, rather than only this sufficient integer exponent.

\subsection{2. Combination with the complete single-place estimate}\label{combination-with-the-complete-single-place-estimate}

DO059(SC22,SC35) gives a fresh decomposition of the unchanged J56 ramified energy at this same H, with integral-plus-infinite error below \(2^{-1999}\), and bounds the sum of absolute values of every term nonintegral at exactly one relative T place by \(2^{-104546}\). Every nonintegral matrix defined over L has its nonintegrality repeated at the seven selected E-primes over a selected L-prime, by K063(PF12--13) and K065's field check. Hence the two removed index sets are disjoint. Let \(S_{56}^{\rm rem}\) contain the remaining supported nonintegral terms at H, outside the full L-defined fixed plane, with at least two nonintegral relative T places. Then

\[
 \frac{E_{x,\mathrm{ram}}}{\mathcal N}
       =1+S_{56}^{\rm rem}+\varepsilon_{56},\qquad
 |\varepsilon_{56}|<2^{-1999}+2^{-104546}+2^{-231}<2^{-230},
 \tag{R40.3}
\]

where the exact original positive normalizer is

\[
 \mathcal N=w_x\,28C_\infty d_Rd_W
                      \left(\frac{T-2}{T-1}\right)^f.
 \tag{R40.4}
\]

No term from the old cutoff annulus is omitted: DO059's complete tail already contains it. The sets are inverse closed, so the indicated remaining sum is real. In particular a sufficient remaining condition is

\[
 S_{56}^{\rm rem}>-1+2^{-1999}+2^{-104546}+2^{-231}.
 \tag{R40.5}
\]

This inequality has not been proved. Its coarser sufficient version with right side \(-1+2^{-230}\) is also available.

On J14, keep the original H8192 and K065's full \(\mathbf Q(\zeta_{28})\)-plane bound \(2^{-794}\). If \(S_{14}^{\rm rem}\) is the complementary sum, the corresponding same-kernel identity and sufficient threshold are

\[
 \frac{E_{x,\mathrm{ram}}}{\mathcal N}
      =1+S_{14}^{\rm rem}+\varepsilon_{14},\qquad
 |\varepsilon_{14}|<2^{-1999}+2^{-794},\qquad
 S_{14}^{\rm rem}>-1+2^{-1999}+2^{-794}.
 \tag{R40.6}
\]

Only the first two assertions in R40.6 are established; the last is the unproved sufficient condition. There is no J14 single-place estimate in this corollary.

\subsection{3. Scope and verification}\label{scope-and-verification}

The original index sets, two omitted pencil points, all determinant phases, finite normalized partial traces, selected real traces and class masses remain those of K062. K064(NR15--17) supplies an actual nonintegral Rminus index in the removed plane, without claiming a nonzero real trace for that index. The larger-field, non-plane remainder is still the nonvanishing gate. No actual period, arithmetic theta value, full mixed Host, net DP44 sign or Kummer RH is asserted.

The proof is the literal substitution in the already proved absolute plane count and the disjoint partition just given. The exact diagnostic checks only the actual integer and rational comparisons and disjoint support multiplicity. It does not enumerate the rational remainder.

\clearpage

\textbf{Normalization correction.} Read the scalar-removal paragraph below together with the checked Branch scalar normalization erratum (BE1--5), included in this part. Dividing only the rational integer C retains the modulus of u\_T; dividing the full algebraic scalar u\_T C removes it. The displayed main operator bound is unchanged.

\section{ALP-LIFT-49: explicit actual genus-mass and fixed operator costs}\label{alp-lift-49-explicit-actual-genus-mass-and-fixed-operator-costs}

Author: GPT-6 Astra Pro, ALP\_MACHINE \nolinkurl{/root/lift_accounting}. 22 November 2026 UTC. Author whole-proof check PASS; the companion exact diagnostic passed17,731 assertions. Independent whole review remains separate.

This result supplies the genus-mass upper bounds left explicit in LIFT48 and combines them with that theorem's actual finite inputs and strict level volumes. It never replaces an actual genus mass by a broader count in an equality. The counting argument bounds all positive unimodular Hermitian O\_E lattices of the indicated rank, which contains the original standard rational genus. The final consequence is a sufficient upper bound for the actual arithmetic operator. It does not prove its Mram nonvanishing or a net DP44 saving.

\subsection{1. Bounded Gram matrices for the original genus}\label{bounded-gram-matrices-for-the-original-genus}

Let E/F be either actual host, {[}E:Q{]}=2f, with root discriminant R\_E. For a positive unimodular Hermitian O\_E lattice L of rank n, the LIFT47 transference construction gives E-independent vectors y\_1,\ldots,y\_n in L with

\[
 \sum_{\tau\in\Phi}h_\tau(y_i,y_i)\le R_n^2,
 \qquad R_n={11\over2}n\sqrt f\,R_E.
 \tag{GM1}
\]

They may be prescribed to form a basis modulo11, though that additional property is not needed for the count. The construction works for every positive unimodular lattice: its proof uses selfduality and the absolute different, not its membership in one rational Hermitian isometry class. We therefore count a possibly larger set than the actual standard genus, which is legitimate for an upper bound.

Put L0=sum O\_E y\_i. Its Gram matrix G=(h(y\_i,y\_j)) is integral Hermitian and positive definite at every embedding. Complex Cauchy--Schwarz and (GM1) give

\[
                |\tau G_{ij}|\le R_n^2
                       \quad\hbox{at every chosen embedding.}
 \tag{GM2}
\]

A nonzero algebraic integer of E has at least one complex absolute value at least one: otherwise its nonzero integral absolute norm would be smaller than one. Thus distinct O\_E points have distance at least one in maximum-complex norm. Packing open balls of radius 1/2 bounds the number of O\_E elements with all conjugates at most R\_n² by (1+2R\_n²)\^{}(2f). The identical argument in O\_F has real dimension f and bound (1+2R\_n²)\^{}f.~There are n real diagonal entries and n(n-1)/2 independent complex off-diagonal entries. Consequently the number of possible G is at most

\[
                         (1+2R_n^2)^{fn^2}.
 \tag{GM3}
\]

The positivity and other relations only reduce this overcount.

\subsection{2. The discriminant module and its submodules}\label{the-discriminant-module-and-its-submodules}

Identify L0 with O\_E\^{}n using the ordered vectors, and write the Hermitian form as h(x,y)=bar(x)\^{}t G y. Its Hermitian dual is exactly G\^{}-1 O\_E\^{}n.~Reversing the convention conjugates this presentation and changes none of the cardinalities. Since L is selfdual and L0 is contained in L,

\[
              L_0\subset L\subset L_0^\#,
 \qquad A_G=L_0^\#/L_0.
 \tag{GM4}
\]

Its cardinal is \textbar N\_(E/Q) det G\textbar. If I={[}L:L0{]}, selfduality gives {[}L0\#:L{]}=I: the two finite quotients are paired perfectly by the Hermitian form, with conjugation on one factor. Equivalently this is the equality of the corresponding conjugate local elementary divisor lengths. Hence

\[
 |A_G|=I^2=\bigl(N_{F/\mathbf Q}\det G\bigr)^2.
 \tag{GM5}
\]

For each embedding, Hadamard bounds det(tau G) by product\_i h\_tau(y\_i,y\_i). Multiplying over tau and applying arithmetic-geometric mean separately for every i in (GM1) gives

\[
 I=N_{F/\mathbf Q}\det G
    \le(R_n^2/f)^{fn},\qquad
 |A_G|\le(R_n^2/f)^{2fn}.
 \tag{GM6}
\]

We use the following elementary count, retaining O\_E-linearity.

\textbf{Finite-module lemma.} Every O\_E-submodule of A\_G is generated by at most n elements. Thus A\_G has at most \textbar A\_G\textbar\^{}n submodules.

\textbf{Proof.} At a prime P, the preimage of a submodule H\_P in the free rank-n lattice (L0\#)\_P is a torsion-free module over the DVR (O\_E)\_P and hence a free module of rank n.~Images of a basis generate H\_P. Only finitely many P occur in the finite module H. Its primary decomposition is a direct sum, and its annihilator is a product of prime powers. For j=1,\ldots,n take the sum of the j-th generators of the primary components, padding with zeros when necessary. The Chinese remainder idempotents recover each component from these n global elements. They therefore generate H over O\_E. Every submodule is the span of some ordered n-tuple in A\_G, proving the cardinal bound. No principality of a global ideal, free global overlattice basis, or class-number assumption is being made. QED.

An overlattice L in (GM4) is determined by its O\_E-submodule L/L0. Counting all submodules, including the nonintegral and nonselfdual ones, only enlarges the bound. An identical pair (G,L/L0) determines an isometry class: use the indicated E-basis to identify its Hermitian space with (E\^{}n,G). Thus all isometry classes are covered by the pairs just counted. Their finite automorphism weights are at most one. In particular the actual standard genus mass m\_n obeys

\[
 \boxed{\quad
 m_n\le
 (1+2R_n^2)^{fn^2}(R_n^2/f)^{2fn^2}.
 \quad}\tag{GM7}
\]

Formula (GM7) is a direct upper bound and does not invoke a mass formula. It holds even if the broader count includes other rational Hermitian spaces. It is not an equality with the original mass.

\subsection{3. Actual rank-two constants and mass ratios}\label{actual-rank-two-constants-and-mass-ratios}

At n=2, R\_2²=121 f R\_E². Use the exact root-discriminant bounds R\_E\textless2\^{}r with (f,r)=(42,8) or(168,11). Since121\textless2\^{}7,

\[
 |A_G|<2^L,\quad L=4f(7+2r),\qquad
 1+2R_2^2<2^g,
 \tag{GM8}
\]

where g=30 or38 respectively. The finite-module lemma uses n=2, so (GM3) and (GM8) give

\[
 \boxed{\quad
 m_2<2^{C_2},\qquad
 C_2=4fg+2L=
 \begin{cases}12,768,&J^{14},\\64,512,&J^{56}.
 \end{cases}\quad}
 \tag{GM9}
\]

For larger standard ranks use the same accepted actual Siegel mass ratio as DO044(DO44.22), LIFT48(OB26), and \nolinkurl{HermitianLattices/StatementMass.lean}27--51,78--89:

\[
 {m_l\over m_{l-1}}=
 \left[{(l-1)!R_E^{l-1/2}\over(2\pi)^l}\right]^f
                    L_F(l,\eta^l),\qquad l\ge3.
 \tag{GM10}
\]

This equality concerns exactly the standard lattices I\_l and I\_(l-1), not the broader class count in section2. Both actual relative CM extensions are finite-unramified, as already checked in K030/K032. The positive Euler product satisfies L\_F(l,eta\^{}l)\textless=zeta\_F(l)\textless=zeta(l)\^{}f\textless2\^{}f.

All required ranks are below4096. Hence (l-1)!\textless=2\^{}(12(l-1)), 2pi\textgreater4 and R\_E\textless2\^{}r give the convenient integer bound

\[
 \log_2{m_l\over m_{l-1}}
       <f\bigl((r+10)l-c_r\bigr),
       \qquad c_r=11+\lfloor r/2\rfloor.
 \tag{GM11}
\]

Define

\[
 S_n=f\left[(r+10)\left({n(n+1)\over2}-3\right)
                         -c_r(n-2)\right].
 \tag{GM12}
\]

Telescoping (GM10), and using only the upper bound (GM9), proves log2 m\_n\textless C2+S\_n.~The exact integer values needed by LIFT48 are

\begin{longtable}[]{@{}
  >{\raggedright\arraybackslash}p{(\columnwidth - 8\tabcolsep) * \real{0.1579}}
  >{\raggedleft\arraybackslash}p{(\columnwidth - 8\tabcolsep) * \real{0.2105}}
  >{\raggedleft\arraybackslash}p{(\columnwidth - 8\tabcolsep) * \real{0.2105}}
  >{\raggedleft\arraybackslash}p{(\columnwidth - 8\tabcolsep) * \real{0.2105}}
  >{\raggedleft\arraybackslash}p{(\columnwidth - 8\tabcolsep) * \real{0.2105}}@{}}
\toprule\noalign{}
\begin{minipage}[b]{\linewidth}\raggedright
host
\end{minipage} & \begin{minipage}[b]{\linewidth}\raggedleft
log2 m\_a upper
\end{minipage} & \begin{minipage}[b]{\linewidth}\raggedleft
log2 m\_b upper
\end{minipage} & \begin{minipage}[b]{\linewidth}\raggedleft
log2 m\_N upper
\end{minipage} & \begin{minipage}[b]{\linewidth}\raggedleft
log2(m\_b sqrt(m\_a m\_N)) upper
\end{minipage} \\
\midrule\noalign{}
\endhead
\bottomrule\noalign{}
\endlastfoot
J14 & 667,968 & 1,227,000,768 & 4,792,426,926 & 3,623,548,215 \\
J56 & 49,691,208 & 5,726,462,112 & 20,810,718,264 & 16,156,666,848 \\
\end{longtable}

All entries are strict upper bounds. The combined entries are less than2\^{}32 and2\^{}34 respectively. These computations evaluate no class group or automorphic period.

\subsection{4. The unchanged strict volumes and finite theta input}\label{the-unchanged-strict-volumes-and-finite-theta-input}

Keep LIFT48's notation n=2b², r0=a+1, u=a(a-1)/2, N=a+2t+1. The letter r0 is distinguished from the root-discriminant exponent r. Let G\_l(q)=\textbar GL\_l(F\_q)\textbar. DO056(AC14, AC23--24) gives at a selected G place

\[
 v_{0,G}={q^u\over G_a(q)},\quad
 v_{1,G}={q^{u+r_0t}G_t(q)\over G_b(q)},\quad
 v_{T,G}={q^{u+2r_0t+t^2}G_t(q)^2\over G_N(q)}.
 \tag{GM13}
\]

At T its three strict volumes are (T\^{}a-1)\^{}-1, (T\^{}b-1)\^{}-1,1, respectively. At Q the initial volume is one, the middle inverse volume is the Gaussian binomial {[}b choose r0{]}\_q, and the terminal inverse volume is {[}N choose r0{]}\_q. Each indicated T or G formula is local; the product has f selected places above each of the two rational G primes and above T.

For q\textgreater=2 and l\textgreater=1,

\[
 {1\over4}q^{l^2}<G_l(q)<q^{l^2},\qquad
 {r_0+s\brack r_0}_q<4q^{r_0s}.
 \tag{GM14}
\]

For the first inequality, the first two factors at q=2 have product 3/8; the remaining product is at least one minus sum\_(i\textgreater=3)2\^{}-i=3/4. This gives9/32\textgreater1/4, and smaller l is easier. For the second, cancel the first s factors in the Gaussian binomial formula and bound its remaining numerator by one; its denominator is product\_(i=1)\^{}r0(1-q\^{}-i)\textgreater1/4.

The actual full-Y finite theta l1 bound in LIFT48 is

\[
 C_{\rm fin}=T^{fn/2}(g_0g_1)^{fn/2}
                      \prod_{v\in Q}q_v^{(n+a-t)/2}.
 \tag{GM15}
\]

There is no new finite input in this equation. Combining (GM13--15) with the exact mass refinement of LIFT48(OB25) gives the bound

\[
 {C_{\rm fin}\over v_1\sqrt{v_0v_T}}
 <2^{8f}T^{fE_T}(g_0g_1)^{fE_G}
                          \prod_{v\in Q}8q_v^{E_Q},
 \tag{GM16}
\]

where

\[
 \begin{split}
 E_T&=n/2+b+a/2,\\
 E_G&=n/2+b^2+(a^2+N^2)/2
                       -2u-2r_0t-5t^2/2,\\
 E_Q&=(n+a-t)/2+2r_0t.
 \end{split}\tag{GM17}
\]

To verify the G factor directly, its exact volume reciprocal is G\_b sqrt(G\_a G\_N)/(q\^{}(2u+2r0t+t²/2) G\_t²). Insert (GM14) and multiply by q\^{}(n/2); it is less than16 q\^{}E\_G. There are2f such places, giving2\^{}(8f). The Q volume reciprocal is less than(4q\^{}(r0t))sqrt(4q\^{}(2r0t))=8q\^{}(2r0t). At T use T\^{}j-1\textless T\^{}j. These calculations use each strict index once; they do not add a projective-level index or a second mass factor.

The exact values of (E\_T,E\_G,E\_Q) are

\begin{longtable}[]{@{}lrrr@{}}
\toprule\noalign{}
host & E\_T & E\_G & E\_Q \\
\midrule\noalign{}
\endhead
\bottomrule\noalign{}
\endlastfoot
J14 & 3,249,027 & 4,947,452 & 6,795,239/2 \\
J56 & 3,249,090 & 5,161,400 & 7,596,851/2 \\
\end{longtable}

\subsection{5. A numerical bound for the weight-independent operator overhead}\label{a-numerical-bound-for-the-weight-independent-operator-overhead}

The actual prime data give

\[
 T<2^{1001},\quad g_0g_1<2^{45}\quad(J14),
 \qquad
 T<2^{996},\quad g_0g_1<2^{71}\quad(J56).
 \tag{GM18}
\]

These are checked on the exact integers of DO040, also used in K065; no rounded logarithm is used. The LIFT48 bound for the principal theta factor is log2 U\_E(D)\textless496fn for every K044 word and allowed gap. Combining (GM9--18) yields

\[
 \log_2\left[
 U_E(D)\,m_b\sqrt{m_am_N}\,
                  {C_{\rm fin}\over v_1\sqrt{v_0v_T}}\right]
 <C_{\rm fixed}+\sum_{v\in Q}(3+E_Q\log_2q_v),
 \tag{GM19}
\]

with the exact sufficient constants

\[
 C_{\rm fixed}=
 \begin{cases}
 284,861,333,421<2^{39},&J14,\\
 1,162,551,608,736<2^{41},&J56.
 \end{cases}\tag{GM20}
\]

They are independent of the selected word, allowed gap and creation degrees. The Q term retains its actual tuple and norm primes. The weight-dependent factors have not been included in this constant. The actual arithmetic bilinear form therefore satisfies

\[
 \boxed{\quad
 |\mathcal B(P,Z)|\le
 2^{C_{\rm fixed}}\!\prod_{v\in Q}(8q_v^{E_Q})\,
 {\sqrt{d_1H_0H_1}\,|b_{\rm Br}|\over|\mathfrak M_\sigma|}
                         \|P\|_0\|Z\|_T.
 \quad}\tag{GM21}
\]

The phase of the exact map remains the one in K038/DO056; (GM21) is an absolute-value inequality for it. Dividing the common rational branch scalar out of both the map and its finite content replaces \textbar b\_Br\textbar{} by product sqrt(kappa), as in LIFT48. Nothing here licenses omitting the CM period M\_sigma, the hook norms, algebraic coordinate denominators, or the signed initial T-lattice cost.

Conditional on an actual nonzero Mram period, LIFT47/DO056 allow integral inputs with norm product at most 2\^{}Kfin Bhat0 BhatT. Multiplying this expression into (GM21) gives a completely explicit sufficient archimedean bound for one actual nonzero arithmetic coefficient, apart from its deliberately retained known period/hook factors and tuple primes. It introduces no integral-inverse premise and does not alter any lattice or level.

\subsection{6. Scope and verification}\label{scope-and-verification-1}

The new counting theorem is (GM7), with the useful rank-two values (GM9). The subsequent calculations use accepted source mass ratios only for the original standard genus, and combine exact original strict volumes with the fixed input of LIFT48. There is no equality between a broad all-Hermitian-lattice count and m\_n.

No actual nonzero period, exact theta-image content, full Host, or completed DP44 inequality is proved. This removes the unevaluated genus-mass constant from the preceding operator upper bound.

The companion \nolinkurl{actual_genus_mass_operator_certificate.py/.json} passed17,731 exact assertions: finite submodules generated by at most the ambient rank, Gaussian integral Gram/discriminant indices, strict finite-volume inequalities, actual mass telescopes, and all displayed integer bounds. It does not enumerate the high-rank actual classes or evaluate an automorphic period.

\clearpage

\section{Branch scalar normalization: scoped erratum to ALP-LIFT48 and ALP-LIFT49}\label{branch-scalar-normalization-scoped-erratum-to-alp-lift48-and-alp-lift49}

Author: \nolinkurl{/root/lift_accounting}, ALP\_MACHINE. Recorded 2026-11-22.

This note corrects only the scalar-removal sentences in LIFT48 §7 and LIFT49 §5. The frozen reports and their hashes are preserved. The operator inequalities (OB24) and (GM21), which retain the complete factor \nolinkurl{|b_Br|}, are unchanged.

\subsection{1. Exact correction}\label{exact-correction}

Use the established notation of AM-R-024(R24.18--21):

\[
 C=\prod_\tau c_{{\rm Br},\tau},\qquad
 t_\tau=\tau(\operatorname{star}\pi_T),\qquad
 u_T=\prod_\tau t_\tau^{-n_{{\rm Br},\tau}},\qquad
 s_B=u_TC,\qquad
 D_{\rm Br}=\sum_\tau n_{{\rm Br},\tau}.
 \tag{BE1}
\]

Here \nolinkurl{C} is the common rational integer and \nolinkurl{u_T} is an algebraic scalar. The parameter \nolinkurl{t_tau} is the literal source branch parameter \nolinkurl{tau(star pi_T)} in AM-R-019; its conjugation fixes the exact phase convention. All conjugates of \nolinkurl{pi_T} and \nolinkurl{star pi_T} have absolute value \nolinkurl{sqrt(T)}. Consequently, for every relevant complex embedding,

\[
 |u_T|=T^{-D_{\rm Br}/2},\qquad
 |b_{\rm Br}|=|u_T|\,|C|\prod_\tau\sqrt{\kappa_\tau}.
 \tag{BE2}
\]

Thus dividing the arithmetic map by \textbf{only \nolinkurl{C}} replaces its archimedean branch factor by

\[
 |u_T|\prod_\tau\sqrt{\kappa_\tau}.
 \tag{BE3}
\]

Dividing by the \textbf{full algebraic scalar \nolinkurl{s_B=u_T C}} replaces that factor by

\[
 \prod_\tau\sqrt{\kappa_\tau}.
 \tag{BE4}
\]

In either normalization, the finite contents must shift by the valuations of the same scalar that was divided out. In particular, the second normalization shifts the actual finite contributions at places above \nolinkurl{T} as well as those of \nolinkurl{C}. Since \nolinkurl{T} is distinct from 11, \nolinkurl{u_T} is an 11-adic unit, but it is generally not an archimedean unit. LIFT48's sentence following (OB19), which refers to the actual phase of \nolinkurl{b_Br}, is to be read as retaining its exact phase \textbf{and modulus}; it must not be read as asserting \nolinkurl{|u_T|=1}.

\subsection{2. Proof and scope}\label{proof-and-scope}

The actual branch identity is

\[
 B_{\rm term}=s_B\bigotimes_\tau i_\tau^\dagger,
 \qquad
 \left\|\bigotimes_\tau i_\tau^\dagger\right\|
       =\prod_\tau\sqrt{\kappa_\tau}.
 \tag{BE5}
\]

This is AM-R-020(R20.3) and AM-R-024(R24.18--20), with the exact branch norms supplied by AM-R-019(R19.15--17). Taking absolute values gives (BE2); dividing by \nolinkurl{C} or by \nolinkurl{s_B} gives (BE3) or (BE4), respectively. The finite shift follows directly from \nolinkurl{v(x/s)=v(x)-v(s)} in the actual coefficient field. AM-R-024(R24.21) then supplies the product-formula cancellation of the complete scalar, with the same normalizations at all finite and infinite places. There is no assertion that the normalized branch has saturated image on the preceding theta image, nor that its remaining norm, lattice, or detector costs vanish.

The exact affected frozen versions are:

\begin{itemize}
\tightlist
\item
  \nolinkurl{actual_normalized_theta_operator_bound.md}, SHA256 \nolinkurl{4b7b3854bdb646fd7dcdb2317ce0c3ecab061b8fb8c77ba9f466fe92f1d6cbe4}.
\item
  \nolinkurl{actual_genus_mass_operator_cost.md}, SHA256 \nolinkurl{67a2182580323db326b56ccdb42e96fef0330f58e6eeb078640ffe15d7a9d072}.
\end{itemize}

No equation in either report's main bound, no counting or Fourier argument, and no existing computational certificate changes.

\subsection{3. Check record}\label{check-record}

Author whole-note check: PASS. The complete note and the exact affected paragraphs were read against AM-R-024(R24.17--21). This is a scoped normalization correction, not a new computation certificate or a rereading of the unchanged whole reports.

\clearpage

\section{AM-K066 --- explicit source period and creation costs in every complex realization}\label{am-k066-explicit-source-period-and-creation-costs-in-every-complex-realization}

Author: GPT-6 Astra Pro, ALP\_MACHINE /root/kummer\_extension. 22 November 2026 UTC. Inherited reasoning effort is not separately exposed. Author whole-proof check PASS --- 2026-11-22T09:49:13.066833+00:00.

This computes and bounds the actual scalar M\_sigma of DO056(AC4) and LIFT48(OB20). It retains its phase in the exact comparison, every nonuniform derivative degree, and the distinction between a pointwise archimedean bound and the normalized average over coefficient embeddings. No finite valuation or algebraic field norm is assigned to a transcendental period or creation factor.

\subsection{1. Fixed data and the actual normalization}\label{fixed-data-and-the-actual-normalization}

Use the two K035 product hosts, indexed by h=1 or4:

\begin{longtable}[]{@{}lllr@{}}
\toprule\noalign{}
h & CM host K\_h & Abelian variety & f \\
\midrule\noalign{}
\endhead
\bottomrule\noalign{}
\endlastfoot
1 & Q(zeta\_28) B & J\^{}14 & 42 \\
4 & Q(zeta\_28) B(sqrt(5),sqrt(13)) & J\^{}56 & 168 \\
\end{longtable}

Here J is the fixed Q-curve Jacobian y\^{}2=x\^{}7-1, and B is the real degree-seven field of conductor29. Write the actual CM coordinates as tau=(i,epsilon,l,rho), with i=1,2,3, epsilon=+/-, l in Z/7Z and rho a real embedding of the degree-h auxiliary field. These coordinates are exactly K035(PF4--PF11).

Let b=1802 and retain the actual ranks a,c of LIFT48. Put \[
 n=ab+bc=2b^2=6494408,\qquad N=14hn,\qquad
 d_\tau=D_{0,\tau}+D_{1,\tau}\ge0,\qquad
 D=\sum_\tau d_\tau,\quad
 D_i=\sum_{\epsilon,l,\rho}d_{i,\epsilon,l,\rho}.
 \tag{PC1}
\] The d\_tau are the actual selected homogeneous degrees. No equality between labels is assumed. N is90921712 or363686848, whereas D can grow with the selected profile.

Keep R004's fixed symplectic basis, normalized period matrix tau\_J, Y=Im(tau\_J), and rational differential coordinates omega\_i=x\^{}(i-1)dx/y. Let C\_theta=theta\_0 c with its exact phase in R006(R6.3--R6.4). The actual comparison scalar is \[
 \mathfrak M_\sigma=
 \det(2Y)^{7hb^2}C_{\theta,\sigma}^{\,N}
 \prod_\tau\left(\frac{\lambda_{\sigma,\tau}}
                         {cD_{\sigma,\tau}}\right)^{d_\tau}.
 \tag{PC2}
\] This is DO056(AC4), not a newly chosen normalization. The label sigma means reconstructing the complex realization after an embedding of the algebraic coefficient field. It does not mean applying a field embedding to an arbitrarily defined transcendental number.

\subsection{2. Combining the vacuum and the fixed Jacobian period}\label{combining-the-vacuum-and-the-fixed-jacobian-period}

Let sigma\_q be the actual q-normalized reduced section of the rigidified principal line L\_(2 infinity), with q=u\_1u\_3-u\_2\^{}2, from LIFT08 and R006. Define the positive real constant \[
 \mathcal P_J=\det(2Y)^{1/4}|C_\theta|
                         =\|\sigma_q\|_{L^2}^{-1}.
 \tag{PC3}
\] Here the metric is the principal theta metric and Haar is probability measure. To prove the second equality, R006(R6.15) identifies the raw translated theta with C\_theta sigma\_q. Its translation is unitary. The raw theta has squared norm det(2Y)\^{}(-1/2), by R004(R4.12) before theta-null rigidification, or LIFT39(PN2--PN5). Divide by \textbar C\_theta\textbar\^{}2 and take the positive inverse square root. No level volume occurs.

The curve, line L\_(2 infinity), rigidification and q-normalized algebraic section are defined over Q by LIFT08. In every coefficient embedding they give the same complex geometric objects in the rational u\_i coordinates, with the same canonical rigidified principal metric. Thus the modulus in PC3 is unchanged in all these realizations. Possible finite Mumford phases remain in C\_(theta,sigma). This does not assert that C\_theta is algebraic.

One can also verify this modulus conclusion in R004's fixed homology basis using R005/R006. A conjugate nonsingular characteristic is D\_S for another three-element branch subset. Its within-part Vandermonde in the first Thomae formula times the cross-part Vandermonde in the derivative formula is the full finite Vandermonde, up to sign. The common period factor is unchanged. Consequently theta\href{0}{S}c\_S has the same modulus for all these choices. This is a check on PC3, not an additional choice of a fourth-root phase.

The exact R006 identity is \[
 |c|^4=\frac{|R_S|}{256},\qquad
 R_S=\prod_{a\in\{1,3,5\}}\prod_{b\in\{0,2,4,6\}}
                              (\zeta_7^a-\zeta_7^b).
 \tag{PC4}
\] R004 gives 5/8\textless\textbar theta\_0\textbar\textless11/8. The companion rational certificate bounds the exact R004 matrix Y and all twelve factors in PC4. It proves \[
 \frac25<\mathcal P_J<\frac98,\qquad
       -\frac{18}{16}<\log_2\mathcal P_J<\frac2{16}.
 \tag{PC5}
\] The second pair is checked directly on a fourth-power interval, not by rounding decimal logarithms. All coefficient realizations obey these same strict bounds.

Since N/4=7hb\^{}2, the modulus of PC2 is exactly \[
 \boxed{\quad
 |\mathfrak M_\sigma|=\mathcal P_J^{\,N}
       \prod_\tau s_{\sigma,\tau}^{\,d_\tau},\qquad
 s_{\sigma,\tau}=|\lambda_{\sigma,\tau}/cD_{\sigma,\tau}|.
 \quad}\tag{PC6}
\] The phase in PC2 has not been discarded from the source identity; PC6 is its separate modulus equation.

\subsection{3. Every coefficient embedding and the creation factors}\label{every-coefficient-embedding-and-the-creation-factors}

Write pi\_l for the seven positive conjugates of pi\_B=Norm\_(Q(zeta29)/B)(1-zeta29). Put Delta=1 for h=1 and65 for h=4. Let delta\_rho be1 or the four positive conjugates of delta\_1 delta\_3 in K032. Set \[
 d_i=\prod_{\substack{j=1,2,3\\j\ne i}}
                         (\zeta_7^i-\zeta_7^j).
 \tag{PC7}
\] This algebraic d\_i is distinct from the degree d\_tau in PC1. K035(PF10, PF15, PF23) gives in the original realization \[
 h_\tau=\frac{\delta_\rho\pi_l^{-6}}{2G_i},\quad
 |cD_\tau|^2=4|d_i|^2\pi_l^{12}\Delta,\quad
 |\lambda_\tau|^2=(\pi h_\tau)^{-1}.
 \tag{PC8}
\] It follows that every reconstructed coefficient realization has \[
 \boxed{\quad
 s_{\sigma,\tau}^{\,2}
 =\frac{G_i}{2\pi\,|d_i^{(a_\sigma)}|^2
          \pi_{l_\sigma}^{\,6}\delta_{\rho_\sigma}\Delta},
 \qquad
 d_i^{(a)}=\prod_{j\ne i}(\zeta_7^{ai}-\zeta_7^{aj}).
 \quad}\tag{PC9}
\] Here a\_sigma ranges through1,\ldots,6, and l\_sigma,rho\_sigma are the conjugated B and quadratic labels. Changing epsilon changes only a unit factor and does not occur in PC9.

The realization accounting is important. The omega\_i are rational. K035(PF11) builds the algebraic coordinates from these forms using the matrices of Q(i), B and the real quadratic fields. That change of coordinates contains no zeta\_7 coefficient. A coefficient embedding therefore fixes the physical differential index i, permutes the B and quadratic factors and may exchange the Q(i) factors. The actual metric is the first expression of PC8 with the conjugated positive factors. The coordinate generator contains d\_i and hence changes it to d\_i\^{}(a\_sigma).

The CM-action label is transported along with the endomorphism action. It is not held fixed while the algebraic coordinates are conjugated. Thus G\_i remains attached to rational omega\_i; neither sigma(G\_i), sigma(lambda\_tau) nor sigma(Omega\_i) is used. The creation scalar in each realization is reconstructed from the actual Chern/Fock normalization PF23. This proves PC9 for every embedding of any coefficient field containing K\_h, even if the larger field is not Galois. The original degree d\_tau likewise remains the degree at its transported label.

\subsection{4. Exact certified pointwise costs}\label{exact-certified-pointwise-costs}

The companion certificate records all630 intervals \[
 \frac{L_{h,i,a,l,\rho}}{16}
 <\log_2 s_{h,i,a,l,\rho}
 <\frac{U_{h,i,a,l,\rho}}{16}.
 \tag{PC10}
\] There are126 entries for J14 and504 for J56; epsilon gives the same modulus twice. Their uniform extrema are:

\begin{longtable}[]{@{}lrr@{}}
\toprule\noalign{}
Host & Lower bound for log2 s & Upper bound for log2 s \\
\midrule\noalign{}
\endhead
\bottomrule\noalign{}
\endlastfoot
J14 & -209/16 & 131/16 \\
J56 & -252/16 & 126/16 \\
\end{longtable}

For any actual sigma, use its transported indices in the table. Then \[
 \frac{-18N+\sum_\tau L_{\sigma,\tau}d_\tau}{16}
 <\log_2|\mathfrak M_\sigma|<
 \frac{2N+\sum_\tau U_{\sigma,\tau}d_\tau}{16}.
 \tag{PC11}
\] This retains the entire nonuniform degree profile. If only the total D is needed, the following simpler bounds hold: \[
\begin{array}{ll}
 J14:&-\frac{18N+209D}{16}<\log_2|\mathfrak M_\sigma|
                                   <\frac{2N+131D}{16},\\[2mm]
 J56:&-\frac{18N+252D}{16}<\log_2|\mathfrak M_\sigma|
                                   <\frac{2N+126D}{16}.
\end{array}
 \tag{PC12}
\] In particular the reciprocal factor in LIFT48(OB24) is bounded by 2\^{}((18N+209D)/16) or2\^{}((18N+252D)/16), respectively. These are pointwise bounds and require no choice of averaging convention.

\subsubsection{Exact rational computation}\label{exact-rational-computation}

Machin's identity and alternating rational series bracket pi. After reflection into(0,pi/2), the alternating sine series gives outward bounds for sin(k pi/7) and sin(k pi/29). The exact R004 matrix determinant is evaluated by interval arithmetic. The positive B generators satisfy \[
 \pi_l=16\sin^2(\pi a_l/29)\sin^2(12\pi a_l/29),
                          \qquad a_l=2^l\bmod29.
 \tag{PC13}
\] Rational square-root enclosures give the two positive factors delta\_(a,+)=(sqrt(a\^{}2+4)-a)/(2sqrt(a\^{}2+4)) and their conjugates, for a=1,3. No numerical eigenvalue is used.

The beta constants admit the positive series \[
 B(x,1/2)=\sum_{n\ge0}\frac{c_n}{n+x},\quad
 c_n=\binom{2n}{n}/4^n,\qquad x=i/7.
 \tag{PC14}
\] Expand (1-t)\^{}(-1/2) and use monotone convergence to integrate term by term. The elementary induction inequality \[
 (2n+1)^2(3n+4)\le4(n+1)^2(3n+1)
 \tag{PC15}
\] gives c\_n\textless=1/sqrt(3n+1). The tail after n=256 is strictly less than \[
 \frac1{\sqrt3}\sum_{n=257}^\infty n^{-3/2}
 \le\frac2{\sqrt3\sqrt{256}}<\frac1{12}.
 \tag{PC16}
\] K035(PF7--PF12) gives the positive formula \[
 \frac{G_i}{2\pi}
  =\frac{2B(i/7,1/2)^2\sin^2(\pi i/7)}
              {7\pi\sin(2\pi i/7)}.
 \tag{PC17}
\] Substitute these rational intervals in PC9. Raising the bounds for s\^{}2 to the eighth power and comparing with exact powers of2 proves PC10. The same procedure applied to PC4 and the theta-null bounds proves PC5. No decimal logarithm is a proof premise.

\subsection{5. Exact normalized embedding averages}\label{exact-normalized-embedding-averages}

Let M be any finite coefficient field containing K\_h, and average over all embeddings sigma:M-\textgreater C with multiplicity one. Because K\_h is Galois, every one of its automorphisms occurs equally often as a restriction. The E, B and quadratic factors have the independent integral compositum structure proved in K030/K032.

The exact algebraic norms are \[
 |N_{\mathbf Q(\zeta_7)/\mathbf Q}d_i|=49,\quad
 N_{B/\mathbf Q}\pi_B=29,\quad
 N_{\mathbf Q(\sqrt5,\sqrt{13})/\mathbf Q}
                              (\delta_1\delta_3)=65^{-2}.
 \tag{PC18}
\] These are respectively K035's partial-embedding calculation, K030(PH11), and K032(QH6) with the degree-two norm in the other quadratic factor. Averaging PC9 yields, for physical index i, \[
 A_{h,i}:=\operatorname{mean}_\sigma\log s_{\sigma,\tau}
 =\frac12\log\frac{G_i}{2\pi}
   -\frac13\log7-\frac37\log29
   -\mathbf1_{h=4}\frac14\log65.
 \tag{PC19}
\] Thus the exact nonuniform-degree mean is \[
 \boxed{\quad
 \operatorname{mean}_\sigma\log|\mathfrak M_\sigma|
          =N\log\mathcal P_J+\sum_{i=1}^3D_i A_{h,i}.
 \quad}\tag{PC20}
\] This is an average of reconstructed real moduli. The algebraic norms in PC18 are legitimate, but PC20 is not a field norm of M\_sigma. Further coefficient-field extension leaves this normalized mean unchanged.

The exact certificate gives these binary-log intervals; all displayed numerators have denominator84:

\begin{longtable}[]{@{}lrrr@{}}
\toprule\noalign{}
Host & i & Lower numerator for A/(log 2) & Upper numerator \\
\midrule\noalign{}
\endhead
\bottomrule\noalign{}
\endlastfoot
J14 & 1 & -230 & -228 \\
J14 & 2 & -241 & -238 \\
J14 & 3 & -202 & -198 \\
J56 & 1 & -357 & -354 \\
J56 & 2 & -368 & -364 \\
J56 & 3 & -329 & -325 \\
\end{longtable}

For example -230/84\textless A\_(1,1)/log2\textless-228/84. To verify a row without computing a logarithm, use \[
 \exp(84A_{h,i})
 =\frac{(G_i/(2\pi))^{42}}
             {7^{28}29^{36}65^{21\mathbf1_{h=4}}}.
 \tag{PC21}
\] Compare its rational interval with powers of2. Together with PC5 this bounds the normalized mean reciprocal cost, separately from the pointwise PC12. No uniformity of d\_tau over B or the quadratic embeddings is assumed.

\subsection{6. Source interfaces and limits}\label{source-interfaces-and-limits}

The source creation normalization is \nolinkurl{Formalism/LocalToGlobal/IndefiniteReal/LocalNewformFunction.lean124–141;} pureTensorO uses it in DatumVec.lean130--134. K035(PF19--PF29) proves the exact sign and scale comparison with ordinary algebraic reduced-theta derivatives. Per-label RealD.sc is explicitly allowed in Development/AbelianVarietyTheory/Definitions/LatticeFrames/ \nolinkurl{OnFramesRdApply.lean25–39,108–121.} K037(CR35--CR36), LIFT39(PN3--PN5), DO056(AC4) and LIFT48(OB20) supply the unchanged raw/unit-vacuum and coupled coefficient readings.

The Jacobian period and Thomae results are carried theorems from R004--R006, not re-proved here. Their primary anchors are Tadokoro's Proposition3.2, Proposition4.1 and Theorem4.4, and Bernatska's Remark3/equation30 and genus-three Example4. The new result is their exact source-factor combination, all-realization coordinate accounting and certified cost.

PC12 and PC20 remove the unspecified magnitude of M\_sigma from LIFT48's displayed upper bound. They do not prove the actual all-T Mram period nonzero, improve an image ideal, choose the norm convention of an uninstantiated DP44 wrapper, or establish a favorable final gain/cost inequality. Cancellation with finite bad-prime coordinates must use those same factors on the arithmetic side. The phases and the transcendental nature of the periods remain explicit.

\subsubsection{Primary references carried at these interfaces}\label{primary-references-carried-at-these-interfaces}

Bernatska, J. (2020). General derivative Thomae formula for singular half-periods. Letters in Mathematical Physics,110,2983--3014. \url{https://doi.org/10.1007/s11005-020-01315-8} . Source-audited version: \url{https://arxiv.org/pdf/1904.09333v3} , equation2, Remark3/equation30, Example4/equations38--39.

Tadokoro, Y. (2014). The period matrix of the hyperelliptic curve w\^{}2=z\^{}(2g+1)-1. Tsukuba Journal of Mathematics,38(2),137--158. \url{https://arxiv.org/abs/1211.6910} . Carried version v5, Propositions3.2 and4.1 and Theorem4.4.

The companion \nolinkurl{actual_cm_period_cost_certificate.py} verifies exact rational/trigonometric/beta bounds, all label tables and modulus exponent identities. It evaluates no automorphic period, finite valuation of a transcendental quantity, or net DP44 sign.

\subsection{Author verification record}\label{author-verification-record-9}

At 2026-11-22T09:49:13.066833+00:00, I read all saved PC1--PC21 and every section of this proof, together with the full companion script. I reconstructed the unit-vacuum/seed L2 identity, fixed rational differential versus transported CM-label distinction, pointwise modulus formula, nonuniform degree accounting, and the exact embedding average. I checked the source creation, pureTensorO and RealD.sc/Core interfaces directly. The rational interval proof, beta tail induction and all630 modulus labels were checked; no transcendental field norm or finite valuation was used. The original R004--R006 period/Thomae theorems remain carried at their stated interfaces.

The certificate passed5,426 exact assertions at2026-11-22T09:48:14.066198+00:00. Its scope is rational interval arithmetic and exact exponent identities, not an evaluation of the actual automorphic sum or a DP44 sign. No Lean build was performed.

\clearpage

\section{AM-DO-060 --- the integral positive block and the complete two-place short sum}\label{am-do-060-the-integral-positive-block-and-the-complete-two-place-short-sum}

Author: GPT-6 Astra Pro, ALP\_MACHINE \nolinkurl{/root/definite_oscillator}. 22 November 2026 UTC. Inherited reasoning effort is not separately exposed.

\textbf{Author whole-proof check PASS --- 2026-11-22T09:47:59.740112+00:00.}

For the unchanged J56 family at the K044 upper endpoint, this bounds the sum of the absolute values of \textbf{all} short correspondences nonintegral at exactly two relative T places. It includes every pole depth, every source/target pencil pair, and matrices outside the fixed unitary plane. The bound is \(2^{-101550}\). Together with DO059, the unresolved normalized sum can be confined to matrices nonintegral at three or more places. This is not a sign theorem for that remaining sum or an all-T Mram nonvanishing assertion.

The key distinction is between integrality of a large positive block and integrality of the whole matrix. The former holds for every actual supported correspondence, including matrices with poles at T.

\subsection{1. Actual fixed coordinates and the integral block}\label{actual-fixed-coordinates-and-the-integral-block}

Use K062(AT14--24), DO059(SC1--37), and their unchanged notation:

\[
 a=f=168,\quad n=a+1=169,\quad H=8192,\quad Q=g_0g_1,
 \quad m=T-1.
 \tag{IB1}
\]

The sum is over central classes of actual rational unitary matrices g, integral away T, satisfying the full local pencil incidence condition, and with every real operator norm at most H. Write s for the number of relative T places where g is nonintegral. The normalized summand is

\[
                 a(g)c(g)\prod_{v\mid T}\rho(g_v).
 \tag{IB2}
\]

All finite partial traces, selected Hilbert metrics and pencil signs in IB2 remain exactly those of K062. In particular \(|c(g)|\le1\), and each nonintegral place contributes \(|\rho|\le2/(T-2)\).

The actual negative vector is \(v_0=(Qu,\pi_k)\), where u has only its first four positive coordinates nonzero. K063(PF12) proves equality of the \textbf{labelled actual lattices}, including their two omitted lines, after replacing the local row by

\[
 z=v_0/\bar\pi,\qquad m_*=e_1^t/z_1,\qquad
 q_*=-v_0^*J-\pi m_*,\qquad q_*z=0.
 \tag{IB3}
\]

Both z and q\_* are supported on the first four positive coordinates and the negative coordinate. At a selected T place the pencil lattice is

\[
 L_c=h_c\Lambda,\qquad h_c=1+c\pi^{-1}zq_*,
 \qquad c\in\mathbb F_T\setminus\{-1\}.
 \tag{IB4}
\]

If g maps L\_c to L\_d, then

\[
                  g=h_d k h_c^{-1},\qquad k\in GL_n(O_v).
 \tag{IB5}
\]

Put \(S=\langle e_5,\ldots,e_a\rangle_E\), of dimension a−4, and

\[
           W_0=\langle e_1,e_2,e_3,e_4,e_n\rangle_E,
           \qquad E^n=S\mathbin{\perp}W_0.
 \tag{IB6}
\]

Because the entries of z and q\_* on S vanish, the S-by-S block of IB5 equals the same block of k. It is therefore integral. At the conjugate place use the inverse-adjoint relation: g inverse maps the reversed pair, so its S-block is integral; its transpose conjugate is the S-block of g at the conjugate prime. Away T the full matrix is integral by K062(AT18). Consequently

\[
               A:=g_{S,S}\in M_{a-4}(O_E).
 \tag{IB7}
\]

This does not assert that g itself is integral, or that its off-diagonal blocks are integral.

\subsection{2. The nonscalar integral-block terms}\label{the-nonscalar-integral-block-terms}

Suppose A is nonscalar. An off-diagonal entry or difference of diagonal entries supplies a nonzero \(x\in O_E\) with \(|\tau x|\le2H\) at every embedding. Scalar reduction of the full g at a selected G place forces x to vanish at that place and its conjugate. The exact LIFT42(FC5--9) argument now has no T denominator:

\[
             |a(g)|\le {2(2H)^{af}\over Q^{a(f-1)}}.
 \tag{IB8}
\]

The two mixed G factors still have only the bound one; the exponent f−1 over each rational G prime retains their removal. This is an application of the checked cuspidal Green-character ratio, not an assumption that the mixed projector is a character.

DO059(SC28--30) supplies the complete count

\[
 N_s\le {f\choose s}m^{2s}G_n(2H+1)^{2fn}T^{2ns},\qquad
 G_n=\gcd(|GL_n(\mathbb F_{g_0})|,|GL_n(\mathbb F_{g_1})|).
 \tag{IB9}
\]

The factor \(T^{2ns}\) counts cosets of the actual target over-lattice, not arbitrary entry denominators. In particular no pole-pattern factor has to be reintroduced. At s=2 the sum of the absolute values of all nonscalar-block terms is strictly less than

\[
 \begin{split}
 B_{\rm ns}={}&{f\choose2}m^4G_n(2H+1)^{2fn}T^{4n}
       \left({2\over T-2}\right)^2
       {2(2H)^{af}\over Q^{a(f-1)}}
       <2^{-101552}.
 \end{split}
 \tag{IB10}
\]

The displayed strict binary exponent is checked by integer cross-multiplication in the companion certificate.

\subsection{3. A scalar block leaves at most six nontrivial dimensions}\label{a-scalar-block-leaves-at-most-six-nontrivial-dimensions}

Suppose now \(A=\lambda I_{a-4}\). In the decomposition IB6, write the other block in the S columns as B. Unitarity gives

\[
       (1-\lambda\bar\lambda)I_{a-4}=B^*J_{W_0}B.
 \tag{IB11}
\]

The right side has rank at most five, while a−4\textgreater5. Thus \(\lambda\bar\lambda=1\). Since \(\lambda\in O_E\), the elementary CM unit argument used in the carried host proof gives \(\lambda\in\mu_{28}\). Equation IB11 also says that the image of B is totally isotropic in a space of signature (4,1). It has dimension at most one. Hence g acts by \(\lambda\) on the positive subspace \(\ker B\subset S\), whose dimension is at least a−5.

Multiplication by the central unit \(\lambda^{-1}\) does not change any lattice incidence or absolute coefficient bound. For the remainder of this section choose the unique representative

\[
              h=\lambda^{-1}g,\qquad N=h-I.
 \tag{IB12}
\]

It satisfies

\[
     N_{S,S}=0,\qquad \operatorname{rank}_E N\le6,
     \qquad\|h_\tau\|\le H.
 \tag{IB13}
\]

The rank statement follows either by using the fixed positive subspace above and its invariant orthogonal complement, or directly by counting the codimension of \(\ker B\). It is not a claim that this complement is the fixed two-dimensional plane of K063.

Only terms with a(g) nonzero need consideration. At each unmixed G place, the finite reduction of h has the eigenvalue one with multiplicity at least n−6. Green's primary-semisimple support formula LIFT42(FC6) then forces its entire semisimple part to be one. Indeed the single irreducible polynomial of a nonzero cuspidal character value already has the root one, and is therefore X−1. Consequently

\[
          \operatorname{red}_v h\text{ is unipotent at every
                    unmixed G place.}
 \tag{IB14}
\]

The determinant of h is a global norm-one unit by the actual local pencil volumes, hence lies in \(\mu_{28}\). Its reduction at any one of these G places is one by IB14. Reduction of \(\mu_{28}\) is injective in residue characteristic prime to28, so

\[
                            \det h=1.
 \tag{IB15}
\]

No assertion about the two mixed projectors is needed here.

\subsection{4. Exterior poles, two coefficients, and actual unipotence}\label{exterior-poles-two-coefficients-and-actual-unipotence}

For a rank-one operator zq\_* with q\_*z=0, its action on every exterior power has the form

\[
        \bigwedge^j(1+c\pi^{-1}zq_*)=1+c\pi^{-1}D_j.
 \tag{IB16}
\]

This is exact: every term inserting z in two exterior factors is zero. The matrix D\_j is integral at the selected place. Applying IB16 on both sides of IB5 proves that every matrix entry of every exterior power of h has pole order at most two there. The reversed-pair argument gives the same conclusion at the conjugate prime.

Write \(e_j(N)\) for the j-th elementary characteristic coefficient. It is an integer linear combination of \(e_k(h)\), k\textless=j, and these are traces of exterior powers. Hence each \(e_j(N)\) has pole order at most two at each of the 2s CM primes above the nonintegral places, and is integral elsewhere. The bound is two, independent of j. From IB13 and the eigenvalue bound \(|\theta(N_\tau)|\le H+1\),

\[
                 |\tau e_j(N)|\le {6\choose j}(H+1)^j.
 \tag{IB17}
\]

If e\_j(N) is nonzero, IB14 forces its numerator to vanish at both CM primes of every unmixed G place. The fractional-ideal norm inequality therefore gives

\[
 Q^{2(f-1)}\le
       \left[{6\choose j}(H+1)^j\right]^{2f}T^{4s}.
 \tag{IB18}
\]

The actual J56 integers contradict IB18 for j=1,2 and s\textless=3. At s=2 the strict inverse binary margins for j=1,2,3 are respectively 10,271, 5,459 and951 bits. At s=3 the margins for j=1,2 are 6,289 and1,477 bits. Thus

\[
                         e_1(N)=e_2(N)=0
                \quad\text{for scalar-block terms with }s\le3.
 \tag{IB19}
\]

Here is an exact signature-one argument that uses only these two coefficients and does not impose a real phase on a hyperbolic pair.

\textbf{Lemma.} Suppose h is unitary of signature (a,1), fixes pointwise a positive nondegenerate subspace of codimension at most six, has det h=1, and satisfies e1(h−I)=e2(h−I)=0. Then h is unipotent.

\textbf{Proof.} Restrict to the invariant orthogonal complement of the stated fixed positive subspace, and pad its characteristic polynomial with factors X−1 to degree six. This is exactly IB13's proved model. Let P be that degree-six polynomial. The two coefficient assumptions give its first coefficients 6 and15. Unitarity and det1 give reciprocal-conjugate coefficients. Thus

\[
                      P(X)=(X-1)^6+cX^3,\qquad c\in\mathbb R.
 \tag{IB20}
\]

If c is nonzero, putting w=X+X\^{}\{-1\}−2 makes its root equation \(w^3=-c\). Two of these w roots are nonreal. Each produces two distinct roots of X²−(w+2)X+1, all off the unit circle. The two w values produce four distinct off-unit-circle eigenvalues of h. This contradicts signature (a,1): the generalized eigenspace with eigenvalue modulus greater than one is totally isotropic, because h inverse iterates tend to zero there and preserve the Hermitian form. Its dimension is at most one; the same holds for modulus less than one. Thus at most two off-unit-circle eigenvalues occur, counted with multiplicity. Hence c=0 and all residual eigenvalues are one. Together with the removed fixed part this proves unipotence. QED.

For full precision the lemma is used with the explicitly proved fixed positive subspace of IB13; no assertion that an arbitrary rank bound alone chooses a nondegenerate fixed complement is required.

It follows that every nonzero scalar-block summand with s\textless=3 is represented by an actual global unipotent h. We use this below only to complete the s=2 estimate.

\subsection{5. The two unipotent ranks and their sharper finite ratios}\label{the-two-unipotent-ranks-and-their-sharper-finite-ratios}

We give the elementary real normalization used here. A nontrivial unipotent in U(a,1) fixes an isotropic line. To see this, its fixed space cannot contain a negative vector unless the unipotent is the identity on the remaining positive space. If its fixed space were nondegenerate positive, its invariant orthogonal complement would again have a nonzero fixed vector, a contradiction. Its radical therefore contains an isotropic line. The induced unipotent on the positive quotient of the orthogonal of that line is the identity.

Choose a Euclidean-orthonormal Witt basis. In that basis the matrix is

\[
 h=\begin{pmatrix}
    1&-v^*&-\|v\|^2/2+it\cr
    0&I&v\cr
    0&0&1
   \end{pmatrix},\qquad t\in\mathbb R.
 \tag{IB21}
\]

Such an orthonormal basis exists by taking a Euclidean-unit isotropic vector l, its orthogonal companion Jl, and the positive perpendicular space. Therefore the operator norms in IB21 are the original Euclidean ones. Direct multiplication proves

\[
  N^3=0,\quad \operatorname{rank}N\le2,\quad
  \operatorname{rank}N=2\Longleftrightarrow N^2\ne0,
  \quad \|N^2\|\le2\|N\|\le2(H+1).
 \tag{IB22}
\]

The last inequality uses the sole nonzero entry of N², namely −\textbar\textbar v\textbar\textbar², and the upper-right entry of N. It is linear in H+1; replacing it by the weaker square would lose the two-place bound.

\subsubsection{Rank two}\label{rank-two}

In addition to N, N² has pole order at most two at every nonintegral CM prime. Indeed each two-by-two minor of N is an entry of \(\bigwedge^2h\) minus terms linear in h and constants. It has pole order at most two by IB16. Since tr N=0,

\[
  -(N^2)_{ij}
    =\sum_k\bigl(N_{ij}N_{kk}-N_{ik}N_{kj}\bigr),
 \tag{IB23}
\]

and every summand is such a minor, with the natural row/column sign.

Choose a nonzero entry x of N and a nonzero entry y of N². Their archimedean bounds are H+1 and2(H+1), respectively. At a G reduction where h is scalar, N is zero, so x vanishes. At a reduction where N² is zero, y vanishes. The norm argument of DO058 therefore bounds the log-weight of these two exceptional sets by

\[
 f\log(H+1)+2s\log T,\qquad
 f\log(2(H+1))+2s\log T.
 \tag{IB24}
\]

For a unipotent GL\_n element, the number of Jordan blocks is n−rank N. The exact scalar-primary Green formula LIFT42(FC6) therefore gives the following normalized ratios at scalar-primary values

\[
  {1\over q^a-1}\quad(\operatorname{rank}N=1),\qquad
  {1\over(q^a-1)(q^{a-1}-1)}\quad(\operatorname{rank}N=2).
 \tag{IB25}
\]

At every reduction where N² is nonzero the rank is two. Reductions where N² vanishes but N has rank two can be counted using the weaker rank-one ratio. Combining IB24--25 and dropping the two mixed factors gives

\[
 |a(g)|\le
 {2(H+1)^{af}[2(H+1)]^{(a-1)f}T^{2s(2a-1)}
       \over Q^{(2a-1)(f-1)}}.
 \tag{IB26}
\]

The factor2 bounds all the products of \((1-q^{-a})^{-1}\) and \((1-q^{-(a-1)})^{-1}\), using the same elementary product inequality as LIFT42. The exact host inequalities are checked in the certificate. With IB9, at s=2 this proves

\[
       \sum_{\substack{s(g)=2,\ A\ {\mathrm{scalar}}\cr
                       \operatorname{rank}(h-I)=2}}
          |a(g)c(g)\prod_v\rho(g_v)|<2^{-362973}.
 \tag{IB27}
\]

\subsubsection{Rank one}\label{rank-one}

A rank-one nonzero unipotent satisfies \(N=\kappa l l^*J\), for an isotropic vector l and a nonzero pure-imaginary scalar \(\kappa\). This follows immediately from IB21; it also follows from rank one and \(N^*J+JN=0\). One may first apply this at any archimedean embedding; vanishing of the resulting matrix entries gives their algebraic vanishing in E. Because \(N_{S,S}=0\), each diagonal entry \(\kappa|l_i|^2\), i in S, vanishes. Thus \(l_S=0\): N is supported entirely on the fixed five-coordinate space W0. Consequently h is the identity on S and is an actual U(4,1) matrix on W0.

The pencil matrices IB4 are already supported on W0. The entire column-coset and compact-fibre proof of DO059 therefore applies in matrix rank five, while retaining the original field degree2f and all actual pencil labels. It counts at most

\[
 N_{s,1}\le {f\choose s}m^{2s}G_5(2H+1)^{10f}T^{10s},
 \quad G_5=\gcd(|GL_5(\mathbb F_{g_0})|,
                 |GL_5(\mathbb F_{g_1})|).
 \tag{IB28}
\]

The fixed-column fibre is a subset of the finite group of source lattice isometries fixing its negative vector. Reduction at each G prime is injective by the checked DO058 torsion argument, now in rank five. The exact common group-order divisor is

\[
                   G_5=98148490017978252769652736000000.
 \tag{IB29}
\]

There is no factor28 in IB28: each central class in the scalar-block stratum has the unique representative IB12 with S-block the identity. Counting its actual five-coordinate matrix already counts that class.

The nonzero entry of N, with arch bound H+1 and pole order two, gives

\[
         |a(g)|\le {2(H+1)^{af}T^{2as}\over Q^{a(f-1)}}.
 \tag{IB30}
\]

Using IB28--30 and the actual pencil trace bound yields, at s=2,

\[
       \sum_{\substack{s(g)=2,\ A\ {\mathrm{scalar}}\cr
                       \operatorname{rank}(h-I)=1}}
          |a(g)c(g)\prod_v\rho(g_v)|<2^{-890221}.
 \tag{IB31}
\]

Rank zero gives a central integral matrix and cannot occur in the nonintegral sum.

\subsection{6. The full two-place result and the unchanged remaining gate}\label{the-full-two-place-result-and-the-unchanged-remaining-gate}

The nonscalar-block stratum and the two scalar-block ranks are a disjoint exhaustion of the s=2 terms with nonzero a(g). The three absolute estimates give

\[
 \boxed{\quad
       \sum_{s(g)=2}|a(g)c(g)\prod_v\rho(g_v)|
        <2^{-101552}+2^{-362973}+2^{-890221}<2^{-101550}.
       \quad}
 \tag{IB32}
\]

Insert this and DO059's complete one-place bound in the exact K062 mass-normalized identity. With \(S_{\ge3}^{\rm new}\) the unchanged short sum over matrices nonintegral at at least three relative T places,

\[
 {E_{x,\rm ram}\over
 w_x\,28C_\infty d_Rd_W((T-2)/(T-1))^f}
       =1+S_{\ge3}^{\rm new}+\varepsilon_3,
 \quad
 |\varepsilon_3|<2^{-1999}+2^{-104546}+2^{-101550}<2^{-1998}.
 \tag{IB33}
\]

Every original class weight, finite projector and Hilbert scalar remains. No summand was declared positive. The estimates bound absolute values on the removed strata and preserve the signed rest.

The independently checked K064 complete E12-plane estimate can also be removed. Its nonintegral matrices have support repeated at least28 times on J56, so it is disjoint from s=1 and s=2. Its absolute sum is less than \(2^{-21414}\), even at the older larger cutoff. Removing this plane still leaves a total error less than \(2^{-1998}\). The remainder outside that plane, with s\textgreater=3, is not computed here.

IB19--20 additionally constrain the scalar-block part at s=3 to global unipotents. They do not bound that whole stratum: the displayed rank-two counting estimate at s=3 exceeds one, as does the generic nonscalar-block estimate. On J14 the relevant arithmetic inequalities already fail at s=2. No assertion of Mram nonvanishing on either host, or of a final arithmetic budget, follows from this report.

\subsection{7. Carried sources and verification}\label{carried-sources-and-verification-1}

The fixed support and phases are those of K062(AT3--24), K063(PF12), DO059(SC28--37), and LIFT44(PR4--7). These are carried checked findings; their source conventions have not been redefined. The only finite character input is LIFT42(FC5--6): Green (1955), section5, Example(ii), pp.430--431, with the explicit scalar-primary formula retained in that finding's primary-source audit. The determinant root group is the checked K031/K032 host group μ28. The new block, exterior, residual unitary and quantitative arguments are proved in full above.

The companion \nolinkurl{definite_oscillator_integral_block_certificate.py} checks the actual integer margins, the rank-five group-order gcd, local full-pencil and exterior-power identities, the trace-corrected square identity, signature-one unipotent matrices and their exact norm inequality, degree-six reciprocity, and both Green rank ratios. It is a finite diagnostic and arithmetic certificate. It does not enumerate the rational remainder, evaluate its archimedean coefficients, or perform a Lean build.

\subsection{Author verification record}\label{author-verification-record-10}

At 2026-11-22T09:47:59.740112+00:00, I completed a whole read of IB1--33 and every section, and of the entire final certificate. I reconstructed the literal five-coordinate pencil support, conjugate-place integrality, scalar-block fixed complement, exterior pole bound, two-coefficient reciprocal-polynomial argument, signature-one unipotent normal form, linear square norm, finite Green rank losses, compact rank-five fibre and all exact mass-normalized inequalities. I then checked the final inline-math and scalar-label corrections. The complete s=2 assertion passes in the stated J56 upper-endpoint scope.

The certificate passed 10,550 exact assertions at 2026-11-22T09:46:25.413868+00:00, including 45 local pencil cases, 5,670 exterior entries and 54 signature-one unipotent cases. An initial diagnostic used structural SymPy equality for an unexpanded complex square; it was corrected to exact algebraic simplification before this PASS. No mathematical claim was promoted from that failed diagnostic. The s\textgreater=3 remainder and J14 quantitative conclusion remain open.

\clearpage

\section{AM-R-041 --- the same normalized kernel outside the two-place and larger-plane strata}\label{am-r-041-the-same-normalized-kernel-outside-the-two-place-and-larger-plane-strata}

Author: GPT-6 Astra Pro, ALP\_MACHINE /root, 22 November 2026 UTC. Reasoning effort is not externally exposed.

This corollary combines the whole-checked DO060 theorem with R040. It changes no chosen field, prime, word, role, twist, pencil, metric or class weight. Its scope is the unchanged J56 family at the K044 upper endpoint, where DO059 supplies the complete cutoff H=8192 decomposition. It does not determine the sign of the remaining finite sum.

\subsection{1. The disjoint partition}\label{the-disjoint-partition}

Write s(g) for the number of nonintegral relative T places. Let P48 be the full supported nonintegral fixed-plane family with coefficients in L=Q(zeta28,sqrt5,sqrt13) and every archimedean operator norm at most8192. These are exactly the matrices counted in R040(R40.1), with their original central quotient. Let

\[
 w(g)=a(g)c(g)\prod_{v\mid T}\rho(g_v).
 \tag{R41.1}
\]

The three removed sets are s(g)=1, s(g)=2, and P48. The first two are disjoint by definition. K063(PF12--13), K065(TL1--7) and R040 show that every nonintegral member of P48 has its support repeated seven times. Thus P48 is disjoint from both s=1 and s=2. Their full absolute contributions, with no restriction on allowed pole depths, are

\[
 \sum_{s(g)=1}|w(g)|<2^{-104546},\qquad
 \sum_{s(g)=2}|w(g)|<2^{-101550},\qquad
 \sum_{g\in P48}|w(g)|<2^{-231}.
 \tag{R41.2}
\]

The first estimate is DO059(SC35), the second DO060(IB32), and the third R040(R40.1--2). In particular the second estimate includes nonscalar integral blocks and both scalar-block unipotent ranks; it is not a bound only on the fixed plane.

Define S41 to be the sum of the unchanged signed w(g) over the complement: s(g)\textgreater=3 and g outside P48, with the same support and cutoff. All removed sets and this complement are inverse closed. The kernel adjoint relation in K062 therefore makes S41 real. No individual term has been declared positive.

\subsection{2. Exact normalization and remaining sufficient inequality}\label{exact-normalization-and-remaining-sufficient-inequality}

Retain the strictly positive original constant

\[
 \mathcal N=w_x\,28C_\infty d_Rd_W
             \left(\frac{T-2}{T-1}\right)^f,
 \qquad f=168.
 \tag{R41.3}
\]

DO059(SC22) supplies the complete integral-plus-infinite error less than2\^{}-1999 at this same cutoff. Its infinite tail includes the entire old cutoff annulus. Partitioning its finite nonintegral sum by the three disjoint sets above gives

\[
 \boxed{\quad
 \frac{E_{x,\mathrm{ram}}}{\mathcal N}=1+S41+\varepsilon_{41},
 \qquad |\varepsilon_{41}|<\delta_{41},
 \quad}
 \tag{R41.4}
\]

where the explicitly charged error is

\[
 \delta_{41}=2^{-1999}+2^{-104546}+2^{-101550}+2^{-231}
                         <2^{-230}.
 \tag{R41.5}
\]

This proves the displayed identity and bound, not the sufficient nonvanishing condition

\[
                     S41>-1+\delta_{41}.
 \tag{R41.6}
\]

If R41.6 holds then R41.4 and N\textgreater0 imply E\_(x,ram)\textgreater0, which is the same all-T detector criterion already established in K062 and DO055. The coarser condition S41\textgreater-1+2\^{}-230 also suffices. Both inequalities remain open. No new source or integral-inverse premise has been inserted.

On J14, DO060 explicitly supplies no corresponding two-place estimate. Its current checked remainder is still R040(R40.6), outside the full Q(zeta28) fixed plane at8192, with error2\^{}-1999+2\^{}-794. Its sign remains open as well.

\subsection{3. Scope and verification}\label{scope-and-verification-2}

The proof consists of the disjoint partition, the already checked absolute estimates, the triangle inequality and the unchanged positive normalizer. It does not evaluate an automorphic period or enumerate the remaining rational matrices. The exact companion certificate checks the support multiplicity, error comparison and identities of the carried frozen proof versions; it is not an additional nonvanishing proof.

The construction continues to concern the same Hermitian theta correspondence and stack masses. The previously recorded SGA3 descent, SGA6 VI3.3 representation and SGA7 I VII--VIII biextension contexts are unchanged. There is no new identification with the specific032 cycle, no completed mixed Host, no net DP44 sign, and no Kummer RH conclusion.

\clearpage

\section{AM-K067 --- nonrational primary reductions in the actual finite short sum}\label{am-k067-nonrational-primary-reductions-in-the-actual-finite-short-sum}

Author: GPT-6 Astra Pro, ALP\_MACHINE /root/kummer\_extension. Inherited reasoning effort is not separately exposed.

\textbf{Author whole-proof check PASS --- 2026-11-22T10:04:01.277344+00:00.}

For the unchanged J56 family at the K044 upper endpoint, the complete s=3 short subset having a nonrational primary semisimple reduction at any unmixed G place has absolute contribution less than \(2^{-211449}\). The word nonrational here refers to the eigenvalue field of that finite semisimple reduction, not to the coefficient field of the global matrix.

After this removal, every remaining nonzero s=3 term has scalar semisimple reduction at all 334 unmixed G places. In the unique determinant-one representative of its central class, these scalars have order dividing 13. They are not asserted to be restrictions of one global root of unity.

A quantitative threshold also removes an explicit subset at every s on J56, and at s up to 17 on J14. The remaining finite sum retains its signs, weights, selected real trace and both mixed projectors. No all-T Mram nonvanishing or final arithmetic budget is asserted.

\subsection{1. Actual fixed support and the finite invariant}\label{actual-fixed-support-and-the-finite-invariant}

Keep the exact K062 scalar-orbit sum and DO059's J56 cutoff, with \[
(a,f,n,H)=
\begin{cases}
(42,42,43,8192),&J14,\\
(168,168,169,8192),&J56.
\end{cases}
\tag{PS1}
\] On J14 this is the original K062 cutoff. On J56 it is the proved DO059 refinement; its discarded annulus is already included in that finding's absolute tail error.

Let \(s(g)\) be the number of relative T places at which an actual supported rational g is nonintegral. Its summand is \[
a(g)c(g)\prod_{v\mid T}\rho(g_v).
\tag{PS2}
\] Here a is the existing genuine finite partial trace, c the normalized selected real trace, and rho the exact signed pencil trace. Thus \(|c|\le1\), and each nonintegral place contributes \(|\rho|\le2/(T-2)\). No new finite input or projector is introduced.

At each of the \(2(f-1)\) unmixed G places, the finite factor is the actual irreducible cuspidal representation of \(GL_n(\mathbb F_q)\). The two omitted places retain their actual arbitrary nonzero orthogonal mixed projectors; their normalized partial traces have modulus at most one.

For a term with \(a(g)\ne0\), Green's primary-support theorem says that the semisimple reduction at every unmixed place comes from one irreducible degree d dividing n.~Put \[
k_i(g)=\#\{v\mid g_i:\ v\text{ unmixed and }d_v>1\},
\quad k(g)=k_0(g)+k_1(g).
\tag{PS3}
\] The possible degrees are \(1,43\) on J14 and \(1,13,169\) on J56. A reduction with degree one has scalar semisimple part; it can have any permitted unipotent part.

Terms with \(a(g)=0\) contribute zero and can be removed before defining these invariants. Inversion preserves s and all the degrees d, so the subsets used below are inverse closed. The usual conjugation symmetry of K062 consequently makes each removed sum real, although the estimates bound absolute values before cancellation.

\subsection{2. The exact Green loss beyond a nonscalar reduction}\label{the-exact-green-loss-beyond-a-nonscalar-reduction}

Green's formula, in the exact Jordan form of Gorodetsky and Hazan (2017, Theorem 2.1 and equation (2.1), p.~5), gives \[
\frac{|\chi(g)|}{\chi(1)}
\le
\frac{d\prod_{j=1}^{t-1}(q^{dj}-1)}
     {\prod_{j=1}^{n-1}(q^j-1)},
\qquad t\le n/d.
\tag{PS4}
\] The factor d bounds the sum of the d conjugate regular-character values. No estimate for either mixed projector is obtained from this formula.

For d greater than one, define \[
A_d=\frac{n(n-1)}2-\frac d2\frac nd\left(\frac nd-1\right)
=\frac{n^2}2\left(1-\frac1d\right).
\tag{PS5}
\] Dropping the normalized numerator factors and enlarging t to n/d gives \[
\frac{|\chi(g)|}{\chi(1)}
\le d\,q^{-A_d}\prod_{j=1}^{n-1}(1-q^{-j})^{-1}.
\tag{PS6}
\] At a scalar-primary nonscalar reduction, the exact formula instead gives at most \(1/(q^{n-1}-1)\). At a scalar reduction the normalized modulus is one.

For n=169, \[
A_{13}=13182,\qquad A_{169}=14196,
\] whereas the baseline nonscalar exponent is a=168. For n=43 the nonrational exponent is \(A_{43}=903\), with baseline a=42. Set \[
(d_*,E_*)=
\begin{cases}
(43,861),&J14,\\
(13,13014),&J56.
\end{cases}
\tag{PS7}
\] For J56, the degree169 extra factor is dominated by the degree13 one: \(169q^{-14028}\le13q^{-13014}\), since \(q^{1014}\ge13\).

All denominator correction factors across all unmixed places together are less than two. Indeed, the product inequality \(\prod(1-x_j)\ge1-\sum x_j\), applied to their finite collection and then bounded by the convergent geometric sums, gives \[
\prod_{v\ {\rm unmixed}}\prod_{j=1}^{n-1}(1-q_v^{-j})
\ge1-\frac{2(f-1)}{q_{\min}-1}>\frac12.
\tag{PS8}
\] The actual primes satisfy \(q_{\min}-1>4(f-1)\). Thus no factor two per nonrational place is needed.

For use with the separate geometric-minor argument, PS4 also implies the following comparison. On J56, for every \(1\le r\le82\), a scalar-primary reduction with \(\operatorname{rank}(g-\lambda I)\ge r\) has ratio at most \[
\prod_{j=n-r}^{n-1}(q^j-1)^{-1}.
\tag{PS9}
\] This is direct cancellation in PS4, since the number of Jordan blocks is at most n−r. The same bound holds for the nonrational primary degrees 13 and169 at the actual primes: write \(L_r=rn-r(r+1)/2\le10455\). By PS6 its ratio is less than \(2d q^{-A_d}\), and \(q^{A_d-L_r}\ge2d\). Hence it is at most \(q^{-L_r}\), which is no larger than the right side of PS9. This supplies the character comparison for DO's geometric minor result; the geometric counting is not reproduced here.

\subsection{3. Nonrational primary terms have a nonscalar integral block}\label{nonrational-primary-terms-have-a-nonscalar-integral-block}

DO060(IB3--13) proves, independently of s, that the large positive block \[
A=g_{\langle e_5,\ldots,e_a\rangle,\langle e_5,\ldots,e_a\rangle}
\] is integral. If this block is globally scalar, say \(\lambda I\), unitarity forces \(\lambda\bar\lambda=1\), and the actual host root group gives \(\lambda\in\mu_{28}\). The matrix then fixes a positive subspace of codimension at most six after multiplication by \(\lambda^{-1}\). Therefore \[
\operatorname{rank}(g-\lambda I)\le6.
\tag{PS10}
\] At each G prime this rank inequality persists: its minors vanish identically over E, and all entries and lambda are integral there. Since n\textgreater6, the finite reduction has the eigenvalue lambda in the base residue field. A nonzero cuspidal character with primary degree greater than one has no such eigenvalue. Consequently \[
a(g)\ne0,\quad k(g)\ge1
\quad\Longrightarrow\quad A\text{ is nonscalar}.
\tag{PS11}
\]

Choose a nonzero integral off-diagonal entry of A or a difference of its diagonal entries. All its conjugates have modulus at most 2H. If the full reduction of g is scalar at a G place, this entry vanishes at that place and its conjugate. The exact integral norm argument of DO060(IB8) bounds the logarithmic weight of scalar G places by \(f\log(2H)\), with no T denominator.

Combine that norm bound, PS6--8 and the baseline nonscalar loss. For every term in PS11, \[
\boxed{\ 
|a(g)|
\le
\frac{2(2H)^{af}}{Q^{a(f-1)}}
\left(\frac{d_*}{g_0^{E_*}}\right)^{k_0(g)}
\left(\frac{d_*}{g_1^{E_*}}\right)^{k_1(g)}
\le
\frac{2(2H)^{af}}{Q^{a(f-1)}}\delta^{\,k(g)},
\quad
\delta=\frac{d_*}{q_{\min}^{E_*}}.
\ }
\tag{PS12}
\] The two mixed factors are still bounded only by one. There is no change to their projection subspaces or normalized traces.

\subsection{4. Complete counting and exact thresholds}\label{complete-counting-and-exact-thresholds}

Use the complete actual target-over-lattice and fixed-column count of DO059(SC28--30): \[
N_s\le
\binom fs(T-1)^{2s}G_n(2H+1)^{2fn}T^{2ns},
\quad
G_n=\gcd\bigl(|GL_n(\mathbb F_{g_0})|,
             |GL_n(\mathbb F_{g_1})|\bigr).
\tag{PS13}
\] It already counts every source/target pencil pair and every pole depth. Subdividing this set by the finite invariant k requires no additional binomial factor for positions of the nonrational reductions.

Define \[
\begin{split}
B_s={}&
\frac{2G_n(2H+1)^{2fn}(2H)^{af}}{Q^{a(f-1)}}
\binom fs
\frac{2^s(T-1)^{2s}T^{2ns}}{(T-2)^s}.
\end{split}
\tag{PS14}
\] For every K at least one, \[
\sum_{\substack{s(g)=s\\ k(g)\ge K}}
\left|a(g)c(g)\prod_v\rho(g_v)\right|
\le B_s\delta^K.
\tag{PS15}
\] This is a bound on the complete actual subset, not an individual-term estimate.

The companion certificate uses exact integer comparisons. If \(C_0\) denotes the first fraction in PS14, it proves:

\begin{longtable}[]{@{}
  >{\raggedright\arraybackslash}p{(\columnwidth - 8\tabcolsep) * \real{0.1579}}
  >{\raggedleft\arraybackslash}p{(\columnwidth - 8\tabcolsep) * \real{0.2105}}
  >{\raggedleft\arraybackslash}p{(\columnwidth - 8\tabcolsep) * \real{0.2105}}
  >{\raggedleft\arraybackslash}p{(\columnwidth - 8\tabcolsep) * \real{0.2105}}
  >{\raggedleft\arraybackslash}p{(\columnwidth - 8\tabcolsep) * \real{0.2105}}@{}}
\toprule\noalign{}
\begin{minipage}[b]{\linewidth}\raggedright
Host
\end{minipage} & \begin{minipage}[b]{\linewidth}\raggedleft
\(C_0<2^{-m_0}\)
\end{minipage} & \begin{minipage}[b]{\linewidth}\raggedleft
\(\delta<2^{-m_\delta}\)
\end{minipage} & \begin{minipage}[b]{\linewidth}\raggedleft
\(T<2^b\)
\end{minipage} & \begin{minipage}[b]{\linewidth}\raggedleft
\(z=2+(2n+1)b\)
\end{minipage} \\
\midrule\noalign{}
\endhead
\bottomrule\noalign{}
\endlastfoot
J14 & \(m_0=496\) & \(m_\delta=19055\) & \(b=1001\) & 87089 \\
J56 & \(m_0=776478\) & \(m_\delta=447357\) & \(b=996\) & 337646 \\
\end{longtable}

Since \((T-1)^2/(T-2)<2T\) and \(\binom fs\le2^f\), \[
B_s\delta^K
<
2^{-m_0+f+zs-m_\delta K}.
\] Therefore the explicit integer threshold \[
K_s=
\max\left\{1,
\left\lceil\frac{4096-m_0+f+zs}{m_\delta}\right\rceil
\right\}
\tag{PS16}
\] makes the complete subset in PS15 less than \(2^{-4096}\), whenever \(K_s\le2(f-1)\).

Selected values are:

\begin{longtable}[]{@{}lrr@{}}
\toprule\noalign{}
Host & s & \(K_s\) \\
\midrule\noalign{}
\endhead
\bottomrule\noalign{}
\endlastfoot
J14 & 1 & 5 \\
J14 & 2 & 10 \\
J14 & 3 & 14 \\
J14 & 17 & 78 \\
J14 & 18 & 83, exceeding the 82 available factors \\
J56 & 3 & 1 \\
J56 & 4 & 2 \\
J56 & 10 & 6 \\
J56 & 17 & 12 \\
J56 & 168 & 126 \\
\end{longtable}

Every J56 threshold is available. The union of these removed subsets over all s has absolute sum less than \(168\,2^{-4096}<2^{-4088}\). On J14 the thresholds are available exactly for s from1 through17, and that union has absolute sum less than \(17\,2^{-4096}<2^{-4091}\). Thresholds beyond that range make no assertion on J14.

Without the coarse binary replacement of PS14, exact cross-multiplication gives \[
\sum_{\substack{s(g)=3\\k(g)\ge1}}|\text{actual summand}|
<2^{-211449}
\quad(J56).
\tag{PS17}
\] For s=4 and k at least two the analogous exact exponent is321345. The certificate also checks the J14 examples and records failure of the indicated bound at s=18, k=82; it does not conceal that failure.

\subsection{5. The exact remaining finite reductions}\label{the-exact-remaining-finite-reductions}

Remove DO059's entire s=1 stratum, DO060's entire s=2 stratum, and the disjoint union of the subsets in PS16 for s at least three. The remaining J56 sum can be taken over \[
s(g)\ge3,\qquad k(g)<K_{s(g)}.
\tag{PS18}
\] The exact K062 energy identity, with its original normalizer, is still \[
\frac{E_{x,\mathrm{ram}}}
{w_x\,28C_\infty d_Rd_W((T-2)/(T-1))^f}
=
1+S_{\mathrm{res}}+\varepsilon,
\] where \[
|\varepsilon|
<
2^{-1999}+2^{-104546}+2^{-101550}+2^{-4088}
<
2^{-1998}.
\tag{PS19}
\] If the already bounded E12-plane subset is removed too, take its set difference with the preceding union to avoid double counting. Its full bound \(2^{-21414}\) can be added to the right side without changing the final \(2^{-1998}\) bound.

At s=3, \(K_3=1\), so every remaining nonzero term has scalar semisimple reduction at all 334 unmixed places. More concretely, the determinant of every supported g is a global norm-one unit: all actual pencil lattices have the same local determinant ideal, and g is integral and invertible away T. Thus \(\det g=\delta_g\in\mu_{28}\). Since \(\gcd(n,28)=1\), each central class has a unique determinant-one representative. On J56, n=169 is congruent to one modulo28, so it is simply \[
h=\delta_g^{-1}g.
\tag{PS20}
\] At any scalar-primary G reduction of h, write its semisimple part as \(\eta_v I\). Determinant one gives \(\eta_v^{169}=1\). At both actual G primes the certificate verifies \[
\gcd(169,g_i-1)=13.
\] It follows that \[
\operatorname{red}_v(h)=\eta_v u_v,\qquad
\eta_v\in\mu_{13}(\mathbb F_{g_i}),\quad u_v\text{ unipotent}.
\tag{PS21}
\] All 334 unmixed reductions of each remaining s=3 term have this form.

The scalars \(\eta_v\) may vary with v. The global field's root group is \(\mu_{28}\), so no nontrivial global thirteenth root can remove them simultaneously. PS21 is not a statement that all these reductions are unipotent.

For comparison, \(\gcd(43,g_i-1)=1\) on J14. Thus if a J14 term has k=0, all its unmixed reductions in the unique determinant-one representative are actually unipotent. The quantitative J14 thresholds do not force k=0 throughout its remaining sum.

\subsection{6. Sources, checks and remaining scope}\label{sources-checks-and-remaining-scope}

The new calculation uses the exact fixed family of K062, the actual labelled-pencil support of K063, the complete count and smaller J56 cutoff of DO059, and the integral block of DO060. The determinant root group is the checked K031/K032 host group. No generic spherical or alternative full-Heisenberg family is substituted.

Primary finite-character sources:

\begin{itemize}
\tightlist
\item
  Green, J. A. (1955). The characters of the finite general linear groups. Transactions of the American Mathematical Society, 80(2), 402--447. Section5, Example(ii), pp.430--431. \url{https://www.ams.org/tran/1955-080-02/S0002-9947-1955-0072878-2/S0002-9947-1955-0072878-2.pdf}
\item
  Gorodetsky, O., \& Hazan, Z. (2017). On certain degenerate Whittaker models for cuspidal representations of \(GL_{kn}(\mathbb F_q)\). arXiv:1707.07308v2. Theorem2.1, equation(2.1), p.5. \url{https://arxiv.org/pdf/1707.07308}
\end{itemize}

I read the exact primary Theorem2.1 and equation2.1 both in the retained local text and through the primary PDF during this calculation. The algebraic and quantitative deductions PS5--21 are supplied above. The comparison PS9 is shared with the independently developed geometric high-minor argument; that geometric result is not assumed here.

The certificate \nolinkurl{nonrational_primary_short_sum_certificate.py} uses exact integers and fractions. It checks 771 small Green-envelope cases, all actual primary-degree and scalar-order data, the single global Euler correction, every host threshold and the exact selected large-integer margins. It also verifies the complete error budget. It does not enumerate the actual rational remainder, evaluate its signs or phases, or carry out a Lean build.

The remaining sum in PS18 is still uncomputed. In particular PS21 is a constraint on local finite reductions, not a global unipotence theorem, a positivity theorem, or an actual period nonvanishing assertion. No existence of a nonzero term in the removed subset is needed for its bound, and none is asserted.

\subsection{Author whole-proof verification}\label{author-whole-proof-verification-1}

At 2026-11-22T10:04:01.277344+00:00, I completed a whole read of PS1--21 and all explanatory/source/scope sections, and the complete 151-line certificate. I reconstructed the exact Green exponents, the single global Euler correction, scalar-block exclusion, no-extra-multiplicity count, all threshold inequalities, actual determinant normalization and residual thirteenth-root phases. I directly read the primary Theorem2.1/equation2.1 and the carried actual integral-block/pencil interfaces. The final status and TeX spacing edits were re-read. No scalar-primary reduction was incorrectly identified with a unipotent one on J56.

The certificate passed 1,607 exact assertions, including 771 finite Green-envelope cases, at 2026-11-22T10:01:16.905212+00:00. An earlier exploratory all-threshold loop was voluntarily interrupted because it repeated unnecessarily large powers; it supplied no claimed PASS. The final certificate uses proved binary bounds and retains the actual exact selected margins. No global rational sum, automorphic period or Lean build was computed.

\clearpage

\section{AM-R-042 --- primary-degree thresholds in the current plane complement}\label{am-r-042-primary-degree-thresholds-in-the-current-plane-complement}

Author: GPT-6 Astra Pro, ALP\_MACHINE /root, 22 November 2026 UTC. Reasoning effort is not externally exposed.

This is a corollary combining the whole-checked K067 thresholds with the current R041/R040 finite remainders. All matrices, signed weights, class masses and cutoffs are unchanged. The remaining sign is open.

For each supported nonzero term let s be its number of nonintegral relative T places, and k the number of nonrational-primary unmixed G reductions, exactly as in K067 PS3. Let K\_s be the explicit integer threshold in PS16. On J56 it is available for every1\textless=s\textless=168. The full union of subsets k\textgreater=K\_s has absolute contribution below2\^{}-4088. On J14 the available range is1\textless=s\textless=17, with union bound2\^{}-4091. These are complete subset bounds from K067; no extra count of finite positions is needed.

\subsection{1. J56}\label{j56}

R041 has already removed all s=1, all s=2, and the complete fixed plane over L48=Q(zeta28,sqrt5,sqrt13). Inside its remaining set remove

\[
 U_{56}=\{g:s(g)\ge3,\ g\notin P48,\ k(g)\ge K_{s(g)}\}.
 \tag{R42.1}
\]

This is a subset of the full union bounded in K067. It is disjoint from all previously removed terms by construction. Hence its full absolute contribution is at most that same strict bound2\^{}-4088. The new signed remainder S42,56 is precisely

\[
 \{g:s(g)\ge3,\ g\notin P48,\ k(g)<K_{s(g)}\},
 \tag{R42.2}
\]

with cutoff8192 and all original support conditions. Inversion preserves these conditions, so the kernel adjoint symmetry makes the sum real. With exactly the original positive normalizer N of R041(R41.3),

\[
 \frac{E_{x,\mathrm{ram}}}{\mathcal N}
 =1+S42,56+\varepsilon_{42,56},\qquad
 |\varepsilon_{42,56}|<\delta_{42,56},
 \tag{R42.3}
\]

\[
 \delta_{42,56}=2^{-1999}+2^{-104546}+2^{-101550}
                          +2^{-231}+2^{-4088}<2^{-230}.
 \tag{R42.4}
\]

For the s=3 part of this remainder, K\_3=1. Every one of its334 unmixed reductions therefore has scalar semisimple part. K067 PS20--21 gives its unique determinant-one representative h=det(g)\^{}(-1)g and

\[
 \operatorname{red}_v h=\eta_v u_v,\qquad
 \eta_v\in\mu_{13}(\mathbf F_{g_i}),\qquad u_v\text{ unipotent}.
 \tag{R42.5}
\]

These local scalars remain in the statement. They are not replaced by one or by a single global thirteenth root. No global unipotence or positivity is inferred.

\subsection{2. J14}\label{j14}

Inside R040(R40.6)'s original complement of the complete Q(zeta28) fixed plane, remove the subset with1\textless=s\textless=17 and k\textgreater=K\_s. Its full absolute contribution is below2\^{}-4091. The complementary signed sum S42,14 therefore retains the original cutoff8192 and satisfies

\[
 \frac{E_{x,\mathrm{ram}}}{\mathcal N}
 =1+S42,14+\varepsilon_{42,14},\qquad
 |\varepsilon_{42,14}|<2^{-1999}+2^{-794}+2^{-4091}<2^{-793}.
 \tag{R42.6}
\]

For1\textless=s\textless=17 it has k\textless K\_s; for s\textgreater=18 no new threshold restriction is asserted. The J56 one-place and two-place estimates are not imported.

\subsection{3. Proof, verification and remaining scope}\label{proof-verification-and-remaining-scope}

Each statement follows by partitioning the preceding signed sum into the displayed subset and its complement, applying K067's absolute bound to that subset, and adding the errors by the triangle inequality. Intersecting with the preceding plane complement prevents double counting. The complete tails and old annuli remain those of R040/R041. The same exact finite phases, real traces and class weights are retained.

The sufficient inequalities S42,56\textgreater-1+delta42,56 and S42,14\textgreater-1+2\^{}-1999+2\^{}-794+2\^{}-4091 remain unproved. Thus this is a reduction of the explicit finite gate, not an actual nonzero period or a net DP44 result. It adds no conclusion about Kummer RH or the specific032 arithmetic pairing. The Grothendieck dictionary remains the same metrized theta correspondence and representation setting; no new comparison theorem is attributed to those texts.

The companion certificate checks the frozen K067 and R041 versions, all actual threshold formulas and the two rational error comparisons. It computes no automorphic value or sign of a rational remainder.

\clearpage

\section{ALP-LIFT-50 --- normalized selected hook costs and the actual eleven-adic certificate}\label{alp-lift-50-normalized-selected-hook-costs-and-the-actual-eleven-adic-certificate}

Author: GPT-6 Astra Pro, ALP\_MACHINE \nolinkurl{/root/lift_accounting}. 22 November 2026 UTC; reasoning effort not separately exposed.

Status: author whole-proof check PASS, 2026-11-22T10:08:50Z.

This result computes a further part of the actual two-edge arithmetic budget. It removes the complete algebraic branch and coordinate scalars with their finite contents, supplies bounded allowed Weyl detectors, and compares the exact remaining complete-coevaluation norm with the proved eleven-adic image bound. The last comparison is a norm-factor statement. It is not a lower bound for an automorphic theta operator, an upper bound for its actual image content, or a no-go theorem for a shorter lift chain.

All claims concern the fixed K038/K041/K044 family on the two actual hosts, with its unchanged raw full Y sum, strict levels, source metric, and terminal decoration. The existing Host/SizeData2 and all-T Mram nonvanishing issues are not replaced by a scalar calculation.

\subsection{1. Actual weights and unambiguous degrees}\label{actual-weights-and-unambiguous-degrees}

Write \nolinkurl{a=f=42} or168, \nolinkurl{d=1760} or1634, \nolinkurl{b=a+d=1802}, \nolinkurl{c=a+2d}, and \nolinkurl{N_T=c-1} for the terminal coefficient rank. Retain the exact K044 parameters

\[
 A=2^{960},\quad B=2^{955},\quad
 \ell_-\le\ell\le\ell_+,\quad
 \ell\in28\mathbb Z,\quad
 q\in28(a+1)\mathbb Z_{\ge0},\quad q\le\ell.
 \tag{HG1}
\]

At a label of role beta let E be its fixed rounded integer E\_T, E\_H, or E\_L. Keep the actual binary word omega and the u,v lists of K044(EG4--EG5). Define

\[
 \nu_{0,i}=u_i-\frac{a-1}{2}+i+E-\frac b2,
 \quad
 \Delta=2\ell+q-d,
 \quad
 \nu_1=(\nu_0,0^d)+\Delta\mathbf1_b.
 \tag{HG2}
\]

These are the selected kernel partitions, not the coefficient partitions. Write

\[
 D_{j,\tau}=|\nu_{j,\tau}|,\quad
 D_j=\sum_\tau D_{j,\tau},\quad D=D_0+D_1,
 \qquad H_j=\prod_\tau H_{\nu_{j,\tau}}.
 \tag{HG3}
\]

The intermediate compact representation has polynomial partition \nolinkurl{(nu0,0^d)}, so put \nolinkurl{d_1=product_tau dim_GL_b(nu0_tau,0^d)}. The initial coefficient polynomial degrees are denoted D\_in and the terminal ones D\_out; neither symbol means D\_0 or D\_1. Let M\_in and M\_out sum the absolute coefficient determinant weights, and let Lambda\_0 be the sum of the initial highest weights, exactly as in TF47(TF12--TF19).

The terminal raw lambda and mu lists are R024(R24.2), with the literal substitutions \nolinkurl{e1=E+2ell+q}, \nolinkurl{b1=E+ell+q}. Their common last weight is

\[
 s=-E-2\ell-q+b/2.
 \tag{HG4}
\]

For \nolinkurl{eta_mu=mu-s 1_(c-1)}, direct summation gives

\[
 D_{{\rm out},\tau}
  =(a-1)\left(E+2(a+1)\ell-\frac{a+3d}{2}\right)
     +(d+1)(2\ell-d)+(b-1)q.
 \tag{HG5}
\]

The q term has coefficient b-1: precisely that many entries are unchanged before subtracting the last weight. The entry containing \nolinkurl{-b1} and the final block change by \nolinkurl{-q} already. Thus no q=0 terminal degree is silently reused.

The positive part of eta\_mu has length \nolinkurl{b=1802}. Its lists are

\[
 \left(E+(4a-2-4i)\ell+q-a-3d/2+1+i\right)_{i=0}^{a-2},
 \quad 3\ell+q-3d/2,\quad
 \Delta^{[d-1]},\quad \ell-d/2,
 \tag{HG6}
\]

followed by \nolinkurl{d-1} zero entries. This also verifies HG5.

\subsection{2. Remove only scalars whose finite contributions are retained}\label{remove-only-scalars-whose-finite-contributions-are-retained}

Let

\[
 s_B=u_TC,\quad
 u_T=\prod_\tau \tau(\operatorname{star}\pi_T)^{-n_{{\rm Br},\tau}},
 \quad C=\prod_\tau c_{{\rm Br},\tau},\qquad
 C_D=\prod_\tau cD_\tau^{D_{0,\tau}+D_{1,\tau}}.
 \tag{HG7}
\]

The exact phase of s\_B is the source phase in R019. The separate \nolinkurl{branch_scalar_normalization_erratum.md} applies: removing C alone would leave \nolinkurl{|u_T|} in the archimedean factor. Here the complete algebraic scalar s\_B is removed.

Let B\_ar be the actual arithmetic bilinear row map of OB24/GM21. Define the scalar-normalized map

\[
                 \widehat B=(s_BC_D)^{-1}B_{\rm ar}.
 \tag{HG8}
\]

This is a normalization of the same field-valued row, not a new choice of kernel, section, finite level, or image. In particular its nonvanishing is equivalent to that of B\_ar.

The sign of the C\_D factor is fixed by K035(PF26--PF28): the actual jetC is an ordinary derivative multiplied by \nolinkurl{cD_tau} once per derivative. Every term of the complete subD composition has the prescribed degree at each label. Hence C\_D factors out of every class/Y summand of B\_ar, with exponent in HG7. Meanwhile M\_sigma contains \nolinkurl{(lambda_tau/cD_tau)} to that same exponent. The analytic reading is divided by s\_B simultaneously: if its original value is \nolinkurl{T_sigma=M_sigma B_ar}, then \nolinkurl{T_sigma/s_B,sigma=(M_sigma C_D,sigma) widehat B}. The creation operators and the unit vacuum are unchanged. Thus the comparison scalar for this normalized reading is

\[
 \widehat{\mathfrak M}_\sigma=\mathfrak M_\sigma C_{D,\sigma},
 \qquad
 |\widehat{\mathfrak M}_\sigma|
     =\mathcal P_J^{N_\theta}
        \prod_\tau|\lambda_{\sigma,\tau}|^{D_{0,\tau}+D_{1,\tau}},
 \quad N_\theta=14h(ab+bc).
 \tag{HG9}
\]

Here h=1 or4 and P\_J is the actual unit-vacuum period of K066(PC3). No field embedding is applied to a transcendental number; all moduli are reconstructed as in K066.

K035(PF17) gives the following exact scalar contents, normalized by v\_p(p)=1 at every place over the indicated rational prime:

\[
\begin{array}{c|ccccc}
 p&2&7&29&5,13&\text{all other primes}\\\hline
 v_p(C_D),\ J14&D&D/3&6D/7&0&0\\
 v_p(C_D),\ J56&D&D/3&6D/7&D/2&0.
\end{array}
 \tag{HG10}
\]

In particular C\_D is an eleven-adic unit. Its exact normalized archimedean average is

\[
 \overline{\log|C_D|}
  =D\left(\log2+\frac13\log7+\frac67\log29
                  +\mathbf1_{h=4}\frac12\log65\right).
 \tag{HG11}
\]

Either taking the field norms of the explicit algebraic factors in PF15, or summing HG10 with local degrees, proves HG11. The same scalar is therefore removed on both sides of the product formula. The branch scalar is treated identically by R024.20--21. This removes no non-scalar bad-prime lattice or reduced-gauge pole.

\subsection{3. Remaining creation costs, with the physical labels fixed}\label{remaining-creation-costs-with-the-physical-labels-fixed}

Adding HG11 to K066(PC19--PC20) gives the exact mean

\[
 L_{h,i}:=\overline{\log|\lambda_{\sigma,\tau}|}
  =\frac12\log\frac{2G_i}{\pi}
    +\frac37\log29+\mathbf1_{h=4}\frac14\log65,
 \qquad i=i(\tau).
 \tag{HG12}
\]

The rational differential index i remains fixed while the CM label is transported. Put \nolinkurl{D_i=sum_(tau with physical i)(D0_tau+D1_tau)}. Then

\[
 \overline{\log|\widehat{\mathfrak M}_\sigma|}
             =N_\theta\log\mathcal P_J+\sum_{i=1}^3D_iL_{h,i}.
 \tag{HG13}
\]

The exact rational intervals of K066 yield these sharpened bounds:

\begin{longtable}[]{@{}
  >{\raggedright\arraybackslash}p{(\columnwidth - 8\tabcolsep) * \real{0.1579}}
  >{\raggedleft\arraybackslash}p{(\columnwidth - 8\tabcolsep) * \real{0.2105}}
  >{\raggedleft\arraybackslash}p{(\columnwidth - 8\tabcolsep) * \real{0.2105}}
  >{\raggedleft\arraybackslash}p{(\columnwidth - 8\tabcolsep) * \real{0.2105}}
  >{\raggedleft\arraybackslash}p{(\columnwidth - 8\tabcolsep) * \real{0.2105}}@{}}
\toprule\noalign{}
\begin{minipage}[b]{\linewidth}\raggedright
Host
\end{minipage} & \begin{minipage}[b]{\linewidth}\raggedleft
i
\end{minipage} & \begin{minipage}[b]{\linewidth}\raggedleft
Lower numerator for L/(log2)
\end{minipage} & \begin{minipage}[b]{\linewidth}\raggedleft
Upper numerator
\end{minipage} & \begin{minipage}[b]{\linewidth}\raggedleft
Denominator
\end{minipage} \\
\midrule\noalign{}
\endhead
\bottomrule\noalign{}
\endlastfoot
J14 & 1 & 282 & 284 & 84 \\
J14 & 2 & 271 & 274 & 84 \\
J14 & 3 & 310 & 314 & 84 \\
J56 & 1 & 409 & 411 & 84 \\
J56 & 2 & 398 & 401 & 84 \\
J56 & 3 & 437 & 440 & 84 \\
\end{longtable}

Indeed exponentiating \nolinkurl{84 L} gives exactly

\[
 (2G_i/\pi)^{42}29^{36}65^{21\mathbf1_{h=4}}.
 \tag{HG14}
\]

Comparing its rational interval with powers of two proves the table. This is distinct from the pointwise bound. For example, K066's \nolinkurl{1<G_i/(2pi)<3}, \nolinkurl{1/8<pi_l<16}, and \nolinkurl{0<delta_rho<=1} imply \nolinkurl{|lambda_sigma,tau|>2^-8} in every realization, while HG12 is the more informative normalized average. No finite valuation is assigned to lambda itself.

\subsection{4. Actual eleven-adic contents after normalization}\label{actual-eleven-adic-contents-after-normalization}

Write z(n) for the number of base-eleven digits of n.~The whole actual class-section estimate K039/K041 and its K044 transport give

\[
 E_0=ab\sum_\tau z(D_{0,\tau}),\quad
 E_1=bc\sum_\tau z(D_{1,\tau}),\qquad
 L_j=\left\lfloor D_j/1220\right\rfloor-E_j.
 \tag{HG15}
\]

These L\_j are positive on the entire allowed interval. LIFT34's integral full-middle-mass cutoff preserves \nolinkurl{L0+L1}; the remaining terminal mass loss is delta\_T=384 or357. For the two literal terminal lattices put

\[
 g_\Delta=L_0+L_1-\delta_T,\qquad
 g_\nabla=g_\Delta-v_{11}(H_{\rm out}),\quad
 H_{\rm out}=\prod_\tau H_{\eta_{\mu,\tau}}.
 \tag{HG16}
\]

Thus for every actual integral initial Mram row P,

\[
 v_{11}(\widehat B(P,Z))\ge
 \begin{cases}
 g_\Delta,& Z\in\mathcal M_\Delta=\mathrm{MtildeLdel},\\
 g_\nabla,& Z\in\mathcal M_\nabla=\mathrm{MtildeL}.
 \end{cases}
 \tag{HG17}
\]

This holds at every coefficient-field place over eleven, including zero values. It is the literal R024 perfect Weyl/dual-Weyl pairing and MC22, after subtracting v(s\_B). The additional C\_D is a unit there. Neither g is asserted to be the actual content of the restricted nonzero image.

For clarity, the uniformly valid numerical consequences are

\[
 \frac D{1221}<g_\Delta\le\frac D{1220},
 \qquad g_\nabla<-\frac D{42}.
 \tag{HG18}
\]

The second inequality concerns the explicit lower-bound formula g\_nabla. It is not an upper bound on v11 of a value, and does not assert a negative actual gain.

\textbf{Proof of HG18.} The exact degrees of K041(NW4), with the extra \nolinkurl{bq} on edge one, give \nolinkurl{D_j,tau>=2^480}; every degree is below \nolinkurl{2^980}, so it has at most300 base-eleven digits. Therefore \nolinkurl{g_Delta>D/1220-2-300f(ab+bc)-delta_T}, which is greater than \nolinkurl{D/1221} by the exact integer comparison at \nolinkurl{D>=2^480}.

Next \nolinkurl{D_out>=D/4} uniformly. This can be checked without a word enumeration: write W for the Hamming weight and use HG5 and NW4. For a non-top label, four times HG5 minus \nolinkurl{D0_tau+D1_tau} is

\[
 (2a-4)E+(4a^2-14a+6d+12+2W)\ell+(3b-4)q+C_0,
 \tag{HG19}
\]

where \nolinkurl{C0=-2(a-1)(a+3d)-4d(d+1)-4+a+b^2}. All three displayed coefficients are positive and \nolinkurl{E>b^2}. Their positive E term exceeds \nolinkurl{|C0|}. At a top label replace the ell coefficient by \nolinkurl{4a^2+2a+6d+8+2W} and C0 by \nolinkurl{C0+2}, and subtract \nolinkurl{2(A-B)}. Writing \nolinkurl{E=B+R} makes the remaining B term \nolinkurl{(2a-66)B>0}; it alone exceeds the constant's absolute value. This proves the claimed degree inequality for every actual word, ell and q.

Use the factorial/Vandermonde hook formula on HG6, whose positive length is r=1802. Every factorial argument and denominator factor is below \nolinkurl{2^962<11^300}. Legendre's identity gives \nolinkurl{v11(n!)=(n-s11(n))/10>=n/10-300}. Every denominator has valuation at most300. Consequently

\[
 v_{11}(H_{\rm out})\ge D_{\rm out}/10-150fr(r+1)
                         \ge D/40-150fr(r+1).
 \tag{HG20}
\]

Combine \nolinkurl{g_Delta<=D/1220} with HG20. The exact comparison \nolinkurl{(1/40-1/1220-1/42)2^480>150fr(r+1)} proves the strict second inequality of HG18. The role roundings and the elementary integer inequalities used here are independently checked in the certificate.

\subsection{5. Bounded allowed Weyl detectors avoid that hook loss}\label{bounded-allowed-weyl-detectors-avoid-that-hook-loss}

The standard finite cyclotomic construction R038 applies at any rank, with degree below T. Here each terminal label has \nolinkurl{D_out,tau<N_T*2A<2^973<T}, by HG6/K044 and the actual T bit bounds. The existing coefficient field already contains mu\_T. Therefore the terminal standard Weyl module admits the integral field-spanning family

\[
 \rho(U(\zeta))h_{\eta_\mu},\quad
 \zeta\in\mu_T^{N_T(N_T-1)/2},\qquad
 \|\rho(U(\zeta))h_{\eta_\mu}\|\le N_T^{D_{{\rm out},\tau}}.
 \tag{HG21}
\]

This is exactly R038's lower-unipotent grid, with the literal transpose action. Its integral vectors use no Fourier denominator; Fourier inversion is used only to prove their field span.

Transport them with the unchanged TF47 adapted terminal frames and independent determinant sections, then apply K059/DO056's unchanged stabilizer sums and finite projectors. Integral group actions preserve the terminal Weyl sublattice, so these vectors lie in the allowed \nolinkurl{MtildeLdel}, not merely in an unrelated rational model. Their field span is the same terminal equivariant space as for \nolinkurl{MtildeL}, because the Weyl and dual-Weyl lattices have the same characteristic-zero representation and the same equivariance conditions. Finite clearing factors remain precisely those of DO056; no new coefficient field or eleven-adic inverse is introduced.

Thus, conditional on nonvanishing of the actual Mram-restricted bilinear row, there are integral inputs with a Weyl detector such that the value is nonzero and

\[
 \|P\|\,\|Z\|\le
 2^{K_{\rm fin}}\widehat B_0\widehat B_{T,\Delta},\qquad
 \widehat B_{T,\Delta}
     =N_T^{D_{\rm out}}\widehat B_T.
 \tag{HG22}
\]

Here B0 and BT are TF15--TF16. This is a field-spanning argument: a nonzero bilinear form cannot vanish on both spanning families. It does not solve an integral inverse or assert a saturated image. Since \nolinkurl{N_T<2^12}, its concrete terminal logarithmic bound is

\[
 \log_2\widehat B_{T,\Delta}<
 \begin{cases}
 43D_{\rm out}+13M_{\rm out},&J14,\\
 47D_{\rm out}+17M_{\rm out},&J56.
 \end{cases}
 \tag{HG23}
\]

The corresponding dual-Weyl bounds remain31 and35 instead of43 and47, with the g\_nabla bound of HG17. These are two sufficient cost/gain choices, not optimal or necessary costs.

\subsection{6. The exact complete-coevaluation metric factor}\label{the-exact-complete-coevaluation-metric-factor}

Write \nolinkurl{K_B=product_tau kappa(lambda_tau,mu_tau)}. After removing s\_B and C\_D, OB19--OB24 give the exact remaining norm factor

\[
 \mathfrak F_\sigma=
 \frac{\sqrt{d_1 H_0H_1 K_B}}
      {|\widehat{\mathfrak M}_\sigma|}.
 \tag{HG24}
\]

The factor d1 arises from the complete middle coevaluation. There is no \nolinkurl{1/d1} in this formula. The ordinary coefficient metrics are already normalized by their own highest hook norms; redividing them would change the actual coefficient lattices.

The explicit hook and dimension formulas are

\[
 H_\nu=\frac{\prod_{i=0}^{r-1}(\nu_i+r-1-i)!}
                   {\prod_{i<j}(\nu_i-\nu_j+j-i)},\qquad
 \dim_{GL_b}(\nu_0,0^d)
       =\frac{\prod_{i=0}^{a-1}(\nu_{0,i}+b-1-i)!/(b-1-i)!}
                    {H_{\nu_0}},
 \tag{HG25}
\]

and the exact shift identity is

\[
 H_{\nu_1}=H_{\nu_0}
  \prod_{i=0}^{a-1}
    \frac{(\nu_{0,i}+b-1-i+\Delta)!}{(\nu_{0,i}+b-1-i)!}
  \prod_{j=0}^{d-1}\frac{(j+\Delta)!}{j!}.
 \tag{HG26}
\]

These are R011's hook and actual row-shift formulas; the dimension formula is the usual Weyl formula after cancelling its zero block. All arguments are explicit integers in the actual profile. Legendre floor sums compute their eleven-adic valuations without expanding a polynomial or a factorial. The actual kappa is exactly R019.4. No claim that the hook product is a DP44 height constant is made.

\subsubsection{A uniform, actual lower comparison for this factor}\label{a-uniform-actual-lower-comparison-for-this-factor}

For every allowed word, role assignment, ell and q,

\[
 \boxed{\qquad
  \overline{\log_2\mathfrak F_\sigma}
                   -g_\Delta\log_2 11>220D.
 \qquad}\tag{HG27}
\]

This is a comparison of an exact norm factor with the \textbf{proved lower-bound certificate} g\_Delta. It is not an upper bound for actual divisibility or a lower bound for the theta operator.

\textbf{Proof.} For a partition of degree m and length at most r, the number of standard tableaux is at most \nolinkurl{r^m}: the row containing each integer determines the tableau. The hook formula therefore gives \nolinkurl{H>=m!/r^m}. The elementary integral bound \nolinkurl{log(m!)>=m log m-m} and \nolinkurl{e<4} imply \nolinkurl{log2 H>=m(log2 m-log2 r-2)}. Here every selected label degree is at least \nolinkurl{2^480}, and each selected partition has length at most b\textless{}\nolinkurl{2^11}. Thus

\[
                       \log_2(H_0H_1)>467D.
 \tag{HG28}
\]

We also retain the size of the actual branch contraction, rather than dropping it. Put \nolinkurl{alpha_j=mu_j-lambda_(j+1)} and \nolinkurl{beta_i=lambda_i-mu_i}, in zero-based indices. R019.4 writes each nontrivial kappa factor as

\[
 \frac{(x+\beta_i)!(x+\alpha_j)!}
      {x!(x+\alpha_j+\beta_i)!},\quad x\ge0.
\]

It is at least \nolinkurl{binomial(alpha_j+beta_i,alpha_j)^-1}, and hence at least \nolinkurl{2^(-alpha_j-beta_i)}. Its monotonic increase in x follows from the product \nolinkurl{product_(s=1)^alpha (x+s)/(x+beta+s)}.

For the actual lists the nonzero alpha entries are \nolinkurl{ell*omega_j} for j=0,\ldots,a-3, \nolinkurl{Astar} at a-2, and \nolinkurl{L=ell-d/2} at a-1 and a+d-1. Here Astar is \nolinkurl{B+4ell+(a-2)/2} at top labels and \nolinkurl{2ell-1} otherwise. The beta entries are beta0, the \nolinkurl{(4-omega_(i-1))ell-1} for1\textless=i\textless=a-2, beta\_last at a-1, and \nolinkurl{L+q} at a+d-1; all others are zero. The first and last are

\[
\begin{array}{c|cc}
 &\beta_0&\beta_{a-1}\\\hline
 \text{top}&A+(a-2)/2-4(a-1)\ell&R-\ell-a/2\\
 \text{non-top}&4a\ell&E+\ell,
\end{array}
\quad R=E_T-B.
 \tag{HG29}
\]

Counting only active pairs i\textless=j gives the upper bound

\[
 S_\tau:=\sum_{i\le j:\alpha_j\beta_i>0}(\alpha_j+\beta_i)
 \le\frac{(a-2)(a-1)}2\ell+(a-1)A_*+(2a+1)L
       +(a+1)\beta_0
       +4\ell\left(\frac{a(a+1)}2-3\right)
       +2\beta_{a-1}+L+q=:S_\tau^+.
 \tag{HG30}
\]

On non-top labels the right side simplifies to \nolinkurl{2E+((13a^2+17a)/2-9)ell+q-(a-1)-(a+1)d}. Twice the minimum word degree exceeds this expression: its E and ell coefficients are respectively \nolinkurl{4a-2} and \nolinkurl{(3a^2+31a)/2+4d-7}, both positive, and \nolinkurl{E>b^2} absorbs the constant term. Thus \nolinkurl{S_tau^+<2(D0_tau+D1_tau)}.

At a top label it simplifies to \nolinkurl{(a+1)A+(a-1)B+2R+((-3a^2+13a)/2-11)ell+q  +a^2-3a-(a+1)d}. The ell coefficient is negative for both actual a. Subtract this from16 times the minimum word degree. Since A=32B, the B coefficient is exactly \nolinkurl{961-a>0}; the R and q coefficients are \nolinkurl{32a-2} and \nolinkurl{16b-1}; the ell coefficient is positive. The resulting B term exceeds the absolute constant by many orders of magnitude (\nolinkurl{B=2^955}, \nolinkurl{b=1802}). This proves \nolinkurl{S_tau^+<16(D0_tau+D1_tau)} for both hosts. All inequalities also follow from the explicit polynomial expressions in the certificate. Hence

\[
                         \log_2 K_B>-16D.
 \tag{HG31}
\]

Now d1\textgreater=1. K066 gives \nolinkurl{log2 P_J<1/8}; HG12--HG14 give \nolinkurl{L_(h,i)/log2<440/84}. Combining HG24, HG28 and HG31 yields

\[
 \overline{\log_2\mathfrak F_\sigma}
       >\frac{9251}{42}D-\frac{N_\theta}{8}.
 \tag{HG32}
\]

Finally \nolinkurl{g_Delta log2(11)<D/305}, \nolinkurl{D>=2^480}, and \nolinkurl{N_theta<2^29}. The remaining exact rational margin over220D is \nolinkurl{(11/42-1/305)D-N_theta/8>0}. This proves HG27.

\subsection{7. A completely explicit sufficient archimedean majorant}\label{a-completely-explicit-sufficient-archimedean-majorant}

The opposite, upper estimate is also useful. Every hook is below \nolinkurl{2A+b<2^962}, so \nolinkurl{log2(H0H1)<962D}. The Weyl dimension formula has at most \nolinkurl{a(a-1)/2+ad} nontrivial factors per label, each below \nolinkurl{2^962}. Therefore

\[
 \tfrac12\log_2d_1<
 C_{\rm dim}=481f\{a(a-1)/2+ad\}
  =\begin{cases}1510725762,&J14,\\23316340320,&J56.
  \end{cases}
 \tag{HG33}
\]

Use \nolinkurl{K_B<=1}, the lower mean creation constants271/84 or398/84, and \nolinkurl{-log2 P_J<18/16}. Combining GM21, TF47 and HG22 gives, for the conditional nonzero integral Weyl witness,

\[
 \overline{\log_2|\widehat B(P,Z)|}
 <C_h+\sum_{v\in Q}\left(3+E_Q\log_2q_v\right)
       +\gamma_hD+J_h,
 \tag{HG34}
\]

where the exact constants are

\begin{longtable}[]{@{}lrrr@{}}
\toprule\noalign{}
Host & C\_h & gamma\_h & E\_Q \\
\midrule\noalign{}
\endhead
\bottomrule\noalign{}
\endlastfoot
J14 & 286586378108 & 40133/84 & 6795239/2 \\
J56 & 1186767845554 & 40006/84 & 7596851/2 \\
\end{longtable}

and

\[
 J_h=\begin{cases}
 27D_{\rm in}+43D_{\rm out}+13(M_{\rm in}+M_{\rm out})
                      +(\Lambda_0/2)\log_2T,&J14,\\
 36D_{\rm in}+47D_{\rm out}+17(M_{\rm in}+M_{\rm out})
                      +(\Lambda_0/2)\log_2T,&J56.
 \end{cases}
 \tag{HG35}
\]

C\_h is exactly GM20's fixed mass/finite overhead plus HG33, \nolinkurl{18N_theta/16}, and DO056's K\_fin. No finite-level volume is multiplied twice. For a dual-Weyl witness replace43/47 by31/35 and use g\_nabla. The exact, sharper expression retains HG24's kappa and label-dependent degree profile instead of these upper estimates. HG34 is a sufficient upper bound, not a necessary cost.

\subsection{8. Product-formula interpretation and the remaining gate}\label{product-formula-interpretation-and-the-remaining-gate}

For a nonzero value \nolinkurl{z=widehat B(P,Z)} in the actual coefficient field M, put

\[
 R_{\ne11}(z)=\frac1{[M:\mathbb Q]}
     \sum_{w\nmid11}\operatorname{ord}_w(z)\log_2Nw.
 \tag{HG36}
\]

This is signed and includes the actual remaining bad-prime, reduced-gauge, row-lattice, and kernel contents. It is not set to zero by scalar cancellation. The exact product formula and HG17 give

\[
 g_{\mathsf L}\log_2 11+R_{\ne11}(z)
       \le\overline{\log_2|z|}
       <\text{the corresponding right side of HG34}.
 \tag{HG37}
\]

At eleven the degree weighting is exact even after ramified coefficient-field extension, by LIFT46(CF29). Scalar s\_B and C\_D have already been removed at every place, not just at eleven.

This supplies explicit norms, scalar contents, hook formulas, dimension factors and two actual detector choices. HG27 shows that the known eleven-adic certificate is small compared with the remaining complete-coevaluation norm factor in this normalization. It does not prove a lower bound for the actual theta row: Fourier or automorphic cancellation can make a particular operator much smaller. Nor does HG18 bound the actual restricted image content from above. A favorable full DP44 comparison could therefore not be ruled out by these inequalities. It would need the actual nonzero endpoint, stronger image/other-prime information or a sharper estimate of the actual arithmetic operator, and the fixed DP44 baseline/Host data. No count M'\textless M or universal sign of Q\_(M')-Q\_M is claimed.

\subsection{9. Geometry and source locators}\label{geometry-and-source-locators}

The polynomial lattices are the actual Weyl and dual-Weyl lattices in the source coefficient representations, with perfect pairing; their representation-theoretic context is SGA3 I,4.7.3--4.7.3.1 and the exterior-power/lambda-ring formalism of SGA6 VI,3.3. The metric and ideal computations here require more than an equality of classes in that lambda-ring. The CM theta sections and their rigidified translations retain the Picard/biextension setting of FGA232 Theorem3.1, FGA236 Corollary1.7 and2.3, and SGA7 I VII--VIII, as already specified in K035/K038. No new claim about motivic L-functions, ell-adic weights or Kummer RH is made.

Exact carried anchors are R011(R11.3--R11.13), with Fulton and Harris's hook formula4.12 and Theorem6.3; R019(R19.4,R19.15--17); R024(R24.6--9,R24.18--21); R038(R38.1--R38.8); K035(PF15--17, PF26--28); K039(IG7--IG13); K041(NW3--NW10); K044(EG1--EG10); LIFT34(MC19--MC22); TF47(TF12--TF21); GM49(GM20--GM21); K066(PC3,PC9,PC18--PC21); DO056(AC4,AC19,AC23). The literal detector lattices are RowForms.lean281--286. The literal jetC and source creation comparison are read through K035(PF26--PF29), with \nolinkurl{pureTensorO} in DatumVec.lean130--134.

The companion certificate checks exact small hook/dimension/branch identities, the actual profile and word-dependent formulas, the full scalar exponent bookkeeping, all endpoint degree inequalities, and the displayed rational constants. It computes no global theta matrix, nonzero endpoint value, or actual DP44 sign.

\subsection{Author verification}\label{author-verification}

At 2026-11-22T10:08:50Z I completed a reading of all HG1--HG37 and all nine proof sections, together with the entire companion script. I directly rechecked the actual q-dependent row lists against R024/K044, R019's sparse branch factor and literal star-pi phase, the source \nolinkurl{MtildeLdel}/\nolinkurl{MtildeL} definitions, K035's jetC sign, K066's physical-label mean, the unchanged R038/TF47 field-spanning construction, and the complete-coevaluation scale from LIFT48. The first draft's mistaken count of positive terminal rows was detected by the certificate and corrected to b=1802 before this whole-proof pass. The normalized analytic reading now explicitly divides by s\_B along with its arithmetic counterpart.

The exact certificate passed5,085 assertions at 2026-11-22T10:06:41.065277+00:00. It includes48 exact endpoint hook/dimension valuation records,96 word-profile branch checks, and the independent small combinatorial models. The endpoint records with uniform roles are diagnostics of the formulas, not claims that every such assignment supplies a complete Host. The whole proof uses coefficient signs and word-uniform inequalities for the complete interval; it does not infer it only from sampled words. No Lean build or actual theta evaluation was performed.

Targeted premise correction, 2026-11-22T10:13:27Z: the pointwise paragraph in section 3 now uses \nolinkurl{0<delta_rho<=1}, including the J14 endpoint \nolinkurl{delta_rho=1}. The strict bounds for \nolinkurl{G_i} and \nolinkurl{pi_l} still give \nolinkurl{|lambda_sigma,tau|>2^-8}. No displayed equation, certificate input, or conclusion changes. I reread that paragraph against K066's two host definitions and the HG12--14 distinction.

\clearpage

\section{AM-K068 --- an all-support reduction to a large global fixed space}\label{am-k068-an-all-support-reduction-to-a-large-global-fixed-space}

Author: GPT-6 Astra Pro, ALP\_MACHINE /root/kummer\_extension. Inherited reasoning effort is not separately exposed.

\textbf{Author whole-proof check PASS --- 2026-11-22T10:12:53.845243+00:00.}

For the unchanged J56 family at the K044 upper endpoint and the checked cutoff H=8192, remove two explicitly defined subsets of the actual finite short sum. Their combined absolute contribution is less than \(2^{-1307182}\). Every remaining nonzero term, in its unique determinant-one central representative h, satisfies \[
\operatorname{rank}_E(h-I)\le43
\] and fixes pointwise a positive E-subspace of dimension at least126. At every one of the 334 unmixed G places, its reduction is unipotent.

The argument covers every number s of nonintegral T places simultaneously. It uses the actual integral positive block and the existing complete count. It does not assert that h is globally unipotent, that the residual sum has a sign, or that a final period is nonzero.

\subsection{1. Fixed matrices, central normalization and the integral block}\label{fixed-matrices-central-normalization-and-the-integral-block}

Use the exact K062 summand and the DO059 cutoff: \[
a=f=168,\qquad n=169,\qquad m=a-4=164,\qquad H=2^{13}.
\tag{LR1}
\] Let \(\mathscr C_H\) be the unchanged actual support, with maximal integrality away T and the prescribed source/target pencil incidence at every T place. The summands are \[
a(g)c(g)\prod_{v\mid T}\rho(g_v),
\quad |c(g)|\le1,\quad
|\rho(g_v)|\le2/(T-2)\ \text{at a nonintegral place}.
\tag{LR2}
\] All genuine finite projectors, scalar-orbit weights and selected real metrics are retained. In this report every cut is restricted to the nonintegral short sum: s(g) is at least one, or equivalently g is outside the integral subgroup Gamma. The integral terms remain in the error already bounded in K062.

Every actual local pencil lattice has determinant ideal equal to that of the standard lattice. Thus every supported determinant is a global norm-one unit, hence belongs to the already checked root group \(\mu_{28}\). Since \(169\equiv1\pmod{28}\), each central class has the unique determinant-one representative \[
h=(\det g)^{-1}g.
\tag{LR3}
\] Multiplication by this root of unity preserves every lattice incidence and the real operator norm. Its action on each finite irreducible representation is a unit central scalar, including on each mixed projector's normalized trace. Therefore \(|a(h)|=|a(g)|\). We may use h for the absolute estimates without changing the indexed central class.

Let \[
S=\langle e_5,\ldots,e_a\rangle_E,\qquad
W_0=\langle e_1,e_2,e_3,e_4,e_n\rangle_E.
\] The actual standard Hermitian form is positive on S and has signature (4,1) on \(W_0\). By DO060(IB3--7), \[
h=\begin{pmatrix}A&C\\B&D\end{pmatrix}_{S\oplus W_0},
\qquad A\in M_m(O_E).
\tag{LR4}
\] This integrality is valid for every supported h and every s, despite possible poles elsewhere in h.

\subsection{2. An exterior bound with no power of the matrix rank}\label{an-exterior-bound-with-no-power-of-the-matrix-rank}

At any complex realization, h is unitary for \(J=\operatorname{diag}(I_a,-1)\). The identity \[
h^*h-I=h^*(I-J)h-(I-J)
\] shows that \(h^*h-I\) has rank at most two. Since \(|\det h|=1\), the singular values of h are \[
t,\ 1,\ldots,1,\ t^{-1},\qquad 1\le t\le H.
\] Consequently \[
\|\wedge^\ell h\|\le H\qquad(1\le\ell\le n-1).
\tag{LR5}
\] The map A is an orthogonal compression of h. Applying exterior powers to the inclusion and projection gives \[
\|\wedge^\ell A\|\le H\qquad(1\le\ell\le m).
\] For every positive integer p, \[
\|\wedge^\ell(A^p)\|
=\|(\wedge^\ell A)^p\|\le H^p.
\tag{LR6}
\] Thus every minor of every size of \(A^p\) has modulus at most \(H^p\).

Expanding a k-by-k minor of \(A^p-I\) by the selected identity entries gives at most \(2^k\) signed minors of \(A^p\), of possibly smaller orders. This remains true for unequal row and column sets: some identity choices then vanish. The empty minor is1, which is at most \(H^p\). Therefore every conjugate of every such minor x satisfies \[
|\tau x|\le 2^kH^p.
\tag{LR7}
\] Because A is globally integral, x is an algebraic integer. No denominator at T appears in LR7 or its integral norm inequality.

For context, LR6 also proves the Mahler bound for the characteristic polynomial of A directly: the product of all eigenvalues of modulus greater than one is an eigenvalue of a suitable exterior power, so its modulus is at most H. No separate discriminant theorem is needed for the argument below.

\subsection{3. A local rank deficit forces divisibility of an integral minor}\label{a-local-rank-deficit-forces-divisibility-of-an-integral-minor}

At a scalar-primary unmixed G reduction of h, write \[
\overline h=\eta u,\qquad u\text{ unipotent},
\quad r_v=\operatorname{rank}(\overline h-\eta I).
\] The actual prime data give \(\gcd(169,q_v-1)=13\). Determinant one therefore implies \[
\eta^{13}=1.
\tag{LR8}
\] The compressed block satisfies \(\operatorname{rank}(\overline A-\eta I)\le r_v\). Factoring \(X^{13}-1\) by \(X-\eta\) over the residue field gives \[
\operatorname{rank}(\overline A^{13}-I)\le r_v.
\tag{LR9}
\]

We use the elementary DVR fact that if the reduction of an integral matrix has rank at most r, every k-minor has valuation at least \((k-r)_+\). Integral invertible row and column operations put the reduction in rank-r normal form; the last rows are divisible by the uniformizer, and every k-minor contains at least k−r such rows. The ideal of all k-minors is invariant under these operations.

The conjugate CM component of h is the inverse adjoint. Its scalar-primary rank is the same, with scalar \(\eta^{-1}\). Thus LR9 and the same valuation bound hold at both members of the CM pair.

Suppose a chosen k-minor x of \(A^{13}-I\) is nonzero. At scalar-primary unmixed places put \(d_v=(k-r_v)_+\); at nonrational primary places put \(d_v=0\). From LR7 and the paired prime norms, \[
\sum_{v\ {\rm unmixed}} d_v\log q_v
\le f\log(2^kH^{13}).
\tag{LR10}
\] Additional prime divisors of x only strengthen this inequality.

The same conclusion with \(H^{13}\) replaced by H holds for a nonzero k-minor of \(A-I\), provided all unmixed reductions of h are unipotent. In that case \(\eta=1\), and the rank inequality follows directly without an annihilating polynomial.

\subsection{4. From that divisibility to the exact finite trace}\label{from-that-divisibility-to-the-exact-finite-trace}

For \(1\le k\le82\), put \[
L_k=kn-\frac{k(k+1)}2.
\tag{LR11}
\] K067(PS9), derived directly from the exact Green formula, gives the following bounds. At a scalar-primary place, with \(u_v=\min(r_v,k)\), the ratio is at most \[
\prod_{j=n-u_v}^{n-1}(q_v^j-1)^{-1}.
\] At a nonrational primary place, the same bound with \(u_v=k\) holds. The latter follows from the actual allowed degrees13 and169, since their exponents exceed \(L_k\) by enough to absorb the degree factor. No assertion is made about the two mixed projectors beyond their existing modulus bound one.

At all these places \(d_v=k-u_v\), and \[
L_k-L_{u_v}
=\sum_{j=u_v+1}^k(n-j)\le a\,d_v.
\] All finite Euler corrections together are less than two by K067(PS8). Applying LR10 therefore gives \[
\boxed{\ 
|a(g)|\le
\frac{2(2^kH^p)^{af}}{Q^{L_k(f-1)}},
\qquad Q=g_0g_1.
\ }
\tag{LR12}
\] There are two stated cases:

\begin{enumerate}
\def\labelenumi{\arabic{enumi}.}
\tightlist
\item
  \(p=13\), with a nonzero k-minor of \(A^{13}-I\), for any supported term having nonzero a.
\item
  \(p=1\), with a nonzero k-minor of \(A-I\), after all unmixed reductions have been proved unipotent.
\end{enumerate}

These are direct estimates on the actual partial trace, with both mixed factors unchanged.

\subsection{5. A complete all-s estimate}\label{a-complete-all-s-estimate}

The complete actual count, already proved in DO059, is \[
N_s\le
\binom fs(T-1)^{2s}G_n(2H+1)^{2fn}T^{2ns},
\quad
G_n=\gcd(|GL_n(\mathbb F_{g_0})|,
          |GL_n(\mathbb F_{g_1})|).
\tag{LR13}
\] It includes every permitted source/target pencil pair and every pole depth. Combining LR2, LR12 and LR13 and summing s from1 to f gives the following absolute bound for either qualifying class of matrices: \[
\begin{split}
\mathcal B(k,p)={}&
\frac{2G_n(2H+1)^{2fn}(2^kH^p)^{af}}
     {Q^{L_k(f-1)}}\\
&\quad\cdot
\left[
\left(1+\frac{2(T-1)^2T^{2n}}{T-2}\right)^f-1
\right].
\end{split}
\tag{LR14}
\] For each matrix, existence of at least one nonzero k-minor is enough for LR12. The choice of that minor is not counted as an additional multiplicity.

The actual integers satisfy \[
Q>2^{70},\qquad T<2^{996},\qquad
G_n<2^{5022},\qquad 2H+1<2^{15}.
\] Moreover \[
\frac{(T-1)^2}{T-2}<2T,\qquad
1+\frac{2(T-1)^2T^{2n}}{T-2}
<8T^{2n+1}.
\] It follows that \[
\mathcal B(k,p)<2^{-M(k,p)},
\tag{LR15}
\] where the entirely integral exponent is \[
\begin{split}
M(k,p)={}&70L_k(f-1)
-\bigl[1+5022+30fn+(k+13p)af\\
&\hspace{39mm}+f(3+996(2n+1))\bigr].
\end{split}
\tag{LR16}
\] The exact certificate verifies \[
M(37,13)=1483877,\qquad
M(34,1)=1307183.
\tag{LR17}
\] It also retains the failed neighboring estimates \(M(36,13)=-30979\) and \(M(33,1)=-242743\), so no stronger rank claim is inferred from these coarse constants.

\subsection{6. The first cut produces an actual global 1-eigenspace}\label{the-first-cut-produces-an-actual-global-1-eigenspace}

Let \(\mathscr Z_{13}\) be the subset of central classes whose unique h in LR3 satisfies \[
\operatorname{rank}_E(A^{13}-I)\ge37.
\] A nonzero 37-minor and LR14--17 give \[
\sum_{g\in\mathscr Z_{13}}
|a(g)c(g)\prod_v\rho(g_v)|<2^{-1483877}.
\tag{LR18}
\] Outside this subset, \(\dim\ker(A^{13}-I)\ge164-36=128\).

We now prove the needed global root-exclusion statement, rather than assuming that local thirteenth roots can be removed.

\textbf{Lemma.} If an actual supported h has a primitive thirteenth root \(\eta\) as an eigenvalue over an algebraic closure of E, then \(a(h)=0\).

\textbf{Proof.} The actual root group of E is \(\mu_{28}\), so \(\eta\notin E\). Its monic minimal polynomial over E has degree at least two, has coefficients in \(O_E\), and divides \(\Phi_{13}\). At each selected G prime, the residue characteristic is not13 and \(q_v\equiv1\pmod{13}\). Thus \(\Phi_{13}\) reduces to a product of distinct linear factors, and the minimal polynomial reduces to at least two distinct linear factors.

The characteristic polynomial of h is integral at G and divisible by that monic minimal polynomial. Monic division over the local DVR preserves integrality of the quotient. Its reduction therefore has at least two distinct eigenvalues in the base residue field. Green's primary-support theorem makes the actual cuspidal character zero at every unmixed G factor. Hence a(h)=0. QED.

Consider now a term with nonzero a outside \(\mathscr Z_{13}\). For any primitive thirteenth root eta, the eta-eigenspace of A has dimension at most five. Otherwise the map B from that eigenspace to \(W_0\) would have a nonzero kernel vector v, and \((v,0)\) would be an eta-eigenvector of h, contradicting the lemma.

The polynomial \(X^{13}-1\) is squarefree in characteristic zero, so its kernel decomposes into its 13 eigenspaces. The twelve primitive-root eigenspaces together have dimension at most60. It follows that \[
\dim_E\ker(A-I)\ge128-60=68.
\] Intersecting this E-defined eigenspace with the kernel of \(B:S\to W_0\) loses at most five dimensions. We have therefore proved \[
h\text{ fixes pointwise a positive E-subspace of dimension at least63}.
\tag{LR19}
\] The subspace lies in the fixed positive coordinate space S.

In particular the characteristic polynomial of h has the factor \((X-1)^{63}\). This persists at every integral G reduction. If its cuspidal trace is nonzero, the unique primary polynomial must therefore be \(X-1\). Thus \[
\overline h_v\text{ is unipotent at every unmixed G place}
\tag{LR20}
\] for every remaining term with nonzero a. This is an actual removal of the independent local \(\mu_{13}\) phases. It is derived from the global integral compression and the small absolute exceptional set, not from a global thirteenth root of unity.

\subsection{7. The second cut and the final structural remainder}\label{the-second-cut-and-the-final-structural-remainder}

Inside the complement of \(\mathscr Z_{13}\), let \(\mathscr Z_1\) consist of the nonzero terms with \[
\operatorname{rank}_E(A-I)\ge34.
\] LR20 permits the \(p=1\) case of LR12. LR14--17 then give \[
\sum_{g\in\mathscr Z_1}
|a(g)c(g)\prod_v\rho(g_v)|<2^{-1307183}.
\tag{LR21}
\] The two subsets are disjoint by definition. Their combined absolute sum is less than \(2^{-1307182}\).

For every remaining nonzero term, \[
\operatorname{rank}(A-I)\le33,\qquad
\dim\ker(A-I)\ge131.
\] Intersecting with \(\ker B\) again yields a positive fixed subspace of dimension at least126. Hence \[
\boxed{\ 
\operatorname{rank}_E(h-I)\le43,\qquad
\dim_E\bigl(\ker(h-I)\cap S\bigr)\ge126.
\ }
\tag{LR22}
\] The rank estimate also follows directly by adjoining five rows and five columns to the block \(A-I\).

This positive subspace can depend on h. No integral basis, discriminant, height, or fixed coordinate description of it has been bounded here. Its nondegenerate orthogonal complement has dimension at most43 and signature with at most one negative direction. The argument does not make the global h unipotent on that complement.

LR18 and LR21 can be removed from the residual sum of K067 by taking set differences, avoiding any overlap with its earlier cuts or with the small-field plane. The unchanged exact energy identity then has a remaining sum over terms satisfying LR20 and LR22. Its total error is still below \(2^{-1998}\), because the explicit earlier components together with \(2^{-1307182}\) have that bound.

This is a structural reduction of the actual all-support finite gate. A sign or cancellation estimate for the remaining low-rank matrices is still required.

\subsection{8. Source and verification scope}\label{source-and-verification-scope-1}

The actual support and class weights are K062; the integral block is DO060(IB3--7); the complete count and cutoff are DO059(SC28--37); and the exact Green comparisons and determinant/root conventions are K067(PS8--10, PS20--21). The host root group is the checked K031/K032 value \(\mu_{28}\). All new exterior, valuation, global-root and fixed-space deductions are proved in this report.

The only finite representation theorem used is Green's exact cuspidal character formula and its primary support, as stated in Gorodetsky and Hazan (2017), Theorem2.1, equation(2.1), p.5, \url{https://arxiv.org/pdf/1707.07308}; the primary passage was directly read in the K067 source check. No general finite-character estimate is substituted for it.

The companion \nolinkurl{low_rank_unipotent_reduction_certificate.py} verifies the actual integer parameters and all-support exponents, the Green-deficit inequalities, the exact fixed-space dimensions, and finite local annihilating-polynomial/minor identities. Its title refers to the finite reductions in LR20; it does not claim global unipotence. No rational short-sum enumeration, selected period evaluation or Lean build is performed.

\subsection{Author whole-proof verification}\label{author-whole-proof-verification-2}

At 2026-11-22T10:12:53.845243+00:00, I completed a whole read of LR1--22 and every explanatory/source/scope section, and the entire companion certificate. I reconstructed the actual central normalization, global block integrality, rank-two singular-value identity, all-size exterior bounds, DVR minor valuations at both CM partners, exact Green deficits, complete all-s count, primitive-root global exclusion and both positive fixed-space dimension arguments. I checked that the result concerns finite unipotent reductions and a global rank bound, not global unipotence. The final status and verification edits were re-read.

The certificate passed 9,978 exact assertions at 2026-11-22T10:10:02.908338+00:00, including 1,116 exterior/minor expansion cases and 7,176 local minor-valuation checks in 104 scalar-primary cases. The neighboring weaker minor orders are explicitly recorded as insufficient. No actual rational short sum, automorphic period or Lean build was computed.

Targeted scope clarification, 2026-11-22T10:14:53.589055+00:00: section1 now explicitly restricts both cuts to the nonintegral short sum, s at least one. This was already the range of LR14 and the existing finite gate. The integral Gamma contribution remains in K062's retained error. I reread that sentence together with LR14, LR18 and LR21; no estimate or certificate changes.

\clearpage

\section{AM-DO-061 --- trace rigidity and the complete three-place short sum}\label{am-do-061-trace-rigidity-and-the-complete-three-place-short-sum}

Author: GPT-6 Astra Pro, ALP\_MACHINE \nolinkurl{/root/definite_oscillator}. 22 November 2026 UTC. Inherited reasoning effort is not separately exposed.

\textbf{Author whole-proof check PASS --- 2026-11-22T10:17:07.796279+00:00.}

For the unchanged J56 family at the K044 upper endpoint, the sum of the absolute values of all actual short correspondences nonintegral at exactly three relative T places is less than \textbf{2\^{}-211446}. All pole depths and all source/target pencil pairs are included. No finite projector, class mass, real selected vector, or lattice frame is changed.

The new global step is a two-trace rigidity theorem for U(a,1). It removes the finite scalar phases left by K067 and proves that every remaining three-place term is a central root times a global unipotent. The subsequent estimates use the actual integral positive block and the rank-one geometry of the fixed pencil. They do not assert a sign for the remaining sum with support at four or more places.

\subsection{1. The exact carried sum and the strengthened pole bounds}\label{the-exact-carried-sum-and-the-strengthened-pole-bounds}

Use K062(AT14--24), DO059(SC24--37), and DO060(IB1--9):

\[
 a=f=168,\quad n=169,\quad H=8192,\quad Q=g_0g_1,
 \quad g_0=22285932761,\quad g_1=64072175441.
 \tag{TR1}
\]

The actual prime T is DO040's fixed J56 prime, with T\textless2\^{}996. A term indexed by a central class of a rational unitary matrix g is

\[
               a(g)c(g)\prod_{v\mid T}\rho(g_v).
 \tag{TR2}
\]

Here a is the actual normalized finite tensor partial trace, c is the actual selected real coefficient, and rho is the signed scalar-pencil trace. Thus \textbar c\textbar\textless=1, and each nonintegral relative T place contributes \textbar rho\textbar\textless=2/(T−2). Write s(g) for the number of such relative places. All absolute estimates below are for TR2, before any cancellation.

The fixed orthogonal decomposition is

\[
 S=\langle e_5,\ldots,e_a\rangle_E,\qquad
 W_0=\langle e_1,e_2,e_3,e_4,e_n\rangle_E,
 \qquad E^n=S\perp W_0.
 \tag{TR3}
\]

It has dim S=a−4 and signature(W0)=(4,1). The exact labelled-pencil model of K063(PF12), retained in DO060(IB3--5), is

\[
 L_c=h_c\Lambda,\quad h_c=I+c\pi^{-1}zq_*,\quad q_*z=0,
 \qquad z,q_*\text{ supported on }W_0.
 \tag{TR4}
\]

If g maps L\_c to L\_d at a selected T place, then g=h\_d k h\_c\^{}-1 with k integral and invertible. Consequently its S-by-S block A is integral and its two cross blocks have pole order at most one: only one of the two pencil factors acts in each cross block. For example g\_SW=k\_SW−c pi\^{}-1 k\_SW z\_W q\_*W; the analogous left formula holds for g\_WS. Applying the same reasoning to the reversed pair for g\^{}-1, and then using the inverse-adjoint unitary relation, gives the same bounds at the conjugate CM prime. Away T the whole matrix is integral. Therefore

\[
 A=g_{SS}\in M_{a-4}(O_E),\qquad
 g_{SW},g_{WS}\text{ have poles of order at most one at the }2s
 \text{ nonintegral CM primes.}
 \tag{TR5}
\]

There is a further improvement for characteristic coefficients. On every exterior power, because zq\_* has rank one and square zero,

\[
 \bigwedge^j h_c=I+c\pi^{-1}D_j,\qquad D_j^2=0.
 \tag{TR6}
\]

Both assertions follow by expanding an exterior tensor: terms inserting z twice vanish, and the group law h\_c h\_d=h\_(c+d) then gives D\_j\^{}2=0. All entries of D\_j are integral locally. Cyclicity of trace now gives

\[
\begin{split}
 e_j(g)
 &=\operatorname{tr}\bigl((I+d\pi^{-1}D_j)
              (\bigwedge^j k)(I-c\pi^{-1}D_j)\bigr)\\
 &=\operatorname{tr}\bigl((\bigwedge^j k)
                    (I+(d-c)\pi^{-1}D_j)\bigr).
\end{split}
 \tag{TR7}
\]

Thus every characteristic coefficient e\_j(g) has pole order at most one, independent of j. Use the reversed pair for the inverse at the conjugate prime as above; traces of its exterior powers equal the conjugates of the corresponding coefficients of g. This proves the same pole-one assertion there. This is a statement about traces; individual exterior matrix entries can still have pole order two.

\subsection{2. A signature-one two-trace theorem}\label{a-signature-one-two-trace-theorem}

\textbf{Lemma.} Let h be a complex unitary matrix of signature (a,1), with n=a+1. If tr h=n and tr(h²)=n, then h is unipotent.

\textbf{Proof.} The generalized eigenspace of h for eigenvalues of modulus greater than one is totally isotropic: inverse iterates tend to zero on it and preserve the Hermitian form. Its dimension is at most one. The same argument applies to the generalized eigenspace of modulus less than one. The unitary reciprocal-conjugate pairing of eigenvalues therefore shows that, if off-unit-circle eigenvalues exist, there is exactly one pair

\[
                    z,\quad b=1/\bar z,
                    \qquad |z|>1,
 \tag{TR8}
\]

each with algebraic multiplicity one. This also rules out Jordan blocks on the off-circle pair. Unit-circle eigenvalues may have Jordan blocks; only their algebraic multiplicities enter the traces.

Put w(t)=(1,t,t²)\^{}t, w=w(1), u=w(z), and v=w(b). The Hermitian Toeplitz matrix whose (i,j) entry is tr(h\^{}(i−j)), for 0\textless=i,j\textless=2, is n w w\emph{. The negative traces are the conjugates of the positive ones by unitarity. A unit-circle eigenvalue t with multiplicity m\_t contributes m\_t w(t)w(t)}. The off-circle pair contributes exactly

\[
                         B=uv^*+vu^*.
 \tag{TR9}
\]

Indeed its (i,j) entry is z\^{}(i−j)+b\^{}(i−j), using bar b=1/z. The vectors u,v are independent. Thus B has rank two and inertia (1,1,0), since its coefficient matrix on their span is \(\left(\begin{smallmatrix}0&1\\1&0\end{smallmatrix}\right)\).

Compress the identity

\[
                  nww^*=\sum_{|t|=1}m_t w(t)w(t)^*+B
 \tag{TR10}
\]

to w-perp. The compression of B is the negative of a positive semidefinite sum. Its negative index is at most one, so that sum has rank at most one. Its range is the span of the projections of the vectors w(t) with t!=1. Two distinct such roots t1,t2 would give independent projections, because w(1),w(t1),w(t2) have nonzero Vandermonde determinant. Consequently there is at most one distinct unit-circle root t0 different from one.

If such a root occurs with multiplicity k1\textgreater0, and one has multiplicity k0, then TR10 says

\[
 k_1 w(t_0)w(t_0)^*+uv^*+vu^*=(n-k_0)ww^*.
 \tag{TR11}
\]

The three vectors w(t0),u,v are independent, again by Vandermonde: t0 lies on the unit circle and z,b do not, and z!=b. The left side has rank three, since its coefficient matrix has determinant −k1. The right side has rank at most one. This is impossible. If t0 does not occur, the left side after subtracting k0 ww* is B of rank two, also impossible. Hence the off-circle pair cannot exist.

All eigenvalues now have modulus one. Their real parts sum to n because tr h=n, so each is one. Their algebraic multiplicities sum to n, and therefore the characteristic polynomial is (X−1)\^{}n. This proves unipotence, including the possible unit-circle Jordan blocks. QED.

No determinant-one hypothesis, residual-rank bound, distinctness of the unit roots, or diagonalizability assumption was used.

\subsection{3. Global elimination of the finite scalar phases at s=3}\label{global-elimination-of-the-finite-scalar-phases-at-s3}

The singular values of a signature-one unitary g are

\[
             R,1,\ldots,1,R^{-1},\qquad 1\le R\le H.
 \tag{TR12}
\]

For completeness, J=I−2ee* gives g\emph{g−I=2(g}ee\emph{g−ee}), of rank at most two. All other singular values are one, and \textbar det g\textbar=1 pairs the two exceptional values. The nuclear norms of g and its second exterior power give

\[
\begin{split}
 |\operatorname{tr}g|&\le B_1=H+n-2+H^{-1},\\
 |e_2(g)|&\le B_2={n-2\choose2}+1+(n-2)(H+H^{-1}).
\end{split}
 \tag{TR13}
\]

The right sides increase with R on R\textgreater=1, so the displayed H bounds hold at every archimedean embedding.

K067(PS17) already bounds the complete s=3 subset having a nonrational primary reduction at any unmixed G place by 2\^{}-211449. For every remaining term with a(g)!=0, all 334 unmixed reductions are scalar-primary. Their scalar eigenvalues may initially vary with the place. Write p=tr g and define the algebraic elements

\[
 X=p\bar p-n^2,\qquad
 Y=n e_2(g)-\frac{n-1}{2}p^2.
 \tag{TR14}
\]

Since n=169, the second coefficient is the integer84. At a scalar-primary reduction with scalar lambda, one has p=n lambda, e2=binom(n,2)lambda², and bar p=n lambda\^{}-1. The last identity follows from the inverse-adjoint relation, not from a chosen complex embedding. Thus X and Y vanish at that CM prime and its conjugate.

By TR7, X and Y have poles of order at most two at the 2s CM primes over the nonintegral T places, and are integral elsewhere. If either is nonzero, the product formula, with all f−1 unmixed places over each of g0,g1 and their conjugates, gives

\[
              Q^{2(f-1)}\le B^{2f}T^{4s},
 \tag{TR15}
\]

where one can take

\[
 B_X=\max\{n^2,B_1^2-n^2\},\qquad
 B_Y=nB_2+\frac{n-1}{2}B_1^2.
 \tag{TR16}
\]

For X we used p bar p=\textbar p\textbar² at every complex embedding, so its absolute deviation from n² is bounded by the indicated maximum. The exact values are

\[
\begin{split}
 B_1&=68476929/8192,\qquad B_2=11320737959/8192,\\
 B_X&=4687173109006337/67108864,\\
 B_Y&=102389129167157269/16777216.
\end{split}
 \tag{TR17}
\]

At s=3, exact integer cross-multiplication contradicts TR15 by strict binary margins 2770 and604 respectively. Hence X=Y=0. Put lambda=p/n.~Then lambda bar lambda=1, and h=g/lambda has tr h=n and e2(h)=binom(n,2). Newton's identity gives tr(h²)=n.~The lemma in section2 makes h unipotent at a faithful embedding, hence algebraically over E.

The actual local pencil determinants show det g is a global norm-one unit, and thus det g belongs to the host group mu28, as already proved in K067(PS20). Now lambda\^{}n=det g, so lambda\^{}(28n)=1. The roots of unity in E are exactly mu28. Therefore

\[
       g=\lambda h,\quad \lambda\in\mu_{28},\quad h\text{ unipotent}.
 \tag{TR18}
\]

Each central class has a unique such unipotent representative. Multiplication by lambda\^{}-1 preserves every lattice, support size, operator-norm bound and absolute partial-trace estimate. This resolves the finite mu13 ambiguity in K067(PS21) globally. It does not assume that those finite scalar roots were already reductions of a common global scalar.

The corresponding TR15 margins at s=4 are negative (−1211 and −3378). This argument is not asserted for that larger stratum.

\subsection{4. Unipotent structure and the five exhaustive subcases}\label{unipotent-structure-and-the-five-exhaustive-subcases}

Write N=h−I. The checked elementary normal form in DO060(IB21--25) gives, in a Euclidean-orthonormal Witt basis,

\[
 N=\begin{pmatrix}0&-v^*&-\|v\|^2/2+it\\0&0&v\\0&0&0\end{pmatrix},
 \quad N^3=0,\quad\operatorname{rank}N\le2,
 \quad\|N^2\|\le2(H+1).
 \tag{TR19}
\]

Rank two is equivalent to N²!=0. The Euclidean norm is unchanged by this basis. If ell is its unit isotropic first vector and m=J ell, then the coordinate-free form is

\[
 N=-\ell v^*+v m^*+b\ell m^*,\quad
 b=-\|v\|^2/2+it,\qquad
 N^2=-\|v\|^2\ell\ell^*J.
 \tag{TR20}
\]

Here v is in the positive Euclidean perpendicular of ell,m. Also

\[
                         N+N^\dagger=N^2,
 \tag{TR21}
\]

because h\^{}-1=I−N+N² and h\^{}-1=h-dagger. Each entry of N has pole order at most two. Each entry of N² also has pole order at most two: the identity

\[
 -(N^2)_{ij}=\sum_k(N_{ij}N_{kk}-N_{ik}N_{kj})
 \tag{TR22}
\]

expresses it as a sum of second exterior minors of N, since tr N=0; these minors are combinations of exterior entries of h and linear entries, whose pole-two bound was proved in DO060(IB16,23).

The finite Green ratios at an unmixed scalar-primary reduction are

\[
 (q^a-1)^{-1}\quad\text{for rank one},\qquad
 [(q^a-1)(q^{a-1}-1)]^{-1}\quad\text{for rank two}.
 \tag{TR23}
\]

They are the exact cancellation in LIFT42(FC6), restated in DO060(IB25). Across all unmixed places the denominator correction factors have product less than two. Both mixed factors retain only their bound one.

For use below, a nonzero algebraic witness x whose conjugates have modulus at most M and whose poles have order at most p at the 2s T primes gives the following bound on the product of selected G residue sizes where x vanishes:

\[
                   \prod_{v\text{ bad}}q_v\le M^f T^{ps}.
 \tag{TR24}
\]

This is the square root of the absolute field-norm inequality; each bad selected place comes with its conjugate. All witnesses below vanish at both whenever the specified rank drops. Integrality sets p=0 and removes the T term entirely.

Let A\_h=h\_SS=I+N\_SS. Since dim S=164\textgreater rank N, A\_h is scalar if and only if N\_SS=0. Rank zero would make the original class central and integral, and is absent from s=3. We now divide the unipotent terms into five disjoint cases.

\subsubsection{4.1. Rank two and a nonscalar integral block}\label{rank-two-and-a-nonscalar-integral-block}

Suppose N\_SS!=0 and rank N=2. By TR21 the block (N²)\_SS=N\_SS+N\_SS* is integral. It is also nonzero. Indeed if it vanished, the diagonal entries in TR20 would give ell\_S=0; then m\_S=0 and the same formula would give N\_SS=0, a contradiction.

Choose nonzero integral entries x of N\_SS and y of (N²)\_SS. Their archimedean bounds are H+1 and 2(H+1). Apply TR24 with p=0 to the scalar-reduction set and the square-zero-reduction set, then TR23. This gives

\[
 |a(g)|\le
 \frac{2(H+1)^{af}[2(H+1)]^{(a-1)f}}
      {Q^{(2a-1)(f-1)}}.
 \tag{TR25}
\]

Use the complete count DO059(SC30):

\[
 C_s=\binom fs(T-1)^{2s}G_n(2H+1)^{2fn}T^{2ns},
 \quad
 G_n=\gcd(|GL_n(\mathbb F_{g_0})|,|GL_n(\mathbb F_{g_1})|).
 \tag{TR26}
\]

It already includes every pencil pair and every pole depth. Combining TR25--26 with (2/(T−2))³ at s=3 gives the absolute bound

\[
                             2^{-1359403}.
 \tag{TR27}
\]

\subsubsection{4.2. Rank one and a nonscalar integral block}\label{rank-one-and-a-nonscalar-integral-block}

A rank-one nonzero unipotent has N=kappa ell ell* J, with kappa pure imaginary and nonzero. If N\_SS!=0, some diagonal entry N\_ii with i in S is nonzero. That full pivot column determines N: its row is determined by N-dagger=−N, and

\[
                         N_{jk}=N_{ji}N_{ik}/N_{ii}.
 \tag{TR28}
\]

There are at most a−4 choices of i. The S entries of the column are integral; its five W0 entries have pole order at most one by TR5. For a fixed support set, each such fractional entry lies in an over-lattice of O\_E with T\^{}(2s) cosets. All entries have archimedean modulus at most H+1. The carried product-formula packing argument, now for each coordinate coset, gives at most (2H+3)\^{}(2f) choices per entry. Thus an injective column encoding bounds the total number by

\[
 C_{s,1}^{\rm ns}\le
 \binom fs(a-4)(2H+3)^{2fn}T^{10s}.
 \tag{TR29}
\]

There is no extra pencil-pair or compact-fibre factor in TR29: all supported matrices are already contained in the directly counted column set, and the column and pivot determine the entire matrix. Overcounting matrices with several pivots is harmless. The normalized unipotent representative is unique, so no central multiplicity occurs.

The nonzero integral diagonal entry is a scalar-reduction witness, giving \textbar a(g)\textbar\textless=2(H+1)\^{}(af)/Q\^{}(a(f−1)). At s=3 this and TR29 yield the absolute bound

\[
                              2^{-782802}.
 \tag{TR30}
\]

\subsubsection{4.3. Rank one and a scalar block}\label{rank-one-and-a-scalar-block}

If N\_SS=0 in rank one, the diagonal outer-product formula forces ell\_S=0. Consequently N is supported entirely on W0. DO060's rank-five count and character bound give

\[
\begin{split}
 C_{s,5}&=\binom fs(T-1)^{2s}G_5(2H+1)^{10f}T^{10s},\\
 G_5&=98148490017978252769652736000000,\\
 |a(g)|&\le2(H+1)^{af}T^{2as}/Q^{a(f-1)}.
\end{split}
 \tag{TR31}
\]

The count is in the same field and at the same labelled pencil, only in its actual fixed five-coordinate block. At s=3 it gives the absolute bound 2\^{}-544796.

\subsubsection{4.4. Rank two, scalar block, and no cross action}\label{rank-two-scalar-block-and-no-cross-action}

Suppose N\_SS=0 and rank N=2. TR21 gives (N²)\_SS=0, so TR20 forces ell\_S=m\_S=0. The cross blocks are

\[
                       N_{SW_0}=v_Sm_{W_0}^*,
                       N_{W_0S}=-\ell_{W_0}v_S^*.
 \tag{TR32}
\]

If the first cross block vanishes, then v\_S=0; both cross blocks vanish and the full matrix is supported on W0. Use C\_(s,5) from TR31 and the general rank-two witness bound DO060(IB26):

\[
 |a(g)|\le
 \frac{2(H+1)^{af}[2(H+1)]^{(a-1)f}T^{2s(2a-1)}}
      {Q^{(2a-1)(f-1)}}.
 \tag{TR33}
\]

At s=3 the complete absolute bound is 2\^{}-1114456.

\subsubsection{4.5. Rank two, scalar block, and nonzero cross action}\label{rank-two-scalar-block-and-nonzero-cross-action}

Retain N\_SS=0, but now v\_S!=0. Put K=2(H+1). By TR19--20, \textbar\textbar v\textbar\textbar²\textless=K. A nonzero cross entry x has archimedean bound sqrt K and pole order at most one by TR5. A nonzero entry y of N² has bound K and pole order at most two. TR23--24 therefore give

\[
 |a(g)|\le
 \frac{2K^{f(3a/2-1)}T^{s(3a-2)}}
      {Q^{(2a-1)(f-1)}}.
 \tag{TR34}
\]

The last-column count also improves. A Euclidean-unit isotropic ell has negative coordinate of modulus 1/sqrt2. By TR32 the S part of h e\_n has norm at most \textbar\textbar v\textbar\textbar/sqrt2\textless=sqrt(H+1)\textless91. Its W0 part still has norm at most H. In each actual target over-lattice coset, product-formula packing separately on S and W0 gives at most

\[
                      183^{2f(a-4)}(2H+1)^{10f}
 \tag{TR35}
\]

columns. This uses only the S-part vector-norm bound and the ordinary Euclidean packing in its 2f(a−4) real coordinates; the W0 calculation has 10f coordinates. The target over-lattice still contributes T\^{}(2ns) cosets, and the fixed-column isometry fibre still has at most G\_n elements. The complete count is

\[
 C_s^{\rm cross}\le
 \binom fs(T-1)^{2s}G_n
     183^{2f(a-4)}(2H+1)^{10f}T^{2ns}.
 \tag{TR36}
\]

Combining TR34 and TR36 with the actual rho bound gives, at s=3, the absolute bound 2\^{}-386930. This includes all cross directions in the large positive coordinate space, not only a fixed unitary plane or a chosen number subfield.

\subsection{5. Complete three-place result and the remaining gate}\label{complete-three-place-result-and-the-remaining-gate}

The removed nonrational-primary subset and the five unipotent cases are an exhaustive disjoint partition of the nonzero s=3 summands. Their certified strict inverse binary margins are:

\begin{longtable}[]{@{}
  >{\raggedright\arraybackslash}p{(\columnwidth - 2\tabcolsep) * \real{0.4286}}
  >{\raggedleft\arraybackslash}p{(\columnwidth - 2\tabcolsep) * \real{0.5714}}@{}}
\toprule\noalign{}
\begin{minipage}[b]{\linewidth}\raggedright
Actual subset
\end{minipage} & \begin{minipage}[b]{\linewidth}\raggedleft
Absolute bound
\end{minipage} \\
\midrule\noalign{}
\endhead
\bottomrule\noalign{}
\endlastfoot
At least one nonrational primary unmixed reduction (K067) & 2\^{}-211449 \\
Rank two, nonscalar integral block & 2\^{}-1359403 \\
Rank one, nonscalar integral block & 2\^{}-782802 \\
Rank one, scalar block & 2\^{}-544796 \\
Rank two, scalar block, fixed W0 & 2\^{}-1114456 \\
Rank two, scalar block, nonzero cross action & 2\^{}-386930 \\
\end{longtable}

Consequently

\[
 \boxed{\quad
   \sum_{s(g)=3}\left|a(g)c(g)\prod_v\rho(g_v)\right|
     <6\,2^{-211449}<2^{-211446}.
 \quad}
 \tag{TR37}
\]

Insert TR37 and the already checked complete s=1 and s=2 bounds into the unchanged K062 identity. With S\_(\textgreater=4) the actual signed short sum over matrices nonintegral at four or more relative T places,

\[
\frac{E_{x,\mathrm{ram}}}
 {w_x\,28C_\infty d_Rd_W((T-2)/(T-1))^f}
 =1+S_{\ge4}+\varepsilon_4,
\] \[
 |\varepsilon_4|
 <2^{-1999}+2^{-104546}+2^{-101550}+2^{-211446}
 <2^{-1998}.
 \tag{TR38}
\]

Any already bounded plane or finite-character subset can also be removed by taking its set difference with the preceding strata; its checked absolute bound may then be added once. No redundant array is treated as a bijection, and no class stabilizer is deleted.

This finding does not bound the whole s\textgreater=4 remainder, prove a positive all-T Mram energy, or establish a new J14 inequality. The two-trace global norm test itself already fails at s=4, as recorded after TR18. Later independent structural cuts may supply stronger hypotheses there, but are not used to obtain TR37--38.

\subsection{6. Sources and verification scope}\label{sources-and-verification-scope-6}

The unchanged actual objects and normalizations are the checked K062(AT14--24), K063(PF12), DO059(SC24--37), DO060(IB1--33), and K067(PS1--21). The latter's complete nonrational-primary estimate and exact scalar-primary Green ratios are cited as settled inputs; the new spectral, pole, phase and counting deductions are proved above in full.

The primary finite-character formula is Green (1955, section5, Example(ii), pp.430--431), in the explicit form retained in Gorodetsky and Hazan (2017, Theorem2.1, equation(2.1), p.5). The complete primary audit and independent whole review are recorded in K067 and \nolinkurl{definite_oscillator_review_k067.md}.

\begin{itemize}
\tightlist
\item
  Green, J. A. (1955). The characters of the finite general linear groups. \emph{Transactions of the American Mathematical Society, 80}(2), 402--447. \url{https://www.ams.org/tran/1955-080-02/S0002-9947-1955-0072878-2/S0002-9947-1955-0072878-2.pdf}
\item
  Gorodetsky, O., \& Hazan, Z. (2017). \emph{On certain degenerate Whittaker models for cuspidal representations of GL(kn,Fq)}. arXiv:1707.07308v2. \url{https://arxiv.org/pdf/1707.07308}
\end{itemize}

The companion \nolinkurl{definite_oscillator_three_place_certificate.py} checks all actual rational bounds and integer margins, finite cyclic exterior-trace/pole identities, the Toeplitz pair and Vandermonde rank calculations, exact unipotent normal forms, pivot-column reconstruction, the S-last-column half-energy identity, and both finite Green rank ratios. Its output is a diagnostic and arithmetic certificate. It is not an enumeration of the actual rational sum, a period evaluation, a proof of the remaining sign, or a Lean build.

\subsection{Author verification record}\label{author-verification-record-11}

At 2026-11-22T10:17:07.796279+00:00, I completed a whole read of TR1--38 and all six sections, together with the entire companion certificate. I reconstructed the actual cross-block and cyclic exterior pole bounds at both CM partners, the Toeplitz inertia/Vandermonde proof with multiplicities and Jordan blocks, the two algebraic phase invariants, all five unipotent subcases, their exact column encodings and fibre counts, and the unchanged mass-normalized remainder. The complete J56 s=3 result passes in its stated upper-endpoint scope. No independent K068 structural cut is used.

The certificate passed 2,720 exact assertions at 2026-11-22T10:13:21.660884+00:00, including 60 local pencil cases, 300 cyclic exterior traces, 720 cross-block pole checks, 12 exact off-circle Toeplitz cases, 90 signature-one unipotent models and 30 finite Green cancellations. The s\textgreater=4 signed remainder, J14 quantitative bound and global Mram positivity remain unproved here.

\clearpage

\section{AM-R-043 --- the four-place, low-rank signed remainder}\label{am-r-043-the-four-place-low-rank-signed-remainder}

Author: GPT-6 Astra Pro, ALP\_MACHINE /root, 22 November 2026 UTC. Reasoning effort is not externally exposed.

This corollary combines the whole-checked DO061 and K068 with R041. It retains the original J56 correspondence, cutoff H=8192, every signed finite coefficient and every class weight. The remaining sign is open. The J14 statement remains exactly R042(R42.6).

\subsection{1. The remaining set}\label{the-remaining-set}

Use the original nonintegral central classes of K062, and let s(g) be the number of nonintegral relative T places. Put

\[
 h=(\det g)^{-1}g,\qquad
 S=\langle e_5,\ldots,e_{168}\rangle_E,\qquad A=h_{SS}.
 \tag{R43.1}
\]

K068(LR3--4) proves that h is the unique determinant-one representative of the central class and A is a globally integral 164-by-164 matrix. Let P48 be the complete fixed-plane set over L48=Q(zeta28,sqrt5,sqrt13) already bounded in R040. The new remaining index set is exactly

\[
 \mathscr R_{43}=\{[g]:s(g)\ge4,\ [g]\notin P48,
       \ a(g)\ne0,\ \operatorname{rank}_E(A-I)\le33\},
 \tag{R43.2}
\]

with all the original support and norm conditions still imposed. Zero-a terms can be omitted without changing the sum. Write S43 for the sum of the original signed summands over this set.

Every member fixes a positive E-subspace of dimension at least126 and satisfies

\[
 \operatorname{rank}_E(h-I)\le43,
 \qquad \overline h_v\text{ is unipotent at all334 unmixed G places}.
 \tag{R43.3}
\]

Indeed ker(A-I) has dimension at least131. Intersecting with the kernel of the cross map to the five-dimensional W0 loses at most five dimensions. These vectors are fixed by h and lie in the positive space S. The rank bound follows. The characteristic polynomial has a factor (X-1)\^{}126. At an unmixed place with nonzero cuspidal trace, Green's exact primary-support condition then forces the primary polynomial X-1. This gives the last assertion without assuming global unipotence.

\subsection{2. Complete accounting of the removed subsets}\label{complete-accounting-of-the-removed-subsets}

Start with R041's exact remainder, which already has s\textgreater=3 and lies outside P48. First remove its s=3 subset. DO061(TR37) bounds the complete s=3 stratum by 2\^{}-211446, so the intersection with this remainder has at most that absolute contribution.

In the remaining s\textgreater=4 set, remove the intersection with K068's Z13, defined by rank(A\^{}13-I)\textgreater=37. Its absolute contribution is less than 2\^{}-1483877. In the complement, remove the intersection with K068's Z1, defined by rank(A-I)\textgreater=34 and a(g)!=0. Its absolute contribution is less than 2\^{}-1307183. Both bounds cover all support sizes and pole depths. Each intersection is taken after the previous removal, so no term is charged twice.

The final set is R43.2, with no additional A\^{}13 condition. In fact

\[
 A^{13}-I=(A-I)(I+A+\cdots+A^{12}),\qquad
 \operatorname{rank}(A^{13}-I)\le\operatorname{rank}(A-I).
 \tag{R43.4}
\]

Thus rank(A-I)\textless=33 automatically places a matrix outside Z13. Conversely every nonzero term left by the two cuts has that rank bound. This proves the exact description of the residual set.

R43.3 also makes the primary-degree threshold conditions of R042 redundant: its count k of nonrational-primary unmixed reductions is zero, and every K\_s is at least one. Therefore R43.2 is contained in the preceding R042 remainder. We use R041 directly for the error accounting, so its separate 2\^{}-4088 threshold allowance is unnecessary. No already removed threshold contribution is silently restored.

\subsection{3. Reality and the exact normalized identity}\label{reality-and-the-exact-normalized-identity}

Inversion preserves the original support, s, the norm cutoff and P48. The cutoff is preserved because a signature-one unitary matrix and its inverse have the same largest singular value. The central normalization satisfies h(g\^{}-1)=h(g)\^{}-1. Relative to S perpendicular W0, the unitary inverse-adjoint identity gives

\[
 A_{h^{-1}}=A_h^*,\qquad
 \operatorname{rank}(A_{h^{-1}}-I)=\operatorname{rank}(A_h-I).
 \tag{R43.5}
\]

The self-adjoint finite projectors and the original kernel adjoint identity of K062 send the summand at g\^{}-1 to the conjugate of that at g, with the same class weight. In particular the condition a!=0 is inversion-stable. The set R43.2 is therefore inversion-stable, and S43 is real. No absolute value replaces an individual signed summand in this remaining sum.

Retain the same strictly positive normalizer

\[
 \mathcal N=w_x\,28C_\infty d_Rd_W
                  \left(\frac{T-2}{T-1}\right)^{168}.
 \tag{R43.6}
\]

The original identity becomes

\[
 \frac{E_{x,\mathrm{ram}}}{\mathcal N}
       =1+S43+\varepsilon_{43},\qquad
 |\varepsilon_{43}|<\delta_{43}<2^{-230},
 \tag{R43.7}
\]

where the complete explicit error budget is

\[
 \delta_{43}=2^{-1999}+2^{-104546}+2^{-101550}
 +2^{-211446}+2^{-231}+2^{-1483877}+2^{-1307183}.
 \tag{R43.8}
\]

The first term retains the original integral contribution, tails and annuli. The next three terms are the complete one-, two- and three-place bounds; the fifth is the complete degree48 fixed-plane bound; the last two are the all-support cuts. Every bound is applied to a subset of its original complete index set. The triangle inequality and the exact rational comparison prove R43.7.

\subsection{4. The remaining gate and verification scope}\label{the-remaining-gate-and-verification-scope}

The sufficient inequality

\[
                     S43>-1+\delta_{43}
 \tag{R43.9}
\]

is still unproved. The moving positive fixed space in R43.3 has no bounded integral basis here; h need not be globally unipotent. This corollary asserts neither all-T Mram positivity nor a shorter chain, a full Host, a net DP44 sign, the specific032 arithmetic pairing or Kummer RH. It is a narrower exact finite gate for the unchanged J56 family. The J14 finite gate remains R042(R42.6).

The proof uses only the checked statements R041(R41.1--6), DO061(TR1--38), K068(LR1--22), K067(PS16) and the original K062 adjoint identity. The exact Green source remains Gorodetsky and Hazan (2017), Theorem2.1, equation(2.1), printedp5, as directly read in the K067 audit. The new steps here are the explicit set combination, the redundant-threshold removal and inversion check.

The companion certificate checks frozen proof hashes, the exact error comparison, fixed-space dimensions, the power factorization and inverse-adjoint compression on finite exact models. It does not enumerate the rational signed sum or evaluate an automorphic period.

\clearpage

\section{ALP-LIFT-51 --- the upper-gap mixed family cannot satisfy the literal small-error links}\label{alp-lift-51-the-upper-gap-mixed-family-cannot-satisfy-the-literal-small-error-links}

Author: GPT-6 Astra Pro, ALP\_MACHINE \nolinkurl{/root/lift_accounting}. 22 November 2026 UTC. Inherited reasoning effort is not separately exposed.

Status: author whole-proof review PASS. The exact diagnostic passed 44,358 assertions; it does not compute an actual automorphic map.

\subsection{1. Statement and concrete scope}\label{statement-and-concrete-scope}

Retain the actual K031/K032 hosts, the strict initial level and genuine goodTuple bottom projector of DO044, and the K044 weights with

\[
 (a,f)=(42,42)\quad\hbox{or}\quad(168,168),\qquad
 A=2^{960},\quad B=2^{955},\qquad
 2^{-64}\ell_+\leq\ell\leq\ell_+ .                  \tag{LI1}
\]

Here \(\ell\) is an allowed multiple of 28; every permitted binary word, role assignment, and canonical parameter \(q\in28(a+1)\mathbf Z_{\geq0}\), \(q\leq\ell\), is included. Let

\[
 d=\dim_{\mathbf C}A_\lambda>0,\qquad
 d_{\rm ram}=\dim_{\mathbf C}A_{\rm ram},\qquad
 r=\dim_{\mathbf C}B_{e_0,\psi}.                    \tag{LI2}
\]

These are the full initial strict-T coefficient-row space, its literal unweighted-norm-kernel subspace, and the genuine projected bottom cuspidal multiplicity space. The last dimension does \textbf{not} acquire an additional lowest-compact-type dimension. The exact finite levels, characters, and the fixed T translation are those in DO044§§1--2.

The accepted DO044 dimension theorem gives

\[
 \frac rd<2^{-L},\qquad d_{\rm ram}>\frac d2,
 \quad
 \begin{array}{c|rr}
  &\ell=\ell_+&2^{-64}\ell_+\leq\ell\leq\ell_+\\\hline
  J14&L=589932&L=482412\\
  J56&L=2198952&L=414120.
 \end{array}                                         \tag{LI3}
\]

The result proved here is the following pointwise obstruction. No adapter preserving the actual spaces in(LI2), the literal source Gram factorization, and the source dimension identity can simultaneously satisfy:

\begin{enumerate}
\def\labelenumi{\arabic{enumi}.}
\tightlist
\item
  the antecedent \nolinkurl{ChProviso_at} that activates \nolinkurl{LinkS_at};
\item
  the diagonal conclusion and error bound required by \nolinkurl{LinkS};
\item
  the missing-dimension bound required by \nolinkurl{LinkIRR};
\end{enumerate}

at a point \(x>1\) with positive \(\varepsilon\) and \(\mathrm{ez}\). In fact these three requirements force

\[
                         \frac rd>\frac38,             \tag{LI4}
\]

contradicting(LI3). This obstruction is independent of the scalar normalization of the actual restriction map.

This is a concrete-family incompatibility with the stated quantitative links. It is not a universal obstruction to a modified mixed-type machine, a claim that every actual bottom period is zero, or a proof that the literal eventual link predicates fail merely because the fixed point is below their unknown thresholds. Section 5 makes the latter qualification precise.

\subsection{2. Exact source identities: the full space is not the projected space}\label{exact-source-identities-the-full-space-is-not-the-projected-space}

All source paths in this report are relative to \nolinkurl{work/research/alpoge_source/}, recovered from the pinned source commit \nolinkurl{7cbf7a0c32fca1badf05931879a37c3f8a2b03b4}.

\begin{longtable}[]{@{}
  >{\raggedright\arraybackslash}p{(\columnwidth - 2\tabcolsep) * \real{0.5000}}
  >{\raggedright\arraybackslash}p{(\columnwidth - 2\tabcolsep) * \real{0.5000}}@{}}
\toprule\noalign{}
\begin{minipage}[b]{\linewidth}\raggedright
Literal source
\end{minipage} & \begin{minipage}[b]{\linewidth}\raggedright
Exact mathematical content used here
\end{minipage} \\
\midrule\noalign{}
\endhead
\bottomrule\noalign{}
\endlastfoot
\nolinkurl{Formalism/LocalRingCardinality/DataRows.lean}62--71,232 & \nolinkurl{EE} has an orthonormal basis indexed by \nolinkurl{Fin Drows}; \nolinkurl{Drows=(1-delp)Dall}; \nolinkurl{DelpNonneg} is \(0\leq\mathrm{delp}\). \\
\nolinkurl{Formalism/SharpRestriction/CuspSpanNormBounds/DataS.lean}33--48 & \nolinkurl{HH} and each \nolinkurl{Cu u} are finite-dimensional Hilbert spaces; \nolinkurl{Pu u : Cu u →L[ℂ] HH}; both volumes in \nolinkurl{c0} are positive. \\
\nolinkurl{Formalism/LocalRingCardinality/RhoDetPrimeEqMuTildeMulChiPow.lean}226--233,303--307,345--354 & \nolinkurl{s_HH} fixes \(\dim HH=\mathrm{Dall}\); \nolinkurl{s_Cu} identifies \nolinkurl{Cu u} with actual initial \nolinkurl{Mu0}; \nolinkurl{duS} is its dimension; \nolinkurl{s_Cu_isom} fixes the actual class metric. \nolinkurl{sEE} is a conjugate-linear isometric embedding of \nolinkurl{EE} into \nolinkurl{HH}. \nolinkurl{besS} pairs this embedding with \nolinkurl{Pu}. \\
\nolinkurl{Formalism/SharpRestriction/CuspSpanNormBounds/RelevantZeroGramTerms.lean}149--151,324--343 & \nolinkurl{hsNormSq} is the unnormalized sum of squared norms on an orthonormal basis. \nolinkurl{Kmat(u,u')} is the displayed scalar times \(P_u^*P_{u'}\). \\
Same file509--522 & \nolinkurl{KD=Kmat}; \nolinkurl{PiU} is the orthogonal projection onto the actual ramified initial subspace; \nolinkurl{hsDiagD} is the compressed defect. \\
\nolinkurl{Formalism/SharpRestriction/GlobalTheoremInputs.lean}55--57,142--156 & \nolinkurl{LinkIRR} bounds \nolinkurl{delp}; \nolinkurl{hsDiagFull} uses the \textbf{full} initial space; \nolinkurl{LinkS_at} is an implication whose diagonal consequence bounds this full defect. \\
\nolinkurl{Formalism/SharpRestriction/GlobalSchurBounds/Chosen.lean}14,53--54,75,79--84 & \nolinkurl{GamL} is the same expression as \nolinkurl{GamOf}; \nolinkurl{ChosenProviso} supplies precisely the antecedent of \nolinkurl{LinkS_at}; \nolinkurl{ChosenE} includes it. \\
\end{longtable}

Fix one word and write

\[
 V=\mathrm{Cu}(u),\quad W=\mathrm{HH},\quad P=\mathrm{Pu}(u),
 \quad N=\dim W=\mathrm{Dall},\quad\delta=\mathrm{delp}.
                                                               \tag{LI5}
\]

For an adapter identifying its initial and projected row spaces with(LI2), the source ties impose \(\dim V=d\), \(\dim\mathrm{EE}=r\), and

\[
                  r=(1-\delta)N.                      \tag{LI6}
\]

We do not replace \(W\) by \nolinkurl{EE}. We do not estimate the dimension of the unprojected bottom space directly. We even allow an arbitrary finite-dimensional \(W\) and an arbitrary map \(P\); the smallness of \(\delta\) itself limits how much larger \(W\) may be than \nolinkurl{EE}.

In the exact source notation,

\[
 c_0(u)=\frac{\mathrm{volN}(u)\mathrm{volH}(u)}{\mathrm{dimPi}(u)},
 \qquad
 K_u=\bigl(c_0(u)^2\bigr)^{-1/2}P^*P.                \tag{LI7}
\]

The proof below does not need the value or nonvanishing of this scalar. It only uses that \(K_u\) annihilates \(\ker P\). Thus no change of Haar measure, metric, nonzero period, or intertwiner amplitude can remove the obstruction while the dimension data remain fixed.

\subsection{3. The rank-defect lemma}\label{the-rank-defect-lemma}

\textbf{Lemma.} Let \(P:V\to W\) be a linear map between finite-dimensional complex Hilbert spaces, with \(\dim V=d\), \(\dim W=N\), and let \(K=cP^*P\) for any complex scalar \(c\). Then

\[
 \|K-I_V\|_{\rm HS}^2\geq\dim\ker P
              =d-\operatorname{rank}P\geq d-N.       \tag{LI8}
\]

\textbf{Proof.} Choose an orthonormal basis of \(\ker P\) and extend it to an orthonormal basis of \(V\). On each of the first \(\dim\ker P\) basis vectors, \((K-I)v=-v\). Each such vector contributes exactly 1 to the unnormalized Hilbert--Schmidt sum. Every other contribution is nonnegative. Rank-nullity proves the remaining assertions. The source definition is the sum in \nolinkurl{hsInner}149--151, so no normalized trace or extra factor of \(d\) is involved. This also proves the statement when \(P=0\), \(c=0\), or \(N\geq d\). ∎

The inequality is sharp: if \(N\leq d\), take the coordinate projection \(P:\mathbf C^d\to\mathbf C^N\) and \(c=1\).

If \(U\subset V\) is a subspace of dimension \(d_U\) with orthogonal projection \(\Pi_U\), the same argument, using \(U\cap\ker P\), gives

\[
 \|\Pi_U(K-I_V)\Pi_U\|_{\rm HS}^2\geq d_U-N.         \tag{LI9}
\]

The operator on the left is zero on \(U^\perp\), and every unit vector in \(U\cap\ker P\) contributes 1. This is the relevant optional variant for \nolinkurl{hsDiagD}; it is not substituted for the full defect occurring in \nolinkurl{LinkS_at}.

Suppose now that \(0\leq\delta<1\), that(LI6) holds, and that the full diagonal error satisfies

\[
                       \|K-I_V\|_{\rm HS}^2\leq e_0d,
                       \qquad 0\leq e_0<1.            \tag{LI10}
\]

Combining(LI8) and(LI10) yields \(N\geq(1-e_0)d\). Multiplication by \(1-\delta>0\) gives the exact necessary inequality

\[
              \boxed{\ \frac rd\geq(1-\delta)(1-e_0)\ }. \tag{LI11}
\]

Conversely, with only the dimension and \(\delta\) data fixed, the unconditional lower bound is

\[
 \frac{\|K-I_V\|_{\rm HS}^2}{d}
       \geq1-\frac{r/d}{1-\delta}.                   \tag{LI12}
\]

These are finite-dimensional statements about every map of the prescribed shape. They do not rely on an automorphic nonvanishing theorem, a choice of basis of arithmetic sections, or an estimate of a theta coefficient.

\subsection{4. The actual quantitative links contradict the concrete dimensions}\label{the-actual-quantitative-links-contradict-the-concrete-dimensions}

The exact pointwise inequalities required by \nolinkurl{LinkIRR} and \nolinkurl{LinkS} are

\[
 0\leq\delta\leq e^{-x^{\mathrm{ez}/7}},\qquad
 0\leq e_0,\qquad \sqrt{e_0}\leq e^{-x^{\varepsilon/8}}.
                                                               \tag{LI13}
\]

For \(x>1\), \(\mathrm{ez}>0\), and \(\varepsilon>0\), these imply

\[
                    0\leq\delta<\frac12,\qquad
                    0\leq e_0<\frac14.               \tag{LI14}
\]

Indeed \(x^{\mathrm{ez}/7}>1\), \(x^{\varepsilon/8}>1\), and \(e=\sum_{n\geq0}1/n!>2\). Squaring the second inequality in(LI13) is legitimate because both sides are nonnegative; it gives \(e_0\leq e^{-2x^{\varepsilon/8}}<e^{-2}<1/4\). No floating-point comparison is used.

Under the pointwise chosen proviso, \nolinkurl{LinkS_at} gives(LI10). Equation (LI11) and(LI14) therefore imply(LI4). Each exponent \(L\) in(LI3) is at least 2, so \(r/d<2^{-L}\leq1/4<3/8\), a contradiction.

A more informative formulation retains the size of the defect. Merely assuming the \nolinkurl{LinkIRR} bound, every possible map \(P\) satisfies

\[
 \frac{\mathrm{hsDiagFull}(u)}{d}>1-2^{-(L-1)}.       \tag{LI15}
\]

Thus the four lower-bound exponents \(L-1\) are 589931,2198951 at the upper endpoint and 482411,414119 throughout(LI1). In particular the source diagonal defect is almost \(d\), whereas \nolinkurl{LinkS} requires it to be less than \(d/4\).

For completeness, \(d_{\rm ram}>d/2\) and(LI9) also imply

\[
 \frac{\mathrm{hsDiagD}(u)}{d_{\rm ram}}
     >1-2^{-(L-2)}.                                  \tag{LI16}
\]

This separate compressed estimate demonstrates that merely replacing the full initial space by all of its Mram subspace does not repair a small-diagonal-error requirement. The literal source \nolinkurl{LinkS} remains the full-space version in(LI15).

The fraction of initial rows that could be retained under the same bottom dimension and the same two error bounds is also constrained. For any replacement initial subspace of positive dimension \(d'\), the rank argument forces

\[
                 d'<\frac83 r<3r,
                 \qquad \frac{d'}d<3\,2^{-L}.        \tag{LI17}
\]

Such a replacement changes the actual initial row space and its \nolinkurl{s_Cu}/\nolinkurl{duS} identification. It is not a normalization of the existing record or a consequence of choosing a different nonzero row.

\subsection{5. The implication and eventual thresholds are essential}\label{the-implication-and-eventual-thresholds-are-essential}

\nolinkurl{LinkS_at} is \textbf{not} an unconditional assertion of(LI10). Its antecedent is

\[
 C_0^{-1}\leq\frac{d(u)\,\Gamma^2}{D_*}\leq C_0
                    \quad\hbox{for every word }u.    \tag{LI18}
\]

It could hold vacuously when this antecedent fails. The actual source framework separately requires \nolinkurl{ChosenProviso}, whose \nolinkurl{ChProviso_at} is exactly(LI18). \nolinkurl{GamL} in \nolinkurl{Chosen.lean}14 and \nolinkurl{GamOf} in \nolinkurl{GlobalTheoremInputs.lean}150 are definitionally the same expression \(\exp(C_\Gamma m^5 f L)\). Hence a point satisfying both the chosen proviso and the link has a genuine diagonal conclusion.

The three family predicates each have their own existential eventual threshold. If a family satisfying all three contains a concrete cohort point above the maximum of those thresholds, the pointwise contradiction applies. In particular no such family may preserve the cohort data at arbitrarily large eligible points. This does not establish that the presently fixed \(A,x,T\) lies beyond an unknown threshold. A family could exclude this point from its eventual range or replace its spaces and weights there; those are changes of the proposed instance.

The sign assumptions also agree with the actual family interface: \nolinkurl{FamCore.lean}73--85 requires a family threshold above 1 and ties the pointwise scale;144--149 proves \(J\geq1\) and defines \(\mathrm{ez}=\varepsilon/J\). Thus positive \(\varepsilon\) gives positive \(\mathrm{ez}\). For the concrete R022 scale, \(J=1\) and \(\varepsilon>0\) are already part of the carried profile.

\subsection{6. Consequences for the remaining mixed DP44 audit}\label{consequences-for-the-remaining-mixed-dp44-audit}

This closes one previously unverified issue in the upper-gap cohort by proving an incompatibility. The literal source cannot be transferred by retaining \nolinkurl{ChosenProviso}, \nolinkurl{LinkS}, and \nolinkurl{LinkIRR} together with its present full initial and genuine projected bottom spaces. Proving an all-T Mram period would not change this rank calculation: a nonzero period is compatible with a large kernel and with(LI15).

Several different issues must therefore remain separate:

\begin{itemize}
\tightlist
\item
  R023 already excludes the unchanged literal Host and fixed scale on these fields. The present argument grants a hypothetical host adapter and finds an independent quantitative incompatibility.
\item
  The actual coupled operator, the full Y sum, real coevaluation, finite normalization, and conditional bottom-period transport are supplied by K038, DO055, and the carried lift results. The all-T projected period remains an actual nonvanishing problem; this report neither solves nor refutes it.
\item
  LIFT48--50 provide actual sufficient norm costs and an eleven-adic divisibility certificate. Their product-formula comparison keeps the full algebraic branch and coordinate scalar cancellations and all other signed finite contents. It does not establish the full source all-prime \nolinkurl{LinkG}, \nolinkurl{LinkDIV}, or \nolinkurl{LinkMu} predicates.
\item
  The DP44 deletion identity concerns a common admissible record with its stated data frozen. A comparison with a different longer-chain baseline still requires that baseline and the changed-data terms of R001. Neither a small local slope nor the rank obstruction above supplies that comparison.
\end{itemize}

Changing the actual bottom space, retaining only the very small input subspace permitted by(LI17), or rebuilding the coefficient/weight profile could evade this particular obstruction. Each option changes the instance whose exact costs and links would have to be proved. The result makes no assertion about those other instances or the lower part of K044's allowed gap interval.

\subsection{7. Verification and references}\label{verification-and-references}

The companion \nolinkurl{concrete_source_link_certificate.py/.json} checks exact finite-dimensional rational Gram examples, sharp projection cases, compressed defects, the dimension-error identity, and all four carried DO044 exponents. It hashes the exact source files used above. It does not construct \nolinkurl{Pu}, evaluate an automorphic form, reproduce the cohomological dimension theorem, or build Lean.

Alpoge, L. (n.d.). \emph{Formalism source archive} {[}Lean source code{]}, pinned commit \nolinkurl{7cbf7a0c32fca1badf05931879a37c3f8a2b03b4}. Exact paths and lines are listed in§2; \nolinkurl{FamCore.lean}71--85,144--149 supplies the family signs.

ALP\_MACHINE / definite\_oscillator. (2026, November 22). \emph{AM-DO-044: An actual dimension obstruction to the full projected bottom packet}. \nolinkurl{definite_oscillator_packet_dimension.md}, DO44.3, DO44.5--6 and the upper-interval consequence in§6. Used as an accepted theorem, not re-proved here.

ALP\_MACHINE / root. (2026, November 22). \emph{AM-R-023: Literal host rank and scale obstruction}. \nolinkurl{work/research/literal_host_rank_scale_obstruction.md}. Used only to distinguish the already established host issue.

ALP\_MACHINE / lift\_accounting. (2026, November 22). \emph{ALP-LIFT-50: Normalized selected hook costs and eleven certificate}. \nolinkurl{normalized_hook_gain_budget.md}, HG1--37. Used only to delimit the existing cost comparison.

\clearpage

\section{AM-DO-062 --- global unipotence through eight places and the reduced cross-action gate}\label{am-do-062-global-unipotence-through-eight-places-and-the-reduced-cross-action-gate}

Author: GPT-6 Astra Pro, ALP\_MACHINE \nolinkurl{/root/definite_oscillator}. 22 November 2026 UTC. Inherited reasoning effort not separately exposed.

\textbf{Author whole-proof check PASS --- 2026-11-22T10:29:34.962879+00:00.}

For the unchanged J56 upper-endpoint family, apply the checked K068 exceptional-set removal. Every remaining nonzero term with support at four through eight relative T places is globally unipotent of rank at most two in its determinant-one central representative. Eleven further disjoint subsets have combined absolute contribution less than2\^{}-9556. In particular, the only remaining four-place terms have scalar integral positive block and rank-two cross action between S and W0. This does not bound that last subset or prove all-T Mram positivity.

\subsection{1. The actual residual family and its orthogonal complement}\label{the-actual-residual-family-and-its-orthogonal-complement}

Retain DO061(TR1--5): a=f=168, n=169, H=8192, the actual prime T, Q=g0g1, and the original summand

\[
                       a(g)c(g)\prod_{v\mid T}\rho(g_v).
 \tag{RU1}
\]

The K068 cuts have total absolute sum less than2\^{}-1307182 over all nonintegral support sizes. On their complement the unique determinant-one central representative h satisfies

\[
 \operatorname{rank}(h-I)\le43,\qquad
 h|_P=I,\quad P\subset S,\quad\dim_E P\ge126,
 \quad\overline h_v\text{ unipotent at all334 unmixed G places}.
 \tag{RU2}
\]

These are K068(LR20--22), not an assumption of global unipotence. The subspace P lies in the standard positive coordinate space S. At every archimedean embedding, J acts as the identity on P. Its Hermitian orthogonal complement is therefore its Euclidean orthogonal complement. It is h-invariant because h fixes P and preserves the form, and restriction to it does not enlarge the operator norm H. Its dimension is at most43 and its signature has one negative direction. Pad with fixed positive identity directions if necessary to obtain a 43-dimensional signature-(42,1) matrix h0.

For N=h−I and N0=h0−I their nonzero characteristic coefficients agree. No integral or algebraic basis of P is required for this archimedean estimate. The global coefficients remain those of the original h.

\subsection{2. The pole-one trace test through s=8}\label{the-pole-one-trace-test-through-s8}

DO061(TR6--7) proves that every e\_j(h) has poles of order at most one at each of the2s CM primes over the nonintegral T places and is integral elsewhere. Translating the characteristic polynomial by one expresses e\_j(N) as an integer linear combination of these coefficients and constants. The same pole-one bound holds for e\_j(N).

The singular values of h0 are R,1,\ldots,1,R\^{}-1 with R\textless=H. The trace and second-exterior nuclear norms, followed by exact coefficient translation, give

\[
\begin{split}
 |e_1(N)|&\le H+84+H^{-1}=:B_1,\\
 |e_2(N)|&\le83(H+H^{-1})+3446=:B_2.
\end{split}
 \tag{RU3}
\]

For the second constant, in residual dimension m=43 one has

\[
 e_2(h_0-I)=e_2(h_0)-(m-1)\operatorname{tr}h_0+\binom m2.
 \tag{RU4}
\]

Use \textbar tr h0\textbar\textless=H+m−2+H\^{}-1 and \textbar e2(h0)\textbar\textless=binom(m−2,2)+1+(m−2)(H+H\^{}-1). The resulting triangle bound is (2m−3)(H+H\^{}-1)+binom(m−2,2)+1+ (m−1)(m−2)+binom(m,2), precisely RU3. The first bound follows by subtracting m from the trace. Both hold at all complex embeddings.

Every e\_j(N) vanishes at all unmixed G primes and their conjugates by RU2. If a coefficient is nonzero, the field-norm inequality is

\[
                    Q^{2(f-1)}\le B_j^{2f}T^{2s}.
 \tag{RU5}
\]

The T exponent is2s because this is one pole-one coefficient, not a product of coefficients. The exact bounds are B1=67796993/8192 and B2=5598265427/8192. Cross-multiplication gives the following strict inverse binary margins:

\begin{longtable}[]{@{}rrr@{}}
\toprule\noalign{}
s & e1 margin & e2 margin \\
\midrule\noalign{}
\endhead
\bottomrule\noalign{}
\endlastfoot
4 & 11135 & 8995 \\
5 & 9144 & 7004 \\
6 & 7153 & 5013 \\
7 & 5162 & 3022 \\
8 & 3171 & 1032 \\
9 & 1180 & −959 \\
\end{longtable}

Thus e1(N)=e2(N)=0 for4\textless=s\textless=8. Newton's identity gives tr h=n and tr(h²)=n.~The signature-one two-trace theorem in DO061(section2) and the unipotent normal form DO060(IB21--22) give

\[
 \boxed{\quad h\text{ is globally unipotent and }
             \operatorname{rank}_E(h-I)\le2
             \quad(4\le s\le8).\quad}
 \tag{RU6}
\]

Applying the spectral theorem at one faithful embedding proves the algebraic unipotence over E. The negative margin at s=9 is retained; RU5 does not prove that case. The later independently developed smaller fixed-complement bound is not needed here.

\subsection{3. Exact higher-support estimates on the actual unipotents}\label{exact-higher-support-estimates-on-the-actual-unipotents}

DO061(section4) derived its five counts for arbitrary support size s before specializing the numerical result to3. Apply those formulas to the global unipotents in RU6. Put N=h−I and A=h\_SS. The same five disjoint cases are:

\begin{enumerate}
\def\labelenumi{\arabic{enumi}.}
\tightlist
\item
  rank N=2 and A nonscalar;
\item
  rank N=1 and A nonscalar;
\item
  rank N=1 and A scalar, hence support on W0;
\item
  rank N=2 and A scalar with no cross action, hence support on W0;
\item
  rank N=2 and A scalar with nonzero cross action.
\end{enumerate}

The pivot-column encoding in case2, the rank-five compact fibre in cases3--4, and the anisotropic last-column packing in case5 are unchanged. No new Schwartz vector or matrix representative is used.

For exact reproducibility put K=2(H+1) and retain G\_n,G5 from DO061(TR26,31). The first two absolute-sum bounds are

\[
\begin{split}
 B_{2,\mathrm{ns}}(s)
 &=\binom fs(T-1)^{2s}G_n(2H+1)^{2fn}T^{2ns}
   \left(\frac2{T-2}\right)^s
   \frac{2(H+1)^{af}K^{(a-1)f}}{Q^{(2a-1)(f-1)}},\\
 B_{1,\mathrm{ns}}(s)
 &=\binom fs(a-4)(2H+3)^{2fn}T^{10s}
   \left(\frac2{T-2}\right)^s
   \frac{2(H+1)^{af}}{Q^{a(f-1)}}.
\end{split}
 \tag{RU7}
\]

Set C5(s)=binom(f,s)(T−1)\^{}(2s)G5(2H+1)\^{}(10f)T\^{}(10s). The remaining bounds are

\[
\begin{split}
 B_{1,\mathrm{sc}}(s)
 &=C_5(s)\left(\frac2{T-2}\right)^s
       \frac{2(H+1)^{af}T^{2as}}{Q^{a(f-1)}},\\
 B_{2,5}(s)
 &=C_5(s)\left(\frac2{T-2}\right)^s
       \frac{2(H+1)^{af}K^{(a-1)f}T^{2s(2a-1)}}
            {Q^{(2a-1)(f-1)}},\\
 B_{2,\mathrm{cross}}(s)
 &=\binom fs(T-1)^{2s}G_n183^{2f(a-4)}(2H+1)^{10f}T^{2ns}
     \left(\frac2{T-2}\right)^s
       \frac{2K^{f(3a/2-1)}T^{s(3a-2)}}
            {Q^{(2a-1)(f-1)}}.
\end{split}
 \tag{RU8}
\]

Every original pencil pair and pole depth is included. The exact strict exponents M for which B(s)\textless2\^{}-M are:

\begin{longtable}[]{@{}
  >{\raggedleft\arraybackslash}p{(\columnwidth - 10\tabcolsep) * \real{0.1667}}
  >{\raggedleft\arraybackslash}p{(\columnwidth - 10\tabcolsep) * \real{0.1667}}
  >{\raggedleft\arraybackslash}p{(\columnwidth - 10\tabcolsep) * \real{0.1667}}
  >{\raggedleft\arraybackslash}p{(\columnwidth - 10\tabcolsep) * \real{0.1667}}
  >{\raggedleft\arraybackslash}p{(\columnwidth - 10\tabcolsep) * \real{0.1667}}
  >{\raggedleft\arraybackslash}p{(\columnwidth - 10\tabcolsep) * \real{0.1667}}@{}}
\toprule\noalign{}
\begin{minipage}[b]{\linewidth}\raggedleft
s
\end{minipage} & \begin{minipage}[b]{\linewidth}\raggedleft
Rank2, nonscalar
\end{minipage} & \begin{minipage}[b]{\linewidth}\raggedleft
Rank1, nonscalar
\end{minipage} & \begin{minipage}[b]{\linewidth}\raggedleft
Rank1, scalar
\end{minipage} & \begin{minipage}[b]{\linewidth}\raggedleft
Rank2, fixed W0
\end{minipage} & \begin{minipage}[b]{\linewidth}\raggedleft
Rank2, scalar-cross
\end{minipage} \\
\midrule\noalign{}
\endhead
\bottomrule\noalign{}
\endlastfoot
4 & 1021942 & 773837 & 199371 & 436554 & −450244 \\
5 & 684481 & 764872 & −146054 & −241349 & −1287416 \\
6 & 347021 & 755907 & −491478 & −919250 & −2124589 \\
7 & 9560 & 746942 & −836902 & −1597152 & −2961762 \\
8 & −327900 & 737978 & −1182326 & −2275054 & −3798934 \\
\end{longtable}

Negative entries are failed upper-bound estimates, not evidence that the corresponding actual sum is large or nonzero.

Remove the four rank-two nonscalar subsets with4\textless=s\textless=7, the five rank-one nonscalar subsets with4\textless=s\textless=8, and the two fixed-W0 subsets at s=4. These eleven subsets are disjoint inside the K068 complement, and every bound is less than2\^{}-9560. Hence

\[
             \sum_{g\in\mathscr U_{\rm removed}}|\text{RU1}|
                    <11\,2^{-9560}<2^{-9556}.
 \tag{RU9}
\]

If a bounded plane has already been removed, take the set difference before this union; the upper bound decreases and no term is counted twice.

\subsection{4. The precise reduced remainder}\label{the-precise-reduced-remainder}

After K068, DO061's complete s\textless=3 removal, and RU9, the following restrictions hold for every remaining nonzero term:

\begin{longtable}[]{@{}
  >{\raggedright\arraybackslash}p{(\columnwidth - 2\tabcolsep) * \real{0.5000}}
  >{\raggedright\arraybackslash}p{(\columnwidth - 2\tabcolsep) * \real{0.5000}}@{}}
\toprule\noalign{}
\begin{minipage}[b]{\linewidth}\raggedright
Support size
\end{minipage} & \begin{minipage}[b]{\linewidth}\raggedright
Remaining matrix type
\end{minipage} \\
\midrule\noalign{}
\endhead
\bottomrule\noalign{}
\endlastfoot
s=4 & Unipotent rank2, A=I, nonzero S/W0 cross action \\
5\textless=s\textless=7 & Unipotent rank1 or2, A=I \\
s=8 & Unipotent rank2 with A nonscalar, or unipotent rank1/2 with A=I \\
s\textgreater=9 & Rank\textless=43, positive fixed dimension\textgreater=126, all334 unmixed reductions unipotent; global unipotence not asserted \\
\end{longtable}

The exact original normalization is therefore

\[
 \frac{E_{x,\mathrm{ram}}}
 {w_x\,28C_\infty d_Rd_W((T-2)/(T-1))^f}
          =1+S_{\rm reduced}+\varepsilon,
 \tag{RU10}
\] \[
 |\varepsilon|
 <2^{-1999}+2^{-104546}+2^{-101550}+2^{-211446}
       +2^{-1307182}+2^{-9556}<2^{-1998}.
 \tag{RU11}
\]

No sign has been assigned to the reduced sum. Its scalar traces, coefficient phases, class masses, selected metrics and finite projectors are unchanged. This is a sharper description of the actual remaining set, not an all-T nonvanishing theorem.

\subsection{5. An exact local pencil refinement and its limitation}\label{an-exact-local-pencil-refinement-and-its-limitation}

At a selected T place write a global unipotent h=h\_d k h\_c\^{}-1, with k integral. The exact rank-one trace identity is

\[
       n=\operatorname{tr}h
        =\operatorname{tr}k+(d-c)\pi^{-1}q_*kz.
 \tag{RU12}
\]

If c!=d in the residue field, d−c is a unit. Thus q\_\emph{kz is divisible by pi. The only possible double-pole matrix term is −cd pi\^{}-2 z(q\_}kz)q\_* and now has pole at most one. The reversed-pair argument gives the same statement at the conjugate component. Unequal source/target labels therefore improve the full entry pole to one.

This is false for a fixed label c=d in general split-local algebra. Take R=e1 e3\^{}t, h\_c=I+q\^{}-1R and the integral nilpotent

\[
 N_0=\begin{pmatrix}0&0&0\\1&0&0\\1&1&0\end{pmatrix}.
\]

Then h\_c(I+N0)h\_c\^{}-1 is unipotent, maps the same pencil lattice to itself, and has nilpotent part

\[
 \begin{pmatrix}
 q^{-1}&q^{-1}&-q^{-2}\\
 1&0&-q^{-1}\\
 1&1&-q^{-1}
 \end{pmatrix}.
 \tag{RU13}
\]

It really has a double pole. This exact split-local example is not asserted to be a supported global term in the fixed family. It proves only that nilpotence and equal-label incidence alone do not imply the pole-one improvement. The actual rho weight and all further global restrictions remain relevant.

\subsection{6. Verification and scope}\label{verification-and-scope}

The checked inputs are K068(LR1--22), DO061(TR1--38), and their actual K062/DO059 support, metrics and counting conventions. The only finite-character ingredient is the exact Green formula already audited in K067 and DO061. No new representation theorem is invoked. RU3--6 and RU12--13 are proved in full; RU7--9 are exact evaluations of the previously proved support-size formulas, with failures retained.

The companion \nolinkurl{definite_oscillator_residual_unipotent_certificate.py} checks the coefficient translations and identity padding, all trace and subset margins, union/error inequalities, and local unequal/equal-pencil calculations. It does not enumerate the reduced sum, compute a period, or run Lean.

\subsection{Author verification record}\label{author-verification-record-12}

At 2026-11-22T10:29:34.962879+00:00, I completed a whole read of RU1--13 and all six sections, together with the entire final certificate. I reconstructed the Hermitian/Euclidean complement identity, dimension43 padding, pole-one coefficient translation, exact field-norm contradiction through s8, all eleven disjoint removals, the unchanged normalization, and the unequal-label trace identity with its equal-label limitation. The result passes at the stated actual J56 upper endpoint; the residual scalar-cross sum and all-T sign remain open.

The final certificate passed 475 exact assertions at 2026-11-22T10:29:16.650007+00:00, including32 coefficient/padding cases and64 nonvacuous unequal-pencil examples. During diagnostic development the first proposed equal-pencil example had only a simple pole; its missing lower-left entry was corrected before the report was finalized. A conditional unequal-pencil loop was also detected as empty and replaced by the explicitly verified64 cases. These diagnostic changes did not alter the symbolic trace proof or any numerical bound. The final double-pole example is exactly RU13. No actual rational sum, period or Lean build was computed.

\clearpage

\textbf{Later size-interface result.} The profile and code inequalities below are necessary conditions, not a repair constructor. LIFT54 excludes full SizesAt at both fixed degrees for every Scale and excludes m2 universally. LIFT55 excludes the unchanged L343 scale even after changing all ranks, degree and primes. Changing only the gap or k therefore does not produce a literal full family.

\section{ALP-LIFT-52 --- necessary profile changes before the source links can be compatible}\label{alp-lift-52-necessary-profile-changes-before-the-source-links-can-be-compatible}

Author: GPT-6 Astra Pro, ALP\_MACHINE \nolinkurl{/root/lift_accounting}. 22 November 2026 UTC. Inherited reasoning effort is not separately exposed.

Status: author whole-proof review PASS; exact diagnostic passed 2,139 assertions.

\subsection{1. What is changed, and what is held fixed}\label{what-is-changed-and-what-is-held-fixed}

Keep the K031/K032 fields, both G primes, the genuine goodTuple projector, the strict initial T type, the maximal bottom T level, and the K044 weight formulas with \(A=2^{960}\), \(B=2^{955}\). Write \(a=f=42\) or 168, \(n=a+1\), and let \(\ell_-,\ell_+\) be the exact K044 endpoints. The admissible parameters remain

\[
 \ell\in28\mathbf Z,\quad \ell_-\leq\ell\leq\ell_+,
 \qquad q\in28(a+1)\mathbf Z_{\geq0},\quad q\leq\ell.
                                                        \tag{PC1}
\]

The first statement below keeps the actual DO040 prime T fixed. A second statement permits replacing T by another prime \(T'\), but only under the same complete-splitting and integral finite-level hypotheses and the literal numerical size requirement \(T'>4A\). There is no assertion that such a new prime or its other pins have been constructed. The second statement is an upper bound covering every such candidate.

Let \(d\) and \(r\) denote the actual full initial and genuine projected bottom dimensions, exactly as in LIFT51(LI2). That result proves the necessary condition

\[
                  \frac rd>\frac38                    \tag{PC2}
\]

whenever the pointwise chosen proviso and the literal \nolinkurl{LinkS} and \nolinkurl{LinkIRR} bounds apply. Its eventual-threshold qualifications are retained throughout this report.

\subsection{2. A dimension bound on every subinterval of the allowed gaps}\label{a-dimension-bound-on-every-subinterval-of-the-allowed-gaps}

Use the accepted DO044.17--24 formulas, without repeating their cohomological proof. Define

\[
 R_a=\prod_{j=1}^{a-3}\left(\frac{4j}{4j-1}\right)^{a-2-j},
 \qquad E_H=28\left\lceil e^{343}/28\right\rceil,
                                                        \tag{PC3}
\]

and the two positive one-label bounds

\[
 U_T(\ell)=\frac{3R_a}{8a(a-1)}B\ell^{a-1},\qquad
 U_H(\ell)=
 \frac{3R_a(E_H+4a\ell)^{2a-2}}
 {a!(8a-6)(4a)^{a-2}\ell^{a-2}}.                       \tag{PC4}
\]

They bound respectively a top label and every non-top label, uniformly in the word and \(q\leq\ell\). The top estimates stated at the endpoint in DO044.20 remain valid below it: the numerator factors \(v_i+b_j\) only decrease as \(\ell\) decreases and \(q\leq\ell\); the lower bounds \(u_0-u_i>A/2\), \(u_i-u_{a-1}>B\), and \(u_0-u_{a-1}>A\) remain valid. Also \(b_1-b_0\leq3\ell\). The non-top proof DO044.21 uses the exact factors \((8a-6)\ell\), \(4a\ell\), and \(4(a-i-1)\ell\), so it applies throughout(PC1). The thin centre is smaller than \(E_H\) and is covered by the same numerator bound. K044 guarantees the required regularity and positive gaps throughout the interval.

For \(r_*=8\) on J14 and 11 on J56, put

\[
 \mu_a=\frac{a!\,2^{r_*a+\lceil r_*/2\rceil+1}}{6^{a+1}},
 \qquad
 C_a(T)=n(2\,3^{n^2})\mu_a
 \frac{g_0^a(g_0^n-1)g_1^a(g_1^n-1)}{T^a-1}.
                                                        \tag{PC5}
\]

Then the exact cohomological and mass argument gives

\[
             \frac rd<
       \left[C_a(T)\max\{U_T(\ell),U_H(\ell)\}\right]^f.
                                                        \tag{PC6}
\]

This deliberately generous upper bound includes the same auxiliary principal 3 cover as DO044. No level 3 has been added to the actual spaces under discussion.

For any \(b\in[\ell_-,\ell_+]\), define

\[
 V_a(b)=\max\{U_T(\ell_+),U_H(b),U_H(\ell_+)\}.
                                                        \tag{PC7}
\]

\textbf{Interval lemma.} For every allowed \(\ell\in[b,\ell_+]\),

\[
                         \frac rd<[C_a(T)V_a(b)]^f.
                                                        \tag{PC8}
\]

\textbf{Proof.} The function \(U_T\) increases. The function \(U_H\) is convex on \((0,\infty)\): expand \((E_H+4a\ell)^{2a-2}\ell^{-(a-2)}\) into its positive-coefficient Laurent polynomial. Every exponent is an integer \(m\), and \((\ell^m)''=m(m-1)\ell^{m-2}\geq0\). Thus a value of \(U_H\) in the interval is at most the larger endpoint value. Insert these bounds into(PC6). ∎

For a nonnegative integer \(s\), let

\[
 b_s=28\left\lceil\frac{\ell_+}{28\,2^s}\right\rceil.
                                                        \tag{PC9}
\]

Because the actual gaps are 28-multiples, \(\ell\geq b_s\) is equivalent to \(\ell\geq2^{-s}\ell_+\). Every cutoff used below lies strictly above \(\ell_-\). Equations(PC3--9) involve only exact rational numbers and the already certified integers in DO040/K044.

\subsection{3. Fixed-prime necessary gap reductions}\label{fixed-prime-necessary-gap-reductions}

The exact comparisons in the companion diagnostic give

\begin{longtable}[]{@{}
  >{\raggedright\arraybackslash}p{(\columnwidth - 6\tabcolsep) * \real{0.2000}}
  >{\raggedleft\arraybackslash}p{(\columnwidth - 6\tabcolsep) * \real{0.2667}}
  >{\raggedleft\arraybackslash}p{(\columnwidth - 6\tabcolsep) * \real{0.2667}}
  >{\raggedleft\arraybackslash}p{(\columnwidth - 6\tabcolsep) * \real{0.2667}}@{}}
\toprule\noalign{}
\begin{minipage}[b]{\linewidth}\raggedright
Host
\end{minipage} & \begin{minipage}[b]{\linewidth}\raggedleft
Cutoff \(b_s\)
\end{minipage} & \begin{minipage}[b]{\linewidth}\raggedleft
\(\lceil\log_2(C_a(T)V_a(b_s))\rceil\)
\end{minipage} & \begin{minipage}[b]{\linewidth}\raggedleft
Uniform bound for \(r/d\)
\end{minipage} \\
\midrule\noalign{}
\endhead
\bottomrule\noalign{}
\endlastfoot
J14 & \(s=353\) & \(-29\) & \(<2^{-1218}\) \\
J56 & \(s=80\) & \(-40\) & \(<2^{-6720}\) \\
\end{longtable}

Thus for the current T and all the other fixed data, compatibility with(PC2) requires

\[
 \boxed{\ \ell<\ell_+/2^{353}\quad(J14),\qquad
          \ell<\ell_+/2^{80}\quad(J56)\ }.             \tag{PC10}
\]

The statement uses the 28-multiple convention in(PC9). These necessary departures are much larger than the 64-bit interval already excluded by DO044 and LIFT51. They do not prove compatibility at a smaller gap.

For an exact check on the limit of this particular upper bound, the next dyadic cutoffs give per-label logarithmic ceilings 11 and 126, respectively. The actual rational bounds at those cutoffs are greater than 1. Consequently this bound cannot exclude all gaps down to those next cutoffs; this is a limitation of the sufficient dimension bound, not evidence that the actual dimension inequality reverses there.

\subsection{4. Lowering T within the source size condition does not fix the endpoint}\label{lowering-t-within-the-source-size-condition-does-not-fix-the-endpoint}

The literal \nolinkurl{Columns/AbelianHostFieldConditions.lean}486--494 defines \nolinkurl{Sizes}; line 490 is \(4a_{\rm size}<T\). For the unchanged R022.2--3 scale, \(a_{\rm size}=A=2^{960}\); here the scalar \(a_{\rm size}\) is kept distinct from the matrix rank \(a\).

For every candidate \(T'>4A\),

\[
 C_a(T')<C_a^{\min}:=
 n(2\,3^{n^2})\mu_a
 \frac{g_0^a(g_0^n-1)g_1^a(g_1^n-1)}{(4A)^a-1}.
                                                        \tag{PC11}
\]

The source proof of the initial dimension formula uses the G-level neatness, which is unchanged, and the strict-T index \((T'^a-1)^f\). The bottom level is maximal at T, so its dimension bound has no compensating T-index. Hence (PC6--8) apply with this larger constant to every candidate satisfying the stated split and level hypotheses. No prime-existence assertion is used to form the upper bound.

At \(\ell=\ell_+\), one obtains

\[
 \frac rd<2^{-521598}\quad(J14),\qquad
 \frac rd<2^{-1255128}\quad(J56)                       \tag{PC12}
\]

for \textbf{every} such \(T'>4A\). Thus replacing the current large prime by a smaller allowed prime alone cannot repair the endpoint profile.

The interval lemma also gives the following uniform necessary gap reductions, even if T may be replaced:

\begin{longtable}[]{@{}
  >{\raggedright\arraybackslash}p{(\columnwidth - 6\tabcolsep) * \real{0.2000}}
  >{\raggedleft\arraybackslash}p{(\columnwidth - 6\tabcolsep) * \real{0.2667}}
  >{\raggedleft\arraybackslash}p{(\columnwidth - 6\tabcolsep) * \real{0.2667}}
  >{\raggedleft\arraybackslash}p{(\columnwidth - 6\tabcolsep) * \real{0.2667}}@{}}
\toprule\noalign{}
\begin{minipage}[b]{\linewidth}\raggedright
Host
\end{minipage} & \begin{minipage}[b]{\linewidth}\raggedleft
Cutoff \(b_s\)
\end{minipage} & \begin{minipage}[b]{\linewidth}\raggedleft
\(\lceil\log_2(C_a^{\min}V_a(b_s))\rceil\)
\end{minipage} & \begin{minipage}[b]{\linewidth}\raggedleft
Uniform bound for \(r/d\)
\end{minipage} \\
\midrule\noalign{}
\endhead
\bottomrule\noalign{}
\endlastfoot
J14 & \(s=313\) & \(-1\) & \(<2^{-42}\) \\
J56 & \(s=46\) & \(-65\) & \(<2^{-10920}\) \\
\end{longtable}

Compatibility therefore requires \(\ell<\ell_+/2^{313}\) or \(\ell<\ell_+/2^{46}\), respectively, under this more permissive prime variation. The next-cutoff rational bounds exceed 1, just as in Section 3. The actual cutoff integers are recorded in the certificate.

\subsection{5. Other necessary changes and the limits of this conclusion}\label{other-necessary-changes-and-the-limits-of-this-conclusion}

With the endpoint gap and initial space fixed, any enlarged projected bottom space satisfying the same small-error links must have dimension \(r'>3d/8\). If the original \(r>0\), this forces

\[
       \frac{r'}r>\frac38\,2^{589932}\quad(J14),\qquad
       \frac{r'}r>\frac38\,2^{2198952}\quad(J56).       \tag{PC13}
\]

If \(r=0\), a positive-dimensional replacement is required and a relative growth factor is not defined. Refining the bottom finite level or changing its projector might enlarge this space, but is a change of the actual \nolinkurl{EE}, \nolinkurl{dat0}, finite kernels, and source ties; neither modification is constructed or certified here.

Alternatively, keeping the current bottom and retaining a positive initial subspace of dimension \(d'\) requires \(d'<8r/3\), as proved in LIFT51. All of Mram is much larger than that, since \(d_{\rm ram}>d/2\). Selecting a single nonzero row can be useful for a different argument, but it does not preserve the literal full \nolinkurl{Mu0} and \nolinkurl{duS} record.

Changing only a nonzero scalar normalization, a Haar measure, a period, or the constants in the chosen proviso cannot change the rank bound. Changing the positive scale exponents cannot help at a point where the links apply, because \(\delta<1/2\) and \(e_0<1/4\) follow from \(x>1\), \(\varepsilon,\mathrm{ez}>0\) alone. Making the proviso false would make \nolinkurl{LinkS_at} vacuous, but would fail the required \nolinkurl{ChosenProviso}.

All these statements are necessary constraints. No assertion here constructs an admissible literal Host, repairs R023's independent Host/Scale obstruction, proves a new all-T period, or shows that any smaller gap satisfies the links. In particular the effective selected packet results established at \(\ell_+\) do not automatically transfer to a radically smaller gap. Such a new profile would need its own actual nonvanishing, source-link, and full gain/cost analysis.

\subsection{6. Verification and source locators}\label{verification-and-source-locators-1}

The companion \nolinkurl{necessary_profile_changes_certificate.py/.json} reconstructs the exact rational functions from the certified DO040/K044 integers, verifies the endpoint and interval comparisons, checks the next cutoffs, and tests the convexity expansion in exact small models. It computes no actual cusp dimension or restriction map.

ALP\_MACHINE / definite\_oscillator. (2026, November 22). \emph{AM-DO-044: An actual dimension obstruction to the full projected bottom packet}. \nolinkurl{definite_oscillator_packet_dimension.md}, DO44.5,12--24 and§6. The cohomological theorem and actual mass formula are carried results.

ALP\_MACHINE / kummer\_extension. (2026, November 22). \emph{AM-K044: An exact larger gap interval inside the same weight windows}. \nolinkurl{enlarged_gap_admissible_interval.md}, EG1--10.

Alpoge, L. (n.d.). \emph{Formalism source archive} {[}Lean source code{]}, pinned commit \nolinkurl{7cbf7a0c32fca1badf05931879a37c3f8a2b03b4}. \nolinkurl{Formalism/LocalToGlobal/Columns/AbelianHostFieldConditions.lean}486--494. For the exact small-error link definitions, see LIFT51§2.

\clearpage

\section{AM-K069 --- the isotropic cross map and a count from the integral block}\label{am-k069-the-isotropic-cross-map-and-a-count-from-the-integral-block}

Author: GPT-6 Astra Pro, ALP\_MACHINE \nolinkurl{/root/kummer_extension}. 22 November 2026 UTC. Inherited reasoning effort is not separately exposed.

\textbf{Author whole-proof check PASS --- 2026-11-22T10:38:55.878995+00:00.}

This report concerns the unchanged J56 finite signed remainder of K062, at the checked cutoff H=8192 and after the K068 cuts. It makes two refinements. Every remaining matrix fixes a positive E-subspace of dimension at least130, and its difference from the identity has rank at most39. For a specified class of matrices whose fixed space is determined by the integral positive block, a count in the original coordinates removes a further set of total absolute contribution less than 2\^{}-4084. No basis of a moving fixed space is assumed integral or bounded. The complementary isotropic cross action remains explicit.

\subsection{1. The unchanged support and notation}\label{the-unchanged-support-and-notation}

Use the determinant-one representative and decomposition of K068:

\[
a=f=168,\quad n=169,\quad m=164,\quad H=8192,
\qquad
h=\begin{pmatrix}A&C\\B&D\end{pmatrix}_{S\oplus W_0}.
\tag{IC1}
\]

Here S is the fixed positive coordinate space spanned by e5 through e168; W0 is the fixed five-coordinate space of signature(4,1); A is globally integral. Put

\[
X=A-I,\qquad r=\operatorname{rank}_E X\le33,
\qquad U=\ker X,\qquad K_A=\ker X\cap\ker X^*.
\tag{IC2}
\]

The star is conjugate transpose for the positive standard form on S. The actual central representative has operator norm at most H at every complex realization. Every one of its334 unmixed G reductions is unipotent. The summand, with both mixed projectors unchanged, is

\[
w(g)=a(g)c(g)\prod_{v\mid T}\rho(g_v),
\quad |c(g)|\le1,\quad
|\rho(g_v)|\le\frac2{T-2}\quad(v\text{ nonintegral}).
\tag{IC3}
\]

All sets below are subsets of the original nonintegral short sum: s(g) is at least1. Integral terms stay in the already charged K062 error. Central normalization preserves the same indexed classes and absolute summands, by K068(LR3).

\subsection{2. The cross map loses at most one fixed dimension}\label{the-cross-map-loses-at-most-one-fixed-dimension}

\textbf{Lemma.} If h is unitary for the form in IC1, then B(U) is totally isotropic in W0. Consequently

\[
\dim_E B(U)\le1,\qquad
\dim_E\{u\in S:h u=u\}\ge m-r-1,
\qquad
\operatorname{rank}_E(h-I)\le r+6.
\tag{IC4}
\]

\textbf{Proof.} For u,v in U one has h(u,0)=(u,Bu). Unitarity gives \(\langle Bu,Bv\rangle_{W_0}=0\). At a complex realization a totally isotropic subspace of a form with one negative coordinate has dimension at most one: projection to that negative coordinate is injective, because a vector in its kernel has positive norm unless it is zero. The E-dimension is unchanged by the realization, so the same bound holds over E. The kernel of B on U has dimension at least m-r-1, and its elements are fixed by h. They lie in the positive space S. Rank-nullity in dimension m+5 proves the final bound. ∎

In the actual residual r\textless=33 this gives a positive fixed space of dimension at least130 and rank(h−I)\textless=39. It does not assert global unipotence, and it does not require an arithmetic basis of that space.

There is also a small, block-determined candidate for the fixed space. The upper-left block of h*Jh=J is

\[
A^*A+B^*J_{W_0}B=I.
\]

For u in U it gives

\[
X^*u=-B^*J_{W_0}Bu.
\tag{IC5}
\]

The right side has rank at most one on U by IC4. Therefore

\[
\operatorname{codim}_S K_A\le r+1.
\tag{IC6}
\]

Define the following actual subset, without changing any coefficient:

\[
\mathscr F=\{[g]:h\text{ fixes }K_A\text{ pointwise}\}.
\tag{IC7}
\]

This condition is substantially more specific than merely having a large fixed space. Its precise complement has the following form.

\textbf{Lemma.} A matrix fails IC7 if and only if

\[
K_A=U,\quad L:=B(U)\ne0,\quad\dim_E L=1,
\quad L\text{ totally isotropic},\quad
\langle L,B(S)\rangle_{W_0}=0.
\tag{IC8}
\]

\textbf{Proof.} If B(U)=0, h fixes U and hence K\_A. Suppose B(U)=L is a nonzero line. On U, IC5 has the same kernel as B unless the linear map \(B^*J_{W_0}\) vanishes on L. If it does not vanish, a vector in K\_A has Bu=0, so h fixes K\_A. If it vanishes, IC5 gives X\emph{U=0. Since ker X and ker X} have the same dimension, they are equal; thus K\_A=U. The map B on that space is nonzero, so h does not fix it. Vanishing of B*J on L is exactly the last orthogonality in IC8. ∎

Thus the uncounted failure of IC7 is one isotropic cross direction inside the common kernel of X and X*. It is not an arbitrary loss of five fixed dimensions. No bound on the number or height of those directions is inferred here.

\subsection{3. Counting the integral block without a moving lattice basis}\label{counting-the-integral-block-without-a-moving-lattice-basis}

Fix r\textgreater=1. Every entry of X has all archimedean moduli at most H+1. For any number field E of degree2f with the current complex embeddings, the number of algebraic integers with that bound is at most

\[
                         (2H+3)^{2f}.
\tag{IC9}
\]

Indeed, embed O\_E in the real space underlying C\^{}f.~A nonzero integral element has absolute field norm, the product of its f squared complex moduli, at least1. The sum of those squared moduli is therefore at least f by the arithmetic-geometric mean inequality. Distinct embedded elements therefore have distance at least sqrt(f). The elements being counted lie in the ball of radius sqrt(f)(H+1). Disjoint balls of radius sqrt(f)/2 around them fit in the ball enlarged by that radius. The ratio of real2f-dimensional volumes is IC9. This argument does not need a discriminant bound.

A rank-r m-by-m matrix is determined by any invertible r-by-r minor, together with its r complete rows and r complete columns. In block notation the remaining block is

\[
X_{I^c,J^c}=X_{I^c,J}(X_{I,J})^{-1}X_{I,J^c}.
\tag{IC10}
\]

There are 2mr−r² stored entries. Taking the union over the possible row and column sets and applying IC9 gives the complete count

\[
N_A(r)\le
\binom mr^2(2H+3)^{2f(2mr-r^2)}.
\tag{IC11}
\]

The determinant of the chosen minor may have nontrivial finite valuation. Formula IC10 is used only for uniqueness of the matrix from the stored original entries; no inverse minor is applied to a lattice or treated as integral.

\subsection{4. The last column uses only r+6 original coordinates}\label{the-last-column-uses-only-r6-original-coordinates}

Now fix an A occurring in the class IC7. The subspace K\_A depends only on A. Since h fixes it and preserves the form, h e\_n is orthogonal to K\_A. Its S-component lies in the positive orthogonal complement of K\_A inside S, a space of dimension at most r+1 by IC6. Choose once for this A a set of r+1 original S coordinates whose projection is injective on that complement. Such a set exists by an ordinary nonzero-minor argument; add unused coordinates if its dimension is smaller. No extra choice factor is needed: use one fixed set for each A.

The full last column is therefore determined by these r+1 entries and its five W0 entries. This is an injective encoding by original matrix coordinates; it introduces no change-of-basis denominator.

Fix a support set of size s and actual ordered source/target pencil points at those places. The labelled pencil acts only in W0. Thus at each CM component its target lattice has the exact decomposition

\[
L'=O_v^m\oplus L'_{W_0},\quad
\varpi O_v^5\subset L'_{W_0}\subset\varpi^{-1}O_v^5,
\quad \operatorname{vol}(L'_{W_0})=\operatorname{vol}(O_v^5).
\tag{IC12}
\]

The actual source contains varpi e\_n.~Hence h e\_n lies in varpi\^{}-1 L'. The r+1 selected S entries lie in varpi\^{}-1 O\_v, while the W0 block lies in varpi\^{}-1 L'\_(W0). The latter contains O\_v\^{}5 and has index T\^{}5 over it, by its determinant valuation. Thus the encoded column lies in an over-lattice of O\_E\^{}(r+6) with exactly

\[
                             T^{2(r+6)s}
\tag{IC13}
\]

cosets. Both conjugate CM components are counted. All other finite places are integral. This is the actual target over-lattice of DO059(SC24--27), restricted to the selected original coordinates.

The encoded column has Euclidean norm at most H at each realization, because it is an orthogonal coordinate projection of h e\_n.~In each coset, the same packing proof as IC9, now in real dimension2f(r+6), gives at most (2H+1)\^{}(2f(r+6)) columns. For a fixed full last column, the actual source-lattice stabilizer fibre has order at most

\[
G_n=\gcd(|GL_n(\mathbf F_{g_0})|,|GL_n(\mathbf F_{g_1})|),
\tag{IC14}
\]

as in DO058/DO059. It is the finite group fixing the negative vector e\_n in the source lattice; no cutoff subset is claimed to be a group. Adding the support and pencil choices and then IC11 yields

\[
\begin{split}
N_{\mathscr F}(r,s)\le{}&
\binom mr^2(2H+3)^{2f(2mr-r^2)}
\binom fs(T-1)^{2s}G_n\\
&\hspace{8mm}\cdot(2H+1)^{2f(r+6)}T^{2(r+6)s}.
\end{split}
\tag{IC15}
\]

Every actual allowed pencil pair and pole depth is included. Counting the same matrix through several pairs or pivot choices only enlarges this upper bound. The argument does not count an unrelated full theta group or replace the labelled pencil.

\subsection{5. Exact finite-character bound and the removable range}\label{exact-finite-character-bound-and-the-removable-range}

For rank X=r, choose a nonzero r-minor. Since all unmixed reductions are unipotent, K068(LR7--12) with p=1 and k=r gives

\[
|a(g)|\le
\frac{2(2^rH)^{af}}{Q^{L_r(f-1)}},
\quad Q=g_0g_1,
\quad L_r=rn-\frac{r(r+1)}2.
\tag{IC16}
\]

The factor2 includes the entire finite Euler correction; both actual mixed projectors retain only their modulus bound1. Combining IC3, IC15 and IC16 proves a complete absolute bound, denoted B(r,s), for the actual intersection of IC7 with the rank-r, support-s stratum. It is exactly the right side of IC15 multiplied by

\[
\left(\frac2{T-2}\right)^s
\frac{2(2^rH)^{af}}{Q^{L_r(f-1)}}.
\tag{IC17}
\]

The following conservative rational exponents make the large range checkable without huge rational-matrix enumeration. The actual primes and integer G\_n satisfy

\[
Q>2^{70},\quad T<2^{996},\quad G_n<2^{5022}.
\]

The exact integer comparisons

\[
(2H+3)^{2048}<2^{14\cdot2048+1},\qquad
(2H+1)^{8192}<2^{14\cdot8192+1}
\tag{IC18}
\]

give strict logarithmic upper bounds14+1/2048 and14+1/8192. Also binom(m,r)\textless=2\^{}m, binom(f,s)\textless=2\^{}f, and

\[
\frac{2(T-1)^2}{T-2}<4T.
\]

Consequently

\[
\log_2 B(r,s)<C_r+b_rs,
\tag{IC19}
\]

where all constants on the right are rational or integral:

\[
\begin{split}
C_r={}&f+2m+5022
 +2f(2mr-r^2)\left(14+\frac1{2048}\right)\\
&+2f(r+6)\left(14+\frac1{8192}\right)
 +1+(r+13)af-70L_r(f-1),\\
b_r={}&2+996(2r+13),\\
s_r={}&\left\lfloor\frac{-4096-C_r}{b_r}\right\rfloor.
\end{split}
\tag{IC20}
\]

For 2\textless=r\textless=33 the resulting complete range is:

\begin{longtable}[]{@{}rrrrrrrr@{}}
\toprule\noalign{}
r & s\_r & r & s\_r & r & s\_r & r & s\_r \\
\midrule\noalign{}
\endhead
\bottomrule\noalign{}
\endlastfoot
2 & 22 & 10 & 104 & 18 & 129 & 26 & 140 \\
3 & 40 & 11 & 108 & 19 & 131 & 27 & 140 \\
4 & 55 & 12 & 112 & 20 & 132 & 28 & 141 \\
5 & 67 & 13 & 116 & 21 & 134 & 29 & 142 \\
6 & 76 & 14 & 119 & 22 & 135 & 30 & 142 \\
7 & 85 & 15 & 122 & 23 & 136 & 31 & 143 \\
8 & 92 & 16 & 125 & 24 & 138 & 32 & 143 \\
9 & 98 & 17 & 127 & 25 & 139 & 33 & 144 \\
\end{longtable}

Thus every such stratum with1\textless=s\textless=s\_r has complete absolute contribution less than2\^{}-4096. There are exactly3677 pairs (r,s) in this table and range; hence their union satisfies

\[
\boxed{\quad
\sum_{\substack{[g]\in\mathscr F,\ 2\le r\le33\\
                 1\le s(g)\le s_r}} |w(g)|
 <3677\,2^{-4096}<2^{-4084}.
\quad}
\tag{IC21}
\]

If the preceding remainder already has s\textgreater=4, only3581 pairs occur. The same upper bound is valid. If other subsets were removed earlier, take their set difference before applying IC21. The displayed s\_r is the largest integer satisfying this particular conservative exponent test; a failed next integer is not a lower bound on the actual sum. The test gives no positive support range at r=1, and no claim for r=0 is made by IC16.

\subsection{6. The remaining structural alternatives and scope}\label{the-remaining-structural-alternatives-and-scope}

After IC21 is charged as an additional error, every remaining matrix of rank2\textless=r\textless=33 either has s\textgreater s\_r, or satisfies exactly IC8. Ranks0 and1 remain separate. All remaining matrices still obey the improved positive fixed dimension130 and global rank39 bound of IC4.

This removal is invariant under inversion. The inverse-adjoint identity gives A\_(h\^{}-1)=A\_h*, so it preserves r and exchanges the two kernels defining K\_A. The inverse fixes that same K\_A exactly when h does. The original support size and cutoff are also preserved, as in R043. Consequently the original adjoint relation still makes the complementary signed sum real; no individual sign has been changed.

The orthogonal complement used for the column count was determined by A, and only its coordinate projection was used. No bounded integral basis or small discriminant of a moving fixed lattice was asserted. The exceptional IC8 line is not covered by that count. The result does not prove the sign of the residual finite sum, all-T Mram positivity, a prescribed arithmetic-image pairing, a source LinkS estimate, a net DP44 inequality, or Kummer RH. In particular the separate upper-gap source-link incompatibility of LIFT51 is unchanged.

The proof uses the already checked actual support and over-lattice identities K062(AT1--3,13--20), DO059(SC24--30), the integral block DO060(IB3--7), and K068(LR1--22). Its finite-character input is precisely K068(LR12), already derived from Green's formula in K067(PS9); no additional finite representation theorem is invoked. The new steps are IC4--8, the injective coordinate count IC11--15, and its specified arithmetic application IC18--21.

The companion certificate checks the dimension identities on exact unitary models, including a nontrivial exceptional cross action and the sharp joint-kernel loss; the rank chart reconstruction; actual target-lattice determinant/containment models; and every rational threshold and strict integer comparison. It does not enumerate the global rational support or evaluate an automorphic period.

\subsection{Author verification record}\label{author-verification-record-13}

At 2026-11-22T10:38:55.878995+00:00, I completed a whole read and reconstruction of IC1--21 and all six sections, together with the entire certificate. The fixed-space and exact exceptional-line lemmas, integral rank-chart encoding, injective original-coordinate column projection, actual target over-lattice index, unchanged finite fibre, paired-G character estimate, all rational support thresholds and inversion stability were checked. The final exact replay passed 4800 assertions at 2026-11-22T10:38:43.455118+00:00, including60 rational unitary models,45 rank charts,52 local pencil models and every one of the3677 range inequalities. These diagnostics supplement the symbolic proof and do not enumerate the global rational support. No theorem outside the stated unchanged finite-family scope is promoted.

\clearpage

\textbf{Later size-interface result.} The profile and code inequalities below are necessary conditions, not a repair constructor. LIFT54 excludes full SizesAt at both fixed degrees for every Scale and excludes m2 universally. LIFT55 excludes the unchanged L343 scale even after changing all ranks, degree and primes. Changing only the gap or k therefore does not produce a literal full family.

\section{ALP-LIFT-53 --- the literal code requirement forces a larger initial rank}\label{alp-lift-53-the-literal-code-requirement-forces-a-larger-initial-rank}

Author: GPT-6 Astra Pro, ALP\_MACHINE \nolinkurl{/root/lift_accounting}. 22 November 2026 UTC. Inherited reasoning effort is not separately exposed.

Status: author whole-proof review PASS; exact diagnostic passed 4,678 assertions.

\subsection{1. Statement and separation of the quantifiers}\label{statement-and-separation-of-the-quantifiers}

For the concrete K031/K032 family, the initial rank and real field degree are \(k=f=42\) or 168, the number of horizontal edges is \(m=2\), and the R022 scale has \(L=343\). These equalities do not change when the K044 gap \(\ell\) or its permitted canonical parameter changes.

The literal \nolinkurl{ChCode_at} condition, with \(C_0\geq1\), \(C_H\geq0\), forces

\[
 C_\Gamma<\frac1{31360},\qquad
 C_H<\frac{140}{1597}\quad(J14),\qquad
 C_H<\frac{581}{1597}\quad(J56).                       \tag{CB1}
\]

In particular it fails at every gap for the current ranks whenever \(C_\Gamma\geq1\) or \(C_H\geq1\). This is an independent obstruction from the dimension estimate in LIFT51--52.

There is no assertion that \nolinkurl{SizeData₂.CGam_pos} by itself imposes \(C_\Gamma\geq1\): it only imposes positivity. The lower bound 1 is instead an allowed request in the source \textbf{every-K} design predicate. Its exact meaning and the eventual-threshold issue are proved in Section 4. The pointwise theorem also applies to any supplied actual constants satisfying these inequalities, independently of that design predicate.

\subsection{2. Exact source code and the logarithmic inequality}\label{exact-source-code-and-the-logarithmic-inequality}

All source locators below refer to the pinned commit \nolinkurl{7cbf7a0c32fca1badf05931879a37c3f8a2b03b4} in \nolinkurl{work/research/alpoge_source/}.

\nolinkurl{Formalism/SharpRestriction/CuspSpanNormBounds/DataS.lean}25--27 contains an injective map

\[
 \operatorname{letter}:\mathrm{Word}\longrightarrow
                  \{0,1\}^{k-2}.
\]

Consequently, for \(k\geq2\),

\[
                     |\mathrm{Word}|\leq2^{k-2}.      \tag{CB2}
\]

This is an upper bound even if only some of the possible binary words are used. Repeating a word with additional labels would violate the same injectivity clause; it cannot enlarge the literal code.

\nolinkurl{Formalism/SharpRestriction/GlobalSchurBounds/Chosen.lean}14,56--57 defines

\[
 \Gamma=\exp(C_\Gamma m^5 fL),\qquad
             C_0\Gamma^2 x^{C_H}\leq|\mathrm{Word}|.
                                                        \tag{CB3}
\]

The source \nolinkurl{ChosenEx} includes \(C_0\geq1\) at lines86--87. We state that hypothesis explicitly when using \nolinkurl{ChosenE} or \nolinkurl{DesignExists}, whose \(C_0\) is an external parameter. The nonnegative \(C_H\) hypothesis is explicit in \nolinkurl{DesignExists}, discussed below.

Combining(CB2--3) and taking logarithms gives the exact necessary bound

\[
 (k-2)\log2\ \geq\ \log C_0+2C_\Gamma m^5 fL+C_H\log x.
                                                        \tag{CB4}
\]

The left side of(CB3) is positive, so the logarithms are legitimate. For \(x>1\), \(C_0\geq1\), \(C_H\geq0\), and \(C_\Gamma>0\), dropping the two nonnegative terms and using \(\log2<7/10\) yields

\[
                  k-2>\frac{20}{7}C_\Gamma m^5 fL.    \tag{CB5}
\]

For \(m=2\), \(L=343\), the rational coefficient is exactly

\[
                  \frac{20}{7}\,2^5\,343=31360.
                                                        \tag{CB6}
\]

Thus(CB5), with \(k=f\), proves the first bound in(CB1).

The second term in(CB4) is strictly positive, so dropping it gives \(C_H<(f-2)\log2/\log x\). R022.2--3 has \(x=2^{960}e^{-343}\); the exact bounds \(69/100<\log2<7/10\) imply

\[
 \log x>960\frac{69}{100}-343=\frac{1597}{5}.
\]

Hence \(C_H<7(f-2)/3194\), which is precisely the remaining two bounds in(CB1). The logarithm inequalities are also certified directly by a positive rational series in the companion diagnostic.

These are code constraints, not estimates of a theta operator. They are unaffected by changing \(\ell\), the canonical parameter, a period, or a scalar normalization while \(k,f,m,L,x\) remain fixed.

\subsection{3. Required rank and the actual G-prime restrictions}\label{required-rank-and-the-actual-g-prime-restrictions}

Suppose the concrete scale is retained but the even initial rank \(k\) may change independently of \(f\). If \(C_\Gamma\geq1\), (CB5--6) and parity force the following necessary lower bounds:

\begin{longtable}[]{@{}
  >{\raggedright\arraybackslash}p{(\columnwidth - 8\tabcolsep) * \real{0.1579}}
  >{\raggedleft\arraybackslash}p{(\columnwidth - 8\tabcolsep) * \real{0.2105}}
  >{\raggedleft\arraybackslash}p{(\columnwidth - 8\tabcolsep) * \real{0.2105}}
  >{\raggedleft\arraybackslash}p{(\columnwidth - 8\tabcolsep) * \real{0.2105}}
  >{\raggedleft\arraybackslash}p{(\columnwidth - 8\tabcolsep) * \real{0.2105}}@{}}
\toprule\noalign{}
\begin{minipage}[b]{\linewidth}\raggedright
Host
\end{minipage} & \begin{minipage}[b]{\linewidth}\raggedleft
\(f\)
\end{minipage} & \begin{minipage}[b]{\linewidth}\raggedleft
Necessary even \(k\)
\end{minipage} & \begin{minipage}[b]{\linewidth}\raggedleft
Necessary \(N_1=k+t+1\)
\end{minipage} & \begin{minipage}[b]{\linewidth}\raggedleft
Necessary \(N_2=k+2(t+1)\)
\end{minipage} \\
\midrule\noalign{}
\endhead
\bottomrule\noalign{}
\endlastfoot
J14 & 42 & 1,317,124 & 2,634,256 & 3,951,388 \\
J56 & 168 & 5,268,484 & 10,536,976 & 15,805,468 \\
\end{longtable}

The exact source \nolinkurl{Formalism/LocalRingCardinality/Params.lean}6--22 requires \(k\) even, \(t\) odd, \(m\geq2\), \(m+3\leq k\), \(k+m-1+5\leq t\), and \(f\leq k\). It does \textbf{not} require \(k=f\). For \(m=2\), parity sharpens \(t\geq k+6\) to

\[
 t\geq k+7,\qquad N_1\geq2k+8,\qquad N_2\geq3k+16.
                                                        \tag{CB7}
\]

This proves the displayed target-rank columns. They are necessary floors, not sufficient choices satisfying(CB4) or all other source conditions. In particular the old common middle rank 1802 cannot be retained under such a rank change.

The source numerical prime condition is \nolinkurl{Columns/AbelianHostFieldConditions.lean}488:

\[
                       2fk<g_j-1\quad(j=0,1).          \tag{CB8}
\]

With the actual primes of K031/K032, the smaller G prime gives the following exact maximum \textbf{even} rank:

\begin{longtable}[]{@{}lrr@{}}
\toprule\noalign{}
Host & Current \(g_0\) & Maximum even \(k\) permitted by(CB8) \\
\midrule\noalign{}
\endhead
\bottomrule\noalign{}
\endlastfoot
J14 & 4,615,409 & 54,944 \\
J56 & 22,285,932,761 & 66,327,180 \\
\end{longtable}

Thus the fixed J14 primes already exclude the \(C_\Gamma\geq1\) necessary floor. The J56 primes do \textbf{not} exclude that floor alone. It would be incorrect to conclude that both pairs must change merely from the table in this section.

For J56, an allowed larger constant request does give an explicit failure: \(C_\Gamma\geq13\) forces \(k\geq68,490,244\), already above 66,327,180. In general larger requested \(C_\Gamma\) forces larger \(k/f\), and any fixed G primes eventually cease to satisfy(CB8).

\subsection{4. What the every-K design condition actually requires}\label{what-the-every-k-design-condition-actually-requires}

The exact source is \nolinkurl{Formalism/SharpRestriction/ParamsHeightRankConstant.lean}:

\begin{itemize}
\tightlist
\item
  lines111--112 define \nolinkurl{cMid(CH)} as a list beginning with \([C_H,C_\Gamma,\ldots]\);
\item
  lines126--130 define \nolinkurl{TransportedAlt} with arbitrary real-valued functions \nolinkurl{t}, \nolinkurl{mid}, and \nolinkurl{last};
\item
  lines138--140 define \nolinkurl{FamP.Large K CH}, which in particular requires \(K.\mathrm{mid}(\mathrm{cOut})\leq c\) for every entry \(c\in\mathrm{cMid}(C_H)\);
\item
  lines187--190 define \nolinkurl{DesignExists(𝒞,C0)} with \textbf{for every} \nolinkurl{K}, there exist a family and \(C_H\geq0\) satisfying \nolinkurl{Large K CH}, \nolinkurl{ChosenE}, and the other displayed predicates.
\end{itemize}

For any positive \(M\), choose the legitimate request

\[
 K_M.\mathrm{pre}=0,\quad K_M.t=0,\quad
 K_M.\mathrm{mid}(l)=M,\quad K_M.\mathrm{last}(l)=0.
                                                        \tag{CB9}
\]

If a family is large for this request, then

\[
                        C_H\geq M,\qquad C_\Gamma\geq M.
                                                        \tag{CB10}
\]

The choice \(M=1\) gives the obstruction(CB1) for every concrete K044 gap at the current ranks. The choice \(M=13\), allowing an enlarged rank but retaining the current scale and G primes, gives the further J56 obstruction in Section 3. No positivity condition has been invented for \nolinkurl{TransportedAlt}: these are valid explicit choices of its unrestricted fields.

\nolinkurl{ChosenE} includes \nolinkurl{ChosenCode}, which is an \textbf{eventual} predicate. The pointwise incompatibility applies only when the point is above its threshold. A fixed concrete point below that threshold is not a counterexample to the predicate. Therefore the conclusion about a design is the following precise one: it cannot meet the indicated requests while preserving the excluded concrete ranks/scale/prime constraints at an eligible point where the code condition applies. If the design is allowed to choose different data in its eventual range, those data are outside this excluded cohort.

There is also a direct asymptotic necessity which does not require the auxiliary request \(M=1\). \nolinkurl{FamCore.lean}113--114 fixes one positive \(C_\Gamma\) throughout a family; its positivity follows from \nolinkurl{SizeData₂.CGam_pos} and the nonempty family at 0. At every point beyond the scale-tie threshold, \(L=x^\varepsilon\), with \(\varepsilon>0\), by \nolinkurl{FamCore.lean}73--85 and \nolinkurl{DataParams₂.lean}9--27. The rank clauses give \(f\geq m\geq2\), so division by \(f\) is legitimate. Since \(m\geq2\), (CB5) implies

\[
                     \frac{k}{f}>
                    \frac{640}{7}C_\Gamma x^\varepsilon.
                                                        \tag{CB11}
\]

The right side tends to infinity. Thus a family satisfying the code condition, with \(C_0\geq1\), \(C_H\geq0\), cannot have a bounded ratio \(k/f\) along arbitrarily large eligible points. To prove this without an implicit nonemptiness assumption, use the unbounded window and its nonempty index set in \nolinkurl{FamCore.lean}78--80, choosing a point beyond both the code threshold and the desired lower bound on \(x^\varepsilon\). In particular an eventual family with \(k=f\) cannot work, even if its standalone \(C_\Gamma\) was chosen very small. This is a constraint on family growth, not an assertion that \(m=2\) is impossible when the other ranks grow.

\subsection{5. What increasing k does and does not solve}\label{what-increasing-k-does-and-does-not-solve}

The literal rank axioms allow increasing \(k\) independently of the fixed CM degree. Therefore there is a numerical direction beyond the excluded \(k=f\) profile: enlarge the initial rank, change \(t\) and both target ranks according to(CB7), and recheck the G primes. The necessary noncompact gap scale then changes. In the two-edge signature the terminal odd rank is \(S=k+1+2t\); the K044 argument for the smallest canonical parameter requires

\[
                         2\ell>S^3.                  \tag{CB12}
\]

The top weight window requires \(\ell\leq E_T-B-k/2-3\), rounded down to a 28-multiple. Those two bounds may still be compatible at million-scale \(k\), but that compatibility and the remaining word, finite-character, and cost conditions must be checked for the new ranks. The existing fixed-rank proofs and numerical norms are not automatically reusable.

This report proves neither a new complete \nolinkurl{Ranks}/Host constructor nor a full new \nolinkurl{ChosenE} instance. In particular satisfying the code-size inequality alone does not establish the separate chosen proviso, dimension-error links, actual nonzero period, or DP44 cost comparison. It identifies an unavoidable further change after the smaller-gap possibility in LIFT52 and keeps the independent R023 Host/Scale obstruction intact.

\subsection{6. Verification and references}\label{verification-and-references-1}

The companion \nolinkurl{chosen_code_rank_certificate.py/.json} checks the rational logarithm bounds, exact parity/rank and G-prime arithmetic, binary injection counts in finite models, and source fingerprints. No family, automorphic form, or complete source record is computed.

Alpoge, L. (n.d.). \emph{Formalism source archive} {[}Lean source code{]}, pinned commit \nolinkurl{7cbf7a0c32fca1badf05931879a37c3f8a2b03b4}. Exact paths and lines are given in Sections 2--4.

ALP\_MACHINE / root. (2026, November 22). \emph{AM-R-022: Narrow-cloud weight profile}. \nolinkurl{work/research/narrow_cloud_weight_profile.md}, R22.2--3.

ALP\_MACHINE / kummer\_extension. (2026, November 22). \emph{AM-K044: An exact larger gap interval inside the same weight windows}. \nolinkurl{enlarged_gap_admissible_interval.md}, EG1--10, for the carried gap and top-window formulas.

\clearpage

\section{ALP-LIFT-54 --- the literal size record excludes two-step chains}\label{alp-lift-54-the-literal-size-record-excludes-two-step-chains}

Author: GPT-6 Astra Pro, ALP\_MACHINE \nolinkurl{/root/lift_accounting}, 22 November 2026 UTC. Reasoning effort is not separately exposed.

\textbf{Status.} Author whole-proof check PASS at 2026-11-22T10:52:45.058831+00:00. Exact certificate: 413 assertions PASS. Independent review is separate. No Lean build is claimed.

\subsection{1. Statement and scope}\label{statement-and-scope-1}

The literal source condition \nolinkurl{SizesAt} cannot hold with \nolinkurl{Ranks.m = 2}, for any choice of initial rank, weights, scale, finite primes, or CM field. This follows from one explicit numerical conjunct, rather than from an assertion about the nonvanishing or size of a theta lift. More generally, every such record necessarily satisfies

\[
 m\ge 2+2(\log f)^2,
 \qquad
 m\ge 2+2\bigl\lceil(\log f)^2\bigr\rceil,
                                                        \tag{SL1}
\]

where the second bound uses the required parity of m and all logarithms are natural. The second inequality is sharp for this single logarithmic condition at fixed f.~It is not sufficient for \nolinkurl{SizesAt}, \nolinkurl{Host}, or any actual theta construction.

At f = 42 and f = 168 the exact necessary even floors from this condition are, respectively, m = 30 and m = 56. Combining the same condition with the accepted AM-R-023 lower bound for a literal Host whose field retains K14 or K56 gives m at least 79,368. If one also retains the full \nolinkurl{SizeM} and \nolinkurl{o_tie} conditions, the stronger bound is

\[
                 m>21^{10371/4}>2^{11388}.
                                                        \tag{SL2}
\]

More strongly, neither fixed degree f = 42 nor f = 168 admits \nolinkurl{SizesAt} at \textbf{any} Scale: its required lower and upper bounds for k contradict each other. At the unchanged L343 scale the explicit necessary bounds are m \textgreater{} 2\^{}231, f \textgreater{} 2\^{}998 and k \textgreater{} 2\^{}2177. These conclusions concern the literal source parameter package. They do not invalidate the independently constructed mixed theta kernels, and they do not show that every modified mixed construction has a negative DP44 difference.

\subsection{2. Exact source clauses}\label{exact-source-clauses}

The source is the recovered archive pinned at commit \nolinkurl{7cbf7a0c32fca1badf05931879a37c3f8a2b03b4}. The relevant definitions have been read directly:

\begin{enumerate}
\def\labelenumi{\arabic{enumi}.}
\item
  \nolinkurl{Formalism/LocalRingCardinality/Params.lean}, lines 6--22 and 28--38, requires m even, 2 \textless= m \textless= f, defines f = 2 ord, and defines the natural number rS = m / 2 - 1.
\item
  \nolinkurl{Formalism/SharpRestriction/ElevenLe.lean}, line 44, makes

  \[
                 \log f\le\sqrt{r_S}
                                                       \tag{SL3}
  \]

  the ninth conjunct of \nolinkurl{SizeKRel}. Lines 134--138 make \nolinkurl{SizeKRel} a conjunct of \nolinkurl{SizesAt}.
\item
  \nolinkurl{Formalism/SharpRestriction/ParConsts.lean}, line 31 and line 44, defines \nolinkurl{FamP} as extending \nolinkurl{FamCore} and requires \nolinkurl{sizes : ∀ x, x1 < x → ∀ i : Ix x, SizesAt (L x i)}.
\item
  \nolinkurl{Formalism/SharpRestriction/FamCore.lean}, lines 73--85, requires x1 \textgreater{} 1, an unbounded set Win, nonempty Ix x for x in Win, and equality of the actual Scale.x with x when x \textgreater{} x1. Its theorem \nolinkurl{FamCore.exists_pt}, lines 138--142, supplies such an eligible point above every proposed threshold.
\end{enumerate}

The numerical \nolinkurl{Sizes} predicate previously used for the finite prime construction is a different object. In particular, its \nolinkurl{g_big} and \nolinkurl{T_gt} clauses in \nolinkurl{AbelianHostFieldConditions.lean}, lines 486--494, do not replace \nolinkurl{SizesAt}.

\subsection{3. The logarithmic obstruction, including natural-number rounding}\label{the-logarithmic-obstruction-including-natural-number-rounding}

Write m = 2h. The source parity and m \textgreater= 2 give an integer h \textgreater= 1. Hence its natural-number division and subtraction satisfy rS = h - 1 exactly; there is no truncation error on passage to real numbers. Since f \textgreater= m \textgreater= 2, log f \textgreater{} 0. Squaring (SL3), whose two sides are nonnegative, gives

\[
             (\log f)^2\le h-1=\frac m2-1.
                                                        \tag{SL4}
\]

This proves the first part of (SL1). Since h - 1 is an integer, (SL4) is equivalent to h - 1 \textgreater= ceil((log f)\^{}2), proving the second part. Conversely that rounded inequality implies (SL3), so the floor is exact for this one condition.

For m = 2, the right side of (SL3) is zero and its left side is strictly positive. This contradiction is independent of all remaining source data. In fact the same source conditions imply the universal necessary bound

\[
                              m\ge20.
                                                        \tag{SL5}
\]

To check this without a numerical guess, use f \textgreater= m in (SL4). For the nine even integers m = 2, 4, \ldots, 18 the exact positive-series bounds described below give (log m)\^{}2 \textgreater{} m/2 - 1. For m = 20 they give (log 20)\^{}2 \textless{} 9. Thus m = f = 20 passes this isolated test, establishing that (SL5) cannot be improved using only parity, 2 \textless= m \textless= f, and (SL3). This last observation asserts no Frame or Host at f = 20.

The field-specific intervals are

\[
 13<(\log42)^2<14,
 \qquad26<(\log168)^2<27.
                                                        \tag{SL6}
\]

Thus m = 28 and m = 54 fail, while m = 30 and m = 56 respectively pass the isolated logarithmic test. Although the actual K14 and K56 fields already fail a separate literal \nolinkurl{HostRaw} condition by AM-R-023, these are useful necessary numerical floors for their degree parameters; no actual source Host on either field is being assumed.

All displayed logarithm intervals use rational arithmetic. For a \textgreater{} 1 set z = (a - 1)/(a + 1). The identity and tail bound

\[
 \log a=2\sum_{j\ge0}\frac{z^{2j+1}}{2j+1},\qquad
 0<\log a-2\sum_{j=0}^{N-1}\frac{z^{2j+1}}{2j+1}
 \le\frac{2z^{2N+1}}{(2N+1)(1-z^2)}
                                                        \tag{SL7}
\]

follow by integrating the geometric series for 1/(1-z\^{}2). The certificate uses a = 2, 3, 5, 7 and N = 64, 64, 96, 128 respectively. Factoring 42, 168 and the small m values reduces every comparison to these four intervals. No floating-point approximation enters a proof or a pass condition.

\subsection{4. Exact family quantifiers and unbounded length}\label{exact-family-quantifiers-and-unbounded-length}

Let Fm be a literal \nolinkurl{FamP eps}. For \textbf{every} x \textgreater{} Fm.x1 and \textbf{every} i : Fm.Ix x, its \nolinkurl{sizes} field supplies (SL3). Consequently its datum at that eligible point cannot have m = 2. This conclusion is immediate after the specified x1, and does not await the distinct eventual thresholds in \nolinkurl{ChosenCode}, \nolinkurl{LinkS}, or \nolinkurl{LinkIRR}.

The qualification on x matters. The \nolinkurl{sizes} field does not assert \nolinkurl{SizesAt} at x \textless= x1, and the nonempty dummy fibre at x = 0 is not an eligible point for this argument. Nor does the argument claim that \nolinkurl{D7}, \nolinkurl{Ranks}, or an actual two-edge theta kernel alone is inconsistent. It excludes their simultaneous extension to the stated size predicate at an eligible family point.

The same source supplies a stronger growth condition. \nolinkurl{ElevenLe.lean}, lines 11 and 27--29, together with \nolinkurl{DataParams₂.lean}, lines 27 and 45, gives

\[
 m\ge x^{-\epsilon/100}\sqrt{x\,x^{-\epsilon/8}}
     =x^{1/2-29\epsilon/400}.
                                                        \tag{SL8}
\]

The Scale clauses x \textgreater{} 1 and 0 \textless{} epsilon \textless{} 1/52 therefore imply

\[
 m>x^{10371/20800}.
                                                        \tag{SL9}
\]

To see the exact coefficient, 1/100 + 1/16 = 29/400 and 1/2 - 29/(400·52) = 10371/20800. Strictness in (SL9) comes from epsilon \textless{} 1/52 and x \textgreater{} 1, not from changing the weak inequality in \nolinkurl{SizeM}.

It follows that no literal FamP can have m bounded on its unbounded eligible window. Given any bound B \textgreater= 1, choose y \textgreater{} max(x1, B\^{}(20800/10371), 1). \nolinkurl{FamCore.exists_pt} supplies an eligible x \textgreater= y and nonempty Ix x, and (SL9) forces m \textgreater{} B there. This conclusion holds even if the initial rank k and the field degree vary. In particular, the fixed floors 30 and 56 are pointwise necessary floors; they are not fixed-length asymptotic-family constructions.

\subsection{5. Neither fixed degree admits the full size predicate at any Scale}\label{neither-fixed-degree-admits-the-full-size-predicate-at-any-scale}

The following strengthens the two-step obstruction without using the \nolinkurl{Host} field axioms. \nolinkurl{SizeN} and epsilon \textgreater{} 0 imply

\[
 x^2<x^{2+799\epsilon/100}\le n=\frac f2.
                                                        \tag{SL10a}
\]

The upper exponent in \nolinkurl{SizeK} is 9/2 + 1001 epsilon/100. It is less than 5, since epsilon \textless{} 1/52 and 9/2 + 1001/5200 \textless{} 5. The first conjunct of \nolinkurl{SizeKRel} has the strictly positive extra multiplier x\^{}(201 epsilon/100) \textgreater{} 1. Hence

\[
 m^5f<k<x^5<\left(\frac f2\right)^{5/2},
 \qquad f^3>32m^{10}.
                                                        \tag{SL10b}
\]

The second statement follows by squaring the first bound on m\^{}5f and canceling f\^{}2 \textgreater{} 0. All inequalities are strict for the reasons just given; no source weak inequality has been silently strengthened.

For f = 42, (SL10a) gives x \textless{} 5, hence k \textless{} 3125. The logarithmic floor m \textgreater= 30 instead gives k \textgreater{} 30\^{}5·42 = 1,020,600,000. For f = 168 it gives x \textless{} 10, hence k \textless{} 100,000, while m \textgreater= 56 gives k \textgreater{} 56\^{}5·168 = 92,522,938,368. Both are contradictions. The rounding, powers, and products are exact integer computations.

Therefore changing ell, k, m, T, or Scale while retaining either of these fixed degree parameters cannot supply the literal \nolinkurl{SizesAt}. One must change the degree parameter or prove an applicable replacement for the size interface. This statement is independent of the separate AM-R-023 \nolinkurl{HostRaw} exclusion.

\subsection{6. Explicit costs at the unchanged L343 scale}\label{explicit-costs-at-the-unchanged-l343-scale}

Use precisely the actual scale from AM-R-022 and the subsequent narrow-cloud profiles:

\[
 L=343,\quad x=2^{960}e^{-343},\quad
 \epsilon=\frac{\log343}{\log x}.
                                                        \tag{SL10}
\]

These values are retained here; varying only the admissible gap ell or the canonical twist does not change them. The lower bounds in \nolinkurl{SizeM}, \nolinkurl{SizeN} and \nolinkurl{SizeK} become

\[
 \begin{aligned}
 \log m&\ge\tfrac12\log x-\tfrac{29}{400}\log343,\\
 \log n&\ge2\log x+\tfrac{799}{100}\log343,\\
 \log k&\ge\tfrac92\log x+\tfrac{999}{100}\log343.
 \end{aligned}
                                                        \tag{SL11}
\]

The coefficients in the last two lines are exactly 8 - 1/100 and 10 - 1/100. With log x = 960 log 2 - 343 and log 343 = 3 log 7, (SL7) gives the following strict rational enclosures after division by log 2:

\begin{longtable}[]{@{}lr@{}}
\toprule\noalign{}
Lower-bound expression & Certified interval \\
\midrule\noalign{}
\endhead
\bottomrule\noalign{}
\endlastfoot
(one half log x - 29 log343/400)/log2 & (231, 232) \\
(2 log x + 799 log343/100)/log2 & (997, 998) \\
(9 log x/2 + 999 log343/100)/log2 & (2177, 2178) \\
\end{longtable}

Thus any full \nolinkurl{SizesAt} at this scale must satisfy

\[
                  m>2^{231},\qquad f=2n>2^{998},\qquad k>2^{2177}.
                                                        \tag{SL12}
\]

In particular no choice of m on either fixed f = 42 or f = 168 field can repair this size record at the unchanged scale. This is stronger than excluding m = 2, and its proof uses no code cardinality, no rank-deficient period map, and no estimate on an automorphic form. These bounds quantify the additional rank growth required at this scale. Changing the scale alone is still insufficient at the two fixed degrees by section 5; any proposed degree-and-scale replacement must also revisit the previously fixed weight windows and quantitative tests.

For clarity, two further independent source clauses also remain unmet by the earlier exploratory profile. \nolinkurl{SizeWords} at line 132 caps the word set by 2\^{}(k-3); the entire binary set of length k-2 has twice that size when k \textgreater= 3. This calls for a subcode, not a failure of every permitted code. \nolinkurl{SizeG} at lines 123--125 requires g\_i \textgreater= (k exp L)\^{}8. At L = 343 this is greater than 2\^{}2744, whereas all four actual g\_i used by the J14/J56 examples are less than 2\^{}36. These observations do not alter the independent literal \nolinkurl{g_big} check previously made for those primes.

\subsection{7. A retained literal Host enlargement has a much larger necessary length}\label{a-retained-literal-host-enlargement-has-a-much-larger-necessary-length}

We use the accepted AM-R-023 result, rather than reprove its field and Frame obstruction. If a literal Host field E contains K14 or K56 and has the source degree {[}E:Q{]} = 4n, the degree tower forces 21 \textbar{} n.~AM-R-023, equations R23.8--10, shows that the Frame and full Host inequalities then give

\[
 o\ge21,\qquad e\ge282,\qquad f=2n\ge42\,2^{282}.
                                                        \tag{SL13}
\]

The hypotheses include the literal Host; in particular its Galois realization is abelian. This is a bound for an enlargement retaining the actual field, not for arbitrary new CM fields or modified Host axioms. AM-R-023 separately states that its abstract group bound does not construct a CM field, polarization or full host, and we preserve that distinction.

The logarithmic condition alone now gives

\[
 \begin{gathered}
 39682<\bigl(\log(42\,2^{282})\bigr)^2<39683,\\
                         m\ge79368.
 \end{gathered}
                                                        \tag{SL14}
\]

The certificate proves the first interval by writing the logarithm as 283 log2 + log3 + log7 and applying (SL7). Monotonicity in f and the exact even rounding in (SL1) then prove the second line. This lower bound is sharp for the isolated logarithmic clause at the numerical lower degree in (SL13), not for the complete source record.

The full record is substantially more restrictive. Its exact \nolinkurl{o_tie} clause at \nolinkurl{DataParams₂.lean}, line 131, says o \textless= x\^{}(epsilon/100), equivalently L \textgreater= o\^{}100. Since log x \textgreater{} 52 log L, (SL8) implies

\[
 \begin{aligned}
 \log m
 &\ge\tfrac12\log x-\tfrac{29}{400}\log L\\
 &>\tfrac{10371}{400}\log L
 \ge\tfrac{10371}{4}\log o.
 \end{aligned}
                                                        \tag{SL15}
\]

Here log L \textgreater{} 0, so the middle inequality is strictly justified. For o \textgreater= 21 this proves the first part of (SL2). The same exact series bounds certify

\[
 11388<\frac{10371\log21}{4\log2}<11389,
                                                        \tag{SL16}
\]

which gives m \textgreater{} 2\^{}11388. Since the source also requires m \textless= f \textless= k, those ranks must grow at least as much. The weaker 79,368 floor should therefore not be used as if it were the complete necessary length for this retained-host repair.

\subsection{8. Consequences for the two main questions}\label{consequences-for-the-two-main-questions}

The concrete two-edge mixed family cannot be a literal \nolinkurl{FamP} size instance at any eligible point, even if all unresolved periods, image lattices, coding choices, and analytic ties were supplied. Increasing k alone cannot repair m = 2. At either fixed degree, even changing every rank and the scale cannot repair the literal size predicate. An enlargement retaining K14 or K56 and the full old Host/Scale/size clauses has the much larger necessary length (SL2).

These are source-admissibility obstructions. The actual mixed theta lifts and their local gains, phase identities, global nonvanishing results in their stated scopes, and conditional operator bounds remain valid. None of the inequalities above compares two complete DP44 budgets or proves that an adapted mixed theorem cannot use fewer lifts. A shorter-chain comparison can still be sought after both candidates satisfy the applicable conditions, or under a separately proved replacement for those conditions. For a candidate retaining the literal interface, its pointwise chain length and the unbounded-window growth must both be checked.

This also distinguishes the result from LIFT51 and LIFT53. Those findings carefully retained their eventual link/code thresholds. The present contradiction uses the explicit \nolinkurl{FamP.sizes} field immediately for x \textgreater{} x1, and the direct \nolinkurl{SizesAt} contradiction does not require any family at all.

\subsection{9. References and verification scope}\label{references-and-verification-scope}

Alpöge, L. (n.d.). \emph{Remark}, step 8, associated formal source archive, pinned commit \nolinkurl{7cbf7a0c32fca1badf05931879a37c3f8a2b03b4}. Exact local files and lines are listed in section 2; the additional scale anchors are \nolinkurl{Formalism/LocalRingCardinality/DataParams₂.lean}, lines 9--17, 27--45, and 131, and \nolinkurl{Formalism/SharpRestriction/ElevenLe.lean}, lines 11, 24--49, 123--138.

ALP\_MACHINE. (2026). AM-R-023, \emph{Literal Host rank and Scale obstructions for the repaired CM fields}, equations R23.8--13; AM-R-022 and AM-K044, the fixed scale and admissible gap profiles. These are accepted internal findings used at their stated scopes. No previously accepted field obstruction or theta theorem is presented as a new proof here.

The accompanying \nolinkurl{literal_size_chain_length_certificate.py} uses exact rational logarithm intervals, natural-number rounding, integer powers, and literal source fingerprints. It does not construct a \nolinkurl{SizesAt} record, a Host, a family, a period, or a new theta image; it does not run Lean. Its source hashes and computation timestamp are recorded in the JSON output. Author whole-proof and independent reviews are separate check events.

\clearpage

\section{AM-DO-063 --- actual fractional-lattice volumes close the four-place stratum}\label{am-do-063-actual-fractional-lattice-volumes-close-the-four-place-stratum}

Author: GPT-6 Astra Pro, ALP\_MACHINE /root/definite\_oscillator. 22 November 2026 UTC. Inherited reasoning effort is not separately exposed.

\textbf{Author whole-proof check PASS --- 2026-11-22T10:50:26.422876+00:00.} The final exact certificate passed 468 assertions; independent reviews are separate events.

This finding retains the J56 family, the actual K062 signed kernel, the upper K044 gap, the cutoff H=8192 and every original finite and real normalization. It bounds the remaining scalar-cross part with exactly four nonintegral relative T places by \(2^{-11674}\). Together with the already proved structural cut and four other cases, the complete four-place stratum has absolute contribution less than \(2^{-11671}\).

The new ingredients are the covolume of the actual fractional target lattice, a covering-radius bound for that lattice, and K060's bound on the full selected real compression. Counting ordinary integral cosets separately would lose the saving. The real coefficient is evaluated at the actual near-origin point, not at the standard origin.

\subsection{1. Fixed objects and the exact subset}\label{fixed-objects-and-the-exact-subset}

Use the notation and normalizations of K062(AT14--24), DO059(SC24--37), DO061(TR19--36) and DO062(RU1--13). In particular,

\[
 a=f=168,\quad n=169,\quad m=a-4=164,\quad s=4,\quad H=8192,
 \qquad Q=g_0g_1,
 \tag{FV1}
\]

where \(g_0=22285932761\), \(g_1=64072175441\), and T is the unchanged 996-bit prime of DO040. Write E for the degree-336 CM field. The trace Euclidean norm is

\[
 \|v\|_{\mathrm{tr}}^2=\sum_{\tau\in\Phi}\|\tau v\|^2,
 \qquad
 \langle v,w\rangle_{\mathrm{tr}}
       =\tfrac12\operatorname{Tr}_{E/\mathbf Q}(v^*w).
 \tag{FV2}
\]

The decomposition is the original coordinate decomposition \(E^n=S\oplus W_0\), with S positive of dimension m and \(W_0\) of signature (4,1). No moving basis of a large fixed space is used.

K068's two cuts have total absolute contribution less than \(2^{-1307182}\), over all nonintegral support sizes. Off those cuts, DO062 proves global unipotence for \(4\le s\le8\). Choose the unique unipotent representative h of each central class under discussion and put \(N=h-I\). The only s=4 case not already bounded in DO062 is

\[
 \operatorname{rank}N=2,\qquad N_{SS}=0,\qquad N_{SW_0}\ne0.
 \tag{FV3}
\]

All finite supports, absolute partial traces and real norms are unchanged by this central normalization. We bound the absolute summand, so the original central phase is not assigned a new value.

The carried unipotent calculation gives

\[
 \|(h e_n)_S\|\le\sqrt{H+1},
 \qquad |a(h)|\le
 \frac{2[2(H+1)]^{f(3a/2-1)}T^{s(3a-2)}}
      {Q^{(2a-1)(f-1)}}.
 \tag{FV4}
\]

Here a(h) is precisely K062's normalized genuine finite partial trace; both mixed projectors retain their modulus bound one. The factor 2 in FV4 is the whole finite Euler correction. The finite pencil factor still obeys \(|\rho_v|\le2/(T-2)\) at each of the s nonintegral places. These are DO061(TR23--24,32--35), cited as settled inputs rather than recalculated with a new finite model.

\subsection{2. The actual target over-lattice and its covolume}\label{the-actual-target-over-lattice-and-its-covolume}

Fix a support set \(\Sigma\) of s relative T places, and actual ordered source/target pencil points at those places. Let

\[
 I_\Sigma=\prod_{\mathfrak P\mid\Sigma}\mathfrak P^{-1},
 \qquad N(I_\Sigma)=T^{-2s}.
 \tag{FV5}
\]

The product runs over both CM partners above each relative place, so \(I_\Sigma=\overline{I_\Sigma}\). All those primes have absolute degree one.

At a CM component the pencil acts only on \(W_0\). Its target lattice is

\[
 \mathcal O_v^m\oplus L'_{v,W_0},\qquad
 \varpi\mathcal O_v^5\subset L'_{v,W_0}
       \subset\varpi^{-1}\mathcal O_v^5,\qquad
 v(\det L'_{v,W_0})=0.
 \tag{FV6}
\]

The partner is the paired Hermitian dual lattice. Assemble the target \(W_0\) lattices at \(\Sigma\), and the standard lattice elsewhere, into a global lattice \(\mathcal L\subset E^5\). It is self-dual for \(J_0=\operatorname{diag}(1,1,1,1,-1)\), has determinant ideal of valuation zero everywhere, and satisfies \(\mathcal L\subset I_\Sigma\mathcal O_E^5\).

DO059(SC25) places the actual last column in the following lattice:

\[
 h e_n\in I_\Sigma^m\oplus\mathcal B,
 \qquad \mathcal B=I_\Sigma\mathcal L.
 \tag{FV7}
\]

At places outside \(\Sigma\), h is integral and its last column is in the standard lattice; no assertion about a different pencil there is needed. At a place in \(\Sigma\), the source contains \(\varpi e_n\), whence \(h e_n\in\varpi^{-1}L'\). This proves FV7 at every completion and hence globally.

The exact host discriminant K032(QH15) gives

\[
 \Delta_E=2^{336}7^{280}29^{288}5^{168}13^{168},\qquad
 D_0:=\operatorname{covol}(\mathcal O_E)
 =\frac{\sqrt{\Delta_E}}{2^f}
 =7^{140}29^{144}5^{84}13^{84}.
 \tag{FV8}
\]

These covolumes use FV2. The zero determinant-valuation target property and the exact ideal norm in FV5 imply

\[
 \begin{split}
 \operatorname{covol}(I_\Sigma^m)&=D_0^m T^{-2ms},\\
 \operatorname{covol}(\mathcal B)&=D_0^5 T^{-10s}.
 \end{split}
 \tag{FV9}
\]

In particular, the two blocks together retain the original exact index \(T^{2ns}\). There is no replacement by an arbitrary entry denominator window.

\subsection{3. Covering radii for these particular fractional lattices}\label{covering-radii-for-these-particular-fractional-lattices}

Write \(\mathfrak D_E\) for the different and \(\operatorname{rd}(E)=\Delta_E^{1/(2f)}=2D_0^{1/f}\). The exact integers in FV8 give

\[
 \operatorname{rd}(E)<1500,\qquad T>2^{995}.
 \tag{FV10}
\]

We use the covering-radius transference theorem in the precise Euclidean normalization already checked in LIFT47(TF7):

\[
                  \mu(L)\lambda_1(L^*)\le\dim_{\mathbf R}(L)/2.
 \tag{FV11}
\]

This is Banaszczyk (1993, Theorem 2.2, p.632). We do not substitute a Hermitian dual for the Euclidean dual.

For the scalar fractional ideal in FV5, trace duality gives

\[
 I_\Sigma^*=2\mathfrak D_E^{-1}I_\Sigma^{-1}.
 \tag{FV12}
\]

If y is a nonzero element of this dual, fractional-ideal norms give \[
 |N_{E/\mathbf Q}y|
 \ge 2^{2f}\Delta_E^{-1}T^{2s}.
\] Apply arithmetic--geometric mean to the f squared complex moduli. It follows that

\[
 \lambda_1(I_\Sigma^*)\ge
       \frac{2\sqrt f}{\operatorname{rd}(E)}T^{s/f},
 \qquad
 \mu(I_\Sigma)\le
       \frac{\sqrt f\,\operatorname{rd}(E)}{2T^{s/f}}.
 \tag{FV13}
\]

The m factors form an orthogonal Euclidean product. Choosing nearest points independently in each factor gives \(\mu(I_\Sigma^m)\le\sqrt m\,\mu(I_\Sigma)\). Since \(s=4\), \(T^{s/f}>2^{23}\), \(\sqrt m<13\), and \(\sqrt{8193}<90.516\), we obtain the strict radius bound

\[
 \sqrt f\sqrt{H+1}+\mu(I_\Sigma^m)
 <\sqrt f\left(90.516+\frac{13\cdot1500}{2\cdot2^{23}}\right)
 <\sqrt f\,90.518.
 \tag{FV14}
\]

For the other block, self-duality of \(\mathcal L\) for \(J_0\) gives the Euclidean dual

\[
 \mathcal B^*
   =2\mathfrak D_E^{-1}I_\Sigma^{-1}J_0\mathcal L
   \subset 2\mathfrak D_E^{-1}\mathcal O_E^5.
 \tag{FV15}
\]

The inclusion uses the actual \(\mathcal L\subset I_\Sigma\mathcal O_E^5\), and \(J_0\) is a diagonal unit matrix. This does not treat the indefinite norm as positive: choose one nonzero coordinate of a vector in FV15, use its absolute field norm, and then arithmetic--geometric mean for the positive Euclidean trace norm. Exactly as in FV13 with s=0,

\[
 \lambda_1(\mathcal B^*)\ge\frac{2\sqrt f}{\operatorname{rd}(E)},
 \qquad
 \mu(\mathcal B)\le\frac52\sqrt f\,\operatorname{rd}(E)
                         <3750\sqrt f.
 \tag{FV16}
\]

The real dimension used here is \(10f\), not \(2f\).

There is also a useful exact bound on the full negative column. For any matrix h in \(U(a,1)\), the identity \(hJh^*=J\) gives

\[
 \|h\|_{\mathrm{HS}}^2-n=2(\|h e_n\|^2-1).
\]

Its singular values are \(R,1,\ldots,1,R^{-1}\), with \(1\le R\le H\), as already proved in DO061(TR12). Consequently

\[
 \|h e_n\|^2=\frac{R^2+R^{-2}}2
 \le\frac{H^2+H^{-2}}2<5793^2.
 \tag{FV17}
\]

The \(W_0\) projection has no larger norm. Combining FV16--17, its trace ball enlarged by the covering radius is contained in the ball of radius \(9543\sqrt f\).

\subsection{4. A volume count for all actual last columns}\label{a-volume-count-for-all-actual-last-columns}

For a full Euclidean lattice L, let \(\mathcal V\) be its closed Dirichlet cell. It has volume \(\operatorname{covol}(L)\), lies in the ball of radius \(\mu(L)\), and its translates are disjoint outside boundaries of measure zero. For the finite set of lattice points in a ball of radius R, their cells therefore lie in the ball of radius \(R+\mu(L)\). Comparing volumes proves

\[
 \#\{v\in L:\|v\|\le R\}
 \le\frac{\operatorname{vol}B_{\dim L}(R+\mu(L))}
          {\operatorname{covol}(L)}.
 \tag{FV18}
\]

This argument counts the entire fractional lattice directly. There is no separate packing loss for each integral coset.

The volume of a real \(2d\)-ball of radius r is \(\pi^d r^{2d}/d!\). Apply FV18 to the two blocks of FV7, with FV9, FV14 and FV16--17. Every actual last column is counted by

\[
 \begin{split}
 V_S&=
 \frac{[\pi f(90.518)^2]^{fm}}{(fm)!\,D_0^m}\,T^{2ms},\\
 V_W&=
 \frac{[\pi f(9543)^2]^{5f}}{(5f)!\,D_0^5}\,T^{10s}.
 \end{split}
 \tag{FV19}
\]

The S and \(W_0\) constraints need not be independent on an actual column; the product of the two full counts remains an upper bound. For a fixed full last column, the original isometry fibre has at most

\[
 G_n=\gcd\bigl(|\operatorname{GL}_n(\mathbf F_{g_0})|,
              |\operatorname{GL}_n(\mathbf F_{g_1})|\bigr)
 \tag{FV20}
\]

elements, exactly as in DO059(SC29). There are \(\binom fs(T-1)^{2s}\) support and ordered-pencil choices. Hence the complete number of matrices in FV3 is at most

\[
 \binom fs(T-1)^{2s}G_n V_S V_W.
 \tag{FV21}
\]

No choice of a nonzero witness entry is counted as an additional matrix: it will be used only for a pointwise estimate below.

\subsection{5. The real selected trace supplies the remaining saving}\label{the-real-selected-trace-supplies-the-remaining-saving}

Let \(k_p\) carry the origin to the actual point p of LIFT41. K062(AT14) defines the real factor by

\[
 c(g)=\frac{1}{d_W}
       \operatorname{Tr}_W
       \bigl(P_W\pi(k_p^{-1}g_\infty k_p)|_W\bigr).
 \tag{FV22}
\]

Every selected component carries the isometric PBW metric. The central normalization from g to h changes this by a unit phase. Write \(h_p=k_p^{-1}h_\infty k_p\).

At every CM label, LIFT41(NP12) gives

\[
 \sinh^2r_p=\frac{77Q^4-14Q^2}{T}<2^{-704}.
 \tag{FV23}
\]

If \(x=\sinh^2r_p<1/8\), then \((\sqrt{1+x}+\sqrt x)^2<2\), by squaring the positive inequality \(2\sqrt{x(1+x)}<1-2x\). Thus

\[
       \|k_p\|\,\|k_p^{-1}\|=e^{2r_p}<2.
 \tag{FV24}
\]

For a unipotent u in \(U(a,1)\), write \(R_u=e^{r(u)}\). Since \(\operatorname{tr}u=n\), the singular-value identity gives

\[
 \|u-I\|_{\mathrm{HS}}^2
   =R_u^2+R_u^{-2}-2=(2\sinh r(u))^2.
\] The operator norm is at most the Hilbert--Schmidt norm, so

\[
                  \|u-I\|\le2\sinh r(u).
 \tag{FV25}
\]

Choose any nonzero entry x of \(N_{SW_0}\). It belongs to \(I_\Sigma\) by DO061(TR5), so the product formula gives

\[
                     \prod_{\tau\in\Phi}|\tau x|\ge T^{-s}.
 \tag{FV26}
\]

At each label, FV24--25 imply

\[
 |\tau x|\le\|h_\tau-I\|
    \le\|k_{p,\tau}\|\,\|k_{p,\tau}^{-1}\|
                     \|h_{p,\tau}-I\|
    <2\,(2\sinh r(h_{p,\tau})).
 \tag{FV27}
\]

No Cartan factor in FV27 is zero: a compact unipotent is the identity, which at a faithful embedding would contradict x nonzero.

K060(RC14--15) is a bound on the complete homogeneous compression, and therefore on the actual selected component:

\[
 \|P_W\pi(h_p)|_W\|
 \le\prod_{\tau\in\Phi}
       \frac{12}{\sqrt{2\sinh r(h_{p,\tau})}
                       (2\ell_++q_s+a)^{1/4}}.
 \tag{FV28}
\]

For the unchanged J56 upper endpoint, \(2\ell_++q_s+a\ge2\ell_++a>2^{485}\). Combining FV26--28 and the normalized trace in FV22 proves

\[
 |c(g)|<
 12^f\,2^{f/2}\,2^{-485f/4}\,T^{s/2}.
 \tag{FV29}
\]

The trace bound uses \(|\operatorname{Tr}A|/d_W\le\|A\|\); no selected-dimension factor has been deleted. The finite factor a(g), the real factor c(g), the original centre and all class masses still have their K062 meanings.

\subsection{6. Exact four-place arithmetic}\label{exact-four-place-arithmetic}

Multiply the count FV21, the character bound FV4, the real bound FV29, and the original pencil bound \((2/(T-2))^s\). Replace \(\pi\) in FV19 by the strict rational upper bound \(22/7\). At s=4 the resulting rational number is

\[
 \begin{split}
 \mathcal U_4={}&
 \binom f4(T-1)^8G_n
 \left(\frac2{T-2}\right)^4
 \frac{[(22/7)f(90.518)^2]^{fm}T^{8m}}
      {(fm)!D_0^m}\\
 &\quad\cdot
 \frac{[(22/7)f(9543)^2]^{5f}T^{40}}
      {(5f)!D_0^5}
 \frac{2[2(H+1)]^{f(3a/2-1)}T^{4(3a-2)}}
      {Q^{(2a-1)(f-1)}}
 \frac{12^f2^{f/2}T^2}{2^{485f/4}}.
 \end{split}
 \tag{FV30}
\]

Every exponent is an integer. Exact integer cross-multiplication in the companion certificate gives

\[
                         \mathcal U_4<2^{-11674}.
 \tag{FV31}
\]

For transparency, omitting only the last real factor leaves a bound between \(2^{6017}\) and \(2^{6018}\); that volume estimate alone does not close the stratum. The last factor is less than \(2^{-17692}\). FV31 is computed from the entire rational number, not by subtracting these rounded exponents.

The other four s=4 unipotent cases have the accepted absolute bounds from DO062:

\begin{longtable}[]{@{}lr@{}}
\toprule\noalign{}
Actual case & Absolute bound \\
\midrule\noalign{}
\endhead
\bottomrule\noalign{}
\endlastfoot
Rank two, nonscalar integral block & \(2^{-1021942}\) \\
Rank one, nonscalar integral block & \(2^{-773837}\) \\
Rank one, scalar block & \(2^{-199371}\) \\
Rank two, scalar block, no cross action & \(2^{-436554}\) \\
Rank two, scalar block, nonzero cross action & \(2^{-11674}\) \\
\end{longtable}

The intersection of the two K068 cuts with s=4 has bound \(2^{-1307182}\). These six sets exhaust the nonzero s=4 contribution. Each has bound at most \(2^{-11674}\), so

\[
 \boxed{\displaystyle
  \sum_{s(g)=4}\left|a(g)c(g)\prod_{v\mid T}\rho(g_v)\right|
       <6\,2^{-11674}<2^{-11671}.}
 \tag{FV32}
\]

This is a complete support-stratum bound, not only a plane, fixed-kernel or chosen-entry subset.

\subsection{7. The remaining actual signed sum}\label{the-remaining-actual-signed-sum}

With \(S_{\ge5}\) denoting the unchanged signed short sum over matrices nonintegral at five or more relative T places, the original K062 identity now becomes

\[
 \frac{E_{x,\mathrm{ram}}}
 {w_x\,28C_\infty d_Rd_W((T-2)/(T-1))^f}
     =1+S_{\ge5}+\varepsilon_5,
 \tag{FV33}
\] \[
 |\varepsilon_5|
 <2^{-1999}+2^{-104546}+2^{-101550}
                 +2^{-211446}+2^{-11671}<2^{-1998}.
 \tag{FV34}
\]

Previously bounded planes, character cuts and fixed-kernel subsets can be removed further by taking their set differences with these four complete strata, then adding their checked absolute bounds. One must not charge an overlapping set twice.

The s\textgreater=5 signed remainder is not bounded in full here. No all-T Mram positivity, actual terminal nonzero value, source-link compatibility or net DP44 gain is claimed. In particular LIFT51--52's independent full-space source-link obstruction is unchanged. This finding changes neither the host, finite primes, projector, upper gap nor the actual class masses. It makes no new J14 assertion.

\subsection{8. Sources and verification scope}\label{sources-and-verification-scope-7}

The actual kernel, supports and metrics are the checked K062(AT14--24), K032(QH15), LIFT41(NP12), LIFT47(TF1--7), K060(RC14--15), DO059(SC24--30), DO061(TR19--36), DO062(RU1--13) and K068(LR1--22). The external transference and selected-compression theorems are cited as established inputs with their exact normalizations; the new ideal-dual, volume, conjugation and counting arguments are proved above.

Banaszczyk, W. (1993). New bounds in some transference theorems in the geometry of numbers. \emph{Mathematische Annalen, 296}, 625--635. \url{https://doi.org/10.1007/BF01445125}. The precise input is Theorem 2.2, printed p.632, with the Euclidean definitions on pp.625--626. The original theorem and its factor 1/2 were directly checked in LIFT47 and its independent reviews.

Haagerup, U., \& Schlichtkrull, H. (2014). Inequalities for Jacobi polynomials. \emph{The Ramanujan Journal, 33}, 227--246. \url{https://doi.org/10.1007/s11139-013-9472-4}. Theorem 1.1, p.228, supplies the uniform constant C\textless12 through the complete graded-ladder argument of K060. That argument, not an unrestricted unitary-operator bound, is the input FV28.

The companion \nolinkurl{definite_oscillator_fractional_volume_certificate.py} checks the exact host integers, both fractional-lattice radii, paired split-pencil dual identities, unipotent Frobenius and negative-column identities, small fractional-ball volume normalizations, the complete rational product FV30 and the retained error. It does not enumerate the actual rational support, evaluate a period, or establish the sign of the remaining sum.

\subsection{Author verification record}\label{author-verification-record-14}

At 2026-11-22T10:50:26.422876+00:00, I completed a whole read and reconstruction of FV1--34 and all eight sections, together with the entire final companion certificate. I checked the actual paired target-lattice decomposition, half-trace covolumes, both different-ideal Euclidean duals and covering-radius dimensions, the Dirichlet-cell volume count, finite isometry fibre, the negative-column and unipotent Frobenius identities, the actual near-point conjugation and normalized selected trace, all exact products and the complete-stratum versus residual-subset bookkeeping. The source-size incompatibility and the s\textgreater=5 signed gate remain separate.

The final replay passed 468 exact assertions at 2026-11-22T10:48:31.642333+00:00, including 52 paired pencil models, 32 unipotent models, 27 fractional-ball count normalizations and the strict 11,674-bit full-product comparison. The preliminary arithmetic replay had already obtained the same main margin; final bookkeeping explicitly included the K068 cut in the complete four-place union and removed two irrelevant numerical assertions. No promoted mathematical bound was changed. No actual rational-support enumeration, period, all-T sign or Lean build is claimed.

\clearpage

\section{AM-K070 --- characteristic poles count only nonfixed unequal-pencil transitions}\label{am-k070-characteristic-poles-count-only-nonfixed-unequal-pencil-transitions}

Author: GPT-6 Astra Pro, ALP\_MACHINE \nolinkurl{/root/kummer_extension}. 22 November 2026 UTC. Inherited reasoning effort is not separately exposed.

\textbf{Author whole-proof check PASS --- 2026-11-22T10:46:32.422947+00:00.}

The unchanged K062 finite kernel distinguishes the total number s of nonintegral T places from a smaller number t defined below. On the K068 residual, t\textless=8 already forces global unipotence, even when s is larger than8. At each of the t exceptional places the exact local factor has an additional denominator T−1. This denominator is not asserted at arbitrary places with no fixed actual lattice.

\subsection{1. The exact local alternatives}\label{the-exact-local-alternatives}

Let O be the local DVR at a selected relative T place, q=T, m=q−1, and retain K062's actual full pencil

\[
\mathscr L=\{L:M\subset L\subset N,\ \dim_k(L/M)=1\}
          \simeq\mathbf P^1(k).
\tag{PT1}
\]

The actual labelled subfamily D omits exactly the two lines infinity and −1, and includes the standard lattice Lambda at the line0. For the actual supported local matrix B put

\[
C(B)=\#\{(d,c)\in D^2:B L_c=L_d\},\qquad
F(B)=\#\{c\in D:B L_c=L_c\},
\quad
\rho(B)=\frac{F(B)-C(B)/m}{m-1}.
\tag{PT2}
\]

The complete support has C(B)\textgreater0. The elementary pencil argument of K062(AT7--10) gives the following three disjoint cases:

\begin{longtable}[]{@{}
  >{\raggedright\arraybackslash}p{(\columnwidth - 4\tabcolsep) * \real{0.3333}}
  >{\raggedright\arraybackslash}p{(\columnwidth - 4\tabcolsep) * \real{0.3333}}
  >{\raggedright\arraybackslash}p{(\columnwidth - 4\tabcolsep) * \real{0.3333}}@{}}
\toprule\noalign{}
\begin{minipage}[b]{\linewidth}\raggedright
Case
\end{minipage} & \begin{minipage}[b]{\linewidth}\raggedright
Exact incidence condition
\end{minipage} & \begin{minipage}[b]{\linewidth}\raggedright
Invariant full O-lattice
\end{minipage} \\
\midrule\noalign{}
\endhead
\bottomrule\noalign{}
\endlastfoot
I & B preserves the full pencil & M and N \\
II & The full pencils differ and F(B)=1 & The unique fixed L\_c \\
III & The full pencils differ and F(B)=0 & None supplied by the incidence alone \\
\end{longtable}

In Cases II and III one has C(B)=1. Two distinct lattices common to the two full pencils would have the same intersection M and sum N, forcing equality of the full pencils. This proves the asserted bound for unequal pencils, without imposing it on Case I.

Case I preserves M because M is the intersection of all members of the full pencil; similarly it preserves N. In Cases I and II, B is therefore an automorphism of a full O-lattice, and all its characteristic coefficients are integral. The inverse has the same property. The conjugate CM component is inverse adjoint and hence also has integral characteristic coefficients. No integrality of B in the standard Lambda frame is inferred.

The exact weights in Cases II and III are

\[
\rho(B)=\frac1{T-1}\quad\text{in Case II},
\qquad
\rho(B)=-\frac1{(T-1)(T-2)}\quad\text{in Case III}.
\tag{PT3}
\]

Case III is necessarily nonintegral in the standard frame: the actual incidence makes det(B) a local unit, so an integral B would fix the included Lambda, making F(B)\textgreater=1. The inverse preserves the case, because it exchanges the unique source and target and preserves F and C.

It would be false to replace Case III by the sole condition F(B)=0. For the full-pencil projective translation c maps to c+b, b!=0, the only fixed line is the omitted infinity. Here F=0, C=T−2 and

\[
                         \rho(B)=-\frac1{T-1}.
\tag{PT4}
\]

The matrix 1+b varpi\^{}-1 zv from K062(AT3), with vz=0, realizes this translation. It moves the standard lattice and is nonintegral. They are nevertheless in Case I and have integral characteristic coefficients. Thus neither a uniform C\textless=2 nor the extra denominator in PT3 holds for all nonintegral matrices.

\subsection{2. The number of possible characteristic-pole pairs}\label{the-number-of-possible-characteristic-pole-pairs}

For a rational matrix in the unchanged K062 support let t(g) be the number of relative T places in Case III. Then

\[
                        0\le t(g)\le s(g)\le f.
\tag{PT5}
\]

DO061(TR6--7) proves that every characteristic coefficient of the actual matrix has pole order at most1 at either CM component of a nonintegral T place. Section1 removes all possible poles in Cases I and II. Away T the matrix is integral. Consequently every coefficient of h, and every coefficient of h−I, is integral away at most2t(g) CM primes and has pole order at most1 at those primes. Translating the polynomial by1 uses integer linear combinations and preserves the assertion. The determinant-one central representative h differs by a global unit, so it has exactly the same local incidences and coefficient-integrality conclusions.

This statement concerns characteristic coefficients. Individual entries may have double poles in Case II. The equal-label example DO062(RU13) demonstrates why coefficient integrality must not be promoted to entry integrality.

The normalized finite product also satisfies the exact improvement

\[
\prod_{v\mid T}|\rho(B_v)|
\le
\left(\frac2{T-2}\right)^{s(g)}
\left(\frac1{2(T-1)}\right)^{t(g)}.
\tag{PT6}
\]

At the nonintegral places outside Case III use K062's original bound 2/(T−2); at each Case III place use PT3. At integral places use the unchanged contraction bound1. No local coefficient phase is discarded in the kernel identity; PT6 is only an absolute upper bound.

\subsection{3. Global unipotence with t\textless=8, at any total support size}\label{global-unipotence-with-t8-at-any-total-support-size}

Retain the actual J56 K068 residual and H=8192. Its positive fixed space has dimension at least126, its global rank is at most43, and all334 unmixed G reductions are unipotent. Use the same dimension43 positive padding as DO062(RU2--4), without requiring a bounded arithmetic basis. For N=h−I this gives, at every realization,

\[
|e_1(N)|\le B_1=\frac{67796993}{8192},\qquad
|e_2(N)|\le B_2=\frac{5598265427}{8192}.
\tag{PT7}
\]

By Section2, a nonzero coefficient would now imply

\[
                  Q^{2(f-1)}\le B_j^{2f}T^{2t(g)}.
\tag{PT8}
\]

This is exactly the DO062 norm comparison with s replaced by t. At t=8 its strict inverse binary margins are3171 for j=1 and1032 for j=2. Since T\textgreater1, both contradictions remain valid for every 0\textless=t\textless=8. Therefore

\[
\boxed{\quad t(g)\le8\ \Longrightarrow\
 h\text{ is globally unipotent and }\operatorname{rank}_E(h-I)\le2.
 \quad}
\tag{PT9}
\]

Indeed e1(N)=e2(N)=0 gives tr(h)=tr(h²)=n by Newton's identity. The already checked signature-one two-trace lemma DO061§2 proves unipotence, and its unitary normal form gives rank at most2. The conclusion is algebraic over E because it holds at a faithful complex realization. No upper bound on the total support s is used in PT9. The existing negative second-coefficient margin at t=9 is retained; this test does not settle that case.

\subsection{4. The exact local weights after global unipotence}\label{the-exact-local-weights-after-global-unipotence}

Suppose h is globally unipotent. At a Case I place its action on N/M is a two-dimensional unipotent linear map. If its projective action is the identity then rho=1, and it fixes Lambda, so the place is integral. Otherwise it has exactly one fixed projective line.

Write E0 for the two omitted lines and \(k=|E_0\cap B E_0|\). Since the full pencil has q+1 members,

\[
                         C=q-3+k=m-2+k.
\tag{PT10}
\]

If the unique fixed line lies in E0, then k=1. The other omitted line cannot map to the fixed line or itself, so this is exact; F=0 and rho=−1/(T−1). If the fixed line lies in D, then F=1 and k is either0 or1. It cannot be2 because a nontrivial unipotent projective map has order T\textgreater2 and cannot permute two nonfixed points. In this case rho is respectively2 or1 divided by (T−1)(T−2).

Together with PT3, the complete nonidentity local table is

\begin{longtable}[]{@{}
  >{\raggedright\arraybackslash}p{(\columnwidth - 2\tabcolsep) * \real{0.4286}}
  >{\raggedleft\arraybackslash}p{(\columnwidth - 2\tabcolsep) * \real{0.5714}}@{}}
\toprule\noalign{}
\begin{minipage}[b]{\linewidth}\raggedright
Actual local situation
\end{minipage} & \begin{minipage}[b]{\linewidth}\raggedleft
rho
\end{minipage} \\
\midrule\noalign{}
\endhead
\bottomrule\noalign{}
\endlastfoot
Unequal pencils, unique fixed edge & 1/(T−1) \\
Unequal pencils, unique off-diagonal edge & −1/{[}(T−1)(T−2){]} \\
Equal pencils, unique fixed projective line omitted & −1/(T−1) \\
Equal pencils, unique fixed projective line actual & (2−k)/{[}(T−1)(T−2){]}, k=0 or1 \\
\end{longtable}

These are signed local identities. In particular an omitted fixed line can still give the larger negative magnitude at arbitrarily many places. Neither PT9 nor the table supplies positivity of the entire global summand, whose G and real factors are unchanged.

\subsection{5. Scope and verification}\label{scope-and-verification-3}

The result refines the actual residual description by the two integers (s,t). It does not alter the set of allowed scalar bases, the genuine G projectors, the two omitted pencil lines, the real centre, class masses, selected metrics, or the long-tail error. The Case III count and all displayed classifications are inverse stable.

No complete absolute bound for the remaining high-s unipotents is proved here. PT6 is an exact additional factor that can be used in a future count with the same incidence classification. It is not a uniform replacement of the older bound without tracking t. No all-T Mram sign, source LinkS repair, final pairing or DP44 inequality is asserted.

The only inputs are the exact actual pencil and partial-permutation proof K062(AT3--12), the coefficient-pole theorem DO061(TR6--7), and the checked trace/norm/unitary arguments DO062(RU2--6). The new steps are the invariant-lattice separation, the count t in PT8--9, and the exact unipotent omitted-line table. The companion finite certificate checks projective-line incidence counts, projective unipotents, the coefficient norm margins, and the strict local factor identity. It does not enumerate the global rational support.

\subsection{Author verification record}\label{author-verification-record-15}

At 2026-11-22T10:46:32.422947+00:00, I completed a whole saved proof read of PT1--10 and all five sections, together with the complete certificate. I checked the exact actual full/unequal-pencil alternatives, characteristic integrality on invariant lattices at both CM components, the only possible coefficient-pole pairs, the strict t\textless=8 norm contradiction, every signed local unipotent weight and inverse invariance. The certificate passed 51923 exact assertions at 2026-11-22T10:45:48.039561+00:00, including15,744 invertible projective matrix models and1,600 nonidentity projective unipotents. The full-pencil translation counterexample is retained explicitly; no uniform C\textless=2 or extra denominator is asserted outside Case III. No global rational sum, period, source link or Lean build was computed.

\clearpage

\section{ALP-LIFT-55 --- universal scale and chain floors in the literal size interface}\label{alp-lift-55-universal-scale-and-chain-floors-in-the-literal-size-interface}

Author: GPT-6 Astra Pro, ALP\_MACHINE \nolinkurl{/root/lift_accounting}, 22 November 2026 UTC. Reasoning effort is not separately exposed.

\textbf{Status.} Author whole-proof check PASS at 2026-11-22T10:56:31.207358+00:00. Exact certificate: 94 assertions PASS. Independent review is separate. No Lean build is claimed.

\subsection{1. Scope and main conclusion}\label{scope-and-main-conclusion}

Let D, L be any data satisfying the literal \nolinkurl{SizesAt} from the source archive pinned at commit \nolinkurl{7cbf7a0c32fca1badf05931879a37c3f8a2b03b4}. Use x, epsilon, L = x\^{}epsilon, b = exp L for its Scale parameters, and m, f, k for its \nolinkurl{Ranks} parameters. Here the Roman L in formulas denotes Scale.L, not the complete \nolinkurl{SizeData₂} record.

The following universal necessary bounds hold:

\[
 L>2^{300},\qquad
 m>2^{31113/4}>2^{7778},\qquad
 f>2^{33598},\qquad
 k>2^{73197}.
                                                        \tag{US1}
\]

In particular, the unchanged L = 343 profile cannot satisfy full \nolinkurl{SizesAt}, even after changing all its ranks, CM degree, and finite primes. These are obstructions to the literal size interface. They do not refute the accepted abstract \nolinkurl{Ranks} construction at smaller even counts, the actual mixed theta kernels already constructed, or a separately proved replacement for these size hypotheses. They also do not assert that \nolinkurl{SizesAt} is inconsistent at all sufficiently large scales.

The weaker elimination in section 2 is included to identify precisely which clauses are responsible. The stronger bounds (US1) use \nolinkurl{SizeG}; no additional arithmetic or analytic theorem is invoked.

\subsection{2. Elimination using SizeM, SizeN, SizeK and the first SizeKRel clause}\label{elimination-using-sizem-sizen-sizek-and-the-first-sizekrel-clause}

The source \nolinkurl{ElevenLe.lean}, lines 11 and 24--37, supplies

\[
 \begin{aligned}
 m&\ge x^{1/2-29\epsilon/400},\\
 f=2n&\ge2x^{2+799\epsilon/100},\\
 k&\le x^{9/2+1001\epsilon/100},\\
 k&\ge x^{\epsilon/100}m^5 fLx^\epsilon.
 \end{aligned}
                                                        \tag{US2}
\]

The first line is the exact \nolinkurl{SizeM} simplification already checked in LIFT54. All factors are positive. Substitute its first two lines into its last line, and use L = x\^{}epsilon:

\[
 k\ge2x^{9/2+(3855/400)\epsilon}.
                                                        \tag{US3}
\]

Indeed the epsilon coefficient is -5·29/400 + 799/100 + 1/100 + 2 = 3855/400. The upper epsilon coefficient in \nolinkurl{SizeK} is 1001/100 = 4004/400. Dividing the two positive inequalities for k gives

\[
                         2\le L^{149/400},
 \qquad L\ge2^{400/149}.
                                                        \tag{US4}
\]

No strictness is asserted in (US4). The factor 2 from f = 2n is essential.

The Scale conditions x \textgreater{} 1 and 0 \textless{} epsilon \textless{} 1/52 give log x \textgreater{} 52 log L \textgreater{} 0. Consequently the three lower size inequalities imply

\[
 \begin{aligned}
 \log m&>\frac{10371}{400}\log L,\\
 \log f&>\log2+\frac{11199}{100}\log L,\\
 \log k&>\frac{24399}{100}\log L.
 \end{aligned}
                                                        \tag{US5}
\]

The last line uses the \textbf{lower} \nolinkurl{SizeK} inequality, with coefficient 999/100, rather than its upper inequality. Thus (US4) alone already yields

\[
 m>2^{10371/149}>2^{69},\qquad
 f>2^{44945/149}>2^{301},\qquad
 k>2^{97596/149}>2^{655}.
                                                        \tag{US6}
\]

These weaker floors are exact diagnostic consequences of the stated subset of clauses, not independently constructed size records.

\subsection{3. The stronger SizeG elimination}\label{the-stronger-sizeg-elimination}

Take i = j = 0 in \nolinkurl{SizeG}, defined at \nolinkurl{ElevenLe.lean}, lines 123--125. It gives a positive integer g with

\[
             (k\exp L)^8\le g,
             \qquad \log g\le x^{\epsilon/100}L=L^{101/100}.
                                                        \tag{US7}
\]

The source \nolinkurl{Ranks} conditions imply k \textgreater= m + 3 \textgreater= 5; in fact k is even and hence at least 6. Thus log k \textgreater{} 0. Taking logarithms of the positive lower bound in (US7) proves

\[
           8(L+\log k)\le\log g\le L^{101/100}.
                                                        \tag{US8}
\]

Divide by L \textgreater{} 0. Since log k \textgreater{} 0, this yields

\[
                L^{1/100}>8,\qquad L>8^{100}=2^{300}.
                                                        \tag{US9}
\]

This uses both sides of \nolinkurl{SizeG}: choosing a much larger auxiliary prime cannot evade its simultaneous logarithmic upper bound. The relation between the two g\_i is not needed.

Substitution in (US5) gives

\[
 \begin{aligned}
 \log_2m&>300\frac{10371}{400}=\frac{31113}{4},\\
 \log_2f&>1+300\frac{11199}{100}=33598,\\
 \log_2k&>300\frac{24399}{100}=73197.
 \end{aligned}
                                                        \tag{US10}
\]

This proves (US1). All constants and strict inequalities are rational or integer consequences of the displayed definitions. No numerical approximation to exp or log is required.

\subsection{4. Family and main-question interpretation}\label{family-and-main-question-interpretation}

\nolinkurl{FamP.sizes} supplies \nolinkurl{SizesAt} for every x \textgreater{} x1 and every actual fibre index, as checked in LIFT54 from \nolinkurl{ParConsts.lean}, line 44. Therefore (US1) applies immediately at all such eligible points. It is not conditional on the later thresholds of the separate \nolinkurl{ChosenCode}, \nolinkurl{LinkS} or \nolinkurl{LinkIRR} links.

In particular, the formerly explored L343 two-edge family cannot become a literal source instance by repairing only its degree, ranks or finite primes. Its scale must change as well. The larger retained-host floor m \textgreater{} 2\^{}11388 in LIFT54 still applies under that finding's additional o \textgreater= 21 and \nolinkurl{o_tie} hypotheses. It is not replaced by the weaker universal bound here.

Neither conclusion computes an actual changed-family DP44 difference. A new mixed construction may require a newly proved size interface, while a construction retaining the literal interface must meet these necessary floors and all its remaining requirements. A comparison between two sufficiently large admissible lengths remains logically open. The purpose of this finding is to exclude an unsupported small-count promotion, not to turn an interface condition into a universal obstruction to theta lifting.

\subsection{5. Exact locators and checks}\label{exact-locators-and-checks}

Alpöge, L. (n.d.). \emph{Remark}, step 8, associated source archive, pinned commit \nolinkurl{7cbf7a0c32fca1badf05931879a37c3f8a2b03b4}:

\begin{itemize}
\tightlist
\item
  \nolinkurl{Formalism/SharpRestriction/ElevenLe.lean}, lines 11, 24--37, 123--138: the precise size inequalities and their inclusion in \nolinkurl{SizesAt}.
\item
  \nolinkurl{Formalism/LocalRingCardinality/DataParams₂.lean}, lines 9--17 and 27--45: Scale, L, b and delta1.
\item
  \nolinkurl{Formalism/LocalRingCardinality/Params.lean}, lines 6--22 and 28--30: rank parity, lower rank and f = 2n.
\item
  \nolinkurl{Formalism/SharpRestriction/ParConsts.lean}, lines 31 and 44: eligible family points.
\end{itemize}

LIFT54 is used for its exact natural-rank and family-quantifier audit; AM-R-023 is used only for the separately identified retained-host consequence. The new proof above is the direct elimination of the displayed size inequalities.

The companion \nolinkurl{universal_size_scale_certificate.py} verifies every rational exponent, exact power, fixed-profile exclusion, and source fingerprint. It does not construct or numerically search a full source record, run Lean, evaluate an automorphic form, or compare DP44 costs. Author whole-proof and independent check events are recorded separately.

\clearpage

\section{AM-K071 --- actual pencil types control nilpotent poles and weighted column counts}\label{am-k071-actual-pencil-types-control-nilpotent-poles-and-weighted-column-counts}

Author: GPT-6 Astra Pro, ALP\_MACHINE \nolinkurl{/root/kummer_extension}. 22 November 2026 UTC. Inherited reasoning effort is not separately exposed.

\textbf{Author whole-proof check PASS --- 2026-11-22T11:06:07+00:00.} The final certificate passed 62,067 exact assertions; independent reviews are separate events.

This finding retains the unchanged J56 upper-endpoint mixed family, its actual K062 signed kernel, H=8192, the actual near-origin point and all finite projectors. It refines the local types of K070 rather than changing the scalar packet. In particular, a full-pencil-preserving unipotent has integral cross blocks and integral nilpotent square. These facts remove large subsets at every total support size, not only at small support.

The remaining five-place terms have a particularly explicit form: every nonintegral place has unequal full pencils, and at least one is a unique fixed incidence. Terms supported on the fixed five-coordinate block must have a unique fixed incidence at all five places. The remaining signed sum is not proved positive.

\subsection{1. Fixed notation and the three local pole bounds}\label{fixed-notation-and-the-three-local-pole-bounds}

Use the actual coordinate decomposition and normalization of DO061(TR3--5), DO063(FV1--3) and K070(PT1--5):

\[
 a=f=168,\quad n=169,\quad m=164,\quad H=8192,\qquad
 E^n=S\oplus W_0,
 \tag{NB1}
\]

where S is positive of dimension m and W0 has signature (4,1). The primes, field and T are exactly those of K032/DO040. At a selected split T place, the full pencil is

\[
 \mathscr L=\{L:\mathcal M\subset L\subset\mathcal N,
                         \ \dim_k L/\mathcal M=1\},
 \qquad \mathcal N/\mathcal M\simeq k^2.
 \tag{NB2}
\]

Its actual subfamily D has T−1 members, obtained by omitting infinity and −1, and includes the standard lattice Lambda. The pencil changes only W0; hence

\[
 S_{\mathcal O}\subset\mathcal M\subset\Lambda
       \subset\mathcal N\subset\varpi^{-1}\Lambda,
 \qquad \mathcal N_S=S_{\mathcal O}.
 \tag{NB3}
\]

At a supported place use the exact disjoint types of K070: I means preservation of the full pencil; II means unequal full pencils with one fixed actual lattice; III means unequal full pencils with a single off-diagonal incidence. Let s be the total number of nonintegral relative T places, u the number of type II places, t the number of type III places and v=u+t. Thus s−v counts the nonintegral type I places. The classification is the same at the conjugate CM component, under the inverse-adjoint paired-lattice identification, and is invariant under inversion.

For now suppose the uniquely normalized central representative h is globally unipotent, and put Z=h−I. The signature-one normal form gives Z\^{}3=0 and rank Z at most2. In particular

\[
                    Z^2=h+h^{-1}-2I.
 \tag{NB4}
\]

\textbf{Lemma NB5.} At either CM component the following are upper bounds for entry pole orders. A zero means integrality in the original standard coordinate frame.

\begin{longtable}[]{@{}lrrr@{}}
\toprule\noalign{}
Actual type & Full Z & Both S/W0 cross blocks of Z & Full Z² \\
\midrule\noalign{}
\endhead
\bottomrule\noalign{}
\endlastfoot
I: full pencil preserved & 1 & 0 & 0 \\
II: unequal pencils, unique fixed incidence & 2 & 1 & 2 \\
III: unequal pencils, unique off-diagonal incidence & 1 & 1 & 1 \\
\end{longtable}

At a standard integral place every entry in the table is integral.

\textbf{Proof.} In type I, h and h\^{}−1 preserve the intersection M and sum N of the full pencil. Thus h Lambda is contained in N, proving the full pole-one bound. Its S projection is integral by NB3, proving integrality of the S-by-W0 block. Also h(S\_O) is contained in M and hence Lambda, which proves integrality of the other cross block.

The nilpotent operator Z acts on the two-dimensional vector space N/M. Its square on that quotient is zero: a nilpotent linear map on a two-dimensional space has square zero. Consequently Z² N is contained in M. Since Lambda is contained in N and M in Lambda, Z² Lambda is contained in Lambda. This proves full entry integrality of Z², not only integrality of its characteristic coefficients.

For type III write the actual local incidence as h=h\_d k h\_c\^{}−1, where c−d is a residue unit and k is integral. The precise trace identity DO062(RU12) is

\[
 n=\operatorname{tr}h
   =\operatorname{tr}k+(d-c)\varpi^{-1}q_*kz.
\]

It forces q\_*kz to be divisible by the uniformizer and removes the sole possible double-pole term of h. Reversing the pair gives the same result for h\^{}−1. NB4 then proves the pole-one bound for Z². The paired inverse-adjoint argument gives both CM components.

In type II the original two pencil factors give full entry pole at most2 for h and h\^{}−1. NB4 gives the same for Z². The cross blocks have only one nontrivial pencil factor, so their pole is at most1 by DO061(TR5). At an integral place the unit determinant makes the inverse integral. These observations prove all entries of the table. QED.

The last column of NB5 cannot in general be improved in type II. In split-local dimension3 put R=e1 e3\^{}t, h\_c=I+varpi\^{}−1 R and

\[
 Z_0=\begin{pmatrix}0&0&0\\1&0&0\\1&1&0\end{pmatrix},
 \qquad h=h_c(I+Z_0)h_c^{-1}.
\]

This h fixes L\_1, is unipotent and has

\[
 (h-I)^2=
 \begin{pmatrix}
 \varpi^{-1}&0&-\varpi^{-2}\\
 0&0&0\\
 1&0&-\varpi^{-1}
 \end{pmatrix}.
 \tag{NB6}
\]

It does not preserve N and hence not the full pencil. It is a genuine type II local example, not an assertion that this matrix occurs as a global supported term in the fixed family.

\subsection{2. Incidence weights cancel most ordered-pair costs}\label{incidence-weights-cancel-most-ordered-pair-costs}

Write b=T−1 in this section only. K070's exact local unipotent table implies that every nonintegral type I factor has modulus at most1/b; type II has rho=1/b and type III has rho=−1/{[}b(b−1){]}. An identity projective action in type I would fix Lambda and would be integral, so it is absent from the nonintegral type I count.

For a prescribed local type, the number of available ordered actual source/target pairs is bounded by b², b, b(b−1), respectively. Multiplying these counts by the corresponding modulus bound gives

\[
 \begin{array}{c|ccc}
 \text{type}&I&II&III\\ \hline
 \text{pair count times }|\rho|&b&1&1.
 \end{array}
 \tag{NB7}
\]

For fixed s,v,u, supports and type sets therefore have the weighted enumeration factor

\[
 W_{s,v,u}
 =\binom fs\binom sv\binom vu (T-1)^{s-v}.
 \tag{NB8}
\]

This is an upper factor for the weighted sum, not a claim that the number of matrices is at most NB8. Fix a support, its type partition and an ordered pair at every place; apply the uniform actual-lattice matrix count to that choice and its pointwise rho bound, then sum. Each actual matrix is represented by at least one such ordered-pair choice, often by more in type I, which only increases the upper bound. Both binomial type-position factors are retained. No new finite projector is used, and NB7 does not assign a sign to the global product.

\subsection{3. Global pole ideals and the genuine G character}\label{global-pole-ideals-and-the-genuine-g-character}

Let I\_U be the product of inverse primes at both CM partners over a set U of relative T places. NB5 gives the following global pole budgets:

\[
 \begin{array}{c|c}
 \text{algebraic entry}&\text{sum of pole orders over relative places}\\
 Z& s+u\\
 Z_{SW_0}, Z_{W_0S}&v\\
 Z^2&v+u=t+2u.
 \end{array}
 \tag{NB9}
\]

The last two lines use integrality at every type I place. Characteristic coefficient poles are still controlled by t alone, as in K070. On the K068 residual, that accepted comparison proves global unipotence whenever t\textless=8, independently of s. Every subset explicitly removed below meets this condition; no new unipotence claim for t\textgreater=9 is made.

The field-norm witness argument DO061(TR23--24) works with nonuniform pole orders: a nonzero entry x, all of whose conjugates have modulus at most M and whose paired pole-order sum is d, bounds the product of bad selected G residue sizes by M\^{}f T\^{}d.~The rank-drop condition vanishes at both CM partners. The exact Green factors and their total Euler correction less than2 are unchanged, including the two mixed partial traces bounded only by1.

Put K=2(H+1). For rank2, A=h\_SS=I and nonzero cross action, the same witnesses as DO061(TR34) have bounds sqrt(K) for a nonzero cross entry of Z and K for a nonzero entry of Z². NB9 therefore gives

\[
 |a(h)|\le
 \frac{2K^{f(3a/2-1)}
                 T^{av+(a-1)(v+u)}}
      {Q^{(2a-1)(f-1)}}.
 \tag{NB10}
\]

For the rank1 or rank2 terms supported entirely on the fixed W0 block, the corresponding bounds are

\[
 \begin{split}
 B_1(s,v,u)&=
 \frac{2(H+1)^{af}T^{a(s+u)}}{Q^{a(f-1)}},\\
 B_2(s,v,u)&=
 \frac{2(H+1)^{af}K^{(a-1)f}
                 T^{a(s+u)+(a-1)(v+u)}}
      {Q^{(2a-1)(f-1)}}.
 \end{split}
 \tag{NB11}
\]

The first bound is independent of t apart from its incidence factor and the condition ensuring global unipotence. These are bounds on the actual normalized G partial trace, not on a newly chosen finite character.

\subsection{4. The S-column lattice loses all type I denominators}\label{the-s-column-lattice-loses-all-type-i-denominators}

Fix a support and type partition, and write U for the v unequal-pencil places. At type I, the S part of h e\_n is integral by NB5. At type II or III it has pole at most1. Thus the actual last-column inclusion improves DO063(FV7) to

\[
                 h e_n\in I_U^m\oplus\mathcal B_\Sigma,
 \tag{NB12}
\]

where Sigma is the full support of size s and B\_Sigma=I\_Sigma L is the very same paired target lattice of DO063. The five-coordinate block retains every type I denominator. No arithmetic change of the large S-coordinate basis is used.

Retain the positive half-trace Euclidean norm, exact

\[
 D_0=7^{140}29^{144}5^{84}13^{84},\qquad
 \operatorname{rd}(E)<1500.
\]

The ideal-dual argument DO063(FV12--16), with s replaced by v only in the S factor, gives its covolume D0\^{}m T\^{}(−2mv) and covering radius at most sqrt(mf) rd(E)/(2 T\^{}(v/f)). Its actual S-column radius is sqrt(f(H+1)). If r\_v is any positive rational with r\_v\^{}f\textless=T\^{}v, then the enlarged radius is strictly less than sqrt(f) R\_v, where

\[
 R_v=\frac{90516}{1000}+\frac{9750}{r_v}.
 \tag{NB13}
\]

Here sqrt(m)\textless13, sqrt(H+1)\textless90.516 and rd(E)\textless1500. Also, in each coset of O\_E\^{}m in I\_U\^{}m, the minimum positive Euclidean distance is at least sqrt(f), by the norm of one nonzero integral coordinate. Ordinary disjoint-ball packing therefore gives fewer than 183\^{}(2fm) points per coset. The exact number of cosets is T\^{}(2mv).

Combining this packing estimate with the fractional-lattice volume estimate gives the actual S-column bound

\[
 \begin{split}
 \#S\text{-columns}&\le A_v T^{2mv},\\
 A_v&=\min\left\{183^{2fm},
  \frac{[(22/7)f R_v^2]^{fm}}{(fm)!D_0^m}\right\}.
 \end{split}
 \tag{NB14}
\]

The bound 22/7 is used only to make the numerical certificate rational. For v=0 the packing alternative is smaller. For v=1,\ldots,5 the checked volume alternative is smaller. The rational choices are

\begin{longtable}[]{@{}rr@{}}
\toprule\noalign{}
v & r\_v \\
\midrule\noalign{}
\endhead
\bottomrule\noalign{}
\endlastfoot
0 & 1 \\
1 & 243/4 \\
2 & 3688 \\
3 & 224256 \\
4 & 13631488 \\
5 & 828375040 \\
\end{longtable}

Each r\_v\^{}f\textless=T\^{}v is checked by exact integer arithmetic. For the W0 factor, DO063's full negative-column radius and covering radius are unchanged. Put

\[
               B=\frac{[(22/7)f\,9543^2]^{5f}}
                               {(5f)!D_0^5}.
 \tag{NB15}
\]

Its column count is at most B T\^{}(10s). The fixed-column isometry fibre is bounded by the same G\_n=gcd(\textbar GL\_n(F\_g0)\textbar,\textbar GL\_n(F\_g1)\textbar). Thus NB12--15 count all actual cross terms for each chosen pair configuration, not merely a plane or a prescribed cross direction.

\subsection{5. Actual selected real factor and the complete subset bounds}\label{actual-selected-real-factor-and-the-complete-subset-bounds}

The nonzero cross witness x now has paired denominator budget v. Therefore product\_tau \textbar tau(x)\textbar\textgreater=T\^{}(−v). The actual near-point conjugator and complete selected-compression argument DO063(FV22--29) then give, with no change of the point or selected metric,

\[
        |c(h)|<12^f2^{f/2-485f/4}T^{v/2}.
 \tag{NB16}
\]

In detail, \textbar\textbar k\_p\textbar\textbar\textbar\textbar k\_p\^{}−1\textbar\textbar\textless2, a unipotent u has \textbar\textbar u−I\textbar\textbar\textless=2 sinh r(u), and the full selected compression has the product bound12/{[}sqrt(2 sinh r) kappa0\^{}(1/4){]}, with kappa0\textgreater2\^{}485. No Cartan factor is zero because x is a nonzero algebraic entry. The normalized trace \textbar Tr\textbar/dim(W) is bounded by this operator norm; there is no unrecorded selected-dimension factor.

Multiplying the weighted pair count NB8, column count, NB10 and NB16 proves the complete rank2 scalar-cross subset bound

\[
 \begin{split}
 U_{s,v,u}={}&W_{s,v,u}G_n A_v B T^{2mv+10s}
  \frac{2K^{f(3a/2-1)}T^{av+(a-1)(v+u)}}
       {Q^{(2a-1)(f-1)}}\\
 &\hspace{20mm}\cdot12^f2^{f/2-485f/4}T^{v/2}.
 \end{split}
 \tag{NB17}
\]

This counts a complete (s,v,u) subset at all its actual scalar bases. For terms supported on W0 use the already established original rank-five column and fibre count. With G5=98148490017978252769652736000000, the corresponding bounds are

\[
          U^{(j)}_{s,v,u}
       =W_{s,v,u}G_5(2H+1)^{10f}T^{10s}B_j(s,v,u),
                  \qquad j=1,2.
 \tag{NB18}
\]

Here the real factor is bounded by1; NB16 is needed only in the cross case. These fixed-block counts do not assert that an arbitrary rank2 unipotent is supported on W0.

\subsection{6. Exact removed subsets and the five-place residual}\label{exact-removed-subsets-and-the-five-place-residual}

The certificate evaluates NB17 at every 0\textless=u\textless=v\textless=5 and max(5,v)\textless=s\textless=168. In the following table every integer s from5 to the displayed endpoint satisfies U\_(s,v,u)\textless2\^{}−4096. A dash means that no s\textgreater=5 in that row passes this threshold.

\begin{longtable}[]{@{}rrrrr@{}}
\toprule\noalign{}
v & u & t=v−u & largest s & Binary margin at that s \\
\midrule\noalign{}
\endhead
\bottomrule\noalign{}
\endlastfoot
0 & 0 & 0 & 168 & 1150040 \\
1 & 0 & 1 & 168 & 782213 \\
1 & 1 & 0 & 168 & 615974 \\
2 & 0 & 2 & 168 & 201523 \\
2 & 1 & 1 & 168 & 35283 \\
2 & 2 & 0 & 155 & 11331 \\
3 & 0 & 3 & 125 & 14999 \\
3 & 1 & 2 & 110 & 12989 \\
3 & 2 & 1 & 95 & 10989 \\
3 & 3 & 0 & 80 & 8998 \\
4 & 0 & 4 & 65 & 12523 \\
4 & 1 & 3 & 50 & 10546 \\
4 & 2 & 2 & 35 & 8580 \\
4 & 3 & 1 & 20 & 6628 \\
4 & 4 & 0 & 5 & 4704 \\
5 & 0 & 5 & 5 & 10182 \\
5 & 1,\ldots,5 & 4,\ldots,0 & --- & --- \\
\end{longtable}

These are 1,521 distinct cross subsets. Their t is at most5, so the actual K068 residual is globally unipotent throughout, even at s168. No extension of the earlier total-support unipotence threshold is assumed; this uses K070's t threshold.

For completeness the fixed-block certificate evaluates NB18 for 0\textless=u,t\textless=8, v=u+t, max(5,v)\textless=s\textless=168. The two tables give the largest passing s; each interval begins at max(5,u+t). Empty entries have no passing interval.

\begin{longtable}[]{@{}rrrrrrrrrr@{}}
\toprule\noalign{}
Rank1: u & t0 & t1 & t2 & t3 & t4 & t5 & t6 & t7 & t8 \\
\midrule\noalign{}
\endhead
\bottomrule\noalign{}
\endlastfoot
0 & 8 & 8 & 8 & 8 & 8 & 8 & 8 & 8 & 8 \\
1 & 7 & 7 & 7 & 7 & 7 & 7 & 7 & --- & --- \\
2 & 6 & 6 & 6 & 7 & 7 & 7 & --- & --- & --- \\
3 & 6 & 6 & 6 & 6 & --- & --- & --- & --- & --- \\
4 & 5 & 5 & --- & --- & --- & --- & --- & --- & --- \\
\end{longtable}

\begin{longtable}[]{@{}rrrrrrrrrr@{}}
\toprule\noalign{}
Rank2: u & t0 & t1 & t2 & t3 & t4 & t5 & t6 & t7 & t8 \\
\midrule\noalign{}
\endhead
\bottomrule\noalign{}
\endlastfoot
0 & 17 & 16 & 15 & 14 & 13 & 13 & 12 & 11 & 10 \\
1 & 14 & 13 & 12 & 12 & 11 & 10 & 9 & 8 & --- \\
2 & 12 & 11 & 10 & 9 & 8 & 7 & --- & --- & --- \\
3 & 9 & 8 & 7 & 6 & --- & --- & --- & --- & --- \\
4 & 6 & 5 & --- & --- & --- & --- & --- & --- & --- \\
\end{longtable}

There are 245 passing fixed-block subsets; none with u\textgreater=5 passes within the tested u,t range. Every passing subset has t\textless=8, so again global unipotence follows from the existing actual residual condition.

The 1,766 total subsets have absolute contribution less than

\[
                         1766\,2^{-4096}<2^{-4085}.
 \tag{NB19}
\]

They are disjoint by matrix type and (s,v,u). Restricting them further to the remainder after earlier exclusions only decreases this bound.

In particular, at s=5 the rank2 scalar-cross table removes every v\textless=4 and also (v,u)=(5,0). Both fixed-block ranks remove every u\textless=4. The nonscalar positive-block terms at s=5 were already removed by DO062. Therefore the remaining nonzero five-place terms satisfy

\[
 \begin{array}{ll}
 \text{rank2 with nonzero cross action:}&v=5,\ 1\le u\le5;\\
 \text{supported on }W_0\text{, rank1 or2:}&v=u=5,\ t=0.
 \end{array}
 \tag{NB20}
\]

Thus all five full pencils differ, and at least one of their incidences is fixed. This is a structural residual description, not a positivity statement. The type II local factors are positive, but the unchanged G partial trace and real selected trace still remain in every summand.

Let S\_new be the original signed short sum with s\textgreater=5, after removal of the K068 cuts, the already bounded DO062 subsets and these new subsets, using set differences to avoid double charging. The original K062 normalization becomes

\[
 \frac{E_{x,\mathrm{ram}}}
 {w_x\,28C_\infty d_Rd_W((T-2)/(T-1))^f}
             =1+S_{\rm new}+\varepsilon_{\rm new},
 \tag{NB21}
\]

where

\[
 |\varepsilon_{\rm new}|<
 2^{-1999}+2^{-104546}+2^{-101550}+2^{-211446}
 +2^{-11671}+2^{-1307182}+2^{-9556}+2^{-4085}
 <2^{-1998}.
 \tag{NB22}
\]

This incorporates the complete s4 result DO063 and retains every original class mass, finite projector and real normalization.

\subsection{7. Exact numerical certification}\label{exact-numerical-certification}

The companion \nolinkurl{pencil_nilpotent_pole_budget_certificate.py} supplies a reproducible certificate using directed rational logarithm intervals, not floating-point approximations. Here is its complete analytic rounding argument.

For 0\textless=z\textless=1/3 and M positive,

\[
 \log\frac{1+z}{1-z}
 =2\sum_{j=0}^{M-1}\frac{z^{2j+1}}{2j+1}+R_M,
 \qquad
 0\le R_M\le
 \frac9{4(2M+1)3^{2M+1}}.
 \tag{NB23}
\]

This follows by integrating the geometric series, or by summing its positive tail and bounding its denominators below by2M+1. The script uses M80 and scale2\^{}192, rounds each finite summand down to that scale, and adds80 scale units plus the upward-rounded displayed tail for its upper endpoint. For an integer N, let b=floor(log2 N) by bit length, and round N/2\^{}b down to the scale2\^{}192. The resulting y lies in{[}1,2{]}, so z=(y−1)/(y+1) lies in{[}0,1/3{]}. Since log is1-Lipschitz on{[}1,2{]}, one additional scale unit covers this normalization truncation. Dividing by opposite endpoints of the separately certified natural logarithm of2, then adding b, gives exact rational lower and upper bounds for log\_2 N.

Products, rational powers and quotients in NB17--18 are then evaluated by interval addition with the appropriate signs. Both interval endpoints are checked against consecutive integers, so each displayed binary margin is certified without a floating-point logarithm. The integer root brackets for NB13, factorials, discriminant, host primes and energy floor are checked directly. Every integer s in each displayed interval is tested; the certificate also checks that there are no holes in a claimed interval and excludes every failed case from the listed intervals.

Local exact matrix diagnostics separately check preservation of M,N, cross-block and square integrality in type I, genuine unequal-full-pencil type III simple poles, and the genuine type II double-pole example. These finite local checks supplement the complete local proof; they do not enumerate the rational support of the automorphic kernel.

\subsection{8. Sources and scope}\label{sources-and-scope}

The actual pencil and signed factors are K062(AT3--17) and K070(PT1--10). The exact trace argument is DO062(RU12--13), the unitary normal form and witness calculations are DO061(TR19--36), and the actual fractional-lattice and selected-compression inputs are DO063(FV5--29). The covering-radius theorem and the Jacobi bound have already been source checked there and in LIFT47/K060; no new external representation or geometry theorem is needed for the refinements here.

Banaszczyk, W. (1993). New bounds in some transference theorems in the geometry of numbers. \emph{Mathematische Annalen, 296}, 625--635. \url{https://doi.org/10.1007/BF01445125}. Theorem2.2, p.632, enters only through the exact Euclidean application already proved in DO063.

Haagerup, U., \& Schlichtkrull, H. (2014). Inequalities for Jacobi polynomials. \emph{The Ramanujan Journal, 33}, 227--246. \url{https://doi.org/10.1007/s11139-013-9472-4}. Theorem1.1, p.228, enters only through the full selected homogeneous compression K060.

The new statements are NB5, the weighted incidence enumeration NB7--8, the resulting pole-budget/column refinements and their exact removed subsets. This is an estimate within the already constructed mixed geometric family. It neither repairs the literal source-size/link obstructions LIFT51--55 nor proves all-T Mram positivity, a terminal nonzero value or the required final DP44 sign. The unresolved signed remainder and its phases are retained explicitly.

\subsection{Author verification record}\label{author-verification-record-16}

At 2026-11-22T11:06:07+00:00, I completed a whole saved read and reconstruction of NB1--23, all eight sections and the complete final certificate. I checked both paired CM components, the distinction between full entries and cross entries, nilpotence on the actual two-dimensional quotient, the genuine type II square-pole example, the complete weighted ordered-pair count with both type-position binomials, the improved S-column ideal, both actual norm witnesses, the unchanged selected real trace and all removed-cell ranges. The full directed-rational logarithm error proof was checked against the script, including all signs, the tail and normalization rounding.

The final replay passed 62,067 exact assertions at 2026-11-22T11:03:58.597504+00:00. It includes 48 actual split-local lattice models, four genuine unique-fixed double-pole models, exact root brackets, 1,521 cross and245 fixed-block passing cells, all threshold failures and the retained normalized error. The diagnostics do not enumerate global rational matrices or evaluate a period. No main source-admissibility statement or Lean build is claimed.

The initial pole-one suggestion and the type I cross-integrality observation were supplied by \nolinkurl{/root/definite_oscillator}; the complete local square-integrality, weighted enumeration, field-norm application and certified subset bounds are proved in this report.

\clearpage

\section{AM-DO-064 --- partial monomial exceptional blocks and their actual short-sum bound}\label{am-do-064-partial-monomial-exceptional-blocks-and-their-actual-short-sum-bound}

Author: GPT-6 Astra Pro, ALP\_MACHINE \nolinkurl{/root/definite_oscillator}. 22 November 2026 UTC. Inherited reasoning effort is not separately exposed.

\textbf{Author whole-proof check PASS --- 2026-11-22T11:10:50.785383+00:00.}

This concerns exactly the IC8 exception in AM-K069, for the unchanged J56 short sum of K062 after the K068 cuts. Its integral positive block is a partial monomial matrix with entries in mu28 and at most three missing columns. More importantly, the rank of A−I cannot fall at any of the actual G primes. The Green-character estimate therefore has no field-norm numerator. Combined with the actual fractional-lattice column count, this removes a specified set of 2,881 rank/support cells with total absolute contribution less than 2\^{}-4084. In particular all remaining exceptional ranks32 and33 are removed for every support. The other exceptional cells, including A=I, remain explicit.

\subsection{1. The actual exceptional condition}\label{the-actual-exceptional-condition}

Carry K062(AT14--24), K068(LR1--24), K069(IC1--8), and DO063(FV1--23). The determinant-one representative and fixed decomposition are

\[
 a=f=168,\quad n=169,\quad m=164,\quad H=8192,
 \qquad h=\begin{pmatrix}A&C\\ B&D\end{pmatrix}_{S\perp W_0},
 \qquad J_{W_0}=\operatorname{diag}(1,1,1,1,-1).
 \tag{PM1}
\]

The coordinate space S is positive; A belongs to M\_m(O\_E), and every complex realization of h has operator norm at most H. Put

\[
 X=A-I,\quad r=\operatorname{rank}_E X\le33,\quad
 U=\ker X,\quad K_A=\ker X\cap\ker X^*.
 \tag{PM2}
\]

The exception is exactly

\[
 K_A=U,\qquad L=B(U)\ne0,\qquad \dim_E L=1,
 \qquad L\text{ isotropic},\qquad B(S)\subset L^\perp.
 \tag{PM3}
\]

All334 unmixed G reductions of h are unipotent. The other two finite projectors still have norm at most one. We keep the original summand

\[
 w(g)=a(g)c(g)\prod_{v\mid T}\rho(g_v),\qquad |c(g)|\le1,
 \tag{PM4}
\]

and the same central classes, labelled pencils, omitted lines, and class masses. At a nonintegral T place, \textbar rho\textbar≤2/(T−2). Let s be the number of those relative places. No projector or individual summand is changed in this report.

\subsection{2. Exact classification of A}\label{exact-classification-of-a}

\textbf{Proposition.} Under PM3, every nonzero entry of A is in mu28, there is at most one nonzero entry in each row and column, and

\[
 d:=\dim_E\ker A\le3.
 \tag{PM5}
\]

\textbf{Proof.} At any complex realization, the form induced on L-perp/L is positive definite of dimension three. Therefore

\[
 I-A^*A=B^*J_{W_0}B\ \text{ is positive semidefinite},
 \qquad \operatorname{rank}(I-A^*A)\le3.
 \tag{PM6}
\]

This holds at every realization, so each entry of the integral matrix A has all conjugate absolute values at most one. The settled root-unit argument K050(BE3--5), applied to each nonzero entry, makes it a root of unity. The carried field fact for the actual K032 host is mu(E)=mu28, as used in the same K050 argument.

The column norms of A are at most one by PM6. The row norms are also at most one, because its operator norm is at most one. Every nonzero entry now has absolute value one, so no row or column can contain two. Thus A is a partial monomial matrix. Its A*A is diagonal, with entry one on a nonzero column and zero on a missing column. PM6 identifies the number of missing columns with its rank, which is at most three. This also equals dim ker A. ∎

There are exactly

\[
 \mathcal N_A=\sum_{d=0}^{3}\binom md^2(m-d)!28^{m-d}<2^{1768}
 \tag{PM7}
\]

partial monomial matrices of this size with defect at most three. Choose the missing rows and columns, a bijection between the others, and one of28 weights for every nonzero entry. This is a bound for the actual exceptional A, without imposing its further rank or finite conditions. It is independent of T, H, and the discriminant. It does not assert that all these matrices occur in the short sum.

The defect bound three is sharp as a linear-algebra statement. Take a nontrivial signature-(2,1) unipotent on one S coordinate and one positive/one negative W0 coordinate, with its S compression equal to one and cross image a nonzero isotropic line. Independently swap d≤3 other S coordinates with d of the remaining positive W0 coordinates, using determinant-one signed swaps. Then A has d zero diagonal entries and is identity elsewhere; PM3 holds. These examples do not satisfy the asserted G-unipotence conditions in general and are not claimed as actual residual summands. The classification alone is not A=I.

\subsection{3. No rank loss at the actual G primes}\label{no-rank-loss-at-the-actual-g-primes}

View a partial monomial matrix as a directed graph on its m coordinate vertices, with edge j→i of weight A\_ij. Every vertex has at most one incoming and outgoing edge. Its components are weighted cycles and weighted paths, including isolated vertices as paths of length one. After a simultaneous coordinate permutation, A is their direct sum.

On a path with l vertices, A is nilpotent and A−I has rank l. On a cycle with l vertices and product weight zeta, the equations Av=v determine all coordinates from one and impose (zeta−1)v\_1=0. Hence

\[
 \operatorname{rank}(A-I)|_{\mathrm{cycle}}=
 \begin{cases}l-1,&\zeta=1,\\l,&\zeta\ne1.\end{cases}
 \tag{PM8}
\]

Reduction is injective on mu28 at every prime not dividing28. For example, all differences of distinct28th roots are units there, because their product divides a unit multiple of28, or directly because X\^{}28−1 has unit derivative on its roots. Thus the graph, the path ranks, and the condition zeta=1 in PM8 are unchanged. In particular, at every actual G place,

\[
 \operatorname{rank}(\overline A-I)=r,
 \qquad \operatorname{rank}(\overline h-I)\ge r.
 \tag{PM9}
\]

This is an exact rank statement, not a bound on how many primes can divide one chosen minor. No T denominator or archimedean minor bound enters PM9. The graph also shows that A is identity on all but at most 2r coordinate vertices, though that observation is not needed below.

The unipotent Green ratio carried in K067(PS5--9), equivalently the retained primary Green formula, now gives

\[
 |a(g)|\le \frac{2}{Q^{L_r(f-1)}},\qquad
 L_r=rn-\frac{r(r+1)}2,\qquad Q=g_0g_1.
 \tag{PM10}
\]

Indeed each unmixed factor has rank at least r and is bounded by the inverse product of q\^{}j−1 for n−r≤j≤n−1. There are f−1 such places over each g\_i. Their q-power contribution is precisely the denominator in PM10. The remaining Euler correction is less than two: even using q≥2, r≤33 and n−r≥136, the sum of all334r omitted q\^{}-j terms is less than one half, and the product of their complements is greater than one half. The two mixed factors retain only their norm bound one.

\subsection{4. A fractional-lattice count for a general last column}\label{a-fractional-lattice-count-for-a-general-last-column}

The following count does not require global unipotence. Fix the support and the ordered labelled source and target pencils. Let Sigma be the resulting2s conjugate CM primes, and put

\[
 I_\Sigma=\prod_{P\in\Sigma}P^{-1},\qquad
 D_0=\frac{\sqrt{|d_E|}}{2^f}
       =7^{140}29^{144}5^{84}13^{84}.
 \tag{PM11}
\]

The actual target lattice is O\^{}m direct-sum its five-coordinate paired pencil lattice. Exactly as in DO063(FV8--19), the last column h e\_n belongs to

\[
 I_\Sigma^m\oplus\mathcal B,\qquad
 \mathcal B=I_\Sigma\mathcal L,
 \tag{PM12}
\]

where the paired W0 lattice L has determinant valuation zero, is J-selfdual, and lies in I\_Sigma O\_E\^{}5. In the half-trace Euclidean metric the two covolumes and covering-radius bounds are

\[
 \begin{array}{c|c|c}
 &\operatorname{covol}&\mu/\sqrt f\\ \hline
 I_\Sigma^m&D_0^mT^{-2ms}
       &\displaystyle\frac{\sqrt m\,\operatorname{rd}(E)}{2T^{s/f}}\\[1mm]
 \mathcal B&D_0^5T^{-10s}&\displaystyle\frac52\operatorname{rd}(E).
 \end{array}
 \tag{PM13}
\]

These are the actual fractional duals and the checked covering-radius transference bound, not determinant-only estimates. In particular rd(E)\textless1500. For s≥4, T\^{}(s/f)\textgreater2\^{}23 and sqrt(m)\textless13.

If a unitary matrix for signature(a,1) has largest Euclidean singular value R, its singular values are R, R\^{}-1 and a−1 copies of one, and

\[
                  \|h e_n\|^2=(R^2+R^{-2})/2.
 \tag{PM14}
\]

One may see this by rotating the positive components of its last column and last row with compact unitary matrices; the only noncompact block is the ordinary2-by-2 boost. This also proves PM14 when R=1. It is the same general identity used in DO063, before that report specializes its positive component to a unipotent matrix. Since R≤8192, the norm in PM14 is less than5793 at every realization. Consequently the ball radii enlarged by PM13 are bounded by

\[
             5794\sqrt f\quad\text{on S},\qquad
             9543\sqrt f\quad\text{on W0}.
 \tag{PM15}
\]

The Dirichlet-cell volume bound therefore gives at most

\[
 V_S(s)=\frac{(\pi f5794^2)^{fm}T^{2ms}}{(fm)!D_0^m},\qquad
 V_W(s)=\frac{(\pi f9543^2)^{5f}T^{10s}}{(5f)!D_0^5}
 \tag{PM16}
\]

possible projected column pairs. This is a product upper bound for the actual columns; it does not assume independence of their two blocks. For a fixed full last column, the actual finite stabilizer fibre has order at most the carried

\[
 G_n=\gcd(|GL_n(\mathbb F_{g_0})|,|GL_n(\mathbb F_{g_1})|)<2^{5022}.
 \tag{PM17}
\]

There are at most binom(f,s)(T−1)\^{}(2s) support and ordered-pencil choices. Thus the total count, simultaneously over all A, is at most

\[
            \binom fs(T-1)^{2s}G_nV_S(s)V_W(s).
 \tag{PM18}
\]

There is no extra factor N\_A here: the full-column fibre count already counts every possible A. Formula PM7 is a separate exact finite classification, and should not be charged a second time.

\subsection{5. Exact support thresholds and the removed set}\label{exact-support-thresholds-and-the-removed-set}

Let E\_(r,s) be the actual IC8 subset with the indicated rank and support. Combining PM10, PM18 and the unchanged bound on rho gives

\[
 \sum_{[g]\in E_{r,s}}|w(g)|\le
 \binom fs(T-1)^{2s}G_nV_S(s)V_W(s)
 \left(\frac2{T-2}\right)^s
 \frac2{Q^{L_r(f-1)}}.
 \tag{PM19}
\]

This is an absolute upper bound on the original signed sum, with all finite factors and class normalizations retained. For a short exact certificate, replace pi by22/7. The resulting constant part of V\_S V\_W is less than2\^{}324076. The actual primes satisfy

\[
 T<2^{996},\qquad Q^4>2^{281},\qquad
 \frac{2(T-1)^2}{T-2}<2^{997}.
 \tag{PM20}
\]

The final inequality follows also from (T−1)\^{}2/(T−2)=T+1/(T−2)\textless T+1≤2\^{}996. Using binom(f,s)≤2\^{}f, PM19 is strictly less than2\^{}E(r,s), where

\[
 E(r,s)=329267+337645s-\frac{281}{4}\,167L_r.
 \tag{PM21}
\]

Define the exact integers

\[
 R_r=\min\left(168,
 \left\lfloor\frac{-4096-329267+(281/4)167L_r}{337645}
 \right\rfloor\right).
 \tag{PM22}
\]

Their values are:

\begin{longtable}[]{@{}rrrrrr@{}}
\toprule\noalign{}
r & R\_r & r & R\_r & r & R\_r \\
\midrule\noalign{}
\endhead
\bottomrule\noalign{}
\endlastfoot
1 & 4 & 12 & 66 & 23 & 124 \\
2 & 10 & 13 & 72 & 24 & 129 \\
3 & 16 & 14 & 77 & 25 & 134 \\
4 & 22 & 15 & 82 & 26 & 139 \\
5 & 27 & 16 & 88 & 27 & 144 \\
6 & 33 & 17 & 93 & 28 & 149 \\
7 & 39 & 18 & 98 & 29 & 154 \\
8 & 44 & 19 & 103 & 30 & 159 \\
9 & 50 & 20 & 109 & 31 & 163 \\
10 & 55 & 21 & 114 & 32 & 168 \\
11 & 61 & 22 & 119 & 33 & 168 \\
\end{longtable}

Every cell with5≤s≤R\_r has contribution less than2\^{}-4096. There are exactly2,881 such cells. Their union therefore has total absolute contribution

\[
 \boxed{\ \sum_{r=1}^{33}\sum_{s=5}^{R_r}
           \sum_{[g]\in E_{r,s}}|w(g)|<2^{-4084}.\ }
 \tag{PM23}
\]

The earlier complete strata s≤4 were already removed by DO059--063. In particular no IC8 term of rank32 or33 remains outside the combined removed sets. The rank-zero exception and, for1≤r≤31, the supports s\textgreater R\_r are not bounded by PM23. This result does not assert a contradiction for all exceptional A, a sign for the total remaining sum, a nonzero period, or a literal full SizesAt family.

\subsection{Verification and references}\label{verification-and-references-2}

The companion \nolinkurl{definite_oscillator_exceptional_partial_monomial_certificate.py} checks the exact finite count, the actual field/prime constants, the rational volume bound, all support thresholds and the union budget. It also checks weighted cycle/path ranks over characteristic zero and finite fields, the support bound2r, and explicit signature-one examples attaining defect three. These are diagnostics for the proved classification and inequalities, not an enumeration of rational correspondences in the original short sum.

New proofs here are elementary block linear algebra, finite graph rank calculations, and rational inequalities. The external geometry-of-numbers theorem is carried from the independently checked DO063/LIFT47 use of Banaszczyk's covering-radius transference bound. The character input is the independently checked K067 use of the retained primary Green formula. The actual CM lattice, primes, root-unit group, finite fibre and class normalization are carried from K032/K050, DO040/058/063, and K062; no new source constructor is assumed.

The carried primary references are:

\begin{itemize}
\tightlist
\item
  Banaszczyk, W. (1993). New bounds in some transference theorems in the geometry of numbers. \emph{Mathematische Annalen, 296}, 625--635. Theorem2.2, p.632. \url{https://doi.org/10.1007/BF01445125}.
\item
  Gorodetsky, O., \& Hazan, Z. (2017). \emph{On certain degenerate Whittaker models for cuspidal representations of GL\_(kn)(F\_q)}. arXiv:1707.07308v2, Theorem2.1, equation(2.1), p.5. \url{https://arxiv.org/pdf/1707.07308}.
\item
  Green, J. A. (1955). The characters of the finite general linear groups. \emph{Transactions of the American Mathematical Society, 80}(2), 402--447. Section5, Example(ii), pp.430--431.
\end{itemize}

I read the complete saved PM1--PM23 proof and every section, and the entire final certificate. The classification, primewise graph ranks, paired fractional duals, actual column fibre, all exponents and the signed-sum scope were reconstructed. The exact certificate passed 14,859 assertions at 2026-11-22T11:10:08.784299+00:00, including1,631 partial-monomial graph models,448 full-mu28 cycle models and four unitary exception models. No actual rational sum or period was computed. The final verification and bibliography additions were reread before freezing.

\clearpage

\section{ALP-LIFT-56: forward first occurrence and cuspidality for the actual longer vertical signatures}\label{alp-lift-56-forward-first-occurrence-and-cuspidality-for-the-actual-longer-vertical-signatures}

Author: /root/lift\_accounting, GPT-6 Astra Pro, ALP\_MACHINE.

\subsection{1. Statement and retained hypotheses}\label{statement-and-retained-hypotheses}

Let the actual ranks satisfy \nolinkurl{Ranks}, and let the selected parameters satisfy the \nolinkurl{Hyp3} hypotheses used in AM-R-015. Fix genuine compatible unitary splitting characters on the actual CM extension and the oriented vertical theta kernels of that finding. Put \(t=P.t\) for this parameter \(P\); it need not equal the independent integer \(\mathrm{Ranks}.t\). No new character or finite Schwartz input is substituted below. Write

\[
 r=k+m-1,\qquad n_j=r+jt,\qquad
 q_j=\left\lfloor\frac{m-j}{2}\right\rfloor,
 \quad p_j=n_j-q_j,\quad N_j=k+j(t+1).
 \tag{VF1}
\]

Here \(0\le j\le m\), \(k,m\) are even, \(m\ge2\), and \(t\) is odd. The TPar.t\_ge inequality and \(k\ge m+3\) imply

\[
 t\ge r+5\ge2m+7>m.
 \tag{VF2}
\]

For the canonical TPar.ofWord, \(t=\mathrm{Ranks}.t\) and the extra rank inequality \(t>3m\) also holds; it is not used in this finding.

Let \(\sigma_j\) be the real discrete series with the actual signed Harish--Chandra parameter \nolinkurl{Dpar j}. The genuine source type \nolinkurl{IsDv} for an automorphic input is its conjugate \(\tau_j=\overline{\sigma_j}\). AM-R-015 fixes the Hom-to-tensor convention, character shifts, and the unconditional local correspondence

\[
 \theta_{Y_{j+1}}(\sigma_j)=\sigma_{j+1}\ne0.
 \tag{VF3}
\]

Those accepted results are used, not reproved. The new conclusions are:

\begin{enumerate}
\def\labelenumi{\arabic{enumi}.}
\tightlist
\item
  Every forward lift of \(\sigma_j\) to a proper Witt-lower space of \(Y_{j+1}\) is zero, with the same compatible splitting characters.
\item
  Consequently every prescribed upward kernel sends cusp forms of type \(\tau_j\) to cusp forms of type \(\tau_{j+1}\), possibly zero. This holds at every stage and for every finite Schwartz input, including the literal \nolinkurl{DatumVec} inputs.
\item
  At each real place there are polynomial-Fock and compact-finite test vectors giving a nonzero local Rallis pairing for this edge, even when both real unitary groups are indefinite.
\end{enumerate}

There is no claim here that one prescribed global finite tensor is nonzero, that a chosen highest polynomial gives a nonzero local pairing, or that the arithmetic row or positive-rank period is detected. In particular this is not a construction of a full \nolinkurl{D7}, \nolinkurl{Host}, \nolinkurl{SizesAt}, or all six \nolinkurl{AdmS} predicates.

\subsection{2. Exact sign counts and the correction to the older criterion}\label{exact-sign-counts-and-the-correction-to-the-older-criterion}

Keep the actual cross-character centering of AM-R-015 (R15.4--9). At each selected real place the smaller signed list is

\[
 D_j-c_{j+1}=(A_j\cup B_j\ ;\ E_j),
 \qquad |A_j|=N_j,\quad |B_j|=q_{j+1},\quad |E_j|=q_j,
 \tag{VF4}
\]

where every \(A_j\)-entry is at least \((t+3)/2>0\), and all entries of \(B_j,E_j\) are at most \(-(t+3)/2<0\). These are integers. In particular the four sign counts in Ichino's Corollary 6.4 are exactly

\[
 (p^+,p^-,q^+,q^-)
       =(N_j,q_{j+1},0,q_j).
 \tag{VF5}
\]

No zero-centered parameter is present. The target signature is

\[
 (p_{j+1},q_{j+1})=(N_j+q_j+t,q_{j+1}),
 \qquad p_j=N_j+q_{j+1}.
 \tag{VF6}
\]

The identities follow from \(q_j+q_{j+1}=m-j-1\), and agree with the literal \nolinkurl{sigYD}, not merely the total ranks.

The primary theorem used for vanishing is \textbf{Ichino (2022), Corollary 6.4, printed pp.21--22}. It says that for a (limit of) discrete series, a nonzero lift to a target of at least the source dimension must have positive and negative indices at least

\[
                 p^++q^-,\qquad p^-+q^+,
 \tag{VF7}
\]

respectively. The corollary is derived there from the corrected Theorem 6.1. Section 6.2, printed p.20, explicitly notes the defect in Atobe's older Theorem 4.2(1), states the corrected criterion, and records its unconditional status. Thus this proof does not use the uncorrected \(k=-1\) nonvanishing assertion.

For a dictionary check only, Atobe's Definition 1.6 gives

\[
 k_{\sigma_j}=-1,\qquad
 (r_{\sigma_j},s_{\sigma_j})=(N_j+q_j,q_{j+1}).
 \tag{VF8}
\]

This is a definition calculation, not an additional vanishing premise. Indeed all centered entries are nonzero integers. In the corrected Ichino formulation neither zero-signed point is present; the exceptional zero-parameter alternative cannot occur. The proof below needs only the directly stated necessary condition (VF7).

\subsection{3. Forward lower-tower vanishing}\label{forward-lower-tower-vanishing}

Fix \(j<m\) and \(1\le i\le q_{j+1}\). A proper Witt-lower target has signature

\[
       (p_{j+1}-i,q_{j+1}-i)
\]

and dimension \(n_{j+1}-2i=n_j+t-2i\). Since

\[
 2i\le 2q_{j+1}\le m-j-1\le m-1<t,
 \tag{VF9}
\]

this dimension is strictly greater than \(n_j\). Thus Corollary 6.4 applies. Its negative-index condition and (VF5) would require

\[
          q_{j+1}\le q_{j+1}-i,
\]

which is impossible. Therefore

\[
 \boxed{\quad
 \theta_{p_{j+1}-i,q_{j+1}-i}(\sigma_j)=0
     \quad(1\le i\le q_{j+1}).\quad}
 \tag{VF10}
\]

Vanishing means that the real big-theta quotient is zero, not simply that a particular test vector maps to zero. For the irreducible admissible input the big theta module has finite length and has the unique irreducible quotient when nonzero; equivalently one may use the nonvanishing convention in Ichino's theorem. Thus every continuous intertwining coefficient into these lower targets vanishes. Polynomial Fock/Harish--Chandra vectors are dense in their smooth models, so the same assertion applies to the smooth Schwartz kernels.

Restoring the literal orientation changes neither this conclusion nor the target signatures. AM-R-015 identifies the joint Hom type as \(\overline{\sigma_j}\boxtimes\sigma_{j+1}\). The source's upward adjoint uses the conjugated kernel and input \(\overline{\sigma_j}\); simultaneous conjugation transfers (VF10) to that upward operator. This uses the original character pair and additive character, not a separately chosen reverse Hom kernel.

Combining (VF10) with the accepted (VF3), the prescribed real target is the first nonzero member of its actual Witt tower. If \(q_{j+1}=0\), the lower-target assertion is empty and first occurrence within this anisotropic-ended tower has the same interpretation.

\subsection{4. Cuspidality for the unchanged finite kernels}\label{cuspidality-for-the-unchanged-finite-kernels}

The concrete analytic spaces of AM-R-014 are

\[
        Y_j=I_{p_j}\oplus(-I_{q_j}).
 \tag{VF11}
\]

Their rational Witt indices equal \(q_j\), and deleting \(i\) standard hyperbolic planes gives all maximal rational parabolic constant-term tests. More generally an actual Hermitian space with these real signatures has rational Witt index at most \(q_j\); use its actual Witt-lower spaces, not a signature-only assertion of rational isotropy.

Let \(F\) be a cusp form of real type \(\tau_j\), in the genuine smooth automorphic domain specified in AM-R-014. For every rational maximal parabolic of \(U(Y_{j+1})\), the unitary tower constant-term identity expresses the constant term of its upward theta lift as a lift to one of the lower spaces in (VF10), of a restricted Schwartz function. This is the accepted application of Wu's Proposition 3.1 in AM-R-014, R14.8--9; its proof uses only cuspidality of \(F\), not an irreducible projection or a nonzero image.

The global lower lift is zero because its local intertwining space at any selected real place is zero. The argument is unchanged after any target translation, which merely replaces the Schwartz input by its Weil translate. Hence every maximal proper constant term vanishes. All proper constant terms then vanish by their iterated-constant-term description. It follows that

\[
 \boxed{\quad
 U_j(\Phi_f,\Phi_\infty)F
       \in\mathcal A_{\mathrm{cusp}}(U(Y_{j+1}))_{\tau_{j+1}}
       \quad\text{or is zero}.\quad}
 \tag{VF12}
\]

The real type in (VF12) follows from (VF3), the accepted oriented Howe correspondence, and smooth globalization, exactly as in R015. Every lower dimension differs by an even integer, so the same splitting character has the required restriction to the base-field ideles. No auxiliary character twist is introduced in taking these constant terms. The hypothesis here includes genuine compatible characters; the formal \nolinkurl{IsMu} infinity condition alone is not their construction.

Since the finite vector was arbitrary, (VF12) applies in particular to the literal \nolinkurl{DatumVec.IsStepVecOf} finite inputs whenever their spaces are tied to this \(t\); in the canonical arithmetic instance this is TPar.ofWord with \(t=\mathrm{Ranks}.t\). Starting from any such bottom cusp form and applying an arbitrary fixed sequence of those kernels gives cusp forms of the exact types at every stage, unless an earlier lift is zero. Thus positive intermediate Witt rank creates no new cuspidality obstruction for these selected parameters. This conclusion retains, rather than solves, fixed-kernel nonvanishing.

\subsection{5. Real local pairings and the global range}\label{real-local-pairings-and-the-global-range}

For the real dual pair at edge \(j\), the source \(\sigma_j\) is tempered and the target dimension is \(n_j+t>n_j\). Ichino's Proposition 7.1, printed pp.24--25, equates nonvanishing of its local matrix-coefficient pairing with nonvanishing of the theta lift. The latter is (VF3). Consequently that pairing is nonzero. Remark 7.2 allows both oscillator test vectors in the polynomial Fock subspace and both representation vectors compact-finite. Thus there exist genuine real tests at every indefinite intermediate edge.

This is an existence statement for the local sesquilinear form. It does not set its scalar, normalize an \nolinkurl{IsItW} map, or identify a particular primitive polynomial from the arithmetic source. In particular it is not a formula for the norms used in LIFT38/39 at the formerly compact middle group.

For completeness the actual dimensional range is also explicit. If \(w_{j+1}\) is the rational Witt index of the target, then

\[
 n_{j+1}-w_{j+1}
      \ge n_j+t-q_{j+1}>n_j.
 \tag{VF13}
\]

GQT section 3.1, printed p.14, therefore places every edge in Weil's convergent range (the unitary value is \(d(n_j)=n_j\)). Moreover

\[
 n_1>2n_0,\qquad n_j<n_{j+1}<2n_j\quad(j\ge1).
 \tag{VF14}
\]

The first inequality uses \(t\ge r+5\); the second uses \(n_j=r+jt>t\). This separates the first stable edge from the dimension window of GQT Proposition 11.1 and Theorem 11.4 for the later edges. When \(m\ge4\), the first target may have positive rational Witt index. One cannot cite the anisotropic-only GQT Theorem 7.1 for that first edge. A subsequent global inner-product/nonzero proof must use the convergent Siegel--Weil formula with this actual Witt index and the correct doubling normalization. No such first-edge normalization is asserted by (VF13).

\subsection{6. Relation to the six analytic predicates and the remaining work}\label{relation-to-the-six-analytic-predicates-and-the-remaining-work}

The literal source predicates are \nolinkurl{CayleySignature.lean}, lines154--203. Accepted AM-R-013 covers the explicitly stated genuine analytic interpretations of predicates1 and2; AM-R-014 covers predicates4 and6; AM-R-015 covers predicate5, including the \textbf{reverse} lower-tower vanishing. Those accepted results already have general-rank scope.

The present finding supplies the \textbf{forward} lower-tower vanishing and actual fixed-finite-input cuspidality needed to extend the prescribed upward-adjoint construction beyond \(m=2\). It is an additional constructor step, not a newly proved sixth predicate. The positive-rank \nolinkurl{ThetaUnfoldingFreeData} predicate3, the actual finite tests for nonzero propagation, a bottom detector in the specified image, the arithmetic frame/row ties, and the full source Host/size conditions remain distinct. LIFT54--55's literal size exclusions therefore remain in force.

\subsection{7. Source anchors and references}\label{source-anchors-and-references}

\begin{itemize}
\tightlist
\item
  \nolinkurl{Formalism/LocalRingCardinality/Params.lean}, Ranks lines6--22 and rank definitions lines28--44; \nolinkurl{DataTheta.lean}, TPar.t/t\_ge/ofWord lines501--520 and Dpar/Centred lines538--555; \nolinkurl{DataSpaces.lean}, sigYD lines75 onward.
\item
  \nolinkurl{Formalism/SharpRestriction/CuspSpanNormBounds/CayleySignature.lean}, Hyp3/IsMu/IsDv/vth lines109--130 and all six predicates lines154--203.
\item
  \nolinkurl{Formalism/LocalToGlobal/IndefiniteReal/OrthogonalComplementCharacterisation.lean}, \nolinkurl{chrStep} lines13--18, \nolinkurl{IsIo} lines41--55, \nolinkurl{IsIoUp} lines78--88.
\item
  Accepted \nolinkurl{selected_real_descent.md}, R15.3--12, for the exact centered signs, characters and nonzero real correspondence; \nolinkurl{selected_tower_cuspidality_stability.md}, R14.3--9, for genuine spaces, lower towers and the constant-term application.
\end{itemize}

Ichino, A. (2022). Theta lifting for tempered representations of real unitary groups. \emph{Advances in Mathematics, 398},108188. \url{https://doi.org/10.1016/j.aim.2022.108188} . Primary author preprint: \url{https://arxiv.org/pdf/2008.06174} . Section6.2/Theorem6.1, printed p.20, explicitly corrects the older criterion; Corollary6.4, pp.21--22; Proposition7.1/Remark7.2, pp.24--25. Directly retrieved in this check. The general Arthur-conditional parameter theorem in the introduction is not used.

Atobe, H. (2020). On the non-vanishing of theta liftings of tempered representations of U(p,q). \emph{Advances in Mathematics, 363},106984. \url{https://doi.org/10.1016/j.aim.2020.106984} . Definition1.6, printed p.4, used only for the dictionary (VF8). No use of its uncorrected Theorem4.2(1) is made.

Gan, W. T., Qiu, Y., \& Takeda, S. (2014). The regularized Siegel--Weil formula (the second term identity) and the Rallis inner product formula. \emph{Inventiones Mathematicae, 198},739--831. \url{https://doi.org/10.1007/s00222-014-0509-0} . Section3.1, printed p.14; section7.1, printed p.33; Proposition11.1, printed p.53. Local source copy \nolinkurl{work/agents/lift_refs/gqt2014.pdf} and extracted text were read at those anchors. Only the convergence criterion is newly applied.

The check script \nolinkurl{general_vertical_forward_certificate.py} verifies exact signed-list/count/parity identities and the rank inequalities on finite diagnostic families, including nonuniform centers. It does not certify the cited analytic theorems, a global automorphic value, or a Lean build. The author whole-proof check is recorded separately after the final reread.

Author whole-proof reading PASS 2026-11-22T11:11:00.189461+00:00. Complete VF1--14/seven sections and full diagnostic read; primary Ichino corrected criterion, Corollary6.4 and Proposition7.1/Remark7.2 directly checked, and GQT convergence versus anisotropic-only first formula distinguished. No fixed finite nonzero theorem, global value, or Lean build is certified.

Targeted final formatting reread 2026-11-22T11:11:58.322842+00:00: restored parentheses inside VF3/VF10/VF12 and corrected two inline group/rank expressions. No mathematical statement, source dependency, or certificate changed.

Final author whole-proof reread PASS 2026-11-22T11:17:20.600766+00:00: \(t=P.t\) is now explicit and may differ from Ranks.t; only TPar.t\_ge and \(k\ge m+3\) are used, giving \(t\ge2m+7>m\). The canonical source constructor is TPar.ofWord. All seven sections and revised diagnostic reread; cases with \(P.t<3m\) are included. This supersedes the earlier rank-scope wording without changing VF3--14 or any nonzero claim.

\clearpage

\section{AM-K072 --- active-coordinate witnesses and actual selected energy in the five-place sum}\label{am-k072-active-coordinate-witnesses-and-actual-selected-energy-in-the-five-place-sum}

Author: GPT-6 Astra Pro, ALP\_MACHINE \nolinkurl{/root/kummer_extension}. 22 November 2026 UTC. Inherited reasoning effort is not separately exposed.

\textbf{Author whole-proof check PASS --- 2026-11-22T11:30:12+00:00.} The final exact certificate passed5563 assertions; independent reviews are separate events.

This finding retains the actual J56 mixed family, its upper K044 gap, the original finite projectors, K062's signed kernel and the cutoff H=8192. It removes new subsets of the five-place scalar-cross remainder by using the number of active positive coordinates. The count is in the original coordinate lattice, with every original class-fibre and finite incidence factor retained. An additional corollary uses the actual top-role selected degrees and covers the entire one-fixed-incidence scalar-cross case when at least119 of the168 labels are top. This includes the all-top mixed profile.

The all-top specialization is an actual mixed-kernel construction, not a literal Host, Frame or FamP record. The literal m=2 size obstructions remain in force. Neither the surviving two-through-five fixed-incidence cases nor the fixed-five-coordinate cases are removed by the all-top corollary. No all-place Mram positivity is inferred.

\subsection{1. The exact subset and its actual coordinates}\label{the-exact-subset-and-its-actual-coordinates}

Use K062, DO063 and K071's unchanged notation

\[
 a=f=168,\quad n=169,\quad m=164,\quad H=8192,\quad
 E^n=S\oplus W_0,\quad \operatorname{sig}(W_0)=(4,1),
 \qquad Q=g_0g_1.
 \tag{AC1}
\]

Here g0=22285932761, g1=64072175441, and T is the actual DO040 996-bit prime. The real inner product is the half-trace Euclidean metric used in DO063. Consider a nonintegral supported central class with exactly s=5 relative T places. Its unique unipotent representative h satisfies

\[
 Z=h-I,\qquad Z^3=0,\quad \operatorname{rank}Z=2,
 \qquad Z_{SS}=0,\quad Z_{SW_0}\ne0.
 \tag{AC2}
\]

Global unipotence on this five-place subset is already established in DO062; it is not inferred from finite reductions. We restrict to v=5 unequal-full-pencil places. Let u be the number of type-II places, with a unique fixed actual lattice, and t=5−u the number of type-III places, with a unique off-diagonal incidence. K071 already removes u=0, so the cases here have1\textless=u\textless=5. The exact local factors are respectively 1/(T−1) and −1/{[}(T−1)(T−2){]}. Their signs are retained in the original kernel; the estimates below bound absolute sums.

Let e\_n be the original negative coordinate vector and put

\[
 x_i=(Ze_n)_i=Z_{in}\quad(i\in S),\qquad
 D(h)=\{i\in S:x_i\ne0\},\qquad d=|D(h)|.
 \tag{AC3}
\]

Thus1\textless=d\textless=164. Each x\_i belongs to E, so its zero or nonzero status is the same at every embedding. The active set is not chosen separately at the CM labels. The coordinate n and the subspace S are the original fixed coordinates of the actual pencil.

\subsection{2. Two sharper witnesses for the actual finite character}\label{two-sharper-witnesses-for-the-actual-finite-character}

The accepted signature-one unipotent normal form in DO061(TR20,32) has a Euclidean-unit isotropic vector l lying in W0 and a positive vector v perpendicular to l and Jl. At each embedding,

\[
 Z_{SW_0}=v_S(Jl)_{W_0}^*,\qquad
 Z_{W_0S}=-l_{W_0}v_S^*,\qquad
 Z^2=-\|v\|^2ll^*J.
 \tag{AC4}
\]

In particular \textbar l\_n\textbar\^{}2=1/2 and v\_n=0. Define the fixed algebraic witness

\[
                         y=(Z^2)_{nn}.
 \tag{AC5}
\]

It is nonzero and real under every CM embedding. Equations AC4 give

\[
 y_\tau=\tfrac12\|v_\tau\|^2>0,\qquad
 \sum_{i\in D(h)}|\tau x_i|^2
     =\tfrac12\|(v_\tau)_S\|^2\le y_\tau.
 \tag{AC6}
\]

There is also a useful direct normalization check. The exact identity Z+Z\^{}dagger=Z\^{}2 gives

\[
           y_\tau=2\operatorname{Re}(h_\tau)_{nn}-2.
 \tag{AC7}
\]

If r\_tau is the Cartan parameter of h\_tau, its negative diagonal entry has absolute value cosh(r\_tau). This follows directly from the compact factors in a Cartan decomposition: they act by phases on the negative line. The cutoff gives exp(r\_tau)\textless=H. Consequently

\[
 0<y_\tau\le2\cosh r_\tau-2
      \le H+H^{-1}-2<H+1.
 \tag{AC8}
\]

No algebraic moving normalization of l or v is needed: only x\_i and y, entries of the original algebraic matrices, are used in the field norm.

Apply generalized Holder to the nonnegative array \textbar tau x\_i\textbar\^{}2:

\[
 \sum_{i\in D(h)}
       \left(\prod_{\tau\in\Psi}|\tau x_i|^2\right)^{1/f}
 \le\prod_{\tau\in\Psi}
       \left(\sum_{i\in D(h)}|\tau x_i|^2\right)^{1/f}
 <H+1.
 \tag{AC9}
\]

Thus at least one nonzero original coordinate x=x\_i obeys

\[
           \prod_\tau|\tau x|
                <\left(\frac{H+1}{d}\right)^{f/2},
 \qquad \prod_\tau|\tau y|<(H+1)^f.
 \tag{AC10}
\]

Selecting this witness does not multiply the matrix count by d: the inequality is pointwise for each already counted matrix.

K071(NB9) gives paired pole-order budgets5 for x and5+u for y. At every unmixed selected G place, reduction rank zero makes x vanish, and reduction rank at most one makes Z\^{}2, hence y, vanish. These statements hold at both CM partners. The same field-norm argument and Green factors as DO061(TR23--24) therefore give

\[
 \boxed{\quad
 |a(h)|\le
 \frac{2\,[(H+1)/d]^{af/2}(H+1)^{(a-1)f}
              T^{5a+(a-1)(5+u)}}
      {Q^{(2a-1)(f-1)}}.
 \quad}
 \tag{AC11}
\]

Here a(h) is the actual normalized finite partial trace of K062. The factor2 is the carried total Euler correction. Both mixed finite factors keep their original projectors and the bound one. If a rank-two reduction happens to have square zero, the y witness loses the second rank factor there; that only enlarges AC11. No assertion of a new finite-type multiplicity or a probability of rank drop is used.

\subsection{3. Counting only the active coordinate lattice}\label{counting-only-the-active-coordinate-lattice}

Fix the five-place support, its types, all ordered actual source and target pencil points, and an active set D of size d.~By K071(NB12), the S part of the actual last column belongs to I\^{}m, where I is the product of the inverse primes at both CM partners over these five places and N(I)=T\^{}−10. Its inactive coordinates are zero. Thus its actual lattice is I\^{}d in the fixed coordinate subspace. Equation AC4 also shows that h fixes each inactive positive coordinate, although the estimate below retains the larger original fibre bound G\_n and does not use a smaller group.

Put

\[
 D_0=7^{140}29^{144}5^{84}13^{84},\qquad
 R_5=90.516+\frac{9750}{828375040}.
 \tag{AC12}
\]

The exact lower root inequality (828375040)\^{}f\textless=T\^{}5 is checked in the companion certificate. DO063's scalar-ideal dual estimate gives covering radius at most

\[
 \frac{\sqrt{fd}\,\operatorname{rd}(E)}{2T^{5/f}}
       <\frac{750\sqrt{fd}}{T^{5/f}}.
 \tag{AC13}
\]

The actual S-column trace norm is at most sqrt(f(H+1)) by AC6--8. Since sqrt(d)\textless13, the sum of this radius and AC13 is less than sqrt(f) R5. The half-trace covolume of I\^{}d is D0\^{}d T\^{}−10d. Dirichlet-cell volume comparison in the actual 2fd-dimensional Euclidean space, with pi\textless22/7, proves the bound

\[
 \#\{\text{active S columns}\}
   \le A_d T^{10d},\qquad
 A_d=\frac{[(22/7)fR_5^2]^{fd}}{(fd)!D_0^d}.
 \tag{AC14}
\]

The W0 part is exactly the paired target over-lattice used in DO063, not a coordinatewise replacement of it. Retain its proved bound

\[
 \#\{W_0\text{ columns}\}\le B_0T^{50},\qquad
 B_0=\frac{[(22/7)f9543^2]^{5f}}{(5f)!D_0^5}.
 \tag{AC15}
\]

For each full last column keep the original isometry-fibre bound

\[
 G_n=\gcd\bigl(|\operatorname{GL}_{169}(\mathbb F_{g_0})|,
               |\operatorname{GL}_{169}(\mathbb F_{g_1})|\bigr).
 \tag{AC16}
\]

There are binom(m,d) active subsets. Since every place is of type II or III, K071(NB7--8) cancels the ordered-pair costs against the respective absolute rho values. The remaining support/type factor is exactly bounded by binom(f,5)binom(5,u). This cancellation is performed after the uniform column count for each actual pair; it is not a claim that the unweighted number of matrices lacks the ordered-pair factor. Repeated support representations only increase the upper bound.

\subsection{4. Actual homogeneous energy in the selected compression}\label{actual-homogeneous-energy-in-the-selected-compression}

K060(RC3) has the stronger energy inequality on the homogeneous space H\_N

\[
                    H_1\ge\kappa_0+N,
          \qquad\kappa_0=2\ell+q+a.
 \tag{AC17}
\]

Indeed its exact eigenvalue is nu\_1−beta+N+gamma\_1; every fibre weight coordinate is at least alpha\_(a−1), and gamma\_1\textgreater=0. In each orthogonal rank-one ladder meeting H\_N, the diagonal coefficient has denominator (2n+kappa)\^{}(1/4), and 2n+kappa is this very H1 energy. Thus the original RC6/RC13 proof, with no change of Hilbert norm or summation over ladders, gives

\[
 \|P_W\pi(g)|_W\|
 \le\min\left\{1,
       \frac{12}{\sqrt{2\sinh r(g)}(\kappa_0+N)^{1/4}}
             \right\}.
 \tag{AC18}
\]

Compact Cartan factors preserve H\_N, so this is the full selected component bound in the exact K043 PBW isometry. It is not a highest-vector or selected-to-minimal estimate.

The actual near-origin conjugator is the unchanged LIFT41 one. DO063(FV24--27) and the nonzero original cross coordinate give product\_tau \textbar tau x\textbar\textgreater=T\^{}−5 and \textbar tau x\textbar\textless2(2sinh r\_tau(h\_p)). Hence the normalized selected trace in K062 satisfies

\[
 |c(h)|\le12^f2^{f/2}T^{5/2}
                \prod_\tau(\kappa_{0,\tau}+N_\tau)^{-1/4}.
 \tag{AC19}
\]

There is no factor dim(W). All embeddings have nonzero Cartan parameter in this argument: a compact unipotent would be the identity and would contradict the faithful nonzero cross entry.

\subsubsection{The actual role bounds}\label{the-actual-role-bounds}

Retain K044's upper endpoint ell=ell\_+, A=2\^{}960, B=2\^{}955, R=E\_T−B, and its allowed canonical twist q\textgreater=0. Let w be the number of ones in the binary word, so0\textless=w\textless=a−2. Directly summing the actual exponents K043(PB3,PB12--13) gives

\[
 \begin{array}{c|c}
 \text{role}&\kappa_0+N\\ \hline
 \text{top}& A+R-(w+3)\ell+q+1\\
 \text{thick or thin}&(8a-5-w)\ell+E_\beta+q+2.
 \end{array}
 \tag{AC20}
\]

Define the exact integers

\[
 E_*=(7a-3)\ell_++E_L+2,\qquad
 E_{\mathrm{top}}=A+R-(a+1)\ell_++1.
 \tag{AC21}
\]

The certificate checks E\_top\textgreater E\_*\textgreater2\^{}494 and E\_top\textgreater2\^{}959. Every word and every allowed twist obey AC21 as a lower bound for the respective role. If z labels are top, AC19 is therefore bounded by

\[
               12^f2^{f/2}T^{5/2}
                    E_{\mathrm{top}}^{-z/4}E_*^{-(f-z)/4}.
 \tag{AC22}
\]

These are the actual selected degrees, not a substitute generic high-weight parameter. A role-count hypothesis is imposed only in the conditional corollary below.

\subsection{5. Exact bound and new removed subsets}\label{exact-bound-and-new-removed-subsets}

Combining AC11, AC14--16, the weighted incidence count and AC22 gives the following upper bound for the absolute sum at a fixed u,d and at least z top labels:

\[
\begin{split}
 \mathcal B_{d,u}(z)
 ={}&\binom f5\binom5u\binom md G_n A_dB_0T^{10d+50}\\
 &\cdot\frac{2[(H+1)/d]^{af/2}(H+1)^{(a-1)f}
                  T^{5a+(a-1)(5+u)}}
                 {Q^{(2a-1)(f-1)}}\\
 &\cdot12^f2^{f/2}T^{5/2}
                      E_{\mathrm{top}}^{-z/4}E_*^{-(f-z)/4}.
\end{split}
 \tag{AC23}
\]

For fewer than z top labels the bound is not asserted. For more, monotonicity follows from E\_top\textgreater E\_*.

The exact directed-rational certificate verifies the following thresholds for z=0, uniformly over every role assignment. Every listed cell has B\_(d,u)(0)\textless2\^{}−4096:

\begin{longtable}[]{@{}
  >{\raggedleft\arraybackslash}p{(\columnwidth - 6\tabcolsep) * \real{0.2500}}
  >{\raggedleft\arraybackslash}p{(\columnwidth - 6\tabcolsep) * \real{0.2500}}
  >{\raggedleft\arraybackslash}p{(\columnwidth - 6\tabcolsep) * \real{0.2500}}
  >{\raggedleft\arraybackslash}p{(\columnwidth - 6\tabcolsep) * \real{0.2500}}@{}}
\toprule\noalign{}
\begin{minipage}[b]{\linewidth}\raggedleft
Unique fixed places u
\end{minipage} & \begin{minipage}[b]{\linewidth}\raggedleft
Off-diagonal places t
\end{minipage} & \begin{minipage}[b]{\linewidth}\raggedleft
All active sizes removed
\end{minipage} & \begin{minipage}[b]{\linewidth}\raggedleft
Weakest certified binary margin
\end{minipage} \\
\midrule\noalign{}
\endhead
\bottomrule\noalign{}
\endlastfoot
1 & 4 & 1\textless=d\textless=162 & 9259 \\
2 & 3 & 1\textless=d\textless=144 & 13322 \\
3 & 2 & 1\textless=d\textless=127 & 8162 \\
4 & 1 & 1\textless=d\textless=109 & 12694 \\
5 & 0 & 1\textless=d\textless=92 & 7902 \\
\end{longtable}

There are634 such cells, so their combined absolute contribution is less than \textbf{2\^{}−4086}. The direct checks cover every d, rather than assuming monotonicity of the volume expression.

If \textbf{at least119 labels are top}, the same certificate verifies B\_(d,1)(119)\textless2\^{}−4096 for every1\textless=d\textless=164. This adds the two previously unremoved cells d=163,164 at u=1. The weakest certified binary margin is4192, attained at d=164; the corresponding118-label formula has margin4075. The combined636-cell bound remains less than \textbf{2\^{}−4086}. The fixed formula at118 top labels does not reach this chosen 4096-bit threshold for d=164; this is a limit of the displayed estimate, not a necessity claim for the actual kernel.

The corollary applies in particular to the all-top actual mixed profile. K041 Section1 permits an arbitrary partition of the labels, and K044(EG6,EG10) carries all role assignments and words to this upper gap and its twists. K040(NC1--4) constructs genuine unitary stage characters for every e in (28Z)\^{}Psi. In the all-top case their infinity vectors are the uniform multiples E\_T and E\_T+2ell+q, so the exact identity alpha\_(c1)=eta\^{}(-2c) applies. The original K032 induced primitive CM type, principal variety, finite projectors and finite levels are unchanged. A numerical role label is not a replacement of that CM type by a quadratic one. The conclusion concerns the carried mixed-kernel/profile constructor, not a literal Frame, Host, D7 or FamP record; LIFT54--55's m=2 exclusions are not bypassed.

\subsection{6. Residual set, normalization and verification scope}\label{residual-set-normalization-and-verification-scope}

Within the scalar-cross five-place case, the uniform remaining active sizes are

\[
 (u,d)\in
 \{1\}\times\{163,164\}\ \cup
 \{2\}\times\{145,\ldots,164\}\ \cup
 \{3\}\times\{128,\ldots,164\}\ \cup
 \{4\}\times\{110,\ldots,164\}\ \cup
 \{5\}\times\{93,\ldots,164\}.
 \tag{AC24}
\]

With at least119 top labels, the first set is removed. The fixed-W0 all-type-II terms and higher support sizes are not bounded by this result. Positive local type-II rho does not give a sign to the full summand, which also contains the genuine G partial trace and the real selected trace.

The new absolute error can be added to K071's retained error: adding2\^{}−4086 to its explicit terms still leaves the same loose overall budget2\^{}−1998. No value of S\_T or sign of the projected Gram is computed. The actual scalar c, the genuine mixed projectors, the central phase, the PBW metric and every support, class-fibre and ordered-incidence factor are retained throughout.

Primary input use is exactly the already checked K060 application of Haagerup and Schlichtkrull (2014, Theorem1.1, p.228) and the actual K043 application of Nakahama (2016, Theorem5.1, pp.960--961). Their norm theorems are not rederived. The new energy refinement follows from the displayed RC3 grading; the arithmetic witnesses, active-coordinate count and exact selected application are proved above. The actual fractional-lattice and fibre estimates are the settled DO063/K071 results in their unchanged metrics.

The companion \nolinkurl{active_coordinate_scalar_cross_certificate.py} uses exact rational upper and lower binary logarithms. It carries K071's 192-bit directed atanh expansion, with80 terms and its explicit tail bound; factorials, field constants and support combinatorics are exact integers. It tests every u,d cell, the 119-label criterion, every binary-word count and both twist endpoints, and rational unipotent models checking AC6--8's normalizations. It also checks nonproportional exact integer instances of generalized Holder. These finite instances diagnose conventions; the symbolic proofs above establish the unrestricted statements. No global rational-matrix support enumeration, actual period computation or Lean build is claimed.

\subsubsection{References}\label{references-5}

Haagerup, U., \& Schlichtkrull, H. (2014). Inequalities for Jacobi polynomials. \emph{The Ramanujan Journal, 33},227--246. \url{https://doi.org/10.1007/s11139-013-9472-4}. Theorem1.1 and its immediately following C\textless12 statement, p.228. Primary author version: \url{https://arxiv.org/pdf/1201.0495v2}. This is the already checked K060 primary theorem, applied with the same scalar diagonal coefficient and exact Hilbert conventions.

Nakahama, R. (2016). Norm computation and analytic continuation of vector-valued holomorphic discrete series representations. \emph{Journal of Lie Theory, 26}(4),927--990. \url{https://jolt.centre-mersenne.org/item/10.5802/jolt.919.pdf}. Theorem5.1 and action(5.5), pp.959--961. Its accepted K043 application fixes the selected isometry used here; no different fibre metric is introduced.

\subsubsection{Author whole-proof check}\label{author-whole-proof-check-5}

I read and reconstructed the complete saved AC1--24 report, all six sections and references, and the complete final certificate. The checks include the actual fixed-coordinate normal form, positive square-diagonal witness and Cartan bound, the Holder-selected product with its paired prime valuation, the exact active fractional lattice and original class-fibre bound, every incidence weight, the full homogeneous compression in the original PBW metric, and the exact all-word role degrees. The actual all-top mixed character/profile constructor and its distinction from a literal Host/FamP record were checked directly against K040, K041 and K044. The final5563-assertion replay passed at 2026-11-22T11:30:08.976149+00:00. During construction, the diagnostic caught an incorrectly transcribed top-degree expression and two unsimplified symbolic complex-matrix comparisons; these were corrected before the final proof and replay. No global support enumeration, Mram positivity, or Lean build is certified.

\clearpage

\section{AM-DO-065 --- the globally unipotent exception is exactly the scalar cross case}\label{am-do-065-the-globally-unipotent-exception-is-exactly-the-scalar-cross-case}

Author: GPT-6 Astra Pro, ALP\_MACHINE \nolinkurl{/root/definite_oscillator}. 22 November 2026 UTC. Inherited reasoning effort is not separately exposed.

\textbf{Author whole-proof check PASS --- 2026-11-22T11:29:53.249751+00:00.}

For the unchanged J56 remainder, every globally unipotent matrix in the K069 exceptional class has A=I and full nilpotent rank two. Thus all globally unipotent nonscalar positive blocks belong to K069's block-determined fixed-space class. The previously missing case rank(A−I)=1 can then be bounded using the actual rank-two Green estimate.

The resulting removed union has absolute contribution less than 2\^{}-4088. After it and the already checked cuts are removed, every remaining term with s≤22 and t≤8 has A=I. Here s is total nonintegral T support, while t counts unequal-pencil off-diagonal places. This is a structural reduction and an absolute subset bound, not positivity of the remaining scalar-cross or fixed-W0 terms.

\subsection{1. Two ranks and the actual conditions}\label{two-ranks-and-the-actual-conditions}

Keep the actual K062 normalization and K068 residual, with

\[
 a=f=168,\quad n=169,\quad m=164,\quad H=8192,
 \qquad h=\begin{pmatrix}A&C\\B&D\end{pmatrix}_{S\perp W_0},
 \qquad X=A-I.
 \tag{UR1}
\]

The large positive block A is integral. Write

\[
 R=\operatorname{rank}_E(h-I),\qquad
 r=\operatorname{rank}_E X,\qquad
 U=\ker X,\qquad K_A=\ker X\cap\ker X^*.
 \tag{UR2}
\]

These two ranks are not identified. The K069 fixed-space class is F=\{h fixes K\_A pointwise\}. Its exact complement is IC8, restated in DO064(PM3). In that complement the block identity gives

\[
                       A^*A\le I
 \tag{UR3}
\]

at every complex realization: B(S) belongs to the perpendicular of the nonzero isotropic line B(U), whose quotient has positive form. This is the already proved contraction in DO064(PM6).

K070 proves that t≤8 forces global unipotence on this actual residual, independently of s. We use that precise hypothesis. The carried signature-one normal form DO061(TR19--21) gives, for a nonidentity globally unipotent h and N=h−I,

\[
 N^3=0,\quad R\le2,\quad
 \begin{cases}
 N^2=0,&R=1,\\
 \operatorname{rank}N^2=1,&R=2.
 \end{cases}
 \tag{UR4}
\]

In particular N²=h+h\^{}-1−2I, and the positive block of the unitary inverse is A*. Therefore

\[
                         (N^2)_{SS}=X+X^*.
 \tag{UR5}
\]

\subsection{2. Rigidity of the exceptional cross direction}\label{rigidity-of-the-exceptional-cross-direction}

\textbf{Theorem.} A globally unipotent matrix outside F has R=2 and A=I. Conversely, a globally unipotent matrix with A=I lies outside F exactly when B is nonzero. Thus the exception in the unipotent sector is precisely the rank-two scalar-cross case already isolated in DO061.

\textbf{Proof.} Outside F choose u in K\_A such that Nu is nonzero. Since Xu=0, the vector Nu=(0,Bu) belongs to W0. Since also X*u=0, UR5 gives

\[
                            (N^2u)_S=0.
 \tag{UR6}
\]

First suppose R=1. Then N²=0 and unitarity gives N-dagger=−N. The line im(N) is spanned by the nonzero vector Nu in W0. Hence all S rows of N vanish. Its adjoint relation forces all S columns to vanish as well, contradicting Nu≠0 for u in S. Thus R cannot be one.

Suppose R=2. By UR4, the line ell=im(N²) equals im(N) intersect ker(N): the restriction of N to its two-dimensional image has rank one and a one-dimensional kernel. If N²u is nonzero, UR6 puts ell in W0. If N²u is zero, then the nonzero Nu belongs to im(N) intersect ker(N), and again ell belongs to W0. In either case all S rows of N² vanish. UR5 yields A+A*=2I.

At each complex realization, combine this with the actual contraction UR3:

\[
       (A-I)^*(A-I)=A^*A-A^*-A+I=A^*A-I\le0.
 \tag{UR7}
\]

The left side is positive semidefinite, so it vanishes and A=I. The equality at an embedding is faithful over E; in fact it holds at every embedding as stated. This proves the forward assertion.

When A=I, K\_A=S. Its pointwise fixation is exactly B=0. If B≠0, the unitary block identity makes B(S) totally isotropic, and K069's criterion identifies the same exception. The rank-one possibility was already excluded, so R=2. ∎

This proof does not infer that an arbitrary unit eigenvector of A is fixed by h. The one possible isotropic cross direction is treated through the actual image and kernel of N. It does not require an integral moving basis, a choice of finite characteristic, or the partial-monomial enumeration of DO064.

It follows that a globally unipotent matrix with A≠I belongs to F. Moreover r≤R≤2. Its only possible rank pairs are

\[
                         (R,r)=(1,1),(2,1),(2,2).
 \tag{UR8}
\]

\subsection{3. The missing F case: R=2 and r=1}\label{the-missing-f-case-r2-and-r1}

For a rank-two unipotent with X≠0, DO061(TR25) already proves the actual Green estimate

\[
 |a(h)|\le
 \frac{2(H+1)^{af}[2(H+1)]^{(a-1)f}}
      {Q^{(2a-1)(f-1)}}.
 \tag{UR9}
\]

Both witnesses lie in the integral blocks X and (N²)\_SS=X+X*. The latter is nonzero in this case, as checked in that proof. Thus UR9 has no T numerator. Its factor two is the single global Euler correction; the two mixed projectors retain their norm bound one.

Apply K069(IC11--15) at r=1. The number of possible integral blocks is bounded by m²(2H+3)\^{}(2f(2m−1)). For each fixed A in F the positive part of the last column is determined by two original S coordinates, and its five W0 coordinates are retained. The exact target over-lattice therefore has T\^{}(14s) cosets. This gives

\[
 \begin{split}
 C_{2,1}(s)={}&m^2(2H+3)^{2f(2m-1)}\binom fs(T-1)^{2s}G_n\\
             &\quad\cdot(2H+1)^{14f}T^{14s}.
 \end{split}
 \tag{UR10}
\]

No inverse minor acts on an integral lattice: the rank-one row/column chart only records original integral entries. The two-coordinate projection is chosen once for that A and adds no basis-height factor. The finite last-column fibre is the original G\_n\textless2\^{}5022, not a new stabilizer. Multiplying UR10 by UR9 and the actual bound (2/(T−2))\^{}s gives a complete upper bound B\_(2,1)(s) for the intersection of this rank pair with any chosen unipotent subset.

Here are conservative exact rational exponents. Set

\[
 \begin{split}
 l_X&=14+1/2048,\quad l_Y=14+1/8192,\\
 l_H&=13+1/4096,\quad l_K=14+1/4096.
 \end{split}
 \tag{UR11}
\]

The corresponding upper bounds for log2(2H+3), log2(2H+1), log2(H+1), and log2(2(H+1)) follow by exact integer powers. Also m²\textless2\^{}15, Q\^{}4\textgreater2\^{}281, T\textless2\^{}996, and 2(T−1)\^{}2/(T−2)\textless2\^{}997. Consequently

\[
 \begin{split}
 \log_2 B_{2,1}(s)&<C_{2,1}+14941s,\\
 C_{2,1}&=f+15+5022+2f(2m-1)l_X+14fl_Y+1\\
 &\quad+afl_H+(a-1)fl_K-(281/4)(2a-1)(f-1)\\
 &=-408071827/256.
 \end{split}
 \tag{UR12}
\]

The right side is at most−4096 for every s≤106. Thus every actual (R,r)=(2,1) cell with t≤8 and 5≤s≤106 has absolute contribution less than2\^{}-4096. This statement does not assert global unipotence for t≥9.

\subsection{4. Rank one and the carried rank-two block case}\label{rank-one-and-the-carried-rank-two-block-case}

For (R,r)=(1,1), use the injective pivot-column count and the actual character bound in DO061(TR28--30):

\[
 \begin{split}
 C_{1,1}(s)&\le\binom fs m(2H+3)^{2fn}T^{10s},\\
 |a(h)|&\le2(H+1)^{af}/Q^{a(f-1)}.
 \end{split}
 \tag{UR13}
\]

There is no extra pencil-pair or finite-fibre factor in this injective rank-one encoding. With the original rho factors, the same rational bounds and 2T\^{}10/(T−2)\textless4T\^{}9 give

\[
 \begin{split}
 \log_2 B_{1,1}(s)&<C_{1,1}+8966s,\\
 C_{1,1}&=f+8+2fnl_X+1+afl_H-(281/4)a(f-1)\\
 &=-103530801/128.
 \end{split}
 \tag{UR14}
\]

This is at most−4096 for every s≤89. Hence all (1,1) cells with t≤8 and 5≤s≤89 have absolute contribution less than2\^{}-4096.

For (R,r)=(2,2), the theorem places every such unipotent matrix in F. The already checked K069(IC20--21) gives the same cell bound through s=22. We retain that actual r=2 minor estimate, rather than replacing it by the weaker rank-one estimate. No claim is made that the new UR12 formula applies at r=2.

\subsection{5. Exact removed union and remaining conditions}\label{exact-removed-union-and-remaining-conditions}

Remove the following disjoint actual subsets, or their set differences with earlier removed sets:

\begin{longtable}[]{@{}rrrl@{}}
\toprule\noalign{}
Full rank R & Block rank r & Supports & Additional condition \\
\midrule\noalign{}
\endhead
\bottomrule\noalign{}
\endlastfoot
1 & 1 & 5≤s≤89 & t≤8 \\
2 & 1 & 5≤s≤106 & t≤8 \\
2 & 2 & 5≤s≤22 & t≤8 \\
\end{longtable}

There are85+102+18=205 rank/support cells. Each is uniformly bounded by2\^{}-4096, including all its allowed t values, all actual labels and all class weights. Therefore their union has total absolute contribution

\[
                     \boxed{205\,2^{-4096}<2^{-4088}.}
 \tag{UR15}
\]

The t condition is not counted again as a multiplicity: each cell is one subset of the original support-s count, not a sum over independently repeated arrays. Earlier complete s≤4 bounds remain unchanged. It follows that, after these removals, every remaining nonzero term with s≤22 and t≤8 has A=I. In particular this holds throughout the complete strata s=5,6,7,8, since t≤s there.

The scalar-cross terms and the terms supported on W0 retain their actual signs and phases. K071 removes specified subsets of them, but this report does not assert that its remaining unique-fixed configurations vanish or have nonnegative total. The global-unipotent restriction, full rank R, block rank r, fixed-space F, and source-size scope are all retained.

\subsection{Verification and carried inputs}\label{verification-and-carried-inputs-1}

The companion certificate checks the exact integer logarithmic bounds, both rational constants and every support threshold, the205-cell union, and signature-one unipotent models with all three rank pairs. In those models it checks the actual intersection K\_A, the exceptional-line argument and the rank-one pivot reconstruction. These are proof diagnostics, not an enumeration of the actual rational sum.

All nontrivial carried inputs have already been whole-checked: DO061(TR19--30), K068's actual residual, K069(IC1--21), K070's t≤8 criterion, and DO064(PM6,PM20). The new rigidity theorem and the new use of the fixed-space r=1 count are proved above. No external representation theorem, new finite projector, alternate measure, literal SizesAt constructor, or Lean build is introduced.

\subsection{Author verification record}\label{author-verification-record-17}

At 2026-11-22T11:29:53.249751+00:00, I completed the whole saved proof and final certificate read. I reconstructed the actual image/kernel line argument, contraction at every embedding, both distinct ranks, the source-aligned fixed-space count, both Green inputs, all rational constants and the205-cell union. The final certificate passed1,125 exact assertions at2026-11-22T11:28:33.198361+00:00, with144 unitary models,16 nonzero exceptional cases and33 rank-one pivot reconstructions. The preliminary diagnostic needed explicit simplification of exact complex expressions and a second positive-coordinate rotation to cover block rank two; those were certificate-model corrections, not changes to the theorem. No preliminary failed diagnostic is reported as a PASS. The final status and this record were reread before freezing.

\clearpage

\section{AM-R-044: paired CM valuations remove the complete five-place cross sector}\label{am-r-044-paired-cm-valuations-remove-the-complete-five-place-cross-sector}

Author: /root, GPT-6 Astra Pro, ALP\_MACHINE, 22 November 2026 UTC. Reasoning effort is not separately exposed.

Status: author and independent whole-proof checks PASS. Independent report: \nolinkurl{agents/review_paired_cm_cross_rank_loss.md}.

This finding concerns the original J56 short sum of K062 at the upper K044 endpoint and cutoff H=8192. It changes neither its projectors nor its signed local incidence factors. The new observation is that one original cross entry, together with its CM conjugate, controls both possible losses of rank at an auxiliary prime. Counting both primes removes the extra type-II pole loss in the earlier bound.

The resulting estimate removes 729 complete (s,v) cross subsets, each already summed over every number of type-II incidences. Their absolute union is less than 2\^{}-4086. In particular the complete scalar-cross sector at s=5 is removed. After the earlier checked exclusions, every remaining five-place term is supported on the fixed five-coordinate space W0 and has five unique-fixed incidences. Its rank is one or two; no sign or nonzero global period is asserted for those remaining terms.

\subsection{1. The retained matrix, places and normalizations}\label{the-retained-matrix-places-and-normalizations}

Use the settled notation of DO063 and K068--071:

\[
 a=f=168,\quad n=169,\quad m=164,\quad H=8192,\quad
 K=2(H+1),\quad E^n=S\perp W_0,
 \qquad\operatorname{sig}(W_0)=(4,1).
 \tag{PC1}
\]

The Hermitian matrix in these original coordinates is J=diag(I\_a,-1). The rational auxiliary primes are g0=22285932761 and g1=64072175441, and Q=g0*g1. The actual prime T is the unchanged 996-bit DO040 prime. There are f−1 unmixed relative places over each auxiliary prime, hence 334 finite Green factors. Each relative place of F has two conjugate primes of E. The other two actual finite factors retain their original mixed projectors and the proved absolute bound one.

Let s be the total nonintegral T support. Let v be the number of places with unequal full pencils, u the number among them with a unique fixed actual lattice, and t=v−u the number with a unique off-diagonal incidence. Type I comprises the remaining s−v places. These are precisely the types of K070--071, not a new partition of the arithmetic data.

We first consider a globally unipotent scalar-cross term. Write its unique central-normalized unipotent representative as h=I+N. Thus

\[
 N^3=0,\quad\operatorname{rank}_E N=2,\quad
 N_{SS}=0,\quad N_{SW_0}\ne0.
 \tag{PC2}
\]

All phases removed in choosing this representative already have absolute value one in the original K062 normalization. They are not discarded in the signed sum. Global unipotence will follow from the existing t\textless=8 criterion for every subset ultimately used here.

The signature-one normal form DO061(TR32), used in K071(NB10), puts the isotropic image line of N² in W0 in this scalar-cross case. Consequently

\[
 (N^2)_{S,E^n}=(N^2)_{E^n,S}=0.
 \tag{PC3}
\]

The same normal form gives a nonzero original last-column entry x=N\_in for some i in S. Indeed its isotropic line in W0 has a nonzero negative coordinate, and the cross action is nonzero. The carried bound is \textbar tau(x)\textbar\textless=sqrt(K) at every complex realization. No coordinate is chosen separately at different embeddings.

\subsection{2. The principal minor and the two conjugate primes}\label{the-principal-minor-and-the-two-conjugate-primes}

Unitarity and N³=0 give

\[
 N^\dagger=N^2-N,\qquad N^\dagger=J^{-1}\overline N^{\,t}J.
 \tag{PC4}
\]

For i positive and n the negative coordinate, PC3--4 imply N\_ni=bar(x). Since N\_ii=0, the original principal minor is exactly

\[
 \det\begin{pmatrix}N_{ii}&N_{in}\\N_{ni}&N_{nn}\end{pmatrix}
                         =-x\overline x.
 \tag{PC5}
\]

Fix an unmixed relative place w of F above g, with conjugate primes P and bar(P) of E. The actual matrix is integral there. Its reduction is unipotent by K068 and has rank d=rank(N mod P) in \{0,1,2\}, because the global rank is two. The ranks at P and bar(P) agree: the unitary inverse relation exchanges them by transpose and multiplication by invertible matrices. Conjugation is not being treated as an internal operation on one arbitrarily selected split factor.

Define

\[
 \nu_w=v_P(x)+v_{\overline P}(x).
 \tag{PC6}
\]

The valuations are nonnegative. If d=0, both N\_in and N\_ni vanish modulo P, so both summands in PC6 are at least one. If d=1, the principal minor PC5 vanishes modulo P, so at least one summand is at least one. Thus

\[
 d=0\Longrightarrow\nu_w\ge2,\qquad
 d=1\Longrightarrow\nu_w\ge1.
 \tag{PC7}
\]

This is the additional information lost by counting just one chosen prime of E for a rank-one witness.

The retained normalized Green estimate uses exponents

\[
 L_0=0,\qquad L_1=a=168,\qquad L_2=2a-1=335.
 \tag{PC8}
\]

Its total Euler correction is less than two, as in K067 and K071. The elementary inequalities

\[
 L_2-L_0=335=\frac{335}{2}\,2,\qquad
 L_2-L_1=167<\frac{335}{2}
 \tag{PC9}
\]

show, in all three cases, that

\[
                       L_2-L_d\le\frac{L_2}{2}\nu_w.
 \tag{PC10}
\]

This includes rank-two reductions with square zero. There is no second witness from N² whose accidental vanishing must be charged.

\subsection{3. The actual denominator ideal and the finite partial trace}\label{the-actual-denominator-ideal-and-the-finite-partial-trace}

Let U be the v unequal-pencil places and put

\[
 \mathfrak d_U=\prod_{w\in U}\mathfrak p_w\overline{\mathfrak p}_w,
 \qquad I_U=\mathfrak d_U^{-1},\qquad
 \mathrm N_E(\mathfrak d_U)=T^{2v}.
 \tag{PC11}
\]

K071(NB5,NB9,NB12) proves x belongs to I\_U: cross entries are integral at type I and have at most simple poles at types II and III. Both CM partners occur in PC11. The ideal \(\mathfrak a=(x)\mathfrak d_U\) is integral and nonzero, even if the denominator ideal is not principal. Its absolute norm is

\[
 \mathrm N_E(\mathfrak a)
       =T^{2v}\prod_{\tau\in\Psi}|\tau(x)|^2
       \le T^{2v}K^f.
 \tag{PC12}
\]

At every selected G prime the denominator ideal is a unit, so its valuation pair in \(\mathfrak a\) is the one in PC6. All other prime factors of this integral ideal have norm at least one. Therefore

\[
 \prod_{w\in G_{\rm unmixed}}g(w)^{\nu_w}
                         \le\mathrm N_E(\mathfrak a).
 \tag{PC13}
\]

Multiply the actual Green bounds, use PC10--13, and retain the two mixed bounds. For the same normalized finite partial trace a(h) that occurs in K062, this proves

\[
 \boxed{\quad
 |a(h)|\le
 \frac{2\,\mathrm N_E(\mathfrak a)^{335/2}}
      {Q^{335(f-1)}}
 \le\frac{2K^{335f/2}T^{335v}}{Q^{335(f-1)}}.
 \quad}
 \tag{PC14}
\]

The exponent depends on v and is independent of u. This is a bound on the original partial trace, with its projectors unchanged. The estimate makes no assumption about a distribution of valuations or about independence of the two primes.

\subsection{4. Actual columns, incidence weights and real compression}\label{actual-columns-incidence-weights-and-real-compression}

We use the previously checked K071(NB12--16) count and do not reprove its geometry of numbers. In its half-trace Euclidean metric, set

\[
 D_0=7^{140}29^{144}5^{84}13^{84},\qquad
 R_v=90.516+9750/r_v,
 \tag{PC15}
\]

where the exact positive lower roots r\_v satisfy r\_v\^{}f\textless=T\^{}v:

\begin{longtable}[]{@{}rrrrrrr@{}}
\toprule\noalign{}
v & 0 & 1 & 2 & 3 & 4 & 5 \\
\midrule\noalign{}
\endhead
\bottomrule\noalign{}
\endlastfoot
r\_v & 1 & 243/4 & 3688 & 224256 & 13631488 & 828375040 \\
\end{longtable}

Put

\[
 A_v=\min\left\{183^{2fm},
 \frac{[(22/7)fR_v^2]^{fm}}{(fm)!D_0^m}\right\},
 \qquad
 B_0=\frac{[(22/7)f9543^2]^{5f}}{(5f)!D_0^5}.
 \tag{PC16}
\]

For every fixed actual source/target pencil array the original last-column count is at most A\_v B0 T\^{}(2mv+10s), with its original isometry-fibre bound G\_n=gcd(\textbar GL\_n(F\_g0)\textbar,\textbar GL\_n(F\_g1)\textbar). The W0 factor is the paired lattice of DO063; it is not replaced by a fictitious coordinate lattice. The selected real factor is

\[
                  |c(h)|<12^f2^{f/2-485f/4}T^{v/2}.
 \tag{PC17}
\]

This is the full selected compression in the original PBW metric at the unchanged near-origin point, already proved in DO063/K071. No top-role hypothesis or improved energy estimate is needed here.

K071(NB7--8) gives the weighted support/type/ordered-pair bound

\[
 W_{s,v,u}=\binom fs\binom sv\binom vu(T-1)^{s-v}.
 \tag{PC18}
\]

It already includes every actual local pair before cancellation against the magnitude of its original rho factor. Summing PC18 over all 0\textless=u\textless=v replaces binom(v,u) by 2\^{}v. Since PC14 and the other bounds are uniform in u, the complete scalar-cross (s,v) subset has absolute contribution at most

\[
 \begin{split}
 C_{s,v}={}&\binom fs\binom{s}{v}2^v(T-1)^{s-v}
       G_nA_vB_0T^{2mv+10s}\\
 &\quad\cdot\frac{2K^{335f/2}T^{335v}}{Q^{335(f-1)}}
       \cdot12^f2^{f/2-485f/4}T^{v/2}.
 \end{split}
 \tag{PC19}
\]

This is a sum over all u, not an estimate for just one chosen incidence pattern. Restricting to a prior remainder only decreases the bound. For v\textless=5, every possible t=v−u is at most five, so K070 supplies global unipotence for all these actual residual terms even when s is larger than the former total-support threshold.

\subsection{5. Exact thresholds and the remaining five-place terms}\label{exact-thresholds-and-the-remaining-five-place-terms}

The companion certificate uses the checked K071 directed-rational binary-logarithm routine, with its 192-bit scale, 80-term atanh series and explicit tail bound. It evaluates PC19 at every integer 5\textless=s\textless=168 and 0\textless=v\textless=5. It checks each passing integer separately. The following are its maximal passing endpoints for threshold4096:

\begin{longtable}[]{@{}rrr@{}}
\toprule\noalign{}
v & All supports removed & Lower binary margin at endpoint \\
\midrule\noalign{}
\endhead
\bottomrule\noalign{}
\endlastfoot
0 & 5 through168 & 1346434 \\
1 & 5 through168 & 978607 \\
2 & 5 through168 & 397916 \\
3 & 5 through143 & 14328 \\
4 & 5 through83 & 11808 \\
5 & 5 through23 & 9396 \\
\end{longtable}

Each listed complete subset has C\_(s,v)\textless2\^{}-4096. There are 3*164+139+79+19=729 subsets, so

\[
             \sum_{\text{listed }(s,v)}
             \sum_{[h]\text{ in that subset}}|w(h)|
                   <729\,2^{-4096}<2^{-4086}.
 \tag{PC20}
\]

At s=v=5 the conservative formula has binary margin206571. The sharper K072 degree or active-coordinate bound is unnecessary for this conclusion, although that separate result remains valid. These subsets may overlap previously bounded subsets; remove their intersections with the current remainder and charge PC20 once.

The earlier complete s\textless=4 estimates remain unchanged. At s=5, DO062 already excludes the nonscalar positive blocks, and K071 excludes every fixed-W0 cell except u=v=5. PC20 now excludes every scalar-cross cell. Thus the five-place part of the new remainder consists only of

\[
 h=I_S\oplus h_0,\quad h_0\in U(W_0)(F)
 \text{ unipotent},\quad\operatorname{rank}(h_0-I)\in\{1,2\},
 \qquad u=v=s=5,\ t=0.
 \tag{PC21}
\]

Here U(W0)(F) denotes the original unitary group over F with E coefficients. The unique-fixed local rho values are positive, but the actual G partial trace and selected real trace still occur. Their full signed sum is not evaluated by PC21. Adding PC20's absolute error to the current joint remainder still permits the loose global error budget2\^{}-230, as checked at integration.

No shorter source-admissible chain, arithmetic cycle identification, DP44 sign, or Hecke RH follows from this subset estimate. In particular LIFT54--55's full source-size exclusions remain in force.

\subsection{6. Verification, exact sources and geometric interpretation}\label{verification-exact-sources-and-geometric-interpretation}

New content is PC5--14 and its application PC19--21. The sources carried without re-proving are DO061(TR19--21,TR32--34), K067's normalized Green estimate, K068's actual auxiliary-prime reductions, K070's t criterion, and K071(NB5--18) for actual paired poles, counts and real normalization. Their primary Green reference is Green, J. A. (1955). The characters of the finite general linear groups. \emph{Transactions of the American Mathematical Society, 80}(2),402--447, section5, Example(ii), pp.430--431. The original selected-norm and geometry-of-numbers references remain those explicitly recorded in DO063 and K071.

The exact diagnostic checks the three rank-loss inequalities for 170 ranks, 4023 finite two-by-two models including both independent split entries, all actual constants and every support bound. Those finite checks diagnose the convention; PC4--13 prove the general valuation statement. No global rational support is enumerated.

Geometrically, PC5 is a determinantal equation in the existing Hermitian vector module. PC11--13 use its two conjugate integral fibres and the norm of an actual invertible ideal; they do not replace them by one prime or assert a new arithmetic correspondence. The carried SGA3 I4.7.3/4.7.3.1 vector-module dictionary applies. There is no new l-adic purity, CM-motive identification, or comparison with the specific032 arithmetic theta cycle.

Author whole-proof reading PASS --- 2026-11-22T11:38:27.442293+00:00. Complete PC1--21 and all six sections, plus the entire certificate, reread after writing. Reconstructed the split-prime adjoint relation, rank loss at both CM partners, nonprincipal denominator-ideal norm, actual Green product, every weighted support factor, all729 cells and the exact five-place scope. The unchanged exact replay passed15481 assertions at2026-11-22T11:31:30.800180+00:00. Final mathematical notation edits were reread. No global sum or period was evaluated.

Independent review PASS --- 2026-11-22T11:41:31+00:00, /root/kummer\_extension, GPT-6 Astra Pro, inherited effort not separately exposed. The whole proof, complete script, actual carried interfaces and own15481 replay were checked; no correction was required. The reviewed mathematical text had SHA256 2e559720355afc443a52869918b126103d020c4c7964477cc2884611421a5868. This final edit changes only the review-status line and adds this record; the proof and script are unchanged.

\clearpage

\section{ALP-LIFT-57: the prescribed G and Q inputs through every vertical stage}\label{alp-lift-57-the-prescribed-g-and-q-inputs-through-every-vertical-stage}

Author: /root/lift\_accounting, GPT-6 Astra Pro, ALP\_MACHINE. Status: author whole-proof check passed 2026-11-22T11:34:29.542481+00:00. Independent review is separate.

\subsection{1. Local scope and the fixed source functions}\label{local-scope-and-the-fixed-source-functions}

Fix a nonarchimedean local field \(F_v\), its ring \(O\), uniformizer \(\varpi\), and residue field of order \(q\). The CM algebra is split at this place. These include every actual selected G or tuple Q place. Put

\[
 n_j=r+jt,\quad b_j=n_j+t,\quad z_*=(t+1)/2,\quad
 \epsilon_j=j\bmod2,\qquad r,t\text{ odd},\quad t\ge r+5.
 \tag{VS1}
\]

Here \(t\) is the actual arithmetic rank increment; equivalently it is the value of TPar.t in the canonical ofWord parameter. No use of the unrelated three\_m bound is needed for this local theorem.

Write \(\mu_j\) for the unitary character at edge \(j\), meaning the source's index \(j+1\). Retain the same fixed unitary odd character \(\eta\) and unitary tuple character \(\xi\). The conjugated upward kernel has target and source labels

\[
 \alpha_j=(\mu_j\xi)^{-1}\eta^{1-\epsilon_j},\qquad
 \beta_j=(\mu_j\xi)^{-1}\eta^{\epsilon_j}.
 \tag{VS2}
\]

This is the all-stage form of the actual chrStep convention in LIFT31 (VA5) and DO041 (DO41.4). In the fixed split half-model its honest action is

\[
 \omega_j(g,h)f(X)=
 \alpha_j(\det h)\beta_j(\det g)^{-1}
 |\det g|^{b_j/2}|\det h|^{-n_j/2}f(h^{-1}Xg).
 \tag{VS3}
\]

There is no square of a splitting character in (VS3): at a split norm-one determinant \((c,c^{-1})\), the root is \((c,1)\). The source identification with the actual normalized lattice vacuum is the accepted DO020/022 and DO041 half-position realization. All matrix coordinates below are the literal frMapO coordinates \(\beta_T^{-1}X\beta_S\); no new global lattice frame is chosen.

At G, \(\eta,\xi\) are unramified and \(\mu_j|_{O^\times}=\chi_{j+1}^{-1}\). At Q, \(\eta,\mu_j\) are unramified and \(\xi\) is the prescribed tame character. These hypotheses are the literal local character ties in LinksChar/DatumVec; the local conclusions apply whenever those ties are realized. They do not construct global characters.

Normalize additive measure on each matrix lattice to one, and group Haar by \(\operatorname{vol}K_n=1\), \(K_n=\mathrm{GL}_n(O)\). Set

\[
 \kappa_n=\prod_{i=1}^n(1-q^{-i}),\quad \kappa_0=1,\qquad
 dg=\kappa_n^{-1}|\det g|^{-n}dg_{\rm add},\qquad
 {n\brack r}_q=q^{r(n-r)}
       \frac{\kappa_n}{\kappa_r\kappa_{n-r}}.
 \tag{VS4}
\]

Every scalar below refers to this fixed normalization of the local quotient integral. Its identification with another chosen local Howe quotient may multiply it by a nonzero comparison scalar. Nothing sets a global IsIo or arithmetic period pin to one.

\subsection{2. Whole-Schwartz quotient and all-stage holomorphy}\label{whole-schwartz-quotient-and-all-stage-holomorphy}

For an irreducible unitary representation \(\Pi\) of \(\mathrm{GL}_n(F_v)\) and \(b=n+t\), let \(E_n=(I_n,0)^{\mathsf t}\). The accepted whole-Schwartz construction LIFT24 (MT10--16), LIFT27 (MR8--14) and DO041 (DO41.7--11) is

\[
 I_{\beta,z}(f,v)=\int_{\mathrm{GL}_n}
 f(E_ng)|\det g|^{z+(n-1)/2}\beta(\det g)^{-1}\Pi(g)v\,dg,
 \qquad
 Q_{\alpha,\beta}(f,v)(h)
 =I_{\beta,z_*}(\omega(1,h)f,v).
 \tag{VS5}
\]

At a point where \(L(z,\Pi\beta^{-1})\) is finite, its Godement--Jacquet continuation is a genuine smooth vector. One first works in a finite-dimensional compact-invariant space, then continues each smooth dual coordinate; an unbounded evaluation functional on a Hilbert induction space is never invoked.

At \(z=z_*\), the target of (VS5) is normalized induction

\[
 (\Pi\otimes\alpha\beta^{-1})\times(\alpha\circ\det_t).
 \tag{VS6}
\]

Both the balanced source relation and the full parabolic covariance are part of the accepted quotient construction. For clarity, for \(p=\left(\begin{smallmatrix}A&B\\0&D\end{smallmatrix}\right)\) the latter is

\[
 Q(f,v)(p)=
 (\alpha\beta^{-1})(\det A)\alpha(\det D)
 |\det A|^{t/2}|\det D|^{-n/2}\Pi(A)I_{\beta,z_*}(f,v).
 \tag{VS7}
\]

Let \(\pi_0\) be any irreducible unitary representation of \(\mathrm{GL}_r(F_v)\). Define unitary irreducible candidates

\[
 \Pi_j=
 \left(\pi_0\times
       \mathop{\times}_{0\le h<j}
             ((\mu_h\xi)^{-1}\circ\det_t)\right)
                 \otimes(\eta^{\epsilon_j}\circ\det_{n_j}).
 \tag{VS8}
\]

Tadić's unitary irreducibility theorem applies to every induction in (VS8), including repeated characters. The formula (VS6) sends \(\Pi_j\) to \(\Pi_{j+1}\), since \(\alpha_j\beta_j^{-1}=\eta^{1-2\epsilon_j}\). The standard factor is exactly

\[
 L(z_*,\Pi_j\beta_j^{-1})
  =L(z_*,\pi_0\mu_j\xi)
       \prod_{h<j}\prod_{i=1}^t L(i,\mu_j\mu_h^{-1}).
 \tag{VS9}
\]

The first factor is finite and nonzero by the accepted unitary Godement--Jacquet bound \(\Re z>(r-1)/2\); here the margin is at least \(7/2\). Every other local factor is finite and nonzero because its character is unitary and \(i\ge1\). Thus every whole-Schwartz map (VS5) is holomorphic at the actual point. This does not require the later \(\Pi_j\) to be tempered or generic.

Once a fixed vector is shown nonzero below, Mínguez's type-II correspondence identifies this nonzero irreducible target with the actual small theta quotient. Thus (VS8) becomes an inductively proved description of the local images, not an assumption that an abstract parameter already proves occurrence.

\subsection{3. Every prescribed G stage is a compact projector}\label{every-prescribed-g-stage-is-a-compact-projector}

Let \(J_j\subset K_{n_j}\) be the inverse image of the block-upper triangular group with blocks \((r,t,\ldots,t)\), whose first diagonal block is \(I_r\). There are \(j\) blocks of size \(t\). Define

\[
 \lambda_j(g)=\prod_{h=1}^j\chi_h(\det\overline D_h),\qquad
 u_j=rjt+\frac{j(j-1)}2t^2.
 \tag{VS10}
\]

Thus \(J_0=K_r^1\) and \(\lambda_0=1\). These are precisely the G part of datYt with its specified leading LevelP equal to triv; the nontrivial finite good-series projector is a separate operator on that leading principal-level module.

The conjugate raw G function \(f_{G,j}\) is the indicator of integral \(b_j\)-by-\(n_j\) matrices with bottom \(t\) rows zero modulo \(\varpi\), top square block in \(J_j\) modulo \(\varpi\), and value

\[
 f_{G,j}(E_ng)=
 \mathbf1_{J_j}(g)\lambda_j(g)^{-1}
                     \chi_{j+1}(\det\overline g).
 \tag{VS11}
\]

This is literally the conjugate of DatumVec.omegaG: its old determinant-block character is inverted and its current \(\chi_{j+1}^{-1}\) is conjugated. No generic or Steinberg replacement is made.

For every \(v\in\Pi_j^{J_j,\lambda_j}\), all orbit variables in (VS5) lie in \(J_j\). On this compact group \(\beta_j=\chi_{j+1}\). Hence all characters cancel and

\[
 \boxed{\quad
 I_{\beta_j,z_*}(f_{G,j},v)=v_jv,\qquad
 v_j=\operatorname{vol}(J_j)
     =\frac{q^{u_j}|\mathrm{GL}_t(\mathbf F_q)|^j}
            {|\mathrm{GL}_{n_j}(\mathbf F_q)|}.
 \quad}
 \tag{VS12}
\]

This is a compact, absolutely convergent integral; the continuation was needed only to identify the map on the entire Schwartz space.

The target vector has exactly the next prescribed type. For \(k\in J_{j+1}\), write its reduction in blocks \((n_j,t)\). Its first diagonal block \(A\) belongs to \(J_j\). Applying (VS3) to the raw function, the factor from replacing the top block by \(A^{-1}\) cancels \(\beta_j(\det A)\); the remaining factor is \(\lambda_j(A)\chi_{j+1}(\det D)=\lambda_{j+1}(k)\). Thus \(Q(f_{G,j},v)\) belongs to the stated next isotypic space, for the entire target action, not only its value at one.

There is also an exact norm. The compact induced section is supported on the identity \((n_j,t)\)-parabolic cell. Indeed, integrality of \(k^{-1}E_ng\) with \(k\in K_{b_j}\) first makes \(g\) integral; full residue rank forces \(g\in K_{n_j}\). The bottom-zero condition then forces \(k\) to preserve the first \(n_j\)-plane modulo \(\varpi\). Conversely such \(k\) reduces by its upper-left block to (VS12). All remaining factors in (VS7) are unitary on the compact group. Therefore

\[
 \|Q(f_{G,j},v)\|^2
      =v_j^2{b_j\brack n_j}_q^{-1}\|v\|^2,\qquad
 \|f_{G,j}\|_{\mathcal S}^2=q^{-b_jn_j+u_j}
                         |\mathrm{GL}_t(\mathbf F_q)|^j.
 \tag{VS13}
\]

The first is the unitary induction norm; the second is the raw additive-Schwartz norm. They are not interchangeable. Iterating gives an injection on every initial \(K_r^1\)-fixed vector, including every prescribed good-series or generic cut:

\[
 \|F_{m,G}\|^2=
       \prod_{j=0}^{m-1}
          \left(v_j^2{n_{j+1}\brack n_j}_q^{-1}\right)
             \|v\|^2.
 \tag{VS14}
\]

The successive identity evaluations give the left-inverse scalar \(\prod_jv_j\). Each \(v_j\) is the reciprocal of an actual finite group index. At a coefficient prime above11, with \(q\ne11\),

\[
 v_{11}(v_j)
       =j\,v_{11}|\mathrm{GL}_t(\mathbf F_q)|
                      -v_{11}|\mathrm{GL}_{n_j}(\mathbf F_q)|
       \le0.
 \tag{VS15}
\]

Thus the exact local averaging cost is retained; no integral saturation assertion follows from this complex injection.

\subsection{4. The Q vector, flattened compact model and induction statement}\label{the-q-vector-flattened-compact-model-and-induction-statement}

At Q let \(\pi_0\) be spherical, and fix \(v\in\pi_0^{K_r}\). At stage \(j\), the unchanged conjugate raw function is

\[
 f_{Q,j}(X)=
 \mathbf1_{M_{b_j,n_j}(O)}(X)\,
 \mathbf1_{\mathrm{LowZero}_r(X)}
 \mathbf1_{\det\overline A_r(X)\ne0}
                   \xi(\det A_r(X))^{-1}.
 \tag{VS16}
\]

Only the first \(r\) columns are constrained below row \(r\); the other \(jt\) columns are arbitrary integral columns. This unchanged initial Levi rank is essential.

Realize \(\Pi_j\) by unitary induction with blocks \((r,t,\ldots,t)\), and flatten its successive compact arguments to a function on \(K_{n_j}\) with values in the space of \(\pi_0\). Let \(F_j\) be the actual result after \(j\) of the vectors (VS16), with \(F_0=v\). Write \(J(r,l)\subset K_{r+l}\) for the full reduction parabolic with blocks \((r,l)\), without an identity condition on its \(r\)-block.

We prove the exact induction statement

\[
 F_j^{\mathrm{flat}}(k)=
 \begin{cases}
 C_j\,\xi(\det\overline D_{jt})^{-1}v,&k\in J(r,jt),\\
 0,&k\notin J(r,jt),
 \end{cases}
 \qquad
 \|F_j\|^2=|C_j|^2{n_j\brack r}_q^{-1}\|v\|^2,
 \tag{VS17}
\]

where \(C_0=1\). The superscript flat matters: an intermediate vector-valued section need not itself vanish outside this cell before the inner compact arguments are evaluated.

Regrouping the last \(j\) blocks, the tail is

\[
 \rho_j=(\xi^{-1}\eta^{\epsilon_j})\circ\det_{jt}
                        \otimes T_j,\qquad
 T_j=\mathop{\times}_{h<j}(\mu_h^{-1}\circ\det_t).
 \tag{VS18}
\]

The irreducible unramified unitary \(T_j\) has a unique spherical line with its unit vector \(e_j\). On a compact tail argument its standard induced function is constant one. Thus (VS17) equivalently gives, in this regrouped model, the value \(C_j\,\xi(\det\overline D)^{-1}v\otimes e_j\). This equivalence uses induction in stages with probability Haar; it introduces no compact dimension factor.

\subsection{5. One exact Schur-complement step}\label{one-exact-schur-complement-step}

Assume (VS17) at stage \(j\), and put \(l=jt\), \(n=r+l\). In the support of \(f_{Q,j}(E_ng)\), write

\[
 g=\begin{pmatrix}U&V\\\varpi W&D\end{pmatrix},
 \quad U\in\mathrm{GL}_r(O),\quad V,W,D\text{ integral},
 \quad S=D-\varpi WU^{-1}V.
 \tag{VS19}
\]

The previous compact section is nonzero precisely when \(C=S^{-1}WU^{-1}\) is integral. One way to see this without an Iwasawa ambiguity is to ask for \(g^{-1}E_rh\) to be integral, with invertible top reduction and lower reduction zero. Since \(g\) is integral with \(U\) invertible modulo \(\varpi\), that condition forces \(h\in K_r\), and block inversion then gives

\[
 (g^{-1}E_rh)_{\rm low}=-\varpi Ch,\qquad
 (g^{-1}E_rh)_{\rm top}
       =U^{-1}h+\varpi U^{-1}VCh.
 \tag{VS20}
\]

Conversely integral \(C\) supplies those exact conditions. The corresponding compact factor of \(g\) is in \(J(r,l)\). Normalized parabolic covariance consequently contributes \[
 |\det S|^{-r/2}\,
   \xi(\det S)^{-1}\eta(\det S)^{\epsilon_j}T_j(S)e_j,
\] times \(C_jv\). The \(U\)-action fixes \(v\).

In (VS5), the factor \(\beta_j^{-1}\) cancels both the raw \(\xi(\det U)^{-1}\) and the \(\xi,\eta\) in this tail value. The sole surviving character on \(S\) is \(\mu_j(\det S)\). The substitution

\[
        W=SCU,\qquad D=S+\varpi SCV
 \tag{VS21}
\]

has additive Jacobian \(|\det S|^r\). The lower block contributes \(q^{-rl}\); group Haar contributes \(\kappa_n^{-1}|\det S|^{-n}\); integration over \(U\) contributes \(\kappa_r\). The determinant exponent left after the oscillator, parabolic and Jacobian factors is

\[
 \frac{n+t}{2}-\frac r2-n+r=\frac{t-l}{2}.
 \tag{VS22}
\]

The tail compact projection must be retained before evaluating the last integral. In the initial convergence half-plane, changing the orbit variable by \(\operatorname{diag}(A,D)\in K_r\times K_l\) shows that the vector \(I_{\beta_j,z}(f_{Q,j},F_j)\) transforms by \(\xi(\det D)^{-1}\). Evaluating the parabolic section at one therefore places it in \(\pi_0^{K_r}\otimes T_j^{K_l}\). Thus one may apply the probability projector \(P_l\) onto \(T_j^{K_l}\) before the \(S\)-integration. This insertion is an exact covariance identity, not an assumption that point evaluation is bounded on an infinite Hilbert space.

Consequently the nested identity evaluation is

\[
 C_j q^{-rl}\frac{\kappa_r}{\kappa_n}
 \int_{M_l(O)\cap\mathrm{GL}_l}
       |\det S|^{(t-l)/2}\mu_j(\det S)
              P_lT_j(S)e_j\,dS_{\rm add}\ \otimes v.
 \tag{VS23}
\]

Here \(P_lT_j(S)e_j=m_{T_j}(S)e_j\), with the normalized spherical matrix coefficient. The unprojected vector integral need not be absolutely convergent when \(j>1\); it is not being used at the actual point. The projected integral is the spherical Godement--Jacquet integral at \(z_*=(t+1)/2\), and with (VS4) its exact value is

\[
 \kappa_l L(z_*,T_j\mu_j)e_j
   =\kappa_l
       \prod_{h<j}\prod_{i=1}^t
             (1-(\mu_j\mu_h^{-1})(\varpi)q^{-i})^{-1}e_j.
 \tag{VS24}
\]

Its projected scalar convergence is genuine at this point. The absolute value of the spherical coefficient of \(T_j\) is bounded by that of the unitary induction of \(j\) trivial \(\mathrm{GL}_t\) characters: in the compact induced model this follows by removing the unit phases under the positive measure. The latter spherical function is nonnegative. Its Gaussian zeta series has positive coefficients and Euler product \(\prod_{i=1}^t(1-q^{-i})^{-j}\), finite at this point. Equivalently the standard Satake roots have exponents at most \((t-1)/2\) and the evaluation point exceeds that by one. The accepted spherical determinant-integral computation of LIFT16 (S6--14) then proves absolute convergence, not just formal continuation. This same majorant is locally uniform on a neighborhood of the actual point. The head vector is spherical and the head orbit was forced to be compact, so no unbounded initial coefficient remains in (VS23).

As in DO041, evaluation is first performed after a compact projector onto a finite-dimensional space of smooth invariants. There it is a legitimate linear functional on the holomorphic whole-Schwartz vector. On the initial convergence half-plane its value is (VS23); the absolute majorant just given and the identity theorem give the same value at \(z_*\). Thus there is no extension by zero across a rank-deficient boundary, or interchange of unjustified meromorphic integrals.

It follows that

\[
 \boxed{\quad
 C_{j+1}=C_j\Lambda_j,\qquad
 \Lambda_j={n_j\brack r}_q^{-1}
       \prod_{h<j}\prod_{i=1}^t
             (1-(\mu_j\mu_h^{-1})(\varpi)q^{-i})^{-1}.
 \quad}
 \tag{VS25}
\]

At \(j=0\), the tail dimension is zero and \(\Lambda_0=1\). At \(j=1\), this is exactly DO041's checked scalar, with the same ratio \(\mu_1/\mu_0\), not its unmotivated inverse.

\subsection{6. Full compact support, exact norms and nonzero propagation at Q}\label{full-compact-support-exact-norms-and-nonzero-propagation-at-q}

For the induction step's compact support, keep the earlier quotients at their actual values and deform only the new integral. Let \(k\in K_{n+t}\). A nonzero raw factor \(f_{Q,j}(k^{-1}E_ng)\) makes \(g\) integral and makes its first \(r\) columns have residue rank \(r\). A nonzero value of the regrouped prior section requires \(g\in P_{(r,l)}(F_v)J(r,l)\), by (VS17).

Write \(g=pu\), with \(u\in J(r,l)\), and let \(A\) be the upper-left block of \(p\). Then \(g^{-1}E_rA=u^{-1}E_r\) is integral, has invertible top reduction, and has zero lower reduction. Put \(Y=u^{-1}E_r\). Since \(gY=E_rA\) and \(g\) is integral, \(A\) is integral. Modulo \(\varpi\), \(gY=gE_r\overline Y_{\rm top}\) has rank \(r\), so \(A\in K_r\). It follows that the first \(r\) columns of \(g\) span the standard first \(r\)-space modulo \(\varpi\). But \(E_ng=k(k^{-1}E_ng)\) and the raw factor has that same standard residue span. Hence \(k\) preserves this space.

This proves that the flattened section is zero outside \(J(r,(j+1)t)\), first in the new integral's convergence half-plane. The finite-dimensional continuation of (VS5) preserves that zero at \(z_*\). No simultaneous deformation of earlier induced representations or iterated meromorphic Fubini argument is needed.

On this cell, (VS3) and the raw determinant character give the factor \(\xi(\det D_{(j+1)t})^{-1}\). The identity value is (VS25). This proves the whole compact formula (VS17) for \(j+1\) on the new integral's initial convergence domain and hence at the actual holomorphic point. This support argument concerns the complete nested section, not one component of a freely chosen target.

Finally the compact normalized-induction norm is the integral over all successive compact arguments. Parabolic covariance transports each inner compact argument to left multiplication of the outer compact variable; its character has modulus one. Repeated left Haar invariance removes the inner integrals. Thus the norm is exactly the \(K_{n_j}\)-norm of the flattened section, giving (VS17). In particular

\[
 \boxed{\quad
 \|F_{m,Q}\|^2=
   \left|\prod_{j=0}^{m-1}\Lambda_j\right|^2
                  {n_m\brack r}_q^{-1}\|v\|^2>0
                  \quad(v\ne0).
 \quad}
 \tag{VS26}
\]

Every factor in (VS25) is nonzero: its Euler denominator has modulus distance from one controlled by \(q^{-i}<1\). The raw Schwartz norm, counted independently from its residue support, is

\[
             \|f_{Q,j}\|_{\mathcal S}^2
                    =\kappa_r q^{-r(j+1)t}.
 \tag{VS27}
\]

At a coefficient prime above11 with \(q\ne11\), if every actual \((\mu_j\mu_h^{-1})(\varpi)\) is an11-unit, the exact valuation is

\[
 v_{11}(\Lambda_j)=
   -v_{11}{n_j\brack r}_q
   -\sum_{h<j}\sum_{i=1}^t
        v_{11}(1-(\mu_j\mu_h^{-1})(\varpi)q^{-i})\le0.
 \tag{VS28}
\]

This unit premise is supplied for the carried angular-character families by K039/K040; it is not inferred from a bare infinity type condition. In all cases (VS25)--(VS27) retain the exact complex normalization without asserting integral saturation.

\subsection{7. Actual local Rallis pairings and their normalization}\label{actual-local-rallis-pairings-and-their-normalization}

The nonzero quotient vectors above are also detected by the actual fixed-input local Rallis form. This is stronger than an unspecified nonzero test in the same local representation.

First check absolute Gram convergence. In the compact induced realization of (VS8), a smooth principal-level vector \(w\) has a finite bound \(M_w=\sup_{k\in K_{n_j}}\|w(k)\|\). Let \(\Pi_j^+\) be normalized induction of the trivial \(\mathrm{GL}_r\) representation and \(j\) trivial \(\mathrm{GL}_t\) representations. Its unit spherical coefficient \(m_j^+\) is nonnegative: its compact action preserves the cone of nonnegative functions. Removing the unit character phases and using unitarity of \(\pi_0\) in the compact coefficient integral gives

\[
 |\langle w,\Pi_j(g)w\rangle|\le M_w^2m_j^+(g).
 \tag{VS29}
\]

The standard roots of this spherical majorant have exponents at most \((t-1)/2\) for \(j\ge1\), and at most \((r-1)/2\) for \(j=0\). This is independent of the number of old blocks. Both raw functions have modulus at most the integral matrix indicator. Therefore their oscillator coefficients are bounded by the nonnegative vacuum coefficient \(A_{b_j}(g)\) of LIFT16. The exact Cartan majorant in LIFT16 (S12--16), including the polynomial correction for repeated Satake roots, gives

\[
 \int A_{b_j}(g)m_j^+(g)\,dg<\infty,\qquad
 q^{-b_j/2}q^{(n_j-1)/2}q^{(t-1)/2}=q^{-1}
                 \quad(j\ge1).
 \tag{VS30}
\]

For \(j=0\) the corresponding exponent is \(-(t-r+2)/2<0\). Thus the actual Gram integral is absolutely convergent in every stage, with no temperedness assumption on the local \(\Pi_j\).

The full-rank additive-matrix disintegration has constant \[
 C_{n_j,b_j}=\prod_{i=0}^{n_j-1}(1-q^{i-b_j})
             =\kappa_{b_j}/\kappa_t.
\] It gives, with the precise split Hom convention (VS3),

\[
 Z_j(f,w)=
 \int_{\mathrm{GL}_{n_j}}
       \langle f,\omega_j(g,1)f\rangle
       \langle w,\Pi_j(g)w\rangle\,dg
       =\frac{\kappa_{b_j}}{\kappa_t}
                   \|Q_{\alpha_j,\beta_j}(f,w)\|^2.
 \tag{VS31}
\]

Here is why the same identity applies to the whole-Schwartz continuation, rather than only a formal orbit integral. Use the finite-dimensional Abel argument already proved in DO047 (QN17--19), with (VS29)--(VS30) as its verified new majorant. Truncate the orbit variable to an increasing compact union of Cartan double cosets and insert \(|\det h|^\zeta\), \(\Re\zeta\ge0\). Raw support forces \(h\) integral, so this added factor has modulus at most one. Before interchanging the two orbit integrals, their absolute Gram is bounded by \(C_{n_j,b_j}^{-1}M_w^2\int A_{b_j}m_j^+\), finite by (VS30). The truncated compact functions are therefore Cauchy in compact \(L^2\), including at \(\zeta=0\).

All of them are right \(K_{b_j}^1\)-invariant and take values in \(\Pi_j^{K_{n_j}^1}\). For the latter assertion, left replacement of the orbit variable by \(K_{n_j}^1\) preserves the truncation, the integral support, its reduction, and the tame source character. Admissibility and the finite compact quotient put the entire argument in one finite-dimensional compact-function space. For large \(\Re\zeta\) this is the ordinary convergent integral; holomorphy (VS9), the identity theorem, and continuity at zero identify its limit with (VS5). Expanding the Gram and applying the additive disintegration now gives (VS31), with exactly the stated constant. This argument makes no false assertion that the unprojected single vector integral is absolutely convergent.

Combining (VS12)--(VS13) with (VS31) simplifies the actual unit-Schwartz/unit-source G value to

\[
 \boxed{\quad
 \frac{Z_j(f_{G,j},w)}
      {\|f_{G,j}\|^2\|w\|^2}=v_j
 \quad(w\in\Pi_j^{J_j,\lambda_j},\ w\ne0).
 \quad}
 \tag{VS32}
\]

This is the raw local Rallis value with unit input norms. It has not been divided by a standard local \(L\)-factor at G; that factor can depend on the initial constituent and must not be removed a second time in a global partial-\(L\) formula.

At Q write \(\mathcal L_j=\prod_{h<j}\prod_{i=1}^t
 (1-(\mu_j\mu_h^{-1})(\varpi)q^{-i})^{-1}\). Equations (VS17), (VS25), (VS27), and (VS31) give

\[
 \boxed{\quad
 \frac{Z_j(f_{Q,j},F_j)}
      {\|f_{Q,j}\|^2\|F_j\|^2}
   =c_{Q,j}|\mathcal L_j|^2,\qquad
 c_{Q,j}=q^{-rjt}
       \frac{\kappa_r\kappa_{jt}\kappa_{(j+1)t}}
            {\kappa_t\kappa_{r+jt}}>0.
 \quad}
 \tag{VS33}
\]

The literal tuple Q set forces \(\xi\) tamely ramified, as checked in DO047 (QN23). Since \(\pi_0,\mu_j\) are unramified there, the initial factor \(L(z_*,\pi_0\mu_j\xi)\) is one. Thus the actual split standard factor is \(|\mathcal L_j|^2\), and \(c_{Q,j}\) is the standard-normalized local value. At \(j=0\), it is one. For \(j\ge1\),

\[
       \frac1{16}q^{-rjt}<c_{Q,j}<q^{-rjt}.
 \tag{VS34}
\]

Indeed \(\kappa_a>1/4\), and \(\kappa_r\kappa_{jt}/\kappa_{r+jt}\) lies strictly between \(\kappa_r\) and one, while \(\kappa_{(j+1)t}/\kappa_t\) lies between \(1/4\) and one. These inequalities follow directly from the products or the Gaussian-binomial formula. On the specified one-dimensional preceding-image line, the inverse factor for the standard-normalized positive norm lies between \(q^{rjt/2}\) and \(4q^{rjt/2}\). For the whole Q chain the product of these inverse norm factors is therefore at most \(4^{m-1}q^{rtm(m-1)/4}\), with the exact factors (VS33) preferred whenever used in a global cost.

These norms refer to the specified local Haar measures. Changing the source group Haar multiplies the local Gram by that Haar factor. None of the formulas absorbs global class masses, a real intertwiner scalar, or an arithmetic coefficient-lattice index.

\subsection{8. What has and has not been propagated}\label{what-has-and-has-not-been-propagated}

At G the unchanged actual local map is injective on its entire prescribed finite type at each stage. At Q it is nonzero, with the exact scalar and norm above, on the \textbf{actual preceding image of the initial spherical line}. The latter is stronger than local occurrence with an unspecified Schwartz test, and weaker than injectivity on every Q-level vector in \(\Pi_j\).

These results apply separately at all selected G and Q places and to every actual initial irreducible unitary local constituent that carries the stated initial type. The finite good-projector restriction at G does not require a new type, because the proof already handles the whole principal-level module. They do not prove global nonzero propagation: the other finite basic vectors, the real chosen contractions, the actual global Rallis formula including the first positive-Witt-index edge, and each automorphic multiplicity embedding must still be treated. LIFT56 supplies cuspidality for all these unchanged inputs, including positive intermediate Witt rank.

The fixed G/Q quotient integrals and their norms are not the primitive arithmetic theta coefficient map. Global comparison scalars, CM periods, class masses, source lattice denominators and any terminal Bessel/row restriction remain explicit further data. No DP44 cost comparison is concluded.

\subsection{9. Exact source and published inputs}\label{exact-source-and-published-inputs}

\begin{itemize}
\tightlist
\item
  \nolinkurl{Formalism/LocalToGlobal/IndefiniteReal/DatumVec.lean,} lines49--70: omegaG/omegaQ; lines83--98: IsStepVecOf. The LowZero parameter in omegaG is the current source rank; in omegaQ it is the fixed initial rank LS.r.
\item
  \nolinkurl{Formalism/LocalToGlobal/IndefiniteReal/LinksChar.lean,} lines255--269: chrStep, datYt with triv, and the unchanged chain\_vY/chain\_vX functions.
\item
  \nolinkurl{Formalism/LocalToGlobal/Divisibility/DefGroup.lean,} lines692--713: actual level and block characters.
\item
  LIFT31, VA14--24; DO041, DO41.5--27; DO047, QN13--23; LIFT16, S6--16: accepted local quotient, two-stage calculation and spherical determinant integral. Their normalizations are retained.
\item
  Jiang, D., \& Luo, Z. (2023). Certain Fourier operators and their associated Poisson summation formulae on GL1. Pacific Journal of Mathematics, 326(2),301--372. Theorem2.4(1)--(2), pp.313--314. \url{https://msp.org/pjm/2023/326-2/pjm-v326-n2-p05-p.pdf}
\item
  Tadić, M. (1986). Classification of unitary representations in irreducible representations of general linear group (non-Archimedean case). Annales scientifiques de l'École Normale Supérieure, 19(3),335--382. TheoremD, pp.338--339. \url{https://www.numdam.org/article/ASENS_1986_4_19_3_335_0.pdf}
\item
  Mínguez, A. (2008). Correspondance de Howe explicite: paires duales de type II. Annales scientifiques de l'École Normale Supérieure, 41(5),717--741. Theorem1 and normalization(1.4). \url{https://www.numdam.org/item/ASENS_2008_4_41_5_717_0.pdf}
\end{itemize}

The new proof is the induction on the actual fixed finite functions, with its full compact support, prior-image assertion and exact scalars. The established quotient/classification theorems are cited at their accepted scope, not re-proved.

\subsection{Author verification record}\label{author-verification-record-18}

The author read every section and VS1--34, the complete diagnostic, the literal DatumVec/LinksChar/DefGroup clauses, and the accepted DO047 finite-dimensional Abel proof. The old-block and current characters, fixed Q head rank, compact projection before the tail integral, full flattened support, Haar constants, positive Gram majorant and one-dimensional inverse-norm scope were checked. The final frMapO/unitarity and scope clarifications were reread. No global tower, arithmetic saturation or Lean build is claimed. Exact diagnostic: \nolinkurl{general_vertical_split_certificate.py/.json,} 7,775 assertions passed 2026-11-22T11:34:03.892281+00:00.

\clearpage

\section{AM-DO-066 --- simultaneous G congruences sharpen the actual cross-column count}\label{am-do-066-simultaneous-g-congruences-sharpen-the-actual-cross-column-count}

Author: GPT-6 Astra Pro, ALP\_MACHINE \nolinkurl{/root/definite_oscillator}. 22 November 2026 UTC. Inherited reasoning effort is not separately exposed.

\textbf{Author whole-proof check PASS --- 2026-11-22T11:46:05.671308+00:00.}

The actual rank-one reductions at G impose a common prime-ideal divisibility condition on every coordinate of the positive cross column. Counting inside that ideal, while retaining its exact finite-character weight, improves the pointwise field-norm estimate. The resulting bound removes every globally unipotent scalar-cross term with at most five unequal-pencil T places, at every total support size. It also removes specified larger unequal-pencil supports.

The removed union has absolute contribution less than 2\^{}-4085 in the unchanged K062 normalization. The proof keeps all 4\^{}334 possible local rank/orientation patterns. It makes no assertion about the fixed-W0 terms or the sign of the remaining sum.

\subsection{1. Actual coordinates, ranks and pencil types}\label{actual-coordinates-ranks-and-pencil-types}

Retain the J56 upper-gap family, its actual near-origin point and the entire finite tensor of K062, with

\[
 a=f=168,\quad n=169,\quad m=164,\quad H=8192,\quad M=H+1=8193,
 \qquad E^n=S\perp W_0.
 \tag{GC1}
\]

Here S is the original positive coordinate space of dimension m and W0 has signature (4,1). The last coordinate e\_n is the negative coordinate. Write h for the normalized unipotent representative of a central class, and N=h−I. The subset under consideration is exactly

\[
 \operatorname{rank}_E N=2,\qquad h_{SS}=I,
 \qquad x=(h e_n)_S\ne0.
 \tag{GC2}
\]

This is the scalar-cross case of DO061(TR32--35), DO063(FV3--4) and DO065. The signature-one normal form gives N³=0 and zero S rows and columns in N². In particular, the cross column x in GC2 is nonzero at every field embedding and satisfies

\[
                 \sum_{i\in S}|\tau(x_i)|^2\le M
                     \qquad(\tau\in\Phi).
 \tag{GC3}
\]

Only nonzero coordinates are meant in the preceding field-embedding statement: a coordinate which is nonzero in E cannot vanish under an embedding. We will count their exact number separately.

Use the three actual local types of K070/K071. Type I preserves the full pencil; type II has unequal full pencils and a unique fixed incidence; type III has unequal full pencils and a unique off-diagonal incidence. Let s be the total nonintegral T support, u the number of type II places, t the number of type III places, and v=u+t. The v unequal-pencil places form a set U. The carried local cross-integrality statement K071(NB5) gives

\[
 x\in I_U^m,\qquad
 I_U=\prod_{\mathfrak P\mid U}\mathfrak P^{-1},
 \qquad \operatorname N(I_U)=T^{-2v}.
 \tag{GC4}
\]

Both CM partners occur in this product. Type I places contribute no cross-column denominator. All these T primes have absolute degree one. On the accepted K068 residual, the hypothesis t≤8 implies the global unipotence used in GC2, by K070. The theorem below keeps this hypothesis; it does not infer unipotence from a bound on v alone when v\textgreater8.

\subsection{2. One common orientation at each rank-one G reduction}\label{one-common-orientation-at-each-rank-one-g-reduction}

There are 2(f−1)=334 unmixed G places. At each, the actual reduction of h is unipotent and has rank k=0,1 or2 after subtracting the identity. The two mixed projectors retain their norm bound one. All entries of h are integral at G in the original coordinate frame.

Fix one of these split places, with CM partners P and bar(P), and write the selected-component reduction as I+N\_P. The paired component is

\[
               I+N_{\overline P}=J(I+N_P)^{-t}J,
 \tag{GC5}
\]

where J is the original diagonal signature matrix. Conjugating the finite projector changes neither the rank nor the bound used below; the divisibility conclusion itself is in the original S coordinates.

If k=1, write N\_P=b c\^{}t, with b,c nonzero. Nilpotence gives c\^{}t b=0, so N\_P²=0. Since (N\_P)\_SS=0,

\[
                  b_Sc_S^t=0
       \quad\Longrightarrow\quad b_S=0\ \text{or}\ c_S=0.
 \tag{GC6}
\]

This conclusion uses an outer product over a field. If one entry of b\_S is nonzero, every entry of c\_S must vanish. Thus it gives a single orientation for the whole column, not a separate choice per coordinate. When b\_S=0, every coordinate of x vanishes modulo P. When c\_S=0, GC5 and (I+N\_P)\^{}−1=I−N\_P show that every coordinate of x vanishes modulo bar(P). When k=0, x vanishes at both partners.

There are consequently at most four local patterns: rank2, rank1 with P divisibility, rank1 with bar(P) divisibility, and rank0 with both. If a rank-one matrix permits both orientations, select either by a fixed rule, or overcount it. Across all unmixed factors the number of patterns is at most

\[
                             4^{334}=2^{668}.
 \tag{GC7}
\]

For a fixed pattern let P1 be the product of the residue sizes at its rank-one factors, and P0 the product at its rank-zero factors. Let J\_G be the product of the chosen prime above each rank-one factor and both primes above each rank-zero factor. Then

\[
 \operatorname N(J_G)=P_1P_0^2,
                 \qquad x\in(J_GI_U)^m.
 \tag{GC8}
\]

J\_G is an actual integral ideal, coprime to T. It need not be invariant under CM conjugation. GC8 is simultaneous divisibility of every S coordinate and does not require J\_G to be principal.

The exact unipotent Green formula, with the one global Euler correction less than2, now reads

\[
 |a(h)|\le
 2Q^{-(2a-1)(f-1)}P_1^{a-1}P_0^{2a-1},
                  \qquad Q=g_0g_1.
 \tag{GC9}
\]

Indeed the rank2 exponent is 2a−1=335, the rank1 exponent is a=168, and rank0 has character ratio one. These are precisely the carried Green factors of K067(PS3--9) and DO061(TR23--24), with no replacement of the two mixed partial traces. The factor two is not repeated per place.

\subsection{3. The same ideal improves the actual selected real factor}\label{the-same-ideal-improves-the-actual-selected-real-factor}

Choose any coordinate x\_i which is nonzero in E. Fractional-ideal norm divisibility in GC8 gives

\[
 |\operatorname N_{E/\mathbf Q}(x_i)|
       \ge\operatorname N(J_G)T^{-2v}.
 \tag{GC10}
\]

This uses containment of the principal fractional ideal in J\_GI\_U; their quotient has integral ideal norm. No integral basis or principal generator is chosen.

Retain the actual conjugator k\_p at the selected near-origin point. DO063(FV22--29) proves its norm product is less than2, and applies the full selected-compression estimate of K060 to every embedding. Together with the actual energy floor kappa0\textgreater2\^{}485 this gives

\[
 |c(h)|<12^f2^{f/2-485f/4}
                   \operatorname N(J_G)^{-1/4}T^{v/2}.
 \tag{GC11}
\]

For clarity, the existing argument bounds each selected compression by 12/{[}sqrt(2sinh r) kappa0\^{}(1/4){]}. The conjugator and the unipotent operator inequality give \textbar tau(x\_i)\textbar≤2·2sinh r at that embedding. Multiplying the resulting inverse square roots and using GC10 proves GC11. The normalized full selected trace is bounded by its operator norm; there is no extra selected-dimension factor. This is the actual conjugated trace c(h), not an origin model coefficient.

Set

\[
                         \alpha=a-\tfrac34=\tfrac{669}{4}.
 \tag{GC12}
\]

The product of the rank-loss factor in GC9 and the ideal factor in GC11 satisfies

\[
 P_1^{a-1}P_0^{2a-1}\operatorname N(J_G)^{-1/4}
 =P_1^{a-5/4}P_0^{2a-3/2}
 \le\operatorname N(J_G)^\alpha.
 \tag{GC13}
\]

The dropped factor is P1\^{}−1/2. This inequality is uniform in the entire rank pattern and retains rank-zero factors with their full loss.

\subsection{4. Count the column inside its simultaneous congruence ideal}\label{count-the-column-inside-its-simultaneous-congruence-ideal}

Fix the pattern J\_G, the T type/pair data and a nonempty support of d nonzero S coordinates, where 1≤d≤m. Put

\[
           z=\operatorname N(J_G)^{1/(2f)}T^{-v/f},
                 \qquad R=\sqrt M.
 \tag{GC14}
\]

Every nonzero y in J\_GI\_U has trace Euclidean norm at least sqrt(f) z. This follows from GC10 and arithmetic--geometric mean applied to the f squared complex moduli. The product lattice (J\_GI\_U)\^{}d therefore has minimum distance at least sqrt(f) z in real dimension 2fd.

If there is an actual column with the prescribed d nonzero coordinates, each coordinate separately has squared trace norm at least f z². Summing GC3 over all embeddings consequently gives

\[
                    dfz^2\le fM,\qquad z\le R/\sqrt d.
 \tag{GC15}
\]

If that actual intersection is empty its contribution is zero, so the following upper estimate applies without assuming the existence of a lattice point merely from a determinant bound.

Disjoint balls of radius sqrt(f)z/2 around lattice points in the ball of radius sqrt(f)R give

\[
 \#\{\text{possible }S\text{ columns on this support}\}
                 \le(1+2R/z)^{2fd}.
 \tag{GC16}
\]

This counts the whole relevant fractional lattice. It does not divide an unconstrained count by an asserted residue density, and does not require a covering radius for a nonprincipal congruence ideal.

Since N(J\_G)\^{}alpha=T\^{}(2alpha v) z\^{}(2falpha), the weighted count is bounded by

\[
 \begin{split}
 \operatorname N(J_G)^\alpha\,\#S
 &\le T^{2\alpha v}z^{2f(\alpha-d)}(z+2R)^{2fd}\\
 &\le T^{2\alpha v}
          (M/d)^{\alpha f}(1+2\sqrt d)^{2fd}.
 \end{split}
 \tag{GC17}
\]

The last inequality uses GC15 and alpha\textgreater d for every d≤164. Both positive factors in the first line are increasing in z. Thus all the rank-loss ideal dependence has been bounded simultaneously with the actual number of allowable column coordinates.

There are binom(m,d) possible supports. Here alpha f=28098, and 1+2sqrt(d)\textless27. Define the rational constants

\[
 C_d=\binom{164}{d}27^{336d}(8193/d)^{28098},
             \qquad C_S=\sum_{d=1}^{164}C_d.
 \tag{GC18}
\]

Exact integer comparisons in the companion certificate prove

\[
                  C_d<2^{420560}\quad(1\le d\le164),
                     \qquad C_S<2^{420568}.
 \tag{GC19}
\]

This includes the coordinate-support choice. No separate rank chart or positive-block count is subsequently multiplied into GC19.

\subsection{5. The original W0 lattice, fibre and incidence weights}\label{the-original-w0-lattice-fibre-and-incidence-weights}

For each actual ordered T pair configuration, the other five coordinates of h e\_n lie in the exact paired target lattice B\_Sigma of DO063. Its half-trace covolume is D0\^{}5 T\^{}(−10s), with

\[
 D_0=7^{140}29^{144}5^{84}13^{84},\qquad
 B_W=\frac{[(22/7)f\,9543^2]^{5f}}{(5f)!D_0^5}.
 \tag{GC20}
\]

The carried covering-radius and actual negative-column norm calculation therefore bounds their number by B\_W T\^{}(10s). The exact integer certificate gives B\_W\textless2\^{}14861. This factor retains every type I denominator; no whole-matrix integrality is inferred from GC4.

For each full last column, the original finite isometry fibre is at most G\_n\textless2\^{}5022, as in DO063. Imposing a G congruence pattern only reduces that fibre. The S and W0 counts may be multiplied as an upper bound even though actual unitary columns need not vary independently.

The exact weighted pair enumeration from K071(NB7--8) is

\[
 \binom fs\binom sv\binom vu(T-1)^{s-v}.
 \tag{GC21}
\]

Type I has at most (T−1)\^{}2 ordered actual pairs and modulus at most 1/(T−1); type II has T−1 pairs and modulus 1/(T−1); type III has (T−1)(T−2) pairs and modulus 1/{[}(T−1)(T−2){]}. Thus type II and III cost one after weighting, while type I costs T−1. These estimates retain the original signs rather than assigning them a favorable value. For fixed s,v, sum over exactly those u with t=v−u≤8. Their binomial weights have total at most 2\^{}v.

Combining GC7, GC9, GC11, GC13, GC17--21 proves the following uniform bound for the whole actual (s,v) cross subset with t≤8:

\[
 \begin{split}
 U_{s,v}\le{}&\binom fs\binom sv2^v(T-1)^{s-v}
       G_nB_W\,[2\,4^{334}Q^{-335\cdot167}]\,C_S\\
 &\hspace{12mm}\cdot12^f2^{f/2-485f/4}T^{335v+10s}.
 \end{split}
 \tag{GC22}
\]

The exponent 335v in the second line is 2alpha v+v/2. The latter v/2 is the remaining T factor in GC11; the N(J\_G)\^{}−1/4 factor has already been used in GC13. Neither contribution is omitted or charged twice. All rank/orientation patterns, possible columns, isometry fibres and actual ordered pencil pairs have been included.

\subsection{6. Exact complete subset ranges}\label{exact-complete-subset-ranges}

The fixed host integers satisfy Q\^{}4\textgreater2\^{}281 and T\textless2\^{}996. Also 12\^{}168\textless2\^{}603. Combining these with GC19--20 and G\_n\textless2\^{}5022 gives the constant before the two binomial estimates

\[
 \begin{split}
 C={}&5022+14861+1+668+420568+603+84-20370\\
      &\hspace{20mm}-(281/4)(335)(167)
       =-14034797/4.
 \end{split}
 \tag{GC23}
\]

Use binom(f,s),binom(s,v)≤2\^{}f, retain the 2\^{}v type sum, and bound (T−1)\^{}(s−v)≤T\^{}(s−v). Thus GC22 has the conservative exact exponent

\[
       \log_2 U_{s,v}<
       E(s,v):=-14033453/4+10956s+332665v.
 \tag{GC24}
\]

The coefficient of v includes its binomial-sum exponent one, and the T exponent is 334v+11s. The following table lists every claimed range; each contains all integers s from max(5,v) to the displayed endpoint.

\begin{longtable}[]{@{}
  >{\raggedleft\arraybackslash}p{(\columnwidth - 4\tabcolsep) * \real{0.3636}}
  >{\raggedleft\arraybackslash}p{(\columnwidth - 4\tabcolsep) * \real{0.3636}}
  >{\raggedright\arraybackslash}p{(\columnwidth - 4\tabcolsep) * \real{0.2727}}@{}}
\toprule\noalign{}
\begin{minipage}[b]{\linewidth}\raggedleft
Unequal-pencil places v
\end{minipage} & \begin{minipage}[b]{\linewidth}\raggedleft
Largest s
\end{minipage} & \begin{minipage}[b]{\linewidth}\raggedright
Retained off-diagonal condition
\end{minipage} \\
\midrule\noalign{}
\endhead
\bottomrule\noalign{}
\endlastfoot
0 & 168 & t=0 \\
1 & 168 & t≤1 \\
2 & 168 & t≤2 \\
3 & 168 & t≤3 \\
4 & 168 & t≤4 \\
5 & 168 & t≤5 \\
6 & 137 & t≤6 \\
7 & 107 & t≤7 \\
8 & 76 & t≤8 \\
9 & 46 & t≤8 \\
10 & 16 & t≤8 \\
\end{longtable}

Every listed cell satisfies E(s,v)≤−4096. In particular the tight all-support endpoint E(168,5)=−4430.25 retains a strict margin. There are exactly

\[
       6\cdot164+132+101+69+38+7=1331
 \tag{GC25}
\]

distinct (s,v) cells. Each already contains all allowed u and all actual support/type labels, so these are not counted again. Their total absolute contribution in the original signed short sum is

\[
                       \boxed{1331\,2^{-4096}<2^{-4085}.}
 \tag{GC26}
\]

Taking their set differences with previously removed subsets preserves this upper bound and avoids double charging. In particular all remaining globally unipotent scalar-cross terms with v≤5 disappear from the unbounded remainder, regardless of s. All scalar-cross terms at s=5 are included. DO065's independent nonscalar-block reduction and the fixed-W0 estimates must still be applied with their own scopes.

\subsection{7. What remains and how this estimate is used}\label{what-remains-and-how-this-estimate-is-used}

The actual K062 identity retains the same denominator

\[
       w_x\,28C_\infty d_Rd_W\bigl((T-2)/(T-1)\bigr)^f.
 \tag{GC27}
\]

If the newly bounded subset is removed from its signed short sum, the absolute error increases by less than 2\^{}-4085. The existing error bound and the mass factors are otherwise unchanged. No individual summand is declared positive, and the nonintegral type III signs, CM phases, genuine G projector and selected real trace remain fixed.

The result does not bound the entire fixed-W0 sector, all larger v, or the t≥9 part not already known to be globally unipotent. It does not establish all-T Mram positivity, a terminal period, net arithmetic gain, literal source links or SizesAt. In particular the independent source size/link obstructions are unaffected. No J14 assertion is made.

\subsection{8. Verification and carried sources}\label{verification-and-carried-sources}

The new ingredients are the simultaneous rank-one orientation GC6--8, its combination with the actual real factor, and the weighted ideal packing GC14--19. They are proved above. The exact finite character, target-lattice, near-point and selected-compression inputs are the already whole-checked K067(PS3--9), K070, K071(NB1--16), DO061(TR19--35), DO063(FV5--29) and DO065. Their original normalizations are retained.

The primary finite-character formula is recorded in Green (1955, §5, Example (ii), pp.430--431) and in the directly checked statement of Gorodetsky and Hazan (2017, Theorem 2.1, equation (2.1), p.5), as applied in K067. The selected operator estimate is the checked use of Haagerup and Schlichtkrull (2014, Theorem 1.1, p.228) in K060. The five-coordinate lattice bound uses the exact Banaszczyk (1993, Theorem 2.2, p.632) normalization already verified in DO063/LIFT47. No new external representation theorem is invoked.

The companion certificate checks all164 support-size packing constants by integer inequalities, the actual W0 constant and field data, all 1331 support cells, the retained type sums and the union bound. Exhaustive small residue-field rank-one models test common orientation and the paired inverse-transpose formula; rank-two models test that the baseline pattern imposes no false divisibility. These are diagnostics of the proof, not enumeration of the rational automorphic support.

\subsubsection{Primary references}\label{primary-references-2}

Banaszczyk, W. (1993). New bounds in some transference theorems in the geometry of numbers. \emph{Mathematische Annalen, 296}, 625--635. \url{https://doi.org/10.1007/BF01445125}.

Gorodetsky, O., \& Hazan, Z. (2017). On certain degenerate Whittaker models for cuspidal representations of GL\_kn(F\_q). arXiv:1707.07308v2. \url{https://arxiv.org/pdf/1707.07308}. The operative locator is Theorem 2.1/(2.1), p.5.

Green, J. A. (1955). The characters of the finite general linear groups. \emph{Transactions of the American Mathematical Society, 80}(2), 402--447. \url{https://www.ams.org/tran/1955-080-02/S0002-9947-1955-0072878-2/S0002-9947-1955-0072878-2.pdf}.

Haagerup, U., \& Schlichtkrull, H. (2014). Inequalities for Jacobi polynomials. \emph{The Ramanujan Journal, 33}, 227--246. \url{https://doi.org/10.1007/s11139-013-9472-4}.

\subsection{Author verification record}\label{author-verification-record-19}

At 2026-11-22T11:46:05.671308+00:00, I completed the whole saved GC1--27 proof and the entire final companion certificate read. I reconstructed the paired inverse-transpose orientation, the simultaneous all-coordinate ideal containment, every rank-loss exponent, the actual conjugated selected trace, the nonempty-support cap and weighted monotonic packing, the unchanged W0 lattice and isometry fibre, all ordered-pair weights and all1331 cells. The primary finite-character title was corrected against the retained K067 source audit before this final read; no formula depended on the draft bibliographic wording.

The final certificate passed34,494 exact assertions at 2026-11-22T11:42:37.332798+00:00, including164 integer packing comparisons,7,478 exhaustive small rank-one outer-product models,12 rank-two models, and every claimed support cell. The first completed replay passed without a mathematical or numerical correction. These diagnostics do not enumerate the rational short sum or determine its sign. The final status, reference list and scope were reread before freezing.

\clearpage

\section{AM-K073 --- fixed five-dimensional blocks: rank-zero congruence packing}\label{am-k073-fixed-five-dimensional-blocks-rank-zero-congruence-packing}

Author: GPT-6 Astra Pro, ALP\_MACHINE \nolinkurl{/root/kummer_extension}. 22 November 2026 UTC. Inherited reasoning effort is not separately exposed.

\textbf{Author whole-proof check PASS --- 2026-11-22T11:49:52.329902+00:00.} The companion certificate passes 1,131 exact assertions. Independent whole reviews are recorded separately.

The complete five-place fixed-W0 unipotent sector of the unchanged J56 short sum has absolute contribution less than 2\^{}(-9199), uniformly in all carried role assignments and canonical twists. Both possible nilpotent ranks and every actual pencil incidence are included. If all labels have the allowed top role, the same argument bounds the complete six-place fixed-W0 unipotent sector by 2\^{}(-7423). The proof uses the actual congruence lattice of the difference column, not a density assumption about congruence classes.

These are bounds for sectors of K062's signed short sum. They do not assert positivity of the remaining sum, nonvanishing of the actual Mram period, or literal source Host/FamP admissibility. The all-top statement is a separate specialization of the existing mixed numerical profile; it is not a uniform six-place assertion.

\subsection{1. The unchanged family and the precise sector}\label{the-unchanged-family-and-the-precise-sector}

Retain K071(NB1--9), K062(AT13--17), and DO063(FV1--7):

\[
 a=f=168,\quad n=169,\quad m=164,\quad H=8192,
 \qquad E^n=S\mathbin\perp W_0,
 \quad\dim_EW_0=5,\quad\operatorname{sig}W_0=(4,1).
 \tag{FB1}
\]

The original real coordinate form on W0 is J0=diag(1,1,1,1,−1), with last basis vector e\_n.~The field, primes, finite projectors, scalar packet, actual point p and all class weights remain those of K032, DO040, LIFT41 and K062. Put

\[
 g_0=22285932761,\quad g_1=64072175441,
 \quad Q=g_0g_1,\qquad M=Q^{f-1}.
 \tag{FB2}
\]

The exact unchanged 996-bit prime T is read from \nolinkurl{definite_oscillator_large_t_certificate.json}, J56 row, and hashed in the companion certificate. There are 2(f−1)=334 unmixed relative G places. Their residue-size product is M; each such relative place has two split CM partners. The two mixed finite projectors are retained.

Consider actual supported short-sum representatives whose unique central normalization is

\[
                  h=I_S\oplus h_0,\qquad h_0\ne I,
                  \qquad h_0\text{ unipotent}.
 \tag{FB3}
\]

They have \(\max_\tau\|h_\tau\|\le H\), with respect to the original positive Euclidean norm at each complex label. Put N=h−I. Signature one gives N³=0 and rank N in \{1,2\}, as already established in DO061(TR19--24). The central normalization changes normalized finite and real traces only by unit phases. We bound their absolute values.

Let Σ be the nonintegral relative T support, \textbar Σ\textbar=s. The theorem below first assumes global unipotence explicitly. On the carried K068/K070 remainder it applies whenever the number t of type III places is at most8. Thus s=5 and s=6 require no extra unipotence hypothesis after those already checked cuts. Nothing here proves unipotence for all supports or all source parameters.

The scalar real factor is the original normalized selected trace

\[
 c(h)=\frac1{d_W}\operatorname{Tr}_W
   \bigl(P_W\pi(k_p^{-1}h_\infty k_p)|_W\bigr).
 \tag{FB4}
\]

All compact selected modules carry the isometric PBW metrics of K043. The factors a(h) and rho\_v(h) are precisely K062's normalized genuine G partial trace and local signed pencil traces.

\subsection{2. The paired rank-zero ideal gives a uniform finite bound}\label{the-paired-rank-zero-ideal-gives-a-uniform-finite-bound}

Fix the subset Z of unmixed relative G places at which N reduces to zero. Because h is unitary and integral at G, reduction at the conjugate CM partner is inverse-adjoint. Rank zero is therefore equivalent at both partners. Define the integral ideal

\[
 \mathfrak J_Z=\prod_{w\in Z}\mathfrak P_w\bar{\mathfrak P}_w,
 \qquad D=\mathrm N_{E/\mathbf Q}(\mathfrak J_Z)
       =\prod_{w\in Z}q_w^2,\qquad 1\le D\le M^2.
 \tag{FB5}
\]

There are at most 2\^{}334 exact subsets Z. This count includes both extreme subsets; it is not an independence or density assertion.

At an unmixed place with nonzero N reduction, its rank is1 or2. The accepted exact Green formula, in the normalization checked in LIFT42(FC6) and K067(PS5--7), gives

\[
 \frac{|\chi(\bar h)|}{\dim\Pi}\le\frac1{q^a-1}.
 \tag{FB6}
\]

This discards the stronger saving available at rank2. At a rank-zero place the ratio is1. At both mixed places the original normalized partial trace has modulus at most1. The full Euler correction is less than2: \(\prod(1-q^{-a})\ge1-\sum q^{-a}>1/2\). Consequently

\[
                    |a(h)|\le \frac{2D^{a/2}}{M^a}.
 \tag{FB7}
\]

The exponent a/2 in FB7 uses the paired ideal norm q². Counting a single CM prime instead would give a different and incorrect bound. No rank1-versus-rank2 distinction will be needed below.

\subsection{3. The actual difference column and a packing bound}\label{the-actual-difference-column-and-a-packing-bound}

Fix a support Σ and every actual ordered source/target pencil pair on that support. Let

\[
 I_\Sigma=\prod_{\mathfrak P\mid\Sigma}\mathfrak P^{-1},
 \qquad \mathrm N(I_\Sigma)=T^{-2s}.
 \tag{FB8}
\]

The product includes both CM partners. At a selected component the actual target W0 lattice L'\_v satisfies

\[
 \varpi\mathcal O_v^5\subset L'_v
           \subset\varpi^{-1}\mathcal O_v^5.
 \tag{FB9}
\]

Assemble the paired target lattices and the standard lattices elsewhere into the actual global lattice \(\mathcal L\subset E^5\). Exactly as in DO063(FV6--7), put

\[
 \mathcal B_\Sigma=I_\Sigma\mathcal L,
 \qquad \mathcal O_E^5\subset\mathcal B_\Sigma
                   \subset I_\Sigma^2\mathcal O_E^5,
 \qquad h e_n\in\mathcal B_\Sigma.
 \tag{FB10}
\]

Both inclusions are local consequences of FB9, including at the conjugate partner. The last assertion holds because the actual source pencil lattice contains \(\varpi e_n\), so its image lies in L'\_v. Away from Σ the actual matrix is integral.

In particular e\_n belongs to \(\mathcal B_\Sigma\). The difference column, rather than just the last column, therefore satisfies

\[
 v:=Ne_n\in\mathcal B_\Sigma,
 \qquad v\in\mathfrak J_Z\mathcal B_\Sigma
               \subset\mathfrak J_Z I_\Sigma^2\mathcal O_E^5.
 \tag{FB11}
\]

The middle membership uses the rank-zero congruence at both CM partners over G and the fact that \(\mathcal B_\Sigma\) is the standard lattice there. All other localizations are unchanged. This is an exact lattice inclusion, not a coset-density replacement.

Use the half-trace Euclidean norm \(\|v\|^2=\sum_{\tau\in\Phi}\|v_\tau\|^2\), of real dimension10f. Every nonzero coordinate x in \(\mathfrak J_Z I_\Sigma^2\) has absolute norm at least \(DT^{-4s}\): its principal ideal is contained in that fractional ideal. Arithmetic--geometric mean gives the lattice minimum

\[
 \lambda_1(\mathfrak J_Z\mathcal B_\Sigma)
       \ge\sqrt f\,z,
 \qquad z=D^{1/(2f)}T^{-2s/f}.
 \tag{FB12}
\]

Every actual difference column has norm at most \(\sqrt f(H+1)\). Disjoint open balls of radius half the lattice minimum, centered at these points, lie in the ball with radius enlarged by that amount. Comparing ordinary Euclidean volumes proves

\[
 \#\{Ne_n\}\le
             \left(1+\frac{2(H+1)}z\right)^{10f}.
 \tag{FB13}
\]

This argument requires neither a covering-radius estimate nor a claim that congruence classes are equidistributed.

For a fixed full last column and the prescribed source/target lattices, the number of W0 isometries is at most the already checked finite stabilizer bound

\[
 G_5=\gcd\bigl(|\operatorname{GL}_5(\mathbf F_{g_0})|,
              |\operatorname{GL}_5(\mathbf F_{g_1})|\bigr)
    =98148490017978252769652736000000.
 \tag{FB14}
\]

This is DO060/DO061's original fixed-W0 fibre bound. It applies to both ranks. The certificate independently recomputes the displayed gcd, without replacing the carried stabilizer argument by that calculation. Since h fixes S pointwise, no larger moving-space fibre is present.

\subsection{4. The same congruence ideal controls the real selected trace}\label{the-same-congruence-ideal-controls-the-real-selected-trace}

A nonidentity unipotent cannot fix a negative vector: its stabilizer in the real unitary group is compact, and a compact unipotent is the identity. Thus v=Ne\_n is nonzero. Choose one nonzero coordinate x of v. FB11 gives

\[
             \prod_{\tau\in\Phi}|\tau x|
                   \ge D^{1/2}T^{-2s}=z^f.
 \tag{FB15}
\]

The coordinate may depend on h; choosing it does not multiply the matrix count. At every embedding it is nonzero.

The actual near-origin transport satisfies \(\|k_{p,\tau}\|\|k_{p,\tau}^{-1}\|<2\), by LIFT41(NP12) and DO063(FV23--24). For a unipotent matrix u in U(a,1), \(\|u-I\|\le2\sinh r(u)\), by the exact singular-value identity in DO063(FV25). Hence, writing \(h_{p,\tau}=k_{p,\tau}^{-1}h_\tau k_{p,\tau}\),

\[
 |\tau x|\le\|N_\tau\|
       <2\bigl(2\sinh r(h_{p,\tau})\bigr).
 \tag{FB16}
\]

Let E\_tau be any valid lower bound for the rank-one compact energy on the actual homogeneous selected module. K060(RC13--15), with the homogeneous refinement K072(AC17--22) when used, bounds its full compression. Multiplying FB16 and using FB15 yields

\[
 |c(h)|\le
 12^f2^{f/2}\left(\prod_\tau E_\tau^{-1/4}\right)z^{-f/2}.
 \tag{FB17}
\]

The normalized trace in FB4 is bounded by the operator norm. There is no omitted factor dim W. Compact phases of h, of the near-point transport and of the selected modules are all retained by this norm inequality. No special highest vector is substituted for W.

\subsection{5. Sum every rank-zero pattern and every actual pencil type}\label{sum-every-rank-zero-pattern-and-every-actual-pencil-type}

For a fixed Z and fixed local ordered pairs, FB7, FB13, FB14 and FB17 bound the weighted matrix sum, before the pencil factors, by

\[
 \frac{2G_5\,12^f2^{f/2}}{M^a}
 \left(\prod_\tau E_\tau^{-1/4}\right)T^{2as}
 z^{(a-1/2)f}
 \left(1+\frac{2(H+1)}z\right)^{10f}.
 \tag{FB18}
\]

The z-function is increasing on z\textgreater0. Indeed, with p=(a−1/2)f, b=10f and C=2(H+1), it equals \(z^{p-b}(z+C)^b\), and p\textgreater b. Since FB5 gives

\[
 z\le z_s:=M^{1/f}T^{-2s/f},
 \qquad z_s^f=M/T^{2s},
 \tag{FB19}
\]

substitution cancels the entire factor \(M^{-a}T^{2as}z_s^{af}\). Thus FB18 is at most

\[
 2G_5\,12^f2^{f/2}
 \left(\prod_\tau E_\tau^{-1/4}\right)
 z_s^{-f/2}\left(1+\frac{2(H+1)}{z_s}\right)^{10f}.
 \tag{FB20}
\]

The argument deliberately allows every rank-zero pattern, including ones with no actual matrices. Multiplication by 2\^{}334 therefore bounds the sum over all actual patterns.

It remains to count the actual ordered pencil pairs with their weights. K071(NB7--8) gives, at a nonintegral place, type I weighted count T−1, type II count1 and type III count1. Consequently

\[
 \sum_{v=0}^s\sum_{u=0}^v
 \binom fs\binom sv\binom vu(T-1)^{s-v}
                    =\binom fs(T+1)^s.
 \tag{FB21}
\]

This is a weighted upper enumeration: each actual matrix is counted through at least one allowed pair, and extra representations only increase the upper bound. The actual local factors may have either sign. All contributions have been replaced by absolute values only after fixing the genuine K062 factors.

\textbf{Theorem FB22.} Let \(\mathcal F_s\) be the absolute contribution of all actual FB3 representatives with support size s. Then

\[
 \boxed{\displaystyle
 \mathcal F_s\le
 2^{335}G_5\binom fs(T+1)^s\,12^f2^{f/2}
 \left(\prod_\tau E_\tau^{-1/4}\right)
 z_s^{-f/2}
 \left(1+\frac{2(H+1)}{z_s}\right)^{10f}.}
 \tag{FB22}
\]

The proof includes both nilpotent ranks, all fixed and off-diagonal incidences, both mixed projectors, the actual finite frame and the actual selected trace. It does not use the sharp rank2 Green bound.

\subsection{6. Exact consequences and the remaining scope}\label{exact-consequences-and-the-remaining-scope}

At the unchanged upper endpoint, K071 carries the uniform floor E\_tau\textgreater2\^{}485. The exact companion certificate brackets z\_s by dyadic rationals using integer powers; directed rational logarithms then give

\[
                         \mathcal F_5<2^{-9199}.
 \tag{FB23}
\]

This is the complete fixed-W0 five-place sector, not only the last unique-fixed cell left by K071. R044 separately removes the complete five-place scalar-cross sector; DO065 removes the nonscalar unipotent sectors. Their errors, the earlier K068 cuts and the preexisting K062 long/integral errors remain separate. No error is silently reset to zero.

For the allowed all-top assignment, K072(AC17--22) supplies E\_tau\textgreater2\^{}959 at every label, with exactly the same CM host, finite projectors, actual point and permitted canonical twists. Hence

\[
            \mathcal F_5<2^{-29107},\qquad
            \mathcal F_6<2^{-7423}\quad\text{(all top only)}.
 \tag{FB24}
\]

The all-top assignment is an actual K040/K041/K044 mixed-profile specialization, as explicitly checked in K072. It changes no primitive CM type. It is not a literal Host/Frame/FamP constructor; the settled LIFT54--55 size obstructions remain.

For transparency, the same upper-bound formula with uniform energy 2\^{}485 gives a number between 2\^{}12484 and2\^{}12485 at s=6. Even the all-top 2\^{}959 substitution gives a number between 2\^{}14481 and2\^{}14482 at s=7. Therefore this estimate proves neither uniform six-place smallness nor all-top seven-place smallness. Failure of an upper bound to be small is not a lower bound on the actual sector.

\subsection{7. Sources and verification scope}\label{sources-and-verification-scope-8}

The only analytic inequality imported here is the already checked K060 compression theorem and its exact K072 homogeneous-energy refinement. Its primary source is:

Haagerup, U., \& Schlichtkrull, H. (2014). Inequalities for Jacobi polynomials. \emph{The Ramanujan Journal, 33}, 227--246. \url{https://doi.org/10.1007/s11139-013-9472-4} Theorem1.1 and the normalized Jacobi definition, printed p.228, with C\textless12. The direct original-PDF check is recorded in K060; this finding applies that accepted result without rederiving it.

The exact finite character input is Green's cuspidal formula as checked in LIFT42(FC6), K067 and the primary source records there. The actual pairing and lattice inputs are K062(AT13--17), K071(NB1--9), DO063(FV5--7,FV23--28), DO060/61's fixed-W0 fibre, and LIFT41(NP12). Their established scopes are retained.

The companion script \nolinkurl{fixed_block_rank_zero_packing_certificate.py} recomputes G5, brackets all displayed roots by exact integer comparisons, checks every reported binary margin with directed rational logarithms, verifies the complete incidence binomial identity on finite diagnostics and checks the symbolic exponent conditions. It does not enumerate the global rational support, numerically evaluate an automorphic kernel, prove an Mram period nonzero or execute Lean.

The entire saved FB1--24 and companion script were read and reconstructed before this author PASS, including both paired CM ideals, the essential inclusion e\_n in BΣ, the actual negative-column fibre, the compression normalization, cap monotonicity and the complete signed-incidence upper enumeration. The final arithmetic replay is a separate event. The accepted Green and Jacobi inputs were carried at their already checked scope; no new direct primary-PDF review is claimed here.

\clearpage

\section{ALP-LIFT-58: the actual longer vertical map is injective}\label{alp-lift-58-the-actual-longer-vertical-map-is-injective}

Author: /root/lift\_accounting, GPT-6 Astra Pro, ALP\_MACHINE. Status: author whole-proof check passed 2026-11-22T11:49:40.776956+00:00. Independent review is separate.

\subsection{1. Statement and genuine input hypotheses}\label{statement-and-genuine-input-hypotheses}

Retain an actual CM extension E/F of the kind constructed in K031/K032, or a real extension with the same following local properties: E/F is unramified at every finite nonsplit place, every dyadic place is relatively split, and the prescribed Hermitian lattices are selfdual. Absolute ramification of F is allowed. No full old Host or SizesAt instance is asserted.

Let m be even and use the actual vertical spaces and characters of the source, with

\[
 r=k+m-1,\quad n_j=r+jt,\quad
 q_j=\lfloor(m-j)/2\rfloor,\quad p_j=n_j-q_j,
 \quad G_j=U(Y_j),\quad 0\le j\le m.
 \tag{VI1}
\]

Here Y\_m denotes the actual positive terminal space X', and t is the arithmetic rank increment, equal to TPar.t in the canonical ofWord realization. Assume the actual Interlace, TParity and Centred hypotheses and \(t\ge r+5\), with r,t odd. At infinity use the genuine discrete series \(\tau_j=\overline{\sigma_j}\) identified by R015, with their actual nonzero real intertwiners. LIFT56's corrected theorem uses only these hypotheses on t, not a confusion of an independent TPar.t with Ranks.t.

Keep the unitary stage characters \(\mu_j\), the common odd character \(\eta\), and the actual tuple character \(\xi\), whose infinity type is zero. They are the source's characters, not freshly chosen local ones. Assume their literal finite ties: \(\eta\) is unramified; \(\mu_j\) is unramified away G with its prescribed character at G; \(\xi\) is tamely ramified at the tuple Q places and unramified elsewhere. Their angular weights are those in R015, in particular \(\operatorname{infTy}\mu_j=-e_j\). The same fixed tuple is used at every edge.

At every finite place use exactly DatumVec.IsStepVecOf and the actual frames. At G the leading principal-level compact and old determinant-block characters are the J\_j,\(\lambda_j\) of LIFT57. At Q the head rank stays r, and its level is the parahoric J(r,jt) with tail character \(\xi^{-1}\). Every other finite place has the actual selfdual lattice vacuum, including T. There is no vertical T incidence factor.

Let \(\mathcal C_j\) be the finite-dimensional cuspidal automorphic multiplicity space with real type \(\tau_j\) and these fixed finite level/character conditions. Precisely, it is the space of smooth \(G_j(F_\infty)\)-intertwiners from \(\tau_j\) to the cuspidal automorphic functions of that finite type. The domain in this Hom space is the full smooth real representation; \(\mathcal C_j\) is its multiplicity space, not the infinite-dimensional representation itself or a purported group-stable space of only K-finite scalar functions. Its finiteness is the usual fixed infinitesimal character and fixed finite level cuspidal finiteness.

Let \(D_j:\mathcal C_{j+1}\to\mathcal C_j\) be the actual downward map with its unbarred kernel and prescribed nonzero real intertwiner. Let \(U_j\) be the upward map with the conjugated kernel and the actual adjoint real intertwiner. The source convention is the one in R015 and LIFT31:

\[
 D=D_0D_1\cdots D_{m-1},\qquad
 U=U_{m-1}\cdots U_1U_0.
 \tag{VI2}
\]

\textbf{Theorem.} Under these genuine analytic and finite-data hypotheses,

\[
 \boxed{\quad U=D^*,\qquad \ker U=0,\qquad
           D(\mathcal C_m)=\mathcal C_0.\quad}
 \tag{VI3}
\]

The assertion is about the unchanged prescribed finite functions on the actual preceding image at every stage. It is not merely local occurrence or freedom to replace a later Schwartz function. For every actual idempotent e on \(\mathcal C_0\), it also gives \(eD(\mathcal C_m)=e\mathcal C_0\). It does not prove that either bottom space is nonzero, nor a nonzero bottom Bessel period, an arithmetic integral inverse, positive-rank free-data predicate3, or a complete source family with full sizes.

\subsection{2. Cuspidality and the exact dimension windows}\label{cuspidality-and-the-exact-dimension-windows}

R015 proves the actual real correspondences in the stated character orientation. R014 proves downward cuspidality for these genuine types. LIFT56 proves that every forward Witt-lower lift is zero and hence that every unchanged upward finite kernel is cuspidal, possibly zero. These facts are used as checked inputs, not re-proved here.

The global Witt index of Y\_(j+1) is q\_(j+1). Its displayed I/-I form gives at least that many hyperbolic planes, while any real embedding gives the reverse inequality. Since \(t>q_{j+1}\),

\[
 n_{j+1}-q_{j+1}>n_j,\qquad s_0=t/2,\qquad
 z_*=s_0+1/2=(t+1)/2.
 \tag{VI4}
\]

Thus every doubled theta integral lies in Weil's convergent range. Here GQT's doubled skew-Hermitian space has dimension 2n\_j and its maximal isotropic dimension is n\_j; one must not replace that n\_j by 2n\_j in the convergence criterion. The first edge has

\[
 n_1=r+t>2r=2n_0,
 \tag{VI5}
\]

while every later edge has

\[
 n_j<n_{j+1}<2n_j\qquad(j\ge1).
 \tag{VI6}
\]

The lower-tower hypothesis in GQT for those later edges is exactly the forward vanishing proved in LIFT56. It is not vacuous at a target of positive rational Witt index.

\subsection{3. Basic finite vectors, including odd source ranks}\label{basic-finite-vectors-including-odd-source-ranks}

LIFT57 supplies the entire prescribed G calculation and the actual preceding Q image. It remains to check the basic vector outside G and Q at every stage. The rank-independent argument of LIFT31 section6 applies with the following inductive bound, so we spell out its use.

At a relative split place the actual normalized basic vector is the integral rectangular-matrix indicator in the frMapO half-model. On a spherical source vector its whole-Schwartz quotient has identity value

\[
 I(f^\circ,v^\circ)=
 L_v(z_*,\Pi_j\beta_j^{-1})v^\circ,\qquad
 \beta_j=(\mu_j\xi)^{-1}\eta^{j\bmod2}.
 \tag{VI7}
\]

This is the direct all-rank spherical integral of LIFT16, also at residue characteristic2. Its value is finite and nonzero by (VI10) below. The resulting target is spherical; the fixed basic vector, not an unspecified spherical test, has therefore survived. This includes the absolutely ramified places above29 on the carried mixed hosts, which are relatively split.

At a relatively nonsplit place the residue characteristic is odd and E/F is unramified. Use the actual selfdual Hermitian lattices and the actual additive symplectic pairing. The trace-zero coefficient is a trace-zero unramified unit times an F-scalar. Move that scalar simultaneously into the additive character; selfduality makes its conductor O. This is a model isomorphism of the same representation and the same normalized vacuum, not a change of its scalar. Absolute ramification of F, including the carried places above7, is allowed.

The canonical doubling transform of the vacuum and its conjugate is the indicator of the doubled selfdual lattice, with value one at zero. Since2 is a unit, the diagonal/antidiagonal lattice change is integral, and selfdual Fourier measure fixes the indicator. The exact spherical Siegel--Weil section assertion is Gong and Grenie (2011), Lemma A.1, as checked in LIFT31. Consequently the fixed local doubling value is

\[
 Z_v(s_0,\delta(f^\circ\otimes\overline{f^\circ}),
                      v^\circ,v^\circ)
       =\frac{L_v(z_*,\phi_j\beta_j^{-1})}
                    {d_v(s_0,\beta_j)}\ne0.
 \tag{VI8}
\]

The multiplication convention in the printed GQT factor is transported by the actual Hom/tensor dictionary. At a nonsplit unramified place the uniformizer is in F and \(\beta_j(\varpi)=\pm1\), so direct and inverse characters give the same factor. At a split place (VI7) fixes the inverse character directly. This does not extrapolate an even-character determinant argument to an unfixed odd splitting.

Nonzero (VI8) implies a nonzero big theta coinvariant pairing, hence local occurrence. Only after that step do we apply Harris (2007), Theorem2.A.1(b): the image of the same lattice vacuum is the tensor of the two spherical lines. His Hecke-algebra argument, Corollary2.A.11 and its last sentence, gives the spherical parameter formula on every point of the actual theta domain. No generic-domain assumption is substituted for this occurrence check.

\subsection{4. Inductive unramified parameter and local holomorphy}\label{inductive-unramified-parameter-and-local-holomorphy}

Write \(\phi_{0,v}\) for the initial unramified standard Satake parameter. At every unramified place the actual successive spherical images satisfy

\[
 \phi_{j,v}=
 \left(\phi_{0,v}\ \oplus\
    \bigoplus_{0\le h<j}(\mu_h\xi)^{-1}
      \bigoplus_{a=0}^{t-1}|\cdot|_E^{(t-1)/2-a}\right)
             \eta^{j\bmod2}.
 \tag{VI9}
\]

This is an identity of semisimple Satake roots. It does not assert a Steinberg constituent or a global automorphic base change. At split G/Q places the corresponding full induced representation is LIFT57 (VS8), proved from the actual quotient maps.

For the initial unitary spherical rank-r component, the standard root moduli are at most \(q_E^{(r-1)/2}\), by the accepted bounded-spherical- coefficient/Satake convexity estimate of LIFT14. Every inserted block has largest exponent \((t-1)/2\). Since t\textgreater r, this gives the exact margins

\[
 z_*-(r-1)/2=(t-r+2)/2>1,\qquad
 z_*-(t-1)/2=1.
 \tag{VI10}
\]

All local standard factors in (VI7)--(VI8) are finite and nonzero. Their normalizing d\_v factors are products of finite nonzero Tate factors at positive arguments. At G the initial component need not be spherical; the unitary Godement--Jacquet bound at the first margin of (VI10) and LIFT57(VS9) give the same holomorphy for that component. At Q, its actual tame xi makes the initial local factor one.

This completes the finite induction without assuming finite temperedness of later theta images.

\subsection{5. Nonzero global partial standard values at every edge}\label{nonzero-global-partial-standard-values-at-every-edge}

Let \(\pi_j\) be an actual cuspidal descendant, and select its initial ancestor \(\pi_0\). This is a constituentwise choice; there need not be one common ancestor for a sum of different irreducible constituents. Let S contain all ramified finite places and the archimedean places. Equation (VI9) gives

\[
 \mathcal L_j^S(z_*)=
 \mathcal L^S(z_*,\phi_{\pi_0}\mu_j\xi)
       \prod_{h<j}\prod_{i=1}^t
                      L_E^S(i,\mu_j/\mu_h).
 \tag{VI11}
\]

The first term denotes only the standard unramified Euler product. Its absolute convergence and nonzero value follow from the initial rank-r root bound and the first margin in (VI10). Thus no global base-change existence or global temperedness premise is hidden.

Literal Interlace orders \(b_h<e_h<b_{h+1}<e_{h+1}\). Hence \(e_j-e_h>0\) at every selected label when h\textless j. The unitary character \(\mu_j/\mu_h\) and its square both have nonzero angular infinity type. The checked boundary lemma LIFT14(E20--22) therefore gives

\[
       0\ne L_E^S(1,\mu_j/\mu_h)\in\mathbb C.
 \tag{VI12}
\]

The lemma uses continuation and the usual positive Euler-product inequality; it does not assert absolute convergence at one. The factors with i\textgreater=2 are absolutely convergent and nonzero. Thus (VI11) is holomorphic and nonzero at the actual point.

Every omitted finite standard factor is finite and nonzero by section4 and LIFT57. At infinity the actual \(\tau_j\) is discrete series and therefore tempered; an angular unitary twist gives gamma arguments z\_* plus nonnegative angular sizes. These factors also are finite and nonzero. Thus the completed doubling standard value needed below is finite and nonzero. The normalizing unramified product d\^{}S is convergent and nonzero for s\_0\textgreater0, with GQT's actual normalization retained.

\subsection{6. Real tests and the prescribed real contraction}\label{real-tests-and-the-prescribed-real-contraction}

At every real place the actual small correspondence \(\tau_j\leftrightarrow\tau_{j+1}\) is R015, in the upward conjugated orientation. Both representations are discrete series, with a finite-dimensional representation when the relevant group is compact.

Ichino (2022), Proposition7.1, applies because n\_(j+1)\textgreater n\_j and \(\tau_j\) is tempered. It gives a nonzero local Rallis pairing on the corresponding theta quotient. Remark7.2 permits polynomial Fock vectors and K-finite representation vectors. Its absolute convergence follows from Lemma7.4 and the Harish--Chandra Schwartz coefficient estimate. This avoids the old incomplete real converse in GQT.

These are existence tests at infinity; no fixed primitive-highest value or formal-degree scalar is inferred from them. That distinction does not obstruct the actual multiplicity map. The full source map allows all real creation polynomials, and a prescribed nonzero IsItW is an intertwiner on the same Harish--Chandra modules. Real Howe multiplicity one identifies it with the same small-quotient contraction up to its actual nonzero scalar. Its adjoint gives the upward contraction. The scalar is retained throughout.

In particular no newly chosen finite Schwartz vector is used to arrange the real nonzero test. One first proves the nonzero global multiplicity map with the allowed polynomial Fock inputs, then uses this one-dimensional intertwiner comparison. There is no claim that every individually selected real highest polynomial has a nonzero pairing.

\subsection{7. The first positive-Witt edge: exact global normalization}\label{the-first-positive-witt-edge-exact-global-normalization}

This edge has target dimension \(n_1>2n_0\), so the dimension window in GQT Theorem11.4 must not be applied. The target need not be anisotropic.

Use the classical convergent formula in Liu (2011), Theorem2.1(1), printed p.859. Its preceding paragraph normalizes the target quotient Haar to volume one. For the doubled source of maximal isotropic dimension n\_0 and target dimension n\_1\textgreater2n\_0, both the Eisenstein series and theta integral converge absolutely, and

\[
 E(s_0,\Phi_\delta)
       =\int_{[G_1]}^{\mathrm{prob}}
                    \theta_\delta(g,h)\,dh.
 \tag{VI13}
\]

The same convergent identity with coefficient one is Yamana (2013), Theorem2, printed p.237; the range m\textgreater2n is explicitly attributed there to Weil. Neither statement imposes anisotropy in this range. Liu's later anisotropic simplification after Remark2.2 is not a premise of Theorem2.1(1).

The doubling transform is the actual isometry with \(\delta(\varphi_1\otimes\overline{\varphi_2})(0)
=\langle\varphi_1,\varphi_2\rangle\), including the determinant character in GQT section11.3. Unfold against the two cuspidal source forms. Absolute convergence of (VI13), moderate growth in the source variables, and cuspidal rapid decay justify the same Fubini and doubling unfolding as GQT sections11.3 and11.5--11.6. For Tamagawa Haar, the target unitary quotient has volume2. Thus

\[
 \langle\Theta_0(\varphi_1,f_1),
                 \Theta_0(\varphi_2,f_2)\rangle
 =2\,Z(s_0,\Phi_\delta,f_1,f_2).
 \tag{VI14}
\]

This proves the first-edge formula at the actual positive-Witt target, with the same factor2 as the later unitary formulas. It does not infer that a factor from probability Haar has vanished. Other specified Haar normalizations are obtained by their actual positive change-of-measure factors.

The partial standard factor is nonzero by section5. At finite G/Q the unchanged first tests have nonzero local Rallis values by LIFT57(VS31--33). At basic places section3 gives the actual test. At real places section6 provides nonzero allowed polynomial tests. The Eulerian identity therefore gives a nonzero first theta map for every nonzero cuspidal automorphic embedding of \(\pi_0\).

\subsection{8. Later global edges and automorphic multiplicities}\label{later-global-edges-and-automorphic-multiplicities}

For j\textgreater=1, (VI6) places the edge in GQT's dimension window. Equation (VI4) makes it the convergent case within that window, even when the target has positive Witt index. LIFT56 supplies the actual lower-tower vanishing. GQT Proposition11.1(ii), Theorem11.4 and section11.6 therefore apply with the actual global characters and the actual target.

All local theta quotients exist by LIFT57, section3 and R015. Their local normalized Rallis forms are nonzero: LIFT57 gives the fixed G/Q tests, section3 the fixed basics, and section6 the real test. The completed standard value is finite and nonzero by section5. It follows that the global theta map on every nonzero automorphic embedding of the preceding representation is nonzero.

This last claim is uniform in an automorphic multiplicity vector. It is not a multiplicity one assertion for the global cuspidal spectrum. For a fixed irreducible source representation, local Howe multiplicity one factors the global map as

\[
          A_j\otimes\bigotimes_v q_{j,v},
 \tag{VI15}
\]

where A\_j acts on its actual automorphic multiplicity space and q\_(j,v) is the actual small quotient, in any chosen consistent unitary realization. This is the checked cuspidal algebraic-span factorization of LIFT26(GS21--29), applied after the forward cuspidality just established.

If A\_j had a nonzero kernel vector, that vector would give a nonzero cuspidal embedding whose theta lift vanished for all Schwartz data. The preceding Rallis argument rules this out. Thus A\_j is injective. Rescaling q\_(j,v) changes A\_j inversely; no global pin is fixed by choosing a convenient local model.

\subsection{9. The fixed finite tensor preserves the entire prior image}\label{the-fixed-finite-tensor-preserves-the-entire-prior-image}

Now insert exactly the finite tensors from DatumVec, at every edge and every place.

At G, LIFT57(VS12--14) is injective on the whole prescribed type, so it applies to every vector of the actual preceding image, including its arbitrary good-series and generic leading components. At Q the initial spherical line has the precise image F\_j of LIFT57(VS17); its next image is the nonzero multiple \(\Lambda_j\) of the next flat section. This is a computation on that prior image, not a generic local occurrence theorem. At all other finite places the same actual spherical line survives by section3.

Tensoring these local injections and the injective A\_j proves injectivity of the actual upward map on every initial irreducible constituent and its entire automorphic multiplicity space. Its real factor is the fixed nonzero contraction, by section6.

Distinct initial constituents cannot cancel after the full tower: equality of terminal irreducibles would imply equality of their preceding local representations by local Howe injectivity, at every place; iteration then identifies the initial representations. Orthogonality of distinct cuspidal irreducibles and finite-dimensional fixed-level decomposition give

\[
                         \ker U=0.
 \tag{VI16}
\]

The argument does not assert that U\_j is injective on the whole independently chosen Q-level space at stage j. The actual preceding line is the necessary and proved assertion.

\subsection{10. The exact adjoint identity when both groups are noncompact}\label{the-exact-adjoint-identity-when-both-groups-are-noncompact}

The proof of LIFT31's adjoint identity extends without a compact quotient assumption because both selected scalar coefficient functions are cuspidal. We give the extra convergence check, instead of importing the formal AdjointLiftPeterssonIdentity predicate.

For one edge, let F be a larger-group multiplicity vector and u a smaller-group one. Let T\_phi and its actual adjoint T\_phi\^{}+ be the prescribed real maps. Use inner products linear in the first variable, as in the source. Since T is nonzero, choose polynomial Fock phi and compact-finite real vectors v',v with

\[
 \beta=\langle T_\phi v',v\rangle
       =\langle v',T_\phi^+v\rangle\ne0.
 \tag{VI17}
\]

Consider

\[
 \mathcal I=
 \int_{[G_j]}\int_{[G_{j+1}]}
  \theta_j(g,h;\phi)\,F(v';h)\,
                         \overline{u(v;g)}\,dh\,dg.
 \tag{VI18}
\]

For fixed Schwartz data the theta function has joint uniform moderate growth on a product of Siegel sets. More explicitly, an ordinary lattice-point estimate after the rational representation maps gives a bound \(C\|g\|^N\|h\|^N\), for some finite C,N. The fixed cuspidal scalar coefficients decay faster than every inverse power on their respective Siegel sets. Reduction theory bounds the Haar densities there by powers as well. Choose the two decay exponents larger than N and those density exponents. This proves absolute convergence of (VI18) on the product quotient, even when both real groups and both automorphic quotients are noncompact. The center is the anisotropic norm-one torus and introduces no noncompact split central direction.

The two orders of integration are now legitimate. The actual multiplicity Schur identities give, respectively,

\[
 \mathcal I=\langle D_jF,u\rangle\,\beta,
 \qquad
 \mathcal I=\langle F,U_ju\rangle\,\beta.
 \tag{VI19}
\]

The second order uses the conjugated kernel exactly as IsIoUp does. The equality of the real scalars is IsTPlus, not a newly chosen isometry. Cancelling beta proves \(U_j=D_j^*\), preserving all Haar, real and Schur factors. Composition yields U=D*. Since the spaces are finite dimensional,

\[
  (\operatorname{im}D)^\perp=\ker D^*
                         =\ker U=0.
 \tag{VI20}
\]

This proves (VI3).

\subsection{11. Costs and exact remaining gates}\label{costs-and-exact-remaining-gates}

For each split G place, the fixed unit-input raw local Rallis factor at edge j is \[
 v_j=\frac{q^{rjt+j(j-1)t^2/2}
                 |\mathrm{GL}_t(\mathbf F_q)|^j}
                {|\mathrm{GL}_{r+jt}(\mathbf F_q)|}.
 \tag{VI21}
\] At each Q place the standard-normalized factor on the actual previous line is \[
 c_{Q,j}=q^{-rjt}
  \frac{\kappa_r\kappa_{jt}\kappa_{(j+1)t}}
       {\kappa_t\kappa_{r+jt}},
 \quad c_{Q,0}=1,\quad
 q^{-rjt}/16<c_{Q,j}<q^{-rjt}\ (j>0).
 \tag{VI22}
\]

These are the exact LIFT57 factors. The Q Euler product belongs in the global standard value exactly once; the raw G factor has not already been divided by its local standard L-factor. All basic d\_v factors, boundary Hecke values, global masses, real comparison scalars and algebraic normalizations remain. There is no uniform numerical inverse bound here. For a fixed actual family the least-norm right inverse is \(D^*(DD^*)^{-1}\), whose existence does not imply an integral right inverse.

The theorem removes an unrestricted finite-kernel nonvanishing premise for this longer vertical tower. It does not prove a nonzero bottom cusp space, an actual projected bottom period, the positive-rank free-data seesaw of predicate3, the remaining same-family arithmetic row identities, or full Host/SizesAt compatibility. In particular, LIFT54--55's literal size exclusions remain unchanged. There is no conclusion about a shorter DP44 chain or its net arithmetic cost.

\subsection{12. References and exact locators}\label{references-and-exact-locators}

Accepted inputs, used at their checked scope: R014/R015 for genuine smooth types and downward cuspidality; LIFT56 for all forward lower towers; LIFT57 for fixed G/Q maps; LIFT31 section6 for the all-rank selfdual basic vector and odd-character dictionary; LIFT14(E20--22) for the boundary Hecke lemma; LIFT26(GS21--29) for the multiplicity-safe factorization.

Literal source, relative to work/research/alpoge\_source: \nolinkurl{Formalism/LocalRingCardinality/DataTheta.lean} 501--560, TPar/ofWord/Interlace/Dpar; Formalism/LocalToGlobal/IndefiniteReal/ \nolinkurl{LinksChar.lean58–71,255–287,} actual IsItW and chain inputs; the same directory's DatumVec.lean49--98; \nolinkurl{OrthogonalComplementCharacterisation.lean} 41--55,74--88,278--298, the two integrals, IsTPlus and adjoint predicate.

Gan, W. T., Qiu, Y., \& Takeda, S. (2014). The regularized Siegel--Weil formula (the second term identity) and the Rallis inner product formula. Inventiones Mathematicae, 198,739--831. Sections3.1,11.3--11.6, Proposition11.1(ii), Theorem11.4. Local lift\_refs/gqt2014.pdf.

Gong, Z., \& Grenie, L. (2011). An inequality for local unitary theta correspondence. Annales de la Faculte des Sciences de Toulouse,20(1),167--202. Lemma A.1, printed199--200. \url{https://doi.org/10.5802/afst.1289} . Local lift\_refs/gong\_grenie2011.pdf.

Harris, M. (2007). Cohomological automorphic forms on unitary groups, II: Period relations and values of L-functions. In Harmonic Analysis, Group Representations, Automorphic Forms and Invariant Theory. Corollary2.2.7 and Appendix2.A, Theorem2.A.1(b), Corollary2.A.11. \url{https://webusers.imj-prg.fr/~michael.harris/periodrelations.pdf} . Local lift\_refs/harris2007\_periodrelations.pdf.

Ichino, A. (2022). Theta lifting for tempered representations of real unitary groups. Advances in Mathematics, 398,108188. Proposition7.1, Remark7.2 and Lemma7.4, printed25--26 in the retrieved author version. \url{https://arxiv.org/pdf/2008.06174} . Directly inspected primary reference turn646674667view2, lines2461--2537.

Liu, Y. (2011). Arithmetic theta lifting and L-derivatives for unitary groups, I. Algebra \& Number Theory, 5(7),849--921. Theorem2.1(1) and its immediately preceding measure normalization, printed859. \url{https://msp.org/ant/2011/5-7/ant-v5-n7-p01-p.pdf} . Local lift\_refs/liu2011\_arithmetic\_theta.pdf and .txt, text519--535. Primary web turn619810063view2, lines678--720.

Yamana, S. (2013). The Siegel--Weil formula for unitary groups. Pacific Journal of Mathematics,264(1),235--256. Theorem2, printed237. \url{https://msp.org/pjm/2013/264-1/pjm-v264-n1-p11-p.pdf} . Local lift\_refs/yamana2013\_siegel\_weil.pdf and .txt. Primary web turn746667115view0, lines137--180.

\subsection{Author verification record}\label{author-verification-record-20}

The author read the complete twelve-section proof/VI1--22 and the complete exact diagnostic. Direct source reads covered Interlace, chrStep, full IsItW, IsIo/IsIoUp/IsTPlus and the actual Schur/adjoint predicate. Primary Liu Theorem2.1(1) and its volume-one normalization, Yamana Theorem2, Ichino Proposition7.1/Remark7.2/Lemma7.4, GQT ranges and Harris all-domain spherical statement were inspected. The finite-dimensional multiplicity-space clarification and tuple infinity type were reread. Actual finite prior-image preservation, every automorphic embedding, the first positive-Witt factor2 and noncompact absolute Fubini were checked. No bottom cusp or period existence, full source constructor, arithmetic inverse or Lean build is claimed. Diagnostic: \nolinkurl{general_vertical_global_certificate.py/.json,} 42012 exact assertions PASS 2026-11-22T11:47:28.369145+00:00.

\clearpage

\section{AM-DO-067 --- nine off-diagonal places are rigid when the integral block has rank at most one}\label{am-do-067-nine-off-diagonal-places-are-rigid-when-the-integral-block-has-rank-at-most-one}

Author: GPT-6 Astra Pro, ALP\_MACHINE \nolinkurl{/root/definite_oscillator}. 22 November 2026 UTC. Inherited reasoning effort is not separately exposed.

\textbf{Author whole-proof check PASS --- 2026-11-22T11:52:20.361856+00:00.}

On the unchanged J56 K068 residual, rank(A−I)≤1 and at most nine unequal-pencil off-diagonal T places force global unipotence. The smaller positive fixed complement makes both original coefficient-norm comparisons strict. The assertion applies at every total support size; it is not a claim for arbitrary higher block rank at nine places.

This also extends specified DO065 and DO066 bounds to their newly eligible t=9 subsets. The added union has absolute contribution less than2\^{}-4088. The remaining fixed-W0 and signed-sum gates are retained.

\subsection{1. The actual small complement}\label{the-actual-small-complement}

Keep the original J56 support, determinant-one representative h, decomposition E\^{}169=S perpendicular W0, and constants

\[
 f=a=168,\quad m=\dim S=164,\quad H=8192,
 \qquad r=\operatorname{rank}_E(A-I)\le1,
 \quad A=h_{SS}.
 \tag{ST1}
\]

We are on the checked K068 residual: every one of the334 unmixed G reductions is unipotent, h is integral away T, and its operator norm at every complex realization is at most H. No global unipotence is assumed here.

K069(IC4) gives an E-subspace P in S, fixed pointwise by h, with

\[
                 \dim_E P\ge m-r-1,
            \qquad k:=169-(m-r-1)=r+6\le7.
 \tag{ST2}
\]

This is the actual kernel of the cross map on ker(A−I); it is not a subspace inferred from a unit eigenvalue alone. P is positive at every embedding and lies in the original positive coordinate space S. Consequently its Hermitian orthogonal complement is also its Euclidean orthogonal complement. Restriction of h to it does not enlarge the operator norm H. As in DO062(RU2), pad by positive identity directions if necessary to obtain a matrix h0 of dimension k and signature(k−1,1).

The first two characteristic coefficients of N=h−I agree with those of h0−I. The coefficients remain algebraic numbers in the original field E. Choosing this complex orthogonal complement does not introduce a denominator or require an integral moving basis.

\subsection{2. The unchanged coefficient estimates in dimension six or seven}\label{the-unchanged-coefficient-estimates-in-dimension-six-or-seven}

The singular values of h0 are R,1,\ldots,1,R\^{}−1, with 1≤R≤H. The precise bounds used in DO062(RU3--4), now in dimension k, are

\[
 \begin{split}
 |e_1(N)|&\le B_1(k):=H+2k-2+H^{-1},\\
 |e_2(N)|&\le B_2(k):=(2k-3)(H+H^{-1})
   +\binom{k-2}{2}+1+(k-1)(k-2)+\binom k2.
 \end{split}
 \tag{ST3}
\]

These are the same trace and second-exterior nuclear-norm estimates, with the exact coefficient translation

\[
        e_2(h_0-I)=e_2(h_0)-(k-1)\operatorname{tr}(h_0)+\binom k2.
 \tag{ST4}
\]

They hold at every CM embedding. In particular,

\[
 \begin{array}{c|cc}
 k&B_1(k)&B_2(k)\\ \hline
 6&67190785/8192&604323849/8192\\
 7&67207169/8192&738705419/8192.
 \end{array}
 \tag{ST5}
\]

The second expressions are respectively 9(H+H\^{}−1)+42 and 11(H+H\^{}−1)+62. The first row is available when r=0, the second uniformly when r≤1.

\subsection{3. The exact global norm support}\label{the-exact-global-norm-support}

Let t count precisely the type III places of K070: the full pencils differ and their unique actual incidence is off-diagonal. At type I, h preserves both the full-pencil intersection and sum lattices. At type II it preserves its unique fixed actual lattice. Thus all its characteristic coefficients, and those of h−I, are integral at both CM components in those two cases. At type III, DO061's cyclic exterior trace argument gives pole order at most one at either component.

Therefore the only possible poles of e\_j(N) are the2t CM primes above type III places, each of absolute norm T and each with pole order at most one. This is K070(PT5--8); no standard-frame entry integrality is asserted at type II.

Each e\_j(N) vanishes modulo every unmixed G prime and its partner, since h reduces to a unipotent matrix there. The actual primes are disjoint from T, and their product of norms is Q\^{}(2(f−1)), where Q=g0g1. A nonzero coefficient would hence satisfy exactly

\[
                 Q^{2(f-1)}\le B_j(k)^{2f}T^{2t}.
 \tag{ST6}
\]

This is the norm of one coefficient, rather than a product of two. The power2t accounts for both CM partners and does not depend on the larger total nonintegral support s.

Exact integer cross-multiplication with the original DO040 T and G primes gives the following certified margins at t=9:

\begin{longtable}[]{@{}
  >{\raggedright\arraybackslash}p{(\columnwidth - 6\tabcolsep) * \real{0.2000}}
  >{\raggedleft\arraybackslash}p{(\columnwidth - 6\tabcolsep) * \real{0.2667}}
  >{\raggedleft\arraybackslash}p{(\columnwidth - 6\tabcolsep) * \real{0.2667}}
  >{\raggedleft\arraybackslash}p{(\columnwidth - 6\tabcolsep) * \real{0.2667}}@{}}
\toprule\noalign{}
\begin{minipage}[b]{\linewidth}\raggedright
Eligible block
\end{minipage} & \begin{minipage}[b]{\linewidth}\raggedleft
Residual dimension
\end{minipage} & \begin{minipage}[b]{\linewidth}\raggedleft
e1 inverse binary margin
\end{minipage} & \begin{minipage}[b]{\linewidth}\raggedleft
e2 inverse binary margin
\end{minipage} \\
\midrule\noalign{}
\endhead
\bottomrule\noalign{}
\endlastfoot
r=0 & 6 & at least1184 & at least119 \\
r≤1 & 7 & at least1184 & at least22 \\
\end{longtable}

For example the last entry means that the left side of ST6 is strictly greater than2\^{}22 times its right side at t=9. For smaller t the right side only decreases. Consequently

\[
             e_1(N)=e_2(N)=0\qquad(r\le1,\ t\le9).
 \tag{ST7}
\]

Newton's identity now gives tr(h)=tr(h²)=169. The already checked signature-one two-trace theorem DO061(section2) implies that h is unipotent at a faithful complex embedding, and hence algebraically unipotent over E. Its signature-one normal form then gives

\[
 \boxed{\ r\le1,\ t\le9\quad\Longrightarrow\quad
        h\text{ globally unipotent},\ (h-I)^3=0,
                 \operatorname{rank}_E(h-I)\le2.\ }
 \tag{ST8}
\]

This result uses the actual all-unmixed-G condition of the K068 residual. It does not follow from the archimedean norm or the pencil alone.

The certificate retains the boundaries of the argument. B2(8), which is the immediate r≤2 bound, fails ST6's strict contradiction at t=9. For both k=6 and k=7, the displayed coefficient tests fail at t=10. These failed estimates are not examples of nonunipotent supported matrices and do not prove a converse to ST8.

\subsection{4. Exactly which checked counting bounds extend}\label{exactly-which-checked-counting-bounds-extend}

For the nonscalar block r=1, ST8 supplies the sole missing hypothesis in DO065's two estimates. Their formulas depend on s and on the full rank R=rank(h−I), but not on t after global unipotence is known. Therefore the new t=9 subsets satisfy

\[
 \begin{array}{c|c|c}
 (R,r)&\text{new total supports}&\text{absolute cell bound}\\ \hline
 (1,1)&9\le s\le89&2^{-4096}\\
 (2,1)&9\le s\le106&2^{-4096}.
 \end{array}
 \tag{ST9}
\]

All actual support labels, class masses and pencil weights remain in the original DO065 counts. There are81+98=179 newly eligible cells. No r=2 counting statement is extended by ST9.

For the scalar block r=0 and nonzero cross action, the same argument extends DO066(GC22--26) to t≤9. Its proof had already bounded the sum of all allowed type coefficients by the full binomial sum2\^{}v, so this extension changes neither that bound nor any numerical constant. The only newly eligible cells in its displayed table are

\[
 \begin{array}{c|c|c}
 (v,t,u)&\text{new total supports}&\text{number of cells}\\ \hline
 (9,9,0)&9\le s\le46&38\\
 (10,9,1)&10\le s\le16&7.
 \end{array}
 \tag{ST10}
\]

Each new subset is bounded by its original entire(s,v) upper bound, less than2\^{}-4096. Their number is45. They are disjoint from ST9, because their positive block is scalar. Thus the newly bounded union has total absolute contribution

\[
                    \boxed{224\,2^{-4096}<2^{-4088}.}
 \tag{ST11}
\]

As usual, take set differences with previously removed subsets before adding errors. Alternatively, replacing DO066's old t≤8 union by its enlarged t≤9 union preserves its full1331-cell bound2\^{}-4085. These are two valid bookkeeping choices, not simultaneous extra charges.

The theorem also proves unipotence for eligible fixed-W0 matrices with t≤9, but it gives them no new absolute-sum bound. The original signed kernel, its CM phases, finite projector, real selected trace and mass normalizer remain unchanged. There is no all-T positivity, terminal nonvanishing, source-size repair or J14 claim.

\subsection{5. Verification and carried inputs}\label{verification-and-carried-inputs-2}

The exact geometry and finite support are K068, K069(IC1--4), K070(PT1--9), and DO062(RU2--6). The spectral two-trace theorem and normal form are the already whole-checked DO061(section2, TR19--21). No new external classification result is introduced. The new step is to use the actual r-dependent complement before applying the unchanged global coefficient-norm inequality.

The companion \nolinkurl{definite_oscillator_small_block_trace_certificate.py} checks every t≤9 comparison by integer cross-multiplication, retains the failed k8/t9 and k6--7/t10 estimates, checks the extended binomial type sums, and tests the original coefficient translation and bounds on96 exact signature-one spectral models. The finite diagnostics do not enumerate the rational support or infer signs of its summands.

\subsection{Author verification record}\label{author-verification-record-21}

At 2026-11-22T11:52:20.361856+00:00, I completed a whole read of ST1--11 and all five sections, with the entire final certificate. I checked the actual r-dependent positive fixed space, the common Hermitian/Euclidean complement, original-field coefficient identity, both paired finite supports, exact norm margins and every enlarged counting scope. Higher block rank, t10, fixed-W0 absolute sums and the signed period remain excluded.

The final certificate passed374 exact assertions at 2026-11-22T11:52:04.526891+00:00, including96 exact spectral models, every t\textless=9 comparison, the deliberately failed neighboring tests, and all224 newly eligible cells. A bookkeeping-only replay added the explicit45-cross-cell/224-union equality after the first373-check PASS. The theorem and coefficient inequalities did not change. Final status and scope were reread before freezing.

\clearpage

\section{AM-R-045 --- the complete five-place sector and the all-top six-place sector}\label{am-r-045-the-complete-five-place-sector-and-the-all-top-six-place-sector}

Author: GPT-6 Astra Pro, ALP\_MACHINE \nolinkurl{/root}. 22 November 2026 UTC. Inherited reasoning effort is not separately exposed.

\textbf{Author whole-proof check PASS --- 2026-11-22T11:59:35+00:00.} The entire five-section proof was read after writing; both complete error sums passed an exact rational comparison. Independent review is recorded separately. The final status and formatting were reread.

This is a corollary of the checked subset estimates, with their original normalization and separate errors retained. It consolidates the current finite nonvanishing gate. It does not evaluate the remaining signed sum.

\subsection{1. Fixed normalization and exact residual}\label{fixed-normalization-and-exact-residual}

Use the J56 upper-endpoint family of K062(AT13--17), K068, and K071, with the unchanged actual point, finite projectors, class masses, canonical twists and selected real trace. The normalized ramified energy has its original central term one. Remove a subset only by intersecting it with the remainder after all earlier removals. Write

\[
             \mathcal M=1+\mathcal S_{\mathcal R}+\varepsilon.
             \tag{JF1}
\]

Here \(\mathcal S_{\mathcal R}\) is the sum of the original signed summands over the retained set \(\mathcal R\). No removed summand is replaced by a positive value. The preexisting cuts are DO059--063, R040, K068--072, DO064--065 and R044, in their recorded scopes. The original finite support is decomposed into type I, type II and type III places as in K070/K071. As before, s is total nonintegral support, u counts type II, t counts type III, and v=u+t.

Append DO066's scalar-cross subset, K073's complete fixed-W0 five-place subset, and DO067's newly eligible t=9 subsets, taking the set difference at each step. A conservative uniform error bound is the sum of \(2^{-e}\) for

\[
\begin{aligned}
e\in(&1999,104546,101550,211446,11671,231,1483877,1307183,\\
     &9556,4084,4085,4084,4086,4088,4086,4085,9199,4088).
\end{aligned}
\tag{JF2}
\]

Repeated exponents denote different conservative subset bounds. An overlap of subsets contributes no summand twice, because the removed sets themselves are disjoint after the successive intersections. The exact rational comparison recorded with this finding gives

\[
                         |\varepsilon|<2^{-230}.
                         \tag{JF3}
\]

This error is deliberately not optimized. Earlier bounds remain in the sum even when a later estimate also covers part of an earlier subset.

\subsection{2. The complete five-place removal}\label{the-complete-five-place-removal}

Consider a retained term with s=5. The K068 residual has all 334 unmixed G reductions unipotent. Since t is at most s, K070 makes the actual global representative unipotent. Its nonzero nilpotent part has rank one or two, by DO061. The identity has no nonintegral support.

For the original positive block A, the K069/DO065 cuts remove every nonscalar possibility at s=5. In detail, the K069 fixed-space ranges remove the higher block ranks; DO065 identifies the unipotent exception with scalar A and nonzero cross action and removes the remaining nonscalar rank pairs. Thus A=I. The scalar cases split into nonzero cross action and the fixed-W0 block. DO066 removes the first case, since v is at most5 and its table includes s=5 for every such v. K073(FB23) removes the second case for both nilpotent ranks, all exact rank-zero G patterns and every actual pencil type. This exhausts the retained s=5 possibilities.

Therefore (JF1) has the stronger proved residual condition

\[
                       \boxed{\mathcal R\subset\{s\ge6\}.}
                       \tag{JF4}
\]

At s=6 the same global-unipotence and nonscalar-block arguments apply. DO066 includes every v at most6 at that support. Consequently every remaining six-place term is a fixed-W0 unipotent term, of rank one or two, still subject to all previous subset cuts. No estimate here makes that entire uniform six-place sector small.

\subsection{3. The allowed all-top specialization}\label{the-allowed-all-top-specialization}

K072(AC17--22) verifies that the all-top role assignment is an allowed specialization of the actual K040/K041/K044 mixed profile, retaining the primitive CM type and the genuine finite data. Its selected energy at every label is greater than \(2^{959}\). The uniform cuts used above continue to apply, since their lower energy bound is weaker.

At s=6 the only remaining cases are precisely the fixed-W0 terms just identified. K073(FB24) bounds their complete sector by \(2^{-7423}\). After its intersection with the previous remainder, (JF1) therefore holds in this specialization with

\[
 \boxed{\mathcal R_{\rm top}\subset\{s\ge7\},\qquad
       |\varepsilon_{\rm top}|<2^{-230}.}
       \tag{JF5}
\]

The last inequality is checked by appending 7423 to (JF2). Uniform five-place error 9199 was retained; the stronger all-top five-place bound is not needed. This is a separate actual profile, and (JF5) is not asserted uniformly over arbitrary role assignments.

\subsection{4. Other surviving restrictions and the open conclusion}\label{other-surviving-restrictions-and-the-open-conclusion}

All earlier residual restrictions remain: outside the complete degree48 fixed plane, nonzero finite coefficient, block rank at most33, a positive fixed subspace of dimension at least130, total rank at most39, and unipotence at all 334 unmixed G reductions. The fixed-space and exceptional-block tables remain effective with their original domains. K070 gives global unipotence when t is at most8; DO067 extends this to t at most9 when the block rank is at most1. DO066's complete table and the additional DO067 t=9 cells are retained, not reduced to (JF4) alone.

The sufficient signed inequality \(\mathcal S_{\mathcal R}>-1+2^{-230}\) is still unproved. Neither (JF4) nor (JF5) asserts that the remaining sum is nonnegative. The actual moved Mram period, its quantitative arithmetic comparison, the specific032 cycle identification and the DP44 cost difference remain open. These analytic mixed profiles do not satisfy the literal full source Host/SizesAt interfaces: the checked R003/R023 and LIFT54--55 obstructions remain. No J14 bound is changed.

\subsection{5. Verification and geometric interpretation}\label{verification-and-geometric-interpretation}

The proof consists of the exhaustive scalar/nonscalar and cross/fixed-W0 case split above and the exact rational error addition. The individual character, lattice, oscillator and finite-support theorems are cited at their checked interfaces, not re-proved. The contributing new proofs are DO066(GC1--27), K073(FB1--24), and DO067(ST1--11).

In the existing Grothendieck dictionary this remains a bound on the finite rational-point sum of the actual Hermitian isometry correspondence, with its oscillator trace and quotient-mass factors. It supplies no new scheme morphism, biextension, algebraic cycle, integral Picard comparison, motivic realization, or purity theorem. The previously recorded SGA3 vector-module and CM-motive locators continue to describe the underlying objects; the strict reduction is analytic and arithmetic counting on those same objects. For the run interfaces it reduces DP\_MULT Q3e and THETA3 G1's conditional period gate, without resolving either row or the CX-ALG-024 Hecke gate.

\clearpage

\section{AM-K074 --- the actual nilpotent square bounds fixed-W0 rank-two supports}\label{am-k074-the-actual-nilpotent-square-bounds-fixed-w0-rank-two-supports}

Author: GPT-6 Astra Pro, ALP\_MACHINE \nolinkurl{/root/kummer_extension}. 22 November 2026 UTC. Inherited reasoning effort is not separately exposed.

\textbf{Author whole-proof check PASS --- 2026-11-22T12:02:22.517102+00:00.} The exact companion certificate passed818,936 assertions at 2026-11-22T12:01:59.976694+00:00.

Retain the exact J56 family, constants, normalized selected trace and actual incidence types of K073(FB1--22). This finding bounds only globally unipotent fixed-W0 matrices of nilpotent rank two. It removes the whole support range5≤s≤10 with an absolute union bound2\^{}(-5183), uniformly in all carried roles and canonical twists. Global unipotence remains an explicit premise; on the carried residual use its separately proved t≤8 criterion (or a separately checked stronger criterion). Rank one and the full signed remainder are not bounded here.

\subsection{1. A second congruence cap from an actual nonzero square}\label{a-second-congruence-cap-from-an-actual-nonzero-square}

Use a=f=168, H=8192, K=2(H+1)=16386, M=Q\^{}(f−1) and G5 exactly as in K073. Put N=h−I. In this sector N³=0, N²≠0 and

\[
                 N^2=h+h^{-1}-2I,\qquad \|N_\tau^2\|\le K.
 \tag{SQ1}
\]

The inverse has the same operator norm as h in signature one, so the last inequality follows from the actual short cutoff. Choose one nonzero algebraic entry y of N². Its choice does not add a matrix count.

Fix K073's paired rank-zero ideal J at the unmixed G factors, and let D=Norm(J). Every entry of N is in J at those primes; every entry of N² is therefore in J² there. At T the actual K071(NB5,NB9) square pole budget is d2=v+u: type I contributes0, type II contributes2, and type III contributes1, at both CM partners. All other finite places are integral. Consequently the product formula gives

\[
       D^2T^{-2d_2}\le |\operatorname N_{E/\mathbf Q}y|
                      \le K^{2f},\qquad D\le K^fT^{d_2}.
 \tag{SQ2}
\]

This uses the actual square-pole bound. Squaring a crude entry pole window would lose the stated improvement. No ideal generator is chosen.

For the parameter z=D\^{}(1/(2f))T\^{}(-2s/f) in K073(FB12), SQ2 implies

\[
 z\le z_2:=\sqrt K\,T^{(d_2/2-2s)/f}
      \le\sqrt K\,T^{-s/f}<1\qquad(s\ge2).
 \tag{SQ3}
\]

The last inequality follows from K\^{}f\textless T\^{}4, checked exactly below.

\subsection{2. Weighted packing and the complete type sum}\label{weighted-packing-and-the-complete-type-sum}

K073(FB18), including the original G5 isometry fibre and the actual full-selected compression with uniform energy floor2\^{}485, is

\[
 \frac{2G_5\,12^f2^{f/2-485f/4}}{M^a}
 T^{2as}z^{(a-1/2)f}(1+K/z)^{10f}.
 \tag{SQ4}
\]

There are at most2\^{}334 exact rank-zero patterns. Since z≤z2\textless1,

\[
 z^{(a-1/2)f}(1+K/z)^{10f}
 =z^{cf}(z+K)^{10f}
 \le z_2^{cf}(K+1)^{10f},\qquad c=a-21/2=315/2.
 \tag{SQ5}
\]

Both powers are positive; the inequality uses the actual nonempty matrix's square cap. Substitution gives the T exponent

\[
             2as+c(d_2/2-2s)=21s+\frac{315}{4}d_2.
 \tag{SQ6}
\]

Let

\[
 C_2=\frac{2^{335}G_5\,12^f2^{f/2-485f/4}
                    (K+1)^{10f}K^{315f/4}}{Q^{a(f-1)}}.
 \tag{SQ7}
\]

For fixed s,v,u multiply C2 T\^{}(21s+315(v+u)/4) by the exact K071 weighted incidence count binom(f,s)binom(s,v)binom(v,u)(T−1)\^{}(s−v). Summing all types gives

\[
 \mathcal R_{2,s}\le C_2\binom fs T^{21s}
       \bigl(T-1+T^{315/4}+T^{315/2}\bigr)^s
 \le C_2\binom fs3^sT^{357s/2}.
 \tag{SQ8}
\]

Every ordered actual pair, sign replaced by its modulus, support, rank-zero pattern and original fibre is included. If only a subset of the types is known to be globally unipotent, its sum is still bounded by the same nonnegative expression. No unproved type is inserted into the mathematical sector being estimated.

\subsection{3. Exact removed range and scope}\label{exact-removed-range-and-scope}

The companion exact rational certificate gives the following inverse binary bounds for the final expression in SQ8:

\begin{longtable}[]{@{}rr@{}}
\toprule\noalign{}
s & Certified inverse binary exponent \\
\midrule\noalign{}
\endhead
\bottomrule\noalign{}
\endlastfoot
5 & 893648 \\
6 & 715955 \\
7 & 538262 \\
8 & 360570 \\
9 & 182878 \\
10 & 5186 \\
\end{longtable}

Therefore the entire globally unipotent fixed-W0 rank-two union over 5≤s≤10 has absolute contribution less than6·2\^{}(-5186)\textless2\^{}(-5183). The s=11 upper expression exceeds1, so that support is not claimed. This is a limit of the estimate, not a lower bound on the actual sum.

In particular all rank-two fixed-W0 terms at s=6 are removed uniformly; K073's rank-one contribution at s=6 still requires its all-top condition unless another independent argument is supplied. The current finding retains the K062 normalization and all previously charged errors. It proves neither all-T positivity, a period, net arithmetic gain nor a literal source-family record. The known size/link obstructions are unaffected.

\subsection{4. Verification and exact dependencies}\label{verification-and-exact-dependencies}

The new arithmetic step is SQ2: the square of the paired G ideal is combined with the independently established actual square pole budget. K073 supplies the actual difference-column packing and selected trace; K071(NB5,NB9) supplies the three local square pole orders. The underlying Jacobi estimate is the already checked Haagerup--Schlichtkrull (2014, Theorem1.1,p.228) application in K060, with its unchanged PBW metric. No new external representation theorem is invoked or rederived.

The certificate uses directed192-bit rational logarithms and exact integer powers. It checks every displayed support margin, all actual s,v,u exponent identities, the square-cap threshold and finite instances of the exact complete-type binomial identity. These are arithmetic and normalization checks, not enumeration of the automorphic support. No Lean build or numerical period computation is claimed.

The entire saved SQ1--8 and complete companion script were read and reconstructed before this author PASS. The final status paragraph was reread. Both CM-prime square valuations, the distinct actual T square poles, the monotone cap, complete type sum and all scope restrictions were checked. Independent reviews are separate records.

\clearpage

\section{AM-K075 --- exact rank-one content and norm equations remove fixed-block supports}\label{am-k075-exact-rank-one-content-and-norm-equations-remove-fixed-block-supports}

Author: GPT-6 Astra Pro, ALP\_MACHINE \nolinkurl{/root/kummer_extension}. 22 November 2026 UTC. Inherited reasoning effort is not separately exposed.

\textbf{Author whole-proof check PASS --- 2026-11-22T12:14:02.173769+00:00.} The final exact certificate passed1,441 assertions at 2026-11-22T12:13:35.424231+00:00. Independent whole reviews are separate records.

This finding replaces the ambient five-coordinate count for rank-one unipotents by their exact algebraic normal form. The coordinate-content ideal and one relative norm equation give a smaller count. The actual CM field has no prime ideals of norm below29, which makes the required ideal-divisor estimate effective without a field-dependent constant. Keeping the original pencil pole orders then removes the complete six-place rank-one fixed-W0 sector uniformly, with bound2\^{}(-8105). For the allowed all-top profile it also removes the seven- and eight-place rank-one fixed-W0 sectors, with bounds2\^{}(-17434) and 2\^{}(-7227), respectively.

All statements concern the unchanged J56 actual short sum of K062. They retain the genuine projectors, selected PBW metric, near-origin point, class masses and original finite packet. The higher-support signed remainder and all-T positivity are not proved here. The actual mixed-profile assertions do not construct literal source Host/FamP records or bypass the settled size/link obstructions.

\subsection{1. Exact field and finite-prime input}\label{exact-field-and-finite-prime-input}

Write E for the actual J56 CM host K* of K032 and F=E\^{}+. Thus

\[
 E=\mathbf Q(\zeta_{28})B\mathbf Q(\sqrt5,\sqrt{13}),
 \qquad [E:\mathbf Q]=336,
 \qquad [F:\mathbf Q]=f=168,
 \tag{RN1}
\]

where B is the real cyclic degree-seven subfield of \(\mathbf Q(\zeta_{29})\). In particular B and \(F_0=\mathbf Q(\zeta_{28})^+\) are subfields of F. Also i belongs to E, and the carried exact root group is \(\mu(E)=\mu_{28}\). These are the actual host fields; no enlarged coefficient field is used in any ideal count below.

\textbf{Lemma RN2.} Every nonzero prime ideal of either E or F has absolute norm at least29. For either field K and every nonzero integral ideal a,

\[
                   d_K(\mathfrak a)
                      \le \mathrm N(\mathfrak a)^{5/24},
 \tag{RN2}
\]

where d\_K counts integral ideal divisors.

\textbf{Proof.} The order-four subgroup fixing B in \((\mathbf Z/29\mathbf Z)^\times\) is \(H_{29}=\{1,12,17,28\}\). Among rational primes below29,

\[
                 2,3,5,7,11,13,17,19,23,
 \tag{RN3}
\]

only17 lies in H29. For every other prime in this list, Frobenius in the degree-seven quotient is nontrivial, hence has order7. Thus every prime of B over it has residue degree7. Any prime of E or F above it has a residue field containing that field, so its norm is at least 2\^{}7\textgreater29. This remains true if a different subextension is ramified at that rational prime: B itself is unramified there.

At17, the Frobenius class in the real cyclotomic field F0 has order6: the powers of17 modulo28 are17,9,13,25,5,1, with none congruent to ±1 before the sixth. Hence every E or F prime above17 has norm at least 17\^{}6\textgreater29. Rational primes at least29 already give that norm bound. The Frobenius assertions here are the elementary cyclotomic ones: raising a reduced root of unity to the p-th power induces the residue Frobenius, and passage to the fixed subfield takes its image in the specified Galois quotient. The certificate checks the complete list.

For an ideal factorization a=product P\^{}e, unique ideal factorization gives d\_K(a)=product(e+1). For e≥0, \(e+1\le2^e\), and \(29^5>2^{24}\). Thus \((e+1)\le\mathrm N(P)^{5e/24}\) for every prime factor. Multiplying proves RN2, including the unit ideal. No asymptotic divisor estimate or unspecified field constant is used. QED.

\subsection{2. The rank-one matrix is determined by its negative column}\label{the-rank-one-matrix-is-determined-by-its-negative-column}

Retain a=f=168, H=8192, the original decomposition \(E^{169}=S\perp W_0\), and the standard form \(J_0=\operatorname{diag}(1,1,1,1,-1)\) on W0. Consider an actual centrally normalized term

\[
 h=I_S\oplus h_0,\qquad N=h_0-I,
 \qquad \operatorname{rank}N=1,\quad N^2=0.
 \tag{RN4}
\]

Global unipotence is explicit in RN4. On the carried remainder it is supplied by the already checked small-block criteria; in particular DO067 covers at most nine type III places when the positive integral block is scalar. No unipotence conclusion is inferred merely from one finite reduction.

Let e be the negative last basis vector of W0 and put

\[
                         x=Ne,\qquad t=x_5=N_{55}.
 \tag{RN5}
\]

Since h is unitary and \(h^{-1}=I-N\), \(N^\dagger=J_0\overline N^{\mathsf t}J_0=-N\). A nonidentity unipotent cannot fix a negative vector, whose real stabilizer is compact. Hence x≠0. Write N=x r\^{}t with r\_5=1. Comparing the fifth column in \(N^\dagger=-N\) shows t≠0 and gives

\[
       \boxed{\quad N=\frac{x\overline x^{\mathsf t}J_0}{t},
          \qquad \bar t=-t,\qquad
          \sum_{j=1}^4x_j\bar x_j=t\bar t.\quad}
 \tag{RN6}
\]

For precision, the diagonal skew-Hermitian identity gives bar(t)=−t. The fifth-column comparison is \(-J_0\bar r\,\bar t=-x\); conjugate-transposition then gives \(r^t=\bar x^tJ_0/t\). Finally N²=0 forces \(\bar x^tJ_0x=0\). Thus every sign and the denominator in RN6 are fixed algebraically. In particular one full column x determines h exactly; there is no G5 isometry-fibre multiplier in the following count.

At every complex label, RN6 gives \(\|x\|^2=2|t|^2\), and the rank-one outer-product norm is \(\|N\|_{\rm HS}=\|x\|^2/|t|=2|t|\). The actual signature-one singular-value identity DO063(FV25) gives \(\|h_0-I\|_{\rm HS}=2\sinh r(h_0)\). The unchanged cutoff \(\|h_{0,\tau}\|\le H\) therefore implies

\[
                |\tau t|\le R:=\frac{H-H^{-1}}2
                   =\frac{67108863}{16384}.
 \tag{RN7}
\]

Also t belongs to the pure-imaginary line iF, and RN6 holds at every embedding. These are genuine constraints on the original algebraic column, not on a chosen analytic normalization.

\subsection{3. Keep the actual full-entry pole budget}\label{keep-the-actual-full-entry-pole-budget}

Fix the s nonintegral relative T places and a type assignment with u unique-fixed type II places. Types I and III each have full-N pole order at most1, whereas type II has order at most2 at both CM partners, by K071(NB5). Let

\[
 \mathfrak I=\prod_{w\in\Sigma}
       (\mathfrak P_w\bar{\mathfrak P}_w)^{-r_w},
 \qquad r_w=1+\mathbf 1_{\{w\text{ type II}\}},
 \qquad \mathrm N(\mathfrak I)=T^{-2(s+u)}.
 \tag{RN8}
\]

Fix the subset of the334 unmixed G places where N reduces to zero, and use the exact paired ideal J of K073(FB5). Its norm D satisfies D≤M², where M=Q\^{}(f−1), Q=g0g1. Every matrix entry of N belongs to

\[
 \mathfrak l=\mathfrak J\mathfrak I,
 \qquad \mathrm N(\mathfrak l)=DT^{-2(s+u)}=z^{2f},
 \qquad z=D^{1/(2f)}T^{-(s+u)/f}.
 \tag{RN9}
\]

At G this is exactly the rank-zero congruence; at T it is the actual full-N pole bound; all other localizations are integral. Both ideals are conjugation invariant. This is a valid scalar over-lattice for the actual entries, without claiming every vector in it occurs.

For any nonzero entry in l, the product formula and arithmetic--geometric mean give half-trace Euclidean norm at least sqrt(f)z. In particular, the nonzero t in RN5 forces z≤R. The lattice l intersect iF has real dimension f; its nonzero vectors have that same minimum bound. Disjoint-ball packing therefore bounds its allowable t values by

\[
                         (1+2R/z)^f.
 \tag{RN10}
\]

Only t≠0 is used in the subsequent argument; including zero in this packing upper bound is harmless.

\subsection{4. Count the coordinate-content ideal before the other coordinates}\label{count-the-coordinate-content-ideal-before-the-other-coordinates}

For a column x let b be the fractional ideal generated by its five coordinates, and let a\_N be the ideal generated by the25 entries of N. Formula RN6 gives the exact ideal identities and inclusions

\[
 (t)\subset\mathfrak b\subset\mathfrak l,
 \qquad \mathfrak b\bar{\mathfrak b}=(t)\mathfrak a_N,
 \qquad \mathfrak a_N\subset\mathfrak l.
 \tag{RN11}
\]

No ideal in this display is assumed principal except the explicitly written (t). The identity holds because all products x\_i bar(x\_j) appear, and the diagonal entries of J0 are units. It implies

\[
 \mathrm N(\mathfrak b)
   \ge |\mathrm N_{E/\mathbf Q}t|^{1/2}z^f.
 \tag{RN12}
\]

For fixed t, mapping b to b l\^{}−1 embeds its possibilities into the integral ideal divisors of (t)l\^{}−1. The latter ideal is integral, and its norm is at most (R/z)\^{}(2f), by RN7 and RN9. Applying RN2 gives

\[
                  \#\{\mathfrak b\mid t\}\le (R/z)^{70},
             \qquad 70=2f\cdot\frac5{24}.
 \tag{RN13}
\]

Here the vertical bar denotes the possibilities for fixed t, not an extra divisibility convention.

Fix t and b. On the original complex embedding space use the explicit positive metric

\[
 \|w\|_t^2=\sum_{\tau\in\Phi}
                         \frac{|\tau w|^2}{|\tau t|}.
 \tag{RN14}
\]

For a nonzero w in b, arithmetic--geometric mean and RN12 give

\[
 \|w\|_t^2\ge
 f\left(\frac{|\mathrm N_Ew|}{|\mathrm N_Et|^{1/2}}\right)^{1/f}
 \ge f\left(\frac{\mathrm N_E\mathfrak b}
                         {|\mathrm N_Et|^{1/2}}\right)^{1/f}
 \ge fz.
\]

For the first three positive coordinates, RN6 gives \(\sum_\tau\sum_{j=1}^3|\tau x_j|^2/|\tau t|
 \le\sum_\tau|\tau t|\le fR\). Thus the product lattice b\^{}3 in this specified metric has minimum at least sqrt(fz), radius at most sqrt(fR), and real dimension6f. Disjoint-ball packing bounds the number of triples by

\[
                       (1+2\sqrt{R/z})^{6f}.
 \tag{RN15}
\]

This is an arithmetic fractional-ideal count, not an unproved assertion about the saturation of a formal or analytic lattice.

\subsection{5. The fourth positive coordinate satisfies one relative norm equation}\label{the-fourth-positive-coordinate-satisfies-one-relative-norm-equation}

Once t,b,x1,x2,x3 are fixed, put

\[
                  \beta=t\bar t-\sum_{j=1}^3x_j\bar x_j\in F.
 \tag{RN16}
\]

Any remaining coordinate w=x4 lies in b and satisfies w bar(w)=beta. If beta=0 then w=0 uniquely. If a solution exists and beta≠0, every embedding of beta is positive and at most \textbar tau(t)\textbar².

For a nonzero solution, A=(w)b\^{}−1 is an integral E-ideal and

\[
 \mathrm N_{E/F}\mathfrak A
       =(\beta)(\mathrm N_{E/F}\mathfrak b)^{-1}
       =:\mathfrak c,
 \qquad \mathfrak c\text{ an integral }F\text{-ideal}.
 \tag{RN17}
\]

In particular \(\mathfrak c\mathcal O_E=(\beta)
 (\mathfrak b\bar{\mathfrak b})^{-1}\) and \(\mathfrak A\bar{\mathfrak A}=\mathfrak c\mathcal O_E\). The count is over divisors in F, not an E-ideal divisor bound. The number of integral E-ideals of relative norm c is at most d\_F(c). Indeed, at an F-prime with exponent e, a split pair gives e+1 choices; an inert prime gives at most one (only when e is even), and a ramified prime also gives at most one. Multiplication over prime factors proves the assertion, without any principality requirement.

For each such ideal A, its product A b has at most28 generators w with w bar(w)=beta. Two generators differ by an algebraic unit of relative norm1. The established CM root-unit statement makes that unit an element of the exact group mu(E)=mu28. Nonprincipal A b contributes no generator, which can only reduce the upper bound.

Finally RN7, RN12 and RN16 imply

\[
 \mathrm N_{F/\mathbf Q}\mathfrak c
   =\frac{|\mathrm N_{F/\mathbf Q}\beta|}{\mathrm N_E\mathfrak b}
   \le\frac{|\mathrm N_Et|}{|\mathrm N_Et|^{1/2}z^f}
   \le(R/z)^f.
 \tag{RN18}
\]

Apply RN2 over F. In both the zero and nonzero beta cases, the number of possible w is at most

\[
                          28(R/z)^{35}.
 \tag{RN19}
\]

For beta=0 this uses z≤R and28≥1. Combining RN10, RN13, RN15 and RN19 gives the complete rank-one matrix count for a fixed G pattern and fixed actual type/pair data:

\[
 \boxed{\quad
 28y^{105}(1+2y)^f(1+2\sqrt y)^{6f},\qquad y=R/z.
 \quad}
 \tag{RN20}
\]

The count includes all content ideals and all norm solutions. The matrix is reconstructed uniquely by RN6. No separate rank chart, compact isometry fibre or ideal generator is chosen.

\subsection{6. The same ideal and the exact selected energy}\label{the-same-ideal-and-the-exact-selected-energy}

The actual finite G trace has K073's bound \(2D^{a/2}M^{-a}\), for rank one as well as rank two. Its two mixed projectors retain modulus at most1, and the single Euler correction less than2 is unchanged. The actual normalized selected real trace, using the nonzero coordinate t, satisfies

\[
 |c(h)|\le12^f2^{f/2}
      \left(\prod_\tau E_\tau^{-1/4}\right)z^{-f/2},
 \tag{RN21}
\]

by the already checked K073(FB15--17) argument in the original PBW metric. The near-point conjugator, compact phases and normalized trace factor are retained.

Before substituting a cap, the z dependence of the product of RN20, RN21 and the finite G bound is, up to positive constants,

\[
 z^{(a-1/2)f-105-4f}
        (z+2R)^f(\sqrt z+2\sqrt R)^{6f}.
 \tag{RN22}
\]

The first exponent is27363\textgreater0, so this expression is increasing for every z\textgreater0. Since D≤M², one may replace z by

\[
                     z_{s,u}=M^{1/f}T^{-(s+u)/f}.
 \tag{RN23}
\]

The finite factor M\^{}−a T\^{}(a(s+u)) z\_\{s,u\}\^{}(af) cancels exactly. Actual nonempty columns have z≤R, but the auxiliary expression RN22 is increasing on the entire positive axis; the larger cap is valid even if z\_\{s,u\}\textgreater R. No existence assertion at that cap is used. There are at most2\^{}334 exact rank-zero G patterns, already as in K073.

\subsection{7. Sum the actual pencil types with their different pole orders}\label{sum-the-actual-pencil-types-with-their-different-pole-orders}

For a fixed number u of type II places, sum over j type III places and s−u−j type I places. K071's weighted incidence formula gives

\[
 \sum_{j=0}^{s-u}\binom fs\binom s{u+j}\binom{u+j}u
                         (T-1)^{s-u-j}
               =\binom fs\binom su T^{s-u}.
 \tag{RN24}
\]

All terms have the same full-entry pole budget s+u, so RN23 is constant through this sum. In particular the original pair multiplicities are not replaced by a single representative.

Define

\[
 \mathcal F(z)=z^{-f/2}(R/z)^{105}
                    (1+2R/z)^f(1+2\sqrt{R/z})^{6f}.
 \tag{RN25}
\]

As u increases by one, y=R/z\_\{s,u\} is multiplied by T\^{}(1/f). The function F(z), as a function of y, is R\^{}(-f/2) times \(y^{189}(1+2y)^f(1+2\sqrt y)^{6f}\). Therefore

\[
 \frac{\binom s{u+1}T^{s-u-1}\mathcal F(z_{s,u+1})}
      {\binom su T^{s-u}\mathcal F(z_{s,u})}
 \ge\frac{T^{1/8}}f>2,
 \qquad 0\le u<s\le f.
 \tag{RN26}
\]

Here189/f=9/8, the binomial ratio is at least1/f, and the exact prime satisfies T\textgreater(2f)\^{}8. Thus the full positive sum over u is less than twice its last term u=s. This is a bound on absolute weighted contributions; it does not assign a sign to the original pencil trace.

\textbf{Theorem RN27.} Let R\_(1,s) be the absolute contribution of every actual globally unipotent rank-one fixed-W0 term with support size s. Put z\_s=M\^{}(1/f)T\^{}(-2s/f), y\_s=R/z\_s. Then

\[
 \boxed{\displaystyle
 R_{1,s}\le2^{336}\,28\binom fs\,12^f2^{f/2}
 \left(\prod_\tau E_\tau^{-1/4}\right)
 z_s^{-f/2}y_s^{105}(1+2y_s)^f
                             (1+2\sqrt{y_s})^{6f}.}
 \tag{RN27}
\]

The powers2\^{}336 consist of the one finite Euler factor2, all2\^{}334 rank-zero patterns and the final geometric-series factor2. The28 counts the possible norm-one generators in RN19. Each factor has a separate origin; no additional isometry fibre is omitted.

\subsection{8. Exact actual-family consequences}\label{exact-actual-family-consequences}

K072(AC20--22), in the actual K044 upper-gap profile, proves the uniform homogeneous-energy floor E\_tau\textgreater2\^{}494 for every role, word and allowed canonical twist. For the allowed all-top assignment the floor is E\_tau\textgreater2\^{}959. Both retain the primitive mixed CM type and the actual finite data; neither is a literal source Host/FamP assertion.

The exact companion certificate brackets z\_s and sqrt(y\_s) by rational numbers and uses directed rational logarithms. Applying RN27 gives

\[
                    R_{1,6}<2^{-8105}
                        \qquad\text{uniformly in the roles},
 \tag{RN28}
\]

and, for the all-top assignment only,

\[
                    R_{1,7}<2^{-17434},\qquad
                    R_{1,8}<2^{-7227}.
 \tag{RN29}
\]

K074 independently bounds the complete globally unipotent rank-two fixed-W0 supports5 through10. Thus these two findings remove both fixed-W0 ranks at s=6 uniformly, and at s=7,8 in the allowed all-top profile. Global unipotence is available throughout these supports from the already checked residual criteria, since the off-diagonal count is at most s≤8. Their errors must be added once, using set differences from earlier cuts.

The same uniform494 formula at s=7 exceeds1; the all-top959 formula at s=9 also exceeds1. These are recorded failures of this upper estimate, not lower bounds on the actual sums or examples of surviving matrices. No larger support, complete S\_T, all-T Mram positivity, terminal period or arithmetic source-gain claim is made here.

\subsection{9. Exact sources and verification scope}\label{exact-sources-and-verification-scope}

The actual fields and inclusions are K032(QH9--11); the norm-one root unit group is the already established K029/K032 fact used explicitly in K050(BE2). The divisor and norm-equation counting lemmas RN2 and RN11--20 are proved above using ideal factorization, not an unquoted asymptotic theorem. The rank-one algebraic formula is also proved here.

The actual finite trace, near-point metric, selected normalization and pencil types are K062, K071 and K073. The energy floor is the exact K072(AC17--22) refinement of K060. Its primary analytic input remains:

Haagerup, U., \& Schlichtkrull, H. (2014). Inequalities for Jacobi polynomials. \emph{The Ramanujan Journal, 33}, 227--246. \url{https://doi.org/10.1007/s11139-013-9472-4}. Theorem1.1 and C\textless12, printed p.228, already directly audited in K060. No accepted theorem is rederived or enlarged here.

The companion certificate checks every rational prime below29, the exact subgroup/Frobenius calculations, the divisor exponent, the actual energy floors, the complete type identity, the geometric ratio, rational root brackets and binary margins. Exact rational complex rank-one models check RN6 and RN7's Hilbert--Schmidt normalization. These are finite diagnostics and exact arithmetic checks, not enumeration of the global rational kernel or a numerical period computation.

\subsection{Author whole-proof verification}\label{author-whole-proof-verification-3}

The complete saved RN1--29 and entire companion script were read and reconstructed before this author PASS. I checked the actual subfield residue degrees at every prime below29, the constant-free ideal-divisor bound, rank-one skew-Hermitian reconstruction and exact negative radius, content-ideal inclusions without principality, the rescaled metric, relative norm fibres over F, all norm-one generators, weighted cap and complete incidence sum. The final notation changes distinguishing the negative entry from the off-diagonal place count, the explicit relative-ideal extension identity, and the status text were reread. The final certificate replay is recorded separately. No new direct primary-PDF read, global support enumeration or Lean build is claimed.

\clearpage

\section{AM-R-046 --- complete six-place removal and the all-top eight-place removal}\label{am-r-046-complete-six-place-removal-and-the-all-top-eight-place-removal}

Author: GPT-6 Astra Pro, ALP\_MACHINE \nolinkurl{/root}. 22 November 2026 UTC. Inherited reasoning effort is not separately exposed.

\textbf{Author whole-proof check PASS --- 2026-11-22T12:18:44.272877+00:00.} The entire final four-section proof was reread after the explicit F/IC8 case-split correction. Both error lists passed exact rational addition; independent review is recorded separately.

This corollary combines the checked rank-one count K075 with the original finite remainder. It retains the original signed summands. All constants and normalizations are those of R045; no claim about the remaining sign or the full source admissibility is made.

\subsection{1. The residual normalization}\label{the-residual-normalization}

Let s be the number of nonintegral relative T places, let u count type II, let t count type III, and put v=u+t, exactly as in K071. Use the actual J56 upper-endpoint family, its primitive mixed CM type, all original finite projectors, selected PBW metrics, the fixed near-origin point and the original class masses. R045 and K074 give

\[
 \mathcal M=1+\mathcal S_{\mathcal R}+\varepsilon.
 \tag{JR1}
\]

Here S\_R is still the sum of the original signed terms in the remainder. Successive removed subsets are intersected with what remains after the previous removals. Thus no rational summand is charged twice.

Append K075's complete uniform rank-one fixed-W0 support-six estimate. The conservative uniform exponent list is

\[
\begin{aligned}
\mathcal E={}&(1999,104546,101550,211446,11671,231,1483877,1307183,\\
 &9556,4084,4085,4084,4086,4088,4086,4085,9199,4088,5183,8105).
\end{aligned}
\tag{JR2}
\]

The first eighteen entries are precisely R045(JF2), 5183 is K074's rank-two union at supports5 through10, and 8105 is K075(RN28). For the separately allowed all-top profile use the list E\_top formed from E by appending 7423,17434,7227. These are the earlier all-top six-place bound and K075(RN29)'s two new rank-one bounds. Retaining overlapping earlier upper bounds is conservative because the actual removed sets were made disjoint. Exact rational addition gives

\[
 \sum_{e\in\mathcal E}2^{-e}<2^{-230},\qquad
 \sum_{e\in\mathcal E_{\rm top}}2^{-e}<2^{-230}.
 \tag{JR3}
\]

No numerical cancellation among the original signed terms is used.

\subsection{2. Exhaustion of all cases at supports six through nine}\label{exhaustion-of-all-cases-at-supports-six-through-nine}

The following elementary case split is useful independently of which rank-one bound is applied. In the inherited remainder consider a term with 6\textless=s\textless=9. Put r=rank(A-I). On the fixed-space class F, the K069 table removes every r from2 through33 at these supports: its smallest threshold is22. On its exceptional complement IC8, the DO064(PM22--23) table removes the same ranks: its smallest threshold for r\textgreater=2 is R\_2=10. These two classes exhaust the possibilities, and the earlier K068 bound already excludes higher ranks. Hence r\textless=1.

If s\textless=8, K070 gives global unipotence because t\textless=s\textless=8. If s=9, the just-established r\textless=1 and t\textless=9 put the term within DO067(ST1--11), so global unipotence again holds. The same signature-one theorem DO061 used there gives N=h-I nonzero of rank one or two. It is nonzero because the identity has no nonintegral support.

For r=1, DO065's full nonscalar ranges, together with DO067's t=9 extension, remove every possibility here. The rank pairs (rank N,r) are (1,1) or (2,1), with thresholds89 and106 respectively, both larger than9. Consequently A=I.

If the scalar positive block has nonzero cross action, DO066 removes it. Indeed v\textless=s\textless=9. For v\textless=5 its table covers every s; for v=6,7,8,9 its respective thresholds are137,107,76,46, all greater than or equal to9. The sole t=9 possibility has u=0,v=9 and is one of the explicitly extended DO067 cells. Thus the old table's original t\textless=8 hypothesis has not been silently dropped.

Every remaining term at these supports is therefore fixed on the whole S block, and its only nontrivial action is on W0. K074 removes rank two throughout supports5 through10, with global unipotence now verified in the needed scope. The only cases still requiring an estimate are precisely the globally unipotent rank-one fixed-W0 terms. This exhausts the inherited remainder at s=6,7,8,9.

\subsection{3. Apply the exact rank-one norm-equation count}\label{apply-the-exact-rank-one-norm-equation-count}

K075(RN28) bounds the entire support-six rank-one fixed-W0 sector by2\^{}(-8105), uniformly in every allowed role, word and canonical twist. Together with the preceding exhaustive case split and R045's earlier support removal, (JR1) now has

\[
 \boxed{\mathcal R\subset\{s\ge7\},\qquad
                  |\varepsilon|<2^{-230}.}
 \tag{JR4}
\]

Its support-seven sector consists only of rank-one fixed-W0 terms, subject to every earlier cut. The same assertion holds for supports eight and nine if those supports are considered separately.

For the allowed all-top assignment the actual energy floor is 2\^{}959 at every label, by K072(AC20--22). This remains an allowed actual mixed profile; it is not a literal Host/FamP record. K075(RN29) bounds the complete rank-one sectors at s=7 and s=8 by2\^{}(-17434) and2\^{}(-7227). Applying the same case split gives

\[
 \boxed{\mathcal R_{\rm top}\subset\{s\ge9\},\qquad
                  |\varepsilon_{\rm top}|<2^{-230}.}
 \tag{JR5}
\]

At support nine only the rank-one fixed-W0 sector remains under this corollary. The unsuccessful K075 estimate there is an upper-bound failure, not a surviving matrix or a lower bound for the signed sum.

\subsection{4. Consequences and exact limits}\label{consequences-and-exact-limits}

All the other inherited restrictions and higher-support subset tables remain in force. The sufficient inequality S\_R\textgreater-1+2\^{}(-230), or its all-top counterpart, is still unproved. Neither (JR4) nor (JR5) supplies the moved Mram period, the specific032 arithmetic cycle identification, the missing full source records or the actual DP44 net cost. The J14 result R042 is unchanged.

The new reduction concerns the same rational points of the Hermitian isometry correspondence and their actual oscillator traces. In the recorded Grothendieck dictionary, the endomorphism module remains the one in the SGA3 I4.7.3/4.7.3.1 vector-module interpretation; the relative ideal norms in K075 are arithmetic data on Spec(O\_E)-\textgreater Spec(O\_F). No new scheme comparison, biextension map, algebraic cycle, motivic realization or purity assertion is introduced. The conditional period gates for DP\_MULT Q3e and THETA3 G1 are strictly reduced; their original questions and CX\_GRH CX-ALG-024 remain open.

The proof uses R045(JF1--5), K069's full threshold table, K070/K071's actual types, DO064(PM22--23), DO065(UR1--15), DO066(GC1--27), DO067(ST1--11), K074(SQ1--8), and K075(RN1--29), at their checked scopes. It re-proves none of their character or oscillator estimates.

\clearpage

\section{AM-K076 --- the actual content-ideal metric sharpens fixed-block rank-one counts}\label{am-k076-the-actual-content-ideal-metric-sharpens-fixed-block-rank-one-counts}

Author: GPT-6 Astra Pro, ALP\_MACHINE \nolinkurl{/root/kummer_extension}. 22 November 2026 UTC. Inherited reasoning effort is not separately exposed.

\textbf{Author whole-proof check PASS --- 2026-11-22T12:22:20.134040+00:00.}

This finding refines the three-coordinate count in K075 using its exact metric covolume and the already accepted Euclidean transference theorem. It proves a uniform rank-one fixed-W0 seven-place bound2\^{}(-2667), and an all-top rank-one fixed-W0 nine-place bound2\^{}(-2067), in the unchanged actual J56 kernel. The absolute different, real trace normalization and orthogonal-product covering radius are retained explicitly.

The sector is globally unipotent and fixed-W0, with the original G/T projectors, scalar packet and selected PBW metric. The result does not bound the full signed remainder, infer Mram positivity or instantiate literal source Host/FamP records.

\subsection{1. The exact positive metric and its covolume}\label{the-exact-positive-metric-and-its-covolume}

Use all notation of frozen K075. In particular E is the actual CM field of degree2f=336, b is a fractional coordinate-content ideal, and 0≠t∈iF is the original negative diagonal entry. Its norm bounds give

\[
 \mathrm N_E\mathfrak b\ge|\mathrm N_Et|^{1/2}z^f,
 \qquad |\tau t|\le R=\frac{67108863}{16384}.
 \tag{MC1}
\]

Set

\[
 z_b=\left(\frac{\mathrm N_E\mathfrak b}
                         {|\mathrm N_Et|^{1/2}}\right)^{1/f}
       \ge z,
 \qquad
 \|w\|_t^2=\sum_{\tau\in\Phi}\frac{|\tau w|^2}{|\tau t|}.
 \tag{MC2}
\]

This is exactly RN14's metric, not an unmentioned integral coordinate change. Let D\_E be the \textbf{absolute} different of E/Q. K032 gives

\[
 \Delta_E=2^{336}7^{280}29^{288}5^{168}13^{168},
 \quad D_0=2^{-f}\sqrt{\Delta_E}
           =7^{140}29^{144}5^{84}13^{84},
 \quad \operatorname{rd}(E)=\Delta_E^{1/(2f)}<1500.
 \tag{MC3}
\]

The unscaled half-trace Euclidean pairing is \(\operatorname{Re}\sum_\tau x_\tau\bar y_\tau
 =\tfrac12\operatorname{Tr}_{E/\mathbf Q}(x\bar y)\). Thus b has ordinary Euclidean covolume D0 Norm(b). To obtain MC2, each complex coordinate is scaled by \textbar tau(t)\textbar\^{}(-1/2). The real Jacobian is product\textbar tau(t)\textbar\^{}(-1)=\textbar Norm(t)\textbar\^{}(-1/2), giving

\[
                    \operatorname{covol}_t(\mathfrak b)
                              =D_0z_b^f.
 \tag{MC4}
\]

No relative different is substituted for the absolute one in MC3. The relative ideal norm used to count the fourth coordinate in K075 is a separate operation and is unchanged.

\subsection{2. The actual dual lattice and the transference constant}\label{the-actual-dual-lattice-and-the-transference-constant}

For the unscaled half-trace Euclidean pairing, the exact dual ideal is

\[
             \mathfrak b^*=2D_E^{-1}\bar{\mathfrak b}^{-1}.
 \tag{MC5}
\]

Indeed, pairing x∈b with y integrally means Tr(x bar(y))∈2Z; the definition of the inverse different gives MC5. Complex conjugation preserves the absolute different. This explains the factor2; deleting it would change the covering estimate.

Let A\_t be the positive diagonal real scaling by \textbar tau(t)\textbar\^{}(-1/2). The dual of A\_t b in the standard Euclidean metric is A\_t\^{}(-1) b\emph{. For nonzero y in b}, its squared norm in that dual is

\[
 \sum_\tau|\tau y|^2|\tau t|
 \ge f\bigl(|\mathrm N_Ey|\,|\mathrm N_Et|^{1/2}\bigr)^{1/f}.
 \tag{MC6}
\]

The exact ideal norm in MC5 gives \(|\mathrm N_Ey|\ge2^{2f}/(\Delta_E\mathrm N_E\mathfrak b)\). Therefore

\[
 \lambda_1\bigl((A_t\mathfrak b)^*\bigr)
       \ge\frac{2\sqrt f}{\operatorname{rd}(E)\sqrt{z_b}}.
 \tag{MC7}
\]

Apply the already checked transference inequality \(\mu(L)\lambda_1(L^*)\le\dim_{\mathbf R}(L)/2\), exactly as in LIFT47(TF7) and DO063(FV11), to the real2f-dimensional lattice A\_t b. It gives

\[
             \mu_t(\mathfrak b)
                 \le\frac{\operatorname{rd}(E)}2\sqrt{fz_b}.
 \tag{MC8}
\]

For the three-coordinate lattice use its orthogonal product structure: choose a nearest lattice vector separately in each coordinate. Hence

\[
             \mu_t(\mathfrak b^3)
                 \le\sqrt3\,\mu_t(\mathfrak b)
                 \le\frac{\operatorname{rd}(E)}2\sqrt{3fz_b}.
 \tag{MC9}
\]

This is the square root of3 product factor; applying the dimension6f bound to a single minimum without using the product would be weaker.

\subsection{3. Exact ball volume for the actual three-coordinate choices}\label{exact-ball-volume-for-the-actual-three-coordinate-choices}

K075(RN6,RN14) puts the first three positive coordinates in the ball of radius sqrt(fR) in the metric MC2. Closed Dirichlet cells have volume equal to the lattice covolume, lie within the covering radius, and have disjoint interiors. Thus their union for these lattice points fits in the ball of radius sqrt(fR)+mu\_t(b³). This is the same exact volume argument already used in DO063(FV18).

The real dimension is6f. Equations MC4 and MC9 give

\[
 \begin{split}
 \#\{x_1,x_2,x_3\}
 &\le\frac{(\pi f)^{3f}}{(3f)!D_0^3}
       \left(\sqrt{R/z_b}+
                    \frac{\sqrt3}{2}\operatorname{rd}(E)\right)^{6f}\\
 &\le A_0\left(\sqrt{R/z}+1300\right)^{6f},
 \qquad A_0=\frac{[(22/7)f]^{3f}}{(3f)!D_0^3}.
 \end{split}
 \tag{MC10}
\]

The first line cancels the complete z\_b\^{}(3f) covolume against the radius scale before estimating. Its right side decreases with z\_b, so MC1 permits replacing z\_b by z. The last line uses pi\textless22/7, rd(E)\textless1500 and 750sqrt(3)\textless1300. Every inequality is for the actual content ideal and the actual metric; no unrelated standard lattice or presumed residue density is substituted.

\subsection{4. The weighted cap and full type sum still apply}\label{the-weighted-cap-and-full-type-sum-still-apply}

All other factors of K075 remain: the pure-imaginary t count, the content-ideal divisor count over E, the fourth-coordinate norm fibres over F and their28 possible unit generators, the genuine finite trace and the actual selected compression. Replace only RN15's factor \((1+2\sqrt{R/z})^{6f}\) by MC10.

For completeness, the z-dependent part before the rank-zero cap is now proportional to

\[
 z^{27363}(z+2R)^f(\sqrt R+1300\sqrt z)^{6f},
 \tag{MC11}
\]

which is increasing for all z\textgreater0. Thus the same exact cap z\_(s,u)=M\^{}(1/f)T\^{}(-(s+u)/f) is valid, including the case that the cap itself does not describe a nonempty actual lattice intersection. As a function of y=R/z, the remaining expression has the positive factor \(y^{189}(1+2y)^f(\sqrt y+1300)^{6f}\). Increasing u by one therefore multiplies the weighted type term by at least T\^{}(1/8)/f\textgreater2, exactly as in RN26. The complete sum over all actual type assignments is still less than twice the all-type-II upper term. No sign or original multiplicity is omitted.

\textbf{Theorem MC12.} With z\_s=M\^{}(1/f)T\^{}(-2s/f) and y\_s=R/z\_s, the actual globally unipotent fixed-W0 rank-one sector obeys

\[
 \boxed{\displaystyle
 R_{1,s}\le2^{336}\,28\binom fs12^f2^{f/2}
       \left(\prod_\tau E_\tau^{-1/4}\right)
       z_s^{-f/2}y_s^{105}(1+2y_s)^f
                             A_0(\sqrt{y_s}+1300)^{6f}.}
 \tag{MC12}
\]

One may use the smaller of MC12 and the frozen K075 bound. The proof never requires the volume expression to be smaller for every possible metric; the exact certificate checks the reported numerical applications.

\subsection{5. New actual support bounds and remaining scope}\label{new-actual-support-bounds-and-remaining-scope}

Retain the actual uniform energy E\_tau\textgreater2\^{}494 and, separately, the allowed all-top E\_tau\textgreater2\^{}959 from K072. Exact rational root brackets and directed logarithms give the following inverse binary margins:

\begin{longtable}[]{@{}rrr@{}}
\toprule\noalign{}
s & Uniform494 floor & All-top959 floor \\
\midrule\noalign{}
\endhead
\bottomrule\noalign{}
\endlastfoot
6 & 9303 & 28833 \\
7 & 2667 & 22197 \\
8 & upper expression exceeds1 & 12270 \\
9 & upper expression exceeds1 & 2067 \\
10 & upper expression exceeds1 & upper expression exceeds1 \\
\end{longtable}

In particular,

\[
           R_{1,7}<2^{-2667}\quad\text{uniformly},
 \qquad R_{1,9}<2^{-2067}\quad\text{in the all-top profile}.
 \tag{MC13}
\]

K074 independently supplies the corresponding rank-two fixed-W0 bounds. DO067 supplies global unipotence for fixed-W0 terms with at most nine type III places, so that premise holds for every support in MC13 after the carried residual cuts. Thus the complete fixed-W0 sector at s=7 is bounded uniformly, and at s=9 under the all-top condition, with the two rank contributions added once.

This statement concerns fixed-W0 terms. It does not assert that every other nine-place integral-block rank is globally unipotent or bounded. The incoming error must be increased by the new absolute bounds, using set differences from previously removed terms. It is not reset to the old numerical value. Nothing here proves the sign of the entire short sum, a terminal period, arithmetic net gain or a literal source-family constructor. The failed s8-uniform and s10-all-top expressions are limits of this estimate, not lower bounds on actual contributions.

\subsection{6. Sources and verification scope}\label{sources-and-verification-scope-9}

The transference result is the precise Euclidean statement already accepted and directly checked in LIFT47(TF7), DO063(FV11):

Banaszczyk, W. (1993). New bounds in some transference theorems in the geometry of numbers. \emph{Mathematische Annalen, 296}, 625--635. \url{https://doi.org/10.1007/BF01445125}. Theorem2.2, printed p.632, in the normalization specified above.

K032(QH15) supplies the actual absolute discriminant; K075(RN11--19) supplies the exact content and relative-norm identities. K060/K072 supply the unchanged selected norm bound and actual homogeneous energy. The present proof derives the metric dual and every factor2 explicitly; it does not enlarge the accepted transference theorem.

The companion certificate verifies the source discriminant, its exact half-trace covolume, the root-discriminant and product-radius caps, all directed root/volume/norm margins and the unchanged cap/type exponents. Independent Gaussian-ideal metric calibrations test the absolute trace factor2 and orthogonal triple radius. No global rational support enumeration, period computation or Lean build is claimed.

\subsection{Author verification record}\label{author-verification-record-22}

At 2026-11-22T12:22:20.134040+00:00, I completed a whole read and reconstruction of MC1--13, every scope statement, and the entire final companion certificate. I checked the half-trace Euclidean dual with its factor2, the absolute different and discriminant, the inverse metric on the dual, the exact metric covolume, the accepted dimension/2 transference constant, the orthogonal-product radius, complete ball-volume cancellation, both weighted monotonicity arguments and all actual rank/type/energy premises. The final 225-check certificate passed at2026-11-22T12:18:04.072195+00:00 before this status update. Its checks are exact arithmetic and convention diagnostics, not a global support enumeration. The final status and this record were reread before freezing.

\clearpage

\section{AM-R-047 --- the remaining support floors are eight and ten}\label{am-r-047-the-remaining-support-floors-are-eight-and-ten}

Author: GPT-6 Astra Pro, ALP\_MACHINE \nolinkurl{/root}. 22 November 2026 UTC. Inherited reasoning effort is not separately exposed.

\textbf{Author whole-proof check PASS --- 2026-11-22T12:25:47.095722+00:00.} All three sections/JC1--4 were read in full. Both enlarged error lists passed exact rational addition. The support-ten limitation was checked against the original t\textless=9 hypothesis; independent review is separate.

\subsection{1. Exact unchanged remainder}\label{exact-unchanged-remainder}

Retain R046(JR1--5), including the actual J56 kernel, its signed summands, class masses, finite G/T projectors, canonical twists, selected PBW metric and point. Let E and E\_top denote exactly its twenty-term and twenty-three-term conservative error lists. Define

\[
 \mathcal E'=\mathcal E\mathbin{\|}(2667),\qquad
 \mathcal E'_{\rm top}=\mathcal E_{\rm top}\mathbin{\|}(2667,2067),
 \tag{JC1}
\]

where the symbol denotes concatenation of lists, with every earlier entry retained. The extra bounds are K076(MC13), first uniformly and then in the separately allowed all-top profile. The uniform support-seven bound may also be retained in the latter error list; its overlap with the earlier all-top removal is harmless after the successive set differences. Exact rational addition verifies

\[
 \sum_{e\in\mathcal E'}2^{-e}<2^{-230},\qquad
 \sum_{e\in\mathcal E'_{\rm top}}2^{-e}<2^{-230}.
 \tag{JC2}
\]

\subsection{2. Complete removal, with the proved case split}\label{complete-removal-with-the-proved-case-split}

R046(section2) already proves that, in its remaining set, every term at each support s=7,8,9 is a globally unipotent rank-one fixed-W0 term. That proof uses both K069's fixed-space table and DO064's exceptional IC8 table, followed by the exact DO067 rank-at-most-one extension at nine type III places. In particular it does not infer global unipotence at arbitrary positive-block rank or at ten type III places.

Apply K076(MC13)'s complete bound2\^{}(-2667) to the support-seven subset of that remainder. The same set-difference convention and (JC2) give

\[
 \boxed{\mathcal M=1+\mathcal S_{\mathcal R'}+\varepsilon',\qquad
 \mathcal R'\subset\{s\ge8\},\qquad |\varepsilon'|<2^{-230}.}
 \tag{JC3}
\]

At supports eight and nine, every surviving term is still a rank-one fixed-W0 unipotent, subject to the earlier cuts. K076's failed uniform support-eight upper estimate supplies neither a surviving matrix nor a lower bound for its contribution.

For the allowed all-top specialization R046 already starts at support nine. Its complete support-nine remainder is exactly of the type to which K076(MC13)'s bound2\^{}(-2067) applies. Consequently

\[
 \boxed{\mathcal M_{\rm top}=1+\mathcal S_{\mathcal R'_{\rm top}}
       +\varepsilon'_{\rm top},\qquad
 \mathcal R'_{\rm top}\subset\{s\ge10\},\qquad
 |\varepsilon'_{\rm top}|<2^{-230}.}
 \tag{JC4}
\]

No exhaustion of the support-ten terms into globally unipotent rank one is asserted. DO067's hypothesis t\textless=9 remains unchanged, and the unsuccessful all-top support-ten volume estimate would not prove such an exhaustion in any case.

\subsection{3. What has and has not been proved}\label{what-has-and-has-not-been-proved}

The nonzero central term and every original phase remain as in R046. The sufficient inequality S\_(R')\textgreater-1+2\^{}(-230), or its all-top analogue, is unproved. All earlier higher-support subset tables remain available; they are not discarded in favor of the support floor alone. J14 is unchanged. The actual moved Mram period, its quantitative arithmetic comparison, the specific032 cycle identification, complete source admissibility and the DP44 net cost remain open.

The new covolume in K076 is the covolume of the same fractional ideal module with its explicit positive metric. Its dual uses the absolute different and the half-trace factor two. In the recorded Grothendieck dictionary this refines arithmetic counting on the existing module over Spec(O\_E), as interpreted through SGA3 I4.7.3/4.7.3.1, and changes neither the biextension nor the CM motive. There is no new comparison morphism, terminal algebraic cycle, motivic weight assertion or Hecke-RH conclusion. DP\_MULT Q3e and THETA3 G1 have strictly smaller conditional period gates; their original requested conclusions are not promoted.

The only inputs are the already checked R046(JR1--5) and K076(MC1--13), with their exact hypotheses and their own underlying source locators. The new proof is the exhaustive inherited case split followed by the two exact error additions. It re-proves no accepted theorem.

\clearpage

\section{ALP-LIFT-59: a genuine goodTuple bottom cusp space at the literal level}\label{alp-lift-59-a-genuine-goodtuple-bottom-cusp-space-at-the-literal-level}

Author: /root/lift\_accounting, GPT-6 Astra Pro, ALP\_MACHINE. Status: author whole-proof check PASS at 2026-11-22T12:19:04.884794+00:00. Independent review is separate. No Lean build is claimed.

\subsection{1. Scope and conclusion}\label{scope-and-conclusion}

This result supplies a nonzero bottom cuspidal space for the longer vertical tower. It uses the actual finite goodTuple projector and the prescribed finite level, without increasing a prime or a congruence depth. It does not supply the prescribed Bessel/Mram period.

Let E/F be a CM extension of degree {[}F:Q{]}=f in the genuine setting of LIFT58. At the bottom use the standard Hermitian space and lattice

\[
 Y_0=(E^r,I_p\oplus(-I_q)),\quad
 r=k+m-1,\quad q=m/2,\quad p=r-q,\quad
 \Lambda=\mathcal O_E^r.                                      \tag{BC1}
\]

Here k,m are even, m\textgreater=2 and k\textgreater=m+3. The actual bottom compact at every finite place outside G is the stabilizer of this selfdual lattice. At the selected degree-one split G places the level is its principal residue kernel. This is exactly the bottom of the fixed finite family in LIFT58: the Q parahoric has zero tail at stage zero, and there is no vertical T incidence condition. Use the source frames to read these groups; a change of integral lattice basis does not change the assertion. The statement concerns this actual standard rational Hermitian space, rather than an unspecified space with the same real signature.

Let tau\_0 be the genuine bottom discrete series from R015 in the conjugated upward convention. Assume the literal Interlace, TParity, Centred and NCGaps conditions. Put

\[
 N=r+mt.
\]

The arithmetic increment t is the actual canonical TPar.t when the theorem is applied to the arithmetic family. As in corrected LIFT56, no inequality for an independent Ranks.t is silently imposed on an arbitrary TPar.

Assume the numerical conditions SizesAt at the point under consideration, including SizeG and the first SizeKRel clause, and the genuine finite splitting/degree-one hypotheses on the two G primes. No existence of a complete old Host/D7/SizesAt record on a mixed CM field is asserted. The theorem is an implication for genuine data satisfying these stated hypotheses. In fact the proof only needs the explicit weaker numerical conditions isolated in section 2.

At each G place the actual list of forbidden character values is the source list chi\_j together with its inverses. Use genuine finite GL regular-series projectors for eSer and hence for e0. Let C\_0 be the bottom cuspidal multiplicity space with real type tau\_0 and the fixed finite principal-G/maximal-away-G level, as in LIFT58.

\textbf{Theorem.} There is an actual goodTuple, with the required aggregate global scalar character, and a nonzero cuspidal automorphic form in its finite isotypic component and real type tau\_0. Consequently

\[
                         e_0 C_0\ne0.                        \tag{BC2}
\]

The genuine generic finite-G cut of this space is also nonzero. In conjunction with LIFT58 this gives a nonzero actual prescribed vertical chain and a nonzero image after the e0 projection. The theorem does not prove that this bottom space has a nonzero prescribed Bessel period, that its restriction has a nonzero Mram component, or that the source's other positive-rank analytic/arithmetic predicates hold. It supplies neither a primitive arithmetic inverse nor a DP44 cost comparison.

The new analytic estimate below is a normalized Flensted--Jensen coefficient bound using only adjacent noncompact HC gaps. The ordinary all-root distance bound in Kassel--Kobayashi is not substituted for this estimate.

\subsection{2. Explicit numerical hypotheses and their literal source implication}\label{explicit-numerical-hypotheses-and-their-literal-source-implication}

The proof uses

\[
 \begin{gathered}
 k\ge64,\quad 3\le f\le k,\quad m\le k-3,\quad r=k+m-1\le2k,\\
 g_i\ge k^8,\quad g_i\ge2^{16},\quad
 r\ge200f\bigl(\log(4r)+1\bigr),\quad
 \delta>N^3\ge r^3,                                      \tag{BC3}
 \end{gathered}
\]

where delta is the minimum over all real places and all adjacent noncompact HC gaps plus one. All logarithms without subscript are natural.

Here is the precise implication from the source sizes. LIFT55 proves, from the literal SizesAt alone, L\textgreater2\^{}300, k\textgreater2\^{}73197 and the other stated universal floors. In particular x\textgreater2\^{}16, k\textgreater=64 and f\textgreater=3. The source rank conditions give f\textless=k and m\textless=k-3. Its SizeG gives g\_i\textgreater=(k exp L)\^{}8\textgreater k\^{}8. The upper SizeK exponent is

\[
 9/2+(1001/100)\epsilon<5\qquad(0<\epsilon<1/52),
\]

so r\textless2x\^{}5. The lower SizeM estimate, exactly simplified in LIFT54/55, is

\[
 m\ge x^{1/2-29\epsilon/400}.
\]

Its fifth power is greater than x\^{}2. The first SizeKRel clause then yields

\[
 \frac r f\ge\frac k f\ge x^{\epsilon/100}m^5Lx^\epsilon>x^2.
                                                               \tag{BC4}
\]

For x\textgreater=2\^{}16, the function x\^{}2/log x is increasing and 2\^{}32/(16 log2)\textgreater2\^{}28\textgreater1800. Moreover

\[
 \log(4r)+1<5\log x+\log8+1<9\log x.
\]

It follows that r/f\textgreater200(log(4r)+1), proving the numerical margin in (BC3). Literal NCGaps is

\[
 (r+mt)^3<|h_i-h_j|\quad\hbox{whenever the two HC signs differ}.
                                                               \tag{BC5}
\]

After sorting a real HC list decreasingly, every adjacent pair with different signs satisfies (BC5). Adding the adjacent rho gap, which equals one, proves the delta statement. Complex conjugation/contragredience reverses the ordering but preserves these absolute gaps.

The large accepted size floors are used only to justify the simple bounds in (BC3). They are not new existence claims for those size records.

\subsection{3. A uniform good-root count for arbitrary m and r}\label{a-uniform-good-root-count-for-arbitrary-m-and-r}

We first extend the cyclic calculation in DO042 to the arbitrary ranks allowed here. Let g be an odd prime, let 4\textless=d\textless=r\textless=2k, and let C be any subset of F\_g\^{}times with \textbar C\textbar\textless=2m. Put n=g\^{}d-1 and L\_l(n)=l\^{}\{v\_l(n)\}.

The literal source primeToPart operation on the cyclic group of order n is its prime-to-l projection. Its kernel has size L\_l(n), and for a fixed c in F\_g\^{}times the equation x\_(l')=c\_(l') has L\_l(n) solutions. These exact source facts were proved in DO042 section 2; we use them as accepted facts.

We need a bound that improves with g even when m is large. If l does not divide g-1, then g-1 divides n/L\_l(n), so

\[
                  L_l(n)\le n/(g-1).                         \tag{BC6}
\]

If l divides g-1, elementary lifting the exponent gives, for odd l,

\[
 v_l(g^d-1)=v_l(g-1)+v_l(d),
\]

hence L\_l(n)\textless=d(g-1). At l=2, odd d gives the same bound; for even d the formula is

\[
 v_2(g^d-1)=v_2(g-1)+v_2(g+1)+v_2(d)-1,
\]

so L\_2(n)\textless=d(g\^{}2-1)/2. Thus the uniform bound d(g\^{}2-1) is valid in this second case. There are fewer than d log\_2 g prime divisors of n.~Combining the two cases gives

\[
 \frac1n\sum_{l\mid n}L_l(n)
 \le\frac{d\log_2g}{g-1}
    +\frac{d^2(g^2-1)\log_2g}{g^d-1}.                       \tag{BC7}
\]

A root is called strongly good if it has degree d over F\_g, every prime-to-l projection for l outside the literal BB(r,g) still has degree d, and it avoids all projected c in C for every l different from g. Full degree excludes collisions for primes not dividing n, since both projections then act as the identity.

The full-degree failures number at most d g\^{}\{d/2\}. For a nonexceptional l, the literal definition of BB gives L\_l(n)\textless=g\^{}\{d/4\}. By the cyclic subgroup count in DO042, all additional nonexceptional degree failures number at most d\^{}2(log\_2 g)g\^{}\{3d/4\}. Character collisions number at most \textbar C\textbar{} times the sum in (BC7). Therefore their total fraction is at most

\[
 \begin{split}
 B(g,d,m)={}&2m\left(\frac{d\log_2g}{g-1}
          +\frac{d^2(g^2-1)\log_2g}{g^d-1}\right)\\
 &+\frac{d g^{d/2}+d^2(\log_2g)g^{3d/4}}{g^d-1}.
 \end{split}                                                \tag{BC8}
\]

This includes exceptional-prime character collisions; only the separate injectivity condition is restricted to nonexceptional primes.

Put u=g\^{}\{1/8\}. The size hypotheses give m\textless=u and d\textless=2u. Since d\textgreater=4, g-1\textgreater=g/2 and g\^{}d-1\textgreater=g\^{}d/2, (BC8) is bounded above by

\[
 16g^{-3/4}\log_2g+16g^{-13/8}\log_2g+4g^{-15/8}.
                                                               \tag{BC9}
\]

For g\textgreater=2\^{}16, log\_2 g\textless=g\^{}\{1/4\}. This follows by equality at 2\^{}16 and the negative derivative of (log\_2 g)g\^{}\{-1/4\} thereafter. Hence (BC9) is at most

\[
 16g^{-1/2}+16g^{-11/8}+4g^{-15/8}
 \le 2^{-4}+2^{-18}+2^{-28}<1/5.                            \tag{BC10}
\]

The displayed last comparison is rational and exact. Thus, uniformly for the actual arbitrary lists C,

\[
 \#\{\hbox{strongly good degree-d roots}\}>\frac45(g^d-1).
                                                               \tag{BC11}
\]

The set is Frobenius stable. The proof requires no factorization of any enormous g\^{}d-1 and no enumeration of its roots.

\subsection{4. A literal goodTuple with the needed scalar character}\label{a-literal-goodtuple-with-the-needed-scalar-character}

By Nagura's theorem, applied to r/2\textgreater25, there is a prime p\_* with

\[
               r/2<p_*<3r/5.
                                                               \tag{BC12}
\]

In particular p\_\emph{+3\textless=r, and both p\_} and r-p\_* lie between 4 and r. At one selected place above each of the two G primes choose strongly good roots of these two degrees and take the product of their irreducible polynomials. At every other selected place choose one strongly good degree-r root and its irreducible polynomial.

The literal six clauses of IsGood now follow exactly as in DO042 section 5, with their full arbitrary-m lists retained. The roots are nonzero and the polynomials monic and squarefree. A nonexceptional projection preserves the degree of every root, so it is injective within each Frobenius orbit; roots from the two factors cannot collide because their preserved degrees differ. All projected character values and their inverses were included in C. Any full-degree place witnesses every no-factor-degree clause a\textless=r/2. The degree p\_* factor witnesses the final prime-degree clause. Thus this is the actual goodTuple predicate, not merely regular semisimplicity.

Let Z=mu(E), and write w=\textbar Z\textbar. Since Q(zeta\_w) is contained in E, phi(w)\textless={[}E:Q{]}=2f. The elementary inequality

\[
 \frac{\phi(w)^2}{w}
   =\prod_{l^a\parallel w}l^{a-2}(l-1)^2\ge1/2
\]

shows

\[
                           w\le8f^2.                         \tag{BC13}
\]

The selected primes split completely in E. Therefore w divides g\_i-1, and reduction is injective on Z at these primes.

We may prescribe the aggregate finite central character on the diagonal copy of Z. Indeed, at a full-degree place, choose a compatible generator to label the regular degree-r torus parameters. Their characters on Z are their exponents modulo w, up to the fixed embedding automorphism. Every such residue class has exactly (g\^{}r-1)/w parameter roots. By (BC11), the set of realized characters has more than 4w/5 elements. Frobenius does not change that residue because g is one modulo w. Thus the same conclusion holds for irreducible-polynomial parameters. Two such realized subsets of the cyclic character group have sum equal to the whole group: for any target omega, the sets A and omega-B have total cardinality greater than w and hence intersect.

There are at least two full-degree places left after the two mixed-degree choices, since f\textgreater=3. Fix every other parameter first and use these two places to realize any desired aggregate character Omega on Z. This is the same exact Green/Deligne--Lusztig central-character dictionary checked in DO043; no scalar twist of a fixed good polynomial is assumed. Every other literal goodTuple condition remains true.

For a regular torus parameter the genuine finite GL series consists of the corresponding generic irreducible representation. At a full-degree place it is cuspidal; at a mixed-degree place it is the generic irreducible representation with those two distinct cuspidal supports. These are the accepted genuine-series facts used in DO042/43. Let Pi be the tensor product over all G places. Then Pi belongs to e0, is generic for either chosen nondegenerate unipotent orientation after contragredience if required, and its central character on diagonal Z can be the prescribed Omega. This proves a nonzero genuine finite projector before the automorphic construction.

\subsection{5. The normalized noncompact-gap coefficient estimate}\label{the-normalized-noncompact-gap-coefficient-estimate}

We prove the analytic estimate used in the Poincare sum. Let G\_R=U(p,q), n=p+q, with J=diag(epsilon\_1,\ldots,epsilon\_n), epsilon\_i in \{+1,-1\}. Reorder coordinates so that the HC parameter h\_1\textgreater\ldots\textgreater h\_n is decreasing. The signs can be interlaced in this ordering. Let

\[
 B=\{j:1\le j<n,\ \epsilon_j\ne\epsilon_{j+1}\},\qquad
 \delta_h=\min_{j\in B}(h_j-h_{j+1}+1).
                                                               \tag{BC14}
\]

The case pq=0 is compact and needs no decay estimate. In the present bottom pq\textgreater0, so B is nonempty.

The regular HC list and its Blattner weight \(\mu_h=h+\rho_n-2\rho_c\) satisfy the dominance and integrality hypotheses required in Kassel--Kobayashi(7.8). Integrality is part of the genuine discrete-series construction in R015 under TParity; at the odd bottom rank the h\_i and rho\_n,i are integral and 2rho\_c has integral coordinates. For dominance, if i\textless l are consecutive indices of one compact-sign block, then rho\_n,i-rho\_n,l=l-i and (2rho\_c)\_i-(2rho\_c)\_l=2. Thus mu\_h,i-mu\_h,l=h\_i-h\_l+(l-i)-2\textgreater=0, since the positive compact HC difference is an integer at least one. In particular NCGaps alone is not being used as a substitute for the weight-lattice condition.

Kassel--Kobayashi's Flensted--Jensen construction, sections 7.1--7.3 and 10.2, therefore gives a K-finite discrete-series coefficient psi\_h with psi\_h(1)=1. We choose its orientation so that its right real representation is the required tau\_0; replacing by the dual reverses the HC list and does not alter the gap estimate. For physical Cartan decomposition g=k\_1 exp(Y)k\_2, Y is Hermitian, JYJ=-Y, and its eigenvalues are plus/minus the nonnegative Cartan parameters and zeros. Put R(g)=\textbar\textbar Y\textbar\textbar\_op.

The exact group-manifold normalization matters. In the dual real form from their section 5.2, physical exp(Y) is represented by the pair

\[
                    (\exp(Y/2),\exp(-Y/2)).                  \tag{BC15}
\]

Under first projection the dual group is GL\_n(C), its maximal compact is U(n), and its Iwasawa projection is ordinary QR decomposition with positive real diagonal. This follows directly from the real Lie-algebra formula in section 5.2: a Hermitian noncompact Y occurs as (Y,-Y), and a compact element as (X,X). Equivalently the dual group is the real graph g maps to (g,(g\^{}*)\^{}\{-1\}).

Here is the precise comparison with the holomorphic graph notation in their section10.2. Let sigma\_R(g)=J(g\^{}*)\^{}\{-1\}J be conjugation for the real form U(p,q), and let theta\_C(g)=JgJ be the holomorphic extension of its Cartan involution. The real dual embedding obtained from section5.2 is \[
 \iota(g)=(g,\theta_C(\sigma_R(g)))=(g,(g^*)^{-1}).
\] The holomorphic map Phi(g)=(g,theta\_C(g)) agrees with iota on the original real group, because sigma\_R(g)=g there. In particular they agree on every exp(-Y/2)ell occurring in (7.8), and on the Cartan slice (BC15). Away from that real group one must keep the real dual embedding and the holomorphic continuation distinct. The QR projection used here is that of the real dual group from section5.2, rather than an asserted identification of those two graphs on all complex matrices. This resolves the half-Cartan factor without changing the parameter dictionary.

The group-manifold HC parameter is (h,-h). On a dual diagonal (D,-D) its value is 2\textless h,D\textgreater, and the restricted roots occur with multiplicity two; rho therefore gives 2\textless rho\_n,D\textgreater, with rho\_n=((n-1)/2,\ldots,(1-n)/2). This is also the parameter dictionary in Kassel--Kobayashi section 10.2. Formula (7.8), together with Lemma 7.5(1), consequently gives

\[
 |\psi_h(k_1\exp(Y)k_2)|
 \le\int_K\prod_{j=1}^{n-1}
 \left\|\bigwedge^j(\exp(-Y/2)\ell)
                       (e_1\wedge\cdots\wedge e_j)\right\|^{-2(h_j-h_{j+1}+1)}d\ell.
                                                               \tag{BC16}
\]

The Haar measure in this formula is probability. For completeness, the QR identity behind (BC16) is that the norm of the leading j-fold wedge is the product of the first j positive QR entries. Summation by parts rewrites the weight 2(h+rho\_n) as the displayed differences. The determinant absolute value is one, so its central term vanishes. These observations fix both the factor two in the exponent and the half argument in (BC15). In rank one they recover the familiar holomorphic coefficient cosh(s)\^{}\{-k\}, before the weaker estimate below is applied.

We now give a uniform lower bound for the relevant wedge norms. By left U(n) invariance we may replace exp(-Y/2)ell by exp(A), where A=-ell\^{}\{-1\}Yell/2 is Hermitian and anticommutes with J. For every prefix highest vector v\_j=e\_1 wedge \ldots{} wedge e\_j, pi\_j(J)v\_j is a scalar of modulus one and pi\_j(J)dpi\_j(A)pi\_j(J)\^{}\{-1\}=-dpi\_j(A). The spectral measure of the self-adjoint dpi\_j(A) in v\_j is therefore symmetric. Hence

\[
 \|\pi_j(e^A)v_j\|^2
  =\langle\cosh(2d\pi_j(A))v_j,v_j\rangle\ge1.              \tag{BC17}
\]

Take the tensor product of the prefix exterior representations for j in B, with unit highest vector v. Its highest weight is lambda\^{}B=sum\_(j in B) varpi\_j, where varpi\_j are the usual fundamental GL\_n weights. For i\textless l of opposite J signs, the interval from i to l crosses a boundary in B; thus lambda\^{}B\_i-lambda\^{}B\_l\textgreater=1. Raising operators annihilate v and lowering vectors of different roots are orthogonal. The standard compact adjoint relation E\_il\^{}*=E\_li gives the exact second moment

\[
 \|d\pi(A)v\|^2
   =\sum_{i<l}|A_{li}|^2(\lambda^B_i-\lambda^B_l)
   \ge\tfrac12\|A\|_F^2.                                  \tag{BC18}
\]

Only opposite-sign entries of A can be nonzero. This proves the inequality without a dimension-dependent constant. The same J symmetry holds in the tensor product. Since x maps to cosh(2sqrt(x)) is convex on x\textgreater=0 by its power series, the spectral theorem and Jensen's inequality yield

\[
 \begin{split}
 \|\pi(e^A)v\|^2
 &\ge\cosh\bigl(2\|d\pi(A)v\|\bigr)
 \ge\cosh(\sqrt2\|A\|_F),\\
 \|\pi(e^A)v\|&\ge\cosh(\|A\|_F/\sqrt2).
 \end{split}                                                \tag{BC19}
\]

The last line uses cosh(2z)\textgreater=cosh(z)\^{}2. The norm on the left is exactly the product of the boundary prefix norms. By (BC17) every other prefix norm is at least one. Since \textbar\textbar A\textbar\textbar\_F=\textbar\textbar Y\textbar\textbar\_F/2 and \textbar\textbar Y\textbar\textbar\_F\textgreater=sqrt2 R(g), equations (BC16--19) prove

\[
 \boxed{\quad |\psi_h(g)|\le
                \cosh(R(g)/2)^{-2\delta_h}.\quad}            \tag{BC20}
\]

This is the asserted noncompact-gap estimate. It holds for every compact factor in the Cartan decomposition and every ordering of the HC signs. It is not an estimate for an arbitrary selected high-PBW coefficient; only this explicitly chosen normalized Flensted--Jensen coefficient is used in the bottom-existence argument.

\subsection{6. Integrability and the actual archimedean type}\label{integrability-and-the-actual-archimedean-type}

Miličić's integrability criterion says that a discrete series is L\^{}1 precisely when every noncompact root alpha satisfies \textbar\textless h,alpha\textgreater\textbar\textgreater k(alpha), where k(alpha)=one quarter of the sum over all roots beta of \textbar\textless alpha,beta\textgreater\textbar. For type A\_(n-1), with roots of squared length two,

\[
 \sum_\beta|\langle\alpha,\beta\rangle|=4(n-1),
 \qquad k(\alpha)=n-1.                                    \tag{BC21}
\]

Indeed alpha and -alpha contribute four, and the four roots involving any one of the other n-2 indices contribute four more. Thus literal NCGaps proves L\^{}1 for every actual indefinite discrete-series stage, not just the bottom. This uses the strict form of the criterion; equality is not substituted. Compact factors and the compact real center do not change integrability.

The Flensted--Jensen function is a finite sum of K-finite matrix coefficients in the actual discrete-series module: Kassel--Kobayashi section 10.2 identifies the group-manifold module with pi tensor pi-dual. Hence its chosen right type is tau\_0, and the product over the f real places is integrable and has that genuine full real type. Define

\[
                  \psi_\infty=\prod_{\tau:F\hookrightarrow\mathbb R}\psi_{h_\tau},
 \qquad \psi_\infty(1)=1.
\]

Every factor has absolute value at most one. If R(g)=max\_tau R(g\_tau), (BC20) and the definition of delta give

\[
                 |\psi_\infty(g)|\le\cosh(R(g)/2)^{-2\delta}.
                                                               \tag{BC22}
\]

Let omega\_infty denote its actual scalar character on diagonal Z=mu(E). Choose the finite Pi in section 4 with aggregate central character Omega=omega\_infty\^{}\{-1\}. This specifies the cancellation at rational scalar elements directly and avoids a guessed left/right central sign.

\subsection{7. Arithmetic packing and the noncompact sum}\label{arithmetic-packing-and-the-noncompact-sum}

Put K\_f\^{}max=Stab(Lambda tensor Zhat) and Gamma=U(Y\_0)(F) intersect K\_f\^{}max. Its matrices have entries in O\_E. Use the Euclidean embedding of M\_r(O\_E) into C\^{}(fr\^{}2) with squared norm sum\_tau \textbar\textbar A\_tau\textbar\textbar\_F\^{}2. A nonzero algebraic integer has sum\_tau \textbar tau(a)\textbar\^{}2\textgreater=f by the arithmetic--geometric mean and its nonzero integral absolute norm. Therefore any two distinct matrices in this lattice have distance at least sqrt(f).

If R(gamma)\textless=B, then \textbar\textbar gamma\_tau\textbar\textbar\_op\textless=exp B and \textbar\textbar gamma\_tau\textbar\textbar\_F\textless=sqrt(r)exp B. The complete matrix vector lies in the Euclidean ball of radius sqrt(fr)exp B. Disjoint balls of radius sqrt(f)/2 about these points give the elementary packing bound

\[
 \#\{\gamma\in\Gamma:R(\gamma)\le B\}
       \le(1+2\sqrt r\,e^B)^{2fr^2}
       \le(4r e^B)^{2fr^2}.                                \tag{BC23}
\]

No covolume or discriminant estimate is used here. It is an upper count in the entire integral matrix lattice, so imposing the actual unitary equations only decreases it.

Let S=Gamma intersect K\_infty. If gamma is outside S, some matrix entry in an off-sign block is a nonzero algebraic integer. At some real embedding its absolute value is at least one. In the U(p,q) Cartan decomposition the operator norm of an off-sign block is sinh R. Consequently

\[
                       R(\gamma)\ge\operatorname{arsinh}(1)>1/2.
                                                               \tag{BC24}
\]

This gap concerns the integral group used here, not a translated scalar packet with T denominators.

Set D=2fr\^{}2 and A=log(4r). For integers j\textgreater=1, use the shell j/2\textless=R\textless(j+1)/2. Since log cosh(z) is convex and vanishes at zero,

\[
 \log\cosh(j/4)\ge j\log\cosh(1/4)>j/40.
\]

The last strict inequality follows from cosh(1/4)\textgreater=33/32 and log(33/32)\textgreater=1/33\textgreater1/40. Equations (BC22--24) give

\[
 \sum_{\gamma\in\Gamma\setminus S}|\psi_\infty(\gamma)|
 \le e^{D(A+1/2)}\sum_{j\ge1}
                    e^{-j(\delta/20-D/2)}.                 \tag{BC25}
\]

By (BC3), delta\textgreater=100D(A+1). The first term on the right is at most exp(-4D(A+1)), and its successive ratio is at most exp(-4D). Since D\textgreater=1 and exp(1)\textgreater2, the right side is strictly less than

\[
                           1/15<1/4.                        \tag{BC26}
\]

This bounds the entire infinite noncompact rational sum. It does not discard a short-correspondence set, an unipotent orbit, or any determinant phase.

\subsection{8. The complete compact origin stabilizer}\label{the-complete-compact-origin-stabilizer}

The compact group S can be described exactly in the standard lattice. Each positive or negative block is an integral positive-definite unitary matrix at all real embeddings. For any column, the sum of the averages of the squared absolute values of its entries is one. Each nonzero entry has average at least one by the same norm/AM--GM argument. Hence there is exactly one nonzero entry in each column, all its conjugates have absolute value one, and it is a root of unity. Therefore

\[
              S\simeq Z^r\rtimes(S_p\times S_q),
              \qquad |S|=w^r p!q!\le w^r r!.                \tag{BC27}
\]

The scalar subgroup Z is retained, not removed by a neatness assertion.

Every s in S minus Z reduces to a nonscalar matrix at every selected G place. A nontrivial permutation has a nonzero off-diagonal root-of-unity entry. A diagonal nonscalar matrix has two distinct roots of unity, whose reductions remain distinct because g\textgreater w. Conjugation by an integral source frame preserves scalar versus nonscalar reduction.

At a full-degree place, the accepted general character estimate LIFT42(FC5) for the genuine finite cuspidal representation rho of GL\_r(F\_g) gives

\[
                 \frac{|\operatorname{tr}\rho(s)|}{\dim\rho}
                    \le\frac1{g^{r-1}-1}\quad(s\notin Z).
                                                               \tag{BC28}
\]

All other normalized finite characters have absolute value at most one by unitarity. Put

\[
 c_f(k)=\prod_{v\mid G}\frac{\operatorname{tr}\Pi_v(\bar k_v)}{\dim\Pi_v}
 \quad(k\in K_f^{max}),\qquad c_f=0\text{ off }K_f^{max}.
\]

Then c\_f(1)=1 and c\_f has the actual principal-G/maximal-away-G level. The noncentral compact contribution is bounded by

\[
 \sum_{s\in S\setminus Z}|\psi_\infty(s)c_f(s)|
 \le\frac{w^r r!}{g^{r-1}-1}
 \le2\frac{(16k^3)^r}{k^{8(r-1)}}
 \le2k^{8-4r}<1/4.                                       \tag{BC29}
\]

Here w\textless=8f\^{}2\textless=8k\^{}2, r\textless=2k, k\textgreater=64 and g\textgreater=k\^{}8 were used. The scalar terms have exact value one:

\[
                 \psi_\infty(z)c_f(z)=1\quad(z\in Z).
                                                               \tag{BC30}
\]

This is why the aggregate central-character choice in section 4 is necessary.

\subsection{9. A nonzero cuspidal Poincare series at the unchanged level}\label{a-nonzero-cuspidal-poincare-series-at-the-unchanged-level}

Consider the actual adelic function

\[
 P(g)=\sum_{\gamma\in U(Y_0)(F)}
       \psi_\infty(\gamma_\infty g_\infty)c_f(\gamma_f g_f).
                                                               \tag{BC31}
\]

At the identity its finite support reduces the sum exactly to Gamma. Equations (BC26), (BC29) and (BC30) prove

\[
                         |P(1)-w|<1/2,
 \qquad\text{in particular }P(1)\ne0.                       \tag{BC32}
\]

This is an actual normalized value estimate. It is not obtained by simply asserting eventual congruence-level nonvanishing.

We verify that the sum is a cuspidal automorphic form. The real function is smooth, Z-finite, K-finite and L\^{}1. The finite function is compactly supported and locally constant. Decompose its restriction to K\_f\^{}max into cosets of the principal-G kernel. On each relevant finite component, (BC31) is a finite sum of congruence Poincare series of left translates of psi\_infty. Muić's Lemma 4-2(i) gives absolute locally uniform convergence and cuspidality for each such series. Its stated semisimple setting applies here by restricting to SU(Y\_0): the determinant of an integral rational unitary element is a norm-one integral unit, hence lies in the finite group Z. Thus the congruence group is a finite union of cosets of its intersection with SU(Y\_0), and the real unitary center is compact. Lift to the finite central cover SU(p,q) times U(1) at every real place and apply the semisimple assertion to the finitely many resulting translates. Every proper unipotent radical lies in the derived group, so its constant terms vanish after this finite extension as well. Equivalently one may use the adelic version of the same Poincare theorem cited there. No noncompact split-center integrability issue occurs.

For a finite component not initially in K\_f\^{}max, either no rational translate meets the support or one chosen translate reduces it to the same calculation with congruence cosets. Hence this argument concerns the entire adelic function, not only its restriction to the identity component. The explicit decay and packing above also directly verify absolute convergence near the identity and on its real compact translates.

Right translation and the summation map commute. The finite coefficient c\_f lies in the Pi-isotypic module for the compact G quotient; its image P therefore lies in that same isotypic module, or is zero. It is not zero by (BC32). At infinity it lies in the specified discrete-series module tau\_0. This proves (BC2).

The finite representation Pi is generic. The compact span of P is a nonzero sum of copies of Pi; it need not be a single copy. The intersection of the kernels of all compact conjugates of the generic projector is a compact submodule, and is zero on every Pi copy. Hence some finite compact translate of P has nonzero projection to the prescribed generic character line. This preserves the principal finite level and real type. It proves the generic-cut assertion. It is not a statement about the initial T dilation characters, Mram, or a restriction to an embedded smaller unitary group.

\subsection{10. Consequence for the actual longer vertical chain, and remaining gate}\label{consequence-for-the-actual-longer-vertical-chain-and-remaining-gate}

Under LIFT58's additional genuine finite unramified/basic-vector and splitting hypotheses, its actual upward map U is injective and its actual downward map D is onto C\_0. Applying its conclusion to the genuine idempotent e0 and (BC2) gives

\[
 e_0D(C_m)=e_0C_0\ne0,\qquad U(e_0C_0)\ne0.                 \tag{BC33}
\]

Every finite kernel here is the prescribed one at its original level. The new input is the nonzero bottom cuspidal space; LIFT58's injection/surjection is cited, not re-proved.

The form in (BC31) is chosen for its normalized value at the integral origin. The prescribed bottom Bessel/Mram detector evaluates a different functional, often after a fixed rational move and with a selected higher PBW type. Neither (BC32) nor finite genericity implies that this functional is nonzero. Establishing that particular period remains necessary before invoking the fixed seesaw as a nonzero arithmetic-row statement. The theorem also does not turn the presently incompatible fixed small-host profiles into full source size instances; the large numerical assumptions here are explicit.

\subsection{11. Primary and exact source anchors}\label{primary-and-exact-source-anchors}

The source archive is pinned at commit 7cbf7a0c32fca1badf05931879a37c3f8a2b03b4.

\begin{itemize}
\tightlist
\item
  \nolinkurl{Formalism/SharpRestriction/CuspSpanNormBounds/CayleySignature.lean110–119:} literal NCGaps and Hyp3.
\item
  \nolinkurl{Formalism/LocalRingCardinality/DataTheta.lean501–560:} TPar, its independent t, Interlace/Centred and Dpar. The canonical arithmetic interpretation of t is retained.
\item
  \nolinkurl{Formalism/SharpRestriction/ElevenLe.lean24–43,123–138:} SizeM, SizeK, first SizeKRel, SizeG and SizesAt.
\item
  \nolinkurl{Formalism/CyclotomicNorms/StatementI.lean10–12} and StatementT.lean9--29: BB, primeToPart and the complete IsGood predicate.
\item
  \nolinkurl{Formalism/SharpRestriction/LevelLetterCocycles/DataN.lean996–1007:} both character orientations, goodTuple and the genuine e0 sum.
\item
  LIFT58 section1 and LIFT57: actual bottom principal-G/maximal-away-G finite family and subsequent prescribed vertical maps.
\end{itemize}

Kassel, F., \& Kobayashi, T. (2016). Poincaré series for non-Riemannian locally symmetric spaces. Advances in Mathematics, 287, 123--236. \url{https://doi.org/10.1016/j.aim.2015.08.029} . Author PDF: \url{https://www.ihes.fr/~kassel/StableSpec.pdf} . The local copy lift\_refs/kassel\_kobayashi2016\_poincare.pdf has 89 pages; exact anchors are section5.2, printed36; formula(7.8), printed50; Lemma7.5(1), printed51; and the group-manifold parameter/module dictionary, printed80--81. Their positive all-root estimate is not the new NC-only estimate (BC20).

Miličić, D. (1977). Asymptotic behaviour of matrix coefficients of the discrete series. Duke Mathematical Journal, 44(1), 59--88. Local primary lift\_refs/milicic1977\_discrete\_series.pdf; strict L\^{}p criterion printed60, with proof and final theorem printed84--85. The factor k(alpha) is computed explicitly in (BC21).

Muić, G. (2015). Fourier coefficients of automorphic forms and integrable discrete series {[}Preprint, arXiv:1505.02263{]}. arXiv. \url{https://doi.org/10.48550/arXiv.1505.02263} . Primary author preprint local lift\_refs/muic2015\_integrable\_poincare.pdf; Lemma4-2(i), printed15, and its stated adelic reference. The eventual-level part(iii) is not used to infer a fixed-level conclusion.

Nagura, J. (1952). On the interval containing at least one prime number. Proceedings of the Japan Academy, 28(4), 177--181. \url{https://doi.org/10.3792/pja/1195570997} . Original primary local lift\_refs/nagura1952\_prime\_interval.pdf, theorem printed180 and conclusion181, giving x\textless p\textless6x/5 for x\textgreater=25.

Accepted findings DO042/43 supply the literal cyclic and genuine-series/central-character dictionary. LIFT42(FC5) supplies the general nonscalar finite cuspidal character estimate. LIFT54/55 supply exact size simplifications. R015 supplies the actual discrete-series types, and LIFT58 supplies actual vertical injection and adjoint surjectivity. None of those accepted results is counted again as a new theorem here.

\subsection{12. Check scope}\label{check-scope-1}

The companion \nolinkurl{general_good_tuple_bottom_cusp_certificate.py/.json} checks the exact rational exponent estimates, finite cyclic projection fibres, boundary-weight second moment, rank-one normalization, and source fingerprints. Its 60,042 exact assertions pass. It does not construct a full Host/SizesAt record, enumerate the enormous finite fields, evaluate an actual automorphic form numerically, or prove a Bessel/Mram period. The displayed proofs, rather than finite testing, establish the universal assertions. Author whole-proof and independent reviews are recorded separately when completed.

\clearpage

\section{AM-K077 --- affine pencil congruences in the rank-one norm-equation count}\label{am-k077-affine-pencil-congruences-in-the-rank-one-norm-equation-count}

Author: GPT-6 Astra Pro, ALP\_MACHINE \nolinkurl{/root/kummer_extension}. 22 November 2026 UTC. Inherited reasoning effort is not separately exposed.

\textbf{Author whole-proof check PASS --- 2026-11-22T12:34:35.693690+00:00.}

This finding retains the exact lattice pencil in K075's rank-one count. At a double pole, the negative diagonal entry determines the leading residue of the entire negative column. For fixed negative entry and coordinate-content ideal, the other coordinates therefore belong to specified affine ideal cosets. Their difference lattice has a larger norm than the scalar over-lattice used in K075. The complete weighted pencil sum admits an exact two-endpoint bound.

In the unchanged J56 kernel, this removes the uniform rank-one fixed-W0 sectors at supports8 and9, by2\^{}(-7319) and2\^{}(-3084), respectively. For the allowed all-top profile, globally unipotent rank-one fixed-W0 sectors are bounded through support14. The higher-support assertion retains global unipotence as a hypothesis. No sign of the full remainder, all-place Mram period, literal Host/FamP constructor or DP44 comparison is inferred.

\subsection{1. The actual negative pencil entries are units}\label{the-actual-negative-pencil-entries-are-units}

Keep K075's notation: a=f=168, H=8192, the original host E/F, fixed positive coordinate space S and five-coordinate space W0 of signature (4,1), the exact G/T primes, selected metric and near-origin packet. Write e for its negative basis vector. LIFT41 and LIFT44 give

\[
 Q=g_0g_1,\quad \pi_k=\pi-7Q^2+2Q^2\sqrt{-7},\quad
 v_0=(w,\pi_k),\quad z=v_0/\bar\pi,
 \quad \pi_k\bar\pi_k=T+W,
 \quad W=77Q^4-14Q^2.
 \tag{AP1}
\]

The actual integers satisfy0\textless W\textless T. Hence at every CM prime over T, both pi\_k and bar(pi\_k) are units. At a selected component pi is a uniformizer and bar(pi) is a unit. The actual full pencil has finite members

\[
 L_c=A_c\Lambda,\qquad A_c=I+c\pi^{-1}zq,\qquad qz=0,
 \quad q=-v_0^*J-\pi m,
 \tag{AP2}
\]

where m is integral, exactly as in LIFT44(PR4--7). It follows that

\[
                 z_n\in\mathcal O_v^\times,
 \qquad q_n=\bar\pi_k-\pi m_n\in\mathcal O_v^\times.
 \tag{AP3}
\]

At the conjugate component the pencil is the inverse-adjoint dual. In split coordinates its pole column and row can be taken as Jq\^{}t and z\^{}tJ, with an irrelevant sign on c.~Their negative entries are again units. Thus every local assertion below holds at both partners. It does not identify the two selected-component matrices. The pencil changes only W0 as a lattice, by K071(NB2--3); using the full matrix in AP2 leads to the same negative-column statements.

\subsection{2. Double poles are detected by the actual negative diagonal}\label{double-poles-are-detected-by-the-actual-negative-diagonal}

Suppose a term is globally unipotent of rank one and fixed on S:

\[
 h=I_S\oplus(I+N),\quad \operatorname{rank}N=1,\quad N^2=0,
 \quad x=Ne,\quad t=x_n=N_{nn}\ne0.
 \tag{AP4}
\]

This is precisely the sector of K075; in particular bar(t)=-t. At a type-II place, h preserves a unique actual lattice L\_c.~A nonintegral term cannot fix L\_0=Lambda, so c is a residue unit. Put K=A\_c\^{}\{-1\}hA\_c, which is integral and invertible. Since(zq)\^{}2=0,

\[
 h=K+c\pi^{-1}(zqK-Kzq)
                  -c^2\pi^{-2}z(qKz)q.
 \tag{AP5}
\]

After multiplication by pi² and reduction modulo pi, the only possible nonzero term is a scalar multiple of zq. Both entries in AP3 are units. Therefore

\[
 \boxed{\ \text{some entry of N has pole2}
                    \iff v_\pi(t)=-2.\ }
 \tag{AP6}
\]

This statement includes the zero leading coefficient: if t has no double pole, the leading scalar is zero, and every entry has pole at most1. For every positive coordinate i, the same leading-column calculation gives

\[
                 x_i-(z_i/z_n)t\in\pi^{-1}\mathcal O_v.
 \tag{AP7}
\]

The quotient z\_i/z\_n is integral. At the conjugate component use its dual pole column, as above, to get the corresponding congruence. No global generator of a fractional ideal is assumed.

At types I and III the full pole-one bound already holds by K071(NB5), and at all other places the entries are integral. Because bar(t)=-t, its double-pole valuation is the same at the two CM partners. Let D be the set of relative places where it is a double pole, and put d=\textbar D\textbar. Necessarily D consists of type-II places. Thus 0\textless=d\textless=s, where s is the full nonintegral relative support.

This differs from simply counting all type-II places: a type-II place whose leading scalar vanishes has only a simple pole and is included among the remaining alternatives.

\subsection{3. The exact affine difference ideal}\label{the-exact-affine-difference-ideal}

Let I\_Sigma be the product of the inverse primes at both CM partners over the s nonintegral places, and I\_D the analogous product over D. Retain the paired rank-zero G ideal J of K073/K075, with

\[
 D_G=\mathrm N_E\mathfrak J\le M^2,\qquad M=Q^{f-1}.
\]

By AP6, every entry of N belongs to

\[
 \mathfrak l=\mathfrak J I_\Sigma I_D,
 \qquad \mathrm N_E\mathfrak l=D_GT^{-2(s+d)}=z_0^{2f},
 \qquad z_0=D_G^{1/(2f)}T^{-(s+d)/f}.
 \tag{AP8}
\]

Use z0 here to avoid confusing this positive real norm parameter with the actual local pole column z of AP1--7. K075's exact coordinate content ideal b satisfies

\[
 (t)\subset\mathfrak b\subset\mathfrak l,\qquad
 \mathrm N_E\mathfrak b\ge|\mathrm N_Et|^{1/2}z_0^f.
 \tag{AP9}
\]

Fix t and b. At a double-pole place AP7 puts each positive coordinate in a specified affine coset of pi\^{}(-1)O\_v. At a simple-pole place this condition is already automatic. All of the other constraints remain membership in b. Therefore the possibilities for each positive coordinate, if nonempty, form a coset of

\[
                    \mathfrak c=\mathfrak b\cap\mathfrak J I_\Sigma.
 \tag{AP10}
\]

Only a containing coset is needed if other matrix conditions cut it down further. The same t and same prescribed local pencil data give this statement separately for every coordinate. CRT identifies the intersection of these local additive conditions with an affine coset of the global fractional ideal; no choice of a principal ideal is used.

At either prime over a place of D, t has valuation-2, all coordinates have valuation at least-2 and one is t itself. Thus b has valuation-2, whereas c has valuation-1. At every other prime b is already contained in J I\_Sigma, so the two valuations agree. Consequently

\[
          \boxed{\quad\mathrm N_E\mathfrak c
                       =\mathrm N_E\mathfrak b\,T^{2d}.\quad}
 \tag{AP11}
\]

This is the exact index recovered from the actual pencil. It is not an assumption that congruence classes have their expected density.

\subsection{4. Retain the content and relative norm equation}\label{retain-the-content-and-relative-norm-equation}

Let R=(H-H\^{}(-1))/2=67108863/16384. The exact normal form K075(RN6--7) continues to give

\[
 |\tau t|\le R,\qquad
 N=x\bar x^tJ_0/t,\qquad
                \sum_{i=1}^{4}x_i\bar x_i=t\bar t.
 \tag{AP12}
\]

The permitted nonzero t values are bounded by(1+2R/z0)\^{}f.~For each t, the content-ideal possibilities are bounded by(R/z0)\^{}70. These are K075(RN10--13) with the sharper actual ideal AP8, not new divisor estimates. They use z0\textless=R on the nonempty actual set.

In the same original positive metric as K075 and K076,

\[
                \|w\|_t^2=\sum_\tau\frac{|\tau w|^2}{|\tau t|},
 \tag{AP13}
\]

any nonzero difference in c has squared norm at least

\[
 f\left(\frac{\mathrm N_E\mathfrak c}
                       {|\mathrm N_Et|^{1/2}}\right)^{1/f}
                    \ge fz_0T^{2d/f}.
 \tag{AP14}
\]

The first three coordinates lie in an affine coset of c³, inside the ball of squared radius fR by AP12. Disjoint-ball packing depends only on pairwise differences, and therefore applies equally to that affine coset. Its real dimension is6f. The number of triples is at most

\[
                 \left(1+2\sqrt{\frac R{z_0T^{2d/f}}}\right)^{6f}.
 \tag{AP15}
\]

For each triple, the fourth relative norm equation is unchanged. K075(RN16--19) uses the original content b, not c: its ideal norm over F has norm at most(R/z0)\^{}f, so the number of fourth coordinates is at most28(R/z0)\^{}35. Keeping the content estimate with b is essential; AP11 is not inserted into that norm equation without justification.

Thus the complete allowable column count for a fixed actual pencil choice, support and marked double-pole set is at most

\[
 28y^{105}(1+2y)^f(1+2\sqrt{y_c})^{6f},
 \qquad y=R/z_0,\quad y_c=R/(z_0T^{2d/f}).
 \tag{AP16}
\]

Each column determines the rank-one matrix exactly. There is no additional isometry-fibre factor. Discarding other pencil conditions only enlarges the upper count.

\subsection{5. The actual weighted type sum and rank-zero cap}\label{the-actual-weighted-type-sum-and-rank-zero-cap}

Retain all2\^{}334 paired rank-zero G patterns. Their finite normalized trace is bounded by2 D\_G\^{}(a/2)/M\^{}a, exactly as in K075. The same nonzero t witness gives the selected real factor

\[
          12^f2^{f/2}\left(\prod_\tau E_\tau^{-1/4}\right)z_0^{-f/2}.
 \tag{AP17}
\]

Fix support Sigma and D inside it. Each place of D must be type II and has weighted ordered-pair cost1. At every other place all three alternatives are permitted: the K071 costs are T-1,1,1 for types I, II,III. Their full sum is T+1. Summing supports and marked sets gives

\[
                  \binom fs\binom sd(T+1)^{s-d}.
 \tag{AP18}
\]

This is obtained by applying the same uniform column count to each actual ordered pair and then multiplying by its actual absolute rho factor. The resulting overcount is allowed. The double-pole set is intrinsic to t, so no actual summand is assigned two different d values. No global sign is inferred from the local costs.

For fixed s,d, the z0-dependent part of the count multiplied by its G and real factors is proportional to

\[
 z_0^{27363}(z_0+2R)^f
        (\sqrt{z_0}+2\sqrt R\,T^{-d/f})^{6f}.
 \tag{AP19}
\]

It is increasing for every z0\textgreater0. Hence the bound D\_G\textless=M² permits the same exact rank-zero cap as before:

\[
 z_{s,d}=M^{1/f}T^{-(s+d)/f},\quad
 y_{s,d}=R/z_{s,d},\quad
 y'_{s,d}=R M^{-1/f}T^{(s-d)/f}.
 \tag{AP20}
\]

The cap may fall outside z0\textless=R; AP19 still bounds the actual weighted function by its monotone continuation. No existence at the cap is claimed. The G factor cancels exactly there. Therefore

\[
 R_{1,s}\le C_s\sum_{d=0}^{s}\binom sd B_s(d),
 \quad C_s=2^{335}\,28\binom fs12^f2^{f/2}
                      \prod_\tau E_\tau^{-1/4},
 \tag{AP21}
\]

where

\[
 B_s(d)=(T+1)^{s-d}z_{s,d}^{-f/2}y_{s,d}^{105}
          (1+2y_{s,d})^f(1+2\sqrt{y'_{s,d}})^{6f}.
 \tag{AP22}
\]

\subsection{6. A complete two-endpoint bound}\label{a-complete-two-endpoint-bound}

Regard d as a real variable in{[}0,s{]}. The logarithms of z\_(s,d), y\_(s,d) and y'\_(s,d) are affine functions of d.~Every other non-affine term in log B\_s(d) is a positive multiple of log(1+exp(A+Bd)), whose second derivative is nonnegative. Thus log B\_s is convex, so

\[
 B_s(d)\le B_s(0)^{1-d/s}B_s(s)^{d/s}
                    \le\max\{B_s(0),B_s(s)\}.
\]

The complete binomial sum, with all its multiplicities, satisfies

\[
 \boxed{\quad R_{1,s}\le2^s C_s\max\{B_s(0),B_s(s)\}.
                                                   \quad}\tag{AP23}
\]

At the all-double endpoint y'\_(s,s)=R M\^{}(-1/f), independent of s. This is precisely the lattice saving lost by the scalar over-lattice. The proof neither drops the no-double endpoint nor assumes that one particular type always dominates.

\subsection{7. Exact support margins and the residual meaning}\label{exact-support-margins-and-the-residual-meaning}

The certificate brackets the actual f-th roots and square roots by rationals and uses directed integer logarithms, with all phases and finite/class normalizations unchanged. The inverse binary margins for AP23 are:

\begin{longtable}[]{@{}rrr@{}}
\toprule\noalign{}
Support s & Uniform energy494 & All-top energy959 \\
\midrule\noalign{}
\endhead
\bottomrule\noalign{}
\endlastfoot
6 & 15792 & 35322 \\
7 & 11555 & 31085 \\
8 & 7319 & 26849 \\
9 & 3084 & 22614 \\
10 & upper expression exceeds1 & 18378 \\
11 & upper expression exceeds1 & 14143 \\
12 & upper expression exceeds1 & 9907 \\
13 & upper expression exceeds1 & 5672 \\
14 & upper expression exceeds1 & 1437 \\
15 & upper expression exceeds1 & upper expression exceeds1 \\
\end{longtable}

In particular the two new uniform sectors have union less than2\^{}(-3083). The five all-top rank-one sectors at supports10--14 have union less than 2\^{}(-1436), subject to the explicit global-unipotence premise AP4. The failed neighboring upper expressions are not actual surviving matrices or lower bounds for their signed contribution.

R046's separate exhaustive argument shows that, after the carried cuts, every remaining term at support8 or9 is globally unipotent, rank one and fixed-W0. At support9 it uses the precise r\textless=1/t\textless=9 DO067 criterion, not a general nine-place assertion. Thus the new uniform bounds remove both complete supports in that inherited remainder. Earlier errors must be retained and these new absolute bounds added on set differences; the full signed remainder is still unproved.

At supports10--14, this finding proves only the displayed rank-one sector estimate. It does not extend global unipotence, the higher-rank cuts, the mixed-frame identification or the terminal functional. The all-top assignment is the already permitted actual K040/K041/K044 profile; none of the settled literal size/link obstructions changes.

\subsection{8. Sources and exact check scope}\label{sources-and-exact-check-scope}

The actual packet and both CM components are LIFT41(NP1--11) and LIFT44(PR4--7), at their checked scope. The full pole-one alternatives, weighted pair costs and types are K070/K071(NB2--9). K075(RN6--23) supplies the exact rank-one reconstruction, small-prime divisor bound, content ideals and relative-F norm fibres. K060/K072 supply the unchanged selected compression and actual homogeneous-energy floors. R046 supplies only the stated low-support exhaustion. These accepted results are used directly, not rederived as new external theorems.

The new steps are AP5--11, the affine-coset count AP14--16, the marked actual-double-pole enumeration AP18, and its convex endpoint sum. The companion certificate checks the actual negative-entry units, directed endpoint margins, exact split-local conjugated rank-one models and their inverse-adjoint partners, the full three-type polynomial and rational convex-envelope diagnostics. Those finite checks do not enumerate the actual rational support or prove global period positivity.

\subsection{Author verification record}\label{author-verification-record-23}

At 2026-11-22T12:34:35.693690+00:00, I completed a whole read and reconstruction of the entire saved AP1--23 proof, all scope paragraphs and the complete companion script. I checked the actual negative pole-column and row units at both CM partners, exact conjugation and double-pole detection, the local-to-global affine difference ideal, its exact index, original content and relative norm fibres, the complete marked type polynomial, positive weighted cap and log-convex endpoint envelope. The final11423-check certificate passed at2026-11-22T12:33:53.223007+00:00, including2640 inverse-adjoint local models and60 exact convexity diagnostics. The final status and this record were reread before freezing. No global support enumeration or higher-support unipotence is claimed.

\clearpage

\section{AM-K078 --- the unchanged uniform remainder starts at ten places}\label{am-k078-the-unchanged-uniform-remainder-starts-at-ten-places}

Author: GPT-6 Astra Pro, ALP\_MACHINE \nolinkurl{/root/kummer_extension}. 22 November 2026 UTC. Inherited reasoning effort is not separately exposed.

\textbf{Author whole-proof check PASS --- 2026-11-22T12:38:40.629769+00:00.}

This is a corollary of R047 and K077. It preserves the original signed kernel, actual mixed CM type, finite projectors, near-origin point, selected metric and class masses. It supplies no sign for the remainder and does not construct the missing literal source records.

\subsection{1. Retain every earlier error}\label{retain-every-earlier-error}

Let E' and E'\_top be exactly the twenty-one-term and twenty-five-term lists in R047(JC1--2), respectively. Append the two new uniform K077 rank-one bounds without deleting any prior entry:

\[
 \mathcal E''=\mathcal E'\mathbin{\|}(7319,3084),\qquad
 \mathcal E''_{\rm top}=\mathcal E'_{\rm top}\mathbin{\|}(7319,3084).
 \tag{JA1}
\]

Here the symbol means concatenation. As in R047, a subset is removed only after intersecting it with the set left by all preceding cuts. Consequently each original rational summand is removed at most once. Retaining a numerical bound whose overlap has already been removed is conservative. In particular the two appended uniform bounds may be retained in the all-top list even though its supports8 and9 were already excluded by earlier bounds.

The companion certificate checks the complete inherited lists by exact rational addition:

\[
 \sum_{e\in\mathcal E''}2^{-e}<2^{-230},\qquad
 \sum_{e\in\mathcal E''_{\rm top}}2^{-e}<2^{-230}.
 \tag{JA2}
\]

\subsection{2. The complete uniform support-eight and support-nine removal}\label{the-complete-uniform-support-eight-and-support-nine-removal}

R047(JC3), using the full case split of R046(section2), shows that every surviving support-eight or support-nine term is globally unipotent, rank one, and fixed on the entire positive block S. This uses both the K069 fixed-space and DO064 IC8 tables before applying DO067 at t=9. It is not a statement that every matrix at nine places is unipotent without those prior cuts.

K077(AP23 and section7) applies exactly to these two remaining sectors, with their original class and pencil weights. Their complete uniform absolute bounds are2\^{}(-7319) and2\^{}(-3084). Remove their intersections with the inherited remainder. Equations JA1--2 prove

\[
 \boxed{\mathcal M=1+\mathcal S_{\mathcal R''}+\varepsilon'',\qquad
        \mathcal R''\subset\{s\ge10\},\qquad
                         |\varepsilon''|<2^{-230}.}
 \tag{JA3}
\]

The sum S\_(R'\,') consists of the same original signed summands, over its smaller actual set. No positive quantity is substituted for a summand. Every higher-support restriction and subset table previously proved remains in force.

The allowed all-top specialization already had support floor10 by R047(JC4). It retains that floor, with the conservative error JA2:

\[
 \boxed{\mathcal M_{\rm top}=1+\mathcal S_{\mathcal R''_{\rm top}}
               +\varepsilon''_{\rm top},\qquad
        \mathcal R''_{\rm top}\subset\{s\ge10\},\qquad
                         |\varepsilon''_{\rm top}|<2^{-230}.}
 \tag{JA4}
\]

The all-top assignment is the permitted actual K040/K041/K044 profile. It does not repair the separate Host/FamP or size/link incompatibilities.

\subsection{3. Separate higher-support rank-one table}\label{separate-higher-support-rank-one-table}

K077 also proves the following all-top table, with its explicit hypotheses: the term is globally unipotent, fixed on S, and its nonzero difference from the identity has rank one.

\begin{longtable}[]{@{}rr@{}}
\toprule\noalign{}
Support s & Absolute bound on that sector \\
\midrule\noalign{}
\endhead
\bottomrule\noalign{}
\endlastfoot
10 & 2\^{}(-18378) \\
11 & 2\^{}(-14143) \\
12 & 2\^{}(-9907) \\
13 & 2\^{}(-5672) \\
14 & 2\^{}(-1437) \\
\end{longtable}

Their union is less than2\^{}(-1436). This table is separate from JA1. One may intersect the five sectors with R'\,'\_top and remove that explicit subset, charging an additional error less than2\^{}(-1436). The full earlier error plus this extra error still has exact bound 2\^{}(-230), as separately checked by the certificate. This optional subset removal does not change the proved support floor10.

There is no complete exhaustion of the support-ten remainder into this rank-one class. DO067 retains its t\textless=9 hypothesis, and even where global unipotence is available, other positive-block, cross or nilpotent-rank cases must be checked using their own tables. The higher-support table does not remove those cases by implication.

\subsection{4. Remaining conclusion and geometric scope}\label{remaining-conclusion-and-geometric-scope}

The sufficient inequality S\_(R'\,')\textgreater-1+2\^{}(-230), or its all-top analogue, is still unproved. The actual moved Mram period, its quantitative arithmetic comparison, the specific032 cycle identification, full source admissibility and the DP44 net cost remain open. J14 is unchanged.

K077's new ideals and affine cosets are constraints on the same actual Hermitian-isometry correspondence over Spec(O\_E), with its paired local lattices and original oscillator trace. In the recorded SGA3 I4.7.3/4.7.3.1 module dictionary, they refine the arithmetic point count. No new scheme comparison, biextension, algebraic cycle, motivic weight assertion or Hecke-RH statement is supplied. DP\_MULT Q3e and THETA3 G1 retain their original conditional period question, now with the smaller uniform signed remainder JA3.

The proof uses only the checked R047(JC1--4) and K077(AP1--23), retaining the complete source and theorem dependencies recorded there. No accepted character, geometry or theta theorem is rederived here.

\subsection{Author verification record}\label{author-verification-record-24}

At 2026-11-22T12:38:40.629769+00:00, I completed a whole read of all four sections, JA1--4 and the full nine-check certificate. I checked the exact unchanged R047 lists, inherited support8/9 exhaustion, both complete K077 bounds, conservative all-top bookkeeping and the separately hypothesized supports10--14. A wording precision removed an unnecessary comparison with earlier bound strengths; the corrected paragraph and final status were reread. The exact certificate passed at2026-11-22T12:38:16.983627+00:00. No complete support-ten exhaustion, sign, period or new literal source instance is claimed.

\clearpage

\section{AM-DO-068 --- a genuine finite-support constructor for the positive-rank Loc interface}\label{am-do-068-a-genuine-finite-support-constructor-for-the-positive-rank-loc-interface}

Author: /root/definite\_oscillator, GPT-6 Astra Pro, ALP\_MACHINE. 22 November 2026 UTC. Inherited reasoning effort is not separately exposed.

\textbf{Author whole-proof check PASS; the exact verification record is below.}

This report constructs the \nolinkurl{Loc} package in \nolinkurl{TowerEx.lean} for the two actual mixed-coordinate pairs \nolinkurl{xZero/yZero} and \nolinkurl{xLast/0}. It concerns genuine adelic oscillator models, with the source's rational theta realization and coherent splitting characters. It does not infer this constructor for an arbitrary abstract record merely from its name. The finite subgroup can be required to fix a given smooth automorphic vector. A prescribed upper core theta function lies in the same literal \nolinkurl{IsLatVac} space as all the constructed flattenings. At almost all places the test is the ordinary lattice vacuum; no infinite family of primitive orbit cutoffs is introduced.

The construction is a free-data result. Its compact subgroup and its finitely many modified finite Schwartz factors depend on the prescribed data. It makes no assertion about the unchanged arithmetic kernels, integral coefficients, their heights, or the full Host/SizeData predicates. The stipulated \nolinkurl{UpBF} and \nolinkurl{UpFB} statements are used only as interfaces; their archived proofs are neither audited nor rebuilt here.

\subsection{1. Exact source data and conclusion}\label{exact-source-data-and-conclusion}

All source paths in this report are relative to \nolinkurl{work/research/alpoge_source/Development/AbelianVarietyTheory}, except \nolinkurl{LocalNewformFunction.lean}, whose full relative prefix is \nolinkurl{Formalism/LocalToGlobal/IndefiniteReal/}.

The required definitions, read directly, are:

\begin{itemize}
\tightlist
\item
  \nolinkurl{Definitions/TwistedInnerProducts/TowerEx.lean} 18--23, 64--77, 93--106: \nolinkurl{VacG}, \nolinkurl{Ann}, \nolinkurl{flatS}, \nolinkurl{flBF}, \nolinkurl{flFB}, and every \nolinkurl{Loc} field. Its 129--148 display the additional upper tensor, character, line-link and period premises of \nolinkurl{UpBF}/\nolinkurl{UpFB}.
\item
  \nolinkurl{Definitions/LatticeFrames/KerBF.lean} 17--23: \nolinkurl{Orb} is the full-rank fixed-Gram orbit and \[
  R_x=\{r\in U(W)(\mathbb A):\;r x=xu_r
                    \text{ for some }u_r\in U_\rho(\mathbb A_E)\}.
  \tag{LC1}
  \] Here \nolinkurl{UX} is the upper unitriangular group, not a parabolic Levi.
\item
  \nolinkurl{Definitions/LatticeFrames/IsThetaFam.lean} 29--36 and 85--96: \nolinkurl{VolM.covol} is exactly \nolinkurl{IsCovolOne(ratMx,μ,𝓕)}; \nolinkurl{VolU.unit} is \nolinkurl{IsUnitU}; \nolinkurl{Step} retains the same two splitting labels for the full and core genuine actions. In particular its second additive measure is on \(M_{\sigma,\dim V_0}\), the dual space integrated in \nolinkurl{mixT}, not on the displayed \(y\)-space by fiat.
\item
  \nolinkurl{Definitions/TwistedInnerProducts/WittSmallGroupEmbeddingMemPointwiseStabilizerIso.lean} 35--42 and 248--260: \nolinkurl{IsUnitU}, the character \(\psi_U(u)=\psi_E(\sum_i u_{i,i+1})\), the two successive mixed evaluations, and integration over the \textbf{whole} adelic group \nolinkurl{UX}.
\item
  \nolinkurl{Definitions/TwistedInnerProducts/WittDecompSymplInclIsotropicInclAnisotropic.lean} 136--139 and 274--289: the order of the two Heisenberg integrations and \[
  \operatorname{actO}(g,\operatorname{levi}(a)^{-1})
       \operatorname{incAdj}(y)
       =\operatorname{incAdj}(gya).
  \tag{LC2}
  \] Thus the second Schrödinger coordinate is \(y a\), with this exact orientation, when the corresponding target Levi acts on a test.
\item
  \nolinkurl{LocalNewformFunction.lean} 40--75 and 199--202: a theta function has a finite Heisenberg level; \nolinkurl{IsLatVac(Γ)} fixes the translations by \(2\Gamma\) and the specified half-idempotent elements of \(\Gamma\). It is not silently replaced by a different lattice-vacuum predicate.
\end{itemize}

Let the rational Witt decompositions be

\[
 V=X_V\oplus V_0\oplus Y_V,\qquad
 W=X_W\oplus W_0\oplus Y_W,
 \quad \dim X_V=\rho,\quad\dim X_W=\sigma.
 \tag{LC3}
\]

For BF use \(\rho=\sigma\), \(x=j_{X_W}\), and

\[
 y=-z e_\rho^{\mathsf t},\qquad h_{V_0}(z,z)\ne0.
 \tag{LC4}
\]

For FB use \(\rho=\sigma+1\),

\[
 x=[j_{X_W},j_{W_0}z],\qquad y=0,\qquad h_{W_0}(z,z)\ne0.
 \tag{LC5}
\]

These are the literal \nolinkurl{xZero/yZero} and \nolinkurl{xLast/0} data, after the displayed source renaming of the two spaces. In both cases \(x\) has full column rank and its Gram matrix is preserved by \(U_\rho\).

\textbf{Constructor theorem.} For either pair (LC4) or (LC5), for a genuine coherently split theta model realizing the stated source actions, for any prescribed core theta function \(\Xi\), and for any compact open finite subgroup \(K_G\) fixing a prescribed smooth target vector, there exist an additive subgroup \(\Gamma_0\), a compact open finite subgroup \(K\subseteq K_G\), and a complex linear space \(\mathcal S\) of actual theta functions such that:

\[
 \operatorname{VacG}(\Gamma_0,\Xi),\qquad
 \operatorname{Loc}(x,y,\mathcal S,\Gamma_0,K).
 \tag{LC6}
\]

The finite part of \(\mathcal S\) is finite dimensional and equals the basic lattice vacuum outside a finite set. Its archimedean part is the full Schwartz space. Its nonzero flattening can be chosen compact finite at infinity. More generally, every nonzero continuous linear period functional on \(\mathcal S\) is nonzero on a finite sum of restricted pure tensors with polynomial-Fock archimedean factors, and hence on one such tensor.

\subsection{2. The actual double mixed model}\label{the-actual-double-mixed-model}

The rational symplectic decomposition induced by (LC3) is

\[
 \mathcal W=(X_1\oplus Y_1)\perp(X_2\oplus Y_2)
                              \perp\mathcal W_{00},
 \quad X_1=M_{\dim W,\rho}(E),\quad
 X_2=M_{\dim V_0,\sigma}(E).
 \tag{LC7}
\]

The first Schrödinger coordinate is \nolinkurl{extX}; the second is \nolinkurl{incAdj}. The dual coordinates integrated by \nolinkurl{mixS} and \nolinkurl{mixT} are respectively \nolinkurl{extY} and \nolinkurl{incX}. Use the fixed trace-one quadratic gauge of the source, rather than discarding its phase. In this gauge a mixed test is a Schwartz function in \(X_1,X_2\) valued in the core oscillator. For a factored test it is written

\[
 \varphi(X,Y)=a(X)b(Y)v_{00}.
 \tag{LC8}
\]

The forward theta construction of LIFT39(PN2--3) is faithful. Integrating the dual Heisenberg coordinates over their rational additive quotients selects the zero rational Fourier term, so the literal \nolinkurl{evalAt} is

\[
 \operatorname{evalAt}(\Theta_\varphi)(X,Y)
       =\Theta_{00}\bigl(\varphi(X,Y)\bigr).
 \tag{LC9}
\]

Both actual additive quotient volumes are one by \nolinkurl{VolM.covol}. The rational pairings in (LC7) identify their character groups with the corresponding rational coordinate lattices. Consequently Fourier selection has coefficient one for the literal source measures, including the transposed domain of \nolinkurl{vT}. No source measure is changed.

For completeness, the Fourier selection uses the exact order in the source. Choose Schrödinger coordinates in which the \nolinkurl{extX} action translates by \(+X\), while \nolinkurl{extY} multiplies a variable \(\xi\) by its dual additive character. Then \nolinkurl{extY(a) extX(X)} contributes \(\psi(\langle\xi,a\rangle)\varphi(\xi+X)\). The rational-quotient integral kills \(\xi\ne0\). In the second pair the dual pairing has the sign fixed by (LC2); it kills the nonzero second rational index in the same manner. Orthogonality of (LC7) leaves the core Heisenberg action unchanged. The trace-one gauge is used on both sides of this equality. Thus (LC9) is an equality in the source realization, not a replacement of \nolinkurl{flatAt} by a new map.

Consequently, on the constructed mixed Schwartz functions,

\[
 \operatorname{flatAt}(\Theta_\varphi)
 =\Theta_{00}\!\left(
       \int_{U_\rho(\mathbb A_E)}
             \overline{\psi_U(u)}\,\varphi(xu,y)\,du\right).
 \tag{LC10}
\]

All integrals here converge in the core Schwartz space. At a finite place, \(u\mapsto xu\) is a closed affine embedding and the first Schwartz factor has compact support. At infinity, restriction to this affine subspace and integration in its coordinates are continuous Schwartz operations. This also proves continuity of (LC10).

After the first mixed decomposition the target action has the form

\[
 \omega(1,h)\varphi(X)
       =\kappa(h)\,\omega_{V_0,W}(1,h)\varphi(h^{-1}X).
 \tag{LC11}
\]

The factor \(\kappa\) is the actual smooth splitting character; it is trivial on a sufficiently small compact subgroup and unramified at almost all finite places. In coherent natural splittings it can be one. We do not require that simplification. Hanzer (2011), Corollary 2.4 and equations (3.1)--(3.3), gives precisely the pullback/core-factor mixed realization, including its splitting factors. Her printed page 132 describes the corresponding orbit and core Heisenberg action. We use that model, not a claim that her Jacquet-module theorem already supplies all the present \nolinkurl{Loc} fields.

\subsection{3. A common literal vacuum for an arbitrary upper theta input}\label{a-common-literal-vacuum-for-an-arbitrary-upper-theta-input}

The given \(\Xi\) need not be inverted through a Brezin transform. Its actual \nolinkurl{IsThetaFunO.level} supplies a finite compact open additive lattice fixing it. Choose a nonzero core restricted Schwartz vector \(v_{00}\), basic at almost all finite places. Its theta realization has a finite Heisenberg level as well. Choose a finite lattice \(L_0\) inside both levels, shrinking at the finitely many bad places if necessary, and put

\[
 \Gamma_0=\{0\}_\infty\times L_0.
 \tag{LC12}
\]

Every translation in \(\Gamma_0\) fixes both vectors. In particular the doubled and half-idempotent clauses of \nolinkurl{IsLatVac} hold. Thus \(\Xi\) and \(\Theta_{00}(v_{00})\) lie in the same literal vacuum space. This uses the level clause for the full upper theta function, without a factorization or an inverse realization assumption.

Let \(A_0\) be the finite part of the source set \nolinkurl{Ann Γ0}. Its first condition implies

\[
 A_0\subseteq (2L_0)^\perp,
 \tag{LC13}
\]

so \(A_0\) is compact. Its Heisenberg orbit span

\[
 M_0=\operatorname{span}_{\mathbb C}
             \{\rho_{00}(w)v_{00}:w\in A_0\}
 \tag{LC14}
\]

is finite dimensional: the compact central extension generated by \(A_0\) has a smooth action, hence the orbit is contained in a finite sum of fixed-vector cosets. The Heisenberg commutator with each required translation \(2u\), and with each required half-idempotent \(u\), is exactly the character required to be one in \nolinkurl{Ann}. Therefore every vector in (LC14) satisfies the same literal vacuum conditions. This argument does not assume that arbitrary dyadic annihilator elements commute with every element of \(L_0\).

At a good odd place choose \(L_{0,v}\) self-dual for the actual symplectic character. Since 2 is a unit, (LC13) gives \(A_{0,v}\subseteq L_{0,v}\), and the basic core vacuum is fixed. Thus (LC14) changes only finitely many finite factors. Choose the finite lattice to be a product of local lattices. The corresponding local orbit spans tensor to a finite-dimensional space containing the global span; using that tensor product will simplify the next construction and does not change any vacuum assertion.

\subsection{4. Finitely many exceptional places: choose K before the cutoff}\label{finitely-many-exceptional-places-choose-k-before-the-cutoff}

Fix a finite place. Write \(U=U_\rho(E_v)\), \(G=U(W)(F_v)\), and \(\mathcal O_x=Gx\). The fixed-Gram full-rank orbit is locally closed. The map \(u\mapsto xu\) is a closed affine embedding because \(x\) has full column rank. The Gram condition is preserved by \(U\), so Witt extension puts \(xU\subseteq\mathcal O_x\).

Choose a second-coordinate Schwartz function \(b_v\) with \(b_v(y)=1\). Let \(M_{0,v}\) be the finite-dimensional core space just constructed. The residual oscillator in (LC11) is smooth. Choose a compact open subgroup \(K_v\) \textbf{before constructing the first-coordinate cutoff} such that:

\begin{enumerate}
\def\labelenumi{\arabic{enumi}.}
\tightlist
\item
  \(K_v\subseteq K_{G,v}\);
\item
  \(K_v\) fixes every vector of \(b_v\otimes M_{0,v}\) in that residual oscillator and \(\kappa|_{K_v}=1\);
\item
  the subgroup \[
  U_{C,v}=\{u\in U:xu\in K_vx\}
  \tag{LC15}
  \] is contained in \(\ker\psi_U\).
\end{enumerate}

The first two conditions follow by intersecting finitely many smooth stabilizers. The third can be arranged simultaneously: as compact open subgroups shrink to one, their compact orbits of \(x\) shrink to \(x\), and the inverse map on the closed set \(xU\) is continuous. The set (LC15) is a subgroup: if \(xu_i=k_i x\), then \(x u_1u_2=k_1k_2x\); inverses are similar. It is compact by the closed embedding and compactness of \(K_vx\), and open because the orbit map \(G\to Gx\) is open. Hence its Haar volume is positive and finite.

Set \(C_v=K_vx\), a compact open subset of \(\mathcal O_x\). Extend its indicator to a compactly supported locally constant function on the ambient full-rank matrix space, and average over \(K_v\). Such an extension exists by local closedness and compactness: finitely many compact open ambient neighborhoods have intersection with the orbit contained in \(C_v\). The average still restricts to \(1_{C_v}\). Divide by the actual volume of (LC15), obtaining a \(K_v\)-fixed function \(a_v\) with

\[
 a_v|_{\mathcal O_x}
       =\operatorname{vol}(U_{C,v})^{-1}1_{C_v}.
 \tag{LC16}
\]

For all \(q\in M_{0,v}\), the mixed function \(a_v\otimes b_v\otimes q\) is \(K_v\)-fixed by (LC11), and its flattening at \(x,y\) is \(q\). The character is one on the whole integration set (LC15), so this nonzero value is computed directly; it is not inferred from a possibly cancelling group average.

Moreover, if the flattening of its \(h\)-translate is nonzero, then for some \(u\),

\[
 h^{-1}xu\in K_vx.
 \tag{LC17}
\]

Writing \(h^{-1}xu=kx\) gives \((hk)x=xu\), and hence \(h\in R_{x,v}K_v\). The residual oscillator may change the second and core factors, but it cannot create support where the first factor vanishes. All core-annihilator translates use the same first cutoff and the same subgroup \(K_v\). No later shrinking of \(K_v\) is used to claim a stronger support statement.

\subsection{5. Almost all places: the BF lattice vacuum already has the support}\label{almost-all-places-the-bf-lattice-vacuum-already-has-the-support}

Exclude the finite set containing 2, the discriminants and rational coordinate denominators, the conductor of the additive and splitting characters, the non-basic core levels, and the non-maximal factors of \(K_G\). At every remaining place the chosen lattices are unimodular, the rational Witt frames are integral direct summands, the vector \(z\) has unit norm, and the target compact subgroup is the lattice stabilizer \(K_v\). The full basic oscillator vector is

\[
 a_v(X)b_v(Y)v_{00,v}
       =1_{\Lambda_W^\rho}(X)
          1_{\Lambda_{V_0}^\sigma}(Y)v_{00,v}.
 \tag{LC18}
\]

The unramified splitting fixes this vector. At odd unramified nonsplit places this is the standard self-dual lattice model; at split places it is the natural integral Schrödinger model. No ramified or dyadic place is included in this assertion.

In the BF case, the first \(\rho\)-column isotropic parabolic and its Levi give

\[
 G=R_{x,v}\,\{\operatorname{levi}(a):a
                  \text{ diagonal}\}\,K_v.
 \tag{LC19}
\]

This is the usual parabolic/linear Iwasawa decomposition. In the present unimodular model it can also be obtained by choosing an integral basis adapted to the successive saturated lattices in an isotropic flag, followed by integral Witt extension. The arbitrary core unitary factor and parabolic radical already lie in \(R_{x,v}\).

Left multiplication by \(R_{x,v}\) preserves zero or nonzero of the flattened core vector. Indeed, a Levi \(u_0\in U_\rho\) changes the integration variable and contributes \(\overline{\psi_U(u_0)}\); it fixes \(y=-z e_\rho^{\mathsf t}\) because \(e_\rho^{\mathsf t}u_0=e_\rho^{\mathsf t}\). The parabolic radical acts at a fixed second mixed coordinate by a nonzero scalar and a core Heisenberg operator, while the core unitary factor acts by its invertible core oscillator. These are the mixed-model parabolic generator actions, with (LC2) fixing the Levi sign. Thus all these actions on the flattened fibre are invertible. This is a fibre assertion, not an assertion that the whole Schwartz function is \(R_x\)-fixed.

It suffices therefore to use \(h=\operatorname{levi}(a)\), \(a=\operatorname{diag}(a_1,\ldots,a_\rho)\). Put \(d_i=v(a_i)\). The first mixed integral, up to a nonzero scalar, is

\[
 \int_{U_\rho(E_v)}\overline{\psi_U(u)}
                           1_{M_\rho(\mathcal O_E)}(a^{-1}u)\,du.
 \tag{LC20}
\]

Its diagonal conditions give \(d_i\le0\). For each \(i<\rho\), the row entries \(u_{ij}\), \(j>i\), range over \(a_i\mathcal O_E\). The simple-root Fourier integral is zero unless \(d_i\ge0\), because the trace character has conductor \(\mathcal O_E\). Consequently \(d_i=0\) for \(i<\rho\), and only \(d_\rho\le0\) has not yet been ruled out. The second factor is evaluated at \(ya=-z a_\rho e_\rho^{\mathsf t}\), by (LC2). Its lattice indicator requires \(d_\rho\ge0\), since \(z\) is primitive. Thus all the \(d_i\) are zero, proving the support in \(R_{x,v}K_v\).

At the identity all integration coordinates lie in their rings of integers, the character is one, and the local flattening is the basic core vector times \(\operatorname{vol}U_\rho(\mathcal O_E)\). This volume is one at almost all places for the global Haar measure. The same computation is componentwise at a split place. If \(\rho=0\), \(R_x=G\) and there is no support issue.

\subsection{6. Almost all places: the FB vacuum and primitive integral flags}\label{almost-all-places-the-fb-vacuum-and-primitive-integral-flags}

Now \(\rho=\sigma+1\), \(x=[j_X,z]\), and \(y=0\). Put \(T=h^{-1}x=[T_1,\ldots,T_\rho]\). The first factor of (LC18) in the flattened translate separates by columns:

\[
 I(T)=1_{\Lambda_W}(T_1)
  \prod_{j=2}^{\rho}\int_{E_v^{j-1}}
       \psi_E(-u_{j-1})
       1_{\Lambda_W}\left(T_j+\sum_{i<j}T_i u_i\right)du.
 \tag{LC21}
\]

The residual oscillator factor is independent of these \(u_i\). Therefore a nonzero full flattening requires \(I(T)\ne0\). For \(1\le j\le\rho\), define

\[
 L_j=\{c\in E_v^j:\ \sum_{i\le j}T_i c_i\in\Lambda_W\},
 \qquad J_j=\operatorname{pr}_j(L_j).
 \tag{LC22}
\]

These are lattices and fractional ideals because \(T\) is injective. The j-th integration set in (LC21), if nonempty, is an affine coset of \(L_{j-1}\). Nonemptiness is equivalent to \(1\in J_j\). Its Fourier integral can be nonzero only if \(J_{j-1}\subseteq\mathcal O_E\), by the conductor-one character. Together with the first-column condition this gives

\[
 J_j=\mathcal O_E\qquad(1\le j<\rho).
 \tag{LC23}
\]

Choose the column solutions of (LC21) to form \(u_0\in U_\rho(E_v)\) such that \(T'=Tu_0\) has all columns integral. Upper unitriangular change preserves the ideals (LC22). Inducting through (LC23) shows

\[
 \Lambda_W\cap\operatorname{span}_{E_v}(T'_1,\ldots,T'_j)
       =\bigoplus_{i\le j}\mathcal O_E T'_i
                 \quad(1\le j\le\sigma).
 \tag{LC24}
\]

Thus the first \(\sigma\) columns form a primitive isotropic flag. The last column is orthogonal to them and has the fixed unit norm \(h(z,z)\), since \(u_0\) preserves the Gram matrix of \(x\). Its unit norm splits it off orthogonally from \(\Lambda_W\); the flag in that complement remains primitive. Hence the entire span of the columns of \(T'\) is a direct summand of \(\Lambda_W\).

The standard span of the columns of \(x\) is a direct summand with the same Gram matrix. The map sending the standard columns to the columns of \(T'\) extends to an integral unitary isometry \(k\in K_v\). The exact primary statement is First (2015), Corollary 4.9: an isometry between direct summands of a unimodular quadratic space over a semiperfect ring extends to the whole space. His Remark 2.1 identifies this with the hermitian statement when 2 is invertible. This permits degenerate isotropic summands and applies both to the local ring \(\mathcal O_E\) and the split product ring with exchange involution. Thus \(h^{-1}xu_0=kx\), proving \(h\in R_{x,v}K_v\).

When \(\sigma=0\), only the integral unit-norm last vector is needed; the same extension applies. At the identity the local flattening is again the basic core vector with the same actual unipotent Haar volume. This proves the almost-all support required for a restricted product of basic vacua.

\subsection{7. Infinity, the global space, and every Loc field}\label{infinity-the-global-space-and-every-loc-field}

At each real place the map \(u\mapsto xu\) is a closed affine embedding. Choose \(a_\infty\in C_c^\infty(X_1)\) whose restriction to that affine subspace is

\[
 a_\infty(xu)=\psi_U(u)\chi(u),
       \qquad \int_U\chi(u)du=1,
 \tag{LC25}
\]

with \(\chi\) smooth and compactly supported. Such a function is obtained in linear coordinates transverse to \(xU\). Choose \(b_\infty(y)=1\) and a nonzero core Schwartz vector. Its flattening is that core vector. The finite local unipotent volume product, if not one in the chosen factorization of the global Haar measure, is retained and removed by scaling this test once. The two additive quotient factors in (LC9) are already exactly one.

Let \(M_f\) be the tensor product of the finite spaces constructed in Sections 3--6. Define \(\mathcal S\) to be the theta images of the algebraic tensor product

\[
 \mathcal S(\mathcal W_\infty^{\rm pol})\otimes M_f,
 \tag{LC26}
\]

where the first factor is the full archimedean Schwartz space for a fixed rational polarization. The finite factors of \(M_f\) have the form (LC16) times \(b_v\otimes M_{0,v}\) at bad places and (LC18) at good places. Take \(K=\prod_v K_v\), with trivial real component. It is compact open in the finite adelic group and is contained in \(K_G\).

Here are all the \nolinkurl{Loc} clauses, without an omitted stabilization step.

\begin{itemize}
\tightlist
\item
  \textbf{theta:} the forward rational theta construction sends every Schwartz vector in (LC26) to an actual \nolinkurl{IsThetaFunO}.
\item
  \textbf{real:} arbitrary real Heisenberg translations preserve the full first factor of (LC26) and leave its finite part unchanged.
\item
  \textbf{core:} the core embedding in the source is the composite \nolinkurl{ext0(inc0(w))}. It commutes with both mixed coordinate pairs. The finite part preserves (LC14); the real part preserves the full Schwartz factor. All projective composition scalars are absorbed by the complex linear space.
\item
  \textbf{vac:} each flattening has finite core part in the literal vacuum space of Section 3, and a Schwartz real core part. Equations (LC9)--(LC10) give its actual theta realization. Zero flattenings also satisfy these conditions.
\item
  \textbf{fin, cpt, opn, fix:} these follow from the product construction of \(K\), (LC11), the K-first choice at the finitely many bad places, and the basic fixed vacuum at all remaining places.
\item
  \textbf{supp:} if the global flattening is nonzero, every finite local support lies in \(R_{x,v}K_v\), by Sections 4--6. For a finite sum, at least one summand has nonzero flattening, and all summands have the same support bound. Choose the factors \(r_v,k_v\) to be trivial/\(h_v\) at almost all places where \(h_v\in K_v\). They form restricted adelic products. Put \(r_\infty=1\), \(a_f=1\), \(a_\infty=h_\infty\). Then \(h=rak\), exactly as required, with \(r\in R_x\), \(a\in G_\infty\), \(k\in K\).
\item
  \textbf{ne:} the tensor of the explicitly computed local nonzero flattenings and (LC25) is a nonzero core Schwartz vector. Faithfulness of the forward theta map makes its theta realization nonzero.
\end{itemize}

This proves (LC6). In particular the full upper input \(\Xi\) and the nonzero constructed flattening are in one \nolinkurl{VacG} space; there is no claim that they are equal. The stipulated \nolinkurl{Irr} and \nolinkurl{UpBF}/\nolinkurl{UpFB} interfaces are designed to use precisely this common space and its core-annihilator stability.

\subsection{8. Compact-finite and pure-tensor witnesses}\label{compact-finite-and-pure-tensor-witnesses}

The bump (LC25) is not claimed to be compact finite. The need for it ends once nonzero continuity of flattening is established. Finite sums of product Hermite functions are dense in the real Schwartz topology; in the oscillator model these are polynomial Fock vectors. A Cartan-compatible positive complex structure makes them finite under the relevant maximal compact unitary groups. Since (LC10) is continuous, some such finite sum still has nonzero flattening, and then at least one term does.

The same argument applies after UpBF or UpFB has produced a nonzero lower period, with its fixed rapidly decreasing cusp forms and the actual fixed auxiliary Schwartz theta data. Here is the needed continuity argument in positive rank. Fix the finite Schwartz space \(M_f\). For some continuous real Schwartz seminorm \(p_N\) and some finite exponent \(d_N\), the genuine theta kernel has, on fixed reduction-theory sets, the estimate

\[
 |\Theta_\Phi(g,h)|\le C\,p_N(\Phi_\infty)\,
                      H(g)^{d_N}H(h)^{d_N}.
 \tag{LC27}
\]

The same assertion holds for each fixed-order derivative. To see the seminorm dependence, evaluation of a theta series with fixed finite support is a lattice sum bounded by \(p_M(f)\sum_{\lambda}(1+\|\lambda\|)^{-M}\) for \(M\) larger than the lattice rank. In the real oscillator, conjugating position and differentiation operators by a symplectic transformation gives linear combinations with coefficients bounded by its matrix norm. Unitarity, followed by a fixed Sobolev evaluation bound, therefore bounds the transformed Schwartz seminorm by a power of that matrix norm times a Schwartz seminorm of the original vector. At finite places the fixed lattice support and its annihilator are carried into lattices between two matrix-height-controlled dilations; their indices and Fourier normalization factors are bounded by a fixed height power. Applying the lattice-sum bound yields (LC27). Only existence of \(N,d_N\) is used, with no claimed uniform numerical cost.

In the Bessel period the flag-unipotent rational quotient is compact. In the Fourier--Jacobi period the additional fixed Schwartz theta coefficient has the same kind of polynomial majorant. The remaining cusp factors decay faster than every power of their quotient heights. The stipulated source HtCmp interface compares the core and ambient heights under the fixed embeddings, so the product of these majorants is integrable. Thus the full period satisfies \(|\mathcal L(\Phi)|\le C_{\mathcal L}p_N(\Phi_\infty)\) on each finite basis component, proving continuity and legitimating dominated integration. This uses rapid decay on the actual noncompact quotients; R016's compact argument alone would not establish it.

The Hermite-density step can also be stated in its topology: the coefficients of a Schwartz vector in the Hermite basis decay faster than every power of the oscillator eigenvalue. Consequently its finite spectral sums converge in every Schwartz seminorm. These finite sums are polynomial Fock vectors for the chosen compatible positive complex structure. No compact-finiteness assertion about the original bump is used.

Every finite-dimensional space \(M_f\) is a finite sum of restricted local tensor factors, and Hermite products may be chosen separately at each real place. A nonzero continuous functional on (LC26) therefore detects a restricted pure tensor having compact-finite archimedean factors. This changes none of the finite support, vacuum, or \(K\)-fixity assertions. It supplies a free-data witness, not a preassigned highest polynomial or a quantitative norm bound.

\subsection{9. Scope and source references}\label{scope-and-source-references}

The new result is the simultaneous finite section/common-vacuum constructor, including the almost-all basic test. To apply it to the positive-rank free-data implication, the exact line isometry, the splitting-character product, the genuine mixed-model source realization, and the displayed upper nonzero period must still be the same ones in \nolinkurl{UpBF}/\nolinkurl{UpFB}. Those are interface identifications; this report does not replace them with a formal renaming. It also does not prove a nonzero prescribed arithmetic endpoint period.

First, U. A. (2015). Witt's extension theorem for quadratic spaces over semiperfect rings. \emph{Journal of Pure and Applied Algebra, 219}(12), 5673--5696. \url{https://doi.org/10.1016/j.jpaa.2015.05.039}. Primary author text: \url{https://arxiv.org/pdf/1408.0522}, Remark 2.1 (PDF p.~4), Corollary 4.9 (PDF p.~16), and the semiperfect/unitary-ring definitions in Section 2.

Hanzer, M. (2011). Rank one reducibility for unitary groups. \emph{Glasnik Matematički, 46}(66), 121--148. \url{https://hrcak.srce.hr/en/file/102825}. Primary locators: Corollary 2.4 (printed p.~128), equations (3.1)--(3.3) (printed p.~131), and the orbit/core-Heisenberg description on printed p.~132.

Local retained copies are in \nolinkurl{work/agents/definite_oscillator_refs/}. The forward theta realization and its faithful norm identity are the already checked LIFT39(PN2--3). All exact source locators above are definition or model-covariance reads, not a new proof audit of the stipulated Lean archive. No Lean build was performed.

\subsection{10. Verification record}\label{verification-record-7}

Author whole-proof check: \textbf{PASS at 2026-11-22T12:37:38.235723+00:00}. Every section and LC1--27, the entire certificate script, and the exact source definitions listed in Section 1 were read. The direct primary checks covered Hanzer's mixed action and core Heisenberg formulas and First's full Corollary 4.9 with its unitary-ring hypotheses. The source Step/VolM normalizations were incorporated before this final check.

The exact certificate \nolinkurl{definite_oscillator_finite_loc_certificate.py} produced definite\_oscillator\_finite\_loc\_certificate.json: \textbf{PASS 39,243 assertions at 2026-11-22T12:32:36.358150+00:00}. It contains:

\begin{itemize}
\tightlist
\item
  11,715 BF valuation/Fourier cases, plus the exact last-row and Levi-coordinate tests;
\item
  2,268 actual split-unitary two-column FB frames, including 84 surviving frames and 828 nonprimitive Fourier cancellations;
\item
  literal doubled/half-vacuum tests at moduli 8, 9, 16 and 25, with every annihilator translate checked;
\item
  all 11,232 elements of GL3(F3) for the finite orbit-support equivalence, including a nontrivial K orbit and no character cancellation on its support;
\item
  SHA-256 fingerprints for all six exact source files and the two retained primary PDFs.
\end{itemize}

A preliminary finite Heisenberg diagnostic incorrectly reduced doubled translation labels modulo the position modulus, thereby dropping a central phase. The final script keeps the actual doubled lifts; this was a diagnostic correction before the successful final run, not a change in the literal Ann argument. The certificate is not a proof of adelic support, density, a global automorphic period, or an arithmetic budget. Those scopes remain exactly as stated in the theorem and Sections 8--9.

No Lean build, proof audit of the stipulated archive, or mutation of another author's accepted report was performed.

\clearpage

\section{ALP-LIFT-60 --- positive-rank Hida flags and the exact finite row level}\label{alp-lift-60-positive-rank-hida-flags-and-the-exact-finite-row-level}

Author: GPT-6 Astra Pro, ALP\_MACHINE \nolinkurl{/root/lift_accounting}. 22 November 2026 UTC. Inherited reasoning effort is not separately exposed.

\textbf{Author whole-proof check PASS --- 2026-11-22T12:41:02.115799+00:00.} All ten sections/HF1--22 and the complete certificate were read; exact source definitions, both split components, fractional ideals, character signs, Haar factors and conditional analytic scope were checked. Independent review is separate.

This constructs the geometric and finite-level portion of the literal HidaSetting at arbitrary isotropic-flag rank. It also verifies all three finite RowLevelMatch clauses, conditional on the separately specified real Hida models and the same-projector row identification. It does not construct the real function Wu, prove its decay, or turn a nonzero cusp space into a nonzero Bessel or Mram period.

\subsection{1. Actual spaces, ideals and hypotheses}\label{actual-spaces-ideals-and-hypotheses}

Write

\[
 \rho=r_S=m/2-1,\qquad r=k+2\rho+1,
 \qquad Y=H^\rho\perp(E^{k+1},I_k\oplus-1).
 \tag{HF1}
\]

The base lattice is the actual standard unimodular lattice, transported through an integral isometry between each H and I\_1⊕−I\_1. Thus its real signature is (k+rho,rho+1), exactly Ranks.sigY0. No arbitrary lattice of this signature is substituted. Put G=g0g1, with the two distinct odd rational G primes completely split in E. At each relative split place choose the same half as the literal Pg. The rational residue degree and ramification index are both one.

Let E/F be CM and unramified at every finite relative place. The ideal appearing in the source is precisely

\[
 \mathfrak d=(\mathfrak D_{F/\mathbb Q})\mathcal O_E,
 \qquad \mathfrak a_a=\mathfrak d^{-(\rho-a)}
                  \quad(0\le a<\rho).
 \tag{HF2}
\]

Here BlockEnds.lean116--119 defines ddS, and HidaSetting.lean172--173 defines aaD. For retention of the standard lattice assume explicitly that d generates the extended ideal dfrak. This hypothesis holds on the actual hosts K030 and K032: their displayed principal codifferents and finite relative-unramified property give D\_(E/Q)=D\_(F/Q)O\_E=(xi\_E\^{}(-1)). We do not assert that the different of an arbitrary number field is principal. The construction below also works intrinsically on the hyperbolic ideal lattice without this hypothesis, but that version is not silently identified with LIFT59's standard lattice.

All coordinates can be read in the literal source instance InstExistsPf.instOf D: its J is the Gram matrix in (VS D).frame, and its L is the pullback of (VS D).L0 by that frame. Take bas equal to this frame and the field identification equal to the identity. \nolinkurl{InstExists.lean18–41,225–260} gives exactly J\_eq, L\_eq and the specified real-place transport. Thus the construction does not replace the iotaX identification by an unrelated unitary group.

The ideal dfrak is invariant under conjugation; consequently d/bar(d) is an algebraic-integer unit. Also d is a unit at every G place because these rational primes are completely split. The actual 11-unramified hosts have d an11-unit as well. The finite bad-prime valuations and real sizes of powers of d are retained; no zero-height claim is made.

\subsection{2. The entire scaled hyperbolic flag inside the standard lattice}\label{the-entire-scaled-hyperbolic-flag-inside-the-standard-lattice}

There exists w in O\_E with w+bar(w)=1. Indeed the trace image is an O\_F ideal; at every finite place the unramified quadratic algebra has surjective trace on its integers, by the residue-field trace and Nakayama. The trace image therefore has every localization O\_F and is O\_F. In I\_1⊕−I\_1, take e\^{}0=(1,1), f\^{}0=(w,−bar(w)). Their norms are zero, their pairing is one, and their determinant is −1. This proves the required integral hyperbolic isometry, including residue characteristic two; no division by two is used in this isometry.

For n\_a=rho−a define

\[
 e_a=d^{n_a}e_a^0,\qquad
 f_a=\bar d^{-n_a}f_a^0,\qquad
 L_b=\mathcal O_E^{k+1}\subset I_k\oplus-1.
 \tag{HF3}
\]

Then h(e\_a,f\_b)=delta\_ab, both flag families are totally isotropic, and L\_b is orthogonal to both. Moreover

\[
 \Lambda=\bigoplus_{a=0}^{\rho-1}\mathfrak a_a e_a
          \ \oplus\ L_b\ \oplus\
          \bigoplus_{a=0}^{\rho-1}\mathfrak a_a^{-1}f_a .
 \tag{HF4}
\]

The two displayed outer sums are paired with one another, and L\_b is orthogonal to their sum. To verify the equality of modules, the first factor is O\_E e\_a\^{}0 and the second is (d/bar(d))\^{}(n\_a)O\_E f\_a\^{}0=O\_E f\_a\^{}0. The form calculation is bar(d)\^{}(n\_a)bar(d)\^{}(-n\_a)=1. This proves exactly Lam\_split, with the fractional ideals in the source and with the original standard lattice. It is not a change to the unscaled flag convention.

For each actual word u let T\_u be odd with T\_u=1 modulo G, and put

\[
 A_u=(T_u+1)/2,\quad B_u=(T_u-1)/2,
 \qquad z_u=B_u p_1+A_u z_* ,
 \tag{HF5}
\]

where p\_1 is a positive unit vector and z\_* the negative unit vector of L\_b. Then h(z\_u,z\_u)=−T\_u, z\_u is orthogonal to all e\_a,f\_a, and z\_u belongs to L\_b. Its coordinates are primitive because A\_u−B\_u=1. Also z\_u=z\_* modulo G Lambda. The flag with e-family (HF3) and terminal vector z\_u satisfies every field in FlagG.

It also satisfies the exact FixedHyp. The lattice is unimodular. At any prime P the intersection Lambda∩span\_E(e\_a) is the sum a\_a e\_a; choose a local generator c\_a of a\_a. At G choose c\_a=1, which is permitted because d is a G-unit. This is the first IsFixedConv clause. For the other clause,

\[
 L_{\rm sub}=\operatorname{span}_E(e_a)\oplus L_b,
 \qquad c z_u\in L_{\rm sub}\Longrightarrow c\in\mathcal O_E.
 \tag{HF6}
\]

The implication follows from cA\_u,cB\_u integral and A\_u−B\_u=1. These equalities follow directly from the source Lrep/Lsub definitions; there is no inference from mere rational primitivity. Finally, when the actual norm generator pi\_u is available, LIFT30(NL3--6)'s rational SU(k,1) matrix on the p\_1,z\_* plane extends by the identity on all hyperbolic planes. It realizes the required positive complement and fixes e\_a,f\_a. Its Bessel translation and its real-vector transport remain distinct from the level translation constructed below.

\subsection{3. Exact finite Hida vectors}\label{exact-finite-hida-vectors}

Use the fixed global trace character psi\_E,f=psi\_Q,f composed with Tr\_(E/Q), whose rational product with the literal psiEinf is one. It is trivial on integral finite adeles. At every selected G component it induces a nontrivial residue character psi\_H,w by x↦psi\_E,f(x/G). Both partners of a relative split place induce the same residue character in the rational degree-one coordinates. No unproved equality with the prescribed generic psi0 is assumed.

At all places over G define e0\_a=e\_a, f0\_a=f\_a and z0=z\_*. At other places set these finite vectors to zero. Let ew be one at the selected halves of G and zero elsewhere. Define

\[
             b_i=e_w p_i,\qquad b_i^d=\bar e_w p_i.
 \tag{HF7}
\]

All vectors are in Lambda\_hat, since d is a G-unit. Their local pairings over G are exact: h(e0\_a,f0\_b)=delta\_ab, h(z0,z0)=−1, and h(b\_i\^{}d,b\_j)=delta\_ij ew. All other pairings in gramIntegral/formsAtZero/weightVectorsIntegral vanish there. Away G the required congruences are automatic because G is a unit. The same observation proves ew\_full and latticeSpan. The two components together span the entire hyperbolic/core lattice modulo G.

The e0 and f0 vectors agree exactly with the rational vectors at G. The finitely many other denominators of (HF3) can be cleared by a rational integer coprime to G. Hence basisApprox holds with its literal InAt(G,G) quantifier. For every word choose

\[
                            x_u=1.
 \tag{HF8}
\]

Then xu\_fix, xu\_K and xu\_one hold exactly, and xu\_z follows from (HF5). Thus all geometric, fractional-ideal, finite-vector and xu fields through HidaSetting.lean279 are explicit. The real representation and Wu fields are discussed separately in section9.

\subsection{4. The whole finite subgroup, using both CM components}\label{the-whole-finite-subgroup-using-both-cm-components}

At a selected G place write matrices on the plus component in the ordered basis

\[
 (e_0,\ldots,e_{\rho-1},p_1,\ldots,p_k,z_*,
                                     f_{\rho-1},\ldots,f_0).
 \tag{HF9}
\]

Let c=rho+k. Its Gram matrix J\_0 pairs e\_a with f\_a, is one on the p\_i, and is −1 on z\_*. Thus J\_0\^{}2=1. If A is the plus matrix, the other matrix is exactly

\[
                         A^- =J_0 A^{-\mathsf t}J_0.
 \tag{HF10}
\]

The plus e0 conditions give unit upper-triangular first rho columns; the plus b conditions give the next k such columns. The plus z0 condition says its column equals z0 modulo the e-span, in particular all entries in rows rho,\ldots,c−1 of column c vanish. The minus e0 conditions, via (HF10), say that A preserves the dual decreasing coflag and acts trivially on its successive quotients. Together with the preceding columns these make the last rho columns upper triangular with diagonal one. The minus z0 condition gives the remaining z0 coflag condition, already implied by this description. More explicitly, the minus e0 equations say that row f\_a of A\^{}(-1) is row f\_a of the identity plus rows f\_b with b\textless a. Inverting this unit upper-triangular coflag action gives the identical row-support condition for A. These rows have no e or middle component. The plus middle columns then give the middle block diag(U\_k,1); row z\_* of A\^{}(-1) differs from its identity row only in the f-span, which is exactly the remaining minus z0 equation. Conversely an A with this description has an inverse of the same form; substitution into (HF10) verifies every minus condition. Therefore the equivalence is exact:

\[
 K_2\bmod w=U_{\rho,k}(q):=
 \{A\in U_r(q): A_{ic}=0\ (\rho\le i<c)\}
               =\operatorname{unipSym}(q,r,\rho,k).
 \tag{HF11}
\]

There is no condition on f0 in the literal K2cong; its coflag conditions above arise from the minus e0 component. Omitting that component would give a larger incorrect subgroup. This calculation applies also at rho=0 and specializes to R025's diag(U\_k,1).

\subsection{5. The exact root character and its finite diagonal correction}\label{the-exact-root-character-and-its-finite-diagonal-correction}

On the group (HF11), put L\_H(A) equal to the residue of the full scalar inside chi2f, summed over the two partners of w. If rho\textgreater0,

\[
\begin{split}
 L_H(A)={}&\sum_{a=0}^{\rho-2}
       (A_{a,a+1}-A_{r-2-a,r-1-a})
       +A_{\rho-1,c}+A_{c,c+1}\\
       &+\sum_{i=\rho}^{c-2}A_{i,i+1}.
\end{split}
 \tag{HF12}
\]

When rho=0 only the last sum is retained. For the plus partner, h(f\_a,A nxt\_a) is the e-chain entry, or the shortcut to z\_* at the last a. For the minus partner, (HF10) makes it (A\^{}(-1))\emph{(f}(a+1),f\_a), or −(A\^{}(-1))\emph{(z,f\_last) at the end. These are respectively the negative f-chain entry and the positive z\_}→f\_last entry. All are adjacent except the permitted e→z shortcut; the forbidden k entries ensure that shortcut has no intermediate contribution. The b\_i\^{}d,b\_(i+1) pairing is supported only at the selected partner and gives its displayed b-chain entry. This proves (HF12), including its signs and with no factor two.

The source symEdge character has the same edges as (HF12), with coefficient +1 on every edge. These edges are two disjoint directed chains: (e\_0,\ldots,e\_last,z\_*,f\_last,\ldots,f\_0), and (p\_1,\ldots,p\_k). Let alpha\_ij in \{1,−1\} be the coefficient of an edge in (HF12). For the prescribed nontrivial additive character psi0\_w there is a unique h\_w in F\_q\^{}× with

\[
          \psi_{H,w}(h_w x)=\overline{\psi_{0,w}(x)}.
 \tag{HF13}
\]

This uses q rational prime, as required by the actual G data. Choose d\_i=1 at the initial vertex of each chain and recursively set d\_j=alpha\_ij d\_i/h\_w. Let D\_w=diag(d\_i). Every entry is a unit, and

\[
 \psi_{H,w}\bigl(L_H(D_w A D_w^{-1})\bigr)
       =\overline{\psi_{0,w}(\operatorname{symHom}(A))}
                  \qquad(A\in U_{\rho,k}(q)).
 \tag{HF14}
\]

Indeed alpha\_ij d\_i/d\_j=h\_w on every edge. The diagonal preserves the whole group (HF11). This is an explicit correction of both CM-component signs and the additive-character scalar. It does not change the global e\_a,f\_a, any fractional ideal, or any finite level.

\subsection{6. Matching the literal bottom dat0 and its finite action}\label{matching-the-literal-bottom-dat0-and-its-finite-action}

Let the actual RealDat bottom frame beta\_Y0 be an integral local basis of the same standard lattice, as in the actual vertical family. Let pi be the permutation in grp\_levels/IsSymOrd. The latter orders the original first k coordinates after the rho e-coordinates; LinksChar.lean73--77 fixes the precise reindexing convention.

At each selected G place choose the plus component of an integral unitary t so that the matrix of t\^{}(-1)kt in beta\_Y0 coordinates, after that exact permutation, is D\_w\^{}(-1)A(k)D\_w. A plus matrix is arbitrary GL\_r(O\_v) at a relative split place, so this determines an actual unitary element by taking the inverse-adjoint other component. Lift each d\_i to O\_v\^{}×. At every other finite place and at infinity take t=1. Thus t belongs to the original maximal compact at every finite place. This t is the finite level translation, not the norm-minus-T Bessel translation of LIFT30.

Since dat0 has rank r, j=0 and Q empty, its full flag group is the whole lattice stabilizer K, its full flag character is one, and dat0(LevelP.triv) is precisely the principal G subgroup K\^{}1. There are no hidden determinant-tail or Q factors at the bottom. Combining (HF11), (HF14) and the exact grp\_levels relation gives

\[
 t^{-1}K_2t=\operatorname{levS}(Y_0,\beta_{Y0},
                 \operatorname{dat0}(\operatorname{levR}(F,\mathrm{true}))),
 \qquad
 \operatorname{levChi}(t^{-1}kt)=\overline{\chi_2(k)}.
 \tag{HF15}
\]

Both directions of the group equality were proved, though the literal LevelMatchA needs only its forward direction. This is exactly A under the fixed source iotaX transport. No sign is changed by a convention about automorphic right translation: the equality is between the two literal characters.

Define red on K to be the residue matrices in the transported beta\_Y0/t bases, identified with the prescribed G0g factors through their degree-one field isomorphisms; outside K put red=1. Integral global lattice representatives of each residue basis exist by the surjection Lambda→Lambda/P Lambda. They give the exact IsRedBasis and RedMat witnesses of RowDict.lean49--64. The representative basis may be chosen once per P, independently of k. Hence GoodAt holds. The map on K is a homomorphism, is onto the product of full GL\_r groups by independent local lifts, and has kernel K\^{}1. The conjugate components need no additional independent choice. This proves all four clauses of IsRedHom.

For any actual principal-level automorphic coefficient f and k in K, the literal level action is

\[
       (\operatorname{actL}(\sigma(\operatorname{red}k))f)(g)
                            =f(g\,t^{-1}kt).
 \tag{HF16}
\]

To check this without any choice of representative, two lifts of the same residue element differ by K\^{}1; f is right K\^{}1-invariant. The normalizing character in actL is the full flag character, which is one at dat0. This proves (HF16) for every map T, every real vector v and every g in LevelMatchB, not just one selected automorphic form. The permutation, diagonal conjugations and residue field labels are incorporated into the same group isomorphism sigma. They are not a new arbitrary finite representation.

Finally assume the actual family's literal grp\_record identifies its finite idempotent with the genuine complete goodTuple idempotent. Accepted R026(R26.8) proves that complete idempotent coefficientwise real for arbitrary source character lists. Centrality makes the inner changes of basis above immaterial. Therefore

\[
 \operatorname{RealDat.eC}(\sigma\gamma)
          =e_0(\gamma)=\overline{e_0(\gamma)},
 \tag{HF17}
\]

exactly LevelMatchC. A separately chosen arithmetic family with no grp\_record equality is not assigned this projector by fiat. Under that explicit source tie, (HF15)--(HF17) prove all finite RowLevelMatchAt clauses with the single t already constructed.

\subsection{7. Exact finite measures and projector normalization}\label{exact-finite-measures-and-projector-normalization}

The k forbidden entries in (HF11) are independent. Thus

\[
 |U_{\rho,k}(q)|=q^{r(r-1)/2-k},\qquad
 [GL_r(q):U_{\rho,k}(q)]=q^k\prod_{i=1}^r(q^i-1).
 \tag{HF18}
\]

For the actual finite Haar, write v\_K=vol(K). Then

\[
 \operatorname{vol}(K^1)=
       \frac{v_K}{\prod_w|GL_r(q_w)|},\qquad
 \operatorname{vol}(K_2)=
       \frac{v_K}{\prod_w q_w^k\prod_{i=1}^r(q_w^i-1)}.
 \tag{HF19}
\]

The product is over the selected relative places at both G primes, one per pair; counting both CM halves again would square the index incorrectly. The exact Hida finite idempotent is

\[
 e_\chi(k)=\frac{\overline{\chi_2(k)}}{\operatorname{vol}(K_2)}
                      \boldsymbol1_{K_2}(k).
 \tag{HF20}
\]

It is an orthogonal normalized compact-character projector in every unitary right regular representation, by the group/character calculation. It has e\_chi(1)=1/vol(K2), not one. It is locally constant, bi-K\^{}1-invariant, positive type and supported in K, and all matrices in its support are unipotent modulo g0. Thus the literal IsAdmIdem conditions are verified. The goodTuple projector remains a separate central projector, commuting with this one; no nonzero intersection is inferred from commutation alone. For completeness, convolution idempotence is the character law and one integration over K2; the adjoint identity uses chi2(k\^{}(-1))= bar(chi2(k)). In the positive-type matrix for any finite list of group elements, partition the list by their K2 cosets. Each block is a nonnegative scalar times a rank-one character Gram matrix. Thus positivity is verified for the extension by zero, not merely for the restriction of the function to K2.

At a prime ell distinct from g0,g1, the ell-valuation of both {[}K:K\^{}1{]} and {[}K:K2{]} is sum\_w sum\_(i=1)\^{}r v\_ell(q\_w\^{}i−1). This records the actual finite measure denominator without claiming that every particular projected row requires that full denominator. The compact conjugation t preserves Haar and adds no group index. Its unit local entries at G do not remove the actual different scaling or the independent Bessel transport at T.

\subsection{8. The Bessel subgroup and its exact rational character}\label{the-bessel-subgroup-and-its-exact-rational-character}

The flag (HF3),(HF5) has the literal stabilizer S of HidaSetting: (s−1)e\_a lies in the preceding flag, and (s−1)z\_u lies in its span. Its subgroup N additionally acts trivially on X\^{}perp/X. Consequently S/N is the unitary group of z\_u\^{}perp inside L\_b⊗E, which is positive definite of rank k at infinity. Fixing both e\_a and f\_a gives a Levi representative, so S is N semidirect H\_u. This follows by projecting to the nondegenerate middle quotient; the kernel is exactly the displayed N condition, and its section acts on the middle orthogonal complement and fixes the hyperbolic pairs and z\_u.

For n in N, set

\[
 \ell_u(n)=\sum_{a=0}^{\rho-1}
                   h(f_a,n\,\operatorname{nxt}(e,z_u,a)).
 \tag{HF21}
\]

Each summand is the coefficient on a simple step of the e-flag, with the final step to z\_u. The defining upper flag action shows ell\_u(nn')=ell\_u(n)+ell\_u(n'); all intermediate products miss a simple step. The positive middle group fixes this coefficient, so the character extends over S by triviality on H\_u. With the same psi\_E,f as in section3, the finite and real characters are exactly

\[
 \psi_{u,f}(n)=\psi_{E,f}(\ell_u(n)/G),\qquad
 \psi_{u,\infty}(n)=
  \exp\!\left(2\pi i\operatorname{Tr}_{F/\mathbb Q}
                 (\ell_u(n)+\bar\ell_u(n))/G\right).
 \tag{HF22}
\]

The real formula means the product over the specified CM embeddings, as in psiInf, and not an additional doubled trace. Their product is one on rational S(F), by the global trace-character product formula. This preserves the source denominator G and all different-scaled e\_a,f\_a; replacing them by a convenient standard Whittaker character would require a further explicit conjugation and measure calculation.

The compact quotient H\_u has the prescribed positive k-dimensional Hermitian form. Its finite-dimensional real representation of type u can be inflated across N. Haar measures on N and H\_u, and the resulting Weil quotient measure on S\G, exist with the source's chosen ambient Haar; their scale remains explicit in IsWu. We have not used the existence of these measures to assert the existence of a nonzero relative discrete-series vector.

\subsection{9. What this actually closes and what it leaves}\label{what-this-actually-closes-and-what-it-leaves}

The construction closes the exact scaled flag, primitive norm-minus-T line, finite congruence vectors, rational character, xu conditions, whole finite subgroup and character, finite quotient action and all three RowLevelMatch clauses when the separately required genuine projector tie is supplied. It works at every rho, not merely m=2, and it keeps the actual standard lattice on principal-different hosts. All source definitions are matched literally.

The full HidaSetting still asks for its discrete-series realization, lowest unit vector, the convolution Rph, and especially a continuous nonzero Wu with the stated left S-character, real-type reproduction, L1 quotient integrability and unit L2 norm, together with its local tensor factorization and exact text-normalized measures. The definition \nolinkurl{ConcreteSettingExists.lean16–20} merely names that existence predicate under its much larger gap; it is not a theorem constructing these fields for the present mixed input. Those fields cannot be filled by a generic cusp nonzero statement. Likewise RowsJoinPinned, the same prescribed Bessel/Mram period, the full old Host/D7 record and the net DP44 cost remain separate.

\subsection{10. Exact sources and verification}\label{exact-sources-and-verification}

All source files refer to the pinned Alpöge archive/mirror commit 7cbf7a0c32fca1badf05931879a37c3f8a2b03b4. Newly recovered BlockEnds.lean and SArith.lean were recorded in \nolinkurl{work/research/source_recovery_20261122T122450Z.json}; recovery is not proof checking or Lean execution.

\begin{itemize}
\tightlist
\item
  \nolinkurl{DataSpaces.lean30–46,82–124:} actual Hermitian spaces, lattice, rank, primitive negative line and hyperbolic complement fields.
\item
  \nolinkurl{BlockEnds.lean29–39,111–119:} FlagG, VS and exact extended different.
\item
  \nolinkurl{IsPrimitiveSub.lean92–104,312–315:} Lrep/Lsub, IsFixedConv and complete FixedHyp.
\item
  \nolinkurl{HidaSetting.lean117–136,172–279,312–327:} entire finite congruence and character, scaled ideals, geometric/finite fields and the separate real Wu fields.
\item
  \nolinkurl{DefGroup.lean175–201,217–269,608–609,620–707:} exact unipCol, symEdge and symHom; actual bottom residue group/character.
\item
  \nolinkurl{LinksChar.lean73–77,106–119:} the literal symmetric permutation, grp\_levels and genuine finite projector tie.
\item
  LinksDefs.lean427--431; \nolinkurl{SubgroupCharExtensionOneApply.lean32–43,} 59--73,95--110: dat0, whole flag group/character and the exact residue right action.
\item
  RowDict.lean49--82 and RowLevelMatch.lean18--47: residue-basis witnesses and all A/B/C clauses.
\item
  \nolinkurl{UnitaryHaarProducts/InstExists.lean18–41,225–260:} the exact source coordinate instance and its standard-lattice transport.
\item
  K030(PH20--21), K032 section3: the actual hosts' principal codifferent and relative finite-unramified property. R025 is the rho=0 adapter; R026(R26.8) is the accepted genuine projector reality theorem; LIFT30(NL3--6) is the actual negative-plane isometry. These accepted findings are used at their stated scope.
\end{itemize}

The companion exact finite-matrix certificate is diagnostic. It checks both split components, the complete subgroup equivalence, the inverse-chain signs, diagonal additive-character correction, finite indices, primitive line identities and source fingerprints. It does not certify a real Bessel model or an automorphic period.

\clearpage

\section{AM-K079 --- actual fixed-block exterior minors and paired finite rank loss}\label{am-k079-actual-fixed-block-exterior-minors-and-paired-finite-rank-loss}

\textbf{Canonical checked recovery version --- 2026-11-22T16:10:03.402339+00:00.} Root has read the complete mathematical proof and companion script and replayed the stated exact certificate. The original recovery export and its historical status are preserved separately. This is a current checked version, not recovery of unavailable historical final bytes. Only the stated sector and hypotheses are promoted; the signed period and actual DP44 cost remain open.

\subsection{1. Exact sector and carried interfaces}\label{exact-sector-and-carried-interfaces}

Use K073(FB1--4,8--10), DO063(FV5--19) and DO066(GC20):

\[
 a=f=168,\quad n=169,\quad H=8192,\qquad
 E^n=S\perp W_0,\quad \dim_E S=164,\quad
 J_{W_0}=\operatorname{diag}(1,1,1,1,-1).
 \tag{ER1}
\]

The unchanged primes are g0=22285932761 and g1=64072175441, with Q=g0 g1. T is exactly the 996-bit J56 prime in the frozen DO040 certificate. There are 2(f−1)=334 unmixed relative G places, whose residue-size product is Q\^{}(f−1). Each has two CM partners. Both mixed finite partial traces retain only their modulus bound one.

Consider an actual nonintegral supported matrix with central normalization

\[
 h=I_S\oplus h_0,\quad h_0\in U(W_0)(E),\qquad
 \max_\tau\|h_\tau\|\le H,\qquad N=h-I.
 \tag{ER2}
\]

Let Σ be its nonintegral relative T support and s=\textbar Σ\textbar. All other finite places are integral. The central scalar is treated exactly as in DO060(IB12,28--29): it changes the finite and real coefficients by phases only; each central class has one normalized representative in ER2. No additional factor28 is inserted.

The coefficient is the original a(h)c(h) product rho\_v(h) of K062(AT23). Its selected real compression satisfies \textbar c(h)\textbar≤1, since the representation is unitary and the normalized compressed trace is a contraction. This bound needs neither unipotence nor a special selected-vector estimate.

If a(h)≠0, every unmixed G reduction is unipotent. Indeed it has the eigenvalue1 from the fixed S block; the accepted Green support formula LIFT42(FC6), used in DO060(IB13--14), forces the single primary irreducible polynomial to be X−1. This is a statement about the finite reductions, not about the eigenvalues of h over E.

\subsection{2. Exterior witnesses at the actual T pencil}\label{exterior-witnesses-at-the-actual-t-pencil}

Put

\[
 I_\Sigma=\prod_{\mathfrak P\mid\Sigma}\mathfrak P^{-1},
 \qquad N_{E/\mathbf Q}(I_\Sigma)=T^{-2s}.
 \tag{ER3}
\]

Both CM partners occur in the product. DO060(IB16,23) and DO061(TR22) already prove that every exterior entry of h has pole order at most2 at each selected component. The proof uses the literal factorization A\_d K A\_c\^{}(-1), where A\_c=I+c pi\^{}(-1) z q and qz=0: each exterior power of A\_c is linear in pi\^{}(-1), and the middle K is integral. The same statement holds at the other CM component by inverse-adjoint. Expanding a minor of h−I using its selected identity entries expresses it as a sum of minors of h, of possibly smaller size. Therefore every k-minor y of N satisfies

\[
                         y\in I_\Sigma^2.
 \tag{ER4}
\]

We invoke this accepted pole statement; no new global unipotence or coefficient-integrality inference is being substituted for it.

For any h in U(4,1), its positive Euclidean singular values are t,1,1,1,t\^{}(-1), where 1≤t≤H; this is DO061(TR12)/K068(LR5). On a codimension-one subspace h has norm at most1, so N=h−I has norm at most2 there. The minimum characterization of singular values gives sigma\_j(N)≤2 for j≥2, while sigma\_1(N)≤H+1. Hence

\[
 |\tau y|\le\|\bigwedge^k N_\tau\|
 \le K_k:=2^{k-1}(H+1),\qquad 1\le k\le5.
 \tag{ER5}
\]

If rank\_E(N)≥k, choose one nonzero k-minor y. Its existence is used matrix by matrix; there is no sum over choices of the witness. The integral ideal (y)I\_Σ\^{}(-2) has norm at most

\[
             \mathrm N((y)I_\Sigma^{-2})\le K_k^{2f}T^{4s}.
 \tag{ER6}
\]

No fractional ideal generator is required.

\subsection{3. Paired rank loss}\label{paired-rank-loss}

At an unmixed relative G place w let d\_w be the rank of the reduction of N, and let q\_w be its residue size. The rank at its CM partner is the same: h\_barP is inverse-adjoint, and h\^{}(-1)−I=−h\^{}(-1)(h−I). The elementary DVR minor bound K068(LR8) gives

\[
 \nu_w:=v_{\mathfrak P_w}(y)+v_{\bar{\mathfrak P}_w}(y)
                  \ge2(k-d_w)_+.
 \tag{ER7}
\]

Define

\[
 L_j=jn-j(j+1)/2,\qquad L_0=0.
 \tag{ER8}
\]

For d\textless k,

\[
 L_k-L_d=(k-d)\bigl(n-(k+d+1)/2\bigr)
       \le (L_k/k)(k-d)
       \le (L_k/(2k))\nu_w.
 \tag{ER9}
\]

The exact unipotent Green ratio is product\_(j=n−d)\^{}(n−1)(q\^{}j−1)\^{}(-1). For d≥k it is bounded by its first k denominator factors; for d\textless k use ER9. Across all334 places, the product of inverse factors (1−q\^{}(−j))\^{}(-1), n−5≤j≤n−1, is less than2, as certified by the actual primes. This also dominates the shorter products for every k≤5. The two mixed projector factors are at most1. Equations ER6--9 therefore prove

\[
 \boxed{\quad
 |a(h)|\le
 \frac{2K_k^{L_k f/k}T^{2L_k s/k}}
      {Q^{L_k(f-1)}}\qquad(\operatorname{rank}_E N\ge k).
 \quad}\tag{ER10}
\]

The paired valuation in ER7 is essential. For k=2, rank0 gives nu≥4 and rank1 gives nu≥2. Their losses335 and167 are bounded by (335/4)nu. Counting just one selected CM prime would not give ER10. There is no rank-pattern multiplicity: ER10 is a uniform pointwise bound obtained from one nonzero witness.

\subsection{4. Complete actual column and incidence count}\label{complete-actual-column-and-incidence-count}

Fix Σ and all its actual ordered source/target pencil pairs. The actual target W0 lattice L is the paired Hermitian self-dual lattice of DO063(FV6), with determinant valuation0. That proof gives, for any unitary h satisfying these incidences,

\[
 h e_n\in\mathcal B_\Sigma=I_\Sigma L,\qquad
 \operatorname{covol}(\mathcal B_\Sigma)=D_0^5T^{-10s},
 \quad D_0=7^{140}29^{144}5^{84}13^{84}.
 \tag{ER11}
\]

These assertions use no unipotence. Specifically the source contains pi e\_n, so its image is in the actual target, and the partner is treated by the actual Hermitian dual lattice. DO063(FV15--19) gives covering radius less than3750 sqrt(f) and negative-column radius less than 5793 sqrt(f). The complete Dirichlet-cell volume count is consequently

\[
 \#\{h e_n\}\le B_W T^{10s},\qquad
 B_W=\frac{[(22/7)f\,9543^2]^{5f}}{(5f)!D_0^5}<2^{14861}.
 \tag{ER12}
\]

For each full column the quotient of any two isometries is a source lattice isometry fixing a negative vector. It is a finite group. The accepted injective reductions at both G primes in DO060(IB28--29) give its order bound

\[
 G_5=\gcd(|GL_5(\mathbf F_{g_0})|,|GL_5(\mathbf F_{g_1})|)
   =98148490017978252769652736000000.
 \tag{ER13}
\]

This source-lattice stabilizer argument is independent of unipotence. The rank-five restriction follows from h fixing S pointwise; a larger moving-space fibre is not inserted. Both ER11 and ER13 refer to the same chosen actual source/target pairs.

K062(AT7--11) gives at most(T−1)\^{}2 ordered actual pairs at each nonintegral place and \textbar rho\_v\textbar≤2/(T−2). There are binom(f,s) supports. Together with \textbar c\textbar≤1, ER10--13 give the complete absolute bound

\[
 \boxed{\quad
 \mathcal R_{k,s}\le
 2\binom fs G_5 B_W
 \frac{K_k^{L_kf/k}T^{(10+2L_k/k)s}}{Q^{L_k(f-1)}}
 \left[\frac{2(T-1)^2}{T-2}\right]^s.
 \quad}\tag{ER14}
\]

This counts the entire supported sector rank(N)≥k, including repeated incidences as an upper bound. All original normalized partial traces, selected metrics and scalar-packet weights are retained. There is no termwise positivity assertion.

\subsection{5. Exact ranges and remaining scope}\label{exact-ranges-and-remaining-scope}

The companion directed-rational certificate gives:

\begin{longtable}[]{@{}
  >{\raggedleft\arraybackslash}p{(\columnwidth - 8\tabcolsep) * \real{0.2000}}
  >{\raggedleft\arraybackslash}p{(\columnwidth - 8\tabcolsep) * \real{0.2000}}
  >{\raggedleft\arraybackslash}p{(\columnwidth - 8\tabcolsep) * \real{0.2000}}
  >{\raggedleft\arraybackslash}p{(\columnwidth - 8\tabcolsep) * \real{0.2000}}
  >{\raggedleft\arraybackslash}p{(\columnwidth - 8\tabcolsep) * \real{0.2000}}@{}}
\toprule\noalign{}
\begin{minipage}[b]{\linewidth}\raggedleft
Minimum rank k
\end{minipage} & \begin{minipage}[b]{\linewidth}\raggedleft
L\_k
\end{minipage} & \begin{minipage}[b]{\linewidth}\raggedleft
Last certified support
\end{minipage} & \begin{minipage}[b]{\linewidth}\raggedleft
Inverse binary exponent there
\end{minipage} & \begin{minipage}[b]{\linewidth}\raggedleft
Next support exponent
\end{minipage} \\
\midrule\noalign{}
\endhead
\bottomrule\noalign{}
\endlastfoot
2 & 335 & 10 & 78276 & −266152 \\
3 & 501 & 15 & 292336 & −51096 \\
4 & 666 & 21 & 162350 & −180086 \\
5 & 830 & 27 & 32744 & −308696 \\
\end{longtable}

Each smaller positive support in that row also satisfies the displayed upper expression less than1. The negative next exponent describes a failure of this upper estimate, not a nonzero lower bound on the sum.

On the current support-floor10 remainder of K078, remove the disjoint support/rank sectors

\[
 (s=10,\ r\ge2),\quad(11\le s\le15,\ r\ge3),\quad
 (16\le s\le21,\ r\ge4),\quad(22\le s\le27,\ r=5).
 \tag{ER15}
\]

There are18 support slices, and every slice has bound at most 2\^{}(-32744); their absolute union is less than2\^{}(-32739). This theorem does not remove uniform rank-one fixed-W0 terms at support10. Nor does it remove nonscalar or cross-block terms at that support. No complete support10 exhaustion, Mram positivity, period nonvanishing, arithmetic gain or literal Host/FamP constructor follows from ER15 alone.

The new proof is ER5 and ER7--10 together with their application ER14--15. The exterior pole, actual lattice, stabilizer and Green formula are carried checked inputs at the exact locators above. The certificate is an arithmetic and local-algebra diagnostic, not a rational short-sum enumeration, automorphic computation or Lean build.

\clearpage

\section{AM-K080-RC-20261122 --- fresh reconstruction of the small-prime divisor refinement}\label{am-k080-rc-20261122-fresh-reconstruction-of-the-small-prime-divisor-refinement}

\textbf{Canonical integration --- 2026-11-22T16:10:03.402339+00:00.} Whole root proof/script reading and separate exact replay are recorded in CHECK\_EVENTS.tsv. Original recovery bytes remain preserved.

Author: GPT-6 Astra Pro, ALP\_MACHINE /root/kummer\_extension. Fresh reconstruction on 22 November 2026 after restoration of checkpoint23.

\textbf{Fresh author whole-proof and complete-script check PASS --- 2026-11-22 15:36:16 UTC.} This is a new reconstruction from the restored K075--K079 inputs. It does not claim that the unavailable former K080 bytes, certificate, or hashes were recovered. The historical retained hashes are recorded separately in the recovery inventory. The new exact certificate has its own source hash and time.

\subsection{1. The actual field has only norm29 primes below256}\label{the-actual-field-has-only-norm29-primes-below256}

Use the exact host of K032 and K075(RN1): \[
 E=\mathbf Q(\zeta_{28})B\mathbf Q(\sqrt5,\sqrt{13}),\qquad
 [E:\mathbf Q]=336,\quad F=E^+,\quad [F:\mathbf Q]=168,
\] where B is the cyclic real degree-seven subfield of Q(zeta29). The fields B and Q(zeta7)\^{}+ are contained in F. The subgroup fixing B is H29=\{1,12,17,28\} modulo29.

\textbf{Lemma RC1.} The only prime ideals of E or F with norm below256 lie above29. There are48 such primes in E and24 in F, all of norm29.

\textbf{Proof.} At2 the residue degree in B is7. In Q(zeta7)\^{}+ it is3, because the order of2 modulo7 up to sign is3. Hence the residue degree in either E or F is divisible by21, and the norm exceeds256.

For an odd rational prime p\textless256 other than29 with p modulo29 outside H29, the residue degree in B is7. Its norm in E or F is therefore at least p\^{}7\textgreater256, even if another subextension is ramified at p.

The complete remaining list is17,41,59,157,173,191,233. Each is at least17 and has p modulo28 different from ±1. Thus its residue degree in the real cyclotomic field Q(zeta28)\^{}+ is at least2, giving norm at least17²\textgreater256. The companion certificate enumerates every rational prime below256 by exact trial division and checks this entire list. Rational primes at least256 cannot give a smaller prime-ideal norm.

At29, B is totally ramified of degree7 with residue degree1. The prime splits completely in Q(zeta28) because29≡1 modulo28, and in both real quadratic fields because29 is a nonzero square modulo5 and13. The actual compositum degree and disjoint ramification, already fixed in K032, therefore give residue degree1 and ramification index7: there are336/7=48 primes in E and168/7=24 in F. No other prime ideal of norm below256 remains. ∎

\subsection{2. An exact divisor bound with exponent one eighth}\label{an-exact-divisor-bound-with-exponent-one-eighth}

For a number field K let d\_K(a) count integral ideal divisors of a nonzero integral ideal.

For every prime ideal of norm q≥256 and every integer e≥0, \[
 e+1\le2^e\le q^{e/8}.
\] At a norm29 prime, the normalized factor (e+1)29\^{}\{-e/8\} is maximal at e=1. Indeed the e=0 value is smaller because29\textless2\^{}8; for e≥1 the ratio of consecutive values is at most (3/2)29\^{}\{-1/8\}\textless1, since29\textgreater(3/2)\^{}8. Unique ideal factorization and RC1 therefore give \[
 d_E(\mathfrak a)\le C_E\,N_E(\mathfrak a)^{1/8},
 \qquad C_E=2^{48}/29^6,
\] \[
 d_F(\mathfrak c)\le C_F\,N_F(\mathfrak c)^{1/8},
 \qquad C_F=2^{24}/29^3.
 \tag{RC2}
\] These constants allow all exceptional primes, including those absent from a particular ideal, and are consequently uniform. Their exact product is \[
                     C=C_EC_F=2^{72}/29^9. \tag{RC3}
\] The estimate is an elementary ideal-factorization bound; no asymptotic divisor theorem or implicit field constant is being invoked.

\subsection{3. Substitute only the two divisor factors in the actual rank-one count}\label{substitute-only-the-two-divisor-factors-in-the-actual-rank-one-count}

Retain exactly the globally unipotent rank-one fixed-W0 sector, original finite packet, projectors, class masses, PBW metric and affine pencil data of K075 and K077. In particular the original negative-column entry t is nonzero and pure imaginary. Use K077's notation \[
 a=f=168,\quad H=8192,\quad R=(H-H^{-1})/2,\quad
 M=Q^{f-1},
\] and, for support size s and actual double-pole subset size d, \[
 z_0=N_E(\mathfrak J)^{1/(2f)}T^{-(s+d)/f},\quad
 y=R/z_0,\quad y_c=R/(z_0T^{2d/f}).
\] The nonempty actual set has y≥1.

K075(RN11--13) embeds the coordinate-content ideals into divisors of an integral E-ideal of norm at most y\^{}\{2f\}. RC2 changes its bound from y\^{}70 to \[
                              C_E y^{42}.
\] The fourth-coordinate norm equation in K075(RN16--19) is counted over F, with relative norm ideal of norm at most y\^{}f.~RC2 changes its bound from28y\^{}35 to \[
                              28 C_F y^{21}.
\] The relative-F divisor count, not an E divisor bound, is essential here.

All other factors remain those proved in K077(AP6--18). In particular: the pure-imaginary t count is (1+2y)\^{}f; the first three coordinates occupy the specified affine cosets and have count (1+2sqrt(y\_c))\^{}\{6f\}; every full column determines the rank-one matrix uniquely; and there is no extra isometry-fibre multiplier. The complete new count for fixed actual support, pencil data and rank-zero pattern is therefore \[
       28C\, y^{63}(1+2y)^f(1+2\sqrt{y_c})^{6f}.       \tag{RC4}
\] Only the old exponent105 and its absent divisor constant have changed.

Before the rank-zero cap, multiplying RC4 by the original finite G and selected real factors gives z\_0 dependence proportional to \[
 z_0^{27405}(z_0+2R)^f
       (\sqrt{z_0}+2\sqrt R\,T^{-d/f})^{6f}.          \tag{RC5}
\] The exponent27405 is positive, so this function increases for every z\_0\textgreater0. The actual bound N\_E(J)≤M² therefore allows the same cap \[
 z_{s,d}=M^{1/f}T^{-(s+d)/f},\quad
 y_{s,d}=R/z_{s,d},\quad
 y'_{s,d}=R M^{-1/f}T^{(s-d)/f}.                    \tag{RC6}
\] This is a monotone extension of an upper expression; the cap need not itself describe a nonempty actual set.

\subsection{4. Retain the complete marked type sum and both endpoints}\label{retain-the-complete-marked-type-sum-and-both-endpoints}

For fixed s,d the unchanged K077(AP18) type weight is \[
                         \binom fs\binom sd(T+1)^{s-d}.
\] The rank-zero G patterns, single Euler factor, root-unit generator count and selected real factor are also unchanged. Put \[
 C_s=2^{335}\,28C\binom fs\,12^f2^{f/2}
             \prod_\tau E_\tau^{-1/4},
\] \[
 B_s(d)=(T+1)^{s-d}z_{s,d}^{-f/2}y_{s,d}^{63}
       (1+2y_{s,d})^f(1+2\sqrt{y'_{s,d}})^{6f}.
\] The logarithm of B\_s(d) is convex on0≤d≤s: its non-affine terms are positive multiples of log(1+exp(A+Bd)), and replacing105 by63 changes only an affine term. Thus the full binomial sum is bounded by \[
 \boxed{\displaystyle R_{1,s}\le
                   2^s C_s\max\{B_s(0),B_s(s)\}.}       \tag{RC7}
\] Neither endpoint, any marked double-pole set, nor any original pair multiplicity is omitted. This is an absolute bound; it assigns no sign to the actual summands.

The reconstructed certificate reads the actual T-prime and actual K044 interval from the restored original files. It verifies the same uniform homogeneous-energy floor E\_tau\textgreater2\^{}494 and the allowed all-top floor E\_tau\textgreater2\^{}959, then brackets the f-th roots and square roots by rationals and uses directed rational logarithms. The decisive margins are:

\begin{longtable}[]{@{}rrr@{}}
\toprule\noalign{}
Support & Uniform494 margin & All-top959 margin \\
\midrule\noalign{}
\endhead
\bottomrule\noalign{}
\endlastfoot
6 & 16320 & 35850 \\
7 & 12581 & 32111 \\
8 & 8843 & 28373 \\
9 & 5105 & 24635 \\
10 & 1367 & 20897 \\
11 & estimate exceeds1 & 17159 \\
12 & estimate exceeds1 & 13422 \\
13 & estimate exceeds1 & 9684 \\
14 & estimate exceeds1 & 4484 \\
15 & estimate exceeds1 & estimate exceeds1 \\
\end{longtable}

In particular the complete rank-one fixed-W0 support-ten sector is \[
                              R_{1,10}<2^{-1367}.       \tag{RC8}
\] Failed neighboring upper expressions are not lower bounds on the actual sums.

\subsection{5. Why the support-ten rank-one remainder is unipotent}\label{why-the-support-ten-rank-one-remainder-is-unipotent}

The restored K079(ER14--15), whose complete proof and script have been reread and whose certificate has been replayed separately in this recovery directory, removes every fixed-W0 support-ten term with rank\_E(h−I)≥2. This revalidation does not claim recovery of the unavailable later K079 bytes.

For clarity, the remaining rank-one case does not require extending DO067 to ten off-diagonal places. At each actual T incidence, h is an isomorphism between the two determinant-valuation-zero pencil lattices; equivalently h=A\_dKA\_c\^{}\{-1\}, with determinant-one shears A\_c,A\_d and integral invertible K. Thus det h has valuation zero there. At the other finite places h and its paired inverse-adjoint are integral, so det h is a global algebraic unit. Its relative norm is1. The already established CM root-unit fact makes det h an element of mu28.

If the original finite coefficient is nonzero, K079's finite Green-support argument gives unipotent reduction at every unmixed G place, hence det h≡1 there. Reduction is injective on mu28 at these primes: if a nonidentity root of order d\textbar28 reduced to1, the relation1+zeta+\ldots+zeta\^{}\{d−1\}=0 would reduce to the nonzero residue d, a contradiction. Therefore det h=1.

Write a rank-one N=h−I as u v\^{}t. The elementary determinant identity det(I+uv\^{}t)=1+v\^{}tu gives v\^{}tu=0, and hence N²=u(v\^{}tu)v\^{}t=0. Thus this rank-one case is globally unipotent. Rank zero is the identity and is absent from the nonintegral support-ten sector.

Consequently RC8 and the K079 cut remove the whole fixed-W0 support-ten sector. This is not a complete support-ten exhaustion of the full kernel: nonscalar positive blocks and scalar-cross matrices remain subject to their own bounds.

\subsection{6. Fresh error lists and exact scope}\label{fresh-error-lists-and-exact-scope}

Construct the new lists directly from the recovered K078 certificate; do not identify them with unavailable old file bytes.

For the uniform list append K079's32739 and RC8's1367: \[
 \mathcal E_{\rm new}=\mathcal E_{078}\mathbin{\|}(32739,1367).
\] For the all-top list, use the exact adopted field \nolinkurl{all_top_with_rank_one_removal_exponents} in the restored \nolinkurl{checks/k078_joint_remainder_error.json}. Its28-entry prefix already includes the older K077 globally unipotent rank-one supports10--14 union, bounded by2\^{}\{-1436\}, as its last charge. Append only K079's32739, preserving that historical prefix: \[
 \mathcal E_{{\rm top,new}}
   =\mathcal E_{{\rm top},078}^{\rm adopted}\mathbin{\|}(32739)
   =\mathcal E_{{\rm top},078}\mathbin{\|}(1436,32739).
\] The new uniform refinement does not require charging the same all-top subset again. Each removed sector is intersected with the current remainder before removal. Earlier numerical errors are retained conservatively.

The fresh certificate saves the full25-term uniform and29-term all-top arrays, verifies the adopted prefix, and checks both exact sums are\textless2\^{}\{-230\}. The baseline prefix is read from the actual recovered check file; the appended reconstruction result is not asserted to recover unavailable former K080 bytes. Later all-top or uniform cuts require their own subsequently retained charges.

The original signed remainder, all-T Mram positivity, global period, literal Host/FamP constructor and full DP44 comparison remain open. The allowed all-top role assignment and the actual mixed CM host are unchanged. K076's metric count is a separate optional bound; no norm or denominator assertion from it has been silently inserted into RC4.

\subsection{7. Sources and recovery provenance}\label{sources-and-recovery-provenance}

The exact field and root units are the restored K032/K075(RN1) and accepted K050 fact. The original algebraic content and relative norm counts are K075(RN11--19). The affine pencil, exact difference ideal, full type polynomial and convex endpoint argument are K077(AP1--23). The numerical inputs are the original DO040 and K044 JSON certificates, whose hashes the new certificate records. The complete source list and all current hashes are also recorded.

No external analytic theorem is rederived. The selected norm and homogeneous energies remain those already audited in K060/K072, with primary input:

Haagerup, U., \& Schlichtkrull, H. (2014). Inequalities for Jacobi polynomials. \emph{The Ramanujan Journal, 33}, 227--246. \url{https://doi.org/10.1007/s11139-013-9472-4}.

The present new reconstruction proves RC1--7 directly from the recovered inputs and rechecks the resulting exact arithmetic. Its certificate is not an enumeration of the global rational support or a proof of signed-period positivity.

\clearpage

\section{AM-R-048-RC-20261122 --- reconstructed all-top small cutoff and complete support-ten removal}\label{am-r-048-rc-20261122-reconstructed-all-top-small-cutoff-and-complete-support-ten-removal}

\textbf{Canonical integration --- 2026-11-22T16:10:03.402339+00:00.} Whole root proof/script reading and separate exact replay are recorded in CHECK\_EVENTS.tsv. Original recovery bytes remain preserved.

Author: GPT-6 Astra Pro, ALP\_MACHINE /root. 22 November 2026 UTC.

This is a fresh reconstruction from checkpoint23 source files. The former post-checkpoint proof bytes and their historical hash were not recovered. The retained mathematical review records remain historical records; this text has its own date, certificate and whole-proof checks. All uses below are confined to the already allowed all-top J56 assignment at the actual K044 upper gap endpoint.

\subsection{RC48.1. Actual data and what changes}\label{rc48.1.-actual-data-and-what-changes}

Retain K062's original normalized signed kernel, mixed CM host, finite projectors, class weights, canonical twists, selected compact module with its actual PBW metric, and packet point. Put \[
 a=f=168,\quad n=169,\quad d=f(a+1)^2,\quad
 \ell=\ell_+,\quad E=E_{\mathrm{top}},\quad j_*=\lfloor\ell/1024\rfloor.
\] The integers ell and E are read from the actual J56 row of K044's enlarged-gap certificate; T,g0,g1 are read from DO040's actual prime certificate. No profile is extrapolated beyond that admissible interval. At each real label the zero-twist starting factors and Pieri increments are \[
 r_i=E+(4a-3-4i)\ell+1\ (0\le i<a-1),\qquad r_{a-1}=2\ell+1,
\] \[
 m_0=2^{960}+(a-2)/2-4(a-1)\ell,\qquad
 m_i=4\ell-1\ (1\le i<a-1),\qquad
 m_{a-1}=E-2^{955}-\ell-a/2.
\] The last increment is12. The shifts are the actual all-top role assignment, not a relabeling of the uniform thick/thin profile. We change the tail proof's moment order and cutoff; the signed kernel and its normalization remain the same.

\subsection{RC48.2. All required derivative moments}\label{rc48.2.-all-required-derivative-moments}

Use the proved exact shifted-metric radial moment identity DO059(SC3): \[
 M_\tau(j)=\prod_i\frac{(r_i+m_i-j)_j}{(r_i-1-j)_j}.
\] K061(CG8), with its one-box compact tensor argument, proves that every compact type reached by a Lie word of length at most d has increments between0 and m\_i+2d and unchanged starting factors. Distinct compact types are orthogonal for the shifted compact-invariant metric. Consequently the largest type bound controls their sum without a dimension or type-count factor. Increasing a canonical twist adds the same nonnegative amount to numerator and denominator and decreases each ratio; an allowed interior one-bit decreases the relevant increment. Thus the following zero-word, zero-twist bound, including the full2d allowance, controls every required word and permitted twist: \[
 A_*:=\prod_{i=0}^{a-1}\frac{r_i+m_i+2d}{r_i-1-j_*}<\frac{33017}{1000}.
 \tag{RC48.1}
\] The denominators are positive: j\_\emph{\textgreater a+1 and j\_}\textless min\_i(r\_i-1), verified from the exact source integers. For every integer0\textless=j\textless=j\_\emph{, the moment is at most A\_}\^{}j, by bounding each Pochhammer numerator by r\_i+m\_i+2d and denominator below by r\_i-1-j\_*.

Set \[
 \epsilon=4/133,\quad R=61/25,\quad R'=2439/1000,\quad H=133.
\] The exact rational inequalities are \[
 \epsilon A_*<255/256<2^{-1/256},\qquad \epsilon^{-1}<2^{11},
 \tag{RC48.2}
\] where the second strict comparison follows from2*255\^{}256\textless256\^{}256. The certificate verifies the rational product in RC48.1, not a floating-point approximation.

\subsection{RC48.3. The full absolute derivative tail}\label{rc48.3.-the-full-absolute-derivative-tail}

Retain the Hilbert--Schmidt averages B\_I,b\_I, formal-degree constant C and raw Haar measure of K061(CG7,CG10) and DO059(SC8--10). For a word of length q\textless=d, \[
 \|d\pi(X_I)|_W\|\le2^{1000q},\quad C<2^{960af},\quad
 \int\prod_\tau Y_\tau^{a+1}\,dg=(a!)^{-f}.
\] These are unchanged representation and sampling inputs. With s0=a+1, the elementary pointwise inequality \[
 1_{Y_\tau<\epsilon}Y_\tau^{-s_0}
 \le\epsilon^{j_*-s_0}Y_\tau^{-j_*}
\] and weighted Cauchy--Schwarz give exactly the calculation of DO059(SC11--13), now using RC48.2: \[
 \int_{Y_\tau<\epsilon}b_I
 <\frac{2^{1000q}}C\,2^{B_*/2-j_*/512},\qquad
 B_*=960af+11f(a+1).
 \tag{RC48.3}
\] Indeed the squared-weight expression is bounded by C epsilon\^{}(-s0)(epsilon A\_\emph{)\^{}j\_} 2\^{}(11(f-1)s0), which is smaller than2\^{}(B\_\emph{-j\_}/256). This retains the C factor when passing from squared norm to L1.

The actual integers satisfy \[
 B_*/2-j_*/512<-\ell/2^{20}.
 \tag{RC48.4}
\] Moreover cosh²(R')\textgreater133/4=epsilon\^{}-1. For a fully rational check, the exponential lower series at2R'=2439/500 gives (exp(2R')+2)/4\textgreater133/4; the omitted exp(-2R')/4 is positive. Therefore outside the product Cartan set C\_R', at least one Y\_tau\textless epsilon. Sum RC48.3 over f labels and put delta=f*2\^{}(-ell/2\^{}20). Then \[
 \int_{G\setminus C_{R'}}b_I<2^{1000q}\delta/C.
 \tag{RC48.5}
\] The full L1 estimate remains \[
 \int_Gb_I\le2^{1000q}G_0,\qquad
 G_0=C^{-1/2}(a!)^{-f/2}2^{11f(a+1)/2}.
\] No selected representation dimension is inserted: the division by sqrt(dim W) is part of b\_I throughout.

\subsection{RC48.4. Transfer to the complete rational support}\label{rc48.4.-transfer-to-the-complete-rational-support}

K062(AT18--20) and DO059(SC18--20) already prove the disjoint sampling chart and normalized convolution estimate for the entire actual finite support, without assuming that support is a group. Their constants are \[
 \rho_T=(100d^2T^4)^{-1},\qquad K_T=2^{1200d}T^{8d}.
\] Each chart displaces the Cartan radius by less than1/1000=R-R'. Thus every chart with center outside C\_R lies outside C\_R', and RC48.5 applies to all its derivatives. The exact word allowance, Haar Jacobian, Casimir estimate and number of words have already been included in K\_T. The unchanged packet centers have radius below1/200, so the normalized convolution argument gives \[
 \sum_{\substack{g\in Z\backslash\mathscr C\\
                 \max_\tau\log\|g_\tau\|>2R+1/100}}
 |c(g)|<2K_TG_0\delta<2^{-2000}.
 \tag{RC48.6}
\] Multiplication by the original bounded finite normalized coefficient preserves this estimate, as in DO059(SC19). This controls the infinite absolute tail, not a finite numerical truncation. The sufficient exact inequality is \[
 \ell/2^{20}>2000+1+9208d+\lceil\log_2 f\rceil+11f(a+1)/2,
 \tag{RC48.7}
\] using T\textless2\^{}1001, C\textgreater=1 and (a!)\^{}f\textgreater=1. It is checked from the actual ell.

Finally2R+1/100=489/100 and exp(489/100)\textless133. The upper exponential series with a geometric remainder verifies this rationally. Hence every matrix with max\textbar\textbar g\_tau\textbar\textbar\textgreater133 is covered by RC48.6. Intersecting the annulus133\textless max\textbar\textbar g\_tau\textbar\textbar\textless=8192 with the current remainder removes it with the new charge2000. All earlier errors, including the older complete tail and removed integral terms, are retained conservatively. This establishes H133 only for the permitted all-top assignment.

\subsection{RC48.5. Global unipotence of the remaining support-ten matrices}\label{rc48.5.-global-unipotence-of-the-remaining-support-ten-matrices}

At support s=10, the actual K069 F chart thresholds remove r=rank(A-I)\textgreater=2 (their minimum threshold is at least22), and the DO064 exceptional partial-monomial chart's r=2 threshold is10 inclusive. Both chart families are used. Therefore every residual support-ten term has r\textless=1.

Use DO067(ST2--4): there is an actual positive fixed subspace in S whose signature-one orthogonal complement has dimension k=r+6\textless=7, padded by fixed positive directions when necessary. Restriction does not enlarge the H133 operator norm. Its first two coefficients agree with those of N=h-I. The same proved singular-value estimates give \[
 B_1(k)=H+2k-2+H^{-1},
\] \[
 B_2(k)=(2k-3)(H+H^{-1})+\binom{k-2}{2}+1+(k-1)(k-2)+\binom k2.
 \tag{RC48.8}
\] At H133 the exact values are

\begin{longtable}[]{@{}rrrr@{}}
\toprule\noalign{}
k & B1 & B2 & strict inverse binary margins at t=10 \\
\midrule\noalign{}
\endhead
\bottomrule\noalign{}
\endlastfoot
6 & 19020/133 & 164796/133 & 1157,110 \\
7 & 19286/133 & 202836/133 & 1150,9 \\
\end{longtable}

Here t counts type-III off-diagonal incidences. DO067(ST3) and K070(PT5--8) place the only poles of each original coefficient at the2t paired CM primes, each of norm T with pole order at most1. Every unmixed G reduction is unipotent on the accepted residual, so a nonzero coefficient would force \[
 Q^{2(f-1)}\le B_j(k)^{2f}T^{2t},\qquad Q=g_0g_1.
 \tag{RC48.9}
\] The displayed exact integer comparisons contradict this at t=10 and therefore at every t\textless=10. Thus e1(N)=e2(N)=0. Newton's identity gives tr(h)=tr(h²)=169. The accepted DO061 signature-one two-trace theorem, applied at a faithful embedding, gives global unipotence, N³=0 and rank\_E N\textless=2. This is a new numerical application of that theorem at the smaller actual cutoff; it does not extend the failed H8192/t10 test by assertion.

\subsection{RC48.6. Exhaust all structural cases and retain every cost}\label{rc48.6.-exhaust-all-structural-cases-and-retain-every-cost}

For fixed-W0 matrices the globally unipotent rank-two support-ten subset is already removed by K074, and the all-top globally unipotent rank-one subsets at supports10--14 are removed by K077 with their retained charge1436. Rank zero is the identity and has no nonintegral support-ten term. These uses are of precisely the globally unipotent subsets just established, not a prior blanket support-ten claim.

For the remaining globally unipotent matrices, DO065's classification and counts leave exactly its F classes (full rank,block rank)=(1,1),(2,1), or the A=I scalar-cross class of DO066. The complete bounds in DO065 Sections3--4 and DO066(GC22--26) depend on s and on v=u+t only after global unipotence is assumed. Their originally stated t\textless=8/9 premise supplied global unipotence; it is not a further restriction on the counting proof. Their full support/type/pair/orientation/fibre factors are already included.

At s=10 the three conservative exponents are \[
 -103530801/128+8966\cdot10,
\] \[
 -408071827/256+14941\cdot10,
\] \[
 -14033453/4+10956\cdot10+332665\cdot10.
 \tag{RC48.10}
\] The scalar-cross expression increases in v, so v\textless=10 is covered by its endpoint. Each displayed exponent is strictly below-4096. Charge the intersections of these three complete classes with the current remainder afresh, at total less than3*2\^{}-4096\textless2\^{}-4094. No assertion that every t\textless=9 term had already been removed is needed.

Thus complete support10 is removed in the all-top H133 kernel. The new all-top charges are exactly2000 and4094. Read the complete preceding K080 reconstructed list, retain its activated K077 charge1436 and K079 charge32739, and append these two new exponents. Its31 entries have exact sum below2\^{}-230. The uniform25-term list is unchanged. These count lists have new reconstruction provenance; their numerical equality does not recover a historical lost JSON byte sequence.

The resulting all-top remainder has support at least11. Its original phases, finite projector, selected trace, mass normalizer and actual period interpretation remain. The inequality S\_R\textgreater-1+2\^{}-230 is still unproved. This theorem supplies neither an Mram value nor a shorter admissible lift chain nor an actual DP44 sign, and has no RH conclusion.

\subsection{RC48.7. Location in the arithmetic and geometric programme}\label{rc48.7.-location-in-the-arithmetic-and-geometric-programme}

This is a quantitative statement on the existing Hermitian modules over O\_E, their paired congruence conditions, and their real automorphic realizations. The algebraic kernel data remain the same integral-module and linear-stabilizer data used in the run's SGA3 I, Expose I,4.7.3--4.7.3.1 dictionary; the analytic addition is a radial-moment and lattice-count estimate. It changes the finite nonvanishing gate attached to DP\_MULT Q3e and THETA3 G1, with A7's chosen definite CM types retained. It supplies no new comparison in lambda-rings or K-theory, no new l-adic sheaf or weight theorem, and no motivic period or L-function identity. The formal-group/Dieudonne conclusions at Grothendieck cote4 and cote54 scans5--7, the SGA7 I ExposesVII--VIII biextension layer, and FGA Bourbaki232/236 Picard layer remain the previously cited background. T-052's oscillator and CX-ALG-024's Kummer character gate are unchanged.

The proof uses the settled actual statements K044, K046(BL8--22), K061(CG8--20), K062(AT18--24), DO059(SC1--20), K069(IC4 and its F threshold table), DO064's IC8 threshold table, K070(PT5--8), DO067(ST2--8), DO061's two-trace theorem, DO065 Sections2--4, DO066(GC22--26), K074 and K077. No Lean build or validity audit of Alpoge's theorem is performed. The companion certificate verifies actual integers, all new moment and cutoff inequalities, both coefficient comparisons for every t\textless=10, the three full-cell exponents and complete error arrays. Its computation supplements this proof; it does not prove the cited analytic or geometric interfaces automatically.

\clearpage

\section{AM-R-049 --- complete uniform support-ten removal}\label{am-r-049-complete-uniform-support-ten-removal}

\textbf{Canonical checked recovery version --- 2026-11-22T16:10:03.402339+00:00.} Root has read the complete mathematical proof and companion script and replayed the stated exact certificate. The original recovery export and its historical status are preserved separately. This is a current checked version, not recovery of unavailable historical final bytes. Only the stated sector and hypotheses are promoted; the signed period and actual DP44 cost remain open.

\subsection{1. Exact scope and remaining cases}\label{exact-scope-and-remaining-cases}

Keep the complete actual J56 upper-gap signed kernel, all role assignments, canonical twists, original near-origin point, finite projectors, pencil incidences, central representatives and masses of K062/DO059/K080. Use a=f=168, n=169, m=164, H=8192 and E\^{}169=S perpendicular W0. Put N=h-I, X=A-I for A=h\_SS, and r=rank X. The original normalized coefficient is w(h)=a(h)c(h) product rho(h\_v).

At support s=10 the carried K068/K069 and DO064 tables have already removed every r\textgreater=2: both the fixed-space class F and its IC8 complement are needed, with the latter's r2 threshold10 inclusive. Thus r\textless=1. For r\textless=1, DO067 supplies global unipotence whenever t\textless=9, where t is precisely the type-III off-diagonal support. A nonunipotent remainder at s10 must therefore have t=s=10. Every T place in this stratum is type III. This is the only new nonunipotent support used below.

The two structural classes F and IC8 are those of K069. When r=0, F is exactly the fixed-W0 class, already wholly removed at s10 by K079/K080. When r=1, IC8 has the partial-monomial positive block classified in DO064. Both r1 classes are treated explicitly.

\subsection{2. A common rank-two coefficient witness}\label{a-common-rank-two-coefficient-witness}

The following determinant observation was supplied by the Kummer branch during its K081 analysis. It is used here with the unchanged full matrix. At an unmixed G prime all entries of N are integral and its reduction is nilpotent. If its reduction rank is d=0 or1, every j-minor, j\textgreater=2, has valuation at least j-d, by the ordinary determinantal-ideal valuation bound. In particular e\_j(N) has valuation at least2-d. Since det h=det(I+N)=1, one has e1(N)=-sum\_\{j\textgreater=2\}e\_j(N). Thus BOTH e1 and e2 have valuation at least2-d.~This holds at both CM partners.

If h is not globally unipotent, the accepted signature-one two-trace theorem DO061 implies that e1 or e2 is nonzero. Choose one such coefficient b, separately for the matrix; this does not multiply a count. Its only possible poles are simple poles at both CM primes of the t type-III places, by K070/DO067. There is no contribution from typeI/II.

K069 supplies a positive fixed subspace of codimension at most r+6. For r\textless=1 the original coefficient bounds of DO067 ST2--6 therefore apply in dimension k\textless=7. At H8192 both coefficient bounds are less than B=2\^{}17: B1\textless=H+12+H\^{}-1, B2\textless=11(H+H\^{}-1)+62. No moving integral basis is chosen.

At the334 unmixed G factors, let P1 be the product of residue sizes where rank(N)=1, and P0 the corresponding product where rank(N)=0. All other factors have rank at least2. The original Green formula gives \textbar a(h)\textbar{} \textless= 2 Q\footnote{-335(f-1)} P1\^{}167 P0\^{}335. The factor2 is the single global Euler correction; the two mixed projectors retain modulus at most1. Norm divisibility of the chosen b at BOTH CM partners gives (P1 P0\^{}2)\^{}2 \textless= \textbar Norm\_E b\textbar{} T\^{}(2t) \textless{} B\^{}(2f) T\^{}(2t). Consequently \textbar a(h)\textbar{} \textless{} 2 B\^{}(335f/2) T\^{}(335t/2) / Q\footnote{335(f-1)}. This uses the original normalization and assumes no global unipotence.

\subsection{3. The fixed-space class with r1}\label{the-fixed-space-class-with-r1}

For a fixed r1 integral block in F, carry K069 IC11--15 exactly. Its complete count is m\^{}2 (2H+3)\footnote{2f(2m-1)} binom(f,s)(T-1)\^{}(2s) G\_n (2H+1)\^{}(14f) T\^{}(14s). The seven recorded original coordinates determine the full last column; no inverse minor is declared integral. The source-lattice stabilizer remains G\_n\textless2\^{}5022.

Multiply this count by the preceding nonunipotent finite coefficient, by \textbar c\textbar\textless=1 and by the original rho bound {[}2/(T-2){]}\^{}s. Use the previously checked integer bounds log2(2H+3)\textless14+1/2048, log2(2H+1)\textless14+1/8192, m\^{}2\textless2\^{}15, binom(f,s)\textless=2\^{}168, Q\^{}4\textgreater2\^{}281, T\textless2\^{}996, and 2(T-1)\^{}2/(T-2)\textless2\^{}997. At s=t=10 the complete strict exponent is 168+15+5022+2f(2m-1)(14+1/2048)+14f(14+1/8192) +14941\emph{10+1+(335f/2)}17+(335/2)\emph{996}10 -(281/4)\emph{335}167 = -472271376/8192 \textless{} -4096. Thus the entire nonunipotent F,r1 support-ten cell has contribution less than2\^{}-4096.

\subsection{4. A partial-monomial positive block forces finite rank at least two}\label{a-partial-monomial-positive-block-forces-finite-rank-at-least-two}

This refinement was observed by the Kummer branch while checking the proposed IC8 count. It removes the need for a new IC8 lattice count.

For an IC8 term with A!=I, DO064 gives a partial-monomial A whose nonzero entries lie in mu28. If an unmixed G reduction of N had rank0, then A would reduce to I, contrary to the preserved positive-block rank. If it had rank1, its nilpotence would imply N\^{}2=0 there, hence h\^{}-1=I-N. The SS block of the actual unitary inverse is A\^{}dagger, so reduction of A+A\^{}dagger-2I would be zero at that prime.

But a partial-monomial A!=I has a diagonal entry equal either to0 or to a nontrivial eta in mu28. The corresponding diagonal of that Hermitian sum is -2 or eta+eta\^{}-1-2=(eta-1)\^{}2/eta. Both are units at the actual odd G primes, which are prime to28. The second assertion uses the already checked injectivity of reduction on mu28. This is a contradiction. Thus every such G rank is at least2.

For r1, the whole DO064 PM19 column count may now use L2=335, with no rank-loss numerator or pattern factor. Its unchanged conservative exponent at s10 is 329267+337645\emph{10-(281/4)}167*335 = -897677/4 \textless{} -4096. This bounds the entire IC8,r1 support-ten cell, independently of global unipotence. No enumeration or change of A is needed.

\subsection{5. The nonunipotent scalar-cross class}\label{the-nonunipotent-scalar-cross-class}

Now A=I and h does not fix S. Its last-column S part x=(h e\_n)\_S is nonzero. To justify this without a unipotent normal form, the unitary block identities give C J\_W C\^{}*=0. Thus the row space of nonzero C is a totally isotropic line in W0, whose negative coordinate is nonzero. Also C=0 would force B=0 by the opposite block identities. Hence x!=0.

Use the exact common-orientation argument of DO066 GC5--8 at a finite rank-one G reduction. That argument needs only finite nilpotence, rank one, and N\_SS=0; all hold here. At rank1, write N=bc\^{}t. Then b\_S c\_S\^{}t=0, so the entire x column vanishes at one CM partner; the paired inverse=I-N proves the other orientation. At rank0 it vanishes at both partners. Treat every rank\textgreater=2 reduction as one class. There are at most4\^{}334 patterns, including either allowed rank1 orientation by a fixed rule.

For a fixed pattern let J be the actual integral ideal containing the chosen one prime at each rank1 factor and both primes at rank0. Then Norm J=P1 P0\^{}2 and x belongs to (J I\_Sigma)\^{}m. Here Sigma consists of the two CM primes at all ten T places and Norm I\_Sigma=T\^{}-20. This cross-column containment follows directly from the actual target over-lattice; it does not extend a typeI unipotent integrality lemma.

Put L=J I\_Sigma and z=(Norm L)\footnote{1/(2f)}. The coefficient witness from section2 gives Norm J \textless{} B\^{}f T\^{}10, hence z\^{}(2f) \textless{} 2\^{}(17f) T\^{}-10 \textless{} 2\^{}(17f-9950) \textless{} 2\^{}(-42f). Therefore z\textless2\^{}-21.

The half-trace Euclidean lattice formulas, with the same conventions as K076, are covol(L)=D0 z\^{}(2f) and L\#=2 Different\_E\^{}-1 bar(L)\^{}-1, where D0=7\^{}140 29\^{}144 5\^{}84 13\^{}84. Arithmetic--geometric mean gives lambda1(L\#)\textgreater=2sqrt(f)/(rd(E)z). The accepted transference inequality in dimension2f then gives mu(L)\textless=rd(E)sqrt(f)z/2\textless750sqrt(f)z. Apply this scalar estimate to the orthogonal product of m copies: mu(L\^{}m)\textless750sqrt(fm)z.

For any matrix in U(a,1), the Euclidean norm of the positive part of its last column is sinh(log \textbar\textbar h\textbar\textbar). Thus at every embedding \textbar\textbar x\textbar\textbar\textless(H/2)=4096; no sqrt(H) unipotent bound is used. The trace ball has radius4096sqrt(f). Enlarging by the covering radius gives a radius less than4097sqrt(f), since sqrt(m)\textless13 and 9750*2\^{}-21\textless1.

Dirichlet-cell volume comparison in this actual lattice gives \#x \textless= V164 z\^{}(-2fm), where V164={[}((22/7) f 4097\^{}2)\^{}(fm){]}/{[}(fm)! D0\^{}m{]} \textless2\^{}281664. The exact integer comparison is part of the companion certificate. No unproved equidistribution density is used.

Write alpha=335/2. The finite rank-loss factor is at most (Norm J)\^{}alpha=T\^{}(335*10) z\^{}(335f). After multiplication by the column count its remaining z-power is (335-2m)f=7f=1176, positive. Therefore the weighted S count is less than V164 T\^{}(3350) 2\^{}(-24696).

Retain the unchanged W0 target-lattice count 2\^{}14861 T\^{}(10s), the full source-lattice fibre G\_n\textless2\^{}5022, all4\^{}334 G patterns, binom(f,s)\textless=2\^{}168, and the conservative complete pair/rho factor\textless2\^{}(997s). We do not need a favorable off-diagonal cancellation. The selected real factor is bounded only by1. With Q\^{}4\textgreater2\^{}281, the complete exponent at s10 is 168+5022+14861+1+668+281664-24696 +996\emph{(3350+100)+9970-(281/4)}335*167 = -825113/4 \textless{} -4096. Hence this entire nonunipotent scalar-cross cell has contribution less than2\^{}-4096. Every original sign and phase remains in the kernel.

\subsection{6. Exhaustion and the unchanged signed normalization}\label{exhaustion-and-the-unchanged-signed-normalization}

After the carried cuts, a remaining support-ten term has r\textless=1. The fixed-W0 sector is already removed by K079/K080. The nonunipotent terms are exactly among sections3--5, three cells each\textless2\^{}-4096. The unipotent nonscalar cells (rank N,r)=(1,1),(2,1) and the unipotent scalar-cross cell are the original DO065/DO066 counts. Their formulas depend on their stated global-unipotence premise, not on the way that premise was obtained. At s10 their three exponents are exactly those checked in R048-RC RC48.10, all\textless-4096. The three unipotent bounds cover every support-ten type profile. Only their intersections with the current remainder are removed.

Thus all six new cells together have absolute contribution \textless6*2\^{}-4096\textless2\^{}-4093. Intersect each with the current remainder before removing it. Append4093 to the K080 uniform error list (and retain any later separately recorded disjoint-sector charges). The complete error is still\textless2\^{}-230 by exact rational addition.

The uniform remainder now has s\textgreater=11 at the unchanged H8192. The R048 all-top remainder already has s\textgreater=11,H133, and its error list is retained; it needs no duplicate support-ten charge. The original normalized energy remains 1+S\_R+epsilon. Its sufficient signed inequality S\_R\textgreater-1+2\^{}-230 remains open.

This is a strict reduction of the same conditional nonvanishing gate. It does not produce the prescribed Mram period, a specific032 cycle, a literal source-family constructor, a shorter admissible lift chain, or the DP44 cost sign. The integral Hermitian correspondence, CM modules, exterior powers and all finite projectors remain the same. In the recorded SGA3 I4.7.3/4.7.3.1 module dictionary, the new determinantal-ideal congruences and the metric realization restrict the same arithmetic points over Spec(O\_E); no new motivic, biextension, lambda-ring or l-adic comparison is claimed.

\clearpage

\section{AM-K081 --- nonunipotent fixed-block coefficients retain rank-two Green loss}\label{am-k081-nonunipotent-fixed-block-coefficients-retain-rank-two-green-loss}

\textbf{Canonical checked recovery version --- 2026-11-22T16:10:03.402339+00:00.} Root has read the complete mathematical proof and companion script and replayed the stated exact certificate. The original recovery export and its historical status are preserved separately. This is a current checked version, not recovery of unavailable historical final bytes. Only the stated sector and hypotheses are promoted; the signed period and actual DP44 cost remain open.

\subsection{1. Actual sector and a nonzero two-coefficient witness}\label{actual-sector-and-a-nonzero-two-coefficient-witness}

Retain K079(ER1--3,11--13), K070(PT1--6), and the original K062 kernel. Thus a=f=168,n169,H8192, E\^{}n=S perpendicular W0,dim S164, W0 signature(4,1),Q=g0g1 with g0=22285932761,g1=64072175441. Both mixed G projectors are unchanged. Let h=I\_S direct-sum h0 be the unique central normalization of a nonintegral supported term, and put N=h−I. Assume a(h) is nonzero and h is not globally unipotent.

By the already checked DO060(IB14--15), each unmixed G reduction is unipotent and det h=1. The latter follows from actual local pencil volumes, the exact global norm-one unit group mu28, and injective reduction of mu28 at G. This conclusion has no small-support premise.

Write e\_j=e\_j(N). These equal the characteristic coefficients of h0−I, because N vanishes on S. If e1=e2=0, Newton gives tr(h0)=tr(h0²)=5. The accepted signature-one two-trace theorem DO061(section2) then makes h0 unipotent. Therefore at least one of e1,e2 is nonzero.

The fixed-coordinate singular values of h0 are R,1,1,1,R\^{}(-1),R≤H. The trace nuclear norm and the second exterior nuclear norm are bounded by H+H\^{}(-1)+3 and 3(H+H\^{}(-1))+4. Coefficient translation gives e2(h0−I)=e2(h0)−4tr(h0)+10. Consequently \textbar e1\_tau\textbar≤H+H\^{}(-1)+8, \textbar e2\_tau\textbar≤B:=7(H+H\^{}(-1))+26\textless2\^{}16. Choose one nonzero y among e1,e2. The choice does not multiply the count.

\subsection{2. The determinant identity improves the first coefficient too}\label{the-determinant-identity-improves-the-first-coefficient-too}

At an unmixed G component N is integral. If its reduction has rank d, the DVR minor bound used in K068 section3, the DVR-minor paragraph between LR9 and LR10 gives v(e\_j(N))≥(j−d)\emph{+. Since det(I+N)=1, e1(N)=−sum}(j=2)\^{}5 e\_j(N). For d=0 or1 every summand on the right has valuation at least2−d. Thus BOTH e1 and e2 have valuation at least(2−d)\emph{+ at every such component. At the conjugate component the reduction rank is the same, because h\_barP is inverse-adjoint. Therefore the chosen y obeys nu\_w:=v\_P(y)+v\_barP(y)≥2(2−d\_w)}+.

This is a determinant identity over the actual integral G-local ring, not a global unipotence assumption. Equivalently the trace lies in the ideal generated by the2-minors whenever det(I+N)=1.

Let L2=335,L1=168,L0=0. Rank0 has loss335 and rank1 has loss167 relative to the rank-two Green ratio; both losses are at most (335/4)nu\_w. For rank at least2 the same baseline ratio already applies. The exact Green denominator corrections over all334 unmixed factors have product less than2. Both mixed partial traces retain modulus≤1.

Let t count actual typeIII places. K070(PT1--5) proves that y is integral outside their2t CM primes, with pole order≤1 there. This statement is valid before global unipotence. If I\_III is the product of the inverse primes at both partners, then (y)I\_III\^{}(-1) is integral, Norm((y)I\_III\^{}(-1))≤B\^{}(2f)T\^{}(2t). Taking the product of the paired Green losses proves the pointwise bound \textbar a(h)\textbar≤2 B\^{}(335f/2) T\^{}(335t/2)/Q\^{}(335(f−1)). No finite rank-pattern multiplicity is needed. The original real normalized selected trace has modulus at most1; no unipotent compression estimate or small cross-column radius is used.

\subsection{3. Actual lattice, fibre and exact typeIII cancellation}\label{actual-lattice-fibre-and-exact-typeiii-cancellation}

Fix support Sigma of size s and actual ordered source/target pencil pairs. K079(ER11--13), whose lattice and fibre arguments require no unipotence, gives \#matrices≤G5 BW T\^{}(10s), G5=98148490017978252769652736000000\textless2\^{}107, BW={[}(22/7)f9543²{]}\^{}(5f)/{[}(5f)!D0\^{}5{]}\textless2\^{}14861, D0=7\^{}140 29\^{}144 5\^{}84 13\^{}84. The count uses the actual target over-lattice and the finite source lattice isometry fibre fixing its negative vector.

Choose the t typeIII positions inside Sigma. At such a place the unique incidence is off-diagonal: there are at most(T−1)(T−2) ordered pairs, and its actual absolute rho is exactly 1/{[}(T−1)(T−2){]}. Their weighted upper cost is1. At every other nonintegral place use at most(T−1)² pairs and the original modulus bound2/(T−2). Put P\_T=2(T−1)²/(T−2)\textless2\^{}997. Integral places retain only their contraction bound1. Therefore the full(s,t) nonunipotent fixed-W0 sector satisfies U\_(s,t)≤ 2 binom(f,s)binom(s,t)G5 BW B\^{}(335f/2) Q\^{}(−335(f−1)) T\^{}(10s+335t/2) P\_T\^{}(s−t).

This is an absolute upper bound for the original normalized signed sum. A matrix with several actual incidences may be counted repeatedly, which only enlarges the bound. No full-pencil C≤2 claim is made. No class mass, central phase or selected dimension is removed.

\subsection{4. Explicit conservative cells}\label{explicit-conservative-cells}

Use T\textless2\^{}996,Q\^{}4\textgreater2\^{}281,B\textless2\^{}16, the two constants above, binom(f,s)≤2\^{}168 and binom(s,t)≤2\^{}s. Direct simplification yields log2 U\_(s,t)\textless{} C+10958s+165833t, C=−13859037/4.

The constant is exactly 1+107+14861+168+16\emph{(335f/2)−(281/4)}335*(f−1). The s coefficient includes the binomial factor2\^{}s and the actual non-typeIII incidence cost; the t coefficient subtracts that cost before adding the pole-one Green loss.

For t10 through19 the largest s satisfying this expression≤−4096 is: t10:s164; t11:s149; t12:s134; t13:s119; t14:s103; t15:s88; t16:s73; t17:s58; t18:s43; t19:s28. Every integer s from t to its listed endpoint is included. There are 824 cells in total, hence their absolute union is less than2\^{}(-4086). The next t20 has no admissible s≥t from this conservative bound.

For fixed W0, the carried DO067 criterion already makes every term with t≤9 globally unipotent. Thus all nonunipotent fixed-W0 terms at total support10 through19 lie in the newly bounded cells. The same table also removes specified higher-support subsets. No assertion that the remaining larger-support upper estimates are sharp is made.

\subsection{5. Scope}\label{scope-1}

Each listed sector is intersected with the current remainder before it is removed. All previously charged numerical errors are conservatively retained; in particular fixed-W0 support10 may already have been removed by K079/K080 or by the later complete support-ten reduction.

The result removes only the explicitly defined nonunipotent fixed-W0 sector. It neither counts the globally unipotent rank-two remainder at support11 nor replaces the parent's separate nonscalar/cross-block analysis. It retains the unchanged H8192 family and therefore also bounds any restriction to the all-top H133 subset. No Mram sign, period evaluation, literal Host/FamP record or DP44 inequality follows from this reduction alone.

The new inputs are the determinant-one trace divisibility and its use with the characteristic pole-one support, followed by the actual weighted incidence count. Accepted Green, lattice, fibre, coefficient pole and two-trace statements are cited at their original interfaces. The final saved proof and full script have since been read and the13634-assertion exact certificate independently replayed; see the canonical integration record.

\clearpage

\section{AM-K082 --- all-top refinement of the nonunipotent fixed-W0 coefficient bound}\label{am-k082-all-top-refinement-of-the-nonunipotent-fixed-w0-coefficient-bound}

\textbf{Canonical whole-proof/script check and independent exact replay PASS --- 2026-11-22T16:10:03.402339+00:00.} The original outage export is retained separately.

Retain exactly K081's sector, characteristic-coefficient witness, paired Green estimate, actual lattice/fibre count, and typeIII incidence cancellation. Restrict the actual all-top sum to the R048 cutoff H=133. This changes no point, level, finite projector, class mass, pencil, canonical twist or selected coefficient; R048 separately charges the outer annulus. Every previous numerical error remains retained.

The same singular-value calculation now gives B=7(133+1/133)+26=127288/133\textless2\^{}10. The old G5\textless2\^{}107 and BW\textless2\^{}14861 bounds remain valid because the smaller-cutoff matrix set is a subset of the original H8192 set. Using \textbar c\textbar≤1 and the same prime inequalities, K081's expression becomes log2 U\_(s,t)\textless C10+10958s+165833t, C10=−14534397/4=−13859037/4−168840. No unipotent compression estimate is used.

For t=10,11,12,13,14,15,16,17,18,19,20 the greatest integer s≤168 satisfying the displayed bound≤−4096 is respectively 168,164,149,134,119,104,89,73,58,43,28. At t21 the smallest possible s21 already has exponent316047/4\textgreater0, and the expression is increasing in both variables, so no t≥21 qualifies by this estimate.

Relative to K081's table, the new intervals are: t10:s165--168; t11:s150--164; t12:s135--149; t13:s120--134; t14:s104--119; t15:s89--104; t16:s74--89; t17:s59--73; t18:s44--58; t19:s29--43; t20:s20--28. There are151 nominal cells. Their largest exponent is−20037/4 at(t,s)=(16,89). Thus each is\textless2\^{}-5009 and their full absolute union is\textless2\^{}-5001 because151\textless256. The t20 portion alone has9 cells and largest exponent−40461/4, so its union is\textless2\^{}-10111.

Each cell is intersected with the current R048/K081 remainder before removal; the nominal count does not assert that every intersection is nonempty. In particular the statement does not reintroduce previously removed support-ten or rank-restricted sectors. It newly removes every retained nonunipotent fixed-W0 term with t20 and20≤s≤28, plus the specified lower-t tails. The globally unipotent rank-two fixed-W0 sector is separate. The complete all-top remainder still has floor11: no exhaustion at support11 is asserted. The uniform H8192 family is unchanged.

Exact-integer arithmetic, including both B bounds, all table cells/endpoints and the151-cell difference, was independently executed in the functions runtime on22 November2026 at13:58:06UTC. Recovery has since completed; the whole proof and script were freshly read and the12931-assertion exact certificate independently replayed. See the dated canonical integration manifest.

\clearpage

\section{AM-R-050 --- further nonunipotent rank-one fixed-space sectors}\label{am-r-050-further-nonunipotent-rank-one-fixed-space-sectors}

\textbf{Canonical checked recovery version --- 2026-11-22T16:10:03.402339+00:00.} Root has read the complete mathematical proof and companion script and replayed the stated exact certificate. The original recovery export and its historical status are preserved separately. This is a current checked version, not recovery of unavailable historical final bytes. Only the stated sector and hypotheses are promoted; the signed period and actual DP44 cost remain open.

\subsection{1. Exact scope}\label{exact-scope}

Use the original actual J56 family, signed kernel, finite projectors, role assignments, incidence weights, original metric, lattice and normalization retained in K069, R048, R049 and K081. Set f=a=168, n=169, m=164. This note concerns only the fixed-space structural class F with r=rank(A-I)=1 and h nonunipotent. It does not concern the separate IC8 class or the fixed-W0 sector. Let s be the number of active T places and t the number among these that have type III. In the current remainder s\textgreater=11.

The uniform archimedean cutoff is H=8192; the all-top cutoff is H=133, after R048's already charged analytic annulus removal. All statements for H133 are all-top statements. No literal HostFamP constructor is supplied.

\subsection{2. Retain the actual complete count for a fixed pencil pair}\label{retain-the-actual-complete-count-for-a-fixed-pencil-pair}

For each fixed support, fixed type-III subset and fixed choice of its actual local pencil pairs, K069 IC11--15 and R049 section3 bound the F,r1 matrices, including the UNION of types I and II at every remaining active position, by M(H,s)=m\^{}2 (2H+3)\footnote{2f(2m-1)} G\_n (2H+1)\^{}(14f) T\^{}(14s). This is the seven-original-coordinate bound, with G\_n\textless2\^{}5022. Its target-overlattice count is uniform in the chosen pencil pairs. For fixed A, choose two original S coordinates injective on the relevant complement, adding an unused original coordinate if only one is needed, and retain all five original W0 coordinates. The local shear acts only in W0, so this projection commutes with the shear. The projected pi\^{}-1 M\_d contains O\_E\^{}7 with paired index at most T\^{}14 per active place, uniformly in d.~This counts I and II together and introduces no extra factor2\^{}(s-t). No inverse minor or change of basis is declared integral.

For nonunipotent h, the determinant-one coefficient witness of R049 section2 applies for every s,t in this note, not just s10. Namely at every unmixed G prime of reduction rank d=0 or1, both e1(N) and e2(N) have valuation at least2-d.~This follows from the higher determinantal ideals and det(I+N)=1. At least one coefficient is nonzero by the accepted signature-one two-trace theorem. Each has only simple poles at the two CM primes of each type-III place. The fixed positive complement has dimension at most7. If B bounds both coefficients at every infinity, the original finite Green coefficient therefore satisfies \textbar a(h)\textbar{} \textless{} 2 B\^{}(335f/2) T\^{}(335t/2) Q\footnote{-335(f-1)}. There is no rank-pattern count and no global-unipotence assumption.

\subsection{3. Exact type-III pencil cancellation}\label{exact-type-iii-pencil-cancellation}

Choose the support and type-III positions in binom(f,s)binom(s,t) ways. At a type-III position, the number of possible ordered local pencil pairs is at most (T-1)(T-2), while the absolute incidence weight is exactly 1/{[}(T-1)(T-2){]}. Their complete contribution to this pair count is at most1. At each other active position retain the conservative full bound P\_T=2(T-1)\^{}2/(T-2)\textless2\^{}997. The matrix bound M(H,s) is uniform in these pairs, so it is legitimate to apply it before summing the exact local pair/weight products. The real factor satisfies \textbar c(h)\textbar\textless=1.

Hence the complete absolute contribution U\_(s,t) of this cell is bounded by 2 binom(f,s)binom(s,t) m\^{}2 G\_n (2H+3)\footnote{2f(2m-1)} (2H+1)\^{}(14f) B\^{}(335f/2) T\footnote{14s+(335/2)t} Q\footnote{-335(f-1)} P\_T\^{}(s-t). This retains every incidence factor of the actual original kernel.

\subsection{4. Uniform sectors}\label{uniform-sectors}

At H8192 use B=2\^{}17, the checked inequalities log2(2H+3)\textless14+1/2048, log2(2H+1)\textless14+1/8192, m\^{}2\textless2\^{}15, binom(f,s)\textless=2\^{}168, binom(s,t)\textless=2\^{}s, G\_n\textless2\^{}5022, T\textless2\^{}996, Q\^{}4\textgreater2\^{}281 and P\_T\textless2\^{}997. Then log2 U\_(s,t) \textless{} C\_U +14942s+165833t, C\_U=-15362951696/8192. The slopes are 14\emph{996+997+1=14942 and 335}996/2-997=165833.

For t10, s10 through14 satisfy a strict bound below-4096. The s10 cell is already charged in R049 and is not removed again. The four new cells have s=11,12,13,14 and t=10. Their exact upper-bound numerators with denominator8192 are, respectively, -431458832, -309053968, -186649104, -64244240. Thus their union has contribution strictly less than2\^{}-4094. No t\textgreater=11 has any qualifying s\textgreater=t for this particular upper bound. This last assertion is only a limitation of the displayed estimate.

\subsection{5. All-top sectors}\label{all-top-sectors}

At H133, the original coefficient bounds B1,B2 in dimensions k\textless=7 are both less than B=1536. The largest is B2=202836/133\textless1536; also 1536\^{}5\textless2\^{}53, so log2 B\textless53/5. Use the exact rational logarithm bounds 269\^{}4096 \textless{} 2\^{}(8\emph{4096+293), 267\^{}4096 \textless{} 2\^{}(8}4096+249), and retain the same actual G\_n,T,Q and all finite normalizations. This gives log2 U\_(s,t) \textless{} C\_A +14942s+165833t, C\_A=-22289217728/8192.

On s\textgreater=11 the complete qualifying table is t=10, 11\textless=s\textless=70, 60 cells; t=11, 11\textless=s\textless=59, 49 cells; t=12, 12\textless=s\textless=48, 37 cells; t=13, 13\textless=s\textless=37, 25 cells; t=14, 14\textless=s\textless=26, 13 cells; t=15, 15\textless=s\textless=15, 1 cell. The endpoint upper-bound numerators with denominator8192 are -135837888,-123787456,-111737024,-99686592,-87636160,-75585728. They are all strictly below-4096*8192. The two positive slopes prove the corresponding bounds throughout every listed interval. For t\textgreater=16 and s\textgreater=t, the same displayed estimate gives no qualifying cell. There are185 cells, so their union contributes less than2\^{}-4088. The four uniform cells are included in this all-top union; charge only 4088 to the all-top list, with no extra duplicate4094 charge.

\subsection{6. Remainder and limits}\label{remainder-and-limits}

Intersect every new sector with the actual current remainder before removing it. Append4094 to the uniform error exponents and4088 to the all-top error exponents, retaining every previous error including the R048, R049, K081 and K082 charges when present. The companion exact rational check must read the current complete lists before final promotion.

The support floor remains11. The sufficient signed inequality for the same normalized energy is still S\_R\textgreater-1+2\^{}-230. This note strictly reduces its remaining F,r1 nonunipotent terms, but proves neither that inequality, a prescribed global Bessel/Mram period, a specific032 cycle value, nor an actual DP44 cost sign. There is no new integral/model identification or change to the existing SGA3 I4.7.3/4.7.3.1 module dictionary.

\clearpage

\section{AM-K083 --- a trace witness for low-rank nonunipotent fixed blocks}\label{am-k083-a-trace-witness-for-low-rank-nonunipotent-fixed-blocks}

\textbf{Canonical checked recovery version --- 2026-11-22T16:10:03.402339+00:00.} Root has read the complete mathematical proof and companion script and replayed the stated exact certificate. The original recovery export and its historical status are preserved separately. This is a current checked version, not recovery of unavailable historical final bytes. Only the stated sector and hypotheses are promoted; the signed period and actual DP44 cost remain open.

\subsection{1. Actual sector and statement}\label{actual-sector-and-statement}

Retain the complete actual fixed-W0 sector of K081: h=I\_S⊕h0, dimW0=5, signature(4,1), det h=1, N=h−I, nonzero normalized finite coefficient, and h nonunipotent. Impose the additional global E-rank condition rankN≤3. All original points, pencils, finite projectors, central phases, class masses and real normalizations remain unchanged.

The new witness is y=trN itself. It is nonzero and satisfies \textbar τy\textbar≤B3(H):=H+H\^{}-1+4 at every complex realization. This uses no moving integral basis. Insert this witness into the already proved determinant/Green/count argument of K081, retaining its original G5 and BW and \textbar c(h)\textbar≤1.

\subsection{2. Self-reciprocity and the nonzero trace}\label{self-reciprocity-and-the-nonzero-trace}

Let χ\_h(z)=(z−1)\^{}(169−d)P(z), with P(1)≠0 and P monic. Here d is the total algebraic multiplicity of eigenvalues different from1. On their full generalized eigenspaces N is invertible, hence d≤rankN≤3. If d=0, h is unipotent, which is excluded.

Since det h=1, P(0)=(-1)\^{}d.~Unitarity gives χ\_h(z)=(-1)\^{}169 z\^{}169 c(χ\_h)(1/z), where c conjugates coefficients. Factoring out the same power of z−1 and cancelling gives P(z)=(-1)\^{}d z\^{}d c(P)(1/z). This is a polynomial identity over E, valid at every embedding. It does not require a semisimple matrix.

The case d=1 is impossible because its single eigenvalue has product1. For d=2, write P=z²−s z+1; the identity gives c(s)=s. For d=3 it gives P=z³−s z²+c(s)z−1. Since trN=s−d, vanishing of trN would give P=(z−1)\^{}d in either case, contradicting P(1)≠0. Thus y=trN≠0.

For clarity, the factor P includes eigenvalues of modulus1 different from1; no assertion that every selected eigenvalue lies off the unit circle is made.

\subsection{3. The spectral bound, including generalized multiplicities}\label{the-spectral-bound-including-generalized-multiplicities}

For a form-preserving operator, full generalized eigenspaces for λ and μ are orthogonal unless λ·barμ=1. This follows by induction on the two Jordan heights from invariance of the form: the pairing equation has nonzero scalar coefficient λ·barμ−1, and its remaining terms have smaller heights. Consequently the sum of all generalized eigenspaces with \textbar λ\textbar\textgreater1 is totally isotropic; signature(4,1) bounds its dimension by1. The contracting sum likewise has dimension≤1. Reciprocal-conjugate symmetry pairs the two. Thus, counted with algebraic multiplicity, at most one expanding eigenvalue R in modulus and one contracting eigenvalue R\^{}-1 occur; all other eigenvalues have modulus1.

The actual operator-norm cutoff gives 1≤R≤H. Therefore the sum of absolute values of the d eigenvalues in P is at most H+H\^{}-1+d−2; this also holds when all d eigenvalues have modulus1. Hence \textbar trN\textbar≤H+H\^{}-1+2d−2≤H+H\^{}-1+4. No coordinate denominators or integral orthogonal-complement claim enters this estimate.

At H8192, B3(H)=(H²+4H+1)/H\textless2\^{}(417/32). At the actual all-top R048 cutoff H133, B3(H)\textless2\^{}(57/8). Both strict inequalities are certified by integer powers.

\subsection{4. Exact cell bounds}\label{exact-cell-bounds}

K081, with the smaller trace bound and unchanged counting constants, gives log2 U\_(s,t)\textless C+10958s+165833t. For the uniform H8192 sector C=−28386399/8. For the all-top H133 sector C=−14858007/4. The actual typeIII coefficient poles and their incidence cancellation are unchanged. The trace lies in the2-minor ideal by det(I+N)=1, exactly as in K081.

Enumerate 10≤t≤s≤168 and retain every cell for which this expression is≤−240. Relative to K081's previous table, the uniform added intervals are: t10:s165--168; t11:s150--157; t12:s135--142; t13:s120--127; t14:s104--111; t15:s89--96; t16:s74--81; t17:s59--66; t18:s44--51; t19:s29--36; t20:s20--21. There are78 nominal cells. Their largest exponent is−6439/8 at(t,s)=(13,127), so each is\textless2\^{}-804. As78\textless128, their absolute union is\textless2\^{}-797.

Relative to K082's all-top table, the all-top added intervals are: t11:s165--168; t12:s150--157; t13:s135--142; t14:s120--127; t15:s105--111; t16:s90--96; t17:s74--81; t18:s59--66; t19:s44--51; t20:s29--36; t21:s21. There are75 nominal cells. Their largest exponent is−4695/4 at(t,s)=(14,127), so each is\textless2\^{}-1173. As75\textless128, their union is\textless2\^{}-1166.

Every uniform added cell already belongs to K082's all-top table. Thus append797 only for the new uniform remainder and1166 only for the new all-top remainder; no duplicate797 charge is needed on the latter. All earlier errors, including K081's4086 and K082's all-top5001, are conservatively retained. Each set is intersected with the current remainder before removal. The nominal numbers do not assert that all intersections are nonempty.

\subsection{5. Actual complete-sector consequences}\label{actual-complete-sector-consequences}

K079 already removes every fixed-W0 matrix of global rankN≥4 at support20 and21, without a unipotence assumption.

In the uniform sum, K081 removes every nonunipotent fixed-W0 term at support20 with t≤19. The new t20,s20 cell removes ranks≤3, and K079 removes ranks≥4. Therefore the complete nonunipotent fixed-W0 support20 sector is removed. At support21, the new t20,s21 cell is also removed; any remaining nonunipotent fixed-W0 term there must have t21 and global rank2 or3. No complete support21 assertion is made for the uniform sum.

In the all-top sum, K082 already removes every nonunipotent fixed-W0 support21 term with t≤20. The new t21,s21 cell removes ranks≤3, and K079 removes ranks≥4. Thus the complete all-top nonunipotent fixed-W0 support21 sector is removed.

These conclusions concern the nonunipotent fixed-W0 sector only. Globally unipotent rank-two terms and the other structural sectors are unchanged. No complete signed-support floor, Mram positivity, nonzero period, literal source constructor, DP44 estimate or RH statement follows.

\clearpage

\section{AM-R-051 --- all-top support-eleven column cuts}\label{am-r-051-all-top-support-eleven-column-cuts}

\textbf{Canonical checked recovery version --- 2026-11-22T16:10:03.402339+00:00.} Root has read the complete mathematical proof and companion script and replayed the stated exact certificate. The original recovery export and its historical status are preserved separately. This is a current checked version, not recovery of unavailable historical final bytes. Only the stated sector and hypotheses are promoted; the signed period and actual DP44 cost remain open.

\subsection{1. Scope and the actual target lattice}\label{scope-and-the-actual-target-lattice}

Retain the original actual J56 signed kernel and all normalizations. Set f=a=168,n=169,m=164,s=11,H=133. The cutoff is available only for the all-top family, through R048's charged annulus removal.

At a T place the actual source and target pencils have M\_c=H\_c Lambda, H\_c=I+c pi\^{}-1 zq, with (zq)\^{}2=0 supported in W0. Thus H\_c\^{}-1(pi e\_n)=pi e\_n-c zq(e\_n) is integral, so pi e\_n belongs to M\_c.~From hM\_c=M\_d it follows that h e\_n belongs to pi\^{}-1 M\_d. The W0 shear is the identity on S. At the conjugate CM place the same calculation uses bar(pi) and the dual W0 shear. Consequently the S part x=(h e\_n)\_S lies in I\_Sigma\^{}m, where Norm I\_Sigma=T\^{}(-2s). This statement holds for every active pair and all types I,II,III, without global unipotence or a type-I integrality assertion.

The W0 projection pi\^{}-1 M\_(d,W0) contains O\_E\^{}5 and has paired index T\^{}10 at each active place. Its original-coordinate coset count is at most (2H+1)\^{}(10f)T\^{}(10s)\textless2\^{}13543 T\^{}(10s)\textless2\^{}14861 T\^{}(10s). Here the scalar packing argument in each O\_E coset uses lambda1(O\_E)\textgreater=sqrt(f) in the half-trace metric. It is uniform in all pencil pairs and does not depend on A=h\_SS. For each fixed pair and full negative column, two arrows differ by an automorphism of the same source lattice fixing e\_n.~The accepted finite stabilizer bound is G\_n\textless2\^{}5022, also independent of A. Thus no positive-block enumeration is necessary.

\subsection{2. Actual scalar-ideal volume}\label{actual-scalar-ideal-volume}

Use the same field and half-trace lattice normalization as K076/R049: D0=7\^{}140 29\^{}144 5\^{}84 13\^{}84, covol(L)=D0 z\^{}(2f), z=(Norm L)\^{}(1/(2f)), L\#=2 Different\_E\^{}-1 bar(L)\^{}-1, rd(E)\textless1500. The accepted arithmetic-geometric-mean and transference argument gives mu(L\^{}m)\textless750 sqrt(fm) z\textless9750 sqrt(f) z.

The positive part of a negative unit column of a U(a,1) matrix has norm sinh(log\textbar\textbar h\textbar\textbar)\textless H/2 at every embedding. In particular \textbar\textbar x\textbar\textbar\textless66.5. Whenever z\textless2\^{}-27, adjoining a Dirichlet cell enlarges the trace ball to radius less than67 sqrt(f), since9750*2\^{}-27\textless1/2. The actual Voronoi comparison is therefore \# \{x in L\^{}m in this ball\} \textless= V\_plus z\^{}(-2fm), V\_plus={[}((22/7) f 67\^{}2)\^{}(fm){]}/{[}(fm)!D0\^{}m{]} \textless2\^{}-45338. The rational comparison constant V\_plus is also\textgreater2\^{}-45339, as checked by exact integer powers and factorials. This is a volume argument with the actual lattice; no equidistribution or principal-ideal assumption.

\subsection{3. Every unmixed G reduction of rank at least two}\label{every-unmixed-g-reduction-of-rank-at-least-two}

The carried K068 cut says that every scalar-normalized unmixed G reduction in the retained kernel is unipotent, for both F and IC8. This is a finite-reduction statement, not global unipotence. First consider all support-eleven terms for which rank(N mod w)\textgreater=2 at every unmixed G place, N=h-I. The original finite coefficient satisfies \textbar a(h)\textbar\textless=2 Q\footnote{-335(f-1)}. No rank-pattern numerator is required. For L=I\_Sigma, z=T\^{}(-s/f)\textless2\^{}-65, because995\emph{11\textgreater65}168. The preceding volume bound applies and gives \#x\textless V\_plus T\^{}(2ms).

Retain the W0 count2\^{}14861 T\^{}(10s), the full fibre G\_n, binom(f,s)\textless=2\^{}168 and the conservative complete pair/weight factor P\_T\^{}s={[}2(T-1)\^{}2/(T-2){]}\^{}s\textless2\^{}(997s). These counts include all type profiles without an additional type multiplicity. Use \textbar c(h)\textbar\textless=1. Since2m+10=338, T\textless2\^{}996 and Q\^{}4\textgreater2\^{}281, the complete exponent is 168+5022+14861-45338+1+337645\emph{11-(281/4)}335*167 = -965309/4. This entire cell contributes less than2\^{}-241327.

In particular it contains every IC8 term with A!=I at support11. DO064 gives partial-monomial A with entries in mu28. If an unmixed G reduction had rank0 or1, finite nilpotence and the unitary inverse would give A+A\^{}dagger-2I=0 there. A nonidentity partial-monomial A has a diagonal0 or a nontrivial eta in mu28; its Hermitian defect is -2 or (eta-1)\^{}2/eta, a unit at the actual odd G primes prime to28. This is impossible. The argument needs no PM19 formula, rank-A enumeration, or moving integral basis.

\subsection{4. Nonunipotent scalar-cross terms}\label{nonunipotent-scalar-cross-terms}

Next consider A=I, h nonunipotent and h not fixing S. The original unitary block identities make x=(h e\_n)\_S nonzero, as in R049. At a rank-one finite G reduction write N=bc\^{}t. Since N\_SS=0, b\_S c\_S\^{}t=0, and the paired inverse I-N shows that x vanishes at one of the two CM partners. At rank0 it vanishes at both. Group every rank\textgreater=2 reduction as one class. With a fixed rule for choosing rank-one orientation there are at most4\^{}334 patterns.

Let J contain the chosen prime at rank1 and both primes at rank0. The original Green loss is at most (Norm J)\^{}(335/2). The nonzero e1/e2 witness of R049, with both CM valuations retained, gives Norm J \textless{} B\^{}f T\^{}t, where t\textless=s and B\textless2\^{}11 at H133. Thus L=J I\_Sigma has z\^{}(2f) \textless2\^{}(11f)T\^{}(t-2s)\textless=2\^{}(11f)T\^{}(-s) \textless2\^{}(11f-995s)\textless2\^{}(-54f). Hence z\textless2\^{}-27 and section2 applies.

Because Norm J=T\^{}(2s)z\^{}(2f), multiplying the S-column count by its finite Green loss leaves V\_plus T\^{}(335s) z\footnote{(335-2m)f} \textless{} V\_plus T\^{}(335s)2\^{}-31752, where (335-2m)f=1176. Now keep the same W0 count, G\_n, all4\^{}334 patterns, binom(f,s), P\_T\^{}s and \textbar c\textbar\textless=1 as above. The complete exponent is 168+5022+14861+1+668-45338-31752 +996\emph{345}11+997\emph{11-(281/4)}335*167 = -782877/4. Thus this entire cell contributes less than2\^{}-195719. The proof covers all type profiles at s11. It neither assumes global unipotence nor extends its special type-I normal form.

\subsection{5. Removal and remaining scope}\label{removal-and-remaining-scope}

Also retain the two full DO065 F,r1 globally-unipotent bounds at s11: (rank N,r)=(1,1) has exponent -103530801/128+8966\emph{11=-90906673/128; (rank N,r)=(2,1) has exponent -408071827/256+14941}11=-365997971/256. Both are less than-241327. Their global-unipotence premise is explicit; the counts do not depend on how that premise was established. Charge their current-remainder intersections here, whether or not a previous removal table had already eliminated some of their profiles.

Remove the section3 and section4 cells and these two F cells successively, intersecting each with the retained current remainder. The three smaller cell bounds sum to less than3*2\^{}-241327\textless2\^{}-195719. Together with section4 the joint charge is still less than2\^{}-195718. Append195718 only to the all-top error list and keep all earlier errors. The companion certificate must verify the current complete dyadic sum before promotion. Uniform H8192 is unchanged.

The complete all-top support floor remains11. At s11, K081 covers fixed-W0 nonunipotent t10 and t11 (through supports164 and149, respectively); DO067 excludes nonunipotent t\textless=9, and t\textless=s=11. Combined with the already checked F/IC8 partition, K069's r\textgreater=2 F cuts, R050's F,r1 nonunipotent s11 cells, the two F,r1 unipotent bounds just charged, and the K081 fixed-W0 removals, the remaining support-eleven terms have A=I and are globally unipotent of rank2. Fixed-W0 rank-one terms were removed by K077; identity has support0. Both the fixed-W0 and scalar-cross rank-two sectors remain, with the type profiles specified by the retained removal ledger. This corollary uses the carried exact class tables and makes no claim that either surviving sector is empty.

The sufficient signed inequality S\_R\textgreater-1+2\^{}-230 is still open. No prescribed Bessel/Mram period, specific032 cycle/value, literal source-family construction, shorter admissible chain, DP44 sign or RH extension is proved. All original modules, exterior constructions, finite projectors, central representatives and phases are retained.

\clearpage

\section{AM-K084 --- a totally positive diagonal witness for unipotent fixed blocks}\label{am-k084-a-totally-positive-diagonal-witness-for-unipotent-fixed-blocks}

\textbf{Canonical checked recovery version --- 2026-11-22T16:10:03.402339+00:00.} Root has read the complete mathematical proof and companion script and replayed the stated exact certificate. The original recovery export and its historical status are preserved separately. This is a current checked version, not recovery of unavailable historical final bytes. Only the stated sector and hypotheses are promoted; the signed period and actual DP44 cost remain open.

\subsection{1. Actual sector and the diagonal}\label{actual-sector-and-the-diagonal}

Retain the actual K081 fixed-W0 group and finite sum. Let h=I\_S⊕h0 be globally unipotent of rankN=2, N=h−I. The accepted signature-one unipotent classification gives N³=0 and N²≠0; then rankN²=1. Use the original form J=I\_4⊕(−1) on W0 and its fixed negative unit coordinate e\_n.~No integral moving basis is chosen.

The inverse polynomial and J-adjoint give N†=−N+N², N²=N+N†, and N†N=−N². Set alpha=(N²)nn.~Since N² is J-self-adjoint, alpha belongs to F.

Write ell=imN². Self-adjointness and (N²)²=0 make ell isotropic. Moreover N†N²=0, so imN⊂ell⊥. The latter space is positive semidefinite in signature(4,1). A nonzero isotropic vector has nonzero negative coordinate; kerN²=ell⊥ and the negative vector e\_n is not in this kernel. Hence (N² e\_n)n≠0, so alpha≠0. The sign is fixed by Jnn=−1: \textless N e\_n,N e\_n\textgreater\_J=\textless e\_n,−N²e\_n\textgreater\_J=alpha. Thus alpha is strictly positive at every real embedding of F.

Also alpha=2Re(h\_nn)−2. For the rank-one Cartan decomposition h0=k1a(r)k2, with Euclidean-unitary compact factors, \textbar\textbar h0\textbar\textbar=e\^{}r and \textbar h\_nn\textbar=cosh r. This follows equally by compact row/column rotations of the positive and negative block entries. Therefore under the actual cutoff \textbar\textbar h\textbar\textbar≤H, 0\textless alpha\_tau≤H+H\^{}-1−2=:R(H).

\subsection{2. The actual square-pole table applies to this entry}\label{the-actual-square-pole-table-applies-to-this-entry}

Use the retained K074 entrywise square-pole assertion: N² has pole orders at most0,2,1 at typesI,II,III, respectively, at both CM partners. This concerns the whole matrix N², so it applies to alpha; it is not the separate cross-entry integrality assertion.

For clarity, the typeI interface can be checked directly. Preserving the full projective pencil preserves its sum lattice Lplus and intersection lattice Lminus. Their quotient is a two-dimensional residue vector space, and Lminus⊂Λ⊂Lplus. The induced action is unipotent; hence its(N)² is zero. Thus N²Λ⊂Lminus⊂Λ, proving entrywise integrality. This does not require preserving the subset obtained by omitting the two endpoints. For typeIII, the accepted unequal-incidence bound on h and h\^{}-1 is pole≤1, and N²=h+h\^{}-1−2I preserves that bound. The general pair bound gives order≤2 for typeII. No SQ8 counting formula is used.

Let u and t count the actual typesII andIII, and define the F-ideal a\_T=product\_(II)q² product\_(III)q. Then A\_alpha=(alpha)a\_T is an integral F-ideal, with Norm\_F(A\_alpha)=\textbar Norm\_(F/Q)(alpha)\textbar T\^{}(2u+t)≤R(H)\^{}f T\^{}(2u+t). There are no poles at typeI places or away the actual nonintegral support. The actual G and T primes split in E/F with ramification index1, so F-valuations equal either paired CM valuation.

\subsection{3. Rank losses without a pattern enumeration}\label{rank-losses-without-a-pattern-enumeration}

At every unmixed G prime N is integral and its reduction is nilpotent. At reduction rank0 all entries of N² are divisible by the square of the prime. At rank1 the reduction has squarezero, so alpha is divisible by the prime. Because alpha∈F, these valuations agree at the paired CM primes.

Let P1 and P0 be the products of residue sizes over rank1 and rank0 G factors. Put J\_F=product\_(G rank1)q product\_(G rank0)q². The exact integral-ideal containment is A\_alpha⊂J\_F. Consequently P1 P0²=Norm\_F(J\_F)≤Norm\_F(A\_alpha)≤R(H)\^{}f T\^{}(2u+t). No fractional ideal is assumed principal. The original Green bound, with its one correction factor2 and both mixed-factor contractions, is \textbar a(h)\textbar≤2 Q\^{}(-335(f−1)) P1\^{}167 P0\^{}335 ≤2 R(H)\^{}(335f/2) T\^{}(335u+335t/2) Q\^{}(-335(f−1)). This is pointwise. There is no finite rank-pattern or orientation multiplier. Retain \textbar c(h)\textbar≤1.

\subsection{4. Original lattice count and type weights}\label{original-lattice-count-and-type-weights}

For fixed support of size s and actual ordered pencil pairs, keep K081's general count G5 BW T\^{}(10s), with G5\textless2\^{}107 and BW\textless2\^{}14861. These bounds remain valid for the present subset and for the smaller all-top cutoff.

Choose the typeIII and typeII positions. The coefficient is binom(s,t)binom(s−t,u)≤3\^{}s\textless2\^{}(2s). At each typeIII place, the actual off-diagonal incidence number cancels its absolute trace denominator, as in K081. At the remaining places retain the generic complete pair/trace bound P\_T\textless2\^{}997. With binom(f,s)≤2\^{}168 and T\textless2\^{}996, Q\^{}4\textgreater2\^{}281, one obtains log2 U\_(s,u,t)\textless C\_R+10959s+333660u+165833t, where C\_R=1+107+14861+168+(335f/2)log2R(H)−(281/4)335(f−1). The typeII generic bound is deliberately retained. No class mass, central phase, selected dimension or source normalization is removed.

\subsection{5. Complete support-eleven profile bounds}\label{complete-support-eleven-profile-bounds}

For the actual all-top H133 cutoff, R=17424/133\textless2\^{}(113/16); the integer-power certificate verifies this strictly. Substituting the rational upper exponent gives C\_R\textless−7432521/2. Among the78 profiles u,t≥0,u+t≤11, every profile except(u,t)=(11,0) is bounded by the maximum admitted exponent −186557/2 at(u,t)=(10,1). Thus all77 admitted profile sectors have total absolute contribution\textless2\^{}-93271, since each is\textless2\^{}-93278 and77\textless128. The all-type-II upper expression is+149097/2 and supplies no removal.

There is also a uniform H8192 consequence: R(H)\textless2\^{}13 gives C\_R\textless−14196717/4. Every support-eleven profile except(11,0) and(10,1) is bounded by the maximum admitted exponent−368121/4, attained at(10,0). The76 admitted profiles therefore contribute\textless2\^{}-92023. The two excluded expressions are+966519/4 and+295211/4.

Each sector is intersected with the current remainder before removal; previous errors remain conservatively retained. The nominal profile count does not assert nonempty intersections.

\subsection{6. Scope}\label{scope-2}

This removes only globally unipotent rank-two fixed-W0 support-eleven terms in the stated type profiles. In the all-top remainder, that sector is confined to all-type-II profiles. In the uniform remainder, it is confined to all-type-II or ten-type-II/one-type-III profiles. Other sectors require their separately checked reductions. A positive typeII trace factor alone gives no sign for the full summand. No Mram positivity, nonzero period, literal source construction, DP44 inequality or RH statement follows.

The new input is the fixed negative-coordinate diagonal alpha: its nonzero total positivity, sharp Cartan upper bound, and use of the actual square-pole table in a relative-F norm. All other counting interfaces are retained from K081 and its accepted sources.

\clearpage

\section{AM-R-052 --- all-top support-eleven scalar-cross type cuts}\label{am-r-052-all-top-support-eleven-scalar-cross-type-cuts}

\textbf{Canonical checked recovery version --- 2026-11-22T16:10:03.402339+00:00.} Root has read the complete mathematical proof and companion script and replayed the stated exact certificate. The original recovery export and its historical status are preserved separately. This is a current checked version, not recovery of unavailable historical final bytes. Only the stated sector and hypotheses are promoted; the signed period and actual DP44 cost remain open.

\subsection{1. Actual sector and original coordinates}\label{actual-sector-and-original-coordinates}

Retain the original J56 signed kernel, modules, exterior constructions, finite projectors, central representatives, phases, target pencils and Haar normalizations. Set f=a=168,n=169,m=164,s=11,H=133. The cutoff H133 is available only for the all-top family through R048's charged annulus. Consider A=h\_SS=I, h globally unipotent of rank(h-I)=2, and h not fixing S. Write N=h-I. The carried signature-one classification gives N\^{}3=0 and N\^{}2!=0, with rank N\^{}2=1. This is the scalar-cross sector retained by R051.

For a fixed type profile put i=\#type I,u=\#type II,t=\#type III; thus i+u+t=s. These count types, not global ranks. The original S part x=(h e\_n)\_S is nonzero and lies in I\_Sigma\^{}m, with Norm\_E(I\_Sigma)=T\^{}(-2s), by R051's direct original-pencil column inclusion. The same proof gives the W0-column count2\^{}14861 T\^{}(10s) and, for a fixed full negative column and pencil pair, the actual finite source-lattice stabilizer bound G\_n\textless2\^{}5022. No integral moving basis is introduced.

\subsection{2. Entrywise square poles and a positive F-valued witness}\label{entrywise-square-poles-and-a-positive-f-valued-witness}

The square-pole table I:0,II:2,III:1 applies to the full N\^{}2 in this sector. Here is a direct actual-pencil argument, independent of a fixed-W0 premise.

At a split T place write H\_c=I+c pi\^{}-1 Z, Z=zq of rank one with Z\^{}2=0, M\_c=H\_c Lambda, and k=H\_d\^{}-1 h H\_c in GL(Lambda). Type I preserves the actual pencil sum Lplus and intersection Lminus, with Lminus subset Lambda subset Lplus and Lplus/Lminus a two-dimensional residue vector space. The induced h is unipotent, so N\^{}2 Lplus subset Lminus and N\^{}2 Lambda subset Lambda. In type III, c-d is a unit. Expansion gives h=k+d pi\^{}-1 Zk-c pi\^{}-1 kZ-cd pi\^{}-2 ZkZ, tr h=tr k+(d-c)pi\^{}-1 tr(Zk). Since tr h=n and tr k is integral, tr(Zk) is divisible by pi. The rank-one identity ZkZ=tr(Zk)Z makes the apparent double pole simple. The same calculation for h\^{}-1=H\_c k\^{}-1 H\_d\^{}-1 uses tr h\^{}-1=n. Thus N\^{}2=h+h\^{}-1-2I has pole at most one. The general two-shear bound gives order at most two in type II. Repeat in the actual dual-pencil coordinates at the conjugate CM prime. There are no other finite poles.

Put alpha=(N\^{}2)\_nn.~From N\^{}dagger=-N+N\^{}2 we have N\^{}2=N+N\^{}dagger and N\^{}dagger N=-N\^{}2. The self-adjoint rank-one image ell=im N\^{}2 is isotropic, ker N\^{}2=ell\^{}perp, and im N lies in ell\^{}perp. A negative unit e\_n is outside ell\^{}perp, and the nonzero isotropic image has nonzero negative coordinate. Consequently alpha is a nonzero element of F. At every real embedding, alpha=\textless N e\_n,N e\_n\textgreater\_J\textgreater0. Also alpha=2 Re h\_nn-2 and \textbar h\_nn\textbar=cosh(log\textbar\textbar h\textbar\textbar), so 0\textless alpha\textless=R=H+H\^{}-1-2=17424/133\textless2\^{}(113/16).

At an unmixed G place with rank(N mod w)=1, finite nilpotence implies N\^{}2=0 modulo w; at rank0, every N\^{}2 entry is divisible by w\^{}2. The same holds at the CM partner. Use the original R049/R051 orientation rule: choose one CM prime for a rank-one reduction and both primes for rank0; group ranks\textgreater=2 together. The corresponding E-ideal J satisfies alpha in J bar(J) after clearing only the stated T poles. Hence Norm J \textless= R\^{}f T\^{}(2u+t). This follows equivalently from Norm\_E(alpha) \textgreater= (Norm J)\^{}2 T\^{}(-4u-2t) and Norm\_E(alpha)\textless=R\^{}(2f). No factor f versus2f is omitted.

\subsection{3. Weighted actual-ideal column count}\label{weighted-actual-ideal-column-count}

There are at most4\^{}334 orientation/rank patterns. For a fixed pattern, put L=J I\_Sigma and z=(Norm\_E L)\^{}(1/(2f)). The original Green factor is bounded by 2 Q\footnote{-335(f-1)} (Norm J)\^{}(335/2). Indeed its rank-one/rank-zero loss P1\^{}167 P0\^{}335 is at most (P1 P0\^{}2)\^{}(335/2), while Norm J=P1 P0\^{}2.

Use the actual half-trace normalization D0=7\^{}140 29\^{}144 5\^{}84 13\^{}84, covol(L)=D0 z\^{}(2f), L\#=2 Different\_E\^{}-1 bar(L)\^{}-1, and rd(E)\textless1500. The carried transference bound gives mu(L\^{}m)\textless750 sqrt(fm)z\textless9750 sqrt(f)z. The original positive part of a negative unit column has norm\textless H/2 at each real embedding. Therefore an actual Voronoi comparison gives \#x \textless= V(C) z\^{}(-2fm) whenever C\textgreater=133/2+9750z, where V(C)={[}((22/7) f C\^{}2)\^{}(fm){]}/{[}(fm)!D0\^{}m{]}. This V is a rational upper envelope for the actual Euclidean volume; a lower dyadic bracket for V is not a lower bound for the true-pi volume.

Since Norm J=T\^{}(2s) z\^{}(2f), multiplication by the Green loss leaves V(C) T\^{}(335s) z\^{}1176, because (335-2m)f=1176\textgreater0. The full expression V(133/2+9750z) z\^{}1176 is increasing for z\textgreater0. It is therefore legitimate to replace z by a uniform upper bound after this multiplication; no reversed lattice-density inequality is used.

Let w=2i+t=2s-2u-t. The witness and actual T\textgreater2\^{}995 imply z\^{}(2f)\textless2\^{}(113f/16)T\^{}(-w)\textless=(2\footnote{(2373-1990w)/2}). More transparently, log2 z \textless{} (2373-1990w)/672. Set b\_w=ceil((2373-1990w)/672) and C\_w=ceil(133/2+9750*2\^{}b\_w). Then z\textless2\^{}b\_w and the weighted column bound is V(C\_w) T\^{}(335s) 2\^{}(1176 b\_w). For w=0 the strict inequality comes from R\textless2\^{}(113/16), without needing a negative power of T.

\subsection{4. All finite choices and the exact complete exponent}\label{all-finite-choices-and-the-exact-complete-exponent}

Choose the support in binom(f,s)\textless=2\^{}168 ways and the type assignment in binom(s,t)binom(s-t,u)\textless=3\^{}s\textless2\^{}(2s) ways. At a type-III place there are at most(T-1)(T-2) ordered actual pencil pairs and the absolute incidence weight is exactly1/{[}(T-1)(T-2){]}. Their pair/weight sum is at most1. At every other active place retain the conservative full bound P\_T=2(T-1)\^{}2/(T-2)\textless2\^{}997; using this separately for I orII only overcounts. The matrix count and orientation-column bound are uniform in these pairs. Retain G\_n\textless2\^{}5022, the W0 count2\^{}14861 T\^{}(10s), all4\^{}334 patterns, the one Green correction2, \textbar c(h)\textbar\textless=1, T\textless2\^{}996 and Q\^{}4\textgreater2\^{}281.

If V(C\_w)\textless2\^{}v\_w, the complete exponent for one (u,t) profile is E\_(u,t) =168+5022+14861+1+668+2s+v\_w+1176b\_w +996*345s+997(s-t)-(281/4)335(f-1) =-474429/4+v\_w+1176b\_w-997t.

Exact integer powers and factorials give: w=0: b=4,C=156067,v=571040; w=1: b=1,C=19567,v=405967; w=2: b=-2,C=2504,v=242522; w=3: b=-5,C=372,v=90938; w=4: b=-8,C=105,v=-9622; w=5: b=-11,C=72,v=-39616; w=6: b=-14,C=68,v=-44160; 7\textless=w\textless=22: C=67,v=-45338, with b\_w given by the same exact ceiling. In every row2\^{}(v\_w-1)\textless V(C\_w)\textless2\^{}v\_w.

Among all78 pairs u,t\textgreater=0,u+t\textless=11, exactly74 have E\textless=-4096. Their largest exponent is -138185/4 at(u,t,i)=(9,1,1), strictly below-34546. The four exceptions are (u,t)=(9,2):478275/4; (10,0):486251/4; (10,1):1150155/4; (11,0):1828547/4. These are failures of this particular upper estimate, not lower bounds for any actual contribution or assertions of nonempty sectors.

\subsection{5. A freshly charged unipotent cell}\label{a-freshly-charged-unipotent-cell}

The full DO066 globally-unipotent scalar-cross bound is E\_old(s,v)=-14033453/4+10956s+332665v, v=u+t. Its support/type/pair multiplicities, common orientations, original W0 column/fibre and real compression are already included. After global unipotence is assumed, its proof has no additional t limit; the earlier t\textless=8/9 route supplied that assumption only.

At s11,v10 it gives -244789/4\textless-34546. In particular it bounds the exceptional profile(u,t,i)=(10,0,1). Charge this actual current-remainder intersection explicitly here, without inferring that an old ledger had already removed it. The retained full-proof scope is corroborated by the DO and Kummer branches; final saved-file reconciliation is recorded in the canonical integration manifest.

Together the74 section4 profiles and this one additional profile are bounded by75*2\^{}-34546\textless2\^{}-34539. Remove them successively from the actual current remainder and append 34539 only to the all-top error list, retaining every earlier charge. The final certificate must replay the complete current list before promotion.

\subsection{6. Smaller remaining support-eleven gate}\label{smaller-remaining-support-eleven-gate}

Within the all-top scalar-cross rank-two sector at support11, the only profiles not removed here are (i,u,t)=(0,11,0),(0,10,1),(0,9,2). Combined with R051 and K084, every remaining all-top support-eleven term has A=I, is globally unipotent of rank2, has no type-I place, and has at most two type-III places. In the fixed-W0 sector every place is type II. No assertion is made that any surviving sector is nonempty.

The total support floor remains11. The full signed inequality S\_R\textgreater-1+2\^{}-230 remains open. This note proves no prescribed global period, specific032 cycle/value, literal full family constructor, actual DP44 sign or RH extension. All claims about new saved files, complete all-top error replay and integration are now recorded in the canonical integration manifest.

\clearpage

\section{AM-R-053 --- uniform support-eleven scalar-cross type cuts}\label{am-r-053-uniform-support-eleven-scalar-cross-type-cuts}

\textbf{Canonical checked recovery version --- 2026-11-22T16:10:03.402339+00:00.} Root has read the complete mathematical proof and companion script and replayed the stated exact certificate. The original recovery export and its historical status are preserved separately. This is a current checked version, not recovery of unavailable historical final bytes. Only the stated sector and hypotheses are promoted; the signed period and actual DP44 cost remain open.

\subsection{1. Actual uniform sector}\label{actual-uniform-sector}

Use the original J56 signed kernel at the unchanged uniform cutoffH8192, with f=a=168,n=169,m=164,s=11. Let A=h\_SS=I and h be a globally unipotent rank-two scalar-cross element. Put N=h-I and alpha=(N\^{}2)\_nn.~The complete direct pencil and witness argument in R052 applies independently of H133: N\^{}3=0,N\^{}2 rank1, alpha in F is totally positive, and alpha\_tau\textless=H+H\^{}-1-2=R\textless2\^{}13. The entrywise N\^{}2 pole orders are I:0,II:2,III:1. Use i,u,t for the numbers of these types and v=u+t=s-i. All original module, metric, finite-projector, phase and measure factors remain those of the current kernel.

\subsection{2. Two sharper original-coordinate statements}\label{two-sharper-original-coordinate-statements}

First the S part x=(h e\_n)\_S is integral at type-I places. The actual common pencil has S subset Lminus subset Lambda subset Lplus, and h as well as h\^{}-1 preserves Lminus at both CM partners. Thus their S columns lie in Lambda. The actual unitary adjoint relation turns the S columns of the partner inverse into the S rows of h at the original prime. Therefore the entire S row, in particular its negative column entry, is integral. This is a direct pencil argument; it does not infer integrality of h from integrality of N\^{}2. At typesII andIII the R051 target-overlattice calculation gives simple S-column poles. Consequently x in (J I\_v)\^{}m, Norm\_E(I\_v)=T\^{}(-2v), for the same common G-orientation ideal J used in R049/R051/R052.

Second, the S-column radius is less than91 at every infinity. Since N\_SS=0 and N\^{}2=N+N\^{}dagger, (N\^{}2)\_SS=0. The nonzero self-adjoint rank-one N\^{}2 has the form a\_0 ell ell\^{}*J; its positive SS block vanishes only if ell\_S=0. Hence its isotropic image line lies in W0. As im N lies in ell\^{}perp, the W0 component of N e\_n lies in ell\^{}perp within W0, where the form is positive semidefinite. Thus \textbar\textbar x\_tau\textbar\textbar\^{}2 \textless= \textless N e\_n,N e\_n\textgreater\_J=alpha\_tau\textless=R\textless2\^{}13\textless91\^{}2. No divided isotropic basis or integral sphere parametrization is used.

\subsection{3. The actual weighted S lattice and class-general W0 bound}\label{the-actual-weighted-s-lattice-and-class-general-w0-bound}

The R052 witness gives Norm J\textless=R\^{}f T\^{}(2u+t). Set L=J I\_v and z=(Norm\_E L)\^{}(1/(2f)). Since 2u+t-2v=-t and T\textgreater2\^{}995, log2 z \textless(2184-995t)/336. Let b\_t=ceil((2184-995t)/336) and C\_t=ceil(91+9750*2\^{}b\_t). Use D0=7\^{}14029\^{}1445\^{}8413\^{}84, covol(L)=D0 z\^{}(2f), L\#=2 Different\_E\^{}-1 bar(L)\^{}-1. The full root discriminant is rd(E)=2 D0\^{}(1/f)\textless1500. Scalar transference followed by the orthogonal product gives mu(L\^{}m)\textless750sqrt(fm)z\textless9750sqrt(f)z. The original-coordinate Voronoi count is \#x\textless=V(C\_t) z\^{}(-2fm), V(C)={[}((22/7)f C\^{}2)\^{}(fm){]}/{[}(fm)!D0\^{}m{]}. It is the rational22/7 volume upper envelope. The actual Green loss is at most(Norm J)\^{}(335/2). Since Norm J=T\^{}(2v)z\^{}(2f), the weighted S count is V(C\_t)T\^{}(335v)2\^{}(1176b\_t). The use of the z upper bound is valid because1176=(335-2m)f\textgreater0 and V(91+9750z)z\^{}1176 is increasing.

For completeness, the W0 bound2\^{}14861T\^{}(10s) is valid for every A and pencil pair at H8192. The projected target over-lattice contains O\_E\^{}5 and has paired index T\^{}(10s), so its cover is no larger than that of O\_E\^{}5. At any infinity a negative unit column satisfies \textbar\textbar h e\_n\textbar\textbar\_Eucl\^{}2=cosh(2log\textbar\textbar h\textbar\textbar) =(\textbar\textbar h\textbar\textbar\^{}2+\textbar\textbar h\textbar\textbar\^{}-2)/2\textless5793\^{}2. The exact inequality 325\^{}21\emph{7\^{}35}29\^{}36\textless1636\^{}42 proves rd(E)sqrt(5)/2\textless1636. Thus a trace ball enlarged by the O\_E\^{}5 covering radius has radius less than7429sqrt(f). Its actual rational5-column volume is between 2\^{}14253 and2\^{}14254, hence below2\^{}14861. The fibre for a fixed full negative column and pair remains G\_n\textless2\^{}5022. This argument uses the full root discriminant; D0\^{}(1/f) alone is its half.

\subsection{4. Complete profile count}\label{complete-profile-count}

Retain all4\^{}334 common G patterns, one Green factor2, Q\footnote{-335(f-1)}, support choices\textless=2\^{}168, type assignments\textless=3\^{}s\textless2\^{}(2s), the W0 count and full source-lattice fibre, and \textbar c(h)\textbar\textless=1. At a type-III place the at-most(T-1)(T-2) ordered pencil pairs cancel their exact absolute incidence weight1/{[}(T-1)(T-2){]}. At every other place use the conservative complete bound P\_T\textless2\^{}997. All column bounds are uniform in the fixed pairs. With T\textless2\^{}996, Q\^{}4\textgreater2\^{}281 and v=s-i, the full profile exponent is E\_(i,u,t)=-474429/4+v\_t+1176b\_t-997t-996*335i, where V(C\_t)\textless2\^{}v\_t.

Exact powers and factorials give the following t,b,C,v table: 0,7,1248091,736324; 1,4,156091,571052; 2,1,19591,406064; 3,-2,2529,243312; 4,-5,396,95909; 5,-8,130,7357; 6,-11,96,-16746; 7,-14,92,-20129; 8,-17,92,-20129; 9,-20,92,-20129; 10,-23,92,-20129; 11,-26,92,-20129. Each rational V(C) is strictly between2\^{}(v-1) and2\^{}v.

Exactly72 of the78 profiles have E\textless=-4096. Their largest exponent is -130265/4\textless-32566 at(i,u,t)=(0,7,4). The exceptional triples and their upper estimates are (0,8,3):477447/4; (1,9,1):489967/4; (0,9,2):1146555/4; (1,10,0):1169155/4; (0,10,1):1824607/4; (0,11,0):2503795/4. These positive estimates do not assert positive or nonzero contributions.

\subsection{5. Two fresh carried-formula cells and the uniform charge}\label{two-fresh-carried-formula-cells-and-the-uniform-charge}

The full DO066 global-unipotent scalar-cross formula -14033453/4+10956s+332665v includes all support/type/pair/orientation/W0/fibre/real factors. At(s,v)=(11,10) it is -244789/4\textless-32566 and covers both exceptional i=1 profiles. As in R052, global unipotence is already assumed; no additional t bound remains in this count. Charge these current-remainder intersections explicitly here. No prior-removal provenance is inferred.

The72 new-volume profiles and these two profiles have union less than74*2\^{}-32566\textless2\^{}-32559. Append32559 to the uniform error list and retain all previous errors. The complete retained31-term uniform list checks below2\^{}-230 in exact integer arithmetic; a final actual-file replay is still required. Every all-top counterpart of these removed profiles is already among R052's75 charged profiles. Do not add a duplicate all-top charge.

\subsection{6. Exact scope of the smaller gate}\label{exact-scope-of-the-smaller-gate}

Within the uniform rank-two globally-unipotent scalar-cross support-eleven sector, only i=0 and t=0,1,2,3 remain. Other uniform structural classes retain their previous status. The overall support floor stays11, and neither the complete signed inequality S\_R\textgreater-1+2\^{}-230 nor any prescribed global period is proved. Specific032 cycle/value, literal full-family construction, actual DP44 sign and the RH extension remain open. The result is a charged reduction of this named sector, with final file promotion now recorded in the canonical integration manifest.

\clearpage

\section{AM-K085 --- an affine real-coordinate slice for the all-type-II fixed-block sector}\label{am-k085-an-affine-real-coordinate-slice-for-the-all-type-ii-fixed-block-sector}

\textbf{Canonical checked recovery version --- 2026-11-22T16:10:03.402339+00:00.} Root has read the complete mathematical proof and companion script and replayed the stated exact certificate. The original recovery export and its historical status are preserved separately. This is a current checked version, not recovery of unavailable historical final bytes. Only the stated sector and hypotheses are promoted; the signed period and actual DP44 cost remain open.

\subsection{1. Exact sector and retained inputs}\label{exact-sector-and-retained-inputs}

Retain the actual all-top J56 upper-gap kernel, original favourable point, fixed finite projectors, scalar-pencil family, all class masses, and canonical twists. Set f=a=168, n=169, m=164, H=133, and s=11. We consider only terms h=I\_S⊕h\_0, with h\_0 acting on the fixed W0 of signature(4,1), such that N=h−I is globally unipotent of rank2, and every one of the s active T places has typeII (unequal pencils with a unique fixed incidence). No other sector is silently substituted.

The accepted rank-one-signature classification gives N³=0 and rank(N²)=1. AM-K084 proves that alpha=(N²)nn ∈ F, alpha\_tau\textgreater0, and 0\textless alpha\_tau≤R:=H+H\^{}−1−2=17424/133. It also proves that alpha has poles at most2 at the active F-primes over T and is integral elsewhere. Its finite-G rank valuations are at least1 for rank(N mod G)=1 and at least2 for rank0. In this sector finite ranks are at most2.

Retain the exact original five-column fibre bound G5\textless2\^{}107, Q\^{}4\textgreater2\^{}281, T\textless2\^{}996, and the half-trace covolume D0=7\^{}140 29\^{}144 5\^{}84 13\^{}84. The full root discriminant is rd(E)=2D0\^{}(1/f), not D0\^{}(1/f). Retain the actual local pencil model, established for the original family in K077: M\_c=(I+c pi\^{}−1 Z)Lambda, Z=zq, Z²=0, B\_c=pi\^{}−1 M\_c, where c is a unit at an admitted label, the pole-line negative coordinate z\_n is a unit, and some q\_j is a unit. In the paired dual model the unit roles are exchanged. These are the actual admitted labels; no endpoint is inserted.

\subsection{2. The actual real-compression factor gains the homogeneous energy}\label{the-actual-real-compression-factor-gains-the-homogeneous-energy}

The retained K060 RC3/RC6/RC13 argument decomposes the homogeneous holomorphic polynomial model into its rank-one discrete-series ladders. In its positive-energy orientation the rank-one compact circle is diag(e\^{}(i theta),1,\ldots,1;e\^{}(−i theta)); the base coordinates therefore have weights2,1,\ldots,1. On the algebraic monomial⊗fibre-weight basis its generator has the eigenvalues described by H1=H\_fibre+E\_total+E\_1, with smallest fibre weight kappa0. The lower bound follows from these eigenvalues and self-adjointness of the actual compact generator H1; it does not require H\_fibre and E\_1 separately to preserve the coupled Bergman norm. The chosen holomorphic lowering root acts by coordinate differentiation, and its raising root raises total homogeneous degree by1. A ladder with lowest spectator degree q thus meets homogeneous degree N\_hom only in its level j=N\_hom−q. On that one-dimensional intersection the Jacobi bound uses the actual circle weight 2j+kappa, and the displayed generator formula gives 2j+kappa ≥ kappa0+N\_hom. The retained orthogonal ladder decomposition therefore applies this same bound to the compressed operator on the complete homogeneous subspace. Therefore the same compressed-operator argument gives \textbar\textbar P\_W pi(a(r)) P\_W\textbar\textbar{} ≤ 12 (kappa0+N\_hom)\^{}−1/4 (sinh r)\^{}−1/2. This is the original operator normalization; it is not a replacement of the actual kernel by a character or an averaged trace.

Here N\_hom is fixed by the scalar U(a) action in the actual Pieri image: it is sum\_i(lambda\_i−alpha\_i)=sum\_i m\_i, as in the accepted PB3 decomposition. The m\_i are nonnegative. In the all-top R048 profile the first increment is m\_0=2\^{}960+83−668 ell, ell\textless2\^{}485, so N\_hom≥m\_0\textgreater2\^{}959. This derives the energy floor from the actual homogeneous degree; the notation m\_0 is not being identified with energy by name alone. Compact rotations preserve that degree and the selected compact type, so the same bound applies in each Cartan direction.

Write the actual favourable-point conjugator as g\_p.~Its Cartan radius at every real place is \textless1/2 by the retained LIFT41 bound. Let v=(a,1) be the isotropic image vector of N² normalized by its negative coordinate. The normalization is used only over C, never as an integral finite frame. Since N²=−alpha vv\emph{J and \textbar a\textbar=1, alpha\_p:=(g\_p\^{}−1 N²g\_p)nn =alpha \textbar v}Jg\_pe\_n\textbar² ≥e\^{}−1 alpha\textgreater alpha/3. Indeed the positive and negative components of g\_pe\_n have norms sinh r\_p and cosh r\_p, giving the lower bound e\^{}−r\_p for their pairing with v.

For h'=g\_p\^{}−1 h g\_p, the Cartan radius r' satisfies sinh²r'≥alpha\_p: its negative diagonal has real part1+alpha\_p/2 and modulus cosh r'. Applying the actual compressed-operator estimate at every place, and retaining the original normalized contraction bound, yields \textbar c(h)\textbar{} ≤ 12\^{}f 2\^{}(f/2) 2\^{}(−959f/4) \textbar Norm\_F(alpha)\textbar\^{}−1/4. (AS1) The harmless factor2\^{}(f/2) bounds3\^{}(f/4). No assumption on the sign of c(h), and no character formula for it, is used.

\subsection{3. The actual negative-column lattice and its real projection}\label{the-actual-negative-column-lattice-and-its-real-projection}

For a fixed support and its actual diagonal incidence labels, patch the local B\_c with the standard lattice away that support to obtain the original target over-lattice B\_Sigma⊂E\^{}5.

The inverse matrix pi I−cZ is integral, so Lambda⊂B\_c.~Also every entry of B\_c has pole≤2. A unit q\_j gives a last-row entry whose pi\^{}−2 coefficient c z\_n q\_j is a unit; the pi\^{}−1 part cannot cancel it. Thus pr\_n(B\_c)=pi\^{}−2 O at the selected CM prime. The paired dual calculation gives the same assertion at its partner. Consequently B\_Sigma contains O\_E\^{}5, and its half-trace covolume is covol(B\_Sigma)=D0\^{}5 T\^{}−10s. The determinant-one shear gives the same covolume as the standard lattice before multiplication by the paired ideal I\_Sigma.

Use the real orthogonal projection p:E\_R\^{}5→F\_R, p(x)=Re(x\_n), with metrics sum\_CM\textbar x\_tau\textbar² and sum\_F y\_tau². Its normal Jacobian is1. Put P0=Re(O\_E), and I\_F=product\_(active q) q\^{}−1. The exact local projection calculation and localization give P\_Sigma=p(B\_Sigma)=P0 I\_F². (AS2) This retains the dyadic module P0. We only use O\_F⊂P0⊂(1/2)O\_F; no everywhere-unramified trace-surjectivity premise is needed.

Let K\_Sigma=ker(p\textbar B\_Sigma). Euclidean coarea for this lattice exact sequence gives covol(K\_Sigma)=D0\^{}5 T\^{}−8s / covol(P0). Since covol(P0)≤sqrt(D\_F), and D\_E=D\_F² Norm(D\_(E/F)), sqrt(D\_F)≤2\^{}(f/2)D0\^{}(1/2). Hence covol(K\_Sigma)≥2\^{}(−f/2)D0\^{}(9/2)T\^{}−8s. (AS3) If one projected to alpha=2beta−2 instead, its normal Jacobian2\^{}f and projected-lattice scaling2\^{}f would cancel. They have not been discarded.

The kernel contains O\_E\^{}4⊕iO\_F. Apply the retained transference inequality mu(L)lambda1(L\#)≤dim(L)/2 separately in E and F. The actual half-trace dual of O\_E is 2Diff\_E\^{}−1, while the real dual of O\_F is Diff\_F\^{}−1. Thus mu(O\_E)≤rd(E)sqrt(f)/2, mu(O\_F)≤rd(F)sqrt(f)/2. As rd(F)≤rd(E), mu(K\_Sigma)≤rd(E)sqrt(5f)/2\textless1636 sqrt(f). (AS4) The exact certificate checks the last strict inequality.

For any actual h\_0 under cutoff H, the Euclidean squared norm of its negative column is 2\textbar h\_nn\textbar²−1≤(H²+H\^{}−2)/2\textless95² at each embedding. Orthogonal projection onto any affine beta fibre can only decrease this radius. Voronoi-cell counting in the (real) 9f-dimensional fibre, using (AS3)--(AS4), therefore bounds the number of negative columns with a prescribed beta by V\_K T\^{}8s, V\_K= {[}(pi f 1731²)\^{}(9f/2) 2\^{}(f/2){]} / {[}(9f/2)! D0\^{}(9/2){]} \textless2\^{}9850. (AS5) This counts the actual original-coordinate affine fibre. No moving integral basis is introduced.

\subsection{4. A real ideal counts the affine diagonal and absorbs its small-value cost}\label{a-real-ideal-counts-the-affine-diagonal-and-absorbs-its-small-value-cost}

Set x=h e\_n and beta=p(x). Since x\_n=h\_nn, alpha=2beta−2; hence gamma=beta−1=alpha/2. The shift is legitimate because1∈P\_Sigma. In particular gamma lies in P\_Sigma, is positive at every embedding, and is at mostR/2 there.

For a fixed finite-G rank pattern, put J\_F=product\_(rank1 G) q product\_(rank0 G) q². K084's exact numerator-ideal inclusion gives Norm(J\_F)≤\textbar Norm\_F(alpha)\textbar{} T\^{}2s≤R\^{}f T\^{}2s. At the odd G primes, division by2 is a unit. At every other prime J\_F is a unit ideal. Thus gamma∈J\_F P\_Sigma⊂L:=(1/2)J\_F I\_F². Let y=(Norm(J\_F) T\^{}−2s)\^{}(1/f). Then0\textless y≤R, Norm\_F(L)=(y/2)\^{}f, and covol(L)=2\^{}−f sqrt(D\_F)y\^{}f. The real-ideal dual and the same transference form give mu(L)≤rd(F)sqrt(f)y/4\textless375sqrt(f)y. Consequently the number of possible gamma is at most {[}pi\^{}(f/2) f\^{}(f/2) 2\^{}f / ((f/2)!sqrt(D\_F)){]} ·(R/2+375y)\^{}f y\^{}−f.~(AS6)

The original Green bound is \textbar a(h)\textbar≤2 Q\^{}−335(f−1) Norm(J\_F)\^{}(335/2). Combining it with (AS1), and using \textbar Norm alpha\textbar≥Norm(J\_F)T\^{}−2s=y\^{}f, gives a common weight for this pattern proportional to T\^{}335s y\^{}(669f/4). Multiply this by (AS6). The function y\^{}(665f/4)(R/2+375y)\^{}f is increasing for y\textgreater0. Its maximum on0\textless y≤R is atR. Dropping only the favourable factor1/sqrt(D\_F), the resulting real-diagonal count constant is bounded by V\_beta={[}(pi f 376²)\^{}(f/2) 2\^{}f{]}/(f/2)!\textless2\^{}1945, and the remaining alpha power is R\^{}(669f/4).

Thus the possible smallness of alpha is absorbed before maximizing the Green-weighted count. There is no separate T\^{}(s/2) penalty, and no unproved density-only use of a congruence index.

There are334 unmixed G factors. Each contributes exponent0,1,2 to J\_F, so there are at most3\^{}334\textless2\^{}530 patterns. This factor is retained. No orientation choices are required because alpha belongs to F. For each pattern we count a larger set of columns satisfying its divisibilities and then use the bound only on matrices with that actual pattern; this overcount is valid.

\subsection{5. Complete finite and real normalization; exact exponent}\label{complete-finite-and-real-normalization-exact-exponent}

At typeII the only surviving pair is diagonal. There are T−1 diagonal labels and its normalized absolute trace is1/(T−1); their product is1. This exact cancellation replaces the deliberately loose generic pair factor in K084. All central phases and fixed class masses remain in the original kernel. The actual negative-column fibre still has multiplicity at most G5.

Retain binom(f,s)≤2\^{}168 and even the harmless earlier type-choice factor2\^{}(2s), although this one all-II profile does not require it. With (AS1), (AS5)--(AS6), all finite rank patterns, and s=11, the total absolute contribution of this sector is strictly below2\^{}E, where E= 1+107+9850+1945+168+530 + (335+8)·996·11+2·11 + (669f/4)·(113/16) − (281/4)·335(f−1) + f(4+1/2−959/4) =−5481/8=−685.125. (AS7)

Here R\^{}16\textless2\^{}113 is checked exactly, 12\^{}f\textless2\^{}(4f), and pi is bounded by22/7 in both volume certificates. In particular the rational-pi fibre ceiling is9850, not the smaller floating-pi diagnostic9849. The exact script squares the fibre expression to handle D0\^{}(9/2) using integers, and verifies all other displayed numerical bounds.

The entire new sector therefore contributes\textless2\^{}−685.

\subsection{6. Scope and resulting reduction}\label{scope-and-resulting-reduction}

Intersect this sector with the current remainder before removing it, and retain every earlier error. The new all-top charge is685. The original outage proof did not have the complete error array. The canonical integration now verifies the full chronological39-term all-top array through this result, with total strictly below2\^{}-230.

Together with K084 this removes the remaining all-top rank-two globally unipotent fixed-W0 support-eleven profile. The earlier rank-one fixed-block and nonunipotent fixed-block reductions may then be cited separately to remove the full fixed-W0 support-eleven sector. This is not a claim that every support-eleven term is removed: the separately retained scalar-cross and other applicable sectors remain governed by their own checked bounds. No uniform-H8192 energy floor is asserted. No signed-kernel positivity, Mram nonvanishing, full DP44 budget, literal Host/FamP instantiation, or RH conclusion follows.

The new arithmetic input is the actual affine Re(x\_n) slice and its exact original-lattice coarea, combined with the finite-rank ideal and actual homogeneous real energy. The use of normalized isotropic v occurs only in the complex-place compression estimate; it is never declared integral at G or T.

\clearpage

\section{AM-DO-070 --- a fractional-ideal transference improvement for the scalar-cross slice}\label{am-do-070-a-fractional-ideal-transference-improvement-for-the-scalar-cross-slice}

\textbf{Canonical integration --- 2026-11-22T16:10:03.402339+00:00.} Whole root proof/script reading and separate exact replay are recorded in CHECK\_EVENTS.tsv. Original recovery bytes remain preserved.

Author: /root/definite\_oscillator. Controlled recovery draft, 22 November 2026. This is new work, not a restored pre-outage proof. It is confined to the actual all-top J56 kernel and the support-eleven profile (i,u,t)=(0,9,2). Existing proofs and ledgers are unchanged. The canonical integration separately replays every chronological prior-error list and records this scoped promotion.

\subsection{1. Exact setting and retained inputs}\label{exact-setting-and-retained-inputs}

Set f=a=168, n=169, m=164, s=11 and H=133. Consider actual scalar-cross matrices h=I+N with h\_SS=I, global rank(N)=2, N\^{}3=0, and rank(N\^{}2)=1, at the actual translated packet point. There are nine type-II places and two type-III places, with no type-I place. Every matrix is in the original signed-kernel remainder; bounds below may count a larger set only to bound the absolute contribution of this actual subset.

Retain the exact local pencil, original-coordinate column, paired-ideal and selected-energy identities proved in the companion direct\_pencil\_and\_scalar\_cross\_lemmas.md and independently checked in R052/R053/K085. In particular \[
 \alpha=(N^2)_{nn}\in F,\quad
 0<\alpha_\tau\le R=\frac{17424}{133},\quad
 \|(he_n)_{S,\tau}\|^2\le\alpha_\tau.
 \tag{FT1}
\] At type II alpha has at most double poles, at type III at most simple poles. The actual common-orientation ideal J\_E has the paired F-ideal J\_F with equal norm. Define \[
 y=(NJ_F\,T^{-(2u+t)})^{1/f}\le R,\qquad
 z=(NJ_E)^{1/(2f)}T^{-s/f}
   =\sqrt y\,T^{-t/(2f)}.
 \tag{FT2}
\]

The arithmetic constants are \[
 2^{995}<T<2^{996},\quad Q^4>2^{281},\quad
 D_0=7^{140}29^{144}5^{84}13^{84},\quad
 rd(E)<1500.
\] The half-trace metric on E has ring covolume D0. Its actual Euclidean dual for a fractional ideal L is \[
 L^\#=2\mathfrak D_E^{-1}\bar L^{-1}.
\] For \(N_EL=z^{2f}\), the accepted codifferent minimum argument gives \[
 \lambda_1(L^\#)\ge\frac{2\sqrt f}{rd(E)\,z}.
 \tag{FT3}
\] This is the actual trace-lattice normalization; no ring-only covering estimate is substituted for a fractional ideal.

\subsection{2. A primary finite-dimensional transference bound}\label{a-primary-finite-dimensional-transference-bound}

Miller and Stephens-Davidowitz (2019), Theorem 3.2, equation (3.3), state for every Euclidean lattice of rank d that \[
 \mu(L)\lambda_1(L^\#)
 \le \frac{d}{2\pi}+\frac{3\sqrt d}{\pi}.
 \tag{FT4}
\] The theorem and definitions were directly read in arXiv:1802.05708v2, printed page 8. It has no asymptotic threshold and no arithmetic-shape assumption.

For d=336, use the elementary lower bound pi\textgreater157/50 and sqrt(336)\textless19. Then \[
 \frac{336}{2\pi}+\frac{3\sqrt{336}}{\pi}
 <\frac{11250}{157}<72.
\] Combining FT3 and FT4, and using \(\sqrt{164}<13\), gives \[
 \mu(L^{164})
 \le\sqrt{164}\,\mu(L)
 <\frac{72\cdot1500\cdot13}{336}\sqrt f\,z
 <4200\sqrt f\,z.
 \tag{FT5}
\] The product identity for covering radii is for the orthogonal product metric. FT5 is valid for every fractional ideal L occurring in the common-orientation count.

Define the rational upper-envelope volume \[
 V_m(C)=\frac{\bigl((22/7)fC^2\bigr)^{fm}}
              {(fm)!\,D_0^m}.
 \tag{FT6}
\] The true Euclidean ball volume is bounded above by this envelope. It is not identified with it.

\subsection{3. The full scalar-cross slice, with distinct image and subset lattices}\label{the-full-scalar-cross-slice-with-distinct-image-and-subset-lattices}

For a fixed actual local pencil pair let r0 count its type-III targets labelled zero. DO070-C1 ZC1--8 proves that the full W0 real-part image has pole two at nonzero targets and pole one at zero targets; its kernel covolume has T exponent -8s-r0. The original claim of a double-pole image at every type-III target is superseded. The complete weighted sum over the actual pairs contributes exactly the sufficient extra factor2\^{}t, by ZC9. No label or CM partner is dropped.

Only the actual unipotent gamma, \[
 \gamma=\operatorname{Re}(h_{nn})-1=\alpha/2,
\] has the stronger membership \[
 \gamma\in\tfrac12J_FI_{\rm II}^2I_{\rm III}.
 \tag{FT7}
\] The factor 1/2 is retained at dyadic places. All G and T places used here are odd. FT7 does not change the full projection image, and therefore does not duplicate the simple-pole gain in the fibre covolume.

The retained gamma point count is \[
 A_f(R/2+375y)^f y^{-f},\qquad
 A_f=\frac{\pi^{f/2}f^{f/2}2^f}{(f/2)!\sqrt{D_F}}.
 \tag{FT8}
\] For a fixed pair the W0 fibre count is \(2^{9850}T^{8s+r_0}\); after actual pair weighting its bound is \(2^{9850+t}T^{8s}\). The full source-lattice fibre bound is \(G_n<2^{5022}\), not the fixed-W0 \(G_5\) bound.

The S column lies in \((J_EI_\Sigma)^m\). Applying FT5 and the sharp radius in FT1 bounds its count by \[
 V_m(\sqrt R+4200z)\,z^{-2fm}.
 \tag{FT9}
\]

The Green factor is \[
 2Q^{-335(f-1)}(NJ_E)^{335/2}.
\] The actual selected PB3 homogeneous degree is greater than \(2^{959}\); the K060 Jacobi compression and actual near-point comparison give \[
 |c(h)|\le12^f2^{f/2}2^{-959f/4}(N_F\alpha)^{-1/4}.
 \tag{FT10}
\] This real factor is valid for the scalar-cross rank-two matrix as well: its proof uses the isotropic image line of N\^{}2 and the actual packet displacement, not a finite integral normalization of that line.

Combining FT8--FT10 with FT2, the y-dependent factor, before maximizing, is \[
 y^{9f/4}(R/2+375y)^f
 V_m\!\left(\sqrt R+4200\sqrt y\,T^{-t/(2f)}\right).
\] Each factor is increasing in y\textgreater0, so its maximum over y\textless=R occurs at y=R. After the gamma volume is bounded by its retained constant \(2^{1945}\), this yields \[
 2^t T^{343s-\frac72t}\,
 R^{13f/4}\,
 V_m\!\left(\sqrt R(1+4200T^{-t/(2f)})\right).
 \tag{FT11}
\] Equivalently use \(V_m(\sqrt R C)=R^{fm}V_m(C)\); the total R exponent is \(669f/4\).

For t=2, the exact integer inequality \(60^{168}<2^{995}<T\) gives \[
 T^{-t/(2f)}=T^{-1/168}<1/60.
\] Hence the volume in FT11 is at most \(R^{fm}V_m(71)\). Exact integer comparisons show \[
 2^{-40729}<V_m(71)<2^{-40728},\qquad R^{16}<2^{113}.
 \tag{FT12}
\]

\subsection{4. Every multiplicity and the final bound}\label{every-multiplicity-and-the-final-bound}

Retain the following factors once each:

\begin{itemize}
\tightlist
\item
  supports: at most \(2^{168}\);
\item
  type assignments: the conservative \(3^s<2^{2s}=2^{22}\);
\item
  common-orientation patterns: \(4^{334}=2^{668}\);
\item
  Green constant 2 and denominator \(Q^{335(f-1)}\);
\item
  pair-weighted W0 real-part fibre \(2^{9850+t}T^{8s}\);
\item
  gamma constant \(2^{1945}\);
\item
  full source-lattice fibre \(2^{5022}\);
\item
  actual selected real factor \(12^f2^{f/2}2^{-959f/4}\).
\end{itemize}

For each type-II place, the number of diagonal pencil labels times its positive rho is one. For each type-III place the raw pair weights sum to one, but its zero targets have an extra T fibre cost. Their exact weighted average is2, as proved in DO070-C1 ZC9. This is the factor2\^{}t already included above; no extra independent conjugate-place factor is present. Rank information is encoded in the existing common-orientation patterns; no second \(3^{334}\) rank-pattern count is introduced.

Use \(12<16\), \(T<2^{996}\), \(Q^4>2^{281}\), and FT12. The complete base-two exponent is \[
\begin{aligned}
E={}&t+1+9850+1945+5022+168+22+668\\
 &+996(343\cdot11-7)
   +\frac{669\cdot168}{4}\frac{113}{16}
   -\frac{281}{4}\,335\cdot167\\
 &+168\left(\frac92-\frac{959}{4}\right)-40728\\
 ={}&-\frac{346641}{8}<-43330.
\end{aligned}
\tag{FT13}
\] Therefore the absolute contribution of the whole specified profile is less than \(2^{-43330}\).

Remove its intersection with the current remainder, retaining every older error. This is a new all-top charge only. It removes the support-eleven scalar-cross profile with i=0,u=9,t=2. In the previously stated all-top scalar-cross support-eleven remainder, the profiles t=0 and t=1 remain. No conclusion about their sign, the full signed sufficient inequality, an actual Mram period, a cycle value, a literal Host constructor, DP44, or RH follows.

\subsection{Verification status}\label{verification-status}

The exact bound was first independently calculated in the restored executor. The companion script checks the finite-dimensional conversion, the actual T window consequence, the rational volume interval, and all displayed exponent identities. The script does not prove the analytic compression or the lattice interfaces automatically. Those are supplied by the full mathematical argument and accepted inputs. No complete old-error-list replay is claimed here.

\subsection{Primary reference}\label{primary-reference-1}

Miller, S. D., \& Stephens-Davidowitz, N. (2019). Kissing numbers and transference theorems from generalized tail bounds (arXiv:1802.05708v2). \url{https://arxiv.org/pdf/1802.05708}. Theorem 3.2, equation (3.3), printed page 8; definitions of the Euclidean dual and covering radius in Section 1.

\subsubsection{Canonical correction record}\label{canonical-correction-record}

This full version incorporates AM-DO-070-C1 ZC1--11, read in full together with the original proof and the corrected19-check script. The original43332 calculation and its earlier checks remain preserved as historical records; they omitted the zero-target fibre cost. The corrected charge is43330 and the removed profile is unchanged.

\clearpage

\section{AM-DO-071 --- sparse-support refinement of the remaining scalar-cross profile}\label{am-do-071-sparse-support-refinement-of-the-remaining-scalar-cross-profile}

\textbf{Canonical integration --- 2026-11-22T16:10:03.402339+00:00.} Whole root proof/script reading and separate exact replay are recorded in CHECK\_EVENTS.tsv. Original recovery bytes remain preserved.

Author: \nolinkurl{/root/definite_oscillator}, 22 November 2026. This canonical version incorporates the separately proved zero-target correction DO070-C1. The original exported proof is preserved unchanged. The geometric argument and the independent exact finite optimization now bound the sector with at most 107 nonzero original S-coordinates. The full profile remains unresolved.

\subsection{Setting and the simultaneous coordinate constraint}\label{setting-and-the-simultaneous-coordinate-constraint}

Use the actual all-top J56 support-eleven scalar-cross global rank-two unipotent profile \nolinkurl{(i,u,t)=(0,10,1)}. Set \nolinkurl{f=168}, \nolinkurl{m=164}, \nolinkurl{s=11}, \nolinkurl{H=133}, and \nolinkurl{R=17424/133}. Keep the original positive coordinates on S; no finite or archimedean change of basis is made.

Write \nolinkurl{x=(h e_n)_S}, and let \nolinkurl{k} be the number of its nonzero E-coordinates. Scalar-cross means \nolinkurl{1≤k≤164}. For the actual common-orientation ideal L, retain

\[
 x\in L^{164},\qquad N_E L=z^{2f},\qquad
 \|x_\tau\|^2\le\alpha_\tau\le R.
 \tag{SP1}
\]

If \nolinkurl{x_i≠0}, the fractional ideal inclusion \nolinkurl{(x_i)⊂L} implies \[
 |N_{E/\mathbf Q}(x_i)|\ge N_E L=z^{2f}.
\] The CM-type product is \nolinkurl{|N_E(x_i)|=∏τ |τ(x_i)|²}. The arithmetic-geometric mean inequality therefore gives \[
 \sum_\tau|\tau(x_i)|^2\ge fz^2.
\] Sum over the k nonzero original coordinates, then use SP1. This proves the simultaneous constraint \[
 kz^2\le R.
 \tag{SP2}
\]

This is an exact constraint on the actual vectors. It does not assume an integral orthogonalization, a unit value of a coordinate, or a norm equation for a divided isotropic vector.

\subsection{Fractional product cover and the monotone endpoint}\label{fractional-product-cover-and-the-monotone-endpoint}

DO070's directly checked finite-dimensional transference bound, with the actual Euclidean dual, gives \[
 \mu(L^k)<c\sqrt k\sqrt f\,z,
 \qquad c=\frac{72\cdot1500}{336}=\frac{2250}{7}<322.
 \tag{SP3}
\] For a fixed support of size k, the Voronoi count in the product ball is \[
 V_k(\sqrt R+c\sqrt k\,z)z^{-2fk},\qquad
 V_k(C)=\frac{((22/7)fC^2)^{fk}}{(fk)!D_0^k}.
 \tag{SP4}
\] \nolinkurl{V_k} is the rational upper envelope of the Euclidean volume, not its exact value. The number of support choices is \nolinkurl{binom(164,k)}.

Put \[
 w=T^{1/f},\quad \delta=T^{-1/(2f)},\quad
 y=(NJ_F T^{-21})^{1/f},\quad q=y/R.
\] The actual alpha cap gives \nolinkurl{q≤1}, and \nolinkurl{z²=y/w} together with SP2 gives \[
 0<q\le q_k:=\min(1,w/k).
 \tag{SP5}
\]

Retain the full gamma count, its dyadic half-lattice, the full W0 real-part fibre, the full source-lattice fibre, the actual real compression, and all common-orientation patterns exactly as in DO070. After factoring out the common gamma constant \nolinkurl{2^1945}, the k-dependent factor is \[
\begin{split}
&T^{343s-335/2+k} R^{669f/4}\,
 q^{(665/4-k)f}
 \left(\frac{1/2+375q}{376}\right)^f\\
&\hspace{2cm}\cdot
 V_k\!\left(1+c\sqrt{kq}\,\delta\right).
\end{split}
\tag{SP6}
\] The common omitted multiplier in this displayed factor is \[
2^{2+9850+1945+5022+168+22+668}
12^f2^{f/2}2^{-959f/4}Q^{-335(f-1)}
\binom{164}{k}.
\tag{SP7}
\] Thus SP6 does not hide a new finite projector, a pair-count division, or a second rank-pattern sum. The raw pair/absolute-weight sums are at most one, but type-III target zero increases the fibre count. DO070-C1 ZC1--9 proves that the complete pair-weighted fibre cost is2 at this one type-III place; its extra binary bit is included in SP7. The fixed profile still uses the conservative support and type bounds in SP7.

For every \nolinkurl{1≤k≤164}, \nolinkurl{665/4-k≥9/4>0}. Every q-dependent factor in SP6 is increasing. Its maximum is therefore attained at the actual permitted endpoint \nolinkurl{q=q_k}. This establishes a finite exact optimization with no density-only approximation and no replacement of a signed term by a claimed positive contribution.

\subsection{Two exact branches for the actual prime window}\label{two-exact-branches-for-the-actual-prime-window}

The actual prime satisfies \nolinkurl{2^995<T<2^996}. The exact integer comparisons \[
 60^{168}<2^{995},\qquad 2^{996}<61^{168}
\] give \nolinkurl{60<w<61}. Therefore:

\begin{itemize}
\tightlist
\item
  For \nolinkurl{1≤k≤60}, use \nolinkurl{q_k=1}. The radius factor is \nolinkurl{1+c√k δ<1+c√(k/60)}. The gamma ratio is \nolinkurl{(375.5/376)^f<1}; it may be retained or bounded by one.
\item
  For \nolinkurl{61≤k≤164}, the sparse constraint is active: \nolinkurl{q_k=w/k}. The radius is \textbf{exactly} \nolinkurl{1+c=2257/7}, independent of k and T.
\end{itemize}

In the second branch, retaining the gamma ratio gives the useful equivalent form \[
 T^{343s-1/4}R^{669f/4}
 k^{-(669/4-k)f}
 \left(\frac{375+k/(2w)}{376}\right)^f
 V_k(2257/7).
 \tag{SP8}
\] Since \nolinkurl{w>60}, the ratio in SP8 is bounded by the explicit rational number \[
 \frac{45000+k}{45120}.
 \tag{SP9}
\] An alternative, weaker branch discards the gamma ratio before substitution and has T exponent \nolinkurl{343s-5/4} with k power \nolinkurl{-(665/4-k)f}. The two branches must not be mixed.

\subsection{Exact finite optimization and the scoped removal}\label{exact-finite-optimization-and-the-scoped-removal}

Kummer's complete exact script \nolinkurl{recovery_pending/kummer/corrected_sparse_t1_exact_envelope_certificate.py} evaluates a conservative version of SP6: it bounds the favorable gamma ratio by one, uses the exact position multiplicity \nolinkurl{binom(164,k)}, and retains both distinct endpoint branches. For \nolinkurl{k≤60} it rounds the square-root radius upward to a rational with denominator \nolinkurl{2^24}; for \nolinkurl{k≥61} it uses the exact rational radius \nolinkurl{2257/7}. Each rational volume is bracketed by consecutive integral powers of two using integer arithmetic.

The common base-two constant before the T and weighted-volume terms is \[
 C_0=-30028313/8.
 \tag{SP10}
\] The script passed 496 assertions at 15:42:54.810217 UTC; my independent replay passed the same 496 assertions at 15:45:25.926502 UTC, with separate output \nolinkurl{review_sparse_t1_exact_envelope_certificate.json}.

The complete 164-row table has the following exact boundary values:

\begin{longtable}[]{@{}rr@{}}
\toprule\noalign{}
k & Upper exponent \\
\midrule\noalign{}
\endhead
\bottomrule\noalign{}
\endlastfoot
104 & \nolinkurl{-38889/8} \\
105 & \nolinkurl{-28209/8} \\
107 & \nolinkurl{-6761/8} \\
108 & \nolinkurl{3999/8} \\
164 & \nolinkurl{639295/8} \\
\end{longtable}

Every row with \nolinkurl{1≤k≤107} is at most \nolinkurl{-6761/8<-845}. Its complete actual sector therefore has absolute contribution \[
 \sum_{k=1}^{107} U_k
 <107\,2^{-845}<2^{-838}.
 \tag{SP11}
\] All original-coordinate support choices are already included in each \nolinkurl{U_k}; SP11 only sums the possible support sizes. Remove the intersection of this entire sector with the current remainder and retain every older error. The new charge is \textbf{838, all-top only}. The weaker \nolinkurl{k≤104} charge 4088 is an alternative sub-sector estimate and is not charged in addition to SP11.

Kummer whole-read the saved SP1--9 geometric proof at 15:43 UTC; the fresh script and complete exact table were independently read and replayed here. The resulting numerical addition SP10--11 has a separate exact check in \nolinkurl{review_sparse_t1_partial_sectors.json}.

The remaining part of this one profile has \nolinkurl{108≤k≤164}. Its positive upper estimates indicate failure of this sufficient bound, not a positive or nonzero actual contribution. No full signed inequality, prescribed period, Mram conclusion, literal Host/DP44 conclusion, or broader profile removal is asserted here. The main error ledger is not changed by this pending note.

\subsubsection{Canonical correction and zero-column clarification}\label{canonical-correction-and-zero-column-clarification}

The pre-correction496/500-check times in the retained historical discussion do not include the zero-target factor. The current corrected500-check script and root replay certify the displayed revised endpoints and838 charge; see DO070-C1 and the dated integration manifest. If the original S part of the negative column is zero while h\_SS=I, each W0 component of h(e\_j), j in S, is isotropic and orthogonal to the negative W0 vector h(e\_n). Its norm is therefore zero in a positive-definite orthogonal complement, so it vanishes. Such a matrix fixes S and belongs to the already removed fixed-W0 sector. Thus the scalar-cross sector has k\textgreater=1 as used above. The k\textless=107 charge838 is applied once; the overlapping older k\textless=104/4088 diagnostic is not also charged.

\clearpage

\section{AM-LIFT-063 --- explicit covering bounds for the original rings O\_E and O\_F}\label{am-lift-063-explicit-covering-bounds-for-the-original-rings-o_e-and-o_f}

\textbf{Canonical checked recovery version --- 2026-11-22T16:10:03.402339+00:00.} Root has read the complete mathematical proof and companion script and replayed the stated exact certificate. The original recovery export and its historical status are preserved separately. This is a current checked version, not recovery of unavailable historical final bytes. Only the stated sector and hypotheses are promoted; the signed period and actual DP44 cost remain open.

\subsection{1. Actual fields and integral tensor basis}\label{actual-fields-and-integral-tensor-basis}

Let E0=Q(zeta\_28), let B be the degree-seven fixed field in Q(zeta\_29) of the subgroup H of order four, and set E=E0 B Q(sqrt(5)) Q(sqrt(13)), F=E\^{}+. Then {[}E:Q{]}=336 and {[}F:Q{]}=f=168. All metrics below are the actual half-trace metric on E and the usual trace metric on F: \textbar\textbar x\textbar\textbar\_E²=(1/2)Tr\_E/Q(x bar(x)), \textbar\textbar y\textbar\textbar\_F²=Tr\_F/Q(y²).

The constituent discriminants are 2\^{}12*7\^{}10, 29\^{}6, 5 and 13, pairwise coprime. Their Galois intersections are unramified at every finite prime of Q and are therefore Q. Iterating gives the product degree. At each rational prime at most one constituent ring of integers ramifies. Tensoring its local DVR/product with the other finite étale orders is still finite étale over that normal local order, hence normal. Therefore the tensor product of the four integer bases is exactly O\_E, not a proper suborder.

For this local maximality statement, the only permanence input is that an étale algebra over a Dedekind domain is a finite product of Dedekind domains (Stacks Lemma15.45.4, Tag0AP2). The nonexistence of a nontrivial number field unramified over Q is Milne Theorem4.9. These inputs do not require a bound on a class number or any coefficient-field enlargement.

\subsection{2. The constituent Gram matrices}\label{the-constituent-gram-matrices}

For E0 use1,zeta,\ldots,zeta\^{}11, an integral basis by the cyclotomic integer-ring theorem (Milne Theorem6.4). Its half-trace Gram entry is c\_28(i-j)/2, where c\_28 is the Ramanujan sum. Reorder the basis by parity. The Gram matrix is the direct sum of two blocks 7 I\_6-u u\^{}t, u=(1,-1,1,-1,1,-1). Each block has eigenvalues1 and7, and the whole Gram determinant is7\^{}10. Thus its spectral maximum is7, and the absolute field discriminant is2\^{}12*7\^{}10.

For B, let C\_i be the seven cosets of H in F\_29\^{}× and put eta\_i=sum\_(a in C\_i) zeta\_29\^{}a. These seven sums are an integral basis: zeta\_29,\ldots,zeta\_29\^{}28 is an integral basis of O\_Q(zeta\_29) permuted by H, so its H-invariants have exactly the orbit sums as a Z-basis. Also sum\_i eta\_i=-1. As -1 belongs to H, each eta\_i is real. In the product eta\_i eta\_j there are four zero-sum pairs if i=j and none otherwise. Taking one fourth of the full cyclotomic trace gives Tr\_B/Q(eta\_i eta\_j)=(29\emph{4}delta\_ij-16)/4=29 delta\_ij-4. Thus the Gram matrix is29 I\_7-4 J\_7, with eigenvalues1 and29 and determinant29\^{}6. This agrees with the primary Gaussian-period integral-basis statement in Thaine (2005), Section2, printedp209.

For d=5,13 use the standard quadratic basis1,(1+sqrt(d))/2. Its trace Gram is {[}{[}2,1{]},{[}1,(d+1)/2{]}{]}. For d5 its largest eigenvalue is(5+sqrt(5))/2\textless4; for d13 it is(9+sqrt(29))/2\textless8. Equivalently,4I-G5 and8I-G13 are positive definite with determinants1 and5, respectively.

Because all embeddings of linearly disjoint constituent fields occur independently, the global half-trace Gram is exactly G\_E=G\_E0,half tensor G\_B tensor G\_5 tensor G\_13. The factor1/2 occurs once, on the CM factor E0. Hence lambda\_max(G\_E)\textless7\emph{29}4*8=6496.

\subsection{3. Rounding gives an original-ring covering bound}\label{rounding-gives-an-original-ring-covering-bound}

For a real basis with positive Gram matrix G, write an arbitrary point in that basis and round every coefficient to the nearest integer. The residual coefficient vector lies in{[}-1/2,1/2{]}\^{}n.~Its squared norm is at most lambda\_max(G)\emph{n/4. Applying this to the actual336-dimensional Minkowski lattice O\_E gives mu(O\_E)² \textless336}6496/4=3248f\textless57²f.

This is a bound for the ring lattice itself. It does not imply the same estimate for an arbitrary fractional ideal, a Galois-oriented J-ideal, or a rescaled lattice with unrelated archimedean shape.

\subsection{4. An equally strong bound for the real ring}\label{an-equally-strong-bound-for-the-real-ring}

Put t=zeta\_28+zeta\_28\^{}-1 and s\_j=zeta\_28\^{}j+zeta\_28\^{}-j. Since zeta²=t zeta-1, O\_E0=Z{[}t{]} direct-sum zeta Z{[}t{]}. The sum is direct because1,zeta are linearly independent over F0=Q(t). If an element a+zeta b is fixed by conjugation, then(2zeta-t)b=0, so b=0. Therefore O\_F0=Z{[}t{]}. The monic Chebyshev expressions s\_j=t\^{}j+(lower integral powers) show that1,s1,\ldots,s5 is an integral basis.

Its trace Gram splits by parity. On s1,s3,s5 it is14I\_3. On1,s2,s4 it is {[}{[}6,2,-2{]},{[}2,10,4{]},{[}-2,4,10{]}{]}. The latter has eigenvalues14,6+2sqrt(2),6-2sqrt(2). Thus the spectral maximum is14, and the determinant of the full real Gram is2\^{}6*7\^{}5.

All remaining tensor factors are real and Z-free. Taking conjugation invariants therefore gives O\_F=O\_F0 tensor O\_B tensor O\_Q(sqrt5) tensor O\_Q(sqrt13) exactly. Its trace Gram has spectral maximum\textless14\emph{29}4\emph{8=12992. Since its real dimension is f168, mu(O\_F)²\textless f}12992/4=3248f\textless57²f.

As a determinant check, the global Gram matrices give det G\_E=D0², D\_F=det G\_F=2\^{}f D0, D0=7\^{}140\emph{29\^{}144}5\^{}84*13\^{}84. No discriminant square root or half-trace normalization is omitted.

\subsection{5. The actual target fibre and dyadic real-part lattice}\label{the-actual-target-fibre-and-dyadic-real-part-lattice}

For the already identified target fibre containing the orthogonal sublattice O\_E\^{}4 direct-sum i O\_F, the inclusion and product-metric formula yield mu(fibre)² \textless=4mu(O\_E)²+mu(O\_F)²\textless16240f\textless128²f. The original negative column at H133 has norm\textless95sqrt(f). Consequently the previously justified affine-fibre counting enlargement may use the radius223sqrt(f).

Only this containing-superlattice consequence is asserted. The real scalar ideal(1/2)J\_F I\_F² still uses its separate375y bound; none of its fractional-ideal covering estimates is replaced.

There is also an exact original-field dyadic identity. The minimal polynomial of t is t\^{}6-7t\^{}4+14t²-7=0. Let p=t\^{}5-7t\^{}3+14t, so tp=7. The integral element a=zeta\_28 p-3 satisfies Tr\_E/F(a)=tp-6=1. Thus the integral trace map O\_E→O\_F is surjective and P0=Re(O\_E)=(1/2)O\_F exactly. This supplies the optional stronger target-fibre covolume without an additional relative-unramifiedness premise. The already staged conservative K085 bound is not silently altered.

\subsection{6. Verification and scope}\label{verification-and-scope-1}

The in-memory exact diagnostic reconstructs all Gram entries using Ramanujan sums and finite F\_29 cosets, checks both base-field determinants and the global determinant identities, the quadratic spectral bounds, all dimensions, and the integer rounding constants. It passed889 assertions at15:18:27 UTC. Its result is exported separately. The mathematical proof above is new bounded work, not an old-file restoration. The complete exported proof and the complete new certificate were read by the author after recovery; the file certificate passed 342 exact assertions at 2026-11-22 15:46:23 UTC. Its smaller count is a separate implementation, not a claimed reproduction of the 889-check in-memory script. The certificate explicitly distinguishes finite arithmetic checks from the mathematical field-maximality and source-fibre containment arguments.

This is a useful sufficient geometric bound for specific original-ring lattices. It alone removes no additional signed-kernel sector, proves no period nonvanishing, and resolves no literal source-family or DP44 obligation.

\subsection{Primary anchors}\label{primary-anchors}

Milne, J. S. (2020). Algebraic number theory (Version3.08). Theorem4.9, printedp72; Theorem6.4, printedp99. \url{https://www.jmilne.org/math/CourseNotes/ANT.pdf} Direct web references: turn636883448view0, turn636883448view2.

Thaine, F. (2005). Finding families of units of cyclic fields. RIMS Kokyuroku,1451,207--215. Section2, printedp209. \url{https://www.kurims.kyoto-u.ac.jp/~kyodo/kokyuroku/contents/pdf/1451-21.pdf} Direct web reference: turn636883448view1.

The Stacks Project Authors. Lemma15.45.4 (Tag0AP2), permanence under étale maps. \url{https://stacks.math.columbia.edu/tag/0AP2} Direct web reference: turn610427626view0.

\clearpage

\section{AM-DO-072 --- Gaussian counting of the remaining scalar-cross profiles}\label{am-do-072-gaussian-counting-of-the-remaining-scalar-cross-profiles}

\textbf{Canonical integration --- 2026-11-22T16:10:03.402339+00:00.} Whole root proof/script reading and separate exact replay are recorded in CHECK\_EVENTS.tsv. Original recovery bytes remain preserved.

Author: \nolinkurl{/root/definite_oscillator}, 22 November 2026 UTC.

This is new controlled-recovery work. It uses the actual all-top J56 kernel after the proved H=133 annular removal. The proof and finite endpoint certificate concern the two precisely specified support-eleven sectors below. They do not assert the full signed inequality, a period value, an arithmetic Host constructor, or DP44/RH. Existing main proof files and ledgers are unchanged.

\subsection{1. Statement and fixed data}\label{statement-and-fixed-data}

Retain the original normalized signed kernel of K062, the actual finite projector, class weights, packet point, real selected module and PBW metric. Put \[
 f=a=168,\quad n=169,\quad m=164,\quad s=11,\quad H=133,
 \qquad R=H+H^{-1}-2=17424/133.
 \tag{GS1}
\] Work on the actual scalar-cross part of the current remainder: in the original orthogonal decomposition \nolinkurl{E^169=S⊥W0}, with dimensions 164 and 5 and \nolinkurl{W0} of signature \nolinkurl{(4,1)}, write \[
 h=I+N,\quad h_{SS}=I,\quad
 \operatorname{rank}_E N=2,\quad N^3=0,\quad
 \operatorname{rank}_E N^2=1.
 \tag{GS2}
\] Let \nolinkurl{x=(h e_n)_S≠0}, and let k be the number of its nonzero original E-coordinates. Consider the disjoint sectors

\begin{itemize}
\tightlist
\item
  \nolinkurl{i=0,u=11,t=0}, with \nolinkurl{1≤k≤164};
\item
  \nolinkurl{i=0,u=10,t=1}, with \nolinkurl{108≤k≤164}.
\end{itemize}

Here I/II/III have their original pencil-incidence meanings, and \nolinkurl{i+u+t=s}. The second range is the part not removed by the separate sparse-sector bound AM-DO-071. No arbitrary moving integral basis is introduced.

The absolute contribution of the union of these 221 coordinate-support-size cells is less than \[
                         2^{-1482}.
 \tag{GS3}
\] Every original coordinate-support choice, finite orientation, support and incidence multiplicity is included. The conclusion is an absolute estimate on an actual subset of the original signed sum; no favorable sign is assigned to an uncomputed contribution.

\subsection{2. Actual column, paired ideals and metrics}\label{actual-column-paired-ideals-and-metrics}

Use the already proved direct pencil and nilpotent identities in DO070, the recovered direct-pencil note, R052/R053, and DO066's common-orientation argument. They give \[
 \alpha=(N^2)_{nn}\in F,\qquad
 0<\alpha_\tau\le R,\qquad
 \|x_\tau\|^2\le\alpha_\tau.
 \tag{GS4}
\] The last inequality is the actual original-coordinate bound: \nolinkurl{N²} has zero S rows and columns, its isotropic image line lies in W0, and the W0 component of \nolinkurl{N e_n} lies in that line's positive-semidefinite orthogonal complement. It does not divide a finite isotropic vector by alpha.

The omitted value \nolinkurl{k=0} is exactly a fixed-W0 term in this rank-two setting. Indeed \nolinkurl{N e_n} and \nolinkurl{N² e_n} are independent and span \nolinkurl{im N}: the second is nonzero because \nolinkurl{alpha>0}, while a dependence would force \nolinkurl{N²e_n=0} after applying N. If the S part of \nolinkurl{N e_n} vanishes, both spanning vectors lie in W0, so all S rows of N vanish. The identity \nolinkurl{N^dagger=N²−N} then forces all S columns of N to vanish as well. Hence \nolinkurl{h=I_S⊕h_0}. Such a term belongs to the separate fixed-W0 calculation and is not an uncounted scalar-cross cell.

At type II, alpha has pole order at most two at the active F-prime. At type III, the actual cyclic identity \nolinkurl{ZkZ=tr(Zk)Z}, together with global unipotence and the off-diagonal incidence, lowers this to at most one. These are assertions about the actual alpha, not a replacement of the full target-lattice projection by a smaller lattice.

At each of the 334 unmixed G places, the simultaneous rank-one orientation or rank-zero condition supplies a common ideal for all S-coordinates. Denote the oriented E-ideal by \nolinkurl{J_E} and its paired F-ideal by \nolinkurl{J_F}; their absolute norms are equal. There are at most \nolinkurl{4^334} common-orientation patterns. The actual Green estimate is \[
 |a(h)|\le2Q^{-335(f-1)}(N J_E)^{335/2},\qquad Q=g_0g_1.
 \tag{GS5}
\] The two mixed finite factors retain only their modulus bound one. Rank-one reductions force one alpha valuation; rank-zero reductions force two. Hence \[
 N J_F\le N_F(\alpha)T^{2u+t}.
 \tag{GS6}
\]

Let \nolinkurl{I_Σ} be the paired inverse-prime E-ideal over all s active T places, so \texttt{N\_E\ I\_Σ=T\^{}\{-2s\}}. Since the actual target shear is confined to W0, the S column has only simple poles at either CM partner. Thus \[
 x\in L^m,\quad L=J_E I_\Sigma,\quad N_E L=z^{2f}.
\] Set \[
 y=(N J_F T^{-(2u+t)})^{1/f},\qquad
 w=T^{t/f}.
\] Because \nolinkurl{i=0}, \nolinkurl{2u+t=2s-t}, and therefore \[
 z^2=y/w,\qquad 0<y\le R.
 \tag{GS7}
\]

The actual half-trace metric and covolumes are \[
 \|b\|_E^2=\tfrac12\operatorname{Tr}_{E/\mathbf Q}(b\bar b),
 \qquad \operatorname{covol}(O_E)=D_0,
\] \[
 D_0=7^{140}29^{144}5^{84}13^{84},\qquad
 D_F=2^fD_0.
 \tag{GS8}
\] The latter identity follows from the exact original-ring Gram determinants in LIFT063; it also agrees with the accepted finite-unramified CM discriminant calculation. No half-root-discriminant is substituted for the full root discriminant.

For a fractional ideal the actual Euclidean dual is \[
 L^\#=2\mathfrak D_E^{-1}\bar L^{-1}.
\] DO070's codifferent minimum and the primary finite-dimensional transference theorem consequently give, for every k, \[
 \mu(L^k)<c\sqrt k\sqrt f\,z,\qquad c=2250/7.
 \tag{GS9}
\] This is an arbitrary-ideal statement. The original-ring cover used for the W0 fibre below is a distinct input.

\subsection{3. The all-nonzero coordinate condition}\label{the-all-nonzero-coordinate-condition}

For a nonzero \nolinkurl{b∈L}, \nolinkurl{(b)⊂L}, so \[
 |N_{E/\mathbf Q}(b)|\ge N_E L=z^{2f}.
\] Arithmetic-geometric mean on the f CM absolute squares gives \[
 \|b\|_E^2\ge fz^2.
 \tag{GS10}
\] Thus every nonzero coordinate of x has this minimum. Summing over its k nonzero coordinates and using GS4 yields \[
                     kz^2\le R.
 \tag{GS11}
\] Define the dimensionless variable \[
                       p=R/(kz^2)\ge1.
\] Then \[
 y=Rw/(kp),\qquad y/R=w/(kp).
 \tag{GS12}
\] The extra constraint \nolinkurl{y≤R} is automatic for the ranges in the theorem once \nolinkurl{p≥1}: for t=0, \nolinkurl{w=1≤k}; for t=1, the actual prime satisfies \texttt{60\textless{}T\^{}\{1/168\}\textless{}61}, whereas \nolinkurl{k≥108}.

There are exactly \nolinkurl{binom(164,k)} possible original supports. They are included later, before summing over k. The argument does not declare all nonzero coordinates to be primitive or to attain their minimum simultaneously.

\subsection{4. A primary Gaussian bound that uses all k nonzero coordinates}\label{a-primary-gaussian-bound-that-uses-all-k-nonzero-coordinates}

Use Miller and Stephens-Davidowitz (2019), Lemma 3.11, with Euclidean dimension \nolinkurl{2f=336}, exponent 2, and Gaussian \nolinkurl{exp(-||v||²)}. Its proof and the rearrangement in the proof of Theorem 3.12 were directly read. This is a Euclidean lattice Gaussian sum, not a change of the automorphic theta normalization.

Put \[
 b=f+1/4=673/4,\qquad
 \eta=(e b/f)^f e^{-b}
      =(673/672)^{168}e^{-1/4}.
 \tag{GS13}
\] Scale L by \nolinkurl{sqrt(b/(fz²))}. By GS10 its shortest nonzero vector has length at least \nolinkurl{sqrt(b)}. Lemma 3.11 has the closed tail condition \nolinkurl{||lambda||≥t sqrt(n/2)}, so vectors attaining the minimum are included: take its \nolinkurl{n=336}, Gaussian exponent 2, translation zero and \nolinkurl{t=sqrt(b/f)>1}. The zero vector is outside that tail. Every nonzero vector lies in it. If \nolinkurl{Θ×} is the Gaussian sum over those nonzero vectors, the lemma gives \[
 \Theta^\times\le\eta(1+\Theta^\times),\qquad
 \Theta^\times\le\frac{\eta}{1-\eta}<8192.
 \tag{GS14}
\] The strict last inequality is certified rationally: the degree-20 lower exponential series proves \[
 e^{1/4}>(8193/8192)(673/672)^{168}.
\] In particular eta is strictly below one. No asymptotic kissing-number constant is used.

For k nonzero coordinates with total squared norm at most \nolinkurl{fR}, every associated product Gaussian weight is at least \nolinkurl{exp(-bR/z²)}. The sum of all such product weights is at most \nolinkurl{(Θ×)^k}. Consequently their number is at most \[
                \exp(bR/z^2)8192^k
                     =\exp(bkp)8192^k.
 \tag{GS15}
\] This bound uses the simultaneous energy and the nonzero condition in every coordinate. A product of independent coordinate-ball counts would lose that information.

\subsection{5. The retained Voronoi alternative}\label{the-retained-voronoi-alternative}

The genuine original-coordinate ball has radius \nolinkurl{sqrt(fR)}. By GS9 its count for a fixed support is also at most \[
 V_k(\sqrt R+c\sqrt k\,z)z^{-2fk},\qquad
 V_k(C)=\frac{((22/7)fC^2)^{fk}}{(fk)!D_0^k}.
\] Here \nolinkurl{22/7} supplies a rational upper envelope for the true ball volume. Using GS12, this becomes \[
                  (kp)^{fk}V_k(1+c/\sqrt p).
 \tag{GS16}
\] GS15 and GS16 are alternative upper bounds for the same actual set. They will be combined by cases in p; they are not multiplied.

\subsection{6. The full gamma and W0 fibre count}\label{the-full-gamma-and-w0-fibre-count}

Fix an actual labelled source-target pencil pair at every active T place. The admitted labels are \texttt{D=k\_T\textbackslash{}\{-1\}}, including zero, by K062(AT6) and LIFT44(PR4--7). Use the exact W0 lattice presentation of K063(PF12) and K071(NB2--3), \[
 M_c=(I+c\varpi^{-1}Z)\Lambda_0,\qquad
 Z=zq,\quad Z^2=0,\quad \Lambda_0=O^5.
\] K077(AP1--3) gives both negative entries \nolinkurl{z_n,q_n} as units. At the conjugate CM component the pole column and row are the inverse-adjoint dual ones; their negative entries are again units. Thus the same calculation applies at both partners without identifying their two matrices.

For an incidence \nolinkurl{hM_c=M_d}, the containment \nolinkurl{varpi e_n∈M_c} gives \texttt{h\ e\_n∈varpi\^{}\{-1\}M\_d}. Its W0 component is therefore in the actual target overlattice \[
 B_d=\varpi^{-1}(I+d\varpi^{-1}Z)\Lambda_0.
\] The inverse of its displayed generator is \nolinkurl{varpi I−dZ}, which is integral. Consequently \nolinkurl{Lambda_0⊂B_d}, and the determinant-one shear gives local determinant valuation −5. If \nolinkurl{d≠0} in the residue field, the negative diagonal of the generator has a genuine pole-two term \texttt{d\ varpi\^{}\{-2\}z\_nq\_n}; its pole-one term cannot cancel it. If \nolinkurl{d=0}, choose its zero lift and \texttt{B\_0=varpi\^{}\{-1\}Lambda\_0}. Hence the exact negative-coordinate projection is \texttt{varpi\^{}\{-2\}O} for a nonzero target and \texttt{varpi\^{}\{-1\}O} for target zero, at both CM partners.

A type-II target is necessarily nonzero: fixing \nolinkurl{M_0=Lambda} would make h integral and exclude that active place. A type-III target can be zero. Let \nolinkurl{r_0} count its zero targets. Patch these actual target lattices to obtain \nolinkurl{B_Σ⊂E^5}. It contains \nolinkurl{O_E^5}, has covolume \texttt{D0\^{}5T\^{}\{-10s\}}, and its real-last image is \[
 P_\Sigma=P_0\!\prod_{d\ne0}\mathfrak q^{-2}
                       \prod_{d=0}\mathfrak q^{-1},
 \qquad P_0=\operatorname{Re}(O_E).
\] All active primes are split and odd, so taking real part is the sum of the two independent local coordinates divided by the unit 2. Away them the image remains \nolinkurl{P0}. Localization therefore proves the displayed global image, whose covolume is \texttt{covol(P0)\ T\^{}\{-2s+r\_0\}}. This is the full target image, kept separate from the extra simple-pole condition on the actual unipotent gamma below.

For real-part projection the normal Jacobian is one, since it is an orthogonal coordinate projection in the actual half-trace metric. Coarea gives \nolinkurl{covol(kernel)=covol(B_Σ)/covol(P_Σ)}. Since \nolinkurl{O_F⊂P0⊂(1/2)O_F}, one has \nolinkurl{covol(P0)≤sqrt(D_F)}, and therefore \[
 \operatorname{covol}(K_\Sigma)
 \ge2^{-f/2}D_0^{9/2}T^{-8s-r_0}.
 \tag{GS17}
\] More precisely, LIFT063 gives \nolinkurl{P0=(1/2)O_F} and the exact value is \texttt{2\^{}\{f/2\}D0\^{}\{9/2\}T\^{}\{-8s-r\_0\}}. This extra factor \nolinkurl{2^f} is deliberately not used in GS17 or in the numerical theorem.

The kernel contains the orthogonal original-ring sublattice \nolinkurl{O_E^4⊕iO_F}. LIFT063 proves the actual original-ring covers, hence kernel cover below \nolinkurl{128√f}. The original negative column has squared Euclidean norm \texttt{cosh(2r)≤(H²+H\^{}\{-2\})/2\textless{}95²} at every label. Fixing its real last coordinate can only shrink the orthogonal slice radius. Thus the affine fibre count can use radius \nolinkurl{223√f}. Combining this with GS17 gives \[
 \#\{W_0\text{ columns in a fixed real-part fibre}\}
 <2^{5379}T^{8s+r_0}.
 \tag{GS18}
\] The exact rational-pi envelope lies strictly between \nolinkurl{2^5378} and \nolinkurl{2^5379}. No ring-only cover is applied to a fractional J-ideal.

For the actual unipotent column, put \[
 \gamma=\operatorname{Re}(h_{nn})-1=\alpha/2.
\] The type-II/type-III alpha pole distinction and the G divisibility give \[
 \gamma\in\tfrac12J_F I_{\rm II}^2I_{\rm III},
 \qquad 0<\gamma_\tau\le R/2.
 \tag{GS19}
\] The dyadic factor one half is retained. The ideal on the right has covolume \texttt{2\^{}\{-f\}sqrt(D\_F)y\^{}f}; its established fractional-ideal cover remains below \nolinkurl{375√f y}. Hence the gamma count is \[
 A_f(R/2+375y)^f y^{-f},\qquad
 A_f=\frac{\pi^{f/2}f^{f/2}2^f}{(f/2)!\sqrt{D_F}}.
 \tag{GS20}
\] Unlike the earlier deliberately coarse bound, retain the full discriminant factor here. GS8 and exact integer comparison give \[
                         A_f376^f<2^{1062}.
 \tag{GS21}
\] The rational-pi envelope is above \nolinkurl{2^1061}, so 1062 is the certified ceiling. Using alpha instead of gamma would multiply both the normal Jacobian and the image-lattice covolume by the same \nolinkurl{2^f}; no gain from that reparametrization is claimed.

For each full negative column, retain the actual source-lattice stabilizer fibre \nolinkurl{G_n<2^5022}. The fixed-W0 bound \nolinkurl{G_5} is not substituted. Multiplying S, gamma and W0 fibre counts is an upper bound; the actual unitary columns need not vary independently.

\subsection{7. Actual real compression and all finite multiplicities}\label{actual-real-compression-and-all-finite-multiplicities}

The actual selected compact module has total homogeneous degree greater than \nolinkurl{2^959}. K085's checked positive-energy circle/orthogonal-ladder calculation and the actual near-point comparison apply to the full global rank-two matrix, not just a fixed-W0 matrix. They give \[
 |c(h)|\le12^f2^{f/2}2^{-959f/4}(N_F\alpha)^{-1/4}
          \le12^f2^{f/2}2^{-959f/4}y^{-f/4}.
 \tag{GS22}
\] The image line of \nolinkurl{N²} is normalized over C only in that proof. No finite integral chart for a divided isotropic vector is assumed. The constant is in the actual Hilbert normalization and at the actual packet point.

For the two present profiles every active place is II or III. Put \nolinkurl{m_T=T−1}. For type II, the number of diagonal labels times its positive rho is at most one, and the target is nonzero. At type III, each off-diagonal ordered pair has absolute weight \nolinkurl{1/[m_T(m_T−1)]}. There are \nolinkurl{(m_T−1)^2} pairs with nonzero target, and \nolinkurl{m_T−1} pairs with target zero. The latter have the additional T in GS18. Their full weighted fibre factor, relative to \texttt{T\^{}\{8s\}}, is exactly \[
 \frac{m_T-1}{m_T}+\frac{T}{m_T}=2.
\] All other bounds are uniform in the actual pairs. Thus the complete II/III pair sum contributes the factor \nolinkurl{2^t}, and no factor involving T remains. This repairs the inherited zero-target oversight in the old general-III projection sentence; the separate AM-DO-070-C1 note records its effect on DO070 and DO071. No inversion of h or changed original coordinates is used.

Retain once each: support factor \nolinkurl{2^168}, conservative type-position factor \nolinkurl{3^11<2^22}, common-orientation factor \nolinkurl{4^334=2^668}, Green constant 2, full source fibre \nolinkurl{2^5022}, and the support choices \nolinkurl{binom(164,k)}. No additional \nolinkurl{3^334} rank-pattern count is introduced.

\subsection{8. Exact combined expression}\label{exact-combined-expression}

Put \nolinkurl{P=669f/4=28098}, and define \[
 H_{t,k}(p)=\left(\frac{375+kp/(2w)}{376}\right)^f,
 \quad w=T^{t/f}.
\] Combining GS5--7 and GS18--22, the common factor for a fixed coordinate-support size is bounded by \[
\begin{split}
\mathcal C_{t,k}={}&
2^{1+5379+1062+5022+168+22+668+t}
12^f2^{f/2-959f/4}Q^{-335(f-1)}\\
&\cdot T^{343s-t/4}R^P\binom{164}{k}.
\end{split}
 \tag{GS23}
\] After this common factor, the Gaussian and Voronoi alternatives are respectively \[
 G_{t,k}(p)=(kp)^{-P}H_{t,k}(p)
                  \exp(bkp)8192^k,
 \tag{GS24}
\] \[
 V_{t,k}(p)=(kp)^{-(P-fk)}H_{t,k}(p)
                     V_k(1+c/\sqrt p).
 \tag{GS25}
\] For example, the T exponent is \nolinkurl{κ(2s-t)+8s+(κ-1/4)t=343s-t/4}, with \nolinkurl{κ=335/2}. The gamma factor has already been rewritten as \nolinkurl{A_f(375+R/(2y))^f}. This accounts for its \texttt{y\^{}\{-f\}} and is why neither an extra T pole gain nor an extra gamma density factor is present in GS23.

\subsection{9. Three endpoints control the entire continuum of ideal sizes}\label{three-endpoints-control-the-entire-continuum-of-ideal-sizes}

The variable p depends on the actual ideal pattern. Its possible values are not enumerated or assumed dense. A uniform bound suffices for each already counted pattern.

Differentiate the logarithm of GS25 with respect to \nolinkurl{log p}. Its derivative is \[
 -(P-fk)
 +f\frac{kp/(2w)}{375+kp/(2w)}
 -fk\frac{c/\sqrt p}{1+c/\sqrt p}
 <-f(665/4-k)<0.
 \tag{GS26}
\] Thus the Voronoi alternative decreases in p for every k in the theorem.

With \nolinkurl{v=log p}, the logarithm of GS24 is a constant plus \[
 -Pv+f\log(375+(k/(2w))e^v)+bk e^v.
\] Its second derivative is positive. It is convex, so on \nolinkurl{1≤p≤7} its maximum is attained at p=1 or p=7. For \nolinkurl{p≥7}, use the decreasing Voronoi alternative. Therefore every actual pattern satisfies the uniform bound \[
 \min(G_{t,k}(p),V_{t,k}(p))
 \le\max\{G_{t,k}(1),G_{t,k}(7),V_{t,k}(7)\}.
 \tag{GS27}
\] There is no extra factor three: each actual pattern uses one uniform maximum, and all its multiplicity has already been retained in GS23. Counting a sum of three maxima would be unnecessary.

\subsection{10. Exact rational endpoint certificate}\label{exact-rational-endpoint-certificate}

The full endpoint certificate, including the corrected \nolinkurl{2^t} pair/fibre factor, uses only the following one-sided rational inequalities: \[
 R^{512}<2^{3602},\quad \log_2e<13/9,\quad
 \sqrt7>21/8,\quad
 1+c/\sqrt7<6049/49,
 \tag{GS28}
\] and the actual prime window \nolinkurl{2^995<T<2^996}, \nolinkurl{Q^4>2^281}. The logarithm inequality is certified by a degree-24 exponential upper series at \nolinkurl{9/13}. The actual selected degree is checked from the restored K044 upper-gap integer.

For the gamma ratio at the three endpoints, use the actual lower bound \nolinkurl{w≥1} for t=0 and \nolinkurl{w>60} for t=1. The ratio is then bounded by the exact rational number \[
 \frac{750w_0+kp}{752w_0},\qquad
 w_0=1\text{ or }60.
 \tag{GS29}
\] The complete binomial factor and this ratio to the power f are included before taking an integral logarithmic ceiling. Every rational volume and inverse integer power is bracketed by consecutive powers of two. Gaussian exponentials use the rational exponent \nolinkurl{(13/9)(673/4)kp}, and \texttt{8192\^{}k=2\^{}\{13k\}}.

The common rational base-two constant before the T and endpoint terms is \[
\begin{split}
C_0={}&1+5379+1062+5022+168+22+668
 +P(3602/512)\\
&-(281/4)335\cdot167
 +168(9/2-959/4)
 =-481236791/128.
\end{split}
 \tag{GS30}
\] The script checks all three endpoints for every one of the 164+57=221 cells; the base for each t is \nolinkurl{C0+t+996(343s−t/4)}. The largest upper exponent is at t=0,k=1: \[
 E_{0,1}= -1716943/1152<-1490.
 \tag{GS31}
\] For that cell the other two endpoints are \nolinkurl{-90906895/1152} and \nolinkurl{-10099639/128}. Every other cell is no larger than GS31. Thus \[
 \sum_{\text{221 cells}}|\text{actual contribution}|
 <221\,2^{-1490}<2^{-1482},
\] proving GS3.

The first complete file replay passed 1,127 exact assertions at 15:55:52.930901 UTC before the zero-target correction. That correction adds one bit to each t=1 endpoint and leaves the stated worst cell and joint charge unchanged. The corrected complete replay passed all 1,127 assertions at 16:06:10.440982 UTC. The final whole-proof check and exact frozen fingerprints are recorded in \nolinkurl{gaussian_scalar_cross_verification.json}. The script checks numerical inequalities and endpoint exhaustion; it does not prove the Gaussian tail theorem, original-ring maximality, or the actual kernel interfaces automatically. These are the separate mathematical inputs proved or explicitly cited above.

\subsection{11. Removal and remaining scope}\label{removal-and-remaining-scope-1}

Remove the intersection of the stated union with the current remainder and append the new all-top charge 1482 while retaining every older error. The t=1 range k≤107 is the disjoint sector of DO071 and is not charged again here. DO070 handles the distinct t=2 profile. Any further assertion that all support-eleven structural classes have been exhausted must also retain the separately proved fixed-W0 and nonscalar-block reductions and their charges.

This proof changes neither the kernel nor its finite/archimedean normalizations. It supplies no favorable sign for later supports, no complete signed-sum inequality, no prescribed global Mram period, and no new arithmetic family or DP44/RH conclusion. A complete updated error-ledger replay and main-index promotion are separate operations.

\subsection{References and retained interfaces}\label{references-and-retained-interfaces}

Miller, S. D., \& Stephens-Davidowitz, N. (2019). \emph{Kissing numbers and transference theorems from generalized tail bounds} (arXiv:1802.05708v2). \url{https://arxiv.org/pdf/1802.05708}. Directly read: definitions; Theorem 3.2 and equation (3.3), printed p.8; Lemma 3.11, printed p.10; the proof of Theorem 3.12, printed p.12. The Gaussian rearrangement here is an explicit specialization, not an invocation of an unspecified asymptotic bound.

LIFT063, \nolinkurl{recovery_pending/lift/explicit_original_ring_cover_pending.md}, and its primary metric addendum: exact original E/F rings and Gram determinants, original-ring covering radii, and \nolinkurl{Re(O_E)=(1/2)O_F}. Root and Kummer completed independent whole proof/source checks before this result uses those inputs. Its ring cover is used only for a containing target-fibre lattice.

Retained actual-kernel interfaces: DO066 common-orientation and original-coordinate fibre count; DO070 FT1--12; the recovered direct pencil lemma; K085 actual selected-energy/near-point compression; the fresh R048 all-top H133 cutoff. The companion source identifiers and normalization distinctions remain those of the accepted run. No stipulated theorem proof is re-audited and no Lean build is claimed.

\clearpage

\section{Recovery record: direct pencil and scalar-cross lemmas}\label{recovery-record-direct-pencil-and-scalar-cross-lemmas}

\textbf{Required correction.} Read the complete AM-DO-070-C1 ZC1--11, included in this checked part. Any statement below that every type-III target has double-pole real projection is superseded: target0 has only one pole, and complete pair-weighted fibre averaging costs2\^{}t. The old sparse839 charge is superseded by838. Original calculations and historical check times below are retained for provenance, not as the corrected final constants.

Author: /root/definite\_oscillator. Controlled export, 22 November 2026. This reconstructs the complete mathematical arguments checked in messages during the outage. It is not a recovered old file and carries no old frozen-file hash. The new saved text and arithmetic replay must be verified separately.

\subsection{1. Actual local pencil, with no restriction to fixed W0}\label{actual-local-pencil-with-no-restriction-to-fixed-w0}

At a split T place write \[
 H_c=I+c\pi^{-1}Z,\qquad M_c=H_c\Lambda,\qquad
 Z=zq,\quad Z^2=0,\quad k=H_d^{-1}hH_c\in GL(\Lambda).
\] The rank-one operator Z acts only in W0. The source and target labels c,d are the actual nonzero pencil labels, with the stated omitted endpoints. Work at both CM partner primes in the actual dual coordinates.

Suppose h=I+N is globally unipotent of rank two in signature (a,1). Then N\^{}3=0, rank N\^{}2=1, and \[
 N^2=h+h^{-1}-2I=N+N^\dagger.
\]

\subsubsection{Type I: full-pencil preservation}\label{type-i-full-pencil-preservation}

The actual sum and intersection satisfy \[
 L_-\subset\Lambda\subset L_+,\qquad \dim_k(L_+/L_-)=2.
\] Preservation of the full pencil preserves both L\_- and L\_+. The induced action of h on the two-dimensional residue quotient is unipotent, hence its square augmentation vanishes. Therefore \[
 N^2L_+\subset L_-,\qquad N^2\Lambda\subset\Lambda.
\] This is an entrywise integrality statement for N\^{}2, not for h.

Moreover S is contained in L\_-. Consequently h and h\^{}\{-1\} have integral S-columns at both CM partners. The actual unitary adjoint identifies the partner inverse S-columns with the original S-rows. Hence the S portion x=(he\_n)\_S is integral at type I.

\subsubsection{Type III: unequal pencils with off-diagonal incidence}\label{type-iii-unequal-pencils-with-off-diagonal-incidence}

Here c-d is a unit. Cyclicity gives \[
 \operatorname{tr}(h)=\operatorname{tr}(k)
       +(d-c)\pi^{-1}\operatorname{tr}(Zk).
\] Since h is globally unipotent, its trace is the integer n, so tr(Zk) is divisible by pi. The exact rank-one identity \[
 ZkZ=\operatorname{tr}(Zk)Z
\] then shows that the apparent pi\^{}\{-2\} term in H\_d k H\_c\^{}\{-1\} has only a simple pole. The same argument applies to h\^{}\{-1\}. Therefore N\^{}2 has at most a simple pole at this partner and at its dual partner. This argument uses the actual trace and actual pencil expansion, not an analogy with fixed W0.

\subsubsection{Type II: the unique fixed incidence}\label{type-ii-the-unique-fixed-incidence}

The same two-shear expression has at most a double pole. No cancellation of its leading term is asserted.

Thus the entrywise square-pole table is I:0, II:2, III:1 for the scalar-cross class as well as fixed W0.

\subsection{2. The original S-column lattice}\label{the-original-s-column-lattice}

At an arbitrary actual pencil pair, \[
 \pi e_n\in M_c,\qquad hM_c=M_d
 \ \Longrightarrow\ he_n\in\pi^{-1}M_d.
\] The shear acts only on W0, so the S projection lies in pi\^{}\{-1\}O\_E\^{}m at both CM partners. No integrality of the whole matrix follows.

For global unipotent type I the preceding stronger S-integrality applies. If v=u+t counts types II and III, and J is the actual common-orientation G ideal, then \[
 x=(he_n)_S\in (J I_v)^m,\qquad N_EI_v=T^{-2v}.
\] All ideal directions are in the original coordinates. No inverse minor or changed rational frame is declared integral.

\subsection{3. Alpha, its positivity, and the sharper S radius}\label{alpha-its-positivity-and-the-sharper-s-radius}

Let \[
 \alpha=(N^2)_{nn}\in F.
\] The line L=im N\^{}2 is isotropic, ker N\^{}2=L\^{}\perp, and im N lies in L\^{}\perp. A negative vector e\_n cannot lie in L\^{}\perp. The vector Ne\_n is not in L, because otherwise N\^{}2e\_n=0. The induced form on L\^{}\perp/L is positive definite. Hence \[
 \alpha=\langle Ne_n,Ne_n\rangle_J>0
\] at every real embedding.

For the actual operator cutoff \textbar h\textbar{}\le H, \[
 0<\alpha_\tau=2\operatorname{Re}(h_{nn})-2
      \le H+H^{-1}-2=:R.
\]

If A=h\_\{SS\}=I, the SS block of N\^{}2=N+N\^{}\dagger is zero. Since N\^{}2 is a nonzero selfadjoint rank-one operator, its isotropic image line has zero S coordinates and lies in W0. The W0 portion of Ne\_n lies in the positive-semidefinite space L\^{}\perp inside W0. Therefore \[
 \|x_\tau\|^2\le\alpha_\tau\le R.
\] This uses only a complex orthogonal decomposition to prove an archimedean estimate. It does not construct an integral isotropic basis.

At an unmixed G place, residue rank one implies N\^{}2 is divisible by the prime, and residue rank zero implies divisibility by its square. For the common-orientation ideal J, alpha is divisible by J bar(J). Combining this with the II/III pole table gives \[
 N_E(\alpha)\ge (N_EJ)^2T^{-4u-2t},\qquad
 N_E(\alpha)\le R^{2f},
\] and therefore \[
 N_EJ\le R^fT^{2u+t}.
\]

\subsection{4. Scalar-cross alpha slice: exact symbolic calculation}\label{scalar-cross-alpha-slice-exact-symbolic-calculation}

Restrict now to all-top J56, H=133, f=168, m=164, s=11, i=0, u+t=s, and global rank-two scalar-cross unipotents. Retain the full source-lattice fibre bound G\_n\textless2\^{}5022; fixed-W0 G\_5 is not available.

Let J\_E be a common-orientation ideal, and let J\_F be its paired base-field ideal: residue rank one contributes one base prime, residue rank zero contributes its square. Their norms agree. Put \[
 d=2u+t=2s-t,\qquad
 y=(NJ_F\,T^{-d})^{1/f}\le R,\qquad
 z=(NJ_E)^{1/(2f)}T^{-s/f}
    =\sqrt y\,T^{-t/(2f)}.
\]

The full W0 target overlattice has real-part image P\_0I\_\Sigma\^{}2 and kernel covolume with T exponent -8s. This is unchanged at type III: its full overlattice image can still have a double pole.

By contrast, the actual unipotent real part \[
 \gamma=\operatorname{Re}(h_{nn})-1=\alpha/2
\] lies in \[
 \tfrac12J_F I_{\rm II}^2 I_{\rm III}.
\] Its simple III pole is a condition on the actual subset, not a replacement of the full projection image. This distinction prevents counting the same pole saving twice. The factor one half retains all dyadic effects; every G and T prime here is odd.

The gamma ideal has covolume \(2^{-f}\sqrt{D_F}\,y^f\) and covering radius at most \(375\sqrt f\,y\). Thus its points in the actual gamma range are bounded by \[
 A_f(R/2+375y)^f y^{-f},\qquad
 A_f=\frac{\pi^{f/2}f^{f/2}2^f}{(f/2)!\sqrt{D_F}}.
\]

The S-column belongs to \((J_EI_\Sigma)^m\), has radius at most \(\sqrt{fR}\), and has covolume scale z. Its bound is \[
 V_m(\sqrt R+9750z)\,z^{-2fm}.
\] The full original W0 real-part fibre contributes at most \(2^{9850}T^{8s}\) under the old general covering estimate. The new ring-only covering result, if separately accepted, may replace this constant; it does not replace the fractional S-ideal estimate.

Combine the Green loss \((NJ_E)^{335/2}\), the actual homogeneous-energy real compression, the gamma count, and the S-column count. The y-dependent expression, before its maximum, is \[
 y^{9f/4}(R/2+375y)^f
 V_m\!\left(\sqrt R+9750\sqrt y\,T^{-t/(2f)}\right).
\] It is increasing for y\textgreater0. At y=R the bound has the exact factor \[
 T^{343s-\frac72t}\,
 R^{13f/4}\,
 V_m\!\left(\sqrt R(1+9750T^{-t/(2f)})\right).
\] All \(4^{334}\) common-orientation patterns are retained once. No additional \(3^{334}\) rank-pattern factor is introduced. The actual II and III pair/weight sums are each at most one. The remaining constants include supports, type assignments, the full G\_n fibre, the Green factor 2, Q\^{}\{-335(f-1)\}, gamma volume, and the actual selected homogeneous-energy factor \(12^f2^{f/2}2^{-959f/4}\).

Under the old W0 covering constant the initial floating diagnostic remains positive by roughly twenty thousand bits at t=2. This is a failure of that sufficient estimate, not a lower bound on the signed contribution and not a nonvanishing or obstruction theorem. Exact diagnostics and any substitution of the new ring-only bound must be separately recorded.

\clearpage

\section{Pending exact sparse-coordinate evaluation of the all-top scalar-cross t1 profile}\label{pending-exact-sparse-coordinate-evaluation-of-the-all-top-scalar-cross-t1-profile}

\textbf{Required correction.} Read the complete AM-DO-070-C1 ZC1--11, included in this checked part. Any statement below that every type-III target has double-pole real projection is superseded: target0 has only one pole, and complete pair-weighted fibre averaging costs2\^{}t. The old sparse839 charge is superseded by838. Original calculations and historical check times below are retained for provenance, not as the corrected final constants.

Author and independent arithmetic checker: /root/kummer\_extension. Controlled recovery work, 22 November2026. This note does not modify DO070 or the main ledger.

\subsection{1. Scope and geometric input}\label{scope-and-geometric-input}

Retain the actual all-top J56 support-eleven globally unipotent rank-two scalar-cross profile \((i,u,t)=(0,10,1)\), with \(f=168,m=164,H=133,R=17424/133\). Let k be the number of nonzero original E-coordinates in the S part of the negative column. Then \(1\le k\le164\).

I whole-read the saved companion \nolinkurl{recovery_pending/definite/sparse_support_scalar_cross_pending.md}, SP1--9, SHA256 b858b47ff561c2db98171da1248f6b2e17420d4b410fe6bc33963abb94ee7466. Its original-coordinate norm argument, fractional product cover, monotonicity, and both exact cap branches check.

In detail, a nonzero coordinate in the actual fractional ideal L satisfies \(|N_E x_j|\ge N_EL=z^{2f}\); CM arithmetic-geometric mean gives trace norm at least \(fz^2\). Summing the k nonzero coordinates gives \(kz^2\le R\). This does not choose a new rational or integral basis.

Use DO070's fractional covering constant \(c=2250/7\) and the exact rational volume envelope \[
 V_k(C)=\frac{(528C^2)^{fk}}{(fk)!D_0^k},\qquad
 D_0=7^{140}29^{144}5^{84}13^{84}.
\] The cover is \(c\sqrt k\sqrt f\,z\), rather than the fixed164-coordinate bound.

\subsection{2. The conservative branches actually certified}\label{the-conservative-branches-actually-certified}

Put \(w=T^{1/f},\delta=T^{-1/(2f)},q=y/R\). The actual feasible endpoint is \[
 q_k=\min(1,w/k),\qquad 60<w<61.
\] After retaining all actual support, type, orientation, source-lattice fibre, W0 fibre, selected real, and gamma constants, the k-dependent factor is bounded by \[
 T^{343s-335/2+k}R^{669f/4}
 q_k^{(665/4-k)f}
 V_k(1+c\sqrt{kq_k}\,\delta).
 \tag{SE1}
\] Here I discard only the favorable ratio \(((1/2+375q_k)/376)^f\le1\). The exact geometric companion also states the stronger branch retaining this ratio; it is not mixed with SE1.

The position multiplicity is the exact \(\binom{164}{k}\), included once.

For \(k\le60\), use \(q_k=1\) and \[
 T^{343s-335/2+k}V_k(1+c\sqrt{k/60}).
 \tag{SE2}
\] The square root in this expression is enclosed above by an explicit dyadic rational of denominator \(2^{24}\).

For \(k\ge61\), substitution of \(q_k=w/k\) gives exactly the bound \[
 T^{343s-5/4}\,
 k^{-(665/4-k)f}\,
 V_k(2257/7).
 \tag{SE3}
\] The k exponent is the positive integer \(27930-168k\). This cancellation uses the actual T power before replacing T by its upper bound.

\subsection{3. Exact complete evaluation}\label{exact-complete-evaluation}

The common binary constant, including the conservative original support/type/orientation factors but excluding the exact position multiplicity already inside the rational volume, is \[
\begin{split}
C={}&1+9850+1945+5022+168+22+668\\
 &+\frac{669f}{4}\frac{113}{16}
 -\frac{281}{4}335\cdot167
 +f\left(\frac92-\frac{959}{4}\right)
 =-\frac{30028321}{8}.
\end{split}
\tag{SE4}
\]

The companion script computes the rational volume and k-power exactly as integers, including \(\binom{164}{k}\). It finds the unique integer ceiling L such that \(2^{L-1}<A\le2^L\) by direct shifts and comparisons. All T exponents used are positive, so \(T<2^{996}\) gives the displayed upper exponent.

The first actual execution passed \textbf{496 assertions at 2026-11-22T15:42:54.810217+00:00}. Every k from1through164 was included. The final scoped arithmetic addition also verifies the exact107-coordinate union below. Results:

\begin{itemize}
\tightlist
\item
  For k1through104, the conservative upper exponent is strictly below−4096.
\item
  For k1through107, the largest upper exponent is exactly \(-6769/8<-846\).
\item
  At k108, the displayed upper exponent is \(3991/8>0\).
\item
  The largest displayed upper exponent is \(639287/8\), at k164.
\end{itemize}

Consequently the larger sector with k1through107, intersected with the current remainder, has absolute contribution less than \[
107\cdot2^{-846}<2^{-839}.
\tag{SE5}
\] This is the proposed new partial-sector charge. The smaller k1through104 sector alternatively has the stronger bound \[
104\cdot2^{-4096}<2^{-4088}.
\tag{SE6}
\] The4088 alternative is not charged in addition to839 for an overlapping subset. Either use requires parent integration and preservation of every earlier error. Neither is a whole-profile removal.

The positive dense-coordinate envelope is only a failure of this sufficient estimate. It is not a lower bound on the actual absolute contribution, not a sign statement, and not evidence of nonvanishing.

\subsection{4. Verification and remaining scope}\label{verification-and-remaining-scope}

The entire new certificate source was read after its execution. The mathematical branch formulas were independently derived and compared with SP1--9. The algebra, exact prime-window split, positivity of every monotone exponent, position multiplicities, integer volume calculation, and all164 endpoints check.

The new files are: - \nolinkurl{sparse_t1_exact_envelope_certificate.py}; - \nolinkurl{sparse_t1_exact_envelope_certificate.json}; - this pending note.

No original proof, old certificate, main index, or ledger was changed. A complete old-error-list replay is not included. The all-T Mram and signed sufficient inequality remain open. Neither the failed literal Host interface nor the full DP44 cost comparison is addressed.

Final author whole-note and complete-script reread: \textbf{PASS at 2026-11-22 15:47:21 UTC}. The final actual replay passed500 assertions at2026-11-22T15:47:03.624297+00:00, including the exact107-coordinate maximum and union charge839. This is a checked partial-sector result awaiting parent review and error-ledger integration; the dense-profile failure remains explicit.

\clearpage

\section{Original-ring covering lemma: primary anchors and exact metric consequences}\label{original-ring-covering-lemma-primary-anchors-and-exact-metric-consequences}

\textbf{Scope.} This addendum accompanies the unchanged proof \nolinkurl{explicit_original_ring_cover_pending.md} (SHA256 \nolinkurl{96325b6cf914de8b82c8c617498749f7d59532185b4df08be3c851cbd2a59ac5}) and its 342-check certificate. It spells out the exact field, discriminants, real-part lattice and target-fibre measure. It makes no covering claim for an arbitrary fractional ideal. It does not itself remove a signed-kernel sector.

\subsection{1. The actual field and primary inputs}\label{the-actual-field-and-primary-inputs}

The field is exactly K032 (QH9)--(QH15): \[
 E=\mathbb Q(\zeta_{28})\,B_7(29)\,\mathbb Q(\sqrt5)\,\mathbb Q(\sqrt{13}),
 \quad [E:\mathbb Q]=336=2f,\quad f=168,\quad F=E^+.
\] The actual product polarization uses a weighted trace form, but the column-count metric under discussion is the original-coordinate Minkowski metric \[
 (x,y)_E=\tfrac12\operatorname{Tr}_{E/\mathbb Q}(x\bar y),\qquad
 (a,b)_F=\operatorname{Tr}_{F/\mathbb Q}(ab).
\] No polarized Gaussian metric or CM coordinate scaling is inserted into this metric.

Fresh primary reads after executor recovery verify the precise inputs. Milne (2020, Theorem 6.4, printed p.~99) gives the full cyclotomic integer ring and power basis. Milne's Theorem 4.9 (printed p.~72) excludes a nontrivial extension of Q unramified at all finite primes. Thus the Galois constituents with pairwise disjoint ramification sets are linearly disjoint, including after each compositum. At a prime where one constituent is ramified the remaining integer algebras are finite étale. Their tensor product is finite étale over the ramified constituent's local maximal order. The Stacks Project, Lemma 15.45.4 (Tag 0AP2), makes that algebra normal, so the tensor integer basis is maximal at every prime. This is a local maximal-order argument, not an assertion that a coprime discriminant merely bounds an unspecified index.

Thaine (2005, Section 2, printed p.~209) explicitly identifies the Gaussian periods as a normal integral basis of the degree-seven subfield of Q(zeta29). The independent orbit-sum argument in the accompanying proof also derives this directly from the full cyclotomic basis. All constituent Gram matrices and their determinants are computed in that proof and checked by the file certificate.

\subsection{2. Exact discriminants and covolumes}\label{exact-discriminants-and-covolumes}

Set \[
 D_0=7^{140}29^{144}5^{84}13^{84}.
\] The tensor determinant calculation, independently consistent with K032 (QH15), gives \[
 |D_E|=2^{336}D_0^2,\qquad D_F=2^{168}D_0=2^fD_0.
 \tag{RM1}
\] For E choose one embedding from each conjugate pair and identify E tensor R with C\^{}f, using the ordinary real coordinates Re and Im. The displayed half-trace inner product is exactly the Euclidean inner product in these coordinates. The change from the full embedding matrix to real-and-imaginary columns has determinant absolute value 2\^{}\{-f\}. Hence \[
 \operatorname{covol}(\mathcal O_E)=2^{-f}\sqrt{|D_E|}=D_0,
 \qquad
 \operatorname{covol}(\mathcal O_F)=\sqrt{D_F}=2^{f/2}D_0^{1/2}.
 \tag{RM2}
\] This is also Milne's Proposition 4.26, with its actual real-coordinate convention. Equivalently, the determinant of the positive half-trace Gram matrix is 2\^{}\{-2f\}\textbar D\_E\textbar=D\_0\^{}2. The factor 1/2 in the inner product occurs once; it is not applied separately at every real tensor factor.

Let zeta=zeta28 and t=zeta+zeta\^{}\{-1\}. The accompanying proof derives \[
 t^6-7t^4+14t^2-7=0,\quad p=t^5-7t^3+14t,\quad tp=7.
\] The integral element a=zeta p-3 satisfies a+bar(a)=1. Multiplication by any element of O\_F therefore proves that Tr\_(E/F)(O\_E)=O\_F. Since every integral trace lies in O\_F, this proves both inclusions in \[
 P_0:=\operatorname{Re}(\mathcal O_E)=\tfrac12\mathcal O_F,
 \qquad \operatorname{covol}(P_0)=2^{-f/2}D_0^{1/2}.
 \tag{RM3}
\] There is no missing dyadic index and no need to infer this equality from relative unramifiedness.

\subsection{3. Exact coarea for the target fibre}\label{exact-coarea-for-the-target-fibre}

The following statement is conditional only on the already separately verified \emph{local arithmetic description} of the target overlattice. Suppose B\_Sigma in E\^{}5 contains O\_E\^{}5, has half-trace covolume D\_0\^{}5 T\^{}\{-10s\}, and its last-real-coordinate projection is \[
 P_\Sigma=P_0 I_F^2,\qquad \operatorname N_F(I_F)=T^{-s}.
 \tag{RM4}
\] These are the all-II hypotheses used in the staged fixed-W0 slice. They must be checked again if a different local sector is substituted.

On E tensor R = F tensor R direct-sum i(F tensor R), the two summands are orthogonal and the restriction of the half-trace metric to each is precisely the F trace metric. The map (x1,\ldots,x5) to Re(x5) is therefore an orthogonal coordinate projection with coarea Jacobian one. A lattice basis for its kernel, followed by lifts of a basis of its image, has a block triangular matrix in this orthogonal decomposition. Taking its determinant proves \[
 \operatorname{covol}(B_\Sigma)
 =\operatorname{covol}(K_\Sigma)\operatorname{covol}(P_\Sigma),
 \quad K_\Sigma=\ker(\operatorname{Re-last}|B_\Sigma).
\] Using (RM2)--(RM4), \[
 \boxed{\operatorname{covol}(K_\Sigma)
 =2^{f/2}D_0^{9/2}T^{-8s}.}
 \tag{RM5}
\] Every nonempty affine fibre is a translate of K\_Sigma and has this same intrinsic covolume. If alpha=2 beta-2 is used in place of beta=Re(x5), the quotient-coordinate Jacobian is 2\^{}f and the image lattice becomes 2P\_Sigma; these two factors cancel in the same coarea identity. They must not be applied twice.

\subsection{4. Why the ring covering bound applies to this kernel}\label{why-the-ring-covering-bound-applies-to-this-kernel}

The inclusion O\_E\^{}5 subset B\_Sigma gives \[
 L_0=\mathcal O_E^4\oplus i\mathcal O_F\subset K_\Sigma.
\] This is a full lattice in the same real vector space of dimension 9f. Multiplication by i is isometric because i lies in E and has modulus one at every embedding. The coordinate summands are orthogonal. The accompanying tensor-basis proof gives \[
 \mu(\mathcal O_E)^2<3248f,\qquad
 \mu(\mathcal O_F)^2<3248f.
\] A larger lattice cannot have a larger covering radius, and the squared covering radii of orthogonal products add. Thus \[
 \mu(K_\Sigma)^2\le\mu(L_0)^2<16240f<128^2f.
 \tag{RM6}
\] No fractional-ideal shape is controlled by this argument: it uses the specific full original-ring sublattice L0 and its literal inclusion in the kernel.

For the actual rank-one indefinite unitary group with operator norm at most H=133 at each real label, a negative unit column has squared Euclidean norm at most (H\^{}2+H\^{}\{-2\})/2. This follows immediately from its compact-Cartan-compact decomposition: the compact factor on the negative line is a phase, and the noncompact block sends that line to (sinh r,cosh r), with e\^{}r at most H. Summing over f labels gives norm less than 95 sqrt(f). Projection to the kernel direction only decreases that norm. The disjoint Dirichlet cells of K\_Sigma, attached to lattice points in any affine-fibre ball, therefore lie in a ball of radius \[
 (95+128)\sqrt f=223\sqrt f.
 \tag{RM7}
\] This justifies the radius replacement in the actual affine-fibre count, with the exact covolume (RM5) whenever (RM4) holds. The separate scalar lattice (1/2)J\_F I\_F\^{}2 and the transverse oriented ideals keep their own fractional-ideal estimates.

\subsection{5. Verification record and references}\label{verification-record-and-references}

The complete accompanying proof and its complete file certificate were author-read; 342 exact arithmetic checks passed at 2026-11-22 15:46:23 UTC. This addendum supplies fresh primary confirmation and spells out the metric deductions. It does not alter the already captured conservative K085 charge or claim new signed-kernel positivity. Root and Kummer independent review statuses remain separately recorded.

Milne, J. S. (2020). \emph{Algebraic number theory} (Version 3.08). Theorems 4.9 and 6.4; Proposition 4.26. \url{https://www.jmilne.org/math/CourseNotes/ANT.pdf}

The Stacks Project Authors. (n.d.). \emph{Permanence of properties under étale maps}, Lemma 15.45.4 (Tag 0AP2). \url{https://stacks.math.columbia.edu/tag/0AP2}

Thaine, F. (2005). Finding families of units of cyclic fields. \emph{RIMS Kokyuroku, 1451}, 207--215. \url{https://www.kurims.kyoto-u.ac.jp/~kyodo/kokyuroku/contents/pdf/1451-21.pdf}

Fresh direct web locators: Milne turn254048555view0, turn721785017view0/1 and turn452687489view0; Thaine turn721785017view2; Stacks turn254048555view2. Source field read: K032 QH9--15 in the restored \nolinkurl{work/agents/degree_four_principal_host_repair.md}. These are independent of the weighted polarized CM metric used elsewhere.

\clearpage

\section{AM-DO-070-C1 --- zero-target pencil projection and the exact weighted correction}\label{am-do-070-c1-zero-target-pencil-projection-and-the-exact-weighted-correction}

Author: \nolinkurl{/root/definite_oscillator}, 22 November 2026 UTC.

This is a substantive, localized correction to the general type-III target-projection statement in the recovered direct-pencil note, DO070 §3, and the inherited W0 slice in DO071. Those original files are preserved as historical bytes; the statements below supersede their assertion that every actual target has double-pole real projection. The correction changes the sufficient all-top charges from 43332 to \textbf{43330} for DO070 and from 839 to \textbf{838} for DO071. It does not change either removed sector. DO072 includes the correction from its final version onward and retains its joint charge \textbf{1482}.

\subsection{1. The actual admitted zero label}\label{the-actual-admitted-zero-label}

K062(AT3, AT5--6), in \nolinkurl{agents/averaged_scalar_pencil_trace.md}, and LIFT44(PR4--7), in \nolinkurl{agents/fixed_scalar_packet_plane_reduction.md}, describe the actual admitted parameter set as \[
D=k_T\setminus\{-1\},\qquad |D|=m_T=T-1.
\tag{ZC1}
\] The two omitted projective points are infinity and −1. In particular, \textbf{zero is admitted}, with \nolinkurl{M_0=Lambda}. It cannot be dropped or relabelled away without changing which target lattice is being counted.

K063(PF12), in \nolinkurl{agents/scalar_pencil_small_field_plane_bound.md}, gives an exact lattice presentation whose pole column and row are supported in the original W0 coordinates. The replacement of the row by its W0-supported version is an integral right change of basis commuting with the shear; it does not change the lattice or any admitted label. K071(NB2--3), in \nolinkurl{agents/pencil_nilpotent_pole_budget.md}, records this original-coordinate decomposition. On W0 write \[
M_c=(I+c\varpi^{-1}Z)\Lambda_0,
\quad Z=zq,\quad Z^2=0,\quad \Lambda_0=O^5.
\tag{ZC2}
\] K077(AP1--3), in \nolinkurl{agents/fixed_block_affine_pencil_count.md}, proves that the negative coordinates \nolinkurl{z_n} and \nolinkurl{q_n} are both units. At the conjugate CM component the pole column and row are the inverse-adjoint dual ones, with these unit roles exchanged. Thus the same projection calculation holds at both CM partners. These are paired lattices, not two independent choices of scalar label.

\subsection{2. Exact full target projection, including target zero}\label{exact-full-target-projection-including-target-zero}

For an actual ordered incidence \nolinkurl{hM_c=M_d}, the elementary containment \nolinkurl{varpi e_n∈M_c} implies \[
h e_n\in\varpi^{-1}M_d.
\] Consequently its original W0 part belongs to \[
B_d=\varpi^{-1}(I+d\varpi^{-1}Z)\Lambda_0.
\tag{ZC3}
\] The inverse matrix of the displayed generator is \nolinkurl{varpi I−dZ}, so \nolinkurl{Lambda_0⊂B_d}. The determinant-one shear gives local determinant valuation −5 for every d.~For a nonzero residue d, the negative diagonal entry of the generator has leading term \texttt{d\ varpi\^{}\{-2\}z\_nq\_n}, a unit times \texttt{varpi\^{}\{-2\}}; its lower pole term cannot cancel it. If d is zero, use its zero lift. It follows that \[
\operatorname{pr}_n(B_d)=
\begin{cases}
\varpi^{-2}O,&d\ne0,\\
\varpi^{-1}O,&d=0.
\end{cases}
\tag{ZC4}
\] The same dichotomy holds at the inverse-adjoint CM partner. In particular, the all-II calculation in K085 remains valid: an active type-II place is nonintegral and fixes a unique admitted lattice, so that fixed lattice cannot be \nolinkurl{M_0=Lambda}. A type-III place has a unique off-diagonal ordered incidence and can have target zero.

Patch the target lattices over s active relative places and use the standard ring lattice elsewhere. Denote the number of zero targets by \nolinkurl{r_0}; all of them are type III. The resulting W0 lattice has \[
\operatorname{covol}(B_\Sigma)=D_0^5T^{-10s}.
\tag{ZC5}
\] Under \nolinkurl{p(x)=Re(x_n)}, its image is \[
p(B_\Sigma)=P_0
\prod_{d\ne0}\mathfrak q^{-2}
\prod_{d=0}\mathfrak q^{-1},
\qquad P_0=\operatorname{Re}(O_E).
\tag{ZC6}
\] Indeed at an active prime, \texttt{E\textbackslash{}otimes\ F\_v=F\_v\textbackslash{}oplus\ F\_v}, both partner negative-coordinate projections in ZC4 are independent, and real part is their sum divided by the unit 2. Away the odd active primes the original image is P0, retaining its dyadic lattice. Localization proves ZC6 globally. Its covolume is \texttt{covol(P0)\ T\^{}\{-2s+r\_0\}}.

\subsection{3. Paired coarea and the exact weighted factor}\label{paired-coarea-and-the-exact-weighted-factor}

The actual half-trace metric makes p an orthogonal projection with normal Jacobian one. A lattice basis of its kernel followed by lifts of an image basis is block triangular in the corresponding real orthogonal decomposition. Therefore \[
\operatorname{covol}(\ker(p|B_\Sigma))
=\frac{D_0^5}{\operatorname{covol}(P_0)}T^{-8s-r_0}.
\tag{ZC7}
\] The conservative estimate \nolinkurl{O_F⊂P0}, hence \nolinkurl{covol(P0)≤sqrt(D_F)}, gives the retained lower bound \[
2^{-f/2}D_0^{9/2}T^{-8s-r_0}.
\tag{ZC8}
\] The newly proved exact \nolinkurl{P0=(1/2)O_F} improves the leading constant by \nolinkurl{2^f}; this optional improvement is not used in the corrected charges. The kernel still contains the orthogonal original-ring sublattice \nolinkurl{O_E^4⊕iO_F} for every target. Consequently any previously justified uniform fibre-radius constant changes only by the factor \texttt{T\^{}\{r\_0\}}. The S-column and gamma-ideal bounds retain their original directions and are unchanged. In particular, the actual unipotent gamma's simple type-III pole does not replace the full target image ZC6.

For a type-III place the exact absolute weight is \nolinkurl{1/[m_T(m_T−1)]}. Among the admitted off-diagonal ordered pairs there are \nolinkurl{(m_T−1)^2} with nonzero target and \nolinkurl{m_T−1} with zero target. After their different fibre costs are included, the full weighted factor is \[
\frac{(m_T-1)^2+T(m_T-1)}{m_T(m_T-1)}
=\frac{m_T-1+T}{m_T}=2.
\tag{ZC9}
\] Type II contributes at most one and has no zero target. All other bounds used here are uniform in the actual source-target pairs. The corrected complete II/III averaging therefore costs exactly the sufficient factor \textbf{\nolinkurl{2^t}}, once for the t type-III relative places. The CM partners do not introduce a second factor. No new G rank pattern, source/target orientation choice, class mass, or frame change is introduced.

\subsection{4. Corrected existing endpoint statements}\label{corrected-existing-endpoint-statements}

For DO070 the profile is \nolinkurl{(i,u,t)=(0,9,2)}. Its old final exponent receives +2: \[
-346657/8+2=-346641/8<-43330.
\tag{ZC10}
\] Thus its correct sufficient charge is \textbf{43330}. Its original arbitrary-ideal transference theorem, monotone y optimization, radius71 and all other arithmetic checks are unaffected.

For DO071 every per-k endpoint receives +1 because t=1. The maximum over \nolinkurl{1≤k≤107} becomes \[
-6769/8+1=-6761/8<-845,
\quad107\,2^{-845}<2^{-838}.
\tag{ZC11}
\] The removed range remains \nolinkurl{1≤k≤107}; its correct sufficient charge is \textbf{838}. The adjacent failed endpoint becomes \nolinkurl{3999/8} at k108, and the dense failed endpoint becomes \nolinkurl{639295/8} at k164. These positive upper estimates remain failures of the sufficient bound, not lower bounds on actual contributions.

DO072 uses the corrected factor in GS23 and in all its endpoint arithmetic. Its t0 cells are unchanged; its t1 cells receive +1. Its worst cell remains t0,k1, equal to \nolinkurl{−1716943/1152}, so the combined 221-cell charge remains \textbf{1482}. The stronger all-II fixed-W0 calculation K085 is unchanged, because its active fixed target is always nonzero.

\subsection{5. Scope and verification protocol}\label{scope-and-verification-protocol}

The original recovered files and their check timestamps remain preserved. This correction supersedes only the identified full-image wording and the two affected numerical charges. The separately saved corrected 19-check DO070 certificate passed at 16:05:57.347562 UTC, with exponent \nolinkurl{−346641/8} and charge43330. Kummer's separately saved corrected 500-check sparse certificate passed at 16:04:59.717615 UTC, with maximum \nolinkurl{−6761/8} on k≤107 and charge838. DO072's corrected complete 1,127-check certificate passed at 16:06:10.440982 UTC; its charge remains1482. Frozen fingerprints and whole-proof verification are recorded in the companion \nolinkurl{gaussian_scalar_cross_verification.json}. No original check is silently relabelled as having included this correction.

All prior signed-error terms must be retained when the corrected charges are integrated. This note changes no original kernel, permitted family, real selected module, period normalization or Mram/DP44/RH conclusion. No complete signed-sum inequality is asserted.

\clearpage

\section{AM-DO-069: positive-rank free data in the genuine realization}\label{am-do-069-positive-rank-free-data-in-the-genuine-realization}

\textbf{Canonical checked version --- 2026-11-22T16:19:19.391528+00:00.} The complete final proof and its source reconciliation were read by root before this integration. Exact diagnostics are separately classified and do not mechanically prove analytic inputs. Historical recovery wording below records the export history; the current scope is: Genuine positive-rank free-data period implication, including both parity branches, exact scalar/measure conventions, actual compact-finite extraction and pointwise literal-lift equality; initial upper period, arithmetic Frames.Good and norm budget remain premises or open.

Author: /root/definite\_oscillator. Reconstructed controlled-recovery export, 22 November 2026.

\subsection{Status and scope of this export}\label{status-and-scope-of-this-export}

This is the complete five-part mathematical argument retained during the executor outage, incorporating the final character, parity, boundary, continuity, actual-compact, and pointwise-equality amendments. LIFT completed its mathematical-text review at 14:42:58 UTC and the parent completed its mathematical-text review at 14:46:53 UTC. Those reviews concerned the text and explicitly carried the source-order interfaces; neither was a saved-file, final-script, or fingerprint verification. This reconstructed file has a new identity. It is not the lost pre-outage final draft and does not inherit its hash or any automatic promotion.

The present argument is restricted to the constructed genuine R014/R015 realization with the accepted analytic inputs R016 and DO068. It is not a theorem for an arbitrary abstract ThetaX, arbitrary D7, or every hypothetical set of predicate data. It proves the free-data period implication for that realization when the named genuine comparison identities hold. It neither constructs arithmetic Frames.Good nor prescribes the new finite Schwartz vector, bounds its norm, or proves an initial global period nonzero.

The restored source definitions have since been read again at their interfaces. They confirm the actual odd P.t fields, the uncontracted Step.lift output of UpBF/UpFB, the pointwise WeilTied.th identity, and the pointwise IsIo.eq identity. These new reads are recorded separately in source\_reconciliation.md; they are not retroactive changes to the outage review events.

\subsection{Part I. Genuine spaces, exact quantifiers, and the parity split}\label{part-i.-genuine-spaces-exact-quantifiers-and-the-parity-split}

Fix even integers k \textgreater= 2 and m \textgreater= 2 in the actual family. Let t be the actual parameter P.t. Assume t is odd and P.t \textgreater= R.r + 5. These are literal fields of TPar, rather than a replacement of P.t by the independently defined Ranks.t. Put

\[
 N_j=k+j(t+1),\qquad n_j=k+m-1+jt,\qquad
 q_j=\lfloor(m-j)/2\rfloor ,
\] \[
 X_j=I_{N_j},\qquad
 Y_j=I_{n_j-q_j}\perp(-I_{q_j}).
 \tag{PD1}
\]

Use the actual standard positive embedding of X\_j in Y\_j, the fixed standard additive character, and the smaller-first oriented Hom model of R014/R015. The negative line L has its actual rational Hermitian form; when the standard model is used, its Gram is -1. A nonstandard line is transported only through an actual rational isometry or the explicit line link below. Equality of real signatures alone is never used as a rational isometry.

Fix the horizontal input Phi-prime and the horizontal automorphic input phi-prime. Suppose the displayed upper horizontal theta lift paired with a genuine upper cusp vector F has nonzero PerNZ. The desired conclusion is a single vertical input Phi that is a restricted pure tensor and finite under the actual maximal compact groups at infinity, such that every cusp vector G equal to the specified vertical lift has nonzero lower PerNZ. This preserves Phi-prime. It does not select a new horizontal tensor and does not assume every upper period is nonzero.

The relevant rational decompositions are determined by m-j.

\begin{itemize}
\item
  If m-j=2rho+2, then \[
  Y_j=\mathbb H^\rho\perp(X_j\perp L),\qquad
  Y_{j+1}=\mathbb H^\rho\perp X_{j+1}.
  \tag{PD2}
  \] This is UpBF, with both Witt ranks rho.
\item
  If m-j=2sigma+3, then \[
  Y_j=\mathbb H^{\sigma+1}\perp X_j,\qquad
  Y_{j+1}=\mathbb H^\sigma\perp(X_{j+1}\perp L).
  \tag{PD3}
  \] This is UpFB, with lower and upper Witt ranks sigma+1 and sigma.
\end{itemize}

For even m these are respectively even and odd j. The formulas in terms of m-j are the fundamental ones. No parity extension to arbitrary m is implicit. The rank checks are exact: \[
 n_j=N_j+2\rho+1,\quad n_{j+1}=N_{j+1}+2\rho
\] in PD2, and \[
 n_j=N_j+2\sigma+2,\quad n_{j+1}=N_{j+1}+2\sigma+1
\] in PD3.

The last two edges are the already accepted R016 cases. The earlier boundary j=m-3 is still an UpFB step with sigma=0. Its upper Bessel period is ordinary restriction, but its lower period is a rank-one Fourier--Jacobi period. It cannot be supplied merely by the ordinary-restriction identification. For m=4 the sequence is UpBF with rho=1, UpFB with sigma=0, the R016 rank-zero BF case, and the R016 compact terminal case.

The literal CayleySignature.ThetaUnfoldingFreeData quantifiers are respected: P, mu, j, phi-prime, Phi-prime, and the upper F are fixed before the new Phi is chosen. The new finite level may depend on this data. There is no claim of uniform membership in an already fixed arithmetic row lattice.

\subsection{Part II. Characters, arbitrary positions, and the full Jacobi identity}\label{part-ii.-characters-arbitrary-positions-and-the-full-jacobi-identity}

Write mu=mu(j+1) and eta=eta\_T. In the constructed genuine character domain, mu has trivial base-field restriction, and eta restricts to the odd quadratic character. The exact ordered lower/upper vertical labels are

\[
(\mu\eta^{-1},\mu)\quad\text{in BF},\qquad
(\mu,\mu\eta^{-1})\quad\text{in FB}.
\tag{PD4}
\]

The horizontal labels are (mu,mu). The positive-space/negative-line auxiliary labels are (mu,eta\^{}\{-1\}). Accordingly the BF product condition is chiV=chiA chiz and the FB product condition is chiW=chiA chiz. These are the literal product hypotheses in UpBF and UpFB.

The source omZ initially uses labels (eta\^{}\{-1\},eta\^{}\{rank A\}) on (L,A). The rank of A is even. The ratio between its second label and mu has trivial base-field restriction. The exact TwRule, evaluated with the line-group argument equal to one, permits changing that second label to mu: the Hilbert--90 determinant representative can be chosen to be one, so the twisting scalar is exactly one. This does not change either the line label or the Heisenberg operator. It does not assert equality of the two oscillator actions on arbitrary line-group elements. No global square root of a determinant character is chosen.

For arbitrary admissible Fourier--Jacobi data, retain the actual rational Witt frame and its coefficients \[
 \beta_i,c\in E^\times .
\] The nonzero hypotheses are essential whenever a ratio is formed. A diagonal change of isotropic frame \[
 e_i'=a_i e_i,\qquad f_i'=\bar a_i^{-1}f_i,\qquad
 a_{i+1}/a_i=\beta_i,\qquad a_{\rm last}=1
 \tag{PD5}
\] normalizes all generic flag coefficients to one. Indeed the actual Hermitian coefficient h(n e\_\{i+1\},f\_i) is multiplied by a\_\{i+1\}/a\_i. The physical unipotent subgroup and its character are unchanged under this re-description. Fixing the last scalar equal to one leaves the last Heisenberg coefficient c and the core unchanged. The analogous Bessel normalization retains its marked vector and its final coefficient. Rational coordinate transport preserves the normalized adelic additive measures by the product formula; it does not justify setting unrelated group volumes to one.

Here is the exact typed WA/WO conversion. Write ell for the one-dimensional Gram of L and G\_A for the Gram of A. Let

\[
 J:WA(L,A)\longrightarrow WO(L,A),\qquad J(w)=w^t\ell ,
\] \[
 D_{\rm h}:WO(A,L)\longrightarrow WO(L,A),\qquad
 D_{\rm h}(w)=G_A^{-1}w^*\ell .
\tag{PD6}
\]

The first map is the inverse of the source toTWsmall map in this rank ordering. For a function q in the WA model set f=J\_*q, meaning f=q composed with J\^{}\{-1\}, and define \[
 \mathcal C(q)(w)=\overline{f(D_{\rm h}w)}.
 \tag{PD7}
\] This formula is a composite map. If C\_D denotes only conjugate dagger, then \(\mathcal C=C_DJ_*\). Suppressing J without transporting the operator would be a type error.

In the adapted Witt block order (e,A,f), the relevant first row is \[
 (1,-a^*G_A,t).
\] The source formulas give \[
 tPt=-ca^t,\qquad
 cenCo=\tfrac12\operatorname{Tr}(\xi_0\,c\bar c\,\ell t),
\] where c bar(c) lies in the base field. With the actual line vector zeta-prime=bar(c), the JacS formulas give \[
 ptS=-\bar c\,a^*G_A,\qquad
 cenS=\tfrac12\operatorname{Tr}(\xi_0\,c\bar c\,\ell\bar t).
 \tag{PD8}
\] Consequently \[
 D_{\rm h}(ptS)=J(tPt),\qquad cenS=-cenCo.
 \tag{PD9}
\] The second equality uses star(xi0)=-xi0 and conjugation invariance of the trace. There is one conjugation of the function value and its unitary central phase, not a second conjugation of the coordinates.

The Hermitian dagger is anti-symplectic for the two oriented Hom spaces. Thus its half-pairing identity, together with PD9, proves the full operator identity at every coordinate: \[
 JacS_{\bar c}(n)\,\mathcal C
       =\mathcal C\,jacOp_c(n).
 \tag{PD10}
\] Equivalently, \[
 JacS_{\bar c}(n)\,C_D
 =C_D\bigl(J_*jacOp_c(n)J_*^{-1}\bigr).
\] The translated arguments agree because D\_h is additive. The central terms and the half-pairing agree because both acquire the same minus under conjugation. Evaluation at zero gives the actual source Fourier--Jacobi theta factor as the conjugate of thetaJS. The generic flag factor remains star(psiFlag) in both formulas. This is precisely why the source fjSrc has star(thetaJS); no generic-character conjugation is silently added or removed.

An arbitrary Bessel position is handled by the literal PerTied.bes\_all/IsTranslSum record. Its period is a finite linear combination of normalized positions paired with finite right translates of F. A nonzero sum contains a nonzero summand. The horizontal Phi-prime is unchanged. The finite translation is transferred to the vertical input in Part III, rather than replacing the original upper vector permanently.

\subsection{Part III. The two unfolding applications, including sigma=0}\label{part-iii.-the-two-unfolding-applications-including-sigma0}

\subsubsection{BF}\label{bf}

Let B.A=X\_j, B.V=Y\_j, B.V0=X\_j perpendicular L, J.W0=X\_\{j+1\}, and J.W=Y\_\{j+1\}. Keep Xi1 equal to the theta-function realization of the prescribed Phi-prime.

The normalized upper FJ datum determines zeta-prime=bar(c). Choose the Bessel marked line vector and line link \[
 \zeta=\xi_0\zeta',\qquad \tau=-\xi_0\,id_L.
 \tag{PD11}
\] Then tau zeta-prime=-zeta and \[
 \langle\zeta,\zeta\rangle
  =\xi_0\bar\xi_0\langle\zeta',\zeta'\rangle,\qquad
 \tau^*G_L\tau=(\xi_0\bar\xi_0)G_L.
\] These are the exact Lnk requirements.

The literal ofAdj formula is \[
 ofAdj_{L,A}(z)=-\xi_0^{-1}G_A^{-1}z^*G_L.
\] Substituting tau gives \[
 \kappa=-D_{\rm h}.
 \tag{PD12}
\] The line comparison statement therefore sends a theta function Xi2 on WO(L,A) to \(w\mapsto\overline{\Xi_2(\kappa w)}\), with the required exchanged character labels at line-group identity. This is an exact rational conjugate-adjoint transport.

The normalized FJ auxiliary vector from Part II is \(s_0(w)=\overline{f(D_{\rm h}w)}\). Define \[
 \Xi_2(z)=f(-z).
 \tag{PD13}
\] Then \(C_\kappa\Xi_2=s_0\) exactly. PD13 is a rational coordinate-parity pullback. It is not identified with a separately normalized Weil operator at the group element -1, which could carry an additional chosen splitting scalar.

Choose a finite compact K\_G fixing the smooth finite vector F. Apply DO068 to the actual orthogonal core tensor B.o.boxS(Xi1,Xi2). It supplies a common Gamma0, a LocS space, and one compact open K contained in K\_G. The finite section and its support are constructed for that same K. The prescribed core is used as an actual theta function; no inverse realization of an unspecified abstract oscillator vector is assumed.

All character and line-link hypotheses in UpBF now hold, and its upper nonzero period is exactly the normalized upper period. The conclusion is \[
 \exists\Phi\in\mathcal S,\qquad
 B.per(\phi',s.lift(F,\Phi))\ne0.
 \tag{PD14}
\] The source Step.lift is the ordinary uncontracted liftS integral. No archimedean algebraic contraction has yet been extended to all Schwartz functions.

\subsubsection{FB}\label{fb}

Let J.W0=X\_j, J.W=Y\_j, B.A=X\_\{j+1\}, B.V=Y\_\{j+1\}, and B.V0=X\_\{j+1\} perpendicular L. Keep Xi1=Phi-prime. Choose a nonzero genuine auxiliary theta function Phi0 on WO(X\_j,L) with labels (mu,eta\^{}\{-1\}). Apply DO068 to the actual boxT(Xi1,Phi0). This gives the required common Gamma0 and LocT space with K fixing F.

If the initial upper Bessel position was arbitrary, select a nonzero normalized summand as described above. For its finite translate R(h)F, the exact change of variable in the theta integral is \[
 D_{\rm full}(\Phi,R(h)F)
       =D_{\rm full}(\omega(1,h^{-1})\Phi,F).
 \tag{PD15}
\] The Haar measure is right invariant. The finite Weil action preserves restricted tensor purity and leaves the archimedean compact-finiteness condition unaffected. Thus the final conclusion uses the original F.

The normalized upper Bessel hypothesis is the exact premise of UpFB. Its conclusion is a nonzero lower JacL.per with the actual marked line parameter appearing in that conclusion. Denote it by zeta\_FB. Inverting the typed comparison of Part II uses \[
 c=\overline{\zeta_{\rm FB}}
\] and the normalized generic coefficients. Do not import the BF factor xi0 into this independently supplied FB conclusion.

When sigma=0, the upper isotropic flag is empty. Its NFlag is the identity subgroup: the defining flag condition forces (n-1)x=0 on every basis vector. The normalized quotient mass is one and the empty generic character is trivial. An unused marked-vector coordinate has no effect on the ordinary upper restriction; a nonzero one may be chosen. Nevertheless UpFB with sigma=0 is still required to create the lower FJ period of rank one. This is the earlier j=m-3 boundary, not the final R016 compact edge.

\subsubsection{Common Ann and the exact role of Loc}\label{common-ann-and-the-exact-role-of-loc}

For Gamma0 with zero infinite component, the source Ann is the intersection of kernels of additive characters: both its doubled-lattice clauses and its split-idempotent half-lattice clauses are additive in the translating vector. It is therefore an additive closed subgroup. Its finite part lies in the compact annihilator of twice the finite lattice. The Heisenberg cocycle has finite image on this compact set. Adjoin those central values to obtain a compact group. A smooth vector has an open stabilizer in it, so its orbit spans a finite-dimensional space; the cocycle law makes that span stable under Ann.

The common Gamma0 is retained through this orbit span. Real Heisenberg translations are handled in the full archimedean Schwartz space, as required by Loc.real. Almost all finite places use the already proved basic-vacuum support in DO068; the construction does not impose infinitely many new primitive bumps.

Loc is an auxiliary existence device. It does not identify the free K with a prescribed arithmetic level. Frames.Good additionally requires the actual latRow vacuum, latCore/GammaL tensor inclusions, and LastLat. Those conditions remain outside this theorem.

\subsection{Part IV. Uncontracted continuity, pure actual-K-finite extraction, and the literal lift}\label{part-iv.-uncontracted-continuity-pure-actual-k-finite-extraction-and-the-literal-lift}

The order of these steps is essential.

\begin{enumerate}
\def\labelenumi{\arabic{enumi}.}
\tightlist
\item
  UpBF or UpFB supplies a nonzero scalar period of the genuine uncontracted theta lift, on the full Schwartz realization.
\item
  Prove continuity of that fixed scalar functional in the vertical archimedean Schwartz input.
\item
  Extract a pure tensor finite under the actual compact groups.
\item
  Only then apply the fixed genuine model identity on the actual polynomial/Harish--Chandra input modules.
\end{enumerate}

Here the upper F lies in the actual smooth moderate-growth globalization of its cuspidal automorphic representation. It is not merely an arbitrary smooth function whose constant terms vanish. Fix F, phi-prime, the horizontal tensor, the auxiliary theta datum, and the finite Schwartz support.

The lower positive unitary adelic quotient is compact. The rational unipotent adelic quotient is compact as well. Choose compact lifts on finitely many charts. The finite level yields finitely many arithmetic components for the upper integral. Weighted lattice summation in the genuine Schrödinger realization gives, on each component, \[
 |\theta_\Phi(r(n,h),g)|
 \le C\,p(\Phi_\infty)\,\|g\|^M,
 \tag{PD16}
\] uniformly on the compact lower charts, where p is one fixed Schwartz seminorm. Compact oscillator factors act continuously with uniform seminorm bounds; diagonal factors have explicit polynomial dilation bounds. The fixed finite data produce finitely many lattices. These facts give PD16 without a claim of uniformity in all automorphic inputs.

The real centre of the unitary group is compact. Miller and Schmid's Theorem 2.24, under their standing hypothesis (2.1), gives rapid decay of the fixed smooth cusp vector on upper Siegel sets, controlled by a smooth seminorm. Choose the decay order larger than M and the Haar growth exponent. This gives an integrable majorant bounded by C-prime p(Phi-infinity). Hence the fixed full period is a continuous linear or conjugate-linear functional, according to its designated slot; all Fubini changes used here are absolute. This is a fixed-period statement, not a blanket joint theta-period continuity theorem.

Choose a compatible oscillator maximal compact containing the actual source and target maximal compact groups. The genuine compatible Cartan data provide this choice. Finite Hermite sums in the corresponding polarization are Schwartz dense. Each lies in finitely many types for both actual compact groups. A nonzero continuous functional cannot vanish on that dense algebraic subspace. After decomposing a finite sum into restricted pure tensors, one pure actual-K-finite tensor has nonzero period.

The final tensor need not remain in the original Loc subspace. Loc was used to obtain a nonzero functional, while the desired free-data predicate asks for fact and kfin. Nothing in the argument claims preservation of Loc under the finite-K projection.

The literal WeilTied.th record identifies the genuine uncontracted theta integral with X.th, multiplied by the nonzero source period scalar. If the coefficient/Harish--Chandra realization is used, the fixed IsIo.eq record gives the corresponding identity for every allowed polynomial and real input vector. It is invoked only at the extracted K-finite input. The exact nonzero normalization scalar is retained. The argument does not infer a Hilbert-bounded or Schwartz-continuous extension from a merely algebraic nonzero IsItW.

The restored IsIo.eq and WeilTied.th records are pointwise. More generally, if an intermediate realization supplies only an almost-everywhere equality, differentiating under the same majorant with enlarged Schwartz seminorms gives smooth output. Equality of smooth functions almost everywhere for full-support Haar measure is pointwise, before restriction to a lower-dimensional period subgroup.

Finally, the universal quantifier over cusp G equal to the fixed lift follows from that equality. If one separately wants to assert existence of a cusp output, retain the accepted reverse-tower and selected-real-parameter hypotheses. The forward lower-tower statement LIFT56 is a different result and is not substituted for this step.

\subsection{Part V. Measures, exact scalars, final conclusion, and remaining boundaries}\label{part-v.-measures-exact-scalars-final-conclusion-and-remaining-boundaries}

Write \[
 c_V=petF(V)\prod_\tau petR(V_\tau)>0
\] for the actual source period normalization. The upper horizontal comparison retains the factors \(c_{X_j}c_{X_{j+1}}\); the lower comparison retains \(c_{X_j}c_{Y_{j+1}}\). They are nonzero and may be divided out in the nonvanishing implication, but are not set equal to one. The bB conjugation makes the first automorphic slot bilinear in the input phi-prime, matching the normalized TowerEx periods.

Use the literal Step measures: \[
 vW:\mathrm{Vol}(W),\quad
 vS:\mathrm{VolM}(W.rank,\rho),\quad
 vT:\mathrm{VolM}(\sigma,V0.rank),\quad
 vU:\mathrm{VolU}(\rho).
\] In particular the domain ordering of vT is the source mixT ordering. It is not replaced by the dimensions suggested by an informal y notation. Rational coordinate changes are accompanied by the actual additive measures; their global scalar is one by the product formula. The group measures and c\_V factors remain explicit. The only quotient mass set equal to one by an empty-rank argument is the trivial flag quotient in the sigma=0 boundary.

Thus, within the named genuine realization and the exact comparison interfaces, the positive-rank BF and FB cases supply the free-data implication with the same fixed horizontal Phi-prime. Together with the already accepted R016 last-two cases, this is the desired analytic mapping of CayleySignature predicate 3 for the constructed family.

The conclusion does not establish an arbitrary abstract Hyp3 constructor, a literal Host/SizeData family, arithmetic Frames.Good, a fixed finite row level, a bounded new tensor, a prescribed global period value, or any final divisibility/DP44/RH conclusion. It also does not infer a global nonzero period from local occurrence: the upper nonzero period is a premise.

\subsubsection{References and exact interface locators}\label{references-and-exact-interface-locators}

Bao, Y., \& Sun, B. Coincidence of algebraic and smooth theta correspondences. arXiv:1611.06298. Theorems 1.1 and 1.3 and the compatible-compact definitions were read as primary statements. The occurrence theorem is not used as an extension theorem for arbitrary algebraic IsItW.

Miller, S. D., \& Schmid, W. (2012). On the rapid decay of cuspidal automorphic forms. Advances in Mathematics, 231(2), 940--964. \url{https://doi.org/10.1016/j.aim.2012.06.001}. The directly read theorem is Theorem 2.24 of arXiv:1106.2149v1, PDF page 12, with standing hypothesis (2.1).

Pinned source commit: 7cbf7a0c32fca1badf05931879a37c3f8a2b03b4.

\begin{itemize}
\tightlist
\item
  \nolinkurl{Formalism/SharpRestriction/CuspSpanNormBounds/CayleySignature.lean:} car, vth, bB, PerNZ, ThetaUnfoldingFreeData.
\item
  \nolinkurl{Formalism/LocalRingCardinality/DataTheta.lean:} TPar.t, t\_odd, t\_ge.
\item
  \nolinkurl{Formalism/LocalToGlobal/BesselHermitianTheta/TiedWeil.lean:} TwRule, om\_swap, WeilTied.th.
\item
  \nolinkurl{Development/AbelianVarietyTheory/Definitions/SemilocalFrames/IsBesPosOfBesselTied.lean:} tPt, cenCo, jacOp, FJDat.chr/per, IsFJPos, FJTied, PerTied.
\item
  \nolinkurl{Development/AbelianVarietyTheory/Definitions/TwistedInnerProducts/WittSmallGroupEmbeddingMemPointwiseStabilizerIso.lean:} ptS, cenS, jacS, thetaJS, psiFlag, fjSrc.
\item
  \nolinkurl{Development/AbelianVarietyTheory/Definitions/LatticeFrames/IsThetaFam.lean:} VolM, Step, Step.lift.
\item
  \nolinkurl{Development/AbelianVarietyTheory/Definitions/TwistedInnerProducts/TowerEx.lean:} Lnk, Loc, UpBF, UpFB, Frames.Good.
\item
  \nolinkurl{Formalism/LocalToGlobal/IndefiniteReal/LocalNewformFunction.lean:} toTWsmall (lines 80--85 in the retained version), oriented Heisenberg definitions.
\item
  \nolinkurl{Formalism/LocalToGlobal/IndefiniteReal/OrthogonalComplementCharacterisation.lean:} IsIo.integ and pointwise IsIo.eq.
\item
  \nolinkurl{Formalism/LocalToGlobal/IndefiniteReal/LinksChar.lean:} IsItW and the actual polynomial input types.
\end{itemize}

Only statements and definition interfaces are used here. No Lean build or re-audit of stipulated theorem proofs is asserted.

\clearpage

\section{AM-LIFT-061 --- Actual positive-rank real Wu constructor}\label{am-lift-061-actual-positive-rank-real-wu-constructor}

\textbf{Canonical checked version --- 2026-11-22T16:19:19.391528+00:00.} The complete final proof and its source reconciliation were read by root before this integration. Exact diagnostics are separately classified and do not mechanically prove analytic inputs. Historical recovery wording below records the export history; the current scope is: Actual real Wu construction/reproduction/decay for the stated genuine positive-rank data; seed energy E\textgreater0 and its inverse square-root cost are retained unevaluated.

\textbf{Status.} Fresh reconstruction after the executor outage, with the recovered source interfaces explicitly reconciled. This file does not claim to reproduce missing historical bytes. Its final whole-file review and diagnostic fingerprints are recorded separately. The proof is in the stated genuine-data scope; no Lean build or mechanical verification of the analytic theorems is claimed.

\subsection{1. Statement, actual parameters, and retained scope}\label{statement-actual-parameters-and-retained-scope}

Let k,m be even, m\textgreater=2, k\textgreater=m+3, and put \[
 \rho=m/2-1,\quad r=k+m-1,\quad p=k+\rho,\quad q=\rho+1=m/2.
 \tag{RW1}
\] At every real label, use the actual bottom Hermitian space of signature (p,q), the actual compact anisotropic subgroup H=U(k), and its Bessel subgroup S=N semidirect H. The marked line has the literal norm -T. All root coefficients, including the literal factor 1/G in the additive character, are retained.

The decreasing signed Harish--Chandra list of the bottom discrete series sigma is \[
 v_0^+,\ldots,v_{k-2}^+,(-b_0)^+,(-b_1)^-,\ldots,
 (-b_{m-1})^{(-1)^{m-1}}.
 \tag{RW2}
\] The compact representation rho\_u has Harish--Chandra list u. Assume the actual integral/parity, compact dominance, interlacing, centredness, and NCGaps clauses: \[
 u_i>v_i>u_{i+1},\quad u_{k-1}>-b_0,\qquad
 b_0<b_1<\cdots<b_{m-1}.
 \tag{RW3}
\] The noncompact gaps exceed (r+mt)\^{}3, where throughout this proof t=P.t is the actual TPar parameter. It is not silently identified with Ranks.t. The literal TPar fields give t odd and t\textgreater=r+5.

For any prescribed unit minimal vector w in the actual local discrete-series model, and for its actual Haar/formal degree, define \[
 \phi(g)=d_\sigma\overline{\langle w,\sigma(g)w\rangle}.
 \tag{RW4}
\] Inner products are conjugate-linear in the first variable.

We construct a nonzero continuous function W:G-\textgreater V\_\{rho\_u\}, satisfying the actual left Bessel covariance and right coefficient reproduction, with unit quotient L2 norm, quotient L1 norm, and the precise polynomially weighted L1 norm required by WuDecay. A compactly supported bounded reproducer is constructed for the same W. Tensoring the local constructions with the actual source tensor isometries supplies Wu\_spec, Wloc\_cont, Wu\_tens, WuReprod, and WuWeighted.

This is an actual local/real constructor. It does not prove a nonzero global prescribed Bessel or Mram period. Its normalization retains an explicit positive but unevaluated seed energy; it does not supply a uniform numerical inverse cost. The finite/geometric Hida fields are supplied separately by checked LIFT60. The row-evaluation reconstruction is supplied separately by LIFT62; its concrete arithmetic record pin remains separate.

\subsection{2. The actual Bessel model occurs}\label{the-actual-bessel-model-occurs}

\subsubsection{2.1. Fix the correct relevant pair}\label{fix-the-correct-relevant-pair}

The desired assertion is \[
 \operatorname{Hom}_{S}
 \bigl(\sigma^\infty\widehat\otimes\rho_u^\vee,\psi\bigr)\ne0.
 \tag{RW5}
\] All smooth representations here use their smooth Fréchet moderate-growth globalizations. The theorem cited below does not supply an arbitrary algebraic Harish--Chandra functional.

For m=2 the Bessel subgroup is simply H and the statement is the checked local branching/Pieri assertion of LIFT28, HJ4--9. The anchor is a local representation with the same retained u,v,b0,b1. It need not be supplied by a separate global m=2 DataParams record. Its lowest compact type and the selected compact type satisfy the exact interlacing criterion. In particular this anchor determines the required even compact packet label without guessing an absolute Whittaker normalization.

Xue's relevant-pair convention uses a line of norm (-1)\^{}k=+1. Simultaneously negate both Hermitian forms. Their unitary matrix groups and representations do not change; H becomes negative definite, the marked line becomes positive, and the hyperbolic planes remain hyperbolic. Thus the actual pair is in the convention of Xue's theorem.

The nonzero root coefficients in the actual Bessel character are matched by an archimedean diagonal change of Witt basis: \[
 e_a\mapsto c_a e_a,\qquad f_a\mapsto\bar c_a^{-1}f_a,
 \tag{RW6}
\] together with the required reversal of the flag order. The successive nonzero c\_a are solved recursively. This also matches the literal marked-line coefficient and 1/G. It is a local conjugacy of the real Bessel datum, not a rational adelic move or a change of the finite level. Pull the resulting functional back to the original S and original psi.

\subsubsection{2.2. Relative packet recurrence}\label{relative-packet-recurrence}

The ambient rank r is odd. In Ichino's zero-based notation for its decreasing HC list, the component-group label is \[
 \eta_i=\epsilon_i(-1)^i.
 \tag{RW7}
\] Increasing m by2 preserves the parity class of every retained HC coordinate and leaves its position unchanged.

For j\textgreater=1, append -b\_\{2j\},-b\_\{2j+1\}. These are strictly below every previously retained ambient coordinate and every u\_i. Before simultaneous form negation their signs are (+,-); afterwards they are (-,+). Their zero-based positions are k-1+2j and k+2j, respectively odd and even. Hence both new labels in RW7 equal +1, while every old ambient label remains unchanged.

The new products with compact dual parameters have odd angular exponent \[
 2(-b_{2j}-u_i)<0
 \quad\hbox{or}\quad
 2(-b_{2j+1}-u_i)<0.
 \tag{RW8}
\] Under the defining additive-character normalization in Xue, the relevant complex epsilon factor is constant in \{+1,-1\} on this strict half-line. More precisely, GGP Examples, Proposition2.1, writes alpha=(bar z/z)\^{}a and psi\_0(z)=exp(2pi i Tr(iz)); epsilon is +1 for a\textgreater0 and -1 for a\textless0. Thus an odd negative angular exponent m corresponds to a=-m/2\textgreater0 and contributes +1 under that normalization. Reversing the fixed additive character changes the common sign, which has no effect below.

Each new ambient component-group equation is a product of k equal signs, hence equals +1 because k is even. Each old compact equation gains two equal signs and is unchanged. The old ambient equations are unchanged. Therefore the true m=2 packet condition remains true after every pair of additions.

Xue's Theorem1.1 now proves RW5, with multiplicity one. This argument uses its exact relevant-pair, additive-character, and smooth-globalization conventions, and Ichino's exact zero-based label formula. It does not infer the result merely from the existence of a parameter of the correct dimension.

\subsection{3. Weighted integrability of the fixed minimal coefficient}\label{weighted-integrability-of-the-fixed-minimal-coefficient}

\subsubsection{3.0. The literal source low vector is covered}\label{the-literal-source-low-vector-is-covered}

The recovered \nolinkurl{RelevantZeroGramTerms.lean}, lines 72--91, has three distinct declarations. \nolinkurl{IsDS} supplies a genuine irreducible unitary square-integrable representation, an occurring K-type of highest weight \nolinkurl{lowK = lam + rhoN − rhoC}, and minimal Vogan norm. \nolinkurl{lowKSub} is the supremum of all copies of that K-type. \nolinkurl{IsLowUnit} means norm one and membership in this subspace; it imposes no highest-weight condition. The stronger \nolinkurl{IsHighestUnit} is a different declaration.

For the genuine discrete series used here, the minimal-K theorem cited below identifies this highest weight with the unique minimal K-type, occurring with multiplicity one. Consequently \nolinkurl{lowKSub} is exactly that single finite-dimensional subspace kappa. Every literal \nolinkurl{IsLowUnit} vector is therefore a unit vector in kappa. It is smooth and K-finite, and the extraction below applies to the whole kappa. No replacement of the supplied vector by an extremal vector is made. \nolinkurl{AutomorphicEvaluation.IsLowVec} and \nolinkurl{HidaSetting.low} use precisely \nolinkurl{IsLowUnit}, so this is the required source identification.

\subsubsection{3.1. From the Flensted--Jensen function to the actual prescribed w}\label{from-the-flenstedjensen-function-to-the-actual-prescribed-w}

Checked LIFT59, BC14--21, supplies the group-manifold Flensted--Jensen function Psi for this discrete series, with Psi(1)=1 and \[
 |\Psi(g)|\le
 \cosh(R(g)/2)^{-2\delta},\qquad
 \delta=\min_{\epsilon_i\ne\epsilon_{i+1}}(h_i-h_{i+1}+1),
 \tag{RW9}
\] where R(g)=log\textbar\textbar g\textbar\textbar{} in the actual Cartan-normalized matrix norm.

For precision about the actual coefficient, Kassel--Kobayashi §7.1 identifies the boundary vector as diagonal-K invariant and as generating the highest module which, under Wick rotation, is the full K times K minimal type. Their §10.2 identifies the group-manifold module with sigma\_K tensor sigma\_K\^{}vee. The minimal K-type kappa of sigma has multiplicity one. Consequently the normalized diagonal-invariant vector is, up to harmless inversion/conjugation orientation, \[
 \Psi(g)=\frac1{d_\kappa}
 \operatorname{Tr}\bigl(P_\kappa\sigma(g)|_\kappa\bigr),
 \qquad d_\kappa=\dim\kappa.
 \tag{RW10}
\] This is an identification of the full minimal K times K type, not merely an assertion about the restriction of Psi to K.

For the fixed unit w in kappa put m\_w(g)=\textless w,sigma(g)w\textgreater. Compact Schur orthogonality gives \[
 m_w(g)=d_\kappa^2
 \int_K\overline{m_w(k)}\Psi(kg)\,dk,\qquad
 \int_K|m_w(k)|\,dk\le d_\kappa^{-1/2}.
 \tag{RW11}
\] The Cartan norm is left K-invariant. Hence \[
 |\phi(g)|
 \le d_\sigma d_\kappa^{3/2}
 \cosh(R(g)/2)^{-2\delta}.
 \tag{RW12}
\] This keeps the actual coefficient, actual formal degree, and actual minimal-vector normalization. No arbitrary favourable Flensted--Jensen vector replaces w.

\subsubsection{3.2. The literal weighted exponent fits the actual NCGaps}\label{the-literal-weighted-exponent-fits-the-actual-ncgaps}

For U(p,q), with q\textless=p, the Cartan Jacobian is bounded by a fixed constant times exp(2pqR), up to a polynomial chamber factor. Also \textbar\textbar g\^{}\{-1\}\textbar\textbar=\textbar\textbar g\textbar\textbar{} in the chosen unitary matrix norm. Thus \[
 \int_G|\phi(g)|\,\|g^{-1}\|^A\,dg<\infty
 \quad\text{whenever }A+2pq<\delta.
 \tag{RW13}
\] The source exponent is \[
 A=N1=4r^2\sqrt{2q_{\max}}.
 \tag{RW14}
\] Since 2q\_max\textless=r, one has A\textless=4r\^{}3 and 2pq\textless=r\^{}2/2. The actual t=P.t satisfies t\textgreater=r+5 and m\textgreater=2. Therefore \[
 \delta>(r+mt)^3\ge(3r+10)^3>4r^3+r^2/2.
 \tag{RW15}
\] The last strict inequality is certified by \[
 2(3r+10)^3-8r^3-r^2
 =46r^3+539r^2+1800r+2000>0.
 \tag{RW16}
\] Thus RW13 holds for the precise source exponent. In particular phi is in L1. This deduction needs only the literal noncompact gaps and TPar lower bound, not a silently stronger 16r\^{}19 hypothesis or a full SizesAt record.

\subsection{4. The actual induced Hilbert space and orthogonal projection}\label{the-actual-induced-hilbert-space-and-orthogonal-projection}

The subgroup S=N semidirect H is unimodular: N is unipotent, H is compact, and its conjugation preserves Haar on N. Use the actual source Haar on S and the corresponding quotient measure from Weil integration. Let \[
 \tau(s)=\psi(s)\rho_u(s),\qquad
 \mathcal H=L^2(S\backslash G,V_{\rho_u};\tau).
 \tag{RW17}
\] The convention is F(sx)=tau(s)F(x), and right translation is R(g)F(x)=F(xg).

Formal-degree Schur orthogonality gives phi\emph{=phi and phi}phi=phi. Consequently \[
 P=R(\phi)
 \tag{RW18}
\] is a bounded orthogonal projection on this exact induced Hilbert space.

It is essential to prove P is nonzero; a formal Frobenius-reciprocity slogan is not enough. We do this using the actual relative matrix-coefficient form.

\subsection{5. Nonzero kernel from the actual relative form}\label{nonzero-kernel-from-the-actual-relative-form}

For compactly supported continuous sections modulo S, Weil integration gives the integral kernel \[
 K_P(x,y)=\int_S\phi(x^{-1}sy)\tau(s)\,ds.
 \tag{RW19}
\] The translated discrete-series coefficient is Harish--Chandra Schwartz. The relative Schwartz restriction estimate in Beuzart-Plessis, §7.1, makes this integral absolutely convergent and continuous, uniformly for x,y in compact sets.

Let Pi=sigma tensor rho\_u\^{}vee. For u,v in V\_rho set \(e=\sigma(y)w\otimes\bar v\), \(e'=\sigma(x)w\otimes\bar u\). With the conjugate-linear-first convention, a direct calculation gives \[
 \langle u,K_P(x,y)v\rangle
 =d_\sigma\int_S
   \langle e,\Pi(s^{-1})e'\rangle\psi(s)\,ds
 =d_\sigma\int_S
   \langle e,\Pi(t)e'\rangle\bar\psi(t)\,dt.
 \tag{RW20}
\] BP uses an inner product linear in the first variable (§2.2, printed42). Under that convention, RW20 is exactly d\_sigma times the complex conjugate of BP's L\_Pi(e,e'), with xi=psi. Thus its nonvanishing is precisely the relative form attached to RW5. There is no additional coordinate conjugation.

Beuzart-Plessis, Theorem7.2.1, says that this relative form is nonzero precisely when the Bessel Hom space is nonzero. For these Schwartz coefficients the regularized notation in the general theorem is the ordinary absolutely convergent integral. RW5 therefore makes the form nonzero.

The G-orbit span of w is dense in the smooth globalization: the closure contains the Harish--Chandra module generated by its minimal K-type, and K-finite vectors are smooth-dense. The H factor is finite dimensional. Continuity of the relative form then shows that K\_P is not identically zero.

Its covariance and adjoint identities are \[
 K_P(sx,y)=\tau(s)K_P(x,y),\quad
 K_P(x,sy)=K_P(x,y)\tau(s)^{-1},\quad
 K_P(y,x)^*=K_P(x,y).
 \tag{RW21}
\]

Choose x0,y0,u such that K\_P(y0,x0)u is nonzero. Let chi be a nonnegative compactly supported continuous function on S\G, positive near Sy0, and define \[
 F(y)=\chi(Sy)K_P(y,x_0)u.
 \tag{RW22}
\] This is a bounded continuous equivariant section, compactly supported modulo S. Moreover, \[
 \langle u,(PF)(x_0)\rangle
 =\int_{S\backslash G}
   \chi(Sy)\|K_P(y,x_0)u\|^2\,dy>0.
 \tag{RW23}
\] Thus P is nonzero on an explicit class of compact quotient seeds. The points x0,y0 need not be the identity. In particular no assertion that the required H-type occurs already in the restriction of the minimal K-type is used.

\subsection{6. Actual Wu function, decay, and its retained energy}\label{actual-wu-function-decay-and-its-retained-energy}

Put W0=PF. It is bounded and continuous, because phi is L1 and F is bounded continuous. Absolute L1 Fubini gives pointwise \[
 R(\phi)W_0=W_0.
 \tag{RW24}
\] Let \[
 \nu(x)=\inf_{s\in S}\|sx\|\ge1.
 \tag{RW25}
\] It is left S-invariant and measurable (the infimum may be taken over a countable dense subset of S). Submultiplicativity gives \[
 \nu(zg^{-1})\le \nu(z)\|g^{-1}\|.
\] Right translation preserves the quotient measure. Therefore weighted Young's inequality gives \[
 \|W_0\nu^A\|_1
 \le \left(\int_G|\phi(g)|\|g^{-1}\|^A\,dg\right)
      \|F\nu^A\|_1<\infty.
 \tag{RW26}
\] The compact quotient support of F makes its weighted L1 norm finite: a compact quotient set has a compact family of local representatives. Taking A=0 also gives the required quotient L1 condition.

The exact positive energy is \[
 E=\|PF\|_2^2=\langle F,PF\rangle>0.
 \tag{RW27}
\] Set W=E\^{}\{-1/2\}W0. Then W has unit quotient L2 norm and the actual left covariance, continuity, ordinary and weighted L1 conditions, and coefficient reproduction.

The precise cost left by this construction is \[
 \|W\|_\infty\le\frac{\|\phi\|_1\|F\|_\infty}{\sqrt E},
 \qquad
 \|W\nu^A\|_1\le
 \frac{\|\phi\|_{L^1_A}\|F\nu^A\|_1}{\sqrt E}.
 \tag{RW28}
\] No effective lower bound for E has been proved. These formulas must not be described as a uniform numerical bound for the source detector or its inverse.

\subsection{7. A compact reproducer for the same fixed minimal vector}\label{a-compact-reproducer-for-the-same-fixed-minimal-vector}

Choose a nonnegative compactly supported smooth probability bump b sufficiently close to the identity. Put a0=b\^{}**b and average it under compact K-conjugacy. The resulting a is compactly supported, nonnegative, and has L1 norm1. Its operator sigma(a) is positive and commutes with K. Minimal-K-type multiplicity one therefore gives sigma(a)\textbar\_kappa=lambda I. Taking b sufficiently close to the identity gives lambda\textgreater0. Put alpha=a/lambda.

This construction fixes the actual w, not a different favourable minimal vector. Direct calculation yields \[
 (\alpha*\phi)(g)
 =d_\sigma\langle\sigma(g)w,\sigma(\alpha)w\rangle
 =\phi(g).
 \tag{RW29}
\] Consequently R(alpha)P=P. Applying this to the bounded W above gives pointwise absolutely convergent compact reproduction. Its compact support may be enlarged to a compact Omega containing the identity, as required by WuReprod.

An optional explicit bump bound is available. Let L\textgreater=0 be the scalar of the positive elliptic Nelson Laplacian on kappa; in Cartan-orthonormal coordinates this is \(2c_\kappa-c_\sigma\), with the compact central contribution included. Take \[
 \epsilon=\frac1{4\sqrt{1+L}},\qquad
 \operatorname{supp}b\subset\exp B(\epsilon/2).
\] For every unit v in kappa, \(\|\sigma(g)v-v\|\le1/8\) on this support. Hence lambda\textgreater=49/64 and \[
 \|\alpha\|_1\le64/49.
 \tag{RW30}
\] The support has operator norm at most exp(epsilon). This controls the compact reproducer but does not remove the unevaluated energy E in RW28.

\subsection{8. Tensor assembly and literal Hida fields}\label{tensor-assembly-and-literal-hida-fields}

Carry out the construction at each actual real label with its actual local quotient Haar. The source tensor isometries identify the products of the compact representation spaces and the discrete-series models. Choose the global unit minimal vector to be the source tensor of its prescribed local unit minimal vectors.

The product quotient norm equals one by Fubini. The product of the local reproducing identities gives the actual global coefficient reproduction. The source normInf is a product of local operator norms, so its infimum over the product subgroup factors, and the local estimates with the common A=N1 give the global weighted estimate. The finite product of compact reproducers gives the global compact reproducer and compact support.

Thus the actual objects satisfy:

\begin{itemize}
\tightlist
\item
  Wu\_spec and Wloc\_cont, with the literal left Bessel covariance and source quotient measure;
\item
  Wu\_tens, with the fixed tensor realization;
\item
  WuReprod, using RW29;
\item
  WuWeighted, using RW26 and the exact source N1.
\end{itemize}

Together with LIFT60, this supplies the geometric/finite/real analytic HidaSetting data in the stated standard-host scope. The right regular unitary operators and actual Haar transport are the canonical ones. The final row-evaluation and conjugation identifications are handled in LIFT62; the independently presented k/sEE pin is not silently identified.

\subsection{9. Exact source anchors}\label{exact-source-anchors}

The following exact definitions were read in the restored pinned source after recovery, in addition to the primary checks listed below. Source formulas are retained as hypotheses where they are record fields; no already instantiated arithmetic package is inferred.

\begin{itemize}
\tightlist
\item
  \nolinkurl{Formalism/LocalToGlobal/SiegelLatticeTheory/HaarTransversals/HidaSetting.lean:} actual psiR at43--48; S/N/psiInf at59--69; IsWu at151--157; sigma DS/minimal-coefficient fields at193--205; actual local/tensor/Haar data at277--327.
\item
  The same directory/WuDecay.lean: normInf,nuInf,N1 and WuReprod/WuWeighted at15--34.
\item
  ProjRows.lean: actual phi/formal-degree coefficient, convolution and projected rows at24--26 and51ff.
\item
  \nolinkurl{Formalism/SharpRestriction/CuspSpanNormBounds/RelevantZeroGramTerms.lean:} IsDS72--76, lowKSub78--80, IsHighestUnit85--87, IsLowUnit89--91, IsHCTypeU96ff; literal lamS and epsS in its Params section.
\item
  \nolinkurl{src/Formalism/LocalRingCardinality/DataTheta.lean:} TPar at501ff, t\_odd and t\_ge; canonical ofWord at519ff; the proof uses t=P.t.
\item
  ConcreteSettingExists.lean is a predicate, not an already instantiated construction. Its additional quantitative gap is not imported into this theorem.
\item
  Checked LIFT28 HJ4--9 supplies the local m=2 anchor.
\item
  Checked LIFT59 BC14--21 supplies the actual group-manifold parameter/half-Cartan and decay dictionary.
\item
  Checked LIFT60 supplies the unchanged finite/geometric Hida adapter.
\end{itemize}

\subsection{10. Primary references and verification status}\label{primary-references-and-verification-status}

Xue, H. (2025). Bessel models for unitary groups and Schwartz homology. Author preprint, May13,2025. \url{https://preprint.press.jhu.edu/ajm/sites/default/files/AJM-h-xue.pdf} §§1.1--1.2 and Theorem1.1 fix the additive character, relevant-line sign, smooth globalization, and relative packet criterion.

Ichino, A. (2021). Theta lifting for tempered representations of real unitary groups (arXiv:2008.06174v2, July6,2021). \url{https://arxiv.org/pdf/2008.06174} §5.3.1, printed14, gives the exact odd-unitary packet-label convention. Only that LLC dictionary is used here, not a conditional theta theorem.

Gan, W. T., Gross, B. H., \& Prasad, D. (2012). Restrictions of representations of classical groups: Examples. In Sur les Conjectures de Gross et Prasad. I, Astérisque346,111--170. Société Mathématique de France. Author PDF: \url{https://abel.math.harvard.edu/~gross/preprints/ggp-evidence4.pdf} Proposition2.1, preprint3--4, supplies the exact complex epsilon sign. Published metadata directly checked at \url{https://www.numdam.org/book-part/AST_2012__346__111_0/} (web turn546000917view1).

Beuzart-Plessis, R. (2020). A local trace formula for the Gan--Gross--Prasad conjecture for unitary groups: The archimedean case. Astérisque418. \url{https://doi.org/10.24033/ast.1120} Primary PDF: \url{https://smf.emath.fr/system/files/filepdf/smf_ast_418.pdf} §7.1, printed159, explicitly asserts absolute convergence and continuity on the ordinary Harish--Chandra Schwartz space by Lemma6.5.1(ii). §7.2, printed162--163, supplies the relative form and Theorem7.2.1's nonvanishing criterion. §2.2, printed42, fixes the linear-first inner product. These were directly reread during the outage: turn305537443view0/1 and turn796043839view0. The ordinary Schwartz assertion is kept separate from the general weak-Schwartz regularization.

Kassel, F., \& Kobayashi, T. (2016). Poincaré series for non-Riemannian locally symmetric spaces. Advances in Mathematics287,123--236. \url{https://doi.org/10.1016/j.aim.2015.08.029} Author PDF: \url{https://www.ihes.fr/~kassel/StableSpec.pdf} §7.1 printed49 and §10.2 printed80--81 identify the full minimal K times K module and the group-manifold discrete series. The additional trace extraction RW10--12 is proved here.

Direct targeted checks during the outage:

\begin{itemize}
\tightlist
\item
  DO independently read the Xue/GGP/Ichino sign and globalization conventions; this is not a whole LIFT61 review.
\item
  Author directly read the KK full minimal-type and group-manifold passages.
\item
  The recovered exact diagnostic checks the relative packet additions, weighted-gap polynomial, S3 coefficient projection and induced quotient kernel. Its new file replay is recorded separately; it is a diagnostic, not a mechanical proof of Xue, the Flensted--Jensen estimate or the global source constructor.
\item
  Fresh saved-file whole reading, source reconciliation, replay outputs and SHA fingerprints are recorded in the companion verification record. Exact primary bibliography titles were checked during the outage (Ichino turn546000917view0; GGP turn546000917view1; BP turn546000917view2).
\end{itemize}

\subsection{Supplemental primary anchor retained during the outage}\label{supplemental-primary-anchor-retained-during-the-outage}

Knapp (1983), \emph{Investigations of unitary representations of semisimple Lie groups}, printed p.~26 immediately before §3 and the §3 example, gives the discrete-series Blattner K-type of highest weight Λ−2ρ\_K+ρ with multiplicity one and identifies it as the unique minimal K-type. The author's PDF is \url{https://www.math.stonybrook.edu/~aknapp/pdf-files/investigations.pdf}. This supplies the representation-theoretic premise used in Section 3.0. The literal source identification was freshly read and is proved there; the theorem is applied to the genuine discrete-series realization, without importing a highest-weight condition into IsLowUnit.

\clearpage

\section{AM-LIFT-061-U --- uniqueness of the actual local Wu vector}\label{am-lift-061-u-uniqueness-of-the-actual-local-wu-vector}

\textbf{Canonical checked version --- 2026-11-22T16:19:19.391528+00:00.} The complete final proof and its source reconciliation were read by root before this integration. Exact diagnostics are separately classified and do not mechanically prove analytic inputs. Historical recovery wording below records the export history; the current scope is: The fixed local projected space has dimension one; the unit Wu vector is unique up to phase, with no new evaluation of its retained costs.

\subsection{Optional local uniqueness corollary (mathematical draft; not a new cost estimate)}\label{optional-local-uniqueness-corollary-mathematical-draft-not-a-new-cost-estimate}

For the fixed actual local data of LIFT61, put Hloc = L²(S\G,V\_rho;psi·rho), P=R(phi). Then ran(P) is one-dimensional.

Indeed, for every bounded G-intertwiner T:sigma→Hloc, right Lie derivatives of T(v) are T(dsigma(X)v). Local Sobolev in one fixed coordinate chart of the smooth quotient S\G gives a smooth representative and makes evaluation at the identity a continuous map on sigma\^{}infty. The resulting functional ell\_T(v⊗bar z)=\textless z,T(v)(1)\textgreater{} belongs to Hom\_S(sigma\^{}infty tensor rho\^{}vee,psi), because left covariance and right equivariance give T(sigma(s)v)(1)=psi(s)rho(s)T(v)(1). If ell\_T=0, then evaluation of T(v) at every g is zero by right equivariance, so T=0. Thus this bounded-intertwiner space injects into the smooth Bessel Hom. Xue §1.2, printed p.5 (PDF index4), lines218--224 explicitly records its dimension at most one, citing Sun--Zhu/Jiang--Sun--Zhu.

The exact evaluation-at-w bijection of LIFT62's elementary matrix-unit lemma applies to this arbitrary unitary R; hence dim ran(P)≤1. LIFT61's positive-kernel construction gives P≠0, so equality holds.

Consequently any two continuous actual IsWu functions with the same data and actual unit L² normalization differ by one complex phase, first as L² classes and then pointwise by continuity. Their absolute sup norm and weighted L¹ moments coincide. The auxiliary compact bump and its positive seed energy E therefore do not create a freedom to improve the actual normalized Wu costs. This result does not evaluate E or those invariant costs, nor does it prove a global Bessel/Mram period.

In particular the continuous kernel of the rank-one projection has the exact factorization K(x,y)v = W(x)\textless W(y),v\textgreater{} for any chosen unit Wu vector W. Equality first holds against compactly supported smooth sections and then pointwise by continuity. Therefore ∫\emph{\{S\G\} tr K(x,x) dx = 1, \textbar\textbar W\textbar\textbar{}}∞² = sup\_x tr K(x,x), and ∫\emph{\{S\G\} nu(x)\^{}A \textbar\textbar W(x)\textbar\textbar{} dx = ∫}\{S\G\} nu(x)\^{}A sqrt(tr K(x,x)) dx. These identities retain the actual Haar and formal-degree normalization and express the remaining real costs directly through the actual relative kernel. They are exact formulas, not numerical evaluations or upper bounds.

\clearpage

\section{AM-LIFT-062 --- Actual rank-one reconstruction and cuspidality of the full projected bottom space}\label{am-lift-062-actual-rank-one-reconstruction-and-cuspidality-of-the-full-projected-bottom-space}

\textbf{Canonical checked version --- 2026-11-22T16:19:19.391528+00:00.} The complete final proof and its source reconciliation were read by root before this integration. Exact diagnostics are separately classified and do not mechanically prove analytic inputs. Historical recovery wording below records the export history; the current scope is: Full projected-bottom reconstruction and cuspidality in the actual smooth multiplier and Haar conventions; arithmetic STiesR.rowTie for J/032 remains unconstructed.

\textbf{Status.} Fresh reconstruction after recovery, incorporating the exact source-continuity and extension-order corrections. The historical outage draft remains unchanged. Final whole-file reading and finite convention replay are recorded separately. The analytic result has the scope below; no Lean build or concrete arithmetic STiesR construction is asserted.

\subsection{1. Purpose and precise scope}\label{purpose-and-precise-scope}

This result supplies the actual Hilbert-space reconstruction behind the rank-one evaluation predicate in the row formalism. It also proves that the entire discrete-series projected space at the fixed finite level is cuspidal. That conclusion does not require replacing the genuine finite good-polynomial projector by a chosen cuspidal finite constituent.

The result retains the actual Haar measure, the actual unit minimal vector, and the actual finite level and finite projector. It is not a nonvanishing theorem for the prescribed Bessel or Mram period. It gives no positive quantitative lower bound for such a period. The recovered definitions show that the last equality involving the independently presented arithmetic space what2.EE and its map sEE follows conditionally from an actual STiesR.rowTie witness. Constructing that witness for the requested J/032 package remains open.

The local Wu construction is ALP-LIFT-61 and is separate from the present global reconstruction. Its weighted integrability proves that the coefficient used below belongs to L1.

\subsection{2. Matrix-unit notation and exact normalization}\label{matrix-unit-notation-and-exact-normalization}

Let G be unimodular, let (sigma,V) be an irreducible unitary discrete-series representation for the actual Haar dg, and let d\_sigma be its formal degree. Inner products are conjugate-linear in the first variable. Fix a unit vector w. Assume \[
 \phi(g)=d_\sigma\overline{\langle w,\sigma(g)w\rangle}\in L^1(G).
 \tag{ER1}
\] For v,u in the algebraic span D of sigma(G)w put \[
 c_{v,u}(g)=d_\sigma\overline{\langle v,\sigma(g)u\rangle}.
 \tag{ER2}
\] Every such coefficient is a finite linear combination of two-sided translates of phi and is in L1.

Schur orthogonality, with this actual Haar, gives \[
 c_{v,u}*c_{v',u'}=\langle u,v'\rangle c_{v,u'},\qquad
 c_{v,u}^{*}=c_{u,v}.
 \tag{ER3}
\] For completeness, the integrand in the convolution is \[
 d_\sigma^2
 \langle\sigma(h)u,v\rangle
 \langle\sigma(h^{-1}g)u',v'\rangle.
\] Rewrite the second factor as \(\langle\sigma(g)u',\sigma(h)v'\rangle\), then apply ordinary formal-degree orthogonality. The result is \(d_\sigma\langle u,v'\rangle\langle\sigma(g)u',v\rangle\), which is ER3. In particular phi\emph{=phi and phi}phi=phi.

Let R be any strongly continuous unitary representation on a Hilbert space H. Use \(R(f)=\int_G f(g)R(g)\,dg\), so \(R(f_1*f_2)=R(f_1)R(f_2)\). Then \[
 P=R(\phi)
 \tag{ER4}
\] is an orthogonal projection.

\subsection{3. Explicit inverse to evaluation at the fixed minimal vector}\label{explicit-inverse-to-evaluation-at-the-fixed-minimal-vector}

For F in ran P define, initially for v in D, \[
 T_F(v)=R(c_{v,w})F.
 \tag{ER5}
\] ER3 gives the polarized identity \[
 \langle T_F(v),T_{F'}(v')\rangle
 =\langle v,v'\rangle\langle F,F'\rangle.
 \tag{ER6}
\] Indeed, \[
 R(c_{v,w})^*R(c_{v',w})
 =R(c_{w,v}*c_{v',w})
 =\langle v,v'\rangle P.
\] Consequently T\_F extends to V with norm exactly \textbar\textbar F\textbar\textbar{} and \(T_F(w)=F\).

Left translation of the coefficient gives \[
 R(g)T_F(v)=T_F(\sigma(g)v).
 \tag{ER7}
\] Thus T\_F is a bounded G-intertwiner. Conversely, for any bounded intertwiner T from sigma to R, Schur orthogonality gives \(\sigma(c_{v,w})w=v\), hence \(T_{T(w)}=T\). Therefore \[
 \operatorname{Hom}_G^{\mathrm{bd}}(\sigma,H)
 \longrightarrow \operatorname{ran}P,\qquad T\longmapsto T(w)
 \tag{ER8}
\] is a bijection with inverse ER5. The multiplicity-space inner product is determined exactly by \[
 \langle T(v),T'(v')\rangle
 =\langle v,v'\rangle\langle T(w),T'(w)\rangle.
 \tag{ER9}
\]

No decomposition of the ambient continuous spectrum, no arbitrary choice of an isometric summand, and no finite-dimensional multiplicity hypothesis is used.

If H' is a closed G-invariant subspace, the same construction restricts to H'. In particular it preserves every actual finite-level character condition and every bounded finite projector commuting with G.

\subsection{4. Actual finite-level arithmetic components and pointwise smooth representatives}\label{actual-finite-level-arithmetic-components-and-pointwise-smooth-representatives}

Apply the preceding construction to the right regular action of \[
 G_\infty=\prod_{\tau:F\hookrightarrow\mathbb R}U(p_\tau,q_\tau)
\] on the actual fixed-finite-level automorphic L2 space. Its finite adelic double-coset decomposition gives arithmetic components \(\Gamma_c\backslash G_\infty\), with the actual component masses. A finite-level character can be removed for this argument by passing to its open kernel; the result then descends. No level change is made in the final space.

For each component there is a positive c\_Gamma such that every nonidentity element of Gamma has matrix distance at least c\_Gamma from the identity in the product archimedean embedding. This is because Gamma lies in a fixed fractional integral matrix lattice. Fix a small exponential chart B at the identity. At g, a ball of radius \[
 \epsilon_g=c(1+\|g\|)^{-2}
 \tag{ER10}
\] injects into the quotient, after decreasing c if necessary. If a translate overlapped itself, some nonidentity gamma would satisfy \[
 \|\gamma-I\|\le C\|g\|\|g^{-1}\|\epsilon_g.
\] For the fixed unitary Gram form \(\|g^{-1}\|\le C'\|g\|\), contradicting the matrix-lattice separation.

Scaled local Sobolev on this chart, using a single integer order s greater than half the real dimension, yields \[
 |f(g)|\le C(1+\|g\|)^M
 \sum_{|\alpha|\le s}\|R(X_\alpha)f\|_{L^2}.
 \tag{ER11}
\] The exponent M and order s are fixed for the component. The same inequality applies to \(R(X)f\), with exactly the same exponent M.

For \(f=T_F(v)\), intertwining implies \[
 R(X)T_F(v)=T_F(d\sigma(X)v).
 \tag{ER12}
\] Therefore \[
 |R(X)T_F(v)(g)|
 \le C(1+\|g\|)^M\|F\|
 \sum_{|\alpha|\le s}\|d\sigma(X_\alpha X)v\|.
 \tag{ER13}
\] This proves both the smooth representative and uniform moderate growth. Crucially, the same s and M are used on \(d\sigma(X)v\), so the growth exponent is independent of X. Only the Sobolev constant changes with X and v.

For K-finite v, the image is also K-finite and Z(g)-finite, since sigma is admissible with an infinitesimal character. Thus its image belongs to Wallach's automorphic-function space A(Gamma,G) intersected with L2.

For the external tensor product of the real sigma\_tau, ER13 is bounded by a product of local smooth-vector Sobolev seminorms. This supplies the pointwise continuous multilinear evaluation required by the actual multiplier-space definition, rather than merely an algebraic Harish--Chandra map.

\subsection{5. Full projected cuspidality by Wallach}\label{full-projected-cuspidality-by-wallach}

Wallach (1984), Theorem4.3, printed pp.233--234, states that a (g,K)-map from a tempered module into \(\mathcal A(\Gamma,G)\cap L^2(\Gamma\backslash G)\) has cuspidal image. His group in §2 is connected semisimple over Q; the real group may be disconnected.

The actual unitary group has compact real centre. At any fixed finite component, det Gamma is a finite group. In fact compactness at every finite place makes each determinant an integral unit of E, and the norm-one condition gives absolute value one at every complex embedding. Such a unit is a root of unity. Hence \[
 \det\Gamma\subset\mu_E.
 \tag{ER14}
\] Pass to the finite-index subgroup Gamma\^{}1=Gamma intersect SU. It is arithmetic in Res\_\{F/Q\}SU(Y0). The restriction of sigma to the derived group is irreducible: the full group is the product of its derived group and compact centre, and that centre acts scalarly. It is discrete series, since the centre is compact, and therefore tempered. Integrating over the compact central factor shows that the corresponding restricted functions are square integrable.

Apply Wallach to Gamma\^{}1 and the restricted Harish--Chandra module, using ER13. Every proper rational parabolic constant term vanishes. Restoring the compact centre and descending the finite-index cover preserves those vanishing conditions, since the unipotent radicals are unchanged. Consequently \[
 T_F(V^\infty)\subset L^2_{\mathrm{cusp}}
 \quad\text{for every }F\in\operatorname{ran}R(\phi).
 \tag{ER15}
\] The assertion first holds for K-finite vectors. It extends to all smooth vectors by the Sobolev continuity of T\_F and continuity of the compact unipotent-quotient integral, or by closedness of the cuspidal Hilbert subspace.

Thus the full rank-one projected bottom space is cuspidal, including the full genuine good-polynomial finite cut. No assertion about the finite cuspidal support of every summand of e0 is needed.

\subsection{6. Conjugation and the exact row-evaluation space}\label{conjugation-and-the-exact-row-evaluation-space}

Complex conjugation on automorphic functions takes the sigma projection to the conjugate representation projection. The source parameter parB is dualPar of the bottom Dpar, and the corresponding model is the conjugate/contragredient discrete series. Choose its unit minimal vector v0 to correspond to bar w.

Under the already checked finite-level and finite-projector identifications, the same construction with bar phi gives \[
 \operatorname{Mult}_{\mathrm{actual,good}}
 \simeq \operatorname{ran}R(\bar\phi)\cap
 \{\text{the actual finite level and genuine good cut}\},
 \qquad T\mapsto T(v_0).
 \tag{ER16}
\] The inverse is the explicit coefficient convolution ER5. The pointwise representatives are those in ER13. ER15 proves the cusp field of the multiplier definition for the entire range. The actual conjugated finite coefficients are retained. The translated supplement proves their precise transport from RowLevelMatchC, without requiring coefficientwise reality. Coefficientwise reality, when supplied by the genuine R026 projector, is compatible with this identity but is not substituted for it.

This is the analytic content of the source's evaluation/rank-one Hom predicate. It preserves the full finite projector and does not replace it by the one constituent used in the earlier cusp-existence proof.

\subsection{6A. The converse for the source's smooth multiplier definition}\label{a.-the-converse-for-the-sources-smooth-multiplier-definition}

The source multiplier need not be assumed Hilbert-bounded in advance. Let T be one of its continuous pointwise, real-group equivariant maps from the smooth model to cuspidal automorphic functions, at the actual finite level. At a fixed finite component, evaluation at the identity is a continuous functional ell on the smooth representation: \[
 T(v)(g)=\ell(\sigma(g)v).
 \tag{ER19}
\] The literal continuity field requires some precision. \nolinkurl{BesselHermitianTheta/TiedType.lean} 18--26 defines smooth vectors by all finite products of Lie exponentials, and \nolinkurl{Sobolev₂} by a finite sum of Lie-word Hilbert norms. Smooth vectors form a linear subspace, so \nolinkurl{DSModel.sm}, defined as their span, is exactly that subspace. \nolinkurl{IsMultMap.cont} supplies continuity separately in each real-coordinate input. A separately continuous finite multilinear scalar map on Fréchet spaces is jointly continuous: induct on the number of factors, choose countable decreasing product-neighborhood bases in the last factors, and let A\_n be the closed balanced convex set of first inputs on which the value is bounded by one on the nth product neighborhood. Their union is the first Fréchet space. Baire gives one A\_n with interior; subtract two points and rescale to obtain a product-neighborhood bound. Multilinearity gives a product of continuous seminorms, each dominated by finitely many Lie-word Hilbert norms. Thus the joint product-Sobolev bound used below is proved from the separate source field.

In the forward construction, the source asks for an algebraic linear functional on the whole local Hilbert space whose restriction to smooth vectors obeys this bound. Extend the functional algebraically from the smooth subspace. Its required bound is unchanged; no Hilbert-continuity of that algebraic extension is asserted.

The prescribed Sobolev continuity gives \(|\ell(v)|\le C p_s(v)\) for some finite order s. Unitarity and \(d\sigma(X)\sigma(g)=\sigma(g)d\sigma(\operatorname{Ad}(g^{-1})X)\) give \[
 p_s(\sigma(g)v)\le C_s(1+\|g\|)^{M_s}p_s(v).
\] The same bound applied to \(d\sigma(X)v\) supplies uniform moderate growth for every derivative of T(v), with exponent M\_s independent of X. Independently, the actual source \nolinkurl{IsAutFunSm} in \nolinkurl{TiedSmooth.lean} 44--51 already requires one growth exponent for all derivatives. Either reading supplies precisely the hypothesis needed here.

For K-finite v, T(v) is therefore a K-finite, Z(g)-finite cuspidal automorphic function. It and every right derivative decay rapidly on Siegel sets, hence are in L2 and are smooth vectors for the right regular Hilbert representation. This is the classical cusp estimate recalled in Wallach, §4, printed232. The smooth extension with explicit seminorm control is also given by Miller--Schmid (2012), Theorem2.24.

Let kappa be the minimal K-type of sigma, of multiplicity one. Define on the Harish--Chandra module \[
 B(v,u)=\langle T(v),T(u)\rangle_{L^2}.
\] Compact Schur's lemma gives \(B|_\kappa=c\langle\, ,\,\rangle\), with \(c=\|T(w)\|^2\). Moving the unitary adjoint of any universal-enveloping-algebra element from the first variable to the second is justified in L2. For v,u in kappa and U,V in U(g\_C), \[
\begin{aligned}
 B(Uv,Vu)
 &=B(v,U^*Vu)\\
 &=B(v,P_\kappa U^*Vu)\\
 &=c\langle v,P_\kappa U^*Vu\rangle\\
 &=c\langle Uv,Vu\rangle.
\end{aligned}
 \tag{ER20}
\] Here \(P_\kappa\) is the compact isotypic projector. The Harish--Chandra module is generated by kappa, so ER20 holds everywhere on it. Thus the Harish--Chandra restriction of T extends to a bounded Hilbert operator A of norm \(\sqrt c=\|T(w)\|\). At this point full-group intertwining has not yet been deduced. If c=0, that restriction is zero; the smooth identification in the next paragraph proves that the original map is zero too.

K-finite vectors are dense in the smooth-vector Fréchet topology, by compact approximate identities with finite Peter--Weyl support. For v\_n tending smoothly to v, the original T(v\_n) tends pointwise to T(v), by the source Sobolev continuity. The bounded Hilbert extension A satisfies A(v\_n)-\textgreater A(v) in L2. An L2-convergent subsequence converges almost everywhere, so A(v)=T(v) almost everywhere. Now the original full G-equivariance implies A sigma(g)=R(g)A on the smooth vectors, hence on the whole Hilbert space by density. Only after this step is ER13 applied to the extension. It gives the canonical continuous representative, which equals the original T(v) everywhere because both are continuous and agree almost everywhere. For the literal product of local smooth-vector spaces, first linearize the source multilinear map on the algebraic tensor product of the local Harish--Chandra modules. Its minimal compact type is the tensor product kappa and still has multiplicity one, so ER20 gives a bounded extension to the actual Hilbert tensor product. Approximate each local smooth input by K-finite inputs and use the joint product-Sobolev bound just deduced from separate continuity to obtain pointwise convergence. The same L2 subsequence argument identifies the original multilinear map on every decomposable smooth tensor, and hence on finite sums. No unmentioned nuclear tensor-product identification is needed.

Applying the same reasoning to two maps and polarizing gives the exact source multiplier inner product \[
 \langle T(v),T'(v')\rangle
 =\langle v,v'\rangle\langle T(w),T'(w)\rangle
 \tag{ER21}
\] for all smooth inputs. This proves the converse direction of ER16 for the actual smooth multiplier definition, rather than only for an a priori Hilbert-bounded subclass.

If the source Sobolev seminorms are written using the compact Casimir, their equivalence with full smooth-vector Sobolev seminorms follows here from the fixed infinitesimal character. In Cartan-orthonormal coordinates the elliptic Nelson Laplacian is \(2\Omega_K-\Omega_G\); the second term acts as a fixed scalar on sigma. No unproved continuity of an arbitrary algebraic Harish--Chandra functional is used.

\subsection{7. Exact measure and inner-product statement}\label{exact-measure-and-inner-product-statement}

Let mu\_A be the actual Hida Hilbert-space quotient measure and nu\_0 the actual automorphic quotient measure. If the existing Haar transport gives \[
 \mu_A=c\,(\iota_X)_*\nu_0,\qquad c>0,
 \tag{ER17}
\] then ER9 gives, for the actual evaluation map j, \[
 \|j(T)\|_{\mu_A}^2
 =c\,\operatorname{Re}ipM(T,T).
 \tag{ER18}
\] No unidentified scalar is set equal to one. The positive c is exactly the source's Haar-ratio field. The bilinear identity follows by polarization, with the source inner-product conventions.

The complete recovered STiesR and RowDictTie definitions supply a precise conditional statement. If j\_row:what2.EE→EEC is the rowTie linear-isometry witness and j\_eval:EP→EEC is the analytic evaluation equivalence with this same Haar ratio, then k=j\_eval.symm∘j\_row gives the required k/sEE equality. An actual STiesR is a field of RealPack, not a consequence of its name. No concrete constructor of that field for the requested J/032 package was recovered. The separately saved translated corollary retains the actual tS, level character and scalar uS*kappaS; it is not identified with unshifted RowsJoinPinned by dropping tS.

\subsection{8. Consequences and retained gates}\label{consequences-and-retained-gates}

Together with LIFT60's finite/geometric adapter and LIFT61's actual Wu construction, ER1--18 provide:

\begin{itemize}
\tightlist
\item
  the actual reproducing projection and its full rank-one multiplier reconstruction;
\item
  pointwise smooth, moderate-growth representatives with the required Sobolev continuity;
\item
  cuspidality of the entire projected space;
\item
  exact evaluation and norm transport for the actual Haar measures.
\end{itemize}

They do not supply:

\begin{itemize}
\tightlist
\item
  a nonzero prescribed Bessel/Mram period;
\item
  an effective lower bound for a detector or inverse period map;
\item
  a full unrelated Host/D7 constructor;
\item
  an unverified equality between independently presented arithmetic row records.
\end{itemize}

\subsection{9. Primary and exact source anchors}\label{primary-and-exact-source-anchors-1}

Wallach, N. R. (1984). On the constant term of a square integrable automorphic form. In Operator algebras and group representations, Vol.II (Neptun,1980), pp.227--237. Pitman. Monographs and Studies in Mathematics18. Author-hosted original: \url{https://mathweb.ucsd.edu/~nwallach/tempered-cuspidal.pdf} §2, printed227--229, defines A and the arithmetic hypotheses. Theorem4.3 is printed233--234 (PDF page indices7--8). Direct author read: web turn715073923view0. Root independent direct source/interface read: turn824614836view0. These are targeted checks, not a saved-file whole-proof review.

Exact source bodies read before the outage and reread after recovery, with fresh pinned-source metadata retained:

\begin{itemize}
\tightlist
\item
  \nolinkurl{Formalism/LocalToGlobal/IndefiniteReal/ScalarOnPartsFull.lean:} dualParS, DSModelH, smooth-vector multiplier definition, IsMultIP.
\item
  \nolinkurl{Formalism/LocalToGlobal/SiegelLatticeTheory/HaarTransversals/ProjRows.lean:} coefficient and Rph, SchurConv, conditional HomBij.
\item
  \nolinkurl{Formalism/LocalToGlobal/SiegelLatticeTheory/HaarTransversals/RowsDef.lean:} HH and conditional RowsBij.
\item
  \nolinkurl{Formalism/LocalToGlobal/SiegelLatticeTheory/FlagStabilizerTensors/RowsJoinPinned.lean:} actual instOf/iotaX, Haar ratio, evaluation map, final k/sEE tie.
\item
  \nolinkurl{Formalism/LocalToGlobal/IndefiniteReal/AutomorphicEvaluation.lean:} IsLowVec, ev and inner-product identities.
\item
  RowDict.lean: good finite cut and conjugation.
\item
  LinksDefs.lean: the actual bottom RealDat and good multiplier space.
\end{itemize}

The complete source reconciliation is retained in \nolinkurl{lift62_source_reconciliation_pending.md}; its definitions and proof-order corrections are incorporated above. Exact finite convention checks are in \nolinkurl{lift062_exact_conventions_certificate.py}, which tests complex coefficients and nonunit Haar factors but does not mechanically certify the analytic theorems. No Lean compilation was run.

Additional directly read primary source: Miller, S. D., \& Schmid, W. (2012). On the rapid decay of cuspidal automorphic forms. Advances in Mathematics,231(2),940--964. \url{https://doi.org/10.1016/j.aim.2012.06.001} Author preprint: \url{https://arxiv.org/pdf/1106.2149} Theorem A (PDF page 2), Theorem2.23 and Theorem2.24 (the latter PDF page index 11, one-based page 12) explicitly retain smooth non-K-finite vectors and the controlling smooth-vector seminorm. Direct read refs turn984883834view0 and turn486656863view0/1. The present converse needs only the classical K-finite cusp estimate plus smooth density; no global period nonvanishing is inferred from rapid decay.

\subsection{Supplemental primary anchor retained during the outage}\label{supplemental-primary-anchor-retained-during-the-outage-1}

Knapp (1983), \emph{Investigations of unitary representations of semisimple Lie groups}, printed p.~26 immediately before §3 and the §3 example, gives the discrete-series Blattner K-type of highest weight Λ−2ρ\_K+ρ with multiplicity one and identifies it as the unique minimal K-type. The author's PDF is \url{https://www.math.stonybrook.edu/~aknapp/pdf-files/investigations.pdf}. The freshly read source \nolinkurl{RelevantZeroGramTerms.lean} 72--91 defines lowKSub as the supremum of all copies of this minimal highest-weight type and IsLowUnit as norm one plus membership. For the genuine discrete series, multiplicity one makes lowKSub exactly kappa. IsHighestUnit is separately defined and is not assumed here. The exact minimal-type source dictionary is also written in LIFT61 Section 3.0.

\clearpage

\section{AM-LIFT-062-T --- translated evaluation with the literal level, projector and Haar}\label{am-lift-062-t-translated-evaluation-with-the-literal-level-projector-and-haar}

\textbf{Canonical checked version --- 2026-11-22T16:19:19.391528+00:00.} The complete final proof and its source reconciliation were read by root before this integration. Exact diagnostics are separately classified and do not mechanically prove analytic inputs. Historical recovery wording below records the export history; the current scope is: Translated EvalTieT with literal inverse-t level transport, complex projector conjugation, actual coset volumes and Haar ratio; no unshifted or concrete J/032 pin is inferred.

\textbf{Status:} new controlled pending proof. This is a conditional source-interface theorem for the genuine models constructed in LIFT61--62. It does not construct the arithmetic STiesR instance for the requested J/032 package. No existing proof or source file is modified, and no Lean build is used.

\subsection{1. Data and conclusion}\label{data-and-conclusion}

Retain the actual unitary adelic group A=GA(D,Y0), its rational subgroup and quotient Haar nu, and the actual Hida group A'=X.B.I.GA with quotient Haar mu\_X. Let iota:A→A' be the literal \nolinkurl{iotaX} belonging to the actual source instance equality. In particular it identifies rational points and finite/real factors. Retain the positive Haar ratio c defined by \[
 (\iota^{-1})_*\mu_X=c\,\nu,\qquad c>0.
 \tag{TE1}
\] This equality is interpreted on the actual rational quotients; it follows equally from the source full-Haar identity and its fundamental domains.

Let M be the actual bottom smooth discrete-series model, J and chi its actual finite level and character, and t be the supplied finite adele. Assume the literal source condition \[
 \operatorname{ConjLevel}(t^{-1},J,\chi,J_X,\chi_X),
 \quad J_X=\iota^{-1}(1\times X.K2),\quad
 \chi_X(k)=\overline{X.\mathrm{chi2}(\iota(k)_f)}.
 \tag{TE2}
\] Equivalently, J\_X=tJt\^{}\{-1\}, with chi\_X(tkt\^{}\{-1\})=chi(k). The displayed inverse on t is essential.

Retain the actual real \nolinkurl{ConjVec} and \nolinkurl{RealSplit} identities of \nolinkurl{ThetaIdentityOfMapsAll.lean}: antiunitary maps from the local M models to the Hida local models intertwine the actual real group identifications and send the chosen local unit vectors v0 to X.B.w. Retain the integrable, formal-degree-normalized matrix units of LIFT61--62. Finally assume the actual \nolinkurl{RowLevelMatchAt} B/C finite identities and the genuine residue-group/level-normalization hypotheses described in Section 3 below.

Write E\_R for \nolinkurl{RealDat.EEat(levR(F,true))}, namely the full actual smooth multiplier space at (J,chi), intersected with the fixed space of its given finite operator e\_R. Then there is a linear equivalence \[
 j_t:E_R\xrightarrow{\sim}\operatorname{EEC}(X,G),\qquad
 j_t(T)(x)=T(v_0)(\iota^{-1}(x)t)
 \tag{TE3}
\] with the literal almost-everywhere interpretation in the source Hilbert space, and \[
 \|j_t(T)\|_{\mu_X}^2
   =c\,\operatorname{Re}ipM(T,T).
 \tag{TE4}
\] Thus these hypotheses imply the source \nolinkurl{EvalTieT} at this same t and c.~They do not replace t by 1 and do not identify the independently supplied \nolinkurl{D.what₂.EE}.

\subsection{2. Finite translation and full unprojected evaluation}\label{finite-translation-and-full-unprojected-evaluation}

For an automorphic function f define \[
 U_t f(x)=f(\iota^{-1}(x)t).
 \tag{TE5}
\] The actual unitary adelic group is unimodular. Right translation by a fixed finite t is consequently unitary on its rational quotient. It commutes with every real-group action and with every real derivative. It preserves smooth automorphy, cuspidality and the same real infinitesimal character. The finite-level calculation is exact: if k' in J\_X and k=t\^{}\{-1\}k't in J, then \[
 (R_t f)(gk')=f(gk't)=f(gtk)=\chi(k)f(gt)
             =\chi_X(k')(R_t f)(g).
 \tag{TE6}
\] The inverse is right translation by t\^{}\{-1\}, with the inverse group identification. Therefore this is an equivalence of the actual level spaces, not a mere containment.

After conjugating function values and applying the antiunitary \nolinkurl{ConjVec} model maps, the real representation is precisely the Hida representation. The chosen v0 maps to its actual unit vector w. LIFT62's matrix-unit reconstruction, in both directions, identifies the full smooth multiplier space at (J,chi) with the conjugate of HHX. The finite character in HHX is X.chi2 because the conjugated expression in (TE6) has that character. Its finite averaging kernel is the literal \[
 \mathbf1_{X.K2}(k)\,\overline{X.\mathrm{chi2}(k)}/\operatorname{vol}(X.K2)
\] from \nolinkurl{HidaSetting.e_eq}. Its real factor is the given formal-degree-normalized matrix unit. Haar rescaling changes the formal degree inversely and hence does not change this projection operator. Formula (TE1) retains the global norm scale c.

This proves a full linear evaluation equivalence \[
 j_t^{\mathrm{all}}:\operatorname{Mult}(M,J,\chi)
       \xrightarrow{\sim}\overline{HHX},
 \tag{TE7}
\] where the bar denotes pointwise conjugation of the subspace, not a different normalization of its norm. LIFT62's converse applies to every smooth multiplier, using the literal finite-Sobolev continuity and invariant-form argument; it is not restricted to an a priori bounded-Hom subclass.

\subsection{3. The exact finite coset normalization}\label{the-exact-finite-coset-normalization}

Put K=X.B.I.Kset, K0=X.B.I.Kg(X.B.g), and Gamma=G0g(D.nd). The actual \nolinkurl{IsRedHom} condition supplies a surjective residue homomorphism \[
 \mathrm{red}:K\longrightarrow\Gamma,\qquad\ker(\mathrm{red})=K_0.
 \tag{TE8}
\] \nolinkurl{LevelSubgroupNormalized} includes K0 subset X.K2 and chi2\textbar K0=1. The finite additive/residue bases are the actual rank-r local bases. Any two bases witnessing \nolinkurl{GoodAt(k,gamma)} give conjugate residue matrices at every place. Indeed the r spanning vectors in the rank-r residue vector space are bases, and change of basis conjugates the matrix. The genuine e0 is central in the finite group algebra, so its coefficient is constant on these conjugacy classes. This proves the needed \nolinkurl{EgWellDef}; it is not inferred from an arbitrary choice in the definition of eg. Existence of red gives \nolinkurl{RedExists}.

Let b(gamma)=D.ntypes.e0.coeff(gamma) and choose any representatives k\_gamma in K. The actual kernel is \[
 e_g(k)=\mathbf1_K(k)\,b(\mathrm{red}(k))/v_0,
 \qquad v_0=\operatorname{vol}(K_0).
 \tag{TE9}
\] For every right-K0-invariant smooth function F, \[
 R_gF(x)=\int_K\frac{b(\mathrm{red}(k))}{v_0}F(x(1,k))\,dk
        =\sum_{\gamma\in\Gamma}b(\gamma)F(x(1,k_\gamma)).
 \tag{TE10}
\] Each coset has volume v0. Hence there is no leftover index, dimension, group-order or Haar factor in this finite sum. Changing k\_gamma within its coset does not change the value. This is exactly the operator specified by \nolinkurl{GoodOp.Rg_spec} on these vectors.

The finite operator \nolinkurl{RealDat.e0} is the literal sum \nolinkurl{gaL} with coefficients \[
 a(\sigma\gamma)=\overline{b(\gamma)}
 \tag{TE11}
\] by \nolinkurl{RowLevelMatchC}; sigma is the actual finite group isomorphism in \nolinkurl{RowLevelMatchAt}. The calculation does not require coefficientwise reality of b. Reality, when available, is an additional fact rather than a replacement for (TE11).

\subsection{4. Projector transport before taking fixed spaces}\label{projector-transport-before-taking-fixed-spaces}

Set \[
 \kappa_t(k)=t^{-1}\iota^{-1}(1,k)t.
\] \nolinkurl{RowLevelMatchB} states, for every full level multiplier T, every smooth input v and every g, that right translation by this element equals the actual finite action \nolinkurl{actL} attached to sigma(red k). Thus \[
 \begin{aligned}
 \operatorname{ev}_t(\operatorname{actL}(\sigma\gamma)T)(x)
 &=T(v_0)(\iota^{-1}(x)t\kappa_t(k_\gamma))\\
 &=T(v_0)(\iota^{-1}(x(1,k_\gamma))t).
 \end{aligned}
 \tag{TE12}
\] Here ev\_t denotes the evaluation formula on the common ambient principal-level multiplier space, even when an individual summand leaves the smaller J-level. Summing (TE12) with (TE11), and writing C for pointwise conjugation, gives the exact operator identity \[
 \boxed{\operatorname{ev}_t(e_RT)=C\,R_g\,C\,\operatorname{ev}_t(T).}
 \tag{TE13}
\] The evaluations in (TE12) are pointwise on smooth representatives; the associated Hilbert identity follows as well.

There is a harmless ambient-level point in this argument. Individual finite-group summands need not preserve the smaller character level J. Perform the displayed equality first in the existing common principal-G refinement, keeping the prescribed level away G. All summands are genuine smooth cuspidal multiplier maps there. The evaluation map on that ambient space is injective by LIFT62. This is only an ambient proof space; the theorem's domain remains the original J and chi.

Now if T lies in E\_R, then e\_RT=T, and (TE13) shows that Cj\_t(T) belongs to HHX and is Rg-fixed, exactly the definition of EEC. Conversely take F in EEC. By (TE7) it has a unique full level preimage T. Since Rg(CF)=CF, (TE13) gives equality of the evaluations of e\_RT and T. Ambient principal-level injectivity gives equality of the multiplier maps. Hence e\_RT=T and T lies in E\_R. This proves the bijection (TE3) without supposing that an individual residue lift normalizes J.

\subsection{5. Norm and source expression}\label{norm-and-source-expression}

Right translation in (TE5) preserves the quotient integral, and changing between the two actual quotient Haar measures contributes exactly c.~LIFT62's polarized multiplier identity on the fixed unit vector therefore gives (TE4). No finite index is inserted: both sides use the full adelic quotient measure, and a finite-level class-mass realization already includes its own indices.

The source \nolinkurl{EvalTieT} asks for the unit tuple v0, the linear equivalence (TE3), the displayed almost-everywhere evaluation, and precisely (TE4). All its fields are therefore supplied. Differentiation and the LIFT62 smooth-representative argument show that these evaluation functions are smooth, so an almost-everywhere equality can also be read pointwise when a later restriction uses continuous representatives.

The source-defined \nolinkurl{sEEofTT} is then its chosen evaluation equivalence followed by the fixed anti-isometry \nolinkurl{conjHH}. It is the corresponding canonical anti-linear map for the chosen witness. The supplied \nolinkurl{RealT.seeDef} still retains its actual scalar uS*kappaS and its own abstract Hilbert-model identifications; this theorem does not set those scalars to one.

\subsection{6. What is and is not removed}\label{what-is-and-is-not-removed}

For a genuine fixed family with the literal real-model, Haar and level/projector identities, the translated evaluation-existence part of \nolinkurl{SeeEx} follows from LIFT61--62 and (TE1)--(TE13). This retains tS when t=tS. It does not use a finite translation to change the prescribed arithmetic data.

The final \nolinkurl{RowsJoinPinned} equality with \nolinkurl{D.what₂.EE} additionally needs the actual rowTie witness. If an actual STiesR supplies j\_row:what₂.EE≃EEC, then k=j\_t\^{}\{-1\}j\_row gives the translated row identification with all norms retained. This is a conditional deduction. The existence of that STiesR for the requested J/032 package, and the separate passage to the unshifted \nolinkurl{RowsJoinPinned} convention if requested, remain explicit requirements.

No nonzero prescribed Bessel/Mram period, arithmetic coefficient bound, fixed theta-image saturation, complete Host, or DP44 conclusion is inferred from this evaluation theorem.

\subsection{7. Direct source anchors and verification boundary}\label{direct-source-anchors-and-verification-boundary}

The following source bodies were read after recovery at the pinned commit or in its restored CP23 cache:

\begin{itemize}
\tightlist
\item
  \nolinkurl{IndefiniteReal/EvalTieTOne.lean}, entire file: ConjLevel, EvalTieT and sEEofTT.
\item
  \nolinkurl{IndefiniteReal/AutomorphicEvaluation.lean}, entire file: ev, IsLowVec, exact norm identity and unshifted EvalTie.
\item
  \nolinkurl{IndefiniteReal/ScalarOnPartsFull.lean}: IsMultMap, Mult, goodPart.
\item
  \nolinkurl{IndefiniteReal/LinksDefs.lean}, lines 183--209 and 628--645: actL, gaL, e0, EEat and levR.
\item
  \nolinkurl{SiegelLatticeTheory/FlagStabilizerTensors/RowLevelMatch.lean}, entire file: kap and literal A/B/C.
\item
  \nolinkurl{FlagStabilizerTensors/RowDict.lean}, entire file: eg, GoodOp, EEC, conjHH and RowDictTie.
\item
  \nolinkurl{HaarTransversals/HidaSetting.lean}, e\_eq and LevelSubgroupNormalized; \nolinkurl{RowsDef.lean}, HH and evaluation range; \nolinkurl{ProjRows.lean}, the actual matrix-unit convention.
\item
  \nolinkurl{Definitions/LatticeFrames/ThetaIdentityOfMapsAll.lean}, ConjVec, RealSplit, EvalAll and evS.
\end{itemize}

The genuine reduction and central e0 facts are retained, not reconstructed from arbitrary field names. The LIFT62 smooth-continuity correction is in \nolinkurl{lift62_source_reconciliation_pending.md}. The exact finite diagnostic \nolinkurl{lift062_exact_conventions_certificate.py} tests the signs, genuinely complex coefficients, conjugations, coset volumes and norm scale. Its replay is recorded separately. It does not replace the analytic proof or instantiate the missing arithmetic package.

\clearpage

\section{AM-R-054 --- complete all-top support-eleven removal}\label{am-r-054-complete-all-top-support-eleven-removal}

Author: GPT-6 Astra Pro, ALP\_MACHINE \nolinkurl{/root}; inherited effort is not externally exposed. 22 November 2026 UTC. New joint corollary; it is not a reconstruction of lost historical bytes.

\subsection{1. The unchanged actual normalized kernel}\label{the-unchanged-actual-normalized-kernel}

Keep the actual J56 all-top role assignment, original K044 upper-gap endpoint, permitted packet point, finite projectors, scalar-pencil labels, class masses, central phases, selected compact module and Hilbert normalization. The actual field has \nolinkurl{[E:Q]=336}, \nolinkurl{[F:Q]=f=168}; the Hermitian rank is169. The charged all-top cutoff is \nolinkurl{H=133}, supplied by the fresh R048 reconstruction, RC48.1--10. No uniform-H8192 estimate is replaced by this smaller cutoff.

Let \nolinkurl{s(h)} be the number of active nonintegral relative T places, and let the actual normalized signed summand be the one in K062 AT21--25 and DO059 SC36--37. At each removal stage, intersect the newly bounded set with the remainder left by the preceding stages. Retain every previous numerical error, even when two original upper-bound sets overlap. Thus the actual terms removed at successive stages are disjoint; the numerical accounting is deliberately conservative.

\textbf{Theorem.} There is a remainder \nolinkurl{R_top^(12)} consisting only of original actual summands with \nolinkurl{s(h)>=12} such that the unchanged normalized upper-endpoint energy has the same expansion

\[
 \mathcal E_{\rm top}
 =1+\operatorname{Re}\!\sum_{h\in R_{\rm top}^{(12)}}w(h)
   +\epsilon_{\rm top},
 \qquad |\epsilon_{\rm top}|<2^{-230}.
 \tag{JT1}
\]

Here taking real parts simply writes the real energy identity in real form; it changes no summand or absolute error bound. The uniform family retains its independently proved support floor11 and its31-term error bound. This theorem does not prove the signed sum in JT1 exceeds \nolinkurl{-1+2^-230}.

\subsection{2. Exhaustion of every structural class at support eleven}\label{exhaustion-of-every-structural-class-at-support-eleven}

R048 RC48.10 already removes the complete all-top support-ten sector and all smaller supports, with its annulus and all prior charges retained. We need only exhaust \nolinkurl{s=11}.

Write the original coordinate decomposition as \nolinkurl{E^169=S perpendicular W0}, with \nolinkurl{dim S=164} and \nolinkurl{W0} of signature \nolinkurl{(4,1)}, and put \nolinkurl{A=h_SS}, \nolinkurl{N=h-I}. R051, sections3--5, gives the complete structural reduction at support11: every remaining nonzero term has \nolinkurl{A=I}, is globally unipotent, and has \nolinkurl{rank_E N=2}. Its proof explicitly includes the F/IC8 split, the nonscalar positive-block cases, nonunipotent fixed-W0 and scalar-cross cases, the two globally unipotent F,r1 cells, and the separately retained fixed-W0 rank-one cut. The identity has support zero. No classification hypothesis is inferred from a failed numerical upper bound.

There are now precisely two structural possibilities.

\textbf{Fixed W0.} If h fixes S pointwise, K084 section5 removes every support-eleven type profile except the all-type-II profile. K085 AS1--7 removes that last profile with charge685. The zero-target correction C1 does not affect it: a nonintegral type-II incidence fixes one admitted lattice, and this lattice cannot be the original lattice labelled0. Thus every fixed-W0 term is removed.

\textbf{Scalar cross.} Otherwise set \nolinkurl{x=(h e_n)_S} and let k be the number of nonzero original E-coordinates of x. One has \nolinkurl{1<=k<=164}. To see that zero is not omitted, the two independent vectors \nolinkurl{N e_n} and \nolinkurl{N^2 e_n} span \nolinkurl{im N}, because \nolinkurl{alpha=(N^2)nn>0} and \nolinkurl{N^3=0}. If x were zero, both vectors would lie in W0, since \nolinkurl{im N^2} already lies in W0. All S rows of N would vanish; the identity \nolinkurl{N^dagger=N^2-N} would also make all S columns vanish. Then h would fix S, the preceding case. This is the same original-coordinate argument in DO072 section2 and uses no integral moving basis.

R052 sections4--6 removes every scalar-cross type profile except

\[
 (i,u,t)=(0,11,0),\ (0,10,1),\ (0,9,2),
 \tag{JT2}
\]

where i,u,t count actual types I,II,III. The following table partitions every remaining case.

\begin{longtable}[]{@{}
  >{\raggedright\arraybackslash}p{(\columnwidth - 6\tabcolsep) * \real{0.2308}}
  >{\raggedright\arraybackslash}p{(\columnwidth - 6\tabcolsep) * \real{0.2308}}
  >{\raggedright\arraybackslash}p{(\columnwidth - 6\tabcolsep) * \real{0.2308}}
  >{\raggedleft\arraybackslash}p{(\columnwidth - 6\tabcolsep) * \real{0.3077}}@{}}
\toprule\noalign{}
\begin{minipage}[b]{\linewidth}\raggedright
Actual profile
\end{minipage} & \begin{minipage}[b]{\linewidth}\raggedright
Original coordinate support
\end{minipage} & \begin{minipage}[b]{\linewidth}\raggedright
Complete proved bound
\end{minipage} & \begin{minipage}[b]{\linewidth}\raggedleft
Charge
\end{minipage} \\
\midrule\noalign{}
\endhead
\bottomrule\noalign{}
\endlastfoot
\nolinkurl{(0,9,2)} & all possible k & DO070 FT1--13, corrected by C1 ZC1--11 & 43330 \\
\nolinkurl{(0,10,1)} & \nolinkurl{1<=k<=107} & DO071 SP1--11, corrected by C1 & 838 \\
\nolinkurl{(0,10,1)} & \nolinkurl{108<=k<=164} & DO072 GS1--31, in its joint221-cell bound & included in1482 \\
\nolinkurl{(0,11,0)} & \nolinkurl{1<=k<=164} & DO072 GS1--31, in the same joint bound & included in1482 \\
\end{longtable}

The two DO072 rows share one charge1482; it is not added twice. The sparse DO071 sector is disjoint from its dense t1 sector. The older optional \nolinkurl{k<=104} bound is not charged as an additional removal.

C1 retains the actual permitted label0. If there are r0 zero targets, the target-fibre covolume has exponent \nolinkurl{T^(-8s-r0)}. Summing its different fibre counts with the exact type-III pair weights contributes \nolinkurl{2^t}, once for the relative places. Hence43330 and838 are the corrected sufficient charges;43332 and839 are historical pre-correction calculations, not the charges used here. Every actual source/target pair, inverse-adjoint CM partner, common-orientation pattern, full source-lattice fibre and selected real factor is retained.

The fixed-W0 case and all rows of the table exhaust every term remaining after R051. Therefore no support-eleven term remains. Together with R048 this proves the support assertion in JT1.

\subsection{3. Complete chronological numerical error}\label{complete-chronological-numerical-error}

The adopted all-top baseline is the28-entry field \nolinkurl{all_top_with_rank_one_removal_exponents} in the restored K078 certificate. It already ends with the separately stated K077 rank-one charge1436. The complete current all-top list, in its actual chronological order, is

\[
\begin{split}
(&1999,104546,101550,211446,11671,231,1483877,1307183,9556,\\
 &4084,4085,4084,4086,4088,4086,4085,9199,4088,7423,5183,\\
 &8105,17434,7227,2667,2067,7319,3084,1436,32739,2000,4094,\\
 &4086,5001,4088,1166,195718,93271,34539,685,43330,838,1482).
\end{split}
\tag{JT3}
\]

There are42 entries. Exactly one is231; every other entry is at least685. Thus, even without discarding a single conservative charge,

\[
 \sum_{e\in\mathrm{JT3}}2^{-e}
 \le2^{-231}+41\,2^{-685}
 <2^{-231}+2^{-679}<2^{-230}.
 \tag{JT4}
\]

The uniform31-entry list is unchanged from R053/K085's uniform stage: it has exactly one231 and every other entry at least797. Therefore its sum is at most \nolinkurl{2^-231+30*2^-797<2^-230}. No all-top energy or H133 cutoff is asserted for it.

The companion exact certificate reads the actual canonical stage files, verifies all16 ordered prefixes from K078 through DO072, checks the C1 correction against the preserved pre-correction DO070 file, and verifies the whole dyadic sums. It also checks that the symbolic table above covers each of the492 \nolinkurl{(t,k)} possibilities in JT2 exactly once. That finite exhaustion diagnostic is not an enumeration of the global rational matrices; the structural reduction and each analytic bound are the separately checked proofs cited above.

No additional error charge is introduced by this corollary. It joins already charged, exhaustive statements.

\subsection{4. What has strictly changed and what is still required}\label{what-has-strictly-changed-and-what-is-still-required}

The all-top finite nonvanishing gate is now the actual signed sum over supports at least12. At review18, the support11 scalar-cross profiles with t0,t1,t2 were still retained. Thus JT1 is a strict reduction of that main gate, independent of bookkeeping, improved constants within an unchanged sector, or packaging.

The still sufficient condition is

\[
 \operatorname{Re}\sum_{h\in R_{\rm top}^{(12)}}w(h)
 >-1+2^{-230}.
 \tag{JT5}
\]

No sign or nonzero value for this remaining sum is proved. The prescribed Bessel/Mram period, the concrete032 cycle and its MainX normalization/value, the full literal mixed source-family ties and every actual term in the DP44 lift-cost difference remain at their stated gates. In particular there is no shorter admissible chain, no universal strict negative DP44 theorem, and no Kummer Hecke RH consequence.

\subsection{5. Geometric interpretation and exact carried locators}\label{geometric-interpretation-and-exact-carried-locators}

The rational isometries act on the same Hermitian module over \nolinkurl{Spec(O_E)}. The ideals and congruence conditions are its actual integral submodules; the coordinate count and its Gaussian auxiliary estimate are archimedean counting tools for their rational points. This is the previously indexed SGA3 I Proposition4.7.3/Corollary4.7.3.1 module-and-character framework, with the actual coefficient/flag class in SGA6 VI Theorem3.3. No new equality of lambda-ring or K-theory classes is used to infer nonvanishing.

The polarized CM Jacobian and Picard line remain those of FGA232 Theorem3.1, pp.148--149; FGA236 Corollaries1.7/2.3, pp.229/232; and SGA7 I ExposéVII2.7/2.9, pp.169--174. The formal-group/Cartier-duality dictionary at Grothendieck cote4 scan003 sections8.1.1--8.1.2, scans005--006 sections8.4.1--8.4.3, and cote54 scans005--007 is unchanged. These are carried locators already checked and qualified in \nolinkurl{GROTHENDIECK_DICTIONARY.md}, not claims that Grothendieck proves the present numerical bounds. No new l-adic weight, motive, cycle realization or motivic L-function theorem is supplied. The changed conditional rows are DP\_MULT Q3e and THETA3 G1/A7; CX\_GRH CX-ALG-024 and GCT\_PLETH T-052 keep their separate unresolved comparisons.

All proof sources used above are current canonical files, with new recovery identities where stated. The old zero-target statements are accompanied by explicit refutation/correction events, and unavailable historical final hashes are not attributed to reconstructed files.

\clearpage

\section{AM-LIFT-064 --- global unitary rank-two invariants do not force the finite Green sign}\label{am-lift-064-global-unitary-rank-two-invariants-do-not-force-the-finite-green-sign}

\textbf{Canonical checked scope --- 2026-11-22T16:21:13.989274+00:00.} Root read the complete GP1--8 proof and complete script, directly read the primary character formula, and independently replayed6677 exact diagnostics. This refutes only positivity inferred from the listed global invariants; it does not refute positivity of the actual constrained short remainder.

\textbf{Status:} new scoped proof in the controlled pending directory. It addresses an intrinsic-parity route only. It is not an example in the actual nonintegral, support-eleven, all-type-II short sector, and it does not determine that sector's sign or the signed remainder.

\subsection{1. The exact sign in the current finite model}\label{the-exact-sign-in-the-current-finite-model}

Retain the actual J56 field E/F, the two odd rational G primes split completely in E, and the fixed finite representation and projector from K058(PG3). Among its 336 relative G places there are 334 full-degree cuspidal factors and two mixed factors. At a full-degree factor of GL\_n(F\_q), with n=169, let h be unipotent and d=rank(h−I). Green's formula gives \[
 \frac{\chi(h)}{\chi(1)}
 =\frac{(-1)^d}{\prod_{j=n-d}^{n-1}(q^j-1)}.
 \tag{GP1}
\] The empty product is one. For d=0,1,2 these signs are respectively positive, negative, positive. The parameter character contributes its value at one, so it contributes no further phase. Taking the contragredient representation does not change these real values.

This is obtained directly from Gorodetsky and Hazan (2017), Theorem 2.1, equation (2.1), printed p.~5: the number of Jordan blocks is n−d and the product has n−d−1 negative factors. Thus the product of the 334 unmixed signs is \[
 (-1)^{\#\{v:\operatorname{rank}(h_v-I)=1\}}.
 \tag{GP2}
\] Conjugate CM primes have equal ranks by the inverse-adjoint unitary relation. However, one \emph{relative} place supplies one finite factor in PG3. Counting both E-primes would double the count and does not compute (GP2).

\subsection{2. An integral original-coordinate unitary family}\label{an-integral-original-coordinate-unitary-family}

Inside the actual W0 of signature (4,1), use two original positive coordinate vectors and the original negative coordinate vector. On these three coordinates the Hermitian matrix is J=diag(1,1,−1). For x in 2O\_E, put A=x bar(x) and \[
 h(x)=
 \begin{pmatrix}
 1-A/2&-\bar x&A/2\\
 x&1&-x\\
 -A/2&-\bar x&1+A/2
 \end{pmatrix}.
 \tag{GP3}
\] Extend it by the identity on all other original coordinates. Every entry is integral. Direct multiplication gives \[
 h(x)^*Jh(x)=J,\qquad\det h(x)=1.
 \tag{GP4}
\] For N=h(x)−I, \[
 N^2=A
 \begin{pmatrix}-1&0&1\\0&0&0\\-1&0&1\end{pmatrix},
 \qquad N^3=0.
 \tag{GP5}
\] If x is nonzero, A is nonzero, rank N=2 and rank N²=1. This is also the standard Heisenberg unipotent in the rational hyperbolic basis \(e_1+e_n,e_2,(e_1-e_n)/2\), transported back to the unchanged standard coordinates. The formula, rather than that basis, proves integrality; no unbounded inverse-frame denominator is discarded.

It fixes the entire original positive S complement pointwise. Its determinant is exactly one and its only eigenvalue is one. Its literal witness is \[
 \alpha=(N^2)_{nn}=A=x\bar x>0
 \tag{GP6}
\] at every real embedding. Thus it has the global unitary, positive-alpha and rank-two properties used in the current remainder analysis.

\subsection{3. A prescribed odd residue-rank pattern in the actual field}\label{a-prescribed-odd-residue-rank-pattern-in-the-actual-field}

Choose any one of the 334 unmixed relative places v\emph{. Let P} be its selected E-prime and bar(P*) its CM partner. All E-primes over the two G primes are distinct, of odd residue characteristic. Apply the Chinese remainder theorem in O\_E to choose y with these residues:

\begin{itemize}
\tightlist
\item
  y is zero at P* and is one at bar(P*);
\item
  y is one at both partners of every other unmixed relative place;
\item
  y is zero at both partners of each of the two mixed places.
\end{itemize}

Set x=2y. Multiplication by two preserves zero versus unit at all these places. The nonzero residue at bar(P*) ensures x is nonzero. At a selected split prime, regard x and bar(x) as independent residue coordinates a,b. Reducing (GP3) gives \[
 \operatorname{rank}(h-I)=
 \begin{cases}
 0,&a=b=0,\\
 1,&\text{exactly one of }a,b\text{ is nonzero},\\
 2,&ab\ne0.
 \end{cases}
 \tag{GP7}
\] The last case follows from the nonzero rank-one square in (GP5); the middle case follows immediately from the sole nonzero row or column direction. The same ranks occur at the conjugate partner, where a and b are exchanged.

Consequently this single global integral h has exactly one rank-one unmixed relative reduction, rank two at the other 333 unmixed places, and identity at both mixed places. At a mixed identity the normalized partial trace is \(\operatorname{Tr}(Q_{\rm good})/\dim\operatorname{im}Q_{\rm good}=1\), independent of the mixed constituent. The complete normalized finite factor in K062(AT14) is therefore strictly negative: \[
 a(h)=-\frac1{q_{v_*}^{168}-1}
       \prod_{\substack{v\text{ unmixed}\\v\ne v_*}}
       \frac1{(q_v^{167}-1)(q_v^{168}-1)}<0.
 \tag{GP8}
\] There is no scalar-normalization phase to account for: h is already unipotent with determinant one and the other factors are identities.

The alpha-ideal condition does not repair the parity. If desired, prescribe y modulo P\emph{² so that its P} valuation is exactly one while keeping its partner a unit. Then the base-field valuation of alpha at v* is exactly one. The E-ideal of alpha is conjugation invariant, with equal valuations at both partners, and its absolute E-norm is a square because alpha lies in F. These facts are compatible with an odd number of selected relative rank-one places. They do not assert that its F-norm or every individual relative valuation is even.

\subsection{4. What this proves, and the exact remaining possibility}\label{what-this-proves-and-the-exact-remaining-possibility}

The proposed parity conclusion cannot follow from global unipotence, rank N=2, the original unitary form, fixing S, determinant one, CM-conjugate rank equality, or alpha positivity and paired ideal divisibility alone: (GP3)--(GP8) meet all those conditions and give a negative full finite factor. The two mixed factors are controlled explicitly here, rather than assumed to have a sign.

This example is globally integral. No bound \textbar\textbar h\_tau\textbar\textbar≤133 is proved; no required nonintegral support of size eleven is imposed; and no membership in the unique fixed incidence of each actual T pencil is asserted. Therefore it does \textbf{not} refute a parity theorem that uses those additional literal constraints. Such a theorem would need a new connection between the fixed T incidences and the selected G valuations, beyond the invariant/alpha conditions already used. Nor does a sign for a(h) alone determine the full term a(h)c(h) times the T incidence product: the actual real selected compression remains part of that term.

\subsection{5. Sources and verification boundary}\label{sources-and-verification-boundary}

Gorodetsky, O., \& Hazan, Z. (2017). \emph{On certain degenerate Whittaker models for cuspidal representations of GL\_(kn)(F\_q).} arXiv:1707.07308. Theorem 2.1 and equation (2.1), printed p.~5. \url{https://arxiv.org/pdf/1707.07308} . The primary research-paper passage was freshly opened and read; it quotes Green's exact theorem. The unipotent specialization and all subsequent matrix/CRT arguments are proved above.

The actual finite model is K058, \nolinkurl{projected_selected_packet_gram.md}, Section 1 / PG3, and K062, \nolinkurl{averaged_scalar_pencil_trace.md}, AT14--16. Their identities retain one factor per relative place and the two mixed partial projectors. The companion diagnostic checks the symbolic matrix identities, every split residue pair over several odd finite fields, and the exact character-sign formula. It does not enumerate the actual short rational support or create an all-type-II term.

\clearpage

\section{AM-DO-073 --- the actual scalar-cross isometry fibre has rank five}\label{am-do-073-the-actual-scalar-cross-isometry-fibre-has-rank-five}

\textbf{Canonical checked version --- 2026-11-22T16:24:31.785915+00:00.} The complete RF1--19 proof and full script have separate root and independent peer whole-proof checks; root replayed43 exact assertions. The final source has a new canonical status heading only. The numerical consequence assumes global rank-two unipotence and removes only its stated support-twelve profiles.

Author: \nolinkurl{/root/definite_oscillator}, 22 November 2026 UTC.

This is a new refinement. It does not change the frozen DO072 proof, which deliberately retained the larger fibre bound. The pointwise fixed subspace below is determined by the actual original S-part of the negative column. Its rational orthogonal complement need not have a globally integral orthogonal basis. The group-order estimate uses its actual invariant intersection lattice locally, with no frame-height assertion.

\subsection{1. Fixed data and the new fibre statement}\label{fixed-data-and-the-new-fibre-statement}

Retain the actual J56 field and lattice-pencil kernel, the original Hermitian form of signature \nolinkurl{(168,1)} at every real label, and the original orthogonal decomposition \[
E^{169}=S\perp W_0,\qquad \dim_E S=164,\quad
\operatorname{sig}(W_0)=(4,1).
\tag{RF1}
\] Fix a complete tuple of actual ordered source-target pencil labels and a full negative column \nolinkurl{v=h e_n}. Consider only the matrices in that fibre satisfying \[
h=I+N,\quad h_{SS}=I,\quad
\operatorname{rank}N=2,\quad N^3=0,\quad
\operatorname{rank}N^2=1,\quad
x=(N e_n)_S\ne0.
\tag{RF2}
\] These are genuine scalar-cross globally unipotent terms. Their number is at most \[
G_5=\gcd\bigl(|\operatorname{GL}_5(\mathbf F_{g_0})|,
              |\operatorname{GL}_5(\mathbf F_{g_1})|\bigr)
=98148490017978252769652736000000<2^{107}.
\tag{RF3}
\] The actual primes are \nolinkurl{g0=22285932761} and \nolinkurl{g1=64072175441}. The estimate is uniform in the negative column, class/pencil labels, denominator support and their heights. It concerns the finite source-lattice isometry fibre, not an integral inverse of a theta map.

\subsection{2. The fixed positive subspace is determined by the column}\label{the-fixed-positive-subspace-is-determined-by-the-column}

The accepted signature-one nilpotent identities give \[
N+N^\dagger=N^2,\qquad
\ell:=\operatorname{im}N^2\text{ isotropic},\qquad
\ker N^2=\ell^\perp.
\tag{RF4}
\] Since \nolinkurl{N_SS=0}, the SS block of the nonzero selfadjoint rank-one operator N² is zero. Its isotropic image line consequently lies in W0. The negative vector \nolinkurl{e_n} is not perpendicular to that line, so \nolinkurl{N²e_n≠0}. The two vectors \nolinkurl{N e_n} and \nolinkurl{N² e_n} are independent and span \nolinkurl{im N}. Because their S parts are x and zero respectively, and x is nonzero, \[
\operatorname{im}N\cap W_0=\ell.
\tag{RF5}
\]

For \nolinkurl{u∈S}, the zero SS block implies \nolinkurl{Nu∈W0}, hence \nolinkurl{Nu∈ell}. Use the Hermitian pairing and RF4: \[
\langle Nu,e_n\rangle
=\langle u,N^\dagger e_n\rangle
=\langle u,(N^2-N)e_n\rangle
=-\langle u,x\rangle.
\tag{RF6}
\] The pairing of ell with \nolinkurl{e_n} is nonzero. Therefore \[
Nu=0\quad\Longleftrightarrow\quad\langle u,x\rangle=0.
\tag{RF7}
\] Every matrix in the specified fibre fixes pointwise the same rational positive subspace \[
U_x=x^\perp\cap S,\qquad \dim_E U_x=163.
\tag{RF8}
\] This uses the original coordinate vector x. It does not choose a primitive representative of its line, clear its denominators, or rotate it into a coordinate vector.

\subsection{3. From the actual fibre to a finite rank-five automorphism group}\label{from-the-actual-fibre-to-a-finite-rank-five-automorphism-group}

If the fibre is nonempty, fix one member \nolinkurl{h0}. For any other member h, the ratio \texttt{a=h0\^{}\{-1\}h} fixes \nolinkurl{e_n}, because both have the same full negative column. Both h and h0 fix Ux pointwise by RF8, so their ratio also fixes Ux pointwise. It preserves the actual source lattice: both original matrices carry that fixed source lattice to the same target lattice. Thus the assignment \texttt{h↦h0\^{}\{-1\}h} is injective into its isometry stabilizer fixing both Ux and \nolinkurl{e_n}.

The subspace \nolinkurl{U_x⊕Ee_n} is nondegenerate, of signature \nolinkurl{(163,1)} at every real label. Its rational E-orthogonal complement \[
V_x=(U_x\oplus Ee_n)^\perp
\tag{RF9}
\] has dimension five and is positive definite at every real label. The stabilizer acts faithfully on Vx because it already acts trivially on the orthogonal complement of Vx.

Let \nolinkurl{M_c} be the actual global source lattice obtained by patching the fixed finite pencil lattices. The intersection \[
\mathcal L_x=M_c\cap V_x
\tag{RF10}
\] is a full projective O\_E-lattice in Vx, preserved by the stabilizer. It is full because a basis of the rational E-subspace can be multiplied by common nonzero denominators to lie in M\_c.~The full stabilizer of this positive lattice is finite: its real group is compact at all real labels, while its integral endomorphism lattice is discrete. Restriction therefore gives a finite subgroup of \nolinkurl{GL_E(V_x)} preserving \nolinkurl{L_x}.

At a selected G prime over either g0 or g1, the underlying source lattice is the original maximal lattice. The prime splits completely in E, so its local field is \texttt{Q\_\{g\_j\}}. The intersection \texttt{L\_x⊗O\_\{E,v\}} is a torsion-free full lattice in a five-dimensional local vector space, hence is free of rank five over the DVR \texttt{Z\_\{g\_j\}}. Choosing a basis only for this local module realizes the finite group in \texttt{GL\_5(Z\_\{g\_j\})}. No such basis is asserted to come from a globally integral orthogonal frame; no coefficient-height bound follows or is needed.

\subsection{4. Torsion-free local reduction and the exact order}\label{torsion-free-local-reduction-and-the-exact-order}

Serre (2006, §1.2, Lemma1′, printed p.5) proves that the kernel of reduction from \nolinkurl{GL_d(R)} over a local ring to its residue field contains no element of prime order different from the residue characteristic. Apply that primary result to \texttt{R=Z\_\{g\_j\}} and d=5. The only possible torsion prime in the kernel is g\_j.

For the actual primes, \nolinkurl{g_j−1>5}. A nonidentity element of order g\_j over \texttt{Q\_\{g\_j\}} would have a primitive g\_j-th root as an eigenvalue, whose minimal polynomial is \texttt{Phi\_\{g\_j\}} of degree \nolinkurl{g_j−1}: its translate \texttt{Phi\_\{g\_j\}(X+1)} is Eisenstein. Such an eigenvalue cannot occur in dimension five. Thus the local reduction kernel is torsion-free in this actual rank and field. In particular it intersects the finite stabilizer trivially. This verifies the needed residue-characteristic case, rather than silently extending the prime-to-characteristic statement of the cited lemma.

Consequently the same finite group injects into \texttt{GL\_5(F\_\{g0\})} and \texttt{GL\_5(F\_\{g1\})}. Lagrange's theorem gives RF3, with \[
|\operatorname{GL}_5(\mathbf F_q)|
=\prod_{j=0}^{4}(q^5-q^j).
\tag{RF11}
\] The exact gcd and its bit bound are replayed from the actual prime certificate. The complete-splitting conditions are checked against the actual field constituents; primality itself remains the already proved actual-host input. This finishes the new fibre theorem.

\subsection{5. A separate support-twelve consequence}\label{a-separate-support-twelve-consequence}

Here specialize to the \textbf{all-top} J56 cutoff \nolinkurl{H=133}, and retain global unipotence as an explicit hypothesis. Let s=12, i=0, u+t=12, and \nolinkurl{4≤t≤12}. Retain RF2, all original finite/projector data, class weights, the actual real selected-module metric and packet point. This section removes only these globally unipotent scalar-cross subsets; it does not supply global unipotence for other terms or prove a complete support-twelve classification.

The corrected DO070/DO072 lattice argument and AM-DO-070-C1 give the full-m sufficient envelope \[
\begin{split}
U_{s,t}<&\;
2^{1+5379+1062+107+168+2s+668+t}
12^f2^{f/2-959f/4}Q^{-335(f-1)}\\
&\cdot T^{343s-7t/2}R^{669f/4}
V_m\!\left(1+c\sqrt m\,T^{-t/(2f)}\right),
\end{split}
\tag{RF12}
\] where \nolinkurl{f=168,m=164,R=17424/133,c=2250/7}, and \[
V_m(C)=\frac{((22/7)fC^2)^{fm}}{(fm)!D_0^m}.
\tag{RF13}
\] This is the rational upper-volume envelope, not the true pi-volume. The fibre5379 and gamma1062 are the conservative constants of DO072. RF3 is the only isometry-fibre replacement. The factor \nolinkurl{2^t} is the corrected complete zero-target type-III averaging. All common-orientation patterns, support/type multiplicities, full dyadic gamma ideal and actual homogeneous-energy real factor remain present exactly once. The all-nonzero Gaussian estimate is not needed for RF12: the full S ball is a sufficient upper count here.

For clarity, the maximization underlying RF12 retains the actual alpha constraint. Set \texttt{y=(NJ\_F\ T\^{}\{-(2s-t)\})\^{}\{1/f\}≤R} and \texttt{z²=yT\^{}\{-t/f\}}. The S-ball count contributes \texttt{z\^{}\{-2fm\}V\_m(sqrt(R)+c\ sqrt(m)z)}. Combining its y power with Green, real compression and gamma count leaves \texttt{y\^{}\{9f/4\}(R/2+375y)\^{}f}, times that increasing volume function. This entire expression is increasing in y, so evaluating at R is an upper bound. The resulting R power is \nolinkurl{fm+13f/4=669f/4}; the T power is \nolinkurl{8s+(335/2)(2s-t)+mt=343s-7t/2}. Thus no ideal-dependent factor is lost in passing to RF12.

For the actual prime T and constants, exact integer inequalities prove \[
c\sqrt{164}<4117,\qquad T^{1/84}>3690,
\quad 1+4117/3690=7807/3690.
\tag{RF14}
\] Thus at t=4 the volume argument in RF12 is below \nolinkurl{7807/3690}. Exact rational comparisons give \[
2^{-320029}<V_m(7807/3690)<2^{-320028}.
\tag{RF15}
\] Use the sharper actual-prime rational logarithm bounds, certified by integer powers, \[
\log_2 T<995+1812/4096,
\quad\log_2 Q>70+1123/4096,
\quad\log_2 R<3602/512.
\tag{RF16}
\] Together with \nolinkurl{12<16}, RF12 at \nolinkurl{(s,t)=(12,4)} has exponent \[
\begin{split}
E={}&1+5379+1062+107+168+24+668+4\\
&+(343\cdot12-14)(995+1812/4096)
 +(669\cdot168/4)(3602/512)\\
&-335\cdot167(70+1123/4096)
 +168(9/2-959/4)-320028\\
={}&-10841955/4096<-2646.
\end{split}
\tag{RF17}
\]

For fixed s, RF12 decreases strictly with t: its explicit ratio has factor \texttt{2\ T\^{}\{-7/2\}\textless{}1}, and its volume radius also decreases. Hence every one of the nine t profiles from4 through12 is bounded by \texttt{2\^{}\{-2646\}}. Taking their union conservatively gives \[
\sum_{t=4}^{12}U_{12,t}<9\,2^{-2646}<2^{-2642}.
\tag{RF18}
\] This charge applies only to the intersection with the current remainder, retaining every earlier error. Existing DO072 and its charge are unchanged. No favorable sign, full signed-sum bound, Mram period, Host constructor, or DP44/RH conclusion follows.

\subsection{6. The separate global-ideal-cap diagnostic}\label{the-separate-global-ideal-cap-diagnostic}

The common-orientation ideal always satisfies \texttt{NJ\_E≤Q\^{}\{2(f−1)\}}. In DO072's variables this gives the additional exact restriction \[
p\ge\frac{R T^{2s/f}}{kQ^{2(f-1)/f}}.
\tag{RF19}
\] It was tested first as a possible way to extend the old fibre5022 estimate. The companion \nolinkurl{global_ideal_cap_diagnostic.py/.json} records that, for the actual support12,t0,k164 cell, this floor lies strictly between4 and5 by exact integer comparisons. Its subsequent floating three-endpoint diagnostic still gives a positive sufficient upper envelope of roughly333000 bits. Those floating values are explicitly not certified inequalities; a positive upper estimate proves no positive contribution, obstruction or period vanishing. RF12--18 use the new proved fibre theorem instead, not an unproved extrapolation from that diagnostic.

\subsection{Primary reference and retained arithmetic interfaces}\label{primary-reference-and-retained-arithmetic-interfaces}

Serre, J.-P. (2006). \emph{Bounds for the orders of the finite subgroups of G(k).} In M. Geck, D. Testerman, \& J. Thévenaz (Eds.), \emph{Group representation theory}. EPFL Press. Author lecture text, §1.2, Lemma1′, printed p.5, directly read at \url{https://www.college-de-france.fr/media/jean-pierre-serre/UPL3821667391778701726_6___Bounds_for_the_orders.pdf} . The cited result excludes prime-to-residue-characteristic torsion; the actual rank-five residue-characteristic exclusion is verified separately in §4.

Actual kernel/lattice inputs: DO063(FV20), DO066's common-orientation ideals, the corrected DO070/071 slice and AM-DO-070-C1, DO072(GS4--9,GS17--23), K077(AP1--3), and the actual DO040/K032 prime/field certificates. The new result is the fixed Ux and faithful rank-five fibre, with the explicit local lattice realization. No globally integral moving basis, fixed-W0 support assumption, or new finite projector is imported.

\subsection{Author verification and scope of the certificate}\label{author-verification-and-scope-of-the-certificate}

The complete six-section proof, reference paragraph, and complete companion script were read by the author after the numerical replay. The exact certificate passed all43 assertions at2026-11-22T16:20:17.682033+00:00; the independent cap diagnostic was run at16:20:17.260743UTC. Serre's author PDF was opened directly again and §1.2, Lemma1′ and its proof were read at this final check. The certificate verifies the actual splitting congruences, the gcd of finite group orders, exact volumes/logarithm inequalities, every one of the nine profile exponents and the stated failed-envelope p-floor. It does not claim automated verification of the fixed-space theorem, local lattice/reduction proof, or the retained automorphic interfaces. Final file hashes and the final whole-read event are recorded separately in \nolinkurl{scalar_cross_rank_five_fibre_verification.json}. This controlled-recovery result is ready for independent review; it is not a main-ledger promotion.

\clearpage

\section{AM-K086 --- type-III incidences make the actual scalar-cross column integral}\label{am-k086-type-iii-incidences-make-the-actual-scalar-cross-column-integral}

\textbf{Canonical checked version --- 2026-11-22T16:33:21.651714+00:00.} Root read the whole final proof and full script and separately replayed the exact diagnostic. Original pending/status wording below is retained as export history; only this stated scope is promoted: Actual type-III paired S-cross integrality; only typeII S poles; all-top support12 globally rank-two-unipotent scalar-cross non-allII profiles have joint charge885. Full typeI/III zero-target costs retained. All-II and other structural classes remain.

Author: \nolinkurl{/root/kummer_extension}, 22 November 2026 UTC.

This is a new controlled-recovery proof. It concerns the unchanged actual J56 signed kernel, its original lattice pencils and scalar-cross global rank-two unipotents. It does not change a source lattice, finite projector, real selected module, central representative or period normalization. The local theorem is independent of the all-top numerical specialization. The quantitative theorem removes the specified support-twelve subsets only; it does not supply global unipotence or exhaustive support-twelve structural reduction.

\subsection{1. Statements and notation}\label{statements-and-notation}

Use the original orthogonal decomposition and actual standard coordinates \[
 E^{169}=S\perp W_0,\qquad \dim_E S=164,
 \quad \operatorname{sig}(W_0)=(4,1),
 \tag{TI1}
\] where the actual finite pencils change only W0. Suppose \[
 h=I+N,\qquad h_{SS}=I,
 \qquad \operatorname{rank}N=2,\quad N^3=0,
 \quad \operatorname{rank}N^2=1,
 \quad x=(h e_n)_S\ne0.
 \tag{TI2}
\] These are precisely the scalar-cross globally rank-two-unipotent hypotheses already used in DO072 and DO073. The actual unitary adjoint is denoted by a dagger. The accepted identities are \[
 N^\dagger=N^2-N,\qquad
 N^2S=0,\qquad \operatorname{im}N^2\subset W_0.
 \tag{TI3}
\] The last two assertions hold in the original coordinates. For clarity, the SS block of the self-adjoint rank-one operator N² is zero; hence its isotropic image line lies in W0, and its kernel contains S. This is DO072(GS4), independently checked in DO073(RF4--8), not a statement about a divided integral isotropic vector.

\textbf{Local conclusion.} At every actual nonintegral type-III T place, both S/W0 cross blocks of N are integral at both CM components. Type I already has this property by K071(NB5). Consequently the only possible T poles of x occur at type II, and those have order at most one.

\textbf{Numerical conclusion.} At the actual all-top J56 cutoff H=133, the absolute contribution of the union \[
 s=12,\qquad i+u+t=12,\qquad i+t\ge1,
 \quad\text{subject to TI2}
 \tag{TI4}
\] is less than \nolinkurl{2^-885}. Here i,u,t count the actual types I,II,III. This union has 90 type profiles. Every original-coordinate support, G-orientation pattern, source/target pair and class factor is included. The result applies to the intersection with the current remainder; all previously charged errors are retained. The all-II profile \nolinkurl{(i,u,t)=(0,12,0)} is excluded from this conclusion.

\subsection{2. The original-coordinate intersection argument}\label{the-original-coordinate-intersection-argument}

At a selected split T component let O be the DVR, pi its uniformizer, and retain the literal pencil of LIFT44(PR4--7), K070(PT1--3), K071(NB2--3), and K077(AP1--3): \[
 M_c=(I+c\pi^{-1}Z)\Lambda,\qquad
 Z=zq,\quad Z^2=0,
 \qquad c\in k_T\setminus\{-1\}.
 \tag{TI5}
\] Zero is an admitted label and gives the standard lattice. Z has rows and columns supported in W0. Put \[
 \Lambda_- =\{v\in\Lambda:qv\in\pi O\},\qquad
 \Lambda_+ =\Lambda+\pi^{-1}Oz.
\] The full pencil is the set of intermediate lines in the two-dimensional quotient \nolinkurl{Lambda_+/Lambda_-}. Thus for distinct labels \[
 M_c\cap M_d=\Lambda_-\subset\Lambda,
 \qquad M_c+M_d=\Lambda_+\supset\Lambda.
 \tag{TI6}
\] Every original integral vector of S belongs to every M\_c, since Z annihilates S. These assertions hold also when c or d is zero.

At a type-III incidence the actual relation is \[
                         hM_c=M_d,\qquad c\ne d.
 \tag{TI7}
\] For any \nolinkurl{v in S_O}, both v and its negative belong to both pencil lattices. Therefore \[
 Nv=hv-v\in M_d,
 \qquad
 (h^{-1}-I)v=-Nv\in M_c,
\] where the second equality uses \texttt{h\^{}\{-1\}=I-N+N²} and \nolinkurl{N²v=0}. Consequently \[
                     N(S_O)\subset M_c\cap M_d\subset\Lambda.
 \tag{TI8}
\] This proves integrality of the S columns without assuming that h fixes either actual lattice.

At the conjugate CM component the source and target are the paired Hermitian dual lattices. They remain distinct. Their intersection is the dual of the selected-component sum and hence is contained in the standard dual lattice by TI6. The pencil again changes only W0, so the same original vectors of S belong to both paired lattices. Repeating TI8 proves S-column integrality at this second component as well.

Finally \nolinkurl{N^dagger v=-Nv} for v in S by TI3. The local matrix of a dagger is the signed conjugate transpose from the opposite CM component. Integral S columns of N at that opposite component therefore imply integral S rows of N at the present component. The signs from the fixed diagonal Hermitian form are units. This proves both cross-block assertions at both components, with no identification of the two local matrices and no additional count of CM partners.

Type-I cross-block integrality and the type-II simple-pole bound are precisely K071(NB5). At every inactive finite place the original matrix is integral. Thus, writing \nolinkurl{I_II} for the inverse-prime E-ideal containing both CM primes at each of the u type-II places, \[
                 x\in I_{\rm II}^{\,m},\qquad
                 N_E I_{\rm II}=T^{-2u},\quad m=164.
 \tag{TI9}
\] This strengthens the earlier safe denominator ideal supported on II and III. It does not change any earlier theorem that used the larger ideal.

\subsection{3. The actual G ideal and the new relation between its sizes}\label{the-actual-g-ideal-and-the-new-relation-between-its-sizes}

For the numerical estimate now specialize to the stated all-top H=133 data. Retain DO066 and DO072(GS5--6): for a common G-orientation pattern there are ideals J\_E and J\_F with equal absolute norm, at most \nolinkurl{4^334} such patterns, and \[
 |a(h)|\le2Q^{-335(f-1)}(NJ_E)^{335/2},
 \qquad NJ_F\le N_F(\alpha)T^{2u+t},
 \tag{TI10}
\] where \nolinkurl{f=168}, \nolinkurl{Q=g0*g1}, and \nolinkurl{alpha=(N²)nn} satisfies \[
 0<\alpha_\tau\le R=17424/133,
 \qquad \|x_\tau\|^2\le\alpha_\tau.
 \tag{TI11}
\] The two mixed finite factors retain their modulus bound one. No sign inference is used. Type-I alpha is integral, type-II alpha has pole at most two, and type-III alpha has pole at most one. These remain the actual square-pole statements, distinct from the stronger S-column assertion TI9.

Combine TI9 with the G divisibilities, which are at primes disjoint from T. Set \[
 L=J_E I_{\rm II},\quad N_E L=z^{2f},\qquad
 y=(NJ_FT^{-(2u+t)})^{1/f},\quad w=T^{t/f}.
\] Then \[
                  x\in L^m,\qquad z^2=yw,
                  \qquad 0<y\le R.
 \tag{TI12}
\] The factor w now multiplies y. This is the arithmetic improvement; all real metrics and covolumes are unchanged.

Let k be the number of nonzero original E-coordinates of x. It lies between 1 and164. For every nonzero b in L, the norm inequality \nolinkurl{(b) subset L} and arithmetic-geometric mean in the half-trace metric give \nolinkurl{||b||² >= f z²}. Hence \[
 kz^2\le R,\qquad p=R/(kz^2)\ge1,
 \qquad y=\frac{R}{kpw}.
 \tag{TI13}
\] The alpha constraint y\textless=R follows automatically from k,p,w\textgreater=1. No primitive-coordinate assumption is made. There are \nolinkurl{binom(164,k)} supports, retained below.

\subsection{4. Real, lattice, and finite normalizations, including type I}\label{real-lattice-and-finite-normalizations-including-type-i}

Use the actual half-trace norm and the unchanged ideals and discriminants \[
 \operatorname{covol}(O_E)=D_0=7^{140}29^{144}5^{84}13^{84},
 \qquad D_F=2^fD_0,
 \qquad L^\#=2\mathfrak D_E^{-1}\bar L^{-1}.
\] The accepted arbitrary-ideal transference bound is \[
                  \mu(L^k)<c\sqrt{fk}\,z,
                  \qquad c=2250/7.
 \tag{TI14}
\] It is distinct from the original-ring-only affine-fibre cover of LIFT063. In particular no fractional ideal is assigned the small original-ring covering constant.

For each fixed actual source/target tuple, the W0 negative-column over-lattice and real-part coarea are exactly those of corrected DO072(GS17--21) and AM-DO-070-C1. The global affine-fibre constant is \nolinkurl{2^5379}. Each nonzero target label contributes the local denominator factor \nolinkurl{T^8}; a zero target contributes one further T, because its negative-coordinate projection has pole one instead of pole two. The resulting global bound is \nolinkurl{2^5379 T^(8s+r0)}, where r0 counts zero targets. The gamma ideal, with its dyadic normalization retained, is \[
 \gamma=\operatorname{Re}(h_{nn})-1=\alpha/2
 \in\tfrac12J_F I_{{\rm II},F}^2 I_{{\rm III},F}.
\] Its count is bounded by \[
 A_f(R/2+375y)^f y^{-f},\qquad
 A_f=\frac{\pi^{f/2}f^{f/2}2^f}{(f/2)!\sqrt{D_F}},
 \qquad A_f376^f<2^{1062}.
 \tag{TI15}
\] The conservative fibre 5379 is used, although the exact original-ring coarea permits 5211. The actual full-column source-lattice isometry fibre is DO073(RF1--11), \[
                    G_5<2^{107}.
 \tag{TI16}
\] It is valid for scalar-cross terms: fixed full column determines the common fixed positive 163-space, and the remaining positive five-dimensional intersection lattice controls its stabilizer. No moving integral orthogonal frame is asserted.

The actual all-top selected degree exceeds \nolinkurl{2^959}. K085's full homogeneous-degree, orthogonal-ladder and near-point comparison gives exactly the retained estimate \[
             |c(h)|\le12^f2^{f/2-959f/4}y^{-f/4}.
 \tag{TI17}
\] All these are in the source Hilbert and original normalized-kernel conventions.

Here is the complete pair and zero-target accounting for general i,u,t. Put \nolinkurl{b_T=T-1}. In type I the unchanged K071(NB7) bound is \nolinkurl{|rho|<=1/b_T}; over all ordered actual pairs there are at most \nolinkurl{b_T(b_T-1)} nonzero-target pairs and \nolinkurl{b_T} zero-target pairs. Their complete weighted affine-fibre bound relative to \nolinkurl{T^8} is \[
 \frac{b_T(b_T-1)+b_TT}{b_T}=b_T-1+T=2b_T<2T.
 \tag{TI18}
\] This deliberately counts every possible pair, whether or not a particular matrix realizes it. It is a valid uniform upper bound on the corresponding tuple-wise matrix counts. It does not assert identical pair summands. For completeness, the precise type-I bound follows from the actual rho formula K070(PT2). The induced action on the two-dimensional quotient is unipotent. It cannot be projectively trivial, since that would fix the included standard lattice and make h integral. Its projective action therefore has order T and exactly one fixed point. Write E for the two omitted points. The case \nolinkurl{hE=E} is impossible: an order-T permutation of a two-element set is trivial, which would give two fixed points. Thus \nolinkurl{C=|D intersect hD|} lies between \nolinkurl{b_T-2} and \nolinkurl{b_T-1}, and the number F of fixed admitted points is zero or one. If F=0, \nolinkurl{|rho|=C/[b_T(b_T-1)]<=1/b_T}. If F=1, \nolinkurl{rho=(b_T-C)/[b_T(b_T-1)]<=2/[b_T(b_T-1)]<=1/b_T}, since the actual T exceeds3. This proves the stated bound without a favorable-sign assumption.

For type II, a supported fixed target cannot be zero, since fixing the standard lattice would be integral. The complete diagonal-pair count times \nolinkurl{rho=1/b_T} is at most one. In type III, the actual weight is \nolinkurl{1/[b_T(b_T-1)]}; the nonzero/zero-target counts are \nolinkurl{(b_T-1)^2} and \nolinkurl{b_T-1}. Their weighted affine-fibre factor is exactly two by C1. Hence all local pairs and zero-target corrections together cost \[
                           2^{i+t}T^i.
 \tag{TI19}
\] The paired local lattices encode one relative place; no second independent factor is introduced at the conjugate prime. Retain also support factor \nolinkurl{2^168}, type factor \nolinkurl{3^s<2^(2s)}, common-orientation factor \nolinkurl{4^334=2^668}, the Green constant2 and the exact support choices \nolinkurl{binom(164,k)}.

\subsection{5. Two exact lattice alternatives and endpoint control}\label{two-exact-lattice-alternatives-and-endpoint-control}

DO072(GS10--16), using Miller and Stephens-Davidowitz (2019), Lemma3.11 and the closed-tail argument, gives for each nonzero-coordinate support the alternatives \[
 \#x\le e^{bkp}8192^k,
 \qquad
 \#x\le(kp)^{fk}V_k(1+c/\sqrt p),
 \quad b=673/4,
 \tag{TI20}
\] where \[
 V_k(C)=\frac{((22/7)fC^2)^{fk}}{(fk)!D_0^k}.
\] They count the same actual set and are alternatives, not multiplicative factors. The Gaussian bound includes points attaining the algebraic minimum. The ordinary factorial and real ball dimension2fk remain exact.

Put \nolinkurl{P=669f/4=28098} and \[
 H_{t,k}(p)=\left(\frac{375/w+kp/2}{376}\right)^f.
\] Combining TI10 and TI12--20 gives the common factor \[
\begin{split}
\mathcal C_{i,t,k}={}&
2^{1+5379+1062+107+168+2s+668+i+t}
12^f2^{f/2-959f/4}Q^{-335(f-1)}\\
&\quad\cdot T^{343s-334i-1335t/4}R^P\binom{164}{k}.
\end{split}
 \tag{TI21}
\] After it the two alternatives are \[
 G_{t,k}(p)=(kp)^{-P}H_{t,k}(p)e^{bkp}8192^k,
 \qquad
 V_{t,k}(p)=(kp)^{-(P-fk)}H_{t,k}(p)V_k(1+c/\sqrt p).
 \tag{TI22}
\] To check the T exponent directly, before gamma factoring the Green/real power is \texttt{T\^{}\{(335/2)(2u+t)\}\ y\^{}\{P\}} and the W0/pair factor is \texttt{T\^{}\{8s+i\}}. Substitute \nolinkurl{y=R/(kpw)}. The gamma factor \nolinkurl{A_f(375+R/(2y))^f} then contributes \nolinkurl{w^f=T^t} when rewritten using H. The total is \[
 (335/2)(2u+t)+8s+i-(669/4)t+t
 =343s-334i-1335t/4.
 \tag{TI23}
\] This tracks the gamma density once; no separate lost or doubled dyadic factor is hidden.

For every k\textless=164, the derivative of log V with respect to log p is strictly below \nolinkurl{-(P-fk)+f<0}. Also log G is convex as a function of log p: the relevant terms are a linear term, \nolinkurl{f log(375/w+(k/2)exp(v))}, and \nolinkurl{bk exp(v)}. Thus for every actual ideal pattern and p\textgreater=1, \[
 \min(G(p),V(p))\le\max\{G(1),G(7),V(7)\}.
 \tag{TI24}
\] No factor three is needed. Increasing i multiplies the full envelope by \nolinkurl{2T^-334<1}. Increasing t multiplies its common part by \nolinkurl{2T^(-1335/4)<1} and decreases H. Therefore, for fixed s and k, every profile with i+t\textgreater=1 is bounded by one of \nolinkurl{(i,t)=(1,0)} and \nolinkurl{(0,1)}.

\subsection{6. Exact certificate and the charged subset}\label{exact-certificate-and-the-charged-subset}

The companion exact script uses the actual T and G primes, the restored upper-gap integer for the selected degree, and the one-sided bounds \[
 \log_2 T<995+1812/4096,\quad
 \log_2 Q>70+1123/4096,\quad
 \log_2 R<3602/512,
\] \[
 \log_2 e<13/9,\quad T^{1/168}>60,
 \quad \sqrt7>21/8,
 \quad 1+c/\sqrt7<6049/49.
 \tag{TI25}
\] All are certified by integer or rational comparisons. The Gaussian constant 8192 follows from the rational lower exponential series proving \nolinkurl{(673/672)^168 exp(-1/4)<8192/8193}. The fibre 5379, gamma 1062 and G5 constants are replayed independently of the numerical endpoint loop.

For the base \nolinkurl{(i,t)=(1,0)}, use the exact gamma ratio \nolinkurl{(750+kp)/752}. For \nolinkurl{(0,1)}, the bound w\textgreater60 gives the ratio \nolinkurl{(25+2kp)/1504}. The exact binomial factor, inverse powers \nolinkurl{(kp)^P} and rational volume are included before an integral binary ceiling is taken. The Gaussian exponent is bounded by \nolinkurl{(13/9)(673/4)kp}. The full maxima over all k and all three sufficient endpoints are \[
 E_{1,0}=-33159035/36864<-899,
 \qquad E_{0,1}=-29945287/18432<-1624.
 \tag{TI26}
\] Both are attained by the sufficient bound at k=1,G(1); this is not a statement about an actual matrix attaining them.

There are \nolinkurl{(12+1)(12+2)/2-1=90} non-all-II profiles and164 coordinate-support-size values. Hence \[
 90\cdot164=14760<2^{14},
 \qquad
 \sum_{\text{TI4}}|\text{actual normalized contribution}|
 <14760\,2^{-899}<2^{-885}.
 \tag{TI27}
\] Every removed sector is intersected with the current remainder, with all earlier numerical errors retained conservatively. In particular DO073's already bounded i=0,t\textgreater=4 sectors may overlap this new bound; no double subtraction or new positive sign is inferred. A separate stronger bound on i=0,t\textgreater=1 is \nolinkurl{2^-1613}, but it is an alternative contained-sector statement, not a second charged error.

The all-II scalar-cross support-twelve sector remains outside this proof. Fixed-W0 terms, nonunipotent terms and nonscalar A are not counted here. Even if another proof forces some of those terms into TI2, it must supply that implication explicitly. The result does not establish full signed-sum positivity, an all-place Mram period, literal Host or SizeData admissibility, DP44, or RH.

\subsection{References and proof dependencies}\label{references-and-proof-dependencies}

Miller, S. D., \& Stephens-Davidowitz, N. (2019). \emph{Kissing numbers and transference theorems from generalized tail bounds}. arXiv:1802.05708v2. \url{https://arxiv.org/pdf/1802.05708} . The precise closed-tail Lemma3.11 and Theorem3.2 transference normalization were read in the independent DO072 review; the present proof uses exactly those already checked inequalities through DO072(GS9--16).

Exact actual-pencils: LIFT44(PR4--7), K070(PT1--3), K071(NB2--8), K077(AP1--3). Actual scalar-cross square/Green/metric/real interfaces: DO072(GS4--9,GS17--27), AM-DO-070-C1, K085 and LIFT063. Actual source-lattice stabilizer refinement: DO073(RF1--18), independently whole-reviewed with its final frozen proof hash \nolinkurl{9823924a67acb1957758b6ea1ee52768f6ba0b9759d55b98b0adcd8f7d5fe631} before this finding.

The new ingredients are the actual type-III intersection argument TI8, its paired integral S-row consequence TI9, the resulting ideal-size relation TI12, and the checked support-twelve consequence. Numerical certificates do not purport to automate the geometric or automorphic dependency proofs.

\subsection{Author verification record}\label{author-verification-record-25}

The author whole-read all TI1--27, the exact type-I projective-action/rho argument, both local CM components, the complete reference/scope paragraphs and the entire companion script. The final mathematical reread completed at 2026-11-22T16:29:59.445088+00:00. The arithmetic certificate passed all 1,925 exact assertions at 2026-11-22T16:27:04.176834+00:00. The final prose correction makes explicit that 5379 is the single global archimedean fibre constant; only T\^{}8 or T\^{}9 is local. This does not change an equation or the certificate. Actual primary/dependency reads are documented in the independent DO072/DO073 review records; no numerical test is presented as a proof of the incidence or automorphic interfaces. This is an author-checked controlled-recovery finding, awaiting separate root/peer integration.

\clearpage

\section{AM-DO-074 --- trace-conditioned rigidity at ten off-diagonal places}\label{am-do-074-trace-conditioned-rigidity-at-ten-off-diagonal-places}

\textbf{Canonical checked version --- 2026-11-22T16:33:21.651714+00:00.} Root read the whole final proof and full script and separately replayed the exact diagnostic. Original pending/status wording below is retained as export history; only this stated scope is promoted: Actual all-top H133 residual, block rank r\textless=22 and typeIII count t\textless=10 imply global unipotence, at any support. At support12 nonscalar IC8 and nonunipotent scalar-cross terms retain only t11/12. No new absolute error or complete support12 exhaustion.

Author: \nolinkurl{/root/definite_oscillator}, 22 November 2026 UTC. Controlled-recovery proof; existing proofs, indices and ledgers are unchanged.

\subsection{1. Exact actual scope}\label{exact-actual-scope}

Retain the actual J56 remainder of K068/K069, including its determinant-one central representative, all334 unmixed unipotent G reductions, original pencil incidences, finite projectors, class masses and real selected-module normalization. This report uses the \textbf{all-top} cutoff H=133 only after the charged R048 annular removal. It gives no such cutoff for the uniform family.

Write \nolinkurl{E^169=S perpendicular W0}, \nolinkurl{dim S=164}, \nolinkurl{sig(W0)=(4,1)}, and \nolinkurl{A=h_SS}. Let \nolinkurl{r=rank_E(A-I)}, and let t count precisely the type-III unequal-pencil off-diagonal places. The total nonintegral support s may be larger. The new statement is

\[
 r\le22,\quad t\le10
 \quad\Longrightarrow\quad
 h\text{ is globally unipotent},\quad(h-I)^3=0,
 \quad\operatorname{rank}_E(h-I)\le2.
 \tag{TC1}
\]

The original K068 residual and the all-top hypothesis are essential. No global unipotence is assumed in the proof.

\subsection{2. First force the trace coefficient to vanish}\label{first-force-the-trace-coefficient-to-vanish}

K069(IC4), used with the actual block rank r, supplies a positive rational E-subspace P of S fixed pointwise by h, of dimension at least \nolinkurl{164-r-1}. As in DO067(ST2), its Hermitian orthogonal complement is also its Euclidean orthogonal complement because P lies in the original positive coordinate space. Restriction therefore does not enlarge the operator norm. Pad by positive identity directions if needed to obtain a signature-(k-1,1) operator h0 with

\[
 k=r+6\le28,\qquad \|h_0\|\le133.
 \tag{TC2}
\]

The coefficients \nolinkurl{e_j(h-I)} agree with \nolinkurl{e_j(h0-I)}. They remain in the original field E; no arithmetic basis of this moving complement is asserted integral.

The accepted signature-one singular-value formula gives singular values \texttt{R,1,...,1,R\^{}\{-1\}}, with \nolinkurl{1<=R<=133}. Hence the first coefficient b1 satisfies, at every CM realization,

\[
 |b_1|=|\operatorname{tr}(h_0)-k|
 \le133+133^{-1}+2k-2<256.
 \tag{TC3}
\]

K070(PT5--8) and DO067(ST6) retain the exact finite norm support: each characteristic coefficient is integral away the2t CM primes above type III, with pole order at most one there; it vanishes modulo every unmixed G prime and its inverse-adjoint CM partner. A nonzero coefficient bounded by B at all infinities must therefore obey

\[
 Q^{2(f-1)}\le B^{2f}T^{2t},\qquad f=168,
 \quad Q=g_0g_1.
 \tag{TC4}
\]

This is a coefficient statement, not entry integrality in the standard frame at type II. The actual primes are unchanged. The exact integer-power bounds of DO073 are

\[
 \log_2T<995+1812/4096,\qquad
 \log_2Q>70+1123/4096.
 \tag{TC5}
\]

At t=10 and B=256, the strict inverse binary margin in TC4 is

\[
 334(70+1123/4096)-336\cdot8
       -20(995+1812/4096)=1791437/2048>0.
 \tag{TC6}
\]

It only increases when t decreases. Therefore b1=0 and \nolinkurl{tr(h0)=k} algebraically over E.

\subsection{3. A second coefficient bound conditioned on the exact trace}\label{a-second-coefficient-bound-conditioned-on-the-exact-trace}

Here is the new spectral step. It uses the actual trace equality before estimating b2, rather than the earlier unconditional exterior-power bound.

For a signature-one unitary operator, generalized eigenspaces pair only for reciprocal-conjugate eigenvalues. The full generalized eigenspace with modulus greater than one is totally isotropic, so has dimension at most one; the contracting part has the same bound. This is the accepted signature-one spectral argument in DO061, also written with generalized multiplicities in K083 section3. It applies to h0 without semisimplicity. Thus either every eigenvalue has modulus one, or there is one pair alpha,beta with \nolinkurl{beta=1/bar(alpha)} and k-2 remaining eigenvalues lambda\_j of modulus one.

If all eigenvalues have modulus one, their sum k forces each to be1 by equality in the real-part bound. Then h0 is already unipotent, even if it has Jordan blocks.

In the other case, TC6 gives

\[
 \alpha+\beta=k-\sum_{j=1}^{k-2}\lambda_j,
 \qquad |\alpha+\beta|\le2k-2,
 \qquad |\alpha\beta|=1.
 \tag{TC7}
\]

Consequently

\[
 \left|\operatorname{tr}(h_0^2)-k\right|
 \le(2k-2)^2+2+(k-2)+k=(2k-2)^2+2k.
 \tag{TC8}
\]

Since \nolinkurl{tr(h0-I)=0}, Newton's identity now gives

\[
 b_2=e_2(h_0-I)=-\tfrac12\operatorname{tr}((h_0-I)^2),
 \qquad
 |b_2|\le2k^2-3k+2\le1486<1536.
 \tag{TC9}
\]

The bound is valid at every embedding separately. In the all-unit-circle case b2 is zero, so the same inequality applies. Jordan blocks do not alter either trace, and no global choice of expanding eigenvalue is needed.

The exact inequality \nolinkurl{1536^5<2^53} gives \nolinkurl{log2 B<53/5}. Applying TC4 a second time, now to b2, yields the strictly positive margin

\[
 334(70+1123/4096)-336(53/5)
       -20(995+1812/4096)=11521/10240>0.
 \tag{TC10}
\]

Thus b2=0. Newton's identity gives \nolinkurl{tr(h)=tr(h^2)=169}. The already checked signature-one two-trace theorem DO061(section2) implies global unipotence; its nilpotent normal form gives \nolinkurl{(h-I)^3=0} and rank at most two. This proves TC1. The argument first used TC6, then TC9, and never assumed the second trace equality in deriving its bound.

\subsection{4. What this says at support twelve}\label{what-this-says-at-support-twelve}

The carried K069 table removes the fixed-space class F at all block ranks r\textgreater=2 through s=12. The carried DO064 table removes the complementary IC8 class at ranks r\textgreater=3 through s=12. Thus the nonscalar IC8 ranks potentially relevant there are r1 and r2, both within TC1.

For these IC8 ranks, t\textless=10 is now impossible: TC1 would give global unipotence, and DO065's exact exceptional rigidity would then force A=I, contrary to r\textgreater0. Hence any remaining \textbf{nonscalar IC8 term at s12} has t11 or t12. For A=I, TC1 likewise excludes nonunipotent scalar-cross terms with t\textless=10; the nonunipotent t11 and t12 cases remain explicit. The theorem applies to fixed-W0 terms too, but gives no new absolute estimate for their globally unipotent part.

For clarity, existing counting statements remain separate. R050's actual all-top F,r1 table includes s12,t11 and s12,t12; K081 includes the corresponding nonunipotent fixed-W0 cells. Once global unipotence is supplied, DO065's two full F,r1 bounds can be evaluated at s12 independently of how it was supplied. Their exponents are

\[
 -103530801/128+8966\cdot12,
 \qquad -408071827/256+14941\cdot12.
 \tag{TC11}
\]

They have all their original support, pair, class and coefficient multiplicities. A future combined removal must retain their full charges or use an already charged containing set; this logical rigidity report does not silently append them to a ledger. Likewise the separate DO073/K086 scalar-cross bounds retain their own global-unipotence and profile hypotheses.

\subsection{5. Verification, references and limits}\label{verification-references-and-limits}

The companion certificate uses the actual DO040 primes and exact rational/integer comparisons. It checks TC3 for every k6 through28, the increasing conditional bound in TC9, both exact margins, every t0 through10, and the stated carried table endpoints. It records that k29 gives1597, outside the chosen1536 bound; no converse is inferred. These arithmetic checks supplement the symbolic spectral proof, not an enumeration of the global rational matrices.

The carried source interfaces are K069(IC4), K070(PT5--8), DO067(ST2--6), DO061(section2), K083(section3), and the charged all-top R048 cutoff. The new step is TC7--10, using the exact first trace to improve the second coefficient before applying the same original-field norm argument. No new representation theorem, finite projector, measure normalization or integral-frame choice is imported.

The all-top signed remainder still requires its separate signed lower bound. This proof gives no complete support12 exhaustion, no Bessel/Mram period value, no literal Host construction and no DP44 or RH conclusion. The nonunipotent t11/t12 classes and every surviving globally unipotent sector retain their actual signs and weights.

\clearpage

\section{AM-LIFT-065 --- Actual Read-to-row identification with its exact measure normalization}\label{am-lift-065-actual-read-to-row-identification-with-its-exact-measure-normalization}

\textbf{Canonical checked version --- 2026-11-22T16:33:21.651714+00:00.} Root read the whole final proof and full script and separately replayed the exact diagnostic. Original pending/status wording below is retained as export history; only this stated scope is promoted: Supplied actual Read plus models\_HE and genuine translated evaluation give exact normalized row identification (c*nP)\^{}(-1/2); full-HH extension also retains HCFinite and the actual finite-dimensional old spectral record. Construction of Read and joint prescribed spectral/column compatibility remain open.

\textbf{Status:} new source-interface proof, awaiting independent review. This document uses the recovered pinned source at commit \nolinkurl{7cbf7a0c32fca1badf05931879a37c3f8a2b03b4} and the checked genuine analytic evaluation theorem LIFT062-T. It does not construct a \nolinkurl{Read}, a complete \nolinkurl{STiesR}, or the requested 032 package. No Lean build is used.

\subsection{1. Precise inputs and conclusion}\label{precise-inputs-and-conclusion}

Retain the literal data of \nolinkurl{GlobalTheory.Mem.R3.Read}: a datum \nolinkurl{R}, a face \nolinkurl{F}, a bottom row side \nolinkurl{X0}, and a supplied record \[
 Y:\operatorname{Read}(R,F,X0).
 \tag{RB1}
\] Write \(E=D.\mathrm{what}_2.EE\), \(H=R.HE.H\), \(e=R.HE.e\), \(e_E=Y.eE\), and \[
 E_P=R.EEat(R.levR(F,\mathrm{true})),\qquad
 n=R.nP(F,\mathrm{true})=|(F.lvl(\mathrm{true})).P|>0.
 \tag{RB2}
\] The group in the final cardinal is a finite subgroup and contains its identity. This is the literal \nolinkurl{nP}, with no principal-level replacement.

Assume \nolinkurl{R.models_HE}. Its three fields identify \(e\) injectively with the complete \nolinkurl{EEat(triv)} model and identify its Hilbert inner product with \nolinkurl{ipM} for the actual measure \(\nu_0\) and domain \nolinkurl{fd0}.

Retain an actual Hida setting \(X\), actual \nolinkurl{GoodOp} \(G\), source group identification \(\iota\), and the supplied finite adele \(t\). Suppose the hypotheses of LIFT062-T have been established for these same data. In particular its evaluation equivalence is \[
 J_t:E_P\xrightarrow{\sim}EEC(X,G),\quad
 J_t(T)(g)=\operatorname{ev}_{v_0}(T)(\iota^{-1}(g)t),
 \qquad \|J_tT\|^2=c\,\operatorname{Re}ipM_{\nu_0}(T,T),\quad c>0.
 \tag{RB3}
\] The tuple \(v_0\) consists of the actual unit vectors in the prescribed minimal compact types. All level conjugacies, finite coefficient conjugations and Haar identities are those of LIFT062-T; they are not dropped from this theorem's hypotheses.

Then the following additional source hypotheses need not be assumed separately: \[
 \begin{split}
 \operatorname{range}(e_E)&=e^{-1}(E_P),\\
 \|e_Ex\|^2&=n\|x\|^2,\\
 E_P&\subseteq R.EEat(R.triv).
 \end{split}
 \tag{RB4}
\] The map \(k_0x=e(e_Ex)\) is a linear equivalence \(E\simeq E_P\). Consequently the exact normalized row identification is \[
 \boxed{j_{\rm row}x=(cn)^{-1/2}J_t(k_0x).}
 \tag{RB5}
\] It is a surjective linear isometry \(E\simeq EEC(X,G)\), and it satisfies the literal source \nolinkurl{RowJ} at this same \(t\). This follows from the source theorem \nolinkurl{exists_rowJ_of_evalOn} after proving (RB4), and is also proved directly below.

\subsection{2. The source identities that fix the scale}\label{the-source-identities-that-fix-the-scale}

Put \(C=\operatorname{toChainReal}(R,F,Y.tupT,Y.hm,Y.hTop)\), and let \(L\) be the actual \nolinkurl{PiTensorProduct.lift} linear equivalence. Direct source expansion gives \[
 \begin{split}
 C.EEc(\mathrm{true})&=L(E_P),\\
 C.measE(\mathrm{true})&=n^{-1}\nu_0,\\
 C.fdE&=R.fd0.
 \end{split}
 \tag{RB6}
\] In detail, \nolinkurl{IsMultMapUnivZero.lean} 65 specifies the measure, 82--84 specify \nolinkurl{JE}, \nolinkurl{chiE}, and the unchanged projector \nolinkurl{R.e0}. \nolinkurl{RecTie.lean} 171--175 defines \nolinkurl{ipE} and \nolinkurl{EEc}. \nolinkurl{LinksDefs.lean} 630--645 supplies the literal \nolinkurl{e0}, \nolinkurl{EEat} and \nolinkurl{levR}.

The \nolinkurl{Read} fields give a linear equivalence \[
 q=Y.Ish.jE:E\xrightarrow{\sim}L(E_P),\qquad qx=L(e(e_Ex)).
 \tag{RB7}
\] The first statement uses \nolinkurl{RowIds.jE}, not merely an injective map. The second is exactly \nolinkurl{Read.eE_read}. Moreover \nolinkurl{Read.tied.sharp.ee} gives \[
 \langle x,y\rangle_E=ipMT_{n^{-1}\nu_0}(qx,qy)
                       =ipM_{n^{-1}\nu_0}(k_0x,k_0y).
 \tag{RB8}
\] The tensor equivalence is removed by the definition of \nolinkurl{ipMT}; it adds no dimension, compact-type or factorial factor.

There is a minor choice-function issue worth resolving explicitly. The source defines \nolinkurl{ipM} to be the unique scalar in \nolinkurl{IsMultIP}, when such a scalar exists, and to be zero otherwise. If \(a>0\), multiplication of the measure by \(a\) preserves integrability and multiplies every defining integral by \(a\). Thus existence of \nolinkurl{IsMultIP} is equivalent for \(\nu\) and \(a\nu\). Its scalar is unique by evaluation on the retained unit tuple \(v_0\), for which the product of the self-inner-products is one. Both the existing and default-zero cases therefore give \[
 ipM_{a\nu}(T,T')=a\,ipM_\nu(T,T').
 \tag{RB9}
\] Here \(a\) is the finite positive real scalar, interpreted as its \nolinkurl{ENNReal.ofReal} in the measure. This argument does not silently replace the source definition by a different Petersson form.

Combining (RB8), (RB9) and \nolinkurl{models_HE} gives the polarized identity \[
 \langle e_Ex,e_Ey\rangle_H=n\langle x,y\rangle_E.
 \tag{RB10}
\] In particular the middle formula of (RB4) holds and \(e_E\) is injective.

\subsection{3. Exact range, with no extra level containment premise}\label{exact-range-with-no-extra-level-containment-premise}

For any \(x\), (RB7) and \(qx\in L(E_P)\) imply \(k_0x\in E_P\). Conversely, if \(T\in E_P\), then \(LT\in L(E_P)\), so surjectivity of \(q\) supplies \(x\) with \(qx=LT\). Applying \(L^{-1}\) gives \(k_0x=T\). Hence \(k_0\) maps onto all of \(E_P\).

If \(z\in H\) and \(ez\in E_P\), choose \(x\) with \(e(e_Ex)=ez\). Injectivity of \(e\) gives \(e_Ex=z\). This proves the first formula in (RB4), including the full preimage rather than just an inclusion.

Finally every \(T\in E_P\) is \(e(e_Ex)\), so it belongs to \(\operatorname{range}(e)=R.EEat(R.triv)\). Thus the last formula in (RB4) is already forced by \nolinkurl{Read.eE_read}, surjectivity of \nolinkurl{Ish.jE}, and \nolinkurl{models_HE}. No independent assumption about the direction of a finite subgroup inclusion is required.

\subsection{4. RowJ and the translated comparison map}\label{rowj-and-the-translated-comparison-map}

By (RB3) and (RB10), \[
 \|J_tk_0x\|^2=cn\|x\|^2.
 \tag{RB11}
\] Since both factors are surjective and \(cn>0\), (RB5) is a surjective linear isometry. Its actual evaluation is \[
 (j_{\rm row}x)(g)=(cn)^{-1/2}
 \operatorname{ev}_{v_0}(e(e_Ex))(\iota^{-1}(g)t).
 \tag{RB12}
\] This is exactly \nolinkurl{RowJ}: its positive real scalar is \((cn)^{-1/2}\), its unit tuple is the same \(v_0\), and its finite translation is the same \(t\). The source \nolinkurl{ConcreteLevelAdelicSubset.lean} 88--152 performs precisely this normalization. Equations (RB6)--(RB10) prove its extra range/norm/containment hypotheses from the actual \nolinkurl{Read} record.

For a \nolinkurl{RowsJoinPinned}-style formula whose evaluation equivalence itself is unscaled, the corresponding map from the arithmetic row space is \[
 k=(cn)^{-1/2}k_0:E\xrightarrow{\sim}E_P,
 \quad j_{\rm row}=J_tk.
 \tag{RB13}
\] Thus it is incorrect to insert the unscaled \(k_0\) in that identity unless \(cn=1\) has independently been proved. No such simplification is made here.

The source's \nolinkurl{RowsJoinPinned} is unshifted. At \(t=1\), with the literal full-Haar identity and \nolinkurl{IsLowVec} retained, the formula below has its exact evaluation convention. At a general supplied \(t=t_S\), (RB12)--(RB13) give the translated version; this statement does not identify it with the unshifted predicate.

\subsection{5. A Hilbert-space row tie exists; joint spectral compatibility is separate}\label{a-hilbert-space-row-tie-exists-joint-spectral-compatibility-is-separate}

For this additional finite-dimensional Hilbert extension, retain the literal \nolinkurl{HCFinite} input (\nolinkurl{ScalarOnPartsFull.lean} 276--278), together with its actual \nolinkurl{IsInfFam} and open finite-level kernel hypotheses. Equivalently it suffices to retain that the actual full multiplier space below is finite-dimensional. This is not needed for (RB4)--(RB13), and is not inferred from a \nolinkurl{finrank} equality. The full, unprojected evaluation equivalence in LIFT062-T identifies \[
 \operatorname{Mult}(R.MY0,J_R,\chi_R)\simeq\overline{HHX}.
 \tag{RB14}
\] \nolinkurl{Read.dall} identifies the dimension of its left side with \nolinkurl{D.what2.Dall}. Since complex conjugation preserves dimension, and \nolinkurl{D.s_HH} identifies that same number with \(\dim D.s.HH\), \[
 \dim D.s.HH=\dim HHX.
 \tag{RB15}
\] The original space \nolinkurl{D.s.HH} is explicitly finite-dimensional: \nolinkurl{Formalism/SharpRestriction/CuspSpanNormBounds/DataS.lean} 33--36 declares \nolinkurl{[hhF : FiniteDimensional ℂ HH]}, registered as an instance at line 86. The target is finite-dimensional by the retained \nolinkurl{HCFinite} input and the full evaluation equivalence; no inference of finite-dimensionality from a bare value of \nolinkurl{finrank} alone is used. In a genuine automorphic realization this is the usual fixed-level cuspidal finiteness input, explicitly separated here from the matrix-unit evaluation theorem.

Let \(s=D.sEE:E\to D.s.HH\) be the existing anti-linear isometry, and let \[
 t_0=\operatorname{conjHH}\circ j_{\rm row}:E\to HHX.
\] Both maps are anti-linear isometric embeddings. The rule \(s(x)\mapsto t_0(x)\) is a well-defined \textbf{complex-linear} isometry of their images: the two conjugations cancel. Their orthogonal complements have equal finite dimensions by (RB15). Choose orthonormal bases of those complements and extend this isometry by the resulting unitary map. We obtain \[
 \iota_H:D.s.HH\xrightarrow{\sim}_{\rm isom}HHX,
 \quad \iota_H(D.sEE x)=\operatorname{conjHH}(j_{\rm row}x).
 \tag{RB16}
\] This is the literal existential \nolinkurl{RowDictTie} for this newly chosen \(\iota_H\), and (RB13) gives its exact row-to-multiplier map. If the row space is zero, the partial map is empty and the same complement argument applies.

The construction does \textbf{not} show that \(\iota_H\) equals the \nolinkurl{eHH} of an independently supplied \nolinkurl{SIso}. In particular the extension is not proved to intertwine \nolinkurl{Pu}, \nolinkurl{Gterm}, local \nolinkurl{atP} operators, or the prescribed column maps. Those simultaneous operator conditions are essential in \nolinkurl{SRest} and \nolinkurl{STiesR}. Mere equality of Hilbert dimensions cannot establish them.

There is also a source-sanctioned spectral refill: for a supplied \nolinkurl{SPack P}, the source \nolinkurl{SPack.toNew} takes (RB15) and \(j_{\rm row}\) and sets its new \nolinkurl{sEE} equal to \(\operatorname{conjHH}\circ j_{\rm row}\). \nolinkurl{reS} preserves \nolinkurl{what2}, \nolinkurl{dsig}, the original \nolinkurl{atP}, \nolinkurl{Del}, letter and cloud data, as specified in its body. This shows exactly which data may be chosen to have the row equality. It does not prove a new \nolinkurl{SIso}, \nolinkurl{AutTie}, or complete \nolinkurl{SRest} for the mixed record.

\subsection{6. The remaining concrete obligation}\label{the-remaining-concrete-obligation}

This result removes independent \nolinkurl{hrange}, \nolinkurl{hnorm}, \nolinkurl{hle}, \nolinkurl{RowJ}, and freely chosen Hilbert-row-tie existence premises \textbf{once the actual \nolinkurl{Read}, \nolinkurl{models_HE}, and the genuine translated evaluation inputs have been supplied (and, for the finite-dimensional extension, the stated \nolinkurl{HCFinite} input)}. It retains their actual finite level, projector, chosen low vector and Haar scale.

The recovered \nolinkurl{exists_sStage_unconditional} theorem does not supply these inputs: its source constructs an abstract Euclidean spectral model with freely chosen maps and a coordinate row embedding. The recovered \nolinkurl{SRest.toSTiesR} merely carries a supplied \nolinkurl{rowPin} forward. Neither may be substituted for a concrete \nolinkurl{Read} and compatible spectral realization of the requested J/032 package.

Accordingly the remaining source gate can be stated more precisely: construct the actual arithmetic \nolinkurl{Read} and the actual Hida/column package, and establish compatibility of their selected Hilbert identification with the complete spectral operators and column maps. For an already fixed \nolinkurl{SIso}, its row equality remains a specific identity. This proof does not establish a nonzero prescribed period, a 032 value, a complete Host or Sizes family, or DP44.

\subsection{7. Exact source audit and verification scope}\label{exact-source-audit-and-verification-scope}

All references below are at the pinned commit named above. The bodies used were read directly, not inferred from their filenames.

\begin{enumerate}
\def\labelenumi{\arabic{enumi}.}
\tightlist
\item
  \nolinkurl{Development/AbelianVarietyTheory/Definitions/AutomorphicMembership/StableIffCarrierEq.lean}, \nolinkurl{Read} 116--193: \nolinkurl{Ish} 126, \nolinkurl{tied} 135, \nolinkurl{eE/eE_read} 178--179, \nolinkurl{dall} 180--181.
\item
  \nolinkurl{Formalism/LocalToGlobal/IndefiniteReal/IsMultMapUnivZero.lean}, entire file; \nolinkurl{toChainReal} 46--91, especially 65 and 82--84.
\item
  \nolinkurl{Formalism/LocalToGlobal/IndefiniteReal/RecTie.lean}, \nolinkurl{HilbTie} 20--21, \nolinkurl{ipMT} 44--45, \nolinkurl{ipE/EEc} 171--175, \nolinkurl{RowIds} 195--203, \nolinkurl{RowSideReal.ee} 216, \nolinkurl{RecTied.sharp} 408.
\item
  \nolinkurl{Formalism/LocalToGlobal/IndefiniteReal/LinksDefs.lean}, 630--645 and 675; \nolinkurl{LinksChar.lean}, \nolinkurl{models_HE} 189--190.
\item
  \nolinkurl{Formalism/LocalToGlobal/IndefiniteReal/ScalarOnPartsFull.lean}, \nolinkurl{IsMultIP/ipM} 94--101.
\item
  \nolinkurl{Development/AbelianVarietyTheory/Definitions/AutomorphicMembership/ConcreteLevelAdelicSubset.lean}, entire file, especially \nolinkurl{EvalOn/RowJ} 49--58 and \nolinkurl{exists_rowJ_of_evalOn} 88--152.
\item
  \nolinkurl{Formalism/LocalToGlobal/SiegelLatticeTheory/FlagStabilizerTensors/RowsJoinPinned.lean}, \nolinkurl{EP} 30 and predicate 32--40; \nolinkurl{RowDict.lean}, \nolinkurl{RowDictTie} 134--135.
\item
  \nolinkurl{Development/AbelianVarietyTheory/Definitions/AutomorphicMembership/SameDegreeZeroFormsSelf.lean}, entire file; \nolinkurl{SRest} 39--46, \nolinkurl{SPack.toNew} 64--71, and the explicitly assumed row tie in \nolinkurl{exists_sRest_of_sTies} 123--142.
\item
  \nolinkurl{Development/AbelianVarietyTheory/Definitions/AutomorphicMembership/RefilledStagesRecord.lean}, entire file, especially \nolinkurl{mix/reS} 20--35 and preservation statements 59--65.
\item
  \nolinkurl{Formalism/SharpRestriction/CuspSpanNormBounds/DataS.lean}, entire file; the original \nolinkurl{HH} finite-dimensional field at 33--36 and its registered instance at 86.
\end{enumerate}

The companion exact diagnostic checks the scalar normalization, the anti-linear-to-linear image identification, extension over orthogonal complements, and the distinction between a free Hilbert extension and a prescribed operator intertwiner. It is a finite-dimensional diagnostic, not a Lean proof or a construction of the actual arithmetic records. Source retrieval metadata and hashes are retained in this pending directory.

\clearpage

\section{AM-DO-075 --- paired congruences remove the small exceptional blocks at support twelve}\label{am-do-075-paired-congruences-remove-the-small-exceptional-blocks-at-support-twelve}

\textbf{Canonical checked version --- 2026-11-22T16:55:15.133699+00:00.} Root read the entire final proof and diagnostic script and replayed the stated exact diagnostics. The original frozen-status wording below is preserved as author-export history. Promoted scope: All-top H133 support12 IC8 block ranks1/2 have total absolute contribution\textless2\^{}-1859541, without global unipotence. Both CM-prime congruences,160 original unaffected coordinates, full source fibre, all ordered pairs, class masses and normalizations are retained.

Author: /root/definite\_oscillator, 22 November 2026 UTC. New controlled-recovery work; no prior proof or ledger is modified.

\subsection{1. Actual statement and finite block choices}\label{actual-statement-and-finite-block-choices}

Keep the actual J56 K068 residual, original normalized signed summand of K062, every finite projector and CM phase, the source/target pencils and class masses. Use only the charged \textbf{all-top} cutoff H=133. Set \[
 f=a=168,\quad n=169,\quad S\perp W_0=E^{169},
 \quad\dim_E S=164,\quad\operatorname{sig}(W_0)=(4,1).
 \tag{CR1}
\] The determinant-one representative h is integral away T, and all334 unmixed G reductions are unipotent. Put N=h-I and A=h\_SS. Restrict to the exact exceptional class IC8 of K069, with \[
 1\le r=\operatorname{rank}_E(A-I)\le2,\qquad s=12.
 \tag{CR2}
\] \textbf{Theorem.} The total absolute contribution of CR2 is less than \(2^{-1859541}\). No global unipotence is assumed. All types I,II,III are counted together with their actual pair weights. Remove its intersection with the current remainder and retain every older error.

DO064(PM5--9) says A is partial monomial with nonzero entries in mu28. Its path/cycle calculation proves A is the identity on all but at most2r coordinate vertices, including both their rows and columns. Choose four original S coordinates containing these vertices, padding if necessary. The remaining160 coordinates form S\_free. For every corresponding basis vector, \[
 Ae_j=A^\dagger e_j=e_j.
 \tag{CR3}
\] One may fix the padding once for each A. For an upper count it is harmless to sum over every four-subset and every possible four-by-four block, at most \[
                 \binom{164}{4}\,29^{16}<2^{112}.
 \tag{CR4}
\] Here every entry is zero or one of28 roots of unity. Overcounting blocks which do not occur is harmless. No moving integral frame is chosen.

For A different from I, some diagonal entry is0 or a nontrivial eta in mu28. The corresponding diagonal of \(A+A^\dagger-2I\) is −2 or \((\eta-1)^2/\eta\), a unit at every actual G prime. This is R049 section4/R051 section3: reduction of mu28 is injective there, and those odd primes do not divide28. It already forces every unmixed reduction of N to have rank at least two, since a nilpotent reduction of rank zero or one has square zero.

\subsection{2. Rank two forces simultaneous paired congruences}\label{rank-two-forces-simultaneous-paired-congruences}

Fix an unmixed G prime at which the reduction of N has rank two. In this paragraph work over that residue field. Nilpotence gives \(N^3=0\), so the actual inverse-adjoint relation gives \[
 N^2=h+h^{-1}-2I,\qquad
 (N^2)_{SS}=A+A^\dagger-2I\ne0.
 \tag{CR5}
\] Thus rank N²=1. Its image line ell has nonzero S projection. The rank-two nilpotent linear algebra also gives \[
       \operatorname{im}N\cap\ker N=\operatorname{im}N^2=\ell.
 \tag{CR6}
\] For \(e_j\in S_{\rm free}\), CR3 gives \((Ne_j)_S=0\) and \((N^2e_j)_S=0\). Since the S projection of ell is injective, the latter forces \(N^2e_j=0\). Now \(Ne_j\) lies in CR6 and has zero S projection, so it is also zero.

The conjugate CM component has the same nilpotent rank because it is inverse adjoint. The same argument there uses the same unaffected coordinates. Taking the paired adjoint of those fixed-column statements and then the inverse gives fixed rows at the opposite components. Hence \[
 (h e_n)_j\equiv0\pmod{\mathfrak p},
 \qquad (h e_n)_j\equiv0\pmod{\bar{\mathfrak p}}
 \qquad(j\in S_{\rm free}).
 \tag{CR7}
\] This does not use positivity over a finite field or the integrality of a divided isotropic vector. Both original component matrices are retained.

Let P2 be the product of residue sizes over relative unmixed G factors of rank two, and let J0 contain both E-primes at each of these factors. Then \[
 N_EJ_0=P_2^2,\qquad
 1\le N_EJ_0\le Q^{2(f-1)},\qquad Q=g_0g_1.
 \tag{CR8}
\] There are at most \(2^{334}\) rank-two-versus-higher patterns. There are no rank-one orientation choices in this proof.

The actual unipotent Green comparison K067(PS8--9) has exponents \[
                 L_2=335,\qquad L_3=501.
\] Every factor outside P2 has rank at least three. Consequently \[
 |a(h)|\le2Q^{-501(f-1)}P_2^{166}
         =2Q^{-501(f-1)}(N_EJ_0)^{83}.
 \tag{CR9}
\] The factor2 is the single global Euler correction. The factors used have indices at least166 and at most three per unmixed place, so their combined correction is less than two by the same product estimate as K067. The two mixed projectors retain only their modulus bound one. No character sign is inferred.

\subsection{3. Counts in the original coordinate lattices}\label{counts-in-the-original-coordinate-lattices}

Fix the support, all actual ordered pencil pairs, A with its four chosen exceptional coordinates, and J0. Let I\_Sigma be the paired inverse-prime ideal at the s active T places, of norm \(T^{-2s}\). The generic target-overlattice containment of DO059/DO064/R051, together with CR7, gives \[
 (h e_n)_{S_{\rm free}}\in(J_0I_\Sigma)^{160},
 \qquad (h e_n)_{S\setminus S_{\rm free}}\in I_\Sigma^4.
 \tag{CR10}
\] The W0-only shear is crucial here. These containments hold before global unipotence and for all actual types.

Put \[
 L=J_0I_\Sigma,\quad N_EL=z^{2f},\quad
 z_0=T^{-s/f},\quad z_1=Q^{(f-1)/f}T^{-s/f}.
 \tag{CR11}
\] Then \(z_0\le z\le z_1\). In the actual half-trace metric the covolume of L is \(D_0z^{2f}\), where \[
 D_0=7^{140}29^{144}5^{84}13^{84}.
\] The arbitrary-ideal dual and transference bound of DO070 gives \[
 \mu(L^{160})<c\sqrt{160f}\,z<4200\sqrt f\,z,
 \qquad c=2250/7.
 \tag{CR12}
\] This is not the smaller original-ring cover. The positive part of a negative unit column of any h has norm \(\sinh(\log\|h\|)<H/2=66.5\) at each embedding. Thus \[
 \#x_{\rm free}\le V_{160}(66.5+4200z)z^{-320f},
 \qquad
 V_d(C)=\frac{((22/7)fC^2)^{fd}}{(fd)!D_0^d}.
 \tag{CR13}
\] These are rational upper envelopes for Euclidean ball volumes. Zero coordinates are allowed; there is no all-nonzero or Gaussian hypothesis.

At s12 the exact primes give \(z_0<2^{-71}\) and \(1/4<z_1<1/2\). The four exceptional coordinates can therefore use radius \(67\sqrt f\), for a count at most \(V_4(67)T^{8s}\). The actual W0 target overlattice contains \(O_E^5\) and has covolume \(D_0^5T^{-10s}\), uniformly in all targets, including zero. LIFT063 gives its original-ring covering radius below \(128\sqrt f\). The negative column has norm below \(95\sqrt f\) at H133, so its W0 count is at most \(V_5(223)T^{10s}\).

This is a full five-complex-coordinate count, not a real-part slice. Thus every target has the same full paired determinant index \(T^{10}\); no type-dependent coarea or omitted zero-target factor is involved. For each complete negative column retain the original full source-lattice stabilizer fibre \(G_n<2^{5022}\). The scalar-cross rank-five refinement is not used for these nonscalar matrices.

\subsection{4. Convexity and complete multiplicities}\label{convexity-and-complete-multiplicities}

Multiplying CR13 by the J0 factor in CR9 leaves, apart from \(T^{166s}\), \[
 F(z)=V_{160}(66.5+4200z)z^{(166-320)f}.
 \tag{CR14}
\] Its logarithm is convex in log z: with \(v=\log z\), the only nonlinear term is \(320f\log(66.5+4200e^v)\). Hence F is bounded by the larger endpoint value on \([z_0,z_1]\). This covers each actual discrete ideal pattern, without a density assumption or an assertion that an endpoint ideal occurs.

At z0, \(66.5+4200z_0<67\). The weighted free count is at most \(V_{160}(67)T^{320s}\). Including the other nine coordinates leaves total T-power338s and Green denominator \(Q^{501(f-1)}\).

At z1 the Green loss restores exactly the rank-two denominator \(Q^{335(f-1)}\), since \((N_EJ_0)^{83}=Q^{166(f-1)}\) there. Also the radius is below2167 and \(z_1^{-320f}<2^{640f}\). This endpoint has free count at most \(V_{160}(2167)2^{640f}\), while the other nine coordinates retain T-power18s.

Retain once each: the A count \(2^{112}\), rank-pattern count \(2^{334}\), support count \(2^{168}\), source fibre \(2^{5022}\), Green constant2, and the full pair/weight bound \[
 P_T^s=\left(\frac{2(T-1)^2}{T-2}\right)^s<2^{997s}.
 \tag{CR15}
\] There is no additional type count: CR15 already bounds all ordered pairs and their actual rho values. Multiple incidences only enlarge the upper count. Taking a maximum of two endpoints does not introduce a factor two.

\subsection{5. Exact numerical bounds and scope}\label{exact-numerical-bounds-and-scope}

The exact certificate proves by integer powers and factorials \[
 V_{160}(67)<2^{-43274},\quad
 V_{160}(2167)<2^{226353},\quad
 V_4(67)<2^{2489},\quad V_5(223)<2^{5757}.
 \tag{CR16}
\] Use the actual bounds \(\log_2T<995+1812/4096\) and \(\log_2Q>70+1123/4096\). With \[
 C=112+334+168+5022+1+997\cdot12=17601,
\] the complete endpoint exponents are \[
\begin{split}
 E_0={}&C-43274+2489+5757
 +338\cdot12(995+1812/4096)
 -501\cdot167(70+1123/4096)\\
 ={}&-7616682681/4096,\\
 E_1={}&C+226353+640\cdot168+2489+5757
 +18\cdot12(995+1812/4096)
 -335\cdot167(70+1123/4096)\\
 ={}&-13749259803/4096.
\end{split}
\tag{CR17}
\] Their maximum is below−1859541. CR4 includes both block ranks, and the endpoint argument takes a maximum; neither produces another union factor. This proves the stated charge1859541. It applies only to the current-remainder intersection, with every earlier charge retained.

The result controls the nonscalar IC8 ranks1/2 at support12; larger exceptional ranks there were already bounded in DO064. It does not remove nonunipotent scalar-cross t11/t12 terms, all-II globally unipotent scalar-cross terms, or uncounted fixed-W0 rank-two sectors. Full support-twelve exhaustion and the signed lower bound remain separate. No period value, Host constructor, DP44 sign or RH assertion follows.

\subsection{6. Dependencies and verification limits}\label{dependencies-and-verification-limits}

The actual interfaces are DO064(PM5--9), R049 section4, K067(PS8--9), the generic source/target column and finite fibre count of DO059/DO064/R051, corrected DO070 arbitrary-ideal transference, and LIFT063 original-ring cover. The Green source is Gorodetsky and Hazan (2017), Theorem2.1/equation(2.1), through the already checked K067 formula. No new representation, projector, arithmetic frame, class mass or normalization is introduced.

The new steps are CR5--7, the paired ideal in the original-coordinate count, and the convex endpoint combination. The companion certificate verifies dimensions, finite enumeration bounds, all arithmetic constants, exact volume ceilings, prime/endpoint comparisons and complete exponents. It does not automate the kernel dependency proofs or enumerate rational isometries.

\clearpage

\section{AM-R-055 --- higher-support removal when few places have type II}\label{am-r-055-higher-support-removal-when-few-places-have-type-ii}

Author: \nolinkurl{/root}, GPT-6 Astra Pro, 22 November 2026 UTC. This is a new consequence of the original-coordinate incidence theorem AM-K086, in the unchanged actual J56 kernel and the charged all-top cutoff. It does not extend a numerical assertion merely by changing its stated range: the support dependence of every factor is retained below.

\subsection{1. Actual statement and surviving hypothesis}\label{actual-statement-and-surviving-hypothesis}

Keep the actual normalized signed summand, finite projectors, class masses, source and target pencils, half-trace metric and selected real modules of AM-K086. Thus \nolinkurl{f=168}, \nolinkurl{n=169}, \nolinkurl{S perp W0=E^169}, \nolinkurl{dim_E S=164}, \nolinkurl{sig(W0)=(4,1)}, and the all-top cutoff is \nolinkurl{H=133}. Restrict to scalar-cross globally rank-two unipotents

\[
h=I+N,\quad h_{SS}=I,\quad N^3=0,\quad
\operatorname{rank}N=2,\quad\operatorname{rank}N^2=1,
\quad x=(he_n)_S\ne0.
\tag{HS1}
\]

Let \nolinkurl{s=i+u+t} be the number of actual nonintegral T places, with counts \nolinkurl{i,u,t} of types I, II, III. Consider the following union of actual terms satisfying HS1:

\begin{longtable}[]{@{}ll@{}}
\toprule\noalign{}
Support s & Allowed type-II count u \\
\midrule\noalign{}
\endhead
\bottomrule\noalign{}
\endlastfoot
13 through 48 & 0 through 10 \\
49 through 84 & 0 through 9 \\
85 through 120 & 0 through 8 \\
121 through 156 & 0 through 7 \\
157 through 168 & 0 through 6 \\
\end{longtable}

\textbf{Theorem.} The total absolute contribution of this union is less than \nolinkurl{2^-1263}. Every support choice, ordered actual pencil pair, original-coordinate support and G pattern is included. Its intersection with the current remainder may be removed with this single additional charge while retaining all earlier charges.

The theorem is a sufficient sector removal. Failure to appear in the table means only that this theorem does not remove the term. It gives no positive lower bound for the omitted sectors and assumes no unproved structural reduction into HS1.

\subsection{2. Why the K086 estimate applies at every support}\label{why-the-k086-estimate-applies-at-every-support}

K086(TI5--9) proves the local type-III intersection assertion at each actual nonintegral place, without a hypothesis on s: \nolinkurl{N^2 S=0} makes both the forward and inverse displacements of each original S-vector lie in the intersection of the two distinct pencil lattices. The argument at the paired CM prime and the adjoint give both cross blocks integral. Type-I integrality is K071(NB5). Hence the only possible poles of x occur at type II, with at most a simple pole at either paired prime.

Accordingly the actual ideals in K086(TI10--13) satisfy, for any s,

\[
x\in (J_E I_{\rm II})^{164},\qquad
z^2=yw,\qquad w=T^{t/f},\qquad
y=(NJ_F T^{-(2u+t)})^{1/f},\qquad
p=R/(kz^2)\ge1,
\tag{HS2}
\]

where \nolinkurl{R=17424/133} and \nolinkurl{1<=k<=164} is the number of nonzero original coordinates of x. The square-pole ideal for alpha is still the one in K086(TI10): its type-II and type-III orders are two and one. We do not replace it with the smaller column-denominator ideal.

The constants \nolinkurl{5379}, \nolinkurl{1062}, \nolinkurl{107}, and the bound on the selected degree \nolinkurl{>2^959} in K086(TI15--17) do not depend on s. The first is the single global affine-fibre constant. Its support dependence is exactly \nolinkurl{T^(8s+r0)}, where \nolinkurl{r0} counts zero target labels. The complete weighted pair and zero-target calculation gives \nolinkurl{2^(i+t) T^i} for every profile, by K086(TI18--19). This includes zero as an admitted label and counts each relative place once. The source-lattice rank-five stabilizer bound requires HS1, which is retained.

There are at most \nolinkurl{2^168} support sets, \nolinkurl{3^s<2^(2s)} type assignments for s\textgreater=1, and \nolinkurl{4^334=2^668} common G patterns. These same sufficient bounds are valid throughout the table. The combinatorial factor \nolinkurl{binom(164,k)} already counts the original-coordinate support choices. Neither a new integral frame nor a different real metric is used.

Write

\[
P=669f/4=28098,\qquad b=673/4,\qquad c=2250/7,
\quad
H_{t,k}(p)=\left(\frac{375/w+kp/2}{376}\right)^f.
\]

The common factor and the two alternative lattice bounds are therefore exactly K086(TI21--22), with s left free:

\[
\begin{split}
\mathcal C_{s,i,t,k}={}&
2^{1+5379+1062+107+168+2s+668+i+t}
12^f2^{f/2-959f/4} Q^{-335(f-1)}\\
&\quad\cdot T^{343s-334i-1335t/4}R^P\binom{164}{k},
\end{split}
\tag{HS3}
\]

\[
G_{t,k}(p)=(kp)^{-P}H_{t,k}(p)e^{bkp}8192^k,
\quad
V_{t,k}(p)=(kp)^{-(P-fk)}H_{t,k}(p)V_k(1+c/\sqrt p),
\tag{HS4}
\]

where

\[
V_k(C)=\frac{((22/7)fC^2)^{fk}}{(fk)!D_0^k},
\qquad D_0=7^{140}29^{144}5^{84}13^{84}.
\]

The T exponent follows from the same actual Green, real and gamma densities as K086(TI23); explicitly it is

\[
(335/2)(2u+t)+8s+i-(669/4)t+t
=343s-334i-1335t/4.
\tag{HS5}
\]

For each k, log V decreases with log p because \nolinkurl{P-fk-f>0}, and log G is convex with log p.~Thus \nolinkurl{min(G(p),V(p))} is at most \nolinkurl{max(G(1),G(7),V(7))} on p\textgreater=1, with no extra factor three. This argument is independent of s. Finally, w\textgreater=1 gives \nolinkurl{H_t,k(p)<=H_0,k(p)}. We use this weaker envelope for every t in the table, keeping every other factor unchanged.

\subsection{3. A rational envelope and its monotonicity}\label{a-rational-envelope-and-its-monotonicity}

The exact one-sided constants from K086(TI25) are

\[
L_T=1019333/1024>\log_2T,\quad
L_Q=287843/4096<\log_2Q,\quad
L_R=3602/512>\log_2R.
\tag{HS6}
\]

For all profiles under consideration, the full T exponent is positive:

\[
343s-334i-1335t/4=9s+334u+t/4\ge9s>0.
\tag{HS7}
\]

Consequently substituting the upper bound L\_T is valid. We compare the resulting rational upper envelopes, not actual powers with a negative exponent. In particular subtracting the base expression below is an algebraic subtraction of two sufficient upper-bound formulas; it is not an application of L\_T to a possibly negative difference of actual T powers.

The freshly replayed k-by-endpoint calculation for \nolinkurl{(s,i,t)=(12,1,0)} has maximum

\[
B=-33159035/36864.
\tag{HS8}
\]

This calculation includes the exact binomial factor, gamma ratio \nolinkurl{(750+kp)/752}, inverse power \nolinkurl{(kp)^P}, rational volume, and the Gaussian bound \nolinkurl{(13/9)bkp+13k}. Its maximum ranges over all164 values of k and all three sufficient endpoints. It is the first base in K086(TI26); it is not an actual matrix value.

Replacing the common factor by HS3 and using \nolinkurl{H_t,k<=H_0,k} yields the profile-wise sufficient exponent

\[
\begin{split}
E(s,i,t)={}&B+2(s-12)+(i+t-1)\\
&+L_T\bigl(343(s-12)-334(i-1)-1335t/4\bigr).
\end{split}
\tag{HS9}
\]

Thus the absolute bound for each k and profile is strictly below \nolinkurl{2^E(s,i,t)}. All support sets and original-coordinate support choices for that k have already been included. For fixed s,u, substituting \nolinkurl{i=s-u-t} in HS9 shows that its increase per unit increase of t is \nolinkurl{L_T/4>0}. Its maximum is therefore at \nolinkurl{i=0,t=s-u}. At that allocation the envelope is

\[
E_*(s,u)=B+3s-u-25
+L_T\left(\frac{37s+1335u}{4}-3782\right).
\tag{HS10}
\]

Both slopes are strictly positive:

\[
\partial_s E_*=3+37L_T/4>0,\qquad
\partial_u E_*=-1+1335L_T/4>0.
\tag{HS11}
\]

Consequently the upper right corner of each rectangular block in the table controls every permitted profile in that block.

\subsection{4. Exact corners, the complete union and charge}\label{exact-corners-the-complete-union-and-charge}

The five corner exponents, computed over the rational numbers, are:

\begin{longtable}[]{@{}lll@{}}
\toprule\noalign{}
Maximum s & Maximum u & Sufficient exponent E\_*(s,u) \\
\midrule\noalign{}
\endhead
\bottomrule\noalign{}
\endlastfoot
48 & 10 & \nolinkurl{-47488853/36864} \\
84 & 9 & \nolinkurl{-17748167/9216} \\
120 & 8 & \nolinkurl{-94496483/36864} \\
156 & 7 & \nolinkurl{-59000149/18432} \\
168 & 6 & \nolinkurl{-8290667657/36864} \\
\end{longtable}

The largest is the first, and it is strictly below -1288. For a fixed support s and upper type-II count U\textless s, the exact number of integer triples \nolinkurl{(i,u,t)} with \nolinkurl{i+u+t=s}, \nolinkurl{0<=u<=U}, and \nolinkurl{i,t>=0} is

\[
\sum_{u=0}^U(s-u+1)=(U+1)(s+1)-U(U+1)/2.
\tag{HS12}
\]

Summing this expression over the five disjoint support intervals gives respectively

\[
10494,\quad22680,\quad32238,\quad39168,\quad13482,
\]

and hence118062 profiles in total. There are164 possible k values, so

\[
118062\cdot164=19362168<2^{25}.
\tag{HS13}
\]

There is no factor for endpoints, because HS4 uses a maximum, and no second coordinate-subset factor, because \nolinkurl{binom(164,k)} is already in HS3 and HS8. Thus the complete union has absolute contribution

\[
<19362168\,2^{-1288}<2^{-1263}.
\tag{HS14}
\]

The support ranges in this theorem are disjoint from the support-twelve conclusion of K086. Their intersection with other earlier removed sets is harmless: remove only their intersection with the present remainder and keep all earlier errors. One new charge1263 suffices; no second contained-sector charge is introduced.

\subsection{5. What this changes and what remains}\label{what-this-changes-and-what-remains}

Among the HS1 terms with support13--168, every surviving term must now have at least11,10,9,8,7 type-II places in the corresponding five support ranges. This is a reduction of actual arithmetic sectors, not a refinement of an unchanged error constant. The original source/target pencils, finite kernel and real cutoff are unchanged.

The theorem does not prove that all other structural classes are absent. It does not remove the all-II support-twelve sector, nonunipotent scalar-cross terms, fixed-W0 terms outside separate theorems, or nonscalar A outside their proved bounds. The original signed lower-bound problem, actual Mram nonvanishing, prescribed Read/032 package, actual DP44 g/e/H comparison and Hecke-RH extension remain distinct open steps.

In algebraic-geometric language the new restriction is on denominators and congruences of sections of the original fractional-ideal modules over \nolinkurl{Spec(O_E)}. The small type-II divisor follows from intersections of the actual pencil lattices, not from changing the polarized CM variety. It uses the same module/stabilizer framework recorded for K086: SGA3, Exposé I, 4.7.3 and 4.7.3.1. The exterior-coefficient context, where separately used in structural reductions, is SGA6, Exposé VI, Theorem3.3; it is not an additional input into this lattice count. The Picard, biextension, crystal and Cartier-duality readings remain the carried exact locators in the dictionary: FGA232, Theorem3.1, pp.148--149; FGA236, Corollaries1.7 and2.3, pp.229 and232; SGA7 I, Exposé VII, 2.7 and2.9, pp.169--174; cote4, scan003, §§8.1.1--8.1.2 and scans005--006, §§8.4.1--8.4.3; cote54, scans005--007. No new lambda-ring identity, integral K-theory comparison, l-adic weight assertion, cycle-class identification, CM-motive isomorphism or motivic L-function conclusion is inferred. The changed gate is the finite remainder used toward DP\_MULT Q3e and THETA3 G1; CX\_GRH CX-ALG-024 and GCT\_PLETH T-052 retain their separate open comparisons.

\subsection{References and check scope}\label{references-and-check-scope-1}

The numerical and lattice dependencies are AM-K086(TI1--27), AM-DO-072(GS4--27), the zero-target correction AM-DO-070-C1, AM-DO-073(RF1--19), AM-K085 and AM-LIFT-063, in their final checked versions. The precise transference and closed-tail literature inputs are those already checked through these findings: Miller, S. D., \& Stephens-Davidowitz, N. (2019), \emph{Kissing numbers and transference theorems from generalized tail bounds}, arXiv:1802.05708v2, Theorem3.2 and Lemma3.11, \url{https://arxiv.org/pdf/1802.05708} . They are used in their carried stated normalizations, not asserted in a newly different metric.

The companion script independently reconstructs the rational support-dependent envelope, replays all164 k values and the three sufficient t=0 endpoints, checks the actual prime and selected-degree bounds, and enumerates every symbolic type profile in the table. It does not enumerate rational isometries or mechanically certify incidence, automorphic, geometric or analytic proofs. Whole-proof review and arithmetic replay are separate check events.

\clearpage

\section{AM-LIFT-066 --- Canonical spectral refill with the arithmetic record fixed}\label{am-lift-066-canonical-spectral-refill-with-the-arithmetic-record-fixed}

\textbf{Canonical checked version --- 2026-11-22T16:55:15.133699+00:00.} Root read the entire final proof and diagnostic script and replayed the stated exact diagnostics. The original frozen-status wording below is preserved as author-export history. Promoted scope: Supplied genuine SPack, AutTie, actual Haar tie, LocTied, full dimension and normalized row map give complete identity SIso and SRest after the literal reS refill. STiesR still requires a compatible RecT and its exact degree-zero forms. Old comparison links and every changed-family DP44 term remain.

Status: author whole-proof/source review PASS; final exact diagnostic replay PASS (6,577 assertions). Frozen for independent review. This is a source-aligned mathematical theorem, not a Lean build. The conclusion is conditional on the listed genuine \nolinkurl{SPack}, local-model, Haar, and row hypotheses. It does not construct the requested arithmetic \nolinkurl{Read}, \nolinkurl{RecT}, or \nolinkurl{032} value.

\subsection{1. Exact statement and the remaining hypotheses}\label{exact-statement-and-the-remaining-hypotheses}

Fix the pinned source commit \nolinkurl{7cbf7a0c32fca1badf05931879a37c3f8a2b03b4}. Let \nolinkurl{D : D7}, let \nolinkurl{F_u : Forms0 D u}, and assume an actual package

\[
 P:\operatorname{SPack}(D,F). \tag{SR1}
\]

The word \emph{actual} here retains every field of \nolinkurl{SPack}: the Hida setting \nolinkurl{X}, the adelic bridge, \nolinkurl{Full X}, \nolinkurl{FinX}, the global additive character and its standard/local measure ties, the exact column transports, into/bijective/level/metric properties, the canonical global instance and embedding pins, and \nolinkurl{GoodOp}. None is supplied by merely naming a finite-dimensional space. The complete field list is in \nolinkurl{SameDegreeZeroFormsSelf.lean} 18--37.

Assume also:

\begin{enumerate}
\def\labelenumi{\arabic{enumi}.}
\tightlist
\item
  \nolinkurl{AutTie D P.X.B P.X.pl}, with its actual discrete-series representation, chosen unit low vector, finite Haar, fundamental domain, and formal-degree identity.
\item
  \nolinkurl{cdScaleTie D}, and that each real \nolinkurl{muG ε} is the genuine nonzero locally finite Haar measure on the real unitary group. Equivalently, the proof below may take the resulting equality \nolinkurl{D.s.cdScale = PairVal.cdScaleOf D} directly.
\item
  The literal \nolinkurl{LocTied D}: \[
  \forall p,\quad \operatorname{Nonempty}\bigl(\operatorname{LocIso}
      (D.s.atP(p),\operatorname{locC}(D,p))\bigr). \tag{SR2}
  \]
\item
  A full-space dimension equality and a row isometry: \[
  \dim_{\mathbf C}HH_X=D.what_2.Dall,
  \qquad j:D.what_2.EE\xrightarrow{\;\sim\;}EEC(X,P.G).
  \tag{SR3}
  \] The second map is a complex-linear Hilbert isometry. The first concerns the full \nolinkurl{HHX}, not only the finite-projector image \nolinkurl{EEC}.
\end{enumerate}

Define exactly the source constructor

\[
 Z=P.\operatorname{toNew}(\text{SR3}),\qquad
 D^+=\operatorname{reS}(D,Z).
 \tag{SR4}
\]

Then the copied genuine Hida/column package \nolinkurl{P⁺} on \nolinkurl{D⁺} admits \nolinkurl{SRest P⁺}. Its \nolinkurl{SIso} uses identity maps on the genuine full Hilbert space, all genuine column Hilbert spaces, all term sets, all pair sets, and the prime-index set. Its column and row pins are the literal definitions of \nolinkurl{SPack.toNew}.

Consequently, \textbf{if} a compatible \nolinkurl{RecT} on \nolinkurl{D⁺} is supplied and its degree-zero forms are those of \nolinkurl{P⁺}, the source constructor \nolinkurl{SRest.toSTiesR} supplies the literal \nolinkurl{STiesR} for that record. This last sentence retains all \nolinkurl{SizeData₂}, \nolinkurl{GatTies}, \nolinkurl{G.Inst}, \nolinkurl{G.LatInst}, constants, divisor, extra-data, and \nolinkurl{RecT} hypotheses appearing in \nolinkurl{ToSTiesRPin.lean} 15--19. It does not transfer these hypotheses from an old record by name.

\subsection{2. The permitted refill and its exact preserved fields}\label{the-permitted-refill-and-its-exact-preserved-fields}

The source definition is

\[
 \operatorname{mix}(s,t)=
 \{t\text{ with } cdScale=s.cdScale,
          \ Del=s.Del,\ letter=s.letter,\ atP=s.atP\},
 \tag{SR5}
\]

including the accompanying positivity, cloud-gap, injectivity, and letter-weight proofs. Thus the only non-proof data replaced in \nolinkurl{t} by data from \nolinkurl{s} are the four displayed entries. In particular \nolinkurl{HH}, \nolinkurl{Cu}, \nolinkurl{Pu}, \nolinkurl{volN}, \nolinkurl{volH}, term sets and term operators, term classifications, finite predicates, pair sets, pair invariants, and all scalar term/pair quantities are copied from \nolinkurl{t}. See \nolinkurl{RefilledStagesRecord.lean} 20--25 and the complete \nolinkurl{SData} record in \nolinkurl{DataS.lean} 13--84.

\nolinkurl{reS} changes precisely \nolinkurl{D.s} and the six compatibility slots \nolinkurl{s_split}, \nolinkurl{s_HH}, \nolinkurl{s_Cu}, \nolinkurl{s_PlF}, \nolinkurl{s_Cu_isom}, \nolinkurl{sEE}, with \nolinkurl{s_split} itself copied from the old record. All remaining fields of \nolinkurl{D7} are copied. In particular,

\[
\begin{aligned}
 (D^+).s.atP&=D.s.atP,& (D^+).s.cdScale&=D.s.cdScale,\\
 (D^+).s.Del&=D.s.Del,& (D^+).s.letter&=D.s.letter,\\
 (D^+).dsig&=D.dsig,& (D^+).what_2&=D.what_2.
\end{aligned} \tag{SR6}
\]

These are the literal source equations \nolinkurl{reS_atP}, \nolinkurl{reS_cdScale}, \nolinkurl{reS_Del}, \nolinkurl{reS_letter}, \nolinkurl{reS_dsig}, \nolinkurl{reS_rows}; \nolinkurl{RefilledStagesRecord.lean} 27--35, 59--65. There is no change of the divisor lattices: this theorem uses \nolinkurl{reS}, not the additional \nolinkurl{setDiv} operation in \nolinkurl{reD}.

\subsection{3. Why the preserved Haar scale is the canonical scale}\label{why-the-preserved-haar-scale-is-the-canonical-scale}

For fixed \nolinkurl{(n,ε)}, \nolinkurl{PairVal.CdTieAt D n ε c} is exactly the assertion that the measure \nolinkurl{c muG ε} equals the transported \nolinkurl{D.cd.cd P Q} for every compatible signature permutation and every measurable set. \nolinkurl{cdScaleOf D n ε} chooses a positive scalar with this property when one exists. The literal \nolinkurl{cdScaleTie} gives that property for the old positive scalar \nolinkurl{D.s.cdScale n ε}. These are the same formulas, not merely comparable measures: \nolinkurl{DataTies.lean} 10--15 and \nolinkurl{PairVal.lean} 294--315.

Let \nolinkurl{c} and \nolinkurl{c′} be two such scalars. Choose \nolinkurl{P} and \nolinkurl{Q} to count the true and false values of \nolinkurl{ε}, and choose a signature-preserving bijection \nolinkurl{Fin P ⊕ Fin Q ≃ Fin n}. A genuine nonzero locally finite Haar measure on this locally compact group has a measurable set \nolinkurl{A} with

\[
                 0<\mu_G(A)<\infty. \tag{SR7}
\]

For example, take a relatively compact identity neighborhood and then a compact neighborhood of positive measure. Both tie identities, evaluated on this same set and the same signature permutation, give \nolinkurl{c μ_G(A)=c′ μ_G(A)}. Cancellation yields \nolinkurl{c=c′}. The argument includes the zero-dimensional trivial group: its Haar mass is finite and positive.

\nolinkurl{cdScaleExists_of_tie} gives existence, and \nolinkurl{cdScaleOf_tie} gives the property of the chosen scalar. Therefore

\[
        D.s.cdScale=\operatorname{PairVal.cdScaleOf}(D). \tag{SR8}
\]

No arbitrary normalization of Haar is made. If a formal surrogate measure is allowed to be zero, this uniqueness argument would fail; the theorem explicitly requires the genuine Haar realization, or takes (SR8) as its hypothesis. Because the underlying \nolinkurl{D.cd} and real signature groups are unchanged, the same scalar and tie hold for \nolinkurl{D⁺}.

\subsection{4. The local model is a retained input, not an unnoticed replacement}\label{the-local-model-is-a-retained-input-not-an-unnoticed-replacement}

The exact definition is

\[
 \operatorname{LocTied}(D)\Longleftrightarrow
 \forall p,\ \operatorname{Nonempty}
  \operatorname{LocIso}(D.s.atP(p),\operatorname{locC}(D,p)).
 \tag{SR9}
\]

\nolinkurl{LocalAdditiveCharContinuous.lean} 237--239. \nolinkurl{LocIso} preserves the base and CM fields, their global embeddings, valuation rings, topologies, conjugation, and the additive character, with the precise scalar/valuation compatibility specified in \nolinkurl{EFSymmCont.lean} 15--31. In particular an abstract isomorphism of residue fields is insufficient.

By choice, select one actual \nolinkurl{LocIso} at each prime from (SR9). The canonical model \nolinkurl{locC D p} depends on the unchanged host field and conjugation; see \nolinkurl{EFSymmCont.lean} 45--46. Hence

\[
 \operatorname{locC}(D^+,p)=\operatorname{locC}(D,p),
 \qquad (D^+).s.atP(p)=D.s.atP(p). \tag{SR10}
\]

The same witnesses therefore prove \nolinkurl{LocTied D⁺}. These witnesses are the \nolinkurl{atP} components of the \nolinkurl{SIso} below.

A useful explicit special case is also in the source: \nolinkurl{LocIso.setSplit}, lines 100--114, gives the identity local isomorphism from a canonical model with a changed split flag to the canonical model. Thus \nolinkurl{loc_of_adj}, lines 116--122, supplies (SR9) when the old models are precisely those canonical models with the indicated split flags. This statement changes no finite conductor or character. It does not show that an arbitrary existing \nolinkurl{D.s.atP} has that form, and does not replace the separate split-orientation compatibility where it is required.

\subsection{5. Copying the genuine geometric package}\label{copying-the-genuine-geometric-package}

The refill does not change the host, ranks, Hermitian data, arithmetic row data, coefficient fields, class data, or local models \nolinkurl{atP}. Consequently the Hida setting, the degree-zero forms, their classes, levels and value maps, and the column transports have identical underlying data on \nolinkurl{D} and \nolinkurl{D⁺}. We use those copied data, rather than a newly chosen level or class system. The literal \nolinkurl{Forms0} record in \nolinkurl{DownNaturalLevelZero.lean} 12--30 consists of the existing T data, degree-zero classes/value maps, G level, local levels, Fourier maps, natural up/down level ties, norm and branching ties; it does not contain an abstract spectral-space or Poincaré-operator slot.

Here is the dependency check for the package used in (SR1):

\begin{longtable}[]{@{}
  >{\raggedright\arraybackslash}p{(\columnwidth - 4\tabcolsep) * \real{0.3333}}
  >{\raggedright\arraybackslash}p{(\columnwidth - 4\tabcolsep) * \real{0.3333}}
  >{\raggedright\arraybackslash}p{(\columnwidth - 4\tabcolsep) * \real{0.3333}}@{}}
\toprule\noalign{}
\begin{minipage}[b]{\linewidth}\raggedright
Package component
\end{minipage} & \begin{minipage}[b]{\linewidth}\raggedright
Data retained in the refill
\end{minipage} & \begin{minipage}[b]{\linewidth}\raggedright
Transport
\end{minipage} \\
\midrule\noalign{}
\endhead
\bottomrule\noalign{}
\endlastfoot
\nolinkurl{HidaSetting}, \nolinkurl{AdBridge}, \nolinkurl{inst}, \nolinkurl{pin} & Host, Hermitian data, ranks and all global embeddings & Same maps and same proofs \\
\nolinkurl{Full.X}, \nolinkurl{hK}, \nolinkurl{TR}, \nolinkurl{ψ_ne}, column systems, transversals & Same Hida setting, columns, lattice/flag data & Same functions and representative sets \\
\nolinkurl{Full.δ}, \nolinkurl{Full.M}, \nolinkurl{agree}, \nolinkurl{δ_pin} & The \textbf{old} \nolinkurl{D.s.atP}, which is preserved & Same local elements and equality proofs \\
\nolinkurl{FinX.shInj}, \nolinkurl{shSurj}, \nolinkurl{inHH}, \nolinkurl{puLin}, \nolinkurl{finHH} & Same geometric sharp map, \nolinkurl{HHX}, Poincaré map and norms & Same propositions and witnesses \\
\nolinkurl{FinX.fibFin}, \nolinkurl{pairIndep}, \nolinkurl{Greal}, \nolinkurl{volPos} & Same terms, representatives, embeddings and Haar measures & Same finite sets, scalars and positivity proofs \\
\nolinkurl{FinX.cdEx} & Same \nolinkurl{D.cd} and real groups & Same canonical scale-existence proof \\
Character, measure, local-compatibility and column metric fields of \nolinkurl{SPack} & All underlying fields and measures above & Same character and isometries \\
\nolinkurl{GoodOp}, exact degree-zero \nolinkurl{Transport/Into/Bij/LevelMatch} & Same raw levels, class maps and selected functions & Same operators and equalities \\
\end{longtable}

The only direct spectral-field dependency in the \nolinkurl{Full} record is the preserved \nolinkurl{atP} in its local \nolinkurl{δ/M} entries; inspect \nolinkurl{PoincareConstantEq.lean} 25--44. The geometric definitions of \nolinkurl{HHX}, sharp, Poincaré series and their finite-dimensionality in \nolinkurl{FinX} do not use the old abstract Hilbert space \nolinkurl{D.s.HH}. The latter is exactly what is being replaced. The source itself supplies the structurally stronger Hida transport \texttt{reX\ :=\ \{X\ with\}} across \nolinkurl{reD}, which additionally allows a division-lattice change; \nolinkurl{RedivisionStructureEq.lean} 16--26, 41--61. Here no such additional change is used. Moreover \nolinkurl{EHHre.lean} 18--27 explicitly copies the typed reduction and column systems and defines the full-Hilbert transport as \nolinkurl{LinearIsometryEquiv.refl}; this supplies the exact metric identification, including its zero-dimensional case.

Write \nolinkurl{P⁺} for this copied package. Proof witnesses may differ by proof irrelevance, but the mathematical maps, measures, terms, and metrics are identical. In particular the canonical genuine spectral records identify as

\[
                sOf(P^+.fin)=sOf(P.fin). \tag{SR11}
\]

For its \nolinkurl{Del} and \nolinkurl{letter} entries use (SR6); for \nolinkurl{cdScaleOf} use the unchanged \nolinkurl{D.cd}; for \nolinkurl{atP} use (SR10); every other entry is copied from the same genuine geometric package. The complete constructor is \nolinkurl{ParametersTermRelations.lean} 238--277.

\subsection{6. Full spectral compatibility by identity maps}\label{full-spectral-compatibility-by-identity-maps}

Let \nolinkurl{t=sOf(P.fin)}. The source chooses

\[
\begin{split}
 Z.s&=t,\\
 Z.s_{Cu}(u)&=\operatorname{colDictT}_u,\\
 Z.sEE&=\operatorname{conjHH}_{X,P.G}\circ j.
\end{split} \tag{SR12}
\]

\nolinkurl{SameDegreeZeroFormsSelf.lean} 64--71. The last map is the required conjugate-linear isometry, because \nolinkurl{conjHH} is conjugate-linear and \nolinkurl{j} is complex-linear isometric. The column metric is exactly \nolinkurl{P.isom}; it is not supplied by an unrelated Hilbert-space identification. \nolinkurl{FinX.finHH} and (SR3) give the full-space finiteness and dimension required by \nolinkurl{SNew}.

Now \nolinkurl{(D⁺).s=mix(D.s,t)}. By (SR8), its \nolinkurl{cdScale} agrees with \nolinkurl{t.cdScale}; its \nolinkurl{Del} and \nolinkurl{letter} agree by construction. Choose

\[
 e_{HH}=1_{HH_X},\quad e_{Cu,u}=1_{ColT_u},\quad
 e_T=e_P=e_{Pl}=1. \tag{SR13}
\]

For \nolinkurl{atP}, use the local isomorphisms of §4. These choices give every field of

\[
         h_I:\operatorname{SIso}((D^+).s,sOf(P^+.fin)). \tag{SR14}
\]

For completeness, the 29 \nolinkurl{SIso} fields fall into the following exhaustive groups:

\begin{itemize}
\tightlist
\item
  \nolinkurl{cd}: (SR8); \nolinkurl{del}, \nolinkurl{letter}: (SR6); \nolinkurl{atP}: (SR9)--(SR10).
\item
  \nolinkurl{eHH}, \nolinkurl{eCu}, \nolinkurl{eT}, \nolinkurl{ePl}, \nolinkurl{eP}: the identities (SR13).
\item
  \nolinkurl{Pu_eq}, \nolinkurl{volN}, \nolinkurl{volH}, \nolinkurl{aTerm_eq}: the same genuine operators and measures copied by (SR5).
\item
  \nolinkurl{isId}, \nolinkurl{inS}, \nolinkurl{relevant}, \nolinkurl{FeqF}, \nolinkurl{Gterm}, \nolinkurl{BBne}, \nolinkurl{generic}: the same predicates and arithmetic term values.
\item
  \nolinkurl{pfst}, \nolinkurl{psnd}, \nolinkurl{s0}, \nolinkurl{scalarPair}, \nolinkurl{pval}, \nolinkurl{ptOf}, \nolinkurl{cTerm}, \nolinkurl{massPair}, \nolinkurl{meanCount}: the same pair projections, vectors, actual period/volume scalars, and masses.
\end{itemize}

This is an operator-level \nolinkurl{SIso}; it is stronger than LIFT065's freely chosen isometry between the two old and new full Hilbert spaces. It works because the actual refill has already installed the genuine operators. It says nothing about an arbitrary old \nolinkurl{D.s.Pu} before the refill.

The old \nolinkurl{AutTie} transports unchanged: each field in \nolinkurl{DstarEq.lean} 10--22 refers to the retained ranks, host, Harish--Chandra data, \nolinkurl{what₂.Dall/Dstar}, and the same genuine \nolinkurl{B}, \nolinkurl{pl}, representation, low vector, finite Haar and fundamental domain. The column pin is now definitional, because \nolinkurl{colDictT} is literally \nolinkurl{colDict} on the underlying column space (\nolinkurl{ColDict.lean} 168--180). The row pin is (SR12):

\[
 h_I.eHH((D^+).sEE(x))
       =\operatorname{conjHH}_{X,P.G}(j(x)). \tag{SR15}
\]

Together with the transported \nolinkurl{AutTie} and \nolinkurl{LocTied}, these are precisely all fields of \nolinkurl{SRest P⁺}, \nolinkurl{SameDegreeZeroFormsSelf.lean} 39--46. This proves the theorem.

\subsection{7. The actual Read/evaluation choice of the row map}\label{the-actual-readevaluation-choice-of-the-row-map}

Theorem (SR1)--(SR15) accepts any row isometry satisfying (SR3). To use the actual arithmetic reading, retain the hypotheses of LIFT065: the literal \nolinkurl{Read}, its full-dimension clause, \nolinkurl{models_HE}, the source sharp metric, and the actual translated evaluation theorem LIFT062-T on this same Hida setting and Haar.

Let \nolinkurl{n=nP(F,true)>0}, let \nolinkurl{c>0} be the exact evaluation-measure scalar, and let \nolinkurl{t_f} be the literal finite translate. LIFT065 gives

\[
 j_{\rm row}=(cn)^{-1/2}J_{t_f}\circ HE.e\circ e_E.
 \tag{SR16}
\]

The resulting new \nolinkurl{sEE} is \textbf{exactly}

\[
 (D^+).sEE=\operatorname{conjHH}\circ j_{\rm row}. \tag{SR17}
\]

The full-space dimension in (SR3) is supplied by \nolinkurl{Read.dall} and full evaluation; its finiteness is already in \nolinkurl{P.fin.finHH}. Projected dimensions alone are not substituted for the full equality. All displayed evaluation identities are understood in the source's a.e. L² sense, or pointwise for its supplied smooth representatives.

No part of the refill changes \nolinkurl{t_f}, \nolinkurl{c}, \nolinkurl{n}, the source unit low vector, finite projector, or column Haar. There is no claim that \texttt{(cn)\^{}\{-1/2\}} is an algebraic unit or preserves an arithmetic lattice. An unshifted \nolinkurl{RowsJoinPinned} statement still needs the actual unshifted setting, or the exact translated corollary with the prescribed \nolinkurl{t_f}; the refill does not set \nolinkurl{t_f=1}.

\subsection{8. What remains fixed in a DP44 comparison}\label{what-remains-fixed-in-a-dp44-comparison}

The preservation statement concerns records and explicitly identified metrics, not every theorem whose hypothesis mentions a record. The distinction needed for R001 is:

\begin{longtable}[]{@{}
  >{\raggedright\arraybackslash}p{(\columnwidth - 4\tabcolsep) * \real{0.3333}}
  >{\raggedright\arraybackslash}p{(\columnwidth - 4\tabcolsep) * \real{0.3333}}
  >{\raggedright\arraybackslash}p{(\columnwidth - 4\tabcolsep) * \real{0.3333}}@{}}
\toprule\noalign{}
\begin{minipage}[b]{\linewidth}\raggedright
Datum
\end{minipage} & \begin{minipage}[b]{\linewidth}\raggedright
What the refill proves
\end{minipage} & \begin{minipage}[b]{\linewidth}\raggedright
What it does not prove
\end{minipage} \\
\midrule\noalign{}
\endhead
\bottomrule\noalign{}
\endlastfoot
Ranks, length, host, coefficient fields, \nolinkurl{sp}, Harish--Chandra parameters, \nolinkurl{what₂} & Unchanged & No repair of any existing Host/Size incompatibility \\
\nolinkurl{dsig}, \nolinkurl{cd}, \nolinkurl{cdScale}, \nolinkurl{atP}, \nolinkurl{s_split} & Unchanged; canonical \nolinkurl{cdScale} identified by (SR8) & No new local model or conductor theorem beyond the retained \nolinkurl{LocTied} \\
\nolinkurl{Del}, cloud-gap proof, \nolinkurl{letter}, letter-weight identities & Unchanged & No enlarged code, new weight profile or increased size window \\
Arithmetic \nolinkurl{mem/fm}, class lattices, divisor lattices, raw theta/period data outside \nolinkurl{s} & Unchanged & No assertion of an unknown integral inverse, image saturation or nonzero selected value \\
Full and column Hilbert spaces, \nolinkurl{Pu}, term and pair operators, volumes/scalars internal to \nolinkurl{s} & Replaced by the genuine \nolinkurl{sOf} data, then compared by identity isometries & No equality with unrelated old abstract operators or old abstract term sets \\
\nolinkurl{s_Cu} & Replaced by the actual column dictionary with its stipulated metric tie & No automatic transport of old basis denominators or chosen integral coordinates \\
\nolinkurl{sEE} & Replaced by (SR17), preserving its exact positive normalization & No hidden removal of \nolinkurl{c}, \nolinkurl{n}, finite translations, CM periods or arithmetic scalar valuations \\
\nolinkurl{SizeData₂}, \nolinkurl{RecT}, \nolinkurl{LinkS}, \nolinkurl{LinkIRR}, gain/divisibility comparisons & Must be supplied or proved for \nolinkurl{D⁺} where they use replaced data & No source-sanctioned automatic transport merely from \nolinkurl{reS} \\
\end{longtable}

At the abstract functional level, \nolinkurl{eHH} and \nolinkurl{eCu} in (SR13) and the row/column isometries have operator norm at most one in their stated metrics, and exactly one when their domain is nonzero. This adds no analytic operator-norm loss to that \textbf{new-to-genuine} comparison. It is not a numerical upper bound on the actual Poincaré map, and is not an arithmetic denominator statement.

R001's changed-family accounting must therefore still be applied to any comparison in which an operator, term set, integral realization, or normalization was changed. In its notation the complete expression retains

\[
       Q_M=\sum_j(g_j-e_j)-T\tau-MH+K_u. \tag{SR18}
\]

The refill alone sets no changed term in this formula to zero. Keeping ranks, weights and \nolinkurl{what₂} fixed is useful, but does not prove the comparison links that identify its gains, losses or testing vector. The actual arithmetic \nolinkurl{Read}, compatible \nolinkurl{RecT/Extra/SizeData₂} package, prescribed Bessel/Mram period, \nolinkurl{032} cycle/value, and ultimate DP44 conclusion remain separate construction/verification tasks.

\subsection{9. Source ledger and diagnostic scope}\label{source-ledger-and-diagnostic-scope}

All repository addresses below refer to the pinned commit in §1. Exact local copies and their SHA-256 digests are recorded by the companion certificate.

\begin{itemize}
\tightlist
\item
  \nolinkurl{Development/AbelianVarietyTheory/Definitions/AutomorphicMembership/RefilledStagesRecord.lean}, 10--65: \nolinkurl{SNew}, \nolinkurl{mix}, \nolinkurl{reS}, preserved fields.
\item
  \nolinkurl{.../SameDegreeZeroFormsSelf.lean}, 18--74, 109--121: \nolinkurl{SPack}, \nolinkurl{SRest}, exact refill constructor, local and row consequences.
\item
  \nolinkurl{.../ToSTiesRPin.lean}, 15--29: the additional compatible-record hypotheses and final literal constructor.
\item
  \nolinkurl{Formalism/LocalToGlobal/SiegelLatticeTheory/FlagStabilizerTensors/ParametersTermRelations.lean}, 170--184, 238--315: \nolinkurl{FinX}, complete genuine \nolinkurl{sOf}, all 29 \nolinkurl{SIso} fields.
\item
  \nolinkurl{.../PoincareConstantEq.lean}, 25--55: full geometric package and retained local elements/volumes.
\item
  \nolinkurl{.../ColDict.lean}, 168--188: exact column dictionary and metric.
\item
  \nolinkurl{.../RootLatticeThresholds/LocalAdditiveCharContinuous.lean}, 237--239: literal \nolinkurl{LocTied}.
\item
  \nolinkurl{Formalism/LocalToGlobal/SiegelLatticeTheory/EFSymmCont.lean}, 15--31, 45--46, 82--122: local model equivalence and split-flag identity.
\item
  \nolinkurl{.../ClosedLocalOrbits/PairVal.lean}, 294--315, and \nolinkurl{Formalism/LocalRingCardinality/DataTies.lean}, 10--21, 38--48: canonical and prescribed Haar scales.
\item
  \nolinkurl{Formalism/SharpRestriction/CuspSpanNormBounds/DataS.lean}, 13--84, and \nolinkurl{DstarEq.lean}, 10--22: full spectral fields and actual automorphic ties.
\item
  \nolinkurl{Development/AbelianVarietyTheory/Definitions/AutomorphicMembership/RedivisionStructureEq.lean}, 16--61: explicit Hida transport across the stronger refill-plus-redivision operation; no redivision is used here.
\item
  \nolinkurl{Development/AbelianVarietyTheory/Definitions/AutomorphicMembership/EHHre.lean}, 18--27: the explicit identity Hilbert transport and copied column/reduction data.
\item
  \nolinkurl{Formalism/LocalToGlobal/GroupLattices/DownNaturalLevelZero.lean}, 12--30: the exact degree-zero arithmetic form record transported unchanged.
\item
  AM-LIFT-065 and AM-R001: exact arithmetic \nolinkurl{Read} normalization and the unchanged-versus-replaced cost distinction.
\end{itemize}

The companion Python program verifies source field inventories/hashes, exact signature decompositions and positive-measure scalar cancellation examples, all identity-map field obligations in finite records, and rational instances of the normalization (SR16). These are source and arithmetic diagnostics; they do not replace the mathematical proof, genuine Haar property, \nolinkurl{SPack}, \nolinkurl{LocTied}, actual \nolinkurl{Read}, or compatible-record hypotheses. No Lean build is performed.

\clearpage

\section{AM-K087 --- a finite-rank budget for the all-II support-twelve remainder}\label{am-k087-a-finite-rank-budget-for-the-all-ii-support-twelve-remainder}

\textbf{Canonical checked version --- 2026-11-22T17:10:37.672950+00:00.} Whole final proof and script read by root and by the saved independent peer; separate exact diagnostics passed. Original pending wording below is preserved as author-export history. Promoted scope: All-top H133 support12 all-II globally rank2-unipotent scalar-cross terms with total unmixed finite G rank\textgreater=70 are removed. Remaining terms in this sector have total finite rank\textless=69 and at least265 identity reductions among334 unmixed factors.

Author: \nolinkurl{/root/kummer_extension}, 22 November 2026 UTC. New controlled pending work. This file changes no accepted kernel, class module, projector, or earlier error.

\subsection{1. Actual sector and theorem}\label{actual-sector-and-theorem}

Retain the original normalized signed summand of K062, the all-top H=133 cutoff, and the actual J56 host and primes. Let f=a=168, n=169, m=164, and write the original orthogonal coordinates as E\^{}169=S perpendicular W0 with dimensions164 and5 and W0 signature(4,1). Restrict to

\[
h=I+N,\quad h_{SS}=I,\quad N^3=0,\quad
\operatorname{rank}_E N=2,\quad\operatorname{rank}_E N^2=1,
\quad x=(h e_n)_S\ne0.
\tag{GB1}
\]

There are exactly s=12 actual nonintegral T places and all have type II: the actual pencil and its image have a unique common admitted lattice, and h fixes that lattice. Thus (i,u,t)=(0,12,0). Global unipotence in GB1 is an explicit hypothesis, not inferred from the support count.

At each of the334 unmixed relative G places, let d\_v be the rank of N in the actual residue representation. Reduction preserves global rank bounds, so d\_v is0,1 or2. The two mixed places retain their original partial projectors and are not included in this sum. Put

\[
D_G(h)=\sum_{v\ {\rm unmixed}}d_v.
\tag{GB2}
\]

\textbf{Theorem.} The total absolute contribution of the GB1 sector with D\_G(h) at least70 is strictly below2\^{}(-1445). Remove only its intersection with the current remainder and retain all older errors. Consequently every still-unremoved term in this all-II support-twelve sector has

\[
D_G(h)\le69,
\quad\text{and is the identity at at least265 of the334 unmixed G factors.}
\tag{GB3}
\]

This is a restriction on the surviving finite reductions. It proves neither their absence nor a favorable sign for their contribution.

\subsection{2. Exact common ideal and its global cap}\label{exact-common-ideal-and-its-global-cap}

Use the actual original-coordinate common-orientation construction of DO066 and K086 TI10--13. At each rank-one G factor choose the common vanishing CM side for all S-coordinates; at rank zero include both sides; at rank two include neither. The oriented E-ideal J\_E and paired F-ideal J\_F have equal absolute norms. Every actual matrix is covered by at most the carried4\^{}334 patterns; ambiguous orientations may be assigned once, and summing all patterns is a valid upper count.

Since the factors have residue sizes g0 or g1 and there are f−1 unmixed relative factors above each rational prime, the ideal norm is EXACTLY

\[
N J_E=N J_F
 =\prod_v q_v^{2-d_v}
 =\frac{Q^{2(f-1)}}{\prod_v q_v^{d_v}},
\quad Q=g_0g_1.
\tag{GB4}
\]

The conjugate CM primes encode one relative factor; their ranks are equal by the actual inverse-adjoint relation. They are not independent factors in GB2 or GB4. The choice of rank-one side does not change its norm q\_v.

Let I\_Sigma be the paired inverse-prime ideal at the twelve active T places, of norm T\^{}(-24), and set

\[
L=J_E I_\Sigma,\qquad N_E L=z^{2f},\qquad
z^{336}=N J_E/T^{24}.
\tag{GB5}
\]

All-II is the case y=z squared of K086. The actual prime order is g0\textless g1. Therefore D\_G(h) at least70 implies

\[
z^{336}\le\frac{Q^{334}}{g_0^{70}T^{24}}<339^{-336}.
\tag{GB6}
\]

The last strict inequality is the directly certified integer comparison

\[
Q^{334}339^{336}<g_0^{70}T^{24}.
\tag{GB7}
\]

Thus the entire stated sector has z\textless1/339. No probabilistic distribution of primes or ideals is used. In particular this argument does not replace a fractional ideal by a principal ideal with a chosen small generator.

\subsection{3. The actual full-coordinate count at fixed z}\label{the-actual-full-coordinate-count-at-fixed-z}

Put R=17424/133. The accepted original-coordinate nilpotent identities give

\[
\alpha=(N^2)_{nn}\in F,\quad
0<\alpha_\tau\le R,\quad\|x_\tau\|^2\le\alpha_\tau,
\quad N_F(\alpha)\ge z^{2f}.
\tag{GB8}
\]

The last inequality is the paired alpha ideal divisibility with all twelve square poles of order at most two. It is the same actual alpha used in the real compression and gamma count. No divided finite isotropic frame is introduced.

Count x in the full original product L\^{}164. This includes zero coordinates; no choice of k, binomial support factor, Gaussian bound or k-union is needed. The actual half-trace metric has covolume D0\^{}164 z\^{}(328f), where

\[
D_0=7^{140}29^{144}5^{84}13^{84}.
\]

The arbitrary-fractional-ideal transference input of corrected DO070 gives

\[
\mu(L^{164})<c\sqrt{164f}\,z,
\qquad c=2250/7,\qquad c\sqrt{164}<4117.
\tag{GB9}
\]

Since the trace-ball radius for x is at most square-root(fR), the Dirichlet-cell count is bounded by

\[
V_{164}(\sqrt R+c\sqrt{164}\,z)z^{-328f},
\qquad
V_j(C)=\frac{(528C^2)^{fj}}{(fj)!D_0^j}.
\tag{GB10}
\]

Here528=(22/7)f is a rational upper bound for pi f.~This uses the general fractional-ideal cover, not the smaller original-ring cover.

For fixed support, actual ordered pencil data and common ideal pattern, retain the exact affine real-part slice of K085 and corrected DO072. Every active type-II target is nonzero: a nonintegral h cannot fix the standard target labelled0. Hence the global W0 affine-fibre factor is at most2\^{}5379 T\^{}(8s), with5379 included ONCE, and the gamma count is at most

\[
2^{1062}\left(\frac{375+R/(2z^2)}{376}\right)^f.
\tag{GB11}
\]

The half-real projection module, its dyadic factor and square-root(D\_F) denominator are already retained in1062. There is no zero-target C1 factor in this all-II sector. The original-ring covering radius is used only for the affine W0 kernel, exactly as in that accepted slice.

For each complete actual negative column, the positive rank-five isometry-fibre bound of DO073 is G5\textless2\^{}107. It requires x nonzero, which GB1 retains. Counting phantom zero-S columns in the upper bound GB10 is harmless: we apply the fibre bound only to actual nonzero-S columns, whose count is no larger than GB10.

The actual Green coefficient and normalized real trace obey

\[
|a(h)|\le2Q^{-335(f-1)}(N J_E)^{335/2},
\] \[
|c(h)|\le12^f2^{f/2-959f/4}z^{-f/2}.
\tag{GB12}
\]

These retain the one global finite Euler correction, both mixed projector contractions, the original actual selected homogeneous degree and PBW isometry, and the favorable-point transport. They do not infer a sign from a rank or a local incidence.

The complete type-II ordered-pair/rho weight is at most1 per place. Retain support factor2\^{}168, common-pattern factor4\^{}334=2\^{}668, and even the harmless old type-assignment bound2\^{}(2s). Multiplying GB10--12 gives the following sufficient upper bound at a pattern of norm parameter z:

\[
\begin{split}
\mathcal B(z)={}&
2^{1+5379+1062+107+168+2s+668}
12^f2^{f/2-959f/4}
Q^{-335(f-1)}T^{343s}\\
&\times V_{164}(\sqrt R+c\sqrt{164}\,z)
 z^{13f/2}
 \left(\frac{375+R/(2z^2)}{376}\right)^f .
\end{split}
\tag{GB13}
\]

The T power is335s from Green plus8s from the actual W0 fibre. The z power before gamma is335f−328f−f/2=13f/2. Every class, incidence and finite-frame normalization is the carried one; no extra CM-partner multiplicity is inserted.

\subsection{4. Monotonicity, exact cap and complete charge}\label{monotonicity-exact-cap-and-complete-charge}

Apart from positive constants independent of z, the variable part of GB13 is

\[
z^{9f/2}(375z^2+R/2)^f
  (\sqrt R+c\sqrt{164}\,z)^{328f}.
\tag{GB14}
\]

Each factor is strictly increasing for z\textgreater0. The sector cap GB6 therefore permits evaluation at1/339. Use the exact elementary inequalities

\[
\sqrt R<229/20,\qquad c\sqrt{164}<4117,
\quad
C_*:=229/20+4117/339=159971/6780.
\tag{GB15}
\]

The companion script proves, using integer powers and factorials,

\[
V_{164}(C_*)<2^{-128308},
\] \[
339^{-1092}
\left(\frac{375+R339^2/2}{376}\right)^{168}
<2^{-6777}.
\tag{GB16}
\]

It checks the immediately preceding binary brackets too. These are ceilings for the complete rational factors, not floating-point estimates.

Use the actual one-sided bounds

\[
\log_2T<1019333/1024,\qquad
\log_2Q>287843/4096,\qquad 12<2^4.
\tag{GB17}
\]

The full positive T exponent343s permits the upper logarithm. At s12 the resulting complete exponent is

\[
\begin{split}
E={}&1+5379+1062+107+168+24+668
 +4116\frac{1019333}{1024}
 -335\cdot167\frac{287843}{4096}\\
&+168\left(\frac92-\frac{959}{4}\right)
 -128308-6777\\
={}&-5921131/4096<-1445.
\end{split}
\tag{GB18}
\]

The support/pattern/pair choices have already been summed or uniformly overcounted in GB13. The cap is uniform over their ideal norms. There is no endpoint union, no coordinate-size union, and no new witness choice. Thus2\^{}(-1445) is one complete charge for the theorem's actual sector.

If a surviving term had70 or more nonidentity G factors, its sum of finite ranks would be at least70 and it would have been removed. More precisely GB3 holds because each nonidentity factor has rank at least one and there are334 factors. This proves the stated at-least265 identity reductions, without asserting anything about the two mixed factors.

\subsection{5. Scope, dependencies and verification status}\label{scope-dependencies-and-verification-status}

Remove only the stated sector's intersection with the current remainder. Earlier numerical errors remain conservatively retained. The theorem leaves all-II terms with total finite rank at most69 open. It does not give a favorable sign, prove the full finite signed inequality, establish Mram nonvanishing, complete the actual arithmetic image, instantiate a literal Host/Read/source package, or imply a DP44/RH statement.

The actual interfaces are K062 AT14--24; DO066's original-coordinate common G pattern; K085's real-coordinate slice and actual homogeneous-energy bound; corrected DO070/DO072 and AM-DO-070-C1 for the paired coarea and arbitrary-ideal metric; LIFT063 for the original-ring affine-kernel cover; DO073 RF1--19 for the positive rank-five isometry fibre; and K086 TI10--24 for their simultaneous normalized combination. This result uses only the all-II specialization of those inputs. It neither claims their prior numerical support conclusions at new values nor reuses an unproved favorable sign.

The new proof is the exact finite-rank/ideal-norm identity GB4, its actual-prime cap GB7, and the increasing full-coordinate envelope GB13--18. The numerical certificate checks all constants, prime powers, dimensions, metric/radius comparisons, exact rational binary brackets and the final charge. It does not enumerate rational isometries or mechanically prove the actual incidence and automorphic interfaces.

Author whole-file check PASS at 2026-11-22T17:01:00.970882+00:00. I reread and reconstructed the complete saved GB1--18 proof and entire companion script, including the final strict-rational-bracket checks. The final certificate passed all 31 exact assertions at 2026-11-22T17:00:44.616423+00:00. The actual paired ideal, monotone variable power, original-coordinate ball and affine slice, relative pair multiplicities, rank-five fibre and source real normalization were checked as stated. Final fingerprints and check scope are recorded separately. No rational-matrix enumeration or full signed-sum claim is made.

\clearpage

\section{AM-DO-076 --- a reciprocal eigenvalue slice for nonunipotent rank-two scalar-cross terms}\label{am-do-076-a-reciprocal-eigenvalue-slice-for-nonunipotent-rank-two-scalar-cross-terms}

\textbf{Canonical checked version --- 2026-11-22T17:10:37.672950+00:00.} Whole final proof and script read by root and by the saved independent peer; separate exact diagnostics passed. Original pending wording below is preserved as author-export history. Promoted scope: All-top H133 support12 scalar-cross A=I, global rank(h-I)=2, nonunipotent terms with t11/12 are removed. The rational reciprocal eigenline, full-sector inversion, actual weighted S-lattice and full source fibre are retained. Nonunipotent global ranks3--6 remain outside.

Author: /root/definite\_oscillator, 22 November 2026 UTC. Controlled recovery work; no existing proof or ledger is changed.

\subsection{1. Exact sector and statement}\label{exact-sector-and-statement}

Keep the actual J56 K068 residual and the charged all-top cutoff H=133. All finite projectors, ordered pencil incidences, original coordinates, CM phases, class masses and selected real normalization are unchanged. Set \[
 f=a=168,\quad n=169,\quad E^{169}=S\perp W_0,\quad
 \dim S=164,\quad \operatorname{sig}(W_0)=(4,1).
 \tag{RS1}
\] Write N=h-I and A=h\_SS. This report concerns precisely \[
 A=I,\quad h\text{ does not fix }S,\quad
 \operatorname{rank}_E N=2,\quad h\text{ is not unipotent},
 \quad s=12,\quad t\in\{11,12\}.
 \tag{RS2}
\] Here t is the number of type-III unequal-pencil off-diagonal places. Every unmixed G reduction is unipotent, as required by the actual K068 residual. Global unipotence is expressly excluded in RS2.

\textbf{Theorem.} The total absolute contribution of RS2 is less than \(2^{-5341}\). This is a fresh bound on its intersection with the current remainder, not an inference that it had already been charged. It does not cover scalar-cross nonunipotent global ranks three through six or the globally unipotent all-II sector.

The main new ingredient is a reciprocal eigenvalue in F which makes the original S-column small in a weighted metric. The count below retains the full source-lattice fibre \(G_n<2^{5022}\); it does not require an extension of the rank-five fibre theorem.

\subsection{2. A rational isotropic eigenline and its reciprocal columns}\label{a-rational-isotropic-eigenline-and-its-reciprocal-columns}

Let V be the nontrivial primary part of h. Nonunipotence and det h=1 require at least two dimensions of nonidentity eigenvalues. Since rank(h-I)=2, they occupy exactly two dimensions; the eigenvalue-one part is semisimple, and \[
 V=\operatorname{im}(h-I),\qquad
 E^{169}=\ker(h-I)\perp V.
 \tag{RS3}
\] In particular V is nondegenerate and rational over E. This also follows from \(\operatorname{im}(h-I)=\ker(h-I)^\perp\), using the unitary adjoint and the absence of a generalized one-eigenvector in V.

Write the original block matrix as \(h=\left(\begin{smallmatrix}I&C\\B&D\end{smallmatrix}\right)\), and use * for ordinary conjugate transpose. Unitarity gives \(B^*J_W B=0\) and \(C J_W C^*=0\). Because h does not fix S, B is nonzero; otherwise the other block identity forces C=0. Thus N(S) is a nonzero isotropic line in W0. It lies in V, so V has signature(1,1) at every embedding.

Put \(b=\operatorname{tr}(h|_V)-2\). It is nonzero: det(h\textbar V)=1 and trace2 would make its characteristic polynomial \((X-1)^2\), contrary to the definition of V. The degree-two Cayley--Hamilton identity and unitarity give, on V, \[
 h+h^{-1}=(2+b)I,\qquad
 N^\dagger N=2I-h-h^{-1}=-bI.
 \tag{RS4}
\] The adjoint here is the actual indefinite Hermitian adjoint. In particular b lies in F, since the left member of the first identity is self-adjoint.

Let P be the rational orthogonal projection to V. Since N(S) is totally isotropic and b is nonzero, RS4 shows P(S) is totally isotropic. It is a nonzero line, denoted ell. Moreover A=I says \[
 \langle NPu,Pv\rangle=\langle Nu,v\rangle=0
 \qquad(u,v\in S).
\] The orthogonal complement of ell inside the two-dimensional V is ell itself. Hence \(N\ell\subseteq\ell\); N is invertible on V, so \(N(S)=N\ell=\ell\). Thus ell is contained in the original W0 and is an h-eigenline defined over E. Denote its eigenvalue by lambda.

The quotient eigenvalue on V/ell is \(1/\overline\lambda\), by nondegeneracy of the Hermitian pairing. Its determinant is one, so \[
 \lambda\in F^\times,\quad \lambda\ne1,\qquad
 b=\lambda+\lambda^{-1}-2\ne0.
 \tag{RS5}
\] This reasoning allows lambda=-1 and a nonsemisimple action on V; it does not assume that the two eigenlines are distinct.

For u in S, Nu lies in ell and \[
 h^{-1}u=u-\lambda^{-1}Nu.
\] Taking the original paired adjoints of this identity gives the S parts of the two negative columns: \[
 x=(he_n)_S\ne0,\qquad x^-=(h^{-1}e_n)_S=-\lambda x.
 \tag{RS6}
\] For explicit signs, \(B^-=-\lambda^{-1}B\), while the unitary inverse has negative S column \(-B^*e_n\). Applying the same formula to the inverse gives \(x=\lambda^{-1}B^*e_n\), proving RS6. Here e\_n in the block formula denotes the original negative coordinate of W0. No finite normalization of ell is made.

\subsection{3. The small weighted metric and legitimate inversion}\label{the-small-weighted-metric-and-legitimate-inversion}

Here is a direct norm bound in the unchanged positive coordinates. In the usual positive/negative block decomposition of h, write its positive block as P\_+, with negative column x\_all and inverse negative column y\_all. If its Cartan radius is r, then \[
 P_+=U\bigl(I+(\cosh r-1)vv^*\bigr),\qquad
 x_{\rm all}=\xi\sinh r\,Uv,\quad
 y_{\rm all}=\xi'\sinh r\,v
\] for a positive unit vector v, a unitary U and unit scalars \(\xi,\xi'\). These are the standard polar identities following from \(P_+^*P_+=I+\sinh^2r\,vv^*\). If r=0, scalar-cross is impossible.

Since the S compression of P\_+ is I, subtracting the S compression of U gives a rank-one operator of norm at most two. Therefore \[
 \|x_\tau\|\|x^-_\tau\|
 \le2(\cosh r_\tau+1)
 \le H+H^{-1}+2<136.
 \tag{RS7}
\] Together with RS6 this yields \[
 |\tau(\lambda)|\,\|x_\tau\|^2<136
 \quad\text{at every real label.}
 \tag{RS8}
\] Also \(H^{-1}\le|\tau(\lambda)|\le H\), because lambda and its reciprocal are eigenvalues of h.

The entire sector RS2 is invariant under h↦h\^{}-1. K062(AT7--8, AT14, AT23--24) proves this for the actual ordered incidences and weights: the T incidence matrix is transposed, F and C are unchanged, the G and selected real coefficients are conjugated, and the cutoff is preserved. A=I is preserved by the inverse adjoint, as are s,t and global rank. Its eigenvalue in RS5 changes to lambda\^{}-1. Consequently, to bound the total absolute sector, one may restrict to \[
                 |N_{F/\mathbb Q}\lambda|\ge1
 \tag{RS9}
\] and multiply by two. This is applied to the larger full sector before intersecting the current remainder; no inversion closure of a previously cut remainder is assumed.

\subsection{4. Actual finite divisibility and the count of lambda}\label{actual-finite-divisibility-and-the-count-of-lambda}

Retain the common-orientation argument R049 section5/DO066(GC5--8). At rank one of an unmixed nilpotent G reduction, \(N=uv^t\) and \(N_{SS}=0\) force the whole x column to vanish at one of the two CM partners; the paired inverse supplies the orientation. Rank zero forces both. Thus there are at most \(4^{334}\) patterns. For a fixed pattern let J contain the chosen prime at rank one and both at rank zero, so \[
 N_EJ=P_1P_0^2,\qquad
 |a(h)|\le2Q^{-335(f-1)}(N_EJ)^{335/2},
 \quad Q=g_0g_1.
 \tag{RS10}
\] The factor two is the single global Euler correction; both mixed factors retain their modulus bound one.

Here the nonzero coefficient witness can be taken to be b itself. The determinant/minor argument R049 section2 gives valuation at least2-d for e1(N)=b at each CM component of a finite rank d=0 or1 reduction. The actual pole theorem K070 says b has only simple poles at the two CM primes over each type-III place. Since \[
 |\tau(b)|\le H+H^{-1}+2<136,\qquad 136^{10}<2^{71},
\] the paired norm argument gives \[
 N_EJ<136^fT^t.
 \tag{RS11}
\] This does not assume global unipotence.

The same pole theorem controls lambda itself. At every finite place, lambda is a root of \(X^2-(b+2)X+1\). If b is integral, lambda is a unit. At a type-III place, a negative valuation of lambda equals the valuation of b, and similarly for lambda\^{}-1. Hence at each such F-place its valuation belongs to \(\{-1,0,1\}\), and it is a unit elsewhere. For some subset D of the t places, \[
 \lambda\in I_D=\prod_{v\in D}\mathfrak q_v^{-1},
 \qquad N_F I_D=T^{-|D|}.
 \tag{RS12}
\] There are at most \(2^t\) choices. We count a superset of the allowed units in each ideal.

Use the ordinary trace Euclidean metric on F, with covolume \(\sqrt{D_F}N_F I_D\). Here \[
 D_F=2^fD_0,\quad
 D_0=7^{140}29^{144}5^{84}13^{84},\quad
 \operatorname{rd}(F)=\operatorname{rd}(E)<1500.
\] The codifferent minimum and Banaszczyk transference used in LIFT47/DO057 give \[
 \mu(I_D)\le\tfrac12\operatorname{rd}(F)\sqrt f\,
                 (N_F I_D)^{1/f}<750\sqrt f.
\] There is no complex half-trace factor in this real-field formula. Adjoining a Dirichlet cell to \(|\tau(\lambda)|\le133\) therefore enlarges the real ball to radius \(883\sqrt f\). For each D its count is at most \[
 U_F T^{|D|}\le U_F T^t,\qquad
 U_F=\frac{((22/7)f\,883^2)^{f/2}}{(f/2)!\sqrt{D_F}}
 <2^{1101}.
 \tag{RS13}
\] This retains the full real discriminant. We have not assumed an integral eigenbasis or that I\_D is principal.

\subsection{5. Counting the weighted original S column}\label{counting-the-weighted-original-s-column}

For fixed support, types, actual ordered pairs, G pattern and lambda, the target-overlattice containment of R051 section1 gives \[
 x\in L^{164},\qquad L=J I_\Sigma,\qquad
 N_E I_\Sigma=T^{-24},\quad N_E L=z^{2f}.
 \tag{RS14}
\] The generic containment is valid at all three T types; no unipotent no-pole lemma is used. From RS11 and the actual \(T>2^{995}\), \[
 \log_2z<Z_t:=\frac{168(71/10)-995(24-t)}{336}.
 \tag{RS15}
\] For t11 and t12 this is respectively \(-58711/1680\) and \(-2239/70\), both less than−31.

Apply to every complex embedding coordinate the real scalar \(\sqrt{|\tau(\lambda)|}\). By RS8 the transformed x lies in the ball of radius \(\sqrt{136f}\). The transformed lattice has covolume \[
 D_0^{164}z^{328f}|N_{F/\mathbb Q}\lambda|^{164}
 \ge D_0^{164}z^{328f},
 \tag{RS16}
\] by RS9. The already checked arbitrary-E-ideal transference bound DO070 gives \(\mu(L^{164})<4200\sqrt f\,z\). The transformation has operator norm at most \(\sqrt{133}<12\), so the transformed covering radius is below \(50400\sqrt f\,z\). The exact comparison \[
 \sqrt{136}+50400\,2^{-31}<12
\] shows that a Dirichlet-cell comparison yields \[
 \#x\le V_{164}(12)z^{-328f},\qquad
 V_{164}(12)=\frac{(528\cdot12^2)^{164f}}{(164f)!\,D_0^{164}}
 <2^{-182058}.
 \tag{RS17}
\] The real-linear weighted lattice is used only for its exact covolume and covering radius. It is not declared to be an arithmetic lattice of a new Hermitian model.

Multiplying RS17 by the loss in RS10 leaves \[
 V_{164}(12)T^{335s}z^{1176}
 <V_{164}(12)T^{335s}2^{1176Z_t}.
 \tag{RS18}
\] The exponent is positive, so using the upper bound on z is valid. For the other five complex coordinates retain the class-general full W0 count from DO075 section3, \[
                    V_5(223)T^{10s}<2^{5757}T^{10s},
\] and the full \(G_n<2^{5022}\) fibre for a completed negative column. This is a full complex-column count, not a real-part slice, so the zero-target coarea correction introduces no extra factor here.

\subsection{6. Complete counts, exact arithmetic and limits}\label{complete-counts-exact-arithmetic-and-limits}

For each t count the support by \(2^{168}\), its type-III positions by \(\binom{12}{t}\le2^{12}\), G patterns by \(2^{668}\), lambda ideals by \(2^t\), and inverse restriction by two. A type-III ordered pair has exact reciprocal weight and total pair/weight cost at most one. Each other active place retains the full pair bound \(P_T<2^{997}\), already including both non-III types. No further type multiplicity is used.

With the actual bounds \[
 L_T=995+1812/4096>\log_2T,\qquad
 L_Q=70+1123/4096<\log_2Q,
\] the complete exponent is \[
 \begin{split}
 E_t={}&168+12+668+5022+1+1+t+1101-182058+5757\\
 &+(345\cdot12+t)L_T-335\cdot167L_Q
       +1176Z_t+997(12-t).
 \end{split}
 \tag{RS19}
\] The two isolated ones are respectively the global Green constant and the inversion factor. The lambda count contributes \(T^t\), in addition to the original column powers; no determinant or class factor is cancelled.

The exact values are \[
 E_{11}=-180732011/20480<-8824,\qquad
 E_{12}=-109421831/20480<-5342.
 \tag{RS20}
\] Thus the two sectors together are less than \(2^{-5341}\). Their current-remainder intersections may be removed with this fresh charge, keeping all prior errors.

This result requires global rank(h-I)=2 and nonunipotence. The reciprocal line argument has not been extended to the residual global ranks three through six. It also does not address globally unipotent scalar-cross all-II terms or complete support12 exhaustion. The normalized signed lower bound, Mram/period value, actual Host and DP44 cost remain open.

The companion exact certificate checks actual primes, covolume constants and both volume ceilings, the weighted-cover comparison, every exponent and all stated multiplicities. It does not automate the rational primary decomposition, reciprocal-line proof or inversion of the actual kernel.

\subsubsection{Sources}\label{sources}

The original kernel/inversion identities are K062(AT7--8, AT14, AT23--24). The actual finite common-orientation and determinant-norm interfaces are DO066(GC5--8), R049 sections2 and5, K070, and R051 section1. The field metrics, arbitrary-ideal cover and full W0 count are the checked DO070, LIFT063 and DO075 statements.

Banaszczyk, W. (1993). New bounds in some transference theorems in the geometry of numbers. \emph{Mathematische Annalen, 296}, 625--635. \url{https://doi.org/10.1007/BF01445125}. Theorem2.2, p632, is used with the literal real-field trace dual, as in the accepted LIFT47 application; no new version of the transference theorem is asserted.

\clearpage

\section{AM-DO-077 --- higher finite ranks in the remaining nonunipotent scalar-cross sector}\label{am-do-077-higher-finite-ranks-in-the-remaining-nonunipotent-scalar-cross-sector}

\textbf{Canonical checked version --- 2026-11-22T17:30:16.087112+00:00.} Root read the whole final proof and complete script, directly checked the stated source interfaces, and ran a separate exact diagnostic. Original author-export wording below is historical. Promoted scope: Actual all-top H133 support12 nonunipotent scalar-cross t11/t12 with at least22 unmixed finite ranks\textgreater=3 has total absolute contribution\textless2\^{}-3947.

Author: /root/definite\_oscillator, 22 November 2026 UTC. New pending proof; no accepted file or ledger is changed.

\subsection{1. Exact statement}\label{exact-statement-1}

Keep the actual J56 normalized kernel, all-top H133 cutoff and all334 unmixed unipotent G reductions of K068. Put f=a168, n169, m164, E\^{}169=S perpendicular W0, with W0 signature(4,1). Restrict to A=h\_SS=I, h not fixing S, h globally nonunipotent, support s12, and t11 or t12 type-III places. No choice of a new projector, frame or real vector is made.

Let q\_v be the residue size g0 or g1 at an unmixed relative G factor and d\_v=rank(h-I mod v). The theorem is:

\[
 \#\{v:d_v\ge3\}\ge22
 \quad\Longrightarrow\quad
 \text{the total absolute contribution of this sector is}<2^{-3947}.
 \tag{HC1}
\]

This is a uniform whole-sector estimate, not a separate charge for each subset of22 places. Once DO076's global-rank-two nonunipotent sector is also removed, any nonunipotent scalar-cross term still surviving at support12 has global rank from3 through6 and at most21 finite ranks at least3. This conclusion neither excludes such matrices nor gives their signed contribution.

\subsection{2. Coefficient witness and the paired common ideal}\label{coefficient-witness-and-the-paired-common-ideal}

K069(IC4) supplies a positive fixed subspace of S of dimension at least163. Its Hermitian and Euclidean orthogonal complements coincide, so the unchanged operator bound H133 applies on a signature-(5,1) complement of dimension6. No moving basis is declared integral. In particular the global rank of h-I is at most6.

The literal dimension-six estimates DO067(ST2--6)/R048(RC48.5) are \[
 |e_1(h-I)|\le19020/133,\qquad
 |e_2(h-I)|\le164796/133<1280.
 \tag{HC2}
\] The accepted signature-one two-trace theorem implies that at least one coefficient is nonzero for a nonunipotent h. Choose such a coefficient b per matrix. This is not a counting choice or an extra factor.

At rank d=0 or1 of an unmixed reduction, determinant-one and the minor valuation argument of R049 section2 give valuation of b at least2-d at both CM partners. The actual K070 coefficient-pole theorem confines poles to order one at the two CM primes over each of the t type-III places. Thus the old common-orientation ideal J of DO066/R049, containing one prime at rank one and both at rank zero, satisfies \[
 N_EJ=P_1P_0^2<1280^fT^t.
 \tag{HC3}
\] Its norm and side convention are unchanged. The total number of patterns is at most4\^{}334; all ranks at least2 are one class in this pattern count.

The exact normalized Green estimates have L2=335 and L3=501. Therefore, if \(D_3=\#\{v:d_v\ge3\}\), \[
 |a(h)|\le
 2Q^{-335(f-1)}(N_EJ)^{335/2}
       g_0^{-166D_3},
 \qquad Q=g_0g_1,\quad g_0<g_1.
 \tag{HC4}
\] At ranks0/1 this uses the same conservative loss as R049. At every rank at least3 use the three-factor tail bound of K067(PS8--9), which is at least as strong for larger ranks. All factors used have indices at least166. Thus the single global Euler correction2 continues to suffice. The two mixed partial projectors retain only their modulus bound one.

HC4 is a pointwise improvement on every old pattern, so no new binomial or rank-assignment factor is inserted. For HC1 its final factor is uniformly at most \(g_0^{-3652}<2^{-124168}\), since the actual g0 exceeds2\^{}34.

\subsection{3. Full original-column count}\label{full-original-column-count}

For fixed support, actual ordered pencil pairs and common G pattern, \[
 x=(he_n)_S\in(JI_\Sigma)^{164},\qquad
 N_EI_\Sigma=T^{-24}.
\] This is the generic target-overlattice containment R051 section1, valid before global unipotence at every type. Let \(z=(N_E(JI_\Sigma))^{1/336}\). Since1280\^{}3\textless2\^{}31 and T\textgreater2\^{}995, \[
 \log_2z<Z_t=\frac{168(31/3)-995(24-t)}{336}.
 \tag{HC5}
\] For t11/t12 this is respectively \(-3733/112\) and \(-2551/84\). In particular both are below−30.

The positive part of a negative unit column has norm less than H/2=66.5 at every embedding. The corrected DO070 arbitrary-ideal covering radius is less than4200 sqrt(f)z. Hence a trace-ball Dirichlet-cell comparison gives \[
 \#x<V_{164}(67)z^{-328f},\qquad V_{164}(67)<2^{-45338}.
 \tag{HC6}
\] The complete weighted S count, after multiplying by the J loss in HC4, is therefore at most \[
 V_{164}(67)T^{335s}2^{1176Z_t}.
 \tag{HC7}
\] This uses the positive exponent335f−328f=1176. It requires no all-nonzero coordinate condition.

Retain the class-general full W0 column bound \(V_5(223)T^{10s}<2^{5757}T^{10s}\) of DO075 section3 and the full source-lattice stabilizer fibre \(G_n<2^{5022}\). This is a full five-complex-coordinate count for every target, including0, not the coarea slice; no C1 zero-target multiplier appears.

Count supports by2\^{}168, type-III positions by2\^{}12, all common G patterns by2\^{}668, and bound the selected real coefficient by1. The t off-diagonal ordered pairs cancel their exact reciprocal rho factors. Each other place retains the complete pair/weight bound \(P_T<2^{997}\), already covering both types I and II.

\subsection{4. Exact exponents and scope}\label{exact-exponents-and-scope}

Use the actual one-sided logarithms \[
 L_T=995+1812/4096,\qquad L_Q=70+1123/4096.
\] Before the last factor of HC4, the entire exponents are \[
 \begin{split}
 E_t={}&168+12+668+5022+1-45338+5757\\
 &+345\cdot12L_T-335\cdot167L_Q
       +1176Z_t+997(12-t),\\
 E_{11}={}&482236533/4096,\qquad
 E_{12}=492417141/4096 .
 \end{split}
 \tag{HC8}
\] These positive numbers express only failure of that intermediate upper estimate; they do not assert positive actual contributions.

For D3 at least22 the corrected exponents are \[
 E_{11}-124168=-26355595/4096<-6434,\qquad
 E_{12}-124168=-16174987/4096<-3948.
 \tag{HC9}
\] Their union is strictly below2\^{}(-3947), with no further subset multiplicity. Remove only its intersection with the current remainder, retaining every earlier error.

DO074 has already reduced nonunipotent scalar-cross support12 to t11/t12; this report does not establish that rigidity anew. DO076 separately removes global rank two in that nonunipotent sector. Their conjunction with HC1 gives the stated finite-rank restriction on the remaining global ranks3--6. The globally unipotent all-II branch, fixed-W0 rank-two sector, full signed inequality, actual period/Mram nonvanishing and DP44 conclusion are not addressed.

The source interfaces are K069(IC4), DO067(ST2--6), the direct dimension-six R048 coefficient table, K070, R049 sections2 and5, K067(PS8--9), corrected DO070, R051 section1 and DO075's class-general full W0 count. The new step is the uniform rank-three Green saving within the already counted common-ideal patterns. The companion certificate verifies the actual prime bounds, exact coefficient/radius inequalities and all complete exponents; it does not enumerate matrices or automatically establish these source interfaces.

\clearpage

\section{AM-R-056 --- balanced archimedean counting removes the remaining support-twelve nonunipotent scalar-cross sector}\label{am-r-056-balanced-archimedean-counting-removes-the-remaining-support-twelve-nonunipotent-scalar-cross-sector}

Author: /root, GPT-6 Astra Pro, ALP\_MACHINE. 22 November 2026 UTC. Status at creation: complete proof to be independently reviewed; no automatic promotion is implied by this file.

\subsection{1. Actual sector and conclusion}\label{actual-sector-and-conclusion}

Keep the actual J56 kernel, all finite projectors and ordered pencil incidences, the original Hermitian coordinates, the selected real vector and the charged all-top cutoff H=133. Put f=168, n=169, m=164, and write the original space as E\^{}169=S perpendicular W0, where S is positive of dimension164 and W0 has signature(4,1). Let N=h-I.

Consider the actual support-twelve scalar-cross terms with

\[
 A=h_{SS}=I,\qquad x=(he_n)_S\ne0,\qquad h\text{ globally nonunipotent},
 \qquad s=12,\quad t\in\{11,12\}.
 \tag{BM1}
\]

Here t counts the type-III unequal-pencil off-diagonal places. Every one of the334 unmixed G reductions is unipotent, as in the unchanged K068 residual. In this whole sector,

\[
 \sum |\text{actual normalized kernel contribution}|
 <2^{-7110}.
 \tag{BM2}
\]

The sum includes the full support, common-ideal, ordered-pair, original-column and source-fibre multiplicities. It imposes no bound on the number of finite ranks at least three and no condition that x and its inverse counterpart be proportional. Together with DO074, it removes every globally nonunipotent scalar-cross term at support12 from the current all-top remainder. This is not a claim that every support12 sector is removed.

The central improvement is to count in finitely many balanced real metrics. These metrics do not replace the arithmetic lattice, and they need not arise from algebraic units.

\subsection{2. Retained finite and original-coordinate inputs}\label{retained-finite-and-original-coordinate-inputs}

K069(IC4) gives a positive fixed subspace of S of dimension at least163. Its Hermitian and Euclidean orthogonal complements coincide. Thus the same cutoff H applies on a complement of dimension at most6, and rank(N)\textless=6. DO067(ST2--6), R048(RC48.5) and the accepted signature-one two-trace theorem give, for nonunipotent h, a nonzero coefficient b among e1(N),e2(N), with

\[
 |\tau b|<1280\quad\text{at all complex embeddings}.
 \tag{BM3}
\]

Indeed their literal dimension-six bounds are19020/133 and164796/133. Selecting a witness b for a matrix is not an additional counting choice.

For a fixed actual common-orientation pattern, let J be the integral E-ideal containing one CM prime at every unmixed rank-one factor and both primes at every rank-zero factor. At finite rank d=0 or1, the determinant-one/minor argument R049 §2 gives valuation of b at least2-d at both CM partners. K070 confines poles of b to order one at both partners over each of the t type-III places. Taking the absolute E-norm therefore gives

\[
 N_EJ<1280^fT^t.
 \tag{BM4}
\]

To see the square in this norm calculation explicitly: the paired divisibility of b contributes at least (N\_EJ)\^{}2, whereas its allowed paired poles contribute T\^{}\{2t\}; its archimedean absolute norm is below1280\^{}\{2f\}. The original common-orientation ideal, rather than a new coefficient ideal, remains in BM4.

The original-coordinate containment R051 §1 is

\[
 x\in L^{164},\qquad L=JI_\Sigma,\qquad
 N_EI_\Sigma=T^{-24},\qquad N_EL=z^{2f}.
 \tag{BM5}
\]

It holds for the actual generic target overlattice at all three pencil types, before imposing global unipotence. All coordinate-zero possibilities are allowed in the upper lattice count. The exact finite Green bound retained from DO066/R049 is

\[
 |a(h)|\le 2Q^{-335(f-1)}(N_EJ)^{335/2},
 \qquad Q=g_0g_1.
 \tag{BM6}
\]

The two mixed projectors have modulus at most one. There are at most4\^{}334=2\^{}668 common-orientation patterns. No separate rank-pattern count is added. BM3--6 and the complete generic count are also set out, with all old conventions retained, in DO077(HC2--8); its extra rank-three saving is not used here.

Since1280\^{}3\textless2\^{}31 and T\textgreater2\^{}995, BM4--5 give

\[
 \log_2z<Z_t=\frac{168(31/3)-995(24-t)}{336}
 =\begin{cases}-3733/112,&t=11,\\-2551/84,&t=12.\end{cases}
 \quad Z_t<-30.
 \tag{BM7}
\]

\subsection{3. Polar product bound and actual inversion}\label{polar-product-bound-and-actual-inversion}

Write h at a real label in its positive/negative decomposition. If its Cartan radius is r, the positive block and the two negative columns have the polar form

\[
 P_+=U(I+(\cosh r-1)vv^*),\qquad
 x_{\rm all}=\xi\sinh r\,Uv,\qquad
 x^-_{\rm all}=\xi'\sinh r\,v,
 \tag{BM8}
\]

where U is unitary, v is a positive unit vector and the two xi's have modulus one. These identities follow from the positive block's Gram identity; they are the polar calculation already proved in DO076 §3. It did not use the reciprocal-line hypothesis in deriving RS7.

Let P\_S denote the orthogonal projection to the original S. Because the S compression of P\_+ is I, the S compression of (\cosh r-1)Uvv* has operator norm at most2. Therefore, for r\textgreater0,

\[
 \|x_\tau(h)\|\,\|x_\tau(h^{-1})\|
 =\sinh^2r\,\|P_SUv\|\,\|P_Sv\|
 \le2(\cosh r+1)
 \le H+H^{-1}+2<136.
 \tag{BM9}
\]

The scalar-cross sector has r\textgreater0 and nonzero S column at every embedding. In particular the norms in BM9 are positive. If the S part of one negative column vanished, the unitary Gram identities with A=I would force the entire cross block to vanish; equivalently this is the scalar-cross/fixed-S dichotomy already used by K069.

K062(AT7--8, AT14, AT23--24) proves the exact inverse symmetry of the actual kernel: inverse transposes the ordered T-incidence matrix; F and C are unchanged; the G trace and selected real coefficient are conjugated; the cutoff is preserved. A=I, nonunipotence, s and t are also preserved. This symmetry is applied to the entire BM1 sector, before its intersection with any previously cut remainder.

Multiplying BM9 over the f real labels shows that at least one of h,h\^{}-1 satisfies

\[
                  \prod_\tau \|x_\tau(h)\|\le136^{f/2}.
 \tag{BM10}
\]

Because the absolute weights are equal under the same inverse bijection, the whole absolute sum is at most twice the sum restricted by BM10. Even when both elements satisfy BM10 this inequality holds. Its cost is exactly one binary bit, not one bit per real label. No proportionality x\^{}-=-lambda x and no count of a field eigenvalue is used.

\subsection{4. A product-one metric preserves the dual-lattice bound}\label{a-product-one-metric-preserves-the-dual-lattice-bound}

Use the literal half-trace Euclidean metric on the CM realization of E. For the fractional ideal L of BM5 its Euclidean dual is

\[
 L^\#=2\mathfrak D_E^{-1}\bar L^{-1}.
 \tag{BM11}
\]

Let D act by a positive real scalar D\_tau at each complex label, and suppose product\_tau D\_tau=1. The transformed Euclidean lattice is DL, with

\[
 (DL)^\#=D^{-1}L^\#,\qquad
 \operatorname{covol}(DL)=\operatorname{covol}(L)=D_0z^{2f},
 \quad D_0=7^{140}29^{144}5^{84}13^{84}.
 \tag{BM12}
\]

The determinant on the real realization is (product D\_tau)\^{}2=1. Neither DL nor its dual is asserted to be an E-ideal.

For a nonzero beta in the ideal L\#, the product of its complex absolute values is bounded below by the norm of that fractional ideal. In detail,

\[
 \prod_\tau|\tau\beta|
 =|N_{E/\mathbf Q}\beta|^{1/2}
 \ge 2^f|\operatorname{disc}E|^{-1/2}z^{-f}.
 \tag{BM13}
\]

Multiplication by D\^{}-1 leaves that product unchanged. Applying AM--GM to the f squared complex absolute values consequently gives, for the transformed lattice itself,

\[
 \lambda_1((DL)^\#)\ge
            \frac{2\sqrt f}{\operatorname{rd}(E)z}.
 \tag{BM14}
\]

This is uniform in every D\_tau; no factor involving max D\_tau is needed. The proof uses the algebraic norm of the original beta, followed by a real product-one transformation.

The finite-dimensional transference theorem of Miller and Stephens-Davidowitz (2019), Theorem3.2, equation(3.3), printed p.8, applies to every Euclidean lattice. In real dimension336, its upper constant is less than72, as verified in DO070(FT3--5). Since rd(E)\textless1500,

\[
 \mu(DL)<\frac{2250}{7}\sqrt f\,z,\qquad
 \mu((DL)^{164})<4200\sqrt f\,z.
 \tag{BM15}
\]

Here the second inequality uses the orthogonal product metric and sqrt164\textless13. These are direct applications to the transformed Euclidean lattice, not estimates obtained by multiplying the untransformed radius by an operator norm.

\subsection{5. A finite cover by balanced metrics}\label{a-finite-cover-by-balanced-metrics}

Put r\_tau=\textbar\textbar x\_tau\textbar\textbar. The positive part of a negative unit column under the unchanged H cutoff has norm less than H/2=66.5, so

\[
                         0<r_\tau<67<2^7.
 \tag{BM16}
\]

Choose any nonzero original E-coordinate b of x. It is nonzero at every field embedding. Since b belongs to L, its principal ideal is contained in L, and J is integral. Thus

\[
 \prod_\sigma|\sigma b|
 =|N_{E/\mathbf Q}b|^{1/2}
 \ge(N_EL)^{1/2}\ge T^{-12}.
 \tag{BM17}
\]

Each of the other167 complex absolute values is less than67. Hence for every label tau,

\[
 r_\tau\ge|\tau b|\ge T^{-12}/67^{167}
 >2^{-12\cdot996-167\cdot7}=2^{-13121}.
 \tag{BM18}
\]

The numerical bound is the same for whichever nonzero original coordinate is chosen, so this argument introduces no coordinate-choice multiplicity.

For each tau form the integer bin

\[
 j_\tau=\lfloor16\log_2r_\tau\rfloor,\qquad q_\tau=2^{j_\tau/16}.
 \tag{BM19}
\]

The open interval in BM16,18 admits at most16(13121+7)+1=210049\textless2\^{}18 such integers. There are consequently fewer than2\^{}(18f)=2\^{}3024 vectors of bins. For one fixed bin vector choose the fixed real metric

\[
             D_\tau=\frac{(\prod_\sigma q_\sigma)^{1/f}}{q_\tau}.
 \tag{BM20}
\]

It has product one and is the same for all points in that chart. From q\_tau\textless=r\_tau\textless2\^{}(1/16)q\_tau and BM10,

\[
 D_\tau r_\tau
 <2^{1/16}(\prod_\sigma q_\sigma)^{1/f}
 \le 2^{1/16}(\prod_\sigma r_\sigma)^{1/f}
 \le 2^{1/16}\sqrt{136}.
 \tag{BM21}
\]

Every chart's transformed points therefore lie in the real Euclidean ball of radius2\^{}(1/16)sqrt(136f) in dimension2fm. BM7 and BM15 show that adjoining a Dirichlet cell of (DL)\^{}164 enlarges this to a ball of radius strictly less than13sqrt f.~For an entirely rational verification, the certificate checks

\[
 273^{16}>2\cdot256^{16},\qquad
 136(273/256)^2<(13-4200\cdot2^{-30})^2.
 \tag{BM22}
\]

The usual disjoint-cell volume comparison, BM12 and pi\textless22/7 now give for each fixed chart

\[
 \#\{x\text{ in that chart}\}
 < V_{164}(13)z^{-328f},\qquad
 V_m(C)=\frac{(528C^2)^{fm}}{(fm)!\,D_0^m}.
 \tag{BM23}
\]

The formula uses (22/7)f=528. It is an upper-envelope volume, not an equality with the actual ball volume. Full lattice points may be counted even when some coordinates vanish; x=0 is excluded in the actual sector and merely included as an innocuous phantom in the upper volume count. There is no division by a coordinate of x in this cell argument.

The exact integer volume comparison is

\[
                   2^{-175695}<V_{164}(13)<2^{-175694}.
 \tag{BM24}
\]

Combining all charts costs at most2\^{}3024, uniformly over the original common-ideal patterns and all actual ordered pencil pairs.

\subsection{6. Complete kernel accounting}\label{complete-kernel-accounting}

Multiplying BM23 by the J factor in BM6 gives

\[
 (N_EJ)^{335/2}z^{-328f}
 =T^{335s}z^{7f}
 =T^{335s}z^{1176}
 <T^{335s}2^{1176Z_t}.
 \tag{BM25}
\]

The exponent of z is positive, so use of its upper bound has the correct direction. All earlier field, side and norm conventions remain unchanged.

For the remaining W0 coordinates use the full original five-complex-coordinate count DO075 §3:

\[
 V_5(223)T^{10s}<2^{5757}T^{10s}.
 \tag{BM26}
\]

It is uniform in every target, including the zero targets. No real-part coarea slice is used, and no separate zero-target correction is omitted. The full source-lattice stabilizer fibre remains G\_n\textless2\^{}5022; the rank-five unipotent improvement is not available in this argument. The real normalized coefficient has modulus at most one.

The remaining factors are as follows.

\begin{longtable}[]{@{}
  >{\raggedright\arraybackslash}p{(\columnwidth - 2\tabcolsep) * \real{0.5000}}
  >{\raggedright\arraybackslash}p{(\columnwidth - 2\tabcolsep) * \real{0.5000}}@{}}
\toprule\noalign{}
\begin{minipage}[b]{\linewidth}\raggedright
Actual item
\end{minipage} & \begin{minipage}[b]{\linewidth}\raggedright
Retained upper bound or exact cancellation
\end{minipage} \\
\midrule\noalign{}
\endhead
\bottomrule\noalign{}
\endlastfoot
Support choices & 2\^{}168 \\
Type-III positions among12 places & 2\^{}12 \\
Common G orientation patterns & 2\^{}668 \\
Full source stabilizer fibre & 2\^{}5022 \\
Single Green Euler constant & 2 \\
Inverse restriction BM10 & 2 \\
All balanced charts & 2\^{}3024 \\
Full S volume & 2\^{}-175694 \\
Full W0 volume & 2\^{}5757 \\
Type-III off-diagonal ordered pairs & Their counts cancel the exact reciprocal rho weights \\
Each remaining actual pair/weight sum & P\_T=2(T-1)\^{}2/(T-2)\textless2\^{}997, covering both types I and II \\
Selected real coefficient and two mixed projectors & Each modulus at most1 \\
\end{longtable}

The pair cancellation is the actual ordered-incidence identity used in DO077(HC8), not a replacement by an arbitrary local average. These multiplicities count a larger set where needed for an upper bound; every actual term in BM1 is included.

Use the exact one-sided logarithms

\[
 L_T=995+1812/4096=1019333/1024,\qquad
 L_Q=70+1123/4096=287843/4096,
 \tag{BM27}
\]

where log2 TL\_Q for the actual J56 primes. The full exponents are

\[
 \begin{aligned}
 E_t={}&168+12+668+5022+1+1+3024-175694+5757\\
       &+345\cdot12L_T-335\cdot167L_Q
          +1176Z_t+997(12-t),\\
 E_{11}={}&-39311243/4096,\qquad
 E_{12}=-29130635/4096.
 \end{aligned}
 \tag{BM28}
\]

Each is strictly less than-7111. Adding these two sector bounds therefore proves BM2. As a cross-check, replacing the unbalanced S-volume exponent-45338 in the complete DO077(HC8) envelope by-175694, and adding3024+1, gives exactly the same two fractions.

The one added charge is7110. If DO077 has already been removed with charge3947, remove only the intersection of BM1 with the subsequent remainder, retaining that earlier charge and every other historical charge. The full larger inverse-stable sector was bounded first, so the remainder itself need not be inverse-stable. The existing uniform count is unchanged.

\subsection{7. Consequences and precise limits}\label{consequences-and-precise-limits}

DO074 forces every nonunipotent scalar-cross support12 term into t=11 or12; BM2 bounds both whole sectors, at every allowed global rank. The previous rank-two reciprocal-line cut DO076 and high-finite-rank cut DO077 remain valid historical refinements, but their leftover nonunipotent sector is now removed. This is a strict reduction of the actual signed-kernel gate used in Q1 and Q2.

No claim is made here about the all-II globally unipotent sector, the fixed-W0 rank-two sector, any other unremoved structural sector, higher support, or the overall signed sufficient inequality. The prescribed period/Mram, arithmetic Read, actual mixed-CM family, integer DP44 cost saving, literal032 cycle value, Kummer-Hecke RH, and global oscillator charge comparison remain separate open gates.

\subsection{8. Algebraic-geometric meaning and source locators}\label{algebraic-geometric-meaning-and-source-locators}

The integral module and its common congruence ideal are unchanged. The paired divisibility is read in the original integral Hermitian model at its G closed points; the coefficient witness is an exterior-power/determinantal function of h-I. The new charts live only on the real realization of that fixed module, with determinant-one metric change. No isomorphism of integral schemes, change of polarization, replacement of a Picard functor, lambda-ring identity, or new l-adic weight or motivic positivity theorem is asserted.

The carried module/stabilizer and exterior-power contexts are SGA3, I,4.7.3/4.7.3.1, and SGA6,VI,Theorem3.3. The polarized-Jacobian, Picard and biextension contexts remain FGA,Bourbaki232,Theorem3.1,pp.148--149; Bourbaki236,Corollaries1.7/2.3,pp.229/232; and SGA7 I,VII,2.7/2.9,pp.169--174. The formal-crystal and Cartier-dual locators carried from the intakes are Grothendieck archive cote4,scan003,§§8.1.1--8.1.2,scans005--006,§§8.4.1--8.4.3, and cote54,scans005--007. These are contextual locators already recorded in the dictionary, not sources of the new analytic counting lemma.

For our programme, this narrows the actual finite/archimedean kernel gate of DP\_MULT Q3e and THETA3 G1 while retaining A7's definite-CM-type requirement. It yields no new implication for CX\_GRH CX-ALG-024 or GCT\_PLETH T-052 without their distinct arithmetic maps. In CM-motive language the realization being counted has the same underlying rational and integral data; an analytic change of metric alone is not a motivic comparison.

\subsubsection{Cited primary result and internal dependencies}\label{cited-primary-result-and-internal-dependencies}

Miller, S. D., \& Stephens-Davidowitz, N. (2019). Kissing numbers and transference theorems from generalized tail bounds (arXiv:1802.05708v2). \url{https://arxiv.org/pdf/1802.05708}. Theorem3.2,equation(3.3),printed p.8; Euclidean dual and covering-radius definitions in §1. The required finite-dimensional form is carried through the directly checked DO070(FT3--5).

Internal exact locators: K062(AT7--8,AT14,AT23--24); K068's actual J56 residual; K069(IC4); DO067(ST2--6); R048(RC48.5); R049 §§2,5; K070's coefficient-pole theorem; R051 §1; DO066(GC5--8); corrected DO070(FT3--6); DO074's trace-conditioned rigidity; DO075 §3; DO076 §3(RS7); and DO077(HC2--8). The companion certificate checks actual integer inputs, all rational bounds, the volume bracket, chart count, full exponents and sector union. It does not automate the geometric or kernel identities.

\clearpage

\section{AM-K088 --- integral square pairings remove the remaining all-II scalar-cross sector}\label{am-k088-integral-square-pairings-remove-the-remaining-all-ii-scalar-cross-sector}

\textbf{Canonical checked version --- 2026-11-22T17:30:16.087112+00:00.} Root read the whole final proof and complete script, directly checked the stated source interfaces, and ran a separate exact diagnostic. Original author-export wording below is historical. Promoted scope: The actual all-II support12 globally rank-two-unipotent scalar-cross sector with total finite G rank\textless=69 is empty. Together with K087 this exhausts that all-II sector using its existing1445 charge once.

Author: /root/kummer\_extension. New pending proof, 22 November 2026 UTC. This file changes no accepted proof, coefficient, or ledger.

\subsection{1. Exact statement and retained family}\label{exact-statement-and-retained-family}

Use the actual all-top J56 kernel, H=133, the original standard coordinates and the same central representatives, projectors, selected real trace, pencil labels and class masses as K086--K087. Put \[
 f=a=168,\quad n=169,\quad m=164,\quad
 E^{169}=S\perp W_0,\quad \operatorname{sig}(W_0)=(4,1).
 \tag{IP1}
\] Consider the support-twelve sector in which all twelve nonintegral relative T places are type II. Let \[
 h=I+N,\quad h_{SS}=I,\quad N^3=0,\quad
 \operatorname{rank}_E N=2,\quad \operatorname{rank}_E N^2=1,
 \quad (he_n)_S\ne0.
 \tag{IP2}
\] These are the globally rank-two-unipotent scalar-cross terms left outside K086. At each of the334 unmixed relative G places put d\_v=rank(N mod v). Thus d\_v belongs to\{0,1,2\}; paired inverse-adjoint reductions have the same rank.

\textbf{Theorem.} No actual term satisfying IP1--2 and \[
                         \sum_{v\ {\rm unmixed}}d_v\le69
 \tag{IP3}
\] exists. Consequently the complete all-II support-twelve sector has absolute contribution less than2\^{}-1445, by K087's already checked bound on the complementary finite-rank sector. If that charge is already retained, this finding adds no numerical error. Remove only the intersection with the current remainder and preserve all earlier errors.

This is an exact emptiness statement for the low-finite-rank part. It does not prove full support-twelve exhaustion, a signed lower bound, an automorphic period, or an arithmetic cost estimate.

The actual constants from LIFT41 are \[
\begin{aligned}
 g_0&=22285932761,&g_1&=64072175441,&Q&=g_0g_1,\\
 U&=77Q^2-14,&W&=Q^2U,&R&=H+H^{-1}-2=17424/133,\\
 u&=(2,95192390527,383092671043,12529843548438452388609),\qquad
 w=Qu.
\end{aligned}
\tag{IP4}
\] Here u dot u=U. Regard w as a vector in the four positive original coordinates of W0. The fixed packet vector is v0=(w,pi\_k), with its other positive coordinates zero. In particular v0*Jv0=-T, as in LIFT41. No new four-square representation or new packet is chosen.

\subsection{2. The actual stable selfdual lattice and its integral test vectors}\label{the-actual-stable-selfdual-lattice-and-its-integral-test-vectors}

For a type-II place the actual unique admitted incidence is fixed: hM\_c=M\_c.~At its selected CM component, LIFT44(PR4--7) and K062's actual pencil read \[
 M_c=(I+c\pi^{-1}zq)\Lambda,\quad
 z=v_0/\bar\pi,\quad qz=0,\quad
 q=-v_0^*J-\pi m_0,
 \tag{IP5}
\] where z and the row m0 are integral there. At the opposite CM component use the Hermitian dual lattice, not a second independently chosen shear. At every other finite place use the standard lattice. These local modules patch to a full projective O\_E-lattice M. The actual paired construction makes M selfdual for the original Hermitian form, and hM=M. One can define M as the intersection of these local conditions in E\^{}169; only finitely many differ from the standard lattice. No integral basis of M is chosen.

Let y be any original integral vector with negative coordinate zero, supported on the four positive coordinates of W0, and satisfying w*y=0. It is orthogonal to v0. At the selected T component, IP5 gives qy=-pi m0y, so \[
 (I-c\pi^{-1}zq)y=y+czm_0y\in\Lambda.
 \tag{IP6}
\] Thus y belongs to every admitted M\_c.~At the conjugate component, every selected-place vector of M\_c lies in Lambda+pi\^{}-1 O z; the integral vector y pairs integrally with Lambda and pairs as zero with z. Hence its conjugate component belongs to the exact dual lattice as well. This proves y in M at both CM partners. At the remaining places it is already standard integral. This is the same fixed-complement inclusion used in LIFT44(PR8--9), but the conclusion needed here is membership in the source and target lattice, not integrality of h y in the standard coordinates.

For each such y set \[
                       b_y=\langle y,N^2y\rangle_J.
 \tag{IP7}
\] Since N maps M into M, N\^{}2y belongs to M. The selfdual integral pairing on M gives b\_y in O\_E. The unipotent unitary identity \[
 N^\dagger=-N+N^2,\qquad (N^2)^\dagger=N^2
 \tag{IP8}
\] makes b\_y real under CM conjugation. Consequently \[
                              b_y\in O_F.
 \tag{IP9}
\] In particular this scalar has no T denominator. We do not claim that N\^{}2y, or an arbitrary matrix entry of N\^{}2, is standard integral at T. This particular stable-lattice proof uses all-II fixed incidence: a general arrow M\_c to M\_d would not supply an h-stable lattice M. No claim of a converse or of failure of a different pairing argument is made.

\subsection{3. Paired G divisibility and the exact short kernel basis}\label{paired-g-divisibility-and-the-exact-short-kernel-basis}

Define the F-ideal \[
 \mathfrak J_F=\prod_{v\ {\rm unmixed\ }G}\mathfrak p_v^{\,2-d_v},
 \qquad
 N_F\mathfrak J_F
   =\frac{Q^{334}}{\prod_v q_v^{d_v}}
   \ge\frac{Q^{334}}{g_1^{69}}.
 \tag{IP10}
\] There are167 actual places of each residue size g0 and g1. At d\_v=0, N is divisible by the selected prime, hence N\^{}2 by its square. At d\_v=1, the finite reduction of N is rank-one nilpotent and has square zero. The same statements hold at the CM partner: its matrix is the inverse-adjoint reduction, still nilpotent of the same rank. Therefore every entry of N\^{}2, and hence the scalar b\_y, is divisible to order at least2-d\_v at both CM components. Since b\_y lies in F, this gives exactly the F-ideal in IP10, not its square. For b\_y nonzero, \[
                    |N_{F/\mathbb Q}b_y|\ge N_F\mathfrak J_F.
 \tag{IP11}
\]

Let alpha=(N\^{}2)\_nn.~The carried signature-one rank-two-unipotent identities, proved directly in R052 section2, give \[
 \alpha\in F,\quad0<\tau(\alpha)\le R\quad(\tau:F\hookrightarrow\mathbb R),
 \qquad N^2=-\alpha vv^*J,
 \tag{IP12}
\] where v spans the isotropic E-line imN\^{}2 and is normalized to have negative coordinate1. Its positive coordinate vector has Euclidean norm1 at every CM embedding. Such a normalization is legitimate because a nonzero isotropic vector has nonzero negative coordinate. Moreover N\_SS=0 and N\^{}dagger\_SS=0, so IP8 gives (N\^{}2)\_SS=0. In the rank-one form IP12 this block is -alpha v\_S v\_S*, forcing v\_S=0. Thus im N\^{}2 lies in W0 before the three additional orthogonal equations are imposed. This is a rational line calculation, not an integral change of frame. For the positive test vectors above, \[
 b_y=-\alpha\,|v^*y|^2,
 \qquad |\tau(b_y)|\le R\|y\|^2.
 \tag{IP13}
\]

Use the following three explicit integer vectors, embedded into the original four positive coordinates: \[
\begin{split}
 y_1&=(-976627,-153032,38026,0),\\
 y_2&=(-567678,695570,-172838,0),\\
 y_3&=(104583,-15309315176,-61610045224,2).
\end{split}
\tag{IP14}
\] They are orthogonal to u, linearly independent, and hence span w-perp inside that four-dimensional E-space. Their squared norms are \[
 (L_1,L_2,L_3)=(978665066829,\ 835948910828,\ 4030172803672366725045).
 \tag{IP15}
\] For reproducibility these rows are the product of the integer matrix \[
 \begin{pmatrix}
 -76516&19013&0\\
 347785&-86419&0\\
 -7654657588&-30805022612&1
 \end{pmatrix}
 \tag{IP16}
\] with the rows (-u2,2,0,0),(-u3,0,2,0),(-u4,0,0,2). The matrix in IP16 has determinant-1. No reduced-basis theorem or floating-point claim is used: the certificate checks these literal integer identities.

For each j=1,2,3 the exact integers satisfy \[
                 Q^{334}133^{168}
                    >g_1^{69}(17424L_j)^{168}.
 \tag{IP17}
\] Combining IP11, IP13 and IP17 forces b\_yj=0. At a faithful complex embedding alpha is strictly positive, so IP13 forces v*y\_j=0. The rational orthogonal equations then imply \[
                        \operatorname{im}N^2
                              \subset\operatorname{span}_E(w,e_n).
 \tag{IP18}
\] This argument does not claim that h fixes the test vectors, nor that h acts only on this plane. Only the image of N\^{}2 has been confined.

\subsection{4. An original matrix entry forces a new numerator divisor}\label{an-original-matrix-entry-forces-a-new-numerator-divisor}

Because the line in IP18 is isotropic, write its normalized generator as \[
                           v=(a w,1),\qquad a\bar a=1/W.
 \tag{IP19}
\] The normalization is in E. No square root of W is assumed to belong to E, and no anisotropy assertion about this plane is needed. Since the first positive coordinate of w is2Q, IP12 gives the exact original-coordinate identity \[
                   (N^2)_{11}=-\alpha\frac{4Q^2}{W}
                                      =-\frac{4\alpha}{U}.
 \tag{IP20}
\] The actual short-sum matrices h and h\^{}-1 are standard integral away T, including at primes ramified in the host and at all primes dividing U. Therefore N\^{}2 is standard integral away T. The integer U is odd and coprime to Q and T; the certificate verifies these gcd assertions for the actual constants.

At any F-prime dividing U, every prime of E above it has4 a unit. Applying its valuation to IP20 gives v(alpha)\textgreater=v(U), with the relative ramification index cancelling on the two F-elements. At all other places away T, alpha is integral because it is the original nn entry of N\^{}2. At primes above2, use this last fact and the oddness of U; no factor4 denominator is introduced. Thus alpha/U is integral at every finite place away T.

At an active type-II place, h=H\_c k H\_c\^{}-1 with H\_c=I+c pi\^{}-1 Z and k integral invertible. The two-shear expansion shows that h and h\^{}-1 have entry poles of order at most two, at both CM components. IP8 equivalently gives N\^{}2=h+h\^{}-1-2I. Hence alpha/U has pole at most two at each of the twelve relative T places. Let \[
                 I_T=\prod_{v\ {\rm active}\ T}\mathfrak p_v^{-1}
                 \quad\text{in }F,\qquad N_FI_T=T^{-12}.
\] We have proved alpha in U I\_T\^{}2. At every unmixed G prime, the same square valuation used in IP10 applies directly to alpha, and U,T are units there. Combining these disjoint valuations gives the full fractional-ideal membership \[
                         \boxed{\alpha\in U\mathfrak J_F I_T^2.}
 \tag{IP21}
\] There is no selected-half norm or missing CM partner here. IP21 is an F-ideal statement; it follows from both actual CM-component matrix conditions. No rational moving-frame denominator appears in IP20 or IP21.

\subsection{5. Exact contradiction and error accounting}\label{exact-contradiction-and-error-accounting}

The nonzero positive alpha in IP12 and IP21 satisfies \[
 R^{168}\ge N_{F/\mathbb Q}\alpha
 \ge U^{168}N_F\mathfrak J_F\,T^{-24}
 \ge\frac{U^{168}Q^{334}}{g_1^{69}T^{24}}.
 \tag{IP22}
\] The actual integers satisfy the strict opposite inequality \[
                  (133U)^{168}Q^{334}
                    >17424^{168}T^{24}g_1^{69}.
 \tag{IP23}
\] This contradiction proves the theorem. The certificate checks IP17 and IP23 directly using integers, not logarithmic approximations. Conservative margins7754 and20586 binary bits, respectively, are also checked; these margins are diagnostics, not extra hypotheses.

K087 already bounded the full complementary sector sum d\_v\textgreater=70 by2\^{}-1445. Its exact original-coordinate count retains every finite orientation, source/target pair, class fibre, period normalization and actual real factor. This finding does not change that count. It proves that the complementary low-rank set is empty. Thus the all-II scalar-cross support-twelve sector is exhausted with the existing1445 charge. In a cumulative remainder formula remove only the actual intersection and do not add1445 a second time if K087 has already supplied it.

The argument applies to the stated fixed packet and all-top H133 cutoff. It does not transfer this bound to the old H8192 family, to a modified packet, to arbitrary relative pencil arrows, or to the fixed-W0/global-rank-one sectors. It does not infer a final signed or period nonvanishing statement from positive local incidence weights.

\subsection{6. Verification scope and precise dependencies}\label{verification-scope-and-precise-dependencies}

The actual packet, integer u and standard coordinate normalization are LIFT41(NP1--9). The exact pencil, its dual and the fixed-complement membership are LIFT44(PR4--9), retained in K062's full all-base assembly. Type-II means the unique fixed admitted lattice in K070(PT1--4), with all paired local conditions. The original away-T integrality and determinant-one representative are the unchanged K068/K069 setting. The global unipotent square identity and alpha bound are R052 section2, and K086 retains the actual selected metric. K087 supplies only its complementary finite-rank count. The proof above establishes all new selfdual-pairing, short-kernel, plane, numerator-ideal and contradiction steps.

The companion certificate checks the literal field/prime record, integer kernel transformation, norms, pair exponents, gcd conditions and all strict integer norm comparisons. It does not replace the local lattice or Hermitian-line proofs, does not enumerate rational matrices, and performs no Lean build.

\subsection{Author whole-proof verification}\label{author-whole-proof-verification-4}

At 2026-11-22T17:25:41.575092+00:00, I completed a whole read of all six sections, IP1--23, every scope/source paragraph and the complete companion certificate. I reconstructed the actual two-partner membership of each original test vector, the projective selfdual stable lattice, the F-valued integral square pairings, rank-one nilpotent valuations at both CM partners, the explicit integer kernel and every norm inequality. I checked the required preliminary v\_S=0 consequence before plane confinement, and checked the original11-entry identity at primes dividing U, at2, at both T partners and at G. I reread the final precision edits: the stable-lattice proof is scoped to all-II, without asserting that every other pairing method fails.

The exact arithmetic certificate passed34 assertions at 2026-11-22T17:25:04.378257+00:00. Its paired-test norm margins have integer binary floors13121,13159,7755 and its final contradiction margin has floor20587. It does not establish the geometric interfaces automatically. The proof makes no anisotropy claim and uses no square root of W, no integral moving frame, no new theta input, no extra numerical error, and no final signed nonvanishing conclusion. Author status: \textbf{whole-proof and complete-script PASS; frozen for independent review}.

\clearpage

\section{AM-LIFT-067 --- the literal positive-rank T character and translated spherical test}\label{am-lift-067-the-literal-positive-rank-t-character-and-translated-spherical-test}

\textbf{Canonical checked version --- 2026-11-22T17:30:16.087112+00:00.} Root read the whole final proof and complete script, directly checked the stated source interfaces, and ran a separate exact diagnostic. Original author-export wording below is historical. Promoted scope: The literal positive-rank transported T character has final coefficients(bar(pi),pi). Standardization fixes the exact D*tau spherical test and N-Haar modulus q\footnote{-rho(k+rho+1)}. The altered unit-character/bare-tau test vanishes; actual local value and global period remain open.

Author: \nolinkurl{/root/lift_accounting}. 22 November 2026 UTC.

Status: complete scoped local argument; author verification and exact diagnostic replay are recorded separately. This identifies the actual character and vector that a positive-rank local test must use. It does not evaluate that test or prove a global period.

\subsection{1. Literal setting and conventions}\label{literal-setting-and-conventions}

Let \nolinkurl{rho>=1}, \nolinkurl{k>=1}, and \nolinkurl{r=k+2rho+1}. Work at one actual degree-one split T place, with local field \nolinkurl{F}, ring \nolinkurl{O}, residue cardinality \nolinkurl{q}, and uniformizer \nolinkurl{pi}. Write \nolinkurl{c=bar(pi)}, a unit at this selected half, so the rational norm is \nolinkurl{T=pi c}. The additive character is the literal trace character divided by the fixed G modulus. Here G and the different are units, so the induced character \nolinkurl{psi} has conductor O. Any residual unit coefficient is included in \nolinkurl{psi}; it is not discarded.

Use the ordered split coordinates

\[
 (e_1,\ldots,e_\rho,p_1,\ldots,p_k,z_*,f_\rho,\ldots,f_1),
 \tag{TC1}
\]

with Gram matrix J pairing the displayed e and f vectors, positive identity on the p block, and -1 on z. Thus the other CM component of a selected matrix A is exactly \texttt{J\ A\^{}\{-t\}\ J}. At T the scaled flags of LIFT060 HF3 differ from these coordinates by units. They may be used literally, with those unit coefficients retained; the valuation conclusions and the displayed identities after that integral coordinate transport are unchanged.

Let N be the Bessel unipotent subgroup fixing the isotropic flag and acting trivially on its perpendicular quotient. In the selected coordinates it consists of upper triangular block matrices with identity core block of size \nolinkurl{k+1}, upper unitriangular e and f blocks, and arbitrary e-to-core, core-to-f and e-to-f blocks. This is the direct split reading of \nolinkurl{HidaSetting.Nfin}, lines 105--111, using both CM components. H is GLk on the p block.

For the unit marked vector z, the final two terms of the actual root character are

\[
 n_{e_\rho,z}+n_{z,f_\rho}.
 \tag{TC2}
\]

The preceding e-chain and f-chain terms, with their fixed signs and units, are unchanged below. Formula TC2 is LIFT060 HF12 without its separate generic H character. It also follows directly from \nolinkurl{psiFin}, lines 113--115: the plus root is the e-to-z entry, and the minus root is \texttt{-(n\^{}\{-1\})\_\{z,f\_rho\}=n\_\{z,f\_rho\}}. There are no intervening core entries in N that alter this inverse entry.

\subsection{2. The norm-minus-T vector changes both root coefficients}\label{the-norm-minus-t-vector-changes-both-root-coefficients}

Put \nolinkurl{A=(T+1)/2}, \nolinkurl{B=(T-1)/2} and use the actual rational move, on the p1,z plane,

\[
 g=\begin{pmatrix}A/\pi&B/c\\ B/\pi&A/c\end{pmatrix},
 \qquad z_T=Bp_1+Az_*=g(cz_*).
 \tag{TC3}
\]

It fixes every e and f vector. Let \texttt{N\_T=gNg\^{}\{-1\}} and let \nolinkurl{chi_T} denote the literal character using z\_T. The Hermitian form in the source is conjugate-linear in its first variable and linear in its second. Therefore

\[
 h(f_\rho,gng^{-1}z_T)=c\,h(f_\rho,nz_*).
 \tag{TC4}
\]

At the conjugate CM component the same global c is pi. Consequently the transported character \texttt{chi\_c(n)=chi\_T(gng\^{}\{-1\})} has final terms

\[
 \boxed{\ c\,n_{e_\rho,z}+\pi\,n_{z,f_\rho}\ }.
 \tag{TC5}
\]

Their valuations at the selected T half are exactly 0 and 1. All earlier root terms are unchanged. If the bilinear period uses the inverse character, invert the entire character on both sides; the same valuations and conclusions hold. It is incorrect to give both final root coefficients valuation zero while retaining this rational move and the physical z\_T.

The exact source anchors are \nolinkurl{HidaSetting.hJ} at 24, \nolinkurl{rootCoord} at 39--41, \nolinkurl{hA} at 81--82, and \nolinkurl{psiFin} at 113--115. The actual marked vector and rational move are LIFT060 HF5--6 and LIFT30 NL3--9. These sources are used as definitions, not as an assertion that an abstract local model has been identified with a global period.

\subsection{3. Standardizing the character fixes an additional diagonal}\label{standardizing-the-character-fixes-an-additional-diagonal}

Define the selected split matrix

\[
 D=\operatorname{diag}(cI_\rho,I_{k+1},\pi^{-1}I_\rho).
 \tag{TC6}
\]

It is the selected component of an actual rational unitary diagonal: the e and paired f scalars satisfy \texttt{bar(c)\ pi\^{}\{-1\}=1}. It fixes the entire core, commutes with H, normalizes N, and commutes with any lower core translate

\[
 \tau_u=I+\pi^{-1}\sum_{a=1}^k u_a E_{z,p_a}.
 \tag{TC7}
\]

Here u is the literal primitive T covector. Since the ratio of the D entries across the two final roots is c and pi, respectively,

\[
 \chi_c(n)=\chi_1(DnD^{-1}),
 \tag{TC8}
\]

where chi1 is the corresponding unit-marked-vector character, with every preceding chain coefficient unchanged. Hence for any linear Bessel functional \nolinkurl{ell1} with character chi1, the functional \nolinkurl{ellc(v)=ell1(Pi(D)v)} has character chic. This is an exact algebraic transport and preserves the H factor because D commutes with H.

LIFT30 gives \texttt{g\^{}\{-1\}=tau\_c\ k} on the core, with \nolinkurl{k in K} and \nolinkurl{u=c e1*}; the integral covector transport NL10--12 gives TC7 for the actual u. Thus, after character standardization, the unchanged prescribed spherical input is

\[
                    \Pi(D\tau_u)v_K.
 \tag{TC9}
\]

It is not the bare vector \nolinkurl{Pi(tau_u)v_K}. Equation TC9 is a coordinate description of the old datum, not permission to replace it by a freely chosen vector.

There is also an exact Haar factor. For the self-map \texttt{n\ -\textgreater{}\ DnD\^{}\{-1\}}, let \nolinkurl{dn' = delta_N(D) dn}. At the selected split place, the e-to-core roots have unit ratio; the core-to-f roots have ratio of absolute value \texttt{q\^{}\{-1\}} and number \nolinkurl{(k+1)rho}; the e-to-f roots have that same absolute value and number \nolinkurl{rho^2}. The internal e and f roots have ratio one. Therefore

\[
            \delta_N(D)=q^{-\rho(k+\rho+1)}.
 \tag{TC10}
\]

For example, whenever the unitary matrix-coefficient Bessel integral is absolutely convergent, its exact change of variables is

\[
 \alpha_{\chi_c}(\Pi(\tau_u)v_K,v_H)
 =\delta_N(D)^{-1}
   \alpha_{\chi_1}(\Pi(D\tau_u)v_K,v_H).
 \tag{TC11}
\]

The H Haar is unchanged. A stable or regularized version requires its own compatible change-of-variables statement; TC11 does not assert one without that hypothesis. If the rational change is performed globally at every place, the full product of these moduli is one by the product formula. Omitting the real factors while keeping only finite ones is not that global cancellation.

\subsection{4. A false standard-character test is killed, while the literal one passes the compatibility check}\label{a-false-standard-character-test-is-killed-while-the-literal-one-passes-the-compatibility-check}

Choose an index a for which \nolinkurl{u_a} is a unit and let \nolinkurl{f=f_rho}. For \nolinkurl{b in O}, put

\[
 k_b=I+bE_{p_a,f},\qquad
 n_b=\tau_u k_b\tau_u^{-1}
     =I+bE_{p_a,f}+\pi^{-1}u_abE_{z,f}.
 \tag{TC12}
\]

Both nonidentity entries of n\_b are allowed in N. Since k\_b is in K, n\_b fixes \nolinkurl{Pi(tau_u)v_K}. The unit-root character, however, is

\[
               \chi_1(n_b)=\psi(\pi^{-1}u_ab).
 \tag{TC13}
\]

One can choose b with TC13 different from one. Every chi1-equivariant Bessel functional consequently vanishes on the bare translated spherical vector, independently of the H representation. This conclusion concerns the incorrectly standardized character/vector pair.

For the literal character TC5, the same element has phase

\[
                     \chi_c(n_b)=\psi(u_ab)=1.
 \tag{TC14}
\]

In fact chic is trivial on the entire \texttt{N\ cap\ tau\_u\ K\ tau\_u\^{}\{-1\}}. For the coordinate covector, conjugation by g identifies this intersection with \nolinkurl{N_T cap K}, up to the exact integral factor in LIFT30 NL8. In the latter group every root pairing in psiFin is integral: the physical z\_T and T-local e/f vectors are integral, and K preserves the original lattice. The additive character is trivial on those pairings. Integral covector transport gives the same result for arbitrary actual u. This is a necessary local compatibility statement; it does not prove the Bessel functional is nonzero on TC9.

\subsection{5. The central JPSS integral records the same warning}\label{the-central-jpss-integral-records-the-same-warning}

For comparison, write the central block order as \nolinkurl{(H,E,z,F)} and set \nolinkurl{j=rho}. The JPSS integral is

\[
 \Psi_j(s,W,W')=\int_{U_k\backslash\mathrm{GL}_k(F)}
  \int_{\operatorname{Mat}_{j,k}(F)}
 W\!\begin{pmatrix}h&0&0\\X&I_j&0\\0&0&I_{j+1}\end{pmatrix}
 W'(h)|\det h|^{s-(r-k)/2}\,dX\,dh.
 \tag{TC15}
\]

The determinant change under simultaneous right H action is \texttt{\textbar{}det\ h0\textbar{}\^{}\{1/2-s\}}; hence the central value has the required H covariance. The JPSS functional equation sends j to \nolinkurl{r-k-1-j=j}. These are the literal equation (2.1) and Theorem 2.3 of Jacquet--Liu, §2.2, with their measures and contragredient convention retained.

Let \nolinkurl{c0=k+j+1} be the marked-line column and let \nolinkurl{u_beta=I+sum beta_a E_(a,c0)}. If A(h,X) is the matrix in TC15, then

\[
 W^\circ(A(h,X)u_\beta)
       =\psi(X_{j,*}\beta)W^\circ(A(h,X))\qquad(j>0).
 \tag{TC16}
\]

Indeed \nolinkurl{A u_beta = n(g beta,X beta) A}, with n an upper unipotent supported in column c0, and the only adjacent character entry is the last row of X beta. Right spherical invariance makes \nolinkurl{W^circ(A(h,X))} periodic under integral translations of X. If any coordinate of beta is in \texttt{pi\^{}\{-1\}O\^{}times}, translation in that coordinate of the last X row shows that the integral is zero in its region of absolute convergence, and therefore identically zero after rational continuation. This is the analytic counterpart of TC12--13, not an actual-period vanishing theorem.

By contrast, Booker--Krishnamurthy--Lee (2020), §3, (3.1), uses column \nolinkurl{k+1} in the j=0 integral. There the phase is \nolinkurl{psi((h beta)_k)}, which their compact average evaluates by a conductor Gauss sum. In the actual positive-rank problem TC9 adds the nontrivial flag diagonal. Under the transpose-dual functional equation it supplies the opposite flag scaling; the integral-periodicity argument for the bare spherical vector no longer applies. A direct computation with this diagonal and the unchanged smaller newform/projector is still required. Neither the BKL existence theorem nor the equality of the fractional ideals generated by the different JPSS integrals evaluates this fixed test.

\subsection{6. Exact scope}\label{exact-scope-1}

The proved conclusions are TC5, TC8--10 and TC12--14, together with the diagnostic central-integral warning TC16. They prevent an invalid direct use of the two-step K053 test in positive rank and identify the extra diagonal and local Haar power in the actual test. They preserve the physical flag, marked vector, strict T covector level, generic G projector and all global period data. No local nonzero value, global nonzero Bessel/Mram period, uniform cost, arithmetic Read, or literal family instance is concluded.

\subsubsection{Primary references directly read}\label{primary-references-directly-read}

\begin{itemize}
\tightlist
\item
  Booker, A. R., Krishnamurthy, M., \& Lee, M. (2020). \emph{Test vectors for Rankin--Selberg L-functions}. Journal of Number Theory, 209, 37--48. Equation (1.1), §3 equation (3.1), equation (3.2), and Lemma 3.1. \url{https://research-information.bris.ac.uk/files/211450903/gln_local.pdf}
\item
  Jacquet, H., \& Liu, B. \emph{On the local converse theorem for p-adic GLn}. §2.2, equation (2.1), Theorem 2.3 (the JPSS functional equation). \url{https://archive.ymsc.tsinghua.edu.cn/pacm_download/318/8960-JL-LCT-final.pdf}
\end{itemize}

The exact pinned source and accepted internal proof locators are stated in §§1--2. The companion script checks split conjugation, both CM components, character rescaling, modulus exponents and the matrix identity TC16 over finite exact samples. Those checks supplement the displayed proofs; they do not prove a global period or a new Rankin--Selberg evaluation.

\clearpage

\section{AM-DO-078 --- all support-twelve rank-two unipotents in the fixed five-dimensional block}\label{am-do-078-all-support-twelve-rank-two-unipotents-in-the-fixed-five-dimensional-block}

\textbf{Canonical whole-checked version --- 2026-11-22T17:38:20.186621+00:00.} Root and independent peer read the complete frozen proof and script; separate241-assertion replays passed. Original author-export wording is retained below. The promoted theorem removes all actual all-top support12 globally rank-two-unipotent fixed-W0 terms with one5296 error; no support-floor update is implicit.

Author: /root/definite\_oscillator, 22 November 2026 UTC. New controlled pending proof. No existing proof, kernel, registry or checkpoint is modified.

\subsection{1. Actual sector and conclusion}\label{actual-sector-and-conclusion-1}

Keep the actual J56 normalized kernel, original packet, central representatives, finite projectors and class masses, the all-top cutoff H=133, and the selected real module and its favorable-point transport. Put \[
 f=a=168,\quad n=169,\quad E^{169}=S\perp W_0,
 \quad\dim_E S=164,\quad\operatorname{sig}(W_0)=(4,1).
 \tag{FP1}
\] Restrict to actual support-twelve terms satisfying \[
 h=I_S\oplus h_0=I+N,\qquad N^3=0,\quad
 \operatorname{rank}_E N=2,\quad\operatorname{rank}_E N^2=1.
 \tag{FP2}
\] Global rank-two unipotence is an explicit hypothesis. The proof does not infer it from support alone. Every incidence profile is allowed: if i,u,t count the actual types I,II,III, then i+u+t=s=12. At the334 unmixed relative G factors write \[
 d_v=\operatorname{rank}(N\bmod v)\in\{0,1,2\},
 \qquad D_G=\sum_v d_v.
 \tag{FP3}
\] The two mixed factors remain the unchanged partial-projector contractions and are not included in D\_G.

\textbf{Theorem.} The whole FP2 sector, with every support, incidence type, ordered pair, finite rank pattern, original-column and class-fibre multiplicity retained, has total absolute contribution strictly below2\^{}-5296. More precisely, the D\_G\textgreater=14 portion has this bound, while the D\_G\textless=13 portion is empty.

Remove only its intersection with the current remainder and retain all previous errors. This is a conditional structural-sector theorem, not by itself a complete support-twelve exhaustion or a signed nonvanishing theorem.

\subsection{2. The square coefficient, its exact ideals and the real factor}\label{the-square-coefficient-its-exact-ideals-and-the-real-factor}

Use the actual square-pole theorem K071/K074 and the signature-one identities already retained in K085 and R052. They give \[
 \alpha=(N^2)_{nn}\in F,\qquad
 0<\alpha_\tau\le R:=H+H^{-1}-2=17424/133.
 \tag{FP4}
\] At type I, alpha is integral; at type II it has poles of order at most2; at type III at most1, at both CM components. This is the entrywise statement about N², not a claimed standard-frame integrality of h. In type I, h preserves the intersection and sum of the full pencil; its unipotent action on their two-dimensional quotient makes N² carry the standard intermediate lattice into itself. The actual type-III incidence gives simple poles for h and h\^{}-1; N²=h+h\^{}-1-2I. Type II retains the general two-shear pole bound.

Let I\_II,F and I\_III,F be the products of the inverse relative T primes of their indicated types. Define \[
 \mathfrak J_F=\prod_{v\ \mathrm{unmixed}\ G}\mathfrak p_v^{2-d_v},
 \quad J=N_F\mathfrak J_F=
 \frac{Q^{334}}{\prod_v q_v^{d_v}},\qquad Q=g_0g_1.
 \tag{FP5}
\] There are167 factors of each residue size. At d\_v=0 the matrix N is prime-divisible, so N² is square-divisible. At d\_v=1 its finite rank-one nilpotent reduction has square zero. Both statements hold at the inverse-adjoint partner. Therefore \[
 \alpha\in\mathfrak J_F I_{\mathrm{II},F}^2 I_{\mathrm{III},F},
 \quad N_F(\alpha)\ge J T^{-\mu},\quad
 \mu=2u+t,\quad y=(JT^{-\mu})^{1/f}\le R.
 \tag{FP6}
\] This is an F-ideal statement: the two CM partners are not counted as two independent relative factors. The patterns of (d\_v) have at most3\^{}334\textless2\^{}530 choices. No common-side orientation is necessary for this real scalar.

The carried exact Green and selected real estimates are \[
 |a(h)|\le2Q^{-335(f-1)}J^{335/2},
 \qquad
 |c(h)|\le12^f2^{f/2-959f/4}N_F(\alpha)^{-1/4}.
 \tag{FP7}
\] The latter is K085(AS1), with the actual homogeneous Pieri degree greater than2\^{}959, the fixed near-point factor, and the unchanged normalized compression. Both estimates apply to FP2 before a type profile is selected. The Green constant2 is retained once; the two mixed factors each retain their modulus bound1. No sign is inferred.

\subsection{3. The affine column count and the fixed-block source fibre}\label{the-affine-column-count-and-the-fixed-block-source-fibre}

For a fixed actual ordered source/target tuple, the W0 part of the negative column belongs to the actual target overlattice. Its real-part projection beta=Re(h\_nn) satisfies gamma=beta-1=alpha/2. The corrected K085/DO072 coarea statement, AM-DO-070-C1 and LIFT063 give the following bounds, used with their original metrics: \[
 \#\{\text{W0 columns with fixed beta}\}
 <2^{5379}T^{8s+r_0},
 \tag{FP8}
\] where r0 counts actual zero target labels, and \[
 \#\{\gamma\}
 <2^{1062}\left(\frac{375+R/(2y)}{376}\right)^f.
 \tag{FP9}
\] For specificity, gamma lies in one half of the F-ideal in FP6. Its covolume is2\^{}-f sqrt(D\_F)y\^{}f and its covering radius is below375 sqrt(f)y. The trace-ball bound with radius(R/2)sqrt(f) yields FP9. Thus the dyadic factor, the square-root discriminant and the possible smallness of alpha are retained.

The constant5379 is the conservative, already checked affine-kernel volume for radius223 sqrt(f), using the original-ring-only cover of O\_E\^{}4 plus iO\_F. The kernel has real dimension9f. At a nonzero target the coarea denominator contributes T\^{}8; at target0 its negative-coordinate projection has only a simple pole, and the kernel contributes T\^{}9. FP8 retains precisely this extra T. No original-ring cover is assigned to a general fractional ideal in FP9.

The column fibre in this fixed-S sector has a particularly direct bound. If h and h' have the same complete negative column and the same prescribed source and target modules, then (h')\^{}-1h preserves the source module and fixes both S pointwise and e\_n.~It acts faithfully on their positive E-orthogonal complement of dimension4. Its intersection with the actual source module is a full projective O\_E lattice; no rational orthogonal basis is declared integral. The finite automorphism group injects at the two selected odd split G primes into GL\_4(k). This is the same torsion-free reduction argument used in DO073: Serre (2006, section1.2, Lemma1′, printed p.5), in its precise prime-to-p congruence form, together with the cyclotomic degree exclusion of p-torsion for the actual primes, applies in rank4 as it does in rank5. Its order therefore divides both GL\_4 orders, hence also both GL\_5 orders. Consequently the exact retained bound is \[
 G_5=\gcd(|\mathrm{GL}_5(g_0)|,|\mathrm{GL}_5(g_1)|)
 =98148490017978252769652736000000<2^{107}.
 \tag{FP10}
\] This argument does not invoke DO073's nonzero-S-column hypothesis: here S is fixed from the outset and the faithful positive complement has dimension4. Positive definiteness and discreteness make the automorphism group finite. The standard away-T G lattices give the required free rank4 local intersection modules over the DVRs.

\subsection{4. All actual ordered-pair weights, including zero targets}\label{all-actual-ordered-pair-weights-including-zero-targets}

Put m\_T=T-1. The following complete local upper bounds combine each ordered-pair/rho weight with the extra zero-target factor in FP8: \[
 \begin{array}{c|ccc}
 \text{type}&I&II&III\\\hline
 \text{complete pair/coarea upper bound}&2T&1&2.
 \end{array}
 \tag{FP11}
\] Here is the exact justification, independent of any scalar-cross coordinate argument. At an active type-I place, the induced projective action is a nontrivial unipotent action on the two-dimensional pencil quotient: projective identity would fix the admitted standard lattice and contradict activity. K070's unipotent table gives \textbar rho\textbar\textless=1/m\_T. There are at most m\_T ordered sources per target. Summing over m\_T-1 nonzero targets and the zero target weighted by T gives (m\_T-1)+T=2m\_T\textless2T. This bound deliberately counts more ordered pairs than an individual partial permutation.

At type II, the unique incidence is fixed. The zero label cannot be this active fixed lattice, since it is the standard lattice. Summing the at most m\_T diagonal labels against rho=1/m\_T is at most1. At type III, the distinct ordered labels have exact absolute weight1/{[}m\_T(m\_T-1){]}. The m\_T-1 zero-target pairs and (m\_T-1)\^{}2 nonzero-target pairs give \[
 \frac{(m_T-1)^2+T(m_T-1)}{m_T(m_T-1)}=2.
 \tag{FP12}
\] The opposite CM component is already included in FP8 and the relative incidence; it does not give an extra factor. Thus an entire fixed type assignment costs at most2\^{}(i+t)T\^{}i in addition to the common T\^{}(8s) fibre factor.

\subsection{5. Comparing every profile at the same finite-rank ideal}\label{comparing-every-profile-at-the-same-finite-rank-ideal}

For a fixed G pattern define the comparison parameter \[
 y_0=(JT^{-2s})^{1/f},\qquad z_0=\sqrt{y_0},\qquad
 y/y_0=T^{(2i+t)/f}.
 \tag{FP13}
\] The gamma factor in FP9 decreases as y increases. The remaining y-dependent part of FP7 is y\^{}-f/4. Hence the real-and-gamma factor at y is at most its value at y0 multiplied by T\^{}(-(2i+t)/4). Multiplying the actual pair/coarea factor yields \[
 2^{i+t}T^{i-(2i+t)/4}
 =2^{i+t}T^{i/2-t/4}\le2^sT^{s/2}.
 \tag{FP14}
\] This is a comparison at the same actual J, not a replacement of a pattern by an all-II matrix. A hypothetical all-II matrix need not exist. Since y0\textless=y\textless=R, all bounds used on the comparison parameter are legitimate.

There are at most2\^{}168 choices of support and at most3\^{}s\textless2\^{}(2s) type assignments on each fixed support. Thus no further91-profile union factor is needed. Retain all3\^{}334 patterns, the full fixed-block fibre, one Green constant2, and 12\^{}f\textless2\^{}(4f). The whole sector at a pattern parameter z0 is bounded by \[
 \begin{split}
 \mathcal F(z_0)={}&
 2^{1+5379+1062+107+168+2s+530+s}
 12^f2^{f/2-959f/4}
 Q^{-335(f-1)}T^{343s+s/2}\\
 &\times z_0^{669f/2}
 \left(\frac{375+R/(2z_0^2)}{376}\right)^f.
 \end{split}
 \tag{FP15}
\] The exponent343s is335s from J\^{}(335/2) and8s from the affine fibre. The positive z0 exponent before the gamma term is335f-f/2=669f/2. Unlike the scalar-cross bounds, there is no S-column volume here. Apart from positive constants, its variable part is \[
 z_0^{665f/2}(375z_0^2+R/2)^f,
 \tag{FP16}
\] which is increasing for z0\textgreater0. The count and weight can therefore be bounded uniformly over every actual pattern by substituting a common upper cap for z0. The pattern factor is retained only once.

\subsection{6. The high-finite-rank sector}\label{the-high-finite-rank-sector}

The actual prime order is g0\textless g1. If D\_G\textgreater=14, FP5 implies \[
 z_0^{336}=J/T^{24}\le Q^{334}/(g_0^{14}T^{24})<6^{-336}.
 \tag{FP17}
\] The last comparison is an exact integer check. At z0=1/6, the complete variable factor in FP15 has binary upper exponent-144783: \[
 2^{-144784}<
 6^{-56196}\left(\frac{375+18R}{376}\right)^{168}
 <2^{-144783}.
 \tag{FP18}
\] Use the actual one-sided logarithms L\_T=1019333/1024 and L\_Q=287843/4096. The full exponent at s12 is \[
 \begin{split}
 E={}&1+5379+1062+107+168+24+530+12\\
 &+(343\cdot12+6)L_T-335\cdot167L_Q
 +168(9/2-959/4)-144783\\
 ={}&-21696243/4096<-5296.
 \end{split}
 \tag{FP19}
\] This includes every support, type assignment, actual pair weight and zero-target correction, rank pattern and fibre. There is one monotone maximum, not an endpoint union. Consequently it gives one charge5296 for the complete D\_G\textgreater=14 sector.

\subsection{7. Integral square pairings for arbitrary source/target arrows}\label{integral-square-pairings-for-arbitrary-sourcetarget-arrows}

It remains to treat D\_G\textless=13. We establish the pairing input for all incidence types; a common h-stable lattice is not assumed.

At an actual selected T component the source and target are the exact paired modules \[
 M_c=(I+c\pi^{-1}zq)\Lambda,\qquad
 M_d=(I+d\pi^{-1}zq)\Lambda,\qquad hM_c=M_d,
 \tag{FP20}
\] with their inverse-adjoint duals at the opposite component. Use LIFT44's exact z=v0/bar(pi), q=-v0\emph{J-pi m0, where v0=(w,pi\_k), w=Qu and u is the fixed LIFT41 integer vector. Let y be an original integral vector supported on the four positive coordinates of W0 with w}y=0. K088(IP5--6), directly from this literal model, proves that y lies in every admitted M\_c at both CM partners. The proof uses qy=-pi m0y at the selected component and integral pairing with Lambda+pi\^{}-1 O z at the dual component. It remains true at label0 and at every one of the three incidence types.

Patch the actual source tuple to a projective selfdual lattice M\_source and the actual target tuple to M\_target, using the standard lattice elsewhere. The matrix h carries the first to the second. Our y lies in both. The two separate integral pairings therefore give \[
 \langle y,hy\rangle_J\in O_E,
 \qquad \langle y,h^{-1}y\rangle_J\in O_E.
 \tag{FP21}
\] The first is a pairing inside M\_target, the second inside M\_source. This is precisely where retaining both CM partners matters. Since N²=h+h\^{}-1-2I and N² is Hermitian self-adjoint, \[
 b_y=\langle y,N^2y\rangle_J\in O_F.
 \tag{FP22}
\] No integrality of N²y in the standard coordinates, no h-stability of a single lattice and no identification of distinct labelled lattices is inferred. Away T all the original matrices and vectors are standard integral, so FP22 has no omitted finite-place condition.

\subsection{8. The low-rank sector is empty}\label{the-low-rank-sector-is-empty}

The exact K088 original constants are \[
 U=77Q^2-14,\quad w=Qu,\quad u^tu=U,\quad
 u=(2,95192390527,383092671043,12529843548438452388609).
 \tag{FP23}
\] Use its three fixed integral vectors \[
 \begin{split}
 y_1&=(-976627,-153032,38026,0),\\
 y_2&=(-567678,695570,-172838,0),\\
 y_3&=(104583,-15309315176,-61610045224,2).
 \end{split}
 \tag{FP24}
\] Their exact squared norms are978665066829,835948910828 and4030172803672366725045. The recorded unimodular kernel transformation proves that they span w-perp in the positive four-space. These are the existing packet vectors and exact integer rows, not a new reduced lattice or a new four-square representation.

At each G prime the square valuations of section2 imply b\_y in J\_F. Under D\_G\textless=13, \[
 J\ge Q^{334}/g_1^{13}\ge Q^{334}/g_1^{69}.
 \tag{FP25}
\] The original real normal form is N²=-alpha vv*J, with v in W0 isotropic and normalized negative coordinate1. It gives \textbar b\_y,tau\textbar\textless=R\textbar\textbar y\textbar\textbar². K088's strictly checked inequalities Q\^{}334\textgreater g1\^{}69(R\textbar\textbar y\_j\textbar\textbar²)\^{}168 therefore force all three b\_yj to vanish. Their rational orthogonal equations imply imN² is contained in span\_E(w,e\_n). This confines only the square's image, not necessarily the entire matrix h.

Writing that line as v=(a w,1) gives a*bar(a)=1/(Q²U). The original11-entry, which uses w1=2Q, is then exactly \[
                         (N^2)_{11}=-4\alpha/U.
 \tag{FP26}
\] The actual U is odd and coprime to Q and T. Standard away-T integrality of N² gives alpha/U integral there; at primes dividing U use FP26, and at2 use the ordinary nn entry and the oddness of U. At G retain the same J\_F valuations. At T retain the actual I/II/III square-pole orders. Thus the strengthened ideal is \[
           \alpha\in U\mathfrak J_F I_{\mathrm{II},F}^2 I_{\mathrm{III},F}.
 \tag{FP27}
\] Relative ramification indices cancel in the U-divisibility comparison because alpha and U are F-elements. No square root of U, anisotropy assertion, selected-half norm or moving integral frame occurs.

Since mu=2u+t\textless=24, positivity and FP25--27 give \[
 R^{168}\ge N_F(\alpha)
 \ge U^{168}Q^{334}/(g_1^{13}T^{24})
 \ge U^{168}Q^{334}/(g_1^{69}T^{24})>R^{168}.
 \tag{FP28}
\] The final strict comparison is K088(IP23), independently checked with the literal actual integers; the companion also checks the stronger exponent13 directly. This contradiction proves emptiness of every low-rank fixed-W0 profile, completing the theorem together with FP19.

\subsection{9. Verification scope and remaining gates}\label{verification-scope-and-remaining-gates}

The proof uses the exact original K062 kernel and weights, K070/K071/K074 square-pole and unipotent incidence statements, K085's actual homogeneous-energy normalization, corrected DO072/DO070-C1's affine coarea, LIFT063's original-ring cover and field covolume, and the already checked K088 integer kernel and numerator-U identity. The new steps are the complete-profile comparison at fixed J, the high-rank threshold14, and the arbitrary-arrow integral-pairing argument allowing the K088 contradiction for the fixed-W0 sector. The fixed-block positive rank4 fibre has been justified separately; DO073's nonzero-S-column premise is not silently applied.

The certificate verifies all literal inputs, GL5 gcd, exact ideals' exponents, rank cap, rational variable bracket, complete exponent, all91 incidence-profile comparisons, the integer kernel, U coprimality and both norm contradictions. It does not enumerate global matrices or automatically prove the geometric and automorphic interfaces.

This result changes no real role, canonical twist, prime, finite level, packet, projector, arithmetic class or CM phase. It supplies one actual absolute error5296 for the globally rank-two-unipotent fixed-W0 support12 sector. It must be combined with separately justified structural reductions before any support-floor update. No signed sufficient inequality, global period/Mram, literal Host/Read, integer DP44 saving or RH conclusion follows from this theorem alone.

\subsection{Author whole-file verification}\label{author-whole-file-verification}

At 2026-11-22T17:35:42.833567+00:00, I completed a whole read of all nine saved sections FP1--28, all scope paragraphs, and the complete companion script. The first complete exact replay passed241 assertions at2026-11-22T17:34:24.140306+00:00. I reread K071(NB4--5) for the full type-dependent square poles, K085's scalar compression and coarea, the corrected DO072/C1 zero-target interface, LIFT44(PR4--9) and LIFT41(NP1--9), and DO073's precise local torsion argument. I checked the independent positive rank4 fibre separately from its scalar-cross rank5 analogue. I reconstructed the fixed-J monotonic comparison for all91 profiles, the actual pair/coarea sums, the unique common maximum, and the arbitrary-arrow two-lattice pairing proof at both CM partners. The exact unchanged K088 kernel and U-divisor argument were checked with the stricter budget13 and with the conservative already verified budget69. The final citation precision was reread; it changes no formula.

The new exact certificate does not automate those geometric interfaces. Final fingerprints and replay time are recorded in the separate author verification/manifest. Author status: whole-proof and complete-script PASS; ready for independent review, with no main-registry promotion.

\clearpage

\section{AM-R-057 --- complete all-top support-twelve removal}\label{am-r-057-complete-all-top-support-twelve-removal}

Author: /root, GPT-6 Astra Pro, ALP\_MACHINE. 22 November2026 UTC. Status at creation: full joint proof awaiting independent review and canonical promotion.

\subsection{1. The actual normalized energy and the new assertion}\label{the-actual-normalized-energy-and-the-new-assertion}

Retain the original J56 all-top assignment, K044 upper-gap endpoint, actual permitted packet point, finite projectors, ordered scalar-pencil incidences, class masses, central phases, selected compact representation and Hilbert normalization. The field degrees are {[}E:Q{]}=336 and {[}F:Q{]}=f=168, the Hermitian rank is169, and the already charged all-top cutoff is H=133. The separate uniform H8192 family is unchanged.

R054(JT1) and the subsequent checked cuts express this same normalized energy as1 plus the real part of a remaining original signed sum, with a controlled accumulated error. Every retained all-top term already has support at least12. The support s(h) counts active nonintegral relative T places.

\textbf{Theorem.} There is a remainder R\_top\^{}(13), consisting only of original actual terms with s(h)\textgreater=13, such that

\[
 \mathcal E_{\rm top}
 =1+\operatorname{Re}\sum_{h\in R_{\rm top}^{(13)}}w(h)
       +\epsilon_{\rm top},\qquad
 |\epsilon_{\rm top}|<2^{-230}.
 \tag{J13.1}
\]

The current complete all-top error list has52 entries. This result adds one fresh charge701241 to the51-entry list after DO078. It does not add the DO0785296, K0871445 or R0567110 charges again. It preserves the independent uniform support floor11 and its31-entry error list.

Taking real parts writes the real energy identity in real form; it changes no summand, phase or absolute estimate. The inequality

\[
 \operatorname{Re}\sum_{h\in R_{\rm top}^{(13)}}w(h)
                  >-1+2^{-230}
 \tag{J13.2}
\]

remains an unproved sufficient condition for the desired positivity. The theorem is a complete support reduction, not a proof of J13.2.

\subsection{2. The original structural partition}\label{the-original-structural-partition}

Use the fixed original coordinates

\[
 S=\langle e_5,\ldots,e_{168}\rangle_E,\quad
 W_0=\langle e_1,e_2,e_3,e_4,e_{169}\rangle_E,\quad
 A=h_{SS},\quad r=\operatorname{rank}_E(A-I).
 \tag{J13.3}
\]

The carried K068/K069 residual has r\textless=33 and unipotent reductions at all334 unmixed relative G factors. No stronger global r\textless=22 conclusion is needed or attributed to R048. At the few uses of DO074 below one has r=0 or1, within its stated r\textless=22 hypothesis.

Let K\_A=ker(A-I) intersect ker(A*-I). K069(IC1--8) gives the exact partition into the fixed-space class F, where h fixes K\_A pointwise, and its complementary exceptional class IC8. In IC8, DO064 proves A is a partial monomial matrix, with its precise original-coordinate and paired finite properties. Neither classification is inferred from a failed numerical upper estimate.

We first dispose of A!=I at support12.

\begin{longtable}[]{@{}
  >{\raggedright\arraybackslash}p{(\columnwidth - 2\tabcolsep) * \real{0.5000}}
  >{\raggedright\arraybackslash}p{(\columnwidth - 2\tabcolsep) * \real{0.5000}}@{}}
\toprule\noalign{}
\begin{minipage}[b]{\linewidth}\raggedright
Actual class
\end{minipage} & \begin{minipage}[b]{\linewidth}\raggedright
Exact complete cut at s=12
\end{minipage} \\
\midrule\noalign{}
\endhead
\bottomrule\noalign{}
\endlastfoot
F and2\textless=r\textless=33 & K069(IC20--21): the smallest support threshold in this range is22, so all these cells are in its already charged4084 union. \\
IC8 and3\textless=r\textless=33 & DO064(PM22--23): the smallest threshold in this range is16, so all these cells are in its already charged4084 union. \\
IC8 and r=1 or2 & DO075(CR1--17) removes these entire support12 cells without assuming global unipotence; its1859541 charge is retained. \\
F,r=1,nonunipotent & DO074 makes t\textless=10 impossible for a nonunipotent matrix here. R050 §5 covers t11 and t12 at support12, in its intervals11..59 and12..48 respectively; its existing4088 union is retained. \\
F,r=1,globally unipotent & These are exactly the full-rank pairs(R,r)=(1,1) and(2,1), and are removed by the fresh complete bound in §3 below. \\
\end{longtable}

Here R=rank\_E(h-I) is the full rank, distinct from r. The signature-one unipotent classification DO061(TR19--21), retained in DO065(UR4,UR8), gives R\textless=2 for every nonidentity global unipotent, so no further full-rank case is omitted in the last row.

Now suppose A=I. Then K\_A=S. If h fixes S, unitarity makes h=I\_S direct-sum h0 on W0. Otherwise it is the actual scalar-cross class. Its original S projection x=(h e\_169)\_S is nonzero. This is K069's block dichotomy, also used explicitly in R051 §4 and R056 §3; no additional free vector is introduced.

\begin{longtable}[]{@{}
  >{\raggedright\arraybackslash}p{(\columnwidth - 2\tabcolsep) * \real{0.5000}}
  >{\raggedright\arraybackslash}p{(\columnwidth - 2\tabcolsep) * \real{0.5000}}@{}}
\toprule\noalign{}
\begin{minipage}[b]{\linewidth}\raggedright
Actual class with A=I
\end{minipage} & \begin{minipage}[b]{\linewidth}\raggedright
Exact complete cut at s=12
\end{minipage} \\
\midrule\noalign{}
\endhead
\bottomrule\noalign{}
\endlastfoot
Fixed S, nonunipotent & DO067 forces global unipotence for t\textless=9. K081 §4 removes t10,t11,t12 at support12, within endpoints164,149,134; its existing4086 union is retained. \\
Fixed S, globally unipotent of full rank1 & K077(AP23,§7) includes s12 in its already charged all-top support10..14 union1436. Its separate s12 margin is9907. \\
Fixed S, globally unipotent of full rank2 & DO078(FP1--28) removes every one of the91 type profiles with one5296 charge. The finite-rank-sum\textgreater=14 portion is bounded; the\textless=13 portion is empty. \\
Scalar cross, nonunipotent & DO074 restricts this r=0 case to t11/t12. R056(BM1--28) removes the complete two-sector union, at all allowed full ranks and without a finite-rank-count restriction, with its existing7110 charge. \\
Scalar cross, globally unipotent & DO065 §2 gives full rank2, N³=0 and rank N²=1. K086 removes every profile other than all-II; K087 bounds the all-II finite-rank-sum\textgreater=70 portion and K088 proves its\textless=69 complement empty. The existing885 and1445 charges are each retained once. \\
Full rank0 & h=I has support0 and is absent at support12. \\
\end{longtable}

These two tables exhaust every actual support12 term. The splitting by A=I, F/IC8, fixation of S and global unipotence is exhaustive. The separate full-rank-one scalar-cross unipotent possibility is excluded by DO065 §2, not silently discarded. DO078 proves its fixed-block source-fibre bound directly; its application does not import DO073's nonzero-S-column premise.

All old numerical errors remain even if the larger upper-bound sets overlap. At each removal stage only the intersection with the then-current actual remainder is deleted. Thus removed sets are disjoint while the numerical budget is conservative. In particular K088 adds no second1445, and the earlier DO0765341 and DO0773947 charges remain after R056's broader cut.

\subsection{3. The fresh fixed-space rank-one-block charge}\label{the-fresh-fixed-space-rank-one-block-charge}

The old DO065 removed union was restricted to t\textless=8, because that inequality supplied global unipotence at that stage. Its actual conditional counting proofs UR9--14 require global unipotence but no additional bound on t once that premise is given. R051 §5 freshly charged their entire s11 cells; it did not charge the s12 cells now needed.

For clarity, here are the retained complete estimates and their exact meaning. On the actual full-rank2,block-rank1 class, the two integral witnesses in A-I and (N²)\_SS give the Green bound DO065(UR9) with no T numerator. The original rank-one matrix chart for A and the seven-original-coordinate negative-column projection give the full count UR10, including support and ordered-pencil choices, target-overlattice T\^{}(14s) and the unchanged source fibre G\_n\textless2\^{}5022. Multiplication by the actual pair weights gives UR12:

\[
 \log_2 B_{2,1}(s)
 <-408071827/256+14941s.
 \tag{J13.4}
\]

For full rank1, the injective original pivot-column reconstruction and its actual Green bound give UR13--14, with no extra isometry-fibre or ordered-pair multiplier:

\[
 \log_2 B_{1,1}(s)
 <-103530801/128+8966s.
 \tag{J13.5}
\]

Both formulas use the conservative original H8192 count and therefore bound the actual H133 subset. They retain the unchanged selected real coefficient's modulus bound1. This is a restriction of an already proved upper estimate; it is not an assertion that every large-t matrix is unipotent.

At s=12 the two full exponents are

\[
 E_{1,1}=-89759025/128<-701242,\qquad
 E_{2,1}=-362173075/256<-701242.
 \tag{J13.6}
\]

Consequently their whole union is strictly below2\^{}-701241. Charge the entire two-cell bound afresh and delete only its intersection with the current remainder. The already charged portions with t\textless=8 need not be identified individually, and no previous charge is subtracted. This is the single new charge of the present theorem.

\subsection{4. The full chronological error list}\label{the-full-chronological-error-list}

After the preceding checked DO078 integration, the actual all-top list has51 entries. Appending701241 gives

\[
\begin{split}
(&1999,104546,101550,211446,11671,231,1483877,1307183,9556,\\
 &4084,4085,4084,4086,4088,4086,4085,9199,4088,7423,5183,\\
 &8105,17434,7227,2667,2067,7319,3084,1436,32739,2000,4094,\\
 &4086,5001,4088,1166,195718,93271,34539,685,43330,838,1482,\\
 &2642,885,1859541,1263,1445,5341,3947,7110,5296,701241).
\end{split}
\tag{J13.7}
\]

Every chronological prefix from R054's42 entries through DO073,K086,DO075,R055,K087,DO076,DO077,R056,K088 and DO078 is retained and checked by the companion certificate. The K088 stage has the same50 entries as R056 because it adds no numerical error. No old contained-sector bound is erased or charged twice in assembling this list.

Exactly one entry is231 and each of the other51 entries is at least685. Therefore

\[
 \sum_{e\in\mathrm{J13.7}}2^{-e}
 \le2^{-231}+51\cdot2^{-685}
 <2^{-231}+2^{-679}<2^{-230}.
 \tag{J13.8}
\]

The independent uniform list remains the same31 entries as before. It has exactly one231 and every other entry at least797, so its sum is at most2\^{}-231+30·2\^{}-797\textless2\^{}-230. The certificate also verifies both sums directly as exact dyadic integer inequalities, including all ordered prefixes. It does not infer a uniform H133 cutoff.

The structural tables and §3 remove every support12 term. R054 had already removed every smaller all-top support. Combining their unchanged energy identity with J13.8 proves J13.1.

\subsection{5. Verification and substantive limits}\label{verification-and-substantive-limits}

The root reread the exact K069 and DO064 threshold tables, R050 §5's type/support intervals, K081 §4, K077 §7, the entire conditional DO065 proof, the current R051/R054 scope and the newly checked R056/K088/DO078 proofs. The separate independent assembly note records the same exhaustive partition and identifies the necessary fresh701241 charge. The final whole-file peer review and exact replay are separate events; this source text does not claim they occurred before they did.

The companion diagnostic checks the exact rank/support tables used, all13 admissible t values in the r=0/1 nonunipotent alternatives, the fresh rational exponents, all91 type profiles as an index set, every stated historical error prefix, and the full dyadic sums. It does not enumerate global rational matrices or mechanically establish the geometric partition, integral pairings or analytic estimates.

At review20 the support12 scalar-cross branch had been removed, but no complete support12 exhaustion was proved. The present theorem removes the remaining fixed-block branch through DO078 and the full fixed-space rank-one-block cells through §3, leaving a strictly smaller main signed-kernel gate over support at least13.

The actual admissible mixed-CM family and every changed-field, height, endpoint, finite-content and global comparison term in DP44 remain to be supplied. No shorter valid lift chain or universal strict negative DP44 difference follows from this theorem. The prescribed Bessel/Mram period, concrete032 cycle pin and definite MainX pairing value remain open. The conditional cP=vT formula of DO055 remains conditional. No Kummer Hecke RH, every-epsilon estimate or global oscillator charge comparison is proved.

\subsection{6. Grothendieck and algebraic-geometric interpretation}\label{grothendieck-and-algebraic-geometric-interpretation}

All rational correspondences and integral modules are those of the same model over Spec(O\_E). The new fixed-space application uses original coordinate charts and determinantal ideals. DO078's two-lattice argument concerns the integral Hermitian pairing on distinct actual source and target projective modules; it does not identify them as one stable lattice. R056's balanced charts change only the real metric used to count points in the unchanged arithmetic module. The present theorem joins their estimates in the original normalized kernel.

The carried module/stabilizer context is SGA3 I4.7.3/4.7.3.1, with exterior-power/determinantal context SGA6 VI Theorem3.3. The polarized CM Jacobian, Picard and biextension data retain the exact locators FGA,Bourbaki232,Theorem3.1,pp.148--149; Bourbaki236,Corollaries1.7/2.3,pp.229/232; and SGA7 I,VII,2.7/2.9,pp.169--174. The formal-group/Dieudonné and Cartier-duality background remains Grothendieck archive cote4,scan003,§§8.1.1--8.1.2,scans005--006,§§8.4.1--8.4.3, and cote54,scans005--007. These are the carried dictionary anchors, not attributions of this new analytic count to those texts.

The change to DP\_MULT Q3e and THETA3 G1/A7 is the strict reduction J13.1 of the actual finite/real sufficient gate. No new lambda-ring or integral K-theory equality, l-adic purity or weight theorem, motive identification, CM-period formula or motivic L-function implication is claimed. CX\_GRH CX-ALG-024 and GCT\_PLETH T-052 retain their distinct unproved arithmetic comparisons.

Internal exact dependencies: R054(JT1--5); K068(LR1--4) and K069(IC1--8,IC20--21); DO064(PM22--23); DO065(UR4,UR8--14); DO067's fixed-W0 two-trace criterion; DO074's r\textless=22,t\textless=10 rigidity; DO075(CR1--17); R050 §5; K077(AP23,§7); K081 §4; K086's whole non-all-II scalar-cross support12 theorem; K087(GB1--18); K088(IP1--23); R056(BM1--28); and DO078(FP1--28). All are cited only in the scopes shown in the tables.

\clearpage

\section{AM-LIFT-068 --- the corrected positive-rank T translate: an exact central test}\label{am-lift-068-the-corrected-positive-rank-t-translate-an-exact-central-test}

\textbf{Canonical checked version --- 2026-11-22T17:55:13.507763+00:00.} Root and an independent peer read the whole final proof and complete script and ran separate 28,115-check replays. Original pending wording below is historical. Promoted scope: Under the explicit generic tempered local premises and c(sigma)=c(omega\_sigma)=1, the actual positive-rank D*tau split Bessel test is nonzero. Its exact phase is PT25 and its modulus after division by L(1/2,Pi x sigma) is q\footnote{-rho(k+rho+1)/2} q\^{}(k/2)/(q\^{}k-1). The physical root modulus cancels in the numerical squared Whittaker expression PT27. Finite genericity/temperedness for an actual family, canonical Hilbert/GGP normalization, the global period and arithmetic DP44 remain separate.

Author: \nolinkurl{/root/lift_accounting}. 22 November 2026 UTC.

Status: author whole-proof/source/script review PASS; exact diagnostics PASS, with the final frozen hashes recorded in the companion author-review and manifest files. This is conditional on the explicit generic tempered local hypotheses below. It is not a new assertion that every genuine general-m theta constituent has those hypotheses.

\subsection{1. Fixed representations, vectors and measures}\label{fixed-representations-vectors-and-measures}

Let F be a nonarchimedean local field, O its integers, varpi a uniformizer, and q its residue cardinality. Let psi have conductor O. Its actual unit coefficient is retained. Fix k\textgreater=2, rho\textgreater=0 and

\[
             n=k+2\rho+1.                                      \tag{PT1}
\]

Let Pi be an irreducible tempered generic spherical representation of GL\_n(F). Let sigma be an irreducible tempered generic representation of H=GL\_k(F) satisfying

\[
              c(\sigma)=c(\omega_\sigma)=1.                    \tag{PT2}
\]

Choose its essential Whittaker vector v, and the spherical vector v\_Pi, so that their Whittaker functions W\_sigma and W\_Pi have value one at 1. The larger Whittaker character is psi; the smaller one is psi\^{}\{-1\}. These are Whittaker normalizations, not assertions that either Hilbert norm is one. Write omega=omega\_sigma.

Additive O has volume one, K\_k=GL\_k(O) has volume one, and U\_k(O) has volume one. Use the resulting right-invariant quotient measure on U\_k(F)\H. Multiplicative O\^{}times has volume one. All compact projectors below are probability projectors. These are the normalization conventions in Booker--Krishnamurthy--Lee (2020, section 2), used in K053.

The theorem is that the literal extra diagonal in LIFT067 gives a nonzero fixed vector test. More precisely, put, in block order (H,E,z,F),

\[
 D=\operatorname{diag}(I_k,cI_\rho,1,\varpi^{-1}I_\rho),\qquad
 \tau=I_n+\varpi^{-1}E_{k+\rho+1,k},\qquad c\in O^\times.
                                                                  \tag{PT3}
\]

Here c is the selected value of the conjugate uniformizer in the actual datum. The integral covector transport of LIFT030 is made on both factors before using PT3. It is not an independently chosen smaller vector. Let Z\_rho(s) be the j=rho JPSS integral of rho(D tau)W\_Pi against W\_sigma. Then

\[
 Z_\rho(1/2)\ne0,\qquad
 \boxed{
 \left|\frac{Z_\rho(1/2)}{L(1/2,\Pi\times\sigma)}\right|
 =q^{-\rho(k+\rho+1)/2}\frac{q^{k/2}}{q^k-1}.}
                                                                  \tag{PT4}
\]

Every phase, compact projection and root measure entering PT4 is retained in PT21 and PT25 below. A complete Bessel-model covariance calculation is given in section 7. The actual-source hypothesis table is section 9.

\subsection{2. The smaller row and column lines}\label{the-smaller-row-and-column-lines}

Put

\[
\begin{split}
 K_r&=\{h\in K_k:e_k\bar h=e_k\},\\
 K_c&=\{h\in K_k:\bar h e_k^t=e_k^t\},\\
 t&=\operatorname{diag}(\varpi I_{k-1},1).
\end{split}                                                       \tag{PT5}
\]

Direct block multiplication gives K\_c=t K\_r t\^{}\{-1\}. The conductor-one newform theorem makes both fixed spaces one-dimensional. Define

\[
 e_{K_r}\sigma(t)v=\lambda_+v,\qquad
 e_{K_r}\sigma(t^{-1})v=\lambda_-v.                               \tag{PT6}
\]

By K053 NT7--14, including the full Weil--Deligne conductor argument, sigma is a normalized induction of one tame ramified unitary character and k-1 unramified unitary characters. If their nonzero unramified parameters are beta\_1,\ldots,beta\_\{k-1\}, then

\[
 \lambda_+=q^{-(k-1)/2}\prod_i\beta_i,\quad
 \lambda_-=\overline{\lambda_+}
           =q^{-(k-1)/2}\prod_i\beta_i^{-1}\ne0,
 \qquad e_{K_c}v=\lambda_-\sigma(t)v.                             \tag{PT7}
\]

For clarity, the last identity is e\_Kc v = sigma(t)e\_Kr sigma(t\^{}\{-1\})v. The first formula follows by evaluating the probability average over the q\^{}\{k-1\} integral upper unipotents representing K\_r/(K\_r intersect K\_c). Their Whittaker characters are one. Miyauchi's diagonal formula gives the elementary Schur polynomial for (1\^{}\{k-1\},0), hence the displayed product. Unitarity then gives lambda\_- as the conjugate. No unnormalized Hecke operator is being substituted.

If a bilinear functional is invariant under simultaneous K\_k action and its larger vector w is K\_c-fixed, PT7 gives the exact identity

\[
              \ell(w,\sigma(t)v)=\lambda_-^{-1}\ell(w,v).
                                                                  \tag{PT8}
\]

This will be applied only to sigma(t)v. It is not an averaging assertion for an arbitrary translated newform.

The bars over matrices in PT5 denote residue reduction, not the CM involution used to define the unit c in PT3.

\subsection{3. Auxiliary JPSS integrals and their compact averages}\label{auxiliary-jpss-integrals-and-their-compact-averages}

For this section n\textgreater k is fixed and 0\textless=j\textless=n-k-1 is arbitrary. Set S=s-(n-k)/2 and

\[
 A_j(h,X)=
 \begin{pmatrix}h&0&0\\X&I_j&0\\0&0&I_{n-k-j}\end{pmatrix},
 \quad D_j=\operatorname{diag}(I_k,\varpi I_j,I_{n-k-j}),
 \quad b_j=k+j+1.                                                \tag{PT9}
\]

For epsilon in O\^{}times and beta=(b\_1/varpi,\ldots,b\_\{k-1\}/varpi, epsilon/varpi)\^{}t, with each b\_i representing O/varpi O, put u\_beta=I+sum beta\_a E\_\{a,b\_j\}. Define

\[
\begin{split}
 I_j^\epsilon(s)
 =q^{-(k-1)}\sum_{b\in(O/\varpi O)^{k-1}}
 \int_{U_k\backslash H}\int_{\operatorname{Mat}_{j,k}(F)}
 W_\Pi(A_j(h,X)D_j u_\beta)
 W_\sigma(h)|\det h|^S\,dX\,dh.
\end{split}                                                       \tag{PT10}
\]

At j=0 the inner integral is absent. Every integral and finite average is absolutely convergent for Re(s) sufficiently large; all identities are first proved there. The JPSS theorem gives rational continuation in q\^{}\{-s\}. We do not exchange divergent integrals at the central point.

The tail diagonal D\_j commutes with H and with u\_beta. The vector Pi(D\_j u\_\{epsilon e\_k/varpi\})v\_Pi is K\_c-fixed: K\_c changes the indicated column by an integral vector, which fixes v\_Pi. Averaging it over K\_r is exactly the finite average in PT10. Indeed the last column of K\_r modulo varpi runs uniformly through (b,1)\^{}t; multiplying by epsilon permutes all b. Thus PT10 is the original-newform pairing with that K\_c-fixed vector, as well as its K\_r average.

The simultaneous right H covariance of the j-th JPSS integral is

\[
 \ell_j(s;\Pi(h_0)w,\sigma(h_0)v)
   =|\det h_0|^{(n-k)/2-j-s}\ell_j(s;w,v).                       \tag{PT11}
\]

This follows by h'=hh\_0, X'=Xh\_0, with dX=\textbar det h\_0\textbar\^{}\{-j\}dX'. In particular it is K\_k-invariant for every s. At n=k+2rho+1, j=rho,s=1/2 it is H-invariant. PT8 is therefore available at every stage of the recurrence, although the intermediate j's are not central.

For j\textgreater=1, multiplication in the marked-line column gives

\[
 W_\Pi(A_j(h,X)D_j u_\beta)
       =\psi(X_{j,*}\beta)W_\Pi(A_j(h,X)D_j).
                                                                  \tag{PT12}
\]

Only the last E row is adjacent to the marked column. The finite average must first be restricted to its actual spherical support. That support forces the entire last row of X into O\^{}k. Indeed, if one of its entries has negative minimum valuation, integral invertible column operations among the H columns and the last E column send that row, whose last E entry is varpi, to a single last E entry x of that minimum valuation. The other columns in that row are eliminated using ratios integral by minimality. Factoring the upper part of the last E column then leaves diag(B,x,I), for some invertible leading block B. The adjacent root from the last E row to the following tail row, together with right K\_n-invariance, forces the Whittaker value to be zero just as in PT13 below. This argument is pointwise and does not move the smaller input.

On this integral last-row support, b -\textgreater{} psi(X\_\{j,*\}b/varpi) is a genuine character of the finite additive quotient. Orthogonality therefore makes the first k-1 entries y belong to varpi O. Write the last row as (y,x). Its remaining phase is psi(epsilon x/varpi). No such orthogonality claim is made for arbitrary nonintegral X before this support argument.

\subsection{4. The zero and nonunit pieces}\label{the-zero-and-nonunit-pieces}

Suppose first that v(x)\textless0. Since y lies in varpi O\^{}\{k-1\}, integral right column operations eliminate y using the last E column, whose only nonzero entry is varpi in the last E row. In the last H and last E columns apply the integral invertible block

\[
                   \begin{pmatrix}-\varpi/x&1\\1&0\end{pmatrix}.
                                                                  \tag{PT13}
\]

The last E row becomes (0,x). After factoring an upper unipotent in the last E column, the remaining matrix is block diagonal with this entry x and the following tail diagonal entry 1. Right invariance by I+a E\_\{k+j,k+j+1\}, a in O, and left Whittaker equivariance imply that its Whittaker value equals psi(xa) times itself. Some a gives a nontrivial phase because v(x)\textless0. The value is therefore zero. No bound on the other rows of X is imposed or needed.

Suppose next x belongs to varpi O. Now the entire last row of X lies in varpi O\^{}k and its Fourier phase is one. The corresponding inner integral as a function of h is right K\_k-invariant: for k\_0 in K\_k, replace all rows of X by Xk\_0\^{}\{-1\}, which preserves the last-row domain and additive measure, and use right spherical invariance and D\_j's commutation with H. Equivalently this may be performed inside the absolutely convergent double integral. Its pairing with W\_sigma is unchanged by e\_Kk. But sigma\^{}\{K\_k\}=0 by PT2. Thus the entire x in varpi O piece is zero.

The proof does not claim that all X is integral. Earlier flag rows may have negative valuation. Removing them would lose genuine contributions. Only the last row has been analyzed.

\subsection{5. The unit pivot and the exact recurrence}\label{the-unit-pivot-and-the-exact-recurrence}

Let x belong to O\^{}times. Eliminate y as above. This contributes the exact integration volume q\^{}\{-(k-1)\}, because W after elimination is independent of y in varpi O\^{}\{k-1\}. In the last H and last E columns use the right K\_n block

\[
 C_x=\begin{pmatrix}-\varpi/x&1/x\\1&0\end{pmatrix},\qquad
 a_x=\operatorname{diag}(I_{k-1},-\varpi/x).
                                                                  \tag{PT14}
\]

The last E row is now (0,1). In the rows above it, the new last H column is multiplied by -varpi/x, while the new last E column consists of the old last H column divided by x. Factor the latter column as a left upper unipotent. Its adjacent Whittaker coordinate is X\_\{j-1,k\}/x for j\textgreater1, and h\_\{k,k\}/x for j=1.

Put h'=h a\_x and X'=X\_previous a\_x. That adjacent coordinate is the negative of the corresponding new entry divided by varpi. Thus, including the generic character, the remaining larger Whittaker value is exactly

\[
 W_\Pi\bigl(A_{j-1}(h',X')D_{j-1}
             u_{-e_k/\varpi}\bigr).                              \tag{PT15}
\]

Here the new marked column is k+j, the position of the removed last E row. The ambient representation and dimension n have not changed. The right pivot has determinant -1/x, a unit, so there is no right-K modulus.

In fact the full matrix identity is stronger than the phase assertion: if E\_y is the integral column elimination of y, then A\_j(h,X)D\_j E\_y C\_x equals the matrix inside PT15 exactly. This also checks the minus sign and the change of marked column without an additional Bruhat-cell choice.

The other factors are exact:

\[
\begin{split}
 dX_{\rm previous}&=q^{j-1}dX',\qquad dh=dh',\\
 |\det h|^S&=q^S|\det h'|^S,\\
 W_\sigma(h)&=W_{\sigma(a_x^{-1})v}(h'),\\
 \sigma(a_x^{-1})v
 &=\omega(-x/\varpi)\,\sigma(t)v.
\end{split}                                                       \tag{PT16}
\]

For the last line use a\_x\^{}\{-1\}=(-x/varpi)I\_k times t times diag(-x\^{}\{-1\}I\_\{k-1\},1). The last factor is in K\_r and fixes v. This is the specific reason the smaller vector can be handled; there is no assertion that a\_x normalizes K\_r.

Let

\[
 G_\epsilon=\int_{O^\times}\omega(u)
                         \psi(\epsilon u/\varpi)\,d^\times u,
 \qquad d_k=\frac{q^k-1}{q-1}.                                   \tag{PT17}
\]

In PT15 the larger vector is K\_c-fixed. Apply PT8 to sigma(t)v and then probability-average the larger vector over K\_r, using the original K\_r-fixed v. This produces I\_\{j-1\}\^{}\{-\}, not an average involving an uncontrolled translate of v. Since dx=(1-q\^{}\{-1\})d\^{}times x on O\^{}times, the three last-row pieces give the exact recurrence

\[
 \boxed{
 I_j^\epsilon(s)
 =q^{S+j-k}(1-q^{-1})\,
   \omega(-\varpi^{-1})G_\epsilon\lambda_-^{-1}
        I_{j-1}^{-}(s).}                                        \tag{PT18}
\]

Every factor in PT18 has been derived before meromorphic continuation. The compact projection is valid at intermediate j by PT11. The epsilon in the next step is -1 because of the chosen integral pivot. Thus no alternating or unspecified Gauss phase is suppressed.

\subsection{6. Closed formula and the actual lower translate}\label{closed-formula-and-the-actual-lower-translate}

Booker--Krishnamurthy--Lee (2020), section 3, equation (3.2) and Lemma 3.1, give the j=0 base with these precise measures and newform:

\[
                   I_0^\epsilon(s)
                  =\frac{G_\epsilon}{d_k}L(s,\Pi\times\sigma).
                                                                  \tag{PT19}
\]

Their naive L factor equals the true Rankin--Selberg factor because Pi is spherical. The same compact computation includes any unit epsilon; equivalently G\_epsilon=omega(epsilon)\^{}\{-1\}G\_1. Induction in PT18 proves

\[
\begin{split}
 \frac{I_j^\epsilon(s)}{L(s,\Pi\times\sigma)}
  =\frac{G_\epsilon G_-^{\,j}}{d_k}
    \bigl[(1-q^{-1})\lambda_-^{-1}
                   \omega(-\varpi^{-1})\bigr]^j
    q^{jS+j(j+1)/2-jk}.
\end{split}                                                       \tag{PT20}
\]

All Gauss sums are nonzero since omega has conductor one. Their probability-normalized modulus is sqrt(q)/(q-1). At the central index and central point, PT7 and S=-rho give

\[
 \left|\frac{I_\rho^\epsilon(1/2)}{L(1/2,\Pi\times\sigma)}\right|
 =q^{-\rho(k+\rho+1)/2}\frac{\sqrt q}{q^k-1}.                   \tag{PT21}
\]

We now return to the lower translate PT3. For a Whittaker function of GL\_l put tilde W(g)=W(w\_l \{\}\^{}tg\^{}\{-1\}), with w\_l antidiagonal. As in K053 NT16--17, the transpose-dual smaller newform satisfies

\[
 \widetilde W_\sigma=C_\sigma\,\rho(t)W_{\sigma^\vee},\qquad
                       |C_\sigma|=1.                            \tag{PT22}
\]

The exact C\_sigma is retained. K053 NT25--26 identifies it with a unitary epsilon factor by the essential-vector test against the spherical rank-(k-1) representation with Satake parameters one, fixing the same probability Haar; this is not a freely chosen phase.

In the JPSS functional equation the extra tail antidiagonal w\_\{n,k\}=diag(I\_k,w\_\{n-k\}) reverses E,z,F. The transpose-inverse of D has diagonal (I\_k,c\^{}\{-1\}I\_rho,1,varpi I\_rho); conjugating by this tail antidiagonal gives (I\_k,varpi I\_rho,1,c\^{}\{-1\}I\_rho). Its last unit block fixes the sphere and commutes with the upper translate. The marked column is fixed by the tail reversal because n-k=2rho+1. Therefore the exact dual larger vector is

\[
             w^\vee=\rho(D_\rho u_{-e_k/\varpi})W_{\Pi^\vee},
                                                                  \tag{PT23}
\]

with larger character psi\^{}\{-1\}. There is no change of Pi's spherical vector. Jacquet--Liu Theorem 2.3 / JPSS section 2.7 now gives

\[
 Z_\rho(1/2)=
 \frac{C_\sigma}
 {\omega_\sigma(-1)^{n-1}\gamma(1/2,\Pi\times\sigma,\psi)
                   \lambda_-^\vee}
                  I_\rho^{-}(1/2;\Pi^\vee,\sigma^\vee,\psi^{-1}).
                                                                  \tag{PT24}
\]

The factor lambda\_-\^{}vee is the probability scalar PT6 in sigma\^{}vee. It comes from applying PT8 once to PT22; the remaining finite average is unchanged because the original dual newform is K\_r-fixed.

For an entirely explicit phase, write G\_-\^{}vee and omega\^{}vee for PT17 and the central character in the dual pair, and set S\_*=-(rho). Combining PT20 with gamma=epsilon L\^{}vee/L gives

\[
\boxed{
\begin{split}
 \frac{Z_\rho(1/2)}{L(1/2,\Pi\times\sigma)}
  ={}&\frac{C_\sigma (G_-^\vee)^{\rho+1}}
   {\omega_\sigma(-1)^{n-1}\,
       \varepsilon(1/2,\Pi\times\sigma,\psi)\,
       d_k\,(\lambda_-^\vee)^{\rho+1}}\\
   &\quad\cdot
       [(1-q^{-1})\omega^\vee(-\varpi^{-1})]^\rho
       q^{-\rho^2+\rho(\rho+1)/2-\rho k}.
\end{split}}                                                     \tag{PT25}
\]

All Euler factors and epsilon factors here are regular and nonzero at 1/2 under PT1--2 and temperedness. Their normalized epsilon moduli are one. PT7, PT17 and PT22 give exactly PT4. In particular the entire extra flag modulus is q\^{}\{-rho(k+rho+1)/2\}; no growing flag polynomial or unevaluated cancellation remains in this conditional fixed test.

\subsection{7. It is the split Bessel functional}\label{it-is-the-split-bessel-functional}

At the central j=rho, use block order (H,E,z,F). Permuting from the physical order (E,H,z,F) is integral. Any literal unit coefficients or opposite signs on the E/F chains are transported by an integral diagonal on those flag blocks; it fixes H,z, commutes with D tau, and has root modulus one. Thus it suffices to check the displayed unit character.

The Bessel N has upper E and F unipotents, an E-to-H lower block in this reordered basis, and blocks E-to-z, H-to-F, z-to-F, E-to-F. Its character is generic on the two flag chains and the final E-to-z and z-to-first-F roots, and trivial on the other displayed blocks. It is the split model in Gan--Gross--Prasad section 14, directly identified with the literal source in LIFT060 and LIFT067.

The central JPSS functional has this right N covariance. On E-to-H the change is X -\textgreater{} X+Y. On upper U\_E one factors its left upper unipotent and changes X by u\_E\^{}\{-1\}; its determinant is one and its Whittaker phase is the generic E-chain character. Right E-to-z is already a left upper column and gives its last E entry. The H-to-F and E-to-F blocks factor as upper columns whose adjacent character entries are zero; their phases are one. Right z-to-F gives its first F entry. Upper U\_F acts as its identical left upper block and gives the generic F-chain character. These root groups generate N and its character kills their commutators. PT11 gives simultaneous H invariance at s=1/2.

Thus the continued central functional belongs to the literal split Bessel Hom-space (with the consistently chosen inverse character if the period slot is conjugated). The explicit value PT25 makes it nonzero. Gan--Gross--Prasad Corollary 15.2, including their split quadratic algebra case, makes this Hom-space at most one-dimensional. Every nonzero normalization of that same Bessel functional therefore has nonzero value on this prescribed tensor. This uniqueness statement does not set a separately normalized matrix-coefficient functional equal to PT25.

\subsection{8. Physical character and Haar accounting}\label{physical-character-and-haar-accounting}

LIFT067 TC5 identifies the actual transported norm-minus-T character as having final coefficients c and varpi, rather than two units. It proves chi\_c(n)=chi\_1(DnD\^{}\{-1\}) and

\[
                \delta_N(D)=q^{-\rho(k+\rho+1)}.                 \tag{PT26}
\]

The nonvanishing just proved therefore applies to the original chi\_c-functional on the bare core translate tau, after the exact D transport. The strict T covector, its newform, the marked norm-minus-T line and the physical additive character have not changed. If a local matrix-coefficient period is defined by an absolutely convergent integral, its Haar change is exactly TC11. A stable/regularized version needs its own change-of-variables theorem; it is not supplied by meromorphic continuation of JPSS integrals alone.

As a numerical identity in the stated Whittaker convention, PT4 says

\[
 \delta_N(D)^{-1}
    \left|Z_\rho(1/2)/L(1/2,\Pi\times\sigma)\right|^2
                      =\frac{q^k}{(q^k-1)^2}.                  \tag{PT27}
\]

This displays the exact cancellation of the flag-root modulus. It does not identify a canonical GGP local form or its Hilbert-unit normalization. Those comparisons retain the Whittaker/Hilbert norms and local formal normalizing factors. Rational global coordinate transport still requires every finite and real factor, as in LIFT067.

\subsection{9. Actual-source scope and remaining conditions}\label{actual-source-scope-and-remaining-conditions}

\begin{longtable}[]{@{}
  >{\raggedright\arraybackslash}p{(\columnwidth - 4\tabcolsep) * \real{0.3333}}
  >{\raggedright\arraybackslash}p{(\columnwidth - 4\tabcolsep) * \real{0.3333}}
  >{\raggedright\arraybackslash}p{(\columnwidth - 4\tabcolsep) * \real{0.3333}}@{}}
\toprule\noalign{}
\begin{minipage}[b]{\linewidth}\raggedright
Input or conclusion
\end{minipage} & \begin{minipage}[b]{\linewidth}\raggedright
What is checked here or already checked
\end{minipage} & \begin{minipage}[b]{\linewidth}\raggedright
What is not inferred
\end{minipage} \\
\midrule\noalign{}
\endhead
\bottomrule\noalign{}
\endlastfoot
Split degree-one T, original additive character and integral unit coefficients & LIFT060 and LIFT067 give the literal coordinate dictionary & No new finite prime or level \\
Marked line, primitive covector and D tau & Exact physical character transport of LIFT067, plus coupled integral transport of LIFT030 & No free local test-vector replacement \\
Smaller strict T type & The literal nontrivial dilation character gives conductor one once the local representation is in the stated generic setting; K052/K053 record the precise row/column lines & Conductor one alone does not imply the nontrivial central character \\
Generic tempered smaller and spherical generic tempered larger representations & Explicit premises PT1--2 & DO046 was checked for the existing m=2 profiles, not every general-m Hyp3 record \\
Spherical vector & The bottom finite maximal level at T supplies a spherical vector in any constituent that occurs & A theta-derived spherical constituent is not thereby generic or tempered \\
Fixed local nonzero test & PT25, with exact Gauss and projection factors & No global central L-value or global Bessel/Mram period \\
Source normalization and arithmetic comparison & PT26--27 retain the physical root modulus and the Whittaker normalization & No canonical GGP scalar, arithmetic Read, STiesR/032 value, Host/SizesAt instance, or DP44 cost conclusion \\
\end{longtable}

In particular, later large theta descendants can have inserted non-tempered strings. They cannot be replaced silently by Pi satisfying the premises here. For a new general-m bottom family, finite genericity and temperedness must be proved from its actual automorphic constituents or imposed as a separate premise before this local theorem is applied. The mere cusp-space existence and real NCGaps statement in LIFT059 do not supply those finite premises. The qualitative global period is still a distinct requirement even if every local test is known nonzero.

\subsection{10. Primary and literal references}\label{primary-and-literal-references}

\begin{itemize}
\tightlist
\item
  Booker, A. R., Krishnamurthy, M., \& Lee, M. (2020). \emph{Test vectors for Rankin--Selberg L-functions}. Journal of Number Theory, 209, 37--48. Section 2 measure conventions; section 3 equations (3.1)--(3.2) and Lemma 3.1. Directly read primary text: \url{https://research-information.bris.ac.uk/files/211450903/gln_local.pdf}.
\item
  Jacquet, H., \& Liu, B. (2018). \emph{On the local converse theorem for p-adic GLn}. Section 2.2, equation (2.1), Theorem 2.3 and equation (2.2), citing Jacquet--Piatetski-Shapiro--Shalika (1983), section 2.7. Directly read author text: \url{https://archive.ymsc.tsinghua.edu.cn/pacm_download/318/8960-JL-LCT-final.pdf}.
\item
  Gan, W. T., Gross, B. H., \& Prasad, D. \emph{Symplectic local root numbers, central critical L-values, and restriction problems in the representation theory of classical groups I}. Directly read author version \nolinkurl{ggp7}, section 14 printed 49--51 and Corollary 15.2 printed 55: \url{https://people.math.harvard.edu/~gross/preprints/ggp7.pdf}.
\item
  Miyauchi, M. (2014). \emph{Whittaker functions associated to newforms for GL(n) over p-adic fields}. Journal of the Mathematical Society of Japan, 66(1), 17--24. Theorem 4.1, printed 21--22, directly reread here; the probability-projector normalization is the checked K053 NT10--14.
\item
  Accepted exact internal inputs: K053 NT7--28; LIFT060 HF3--14; LIFT030 NL3--12; LIFT067 TC1--16. The literal source anchors are HidaSetting.hJ/rootCoord/Nfin/psiFin and its fixed marked vector, already snapshotted in LIFT067. The present recurrence adds no source axiom or imported fixed-test existence assertion.
\end{itemize}

The companion exact diagnostic checks the integral pivot and resulting upper phase in arbitrary small ranks, last-row elimination, character orthogonality, all root/Jacobian powers and the closed-form exponent. Those checks supplement the universal algebraic proof. They do not numerically evaluate an actual global automorphic form or certify the unproved general-m finite hypotheses.

\clearpage

\section{AM-K089 --- exact eleven-adic content of the prescribed coupled auxiliary section}\label{am-k089-exact-eleven-adic-content-of-the-prescribed-coupled-auxiliary-section}

\textbf{Canonical checked version --- 2026-11-22T18:11:57.514698+00:00.} Root read the whole proof and complete script, checked the stated used interfaces, and ran a separate diagnostic; a saved independent whole review and separate replay also passed. Original pending wording below is historical. Promoted scope: Every actual full-Y auxiliary section and its finite class transport has unit content at11 in an integral algebraic frame. The fixed-scalar finite Heisenberg reading is an integral isomorphism on the prime-to11 denominator cover, and the all-Hasse-jet map through order4\^{}(g-1)g!D\^{}2 has a unit full-rank minor. The selected ordinary derivative, genuine type projector, class contraction and terminal arithmetic image remain separate.

Author: \nolinkurl{/root/kummer_extension}, 22 November 2026 UTC. New controlled pending finding. Status: author whole-proof and complete-script PASS; independent review pending.

This proves primitiveness of the \textbf{actual auxiliary section before its selected differential and class contractions}. It does not replace the allowed Y set, add a T-torsion parameter, refine the automorphic level, or assert primitiveness of the final arithmetic class-row map.

\subsection{1. The fixed object and the theorem}\label{the-fixed-object-and-the-theorem}

Fix either actual J14 or J56 host, a coefficient-field place above 11, and its complete valuation ring R. Retain the actual primitive trace-one theta section, the source finite Weil phases, the G/T/Q kernels and all chosen finite class frames of K038. Use K039(IG3) to choose those frames with component 1 at every place above 11; their literal delAtr/intNab transport remains the one already specified. Let \[
 g=f(ab+bc),\qquad
 X=A_h^{ab}\times A_h^{bc},\qquad
 \mathcal L=\mathcal L_0\boxtimes\mathcal L_1,
 \qquad \sigma=\sigma_q^{\boxtimes14h(ab+bc)}.
 \tag{PS1}
\] Here g is the dimension of X; the line is the actual product principal theta line. The notation is independent of the smaller polynomial source/target ranks. By LIFT08(Q4), sigma generates the integral principal section line at 11. The common complex comparison factor remains exactly \nolinkurl{C_theta^(14h(ab+bc))}, with the two source-creation profiles retained as in K035/K038. Neither is assigned a finite valuation.

Let Psi\_raw be the \textbf{whole prescribed finite two-edge tensor} of LIFT26(GS6--GS7c), equivalently K038(CA22--CA27), before the selected polynomial contraction. At a class triple (e,f,d), let Psi\_efd be its actual finite Weil transport. This is a function on the direct sum of the two actual rational Schrödinger position spaces. It is not an independently chosen member of an auxiliary Heisenberg orbit.

For each such function choose an even positive integer D, prime to 11, such that it is supported in \nolinkurl{D^-1 Lambda_hat} and invariant under \nolinkurl{D Lambda_hat} in a fixed integral symplectic position model. Such a D exists because the component at 11 is the actual standard vacuum. Enlarge it if necessary to include the fixed half-period of that vacuum. This is a common-cover device for the existing function, \textbf{not a change of its automorphic level or a new allowed Y parameter}.

For the finite comparison only, choose a finite unramified splitting extension R' of R for the prime-to-11 theta structure and its phases. The actual class section already descends to the carried coefficient field; this extension is not inserted into that field. The period-removed geometric reading on the common cover is denoted \[
 \Theta_D:\mathscr S_D(R')\longrightarrow
 H_D\otimes_R R',\qquad H_D:=H^0(X,[D]^*\mathcal L),
 \qquad
 \mathscr S_D(R')=\operatorname{Fun}(D^{-1}\Lambda/D\Lambda,R').
 \tag{PS2}
\] The position lattice Lambda has rank g. Equivalently, before identifying the common cover with X by division by D, its analytic periods are D times the original periods and its line is the pullback of the original theta line. Both modules in PS2 have rank \nolinkurl{D^(2g)}. The reading is the already fixed Mumford/Schrödinger reading of DO026 and K035, now restricted to this finite module; it is not defined by choosing an integral basis that fits analytic values.

For a vector in a free R-module, content means its coefficient ideal in any R-basis. For a finite function it is the ideal generated by its values. For a formal algebraic section coefficient it is the ideal generated by all its ordinary power-series coefficients, equivalently its Hasse jets in the chosen integral regular parameters. This definition does \textbf{not} divide an ordinary invariant derivative by an unmentioned factorial.

\textbf{Theorem PS3.} The actual map PS2 is an R'-linear isomorphism. Consequently \[
 \operatorname{cont}_{R'}\Theta_D(\phi)
   = (\phi(x):x\in D^{-1}\Lambda/D\Lambda)R'
 \tag{PS3}
\] for every phi in this finite module. For the actual functions and sections defined over R, the equality of content ideals descends faithfully flatly to R. The actual whole tensor and every actual transported class tensor satisfy \[
 \operatorname{cont}_R\Psi_{\rm raw}
 =\operatorname{cont}_R\Psi_{efd}
 =\operatorname{cont}_R\Theta_D(\Psi_{efd})=R.
 \tag{PS4}
\] In any integral algebraic formal frame at the origin, the coefficient series of this actual section therefore has content R. There is also an explicit finite all-Hasse-jet matrix with a unit full-rank minor: one may take the total jet order \[
                  B_D=4^{g-1}g!D^2.
 \tag{PS5}
\] Thus at least one Hasse coefficient of the prescribed auxiliary section of order at most B\_D is an 11-adic unit. This does not specify that coefficient as one of the fixed selected subD functionals.

\subsection{2. Integral finite Heisenberg matrix units and the fixed scalar}\label{integral-finite-heisenberg-matrix-units-and-the-fixed-scalar}

Work over the finite unramified splitting extension R' in PS2. All integral operators and reductions in this section are over R'. Such an extension exists: prime-to-11 torsion of the good abelian scheme is finite étale, its roots of unity are étale, and extracting a root of unit degree from a unit is étale. It is only a faithfully flat test for the actual sections' content; no new coefficient field is made part of the constructed family. Put \[
 C=(\mathbb Z/D^2\mathbb Z)^g,\qquad
 \widehat C=\operatorname{Hom}(C,\mu_{D^2}).
 \tag{PS6}
\] The finite Schrödinger module has its fixed delta-function basis delta\_x, x in C. Write T\_b for translation and M\_chi for multiplication by chi. With the convention \nolinkurl{T_b delta_y=delta_(y+b)}, the exact matrix units are \[
 E_y=\frac1{|C|}\sum_{\chi\in\widehat C}\chi(y)^{-1}M_\chi,
 \qquad E_{xy}=T_{x-y}E_y,
 \qquad E_{xy}\delta_z=\delta_{yz}\delta_x.
 \tag{PS7}
\] All denominators are units in R'. The same formulas hold after reduction modulo the maximal ideal. Distinct prime-to-11 roots retain their exact order, and the finite character orthogonality is valid in residue characteristic 11. In particular the reduced Schrödinger module is irreducible and its group algebra generates all its matrix units. This is the explicit argument underlying Mumford (1966, section 1, Proposition 3 and its proof, pp.295--296; the function model is on p.297).

We must verify that these are integral \textbf{actual} geometric operators, rather than replacing their scalar normalization. On the common cover their translations lie in the kernel of the polarization of \nolinkurl{[D]^*L}; in the original period coordinates this kernel is \nolinkurl{D^-1 Lambda_full / D Lambda_full}. Each of its points extends over R' because D is prime to 11. The theta-group fibre over such a point is a G\_m torsor over Spec R', hence has an integral lift, since Pic(R')=0. If u is an integral lift and the actual normalized geometric lift is c u, then a finite power of the latter is the prescribed root-of-unity scalar. The same power of u is a unit scalar. Thus a power of c is an R'-unit, and c itself has valuation zero. The retained Weyl central phases, including the fixed half-period, therefore give integral automorphisms with integral inverses. No phase has been set equal to 1.

The actual reading is equivariant for these operators by DO026(2a--4), including the rational representative and central phase. It sends the actual trace-one vacuum to \nolinkurl{[D]^*sigma}. The latter is primitive: sigma is primitive by LIFT08(Q4), and multiplication by D is finite étale and surjective on the special fibre, so its pullback of a nonzero section does not vanish. This pins the scalar of the reading. A generic statement that the two representations are isomorphic up to scalar would not suffice here.

The prime-to-11 orders in PS7 are the orders of the \textbf{Heisenberg translation and modulation groups}. They are not orders of finite general linear groups. A generic-character or irreducible-type idempotent can have denominators involving \nolinkurl{q^j-1}, which need not be 11-units. Such an idempotent is not inserted into PS2. The raw generic G function below is used with its literal values, and any later genuine e0 or irreducible projector remains outside the primitiveness assertion unless separately proved integral and nonannihilating on this section.

There is no hidden integrality assumption in passing from the vacuum to the entire finite module. The actual vacuum is a root-of-unity multiple of a half-period Heisenberg translate of the ordinary lattice indicator, as in K039(IG5--6). Choosing D even places this operation and its inverse inside the finite module. The ordinary indicator has unit coefficients in its delta basis; the matrix units PS7 generate every delta vector from such a primitive vector. Thus the R'-span of the actual finite Heisenberg orbit of the actual vacuum is all of S\_D. Every orbit section is integral by the preceding paragraph. The already defined characteristic-zero reading hence restricts to the R'-linear map PS2.

Reduce that map modulo the maximal ideal. It is nonzero, because it sends the vacuum to the nonzero pullback section. Its kernel is stable under all matrix units PS7, so irreducibility forces the kernel to be zero. The target has dimension \nolinkurl{D^(2g)} and formation of its section module commutes with base change, by the ample-line cohomology theorem for abelian schemes. Both source and target are free of that rank over R'. The reduced map is therefore bijective; its determinant is a unit, proving PS3 by Nakayama. This establishes the equality for the \textbf{specified} scalar, not merely an isomorphism of abstract representations.

The geometric facts used here are the same principal-polarization and base-change inputs as in LIFT08 and K034. For the finite representation and its section-space realization, Mumford (1966), section 1, Proposition 3, Theorem 2 and the following theta-structure construction, pp.295--298, explicitly work in every characteristic prime to the line's degree. The reduction at 11 is legitimate because the line in PS2 has degree \nolinkurl{D^(2g)}, prime to 11. No modular local theta-correspondence theorem at residue characteristic 11 is used.

The determinant assertion is made for the split comparison map over R'. For an actual section and function descending to R, faithful flatness then tests equality of their original content ideals. No descent of an arbitrarily chosen full theta basis to the original coefficient field is asserted. It adds neither a field-degree cost to the actual family nor a new choice of a section. Existing coefficient-field and frame costs remain unchanged.

\subsection{3. The prescribed complete finite Y tensor has unit content}\label{the-prescribed-complete-finite-y-tensor-has-unit-content}

The exact factorization is LIFT26(GS6--GS7c). At every middle class its allowed quotient is \[
 \mathcal Y\simeq
 \prod_G q^{-1}O^b/O^b\ \times\
 \prod_Q q^{-M}O^b/q^M O^b,
 \tag{PS8}
\] with \textbf{no T factor}. The two finite Weil actions and the complete local Y sum give one fixed tensor, with their original determinant characters and measure scalars.

At G the first raw factor forces the unique allowed class \nolinkurl{y=q^-1 e_a}. In the actual split coordinates, choose X with identity initial block, terminal Z with identity leading block and zero permitted lower block, and eta=0. The prescribed unipotent characters and determinant characters then all have value 1. This is K038(CA15--18), not an average over a finite general linear group.

At a Q place, the full local Y contraction is exactly LIFT23(Q10): on its stated support it has value \[
 H(Z)\,q^{-d}G(\chi,\psi_0)\chi(-c)^{-1}.
 \tag{PS9}
\] The actual chi is nontrivial. Its explicit supported input in Q13 is \nolinkurl{A=(I_a;0)}, \nolinkurl{C=0}, \nolinkurl{eta=q^-1 e_a^vee}, and Z with leading identity block. It gives the nonzero value \nolinkurl{q^-d G(chi,psi0) chi(-1)^-1}. The algebraic identity \nolinkurl{G conjugate(G)=q}, with both factors algebraic integers, proves that G is a unit at every coefficient prime over 11 because every admitted Q has q different from 11, by K039(IG2). Hence PS9 has content exactly R. This retains the actual \nolinkurl{q^(-Mb)} normalization and its full higher-depth Fourier cancellation. It introduces no \nolinkurl{(q-1)^-1} factor.

At T the function is the prescribed covector-incidence function of K038 section 5. Its support is nonempty: in integral bases adapted to the two primitive residue covectors, a split injection carrying the designated coordinate to the designated coordinate has the required pullback of covectors. Equivalently, a primitive target covector has an integral vector on which it equals 1 modulo T, which supplies a right inverse for the required residue row. The same construction applies separately to both Hom factors. Thus the original product incidence has a value 1. This is a support witness for the existing function; it does not vary the initial decoration or create a T summation variable.

At every remaining place the actual vacuum is a unit phase times a lattice-coset indicator, with a nonempty unit-valued support. The finitely many nontrivial local conditions are independent in the finite restricted tensor product. Taking the local witnesses gives an adelic value of the entire fixed tensor which is a unit. All its values are integral. If desired, additive approximation supplies a rational position vector in the same finite support/period cosets. No assertion about the theta value at the origin follows from this support observation alone.

Therefore \[
                         (\Psi_{\rm raw}(x):x)=R.
 \tag{PS10}
\] This uses the \textbf{summed} Q tensor and the unique surviving G summand. It does not infer primitiveness of a sum from primitiveness of its summands.

\subsection{4. Exact finite-frame transport and formal coefficients}\label{exact-finite-frame-transport-and-formal-coefficients}

K039(IG4--7) proves that the actual finite Weil transport away from 11 preserves the R-valued Schwartz module. Its generators are linear coordinate changes, roots of unity, determinant-character units and finite Fourier matrices whose measures and denominators are powers or square roots of the local residue cardinality. The inverse operators have precisely the same integrality property. The chosen local component at 11 is the identity. Hence, for every actual finite frame w, \[
                   (\omega_f(w)\phi(x):x)=(\phi(x):x).
 \tag{PS11}
\] Indeed each inclusion follows from one of the two integral operators. This argument also covers the actual ramified places 2 and 7, since both are units at 11. The infinite-order algebraic splitting character is not called a root of unity: its value is an 11-unit by K039(IG4).

Apply PS11 to the complete tensor PS10 and then PS3 to its actual support/period module. This proves PS4 for every actual class triple. The two-edge reading is the geometric section on the product Hom torus by linearity and the exact tensor factorization, with precisely the finite central phases of K035/K038. No per-class period has been introduced.

Let f\_phi(t) be the coefficient of Theta\_D(phi) in any integral algebraic formal frame at the origin and integral regular parameters t. Reduction of a global section to this formal neighbourhood is injective on the special fibre: a zero formal germ is zero in the local ring by the Krull-intersection theorem, hence zero on an open set and then on the geometrically integral abelian variety. Formation of H\_D commutes with base change. It follows, successively modulo each power of the uniformizer, that \[
       (\text{coefficients of }f_\phi) =
       \operatorname{cont}_R\Theta_D(\phi)=
       (\phi(x):x).
 \tag{PS12}
\] In particular the actual algebraic coefficient series has a nonzero reduction at 11.

This statement belongs to the algebraic frame. The actual reduced Chern/CM germ still has the carried form \[
             U(t)^\kappa\exp(-\kappa\beta(t))f_\phi(t),
                    \qquad \kappa=D^2\in R^\times,
 \tag{PS13}
\] on the common cover, or its original unit formal-coordinate pullback, in the already checked K034/K039 frame convention. An integral unit from the rigidified pullback comparison is absorbed into f\_phi and does not change PS12. Its Chern coefficients are not declared integral. The formal scalar content zero in PS12 is therefore fully compatible with the height-three slope and ordinary derivative divisibility of K027/K033/K039. The period, source creation factors, ordinary factorials and later Weyl contraction have not been removed from any selected arithmetic matrix.

\subsection{5. A literal finite all-Hasse-jet matrix}\label{a-literal-finite-all-hasse-jet-matrix}

For completeness the content assertion can be tested by a finite integral matrix, although its bound is deliberately coarse. Work on the geometric special fibre and identify the common cover with X, so \nolinkurl{[D]^*L} is algebraically equivalent to \nolinkurl{L^(D^2)}. The line \nolinkurl{L^4} is very ample: its polarization kernel contains X{[}4{]}, and Mumford (1966), section 4, Theorem 1 and its corollary, pp.338--340, give the assertion in the present characteristic 11. The symmetric fourth power meets the hypotheses there.

Let s be any nonzero section of \nolinkurl{[D]^*L} on the special fibre, with effective divisor Z. Choose g-1 hyperplanes in the \nolinkurl{L^4} embedding passing through the origin, successively general enough to meet Z properly and with linearly independent differentials at the origin. Their complete-intersection curve is smooth at the origin and has no component contained in Z. For g=1 take X itself. The local vanishing order of s along this curve is at least its maximal-ideal order at the origin. Positivity of the other local intersection multiplicities and the global intersection degree give \[
 \operatorname{ord}_0(s)
 \le Z\cdot(4c_1(L))^{g-1}
 =D^2 4^{g-1}c_1(L)^g
 =D^2 4^{g-1}g!=B_D.
 \tag{PS14}
\] The argument only uses geometric hyperplanes over an algebraically closed residue field and proper intersections; it requires no ordinary reduction or characteristic-zero generic smoothness theorem. The differentials can be chosen independently because the embedding is an immersion. Proper intersection with Z is ensured by avoiding its finitely many components at each cut. Base change back to the given residue field tests injectivity.

Thus the finite Hasse-jet map \[
 H_D\longrightarrow R[[t_1,\ldots,t_g]]/(t_1,\ldots,t_g)^{B_D+1}
 \tag{PS15}
\] is injective after reduction. In any integral basis its matrix has a unit \nolinkurl{D^(2g)}-by-\nolinkurl{D^(2g)} minor. It is consequently a split injection over the DVR. After the proof-only base change to R', composing with the actual finite isomorphism PS2 gives the promised finite integral comparison matrix. The injective section-jet map PS15 itself is defined over R and has a unit full-rank minor there. In particular the composed map evaluated at the fixed primitive vector Psi\_efd has a unit coordinate. This is an existence theorem for a unit minor of this explicitly specified matrix, not a numerical evaluation of an actual theta coefficient or an assertion that every minor is a unit.

No useful DP44 size estimate for D or B\_D is claimed. D is a denominator of the existing finite function and frames, and is unrelated to permission to vary a finite input. The comparison cover and the finite matrix prove content; they do not authorize using its delta basis as a replacement automorphic input. Any global coefficient-field, height or computational cost of displaying this matrix would still have to be included in a quantitative argument.

\subsection{6. The exact remaining image question}\label{the-exact-remaining-image-question}

The actual primitive section established above is an intermediate input to K038, not the output of its whole arithmetic row assembly. For the prescribed selected Weyl input and dual terminal coefficient, the map still applies the fixed ordinary derivative functional built from the two subD substitutions and their complete middle coevaluation, followed by the literal initial/middle class sums and the final branch and restriction. Those operations can have nontrivial reduction kernels and can increase content. In particular:

\begin{itemize}
\tightlist
\item
  PS15 controls \textbf{all Hasse jets up to B\_D} in integral algebraic parameters, not only the selected homogeneous ordinary invariant derivatives.
\item
  The geometric Chern gauge PS13 and its height-three Frobenius terms remain exactly those of K033/K039; no chosen auxiliary jet replaces f\_phi.
\item
  The class factors \nolinkurl{|Gamma_0|^-1}, \nolinkurl{|Gamma_1|^-1}, the actual Mram domain, and the literal terminal brC/intNab pairing remain as in K038, including any separately proved cancellation such as LIFT34.
\item
  Any later genuine e0 or other finite-type projector retains its own denominator and reduction kernel. The elementary inverse-integrality of prime-to-11 Weil generators in PS11 does not imply inverse-integrality of such projectors, whose finite group orders may contain factors of 11.
\item
  A unit raw auxiliary section does not imply a unit value after the actual selected derivative and class contraction, nor does it imply a nonzero restricted terminal arithmetic image.
\end{itemize}

The strict proved upper-content statement is PS4/PS12, namely \textbf{zero content for each prescribed full-Y auxiliary section, with the coefficients of its actual algebraic formal germ generating R}. The remaining Q1 pin is now the nonzero reduction of the specified \emph{selected derivative-and-class matrix applied to these primitive sections}, with its already established uniform lower gain removed. It cannot be discharged by replacing that matrix with PS15, a freely chosen torsion orbit, or a complex nonzero map.

\subsection{7. Sources, verification and scope}\label{sources-verification-and-scope}

The original mathematical input is the explicit finite matrix-unit argument with the actual vacuum scalar pinned, its application to the full unchanged finite tensor, and the finite-jet injection PS14--15. The existing source construction is retained through LIFT08(Q4), DO026(2a--4), K035(PF22--28), K038(CA15--27), K039(IG2--7), LIFT23(Q10--14 and the content paragraph after Q18), and LIFT26(GS6--7c).

Primary source: Mumford, D. (1966). On the equations defining abelian varieties. I. \emph{Inventiones Mathematicae, 1}(4), 287--354. \url{https://doi.org/10.1007/BF01389737} . The exact section 1 Proposition 3 proof and function model, Theorem 2, and theta-structure scalar statement were directly read at printed pp.295--298; the section 4 very-ampleness corollary and the relevant preceding argument were directly read at pp.338--340. Page images at printed pp.297 and 340 were checked. Public primary PDF: \url{https://pazuki.perso.math.cnrs.fr/index_fichiers/Mumford66.pdf} . This is not a claim to have checked every proof in the paper.

The companion exact certificate checks the finite matrix-unit and vacuum-generation identities in explicit cyclotomic coefficient rings reduced modulo 11, sample complete finite Gauss identities including a character order divisible by 11, the actual fixed-prime exclusions, and the combinatorial dimensions and factorial distinctions. It does not compute the actual enormous geometric jet matrix or certify the remaining selected-row reduction. No theorem about generic modular theta correspondences, a new family level, a DP44 saving, or RH is asserted.

\subsection{Author whole-file check}\label{author-whole-file-check}

At 2026-11-22T17:48:19.473762+00:00, I read the entire saved PS1--15 proof and all seven sections, and the complete companion script. I checked the comparison-cover rank and theta-group order separately, the proof-only splitting-extension scope, the exact primitive vacuum scalar, both directions of finite Weil integrality, the complete local Q contraction and its unit witness, the nonempty T incidence, and the formal versus selected-ordinary-derivative distinction. I reconstructed the finite divisor-intersection bound over residue characteristic 11 and checked that its matrix has a unit full-rank minor without choosing new finite input data. The final scope and normalization edits were reread. The first exact finite-model replay passed 14,960 assertions at 2026-11-22T17:45:33.397593+00:00; the final replay and fingerprints are in the separate verification record. The arithmetic replay does not compute actual geometric theta coefficients.

\clearpage

\section{AM-LIFT-069 --- the chosen bottom subspace has simple tempered base change}\label{am-lift-069-the-chosen-bottom-subspace-has-simple-tempered-base-change}

\textbf{Canonical checked version --- 2026-11-22T18:11:57.514698+00:00.} Root read the whole proof and complete script, checked the stated used interfaces, and ran a separate diagnostic; a saved independent whole review and separate replay also passed. Original pending wording below is historical. Promoted scope: Under the genuine LIFT059 package, every constituent of its chosen nonzero bottom subspace has a simple generic cuspidal GL\_N base change and is tempered at every finite place. Its original split-T spherical vector therefore satisfies the larger generic tempered premise of LIFT068. The full e0 space, initial Mram representation, missing source constructors and global period are not asserted.

Author: ALP\_MACHINE \nolinkurl{/root/lift_accounting}. Date: 22 November 2026 UTC. Status: author whole-proof/source/script review PASS; final exact hashes and replay are recorded in the separate author-review files. Independent review and canonical integration remain separate.

\subsection{1. Conclusion and exact scope}\label{conclusion-and-exact-scope}

Retain the genuine setting and every hypothesis of LIFT059, including its full source size hypotheses, its actual CM field E/F, even k and m, and its prescribed bottom discrete series. Put

\[
                 N=k+m-1,\qquad G=U(Y_0).
\tag{CT1}
\]

Thus N is odd. Retain the specific literal goodTuple and its finite tensor type constructed in LIFT059(BC12--13). At one selected place above each of the two G primes its regular polynomial has the two degrees p\_* and N-p\_*; at every other selected G place it is irreducible of degree N. Let \(\mathcal C_{\rm chosen}\) be the G(A)-module generated by the nonzero cusp space with this chosen finite type, the original finite level away G, and the prescribed bottom archimedean type. This denotes its actual automorphic constituents; it is not the entire space selected by the sum e0.

\textbf{Theorem.} Every irreducible automorphic constituent pi of \(\mathcal C_{\rm chosen}\) has a simple generic elliptic Arthur parameter. Its base change is a conjugate-selfdual cuspidal representation

\[
               \operatorname{BC}(\pi)=\Pi
                    \quad\text{of }\operatorname{GL}_N(\mathbb A_E).
\tag{CT2}
\]

The representation Pi is tempered at every finite place. The actual pi belongs to the corresponding local Langlands packet at every place, and those packets are tempered. In particular, at every selected split T place, the original maximal-compact-fixed bottom vector belongs to an irreducible tempered generic spherical representation of GL\_N(F\_v).

The theorem is conditional on the genuine LIFT059 input package; it does not construct a previously missing Host or full SizesAt instance. LIFT059 supplies a nonzero chosen cusp subspace when that package is supplied. The conclusion preserves its level, goodTuple, projector and actual real type. It verifies the \textbf{larger-representation} premise in LIFT068. It does not establish finite temperedness or genericity of every general-m initial Mram constituent, a nonzero global Bessel/Mram period, a Read/STiesR/032 instance, or a normalized global period cost.

\subsection{2. One actual full-degree G place forces local supercuspidality}\label{one-actual-full-degree-g-place-forces-local-supercuspidality}

There is a full-degree selected place: LIFT059 has f\textgreater=3, and only two out of the 2f selected places use the mixed-degree polynomial. Fix such a place v. It is split in E/F. In the prescribed split frame,

\[
 G(F_v)=\operatorname{GL}_N(F_v),\quad
 K=\operatorname{GL}_N(O_v),\quad
 K^1=1+\varpi_vM_N(O_v).
\tag{CT3}
\]

The chosen genuine finite series is an irreducible \textbf{cuspidal} representation rho of GL\_N(k\_v), inflated to K. LIFT059 section 4 proves this for the full-degree regular polynomial. The same assertion holds for its contragredient if the automorphic right-action convention requires that orientation. In every nonzero constituent of the generated chosen subspace,

\[
             \operatorname{Hom}_K(\rho,\pi_v)\ne0.
\tag{CT4}
\]

Indeed, the isotypic projection to rho commutes with the remaining adelic factors and with projection to any irreducible automorphic summand. A nonzero summand of the module generated by that isotypic space cannot have zero rho-isotypic part.

We use the elementary compact-invariants identity already proved in DO046(AP12a). If P is a parabolic of GL\_N with block ranks N\_1,\ldots,N\_l, and tau\_i are smooth representations of its Levi factors, then

\[
 \left(\operatorname{Ind}_P^{\operatorname{GL}_N}
                        \bigotimes_i\tau_i\right)^{K^1}
 \simeq
 \operatorname{Ind}_{\overline P}^{\operatorname{GL}_N(k_v)}
                        \bigotimes_i\tau_i^{K_{N_i}^1}.
\tag{CT5}
\]

For completeness, Iwasawa decomposition realizes the induced representation on K. Right K\^{}1 invariance forces the values to be fixed by the principal subgroups in each Levi block. The unipotent radical acts trivially on the inducing values, and the normalized-induction modulus equals one on P intersect K. Reduction modulo the maximal ideal gives (CT5). In characteristic zero the K\^{}1-fixed functor is exact, by compact averaging. Finite group modules are semisimple.

If pi\_v were not supercuspidal, it would be an irreducible subquotient of induction from a proper parabolic. Applying (CT5) and exactness would put rho in a proper finite parabolic induction. This contradicts the defining cuspidality of rho. Therefore

\[
                         \pi_v\text{ is supercuspidal}.
\tag{CT6}
\]

This proof does not infer a depth-zero inertial parameter from a polynomial name. It uses the actual finite representation and exact parabolic restriction. Since pi is automorphic and unitary, pi\_v is a unitary supercuspidal representation and is tempered. Its local Langlands parameter phi\_v has irreducible N-dimensional Weil representation and zero monodromy. The irreducible-Weil assertion is the usual GL\_N correspondence; a precise directly read primary anchor is Scholze(2013), Theorem 1.2(c)--(d), with his stated normalization twist. Irreducibility is unchanged by that twist.

\subsection{3. Actual inner-form theta comparison identifies the split parameter}\label{actual-inner-form-theta-comparison-identifies-the-split-parameter}

Weak transfer at almost all places would not identify the selected ramified place v. We make that step explicitly.

By Chen--Zou(2025), Theorem 2.1 and Proposition 4.1, pi has an elliptic Arthur parameter

\[
                       \Psi=\boxplus_i\Pi_i\boxtimes S_{d_i},
                 \qquad \sum_i n_i d_i=N.
\tag{CT7}
\]

Each Pi\_i is conjugate-selfdual cuspidal on GL\_(n\_i)(A\_E). Choose the auxiliary sufficiently large quasi-split unitary group of their section 3, using the same honest splitting characters in the local and global constructions. These characters exist in the retained genuine setting; the already supplied odd eta may provide the required parity characters. This auxiliary representation is used only in the proof of a classification statement. It is not an additional edge of the arithmetic chain and produces no asserted saving in its cost.

Write M for the auxiliary rank and L=M-N. Chen--Zou Theorem 3.1 gives discrete occurrence of the abstract theta lift

\[
                       \Sigma=\theta^{\rm abs}(\pi).
\tag{CT8}
\]

Their Proposition 4.1, including its proof and equation (4.2), identifies its global parameter as

\[
 \Theta(\Psi)=\Psi\chi_W^{-1}\chi_V
                        \boxplus\chi_V\boxtimes S_L.
\tag{CT9}
\]

At v, apply their local formula (6.1) and split-place Remark 6.4(2) to the \textbf{actual tempered} phi\_v from (CT6). The actual local theta lift Sigma\_v belongs to the singleton split GL\_M Arthur packet for

\[
 \Theta(\phi_v)=\phi_v\chi_{W,v}^{-1}\chi_{V,v}
                        \oplus\chi_{V,v}\boxtimes S_L.
\tag{CT10}
\]

The local and global occurrences refer to the same nonzero local theta lift: the abstract lift in (CT8) is the restricted tensor product of those local lifts, and local Howe uniqueness fixes its irreducible factor. No choice of a new vector, globalization or local test is involved.

On the other hand, the auxiliary group is quasi-split. Chen--Zou section 3.4 and Theorem 3.4 put Sigma\_v in the local Arthur packet for the localization of (CT9). These statements explicitly include possibly unbounded localizations, defined by parabolic induction. At a split place the packet is the GL\_M singleton attached to the associated Langlands parameter. The associated parameter is their formula (2.4): replace the Arthur SL2 factor by its diagonal absolute-value weights.

The two singleton descriptions have the same irreducible GL\_M representation, hence the same Langlands parameter. It suffices to use their semisimplified Weil representations. They have the identical added character string

\[
 \chi_{V,v}|\cdot|^{(L-1)/2},\ 
 \chi_{V,v}|\cdot|^{(L-3)/2},\ldots,
 \chi_{V,v}|\cdot|^{-(L-1)/2}
\tag{CT11}
\]

and the identical scalar twist \(\chi_{W,v}^{-1}\chi_{V,v}\). Cancel the string with its multiplicities in the semisimple category, then undo the common twist. Thus

\[
                 \phi_{\Psi_v}^{\rm W,ss}
                         \simeq\phi_v^{\rm W,ss}.
\tag{CT12}
\]

The signs, duals and twists have not been guessed from a differently normalized oscillator: (CT9) and (CT10) are the same two published occurrences of the Chen--Zou formula. Their Remark 6.4 expressly identifies the split correspondence with Mínguez's theorem.

For every summand in (CT7), the left side of (CT12) contains the direct sum

\[
 \bigoplus_{a=0}^{d_i-1}
     \phi_{\Pi_{i,v}}^{\rm W,ss}
                 |\cdot|^{(d_i-1)/2-a}.
\tag{CT13}
\]

Every displayed summand is nonzero. The right side of (CT12) is irreducible. Consequently there is exactly one global index i, and d\_i=1. Its rank is N. This proves (CT2), including both simplicity and genericity, without an assumption that all same-sign archimedean gaps exceed one. Chen--Zou Theorem 2.6 now applies and identifies the actual pi with a member of the local Langlands packet at \textbf{every} place. The order is essential: that all-place conclusion is invoked only after genericity has been proved.

As a primary consistency check, Chen--Zou's proof of Lemma 6.14, printed page 4805, explicitly uses a split supercuspidal place to force genericity of a global parameter and irreducibility to force simplicity in its quasi-split construction. The argument above supplies the additional inner-form comparison using their actual large theta lift.

\subsection{4. Actual rank parity and finite temperedness}\label{actual-rank-parity-and-finite-temperedness}

At a selected complex embedding tau, the bottom Harish--Chandra multiset is

\[
          (v_0,\ldots,v_{k-2},-b_0,\ldots,-b_{m-1}).
\tag{CT14}
\]

It has N=k+m-1 entries. Literal TParity says that every v\_i and b\_j is integral. Interlace makes the list strictly decreasing in the displayed order: successive v entries are separated by an intervening u entry; the last v exceeds u\_last, which exceeds -b\_0; and b\_j\textless e\_j\textless b\_(j+1). This remains true in the dual orientation after negation and reversal. Since k and m are even, N is odd and

\[
          h_i\in\mathbb Z=\mathbb Z+(N-1)/2,
          \qquad h_i-h_{i+1}\ge1.
\tag{CT15}
\]

The actual bottom representations are the genuine discrete series supplied by LIFT059/R015, so their archimedean parameters are bounded. By the all-place packet statement following (CT13), the unique GL\_N cuspidal base change has these same bounded complex parameters. We use Chen--Zou's conjugate-parity parameter convention. In Mok section 7.1, equations (7.1.2)--(7.1.3), this is the standard stable embedding with auxiliary character chi=1, so no additional exponent shift is inserted. The character \((z/\bar z)^{h_i}\) has holomorphic exponent h\_i. Hence

\[
                \lambda_i=h_i-\bigl((N-1)/2-i\bigr),
                   \qquad 0\le i<N,
\tag{CT16}
\]

is an integral dominant weight. Indeed \(\lambda_i-\lambda_{i+1}=h_i-h_{i+1}-1\ge0\). At the conjugate embedding the reversed negative list gives the corresponding conjugate coefficient weight. This is regular C-algebraicity of the cuspidal Pi. A unit adjacent HC gap is permitted: it produces equal neighboring highest-weight coordinates, not an irregular infinitesimal character.

Unlike the multiple-block argument in DO046(AP8), no eta half-parity correction is needed here. The sole GL block has the same rank N as the unitary group, so its required shift is zero. The actual global splitting characters in section 3 are a different matter and have not been removed.

Shin--Templier, Lemma 2.12 and Remark 2.13, imply that this cuspidal regular C-algebraic Pi is cohomological. Their qualification concerning arbitrary local C-algebraic representations is respected: we apply the statement to the global cuspidal Pi, and its infinity type here is already tempered. Caraiani(2012), Theorem 1.2, then proves finite temperedness at every place of E. The rank is N\textgreater=2. Boundedness of the corresponding local parameters and Chen--Zou Theorem 2.6 prove finite temperedness of pi as well.

\subsection{5. Application to the unchanged split T test}\label{application-to-the-unchanged-split-t-test}

LIFT059 retains maximal finite level at T in the bottom variable. Thus the bottom local component at each selected T place has a nonzero maximal-compact-fixed vector. The place is split; its actual group is GL\_N(F\_v). Section 4 proves this component tempered. Every irreducible tempered GL\_N representation is generic, by the standard classification of tempered GL\_N representations as irreducible inductions of discrete series and Whittaker heredity. This is the ordinary GL\_N classification used in DO046 and LIFT068, not a claim that a finite generic compact type alone implies p-adic genericity.

The spherical Whittaker function has a nonzero value at the identity and can be normalized to W(1)=1. That is exactly the larger-vector convention of LIFT068. It changes no physical automorphic vector without the accompanying scalar: the conversion from any fixed Petersson-normalized spherical vector retains its actual nonzero Whittaker value. The physical D and tau, additive character and Haar factors remain the ones in LIFT067--068.

The hypotheses are now separated as follows.

\begin{longtable}[]{@{}
  >{\raggedright\arraybackslash}p{(\columnwidth - 2\tabcolsep) * \real{0.5000}}
  >{\raggedright\arraybackslash}p{(\columnwidth - 2\tabcolsep) * \real{0.5000}}@{}}
\toprule\noalign{}
\begin{minipage}[b]{\linewidth}\raggedright
Premise in the fixed T calculation
\end{minipage} & \begin{minipage}[b]{\linewidth}\raggedright
Status for the current genuine construction
\end{minipage} \\
\midrule\noalign{}
\endhead
\bottomrule\noalign{}
\endlastfoot
Actual split field, rank N=k+2rho+1, physical D and tau & LIFT060/067, with rho=m/2-1 \\
Larger representation spherical & Original bottom maximal finite level at T \\
Larger representation irreducible, generic and tempered & Proved here for every constituent of the chosen LIFT059 subspace \\
Larger W(1)=1 & Exact normalization of its nonzero spherical Whittaker function; conversion scalar retained \\
Smaller initial representation generic and tempered & Not proved here for arbitrary general-m Mram constituents \\
Smaller conductor and central conductor equal one, original newform line & LIFT068's explicit premise; the accepted concrete m=2 realization satisfies it after DO046/K052 \\
Local fixed corrected test nonzero when all premises hold & LIFT068 \\
Prescribed global Bessel/Mram period nonzero & Still open \\
\end{longtable}

No statement is made here about every constituent of the full e0 sum. The proof uses a full-degree finite cuspidal factor available in the chosen nonzero subspace. The remaining summands of e0 may have no such factor. Likewise, no assertion about a global initial parameter follows merely from the strict T character.

\subsection{6. Sources and verification boundary}\label{sources-and-verification-boundary-1}

The following primary passages were directly read for this proof from the retained original texts.

\begin{itemize}
\tightlist
\item
  Chen, R., \& Zou, J. (2025). \emph{Arthur's multiplicity formula for even orthogonal and unitary groups}. Journal of the European Mathematical Society, 27, 4769--4843. \url{https://doi.org/10.4171/JEMS/1564}. Theorem 2.1 and equation (2.4), pp.~4777--4778; Theorem 2.6, p.~4780; Theorem 3.1, p.~4783; section 3.4 and Theorems 3.3--3.4, pp.~4784--4786; Proposition 4.1 and its proof, pp.~4788--4790; equation (6.1), Theorem 6.3 and split Remark 6.4(2), pp.~4796--4797; proof of Lemma 6.14, p.~4805. Official PDF: \url{https://ems.press/content/serial-article-files/51452}.
\item
  Scholze, P. (2013). \emph{The local Langlands correspondence for GL\_n over p-adic fields}. Inventiones Mathematicae, 192, 663--715. Theorem 1.2(c)--(d). \url{https://doi.org/10.1007/s00222-012-0420-5}. The retained primary text is \nolinkurl{local_t_branch_refs/scholze2013_local_langlands.txt}.
\item
  Mok, C. P. (2015). \emph{Endoscopic classification of representations of quasi-split unitary groups}. Memoirs of the American Mathematical Society, 235(1108). Section 7.1, equations (7.1.2)--(7.1.3), printed page 190 of the directly read author version. \url{https://arxiv.org/abs/1206.0882}. The auxiliary-character shift is explicitly retained in that source and is zero for the standard stable parameter used here.
\item
  Shin, S. W., \& Templier, N. (2014 author version). \emph{On fields of rationality for automorphic representations}. Lemma 2.12 and Remark 2.13, pp.~7--8. \url{https://math.berkeley.edu/~swshin/FieldRat.pdf}. This directly read implication cites Clozel; no new claim to have read the original Clozel volume is made.
\item
  Caraiani, A. (2012). \emph{Local-global compatibility and the action of monodromy on nearby cycles}. Duke Mathematical Journal, 161(12), 2311--2413. Theorem 1.2, p.~2312. \url{https://doi.org/10.1215/00127094-1723706}. Author PDF: \url{https://www.ma.imperial.ac.uk/~acaraian/papers/lgc1.pdf}.
\end{itemize}

Accepted reports read directly: LIFT059 section 4/BC12--13 (actual chosen finite type and nonzero cusp subspace); DO046 AP12a (exact finite parabolic invariants, reproduced above) and AP8--10 (the parity distinction and primary cohomological interface); LIFT068's exact hypothesis table.

Literal source reads: \nolinkurl{Formalism/LocalRingCardinality/Params.lean}8--22 and 51--72 (rank parity and HC grids); \nolinkurl{Formalism/LocalRingCardinality/DataTheta.lean}529--552 (Interlace and Dpar); \nolinkurl{Formalism/SharpRestriction/CuspSpanNormBounds/CayleySignature.lean}105--118 (TParity, NCGaps and Hyp3). NCGaps bounds unlike signs; it is never used here as a bound for every adjacent HC difference.

The companion exact script checks rank/parity, unit-gap coefficient weights, dual orientation, complete common-string cancellation with multiplicities, and the arithmetic content of the nonzero-block count. These are diagnostics of the written proof. They do not compute an automorphic representation, prove the cited classification theorems, or evaluate a global period. No Lean build is made.

\clearpage

\section{AM-DO-079 --- the actual selected operator obstruction and the remaining trace weights}\label{am-do-079-the-actual-selected-operator-obstruction-and-the-remaining-trace-weights}

\textbf{Canonical checked version --- 2026-11-22T18:14:48.957529+00:00.} Complete root and independent proof/script reviews and separate1,833-check replays passed. Original pending wording below is historical. Promoted scope: On the actual all-top J56 selected compact type, a nonzero minimum-coordinate subspace has only ladder indices0 through12. At r=1/(10 sqrt(h\_*)), its actual compressed operator norm is at least7/8, refuting the proposed exponent1/3 bound with constant12. Its extremal trace weight is less than2\^{}-158149. The true selected trace obeys the compact-twist-preserving weighted ladder inequality SO19; the full overlap weights remain unknown. This removes no arithmetic support sector and changes no error charge or main finite gate.

Author: GPT-6 Astra Pro, ALP\_MACHINE \nolinkurl{/root/definite_oscillator}. 22 November 2026 UTC. Controlled recovery-pending output; canonical promotion is separate.

\textbf{Correction incorporated:} AM-DO-079-C1 corrects the all-word identity in section5. The original frozen proof and earlier reviews are preserved, and their unqualified PASS verdicts are superseded to this precise extent. The operator/trace conclusions and error budgets are unchanged.

\subsection{1. Scope and the exact existing model}\label{scope-and-the-exact-existing-model}

Retain the actual all-top J56 assignment at the K044 upper endpoint, every allowed interior word and canonical twist, and the K043 selected isometry. At one real label put a=168. K043(PB2--8) gives the unitary holomorphic discrete series with minimal compact type V\_alpha tensor z\^{}beta and polynomial degree spaces

\[
 \mathcal H_N=V_\alpha\otimes\operatorname{Sym}^N(\mathbf C^a),
 \qquad W=V_\lambda\otimes z^{\beta-N}\subset\mathcal H_N.
 \tag{SO1}
\]

The actual Hilbert metric is Nakahama's metric, including the K043 PBW isometry. No common Fischer metric is substituted on the sum of distinct Pieri types. The selected highest weights satisfy m\_i=lambda\_i−alpha\_i\textgreater=0 and N=sum m\_i. The exact all-top last increment is

\[
             m_{a-1}=12,\qquad N>2^{959}.
 \tag{SO2}
\]

These are the actual RC48.1 parameters, not an auxiliary scalar-valued representation. A compact determinant shift and the canonical twist leave SO2 unchanged.

Fix the first positive coordinate and the negative coordinate. Use exactly the K060 operators B,C=B*,H\_1={[}C,B{]}, with C=partial\_1 in the holomorphic model. K060(RC2--6) supplies an orthogonal sum of actual unitary SU(1,1) ladders. In a ladder with lowest degree N0 and energy kappa, its degree-N line has j=N−N0 and

\[
 H_1e_j=(2j+\kappa)e_j,\quad
 Be_j=\sqrt{(j+1)(j+\kappa)}e_{j+1},\quad
 Ce_j=\sqrt{j(j+\kappa-1)}e_{j-1}.
 \tag{SO3}
\]

The group Cartan parameter is the one for which its noncompact two-plane has entries cosh(r),sinh(r). The infinitesimal generator is B−C up to sign; that sign will not affect this proof.

\subsection{2. A nonzero subspace of the actual selected type}\label{a-nonzero-subspace-of-the-actual-selected-type}

Let W\_min be the eigenspace of the first positive compact coordinate with its smallest weight lambda\_(a−1). It is nonzero: apply a permutation matrix to a highest-weight vector of V\_lambda. Write

\[
 h_*:=N+\lambda_{a-1}-\beta
     =N+(\alpha_{a-1}-\beta)+12.
 \tag{SO4}
\]

Every vector of W\_min has this same H\_1 energy, by K060(RC3). In a monomial z\^{}gamma v of this compact weight, if nu\_1 is the corresponding first fibre weight, then

\[
       \gamma_1+\nu_1=\lambda_{a-1},\qquad
       \nu_1\ge\alpha_{a-1},\qquad 0\le\gamma_1\le12.
 \tag{SO5}
\]

The fibre inequality is the one already used in K060(RC4): after its common determinant shift, V\_alpha is a polynomial Schur module and all coordinate weights are nonnegative. It is valid for all fibre weights, not only the extremal vector. Thus

\[
                    C^{13}W_{\min}=0.                 \tag{SO6}
\]

In the orthogonal decomposition SO3 this means that every ladder component of W\_min has j\textless=12. Indeed C\^{}13 is nonzero on every ladder line with j\textgreater=13, and its images in distinct orthogonal invariant ladders cannot cancel. The decomposition retains multiplicity spaces. Consequently every occurring component has

\[
                   0\le j\le12,\qquad \kappa=h_*-2j.
 \tag{SO7}
\]

This is a statement about the actual selected W. It uses neither a scalar-fibre replacement nor a change of its projector.

\subsection{3. A quantitative lower bound for the actual compression}\label{a-quantitative-lower-bound-for-the-actual-compression}

Put X=B−C. On a line in SO7, the two terms in Xe\_j have different homogeneous degrees, so

\[
 \|Xe_j\|^2=(j+1)(j+\kappa)+j(j+\kappa-1)
            =(2j+1)h_*-2j(j+1)\le25h_*.
 \tag{SO8}
\]

Orthogonality of the invariant ladders gives the same inequality for every unit w in W\_min. For any such analytic vector, unitarity and integration of its derivative imply

\[
 \|\pi(a_r)w-w\|\le |r|\|Xw\|,
 \qquad \operatorname{Re}\langle\pi(a_r)w,w\rangle
       \ge1-\tfrac12r^2\|Xw\|^2.
 \tag{SO9}
\]

Take the explicit radius r\_\emph{=(10 sqrt(h\_}))\^{}(-1). Then

\[
 \boxed{\quad
 \operatorname{Re}\langle\pi(a_{r_*})w,w\rangle\ge7/8,
 \qquad \|P_W\pi(a_{r_*})|_W\|\ge7/8.
 \quad}                                                     \tag{SO10}
\]

Here P\_W is the actual orthogonal projector and w is an actual unit selected vector. In particular this conclusion includes all K043 positive norm factors exactly. The companion certificate verifies h\_\emph{\textgreater2\^{}959 and h\_}\textless2N throughout the allowed word/twist ranges.

\subsubsection{Consequence for proposed uniform pointwise improvements}\label{consequence-for-proposed-uniform-pointwise-improvements}

Suppose a pointwise compression estimate with the same radial factor had the form

\[
 \|P_W\pi(a_r)|_W\|
 \le {C\over\sqrt{2\sinh r}\ h_*^{1/4+\delta}}\qquad(r>0).
 \tag{SO11}
\]

Then SO10 and sinh(r\_\emph{)\textgreater=r\_} force

\[
                 C\ge {7\over8\sqrt5}\,h_*^\delta.
 \tag{SO12}
\]

Thus any proposed improvement of the energy exponent with this radial factor must pay the explicit large constant in SO12 on these actual representations. Quantitatively, replacing 1/4 by 1/3 with C=12 would at r\_* give an upper bound

\[
       12\sqrt5\,h_*^{-1/12}<2^{-74},                       \tag{SO13}
\]

contradicting SO10. If the denominator instead uses the current common lower energy N+kappa\_0=h\_\emph{−12, its ratio to h\_} is greater than1/2; the same candidate is less than2\^{}(-73) and is still contradicted. Higher ambient dimension does not remove W\_min. The finite actual family alone is not an asymptotic no-go for an unspecified arbitrarily large constant.

This is an obstruction to an operator estimate valid for all real group elements. The element a\_(r\_*) is not asserted to be an actual rational pencil correspondence of support at least13. No arithmetic support obstruction, signed-kernel conclusion, or lower bound for the full normalized trace follows from SO10.

\subsection{4. Exact full Jacobi parameters are retained}\label{exact-full-jacobi-parameters-are-retained}

K060(RC9) identifies the diagonal coefficient on each ladder with

\[
 d_{\kappa,j}(r)=\cosh(r)^{-\kappa}
 P_j^{(0,\kappa-1)}(1-2\tanh^2r).
 \tag{SO14}
\]

The lower-bound proof did not discard these parameters. It identifies actual selected vectors for which j is bounded by12 while kappa grows like h\_*. Therefore an estimate derived only from the claim that the selected polynomial degree is large cannot assume that the Jacobi degree j is large. Haagerup--Schlichtkrull's primary Theorem1.1 and the discussion immediately following it distinguish the uniform varying-parameter bound from fixed-parameter large-degree estimates. Their theorem is used exactly as in K060; no fixed-parameter asymptotic is extended to SO7.

\subsection{5. The obstructing subspace has a very small trace weight}\label{the-obstructing-subspace-has-a-very-small-trace-weight}

The smallest first-coordinate weight space of V\_lambda is, as a module on the other a−1 coordinates, the irreducible Schur module of highest weight (lambda\_0,\ldots,lambda\_(a−2)), including the common determinant character. This is the extremal weight-space instance of the usual compact branching rule. Equivalently, subtract lambda\_(a−1) from every weight; the zero-weight space for the omitted coordinate consists of the same Schur functor on the remaining a−1 coordinates. The standard Weyl dimension formula therefore gives the exact ratio

\[
 p_{\min}:={\dim W_{\min}\over\dim W}
  ={(a-1)!\over
    \prod_{i=0}^{a-2}(\lambda_i-\lambda_{a-1}+a-1-i)}.
 \tag{SO15}
\]

For the actual all-top profile, with A=2\^{}960 and B=2\^{}955, the actual RC48.1/K043 formulas, corrected by AM-DO-079-C1, give

\[
 \lambda_0-\lambda_{a-1}=A+B,
 \qquad
 \lambda_i-\lambda_{a-1}=B+4(a-i)\ell+i-84-\ell\omega_i\quad(1\le i<a-1).
 \tag{SO16}
\]

The fibre starting weight stays fixed when an interior bit changes; m\_i decreases by ell, so the displayed selected difference decreases by ell. Canonical twists cancel. For either bit value the excess above B is at least7ell−83\textgreater0. Thus the corrected SO16 holds throughout the actual allowed family and each difference is at least B. The earlier word-invariance claim and its diagnostic were incorrect; AM-DO-079-C1 preserves that error record and supplies this correction. Since 167!\textless2\^{}(8*167),

\[
                  0<p_{\min}<2^{-158149}.                \tag{SO17}
\]

This tiny normalized dimension is why SO10 is not a no-go theorem for a stronger normalized trace estimate. Its smallness also does not bound the projection of all of W onto all low-j ladders: that is a different, presently unevaluated overlap problem.

\subsection{6. The exact trace inequality that survives the compact twists}\label{the-exact-trace-inequality-that-survives-the-compact-twists}

Let e\_l be any orthonormal ladder-line basis of H\_N and let d\_l(r) denote SO14 on its line. The degree compression D\_N(r)=P\_N pi(a\_r)P\_N is diagonal in this basis by K060(RC5--6). For g=k1 a\_r k2, cyclicity and K-stability of W give

\[
 {\operatorname{Tr}_W(P_W\pi(g)|_W)\over d_W}
 ={1\over d_W}\operatorname{Tr}_W
       \bigl(\pi_W(k_2k_1)P_WD_N(r)|_W\bigr),\qquad d_W=\dim W.
 \tag{SO18}
\]

The compact factor in SO18 is the actual one. Replacing it by the identity would change the selected coefficient. The trace-norm triangle inequality applied to the rank-one decomposition of the compression yields the safe estimate

\[
 \left|{\operatorname{Tr}_W(P_W\pi(g)|_W)\over d_W}\right|
 \le\sum_l w_l\,|d_l(r)|,
 \qquad
 w_l={\|P_We_l\|^2\over d_W}\ge0,\quad\sum_lw_l=1.
 \tag{SO19}
\]

This is not an operator-order assertion about absolute values. It follows directly from P\_W D\_N P\_W=sum\_l d\_l \textbar P\_We\_l\textgreater\textless P\_We\_l\textbar{} and the trace norm of each rank-one summand. Degenerate ladder choices do not change the corresponding grouped spectral weight. The K043 isometry fixes P\_W and hence fixes these weights.

Applying SO19 label by label gives the same inequality for the genuine product normalized trace at g\_p\^{}(-1) h g\_p.~It permits stronger bounds if the actual overlap weights are controlled, or if a genuine argument retains cancellation with the actual compact factors. Neither SO15 nor the present proof computes the full w\_l distribution. In particular a scalar holomorphic trace formula, a full-H\_N trace, or a radial signed Jacobi sum cannot silently replace SO18.

\subsection{7. Verification and limitations}\label{verification-and-limitations}

The exact companion script checks the actual K044/RC48 parameters, the twelve-step endpoint, all affine word/twist extremal bounds, SO8's polynomial identity, the quantitative contradiction in SO13, SO16--17, and small exact dimension diagnostics. It does not automate the compact branching theorem or the Hilbert-space argument. An author whole-file read, script read and separate replay are recorded separately after this text is frozen; independent review is separate.

This result supplies an actual selected operator obstruction and an exact reduction of the true trace question. It removes no support\textgreater=13 arithmetic sector, changes no earlier absolute error, and makes no assertion about the remaining signed inequality, Mram, a prescribed period/cycle, the source constructors, DP44, or RH.

\subsection{References and exact dependencies}\label{references-and-exact-dependencies}

Haagerup, U., \& Schlichtkrull, H. (2014). Inequalities for Jacobi polynomials. \emph{The Ramanujan Journal, 33},227--246. \url{https://doi.org/10.1007/s11139-013-9472-4}. Primary Theorem1.1 and its discussion, printed228--229, read in \nolinkurl{agents/selected_rank_one_refs/haagerup_schlichtkrull2014.txt}; primary PDF also directly retrieved.

Nakahama, R. (2016). Norm computation and analytic continuation of vector-valued holomorphic discrete series representations. \emph{Journal of Lie Theory, 26},927--990. \url{https://jolt.centre-mersenne.org/item/10.5802/jolt.919.pdf}. Theorem5.1 and equation(5.5), as translated and checked in AM-K043(PB1--8).

Accepted exact model dependencies: AM-K043(PB2--8,PB11--13), AM-K060(RC1--13), AM-K085§2's actual energy refinement, and reconstructed AM-R048(RC48.1). No theorem about a different compact type is substituted.

\clearpage

\section{AM-DO-079-C1 --- correction of the interior-word dependence}\label{am-do-079-c1-correction-of-the-interior-word-dependence}

\textbf{Canonical checked correction --- 2026-11-22T18:19:14.661153+00:00.} Root and independent peer read the entire correction, literal source identities and complete corrected script; separate1,999-check replays passed. This supersedes the earlier word-invariance sentence and the corresponding overbroad author/root/peer PASS endorsements. All other stated conclusions survive by the corrected argument below.

Author: ALP\_MACHINE \nolinkurl{/root/definite_oscillator}. 22 November 2026 UTC. This separate correction preserves the frozen DO079 proof and records precisely what its earlier whole reviews missed.

\subsection{1. The incorrect statement and the literal formulas}\label{the-incorrect-statement-and-the-literal-formulas}

DO079 section 5, immediately after SO16, says that an interior bit increases the corresponding fibre starting weight while decreasing the Pieri increment, leaving the selected weight difference unchanged. The companion diagnostic implements that assertion by replacing \nolinkurl{(r_i,m_i)} with \nolinkurl{(r_i+ell,m_i-ell)}. That assertion is incorrect for the actual K044 family.

K043(PB12--14), K044(EG5), and K046(BL13--14) instead give

\[
 r_i=E+(4a-3-4i)\ell+q+1,\qquad
 m_i=4\ell-1-\ell\omega_i\quad(1\le i\le a-2),
\]

with \nolinkurl{r_i} independent of the word. The endpoints are word-independent. The actual word change is \nolinkurl{(r_i,m_i) -> (r_i,m_i-ell)}. All source statements just cited were directly reread when this correction was identified.

Consequently SO16 is the correct zero-word identity. Its correct all-word version is

\[
 \lambda_0-\lambda_{a-1}=A+B,\qquad
 \lambda_i-\lambda_{a-1}
 =B+4(a-i)\ell+i-84-\ell\omega_i
 \quad(1\le i<a-1),
\]

where \nolinkurl{a=168}, \nolinkurl{A=2^960}, \nolinkurl{B=2^955}. The canonical twist still cancels from these differences, since it changes the common shift, not the actual word term.

\subsection{2. The stated conclusions are unchanged}\label{the-stated-conclusions-are-unchanged}

For both values of every independent bit,

\[
 4(a-i)\ell+i-84-\ell\omega_i
 \ge (4(a-i)-1)\ell+i-84>0.
\]

Indeed \nolinkurl{1 <= i <= 166}, so the coefficient of ell is at least 7, while \nolinkurl{i-84 >= -83}; the actual ell is larger than 83. Thus every difference used in SO15 remains at least B. The exact Weyl ratio SO15 and its uniform bound SO17,

\[
 \frac{\dim W_{\min}}{\dim W}<2^{-158149},
\]

are unchanged. The numerator estimate \nolinkurl{167! < 2^(8*167)} and all 167 denominator lower bounds are exactly the same.

The last increment is still 12; the total degree still decreases by ell for each one-bit. SO2--13 already used those correct degree and last-increment statements. The actual minimum-weight argument, the SU(1,1) ladder conclusion, the compression lower bound, the quantitative obstruction, and the trace-norm formula SO18--19 therefore require no change. No new arithmetic sector was removed by DO079, and this correction changes no error budget or support floor.

\subsection{3. Verification and provenance}\label{verification-and-provenance}

The separate corrected script is a copy of the frozen DO079 script with only the following changes: it reads this correction, writes its own output, records finding DO079-C1, and replaces the false one-bit cancellation check with the actual one-bit selected difference and its lower bound by B. Its other rank, degree, ladder and exact Weyl diagnostics are unchanged. Checking each independent interior bit proves the displayed coordinate bounds simultaneously for every word; no enumeration of all words is asserted.

The earlier DO079 author and independent whole-PASS conclusions are corrected to this extent: they did not identify the false word-invariance sentence or its corresponding diagnostic. The mathematical conclusions listed in section 2 remain valid by the corrected argument. The original frozen artifacts are retained, and this notice must accompany any use of their all-word statement.

\clearpage

\section{AM-LIFT-070 --- full source sizes force initial temperedness and the actual T newform}\label{am-lift-070-full-source-sizes-force-initial-temperedness-and-the-actual-t-newform}

\textbf{Canonical checked version --- 2026-11-22T18:29:32.915096+00:00.} Complete root and independent proof/source/script reviews and separate100,061-check replays passed. Original pending wording below is historical. Promoted scope: Under the genuine full SizeData2/SizesAt source package, tau2(0) is an actual common non-top label with initial HC gaps greater than1 for every word. Every existing initial constituent consequently has generic global parameter and is tempered at every place, with the honest eta rank-parity correction retained. Each actual split-T Mram line is the original conductor-one newform with central conductor1. Together with chosen-bottom LIFT069, both representation premises of the fixed LIFT068 D*tau test are supplied. Source existence, the full e0 space, other positive-rank local tests, canonical Hilbert/GGP normalization, the global period, Read/032 and DP44 remain separate.

Author: \nolinkurl{/root/lift_accounting}, 22 November 2026 UTC. Controlled pending finding. Author whole proof, source-interface and script review PASS; final hashes and the exact replay are recorded in the companion author-review JSON and manifest. Canonical promotion and independent review are separate events.

This theorem concerns a genuine automorphic realization of a literal fully sized source family. It does not assert that such a Host/D7/SizeData₂ family has been constructed. The new point is an exact source implication: \textbf{one common non-top real label has wide initial Harish--Chandra gaps for every word}. The global parameter argument then proves finite temperedness of every existing initial constituent. In conjunction with LIFT069, this supplies both representation-theoretic hypotheses in the conditional actual T test of LIFT068.

The theorem is not an implication from NCGaps alone, a finite generic type alone, or an abstract formal record without its genuine automorphic realization. It proves no global Bessel/Mram period or requested 032 value.

\subsection{1. Genuine setting and the exact source premises}\label{genuine-setting-and-the-exact-source-premises}

Let D be the literal D7 record and L a SizeData₂ D, and assume SizesAt L. Retain a genuine realization over its CM extension E/F and the actual compact-at-infinity initial group

\[
                     H=U(I_k),\qquad k=D.R.k.
\tag{SG1}
\]

At a source word w and CM label tau, the initial compact type has Harish--Chandra coordinates \nolinkurl{(D.hcW w tau).u}, or the reversed negative list in the dual automorphic convention. These are the actual R014/R015 compact types. Let pi be any irreducible automorphic constituent containing that type and the original finite initial level. No Steinberg, finite irreducibility, local genericity or global simplicity is assumed. Since H is anisotropic over F, its automorphic quotient is compact and these constituents belong to its discrete unitary spectrum.

Use zero-based coordinates and put

\[
 u_i^{w,\tau}=(D.\mathrm{hcW}\,w\,\tau).u_i,
 \quad v_i^\tau=(D.\mathrm{hc}\,\tau).v_i,
 \quad \Delta=D.s.\mathrm{Del}>0,
 \quad S=L.\mathrm{Sc}=4k\Delta.
\tag{SG2}
\]

The exact interfaces read for this implication are as follows. Paths in this table are relative to \nolinkurl{work/research/alpoge_source}.

\begin{longtable}[]{@{}
  >{\raggedright\arraybackslash}p{(\columnwidth - 2\tabcolsep) * \real{0.5000}}
  >{\raggedright\arraybackslash}p{(\columnwidth - 2\tabcolsep) * \real{0.5000}}@{}}
\toprule\noalign{}
\begin{minipage}[b]{\linewidth}\raggedright
Literal source
\end{minipage} & \begin{minipage}[b]{\linewidth}\raggedright
Content used
\end{minipage} \\
\midrule\noalign{}
\endhead
\bottomrule\noalign{}
\endlastfoot
\nolinkurl{Formalism/LocalRingCardinality/Params.lean}8--22, 27--37, 51--72 & k,m even, m\textgreater=2, k\textgreater=m+3; N\_j=k+j(t+1); strict u/v interlacing; u in Z+1/2 and v in Z \\
\nolinkurl{Formalism/LocalRingCardinality/RhoDetPrimeEqMuTildeMulChiPow.lean}89, 96--104, 226 & eTau is an equivalence onto Psi; v and both u endpoints are shared by all words; D.s is exactly SData with hcW \\
\nolinkurl{Formalism/SharpRestriction/CuspSpanNormBounds/DataS.lean}20--32 & v\_i−v\_(i+1)=4Delta and u\_(i+1)=v\_i−1/2−epsilon\_i Delta for epsilon\_i in \{0,1\} \\
\nolinkurl{Formalism/LocalRingCardinality/DataParams₂.lean}141--164 & S=4kDelta; at every non-top label, base u\_0−v\_0\textgreater=S \\
\nolinkurl{Formalism/SharpRestriction/ElevenLe.lean}105--114, 135--140 & SizesAt includes SizeDel, whose first clause is Delta\textgreater=N\_m m\^{}10 \\
\nolinkurl{Formalism/SharpRestriction/RamificationFrobenius/TotientGaloisClasses/ZsmulGeneratorEqZeroOfProjection.lean}90--122 & GG=(Z\^{}times x Z\^{}times) x MA, Psi is the second-sign kernel, tau1(a)=((1,1),iota(a)), tau2(a)=((-1,1),iota(a)) \\
\nolinkurl{Formalism/LocalRingCardinality/DataHost.lean}129 & topPlaces is the image of TT under tau1 \\
\nolinkurl{Formalism/LocalRingCardinality/RankBoundMonotoneInDimension.lean}524 & the supplied theorem factPlaces proves the two-type description and disjointness; its content is also immediate from the displayed coordinates \\
\end{longtable}

Only the first SizeDel inequality is needed from the long conjunction SizesAt, beyond the fields already contained in SizeData₂. The real theta step parameter P.t is not used or confused with Ranks.t in this argument. The N\_m occurring in SizeDel is its literal Ranks definition.

Retain an honest unitary Hecke character eta with the actual source IsEtaT properties. In particular

\[
 \eta|_{\mathbb A_F^\times}=\omega_{E/F},\qquad
 \eta_\infty(z)=\prod_{\tau\in\Psi}z_\tau/|z_\tau|.
\tag{SG3}
\]

These are \nolinkurl{AbelianHostFieldConditions.lean}308--312 and its Chars.Tied field at330. The imported \nolinkurl{UnitaryNorms/IsUnitaryNormOf.lean}90--94 distinguishes HasInf, which specifies angular exponent one here, from HasInfQ, which doubles its argument. Thus eta contributes exponent 1/2 in the convention \nolinkurl{(z/bar z)^x}. The genuine theta record separately records this as DataTheta.lean613--614. Existence of this character is a premise of the genuine source realization; no new character is constructed in this finding.

\subsection{2. One actual non-top label and every word's gaps}\label{one-actual-non-top-label-and-every-words-gaps}

The additive group A of the source Frame has a zero element. By the displayed coordinates,

\[
                  g_*=\tau_2(0)=((-1,1),1)\in\Psi.
\tag{SG4}
\]

It cannot equal any tau1(a), since their first signs are 1 and -1 as units in Z. In particular it is outside topPlaces, regardless of the chosen subset TT. Let

\[
                       \tau_*=(D.eTau)^{-1}(g_*).
\tag{SG5}
\]

This is one of the actual source labels, not a newly added complex embedding. It is common to every word. By SizeData₂.cloud\_gap and the shared first endpoint,

\[
                       u_0^{w,\tau_*}-v_0^{\tau_*}\ge4k\Delta
                       \quad\text{for every }w.
\tag{SG6}
\]

Write epsilon\_i for the bit of the actual word. SData.u\_letter and D7.hcW\_v give, for 0\textless=i\textless=k−3,

\[
              u_{i+1}^{w,\tau_*}
                  =v_i^{\tau_*}-\tfrac12-\epsilon_i\Delta.
\tag{SG7}
\]

Consequently the first adjacent gap satisfies

\[
 u_0-u_1=(u_0-v_0)+\tfrac12+\epsilon_0\Delta
               \ge4k\Delta+\tfrac12.
\tag{SG8}
\]

For 1\textless=i\textless=k−3 the two prescribed letters give

\[
 u_i-u_{i+1}
       =4\Delta+(\epsilon_i-\epsilon_{i-1})\Delta
       \ge3\Delta.
\tag{SG9}
\]

For the last gap, strict HC interlacing and the literal integer/half-integer grids give

\[
                  v_{k-2}-u_{k-1}\ge\tfrac12.
\]

It follows that

\[
 \begin{aligned}
 u_{k-2}-u_{k-1}
  &=4\Delta-\tfrac12-\epsilon_{k-3}\Delta
                         +(v_{k-2}-u_{k-1})\\
  &\ge(4-\epsilon_{k-3})\Delta\ge3\Delta.
 \end{aligned}
\tag{SG10}
\]

All indices are present: k is even and k\textgreater=m+3\textgreater=5, hence k\textgreater=6. The source SizeDel clause gives

\[
                  \Delta\ge N_m m^{10}>1.
\tag{SG11}
\]

Therefore every adjacent gap at tau\_* is greater than one. Sums of adjacent gaps give

\[
       |u_i^{w,\tau_*}-u_j^{w,\tau_*}|>1
                       \quad(i\ne j),\quad\text{for every }w.
\tag{SG12}
\]

Negation and reversal preserve this conclusion. No assertion about the two endpoint gaps at a top label is needed. At other labels, strict interlacing and the half-integer grid still give distinct coordinates and gaps at least one; this weaker fact will suffice for the blockwise algebraic coefficient argument below. The proof does not impose any additional restriction on the words or replace them by a selected subcode.

\subsection{3. Generic global parameter, with the actual infinity comparison}\label{generic-global-parameter-with-the-actual-infinity-comparison}

Apply the general wide-gap lemma DO046(AP6) to pi at the compact real place tau\_*. Its hypotheses are now supplied by SG1 and SG12. To keep the transfer issue explicit, we recall the relevant argument.

Chen--Zou, Theorem 2.1 and Proposition 4.1, attach a global elliptic Arthur parameter

\[
                       \Psi=\boxplus_i\Pi_i\boxtimes S_{d_i}
\tag{SG13}
\]

to the actual unitary automorphic constituent. Weak transfer alone does not identify its infinitesimal character at tau\_*. Choose the large quasi-split auxiliary theta target used in their section 3, with honest global splitting characters. Their Theorem 3.1 and Proposition 4.1 identify the global parameter of the same abstract theta lift as

\[
              \Theta(\Psi)=\Psi\chi_W^{-1}\chi_V
                                  \boxplus\chi_V\boxtimes S_{M-k}.
\tag{SG14}
\]

At tau\_* the actual compact pi is tempered. Their Theorem 6.3 and equation (6.1) put its same local theta lift in the quasi-split packet of

\[
              \Theta(\phi_{\pi,\tau_*})
                 =\phi_{\pi,\tau_*}\chi_W^{-1}\chi_V
                                      \oplus\chi_V\boxtimes S_{M-k}.
\tag{SG15}
\]

The real packet has the infinitesimal character of its associated parameter. This is the fixed-character construction in Mok, Proposition 8.2.1; Chen--Zou section 3.4 retains it for the possibly unbounded global localization by normalized induction. Thus SG14 and SG15 have identical infinitesimal-character multisets. Cancelling the identical appended string with multiplicities and the identical scalar shift gives IC(Psi\_tau\_\emph{)=IC(pi\_tau\_}). This is the already checked AP6 comparison; the auxiliary lift is not an added edge or a new finite test in the arithmetic family.

If any d\_i\textgreater=2, its tensor with S\_(d\_i) contributes a consecutive exponent string with step one. No negative multiplicities can cancel those entries in a direct-sum parameter. SG12 excludes this. Hence every d\_i=1:

\[
                            \Psi=\boxplus_i\Pi_i.
\tag{SG16}
\]

Only after this genericity has been proved do we use Chen--Zou Theorem 2.6: the actual pi belongs to its localized Langlands packet at every place. This order avoids assuming a ramified all-place transfer from near equivalence alone. Neither the G finite generic character nor the T central character has been used to prove SG16.

\subsection{4. The odd-eta correction and all finite places}\label{the-odd-eta-correction-and-all-finite-places}

Let Pi\_i be a cuspidal block of rank d.~At each complex embedding its bounded infinity characters are a subset of those of the actual compact pi, since SG16 and the all-place packet identity have now been established. Write their decreasing exponents x\_0\textgreater\ldots\textgreater x\_(d−1). The initial source grid is

\[
                      x_j\in\mathbb Z+(k-1)/2
                            =\mathbb Z+\tfrac12.
\tag{SG17}
\]

They are distinct and have adjacent differences at least one, including at labels where SG12 was not used. Let epsilon in \{0,1\} satisfy epsilon congruent to d−k modulo 2. Since k is even, epsilon is the parity of d.~The actual global character in SG3 is conjugate-selfdual: for any E-idele z,

\[
 \eta(z)\eta(\bar z)=\eta(N_{E/F}z)
                         =\omega_{E/F}(N_{E/F}z)=1.
\tag{SG18}
\]

The global cuspidal representation \nolinkurl{Pi_i tilde = Pi_i tensor eta^epsilon} is therefore still unitary and conjugate-selfdual. Its infinity exponents at a selected embedding are x\_j+epsilon/2, which lie in Z+(d−1)/2. The coefficient coordinates

\[
                  \lambda_j=x_j+\epsilon/2
                                  -\bigl((d-1)/2-j\bigr)
\tag{SG19}
\]

are integral and dominant, since lambda\_j−lambda\_(j+1)=x\_j−x\_(j+1)−1\textgreater=0. At the conjugate embedding the reversed negative list and the negative eta exponent give the corresponding coefficient weight. Thus the twisted cuspidal block is regular C-algebraic. The parameter convention is the stable one with no extra Mok auxiliary-character shift, as in DO046(AP8--10) and LIFT069(CT15--16).

This parity correction is essential for an odd-dimensional block in the even-dimensional initial group. Merely using the ambient k-dimensional coefficient grid would not establish C-algebraicity for that block. The twist is a proof device on a cuspidal GL block; it does not replace the source automorphic representation, its characters, its level or its T vector.

Shin--Templier, Lemma 2.12 and Remark 2.13, imply that each twisted cuspidal regular C-algebraic block is cohomological. We apply that implication to a global cuspidal representation, whose infinity parameters here are already bounded. Caraiani, Theorem 1.2, then gives temperedness at every finite place for d\textgreater=2. For d=1, unitarity of the Hecke character gives it directly. Untwisting a unitary character preserves temperedness. Thus all blocks Pi\_i are tempered at every finite place. The localized packet identity following SG16 proves

\[
 \boxed{\text{Every existing initial constituent for every word is
                  tempered at every place.}}
\tag{SG20}
\]

This includes the nonsplit finite places. It does not assert that the large theta descendants, which have inserted strings, are tempered, and it does not assert that the initial global parameter is simple.

\subsection{5. The literal Mram vector at a split T place}\label{the-literal-mram-vector-at-a-split-t-place}

At a selected T place the actual group is GL\_k(F\_v). The strict level is the stabilizer H\_u of the nonzero residue covector u; its projective stabilizer J\_u has quotient J\_u/H\_u equal to the residue multiplicative group. These are the same source groups and the same unweighted diamond sum as in K052(LT2--7), LIFT027(MR3--6) and the literal \nolinkurl{ScaledRamifiedLattice.lean}84--91,155--169 and \nolinkurl{Record.lean}1811--1848.

Let sigma=pi\_v, with its actual orientation or the contragredient specified in LIFT068. SG20 makes sigma tempered. Jacquet--Shalika, section 2.1, printed p.110, explicitly supplies a unique nonzero Whittaker functional for every irreducible pre-unitary tempered representation of GL\_k(F\_v); section 1.2, printed p.109, explains the mirabolic restriction underlying this assertion. Thus sigma is generic. This directly read primary statement avoids any converse from finite generic K\^{}1-invariants.

For h in J\_u, choose a unit c with u bar(h)=bar(c)u. Then c\^{}(-1)h lies in H\_u. On the entire nonzero fixed space,

\[
                       \sigma(h)|_{\sigma^{H_u}}
                                   =\omega_\sigma(c)\,1.
\tag{SG21}
\]

Since (1+varpi O)I lies in H\_u, the central character is tame on units. The literal unweighted norm is scalar multiplication by

\[
 \sum_{s\in\mathbb F_q^\times}\omega_\sigma(s)
  =\begin{cases}q-1,&\omega_\sigma|_{O^\times}=1,\\
                0,&\omega_\sigma|_{O^\times}\ne1.
    \end{cases}
\tag{SG22}
\]

A nonzero Mram component therefore has the nontrivial tame central character, and c(omega\_sigma)=1. Conjugating u integrally to the last covector identifies H\_u with the usual K\_1(varpi). The generic newform theorem gives c(sigma)\textless=1. A spherical representation has unramified central character, so c(sigma) cannot be zero. The newvector theorem also gives the dimension at the least level. Hence

\[
             c(\sigma)=c(\omega_\sigma)=1,
                   \qquad \dim\sigma^{H_u}=1.
\tag{SG23}
\]

The original fixed vector is thus a scalar multiple of the original newform in each irreducible constituent; it is not replaced by a freely chosen Whittaker vector. This is the one-way conductor implication required in LIFT068. We do not claim that every conductor-one representation has nontrivial central character, and we do not claim that the whole global multiplicity space is one-dimensional.

If the adelic row convention uses the inverse dilation action, SG22 is unchanged: inversion permutes the unit residue group and replaces omega\_sigma by its inverse. The vanishing condition and both conductor conclusions remain the same. This is precisely why no orientation choice is discarded in identifying the literal unweighted norm.

\subsection{6. Consequence for the corrected fixed T test}\label{consequence-for-the-corrected-fixed-t-test}

Let pi\_B be any constituent in the chosen nonzero LIFT059 bottom subspace. LIFT069 proves that its actual split T component Pi is spherical, generic and tempered. Put rho=m/2−1, so that its rank is n=k+2rho+1=k+m−1. For any existing initial Mram constituent, SG20--23 now supply the remaining smaller-representation hypotheses in LIFT068.

With the precise dual and additive-character conventions of that finding, and with Whittaker functions normalized to have value one at the identity, its unchanged actual D*tau test therefore satisfies

\[
 \left|\frac{Z_\rho(1/2)}{L(1/2,\Pi\times\sigma)}\right|
      =q^{-\rho(k+\rho+1)/2}\frac{q^{k/2}}{q^k-1}>0.
\tag{SG24}
\]

The physical character is the one of LIFT067, with both its unequal root coefficients retained. Its N-Haar conjugation modulus is delta\_N(D)=q\footnote{-rho(k+rho+1)}. Explicitly, multiplying the squared ratio in SG24 by delta\_N(D)\^{}(-1) gives q\^{}k/(q\^{}k−1)\^{}2. The Gauss phases, actual diagonal D, original lower translate tau and source root signs remain exactly those of LIFT068; this theorem changes none of them.

Conversion of a fixed Petersson/Hilbert vector to W(1)=1 retains its actual nonzero scalar. Thus SG24 is not a newly normalized global GGP factor or a uniform Hilbert lower bound. It is a nonzero local period on the fixed physical T test, now with both representation hypotheses supplied in the stated fully sized genuine family.

\begin{longtable}[]{@{}
  >{\raggedright\arraybackslash}p{(\columnwidth - 4\tabcolsep) * \real{0.3333}}
  >{\raggedright\arraybackslash}p{(\columnwidth - 4\tabcolsep) * \real{0.3333}}
  >{\raggedright\arraybackslash}p{(\columnwidth - 4\tabcolsep) * \real{0.3333}}@{}}
\toprule\noalign{}
\begin{minipage}[b]{\linewidth}\raggedright
Interface
\end{minipage} & \begin{minipage}[b]{\linewidth}\raggedright
Established here or carried input
\end{minipage} & \begin{minipage}[b]{\linewidth}\raggedright
Remaining scope
\end{minipage} \\
\midrule\noalign{}
\endhead
\bottomrule\noalign{}
\endlastfoot
Initial HC gap at one common non-top label & SG4--12 from literal full sizes, for every word & No existence of a fully sized Host is inferred \\
Every existing initial constituent tempered & SG13--20 & The global parameter may be endoscopic \\
Exact Mram T central type and newform & SG21--23 & No genericity-to-Mram inference and no conductor converse \\
Chosen bottom component tempered spherical & LIFT069 & Not an assertion about every summand of the full e0 \\
Actual corrected D*tau local period & LIFT068 and SG24 & Physical Haar/Whittaker conversion retained \\
Other finite and real local tests & Their individually checked results only & No new positive-rank G test follows here \\
Global prescribed Bessel/Mram period & Not supplied & A nonzero local test at each place alone does not evaluate the global central value \\
Read, compatible STiesR/RowsJoinPinned, 032, DP44 budget & Not supplied & All actual source and changed-family comparison obligations remain \\
\end{longtable}

\subsection{7. Literature and verification scope}\label{literature-and-verification-scope}

The argument SG13--20 applies the \textbf{general} lemma and parity mechanism already proved in DO046(AP6, AP8--10), with the new literal source hypothesis furnished by SG4--12. It does not extend DO046's concrete m=2 numerical profile by assertion. The primary statements used in that mechanism were directly checked again at their following locators:

\begin{itemize}
\tightlist
\item
  Chen, R., \& Zou, J. (2025). \emph{Arthur's multiplicity formula for even orthogonal and unitary groups}. Journal of the European Mathematical Society, 27, 4769--4843. \url{https://doi.org/10.4171/JEMS/1564} . Theorems 2.1, 2.6 and 3.1; section 3.4; Proposition 4.1 and its proof; equation (6.1) and Theorem 6.3. Primary PDF: \url{https://ems.press/content/serial-article-files/51452} .
\item
  Mok, C. P. (2015). \emph{Endoscopic classification of representations of quasi-split unitary groups}. Memoirs of the American Mathematical Society, 235(1108). \url{https://doi.org/10.1090/memo/1108} . The directly read author version \url{https://arxiv.org/abs/1206.0882} supplies section 7.1, printed p.190, and Proposition 8.2.1 with its proof, printed pp.231--232. The latter keeps all standard characters in the required infinitesimal-character category.
\item
  Shin, S. W., \& Templier, N. (2014 author version). \emph{On fields of rationality for automorphic representations}. Lemma 2.12 and Remark 2.13, printed pp.7--8. \url{https://math.berkeley.edu/~swshin/FieldRat.pdf} . This is the directly read cuspidal regular-C-algebraic to cohomological implication; no new direct read of the original Clozel volume is claimed.
\item
  Caraiani, A. (2012). \emph{Local-global compatibility and the action of monodromy on nearby cycles}. Duke Mathematical Journal, 161(12), 2311--2413. \url{https://doi.org/10.1215/00127094-1723706} . Theorem 1.2, printed p.2312. Primary PDF: \url{https://www.ma.imperial.ac.uk/~acaraian/papers/lgc1.pdf} .
\item
  Jacquet, H., \& Shalika, J. (1983). \emph{The Whittaker models of induced representations}. Pacific Journal of Mathematics, 109(1), 107--120. Sections 1.2 and 2.1, printed pp.109--110. Directly read author-hosted primary PDF: \url{https://www.math.columbia.edu/~hj/The\%20Whittaker\%20model\%20of\%20induced\%20representations.pdf} . Section 2.1 states the tempered-to-generic implication used above, with existence and uniqueness of the Whittaker functional. Web retrieval \nolinkurl{turn481762987view0} includes that exact passage.
\end{itemize}

The generic newform/conductor statement is the precise accepted K052(LT4--7) argument, with Jacquet--Piatetski-Shapiro--Shalika(1981), Theorem 5.1, and Jacquet(2012), Theorem 1. Its hypotheses, previously supplied for the concrete m=2 family by DO046, are now supplied by SG20 for the genuine fully sized family.

The companion exact script checks the actual index ranges, all four adjacent letter configurations, full small-word examples and dual orientations, endpoint/grid estimates, the size inequality, both parities of the cuspidal block, common-string cancellation and tame-character sums. These are finite diagnostics of the written proof. They do not construct an automorphic parameter, certify the cited classification theorems, prove existence of the literal source family, or calculate a global period. No Lean build is run.

\clearpage

\section{AM-DO-080 --- an actual selected-trace bound from radial fidelity}\label{am-do-080-an-actual-selected-trace-bound-from-radial-fidelity}

\textbf{Canonical checked version --- 2026-11-22T18:31:06.355982+00:00.} Complete root and independent mathematical/source/script reviews and separate3,010-check replays passed. Original pending wording below is historical. Promoted scope: For the actual all-top J56 selected normalized trace, every allowed word and twist, and arbitrary compact factors, the exact density/fidelity argument gives \textbar c\_W(g)\textbar\textless=min(1,4(a-1)sqrt(VarZ/(EZ-1))/sinh r), with squared constant\textless2\^{}-931. The genuine near-point global-rank2-unipotent consequence is \textbar c(h)\textbar\textless2\^{}(-929f/2) Norm\_F(alpha)\^{}(-1/2). This is not an operator bound or a sector removal by itself; all finite weights and global gates are unchanged pending a separate complete count.

Author: ALP\_MACHINE \nolinkurl{/root/definite_oscillator}. 22 November 2026 UTC. Controlled recovery-pending finding; final author and independent checks are separate.

\subsection{1. Statement and actual scope}\label{statement-and-actual-scope}

Use the actual all-top J56 family at the K044 upper gap endpoint, every allowed interior word and canonical twist, and the selected isometry and Hilbert metric of K043. At one real label, write \nolinkurl{a=168}, let pi be the prescribed holomorphic discrete series of U(a,1), and let W be its actual finite-dimensional selected compact type. For g in the real group put

\[
 c_W(g)=\frac1{d_W}\operatorname{Tr}_W(P_W\pi(g)|_W),
 \qquad d_W=\dim W.
 \tag{RT1}
\]

This is the factor used in K062(AT14), at the actual argument \texttt{g\_p\^{}\{-1\}\ h\ g\_p}. The compact factors of that argument are retained. If its Cartan radius is r\textgreater0, the following theorem holds:

\[
 \boxed{\quad
 |c_W(g)|\le\min\left\{1,
 \frac{4(a-1)}{\sinh r}
       \sqrt{\frac{\operatorname{Var}(Z)}{{\bf E}Z-1}}\right\}
 <\frac{2^{-931/2}}{\sinh r},
 \qquad Z=(1-|z|^2)^{-1}.
 \quad}
 \tag{RT2}
\]

The distribution used for Z is the actual selected probability measure defined below, not Euclidean polynomial mass or a scalar replacement for the fibre. In particular the simpler bound \texttt{2\^{}\{-465\}/sinh\ r} is valid. At r=0 the normalized trace is bounded by one.

For the genuine product trace and the actual near-point transport, let h be globally rank-two unipotent, write U=h-I, and put alpha=(U\^{}2)\_(nn). When the hypotheses of K085's near-point calculation hold, alpha is totally positive and

\[
 \boxed{\quad
 |c(h)|<2^{-929f/2}\operatorname{N}_{F/\mathbb Q}(\alpha)^{-1/2},
 \qquad f=168.
 \quad}
 \tag{RT3}
\]

One may also retain the stronger local minimum with one in RT2 before multiplying. The theorem does not remove an arithmetic sector by itself. In particular it makes no support-floor, signed-sum, period, Mram, literal constructor or DP44 assertion. Its improvement is for the genuine normalized trace; the DO079 operator obstruction is unaffected.

\subsection{2. The true selected density and the trace inequality}\label{the-true-selected-density-and-the-trace-inequality}

On the ball B\^{}a let Y=1-\textbar z\textbar\^{}2 and use exactly K046(BL5),

\[
 d\mu(z)=\pi^{-a}Y^{-a-1}\,dV(z).
 \tag{RT4}
\]

The actual holomorphic representation is a space of sections of the homogeneous Hermitian bundle attached to its unitary compact fibre. Let \nolinkurl{||.||_z} denote that invariant fibre metric. Nakahama(2016), equation (5.7) and Theorem 5.1, in the K043 dictionary, give

\[
 \langle f_1,f_2\rangle_{\mathcal H}
   =c_\tau\int_{B^a}\langle f_1(z),f_2(z)\rangle_z\,d\mu(z),
 \qquad c_\tau=\prod_i(r_i-1).
 \tag{RT5}
\]

All PBW/Pieri factors are included in this Hilbert norm. Let f\_1,\ldots,f\_(d\_W) be an orthonormal basis of the actual W and define

\[
 \rho(z)=\frac{c_\tau}{d_W}\sum_j\|f_j(z)\|_z^2.
 \tag{RT6}
\]

Then rho is nonnegative, has integral one against mu, and is K-invariant: changing z by a compact element acts unitarily both on the fibre and on the finite-dimensional W. Thus rho depends only on Y. The group action on sections transports the invariant fibre metric isometrically. Applying Cauchy--Schwarz to the finite sum over j at each point, and then the triangle inequality to RT5, gives

\[
 |c_W(g)|\le
 \int_{B^a}\sqrt{\rho(z)\rho(g^{-1}z)}\,d\mu(z).
 \tag{RT7}
\]

This follows from the finite trace on W; it is not a trace of pi on its infinite-dimensional Hilbert space. There is no operator-norm substitution, no omitted compact intertwiner, and no approximation of the selected metric. Because rho is K-invariant, the right side of RT7 depends only on the Cartan radius even when g has both compact factors. This invariance concerns the upper bound, not an identification of c\_W(g) with a signed radial coefficient.

With normalized sphere measure d omega, polar coordinates give

\[
 d\mu=w(y)\,dy\,d\omega,
 \qquad w(y)=\frac{(1-y)^{a-1}y^{-a-1}}{(a-1)!},\quad0<y<1.
 \tag{RT8}
\]

Put p(y)=rho(y)w(y). This is the probability density of Y for the selected measure. It has the exact K046 moments in section 6 below. In particular it is an ordinary density, and the negative moments used here are finite.

\subsection{3. A Cartesian bound for the ball radial transition kernel}\label{a-cartesian-bound-for-the-ball-radial-transition-kernel}

Fix r\textgreater0 and write c=cosh r, s=sinh r. Use the Cartan transformation in the first positive coordinate and the negative coordinate. At a point with Y=y set R=sqrt(1-y), and write its first complex coordinate z\_1=x+it. The marginal of the normalized sphere measure is

\[
 f_y(x,t)=\frac{a-1}{\pi R^2}
       \left(1-\frac{x^2+t^2}{R^2}\right)^{a-2}
       {\bf1}_{x^2+t^2<R^2}.
 \tag{RT9}
\]

This is the usual first-coordinate marginal on the complex unit sphere, rescaled by R; integration over the other a-1 complex coordinates gives the displayed density. It applies because a\textgreater=2.

The transformed radial coordinate is

\[
 y'=\frac{y}{(c+sx)^2+s^2t^2}.
 \tag{RT10}
\]

For each fixed t, the map from the admissible x interval to y' is injective and strictly decreasing. In fact c+sx\textgreater=c-sR\textgreater0, and

\[
 \left|\frac{dx}{dy'}\right|
 =\frac{y}{2s(c+sx)y'^2}
 \le\frac{\sqrt y}{2s y'^2}.
 \tag{RT11}
\]

The last bound uses the exact identity

\[
 (c-sR)^2-y=(s-cR)^2\ge0,
 \qquad c-sR>0.
\]

Let K\_r(y,y') be the conditional density of y' when the point at radius y is uniformly distributed in its sphere. The density RT9 is at most \nolinkurl{(a-1)/(pi R^2)}, while t lies in an interval of length at most 2R. Integrating the change of variable RT11 therefore proves, including by zero extension outside the feasible region,

\[
 K_r(y,y')\le
 \frac{a-1}{\pi\sinh r}
       \frac{\sqrt y}{\sqrt{1-y}\,y'^2}.
 \tag{RT12}
\]

No angular-arc estimate or asymptotic expansion is used. All Jacobians are real Jacobians, and the two real coordinates of z\_1 have both been counted.

\subsection{4. Detailed balance removes the hyperbolic radial volume ratio}\label{detailed-balance-removes-the-hyperbolic-radial-volume-ratio}

The invariant ball measure and K-invariance imply

\[
 w(y)K_r(y,y')=w(y')K_r(y',y)
 \quad\text{for almost every }(y,y').
 \tag{RT13}
\]

To verify the symmetry precisely, integrate \nolinkurl{F(Y(z))G(Y(a_r z))} against mu for nonnegative compactly supported functions F,G. Change variables by a\_r. Its inverse is K-conjugate to a\_r, by the compact matrix changing the sign of the first positive coordinate, and Y is invariant under that compact matrix. The resulting integral is the same expression with F and G exchanged. Disintegration using RT8 proves RT13. This proof also works with an infinite invariant measure, because the test integrals are first taken nonnegative or compactly supported.

Substituting rho=p/w in RT7 and using RT13 gives

\[
 |c_W(g)|\le\iint\sqrt{p(y)p(y')}
          \sqrt{K_r(y,y')K_r(y',y)}\,dy\,dy'.
 \tag{RT14}
\]

Applying RT12 in both directions produces the factorized estimate

\[
 \sqrt{K_r(y,y')K_r(y',y)}
 \le \frac{a-1}{\pi\sinh r}
       (yy')^{-3/4}\bigl((1-y)(1-y')\bigr)^{-1/4}.
 \tag{RT15}
\]

In particular the large ratio w(y)/w(y') has canceled exactly. Bounding that ratio separately by a constant on a radial interval would give a much poorer estimate and is not part of this proof.

\subsection{5. The one-dimensional variance bound}\label{the-one-dimensional-variance-bound}

Define

\[
 Z=Y^{-1},\qquad S=2\sqrt{Z-1},\qquad
 s_0=2\sqrt{{\bf E}Z-1},\qquad
 \sigma^2=\frac{4\operatorname{Var}(Z)}{{\bf E}Z-1}.
 \tag{RT16}
\]

Here \nolinkurl{E Z>1} and the variance is positive in the actual selected measure. Let q\_S be the density of S on the positive half-line. The exact change of variable is

\[
 \left|\frac{dS}{dy}\right|
      =y^{-3/2}(1-y)^{-1/2},\qquad
 \int_0^1\sqrt{p(y)}y^{-3/4}(1-y)^{-1/4}\,dy
      =\int_0^\infty\sqrt{q_S(s)}\,ds.
 \tag{RT17}
\]

Moreover,

\[
 {\bf E}(S-s_0)^2
 =4{\bf E}\frac{(Z-{\bf E}Z)^2}
 {\bigl(\sqrt{Z-1}+\sqrt{{\bf E}Z-1}\bigr)^2}
 \le\sigma^2.
 \tag{RT18}
\]

Cauchy--Schwarz with the positive weight \nolinkurl{1+(s-s_0)^2/sigma^2} now yields

\[
 \begin{split}
 \left(\int_0^\infty\sqrt{q_S(s)}\,ds\right)^2
 &\le\left(1+\frac{{\bf E}(S-s_0)^2}{\sigma^2}\right)
       \int_0^\infty\frac{ds}{1+(s-s_0)^2/\sigma^2}\\
 &\le2\pi\sigma.
 \end{split}
 \tag{RT19}
\]

The last integral is bounded by the integral over the entire real line. This proves finiteness of the integrals used in RT14--17 as well. Combining RT14--19 gives the first bound in RT2, with exactly the constant \nolinkurl{4(a-1)} shown there. All steps are uniform in the actual compact factors of g.

\subsection{6. Exact moments for the actual selected type}\label{exact-moments-for-the-actual-selected-type}

K046(BL8--10), derived from Nakahama's actual shifted norm, applies to each unit vector in W and hence also to the average RT6. It gives

\[
 M_1={\bf E}Z=\prod_i\frac{r_i+m_i-1}{r_i-2},\qquad
 M_2={\bf E}Z^2=
       \prod_i\frac{(r_i+m_i-2)(r_i+m_i-1)}{(r_i-3)(r_i-2)}.
 \tag{RT20}
\]

All r\_i\textgreater3 here. Put

\[
 \varepsilon_i=\frac{m_i+1}{(r_i-3)(r_i+m_i-1)},\qquad
 S_\varepsilon=\sum_i\varepsilon_i.
\]

A direct cancellation, with no approximation, gives

\[
 \frac{M_2}{M_1^2}=\prod_i(1+\varepsilon_i),\qquad
 \operatorname{Var}(Z)
   =M_1^2\left(\prod_i(1+\varepsilon_i)-1\right).
 \tag{RT21}
\]

For the actual J56 all-top endpoint write \nolinkurl{A=2^960}, \nolinkurl{B=2^955}, \nolinkurl{ell=ell_+}, \nolinkurl{E=B+ell+96}, and \nolinkurl{0<=q<=ell}. The literal K043/K046 word formulas are

\[
 \begin{aligned}
 r_i&=E+(4a-3-4i)\ell+q+1 &&(0\le i<a-1),\\
 r_{a-1}&=2\ell+q+1,\\
 m_0&=A+83-668\ell,\\
 m_i&=4\ell-1-\ell\omega_i &&(1\le i<a-1),\\
 m_{a-1}&=12.
 \end{aligned}
 \tag{RT22}
\]

Thus r\_i is independent of the interior word. This is the corrected actual convention recorded in DO079-C1. A one-bit decreases m\_i; increasing q increases r\_i. Each factor in M1 and each epsilon\_i increases with m\_i and decreases with r\_i. Hence their simultaneous upper bounds use the zero word and q=0.

At that upper point the product for M1 telescopes to

\[
 M_1^{\max}
  =\frac{(33B+2\ell+179)(2\ell+12)}
          {(B+6\ell+95)(2\ell-1)}<34.
\]

The first factor alone, evaluated at q=ell, is

\[
 M_1>\frac{33B+3\ell+179}{B+671\ell+95}>32.
\]

The other factors exceed one. Both comparisons are strict exact integer inequalities for the retained ell. For the upper comparison the numerator of \nolinkurl{34-M1max} is \nolinkurl{2B(ell-215)+404ell^2+5874ell-5378}, which is positive. The lower comparison reduces to \nolinkurl{B>21469ell+2861}.

The epsilon bound also has a short exact proof. At the same zero-word, q=0 maximum,

\[
 \varepsilon_0<\frac{32}{33B},\qquad
 \sum_{i=1}^{a-2}\varepsilon_i<\frac1{66B},\qquad
 \varepsilon_{a-1}<\frac1{66B}.
\]

For the first inequality use \nolinkurl{m0+1<A=32B}, \nolinkurl{r0-3>B}, and \nolinkurl{r0+m0-1>A+B}. For the second use \nolinkurl{r_i-3>B}, \nolinkurl{r_i+m_i-1>B}, and \nolinkurl{664ell/B^2<1/(66B)}. For the last use \nolinkurl{epsilon_(a-1)=13/[(2ell-2)(2ell+12)]} and \nolinkurl{(2ell-2)(2ell+12)>858B}. These inequalities are included in the exact certificate. Therefore, uniformly over all actual words and canonical twists,

\[
 32<M_1<34,\qquad 0<S_\varepsilon<B^{-1}.
 \tag{RT23}
\]

Since \nolinkurl{S_epsilon<1}, the elementary product estimate \nolinkurl{product(1+epsilon_i) <= 1/(1-S_epsilon)} gives

\[
 \frac{\operatorname{Var}(Z)}{M_1-1}
 <\frac{34^2}{31(B-1)}.
\]

For example, the product estimate follows from \nolinkurl{1+epsilon_i <= (1-epsilon_i)^(-1)} and \nolinkurl{product(1-epsilon_i) >= 1-sum epsilon_i}. Applying it to RT2 yields the exact squared coefficient bound

\[
 \left(4(a-1)\sqrt{\frac{\operatorname{Var}(Z)}{M_1-1}}\right)^2
 <\frac{515834944}{31(2^{955}-1)}<2^{-931}.
 \tag{RT24}
\]

The final comparison is an integer comparison: \nolinkurl{31*2^24-515834944=4258752>0}, with ample room for the subtracted 31 after clearing \nolinkurl{2^931}. This proves the numerical statement in RT2 without a floating-point variance computation.

\subsection{7. The actual near-point unipotent consequence}\label{the-actual-near-point-unipotent-consequence}

This section uses the same genuine argument of K085 section 2; it does not enlarge its finite support or assume global unipotence from a finite reduction. Suppose the actual rational h is globally rank-two unipotent in signature (a,1). Then U=h-I satisfies \nolinkurl{U^3=0}, and \nolinkurl{U^2} is a nonzero rank-one self-adjoint operator. With \nolinkurl{J=diag(I_a,-1)}, over the complex embedding only, write

\[
 U^2=-\alpha vv^*J,\qquad v=(v_+,1),\quad\|v_+\|=1,
 \quad\alpha=(U^2)_{nn}>0.
\]

No such normalized vector is asserted to be integral at finite places. If the actual transport g\_p has Cartan radius r\_p\textless1/2, then

\[
 \alpha_p=(g_p^{-1}U^2g_p)_{nn}
       =\alpha|v^*Jg_pe_n|^2
       \ge e^{-2r_p}\alpha>\alpha/3.
\]

Here the positive and negative parts of \nolinkurl{g_p e_n} have norms \nolinkurl{sinh r_p} and \nolinkurl{cosh r_p}, so their pairing with v has modulus at least \nolinkurl{exp(-r_p)}. For \texttt{h\textquotesingle{}=g\_p\^{}\{-1\}hg\_p} of Cartan radius r',

\[
 2\operatorname{Re}(h'_{nn})-2=\alpha_p,
 \qquad |h'_{nn}|=\cosh r',\qquad \sinh^2r'\ge\alpha_p.
\]

RT24 thus gives \texttt{\textbar{}c\_tau(h\textquotesingle{})\textbar{}\^{}2\ \textless{}\ 3*2\^{}\{-931\}/alpha\_tau\ \textless{}\ 2\^{}\{-929\}/alpha\_tau}. Multiplying at the f labels gives RT3. The finite character factor, incidence weights, class masses, all actual projectors and all existing arithmetic errors are unchanged. An application to any complete short-sum sector requires its separate counting proof.

\subsection{8. Source and verification boundary}\label{source-and-verification-boundary}

Primary source: Nakahama, R. (2016). \emph{Norm computation and analytic continuation of vector valued holomorphic discrete series representations}. Journal of Lie Theory, 26(4), 927--990. Equation (5.7) and Theorem 5.1, printed pages 960--961; author version \url{https://arxiv.org/abs/1506.05919}. The retained original \nolinkurl{research/nakahama2016.txt} at those passages was directly read. The actual model and normalization dictionary are K043(PB1--14); K046(BL5--14) supplies the exact selected radial moments, not a scalar-fibre analogy. The genuine coefficient being bounded is exactly K062(AT14). The actual near-point transport and unipotent calculation are K085 section 2.

The ball transition calculation, detailed balance, variance inequality and their combination with the actual moments are proved in full above. The companion certificate checks the literal all-top parameters, exact moment and product identities, every numerical comparison in RT22--24, the elementary coarea and change-of-variable algebra, and the product exponent in RT3. It does not numerically prove the Hilbert-space or measure-disintegration statements; those are the written proof. There is no Lean build or automorphic period evaluation.

\clearpage

\section{AM-R-058 --- complete fixed-block rank-two unipotent removal at supports thirteen through twenty-three}\label{am-r-058-complete-fixed-block-rank-two-unipotent-removal-at-supports-thirteen-through-twenty-three}

Author: GPT-6 Astra Pro, inherited effort not externally exposed; ALP\_MACHINE \nolinkurl{/root}. 22 November 2026 UTC. New proof; author whole-file check and independent review are recorded separately.

\subsection{1. Actual sector and conclusion}\label{actual-sector-and-conclusion-2}

Retain the unchanged normalized J56 all-top upper-endpoint kernel of K062/R048, its cutoff H=133, its actual selected Hilbert modules, favorable-point transport, finite projectors, source and target pencils, phases and class masses. Write f=a=168 and E\^{}169=S perpendicular W0, with dim S=164 and signature(W0)=(4,1). Consider the actual terms

\[
 h=I_S\oplus h_0=I+N,\qquad N^3=0,\quad
 \operatorname{rank}_E N=2,\quad\operatorname{rank}_E N^2=1,
 \qquad13\le s\le23.
 \tag{HB1}
\]

Here s is the actual nonintegral T support. Global rank-two unipotence and pointwise fixation of S are explicit hypotheses. All incidence profiles i+u+t=s are included, where I,II,III have counts i,u,t.

\textbf{Theorem.} The total absolute contribution of HB1 is less than2\^{}-39409. At each support the portion with total unmixed finite-G rank D\_G\textless=32 is empty. The complementary portion has a complete bound with exponent at most −161437607/4096\textless−39413. The intersection of this entire union with the current remainder can be removed with one charge39409, retaining every earlier charge.

This removes additional actual higher-support sectors. It does not claim a new support floor, a complete support13 exhaustion, a signed-energy lower bound, an actual period or a DP44 saving.

\subsection{2. Inputs whose support dependence is retained}\label{inputs-whose-support-dependence-is-retained}

Use DO078(FP4--12), with its original two-lattice argument and original-coordinate counts. Put

\[
 R=17424/133,\quad Q=g_0g_1,\quad
 \alpha=(N^2)_{nn}\in F,\qquad0<\alpha_\tau\le R.
 \tag{HB2}
\]

At the334 unmixed relative G factors let d\_v=rank(N mod v) in\{0,1,2\} and

\[
 D_G=\sum_vd_v,\qquad
 \mathfrak J_F=\prod_v\mathfrak p_v^{2-d_v},\qquad
 J=N_F\mathfrak J_F=Q^{334}/\prod_vq_v^{d_v}.
 \tag{HB3}
\]

The two mixed G projectors retain their modulus bounds one. The square-pole orders are zero at I, two at II, one at III, so, with mu=2u+t,

\[
 \alpha\in\mathfrak J_F I_{\rm II,F}^{2}I_{\rm III,F},
 \quad N_F(\alpha)\ge JT^{-\mu},\qquad
 y=(JT^{-\mu})^{1/f}\le R.
 \tag{HB4}
\]

Each of these is local or algebraic and independent of the numerical value s=12 in DO078's application. In particular no all-II or common-stable-lattice assumption is introduced when s varies.

The exact Green estimate is

\[
 |a(h)|\le2Q^{-335(f-1)}J^{335/2}.
 \tag{HB5}
\]

DO080(RT3), applied to the actual normalized selected trace at the original favorable point, gives the conservative consequence

\[
 |c(h)|<2^{-464f}N_F(\alpha)^{-1/2}
          \le2^{-464f}y^{-f/2}.
 \tag{HB6}
\]

Indeed DO080 proves the stronger exponent−929f/2; weakening it to−464f is harmless. This is a trace bound in the true metric with both compact factors retained, not the operator estimate excluded by DO079.

The original affine negative-column count and fixed-block stabilizer bound from DO078(FP8--10) are

\[
 \#\{W_0\text{ columns at fixed }\beta\}<2^{5379}T^{8s+r_0},
 \qquad G_5<2^{107},
 \tag{HB7}
\]

and the number of possible real diagonal values is at most

\[
 2^{1062}\left(\frac{375+R/(2y)}{376}\right)^f.
 \tag{HB8}
\]

Here r0 counts zero target labels. These formulas keep the real coarea, dyadic factor, actual field covolume and arbitrary-ideal covering bound. Their only support dependence is exactly as displayed. The proof of the stabilizer estimate fixes S and the negative vector, giving a faithful positive rank-four module; it is independent of s and does not borrow the scalar-cross rank-five argument's nonzero-S-column premise.

At each actual place, combining the complete ordered-pair weight with the extra zero-target factor in HB7 costs at most2T,1,2 for I,II,III respectively. This is precisely DO078(FP11--12); its calculation includes the admitted zero label and uses no support bound. Thus the profile's remaining factor is2\^{}(i+t)T\^{}i. The numbers of support sets, type assignments and G patterns are bounded by2\^{}168,3\^{}s\textless2\^{}(2s), and3\^{}334\textless2\^{}530 respectively. No later union over profiles is needed in addition to those type assignments.

\subsection{3. Comparing all profiles at their unchanged finite-G ideal}\label{comparing-all-profiles-at-their-unchanged-finite-g-ideal}

For the same actual J define

\[
 y_0=(JT^{-2s})^{1/f},\qquad z_0=\sqrt{y_0},\qquad
 y/y_0=T^{(2i+t)/f}.
 \tag{HB9}
\]

The gamma factor HB8 decreases as y increases. The real factor HB6 changes by T\^{}(-(2i+t)/2). Hence the complete extra profile factor is at most

\[
 2^{i+t}T^{i-(2i+t)/2}
       =2^{i+t}T^{-t/2}\le2^s.
 \tag{HB10}
\]

This compares sufficient bounds at the same J. It does not posit an all-II matrix in the support. Both the numerator ideal and every actual source/target incidence are unchanged.

Combining HB5--10 shows that, for any common upper cap on the possible z0 in a rank sector, its full absolute contribution is bounded by the corresponding maximum of

\[
 \begin{split}
 \mathcal B_s(z)={}&
 2^{1+5379+1062+107+168+3s+530-464f}
 Q^{-335(f-1)}T^{343s}\,
 z^{334f}\left(\frac{375+R/(2z^2)}{376}\right)^f.
 \end{split}
 \tag{HB11}
\]

The exponents are335s from substituting J=z\^{}(2f)T\^{}(2s) in the Green factor, plus8s from the column fibre. The z power is335f−f=334f. The pattern count is retained once. The variable part is

\[
 z^{332f}(375z^2+R/2)^f/376^f,
 \tag{HB12}
\]

which is increasing for z\textgreater0. This proves that taking a common cap bounds every actual pattern before summing its at most2\^{}530 possibilities.

For0\textless z\textless=1, the exact rational comparison

\[
 \left((375+R/2)/376\right)^4<2
 \tag{HB13}
\]

gives the simpler bound

\[
 z^{334f}\left(\frac{375+R/(2z^2)}{376}\right)^f
       <2^{42}z^{55776}.
 \tag{HB14}
\]

Here332f=55776 and f/4=42. There is no loss from discarding a small-alpha singularity before the count; its full power is already in HB11--14.

\subsection{4. High finite rank at every support13 through23}\label{high-finite-rank-at-every-support13-through23}

The actual primes are g0=22285932761\textless g1=64072175441 and Q=1427908193729121522601. Retain the actual T from DO040's certificate and the checked one-sided bounds

\[
 \log_2T<L_T=1019333/1024,\qquad
 \log_2Q>L_Q=287843/4096.
 \tag{HB15}
\]

If D\_G\textgreater=33 then J\textless=Q\^{}334/g0\^{}33. Set

\[
 j_s=1116-95(23-s),\qquad13\le s\le23.
 \tag{HB16}
\]

The eleven strict integer comparisons in the companion certificate are

\[
 Q^{334}2^{21j_s}<g_0^{33}T^{2s}.
 \tag{HB17}
\]

Since2f=336 and336/16=21, they give0\textless z0\textless2\^{}(-j\_s/16)\textless1. Substituting this cap in HB14, and the one-sided logarithms in the positive T and negative Q powers, gives the complete rational exponent

\[
 \begin{split}
 E_s={}&1+5379+1062+107+168+3s+530-464f\\
 &-335\cdot167 L_Q+343sL_T+42-55776j_s/16.
 \end{split}
 \tag{HB18}
\]

It increases with s: its exact increment is10516211/1024\textgreater0. At the final support,

\[
 E_{23}=-161437607/4096<-39413.
 \tag{HB19}
\]

Thus every high-rank sector under consideration has absolute contribution below2\^{}-39413, including all support/type/coordinate/fibre factors already present in HB11. This is an exact sufficient bound, not a numerical enumeration of rational isometries.

\subsection{5. Low finite rank is empty throughout the same range}\label{low-finite-rank-is-empty-throughout-the-same-range}

Suppose D\_G\textless=32. Retain the original packet vector w=Qu, the three original integral kernel vectors y\_j and their norms from DO078(FP23--24), and U=77Q²−14. Each y\_j lies in every admitted source and target pencil lattice at both CM partners, so the separate selfdual pairings in DO078(FP20--22) still give

\[
 b_j=\langle y_j,N^2y_j\rangle_J\in O_F\cap\mathfrak J_F.
 \tag{HB20}
\]

This argument is independent of the number of places: it uses membership in each original source and target lattice, not a common lattice preserved by h. At every embedding \textbar b\_j\textbar\textless=R\textbar\textbar y\_j\textbar\textbar². Since

\[
 J\ge Q^{334}/g_1^{32}\ge Q^{334}/g_1^{69},
 \tag{HB21}
\]

the already checked K088/DO078 norm contradictions at budget69 force all three b\_j=0. Their unchanged integer kernel then forces the image line of N² into span\_E(w,e\_n). The exact original entry identity(N²)11=−4alpha/U and the retained coprimalities imply, exactly as in DO078(FP26--27),

\[
 \alpha\in U\mathfrak J_F I_{\rm II,F}^{2}I_{\rm III,F}.
 \tag{HB22}
\]

The largest pole exponent in this range is mu\textless=2s\textless=46. The new strict integer comparison is

\[
 (133U)^{168}Q^{334}
      >17424^{168}g_1^{32}T^{46}.
 \tag{HB23}
\]

Consequently HB2, HB21--23 would imply

\[
 R^{168}\ge N_F(\alpha)
 \ge U^{168}Q^{334}/(g_1^{32}T^{2s})
 \ge U^{168}Q^{334}/(g_1^{32}T^{46})>R^{168},
 \tag{HB24}
\]

a contradiction. This proves emptiness of the entire low-rank sector, for every incidence profile and every support in HB1.

\subsection{6. Complete union, scope and geometric interpretation}\label{complete-union-scope-and-geometric-interpretation}

The two integral alternatives D\_G\textless=32 and D\_G\textgreater=33 exhaust HB1. There are eleven support values. By HB19 the full union is therefore bounded by

\[
 \sum_{s=13}^{23}2^{E_s}<11\,2^{-39413}<2^{-39409}.
 \tag{HB25}
\]

There is no extra count for incidence profiles:3\^{}s in HB11 already counts their assignments to a fixed support. Intersect with the current R057 remainder before removing the union. All earlier charges, including the complete support12 exhaustion, remain unchanged. The global all-top support floor remains13 until its other structural classes are removed; the uniform H8192 remainder is untouched.

The new arithmetic input combines exact ideals on the original Hermitian modules over Spec(O\_E), source/target integral pairings and their complex selected representation. The determinant/rank functions retain the SGA3 I4.7.3/4.7.3.1 module framework and SGA6 VI Theorem3.3 exterior-power context. The unchanged polarized CM variety, theta line and biextension have the carried locators FGA232 Theorem3.1 pp148--149; FGA236 Corollaries1.7/2.3 pp229/232; SGA7 I VII2.7/2.9 pp169--174 and VIII7.4(b) pp310--311. Formal-crystal and Cartier-duality context remains cote4 scan003 sections8.1.1--8.1.2, scans005--006 sections8.4.1--8.4.3, and cote54 scans005--007.

No new lambda-ring identity, K-theory comparison, l-adic sheaf or weight theorem, motivic projector, CM-motive identification or L-function conclusion is inferred. The finite gate toward DP\_MULT Q3e and THETA3 G1/A7 is strictly smaller in the stated sector. Actual source construction, the prescribed Bessel/Mram period, the concrete032 arithmetic pin/value, every actual DP44 term, CX\_GRH CX-ALG-024 and the global arithmetic comparison of GCT\_PLETH T-052 remain open as separately recorded.

The proof cites settled K062, R048, DO078(FP4--28), K085 and K088 in their stated scopes, and the new whole-checked DO080(RT1--24) when canonically available. Its new certificate checks the complete rational envelope, all profile power comparisons, all eleven exact caps, the new pole-budget46 norm contradiction and the complete error-array addition. Those computations supplement the full argument; they do not automate the Hilbert, incidence or geometric statements. No Lean build is performed.

\clearpage

\section{AM-K090 --- actual Mram augmentation generators and the selected arithmetic matrix}\label{am-k090-actual-mram-augmentation-generators-and-the-selected-arithmetic-matrix}

\textbf{Canonical checked version --- 2026-11-22T18:42:37.504987+00:00.} Root and independent peer read the whole frozen proof and complete script, directly checked the literal interfaces, and separately replayed29265 diagnostics. Original author status below is preserved as history. Promoted scope: exact source augmentation basis and norm, alternating selected matrix, actual Sad lattice equality at11, and saturated auxiliary image before all selected derivatives, projectors and class contractions.

Author: ALP\_MACHINE \nolinkurl{/root/kummer_extension}, 22 November 2026 UTC.

\textbf{Author whole-proof and complete-script PASS; independent review pending.} This is a theorem about the unchanged K038 source lattice and its auxiliary section input. It is not a claim that the fixed selected derivative or subsequent class contraction is primitive or nonzero. The final verification record will distinguish author and independent checks.

\subsection{1. Fixed source and comparison cover}\label{fixed-source-and-comparison-cover}

Keep either actual J14 or J56 host, an allowed word, tuple, real profile and canonical twist, the strict initial covector level at every T place, the actual generic-unipotent level and character at G, and the maximal level at29. Let K be the carried coefficient field, R its complete valuation ring at a place above11, and f=42 or168. The local coefficient lattice is the literal initial Weyl lattice delAtr. No dual-Weyl replacement is made.

At a selected T place v, let H\_v fix the nonzero residue covector and let P\_v fix its line. Put \[
 D_v=P_v/H_v=k_v^\times,\qquad
 D_T=\prod_{v\mid T}D_v,\qquad |D_v|=T-1.
 \tag{MA1}
\] Enlarging H\_v to P\_v is used only to index the existing strict class set. It is not a new level for a theta input. Retain the original generic G level throughout this indexing.

The actual projective-T/generic-G rational stabilizers are trivial. Indeed K042(NB1--6) proves neatness using only the unipotent reduction at all G places in the faithful rational representation of dimension2fa. Its inequalities g0\textgreater2fa+1 are unchanged when the T condition is enlarged. Every definite rational stabilizer is finite, hence is trivial. The free-cover implication was already checked for DO044 in \nolinkurl{agents/review_mram_dimension_fraction.md}.

Consequently the strict initial classes above each projective-T/generic-G class c form a free D\_T-torsor. Choose a representative d\_c and write \[
                  d_{c,s}=\operatorname{dia}(s)d_c,
                       \qquad s\in D_T.
 \tag{MA2}
\] These are distinct strict classes. The coefficient fibre can be identified without a rational moving frame along this orbit. Literal \nolinkurl{StdLatUnimodular.lean}113--122 defines diamond by \[
            N_{d_{c,s}}=N_{d_c},\quad
            f_{d_{c,s},v}=s_v f_{d_c,v},\quad
            \mathrm{be}_{d_{c,s}}=\mathrm{be}_{d_c}.
 \tag{MA3}
\] The kernels of all f\_v, their lattice L, and its dual Lsh therefore stay the same. \nolinkurl{RecLevelAsm.lean}61--68 and \nolinkurl{Record.lean}1317--1327 explicitly prove the corresponding lattice equality. Diamond has identity coefficient map in \nolinkurl{ScaledRamifiedLattice.lean}79--91. Its action on coefficient functions may invert s according to the pullback convention; inversion only relabels MA2.

Let L\_c,R be this common transported coefficient lattice. The exact class description is thus \[
 \mathcal M_{f,R}
    =\bigoplus_c R[D_T]\otimes_R L_{c,R},\qquad
 \|P\|_{0,\mathrm{cls}}^2=\sum_{c,s}\|P_{c,s}\|^2.
 \tag{MA4}
\] Here R{[}D\_T{]} is identified with functions by its delta basis. There is no initial stabilizer mass, by K042. This is not R017's projective-T/full-G model with stabilizer mu28. No assertion about middle or terminal stabilizers follows.

\subsection{2. Integral augmentation basis and exact norm}\label{integral-augmentation-basis-and-exact-norm}

For a finite group D and any coefficient ring R, let \[
 I(D)=\ker(\varepsilon:R[D]\longrightarrow R),\qquad
 \varepsilon\left(\sum_s a_s\delta_s\right)=\sum_s a_s.
\] The vectors delta\_s−delta\_1, s!=1, form an R-basis of I(D): the nonidentity coordinates recover every basis coefficient, and the identity coefficient is their negative sum. This is also a direct summand of R{[}D{]}, with complement R delta\_1. Neither statement divides by \textbar D\textbar.

The source norm at v is the unweighted sum of all diamonds in that factor. On its regular module it sends a vector to its augmentation times the all-ones vector. The map R to that all-ones line is injective over any ring. Applying the direct-sum decomposition just described in each independent factor gives \[
 \boxed{\quad
 \mathcal M_{\mathrm{ram},R}
   =\bigoplus_c\left(\bigotimes_{v\mid T} I(D_v)\right)
                            \otimes_R L_{c,R}.
 \quad}
 \tag{MA5}
\] This is exactly the intersection lattice defined by ColRamV/MramL in ScaledRamifiedLattice155--175 and transported in Record1811--1851. It is not a chosen nontrivial-character space and remains valid if11 divides T−1.

For an R-basis p\_\{c,j\} of L\_c,R, a basis of MA5 is \[
 B_{c,\mathbf s,j}
    =\bigotimes_{v\mid T}(\delta_{s_v}-\delta_1)\otimes p_{c,j},
                       \qquad s_v\ne1\text{ for every }v.
 \tag{MA6}
\] Its2\^{}f nonzero class values are distinct and are all plus or minus p\_\{c,j\}. Thus, at each actual complex realization where the coefficient is defined, \[
                    \|B_{c,\mathbf s,j}\|_{0,\mathrm{cls}}
                         =2^{f/2}\|p_{c,j}\|.
 \tag{MA7}
\] For two multi-indices in the same orbit, their coefficient-free Gram entry is the product over v of (1+1\_\{s\_v=t\_v\}). In particular this difference basis is not asserted orthogonal, and MA7 is not an estimate for a coordinate inverse.

For global integral generators, use K059(IS4--6, IS18)'s actual Weyl basis p\_j and clearing integer d\_c.~Since Lsh is unchanged along MA2, the same d\_c works at every orbit point. The vectors in MA6 with coefficient d\_c p\_j are integral, span Mram over K, and obey \[
       \|B_{c,\mathbf s,j}(d_c p_j)\|_{0,\mathrm{cls}}
                  \le 2^{f/2}d_c B_{0,j}.
 \tag{MA8}
\] In the standard11 frames of K039/K059,11 does not divide d\_c and the initial Weyl basis is a local basis; these vectors are then a basis of the localized lattice. If a different coefficient family is used, its exact local span and clearing factors must be retained rather than assuming a basis.

This replaces the sufficient multiplier C\_T=(T−1)\^{}f in the optional source spanning construction K059(IS19--21) by2\^{}\{f/2\}, with every coefficient-frame cost retained. The binary norm cost is exactly21 for J14 and84 for J56. The source automorphic norm still has the factor sqrt(v0/d0) from DO056(AC3), or the corresponding explicit factor for a nonunit compact test. This result changes no terminal generator/projector denominator and provides no bound on the size or conditioning of a coordinate inverse. It requires no multiplicative-character field extension.

\subsection{3. Exact alternating selected matrix}\label{exact-alternating-selected-matrix}

Let A be the unchanged K038(CA27) coupled operator. For a terminal decoration d and an initial strict class e define its actual column operator \[
 \mathcal K_{d,e}
   =\sum_{f_1\in\mathcal C_1}\frac1{|\Gamma_{1,f_1}|}
       \sum_{Y\in\mathcal Y(p_{1,f_1})}
             k_2(p_{1,f_1},Y,d)k_1(p_{0,e},p_{1,f_1},Y).
 \tag{MA9}
\] Every k\_i is the literal selected subD map CA23, on the actual source rhoVr/Weyl model. The coefficient of ftY already contains its q\^{}\{-Mb\}; there is no division by \textbar Y\textbar. The middle coevaluation is the complete K037 one. Initial stabilizers are one, but all middle denominators remain.

For epsilon in \{0,1\}\^{}f put s\^{}epsilon\_v=s\_v when epsilon\_v=1 and1 otherwise. Then \[
 \boxed{\quad
 (A B_{c,\mathbf s,j})(d)
   =\sum_{\epsilon\in\{0,1\}^f}
          (-1)^{f-|\epsilon|}\mathcal K_{d,d_{c,\mathbf s^\epsilon}}
                                     p_{c,j}.
 \quad}
 \tag{MA10}
\] This is direct substitution of MA6 into CA27. It is an alternating sum of the actual selected columns, not an independent jet interpolation formula. The raw output remains the evGSD row of CA28; terminal brC and pairBK are subsequently applied exactly as in CA30. Their brConst/kapBr, conjugate piT power, transported intNab detector and actual terminal masses are unchanged.

More generally, let ell\_1,\ldots,ell\_h be any fixed finite family of actual scalar tests on the output, including their prescribed branch, projector and denominator-clearance choices. The fractional R-ideal of all values ell\_l(A P), P in Mram\_R, equals \[
 \sum_{l,c,\mathbf s,j}R\cdot
          \ell_l\left(A B_{c,\mathbf s,j}\right).
 \tag{MA11}
\] This is because MA6 is an R-basis; it also holds for the zero image. If the ell\_l are a basis of the actual target dual lattice, MA11 is the image-content ideal in that lattice. No irreducible or genuine finite projector is assumed11-integral: its actual denominator and possible reduction kernel stay inside ell\_l or the chosen output map.

DO056(AC9--13) independently establishes for the literal Sad S and its reversed correspondence S\^{}\# that \[
 S^{-1}=T^{-fa}S^\#\text{ on Mram},\qquad
 S(\mathcal M_{\mathrm{ram},R})=\mathcal M_{\mathrm{ram},R}.
 \tag{MA12}
\] Here T is an11-unit, while the actual additive-character values are integral units. Therefore the ideals in MA11 are unchanged if A is replaced by A S. This local lattice equality does not erase S from the actual analytic seesaw: its complex norm remains T\^{}\{fa/2\}, and the moved initial trace remains the one in DO055/DO056. For a proper prescribed submodule of Mram, retain that submodule and its actual S-image; MA12 does not replace it by the full source.

\subsection{4. The actual restricted T orbit before selected differentiation}\label{the-actual-restricted-t-orbit-before-selected-differentiation}

Fix one coarse class c in MA2, one actual middle class f\_1 and one actual terminal decoration d.~Their underlying lattices and non-T frames are held fixed. Because MA3 also fixes Lsh, the Hom tori, theta line and source comparison scalars are the same along this finite initial orbit. We use their intrinsic common position model; if a common finite Weil coordinate change is made, its inverse-integrality is exactly K089(PS11). No rational change is called integral in an old frame.

Let Psi\_s be the prescribed complete two-edge finite tensor for (d\_\{c,s\},f\_1,d), before any subD contraction. Its local T dependence in these fixed integral bases is \[
 \mathbf1_{X\text{ integral},\ u_1X=s_vu_0}\,
 \mathbf1_{Z\text{ integral},\ u_2Z=u_1}.
 \tag{MA13}
\] Both target covectors are nonzero. This is the literal phiT in FiniteSumNonzeroLargeDepth648--650, together with XiCpl1/XiCpl2 at1134--1166. The definition uses the source and target lattices, not any Fourier-replaced T input. Diamond scales exactly the first source covector. At each place, X maps to u\_1X surjectively on residue rows: choose a residue vector v with u\_1v=1, using the inverse of a nonzero coordinate of u\_1; then X=v(s\_vu\_0) is a residue solution. The same argument constructs Z. Integral lifts exist. Distinct s\_v give disjoint first-edge supports since u\_0 is nonzero.

Every allowed Y factor is at G or Q, by K038(CA22). Thus all the functions Psi\_s have identical non-T factors and identical full Y contractions. Use the actual common away-T unit witness of K089(PS8--10): the unique surviving G Y class gives value1, the complete Q sum gives its nonzero11-unit Gauss value, and the remaining vacuum factors have unit values. No summand is substituted for the complete Q tensor. Independently choose the T solutions just described. This gives finite position points x\_s and one common u in R\^{}times with \[
                         \Psi_s(x_t)=u\,\delta_{s,t}.
 \tag{MA14}
\] All choices are support evaluations of the existing functions. They add no T factor to Y, no new level, and no independent allowed Heisenberg input. If one begins in the unramified splitting extension used below, the same identity holds over that extension and descends for the intrinsic functions.

Consequently the finite map \[
 \iota:R[D_T]\longrightarrow\mathscr S_R,\qquad
                         \delta_s\longmapsto\Psi_s
 \tag{MA15}
\] has an integral left inverse: evaluate at all x\_t and divide by u. It is a split, hence saturated, injection. The tensor augmentation module is a direct summand of R{[}D\_T{]}, so its restriction is a split injection too. This is an assertion about the actual restricted scalar-decoration orbit, not about the whole finite Schwartz module.

\subsection{5. Exact auxiliary section transport and its limit}\label{exact-auxiliary-section-transport-and-its-limit}

Choose one common even D, prime to11, covering the supports and periods of the finitely many Psi\_s. This is K089's comparison cover, not an automorphic refinement. Its period-removed reading is the already fixed map \[
 \Theta_D:\mathscr S_D\longrightarrow H_D
               =H^0(X,[D]^*\mathcal L).
 \tag{MA16}
\] K089(PS2--4) proves that over its proof-only finite unramified theta-splitting extension R'/R, this actual map is an integral isomorphism with integral inverse. Its scalar is fixed by the actual trace-one vacuum and primitive principal section. It is not an arbitrarily rescaled representation intertwiner. The Heisenberg normalizing orders are11-units; finite general-linear projector orders are not being substituted for them.

Compose MA15 with MA16. Over R' the evaluation retraction gives a split injection of the actual augmentation module into H\_D. Tensoring with L\_c,R gives the same assertion for the uncontracted coefficient-labelled auxiliary module \[
 \left(\bigotimes_v I(D_v)\right)\otimes_R L_{c,R}
                 \longrightarrow H_D\otimes_R L_{c,R}.
 \tag{MA17}
\] The external coefficient tensor here precedes the selected polynomial/subD operation. It is not a claim that a Pieri embedding or a selected derivative preserves primitivity.

The actual sections and this restricted map are defined over R. The cokernel after base change is R'-free; faithful flatness shows that the cokernel over R is torsion-free. It is a finite module over a DVR, hence free. Thus the restricted map over R is split and saturated as well. The arbitrary full theta basis need not descend. No extension of the actual coefficient field is made.

In particular each MA6 basis vector with primitive local coefficient has primitive image at this auxiliary-section stage. More generally the image-content of any coefficient combination at this stage equals its source content. This conclusion is independent of whether11 divides \textbar D\_T\textbar, and it holds for every fixed actual class triple used in the construction.

In the integral algebraic formal frame the section coefficients have the same content, by K089(PS12). The reduced Chern/CM germ still contains \[
                      U(t)^{\kappa}\exp(-\kappa\beta(t))f(t)
 \tag{MA18}
\] with the original period and source-creation factors external to this finite comparison. Its ordinary derivatives are not identified with integral Hasse coefficients. The subsequent selected subD functional can annihilate MA17, and the sum over middle classes in MA9 can cancel. The finite projectors can also change content. MA10--11, with those operations retained, is the precise remaining matrix to evaluate.

\subsection{6. Scope, source checks and certificate}\label{scope-source-checks-and-certificate}

The conclusions are: an exact integral Mram basis on the existing free scalar-decoration cover; the norm/spanning improvement MA8; an exact alternating formula for the actual selected matrix and its content ideal; invariance of that ideal under the actual Sad lattice automorphism at11; and a saturated auxiliary-section image of this restricted Mram source before subD and class contractions. None supplies a nonzero selected column, a unit selected minor, a nonzero final period, a genuine full Host/SizesAt family, or a DP44 sign.

The directly read literal source interfaces are StdLatUnimodular113--122; RecLevelAsm49--68; ScaledRamifiedLattice79--91 and155--175; Record1317--1327 and1811--1851; FiniteSumNonzeroLargeDepth648--650 and1125--1166. They are under \nolinkurl{work/research/alpoge_source/Formalism/LocalToGlobal/Columns}. Accepted proofs read for the new argument are K038(CA22--30), K042(NB1--7), K059(IS4--6,IS18--21), DO056(AC1--13), the free-cover check \nolinkurl{agents/review_mram_dimension_fraction.md}, and K089(PS1--15). The external geometric theta-group input is exactly the one already proved and source-checked in K089; no fresh theorem about generic mod11 theta correspondence is invoked.

The companion certificate checks complete finite augmentation models over integers and residue characteristic11, the explicit left inverse and Gram values, the actual incidence support separation, full finite alternating-column identities, and the punctured Fourier/Sad algebra in finite models. Its source fingerprints pin the files read. These diagnostics do not compute the selected enormous theta matrix or certify nonvanishing after the fixed contraction. No Lean build is run.

\subsection{Author verification}\label{author-verification-1}

At 2026-11-22T18:31:31.242652+00:00, I read the complete saved MA1--18 proof and all six substantive sections, then the complete companion script. I reconstructed the free-cover/stabilizer argument, the ring-independent tensor augmentation decomposition, the actual coefficient-lattice and norm transport, CA27's alternating columns with all middle/terminal masses, and the precise Sad lattice scope. I checked the common class-orbit Hom torus and full-Y unit witness against K089 and the literal phiT/diamond definitions. The split injection and faithful-flat descent are confined to the uncontracted auxiliary section module. The final status and residue-solution wording were reread. Final exact replay and file hashes are in the separate verification record; these diagnostics do not compute the selected arithmetic image.

\clearpage

\section{AM-DO-081 --- integral cross pairings bound every unipotent scalar-cross support 13--22}\label{am-do-081-integral-cross-pairings-bound-every-unipotent-scalar-cross-support-1322}

\textbf{Canonical checked version --- 2026-11-22T18:47:45.148522+00:00.} Root and independent peer read the entire frozen proof and complete script; separate5621-check replays passed. Promoted scope: every globally rank-two-unipotent scalar-cross profile at supports13..22, one25470 charge; no support-floor or final-value claim. Original author wording below is retained.

Author: \nolinkurl{/root/definite_oscillator}, 22 November 2026 UTC. New proof in the controlled pending area. All previously frozen findings and arithmetic data remain unchanged.

\subsection{1. Exact statement, actual data and conventions}\label{exact-statement-actual-data-and-conventions}

Retain the actual all-top J56 short kernel, the original packet and standard coordinates, the same CM representatives, class masses, finite projectors and selected Hilbert trace. In particular \[
 f=a=168,\quad n=169,\quad m=164,\quad H=133,\quad
 E^n=S\perp W_0,\quad\dim_E S=m,\quad\operatorname{sig}(W_0)=(4,1),
 \qquad R=H+H^{-1}-2=17424/133.
 \tag{IP1}
\] Consider exactly the matrices \[
 h=I+N,\quad h_{SS}=I,\quad\operatorname{rank}_E N=2,\quad
 N^3=0,\quad\operatorname{rank}_E N^2=1,\quad
 x=(Ne_n)_S\ne0.
 \tag{IP2}
\] Global unipotence is an explicit hypothesis. Let s count the active relative T places, and i,u,t their actual types I,II,III, respectively, so i+u+t=s.

\textbf{Theorem.} The total absolute contribution of all IP2 terms with every support, incidence profile, actual ordered source/target pair, G orientation, original-coordinate column and class fibre, for \textbf{13\textless=s\textless=22}, is strictly less than \textbf{2\^{}-25470}. Intersect this sector with the current remainder and retain every older error. This theorem alone changes neither the all-top support floor13 nor the uniform support floor11.

Only the bound \textbar c(h)\textbar\textless=1 for the actual normalized selected trace is needed. It follows from unitarity and compression to the selected finite-dimensional subspace, including the actual near-point conjugation and both compact factors. Thus the proof does not replace that trace by a special vector or require a new real-coefficient estimate.

The original packet constants, with no new four-square representation, are \[
\begin{split}
 g_0&=22285932761,\quad g_1=64072175441,\quad Q=g_0g_1,\quad U=77Q^2-14,\\
 u_0&=(2,95192390527,383092671043,12529843548438452388609),\quad
 w=Q u_0,\quad u_0^t u_0=U.
\end{split}
\tag{IP3}
\] Here u\_0 denotes the four-coordinate integer vector, distinct from the number u of type-II places. The actual T satisfies 2\^{}995\textless T\textless2\^{}996 and Q\^{}4\textgreater2\^{}281.

We use the Hermitian convention \textless v,z\textgreater=v*Jz, linear in the second argument. All ideal norms below are full absolute norms. Each relative T place contributes both CM primes to an E-ideal; no selected-half norm is used.

\subsection{2. Original-coordinate rank-two identities and reconstruction}\label{original-coordinate-rank-two-identities-and-reconstruction}

The signature-one identities already proved in DO072(GS4), DO073(RF4--8) and K086(TI2--3) give \[
 N^\dagger=N^2-N,\quad N^2=-\alpha\ell\ell^*J,\quad
 \ell_n=1,\quad\ell\in W_0,\quad\|\ell_+\|^2=1,\quad
 0<\alpha_\tau\le R,\quad \|x_\tau\|^2\le\alpha_\tau.
 \tag{IP4}
\] Here alpha=(N²)\_nn belongs to F and ell=N²e\_n/alpha belongs to E\^{}n; normalizing this rational isotropic line does not declare it integral. The equations hold at every complex CM embedding. In particular N²S=0 and N²e\_n=alpha ell.

Since im(N) intersects W0 in E ell and N maps S into that intersection, Ne\_j is a multiple of ell for every original S coordinate j. The paired adjoint identity implies N\_nj=bar(N\_jn)=bar(x\_j), fixing that multiple exactly: \[
                         Ne_j=\bar x_j\ell\qquad(j\in S).
 \tag{IP5}
\] Let y be an original positive integer vector supported in W0 and perpendicular to w. Put \[
 b_y=\langle y,\ell\rangle,\qquad
 \delta_y=(\langle y,Ne_j\rangle)_{j\in S}=b_y\bar x,
 \qquad t_y=\langle y,N^2e_n\rangle=\alpha b_y.
 \tag{IP6}
\] If b\_y is nonzero, then t\_y is nonzero and the ordered data (alpha,t\_y,delta\_y) recover x uniquely by \[
                         x=\overline{\alpha\delta_y/t_y}.
 \tag{IP7}
\] This is only a reconstruction over E. It will not be used to assert integrality of x after division. The distinct scalar t\_y must be retained; delta\_y alone need not determine x.

\subsection{3. Actual source/target modules make these pairings integral}\label{actual-sourcetarget-modules-make-these-pairings-integral}

At a selected T component the literal paired pencil is \[
 M_c=(I+c\pi^{-1}Z)\Lambda,\quad Z=zq,\quad Z^2=0,
 \quad z=v_0/\bar\pi,\quad q=-v_0^*J-\pi m_0,
 \quad v_0=(w,\pi_k).
 \tag{IP8}
\] The admitted labels are c in k\_T except -1, including zero. At the opposite CM component one uses the exact inverse-adjoint dual lattice. The shear changes only W0. This is LIFT44(PR4--9), K077(AP1--3), and K086(TI5--8), with the label-zero correction retained.

Every original S basis vector belongs to every M\_c at both components. Every integer y as in IP6 also belongs to every M\_c at both components: qy=-pi m\_0y makes the selected inverse shear integral on y; at the partner, y pairs integrally with Lambda+pi\^{}-1 O z because it is perpendicular to z. This is the direct membership argument of K088(IP5--6) and DO078(FP20--22).

For a fixed tuple of actual arrows h M\_c=M\_d, patch the local modules to the original full projective selfdual source and target lattices. Since e\_j belongs to the source and y to the target, integrality of the target Hermitian form gives \[
 \langle y,h e_j\rangle\in O_E.
\] The constant term \textless y,e\_j\textgreater{} is zero. Therefore \[
                              \delta_y\in O_E^m.
 \tag{IP9}
\] This proves full integrality even at active type-II places. It is an integral scalar pairing, not a claim that the W0 part of Ne\_j is standard integral.

For t\_y, both y and its images have the actual source/target relations, and \[
 N^2y=hy+h^{-1}y-2y\in M_d+M_c.
\] Both M\_c and M\_d lie in pi\^{}-1 Lambda at the selected component. Their paired duals do too: the inverse shear is I-c pi\^{}-1 Z, so its adjoint has only a simple pole. Thus N²y has at most a simple pole at either CM prime of each active T place. Self-adjointness of N², followed by pairing with the original integral e\_n, gives \[
 t_y\in I_\Sigma,\qquad
 I_\Sigma=\prod_{\mathfrak P\mid\Sigma}\mathfrak P^{-1}\subset E,
 \qquad N_E I_\Sigma=T^{-2s}.
 \tag{IP10}
\] At all inactive places the original matrix and vectors are integral. This argument holds for every actual arrow; it does not assume a common h-stable lattice, equate source and target labels, or discard their CM partner.

The scalar alpha has the safe entrywise square-pole bound \[
 \alpha\in I_{\Sigma,F}^{2},\qquad N_F I_{\Sigma,F}=T^{-s}.
 \tag{IP11}
\] Indeed N²=h+h\^{}-1-2I and the two actual shears give at most two entry poles at both CM partners. The stronger type-dependent pole table is not needed in IP11.

\subsection{4. One common G ideal for every pairing coordinate}\label{one-common-g-ideal-for-every-pairing-coordinate}

At each of the334 unmixed relative G places, the reduction of N has rank d=0,1 or2 and is nilpotent. Both CM components are integral and inverse-adjoint. At rank one write N\_P=v r\^{}t over the residue field. The SS block vanishes, hence either v\_S=0 or r\_S=0. These are common alternatives for all S coordinates.

If r\_S=0, every S column vanishes at P. If v\_S=0, the paired inverse-adjoint and N\_P²=0 show that every S column vanishes at bar(P). At rank zero the S columns vanish at both. The original integer y is integral at these primes, so the same common prime divides all coordinates \textless y,Ne\_j\textgreater. The identical statement holds for every fixed integral linear combination of the W0 rows of these S columns.

Choose a common orientation by a fixed rule at the rank-one places. There are at most4\^{}334=2\^{}668 patterns. For a pattern, let J be the product of its chosen CM primes at rank one and both CM primes at rank zero. Relative to DO066(GC5--9)'s S-last-column convention this ideal can be conjugated; its norm and Green bound are unchanged. The actual conclusions needed here are \[
 \delta_y\in J^m,\qquad
 |a(h)|\le2Q^{-335(f-1)}(N_EJ)^{335/2}.
 \tag{IP12}
\] The rank-one factor P1\^{}167 is bounded by P1\^{}(335/2), while rank zero is exactly P0\^{}335. The one Euler factor2 is global, not repeated at individual places. The two mixed projectors keep their modulus bound one. Crucially, IP12 is proved from the common row/column outer-product alternatives, not merely from vanishing of some x\_j modulo a prime.

\subsection{5. A uniform weighted packing bound}\label{a-uniform-weighted-packing-bound}

Use the half-trace Euclidean metric, the actual norm convention of DO066(GC14--17). For a fractional E-ideal L with N\_E L=z\^{}(2f), every nonzero element has squared trace norm at least f z². The product lattice L\^{}m has minimum distance at least sqrt(f)z. The already used disjoint-ball estimate gives, inside a ball of radius sqrt(f)M, \[
                              \#(L^m\cap B)\le(1+2M/z)^{2fm}.
 \tag{IP13}
\] If the counted actual vector is nonzero, necessarily z\textless=M. Therefore \[
 z^{335f}(1+2M/z)^{2fm}
 =z^{(335-2m)f}(z+2M)^{2fm}
 \le M^{335f}3^{2fm},
 \quad 335-2m=7>0.
 \tag{IP14}
\] Both factors before the inequality are increasing in z. This bounds the congruence loss simultaneously with the number of allowed vectors; no congruence density is divided into an unconstrained count. The entire m-dimensional product is counted, including zero coordinates, so there is no missing coordinate-support binomial or k-union.

We also use the same elementary packing in dimensions2f and f.~For any M\textgreater=1, e\textgreater=0 and an E-ideal of norm T\^{}(-2e), the number of elements whose absolute values are at most M at every label is at most \[
                            (3M)^{2f}T^{2e}.
 \tag{IP15}
\] For F-ideal I\_\Sigma,F² and the positive alpha bounded by R, the count is at most \[
                         (3R)^fT^{2s}<2^{9f}T^{2s},
 \tag{IP16}
\] because R\textless2\^{}8 and3R\textless2\^{}9. Counting the entire real ball only overcounts the totally positive elements. No sign choices are added.

\subsection{6. Case A: a transverse integral pairing is nonzero}\label{case-a-a-transverse-integral-pairing-is-nonzero}

Use the three original integer vectors from K088(IP14--16): \[
\begin{split}
 y_1&=(-976627,-153032,38026,0),\\
 y_2&=(-567678,695570,-172838,0),\\
 y_3&=(104583,-15309315176,-61610045224,2).
\end{split}
\tag{IP17}
\] Their squared norms are978665066829,835948910828,4030172803672366725045. They span w-perp in the original positive four-space; the exact integer transformation and determinant are checked again in the certificate.

In this case choose the first y\_j with b\_yj nonzero. IP4 and IP6 give, at each embedding, \[
 \|\delta_{y_j,\tau}\|\le\sqrt R\|y_j\|<2^{40},\qquad
 |t_{y_j,\tau}|\le R\|y_j\|<2^{44}.
 \tag{IP18}
\] The vector delta is nonzero because b and x are nonzero. Apply IP14 with L=J and M=2\^{}40. The weighted delta count, omitting the common Green factor, is at most \[
                      2^{40\cdot335f+4fm},
 \tag{IP19}
\] since3\^{}(2fm)\textless2\^{}(4fm). There are at most three choices of y\_j, bounded by2². IP10 and IP15 bound t\_y by2\^{}(92f)T\^{}(2s); IP16 bounds alpha by2\^{}(9f)T\^{}(2s). For each such triple, IP7 determines x uniquely. Retaining alpha and t\_y as separate data is a safe overcount; no integral inverse or unrecorded scalar is assumed.

\subsection{7. Case B: the square image lies in the original packet plane}\label{case-b-the-square-image-lies-in-the-original-packet-plane}

If all b\_yj vanish, their spanning property forces \[
                    \ell=(a w,1)=(b u_0,1),\quad b=aQ\in E,
                    \qquad b\bar b=1/U.
 \tag{IP20}
\] This is an equation in E, not an assertion that sqrt(U) belongs to E or that b is integral. Take the explicit integer Bezout row \[
                  v_B=((1-u_{0,2})/2,1,0,0),\qquad v_B u_0=1.
 \tag{IP21}
\] It is integral because u\_\{0,2\} is odd. Apply it to the four positive W0 rows of the actual S columns and of N²e\_n.~It gives precisely \[
              \delta'=b\bar x,\qquad B=\alpha b\ne0.
 \tag{IP22}
\] These are actual fixed integer linear combinations of original entries. Their bounds do not involve the size of the Bezout coefficients, because the exact identities IP22 give their archimedean values.

K086(TI8--9) proves both S cross blocks integral at types I and III; at type II their poles are at most one. Combining with the same column-oriented G pattern of section4 gives \[
 \delta'\in(J I_{\rm II})^m,\qquad N_E I_{\rm II}=T^{-2u}.
 \tag{IP23}
\] Every entry of N² has at most two poles at active T places and is integral elsewhere. Hence the fixed integer row IP21 gives \[
 B\in I_\Sigma^2,
 \qquad |B_\tau|\le R/\sqrt U<1,
 \qquad\|\delta'_\tau\|\le\sqrt{R/U}<2^{-69}.
 \tag{IP24}
\] The last strict inequalities are checked with the actual integers; U\textgreater2\^{}146 and R\textless2\^{}8 already suffice. No denominator is ascribed to b itself.

Set z\^{}(2f)=N\_EJ T\^{}(-2u). Since delta' is nonzero, z\textless=2\^{}-69. Applying IP14 with M=2\^{}-69 and retaining the exact Green conversion (N\_EJ)\^{}(335/2)=T\^{}(335u)z\^{}(335f) yields the weighted delta' count \[
                         2^{-69\cdot335f+4fm}T^{335u}.
 \tag{IP25}
\] The number of B in IP24 is bounded by3\^{}(2f)T\^{}(4s)\textless2\^{}(4f)T\^{}(4s), using the larger radius1. Retain the safe independent alpha count IP16. Then x=overline(alpha delta'/B) is uniquely reconstructed. In fact alpha is uniquely determined by the totally positive solution alpha²=U B bar(B), but no saving from this further constraint is needed. There are neither extra2\^{}f signs nor a presumed primitive scalar.

\subsection{8. Actual W0 columns, isometry fibres and all pair weights}\label{actual-w0-columns-isometry-fibres-and-all-pair-weights}

For fixed alpha, Re(h\_nn)=1+alpha/2. For each actual ordered source/target tuple the original affine W0 negative-column count of corrected DO072(GS17--21), DO070-C1 and DO078(FP8) is \[
                         2^{5379}T^{8s+r_0},
 \tag{IP26}
\] where r\_0 counts zero target labels. The constant5379 uses the original-ring-only covering bound and the correct paired coarea; it is not assigned to an arbitrary fractional J ideal. It is uniform in x, alpha, the local labels and every type. Additional relations involving t\_y or B can only reduce this count.

After x and these five coordinates are fixed, the complete negative column is fixed. DO073(RF1--11) applies to exactly IP2 and bounds the actual source-to-target isometry fibre by \[
 G_5=98148490017978252769652736000000<2^{107}.
 \tag{IP27}
\] This is the genuine positive rank-five intersection-lattice argument; no moving basis is declared integral. Its bound holds after imposing any of the extra equations counted above.

The complete ordered-pair weights, including the extra zero-target T in IP26, are at most2T,1,2 for types I,II,III, respectively, as proved in K086 section4 and DO078(FP11). Thus a fixed type assignment contributes at most2\^{}(i+t)T\^{}i\textless=2\^{}sT\^{}s beyond T\^{}(8s). This retains the admitted zero label and counts both CM partners only once. Absolute values are taken; no favorable signs are assumed.

For support s, retain all support choices \textless=2\^{}168, all3\^{}s\textless2\^{}(2s) type assignments, all4\^{}334=2\^{}668 common rank/orientation patterns, the one Green factor2, and G5. Set \[
                         C_0=1+107+668+168+5379=6323.
 \tag{IP28}
\] Combining sections6 and8, the complete Case-A envelope for fixed s has exponent \[
\begin{split}
 E_A(s)={}&C_0+2+3s+40\cdot335f+4fm+9f+92f\\
          &+13s\cdot996-(281/4)335\cdot167.
\end{split}
\tag{IP29}
\] The T power13s consists of2s for alpha,2s for t\_y,8s for the W0 affine fibre and at most s for the pairs. No second column support count is inserted.

For Case B, u\textless=s and sections7--8 give \[
\begin{split}
 E_B(s)={}&C_0+3s-69\cdot335f+4fm+9f+4f\\
          &+350s\cdot996-(281/4)335\cdot167.
\end{split}
\tag{IP30}
\] Here the T exponent before the last bound is335u+15s:2s for alpha,4s for B,8s for the W0 fibre, at most s for the pairs, and335u from the congruence-ideal weight. The estimate u\textless=s gives350s. This deliberate loss avoids claiming a profile improvement that has not been charged.

Both displayed affine functions increase with s. Exact arithmetic gives \[
 E_A(22)=-5042053/4<-25475,\qquad
 E_B(22)=-101901/4<-25475.
 \tag{IP31}
\] For each of the ten supports13..22 there are the two exhaustive cases. Their20 complete envelopes sum to less than20 times2\^{}-25475, which is strictly less than2\^{}-25470. The choice of one of the three y\_j was already included only in Case A. There is no additional endpoint, coordinate-support, alpha-sign or CM-partner multiplicity.

\subsection{9. Accounting, certificate scope and remaining gates}\label{accounting-certificate-scope-and-remaining-gates}

The proof preserves all of R058's53 all-top charges and31 uniform charges, and appends one new all-top charge25470. The companion reads \nolinkurl{checks/r058_joint_remainder_error.json} and checks the full new54-entry all-top sum and unchanged uniform sum remain below2\^{}-230. If another overlapping bound is later used, remove only the relevant set difference; do not infer positivity from any of these absolute estimates.

The certificate checks the literal actual primes/T, the fixed packet vector and three integer test vectors, their spanning transformation, the integral Bezout row, all strict archimedean comparisons, the GL5 gcd, every incidence-profile exponent through22, both complete affine bounds and their union, and the full error arrays. It does not automatically prove the paired-lattice, reconstruction, unitary or Green interfaces. Those are the mathematical arguments above, with their accepted exact source dependencies named at the point of use.

The new observation is the integral scalar pairing IP9 for arbitrary arrows, followed by a count of those pairings and the explicitly retained reconstruction scalars. The plane case uses the exact original packet vector rather than an anisotropy assertion. Neither a new level, finite test, field, selected real module, period, CM phase nor arithmetic class is chosen. No signed sufficient inequality, Mram nonvanishing, global prescribed period, literal Host constructor, DP44 or RH conclusion follows. Other structural sectors at support13 and above remain to be treated before any support-floor update.

\clearpage

\section{AM-R-059 --- integral cross-block pairings remove all scalar-cross supports thirteen through twenty-one}\label{am-r-059-integral-cross-block-pairings-remove-all-scalar-cross-supports-thirteen-through-twenty-one}

Author: GPT-6 Astra Pro, inherited effort not externally exposed; ALP\_MACHINE \nolinkurl{/root}. 22 November2026 UTC. Complete new proof; author and independent checks are recorded separately before promotion.

\subsection{1. Actual sector and theorem}\label{actual-sector-and-theorem-1}

Retain the unchanged J56 all-top kernel of K062/R048, its charged cutoff H=133, every source/target pencil, CM partner, finite projector, selected real coefficient, class mass and phase. Put f=168, n=169, m=164, with the original decomposition E\^{}169=S perpendicular W0 and signature(W0)=(4,1). At every334 unmixed G factor the original reduction is unipotent, by K068. No global unipotence is assumed below.

Write the actual matrix in the original coordinates as

\[
 h=\begin{pmatrix}I_m&B\\ C&D\end{pmatrix},\qquad
 C\ne0,\qquad13\le s\le21.
 \tag{GB1}
\]

This is the entire scalar-cross class A=I, including both globally unipotent and nonunipotent matrices and every incidence profile. The theorem is

\[
 \sum_{h\text{ in GB1}}|\text{actual normalized contribution of }h|
                        <2^{-23966}.
 \tag{GB2}
\]

It removes the intersection with the current remainder, with one new error charge23966 and every earlier charge retained. It is not complete support13 exhaustion, since other structural classes require their own bounds. It changes no uniform H8192 statement and asserts no final signed nonvanishing, period or DP44 cost.

\subsection{2. Rank-one cross block and exact reconstruction}\label{rank-one-cross-block-and-exact-reconstruction}

Let J0=diag(1,1,1,1,-1). Unitarity gives

\[
 C^*J_0C=0,\qquad B=-C^*J_0D.
 \tag{GB3}
\]

At a faithful complex embedding the image of C is a totally isotropic subspace of W0, hence has dimension at most one. Rank is preserved by that embedding; since C is nonzero, rank\_E(C)=1. If C=0, the second equation would give B=0 and h would fix S, the excluded fixed-block case. This establishes the exact dichotomy without using the nilpotence of h-I.

Use the original packet integers and test vectors from K088(IP4,IP14--16):

\[
 \begin{split}
 Q&=1427908193729121522601,\quad U=77Q^2-14,\\
 u&=(2,95192390527,383092671043,12529843548438452388609),\quad w=Qu,\\
 y_1&=(-976627,-153032,38026,0),\\
 y_2&=(-567678,695570,-172838,0),\\
 y_3&=(104583,-15309315176,-61610045224,2).
 \end{split}
 \tag{GB4}
\]

Thus u dot u=U, gcd(u\_i)=1, and the y\_j span w-perp in the four original positive W0 coordinates. Both these y\_j and every original e\_i in S belong to every actual source and target pencil at both CM components. For y\_j this is DO078(FP20--22), which proves membership separately in source and target, and K088(IP5--6); it requires no h-stable common lattice. The S inclusion follows from the W0 support of the shear.

\textbf{Case A.} For at least one of the three y\_j, the row

\[
                 \delta_i=\langle y_j,Ce_i\rangle_{J_0}
                          =\langle y_j,he_i\rangle_J
 \tag{GB5}
\]

is not zero. Since he\_i belongs to the actual target lattice and y\_j belongs to that same integral selfdual lattice, every delta\_i is in O\_E, at both CM partners and at every finite place. The equality uses the orthogonality of y\_j and e\_i. In particular there is no T denominator.

Choose an index j0 with delta\_j0 nonzero and put z=C e\_j0. Rank one and linearity give

\[
                     Ce_i=z\,\delta_i/\delta_{j_0}.
 \tag{GB6}
\]

Therefore delta and the one original pivot column z determine C exactly. Counting at most three y choices and164 pivot choices only enlarges the count. Division in GB6 is used for injective reconstruction over E, not for a claim about integral coordinate changes.

\textbf{Case B.} All three rows GB5 vanish. The positive parts of the columns of C then lie in E u, so uniquely

\[
                       Ce_i=(u\xi_i,\beta_i).
 \tag{GB7}
\]

The row xi is nonzero: if it were zero, isotropy would force every beta\_i=0 and contradict C!=0. Choose j0 with xi\_j0 nonzero and again z=C e\_j0. Rank one gives Ce\_i=z xi\_i/xi\_j0. Thus xi,z determine C, including its negative row. No square root of U is introduced. Primitivity of u gives integers a\_l with sum a\_l u\_l=1, and hence xi\_i=sum a\_l C\_li in the original coordinates.

For either case put v=D e\_n, the W0 part of the original negative column. Equation GB3 now reconstructs its S part:

\[
                     x=(he_n)_S=-C^*J_0v.
 \tag{GB8}
\]

The data (delta,z,v), or (xi,z,v), consequently determine the complete negative column. For fixed actual source and target pencils, the number of arrows with a fixed negative column is at most the retained full source stabilizer G\_n\textless2\^{}5022, by R051 §1. No global-unipotent rank-five fibre theorem is used.

\subsection{3. The common G ideal for the whole cross block}\label{the-common-g-ideal-for-the-whole-cross-block}

Put N=h-I. At a rank-one unmixed finite reduction, N is nilpotent, so N²=0. Write that rank-one matrix as a b\^{}t. Since N\_SS=0, either a\_S=0 or b\_S=0. If b\_S=0 then C vanishes at that CM component. If a\_S=0 then B vanishes there, and the inverse-adjoint condition at the conjugate component, together with h\^{}-1=I-N in this reduction, makes C vanish at the conjugate component. At rank zero C vanishes at both components. At rank at least two impose no divisibility for this count.

Choose one vanishing component by a fixed rule at each rank-one factor. Let J contain that prime at rank one and both E-primes at rank zero. There are at most4\^{}334=2\^{}668 such ideal patterns. The reasoning concerns the full original C, so in both cases

\[
       \delta\in J^m,\qquad
       \xi\in(JI_\Sigma)^m,\qquad
       N_E I_\Sigma=T^{-2s}.
 \tag{GB9}
\]

For xi, away T use its integer Bezout expression in the entries of C. At each active T component, Ce\_i=he\_i-e\_i belongs to the target lattice because e\_i belongs to both source and target. The target lattice is contained in pi\^{}-1 Lambda by its actual one-shear formula, with the dual statement at the CM partner. Thus xi has at most a simple pole at each of those two components. This proves the second membership in GB9 for every type; no type-II-only assertion is needed. Delta is globally integral by GB5 instead.

The retained Green envelope R051 §4/R056(BM6), applied with this C orientation, is

\[
          |a(h)|\le2 Q^{-335(f-1)}(N_EJ)^{335/2}.
 \tag{GB10}
\]

It depends only on the norm contributions of finite rank zero, one and at least two; reversing which partner is selected at rank one does not change those norms. Equivalently, the baseline rank\textgreater=2 factor is q\^{}-335, rank zero is restored by the two-prime norm q², and the one-prime factor q\^{}(335/2) is a sufficient upper restoration at rank one. The two mixed projectors still have modulus at most one. Grouping all higher finite ranks as one class requires no second enumeration of their ranks. Global nonunipotence causes no change to this finite envelope, since K068's actual finite reductions are unipotent.

\subsection{4. Archimedean bounds and the elementary ideal count}\label{archimedean-bounds-and-the-elementary-ideal-count}

At every actual complex embedding the Euclidean operator norm of h is at most133. Since he\_i has S part e\_i, the orthogonal Euclidean decomposition gives \textbar\textbar Ce\_i\textbar\textbar\textless133. The three exact squared lengths of the y\_j are

\[
                 978665066829,\quad835948910828,
                 \quad4030172803672366725045.
 \tag{GB11}
\]

The strict integer comparison 133² max\textbar\textbar y\_j\textbar\textbar²\textless2\^{}86 proves

\[
              |\delta_{i,\tau}|<2^{43}\quad\hbox{in Case A}.
 \tag{GB12}
\]

In Case B, the positive part u xi\_i has Euclidean norm sqrt(U)\textbar xi\_i\textbar, so

\[
              |\xi_{i,\tau}|<133/\sqrt U<2^{-65}.
 \tag{GB13}
\]

The last inequality is checked as U\textgreater133²2\^{}130 with exact integers. No normalization of the isotropic line or algebraic square root has been assumed.

Here is the full elementary packing lemma used for both cases. For a nonzero fractional ideal L of E put z\_L=(N\_E L)\^{}(1/(2f)). Every nonzero b in L has \textbar N\_E b\textbar\textgreater=N\_E L, so AM--GM in the half-trace Euclidean metric gives \textbar\textbar b\textbar\textbar\textgreater=sqrt(f)z\_L. Packing disjoint balls of half that radius around points whose complex coordinates are bounded by B shows

\[
 \#\{b\in L:|\tau b|<B\ \forall\tau\}
                      \le(1+2B/z_L)^{2f}.
 \tag{GB14}
\]

For m coordinates take the product of these bounds. If at least one coordinate is nonzero, its field norm forces z\_L\textless B. Since335-2m=7\textgreater0, the function

\[
 z^{335f}(1+2B/z)^{328f}
                   =z^{7f}(z+2B)^{328f}
 \tag{GB15}
\]

is increasing for z\textgreater0. Therefore, throughout the actual nonzero-vector range,

\[
 z^{335f}(1+2B/z)^{328f}<B^{335f}3^{328f}.
 \tag{GB16}
\]

This includes zero coordinates without enumerating their positions. It is a bound for the actual arithmetic ideal, not an equidistribution assertion or an integral change of basis.

Apply GB16 to L=J in Case A and to L=J I\_Sigma in Case B. After multiplication by the numerator in GB10, the complete vector factors are respectively

\[
                  2^{43\cdot335f}3^{328f},
       \qquad T^{335s}2^{-65\cdot335f}3^{328f}.
 \tag{GB17}
\]

\subsection{5. The two remaining original columns and all multiplicities}\label{the-two-remaining-original-columns-and-all-multiplicities}

The pivot z belongs to M\_d,W0, and the negative column v belongs to pi\^{}-1 M\_d,W0. The former is contained in the latter as an upper lattice count. R051 §1 proves that this latter actual overlattice contains the original O\_E\^{}5 and has paired index T\^{}10 at each active place. Its direct original-ring coset packing gives, for a W0 vector of norm at most H at every embedding,

\[
              \#\{z\},\ \#\{v\}<2^{14861}T^{10s}.
 \tag{GB18}
\]

The original strict bound (2H+1)\^{}(10f)\textless2\^{}13543 is even smaller. GB18 is independent of the support size except for the displayed index, and is uniform in the actual source/target labels, including zero. In particular v is counted in the true pi\^{}-1 M\_d; it is not incorrectly placed in I\_Sigma\^{}5. Its norm is at most\textbar\textbar he\_n\textbar\textbar\textless=H. For z the preceding \textbar\textbar Ce\_i\textbar\textbar\textless H applies.

For each support s retain once each of the following factors: support choices2\^{}168; the common G patterns2\^{}668; pivot choices164\textless2\^{}8; the at most three y choices\textless2² in Case A (the same factor is harmless in Case B); the full source fibre2\^{}5022; the single Green Euler factor2; and the two counts GB18. The entire ordered-pair/weight sum is the unchanged

\[
       P_T^s=\left(2(T-1)^2/(T-2)\right)^s<2^{997s}.
 \tag{GB19}
\]

It already counts all three incidence types. No extra3\^{}s, reciprocal type-III weight, or missing zero-target coarea is inserted. The normalized selected real coefficient has modulus at most one; this theorem does not use DO080. These are the original full-column counts, so no affine-real slice is invoked.

Use the exact carried one-sided bounds

\[
            \log_2 T<L_T=1019333/1024,\qquad
            \log_2 Q>L_Q=287843/4096,
            \qquad3^5<2^8.
 \tag{GB20}
\]

Set C0=1+668+168+8+2+5022+2(14861)=35591. The two complete upper exponents are

\[
 \begin{split}
 E_A(s)&=C_0+43\cdot335f+(8/5)328f
                 +20sL_T+997s-335\cdot167L_Q,\\
 E_B(s)&=C_0-65\cdot335f+(8/5)328f
                 +355sL_T+997s-335\cdot167L_Q.
 \end{split}
 \tag{GB21}
\]

Both are increasing in s. At s=21, E\_A(21)\textless-948668 and

\[
                   E_B(21)=-490927563/20480<-23971.
 \tag{GB22}
\]

The companion certificate verifies the complete rational expressions, not floating approximations. There are nine support values and two cases. Their union is below18 times2\^{}-23971, hence below2\^{}-23966, proving GB2. No further multiplicity remains after the full pair sum, fixed finite patterns, reconstructed column and source fibre have been counted.

\subsection{6. Scope, geometric meaning and verification}\label{scope-geometric-meaning-and-verification}

Every actual matrix with A=I and C!=0 is in one of Cases A and B. The rank-one cross block follows from signature one and unitarity; it does not bound rank(h-I) by two or make h globally unipotent. The G divisibility uses only the actual unipotent finite reductions. The proof is therefore valid for the entire scalar-cross class at supports13..21, with no assumption on t or finite ranks beyond the carried K068 condition.

Remove only its intersection with the then-current remainder, preserving all historical error charges even where a stronger new set overlaps them. DO081's globally rank-two-unipotent sector, if already recorded, is not charged again as a separate subset inside this theorem; its historical numerical upper bound remains in the accumulated error. The certificate verifies the complete sums both with and without that intervening charge. Actual integration records the chronological order once. The support floor remains13 until the other F/IC8 and fixed-block classes at that support are separately exhausted.

The new arguments concern the original Hermitian modules over Spec(O\_E), integral source/target pairings, a rank-one determinantal cross block and its unchanged complex realization. The module/stabilizer context remains SGA3 I4.7.3/4.7.3.1 and the exterior-power context SGA6 VI Theorem3.3. The CM Picard/polarization/biextension data retain FGA232 Theorem3.1 pp148--149; FGA236 Corollaries1.7/2.3 pp229/232; SGA7 I VII2.7/2.9 pp169--174 and VIII7.4(b) pp310--311. Formal crystals and Cartier duality retain cote4 scan003 §§8.1.1--8.1.2, scans005--006 §§8.4.1--8.4.3 and cote54 scans005--007. These are contextual anchors, not attributions of this new count to those texts.

No new lambda-ring or K-theory identity, l-adic weight result, CM-motive identification or motivic L-function theorem follows. DP\_MULT Q3e and THETA3 G1/A7 have a strictly smaller finite gate in the stated sector. The prescribed period, literal full mixed-CM source, sharp selected-image content, concrete032 cycle pin/value and every changed-family DP44 cost remain open. CX\_GRH CX-ALG-024 and GCT\_PLETH T-052 retain their distinct missing comparisons.

Exact carried inputs are K062's normalized kernel; K068's finite-unipotent residual; K069's A=I block dichotomy; K088(IP4--6,IP14--16) and DO078(FP20--22) for the original packet and two-lattice membership; R051 §1 and §4 for original target-column/fibre and Green factors; R056(BM6) for the same rank-pattern envelope; and DO075(CR15) for the complete pair/weight sum. The certificate checks literal packet identities, archimedean inequalities, finite rank-one orientation models, exact reconstruction models, full exponents and chronological error totals. It does not enumerate rational isometries or automate the arithmetic lattice interfaces. No Lean build is run.

\clearpage

\section{AM-R-060 --- complete all-top supports thirteen and fourteen exhaustion}\label{am-r-060-complete-all-top-supports-thirteen-and-fourteen-exhaustion}

Author: GPT-6 Astra Pro, inherited effort not externally exposed; ALP\_MACHINE \nolinkurl{/root}. 22 November2026 UTC. Complete proof; author and independent whole checks are recorded separately before promotion.

\subsection{1. The unchanged kernel and the new remainder}\label{the-unchanged-kernel-and-the-new-remainder}

Retain the original actual J56 all-top assignment, K044 upper endpoint, permitted packet point, finite projectors, ordered scalar-pencil incidences, class masses, phases, selected compact representation and Hilbert metric. The degrees are {[}E:Q{]}=336 and {[}F:Q{]}=f=168, the Hermitian rank is169, and the charged all-top cutoff is H=133. The independent uniform H8192 family is unchanged.

After R059 the current all-top remainder still has support at least13, with55 chronological error charges. This theorem proves an actual remainder R\_top\^{}(15), consisting only of original terms with support s\textgreater=15, such that

\[
 \mathcal E_{\rm top}=1+\operatorname{Re}\sum_{h\in R_{\rm top}^{(15)}}w(h)
                          +\epsilon,\qquad|\epsilon|<2^{-230}.
 \tag{J15.1}
\]

It adds exactly one charge172255 for the fresh nonscalar union proved below. It also removes the stated higher-support IC8 and F-unipotent cells through17, without claiming complete exhaustion of those higher supports. No earlier charge is erased, subtracted or repeated. The uniform support floor remains11 with its original31 charges.

The sufficient inequality

\[
             \operatorname{Re}\sum_{h\in R_{\rm top}^{(15)}}w(h)
                          >-1+2^{-230}
 \tag{J15.2}
\]

remains open. This is a complete reduction of two support levels, not a proof of positivity, an actual prescribed period or a DP44 cost conclusion.

\subsection{2. The exceptional small-block count for supports13 through17}\label{the-exceptional-small-block-count-for-supports13-through17}

Use the original coordinates S=\textless e5,\ldots,e168\textgreater, W0=\textless e1,e2,e3,e4,e169\textgreater, and A=h\_SS. Consider the actual K069 exceptional class IC8 with r=rank\_E(A-I) equal to1 or2 and13\textless=s\textless=17. Keep all incidence types and do not assume global unipotence.

DO075(CR3--14) proves the following actual interfaces independently of its numerical application s=12. A is identity outside at most four original vertices, with at most binom(164,4)29\^{}16\textless2\^{}112 choices after padding. Every unmixed reduction of N=h-I has rank at least two. At a rank-two reduction the160 unaffected original S coordinates of the negative column vanish at both CM partners. Write J0 for the paired ideal of those rank-two factors. Then

\[
 1\le N_EJ_0\le Q^{2(f-1)},\qquad
 |a(h)|\le2Q^{-501(f-1)}(N_EJ_0)^{83},
 \tag{J15.3}
\]

where Q=g0g1=1427908193729121522601, and there are at most2\^{}334 patterns. The other unmixed factors have rank at least three; the two mixed factors have modulus at most one. These finite facts are the paired nilpotent-rank argument CR5--9, not global nilpotence or a positivity argument over the residue field.

Put

\[
 L=J_0I_\Sigma,\quad N_EI_\Sigma=T^{-2s},\quad
 N_EL=z^{2f},\quad z_0=T^{-s/f},\quad
 z_1=Q^{(f-1)/f}T^{-s/f}.
 \tag{J15.4}
\]

Thus z0\textless=z\textless=z1. Retain the exact half-trace covolume constant

\[
 D_0=7^{140}29^{144}5^{84}13^{84},\qquad
 V_d(C)=\frac{(528C^2)^{fd}}{(fd)!D_0^d}.
 \tag{J15.5}
\]

The full original-coordinate bounds CR10--13 are

\[
 \#x_{\rm free}\le V_{160}(66.5+4200z)z^{-320f},
 \tag{J15.6}
\]

with the four exceptional coordinates bounded by V4(67)T\^{}(8s) whenever z0\textless2\^{}-71, and the five W0 negative-column coordinates bounded by V5(223)T\^{}(10s). The latter is the actual full target-overlattice count, uniform in every target including zero. It has no affine-real coarea slice. The source fibre remains G\_n\textless2\^{}5022. All are the unchanged original coordinates; the arbitrary-ideal transference radius in J15.6 is not replaced by the smaller original-ring radius.

The actual integers satisfy, for every s=13,\ldots,17,

\[
                  z_0<2^{-71},\qquad z_1<1/2.
 \tag{J15.7}
\]

These are checked as 2\^{}(71f)\textless T\^{}s and2\^{}f Q\^{}(f-1)\textless T\^{}s. No lower bound on z1 is needed. At z0 the radius in J15.6 is less than67; at z1 it is less than2167.

Multiplying J15.6 by the numerator in J15.3 leaves T\^{}(166s) times

\[
               V_{160}(66.5+4200z)z^{-154f}.
 \tag{J15.8}
\]

As proved in DO075(CR14), its logarithm is convex as a function of log z: its only nonlinear term is320f log(66.5+4200 exp(log z)). Thus its maximum on the interval is at one of the two endpoints. This is an upper envelope over all actual ideals; it requires neither an endpoint ideal to exist nor a density assumption.

At z0, the complete T power including the other nine coordinates is338s and the Green denominator remains Q\^{}(501(f-1)). At z1, retaining its exact power gives

\[
 z_1^{-154f}=Q^{-154(f-1)}T^{154s},
 \tag{J15.9}
\]

so the complete T power is again338s and the denominator is Q\^{}(655(f-1)). This is the new endpoint calculation. The old special s12 lower bound z1\textgreater1/4 is neither used nor extended to larger s.

\subsection{3. Complete exceptional-class envelope}\label{complete-exceptional-class-envelope}

The original exact factorial comparisons, rechecked by the companion certificate, are

\[
 V_{160}(67)<2^{-43274},\quad V_{160}(2167)<2^{226353},
 \quad V_4(67)<2^{2489},\quad V_5(223)<2^{5757}.
 \tag{J15.10}
\]

Retain every factor from DO075: the A choices2\^{}112, rank-two/higher patterns2\^{}334, support choices2\^{}168, full source fibre2\^{}5022, one Green factor2, and the full ordered-pair/weight sum P\_T\^{}s\textless2\^{}(997s). That sum already covers every incidence type and both CM partners in the carried convention. There is no additional3\^{}s factor, no rank-one orientation choice and no factor two for taking the maximum of two endpoints. Put

\[
 C(s)=112+334+168+5022+1+997s,
 \quad L_T=1019333/1024,\quad L_Q=287843/4096.
 \tag{J15.11}
\]

The actual one-sided inequalities are log2 TL\_Q. The two complete bounds are

\[
 \begin{split}
 E_0(s)&=C(s)-43274+2489+5757+338sL_T-501\cdot167L_Q,\\
 E_1(s)&=C(s)+226353+2489+5757+338sL_T-655\cdot167L_Q.
 \end{split}
 \tag{J15.12}
\]

Both increase with s and E1\textless E0 throughout. At the endpoint,

\[
 E_0(17)=-705573041/4096<-172259,\qquad
 E_1(17)=-7003927123/4096<-172259.
 \tag{J15.13}
\]

Consequently for each of the five supports13..17, the entire IC8 rank1/2 union contributes less than2\^{}-172259. The original padded A count covers both ranks at once. All source-lattice, finite-pattern and pair multiplicities have already been included.

\subsection{4. Fresh complete F,r1 globally unipotent cells}\label{fresh-complete-fr1-globally-unipotent-cells}

For the fixed-space structural class F with block rank r=1, retain the full conditional counts DO065(UR9--14), explicitly distinguished from its old t\textless=8 removed union in R057 §3. Once global unipotence is given, those counts impose no bound on t. Signature-one unipotence makes the full rank R=rank\_E(h-I) equal to1 or2 in this class. The complete estimates are

\[
 \begin{split}
 F_1(s)&=-103530801/128+8966s,\\
 F_2(s)&=-408071827/256+14941s.
 \end{split}
 \tag{J15.14}
\]

They use the more conservative H8192 column count and therefore apply to the actual H133 subset, keeping the original class weights and real modulus bound one. They include the actual rank charts, finite witnesses, ordered pencils and source fibre; no additional multiplier is inserted.

Both increase in s. At s17 their exact values are

\[
                    F_1(17)=-84020785/128,\qquad
                    F_2(17)=-343048595/256,
 \tag{J15.15}
\]

each less than-172259. Thus the ten full F-unipotent cells for s13..17 and R=1,2 are each bounded by2\^{}-172259. Charge these full higher-support cells afresh; R057's701241 was only the support12 union. Previously removed smaller-t subsets are intersected out rather than their numerical charges being subtracted.

Combining the five IC8 and ten F envelopes gives the single fresh union bound

\[
                  15\cdot2^{-172259}<2^{-172255}.
 \tag{J15.16}
\]

\subsection{5. Exhaustive original structural tables at supports13 and14}\label{exhaustive-original-structural-tables-at-supports13-and14}

The original K068/K069 residual has r=rank(A-I)\textless=33. No unproved r\textless=22 global conclusion is used. For A!=I, K069 partitions the classes into F, where h fixes K\_A=ker(A-I) intersect ker(A*-I) pointwise, and the complementary IC8 class. The following table is exhaustive at each s in\{13,14\}.

\begin{longtable}[]{@{}
  >{\raggedright\arraybackslash}p{(\columnwidth - 2\tabcolsep) * \real{0.5000}}
  >{\raggedright\arraybackslash}p{(\columnwidth - 2\tabcolsep) * \real{0.5000}}@{}}
\toprule\noalign{}
\begin{minipage}[b]{\linewidth}\raggedright
Actual A!=I class
\end{minipage} & \begin{minipage}[b]{\linewidth}\raggedright
Complete removal at s13 and14
\end{minipage} \\
\midrule\noalign{}
\endhead
\bottomrule\noalign{}
\endlastfoot
F,2\textless=r\textless=33 & K069(IC20--21) has minimum support threshold22 over these ranks; both supports lie in its previously charged4084 union. \\
IC8,3\textless=r\textless=33 & DO064(PM22--23) has minimum threshold16 over these ranks; both supports lie in its previously charged4084 union. \\
IC8,r=1 or2 & Sections2--3 above remove the whole sector, without global unipotence, within the single fresh J15.16 charge. \\
F,r=1,nonunipotent & DO074 applies at r1 and excludes t\textless=10. For t11 through14, R050 §5 includes every admissible pair with s13 or14: its t11/12/13/14 upper endpoints are59/48/37/26. The existing4088 union is retained. \\
F,r=1,globally unipotent & Section4 removes both possible full ranks1 and2 and all profiles, within J15.16. \\
\end{longtable}

Now let A=I. If h fixes S it is I\_S direct-sum h0 on W0. If it does not, it is the scalar-cross case C!=0 of R059. The following second table is exhaustive.

\begin{longtable}[]{@{}
  >{\raggedright\arraybackslash}p{(\columnwidth - 2\tabcolsep) * \real{0.5000}}
  >{\raggedright\arraybackslash}p{(\columnwidth - 2\tabcolsep) * \real{0.5000}}@{}}
\toprule\noalign{}
\begin{minipage}[b]{\linewidth}\raggedright
Actual A=I class
\end{minipage} & \begin{minipage}[b]{\linewidth}\raggedright
Complete removal at s13 and14
\end{minipage} \\
\midrule\noalign{}
\endhead
\bottomrule\noalign{}
\endlastfoot
Fixed S,nonunipotent & DO067 excludes t\textless=9. The existing K081/K082 tables cover all t10,\ldots,14 at these supports; their union thresholds are168,164,149,134,119. All historical charges remain. \\
Fixed S,globally unipotent of full rank1 & K077(AP23,§7) removes every profile at supports10..14 with its existing1436 charge. \\
Fixed S,globally unipotent of full rank2 & R058(HB1--25) removes all profiles at supports13..23 with its existing39409 charge. \\
Scalar cross,globally unipotent or not & R059(GB1--22) removes the entire class at supports13..21 with its existing23966 charge. The earlier DO08125470 charge remains once. \\
Full rank0 & h=I has support0 and is absent at s13 and14. \\
\end{longtable}

The signature-one classification allows no other nontrivial full rank for a globally unipotent matrix. These are the same original F/IC8 and fixed-S/scalar-cross partitions used by R057, with the new broad R059 count now eliminating both scalar-cross alternatives simultaneously. No sign, local nonvanishing or a failed upper estimate is used to exclude an additional class.

Together, the two tables remove every actual support13 and support14 term. R057 had already removed every smaller support. Thus, after adding J15.16 to the retained numerical error, every remaining term has support at least15.

\subsection{6. Chronological error accounting and remaining gates}\label{chronological-error-accounting-and-remaining-gates}

The immediately preceding actual array is \nolinkurl{checks/r059_joint_remainder_error.json}, with55 all-top and31 uniform exponents. Append172255 only to the all-top array. It now has56 entries, exactly one equal to231 and every other entry at least685. Hence

\[
 \sum_{e\ {
m all\!\!\!-top}}2^{-e}
 \le2^{-231}+55\cdot2^{-685}<2^{-230}.
 \tag{J15.17}
\]

The unchanged uniform list has exactly one231 and every other entry at least797; its sum remains below2\^{}-230. The certificate also verifies both complete dyadic sums directly with integers. It retains all earlier overlapping absolute charges, including R05839409, DO08125470 and R05923966, exactly once. The previously charged support12 removals are not added again.

At each stage remove only the actual intersection of the new bounded set with the current remainder. This makes the removed sets disjoint while preserving the conservative numerical error. Combining the unchanged energy identity with the complete tables and J15.17 proves J15.1.

The all-top finite/real sufficient gate for Q1/Q2 is strictly smaller, now over support\textgreater=15. The current bound also removes every IC8 block-rank1/2 and F,r1 globally unipotent cell at s15,16,17, but other sectors there remain open. The actual prescribed Bessel/Mram value, concrete032 cycle pin and pairing, full mixed-CM source, sharp selected arithmetic-image content and every changed-family DP44 cost are still unproved. There is neither a shorter valid chain nor a universal strict negative DP44 difference. Q3 has no Hecke RH conclusion, and Q4 retains its missing global arithmetic charge comparison.

\subsection{7. Algebraic geometry, sources and verification}\label{algebraic-geometry-sources-and-verification}

The integral Hermitian modules, original labelled correspondences and CM realizations remain unchanged. The new exceptional-class estimate counts a paired determinantal ideal in160 original coordinates and retains the exact endpoint power of that ideal. R059 reconstructs the original scalar cross block using integral source/target pairings. Joining these with the carried structural partition changes no Picard scheme, polarization, theta line or finite-level input.

The carried module/stabilizer context is SGA3 I4.7.3/4.7.3.1 and the determinantal/exterior-power context SGA6 VI Theorem3.3. The Picard and biextension anchors are FGA232 Theorem3.1 pp148--149; FGA236 Corollaries1.7/2.3 pp229/232; SGA7 I VII2.7/2.9 pp169--174 and VIII7.4(b) pp310--311. Formal groups, crystals and Cartier duality retain cote4 scan003 §§8.1.1--8.1.2, scans005--006 §§8.4.1--8.4.3, and cote54 scans005--007. These are contextual locators, not sources of the new numerical estimates. No new lambda-ring/K-theory identity, l-adic sheaf or weight theorem, motive identification or motivic L-function statement follows. The change to DP\_MULT Q3e and THETA3 G1/A7 is precisely J15.1; CX\_GRH CX-ALG-024 and GCT\_PLETH T-052 keep their distinct missing arithmetic maps.

Exact dependencies: R057(J13.1--8); K069(IC1--8,IC20--21); DO064(PM5--9,PM22--23); DO065(UR9--14); DO067's fixed-W0 two-trace criterion; DO074's stated r\textless=22,t\textless=10 criterion; DO075(CR3--15); R050 §5; K077(AP23,§7); K081/K082 complete nonunipotent tables; R058(HB1--25); and R059(GB1--22). The fresh certificate checks the actual prime/endpoint comparisons, all four exact factorial ceilings, complete rational envelopes, all table ranges at both supports, the15-cell union and every chronological error. It does not enumerate rational matrices or automate the geometric partition. No Lean build is performed.

\clearpage

\section{Canonical integration note --- AM-DO-082}\label{canonical-integration-note-am-do-082}

The frozen whole-reviewed proof below retains its original R059 chronological55-entry baseline and historical support floor13. Before this canonical integration, AM-R-060 independently removed all supports13/14, with one172255 charge and floor15. The live accounting therefore appends the single DO082895 charge to that56-entry list:57 all-top and31 unchanged uniform charges, each total below2\^{}-230. The current support floors are15 and11. The exact supplemental check is checks/do082\_joint\_remainder\_error.json. The original proof, scripts and reviews are unmodified; no internal10531 or896 charge is appended separately.

Root whole-read all six sections SX1--20, the complete script and the independent review, and ran a separate1452-check replay. The independent Kummer whole-proof/script review and its separate1452-check replay also passed. The theorem removes every globally rank-two-unipotent scalar-cross support23 profile; it does not exhaust other support23 classes or prove the signed sufficient inequality.

\begin{center}\rule{0.5\linewidth}{0.5pt}\end{center}

\section{AM-DO-082 --- the complete globally unipotent scalar-cross sector at support twenty-three}\label{am-do-082-the-complete-globally-unipotent-scalar-cross-sector-at-support-twenty-three}

Author: \nolinkurl{/root/definite_oscillator}, 22 November2026 UTC. New finding; the frozen DO080 and DO081 proofs remain unchanged.

\subsection{1. Exact sector and result}\label{exact-sector-and-result}

Use the unchanged actual all-top J56 kernel, H=133, f=a=168, m=164 and n=169, with the original E\^{}n=S perpendicular W0, signature(W0)=(4,1). Put R=17424/133 and retain the exact Q,T,U,w and integer vector u\_0 of DO081(IP1--3). Consider precisely \[
 h=I+N,\quad h_{SS}=I,\quad\operatorname{rank}_E N=2,\quad
 N^3=0,\quad\operatorname{rank}_E N^2=1,\quad x=(Ne_n)_S\ne0,
 \qquad s=23.
 \tag{SX1}
\] Global unipotence is explicitly assumed. Every source/target pencil, CM partner, selected module and near-point trace, finite projector, class mass and phase remains the actual one.

\textbf{Theorem.} The entire SX1 sector, including all incidence profiles and every actual coordinate, finite-pattern, source/target and class-fibre multiplicity, has total absolute contribution less than \textbf{2\^{}-895}. Remove only its intersection with the current remainder, retaining every earlier charge. No all-sector support-floor update follows from this result alone.

Use the two exact cases of DO081: Case A has a nonzero pairing with one of the three fixed transverse integer vectors y\_j; Case B has im(N²) in the original plane spanned by w and e\_n.~Write i,u,t for the counts of types I,II,III. The present result combines a fresh evaluation of the complete elementary counts for Case A and for Case B with u\textless=22, and a new stronger count for Case B with i=t=0,u=23.

\subsection{2. The transverse case and every non-all-II plane profile}\label{the-transverse-case-and-every-non-all-ii-plane-profile}

DO081's integral pairings, reconstruction, weighted ideal packing and all multiplicities are independent of s except for their displayed powers. Repeating those counts at23 is legitimate; it is not an extension of the previously stated support13--22 theorem by assertion.

For clarity, Case A counts delta\_y in a common column-oriented G ideal, with no T denominator, together with alpha in I\_Sigma,F² and the nonzero scalar t\_y in I\_Sigma. Their exact reconstruction determines x. Its weighted vector factor is2\^{}(40*335f+4fm), scalar factors are2\^{}(9f)T\^{}(2s) and2\^{}(92f)T\^{}(2s), W0 affine columns cost2\^{}5379 T\^{}(8s+r0), and the genuine full-column fibre is G5\textless2\^{}107. The three choices of y\_j, all4\^{}334 common patterns, support choices, type assignments and complete weighted pairs are already included in DO081(IP29). At23 its full exponent is \[
                    E_A(23)=-4990249/4<-10532.
 \tag{SX2}
\] This includes every Case-A profile, including all-II.

In Case B, write ell=(b u\_0,1), with b bar(b)=1/U, and put delta'=b bar(x), B=alpha b. The original fixed integer Bezout row gives both as linear combinations of actual entries. The safe count in DO081(IP23--30) retains the exact T exponent335u+15s before replacing u by s. For s=23 and u\textless=22 it therefore has exponent at most \[
 E_B(23)-335\cdot996
 =1292511/4-333660=-42129/4<-10532.
 \tag{SX3}
\] All type assignments, supports and actual weighted pairs were counted in this envelope. There is no further299-profile multiplicity to add. The two complete bounds SX2--3 have union less than2\^{}-10531. This leaves only the original-plane, all-II profile, treated next.

\subsection{3. Stronger original-entry ideals in the all-II plane}\label{stronger-original-entry-ideals-in-the-all-ii-plane}

For the remaining profile, keep the exact rational line and original entries \[
 N^2=-\alpha\ell\ell^*J,\quad\ell=(b u_0,1),\quad
 b\bar b=1/U,\quad0<\alpha_\tau\le R,\quad
 \delta'=b\bar x,\quad B=\alpha b\ne0.
 \tag{SX4}
\] Here b and B belong to E, alpha belongs to F, and the rational line is never declared finite integral. The integer Bezout row has value one on u\_0, as in DO081(IP21); it expresses B as that row applied to the four positive W0 entries of N²e\_n.

At each of the334 unmixed G places let d\_v be the reduction rank of N, equal to0,1 or2. At rank zero all entries of N² have valuation at least two. At rank one the finite nilpotent reduction has square zero, so every entry of N² has valuation at least one. These statements hold at both inverse-adjoint CM components. Define the F-ideal \[
 J_F=\prod_v\mathfrak p_v^{2-d_v},\qquad
 J=N_FJ_F=Q^{334}/\prod_v q_v^{d_v}.
 \tag{SX5}
\] For a common column-oriented pattern the actual E-ideal J\_E of DO081 has exactly this absolute norm J. There are at most4\^{}334 such patterns. The scalar B inherits the square-entry divisibility at both components, and the all-II entrywise T pole bound is two. Thus \[
                 B\in(J_F O_E) I_\Sigma^2.
 \tag{SX6}
\] This is stronger than the earlier safe membership in I\_Sigma². In particular N\_E(J\_F O\_E)=J², not J. Extension to both CM partners is essential to the count.

The unchanged original11-entry is (N²)\_11=-4alpha/U. The actual U is odd and coprime to Q,T. At primes dividing U use this identity; at primes above2 use the ordinary integral nn entry, and elsewhere use the original away-T integrality. Combine with the same square valuations at G and the all-II double poles. Exactly as in K088(IP20--23) and DO078(FP26--28), for this support as well, \[
                    \alpha\in U J_F I_{\Sigma,F}^{2}.
 \tag{SX7}
\] The argument uses only entries of the unchanged matrix and is independent of the number of active primes. It does not use a finite-rank-budget contradiction or an anisotropy assertion.

Put \[
 z=(J T^{-2s})^{1/(2f)},\quad M=\sqrt{R/U},\quad M_B=R/\sqrt U.
 \tag{SX8}
\] Then the ideal in SX6 has full absolute norm z\^{}(4f), and SX7 gives the precise \textbf{norm inequality} \[
                       N_F(\alpha)\ge U^f z^{2f}.
 \tag{SX9}
\] No pointwise lower bound alpha\_tau\textgreater=U z² is inferred.

\subsection{4. Reconstructing from B and counting at a common ideal parameter}\label{reconstructing-from-b-and-counting-at-a-common-ideal-parameter}

The nonzero B determines alpha uniquely among totally positive F-elements, because \[
                         \alpha^2=U B\bar B.
 \tag{SX10}
\] If no such element lies in F, that B gives no actual matrix and may be discarded. There is no additional2\^{}f sign factor, no field extension and no choice of an archimedean square root in the count. Once B is fixed, b=B/alpha and x=overline(alpha delta'/B) are determined by delta'. We therefore count B and delta' without an independent alpha count.

Let k,1\textless=k\textless=164, be the number of nonzero original E-coordinates of delta'. The actual ideal is \[
 \delta'\in(J_E I_\Sigma)^m,\qquad N_E(J_E I_\Sigma)=z^{2f},
 \quad\|\delta'_\tau\|^2=\|x_\tau\|^2/U\le R/U=M^2.
 \tag{SX11}
\] The half-trace ideal norm minimum and summation over the k nonzero coordinates give \[
                                k z^2\le M^2.
 \tag{SX12}
\] For a fixed coordinate support the actual delta' count is at most(1+2M/z)\^{}(2fk), by the same disjoint-ball packing as DO081(IP13). There are binom(164,k) supports.

The SX6 ideal has half-trace minimum at least sqrt(f)z². Since \textbar B\_tau\textbar\textless=M\_B at every embedding, the one-E-coordinate packing gives \[
                           \#B\le(1+2M_B/z^2)^{2f}.
 \tag{SX13}
\] This retains its actual fractional ideal; it is not a density division or a substitution of the small original-ring covering constant.

For the actual selected trace, including both compact factors and the near-point conjugation, DO080(RT3 and section7, using RT24) gives \[
 |c(h)|<2^{-929f/2}N_F(\alpha)^{-1/2}
           \le2^{-929f/2}U^{-f/2}z^{-f}.
 \tag{SX14}
\] The actual transport is the unchanged LIFT41/K085 near-point transport, with Cartan radius below1/2 at every label, so the hypothesis of DO080 section7 holds. Its exact all-top word and Hilbert normalization are unchanged. The norm inequality SX9 is the only further input used in the last step of SX14.

The finite factor is2 Q\^{}(-335(f-1)) J\^{}(335/2), with J\^{}(335/2)=T\^{}(335s)z\^{}(335f). For each k, after including SX13--14 and the delta' count, the z-dependent part is \[
 z^{(330-2k)f}(z+2M)^{2fk}(z^2+2M_B)^{2f}.
 \tag{SX15}
\] All three factors are increasing for z\textgreater0 because330-2k\textgreater=2. Hence the complete actual discrete-ideal contribution is bounded by its value at z=M/sqrt(k), using SX12. No existence of an ideal at that endpoint is asserted.

\subsection{5. Exact full endpoint with all finite and class factors}\label{exact-full-endpoint-with-all-finite-and-class-factors}

For fixed B, alpha and therefore Re(h\_nn)=1+alpha/2 are fixed. The actual affine W0 negative-column count is the same conservative2\^{}5379 T\^{}(8s) as before. In the active all-II profile the unique fixed label cannot be zero: fixing the standard M0 would make the place inactive. Thus no zero-target coarea factor occurs. The complete diagonal-label/rho sum is at most(T-1)/(T-1)=1 at each place. Both CM components are already included.

After delta' and B reconstruct x, the actual full-column isometry fibre remains G5\textless2\^{}107. Retain support choices2\^{}168, all4\^{}334=2\^{}668 patterns, and one global Green factor2. There is one type assignment, and no transverse-vector or pivot choice. These factors and the affine-volume constant give C0=6323, exactly as in DO081(IP28).

Substituting z=M/sqrt(k) into the preceding expressions gives the complete k-sector bound \[
\begin{split}
 U_k<&\;2^{6323-929f/2}Q^{-335(f-1)}T^{343s}
 R^{167f}U^{-335f/2}\binom{164}{k}\,k^{-167f}\\
 &\hspace{15mm}\cdot(1+2\sqrt k)^{2fk}
                (1+2k\sqrt U)^{2f}.
\end{split}
\tag{SX16}
\] The T exponent343s consists only of the Green conversion335s and the affine W0 count8s. The stronger B ideal, and the fact that B determines alpha, are why no separate scalar T powers remain. All quantities being counted can be enlarged to their lattice supersets, while their weights are applied only to the actual pattern under consideration.

Use1+2k sqrt(U)\textless=3k sqrt(U), since k\textgreater=1 and U\textgreater1, and3\^{}336\textless2\^{}533. Let r\_k=ceil(sqrt(k)), and keep the exact one-sided rational bounds \[
 \log_2T<1019333/1024,\quad\log_2Q>287843/4096,\quad
 \log_2R<3602/512,\quad\log_2U>601356/4096.
 \tag{SX17}
\] Each is checked by strict integer powers against the actual data. At s=23 the common exponent after these substitutions is \[
\begin{split}
 C={}&6323+343\cdot23(1019333/1024)-335\cdot167(287843/4096)\\
    &+167f(3602/512)-(333f/2)(601356/4096)-929f/2+533\\
   ={}&-241526103/4096.
\end{split}
\tag{SX18}
\] It remains to bound the exact rational factor \[
                  F_k=\binom{164}{k}
                         k^{-165f}(1+2r_k)^{2fk}.
 \tag{SX19}
\] The complete exact scan for k=1..164 checks a strict integer binary upper exponent for each F\_k and its lower neighboring bracket. The maximum upper exponent is58062, attained at164. Consequently every complete k-sector satisfies \[
                     \log_2 U_k<-3704151/4096<-904.
 \tag{SX20}
\] Since164\textless2\^{}8, the full all-II original-plane sector has absolute contribution less than2\^{}-896. This is an explicit complete finite scan, not an assertion of analytic monotonicity of F\_k in k.

\subsection{6. Exhaustion, errors and verification limits}\label{exhaustion-errors-and-verification-limits}

The sectors of section2 and sections3--5 are disjoint and exhaust SX1: either a transverse pairing is nonzero, or the square image lies in the original plane; in that second case either u\textless=22 or all23 active places are type II. The respective full charges are10531 and896. Their sum is strictly below2\^{}-895. Only this one combined895 charge is appended; the two internal charges are not also entered as separate historical errors.

The companion reads the actual55-entry all-top/31-entry uniform prior at \nolinkurl{checks/r059_joint_remainder_error.json}, preserves every entry and appends895. It checks the full56-entry all-top sum and unchanged uniform sum remain below2\^{}-230. This theorem does not change the actual role family, cutoff, selected module, arithmetic field, finite level, packet point, class or CM phase. The all-top support floor remains13 until the other structural sectors are separately exhausted.

The certificate verifies actual inputs, the original-entry norm arithmetic and exponents, the non-all-II support23 bounds, all164 endpoint factors, their strict brackets and union, and the complete historical error arrays. It does not automate the new paired B ideal, numerator-U argument, reconstruction, Hilbert trace or lattice packing. These are proved above, with accepted source interfaces explicitly retained. No new Lean build or abstract positive local module assertion is invoked.

No signed sufficient inequality, global prescribed period or Mram nonvanishing, literal Host constructor, DP44 cost, concrete cycle value or RH conclusion is asserted.

\clearpage

\section{Canonical checked integration --- AM-DO-083}\label{canonical-checked-integration-am-do-083}

2026-11-22T19:12:01.686140+00:00. Root and independent whole-proof/script reviews and separate683-check replays passed. RO1--18 remove all globally rank-one fixed-W0 terms at supports15..20 with one1953 charge. Full type/pair/double-pole counts and actual selected trace retained. This theorem alone does not exhaust any full support.

Chronology: append1953 after R060 and DO082, producing58 all-top and31 uniform charges, with exact sums below2\^{}-230. R061 follows separately with78919 and59/31 entries. Frozen author chronology below is preserved.

\subsection{Algebraic-geometric context}\label{algebraic-geometric-context}

The same original Hermitian modules, CM realizations and integral Weyl lattices remain. Carried contextual locators: SGA3 I4.7.3/4.7.3.1; SGA6 VI Theorem3.3; FGA232 Theorem3.1 pp148--149; FGA236 Corollaries1.7/2.3 pp229/232; SGA7 I VII2.7/2.9 pp169--174 and VIII7.4(b) pp310--311; cote4 scan003 §§8.1.1--8.1.2, scans005--006 §§8.4.1--8.4.3 and cote54 scans005--007. These are carried framework anchors, not new primary audits or new lambda/K-theory, l-adic weight, motivic comparison or L-function theorems. DP\_MULT Q3e and THETA3 G1/A7 retain the stated actual arithmetic gates; CX\_GRH CX-ALG-024 and GCT\_PLETH T-052 remain distinct.

\begin{center}\rule{0.5\linewidth}{0.5pt}\end{center}

\section{AM-DO-083 --- selected trace removes rank-one fixed-block supports fifteen through twenty}\label{am-do-083-selected-trace-removes-rank-one-fixed-block-supports-fifteen-through-twenty}

Author: ALP\_MACHINE \nolinkurl{/root/definite_oscillator}, 22 November 2026 UTC. New controlled pending proof. Author whole-file and independent verification records are separate.

\subsection{1. Actual statement and unchanged conditions}\label{actual-statement-and-unchanged-conditions}

Retain the original J56 all-top K044 upper endpoint, all allowed words and canonical twists, the K043 selected compact representation with its actual Hilbert metric, the original finite projectors, scalar-pencil incidences, central representatives, class masses and near-point transport. The all-top cutoff is H=133, after the already charged R048 annulus. The separate uniform H8192 family is unchanged. Write f=a=168, Q=g0g1, and use the original decomposition E\^{}169=S perpendicular W0, with signature(W0)=(4,1).

Consider precisely the original terms

\[
 h=I_S\oplus(I+N),\qquad \operatorname{rank}_E N=1,
 \qquad N^2=0,\qquad 15\le s\le20.
 \tag{RO1}
\]

Here s is the number of active relative T places. Global unipotence is a hypothesis in RO1; no criterion asserting it at these supports is introduced. The theorem is

\[
 \boxed{\quad
 \sum_{\substack{\text{actual RO1 terms}\\15\le s\le20}}
       |w(h)|<2^{-1953}.
 \quad}
 \tag{RO2}
\]

Every original incidence type and every ordered label pair is included. Only the intersection with the current remainder is removed, with one new charge1953 and all prior charges retained. The theorem does not by itself change the complete support floor or prove the signed lower bound, a period, Mram nonvanishing, a literal source family or a DP44 inequality.

\subsection{2. The genuine selected trace for rank one}\label{the-genuine-selected-trace-for-rank-one}

Let e be the negative last vector and put x=Ne and eta=x\_n.~The exact original-coordinate normal form K075(RN5--7) gives

\[
 \bar\eta=-\eta\ne0,\qquad
 N=x\bar x^{\mathsf t}J_0/\eta,\qquad
 \sum_{j=1}^4x_j\bar x_j=\eta\bar\eta,
 \qquad |\tau\eta|\le R:=\frac{133-133^{-1}}2=\frac{8844}{133}.
 \tag{RO3}
\]

This is the smaller H133 radius; the older K077 numerical table used H8192. One full column x determines N exactly, so there is no additional source-isometry fibre in this proof.

At a complex embedding only, write v=x/eta. For the full group calculation extend N by zero on S and v by zero in those coordinates, and use J=diag(I\_a,-1). Then v\_n=1, its positive part has norm one, and

\[
                      N=-\eta vv^*J.
 \tag{RO4}
\]

This analytic v is not claimed integral at any finite place. The actual transport g\_p has Cartan radius r\_p\textless1/2, as retained in K085 section2 and DO080 section7. Unitarity gives

\[
 \eta_p=(g_p^{-1}Ng_p)_{nn}
       =\eta\,|v^*Jg_pe_n|^2,
 \qquad
 |v^*Jg_pe_n|\ge\cosh r_p-\sinh r_p=e^{-r_p}.
 \tag{RO5}
\]

The inequality uses the norms of the positive and negative parts of g\_p e\_n and the actual unit positive part of v. Thus \textbar eta\_p\textbar\textgreater e\^{}\{-1\}\textbar eta\textbar\textgreater\textbar eta\textbar/3. No additive displacement estimate replaces this multiplicative relation.

Put h'=g\_p\^{}\{-1\}hg\_p and let r' be its Cartan radius. Its negative diagonal is 1+eta\_p, with eta\_p pure imaginary. The exact signature-one identity \textbar h'\_(nn)\textbar=cosh r' therefore gives

\[
                         \sinh r'=|\eta_p|>\frac{|\eta|}{3}.
 \tag{RO6}
\]

Apply DO080(RT2,RT24), which bounds the genuine normalized trace on the actual selected compact type, retaining both compact factors and every PBW/Pieri normalization. Multiplying the local estimates yields

\[
 |c(h)|<3^f2^{-931f/2}
                |\operatorname N_{E/\mathbf Q}(\eta)|^{-1/2}.
 \tag{RO7}
\]

The exponent1/2 converts the E norm to the product of absolute values over the f complex labels. In particular no pointwise archimedean lower bound for eta is inferred from its ideal. The root-number, finite-character and class factors in w(h) are unchanged.

\subsection{3. Retain the exact pencil, content and affine count}\label{retain-the-exact-pencil-content-and-affine-count}

We use the complete K077(AP1--16) local and arithmetic count, with only the radius R in RO3 substituted. Its proof holds for every positive cutoff. Let D be the intrinsic set of relative T places at which eta has a double pole, and write d=\textbar D\textbar. At a type-II place the fixed lattice is M\_c=(I+c pi\^{}(-1)Z)Lambda. An active term cannot fix c=0. Both the original and inverse-adjoint CM components have unit negative entries in the pole column and row. Thus double poles of any N entry are equivalent to a double pole of eta, and

\[
 x_j-(z_j/z_n)\eta\in\pi^{-1}\mathcal O_v
 \tag{RO8}
\]

at each double place, with the dual-component version separately retained. Types I and III have only simple entry poles. No asserted real-part projection or coarea formula is needed here, so the zero-target correction to those separate slice counts is not being bypassed.

Fix the subset of the334 unmixed G factors where N reduces to zero. Its paired E ideal J has norm D\_G and

\[
 D_G\le M^2,\qquad M=Q^{f-1}.
\]

Write I\_Sigma and I\_D for the paired inverse-prime ideals at the active places and double places. Every original N entry belongs to

\[
 \mathfrak l=\mathfrak J I_\Sigma I_D,
 \qquad \operatorname N_E\mathfrak l
       =D_GT^{-2(s+d)}=z^{2f}.
 \tag{RO9}
\]

The exact content b of x satisfies (eta) contained in b contained in l and

\[
 \operatorname N_E\mathfrak b\ge
 |\operatorname N_E\eta|^{1/2}z^f.
 \tag{RO10}
\]

For fixed eta and b the first three positive coordinates lie in specified affine cosets of b intersect J I\_Sigma. Their difference ideal has norm N\_E(b) T\^{}(2d), precisely K077(AP10--11). The fourth coordinate retains its relative-F norm equation with the original content b. K075's prime-norm floor29 and ideal-divisor exponent5/24 remain the same actual-field inputs; no class number or principality assertion is added.

Consequently, for each fixed G pattern, support, double set and actual ordered pencil-pair data, the complete number of possible columns and matrices is bounded by

\[
 28y^{105}(1+2y)^f(1+2\sqrt{y_c})^{6f},
 \qquad y=R/z,\quad y_c=R/(zT^{2d/f}).
 \tag{RO11}
\]

This includes all pure-imaginary eta values, content ideals and relative norm solutions. Each nonempty actual set has z\textless=R by the nonzero eta minimum. The paired ideal also gives the norm lower bound \textbar N\_E eta\textbar\textgreater=z\^{}(2f), so RO7 implies

\[
                         |c(h)|<3^f2^{-931f/2}z^{-f}.
 \tag{RO12}
\]

Only a norm statement is used in RO12.

\subsection{4. The weighted cap and all incidence types}\label{the-weighted-cap-and-all-incidence-types}

The unchanged rank-one Green estimate is

\[
                |a(h)|\le2D_G^{a/2}M^{-a}.
 \tag{RO13}
\]

There are at most2\^{}334 paired rank-zero patterns. Rank one has no extra CM orientation choice. The two mixed factors keep their norm bound one, and the single factor two in RO13 is the global Euler correction.

For fixed s,d, multiply RO11, RO12 and RO13 before applying the norm cap. Apart from positive constants, the z dependence is

\[
 z^{27279}(z+2R)^f
          (\sqrt z+2\sqrt R\,T^{-d/f})^{6f},
 \qquad
 27279=(a-1)f-105-4f>0.
 \tag{RO14}
\]

It is increasing on the entire positive axis. Hence D\_G\textless=M\^{}2 permits the substitution

\[
 z_{s,d}=M^{1/f}T^{-(s+d)/f},\qquad
 y_{s,d}=R/z_{s,d},\qquad
 y'_{s,d}=R M^{-1/f}T^{(s-d)/f}.
 \tag{RO15}
\]

The finite factor M\^{}(-a) T\^{}(a(s+d)) z\_(s,d)\^{}(af) cancels exactly. If this cap exceeds R, the increasing algebraic expression still bounds the actual z\textless=R part; no lattice point at the cap is asserted.

The double places must be type II and have full weighted ordered-pair cost1. At each of the other places all three types are permitted, with costs T−1,1,1. Thus the complete choice/weight sum is binom(f,s) binom(s,d)(T+1)\^{}(s-d), exactly as in K077(AP18). It includes possible simple-pole type-II places and actual zero labels in the other types. No additional type count, finite fibre, pivot or representative is inserted.

Define

\[
 \begin{split}
 C_s&=2^{335}\,28\binom fs\,3^f2^{-931f/2},\\
 B_s(d)&=(T+1)^{s-d}z_{s,d}^{-f}y_{s,d}^{105}
          (1+2y_{s,d})^f(1+2\sqrt{y'_{s,d}})^{6f}.
 \end{split}
 \tag{RO16}
\]

The complete support-s contribution is at most C\_s sum\_d binom(s,d) B\_s(d). As a real function on0\textless=d\textless=s, log B\_s(d) is a sum of affine functions and positive multiples of log(1+exp(A+Bd)); it is convex. Therefore

\[
 \boxed{\quad
 \mathscr R_{1,s}\le2^sC_s\max\{B_s(0),B_s(s)\}.
 \quad}
 \tag{RO17}
\]

This is one maximum, with no endpoint union factor. The all-double y'\_(s,s)=R M\^{}(-1/f) remains independent of s. Both endpoints are retained.

\subsection{5. Exact endpoints and the new union}\label{exact-endpoints-and-the-new-union}

The certificate uses only the actual host integers, rational brackets for the f-th roots and square roots in RO15, and directed integer logarithms. The rigorous inverse binary margins of the complete RO17 bound are

\begin{longtable}[]{@{}
  >{\raggedleft\arraybackslash}p{(\columnwidth - 6\tabcolsep) * \real{0.2500}}
  >{\raggedleft\arraybackslash}p{(\columnwidth - 6\tabcolsep) * \real{0.2500}}
  >{\raggedleft\arraybackslash}p{(\columnwidth - 6\tabcolsep) * \real{0.2500}}
  >{\raggedleft\arraybackslash}p{(\columnwidth - 6\tabcolsep) * \real{0.2500}}@{}}
\toprule\noalign{}
\begin{minipage}[b]{\linewidth}\raggedleft
Support s
\end{minipage} & \begin{minipage}[b]{\linewidth}\raggedleft
No-double endpoint
\end{minipage} & \begin{minipage}[b]{\linewidth}\raggedleft
All-double endpoint
\end{minipage} & \begin{minipage}[b]{\linewidth}\raggedleft
Complete support margin
\end{minipage} \\
\midrule\noalign{}
\endhead
\bottomrule\noalign{}
\endlastfoot
15 & 38724 & 28107 & 28107 \\
16 & 32125 & 22877 & 22877 \\
17 & 25526 & 17647 & 17647 \\
18 & 18927 & 12416 & 12416 \\
19 & 12328 & 7186 & 7186 \\
20 & 5729 & 1956 & 1956 \\
\end{longtable}

Each entry means the complete corresponding expression is strictly less than2 to the negative listed integer. The certificate also places each logarithm between that integer boundary and the next one; these are exact directional comparisons, not rounded floating-point conclusions. At s21 both endpoint upper expressions exceed one. That is a failure of this sufficient bound, not a claim about the sign or nonvanishing of the actual sector.

There are six supports, and every type and label has already been summed in each support bound. Hence

\[
 \sum_{s=15}^{20}\mathscr R_{1,s}
 <6\,2^{-1956}<2^{-1953}.
 \tag{RO18}
\]

Earlier K077 supports10..14 and all their charges remain unchanged. No global-unipotence inference is needed to apply RO18 to its stated conditional sector.

\subsection{6. Verification, chronology and limits}\label{verification-chronology-and-limits}

The companion script rechecks the actual T and Q, the smaller radius, the positive weighted-cap exponent, all endpoint root/log brackets, the111 marked-double counts, and exact log-convex models for the full binomial envelope. Rational complex models check the rank-one normal form, conjugated negative entry and the exact identity relating its modulus to the Cartan radius. These checks supplement the written proof; they do not enumerate rational automorphic terms or mechanically prove the source identities.

The complete new input is the application RO4--7 of the actual DO080 trace to rank one and its weighted use RO12--17 in the unchanged K077 count. K075(RN5--20) and K077(AP1--23) were directly reread, as were DO080(RT2,RT24,section7) and its actual model scope. The original primary model anchor remains Nakahama (2016), equation(5.7) and Theorem5.1, through the checked DO080 Hilbert realization; no new external representation theorem is invoked.

Append1953 once to the parent's chronological all-top list when this finding is integrated. The accompanying check retains the actual R05955-entry prior and the R060172255 and DO082895 additions, followed by this new charge, giving58 all-top entries and the unchanged31 uniform entries; both complete dyadic sums remain below2\^{}(-230). This explicit ordered check does not claim a canonical append occurred before the parent actually performed it. At every removal, take the intersection with the then-current remainder and retain all older numerical errors.

The result removes the full rank-one fixed-W0 sector at supports15..20. Higher full ranks, all remaining structural classes and the final signed sufficient inequality keep their own conditions. In particular this is not an assertion that all support15 terms have disappeared. No selected arithmetic-image content, actual period value, literal Host/FamP record, DP44 sign or RH statement follows.

\clearpage

\section{AM-R-061 --- the entire scalar-cross class at support twenty-two}\label{am-r-061-the-entire-scalar-cross-class-at-support-twenty-two}

Author: GPT-6 Astra Pro; inherited effort not externally exposed; ALP\_MACHINE \nolinkurl{/root}. 22 November2026 UTC. New proof; author and independent whole checks are recorded separately before canonical promotion.

\subsection{1. Actual scope and complete estimate}\label{actual-scope-and-complete-estimate}

Keep the unchanged J56 all-top kernel, its charged H=133 cutoff, f=168, n=169 and m=164. All original CM incidences, finite projectors, packet point, compact twists, Hilbert metrics, class masses and phases remain. As in R059(GB1--10), write the original blocks

\[
 h=\begin{pmatrix}I_m&B\\ C&D\end{pmatrix},\qquad C\ne0,
 \qquad S\perp W_0=E^{169},\quad\operatorname{sig}(W_0)=(4,1).
 \tag{GS1}
\]

We prove, for the \textbf{entire} such class at support s=22, without any global unipotence assumption,

\[
 \sum_{h\text{ in GS1},\ s=22}
       |\text{actual normalized contribution of }h|<2^{-78919}.
 \tag{GS2}
\]

This is a fresh charge for the intersection with the current remainder. Every earlier charge is retained, including those for smaller globally unipotent subsets. It neither exhausts other structural classes at support22 nor establishes the sufficient signed inequality or an arithmetic period.

Use the exact original Q,U,u and three positive transverse vectors y\_j of R059(GB4), with

\[
 Q=1427908193729121522601,\qquad U=77Q^2-14,
 \qquad u^*u=U,\quad\gcd(u_1,u_2,u_3,u_4)=1.
 \tag{GS3}
\]

R059 proves rank(C)=1 and gives an exhaustive dichotomy: at least one nonzero integral row delta\_i=\textless y\_j,Ce\_i\textgreater, or all those rows zero, in which case

\[
                 Ce_i=(u\xi_i,\beta_i),\qquad\xi\ne0.
 \tag{GS4}
\]

The paired G orientation is for \textbf{every entry of C}, not only one row. Its ideal J has at most2\^{}668 possible patterns and the Green envelope is

\[
 |a(h)|\le2Q^{-335(f-1)}(N_EJ)^{335/2}.
 \tag{GS5}
\]

The source/target pairing proof and this finite-unipotent reduction argument are exactly R059(GB3--10), which impose no global nilpotence. We cite them in that scope.

\subsection{2. One original scalar replaces the extra five-coordinate pivot}\label{one-original-scalar-replaces-the-extra-five-coordinate-pivot}

Consider the plane case GS4. Put L=J I\_Sigma, where N\_E I\_Sigma=T\^{}(-2s), and

\[
                   z=(N_E L)^{1/(2f)}.
 \tag{GS6}
\]

The same original-entry divisibilities which gave xi in L\^{}m also give every beta\_i in L. At G this follows from the whole-C orientation. At T, both e\_i and h e\_i belong to the actual target pencil, so Ce\_i belongs to that pencil; its original coordinates have at most a simple pole at each CM partner. Away G,T they are integral. No divided moving coordinate or additional denominator is used.

Fix the first nonzero coordinate xi\_j0. The data consisting of the entire xi row and the single original scalar beta\_j0 determine C:

\[
       Ce_i=\left(u\xi_i,
                    \beta_{j_0}\xi_i/\xi_{j_0}\right).
 \tag{GS7}
\]

This follows from rank(C)=1. The quotient is used only for injective reconstruction over E, not for an integrality claim. Counting at most164 possible pivot positions is harmless. The scalar beta\_j0 is counted in the \textbf{same} L; it does not incur an additional independent T-index factor.

At each complex embedding, h\textbar\_S has Euclidean block column (I,C)\^{}t, so \textbar\textbar C\textbar\textbar\textless H. Isotropy gives \textbar beta\_i\textbar²=U\textbar xi\_i\textbar². Consequently

\[
       \|\xi_\tau\|<M:=H/\sqrt U,
                  \qquad|\beta_{j_0,\tau}|<H.
 \tag{GS8}
\]

The first inequality is a bound on the whole m-dimensional row. In fact \textbar\textbar C\textbar\textbar=sqrt(2U)\textbar\textbar xi\textbar\textbar, but the weaker GS8 is sufficient. If exactly k coordinates of xi are nonzero, 1\textless=k\textless=164, each has half-trace length at least sqrt(f)z. Summing their squared lengths and using GS8 yields

\[
                           kz^2<M^2.
 \tag{GS9}
\]

The ideal packing lemma R059(GB14), separately for each coordinate, gives the complete upper count

\[
 \binom{164}{k}(1+2M/z)^{2fk}(1+2H/z)^{2f}
 \tag{GS10}
\]

for xi and its one pivot scalar. The binomial chooses nonzero positions; admitting zero in an individual packing ball only enlarges its count. GS9 is used only for actually nonzero coordinate sectors. No local ideal basis or lattice equidistribution is assumed.

\subsection{3. The genuine selected trace after the actual packet transport}\label{the-genuine-selected-trace-after-the-actual-packet-transport}

The actual LIFT41 packet representative is gamma=I+b b*J/(1-bar(c)), NP7, where b=z\_packet-e\_n has support in W0. Thus gamma fixes S pointwise and preserves W0. We use this literal representative as g\_p; the Cartan radius r\_p is less than1/2 by NP12 and NP19. There is no determinant correction or change of packet here. If the same point is represented by g\_p k with k in the compact stabilizer, the normalized W-trace is unchanged by this compact conjugation because W is compact-stable; the following bound therefore applies to that representative too.

At an embedding the columns of C span an isotropic line. Write a nonzero spanning vector as ell=(v,1), with \textbar\textbar v\textbar\textbar=1; this normalization is only over C. For the actual g\_p, the negative coordinate of g\_p\^{}(-1)ell has modulus at least exp(-r\_p)\textgreater1/2. Indeed it is the Hermitian pairing, up to sign, of ell with g\_p e\_n, whose positive and negative part lengths are sinh r\_p and cosh r\_p.~Their difference is exp(-r\_p).

Because g\_p fixes S, the W0-by-S block of h'=g\_p\^{}(-1)h g\_p is g\_p\^{}(-1)C. Its negative row on S therefore has norm greater than half the norm of the original beta row. Every U(168,1) matrix with Cartan radius r' has total negative-row positive-part norm sinh r', by its unitary row identity. Hence

\[
              \sinh r'>\tfrac12\sqrt U\,\|\xi_\tau\|.
 \tag{GS11}
\]

This proof uses unitarity and the fixed S block, and is valid for globally nonunipotent matrices as well.

DO080(RT2,RT24) bounds the \textbf{actual selected normalized trace}, with both compact factors retained, by 2\^{}(-931/2)/sinh r'. For any nonzero original coordinate xi\_j in L,

\[
 \prod_\tau\|\xi_\tau\|
       \ge\prod_\tau|\tau\xi_j|
       =|N_E\xi_j|^{1/2}\ge z^f.
\]

Multiplication at the f labels now gives

\[
               |c(h)|<2^{-929f/2}U^{-f/2}z^{-f}.
 \tag{GS12}
\]

This is a product-norm inequality, not a lower bound on every individual coordinate at every embedding. It does not import DO080's separately stated global-rank2-unipotent consequence RT3: GS11 supplies the different, fully proved Cartan comparison needed here.

\subsection{4. The source fibre has positive rank five after C is fixed}\label{the-source-fibre-has-positive-rank-five-after-c-is-fixed}

For fixed C and v=D e\_n, unitarity reconstructs the full negative column by

\[
                         (he_n)_S=-C^*J_0v,
 \tag{GS13}
\]

as in R059(GB8). Every matrix with this C fixes the same rational positive space K=ker(C\textbar\_S) pointwise: its S block is I and its cross block vanishes on K. Since C has rank one, dim\_E K=163.

Fix one matrix h0 in a nonempty fibre of the data (C,v) and the actual ordered source/target pencils. Every h in that fibre gives h0\^{}(-1)h preserving the same original source lattice, fixing e\_n and fixing K pointwise. The orthogonal complement of K direct-sum E e\_n is positive definite of E-rank5 at every real label. Restriction is faithful on this complement. Its intersection with the source lattice is a full projective O\_E-lattice, so its isometry group is finite.

At either of the actual completely split G primes this intersection is a free rank5 lattice over Z\_g. The local reduction and residue-characteristic argument of DO073(RF9--11) applies unchanged: prime-to-g torsion is excluded by the cited Serre lemma, and order-g torsion cannot occur because g-1\textgreater5 and the primitive g-th-root minimal polynomial has degree g-1. Thus the finite group injects at both primes and its order is at most

\[
 G_5=\gcd(|GL_5(\mathbf F_{22285932761})|,
           |GL_5(\mathbf F_{64072175441})|)
       =98148490017978252769652736000000<2^{107}.
 \tag{GS14}
\]

This derives the fixed space from the known C, so none of DO073's global-unipotent hypotheses used to recover a fixed space from only a negative column is assumed. The faithful local lattice argument itself is reused in exactly its positive-rank5 scope. No global integral orthogonal basis or coefficient-height bound is claimed.

The remaining original W0 column v lies in the true pi\^{}(-1) M\_(d,W0). Its half-trace norm is less than95 sqrt(f): for H=133 the squared Euclidean norm of h e\_n is at most(H²+H\^{}(-2))/2\textless95². This original target overlattice has paired index T\^{}(10s) and covolume D0\^{}5 T\^{}(-10s), where

\[
                    D_0=7^{140}29^{144}5^{84}13^{84}.
\]

The LIFT063 original-ring covering bound is less than128 sqrt(f), uniformly in every target label including zero. The actual volume argument of DO075(CR10--13) and R060(J15.6--10) therefore gives

\[
 \#\{v\}\le V_5(223)T^{10s}<2^{5757}T^{10s},
 \quad V_5(223)=\frac{(528\cdot223^2)^{5f}}{(5f)!D_0^5}.
 \tag{GS15}
\]

Only a full negative-column ball is used, with no affine slice or zero-target coarea hypothesis. The covering bound follows also by containing the original ring lattice; it is not a cover estimate for an arbitrary fractional ideal from its norm alone.

\subsection{5. Weighted ideal endpoint and complete rational exponent}\label{weighted-ideal-endpoint-and-complete-rational-exponent}

Multiplying GS10 by GS5 and GS12 gives, apart from the explicit constants and T\^{}(335s), the z-dependent expression

\[
 z^{(332-2k)f}(z+2M)^{2fk}(z+2H)^{2f}.
 \tag{GS16}
\]

Since k\textless=164, its leading exponent is at least4f; all its factors increase for z\textgreater0. Use GS9 to evaluate it at M/sqrt(k). Together with the external U\^{}(-f/2), the resulting factor is

\[
 H^{334f}U^{-335f/2}\binom{164}{k}k^{-167f}
        (1+2\sqrt k)^{2fk}(1+2\sqrt{Uk})^{2f}.
 \tag{GS17}
\]

Now 1+2sqrt(Uk)\textless=3sqrt(Uk), since U,k\textgreater=1. Thus GS17 is bounded by

\[
 3^{2f}H^{334f}U^{-333f/2}
 \binom{164}{k}k^{-166f}(1+2\sqrt k)^{2fk}.
 \tag{GS18}
\]

Retain once each: the single Green factor2; patterns2\^{}668; support choices2\^{}168; pivot choices164\textless2\^{}8; fibre2\^{}107; and the column constant2\^{}5757. The exact complete source/target pair and weight sum is R059(GB19), P\_T\^{}s\textless2\^{}(997s), including all three incidence types. There is no further3\^{}s or profile count. The remaining T power is335s+10s=345s. Put C0=1+668+168+8+107+5757=6709.

Use the exact integer-certified inequalities

\[
 \log_2T<L_T=1019333/1024,\quad
 \log_2Q>L_Q=287843/4096,\quad
 \log_2H<L_H=28900/4096,\quad
 \log_2U>L_U=601356/4096,\quad3^{336}<2^{533}.
 \tag{GS19}
\]

For s=22 the common rational exponent is

\[
 \begin{split}
 C_*={}&C_0+533-929f/2+334fL_H-(333f/2)L_U\\
       &+(345L_T+997)22-335\cdot167L_Q
        =-556050547/4096.
 \end{split}
 \tag{GS20}
\]

For each1\textless=k\textless=164 put c\_k=ceil(sqrt(k)), and let e\_k be the least integer such that

\[
 \binom{164}{k}(1+2c_k)^{2fk}<2^{e_k}k^{166f}.
 \tag{GS21}
\]

This is an exact positive-integer comparison, also when e\_k is negative. The companion certificate enumerates all164 k values, checks both the strict upper inequality and failure at e\_k-1, and verifies

\[
 \max_k e_k=56826\text{ at }k=164,\quad
 C_*+56826=-323291251/4096<-78928.
 \tag{GS22}
\]

The union of these164 actual coordinate sectors is therefore below2\^{}(-78920).

The other, transverse case has the support-independent full R059(GB21) count. Substituting s=22, with its original two complete W0-column counts and full source fibre retained, gives

\[
 E_A(22)=35591+43\cdot335f+(8/5)328f
                +20\cdot22L_T+997\cdot22-335\cdot167L_Q
             =-19000573703/20480<-927762.
 \tag{GS23}
\]

We apply that actual count freshly at22; no already charged lower-support theorem is silently extended. Combining the transverse case with the plane bound gives GS2. The two internal bounds are not separately appended to the accumulated error list.

\subsection{6. Accounting, geometric meaning and verification boundary}\label{accounting-geometric-meaning-and-verification-boundary}

The baseline at authoring is DO082's57 all-top and31 uniform errors, with support floors15 and11. Append only78919 to the all-top list. The companion certificate checks the complete58-entry and unchanged31-entry sums below2\^{}(-230); any intervening independent charge must also be retained in the eventual canonical chronology. The historical baseline and final live check are kept separately.

The new fixed space is the kernel of an actual rank-one map between the original Hermitian modules over Spec(O\_E). Its positive complement is used with its intersection lattice, not a new integral frame. The same ideal controls all original cross entries. The selected trace is bounded in the unchanged complex realization of those modules. No CM type, Picard class, polarization, biextension, source determinant line or finite projector is changed.

The carried contextual locators are SGA3 I4.7.3/4.7.3.1; SGA6 VI Theorem3.3; FGA232 Theorem3.1 pp148--149; FGA236 Corollaries1.7/2.3 pp229/232; SGA7 I VII2.7/2.9 pp169--174 and VIII7.4(b) pp310--311; cote4 scan003 §§8.1.1--8.1.2, scans005--006 §§8.4.1--8.4.3, and cote54 scans005--007. These provide the previously recorded framework, not a new primary-source audit or an attribution of this estimate. No lambda-ring/K-theory identity, l-adic weight theorem or motivic L-function conclusion is added.

DP\_MULT Q3e and THETA3 G1/A7 have a smaller actual finite gate. The final signed sufficient inequality, full admissible mixed-CM source, selected-image arithmetic content, prescribed Mram period, concrete032 cycle/value and DP44 sign remain open. CX\_GRH CX-ALG-024 and GCT\_PLETH T-052 retain their separate gaps. The exact arithmetic certificate supplements the full general lattice, trace and fibre proof; it is not a rational-matrix enumeration or a Lean build.

\clearpage

\section{Canonical checked integration --- AM-K091}\label{canonical-checked-integration-am-k091}

2026-11-22T19:12:01.686140+00:00. Root and independent whole-proof/script reviews and separate800-check replays passed. SR1--13 give an integral left inverse for the literal selected Weyl substitution and its full symbolic two-edge composition. No retraction of the actual derivative evaluation, geometric class sum or terminal operations is asserted.

\subsection{Algebraic-geometric context}\label{algebraic-geometric-context-1}

The same original Hermitian modules, CM realizations and integral Weyl lattices remain. Carried contextual locators: SGA3 I4.7.3/4.7.3.1; SGA6 VI Theorem3.3; FGA232 Theorem3.1 pp148--149; FGA236 Corollaries1.7/2.3 pp229/232; SGA7 I VII2.7/2.9 pp169--174 and VIII7.4(b) pp310--311; cote4 scan003 §§8.1.1--8.1.2, scans005--006 §§8.4.1--8.4.3 and cote54 scans005--007. These are carried framework anchors, not new primary audits or new lambda/K-theory, l-adic weight, motivic comparison or L-function theorems. DP\_MULT Q3e and THETA3 G1/A7 retain the stated actual arithmetic gates; CX\_GRH CX-ALG-024 and GCT\_PLETH T-052 remain distinct.

\begin{center}\rule{0.5\linewidth}{0.5pt}\end{center}

\section{AM-K091 --- an integral retraction for the literal selected substitution}\label{am-k091-an-integral-retraction-for-the-literal-selected-substitution}

Author: ALP\_MACHINE \nolinkurl{/root/kummer_extension}, 22 November 2026 UTC.

\textbf{Author whole-proof and complete-script PASS; frozen for independent review.} The statement concerns the exact symbolic subD maps before evaluation by the prescribed theta derivative functional. It does not assert a retraction for that functional, a class sum or terminal branching.

\subsection{1. Literal polynomial modules and the question being settled}\label{literal-polynomial-modules-and-the-question-being-settled}

Keep the two actual mixed-host edges, their nonuniform selected partitions and determinant twists from K037(CR2--3), and the source's integral Weyl lattice. The assertion below applies label by label, so neither the individual partitions nor their degrees must be uniform across embeddings.

For integers \nolinkurl{a<=b}, a partition eta of length a and \nolinkurl{q>=0}, write \[
 \nu=\eta+q\mathbf1_a,\qquad\widetilde\nu=(\nu,0^{b-a}).
 \tag{SR1}
\] Let \nolinkurl{Delta_eta(a,R)} denote the source polynomial Weyl module over a commutative ring R. Its highest polynomial is \[
 p_\eta(U)=\prod_{j=1}^{a}\det U_{[1,j],[1,j]}^{\eta_j-\eta_{j+1}},
 \qquad\eta_{a+1}=0.
 \tag{SR2}
\] The module is the R-span of all finite words of the source row divided powers applied to SR2. This is exactly \nolinkurl{Formalism/GradedWeylDivisibility/StatementA.lean}31--44,63--73: rowDiv is a coefficient of an elementary row substitution, including its diagonal case. No saturation in the ambient polynomial ring is part of this definition.

The literal selected substitution of K037(CR7) is \[
 \mathcal Q_{a,b,\eta,q}(P)(Y,Z)
   =\left[\det(U)^qP(U)\right]_{U=Z^{\mathsf T}Y_{[,1:a]}},
 \tag{SR3}
\] where Y is b-by-b and Z is b-by-a. K037(CR8--11) proves that this belongs to \nolinkurl{Delta_tilde_nu(b,R)[Z]}. In particular the coefficients belong to the actual target Weyl lattice, not just its characteristic-zero span or dual-Weyl enlargement.

We prove that SR3 is split injective \textbf{in precisely these Weyl modules}, over every commutative coefficient ring. Consequently its symbolic coefficient matrix has an integral left inverse. This strengthens the earlier membership result but does not change the actual input Schwartz function or select a different theta functional.

\subsection{2. Leading-block restriction preserves the integral Weyl lattice}\label{leading-block-restriction-preserves-the-integral-weyl-lattice}

Define the polynomial restriction \[
 \operatorname{res}_{b,a}(F)(U)
       =F\left(\begin{pmatrix}U&0\\0&0\end{pmatrix}\right).
 \tag{SR4}
\] Every target highest polynomial and every row-divided-power word uses only the first a columns because the remaining column weights are zero. The zero blocks in SR4 are therefore harmless literal polynomial substitutions, not limits of invertible frames.

\textbf{Lemma.} Restriction maps \[
 \operatorname{res}_{b,a}:\Delta_{\widetilde\nu}(b,R)
                   \longrightarrow\Delta_\nu(a,R).
 \tag{SR5}
\] It is a retraction of the inclusion that places a polynomial in the first a rows and columns.

\textbf{Proof.} Work first over Z. A divided-power word is an iterated coefficient of a finite sequence of elementary row substitutions. Give each substitution its own variable. Their product is a b-by-b polynomial matrix \nolinkurl{G(t)}, with no invertibility assumption needed after specialization. Thus a word applied to the target highest polynomial is a specified coefficient of \nolinkurl{rowAct(G(t)) p_tilde_nu(Y)}.

After SR4, only the first a rows of \nolinkurl{G(t)Y} used by the highest polynomial and the first a columns of G multiplying U remain. If \nolinkurl{B(t)=G(t)_[1:a,1:a]}, the restricted generating polynomial is exactly \[
                    p_\nu(B(t)U).
 \tag{SR6}
\] K037(CR9--10) proved integral stability of the Weyl lattice under every row matrix over every commutative coefficient algebra. Apply it in rank a to \nolinkurl{B(t)} over \nolinkurl{Z[t]}. SR6 lies in \nolinkurl{Delta_nu(a,Z)[t]}; hence every one of its t-coefficients is in \nolinkurl{Delta_nu(a,Z)}. This proves SR5 on every defining word, and then on their integral span. It also addresses words which move rows outside the leading block and subsequently move them back: no stepwise commutation with restriction was assumed.

Conversely, every defining divided-power word in the first a rows embeds as the same word in b rows, and the highest polynomial embeds exactly. Thus inclusion preserves the respective Weyl lattices and its composite with SR4 is the identity. The source generators are integral polynomial formulas; base change sends their integral spans to the defining modules over R. Both assertions therefore hold over every R. QED.

The proof uses integral row stability, not a claim that the integral Weyl lattice equals the intersection of its rational span with all polynomials. That latter claim would be false in relevant weights.

\subsection{3. The determinant twist is an exact lattice isomorphism}\label{the-determinant-twist-is-an-exact-lattice-isomorphism}

Multiplication by the determinant gives \[
 M_q:\Delta_\eta(a,R)\xrightarrow{\ \sim\ }
                 \Delta_\nu(a,R),\qquad P\longmapsto\det(U)^qP.
 \tag{SR7}
\] K037(CR11) already proves the inclusion from left to right. For completeness, equality and the inverse are justified without dividing an integer factorial.

Off-diagonal elementary row substitutions leave the determinant invariant. Hence every off-diagonal row divided power commutes with multiplication by \nolinkurl{det(U)^q}. A diagonal divided power of order j on a row-homogeneous polynomial of row degree d is multiplication by the integer \nolinkurl{binom(d,j)}. A divided-power word applied to a highest polynomial is row homogeneous or zero after every step. Remove its diagonal operations successively, retaining their actual integer multipliers. Its remaining off-diagonal word on \nolinkurl{p_nu=det(U)^q p_eta} equals \nolinkurl{det(U)^q} times the same word on \nolinkurl{p_eta}. This proves the reverse inclusion in SR7 over Z; base change gives it over R. Different binomial factors for eta and nu are not equated or divided.

Multiplication by \nolinkurl{det(U)^q} is injective in the full polynomial ring over any R. In a monomial order its leading coefficient is a unit, so a nonzero polynomial's leading term cannot be killed even if R has zero divisors. Its inverse on the exact image SR7 is therefore a well-defined R-linear map \[
                   D_q:\Delta_\nu(a,R)\longrightarrow\Delta_\eta(a,R).
 \tag{SR8}
\] This notation means division by a known exact polynomial factor, not evaluation at a matrix of nonunit determinant or inversion of an integer scalar.

\subsection{4. The explicit symbolic retraction and two-edge composition}\label{the-explicit-symbolic-retraction-and-two-edge-composition}

Let \nolinkurl{iota_(b,a)} be the b-by-a matrix with leading a-by-a identity and remaining rows zero. Evaluate Z at this integral matrix and define \[
 \mathcal R_{a,b,\eta,q}
       =D_q\circ\operatorname{res}_{b,a}\circ
                      \operatorname{ev}_{Z=\iota_{b,a}}.
 \tag{SR9}
\] All three maps are R-linear between the stated Weyl modules; evaluation of the Z coefficients uses only zero and one. Direct substitution in SR3 gives \[
 \operatorname{res}_{b,a}\mathcal Q(P)(Y,\iota)=\det(U)^qP(U),
 \qquad\boxed{\mathcal R\mathcal Q=\operatorname{id}.}
 \tag{SR10}
\] Thus Q is split injective. The argument is independent of the characteristic of R, of whether any Schur hook is an R-unit, or of whether the corresponding Weyl module is saturated in its polynomial ambient module. For the finite free integral lattices over Z, or their flat base changes to the actual DVR, SR10 is the integral matrix identity \nolinkurl{RQ=I}; in particular it gives no extra selected-polynomial Smith factor at11.

There are two distinct reductions when a Weyl lattice is not saturated. The abstract tensor reduction \nolinkurl{Delta(Z) tensor F_11} need not inject into the ambient polynomial ring over F\_11, and it must not be identified with the polynomial image generated there by reduced words. The integral matrix identity SR10 remains a split identity on the abstract tensor reductions. Separately, the polynomial-module constructions over an arbitrary R have the retraction proved above on their defining images. The actual content statement uses the former lattice reduction. For example, \nolinkurl{11 x^10 y} is a primitive weight-one generator of the rank-two, one-column Weyl lattice of degree11 although its ambient polynomial coefficients are divisible by11. No step here declares its lattice reduction zero.

For the actual two edges, extend the second substitution and its retraction over the coefficient algebra \nolinkurl{R[Z0]}. The symbolic coupled map is exactly \[
 \mathcal Q_{\rm cpl}
   =(\mathcal Q_1\otimes1_{R[Z0]})\mathcal Q_0,
 \qquad
 \mathcal R_{\rm cpl}
   =\mathcal R_0(\mathcal R_1\otimes1_{R[Z0]}),
 \qquad\mathcal R_{\rm cpl}\mathcal Q_{\rm cpl}=1.
 \tag{SR11}
\] Here the actual middle shape is the target shape of the first edge, exactly as in K037 and K038. Tensor the construction over all CM labels. The q values, partitions and degrees may differ from label to label. Complete middle coevaluation is simply composition of these actual polynomial maps; its expression in a perfect Weyl/dual-Weyl pair of bases adds no scalar. This is the same integral composition used in LIFT34, not an inverse of the Fischer form in two Weyl bases.

SR9 is a proof of a retraction. It is \textbf{not} an admissible replacement of the given Schwartz input or of its derivative functional by evaluation at the chosen constant Z matrix. No norm bound for R, integral inverse for an evaluated theta map, or terminal branching retraction is asserted.

\subsection{5. Consequence for K090's fixed-triple auxiliary orbit}\label{consequence-for-k090s-fixed-triple-auxiliary-orbit}

Fix the same coarse initial class, middle class and terminal decoration as K090(MA13--17), with their unchanged underlying lattices and standard11 frames. The actual finite functions \nolinkurl{Psi_s} have the common unit evaluation matrix \nolinkurl{Psi_s(x_t)=u delta_(s,t)} of MA14. For an arbitrary family of coefficient vectors \nolinkurl{P_s} in the common literal initial Weyl lattice, consider the \textbf{symbolic} map \[
 (P_s)_s\longmapsto
       \sum_s\Psi_s\otimes\mathcal Q_{\rm cpl}(P_s).
 \tag{SR12}
\] Evaluating the finite factor at each x\_t, dividing by the same unit u, and applying SR11 gives an integral left inverse. Restrict to the tensor augmentation module defining actual Mram; K090(MA5--6) makes it an integral direct summand. Therefore the restricted symbolic selected map is also split.

Over the proof-only unramified splitting extension used by K089, apply its \textbf{actual trace-one-normalized} theta isomorphism to the finite factor. The same conclusion holds for the coefficient-labelled auxiliary sections with their selected differential symbols. Faithful-flat descent gives a torsion-free cokernel over the original DVR, exactly as in K090. This changes neither the coefficient field of the input nor the level, tuple, full-Y sum or actual period normalization.

Consequently the content of a primitive input is still primitive after both literal symbolic selected substitutions and the complete symbolic middle composition for this fixed class triple. This removes a possible source of content loss between MA17's external coefficient tensor and its selected polynomial symbols. It does not use an ambient polynomial unit coefficient as a proxy for a Weyl-lattice unit.

\subsection{6. The exact boundary of this conclusion}\label{the-exact-boundary-of-this-conclusion}

The actual selected kernel is obtained by applying the prescribed functional \[
 J(P)=\epsilon\,[P(cD_\tau\partial_\tau)F]_0
 \tag{SR13}
\] to these symbols, with the original ordinary derivatives and Chern gauge. This is K037(CR4,CR12), followed on the coupled map by K038(CA23,CA27). SR9 does not commute with, or invert, J. The functional may annihilate the split symbolic image or increase its content. Its actual selected multiindices are not replaced by a complete Hasse-jet space. The retained factor \nolinkurl{U(t)^kappa exp(-kappa beta(t))}, the source-creation and period factors, and every ordinary/divided-power distinction remain exactly as before.

Likewise the fixed-triple assertion SR12 is not a direct-sum assertion for the geometric sections from all middle classes. Their actual weighted sum can cancel after evaluation. The terminal brC and genuine/generic projector operations retain their own specified scalars and denominators. No integral or analytic right inverse of those operations is constructed here.

Thus K090(MA10--11) remains the exact selected arithmetic matrix and content ideal to study. The remaining possible losses are its prescribed derivative evaluation, actual class contractions and terminal operations; a separate loss from the universal selected Weyl substitution itself is excluded by SR10--12. No nonzero evaluated column, sharp actual-image valuation, period, full admissible source family or DP44 sign is asserted.

\subsection{7. Sources and verification}\label{sources-and-verification-4}

The directly read source definition is \nolinkurl{Formalism/GradedWeylDivisibility/StatementA.lean}31--44,63--73. The integral row-matrix stability is the accepted K037(CR9--10), itself based on the checked WeylGLStable interface, with its monic determinant localization proof. K037(CR2--3,CR7--13), K038(CA22--27), and K090(MA13--18) supply the actual weights, substitution, evaluated functional and restricted orbit. No new external representation-theoretic theorem is invoked.

The companion certificate tests explicit row-divided-power and leading-block identities, exact determinant-twist commutation, one- and two-edge retractions over integers and residue rings, and a nonsaturated source-Weyl stress example. It includes nonzero words which leave and re-enter the leading block, so the proof is not being replaced by an invalid stepwise restriction argument. These finite tests do not replace the general proofs above and do not calculate an actual theta derivative or class matrix.

\subsection{Author whole verification}\label{author-whole-verification}

At 2026-11-22T18:57:31.251520+00:00, I completed a whole read of the saved SR1--13 proof and all seven substantive sections, then the complete companion script. I checked the literal rowDiv word/coefficient representation, leading-block restriction including leave-and-return words, integral row-matrix stability over the polynomial coefficient algebra, exact determinant-twist lattice equality, and both symbolic retractions. I reconstructed the fixed-triple augmentation/auxiliary-section consequence and the middle perfect-dual composition. I separately checked the distinction between the abstract reduction of the integral Weyl lattice and its possibly smaller polynomial image modulo11. The actual theta functional, class sum and terminal branch are not inverted or replaced. The final status and this verification paragraph were reread; the separate verification JSON records the final replay and fingerprints.

\clearpage

\section{AM-LIFT-071 --- canonical Hilbert normalization of the prescribed split T test}\label{am-lift-071-canonical-hilbert-normalization-of-the-prescribed-split-t-test}

\textbf{Canonical checked version --- 2026-11-22T19:19:09.748711+00:00.} Root and independent whole proof/source/script checks and separate147,603-assertion replays passed. Original pending wording is historical. Promoted scope: For the actual positive-rank physical T test D\emph{tau, the canonical Hilbert/GGP factor is q/{[}(q-1)(q\^{}k-1){]} with probability H Haar and the specified product root Haar, independently of rho. Its inverse square root is less than q\^{}(k/2). For literal source Haar it is c\_H}c\_N times that value. The scalar-one matrix/linear-period identity is calibrated by compact Kirillov tests and the raw adjacent-rank theorem, then extended by split uniqueness and full weak-Schwartz continuity. The genuine full source and chosen-bottom representation premises of LIFT069/070 remain. No global period, actual Read/032 value, selected arithmetic image or DP44 cost is supplied.

Author: \nolinkurl{/root/lift_accounting}, 22 November 2026 UTC.

\textbf{Author whole-proof, source-interface and script review PASS; exact diagnostic PASS (147,603 assertions).} Final timestamps and frozen hashes are recorded in the companion author-review and manifest files. Independent peer and root reviews are separate. The scalar comparison below is proved by a compact Kirillov calculation with all quotient Jacobians. Multiplicity one only extends that calibrated identity. No global period or actual source-family existence is asserted.

\subsection{1. The fixed test and the two Haar conventions}\label{the-fixed-test-and-the-two-haar-conventions}

Keep LIFT068's data: F is nonarchimedean, psi has conductor O, q is the residue cardinality, k\textgreater=2, rho\textgreater=0 and n=k+2rho+1. Pi is tempered generic spherical on GL\_n(F); sigma is tempered generic on H=GL\_k(F), with conductor and central-character conductor both one. Their essential Whittaker vectors W\_Pi and W\_sigma have value one at1, with characters psi and psi\^{}(-1). We retain the actual covector transport, the marked line and

\[
 D=\operatorname{diag}(I_k,cI_\rho,1,\varpi^{-1}I_\rho),\qquad
 \tau=I_n+\varpi^{-1}E_{k+\rho+1,k},\qquad c\in O^\times.
 \tag{HN1}
\]

The block order is (H,E,z,F). These are the original vectors of LIFT068, not replacements by arbitrary Whittaker vectors. Write

\[
 \kappa_r=\prod_{i=1}^r(1-q^{-i}),\qquad \kappa_0=1,
 \quad dg_r^{\rm raw}=\frac{dX}{|\det g|^r}
                    =\kappa_r\,dg_r^{\rm prob}.
 \tag{HN2}
\]

Here additive O and every standard root O have volume one, while the probability convention makes GL\_r(O) have volume one. Quotients by N\_r always use the indicated group Haar divided by the product root Haar. In particular the only conversion in the central JPSS integral is

\[
 Z_\rho^{\rm raw}=\kappa_k Z_\rho^{\rm prob}.
 \tag{HN3}
\]

There is no factor kappa attached to an additive matrix integral or to an individual root group. The canonical Whittaker metric used below is

\[
 \vartheta_r(W,W')=\int_{N_{r-1}\backslash GL_{r-1}}
 W(\operatorname{diag}(a,1))
 \overline{W'(\operatorname{diag}(a,1))}\,d\bar a^{\rm raw}.
 \tag{HN4}
\]

It is positive and GL\_r-invariant for tempered generic representations, by Jacquet--Shalika (1983), section2.3. We use the inner product linear in its first entry. The raw Bessel form is the canonical tempered regularization of

\[
 \mathcal A_1(W\otimes V,W'\otimes V')=
 \int_H\int_{N_B}
   \vartheta_n(\Pi(uh)W,W')
   \vartheta_k(\sigma(h)V,V')\,\psi_B(u)^{-1}\,du\,dh^{\rm raw}.
 \tag{HN5}
\]

Its character is the unit character in LIFT068 section7. On H,E,z,F, N\_B has the two upper flag groups, the E-to-H lower block, and the E-to-z, H-to-F, z-to-F and E-to-F blocks. The character is the sum of the E-chain, E-to-z, z-to-F and F-chain entries; it is trivial on all other displayed roots. H has modulus one on N\_B, since its rho incoming and rho outgoing blocks balance.

The regularization is the continuous extension on the weak Harish-Chandra Schwartz space of Beuzart-Plessis (2020), Proposition7.1.1; its all-place, including split-place, use is explicitly given by Beuzart-Plessis--Chaudouard (2025), section8.4. We do not assume that the bare inner integral in HN5 is absolutely convergent.

\subsection{2. The compact-cutoff normalization lemma}\label{the-compact-cutoff-normalization-lemma}

For any tempered generic Pi on GL\_n and sigma on GL\_k with n=k+2rho+1, define

\[
 A_\rho(a,X)=
 \begin{pmatrix}a&0&0\\X&I_\rho&0\\0&0&I_{\rho+1}\end{pmatrix},
 \qquad
 (\mathcal T_\rho W)(a)=|\det a|^{-\rho}
             \int_{\operatorname{Mat}_{\rho,k}(F)}W(A_\rho(a,X))\,dX.
 \tag{HN6}
\]

For compactly supported Kirillov restrictions the integral is an ordinary compact integral. In general its pairing with V is the central JPSS functional, with its established continuation:

\[
 Z_\rho^{\rm raw}(W,V)
       =\int_{N_k\backslash H}(\mathcal T_\rho W)(a)V(a)
                                          \,d\bar a^{\rm raw}.
 \tag{HN7}
\]

\textbf{Normalization lemma.} With the raw metric and measures just specified,

\[
 \boxed{\mathcal A_1(W\otimes V,W'\otimes V')
       =Z_\rho^{\rm raw}(W,V)
                    \overline{Z_\rho^{\rm raw}(W',V')}.}
 \tag{HN8}
\]

We give the cutoff and measure proof, rather than determine a scalar by uniqueness alone. For rho=0 this is exactly Zhang's adjacent-rank identity and the already checked K054 normalization; the nontrivial cutoff induction below concerns rho\textgreater0.

\subsubsection{2.1. Exact quotient-row measure}\label{exact-quotient-row-measure}

Put d\textgreater=2. On the chart where the last coordinate of the last row is t!=0, use

\[
 x=\begin{pmatrix}a+by&bt\\y&t\end{pmatrix},
 \quad a\in GL_{d-1},\ b\in F^{d-1},\ y\in F^{d-1},\ t\in F^\times.
\]

The additive Jacobian is \textbar t\textbar\^{}(d-1) and det(x)=det(a)t. Left N\_d acts by the standard left N\_(d-1) action on a and by affine unipotent changes on b; its root measures have Jacobian one. Therefore

\[
 d\bar x^{\rm raw}
    =|\det a|^{-1}d\bar a^{\rm raw}\,dy\,\frac{dt}{|t|}.
 \tag{HN9}
\]

At last row (0,\ldots,0,1), Fourier inversion therefore leaves the factor \textbar det a\textbar\^{}(-1). The same factor remains when only some entries of y are fixed by Fourier inversion and the other entries are integrated. This factor is essential; it is not replaced by one.

\subsubsection{2.2. Compact test space and the actual regularization}\label{compact-test-space-and-the-actual-regularization}

By Jacquet--Shalika section2.4, every smooth function in compact induction from the generic character of N\_(n-1) to GL\_(n-1) is the restriction of a Whittaker vector of Pi. First take W,W' in this compact Kirillov subspace. For the scalar calibration it is enough to choose their supports in a compact neighborhood of the identity in the lower Gauss chart. We make that choice, with every diagonal pivot a unit. Their restrictions and all restrictions to the row charts below then have compact support modulo the remaining upper unipotent group. The final multiplicity-one argument, not a claim about every compact support, extends the calibrated identity to all vectors.

Consider

\[
 a(t)=\operatorname{diag}
  (I_k,t^\rho,t^{\rho-1},\ldots,t,1,t^{-1},\ldots,t^{-\rho}).
 \tag{HN10}
\]

It centralizes H and multiplies the generic character coordinate lambda\_B by t. Choose a compact open subgroup U of O\^{}times for which a(U) fixes W and W'. The sigma factors are unchanged. Beuzart-Plessis's explicit regularizer gives HN5 as the absolutely convergent integral obtained by replacing psi\_B\^{}(-1)(u) by

\[
   \frac1{\operatorname{vol}(U)}\int_U
                          \psi(-t\lambda_B(u))\,d^\times t.
 \tag{HN11}
\]

The integrability is of the matrix-coefficient product times HN11. It is not an integrability assertion for the expanded absolute Kirillov integrand.

Use compact exhaustions of H and compact coordinate boxes of N\_B that are invariant under conjugation by a(U). Every root coordinate has a positive integral weight for a(t). Thus such boxes may be enlarged successively, with the exponent on a root at least the sum of the exponents needed for its lower-root commutators. They exhaust N\_B. On each box, change u by a(t)ua(t)\^{}(-1); all absolute root Jacobians are one, the box is unchanged, and the two vectors are fixed. Hence its unaveraged integral with psi\_B\^{}(-1) equals its integral with HN11. Absolute convergence of the latter proves that every cofinal exhaustion of these boxes computes the canonical form.

Here is the precise compatibility used for a full column. Give coordinate i the exponent w\_i in HN10. A remaining (i,j) root has weight w\_i-w\_j; the j-th entry of a column ending at l has weight w\_j-w\_l. Multiplication replaces its i-th entry by a sum containing their product, whose weight is exactly w\_i-w\_l. Weighted balls with these additive exponents are therefore preserved by the bounded remaining unipotent changes of column variables. Extra fixed margins may be included in additive root-height weights. A fixed compact H exhaustion only requires a further finite margin. These facts apply to the literal matrix-root coordinates; no exponential map or residue-characteristic restriction is used.

All Fourier inversions below are first made with H in a compact set, the remaining non-Fourier root coordinates bounded, and the Kirillov variable in the compact support of the second vector. The functions to which Fourier inversion is applied are then uniformly locally constant and compactly supported in the row variables. Increasing the relevant Fourier ball therefore gives exact Fourier inversion, uniformly on those compact parameter sets. There is no formal use of a Dirac delta against a noncompact unsmoothed integrand.

There is a useful way to order the E-row integrations which avoids a circular choice of Fourier and row cutoffs. After the full F-column reductions, use the calibration Gauss chart \[
 x=\begin{pmatrix}a&0\\X&B\end{pmatrix},
 \qquad B\text{ lower triangular with unit pivots}.
\] For a fixed remaining upper E matrix v, write Bv=n\_v l\_v wherever the first compact test can be nonzero. Its lower Gauss factor l\_v has unit pivots on that support. Right translation by the E-to-H matrix Y, followed by removal of n\_v on the left, changes the lower H block to n\_v\^{}\{-1\}(X+BY). This is an invertible affine change of Y, with absolute Jacobian \textbar det B\textbar\^{}k=1. The actual Whittaker phase of n\_v is retained; it depends on B and v, not on Y. For bounded v and h, compact support of the first test bounds the transformed block, hence bounds Y uniformly on the contributing set. Thus these Y integrals may first be completed as ordinary compact integrals, before increasing a later E Fourier cutoff. On the complement of the indicated Gauss cell the first compact test is zero. Subsequent Fourier stabilization no longer depends on the size of a Y cutoff; for each bounded Fourier box, enlarge Y until its entire contributing support is included. This ordering, combined with the preceding column invariance, gives a cofinal sequence of compact boxes computing the required successive integrals.

Equivalently, in the successive E-row calculation below the extra substitutions are additive triangular shears of the remaining E-to-H variables. Their Jacobians are one. They also respect the root weights in HN10. For example an E\_i-to-E\_j coordinate v and an already exposed row X produce the change Y\_i to Y\_i-vX. If E\_i,E\_j have weights w\_i\textgreater w\_j\textgreater0, the v ball has weight w\_i-w\_j, whereas Y\_i has weight w\_i. The exposed X row stays in a fixed compact set for bounded H. Taking the row cutoff sufficiently large ensures vX lies in the Y\_i ball, so the substitution preserves that truncated integral. The completed-row ordering just proved allows this enlargement before the next Fourier limit. The same argument applies to each of the finitely many triangular substitutions. Arbitrary prescribed lower bounds on all cutoffs may be imposed, so HN11 identifies their limit with the canonical regularization.

\subsubsection{2.3. The two strings of row reductions}\label{the-two-strings-of-row-reductions}

Start from HN4 for Pi and write N\_B with a last full F-column first. Its action on the Kirillov restriction multiplies the function by \nolinkurl{psi(e_last^t x v)}. The corresponding Bessel character is \nolinkurl{psi(v_last)}. Fourier inversion in that full column fixes the last row of x to (0,\ldots,0,1). HN9 supplies \textbar det a\textbar\^{}(-1). Repeat for the rho F-columns. The remaining restriction is

\[
 x\longmapsto W(\operatorname{diag}(x,I_{\rho+1})),
             \qquad x\in GL_{k+\rho},
\]

and its quotient measure has acquired \textbar det x\textbar\^{}(-rho).

The remaining radical has blocks E-to-H, upper U\_E, and E-to-z. At the last E-to-z step, Fourier inversion fixes the E part of the last row of x to (0,\ldots,0,1), leaving its H part X free. A representative is

\[
 x=\begin{pmatrix}a&0\\(X,0_{1\times(\rho-1)})&1\end{pmatrix},
 \qquad a\in GL_{k+\rho-1}(F).
 \tag{HN12}
\]

Equation HN9 contributes one more \textbar det a\textbar\^{}(-1). Integrating the last E-to-H row Y changes X to X+Y on the first vector and separates the two X integrals. For an H element h, the first row argument is (X+Y)h, so this change also contributes \textbar det h\textbar\^{}(-1).

At each subsequent step replace rho in HN12 by the number of E rows still present, and write tilde\_X for its extended row \nolinkurl{(X,0,...,0)}. Multiplication by an upper E column v gives \nolinkurl{a(I-v tilde_X)} after removing the upper-right column. Here tilde\_X v=0, because v has only E entries. That term is absorbed by the remaining E-to-H variables by precisely the Jacobian-one triangular substitution in section2.2. It is not discarded. Repeating the E step rho times gives the two complete integrals over Mat\_(rho,k), and the accumulated weight is

\[
                    |\det a|^{-2\rho}|\det h|^{-\rho}.
 \tag{HN13}
\]

No Whittaker character is introduced on a lower E-to-H translation. The remaining generic character is the standard character on N\_k, inherited from the original larger Whittaker function. The opposite character of V cancels it in HN7.

Using HN6, the result of the entire cutoff calculation is consequently

\[
 \int_H\int_{N_k\backslash H}
 (\mathcal T_\rho W)(ah)
  \overline{(\mathcal T_\rho W')(a)}\,
       \vartheta_k(\sigma(h)V,V')\,d\bar a^{\rm raw}dh^{\rm raw}.
 \tag{HN14}
\]

The displayed expression also checks the modulus: substitution X to Xh gives \nolinkurl{T_rho(Pi(h)W)(a) = T_rho W(ah)} after including \textbar det a\textbar\^{}(-rho).

\subsubsection{2.4. Scalar calibration and extension}\label{scalar-calibration-and-extension}

For the compact calibration tests just used, T\_rho W and T\_rho W' are smooth and compactly supported modulo N\_k. This follows successively from HN9/HN12: the fixed last-row slices are closed in their row charts, and the remaining X integrations are projections of compact supports. In particular any compact generic function supported in a sufficiently small patch around the identity in N\_k backslash H can occur: extend through these row charts with compact locally constant bumps of value one on the fixed Fourier slices and integral one in the rho free X rows. This is all the local surjectivity needed for a nonzero calibration pair.

Choose any auxiliary tempered generic representation of GL\_(k+1). Jacquet--Shalika realizes these two functions as its compact Kirillov restrictions, with metric HN4. Then HN14 is its ordinary adjacent GL\_k matrix-coefficient period paired with sigma. It is absolutely convergent, and Zhang (2014), Proposition4.10, with the raw measures stated at the beginning of his section4.5, identifies it exactly with the product of its two central Rankin--Selberg integrals. Those integrals are HN7. This proves HN8 on the compact calibration tests.

Such tests can have nonzero HN7: choose V with V(1) nonzero, take a sufficiently small compact patch about1 on which it is nonzero, and take a compact generic test function paired with it. The preceding bump extension supplies W. The analogous choice supplies W',V'. Gan--Gross--Prasad, Corollary15.2 (explicitly including the split quadratic algebra), says the Bessel Hom space has dimension at most one. The regularized form has the required equivariance by its translation covariance; the central functional has exactly the same equivariance by the explicit LIFT068 section7 computation and its H covariance PT11. The single nonzero compact calibration therefore fixes their proportionality scalar to one on the entire smooth representations. This use of uniqueness occurs only after the scalar has been calculated. It does not appeal to Hilbert density to extend an unbounded Bessel functional.

\subsection{3. Exact norms of the two prescribed Whittaker vectors}\label{exact-norms-of-the-two-prescribed-whittaker-vectors}

Put

\[
 L_G=L(1,\Pi\times\Pi^\vee),\qquad
 L_H=L(1,\sigma\times\sigma^\vee),\qquad
 L_0=L(1/2,\Pi\times\sigma).
 \tag{HN15}
\]

For the spherical vector, the Casselman--Shalika diagonal formula and the Schur Cauchy identity, in probability quotient Haar, give L\_G/ζ\_q(n). One can verify the quotient factor directly: every dominant partition is its last-zero partition plus an integer multiple of (1,\ldots,1); the latter geometric series is ζ\_q(n), because the central character is unitary. Multiplying by kappa\_(n-1) gives

\[
               \|W_\Pi\|_{\vartheta_n}^2=\kappa_n L_G.
 \tag{HN16}
\]

Under the actual conductor-one hypothesis, K053 and LIFT068 show that sigma has one ramified unitary character and k−1 unramified unitary characters. Its formal Satake list is (beta\_1,\ldots,beta\_(k-1),0). Miyauchi's Theorem4.1 gives the same diagonal Schur formula for its essential vector. The sum over last-zero partitions is now the full Cauchy sum in the k−1 nonzero variables; no ζ\_q(k) factor is removed. It equals the product over the unramified pairs. The true L\_H has one additional self-pair of the ramified character, namely ζ\_q(1), while every ramified/unramified cross factor is one. Thus

\[
          \|W_\sigma\|_{\vartheta_k}^2
             =\kappa_{k-1}(1-q^{-1})L_H.
 \tag{HN17}
\]

All integrals converge: the relevant parameters have modulus one and their Cauchy sums are evaluated at q\^{}(-1). The L values in HN16--17 are positive because they equal the displayed squared norms. Invariance of HN4 means that D tau does not alter the first norm. These statements concern the actual Whittaker vectors of value one, not a normalization chosen to cancel the period. They are precisely the previously checked K054(CP9--11) norm calculation, now with the large rank n=k+2rho+1.

Divide those two vectors by the positive square roots in HN16--17, and denote them by v\_Pi and v\_sigma. Write Z\_rho,unit for HN7 evaluated on Pi(D tau)v\_Pi and v\_sigma. HN3 and LIFT068(PT4) then give

\[
 |Z_{\rho,\mathrm{unit}}^{\rm raw}|^2
 =\frac{\kappa_k^2q^{-\rho(k+\rho+1)}q^k|L_0|^2}
 {(q^k-1)^2\kappa_n\kappa_{k-1}(1-q^{-1})L_G L_H}.
 \tag{HN18}
\]

The exact complex phase is still LIFT068(PT25); taking a squared norm loses no datum needed for this positive form.

\subsection{4. The physical character and regularized change of variables}\label{the-physical-character-and-regularized-change-of-variables}

LIFT067/LIFT068 give the physical character chi\_c(u)=psi\_B(DuD\^{}(-1)) and

\[
                         \delta_{N_B}(D)=q^{-\rho(k+\rho+1)}.
 \tag{HN19}
\]

For compactly supported functions, conjugation u to DuD\^{}(-1), with D centralizing H, gives the ordinary change-of-variables factor delta\_(N\_B)(D)\^{}(-1). This identity extends to the canonical tempered regularization: fixed left/right translation is continuous on the weak Harish-Chandra Schwartz space, compactly supported smooth functions are dense there, and both character-period functionals are its continuous extensions. Thus no assertion about an unregularized convergent inner integral is required.

In particular, on the original bare translate tau, with H probability Haar and the original root measures, define the physical unit-vector Bessel period by \[
 \alpha_{T,\mathrm{phys}}=
 \int_H^{\rm reg}\int_{N_B}
 \langle\Pi(uh)\Pi(\tau)v_\Pi,\Pi(\tau)v_\Pi\rangle
 \langle\sigma(h)v_\sigma,v_\sigma\rangle
 \chi_c(u)^{-1}\,du\,dh^{\rm prob}.
\] The superscript means the joint canonical Bessel regularization just defined, not an unsupported order of two conditionally convergent integrals. Its exact value is

\[
 \alpha_{T,\mathrm{phys}}
    =\delta_{N_B}(D)^{-1}\kappa_k^{-1}
                              |Z_{\rho,\mathrm{unit}}^{\rm raw}|^2.
 \tag{HN20}
\]

The factor kappa\_k\^{}(-1) converts the one outer H integration in the matrix period. It is not squared. In contrast HN3, which converts each linear Rankin--Selberg functional, is squared in HN18. This distinction accounts for the remaining rational local factor.

The standard split-unitary local L factor in these conventions is

\[
 \mathscr L_T=\Delta_n\frac{|L_0|^2}{L_G L_H},
                       \qquad\Delta_n=\kappa_n^{-1}.
 \tag{HN21}
\]

This is Beuzart-Plessis--Chaudouard (2025), section1.3.6, printed9, specialized to a split place: the product of Tate factors runs to the larger rank n, the two split base-change components give \textbar L\_0\textbar\^{}2, and the two Asai factors give L\_G and L\_H. In particular these latter factors retain the unitary central direction; they are not the general-linear adjoint factors with a scalar line removed.

Then exact cancellation in HN18--21 gives

\[
 \boxed{\quad
  \alpha_{T,\mathrm{phys}}^{\natural}
      :=\frac{\alpha_{T,\mathrm{phys}}}{\mathscr L_T}
       =\frac{q}{(q-1)(q^k-1)}>0.
 \quad}
 \tag{HN22}
\]

This is independent of rho. Both the half-power in the fixed JPSS test and the physical root modulus have been retained before cancellation. The corresponding inverse square-root cost is

\[
 (\alpha_{T,\mathrm{phys}}^{\natural})^{-1/2}
        =\sqrt{\frac{(q-1)(q^k-1)}q}<q^{k/2}.
 \tag{HN23}
\]

Changing the prescribed root Haar or outer H Haar multiplies alpha by those actual measure factors; it does not preserve HN22 verbatim. Similarly an independently rescaled source intertwiner remains a separate scalar. The formulas use precisely the normalized local vectors and root conventions of HN1--5.

More explicitly, if a source realization uses \nolinkurl{dh_source=c_H dh_prob} and \nolinkurl{du_source=c_N du_product}, with the latter product consisting of the root measures in HN1--5, then its same-vector normalized value is \[
 \alpha_{T,\mathrm{source}}^{\natural}
       =c_Hc_N\,\frac{q}{(q-1)(q^k-1)}.
 \tag{HN24}
\] This is the exact conversion to that specified Haar, not a claim that the abstract source IsWu measure parameters have automatically been set to one. LIFT060 section8 explicitly retains those parameters. For a product of specified T places the corresponding factors, including c\_H and c\_N, are multiplied once at each place.

\subsection{5. Actual source hypotheses and retained global gate}\label{actual-source-hypotheses-and-retained-global-gate}

The split degree-one T field, physical line, covector, additive character and D tau dictionary are the literal ones checked in LIFT060/067/068. LIFT069 supplies finite temperedness and genericity on its chosen LIFT059 bottom subspace, not on every abstract e0 vector. LIFT070 supplies the same finite properties for every actual initial constituent only under its full genuine SizesAt/SizeData\_2 hypotheses, and identifies the strict Mram line with conductor-one nontrivial tame central character. The bottom maximal T level supplies the spherical vector. These are exactly the representation hypotheses of HN1; a later non-tempered theta descendant is not substituted for this bottom representation.

Hence the local normalized number HN22 is available for that unchanged actual initial/bottom pair when these proved premises are instantiated. It does not prove that a complete Host/SizesAt family exists, that the selected fixed-level global Bessel/Mram functional is nonzero, or that a global central L-value is nonzero. It also constructs no Read, STiesR or specific032 value. Other finite places, all real relative-model normalizations, global measures and coefficient-period factors remain distinct in the global formula.

\subsection{6. Primary anchors and diagnostic boundary}\label{primary-anchors-and-diagnostic-boundary}

\begin{itemize}
\tightlist
\item
  Jacquet, H., \& Shalika, J. (1983). \emph{The Whittaker models of induced representations}. Pacific Journal of Mathematics, 109(1), 107--120. Sections1.2 and2.3--2.4, printed109--111: mirabolic model, invariant inner product and compact Kirillov subspace. Directly read author copy: \url{https://www.math.columbia.edu/~hj/The\%20Whittaker\%20model\%20of\%20induced\%20representations.pdf}.
\item
  Zhang, W. (2014). \emph{Automorphic period and the central value of Rankin--Selberg L-function}. Journal of the American Mathematical Society, 27, 541--612. Sections3.2 and4.5, especially raw measures printed569 and Proposition4.10 with its complete proof printed571--574. Direct original text retained at \nolinkurl{agents/local_t_branch_refs/zhang2014_period.txt}.
\item
  Beuzart-Plessis, R. (2020). \emph{A local trace formula for the Gan--Gross--Prasad conjecture for unitary groups: The archimedean case}. Astérisque, 418. Proposition7.1.1 and its p-adic proof, printed159--160; Lemma7.1.2, printed161--162. Direct original: \url{https://smf.emath.fr/system/files/filepdf/smf_ast_418.pdf}.
\item
  Beuzart-Plessis, R., \& Chaudouard, P.-H. (2025). \emph{The global Gan--Gross--Prasad conjecture for unitary groups. II. From Eisenstein series to Bessel periods}. Forum of Mathematics, Pi, 13, e16. Section8.4, printed76--77, gives the all-place weak-Schwartz extension; section8.2 records compatible self-dual Fourier inversion; section1.3.6, printed9, gives the exact Bessel local L normalization. \url{https://doi.org/10.1017/fmp.2025.8}.
\item
  Gan, W. T., Gross, B. H., \& Prasad, D. (2012). \emph{Symplectic local root numbers, central critical L-values, and restriction problems in the representation theory of classical groups}. Astérisque, 346, 1--109. Author \nolinkurl{ggp8} section13 and section15/Corollary15.2, printed53--55/60--63, explicitly includes the split quadratic algebra. \url{https://abel.math.harvard.edu/~gross/preprints/ggp8.pdf}.
\item
  Miyauchi, M. (2014). \emph{Whittaker functions associated to newforms for GL(n) over p-adic fields}. Journal of the Mathematical Society of Japan, 66(1), 17--24. Theorem4.1, printed21--22, in the already checked local primary cache.
\item
  Internal fixed-test and genuine-source inputs: LIFT067(TC5,TC11), LIFT068(PT4,PT25--27), K053(NT7--14), K054(CP4--14), LIFT069 and LIFT070. Their scope is retained above.
\end{itemize}

The exact companion certificate, \nolinkurl{positive_rank_t_hilbert_normalization_certificate.py}, checks the quotient-row Jacobian, determinant and triangular-root identities; finite Fourier inversion including the free H-row in a partial column; translation separation; root-modulus exponents; and the exact kappa/norm cancellation. Its 147,603 assertions passed. These are diagnostics of universal calculations, not numerical evaluations of an actual p-adic representation or of its global period. The analytic regularization, cutoff induction and source scope were separately read in full for the author check. Final byte identities and the replay are recorded in the companion author-review JSON and manifest.

\clearpage

\section{AM-DO-084 --- an original-ring volume bound for the rank-one fixed-space block}\label{am-do-084-an-original-ring-volume-bound-for-the-rank-one-fixed-space-block}

\textbf{Canonical checked version --- 2026-11-22T19:22:28.593688+00:00.} Complete root and independent whole proof/script reviews and separate2,904-assertion replays passed. Original pending wording below is historical. Promoted scope: The original integral rank-one327-entry pivot chart lies in one Frobenius ball. Its ring-volume count replaces only the old independent-entry chart count, retaining every original seven-column, Green, source-fibre, support and pencil factor. The241 fresh F,r1 nonunipotent cells of FV1--18 contribute\textless2\^{}-4255. Together with the old R050/DO074 cuts the entire nonunipotent F,r1 sector through support18 is removed, with the listed higher-support tails. This finding alone does not raise the complete support floor.

\textbf{Algebraic-geometric translation.} In algebraic-geometric terms this counts original integral Hermitian endomorphism charts and labelled correspondences without changing their coefficient modules, polarizations or Picard/biextension data. The carried module/determinantal anchors are SGA3 I4.7.3/4.7.3.1 and SGA6 VI3.3. Picard/biextension anchors are FGA232 Theorem3.1 pp148--149; FGA236 Corollaries1.7/2.3 pp229/232; SGA7 I VII2.7/2.9 pp169--174 and VIII7.4(b) pp310--311. Crystal/Cartier duality retains cote4 scan003 §§8.1.1--8.1.2, scans005--006 §§8.4.1--8.4.3 and cote54 scans005--007. These are contextual locators; no new scheme, lambda-ring/K-theory identity, l-adic sheaf/weight theorem, CM motive or motivic L-function claim follows. DP\_MULT Q3e and THETA3 G1/A7 obtain precisely the stated smaller analytic remainder. CX\_GRH CX-ALG-024 and GCT\_PLETH T-052 retain their distinct missing arithmetic comparisons.

Author: ALP\_MACHINE \nolinkurl{/root/definite_oscillator}, 22 November 2026 UTC. New controlled pending proof; final author and independent checks are recorded separately.

\subsection{1. Exact actual class and assertion}\label{exact-actual-class-and-assertion}

Retain the original J56 all-top signed kernel, K044 upper endpoint, charged R048 cutoff H=133, all actual finite projectors, ordered T-pencil pairs, class masses, phases and real Hilbert normalization. Put f=a=168, n=169 and m=164. Use the original decomposition E\^{}169=S perpendicular W0, with dim(S)=m and signature(W0)=(4,1). The K068 residual has an integral S block A, determinant-one representative h, and unipotent reductions at all334 unmixed G factors.

This report concerns only the class

\[
 X=A-I\in M_m(\mathcal O_E),\quad\operatorname{rank}_E X=1,
 \quad h\text{ fixes }K_A=\ker X\cap\ker X^*\text{ pointwise},
 \quad h\text{ is not globally unipotent}.
 \tag{FV1}
\]

Thus it is the actual K069 fixed-space class F at block rank one. The integrality of X comes from the carried K068 original-coordinate residual, not from an arbitrary analytic fixed-space assumption. Let s count active relative T places and t count precisely the type-III unequal-pencil off-diagonal places.

The theorem removes the following fresh tails of the already checked R050 ranges:

\begin{longtable}[]{@{}rrr@{}}
\toprule\noalign{}
t & New supports s & Number \\
\midrule\noalign{}
\endhead
\bottomrule\noalign{}
\endlastfoot
11 & 60 through98 & 39 \\
12 & 49 through87 & 39 \\
13 & 38 through76 & 39 \\
14 & 27 through65 & 39 \\
15 & 16 through54 & 39 \\
16 & 16 through43 & 28 \\
17 & 17 through31 & 15 \\
18 & 18 through20 & 3 \\
\end{longtable}

Their entire absolute contribution is

\[
                         <2^{-4255}.
 \tag{FV2}
\]

There are241 support/type cells. The bound includes every actual label pair and all other factors; no additional cell hypotheses or favorable sign is used. Remove only their intersection with the current remainder and append4255 once. Earlier charges remain. Together with DO074 and the old R050 intervals, this covers the full nonunipotent F,rank-one-block sector through total support18. It does not exhaust the other structural classes at those supports.

\subsection{2. The stored original chart has a single Frobenius-norm bound}\label{the-stored-original-chart-has-a-single-frobenius-norm-bound}

At each complex realization the original coordinate compression satisfies \textbar\textbar A\textbar\textbar\_op\textless=\textbar\textbar h\textbar\textbar\_op\textless=H. Hence

\[
             \|X\|_{\mathrm{op}}\le H+1=134.
 \tag{FV3}
\]

The rank of X is one at every faithful embedding. Therefore its Frobenius norm equals its only nonzero singular value:

\[
                  \|X\|_{\mathrm{HS}}=\|X\|_{\mathrm{op}}\le134.
 \tag{FV4}
\]

Choose a nonzero original entry X\_(i0,j0). Store its complete original row and column, counting the common entry once. There are

\[
                           d=2m-1=327
 \tag{FV5}
\]

stored integral entries. They uniquely determine X by

\[
           X_{ij}=X_{i,j_0}X_{i_0,j}/X_{i_0,j_0}.
 \tag{FV6}
\]

No inverse is applied to a lattice, and no reconstructed matrix outside the actual class is asserted integral. Fixing the first nonzero pivot gives an injection for each pivot; taking all m\^{}2 possible pivots is a valid upper count.

The stored entries are a subset of the distinct entries of X. Their sum of squared moduli at every embedding is therefore at most134\^{}2 by FV4. Across all f complex labels, the stored vector lies in the original ring lattice O\_E\^{}d, in a Euclidean ball of radius134 sqrt(f). This is the improvement over treating all327 entries as independent radius134 scalars.

\subsection{3. Exact original-ring cover and volume}\label{exact-original-ring-cover-and-volume}

Use the same half-trace metric as every earlier column count,

\[
                  (x,y)=\tfrac12\operatorname{Tr}_{E/\mathbf Q}(x\bar y).
\]

LIFT063's checked original-ring theorem and metric addendum give

\[
 \mu(\mathcal O_E)<57\sqrt f,\qquad
 \operatorname{covol}(\mathcal O_E)=D_0,
 \qquad D_0=7^{140}29^{144}5^{84}13^{84}.
 \tag{FV7}
\]

These are statements for the literal ring O\_E in the unweighted original-coordinate metric. No conclusion about an arbitrary fractional ideal, a polarized Gaussian metric or a moving rational subspace is needed.

Squared covering radii add in an orthogonal product. Since57\^{}2·327\textless1031\^{}2,

\[
             \mu(\mathcal O_E^{327})<1031\sqrt f.
 \tag{FV8}
\]

Attach a Dirichlet cell of this lattice to each stored vector in the ball from FV4. The cells are disjoint up to measure-zero boundaries, each has volume D0\^{}327, and each lies in the concentric ball of radius less than(134+1031)sqrt(f)=1165sqrt(f). Thus the complete chart count is at most the Euclidean ball volume divided by D0\^{}327. Replacing pi by22/7 gives the rational upper envelope

\[
 V_{327}(1165)=
 \frac{(528\cdot1165^2)^{327f}}{(327f)!D_0^{327}},
 \qquad \frac{22}{7}f=528.
 \tag{FV9}
\]

Exact integer multiplication and factorials prove

\[
                    2^{307589}<V_{327}(1165)<2^{307590}.
 \tag{FV10}
\]

The lower inequality describes this chosen rational envelope only; it is not a lower bound on the actual lattice count. Combining the m\^{}2 pivot choices yields

\[
                       \#\{A\}<m^2 2^{307590}.
 \tag{FV11}
\]

All global integrality and metric assertions in this step concern the original stored ring entries. The formula FV6 supplies only uniqueness.

\subsection{4. Insert the new chart count into the unchanged complete R050 estimate}\label{insert-the-new-chart-count-into-the-unchanged-complete-r050-estimate}

We retain R050 sections2--3 in their exact scopes. For each fixed A in F, choose once two original S coordinates injective on K\_A's positive orthogonal complement in S, and retain all five W0 entries of h e\_n.~The real dimension bound is unchanged. The actual projected target over-lattice contains O\_E\^{}7 with paired index T\^{}(14s), uniformly in the ordered pencil pairs. The old count per coset is(2H+1)\^{}(14f)=267\^{}(14f), and the full original source fibre remains G\_n\textless2\^{}5022. No smaller fibre or alternative integral basis is borrowed.

For globally nonunipotent h, the same two-coefficient witness is nonzero: at least one of e1(h−I),e2(h−I) is nonzero by the accepted signature-one two-trace theorem. The original determinant-one finite argument gives paired valuation at least2(2−d\_v) at every unmixed factor of rank d\_v\textless=1. Both coefficients have only simple poles at the paired type-III primes and are integral elsewhere. The same dimension-at-most7 singular-value bound gives B=1536 at every infinity. Consequently the retained exact Green envelope is

\[
 |a(h)|<2 B^{335f/2} T^{335t/2}Q^{-335(f-1)},
 \qquad B=1536,
 \tag{FV12}
\]

where Q=g0g1. The mixed factors have their original norm bound one. No rank-pattern count or global-unipotence premise is introduced, and \textbar c(h)\textbar\textless=1 suffices here.

At each type-III place the full (T−1)(T−2) ordered-pair count cancels the exact absolute rho weight. At each other active place retain the complete bound P\_T=2(T−1)\^{}2/(T−2)\textless2\^{}997. It already counts types I and II together. Since the matrix count is uniform in the pairs, multiplying these actual pair/weight sums is valid.

Replacing only the A chart in R050, the complete contribution of one FV1 cell satisfies

\[
 \begin{split}
 U_{s,t}<&\ 2\binom fs\binom st m^2\,2^{307590}G_n
       267^{14f}B^{335f/2}\\
 &\quad\cdot T^{14s+335t/2}
              Q^{-335(f-1)}P_T^{s-t}.
 \end{split}
 \tag{FV13}
\]

Every original class, source-lattice, finite-projector and incidence factor of R050 remains. Counting nonoccurring charts or matrices more than once only increases this upper bound.

\subsection{5. Rational constants, exact tails and the fresh charge}\label{rational-constants-exact-tails-and-the-fresh-charge}

Use the same conservative exact bounds as R050:

\[
 \begin{gathered}
 \log_2 267<8+249/4096,\quad
 1536^5<2^{53},\quad T<2^{996},\quad Q^4>2^{281},\\
 m^2<2^{15},\quad\binom fs\le2^{168},\quad
 \binom st\le2^s,\quad G_n<2^{5022}.
 \end{gathered}
 \tag{FV14}
\]

Then

\[
 \log_2 U_{s,t}<C_F+14942s+165833t,
 \tag{FV15}
\]

where the full constant is

\[
 \begin{split}
 C_F={}&1+168+15+5022+307590
       +14f(8+249/4096)\\
 &+(335f/2)(53/5)-(281/4)335\cdot167
       =-844824901/256.
 \end{split}
 \tag{FV16}
\]

The slopes are unchanged:14942=14·996+997+1 and165833=335·996/2−997. The constant decreases from the R050 value by exactly148286847/256; no omitted factor accounts for that decrease.

For t11,\ldots,18 the greatest s\textless=168 satisfying the displayed exponent\textless=−4096 is respectively98,87,76,65,54,43,31,20. R050 had already charged all s from t through its old endpoints59,48,37,26,15 at t11 through15. Subtracting only these sets gives exactly the fresh table in section1. The old numerical charges are not subtracted. At t16 through18 the old table was empty.

The complete endpoint exponents of the eight new intervals, in table order, have numerator

\[
 (-2974277,-2597701,-2221125,-1844549,
   -1467973,-1091397,-4539973,-4163397)/256.
 \tag{FV17}
\]

The largest is−1091397/256\textless−4263. Positivity of the s slope proves the same bound throughout each interval. All241 cells therefore satisfy

\[
             \sum_{\text{fresh cells}}U_{s,t}
                    <241\,2^{-4263}<2^{-4255}.
 \tag{FV18}
\]

There is no second union over pencils or types: those multiplicities are already in FV13. The displayed upper estimate has no qualifying s\textgreater=t at t\textgreater=19; this is only a limitation of this estimate. It does not imply a positive contribution or the existence of an unremoved rational matrix.

DO074 excludes t\textless=10 for globally nonunipotent FV1 terms. The old R050 intervals and the newly proved intervals then cover every possible t at each support through18. In particular at support16 both formerly missing types t15 and t16 are now covered. Other A classes require their separate already checked or future counts.

\subsection{6. Chronology, sources and verification boundary}\label{chronology-sources-and-verification-boundary}

The script reads the actual current R06159-entry all-top and31-entry uniform prior. It also checks that the current prior is exactly the historical DO08257-entry all-top list followed, in the recorded integration order, by DO0831953 and R06178919. Appending only the present4255 gives60 all-top entries, with the complete exact dyadic sum below2\^{}(-230). The uniform list is unchanged and its full sum is checked separately. This staging check does not claim that the present finding has already been canonically integrated.

The exact dependency reads are K069(IC9--15), the full R050 proof, LIFT063's original-ring covering theorem and metric addendum, and the actual H133/K068 residual. The new argument is FV3--11: the original rank-one chart has one small Frobenius ball, so the explicit ring covolume can replace the old independent-entry packing. The Green, pencil, source fibre and nonunipotence tests are retained in their original scopes rather than re-proved or enlarged.

The companion certificate verifies the original field/prime constants, the exact factorial ceiling, the complete rational constant, every new table cell and threshold, and full dyadic sums. Exact integral rank-one charts test the nonzero-pivot reconstruction and distinct-entry Frobenius inequality, including nondiagonal pivots. These diagnostics do not enumerate the actual rational kernel or automatically prove the geometric estimates. The inherited primary arithmetic anchors for the actual ring are Milne(2020), Theorem6.4 and Proposition4.26, and Thaine(2005), section2, through the checked LIFT063 proof; no new source classification theorem is imported.

The theorem removes only the indicated nonunipotent F,r1 tails. It gives no fixed-W0 unipotent statement, no full support-floor theorem by itself, and no signed positivity, prescribed Bessel/Mram value, selected arithmetic-image content, literal Host/FamP construction, DP44 sign or RH conclusion.

\clearpage

\section{AM-R-062 --- complete all-top support fifteen through seventeen exhaustion}\label{am-r-062-complete-all-top-support-fifteen-through-seventeen-exhaustion}

Author: GPT-6 Astra Pro; inherited effort not externally exposed; ALP\_MACHINE \nolinkurl{/root}. 22 November 2026 UTC. Complete new proof; author and independent whole checks are recorded separately before promotion.

\subsection{1. Exact kernel and assertion}\label{exact-kernel-and-assertion}

Keep the actual J56 all-top K044 upper-endpoint kernel, its permitted packet point, original coefficient representation and Hilbert metric, paired CM pencils, finite projectors, source class masses, phases, and real normalization. Write f=a=168, n=169, m=164 and use only the already charged H=133 archimedean cutoff of R048. Let s be the number of active relative T places. R060 has removed every support below15; the subsequent DO082, DO083 and R061 cuts and all their errors remain.

The new theorem gives the same normalized energy as

\[
 \mathcal E_{\rm top}=1+\operatorname{Re}\sum_{h\in R_{\rm top}^{(18)}}w(h)+\epsilon,
 \qquad |\epsilon|<2^{-230},\qquad s(h)\ge18.
 \tag{J18.1}
\]

This follows by removing all support15,16,17 terms using DO084 and the one fresh support17 cell below. The fresh charge of this finding alone is177258. DO084 contributes its separately proved4255 once. No previous charge is subtracted or repeated. The independent uniform H8192 support floor remains11.

The sufficient signed inequality

\[
 \operatorname{Re}\sum_{h\in R_{\rm top}^{(18)}}w(h)>-1+2^{-230}
 \tag{J18.2}
\]

remains open. In particular this theorem computes no prescribed global period or actual DP44 cost.

\subsection{2. A full-column support17 count at finite rank at least three}\label{a-full-column-support17-count-at-finite-rank-at-least-three}

Consider every actual support17 matrix for which the reduction of N=h−I has rank at least3 at each of the334 unmixed G factors. All these reductions are unipotent by the retained K068 cut. No global nilpotence, condition on A=h\_SS, finite rank-pattern choice, or type profile is required.

The Green bound K067(PS5--9)/DO064(PM10), with r=3, gives

\[
 |a(h)|\le 2Q^{-501(f-1)},\qquad
 Q=g_0g_1=1427908193729121522601,
 \quad501=3n-3\cdot4/2.
 \tag{J18.3}
\]

Both mixed G factors retain modulus at most one. The single factor2 already bounds the complete Euler correction at all unmixed factors. Even when a finite rank exceeds33, the first three factors of the inverse Green product suffice for J18.3; no upper-rank hypothesis is needed for this envelope.

For each fixed support and actual ordered pencil pairs, the original source contains pi e\_n and the target is M\_d.~Thus v=h e\_n belongs to pi\^{}−1 M\_d.~The original W0 shear acts trivially on S, so

\[
 x=v_S\in I_\Sigma^{164},\quad N_EI_\Sigma=T^{-2s},
 \qquad v_{W_0}\in B_\Sigma\supset\mathcal O_E^5,
 \quad[B_\Sigma:\mathcal O_E^5]=T^{10s}.
 \tag{J18.4}
\]

These are precisely R051 §1 and K069(IC12--14), at both CM partners and for every incidence type. They hold independently of A and count every target, including zero. For a fixed full negative column the original source-lattice stabilizer fibre is bounded by G\_n\textless2\^{}5022. Counting a full column is sufficient; no count of A is multiplied in.

Use the half-trace Euclidean metric and

\[
 D_0=7^{140}29^{144}5^{84}13^{84},\qquad
 V_d(C)=\frac{(528C^2)^{fd}}{(fd)!D_0^d}.
 \tag{J18.5}
\]

For z=(N\_EI\_\Sigma)\^{}(1/(2f))=T\^{}−s/f, R051 §2's arbitrary-ideal transference argument gives covering radius less than9750 sqrt(f)z for I\_\Sigma\^{}164. The positive part of a negative unit column satisfies \textbar\textbar x\_tau\textbar\textbar\textless H/2=66.5 at each embedding. At s17 the exact prime bound gives z\textless2\^{}−27 and9750·2\^{}−27\textless1/2. Attaching disjoint Dirichlet cells therefore gives

\[
 \#x\le V_{164}(67)T^{328s}<2^{-45338}T^{328s}.
 \tag{J18.6}
\]

This uses the actual fractional ideal and its covolume. It does not replace its covering radius by that of the original ring.

For the W0 coordinates, the full negative column has norm less than95 sqrt(f). LIFT063 and its metric addendum give mu(O\_E)\textless57 sqrt(f). Consequently the product ring O\_E\^{}5 has covering radius less than57 sqrt(5f)\textless128 sqrt(f). The containing target over-lattice B\_\Sigma cannot have a larger covering radius. Its exact half-trace covolume is D\_0\^{}5 T\^{}−10s. Full-ball Dirichlet comparison, without fixing a real coordinate or using an affine coarea slice, yields

\[
 \#v_{W_0}\le V_5(223)T^{10s}<2^{5757}T^{10s}.
 \tag{J18.7}
\]

The two coordinate sets are bounded independently. Their product is a valid upper count even though a true column also satisfies the Hermitian equation.

\subsection{3. Complete exponent and single fresh charge}\label{complete-exponent-and-single-fresh-charge}

The full support choices contribute binom(f,s)\textless=2\^{}168. The original pair/absolute-incidence sum, including all types at each active position, is P\_T\^{}s with P\_T=2(T−1)\^{}2/(T−2)\textless2\^{}997. There is no additional type multiplicity. The real factor is bounded by one. Combining J18.3, J18.6 and J18.7 with the complete fibre gives

\[
 U_{17}<2^{1+5022+168-45338+5757}
        T^{338\cdot17}2^{997\cdot17}Q^{-501\cdot167}.
 \tag{J18.8}
\]

With the checked strict rational bounds log2(T)\textless1019333/1024 and log2(Q)\textgreater287843/4096, its exponent is

\[
 E_{17}=1+5022+168-45338+5757
 +(338\cdot1019333/1024+997)17
 -501\cdot167\cdot287843/4096
 =-726048945/4096<-177258.
 \tag{J18.9}
\]

Thus the entire class in §2, with every original pair, profile and A, contributes less than2\^{}−177258. Remove its intersection with the current remainder. There is exactly one new charge177258; no union over A, ranks or endpoints remains. Applying this same estimate at support18 gives a positive exponent and proves no bound there.

\subsection{4. Exhaustive structural tables at supports15,16,17}\label{exhaustive-structural-tables-at-supports151617}

The unchanged K068/K069 structural partition has r=rank\_E(A−I)\textless=33. For A!=I its two alternatives are the fixed-space class F and the exceptional isotropic-cross class IC8. The following table covers every such term at the three stated supports.

\begin{longtable}[]{@{}
  >{\raggedright\arraybackslash}p{(\columnwidth - 2\tabcolsep) * \real{0.5000}}
  >{\raggedright\arraybackslash}p{(\columnwidth - 2\tabcolsep) * \real{0.5000}}@{}}
\toprule\noalign{}
\begin{minipage}[b]{\linewidth}\raggedright
Actual A!=I class
\end{minipage} & \begin{minipage}[b]{\linewidth}\raggedright
Complete removal at s15,16,17
\end{minipage} \\
\midrule\noalign{}
\endhead
\bottomrule\noalign{}
\endlastfoot
F,2\textless=r\textless=33 & K069(IC20--21) has minimum support threshold22 over these ranks. Its previously charged full union covers all three supports. \\
IC8,3\textless=r\textless=33, s15 or16 & DO064(PM22--23) has minimum threshold16 over these ranks. Its original full union applies. \\
IC8,3\textless=r\textless=33, s17 & DO064(PM9) proves rank(N mod w)\textgreater=r\textgreater=3 at every unmixed G factor. Hence the entire class lies in §§2--3, with one177258 charge. \\
IC8,r1 or2 & R060(J15.3--16) removes the full classes through17, without global unipotence, with its retained172255 union charge. \\
F,r1,globally unipotent & R060(J15.14--16) removes both possible full ranks1 and2 and every profile through17. The same172255 charge is retained once. \\
F,r1,nonunipotent & DO074 excludes t\textless=10. R050 §5 and DO084(FV13--18) together cover every11\textless=t\textless=s through18; in particular the new t15,16,17 tails cover their former gaps. DO084's separate4255 charge is retained once. \\
\end{longtable}

For the last row, at s15 the old R050 endpoints already cover t11,\ldots,15. At s16 the old t11,\ldots,14 ranges and DO084 t15/16 apply. At s17 the same old ranges and DO084 t15/16/17 apply. No favorable contribution or failed upper bound is being used to exclude a class.

For A=I, either h fixes S pointwise, so h=I\_S direct-sum h0 on W0, or it is the scalar-cross case of R059. The full complementary table is

\begin{longtable}[]{@{}
  >{\raggedright\arraybackslash}p{(\columnwidth - 2\tabcolsep) * \real{0.5000}}
  >{\raggedright\arraybackslash}p{(\columnwidth - 2\tabcolsep) * \real{0.5000}}@{}}
\toprule\noalign{}
\begin{minipage}[b]{\linewidth}\raggedright
Actual A=I class
\end{minipage} & \begin{minipage}[b]{\linewidth}\raggedright
Complete removal at s15,16,17
\end{minipage} \\
\midrule\noalign{}
\endhead
\bottomrule\noalign{}
\endlastfoot
Fixed S,nonunipotent & DO067 excludes t\textless=9; K081/K082 cover t10,\ldots,17 at these supports. Their corresponding support ceilings are168,164,149,134,119,104,89,73, all at least17. \\
Fixed S,globally unipotent,full rank1 & DO083 removes every profile at supports15..20 with its retained1953 charge. \\
Fixed S,globally unipotent,full rank2 & R058 removes every profile at supports13..23 with its retained39409 charge. \\
Scalar cross,globally unipotent or not & R059 removes the entire class through21 with its retained23966 charge. R061's later support22 charge is retained too. \\
Full rank0 & h=I has support0 and is absent here. \\
\end{longtable}

Signature-one global unipotence permits only the stated nonzero ranks. The two tables therefore exhaust all original terms at supports15,16,17. R060 already exhausted every smaller support. Their removal proves that the remainder in J18.1 has support at least18.

\subsection{5. Exact chronological error and remaining main questions}\label{exact-chronological-error-and-remaining-main-questions}

Start from the actual R061 array in \nolinkurl{checks/r061_joint_remainder_error.json}:59 all-top exponents and31 uniform exponents. Append DO084's4255 once and the present177258 once, in that order. The resulting61-entry all-top list has one exponent231 and every other exponent at least685, so

\[
 \sum_e2^{-e}\le2^{-231}+60\cdot2^{-685}<2^{-230}.
 \tag{J18.10}
\]

The31-entry uniform list is unchanged and has one231 with every other exponent at least797. Its complete sum is also below2\^{}−230. The companion certificate verifies both complete dyadic sums with integers. All previously removed-set overlaps are handled by intersection with the current remainder, without deleting historical absolute-error charges. Combining these sums with the unchanged energy identity proves J18.1.

Q1/Q2 have a strictly smaller finite/real sufficient gate, namely J18.2 over support\textgreater=18 with all other retained cuts. The actual mixed-CM source, prescribed Bessel/Mram period, full arithmetic ties and Read, selected arithmetic-image content, concrete032 pairing/value and every changed-family DP44 term remain unproved. This is neither a shorter valid chain nor a universal negative cost difference. Q3 has no Kummer-Hecke RH proof. Q4's global arithmetic charge comparison is unchanged.

\subsection{6. Algebraic-geometric language and exact carried locators}\label{algebraic-geometric-language-and-exact-carried-locators}

The count uses original integral Hermitian modules and their labelled theta correspondences. It changes neither their schemes, coefficient lattices, polarizations, Picard/biextension data nor finite level. The finite congruence condition is a rank condition on the reduction of the original endomorphism N; DO064 supplies its determinantal comparison with A−I. Counting full negative columns uses the genuine source stabilizer, with no moving lattice basis or new integral inverse.

The carried module/stabilizer context is SGA3 I4.7.3/4.7.3.1 and the exterior/determinantal context SGA6 VI Theorem3.3. The Picard and biextension anchors are FGA232 Theorem3.1 pp148--149; FGA236 Corollaries1.7/2.3 pp229/232; SGA7 I VII2.7/2.9 pp169--174 and VIII7.4(b) pp310--311. Formal groups, crystals and Cartier duality retain cote4 scan003 §§8.1.1--8.1.2, scans005--006 §§8.4.1--8.4.3 and cote54 scans005--007. These locators provide the inherited context, not a source for this new numerical bound. No new lambda-ring/K-theory identity, l-adic sheaf or weight theorem, CM motive comparison or motivic L-function statement is inferred. DP\_MULT Q3e and THETA3 G1/A7 gain precisely the smaller sufficient signed gate. CX\_GRH CX-ALG-024 and GCT\_PLETH T-052 retain their distinct open arithmetic comparisons.

Exact used results: K067 PS5--9; K068 original residual; K069 IC1--8,IC12--15,IC20--21; DO064 PM9--10,PM22--23; R051 §§1--2; LIFT063 §§3--5 and metric addendum RM1--2; R050 §5; DO067/DO074 stated nonunipotence criteria; K081/K082 full type tables; R058; R059; R060; DO083; DO084. Their already checked proofs are cited, not re-derived. The new certificate checks exact prime logarithms, the two volume ceilings, the complete new exponent, every structural table index and both full historical sums. It does not enumerate the rational kernel or automatically prove the structural partition. No Lean build or archive audit is used.

\clearpage

\section{AM-DO-085 --- minimal finite rank forces paired congruences on the untouched exceptional coordinates}\label{am-do-085-minimal-finite-rank-forces-paired-congruences-on-the-untouched-exceptional-coordinates}

\textbf{Canonical checked version --- 2026-11-22T19:29:33.467934+00:00.} Root and independent whole proof/script checks and separate976-assertion replays passed, including15,653 finite projection models and5,299 kernel tests. Original pending wording is historical. In the actual IC8 class at block ranks3..31, equality of full and compressed finite rank forces every untouched original column and the paired row to be fixed. The resulting paired ideal and exact Green jump jointly improve both endpoint counts. All182 listed fresh rank/support cells contribute\textless2\^{}-7548, including the full rank3 class through support23. Every original incidence, source fibre and error remains; the complete all-top support floor stays18.

\textbf{Algebraic-geometric translation.} This is a determinantal rank statement on the reductions of the original integral Hermitian module, coupled to counting its original-coordinate columns. The carried module/determinantal anchors are SGA3 I4.7.3/4.7.3.1 and SGA6 VI3.3. The Picard/biextension context is FGA232 Theorem3.1 pp148--149; FGA236 Corollaries1.7/2.3 pp229/232; SGA7 I VII2.7/2.9 pp169--174 and VIII7.4(b) pp310--311. Formal groups/crystals/Cartier duality retain cote4 scan003 §§8.1.1--8.1.2, scans005--006 §§8.4.1--8.4.3 and cote54 scans005--007. No new scheme, lambda-ring/K-theory identity, l-adic sheaf or weight theorem, CM motive comparison or motivic L-function statement is inferred. DP\_MULT Q3e and THETA3 G1/A7 gain exactly these smaller remaining sets; CX\_GRH CX-ALG-024 and GCT\_PLETH T-052 retain their distinct global arithmetic gaps.

Author: ALP\_MACHINE /root/definite\_oscillator, 22 November 2026 UTC. New controlled pending proof. Final author, exact diagnostic and independent whole checks are recorded separately.

\subsection{1. Exact sector and assertion}\label{exact-sector-and-assertion}

Keep the original J56 all-top signed kernel after the already charged H=133 cutoff, the actual K068 residual, every original finite projector, both mixed factors, all ordered T-pencil incidences, source class masses, CM phases and real Hilbert normalization. Put \[
 f=a=168,\qquad n=169,\qquad m=164,\qquad
 E^n=S\perp W_0,\quad \operatorname{sig}(W_0)=(4,1).
 \tag{MR1}
\] The determinant-one representative h is integral away T, A=h\_SS is globally integral, and every one of the334 unmixed G reductions of h is unipotent.

This report concerns exactly the exceptional class IC8 of K069, with \[
              r=\operatorname{rank}_E(A-I),\qquad 3\le r\le31.
 \tag{MR2}
\] No global unipotence is assumed. Let s count active relative T places. Let R\_r be the already charged DO064 support threshold at rank r. The following table gives a new threshold R'\_r:

\begin{longtable}[]{@{}rrr@{}}
\toprule\noalign{}
r & Old R\_r & New R'\_r \\
\midrule\noalign{}
\endhead
\bottomrule\noalign{}
\endlastfoot
3 & 16 & 23 \\
4 & 22 & 28 \\
5 & 27 & 34 \\
6 & 33 & 40 \\
7 & 39 & 45 \\
8 & 44 & 51 \\
9 & 50 & 56 \\
10 & 55 & 62 \\
11 & 61 & 67 \\
12 & 66 & 73 \\
13 & 72 & 78 \\
14 & 77 & 84 \\
15 & 82 & 89 \\
16 & 88 & 94 \\
17 & 93 & 99 \\
18 & 98 & 105 \\
19 & 103 & 110 \\
20 & 109 & 115 \\
21 & 114 & 120 \\
22 & 119 & 125 \\
23 & 124 & 130 \\
24 & 129 & 135 \\
25 & 134 & 140 \\
26 & 139 & 145 \\
27 & 144 & 150 \\
28 & 149 & 155 \\
29 & 154 & 160 \\
30 & 159 & 164 \\
31 & 163 & 168 \\
\end{longtable}

\textbf{Theorem.} The entire absolute contribution of the182 indicated cells \[
       3\le r\le31,\qquad R_r<s\le R'_r
\] is less than \[
                              2^{-7548}.
 \tag{MR3}
\] Remove only their intersections with the current remainder and retain every earlier charge. Some portions, including the rank-three support17 cell, may already have been removed by later findings; the theorem does not assert those intersections nonempty. The182-cell bound is a conservative containing union.

This removes IC8 rank three through support23, for example. IC8 ranks1/2 and the other structural classes keep their separate scopes. In particular no full support-floor update follows from this finding alone.

\subsection{2. The original exceptional block and its finite rank}\label{the-original-exceptional-block-and-its-finite-rank}

Use the already proved DO064(PM5--9). The actual A is partial monomial, every nonzero entry lies in mu28, and A is the identity outside at most2r original coordinate vertices, in both its rows and columns. At every actual G component, \[
                    \operatorname{rank}(\overline A-I)=r.
 \tag{MR4}
\] This is the exact weighted path/cycle rank statement, using injective reduction of mu28 at the primes in question. No archimedean numerator or T denominator is involved.

For each A choose2r original vertices containing its nonidentity support, padding once if necessary. The untouched coordinate space S\_free has \[
                           k=m-2r
 \tag{MR5}
\] basis vectors. For each of them, A e\_j=A\^{}\dagger e\_j=e\_j. A valid upper count for all the block and coordinate choices is \[
                  \mathcal A_r=\binom m{2r}29^{4r^2}.
 \tag{MR6}
\] Every block entry is zero or one of28 roots of unity. This counts many matrices which are not partial monomial or do not occur; that only enlarges the bound. No further class of rational coordinate frames is chosen.

\subsection{3. The new finite-rank lemma and both CM components}\label{the-new-finite-rank-lemma-and-both-cm-components}

Here is the elementary linear statement that improves the previous uniform-rank estimate.

\textbf{Lemma.} Let N be an endomorphism of a vector space V=S direct-sum W over any field. If \[
                 \operatorname{rank}N
                    =\operatorname{rank}(N_{SS})=r,
\] then projection to S is injective on im(N). Consequently N v=0 for every v in S with N\_SS v=0.

\textbf{Proof.} The image of the map P\_S N has dimension at least r because its restriction to S already has rank r. It has dimension at most rank(N)=r. Thus its restriction P\_S:im(N) to S has rank r on the r-dimensional domain im(N), and is injective. If N\_SS v=0, then Nv lies in this image and projects to zero, so Nv=0. This proves the assertion without a nilpotence or positivity assumption. \(\square\)

Apply the lemma at an unmixed G component where \[
                        \operatorname{rank}(\overline h-I)=r.
 \tag{MR7}
\] By MR4 the compressed rank is also r. For every original untouched S basis vector, the lemma gives (\overline h-I)e\_j=0.

The conjugate CM component has the same full rank: its matrix is the inverse adjoint, and rank(h\^{}\{-1\}-I)=rank(h-I). Its compressed rank is again r by MR4. The same argument gives fixed columns there on the same untouched original coordinates. A fixed column is fixed also by the inverse; taking the paired adjoint turns that statement into a fixed row at the opposite component. Hence the negative column satisfies simultaneously \[
 (h e_n)_j\equiv0\pmod{\mathfrak p},\qquad
 (h e_n)_j\equiv0\pmod{\bar{\mathfrak p}}
               \quad (j\in S_{\rm free}).
 \tag{MR8}
\] The deduction uses the two actual original-component matrices. It does not use positivity over a finite field, an integral orthogonal basis of a moving space, or a divided isotropic vector.

Let J contain both E-primes at precisely those unmixed relative G factors with full finite rank r. Then \[
        N_EJ=P_r^2,\qquad 1\le N_EJ\le Q^{2(f-1)},
 \tag{MR9}
\] where P\_r is the product of their residue sizes and Q=g0g1. There are at most2\^{}334 choices of this rank-r-versus-higher pattern. No rank-one orientation is required.

\subsection{4. Exact Green loss}\label{exact-green-loss}

Write \[
 L_r=rn-\frac{r(r+1)}2,\qquad
 d_r=L_{r+1}-L_r=n-r-1=168-r.
 \tag{MR10}
\] Every unmixed full rank is at least r by MR4. Those outside J have rank at least r+1. The retained unipotent Green formula K067(PS5--9), used with its first r or r+1 inverse factors, gives \[
 |a(h)|\le2Q^{-L_{r+1}(f-1)}P_r^{d_r}
        =2Q^{-L_{r+1}(f-1)}(N_EJ)^{d_r/2}.
 \tag{MR11}
\] The factor2 is the single global Euler correction for all334 factors. Here r+1\textless=32, so every retained denominator index is at least137. The same accepted product bound applies uniformly, even when an actual finite rank is larger. Both mixed projectors keep modulus at most one. We use only \textbar c(h)\textbar\textless=1 for the original real normalized trace.

Thus the finite-rank pattern produces both a Green loss and the paired original-coordinate divisibility MR8. Counting these together is the new step.

\subsection{5. Actual negative-column lattices and their volumes}\label{actual-negative-column-lattices-and-their-volumes}

Fix r, the original exceptional coordinate set and A, support, all ordered pencil pairs, and J. The generic actual source/target containment is R051 section1 and DO064(PM11--18): with I\_Sigma the paired inverse-prime ideal of norm T\^{}\{-2s\}, \[
 \begin{split}
 (h e_n)_{S_{\rm free}}&\in (J I_\Sigma)^k,\\
 (h e_n)_{S\setminus S_{\rm free}}&\in I_\Sigma^{2r},\\
 (h e_n)_{W_0}&\in \mathcal B_\Sigma\supset\mathcal O_E^5,\quad
          [\mathcal B_\Sigma:\mathcal O_E^5]=T^{10s}.
 \end{split}
 \tag{MR12}
\] The W0-supported shear preserves the original S-coordinate projection. All these statements hold for every I/II/III pair, including zero target labels, without global unipotence. The W0 count below is a full complex-coordinate count, not a real-part coarea slice.

Put \[
 L=J I_\Sigma,\quad N_EL=z^{2f},\quad
 z_0=T^{-s/f},\quad z_1=Q^{(f-1)/f}T^{-s/f}.
 \tag{MR13}
\] Thus z0\textless=z\textless=z1. In the original half-trace metric covol(L)=D0 z\^{}\{2f\}, with D0=7\^{}\{140\}29\^{}\{144\}5\^{}\{84\}13\^{}\{84\}. The arbitrary-ideal transference result DO070, with its actual dual normalization, gives \[
 \mu(L^k)<(2250/7)\sqrt{kf}\,z<4200\sqrt f\,z.
 \tag{MR14}
\] This is the fractional-ideal bound, not the smaller ring-only cover.

The positive part of a negative unit column has norm below H/2=66.5 at each embedding. Set \[
                  V_d(C)=\frac{(528C^2)^{fd}}{(fd)!D_0^d},
                  \qquad 528=(22/7)f.
 \tag{MR15}
\] Dirichlet-cell comparison gives the free-coordinate count \[
                V_k(66.5+4200z)\,z^{-2kf}.
 \tag{MR16}
\] Every cell in MR3 has s\textgreater=17, so the exact actual-prime bounds at s12 imply z0\textless2\^{}\{-71\} and z1\textless1/2 throughout the range. The other2r positive coordinates have enlarged radius below67 sqrt(f), giving V\_\{2r\}(67)T\^{}\{4rs\}. For W0 the literal ring cover of LIFT063 is below128 sqrt(f); the full negative-column radius is below95 sqrt(f). Therefore the full W0 count is at most \[
                         V_5(223)T^{10s},
                 \qquad V_5(223)<2^{5757}.
 \tag{MR17}
\] The paired determinant index T\^{}\{10s\} and the containing O\_E\^{}5 hold uniformly for every target. No type-dependent zero-target correction is relevant to this full-coordinate count.

For each entire negative column, retain the original source-lattice stabilizer fibre G\_n\textless2\^{}\{5022\}. The scalar-cross rank-five fibre is not used. Counting the free, remaining positive and W0 coordinates separately enlarges the true Hermitian column set and is valid. The chosen A count may also overcount fibres, which is harmless.

\subsection{6. Convex endpoint calculation and exact constants}\label{convex-endpoint-calculation-and-exact-constants}

Multiplying MR16 by the J factor in MR11 gives, apart from T\^{}\{d\_r s\}, \[
                 F_r(z)=V_k(66.5+4200z)\,
                         z^{(d_r-2k)f}.
 \tag{MR18}
\] Its logarithm is convex in log z: the only nonlinear term is 2kf times log(66.5+4200 exp(v)), with v=log z. Thus the maximum on {[}z0,z1{]} is at an endpoint. This proves a bound for every actual ideal pattern; it asserts neither continuous density nor occurrence of an endpoint ideal.

At z0 the radius is below67. The weighted free count has T-power2ks. After the other coordinates are included, the full T-power is \[
                         2k+4r+10=338.
 \tag{MR19}
\] The Green denominator is Q\^{}\{L\_\{r+1\}(f-1)\}.

At z1 the radius is below2167, because66.5+4200/2=2166.5. Keep its exact power, rather than using a lower bound on z1. The Q-power is \[
  Q^{-[L_{r+1}+2k-d_r](f-1)}
            =Q^{-[L_r+2k](f-1)},
 \tag{MR20}
\] and the full T-power is again338s. No lower bound for z1 at large support is needed.

Let b\_r be the least strict integer upper exponent for the exact count A\_r, and let v\_d(C) denote the exact strict integer upper exponent for the rational volume V\_d(C). Put \[
 C_r=b_r+334+168+5022+1+v_{2r}(67)+5757,
 \quad L_T=1019333/1024,\quad L_Q=287843/4096.
 \tag{MR21}
\] The full two endpoint exponents are \[
 \begin{split}
 E_{r,0}(s)&=C_r+v_k(67)+(338L_T+997)s
                      -L_{r+1}(f-1)L_Q,\\
 E_{r,1}(s)&=C_r+v_k(2167)+(338L_T+997)s
                      -(L_r+2k)(f-1)L_Q.
 \end{split}
 \tag{MR22}
\] Here334 counts all rank patterns,168 counts all supports,5022 retains the full source fibre,1 is the global Green constant, and997s bounds the complete actual ordered-pair/absolute-incidence product P\_T\^{}s. This product already contains every type. There is no extra type multiplicity, rank orientation or endpoint-union factor.

Every bound in MR21--22 is proved by integer powers and factorials in the companion certificate. In all29 ranks E\_\{r,0\} exceeds E\_\{r,1\}; both endpoints are nevertheless checked. The new threshold R'\_r is the greatest s\textless=168 for which their maximum is at most−4096. The positive common slope then proves the estimate throughout the specified interval. The old thresholds are read from the literal DO064 table and checked against its exact formula.

There are182 cells R\_r\textless s\textless=R'\emph{r. Their largest exact exponent is \[
          E_{14,0}(84)=-30953339/4096<-7556.
 \tag{MR23}
\] For example the rank-three endpoint is E}\{3,0\}(23)=-168236325/2048. The full interval union is bounded by \[
                        182\,2^{-7556}<2^{-7548},
 \tag{MR24}
\] which proves MR3. Neither endpoint max nor actual type choices incur an additional union factor.

\subsection{7. Chronology, verification and limits}\label{chronology-verification-and-limits}

The companion reads the actual R06261-entry all-top prior and31-entry uniform prior. It appends7548 once to the former, retaining every existing charge and leaving the latter unchanged. Both complete integer dyadic sums remain below2\^{}\{-230\}. This is a staging check, not a claim that the present finding has already been integrated. The existing all-top support floor18 is unchanged by this report alone.

The new proof is MR7--8 and its simultaneous use with MR11 in the original-coordinate count. The partial-monomial classification and exact primewise ranks are cited from DO064, without re-proving their settled root-unit step. The actual pencil, volume and fibre arguments are carried from R051, DO070 and the checked LIFT063 metric theorem. The Green input is Gorodetsky and Hazan (2017), Theorem2.1/equation(2.1), through K067; the arbitrary-ideal transference input is Miller and Stephens-Davidowitz (2019), Theorem3.2/equation(3.3), through DO070.

The certificate verifies the exact field/prime constants, every factorial ceiling, the two endpoints, the literal old thresholds, all182 cell inequalities, and complete chronological sums. Small finite-field matrix calculations test the characteristic-free projection lemma; they do not enumerate the original rational kernel or replace its symbolic proof. No moving integral frame, new local projector, favorable sign, prescribed global period, selected arithmetic image, literal Host constructor, DP44 sign or RH conclusion is asserted.

\clearpage

\section{AM-DO-086 --- the complete rank-one fixed-space sector from an integral entry witness}\label{am-do-086-the-complete-rank-one-fixed-space-sector-from-an-integral-entry-witness}

\textbf{Canonical checked version --- 2026-11-22T19:37:51.965592+00:00.} Root and independent whole proof/source/script checks and separate187-assertion replays passed; their exact pre-checkpoint events are retained. Original pending wording is historical. Promoted scope: The full actual F class with block rank1 at every support18..94 has total absolute contribution\textless2\^{}-8622, without global-unipotence or type-III restrictions. One original nonzero integral entry controls paired finite rank-zero loss; DO084 original chart geometry and every genuine seven-coordinate/fibre/pair factor are retained. This removes the entire F,r1 class through94 together with the old complete floor18, but supplies no full support-floor theorem by itself.

\textbf{Algebraic-geometric translation.} The original integral Hermitian modules, coefficient lattices and labelled correspondences remain unchanged. One original integral endomorphism entry or an untouched rank-one isotropic image is counted; no moving integral basis or arithmetic theta-image inverse is introduced. The carried module/determinantal anchors are SGA3 I4.7.3/4.7.3.1 and SGA6 VI3.3. Picard/biextension anchors remain FGA232 Theorem3.1 pp148--149; FGA236 Corollaries1.7/2.3 pp229/232; SGA7 I VII2.7/2.9 pp169--174 and VIII7.4(b) pp310--311. Crystals/Cartier duality retain cote4 scan003 §§8.1.1--8.1.2, scans005--006 §§8.4.1--8.4.3 and cote54 scans005--007. These are contextual locators; no new scheme identification, lambda-ring/K-theory identity, l-adic sheaf or weight theorem, CM motive or motivic L-function assertion follows. DP\_MULT Q3e and THETA3 G1/A7 gain precisely the stated smaller sufficient remaining sets; all actual source, selected arithmetic image, prescribed period,032 value and DP44 costs remain open. CX\_GRH CX-ALG-024 and GCT\_PLETH T-052 retain their distinct arithmetic gaps.

Author: ALP\_MACHINE /root/definite\_oscillator, 22 November 2026 UTC. New controlled pending proof. Final author, exact replay and independent whole checks are recorded separately.

\subsection{1. Actual assertion without a global-unipotence restriction}\label{actual-assertion-without-a-global-unipotence-restriction}

Retain the unchanged J56 all-top kernel, the charged H=133 cutoff, f=a=168, n=169, m=164, and the full K068/K069 residual. In particular the original S block A is globally integral and all334 unmixed G reductions of h are unipotent. Both mixed projectors, every actual T-pencil incidence, original class masses, phases and real normalization are unchanged.

Let X=A-I and K\_A=ker(X) intersection ker(X\^{}*). Restrict to the exact fixed-space class \[
       \operatorname{rank}_E X=1,\qquad
                    h\text{ fixes }K_A\text{ pointwise}.
 \tag{FR1}
\] There is no assumption that h is globally unipotent or nonunipotent, and no restriction on the type-III count.

\textbf{Theorem.} The entire absolute contribution of FR1 at total supports \[
                              18\le s\le94
\] is less than \[
                              2^{-8622}.
 \tag{FR2}
\] Remove only the intersection with the current remainder and retain all prior charges. The theorem covers both the globally unipotent and nonunipotent parts of this class. It leaves all other structural classes separate; it does not alone increase the full support floor.

\subsection{2. One actual integral entry controls all finite rank-zero losses}\label{one-actual-integral-entry-controls-all-finite-rank-zero-losses}

Choose a nonzero entry x of the original integral rank-one matrix X. At every complex embedding, \[
                         |x_\tau|\le H+1<2H=266.
 \tag{FR3}
\] Consequently its absolute field norm satisfies \[
                         |N_{E/\mathbf Q}(x)|<(2H)^{2f}.
 \tag{FR4}
\] This is the literal integral element in the original field. No denominator clearing or moving lattice frame occurs.

Let P0 be the product of residue sizes over the unmixed relative G factors at which h-I has reduction rank zero. At each such factor, every entry of X is divisible by its prime; the actual inverse-adjoint companion has rank zero as well, so x is divisible at both CM primes. Since x is integral everywhere, all additional prime divisors contribute nonnegative valuations. Thus \[
                     P_0^2\le |N_{E/\mathbf Q}(x)|,
                     \qquad P_0<(2H)^f.
 \tag{FR5}
\] The same chosen global entry is used at every factor.

The retained Green formula has baseline rank-one decay q\^{}\{-a\}. Every nonzero reduction rank has at least that decay, with its accepted denominator correction; rank zero loses q\^{}a relative to that baseline. Over all334 unmixed factors the total Euler correction is less than two. Therefore \[
 |a(h)|\le 2Q^{-a(f-1)}P_0^a
           <2(2H)^{af}Q^{-a(f-1)},\qquad Q=g_0g_1.
 \tag{FR6}
\] The two mixed factors retain modulus at most one. This is exactly K068(LR10--12), p=k=1, and K069(IC16) in the present H133 subset, written with the actual integral entry and paired prime norms explicit. It applies before any assertion of global unipotence. The one entry is a witness for a pointwise bound; no choice of entries or finite rank patterns is summed.

\subsection{3. The original integral chart and unchanged source-column count}\label{the-original-integral-chart-and-unchanged-source-column-count}

The geometric estimate DO084(FV3--11) used only integrality and rank(X)=1. Its proof did not use the nonunipotence restriction in that report's final application. We retain that exact geometric statement: store the original nonzero pivot row and column, counting their common entry once. Their327 distinct integral entries lie in one Frobenius ball of radius134 sqrt(f). With the original-ring cover and actual half-trace covolume, \[
                   \#A<m^2\,2^{307590}.
 \tag{FR7}
\] The pivot reconstruction uses division only for uniqueness, never to transport an integral basis.

For each fixed A, K069(IC6,IC12--15) gives codimension at most two for K\_A in S. The genuine negative column h e\_n lies in K\_A's Hermitian orthogonal complement. Choose once two original S coordinates injective on its S part, padding with an unused original coordinate if necessary, and retain all five W0 coordinates. This is the original coordinate projection; no inverse minor is declared integral.

The actual projected target over-lattice contains O\_E\^{}7 and has paired index at most T\^{}\{14s\}, uniformly in every ordered source/target pencil pair. The shear acts only in W0 and therefore commutes with this projection. The old integral-coset bound is267\^{}\{14f\}. For a complete original column the genuine source-lattice stabilizer remains G\_n\textless2\^{}\{5022\}. These are the same K069 counts for every F matrix, independently of global unipotence: \[
               \#h\text{ for fixed pairs}
               <m^2\,2^{307590}G_n\,267^{14f}T^{14s}.
 \tag{FR8}
\] All matrix and class multiplicities have their original meaning. The estimate does not replace the full fibre by a smaller scalar-cross fibre.

Choose supports in binom(f,s) ways. The complete actual ordered-pair/absolute-incidence product at every active place is bounded by \[
                    P_T=\frac{2(T-1)^2}{T-2}<2^{997}.
\] It already includes all types I,II,III. There is no specified type-III subset to choose and no extra binom(s,t),3\^{}s or rank-pattern factor. Use only \textbar c(h)\textbar\textless=1 for the true real normalized trace. Combining FR6 and FR8 gives the complete support-s bound \[
 U_s<
  2\binom fs\,m^2\,2^{307590}G_n\,
 267^{14f}266^{af}T^{14s}P_T^sQ^{-a(f-1)}.
 \tag{FR9}
\]

\subsection{4. Exact support range and error}\label{exact-support-range-and-error}

Use the literal prime bounds \[
 \log_2T<1019333/1024,\qquad
 \log_2Q>287843/4096,
\] together with m\^{}2\textless2\^{}\{15\}, binom(f,s)\textless=2\^{}\{168\}, G\_n\textless2\^{}\{5022\}, and \[
             266<267,\qquad267^{4096}<2^{8\cdot4096+249}.
\] Then FR9 is strictly bounded by2\^{}\{C+bs\}, where \[
 \begin{split}
 C={}&1+168+15+5022+307590
        +(14+a)f(8+249/4096)-a(f-1)(287843/4096)\\
    ={}&-723122875/512,\\
 b={}&14(1019333/1024)+997=7645795/512.
 \end{split}
 \tag{FR10}
\] The positive slope gives the largest exponent on18\textless=s\textless=94 at s94: \[
                       C+94b=-4418145/512<-8629.
\] There are77 supports, so \[
                       \sum_{s=18}^{94}U_s
                         <77\,2^{-8629}<2^{-8622}.
\] This proves FR2 with no further union over global unipotence or types. At s95 the same displayed exponent is positive; that failed sufficient bound is not evidence of an actual remaining matrix or a positive summand.

The complete F rank-one sector at smaller supports has already been removed by the current support18 remainder. New subsets may overlap earlier individual-sector bounds at larger supports; only their intersection with the current remainder is removed and all historical errors remain.

\subsection{5. Dependencies, exact verification and scope}\label{dependencies-exact-verification-and-scope}

The proof uses the actual original integral-entry witness and pairwise rank-zero divisibility from K068(LR10--12)/K069(IC16), the source-coordinate and class-fibre count K069(IC6,IC12--15), and DO084's rank-one chart geometry FV3--11. The latter is explicitly applied only at the scope proved in its geometric sections; its previous nonunipotent coefficient estimate is not used. The Green input is the already checked K067 use of Gorodetsky and Hazan(2017), Theorem2.1/equation(2.1). No new primary classification theorem is assumed.

The companion reads the actual canonical R06261-entry all-top/31-entry uniform prior, checks the independently frozen DO085 staged list with7548 appended once, and stages only the present8622 after it. The resulting63-entry all-top list and unchanged31-entry uniform list have exact sums below2\^{}\{-230\}. This is an ordered staging check, not an assertion that either pending result has already been integrated.

The exact diagnostics check every prime bound, coefficient exponent and support endpoint, all77 complete support counts and the historical sums. Finite arithmetic examples test that the product of paired rank-zero prime factors divides the norm of the same integral entry; they do not enumerate the actual rational kernel or prove automorphic occurrence.

No IC8 bound, fixed-W0 estimate, scalar-cross theorem, sign of the remaining total, prescribed Bessel/Mram value, integral selected image, literal Host constructor, DP44 inequality or RH conclusion is inferred. The entire-class improvement is the removal of a nonunipotent coefficient-pole cost: the chosen original block entry is already integral.

\clearpage

\section{AM-R-063 --- untouched-coordinate cross blocks in the small exceptional class}\label{am-r-063-untouched-coordinate-cross-blocks-in-the-small-exceptional-class}

\textbf{Canonical checked version --- 2026-11-22T19:37:51.965592+00:00.} Root and independent whole proof/source/script checks and separate214-assertion replays passed; their exact pre-checkpoint events are retained. Original pending wording is historical. Promoted scope: The entire actual IC8 block-rank1/2 class at supports18..21 has total absolute contribution\textless2\^{}-52215. The cross map on160 untouched original coordinates has rank0/1; its zero/transverse/plane alternatives and nine retained negative-column coordinates reconstruct the whole column. The original source fibre counts every A without extra block enumeration. All original targets, paired primes, incidence weights and finite Green factors remain. This is an entire-class bound without global unipotence, not a standalone full support-floor theorem.

\textbf{Algebraic-geometric translation.} The original integral Hermitian modules, coefficient lattices and labelled correspondences remain unchanged. One original integral endomorphism entry or an untouched rank-one isotropic image is counted; no moving integral basis or arithmetic theta-image inverse is introduced. The carried module/determinantal anchors are SGA3 I4.7.3/4.7.3.1 and SGA6 VI3.3. Picard/biextension anchors remain FGA232 Theorem3.1 pp148--149; FGA236 Corollaries1.7/2.3 pp229/232; SGA7 I VII2.7/2.9 pp169--174 and VIII7.4(b) pp310--311. Crystals/Cartier duality retain cote4 scan003 §§8.1.1--8.1.2, scans005--006 §§8.4.1--8.4.3 and cote54 scans005--007. These are contextual locators; no new scheme identification, lambda-ring/K-theory identity, l-adic sheaf or weight theorem, CM motive or motivic L-function assertion follows. DP\_MULT Q3e and THETA3 G1/A7 gain precisely the stated smaller sufficient remaining sets; all actual source, selected arithmetic image, prescribed period,032 value and DP44 costs remain open. CX\_GRH CX-ALG-024 and GCT\_PLETH T-052 retain their distinct arithmetic gaps.

Author: \nolinkurl{/root}, GPT-6 Astra Pro, inherited effort not externally exposed, ALP\_MACHINE.22 November2026 UTC. \textbf{New complete draft; not yet canonically promoted.} Whole proof/script checks and exact replay remain separate requirements.

\subsection{1. Actual class and assertion}\label{actual-class-and-assertion}

Retain the unchanged J56 all-top kernel, charged H133 cutoff, original paired CM pencils, finite projectors, packet, coefficient metric, source masses and phases. Put f=168,n=169,m=164. In the original decomposition S perpendicular W0, signature(W0)=(4,1), consider the K069 exceptional class IC8 with r=rank\_E(A−I) in\{1,2\}, where A=h\_SS. No global unipotence is assumed. All334 unmixed G reductions are unipotent by K068.

The theorem bounds the entire class at supports18,19,20,21 by

\[
             \sum_{18\le s\le21,\ r\in\{1,2\},\ IC8}|w(h)|
                              <2^{-52215}.
 \tag{UC1}
\]

Only its intersection with the current remainder is removed. All preceding error charges remain and52215 is appended once after final checks. This is not a complete support-floor statement by itself. Its genuinely new step is to count the cross block on the many untouched original S coordinates, and use it to reconstruct most of the negative column.

\subsection{2. Original untouched coordinates and finite Green factor}\label{original-untouched-coordinates-and-finite-green-factor}

DO064(PM5--9), in the exact form DO075(CR3), makes A identity outside at most four original vertices in both rows and columns. Choose a padded four-subset B containing these vertices, once for each occurring A; denote its coordinate space by S\_B and its160-coordinate complement by S\_0. The total number of possible B is

\[
                         \binom{164}{4}<2^{25}.
 \tag{UC2}
\]

There is no separate A count below. The full negative column and the original source stabilizer will already bound all arrows, including their different A blocks.

DO075's argument immediately preceding CR5 is independent of its support12 application: a nonidentity partial-monomial A has a diagonal0 or a nontrivial root eta in mu28. The corresponding entry of A+A\^{}dagger−2I is−2 or(eta−1)\^{}2/eta, a unit at the actual odd G primes prime to28. A rank-zero or rank-one unipotent reduction N=h−I would have N²=0 and h+h\^{}−1−2I=0, contradicting that unit. Thus every unmixed finite rank is at least two. The original Green bound therefore gives, for the whole class,

\[
                 |a(h)|\le2Q^{-335(f-1)},
 \quad Q=g_0g_1=1427908193729121522601.
 \tag{UC3}
\]

This is one full-class bound. No finite-rank pattern, orientation or additional rank1/rank2 union is counted. Both mixed factors retain modulus at most one.

\subsection{3. The untouched cross block and exact column coding}\label{the-untouched-cross-block-and-exact-column-coding}

For each original e\_j in S\_0 write

\[
                    he_j=e_j+C_0e_j,\qquad C_0:S_0\longrightarrow W_0.
 \tag{UC4}
\]

There is no S\_B component because the corresponding original column of A is exactly e\_j. Unitarity gives C\_0\^{}*J\_0C\_0=0, where J\_0=diag(1,1,1,1,−1). Hence either C\_0=0 or rank\_E C\_0=1: after one faithful complex embedding its image is totally isotropic in a signature-(4,1) space and has dimension at most one, and rank is preserved by that embedding.

Let v=h e\_n, write v\_B and v\_W for its S\_B and W0 parts. Orthogonality of he\_j and he\_n gives

\[
                         v_0=-C_0^*J_0v_W.
 \tag{UC5}
\]

Thus C\_0,v\_B,v\_W determine the entire negative column. For every fixed actual source/target pair, at most G\_n\textless2\^{}5022 arrows have a given full negative column, by R051 §1. This source-lattice fibre is the original one, independent of A. Consequently a count of data determining C\_0 and these nine original column coordinates bounds all arrows after multiplying by G\_n; no enumeration of A is needed.

If C\_0=0, UC5 says v\_0=0 and the same coding works with no cross data. If C\_0 is nonzero, use the original packet integers of R059(GB4):

\[
 \begin{split}
 U&=77Q^2-14,\\
 u&=(2,95192390527,383092671043,12529843548438452388609),\\
 y_1&=(-976627,-153032,38026,0),\\
 y_2&=(-567678,695570,-172838,0),\\
 y_3&=(104583,-15309315176,-61610045224,2).
 \end{split}
 \tag{UC6}
\]

They satisfy u dot u=U, gcd(u\_i)=1, and the y\_j span u-perp in the positive W0 coordinates. Each y\_j belongs separately to every actual source and target integral selfdual lattice at both CM components, by DO078(FP20--22)/K088(IP5--6). Each original e\_i of S\_0 also belongs to both because the pencil shear is supported on W0.

\textbf{Transverse case.} At least one row delta\_i=\textless y\_j,C\_0e\_i\textgreater=\textless y\_j,he\_i\textgreater{} is nonzero. It lies in O\_E\^{}160: he\_i and y\_j both belong to the actual target integral lattice at every finite place. Fix a nonzero pivot delta\_i0 and z=C\_0e\_i0. Rank one gives

\[
                    C_0e_i=z\,\delta_i/\delta_{i_0}.
 \tag{UC7}
\]

The data delta,z reconstruct C\_0. At most three y choices and160 pivot choices are allowed. Division proves only uniqueness over E, not integrality of an inverse.

\textbf{Plane case.} All three rows vanish. Then uniquely C\_0e\_i=(u xi\_i,beta\_i). Isotropy implies that xi is nonzero; otherwise every column would be supported only on a negative line and hence zero. With any nonzero xi\_i0 and z=C\_0e\_i0, rank one reconstructs C\_0 by z xi\_i/xi\_i0. The integer Bezout relation for u gives xi\_i as an integral linear combination of the original positive cross entries.

Let I\_Sigma be the paired inverse-prime ideal at the active T places, with N\_E I\_Sigma=T\^{}−2s. Since e\_i lies in both source and target, he\_i−e\_i lies in the target, whose W0 projection lies inside pi\^{}−1 O\^{}5. The same statement holds at the conjugate component. Away T the entries are integral. Thus

\[
                   \delta\in O_E^{160},\qquad
                   \xi\in I_\Sigma^{160}.
 \tag{UC8}
\]

Unlike the earlier A=I argument, no common G ideal is required here: UC3 already supplies the rank-two denominator uniformly.

\subsection{4. The actual vector and target-lattice counts}\label{the-actual-vector-and-target-lattice-counts}

At each embedding \textbar\textbar C\_0e\_i\textbar\textbar\textless H=133, by the orthogonal Euclidean decomposition in UC4. R059(GB11--12) gives133²max\textbar\textbar y\_j\textbar\textbar²\textless2\^{}86, so every \textbar delta\_i\textbar\textless2\^{}43. The elementary original-ideal packing bound R059(GB14), applied to O\_E of norm1, gives

\[
                 \#\delta\le(1+2^{44})^{320f}<2^{45\cdot320f}.
 \tag{UC9}
\]

In the plane case, \textbar\textbar u xi\_i\textbar\textbar\textless133. The exact packet comparison U\textgreater133²2\^{}132 gives \textbar xi\_i\textbar\textless M=2\^{}−66 at every embedding. Set z\_s=T\^{}−s/f.~At s\textgreater=18, the actual prime comparison T\^{}18\textgreater2\^{}(66f) gives z\_s\textless M. Therefore the same fractional-ideal packing lemma implies

\[
 \#\xi\le(1+2M/z_s)^{320f}
                  <3^{320f}2^{-66\cdot320f}T^{320s}.
 \tag{UC10}
\]

This count includes zero coordinates and is uniform in the original ideal. It does not assume a principal generator, an integral isotropic line or an extra field containing sqrt(U).

The pivot z lies in the W0 projection of the true target M\_d.~The negative part v\_W lies in pi\^{}−1 M\_d,W0. The former is contained in the latter, so R051's full original-coordinate coset count, valid for every target including zero, gives separately

\[
                    \#z,\#v_W<2^{14861}T^{10s}.
 \tag{UC11}
\]

Both vectors have norm at mostH at each embedding. The target over-lattice contains O\_E\^{}5 with paired indexT\^{}10 per active place. This is a full complex count with no real-coordinate coarea slice.

The four entries v\_B lie in I\_Sigma\^{}4 by the actual source/target containment R051 §1. This over-lattice of O\_E\^{}4 has paired indexT\^{}(8s). Its per-coset ball count gives

\[
                   \#v_B\le267^{8f}T^{8s}<2^{10835}T^{8s}.
 \tag{UC12}
\]

The estimates count the three data blocks independently, enlarging the true Hermitian column set. In the zero-cross case there is no pivot z, but multiplying its bound, which exceeds one, remains a valid upper bound with the same common factors.

\subsection{5. Every multiplicity and the complete exponent}\label{every-multiplicity-and-the-complete-exponent}

Retain the support choices2\^{}168, padded four-subsets2\^{}25, at most160 pivots\textless2\^{}8, at most three transverse tests\textless2², original full fibre2\^{}5022, single Green factor2, two W0 counts and four positive-coordinate count. Using the same harmless factors for all three cases gives the common constant

\[
                    C=1+168+25+8+2+5022+2(14861)+10835=45783.
 \tag{UC13}
\]

The actual complete ordered-pair/absolute-incidence sum is P\_T\^{}s\textless2\^{}(997s), including all types and both CM components. The real coefficient has modulus at most one. There is no further type, A or finite-pattern multiplier. Put L\_T=1019333/1024 and L\_Q=287843/4096, with log2TL\_Q. From UC3 and UC9--12 the three entire case envelopes are

\[
 \begin{split}
 E_0(s)&=C+(28L_T+997)s-335\cdot167L_Q,\\
 E_A(s)&=E_0(s)+45\cdot320f,\\
 E_B(s)&=C-66\cdot320f+(8/5)320f
                       +(348L_T+997)s-335\cdot167L_Q,
 \end{split}
 \tag{UC14}
\]

where3\^{}5\textless2\^{}8 was used. Each increases with s. Exact rational evaluation at s21 gives

\[
 E_A(21)=-3523577099/4096,\qquad
 E_B(21)=-213891083/4096<-52219,
 \quad E_0(21)<E_A(21).
 \tag{UC15}
\]

All twelve case/support cells are consequently below2\^{}−52219, so their union is less than12·2\^{}−52219\textless2\^{}−52215, proving UC1. The four-subset choice already covers both original ranks simultaneously. Possible overcounting of a matrix through several subsets or pivots only increases this bound.

\subsection{6. Chronology, language and limits}\label{chronology-language-and-limits}

After its whole checks, this theorem appends exactly52215 to the actual complete chronological all-top array then current, retaining all intervening proved cuts and leaving the uniform array unchanged. Before promotion the companion must verify that complete sum. Remove only the intersection of the three cases with the actual remainder. The new theorem alone gives no full floor beyond18, because the F and fixed-W0 classes require their own checked estimates.

This is a statement about rank-one isotropic images of an original integral Hermitian-module map and its integral target pairings. The identity UC5 reconstructs coordinates over E; it is not an integral inverse or arithmetic theta-image retraction. The original coefficient lattices, correspondences, Picard data and polarizations remain unchanged. The carried module/determinantal anchors are SGA3 I4.7.3/4.7.3.1 and SGA6 VI3.3; Picard/biextension anchors are FGA232 Theorem3.1 pp148--149, FGA236 Corollaries1.7/2.3 pp229/232, SGA7 I VII2.7/2.9 pp169--174 and VIII7.4(b) pp310--311. Formal crystals/Cartier duality retain cote4 scan003 §§8.1.1--8.1.2, scans005--006 §§8.4.1--8.4.3 and cote54 scans005--007. These are contextual locators, not attributions of the new count to those texts. No new lambda-ring/K-theory identity, l-adic sheaf or weight theorem, CM motive or motivic L-function statement is asserted.

DP\_MULT Q3e and THETA3 G1/A7 gain only this smaller sufficient signed remainder. The signed inequality, fully admissible mixed-CM family, selected arithmetic image, prescribed global Bessel/Mram period, actual Read/032 pin/value and changed-family DP44 cost remain open. CX\_GRH CX-ALG-024 has no Hecke RH consequence, and GCT\_PLETH T-052 keeps its missing global arithmetic comparison. Exact carried mathematical inputs are DO064 PM5--9; DO075 CR3 and its pre-CR5 rank-two argument; R051 §1; R059 GB4--8,GB11--14,GB18--20; DO078 FP20--22 and K088 IP5--6. No Lean build or archive audit is performed.

\clearpage

\section{AM-R-064 --- complete all-top support eighteen through twenty exhaustion}\label{am-r-064-complete-all-top-support-eighteen-through-twenty-exhaustion}

Author: \nolinkurl{/root}, GPT-6 Astra Pro, inherited effort not externally exposed, ALP\_MACHINE.22 November2026 UTC. Complete new assembly proof. Whole proof/script checks are recorded separately before promotion.

\subsection{1. Same energy, strictly smaller remainder}\label{same-energy-strictly-smaller-remainder}

Keep the original J56 all-top K044 endpoint, permitted packet, selected compact representation and Hilbert norm, every paired CM incidence, finite projector, source class mass and phase. The actual dimensions are f=a=168,n=169,m=164. Use only the already charged H133 cutoff. R062 gives the exact energy identity with all remaining supports at least18. DO085, DO086 and R063 are now checked additional containing-set removals with all their absolute errors retained.

This theorem gives an actual remaining subset R\_top\^{}(21) such that

\[
 \mathcal E_{\rm top}
  =1+\operatorname{Re}\sum_{h\in R_{\rm top}^{(21)}}w(h)+\epsilon,
 \qquad |\epsilon|<2^{-230},\qquad s(h)\ge21.
 \tag{J21.1}
\]

It introduces no new numerical charge. Its content is the complete exhaustion of every original structural class at s18,19,20. The sufficient signed inequality

\[
             \operatorname{Re}\sum_{h\in R_{\rm top}^{(21)}}w(h)
                                      >-1+2^{-230}
 \tag{J21.2}
\]

remains open. The independent uniform H8192 support floor remains11.

\subsection{2. Exhaust every nonidentity positive compression}\label{exhaust-every-nonidentity-positive-compression}

Write A=h\_SS in the unchanged original decomposition S perpendicular W0. The actual residual has r=rank\_E(A−I)\textless=33. For A!=I, K069(IC1--8) partitions the original matrices into the fixed-space class F, in which h fixes K\_A=ker(A−I) intersect ker(A*−I) pointwise, and its exceptional isotropic-cross complement IC8. This partition is made over the original field; no further global-unipotence restriction is imposed.

\begin{longtable}[]{@{}
  >{\raggedright\arraybackslash}p{(\columnwidth - 2\tabcolsep) * \real{0.5000}}
  >{\raggedright\arraybackslash}p{(\columnwidth - 2\tabcolsep) * \real{0.5000}}@{}}
\toprule\noalign{}
\begin{minipage}[b]{\linewidth}\raggedright
Actual class
\end{minipage} & \begin{minipage}[b]{\linewidth}\raggedright
Complete containing-set bound at each s18,19,20
\end{minipage} \\
\midrule\noalign{}
\endhead
\bottomrule\noalign{}
\endlastfoot
F,r=1 & DO086(FR1--10) covers the entire class at every support18..94, with no type-III or global-unipotence restriction. Its single8622 charge is already retained. \\
F,2\textless=r\textless=33 & The original K069(IC20--21) table has minimum support threshold22 over these ranks. All three supports are inside its already charged full union. \\
IC8,r=1 or2 & R063(UC1--15) covers the complete class at every support18..21, including all zero/transverse/plane untouched cross maps and every incidence profile. Its single52215 charge is already retained. \\
IC8,3\textless=r\textless=31 & DO064's old ranges together with DO085(MR1--24) cover every support through at least23 at every rank in this range. DO085's smallest new ceiling is23 at rank3. Its single7548 charge is already retained. \\
IC8,r=32 or33 & DO064(PM22--23) already removes every support through168. \\
\end{longtable}

For the fourth row, DO085's fresh ranges start immediately after each old DO064 ceiling; there is no intervening gap. Later overlaps, including R062's support17 class, are handled by actual remainder intersection, not by removing a historical numerical charge. These five rows exhaust every possibility with A!=I and s in\{18,19,20\}.

\subsection{3. Exhaust the identity compression}\label{exhaust-the-identity-compression}

When A=I, the original unitary block identities give the exact dichotomy of R059(GB1--3): either h fixes S and is I\_S direct-sum h0 on W0, or its nonzero scalar-cross block has rank one. This does not imply that h itself is globally unipotent. Treat every alternative explicitly.

\begin{longtable}[]{@{}
  >{\raggedright\arraybackslash}p{(\columnwidth - 2\tabcolsep) * \real{0.5000}}
  >{\raggedright\arraybackslash}p{(\columnwidth - 2\tabcolsep) * \real{0.5000}}@{}}
\toprule\noalign{}
\begin{minipage}[b]{\linewidth}\raggedright
Actual A=I class
\end{minipage} & \begin{minipage}[b]{\linewidth}\raggedright
Complete containing-set bound at each s18,19,20
\end{minipage} \\
\midrule\noalign{}
\endhead
\bottomrule\noalign{}
\endlastfoot
Fixed S,nonunipotent & DO067 excludes t\textless=9. For t10,\ldots,20, K081/K082 cover every support from t to their respective combined ceilings168,164,149,134,119,104,89,73,58,43,28. Every one of these ceilings is at least20. Thus every possible t\textless=s is covered. \\
Fixed S,globally unipotent,full rank1 & DO083(RO1--18) covers every profile at supports15..20 with its retained1953 charge. \\
Fixed S,globally unipotent,full rank2 & R058(HB1--25) covers every profile at supports13..23 with its retained39409 charge. \\
Scalar cross,globally unipotent or not & R059(GB1--22) covers the entire class at supports13..21 with its retained23966 charge; R061 additionally covers22, with its charge retained. \\
Full rank0 & h=I has support0 and is absent at the three positive supports under consideration. \\
\end{longtable}

The accepted signature-one unipotent classification allows no other nonzero full rank for a globally unipotent h. The first row sums over all possible type-III counts; it does not require that they be small. The scalar-cross row has no global-unipotence hypothesis. Hence these five rows exhaust every matrix with A=I at supports18,19,20.

Combining both tables removes every actual term at these supports. R062 had already removed every smaller support. This proves that the remaining subset has support at least21. No assertion about complete support21 exhaustion is made.

\subsection{4. Exact retained errors and final identity}\label{exact-retained-errors-and-final-identity}

The actual chronological array immediately preceding this assembly is \nolinkurl{checks/r063_joint_remainder_error.json}. It contains64 all-top charges and31 uniform charges. The last five all-top exponents, in actual integration order, are

\[
                      4255,\quad177258,\quad7548,\quad8622,\quad52215,
 \tag{J21.3}
\]

from DO084, R062, DO085, DO086 and R063. None is deleted, combined with a smaller overlapping set or appended again by this theorem. The all-top list has exactly one231 and every other exponent at least685, so

\[
 \sum_{e\in\mathcal L_{\rm top}}2^{-e}
          \le2^{-231}+63\cdot2^{-685}<2^{-230}.
 \tag{J21.4}
\]

The unchanged uniform list has exactly one231 and every other exponent at least797, and its sum is below2\^{}−230. The companion verifies both complete sums with integers and verifies that this assembly changes neither list.

At every removal step, remove only the intersection of the new proved containing set with the actual current remainder. The removed subsets are then disjoint, while their conservative historical numerical charges need not be. The triangle inequality bounds the total discarded contribution by J21.4. The original energy identity, this error bound and the two exhaustive tables prove J21.1.

Q1/Q2 have a strictly smaller sufficient finite/real gate, now J21.2. The actual fully admissible mixed-CM family, prescribed nonzero global Bessel/Mram period, selected arithmetic-image content, full Read/comparison ties, specific032 pairing/value and every changed-family DP44 term remain unresolved. Neither a shorter valid chain nor a universal strict negative DP44 difference is proved. Q3 has no Kummer-Hecke RH conclusion, and Q4 retains its missing global arithmetic charge comparison.

\subsection{5. Algebraic geometry and verification boundary}\label{algebraic-geometry-and-verification-boundary}

The assembly combines rank conditions on the reductions of original integral Hermitian modules, the fixed-space/isotropic-image partition, and their original labelled theta correspondences. It changes no lattice, coefficient module, polarization, finite level, Picard scheme, biextension or CM realization. The root-coordinate reconstructions in the cited counts are over E and are not integral inverses of the arithmetic theta map.

The carried module/stabilizer context is SGA3 I4.7.3/4.7.3.1 and the exterior/determinantal context is SGA6 VI Theorem3.3. The Picard/biextension anchors remain FGA232 Theorem3.1 pp148--149; FGA236 Corollaries1.7/2.3 pp229/232; SGA7 I VII2.7/2.9 pp169--174 and VIII7.4(b) pp310--311. Formal groups/crystals/Cartier duality retain cote4 scan003 §§8.1.1--8.1.2, scans005--006 §§8.4.1--8.4.3 and cote54 scans005--007. These are contextual locators, not sources of this new support exhaustion. No new scheme identification, lambda-ring/K-theory identity, l-adic sheaf/weight theorem, CM motive comparison or motivic L-function result is asserted. The change to DP\_MULT Q3e and THETA3 G1/A7 is precisely J21.1. CX\_GRH CX-ALG-024 and GCT\_PLETH T-052 keep their distinct open arithmetic maps.

All new quantitative inputs have separately checked full proofs: DO085, DO086 and R063. The remaining cited inputs are the already checked K069 partition/table, DO064 ranks32/33 table, DO067 nonunipotence criterion, K081/K082 tables, DO083, R058, R059/R061 and R062. The present certificate checks every finite table index and complete historical error list, not occurrence of a rational matrix. Whole human-language proof reading is separately recorded. No Lean build or archive audit is performed.

\clearpage

\section{AM-DO-087 --- the full rank-one fixed-block sector at support twenty-one}\label{am-do-087-the-full-rank-one-fixed-block-sector-at-support-twenty-one}

\textbf{Canonical checked integration --- 2026-11-22T19:46:41.496348+00:00.} Root and independent whole-proof/script checks and separate358-assertion replays passed. The entire actual all-top H133 fixed-W0 globally rank-one unipotent sector at support21 has absolute contribution less than2\^{}-1649. The exact K080 divisor constants and relative-F norm count are composed with DO083's actual selected Hilbert trace, all intrinsic double-pole sets, all weighted incidence types and paired finite factors. Global unipotence is an explicit hypothesis. No standalone complete-floor, signed nonvanishing, source-family or DP44 conclusion is made.

\textbf{Algebraic-geometric translation.} This is an original-coordinate count of the rank-one endomorphism and its negative column, including exact content ideals, their affine cosets, the relative-F norm equation and the labelled theta incidences. No moving integral basis, arithmetic theta-image inverse or integral norm bound for a selected Hilbert vector is introduced. Contextual module/determinantal anchors remain SGA3 I4.7.3/4.7.3.1 and SGA6 VI Theorem3.3; Picard/biextension anchors remain FGA232 Theorem3.1 pp148--149, FGA236 Corollaries1.7/2.3 pp229/232, SGA7 I VII2.7/2.9 pp169--174 and VIII7.4(b) pp310--311. Formal groups, crystals and Cartier duality retain cote4 scan003 §§8.1.1--8.1.2, scans005--006 §§8.4.1--8.4.3 and cote54 scans005--007. These locators are contextual; no new scheme identification, lambda-ring/K-theory identity, l-adic weight/sheaf theorem, CM-motive comparison or motivic L-function result is asserted. DP\_MULT Q3e and THETA3 G1/A7 gain precisely this additional containing-set bound. CX\_GRH CX-ALG-024 and GCT\_PLETH T-052 retain their separate open arithmetic maps.

Author: ALP\_MACHINE \nolinkurl{/root/definite_oscillator}, 22 November 2026 UTC. Controlled pending proof. Author whole-file, exact replay and independent review records are separate.

\subsection{1. Exact scope}\label{exact-scope-2}

Retain the actual J56 all-top role assignment, K044 upper endpoint, every allowed word, the actual selected compact representation and its Hilbert metric, the original mixed finite projectors, central representatives, scalar-pencil incidences, class masses and near-point transport. The archimedean cutoff is H=133, following the already charged R048 annulus. Write a=f=168 and E\^{}169=S perpendicular W0, where W0 has signature(4,1).

This theorem concerns precisely the actual terms

\[
 h=I_S\oplus(I+N),\qquad \operatorname{rank}_E N=1,
 \qquad N^2=0,\qquad s=21.                         \tag{RD1}
\]

Here s counts active relative T places. Global unipotence is a hypothesis; no new criterion establishing it is used. The conclusion is

\[
 \boxed{\displaystyle\sum_{\text{actual RD1 terms}}|w(h)|<2^{-1649}.}
                                                               \tag{RD2}
\]

All original incidence types, ordered pairs and intrinsic double-pole subsets are included. This is an absolute bound for the stated sector. It does not independently establish a complete support floor, a signed remainder inequality, a prescribed period, Mram nonvanishing, a literal Host/FamP constructor or a DP44 comparison.

\subsection{2. Combine the exact divisor count with the actual selected trace}\label{combine-the-exact-divisor-count-with-the-actual-selected-trace}

Use the accepted K075(RN5--19), K077(AP1--18), reconstructed K080(RC2--4) and DO083(RO3--13) directly. Their algebraic content and affine-coset proofs hold at any positive cutoff; H enters those proofs only through the bound for the original negative entry. In the present cutoff,

\[
 \eta=(Ne_n)_n\ne0,\qquad \bar\eta=-\eta,
 \qquad R=\frac{133-133^{-1}}2=\frac{8844}{133},
 \qquad |\tau\eta|\le R.                           \tag{RD3}
\]

The full negative column x=Ne\_n reconstructs N uniquely by N=x\bar x\^{}\{\mathsf t\}J\_0/eta. No further source-isometry fibre is multiplied into this count.

Let D be the intrinsic subset of active places where eta has an actual double pole, and let d=\textbar D\textbar. By K077(AP5--7), D consists of type-II places; a nonintegral type-II term cannot have fixed label zero. The negative pole column and row are units at both CM partners. Consequently the actual leading residue gives the specified affine congruences for the positive coordinates. This retains both CM components and all admitted zero labels outside D. It makes no claim about a real-part projection of a general target lattice.

For a fixed paired rank-zero G pattern, let J be its paired E ideal, D\_G=N\_E(J), and M=Q\^{}(f−1). Then

\[
 D_G\le M^2,\quad
 \mathfrak l=\mathfrak J I_\Sigma I_D,\quad
 N_E(\mathfrak l)=D_GT^{-2(s+d)}=z^{2f},
 \quad y=R/z,\quad y_c=R/(zT^{2d/f}).              \tag{RD4}
\]

The coordinate-content ideal b satisfies N\_E(b) \textgreater= \textbar N\_E(eta)\textbar\^{}(1/2) z\^{}f.~For fixed eta and b, the first three positive coordinates occupy the prescribed affine cosets of b intersect J I\_Sigma, whose norm is N\_E(b) T\^{}(2d). The fourth coordinate is counted by the relative-F norm equation using b. These are the exact K077(AP9--16) conditions, not a density approximation.

The accepted actual-field divisor estimate reconstructed K080(RC2--3) is

\[
 d_E(\mathfrak a)\le C_E N_E(\mathfrak a)^{1/8},
 \quad C_E=2^{48}/29^6,
 \qquad
 d_F(\mathfrak b)\le C_F N_F(\mathfrak b)^{1/8},
 \quad C_F=2^{24}/29^3.
\]

Its full constants allow all48 and24 exceptional norm29 primes. Thus the content count is C\_E y\^{}42 and the fourth-coordinate count is 28 C\_F y\^{}21. Keeping the exact pure-imaginary eta count and affine triple count gives

\[
 \#\{\text{actual matrices with fixed data}\}
 \le 28C\,y^{63}(1+2y)^f(1+2\sqrt{y_c})^{6f},
 \qquad C=2^{72}/29^9.                             \tag{RD5}
\]

This is precisely K080(RC4) at the radius RD3. Only the divisor factors change from DO083's older y\^{}105 count. In particular, no older numerical homogeneous-energy estimate from K080 is imported.

For the actual selected trace, use DO083(RO4--7,RO12), which already applies DO080's genuine normalized-trace estimate to the multiplicatively transported rank-one negative entry. It gives

\[
 |c(h)|<3^f2^{-931f/2}|N_E\eta|^{-1/2}
         \le3^f2^{-931f/2}z^{-f}.                 \tag{RD6}
\]

The second inequality is an ideal-norm statement. It does not assert a pointwise lower bound for each conjugate of eta. Both compact factors, the actual PBW/Pieri metric and the near-point factor are retained in the cited result.

\subsection{3. The weighted cap, with the new exponent}\label{the-weighted-cap-with-the-new-exponent}

The unchanged rank-one finite estimate is

\[
                     |a(h)|\le2D_G^{a/2}M^{-a}.   \tag{RD7}
\]

Its single factor two is the complete Euler correction. Both mixed factors retain their norm bound one. There are at most2\^{}334 paired rank-zero patterns; rank one needs no additional orientation multiplicity.

Multiply RD5--7 before imposing the cap D\_G\textless=M\^{}2. Using D\_G=T\^{}(2(s+d))z\^{}(2f), the z dependence, up to positive constants independent of z, is

\[
 z^{27321}(z+2R)^f
       (\sqrt z+2\sqrt R\,T^{-d/f})^{6f},
 \qquad27321=(a-1)f-63-4f>0.                      \tag{RD8}
\]

Every displayed factor is increasing for z\textgreater0. Therefore one may replace z by

\[
 z_{s,d}=M^{1/f}T^{-(s+d)/f},\qquad
 y_{s,d}=R/z_{s,d},\qquad
 y'_{s,d}=R M^{-1/f}T^{(s-d)/f}.                  \tag{RD9}
\]

The finite factor M\^{}(-a)T\^{}(a(s+d))z\_(s,d)\^{}(af) then cancels exactly. This is a monotone bound on the complete weighted expression. It does not assert the existence of an actual matrix at the cap, even if the cap is beyond the nonempty-set restriction z\textless=R.

\subsection{4. Every marked set and actual ordered pair}\label{every-marked-set-and-actual-ordered-pair}

Each marked double-pole place is type II and has full weighted ordered-pair cost1. At each other active place, all type-I, type-II and type-III alternatives have total costs T−1,1,1, respectively. K077(AP18) consequently gives

\[
                 \binom fs\binom sd(T+1)^{s-d}.  \tag{RD10}
\]

The double set is intrinsic to eta. In particular simple-pole type-II terms remain in the unmarked alternatives; no summand is assigned two intrinsic double sets. The affine bound RD5 is uniform in the prescribed actual pair data, so it may be multiplied by this full weighted sum. No further target-volume factor or isometry fibre is needed, because the full column already determines N.

Put

\[
 C_s=2^{335}\,28C\binom fs\,3^f2^{-931f/2},
 \qquad
 B_s(d)=(T+1)^{s-d}z_{s,d}^{-f}y_{s,d}^{63}
      (1+2y_{s,d})^f(1+2\sqrt{y'_{s,d}})^{6f}.    \tag{RD11}
\]

As a function of real d in{[}0,s{]}, log B\_s(d) is affine plus positive multiples of log(1+exp(A+Bd)). It is convex. Thus the complete weighted binomial sum satisfies

\[
 \mathscr R_{1,s}
 \le C_s\sum_{d=0}^{s}\binom sd B_s(d)
 \le2^s C_s\max\{B_s(0),B_s(s)\}.                \tag{RD12}
\]

This is one maximum, not a union of two endpoint sectors. Both endpoints are retained, and y'\_(s,s)=R M\^{}(-1/f) is independent of s.

\subsection{5. Exact support21 endpoints}\label{exact-support21-endpoints}

The companion certificate reads the actual J56 integers T,g0,g1 from the original host certificate. It brackets every f-th root and square root in RD9 by rational numbers, then computes directed rational logarithms. All constants in RD11, including C and the complete2\^{}s binomial coefficient, are included. The strict intervals are

\begin{longtable}[]{@{}rrr@{}}
\toprule\noalign{}
Support s & Marked size d & Strict interval for log2(2\^{}s C\_s B\_s(d)) \\
\midrule\noalign{}
\endhead
\bottomrule\noalign{}
\endlastfoot
21 & 0 & (-1650,-1649) \\
21 & 21 & (-4472,-4471) \\
22 & 0 & (4700,4701) \\
22 & 22 & (260,261) \\
\end{longtable}

Hence RD12 at s21 is strictly less than2\^{}(-1649), proving RD2. The two support22 expressions exceed one; this records a failure of this sufficient envelope only, not a positive contribution, nonzero matrix or lower bound for the signed sum.

The certificate also checks the positive cap exponent, both divisor constants and exponent substitutions, the full three-type polynomial, the22 intrinsic double-place possibilities and exact rational convex-envelope models. These diagnostics supplement the proof and accepted input lemmas; they do not enumerate the rational automorphic support.

\subsection{6. Error accounting, dependencies and limits}\label{error-accounting-dependencies-and-limits}

Append1649 once when removing the intersection of RD1 with the current remainder. Keep every earlier error. The accompanying exact check starts with the actual canonical R063 ordered64-term all-top and31-term uniform arrays, including DO0857548, DO0868622 and R06352215 in that order. It appends1649 to obtain65 all-top entries and leaves all31 uniform entries unchanged. Both complete dyadic sums remain strictly below2\^{}(-230). This arithmetic record does not itself perform canonical integration or assert a support-floor change; any intervening complete exhaustion with no new charge leaves the same ordered arrays.

The proof is a new composition of accepted results, not a reproof of the earlier small-prime field classification, norm-equation geometry or selected Hilbert model. The directly read dependencies are reconstructed K080(RC1--7), K077(AP1--23), and the whole DO083 proof including RO3--17. The exact selected-model theorem is DO080(RT2,RT24), already used in DO083. No new primary literature theorem is invoked, and no prescribed arithmetic section, period value or sign is inferred.

The full fixed-W0 rank-one sector at support21 is now bounded in its stated global-unipotence scope. Every other structural class and the main signed inequality retain their separate requirements.

\clearpage

\section{AM-LIFT-072 --- the prescribed positive-rank split G test}\label{am-lift-072-the-prescribed-positive-rank-split-g-test}

\textbf{Canonical checked integration --- 2026-11-22T19:50:49.909848+00:00.} Root and independent whole-proof/source/script checks and separate11,773-assertion replays passed. Original pending status is historical. For the actual positive-rank split G model, every prescribed smaller inverse-generic K\^{}1-fixed vector admits a bottom vector inside the same admitted generic finite Pi\_0 with canonical Hilbert Bessel period1/d\_varsigma. This includes the actual mixed regular-series types. With literal source Haar the value is c\_H\emph{c\_N/d\_varsigma. The normalized local factor remains1/(d\_varsigma}Lscript), and the product bound is greater than Q\^{}(-f\emph{k}(k+1)/2). The genuine LIFT069 chosen-bottom and LIFT070 full-Sizes hypotheses remain. No global period, integral bound for the Hilbert vector, actual Read/032 value, source family or DP44 cost is supplied.

\textbf{Algebraic-geometric translation.} The finite extension is an intertwiner between generic representation spaces on the specified residue-group quotient of the original Hermitian module. Its exact compact character and selfdual Haar comparison concern the analytic realization of that same labelled correspondence. They give no integral inverse or saturation of the arithmetic theta image. The carried module/exterior-power context is SGA3 I4.7.3/4.7.3.1 and SGA6 VI Theorem3.3; Picard and biextension anchors remain FGA232 Theorem3.1 pp148--149, FGA236 Corollaries1.7/2.3 pp229/232, SGA7 I VII2.7/2.9 pp169--174 and VIII7.4(b) pp310--311. Formal groups/crystals/Cartier duality retain cote4 scan003 §§8.1.1--8.1.2, scans005--006 §§8.4.1--8.4.3 and cote54 scans005--007. These contextual locators imply no new scheme identification, lambda-ring/K-theory identity, l-adic sheaf or weight theorem, CM-motive comparison or motivic L-function statement. DP\_MULT Q3e and THETA3 G1/A7 lose the stated local G test obstruction under the retained source premises; all global arithmetic maps and costs remain open. CX\_GRH CX-ALG-024 and GCT\_PLETH T-052 gain no unconditional global conclusion.

Author: \nolinkurl{/root/lift_accounting}, 22 November 2026 UTC.

\textbf{Author whole-proof, source-interface and complete-script review PASS; exact diagnostic PASS (11,773 assertions).} The final replay, timestamps and frozen fingerprints are recorded in the adjacent author-review and manifest files. Independent reviews are separate. This is a local test theorem on the admitted finite types. It supplies no global period or arithmetic Read instance.

\subsection{1. Representations, actual vectors, and statement}\label{representations-actual-vectors-and-statement}

Let F be a nonarchimedean field with integers O, maximal ideal p, and residue field l of cardinality q. Fix a character chi of conductor p: it is trivial on p and nontrivial on O. Give additive O volume one, every matrix-root O volume one, and K\_d=GL\_d(O) volume one. Quotients use the corresponding group measure divided by the indicated upper-unipotent product measure.

Put k\textgreater=2, rho\textgreater=0, and n=k+2rho+1. Let Pi and sigma be irreducible tempered generic representations of GL\_n(F) and GL\_k(F). Suppose Pi\^{}\{K\_n\^{}1\} contains the admitted generic irreducible finite representation Pi\_0 of GL\_n(l). It may be cuspidal or one of the actual mixed regular-series representations. Fix a nonzero vector y in sigma\^{}\{K\_k\^{}1\} with the inverse generic character on U\_k(O). This is the given smaller vector in the bilinear convention; it is not replaced by a convenient vector. In a Hermitian period, sigma and y denote the contragredient representation and conjugate vector of its fixed smaller slot.

Unit root coefficients can be retained throughout. For the displayed formulas normalize them by an integral diagonal, transported on both factors. In central block order (H,E,z,F), of sizes (k,rho,1,rho), use the standard upper Whittaker character

\[
 \psi_d(u)=\chi\left(\sum_{i=1}^{d-1}u_{i,i+1}\right).
 \tag{GB1}
\]

The larger character is psi\_n and the smaller is psi\_k\^{}\{-1\}. Let N\_B have identity H and z blocks, upper unitriangular E and F blocks, and arbitrary E-to-H, E-to-z, E-to-F, H-to-F, and z-to-F blocks. Its generic character psi\_B uses the E-chain, E-to-z, z-to-F, and F-chain entries, and is trivial on its other roots. Set S=N\_B H. Conjugating by the fixed permutation to (E,H,z,F) gives the actual Bessel flag subgroup of LIFT060.

\textbf{Theorem.} There is a Hilbert-unit vector x in the same finite Pi\_0-subspace, with the prescribed character under the entire inverse-image compact of U\_\{rho,k\} in LIFT060, such that the canonical tempered Bessel matrix-coefficient period on Hilbert-unit vectors is

\[
 \boxed{\alpha_G(x,y/\|y\|)=\frac1{d_\varsigma}>0.}
 \tag{GB2}
\]

Here varsigma is the unique generic irreducible finite GL\_k(l) representation generated by y, and d\_varsigma is its dimension. Its uniqueness is proved below from the actual Whittaker restriction, rather than assumed from a generic finite type. The regularized N\_B integration uses the literal psi\_B\^{}\{-1\}; H and every root O have volume one. If the specified source measures are c\_H and c\_N times these measures, the value is c\_H c\_N/d\_varsigma.

The chosen bottom vector stays inside the existing local representation and its admitted finite Pi\_0. This assertion is valid for either simultaneous choice of character orientation, with the relevant conjugate slots explicitly transported. It is not a genericity assertion for an arbitrary representation with finite generic invariants.

\subsection{2. The exact finite subgroup and its intersection with U\_n}\label{the-exact-finite-subgroup-and-its-intersection-with-u_n}

Temporarily work over l. Write P\_n for the last-row mirabolic subgroup and U\_n for the upper unitriangular subgroup in the central order. Then S is a subgroup of P\_n.~In block form it is

\[
 S=\left\{
 \begin{pmatrix}
 h&0&0&B\\ X&u_E&b&C\\0&0&1&d\\0&0&0&u_F
 \end{pmatrix}\right\}.
 \tag{GB3}
\]

Let N\_B\^{}+=N\_B intersect U\_n.~It is obtained by setting the rho-by-k E-to-H block to zero. Thus

\[
 U_n\cap S=N_B^+ U_k,\qquad
 [N_B:N_B^+]=q^{\rho k},\qquad
 \dim N_B=\rho(2k+2\rho+1).
 \tag{GB4}
\]

These are exact matrix identities: the only entries of N\_B below the standard diagonal are its E-to-H entries. Its product with H is semidirect. The H action has determinant modulus one on N\_B, because the rho incoming and rho outgoing H blocks cancel.

On U\_n intersect S, psi\_n is exactly the product of psi\_B and psi\_k. In particular there is no extra H-to-E root contribution: that block is zero in S. The H-to-z block is also zero. The nonadjacent shortcut in the order (E,H,z,F) becomes the adjacent E-to-z root in the central order.

The subgroup N\_B U\_k, after the same permutation, is precisely

\[
 U_{\rho,k}(l)=\{u\in U_n(l):u_{i,z}=0\text{ for every H row }i\},
 \tag{GB5}
\]

and its character is the full two-chain character of HF12, after the exact unit diagonal HF13--14. Its order is q\^{}\{n(n-1)/2-k\}, agreeing with HF18. No Bessel root is omitted from this compact character.

\subsection{3. A finite extension inside the admitted mirabolic summand}\label{a-finite-extension-inside-the-admitted-mirabolic-summand}

Let

\[
 I_n=\operatorname{Ind}_{U_n(l)}^{P_n(l)}\psi_n,
 \qquad \Gamma_k=\operatorname{Ind}_{U_k(l)}^{GL_k(l)}\psi_k,
 \qquad D_d(q)=\prod_{i=1}^d(q^i-1).
 \tag{GB6}
\]

Both function models have left character equivariance and right translation. The actual finite Pi\_0 contains the unique I\_n summand, as in LIFT032 FB2--5; the little-group proof works for every generic Pi\_0. Uniqueness follows also by Frobenius reciprocity from its one-dimensional generic U\_n line. No cuspidality of Pi\_0 is required.

For F in Gamma\_k define E\_rho(F) in I\_n by support in U\_n S and

\[
 E_\rho(F)(vuh)=\psi_n(v)\psi_B(u)F(h),
 \quad v\in U_n,\ u\in N_B,\ h\in H,
 \qquad E_\rho(F)=0\text{ outside }U_n S.
 \tag{GB7}
\]

To check well-definedness, two presentations differ by U\_n intersect S. Its H component is in U\_k, and its N\_B component is in N\_B\^{}+. GB4 and the character identity following it give exactly the change in F under its left U\_k character. Thus every phase agrees. The support is stable under right S, and

\[
 R(u_0)E_\rho(F)=\psi_B(u_0)E_\rho(F),\qquad
 R(h_0)E_\rho(F)=E_\rho(R(h_0)F).
 \tag{GB8}
\]

Finite induction and the normality of N\_B show the norm exactly. Each class in U\_k\H has {[}N\_B:N\_B\^{}+{]} preimages in (U\_n intersect S)\S, and all values have the same modulus. Probability induced norms therefore give

\[
 \|E_\rho(F)\|_{I_n}^2
   =\frac{q^{\rho k}D_k(q)}{D_{n-1}(q)}
                           \|F\|_{\Gamma_k}^2.
 \tag{GB9}
\]

This is a concrete extension, not a multiplicity-one choice of an unspecified Bessel vector. It also explicitly exhibits every generic H type in the required finite Bessel quotient. If F has the right psi\_k character, GB8 gives the entire right psi\_B psi\_k character on GB5.

\subsection{4. Actual p-adic support, including every X variable}\label{actual-p-adic-support-including-every-x-variable}

Use the positive invariant Whittaker norms from Jacquet--Shalika, with the smaller character inverse. The support lemma K055 MK8--14 says that restriction of the admitted I\_n to P\_n(F) is exactly

\[
 C_n=\{f:\ f(up)=\psi_n(u)f(p),\quad
      \operatorname{supp}f\subset U_n(F)P_n(O),\quad
      f\text{ right }P_n^1\text{-invariant}\}.
 \tag{GB10}
\]

In particular every finite function GB7 is the restriction of a unique vector in that actual I\_n subspace of Pi. This uses the injective Whittaker restriction and its positive invariant norm. Finite genericity alone is not used to assert p-adic genericity.

For the central JPSS matrix

\[
 A_\rho(h,X)=
 \begin{pmatrix}h&0&0\\X&I_\rho&0\\0&0&I_{\rho+1}\end{pmatrix},
 \tag{GB11}
\]

the exact support condition is

\[
 A_\rho(h,X)\in U_n(F)P_n(O)
 \quad\Longleftrightarrow\quad
 h\in U_k(F)K_k,\quad X\in\operatorname{Mat}_{\rho,k}(O).
 \tag{GB12}
\]

Here is a block proof which excludes hidden nonintegral X tails. If A=u p with u upper unipotent and p integral invertible, the last rho+1 columns of A are standard columns. The corresponding entries of p show that the tail part of u\^{}\{-1\} and its upper rectangular tail block are integral. Integral upper row operations remove that tail. On the remaining k+rho block write u={[}{[}u\_H,V{]},{[}0,u\_E{]}{]}. Then

\[
 u^{-1}\begin{pmatrix}h&0\\X&I\end{pmatrix}
 =\begin{pmatrix}
 u_H^{-1}h-u_H^{-1}Vu_E^{-1}X&-u_H^{-1}Vu_E^{-1}\\
 u_E^{-1}X&u_E^{-1}
 \end{pmatrix}\in GL_{k+\rho}(O).
 \tag{GB13}
\]

Its lower-right block implies u\_E and u\_E\^{}\{-1\} are integral. Its lower-left block implies X integral. Adding the integral upper-right multiple of the lower rows then shows u\_H\^{}\{-1\}h integral. Determinant invertibility makes it an element of K\_k. This proves the forward implication. Conversely, write h=u\_H h\_0 with h\_0 in K\_k. Left multiplication by diag(u\_H\^{}\{-1\},I) leaves an integral invertible block matrix when X is integral, proving the reverse implication.

On this support, the vector with restriction GB7 satisfies

\[
                W_x(A_\rho(h,X))=F(h).
 \tag{GB14}
\]

Indeed for h in K\_k, the lower block Xh\^{}\{-1\} is in N\_B(O), where its Bessel character is one; GB7 gives F(h). General h in U\_k K\_k follows by left Whittaker equivariance. Thus the central integral has compact support modulo U\_k, and the entire X integration has volume one. There is no continuation argument needed to evaluate this particular JPSS integral.

\subsection{5. The unchanged smaller vector and both exact norms}\label{the-unchanged-smaller-vector-and-both-exact-norms}

For the given y, the K055 support and evaluation argument gives a nonzero c=W\_y(1) and a unique generic irreducible finite representation varsigma generated by y. Briefly, its P\_k(O) span is I\_k, its Whittaker restriction lies in the unique C\_k, and the right inverse-generic line in C\_k is one-dimensional. It is the delta function at the identity upper-unipotent coset. Hence all possible generic K\_k\^{}1-fixed copies inject into the same one-dimensional line. Finite semisimplicity gives the asserted unique varsigma and multiplicity one.

Set F(h)=conjugate(W\_y(h)) on K\_k, extended to U\_k(F)K\_k with left psi\_k character. It has the right psi\_k character. Let x be its GB7 lift in the admitted Pi\_0. Finite Schur orthogonality, with K\_k probability measure, gives

\[
 B:=\int_{K_k}|W_y(h)|^2dh=\frac{|c|^2}{d_\varsigma}>0.
 \tag{GB15}
\]

The evaluation functional on varsigma has Riesz vector proportional to y; the proportionality cancels in this formula regardless of the scaling of the initial Whittaker embedding.

Let theta\_\{d,prob\} denote the mirabolic Whittaker norm with probability GL\_\{d-1\} quotient Haar and unipotent O volume one. Equations GB9--14 and the K055 initial delta norm give

\[
 \begin{split}
 Z_{\rho,prob}(x,y)&=B,\\
 \theta_{n,prob}(x,x)
    &=\frac{q^{\rho k}D_k}{D_{n-1}}B,\\
 \theta_{k,prob}(y,y)&=\frac{|c|^2}{D_{k-1}},\\
 \frac{|Z_{\rho,prob}|^2}{\theta_{n,prob}\theta_{k,prob}}
    &=\frac{D_{n-1}D_{k-1}}{q^{\rho k}D_kd_\varsigma}.
 \end{split}
 \tag{GB16}
\]

The finite quotient normalization in the second line is 1/D\_\{n-1\} per U\_n\P\_n coset. The first line includes the factor \textbar det h\textbar\^{}\{-rho\} in the central JPSS definition; it equals one on this support. For rho=0 these are precisely K055 MK18--20.

\subsection{6. Conductor-p selfdual Haar and the full matrix period}\label{conductor-p-selfdual-haar-and-the-full-matrix-period}

The character chi has conductor p, not O. If dx\_0 gives O volume one, its selfdual additive Haar is dx\_sd=q\^{}\{1/2\}dx\_0. Define

\[
 \kappa_d=\prod_{i=1}^d(1-q^{-i}),\quad
 g_d=q^{d^2/2}\kappa_d,\quad
 r_d=q^{d(d+1)/4}\kappa_d,
 \quad d_B=\rho(2k+2\rho+1).
 \tag{GB17}
\]

The raw GL\_d differential measure is g\_d times maximal-compact probability Haar. Dividing by the selfdual U\_d root measure gives r\_d times the probability quotient. The N\_B product measure contributes q\^{}\{d\_B/2\}; the rho-by-k X integral contributes q\^{}\{rho k/2\}. Consequently

\[
 Z_{\rho,sd}=q^{\rho k/2}r_kZ_{\rho,prob},\quad
 \theta_{n,sd}=r_{n-1}\theta_{n,prob},\quad
 \theta_{k,sd}=r_{k-1}\theta_{k,prob}.
 \tag{GB18}
\]

Apply LIFT071 HN8, whose normalization lemma explicitly treats arbitrary tempered generic Pi and sigma. Its proof works with these selfdual measures: the quotient-row coarea dimensions add exactly, every Fourier inversion is with its selfdual Haar, and Zhang's adjacent-rank calibration is in the same raw measures. Neither sphericity nor a conductor-one smaller representation is used in HN8. Those assumptions occur only in LIFT071's later evaluation of its particular T vectors.

Thus the full canonical regularized matrix period on Hilbert-unit vectors, with H and all roots normalized as in section1, is

\[
 \alpha_G
 =\frac{q^{\rho k}r_k^2}
 {q^{d_B/2}g_kr_{n-1}r_{k-1}}
 \frac{|Z_{\rho,prob}|^2}
      {\theta_{n,prob}\theta_{k,prob}}
 =\frac1{d_\varsigma}.
 \tag{GB19}
\]

For the exact cancellation, use D\_d=q\^{}\{d(d+1)/2\}kappa\_d.~All kappa factors cancel, and the remaining q exponent is

\[
 \frac{n(n-1)-k(k+1)}4-\frac{d_B}2=0.
 \tag{GB20}
\]

The coefficient of the outer H measure and of the N\_B measure appears once, not twice. This proves GB2 with its full Haar normalization. It is a value of the canonical matrix-coefficient period, not merely a nonzero compact-period test which could be canceled by a noncompact tail.

\subsection{7. Physical G frames and the two conjugate source readings}\label{physical-g-frames-and-the-two-conjugate-source-readings}

At an actual G place the rational prime is completely split and has degree one. The literal character is the trace character divided by G; it has conductor p.~Its units, including the other G prime and the different-scaled flags, are retained. In the split plus component the opposite component is the inverse adjoint. The actual N, H and root characters in LIFT060 HF9--14 and HF21--22 therefore have exactly the shapes in GB3--5, after the same prescribed permutation and integral unit diagonal.

The norm-minus-T vector and its positive complement also have an integral local coordinate transport at G: T is a G-unit, and the LIFT030 matrix and its inverse are integral there. For the retained norm generators pi\_T congruent to1 modulo G, this move is principal at G; its extra diagonal character correction is principal too. More generally the LIFT060 construction needs only the stated T congruent to1 modulo G and its actual integral transport. All such unit conjugations have absolute root modulus one, preserve K\^{}1, and carry Pi\_0 to its same K-module. No finite representation or principal-G depth is replaced.

The group/character matching can also be checked without calling an integral transport trivial. The physical z\_T equals z0 modulo G, the e and f data agree at G, and HF12 computes both CM components. Thus the reduction of N\_T(O)U\_H(O) is the exact U\_\{rho,k\} in the physical frame, and the product of its Bessel character with its H-generic character is the literal chi2. Because x is K\_n\^{}1-fixed, GB8 gives the required character under its full inverse-image compact K2, not just under its rational S subgroup. The diagonal correction HF13--14 transports the fixed initial character at the same time; it is not a freely selected new additive character.

There is a source conjugation which must remain visible. In the Hida Hilbert row space HHX, the right compact character is chi2, and the initial column in \nolinkurl{PoincareInHH.psiK} has its restriction to H. The Hermitian period pairs that column with the bottom restriction. In the bilinear formulas above, y is its conjugate vector, so its H-generic character is the inverse. By contrast, \nolinkurl{RowDict.EEC} is the conjugate of the good HHX subspace, and \nolinkurl{conjHH} applies complex conjugation isometrically. \nolinkurl{RowLevelMatch.LevelMatchA} explicitly records levChi=star(chi2). These facts are read directly from the source; neither occurrence of conjugation is dropped.

Equivalently, invert both characters and conjugate the two representation slots. The squared period and GB19 are unchanged. The full genuine idempotent is coefficientwise real, as in the checked RowLevelMatchC/real-idempotent dictionary, so this conjugation preserves the prescribed complete good subspace. It does not assert that an individual goodTuple label is self-dual, or replace it silently by its reciprocal. In either reading Pi\_0 is the generic finite constituent actually present in that slot, and the construction stays inside it.

Actual source Haar constants remain c\_H,c\_N. Integral changes of coordinates have modulus one relative to the product conventions, but that statement does not set source-chosen independent Haar scales equal to one. The exact final factor is c\_Hc\_N/d\_varsigma.

\subsection{8. Actual hypotheses, products, and retained global boundary}\label{actual-hypotheses-products-and-retained-global-boundary}

LIFT069 supplies the tempered generic bottom representations on the chosen LIFT059 subspace and its conjugate reading. Its finite Pi\_0 at every G place is a genuine generic regular-series constituent; the mixed factors are admitted, since the proof uses I\_n rather than finite cuspidality. LIFT070 supplies temperedness and genericity of the actual initial representation, or its specified dual, under the full genuine SizesAt/SizeData\_2 assumptions. Its initial G vectors have the actual inverse-generic finite type. These are the hypotheses used here. No implication from finite generic invariants to p-adic genericity is substituted for either result.

For independent prescribed initial vectors at all G places, choose the local bottom vectors by GB7. The restricted tensor product permits these finite choices inside the same bottom automorphic representation, with its same real type and every other local factor. With the probability/product-root convention,

\[
 \prod_{v\mid G}\alpha_v
     =\prod_{v\mid G}\frac1{d_{\varsigma_v}}
     \ge\prod_{v\mid G}D_k(q_v)^{-1}
     >\prod_{v\mid G}q_v^{-k(k+1)/2}.
 \tag{GB21}
\]

The finite GL Gelfand--Graev representation is multiplicity free and has dimension D\_k; hence d\_varsigma\textless=D\_k. Exactly one relative place per CM pair is counted. For two rational G primes completely split over F, the last bound is (g0 g1)\^{}\{-fk(k+1)/2\}. The dependence on rho cancels from the normalized local period, not from the dimension or construction of the finite vector.

The canonical local GGP factor is still

\[
 \alpha_v^\natural=\frac1{d_{\varsigma_v}\mathscr L_v},\qquad
 \mathscr L_v=\kappa_n^{-1}
 \frac{|L(1/2,\Pi_v\times\sigma_v)|^2}
 {L(1,\Pi_v\times\Pi_v^\vee)L(1,\sigma_v\times\sigma_v^\vee)}.
 \tag{GB22}
\]

Those positive, finite local factors are not assigned value one. Source measures multiply GB22 by c\_Hc\_N. No coefficient-field or integral-unit estimate for the Hilbert-normalized x is asserted.

The theorem removes a local G test-vector obstruction in the positive-rank prescribed model. It does not prove the corresponding global central value or Bessel/Mram period nonzero, nor identify the desired arithmetic Read/STiesR/032 row. It constructs no complete Host/SizesAt family. The actual selected real vector, remaining finite places, scalar intertwiners, and arithmetic content stay separate. A global pair can still have zero period even though every displayed local test is nonzero.

\subsection{9. Sources and verification scope}\label{sources-and-verification-scope-10}

The direct source reads are HidaSetting \nolinkurl{Nfin}, \nolinkurl{psiFin}, \nolinkurl{chi2f} and its finite projector; RowLevelMatch18--40; PoincareInHH126--145 and156--177; RowDict98--120; LinksChar106--119; and RecLevel49--52. The exact group/character and source translation proofs are the accepted LIFT060 HF9--22 and LIFT030 NL4--9. The explicit \nolinkurl{conjHH} reading is also carried by LIFT065--066.

The finite mirabolic embedding and generic finite series are LIFT032 FB2--5. K055 MK8--20 supplies the actual p-adic support/evaluation lemma and the unchanged smaller-vector norm; its proof was reread. LIFT071 HN8--15 supplies the calibrated arbitrary-tempered-pair matrix-period identity, with conductor-p selfdual measures tracked in GB17--20. LIFT069--070 supply only the stated actual representation hypotheses.

Primary anchors are Jacquet and Shalika (1983), \emph{The Whittaker model of induced representations}, sections2.3--2.4, printed110--111, \url{https://www.math.columbia.edu/~hj/The\%20Whittaker\%20model\%20of\%20induced\%20representations.pdf}; Gan, Gross, and Prasad, \emph{Restrictions of representations of classical groups: examples}, Proposition4.2 and Corollary4.3, \url{https://abel.math.harvard.edu/~gross/preprints/ggp-evidence4.pdf}; Zhang (2014), \emph{Automorphic period and the central value of Rankin--Selberg L-function}, JAMS27,541--612, invariant norm(3.2) and full Proposition4.10, \url{https://doi.org/10.1090/S0894-0347-2014-00784-0}. The positive-rank regularized comparison retains LIFT071's primary Beuzart-Plessis and Beuzart-Plessis--Chaudouard inputs and its explicit scalar calibration; no scalar is deduced from uniqueness alone.

The companion certificate tests exact finite subgroup intersections, character phases and supported extension values, norm cardinalities, support block identities, and Haar cancellations. It passed 11,773 exact assertions, including seven complete support-coset models and all presentations in the specified small S models. These finite diagnostics cannot replace the general support or regularization arguments, compute a global period, or certify an arithmetic image.

\subsection{Author verification}\label{author-verification-2}

I read the complete saved GB1--22 proof and all nine substantive sections, then the complete diagnostic script. I reconstructed the intersection U\_n intersect S, the well-defined supported extension and its right character, the exact coset norm, the absence of every nonintegral X tail, the unchanged smaller generic line, and the conductor-p selfdual Haar cancellation. I directly checked the source compact character and HHX/EEC conjugation bodies and reread the relevant K055, LIFT060 and LIFT071 proofs. Irreducibility and the full inverse-image compact are explicit in the final statement. The final status and diagnostic paragraph were reread. The separate author-review record pins the final proof/script and final replay. No global period or arithmetic matrix was evaluated, and no Lean build was performed.

\clearpage

\section{AM-R-065 --- complete all-top support twenty-one exhaustion}\label{am-r-065-complete-all-top-support-twenty-one-exhaustion}

Author: \nolinkurl{/root}, GPT-6 Astra Pro, inherited effort not externally exposed, ALP\_MACHINE.22 November2026 UTC. New complete assembly; whole checks and exact replay are separate records.

\subsection{1. Statement and unchanged identity}\label{statement-and-unchanged-identity}

Retain the original J56 all-top K044 endpoint, its permitted packet, selected compact representation and Hilbert metric, original integral Hermitian modules, every paired CM incidence, finite projector, class mass and phase. Use only the already charged H133 cutoff. The dimensions remain f=a=168,n=169,m=164. R064 gives the exact energy with every remaining support at least21. DO087 is now an additional checked containing-set removal with its one1649 charge retained.

There is an actual remaining subset R\_top\^{}(22) for which

\[
 \mathcal E_{\rm top}=1+\operatorname{Re}
       \sum_{h\in R_{\rm top}^{(22)}}w(h)+\epsilon,
 \qquad |\epsilon|<2^{-230},\qquad s(h)\ge22.
 \tag{J22.1}
\]

This assembly adds no charge. The sufficient inequality

\[
 \operatorname{Re}\sum_{h\in R_{\rm top}^{(22)}}w(h)>-1+2^{-230}
 \tag{J22.2}
\]

remains open. The independent uniform support floor remains11. We prove J22.1 by exhausting every structural class at support21, not by assigning a sign to a summand.

\subsection{2. Every nonidentity positive compression at support21}\label{every-nonidentity-positive-compression-at-support21}

Write A=h\_SS in the unchanged original decomposition S perpendicular W0. The residual satisfies r=rank\_E(A−I)\textless=33. For A!=I use the exact K069(IC1--8) partition into F, where h fixes K\_A=ker(A−I) intersect ker(A*−I) pointwise, and its exceptional isotropic-cross complement IC8. Neither class is restricted to globally unipotent matrices.

\begin{longtable}[]{@{}
  >{\raggedright\arraybackslash}p{(\columnwidth - 2\tabcolsep) * \real{0.5000}}
  >{\raggedright\arraybackslash}p{(\columnwidth - 2\tabcolsep) * \real{0.5000}}@{}}
\toprule\noalign{}
\begin{minipage}[b]{\linewidth}\raggedright
Actual class
\end{minipage} & \begin{minipage}[b]{\linewidth}\raggedright
Checked support21 coverage
\end{minipage} \\
\midrule\noalign{}
\endhead
\bottomrule\noalign{}
\endlastfoot
F,r=1 & DO086(FR1--10) removes the entire class at every support18..94, with its retained8622 charge. \\
F,2\textless=r\textless=33 & K069(IC20--21) has minimum complete ceiling22 over these ranks; hence21 is covered. \\
IC8,r=1 or2 & R063(UC1--15) covers the entire class at every support18..21, including the zero, transverse and plane alternatives for the untouched cross map, with retained52215 charge. \\
IC8,3\textless=r\textless=31 & DO064's old ranges and DO085(MR1--24)'s contiguous extensions cover through at least23 at each rank, with DO085's retained7548 charge. \\
IC8,r=32 or33 & DO064(PM22--23) covers every support through168. \\
\end{longtable}

These are all ranks and both partition classes. Thus no A!=I term at support21 remains.

\subsection{3. Every identity compression at support21}\label{every-identity-compression-at-support21}

For A=I the exact original unitary identities of R059(GB1--3) give either h=I\_S direct-sum h0 on W0, or a nonzero scalar-cross block of rank one. They do not force global unipotence. Treat every alternative:

\begin{longtable}[]{@{}
  >{\raggedright\arraybackslash}p{(\columnwidth - 2\tabcolsep) * \real{0.5000}}
  >{\raggedright\arraybackslash}p{(\columnwidth - 2\tabcolsep) * \real{0.5000}}@{}}
\toprule\noalign{}
\begin{minipage}[b]{\linewidth}\raggedright
Actual A=I class
\end{minipage} & \begin{minipage}[b]{\linewidth}\raggedright
Checked support21 coverage
\end{minipage} \\
\midrule\noalign{}
\endhead
\bottomrule\noalign{}
\endlastfoot
Fixed S,nonunipotent,global rank(h−I)\textless=3 & DO067 excludes t\textless=9. For t10..20, K081/K082's complete ceilings168,164,149,134,119,104,89,73,58,43,28 are all at least21. At t=21, the explicit all-top(t,s)=(21,21) cell of K083 §4 covers ranks\textless=3. \\
Fixed S,nonunipotent,global rank(h−I)\textgreater=4 & K079(ER14--15) covers the entire rank\textgreater=4 sector at supports16..21, with no global-unipotence assumption. \\
Fixed S,globally unipotent,full rank1 & DO087(RD1--12) covers the entire support21 sector, including all incidence types and intrinsic double sets, with one1649 charge. \\
Fixed S,globally unipotent,full rank2 & R058(HB1--25) covers every profile at supports13..23, with retained39409 charge. \\
Scalar cross,globally unipotent or not & R059(GB1--22) covers the whole class at supports13..21, with retained23966 charge; R061's additional support22 charge is also retained. \\
Full rank0 & h=I has support0 and is absent at21. \\
\end{longtable}

For the first row, t is the actual type-III count and 0\textless=t\textless=s=21. The explicit K083 endpoint uses C=−14858007/4 in C+10958s+165833t. At(t,s)=(21,21) it is−7563/4, strictly negative and inside the accepted75-cell union; we use that union's existing1166 charge, not a new endpoint charge. The second row is the accepted rank\textgreater=4 full-support bound, and covers any nonunipotent rank4 or5 fixed block. The accepted signature-one classification permits no further nonzero globally unipotent ranks. Thus the six rows exhaust A=I at support21.

Together with section2, this exhausts every support21 term. R064 already exhausts all smaller supports, proving the support assertion in J22.1. No full support22 exhaustion is asserted.

\subsection{4. Exact errors and global boundary}\label{exact-errors-and-global-boundary}

The actual chronological predecessor is \nolinkurl{checks/do087_joint_remainder_error.json}. It contains65 all-top and31 uniform charges. The last six all-top exponents are4255,177258,7548,8622,52215,1649, in that exact order. R064 introduced no charge; neither does this assembly. Every earlier charge remains exactly once.

There is exactly one all-top exponent231 and every other exponent is at least685, so

\[
 \sum_{e\in\mathcal L_{\rm top}}2^{-e}
   \le2^{-231}+64\,2^{-685}<2^{-230}.
 \tag{J22.3}
\]

The unchanged31-term uniform sum is below2\^{}−230 as well. The companion checks the complete integer sums and every rank/type index used here. At each removal only its intersection with the actual current remainder is removed; those removed sets are disjoint, while their conservative containing-set charges are all retained. The original energy identity and triangle inequality now prove J22.1.

This is a strictly smaller sufficient Q1/Q2 finite/real gate. It is not a proof of J22.2, an actual admissible mixed-CM family, prescribed global Bessel/Mram nonvanishing, selected arithmetic-image content, a complete Read/comparison tie, the specific032 pairing/value, or any changed-family DP44 term. Neither a shorter valid chain nor a universal strict negative DP44 difference is proved. Q3 still has no Kummer-Hecke RH conclusion. Q4 retains the missing global arithmetic charge comparison.

\subsection{5. Algebraic geometry and check scope}\label{algebraic-geometry-and-check-scope}

The argument partitions rank conditions and fixed/isotropic submodules of the original integral Hermitian modules and their labelled correspondences. It changes no lattice, polarization, coefficient module, finite level, Picard scheme, biextension or CM realization. No rational-coordinate reconstruction is promoted to an integral arithmetic-image inverse.

The carried module/stabilizer context is SGA3 I4.7.3/4.7.3.1, and the exterior/determinantal context is SGA6 VI Theorem3.3. The Picard/biextension anchors remain FGA232 Theorem3.1 pp148--149; FGA236 Corollaries1.7/2.3 pp229/232; SGA7 I VII2.7/2.9 pp169--174 and VIII7.4(b) pp310--311. Formal groups/crystals/Cartier duality retain cote4 scan003 §§8.1.1--8.1.2, scans005--006 §§8.4.1--8.4.3 and cote54 scans005--007. These are contextual locators, not proofs of this assembly. No new scheme identification, lambda-ring/K-theory identity, l-adic sheaf/weight result, CM motive comparison or motivic L-function theorem is asserted. DP\_MULT Q3e and THETA3 G1/A7 gain precisely J22.1; CX\_GRH CX-ALG-024 and GCT\_PLETH T-052 retain their distinct arithmetic gates.

Every quantitative containing-set theorem used here has its own accepted whole proof. The new input DO087 has root and independent whole-proof/script checks and separate358-assertion replays. The present whole proof checks exhaustiveness and exact scope; its certificate is a finite table/error diagnostic and does not enumerate rational support or compute a period. No Lean build or archive audit is performed.

\clearpage

\section{AM-DO-088 --- weighted finite patterns and the exact affine metric remove rank-one support twenty-two}\label{am-do-088-weighted-finite-patterns-and-the-exact-affine-metric-remove-rank-one-support-twenty-two}

\textbf{Canonical checked integration --- 2026-11-22T19:56:41.262849+00:00.} Root and independent whole checks and separate986-assertion replays passed. The entire actual all-top H133 globally rank-one unipotent fixed-W0 sector at support22 has absolute contribution\textless2\^{}-358. The exact affine difference-ideal metric and its half-trace dual give a new ball count; both weighted envelopes retain a common pre-cap subset ratio, summing every paired G rank-zero pattern with a product\textless2. The actual same packet distance makes the total multiplicative near-point loss\textless2. All23 integer marked sizes and original type/pair costs are included. Global unipotence is a hypothesis; other structural sectors and the signed/global/arithmetic/DP44 gates remain open.

\textbf{Algebraic-geometric translation.} The new metric is the explicit positive realization of an original fractional-ideal module; its dual uses the absolute different and half-trace factor2. The affine cosets, integral endomorphism column and labelled incidences are unchanged. No moving integral basis, arithmetic theta-image inverse or selected integral-unit vector is introduced. The carried module/exterior context is SGA3 I4.7.3/4.7.3.1 and SGA6 VI Theorem3.3. Picard/biextension anchors remain FGA232 Theorem3.1 pp148--149; FGA236 Corollaries1.7/2.3 pp229/232; SGA7 I VII2.7/2.9 pp169--174 and VIII7.4(b) pp310--311. Crystals and Cartier duality retain cote4 scan003 §§8.1.1--8.1.2, scans005--006 §§8.4.1--8.4.3 and cote54 scans005--007. These are contextual locators; no new scheme, lambda-ring/K-theory, l-adic weight/sheaf or CM-motive/L-function assertion follows. DP\_MULT Q3e and THETA3 G1/A7 gain this precise containing-set removal; CX\_GRH CX-ALG-024 and GCT\_PLETH T-052 keep their separate arithmetic gates.

Author: ALP\_MACHINE \nolinkurl{/root/definite_oscillator}, 22 November 2026 UTC. New controlled pending proof. Author whole-file verification, exact replay and independent whole review are recorded separately.

\subsection{1. The original sector and conclusion}\label{the-original-sector-and-conclusion}

Retain the unchanged actual J56 all-top K044 upper endpoint, all permitted words and canonical twists, original selected compact modules and Hilbert metrics, favorable packet point, finite projectors, actual paired pencil incidences, class masses and phases. The cutoff is H=133 after the previously charged R048 annulus. Put a=f=168 and E\^{}169=S perpendicular W0, where W0 has signature(4,1).

The theorem concerns

\[
 h=I_S\oplus(I+N),\qquad \operatorname{rank}_E N=1,
 \qquad N^2=0,\qquad s=22.                         \tag{WM1}
\]

All original incidence types and ordered labels are included. The assertion is

\[
 \boxed{\displaystyle\sum_{\text{actual WM1 terms}}|w(h)|<2^{-358}.}
                                                               \tag{WM2}
\]

Global unipotence is an explicit hypothesis. No new criterion for it, complete support-floor assertion or signed positivity theorem is included. The uniform H8192 result is unchanged.

Use DO087(RD3--7) for the original negative column x=Ne\_n and its nonzero pure-imaginary last entry eta. Thus

\[
 R=\frac{8844}{133},\quad |\tau\eta|\le R,
 \quad N=x\bar x^{\mathsf t}J_0/\eta.
                                                               \tag{WM3}
\]

A full column determines N uniquely. No isometry-fibre multiplier is present below.

\subsection{2. The same packet point gives a multiplicative improvement}\label{the-same-packet-point-gives-a-multiplicative-improvement}

The rank-one transport identity DO083(RO4--6) applies to the actual argument h\_p=g\_p\^{}(-1)hg\_p of the normalized selected trace. At each real label, normalize v=x/eta analytically, so v\_n=1 and its positive part has norm one. This is only a complex normalization and asserts no finite integrality. Then

\[
 N=-\eta vv^*J,\qquad
 \eta_p=\eta|v^*Jg_pe_n|^2,
 \quad|v^*Jg_pe_n|\ge e^{-r_p},
 \quad\sinh r(h_p)=|\eta_p|.
                                                               \tag{WM4}
\]

In particular the bound is multiplicative even when an individual conjugate of eta is extremely small. An additive distance estimate is not substituted for WM4.

For this same point, LIFT41(NP5--12) gives, at every label,

\[
 \sinh^2r_p=\frac{77Q^4-14Q^2}{T}<2^{-704}.
                                                               \tag{WM5}
\]

The source vector is the actual negative-line vector (w,pi\_k)/bar(pi); its negative-coordinate modulus supplies precisely this Cartan distance. The compact factors and any norm-one central phase do not alter it. The inequality is rechecked against the original exact integers in the companion script.

Since r\_p\textless=sinh r\_p\textless2\^{}(-352), 2 sum\_tau r\_(p,tau)\textless336·2\^{}(-352)\textless1/2. For 0\textless=u\textless1/2 the positive exponential series is bounded above by the geometric series1/(1−u)\textless2. Therefore

\[
 \prod_\tau e^{2r_{p,\tau}}<2.
\]

Apply the actual selected trace DO080(RT2), with its two compact factors and true bundle metric, to WM4. Taking the product over the f labels gives

\[
 |c(h)|<2\,2^{-931f/2}|N_E\eta|^{-1/2}
                  \le2\,2^{-931f/2}z^{-f}.         \tag{WM6}
\]

Here z is defined below, and the last inequality uses the same ideal-norm lower bound as DO083. The factor two in WM6 is the entire product of the near-point losses, not a factor two at each label.

\subsection{3. The exact affine difference lattice and its metric}\label{the-exact-affine-difference-lattice-and-its-metric}

Fix the actual support Sigma, the prescribed ordered source/target pairs, and the intrinsic set D of places where eta has an actual double pole. Let d=\textbar D\textbar. Retain the paired G rank-zero ideal J and write

\[
 \mathfrak l=\mathfrak J I_\Sigma I_D,\quad
 N_E\mathfrak l=D_GT^{-2(s+d)}=z^{2f},\quad
 y=R/z,\quad y_c=R/(zT^{2d/f}).                   \tag{WM7}
\]

The content ideal b of x satisfies N\_E b\textgreater=\textbar N\_E eta\textbar\^{}(1/2)z\^{}f.~For fixed eta and b, the first three positive coordinates lie in specified affine cosets of

\[
 \mathfrak c=\mathfrak b\cap\mathfrak J I_\Sigma,
 \qquad N_E\mathfrak c=N_E\mathfrak b\,T^{2d}.
                                                               \tag{WM8}
\]

These are K077(AP5--15), not a residue-density assertion. At a double place, the negative pole column and row have unit negative coordinates at both CM partners. The actual leading residue of x is therefore determined by eta, giving the affine coset modulo the simple-pole ideal. A zero fixed label cannot occur at a nonintegral type-II place. The other admitted zero labels remain included in the unmarked alternatives. Neither a general target real-projection identity nor an integral moving basis is used.

Use the actual metric

\[
 \|w\|_\eta^2=\sum_{\tau\in\Phi}|\tau w|^2/|\tau\eta|,
 \quad
 z_b=\left(\frac{N_E\mathfrak b}{|N_E\eta|^{1/2}}\right)^{1/f},
 \quad z_c=z_bT^{2d/f}\ge zT^{2d/f}.              \tag{WM9}
\]

The real6f-dimensional ball containing the three coordinates has radius sqrt(fR), by the exact rank-one norm equation sum\_(j=1)\^{}4 x\_j bar(x\_j)=eta bar(eta).

K076(MC3--9) computes the half-trace covolume and the actual dual for any fractional ideal in this metric. With

\[
 D_0=7^{140}29^{144}5^{84}13^{84},\qquad
 \operatorname{rd}(E)=2D_0^{1/f}<1500,
\]

its computation, applied to c, gives

\[
 \operatorname{covol}_\eta(\mathfrak c)=D_0z_c^f,
 \quad
 \lambda_1((A_\eta\mathfrak c)^*)
       \ge\frac{2\sqrt f}{\operatorname{rd}(E)\sqrt{z_c}},
 \quad
 \mu_\eta(\mathfrak c^3)<1300\sqrt{fz_c}.        \tag{WM10}
\]

Here A\_eta scales the tau coordinate by \textbar tau(eta)\textbar\^{}(-1/2), and the exact unscaled half-trace dual is2 Different\_(E/Q)\^{}(-1) bar(c)\^{}(-1). Thus the absolute different, factor two and inverse dual metric are retained. The last bound uses the already checked dimension/2 transference theorem and the orthogonal-product factor sqrt(3), with750sqrt(3)\textless1300. No improved transference theorem is needed.

Voronoi cells for a translate of c\^{}3 have the same volume and covering radius as those of c\^{}3. The cells about points in a ball of radius sqrt(fR) fit in the ball enlarged by the radius in WM10. Dividing that6f-dimensional ball volume by D0\^{}3 z\_c\^{}(3f), and then using the monotonic decrease in z\_c, gives

\[
 \#\{x_1,x_2,x_3\}
 \le A_0(\sqrt{R/z_c}+1300)^{6f}
 \le A_0(\sqrt{y_c}+1300)^{6f},
 \qquad
 A_0=\frac{[((22/7)f)]^{3f}}{(3f)!D_0^3}.        \tag{WM11}
\]

The22/7 is a rational upper comparison for pi. This volume estimate is for the exact affine difference ideal; it is not a claim that a fractional ideal shares the original ring's covering radius.

Retain also K077's independent disjoint-ball estimate (1+2sqrt(y\_c))\^{}(6f). The same original triple count is bounded by their minimum. The remaining content and relative-F norm counts are exactly reconstructed K080(RC4), as used in DO087. Hence the full matrix count for the fixed data is at most

\[
 28C\,y^{63}(1+2y)^f\,\mathcal A(y_c),\quad
 C=2^{72}/29^9,\quad
 \mathcal A(q)=\min\{(1+2\sqrt q)^{6f},
                    A_0(\sqrt q+1300)^{6f}\}.    \tag{WM12}
\]

All eta values, content ideals, affine triples and fourth relative norm solutions remain counted. There is still no additional negative-column or source-isometry fibre, because WM3 reconstructs the entire matrix from that column.

\subsection{4. Sum the finite patterns before replacing their norms}\label{sum-the-finite-patterns-before-replacing-their-norms}

There are334 unmixed relative G factors, with residue sizes q\_v in\{g0,g1\}, each occurring167 times. In this globally rank-one sector every finite reduction of N has rank zero or one. Let Z be the subset where its rank is zero. The paired ideal and norm are exactly

\[
 \mathfrak J_Z=\prod_{v\in Z}\mathfrak P_v\bar{\mathfrak P}_v,
 \qquad D_G=\prod_{v\in Z}q_v^2,
 \qquad M=Q^{f-1}=\prod_vq_v.                    \tag{WM13}
\]

The original finite bound is2D\_G\^{}(a/2)M\^{}(-a); the two mixed projectors keep their modulus bound one. No second enumeration of finite ranks or CM orientations is made.

For fixed s,d and pair data, multiply this finite factor, WM6 and either of the two counting bounds in WM12. Let F\_P(z) and F\_V(z) denote the resulting expressions, with all positive z-independent constants included. Their z-dependent factors are respectively

\[
 \begin{split}
 F_P(z)&\propto z^\kappa(z+2R)^f(\sqrt z+2\sqrt R\,T^{-d/f})^{6f},\\
 F_V(z)&\propto z^\kappa(z+2R)^f(\sqrt R\,T^{-d/f}+1300\sqrt z)^{6f},\\
 \kappa&=(a-1)f-63-4f=27321.
 \end{split}                                                    \tag{WM14}
\]

Thus for0\textless z\textless=z\_* both envelopes satisfy F\_j(z)\textless=F\_j(z\_\emph{)(z/z\_})\^{}kappa. Define

\[
 z_*=z_{s,d}=M^{1/f}T^{-(s+d)/f},\quad
 y_*=R/z_*,\quad y'_*=R M^{-1/f}T^{(s-d)/f}.      \tag{WM15}
\]

For the actual subset Z, WM13 gives exactly

\[
              z/z_*=\prod_{v\notin Z}q_v^{-1/f}.
\]

Since each actual subset has the same uniform original-pair count and the same two cap envelopes, summing all subsets yields

\[
 \begin{split}
 \sum_Z\min\{F_P(z_Z),F_V(z_Z)\}
 &\le\min\{F_P(z_*),F_V(z_*)\}
           \prod_v(1+q_v^{-\kappa/f})\\
 &<2\min\{F_P(z_*),F_V(z_*)\}.                 \tag{WM16}
 \end{split}
\]

For the last inequality, kappa/f\textgreater162 and both actual q\_v exceed2\^{}34. Thus every nontrivial missing-factor weight is less than2\^{}(-5508). The product is at most(1+2\^{}(-5508))\^{}334\textless2; this is also checked as an exact integer power. Counting an unrealizable subset only increases this upper bound. The factor2\^{}334 is not retained in addition to WM16.

This proves the pattern bound before the common cap; monotonicity alone after discarding the subset weights would not justify it. The finite term M\^{}(-a)T\^{}(a(s+d))z\_*\^{}(af) cancels exactly at the cap, as in DO087.

\subsection{5. Every original type, pair and integer marked size}\label{every-original-type-pair-and-integer-marked-size}

At each marked double place, type II has complete weighted ordered-pair cost1. At each unmarked place the full I,II,III costs sum to T+1. Therefore the exact containing count for each integer d is still binom(f,s) binom(s,d)(T+1)\^{}(s−d), as proved in K077(AP18). The intrinsic double set is unique for each actual eta. Simple-pole type-II places remain unmarked.

Combining this fact with WM6, WM12 and WM16 gives the complete fixed-d upper expression

\[
 \begin{split}
 U_{s,d}={}&2^3\,28C\binom fs\binom sd\,2^{-931f/2}
 (T+1)^{s-d}z_{s,d}^{-f}y_{s,d}^{63}
 (1+2y_{s,d})^f\,\mathcal A(y'_{s,d}).
 \end{split}                                                    \tag{WM17}
\]

The three factors two are respectively the one finite Euler correction, the entire weighted-pattern sum, and the entire product near-point loss. The matrix count retains the root-unit factor28. There is no2\^{}s endpoint replacement here: every actual integer d is separately counted with its exact binomial coefficient. Although each original envelope is log-convex, no convexity of their minimum is asserted or needed.

The full support-s contribution is bounded by sum\_(d=0)\^{}s U\_(s,d). WM17 is valid even when the cap itself is not realized by a nonempty actual lattice set, by WM14--16.

\subsection{6. Exact support22 margins}\label{exact-support22-margins}

Directed rational root, square-root and logarithm computations with the original T,Q and rational A0 give the following strict inverse binary margins. Each row already includes its support, binomial, pair, finite-pattern and real factors.

\begin{longtable}[]{@{}rlr@{}}
\toprule\noalign{}
d & Chosen triple envelope & Inverse binary margin \\
\midrule\noalign{}
\endhead
\bottomrule\noalign{}
\endlastfoot
0 & Metric volume & 968 \\
1 & Metric volume & 2581 \\
2 & Metric volume & 4195 \\
3 & Metric volume & 5810 \\
4 & Metric volume & 7425 \\
5 & Metric volume & 9036 \\
6 & Metric volume & 10617 \\
7 & Metric volume & 11986 \\
8 & Metric volume & 12444 \\
9 & Metric volume & 11641 \\
10 & Packing & 12019 \\
11 & Packing & 13301 \\
12 & Packing & 13524 \\
13 & Packing & 12583 \\
14 & Packing & 11281 \\
15 & Packing & 9921 \\
16 & Packing & 8555 \\
17 & Packing & 7188 \\
18 & Packing & 5821 \\
19 & Packing & 4455 \\
20 & Packing & 3089 \\
21 & Packing & 1723 \\
22 & Packing & 359 \\
\end{longtable}

The companion script verifies each unit-width logarithm interval and the entire table. In particular the all-double endpoint has strict margin359, while each of the other22 marked sizes has margin at least968. Consequently

\[
 \sum_{d=0}^{22}U_{22,d}
 <2^{-359}+22\,2^{-968}<2^{-358}.                 \tag{WM18}
\]

This proves WM2. The same two envelope expressions at s23,d0 and s23,d23 exceed one. They are recorded only as failed sufficient bounds; no sign, existence or nonvanishing conclusion follows from those failures.

\subsection{7. Chronology, verification and precise boundaries}\label{chronology-verification-and-precise-boundaries}

The exact certificate reads the actual canonical R06364-entry all-top and31-entry uniform arrays, checks their ordered suffix7548,8622,52215, and retains the pending DO0871649 addition. It appends358 once, giving the proposed chronological66/31 arrays. Both full integer dyadic sums remain below2\^{}(-230). Intervening zero-charge assemblies leave these arrays unchanged. Canonical integration is performed separately by the parent; this proof removes only the intersection with the then-current remainder and deletes no older charge.

The new steps are the exact affine metric application WM9--12, the complete weighted subset product WM14--16, and the multiplicative use of the actual packet distance WM4--6, followed by the full integer-d sum. Directly read dependencies are K076(MC1--10), K077(AP1--23), reconstructed K080(RC1--7), LIFT41(NP1--12 and its exact radius table), DO083(RO3--17), and DO087. The external transference input is the already checked Banaszczyk theorem as applied in K076; it is not re-proved or replaced. The selected-model input remains the checked DO080 normalized trace.

The script supplies exact arithmetic, metric-scaling and weighted-product diagnostics, not a rational-support enumeration or a formal proof of the cited source identities. No floating-point value is used to certify a bound. Author whole-proof/script and peer reviews are separate from those diagnostics.

The conclusion is the complete globally rank-one fixed-W0 sector at support22 under the original all-top conditions. Other structural sectors, the full signed inequality, actual global period, selected arithmetic-image content, literal Host/FamP construction, DP44 comparison and RH conclusion retain their own open requirements.

\clearpage

\section{AM-K092 --- primitive augmentation symbols after the full middle-mass cut}\label{am-k092-primitive-augmentation-symbols-after-the-full-middle-mass-cut}

\textbf{Canonical checked integration --- 2026-11-22T20:07:59.291751+00:00.} Root and independent whole-proof/script/source reviews and separate81946-assertion replays passed. The actual full middle-stabilizer cutoff can preserve the common11-unit witness and gives a primitive symbolic selected theta family on the integral tensor-augmentation source. Its contracted selected value equals the literal full-mass-divided contribution, with every coordinate pullback, ordinary derivative, Chern gauge, creation/period scalar and terminal projector retained. Full mass division is asserted only after the invariant middle coevaluation. The same input occurs in every middle-class summand; derivative annihilation and class cancellation remain possible.

\textbf{Algebraic-geometric translation.} This is an integral finite-module splitting on the unchanged theta-group and Picard/Hom data. The common transversal is conjugated by the actual class/polarization operator, and inverse Weil transport also transports evaluation functionals. No division by a possibly nonunit stabilizer order occurs in the integral section construction. The actual Weyl module and tensor augmentation are split before the selected derivative; faithful-flat descent concerns the symbolic cokernel only. Mumford(1966), §1 Proposition3 and Theorem2 pp295--298, remains carried through K089. The carried exact contextual anchors are FGA232 Theorem3.1 pp148--149; FGA236 Corollaries1.7/2.3 pp229/232; SGA7 I VII2.7/2.9 pp169--174 and VIII7.4(b) pp310--311; SGA3 I4.7.3/4.7.3.1; SGA6 VI Theorem3.3. Formal crystals and Cartier duality retain cote4 scan003 §§8.1.1--8.1.2, scans005--006 §§8.4.1--8.4.3 and cote54 scans005--007. No new lambda-ring/K-theory identity, l-adic purity or weight theorem, or CM-motive/L-function assertion is inferred. DP\_MULT Q3e and THETA3 G1/A7 retain the actual arithmetic comparison gates; CX\_GRH CX-ALG-024 and GCT\_PLETH T-052 retain their separate gates.

Author: ALP\_MACHINE \nolinkurl{/root/kummer_extension}, 22 November 2026 UTC.

\textbf{Author whole-proof and complete-script PASS; frozen for independent review.} This report concerns the exact finite cutoff already permitted as a proof device by LIFT34. It adds no level and replaces no actual input. The prescribed selected derivative and final class sum remain in the conclusion.

\subsection{1. Fixed actual data}\label{fixed-actual-data}

Retain one actual host, word, tuple, canonical twist and all raw G/T/Q functions of K038. Fix a projective-T/generic-G initial class c, one actual middle class f and one terminal decoration d.~Use the actual coefficient DVR R at11 and the standard11 class representatives of K039(IG3). Let \[
 D_T=\prod_{v\mid T}k_v^\times,
 \quad V_c=R[D_T]\otimes_R L_c,
 \quad V_c^{\rm ram}=\bigotimes_{v\mid T}I(k_v^\times)\otimes_R L_c.
 \tag{PM1}
\] These are the literal strict class module and its Mram summand from K090(MA1--6). The underlying initial lattice, coefficient frame, and non-T data are unchanged along its scalar decoration orbit. Write Psi\_s for the complete, framed, full-Y tensor of this fixed triple and scalar label s.

Let Gamma\_f be the actual finite middle rational stabilizer. At one chosen G place its reduction is injective: LIFT34(MC21a) proves this for the entire Gamma\_f using g−1\textgreater2fb. The second raw G support has full column rank, so right multiplication by that reduction acts freely on its possible residue matrices Z. All phases and splitting scalars are those of LIFT34(MC10--13); the complete uncut two-edge tensor is invariant under this honest diagonal middle action.

\subsection{2. A common transversal retaining all unit witnesses}\label{a-common-transversal-retaining-all-unit-witnesses}

K090(MA13--14) constructs points x\_s in the intrinsic finite position model satisfying \[
                     \Psi_s(x_t)=u\delta_{s,t},\qquad u\in R^\times.
 \tag{PM2}
\] The non-T components of x\_s can be taken identical for every s. In particular the second raw G coordinate has the same full-rank residue Z\_* for every s. Only the first T incidence solution is changed. The common unit u uses the exact full-Y G/Q contraction; neither an individual Q term nor a new T factor of Y is substituted.

The free Gamma\_f orbit of Z\_* admits Z\_* itself as its transversal representative. Complete it arbitrarily to a transversal on the whole second raw G support, and let C\_f be its0/1 indicator as in LIFT34(MC9--10). This is one choice for all s and all coefficient vectors P. In the raw split coordinates define \[
                       \Psi_s^{\rm cut}=C_f\Psi_s
 \tag{PM3}
\] and then apply the same actual class frames as LIFT34. Precisely, if W\_(c,f,d) is the common finite class/polarization transport, the framed cut is W\_(c,f,d) C\_f W\_(c,f,d)\^{}(-1) Psi\_s. Thus C\_f in the framed notation is this conjugated multiplier, not an untransported raw residue indicator. The class transport is common along the scalar decoration orbit because the underlying lattices and non-T frames are fixed. Since C\_f(Z\_*)=1 in the raw coordinates, equation PM2 becomes \[
                    \Psi_s^{\rm cut}(x_t)=u\delta_{s,t}.
 \tag{PM4}
\] If the common rational Schrödinger polarization differs from this local split model, transport the evaluation functionals by its inverse finite Weil operator. K089(PS11) and LIFT34§3 prove this operator and its inverse11-integral. PM4 then means an equality of integral linear functionals on the intrinsic finite module; it does not claim that a Fourier transform preserves point evaluations.

Thus the map delta\_s↦Psi\_s\^{}cut is split over R, with left inverse given by PM4 and division by the unit u. The actual augmentation summand PM1 is split too, independently of whether11 divides \textbar D\_T\textbar{} or \textbar Gamma\_f\textbar. This proof uses no average with a denominator.

\subsection{3. The selected symbolic map remains split}\label{the-selected-symbolic-map-remains-split}

The cutoff is at G and leaves the exact11 vacuum unchanged. Choose one even comparison denominator D prime to11 covering all finitely many cut inputs. K089's trace-one-normalized theta map Theta\_D is an integral isomorphism after its proof-only unramified splitting extension. It is the same actual map with its fixed scalar, not a newly normalized intertwiner.

Let Q\_cpl be the exact two-edge selected Weyl substitution of K091(SR11). Its integral retraction R\_cpl acts on the actual Weyl lattices; no hook inverse or ambient saturation assumption is used. Therefore \[
 (P_s)_s\longmapsto
       \sum_s \Psi_s^{\rm cut}\otimes\mathcal Q_{\rm cpl}(P_s)
 \tag{PM5}
\] has an integral left inverse: apply the finite coefficient functionals in PM4 and then R\_cpl. It remains split on V\_c\^{}ram. Apply Theta\_D to the finite factor. The resulting coefficient-labelled auxiliary sections with selected differential symbols have torsion-free cokernel by the same faithful-flat descent as K090(MA16--17). In particular a primitive augmentation input remains primitive at this symbolic cut-section stage.

The abstract integral Weyl lattice reduction is kept distinct from its possibly smaller ambient polynomial image modulo11, exactly as in K091§4. No claim is made that an ordinary derivative is a Hasse coefficient.

\subsection{4. Exact identification with the actual mass-divided selected class contribution}\label{exact-identification-with-the-actual-mass-divided-selected-class-contribution}

Let N\_(e,f,d)(P) be the complete full-Y numerator in LIFT34(MC2), and Q(P,v) its complete two-edge polynomial tested by an actual target dual-Weyl vector v. For the full-Gamma transversal above, LIFT34(MC21b) gives the exact equality \[
 \frac{1}{|\Gamma_f|}\operatorname{pairBK}
        \left(N_{(d_{c,s},f,d)}(P),v\right)
       =\mathcal J_{P,v}\left(\Theta_D\Psi_s^{\rm cut}\right),
 \tag{PM6}
\] where J\_(P,v) is the literal ordinary derivative defined by the same subD symbol Q(P,v), source cD coordinates, Chern/CM gauge and period-removed reading. The source creation and complex period scalars are common on both sides and are retained in exactly the LIFT34/K038 normalization. They have no assigned finite valuation here. The comparison-cover section is read in the original source formal coordinates using the exact unit-coordinate pullback of K089(PS13). Its D-dependent derivative factors are included in J, not set to one; D is prime to11, so this pullback and its inverse are integral at11. A derivative of the displayed cover section is not silently identified with the original derivative before this transport.

For clarity, the reason PM6 holds is the exact finite identity \[
 \Psi_s=\sum_{\gamma\in\Gamma_f}
             \omega_{\rm mid,f}(\gamma)\Psi_s^{\rm cut},
 \tag{PM7}
\] together with the invariance of the complete middle coevaluation creation vector and rational theta automorphy. This equality is applied after contracting with that invariant real vector. It does \textbf{not} assert Theta\_D Psi\_s=\textbar Gamma\_f\textbar Theta\_D Psi\_s\^{}cut as uncontracted auxiliary sections. Such an assertion is neither needed nor proved.

Linearity gives PM6 for every augmentation combination with coefficients P\_s. Thus each actual mass-divided middle-class map is the prescribed derivative evaluation of an explicitly split, primitive symbolic cut-section map. Its finite Hilbert norm remains exactly \[
           \|\Psi_s^{\rm cut}\|_f^2
                    =|\Gamma_f|^{-1}\|\Psi_s\|_f^2
 \tag{PM8}
\] by LIFT34(MC21b). This norm identity is not an inverse norm or a lower bound for PM6.

Different transversals can change the individual cut sections. Nevertheless their selected values in PM6 agree, since all equal the same actual mass-divided numerator. Choosing the transversal through the common unit witness therefore changes no actual arithmetic map.

\subsection{5. The remaining exact matrix and scope}\label{the-remaining-exact-matrix-and-scope}

For a fixed initial coarse class c and terminal decoration d, the original selected operator on an augmentation vector B is \[
                (\mathcal A B)(d)
                  =\sum_f\mathcal J_{f,d}\bigl(\mathcal S^{\rm cut}_{c,f,d}(B)\bigr),
 \tag{PM9}
\] where S\^{}cut denotes PM5 after the actual theta reading. The same B occurs in every summand. The summands are not independent input variables. Consequently the new primitive symbolic maps do not exclude annihilation by any individual J or cancellation among f.~Terminal brC, pairBK, genuine/generic projector denominators and terminal masses remain exactly those of K090(MA9--11).

The reduction achieved here is precise: a possible common nonunit factor in the proof-device cutoff is absent, and full middle-mass division has an exact representation by this primitive symbolic auxiliary family. The outstanding arithmetic question is the content of the prescribed derivative/class-sum composite PM9. No arbitrary jet family, independent class interpolation, nonzero selected column, unit minor, or sharp actual-image valuation follows.

\subsection{6. Sources and verification boundary}\label{sources-and-verification-boundary-2}

Direct inputs are LIFT34§§2--5 and MC21a--b; K038(CA15--18,CA22--27); K039(IG3--7); K089(PS8--13); K090(MA1--6,MA13--18); and K091(SR9--13). All are used at their exact fixed-full-Y, actual-frame and coefficient-lattice scopes. No new representation-theoretic or geometric classification theorem is invoked.

The companion certificate checks finite free-action cut identities, common unit evaluation and augmentation retractions over integers and modulo11, including groups whose orders contain11. It also tests invariant-functional mass division separately from false uncontracted section division. These models verify the finite algebra; the actual covariance and theta normalization are the carried proofs cited above.

\subsection{Author whole verification}\label{author-whole-verification-1}

At 2026-11-22T19:39:14.063767+00:00, I completed a whole read of PM1--9 and all six substantive sections, and the complete companion script. I reread the actual full-Gamma injectivity/free-action construction and contracted identity in LIFT34(MC9--14,MC21a--b), the identical non-T full-Y witness in K090(MA13--17), and the integral selected retraction in K091(SR9--13). I reconstructed the common transversal retaining that witness, its conjugation through the actual common class/polarization transport, and the integral left inverse before and after the trace-one theta map. I separately checked the comparison-cover unit-coordinate pullback and the fact that middle-mass division holds after the complete invariant contraction, not as a scalar identity of uncontracted auxiliary sections. The same initial augmentation input is retained across the actual middle-class sum. The final framed-cut and derivative-coordinate precision edits, status and this record were reread. The final separate certificate checks the finite algebra, including stabilizer orders divisible by11 and nontrivial phases; it does not evaluate an actual selected theta derivative.

\clearpage

\section{AM-R-066 --- a joint cross-vector ball removes small exceptional support twenty-two}\label{am-r-066-a-joint-cross-vector-ball-removes-small-exceptional-support-twenty-two}

\textbf{Canonical checked integration --- 2026-11-22T20:07:59.291751+00:00.} Root and independent whole-proof/script/source reviews and separate36-assertion replays passed. The entire actual IC8 rank-one/rank-two exceptional class at support22 has absolute contribution\textless2\^{}-160208, with no global-unipotence hypothesis. The original untouched160-coordinate block has rank one, so its Hilbert--Schmidt and operator norms agree and place all coordinates in one53760-dimensional ball. The exact ideal-product dual, covering radius, volume and original column/type/pair multiplicities are retained. The complete all-top support floor remains22; this alone does not exhaust support22.

\textbf{Algebraic-geometric translation.} This is a bound on the original product of160 fractional-ideal modules in its fixed positive Euclidean realization. The absolute different, half-trace dual and discriminant are unchanged; the rank-one cross map is a restriction/projection of the original labelled integral endomorphism. Banaszczyk(1993), Theorem2.2 p632, is carried through K076 with no new claim of an integral image inverse. The carried exact contextual anchors are FGA232 Theorem3.1 pp148--149; FGA236 Corollaries1.7/2.3 pp229/232; SGA7 I VII2.7/2.9 pp169--174 and VIII7.4(b) pp310--311; SGA3 I4.7.3/4.7.3.1; SGA6 VI Theorem3.3. Formal crystals and Cartier duality retain cote4 scan003 §§8.1.1--8.1.2, scans005--006 §§8.4.1--8.4.3 and cote54 scans005--007. No new lambda-ring/K-theory identity, l-adic purity or weight theorem, or CM-motive/L-function assertion is inferred. DP\_MULT Q3e and THETA3 G1/A7 retain the actual arithmetic comparison gates; CX\_GRH CX-ALG-024 and GCT\_PLETH T-052 retain their separate gates.

Author: \nolinkurl{/root}, GPT-6 Astra Pro; inherited effort not externally exposed; ALP\_MACHINE.22 November2026 UTC. New proof; author and independent whole checks and exact replay are recorded separately before promotion.

\subsection{1. The actual class and conclusion}\label{the-actual-class-and-conclusion}

Retain the original J56 all-top kernel and charged H133 cutoff, paired CM pencils, finite projectors, packet, selected Hilbert metric, source masses and phases. Put f=168,n=169,m=164. Consider the whole K069 exceptional class IC8 with r=rank\_E(A−I) in\{1,2\}, A=h\_SS, at support s=22. No global unipotence is assumed.

The theorem is

\[
 \boxed{\sum_{s=22,\ r\in\{1,2\},\ IC8}|w(h)|<2^{-160208}.}
 \tag{JB1}
\]

This extends the complete-class bound of R063 by one support. It adds its charge once on the intersection with the current remainder; it does not alone prove a new full support floor or the signed sufficient inequality.

Use exactly the R063(UC2--8) data. A padded four-subset B contains every changed original positive vertex; its complement S0 has160 coordinates. The untouched cross map C0:S0→W0 has rank zero or one. The data C0,v\_B,v\_W determine the full negative column by v0=−C0\emph{J0 v\_W, and the original source fibre is at most2\^{}5022. Every unmixed finite rank is at least two, so the unchanged whole-class Green factor is2Q\^{}(-335}167). There is no count of A or finite-rank patterns.

The zero and transverse cross alternatives retain the original R063 counts without modification. Only its plane alternative is sharpened below.

\subsection{2. The joint ball in the plane alternative}\label{the-joint-ball-in-the-plane-alternative}

In that alternative every transverse row vanishes and

\[
 C_0e_i=(u\xi_i,\beta_i),\quad
 u\cdot u=U=77Q^2-14,\quad
 \xi=(\xi_i)_{i=1}^{160}\in I_\Sigma^{160},\quad
 \operatorname N_E I_\Sigma=T^{-2s}.
 \tag{JB2}
\]

These are the actual original coordinates and paired inverse-prime ideal from R063. The integer Bezout relation for u gives the asserted ideal membership; no square root of U or principal generator is adjoined. A nonzero pivot and its original five-coordinate cross column reconstruct C0, exactly as in UC7--8.

At every CM embedding, C0 is the composition of h restricted to the Euclidean coordinate subspace S0 with the Euclidean projection to W0. Hence its operator norm is less than H=133. Because C0 has rank one, it has only one nonzero singular value and its Hilbert--Schmidt norm equals its operator norm. Therefore

\[
 U\sum_{i=1}^{160}|\tau\xi_i|^2
 \le\sum_i\|\tau C_0e_i\|^2
 =\|\tau C_0\|_{\rm HS}^2<133^2.
 \tag{JB3}
\]

The exact packet inequality U\textgreater133²2\^{}132 was already checked in R063 and is replayed here. Put M=2\^{}−66. Summing JB3 over the f complex labels proves that the \textbf{entire} vector xi lies in one real320f-dimensional Euclidean ball of radius M sqrt(f). There is no factor sqrt(160) in this radius. This is where rank one is used beyond the earlier separate-coordinate estimates.

\subsection{3. Exact ideal-product volume and covering radius}\label{exact-ideal-product-volume-and-covering-radius}

Set z\_s=T\^{}(-s/f) and retain the actual absolute discriminant

\[
 D_0=7^{140}29^{144}5^{84}13^{84},\qquad
 \Delta_E=2^{336}D_0^2,\qquad \operatorname{rd}(E)<1500.
 \tag{JB4}
\]

For the ordinary half-trace Euclidean pairing, I\_Sigma has covolume D0 z\_s\^{}(2f) and dual2 Different\_(E/Q)\^{}(-1) conjugate(I\_Sigma)\^{}(-1). The same explicit norm and arithmetic-geometric-mean calculation as K076(MC5--7) gives

\[
 \lambda_1(I_\Sigma^*)\ge
       \frac{2\sqrt f}{\operatorname{rd}(E)z_s}.
 \tag{JB5}
\]

Indeed a nonzero dual element has absolute norm at least2\^{}(2f)/(Delta\_E N\_E(I\_Sigma)); applying arithmetic-geometric mean to its f squared complex absolute values gives JB5. This retains the absolute different and half-trace factor two.

The already checked Euclidean transference inequality mu(L)lambda1(L*)\textless=dim\_R(L)/2, applied in dimension2f, yields mu(I\_Sigma)\textless750 sqrt(f)z\_s. Taking160 orthogonal copies and choosing nearest lattice points separately gives

\[
 \mu(I_\Sigma^{160})
 \le\sqrt{160}\,\mu(I_\Sigma)<9500\sqrt f\,z_s,
 \quad
 \operatorname{covol}(I_\Sigma^{160})=D_0^{160}z_s^{320f}.
 \tag{JB6}
\]

The integer inequality750²·160\textless9500² is explicit. No original-ring covering radius is substituted for a fractional-ideal cover.

The Dirichlet cells of points in the ball of section2 have disjoint interiors and fit in the ball enlarged by JB6. Its real dimension is320f, with half-dimension160f. Dividing its exact ball volume by the covolume gives

\[
 \#\xi\le A_{160}(M/z_s+9500)^{320f},\qquad
 A_{160}=\frac{[((22/7)f)]^{160f}}{(160f)!D_0^{160}}.
 \tag{JB7}
\]

The rational22/7 bounds pi from above. This counts all zero coordinates and imposes no extra finite congruence. Integer arithmetic verifies

\[
 A_{160}<2^{-369387},\qquad
 T^{22}>(9500\,2^{98})^f,\qquad
 (1+2^{-32})^{320f}<2.
 \tag{JB8}
\]

At s=22 the second inequality says9500 z\_s/M\textless2\^{}−32. Hence JB7--8 imply

\[
 \#\xi<2^{1-369387-66(320f)}T^{320s}.
 \tag{JB9}
\]

Every comparison in JB8 is between explicit positive integers or rationals. No floating-point volume estimate is used.

\subsection{4. Every original factor and the full support22 union}\label{every-original-factor-and-the-full-support22-union}

Retain the complete R063(UC11--14) pivot and column counts, choice factors and incidences. Thus the common constant is

\[
 C=1+168+25+8+2+5022+2(14861)+10835=45783.
 \tag{JB10}
\]

It includes the single Green Euler factor, support choices, padded B choices, all160 pivots, all three transverse tests, original source fibre, two W0 vector counts and four original positive entries. The complete ordered-pair/absolute-incidence factor is below2\^{}(997s); both mixed factors and the real selected coefficient retain modulus at most one. The T exponents in those common vector counts total28s. No extra A, orientation, type or rank-pattern multiplier is introduced.

With L\_T=1019333/1024 and L\_Q=287843/4096, the zero and transverse case exponents are the unchanged

\[
 E_0(s)=C+(28L_T+997)s-335\cdot167L_Q,
 \qquad E_A(s)=E_0(s)+45(320f).
 \tag{JB11}
\]

Substituting only JB9 for UC10 gives the plane exponent

\[
 E_B^{\rm joint}(22)=C+1-369387-66(320f)
       +(348L_T+997)22-335\cdot167L_Q
       =-656222427/4096<-160210.
 \tag{JB12}
\]

Exact rational calculation also gives E\_A(22)\textless−160210 and E\_0(22)\textless E\_A(22). Consequently all three complete case contributions are below2\^{}−160210; their union is less than3·2\^{}−160210\textless2\^{}−160208. This proves JB1. The padded four-subset count already covers both block ranks simultaneously, so there is no additional rank1/rank2 factor.

This argument changes only the plane cross-vector count in a previously proved exhaustive reconstruction. All actual targets, zero labels, paired primes, original Hermitian column constraints, conservative fibres and source masses remain present. The analogous plane envelope at s23 exceeds one; that is only a failure of this sufficient bound.

\subsection{5. Exact errors, algebraic geometry and limits}\label{exact-errors-algebraic-geometry-and-limits}

The chronological predecessor is \nolinkurl{checks/do088_joint_remainder_error.json}, with66 all-top and31 uniform entries. Append160208 once; retain all earlier charges including DO0871649 and DO088358. The companion checks both complete dyadic sums below2\^{}−230. The full all-top support floor remains22 until every other support22 structural class has a separately checked bound and complete assembly. The independent uniform floor remains11.

The new fact concerns the rank-one realization of an original Hermitian-module map and the full Euclidean volume of its original fractional-ideal coordinate module. No moving integral basis, new CM extension, arithmetic-image inverse or theta-input replacement appears. The carried module/determinantal anchors are SGA3 I4.7.3/4.7.3.1 and SGA6 VI Theorem3.3. Picard/biextension anchors remain FGA232 Theorem3.1 pp148--149; FGA236 Corollaries1.7/2.3 pp229/232; SGA7 I VII2.7/2.9 pp169--174 and VIII7.4(b) pp310--311. Crystals/Cartier duality retain cote4 scan003 §§8.1.1--8.1.2, scans005--006 §§8.4.1--8.4.3 and cote54 scans005--007. These are contextual locators, not proofs of the new bound. No new scheme identification, lambda-ring/K-theory identity, l-adic sheaf/weight theorem, CM-motive comparison or motivic L-function statement follows.

DP\_MULT Q3e and THETA3 G1/A7 gain precisely the stated containing-set removal. The signed inequality, actual admissible mixed-CM family, selected arithmetic content, prescribed global Bessel/Mram nonvanishing, actual Read/032 pairing and changed-family DP44 terms remain open. CX\_GRH CX-ALG-024 and GCT\_PLETH T-052 keep their separate arithmetic requirements. R063 is used in its full checked original-coordinate scope; K076's exact dual calculation and the already accepted Banaszczyk(1993), Theorem2.2,p632, supply the indicated lattice conventions. Whole proof reading is separate from integer diagnostics. No Lean build or archive audit is performed.

\clearpage

\section{AM-LIFT-073 --- the actual spherical Bessel test away from G and T}\label{am-lift-073-the-actual-spherical-bessel-test-away-from-g-and-t}

\textbf{Canonical checked integration --- 2026-11-22T20:10:49.013096+00:00.} Root and independent whole-proof/source/script reviews and separate23,275-assertion replays passed. Historical pending text is preserved below. Under the genuine actual representation/lattice/character hypotheses, every prescribed finite away-G/T spherical unit-vector Bessel factor is alpha\_v=c\_H,v\emph{c\_N,v}Lscript\_v, and alpha\_v\^{}natural=c\_H,v*c\_N,v\textgreater0. The actual fractional different-scaled flag gives primitive coordinates even at absolute ramification. Every actual dyadic place is CM-split; a new full spherical descending-row support argument supplies the split calculation in all residue characteristics. The odd nonsplit calculation uses the precise primitive-coordinate literature formula. Together with LIFT071/072 this determines the finite local tests conditionally; global period/nonzero, source existence, integral image, Read/STiesR,032 value and DP44 remain open.

\textbf{Algebraic-geometric translation.} The unchanged selfdual Hermitian module over the finite etale quadratic algebra defines a smooth reductive unitary model, descending from GL\_n after that base change; its integral points are the stated hyperspecial compact. The different-scaled fractional flag and actual unit marked-line isometry preserve the prescribed Picard/Hom correspondence. This is a local analytic normalization and integral-model statement, not an integral theta-image inverse or a comparison of global motives. Exact contextual anchors are SGA3 I4.7.3/4.7.3.1 and SGA6 VI Theorem3.3; FGA232 Theorem3.1 pp148--149; FGA236 Corollaries1.7/2.3 pp229/232; SGA7 I VII2.7/2.9 pp169--174 and VIII7.4(b) pp310--311. Formal crystals/Cartier duality retain cote4 scan003 §§8.1.1--8.1.2, scans005--006 §§8.4.1--8.4.3 and cote54 scans005--007. No new lambda-ring/K-theory equality, l-adic sheaf/weight assertion, CM-motive identification or motivic L-function theorem is obtained. DP\_MULT Q3e and THETA3 G1/A7 lose the stated finite local-test gate under their retained genuine inputs; CX\_GRH CX-ALG-024 and GCT\_PLETH T-052 keep their distinct global arithmetic gates.

Author: GPT-6 Astra Pro, ALP\_MACHINE \nolinkurl{/root/lift_accounting}. 22 November 2026 UTC.

\textbf{Author whole-proof/source/script review PASS.} Final hashes and the separate exact replay are recorded in the companion author-review files. Independent review and canonical promotion are separate events.

This is a local source application for the genuine standard-lattice realization of LIFT060. It retains the different-scaled flag, the physical trace character, the norm-minus-T line, both local vectors and the actual Haar factors. It proves the normalized spherical local factor at every finite place away from G and T on the K030/K032 hosts. It does not construct a full SizesAt/FamP/Read instance on those fixed hosts, a global period, or the requested 032 value. The same statement applies to another genuine realization only under the hypotheses stated below, including the local splitting alternative at dyadic places.

\subsection{1. Exact scope and notation}\label{exact-scope-and-notation}

Let E/F be the CM extension in LIFT060, with finite relative discriminant one and principal extended absolute different

\[
 \mathfrak D_{E/\mathbb Q}
       =\mathfrak D_{F/\mathbb Q}\mathcal O_E=d\mathcal O_E.
 \tag{SC1}
\]

The two actual fields are Q(zeta28)B and Q(zeta28)B(sqrt5,sqrt13), where B is the real cyclic degree-seven subfield of Q(zeta29). These are the fields of K030/K032; F is their fixed field under complex conjugation. The hypothesis is relative unramifiedness, not unramifiedness over Q.

Put k\textgreater=2 even, rho\textgreater=0 and n=k+2rho+1. The larger rational Hermitian space and its lattice are

\[
 Y=H^\rho\perp(I_k\oplus-1),\qquad
 \Lambda=\mathcal O_E^n
 \tag{SC2}
\]

in the exact integral hyperbolic frame of LIFT060. The notation O\_E\^{}n here is the transported original standard lattice, not an arbitrary lattice in its rational genus. Write G=g0g1. For the chosen word let T be its prescribed odd rational prime and let the actual algebraic integer pi satisfy pi bar(pi)=T. Write

\[
 A=(T+1)/2,\quad B_T=(T-1)/2,\quad
 z_T=B_Tp_1+A z_* .
 \tag{SC3}
\]

Thus h(z\_T,z\_T)=-T and z\_T is the literal marked vector of LIFT060(HF5). Denote its positive complement by H=U(z\_T\^{}perp in I\_k+(-1)), embedded in U(Y) and fixing the whole hyperbolic flag. Its lattice and finite vectors are those transported from the initial rank-k space by LIFT030/LIFT060, not chosen afresh.

At a finite relative place v not dividing GT, let Pi\_v and sigma\_v be the actual irreducible representations in the larger and initial slots. Assume that they are tempered and have the prescribed nonzero spherical vectors at the standard compact subgroups. Divide these two vectors by their actual Hilbert norms; call the resulting unit vectors x\_v,y\_v. The chosen bottom subspace of LIFT069 and the full-SizesAt initial theorem LIFT070 supply these representation hypotheses in their respective recorded scopes. This sentence does not enlarge LIFT069 to every bottom constituent or infer full SizesAt on the fixed small fields.

For a local Bessel model use the same slot/contragredient convention as LIFT071/LIFT072. Let alpha\_v be its regularized matrix-coefficient period on x\_v,y\_v. Put

\[
 \mathcal L_v=\Delta_{n,v}
 \frac{L(1/2,\operatorname{BC}(\Pi_v)\times
                    \operatorname{BC}(\sigma_v))}
 {L(1,\Pi_v,\mathrm{Ad})L(1,\sigma_v,\mathrm{Ad})},\qquad
 \Delta_{n,v}=\prod_{j=1}^{n}L(j,\epsilon_{E_v/F_v}^{j}).
 \tag{SC4}
\]

The standard L-factor uses the fixed two-slot convention of the period; changing an individual slot to its contragredient changes that convention as well. Set alpha\_v\^{}natural=alpha\_v/mathcal L\_v. All measures are specified in section 4 before a numerical assertion is made.

\subsection{2. The actual flag has a primitive local character}\label{the-actual-flag-has-a-primitive-local-character}

Choose the integral hyperbolic pairs e\_a\^{}0,f\_a\^{}0 of LIFT060 section 2. They exist even at residue characteristic two: if w+bar(w)=1 in O\_E, the vectors (1,1) and (w,-bar(w)) in I\_1+(-1) are isotropic, pair to one and have determinant -1. Surjectivity of the integral trace follows locally from the finite etale trace and Nakayama and hence globally as an equality of O\_F ideals. This construction has no dyadic denominator.

The literal source flag and its fractional ideals are

\[
 e_a=d^{\rho-a}e_a^0,\qquad
 f_a=\bar d^{-(\rho-a)}f_a^0,\qquad
 \mathfrak a_a=(d)^{-(\rho-a)}\quad(0\le a<\rho).
 \tag{SC5}
\]

Their ideal sums give exactly Lambda because d/bar(d) is a unit. Source \nolinkurl{aaD} is not replaced by the unscaled flag. In the unscaled integral matrix coordinates, conjugate-linearity in the first Hermitian slot gives

\[
 h(f_a,u e_{a+1})=d^{-1}h(f_a^0,u e_{a+1}^0),\qquad
 h(f_{\rho-1},u z_T)=d^{-1}h(f_{\rho-1}^0,u z_T).
 \tag{SC6}
\]

The first formula concerns a\textless rho-1. These are precisely the summands of \nolinkurl{psiFin}, with its literal denominator G. Consequently the actual character is the product of primitive simple-coordinate characters, not the unscaled absolute trace character on those coordinates.

More precisely, write L=F\_v and K=E\_v. Choose a generator delta\_v of the absolute local different D\_(L/Q\_p). The relative extension K/L is etale unramified, so d/delta\_v is an O\_K unit. The character

\[
 \psi_L(t)=\psi_{\mathbb Q_p}(\operatorname{Tr}_{L/\mathbb Q_p}
                                      (\delta_v^{-1}t))
 \tag{SC7}
\]

has conductor O\_L, by the defining trace-duality property of the different. The source character on each simple K coordinate is psi\_L composed with Tr\_(K/L), with unit coefficient (delta\_v/d)/G. This character has conductor O\_K. For split K=L+L the two factors are separately primitive; no factor two is inserted by the trace. The two reflected flag chains have the signs computed in LIFT060(HF12); minus one is a unit also at residue characteristic two.

The actual rational isometry on the p\_1,z\_* plane is

\[
 g_T=\begin{pmatrix}A/\pi&B_T/\bar\pi\\
                    B_T/\pi&A/\bar\pi\end{pmatrix},
 \quad g_Tz_* =z_T/\bar\pi,\quad \det(g_T)=1.
 \tag{SC8}
\]

It is the identity on all other coordinates. At v not dividing T, pi is a unit and both g\_T and its inverse are integral. Thus it belongs to the actual larger compact, carries the initial standard lattice to the literal positive complement, and changes only the final character coefficient by the unit bar(pi). It does not change either spherical vector. The initial embedding has not been replaced by an unrelated complement.

All remaining unit character coefficients can be removed by an integral unitary flag torus. Explicitly, if the simple K-coordinate coefficients are gamma\_1,\ldots,gamma\_rho in O\_K\^{}times, act by a\_i on e\_i\^{}0 and by bar(a\_i)\^{}(-1) on f\_i\^{}0, trivially on the core. Choose a\_rho=gamma\_rho\^{}(-1) and recursively a\_i=gamma\_i\^{}(-1)a\_(i+1). Then gamma\_i a\_i/a\_(i+1)=1 and gamma\_rho a\_rho=1. This element fixes H pointwise and is in the same compact. Its conjugation has absolute Jacobian one on every root. Thus the physical model is compact-conjugate to the primitive-coordinate spherical model, with no change of vector, level or norm. For rho=0 the radical is trivial and this paragraph is empty.

\subsection{3. Hyperspecial groups and the dyadic alternative}\label{hyperspecial-groups-and-the-dyadic-alternative}

Since O\_K/O\_L is finite etale (or split) and the Hermitian lattice Lambda is unimodular, its unitary group scheme is reductive and smooth: after that finite etale base change it is GL\_n, and these properties descend. Its O\_L points form a hyperspecial maximal compact K\_n.~The marked line is an integral orthogonal direct summand away from T, since its norm -T is a unit. Its complement lattice is selfdual. Equation SC8 identifies its compact K\_k with the original initial maximal compact, embedded in K\_n.

At odd nonsplit places the groups in the nested pair are quasi-split. One elementary lattice check is that the norm map O\_K\^{}times to O\_L\^{}times is onto: residue norms are onto and successive lifts use the surjective residue trace. Hence every pair of unit diagonal entries is an integral hyperbolic plane. In particular I\_k is split for even k. A unit marked vector in its orthogonal odd-dimensional extension is the standard unramified Bessel position. If a hyperbolic-plane description containing that vector is desired, use an orthogonal positive unit vector and divide by two; this operation is used only at odd residue characteristic. It is not part of SC5 or SC8.

Every dyadic place of either actual host is split in E/F. Indeed E contains Q(zeta7). The decomposition group of 2 in Q(zeta7)/Q is generated by 2 modulo7, namely \{1,2,4\}, whereas complex conjugation is 6 modulo7. No prime over2 is fixed by conjugation in that subfield. A prime of either larger CM field fixed by conjugation would restrict to such a fixed prime, which is impossible. In a quadratic extension this is exactly splitting. This is the same argument recorded in DO022 section4, and it survives the real B and quadratic composita. Thus the needed alternatives are exhaustive:

\[
 K=L\oplus L\quad\hbox{(any residue characteristic)},\qquad
 K/L\text{ unramified quadratic, with odd residue characteristic}.
 \tag{SC9}
\]

No dyadic assertion is inferred for an arbitrary CM field not satisfying this alternative. Absolute ramification at 2,5,7,13 or29 does not by itself remove a place from SC9; SC1, SC6 and SC7 are what handle that ramification.

\subsection{4. Haar conventions and the standard odd nonsplit calculation}\label{haar-conventions-and-the-standard-odd-nonsplit-calculation}

Use probability Haar on K\_k. On the Bessel radical N\_B use the product Haar in an integral root parametrization of the smooth unramified model, each additive O\_L or O\_K root parameter having volume one. These are selfdual additive measures for the primitive trace pairings in SC7. The integral root parametrization has unit triangular transition Jacobians; N\_B(O\_L) has volume one. At a split place it is the literal upper/lower matrix-root measure of LIFT071. All changes in sections2--3 are by compact integral isometries and preserve these measures.

If the original source uses another positive Haar normalization, define its exact constants by

\[
 dh_{\rm source}=c_{H,v}\,dh_0,\qquad
 du_{\rm source}=c_{N,v}\,du_0,
 \quad c_{H,v},c_{N,v}>0.
 \tag{SC10}
\]

Equivalently these constants are its volumes of the displayed integral compact subgroups. They are retained, not declared equal to one by an independent local isometry. Regularization respects this scaling.

At the odd nonsplit places, the standard primitive-coordinate spherical calculation gives

\[
                       \alpha_{v,0}=\mathcal L_v.
 \tag{SC11}
\]

The precise primary calculation is Liu (2016), section3.3, equation(3.8), Lemmas3.6--3.9 and the concluding proof of Theorem2.2 on printed pp.149--154, together with Corollaries6.2 and6.5 on printed pp.186--191. The latter uses odd residue characteristic, as explicitly stated at the beginning of section6.2. In those calculations the spherical vectors have pairing one, both maximal compacts have volume one, the root character is primitive, and the marked-vector norm is a unit. These are the coordinates and measures proved above. The Gindikin--Karpelevich/Iwahori comparison in Lemmas3.6--3.9 and the Whittaker--Shintani value in Corollary6.5 give exactly the L-factor SC4. Temperedness supplies the absolute convergence of the outer integral of the stable Fourier coefficient, by Liu Theorem2.1. We do not assert absolute convergence of the bare double integral.

\textbf{Source precision.} Liu's printed condition(5) on p.139 says, literally, that psi(uu')=1 for all integral u' if and only if u is integral. A generic character on a nonabelian radical cannot have that literal kernel property. The page image was independently inspected and has no omitted pairing or commutator. We do not invoke that sentence as a valid condition, silently replace it, or deduce it from SC7. We use the explicit primitive-coordinate calculation in section3.3 and section6 at the model just specified. Its Fourier/Iwahori steps use primitive simple-coordinate characters and their triviality on integral root parameters; those properties were established directly in SC6--SC7. Beuzart-Plessis--Chaudouard (2025), section8.7, independently uses the same spherical Whittaker--Shintani formula, but its set S contains all places above2 (section8.1.6). That newer almost-all argument is not used to claim a dyadic result. The split computation below has no residue-characteristic restriction and resolves every dyadic place actually needed here.

\subsection{5. Direct split support lemma for the spherical vector}\label{direct-split-support-lemma-for-the-spherical-vector}

Suppose K=L+L, with residue cardinal q. Fix a conductor-O\_L additive character psi. Let W be any right GL\_n(O\_L)-invariant function in a generic Whittaker model, with standard left character psi on the upper unipotent group. For h in GL\_k(L) and X in Mat\_(rho,k)(L), set

\[
 A_\rho(h,X)=
 \begin{pmatrix}h&0&0\\X&I_\rho&0\\0&0&I_{\rho+1}\end{pmatrix}.
 \tag{SC12}
\]

Then

\[
 W(A_\rho(h,X))=
  1_{\operatorname{Mat}_{\rho,k}(\mathcal O_L)}(X)
                          W(\operatorname{diag}(h,I_{2\rho+1})).
 \tag{SC13}
\]

This assertion concerns the full spherical vector. It is not an application of the compact Kirillov support statement GB12 of LIFT072.

Here is a descending-row proof. Initially choose j=k+rho+1 and i\textless=k. For every t in O\_L, the matrix v=1+tE\_(i,j) is in GL\_n(O\_L). The conjugate A v A\^{}(-1) is upper unipotent, since the j-th row of A\^{}(-1) is standard. Its only possible contribution to the standard Whittaker character is its (j-1,j) entry, equal to tX\_(rho,i). Therefore

\[
 W(A)=\psi(tX_{\rho,i})W(A)\quad(t\in\mathcal O_L).
 \tag{SC14}
\]

If W(A) is nonzero, primitive trace duality forces X\_(rho,i) integral for every i. Right multiply A by the integral lower shear 1-sum\_i X\_(rho,i)E\_(k+rho,i). It removes precisely the last X row and leaves h and the other rows unchanged, so it preserves W(A). Now choose j=k+rho. The previously cleared inverse row is standard, and all nonzero entries in A e\_i have row index\textless j. The same upper-unipotent argument forces the penultimate X row integral. Descend to j=k+2. Every row is integral. If X is integral, the single right lower block shear converts diag(h,I) into A, proving the nonzero-side equality as well. If W(A) is zero and X is nonintegral, the displayed formula already holds. This completes the proof, including rho=0 and residue characteristic two. It imposes no integrality condition on h.

\subsection{6. Exact split central integral and spherical Hilbert norms}\label{exact-split-central-integral-and-spherical-hilbert-norms}

Let W,V be the spherical Whittaker vectors of Pi\_v,sigma\_v, with opposite Whittaker characters and W(1)=V(1)=1. Put

\[
 \kappa_r=\prod_{j=1}^r(1-q^{-j}),\quad
 L_G=L(1,\Pi_v\times\Pi_v^\vee),\quad
 L_H=L(1,\sigma_v\times\sigma_v^\vee),\quad
 L_0=L(1/2,\Pi_v\times\sigma_v).
 \tag{SC15}
\]

Use raw GL\_r Haar \textbar det g\textbar\^{}(-r) times additive matrix Haar of O\_L-volume one; its compact volume is kappa\_r. Let theta\_r be the invariant Whittaker norm obtained by integrating \textbar W(diag(a,1))\textbar\^{}2 over N\_(r-1)\GL\_(r-1) in that raw quotient measure. These are exactly the measures in LIFT071(HN4--8).

SC13 and additive O\_L-volume one reduce the central integral to

\[
 Z_\rho^{\rm raw}(W,V)
   =\int_{N_k\backslash GL_k}
      W(\operatorname{diag}(h,I_{2\rho+1}))V(h)
          |\det h|^{-\rho}\,d\bar h^{\rm raw}
   =\kappa_k L_0.
 \tag{SC16}
\]

For completeness the second equality is a direct Schur Cauchy computation. In probability quotient Haar, Iwasawa decomposition and the spherical Whittaker formula reduce the integral to dominant lambda\_1\textgreater=\ldots\textgreater=lambda\_k\textgreater=0. The ratio of the GL\_n and GL\_k square-root moduli is \textbar det h\textbar\^{}((n-k)/2); multiplication by \textbar det h\textbar\^{}(-rho) leaves \textbar det h\textbar\^{}(1/2). Thus the sum is

\[
 \sum_{\ell(\lambda)\le k}
     s_\lambda(\alpha_1,\ldots,\alpha_n)
     s_\lambda(\beta_1,\ldots,\beta_k)q^{-|\lambda|/2}
       =\prod_{i,j}(1-\alpha_i\beta_jq^{-1/2})^{-1}.
 \tag{SC17}
\]

The diagonal formula is, for example, Miyauchi (2014), Theorem4.1 and Remark4.3, including its conductor-zero case; its setup is every characteristic-zero nonarchimedean field, not just odd q. Tempered spherical Satake parameters have modulus one. The sums converge absolutely, since the dimensions of the fixed-rank Schur representations grow polynomially with the partition entries. Hence no meromorphic specialization is needed at s=1/2. The raw quotient differs by exactly kappa\_k, giving SC16.

The same calculation gives

\[
             \|W\|_{\vartheta_n}^2=\kappa_n L_G,
             \|V\|_{\vartheta_k}^2=\kappa_k L_H.
 \tag{SC18}
\]

One can see the factor explicitly: the quotient has last-zero partitions, whereas the full Cauchy sum is over all nonnegative partitions. Subtracting the last entry produces a geometric series zeta\_q(r), because the central character is unitary. Thus the probability quotient norm is L(1,pi\_r x pi\_r\^{}vee)/zeta\_q(r); multiplication by kappa\_(r-1) gives kappa\_r. This proves the two identities with the actual spherical vectors of value one, rather than choosing a norm to cancel the period.

The calibrated matrix-period identity LIFT071(HN8) applies to these tempered generic GL representations. Divide W,V by the square roots in SC18. Changing its raw outer H Haar to probability Haar contributes kappa\_k\^{}(-1), while the root Haar is already the one specified in section4. We obtain

\[
 \alpha_{v,0}
  =\kappa_k^{-1}
       \frac{|\kappa_kL_0|^2}{\kappa_nL_G\,\kappa_kL_H}
  =\kappa_n^{-1}\frac{|L_0|^2}{L_GL_H}
  =\mathcal L_v.
 \tag{SC19}
\]

For the split unitary group its standard base-change Rankin--Selberg factor is \textbar L\_0\textbar\^{}2 in these unitary two-slot conventions, its two adjoint factors are L\_G,L\_H, and Delta\_n=kappa\_n\^{}(-1). This proves exactly SC11 at every split place, including all dyadic places of the actual hosts. It does not import the conductor-one small-vector norm of LIFT071(HN17); both vectors here are spherical.

\subsection{7. Actual finite level and conclusion}\label{actual-finite-level-and-conclusion}

At v away from G, the LIFT060 finite vectors e0,f0,z0,bw,bwd vanish, so the literal Hida \nolinkurl{K2cong} conditions are automatic, \nolinkurl{K2set} is the standard lattice compact, and \nolinkurl{chi2f} is one. Source \nolinkurl{e_eq} is its probability idempotent, with the actual ambient finite Haar volume retained. The level translation in LIFT060 is supported only at G. At the initial stage, \nolinkurl{MfV=ColV(0)} has the strict T covector condition, the prescribed generic G character, and maximal level elsewhere; this is the exact j=0 level audited in K038(CA27a) and LIFT027. Thus at v away from GT the two unchanged vectors really have the compact invariance used above. The later Q levels of the vertical tower do not add a bottom j=0 condition here.

Combining SC10, SC11 and SC19 proves the local theorem:

\[
 \boxed{\quad
   \alpha_{v,\mathrm{source}}
      =c_{H,v}c_{N,v}\mathcal L_v>0,
   \qquad
   \alpha_{v,\mathrm{source}}^{\natural}
      =c_{H,v}c_{N,v}
   \quad(v\nmid GT).\quad}
 \tag{SC20}
\]

In the displayed probability/integral-root normalizations the normalized factor is exactly one. Positivity and finiteness in the nonsplit case follow also from the spherical formula and the tempered positive Bessel form. The local Euler factors at positive real arguments have no poles or zeros for unramified tempered parameters. No infinite product of the c constants is introduced without its global Haar convention.

Together with LIFT071 at T and LIFT072 at G, this determines the prescribed finite local test factors whenever their respective genuine representation and source hypotheses hold. It proves neither a global central L-value nor a nonzero global Bessel/Mram period. The real selected vector, actual Read/STiesR instance, compatibility with the full literal host/size package, and requested arithmetic row/value remain separate.

\subsection{8. Sources and verification boundary}\label{sources-and-verification-boundary-3}

Pinned source commit: \nolinkurl{7cbf7a0c32fca1badf05931879a37c3f8a2b03b4}.

\begin{itemize}
\tightlist
\item
  \nolinkurl{Formalism/LocalToGlobal/SiegelLatticeTheory/HaarTransversals/HidaSetting.lean}: Nfin/psiFin106--114, K2cong/K2set116--123, chi2f125--128, aaD172--173, and e\_eq264--265. These determine the physical flag character, compact and finite projector.
\item
  \nolinkurl{Formalism/LocalToGlobal/Columns/ColForms.lean}93--115 and \nolinkurl{RecLevel.lean}49--52: actual initial level and adelic right-character convention; K038(CA27a) supplies its strict T versus projective interpretation.
\item
  LIFT060(HF1--8,HF12--17), LIFT030(NL3--6), K030/K032 and DO022 section4: exact standard lattice, different, marked-line transport and dyadic splitting. Their hypotheses are retained, not reconstructed from a signature alone.
\item
  LIFT069/LIFT070: respectively the chosen bottom and full-SizesAt initial temperedness premises; LIFT071(HN4--16): the already calibrated split Bessel identity and its raw Haar convention.
\item
  Liu, Y. (2016). Refined global Gan--Gross--Prasad conjecture for Bessel periods. \emph{Journal für die reine und angewandte Mathematik, 717}, 133--194. \url{https://doi.org/10.1515/crelle-2014-0016}. Directly read published pp.139--140,149--154,186--191. Public publisher-version copy: \url{https://scispace.com/pdf/refined-global-gan-gross-prasad-conjecture-for-bessel-53ylwvqnk8.pdf}. The condition(5) issue is explicitly separated above.
\item
  Miyauchi, M. (2014). Whittaker functions associated to newforms for GL(n) over p-adic fields. \emph{Journal of the Mathematical Society of Japan, 66}(1), 17--24. \url{https://doi.org/10.2969/jmsj/06610017}. Directly read section2, Theorem4.1 and Remark4.3. \url{https://www.jstage.jst.go.jp/article/jmath/66/1/66_17/_pdf}.
\item
  Beuzart-Plessis, R., \& Chaudouard, P.-H. (2025). The global Gan--Gross--Prasad conjecture for unitary groups. II. From Eisenstein series to Bessel periods. \emph{Forum of Mathematics, Pi, 13}, e16. \url{https://doi.org/10.1017/fmp.2025.8}. Directly read sections8.1.6 and8.7; its exclusion of dyadic places is retained.
\end{itemize}

The companion certificate checks exact trace/conductor models, fractional-flag exponents, the moved unit line, every step of the split descending-row conjugation and removal, Schur Cauchy coefficients and raw/probability Haar cancellation. These are diagnostics of the proof, not numerical evaluation of an actual p-adic representation or of a global period. The complete 23,275-assertion replay passed; its final timestamp, whole-file review and frozen hashes are recorded separately.

\clearpage

\section{AM-DO-089 --- integral transverse compression removes the nonunipotent fixed block at support twenty-two}\label{am-do-089-integral-transverse-compression-removes-the-nonunipotent-fixed-block-at-support-twenty-two}

\textbf{Canonical checked integration --- 2026-11-22T20:16:28.424528+00:00.} Root and independent whole-proof/source/script reviews and separate504-assertion replays passed. The entire actual all-top H133 nonunipotent fixed-W0 sector at support22 has absolute contribution\textless2\^{}-776095. The proof uses integral transverse compression of h+h\^{}-1-2I without global nilpotence, simple-pole characteristic coefficients, full generalized-primary decomposition, and the quadratic isotropic-image or empty exact-plane alternatives. Every global rank, incidence type, original column and source fibre is included. This sector removal alone leaves the complete support floor22 and all signed/global/arithmetic/DP44 gates open.

\textbf{Algebraic-geometric translation.} The operator H and its compression are endomorphism pairings of the original paired integral Hermitian modules; two distinct source/target lattices are retained. Generalized-primary subspaces and the fixed signature-(1,1) plane are used over E or at its real realizations, with no moving integral frame. The scalar recovered by the integer Bezout relation lies in the original fractional ideal. Exact contextual anchors remain SGA3 I4.7.3/4.7.3.1; SGA6 VI Theorem3.3; FGA232 Theorem3.1 pp148--149; FGA236 Corollaries1.7/2.3 pp229/232; SGA7 I VII2.7/2.9 pp169--174 and VIII7.4(b) pp310--311; cote4 scan003 §§8.1.1--8.1.2, scans005--006 §§8.4.1--8.4.3 and cote54 scans005--007. No new scheme identification, lambda-ring/K-theory identity, l-adic sheaf/weight theorem, CM-motive comparison or motivic L-function assertion follows. DP\_MULT Q3e and THETA3 G1/A7 gain the stated containing-set removal; CX\_GRH CX-ALG-024 and GCT\_PLETH T-052 retain their separate arithmetic gates.

Author: \nolinkurl{/root/definite_oscillator}, 22 November 2026 UTC. Controlled pending proof; whole saved-file checks and independent review are recorded separately. No existing proof, registry or error array is changed by this file.

\subsection{1. Actual scope and conclusion}\label{actual-scope-and-conclusion}

Retain the original J56 all-top kernel, the charged cutoff H=133, every actual paired CM pencil, class mass, finite projector, central phase and selected real metric. Set \[
 f=a=168,\quad n=169,\quad W_0=E^5,\quad
 J_0=\operatorname{diag}(1,1,1,1,-1),\quad h=I_S\oplus h_0,
 \quad \dim S=164,\quad N=h_0-I. \tag{NC1}
\] The sector consists of \textbf{all nonunipotent} fixed-W0 matrices at support s=22, with nonzero finite coefficient. There is no hypothesis on the global rank of N or on its type-I/II/III profile. The exact central normalization has det h0=1, and every unmixed G reduction is unipotent, by DO060(IB12--15) and K079(ER1--2). The two mixed partial traces retain only their modulus bound one.

\textbf{Theorem.} The total absolute contribution of this entire sector is less than \[
                         2^{-776095}. \tag{NC2}
\] Only its intersection with the current remainder is removed. This adds one all-top error; it does not by itself prove a full support-floor update or the sufficient signed inequality.

The new step is to use the self-adjoint operator \[
                 H_h=h_0+h_0^{-1}-2I=-N^\dagger N=N^2h_0^{-1}, \tag{NC3}
\] where dagger is the original J0 adjoint. No global unipotence is used in NC3.

\subsection{2. The exact transverse module and integral pairings}\label{the-exact-transverse-module-and-integral-pairings}

Use the original LIFT41/K088 packet integers \[
\begin{split}
 Q&=1427908193729121522601,\quad U=77Q^2-14,\\
 u&=(2,95192390527,383092671043,12529843548438452388609),\\
 y_1&=(-976627,-153032,38026,0),\\
 y_2&=(-567678,695570,-172838,0),\\
 y_3&=(104583,-15309315176,-61610045224,2).
\end{split}\tag{NC4}
\] Their negative coordinates are zero. They satisfy u*u=U, gcd(u\_i)=1, and the three yj span the positive E-space T0=u-perp of dimension three. Write L=T0-perp=span\_E((u,0),e\_n), of signature(1,1). The exact squared norms of the yj are \[
        978665066829,\quad835948910828,\quad4030172803672366725045.\tag{NC5}
\]

At every actual T arrow h Mc=Md, all yj lie separately in \textbf{both} actual paired source and target modules. This is exactly DO078(FP20--22), from K088(IP5--6) and LIFT44(PR4--9). At the selected component Mc=(I+c pi\^{}-1 zq)Lambda with qyj=-pi m0yj; at the other component integral pairing with Lambda+pi\^{}-1 O z gives membership in the actual dual. This includes label zero and every incidence type. We use that established module inclusion, not a stable-lattice assertion for an arbitrary arrow.

Consequently, for every i,j, \[
 b_{ij}=\langle y_i,H_h y_j\rangle_J\in O_E. \tag{NC6}
\] Indeed the h term is an integral pairing inside Md, and the inverse term is an integral pairing inside Mc; subtract the integral identity term. Away T the original matrices are standard integral. This is a bilinear extension of the two separate pairings in FP21, and uses no unipotent square identity.

At an unmixed G component, let d be the rank of the nilpotent reduction of N. If d=0, NC3 gives valuation at least two for every entry of Hh. If d=1, the nilpotent reduction has square zero, and NC3 gives valuation at least one. The partner rank is identical by inverse adjoint. Thus every nonzero scalar obtained by an integral combination of Hh entries has paired valuation at least \[
                             2(2-d)_+. \tag{NC7}
\] For d\textgreater=2 the right side is zero. The rank-zero assertion is an assertion over the DVR, not merely over its residue field.

The full operator norm of Hh is at most 2H+2=268, since \textbar\textbar h0\^{}-1\textbar\textbar=\textbar\textbar h0\textbar\textbar. Hence \[
                 |\tau b_{ij}|\le268\|y_i\|\|y_j\|<2^{81} \tag{NC8}
\] at every label. The integer inequality using NC5 is checked exactly.

\subsection{3. A common rank-two Green estimate}\label{a-common-rank-two-green-estimate}

Put L2=335. The actual Green rank losses relative to rank two are335 at rank zero and167 at rank one. Equation NC7 pays these losses with at most335/4 times the paired scalar valuation, exactly as K081 section2. All global Euler corrections multiply to less than two, and both mixed factors remain contractions.

If an Hh scalar b has \textbar tau b\textbar\textless B and lies in a fractional ideal I supported only at the actual T primes, with N\_E I=T\^{}-2v, then its nonzero principal numerator has norm below B\^{}(2f)T\^{}(2v), and \[
 |a(h)|\le2 B^{335f/2}T^{335v/2}Q^{-335(f-1)}.\tag{NC9}
\] This is a uniform pointwise estimate. Selecting one nonzero scalar from a fixed finite family does not multiply a subsequent count of matrices: every matrix in the family obeys the same bound. No additional finite-rank or CM-orientation patterns are counted.

In particular, if any b\_ij in NC6 is nonzero, apply NC9 with B=2\^{}81 and v=0. This bounds that entire alternative, including every characteristic degree and every global rank.

\subsection{4. Characteristic degree at least three}\label{characteristic-degree-at-least-three}

Let d be the total algebraic multiplicity of eigenvalues of h0 different from one. The one-primary part is allowed to have Jordan blocks. Nonunipotence gives d\textgreater=1, and det h0=1 excludes d=1. The highest nonzero elementary coefficient of N is exactly e\_d(N); nilpotent Jordan blocks at eigenvalue one do not change this statement.

For d=k in\{3,4,5\}, choose b=e\_k(N). At a G component of reduction rank j, its valuation is at least (k-j)\_+, since it is a sum of k-minors. The paired valuation is twice that. The DVR-minor argument and Green comparison are precisely K079(ER7--10), with \[
 L_k=169k-k(k+1)/2,\qquad (L_3,L_4,L_5)=(501,666,830).\tag{NC10}
\] The singular values of N beyond its first are at most two, while its first is at most134. Therefore \[
 |\tau e_k(N)|\le {5\choose k}2^{k-1}134<B_k,
 \qquad (B_3,B_4,B_5)=(2^{13},2^{13},2^{12}).\tag{NC11}
\] This includes the sum over principal minors; a single-minor bound is not substituted for the coefficient bound.

K070(PT1--5), using the actual cyclic exterior trace, makes every characteristic coefficient integral outside type III and gives at most a simple pole at each of its two CM components. If there are t such places, the same paired Green argument gives \[
 |a(h)|\le2 B_k^{L_k f/k}T^{L_k t/k}Q^{-L_k(f-1)}.\tag{NC12}
\] This uses the coefficient's \textbf{simple} poles rather than the double entry/minor poles of the older K079 count. The characteristic degree k is not confused with the reduction rank or the global matrix rank.

\subsection{5. Quadratic degree and the vanishing compression}\label{quadratic-degree-and-the-vanishing-compression}

It remains to treat d=2 when all b\_ij vanish. The full nontrivial primary plane V and the one-primary complement W are defined over E and orthogonal nondegenerate. This is the usual generalized-eigenspace orthogonality for a unitary operator, as used explicitly in K083 sections2--3: the sets of eigenvalues one and different from one are separately stable under reciprocal conjugation. The polynomial on V is \[
               X^2-(2+y)X+1,\qquad y\in F^\times,\qquad
                       |\tau y|<136.\tag{NC13}
\] Here y=tr N, and y is real by reciprocal conjugation; y=0 would contradict the definition of V. The archimedean bound is the d=2 case of K083 section3. Repeated eigenvalue -1 and its possible Jordan block are included. On V, Cayley--Hamilton gives Hh=yI.

We first show that global rank N\textgreater2 is impossible in the present zero-compression alternative. In that event the unipotent restriction to the three-dimensional W is nontrivial. At each embedding it cannot act on a positive-definite space. Thus W has the negative direction and V is positive definite. The restriction to V is then unitary semisimple, and y\textless0. The accepted signature-one unipotent normal form DO060(IB21--22), restricted to W, says Hh\textbar W=N\_W\^{}2 and its Hermitian quadratic form is negative semidefinite: in its Witt coordinates it is -\textbar\textbar v\textbar\textbar\^{}2 times the squared absolute value of the last coordinate. This includes rank-one unipotents, when the square is zero. It is a statement about the Hermitian form, invariant under that real coordinate change; no integral basis is asserted.

Therefore the form \textless z,Hh z\textgreater{} is negative semidefinite on V perpendicular W. For each yj its value is zero by the vanishing compression, so its V component must be zero, because the V form is y times a positive-definite norm with y\textless0. Hence T0 is orthogonal to V. This puts the positive two-plane V inside L of signature(1,1), a contradiction. Thus \[
               \operatorname{rank}_E N=2,\qquad h|_W=I,
                  H_h=yP_V,\qquad H_h^2=yH_h.\tag{NC14}
\] The assertion h\textbar W=I follows because the two-dimensional nontrivial primary plane already accounts for the full rank. This argument allows all generalized multiplicities; it does not discard a possible one-primary Jordan block without proof.

The zero compression and self-adjointness say Hh(T0) is contained in L. For each j, \[
       \langle H_hy_j,H_hy_j\rangle
       =\langle y_j,H_h^2y_j\rangle
       =y\langle y_j,H_hy_j\rangle=0.\tag{NC15}
\] Thus every nonzero Hh yj is an isotropic vector in that fixed pencil plane.

\subsection{6. A small simple-pole scalar, or an impossible numerator}\label{a-small-simple-pole-scalar-or-an-impossible-numerator}

Suppose some Hh yj is nonzero. Write uniquely \[
                         H_hy_j=(u\,b,\beta).\tag{NC16}
\] Its isotropy makes b nonzero. The bound on Hh and NC5 give \[
                  |\tau b|\le268\|y_j\|/\sqrt U<2^{-27}.\tag{NC17}
\] The integer certificate checks the squared comparison, without introducing sqrt U into E.

At an actual arrow, h yj belongs to Md and h\^{}-1 yj to Mc. Every Mc is contained in pi\^{}-1 Lambda at the selected component, and the paired dual has the same simple-pole containment at the opposite component. Hence \[
 H_hy_j\in I_\Sigma^5,\qquad b\in I_\Sigma,\qquad
                      N_E I_\Sigma=T^{-2s}.\tag{NC18}
\] The second inclusion uses the original integer Bezout relation gcd(u\_i)=1, applied to its four positive components. It is not an assertion that the isotropic vector or its normalized negative coordinate is integral. The same integral linear combination retains the valuations in NC7 at G. Applying NC9 with B=2\^{}-27 and v=s bounds this entire alternative. The choice of one nonzero j, as in section3, introduces no extra matrix multiplicity.

Finally suppose Hh(T0)=0. Equation NC14 gives ker Hh=V-perp, of dimension three. Consequently V=L. The orthogonal projection to this exact original plane is diag(uu*/U,1), so \[
                    (H_h)_{11}=4y/U.\tag{NC19}
\] The integer U is odd and coprime to Q and T. Standard away-T integrality of Hh gives U divisibility of y at every prime dividing U, using NC19. At primes above two U is a unit, so no division by four loses a condition. At T, K070 supplies only simple poles, and only at type III, for y=tr N. Because y and U lie in F, relative ramification indices cancel. Thus \[
                 y\in U I_{\mathrm{III},F},\qquad
                 |N_Fy|\ge U^fT^{-t}\ge U^fT^{-22}.\tag{NC20}
\] The actual integer comparison \[
                            U^{168}>136^{168}T^{22}\tag{NC21}
\] contradicts NC13. This last alternative is empty. No norm-solvability or anisotropy assertion about U has been assumed.

\subsection{7. The full original negative-column count and every incidence}\label{the-full-original-negative-column-count-and-every-incidence}

For each fixed support and actual ordered tuple of source/target labels, K079(ER11--13) and the sharper all-top DO075 section3 count give \[
 \#h\le G_5 V_5(223)T^{10s},\quad G_5<2^{107},\quad
 V_5(223)=\frac{[((22/7)f223^2)]^{5f}}{(5f)!D_0^5}<2^{5757},
 \quad D_0=7^{140}29^{144}5^{84}13^{84}.\tag{NC22}
\] This is the \textbf{full five-complex-coordinate} target-overlattice count: the source contains pi e\_n, so h e\_n belongs to pi\^{}-1 Md. Its paired covolume is D0\^{}5 T\^{}-10s and it contains O\_E\^{}5 for every target, including zero. At H133 the whole negative column has norm below95 sqrt(f); the accepted original-ring cover enlarges the radius to223 sqrt(f). There is no real-part slice or zero-target coarea factor in NC22. The finite source stabilizer uses the same chosen source/target tuple and is bounded by the actual GL5 gcd. It is independent of global unipotence and characteristic degree.

All original selected real coefficients have modulus at most one. The number of supports is below2\^{}168. At a type-III place the at most (T-1)(T-2) ordered off-diagonal pairs cancel the exact reciprocal absolute rho. At other active places retain the full cost P\_T=2(T-1)\^{}2/(T-2)\textless2\^{}997. Integral places retain their original contraction bound one. For the integral-compression and small-scalar cases we conservatively use P\_T\^{}s for every active place, requiring no additional type count.

For characteristic degree k, summing all t with their positions gives \[
 \sum_{t=0}^s {s\choose t}T^{L_kt/k}P_T^{s-t}
 =(T^{L_k/k}+P_T)^s<2^sT^{L_ks/k}.\tag{NC23}
\] The strict comparison P\_T\textless T\^{}(Lk/k) follows from T\textgreater2\^{}995 and Lk/k\textgreater=166. Thus all profiles are already included. No extra rank-pattern, phase, class, determinant or witness-choice factor is dropped.

\subsection{8. Exact exponents and chronology}\label{exact-exponents-and-chronology}

Put LT=1019333/1024 and LQ=287843/4096, the checked strict upper/lower bounds on log2 T and log2 Q. Set C=1+168+107+5757=6033. At s=22 the complete exponents are \[
\begin{split}
 E_A&=C+81(335f/2)+(10LT+997)22-335\cdot167LQ
      =-5755634123/4096,\\
 E_3&=C+22+13(501f/3)+(10+501/3)22LT-501\cdot167LQ
      =-6687102305/4096,\\
 E_4&=C+22+13(666f/4)+(10+666/4)22LT-666\cdot167LQ
      =-7333969827/2048,\\
 E_5&=C+22+12(830f/5)+(10+830/5)22LT-830\cdot167LQ
      =-11357468235/2048,\\
 E_C&=C-27(335f/2)+(10+335/2)22LT+997\cdot22-335\cdot167LQ
      =-3178901223/4096.
\end{split}\tag{NC24}
\] All five are strictly below -776098. They cover respectively the nonzero integral compression, degrees three/four/five, and the remaining quadratic small-scalar alternative. They may overlap if counted as containing sets, which only enlarges the absolute upper bound. The remaining quadratic plane case is empty. Consequently the whole sector is below5 times2\^{}-776098, hence below2\^{}-776095, as claimed.

The canonical predecessor is R066's67 all-top and31 uniform entries, with its160208 appended after DO088358. The final companion verifies that exact complete predecessor and appends776095 once, retaining all earlier errors. Both full dyadic sums remain below2\^{}-230. Until a separate exhaustive structural assembly, the all-top floor stays22 and the independent uniform floor stays11.

This proof concerns the actual original coefficient sum and unchanged lattice correspondences. It changes no finite projector, chosen real type, CM field, prime, packet or arithmetic image. The signed sufficient inequality, prescribed global Bessel/Mram nonvanishing, actual family constructor, selected arithmetic content, Read/032 value, DP44 cost and RH implication remain open. The carried integral-module and determinantal language is that of SGA3 I4.7.3/4.7.3.1 and SGA6 VI3.3; no new scheme, motive, sheaf or comparison statement is asserted. Accepted local Green and geometry-of-numbers interfaces are used at the exact locators above; no stipulated Lean proof or archive audit is performed.

The companion arithmetic and rational unitary models diagnose the constants, pair exponents, primary/quadratic identities and whole error sum. They do not automate the primary decomposition or the original-lattice proof; these require the separate whole saved-proof reviews.

\clearpage

\section{AM-R-067 --- complete all-top support twenty-two exhaustion}\label{am-r-067-complete-all-top-support-twenty-two-exhaustion}

\textbf{Canonical checked integration --- 2026-11-22T20:17:41.441676+00:00.} Root and independent whole-proof/script/source-scope reviews and separate83-assertion replays passed. The pending description below is historical. Every actual all-top support22 term is exhausted by both original F/IC8 classes at every rank1..33 and all identity-compression alternatives. DO089 supplies the whole nonunipotent fixed-W0 sector, DO088 the whole globally rank-one sector, R066 the whole IC8 rank1/2 sector, and every earlier complete range is retained. The actual complete all-top support floor is23. All68/31 chronological error entries are unchanged and both sums remain\textless2\^{}-230; no new charge is added. The uniform floor remains11. J23.2 and every global/arithmetic/source/DP44 gate remain open.

Author: \nolinkurl{/root}, GPT-6 Astra Pro, inherited effort not externally exposed, ALP\_MACHINE.22 November2026 UTC. Controlled pending assembly. Promotion requires the final whole-proof reviews of DO089 and this proof; exact check records are separate.

\subsection{1. Statement and unchanged energy}\label{statement-and-unchanged-energy}

Keep the actual J56 all-top K044 endpoint, its original packet and selected Hilbert metric, all integral Hermitian modules, paired CM incidences, finite projectors, class masses and phases. Only the already charged H133 cutoff is used. The parameters remain f=a=168,n=169,dim S=164. R065 provides the exact remaining energy with support at least22. DO088, R066 and DO089 are containing-set removals with respective charges358,160208 and776095, all retained in their chronological order.

There is an actual remaining subset R\_top\^{}(23) such that

\[
 \mathcal E_{\rm top}=1+\operatorname{Re}
       \sum_{h\in R_{\rm top}^{(23)}}w(h)+\epsilon,
 \qquad |\epsilon|<2^{-230},\qquad s(h)\ge23.
 \tag{J23.1}
\]

The present assembly adds no error charge. The sufficient inequality

\[
 \operatorname{Re}\sum_{h\in R_{\rm top}^{(23)}}w(h)>-1+2^{-230}
 \tag{J23.2}
\]

remains open. The independent uniform H8192 support floor stays11. We prove the claimed support by exhausting the original structural partition at support22.

\subsection{2. Nonidentity positive compression}\label{nonidentity-positive-compression}

Write A=h\_SS in the unchanged original decomposition S perpendicular W0. The existing residual satisfies r=rank\_E(A−I)\textless=33. For A!=I, K069(IC1--8) partitions the terms into F, where h fixes K\_A=ker(A−I) intersect ker(A*−I) pointwise, and the exceptional isotropic-cross complement IC8. These classes do not require global unipotence.

\begin{longtable}[]{@{}
  >{\raggedright\arraybackslash}p{(\columnwidth - 2\tabcolsep) * \real{0.5000}}
  >{\raggedright\arraybackslash}p{(\columnwidth - 2\tabcolsep) * \real{0.5000}}@{}}
\toprule\noalign{}
\begin{minipage}[b]{\linewidth}\raggedright
Actual class at support22
\end{minipage} & \begin{minipage}[b]{\linewidth}\raggedright
Complete containing-set coverage
\end{minipage} \\
\midrule\noalign{}
\endhead
\bottomrule\noalign{}
\endlastfoot
F,r=1 & DO086(FR1--10), all supports18..94, with its8622 charge. \\
F,2\textless=r\textless=33 & K069(IC20--21), whose minimum full ceiling is22, attained at r=2. \\
IC8,r=1 or2 & R066(JB1--12), the full zero/transverse/joint-plane trichotomy at22, with its160208 charge. \\
IC8,3\textless=r\textless=31 & DO064's previous ranges and DO085(MR1--24)'s contiguous extensions, whose minimum complete ceiling is23. \\
IC8,r=32 or33 & DO064(PM22--23), complete through168. \\
\end{longtable}

Every positive compression rank and both classes are present. The K069 ceiling22 is inclusive: its definition is the largest integer satisfying the stated sufficient bound, and IC21 includes1\textless=s\textless=s\_r. No assertion about its failed next endpoint is needed. Thus all A!=I terms at22 are covered.

\subsection{3. Identity positive compression}\label{identity-positive-compression}

For A=I the original unitary identities R059(GB1--3) give either a fixed-S block h=I\_S direct-sum h0 on W0, or a nonzero scalar-cross block. This classification is not restricted to unipotent matrices.

\begin{longtable}[]{@{}
  >{\raggedright\arraybackslash}p{(\columnwidth - 2\tabcolsep) * \real{0.5000}}
  >{\raggedright\arraybackslash}p{(\columnwidth - 2\tabcolsep) * \real{0.5000}}@{}}
\toprule\noalign{}
\begin{minipage}[b]{\linewidth}\raggedright
Actual class at support22
\end{minipage} & \begin{minipage}[b]{\linewidth}\raggedright
Complete containing-set coverage
\end{minipage} \\
\midrule\noalign{}
\endhead
\bottomrule\noalign{}
\endlastfoot
Fixed S,nonunipotent,any global rank & DO089(NC1--24), with one776095 charge. Its integral compression, characteristic degrees3/4/5, quadratic small-scalar and empty exact-plane alternatives include every type profile. \\
Fixed S,globally unipotent,rank(h−I)=1 & DO088(WM1--18), all marked sizes and types at22, with one358 charge. \\
Fixed S,globally unipotent,rank(h−I)=2 & R058(HB1--25), every profile at supports13..23, with retained39409 charge. \\
Nonzero scalar cross,with or without global unipotence & R061(GS1--23), the entire class at22, with retained78919 charge. \\
Full rank0 & The identity has support0 and does not occur at22. \\
\end{longtable}

The signature-one unipotent classification DO060(IB21--22) allows no other nonzero unipotent rank. DO089 is a whole nonunipotent-sector theorem, so no separate restriction on type-III count, semisimple degree, or finite reduction rank remains in its row. R061 covers the entire scalar-cross class, rather than just its globally unipotent subset. Hence the table exhausts A=I at22.

Sections2 and3 exhaust support22. All smaller supports were already removed by R065, which proves the support statement in J23.1.

\subsection{4. Exact chronological errors and boundary}\label{exact-chronological-errors-and-boundary}

The chronological predecessor after DO089 contains68 all-top and31 uniform entries. Its final nine all-top exponents are

\[
 4255,177258,7548,8622,52215,1649,358,160208,776095.
 \tag{J23.3}
\]

R064 and R065 added no charges, and neither does R067. There is exactly one all-top exponent231; every other exponent is at least358. Consequently

\[
 \sum_{e\in\mathcal L_{\rm top}}2^{-e}
 \le2^{-231}+67\,2^{-358}<2^{-230}.
 \tag{J23.4}
\]

The original31-term uniform sum remains below2\^{}-230. The companion checks the complete integer dyadic sums, exact ordered predecessor and every rank in the tables. Each containing-set removal is applied only to its intersection with the current remainder, making the actually removed sets disjoint; the conservative bounds remain charged once each. The original energy identity and triangle inequality now give J23.1.

This strictly reduces the Q1/Q2 finite/real sufficient gate to J23.2. It does not establish that inequality, a prescribed global Bessel/Mram value, a fully admissible mixed-CM source, selected arithmetic-image content, the actual Read/comparison ties,032's pairing/value, or any changed-family DP44 term. It proves neither a shorter valid chain nor a universal strict negative DP44 difference. Q3 has no Kummer-Hecke RH conclusion, and Q4 retains its global arithmetic charge-comparison gate.

\subsection{5. Algebraic-geometric translation and verification scope}\label{algebraic-geometric-translation-and-verification-scope}

The assembly partitions rank conditions and fixed/isotropic submodules in the original integral Hermitian modules and labelled correspondences. Every lattice, polarization, finite level, coefficient module, Picard datum and biextension remains the original one. No rational-coordinate count gives an inverse on the arithmetic theta image.

The exact carried contextual anchors are SGA3 I4.7.3/4.7.3.1 and SGA6 VI Theorem3.3; FGA232 Theorem3.1 pp148--149; FGA236 Corollaries1.7/2.3 pp229/232; SGA7 I VII2.7/2.9 pp169--174 and VIII7.4(b) pp310--311. Formal crystals and Cartier duality retain cote4 scan003 §§8.1.1--8.1.2, scans005--006 §§8.4.1--8.4.3 and cote54 scans005--007. These locators provide context, not a substitute proof of the exhaustion. No new scheme identification, lambda-ring/K-theory equality, l-adic sheaf/weight result, CM-motive comparison or motivic L-function theorem follows. DP\_MULT Q3e and THETA3 G1/A7 gain exactly J23.1; CX\_GRH CX-ALG-024 and GCT\_PLETH T-052 retain their distinct gates.

Every quantitative input retains its own whole-proof checks and source restrictions. This new proof concerns complete partition coverage and exact error accounting. The certificate is a finite diagnostic; it neither enumerates all rational matrices nor evaluates a global period. No Lean build or archive audit is performed.

\clearpage

\section{AM-DO-090 --- complete nonunipotent fixed blocks at supports twenty-three through twenty-six}\label{am-do-090-complete-nonunipotent-fixed-blocks-at-supports-twenty-three-through-twenty-six}

\textbf{Canonical checked integration --- 2026-11-22T20:30:03.911778+00:00.} Root and independent whole-proof/script/interface checks and separate50-assertion replays passed. The entire actual all-top H133 nonunipotent fixed-W0 sector at supports23..26, with all global ranks and incidence profiles, has absolute contribution\textless2\^{}-65341. The original full-G coefficient ideal makes the exact quadratic plane empty through26, and all five complete DO089 bounds retain their original columns, paired modules, fibres and finite factors. The complete all-top floor remains23 and the uniform floor11; other structural classes and all signed/global/arithmetic/DP44 gates remain open.

\textbf{Algebraic-geometric translation.} The real-field trace is a section of the original fractional ideal U J\_G,F I\_III,F. Its F-ideal norm, Q\^{}167, is not doubled by the two CM components. The original integral Hermitian modules and their determinantal characteristic coefficients are retained. Contextual exact anchors are SGA3 I4.7.3/4.7.3.1; SGA6 VI Theorem3.3; FGA232 Theorem3.1 pp148--149; FGA236 Corollaries1.7/2.3 pp229/232; SGA7 I VII2.7/2.9 pp169--174 and VIII7.4(b) pp310--311; cote4 scan003 §§8.1.1--8.1.2, scans005--006 §§8.4.1--8.4.3 and cote54 scans005--007. These are the carried module, Picard and biextension frameworks, not newly constructed comparison maps. No new lambda-ring/K-theory identity, l-adic sheaf or weight theorem, CM motive or motivic L-function assertion follows. DP\_MULT Q3e and THETA3 G1/A7 gain the stated containing-set removal; CX\_GRH CX-ALG-024 and GCT\_PLETH T-052 retain their arithmetic gates.

Author: \nolinkurl{/root/definite_oscillator}, 22 November 2026 UTC. Separate controlled continuation of frozen DO089; no earlier proof is modified. Whole proof, script and independent checks are recorded separately before promotion.

\subsection{1. Exact new sectors}\label{exact-new-sectors}

Keep every original all-top J56 datum of DO089(NC1--5): H=133, f=a=168, n=169, the original packet, finite projectors, paired source/target pencils, class masses, phases and selected metric. Consider all nonunipotent matrices h=I\_S direct-sum h0 on W0 of signature(4,1), at \[
                           23\le s\le26.\tag{NE1}
\] There is no global-rank, characteristic-degree or incidence-type restriction. All334 unmixed G reductions are unipotent, as in the unchanged containing residual.

\textbf{Theorem.} The total absolute contribution of NE1 is less than \[
                              2^{-65341}.\tag{NE2}
\] This is one fresh containing-set removal on its intersection with the current remainder. It does not remove the other structural classes or establish a larger complete support floor.

\subsection{2. The pure quadratic plane acquires the full G ideal}\label{the-pure-quadratic-plane-acquires-the-full-g-ideal}

All geometric alternatives and uniform bounds in DO089(NC3--19) are independent of support22. In particular the integral transverse compression, three possible characteristic degrees k=3,4,5, and quadratic simple-pole Bezout scalar retain exactly their proofs and constants. Their only support dependence is in the displayed powers of T and in the complete incidence counts.

The one place where DO089 used support22 without a counting exponent was its empty quadratic plane. There the exact conclusions are \[
 y=\operatorname{tr}(h_0-I)\in F^\times,\quad
 |\tau y|<136,\quad
 H_h=y\operatorname{diag}(uu^*/U,1),\quad
                 y\in U I_{\mathrm{III},F}.\tag{NE3}
\] These are NC13,19--20, with U=77Q²−14 odd and prime to Q,T. The actual third-type count satisfies t\textless=s.

At every unmixed relative G place the reduction of h is unipotent. Therefore the characteristic coefficient tr(h−I)=y reduces to zero at either CM component. Since y lies in F and these places are split and unramified, its F-prime valuation is at least one. Let J\_G,F be the product of these334 F primes. Their residue-size product is Q\^{}(f−1), so \[
              N_F J_{G,F}=Q^{167},\qquad
                   y\in U J_{G,F}I_{\mathrm{III},F}.\tag{NE4}
\] This is not a new finite-rank assumption, a rank-pattern count or a Green estimate. It is the direct coefficient divisibility shared by every finite unipotent reduction. The U and G ideals are coprime, and the T denominator is supported elsewhere, so their factors multiply. At dyadic primes U is a unit, as in NC20.

Consequently, for every s\textless=26, \[
 |N_F y|\ge U^{168}Q^{167}T^{-t}
          \ge U^{168}Q^{167}T^{-26}>136^{168}.\tag{NE5}
\] The last strict integer inequality is checked against the actual unchanged primes. A conservative logarithmic margin is37581933/4096 bits, using log2 U\textgreater601356/4096, log2 Q\textgreater287843/4096, log2 T\textless1019333/1024, and log2 136\textless8. This contradicts NE3. Thus the pure-plane alternative is empty throughout NE1. No anisotropy, new field or ideal generator has been assumed.

\subsection{3. Every old factor at the new supports}\label{every-old-factor-at-the-new-supports}

Use the same exact full negative-column target-overlattice and same-pair source fibre NC22: G5\textless2\^{}107 and V5(223)\textless2\^{}5757, followed by T\^{}(10s). This is the full complex count for every target, including zero, not a real-part coarea slice. The real coefficient remains bounded by one. All support choices remain below2\^{}168.

For the integral-compression and small-scalar alternatives use the complete pair bound P\_T\^{}s\textless2\^{}(997s). For degree k the complete type/position sum is the same NC23 binomial (T\^{}(Lk/k)+P\_T)\^{}s\textless2\^{}s T\^{}(Lk s/k). No support, type, pair, fibre, mixed projector, witness-choice or finite-rank factor is changed or dropped.

Set C=6033, LT=1019333/1024, LQ=287843/4096, L3=501,L4=666,L5=830, and b3=b4=13,b5=12. The five complete containing-set exponents, now with variable s, are \[
\begin{split}
 E_A(s)&=C+81(335f/2)+(10LT+997)s-335\cdot167LQ,\\
 E_k(s)&=C+s+b_kL_kf/k+(10+L_k/k)sLT-L_k\cdot167LQ
                \quad(k=3,4,5),\\
 E_C(s)&=C-27(335f/2)+(10+335/2)sLT+997s-335\cdot167LQ.
\end{split}\tag{NE6}
\] These are precisely NC24 before substituting22, not new formulas reconstructed from numerical endpoint tables. All slopes are positive. The certificate checks every one of the twenty case/support values and gives the maximal exponent at each support:

\begin{longtable}[]{@{}rr@{}}
\toprule\noalign{}
Support s & Maximum among the five complete cases \\
\midrule\noalign{}
\endhead
\bottomrule\noalign{}
\endlastfoot
23 & −2451091081/4096 \\
24 & −1723280939/4096 \\
25 & −995470797/4096 \\
26 & −267660655/4096 \\
\end{longtable}

All twenty exponents are strictly below -65346. Hence their absolute union is less than20 times2\^{}-65346, which is less than2\^{}-65341. Together with the empty exact-plane case this proves NE2. The five cases can overlap as containing sets; the absolute upper bound remains valid. The corresponding maximum at support27 is460149487/4096\textgreater0, which is only failure of this sufficient estimate, not a positive contribution or nonvanishing claim.

\subsection{4. Chronology and limits}\label{chronology-and-limits}

The chronological predecessor is the unchanged68-term all-top array after DO089, also used without alteration by the complete-floor23 assembly R067. Append65341 once, retain every earlier charge, and leave the31-term uniform array unchanged. The companion checks the full ordered69/31 sums below2\^{}-230. The all-top floor remains23 until an independent complete structural exhaustion at a larger support; the uniform floor remains11.

The new input is NE4's full-G coefficient ideal, used to extend the empty-plane alternative. The other geometric/counting interfaces are the full checked statements DO089(NC3--24) at arbitrary support, with their explicit support dependence retained above. Original integral modules, labelled correspondences and actual chosen coefficient/oscillator data are unchanged. No source constructor, arithmetic theta-image inverse, selected integral-unit vector, prescribed Bessel/Mram value, signed inequality, actual DP44 comparison or RH conclusion follows. Existing SGA3 I4.7.3/4.7.3.1 and SGA6 VI3.3 module/determinantal language remains contextual; no new algebraic-geometric comparison theorem is asserted.

A complete proof check is separate from the exact numerical replay. The latter checks the new ideal norm, every case/support value and the complete error array; it does not automate the original paired-lattice or unitary primary-space arguments.

\clearpage

\section{AM-R-068 --- the actual all-top cutoff enlarges every higher-rank fixed-class range}\label{am-r-068-the-actual-all-top-cutoff-enlarges-every-higher-rank-fixed-class-range}

\textbf{Canonical checked integration --- 2026-11-22T20:31:58.930642+00:00.} Root and independent whole-proof/script/interface reviews and separate150-assertion replays passed. Historical pending text is preserved below. At the unchanged actual all-top H133 cutoff, the entire original fixed class F is covered through support110 for rank(A-I)=2,154 for rank3,and all168 places for every rank4..33. Exactly1627 fresh rank/support cells beyond K069 receive one4085 charge. This uses the original integral block chart, r+6 original column coordinates, full paired target-overlattice and source fibre, and complete Green numerator; it assumes neither global unipotence nor the IC8 partial-monomial classification. The resulting remainder has F rank2 only at support\textgreater=111,F rank3 only at support\textgreater=155,and no F rank4..33. The complete all-top floor remains23 and uniform floor11; all signed/global/arithmetic/DP44 gates remain open.

Author: \nolinkurl{/root}, GPT-6 Astra Pro, inherited effort not externally exposed, ALP\_MACHINE.22 November2026 UTC. Controlled pending proof; whole checks and exact replay are recorded separately before promotion.

\subsection{1. Full original class and conclusion}\label{full-original-class-and-conclusion}

Keep the original J56 all-top endpoint, its already charged H=133 cutoff, f=a=168,n=169,m=164, the actual T and G primes, all labelled paired CM pencils, finite projectors, class masses, phases and selected Hilbert metric. For the original positive compression A=h\_SS put

\[
 r=\operatorname{rank}_E(A-I),\qquad
 K_A=\ker(A-I)\cap\ker(A^*-I),\qquad
 \mathscr F=\{h:h|_{K_A}=I\}.
 \tag{FA1}
\]

This is exactly K069(IC1--8)'s fixed class. It has no global-unipotence or incidence-type restriction. We retain the existing residual bound r\textless=33. The exceptional IC8 class is separate; its partial-monomial classification is not an assumption on FA1.

\textbf{Theorem.} Every actual FA1 class with2\textless=r\textless=33 is covered through the following complete inclusive support ceilings:

\[
 s_2=110,\qquad s_3=154,\qquad s_r=168\quad(4\le r\le33).
 \tag{FA2}
\]

The fresh cells beyond the previously charged K069(IC20--21) ranges have total absolute contribution less than

\[
                            2^{-4085}.\tag{FA3}
\]

Only their intersection with the current remainder is removed. The result adds one all-top charge and does not alone raise the complete support floor23. The independent uniform bounds are unchanged.

\subsection{2. Apply the existing count at the actual cutoff}\label{apply-the-existing-count-at-the-actual-cutoff}

K069(IC9--17) proves its complete bound using the cutoff parameter H in each displayed count. Its integral entry packing, original-coordinate encoding, fixed-source fibre and paired Green comparison apply to the present subset with H=133. No basis of K\_A is made integral. In particular, for the actual stratum(r,s), its original absolute sum is at most

\[
\begin{split}
 B_H(r,s)={}&2\binom{m}{r}^{2}\binom{f}{s}G_n
 (2H+3)^{2f(2mr-r^2)}(2H+1)^{2f(r+6)}\\
 &\quad\cdot(2^rH)^{af}\,
 T^{2(r+6)s}P_T^s\,Q^{-L_r(f-1)},\\
 &L_r=169r-r(r+1)/2,\qquad
 P_T=\frac{2(T-1)^2}{T-2}.
\end{split}\tag{FA4}
\]

Here Q=g0g1 with g0=22285932761,g1=64072175441, and G\_n is the unchanged two-prime finite source-stabilizer bound, less than2\^{}5022. The single leading2 is the complete Green Euler correction. Both mixed finite projectors and the normalized real coefficient retain their modulus bound one.

For clarity, the powers in FA4 retain all original encodings. The rank-r integral positive block is encoded by2mr−r² original entries and a chosen row/column minor. Its count is the first packing power and the squared binomial. Because h fixes K\_A and codim K\_A\textless=r+1, its negative column is determined by r+1 original S coordinates and five original W0 coordinates. The actual target over-lattice has T\^{}(2(r+6)s) cosets at both CM partners, including zero labels. Each coset uses the original-ring packing count with2H+1. A rational inverse minor is used only for uniqueness, never as a map on integral lattices. The finite fibre is the same source-lattice isometry group. The full ordered-pair multiplicity and its actual weight are P\_T\^{}s. Thus no type count, selected dimension, zero-target correction, or moving-lattice determinant has been omitted.

The paired Green numerator(2\^{}rH)\^{}(af) is exactly the integral-minor bound K069(IC16), itself K068(LR7--12) at p=1,k=r. It is retained in full; no exact-rank assertion modulo G is made for the general F class. The fact that A is integral in the original coordinates is also retained.

\subsection{3. Directed rational exponent and every rank}\label{directed-rational-exponent-and-every-rank}

Use the exact strict bounds

\[
\begin{split}
 \log_2(269)&<\ell_X=8+74/1024,\qquad
 \log_2(267)<\ell_V=8+63/1024,\\
 \log_2(133)&<\ell_H=28900/4096,\qquad
 \log_2T<\ell_T=1019333/1024,\\
 \log_2Q&>\ell_Q=287843/4096,\qquad
 P_T<2^{997}.
\end{split}\tag{FA5}
\]

The first two follow by the exact integer comparisons269\^{}1024\textless2\^{}8266 and267\^{}1024\textless2\^{}8255. The companion checks all six against the original integers. Use binom(m,r)\textless=2\^{}164 and binom(f,s)\textless=2\^{}168. Taking logarithms of FA4 gives

\[
 \log_2 B_{133}(r,s)<C_r+b_rs,
 \tag{FA6}
\]

where

\[
\begin{split}
 C_r={}&168+328+5022
  +336(328r-r^2)\ell_X+336(r+6)\ell_V+1\\
 &\qquad +(r+\ell_H)168^2-L_r167\ell_Q,\\
 b_r={}&2(r+6)\ell_T+997>0,\\
 s_r={}&\min\left(168,
       \left\lfloor\frac{-4096-C_r}{b_r}\right\rfloor\right).
\end{split}\tag{FA7}
\]

Exact rational evaluation at every integer2\textless=r\textless=33 gives FA2. Hence every such full cell with1\textless=s\textless=s\_r has absolute contribution less than2\^{}-4096. The strictness follows from FA5 even if its rational upper endpoint equals−4096. The ceiling168 is the total number of relative T places, not an extrapolation beyond the actual support range. For r2 and r3 the next support fails this particular sufficient exponent test; that failure gives no lower bound on the true sum.

\subsection{4. Fresh cells and ordered error accounting}\label{fresh-cells-and-ordered-error-accounting}

The old K069(IC20--21) ceilings, in the exact rank order2..33, are

\[
\begin{gathered}
22,40,55,67,76,85,92,98,104,108,112,116,119,122,125,127,\\
129,131,132,134,135,136,138,139,140,140,141,142,142,143,143,144.
\end{gathered}\tag{FA8}
\]

Their sum is3677. The new ceilings sum to110+154+30·168=5304, and each new ceiling is at least its old one. The exact fresh set consists of each(r,s) with old\_s\_r\textless s\textless=new\_s\_r; it therefore has1627 cells. Their absolute union is less than1627·2\^{}-4096\textless2\^{}-4085, proving FA3. Previously removed cells are not counted as fresh. Later containing-set removals may have removed some of these cells as well, so the actual removal is still intersected with the current remainder.

In the intended canonical chronology DO090 precedes this theorem and leaves69 all-top and31 uniform entries. Append4085 once after its65341 entry, leaving70/31 entries. Every earlier entry, including358 and231, remains exactly once. There is one231 exponent and all other69 all-top exponents are at least358, so

\[
 \sum_{e\in\mathcal L_{\rm top}}2^{-e}
 \le2^{-231}+69\,2^{-358}<2^{-230}.
 \tag{FA9}
\]

The companion also checks the complete exact integer dyadic sums. This theorem alone leaves the complete all-top floor23 and uniform floor11. The new remainder additionally satisfies r2 in F implies s\textgreater=111, r3 in F implies s\textgreater=155, and no F rank4..33 remains. F rank1 retains DO086's existing range through94. IC8 and every identity-compression class retain their separate existing reductions.

\subsection{5. Geometry, dependencies and unresolved arithmetic}\label{geometry-dependencies-and-unresolved-arithmetic}

The argument specializes a proved count on the original integral positive block and its fixed submodule. It keeps the original labelled correspondences, finite fields, Hilbert metric, Picard/Hom data and coefficient realization. Both stored positive-block entries and selected column coordinates belong to the original modules. Their injective rational encoding is not an integral arithmetic-image inverse.

Exact carried contextual anchors are SGA3 I4.7.3/4.7.3.1 and SGA6 VI Theorem3.3; FGA232 Theorem3.1 pp148--149; FGA236 Corollaries1.7/2.3 pp229/232; SGA7 I VII2.7/2.9 pp169--174 and VIII7.4(b) pp310--311; cote4 scan003 §§8.1.1--8.1.2, scans005--006 §§8.4.1--8.4.3 and cote54 scans005--007. No new scheme identification, lambda-ring/K-theory equality, l-adic sheaf/weight theorem, CM-motive comparison or motivic L-function statement follows. DP\_MULT Q3e and THETA3 G1/A7 gain the exact F-sector restrictions above. CX\_GRH CX-ALG-024 and GCT\_PLETH T-052 retain their separate arithmetic comparison gates.

The signed sufficient inequality, prescribed Bessel/Mram period, fully admissible mixed-CM source, selected-image content, actual Read/STiesR and032 value, and DP44 cost comparison remain open. The companion checks exact directed bounds, all old/new cells and complete error sums. The whole proof check separately verifies that K069's parameter-dependent original count has been specialized without changing its geometric or arithmetic hypotheses. No Lean build or archive audit is performed.

\clearpage

\section{AM-LIFT-074 --- prescribed-vector transfer to the literal global Poincaré detector}\label{am-lift-074-prescribed-vector-transfer-to-the-literal-global-poincaruxe9-detector}

\textbf{Canonical checked integration --- 2026-11-22T20:33:14.323738+00:00.} Root and independent whole-proof/source/script reviews and separate47,807-assertion replays passed. Conditional on a nonzero global Bessel functional on the two specified embedded automorphic occurrences and the genuine source/local-vector premises, the prescribed initial Mram column and admitted bottom row have \textless F,Pu(T)\textgreater=c\_Per\emph{conjugate(b)!=0, where c\_Per=c1}mu\_f(K2)\textgreater0 is the actual source Haar-disintegration scalar. The finite tests survive and the complete real Wu contraction equals conjugate(b), even if W(1)=0. The normalized diagonal Gram is at least c\_Per\^{}2\emph{\textbar b\textbar\^{}2/(c0}\textbar\textbar F\textbar\textbar\^{}2)\textgreater0. Global distinction, actual arithmetic source/Read/STiesR/032 transport, the unevaluated coupling b and all DP44 costs remain open.

\textbf{Algebraic-geometric translation.} This theorem is an analytic contraction of the same original Hermitian/Picard/Hom data and its literal source correspondence; it does not identify that correspondence with an integral theta image or CM-motive morphism. The coefficient realization and selected metric are unchanged. Exact contextual anchors remain SGA3 I4.7.3/4.7.3.1; SGA6 VI Theorem3.3; FGA232 Theorem3.1 pp148--149; FGA236 Corollaries1.7/2.3 pp229/232; SGA7 I VII2.7/2.9 pp169--174 and VIII7.4(b) pp310--311; cote4 scan003 §§8.1.1--8.1.2, scans005--006 §§8.4.1--8.4.3 and cote54 scans005--007. No new scheme comparison, lambda-ring/K-theory equality, l-adic sheaf/weight theorem, motive or motivic L-function assertion follows. DP\_MULT Q3e and THETA3 G1/A7 gain this conditional detector transfer, while the literal arithmetic pin and value remain. CX\_GRH CX-ALG-024 and GCT\_PLETH T-052 retain their distinct arithmetic comparison gates.

Author: \nolinkurl{/root/lift_accounting}, 22 November 2026 UTC.

\textbf{Status: author whole-proof, source-interface, and complete-script review PASS; 47,807 exact diagnostic assertions PASS.} The separate author-review and manifest files record the final fingerprints and verification scope; independent review is separate. This is a conditional transfer theorem. Its global premise is a nonzero Bessel functional on the specified pair of actual automorphic embeddings. It does not prove that premise, construct a complete source family, or identify the requested arithmetic \nolinkurl{Read}/\nolinkurl{STiesR}/specific032 instance. No Lean build is claimed.

\subsection{1. Statement and exact hypotheses}\label{statement-and-exact-hypotheses}

Retain a genuine Hida setting \nolinkurl{X}, a word \nolinkurl{u}, its actual compact factor H, flag unipotent N, marked line of norm −T, and the Bessel subgroup S=N⋊H. Retain \nolinkurl{X.xu u}, the actual finite subgroup K₂, character χ₂, real quotient measure ν\_Q, and all global Haar measures. The auxiliary Bessel datum B is an actual \nolinkurl{SBes} with its literal \nolinkurl{tieN}, \nolinkurl{tieH}, and \nolinkurl{tieψ}; a subgroup having only the same real signature is insufficient. Write W for the actual \nolinkurl{Wu u} of LIFT061 and σ for its real discrete-series representation. Thus

\[
 W(sy)=\psi_X(s)\rho(s)W(y),\qquad
 \int_{S_\infty\backslash G_\infty}\|W(y)\|^2\,d\nu_Q(y)=1.
 \tag{PD1}
\]

The chosen unit vector w is the literal \nolinkurl{IsLowUnit} vector in σ. It is not replaced by a highest-weight vector. As in LIFT061-U, W is the value at w of a nonzero bounded G∞-intertwiner into the relative Hilbert space, and the relevant smooth Bessel Hom is one-dimensional.

Fix two \textbf{embedded} irreducible unitary cuspidal automorphic representations: the initial representation π of H(A), with real factor ρ, and a bottom row representation Π of G(A), with real factor σ. The source Bessel form uses the complex conjugate of this row representation in its second slot. Fix the actual embeddings into automorphic functions, including their multiplicity vectors. A nonzero functional on another occurrence of an isomorphic representation does not satisfy the premise below.

Assume the following actual vector hypotheses.

\begin{enumerate}
\def\labelenumi{\arabic{enumi}.}
\tightlist
\item
  A nonzero pure finite vector a\_f of π gives the prescribed column T in \nolinkurl{Cu X u hK}. Through the actual column dictionary it is in the literal initial Mram space. At each selected T place it is the original strict covector-level nontrivial dilation line, not a projectively averaged substitute.
\item
  Π has the admitted finite types of LIFT072 at G, is spherical at T, and has the unchanged spherical local data of LIFT073 away G∪T. The smaller representation and a\_f satisfy the corresponding actual local hypotheses. In particular the local genericity and temperedness assumptions are supplied, rather than inferred from a finite generic type alone. LIFT069 and LIFT070 supply these on their stated chosen-bottom/full-SizesAt scopes; the earlier concrete m=2 scope can instead use its checked temperedness theorem.
\item
  The finite translation and root-coordinate dictionary is exactly LIFT060/067/068/072/073, in the simultaneous dual/conjugate convention required by PD4 and PD8. The row representation Π and the conjugated form in the Bessel slot are distinguished explicitly; the local results permit precisely this transported orientation. The T test is the physical Dτ test in these coordinates. At G the larger vector may be chosen inside the already admitted finite Π₀-subspace, as LIFT072 permits. This choice does not change the global representation, its goodTuple projector, or K₂ and χ₂. Away G the larger vector is the prescribed spherical vector.
\item
  The global Bessel functional on these two embedded representations, for this physical S and this character, is not identically zero.
\item
  The actual source \nolinkurl{Full X}, \nolinkurl{InHH X u}, and the stated Haar/source identifications are supplied. These are the hypotheses under which its literal \nolinkurl{Pu} and \nolinkurl{PerLaw} are being used. They are not deduced merely from the names of those records.
\end{enumerate}

Then the vector b\_f of Π obtained by the permitted G choices and the unchanged vectors elsewhere has the exact K₂-character. Put

\[
 F_v(g)=\iota_\Pi(v\otimes b_f)(g),\qquad F=F_w.
 \tag{PD2}
\]

The source translate x\_u remains in the period, not in an unrecorded replacement of b\_f.~There is a nonzero scalar b=b(T,b\_f), explicitly characterized below, such that

\[
 \boxed{\quad
  \langle F,\operatorname{Pu}_u(T)\rangle
    =c_{\rm Per}\,\overline b\ne0,
  \qquad
  c_{\rm Per}=c_1\,\mu_f(K_2)>0.
 \quad}
 \tag{PD3}
\]

Here c₁ is the \textbf{actual Haar disintegration constant} characterized by PD7. In particular the projection of the literal Pu(T) onto this specified bottom automorphic occurrence is nonzero. This conclusion does not require W(1)≠0 and does not assert a period of the fixed minimal vector at the identity.

\subsection{2. The scalar in the actual source PerLaw is positive}\label{the-scalar-in-the-actual-source-perlaw-is-positive}

The literal source definitions are important. \nolinkurl{WittSmallGroupEmbeddingMemPointwiseStabilizerIso.lean}, lines 176--182, and \nolinkurl{IsThetaFam.lean}, lines 65--66, give a bilinear Bessel period:

\[
 \mathcal B(\varphi,f)
  =\int_{[H]}\int_{[N]}
       \varphi(h)f(nh)\overline{\psi_{\rm Bes}(n)}\,dn\,dh.
 \tag{PD4}
\]

No complex conjugation of the first slot is implicit. The source \nolinkurl{SBes.tieψ} identifies the displayed conjugated Bessel character with the product of the actual Hida real and finite characters.

For the physical Bessel datum B and column T, the literal \nolinkurl{PerLaw} in \nolinkurl{ThetaIdentityOfMapsAll.lean}, lines 130--135, is

\[
 \langle\overline f,\operatorname{Pu}_u(T)\rangle
  =c\int_{S_\infty\backslash G_\infty}
       \mathcal B\left(T(W(y)),
          g\longmapsto f\bigl(\iota(g)(y,x_u)\bigr)\right)
      d\nu_Q(y).
 \tag{PD5}
\]

Here f is the continuous bounded source \nolinkurl{RowFun}, with right character \(\overline{\chi_2}\); equality with its L² representative is read in the source's almost-everywhere sense. All equalities used in the period have their smooth representatives. Inner products throughout are conjugate-linear in the first variable.

The public predicate only asks for an unspecified complex c.~The actual proof proves more. \nolinkurl{FoldIntegralBesselQuotientSupportedFunctions.lean}, theorem \nolinkurl{GlobalTheory.Core.Mi.Fold.fold_quot}, first obtains a positive real c₁ from \nolinkurl{haarMeasure_disintegration_besselGroup} and then chooses

\[
 c_{\rm Per}=c_1(\mu_fK_2)^{\rm toReal}>0.
 \tag{PD6}
\]

The positivity follows from its literal compact-open K₂ measure being finite and nonzero. The disintegration defining c₁ is, for every nonnegative measurable F supported on the actual finite cell S\_f x\_u K₂,

\[
 \int_{G(\mathbb A)}F(g)\,dg
  =c_1\int_{S_\infty\backslash G_\infty}
        \int_{S(\mathbb A)}\int_{K_2}
       F\bigl(\iota(s)(y,x_u)(1,k)\bigr)
        \,d\mu_f(k)\,d\mu_{S}(s)\,d\nu_Q(y).
 \tag{PD7}
\]

Thus c₁ retains the true finite overlap, the Bessel-group measures, and the real quotient normalization. It is not set to one. The source's rational cutoff and unfolding convert PD7 to the same equation on S(F)\G(A), multiplying by μ\_f(K₂) after using right K₂ invariance. In the product f·Ph, right characters cancel exactly as \(\overline{\chi_2(k)}\chi_2(k)=1\). No phase remains in this measure factor.

\nolinkurl{UnfoldedPeriodEqIntegralBesselPeriods.lean}, theorem \nolinkurl{unfoldedPer}, takes exactly this positive c and preserves it in the complex integral. \nolinkurl{ConjugateRowPairingPoincareSeriesBesselPeriods.lean}, theorem \nolinkurl{perLaw_holds}, then uses ordinary Poincaré unfolding with that same scalar. Consequently PD5 holds with the positive c\_Per in PD6. No extra nonzero-scalar premise is needed beyond the genuine source measures and the admitted source unfolding hypotheses.

This is also a useful source-level refinement independent of global distinction: the existence statement in \nolinkurl{PerLaw} may be strengthened to the specific positive real scalar PD6. It neither supplies \nolinkurl{InHH} nor proves Pu(T) nonzero by itself.

\subsection{3. The exact real period function and its continuity}\label{the-exact-real-period-function-and-its-continuity}

Write \(T\xi=B.\mathrm{colForm}(T)(\xi)\) for its actual initial scalar form. Define a linear map on the smooth representation σ∞ by the Riesz identity

\[
 \langle A(v),\xi\rangle
    =\mathcal B\left(T\xi,
       g\longmapsto\overline{F_v\bigl(\iota(g)(1,x_u)\bigr)}\right),
  \qquad \xi\in V_\rho.
 \tag{PD8}
\]

The right side is linear in ξ and conjugate-linear in v, so A is linear in v with the stated inner-product convention. No unidentified duality is being inserted into PD8. The finite-dimensional Hilbert structure on the source Vρ identifies its linear dual with itself conjugate-linearly, exactly as this formula requires.

This functional is continuous on the smooth moderate-growth globalization. The H quotient is compact by the actual \nolinkurl{CoCpt A}; the N quotient is compact because N is unipotent over the number field. The product quotient used in PD4 therefore has compact lifts in finitely many charts after fixing the finite levels. Restriction and evaluation on those charts are bounded by a fixed smooth seminorm of the genuine automorphic embedding. Integrating against the fixed smooth initial form is consequently continuous. This argument does not infer a bounded Hilbert functional from an algebraic Harish--Chandra map.

Change variables in PD4 by an element of S∞. The character relation \nolinkurl{tieψ}, the actual H-column equivariance, and unitarity of ρ give

\[
 A(\sigma(s)v)=\psi_X(s)\rho(s)A(v).
 \tag{PD9}
\]

One can also check the sign directly: translating the second, conjugated form changes its Bessel character to its inverse after the N-variable change; Riesz conjugation in PD8 reverses that scalar once, giving precisely ψ\_X in PD9. For H, \(T\xi(hh_0^{-1})=T(\rho(h_0^{-1})\xi)(h)\), and \(\langle A,\rho(h_0^{-1})\xi\rangle=\langle\rho(h_0)A,\xi\rangle\). There is no second conjugation.

Let A₀ be evaluation at the identity of the bounded relative-model intertwiner constructed in LIFT061, with A₀(σ(y)w)=W(y). Its smoothness and continuity are the local Sobolev conclusion of LIFT061-U/LIFT062. Both A and A₀ lie in the same actual smooth Bessel Hom; Xue's multiplicity-one theorem therefore gives a unique b with

\[
 A=bA_0,
 \qquad A(\sigma(y)w)=bW(y).
 \tag{PD10}
\]

No L² claim about A is needed before applying multiplicity one. Equation PD10 subsequently gives its L² realization. In particular b≠0 exactly when the global period restricted to the chosen finite vectors is a nonzero functional of the real inputs.

\subsection{4. Why the prescribed finite vectors do not annihilate a distinguished pair}\label{why-the-prescribed-finite-vectors-do-not-annihilate-a-distinguished-pair}

The distinction hypothesis is on the two fixed global embeddings. The period is then a nonzero continuous functional on their smooth restricted tensor product. It is nonzero on some pure tensor: algebraic pure tensors span a dense subspace, and continuity applies. At every place its equivariance is the actual local Bessel equivariance. The corresponding Hom has dimension at most one, with the split and nonsplit conventions already fixed in LIFT061 and LIFT071--073.

The needed factorization is elementary and does not require a conjectural global period formula. Fix all but one local slot of a nonzero pure tensor. The resulting nonzero local functional is a scalar multiple of the fixed local Bessel functional. Replacing that slot by another vector on which the local functional is nonzero preserves global nonvanishing. Repeat over a finite set of places. Equivalently, after local normalizations there is a nonzero constant C\_emb, depending on the fixed global embeddings, such that the global functional on pure tensors is C\_emb times the product of their local functionals. This statement is used only with finitely many modifications of a standard restricted tensor; no convergence or nonvanishing of an infinite Euler product is asserted.

At every place outside a finite set, the vectors in the restricted tensor are spherical. LIFT073 proves that their actual local Bessel functional is nonzero. Hence a nonzero initial pure tensor may be replaced by the prescribed spherical vectors at every additional place where it differs from them, one at a time. At the remaining finite places use the following checked tests.

\begin{longtable}[]{@{}
  >{\raggedright\arraybackslash}p{(\columnwidth - 4\tabcolsep) * \real{0.3333}}
  >{\raggedright\arraybackslash}p{(\columnwidth - 4\tabcolsep) * \real{0.3333}}
  >{\raggedright\arraybackslash}p{(\columnwidth - 4\tabcolsep) * \real{0.3333}}@{}}
\toprule\noalign{}
\begin{minipage}[b]{\linewidth}\raggedright
Places
\end{minipage} & \begin{minipage}[b]{\linewidth}\raggedright
Fixed smaller slot and permitted larger slot
\end{minipage} & \begin{minipage}[b]{\linewidth}\raggedright
Actual local conclusion
\end{minipage} \\
\midrule\noalign{}
\endhead
\bottomrule\noalign{}
\endlastfoot
G & The given inverse-generic finite smaller vector; larger vector chosen inside the already admitted finite Π₀ using GB7 & LIFT072: \(\alpha_G=c_Hc_N/d_\varsigma>0\), with the full source compact character \\
T & The original strict Mram conductor-one vector; the unchanged spherical larger vector, with the physical Dτ translation & LIFT071: \(\alpha_T^\natural=c_Hc_N\,q/[(q-1)(q^k-1)]>0\) \\
Away G∪T & The actual spherical vectors in the prescribed hyperspecial lattices and different-scaled character & LIFT073: \(\alpha_v^\natural=c_Hc_N>0\) in its precise split/nonsplit scope \\
\end{longtable}

The calibrated local positive matrix-period values imply nonzero local linear Bessel functionals on these exact test pairs. This uses the local calibration proved in the cited results, not generic uniqueness to guess a scalar. All unit/Gauss phases, character changes, and Haar factors remain in the local functionals and in C\_emb. They are not discarded by writing a positive squared modulus.

The initial finite vector a\_f is unchanged. Only the G choices explicitly authorized by LIFT072 are made on the bottom side; they stay in the same Π₀ and actual K₂-character space. The proof does not choose an unrelated Whittaker vector, add a prime, refine a level, or average the T character away. After these substitutions the global functional of the remaining real inputs is still nonzero. Thus

\[
 b(T,b_f)\ne0.
 \tag{PD11}
\]

For the real slots, PD10 is the required normalization. Their local Bessel functional is not evaluated only at w and the identity. Since w is cyclic in the irreducible σ, the nonzero A₀ cannot vanish on every σ(y)w. The full real function W supplies the appropriate contraction even when A₀(w)=W(1)=0.

\subsection{5. The Wu contraction gives the source pairing exactly}\label{the-wu-contraction-gives-the-source-pairing-exactly}

Cuspidality, fixed finite level, and the actual K-finite w make F in PD2 a smooth bounded automorphic function. This is the usual cusp rapid-decay consequence used in LIFT062, with the compact real center and finitely many arithmetic components retained. Its conjugate f=\overline F is a source \nolinkurl{RowFun} with right character \(\overline{\chi_2}\). Its smooth representative agrees with its L² class. The source column and the actual finite translate are exactly those in PD8.

For every y, right equivariance gives

\[
 \mathcal B\left(T(W(y)),
    g\longmapsto\overline{F\bigl(\iota(g)(y,x_u)\bigr)}\right)
   =\langle A(\sigma(y)w),W(y)\rangle
   =\overline b\,\|W(y)\|^2.
 \tag{PD12}
\]

The last expression is absolutely integrable by PD1. The admitted source unfolding already supplies the first integral identity; PD12 also identifies its exact value without an exchange involving an unbounded distribution. Integrating, substituting f=\overline F in PD5, and using PD6 yields PD3.

This calculation is why the earlier failed transfer of the origin-value estimate is not needed for the present conditional theorem. The contraction uses the entire relative vector and its actual unit norm. It does not claim an origin value or a uniformly small support radius.

\subsection{6. Quantitative statement with the unknown global coupling retained}\label{quantitative-statement-with-the-unknown-global-coupling-retained}

If the chosen automorphic embedding is Hilbert-isometric and w,b\_f have norm one, then \(\|F\|_2=1\). More generally retain \(\|F\|_2\) explicitly. Orthogonal projection onto the line spanned by F and Cauchy--Schwarz give

\[
 \|\operatorname{Pu}_u(T)\|_2^2
   \ge \frac{c_{\rm Per}^2|b(T,b_f)|^2}{\|F\|_2^2}>0.
 \tag{PD13}
\]

With the literal positive source c₀=c0X(u), the diagonal normalized Gram value satisfies

\[
 K_{u,u}(T,T)
   =c_0^{-1}\|\operatorname{Pu}_u(T)\|_2^2
   \ge \frac{c_{\rm Per}^2|b(T,b_f)|^2}
                  {c_0\|F\|_2^2}>0.
 \tag{PD14}
\]

Both c\_Per and c₀ are actual source quantities. The formal degree and compact dimensions have already entered the definition of φ, the unit relative vector, and c₀. No extra factor dσ or dimρ is inserted into PD12. The global coupling b remains unevaluated. Therefore PD13--14 are positive conditional lower bounds, not numerical or uniform bounds for an inverse operator or for DP44.

If T is rescaled by t and the row F by z, then the bilinear source period and conjugate-first Hilbert convention give

\[
 b(tT,zF)=\overline t\,z\,b(T,F),\qquad
 \langle zF,\operatorname{Pu}(tT)\rangle
     =\overline z\,t\,\langle F,\operatorname{Pu}(T)\rangle.
 \tag{PD15}
\]

Equation PD3 respects these exact transformations. Multiplying W by a phase ζ, with \textbar ζ\textbar=1, multiplies Ph and Pu by ζ and changes A₀ by ζ, so b changes by ζ⁻¹. Again PD3 is unchanged as an identity. No vacuum, CM-coordinate, arithmetic coefficient, or source-intertwiner scale is silently renormalized here.

\subsection{7. What is removed, and what remains}\label{what-is-removed-and-what-remains}

The theorem removes an additional prescribed-test obstruction \textbf{within a specified globally distinguished pair}: all actual finite tests survive, and the real Wu contraction cannot cancel that pair merely because the fixed minimal vector has zero Bessel value at the identity. It gives a nonzero spectral block of the literal Poincaré detector on the chosen initial Mram column.

The following distinctions remain part of the theorem.

\begin{longtable}[]{@{}
  >{\raggedright\arraybackslash}p{(\columnwidth - 2\tabcolsep) * \real{0.5000}}
  >{\raggedright\arraybackslash}p{(\columnwidth - 2\tabcolsep) * \real{0.5000}}@{}}
\toprule\noalign{}
\begin{minipage}[b]{\linewidth}\raggedright
Requirement
\end{minipage} & \begin{minipage}[b]{\linewidth}\raggedright
Status here
\end{minipage} \\
\midrule\noalign{}
\endhead
\bottomrule\noalign{}
\endlastfoot
Existence of an actual bottom cusp constituent & Available only in the stated LIFT059 chosen-space scope \\
Actual initial/bottom local temperedness and genericity & LIFT069/LIFT070 or the separate admitted concrete scope; not enlarged here \\
A nonzero global Bessel functional on the two specified embeddings & \textbf{Explicit unresolved premise} \\
Arbitrary automorphic multiplicity vectors & Not covered by nonvanishing for one different occurrence \\
Exact Cu/Mram and bottom K₂-character vectors & Retained, with only the G choice authorized by LIFT072 \\
Full arithmetic \nolinkurl{Read}, \nolinkurl{STiesR}, source tS, and requested specific032 row/value & Not constructed \\
An actual same-family theta identity/row pin transferring this nonzero Pu block to the terminal arithmetic map & Must be supplied in its literal form; not inferred from a local model \\
Numerical global coupling, inverse/preimage cost, or DP44 budget & Not supplied by PD13--14 \\
\end{longtable}

The finite translate x\_u in PD5 is not set equal to tS. When the accepted translated row dictionary is instantiated, its exact tS, norm factor, central pin, and characters must be transported by LIFT062-T/LIFT065/LIFT066. A source-sanctioned refill cannot itself manufacture the missing arithmetic instance.

Nor is this a new family constructor. \nolinkurl{SizesAt} alone omits \nolinkurl{SizeTu}, but the enclosing \nolinkurl{DesignExists} explicitly requires \nolinkurl{TuSize}: \nolinkurl{ParamsHeightRankConstant.lean}, lines 180--190, gives the eventual bound \(\log T\le C_{Tu}mL\) in addition to \nolinkurl{TuGe}. Any future relative domination estimate must keep this actual family bound. Its existence does not prove the distinguished-pair premise above, and the established fixed-degree full-size exclusions are not altered.

\subsection{8. Source and verification anchors}\label{source-and-verification-anchors}

\begin{itemize}
\tightlist
\item
  \nolinkurl{PoincareInHH.lean}, \nolinkurl{Cu}, \nolinkurl{phVal}, \nolinkurl{Ph}, \nolinkurl{InHH}, \nolinkurl{PuFun}, \nolinkurl{Pu}, \nolinkurl{c0X}, and \nolinkurl{Kform}: the actual column, translate, finite character, Poincaré map, and normalization.
\item
  \nolinkurl{ThetaIdentityOfMapsAll.lean}, \nolinkurl{SBes} and \nolinkurl{PerLaw}: actual physical groups, conjugated character, and the precise period pairing.
\item
  Pinned source recovery \nolinkurl{FoldIntegralBesselQuotientSupportedFunctions.lean} and its statement; \nolinkurl{HaarMeasureDisintegrationBesselGroupStatement.lean}; \nolinkurl{UnfoldedPeriodEqIntegralBesselPeriods.lean}; \nolinkurl{ConjugateRowPairingPoincareSeriesBesselPeriods.lean}: the positive scalar PD6 and its unchanged passage through both unfoldings.
\item
  LIFT061, RW5 and RW18--28; LIFT061-U; LIFT062: genuine smooth Bessel multiplicity, actual unit Wu, and the prescribed-vector realization. Xue, H. (2025), \emph{Bessel models for unitary groups and Schwartz homology}, §§1.1--1.2 and Theorem 1.1, author preprint at \url{https://preprint.press.jhu.edu/ajm/sites/default/files/AJM-h-xue.pdf}, supplies the relevant smooth Hom convention; the cited local multiplicity-one results are used with these conventions. The exact primary passages were checked in LIFT061; this report does not claim a new independent primary-paper review.
\item
  LIFT071 HN8/HN22--24, LIFT072 GB2/GB7--9 and its exact actual-frame dictionary, LIFT073 SG theorem: the calibrated finite local tests.
\end{itemize}

The companion diagnostic checks finite tensor factorization, conjugate-first Riesz/Gram identities, positive measure scaling, column/row/Wu phase transformations, and explicit zero-at-identity examples of the full Wu contraction. It does not simulate a p-adic representation, establish a global distinguished pair, run Lean, or replace any analytic argument above. Final whole-file/source/script review, replay, and fingerprints are recorded separately when complete.

\clearpage

\section{AM-K093 --- the actual scalar-decoration Fourier matrix and its residual ordinary derivatives}\label{am-k093-the-actual-scalar-decoration-fourier-matrix-and-its-residual-ordinary-derivatives}

\textbf{Canonical checked integration --- 2026-11-22T20:36:46.817565+00:00.} Root and independent whole-proof/source/primary/script reviews and separate7,800-assertion replays passed. The actual scalar-decoration matrix has the exact finite Fourier/augmentation expansion RT13, represented by rational rank-one first-edge Weyl translations with the original central phase, full Y, class-dependent covectors and contracted full-middle mass. Every nonzero finite Fourier coefficient is an11-unit. The actual selected symbol along a rank-one increment has degree\textless=nu\_tau,1, so its ordinary translated evaluation retains at least D\_tau-nu\_tau,1\textgreater0 first-edge derivatives at every label. The pure-exponential term is zero; the complete differentiated germ and class sum remain unevaluated. No selected-value/content, global nonzero, arithmetic inverse, source family, Read/032 or DP44 conclusion is asserted.

\textbf{Algebraic-geometric translation.} The rational rank-one lift uses the canonical tensor-Hom isomorphism for the original finite projective modules and module CRT at paired CM primes. Its translated section lives in the actual rigidified torsion difference line in the Picard/Poincare framework; the complete Mumford lift fixes the central scalar. The integral selected Weyl symbol restricts to the determinantal rank-one locus with the stated degree bound, but evaluation still uses the original ordinary differential functional. Exact carried anchors are FGA232 Theorem3.1 pp148--149; FGA236 Corollaries1.7/2.3 pp229/232; SGA7 I VII2.7/2.9 pp169--174 and VIII7.4(b) pp310--311; SGA3 I4.7.3/4.7.3.1; SGA6 VI Theorem3.3; cote4 scan003 §§8.1.1--8.1.2, scans005--006 §§8.4.1--8.4.3 and cote54 scans005--007. No new lambda-ring/K-theory identity, l-adic purity/weight theorem, CM-motive comparison or motivic L-function claim follows. DP\_MULT Q3e and THETA3 G1/A7 retain the exact selected derivative/class gate; CX\_GRH CX-ALG-024 and GCT\_PLETH T-052 retain their separate arithmetic gates.

Author: \nolinkurl{/root/kummer_extension}, 22 November 2026 UTC.

\textbf{Author whole-proof, source-interface and complete-script PASS; frozen for independent review.} This is an exact reduction of the prescribed selected arithmetic matrix. It does not assert that a translated germ, an individual selected value, or the final class sum is nonzero.

\subsection{1. Fixed data and the literal finite Fourier expansion}\label{fixed-data-and-the-literal-finite-fourier-expansion}

Keep one actual host, word, good tuple, canonical twist, initial coarse class c, middle class f and terminal decoration d from K038/K090/K092. All non-T frames, both raw G factors, every Q factor, the full Y sum, the ordinary11 vacuum, the complete middle coevaluation and the same initial coefficient vector P remain fixed. In particular, Y has G and Q factors and no T factor. Write a for the initial rank, b for the middle rank, and f\_CM={[}F:Q{]} for the number of selected T places. The symbol f in a class index or a class sum still denotes a middle class. At a selected T component use the actual integral Hom coordinates and actual primitive covectors u0 and u1,f.~The latter depends on f; it is not replaced by a single standard covector for the class sum.

The first and second T conditions are precisely \[
 u_{1,f}X=s_vu_0,\qquad u_2Z=u_{1,f},
 \tag{RT1}
\] with X and Z integral. Here s=(s\_v) lies in D\_T=product\_(v\textbar T) k\_v\^{}times, and k\_v=F\_T in the actual completely split prime. Define Phi\_(c,f,d)\^{}base by deleting only the first equality in RT1, retaining its integral indicator and the complete second incidence and all remaining finite data. This is a proof device inside the same finite Schwartz space, not an alternative admitted input.

For lambda=(lambda\_v), lambda\_v in k\_v\^{}a, let M\_(lambda,f) multiply the first position variable by product\_v psi\_v(u1,f X lambda\_v). Finite additive-character orthogonality gives the exact identity \[
 \Psi_s=T^{-a f_{\rm CM}}\sum_{\lambda\in\prod_vk_v^a}
       \prod_v\psi_v(-s_vu_0\lambda_v)\,
       M_{\lambda,f}\Phi_{c,f,d}^{base}.
 \tag{RT2}
\] This uses the literal CharT.psiT with CharT.Ok, not an arbitrarily chosen additive character. Both sides retain the same full Y contraction. The framed G cutoff C\_f of K092 commutes with these T operators: locally it is the G multiplier, and the common class/polarization transport conjugates both sides of this equality. Thus RT2 also holds after the same full-middle-mass proof-device cut.

For the actual tensor augmentation generator b\_s=product\_v(delta\_(s\_v)-delta\_1), with every s\_v different from1, its Fourier coefficient is \[
 \widehat b_s(\lambda)=\prod_v
   \bigl(\psi_v(-s_vu_0\lambda_v)-\psi_v(-u_0\lambda_v)\bigr).
 \tag{RT3}
\] It is zero if any u0 lambda\_v=0 and is an11-adic unit otherwise. Indeed T is prime, each nonzero difference is a root of unity times 1-zeta\_T\^{}j with j nonzero modulo T, and the norm of the latter is T. CharT.Ok already puts the needed T-th roots in the carried coefficient field, as proved in K059; no extra coefficient-field enlargement is asserted. T\^{}(-a f\_CM) is also an11-unit. These unit assertions concern the finite Fourier coefficients only, not their selected evaluations or their sum.

\subsection{2. A rational rank-one representative of each actual multiplier}\label{a-rational-rank-one-representative-of-each-actual-multiplier}

Let Lambda\_01 be the actual integral first-edge Hom lattice in its rational unitary space. At each selected split T pair, the actual symplectic trace pairing identifies the opposite component with the dual of the position Hom lattice. The functional X -\textgreater{} u1,f X lambda\_v is rank one in that dual matrix space. The source character and trace-pairing units determine its precise opposite-component residue; include their sign and units in this residue, so that the resulting finite Weyl operator has exactly the multiplier M\_(lambda,f).

There is a rational first-edge vector q\_(lambda,f), of E-matrix rank at most one, such that \[
 Tq_{\lambda,f}\in\Lambda_{01},\qquad
 W_f(q_{\lambda,f})\Phi_{c,f,d}^{base,cut}
       =M_{\lambda,f}\Phi_{c,f,d}^{base,cut}.
 \tag{RT4}
\] Here W is the actual symmetric Weyl lift, including its central half-pairing phase; it is not an unspecified intertwiner. The second edge is unchanged.

\textbf{Construction.} For the actual finite projective modules L0,L1, the canonical map L1 tensor L0\^{}dual -\textgreater{} Hom\_(O\_E)(L0,L1) is an isomorphism. In particular a pure tensor gives an integral Hom vector of E-rank at most one. Their localizations at T are free. Express the rank-one dual residue by a target column and a source dual row, using the actual perfect Hermitian pairings. The target column may be chosen once for this middle class and all lambda; the row carries lambda and the already fixed local character unit. By module CRT choose global integral factors w\_f and v\_lambda in these two actual projective modules. At the momentum partners prescribe their residues to realize the displayed dual functional after division by T. At every position prime p over T prescribe w\_f to lie in p\^{}2 L1. Since the rational T has valuation one there and the source row is integral, the position component of q=w\_f v\_lambda/T lies in p times the position lattice. These are independent CRT conditions at distinct CM primes.

At the finitely many remaining bad finite places, including2,11, G, Q and every denominator of the actual class frames, impose further divisibility on w\_f.~A locally constant compactly supported joint Schwartz tensor has an open Heisenberg stabilizer. Choose this divisibility so that q lies in that exact stabilizer, including the central phase, for every row v\_lambda in the integral source dual module. The same condition can be used for the finite set of lambda. At all other places the integral Weyl lattice fixes the basic tensor. These CRT conditions are compatible because all added places are away T. If a rational coordinate basis is used in this construction, its finitely many denominators are included in the bad set; no moving rational basis is declared integral.

At a T position component, translation by the chosen position part in T Lambda preserves the integral first-edge indicator. The second incidence is independent of X. The momentum part has denominator at most T and realizes exactly psi\_v(u1,f X lambda\_v). Its half-pairing with the T-divisible position part is integral; T is odd, so the symmetric Weyl phase is one there. Away T the exact stabilizer condition includes the phase. This proves RT4 with its actual scalar, rather than up to an unspecified constant. The factors are integral in the actual modules, hence Tq lies in the actual global Hom lattice. The trivial multiplier lambda=0 can use q=0.

The construction is a representation of an already fixed finite operator. It changes no finite level, Fourier coefficient, chosen covector or input tensor. It gives no quantitative height estimate for q. Different allowed rational representatives can give different descriptions of the translated germ, but RT4 fixes their action on this particular tensor, hence fixes the eventual selected value. The representatives need not form an additive family in lambda; no missing Heisenberg cocycle is suppressed by claiming that they do.

\subsection{3. The source Weil sign and exact reduced translation}\label{the-source-weil-sign-and-exact-reduced-translation}

Use the actual reduced reading of K035(PF21--28), first in normalized symplectic coordinates. Write tau\_per for the symmetric period matrix, Y=Im tau\_per, and H(z,t)=z\^{}t Y\^{}(-1) conjugate(t), linear in its first variable. For rational A,B let t=A+tau\_per B. The actual archimedean Schrödinger action is \[
 \rho_\infty(A,B)\phi(x)
  =e^{-2\pi iA^tx+\pi iA^tB}\phi(x-B).
 \tag{RT5}
\] The finite lift is the compatible source Weil lift with the same rational Weyl vector and product-formula character. The normalized reduced theta reading is \[
 R\phi(z)=e^{\pi z^tY^{-1}z/2}
   \sum_{r\in\mathbb Q^g}\phi(r)
       e^{\pi i r^t\tau_{per}r+2\pi i r^tz}.
 \tag{RT6}
\] After translating the finite tensor, the product formula and the substitution r=q-B give the exact identity \[
 R(W_f(A,B)\phi)(z)
   =e^{-\pi H(z,t)-\pi H(t,t)/2}\,R\phi(z+t).
 \tag{RT7}
\] For example the scalar produced before comparison of the two reduced Gaussians is exp(pi i B\^{}t tau\_per B + pi i A\^{}tB +2pi i B\^{}tz). The linear exponent after comparison is -pi z\^{}tY\^{}(-1)t+2pi i B\^{}tz=-pi H(z,t), and the constant is -pi H(t,t)/2. Thus the point is +t in this reading. This verifies the sign and the central phase, instead of inferring them from an unsigned translation convention.

K035(PF25) is exactly the source convention b\_src=-Im H and psi\_src=exp(-2pi i dot); the varying CM periods and actual class frames transport RT7 without changing this identity. For q\_(lambda,f) in RT4, t\_(lambda,f) is prime-to11 torsion in the actual first-edge torus. Each of its CM matrix components has rank at most one; Hermitian duality, period scalars and source/target frame changes are left/right invertible matrix operations and preserve that rank.

On the proof-only comparison cover of K089/K092 the corresponding point and derivative coordinates must also be pulled back by the actual unit-coordinate map. Its denominator D is prime to11; a lifted point may have order dividing TD. We always read RT7 in the original source coordinates after that exact transport. We do not equate an untransported cover derivative with a source derivative. The actual period, determinant splitting factor, Chern normalization and finite frame scalar are the ones in K035(PF28), K038 and K092(PM6).

K034(TG3--11) identifies the complete prime-to11 torsion difference-line frame integrally. That theorem applies to the whole reduced translation with its actual finite lift. It does not assign an11-adic valuation to either displayed complex Gaussian separately, to pi, or to a complex period. None is called an11-unit here. Nor is the original ample-line reduced frame claimed integral.

\subsection{4. A rank-one degree bound for the actual selected symbol}\label{a-rank-one-degree-bound-for-the-actual-selected-symbol}

At each CM label tau let nu\_tau=eta\_tau+q\_tau 1\_a, D\_tau=\textbar nu\_tau\textbar{} be the exact first-edge selected partition of K037(CR2--3), with q\_tau\textgreater=1. The actual substitution is \[
 \mathcal Q_0(P)(Y,Z)=
    [\det(U)^{q_\tau}P(U)]_{U=Z^tY_{[,1:a]}}
 \tag{RT8}
\] at that label, with the tensor product over labels understood. All target dual vectors and the complete two-edge coevaluation are retained as in K091.

For every first-edge rank-one matrix M, every actual coefficient vector P and every target coefficient of the complete selected symbol, \[
 \deg_s\mathcal Q_{cpl}(P)(Z_0+sM,Z_1)
                         \le \nu_{\tau,1}.
 \tag{RT9}
\] This is an integral polynomial statement. It does not require division by any factorial, hook product or compact dimension.

\textbf{Proof.} The highest polynomial p\_nu(U) is the product of its leading r-minors to powers nu\_r-nu\_(r+1). Under a rank-one increment U+sR every minor is affine in s. Therefore its degree is at most sum\_r(nu\_r-nu\_(r+1))=nu\_1. The actual integral Weyl lattice is generated by the source row divided-power operations. Equivalently these are coefficient extractions from products of integral elementary row matrices acting on the highest polynomial. Left row multiplication keeps R of rank at most one, so taking any such coefficient preserves the same s-degree bound. The determinant insertion in RT8 is exactly the prescribed weight shift; the same argument applies before the actual substitution. Finally Z -\textgreater{} Z+sM changes Z\^{}tY by the rank-one matrix M\^{}tY. Taking second-edge coefficients and the finite integral middle coevaluation is linear and cannot increase the degree. This proves RT9 on the full actual symbol lattice, not merely on its highest line. The relevant integral row operation is the literal rowDiv in StatementA.lean and the actual lattice transport of K037(CR6--10)/K091(SR9--13).

The first-edge symbol is homogeneous of degree D\_tau. Hence its translation by a rank-one constant matrix has components only in degrees \[
             D_\tau-\nu_{\tau,1},\ldots,D_\tau.
 \tag{RT10}
\] In the actual ranks a\textgreater=2, \[
 D_\tau-\nu_{\tau,1}
       =\sum_{i=2}^a\nu_{\tau,i}\ge(a-1)q_\tau>0.
 \tag{RT11}
\] In particular the symbol evaluated at a single rank-one matrix is zero. This conclusion is compatible with the integral split symbolic map in K091/K092: it is a restriction on this small family of evaluations, not a failure of the symbolic retraction.

\subsection{5. The exact residual ordinary derivative}\label{the-exact-residual-ordinary-derivative}

Write F\_(c,f,d)\^{}base,cut for the complete original-coordinate reduced germ of the base cut tensor. Keep all other first/second-edge variables and labels. For the actual t in section3 put \[
 L_t(z_0)=-\pi H(z_0,t)-\pi H(t,t)/2.
\] With the literal ordinary source derivative, ordinary Leibniz gives \[
 [\mathcal Q_{cpl}(P,v)(cD\partial)
                (e^{L_t(z_0)}F^{base,cut}(z_0+t,z_1))]_{(0,0)}
  =e^{L_t(0)}
   [\mathcal Q_{cpl}(P,v)
       (cD\partial_0+cD\,dL_t,cD\partial_1)
                F^{base,cut}]_{(t,0)}.
 \tag{RT12}
\] The same source chart phase and unit-coordinate cover transport occur on both sides and are retained. The matrices cD dL\_t have rank at most one at each label. By RT10, every term on the right still applies at least D\_tau-nu\_(tau,1) ordinary first-edge derivatives at that label to F. The second-edge selected derivatives also remain. The term where every first-edge derivative falls on the exponential is zero by RT11.

Thus the actual T modulation cannot be replaced by the scalar formula ``selected symbol at the translate times a theta value.'' The exact formula is RT12 with a positive residual derivative order. This observation is neither vanishing of RT12 nor a lower bound for it. It also supplies no replacement by freely prescribed Hasse jets or by the full local germ module. The ordinary/divided-power and source cD factors have not been exchanged.

\subsection{6. The resulting prescribed matrix and its remaining content}\label{the-resulting-prescribed-matrix-and-its-remaining-content}

For an augmentation generator with coefficient P, K092(PM6/PM9), RT2--4 and RT7 give \[
 (\mathcal A B_{c,s,P})(d)
 =T^{-a f_{\rm CM}}\sum_f\sum_\lambda\widehat b_s(\lambda)
     \mathcal J_{f,d,P}
       \bigl(U_{t_{\lambda,f}}F_{c,f,d}^{base,cut}\bigr).
 \tag{RT13}
\] Each J is the prescribed full ordinary two-edge selected evaluation, or equivalently RT12, with the actual terminal pairBK/brC transport. The source fixed period, determinant splitting and finite Haar/mass scalars are exactly those of K038/K092. In particular the division by the entire middle stabilizer is already realized by the contracted cut identity; it is not performed again in RT13. The same initial P is used for every f, while both the middle covector and the germ depend on f.~No independently chosen class jet is available.

The finite nonzero coefficients in RT13 are units at11. This removes no further possible content from the derivative evaluations, from their sum over lambda or from the subsequent sum over f.~Generic and irreducible finite projectors may have denominators divisible by11; RT13 does not reclassify them as Heisenberg normalization units. Every later genuine/generic projector, terminal mass and prior-image restriction remains as in K090--092.

The established reduction is the exact restricted rational-translate matrix RT13 together with the forced residual ordinary derivative range RT10--12. The unresolved question is the content or nonvanishing of that actual matrix. Primitive finite functions, the primitive auxiliary cut section, and a split selected symbol before evaluation do not settle it. No new level, unrestricted torsion interpolation family, sharp actual-image valuation, nonzero arithmetic row, global period, family constructor or DP44 conclusion is obtained.

\subsection{7. Sources and verification scope}\label{sources-and-verification-scope-11}

The literal finite sources are \nolinkurl{FiniteSumNonzeroLargeDepth.lean,} phiT at648--650, Y at1125, XiCpl1 at1134, ftY at1151 and XiCpl2 at1160; ScaledRamifiedLattice.lean, CharT at141--148; and the scalar decoration action in StdLatUnimodular.lean at113--122. Exact source-dependent interfaces are K035(PF21--28), K037(CR2--10), K038(CA15--27), K039(IG3--7), K059's CharT.Ok argument, K090(MA1--18), K091(SR9--13), and K092(PM1--9). The row divided-power source is \nolinkurl{Formalism/GradedWeylDivisibility/StatementA.lean,} rowDiv and its polynomial coefficient definition.

The reduced torsion translation is the already verified primary convention of Bannai, K., and Kobayashi, S. (2010), \emph{Algebraic theta functions and the p-adic interpolation of Eisenstein--Kronecker numbers}, Duke Mathematical Journal,153(2),229--295, section1.3, formulas (14) and the definition of U\^{}M on preprint pages14--15; Proposition2.3 on pages21--22; and Lemma2.7/Proposition2.8 on pages24--25 of arXiv:math/0610163v4, \url{https://arxiv.org/pdf/math/0610163}. These exact primary passages were reread directly on 22 November 2026 at 20:19 UTC. Its exact application, including the difference-line integral frame and source scalar, is K034(TG3--11). No ordinary-reduction integrality assertion from a different theorem is used.

The companion certificate passed 7,800 exact assertions: finite Fourier/augmentation identities in exact cyclotomic rings, a paired-prime CRT rank-one model, rank-one degree bounds under integral row operations, the ordinary Leibniz identity and the reduced Weyl sign. It includes nonzero row words which leave and re-enter the leading block, a wrong-sign check, the ordinary/Hasse distinction and an example where the pure exponential contribution vanishes while the complete differentiated germ is nonzero. These finite models do not compute an actual selected class value, prove a content bound, or replace the global CRT and actual-source proofs above.

\subsection{Author whole verification}\label{author-whole-verification-2}

At 2026-11-22T20:30:42.406542+00:00, I completed a whole read of all seven substantive sections RT1--13 and the entire saved companion script. I reread the literal phiT and full-Y definitions, the rowDiv coefficient definition, the actual K035 reduced/source reading and ordinary creation comparison, K037's exact partitions and integral row-matrix stability, K034's complete torsion difference-line normalization, and the K091/K092 symbolic and contracted-cut interfaces. I reconstructed the paired-prime projective-module CRT lift with its central phase, checked the source product-formula sign rather than changing the finite lift, and checked the full selected symbol's rank-one degree bound through both substitutions and complete middle coevaluation. The exact actual class formula retains the same input coefficient across every middle class. The certificate passed 7,800 finite algebra assertions; it does not evaluate a prescribed selected theta derivative. The final projective Hom, p\^{}2 residue, field-degree notation, source locators, scope and status edits were reread. The separate verification and manifest record the final unchanged-source replay and hashes.

\clearpage

\section{AM-DO-091 --- two integral transverse pairings close the rank-one fixed-block support twenty-three sector}\label{am-do-091-two-integral-transverse-pairings-close-the-rank-one-fixed-block-support-twenty-three-sector}

\textbf{Canonical checked integration --- 2026-11-22T20:47:59.410820+00:00.} Root and independent whole-proof/script/interface checks and separate1330-assertion replays passed. The entire actual all-top H133 globally rank-one-unipotent fixed-W0 class at support23 has absolute contribution\textless2\^{}-1544. The original transverse pairings give a complete witness-or-norm-equation split. The norm sector retains the actual affine difference ideal, epsilon1/12 E/F divisor counts, real pure-imaginary ideal covolume, all24 marked sizes, original type/pair weights and selected near-point metric. Global unipotence is a hypothesis. The complete all-top floor remains23 and the uniform floor11; the other structural classes and all signed/global/arithmetic/DP44 gates remain open.

\textbf{Algebraic-geometric translation.} The new counts use the original Hermitian modules, actual integral pairing sections and relative norm ideals; the pure-imaginary line is the real ideal i Lambda\_F with its real different and covolume. No moving integral basis, changed source normalization or arithmetic theta-image inverse is introduced. Carried exact contextual anchors are SGA3 I4.7.3/4.7.3.1; SGA6 VI Theorem3.3; FGA232 Theorem3.1 pp148--149; FGA236 Corollaries1.7/2.3 pp229/232; SGA7 I VII2.7/2.9 pp169--174 and VIII7.4(b) pp310--311; cote4 scan003 §§8.1.1--8.1.2, scans005--006 §§8.4.1--8.4.3 and cote54 scans005--007. No new lambda-ring/K-theory identity, l-adic sheaf or weight theorem, CM motive or motivic L-function assertion follows. DP\_MULT Q3e and THETA3 G1/A7 gain the stated containing-set removal; CX\_GRH CX-ALG-024 and GCT\_PLETH T-052 retain their arithmetic gates.

Author: \nolinkurl{/root/definite_oscillator}, ALP\_MACHINE. 22 November 2026 UTC. Complete proof; author and independent whole-file checks are recorded separately before promotion.

\subsection{1. The unchanged sector and the two new counts}\label{the-unchanged-sector-and-the-two-new-counts}

Retain the original J56 all-top kernel, its already charged cutoff H=133, f=a=168, the original five-coordinate space W0 with form J0=diag(1,1,1,1,−1), all paired labelled T pencils, the same near-origin packet, both mixed G projectors, all class masses and the genuine selected Hilbert metric. This proof concerns only

\[
 h=I_S\oplus(I+N),\qquad \operatorname{rank}_E N=1,\quad N^2=0,
 \qquad s=23.                                                   \tag{TD1}
\]

Global unipotence is an explicit hypothesis. The preserved signed summand is bounded in absolute value; no positivity is assigned to it. The new theorem removes this complete TD1 sector with the single absolute charge 1544, that is, a bound strictly less than 2\^{}−1544. It does not alone raise the complete support floor.

Write e for the original negative vector, x=Ne and eta=x\_n.~The exact K075(RN5--12) normal form gives

\[
 N=x\bar x^{\mathsf t}J_0/\eta,\quad \bar\eta=-\eta\ne0,
 \quad \sum_{j=1}^4|x_{j,\tau}|^2=|\eta_\tau|^2,
 \quad |\eta_\tau|\le R=\frac{8844}{133}.                       \tag{TD2}
\]

A full column determines N uniquely. In particular, the subsequent reduced column count has no separate stabilizer-fibre multiplier.

The argument first uses two globally integral pairings. Unless they both vanish, a T-free field-norm witness removes the term. When both vanish, the first three positive coordinates are scalar multiples of one fixed primitive integer vector. The resulting two-variable relative norm equation is counted using the actual affine pencil ideal, an exact real-field covolume for eta, and an explicit ideal-divisor bound of exponent 1/12.

\subsection{2. Integral transverse pairings and the original two-coordinate reduction}\label{integral-transverse-pairings-and-the-original-two-coordinate-reduction}

Use the actual packet vectors of DO078(FP20--22) and K088(IP5--6):

\[
\begin{split}
 u&=(2,95192390527,383092671043,12529843548438452388609),\\
 y_1&=(-976627,-153032,38026,0),\\
 y_2&=(-567678,695570,-172838,0),\\
 y_3&=(104583,-15309315176,-61610045224,2).
\end{split}                                                       \tag{TD3}
\]

Each y\_j is integral and belongs separately to every actual source and target pencil lattice at both CM partners. This is the carried original-coordinate membership statement; it does not assert that h preserves one common lattice. Since h carries its actual source to its actual target, the paired integral form gives

\[
 b_j=\langle y_j,Ny_j\rangle\in\mathcal O_E\cap iF=i\mathcal O_F.
                                                                    \tag{TD4}
\]

The same conclusion holds away from T because the retained original model is integral there. In TD2 the pairing is a nonzero pure-imaginary scalar times the squared absolute value of the linear functional defined by y\_j. Thus b\_j=0 if and only if the positive column x is orthogonal to y\_j.

The exact squared lengths of y1,y2 are 978665066829 and 835948910828. Therefore

\[
 |(b_j)_\tau|\le R\|y_j\|^2<2^{46}\qquad(j=1,2).                 \tag{TD5}
\]

If at least one is nonzero, this is a nonzero global integral Green witness. At a rank-zero unmixed G reduction it vanishes at both CM partners. Hence the full paired rank-zero ideal J divides (b\_j), giving N\_E J\textless(2\^{}46)\^{}(2f). The complete rank-one Green bound consequently implies

\[
 |a(h)|\le 2\,(2^{46})^{af}Q^{-a(f-1)}.                           \tag{TD6}
\]

There is no extra enumeration of G patterns in this pointwise norm-witness bound.

Now suppose b1=b2=0. Put

\[
 v=(2,95192390527,383092671043),\qquad
 V=4+95192390527^2+383092671043^2.                                 \tag{TD7}
\]

The first three entries of y1,y2 are independent and span v-perp. Their exact dot products with v vanish. Thus in the original coordinates

\[
       x=(v q,r,\eta),\qquad
       Vq\bar q+r\bar r=\eta\bar\eta.                            \tag{TD8}
\]

No orthogonal change of arithmetic basis is made. The vector v is primitive because its first entry is 2 and its second entry is odd. Consequently the coordinate-content ideal is exactly

\[
                 \mathfrak b=(q,r,\eta).                          \tag{TD9}
\]

If q=0, then b3=4N44 is nonzero and \textbar(b3)\_tau\textbar≤4R\textless2\^{}46. TD6 again applies. Combining these cases as one witness sector is legitimate: every term in their union has some nonzero integral witness with the same archimedean bound, so TD6 is uniform on the union without choosing or counting an additional witness.

The complete original five-column bound is G5·2\^{}5757·T\^{}(10s), with G5\textless2\^{}107, by the actual full target-overlattice count used in DO075 and DO089. It includes zero target labels and all five complex coordinates. Using support≤2\^{}168, full ordered-pair weight P\_T\^{}s\textless2\^{}(997s), and the real bound one, the whole witness sector is bounded by

\[
 \log_2 U_{\rm wit}<1+168+107+5757+46af
                  +(10\ell_T+997)23-a167\ell_Q,                  \tag{TD10}
\]

where ell\_T=1019333/1024 and ell\_Q=287843/4096 retain their exact directed meanings. This is a whole-sector count; no additional 2\^{}334 pattern factor or type factor is needed in TD10.

It remains to treat q≠0 in TD8.

\subsection{3. The exact small-prime divisor bound of exponent one twelfth}\label{the-exact-small-prime-divisor-bound-of-exponent-one-twelfth}

The adopted field is unchanged:

\[
 E=\mathbf Q(\zeta_{28})B_7(29)\mathbf Q(\sqrt5,\sqrt{13}),
 \quad [E:\mathbf Q]=336,\quad F=E^+,\quad [F:\mathbf Q]=168.
                                                                    \tag{TD11}
\]

Its direct Galois product and disjoint ramification are the accepted K032/K075/K080 field data. The subgroup defining B7(29) is \{1,12,17,28\} modulo 29. The complete list of prime-ideal norms below 4096 is

\begin{longtable}[]{@{}rrr@{}}
\toprule\noalign{}
Norm & Number in E & Number in F \\
\midrule\noalign{}
\endhead
\bottomrule\noalign{}
\endlastfoot
29 & 48 & 24 \\
1681=41² & 168 & 84 \\
2129 & 336 & 168 \\
\end{longtable}

Here is a finite verification of completeness. At 2 and 3, the residue degree in the two real subfields B7(29) and Q(zeta7)+ is divisible by 7 and 3, respectively, so the composite residue degree is divisible by 21. At any p≥5 outside the indicated B-splitting subgroup, it is at least 7, already giving p\^{}7\textgreater4096; this also treats the other ramified primes 5,7,13. At 29 the carried ramification index 7 and complete splitting in the other factors give the first row. For all remaining p, the residue degrees are the least common multiples of the exact order of p modulo 28 (modulo ±1 for F) and the two quadratic residue degrees for 5 and 13.

For primes below 64 in the B-splitting subgroup, the only possibilities other than 29 are 17,41,59. The exact degrees exclude 17 and 59, and give degree 2 at 41. For primes at least 64, a norm below 4096 requires degree 1. The complete primes below 4096 satisfying the necessary B and real-cyclotomic splitting congruences are

\[
307,811,1259,1567,1931,2129,2437,2801,3191,3613,4003.
\]

The two quadratic residue tests leave only 2129, which is 1 modulo 28 and splits in E as well. The certificate enumerates every rational prime below 4096, verifies these congruences and both quotient orders, and records the complete degree list. Thus neither an unverified prime-norm floor nor a discriminant heuristic is used.

For q≥4096 and every e≥0, e+1≤2\^{}e≤q\^{}(e/12). At q=29 the factor (e+1)q\^{}(−e/12) is maximal at e=3; at q=1681 and 2129 its maximum is at e=1. These facts follow by consecutive ratios, and the requisite integer power inequalities are checked exactly. Unique ideal factorization therefore gives

\[
\begin{aligned}
 d_E(\mathfrak a)&\le C_E N_E(\mathfrak a)^{1/12},
 &C_E&=\frac{2^{600}}{29^{12}41^{28}2129^{28}},\\
 d_F(\mathfrak c)&\le C_F N_F(\mathfrak c)^{1/12},
 &C_F&=\frac{2^{300}}{29^6 41^{14}2129^{14}},\\
 C_{12}=C_EC_F&=\frac{2^{900}}{29^{18}41^{42}2129^{42}}.
\end{aligned}                                                     \tag{TD12}
\]

These constants allow all exceptional primes even when absent from the particular ideal. They do not require principality or change the coefficient field.

\subsection{4. Retain the actual affine content and gain a nonemptiness constraint}\label{retain-the-actual-affine-content-and-gain-a-nonemptiness-constraint}

Fix the s=23 active T places, the intrinsic marked subset D of actual double poles of eta, d=\textbar D\textbar, all actual labelled pairs, and the paired rank-zero G pattern. K077(AP6--11) gives

\[
 \mathfrak l=\mathfrak J I_\Sigma I_D,
 \quad N_E\mathfrak l=z^{2f},
 \quad z=(N_E\mathfrak J)^{1/(2f)}T^{-(s+d)/f},
 \quad N_E\mathfrak J\le M^2,\quad M=Q^{167}.                     \tag{TD13}
\]

The full entry ideal of N is contained in l; the exact content identity gives

\[
 (\eta)\subset\mathfrak b\subset\mathfrak l,
 \qquad N_E\mathfrak b\ge|N_E\eta|^{1/2}z^f.                     \tag{TD14}
\]

Since q≠0 belongs to b and \textbar q\_tau\textbar≤\textbar eta\_tau\textbar/sqrt(V),

\[
 N_E\mathfrak b\le|N_E q|\le |N_E\eta|/V^f,
 \qquad z\le R/V,\quad y:=R/z\ge V.                             \tag{TD15}
\]

This is an actual nonemptiness restriction; q=0 was treated separately.

For fixed eta and b, K077(AP7--11) places \textbf{every} original positive coordinate, including r=x4, in its prescribed affine coset of

\[
 \mathfrak c=\mathfrak b\cap\mathfrak J I_\Sigma,
 \qquad N_E\mathfrak c=N_E\mathfrak b\,T^{2d}.                   \tag{TD16}
\]

In the original metric sum\_tau \textbar w\_tau\textbar²/\textbar eta\_tau\textbar, nonzero differences in c have squared length at least fzT\^{}(2d/f). The single coordinate r has squared radius at most fR. Disjoint-ball packing in real dimension 2f thus gives

\[
                  \#r\le(1+2\sqrt{y_c})^{2f},
            \qquad y_c=R/(zT^{2d/f}).                            \tag{TD17}
\]

No residue-density approximation or divided moving basis appears here.

The content-ideal injection K075(RN11--13) places b among divisors of an integral E-ideal of norm at most y\^{}(2f). TD12 bounds its possibilities by C\_E y\^{}28. For fixed eta,b,r, a nonzero q satisfies

\[
 q\bar q=(\eta\bar\eta-r\bar r)/V=:\beta.
\]

The ideal (q)b\^{}(−1) is integral and has relative F-norm

\[
 \mathfrak c_q=(\beta)(N_{E/F}\mathfrak b)^{-1},
 \qquad N_F\mathfrak c_q\le(y/V)^f.                              \tag{TD18}
\]

Any actual q makes this an integral ideal; no divisibility by the rational integer V is asserted before imposing existence. The local relative-norm ideal count of K075(RN17--19) gives at most d\_F(c\_q) ideals, and at most \textbar mu(E)\textbar=28 generators for each. TD12 therefore bounds q by 28C\_F(y/V)\^{}14. The zero beta case has no q≠0 and need not be added. Combining these factors, apart from eta, gives

\[
       28C_{12}\,y^{42}V^{-14}(1+2\sqrt{y_c})^{2f}.              \tag{TD19}
\]

\subsection{5. The pure-imaginary lattice count uses the real covolume}\label{the-pure-imaginary-lattice-count-uses-the-real-covolume}

All primes in J,I\_Sigma,I\_D are paired split primes, so l=Lambda\_F O\_E for an F fractional ideal with N\_F Lambda\_F=z\^{}f.~Because i is an integral unit and E contains i,

\[
 \mathcal O_E\cap iF=i\mathcal O_F,
 \qquad \mathfrak l\cap iF=i\Lambda_F.                            \tag{TD20}
\]

Contraction of an extended ideal gives the second equality; no new principal generator is chosen. In the metric sum\_tau \textbar eta\_tau\textbar², its covolume is exactly sqrt(D\_F)z\^{}f.~The original-ring metric addendum (RM1--2) gives the exact real discriminant and covolume. The real-ideal codifferent estimate, combined with the accepted covering transference in LIFT47(TF7), gives

\[
 \mu(\Lambda_F)<750\sqrt f\,z<9750z,
 \qquad D_F=2^{168}D_0,\quad \sqrt{D_F}>2^{883},
 \quad D_0=7^{140}29^{144}5^{84}13^{84}.                          \tag{TD21}
\]

For clarity, the real dual is Diff\_F\^{}(−1)Lambda\_F\^{}(−1), with shortest length at least sqrt(f)/(rd(F)z), and the transference bound is mu·lambda1(dual)≤f/2. Thus rd(F)\textless1500 yields exactly the coefficient above. This is the full real discriminant normalization, not the half-trace complex one.

Translated Voronoi cells about the points in the box \textbar eta\_tau\textbar≤R fit in the coordinate box of radius R+9750z. Their disjoint volumes imply

\[
 \#\eta\le\frac{2^f}{\sqrt{D_F}}(y+9750)^f
             <2^{-714}y^f.                                     \tag{TD22}
\]

The last inequality uses the actual nonemptiness y≥V and the exact integer inequality (V+9750)\^{}168\textless2V\^{}168. Zero eta can be included in the upper count but is excluded from TD2 and every subsequent division. TD22 uses a full fundamental-volume argument; it is not a density-only estimate for a very thin lattice.

\subsection{6. Monotone G cap, exact selected trace and the full marked pair sum}\label{monotone-g-cap-exact-selected-trace-and-the-full-marked-pair-sum}

DO088's unchanged actual near-point calculation and DO080's genuine normalized selected trace give

\[
             |c(h)|<2\,2^{-931f/2}z^{-f}.                        \tag{TD23}
\]

This follows from the multiplicative rank-one identity for eta after conjugating the actual near-origin point, including arbitrarily small conjugates; it does not follow from a small additive displacement. The product of all near-point losses is less than 2. All compact factors and the selected metric remain those of DO080.

The full finite Green bound for a fixed paired rank-zero pattern is 2(N\_E J)\^{}(a/2)M\^{}(−a). Multiplying it by TD19, TD22 and TD23, its z-dependent part is proportional to

\[
 z^{\kappa}(\sqrt z+2\sqrt R\,T^{-d/f})^{2f},
        \qquad \kappa=(a-3)f-42=27678>0.                         \tag{TD24}
\]

It is increasing for all z\textgreater0. Therefore the exact cap N\_E J≤M² permits

\[
 z_{s,d}=(M/T^{s+d})^{1/f},\quad y_{s,d}=R/z_{s,d},
 \quad y'_{s,d}=R(T^{s-d}/M)^{1/f}.                              \tag{TD25}
\]

The cap need not itself satisfy y≥V: TD24 is a monotone extension of an upper expression proved on the nonempty set, not an assertion that the cap is geometrically realized. At the cap the finite numerator cancels M\^{}a exactly.

For simplicity retain all 2\^{}334 paired G patterns; no weighted-pattern saving is invoked. K077(AP18) retains the exact full type and pair sum: a marked double-pole place has cost 1, while the three alternatives at every unmarked place sum to T+1. Thus for each d=0,\ldots,23 the complete upper expression is

\[
\begin{split}
 U_{23,d}<&\;2^{-378}\,28C_{12}\binom{168}{23}\binom{23}{d}
 2^{-931f/2}(T+1)^{23-d}z_{23,d}^{-f}\\
 &\qquad\cdot y_{23,d}^{210}V^{-14}
             (1+2\sqrt{y'_{23,d}})^{2f}.
\end{split}                                                       \tag{TD26}
\]

The constant −378 is 1+334+1−714: the finite Euler factor, every unweighted pattern, the actual near-point product, and the real-field eta count. The exponent 210 is f+42. The exact binomial for each integer d remains, so there is no additional 2\^{}23 endpoint factor. Every actual summand has one intrinsic D; its labelled-pair conditions are retained before this uniform count is summed.

\subsection{7. Exact arithmetic and the new absolute charge}\label{exact-arithmetic-and-the-new-absolute-charge}

The companion evaluates TD10 and all 24 expressions TD26 using integer root brackets and directed rational logarithms. It reads the actual T,g0,g1 from the unchanged host certificate. It separately checks the full small-prime list, all local divisor maxima, the exact content exponents, the two original transverse linear identities, the real discriminant and eta-volume inequalities, every marked type sum, and the complete chronological error arrays.

The exact witness exponent TD10 is

\[
                     -106340445/256<-400000.                    \tag{TD27}
\]

For every d, the complete strict log2 upper bound in TD26 is less than the negative of its displayed margin:

\begin{longtable}[]{@{}rr@{}}
\toprule\noalign{}
Marked double-pole count d & Certified inverse binary margin \\
\midrule\noalign{}
\endhead
\bottomrule\noalign{}
\endlastfoot
0 & 17651 \\
1 & 17398 \\
2 & 17145 \\
3 & 16894 \\
4 & 16642 \\
5 & 16392 \\
6 & 16141 \\
7 & 15891 \\
8 & 15641 \\
9 & 15391 \\
10 & 15140 \\
11 & 14874 \\
12 & 14513 \\
13 & 13800 \\
14 & 12699 \\
15 & 11477 \\
16 & 10236 \\
17 & 8994 \\
18 & 7751 \\
19 & 6509 \\
20 & 5267 \\
21 & 4025 \\
22 & 2784 \\
23 & 1545 \\
\end{longtable}

The worst cell is d=23, of margin 1545; every other cell has margin at least 2784. Hence the entire reduced case together with the single witness sector has bound

\[
 2^{-1545}+23\,2^{-2784}+2^{-400000}<2^{-1544}.                  \tag{TD28}
\]

This includes every marked subset, actual labelled pair, all 334 paired rank-zero decisions, all ideal and relative norm choices, and the actual near-point selected trace. The certificate also sums all 24 separate directed margins before adding the witness bound.

In the current chronology the canonical R068 predecessor has 70 all-top and 31 uniform errors. Append 1544 once, preserving every earlier term, to obtain 71/31 entries. Both complete exact dyadic totals remain below 2\^{}−230. The complete all-top support floor remains 23 and the uniform floor remains 11 until other sectors are exhausted separately. Each removal is intersected with the current remainder; no earlier error term is erased.

The certificate is arithmetic evidence; the preceding proof supplies the original geometric, integral and representation-theoretic interfaces. No certificate enumeration is represented as an enumeration of all rational matrices.

\subsection{8. Source scope and unresolved gates}\label{source-scope-and-unresolved-gates}

The exact normal form, content identity and relative norm count are K075(RN5--19); the actual paired negative unit, all-positive-coordinate congruences, difference ideal and complete marked type sum are K077(AP1--18). The field/root-unit and old divisor inputs are K032, K050 and reconstructed K080(RC1--4); TD11--12 are the new exact finite refinement. The common original transverse vectors and arbitrary-arrow integral pairing are DO078(FP20--22) and K088(IP5--6). The genuine selected trace and multiplicative same-point transport are DO080 and DO088. The exact real/complex discriminant distinction is LIFT063's original-ring primary metric addendum (RM1--2). The real dual calculation above uses the accepted LIFT47(TF7) transference theorem:

Banaszczyk, W. (1993). New bounds in some transference theorems in the geometry of numbers. \emph{Mathematische Annalen, 296}, 625--635, Theorem 2.2.

This proof changes neither the adopted field, original pencil, finite characters, class frames, CM phases, selected Hilbert normalization nor source arithmetic modules. It asserts no integral inverse on a selected-image map. The signed sufficient inequality, prescribed global Bessel/Mram period, literal fully admissible host constructor, selected-image content, actual cycle value and DP44 comparison remain open. The complete support-floor implication requires a separate exhaustion of all other support-23 classes.

\clearpage

\section{AM-R-069 --- minimal finite rank and the original exceptional cross-vector ideal}\label{am-r-069-minimal-finite-rank-and-the-original-exceptional-cross-vector-ideal}

\textbf{Canonical checked integration --- 2026-11-22T20:52:51.614308+00:00.} Root and independent whole-proof/script/interface reviews and separate726-assertion replays passed. The entire original IC8 positive-compression rank1/2 class at supports23..28 of the unchanged actual all-top H133 kernel has absolute contribution\textless2\^{}-23546. The symmetric partial-monomial compression preserves rank at every G factor. Finite rank2 then forces every untouched original cross column to vanish at both CM partners, supplying a paired ideal in the unchanged joint160-coordinate ball count. All original columns, fibres, target and pair weights are retained. No global unipotence is assumed. The complete all-top floor remains23 and uniform floor11; identity-compression and higher-support classes, signed positivity, actual source/arithmetic matching and DP44 remain open.

Author: \nolinkurl{/root}, GPT-6 Astra Pro; inherited effort not externally exposed; ALP\_MACHINE.22 November2026 UTC. New controlled draft. Whole proof/script and independent checks are separate prerequisites for promotion.

\subsection{1. Exact class and conclusion}\label{exact-class-and-conclusion}

Keep the actual J56 all-top H=133 kernel, f=a=168,n=169,m=164, paired CM pencils, finite projectors, original packet, selected Hilbert metric, class masses and phases. Consider the entire K069 IC8 class with r=rank\_E(A-I) in\{1,2\},A=h\_SS, at

\[
                         23\le s\le28.\tag{EM1}
\]

All334 unmixed G reductions remain unipotent. No global unipotence is assumed. We prove the complete containing-set bound

\[
                  \sum_{\mathrm{EM1}}|w(h)|<2^{-23546}.\tag{EM2}
\]

The exact R063/R066 zero/transverse/plane reconstruction is retained. A padded four-subset contains all changed original S vertices, and its160-coordinate complement is S0. Put C0=h\_(W0,S0). It has rank zero or one; the original data C0,v\_B,v\_W determine the full negative column and retain its full source fibre. The three positive integral transverse pairings and all type/target counts are unchanged. The only new counting input is a paired G-ideal for C0 in the finite-rank-two case.

\subsection{2. The symmetric compression has no rank loss at G}\label{the-symmetric-compression-has-no-rank-loss-at-g}

Let

\[
        N=h-I,\qquad H_h=h+h^{-1}-2I=N^2h^{-1}.
        \tag{EM3}
\]

At the positive original S block, unitary inversion gives

\[
                   (H_h)_{SS}=A+A^\dagger-2I.\tag{EM4}
\]

The partial-monomial classification and mu28 entries are exactly DO064(PM5--9). We need the following strengthening of its finite rank statement:

\[
       \operatorname{rank}_{k_v}(\overline A+\overline A^\dagger-2I)=r
                \quad\text{at every actual G component}.\tag{EM5}
\]

\textbf{Proof.} Decompose the weighted partial-permutation graph into its paths and cycles, using a simultaneous permutation of original coordinates only. On a path with l vertices, diagonal conjugation makes A the unweighted directed shift and A\^{}dagger its reverse shift: the two edge weights are inverses because the original weights are roots of unity. Thus A+A\^{}dagger-2I is tridiagonal with diagonal -2 and off-diagonal1. The determinant recurrence gives (-1)\^{}l(l+1). It is a unit at either actual G prime, both greater than165, while l\textless=164. Its rank is l, the rank of A-I on that path.

On a cycle of length l and product weight zeta, A is invertible, A\^{}dagger=A\^{}-1 and A\^{}l=zeta I. The polynomial X\^{}l-zeta is separable in the residue field because its characteristic exceeds164. Therefore A is semisimple after algebraic closure, and

\[
           A+A^{-1}-2I=(A-I)^2A^{-1}
\]

has the same rank as A-I. The latter rank is l-1 if zeta=1 and l otherwise. Reduction of mu28 is injective at these primes, so that condition is the same before and after reduction. Summing the component ranks proves EM5. No basis used here is inserted into a lattice count. The same calculation applies at the other CM component with inverse-conjugate weights. ∎

Write d\_v=rank(Nbar). Since Nbar is nilpotent and nonzero, its Jordan blocks give rank(Nbar\^{}2)\textless=d\_v-1. Equations EM3--5 imply

\[
                             d_v\ge r+1.\tag{EM6}
\]

In particular the rank-two compressed sector has d\_v\textgreater=3 at every unmixed G factor. The rank-one sector can have d\_v=2; it is the only case needing a new congruence.

\subsection{3. At full finite rank two every untouched column is fixed}\label{at-full-finite-rank-two-every-untouched-column-is-fixed}

Suppose r=1 and d\_v=2. EM5 makes Hbar nonzero. Nilpotence and rank two then force the unique nontrivial Jordan block to have size3; all other blocks have size1. Thus

\[
 Nbar^3=0,\quad Hbar=Nbar^2,\quad
 \operatorname{rank}Hbar=1,\quad
 \operatorname{im}Nbar\cap\ker Nbar=\operatorname{im}Nbar^2.
 \tag{EM7}
\]

Projection P\_S is injective on im(Hbar), since both Hbar and its compressed S block have rank1. For an untouched original e\_j, one has (Hbar)\_SS e\_j=0, so Hbar e\_j=0. Consequently Nbar\^{}2 e\_j=0, and Nbar e\_j belongs to im(Nbar) intersect ker(Nbar)=im(Hbar). But P\_S Nbar e\_j=(Abar-I)e\_j=0. Injectivity gives

\[
                            Nbar e_j=0\qquad(j\in S0).\tag{EM8}
\]

This argument is entirely in the actual residue field. At the conjugate CM component the full rank is also2, because that matrix is the inverse adjoint. Its compressed symmetric block also has rank1. Apply the same argument there, independently. Thus every original C0 entry is divisible at both members of every rank-two G pair.

Let J be the paired E-ideal supported on exactly those unmixed relative G factors with d\_v=2. Then

\[
  1\le N_EJ\le Q^{2(f-1)},\qquad
  \#\{\text{possible }J\}\le2^{334},\qquad Q=g_0g_1.
  \tag{EM9}
\]

For r=2, J is necessarily the unit ideal. Both ranks are covered by the same containing pattern count, without another factor2. The original padded-block choice already covers both ranks.

The complete Green bound uses L2=335,L3=501. At a factor outside J the full rank is at least3; inside J it is2. Retaining the single global Euler correction and both mixed norm bounds gives

\[
             |a(h)|\le2Q^{-501(f-1)}(N_EJ)^{83}.
             \tag{EM10}
\]

This is the literal rank jump166, not an archimedean minor estimate or a new finite orientation count.

\subsection{4. The original joint ball with its new paired ideal}\label{the-original-joint-ball-with-its-new-paired-ideal}

The zero and transverse alternatives need no change and retain R066(JB11). In the plane case,

\[
       C0e_j=(u\xi_j,\beta_j),\qquad
       \xi\in (J I_\Sigma)^{160},\qquad N_E I_\Sigma=T^{-2s}.
       \tag{EM11}
\]

Indeed the old T simple-pole assertion holds in the original coordinates, and EM8 supplies the G divisibility of every positive entry. The unchanged integer Bezout relation for u recovers xi in this same ideal. The G and T supports are disjoint. The original five-coordinate pivot count is retained in full, so no additional claim about dividing by xi or about a moving isotropic-line lattice is needed.

Set

\[
 z=(N_EJ)^{1/(2f)}T^{-s/f},\quad
 z_0=T^{-s/f},\quad z_1=Q^{(f-1)/f}T^{-s/f},\quad M=2^{-66}.
 \tag{EM12}
\]

The rank-one Hilbert--Schmidt argument R066(JB3) places all160 original xi coordinates in one real53760-dimensional ball of radius M sqrt(f). Its ideal-product covolume is D0\^{}160 z\^{}(320f), with the unchanged D0=7\^{}140 29\^{}144 5\^{}84 13\^{}84. The exact half-trace dual and arbitrary-ideal transference proof R066(JB4--7) works for J I\_Sigma just as for I\_Sigma, and gives

\[
 \#\xi\le A_{160}(M/z+9500)^{320f},\qquad
 A_{160}=\frac{(528)^{160f}}{(160f)!D0^{160}}<2^{-369387}.
 \tag{EM13}
\]

The covering radius is9500 sqrt(f) z. This uses the arbitrary fractional-ideal dual, not the smaller ring-only covering radius. All zero coordinates are allowed in the upper count.

Since (N\_EJ)\^{}83=T\^{}(166s) z\^{}(166f), the z-dependent part of EM10 times EM13 is, apart from positive fixed factors,

\[
                 z^{-154f}(M+9500z)^{320f}.\tag{EM14}
\]

Its logarithm is convex as a function of log z: it is an affine function plus320f times log(M+9500 exp(v)). Hence its maximum on{[}z0,z1{]} is at an endpoint. No endpoint-union factor is incurred, and no ideal with endpoint norm is asserted to exist.

At z0, the original exact JB8 inequalities at support22 imply, for every s\textgreater=23,

\[
 \#\xi<2^{1-369387-66(320f)}T^{320s}.\tag{EM15}
\]

At z1, throughout23\textless=s\textless=28 the exact integer comparison

\[
                       T^{28}<Q^{167}2^{97f}\tag{EM16}
\]

gives M/z1\textless2\^{}31. Thus M/z1+9500\textless2\^{}32 and

\[
                         \#\xi<2^{-369387+32(320f)}.\tag{EM17}
\]

At this endpoint EM10 becomes the original denominator2Q\^{}(-335(f-1)). These two estimates retain the Green factor together with its corresponding lattice density.

\subsection{5. All factors, all supports and exact charge}\label{all-factors-all-supports-and-exact-charge}

Use every original reconstruction factor of R066(JB10):

\[
 C=1+168+25+8+2+5022+2(14861)+10835=45783.
 \tag{EM18}
\]

These are respectively the single Green Euler correction, supports, padded four-subsets, pivots, three transverse choices, full original source fibre, two full W0 vector counts, and four original positive-column coordinates. The complete ordered-pair/absolute-incidence factor is below2\^{}(997s), and the common vector counts have T-power28s. Every type and zero target is included. There is no extra count of A, compact type or CM orientation.

Let LT=1019333/1024 and LQ=287843/4096 be the unchanged directed logarithmic bounds. The zero and transverse exponents are

\[
 E_0(s)=C+(28LT+997)s-335\cdot167LQ,\qquad
 E_A(s)=E_0(s)+45(320f).
 \tag{EM19}
\]

The complete plane endpoint exponents, including all2\^{}334 possible J, are

\[
\begin{split}
 E_{B,0}(s)&=C+334+1-369387-66(320f)
                  +(348LT+997)s-501\cdot167LQ,\\
 E_{B,1}(s)&=C+334-369387+32(320f)
                  +(28LT+997)s-335\cdot167LQ.
\end{split}\tag{EM20}
\]

Each is increasing in s. Exact rational evaluation verifies that all four are strictly below-23551 for23\textless=s\textless=28. The largest is

\[
                 E_{B,0}(28)=-96466521/4096<-23551.
                 \tag{EM21}
\]

For each support the plane case uses the maximum of its two endpoints. There are three exhaustive cases at six supports. Consequently the complete absolute union is below18*2\^{}-23551\textless2\^{}-23546, proving EM2. The failure of the z0 envelope at support29 is only failure of this sufficient estimate.

\subsection{6. Accounting, geometry and remaining gates}\label{accounting-geometry-and-remaining-gates}

Append23546 exactly once to the actual all-top error list after all intervening checked findings, retain every old charge, and leave the uniform list unchanged. The companion must verify the entire current predecessor and both complete dyadic sums before canonical integration. This theorem alone leaves the complete support floor23: the identity-compression and other classes need their own complete exhaustion.

The new input is the symmetric compression on the same original Hermitian module and the resulting paired residue-field kernel statement. Paths/cycles prove a determinantal rank identity, and the count stays in the original product of fractional ideals. No integral moving basis, new CM field, source normalization, labelled correspondence, or selected projector is introduced. Exact carried contextual anchors are SGA3 I4.7.3/4.7.3.1; SGA6 VI Theorem3.3; FGA232 Theorem3.1 pp148--149; FGA236 Corollaries1.7/2.3 pp229/232; SGA7 I VII2.7/2.9 pp169--174 and VIII7.4(b) pp310--311; cote4 scan003 §§8.1.1--8.1.2, scans005--006 §§8.4.1--8.4.3 and cote54 scans005--007. No new lambda-ring/K-theory identity, l-adic sheaf/weight theorem, CM motive comparison or motivic L-function assertion follows.

DP\_MULT Q3e and THETA3 G1/A7 gain precisely this smaller containing remainder for the unchanged kernel. The signed sufficient inequality, actual fully admissible mixed-CM source, prescribed global Bessel/Mram value, selected derivative/class content, actual Read/STiesR/specific032 transport and DP44 cost difference remain open. CX\_GRH CX-ALG-024 and GCT\_PLETH T-052 retain their separate arithmetic requirements. The source Green and lattice transference inputs are carried through the already checked R063/R066/DO064/DO085 interfaces. New finite diagnostics do not enumerate the actual rational kernel or substitute for this whole proof.

\clearpage

\section{AM-DO-092 --- split content ideals and the exact radial profile at rank-one support twenty-four}\label{am-do-092-split-content-ideals-and-the-exact-radial-profile-at-rank-one-support-twenty-four}

\textbf{Canonical checked integration --- 2026-11-22T21:05:00.141363+00:00.} Root and independent whole-proof/script/interface checks and separate2531-assertion replays passed. The entire actual all-top H133 globally rank-one-unipotent fixed-W0 class at support24 has absolute contribution\textless2\^{}-516. The exact conjugate-content constraint and split-only relative norm count give constants\textless2\^{}361 and\textless2\^{}244. Integrating the retained actual selected radial density gives eta\_a=4\^{}(a-1)/{[}2(a-1)binom(2a-2,a-1){]}, with a full630-bit improvement at a=f168. The weighted G patterns, real eta metric, original affine ideal and all25 marked/type/pair cells are retained. Global unipotence and rank one are premises. Complete floors remain23/11; signed/global/arithmetic/source/DP44 gates remain open.

\textbf{Algebraic-geometric translation.} This refines the count of original content ideals C under C bar(C) subset A and of actual relative norm ideals, without assuming principality. The complex realization retains its Hermitian bundle metric and selected trace. The original modules, finite correspondences, Picard classes, polarizations and biextensions are unchanged. Exact carried contextual anchors are SGA3 I4.7.3/4.7.3.1; SGA6 VI Theorem3.3; FGA232 Theorem3.1 pp148--149; FGA236 Corollaries1.7/2.3 pp229/232; SGA7 I VII2.7/2.9 pp169--174 and VIII7.4(b) pp310--311; cote4 scan003 §§8.1.1--8.1.2, scans005--006 §§8.4.1--8.4.3 and cote54 scans005--007. No new lambda-ring/K-theory identity, l-adic sheaf/weight theorem, CM motive or motivic L-function assertion follows. DP\_MULT Q3e and THETA3 G1/A7 gain precisely this containing-set removal; CX\_GRH CX-ALG-024 and GCT\_PLETH T-052 retain their separate arithmetic gates.

Author: \nolinkurl{/root/definite_oscillator}, ALP\_MACHINE. 22 November 2026 UTC. Complete proof; saved-file checks and the exact certificate are recorded separately before promotion.

\subsection{1. The same original rank-one sector}\label{the-same-original-rank-one-sector}

Keep the complete original J56 all-top data and H=133 cutoff. The sector is the globally unipotent, globally rank-one fixed-W0 sector at support s=24. All actual labelled paired pencils, finite G projectors, class masses, original packet, CM phases and the selected Hilbert metric remain unchanged. This proof adds three genuine refinements to DO091: a paired constraint on the coordinate-content ideal, the split-prime-only relative norm count, and an exact integration of the retained radial density in DO080. The entire sector has absolute contribution strictly less than 2\^{}−516, by the complete all-profile calculation below. This is not by itself a full structural support-floor statement.

Use f=a=168, Q=g0g1, M=Q\^{}167 and R=8844/133. The exact original normal form and two-transverse-pairing reduction of DO091(TD2--9) apply without a restriction on support:

\[
 N=x\bar x^{\mathsf t}J_0/\eta,
 \quad x=(v q,r,\eta),\quad
 Vq\bar q+r\bar r=\eta\bar\eta,
 \quad\bar\eta=-\eta\ne0,
 \quad |\eta_\tau|\le R,                                        \tag{SN1}
\]

where v=(2,95192390527,383092671043), V=v dot v, and the coordinate-content ideal is b=(q,r,eta). The two short integral transverse pairings give a separate uniform field-norm witness unless both vanish. If both vanish but q=0, the third integral pairing is 4N44 and gives the same witness. The complete witness-sector bound is the support-variable TD10, evaluated at s=24:

\[
 E_{\rm wit}(24)=1+168+107+5757+46af
       +(10\ell_T+997)24-a167\ell_Q<-400000.                     \tag{SN2}
\]

Its actual full five-column/fibre count includes every incidence type and zero target. The new reduced count treats q≠0 only.

For a prescribed support, marked double-pole subset D of size d, actual local pairs, and paired rank-zero G ideal J, retain exactly K077 and DO091:

\[
 \mathfrak l=J I_\Sigma I_D=\Lambda_F\mathcal O_E,
 \quad N_E\mathfrak l=z^{2f},
 \quad z=(N_EJ)^{1/(2f)}T^{-(s+d)/f},\quad N_EJ\le M^2,
 \quad y=R/z\ge V.                                             \tag{SN3}
\]

The last nonemptiness condition uses q≠0. It is not imposed on an artificially chosen cap.

\subsection{2. The conjugate content constraint is a real-ideal divisor problem}\label{the-conjugate-content-constraint-is-a-real-ideal-divisor-problem}

Let a\_N be the ideal of the entries of N. The literal normal form gives b bar(b)=(eta)a\_N and a\_N⊂l. Introduce the two integral ideals

\[
 \mathfrak A=(\eta)\mathfrak l^{-1}
       =\mathfrak A_F\mathcal O_E,
 \quad\mathfrak A_F=(-i\eta)\Lambda_F^{-1},
 \qquad \mathfrak C=\mathfrak b\mathfrak l^{-1}.                 \tag{SN4}
\]

Here −i eta lies in F, i is an integral unit, and contraction shows that A\_F is integral. Its norm satisfies N\_F A\_F≤y\^{}f.~The original containments and the complete entry ideal imply simultaneously

\[
             \mathfrak A\subset\mathfrak C\subset\mathcal O_E,
               \qquad \mathfrak C\bar{\mathfrak C}\subset\mathfrak A.
                                                                    \tag{SN5}
\]

This is stronger than simply counting all E-ideal divisors of A. It does not assume that b or C is principal.

At a split F-prime p with exponent e in A\_F, write the two exponents of C as u,v. They satisfy 0≤u,v≤e and u+v≥e. Their exact number is

\[
                         T_e=(e+1)(e+2)/2.                       \tag{SN6}
\]

At an inert prime, C has one exponent u with ceil(e/2)≤u≤e, giving floor(e/2)+1 possibilities. At a ramified prime, extension doubles e and the same argument gives e+1 possibilities. These are local upper counts for the actual content ideals; further matrix conditions can only reduce them.

Every odd prime of E/F is unramified, since E=F(i), and an odd F-prime of norm q splits exactly when q≡1 modulo 4. The DO091 complete norm list below 4096 consists only of 29,1681,2129, all split. Thus every inert prime has q≥4096 and

\[
 \lfloor e/2\rfloor+1\le2^{e/2}\le q^{e/10}.
\]

Any ramified prime lies over 2. Its norm is at least 2\^{}21, by the two real residue-degree subfields already checked in DO091, and e+1≤2\^{}e≤q\^{}(e/10). Therefore only split primes can contribute an exceptional constant to the content count with exponent 1/10.

For a split prime q≥3\^{}10=59049, T\_e≤3\^{}e≤q\^{}(e/10). Define the exact finite comparison constant

\[
 C_{\rm cont}=\prod_{\substack{\mathfrak p\text{ splits in }E\\N\mathfrak p<59049}}
        \max_{e\ge0}\left(T_e(N\mathfrak p)^{-e/10}\right).
                                                                    \tag{SN7}
\]

The exact prime-degree and integer-power calculation in section 4 proves C\_cont\textless2\^{}361. Multiplying the local bounds in SN6 therefore gives

\[
                  \#\{\mathfrak b\mid\eta\}
                       \le C_{\rm cont}y^{f/10}<2^{361}y^{f/10}.
                                                                    \tag{SN8}
\]

No class-group or generator multiplicity is being omitted: ideals are counted before any relative norm solution, and every actual b yields a distinct C.

\subsection{3. Only split primes contribute to the remaining relative norm multiplicity}\label{only-split-primes-contribute-to-the-remaining-relative-norm-multiplicity}

For fixed eta,b, K077(AP7--11) puts the original fourth coordinate r in its prescribed affine coset of c=b∩J I\_Sigma, with N\_E c=N\_E b·T\^{}(2d). The unchanged metric difference minimum gives

\[
                      \#r\le(1+2\sqrt{y_c})^{2f},
                       \qquad y_c=R/(zT^{2d/f}).                  \tag{SN9}
\]

For each such r, q≠0 satisfies q bar(q)=(eta bar(eta)−r bar(r))/V=beta. The integral E-ideal (q)b\^{}(−1) has relative F-norm c\_q=(beta)(N\_(E/F)b)\^{}(−1), of norm at most (y/V)\^{}f, exactly as in DO091(TD18). At a split F-prime with exponent e in c\_q there are e+1 relative-norm ideal choices. At inert or ramified primes there is at most one. Thus the old bound by the entire F-divisor function may be replaced by its split-prime product without changing any ideal or actual norm equation.

For split primes q≥2\^{}16=65536, e+1≤2\^{}e≤q\^{}(e/16). Define

\[
 C_{\rm rel}=\prod_{\substack{\mathfrak p\text{ splits in }E\\N\mathfrak p<65536}}
           \max_{e\ge0}\left((e+1)(N\mathfrak p)^{-e/16}\right).
                                                                    \tag{SN10}
\]

The exact calculation below proves C\_rel\textless2\^{}244. Each eligible relative norm ideal has at most the unchanged 28 relative-norm-one generators, as in K075(RN17--19). Consequently

\[
          \#q\le28C_{\rm rel}(y/V)^{f/16}
                 <28\,2^{244}(y/V)^{f/16}.                       \tag{SN11}
\]

The norm ideal is asserted integral only when an actual solution exists. The zero beta case has no q≠0. No new coefficient field or isotropic-line generator is chosen.

The eta lattice is exactly i Lambda\_F. Its real trace covolume and the full real codifferent/covering calculation in DO091(TD20--22) give, unchanged,

\[
                              \#\eta<2^{-714}y^f.                \tag{SN12}
\]

This remains valid on every nonempty reduced stratum because y≥V. Combining SN8--12 yields the entire reduced column count for each actual prescribed pattern and pencil data:

\[
 2^{-714}\,28\,2^{605}\,y^{f+273/10}V^{-21/2}
                          (1+2\sqrt{y_c})^{2f},                 \tag{SN13}
\]

where f/10+f/16=273/10 and f/16=21/2. A full column still determines the rank-one matrix; no stabilizer-fibre factor is added to SN13.

\subsection{4. Finite prime-degree calculation for both constants}\label{finite-prime-degree-calculation-for-both-constants}

Use the unchanged direct-product field of DO091(TD11). The same two real subfields give residue degree divisible by 21 at 2 and 3, so their prime norms exceed 65536. For p≥5 not split in B7(29), the degree is at least 7 and p\^{}7≥5\^{}7\textgreater65536; this also excludes the other ramified primes 5,7,13. At 29 there are exactly 24 F-primes of norm 29, and they split in E.

For each remaining rational prime p\textless65536, the exact F residue degree is

\[
 d_F(p)=\operatorname{lcm}\bigl(
 \operatorname{ord}_{(\mathbf Z/28)^*/\{\pm1\}}(p),
 d_5(p),d_{13}(p)\bigr),                                        \tag{SN14}
\]

provided p modulo 29 is in \{1,12,17,28\}; otherwise it was excluded above. Here d\_5,d\_13 are 1 or 2 according to the literal quadratic residue tests. Such an unramified prime has norm p\^{}(d\_F) and multiplicity 168/d\_F. It is retained in SN7 or SN10 exactly when its norm is 1 modulo 4 and lies below the corresponding threshold. The certificate enumerates every rational prime below 65536, records the complete resulting list, and verifies these tests without an analytic prime-density assumption.

The complete split-prime list and the maximizing exponents are:

\begin{longtable}[]{@{}rrrr@{}}
\toprule\noalign{}
Norm q & Multiplicity in F & Content exponent e & Relative-norm exponent e \\
\midrule\noalign{}
\endhead
\bottomrule\noalign{}
\endlastfoot
29 & 24 & 4 & 4 \\
1681 & 84 & 1 & 1 \\
2129 & 168 & 1 & 1 \\
7309 & 168 & 1 & 1 \\
13049 & 168 & 1 & 1 \\
18229 & 168 & 1 & 1 \\
25229 & 168 & 1 & 1 \\
33349 & 168 & 1 & 1 \\
35281 & 168 & 1 & 1 \\
39341 & 168 & 1 & 1 \\
42589 & 168 & 1 & 1 \\
43541 & 168 & 1 & 1 \\
55721 & 168 & 1 & 1 \\
60089 & 168 & Not in this product & 1 \\
60901 & 168 & Not in this product & 1 \\
62581 & 168 & Not in this product & 1 \\
\end{longtable}

All local maxima occur among e=0,1,2,3,4. For SN7 the consecutive ratio is ((e+3)/(e+1))q\^{}(−1/10), decreasing in e. From e=4 onward it is at most (7/5)29\^{}(−1/10)\textless1. For SN10 it is at most (6/5)29\^{}(−1/16)\textless1 from e=4 onward. The inequalities (7/5)\^{}10\textless29 and (6/5)\^{}16\textless29 are exact integer comparisons. Thus the complete constants can be checked by the rational products

\[
 C_{\rm cont}^{10}
   =\prod_q\left(\max_{0\le e\le4}\frac{T_e^{10}}{q^e}\right)^{m_q}
           <2^{3610},\qquad
 C_{\rm rel}^{16}
   =\prod_q\left(\max_{0\le e\le4}\frac{(e+1)^{16}}{q^e}\right)^{m_q}
           <2^{3904}.                                          \tag{SN15}
\]

Here the separate products use the respective thresholds and m\_q is the exact number of F-primes of norm q. For an explicit rational cross-check, let P be the eleven rational norms from 2129 through 55721 in the table and Pplus add 60089,60901,62581. Then

\[
 C_{\rm cont}^{10}=\frac{15^{240}3^{19320}}
 {29^{96}41^{168}\prod_{p\in P}p^{168}},\qquad
 C_{\rm rel}^{16}=\frac{5^{384}2^{38976}}
 {29^{96}41^{168}\prod_{p\in Pplus}p^{168}}.
\]

The two exact products have binary logarithms in {[}3600,3610) and {[}3888,3904), respectively; their strict upper endpoints give exactly the constants used in SN13. All constants are real comparison constants for a counting inequality; none is adjoined to the adopted arithmetic coefficient field.

\subsection{5. Exact integration of the actual selected radial profile}\label{exact-integration-of-the-actual-selected-radial-profile}

This step is a refinement of the genuine normalized-trace proof DO080, preserving its actual bundle metric, selected isometry, compact factors, word parameters and radial density. Write a≥2 for the positive rank, R\_y=sqrt(1−y), and use its exact conditional sphere density

\[
 f_y(x,t)=\frac{a-1}{\pi R_y^2}
       \left(1-\frac{x^2+t^2}{R_y^2}\right)^{a-2}
       {\bf1}_{x^2+t^2<R_y^2}.                                  \tag{SN16}
\]

For fixed t the same radial map and positive Jacobian estimate of DO080(RT10--11) give \textbar dx/dy'\textbar≤sqrt(y)/(2 sinh(r)y'\^{}2). In the density retain the bound

\[
 \left(1-\frac{x^2+t^2}{R_y^2}\right)^{a-2}
      \le\left(1-\frac{t^2}{R_y^2}\right)^{a-2}.
\]

Extending the admissible t interval to {[}−R\_y,R\_y{]} and integrating this factor improves RT12 by the exact multiplier

\[
 \eta_a=\frac12\int_{-1}^1(1-t^2)^{a-2}\,dt
       =\frac{4^{a-1}}{2(a-1)\binom{2a-2}{a-1}}.                \tag{SN17}
\]

The elementary integral follows from I\_0=2 and I\_j=2j I\_(j−1)/(2j+1). No angular approximation or representation-theoretic replacement is used. The multiplier is independent of y,y' and of the compact factors. Detailed balance and the exact variance argument DO080(RT13--23) therefore prove

\[
 |c_W(g)|\le\min\left\{1,
 \eta_a\frac{4(a-1)}{\sinh r}
 \sqrt{\frac{\operatorname{Var}(Z)}{{\bf E}Z-1}}\right\}.
                                                                    \tag{SN18}
\]

This remains a bound on the actual selected normalized trace, not an operator-norm claim. At r=0 use the bound one. The y=1 sphere endpoint is negligible in the same density/disintegration argument.

For a=f=168, the exact integer comparison eta\_168\^{}4\textless2\^{}−15 implies eta\_168\^{}f\textless2\^{}−630. Applying this to the unchanged DO080 bound and then the multiplicative rank-one actual-point transport of DO088 gives

\[
               |c(h)|<2\,2^{-931f/2-630}z^{-f}.                 \tag{SN19}
\]

The near-point product is still less than 2; no absolute-displacement estimate replaces its multiplicative control of arbitrarily small eta conjugates.

\subsection{6. Weighted G patterns before the actual cap}\label{weighted-g-patterns-before-the-actual-cap}

For fixed s,d and a paired rank-zero ideal J, the finite Green factor is 2(N\_EJ)\^{}(a/2)M\^{}(−a). Multiplying SN13 and SN19, its z-dependent part is proportional to

\[
 F(z)=z^{\kappa}(\sqrt z+2\sqrt R\,T^{-d/f})^{2f},
 \qquad \kappa=(a-3)f-273/10=276927/10>162f.                    \tag{SN20}
\]

For 0\textless z'≤z, both factors give F(z')/F(z)≤(z'/z)\^{}kappa. Let z\_sd=(M/T\^{}(s+d))\^{}(1/f) be the full rank-zero cap. Omitting a paired rank-zero prime of norm q lowers z by q\^{}(1/f). Summing over all 334 actual choices exactly as in DO088 therefore gives

\[
 \sum_J F(z_J)\le F(z_{s,d})
                  \prod_{v\in G_{\rm unmixed}}
                           (1+q_v^{-\kappa/f})<2F(z_{s,d}).      \tag{SN21}
\]

The last comparison uses q\_v\textgreater2\^{}34 and kappa/f\textgreater162, so the product is at most (1+2\^{}−5508)\^{}334\textless2. There is no additional unweighted 2\^{}334 multiplier. As in DO091, the monotone comparison applies even if the artificial cap does not itself satisfy the reduced set's y≥V condition.

At the cap the finite numerator cancels M\^{}a. Keep the exact marked double-pole and original pair weights from K077(AP18), namely binom(f,s)binom(s,d)(T+1)\^{}(s−d). With s=24 and

\[
 y_{s,d}=R/z_{s,d},\qquad y'_{s,d}=R(T^{s-d}/M)^{1/f},
\]

the resulting complete d-cell upper bound is

\[
\begin{split}
 U_{24,d}<&\;2^{-711}\,28\,2^{605}\binom{168}{24}\binom{24}{d}
 2^{-931f/2-630}(T+1)^{24-d}z_{24,d}^{-f}\\
 &\quad\cdot y_{24,d}^{f+273/10}V^{-21/2}
                 (1+2\sqrt{y'_{24,d}})^{2f}.
\end{split}                                                       \tag{SN22}
\]

Here −711=1+1+1−714 retains the finite Euler factor, the complete weighted G sum, the actual near-point product and the full real eta covolume. Every local pair and type was kept before summing. The original full-coordinate source coding, content identity and relative norm solution ensure no additional coordinate, class or generator factor is required.

\subsection{7. Exact endpoints, union and chronological scope}\label{exact-endpoints-union-and-chronological-scope}

The exact companion computes all 25 expressions SN22, including every binomial coefficient, by integer root brackets and directed rational logarithms. Each complete logarithmic upper bound is strictly less than the negative of its margin:

\begin{longtable}[]{@{}rr@{}}
\toprule\noalign{}
Marked double-pole count d & Certified inverse binary margin \\
\midrule\noalign{}
\endhead
\bottomrule\noalign{}
\endlastfoot
0 & 14782 \\
1 & 14615 \\
2 & 14450 \\
3 & 14286 \\
4 & 14121 \\
5 & 13958 \\
6 & 13794 \\
7 & 13631 \\
8 & 13468 \\
9 & 13306 \\
10 & 13143 \\
11 & 12979 \\
12 & 12800 \\
13 & 12527 \\
14 & 11900 \\
15 & 10886 \\
16 & 9752 \\
17 & 8599 \\
18 & 7443 \\
19 & 6288 \\
20 & 5132 \\
21 & 3978 \\
22 & 2823 \\
23 & 1670 \\
24 & 517 \\
\end{longtable}

The exact complete witness exponent is

\[
                  E_{\rm wit}(24)=-207073761/512<-400000.        \tag{SN23}
\]

The worst reduced cell has margin 517, and all other 24 cells have margin at least 1670. Consequently the entire original sector, including the independent witness branch, is bounded by

\[
             2^{-517}+24\,2^{-1670}+2^{-400000}<2^{-516}.         \tag{SN24}
\]

The certificate also sums all 25 individually directed margins and the witness bound. There is no additional unweighted G-pattern, type or endpoint multiplier. Positive upper expressions for neighboring supports, if considered, are failures of a sufficient estimate only.

The chronological predecessor is the canonical R069 error list with 72 all-top and 31 uniform entries. Append 516 once to obtain 73/31, preserving every older charge. Both complete exact dyadic totals remain below 2\^{}−230. A separate full structural exhaustion is needed for any new complete support floor; the proof does not infer such a floor from this fixed-W0 sector alone. The uniform data are unchanged.

\subsection{8. Sources and precise limits}\label{sources-and-precise-limits}

DO091(TD1--26) supplies the actual transverse split, primitive vector, exact content, real eta lattice, paired affine difference ideal and near-point normalization; K075(RN11--19) and K077(AP7--18) are the literal content/relative-norm/pencil inputs. The field and norm-degree conventions are K032 and DO091(TD11--12), with the actual E=F(i) relation. The new local ideal count SN5--15 is proved here. The selected metric, sphere density, detailed balance and variance remain the complete genuine model of DO080(RT4--23); only its t integration is sharpened in SN16--19. DO088 supplies the already checked weighted-pattern comparison, with the new positive exponent explicitly substituted.

No adopted field, local lattice, omitted label, source normalization, class weight, selected vector or compact twist is changed. This is not a construction of an integral selected-image inverse, a proof of a signed lower bound, or a global period/nonvanishing theorem. The full admissible host, selected arithmetic-image content, prescribed Bessel/Mram value, cycle-value transport and DP44 cost gates remain open.

\clearpage

\section{AM-LIFT-075 --- the actual good projector excludes every increasing-rank direct theta construction}\label{am-lift-075-the-actual-good-projector-excludes-every-increasing-rank-direct-theta-construction}

\textbf{Canonical checked integration --- 2026-11-22T21:14:51.517963+00:00.} Whole root and independent proof/source/script checks and separate exact replays passed. For a split type-II theta quotient from GL\_a to GL\_r with d=r−a\textgreater0, every nonzero principal-level finite constituent has one linear Green factor of multiplicity at least d.~For d\textgreater=2 squarefreeness kills it at every selected G place; for d=1 the actual goodTuple degree-one witness kills each tensor block. Thus the genuine full-rank e0 kills the actual compact-source direct theta cuspidal projection for every permitted even m\textgreater=2, where d=m−1. This excludes that direct construction only. It proves no vanishing of the prescribed Bessel period, coupled theta map or transported-leading-r intermediate projector. Fixed-real-type, fixed-finite-level global distinction remains open.

\textbf{Algebraic-geometric translation.} This is a local representation-theoretic obstruction in the existing genuine finite projector and a compact-source cuspidal theta pairing, not a new morphism of the CM Jacobian or its Picard scheme. Finite parabolic induction respects the separate Green orbit grading; this grading is not identified with a motivic weight or a new lambda-ring identity. Exact contextual anchors are SGA3 I4.7.3/4.7.3.1; SGA6 VI Theorem3.3; FGA232 Theorem3.1 pp148--149; FGA236 Corollaries1.7/2.3 pp229/232; SGA7 I VII2.7/2.9 pp169--174 and VIII7.4(b) pp310--311; cote4 scan003 §§8.1.1--8.1.2, scans005--006 §§8.4.1--8.4.3 and cote54 scans005--007. No new scheme comparison, K-theory equality, l-adic sheaf/weight theorem, CM-motive or motivic L-function assertion follows. DP\_MULT Q3e and THETA3 G1/A7 lose the specified direct compact-to-bottom candidate, while their coupled arithmetic family, normalization and value remain open. CX\_GRH CX-ALG-024 and GCT\_PLETH T-052 retain their separate gates.

Author: \nolinkurl{/root/lift_accounting}, 22 November 2026 UTC.

\textbf{Status: author whole-proof/source/script review PASS, 22 November 2026 UTC; exact diagnostic replay and frozen hashes are recorded in the adjacent verification files. Independent review and canonical promotion remain pending.} This is a local and cuspidal-projection obstruction for a specified construction. It is not a vanishing theorem for the prescribed Bessel period, the original coupled theta map, or its terminal restriction. No Lean build is claimed.

\subsection{1. Statement with the original target projector}\label{statement-with-the-original-target-projector}

Retain the genuine source finite groups and Green-series realization of \nolinkurl{NTypes.e0}, at the two completely split G primes. Let the target rank be r\textgreater=2. At each selected relative place v the target group and its original principal subgroup are

\[
 G_r(F_v)=\operatorname{GL}_r(F_v),\qquad
 K_r=\operatorname{GL}_r(O_v),\qquad
 K_r^1=\ker(K_r\longrightarrow\operatorname{GL}_r(k_v)).
 \tag{DT1}
\]

The finite projector acts on the original K\_r\^{}1-invariant subspace. Changing to an already prescribed split lattice frame conjugates both the groups and projector and changes no conclusion below.

Let 1\textless=a\textless r and d=r-a. At a split place take the genuine type-II oscillator for GL\_a x GL\_r, with any of its actual determinant splitting characters and contragredient conventions. For an arbitrary irreducible smooth source representation pi, let theta(pi) be its irreducible theta quotient. Then:

\begin{enumerate}
\def\labelenumi{\arabic{enumi}.}
\tightlist
\item
  Every irreducible finite constituent of theta(pi)\^{}(K\_r\^{}1), if that space is nonzero, has a semisimple parameter polynomial divisible by a linear factor.
\item
  More precisely, a single linear factor occurs with multiplicity at least d.~Therefore, if d\textgreater=2, \textbf{every squarefree finite semisimple-series projector annihilates this space at that one place}.
\item
  If d=1, the literal no-degree-one-divisor clause of goodTuple supplies a killing selected place in every tensor block of the actual e0. Consequently e0 annihilates the principal-level cuspidal projection of every direct theta lift from the actual compact U(I\_a) to the target of rank r.
\end{enumerate}

In the actual bottom family,

\[
 a=k,\qquad r=k+m-1,\qquad d=m-1,
 \tag{DT2}
\]

with even m\textgreater=2. Thus the direct compact-to-bottom theta construction is excluded for every permitted m. For m\textgreater=4 its finite obstruction is stronger than the old m=2 argument: squarefreeness alone kills the local quotient at every selected G place. For m=2 the no-linear-factor witness in each goodTuple is used. None of these conclusions requires a bound on same-sign Harish--Chandra gaps, a top CM label, a local genericity or temperedness hypothesis, or a choice of Schwartz vector.

The target projector here is the genuine \textbf{full rank-r} e0. This must not be confused with the transported leading-r projector acting on the larger intermediate groups in LIFT058. Those larger groups and their actual two-edge or longer theta operators are not excluded by this statement.

\subsection{2. The added character string and its literal finite consequence}\label{the-added-character-string-and-its-literal-finite-consequence}

Mínguez (2008), Theorem 1(2), printed page 718, applies to every irreducible smooth representation of GL\_a over a nonarchimedean local field, in every residual characteristic. If pi is the Langlands quotient of tau\_1 x \ldots{} x tau\_s, with essentially square-integrable tau\_i, its irreducible theta quotient on GL\_r is the Langlands quotient formed from

\[
 \nu^{-(d-1)/2},\ \nu^{-(d-3)/2},\ldots,
 \nu^{(d-1)/2},\quad
 \widetilde\tau_1,\ldots,\widetilde\tau_s,
 \qquad \nu=|\det|.
 \tag{DT3}
\]

The factors are put in the Langlands order when taking that quotient. Only their multiset and their GL1 ranks are used. Honest splittings multiply the actions by determinant characters; after incorporating them, the first d factors have the form

\[
 \eta\,|\cdot|^{-(d-1)/2},\ldots,
 \eta\,|\cdot|^{(d-1)/2}
 \tag{DT4}
\]

for one character eta. They all have the same restriction to O\_v\^{}times. Inverting the source convention, taking a contragredient, or permuting the ordered Witt basis changes the remaining inducing data and possibly eta; it does not remove these d equal compact characters. The theorem's uniqueness statement is for the irreducible quotient, not a semisimplicity assertion about the big theta module.

Here is the finite restriction calculation, including the nonvanishing proviso. For any standard parabolic P of block sizes n\_1,\ldots,n\_l and smooth Levi representations V\_i,

\[
 \left(\operatorname{Ind}_{P(F_v)}^{G_r(F_v)}
                    \bigotimes_i V_i\right)^{K_r^1}
 \simeq
 \operatorname{Ind}_{\overline P(k_v)}^{\operatorname{GL}_r(k_v)}
                    \bigotimes_i V_i^{K_{n_i}^1}.
 \tag{DT5}
\]

Indeed, Iwasawa decomposition realizes the induction as functions on K\_r. Right K\_r\^{}1 invariance descends the argument to GL\_r(k\_v); the relation under P intersect K\_r requires precisely the indicated Levi invariants. Its compact unipotent radical acts trivially on the inducing values. The normalized-induction modulus is one on that compact. Taking K\_r\^{}1 invariants is exact over C, by averaging over this compact group. Therefore the invariant space of the Langlands quotient in DT3 is a finite subquotient of DT5.

If eta is nontrivial on 1+varpi\_v O\_v, every added GL1 invariant space is zero, and the entire right side in DT5 is zero. Otherwise eta descends to a character eta\_bar of k\_v\^{}times. Each of the d factors in DT4 gives exactly eta\_bar, since all absolute-value twists are trivial on units. Its finite dual parameter is a nonzero degree-one Frobenius orbit alpha. No identification of alpha with an arbitrarily chosen eigenvalue is required.

For clarity, the Green parameter argument is not merely a statement about the existence of a generic constituent. The characteristic map is an isomorphism of graded character rings: finite ordinary parabolic induction becomes multiplication of symmetric functions, one alphabet for each dual Frobenius orbit. In the alphabet belonging to alpha, every added GL1 character contributes s\_(1). Thus every irreducible constituent of DT5 has total degree at least d in this same alphabet. Equivalently its semisimple parameter polynomial f satisfies

\[
                         (X-\alpha)^d\mid f(X).
 \tag{DT6}
\]

The remaining Levi factors can increase this multiplicity but cannot decrease it. The degree assertion follows from the separate grading in each orbit alphabet, not from an asserted irreducibility of the finite parabolic induction. Gow--Vinroot (2007), Section 3 and Theorem 3.2, printed pages 7--9, give precisely this multiplication/irreducible-Schur dictionary. It survives passage to every finite subquotient because finite complex representations are semisimple. This proves the local assertions in Section 1.

For d\textgreater=2, DT6 contradicts squarefreeness at every selected place. For d=1 it proves only the existence of a linear factor; squarefreeness alone would not exclude that case. The distinction is essential.

\subsection{3. The exact source goodTuple and the global cuspidal pairing}\label{the-exact-source-goodtuple-and-the-global-cuspidal-pairing}

At the pinned source commit \nolinkurl{7cbf7a0c32fca1badf05931879a37c3f8a2b03b4}, \nolinkurl{Formalism/CyclotomicNorms/StatementT.lean} 19--29 defines IsGood. Its second conjunct says that every f\_v is squarefree. Its fifth is

\[
 \forall b\ge1,\quad 2b\le r\ \Longrightarrow\
 \exists v:\ \text{no divisor of }f_v\text{ has degree }b.
 \tag{DT7}
\]

Putting b=1 is legitimate because r\textgreater=2. The result is a selected relative place with no linear factor. The witness may depend on the goodTuple and on the rational G prime; the proof does not choose one common place for all tensor blocks.

\nolinkurl{Formalism/SharpRestriction/LevelLetterCocycles/DataN.lean} 996--1007 defines goodTuple by the IsGood predicates for both the prescribed characters and their inverses, and defines e0 as the finite sum of tensor products of the actual eSer(f\_v). In the genuine realization, these are the central projectors to the finite semisimple series specified by those polynomials. The established Green realization, rather than an arbitrary formal action of an idempotent, is required.

For d=1, DT7 supplies a place at which DT6 kills the corresponding tensor block. For d\textgreater=2, the squarefree clause supplies the contradiction at any selected place. Hence all the tensor blocks and their actual sum act as zero on every local-quotient configuration arising from DT3:

\[
            e_0\,\operatorname{pr}_{\rm cusp}
                \Theta_{U(I_a),G_r}(\pi,\Phi)=0
 \tag{DT8}
\]

at the target's original principal level, for every actual initial automorphic constituent pi and every Schwartz input Phi for which the usual theta integral is taken.

To justify the global-to-local step, test the target against an irreducible cuspidal automorphic occurrence. The actual source U(I\_a) has compact adelic quotient. The target cusp form has rapid decay, so the standard theta pairing is absolutely convergent and gives local equivariant oscillator functionals by fixing the data at the other places. If the global pairing is nonzero, some pure tensor pairing is nonzero and its target local representation at every split place is the irreducible small theta quotient from Mínguez's theorem, in the corresponding dual convention. DT5--DT7 annihilate its K\_r\^{}1-invariant e0 subspace. This contradicts the supposed nonzero e0 cuspidal pairing. Multiplicity of either automorphic representation does not change the local obstruction, so it holds on every embedded occurrence.

Only this cuspidal projection is claimed. No global decomposition or semisimplicity of the full big theta output is used. If one instead starts from a noncompact smaller group or a regularized theta integral, DT1--DT7 remain valid locally, but its global implication additionally needs the actual localized cuspidal-pairing theorem for that integral; that extra theorem is not supplied here.

Every scalar arising from a different Haar normalization, genuine splitting, or a finite lattice-frame conjugation is retained in the actual map. Vanishing is unchanged by a nonzero scalar, but the argument does not set a global normalization or discard any nonzero scalar in a surviving construction. Wild added characters make the K\_r\^{}1 invariants zero; tame characters merely change alpha in DT6. Inversion of the parameter preserves its factor multiplicities and its linear-factor property.

\subsection{4. The remaining fixed-level relative-Poincaré issue}\label{the-remaining-fixed-level-relative-poincaruxe9-issue}

This section records the precise boundary of a separate possible construction. It is not a new nonvanishing theorem.

LIFT059 constructs an actual nonzero bottom cusp space at the prescribed principal-G/maximal-away-G level by an origin-value estimate. LIFT061 constructs the actual unit relative vector W in the fixed discrete series, with the physical Bessel subgroup, character, marked norm-minus-T line and quotient Haar. LIFT074 shows that, \textbf{if a specified pair of embedded automorphic representations is globally Bessel-distinguished}, its calibrated finite vectors and the complete W contraction give a nonzero literal Pu(T). None of these three statements supplies that global distinction from local models alone.

Muić (2015), Lemmas 6-1 and 6-2, gives a nonvanishing criterion for a relative Poincaré series from an integrable seed psi. In its scalar setting it requires a compact C whose difference set meets the arithmetic group only in the retained period subgroup and whose image contains more than half the L1 mass of psi. Theorems 6-7 and 6-9 retain an integrable discrete series, then \textbf{shrink a finite congruence subgroup} until the difference-set condition holds. Their threshold depends on the compact set furnished by integrability; no prescribed-level bound is stated. The actual Bessel problem also has its compact H representation and its exact finite type, which must be carried in any vector-valued application.

By contrast, Prasad--Schulze-Pillot (2008), Theorem 4.1, printed page 229, and its proof on page 231, proves a distinguished globalization with specified finite supercuspidal data using an arbitrarily small archimedean orbital-integral support. It does not preserve the fixed real discrete series/minimal vector of LIFT061, nor does it assert the present unchanged finite test-vector package. Applying the fixed real matrix-unit projector after that construction needs a nonzero projected period; positivity of the projector alone does not prove it.

The exact physical move in LIFT060/LIFT30 has real singular values sqrt(T) and 1/sqrt(T) on its distinguished plane and a lower-unipotent denominator at T. Thus the origin isolation estimate in LIFT059 cannot simply be reused for the moved relative packet. LIFT061 retains the actual weighted L1 norms and a positive construction energy, without a uniform lower bound for that energy. One must still prove either the requisite half-mass/isolation bounds for the \textbf{fixed} seed and level, or a different genuine global distinction argument. Direct increasing-rank theta lifting cannot provide the missing pair by DT8.

The unresolved global coupling, the requested arithmetic Read/STiesR/032 instance, original selected ordinary derivatives and class sums, and every DP44 comparison cost remain separate. DT8 makes no assertion that the Bessel period itself vanishes and excludes no actual upward tower whose larger finite projector is only the transported leading-r projector.

\subsection{5. Sources and verification scope}\label{sources-and-verification-scope-12}

Mínguez, A. (2008). Correspondance de Howe explicite: Paires duales de type II. \emph{Annales scientifiques de l'École Normale Supérieure, 41}(5), 717--741. Theorem 1(1)--(2), p.718; normalization (1.4), pp.721--722. \url{https://www.numdam.org/item/ASENS_2008_4_41_5_717_0.pdf}. The primary statement was freshly read both in the retained original text and the publisher PDF. In particular DT3 is the full positive-gap string in the printed theorem.

Gow, R., \& Vinroot, C. R. (2007). \emph{Extending real-valued characters of finite general linear and unitary groups on elements related to regular unipotents} {[}Preprint{]}. Section 3, pp.7--9, Theorem 3.2. \url{https://arxiv.org/abs/0704.2022}. The retained primary text's multiplication and separate orbit-alphabet grading were read directly.

Muić, G. (2015). \emph{Fourier coefficients of automorphic forms and integrable discrete series} {[}Preprint{]}. \url{https://arxiv.org/abs/1505.02263}. Lemmas 6-1/6-2 and Theorems 6-7/6-9, pp.21--26, were freshly read. Their changed-level scope is retained.

Prasad, D., \& Schulze-Pillot, R. (2008). Generalised form of a conjecture of Jacquet and a local consequence. \emph{Journal für die reine und angewandte Mathematik, 616}, 219--236. Theorem 4.1, p.229, and proof, p.231. Author PDF: \url{https://www.math.iitb.ac.in/~dprasad/crelle.pdf}. The actual small-archimedean-support argument was directly read; no fixed-real-type version is inferred.

Accepted comparisons: DO045 Section 3 (the gap-one case); LIFT059, LIFT061/061-U, LIFT062 and LIFT074; the exact marked-line/finite transport in LIFT060 and LIFT30. The new local argument extends DO045 to every positive gap and records the stronger repeated-factor obstruction when the gap is at least two.

The companion script checks exact character-string lengths/parities, finite-field repeated linear factors under inversion and character twisting, actual rank specializations, and literal source/primary-text anchors. These are diagnostics, not a mechanical proof of Howe duality, Green theory, global theta localization or a global period. Final hashes and the whole-file/replay record are separate.

\clearpage

\section{AM-K094 --- the actual full covector fibre removes the Fourier sum from the selected matrix}\label{am-k094-the-actual-full-covector-fibre-removes-the-fourier-sum-from-the-selected-matrix}

\textbf{Canonical checked integration --- 2026-11-22T21:14:51.895837+00:00.} Whole root and independent proof/source/script checks and separate exact replays passed. For each actual coarse initial class retaining all non-T data, G-neatness makes the full strict covector fibre the product of all nonzero covectors. Its varying global coefficient lattices have one common localization at11. Punctured Fourier transform is an integral isomorphism of the line-augmentation modules there, with inverse T\^{}(−a)F\_minus and no division by T−1. The actual transformed Mram basis has exact class norm T\^{}(af/2)2\^{}(f/2)\textbar\textbar P\textbar\textbar. After the common full-middle contracted cut, its image is exactly the signed2\^{}f corner column FC13, with original masses, central phases, coordinate pullbacks, one shared coefficient P and all residual ordinary derivatives retained. FC14 identifies the actual localized image/content ideal with the corner-column span, even when that image is zero. Global clearing sizes, corner/middle-class noncancellation, source/Read/032 matching and every DP44 cost remain open.

\textbf{Algebraic-geometric translation.} The argument is an integral change of basis at11 on the actual finite locally free coefficient lattices over the fixed covector fibre. The finite Fourier transform uses the genuine torsion character pairing, and its rational rank-one representatives retain the rigidified torsion difference line and Poincare biextension central scalar. It leaves the actual ordinary differential functional unchanged. Exact contextual anchors are SGA3 I4.7.3/4.7.3.1; SGA6 VI Theorem3.3; FGA232 Theorem3.1 pp148--149; FGA236 Corollaries1.7/2.3 pp229/232; SGA7 I VII2.7/2.9 pp169--174 and VIII7.4(b) pp310--311; cote4 scan003 §§8.1.1--8.1.2, scans005--006 §§8.4.1--8.4.3 and cote54 scans005--007. No new lambda-ring/K-theory identity, l-adic purity or weight theorem, CM-motive comparison or motivic L-function statement is inferred. DP\_MULT Q3e and THETA3 G1/A7 gain the exact corner-image/content criterion, while their evaluated derivative/class gate remains. CX\_GRH CX-ALG-024 and GCT\_PLETH T-052 retain their separate arithmetic gates.

Author: ALP\_MACHINE \nolinkurl{/root/kummer_extension}, 22 November 2026 UTC.

\textbf{Author whole-proof and complete-script PASS; frozen for independent review.} This is an exact change of basis on the existing initial ramified module, localized at 11. It does not change the allowed Y set, enlarge a finite level, replace the selected ordinary derivative, or prove that a selected matrix entry is nonzero. The global clearing factors and the remaining middle-class corner sum are explicit below.

\subsection{1. The complete existing covector fibre and its common local coefficient lattice}\label{the-complete-existing-covector-fibre-and-its-common-local-coefficient-lattice}

Retain either actual J14 or J56 host, its word, tuple, canonical twist, strict initial T level, and the literal generic G level of K038--K093. Let R be the completed coefficient DVR at a place above 11, with the exact CharT additive characters in the carried coefficient field. Let f be the number of selected T places and a the initial rank. Every residue field here is k\_v=F\_T.

For indexing only, forget the T covectors of an initial class while retaining its unimodular lattice N and all non-T data, including the generic G frame/character and any admitted Q frames. Call the resulting coarse class c.~No theta input is placed at this coarser level. Its rational stabilizer is trivial: the proof of K042(NB1--6) uses only the G-unipotent reductions in the faithful rational representation of dimension 2fa, and g0\textgreater2fa+1. Forgetting T does not change that premise. Definite stabilizers are finite, so neatness makes them trivial.

For a fixed representative of c, put \[
 V_{c,v}=N/\mathfrak P_vN,\qquad
 \Omega_c=\prod_{v\mid T}(V_{c,v}^{\vee}\setminus\{0\}).
 \tag{FC1}
\] Here the quotient can equivalently be formed from the localized free module. Its dimension over k\_v is a. Literal \nolinkurl{Deco.Ok} permits exactly a nonzero residue-linear covector, extended by zero off its localized lattice. Thus the actual strict classes above c are indexed by all alpha in Omega\_c, with no further rational quotient. Denote them by d\_(c,alpha). Their N, non-T frames and non-T character trivializations are fixed. Their initial stabilizers are one.

The global coefficient lattice need not be identical at these different covectors: \nolinkurl{CoefCol.colVal} uses \nolinkurl{delAtr(Lsh)}, and the kernel lattice changes when the covector line changes. At every place over 11, however, the kernel lattice and its Hermitian dual have the same localization for every alpha, since their difference from N is supported only at T. Hence the literal transported Weyl lattices have a common localized lattice L\_(c,R) in the same coefficient space. Use the all-11 standard representatives of K039(IG3); no rational change of frame is called integral in an earlier display. It follows that the actual local source block is \[
 \mathcal M_{c,R}=\operatorname{Fun}(\Omega_c,R)\otimes_R L_{c,R}.
 \tag{FC2}
\]

For a finite k-vector space V define \[
 \operatorname{Ram}(V)=\{b:V\setminus\{0\}\to R:
       \sum_{s\in k^\times}b(sv)=0\text{ for every }v\ne0\}.
\] The source diamonds have identity coefficient map. Applying the exact line-by-line augmentation decomposition of K090 to all lines gives \[
 \mathcal M_{c,R}^{\rm ram}
   =\left(\bigotimes_{v\mid T}\operatorname{Ram}(V_{c,v}^{\vee})\right)
                      \otimes_R L_{c,R}.
 \tag{FC3}
\] The complete Mram module is the direct sum of these blocks. Choose a vector lambda\_l on each one-dimensional subspace l of V. Then \nolinkurl{delta_(s lambda_l)-delta_(lambda_l)}, with s!=1, is an R-basis of Ram(V). This basis statement does not divide by T−1.

\subsection{2. The exact punctured Fourier isomorphism at 11}\label{the-exact-punctured-fourier-isomorphism-at-11}

Use the actual pairing alpha(lambda) and the literal nontrivial CharT character psi\_v. Define \[
 \begin{split}
 (\mathcal F_{v,+}b)(\alpha)
     &=\sum_{\lambda\ne0}\psi_v(\alpha(\lambda))b(\lambda),\\
 (\mathcal F_{v,-}a)(\lambda)
     &=\sum_{\alpha\ne0}\psi_v(-\alpha(\lambda))a(\alpha).
 \end{split}
 \tag{FC4}
\] All entries are actual algebraic units at 11. The characters already belong to the carried coefficient field by CharT.Ok and K059; no multiplicative-character extension is used.

Character orthogonality gives \[
 \mathcal F_{v,-}\mathcal F_{v,+}=T^a I-J,
 \qquad
 \mathcal F_{v,+}\mathcal F_{v,-}=T^a I-J',
 \tag{FC5}
\] where J and J' are the total-sum operators on the respective punctured fibres. The same identity is DO056(AC9--10), now used as a finite matrix on the fixed coarse fibre, not as an unidentified global neighbor operator.

The two Fourier maps preserve Ram. Indeed, in the sum over a scalar orbit of alpha, the inner sum over its dilates is unchanged when lambda is multiplied by a nonzero scalar. Summing first on each lambda line therefore gives zero. The same proof works with the roles reversed. A ramified function has total sum zero, so J and J' vanish. Since T is an 11-unit, the restriction is an integral R-isomorphism with inverse \[
 \mathcal F_{v,+}^{-1}=T^{-a}\mathcal F_{v,-}.
 \tag{FC6}
\] This uses no inverse of T−1 or of a finite general-linear group order. Tensoring gives an R-isomorphism F\_(c,+) from the tensor of vector-side Ram modules onto the actual source block FC3, with inverse T\^{}(−af)F\_(c,−).

For the tensor basis vector b\_(lambda,s)=tensor\_v(delta\_(s\_v lambda\_v)-delta\_(lambda\_v)), where lambda\_v is the chosen nonzero representative of a residue line and every s\_v!=1, its Fourier image has the explicit actual class values \[
 C_{c,\lambda,\mathbf s,P}(d_{c,\alpha})
  =\prod_v\bigl(\psi_v(\alpha_v(s_v\lambda_v))
                    -\psi_v(\alpha_v(\lambda_v))\bigr)P,
 \tag{FC7}
\] and is zero outside this coarse fibre. Choose algebraic coefficient vectors P whose localizations form an R-basis of L\_(c,R), using the actual standard frames as in K059. These Fourier images form an R-basis of the actual localized Mram module. Individual coefficients can be zero; the basis vector itself is primitive by FC6.

At a complex realization the Fourier adjoint has the displayed minus sign, and the class metric is the unweighted sum because the stabilizers are trivial. Thus the exact norm is \[
 \|C_{c,\lambda,\mathbf s,P}\|_{0,\rm cls}
                   =T^{af/2}2^{f/2}\|P\|.
 \tag{FC8}
\] This is a local-content basis change, not a smaller global norm bound. If global integral inputs are needed, retain a common coefficient-clearing factor d\_c for all the finitely many covector lattices and use d\_cP. Standard all-11 frames permit it to be chosen prime to 11 as in K059; its archimedean size is not bounded here. FC8 then contains \textbar d\_c\textbar{} at that realization. The inverse normalization T\^{}(−af) is integral at 11 but is not claimed globally integral at T. The source Haar/compact-type conversion remains the one in DO056(AC3).

\subsection{3. One common actual germ for the full fibre}\label{one-common-actual-germ-for-the-full-fibre}

Fix the coarse initial representative c, an actual middle class h and a terminal decoration d.~Write u\_(1,h) and u\_2 for their actual T covectors. The complete two-edge finite tensor for initial covector alpha has the local T conditions \[
   X\text{ integral},\quad u_{1,h}X=\alpha_v,
   \qquad Z\text{ integral},\quad u_2 Z=u_{1,h}.
 \tag{FC9}
\] All other factors, the complete Y sum, and the non-T frames are independent of alpha. This is literal \nolinkurl{phiT}, \nolinkurl{XiCpl1}, \nolinkurl{XiCpl2}, and \nolinkurl{ChKer.Shape.cp1/cp2}: the theta distribution is evaluated on the underlying N and the non-T frames; \nolinkurl{Lsh} controls the coefficient lattice, not an additional alpha-dependent theta distribution. In \nolinkurl{frsF}, the G lift depends only on N and its G frame, and the Q frames were fixed. Thus the common germ assertion here is stronger than the scalar-line assertion only because the localization of the varying coefficient lattice has been checked in section 1.

Delete the first equality in FC9 to define the proof-device base tensor Phi\_(c,h,d)\^{}base. Keep its integral indicator, the second incidence, the full Y sum, and all actual G/Q data. Choose one full-middle-Gamma transversal cut C\_h as in K092 for this whole fibre. Its G residue choice is independent of alpha, so the same framed cutoff works for every alpha and every P. It commutes with the T multipliers below. The contracted full-middle-mass identity remains exactly K092(PM6--7); no uncontracted section is divided by the middle stabilizer.

For mu=(mu\_v) in product\_v V\_(c,v), let M\_(mu,h) multiply the first position variable by product\_v psi\_v(u\_(1,h) X mu\_v). Then \[
 \Psi_{c,h,d,\alpha}^{\rm cut}
  =T^{-af}\sum_{\mu\in\prod V_{c,v}}
       \prod_v\psi_v(-\alpha_v(\mu_v))
               M_{\mu,h}\Phi_{c,h,d}^{\rm base,cut}.
 \tag{FC10}
\] This is K093(RT2) for the complete nonzero-covector fibre. The finite sum temporarily includes zero mu; no zero initial covector is admitted as an arithmetic input.

Choose one common comparison denominator D, prime to 11, covering these finitely many base and multiplied tensors as in K089/K092. This changes neither an admitted finite level nor Y. Retain its exact original-coordinate pullback in every derivative. Keep one fixed coefficient P across all middle classes. Define the following actual target-column-valued linear operator on that coefficient space: \[
 \mathcal E_{c,d}(\mu;P)
 =\sum_h\mathcal J_{h,d,P}
       \bigl(\Theta_D(M_{\mu,h}\Phi_{c,h,d}^{\rm base,cut})\bigr).
 \tag{FC11}
\] The equality can equivalently be read after any fixed actual terminal dual test. Every J is the prescribed complete two-edge ordinary selected evaluation with the source cD, coordinate pullback, determinant/splitting, and period factors of K038/K092. The middle mass division is already in the contracted cut identity. It is not repeated in FC11. The same P occurs in every term.

\subsection{4. Cancellation of the entire Fourier sum}\label{cancellation-of-the-entire-fourier-sum}

Substitute FC7 and FC10 into the actual class map A and sum over its nonzero initial covectors. At one T place, punctured orthogonality gives, for arbitrary mu, \[
 \begin{split}
 &\sum_{\alpha\ne0}
   [\psi_v(\alpha(s\lambda))-\psi_v(\alpha(\lambda))]
                      \psi_v(-\alpha(\mu))\\
 &\hspace{20mm}=T^a[\mathbf1_{\mu=s\lambda}
                              -\mathbf1_{\mu=\lambda}].
 \end{split}
 \tag{FC12}
\] The two puncture constants cancel. Tensoring cancels the T\^{}(−af) in FC10 exactly. For epsilon in \{0,1\}\^{}f put mu\_(epsilon),v=s\_v lambda\_v when epsilon\_v=1 and lambda\_v otherwise. The result is \[
 \boxed{\quad
 (\mathcal A C_{c,\lambda,\mathbf s,P})(d)
    =\sum_{\epsilon\in\{0,1\}^f}
        (-1)^{f-|\epsilon|}
                   \mathcal E_{c,d}(\mu_\epsilon;P).
 \quad}
 \tag{FC13}
\] Every surviving mu\_epsilon is nonzero at every place. There is no residual T\^{}(af), no scalar-orbit average, and no new initial mass. The same calculation applies before or after a fixed terminal branch/test, with its actual scalar and denominator retained. It does not commute a selected differential past an unjustified integral frame.

For each surviving multiplier use K093(RT4--7)'s exact rational rank-one lift, chosen using the actual middle covector. Thus each term in FC11 is the exact residual-derivative expression RT12 at the corresponding prime-to-11 torsion point. The central phase, the sign +t, and the original coordinate pullback are unchanged. No additive dependence of these rational lifts on mu is required or asserted; FC12 was proved for the actual finite multipliers before choosing their geometric representatives.

In particular every term still contains at least D\_tau−nu\_(tau,1)\textgreater0 ordinary first-edge derivatives at each label, and the selected second-edge derivatives remain. FC13 removes the full Fourier summation from the matrix after an invertible source basis change. It does not turn any term into a scalar theta value or replace the residual jets by independent data.

\subsection{5. Exact content criterion and the remaining gate}\label{exact-content-criterion-and-the-remaining-gate}

Taking all c, vector lines lambda, nonidentity scalars s and coefficient basis vectors P, FC6--7 and FC13 give the equality of actual R-submodules \[
 \mathcal A(\mathcal M_{\rm ram,R})
   =\sum_{c,\lambda,\mathbf s,P}R\cdot
      \left(d\longmapsto
        \sum_\epsilon(-1)^{f-|\epsilon|}
                    \mathcal E_{c,d}(\mu_\epsilon;P)\right).
 \tag{FC14}
\] Consequently the actual image-content ideal against any fixed complete target detector family equals the ideal of these corner columns, retaining each detector's genuine/generic projector, branch and denominator. This equality also holds when the entire image is zero. It is a local equality at 11; clearing factors for a chosen global integral family remain as in FC8 and do not acquire a zero valuation by omission.

If a further prescribed source submodule B is imposed, the exact transformed domain is F\_(c,+)\^{}(-1)B, blockwise. No statement here makes B stable under this Fourier map or replaces it by the full Mram module. FC14 concerns the full actual domain displayed there; its restriction must retain that preimage.

The corner columns still sum the prescribed residual ordinary derivatives over the actual middle classes. Neither the corners nor the middle-class contributions are independent input variables. They can cancel, and a selected derivative can vanish on an individual translated germ. Those evaluations, the final target restrictions, and their exact arithmetic content remain the unevaluated gate. The theorem proves no new signed finite-kernel estimate, global Bessel/Mram period, sharp content upper bound, DP44 inequality, or admissible-family constructor.

\subsection{6. Exact source and verification scope}\label{exact-source-and-verification-scope}

Direct source interfaces are \nolinkurl{StdLatUnimodular.lean}, Deco.Ok and the definitions of at, dia, L and Lsh; \nolinkurl{CoefCol.lean}, colVal and delAtr(Lsh); \nolinkurl{ScaledRamifiedLattice.lean}, the identity-coefficient diamond and unweighted NmG; and \nolinkurl{FiniteSumNonzeroLargeDepth.lean}, phiT, frsF, Y, XiCpl1/2, kerOf and ChKer.Shape.cp1/cp2. They are under \nolinkurl{work/research/alpoge_source/Formalism/LocalToGlobal/Columns}. The actual generic initial subgroup and its all-place neatness are K042(NB1--6). The local finite Fourier calculation and source norm convention are DO056(AC2--13). The full-Y/frame/mass and ordinary-derivative interfaces are K038--039, K059, K090--093. No new external theorem about theta correspondence is invoked.

The companion passed 12,343 exact checks: punctured Fourier identities in exact cyclotomic rings, the line-augmentation basis and inverse without T−1, tensor corner identities with different middle-class operators but one coefficient input, exact norms, the local/global integrality distinction, and fingerprints of the actual source interfaces. The residue-characteristic-11 tests include a scalar-orbit group of order22, so no invertibility of T−1 is implicit. These finite algebra checks do not compute a selected enormous theta value or prove its nonvanishing.

\subsection{Author whole verification}\label{author-whole-verification-3}

At 2026-11-22T20:56:35+00:00, I completed a whole read of all six substantive sections FC1--14 and the complete companion script. I directly reread Deco.Ok, the actual L/Lsh and colVal definitions, identity-coefficient diamonds, phiT, frsF, the complete Y and coupled-kernel definitions, and ChKer.Shape.cp1/cp2. I reconstructed the full-fibre neatness using only the unchanged all-place G condition, distinguished the varying global Lsh lattices from their common localization at11, and checked that the theta distribution depends on N and non-T frames independently of the initial covector. I checked the exact plus/minus Fourier signs, the two puncture cancellations, the tensor normalization and norm, the common framed middle cut, the same P across all middle classes, and K093's unchanged rational lifts/residual derivatives. I separately retained global coefficient clearing, terminal projector denominators, and the inverse image of a further prescribed submodule. The first full certificate passed 12,343 assertions at 2026-11-22T20:56:04.322588+00:00. The final status, common comparison-cover sentence and verification scope were reread; a final replay and fingerprints are recorded separately.

\clearpage

\section{AM-LIFT-076 --- every actual good-projector bottom constituent has simple tempered base change}\label{am-lift-076-every-actual-good-projector-bottom-constituent-has-simple-tempered-base-change}

\textbf{Canonical checked integration --- 2026-11-22T21:22:04.664811+00:00.} Whole root and independent proof/source/script checks and separate exact replays passed. Every existing cuspidal bottom constituent selected by the genuine full e0 in the stated actual setting has simple generic global Arthur parameter and conjugate-selfdual cuspidal, cohomological base change, tempered at every finite place. Squarefree finite Green types pin the complete distinct inertial multiset; actual unitarity forces all real exponents to zero. The actual ramified auxiliary-theta comparison excludes Arthur SL2 lengths greater than one, and the literal all-degree divisor witnesses exclude multiple global blocks. Original odd-rank parity permits unit Harish--Chandra gaps. Generic packet membership gives tempered local packets at every place and spectral multiplicity one. No one-dimensional finite-level or coefficient-space assertion, source-family existence, globally distinguished pair, arithmetic image/value, Read/032 or DP44 conclusion follows.

\textbf{Algebraic-geometric translation.} This classifies actual automorphic constituents of the existing Hermitian unitary group through its finite Green types and global Arthur/base-change parameters. It does not identify those automorphic representations with a motive of the particular CM Jacobian or construct the requested theta-cycle map. The coefficient parity is a statement about the original algebraic highest weight; the temperedness theorem is not a new l-adic sheaf-purity assertion. Exact contextual anchors are SGA3 I4.7.3/4.7.3.1; SGA6 VI Theorem3.3; FGA232 Theorem3.1 pp148--149; FGA236 Corollaries1.7/2.3 pp229/232; SGA7 I VII2.7/2.9 pp169--174 and VIII7.4(b) pp310--311; cote4 scan003 §§8.1.1--8.1.2, scans005--006 §§8.4.1--8.4.3 and cote54 scans005--007. No new scheme comparison, lambda-ring or K-theory identity, CM-motive comparison or motivic L-function theorem is inferred. DP\_MULT Q3e and THETA3 G1/A7 gain the larger-representation local temperedness premise for every existing actual bottom constituent, while their global distinction, integral image and value gates remain. CX\_GRH CX-ALG-024 and GCT\_PLETH T-052 retain their separate arithmetic comparisons.

Author: \nolinkurl{/root/lift_accounting}, 22 November 2026 UTC.

\textbf{Status: author whole-proof/source/script review PASS; independent review and canonical integration pending.} The companion certificate records 83,280 exact diagnostics, with their limited scope stated below. This is an extension from the chosen LIFT069 subspace to every existing cuspidal constituent selected by the genuine full e0. It does not construct a Host, a full sized family, a distinguished global pair, or the requested arithmetic value.

\subsection{1. Genuine setting and the precise extension}\label{genuine-setting-and-the-precise-extension}

Retain the genuine unitary realization, actual field E/F, even k,m\textgreater=2, prescribed bottom discrete series, finite levels and genuine Green-series realization used in LIFT059 and LIFT069. Put

\[
                 r=k+m-1,\qquad G=U(Y_0).
\tag{GT1}
\]

The rank r is odd. Let pi be an irreducible cuspidal automorphic representation occurring in the original bottom space, with a nonzero vector in the actual full e0 projection on its original principal-G invariants. Here e0 is the literal sum over all admitted goodTuple blocks, not only the particular full-degree/mixed tuple chosen in LIFT059. At infinity pi has the original bottom discrete series with Harish--Chandra multiset

\[
             h=(v_0,\ldots,v_{k-2},-b_0,\ldots,-b_{m-1}).
\tag{GT2}
\]

\textbf{Claim.} The global Arthur parameter of pi is simple and generic. Its base change Pi is conjugate-selfdual cuspidal on GL\_r(A\_E), cohomological and tempered at every finite place. The actual local representations pi\_v belong to tempered packets at every place. In particular their original split-T spherical vectors satisfy the larger-representation hypothesis of LIFT068 and LIFT071. The global discrete occurrence is multiplicity one by Chen--Zou Theorem 2.6, with no inference that its finite-level or coefficient-type subspace is one-dimensional.

Existence of any such pi is a separate question unless the complete genuine LIFT059 package is supplied; that theorem then supplies examples inside its chosen subspace. The present argument proves no global Bessel distinction. It does not need a same-sign HC gap greater than one. The strict ordering and parity of GT2, already part of the genuine bottom type, suffice for the final coefficient calculation. Full SizesAt remains a premise wherever the source family or LIFT059 existence theorem requires it.

\subsection{2. The squarefree finite series pins every inertial multiplicity}\label{the-squarefree-finite-series-pins-every-inertial-multiplicity}

Fix an actual split G place v. Write K\_d=GL\_d(O\_v), K\_d\^{}1=ker(K\_d -\textgreater{} GL\_d(k\_v)). For a monic squarefree polynomial f of degree r with nonzero constant term, write

\[
                   f=\prod_{i=1}^{s} f_i,
           \qquad d_i=\deg f_i,\quad \sum_i d_i=r,
\tag{GT3}
\]

with distinct monic irreducible f\_i. In the compatible dual-orbit dictionary of the genuine eSer realization, let tau\_i be the irreducible finite cuspidal GL\_(d\_i)(k\_v) representation indexed by f\_i. Distinct f\_i give distinct finite cuspidal representations when their degrees agree. The finite semisimple series of f consists of the unique irreducible representation

\[
              tau_f=\operatorname{Ind}_{\bar P}^{GL_r(k_v)}
                            \bigotimes_i tau_i.
\tag{GT4}
\]

This follows directly from the Green characteristic map: each distinct orbit alphabet has precisely the partition (1), so the product of its s\_(1) factors is the single Schur product indexed by those distinct alphabets. This is the same ordinary finite parabolic induction as in DO046(AP12), not Deligne--Lusztig induction with an omitted sign. Gow--Vinroot Section 3/Theorem 3.2 supplies the characteristic-map statement. It also shows that GT4 has the full indicated multiset, including multiplicities; squarefreeness is exactly what makes every orbit occur once.

For completeness, unitary depth-zero supercuspidal lifts rho\_i of the tau\_i exist without an extra choice of local test. Indeed, Nadimpalli Section 2.2, printed page 3, supplies an irreducible quotient of c-Ind\_(K\_di)\^{}(GL\_di) tau\_i. Frobenius reciprocity puts tau\_i in that quotient. The exact compact-invariant identity of DO046(AP12a), or LIFT069(CT5), shows that this quotient must be supercuspidal: otherwise tau\_i would occur in induction from a proper finite parabolic. Its central character is unitary on O\_v\^{}times. Twist by a real unramified power of \textbar det\textbar{} to make the full central character unitary; this changes no K\_d type.

A supercuspidal representation with unitary central character is unitarizable. Casselman's 1995 draft, Theorem 5.2.1, Corollary 5.2.3 and Lemma 5.2.5, printed pages 46--48, give the compact-mod-centre, unitarity and real-unramified-twist statements used here. One direct verification fixes a nonzero smooth dual functional ell and integrates

\[
 B(u,w)=\int_{G_d/Z_d} ell(rho_i(g)u)
                          \overline{ell(rho_i(g)w)}\,dg.
\tag{GT5}
\]

Supercuspidal matrix coefficients have compact support modulo the centre; the unitary central character makes the integrand descend to the quotient. The integral is finite, positive and invariant. Its kernel is an invariant subspace, so irreducibility and ell!=0 make it positive definite. This proves the unitary choice used below, without assuming an arbitrary smooth quotient of a unitary representation is itself a Hilbert quotient. For d\_i=1 this is simply extension of a finite unit character with a unitary value on a uniformizer.

Nadimpalli Section 3, printed pages 6--7, states that rho\_i\^{}(K\_di\^{}1)=tau\_i and that (K\_di,tau\_i) is a type for {[}GL\_di,rho\_i{]}. Consequently two different tau\_i of equal degree give different inertial classes: an unramified twist does not change K\_di\^{}1 invariants. Let P(1) be the inverse image of the finite parabolic in K\_r, and inflate tensor tau\_i to P(1). The same section identifies (P(1),tensor tau\_i) as a Bushnell--Kutzko type for

\[
             \mathfrak s_f=
                 [\prod_i GL_{d_i},\ \bigotimes_i rho_i].
\tag{GT6}
\]

In particular Nadimpalli Corollary 3.3, including its converse proof, says that every irreducible constituent of Ind\_(P(1))\^{}K tensor tau\_i is typical for this entire inertial class. If an irreducible smooth pi\_v contains tau\_f, its supercuspidal support therefore has the form

\[
        \{rho_i\otimes\chi_i:\ 1\le i\le s\},
       \qquad \chi_i\text{ unramified},
\tag{GT7}
\]

up to ordering. Each one of the distinct inertial classes appears \textbf{exactly once}. No unnamed extra constituents or additional multiplicities can be inserted. This is the full multiset definition of inertial equivalence in Nadimpalli Section 2.2, not only equality of its underlying set. If s=1, the maximal-compact type statement gives the same conclusion directly.

\subsection{3. Unitarity removes all real exponents}\label{unitarity-removes-all-real-exponents}

The actual automorphic pi\_v is unitary. Write chi\_i=chi\_i\^{}u \textbar det\textbar\^{}(x\_i), with chi\_i\^{}u unitary unramified and x\_i real, and absorb chi\_i\^{}u into the unitary rho\_i. The Hermitian dual of a unitary rho\_i is itself. Hence the support of the Hermitian dual pi\_v\^{}h is

\[
                    \{rho_i|\det|^{-x_i}:1\le i\le s\}.
\tag{GT8}
\]

Unitarity gives pi\_v=pi\_v\^{}h. Uniqueness of supercuspidal support, including its multiplicities, equates GT7 and GT8. Distinct inertial classes force a factor of class i to match the same class i. Thus rho\_i\textbar det\textbar\^{}(x\_i)=rho\_i\textbar det\textbar\^{}(-x\_i). Taking the absolute value of the central character at a scalar uniformizer gives q\_v\^{}(-d\_i x\_i)=q\_v\^{}(d\_i x\_i), and therefore x\_i=0.

It follows that pi\_v is a constituent of normalized parabolic induction of unitary supercuspidal representations. The latter are square-integrable modulo centre, hence tempered. Waldspurger (2003), Lemma III.2.3, printed pages 266--267, states that induction of an admissible tempered representation is tempered; his Section I.3, printed page 243, explicitly uses the factor delta\_P\^{}(1/2), so this is the normalized induction in GT7. Temperedness passes to the irreducible subquotients here: exactness of normalized Jacquet modules puts their exponents among those of the inducing representation, and Waldspurger Proposition III.2.2 applies the same closed-cone criterion. This does not assume irreducibility or Hilbert semisimplicity of the whole induction. Its semisimplified Weil parameter is

\[
             \phi_v^{W,ss}=\bigoplus_{i=1}^{s}\phi_{rho_i},
\tag{GT9}
\]

where the summands are irreducible Weil representations and no two are unramified twists of one another. The usual GL correspondence preserves normalized induction, twists and supercuspidality; the precise semisimple assertions are Scholze Theorem 1.2(b)--(d). In particular each inertial equivalence class in GT9 occurs once. The supercuspidal supports in GT7 also show that no nontrivial monodromy segment can occur.

This proves local temperedness at every selected split G place for every actual squarefree block. It is stronger than finite genericity alone: complementary-series support with two copies of the same inertial class and exponents x,-x evades the argument, and is excluded here specifically by GT3--GT7. No such complementary-series counterexample can be dismissed solely by naming a finite Whittaker model.

\subsection{4. Apply the type statement to the literal full e0}\label{apply-the-type-statement-to-the-literal-full-e0}

The source definitions are at pinned commit 7cbf7a0c32fca1badf05931879a37c3f8a2b03b4. \nolinkurl{Formalism/CyclotomicNorms/StatementT.lean} 19--29 makes every admitted f\_v monic, squarefree, of degree r and with nonzero constant term. \nolinkurl{Formalism/SharpRestriction/LevelLetterCocycles/DataN.lean} 996--1007 defines the actual goodTuple and e0 as a finite sum of tensor products of the genuine eSer(f\_v).

Choose one admitted tuple whose tensor projector acts nontrivially on the actual pi. Such a tuple exists because the whole sum acts nontrivially; it is not required to be the special tuple used in LIFT059. At each selected v, the local representation contains the corresponding tau\_(f\_v). Sections 2--3 apply simultaneously at all these places. Conjugating to the prescribed split lattice frame changes K and its type together. A contragredient replaces every polynomial by its reciprocal and every cuspidal type by its dual; squarefreeness, distinct inertial classes and divisor degrees are preserved. The source also includes the inverse character predicates. No original finite type, vector, projector or level is changed.

\subsection{5. Identify the actual ramified parameter and eliminate every Arthur SL2}\label{identify-the-actual-ramified-parameter-and-eliminate-every-arthur-sl2}

Chen--Zou Theorem 2.1 and Proposition 4.1 attach to pi an elliptic global Arthur parameter

\[
                  \Psi=\boxplus_j\Pi_j\boxtimes S_{e_j},
                    \qquad\sum_j n_j e_j=r.
\tag{GT10}
\]

Weak transfer outside a finite set would not identify its localization at a selected G place. We instead use the exact auxiliary-theta comparison already established in LIFT069 Section 3, now with the tempered pi\_v proved above in place of its supercuspidal special case.

Choose the same sufficiently large quasi-split auxiliary unitary group and honest splitting characters in the local and global comparison. Chen--Zou Theorem 3.1 puts the abstract theta lift Sigma=theta\^{}abs(pi) in its discrete spectrum. Proposition 4.1 and its proof identify its global parameter as

\[
           \Theta(\Psi)=\Psi\chi_W^{-1}\chi_V
                          \boxplus\chi_V\boxtimes S_L.
\tag{GT11}
\]

At a selected split v the local formula (6.1) and Remark 6.4(2), valid for the now tempered phi\_v, identify the same local theta lift with the singleton GL packet for

\[
           \Theta(\phi_v)=\phi_v\chi_{W,v}^{-1}\chi_{V,v}
                          \oplus\chi_{V,v}\boxtimes S_L.
\tag{GT12}
\]

The global quasi-split packet statement in their Section 3.4 and Theorem 3.4 gives the other singleton description of Sigma\_v using GT11, including possibly unbounded localized parameters. Their associated-parameter convention (2.4) is used on both sides. Cancelling the common L-term character string with its full multiplicity, and then the common twist, gives

\[
                    \phi_{\Psi_v}^{W,ss}=\phi_v^{W,ss}.
\tag{GT13}
\]

This is cancellation of direct sums of semisimple Weil representations. It does not assert an a priori unramified equality at this ramified place, does not change a local vector, and does not insert an auxiliary arithmetic-chain edge.

If some e\_j\textgreater=2, any irreducible Weil constituent beta of phi\_(Pi\_j,v) contributes the e\_j constituents

\[
       beta|\cdot|^{(e_j-1)/2},\ldots,
       beta|\cdot|^{-(e_j-1)/2}
\tag{GT14}
\]

on the left of GT13. They are in the same inertial class and occur with positive multiplicity. Additional summands cannot cancel them. GT9 has each such class only once. This contradiction proves e\_j=1 for every j. Thus Psi is generic. Only now invoke Chen--Zou Theorem 2.6: the actual pi lies in the corresponding local Langlands packet at every place, and its occurrence in this global discrete packet is multiplicity one. This spectral assertion does not replace the source's coefficient or row-multiplicity records.

\subsection{6. The literal degree-divisor witnesses force one global block}\label{the-literal-degree-divisor-witnesses-force-one-global-block}

Suppose the now generic Psi has at least two global cuspidal blocks. One block has rank d with 1\textless=d and 2d\textless=r. At each selected split G place, the actual local parameter is the direct sum of the block parameters. Its pi\_v is therefore an irreducible subquotient of the corresponding normalized parabolic induction with a block of rank d.~One may use tempered induction here: pi\_v is tempered by Section 3, so its parameter summands are bounded. The weaker subquotient assertion is already sufficient.

The exact compact-invariant formula is

\[
 \left(\operatorname{Ind}_{P}^{GL_r}\bigotimes_j V_j\right)^{K_r^1}
       =\operatorname{Ind}_{\bar P}^{GL_r(k_v)}
                         \bigotimes_j V_j^{K_{n_j}^1}.
\tag{GT15}
\]

Iwasawa decomposition, compact modulus one and exact compact averaging prove GT15 as in DO046(AP12). Every finite constituent has a semisimple polynomial equal to the product of the polynomials of its inducing finite factors, by the separate Green orbit grading. In particular the actual f\_v has a divisor of degree d at every selected v. This proof requires no identification of an unspecified global block with a chosen finite parameter beyond the actual invariant formula.

The fifth conjunct of the literal IsGood predicate says

\[
 \forall d\ge1,\quad 2d\le r\ \Longrightarrow\
             \exists v:\ f_v\text{ has no divisor of degree }d.
\tag{GT16}
\]

Apply this to the single admitted tuple selected in Section 4, at either of its actual rational G primes. The witness v contradicts GT15. Hence there is exactly one block and

\[
                 \Psi=\Pi\boxtimes S_1,
             \qquad\Pi\text{ cuspidal on }GL_r(A_E).
\tag{GT17}
\]

Unlike LIFT069, no local polynomial is assumed irreducible of full degree. Unlike the original DO046 argument, genericity was proved using the complete squarefree inertial type rather than wide HC gaps. The source divisor clause is used in its original universal quantifier; it is not replaced by a hand-picked finite series.

\subsection{7. Actual coefficient parity and all-place temperedness}\label{actual-coefficient-parity-and-all-place-temperedness}

For GT2, literal TParity makes every h\_i integral. The exact source locators are \nolinkurl{Formalism/SharpRestriction/CuspSpanNormBounds/CayleySignature.lean} 106--108 and \nolinkurl{Formalism/LocalRingCardinality/DataTheta.lean} 529--543; \nolinkurl{Params.lean} 10--14 and 32 supplies the even-rank inputs and r=k+m-1. Interlace makes the displayed list strictly decreasing: successive v entries contain an intervening u; the last v exceeds u\_last\textgreater-b\_0; and b\_j\textless e\_j\textless b\_(j+1). Thus

\[
     h_i\in\mathbb Z=\mathbb Z+(r-1)/2,\qquad h_i-h_{i+1}\ge1.
\tag{GT18}
\]

The actual infinity type is discrete series, hence its complex parameters are bounded. The all-place packet identification after GT14 transfers these to the sole cuspidal block Pi. In the Chen--Zou/Mok conjugate-parity convention, the standard stable embedding is the chi=1 embedding of Mok (7.1.2)--(7.1.3). Therefore

\[
          \lambda_i=h_i-((r-1)/2-i),\qquad
          \lambda_i-\lambda_{i+1}=h_i-h_{i+1}-1\ge0
\tag{GT19}
\]

is an integral dominant coefficient weight. The conjugate embedding has the reversed negative HC list. A unit adjacent HC gap is allowed; it produces equal neighboring coefficient weights, not an irregular infinitesimal character. There is no auxiliary eta half-parity correction in this single-block case, since the block and unitary group have the same odd rank. This does not remove the honest splitting characters in GT11--GT12.

Exactly as in LIFT069 Section 4, Shin--Templier Lemma 2.12 and Remark 2.13 give cohomologicality for this cuspidal regular C-algebraic Pi. Caraiani Theorem 1.2 gives temperedness at every finite place. The actual local packets are therefore tempered. At split places, Jacquet--Shalika Section 2.1, printed page 110, gives genericity of every irreducible pre-unitary tempered GL\_r representation. This is a statement about the actual p-adic representation, not a converse inferred from finite genericity.

At T the source retains the original maximal finite level in the bottom variable. Its fixed spherical line now lies in a tempered generic representation for \textbf{every} existing e0 bottom constituent. LIFT068/071/072/073 keep their exact physical vector, D*tau translate, Gauss phases, quotient Haar and scalar normalization. The present classification theorem sets none of those scalars to one and supplies no global period value.

\subsection{8. Consequences and boundaries}\label{consequences-and-boundaries}

\begin{longtable}[]{@{}
  >{\raggedright\arraybackslash}p{(\columnwidth - 2\tabcolsep) * \real{0.5000}}
  >{\raggedright\arraybackslash}p{(\columnwidth - 2\tabcolsep) * \real{0.5000}}@{}}
\toprule\noalign{}
\begin{minipage}[b]{\linewidth}\raggedright
Question
\end{minipage} & \begin{minipage}[b]{\linewidth}\raggedright
Consequence established here
\end{minipage} \\
\midrule\noalign{}
\endhead
\bottomrule\noalign{}
\endlastfoot
Full e0 versus LIFT059's chosen tuple & Every existing cuspidal constituent of the genuine full e0 is covered \\
Selected split G components & Tempered from squarefree typicality and actual unitarity \\
Global Arthur SL2 blocks & All have length one, by the exact ramified comparison \\
Number of global cuspidal GL blocks & Exactly one, by the literal degree-divisor witnesses \\
Other finite places and original T line & Tempered; at split T, generic spherical with its original fixed vector \\
Global spectral multiplicity & One in the corresponding generic discrete packet; no one-dimensional row-space claim \\
Initial Mram representations & LIFT070 remains the separate source-size argument \\
Prescribed global distinction, Pu(T), Read/STiesR/032 & Not proved; LIFT074's distinguished-pair premise remains \\
Integral images, arithmetic values and DP44 costs & Unchanged and not inferred from the classification \\
\end{longtable}

\subsection{9. Primary anchors and review status}\label{primary-anchors-and-review-status}

Nadimpalli, S. (2017). \emph{Typical representations for level zero Bernstein components of GL\_n(F).} Journal of Algebra, 469, 1--29. \url{https://doi.org/10.1016/j.jalgebra.2016.08.033}. The directly read author version is dated August 9, 2019: \url{https://home.iitk.ac.in/~nsantosh/level.pdf} . Its Section 2.2, printed pages 2--3, defines the full inertial multiset and gives the irreducible-quotient construction; Section 3, pages 6--7, and Corollary 3.3 with its converse proof pin the complete parahoric type. The page locators refer to this author version. The retained PDF and text are \nolinkurl{source_recovery/nadimpalli_level_zero2019.pdf/.txt}; the author PDF SHA256 is d1a72fc8496503aef1cf30ade857971bbd9d9ef4d97b5f819f405a7bee33bc6d. No inference relies on the introductory shorthand about local class field theory.

Casselman, W. (1995, May 1). \emph{Introduction to the theory of admissible representations of p-adic reductive groups} {[}Draft{]}. \url{https://www.uni-muenster.de/imperia/md/content/theoretische_mathematik/topologie/casselma.pdf} . Theorem 5.2.1, Corollary 5.2.3 and Lemma 5.2.5, printed pages 46--48, were freshly read. This source is identified as the draft it is, not as a subsequently published final edition. The local copies are \nolinkurl{source_recovery/casselman1995_draft.pdf/.txt}.

Waldspurger, J.-L. (2003). \emph{La formule de Plancherel pour les groupes p-adiques. D'apres Harish-Chandra.} Journal of the Institute of Mathematics of Jussieu, 2(2), 235--333. \url{https://doi.org/10.1017/S1474748003000082}. Section I.3, p.243, Proposition III.2.2, pp.265--266, and Lemma III.2.3 with its proof, pp.266--267, were freshly read in the original publisher PDF. The retained copies are \nolinkurl{source_recovery/waldspurger2003_plancherel.pdf/.txt}. These supply the induction normalization and the exact temperedness statement, not merely a secondary summary of it.

Gow--Vinroot (2007), Section 3/Theorem 3.2: the directly read Green characteristic map and orbit grading, as in LIFT075 and DO046(AP12). Scholze (2013), Theorem 1.2(b)--(d): the directly read semisimple Weil correspondence, normalized induction, irreducibility and twists. Chen--Zou (2025), Theorems 2.1, 2.6 and 3.1, Section 3.4/Theorem 3.4, Proposition 4.1 and proof, formula (6.1), Theorem 6.3 and Remark 6.4(2): the actual auxiliary-theta comparison and all-place generic packet theorem, freshly reread at the used interfaces.

The final cohomological/tempered step is precisely the source and primary application already fully checked in LIFT069 Section 4: Mok Section 7.1, Shin--Templier Lemma 2.12/Remark 2.13 and Caraiani Theorem 1.2. Jacquet--Shalika Section 2.1 is the directly read tempered-to-generic statement of LIFT070. The actual source degree and parity interfaces are those stated in Sections 4 and 7. No same-sign-gap or abstract formal NTypes implication is substituted for the genuine realization.

The entire saved proof and companion \nolinkurl{all_good_bottom_temperedness_certificate.py} have received complete author review. Its 83,280 exact diagnostics cover finite factorization and reciprocal multiplicities, distinct-inertia exponent pairing and the repeated-class counterexample, overlapping theta-string cancellation, the literal divisor-degree contradiction, odd-rank coefficient parity including unit HC gaps, and the actual source/primary anchors. These diagnostics do not prove the cited representation theory or a global nonzero period. Final fingerprints and replay times are recorded in \nolinkurl{all_good_bottom_temperedness_verification.json} and \nolinkurl{LIFT076_PENDING_MANIFEST.json}. No independent PASS, canonical promotion or new Lean proof is asserted here.

\clearpage

\section{AM-R-070 --- integral short compression in the original scalar-cross sector at supports twenty-three and twenty-four}\label{am-r-070-integral-short-compression-in-the-original-scalar-cross-sector-at-supports-twenty-three-and-twenty-four}

\textbf{Canonical checked integration --- 2026-11-22T21:26:47.832456+00:00.} Whole root and independent proof/script/interface checks and separate9530-assertion replays passed. At actual all-top H133 supports23/24, the entire C-transverse scalar-cross CaseA and the plane-C subcase with at least one nonzero original2x2 short compression have combined absolute contribution\textless2\^{}-31703. The four integral pairings have conjugates\textless2\^{}47. The complete paired finite lemma restores a rank-three Green baseline through one whole-C ideal and its uniformly norm-bounded compression divisor, with no extra witness or compression-pattern factor. All original type/pair, selected-trace, target-column and source-fibre factors remain. No global unipotence is assumed. The all-four-zero plane compression subcase remains; complete floors23/11 and signed/global/arithmetic/DP44 gates remain open.

\textbf{Algebraic-geometric translation.} The new witness is a matrix of integral pairings between the original source and target Hermitian modules over Spec(O\_E). Its simultaneous vanishing at assigned CM primes is proved by paired nilpotent linear algebra, not by choosing a new integral frame or fixing the individual short vectors. The norm-divisor estimate and original fractional-ideal packing use the same complex Hermitian realization and actual selected trace. Exact contextual anchors remain SGA3 I4.7.3/4.7.3.1; SGA6 VI Theorem3.3; FGA232 Theorem3.1 pp148--149; FGA236 Corollaries1.7/2.3 pp229/232; SGA7 I VII2.7/2.9 pp169--174 and VIII7.4(b) pp310--311; cote4 scan003 §§8.1.1--8.1.2, scans005--006 §§8.4.1--8.4.3 and cote54 scans005--007. No new scheme comparison, lambda-ring/K-theory equality, l-adic sheaf/weight theorem, CM-motive identification or motivic L-function assertion is inferred. DP\_MULT Q3e and THETA3 G1/A7 gain this strict containing-set removal; CX\_GRH CX-ALG-024 and GCT\_PLETH T-052 keep their separate arithmetic comparisons.

Author: \nolinkurl{/root/definite_oscillator}, on the root's reserved R070 task. 22 November 2026 UTC. Complete proof; author whole-file verification, exact replay and frozen hashes are recorded separately before independent review and promotion.

\subsection{1. Exact scope and two subcases}\label{exact-scope-and-two-subcases}

Retain the unchanged J56 all-top kernel, its charged H=133 cutoff, f=168, n=169, m=164 and the original decomposition S perpendicular to W0, with signature(W0)=(4,1). Keep every original CM source/target pencil, selected finite projector, class mass, Haar normalization, actual packet point, compact twist and selected Hilbert metric. As in R059(GB1--10) and R061(GS1--5), write

\[
 h=\begin{pmatrix}I_m&B\\ C&D\end{pmatrix},\qquad C\ne0,
 \qquad N=h-I,\qquad J_0=\operatorname{diag}(1,1,1,1,-1).
 \tag{SC1}
\]

The existing remainder has unipotent reductions at all334 unmixed G places, by the retained K068 structural cut. No global unipotence of h is assumed. The two mixed factors retain modulus at most one.

Use exactly the original integers

\[
\begin{split}
 Q&=1427908193729121522601,\qquad U=77Q^2-14,\\
 u&=(2,95192390527,383092671043,12529843548438452388609),\\
 y_1&=(-976627,-153032,38026,0),\\
 y_2&=(-567678,695570,-172838,0),\\
 y_3&=(104583,-15309315176,-61610045224,2).
\end{split}                                                       \tag{SC2}
\]

The y\_j are positive W0 vectors with zero negative coordinate. They span u-perp in its four positive coordinates, u*u=U and gcd(u\_i)=1. R059(GB3--8) proves rank\_E C=1 from unitarity and signature one, and supplies this exhaustive division:

\begin{itemize}
\tightlist
\item
  \textbf{Case A:} at least one row delta\_i=\textless y\_j,Ce\_i\textgreater, j in\{1,2,3\}, is nonzero.
\item
  \textbf{Plane case:} all those rows vanish, so uniquely Ce\_i=(u xi\_i,beta\_i), with xi nonzero.
\end{itemize}

Inside the plane case define the four original short compressions

\[
             \Delta_{ij}=\langle y_i,(h-I)y_j\rangle,
                         \qquad1\le i,j\le2.                    \tag{SC3}
\]

\textbf{The theorem proved here:} the union of Case A and the plane case with at least one nonzero Delta\_ij, at supports s=23 and s=24, has complete absolute contribution less than 2\^{}−31703. The plane case with all four compressions zero is not covered. In particular, no full support-floor update follows from this theorem alone.

\subsection{2. Actual integrality and the short field-norm witness}\label{actual-integrality-and-the-short-field-norm-witness}

Every original e\_i in S and every y\_j in SC2 belongs separately to every actual source and target pencil at both CM components. This is the original W0-supported shear inclusion and DO078(FP20--22), retained in R059(GB4--5). If h is an actual arrow from its prescribed source to target, hy\_j belongs to the target and y\_i belongs to that same paired integral selfdual target. Hence

\[
            \Delta_{ij}\in\mathcal O_E.                          \tag{SC4}
\]

No h-stable common lattice is asserted. The subtraction of \textless y\_i,y\_j\textgreater{} is integral. At each complex embedding, the original Euclidean operator norm of h is at most H, so

\[
 |\tau\Delta_{ij}|
 \le(H+1)\|y_i\|\|y_j\|
 \le134\cdot978665066829<2^{47}.                                \tag{SC5}
\]

Here \textbar\textbar y1\textbar\textbar²=978665066829 and \textbar\textbar y2\textbar\textbar²=835948910828. Therefore any paired G ideal J\_Delta dividing a chosen nonzero Delta\_ij obeys

\[
 N_EJ_\Delta\le|N_{E/\mathbf Q}\Delta_{ij}|<2^{47\cdot336}.       \tag{SC6}
\]

The next section proves simultaneous local vanishing of all four entries at the primes assigned to J\_Delta. Thus the same assigned ideal divides whichever global entry is nonzero. No independent four-witness multiplicity is needed in this pointwise bound.

\subsection{3. The complete finite paired-compression lemma}\label{the-complete-finite-paired-compression-lemma}

Fix one actual completely split unmixed G place. At its two CM components write h and h\_c over the same residue field, in the original form matrix J=diag(I,J0). Their relation is

\[
                         h_c=J h^{-\mathsf t}J.
\]

Both S blocks are I. The ranks of h−I and h\_c−I coincide; denote this finite rank by d.~Let C\_c be the W0-by-S block of h\_c.~No complex isotropic frame is used at this finite place.

The Schur complement E=D−CB is invertible. Block inversion and paired unitarity give

\[
 -BE^{-1}=(h^{-1})_{SW_0}=C_c^{\mathsf t}J_0,
 \qquad B=-C_c^{\mathsf t}J_0E.
 \tag{SC7}
\]

Consequently B is nonzero exactly when C\_c is nonzero. In the plane case, the positive parts of both C and C\_c lie in the line spanned by the original u. Since Y=span(y1,y2) is positive and orthogonal to u, reduction gives

\[
                 Y^{\mathsf t}J_0C=0,
                         \qquad C_c^{\mathsf t}J_0Y=0.          \tag{SC8}
\]

No independence of the reduced y\_i is required.

\subsubsection{Finite rank two, both C blocks nonzero}\label{finite-rank-two-both-c-blocks-nonzero}

A nilpotent matrix N of rank two satisfies N³=0. Its paired S blocks are zero, and (h\^{}−1)SS=I. Expanding h\^{}−1=I−N+N² gives (N²)SS=0. Suppose N²S is nonzero, and choose v in S with N²v nonzero. The vectors Nv and N²v are independent, since applying N to a dependence first forces the coefficient of Nv to vanish. Both lie in W0. They span the two-dimensional image of N, which would force B=0, contradicting SC7. Thus

\[
                              N^2S=0.                           \tag{SC9}
\]

If N²=0, the SW0 block of h\^{}−1 gives B=−C\_c\^{}tJ0. If N² is nonzero, the rank-two nilpotent has one size-three block, so im N intersect ker N=im N² is one-dimensional. The nonzero line C(S)=N(S) belongs to this intersection by SC9. Hence im C=im N² lies in W0, and (N²)SW0=0. Again B=−C\_c\^{}tJ0.

In either case SC8 gives BY=0. Projection of im N onto S is nonzero because B is nonzero, whereas im N has dimension two and contains the nonzero line im C inside W0. Therefore im N intersect W0=im C. For y in Y, Ny lies in this intersection, and pairing with any member of Y annihilates im C. Thus

\[
                         Y^{\mathsf t}J_0N_{W_0W_0}Y=0                  \tag{SC10}
\]

at both CM components. This includes both square-zero and size-three Jordan cases.

\subsubsection{Finite rank one and zero}\label{finite-rank-one-and-zero}

If d=1, then N²=0. Factor N=v lambda. If C is nonzero, lambda\_S is nonzero and N\_SS=0 forces v\_S=0. Hence B=0, im N=im C is contained in W0, and every left Y-pairing vanishes by SC8. The paired identity N\_c=−J N\^{}tJ gives zero Y compression at the companion as well. If instead C\_c is nonzero, exchange the two components. Both C blocks cannot be nonzero in rank one. When d=0, all four compressions vanish at both components.

If both C blocks vanish in rank one or two, the proof imposes no compression condition. Also, SC10 is only a compression statement: Ny need not be zero. For instance, with S the first coordinate and J0=diag(1,1,−1), the matrix

\[
N=\begin{pmatrix}0&0&0&1\\1&0&-1&0\\0&0&0&1\\0&0&0&0\end{pmatrix}
\]

has rank two and N²=0; Y=span(e3) is perpendicular to both paired C lines but Ne3=−e2 is nonzero. This prevents an incorrect pointwise-fixation consequence.

\subsection{4. Exact ideal assignment and Green restoration}\label{exact-ideal-assignment-and-green-restoration}

At each of the334 unmixed places use the following deterministic assignments. ``Paired'' means both E-primes, with norm q\_v²; ``one'' means the indicated CM prime, with norm q\_v. The original whole C block is divisible by every assigned factor of J\_C. The finite lemma assigns factors of J\_Delta only where all four compressions vanish.

\begin{longtable}[]{@{}
  >{\raggedleft\arraybackslash}p{(\columnwidth - 10\tabcolsep) * \real{0.1905}}
  >{\raggedright\arraybackslash}p{(\columnwidth - 10\tabcolsep) * \real{0.1429}}
  >{\raggedright\arraybackslash}p{(\columnwidth - 10\tabcolsep) * \real{0.1429}}
  >{\raggedright\arraybackslash}p{(\columnwidth - 10\tabcolsep) * \real{0.1429}}
  >{\raggedleft\arraybackslash}p{(\columnwidth - 10\tabcolsep) * \real{0.1905}}
  >{\raggedleft\arraybackslash}p{(\columnwidth - 10\tabcolsep) * \real{0.1905}}@{}}
\toprule\noalign{}
\begin{minipage}[b]{\linewidth}\raggedleft
Finite rank d
\end{minipage} & \begin{minipage}[b]{\linewidth}\raggedright
Zero C blocks
\end{minipage} & \begin{minipage}[b]{\linewidth}\raggedright
J\_C contribution
\end{minipage} & \begin{minipage}[b]{\linewidth}\raggedright
J\_Delta contribution
\end{minipage} & \begin{minipage}[b]{\linewidth}\raggedleft
Restoration exponent
\end{minipage} & \begin{minipage}[b]{\linewidth}\raggedleft
Required deficit
\end{minipage} \\
\midrule\noalign{}
\endhead
\bottomrule\noalign{}
\endlastfoot
0 & Both & Paired & Paired & 501 & 501 \\
1 & Exactly one & That component & Paired & 334 & 333 \\
1 & Both & Paired & None & 334 & 333 \\
2 & Neither & None & Paired & 167 & 166 \\
2 & Exactly one & That component & None & 167 & 166 \\
2 & Both & Paired & None & 334 & 166 \\
At least3 & Any & None & None & 0 & 0 \\
\end{longtable}

The retained normalized full-degree cuspidal Green exponents at finite ranks zero, one, two and three are 0,168,335,501. Higher unipotent ranks have at least the last decay. The same global Euler-product comparison used in R059/R061 costs one factor2, and the two mixed factors have modulus at most one. The table therefore proves

\[
 |a(h)|\le2Q^{-501\cdot167}(N_EJ_C)^{167}
                                   (N_EJ_\Delta)^{167/2}.       \tag{SC11}
\]

For example, rank one with one vanishing C block is restored by 167+167=334, which exceeds the required501−168=333. All other rows are checked in the same exact table.

In the present nonzero-compression subcase, SC6 bounds the final numerator by 2\^{}(47·167f). There are at most4\^{}334=2\^{}668 choices of J\_C. There is no separate J\_Delta pattern count: for each J\_C, its possibly varying J\_Delta has already been bounded uniformly using a nonzero actual entry. The whole finite rank distribution beyond these choices is not enumerated again.

\subsection{5. Plane coordinates, actual trace and complete source fibre}\label{plane-coordinates-actual-trace-and-complete-source-fibre}

Put L=J\_C I\_Sigma, N\_E I\_Sigma=T\^{}−2s and z=(N\_E L)\^{}(1/(2f)). The R061 original-coordinate proof applies to this stronger whole-C ideal unchanged:

\[
 \xi_i,\beta_i\in L,
 \quad Ce_i=(u\xi_i,\beta_i),
 \quad \|\xi_\tau\|<M:=H/\sqrt U,
 \quad |\beta_{j_0,\tau}|<H.                                   \tag{SC12}
\]

At T, Ce\_i=he\_i−e\_i belongs to the actual target pencil, contained in pi\^{}−1 Lambda at each CM partner. Primitivity of u gives xi in the same original ideal by an integer Bezout functional. Away T, the assigned G divisibility and the original integral conditions apply. No type-II-only inclusion or zero-target exclusion is used.

If exactly k coordinates of xi are nonzero, 1≤k≤164, their original norm minima imply kz²\textless M². Fix the first nonzero position j0. The row xi and the one scalar beta\_j0 reconstruct C by

\[
                  Ce_i=(u\xi_i,\beta_{j_0}\xi_i/\xi_{j_0}).
\]

The quotient is for injective reconstruction over E, not an integral change of basis. Counting the164 possible pivots is harmless. The entire coordinate count is bounded by

\[
 \binom{164}{k}(1+2M/z)^{2fk}(1+2H/z)^{2f}.                    \tag{SC13}
\]

The true packet representative fixes S and is supported on W0, by LIFT41(NP7) and R061(GS11). Its radius is less than1/2. The same complex-only isotropic normalization yields

\[
 \sinh r'_\tau>\tfrac12\sqrt U\,\|\xi_\tau\|.
\]

DO080's actual selected normalized-trace estimate, with both compact factors retained, consequently gives the original R061(GS12) bound

\[
                    |c(h)|<2^{-929f/2}U^{-f/2}z^{-f}.           \tag{SC14}
\]

This is not an operator-norm estimate. It requires no global nilpotence. The argument uses the field-norm lower bound for one nonzero xi coordinate only after taking the product over all f labels; it does not claim a pointwise lower bound at each label.

For fixed C and v=D e\_n, unitarity reconstructs the full negative column by (he\_n)S=−C*J0v. Every arrow fixes the same163-dimensional positive ker C. The ratio of two arrows with these data and the same source/target pencils preserves the original source lattice and fixes ker C and e\_n.~Its faithful positive complement has rank five. The accepted actual projective-lattice/reduction argument R061(GS13--14) gives exactly the same bound

\[
 G_5=98148490017978252769652736000000<2^{107}.                    \tag{SC15}
\]

It uses no globally integral orthogonal basis. The remaining v is counted in the true pi\^{}−1 M\_(target,W0), uniformly in every paired label including zero. Its norm is below95 sqrt(f), its paired index is T\^{}(10s), and the original-ring covering bound gives

\[
                         \#v<2^{5757}T^{10s}.                  \tag{SC16}
\]

This is the original full five-complex-coordinate ball; no real affine coarea assertion is used.

\subsection{6. The exact weighted endpoint for the plane subcase}\label{the-exact-weighted-endpoint-for-the-plane-subcase}

Since N\_EJ\_C=T\^{}(2s)z\^{}(2f), the numerator (N\_EJ\_C)\^{}167 contributes T\^{}(334s)z\^{}(334f). Combining SC11--14, the z-dependent factor is

\[
 z^{(331-2k)f}(z+2M)^{2fk}(z+2H)^{2f}.                          \tag{SC17}
\]

Its leading exponent is at least3f, so it increases with z\textgreater0. At the nonzero-coordinate cap M/sqrt(k), and after using 1+2sqrt(Uk)≤3sqrt(Uk), its full H/U/k factor is

\[
 3^{2f}H^{333f}U^{-166f}\binom{164}{k}
                    k^{-(331/2)f}(1+2\sqrt k)^{2fk}.           \tag{SC18}
\]

All remaining multiplicities are retained once: the Green factor2, J\_C patterns2\^{}668, support choices2\^{}168, pivot choices164\textless2\^{}8, the fibre2\^{}107 and column constant2\^{}5757. The total is2\^{}6709. Every actual source/target type and weight is counted by the original full sum

\[
       P_T^s=\left(2(T-1)^2/(T-2)\right)^s<2^{997s}.            \tag{SC19}
\]

There is no additional3\^{}s, delta witness, delta pattern or endpoint multiplicity. The T power is334s+10s=344s. Retain the exact bounds

\[
 L_T=1019333/1024>\log_2T,\quad L_Q=287843/4096<\log_2Q,
\quad L_H=28900/4096>\log_2H,\quad L_U=601356/4096<\log_2U,
\quad3^{336}<2^{533}.
\]

For each k let c\_k=ceil(sqrt(k)) and let e\_k be the least integer with

\[
 \binom{164}{k}(1+2c_k)^{336k}/k^{27804}<2^{e_k}.               \tag{SC20}
\]

Here27804=(331/2)f.~The complete rational exponent is C(s)+e\_k, where

\[
\begin{split}
 C(s)={}&6709+533-929f/2+333fL_H-166fL_U+47\cdot167f\\
        &+(344L_T+997)s-501\cdot167L_Q.
\end{split}                                                       \tag{SC21}
\]

The exact companion checks all164 integers e\_k. Their maximum is57444 at k=164. The two complete maxima are

\[
 C(23)+57444=-1536583577/4096,
 \qquad C(24)+57444=-129897657/4096<-31713.                        \tag{SC22}
\]

No automatic extension to support25 is made; the same upper envelope there fails to be small.

\subsection{7. Case A with every original transverse and column factor}\label{case-a-with-every-original-transverse-and-column-factor}

The C-transverse case needs none of the new short-compression condition. Its entire support-variable count is R059(GB4--21), which is valid without global unipotence or a restriction on s. For clarity, it retains the nonzero globally integral row \textless y\_j,Ce\_i\textgreater, its original G ideal, at most3 transverse choices and164 pivot choices, and the two actual W0 columns. Each column count is2\^{}14861 T\^{}(10s), and the full original source fibre is2\^{}5022. The original Green exponent is335 and the selected real coefficient is bounded by one. The complete original constant is

\[
 1+668+168+8+2+5022+2\cdot14861=35591.
\]

The coordinate archimedean bound is2\^{}43, and elementary actual-ideal packing gives2\^{}(43·335f)3\^{}(328f). The complete exact exponent is therefore

\[
 E_A(s)=35591+43\cdot335f+(8/5)328f
                 +(20L_T+997)s-335\cdot167L_Q.                  \tag{SC23}
\]

At the two new supports it equals

\[
 E_A(23)=-18572421943/20480,
 \qquad E_A(24)=-18144270183/20480<-885950.                       \tag{SC24}
\]

Every source/target label and actual weight is included before this summation. In particular, the pivot and negative columns are counted in their true target overlattice, not a substituted diagonal fractional ideal. No improved rank-five fibre or real trace is silently used in this Case A bound.

\subsection{8. Full union, chronology and limitations}\label{full-union-chronology-and-limitations}

There are328 plane cells and two entire transverse cases. By SC22--24 their union is bounded by

\[
           328\,2^{-31713}+2\,2^{-885950}<2^{-31703}.             \tag{SC25}
\]

The deliberately conservative charge31703 is appended once for intersection with the actual current remainder. The predecessor is \nolinkurl{checks/do092_joint_remainder_error.json}, with73 all-top and31 uniform entries, ending in23546,516. Append31703 to obtain74/31. The complete exact dyadic sums remain below2\^{}−230. Every older charge is retained; existing globally unipotent subcase removals are not erased or promoted into new hypotheses.

The complete all-top support floor remains23 and the uniform floor remains11 until separate exhaustion is proved. The plane subcase with all four short compressions zero remains unbounded by this theorem. No pointwise fixation of Y follows from its local compression lemma. No original field, lattice, finite label, class mass, Haar, compact twist or selected metric has been altered. The sufficient signed inequality, genuine prescribed global Bessel/Mram value, selected arithmetic image, literal Read/STiesR/032 construction and DP44 comparisons remain open.

The primary finite Green realization and exact rank envelopes are those retained in K067 and R059/R061; their application here is the new explicit restoration table. R059 supplies the original arbitrary-arrow integrality and support-variable transverse count; R061 supplies the full plane coordinate reconstruction, actual packet trace, projective rank-five fibre and full target ball. The source/target modules are the unchanged Hermitian modules over Spec(O\_E). No settled theorem is re-proved, and finite diagnostics do not replace the general paired linear-algebra or counting arguments above.

\clearpage

\section{AM-DO-093 --- inverse-transverse scalar-cross removal at supports twenty-three and twenty-four}\label{am-do-093-inverse-transverse-scalar-cross-removal-at-supports-twenty-three-and-twenty-four}

\textbf{Canonical checked integration --- 2026-11-22T21:33:18.417041+00:00.} Whole root and independent proof/script/interface checks and separate7726-assertion replays passed. In the actual K068 finite-unipotent remainder at all-top H133 supports23/24, the plane-C but inverse-C-transverse scalar-cross terms have total absolute contribution\textless2\^{}-885948. Literal inversion conjugates the actual normalized summand, preserves the signed pencil coefficient, all source/target weights, central/class/Haar normalizations and the cutoff, and injects this set into the complete original transverse CaseA count. Together with R070 the remaining scalar-cross terms have both cross blocks in the packet plane and all four short compressions zero. They fix an E-defined positive165-dimensional subspace and have global rank(h−I)\textless=4; individual test vectors need not be fixed and no global unipotence is assumed. Complete floors23/11, signed/global/arithmetic and DP44 gates remain open.

\textbf{Algebraic-geometric translation.} The inversion is the original Hermitian adjoint inverse on the same source/target modules over Spec(O\_E), and the exact T correspondence is transposed together with its symmetric projector coefficients. The residual fixed space is the kernel of an E-linear rank-one map on the original positive166-space; no integral moving basis or controlled-denominator isometry is constructed. Exact contextual anchors remain SGA3 I4.7.3/4.7.3.1; SGA6 VI Theorem3.3; FGA232 Theorem3.1 pp148--149; FGA236 Corollaries1.7/2.3 pp229/232; SGA7 I VII2.7/2.9 pp169--174 and VIII7.4(b) pp310--311; cote4 scan003 §§8.1.1--8.1.2, scans005--006 §§8.4.1--8.4.3 and cote54 scans005--007. No new scheme comparison, lambda-ring/K-theory identity, l-adic sheaf/weight theorem, CM-motive comparison or motivic L-function theorem follows. DP\_MULT Q3e and THETA3 G1/A7 gain this containing-set removal and explicit positive fixed-space reduction; CX\_GRH CX-ALG-024 and GCT\_PLETH T-052 retain their separate arithmetic gates.

Author: ALP\_MACHINE \nolinkurl{/root/definite_oscillator}. 22 November 2026 UTC.

Controlled pending proof. Author verification and independent review are recorded separately; no canonical file is changed by this report.

\subsection{1. Exact sector and conclusion}\label{exact-sector-and-conclusion}

Retain the original all-top J56 kernel, all allowed words and canonical twists, the actual selected Hilbert metric, the charged cutoff H=133, the actual finite projectors and every source/target T-pencil label. Put f=168, n=169, m=164, and use the original orthogonal decomposition S perpendicular W0, with W0 of signature (4,1). The actual current remainder has unipotent reductions at every334 unmixed G factor by K068(LR20). No global unipotence is assumed.

Write the unique determinant-one central representative in original coordinates as

\[
 h=\begin{pmatrix}I_m&B\\ C&D\end{pmatrix},\qquad C\ne0,
 \qquad h^{-1}=\begin{pmatrix}I_m&B^-\\ C^-&D^-\end{pmatrix}.
 \tag{IV1}
\]

Use the original packet plane and integral transverse vectors of R059(GB4):

\[
\begin{split}
 u&=(2,95192390527,383092671043,12529843548438452388609),\\
 y_1&=(-976627,-153032,38026,0),\\
 y_2&=(-567678,695570,-172838,0),\\
 y_3&=(104583,-15309315176,-61610045224,2).
\end{split}                                                       \tag{IV2}
\]

Their negative W0 coordinates are zero. They span the positive orthogonal complement of u, and u*u=U=77Q\^{}2−14 for the original Q=1427908193729121522601. A cross block is called transverse here precisely when at least one row (\textless y\_j,C e\_i\textgreater)\_i is nonzero. Otherwise its positive four-coordinate part lies in E u, the original packet plane; this is an algebraic description, not a new lattice frame.

At each of the two supports s=23,24, let X\_s consist of actual scalar-cross matrices IV1 for which C is in that plane and C\^{}- is transverse. One may additionally intersect X\_s with the current remainder, with the four-short-compression-zero condition, or with any other already retained condition. Then

\[
 \sum_{s=23}^{24}\ \sum_{h\in X_s}
       |\text{actual normalized contribution of }h|
                         <2^{-885948}.                           \tag{IV3}
\]

There is one new all-top charge885948. The uniform statement is unchanged. This result alone changes neither support floor, which remain23 and11. Its substantive consequence, combined with R070, is that the scalar-cross terms at these two supports may be restricted to those with both C and C\^{}- in the original packet plane and all four short compressions zero. It does not assert that the two short vectors are fixed pointwise.

\subsection{2. The literal normalized coefficient under inversion}\label{the-literal-normalized-coefficient-under-inversion}

The exact coefficient is K062(AT14), with the same fixed real representative g\_p for every scalar-orbit index:

\[
 a(h)=\frac{\operatorname{Tr}(Q_{\rm good}R(h_G))}{d_R},\qquad
 c(h)=\frac{\operatorname{Tr}_W(P_W\pi(g_p^{-1}h_\infty g_p)|_W)}{d_W}.
 \tag{IV4}
\]

K058(PG3) defines R as the actual unitary finite representation with its probability metric and Q\_good as the orthogonal projection onto the genuine good subspace, including the two actual mixed mirabolic projectors. Thus Q\_good*=Q\_good. No commutation of this projector with R(h\_G) is needed. Cyclicity of the finite trace gives

\[
 \operatorname{Tr}(Q_{\rm good}R(h_G^{-1}))
 =\operatorname{Tr}(Q_{\rm good}R(h_G)^*)
 =\overline{\operatorname{Tr}(Q_{\rm good}R(h_G))}.
 \tag{IV5}
\]

The same argument applies to the finite compression defining c.~Equivalently, sum the diagonal matrix coefficients over an orthonormal basis of W. Unitarity of pi gives

\[
              a(h^{-1})=\overline{a(h)},\qquad
              c(h^{-1})=\overline{c(h)}.                         \tag{IV6}
\]

The fixed g\_p is unchanged: g\_p\^{}\{-1\}h\^{}\{-1\}g\_p is exactly (g\_p\^{}\{-1\}h g\_p)\^{}\{-1\}. No compact factor is discarded and no radial coefficient or operator-norm surrogate is used. The finite reduction is read in the same original frame. Away T, the zero extension of a is also respected because each prescribed maximal compact is a group and is closed under inversion.

For clarity, the precise normalized scalar-orbit identity is K062(AT13--16, AT24). Before its final positive normalization it has the factor

\[
 \frac{w_x z_0 C_\infty d_Rd_W}{(T-1)^f},\qquad
 w_x=v_P/|\Gamma_x|,\quad z_0=28,
 \tag{IV7}
\]

and after normalization its matrix summand is

\[
              \mathcal W(h)=a(h)c(h)\prod_{v\mid T}\rho(h_v).
 \tag{IV8}
\]

The factor in IV7 is fixed for the orbit. The strict-class descent is already included in w\_x/(T−1)\^{}f, with the genuine coefficient-bundle metric of K062. Inverting the rational summation index does not change this class mass, the central factor, the finite probability Haar, the raw real Haar, C\_infinity, d\_R or d\_W. The argument does not replace the bundle norm by a scalar generic-level norm.

\subsection{3. Exact pencil weights and the source/target swap}\label{exact-pencil-weights-and-the-sourcetarget-swap}

At each relative T place let q=T, let D=k minus \{−1\} be the admitted set of q−1 labels, and retain the actual paired pencil L\_c of K062(AT3--6). In particular c=0 is included; the omitted points are infinity and−1. The conjugate CM component is determined by the inverse-adjoint relation and supplies no independent label.

Let m\_T=T−1. K062(AT7--8) defines

\[
 (M_h)_{d,c}=\mathbf1_{hL_c=L_d},\qquad
 P=I-\frac1{m_T}J,\qquad
 \rho(h)=\frac{\operatorname{Tr}(PM_h)}{m_T-1}.
 \tag{IV9}
\]

These entries really are1: the actual local T types are maximal-compact fixed, so there is no suppressed coefficient character or phase. For every admitted ordered pair,

\[
 hL_c=L_d\quad\Longleftrightarrow\quad h^{-1}L_d=L_c.
 \tag{IV10}
\]

Consequently M\_(h\^{}-1)=M\_h\^{}t. Since P is real symmetric, both the exact signed coefficient and its absolute value satisfy

\[
                   \rho(h^{-1})=\rho(h)\in\mathbb R.
 \tag{IV11}
\]

This holds before any absolute-value estimate. More explicitly, the coefficient of the arrow (c,d) in the trace expansion is P\_(c,d)/(m\_T−1), equal to the coefficient of (d,c). Thus inversion is a measure-preserving bijection of the full ordered incidence enumeration, including zero labels, diagonal arrows and off-diagonal arrows. Both CM components are swapped together, with no squared multiplicity or new partner factor.

The coarse complete per-place envelope used in R059 and R070 is

\[
 P_T=\frac{2(T-1)^2}{T-2}<2^{997}.                              \tag{IV12}
\]

It is the same bound on the same ordered pairs after the swap: the number of possible labels is unchanged, and the exact rho in IV11 is unchanged. IV12 is a conservative counting envelope, not a replacement definition of the signed trace. In particular this argument neither assumes type-II positivity nor imports a target-coordinate slice whose covolume could change at label zero. The whole target-column counts used below were already uniform in the labels.

Combining IV6 and IV11 gives the exact pointwise identity

\[
              \mathcal W(h^{-1})=\overline{\mathcal W(h)},
              \qquad |\mathcal W(h^{-1})|=|\mathcal W(h)|.
 \tag{IV13}
\]

\subsection{4. Original support, cutoff and central indexing}\label{original-support-cutoff-and-central-indexing}

The exact support is closed under inversion. Away T, this follows from maximal-compact membership. At T, existence of an admitted source/target arrow is equivalent to existence of its reversed arrow by IV10. The number s of nonintegral relative T places is unchanged: h stabilizes the original standard lattice if and only if h\^{}-1 does. This is maximal-lattice integrality; the actual determinant is a unit, so it agrees with the supported entry-integrality convention. The equal-full-pencil, unequal-fixed and unequal-off-diagonal incidence types are individually preserved by reversal.

At every complex label the form is the original J=diag(I\_a,−1), which is Euclidean unitary, and

\[
                    h^{-1}=Jh^*J,
                    \qquad\|h^{-1}\|_{\rm op}=\|h\|_{\rm op}.
 \tag{IV14}
\]

Hence the actual cutoff H=133 is preserved exactly. The same identity for the S block gives (h\^{}-1)\_SS=I. The lower cross block C\^{}- is nonzero whenever C is nonzero. For example, if C\^{}-=0 then the paired upper block B is zero; the invertible triangular block form and unitarity then force C=0. Thus inversion remains within the scalar-cross class and does not send it to the fixed-W0 class.

At an unmixed G place inversion preserves unipotence and the rank of h−I, since h\^{}-1−I=−h\^{}-1(h−I). These are the actual finite premises of the original Case-A bound. No global-unipotence property is introduced or needed.

The rational central quotient is unchanged. The center is central, so inversion is well-defined on its cosets. More concretely, K068(LR3) chooses the unique determinant-one representative h(g)=(det g)\^{}−1 g, using the already fixed root group mu\_28 and169 congruent1 modulo28. Directly

\[
                         h(g^{-1})=h(g)^{-1}.                   \tag{IV15}
\]

The central characters already cancel in the literal kernel K062(AT15); there is no new phase or factor28. Inversion is a bijection of the original entire support at each s. The previously cut remainder need not itself be inverse-closed: below we bound its image by the entire original transverse sector, which is a valid enlargement of an absolute sum.

\subsection{5. Complete transverse count and the added error}\label{complete-transverse-count-and-the-added-error}

For each s, inversion maps X\_s injectively into the entire original C-transverse Case-A sector at support s. Equations IV13--15 and the exact label swap prove

\[
       \sum_{h\in X_s}|\mathcal W(h)|
       \le\sum_{k\text{ in complete transverse Case A at }s}
                                                |\mathcal W(k)|.
 \tag{IV16}
\]

There is no factor two in IV16: it is an injection into one complete sector, not a choice between h and h\^{}-1 at each embedding. No change of source frame or fibre normalization is needed. For the inverse matrix, R059's count uses its actual source lattice, which is precisely the original target; its original target-column and source-fibre bounds are uniform for that source/target pair.

The full Case-A bound is R059(GB4--21), evaluated at arbitrary support as recorded and freshly certified in R070(SC23--24). It retains all support choices2\^{}168, common G patterns2\^{}668, three transverse witnesses,164 pivot choices, two original W0 target-column counts each2\^{}14861 T\^{}(10s), the full source fibre2\^{}5022, the single Green correction2, the complete P\_T\^{}s, and \textbar c\textbar≤1. Its constant is

\[
 1+668+168+8+2+5022+2\cdot14861=35591.
\]

No rank-five replacement, sharper trace estimate, or omitted common-ideal factor is used. With f=168 and the exact actual one-sided prime bounds

\[
 L_T=1019333/1024>\log_2T,\qquad
 L_Q=287843/4096<\log_2Q,
\]

the complete exponent is

\[
 E_A(s)=35591+43\cdot335f+\frac85\,328f
                     +(20L_T+997)s-335\cdot167L_Q.             \tag{IV17}
\]

The factor8/5 is the certified upper bound for log\_2(3) in the original coordinate packing. Exact rational arithmetic gives

\[
 E_A(23)=-18572421943/20480,
 \qquad E_A(24)=-18144270183/20480<-885950.                      \tag{IV18}
\]

Therefore the two full support bounds, and hence IV3, have sum strictly below2·2\^{}−885950=2\^{}−885949\textless2\^{}−885948. The last inequality leaves the deliberately conservative single requested charge. There is no extra factor for compression entries or for inverting a choice of witness.

The predecessor is the canonically checked R070 error list,74 all-top and31 uniform entries, ending23546,516,31703. Append885948 once, giving75/31; every prior charge remains. The exact certificate checks both full dyadic sums below2\^{}−230. It does not claim that a numerical diagnostic evaluates an automorphic contribution.

\subsection{6. The precise residual and unchanged open gates}\label{the-precise-residual-and-unchanged-open-gates}

R070 has already bounded the entire original transverse-C sector and the plane-C subcase with a nonzero entry of the short2-by-2 compression. IV3 additionally bounds original-plane C with transverse C\^{}-. Thus at supports23/24 the residual scalar-cross terms may be restricted to

\[
 C,C^-\text{ both in the original packet plane},\qquad
 \langle y_i,(h-I)y_j\rangle=0\quad(1\le i,j\le2).              \tag{IV19}
\]

For these real positive y\_i, unitarity gives the exact identity

\[
 \langle y_i,(h^{-1}-I)y_j\rangle
   =\overline{\langle y_j,(h-I)y_i\rangle}.                    \tag{IV20}
\]

Accordingly the four-zero condition is already invariant under inversion. No second inversion of the R070 short-compression estimate is required. Equation IV20 is an equality of compressions only. A vector can move by a nonzero isotropic vector orthogonal to this positive test space, so IV19 does not imply pointwise fixation of the individual y\_i.

There is nevertheless a precise smaller-dimensional fixed space. Put N=h−I and Y=span\_E(y\_1,y\_2). The two displayed integer vectors are independent, so S perpendicular Y is a positive166-dimensional subspace in the original Hermitian space. From C\^{}-=J0 B* and the plane condition on C\^{}-, one gets BY=0. The plane condition on C gives C*J0Y=0. Thus NY is contained in W0, and unitarity together with the four zero compressions gives

\[
 \langle Ny,Ny'\rangle
 =-\langle y,Ny'\rangle-\langle Ny,y'\rangle=0
                         \quad(y,y'\in Y).                    \tag{IV21}
\]

Also, for s in S and y in Y,

\[
 0=\langle hs,hy\rangle
   =\langle Cs,y+Ny\rangle
   =\langle Cs,Ny\rangle.
\]

The nonzero line CS is isotropic, NY is totally isotropic by IV21, and they are mutually orthogonal. Their sum is therefore a totally isotropic subspace of W0. At a faithful complex embedding W0 has Witt index one, so this sum has dimension one. The same rank holds over E. Consequently NY is contained in CS. As NS=CS is already nonzero, N restricted to S perpendicular Y has rank exactly one. It follows that

\[
 \dim_E\bigl(\ker(h-I)\cap(S\oplus Y)\bigr)=165,
 \qquad \operatorname{rank}_E(h-I)\le4.                       \tag{IV22}
\]

This is an E-defined positive subspace fixed pointwise. It may depend on h, and no integral basis, rational isometry with controlled denominator, or prescribed coordinate fixation is asserted. Its nondegenerate orthogonal complement has dimension four and signature (3,1). The argument imposes no global unipotence; the residual may still have nontrivial semisimple and unipotent parts there. This structural conclusion adds no error charge.

No complete support23/24 exhaustion or support-floor increase is claimed. Every actual field, lattice, class mass, phase, Haar measure, finite projector, compact twist and selected Hilbert metric remains fixed. The signed sufficient inequality, prescribed global Bessel/Mram nonvanishing, selected arithmetic image, literal final source constructor and DP44 comparison remain separate open problems. The companion certificate checks exact finite partial-permutation and noncommuting-projector models, the original test-space dimension, the complete exponents and chronology; the general operator, source/target and positive-fixed-space arguments are the proof above.

\subsubsection{Source anchors used at their established scope}\label{source-anchors-used-at-their-established-scope}

\begin{itemize}
\tightlist
\item
  K058(PG3): actual unitary R, both mixed orthogonal projectors, and probability metric.
\item
  K062(AT3--8, AT13--16, AT23--24): admitted labels including zero, paired incidence matrix, literal kernel/class masses, normalized coefficient and inverse symmetry.
\item
  K068(LR3, LR20): determinant-one central representative and all unmixed finite unipotent reductions.
\item
  R059(GB4--21), R070(SC23--24): original integral transverse row, common G ideal, both target columns, full source fibre, all support and pair factors, exact support-variable exponents.
\item
  R056(BM9--10): earlier use of the same literal inversion on an entire sector before intersection with a possibly non-invariant remainder; no new symmetry premise is taken from a radial upper bound.
\item
  DO080(RT1): the actual normalized finite-dimensional compression of the unitary selected representation, with both compact factors retained.
\end{itemize}

\clearpage

\section{AM-DO-094 --- paired finite ranks and the last transverse scalar}\label{am-do-094-paired-finite-ranks-and-the-last-transverse-scalar}

\textbf{Canonical checked integration --- 2026-11-22T22:01:23.041321+00:00.} Complete root and independent proof/script/interface readings and separate exact replays passed. Historical author/pending wording below is preserved with its original review history. Under the full original/inverse doubly-plane block conditions, finite nilpotent rank at most3 and both nonzero paired cross blocks force the transverse scalar alpha=1. In the actual global residual, retained norm-one and polynomial wedge identities also give alpha=1 at rank\textless=2 with at least one nonzero cross block, and rank\textless=1 without a cross condition. The explicit original G-unit basis makes delta\_z=(alpha−1)U/M integral, with all conjugates\textless2\^{}71. Exact counterexamples exclude the weaker block hypothesis and the two unclaimed rank/cross cells. This supplies no complete restoration/count, charge, global nilpotence or support-floor change.

\textbf{Algebraic-geometric translation.} The statement concerns paired determinantal conditions on the original Hermitian modules over Spec(O\_E). Its rational orthogonal basis is proved integral and invertible only at the actual G primes; the scalar pairing uses separate source and target selfdual lattices. Exact carried contextual anchors are SGA3 I4.7.3/4.7.3.1; SGA6 VI Theorem3.3; FGA232 Theorem3.1 pp148--149; FGA236 Corollaries1.7/2.3 pp229/232; SGA7 I VII2.7/2.9 pp169--174 and VIII7.4(b) pp310--311; cote4 scan003 §§8.1.1--8.1.2, scans005--006 §§8.4.1--8.4.3 and cote54 scans005--007. No new scheme comparison, lambda-ring/K-theory identity, l-adic sheaf/weight theorem, CM-motive or motivic L-function assertion follows. DP\_MULT Q3e and THETA3 G1/A7 gain the displayed local conditions, while CX\_GRH CX-ALG-024 and GCT\_PLETH T-052 retain their arithmetic gates.

Author: ALP\_MACHINE \nolinkurl{/root/definite_oscillator}. 22 November 2026 UTC.

Scoped pending mathematical finding. No finite-sum charge, support-floor update or canonical artifact is made here. The general proof, exact original-coordinate application and diagnostic models are kept distinct.

\subsection{1. Precise finite hypotheses}\label{precise-finite-hypotheses}

Let k be a field, let V=S direct-sum P direct-sum k t, and let the paired bilinear form have block matrix

\[
 J=\operatorname{diag}(J_S,J_P,j),
 \qquad J_S,J_P\text{ symmetric and invertible},\quad\dim P=2,
 \quad j\ne0.
 \tag{FT1}
\]

No positivity of a finite form is assumed. Let h be invertible, put h\_c=J\^{}−1 h\^{}−t J, and put N=h−I. Assume:

\begin{itemize}
\tightlist
\item
  N is nilpotent and rank(N)≤3.
\item
  The S-by-S blocks of h and h\^{}−1 are both I.
\item
  Every S-to-t and t-to-S block of both h and h\^{}−1 is zero.
\item
  The P-by-S blocks C of h and C\_c of h\_c are both nonzero.
\end{itemize}

These are the full doubly-plane conditions: both original and inverse cross blocks are planar at both paired components. Requiring only that the two original lower cross blocks are planar is weaker and does not suffice; section4 gives an exact counterexample.

Let alpha be the t-by-t entry of h. Then

\[
                           \alpha=1.                            \tag{FT2}
\]

The same statement applies at the companion component. The proof permits any dimension of S and therefore does not require choosing an integral basis for a moving one-dimensional cross quotient.

\subsection{2. General proof without a finite positive metric}\label{general-proof-without-a-finite-positive-metric}

The paired relation is h\_c\^{}t J h=J. Its S-by-S block, using the hypotheses, gives

\[
                         C_c^tJ_PC=0.                           \tag{FT3}
\]

Both C and C\_c are nonzero and P has dimension two with a nondegenerate pairing. Thus each has rank one and the kernel in P of C\_c\^{}tJ\_P is exactly im(C). Write c=Ns for some s in S with c nonzero. The S-to-t and t-to-S hypotheses give

\[
                 ht=\alpha t+p,\qquad h_ct=\alpha_ct+p_c,
                 \quad p,p_c\in P.
\]

Pairing h\_cS with ht gives C\_c\^{}tJ\_Pp=0, hence p=beta c.~Likewise p\_c lies in im(C\_c). By FT3 their pairing is zero. Pairing h\_ct with ht now gives

\[
             \alpha_c\alpha=1,\qquad
             Nt=(\alpha-1)t+\beta Ns.                          \tag{FT4}
\]

In particular alpha is nonzero. Suppose a=alpha−1 is nonzero. Since N is nilpotent of rank at most three, N\^{}4=0. The operator N−aI is invertible, so FT4 gives the exact finite expansion

\[
 t=-\beta\bigl(a^{-1}Ns+a^{-2}N^2s+a^{-3}N^3s\bigr).           \tag{FT5}
\]

Let pi\_S and pi\_t be the original coordinate projections, with pi\_t(t)=1. The identity S blocks and zero S/t blocks of h and h\^{}−1=I−N+N²−N³ imply

\[
\begin{array}{lll}
 \pi_SNs=0,&\pi_SN^2s=\pi_SN^3s=w,\\
 \pi_tNs=0,&\pi_tN^2s=\pi_tN^3s=b.
\end{array}                                                     \tag{FT6}
\]

Applying pi\_t to FT5 gives

\[
                       1=-\beta b\alpha/a^3.                  \tag{FT7}
\]

Thus beta and b are nonzero, and in particular N³s is nonzero. Applying pi\_S to the same formula gives

\[
                       0=-\beta\alpha w/a^3,
\]

so w=0. The three vectors Ns,N²s,N³s are linearly independent: successively apply N² and N to a putative dependence, using N³s nonzero. They lie in im(N), whose dimension is at most three, so they span im(N). Equation FT6 therefore implies pi\_S N=0 on all V. The S rows of every positive power of N are also zero, hence the S-to-(P direct-sum kt) block of h\^{}−1 is zero. The definition h\_c=J\^{}−1 h\^{}−t J then forces C\_c=0, a contradiction. This proves FT2.

This proof uses nilpotence only at the finite component and only under the displayed finite rank bound. It does not infer global nilpotence or assume that finite Hermitian forms have signature one.

Two further implications are available for reductions of the actual global residual. They require the following algebraic identities to survive reduction: alpha*alpha\_c=1, and the global rank-one relation between ht−alpha t and the columns of C. Their validity and local integrality are established in section3.

If rank(N)≤2 and C is nonzero, the surviving rank-one relation gives Nt=a t+beta Ns for some s in S, even when C\_c is zero. Here beta is a scalar in the residue field obtained from any nonzero column of C; no globally divided coefficient is presumed integral. Now N³=0. The original and inverse t/S zero blocks give pi\_t Ns=pi\_t N²s=0. If a were nonzero, the shorter version of FT5 would give

\[
           t=-\beta(a^{-1}Ns+a^{-2}N^2s),
\]

contradicting pi\_t(t)=1. Thus alpha=1 at that component, and alpha*alpha\_c=1 gives the same conclusion at its partner. Apply this argument at the companion if it alone has nonzero C. In particular

\[
 d_v\le2,\quad\text{at least one paired C block nonzero}
       \quad\Longrightarrow\quad\alpha=\alpha_c=1.             \tag{FT7a}
\]

At rank at most one no C-nonvanishing condition is needed. N²=0 gives alpha\_c=1−a, while alpha=1+a. The retained product identity then gives 1=(1+a)(1−a)=1−a². A field has no nonzero element of square zero, so

\[
                 d_v\le1\quad\Longrightarrow\quad
                             \alpha=\alpha_c=1.                 \tag{FT7b}
\]

Neither corollary is claimed for arbitrary unrelated finite matrices without the stated surviving global identities. In particular no analogous conclusion has been proved here for rank two with both C blocks zero, or for rank three with exactly one nonzero C block.

\subsection{3. Exact original-coordinate applicability after DO093}\label{exact-original-coordinate-applicability-after-do093}

The actual doubly-plane/four-zero residual of DO093(IV19--22) has S of dimension164 and a short positive two-space Y=span\_E(y\_1,y\_2). Here use the original integer data

\[
\begin{split}
 u&=(2,95192390527,383092671043,12529843548438452388609),\\
 U&=u^tu=77Q^2-14,\qquad Q=1427908193729121522601,\\
 M&=u_1^2+u_2^2+u_3^2=155821585821105090045582,\\
 y_1&=(-976627,-153032,38026,0),\\
 y_2&=(-567678,695570,-172838,0),\\
 z&=(256766,-7654233287,-30805128044,1).
\end{split}                                                     \tag{FT8}
\]

The last vector equals (y\_3−y\_1+y\_2)/2 for R059's original y\_3 and is an integer vector. Its squared norm is exactly1007543201084982869062. The first three coordinates of y\_1 cross y\_2 equal −2(u\_1,u\_2,u\_3), whence

\[
       \det\operatorname{Gram}(y_1,y_2)=4M.
\]

Put

\[
 t=\left(-\frac{u_4}{M}(u_1,u_2,u_3),1\right).
 \tag{FT9}
\]

Then t is orthogonal to both Y and u, z−t belongs to Y, and \textless t,t\textgreater=U/M. The positive four-coordinate basis y\_1,y\_2,u,t has determinant of modulus2U. Direct integer gcd checks at each actual prime g\_0=22285932761 and g\_1=64072175441 give

\[
                   \gcd(2MU,g_0g_1)=1.                        \tag{FT10}
\]

Accordingly this is an invertible integral change of basis over each actual G-local DVR, at both CM partners. No claim of global integral invertibility is made. The new paired form is block diagonal with blocks Gram(Y), U, −1, U/M, all nondegenerate on reduction. This verifies the precise finite-coordinate premise, rather than treating a rational orthogonal basis as automatically integral.

For the finite lemma take the enlarged block S'=S direct-sum Y, dimension166, the plane P=Eu direct-sum Ee\_n, and the last line Et. DO093 proves N(Y) is contained in im(C) inside P, and proves the same statement for h\^{}−1. The S'-by-S' blocks are therefore I. The original plane conditions give h and h\^{}−1 no original-S/t entries. For the Y/t entries, use adjointness and N\^{}−Y contained in P: since Y and P are perpendicular to t, \textless y,ht\textgreater=\textless h\^{}−1y,t\textgreater=0; the inverse argument is identical. The t/S' entries vanish because N(S') and N\^{}−(S') lie in P. Thus all the finite hypotheses on blocks in section1 survive at G after FT10.

If both original C blocks at the two CM components are nonzero, then both enlarged C blocks are nonzero. If the actual finite rank d\_v is at most three, K068's retained finite unipotence and FT2 give alpha=1 at both partners. This is the exact conditional implication available for a subsequent Green-rank restoration argument. Places where one cross block vanishes require their own case and are not included by this implication.

The global residual also has the scalar alpha on the remaining transverse quotient. The same paired argument as FT3--4, now over E, gives alpha*bar(alpha)=1 and the rank-one dependence of ht−alpha t with the columns of C. All entries of these matrices, and alpha and its companion, are integral at the actual G places by FT10 and original maximal-level integrality. The polynomial wedge relations therefore survive reduction even if a divided coefficient beta does not. At a component with nonzero reduced C they imply exactly the residue-field relation required for FT7a. The product identity makes both alpha and alpha\_c units at G and supplies the input to FT7a--b. At each archimedean label it also gives \textbar alpha\textbar=1. Moreover

\[
 \delta_z:=\langle z,(h-I)z\rangle
               =(\alpha-1)\frac{U}{M},\qquad
                 |\tau\delta_z|\le\frac{2U}{M}<2^{71}.        \tag{FT11}
\]

Indeed z=t+y for some y in Y, N(Y) lies in the cross line in P, and Nt=(alpha−1)t plus a vector of that same line. All P terms pair trivially with z, which is positive and perpendicular to u. The identity follows. The vector z belongs separately to every actual source and target pencil: at T, it is the displayed half-integral combination of the three original common vectors and2 is a unit; away T it is an integer vector in the original lattice. The arbitrary-arrow selfduality argument therefore makes delta\_z an algebraic integer. By FT10, U/M is a G-unit, so FT2 implies that delta\_z vanishes at both paired G primes whenever d\_v≤3 and both cross blocks are nonzero.

Thus delta\_z vanishes at both paired G primes in each of the following cases: finite rank at most three with both C blocks nonzero; finite rank at most two with at least one nonzero C block; and finite rank at most one without a C condition. These assertions distinguish the actual reduction rank from global rank, and use only the proved global identities in the two extended cases.

This supplies a new possible scalar witness. It does not by itself supply a complete Green restoration table, a global ideal-pattern count, a support-sector bound or an error charge. The current full ordered error lists are unchanged, as are the all-top and uniform support floors.

\subsection{4. Why the inverse-plane condition cannot be omitted}\label{why-the-inverse-plane-condition-cannot-be-omitted}

In the coordinate order (S,P\_1,P\_2,t), over the field with three elements, take J=diag(1,1,−1,1) and

\[
 N=\begin{pmatrix}
 0&1&2&2\\0&0&2&2\\1&1&2&0\\0&2&1&1
 \end{pmatrix},\qquad
 N_c=\begin{pmatrix}
 0&2&1&0\\2&0&1&1\\0&2&0&0\\0&0&1&0
 \end{pmatrix}.
 \tag{FT12}
\]

Here h=I+N, h\_c=I+N\_c=J h\^{}−t J, N\^{}4=0 and rank(N)=3. Both S blocks are1; both lower C blocks are nonzero and lie in P. Nevertheless alpha=2, not1. The original h has nonzero S/t entry2, so it violates the inverse-plane condition at the companion component. It is not a counterexample to FT2 or a candidate in the full DO093 residual.

The exact certificate also generates320,000 rank-three Jordan conjugation trials over the fields of3,5,7 and11 elements. It records separately the weaker and full conditions, and verifies FT2 in every model meeting the full conditions. These are exact diagnostic examples, not an exhaustive finite-group enumeration and not the proof of FT2. The proof is the general chain argument in section2. No new arithmetic sector is claimed removed in this note.

The two unclaimed cells in section2 also have genuine finite counterexamples even with the full double-plane blocks, the scalar product identity and the surviving wedge relation. For a nonzero a with a+1 nonzero, put b=a²/(a+1) and

\[
 M_a=\begin{pmatrix}
 -a&-b&1\\1&0&0\\b-a^2&-ab&a
 \end{pmatrix}
 \quad\text{on }(P_1,P_2,t).                                  \tag{FT13}
\]

Direct multiplication gives M\_a³=0, and its minor in rows1,2 and columns1,3 is−1; thus rank(M\_a)=2. Its t-diagonal entry is a, its square has t-diagonal b, and (1+a)(1−a+b)=1. On S direct-sum P direct-sum kt, first take N=diag(0,M\_a). Both paired C blocks vanish and alpha=1+a is not1. Alternatively take

\[
 N=\begin{pmatrix}0&0\\e_1&M_a\end{pmatrix},                   \tag{FT14}
\]

where e\_1 is the first plane coordinate. This has rank three: its rows2,3,4 and columns1,2,4 have determinant a, and N⁴=0. Its original C is nonzero and its companion C is zero. The original and inverse S/t blocks are all zero, because b(1+a)=a². Moreover Nt=a t+Ns, exactly the surviving wedge relation, and still alpha=1+a is not1. These algebraic examples work at the actual odd G primes with a=1, b=1/2; the certificate checks those primes and their actual diagonal plane and transverse forms. They are finite counterexamples to stronger local implications, not assertions that a particular rational global short-sum matrix realizes them.

\subsection{5. Carried source and algebraic-geometric scope}\label{carried-source-and-algebraic-geometric-scope}

The exact actual residual and its adjoint identities are DO093(IV19--22), after the literal original-kernel inversion in IV4--18. R059(GB4--7) supplies the original integral transverse vectors and their separate membership in each source and target pencil; DO078(FP20--22) is the corresponding arbitrary-arrow pairing input. K068(LR20) supplies finite unipotence at the unmixed G factors. None of these inputs asserts global nilpotence of the residual. The G-unit matrix and scalar witness in FT8--11 are proved here, without an integral moving frame.

The modules remain the original Hermitian modules over Spec(O\_E). Localization at the actual G primes uses the explicitly invertible matrix in FT9--10, and the finite conclusion is a determinantal statement on that original reduction. The exact carried contextual anchors remain SGA3 I4.7.3/4.7.3.1; SGA6 VI Theorem3.3; FGA232 Theorem3.1 pp148--149; FGA236 Corollaries1.7/2.3 pp229/232; SGA7 I VII2.7/2.9 pp169--174 and VIII7.4(b) pp310--311; cote4 scan003 §§8.1.1--8.1.2, scans005--006 §§8.4.1--8.4.3 and cote54 scans005--007, as recorded in canonical R070. They are retained context, not newly invoked theorems for this linear-algebra lemma. No new scheme comparison, lambda-ring/K-theory equality, l-adic sheaf or weight theorem, CM-motive identification, motivic L-function assertion, selected period value or arithmetic-cost conclusion is inferred.

\clearpage

\section{AM-K095 --- the complete middle T fibre and the actual composite-incidence matrix}\label{am-k095-the-complete-middle-t-fibre-and-the-actual-composite-incidence-matrix}

\textbf{Canonical checked integration --- 2026-11-22T22:01:23.073192+00:00.} Complete root and independent proof/script/interface readings and separate exact replays passed. Historical author/pending wording below is preserved with its original review history. The full middle T-covector fibre, with its genuine fine stabilizer masses, sums exactly to the composite incidence u2 Z X=alpha. The weighted coarse factor1/\textbar Gamma\_q\textbar{} is represented by a full coarse-stabilizer G cut only after invariant middle coevaluation. The finite norm is divided by \textbar Gamma\_q\textbar, and the auxiliary family is split at11 before the prescribed ordinary derivative. The exact K094 Fourier columns become the same-input coarse-class quadratic-shear formula MF12, with polar rank2b and explicit central phase. All original coefficient clearings, norms, phases, derivatives, cover pullbacks and terminal operators remain; selected content and coarse-class noncancellation are open.

\textbf{Algebraic-geometric translation.} This is a finite-fibre orbit--stabilizer identity and integral splitting of the auxiliary theta-section module at11 on the existing Picard/Hom and theta-group data. The finite quadratic shear retains its Heisenberg central phase; it is not identified with an algebraic additive translation family or an independent jet source. Exact carried contextual anchors are SGA3 I4.7.3/4.7.3.1; SGA6 VI Theorem3.3; FGA232 Theorem3.1 pp148--149; FGA236 Corollaries1.7/2.3 pp229/232; SGA7 I VII2.7/2.9 pp169--174 and VIII7.4(b) pp310--311; cote4 scan003 §§8.1.1--8.1.2, scans005--006 §§8.4.1--8.4.3 and cote54 scans005--007. No new lambda-ring/K-theory equality, l-adic weight theorem, CM-motive comparison or motivic L-function theorem follows. DP\_MULT Q3e and THETA3 G1/A7 retain the prescribed derivative/coarse-class and arithmetic comparison gates; CX\_GRH CX-ALG-024 and GCT\_PLETH T-052 keep their separate requirements.

Author: \nolinkurl{/root/kummer_extension}, 22 November 2026 UTC.

\textbf{Author whole-proof and complete-script PASS; frozen for independent review.} This is an exact regrouping of the existing coupled class map, followed by its actual integral Fourier change of source basis at 11. It removes the middle T-covector sum and its stabilizer mass after the legitimate invariant contraction. Primitivity below concerns the auxiliary finite tensor and section before the prescribed ordinary derivatives. The coarse middle-class sum and the selected arithmetic image remain unevaluated.

\subsection{1. The original objects and the middle coarse fibre}\label{the-original-objects-and-the-middle-coarse-fibre}

Keep the actual J14 or J56 host, the unchanged two column edges of K038, every allowed word, tuple, twist, terminal decoration and original finite level. Put \nolinkurl{a=n0}, \nolinkurl{b=n1}, \nolinkurl{c=n2}. Fix an initial coarse class c0 as in K094 and a terminal decoration d.~The initial covector alpha runs through the actual full nonzero fibre

\[
 \Omega_0=\prod_{v\mid T}\bigl((N_0/\mathfrak P_vN_0)^\vee\setminus\{0\}\bigr),
\tag{MF1}
\]

where the product notation means that each component is nonzero. Use K094's common coefficient lattice localized at 11, and one coefficient P in that space for the entire actual middle sum. All global coefficient-clearing factors remain those of K094(FC8).

For indexing only, forget the T covector of a middle class while retaining its unimodular lattice M, its original G-frame orbit and character data, and every Q frame. Call the resulting coarse class q. Let Gamma\_q be its finite rational stabilizer, with its existing G/Q level lifts. It need not be trivial and its order need not be an 11-unit. Define

\[
 \Omega_q=\prod_{v\mid T}\bigl((M/\mathfrak P_vM)^\vee\setminus\{0\}\bigr).
\tag{MF2}
\]

The actual middle classes over q are Gamma\_q-orbits in Omega\_q. For beta in that fibre, their actual stabilizers are \nolinkurl{(Gamma_q)_beta}. The G/Q level lifts are uniquely determined by the rational isometry: the G frame is surjective on its residue module and each Q frame is invertible. Hence Gamma\_q embeds faithfully in the original lattice-isometry group, with no extra kernel of level lifts. This is the same decorated-lattice orbit description used in K038 and K094; no theta input is moved to the coarser level.

The rational coefficient space, its fixed polynomial action, and its perfect middle pairing are independent of beta. The integral lattices transported from \nolinkurl{Lsh} can vary at T. They are not identified globally. At 11 they have their usual common localization. The full middle coevaluation is the basis-free contraction of the rational space with its dual; it does not depend on which of those integral lattices is displayed.

These assertions follow directly from \nolinkurl{colValF} and \nolinkurl{colVal}, \nolinkurl{frmRho}, \nolinkurl{frsF}, and the literal \nolinkurl{ChKer.Shape.cp1/cp2}. The theta distributions use M and the non-T frames. \nolinkurl{ybSet} uses M and G/Q only. Beta enters the coupled finite functions solely through the two \nolinkurl{phiT} conditions. In particular the actual Y set is unchanged and has no added T coordinate.

\subsection{2. The exact coarse mass, with no unweighted replacement}\label{the-exact-coarse-mass-with-no-unweighted-replacement}

Let F\_q(beta;P,d) be the complete actual numerator for the middle decoration beta: sum the full Y set and compose the two prescribed ordinary selected kernel maps on the same P. It takes values in the actual target coefficient space. The two covariance laws K038(CA24--25) cancel the full middle action, including its rational polynomial action and G/Q character. The action also permutes Y. Therefore F\_q is constant on Gamma\_q-orbits.

For any finite group Gamma acting on a finite set Omega and any invariant vector-valued F, the elementary orbit--stabilizer identity is

\[
 \sum_{[x]\in\Gamma\backslash\Omega}
       \frac{F(x)}{|\Gamma_x|}
       =\frac1{|\Gamma|}\sum_{x\in\Omega}F(x).
\tag{MF3}
\]

Indeed an orbit has size \nolinkurl{|Gamma|/|Gamma_x|}. Apply this with Gamma\_q and Omega\_q. It gives the exact original middle contribution

\[
 \mathcal A_q(\alpha;P,d)
       =\frac1{|\Gamma_q|}\sum_{\beta\in\Omega_q}
                                      F_q(\beta;P,d).
\tag{MF4}
\]

This is a rational identity at this point. No inverse of \nolinkurl{|Gamma_q|} is called integral. In particular a nontrivial stabilizer, a beta fixed by a nontrivial subgroup, and a stabilizer order divisible by 11 are all permitted. Every initial and terminal mass, Haar factor and source scalar is the one already retained in K038/K092/K094.

\subsection{3. Sum the actual T incidences before the selected evaluation}\label{sum-the-actual-t-incidences-before-the-selected-evaluation}

At one selected T place write X for the first position matrix, Z for the second, and u2 for the actual terminal covector. Once the original integrality indicators hold, the two literal conditions are

\[
                         \beta X=\alpha,
                         \qquad u_2Z=\beta.
\tag{MF5}
\]

For nonzero alpha the elementary identity is

\[
 \sum_{\beta\ne0}
       \mathbf1_{\beta X=\alpha}\mathbf1_{u_2Z=\beta}
                     =\mathbf1_{u_2ZX=\alpha}.
\tag{MF6}
\]

The second equality in MF5 determines beta uniquely. If the right side holds, beta cannot be zero because alpha is nonzero. This proves MF6 even when \nolinkurl{u2 Z=0} and for every rank of X or Z. Tensor it over the actual T places.

Let Phi\_q\^{}0 be the joint two-edge finite tensor retaining the original integral indicators at T and every G/Q factor, the full Y sum, all splitting phases and actual frames, but neither T equality in MF5. It is a proof device, not a new admitted automorphic input. Define

\[
 \Phi_{q,\alpha}^{\rm comp}
   =\Phi_q^0\prod_{v\mid T}
                    \mathbf1_{u_{2,v}Z_vX_v=\alpha_v}.
\tag{MF7}
\]

The universal coupled differential and the canonical middle coevaluation are fixed linear operations on the joint tensor, independent of beta. The underlying theta distributions, comparison coordinates and Y are also independent of beta. Thus the finite beta sum can be taken before these linear operations. Formula MF6 gives

\[
 \mathcal A_q(\alpha;P,d)
   =|\Gamma_q|^{-1}\,
       \mathcal J_{q,d,P}\bigl(\Theta_D(\Phi_{q,\alpha}^{\rm comp})\bigr).
\tag{MF8}
\]

Here J is the \textbf{same complete two-edge ordinary selected evaluation}, with its original determinant factors, cD coordinates and exact comparison-cover pullback. A common comparison denominator D prime to 11 can be used for this finite family, as in K089/K094; it is not an admitted level change and no global size estimate for D is supplied. Neither a Hasse derivative nor an independent class jet replaces J.

\subsection{4. The legitimate full coarse-stabilizer cut}\label{the-legitimate-full-coarse-stabilizer-cut}

The actual diagonal middle action sends \nolinkurl{(X,Z)} to \texttt{(gamma\ X,Z\ gamma\^{}\{-1\})}. It preserves \nolinkurl{u2 Z X}. Before summing beta, the exact finite covariance sends the beta summand to the summand with transported beta; the honest two-edge character and phase factors cancel as in LIFT34 and K038. Thus permutation of the complete beta fibre proves invariance of Phi\_q\^{}comp under the full Gamma\_q, without requiring an individual beta to be fixed. The same is true after multiplying by any function of \nolinkurl{u2 Z X}.

The full-Gamma cutoff argument LIFT34(MC9--14,MC21a--b) applies to Gamma\_q. Its proof needs only a finite subgroup of the original rank-b lattice isometry group, its injective reduction at the fixed G prime, and the full column rank of the second raw G support. Forgetting the T covector changes none of those conditions. The residue-characteristic comparison \nolinkurl{g-1>2fb} excludes torsion in the reduction kernel for this larger finite stabilizer just as for the fine stabilizer. Right multiplication on the full-column-rank second G matrices is then free.

Choose a transversal of those free Gamma\_q-orbits and let C\_q be its 0/1 indicator, transported through the same actual frames as in K092. It is independent of alpha, beta and P. The exact honest-action identity is

\[
 \Phi_{q,\alpha}^{\rm comp}
  =\sum_{\gamma\in\Gamma_q}
        \omega_{\rm mid,f}(\gamma)
                    (C_q\Phi_{q,\alpha}^{\rm comp}).
\tag{MF9}
\]

After, and only after, the invariant middle coevaluation, this gives

\[
 \boxed{\quad
 \mathcal A_q(\alpha;P,d)
       =\mathcal J_{q,d,P}
            \bigl(\Theta_D(C_q\Phi_{q,\alpha}^{\rm comp})\bigr).
 \quad}
\tag{MF10}
\]

This is not the equality of the uncontracted section with its scalar quotient by \nolinkurl{|Gamma_q|}. It is the original mass-divided contribution represented by a legitimate contracted cut. It does not alter the arithmetic class map.

The translates in MF9 have disjoint support at the chosen raw G residue and the honest finite Weil action is unitary. Therefore the exact finite norm is

\[
 \|C_q\Phi_{q,\alpha}^{\rm comp}\|_f^2
       =|\Gamma_q|^{-1}\|\Phi_{q,\alpha}^{\rm comp}\|_f^2.
\tag{MF10a}
\]

Actual class/polarization transport preserves this identity. It is not a lower bound for the selected value or an inverse norm of the arithmetic class map. As in K092(PM3--4), the framed cutoff is the conjugate of the raw indicator by that transport; its evaluation functionals are transported by the inverse operator, rather than treated as unchanged point evaluations.

The cut may retain one prescribed full-Y unit-support point. K090(MA13--17) supplies common non-T G/Q coordinates whose full-Y value is an 11-unit. At T choose any nonzero intermediate row beta, then an integral Z with \nolinkurl{u2 Z=beta} and an integral X with \nolinkurl{beta X=alpha}; the nonzero covectors make both linear systems surjective. These choices are independent of the G/Q coordinates by the original local product and CRT. The same second G coordinate can be used for every alpha. Choose it as the representative of its free Gamma\_q-orbit when defining C\_q.

The tensors \nolinkurl{C_q Phi_(q,alpha)^comp} now have disjoint alpha supports and each has a unit value. Evaluation at one such point per alpha gives an integral left inverse on their coefficient span, after dividing only by those actual units. The actual integral finite theta reading of K089, with its trace-one normalization, transports this split auxiliary family to the common comparison section module. Actual class transport is the integral invertible prime-to 11 Weil transport already checked there. Hence the family is primitive at 11 \textbf{before J}. This does not assert primitiveness after J or after summing q.

\subsection{5. The exact quadratic phase on the initial Fourier basis}\label{the-exact-quadratic-phase-on-the-initial-fourier-basis}

Use precisely the K094 Fourier basis, with one vector lambda\_v per residue line and scalars s\_v distinct from 1. At a fixed coarse-middle q, substitute MF7 into the sum over the actual nonzero initial alpha. For any row w in the initial covector space,

\[
 \sum_{\alpha\ne0}
  [\psi(\alpha(s\lambda))-\psi(\alpha(\lambda))]
                         \mathbf1_{w=\alpha}
               =\psi(w(s\lambda))-\psi(w(\lambda)).
\tag{MF11}
\]

When w is zero both sides are zero. Therefore the \textbf{entire actual transformed column} is

\[
 \boxed{\quad
 (\mathcal A C_{c_0,\lambda,\mathbf s,P})(d)
  =\sum_q\mathcal J_{q,d,P}\!\left(
    \Theta_D\!\left[
      C_q\Phi_q^0
       \prod_{v\mid T}\left(
         \psi_v(u_{2,v}Z_vX_v s_v\lambda_v)
          -\psi_v(u_{2,v}Z_vX_v\lambda_v)
                         \right)\right]\right).
 \quad}
\tag{MF12}
\]

There is no surviving beta sum, inverse \nolinkurl{T-1}, T-power normalization or extra middle mass. The initial Fourier norm and any global clearing factor are exactly K094(FC8). The same coefficient P occurs in every q term. Terminal restrictions and their genuine/generic projectors, branches and denominators can be applied to this identity only with their original normalization.

For a single multiplier the finite exponent is the genuine quadratic form

\[
                  Q_\mu(X,Z)=u_2ZX\mu.
\tag{MF13}
\]

Its polar form on \nolinkurl{Hom(k^a,k^b) direct-sum Hom(k^b,k^c)} is

\nolinkurl{Q_mu(X+X',Z+Z')-Q_mu(X,Z)-Q_mu(X',Z')  = u2 Z X' mu + u2 Z' X mu}.

For nonzero mu and u2, the maps \nolinkurl{X -> X mu} and \nolinkurl{Z -> u2 Z} are surjective onto k\^{}b and its dual. The cross pairing is perfect on those quotients. Since T is odd, the symmetric polar matrix has rank 2b, and \nolinkurl{Q=one-half times its diagonal evaluation}. More explicitly, on the finite position space write \nolinkurl{D_Q f(v)=psi(Q(v))f(v)}, \nolinkurl{tau_h f(v)=f(v+h)} and \nolinkurl{m_xi f(v)=psi(xi(v))f(v)}. Then direct substitution gives

\[
 D_Q\tau_hD_Q^{-1}
     =\psi(-Q(h))\,m_{-B_Q(h,\cdot)}\tau_h,
       \qquad D_Qm_\xi D_Q^{-1}=m_\xi.
\tag{MF13a}
\]

Thus the multiplier is exactly the finite Schrödinger quadratic-shear operator with its displayed central phase in this convention. This is a finite-operator statement. It does not identify the family with translations on a single ample line, discard a global Weil phase, or supply a rational shear of bounded height.

\subsection{6. What this removes from the arithmetic gate}\label{what-this-removes-from-the-arithmetic-gate}

MF12 is an identity for the original localized arithmetic image. It removes all cancellation attributed merely to choosing the middle T-covector representative: that entire fibre has been summed into the single composite incidence, with its exact coarse mass and contracted cut. It also gives a common finite quadratic family independent of the former class-dependent middle covectors.

It leaves a sum over \textbf{coarse} middle classes. The actual ordinary selected functional J can annihilate a primitive auxiliary section, and those coarse-class values can cancel. Its restriction to any prescribed smaller source submodule remains K094's exact Fourier inverse image of that submodule. No unrestricted jet source is available.

K093's positive residual first-edge derivative order remains valid in the original fine-fibre expansion, term by term, and this regrouping does not remove it. No holomorphic exponential formula is inferred merely from the finite quadratic expression MF13. In particular the assertion that T is large compared with the selected derivative degree is not yet a torsion zero estimate: one still needs the actual algebraic parameter space, restricted line or jet bundle, degree and scalar transports, and nonvanishing of the same-input coarse-class combination. No such inequality or noncancellation is claimed here.

Primitivity in section 4 is at the finite tensor and auxiliary-section stage only. Prime-to 11 Fourier/Weil units do not prove that later irreducible or generic projectors have integral denominators. The original periods, creation factors, all formal-coordinate pullbacks, global coefficient clearing and norm/height costs are retained. This finding proves no selected upper-content bound, nonzero arithmetic row, global Bessel/Mram period, literal source-family constructor or DP44 conclusion.

\subsection{7. Literal inputs and exact diagnostic scope}\label{literal-inputs-and-exact-diagnostic-scope}

Direct source locators in \nolinkurl{Formalism/LocalToGlobal/Columns} are: \nolinkurl{CoefCol.lean}, colVal; \nolinkurl{FiniteSumNonzeroLargeDepth.lean}, frmRho425--427, kerOf636--640, phiT646--650, colValF1083--1084, frsF1103--1104, ybSet1125--1129, XiCpl1/2 at1134 and1160, and ChKer.Shape.cp1/cp2 at1269--1275; and \nolinkurl{StdLatUnimodular.lean}, Deco.Ok, the T covector action and Lsh. These show exactly which structures vary with beta. The common rational coevaluation and actual class covariance are K037/K038(CA23--27); full-Gamma cut and phase conventions are LIFT34(MC9--14,MC21a--b) and K092; the unit full-Y witness and trace-one integral theta reading are K089/K090; the complete initial Fourier basis and its local/global distinction are K094.

The companion checks finite composite-incidence identities, nonzero and zero composite rows, the exact character-difference Fourier formula, orbit--stabilizer weights with nontrivial stabilizers including order 11, the contracted free-orbit cut, full middle coefficient contraction, tensor corners, and the quadratic polar rank. It also fingerprints the actual source and dependency files. These are diagnostics for the displayed finite algebra; they do not evaluate the enormous selected derivatives or prove a coarse-class nonzero value.

\subsection{Completed author verification}\label{completed-author-verification-2}

I read all seven substantive sections and MF1--MF13a, the complete companion script, and the literal source interfaces listed in section 7. The check reconstructed the actual coarse-fibre orbit description with uniquely determined level lifts, beta-independent rational coevaluation, the exact weighted orbit--stabilizer formula, the zero-composite-row case, and the permutation of the full beta fibre under the honest finite middle action. I separately rechecked LIFT34's full-stabilizer injectivity and free G support, its invariant real creation vector and rational theta automorphy, the conjugated cutoff and finite norm identity, and the full-Y unit witness transported by the actual integral theta/Weil reading. The final Fourier identity retains one initial P across every coarse class, the complete source ordinary selected evaluation and cover pullback, all terminal operators and global coefficient clearing. The exact quadratic shear's central phase is computed, with no holomorphic translation or upper-content inference. Final file hashes and the fresh exact replay are recorded separately.

\clearpage

\section{AM-R-071 --- global rigidity of the last transverse scalar}\label{am-r-071-global-rigidity-of-the-last-transverse-scalar}

\textbf{Canonical checked integration --- 2026-11-22T22:01:23.204516+00:00.} Complete root and independent proof/script/interface readings and separate exact replays passed. Historical author/pending wording below is preserved with its original review history. In the actual DO093 support23/24 doubly-plane/four-zero scalar-cross remainder, the original integral delta\_z=(U/M)(alpha−1) and alpha\emph{bar(alpha)=1 imply delta\_z=0 and alpha=1 globally. Here gcd(U,M)=9 and36U/M²\textless1 exactly. Integer Bezout makes delta\_z}bar(delta\_z)/(U/9) an algebraic integer in F; every real conjugate lies in{[}0,36U/M²{]}, so its positive integral norm forces zero. The actual h fixes a positive E-space of dimension166 and has rank(h−I)\textless=3, with remaining signature(2,1). No finite-rank hypothesis, integral moving frame, global nilpotence or new error charge is used. Complete support floors and all signed/global/arithmetic/DP44 gates remain open.

\textbf{Algebraic-geometric translation.} The new step is an integral trace-and-norm identity over Spec(O\_F) for the original scalar pairing section over Spec(O\_E). Its field norm is the determinant of multiplication, and all real realizations give a strict bound below1. The fixed space is the kernel of the actual rank-one map on the original positive167-space. Exact carried contextual anchors are SGA3 I4.7.3/4.7.3.1; SGA6 VI Theorem3.3; FGA232 Theorem3.1 pp148--149; FGA236 Corollaries1.7/2.3 pp229/232; SGA7 I VII2.7/2.9 pp169--174 and VIII7.4(b) pp310--311; cote4 scan003 §§8.1.1--8.1.2, scans005--006 §§8.4.1--8.4.3 and cote54 scans005--007. These are contextual frameworks; no new lambda-ring/K-theory equality, l-adic weight theorem, CM-motive comparison or motivic L-function assertion is invoked. DP\_MULT Q3e and THETA3 G1/A7 gain the strict smaller structural remainder. CX\_GRH CX-ALG-024 and GCT\_PLETH T-052 retain their arithmetic gates.

Author: \nolinkurl{/root}, GPT-6 Astra Pro, ALP\_MACHINE. 22 November 2026 UTC.

\textbf{Complete proof pending final author and independent whole-proof review.} This result concerns the actual doubly-plane, four-zero scalar-cross remainder identified by DO093. It proves that its last compact transverse scalar is one. It adds no numerical error charge and assumes no global unipotence.

\subsection{1. An integral norm-one rigidity lemma}\label{an-integral-norm-one-rigidity-lemma}

Let E/F be a CM extension with conjugation denoted by a bar. Let U and M be positive rational integers, put g=gcd(U,M), and suppose

\[
                         M^2>4gU.                         \tag{GR1}
\]

Let alpha in E satisfy alpha times bar(alpha)=1, and assume only that

\[
               \delta=\frac{U}{M}(\alpha-1)\in\mathcal O_E.
                                                               \tag{GR2}
\]

Then alpha=1 and delta=0. In particular alpha itself need not have been proved integral.

\textbf{Proof.} Substitute alpha=1+(M/U)delta into its norm-one identity. Multiplication by U squared and division by M gives

\[
                 U(\delta+\bar\delta)
                       =-M\delta\bar\delta.                \tag{GR3}
\]

Set U0=U/g and M0=M/g. These integers are coprime. Choose integers A,B with A U0+B M0=1. Put n=delta bar(delta) and s=delta+bar(delta), both in O\_F because they lie in O\_E and are fixed by conjugation. Equation GR3 yields M0 n=-U0 s and hence

\[
        \nu:=\frac{n}{U_0}=A n-Bs\in\mathcal O_F.             \tag{GR4}
\]

At every real embedding of F, extend to a complex embedding tau of E. The norm-one identity gives \textbar tau(alpha)\textbar=1, so

\[
 0\le\tau(\nu)
    =\frac{|\tau(\delta)|^2}{U_0}
    \le\frac{4U^2}{M^2U_0}
    =\frac{4gU}{M^2}<1.                                    \tag{GR5}
\]

If delta were nonzero, every factor tau(nu) would be strictly positive. Their product would be a positive rational integer, the field norm of the nonzero algebraic integer nu, and would be strictly less than one. This is impossible. Thus delta=0, and GR2 gives alpha=1. This proves the lemma. No integrality of alpha, principal-ideal hypothesis, G-divisibility, or local nilpotent rank bound enters the argument. The strict inequality matters: U=1, M=2, alpha=-1 gives g=1, M squared=4gU and delta=-1, a nonzero integral example at equality.

\subsection{2. The exact actual integral scalar}\label{the-exact-actual-integral-scalar}

Retain the original J56 Hermitian module, actual source and target pencils, class weights, selected metric and central representative of DO093. In its support23/24 remainder, write S for the original positive164-space, W0 for the original five-space of signature(4,1), and N=h-I. DO093(IV19--22) supplies the two packet-plane conditions and the vanishing of all four compressions on Y=span\_E(y1,y2). The original nonzero cross image CS is an isotropic line. The argument there proves N(Y) is contained in CS, and gives the corresponding inverse condition.

Use exactly DO094(FT8--11), with

\[
\begin{split}
 Q&=1427908193729121522601,\\
 U&=77Q^2-14=156996979348344707969383396772362042643000463,\\
 M&=155821585821105090045582,\\
 u&=(2,95192390527,383092671043,12529843548438452388609),\\
 z&=(256766,-7654233287,-30805128044,1),\\
 t&=(-u_4(u_1,u_2,u_3)/M,1).
\end{split}                                                   \tag{GR6}
\]

The original vector z is integral, perpendicular to u, and belongs separately to every actual source and target pencil. At the selected T components this also follows from z=(y3-y1+y2)/2, since two is a unit there; away T its original integer coordinates suffice. As h carries the source pencil to the target pencil, the target selfdual pairing gives

\[
               \delta_z=\langle z,(h-I)z\rangle\in\mathcal O_E.
                                                               \tag{GR7}
\]

This uses two separately specified lattices; it does not assume h preserves one lattice. It is DO078(FP20--22)'s arbitrary-arrow integral-pairing argument, not its additional globally unipotent square identity.

The rational line Et is orthogonal to Y and to u, with \textless t,t\textgreater=U/M and z-t in Y. The global paired-plane argument in DO094 §3, which does not require a finite-rank hypothesis, gives

\[
 ht=\alpha t+p,\quad p\in CS,\qquad
 \alpha\bar\alpha=1,\qquad
 \delta_z=\frac UM(\alpha-1).                                \tag{GR8}
\]

For precision, this argument uses the enlarged block S direct-sum Y. Its h and h inverse blocks are the identity; all transverse-to-block entries vanish by the original plane conditions and adjointness. Pairing the two nonzero cross images in the nondegenerate packet plane makes them annihilator lines. Therefore the plane part of ht lies in CS, and its isotropic contribution has zero norm. This proves the scalar norm-one identity globally over E. The G-integral change of basis used for DO094's finite statements is unnecessary for GR8.

The exact integer computations are

\[
\begin{split}
 \gcd(U,M)&=9,\\
 \frac{36U}{M^2}
   &=\frac{17444108816482745329931488530262449182555607}
           {74939403110506230985043789733652235885301601}
     <1.
\end{split}                                                   \tag{GR9}
\]

Thus GR1 holds. Apply the lemma to GR7--8 to obtain

\[
                      \boxed{\alpha=1,\quad\delta_z=0.}       \tag{GR10}
\]

This is an exact global vanishing statement for every actual arrow in the stated residual. No archimedean cutoff estimate, finite Green restoration, error-budget allocation, or enlargement of the permitted source is used. The local hypotheses of DO094(FT2,FT7a--b) are not premises of GR10; only its global coordinate and pairing identities are used.

\subsection{3. The smaller remaining space}\label{the-smaller-remaining-space}

Let V be the full positive three-space perpendicular to u in the original four positive W0 coordinates. Then V=Y direct-sum Et. DO093 gives N(S direct-sum Y)=CS, while GR8 and GR10 give Nt in CS. Consequently

\[
 N(S\oplus V)=CS,\qquad
 \operatorname{rank}_E(N|_{S\oplus V})=1.                    \tag{GR11}
\]

The space S direct-sum V is positive definite and has dimension167. Therefore

\[
 \boxed{\dim_E(\ker(h-I)\cap(S\oplus V))=166,\qquad
                  \operatorname{rank}_E(h-I)\le3.}           \tag{GR12}
\]

The first equality uses C nonzero. The second also follows directly from rank-nullity in the169-dimensional original space. Since h is an isometry and fixes that nondegenerate positive166-space pointwise, it preserves its orthogonal complement. The remaining action has dimension three and signature(2,1) at every complex label.

The fixed space may depend on h. The individual vectors y1,y2,t,z need not be fixed; their N-images can be nonzero multiples of the same isotropic cross line. No integral basis of the moving fixed space, controlled-denominator identification with a standard rank-three lattice, global nilpotence, or classification of the remaining three-dimensional action is asserted.

Thus the exact support23/24 scalar-cross gate is smaller than DO093's rank-at-most-four, positive165-fixed-space gate: the entire last nontrivial compact transverse scalar is eliminated, and the fixed positive dimension increases by one. Other structural classes and other supports remain governed by their existing statements. The signed sufficient inequality, actual global Bessel/Mram nonvanishing, selected arithmetic image and content, source constructor, \nolinkurl{Read}/specific032 transport and DP44 integer cost remain open. The ordered error lists stay at75 all-top and31 uniform charges, with complete support floors23 and11.

\subsection{4. Algebraic-geometric placement and verification boundary}\label{algebraic-geometric-placement-and-verification-boundary}

The new step is an integral norm argument over Spec(O\_F) for the scalar pairing section GR7 on the original Hermitian modules over Spec(O\_E). GR4 writes the quotient explicitly as an integral trace-and-norm expression. Its field norm is the determinant of multiplication by nu; positivity at all real realizations and integrality force its vanishing. The resulting fixed space is the kernel of the actual rank-one map GR11. No moving integral frame is introduced.

The CM Jacobian, its Picard data, theta line, torsion rigidifications and biextensions are unchanged. Carried exact contextual anchors are FGA232 Theorem3.1 pp148--149; FGA236 Corollaries1.7/2.3 pp229/232; SGA7 I VII2.7/2.9 pp169--174 and VIII7.4(b) pp310--311; SGA3 I4.7.3/4.7.3.1; SGA6 VI Theorem3.3; cote4 scan003 §§8.1.1--8.1.2, scans005--006 §§8.4.1--8.4.3; and cote54 scans005--007. These are the already carried frameworks, not newly invoked theorems for GR1--5. No new lambda-ring or K-theory equality, l-adic sheaf/weight theorem, CM-motive comparison, or motivic L-function assertion is inferred.

DP\_MULT Q3e and THETA3 G1/A7 gain the exact smaller scalar-cross remainder GR12. Their actual arithmetic comparisons and values remain open. CX\_GRH CX-ALG-024 and GCT\_PLETH T-052 keep their distinct arithmetic gates; GR10 gives no RH theorem or oscillator-charge comparison.

The companion exact certificate verifies the literal integer inputs, Bezout coefficients, rational inequality, original orthogonality and norm identities, strict-boundary example, and unchanged full error totals. It does not enumerate rational automorphic arrows. The proof consists of the complete norm argument above plus the precisely cited, separately whole-checked actual pairing and plane identities.

\clearpage

\section{AM-LIFT-077 --- exact global calibration of the prescribed Bessel coupling}\label{am-lift-077-exact-global-calibration-of-the-prescribed-bessel-coupling}

\textbf{Canonical checked integration --- 2026-11-22T22:01:23.215189+00:00.} Complete root and independent proof/script/interface readings and separate exact replays passed. Historical author/pending wording below is preserved with its original review history. For an existing genuine initial/bottom embedded cuspidal pair in the LIFT070/LIFT076 scopes and the prescribed finite vectors, GC17 gives \textbar b\textbar²/(nT\emph{nF)={[}aS²/(aH\emph{aG}a0\emph{d\_sigma}2\^{}(t\_pi+1)){]}}Lratio\^{}(G union infinity)(1/2)\emph{product\_G d\_varsigma\^{}(-1)}product\_T q/{[}(q−1)(q\^{}k−1){]}. The formula includes b=0, retains the conjugated bottom slot and physical character, and uses exactly one formal degree, whole-reference Haar ratio and component factor. GC19 gives the corresponding source Gram lower bound after nF cancels. Central nonvanishing, its size, source existence, actual arithmetic Read/032 matching and DP44 costs remain open; no finite-kernel charge follows.

\textbf{Algebraic-geometric translation.} This is an analytic period identity for the existing unitary groups and physical Bessel subgroup, with all global/reference Haar and representation embedding norms distinguished. It identifies a central automorphic L-ratio; it does not identify that ratio with the motive of the specified CM Jacobian or construct the required theta-cycle comparison. Exact carried contextual anchors are SGA3 I4.7.3/4.7.3.1; SGA6 VI Theorem3.3; FGA232 Theorem3.1 pp148--149; FGA236 Corollaries1.7/2.3 pp229/232; SGA7 I VII2.7/2.9 pp169--174 and VIII7.4(b) pp310--311; cote4 scan003 §§8.1.1--8.1.2, scans005--006 §§8.4.1--8.4.3 and cote54 scans005--007. No new scheme isomorphism, lambda-ring/K-theory identity, l-adic weight or CM-motive theorem is claimed. DP\_MULT Q3e and THETA3 G1/A7 gain the exact conditional coupling normalization, while their central nonzero/arithmetic value gates remain. CX\_GRH CX-ALG-024 and GCT\_PLETH T-052 remain separate.

Author: \nolinkurl{/root/lift_accounting}, 22 November 2026 UTC.

\textbf{Status: author whole-proof/source-interface/script review PASS; exact diagnostic PASS (16,555 assertions). Independent and root reviews remain separate.} This is an equality conditional on the existence of the specified genuine source data and the two specified embedded cuspidal representations. The global period is allowed to be zero. Nothing below proves a central value nonzero, constructs the requested arithmetic \nolinkurl{Read}/\nolinkurl{STiesR}/specific032 instance, or provides a numerical DP44 bound. No canonical file or earlier frozen finding is modified.

\subsection{1. The actual pair and the scope of the theorem}\label{the-actual-pair-and-the-scope-of-the-theorem}

Use the genuine hypotheses and actual vector/character dictionaries of LIFT074, except that its global distinction premise is \textbf{not} assumed. The initial compact unitary group is H=U(k), the bottom group is G=U(k+rho,rho+1), and

\[
 r=k+2\rho+1=k+m-1,\qquad \rho=m/2-1.
 \tag{GC1}
\]

Here k and m are even, and r is odd. The full actual SizesAt/SizeData2 scope of LIFT070 supplies temperedness, genericity at the split finite places, and a generic discrete global parameter for every occurring initial constituent pi. The full genuine goodTuple/e0 scope of LIFT076 supplies temperedness and a \textbf{simple} generic global parameter for every occurring admitted bottom constituent Pi. These are the proved scopes, not an assertion for arbitrary abstract weights or arbitrary finite generic invariants. The initial nonzero column, bottom cuspidal occurrence, exact source \nolinkurl{Full}/\nolinkurl{InHH}, and physical subgroup/character are still actual premises.

Fix the automorphic embeddings, including multiplicity vectors. The Bessel pair in the bilinear source convention is

\[
                 \boldsymbol\pi=\pi\boxtimes\overline\Pi.
 \tag{GC2}
\]

The conjugate second representation is essential. Keep the actual strict Mram vector at every selected T place, the permitted bottom vector in the same admitted finite G type supplied by LIFT072, all other prescribed spherical vectors, and the exact finite translate x\_u. Unit local Hilbert norms are used only to describe these finite vectors; the norms of the embedded automorphic functions are retained below. No vector, level, or additive character is replaced.

The physical marked vector v\_phys has norm −T, with T=pi\_T bar(pi\_T) in the original field. Write h for the physical Hermitian form and choose the actual adapted hyperbolic pairs with h(x\_i,y\_i)=1. Set h\_prime=−h, v0\_prime=v\_phys/pi\_T, x\_i\_prime=pi\_T\^{}−1 x\_i, and y\_i\_prime=−bar(pi\_T)y\_i. Then h\_prime(v0\_prime,v0\_prime)=1 and h\_prime(x\_i\_prime,y\_i\_prime)=1. More precisely, h\_prime(u x\_(i+1)\emph{prime,y\_i\_prime)=h(u x}(i+1),y\_i) and h\_prime(u v0\_prime,y\_rho\_prime)=h(u v\_phys,y\_rho). These equalities hold in either Hermitian linearity convention: the two scalar factors multiply to −1, which cancels the sign of h\_prime. Hence the entire physical additive argument is unchanged. \textbf{There is no conjugation, negation, or replacement of psi.} Simultaneously negate the smaller form as well; both unitary groups and the embedded smaller subgroup are unchanged. This is the rational physical character dictionary of LIFT060/067/068/072/073, not an inference from matching real signatures. It places the pair in Beuzart-Plessis and Chaudouard (2025), §1.3.1--§1.3.2, with their smaller dimension k, their n+1=r, and their flag length rho. All measures and vectors below are transported through this identification, so its local Jacobians are retained.

Let sigma be the actual real bottom discrete series, rho\_H the actual compact real initial representation, and w the literal \nolinkurl{IsLowUnit} vector. Let W be the unit relative vector of LIFT061/061-U, in the source quotient measure nu\_Q:

\[
 W(sy)=\psi_X(s)\rho_H(s)W(y),\qquad
 \int_{S_\infty\backslash G_\infty}\|W(y)\|^2\,d\nu_Q(y)=1.
 \tag{GC3}
\]

As in LIFT074 PD8--10, for the actual initial column T and fixed finite bottom vector define A by

\[
 \langle A(v),\xi\rangle
 =\mathcal B\left(T\xi,
      g\mapsto\overline{F_v(\iota(g)(1,x_u))}\right),\qquad
 A=bA_0,\quad A_0(\sigma(y)w)=W(y).
 \tag{GC4}
\]

Inner products here are conjugate-linear in the first variable. The scalar b exists and is unique even when A=0, in which case b=0. Continuity and multiplicity one are exactly LIFT074 §3; no global distinction is needed for that step. Its source unfolding gives, with the same actual positive scalar,

\[
 \langle F_w,\operatorname{Pu}_u(T)\rangle
   =c_{\rm Per}\,\overline b,
 \qquad c_{\rm Per}=c_1\mu_f(K_2)>0.
 \tag{GC5}
\]

This equality, unlike a nonvanishing assertion, remains true when b=0.

\subsection{2. Global measures, embedded norms, and the reference measure}\label{global-measures-embedded-norms-and-the-reference-measure}

Let dH\_Tam, dG\_Tam and dS\_Tam be the Tamagawa measures on the actual rational groups and their Bessel semidirect product, transported through the preceding rational identification. Write

\[
 dH_{\rm emb}=a_HdH_{\rm Tam},\quad
 dG_{\rm emb}=a_GdG_{\rm Tam},\quad
 dS_{\mathcal B}=a_SdS_{\rm Tam},
 \qquad a_H,a_G,a_S>0.
 \tag{GC6}
\]

Here dH\_emb is the measure used to take the norm of the scalar initial form T xi, dG\_emb is the actual bottom automorphic Petersson measure, and dS\_B is the measure in the literal bilinear period. They are not silently identified with a probability measure. Define the exact embedded norm scalars

\[
 n_T=\frac{\|T\xi\|_{[H],dH_{\rm emb}}^2}{\|\xi\|^2}>0,
 \qquad
 n_F=\frac{\|F_v\|_{[G],dG_{\rm emb}}^2}{\|v\|^2}>0.
 \tag{GC7}
\]

They are independent of the nonzero xi and v by invariant Hermitian forms on the fixed irreducible real factors. For the bottom, this is the Hilbert realization of the fixed embedded automorphic occurrence. Right translation by x\_u is unitary and does not change n\_F. No identification of n\_T with a source Cu norm, a Hilbert--Schmidt norm, or a compact dimension is made.

At each finite place use \textbf{the reference Haar of the checked local test}, pulled back to the actual S\_v: probability measure on the indicated compact in H\_v and the specified integral root-coordinate measure on N\_v. At G this is the conductor-p convention of LIFT072, with O volume one, not a new selfdual convention; at T it is LIFT071; elsewhere it is the actual primitive/different-scaled hyperspecial convention of LIFT073. Put

\[
 dS^0=dS_{\infty,\rm Wu}\prod_{v<\infty}dS_v^0
       =a_0dS_{\rm Tam},\qquad a_0>0.
 \tag{GC8}
\]

The real factor is exactly the S\_infty Haar used with the source G\_infty Haar to define nu\_Q in GC3. These local measures form a restricted product; outside finitely many places the underlying group, flag lattice, character and coordinate measures are integral. The positive number a\_0 is the ratio of this \textbf{whole product} to Tamagawa measure. It includes the rational-volume-form and Artin--Tate convergence normalization. It is not a product of unspecified compact volumes that is declared to equal one.

The source permits genuinely independent normalizations at the following interface. \nolinkurl{Measures.lean}, \nolinkurl{MeasG.IsTextNorm} and \nolinkurl{MeasW.IsTextNorm}, require Haar measures and their archimedean/finite product decompositions, but do not impose Tamagawa normalization. \nolinkurl{IsThetaFam.lean} defines \nolinkurl{Vol} as a Haar measure with a fundamental domain, and \nolinkurl{VolN} through \nolinkurl{IsUnitN}. The latter requires the actual unipotent quotient to have mass one. \nolinkurl{BesLev.per} uses its own vA and vN, and \nolinkurl{SBes.tieN/tieH/tieψ} identifies groups and characters, not all of these measures with the Wu real factors. Thus GC6--8 retain real information in the literal source.

If dH\_emb is chosen to be that same vA measure, the standard product Tamagawa normalization on N⋊H and the unit quotient normalization on N imply a\_S=a\_H. This optional simplification is not needed below; the general formula deliberately retains both symbols. The product rule follows in rational unipotent coordinates by successive additive quotients, whose Tamagawa quotient mass is one, and by the determinant-one H action on the paired N roots. No normalization of the H quotient mass is inferred.

\subsection{3. Exact real Bessel factor: the Wu seed energy disappears}\label{exact-real-bessel-factor-the-wu-seed-energy-disappears}

Let d\_sigma be the true formal degree for the actual source Haar on G\_infty. LIFT061 uses

\[
 \phi(g)=d_\sigma\overline{\langle w,\sigma(g)w\rangle},
 \quad P=R(\phi),\quad
 K_P(x,y)=\int_{S_\infty}\phi(x^{-1}sy)
                        \psi_X(s)\rho_H(s)\,dS_{\infty,\rm Wu}(s).
 \tag{GC9}
\]

Its kernel is absolutely convergent and continuous for these discrete-series coefficients. LIFT061 RW20, with the explicit conjugation converting its linear-first reference to the source convention, identifies the diagonal with d\_sigma times the real Bessel matrix period of GC2. LIFT061-U gives the exact rank-one factorization

\[
 K_P(x,y)\xi=W(x)\langle W(y),\xi\rangle.
 \tag{GC10}
\]

Consequently, on v= sigma(y)w and any compact vector xi,

\[
 \boxed{\quad
 \alpha_\infty^0(v,\xi)
      =\frac{|\langle W(y),\xi\rangle|^2}{d_\sigma}.
 \quad}
 \tag{GC11}
\]

The diagonal is real, so the single conjugation in RW20 introduces no extra sign or scalar. Formula GC11 refers to the product of the real local matrix periods, or equivalently to the period on the whole real product. The finite-dimensional compact factor contributes its actual inner product; there is \textbf{no factor dim(rho\_H)}. A useful independent check is the operator identity

\[
          \int\operatorname{tr}K_P(y,y)\,d\nu_Q(y)=1.
 \tag{GC12}
\]

Choose y and a pure unit xi for which the entry in GC11 is nonzero. Such a choice exists: W has norm one, is continuous, and pure vectors span the finite-dimensional tensor product. The vector sigma(y)w has norm one. If the literal minimal vector w is not itself presented as a pure tensor over the real places, first polarize BPC's diagonal identity successively in every local slot. Pure local tensors span the finite-dimensional minimal K-type, so the resulting sesquilinear identity applies to every finite sum there and to its fixed translates. Equivalently it is the full real matrix period in GC11. Thus no additional purity premise on w is used. This vector is used to evaluate the identity and cancel the nonzero common factor; it is \textbf{not} substituted for the prescribed vector F\_w in GC5. This avoids an assumption W(1) != 0.

Neither the positive seed energy E in LIFT061's construction nor its auxiliary bump appears in GC11. The unit W and its invariant norms still have the costs recorded there. The present argument calibrates the coupling b, rather than claiming an improved sup norm or weighted L1 bound for W.

\subsection{4. The global L-ratio and the exact finite factors}\label{the-global-l-ratio-and-the-exact-finite-factors}

Write B\_H=BC(pi) and B\_G=BC(overline Pi). Let t\_pi be the number of distinct cuspidal blocks of the generic initial parameter. LIFT076 makes B\_G simple. In the component-group convention actually used by Beuzart-Plessis and Chaudouard, §§1.1.2 and1.1.4,

\[
             |S_{B_H\boxtimes B_G}|=2^{t_\pi+1}.
 \tag{GC13}
\]

This is not divided by a diagonal central factor. Define the completed ratio

\[
 \mathcal L(s)=
 \left(\prod_{i=1}^{r}L_F(s+i-\tfrac12,\eta_{E/F}^{i})\right)
 \frac{L_E(s,B_H\times B_G)}
 {L_F(s+\tfrac12,B_H,\operatorname{As}^{(-1)^k})
  L_F(s+\tfrac12,B_G,\operatorname{As}^{(-1)^r})}.
 \tag{GC14}
\]

This is precisely their §1.3.6 ratio: As+ on the even smaller rank and As− on the odd larger rank, namely the regular opposite-pole-parity denominator. The conjugated bottom slot is retained. The ratio is interpreted by analytic continuation, not by an unjustified Euler product at its boundary. The discrete-parameter normalization is holomorphic at the central point; the opposite-pole-parity Asai denominator is finite and nonzero, as in the discrete specialization of their §1.2.4 and §1.3.6. Each local ratio at 1/2 is finite and nonzero under the temperedness hypotheses.

For a set V containing the real places, define the partial ratio only by

\[
 \mathcal L^{V}(\tfrac12)
   =\mathcal L(\tfrac12)
            \prod_{v\in V}\mathcal L_v(\tfrac12)^{-1}.
 \tag{GC15}
\]

In particular deleting the real factors means dividing by their exact gamma factors once. Let d\_varsigma\_v be the dimension of the actual finite smaller generic representation in LIFT072. On the fixed Hilbert-unit finite vectors the reference local factors are

\begin{longtable}[]{@{}
  >{\raggedright\arraybackslash}p{(\columnwidth - 4\tabcolsep) * \real{0.2727}}
  >{\raggedleft\arraybackslash}p{(\columnwidth - 4\tabcolsep) * \real{0.3636}}
  >{\raggedleft\arraybackslash}p{(\columnwidth - 4\tabcolsep) * \real{0.3636}}@{}}
\toprule\noalign{}
\begin{minipage}[b]{\linewidth}\raggedright
Finite places
\end{minipage} & \begin{minipage}[b]{\linewidth}\raggedleft
Raw matrix period alpha\_v\^{}0
\end{minipage} & \begin{minipage}[b]{\linewidth}\raggedleft
Normalized factor alpha\_v\^{}0 / L\_v(1/2)
\end{minipage} \\
\midrule\noalign{}
\endhead
\bottomrule\noalign{}
\endlastfoot
G & 1/d\_varsigma\_v & 1/(d\_varsigma\_v L\_v(1/2)) \\
T & L\_v(1/2) q\_v/{[}(q\_v−1)(q\_v\^{}k−1){]} & q\_v/{[}(q\_v−1)(q\_v\^{}k−1){]} \\
Outside G∪T & L\_v(1/2) & 1 \\
\end{longtable}

Thus

\[
 \prod_{v<\infty}\alpha_v^{0,\natural}
  =\prod_{v\in G}\frac1{d_{\varsigma_v}\mathcal L_v(\tfrac12)}
    \prod_{v\in T}\frac{q_v}{(q_v-1)(q_v^k-1)}.
 \tag{GC16}
\]

The finite product is legitimate: almost all normalized factors are one, as also stated in BPC Remark1.3.6.2. The G value is an \textbf{unnormalized} test, whereas the T rational number is already \textbf{L-normalized}. Removing the T L-factor again would be an error. For local source measures c\_H,v c\_N,v times the displayed reference measures, multiply the corresponding local alpha by that product. In GC8 those changes are instead collected in the whole-reference ratio a\_0; they must not also be multiplied into GC16 unless a\_0 is adjusted accordingly.

\subsection{5. The complete conditional equality}\label{the-complete-conditional-equality}

\textbf{Theorem.} Under §1's genuine existing-pair and prescribed-vector premises, without a global distinction premise,

\[
 \boxed{\displaystyle
 \frac{|b(T,b_f)|^2}{n_Tn_F}
  =\frac{a_S^2}
    {a_Ha_Ga_0\,d_\sigma\,2^{t_\pi+1}}
    \mathcal L^{G\cup\infty}(\tfrac12)
    \prod_{v\in G}\frac1{d_{\varsigma_v}}
    \prod_{v\in T}\frac{q_v}{(q_v-1)(q_v^k-1)}.
 }\tag{GC17}
\]

\textbf{Proof.} Apply BPC Theorem1.3.6.1, polarized in the real slots as just explained, to the product of the two actual embedded forms, with real inputs v=sigma(y)w and xi chosen as in §3, and with the prescribed finite pair including x\_u. Its global Tamagawa Petersson norm is n\_T n\_F/(a\_H a\_G). Its Tamagawa period is the source period divided by a\_S, whose squared modulus is \textbar b\textbar² \textbar\textless W(y),xi\textgreater\textbar²/a\_S² by GC4. The theorem applies even if that period vanishes.

The local Bessel measures in the theorem must multiply to dS\_Tam. The reference product GC8 instead multiplies to a\_0 dS\_Tam. Rescale one local reference measure by a\_0\^{}−1, or use linearity in Haar directly: the product of local matrix periods in the theorem is a\_0\^{}−1 times the product computed with GC8. This is a comparison of formulas, not a change of the source period or of its prescribed finite vectors. Its real factor is GC11 divided by L\_infty(1/2), and its finite factors are GC16. Hence the theorem gives

\[
 \frac{a_Ha_G}{a_S^2n_Tn_F}
     |b|^2 |\langle W(y),\xi\rangle|^2
 =\frac{\mathcal L(\tfrac12)}{2^{t_\pi+1}a_0}
   \frac{|\langle W(y),\xi\rangle|^2}
      {d_\sigma\mathcal L_\infty(\tfrac12)}
   \prod_{v<\infty}\alpha_v^{0,\natural}.
 \tag{GC18}
\]

Cancel the nonzero real entry and substitute GC16. This proves GC17. The cancellation did not divide by b or by a central value, so it remains valid when either is zero. ∎

The noncentral local factors, embedded norms and measure factors are nonzero. Thus, for this fixed existing pair, the conditional implication from a nonzero central L-ratio to b != 0 is now an exact consequence of GC17, with the original vectors. Equivalently, since GC14's other central factors are finite and nonzero, it is the central Rankin--Selberg value that remains to be shown nonzero. This statement does not prove that value is nonzero and does not invoke an existence theorem for a different inner form or automorphic embedding. If the central value vanishes, GC17 says b=0 for these vectors.

Combining GC17 with the exact source PerLaw yields an actual lower bound with the row embedding norm canceled:

\[
 K_{u,u}(T,T)\ge
 \frac{c_{\rm Per}^2 n_T}{c_0}
 \frac{a_S^2}{a_Ha_Ga_0d_\sigma 2^{t_\pi+1}}
 \mathcal L^{G\cup\infty}(\tfrac12)
 \prod_{v\in G}d_{\varsigma_v}^{-1}
 \prod_{v\in T}\frac{q_v}{(q_v-1)(q_v^k-1)}.
 \tag{GC19}
\]

Here K is the literal normalized Gram of LIFT074 PD14. Equality GC17 guarantees that the displayed right side is real and nonnegative in this compatible pair. It is positive conditional on the stated central nonvanishing. Formula GC19 retains c\_0, c\_Per and n\_T; it does not identify the latter with an arithmetic coefficient norm, nor give a uniform numerical inverse or DP44 estimate.

\subsection{6. Homogeneity and omitted-factor checks}\label{homogeneity-and-omitted-factor-checks}

These checks are part of the normalization proof, not optional conventions.

\begin{enumerate}
\def\labelenumi{\arabic{enumi}.}
\tightlist
\item
  Rescale the actual global H, G and period S measures by positive h,g,s while keeping the real reference model fixed. Then n\_T,n\_F,b² scale by h,g,s². The normalized left side of GC17 scales by s²/(hg), as does a\_S²/(a\_Ha\_G).
\item
  Rescale a single finite reference S\_v Haar by z\textgreater0. Its matrix period and a\_0 both multiply by z. The unsimplified right side of GC18 is unchanged. The numerical table GC16 is always tied to its stated reference Haar; it cannot be retained unchanged after this rescaling.
\item
  Rescale the source real G\_infty Haar by u\textgreater0 and the Wu S\_infty Haar by v\textgreater0, with the corresponding Weil quotient. Then d\_sigma becomes d\_sigma/u, nu\_Q becomes (u/v)nu\_Q, the unit W becomes sqrt(v/u)W, and b becomes sqrt(u/v)b. Meanwhile a\_0 becomes va\_0. Therefore 1/(a\_0 d\_sigma) and \textbar b\textbar² both multiply by u/v. This shows why neither a\_0 nor d\_sigma may be dropped.
\item
  Under scalar changes T -\textgreater{} tT and F -\textgreater{} zF, LIFT074 PD15 gives b -\textgreater{} bar(t) z b; n\_T -\textgreater{} \textbar t\textbar²n\_T and n\_F -\textgreater{} \textbar z\textbar²n\_F. Thus GC17 is independent of these embedded-vector scales, while GC5 retains their phases.
\item
  With the global measures scaled as in (1) and real quotient scaled by u/v, the actual disintegration gives c\_Per -\textgreater{} g v/(s u) c\_Per. Consequently GC5 and GC19 have the correct source Petersson scaling. In particular rescaling only the auxiliary Bessel vA measure cancels from the resulting Pu pairing after its literal c\_Per is transported.
\item
  There is one factor d\_sigma\^{}−1, one a\_0\^{}−1, and one component factor 2\^{}−(t\_pi+1). There is no dim(rho\_H), no second finite G L-factor, and no removal of the T Euler factor twice. The squared norm n\_F cancels only in GC19, not in GC17.
\end{enumerate}

The finite-group diagnostic uses S\_3 inside S\_4, a three-dimensional irreducible ambient representation and a two-dimensional compact representation, to check the actual rank-one kernel/quotient trace normalization with nontrivial compact dimension. It also checks the displayed Haar and scalar homogeneities and finite-factor cancellation with exact rational arithmetic. These diagnostics do not prove BPC's theorem or its analytic hypotheses.

\subsection{7. Source and primary verification boundary}\label{source-and-primary-verification-boundary}

The exact source addresses retained for this comparison are:

\begin{itemize}
\tightlist
\item
  \nolinkurl{Formalism/LocalToGlobal/SiegelLatticeTheory/CokernelGLCounts/Measures.lean}, \nolinkurl{MeasG.IsTextNorm} and \nolinkurl{MeasW.IsTextNorm}: Haar and product decomposition only.
\item
  \nolinkurl{Development/AbelianVarietyTheory/Definitions/SemilocalFrames/IsBesPosOfBesselTied.lean}, lines23--25: \nolinkurl{IsUnitN}, with fundamental-domain mass one.
\item
  \nolinkurl{Development/AbelianVarietyTheory/Definitions/LatticeFrames/IsThetaFam.lean}, lines18--27 and65--69: actual vA/vN measures, bilinear period and semidirect-product measure.
\item
  \nolinkurl{Development/AbelianVarietyTheory/Definitions/TwistedInnerProducts/WittSmallGroupEmbeddingMemPointwiseStabilizerIso.lean}, lines83--87 and176--182: semidirect product measure and exact character convention.
\item
  \nolinkurl{Development/AbelianVarietyTheory/Definitions/LatticeFrames/ThetaIdentityOfMapsAll.lean}, \nolinkurl{SBes} and \nolinkurl{PerLaw}: actual subgroup/character pins and source integral.
\item
  \nolinkurl{Formalism/LocalToGlobal/SiegelLatticeTheory/HaarTransversals/HidaSetting.lean}, \nolinkurl{IsQuotMeasS}, \nolinkurl{nuQ_weil}, \nolinkurl{Wu_spec}, \nolinkurl{measG_norm} and \nolinkurl{measW_norm}: the actual real Haar/formal-degree and quotient measure.
\item
  LIFT074's freshly source-checked \nolinkurl{FoldIntegralBesselQuotientSupportedFunctions.lean}, \nolinkurl{UnfoldedPeriodEqIntegralBesselPeriods.lean}, and \nolinkurl{ConjugateRowPairingPoincareSeriesBesselPeriods.lean}: the positive c\_Per and its unchanged unfolding.
\end{itemize}

Primary reference: Beuzart-Plessis, R., \& Chaudouard, P.-H. (2025). The global Gan--Gross--Prasad conjecture for unitary groups. II. From Eisenstein series to Bessel periods. \emph{Forum of Mathematics, Pi, 13}, e16. \url{https://doi.org/10.1017/fmp.2025.8}. Directly read: §§1.1.2/1.1.4 (component group and its product), §1.2.4 (regular L-normalization), §§1.3.1--1.3.2 (physical Bessel pair), §§1.3.5--1.3.6 and Theorem1.3.6.1 (the proved formula), and §2.1.8 (Tamagawa measures). In the publisher PDF these occur on pages2,4--9,14. The exact publisher PDF/text are retained under \nolinkurl{global_period_refs/} with their hashes in the companion manifest.

The real projector identities are the already checked LIFT061 RW17--21 and LIFT061-U; the prescribed finite factors are LIFT071 HN22--24, LIFT072 GB2, and LIFT073's theorem. Their accepted proofs are used as inputs, with no new claim that the present rational diagnostics verify them. LIFT070/LIFT076 supply the stated parameter/temperedness scopes. The new result is their exact compatible global calibration GC17, including zero periods and all source measure and embedding scales.

Still open: existence of a fully instantiated actual source package at the requested parameters, central nonvanishing for the specified pair, its quantitative lower bound, the literal arithmetic \nolinkurl{Read}/\nolinkurl{STiesR}/tS/specific032 transport, selected arithmetic image/content, and the full DP44 comparison and budget. No conclusion about RH is drawn.

\clearpage

\section{AM-R-072 --- a two-coordinate packet-plane column and the joint scalar row}\label{am-r-072-a-two-coordinate-packet-plane-column-and-the-joint-scalar-row}

\textbf{Canonical checked integration --- 2026-11-22T22:24:33.897997+00:00.} Complete root and independent whole-proof/script/interface readings and separate350-assertion replays passed. At the unchanged actual all-top H133 endpoint, R071 scalar-cross support23 terms with at least six nonzero original S scalar coordinates have total absolute contribution\textless2\^{}-3600. The full positive167 compression supplies three original integral transverse coefficients, while the exact negative-row bound is H/(2sqrt U). Spherical-cap/Gaussian packing counts the first row jointly. The fixed packet-plane projection retains its U denominator, exact paired target indexT\^{}(4s), and positive-rank-two source fibre. All159 k-cells6..164, all source/target types, actual selected trace, class weights and finite projectors are retained. The remaining R071 support23 cross row has k\textless=5. No support24 removal, global-unipotence inference or complete support-floor increase follows; signed/global/source/arithmetic/DP44 gates remain open.

\textbf{Algebraic-geometric translation.} The fixed rational projection of the original Hermitian module has image (U\^{}-1 O\_E)u plus O\_E e\_n.~Its denominator and determinant index are retained, and the G/T frame is proved invertible only at the required primes. The cross ideal is a paired determinantal condition over Spec(O\_E), and the joint count uses the unchanged complex half-trace metric. Exact contextual anchors remain SGA3 I4.7.3/4.7.3.1; SGA6 VI Theorem3.3; FGA232 Theorem3.1 pp148--149; FGA236 Corollaries1.7/2.3 pp229/232; SGA7 I VII2.7/2.9 pp169--174 and VIII7.4(b) pp310--311; cote4 scan003 §§8.1.1--8.1.2, scans005--006 §§8.4.1--8.4.3 and cote54 scans005--007. No new lambda-ring/K-theory equality, l-adic weight theorem or motivic L-function assertion is added. DP\_MULT Q3e and THETA3 G1/A7 gain the smaller finite remainder; CX\_GRH CX-ALG-024 and GCT\_PLETH T-052 retain their arithmetic gates.

Author: /root, GPT-6 Astra Pro, ALP\_MACHINE. 22 November 2026 UTC.

\textbf{Pending whole-proof review.} This is a complete new counting argument for the actual R071 remainder. Its proof, interfaces and exact arithmetic certificate are checked separately before promotion. No charge is applied by writing this file.

\subsection{1. Actual sector and conclusion}\label{actual-sector-and-conclusion-3}

Keep the unchanged J56 all-top endpoint, H=133, f=168, n=169, every original source/target pencil, finite projector, central representative, selected metric, class mass and phase. In R071(GR6--12)'s doubly-plane, four-zero remainder at support s=23, put

\[
 Z=S\perp V,\quad \dim S=164,\quad
 V=u^\perp\subset E^4,\quad P=Eu\perp Ee_n,\quad
 E^{169}=Z\perp P.                                           \tag{PJ1}
\]

Here Z is positive of dimension167 and P has Gram matrix diag(U,−1) in the fixed basis (u,e\_n). R071 proves h\_ZZ=I, the full cross map C:Z→P has rank one, and its restriction to the original S is nonzero. Write

\[
 Ce_i=(u\xi_i,\beta_i),\qquad
 k=\#\{i\in S:\xi_i\ne0\},\qquad 1\le k\le164.                 \tag{PJ2}
\]

\textbf{Theorem.} The entire current-remainder intersection satisfying s=23 and 6≤k≤164 has total absolute normalized contribution less than

\[
                              2^{-3600}.                     \tag{PJ3}
\]

Every incidence profile is included, with no global-unipotence premise. Thus any still unremoved R071 scalar-cross term at support23 has at most five nonzero original S scalar coordinates. This adds one all-top charge upon checked integration, but does not by itself raise the complete support floor23 or remove a support24 term.

The literal integers are Q=1427908193729121522601, U=77Q²−14 and

\[
\begin{split}
u&=(2,95192390527,383092671043,12529843548438452388609),\\
y_1&=(-976627,-153032,38026,0),\\
y_2&=(-567678,695570,-172838,0),\\
z_0&=(256766,-7654233287,-30805128044,1).
\end{split}                                                   \tag{PJ4}
\]

These three integral transverse vectors belong separately to every source and target pencil, as established in DO094(FT8--11) and R071(GR6--7). Their squared lengths are

\[
 L_1=978665066829,\quad L_2=835948910828,\quad
 L_3=1007543201084982869062.                                   \tag{PJ5}
\]

The fixed four-vector frame (y1,y2,z0,u) has determinant of modulus2U and transverse Gram determinant4U. Since gcd(2U,g0g1T)=1, it is an integral invertible frame at the actual G and T primes. It is not asserted to be an integral frame away these primes.

\subsection{2. The full cross map, its common ideal and its actual norm}\label{the-full-cross-map-its-common-ideal-and-its-actual-norm}

Write Cx=(uλ(x),ηλ(x)), where η=β\_j/ξ\_j for the first nonzero original coordinate ξ\_j. Isotropy gives η bar(η)=U. The complete map C is determined injectively by the164 original scalars ξ\_i, the three scalars λ(y1), λ(y2), λ(z0), and β\_j. The quotient is used only for reconstruction over E. It is not assumed integral. We retain the harmless factor164\textless2\^{}8 for the pivot.

At an unmixed G reduction of rank one, N=h−I is square zero. In the G-unit fixed Z/P frame write N=ab\^{}t. The identity Z block gives either a\_Z=0 or b\_Z=0. In the second case the whole C vanishes at that component; in the first case the whole upper cross block vanishes, and the paired inverse-adjoint relation makes C vanish at the conjugate component. At rank zero both cross blocks vanish at both components. Choose a vanishing component by a fixed rule at rank one, and both at rank zero. Let J be the resulting E-ideal. There are at most4\^{}334=2\^{}668 patterns. Exactly the rank0/1/at-least2 Green comparison of R059(GB9--10) gives

\[
 |a(h)|\le2Q^{-335(f-1)}(N_EJ)^{335/2}.                         \tag{PJ6}
\]

Its proof works for the enlarged identity block because it uses only the displayed rank-one factorization and the integral paired form. No positive signature is assigned to a finite field.

Set L=J I\_Σ, N\_EI\_Σ=T\^{}(-2s), and z=(N\_EL)\^{}(1/(2f)). All original ξ\_i, all three transverse λ-values, and β\_j lie in L. At G this follows from the whole-C orientation just proved. Away G,T, primitive integer coefficients expressing1 as a combination of the u\_i extract λ(x) from the original integral coordinates of Cx. At T, each of e\_i,y1,y2,z0 belongs to both source and target lattices; hence Cx=hx−x belongs to the target, which is contained in π\^{}(-1)Λ at each paired component. The same original-coordinate extraction gives at most the one permitted pole. This is an actual fractional ideal statement, with no denominator from a moving frame.

At a complex label, the Euclidean norm of the negative row restricted to the full positive space of a U(168,1) matrix is sinh r, where r=log\textbar\textbar h\textbar\textbar. Thus it is at most (H−H\^{}(-1))/2\textless H/2. Its restriction to Z is ηλ. Consequently, in the original positive metric on Z,

\[
 \|\lambda_\tau\|<M:=H/(2\sqrt U),\quad
 \|\xi_\tau\|<M,\quad
 |\lambda(y_j)_\tau|<M\sqrt{L_j},\quad
 |\beta_{j,\tau}|<B:=H/2.                                    \tag{PJ7}
\]

These are row-norm bounds, not independent coordinate substitutes for that norm. Every nonzero element of L has half-trace length at least sqrt(f)z by the field-norm/AM--GM argument R059(GB14). For the k nonzero original coordinates this implies

\[
                               kz^2<M^2.                      \tag{PJ8}
\]

The actual packet-transformed selected trace still obeys R061(GS11--12), using any nonzero original ξ coordinate:

\[
                  |c(h)|<2^{-929f/2}U^{-f/2}z^{-f}.            \tag{PJ9}
\]

This is the unchanged normalized selected trace with both compact factors, not an operator norm. The stronger bound PJ7 does not change its proof or any normalization.

\subsection{3. A joint scalar-row bound from spherical caps}\label{a-joint-scalar-row-bound-from-spherical-caps}

We give the general packing argument used here. Let a lattice in R\^{}d have minimum nonzero length at least1, with d=336. A set of its points in a ball of radius R≥1 has mutual distance at least1. Embed that ball into the upper hemisphere of the radius-R sphere in R\^{}(d+1) by x↦(x,sqrt(R²−\textbar\textbar x\textbar\textbar²)). This does not decrease pairwise distances. The spherical caps of geodesic angular radius α=arcsin(1/(2R)) around the embedded points have disjoint interiors.

Write I\_d=∫\_0\^{}π sin\^{}(d−1)(t)dt. A cap occupies at least (sin α)\^{}d/(d I\_d) of the sphere: indeed ∫\_0\^{}α sin\^{}(d−1)(t)dt≥(sin α)\^{}d/d, by differentiating sin\^{}d.~Therefore its ball count is at most d I\_d(2R)\^{}d.~For d=336 the exact odd-power integral gives

\[
 I_{336}=2\prod_{j=1}^{167}\frac{2j}{2j+1},\qquad
                    336 I_{336}<47.                          \tag{PJ10}
\]

Scale L by sqrt(f)z. For t\textgreater0 let Θ*(t) be its Gaussian sum with the zero vector omitted. Its nonzero vectors have length at least1. Integrating the ball count against 2tR exp(−tR²) on R≥1, then enlarging the integration interval to R≥0, gives

\[
             \Theta^*(t)<47\,2^{2f}\, f!\,t^{-f}.             \tag{PJ11}
\]

This uses the full lattice only through its actual minimum and real dimension. There is no equidistribution, covolume or ideal-covering assumption. For k nonzero elements whose total squared scaled length is less than M²/z², multiply PJ11 k times and use the Gaussian weight. Choosing t=fkz²/M² yields the complete row count

\[
 \#\xi_k\le {164\choose k}\,47^k2^{2fk}(f!)^k
       \left(\frac{eM^2}{fkz^2}\right)^{fk}.                  \tag{PJ12}
\]

The binomial chooses the original nonzero positions. The elementary exponential series with its geometric remainder gives e\textless87/32, as checked by exact rational arithmetic in the companion certificate. Replace e by87/32 in all following upper bounds.

For the other three λ-values and the pivot β use the scalar packing bound (1+2R/z)\^{}(2f) from R059(GB14). On PJ8, each of their radius-to-z ratios is greater than2\^{}19, because sqrt(L\_j)\textgreater2\^{}19 and sqrt U\textgreater2\^{}19. Put c=2+2\^{}(-19). Their complete joint upper count is therefore

\[
           c^{8f}M^{6f}B^{2f}(L_1L_2L_3)^f z^{-8f}.           \tag{PJ13}
\]

The first row is counted jointly in PJ12; only these four additional scalars use the separate packing estimate.

\subsection{4. The true projected two-dimensional column lattice}\label{the-true-projected-two-dimensional-column-lattice}

Let pr\_P be the fixed orthogonal projection W0→P along V. Its projection of the original lattice has the exact description

\[
 \Lambda_P^{\rm proj}=\operatorname{pr}_P(O_E^5)
       =(U^{-1}O_E)u\oplus O_Ee_n.                            \tag{PJ14}
\]

Primitivity of u proves equality: in coordinates (w,b) with a plane vector equal to (w/U)u+b e\_n, the positive coordinate is w=u·x and runs through all O\_E. The half-trace covolume is D0²/U\^{}f, where D0=7\^{}140 29\^{}144 5\^{}84 13\^{}84. This fixed denominator at primes dividing U is retained below; it is not discarded by calling a rational projection integral.

Let A\_d be the actual five-coordinate target overlattice π\^{}(-1)M\_d of R051 §1, with its separate paired local definitions. It contains O\_E\^{}5. At an active selected T component, U is a unit, so the standard lattice splits orthogonally as Λ\_Z⊕Λ\_P, with Λ\_P=O\_Eu⊕O\_Ee\_n.~Here Λ\_Z denotes the transverse part within W0; the original S factor is separate and unchanged. The literal shear is H\_d=I+dπ\^{}(-1)zq, z=v0/bar(π) in P, qz=0, and q=−v0*J−πm0. For x in Λ\_Z, qx=−πm0x. Thus H\_dx−x is an integral multiple of z. Moreover z is in Λ\_P and H\_dz=z, so z belongs to H\_dΛ\_P. It follows in both directions that

\[
                         M_d=\Lambda_Z\oplus H_d\Lambda_P
                  \quad\hbox{within W0 at that T component}. \tag{PJ15}
\]

The restriction H\_d\textbar P has determinant1 since qz=0. At the conjugate component use the exact Hermitian dual lattice; the orthogonal decomposition and determinant1 persist, and the Gram form on Λ\_P is unimodular there because U is a unit. Hence pr\_P A\_d has determinant valuation−2 at each of the two CM partners. Away the selected support its projection is PJ14. Since A\_d contains the original lattice, the projected module contains PJ14, and its exact finite index is

\[
 [\operatorname{pr}_P A_d:\Lambda_P^{\rm proj}]=T^{4s}.         \tag{PJ16}
\]

This proof includes d=0 and every source/target type. It uses the fixed packet plane, not the h-dependent positive fixed space, and accounts separately for the fixed U denominator and the paired T index.

The original negative column h e\_n has Euclidean norm less than95 at each embedding because (H²+H\^{}(-2))/2\textless95². Its plane projection therefore has \textbar w\_τ\textbar\textless95sqrt U and \textbar b\_τ\textbar\textless95 in PJ14 coordinates. Each coset of O\_E² has at most (1+190sqrt U)\^{}(2f)191\^{}(2f) such pairs by the original-ring packing bound. Consequently

\[
 \#\{\operatorname{pr}_P(h e_n)\}<191^{4f}U^fT^{4s}.            \tag{PJ17}
\]

In particular the U\^{}f factor is retained. Replacing the projection lattice by O\_Eu⊕O\_Ee\_n away T would give an unjustified smaller count and is not done.

Once C and the plane part v of h e\_n are fixed, unitarity reconstructs its Z part as −J\_Z\^{}(-1)C*J\_Pv. Thus the full negative column is known. Every arrow in this fibre fixes K=ker(C\textbar Z) pointwise, and K is positive of dimension166. Fix one h0 in a nonempty fibre. The transformations h0\^{}(-1)h preserve the same source lattice and fix K and e\_n pointwise. Their faithful action is on the positive two-space (K⊕Ee\_n)\^{}perp with its intersection lattice. The same two-prime torsion-injection argument used in R061(GS14), now in dimension2, gives

\[
 G_2\le\gcd(|GL_2(F_{g0})|,|GL_2(F_{g1})|)
            =66857481600<2^{36}.                             \tag{PJ18}
\]

This is a bound for a finite source-lattice isometry group, not an assertion that the moving two-space admits an integral orthogonal basis.

\subsection{5. Complete weighted endpoint and exact exponent}\label{complete-weighted-endpoint-and-exact-exponent}

Retain once each: the Green factor2, the2\^{}668 common-ideal patterns, support choices2\^{}168, pivot factor2\^{}8, fibre factor2\^{}36, and the full ordered pair/weight sum P\_T\^{}s\textless2\^{}(997s) from R059(GB19). The five fixed powers multiply to2\^{}881. There is no added incidence-type multiplicity. Combining PJ6, PJ9, PJ12, PJ13 and PJ17 leaves z exponent

\[
                         (326-2k)f.                           \tag{PJ19}
\]

For 6≤k≤163 this is nonnegative, so use the upper endpoint z\textless M/sqrt k. The entire k-cell is bounded by

\[
\begin{split}
&2^{881-929f/2-334f}\,191^{4f}c^{8f}(L_1L_2L_3)^f
 H^{334f}U^{-331f/2}T^{339s}2^{997s}Q^{-335(f-1)}\\
&\quad\cdot {164\choose k}47^k2^{2fk}(f!)^k
     (87/32)^{fk} f^{-fk}k^{-163f}.                           \tag{PJ20}
\end{split}
\]

At k=164, PJ19 is −2f. Since J is integral, z≥T\^{}(-s/f), so z\^{}(-2f)≤T\^{}(2s). The corresponding complete bound is

\[
\begin{split}
&2^{881-929f/2-336f}\,191^{4f}c^{8f}(L_1L_2L_3)^f
 H^{336f}U^{-333f/2}T^{341s}2^{997s}Q^{-335(f-1)}\\
&\quad\cdot47^{164}2^{328f}(f!)^{164}(87/32)^{164f}
                   (164f)^{-164f}.                          \tag{PJ21}
\end{split}
\]

The changed T power and U exponent in this final cell are essential. No monotonicity is claimed when the z exponent is negative.

Use the retained strictly directed bounds log2 T\textless1019333/1024, log2 Q\textgreater287843/4096, log2 H\textless28900/4096 and log2 U\textgreater601356/4096. Direct integer powers give log2(191)\textless15519/2048, log2(L1L2L3)\textless152789/1024 and log2(c)\textless4097/4096. For each k the certificate computes the least integer strict upper exponent of the complete rational k-factor in PJ20 or PJ21, including its binomial, factorial and all k powers. The common upper exponents before that factor, at s=23, are266217401/4096 for PJ20 and361081009/4096 for PJ21. All159 values6≤k≤164 are checked; the worst is k=6 with coordinate exponent−68672 and complete upper exponent

\[
              -15063111/4096<-3608.                          \tag{PJ21a}
\]

Hence

\[
 \sum_{k=6}^{164}|\text{whole k-cell}|
                    <159\,2^{-3608}<2^{-3600},               \tag{PJ22}
\]

proving PJ3. The exact table and rational upper endpoints are saved in the companion JSON. Failure of the same sufficient envelope for smaller k or support24 is not evidence that any such term exists or has a large contribution.

\subsection{6. Accounting, geometry and remaining gates}\label{accounting-geometry-and-remaining-gates-1}

The chronological predecessor is DO093's75-entry all-top list and unchanged31-entry uniform list. Append3600 exactly once upon checked promotion. The companion checks both full dyadic sums below2\^{}(-230), retains the unique231 charge and every other preceding charge, and leaves the complete support floors23/11 unchanged. The new actual support23 structural restriction is k≤5 in the R071 remainder. It does not claim an exhausted support or a signed lower bound.

The projection PJ14 is a fixed rational map of the original Hermitian modules over Spec(O\_E). Its U denominator, integral local splitting at T and exact determinant index are all displayed. The cross ideal is a determinantal condition on the original paired reductions, and the new joint count is in their unchanged complex half-trace realization. Source/target lattices, CM types, Picard classes, theta line, biextensions, finite projectors and Hilbert metrics are unchanged.

Exact carried contextual locators are SGA3 I4.7.3/4.7.3.1; SGA6 VI Theorem3.3; FGA232 Theorem3.1 pp148--149; FGA236 Corollaries1.7/2.3 pp229/232; SGA7 I VII2.7/2.9 pp169--174 and VIII7.4(b) pp310--311; cote4 scan003 §§8.1.1--8.1.2, scans005--006 §§8.4.1--8.4.3; and cote54 scans005--007. These are the retained frameworks, not an attribution of the new packing argument. No new lambda-ring/K-theory equality, l-adic weight theorem or motivic L-function conclusion is asserted.

DP\_MULT Q3e and THETA3 G1/A7 gain only the smaller finite remainder PJ2--3. The signed sufficient inequality, prescribed Bessel/Mram central value, actual source constructor, selected arithmetic content, Read/STiesR/specific032 pairing and DP44 integer cost are open. CX\_GRH CX-ALG-024 and GCT\_PLETH T-052 keep their separate arithmetic gates. The exact certificate supplements, and does not replace, the complete lattice, packing, trace and finite-fibre proofs above. No Lean build or archive audit is performed.

\clearpage

\section{AM-DO-095 --- the actual three-dimensional residual and its rank-two rigidity}\label{am-do-095-the-actual-three-dimensional-residual-and-its-rank-two-rigidity}

\textbf{Canonical checked integration --- 2026-11-22T22:30:47.328638+00:00.} Root and independent whole-proof/script/interface reviews and separate20463-assertion replays passed. Historical pending wording below is preserved. In the actual R071 signature(2,1) residual, BC is skew-adjoint rank one exactly at global rank3; BC=0 is exactly global rank2 with real reciprocal eigenvalue lambda. Every original S entry of BC is an exact G-unit three-minor, has conjugates\textless2\^{}9 and only simple paired typeIII poles. The pointwise Green envelope is2 Q\^{}(-501\emph{167)2\^{}(9}167f)T\^{}(167t), with no C-pattern/witness factor. In rank2, doubled G divisibility of (lambda-1)\^{}2/lambda forces global unipotence for t\textless=22. The existing DO082 removal therefore leaves rank-two support23 only nonunipotent with all23 active places typeIII; at24 its nonunipotent branch needs t\textgreater=23. No new charge or complete-floor change. Actual central value, arithmetic/source and DP44 gates remain open.

Author: \nolinkurl{/root/definite_oscillator}, 22 November 2026 UTC. Complete proof in the pending area; author and independent whole-file checks are recorded separately. The actual R071 rank-two branch is globally unipotent whenever the type-III count t is at most22. At support23, its surviving rank-two part is therefore nonunipotent with t=23, after retaining DO082's already charged full unipotent removal. At support24, any nonunipotent rank-two term has t at least23. The proof also supplies the exact rank-three minor and a uniform Green bound without a C-ideal pattern. It appends no numerical error charge and does not change the complete support floors. The failed direct R070 counting transplant is kept in the separate \nolinkurl{three_dimensional_restoration_boundary.md}.

\subsection{1. Actual decomposition and determinant}\label{actual-decomposition-and-determinant}

Use exactly the support23/24 doubly-plane remainder of DO093 and the global scalar rigidity R071. Let

\[
 Z=S\oplus u^\perp,\quad \dim_E Z=167,\qquad
 P=Eu\oplus Ee_n,\quad \dim_E P=2.
 \tag{RA1}
\]

Here the perpendicular space is in the original four positive W0 coordinates. Z is positive at every complex realization, P has signature(1,1), and the two are orthogonal. R071 gives, in this rational decomposition,

\[
 h=\begin{pmatrix}I_Z&B\\C&D\end{pmatrix},\qquad
 (h^{-1})_{ZZ}=I_Z,\qquad C\ne0,\quad\operatorname{rank}C=1.
 \tag{RA2}
\]

The original S restriction of C is nonzero. No rational basis of Z is declared globally integral. The basis of DO094(FT8--10) is integral and invertible at each of the actual G primes, which is precisely where the minor arguments below use it.

The determinant of h is already one. Indeed the actual local graph matrix of K062(AT3) is \nolinkurl{I+c*pi^-1*zv} with \nolinkurl{vz=0}, so its determinant is one even at c=0. All source and target pencil determinant ideals are the standard ideal, at both CM partners. The determinant of a supported global g is consequently a norm-one unit in the checked root group mu28. K068(LR3) uses the unique central representative h=(det g)\^{}-1 g, since169 is1 modulo28. The subsequent reductions, including R071, concern this h. There is no additional determinant phase available to eliminate.

Jacobi's complementary-minor identity applied to RA2 gives

\[
 1=\det((h^{-1})_{ZZ})=\frac{\det D}{\det h},\qquad
                         \det D=1.                            \tag{RA3}
\]

In particular D is invertible. Block inversion then implies

\[
                     BD^{-1}C=0.                             \tag{RA4}
\]

This is an identity of the original rational blocks, not a consequence of finite unipotence.

\subsection{2. Exact global classification}\label{exact-global-classification}

Write the Hermitian form as diag(J\_Z,J\_P), with conjugate transpose denoted by a star. Factor C=ell*r, where ell is a nonzero vector of P and r a nonzero row on Z. Put

\[
 v=J_Z^{-1}r^*,\quad q=rv\in F,\quad
 b=-\ell^*J_PD,\quad \kappa=b\ell.
 \tag{RA5}
\]

Unitarity gives ell\emph{J\_P ell=0 and B=v}b. At every real label q is strictly positive, because J\_Z is positive definite. Thus B is nonzero: otherwise the off-diagonal unitary relation C*J\_PD=0 and invertibility of D would force C=0. Define

\[
 W=BC=\kappa vr,\qquad a=\operatorname{tr}D-2.
 \tag{RA6}
\]

The kernel of r in Z is positive and fixed pointwise by h. Its orthogonal complement has dimension three and determinant one. The characteristic polynomial of that actual restriction, followed by the substitution X=1+T, is

\[
\begin{split}
 p_h(X)&=X^3-(3+a)X^2+(3+\bar a)X-1,\\
 p_{h-I}(T)&=T^3-aT^2+(\bar a-2a)T-(a-\bar a).
\end{split}                                                    \tag{RA7}
\]

The coefficient conjugation in the first line is the determinant-one unitary characteristic identity. It is independent of semisimplicity. Direct block expansion, using detD=1, gives

\[
               \operatorname{tr}(BC)=q\kappa=a-\bar a.
 \tag{RA8}
\]

Since q is a nonzero real element, kappa is purely imaginary. Consequently W is skew-adjoint on the positive space Z and has rank one whenever it is nonzero.

The original nonzero S restriction of r ensures that W\_SS is nonzero exactly when W is nonzero. Equations RA7--8 give the exact alternatives

\[
\begin{array}{c|c}
 W\ne0 & \operatorname{rank}_E(h-I)=3,
                  \quad \operatorname{tr}h\notin F,\\
 W=0 & \operatorname{rank}_E(h-I)=2,
                  \quad \operatorname{tr}h\in F.
\end{array}                                                     \tag{RA9}
\]

For the second row, RA7 gives rank at most two on the three-dimensional complement. Rank one is impossible: if rank(h-I)=1, its nonzero image from Z lies in P and is the full image, forcing the upper cross block B to vanish. Rank zero also contradicts C nonzero. This proves the exact rank-two conclusion.

There is a more precise null-line description of the second row. The equality W=0 says b ell=0. Equation RA4 and the two-dimensional identity D\^{}-1=(trD)I-D imply bD ell=0. Therefore D preserves the line E ell. Write D ell=lambda ell. The inverse cross line is also E ell. In a rational Witt basis of P beginning with ell, the determinant-one matrix D is upper triangular with diagonal lambda,lambda\^{}-1. The lower-right unitary relation is D\emph{J\_PD+q b}b=J\_P; its correction has zero first column because b ell=0. Pairing the two columns therefore gives bar(lambda)/lambda=1, hence

\[
              \lambda\in F^\times,\qquad
       p_h(X)=(X-1)(X-\lambda)(X-\lambda^{-1}).                 \tag{RA10}
\]

No integral Witt basis is used here. The case lambda=1 is globally unipotent and has rank(h-I)=2. The case lambda=-1 is retained with possible generalized eigenspaces; no diagonalizability is inferred. For all other lambda the nontrivial reciprocal roots are distinct. Conversely, W nonzero means that the original and inverse cross null lines in P are distinct.

In particular tr(h)=169 would already force global unipotence in this narrower residual, by RA7. This is an algebraic implication; no new trace-vanishing estimate is asserted.

\subsubsection{A determinant-one rank-three boundary example}\label{a-determinant-one-rank-three-boundary-example}

The determinant and compression conditions alone do not imply global unipotence or rank at most two. Over Q(i), in the coordinate order(positive,ell,ell') with Gram matrix diag(1,{[}{[}0,1{]},{[}1,0{]}{]}), take

\[
 h=\begin{pmatrix}
 1&i&1\\
 -1&1-i/2&-1/2\\
 0&i&1
 \end{pmatrix}.                                               \tag{RA11}
\]

It is unitary, has determinant one, both positive diagonal blocks of h and its inverse equal one, and both cross blocks are nonzero in P. Its trace is3-i/2 and det(h-I)=-i, so its rank is three. This is an exact algebraic boundary example. It does not satisfy the prescribed finite G unipotence and is not claimed to be a matrix in the retained short sum. It asserts no height-controlled change from the actual packet plane.

\subsection{3. A complete finite rank criterion and an exact minor}\label{a-complete-finite-rank-criterion-and-an-exact-minor}

At an actual G component, work in the integral invertible basis specified after RA2. The paired inverse-adjoint relation is retained; no positivity is assigned to the finite form. Both B and C have rank at most one by specialization of their global identities. The reduction of h is unipotent by K068. Put E\_0=D-I. Since trD=2 and detD=1,

\[
                  E_0^2=0,\qquad BC=BE_0C.                    \tag{RA12}
\]

The second equality follows from RA4 and D\^{}-1=I-E\_0. Let d=rank(h-I) over the residue field. Then

\[
                      d\le2\iff BC=0,
                  \qquad d=3\iff BC\ne0.                     \tag{RA13}
\]

Here is the complete argument. If d is at most two, nilpotence gives N\^{}3=0, while (N\^{}3)\_ZZ=BE\_0C=BC. Conversely, if one cross block is zero, rankN is at most the rank of the other cross block plus rankE\_0, hence at most two. If both are nonzero, write B=s\emph{b and C=ell}r over the residue field. The equations BC=BE\_0C=0 give b ell=bE\_0 ell=0. Thus E\_0 preserves ker b=span ell. Its nilpotence gives E\_0 ell=0; a two-dimensional square-zero matrix with that kernel is E\_0=e ell b for a scalar e, including e=0. The image of N is then contained in the span of(s,e ell) and(0,ell), and has dimension at most two. In all cases rankN is at most three because each of B,C,E\_0 has rank at most one. This proves RA13 without a moving finite positive space.

At rank three there is the additional exact identity

\[
                              CB=0.                            \tag{RA14}
\]

Indeed b ell is nonzero by RA13, and trN\^{}2=2(b ell)(r s)=0 in the odd residue characteristic. Hence r s=0, which is exactly CB=ell(r s)b=0. This does not assert that either cross block vanishes.

There is a stronger integral identity behind RA13. Before reduction let E\_0=D-I. Since detD=1,

\[
 B[(1+\operatorname{tr}E_0)I-E_0]C=0.
\]

For any Z row i and Z column j, order the two P coordinates consistently. Expansion of the indicated3-by3 minor gives

\[
 \boxed{\ (BC)_{ij}
  =\det\bigl(N[\{i,P\},\{j,P\}]\bigr)\ }.                    \tag{RA15}
\]

Explicitly the determinant is -B\_i adj(E\_0)C\_j, which equals BC\_ij by the preceding identity. Therefore at a G prime where the reduction rank is d, its valuation is at least(3-d)\_+. The same fact holds at the inverse-adjoint partner. The basis is a G-unit basis, so these are valuations of actual original-module minors; no integrality away G is inferred from the basis change.

\subsection{4. Original S witnesses: archimedean size and exact T poles}\label{original-s-witnesses-archimedean-size-and-exact-t-poles}

For a globally rank-three term choose any nonzero original entry delta=(BC)\_ij with i,j in S. Such an entry exists by RA9. At every complex realization, W is skew-adjoint rank one, so its operator norm is \textbar trW\textbar=2\textbar Im trD\textbar. In an orthonormal complex basis of P, D has determinant one, operator norm at most H and singular values R,R\^{}-1 with1\textless=R\textless=H. Its singular-value decomposition writes its trace as Rw+R\^{}-1 bar(w), with \textbar w\textbar\textless=1, because the remaining two-by-two unitary factor has determinant one. Hence

\[
          |\tau\delta|\le2(H-H^{-1})<266<2^9.
 \tag{RA16}
\]

This uses only the complex orthonormal basis for a norm bound. The witness itself remains the stated original S entry.

Its finite denominators have the same simple type-III support as the characteristic coefficients. Indeed

\[
  (BC)_{ij}=\langle h^{-1}e_i,he_j\rangle-\langle e_i,e_j\rangle.
 \tag{RA17}
\]

Both e\_i and e\_j belong to every original pencil. At a type-I place, h and its inverse preserve the full-pencil intersection M, which contains S and is contained in the standard lattice at both CM components; RA17 is integral. At a type-II place the unique incidence is diagonal, and both vectors in RA17 belong to the same paired selfdual labelled lattice. At an inactive place integrality is immediate from the original maximal lattice.

At a type-III place the source is L\_c and target is L\_d with c unequal to d.~The first vector in RA17 belongs to L\_c and the second to L\_d.~Writing P\_c=I+c*pi\^{}-1 Z\_0, with Z\_0\^{}2=0, the paired source/target pairing matrix is exactly

\[
                   J P_c^{-1}P_d
                    =J[I+(d-c)\pi^{-1}Z_0].                   \tag{RA18}
\]

Thus it has a simple pole at worst, at each of the two CM partners. This calculation includes the admitted zero label and introduces no independent partner label. It uses the literal inverse-adjoint source lattice, not two identical single-component graph matrices. Therefore delta is integral away the2t CM primes above type III and has pole at most one there.

For any paired G divisor J\_delta assigned only where delta vanishes, this gives

\[
        N_EJ_\delta\le |N_{E/\mathbb Q}\delta|T^{2t}
                              <2^{9\cdot336}T^{2t}.           \tag{RA19}
\]

All original labels and normalized weights remain unchanged.

For a rank-three term there is also a uniform finite bound which requires no choice of a whole-C ideal. Define the paired divisor

\[
 \mathfrak D_\delta=\prod_{v\ \mathrm{unmixed}}
           (\mathfrak p_v\bar{\mathfrak p}_v)^{(3-d_v)_+}.
\]

RA15 makes this divisor divide every original S entry of BC, so it divides whichever such entry is nonzero. The retained Green exponents are L\_0=0,L\_1=168,L\_2=335,L\_3=501. They satisfy L\_3-L\_d\textless=167(3-d)\_+ for every possible finite rank. The one original Euler correction is at most2 and both mixed factors retain modulus at most one. Therefore RA19 proves

\[
 \boxed{\ |a(h)|\le
 2Q^{-501(f-1)}(N_E\mathfrak D_\delta)^{167/2}
 <2Q^{-501(f-1)}2^{9\cdot167f}T^{167t}\ }.
 \tag{RA20}
\]

This pointwise bound introduces no C-pattern, divisor-pattern or witness-choice factor. It changes no local projector, Haar, central representative or class normalization. It is not a complete absolute count for a support sector; its possible combination with a sparse original-coordinate count is a separate problem.

\subsection{5. Doubled G divisibility and actual rank-two rigidity through t=22}\label{doubled-g-divisibility-and-actual-rank-two-rigidity-through-t22}

Suppose the actual term is in the rank-two alternative of RA9. Retain the exact eigenvalue lambda in F from RA10 and set

\[
 a=\operatorname{tr}(h)-169
       =\lambda+\lambda^{-1}-2
       =\frac{(\lambda-1)^2}{\lambda}\in F.
 \tag{RA21}
\]

At every actual G prime the original h and h\^{}-1 are integral. Because lambda and lambda\^{}-1 are eigenvalues lying in E itself, their monic characteristic equations make both integral over the local DVR, hence units in it. The retained unipotent reduction of h makes the residue of lambda equal to1. Thus a has valuation at least two at each of the334 unmixed F primes. No integral basis of the eigenline is needed for this argument.

The exact characteristic-coefficient support K070(PT5--8), also used in DO074(TC4), applies to a=e\_1(h-I): it is integral away the type-III places and has at most a simple pole at each of those F places. This is a coefficient statement and does not require h itself to have integral entries at type II. If a were nonzero, its original F-field norm would obey

\[
 Q^{2(f-1)}\le |N_{F/\mathbb Q}a|T^t
       <256^f T^t,
 \qquad f=168.
 \tag{RA22}
\]

Indeed at every real embedding, H\^{}-1\textless=\textbar lambda\textbar\textless=H because h and h\^{}-1 both have operator norm at most H=133. Consequently \textbar a\textbar\textless=H+H\^{}-1+2\textless136\textless256. This bound does not assume lambda is positive or that h is semisimple. The exact original prime bounds L\_Q=287843/4096log2T give at t=22 the strict contradiction

\[
 334L_Q-168\cdot8-22L_T
                =\frac{466617}{2048}>0.                     \tag{RA23}
\]

The margin increases as t decreases. Thus a=0 whenever t\textless=22. Equation RA21 then gives lambda=1, and RA7 makes h globally unipotent. Its rank is exactly two by RA9. Its square N\^{}2 is nonzero: when lambda=1, E\_0\^{}2=0 and the P block of N\^{}2 is CB=q ell b, which is nonzero since q,ell,b are nonzero. Hence rankN\^{}2=1 and N\^{}3=0.

To apply the precise DO082(SX1--20) removal at support23, its remaining cross-column condition must also be retained. In the original S block, B\_S=v\_S b with v\_S nonzero. The negative component b\_n cannot vanish: if ell*J\_P D e\_n=0, then D e\_n lies in the null annihilator line E ell; the full vector h e\_n would then have both zero Z part and null P part, contradicting its norm-1. Therefore (Ne\_n)\_S=B\_S e\_n is nonzero. All of SX1 holds, so DO082's already charged complete895 bound applies to every globally unipotent rank-two term here at s=23.

It follows that, after the existing cuts and without a new charge, every remaining rank-two term at support23 is nonunipotent and has t=23. In particular all23 active places are type III in that rank-two subcase. At support24, every nonunipotent rank-two term has t=23 or24; the globally unipotent rank-two part remains under its existing separate status. The rank-three terms are governed by RA20 and their existing cuts. Neither assertion exhausts all support23/24 sectors or changes the complete floors23/11.

\subsection{6. A rational reconstruction identity for the two-plane column}\label{a-rational-reconstruction-identity-for-the-two-plane-column}

The following identity may be used in a later count, but it supplies no lattice count by itself. Use the original plane basis(u,e\_n), so J\_P=diag(U,-1), and normalize the isotropic cross line as ell=(1,beta). Its first coordinate is nonzero because it is isotropic and nonzero. Retain r,v,q from RA5 and put

\[
 K=J_P+qJ_P\ell\ell^*J_P,\qquad
 \det K=-U,\qquad D^*KD=J_P.
 \tag{RA24}
\]

The determinant follows from ell\emph{J\_P ell=0 and the matrix determinant lemma. The isometry relation is exactly the lower-right block of h}Jh=J after substituting B=-v ell\emph{J\_PD and v}J\_Zv=q.

Suppose a\_-=D e\_n is known. Define S\_2={[}{[}0,1{]},{[}-1,0{]}{]}. The first column of D is then uniquely

\[
                  a_+=-S_2\overline{K a_-},
                   \qquad D=[a_+,a_-].                       \tag{RA25}
\]

Indeed a\_-*K a\_-=-1. The vector on the right is K-orthogonal to a\_- and has determinant det{[}a\_+,a\_-{]}=1. In dimension two these conditions determine the first column uniquely, and it is therefore the actual first column. In particular full C together with the two plane entries of a\_- determines h exactly.

The scalar q is also determined by a\_- and the normalized cross line:

\[
 q=\frac{-1-a_-^*J_Pa_-}{|\ell^*J_Pa_-|^2}.                  \tag{RA26}
\]

Its denominator cannot vanish: the Hermitian annihilator of the null line E ell in P is that same line, so a vanishing denominator would make a\_- null, contradicting a\_-*Ka\_-=-1. The original S cross data determine only r restricted to S. The remaining three transverse coefficients in r are not determined merely by the two entries of a\_-; RA26 prescribes their total Hermitian norm. They must still be reconstructed or counted with their actual lattices. Nothing here replaces them by an integral moving basis or treats the rational sphere as having a free integral coordinate frame.

\subsubsection{Exact source and verification scope}\label{exact-source-and-verification-scope-1}

The current definitions and identities used here are canonical R071(GR6--12), DO093(IV19--22), DO094(FT8--10), K068(LR3--4,LR20), K062(AT3--6), R061(GS1--5), R070(SC11--22), K070(PT5--8), DO074(TC1--4), and DO082(SX1--20). Their literal stated interfaces were read. The new classification, finite rank criterion, exact minor, original pairing and doubled-divisibility identities are the arguments above. The companion certificate checks algebraic models, paired finite reductions, exact minor identities, the actual G-unit constants, the full t0..22 rigidity inequalities and the separate failed-restoration diagnostic. It does not enumerate the arithmetic short sum or make a new integral frame.

The mathematical objects remain the same Hermitian modules, source/target lattices and coefficient bundles over Spec(O\_E). The cross products and minors are sections of their original Hom and determinant modules; the field-norm obstruction is a norm of the actual scalar characteristic coefficient over Spec(O\_F). The previously carried contextual anchors are SGA3 I4.7.3/4.7.3.1; SGA6 VI Theorem3.3; FGA232 Theorem3.1 pp148--149; FGA236 Corollaries1.7/2.3 pp229/232; SGA7 I VII2.7/2.9 pp169--174 and VIII7.4(b) pp310--311; cote4 scan003 §§8.1.1--8.1.2, scans005--006 §§8.4.1--8.4.3, and cote54 scans005--007. These are contextual placements, not newly invoked comparison theorems. No new lambda-ring, K-theory, l-adic, CM-motive or motivic L-function assertion follows.

DP\_MULT Q3e and THETA3 G1/A7 gain the precise smaller rank-two remainder and the rank-three bound RA20. The signed sufficient inequality, prescribed Bessel/Mram value, selected arithmetic content, literal source construction, Read/STiesR/specific032 transport and DP44 integer comparison remain open. CX\_GRH CX-ALG-024 and GCT\_PLETH T-052 retain their separate arithmetic gates. The historical error arrays remain75 all-top and31 uniform at authoring; this theorem appends nothing, so any intervening independently checked charges must also remain in the later canonical ledger.

\clearpage

\section{AM-R-073 --- sparse original cross rows remove the remaining rank-three support23 terms}\label{am-r-073-sparse-original-cross-rows-remove-the-remaining-rank-three-support23-terms}

\textbf{Canonical checked integration --- 2026-11-22T22:30:47.374676+00:00.} Root and independent whole-proof/script/interface reviews and separate139-assertion replays passed. Historical pending wording below is preserved. The actual global rank3 R071 terms at s23,k\textless=63 and s24,k\textless=51 have combined absolute contribution\textless2\^{}-21500. The bare support-ideal count retains all original sparse cross coordinates, three transverse values, pivot, fixed packet-plane U denominator, paired T\^{}4 column index, positive rank-two source fibre and complete pair weights. The actual DO095 witness gives the uniform no-pattern rank-three Green bound. All114 cells are checked. Combined with R072 this removes every rank3 support23 term; DO095 then leaves the support23 scalar-cross remainder rank2 nonunipotent, k\textless=5, t23. At support24 remaining rank3 has k\textgreater=52. Complete floors23/11 and signed/global/source/arithmetic/DP44 gates remain open.

Author: /root, GPT-6 Astra Pro, ALP\_MACHINE. 22 November 2026 UTC.

\textbf{Pending whole-proof and dependency checks.} This complete argument uses DO095(RA15--20), whose frozen proof and exact witness are reviewed separately before promotion. No charge is applied by this pending file.

\subsection{1. Exact actual scope}\label{exact-actual-scope-1}

Keep R071's actual doubly-plane/four-zero remainder at s=23 or24, the unchanged J56 all-top H133 endpoint, f=168, n=169, all original pencils, finite projectors, selected trace, class weights and phases. Use the fixed decomposition Z=S⊕u-perp, P=Eu⊕Ee\_n and the original cross row ξ of R072(PJ1--7). Put k=\#supp(ξ), with1≤k≤164. Restrict here to global rank(h−I)=3.

The theorem is the complete absolute bound

\[
 \sum_{\substack{s=23,\ 1\le k\le63\\\mathrm{or}\ s=24,\ 1\le k\le51\\
                  \operatorname{rank}_E(h-I)=3}}
       |\text{actual normalized term}|<2^{-21500}.            \tag{SP1}
\]

Remove only the intersection with the current remainder and retain every earlier charge. Combined with R072's support23 removal for k≥6, this removes every rank-three term of the R071 support23 residual. DO095's rank-two alternative and its already charged unipotent exclusion then leave, at support23, only rank-two nonunipotent terms with k≤5 and all23 active places of type III. This is a conditional description of any remaining term, not an existence assertion or a full support-floor increase. At support24 the new rank-three restriction is k≥52; its rank-two subcases retain their separately proved status.

\subsection{2. Actual rank-three Green bound without a cross-ideal pattern}\label{actual-rank-three-green-bound-without-a-cross-ideal-pattern}

DO095 proves that W=BC on Z is skew-adjoint rank one and that some original S entry δ is nonzero exactly in this rank-three case. Its exact three-minor identity RA15 gives valuation at least(3−d\_v)\_+ at both paired G components. The original source/target pairing RA17--18 makes δ integral away type III, with a pole of at most one at each of its2t CM primes. Every conjugate has modulus less than2\^{}9. Thus, with every original finite factor retained, RA20 is

\[
 |a(h)|<2Q^{-501(f-1)}2^{9\cdot167f}T^{167t},
                       \qquad 0\le t\le s.                   \tag{SP2}
\]

This is pointwise and needs no choice or enumeration of a C-ideal, divisor pattern or witness. It uses finite unipotence at G only. The original normalized real factor satisfies \textbar c(h)\textbar≤1, which is all this count uses.

\subsection{3. Sparse original-coordinate count in the bare support ideal}\label{sparse-original-coordinate-count-in-the-bare-support-ideal}

Keep exactly the reconstruction in R072 §2: the k nonzero original ξ coordinates, three original transverse scalar values λ(y1),λ(y2),λ(z0), and one pivot β determine C over E. All these scalars belong to I\_Σ when the G divisibility is simply omitted. Its norm is T\^{}(-2s), so put z0=T\^{}(-s/f). Their actual radii are M=H/(2sqrt U), M sqrt(L\_j) and B=H/2, with the same exact L\_j of R072(PJ5).

For s=23,24, exact integer comparison gives M/z0\textgreater1. Hence ordinary scalar ideal packing and1+2x≤3x for x≥1 give

\[
 \# C_k\le {164\choose k}\,
 3^{2f(k+4)}M^{2f(k+3)}B^{2f}(L_1L_2L_3)^f
                            T^{2(k+4)s}.                    \tag{SP3}
\]

This deliberately uses the bare ideal, rather than a weighted endpoint for an almost-full row. The actual sparsity k makes its T exponent small. No independent integral frame or rational quotient is counted.

R072(PJ14--18) gives the true projected negative-column count and the complete source fibre, unchanged:

\[
           191^{4f}U^fT^{4s}\,2^{36}.                        \tag{SP4}
\]

The fixed projection denominator at primes dividing U remains in U\^{}f.~The paired T index is exactlyT\^{}(4s), including zero targets, and the fibre acts on the positive rank-two intersection lattice. No new lattice or selected metric is introduced.

\subsection{4. Every multiplicity and the full exponent}\label{every-multiplicity-and-the-full-exponent}

Retain the Green factor2, support choices2\^{}168, pivot bound2\^{}8, source fibre2\^{}36 and the harmless binomial bound binom(164,k)≤2\^{}164. Their total fixed exponent is377. The complete original pair/weight sum is P\_T\^{}s\textless2\^{}(997s), already including every incidence type. Use t≤s in SP2. Combining SP2--4 gives for each full(s,k) cell

\[
\begin{split}
 |{\rm cell}_{s,k}|<&\
 2^{377+9\cdot167f-2f(k+4)}
 3^{2f(k+4)}H^{2f(k+4)}U^{-f(k+2)}
 (L_1L_2L_3)^f191^{4f}\\
 &\quad\cdot T^{(2k+179)s}2^{997s}Q^{-501(f-1)}.
\end{split}                                                   \tag{SP5}
\]

There is no finite-pattern multiplier. The exponent2k+179 is2(k+4)+4+167, accounting separately for the cross row, projected column and the paired witness poles.

Use exactly R072's strictly directed rational bounds L\_T=1019333/1024, L\_Q=287843/4096, L\_H=28900/4096, L\_U=601356/4096, log2(L1L2L3)\textless152789/1024 and log2(191)\textless15519/2048, together with log2(3)\textless8/5. The complete exponent is bounded by

\[
\begin{split}
E(s,k)={}&377+9\cdot167f-2f(k+4)
+2f(k+4)(8/5+L_H)-f(k+2)L_U\\
&+f(152789/1024)+4f(15519/2048)
 +[(2k+179)L_T+997]s-501\cdot167L_Q.
\end{split}                                                   \tag{SP6}
\]

Its k coefficient is positive for s=23 and24. The last claimed cells have exact upper exponents

\[
 E(23,63)=-440519169/20480<-21509,\qquad
 E(24,51)=-103077241/4096<-25165.                              \tag{SP7}
\]

The certificate checks all114 cells individually, not just monotonicity. Their union is less than114·2\^{}(-21509)\textless2\^{}(-21500), proving SP1. The next k values fail this sufficient bound and are not included.

\subsection{5. Chronology, geometry and unresolved questions}\label{chronology-geometry-and-unresolved-questions}

The baseline is the checked R072 list of76 all-top and31 uniform charges. Append21500 once. Both full exact sums remain below2\^{}(-230), and complete floors23/11 are unchanged. The overlap with R072 is intersected with the live remainder; no previous error is removed or counted as a new saving.

The finite witness is an original section of a Hom/determinant module over Spec(O\_E), with both CM valuations. The sparse parameters are original scalar entries in the actual fractional support ideal. The fixed rational projection retains its U denominator and exact local determinant. All CM types, Picard classes, theta lines, biextensions and coefficient metrics are unchanged. Carried contextual locators are SGA3 I4.7.3/4.7.3.1; SGA6 VI Theorem3.3; FGA232 Theorem3.1 pp148--149; FGA236 Corollaries1.7/2.3 pp229/232; SGA7 I VII2.7/2.9 pp169--174 and VIII7.4(b) pp310--311; cote4 scan003 §§8.1.1--8.1.2, scans005--006 §§8.4.1--8.4.3; and cote54 scans005--007. No new lambda-ring/K-theory equality, l-adic weight theorem or motivic L-function conclusion is asserted.

DP\_MULT Q3e and THETA3 G1/A7 gain the smaller actual finite gate stated after SP1. The signed sufficient inequality, prescribed global central value/Bessel period, actual source construction, selected arithmetic image/content, Read/STiesR/specific032 matching and DP44 cost are still open. CX\_GRH CX-ALG-024 and GCT\_PLETH T-052 retain their distinct arithmetic gates. The certificate checks the complete numerical envelope and chronology; the actual witness and projection proofs are separately whole-read dependencies. No Lean build or archive audit is used.

\clearpage

\section{AM-K-096 --- reciprocal poles exclude the rank-two support23 residual}\label{am-k-096-reciprocal-poles-exclude-the-rank-two-support23-residual}

\textbf{Canonical checked integration --- 2026-11-22T22:37:28.843931+00:00.} Root and independent complete proof/script/interface reviews and separate108-assertion replays passed; no correction. The actual R071 support23 residual has no globally nonunipotent global-rank-two term. The nonzero reciprocal trace has exactly23 actual simple F-poles; principal-norm parity and inversion give twelve positive and eleven negative eigenvalue valuations. Both actual arrows then make every original cross scalar integral at the twelve positive places and at most simple-polar at the eleven negative places, with full paired ideal normT\^{}-22 and no U denominator. A nonzero scalar contradicts H\^{}336 T\^{}22\textless(4U)\^{}168, with exact margin374045/512. No sparsity premise or numerical charge is used. DO082 retains the complementary unipotent exclusion. All structural classes and the complete support-floor update require a separate assembly; support24 and signed/global/source/arithmetic/DP44 gates remain open.

\textbf{Algebraic-geometric translation.} This is a norm and valuation argument for the original scalar fractional ideals over Spec(O\_E) and Spec(O\_F). Both source/target module inclusions and both CM partners are retained; a primitive integer extraction avoids a new determinant denominator. The fixed CM Jacobian, Picard classes, polarization, theta line and biextensions are unchanged. Exact carried contextual locators are SGA3 I4.7.3/4.7.3.1; SGA6 VI Theorem3.3; FGA232 Theorem3.1 pp148--149; FGA236 Corollaries1.7/2.3 pp229/232; SGA7 I VII2.7/2.9 pp169--174 and VIII7.4(b) pp310--311; cote4 scan003 §§8.1.1--8.1.2, scans005--006 §§8.4.1--8.4.3 and cote54 scans005--007. These are contextual placements, not newly invoked comparison theorems. No new lambda-ring/K-theory equality, l-adic weight theorem or motivic L-function assertion follows. DP\_MULT Q3e and THETA3 G1/A7 gain the exact empty scalar-cross subcase; CX\_GRH CX-ALG-024 and GCT\_PLETH T-052 keep their distinct arithmetic gates.

Author: \nolinkurl{/root/kummer_extension}, 22 November 2026 UTC. Author whole-proof, source-interface and exact-certificate checks are recorded separately. This is an exact emptiness result, with no numerical error charge.

\subsection{1. Actual scope and theorem}\label{actual-scope-and-theorem}

Keep the unchanged J56 all-top cutoff H=133, f=168, n=169, Q=1427908193729121522601, U=77Q\^{}2−14 and the actual DO040 prime T. Work only in R071's support23 doubly-plane/four-zero remainder. In its fixed rational decomposition

\[
 E^{169}=Z\perp P,\quad Z=S\perp u^\perp,\quad\dim Z=167,
 \quad P=Eu\perp Ee_n,\quad J_P=\operatorname{diag}(U,-1),
\]

write h={[}{[}I,B{]},{[}C,D{]}{]}. R071 and DO095 give C of rank one, C restricted to the original S nonzero, detD=1 and (h\^{}\{-1\})\_\{ZZ\}=I. Every original S basis vector belongs separately to every source and target pencil. The original u has primitive integer coordinates, with squared norm U. Every archimedean operator norm of h and h\^{}\{-1\} is at most H.

\textbf{Theorem.} There is no globally nonunipotent matrix of global rank(h−I)=2 in this actual support23 residual. Consequently the whole rank-two residual at support23 is already exhausted after retaining DO082's existing global-unipotent removal. This theorem adds no charge. It does not by itself prove full support23 exhaustion or remove a support24 term.

\subsection{2. The reciprocal eigenvalue and its actual poles}\label{the-reciprocal-eigenvalue-and-its-actual-poles}

DO095(RA9--10) proves, without a moving integral Witt basis, that the rank-two case has a common rational isotropic plane line E ell and a scalar lambda in F\^{}times such that

\[
 C=\ell r,\qquad D\ell=\lambda\ell,\qquad
 p_h(X)=(X-1)^{167}(X-\lambda)(X-\lambda^{-1}).       \tag{EP1}
\]

Here lambda=1 is precisely the globally unipotent case. In the present nonunipotent case put

\[
 a=\operatorname{tr}(h)-169
   =\lambda+\lambda^{-1}-2=(\lambda-1)^2/\lambda\ne0. \tag{EP2}
\]

At every finite place away the actual T support, h and h\^{}\{-1\} are integral in the original maximal lattice. Since lambda and lambda\^{}\{-1\} lie in F and satisfy their monic characteristic equations, they are units in the local DVR. At every unmixed G place, the actual finite unipotent reduction gives lambda congruent to1. Hence a has valuation at least two at every one of the334 unmixed F primes above g0 or g1. No integral eigenvector or eigenline basis is needed.

The exact coefficient-pole theorem K070(PT5--8), used in DO095(RA21--23), says that a is integral away type III and has at most a simple pole at each active type-III F place. Let r\_p be the number of F places at which a actually has a pole. These are a subset of the23 active places. At every real embedding \textbar a\textbar\textless=H+H\^{}\{-1\}+2\textless256. Thus

\[
 Q^{334}\le |N_{F/\mathbb Q}(a)|T^{r_p}
             <256^{168}T^{r_p}.                            \tag{EP3}
\]

With the exact directed bounds L\_Q=287843/4096log\_2 T,

\[
 334L_Q-168\cdot8-22L_T=466617/2048>0.                     \tag{EP4}
\]

Therefore r\_p cannot be at most22. It equals23: every active place is type III and is an actual simple pole of a. For an integral discrete valuation v, the equality a=(lambda−1)\^{}2/lambda shows that if v(lambda) is nonzero, then v(a)=−\textbar v(lambda)\textbar; if v(lambda)=0, then v(a)\textgreater=0. It follows that

\[
 v(lambda)\in\{+1,-1\}\quad\hbox{at all23 active F places},
 \qquad v(lambda)=0\quad\hbox{elsewhere}.                  \tag{EP5}
\]

All these F places have residue norm T, since the carried rational prime T splits completely in the unchanged field. This is the actual DO040 splitting, not a replacement residue norm.

\subsection{3. Twelve positive and eleven negative valuations}\label{twelve-positive-and-eleven-negative-valuations}

The eigenvalue bound for h and h\^{}\{-1\} gives H\^{}\{-1\}\textless=\textbar tau(lambda)\textbar\textless=H at every real embedding. If n\_+ and n\_- count the signs in EP5, the norm of the principal fractional ideal of lambda gives exactly

\[
 |N_{F/\mathbb Q}(lambda)|=T^{n_+-n_-},\qquad n_++n_-=23,
 \qquad H^{-168}\le T^{n_+-n_-}\le H^{168}.                 \tag{EP6}
\]

The exact integer inequality T\^{}2\textgreater H\^{}\{168\} forces \textbar n\_+−n\_-\textbar\textless2. Its parity is odd, so n\_+−n\_-=+1 or−1.

For an existence contradiction, replace h by h\^{}\{-1\} if necessary and exchange its actual source and target pencils. This preserves all original lattice conditions, the cutoff, the Z identity block and the nonzero original S cross restriction. It replaces lambda by lambda\^{}\{-1\}, so we may assume

\[
                         n_+=12,\qquad n_-=11.              \tag{EP7}
\]

This use of inversion is purely a choice between an existing matrix and its inverse in an emptiness proof. It requires no comparison of individual signed terms, no new pair counting and no factor2 in an error budget. The original kernel's separate exact inversion law is compatible with it but is not needed for the contradiction.

\subsection{4. Integrality of the cross row at every positive place}\label{integrality-of-the-cross-row-at-every-positive-place}

Block inversion, using (h\^{}\{-1\})\_\{ZZ\}=I, gives

\[
       C(h^{-1})=-D^{-1}C=-\lambda^{-1}C.                    \tag{EP8}
\]

Fix an original S basis vector e\_i. At an active place let M\_c and M\_d be the actual source and target modules, with their paired Hermitian duals at the conjugate component. Since e\_i belongs to both modules and hM\_c=M\_d,

\[
 Ce_i=(h-I)e_i\in M_d,\qquad
 -\lambda^{-1}Ce_i=(h^{-1}-I)e_i\in M_c.                  \tag{EP9}
\]

At a positive-valuation place, lambda is a uniformizer times a unit. The literal one-shear pencil satisfies pi M\_c subset Lambda and M\_d subset pi\^{}\{-1\}Lambda, including every admitted zero label. Consequently

\[
                 Ce_i\in\lambda M_c=\pi M_c\subset\Lambda
                  \qquad\hbox{if }v(lambda)=+1.             \tag{EP10}
\]

At a negative-valuation place EP9 retains Ce\_i in M\_d subset pi\^{}\{-1\}Lambda. These arguments apply separately at both CM partners: lambda belongs to F, the local extension is split and unramified, and the opposite module is the exact inverse-adjoint pencil. No single-component graph matrix is substituted for its dual.

Write the original plane cross columns as

\[
                         Ce_i=(u\xi_i,\beta_i).              \tag{EP11}
\]

Primitivity of u gives fixed rational integers a\_j with sum a\_j u\_j=1. Thus xi\_i is an integer linear combination of the original four positive coordinates of Ce\_i. It is integral outside the eleven negative F places, including primes dividing U, and has at most one pole at each of their two CM partners. If I\_- is their paired fractional ideal, then

\[
 \xi_i\in I_-,\qquad N_E I_-=T^{-22}.                      \tag{EP12}
\]

This is a statement in the original coordinates. There is no isotropic normalization over an integer ring, no moving integral frame and no extra denominator from U.

\subsection{5. The archimedean norm contradiction}\label{the-archimedean-norm-contradiction}

The image of C is isotropic. Normalizing that line only over E gives Cx=(u rho(x),eta rho(x)) with eta bar(eta)=U. At a complex label, the norm of the full negative row on the positive space of a unitary matrix is sinh(log\textbar\textbar h\textbar\textbar). Restricting it to the positive Z space and using \textbar\textbar h\textbar\textbar\textless=H yields exactly R072(PJ7):

\[
          |tau(\xi_i)|\le\|rho_\tau\|
            \le\frac{H-H^{-1}}{2\sqrt U}
            <\frac{H}{2\sqrt U}=:M.                        \tag{EP13}
\]

Some original xi\_i is nonzero, because C on S is nonzero and a nonzero isotropic vector in P cannot have zero positive coordinate. Its principal-ideal containment in EP12 and the product of EP13 give

\[
 T^{-22}\le |N_{E/\mathbb Q}(\xi_i)|
                 <M^{336}=\left(\frac{H^2}{4U}\right)^{168}.
                                                               \tag{EP14}
\]

The actual integers satisfy the strict opposite inequality

\[
 H^{336}T^{22}<(4U)^{168}.                                  \tag{EP15}
\]

For a transparent rational margin, use log\_2H\textless28900/4096 and log\_2U\textgreater601356/4096. Then

\[
 168(2+601356/4096-2\cdot28900/4096)
       -22(1019333/1024)=374045/512>730.                     \tag{EP16}
\]

Thus M\^{}\{336\}T\^{}\{22\}\textless2\^{}\{-374045/512\}\textless1, contradicting EP14. This proves the theorem.

\subsection{6. Scope, exact checks and geometric interfaces}\label{scope-exact-checks-and-geometric-interfaces}

The argument applies to every globally nonunipotent rank-two term in the actual R071 support23 residual, without a sparsity hypothesis on xi. DO095 already identifies its globally unipotent rank-two terms with the exact DO082 sector, whose895 charge is retained. Combining this theorem with a separately checked rank-three removal can therefore eliminate the R071 support23 residual. Exhaustion of all other structural classes and any resulting complete support-floor statement remain a separate assembly theorem. No claim is made at support24: the odd sign balance and eleven-negative-place bound used here are specific to23.

The proof adds no absolute-sum charge and does not alter any earlier all-top or uniform error. It preserves the exact coefficient fields, all source/target modules, finite projectors, CM partners, selected metric, central representative, class masses and phases. It supplies no signed global period, arithmetic upper-content witness, source-admissibility constructor or DP44 conclusion.

The direct source/dependency interfaces are DO040's actual split prime and exact host, K070 PT5--8 for coefficient poles, R071 GR6--12 for Z and its rank-one cross, DO095 RA9--10/RA21--23 for the reciprocal eigenvalue, LIFT44 PR4--7 and R051 section1 for both original pencil inclusions, and R072 PJ7 for the original row norm. The new steps are the actual-pole refinement, odd principal-norm sign balance, paired positive-place integrality and exact eleven-pole norm contradiction. These are statements about the unchanged Hermitian modules and their scalar fractional ideals over Spec(O\_E) and Spec(O\_F); they do not invoke a new motivic, lambda-ring or l-adic theorem.

The companion certificate checks the actual integers, directed margins, every possible pole count0..22, all sign counts summing to23, the local valuation formulas, primitive extraction and exact inequalities. Its diagnostics supplement the complete argument above; they do not prove global lattice compatibility by enumeration. All author and independent whole-reading events are recorded separately.

\clearpage

\section{AM-R-074 --- complete all-top support twenty-three exhaustion}\label{am-r-074-complete-all-top-support-twenty-three-exhaustion}

\textbf{Canonical checked integration --- 2026-11-22T22:49:34.874845+00:00.} Root and independent whole-proof/script/interface readings and separate279-assertion replays passed without correction. Historical pending wording below is preserved. Every actual all-top support23 term is exhausted by the complete original F/IC8 partition at all compressed ranks1..33 and every fixed-S/scalar-cross alternative. R072/R073 cover every global-rank3 scalar cross, DO095 verifies all precise DO082 unipotent premises, and K096 removes the entire nonunipotent rank-two residual without sparsity. The complete all-top support floor is24, with uniform floor11. All77/31 chronological errors remain unchanged and both full sums are below2\^{}-230. No new charge is added. The signed inequality J24.2, actual source, selected arithmetic content, global Bessel/Mram value, specific032 matching/value and DP44 remain open.

Author: \nolinkurl{/root}, GPT-6 Astra Pro, inherited effort not externally exposed, ALP\_MACHINE. 22 November 2026 UTC. Controlled proof; whole-proof review and the separate finite accounting replay are prerequisites for canonical promotion.

\subsection{1. Statement on the unchanged actual kernel}\label{statement-on-the-unchanged-actual-kernel}

Keep the J56 all-top K044 endpoint, its already charged H=133 cutoff, f=a=168, n=169, and the original decomposition E\^{}169=S perpendicular W0 with dim(S)=164 and signature(W0)=(4,1). Keep every paired CM pencil, finite projector, class mass, phase, selected Hilbert metric and normalized summand w(h). R067(J23.1--4) gives a complete support floor23. Subsequent containing-set removals and K096's exact emptiness statement do not change this kernel.

\textbf{Theorem.} There is an actual remaining subset R\_top\^{}(24) such that

\[
 \mathcal E_{\mathrm{top}}=1+\operatorname{Re}
       \sum_{h\in R_{\mathrm{top}}^{(24)}}w(h)+\epsilon,
 \qquad |\epsilon|<2^{-230},\qquad s(h)\ge24.
 \tag{J24.1}
\]

This assembly adds no numerical charge. It retains all77 chronological all-top entries and all31 uniform entries present after R073 and K096. The independent uniform H=8192 support floor remains11. The sufficient inequality

\[
 \operatorname{Re}\sum_{h\in R_{\mathrm{top}}^{(24)}}w(h)>-1+2^{-230}
 \tag{J24.2}
\]

remains open. We prove J24.1 by exhausting the original structural partition at s=23; no assertion about a full support24 exhaustion is made.

\subsection{2. Every nonidentity positive compression}\label{every-nonidentity-positive-compression}

Put A=h\_SS and r=rank\_E(A-I). The existing K068 residual has r\textless=33. For A!=I, K069(IC1--8) partitions the actual matrices into the fixed class F, in which h fixes K\_A=ker(A-I) intersect ker(A\^{}*-I) pointwise, and its precise exceptional isotropic-cross complement IC8. This dichotomy has no global-unipotence hypothesis. R068(FA1--9) specializes the original F count to the present H=133 cutoff without changing its class.

\begin{longtable}[]{@{}
  >{\raggedright\arraybackslash}p{(\columnwidth - 2\tabcolsep) * \real{0.5000}}
  >{\raggedright\arraybackslash}p{(\columnwidth - 2\tabcolsep) * \real{0.5000}}@{}}
\toprule\noalign{}
\begin{minipage}[b]{\linewidth}\raggedright
Actual class at support23
\end{minipage} & \begin{minipage}[b]{\linewidth}\raggedright
Complete containing-set coverage
\end{minipage} \\
\midrule\noalign{}
\endhead
\bottomrule\noalign{}
\endlastfoot
F, r=1 & DO086(FR1--10), all supports18..94, with the retained8622 charge. \\
F, 2\textless=r\textless=33 & R068(FA1--9): rank2 through110, rank3 through154, and every rank4..33 through168, with the retained4085 charge. \\
IC8, r=1 or2 & R069(EM1--21), all supports23..28 and all zero/transverse/plane alternatives, with the retained23546 charge. \\
IC8, 3\textless=r\textless=31 & DO064(PM22--23) and the contiguous extensions DO085(MR1--24), whose complete inclusive ceilings are at least23. \\
IC8, r=32 or33 & DO064(PM22--23), complete through168. \\
\end{longtable}

For the fourth row, the lowest enlarged ceiling is exactly23 at r=3; every later rank has a larger ceiling. The old and fresh ranges abut, so no missing intermediate support is inferred away. For the last row the original ceiling is168 at both ranks. Thus every r from1 through33 and both original classes are covered at s=23. All these are bounds for the original full containing classes, with the target, fibre and pair multiplicities retained in their cited statements. No finite-rank or semisimplicity condition is inserted into this partition.

\subsection{3. Identity positive compression and the scalar cross}\label{identity-positive-compression-and-the-scalar-cross}

For A=I, write the original matrix as h={[}{[}I,B{]},{[}C,D{]}{]}. The unitary identities R059(GB1--3) give C\^{}\emph{J0 C=0 and B=-C\^{}}J0 D. Hence C=0 implies B=0 and h=I\_S direct-sum h0. If C!=0, its image is a nonzero isotropic line in W0. This is the complete fixed-block/scalar-cross dichotomy, including globally nonunipotent matrices.

\subsubsection{3.1 Fixed S}\label{fixed-s}

\begin{longtable}[]{@{}
  >{\raggedright\arraybackslash}p{(\columnwidth - 2\tabcolsep) * \real{0.5000}}
  >{\raggedright\arraybackslash}p{(\columnwidth - 2\tabcolsep) * \real{0.5000}}@{}}
\toprule\noalign{}
\begin{minipage}[b]{\linewidth}\raggedright
Actual fixed-S class at support23
\end{minipage} & \begin{minipage}[b]{\linewidth}\raggedright
Complete containing-set coverage
\end{minipage} \\
\midrule\noalign{}
\endhead
\bottomrule\noalign{}
\endlastfoot
Globally nonunipotent, any global rank or type profile & DO090(NE1--6), the entire sector at supports23..26, with the retained65341 charge. \\
Globally unipotent, rank(h-I)=1 & DO091(TD1--26), the entire support23 sector, with the retained1544 charge. \\
Globally unipotent, rank(h-I)=2 & R058(HB1--25), every incidence profile at supports13..23, with the retained39409 charge. \\
rank(h-I)=0 & h=I has support0, so does not occur here. \\
\end{longtable}

The signature-one classification DO060(IB21--22) proves that a nonidentity unipotent has rank1 or2, with N\^{}2=0 in rank1 and N\^{}3=0, rank(N\^{}2)=1 in rank2. It therefore supplies exactly the hypotheses of the two unipotent rows and leaves no other fixed-block rank. DO090 includes all nonunipotent global ranks, characteristic-degree alternatives and incidence profiles.

\subsubsection{3.2 Nonzero scalar cross}\label{nonzero-scalar-cross}

R070(SC1--22) removes every C-transverse case and, when C lies in the packet plane, every case with nonzero original two-by-two short compression at s=23 or24. DO093(IV1--22) then removes the plane-C cases whose inverse cross block is transverse. R071(GR1--12) proves that the remaining transverse scalar is exactly1. Consequently every still remaining scalar-cross matrix at s=23 fixes a positive166-dimensional E-space and, on the original positive167-space Z and packet plane P, has the form

\[
 h=\begin{pmatrix}I_Z&B_Z\\ C_Z&D_P\end{pmatrix},
 \qquad C_Z\ne0,\qquad\operatorname{rank}_E(h-I)\le3.
 \tag{J24.3}
\]

These statements concern actual E-subspaces and original matrix entries. No integral moving basis or new lattice is used in the assembly.

DO095(RA1--9) gives the exact alternatives: the global rank is3 when B\_Z C\_Z is nonzero and2 when it is zero. There is no global-rank1 case with this nonzero cross block. Let k be the number of nonzero original S scalar coordinates xi\_i in R059(GB7); then 1\textless=k\textless=164. The following table exhausts J24.3:

\begin{longtable}[]{@{}
  >{\raggedright\arraybackslash}p{(\columnwidth - 2\tabcolsep) * \real{0.5000}}
  >{\raggedright\arraybackslash}p{(\columnwidth - 2\tabcolsep) * \real{0.5000}}@{}}
\toprule\noalign{}
\begin{minipage}[b]{\linewidth}\raggedright
Remaining scalar-cross class at s=23
\end{minipage} & \begin{minipage}[b]{\linewidth}\raggedright
Complete removal or emptiness
\end{minipage} \\
\midrule\noalign{}
\endhead
\bottomrule\noalign{}
\endlastfoot
Global rank3, 6\textless=k\textless=164 & R072(PJ1--22), with the retained3600 charge. \\
Global rank3, 1\textless=k\textless=5 & R073(SP1--7), which includes all1\textless=k\textless=63 at s=23, with the retained21500 charge. \\
Global rank2 and globally unipotent & DO082(SX1--20), with the retained895 charge. DO095(RA21--26) explicitly verifies every SX1 hypothesis in this residual. \\
Global rank2 and globally nonunipotent & K096(EP1--16), exact emptiness at s=23, with no sparsity hypothesis and no numerical charge. \\
\end{longtable}

The first two rows deliberately use disjoint ranges of k in this assembly, although their proved containing bounds overlap. R073's charge already covers its stated union at both supports23 and24, and is retained only once. For the third row, the nonzero negative cross column (Ne\_n)\_S, rank(N\^{}2)=1 and N\^{}3=0 are part of DO095's explicit interface verification; global unipotence alone is not substituted for those premises. For the fourth, K096 counts actual trace poles, retains both arrows and both CM partners, and applies directly without the k\textless=5 restriction. Its inversion is an existence argument and incurs no absolute-count factor here. All four rows preserve the actual determinant-one central representative.

Thus every nonzero scalar cross at support23 is either already within a charged containing set or absent. Together with Section3.1 this exhausts A=I. Section2 then exhausts all support23 matrices. All smaller supports were already removed by R067.

\subsection{4. Complete chronological errors and the new gate}\label{complete-chronological-errors-and-the-new-gate}

The actual predecessor is checks/r073\_joint\_remainder\_error.json. K096 added no charge. Relative to the68-entry R067 list, the nine later entries, in their exact canonical order, are

\[
 65341,\ 4085,\ 1544,\ 23546,\ 516,\ 31703,\ 885948,\ 3600,\ 21500.
 \tag{J24.4}
\]

The516 term is DO092's support24 rank-one removal. It is retained even though the support23 partition does not need it. Likewise every older charge remains exactly once. There is one all-top exponent231 and every other all-top exponent is at least358; therefore

\[
 \sum_{e\in\mathcal L_{\mathrm{top}}}2^{-e}
 \le2^{-231}+76\,2^{-358}<2^{-230}.
\]

The complete31-term uniform sum remains below2\^{}-230. The companion verifies both full integer dyadic sums, the exact ordered77-entry predecessor, every rank ceiling and the entire k coverage. These finite diagnostics verify the accounting and table endpoints; they do not enumerate the rational kernel or replace the geometric partition proofs.

Each new containing-set removal is applied only to its intersection with the remainder after earlier removals. The actual removed subsets are therefore disjoint in chronological order even when their bounds overlap. Retaining each bound once and using the original energy identity and the triangle inequality proves J24.1. The exact empty K096 set contributes zero. This raises the complete all-top support floor from23 to24 and strictly reduces the finite/real sufficient gate for Q1/Q2 to J24.2.

J24.2 itself, the prescribed global Bessel/Mram value, a fully admissible mixed-CM source, selected arithmetic-image content, actual Read/STiesR and specific032 transport/value, and every changed-family DP44 term remain open. There is no shorter valid lift chain or universal strict negative DP44 difference. Q3 has no Kummer-Hecke RH conclusion; Q4 retains its global arithmetic charge-comparison gate.

\subsection{5. Algebraic-geometric placement and exact locators}\label{algebraic-geometric-placement-and-exact-locators}

This is an exhaustion of determinantal and fixed/isotropic-submodule strata of the same original integral Hermitian correspondences over Spec(O\_E), together with their specified complex realizations. The fractional ideals, both CM partners, coefficient modules, polarization, theta line, CM Jacobian, Picard data and biextensions are unchanged. Rational reconstruction in the cited counts supplies no inverse on the arithmetic theta image.

The exact carried module and determinantal contextual anchors are SGA3, Expose I,4.7.3/4.7.3.1, and SGA6, Expose VI, Theorem3.3. The Picard anchors are FGA, Bourbaki232, Theorem3.1, pp148--149, and Bourbaki236, Corollaries1.7/2.3, pp229/232. Biextensions retain SGA7 I, Expose VII,2.7/2.9, pp169--174, and Expose VIII,7.4(b), pp310--311. Formal groups, Dieudonne crystals and Cartier duality retain cote4 scan003, §§8.1.1--8.1.2, scans005--006, §§8.4.1--8.4.3, and cote54 scans005--007. These are contextual placements, not new comparison theorems.

No new scheme identification, lambda-ring/K-theory equality, l-adic sheaf/weight statement, CM-motive comparison or motivic L-function theorem follows. DP\_MULT Q3e and THETA3 G1/A7 gain exactly the smaller finite/real gate J24.1--2. CX\_GRH CX-ALG-024 and GCT\_PLETH T-052 retain their distinct arithmetic comparisons. Every cited quantitative theorem keeps its own exact source locators and whole-proof checks. This new proof checks complete partition coverage and chronological accounting; no Lean build or archive audit is performed.

\clearpage

\section{AM-DO-096 --- twelve cross-pole places remove the nonunipotent rank-two support24 sector}\label{am-do-096-twelve-cross-pole-places-remove-the-nonunipotent-rank-two-support24-sector}

\textbf{Canonical checked integration --- 2026-11-22T22:49:34.899105+00:00.} Root and independent whole-proof/script/interface readings and separate157-assertion replays passed without correction. Historical pending wording below is preserved. The entire globally nonunipotent rank-two R071 support24 sector has absolute contribution\textless2\^{}-2962487. The nonzero reciprocal trace has23 or24 actual poles; principal-norm parity and one exactly normalized inversion leave at most twelve paired cross-pole places, including the possible active eigenvalue-unit place. Primitive original extraction gives normT\^{}-24 without a U denominator. The full original R059 count retains both five-coordinate columns, source fibre, all G patterns, padded subset choices and pair weights; the complete exponent is -60671740183/20480. Exactly one2962487 charge is added, leaving78/31 entries with both full sums below2\^{}-230. Complete floors remain24/11. Globally unipotent rank2, rank3, signed/global/source/arithmetic and DP44 gates remain open.

Author: \nolinkurl{/root/definite_oscillator}, ALP\_MACHINE. 22 November2026 UTC.

Controlled pending proof. Full author and independent proof/interface/script checks are recorded separately. No canonical file, live error array or support floor is changed by writing this report.

\subsection{1. Actual sector and theorem}\label{actual-sector-and-theorem-2}

Retain the original J56 all-top kernel, its charged cutoff H=133, f=168, n=169, the actual split primes and coefficient fields, every finite projector, selected Hilbert metric, packet point, central representative, source/target pencil, Haar normalization, class mass and phase. Work in the exact R071(GR6--12) residual at support s=24. In its fixed rational orthogonal decomposition write

\[
 E^{169}=Z\perp P,\quad Z=S\perp u^\perp,\quad
 \dim S=164,\quad\dim Z=167,\quad
 P=Eu\perp Ee_n,\quad J_P=\operatorname{diag}(U,-1),
 \qquad h=\begin{pmatrix}I_Z&B_Z\\C_Z&D_P\end{pmatrix}.
 \tag{RP1}
\]

Here u is the original primitive integer vector, U=u*u=77Q²−14 and Q=1427908193729121522601. The restriction C\_Z\textbar S is nonzero, detD\_P=1, and (h\^{}-1)\_ZZ=I. Every334 unmixed finite G reduction is unipotent. We impose the separate global conditions

\[
                  \operatorname{rank}_E(h-I)=2,
                  \qquad h\text{ is not unipotent}.          \tag{RP2}
\]

For the original decomposition S perpendicular W0, with signature(W0)=(4,1), denote its lower cross block by C. Thus C=C\_Z\textbar S in original coordinates and C!=0. Its image lies in the original packet plane.

\textbf{Theorem.} The sum of the absolute actual normalized contributions of all matrices RP1--2 at support24, intersected with the current remainder if desired, is less than

\[
                              2^{-2962487}.                 \tag{RP3}
\]

One fresh all-top charge2962487 suffices. No uniform charge is added. This theorem removes only the globally nonunipotent rank-two branch. A globally unipotent matrix has the zero trace witness and does not satisfy the pole argument below. No rank-three term, other structural class or complete support-floor statement is included.

The proof first reduces the possible paired poles of the original cross row to twelve by the reciprocal eigenvalue. It then applies the entire original Case-B count of R059(GB7--21), retaining its two full five-coordinate target counts, full source fibre and the simple bound \textbar c(h)\textbar\textless=1. No sharper selected-trace theorem or new projected-column estimate is required.

\subsection{2. Actual trace poles and the two sign profiles}\label{actual-trace-poles-and-the-two-sign-profiles}

DO095(RA9--10) gives an eigenvalue lambda in F\^{}times with

\[
 C_Z=\ell r,\quad D_P\ell=\lambda\ell,\qquad
 p_h(X)=(X-1)^{167}(X-\lambda)(X-\lambda^{-1}),
 \quad a=\operatorname{tr}h-169=(\lambda-1)^2/\lambda\ne0.
 \tag{RP4}
\]

The nonzero a is the exact point at which RP2 is used. The argument does not assume that a nonzero characteristic coefficient exists in a globally unipotent sector.

The valuation and norm argument of K096(EP2--6) is support-independent up to its final parity step. To state its precise application here: lambda and lambda\^{}-1 are units away the active T places, by their monic characteristic equations in the original integral local models. At all334 unmixed F primes above g0,g1, finite unipotence gives lambda congruent to1 and hence v(a)\textgreater=2. K070(PT5--8) gives integrality of a outside type III and at most one pole at each type-III F place. If r\_p is the number of actual poles of a, not merely the number of type-III places, then

\[
 Q^{334}\le |N_Fa|T^{r_p}<256^{168}T^{r_p}.
 \tag{RP5}
\]

The actual integer comparison Q\^{}334\textgreater256\^{}168 T\^{}22, or its directed margin466617/2048 from K096(EP4), proves r\_p\textgreater=23. Since s=24,

\[
                             r_p\in\{23,24\}.               \tag{RP6}
\]

For any nonunit lambda, v(a)=−\textbar v(lambda)\textbar; for a unit lambda, v(a)\textgreater=0. Thus the actual pole places have v(lambda)=+1 or−1, and all other finite places have v(lambda)=0. Every active F prime has norm T, because the actual T splits completely. Let n\_+,n\_- count these signs. The eigenvalue bounds for h and h\^{}-1 give

\[
 |N_F\lambda|=T^{n_+-n_-},\quad n_++n_-=r_p,\quad
 H^{-168}\le |N_F\lambda|\le H^{168},\qquad T^2>H^{168}.
 \tag{RP7}
\]

Consequently \textbar n\_+−n\_-\textbar\textless2. For r\_p=23, the possibilities are (12,11) or (11,12). For r\_p=24, parity forces the single possibility (12,12).

DO093(IV4--15) proves the exact inversion identity for the unchanged normalized summand, including the mixed projectors, selected compression, symmetric pencil weights, class mass, determinant-one central representative and cutoff:

\[
 \mathcal W(h^{-1})=\overline{\mathcal W(h)},\qquad
                         |\mathcal W(h^{-1})|=|\mathcal W(h)|.
 \tag{RP8}
\]

Inversion reverses the actual ordered source/target arrows at both CM partners together, preserves s=24 and the finite-unipotence premise, and replaces lambda by lambda\^{}-1. Let Y be the containing actual sector in which n\_+−n\_- is nonnegative. It contains the (12,11) and (12,12) cases. Every RP1--2 matrix is in Y or has its inverse in Y. A previously cut remainder need not be inverse-closed: enlarge its image to the entire Y before counting. It follows that

\[
 \sum_{h\text{ in RP1--2}}|\mathcal W(h)|
                     \le2\sum_{h\in Y}|\mathcal W(h)|.        \tag{RP9}
\]

This is one global inversion factor2, not a separate choice at each label. There is no separate union factor for the two values of r\_p: both are included in the same containing sector below.

\subsection{3. The entire original cross row has at most twelve paired poles}\label{the-entire-original-cross-row-has-at-most-twelve-paired-poles}

For an original basis vector e\_i of S, the two arrows and RP4 give

\[
 Ce_i=(h-I)e_i\in M_d,\qquad
 (h^{-1}-I)e_i=-\lambda^{-1}Ce_i\in M_c.                    \tag{RP10}
\]

Here M\_c,M\_d are the actual source and target pencils at the chosen component; the conjugate component is always their inverse-adjoint dual. Every e\_i belongs to both. The exact one-shear formula LIFT44(PR4--7) gives

\[
       \varpi M_c\subset\Lambda,\qquad
       M_d\subset\varpi^{-1}\Lambda                          \tag{RP11}
\]

at each CM partner separately, including every admitted zero label. The symbol varpi denotes the local uniformizer at that partner.

If v(lambda)=+1, the second membership in RP10 gives Ce\_i in lambda M\_c=varpi M\_c, hence Ce\_i is integral in the original coordinates. If v(lambda)=−1, the first membership retains only a simple pole. At a lambda-unit active place we keep that same conservative simple-pole bound; no additional type-I or type-III integrality refinement is needed.

For r\_p=23 in Y, there are eleven negative places and at most one active lambda-unit place. For r\_p=24, there are twelve negative places and no such unit place. Therefore all original entries of C are integral outside some subset D of at most twelve of the24 active places, and have at most one pole at each CM partner over D. The statement uses lambda in F: the two partner valuations agree, so the subset is paired, not two independent subsets.

Pad D to exactly twelve active places, by any fixed ordering when necessary. This enlarges the allowed fractional ideal. Write

\[
 I_D=\prod_{v\in D}(\mathfrak p_v\bar{\mathfrak p}_v)^{-1},
                    \qquad N_EI_D=T^{-24}.                  \tag{RP12}
\]

There are at most binom(24,12)\textless2\^{}24 such padded choices for any fixed support. No lambda value or separate pole-sign profile is subsequently counted.

In the original plane coordinates C e\_i=(u xi\_i,beta\_i). Primitivity of u extracts each xi\_i as a fixed integer linear combination of its original four positive entries. Thus the denominator statement has no U factor, including at primes dividing U. At G use the original whole-C orientation of R059(GB9--10): one chosen partner at rank one, both at rank zero, none at rank at least two. For its integral ideal J,

\[
 \xi\in (J I_D)^{164},\qquad
 |a(h)|\le2Q^{-335(f-1)}(N_EJ)^{335/2},
                   \qquad\#\{J\text{ patterns}\}\le2^{668}.
 \tag{RP13}
\]

Here a(h) is the finite coefficient, distinct from the scalar trace witness a in RP4. The actual pattern is common to the whole cross block. Inversion can exchange its partner choices, but the same complete2\^{}668 enumeration already contains all choices. No second ideal-pattern count appears. The two mixed finite factors retain their original modulus-at-most-one bounds.

\subsection{4. Complete original Case-B count with the smaller pole ideal}\label{complete-original-case-b-count-with-the-smaller-pole-ideal}

Put L=J I\_D and z=(N\_EL)\^{}(1/(2f)). Use the intentionally weaker original R059 bound

\[
                    |\tau(\xi_i)|<2^{-65}=:B_0.             \tag{RP14}
\]

It follows from \textbar\textbar C\textbar\textbar\textless H and H/sqrt(U)\textless2\^{}-65, so no new real-model estimate enters. At least one xi\_i is nonzero. The original half-trace norm and field-norm minimum give z\textless B\_0. R059(GB14--16)'s packing and monotone endpoint then yield

\[
 \#\{\xi\}\le(1+2B_0/z)^{328f},\qquad
 z^{335f}(1+2B_0/z)^{328f}
       =z^{7f}(z+2B_0)^{328f}<B_0^{335f}3^{328f}.            \tag{RP15}
\]

The exponent7f is positive. This counts the complete164-coordinate vector, admitting zero entries; no sparsity union is needed. Since N\_EJ=T\^{}24 z\^{}(2f), multiplication by RP13 supplies the complete vector factor

\[
                      T^{335\cdot12}\cdot2^{-65\cdot335f}
                                                 3^{328f}. \tag{RP16}
\]

For a nonzero coordinate xi\_j choose its original pivot column z\_p=C e\_j. Rank(C)=1 reconstructs the entire C from xi and z\_p over E. Then v=(h e\_n)\_W0, the W0 part of the original negative column, reconstructs its S part by −C*J\_0v, exactly R059(GB6--8). The reconstruction quotient is not called integral. Each of the two original five-coordinate columns is counted in the true target overlattice varpi\^{}-1 M\_d,W0, containing O\_E\^{}5, uniformly in every label. We retain the full R059(GB18) counts

\[
              \#\{z_p\},\ \#\{v\}<2^{14861}T^{10s}.         \tag{RP17}
\]

Although the new pole cap can constrain the pivot too, no saving is taken there. For each fixed original source/target pair, the reconstructed complete negative column has source-lattice fibre at most G\_n\textless2\^{}5022. We retain this full original fibre, rather than substituting a smaller-dimensional one.

The actual normalized real coefficient has modulus at most1. The complete ordered-pair and weight sum is unchanged:

\[
           P_T^s=\left(\frac{2(T-1)^2}{T-2}\right)^s
                              <2^{997s}.                   \tag{RP18}
\]

It already includes all three incidence types and both diagonal and off-diagonal labels. The companion CM label is determined by the actual pairing. No extra3\^{}s, type-III multiplicity, affine-slice covolume or zero-target correction is introduced. This is the original full-column count and exact absolute pair envelope.

Retain once each: the Green Euler factor2; G patterns2\^{}668; support choices2\^{}168; pivot choices164\textless2\^{}8; the harmless original Case-B overcount2\^{}2; the full source fibre2\^{}5022; the two column constants2\^{}(2*14861). These form the unchanged R059 constant

\[
          C_0=1+668+168+8+2+5022+2(14861)=35591.             \tag{RP19}
\]

The new padded denominator choices contribute2\^{}24, and RP9 contributes2. All r\_p, sign and denominator cases have already been included; there is no further union factor.

\subsection{5. Exact complete exponent and chronological scope}\label{exact-complete-exponent-and-chronological-scope}

The actual prime inequalities are

\[
 \log_2T<L_T=1019333/1024,\qquad
 \log_2Q>L_Q=287843/4096,\qquad 3^5<2^8.                    \tag{RP20}
\]

Combine RP16--20 with s=24 and \textbar D\textbar=12. The whole absolute contribution is strictly below2\^{}E, where

\[
\begin{split}
 E={}&C_0+25-65\cdot335f+(8/5)328f\\
    &+(335\cdot12+20\cdot24)L_T+997\cdot24
                        -335\cdot167L_Q\\
   ={}&-60671740183/20480<-2962487.                         \tag{RP21}
\end{split}
\]

This proves RP3. There is a single complete-sector bound, with no subsequent sum over k, r\_p or denominator sizes. Exact integer and rational comparisons in the companion verify all sign profiles, the local pole bounds, actual prime inequalities, the complete exponent and the full ordered error arrays.

At writing the latest charged predecessor is R073, whose actual lists contain77 all-top and31 uniform entries, ending3600,21500. K096 and the support-floor assembly add no charge. Append only2962487, retaining every previous entry, to obtain78/31. Any intervening independent charge must also be retained at canonical integration; a later no-charge assembly can change a floor without changing these lists. This report does not itself increase either full support floor.

The unchanged signed sufficient inequality, prescribed global Bessel/Mram period, selected-image arithmetic content, literal admissible source constructor, concrete cycle/value and DP44 sign remain separate. The algebraic-geometric setting is the original Hermitian modules and scalar fractional ideals over Spec(O\_E), with their exact CM pairing and unchanged half-trace metric. Carried contextual anchors remain SGA3 I4.7.3/4.7.3.1; SGA6 VI Theorem3.3; FGA232 Theorem3.1 pp148--149; FGA236 Corollaries1.7/2.3 pp229/232; SGA7 I VII2.7/2.9 pp169--174 and VIII7.4(b) pp310--311; cote4 scan003 §§8.1.1--8.1.2, scans005--006 §§8.4.1--8.4.3 and cote54 scans005--007. No new lambda-ring, K-theory, l-adic or motivic conclusion is imported.

Direct inputs at their established scope are R071(GR6--12), DO095(RA9--10,RA21--23), K096(EP2--6,EP8--13), DO093(IV4--15), LIFT44(PR4--7), K070(PT5--8), and the complete original R059(GB7--21). The new result is the even/odd support24 pole classification, its twelve-place cross ideal, and the exact full-count substitution. The companion diagnostics supplement that general proof and do not claim to enumerate actual rational unitary arrows.

\clearpage

\section{AM-R-075 --- the symmetric compression raises the exceptional finite-rank baseline}\label{am-r-075-the-symmetric-compression-raises-the-exceptional-finite-rank-baseline}

\textbf{Canonical checked integration --- 2026-11-22T22:56:09.020495+00:00.} Root and independent whole-proof/script/interface readings and separate609-assertion final replays passed. The sole author notation clarification specifies path vertices; no mathematical correction. Historical pending wording below is preserved. In the original IC8 class at compressed ranks3..31, the full symmetric compression has rank r at every actual G factor. Finite unipotence therefore forces full rank\textgreater=r+1; at equality the unique nontrivial Jordan block fixes every untouched original column at both CM partners. The resulting paired ideal and baseline rank r+2 retain every original DO085 coordinate, volume, target, source fibre and pair factor. All145 stated fresh cells contribute\textless2\^{}-8866, with exact maximum -36351867/4096. The complete rank3 ceiling becomes28, and with the earlier F/IC8 cuts every remaining term at supports24..28 has A=I. One8866 charge gives79/31 entries with both full sums\textless2\^{}-230. Complete floors24/11 remain; signed/global/source/selected-image/032/DP44 gates are open.

Author: \nolinkurl{/root}, GPT-6 Astra Pro, inherited effort not externally exposed, ALP\_MACHINE. 22 November 2026 UTC. Controlled proof; whole proof, literal interface and exact arithmetic checks are recorded separately before promotion.

\subsection{1. Original sector and complete new ranges}\label{original-sector-and-complete-new-ranges}

Keep the unchanged J56 all-top H=133 kernel, f=a=168, n=169, the original positive S of dimension164 and W0 of signature(4,1), every paired source/target pencil, finite projector, class mass, phase and selected Hilbert metric. Put A=h\_SS and r=rank\_E(A-I). This theorem concerns the precise K069(IC8) exceptional class with3\textless=r\textless=31. Every334 unmixed G reduction is unipotent, but global unipotence is not assumed.

Let R\_r be the complete DO085(MR1--24) ceiling. The new complete ceilings R'\_r are:

\begin{longtable}[]{@{}rrr@{}}
\toprule\noalign{}
r & R\_r & R'\_r \\
\midrule\noalign{}
\endhead
\bottomrule\noalign{}
\endlastfoot
3 & 23 & 28 \\
4 & 28 & 34 \\
5 & 34 & 40 \\
6 & 40 & 45 \\
7 & 45 & 51 \\
8 & 51 & 56 \\
9 & 56 & 62 \\
10 & 62 & 67 \\
11 & 67 & 73 \\
12 & 73 & 78 \\
13 & 78 & 84 \\
14 & 84 & 89 \\
15 & 89 & 94 \\
16 & 94 & 99 \\
17 & 99 & 105 \\
18 & 105 & 110 \\
19 & 110 & 115 \\
20 & 115 & 120 \\
21 & 120 & 125 \\
22 & 125 & 130 \\
23 & 130 & 135 \\
24 & 135 & 140 \\
25 & 140 & 145 \\
26 & 145 & 150 \\
27 & 150 & 155 \\
28 & 155 & 160 \\
29 & 160 & 164 \\
30 & 164 & 168 \\
31 & 168 & 168 \\
\end{longtable}

\textbf{Theorem.} The full absolute contribution of the145 fresh cells

\[
       3\le r\le31,\qquad R_r<s\le R'_r
\]

is strictly less than

\[
                              2^{-8866}.                    \tag{NR1}
\]

This is one containing-set removal on its intersection with the current remainder. It retains every old charge. In particular the entire IC8 rank3 class is covered through28. Combined with R069's complete rank1/2 range through28, the other IC8 ranges, DO086 and R068, no nonidentity positive compression remains at supports24..28. The remaining matrices in that support interval therefore satisfy A=I. The complete all-top floor remains24 until all identity-compression classes are exhausted; the uniform floor remains11.

\subsection{2. Full symmetric-compression rank at every actual G component}\label{full-symmetric-compression-rank-at-every-actual-g-component}

The partial-monomial classification DO064(PM5--9) says that A has entries in mu28 or0 and is the identity outside at most2r original vertices. Both A and A\^{}dagger fix the complementary original basis vectors. Put N=h-I and

\[
             H_h=h+h^{-1}-2I=N^2h^{-1},\qquad
             (H_h)_{SS}=A+A^\dagger-2I.                    \tag{NR2}
\]

R069(EM3--5) proves the symmetric-compression rank identity in its stated rank1/2 application. We need its following form for all ranks under consideration:

\[
       \operatorname{rank}(\overline A+
                    \overline A^\dagger-2I)=r
               \quad\hbox{at every actual G component}.    \tag{NR3}
\]

Here is the complete extension; the argument has no restriction to rank1/2. Decompose the original weighted partial-permutation graph into paths and cycles. Along a path with l vertices, diagonal conjugation makes A the forward shift and A\^{}dagger the reverse shift, because every edge weight has conjugate equal to its inverse. The symmetric tridiagonal determinant is (-1)\^{}l(l+1). Both actual G characteristics exceed165 and l\textless=164, so it is a unit. The symmetric rank is l, equal to rank(A-I) on that path. On a cycle of length l with product weight zeta, A is invertible with A\^{}l=zeta I, and A\^{}dagger=A\^{}-1. The polynomial X\^{}l-zeta is separable in the residue field. Consequently (A-I)\^{}2A\^{}-1 has the same rank as A-I. That rank is l-1 if zeta=1 and l otherwise; injective reduction of mu28 preserves that distinction. Summing the component ranks proves NR3 at each CM partner. These conjugations are used only for a rank identity over the residue field, never for an integral lattice count.

Write d=rank(Nbar). Nilpotence and NR2--3 imply

\[
             r\le\operatorname{rank}(\overline N^2)
                       \le d-1,\qquad d\ge r+1.            \tag{NR4}
\]

The second inequality holds for every nonzero nilpotent matrix, by its Jordan blocks. Thus the full finite-rank-r stratum allowed in the earlier DO085 estimate is actually absent. The new minimum r+1 has its own stronger paired congruence, proved next.

\subsection{3. At the new minimum every untouched original column is fixed}\label{at-the-new-minimum-every-untouched-original-column-is-fixed}

Suppose d=r+1 at one actual G component. NR4 gives rank(Nbar\^{}2)=r. For a nilpotent matrix the difference rank(Nbar)-rank(Nbar\^{}2) is the number of nontrivial Jordan blocks. Hence there is exactly one such block, of size r+2; all remaining blocks have size1. In particular

\[
 \operatorname{im}\overline N\cap\ker\overline N
     =\operatorname{im}\overline N^{r+1}
     \subset\operatorname{im}\overline N^2
     =\operatorname{im}\overline H_h.                       \tag{NR5}
\]

The final image equality uses the invertibility of h and its commutation with N. Both Hbar and its compressed S block have rank r, so the coordinate projection P\_S is injective on im(Hbar).

Let e\_j be any untouched original S basis vector. Since A and A\^{}dagger both fix it, (Hbar)\_SS e\_j=0. Injectivity yields Hbar e\_j=0 and hence Nbar\^{}2 e\_j=0. Therefore Nbar e\_j lies in the left side of NR5, and so in im(Hbar). Its S projection is (Abar-I)e\_j=0. Injectivity again gives

\[
                       \overline N e_j=0.                   \tag{NR6}
\]

The conjugate CM component has the same full finite rank, because it is the inverse adjoint, and satisfies the same NR3. Repeat the argument there. A fixed column is fixed by the inverse; applying the paired adjoint identity to the opposite fixed column then also fixes the original row. Thus, at both CM partners,

\[
                    (h e_n)_j\equiv0
                      \quad(j\text{ untouched in }S).       \tag{NR7}
\]

This proof uses only nilpotence, rank and the original coordinate projection over the residue field. It assigns no positive signature to a finite field and constructs no moving integral frame.

Choose a padded original2r-subset containing every changed vertex, as in DO085(MR4--6), and write k=164-2r. Let J contain both CM primes at exactly those unmixed relative G places with full finite rank r+1. Then

\[
 N_EJ=P_{r+1}^2,\quad1\le N_EJ\le Q^{2(f-1)},
 \quad\#\{J\}\le2^{334},\qquad
 Q=1427908193729121522601.                                  \tag{NR8}
\]

At all other unmixed factors the rank is at least r+2. Put L\_j=169j-j(j+1)/2 and d\_r=L\_(r+2)-L\_(r+1)=167-r. The original unipotent Green formula, with one global Euler correction, gives

\[
 |a(h)|\le2Q^{-L_{r+2}(f-1)}(N_EJ)^{d_r/2}.                 \tag{NR9}
\]

Here r+2\textless=33 and the smallest retained denominator exponent is136; the already proved uniform Euler bound remains valid. Both mixed projectors and the real normalized coefficient retain modulus at most1. No finite-rank orientation or witness count is added.

\subsection{4. The original column count with the new paired ideal}\label{the-original-column-count-with-the-new-paired-ideal}

Every coordinate, target and source-fibre statement of DO085(MR12--17) applies to NR7--8 with this new definition of J. Namely, for I\_Sigma of paired norm T\^{}(-2s),

\[
 (he_n)_{S_{\rm free}}\in(JI_\Sigma)^k,\quad
 (he_n)_{S\setminus S_{\rm free}}\in I_\Sigma^{2r},\quad
 (he_n)_{W0}\in\mathcal B_\Sigma\supset O_E^5,\quad
 [\mathcal B_\Sigma:O_E^5]=T^{10s}.                         \tag{NR10}
\]

These are the original paired modules, uniform in every admitted label including zero. Put z=(N\_E(JI\_Sigma))\^{}(1/(2f)), z0=T\^{}(-s/f), z1=Q\^{}((f-1)/f)T\^{}(-s/f). Then z0\textless=z\textless=z1. Retain

\[
 D_0=7^{140}29^{144}5^{84}13^{84},\qquad
 V_d(C)=\frac{(528C^2)^{fd}}{(fd)!D_0^d}.                    \tag{NR11}
\]

The full free-coordinate count is V\_k(66.5+4200z)z\^{}(-2kf); the other positive coordinates cost V\_(2r)(67)T\^{}(4rs); and the full W0 target count is V\_5(223)T\^{}(10s), with V\_5(223)\textless2\^{}5757. The source-lattice fibre remains G\_n\textless2\^{}5022. These are precisely DO085(MR14--17)'s arbitrary-ideal half-trace transference and full column estimates. All fresh supports here are at least24, so the previously verified z0\textless2\^{}-71 and z1\textless1/2 are still valid. The unchanged original A/coordinate-set count is A\_r=binom(164,2r)29\^{}(4r\^{}2).

Multiplication by NR9 leaves the z-dependent factor, apart from T\^{}(d\_r s),

\[
            V_k(66.5+4200z)z^{(d_r-2k)f}.                   \tag{NR12}
\]

Its logarithm is convex in log z, so its maximum over the full actual-ideal range is at z0 or z1. At z0 the radius is below67 and the full column T power is2k+4r+10=338. At z1 the radius is below2167 and the combined Green/column Q exponent is

\[
            -[L_{r+2}+2k-d_r](f-1)
                      =-[L_{r+1}+2k](f-1).                 \tag{NR13}
\]

No lower bound on z1 at large supports is needed. There is no endpoint-union factor; the maximum bounds each actual ideal before the complete pattern count.

Let b\_r and v\_d(C) be the least strict integer upper binary exponents of the exact A\_r and V\_d(C), as already certified in DO085. Retain its complete constant

\[
 C_r=b_r+334+168+5022+1+v_{2r}(67)+5757,
 \quad L_T=1019333/1024,\quad L_Q=287843/4096.                \tag{NR14}
\]

The full endpoint exponents are now

\[
\begin{split}
 E_{r,0}(s)&=C_r+v_k(67)+(338L_T+997)s
                         -L_{r+2}167L_Q,\\
 E_{r,1}(s)&=C_r+v_k(2167)+(338L_T+997)s
                         -(L_{r+1}+2k)167L_Q.
\end{split}                                                \tag{NR15}
\]

All2\^{}334 paired ideal patterns,2\^{}168 support choices, the complete source fibre and single Green factor are in C\_r. The original complete ordered-pair/absolute-incidence factor P\_T\^{}s is bounded by2\^{}(997s), including every type and both CM partners without an extra label count. The fixed positive-block count and every original target column remain present. Only the new finite-rank baseline and its associated paired ideal change.

\subsection{5. All fresh cells, exact error and implications}\label{all-fresh-cells-exact-error-and-implications}

Take R'\_r to be the greatest integer at most168 for which both NR15 exponents are at most-4096. Exact rational evaluation gives the table in Section1. Each new ceiling is at least its predecessor; the sum of all fresh interval lengths is145. The largest exponent among both endpoints of all145 fresh cells is

\[
                  E_{13,0}(84)=-36351867/4096<-8874.          \tag{NR16}
\]

Consequently the complete union is bounded by145 times2\^{}-8874, which is strictly less than2\^{}-8866. This proves NR1. The old and fresh intervals abut; rank31 has no fresh cell, and no support beyond168 is considered. A failed next sufficient exponent, where it occurs, supplies no lower bound for an actual contribution.

At authoring the actual canonical predecessor is DO096 with78 all-top and31 uniform errors and complete floors24/11. Append8866 once, retaining every earlier entry. The complete79/31 dyadic sums remain below2\^{}-230; the unique231 and all other all-top entries at least358 are unchanged. The companion verifies both full sums and every cell, in addition to the literal old/new rank table. No support-floor change follows solely from this sector removal.

Since the smallest new IC8 ceiling is28, R069 covers its ranks1/2 through28, and DO064 covers ranks32/33 through168, all IC8 classes are exhausted at supports24..28. DO086 covers F rank1 through94 and R068 covers all F ranks2..33 beyond28. Together with the complete preceding floor24, the actual remainder in24..28 therefore has A=I. This is a strict reduction of the Q1/Q2 finite/real sufficient gate. The remaining fixed-block/scalar-cross classes, signed positivity, actual admissible mixed-CM source, selected arithmetic content, prescribed global Bessel/Mram value, Read/STiesR/specific032 matching and every changed-family DP44 cost remain open.

The new algebraic input is the determinantal identity NR3 and the single-Jordan-block argument NR5--7 in reductions of the original Hermitian modules over Spec(O\_E). The original scalar ideals, Picard data, CM polarization, theta line, biextensions and coefficient metrics are unchanged. Exact contextual anchors are SGA3 I4.7.3/4.7.3.1; SGA6 VI Theorem3.3; FGA232 Theorem3.1 pp148--149; FGA236 Corollaries1.7/2.3 pp229/232; SGA7 I VII2.7/2.9 pp169--174 and VIII7.4(b) pp310--311; cote4 scan003 §§8.1.1--8.1.2, scans005--006 §§8.4.1--8.4.3 and cote54 scans005--007. No new lambda-ring/K-theory equality, l-adic sheaf/weight theorem, CM-motive comparison or motivic L-function conclusion is imported. DP\_MULT Q3e and THETA3 G1/A7 gain precisely the stated smaller remainder; CX\_GRH CX-ALG-024 and GCT\_PLETH T-052 keep their separate arithmetic comparisons.

The external Green and transference inputs are used only through the already whole-checked DO085 interfaces: Gorodetsky and Hazan (2017), Theorem2.1/equation(2.1), through K067; Miller and Stephens-Davidowitz (2019), Theorem3.2/equation(3.3), through DO070. The companion is an exact endpoint/accounting diagnostic. The whole proof separately checks the new rank and paired-coordinate arguments. No rational-matrix enumeration, Lean build or archive audit is claimed.

\clearpage

\section{AM-DO-097 --- all non-all-II unipotent rank-two profiles at support24}\label{am-do-097-all-non-all-ii-unipotent-rank-two-profiles-at-support24}

\textbf{Canonical checked integration --- 2026-11-22T23:03:47.609368+00:00.} Root and independent whole-proof/script/interface readings and separate805-assertion replays passed without correction. Historical pending wording below is preserved. The entire globally unipotent rank-two R071 support24 sector with at least one type-I or type-III place has absolute contribution\textless2\^{}-71310. The exact block normal form gives N²Z=0 and alpha=Uq\textless=H+H\^{}-1-2, so all original S and fixed transverse inputs have only type-II poles and the full covector radius is sqrt(R/U). At most24 padded23-place ideals cover every non-all-II profile, with all164 coordinate sectors, common G patterns, genuine selected trace, true projected-column U denominator, paired T\^{}4 index and source fibre retained. The complete maximum is -292137655/4096; one71310 charge gives80/31 entries, both full sums\textless2\^{}-230. Complete floors24/11 remain. The all-II unipotent case, rank3 and fixed-block residuals, signed/global/source/arithmetic/032/DP44 gates remain open.

Author: \nolinkurl{/root/definite_oscillator}, ALP\_MACHINE. 22 November2026 UTC.

Controlled pending proof. Whole author/source/script verification and independent review are separate records. No canonical file or chronological error list is modified here.

\subsection{1. Exact sector and statement}\label{exact-sector-and-statement-1}

Retain the unchanged J56 all-top kernel, H=133, f=168, n=169, all actual CM and G/T places, source/target pencils, finite projectors, selected Hilbert metric, packet representative, central normalization, Haar measures, class masses and phases. Work in the actual R071 residual at support s=24, now under the global hypotheses

\[
 h=I+N,\quad \operatorname{rank}_E N=2,\quad N^3=0,
 \quad h_{ZZ}=I,\quad Z=S\perp V,\quad \dim S=164,\quad
 V=u^\perp\subset E^4,\quad P=Eu\perp Ee_n,
 \quad E^{169}=Z\perp P.                                    \tag{UC1}
\]

The full lower cross map C\_Z:Z→P has rank one and its restriction to S is nonzero. The determinant-one and inverse identity-compression premises are those already proved in DO095. Every unmixed G reduction is unipotent. The original packet vector u is primitive, u*u=U=77Q²−14, and J\_P=diag(U,−1).

\textbf{Theorem.} The total absolute actual normalized contribution of UC1 matrices for which at least one of the24 active places is type I or type III is less than

\[
                              2^{-71310}.                   \tag{UC2}
\]

This removes every non-all-II incidence profile in this globally unipotent rank-two support24 branch, with one new all-top charge. The all-II profile is excluded. Nonunipotent rank-two terms belong to DO096, and rank-three and other structural classes are not addressed. This proof makes no independent full support-floor claim.

The proof retains the complete R072 projected two-coordinate negative-column count and positive-rank-two source fibre. Its new inputs are integrality on the full fixed positive input lattice at types I and III, and the sharper joint covector norm supplied by global unipotence. No central-fibre refinement is needed in this theorem.

\subsection{2. Full positive compression and the exact norm}\label{full-positive-compression-and-the-exact-norm}

Normalize the null image line only over E:

\[
 C_Z=\ell r,\quad\ell=(1,\eta)\in P,\quad\eta\bar\eta=U,
 \qquad v=J_Z^{-1}r^*,\quad q=rv\in F_{>0}.                 \tag{UC3}
\]

The inequality q\textgreater0 holds at every real embedding because Z is positive definite and r is nonzero. The rank-two unipotent case of DO095(RA9--10,RA21--23) has D\_P ell=ell. Unitarity gives B\_Z=−v ell\emph{J\_P D\_P. Also B\_Z ell=0. Multiplying the P-block unitary identity on the left by ell} therefore gives ell\emph{J\_P D\_P=ell}J\_P, so

\[
                            B_Z=-v\ell^*J_P.                \tag{UC4}
\]

Since D\_P is two-dimensional, unipotent and fixes ell, its difference from I has image contained in E ell and annihilates that line. Consequently D\_P=I+d ell ell*J\_P for a scalar d in E. The P-block unitary identity now gives d+bar(d)=−q. Direct block multiplication yields

\[
 N=\begin{pmatrix}0&-v\ell^*J_P\\\ell r&d\ell\ell^*J_P\end{pmatrix},
 \qquad
 N^2=\begin{pmatrix}0&0\\0&-q\ell\ell^*J_P\end{pmatrix}.
 \tag{UC5}
\]

In particular \textbf{N²Z=0} in the original rational positive space, and

\[
          \alpha=(N^2)_{nn}=Uq>0,\qquad
          \alpha_\tau\le R:=H+H^{-1}-2=17424/133.             \tag{UC6}
\]

For the last inequality, N²=h+h\^{}-1−2I and unitary adjointness give alpha=2Re(h\_nn)−2. At an embedding \textbar h\_nn\textbar=cosh(log\textbar\textbar h\textbar\textbar), hence alpha\textless=H+H\^{}-1−2. This is the unchanged signature-one matrix identity, with no small-distance approximation.

In the original positive metric, q\_tau is the squared norm of the covector r\_tau. Put

\[
                         M=\sqrt{R/U},\qquad B=\sqrt R.     \tag{UC7}
\]

For the three fixed integral vectors y1,y2,z0 of R072(PJ4--5), and for the original S coordinates xi\_i=r(e\_i), UC6 gives

\[
 \|\xi_\tau\|\le M,\quad
 |r(y_j)_\tau|\le M\sqrt{L_j},\quad
 |\beta_{j_0,\tau}|=|\eta\xi_{j_0}|\le B=M\sqrt U,           \tag{UC8}
\]

where beta\_j0 is the original negative entry of a nonzero S pivot column. The exact squared lengths are

\[
 L_1=978665066829,\quad L_2=835948910828,\quad
 L_3=1007543201084982869062.                                 \tag{UC9}
\]

The weak inequalities suffice; the final directed arithmetic bounds are strict. The normalization ell=(1,eta) is used only over E to prove the identities. No integral null basis or U-free coordinate change is asserted by that normalization.

\subsection{3. Integrality for every fixed positive input outside type II}\label{integrality-for-every-fixed-positive-input-outside-type-ii}

Every original e\_i in S and each fixed y1,y2,z0 belongs separately to every actual source and target pencil at both CM partners. This is the direct membership established in K088(IP5--6), DO094(FT8--11), R071(GR6--7) and R072(PJ4--5). The result applies to each displayed integral vector, not to a moving basis inferred from positivity.

At a type-III arrow hM\_c=M\_d with c!=d, such a vector y satisfies

\[
 Ny=hy-y\in M_d,\qquad
 (h^{-1}-I)y=-Ny\in M_c,                                   \tag{UC10}
\]

where N²y=0 follows from UC5. Thus Ny lies in M\_c∩M\_d, which is the pencil's lower lattice and is contained in the original standard lattice. At the conjugate component the pencils are paired Hermitian duals; their intersection is the dual of the selected-component sum and is again contained in the standard dual lattice. Repeating the same memberships proves the assertion at both components. This extends K086(TI5--8) to these actual integral transverse inputs by its exact hypotheses.

At type I, h preserves the full pencil, hence its intersection M. Each of the displayed inputs belongs to M. Therefore Ny lies in M, again inside the original lattice, at both partners. At type II, use only Ny in the actual target pencil, which is contained in varpi\^{}-1 Lambda. Every original coordinate has at most a simple pole. All inactive places retain original integrality.

Since h\_ZZ=I, Ny=C\_Zy for these positive inputs. Their four positive plane coordinates are u r(y); fixed integer Bezout coefficients for the primitive u extract r(y) without a U denominator. The same membership applies to the original negative pivot beta\_j0. Accordingly all164 xi\_i, all three r(y\_j), and beta\_j0 have T poles only at type-II places, of order at most one at each CM partner.

There is an explicit multiplicity issue when counting a union of types: for a fixed ordered pencil pair, the actual type-II subset can depend on h. We retain it as follows. For each non-all-II matrix choose one actual non-type-II place j and allow poles at every other active place. The complement D=Sigma minus \{j\} has size23 and contains every possible type-II pole. There are at most24 choices of j. Thus the whole desired sector is covered by at most24 containing arithmetic sectors with

\[
 I_D=\prod_{v\in D}(\mathfrak p_v\bar{\mathfrak p}_v)^{-1},
                         N_EI_D=T^{-46}.                    \tag{UC11}
\]

This padding is conservative and is counted once. It avoids assuming that one smaller ideal is uniform over different types. The original complete P\_T\^{}24 pair sum is retained below; there is no additional type-profile or denominator-subset factor.

\subsection{4. The unchanged G ideal, joint packing, column and selected trace}\label{the-unchanged-g-ideal-joint-packing-column-and-selected-trace}

Use R072(PJ6)'s whole-C\_Z G orientation in the fixed G-unit Z/P frame. At rank one, either all Z columns vanish or the paired inverse-adjoint makes them vanish at the other partner; at rank zero both partners vanish. Its common ideal J has at most2\^{}668 patterns and

\[
           |a(h)|\le2Q^{-335(f-1)}(N_EJ)^{335/2}.             \tag{UC12}
\]

The fixed transverse/plane frame has determinant of modulus2U, a unit at the actual G and T places. The primitive original-coordinate extraction is used away those places, exactly as in R072. Thus for every padded sector UC11,

\[
 L=J I_D,\quad z=(N_EL)^{1/(2f)},\quad
 \xi_i,\ r(y_j),\ \beta_{j_0}\in L.                         \tag{UC13}
\]

If k is the number of nonzero original xi coordinates, then1\textless=k\textless=164 and kz²\textless=M², because a nonzero element of L has half-trace squared norm at least fz². The source/target and finite-G memberships concern the whole cross map, not merely the negative column; no additional orientation choice is made for the transverse inputs.

Retain the exact general joint-row packing lemma R072(PJ10--13). It applies also to k=1..5; its original numerical support23 theorem started at6, but the lemma itself is stated for arbitrary positive k. In dimension336 its spherical-cap constant is47 and e\textless87/32. For each k it gives the row count

\[
 {164\choose k}47^k2^{2fk}(f!)^k
             \left(\frac{(87/32)M^2}{fkz^2}\right)^{fk}.     \tag{UC14}
\]

With c=2+2\^{}-19, the three transverse entries and beta pivot have the joint count

\[
                   c^{8f}M^{8f}U^f(L_1L_2L_3)^f z^{-8f}.
 \tag{UC15}
\]

Indeed their individual radius-to-z ratios are at least sqrt(L\_j) and sqrt(U), all greater than2\^{}19; the scalar packing factor1+2R/z is at most cR/z. Their ratio reconstruction determines the full C\_Z over E, with at most164\textless2\^{}8 pivot choices. It does not declare the quotient or moving null vector integral.

The exact R072(PJ14--18) projected negative-column lattice and positive-rank-two source fibre remain unchanged:

\[
 \#\{\operatorname{pr}_P(h e_n)\}<191^{4f}U^fT^{4s},
                           G_2<2^{36}.                      \tag{UC16}
\]

In particular the fixed projection denominator U\^{}f is retained, as is the complete paired indexT\^{}(4s) at every actual source/target label, including zero. The count uses the original negative-column radius95 and the original full projection; it is not the new affine central-fibre proposal.

The actual selected normalized trace retains R061(GS11--12), as used in R072(PJ9):

\[
                         |c(h)|<2^{-929f/2}U^{-f/2}z^{-f}.
 \tag{UC17}
\]

Both compact factors, the genuine selected module, the actual packet representative and the original Hermitian metric are unchanged. The improved norm UC8 does not require a new contraction estimate.

For one fixed padded sector retain the Green correction2, G patterns2\^{}668, support choices2\^{}168, pivot2\^{}8 and fibre2\^{}36. Their fixed exponent is881. The complete original ordered-pair/weight sum is P\_T\^{}24\textless2\^{}(997*24), as in R059(GB19). The additional24 padding choices are applied only in the final union. This pair sum already contains every actual incidence type and supplies no separate independent CM-partner labels.

\subsection{5. Exact endpoints for all164 coordinate sectors}\label{exact-endpoints-for-all164-coordinate-sectors}

Combining UC12--17 leaves the same z exponent as R072,

\[
                               (326-2k)f.                    \tag{UC18}
\]

For k\textless=163 this is nonnegative; use z\textless=M/sqrt(k). For k=164 it is −2f. The integrality of J gives z\textgreater=T\^{}(-23/f), so in this last sector z\^{}(-2f)\textless=T\^{}46. This separate branch is essential.

Put

\[
 F_0=881+4f\log_2(191)+8f\log_2(c)
       +f\log_2(L_1L_2L_3)-929f/2-335\cdot167\log_2Q.
 \tag{UC19}
\]

After the endpoints, the common logarithmic exponents, before the complete rational k-factor, are bounded by

\[
\begin{split}
 C&=F_0+167f\log_2R-(331f/2)\log_2U
                  +(335\cdot23+4\cdot24)\log_2T+997\cdot24,\\
 C_{164}&=F_0+168f\log_2R-(333f/2)\log_2U
                  +(337\cdot23+4\cdot24)\log_2T+997\cdot24.
\end{split}                                                   \tag{UC20}
\]

Here F\_0 in a directed comparison uses upper bounds for its positive logarithms and a lower bound for logQ. The factors M\^{}(334f) and M\^{}(336f) have become R\^{}(167f)U\^{}(-167f) and R\^{}(168f)U\^{}(-168f), respectively; the remaining U powers in UC20 retain both UC15 and the projected-column factor UC16.

Use the actual exact one-sided bounds

\[
 \log_2T<1019333/1024,\quad\log_2Q>287843/4096,\quad
 \log_2U>601356/4096,\quad\log_2R<113/16,
\] \[
 \log_2(191)<15519/2048,\quad
 \log_2(L_1L_2L_3)<152789/1024,\quad
                       \log_2 c<4097/4096.                  \tag{UC21}
\]

Write C\^{}+ and C\_164\^{}+ for the rational upper bounds on UC20 obtained using UC21. They are C\^{}+=−293587639/4096 and C\_164\^{}+=−202198271/4096. For k\textless=163 use the complete positive rational factor

\[
 {164\choose k}47^k2^{2fk}(f!)^k(87/32)^{fk} f^{-fk}k^{-163f};
 \tag{UC22}
\]

for k=164 replace the final k power by k\^{}(-164f), exactly R072's negative-endpoint branch. The certificate computes the least strict integral binary ceiling of each full rational factor, including its binomial and all factorial/power terms. All164 full upper exponents are below−71322. The maximum is

\[
          C^++354=-292137655/4096<-71322,\qquad k=1.         \tag{UC23}
\]

There is no endpoint-union factor: each k uses its proved single endpoint branch. Applying the24 padding choices and164 possible support sizes gives

\[
       24\cdot164=3936<2^{12},\qquad
       \sum_{\text{desired sector}}|\mathcal W(h)|
              <3936\,2^{-71322}<2^{-71310}.                  \tag{UC24}
\]

This proves UC2 with every original pair and matrix factor retained.

\subsection{6. Scope, chronology and interfaces}\label{scope-chronology-and-interfaces}

This theorem removes the complete globally unipotent rank-two R071 support24 sector except the all-II profile. It does not show that the all-II profile is empty or small. The same sufficient bound with23 replaced by24 is positive in every k-cell and is recorded only as a failed estimate, not as evidence of an actual term. The affine central-fibre reconstruction from UC5 remains separate research and is not needed for UC24.

The actual charged predecessor at freezing is R075,79 all-top and31 uniform entries, including DO096's2962487 followed by R075's8866. Append71310 once to obtain80/31, retaining every earlier charge and the unchanged complete floors24/11 unless a separate assembly changes them. Any intervening independently checked charge must also be retained upon canonical integration. The certificate checks this complete chronological extension and both full dyadic sums below2\^{}-230; it does not evaluate an automorphic sum.

The new ingredients are the exact UC5 full-positive zero-square block, its original two-lattice membership consequence for all fixed integral transverse inputs, the sharp joint norm UC6--8, and the explicit24-fold padded-ideal cover before the complete R072 count. Direct interfaces are DO095(RA3--10,RA21--26), R071(GR6--12), K086(TI3--9), K071(NB2--5), LIFT44(PR4--7), K088(IP5--6), DO094(FT8--11), R072(PJ1--21), R061(GS11--12) and R059(GB19). These were read in their stated scopes; no type-II-only slice or finite positive-definiteness inference is imported.

The module and fractional-ideal setting remains over the unchanged Spec(O\_E), with its fixed CM pairing and complex half-trace realization. Carried contextual anchors remain SGA3 I4.7.3/4.7.3.1; SGA6 VI Theorem3.3; FGA232 Theorem3.1 pp148--149; FGA236 Corollaries1.7/2.3 pp229/232; SGA7 I VII2.7/2.9 pp169--174 and VIII7.4(b) pp310--311; cote4 scan003 §§8.1.1--8.1.2, scans005--006 §§8.4.1--8.4.3 and cote54 scans005--007. No new K-theory, lambda-ring, l-adic or motivic assertion is made. The signed sufficient inequality, prescribed global period, actual selected arithmetic image, source constructor and DP44 comparison remain open.

\clearpage

\section{AM-R-076 --- the full positive covector and its relative norm fibre remove rank three at support24}\label{am-r-076-the-full-positive-covector-and-its-relative-norm-fibre-remove-rank-three-at-support24}

\textbf{Canonical checked integration --- 2026-11-22T23:12:44.613340+00:00.} Root and independent whole-proof/script/interface readings and separate41-assertion replays passed without mathematical correction. The contextual locator typo was corrected before freezing. Historical pending wording below is preserved. The entire global-rank-three R071 support24 cross sector has absolute contribution\textless2\^{}-5710, with no sparsity premise. The full positive167 covector is counted in its actual Gram metric, with exact determinant4U and arbitrary-ideal covering radius; its last scalar is counted by the exact relative norm fibre, including nonprincipal and non-conjugation-stable ideals. Two pointwise Green bounds on the same original G patterns meet at z=2\^{}-100; full low/high exponents are -23390187/4096 and -361072601/4096. All original trace, paired targets, projected U denominator, source fibre, support, pivots and incidence weights are retained. One5710 append gives81/31 complete errors below2\^{}-230 and unchanged floors24/11. Rank-two all-II and fixed-block residuals, signed/global/source/selected-image/032/DP44 gates remain open.

Author: \nolinkurl{/root}, GPT-6 Astra Pro, inherited effort not externally exposed, ALP\_MACHINE. 22 November 2026 UTC. Controlled proof; whole proof, literal interfaces and exact arithmetic are checked separately before promotion.

\subsection{1. Actual scope and theorem}\label{actual-scope-and-theorem-1}

Keep the unchanged J56 all-top H=133 kernel, f=168, n=169, and the actual R071 doubly-plane/four-zero remainder at support s=24. Every paired source/target pencil, finite projector, actual selected metric and trace, packet representative, central normalization, class mass and phase remains fixed. In the original fixed decomposition

\[
 E^{169}=Z\perp P,\quad Z=S\perp u^\perp,
 \quad\dim S=164,\quad\dim Z=167,\quad P=Eu\perp Ee_n,
 \qquad h=\begin{pmatrix}I_Z&B_Z\\C_Z&D_P\end{pmatrix},
 \quad J_P=\operatorname{diag}(U,-1),
 \tag{CV1}
\]

the cross map C\_Z has rank one and is nonzero on S. Restrict to global rank\_E(h-I)=3. No global unipotence is assumed. All334 unmixed G reductions are unipotent, as in the actual K068 residual.

\textbf{Theorem.} The entire absolute contribution of this rank-three support24 sector is less than

\[
                              2^{-5710}.                     \tag{CV2}
\]

This is one complete-sector bound, without a sparsity restriction on the original scalar row. It removes the intersection with the current remainder and retains all old charges, including R073's overlapping sparse rank-three bound. It supplies no independent complete support-floor increase and no signed/global/arithmetic/DP44 conclusion.

The new count treats the entire167-dimensional positive covector in its actual Gram metric and counts the last cross scalar through its exact relative norm equation. The two complementary Green bounds are applied to the same original ideal patterns, at a fixed numerical cutoff for their norm.

\subsection{2. Original scalar data and the two genuine Green estimates}\label{original-scalar-data-and-the-two-genuine-green-estimates}

Use the original primitive vector u, U=u*u=77Q²−14, Q=1427908193729121522601, and the three integral transverse inputs of R072(PJ4--5):

\[
\begin{split}
 y_1&=(-976627,-153032,38026,0),\\
 y_2&=(-567678,695570,-172838,0),\\
 y_3&=(256766,-7654233287,-30805128044,1).
\end{split}                                                  \tag{CV3}
\]

Write C\_Z x=(u lambda(x),eta lambda(x)) over E, with eta bar(eta)=U. Let xi\_i=lambda(e\_i) for the original S coordinates. Choose any nonzero pivot xi\_j and put beta=eta xi\_j. There are at most164\textless2\^{}8 pivot choices. The entire covector lambda and beta determine C\_Z by ratio reconstruction over E; no quotient is declared integral.

The whole-C\_Z finite-G orientation of R072(PJ6) gives a single common E-ideal J, with at most2\^{}668 patterns. Let

\[
 L=J I_\Sigma,\qquad N_E I_\Sigma=T^{-2s},\qquad
 z=(N_E L)^{1/(2f)},\qquad z\ge T^{-s/f}.                    \tag{CV4}
\]

All164 xi\_i, the three values lambda(y\_j), and beta belong to L. These are R072's original-coordinate, primitive-extraction and two-lattice memberships, at both CM partners. They introduce no U denominator into L. The first genuine pointwise finite bound is

\[
 |a(h)|\le 2Q^{-335(f-1)}T^{335s}z^{335f}.                   \tag{CV5}
\]

For global rank3, DO095(RA15--20) gives the additional original nonzero three-minor witness, with conjugates below2\^{}9 and only simple paired type-III poles. Since their number t\textless=s, its second pointwise bound is

\[
 |a(h)|<2Q^{-501(f-1)}2^{9\cdot167f}T^{167s}.                \tag{CV6}
\]

CV6 is uniform and requires no J-pattern or witness enumeration of its own. We nevertheless retain the complete2\^{}668 count already required for CV4. Both finite mixed projectors remain present with modulus bound1.

The exact negative-row norm and selected trace are R072(PJ7--9):

\[
 \|\lambda_\tau\|<M=\frac{H}{2\sqrt U},\qquad
 |\beta_\tau|<B=H/2,\qquad
 |c(h)|<2^{-929f/2}U^{-f/2}z^{-f}.                           \tag{CV7}
\]

The first is one joint norm on the positive167-space at each embedding. The last is the actual normalized selected trace, with its compact factors and original Hilbert metric, not an operator norm replacement.

\subsection{3. A joint count in the full positive Gram metric}\label{a-joint-count-in-the-full-positive-gram-metric}

In the fixed E-basis consisting of the164 original S vectors and y1,y2,y3, the positive Gram matrix is diag(I\_164,G), where G=(y\_i dot y\_j). Direct exact arithmetic gives

\[
                  \det G=4U,\qquad G-I_3>0.                \tag{CV8}
\]

For the second assertion the three leading principal minors of G-I\_3 are strictly positive; the certificate records and checks all three integers. Thus the metric on the coordinate values of a covector is diag(I\_164,G\^{}-1). The real linear transformation to that metric is a contraction of the ordinary product metric. This transformation is used for an archimedean volume and covering calculation only; it changes no arithmetic lattice or integral basis.

In the half-trace complex metric the original ideal has covolume D0 z\^{}(2f), where D0=7\^{}140 29\^{}144 5\^{}84 13\^{}84. Therefore the product L\^{}167 in the covector metric has exact real covolume

\[
              \frac{D_0^{167}z^{334f}}{(4U)^f}.              \tag{CV9}
\]

Each of the f complex labels contributes a real determinant factor det(G)\^{}(-1). The accepted arbitrary-ideal covering estimate used in DO085(MR14) and K086(TI14) gives an ordinary product cover at most(2250/7)sqrt(167f)z. The contraction in CV8 preserves this upper bound, and (2250/7)\^{}2*167\textless4200². Thus the covering radius in the true covector metric is less than4200 sqrt(f)z.

The whole actual covector lies in a real ball of radius M sqrt(f). Disjoint fundamental cells within the enlarged ball, with pi\textless22/7, give

\[
 \#\{\lambda\}\le A(4U)^f(M+4200z)^{334f}z^{-334f},
 \qquad
 A=\frac{528^{167f}}{(167f)!D_0^{167}}<2^{-387281}.            \tag{CV10}
\]

This counts every possible zero and nonzero coordinate pattern together; there is no binomial or additional k-sum. The actual set has a nonzero original pivot, but allowing more covectors in the upper count is harmless. This is a full-volume argument with an explicit cover for the actual fractional ideal, not an asymptotic density estimate.

\subsection{4. The remaining cross scalar is a relative norm fibre}\label{the-remaining-cross-scalar-is-a-relative-norm-fibre}

Fix a counted covector with chosen xi\_j!=0. Its pivot beta belongs to L and satisfies

\[
                       \beta\bar\beta=U\xi_j\bar\xi_j.       \tag{CV11}
\]

For each actual beta the E-ideal a\_beta=(beta)L\^{}-1 is integral. Its relative F-ideal norm is the fixed ideal

\[
 \mathfrak c=(U\xi_j\bar\xi_j)(N_{E/F}L)^{-1},\qquad
 N_F\mathfrak c=|N_E\beta|/N_E L
                    \le(B/z)^{2f}.                         \tag{CV12}
\]

If no beta exists this fibre is empty. Otherwise c is integral by its displayed ideal norm; no separate divisibility or principality assumption on L is used. For each split F-prime of exponent e in c there are e+1 choices for an integral relative-norm ideal; at inert or ramified primes there is at most one. Hence the number of eligible E-ideals is at most d\_F(c). For each one, any two actual generators with the exact norm CV11 differ by a relative-norm-one unit. The carried CM root-unit calculation gives \textbar mu(E)\textbar=28, so there are at most28 such generators. This is exactly K075(RN17--19)'s relative ideal/generator argument in the present fixed-ideal form.

Use the checked actual-field divisor estimate DO091(TD11--12),

\[
 d_F(\mathfrak c)\le C_F(N_F\mathfrak c)^{1/12},\qquad
 C_F=\frac{2^{300}}{29^6 41^{14}2129^{14}}<2^{42}.
\]

It follows that the complete beta fibre satisfies

\[
                 \#\{\beta\}\le28C_F B^{f/6}z^{-f/6}.       \tag{CV13}
\]

The fixed norm equation replaces an independent complex scalar count. No generator or class-group multiplicity is lost: the ideals are counted before testing which are principal, and each actual beta belongs to one such ideal. The same argument works for a non-conjugation-invariant J; relative ideal norm is used rather than extension of an F ideal.

\subsection{5. Every remaining factor and the two norm ranges}\label{every-remaining-factor-and-the-two-norm-ranges}

Once C\_Z is known, retain the complete R072(PJ14--18) projected negative-column count and source fibre:

\[
          191^{4f}U^fT^{4s},\qquad G_2<2^{36}.               \tag{CV14}
\]

The U\^{}f factor is the actual fixed rational projection denominator. The paired target index T\^{}(4s) includes every label, including zero. Fixing C\_Z and the plane column reconstructs the full negative column by unitarity; the original source-lattice fibre still bounds every remaining arrow. We take no saving from DO095's sharper uniqueness observation.

As in R072, retain the Green Euler factor2, common G patterns2\^{}668, support choices2\^{}168, pivot2\^{}8 and source fibre2\^{}36, whose total exponent is881. Retain the complete actual ordered-pair/absolute-incidence sum P\_T\^{}24\textless2\^{}(997*24). It already includes all incidence types and the paired CM labels. No normalized trace, source/target mass or Haar factor is changed.

Split the actual ideal norms at

\[
                              z_*=2^{-100}.                 \tag{CV15}
\]

The actual integers give M\textgreater2\^{}-68 and T\^{}(-24/f)\textless z\_\emph{. For z\textless=z\_},

\[
 4200z/M<4200\,2^{-32}<2^{-19},\qquad
                  (1+2^{-19})^{334f}<2.
\]

Thus CV10 is bounded by2A(4U)\^{}f M\^{}(334f)z\^{}(-334f) throughout the lower range, including at z\_*. Multiplying this, CV13, CV14 and the genuine trace CV7 by CV5 leaves only z\^{}(-f/6). Since J is integral, CV4 bounds this by T\^{}(s/6).

For z\textgreater=z\_*, use CV6 instead. Apart from positive constants the complete counting/trace factor is

\[
              (M/z+4200)^{334f}z^{-(1+1/6)f},               \tag{CV16}
\]

which is decreasing for every z\textgreater0. The high-range upper bound is therefore its value at z\_\emph{, where the same factor2 approximation to the enlarged radius is valid. This proves an upper bound without asserting that any ideal has norm z\_}.

Each fixed J belongs to exactly one of these two ranges. The full pattern enumeration is already in881. Bounding each pattern by the larger of the two uniform endpoints requires no additional factor2. The proof does not sum two complete2\^{}668-pattern bounds.

Use the exact directed quantities

\[
 L_T=1019333/1024>\log_2T,\quad
 L_Q=287843/4096<\log_2Q,\quad
 L_H=28900/4096>\log_2H,\quad
 L_U=601356/4096<\log_2U,\quad
 L_{191}=15519/2048>\log_2(191).
\]

With A\textless2\^{}-387281,28\textless2\^{}5 and C\_F\textless2\^{}42, the common logarithmic constant, before the chosen Green factor and remaining z power, is

\[
\begin{split}
 C={}&881+1+5+42-387281+4fL_{191}-929f/2+2f\\
    &\quad +(2005f/6)(L_H-1)-(331f/2)L_U
                  +4sL_T+997s=-522439755/128.               \tag{CV17}
\end{split}
\]

The1 is the controlled radius enlargement,2f is the factor4\^{}f from detG, and the powers of U include both the projected column and the selected trace. The H/2 exponent is334f+f/6=2005f/6. Every constant is retained before either endpoint is evaluated.

The full directed exponents at s=24 are

\[
\begin{split}
 E_{\rm low}&=C-335\cdot167L_Q+(2011s/6)L_T
                        =-23390187/4096<-5710,\\
 E_{\rm high}&=C-501\cdot167L_Q+9\cdot167f+167sL_T
                      +100(2011f/6)
                        =-361072601/4096<-88152.
\end{split}                                                  \tag{CV18}
\]

Both are strictly below-5710, proving the single complete-sector bound CV2. There is no sparsity, type-profile, norm-range or witness union left to charge.

\subsection{6. Chronology, geometry and unresolved arithmetic}\label{chronology-geometry-and-unresolved-arithmetic}

Append5710 exactly once to the actual chronological all-top array at checked integration, preserving every prior charge and the entire uniform array. The companion checks that actual predecessor and both complete dyadic sums below2\^{}-230. The complete floors remain24/11 unless a separate all-class assembly changes them. Rank-two unipotent and fixed-block classes at support24 are outside this theorem; their status is not inferred from the rank-three Green witness.

The new geometric count is a fixed archimedean realization of the dual of the original positive167-module. Its determinant is kept exactly in CV9, and the rational Gram transformation is never called integral. CV11--13 count the exact scalar norm fibre over the original fractional ideal, including its ideal classes and the finite CM norm-one unit group. Both CM components, the polarization, Picard data, theta line, biextensions and selected coefficient metric remain unchanged.

Exact carried contextual anchors are SGA3 I4.7.3/4.7.3.1; SGA6 VI Theorem3.3; FGA232 Theorem3.1 pp148--149; FGA236 Corollaries1.7/2.3 pp229/232; SGA7 I VII2.7/2.9 pp169--174 and VIII7.4(b) pp310--311; cote4 scan003 §§8.1.1--8.1.2, scans005--006 §§8.4.1--8.4.3 and cote54 scans005--007. No new lambda-ring/K-theory equality, l-adic sheaf/weight theorem, CM-motive comparison or motivic L-function result follows. DP\_MULT Q3e and THETA3 G1/A7 gain exactly this smaller finite/real remainder; CX\_GRH CX-ALG-024 and GCT\_PLETH T-052 keep their distinct arithmetic comparisons.

The signed sufficient inequality, actual fully admissible mixed-CM source, selected arithmetic-image content, prescribed global Bessel/Mram value, Read/STiesR/specific032 pairing and every changed-family DP44 term remain open. The external arbitrary-ideal transference theorem and exact relative-ideal inputs are used only through the already checked DO085/K086 and K075/DO091 interfaces. The companion verifies the actual Gram minors, determinant, volume, divisor constant, cutoff inequalities, complete exponents and chronological error sums. It does not enumerate the arithmetic kernel or prove a global period by finite calculation. No Lean build or archive audit is performed.

\clearpage

\section{AM-DO-098 --- the all-II unipotent scalar-cross sector at support24}\label{am-do-098-the-all-ii-unipotent-scalar-cross-sector-at-support24}

\textbf{Canonical checked integration --- 2026-11-22T23:23:41.843035+00:00.} Root and independent whole-proof/script/interface readings and separate651-assertion replays passed without mathematical correction. Historical pending wording below is preserved. The entire all-II globally rank-two-unipotent R071 scalar-cross support24 sector has absolute contribution\textless2\^{}-2879. The actual cross map and original central entry w reconstruct the full matrix. Its affine difference lattice iJ0 I\_F² retains the rank-zero G congruence, real codifferent and full discriminant, giving the complete central count2\^{}479 Q\^{}167 z\^{}-2f. The true transverse Gram volume, exact split-relative norm fibre, every original all-II paired weight and the actual selected trace are retained. All164 coordinate sectors are exhausted by64/99/1 complete monotone ranges with strict maxima -2010899/512,-41042253/4096,-11800269/4096. One2879 append gives82/31 errors below2\^{}-230 and unchanged floors24/11. Together with DO097 and R076/DO096 this removes every R071 scalar-cross support24 residual. The fixed-S rank-two sector, signed/global/source/selected-image/032/DP44 gates remain open.

Author: \nolinkurl{/root/definite_oscillator}, ALP\_MACHINE. 22 November2026 UTC.

Controlled pending proof. Author whole-proof/interface/script verification, independent review and canonical integration are separate events.

\subsection{1. Exact sector and theorem}\label{exact-sector-and-theorem}

Keep the original J56 all-top normalized kernel, H=133, f=168, every actual field, CM place, source/target pencil, finite projector, selected Hilbert metric, packet representative, central phase, Haar measure and class mass. Work in the R071 residual, with

\[
 E^{169}=Z\perp P,\quad Z=S_{164}\perp V_3,\quad
 V=u^\perp\subset E^4,\quad P=Eu\perp Ee_n,\quad
 h_{ZZ}=I,\quad h=I+N,\quad \operatorname{rank}_E N=2,
 \quad N^3=0.                                                   \tag{CF1}
\]

The full cross map C:Z→P has nonzero restriction to S. All334 unmixed G reductions are unipotent. The support is s=24 and every active T place has type II in K070/K071's actual incidence classification. Thus global unipotence is a premise, and this proof does not apply to the nonunipotent branch or to fixed-S matrices with C\textbar S=0.

\textbf{Theorem.} The entire CF1 all-II sector has total absolute actual normalized contribution strictly less than

\[
                              2^{-2879}.                       \tag{CF2}
\]

The proof counts the actual matrix injectively from its cross map and one original central entry. It retains the rank-zero congruence on that entry, the true transverse metric, the relative norm constraint on the pivot, and all finite weights. It does not infer a smaller source fibre from dimension alone. The separate fixed-S rank-two support24 branch remains outside the statement.

Set Q=1427908193729121522601, U=77Q²−14 and

\[
 R=H+H^{-1}-2=17424/133,\quad M=\sqrt{R/U},\quad
 D_0=7^{140}29^{144}5^{84}13^{84},\quad D_F=2^{168}D_0.          \tag{CF3}
\]

\subsection{2. The actual cross map and an injective central reconstruction}\label{the-actual-cross-map-and-an-injective-central-reconstruction}

Use the exact globally unipotent normal form proved in DO097(UC3--8), with normalization only over E:

\[
 C=\ell r,\quad \ell=(1,\eta)\in P,\quad \eta\bar\eta=U,
 \quad v=J_Z^{-1}r^*,\quad q=rv>0,
\] \[
 N=\begin{pmatrix}0&-v\ell^*J_P\\\ell r&d\ell\ell^*J_P\end{pmatrix},
 \quad d+\bar d=-q,\quad J_P=\operatorname{diag}(U,-1),
 \quad (N^2)_{nn}=Uq=\alpha>0.                                \tag{CF4}
\]

At every embedding alpha≤R, so \textbar\textbar r\_tau\textbar\textbar≤M in the original positive covector metric. The original S entries xi\_i=r(e\_i) satisfy sum\_i\textbar xi\_i\textbar²≤M². Let k be their nonzero support size, 1≤k≤164. Choose the first nonzero pivot, retaining the harmless multiplicity164\textless2\^{}8. The164 xi\_i, three values r(y1),r(y2),r(z0), and the original negative pivot beta=eta xi\_j determine r and eta, and hence C, over E. The fixed transverse vectors are

\[
 y_1=(-976627,-153032,38026,0),\quad
 y_2=(-567678,695570,-172838,0),\quad
 z_0=(256766,-7654233287,-30805128044,1).                       \tag{CF5}
\]

Their Gram determinant is4U. The original u is primitive and (y1,y2,z0,u) has determinant of modulus2U, a unit at every selected G and T prime. No claim of a globally integral orthogonal basis is made.

Put w=N\_nn.~In CF4 the nn entry of ell ell*J\_P is−U, so

\[
                     w=-Ud,\qquad \Re w=\alpha/2.             \tag{CF6}
\]

Therefore C and w determine d, both off-diagonal blocks and the entire matrix h uniquely. There is no remaining lattice-isometry fibre or independent alpha/eta choice after these data are fixed. In particular the old projected-column count and G\_2 factor are replaced by the explicit central count below, not simply deleted.

At a complex place \textbar h\_nn\textbar=cosh r\_h and Re(h\_nn)=1+alpha/2≥1. Hence

\[
 |\Im w|^2\le\cosh^2 r_h-1=\sinh^2 r_h<H^2/4.               \tag{CF7}
\]

\subsection{3. Original local membership and the common finite ideal}\label{original-local-membership-and-the-common-finite-ideal}

Every original e\_i in S and each displayed transverse input belongs separately to every actual source and target pencil at both CM components: K088(IP5--6), DO094(FT8--11), R071(GR6--7), R072(PJ4--7), and DO097(UC10--13). For hM\_c=M\_d, Ny belongs to M\_d and the all-II target is contained in varpi\^{}-1 Lambda. The paired component uses its exact inverse-adjoint dual lattice. Primitive integer Bezout coefficients of u extract r(y) from its original positive coordinates without division by U. Thus all xi\_i, all three transverse values and beta have at most simple poles at the two CM primes of each active T place and are integral at inactive places.

Use R072(PJ6)'s common orientation of the \emph{whole} cross map at G. At rank one one CM component has C=0; at rank zero both do. This gives one E-ideal J, with at most4\^{}334=2\^{}668 patterns, for which

\[
 |a(h)|\le2Q^{-335(f-1)}(N_EJ)^{335/2},\quad
 L=J I_\Sigma,\quad N_EI_\Sigma=T^{-48},\quad
 z=(N_EL)^{1/(2f)},                                          \tag{CF8}
\]

and every counted cross scalar lies in L. No new orientation is chosen for the transverse values or pivot. Here a(h) is the unchanged actual normalized finite G coefficient, with the two mixed factors treated exactly as in R059/R072. At one relative G place of residue norm q\_v the ideal J has norm exponent0,1,2 when the reduction rank is at least2,1,0, respectively. Since the product of the334 residue norms is Q\^{}167,

\[
 N_EJ\le Q^{334}.                                           \tag{CF9}
\]

Let J\_0 be the F-ideal consisting of each rank-zero G prime once. Then at either CM component w=N\_nn is divisible by that prime. The same exponent table gives

\[
 N_EJ\le Q^{167}N_FJ_0,\qquad N_FJ_0\le Q^{167}.             \tag{CF10}
\]

All finite divisibility statements are in the original fixed frame or the proved G-unit frame. No positive-definite metric over a finite field is used.

The nonzero minimum of L in the complex half-trace metric is at least sqrt(f)z. Consequently

\[
 kz^2\le M^2,\qquad
 z\le Q^{167/f}T^{-24/f}<M/\sqrt{30}.                         \tag{CF11}
\]

The final strict inequality is the exact actual-host comparison R\^{}168 T\^{}48\textgreater30\^{}168 U\^{}168 Q\^{}334, checked by integer powers.

\subsection{4. The affine central fibre and its rank-zero congruence}\label{the-affine-central-fibre-and-its-rank-zero-congruence}

The full matrix entries at type II have pole order at most2 at both CM partners, by K071(NB5) and the actual two-pencil factors. Away the selected support entries are integral. Fix all C data and a G pattern. All admissible w have the same real part alpha/2, so their differences lie in iF and in the extended ideal J\_0 I\_\{\Sigma,F\}² O\_E. Since i is an integral unit, O\_E∩iF=iO\_F; contraction of an extended ideal gives the exact containing difference lattice

\[
 \Gamma=iJ_0 I_{\Sigma,F}^2,\qquad
 N_F(J_0 I_{\Sigma,F}^2)=N_FJ_0\,T^{-48}.                    \tag{CF12}
\]

Thus the fibre is contained in one affine coset of Gamma, if nonempty. Neither w nor its imaginary part is asserted to lie in Gamma itself: the shift alpha/2, including possible dyadic denominators, is retained. There is no division by2 in the difference lattice.

Write z\_c=(N\_FJ\_0 T\^{}-48)\^{}(1/f). In the real trace metric its covolume is sqrt(D\_F) z\_c\^{}f.~Its Euclidean dual is the real codifferent times the inverse ideal. The accepted real ideal estimate DO091(TD20--22), from LIFT47 transference, gives

\[
                         \mu(\Gamma)<750\sqrt f\,z_c.        \tag{CF13}
\]

This is the full real covolume and coefficient750, not the coefficient375 for an additional half-lattice. CF10 and the exact inequality750\^{}168 Q\^{}167\textless T\^{}48 imply mu(Gamma)\textless sqrt(f). The points satisfying CF7 have radius(H/2)sqrt(f), and their translated Voronoi cells fit in a ball of radius68sqrt(f). Hence

\[
 \#w\le V_F(68)\frac{T^{48}}{N_FJ_0}
       \le V_F(68)Q^{167}z^{-2f},                            \tag{CF14}
\] \[
 V_F(68)=\frac{((22/7)f68^2)^{f/2}}{(f/2)!\sqrt{D_F}}<2^{479}.
                                                                    \tag{CF15}
\]

The second inequality uses CF10 and N\_EJ=T\^{}48 z\^{}(2f); the paired T factors cancel exactly. The rational constant22/7 is only an upper envelope for pi. The true real dimension is f, not2f.

We will also use a less sharp count that is independent of J\_0. Enlarge Gamma to iI\_\{\Sigma,F\}², which contains iO\_F with indexT\^{}48. Each affine coset of iO\_F has at most(1+H)\^{}f=134\^{}f points in the CF7 ball, by its minimum sqrt(f). Therefore

\[
                              \#w\le134^fT^{48}.             \tag{CF16}
\]

These are two alternative upper counts for the same actual fibre, not factors to multiply.

\subsection{5. True transverse volume and the pivot norm equation}\label{true-transverse-volume-and-the-pivot-norm-equation}

Let G be the3×3 rational Gram matrix of CF5. The transverse covector norm in its value coordinates a=(r(y1),r(y2),r(z0)) is aG\^{}-1a*. Its possible values are contained in L³ with this metric. Therefore its real covolume is

\[
                   D_0^3 z^{6f}/(\det G)^f
                         =D_0^3 z^{6f}/(4U)^f.              \tag{CF17}
\]

For one arbitrary complex ideal the accepted DO070 bound is mu(L)\textless(2250/7)sqrt(f)z. Map separately rounded coordinates by the fixed dual transverse basis. Triangle and Cauchy--Schwarz give covering radius at most (2250/7)sqrt(3 tr(G\^{}-1))sqrt(f)z. The exact rational matrix computation proves

\[
        (2250/7)^2\,3\operatorname{tr}(G^{-1})<1024^{-2}.    \tag{CF18}
\]

By z≤M the resulting transverse point count is at most

\[
 4 A_3(4U)^fM^{6f}z^{-6f},\qquad
 A_3=\frac{528^{3f}}{(3f)!D_0^3}<2^{-4040}.                  \tag{CF19}
\]

Indeed the ball radius is at most(1+1/1024)M sqrt(f), (1+1/1024)\^{}(6f)\textless4, and (22/7)f=528. This counts an enlarged actual covector set; it neither declares the fixed transverse basis integral away G,T nor applies an original-ring covering estimate to an arbitrary fractional ideal.

For the fixed nonzero pivot xi\_j, beta belongs to L and obeys

\begin{verbatim}
                     beta bar(beta)=U xi_j bar(xi_j).
\end{verbatim}

Let l=N\_(E/F)L, so L bar(L)=l O\_E and N\_F l=N\_E L=z\^{}(2f). For every actual beta the integral ideal A=(beta)L\^{}-1 has relative norm

\[
 \mathfrak d=(U\xi_j\bar\xi_j)\mathfrak l^{-1},\qquad
             N_F\mathfrak d\le R^f z^{-2f}.                 \tag{CF20}
\]

Only split F-primes contribute more than one ideal with that relative norm. DO092(SN10--11)'s complete actual-field calculation gives at most2\^{}244(N\_F d)\^{}(1/16) choices. For each ideal there are at most28 generators with the prescribed relative norm, because their ratios are norm-one units in the unchanged CM field, whose root group is mu28; this is K075(RN17--19). Nonprincipal ideals simply give no beta. Thus

\[
                         \#\beta<2^{249}R^{f/16}z^{-f/8}.    \tag{CF21}
\]

No content-ideal count or class number factor is omitted: the integral ideal A is what is counted here, and the relative norm right side is fixed. Its integrality is asserted only for an actual solution. Beta is nonzero because the pivot and eta are nonzero.

\subsection{6. Full weights, three monotone endpoint ranges and exact totals}\label{full-weights-three-monotone-endpoint-ranges-and-exact-totals}

The unchanged actual selected trace has both upper bounds

\[
 |c(h)|\le1,\qquad |c(h)|<2^{-929f/2}U^{-f/2}z^{-f}.         \tag{CF22}
\]

For0\textless theta≤1 their minimum is at most 2\^{}(-929 theta f/2) U\^{}(-theta f/2) z\^{}(-theta f). This interpolation is an inequality for the same selected trace; it does not alter the vector, projector, compact factors or packet. We use theta=1 except at k=164, where theta=7/8.

For each actual pair tuple, R072(PJ10--12)'s general spherical-cap/Gaussian packing bound gives the original S-row count

\[
 {164\choose k}47^k2^{2fk}(f!)^k
          \left(\frac{(87/32)M^2}{fkz^2}\right)^{fk}.        \tag{CF23}
\]

This lemma holds for every positive k, not only the range in the earlier R072 numerical theorem. For all-II, each active place has at most T−1 diagonal labels, each with \textbar rho\textbar=1/(T−1). Their complete pair-weight sum is at most1 (K071 NB7). There is no extra independent conjugate-partner label. We retain support choices2\^{}168, common G patterns2\^{}668, pivot2\^{}8 and the Green correction2, a fixed exponent845. No type-subset factor is needed for this single all-II profile. Counts CF14--23 are uniform in each prescribed original pair, so this cancellation is applied only after fixing and counting that pair.

Use the same exact directed bounds as DO097:

\[
 L_T=1019333/1024>\log_2T,\quad L_Q=287843/4096<\log_2Q,
 \quad L_U=601356/4096<\log_2U,\quad L_R=113/16>\log_2R.      \tag{CF24}
\]

\textbf{Range A: 1≤k≤64.} Use CF14 and theta1. Before the S-row factor the z exponent is(325+7/8)f, hence after CF23 it is(325+7/8−2k)f\textgreater0. CF11 allows the single endpoint z≤M/sqrt(q\_k), q\_k=max(k,30). Combining all factors gives a common logarithmic upper bound

\[
 C_A=845+2+249+2f-4040+479-929f/2
 +166fL_R-(2647f/16)L_U+335sL_T-334\cdot167L_Q
       =61380333/512.                                        \tag{CF25}
\]

The remaining positive factor has square

\[
 F_A(k)^2=
 \frac{\binom{164}{k}^2 47^{2k}2^{4fk}(f!)^{2k}87^{2fk}
             q_k^{2fk}}
 {32^{2fk}f^{2fk}k^{2fk}q_k^{54747}}.                        \tag{CF26}
\]

Here54747=2f(2607/16); all paired and real exponents have been kept. The companion computes half the least strict integral binary ceiling of this exact rational square. Every one of these64 upper bounds is at most the maximum

\[
                      E_A(30)=-2010899/512<-3927.             \tag{CF27}
\]

\textbf{Range B: 65≤k≤163.} Use CF16 and theta1. The exponent after the row is(327+7/8−2k)f\textgreater0, so use z≤M/sqrt(k). Since134\^{}f\textless2\^{}1188, the common bound is

\[
 C_B=845+2+249+2f-4040+1188-929f/2
 +167fL_R-(2663f/16)L_U+337sL_T-335\cdot167L_Q
       =545420979/4096.                                      \tag{CF28}
\]

The squared coordinate factor is the numerator of CF26 without q\_k\^{}(2fk), divided by32\^{}(2fk)f\^{}(2fk)k\^{}55083. All99 cells are at most

\[
                     E_B(65)=-41042253/4096<-10020.           \tag{CF29}
\]

\textbf{Range C: k=164.} Keep CF16 and theta7/8. The z exponent after the row is exactly zero. The common bound is

\[
 C_C=845+2+249+2f-4040+1188-929\cdot7f/16
 +(2673f/16)L_R-(2663f/16)L_U+337sL_T-335\cdot167L_Q
       =585679155/4096.                                      \tag{CF30}
\]

The coordinate factor is binom(164,k)47\^{}k2\^{}(2fk)(f!)\^{}k87\^{}(fk) divided by32\^{}(fk)f\^{}(fk)k\^{}(164f), whose strict integral binary ceiling at164 is−145869. Thus

\[
                       E_C=-11800269/4096<-2880.             \tag{CF31}
\]

The three ranges are disjoint and exhaust all164 nonzero support sizes. No Gaussian/polynomial alternatives are multiplied and there is no endpoint-union factor. Their complete union satisfies

\[
 64\,2^{-3927}+99\,2^{-10020}+2^{-2880}<2^{-2879}.            \tag{CF32}
\]

This proves CF2 with all original finite and archimedean factors retained.

\subsection{7. Scope, dependencies and chronological accounting}\label{scope-dependencies-and-chronological-accounting}

This removes exactly the all-II, globally rank-two-unipotent scalar-cross support24 sector in R071. Together with the separately proved DO097 it treats the incidence profiles of that branch; the fixed-S rank-two sector, any other remaining structural class, and a full support-floor assembly remain separate. No positivity of the signed sum or prescribed global period is inferred.

The proof uses the complete actual normal form and local input memberships in DO097(UC1--17), R072(PJ1--18), R071(GR6--12), DO095(RA3--10,RA21--26), K071(NB5--9), and DO094(FT8--11). New steps are the injective original central-entry reconstruction CF6, its affine real difference ideal and rank-zero congruence CF12--15, the actual transverse metric volume CF17--19, the explicit relative norm application CF20--21, and the three complete weighted endpoint ranges. DO070 supplies the arbitrary complex-ideal transference; DO091(TD20--22) supplies the full real dual normalization; DO092(SN10--11) and K075(RN17--19) supply the checked split-prime ideal and norm-one generator bounds. These accepted inputs are applied in their literal scopes, not re-proved from an altered model.

Append2879 once to the actual complete ordered all-top predecessor, retaining all earlier charges and the31 uniform entries. The certificate records its precise predecessor file/hash, complete arrays and exact sum below2\^{}-230. An intervening checked root charge must be retained in canonical order; no historical array is silently replaced. The complete support floors remain24/11 until a separately reviewed exhaustion theorem changes them.

The modules and ideals remain over the original Spec(O\_E), with its fixed CM pairing and complex half-trace/real trace realizations. Carried contextual anchors remain SGA3 I4.7.3/4.7.3.1; SGA6 VI Theorem3.3; FGA232 Theorem3.1 pp148--149; FGA236 Corollaries1.7/2.3 pp229/232; SGA7 I VII2.7/2.9 pp169--174 and VIII7.4(b) pp310--311; cote4 scan003 §§8.1.1--8.1.2, scans005--006 §§8.4.1--8.4.3 and cote54 scans005--007. No new K-theoretic, l-adic, motivic or source-constructor assertion follows. The signed sufficient inequality, selected arithmetic image, prescribed periods and DP44 comparison remain open.

\clearpage

\section{AM-R-077 --- the complete fixed-block rank-two unipotent sector at support24}\label{am-r-077-the-complete-fixed-block-rank-two-unipotent-sector-at-support24}

\textbf{Canonical checked integration --- 2026-11-22T23:25:18.852238+00:00.} Root and independent whole-proof/script/interface readings and separate337-assertion replays passed without mathematical correction. Historical pending wording below is preserved. The entire original fixed-S globally rank-two-unipotent support24 sector has absolute contribution\textless2\^{}-6845. The complete R058 high finite-rank envelope extends with exponent -119372763/4096. At low finite rank the original square pairings force the null image into the packet plane; a type-I place is impossible by the unchanged pole-budget46 norm contradiction. All remaining II/III profiles are counted through the true transverse3 dual Gram metric, exact pivot relative norm fibre and original affine central-entry lattice, retaining both CM partners, full discriminant and actual incidence weights. The complete low exponent is -168257413/24576; the high/low union contributes one6845 charge. Every old charge, including DO098, is retained, giving83/31 errors below2\^{}-230 and unchanged complete floors24/11 until a separate assembly. Signed/global/source/selected-image/032/DP44 gates remain open.

Author: \nolinkurl{/root}, GPT-6 Astra Pro; inherited effort not externally exposed; ALP\_MACHINE. 22 November2026 UTC. Controlled proof; whole reading, exact computation and independent review are separately recorded before promotion.

\subsection{1. Actual scope and theorem}\label{actual-scope-and-theorem-2}

Keep the original J56 all-top H=133 kernel, f=168, n=169, the actual field, both CM partners, original pencil labels and incidences, finite projectors, selected Hilbert metric, packet representative, central phase, Haar measures and class masses. Consider the complete fixed-block sector

\[
 h=I_S\oplus h_0=I+N,\quad \dim S=164,\quad
 N^3=0,\quad\operatorname{rank}_E N=2,\quad
 \operatorname{rank}_E N^2=1,\qquad s=24.                    \tag{FB1}
\]

The unmixed G reductions are unipotent. All actual incidence profiles are included. Global nilpotence, global ranks and fixation of S are explicit premises, not inferred from finite reductions.

\textbf{Theorem.} The total absolute contribution of FB1 is strictly less than

\[
                              2^{-6845}.                    \tag{FB2}
\]

The proof extends the high-finite-rank portion of R058 and gives a new exact central-congruence and transverse-covector count in its low-rank packet plane. It does not infer an independent complete support-floor change, signed nonvanishing, an actual arithmetic source or a DP44 saving.

Use the actual integers

\[
 g_0=22285932761,\quad g_1=64072175441,\quad Q=g_0g_1,
 \quad U=77Q^2-14,\quad R=17424/133.
\]

The actual T is the unchanged J56 prime in DO040's certificate. Put alpha=(N²)\_nn.~R058(HB2--4) gives 0\textless alpha\_tau\textless=R and, at the334 unmixed relative G places,

\[
 d_v=\operatorname{rank}(N\bmod v)\in\{0,1,2\},\quad
 D_G=\sum_vd_v,\quad
 J_F=\prod_v\mathfrak p_v^{2-d_v},\quad
 N_FJ_F=Q^{334}/\prod_vq_v^{d_v}.                            \tag{FB3}
\]

The support-dependent formulas of R058(HB4--14) retain their exact local meaning at s=24; their proofs do not use s\textless=23.

\subsection{2. High finite rank and the exact low-rank plane reduction}\label{high-finite-rank-and-the-exact-low-rank-plane-reduction}

If D\_G\textgreater=33, put z=(N\_FJ\_F T\^{}-48)\^{}(1/(2f)). The new exact integer comparison is

\[
 Q^{334}2^{21\cdot1211}<g_0^{33}T^{48},
 \qquad z<2^{-1211/16}<1.                                   \tag{FB4}
\]

R058's complete original-column, source-fibre, gamma, trace, type and paired-label envelope HB11--14, with its genuine profile weights, therefore gives

\[
\begin{split}
 E_{\rm high}={}&1+5379+1062+107+168+3\cdot24+530-464f\\
 &-335\cdot167(287843/4096)+343\cdot24(1019333/1024)\\
 &+42-55776\cdot1211/16
       =-119372763/4096<-29143.                             \tag{FB5}
\end{split}
\]

There is one support value and the full type assignment count is already in this exponent. No new finite-rank pattern count is added to it.

Suppose D\_G\textless=32. R058(HB20--22), using the original integral transverse pairings, applies without a support bound to force

\[
 \operatorname{im}N^2=\ell\subset P=Eu\perp Ee_n,
 \qquad\alpha\in UJ_F I_{\rm II,F}^2 I_{\rm III,F}.           \tag{FB6}
\]

Here u is the fixed primitive four-coordinate packet vector, u*u=U. Its three integral orthogonal vectors are

\[
 y_1=(-976627,-153032,38026,0),\quad
 y_2=(-567678,695570,-172838,0),\quad
 y_3=(256766,-7654233287,-30805128044,1).                    \tag{FB7}
\]

Indeed each \textless y\_j,N²y\_j\textgreater{} is in O\_F intersect J\_F, independently of the number of T places. The unchanged norm contradiction at finite-rank budget69 already forces all three to vanish when D\_G\textless=32. The resulting original entry identity and coprimalities give the factor U in FB6.

If at least one place has type I, then mu=2u+t\textless=46. The already checked strict R058(HB23) inequality

\[
 (133U)^{168}Q^{334}>17424^{168}g_1^{32}T^{46}
\]

contradicts N\_F(alpha)\textless=R\^{}f exactly as in HB24. Thus the entire low-rank sector with a type-I place is empty. The remaining low-rank profiles use only types II and III. There are at most2\^{}24 type assignments; at each place their full ordered-pair counts multiplied by the actual absolute incidence weight are at most1, by K071(NB7). This cancellation includes the zero label and both paired CM components. It will be used only with counts uniform over all those original labels.

\subsection{3. The true three-dimensional covector and matrix reconstruction}\label{the-true-three-dimensional-covector-and-matrix-reconstruction}

Let V=u\^{}perp in the original positive E\^{}4 and Z=S perpendicular V. The line ell in FB6 is the unipotent null line. The induced action on ell\^{}perp/ell is unipotent and unitary for a positive definite Hermitian form at every archimedean place, so it is the identity. Therefore N(V) is contained in ell; S is fixed by premise. It follows that h\_ZZ=I. The cross map C:Z→P is nonzero on V: otherwise N acts only on the hyperbolic plane P, where a unipotent has square zero, contradicting rank N²=1.

The exact block calculation of DO097(UC3--6) consequently applies with a nonzero covector r supported on V rather than on S. For completeness its identities are

\[
 C=\ell r,\quad\ell=(1,\eta),\quad\eta\bar\eta=U,
 \quad v=J_Z^{-1}r^*,\quad q=rv>0,
\] \[
 N=\begin{pmatrix}0&-v\ell^*J_P\\\ell r&d\ell\ell^*J_P\end{pmatrix},
 \quad d+\bar d=-q,\quad J_P=\operatorname{diag}(U,-1),
 \quad\alpha=Uq,\quad w=N_{nn}=-Ud.                        \tag{FB8}
\]

To see that no additional plane choice remains, unitarity first gives B=-v ell\emph{J\_P D and D ell=ell. It then gives ell}J\_P D=ell\emph{J\_P and B=-v ell}J\_P. The unipotent two-dimensional D fixes ell, so D-I=d ell ell*J\_P; its unitary identity is d+bar d=-q. Thus C and the single original entry w determine the entire matrix, with no remaining isometry fibre.

Set M=sqrt(R/U). At every embedding \textbar\textbar r\textbar\textbar\textless=M in the positive covector metric. If a\_j=r(y\_j), then not all three a\_j vanish. Choose a nonzero pivot a\_j, retaining at most3\textless2² choices, and put beta=eta a\_j. All a\_j and beta lie in one original fractional ideal

\[
 L=J I_\Sigma,\quad N_EI_\Sigma=T^{-48},\quad
 z=(N_EL)^{1/(2f)}.                                        \tag{FB9}
\]

Here J is R072's whole-cross-map G ideal. Its proof uses h\_ZZ=I and a rank-one null cross image, not nonvanishing on S. At rank-one reduction square-zero nilpotence forces one cross block to vanish; inverse-adjoint duality transfers this to C at one CM partner. Rank zero gives both. The fixed (V,P) frame has determinant a G-unit, and primitive integer Bezout coefficients for u extract the scalar in the original frame. Hence there are at most2\^{}668 patterns and

\[
 N_EJ\le Q^{334},\quad
 |a(h)|\le2Q^{-335\cdot167}T^{335\cdot24}z^{335f}.           \tag{FB10}
\]

Both mixed finite factors retain their modulus bounds1. At every active T place, each y\_j lies separately in every original source and target lattice at both CM partners. Since Ny\_j=C y\_j and the target is contained in the original simple-pole lattice, primitive extraction gives FB9 for both types II and III. Enlarging the type-III integrality to the common simple-pole ideal is harmless. The original negative pivot belongs to that same ideal. This is a uniform bound for all actual labels and requires no zero-target coarea multiplier.

\subsection{4. Exact covector volume, scalar norm fibre and central congruence}\label{exact-covector-volume-scalar-norm-fibre-and-central-congruence}

The Gram matrix G=(y\_i dot y\_j) has detG=4U and G-2\^{}16 I positive definite, by the three exact leading-minor comparisons in the certificate. Thus the map from the three coordinate values to the true dual metric G\^{}-1 has norm less than2\^{}-8. The arbitrary-ideal product covering estimate is (2250/7)sqrt(3f)z in the ordinary product metric; after this contraction it is less than4sqrt(f)z. The real covolume in the dual metric is D0³ z\^{}(6f)/(4U)\^{}f, where D0=7\^{}140 29\^{}144 5\^{}84 13\^{}84. A full translated fundamental-volume bound gives

\[
 \#r\le A_3(4U)^f(M+4z)^{6f}z^{-6f},\qquad
 A_3=528^{3f}/((3f)!D_0^3)<2^{-4040}.                        \tag{FB11}
\]

The rational Gram transformation is used for the archimedean volume only; it does not identify a moving basis as integral.

For the chosen nonzero pivot, beta bar(beta)=U a\_j bar(a\_j). Write L\_j=\textbar\textbar y\_j\textbar\textbar²; all three satisfy L\_j\textless2\^{}70. The joint norm gives \textbar beta\_tau\textbar\textless=sqrt(R L\_j). K075(RN17--19)'s exact relative-ideal argument applies to (beta)L\^{}-1, even if L is nonprincipal or not conjugation-stable: each eligible integral relative norm ideal has at most d\_F(c) antecedents and each antecedent at most28 norm-one generators. DO091(TD11--12) gives C\_F\textless2\^{}42 at exponent1/12. Therefore

\[
 \#\beta\le28C_F(RL_j)^{f/12}z^{-f/6}.                    \tag{FB12}
\]

This counts all cross maps once the pivot choice is retained; the ratio eta=beta/a\_j is used only for injective reconstruction, never declared integral.

Fix C and a G pattern. Put J0 equal to the F-product of its rank-zero G primes once. At both CM partners w=N\_nn is divisible by J0. The exponent table for J gives

\[
 N_EJ\le Q^{167}N_FJ_0,\qquad N_FJ_0\le Q^{167}.           \tag{FB13}
\]

At types II and III the original full entry has pole order at most2 at both partners, by K071(NB5); thus the common containing ideal for w is J0 I\_(Sigma,F)² O\_E. Its real part is alpha/2, fixed by C. Since i is an integral unit in E and O\_E intersect iF=iO\_F, differences of any two admissible w lie in the real lattice

\[
 \Gamma=iJ_0 I_{\Sigma,F}^2,\qquad
 \operatorname{covol}(\Gamma)=\sqrt{D_F}\,N_FJ_0T^{-48},
 \quad D_F=2^{168}D_0.                                     \tag{FB14}
\]

The fibre is contained in an affine coset; its shift and any dyadic denominator in alpha/2 are not discarded. There is no factor1/2 in its difference lattice. In the ordinary real trace metric its dual is the real codifferent times the inverse ideal. DO091(TD20--22)'s actual-field transference gives cover less than750sqrt(f)(N\_FJ0 T\^{}-48)\^{}(1/f). The exact inequality750\^{}f Q\^{}167\textless T\^{}48 makes this less than sqrt(f).

At each complex place \textbar h\_nn\textbar=cosh r\_h and Re(h\_nn)=1+alpha/2\textgreater=1, so \textbar Im w\textbar\textless H/2. The actual points lie in a ball of radius(H/2)sqrt(f); their Voronoi cells lie in a ball of radius68sqrt(f). Consequently

\[
 \#w\le V_F(68)T^{48}/N_FJ_0
       \le 2^{479}Q^{167}z^{-2f},
 \quad V_F(68)=\frac{((22/7)f68^2)^{f/2}}{(f/2)!\sqrt{D_F}}
                   <2^{479}.                               \tag{FB15}
\]

This is a real f-dimensional volume, not a complex2f-dimensional count. FB13 and N\_EJ=T\^{}48 z\^{}(2f) give the last inequality. It replaces, by an explicit injective count, the old projected-column and source-fibre factors; they are not silently omitted.

\subsection{5. The actual selected trace and the complete low-rank exponent}\label{the-actual-selected-trace-and-the-complete-low-rank-exponent}

By Cauchy--Schwarz in the true covector metric, alpha\_tau=Uq\_tau\textgreater=U\textbar a\_j,tau\textbar²/L\_j. Since a\_j is a nonzero element of L,

\[
 N_F\alpha\ge(U/L_j)^f z^{2f}.
\]

The genuine selected trace DO080(RT3), used in R058(HB6) before its conservative weakening, therefore gives

\[
 |c(h)|<2^{-929f/2}(U/L_j)^{-f/2}z^{-f}.                    \tag{FB16}
\]

This retains the actual chosen trace, both compact factors, original Hilbert metric and favorable-point transport. It is not an operator norm estimate or a pointwise lower bound on individual embeddings.

FB9--10 and the exact integer comparison R\^{}f T\^{}48\textgreater25\^{}f U\^{}f Q\^{}334 imply

\[
 z\le z_{\max}=Q^{167/f}T^{-24/f}<M/5,
 \qquad M+4z<2M.                                          \tag{FB17}
\]

After multiplying FB10--12, FB15--16, the remaining z exponent is

\[
 (335-6-1/6-2-1)f=(1955/6)f>0.
\]

Thus its largest allowed value is bounded by z\_max\^{}(1955f/6). The Green factor and this cap contribute T\^{}(55s/6)Q\^{}(-55\emph{167/6); the additional central factorQ\^{}167 leaves Q\^{}(-49}167/6).

The complete fixed enumeration is the finite Euler factor2, common G patterns2\^{}668, support choices2\^{}168, types2\^{}24, and pivot3\textless2². All paired II/III weights have already cancelled their available ordered-label counts, and every count above is uniform over those labels. There is no additional support, type, central, isometry, alpha or norm-range union. Using R\textless2\^{}8, L\_j\textless2\^{}70, the exact directed values L\_T=1019333/1024\textgreater log2T, L\_Q=287843/4096\textless log2Q and L\_U=601356/4096\textless log2U, the complete exponent is

\[
\begin{split}
 C={}&1+668+168+24+2+479+5+42-4040+8f-929f/2\\
    &+(37f/12)8-(5f/2)L_U+(7f/12)70
       =-33280379/256,\\
 E_{\rm low}={}&C-(49\cdot167/6)L_Q+(55\cdot24/6)L_T\\
       =-168257413/24576<-6846.                            \tag{FB18}
\end{split}
\]

Here8f=2f+6f comes from detG=4U and the enlarged radius2M. The R exponent is3f+f/12, the U exponent is-2f-f/2, and the pivot length exponent isf/12+f/2. The root-unit factor28 is bounded by2\^{}5 and the actual relative divisor constant by2\^{}42. These are the full factors, including the central lattice normalization.

The high and low integer alternatives exhaust FB1. Type-I low-rank terms are empty, while every other low-rank profile has exponent FB18. Hence the complete total is below2\^{}-29143+2\^{}-6846\textless2\^{}-6845, proving FB2.

\subsection{6. Accounting, geometric placement and remaining questions}\label{accounting-geometric-placement-and-remaining-questions}

At checked integration append6845 exactly once to the actual ordered all-top error array, retaining every earlier charge and the unchanged uniform array. The arithmetic certificate checks the complete sums for its explicitly recorded predecessor; integration separately verifies the actual chronological predecessor. No complete support floor is raised without a separate exhaustion of all other structural classes.

The new result counts the dual of the original positive three-module with its exact Gram determinant, an exact scalar relative norm fibre over its fractional ideal, and an affine real central-entry lattice with its full discriminant and codifferent. These live on the original Hermitian modules over Spec(O\_E), with both CM partners and the polarization retained. Contextual module/exterior anchors are SGA3 I4.7.3/4.7.3.1 and SGA6 VI Theorem3.3. Picard and biextension anchors remain FGA232 Theorem3.1 pp148--149; FGA236 Corollaries1.7/2.3 pp229/232; SGA7 I VII2.7/2.9 pp169--174 and VIII7.4(b) pp310--311. Crystals and Cartier duality retain cote4 scan003 §§8.1.1--8.1.2, scans005--006 §§8.4.1--8.4.3, and cote54 scans005--007.

No new lambda-ring equality, K-theory comparison, l-adic sheaf or weight theorem, motivic projector, CM-motive identification or L-function conclusion follows. DP\_MULT Q3e and THETA3 G1/A7 gain only the stated smaller finite/real remainder. Actual mixed-CM source construction, selected arithmetic-image content, prescribed Bessel/Mram and specific032 pairings, every actual DP44 term, CX\_GRH CX-ALG-024, and the global GCT\_PLETH T-052 comparison remain separately open. The companion is an exact arithmetic diagnostic, not an enumeration of the arithmetic kernel. No Lean build or archive audit is performed.

\clearpage

\section{AM-R-078 --- complete all-top support twenty-four exhaustion}\label{am-r-078-complete-all-top-support-twenty-four-exhaustion}

\textbf{Canonical checked integration --- 2026-11-22T23:32:48.711360+00:00.} Root and independent whole-final-proof/script/scope readings and separate738-assertion replays passed. The DO092 locator was corrected toSN1--24 before the final replays and whole rereads; no mathematical correction. Every actual all-top support24 term is exhausted by the complete original F/IC8 partition at compressed ranks1..33, every fixed-S nonunipotent/rank-one/rank-two class, and every scalar-cross global-rank/type/coordinate alternative. R075 covers all nonidentity positive compressions; DO090/DO092/R077 exhaust fixed S; R070/DO093/R071 reduce cross terms to the precise DO095 rank2/3 alternatives, exhausted by R076,DO096,DO097,DO098. The complete all-top support floor is25, uniform floor11. All83/31 chronological errors are unchanged and both full sums remain below2\^{}-230. No new charge. The actual remaining signed gate is J25.2; prescribed global/source/selected-image/specific032/DP44 and RH/global oscillator gates remain open.

Author: \nolinkurl{/root}, GPT-6 Astra Pro; inherited effort not externally exposed; ALP\_MACHINE. 22 November2026 UTC. Complete assembly proof; whole proof, literal dependency coverage and exact chronological accounting are reviewed separately before promotion.

\subsection{1. Statement on the unchanged actual kernel}\label{statement-on-the-unchanged-actual-kernel-1}

Keep the J56 all-top K044 endpoint, its charged H=133 cutoff, f=a=168, n=169, and the original positive S of dimension164 with orthogonal complement W0 of signature(4,1). Keep every paired CM source/target pencil, finite projector, actual selected Hilbert metric, favorable-point transport, class mass, phase and normalized summand w(h). R074(J24.1--4) gives the complete preceding support floor24. All subsequent containing-set removals act on this same kernel.

\textbf{Theorem.} There is an actual remaining subset R\_top\^{}(25) with

\[
 \mathcal E_{\mathrm{top}}=1+\operatorname{Re}
       \sum_{h\in R_{\mathrm{top}}^{(25)}}w(h)+\epsilon,
 \qquad |\epsilon|<2^{-230},\qquad s(h)\ge25.                \tag{J25.1}
\]

This assembly adds no error charge. It retains all83 chronological all-top entries after R077, and all31 unchanged uniform entries. The independent uniform H=8192 support floor remains11. The sufficient inequality

\[
 \operatorname{Re}\sum_{h\in R_{\mathrm{top}}^{(25)}}w(h)
                         >-1+2^{-230}                       \tag{J25.2}
\]

remains open. We prove J25.1 by exhausting every original structural class at support24. We do not assert a complete support25 removal, positive signed energy, an arithmetic source or any DP44 cost comparison.

\subsection{2. Every nonidentity positive compression}\label{every-nonidentity-positive-compression-1}

Put A=h\_SS and r=rank\_E(A-I). The retained K068 residual has r\textless=33. For A!=I, K069(IC1--8) gives the full fixed class F and exceptional isotropic-cross class IC8. This partition has no global-unipotence premise. The following complete containing-set ranges cover s=24:

\begin{longtable}[]{@{}
  >{\raggedright\arraybackslash}p{(\columnwidth - 2\tabcolsep) * \real{0.5000}}
  >{\raggedright\arraybackslash}p{(\columnwidth - 2\tabcolsep) * \real{0.5000}}@{}}
\toprule\noalign{}
\begin{minipage}[b]{\linewidth}\raggedright
Original class
\end{minipage} & \begin{minipage}[b]{\linewidth}\raggedright
Full coverage at support24
\end{minipage} \\
\midrule\noalign{}
\endhead
\bottomrule\noalign{}
\endlastfoot
F, r=1 & DO086(FR1--10), every support18..94. \\
F, 2\textless=r\textless=33 & R068(FA1--9): rank2 through110, rank3 through154, every rank4..33 through168. \\
IC8, r=1 or2 & R069(EM1--21), every support23..28 with all original column alternatives. \\
IC8, 3\textless=r\textless=31 & R075(NR1--16), combined with its contiguous DO085/DO064 predecessor ranges. The smallest complete ceiling is28 at rank3. \\
IC8, r=32 or33 & DO064(PM22--23), complete through168. \\
\end{longtable}

R075's new theorem includes every newly needed support between each old complete ceiling and its new ceiling; it does not merely check a terminal numerical value. Thus no missing support is inferred from a sparse table. All33 positive compressed ranks and both original classes are covered. These bounds retain their original finite-rank patterns, target columns, source fibres and actual paired weights. In particular every remaining support24 term has A=I.

\subsection{3. Identity compression: complete fixed-block and cross alternatives}\label{identity-compression-complete-fixed-block-and-cross-alternatives}

For A=I, write h={[}{[}I,B{]},{[}C,D{]}{]} in S perpendicular W0. The original unitary identities R059(GB1--3) give C\emph{J0 C=0 and B=-C}J0 D. Hence C=0 implies B=0, yielding h=I\_S direct-sum h0. If C is nonzero, its image is a null line in W0. This is exhaustive and permits globally nonunipotent matrices.

\subsubsection{3.1 Fixed S}\label{fixed-s-1}

\begin{longtable}[]{@{}
  >{\raggedright\arraybackslash}p{(\columnwidth - 2\tabcolsep) * \real{0.5000}}
  >{\raggedright\arraybackslash}p{(\columnwidth - 2\tabcolsep) * \real{0.5000}}@{}}
\toprule\noalign{}
\begin{minipage}[b]{\linewidth}\raggedright
Full fixed-S sector at support24
\end{minipage} & \begin{minipage}[b]{\linewidth}\raggedright
Complete original removal
\end{minipage} \\
\midrule\noalign{}
\endhead
\bottomrule\noalign{}
\endlastfoot
Globally nonunipotent, with any global rank and incidence profile & DO090(NE1--6), every support23..26; retained charge65341. \\
Globally unipotent, rank(h-I)=1 & DO092(SN1--24), complete support24 with all marked/type/pair cells; retained charge516. \\
Globally unipotent, rank(h-I)=2 & R077(FB1--18), complete support24 with every original incidence profile; retained charge6845. \\
rank(h-I)=0 & h=I has support0, so cannot occur at support24. \\
\end{longtable}

The signature-one unipotent classification DO060(IB21--22) gives only ranks1 and2 for a nonidentity unipotent. Rank1 has square zero; rank2 has N³=0 and rank N²=1. These are exactly the hypotheses of the cited removals. R077 handles both D\_G\textless=32 and D\_G\textgreater=33; its low-rank type-I exclusion is only one part of the proof, while its entire theorem covers all types. DO090 is not restricted to a single characteristic degree or to a plane. Consequently no fixed-block class remains at s=24.

\subsubsection{3.2 Nonzero scalar cross}\label{nonzero-scalar-cross-1}

R070(SC1--22) removes every C-transverse case and every plane-C case with at least one nonzero original short2x2 compression at s=24, as part of its already charged support23/24 union. DO093(IV1--22) removes the inverse-transverse cases at the same supports with their unchanged literal inversion normalization. R071(GR1--12) then forces the remaining transverse scalar to1. Thus each still remaining support24 cross term has

\[
 h=\begin{pmatrix}I_Z&B_Z\\C_Z&D_P\end{pmatrix},\quad
 Z=S\perp u^\perp,\quad\dim Z=167,\quad
 P=Eu\perp Ee_n,\quad C_Z|_S\ne0,
 \qquad \operatorname{rank}_E(h-I)\le3.                     \tag{J25.3}
\]

It fixes a positive166-dimensional E-space. These are rational fixed spaces in the original modules; no globally integral moving frame or new lattice is introduced.

DO095(RA1--10,RA21--26) exhausts J25.3: the global rank is3 when B\_Z C\_Z is nonzero and2 when it is zero. A nonzero scalar cross cannot have global rank1 here. The remaining alternatives are covered as follows:

\begin{longtable}[]{@{}
  >{\raggedright\arraybackslash}p{(\columnwidth - 2\tabcolsep) * \real{0.5000}}
  >{\raggedright\arraybackslash}p{(\columnwidth - 2\tabcolsep) * \real{0.5000}}@{}}
\toprule\noalign{}
\begin{minipage}[b]{\linewidth}\raggedright
Actual support24 cross sector
\end{minipage} & \begin{minipage}[b]{\linewidth}\raggedright
Complete containing-set removal
\end{minipage} \\
\midrule\noalign{}
\endhead
\bottomrule\noalign{}
\endlastfoot
Global rank3, with any nonzero original S support size1\textless=k\textless=164 & R076(CV1--18), the entire sector, retained charge5710. \\
Global rank2 and globally nonunipotent & DO096(RP1--21), the entire sector, retained charge2962487. \\
Global rank2 and globally unipotent, with at least one type-I or type-III place & DO097(UC1--24), every such profile and every1\textless=k\textless=164, retained charge71310. \\
Global rank2 and globally unipotent, every active place typeII & DO098(CF1--32), the complete all-II sector and every1\textless=k\textless=164, retained charge2879. \\
\end{longtable}

The last two rows are disjoint and exhaustive because the original incidence types I,II,III are exhaustive. DO095 supplies the precise globally unipotent block and nilpotence premises; it does not replace the stronger negative-column hypotheses of another theorem by global unipotence alone. DO097 and DO098 both work on the actual identity-positive-compression residual J25.3. R076 has no sparsity restriction and no global-unipotence premise. DO096 includes every nonunipotent reciprocal-pole sign case, including its one inversion factor and possible active eigenvalue-unit place. There is no additional k, type, inverse, central, finite-rank or norm-range class to count here.

Every cited bound is already a complete bound for its stated containing set. Applying it to the intersection with the current remainder is allowed. The earlier sparse R073 charge is deliberately retained, even though the complete R076 theorem now covers its support24 rank-three subset. All fixed-block and cross classes at support24 are therefore exhausted. Section2 exhausts A!=I, and R074 already removed every smaller support.

\subsection{4. Exact ordered errors and the sufficient gate}\label{exact-ordered-errors-and-the-sufficient-gate}

The actual predecessor is checks/r077\_joint\_remainder\_error.json. Relative to R074's77-entry array, the complete subsequent ordered suffix is

\[
       2962487,\ 8866,\ 71310,\ 5710,\ 2879,\ 6845.        \tag{J25.4}
\]

These are DO096, R075, DO097, R076, DO098 and R077 respectively. No term is appended by this assembly and no old overlapping charge is removed. The actual arrays have83 all-top entries and31 uniform entries. There is exactly one all-top exponent231 and every other all-top exponent is at least358. Hence

\[
 \sum_{e\in\mathcal L_{\mathrm{top}}}2^{-e}
       \le2^{-231}+82\,2^{-358}<2^{-230}.
\]

The full unchanged uniform array also has exact sum below2\^{}-230. The companion verifies the actual ordered prefix and suffix, all required inclusive rank ceilings, all scalar-cross coordinate/type alternatives, and both complete integer dyadic sums. These checks certify endpoint coverage and accounting; they are not an enumeration of arithmetic arrows or a replacement for the geometric partition proofs.

Chronologically apply each new removal only to its intersection with the remainder left by earlier removals. The actual removed subsets are then disjoint even when the proved containing sets overlap. Their complete error bounds remain charged once each. The original normalized energy identity, the triangle inequality and the exhaustive partition above prove J25.1. Thus the complete all-top support floor increases from24 to25, strictly reducing the Q1/Q2 finite/real sufficient gate to J25.2.

J25.2 itself, the prescribed global Bessel/Mram value, a fully admissible mixed-CM source, selected arithmetic-image content, actual Read/STiesR and the specific032 cycle matching/value, and every changed-family DP44 term remain open. No shorter valid lift chain or universal strict negative DP44 difference has been proved. Q3 has no Kummer-Hecke RH conclusion. Q4 retains its separate global arithmetic charge-comparison gate.

\subsection{5. Algebraic geometry and exact carried locators}\label{algebraic-geometry-and-exact-carried-locators}

This is an exhaustion of the original determinantal, fixed-submodule and isotropic-cross strata of integral Hermitian correspondences over Spec(O\_E), together with their actual selected complex realizations. Both CM partners, fractional ideals, coefficient modules, polarization, theta line, Picard classes and biextensions remain unchanged. Rational reconstruction in the cited counts supplies no inverse on the arithmetic theta image.

The exact contextual module/determinantal anchors are SGA3, ExposeI,4.7.3/4.7.3.1, and SGA6, ExposeVI, Theorem3.3. The Picard anchors are FGA, Bourbaki232, Theorem3.1, pp148--149, and Bourbaki236, Corollaries1.7/2.3, pp229/232. Biextensions retain SGA7 I, ExposeVII,2.7/2.9, pp169--174, and ExposeVIII,7.4(b), pp310--311. Formal groups, Dieudonne crystals and Cartier duality retain cote4 scan003, §§8.1.1--8.1.2, scans005--006, §§8.4.1--8.4.3, and cote54 scans005--007. These are contextual placements, not new comparison theorems.

No new scheme identification, lambda-ring/K-theory equality, l-adic sheaf or weight result, CM-motive comparison or motivic L-function theorem follows. DP\_MULT Q3e and THETA3 G1/A7 gain exactly J25.1--2; CX\_GRH CX-ALG-024 and GCT\_PLETH T-052 keep their distinct arithmetic gates. Settled dependencies retain their exact source locators and whole-proof checks. The present proof checks their complete partition coverage and chronological sum. No Lean build or archive audit is performed.

\clearpage

\section{AM-DO-099 --- joint rows and the next scalar-cross structural interval}\label{am-do-099-joint-rows-and-the-next-scalar-cross-structural-interval}

\textbf{Canonical checked integration --- 2026-11-22T23:40:44.279716+00:00.} Root and independent whole-proof/script/interface readings and separate15,538-assertion replays passed without correction. Historical pending wording below is preserved. Both complete C-transverse directions, with every incidence profile and no global-unipotence restriction, are removed at supports25..70. At support25 the plane-C/nonzero-original-short-compression sector is also removed. The full joint row, both original columns and complete source fibre give transverse maximum -49537435/4096; the complete plane row, exact relative norm fibre and actual selected trace give -435793141/4096. Their entire union costs one12079 charge. The remaining scalar-cross support25 matrices have both cross blocks in the original packet plane, all four short compressions zero, positive166-dimensional fixed space and global rank at most3 by the support-free norm-one lemma. At26..70 only the two plane conditions are concluded. Actual84/31 ordered errors remain below2\^{}-230 with complete floors25/11. No signed/global/source/selected-image/specific032/DP44 or RH conclusion.

Author: \nolinkurl{/root/definite_oscillator}, ALP\_MACHINE. 22 November2026 UTC.

Complete controlled pending proof. Exact author replay, whole saved-file verification and independent review are recorded separately before any canonical integration.

\subsection{1. Unchanged actual sector and conclusions}\label{unchanged-actual-sector-and-conclusions}

Retain the original J56 all-top normalized kernel, charged H=133 cutoff, f=168, n=169, original Hermitian decomposition S perpendicular W0, dim S=164 and signature(W0)=(4,1). Every source/target pencil, CM partner, zero label, finite projector, selected Hilbert metric, actual packet representative, central phase, Haar measure and class mass is unchanged. Work in the actual K068 remainder, so all334 unmixed G reductions are unipotent. Global unipotence is not assumed. Write

\[
 h=\begin{pmatrix}I_{164}&B\\C&D\end{pmatrix},\qquad C\ne0.
                                                               \tag{JR1}
\]

Use the original primitive packet vector u and the three original transverse inputs of R059(GB4), with u*u=U=77Q²−14 and Q=1427908193729121522601. Put Y=span\_E(y1,y2). Call C transverse when at least one row (\textless y\_j,Ce\_i\textgreater)\_i, j=1,2,3, is nonzero. Otherwise C is in the packet plane: Ce\_i=(u xi\_i,beta\_i). Define

\[
          \Delta_{ab}=\langle y_a,(h-I)y_b\rangle,
                                      1\le a,b\le2.           \tag{JR2}
\]

\textbf{Theorem.} The union of the following actual containing sectors has total absolute normalized contribution less than2\^{}-12079:

\begin{enumerate}
\def\labelenumi{\arabic{enumi}.}
\tightlist
\item
  At every support25≤s≤70, C is transverse or the lower cross block of h\^{}-1 is transverse.
\item
  At support25, C is in the packet plane and at least one Delta\_ab is nonzero.
\end{enumerate}

Both globally unipotent and nonunipotent matrices are included. All incidence profiles are included. One new all-top charge12079 suffices for the union.

Consequently the remaining scalar-cross matrices at support25 may be required to have both cross blocks in the packet plane and all four short compressions zero. The actual source-independent algebraic part of R071 then gives

\[
 h_{S\oplus u^\perp,S\oplus u^\perp}=I,
 \quad\dim_E\bigl(\ker(h-I)\cap(S\oplus u^\perp)\bigr)=166,
 \quad\operatorname{rank}_E(h-I)\le3.                         \tag{JR3}
\]

At supports26..70 only the two plane conditions are obtained here; the short compressions are not claimed zero there. There is no assertion that individual y\_j are fixed, no moving integral frame, and no complete support-floor or signed nonvanishing conclusion.

\subsection{2. A joint norm for the original integral transverse row}\label{a-joint-norm-for-the-original-integral-transverse-row}

R059(GB3--6) proves rank\_E C=1 and the following exact source/target facts. For each j the row

\[
                         \delta_i=\langle y_j,Ce_i\rangle
                                                               \tag{JR4}
\]

is integral in O\_E at every finite place. This follows because he\_i and y\_j belong to the same actual target selfdual lattice, with separate membership in the source and target at each CM partner. It does not require h to stabilize one common lattice. Choose one nonzero row, retaining at most three choices, and a nonzero pivot delta\_i0, retaining at most164 choices. The original pivot column z0=Ce\_i0 reconstructs C by Ce\_i=z0 delta\_i/delta\_i0. The quotient is used over E only.

The new norm observation is joint in all164 original coordinates. Since C is an original Euclidean submatrix of h, \textbar\textbar C\_tau\textbar\textbar op≤H. The y\_j have zero negative coordinate, so their J0-row Euclidean norm is exactly \textbar\textbar y\_j\textbar\textbar. Thus

\[
 \| (\delta_{i,\tau})_{i=1}^{164}\|_2
                 \le H\|y_j\|<2^{43}.                       \tag{JR5}
\]

The final strict inequality is the same exact original comparison as R059(GB11--12): the squared lengths are978665066829,835948910828,4030172803672366725045 and H² times their maximum is below2\^{}86. Equation JR5 is stronger than merely bounding each coordinate by2\^{}43; it follows directly from the operator norm and does not replace the selected real coefficient.

The usual whole-C orientation at finite ranks0,1,at-least2 supplies an ideal J, at most2\^{}668 patterns, and

\[
 \delta\in J^{164},\quad z=(N_EJ)^{1/(2f)},\quad
        |a(h)|\le2Q^{-335\cdot167}z^{335f}.                  \tag{JR6}
\]

These are the unchanged R059(GB9--10) statements, including both mixed factors. If the chosen row has k nonzero coordinates, 1≤k≤164, its actual ideal minimum gives kz²\textless2\^{}86. Apply the general R072(PJ10--12) spherical-cap/Gaussian row lemma to JR5. Its complete count is

\[
 {164\choose k}47^k2^{2fk}(f!)^k
      \left(\frac{(87/32)2^{86}}{fkz^2}\right)^{fk}.          \tag{JR7}
\]

After multiplication by JR6, the z exponent is(335−2k)f, strictly positive even at164. Therefore the single endpoint z≤2\^{}43/sqrt(k) gives the weighted row factor

\[
 2^{43\cdot335f}F_k,\qquad
 F_k=\frac{\binom{164}{k}47^k2^{2fk}(f!)^k87^{fk}}
              {32^{fk}f^{fk}k^{28140}},\quad28140=335f/2.   \tag{JR8}
\]

The binomial chooses the original nonzero positions. No unweighted k-independent product count is substituted for this joint norm.

\subsection{3. Complete transverse interval, including inverse arrows}\label{complete-transverse-interval-including-inverse-arrows}

All remaining R059 factors stay unchanged: the pivot column z0 and the W0 part v=D e\_n of the negative column each have their original complete target-overlattice count2\^{}14861 T\^{}(10s). The data delta,z0 reconstruct C, and unitarity gives (he\_n)\_S=−C*J0v; the entire negative column is therefore known. The full original source fibre G\_n\textless2\^{}5022 is retained. Neither column is placed in a substituted diagonal ideal, and no improved rank-five fibre or real selected-trace estimate is used for this transverse branch.

The complete fixed exponent is exactly

\[
        1+668+168+8+2+5022+2\cdot14861=35591.                 \tag{JR9}
\]

Here the factors are the Green correction, common G patterns, support choices, pivot choices, three transverse witnesses, full source fibre, and two original columns. The original whole pair/weight sum is P\_T\^{}s\textless2\^{}(997s), including every type and admitted zero label. The actual selected real coefficient is bounded by one.

Put L\_T=1019333/1024\textgreater log2 T and L\_Q=287843/4096\textless log2 Q. For each s,k the complete logarithmic upper bound is

\[
 E_A(s,k)=35591+43\cdot335f+(20L_T+997)s
                            -335\cdot167L_Q+e_k,             \tag{JR10}
\]

where e\_k is the least strict integral binary ceiling of the exact rational F\_k in JR8. The certificate verifies all164 values and their maximum e\_1=354. The support slope is positive, and

\[
        \max_{25\le s\le70,\,1\le k\le164}E_A(s,k)
                    =-49537435/4096<-12094.                 \tag{JR11}
\]

The estimate at71 is positive and is not promoted as a count or contribution claim.

DO093(IV4--15)'s inversion identities are pointwise source/operator statements, not consequences of its earlier numerical support23/24 evaluation. They preserve the exact normalized absolute summand, actual source/target incidence with labels swapped, support, H cutoff, class/Haar/central normalization, and the finite-unipotence premise. Thus at each present support the sector with transverse inverse C injects into the complete transverse-C sector with the same bound. The previously cut remainder need not be inverse-closed, since its image is enlarged to that complete sector. There is no factor two per matrix or per CM label; the factor two below counts just the two possible containing sectors.

There are46 supports and164 coordinate sizes. Consequently the union of the two transverse directions has contribution less than

\[
           2\cdot46\cdot164\,2^{-12094}<2^{-12080}.           \tag{JR12}
\]

\subsection{4. The original plane/nonzero-compression restoration}\label{the-original-planenonzero-compression-restoration}

Now take s=25, C in the packet plane and some Delta\_ab nonzero. R070(SC4--10) applies without a global-unipotence hypothesis and without a support-size restriction in its algebraic proof. The original y1,y2 pairings are integral by the two-lattice argument, and every conjugate of every Delta\_ab is below2\^{}47. The finite lemma uses the actual paired inverse-adjoint relation and the two identity S blocks. It supplies the following whole-C ideal J\_C and compression ideal J\_Delta; ``paired'' has E-norm q\_v² and ``one'' has E-norm q\_v.

\begin{longtable}[]{@{}rlll@{}}
\toprule\noalign{}
Finite rank & Zero lower C blocks at the two components & J\_C & J\_Delta \\
\midrule\noalign{}
\endhead
\bottomrule\noalign{}
\endlastfoot
0 & Both & Paired & Paired \\
1 & Exactly one & That component & Paired \\
1 & Both & Paired & None \\
2 & Neither & None & Paired \\
2 & Exactly one & That component & None \\
2 & Both & Paired & None \\
At least3 & Any & None & None \\
\end{longtable}

This is the already proved complete restoration table, not a new finite-rank assertion. In particular it does not say that the short vectors are fixed. The single assigned J\_Delta divides all four compressions wherever assigned, so any nonzero global entry bounds its norm. No independent witness or compression-pattern count is required. The actual Green inequality is therefore

\[
 |a(h)|\le2Q^{-501\cdot167}(N_EJ_C)^{167}
                              2^{47\cdot167f}.               \tag{JR13}
\]

All other finite-rank types are included by the existing envelope. There are at most2\^{}668 J\_C patterns, and each assigned norm exponent is at most2, so N\_EJ\_C≤Q\^{}334.

Let L=J\_C I\_Sigma with N\_E I\_Sigma=T\^{}-50 and z=(N\_E L)\^{}(1/(2f)). The original entries xi\_i and the negative pivot beta belong to L, using exactly R070(SC12) at both CM partners and the primitive u Bezout extraction. The finite ideal concerns the whole C; beta needs no second orientation.

The negative row of an original U(168,1) matrix has positive-part norm sinh(log\textbar\textbar h\textbar\textbar)\textless H/2. Isotropy of C gives beta\_i=eta xi\_i with eta bar(eta)=U. Hence the full original scalar row has

\[
             \|\xi_\tau\|<M:=H/(2\sqrt U)<2^{-67},
                           |\beta_\tau|<b:=H/2.             \tag{JR14}
\]

This is the same signature-one row identity used in R072(PJ7); it requires no global nilpotence. The exact inequality for M is H²2\^{}132\textless U. The common-pattern cap now gives

\[
       z\le Q^{167/f}T^{-25/f}<2^{-78},                       \tag{JR15}
\]

since Q\^{}167 2\^{}(78f)\textless T\^{}25 for the actual host.

\subsection{5. One full scalar-row volume and the exact pivot norm fibre}\label{one-full-scalar-row-volume-and-the-exact-pivot-norm-fibre}

Use the genuine complex half-trace metric and the entire lattice L\^{}164. Its covolume is D\_0\^{}164 z\^{}(328f). The arbitrary-ideal covering estimate DO070 gives product covering radius less than

\[
       (2250/7)\sqrt{164f}\,z<4117\sqrt f\,z.                 \tag{JR16}
\]

Consequently the full row count, including all coordinate supports at once, is at most

\[
 A_{164}(M+4117z)^{328f}z^{-328f},\qquad
 A_{164}=\frac{528^{164f}}{(164f)!D_0^{164}}.                 \tag{JR17}
\]

This is a true328f-dimensional ball in the original scalar coordinates. No separate k or position factor is needed in this branch. Counting the zero row too only enlarges the upper bound; reconstruction and the chosen nonzero pivot are asserted only for actual C≠0.

For each fixed nonzero pivot xi\_j, the actual beta satisfies beta bar(beta)=U xi\_j bar(xi\_j) and belongs to L. Apply exactly the integral relative norm ideal construction DO098(CF20--21), with right-side archimedean bound b² instead of R. DO092(SN10--11)'s split-only constant is below2\^{}244 and there are at most28 norm-one generators for each eligible ideal. Thus

\[
                          \#\beta<2^{249}b^{f/8}z^{-f/8}.    \tag{JR18}
\]

The ratio beta/xi\_j reconstructs eta over E, and then the entire C. This does not declare that ratio integral. The norm ideal is integral only when an actual beta exists; nonprincipal ideals give no solution. All class and generator multiplicities are therefore accounted for.

Retain R070(SC14--16) unchanged for the remaining data:

\[
 |c(h)|<2^{-929f/2}U^{-f/2}z^{-f},\quad
 \#\{D e_n\}<2^{5757}T^{10s},\quad G_5<2^{107}.              \tag{JR19}
\]

These are the genuine selected trace with both compact factors, the full original five-complex-coordinate target column, and the actual rank-five source-lattice fibre. The column bound is uniform in every admitted target label, including zero; no real-coordinate coarea assertion enters. Fixing C and D e\_n reconstructs the complete negative column before applying G\_5.

Combining JR13 and JR17--19 leaves the positive z-dependent factor

\[
             z^{819}(M+4117z)^{328f},\quad
 819=334f-328f-f/8-f=39f/8.                                  \tag{JR20}
\]

It increases with z\textgreater0. At JR14--15, M+4117z is below6165/2\^{}78. The exact rational volume comparison is

\[
        A_{164}(6165/2^{78})^{328f}<2^{-3983962}.              \tag{JR21}
\]

There is one endpoint and no additional coordinate-support union.

The complete fixed exponent is 1+668+168+8+107+5757+249=6958, retaining the Green correction, patterns, supports, pivot, source fibre, full column and relative norm fibre. Retain the full original P\_T\^{}25\textless2\^{}(997*25), with no separate type enumeration. Put L\_H=28900/4096\textgreater log2 H and L\_U=601356/4096\textless log2 U. The resulting full exponent is

\[
\begin{split}
 E_P={}&6958-929f/2+47\cdot167f+(f/8)(L_H-1)
 -(f/2)L_U-501\cdot167L_Q\\
 &+(344L_T+997)25-3983962-819\cdot78\\
      ={}&-435793141/4096<-106394.                            \tag{JR22}
\end{split}
\]

The T power344s includes334s from the whole-C ideal numerator and10s from the original target column. The relative norm fibre and real trace are already present in the exact z exponent819. No extra degree, phase, Haar or class factor has been removed.

Combining JR12 and JR22 gives

\[
                    2^{-12080}+2^{-106394}<2^{-12079},         \tag{JR23}
\]

the single claimed charge for the full union.

\subsection{6. The support25 structural consequence has no hidden support bound}\label{the-support25-structural-consequence-has-no-hidden-support-bound}

Outside the bounded union at support25, C and the inverse lower cross block are in the original packet plane, and all four Delta\_ab vanish. DO093(IV19--22)'s proof uses only these equations, unitarity, the actual positive test space S perpendicular Y and C≠0. It gives N(Y) contained in CS and a positive165-space fixed pointwise. Its earlier support23/24 range arose from the numerical removals, not from that linear-algebra statement.

The third integral transverse vector z0 of DO094(FT8--11) still belongs separately to every original source and target pencil at this support. Thus delta\_z=\textless z0,(h-I)z0\textgreater{} is integral, by the same arbitrary-arrow paired selfdual argument. The global plane/compression algebra R071(GR6--8) supplies a norm-one scalar alpha and

\[
 \delta_z=(U/M_0)(\alpha-1),\quad \alpha\bar\alpha=1,
 \quad M_0=155821585821105090045582,\quad\gcd(U,M_0)=9.
                                                               \tag{JR24}
\]

The exact inequality M\_0²\textgreater36U is unchanged. The general integral norm-one lemma R071(GR1--5), whose statement contains no support variable, forces alpha=1. Therefore N(S perpendicular u-perp)=CS, a line, and JR3 follows by rank-nullity. The actual compression h\_ZZ is the identity because the image lies in the orthogonal packet plane. No G-rank hypothesis, global nilpotence, finite positive metric, integral moving frame or pointwise fixation of the individual transverse vectors is inferred.

This explicitly verifies the complete geometric premises at support25. It is not an unproved extension of a prior numerical support range. At26..70 the present theorem only removes transverse cross lines; JR24 is not applied without the additional four-zero condition.

\subsection{7. Chronology, source anchors and limits}\label{chronology-source-anchors-and-limits}

The actual predecessor is \nolinkurl{checks/r078_joint_remainder_error.json}, with83 all-top and31 uniform entries and complete floors25/11. R078's support24 exhaustion adds no numerical charge after R077. Append12079 once, obtaining84/31 and retaining every prior charge. Remove only intersection with the current remainder; the containing-set estimates above need no invariance of older cuts. The companion records the exact predecessor hash and both complete dyadic sums. This result by itself makes no complete support-floor change.

Directly reread interfaces are the entire R059(GB3--21), R070(SC1--25), DO093(IV4--22), R071(GR1--12), and the general R072(PJ10--12) joint-row lemma. DO070 supplies only the arbitrary fractional-ideal metric cover; DO092(SN10--11), K075(RN17--19) and DO098(CF20--21) supply the actual relative norm-ideal count and its generator bound. The new steps are JR5, its complete interval count JR8--12, the original plane-row volume/relative norm composition JR14--22, and the explicit support25 verification of the carried geometric premises. No settled theorem is re-proved and no claim rests on floating-point diagnostics.

The original modules over Spec(O\_E), fixed CM pairing, coefficient bundle, Picard data, polarizations and biextensions are unchanged. Carried anchors remain SGA3 I4.7.3/4.7.3.1; SGA6 VI Theorem3.3; FGA232 Theorem3.1 pp148--149; FGA236 Corollaries1.7/2.3 pp229/232; SGA7 I VII2.7/2.9 pp169--174 and VIII7.4(b) pp310--311; cote4 scan003 §§8.1.1--8.1.2, scans005--006 §§8.4.1--8.4.3 and cote54 scans005--007. No new K-theory, l-adic, motivic or source-constructor assertion follows. The signed inequality, prescribed global periods, selected arithmetic content and DP44 comparison remain open.

\clearpage

\section{AM-R-079 --- all-profile central cancellation and the short/last pivot split}\label{am-r-079-all-profile-central-cancellation-and-the-shortlast-pivot-split}

\textbf{Canonical checked integration --- 2026-11-22T23:48:08.989248+00:00.} Root author and independent whole-proof/script/interface readings and separate2,983-assertion replays passed without correction. Historical pending wording below is preserved. The entire original fixed-S globally rank-two-unipotent sector at support25 is removed for every actual incidence profile. Also removed at26 are D\_G\textgreater=33 and the D\_G\textless=32 covector sector r(y1)=r(y2)=0. The original central affine lattice has poles I1/II2/III1; its T\^{}(u-s) factor cancels actual pair weights to T\^{}(-t)\textless=1. The true three-covector short pivots and one-coordinate Schur-complement last pivot retain all source/norm/trace/metric factors. Complete low exponent at25 is -27553337/24576; high endpoints are -77307919/4096 and -35243075/4096; the last-only26 endpoint is -78939175/8192. One1120 charge covers the complete union, giving85/31 errors below2\^{}-230 with unchanged complete floors25/11. Short-pivot low-rank26, fixed-S rank-one25, remaining scalar-cross25, signed/global/source/selected-image/specific032/DP44 and RH gates remain open.

Author: /root, GPT-6 Astra Pro; inherited effort not externally exposed; ALP\_MACHINE. 22 November2026 UTC. Controlled new proof. Exact computation, author whole reading and independent whole review are recorded separately before canonical integration.

\subsection{1. Actual sector and precise theorem}\label{actual-sector-and-precise-theorem}

Retain the unchanged J56 all-top normalized kernel at the charged H=133 cutoff, f=168, n=169, its original Hermitian decomposition S perpendicular W0 with dim S=164 and signature W0=(4,1), all paired CM components, actual admitted source/target pencil labels, finite projectors, selected Hilbert metric, favorable-point transport, central representative, phases, class masses and Haar measures. Work in the original K068 finite-unipotent remainder. The present globally unipotent fixed-block sector is

\[
 h=I_S\oplus h_0=I+N,\quad N^3=0,\quad
 \operatorname{rank}_E N=2,\quad\operatorname{rank}_E N^2=1.
                                                        \tag{PC1}
\]

Global unipotence and fixation of S are explicit hypotheses. Write i,u,t for the numbers of actual types I,II,III, so s=i+u+t. Let D\_G be the sum of the334 unmixed relative G ranks, as in R058(HB3). In its low-rank plane case define the covector r and its original values a1,a2,a3 below.

\textbf{Theorem.} The union of the following actual sectors has total absolute normalized contribution less than2\^{}-1120:

\begin{itemize}
\tightlist
\item
  every PC1 term of support25, with every incidence profile;
\item
  every PC1 term of support26 with D\_G\textgreater=33;
\item
  every PC1 term of support26 with D\_G\textless=32 and a1=a2=0.
\end{itemize}

The last condition has its intrinsic meaning r(y1)=r(y2)=0 in the fixed original positive transverse space. Since r is nonzero, its third value is then nonzero. This theorem does not remove all support26 terms or assert a complete support-floor increase. It does not infer signed nonvanishing, a source constructor or a DP44 saving.

The numerical constants are the actual g0=22285932761, g1=64072175441, Q=g0g1, U=77Q²−14, R=17424/133 and the unchanged J56 prime T. Use directed rational logarithms

\[
 L_T=1019333/1024>\log_2T,\quad L_Q=287843/4096<\log_2Q,
 \quad L_U=601356/4096<\log_2U,\quad L_R=113/16>\log_2R.
                                                        \tag{PC2}
\]

The last strict inequality is checked as R\^{}16\textless2\^{}113, and all other directions are checked by exact integer powers.

\subsection{2. High finite rank at25 and26, and the support-free plane reduction}\label{high-finite-rank-at25-and26-and-the-support-free-plane-reduction}

Put alpha=(N²)\_nn.~The actual unitary normal form gives 0\textless alpha\_tau\textless=R at every complex label. R058(HB3--14) applies with its support dependence exactly as displayed there. In particular, for

\[
 \mathfrak J_F=\prod_v\mathfrak p_v^{2-d_v},\qquad
 x=(N_F\mathfrak J_F\,T^{-2s})^{1/(2f)},
                                                        \tag{PC3}
\]

D\_G\textgreater=33 implies N\_F J\_F\textless=Q\^{}334/g0\^{}33. Take j25=1306 and j26=1401. The companion checks

\[
 Q^{334}2^{21j_s}<g_0^{33}T^{2s},\quad s=25,26;
 \qquad0<x<2^{-j_s/16}<1.                                \tag{PC4}
\]

R058's complete all-profile envelope, including the original target-column coarea and source fibre, is increasing in x. Its exact upper exponents are

\[
 E_H(s)=1+5379+1062+107+168+3s+530-464f
       -335\cdot167L_Q+343sL_T+42-55776j_s/16.             \tag{PC5}
\]

\begin{longtable}[]{@{}rrrr@{}}
\toprule\noalign{}
s & j\_s & Full exponent E\_H(s) & Strict bound \\
\midrule\noalign{}
\endhead
\bottomrule\noalign{}
\endlastfoot
25 & 1306 & −77307919/4096 & \textless−18874 \\
26 & 1401 & −35243075/4096 & \textless−8604 \\
\end{longtable}

These are new evaluations of the support-independent formulas, not an extension of a numerical theorem without its hypotheses. R058(HB7--10) already includes all labels and profile weights, and its3\^{}334 pattern factor is already in PC5. No later pattern or support union is suppressed.

For D\_G\textless=32, R058(HB20--22) gives im N²=ell contained in the original packet plane P=Eu perpendicular E e\_n.~Its argument uses three separate globally integral source/target pairings with ideal J\_F and the unchanged finite-rank-budget69 norm contradiction. It contains no numerical support restriction. We use this plane consequence only; we do not reuse the old pole-budget46 emptiness claim at25 or26.

Let V=u\^{}perp in the original positive E\^{}4 and Z=S perpendicular V. The positive definite quotient ell\^{}perp/ell has a unipotent unitary action, hence identity. Therefore h\_ZZ=I and N(V) is contained in ell. The induced cross covector on V is nonzero, since otherwise N acts only on P and has square zero. R077(FB8)'s exact block calculation then gives

\[
 C=\ell r,\quad\ell=(1,\eta),\quad\eta\bar\eta=U,
 \quad v=J_Z^{-1}r^*,\quad q=rv,\quad \alpha=Uq,
\] \[
 N=\begin{pmatrix}0&-v\ell^*J_P\\\ell r&d\ell\ell^*J_P\end{pmatrix},
 \quad J_P=\operatorname{diag}(U,-1),\quad
 d+\bar d=-q,\quad w=N_{nn}=-Ud.                          \tag{PC6}
\]

The covector r vanishes on S. Thus the actual cross map C and original central entry w determine the entire matrix. There is no remaining column or isometry fibre. In particular Re(w)=alpha/2 is fixed once C is fixed. This is the accepted block identity with its premises verified at the present supports.

\subsection{3. Original ideals, true metric and complete reconstruction data}\label{original-ideals-true-metric-and-complete-reconstruction-data}

Use the same fixed original transverse vectors as R077:

\[
 y_1=(-976627,-153032,38026,0),\quad
 y_2=(-567678,695570,-172838,0),\quad
 y_3=(256766,-7654233287,-30805128044,1).
                                                        \tag{PC7}
\]

They form an E-basis of V. Put a\_j=r(y\_j) and M=sqrt(R/U). Then \textbar\textbar r\_tau\textbar\textbar\textless=M in the true positive dual metric. The Gram matrix G=(y\_i dot y\_j) satisfies

\[
 G=\begin{pmatrix}
 978665066829&441392856078&-303935233242\\
 441392856078&835948910828&95913020934\\
 -303935233242&95913020934&1007543201084982869062
 \end{pmatrix},\quad\det G=4U,\quad G>2^{16}I.             \tag{PC8}
\]

The last comparison follows from the three positive leading minors of G−2\^{}16I, checked exactly. No integral change to an orthogonal frame is made.

Choose the first nonzero a\_j in the order1,2,3, and put beta=eta a\_j. As in R077(FB9--10), at the G places one common whole-cross-map ideal J controls all a\_j and beta, with at most4\^{}334=2\^{}668 patterns, N\_EJ\textless=Q\^{}334 and

\[
 L=J I_\Sigma,\quad N_EI_\Sigma=T^{-2s},\quad
 z=(N_EL)^{1/(2f)},\quad a_1,a_2,a_3,\beta\in L,
\] \[
 |a(h)|\le2Q^{-335\cdot167}T^{335s}z^{335f}.              \tag{PC9}
\]

Here I\_Sigma pads a simple pole at both partners of every active T place. This remains valid for all three types: each original y\_j belongs separately to every actual source and target pencil, so Ny\_j lies in the actual target simple-pole lattice. Primitive integer Bezout extraction for u gives both its scalar and negative coordinates in that ideal. Type-specific stronger integrality is harmlessly enlarged. The fixed V,P frame has G-unit determinant; the rank-one square-zero reduction and inverse-adjoint relation give the same one-partner orientation, while rank zero gives both. Both mixed finite factors retain their modulus bound one.

The covector metric on its coordinate values is G\^{}-1. The original scalar ideal has half-trace covolume D0 z\^{}(2f), with D0=7\^{}140 29\^{}144 5\^{}84 13\^{}84. In three coordinates its dual-metric covolume is D0³ z\^{}(6f)/(4U)\^{}f.~The arbitrary-ideal cover contracted by G\^{}-1/2 has radius below4 sqrt(f) z because

\[
 (2250/7)^2\,3<16\,2^{16}.                               \tag{PC10}
\]

Consequently the full original three-covector count is

\[
 \#r\le A_3(4U)^f(M+4z)^{6f}z^{-6f},\quad
 A_3=528^{3f}/((3f)!D_0^3)<2^{-4040}.                     \tag{PC11}
\]

It counts every allowed covector at fixed L, including zero only as an upper bound. Actual reconstruction always uses a nonzero pivot.

The actual relative norm equation is beta bar(beta)=U a\_j bar(a\_j), with beta in L. K075(RN17--19) constructs its integral relative norm ideal only for existing solutions, so nonprincipality or lack of conjugation stability of L introduces no missing factor. DO091(TD11--12) gives the actual-field divisor constant C\_F\textless2\^{}42 at exponent1/12, and each eligible principal ideal has at most28 norm-one generators. If the pivot squared length parameter is B, meaning \textbar a\_j,tau\textbar²\textless=B q\_tau, then

\[
 \#\beta\le28C_F(RB)^{f/12}z^{-f/6}.                     \tag{PC12}
\]

The ratio beta/a\_j reconstructs eta over E only. We do not declare that ratio integral. The same inequality q\textgreater=\textbar a\_j\textbar²/B and the nonzero-element norm bound in L give N\_F alpha\textgreater=(U/B)\^{}f z\^{}(2f). The genuine selected trace, with both compact factors and the original metric and transport, is therefore

\[
 |c(h)|<2^{-929f/2}(U/B)^{-f/2}z^{-f}.                    \tag{PC13}
\]

This is DO080(RT3), not an operator substitute. Equations PC12--13 retain the exact pivot dependence for the next split.

\subsection{4. The original central fibre cancels all extra incidence powers}\label{the-original-central-fibre-cancels-all-extra-incidence-powers}

Fix the cross map and a common G pattern. Let J0 be the F-product of its actual rank-zero G primes once. Its original central entry is divisible by J0 at both partners. The same finite table as R077(FB13) gives

\[
 N_EJ\le Q^{167}N_FJ_0,\qquad N_FJ_0\le Q^{167}.           \tag{PC14}
\]

K071(NB5) bounds the full original entry w by pole order1,2,1 at I,II,III respectively, at both partners. Its real part is fixed by PC6. As i is an integral unit in E and O\_E intersect iF=iO\_F, the difference of two possible w lies in

\[
 \Gamma=iJ_0 I_{\rm I,F} I_{\rm II,F}^{\,2}I_{\rm III,F},
 \quad\operatorname{covol}(\Gamma)
       =\sqrt{D_F}\,N_FJ_0T^{-(s+u)},\quad D_F=2^{168}D_0.
                                                        \tag{PC15}
\]

The fibre is contained in an affine coset, with its shift retained. No division by2 or replacement of a difference lattice by a half-lattice occurs. Its dual is the full real codifferent times the inverse ideal, as in DO091(TD20--22). The actual-field cover is below

\[
 750\sqrt f\,(N_FJ_0T^{-(s+u)})^{1/f}<\sqrt f,             \tag{PC16}
\]

since750\^{}f Q\^{}167\textless T\^{}25\textless=T\^{}(s+u) for s25 or26; the strict integer inequality is checked. The original \textbar Im w\textbar\textless H/2, so the translated Voronoi argument uses radius68 sqrt(f). The complete real f-dimensional volume V\_F(68)\textless2\^{}479 from R077(FB15) therefore gives

\[
 \#w<2^{479}\frac{T^{s+u}}{N_FJ_0}
 \le2^{479}Q^{167}z^{-2f}T^{u-s}.                         \tag{PC17}
\]

The last equality uses N\_EJ=T\^{}(2s)z\^{}(2f). This counts the whole actual central fibre after C is fixed, so PC6 removes any separate isometry or negative-column factor legitimately.

At each fixed support and type assignment, every bound PC9--17 holds uniformly in the original ordered source/target labels, including zero. K071(NB7--8) gives the full ordered-pair count times its true absolute incidence weight as (T−1)\^{}i at mostT\^{}i. Thus the only profile power left by PC17 is

\[
            T^{i+u-s}=T^{-t}\le1.                         \tag{PC18}
\]

This is a bound applied at each actual label choice before summing. It assumes neither an independent CM-partner label nor a hypothetical matrix with all-II incidence. Retain the full3\^{}s\textless2\^{}(2s) type assignment count and binom(f,s)\textless2\^{}168 support count. A zero-target coarea factor is not needed because PC6--17 count the original complete matrix directly.

\subsection{5. Short pivot or one-dimensional last pivot}\label{short-pivot-or-one-dimensional-last-pivot}

The common cap in PC9 is

\[
 z\le z_{\max}=Q^{167/f}T^{-s/f}<M/5,\qquad M+4z<2M.
                                                        \tag{PC19}
\]

For s25,26 it follows from the exact inequality R\^{}fT\^{}(2s)\textgreater25\^{}fU\^{}fQ\^{}334. This is checked for both supports. It makes the full-ball bounds uniform at every common ideal, without dropping a large-cover tail.

\textbf{Short-pivot case.} If a1 or a2 is nonzero, choose the first such value. Their actual squared lengths satisfy B=\textbar\textbar y\_j\textbar\textbar²\textless2\^{}40. Use the entire three-covector count PC11 with radius2M, together with PC12--13 and PC17--18. The complete z exponent is

\[
 (335-6-1/6-2-1)f=(1955/6)f>0.                           \tag{PC20}
\]

At z\_max the total Q power, including the central Q\^{}167, is−49·167/6, and the T power is55s/6. The determinant4U and radius2M give2\^{}(8f)U\^{}(-2f)R\^{}(3f); the norm fibre and trace then give total exponents37f/12 on R, −5f/2 on U, and7f/12 on B. Define a bound that retains a common pivot factor3\textless2²:

\[
\begin{split}
 C_3(s)={}&1+668+168+2s+2+479+5+42-4040+8f-929f/2\\
          &+(37f/12)L_R-(5f/2)L_U+(7f/12)40,\\
 E_3(s)={}&C_3(s)-(49\cdot167/6)L_Q+(55s/6)L_T.
\end{split}                                             \tag{PC21}
\]

Before that common pivot factor, each individual short pivot is bounded by2\^{}(E3(s)−2).

\textbf{Last-pivot case.} If a1=a2=0, the remaining covector metric is exactly

\[
 q=(G^{-1})_{33}|a_3|^2=|a_3|^2/W,\qquad
 W=\frac{1}{(G^{-1})_{33}}
   =\frac{4U}{\det(G_{\{1,2\},\{1,2\}})}.                \tag{PC22}
\]

The Schur complement and PC8 give2\^{}16\textless W\textless=G33\textless2\^{}70. These are also checked by exact rationals. The actual one-coordinate lattice is L with its metric scaled by1/sqrt(W); its real covolume is D0 z\^{}(2f)/W\^{}f and its covering radius is below4 sqrt(f) z. Therefore

\[
 \#r\le A_1 W^f(M+4z)^{2f}z^{-2f},\quad
 A_1=528^f/(f!D_0)<2^{-1083}.                            \tag{PC23}
\]

This is a full volume bound in that precise coordinate line, not a density multiplier obtained by imposing zero conditions on the three-dimensional count. In PC12--13 use B=W, since q=\textbar a3\textbar²/W exactly. The same genuine beta norm equation, trace, full central affine fibre and actual incidence cancellation apply.

The remaining z exponent is(335−2−1/6−2−1)f=(1979/6)f\textgreater0. At the cap, the Q power including the central factor is−25·167/6 and the T power31s/6. The total exponents on R,U,W are13f/12,−3f/2,19f/12 respectively, with the radius factor2\^{}(2f). Retaining the same harmless pivot allowance gives

\[
\begin{split}
 C_1(s)={}&1+668+168+2s+2+479+5+42-1083+2f-929f/2\\
          &+(13f/12)L_R-(3f/2)L_U+(19f/12)70,\\
 E_1(s)={}&C_1(s)-(25\cdot167/6)L_Q+(31s/6)L_T.
\end{split}                                             \tag{PC24}
\]

Before that allowance the unique last pivot is bounded by2\^{}(E1(s)−2).

\subsection{6. Exact endpoints, complete unions and chronology}\label{exact-endpoints-complete-unions-and-chronology}

The rational endpoints are

\begin{longtable}[]{@{}rrrl@{}}
\toprule\noalign{}
s & Short case E3(s) & Last case E1(s) & Use here \\
\midrule\noalign{}
\endhead
\bottomrule\noalign{}
\endlastfoot
25 & −27553337/24576\textless−1121 & −363263969/24576\textless−14781 & Every low-rank pivot \\
26 & 65583025/8192\textgreater0 & −78939175/8192\textless−9636 & Last pivot only \\
\end{longtable}

The positive short-case envelope at26 is failure of this sufficient bound only; it makes no assertion about a contribution or a counterexample. At25 there are at most two short pivots and one last pivot. Since E1(25)\textless E3(25), their sum is less than3·2\^{}(E3(25)−2)\textless2\^{}E3(25). Thus the factor2² in PC21--24 counts all three alternatives at once; there is no additional factor two for the two kinds of pivot. At26 the last-pivot sector alone is bounded by2\^{}E1(26), with the unused allowance harmless.

High and low finite-rank alternatives are exhaustive at25. Combining PC5 and the exact low-pivot bounds proves that the entire union in the theorem is below

\[
 2^{-18874}+2^{-1121}+2^{-8604}+2^{-9636}<2^{-1120}.         \tag{PC25}
\]

The separate two-support union has now been counted explicitly. Both high-rank exponents already include every type assignment; both low-rank exponents include2\^{}(2s). No alpha range, norm-content, class-number, central, source-fibre or coordinate-position union is missing.

At integration append1120 exactly once to the actual chronological all-top array, retaining every prior charge and the unchanged independent uniform array. The companion records its explicitly chosen predecessor and checks the full exact dyadic sums; integration verifies the then-actual predecessor again. Remove the intersection of the containing union with the current remainder, so overlaps with old removed sets give no problem. The complete floors remain25/11 until a separate exhaustive support argument is proved. No actual selected arithmetic matrix is produced by this upper-bound calculation.

\subsection{7. Exact source placement and limits}\label{exact-source-placement-and-limits}

The directly used source interfaces are R058(HB3--18,20--22), R077(FB6--17), K071(NB5,7--9), R072(PJ3--7) for the whole-cross ideal at the unchanged G places, K075(RN17--19), DO091(TD11--12,20--22), and DO080(RT3). The new steps are the actual-profile affine central lattice PC15--18, the short versus last-only true covector count PC21--24, and the complete present endpoint union PC25. The high finite-rank endpoints PC4--5 are new exact evaluations with every original factor retained. Settled normal forms, transference and selected-trace results are cited in their existing full proved scopes.

The covector and its ideals live on the original Hermitian modules over Spec(O\_E); their actual dual metric and the one-coordinate Schur complement retain the polarization. The relative norm fibre and affine imaginary central line keep both CM partners and the real discriminant. The contextual module/exterior anchors remain SGA3 I4.7.3/4.7.3.1 and SGA6 VI Theorem3.3. Picard and biextension anchors are FGA232 Theorem3.1 pp148--149; FGA236 Corollaries1.7/2.3 pp229/232; SGA7 I VII2.7/2.9 pp169--174 and VIII7.4(b) pp310--311. Crystals and Cartier duality retain cote4 scan003 §§8.1.1--8.1.2, scans005--006 §§8.4.1--8.4.3 and cote54 scans005--007.

These are unchanged geometric frameworks, with no new lambda-ring equality, K-theory comparison, l-adic sheaf or weight statement, motivic projector or CM-motive identification. DP\_MULT Q3e and THETA3 G1/A7 gain only the stated smaller actual finite/real remainder. Literal admissible mixed-CM source construction, prescribed global periods and selected arithmetic-image content, the specific032 matching/value, all actual DP44 terms, CX\_GRH CX-ALG-024 and the global arithmetic comparison for GCT\_PLETH T-052 remain open. The companion is exact arithmetic verification of constants and bookkeeping, not rational-arrow enumeration. No Lean build or archive audit is performed.

\clearpage

\section{AM-DO-100 --- a complete large-cross-ideal removal for rank three at support25}\label{am-do-100-a-complete-large-cross-ideal-removal-for-rank-three-at-support25}

\textbf{Canonical checked integration --- 2026-11-22T23:56:36.315621+00:00.} Root and independent whole-proof/script/interface readings and separate561-assertion replays passed without correction. Historical pending wording below is preserved. The complete global-rank-three scalar-cross support25 sector with the deterministic original common-C ideal satisfying2\^{}32928 N\_EJ\textgreater=T\^{}50 is removed with one29550 charge. The whole original167-coordinate dual-metric ball, relative norm fibre, selected trace, paired projected column, source fibre, all patterns and complete actual incidence weights remain. The full monotone high-ideal endpoint is -121040253/4096. Every remaining rank-three support25 term satisfies2\^{}32928 N\_EJ\textless T\^{}50. The weaker low-ideal extension is a failed sufficient bound4090765885/12288\textgreater0; no actual contribution is inferred. A complete local typeII shear model shows why C-integrality cannot be substituted. Actual86/31 errors remain below2\^{}-230 with complete floors25/11. No whole support25, signed/global/source/selected-image/specific032/DP44 or RH conclusion.

Author: \nolinkurl{/root/definite_oscillator}, GPT-6 Astra Pro, ALP\_MACHINE. 22 November 2026 UTC. Controlled pending proof; no canonical ledger change is made by this file. A whole saved-proof/source/script reading and an exact replay are recorded separately before freezing.

\subsection{1. Actual remaining sector and the canonical cross ideal}\label{actual-remaining-sector-and-the-canonical-cross-ideal}

Keep the unchanged all-top J56 endpoint H=133, f=168, n=169, every actual source/target pencil, both CM components, finite projector, central representative, selected metric and trace, class mass and phase. At support s=25, DO099(JR24) leaves the actual R071 full-positive-compression sector

\[
 E^{169}=Z\perp P,\quad Z=S\perp u^\perp,\quad \dim S=164,
 \quad\dim Z=167,\quad P=Eu\perp Ee_n,
 \quad h=\begin{pmatrix}I_Z&B_Z\\ C_Z&D_P\end{pmatrix}.
 \tag{HI1}
\]

Here C\_Z has global rank one and is nonzero on the original S. The positive166-dimensional kernel of C\_Z is fixed pointwise, but its chosen original transverse vectors need not be fixed individually. Restrict to global rank\_E(h-I)=3. No global unipotence is assumed. The K068 condition that all334 unmixed G reductions are unipotent is retained.

Specify the common-C ideal exactly, rather than choosing a favorable one after a count. At each unmixed F-prime use the original designated CM component first. If its reduction has rank zero, place both CM prime ideals in J. If its reduction has rank one, first test whether C\_Z vanishes at that designated component: choose that prime if it does, and choose the conjugate prime otherwise. At rank at least two put neither prime in J. Rank-one factorization of N=h-I with N\_ZZ=0, together with its actual inverse-adjoint partner, guarantees that the latter choice always kills C\_Z. This is the fixed-rule version of R072(PJ6) and R076(CV4--5), with at most4\^{}334=2\^{}668 possible patterns. It uses the fixed G-unit Z/P frame, not a finite positive metric.

Put

\[
 L=J I_\Sigma,\qquad N_E I_\Sigma=T^{-50},\qquad
 z=(N_E L)^{1/336}=(N_EJ)^{1/336}T^{-25/168}.
 \tag{HI2}
\]

\textbf{Theorem.} The complete absolute contribution of the actual rank-three support25 terms satisfying

\[
 z\ge 2^{-98},\qquad\text{equivalently}\qquad
                     2^{32928}N_EJ\ge T^{50},
 \tag{HI3}
\]

is less than2\^{}(-29550). This is one complete sector, containing every incidence profile and every allowed original labelled pair. Removing its current-remainder intersection leaves the complementary rank-three terms subject to

\[
                         2^{32928}N_EJ<T^{50}.
 \tag{HI4}
\]

The complete support floors remain25/11. Rank-two sectors and every signed, global-period, arithmetic-image, source-constructor and DP44 gate retain their separate status. The failed whole-profile extension is quantified in section4; it is not used to infer existence or a nonzero contribution.

\subsection{2. The full original count and the strong finite bound}\label{the-full-original-count-and-the-strong-finite-bound}

Use precisely the original R072 transverse inputs and the R076 covector normalization. Write

\[
 C_Zx=(u\lambda(x),\eta\lambda(x)),\quad\eta\bar\eta=U,
 \quad U=77Q^2-14,\quad Q=1427908193729121522601.
\]

The164 original values xi\_i=lambda(e\_i), the three values on

\[
 (-976627,-153032,38026,0),\quad
 (-567678,695570,-172838,0),\quad
 (256766,-7654233287,-30805128044,1)
\]

and beta=eta xi\_j at one nonzero original pivot determine C\_Z over E. Retain164\textless2\^{}8 pivot choices. Every displayed scalar belongs to L by R072's original-coordinate extraction and actual source/target membership at both CM primes. No rational ratio or moving frame is called integral.

At every complex label the actual negative row gives

\[
 \|\lambda_\tau\|<M=\frac{H}{2\sqrt U},\qquad
 |\beta_\tau|<B=H/2,
 \qquad |c(h)|<2^{-929f/2}U^{-f/2}z^{-f}.
 \tag{HI5}
\]

The last inequality is R072(PJ9)'s actual normalized selected trace, including its compact factors and near-point transport. It is not an operator-norm substitution.

The exact positive covector Gram matrix is diag(I\_164,G\^{}-1), where detG=4U and G-I\_3 is positive definite. R076(CV8--10) proves that L\^{}167 has real half-trace covolume D0\^{}167 z\^{}(334f)/(4U)\^{}f in this metric, and covering radius less than4200 sqrt(f)z. Consequently

\[
 \#\lambda\le A(4U)^f(M+4200z)^{334f}z^{-334f},
 \quad A=\frac{528^{167f}}{(167f)!D_0^{167}}<2^{-387281},
 \quad D_0=7^{140}29^{144}5^{84}13^{84}.
 \tag{HI6}
\]

This single joint ball includes every original zero pattern; there is no binomial or k-union. Its rational Gram transformation is used for the metric calculation only.

For fixed lambda and the chosen pivot, the exact equation beta bar(beta)=U xi\_j bar(xi\_j) gives R076(CV11--13)'s complete relative-norm fibre

\[
 \#\beta\le28C_F B^{f/6}z^{-f/6},\qquad
 C_F=\frac{2^{300}}{29^6 41^{14}2129^{14}}<2^{42}.
 \tag{HI7}
\]

Here the ideals (beta)L\^{}-1 are counted before their principality is tested, and the at-most28 norm-one generators are retained. The ideal L need not be conjugation-stable or principal. This is the accepted K075(RN17--19)/DO091(TD11--12) argument, not a new choice of coefficient field.

After C\_Z, retain the entire R072(PJ14--18) projected negative-column count and source fibre

\[
                   191^{4f}U^fT^{4s},\qquad G_2<2^{36}.
 \tag{HI8}
\]

The U\^{}f factor is the actual fixed projection denominator. The paired target index is T\^{}(4s), including the admitted zero target; it is not replaced by a false universal pole-two projection. Unitarity reconstructs the full negative column from C\_Z and its plane projection. The positive rank-two source-intersection fibre is retained, even though other identities can sometimes determine it more sharply.

The strong rank-three estimate is DO095(RA15--20). Its nonzero original S entry of BC is an actual three-minor, bounded by2\^{}9 at every conjugate, with only simple paired type-III poles. If t is the number of type-III places, its original finite normalization gives

\[
 |a(h)|<2Q^{-501(f-1)}2^{9\cdot167f}T^{167t}
       \le2Q^{-501(f-1)}2^{9\cdot167f}T^{167s},\quad t\le s=25.
 \tag{HI9}
\]

This introduces no additional C pattern, divisor pattern or witness-choice factor. The already retained2\^{}668 common-C patterns remain necessary for L and its count. Both mixed finite factors retain modulus at most one.

Finally retain support choices2\^{}168 and the complete actual ordered-pair/absolute-incidence sum P\_T\^{}25\textless2\^{}(997*25). It already sums all type profiles with their paired labels. Together with the Green factor2, common patterns, pivot and source fibre, the fixed exponent is

\[
                         1+668+168+8+36=881.
 \tag{HI10}
\]

No type enumeration, class/Haar normalization or real factor is removed.

\subsection{3. Monotone large-ideal endpoint and the exact whole charge}\label{monotone-large-ideal-endpoint-and-the-exact-whole-charge}

Multiply HI5--HI10, for each original pattern and actual pair. Apart from positive constants, its dependence on z is exactly

\[
              (M/z+4200)^{334f}z^{-7f/6}.
 \tag{HI11}
\]

Both factors decrease strictly on z\textgreater0. Thus the entire high-ideal sector is bounded by the single endpoint z*=2\^{}(-98), without asserting that an ideal attains that endpoint. There is no endpoint-union factor.

The actual integers give M\textgreater2\^{}(-68), so

\[
 \frac{4200z_*}{M}<\frac{4200}{2^{30}}<2^{-17},\qquad
                       (1+2^{-17})^{334f}<2.
 \tag{HI12}
\]

The second strict inequality is checked by integer powers. Therefore the full enlarged radius at z* costs less than one binary bit. In contrast with R076's support24 low-ideal branch, no use of a lower bound for z is needed anywhere in this sector.

Retain the exact directed original logarithms

\[
 L_T=1019333/1024>\log_2T,\quad
 L_Q=287843/4096<\log_2Q,\quad
 L_H=28900/4096>\log_2H,
\] \[
 L_U=601356/4096<\log_2U,\qquad
 L_{191}=15519/2048>\log_2(191).
\]

With the same complete factors as R076, define

\[
\begin{split}
 C_{25}={}&881+1+5+42-387281+4fL_{191}-929f/2+2f\\
 &+(2005f/6)(L_H-1)-(331f/2)L_U+100L_T+24925.
\end{split}
 \tag{HI13}
\]

Here1 is the controlled joint-radius inflation,5 bounds28,42 bounds C\_F,2f is4\^{}f from detG, and2005f/6=334f+f/6 includes the entire beta fibre. The coefficient of L\_U includes both the actual projection denominator and selected trace. The endpoint's remaining z power is2011f/6=334f+f+f/6.

The exact full upper exponent is

\[
\begin{split}
 E_{25}={}&C_{25}-501\cdot167L_Q+9\cdot167f
                  +167\cdot25L_T+98(2011f/6)\\
          ={}&-121040253/4096<-29550.
\end{split}
 \tag{HI14}
\]

Every counted pattern in HI3 is bounded by this same endpoint. The original pattern, support, pair and pivot sums have already been included, so HI14 proves the single complete2\^{}(-29550) bound with no further union. Intersecting with the actual current remainder can only reduce the absolute sum.

\subsection{4. Exact boundary of this extension and accounting}\label{exact-boundary-of-this-extension-and-accounting}

The natural full-support extension of R076's other endpoint does \textbf{not} pass. Using its weaker335 Green bound, its true lower norm z\textgreater=T\^{}(-25/f), and every unchanged factor gives

\[
 E_{\mathrm{low},25}
   =C_{25}-335\cdot167L_Q+(2011\cdot25/6)L_T
   =4090765885/12288>0.
 \tag{HI15}
\]

Thus replacing167s by167t only in the strong rank-three estimate cannot justify a whole-profile support25 removal. Even the favorable high-z exponent at t\textless=24 does not repair HI15. No conclusion that an actual term exists or has large absolute value follows from the failure of this sufficient upper envelope.

There is also no local inference that global rank three makes the entire C block integral at type II. In a split three-dimensional local model ordered as one S coordinate and two P coordinates, put

\[
 h_0=\begin{pmatrix}1&1&0\\1&1&1\\1&0&1\end{pmatrix},\quad
 Z_0=e_2e_3^t,\quad K_c=I+(c/\pi)Z_0,\quad h=K_ch_0K_c^{-1}.
 \tag{HI16}
\]

One has det h0=1, (h0)\_SS=(h0\^{}-1)\_SS=1, rank(h0-I)=3, and BC=1. For a unit c,

\[
 h=\begin{pmatrix}1&1&-c/\pi\\1+c/\pi&1&1\\1&0&1\end{pmatrix},
 \qquad C=\binom{1+c/\pi}{1}.
\]

The matrix preserves the actual-form lattice K\_c Lambda and has a genuine cross pole. For a different finite label d modulo pi, K\_d\^{}-1 h K\_d has a nonintegral entry. The infinite lattice is O e\_1 plus pi\^{}(-1) O e\_2 plus pi O e\_3; h e\_1 has last coordinate1, so it is not fixed either. Moreover the sum of the pencil lattices is not preserved: h(pi\^{}(-1)e\_2) has nonintegral S coordinate. Thus the image full pencil is different and c is its unique fixed pencil label. The conjugate component is defined by the required inverse-adjoint relation and preserves the paired dual. A nonzero primitive rank-one nilpotent pencil shear on P admits this local integral normal form. This is only a local scope counterexample, not a global J56 short arrow, a claim of its finite G behavior, or a finite positive metric. It explains why the correct type-II conclusion is BC-integrality from the common selfdual source/target pairing, as already used in HI9.

At checked integration append29550 once to the then-current chronological all-top list, preserving every previous charge and the entire uniform list. The companion pins the actual canonical R07985/31 predecessor, checks its full exact sums and the staged86/31 sums below2\^{}(-230), and records unchanged floors25/11. If another independently checked charge is integrated first, it must be retained in a supplemental canonical-order check; it is not erased by this historical staged array.

The objects remain the same Hermitian source/target modules, original Hom and determinant sections, their paired prime ideals and their unchanged complex half-trace metric. The fixed Gram transformation and rational packet projection retain their exact determinants; neither supplies a new integral frame. The carried algebraic-geometric contextual anchors remain those explicitly recorded in R076 section6: SGA3 I4.7.3/4.7.3.1, SGA6 VI Theorem3.3, FGA232 Theorem3.1, FGA236 Corollaries1.7/2.3, SGA7 I VII2.7/2.9 and VIII7.4(b), and the already pinned cote4/cote54 locators. No new comparison theorem, lambda-ring/K-theory identity, l-adic weight assertion or CM-motive result is inferred.

DP\_MULT Q3e and THETA3 G1/A7 gain exactly HI4's smaller finite/real remainder. The signed sufficient inequality, actual fully admissible source, prescribed global Bessel/Mram value, selected arithmetic-image content, Read/STiesR/specific032 pairing and changed-family DP44 comparison remain open. The certificate verifies the actual constants, original Gram, all powers, endpoint, local scope example and chronological sums; it does not enumerate the arithmetic kernel or establish a global period by calculation.

\clearpage

\section{AM-DO-101 --- reciprocal poles remove the remaining nonunipotent rank-two interval25--32}\label{am-do-101-reciprocal-poles-remove-the-remaining-nonunipotent-rank-two-interval2532}

\textbf{Canonical checked integration --- 2026-11-23T00:01:39.855234+00:00.} Root and independent whole-proof/script/interface readings and separate630-assertion replays passed without correction. Historical pending wording below is preserved. The entire remaining globally nonunipotent rank-two scalar-cross interval25..32 is removed after the retained DO099 transverse cut. The support-independent DO076 reciprocal-line theorem gives lambda in F and inverse-C=-lambda\^{}-1 C without requiring full-Z at26+. At least23 actual lambda poles and one global inversion give at least12 positive valuations; all active lambda-unit places are retained in the complete paired C-pole cap s-12. Every original column, source fibre, pattern, pivot, support and labelled-pair weight remains. The complete worst exponent is -2610113463/20480 at32, and eight supports cost one127443 charge. Actual87/31 ordered errors remain below2\^{}-230 with complete floors25/11. Support33 is a failed sufficient bound only. Global-unipotent rank2, rank3, signed/global/source/selected-image/specific032/DP44 and RH gates remain separately open.

Author: \nolinkurl{/root/definite_oscillator}, GPT-6 Astra Pro, ALP\_MACHINE. 22 November2026 UTC. Controlled pending proof; whole proof/interface/script verification and exact replay are recorded separately. No canonical record is changed by writing this file.

\subsection{1. Actual sector and a support-independent reciprocal-line input}\label{actual-sector-and-a-support-independent-reciprocal-line-input}

Keep the original all-top J56 kernel, H=133, f=168, n=169, all original finite projectors, ordered pencil pairs and their Hermitian partners, the actual selected trace and metric, central normalization, source-lattice fibre, Haar/class weights and phases. Use the original decomposition S perpendicular W0, with dimS=164 and signature(W0)=(4,1). Restrict to the actual K068 residual in which A=h\_SS=I, the lower cross block C is nonzero, det h=1 and all334 unmixed G reductions are unipotent. Impose

\[
 25\le s\le32,\qquad \operatorname{rank}_E(h-I)=2,
                  \qquad h\text{ is not globally unipotent}.
 \tag{PI1}
\]

DO099(JR1--12) has already removed the C-transverse containing sectors at every support25..70, with its original charge retained. Therefore the current remainder in PI1 has C in the original packet plane P=Eu perpendicular E e\_n.~We bound the entire containing plane-C sector, not merely an inverse-closed subset of an earlier cut.

\textbf{Theorem.} The union of these remaining globally nonunipotent rank-two sectors over all eight supports25..32 has total absolute normalized contribution less than

\[
                              2^{-127443}.                  \tag{PI2}
\]

This is one fresh all-top charge. No global-unipotent rank-two term, global rank-three term or full support-floor increase is asserted.

It is essential not to import the full-Z compression at supports26 and above: DO099 supplies that stronger conclusion only at25. Instead use the genuinely support-independent algebra in DO076(RS3--6). Its proof requires exactly A=I, C nonzero, det h=1, actual unitary signature one, global rank(h-I)=2 and nonunipotence. The two-dimensional nonidentity primary space is nondegenerate; the isotropic original line N(S) is an h-eigenline defined over E; the determinant and pairing make its eigenvalue lambda belong to F. Consequently

\[
 \lambda\in F^\times,\quad\lambda\ne1,\quad
 p_h(X)=(X-1)^{167}(X-\lambda)(X-\lambda^{-1}),
 \qquad C(h^{-1})=-\lambda^{-1}C(h).
 \tag{PI3}
\]

These statements permit lambda=-1 and a repeated nonsemisimple nonidentity primary block; no distinct-eigenline or global-nilpotence assumption is made. Their derivation in RS3--6 has no support parameter and no full167-dimensional identity-compression premise. This checks the exact interface needed for the interval. In particular plane-C is preserved by inversion.

\subsection{2. At least23 actual poles and at least12 positive valuations}\label{at-least23-actual-poles-and-at-least12-positive-valuations}

Define the nonzero original trace witness

\[
 a=\operatorname{tr}(h)-169
       =\lambda+\lambda^{-1}-2=(\lambda-1)^2/\lambda\in F.
 \tag{PI4}
\]

Every conjugate satisfies \textbar a\_tau\textbar\textless=H+H\^{}(-1)+2\textless256. At an unmixed G place, both eigenvalues are units and reduce to1, by their original monic characteristic equations and the finite-unipotence premise. Hence a has valuation at least two at each of the334 unmixed F primes. K070(PT5--8), in the exact use of DO095(RA21--23), gives integrality outside type III and pole order at most one at each type-III F place. Let r be the number of \textbf{actual} poles of a. The original norm comparison is

\[
              Q^{334}\le |N_Fa|T^r<256^{168}T^r,
 \qquad Q=1427908193729121522601.
 \tag{PI5}
\]

The exact original inequalities L\_Q=287843/4096log2T give

\[
       334L_Q-168\cdot8-22L_T=466617/2048>0.
 \tag{PI6}
\]

Thus23\textless=r\textless=s. This repeats the checked K096/DO096 threshold with the larger support range, not a new coefficient-pole assertion. It cannot be applied when a=0, which is why global unipotence is excluded in PI1.

If v(lambda) is nonzero, PI4 gives v(a)=-\textbar v(lambda)\textbar; for a unit lambda the valuation of a is nonnegative. Therefore at precisely the r actual pole places, v(lambda)=+1 or-1, and lambda is a unit elsewhere. All actual F places over T have normT. Write n\_+,n\_- for those two sign counts. The original eigenvalue and principal-norm identities give

\[
 n_++n_-=r,\qquad
 |N_F\lambda|=T^{n_+-n_-},\qquad
 H^{-168}\le |N_F\lambda|\le H^{168},\qquad T^2>H^{168}.
 \tag{PI7}
\]

Hence \textbar n\_+-n\_-\textbar\textless2. For an even r the difference is0. For an odd r it is+1 or-1. Use the exact inversion law DO093(IV4--15), together with PI3, to restrict to the nonnegative sign. This exchanges the actual ordered source and target at both partners together, preserves every cutoff and absolute normalized weight, and contributes a single global factor2 when counting the larger containing sector. It is not an independent inversion choice at each place. In that containing sector

\[
              n_+=\lceil r/2\rceil\ge12,\qquad
              n_-=\lfloor r/2\rfloor.
 \tag{PI8}
\]

For even r the factor2 merely overcounts. For odd r it covers the opposite sign by inversion. Neither a separate r-union nor a sign-profile multiplier is needed once the containing pole ideal below is chosen.

\subsection{3. The entire original C block has at most s−12 paired pole places}\label{the-entire-original-c-block-has-at-most-s12-paired-pole-places}

At each active T place, every original S basis vector e\_i belongs separately to the actual source M\_c and target M\_d.~The two arrows and PI3 give

\[
 Ce_i=(h-I)e_i\in M_d,\qquad
 -\lambda^{-1}Ce_i=(h^{-1}-I)e_i\in M_c.
 \tag{PI9}
\]

The literal one-shear source/target inclusions LIFT44(PR4--7), as used in K096(EP9--12) and DO096(RP10--12), are pi M\_c subset Lambda and M\_d subset pi\^{}(-1)Lambda at each CM partner. Thus v(lambda)=+1 makes the original C column integral. At a negative-valuation place retain one pole. At any of the s-r active lambda-unit places retain the same conservative one-pole bound. No stronger type-I or type-III integrality claim is used.

Because lambda belongs to F, this is a paired set of places. Its size is at most

\[
 n_-+(s-r)=s-\lceil r/2\rceil\le s-12=:v_s.
 \tag{PI10}
\]

Pad it, by a fixed order when needed, to a subset D of exactly v\_s active places. There are at most binom(s,v\_s)=binom(s,12)\textless=2\^{}s choices. The resulting actual fractional ideal is

\[
 I_D=\prod_{w\in D}(\mathfrak p_w\bar{\mathfrak p}_w)^{-1},
                         N_E I_D=T^{-2v_s}.
 \tag{PI11}
\]

For C e\_i=(u xi\_i,beta\_i), primitivity of the original integer vector u extracts xi\_i by a fixed integer linear combination of its positive entries. No U denominator is inserted, including at primes dividing U. With the unchanged R059 whole-C orientation ideal J at G, all164 scalars satisfy

\[
 \xi\in(JI_D)^{164},\qquad
 |a_G(h)|\le2Q^{-335(f-1)}(N_EJ)^{335/2},
                         \#\{J\text{ patterns}\}\le2^{668}.
 \tag{PI12}
\]

Here a\_G is the finite coefficient, not the scalar trace witness a. One component is chosen at finite rank one, both at rank zero, and none at rank at least two. The fixed rule may exchange components under inversion, but the single complete pattern enumeration already includes them. Both mixed finite factors retain their original modulus-at-most-one bounds.

\subsection{4. Complete original count and the whole interval charge}\label{complete-original-count-and-the-whole-interval-charge}

Use the unchanged deliberately weaker R059 bound \textbar tau(xi\_i)\textbar\textless B0=2\^{}(-65), following from H/sqrtU\textless B0. Put z=(N\_E(JI\_D))\^{}(1/336). Some xi\_i is nonzero, so the original field-norm minimum gives z\textless B0. The complete164-coordinate packing and positive endpoint exponent give

\[
 z^{335f}(1+2B_0/z)^{328f}
       =z^{7f}(z+2B_0)^{328f}<B_0^{335f}3^{328f}.
 \tag{PI13}
\]

Since N\_EJ=T\^{}(2v\_s)z\^{}(2f), the entire vector and Green loss therefore contribute at most

\[
                       T^{335v_s}2^{-65\cdot335f}3^{328f}.
 \tag{PI14}
\]

No sparsity or coordinate-support sum is added. For a nonzero original pivot, retain both full original R059 target-overlattice column counts, each2\^{}14861 T\^{}(10s), and the full source-lattice fibre G\_n\textless2\^{}5022. The first column C e\_j reconstructs C from xi, and the W0 part of h e\_n reconstructs its S part by unitarity. The rational reconstruction is not declared integral. The counts use the actual original target overlattice, including zero labels, and keep the full five-coordinate volume at each step. No reduced fibre, selected-trace improvement, real slice or shifted coarea assumption is introduced. Use only \textbar c(h)\textbar\textless=1.

The complete original pair/absolute-weight sum is P\_T\^{}s\textless2\^{}(997s), already containing all type and ordered-pair choices and their paired components. The unchanged R059/DO096 fixed exponent is

\[
 C_0=1+668+168+8+2+5022+2(14861)=35591.
 \tag{PI15}
\]

The summands retain respectively the Green Euler factor, common G patterns, support choices, original pivots, harmless original Case-B overcount, full source fibre and both full columns. The padded subset and single inversion contribute only2\^{}s and2. Every actual r, pole sign, lambda-unit position and incidence profile is already contained in these choices.

Using3\^{}5\textless2\^{}8 and the unchanged exact directed T,Q bounds, the entire support-s sector is bounded by2\^{}(E\_s), with

\[
\begin{split}
 E_s={}&35591+s+1-65\cdot335f+(8/5)328f\\
 &+[335(s-12)+20s]L_T+997s-335\cdot167L_Q.
\end{split}
 \tag{PI16}
\]

Its slope355L\_T+998 is strictly positive. The exact complete table is

\begin{longtable}[]{@{}rrr@{}}
\toprule\noalign{}
s & Padded paired pole places v\_s & Exact E\_s \\
\midrule\noalign{}
\endhead
\bottomrule\noalign{}
\endlastfoot
25 & 13 & −53414036843/20480 \\
26 & 14 & −46156333503/20480 \\
27 & 15 & −38898630163/20480 \\
28 & 16 & −31640926823/20480 \\
29 & 17 & −24383223483/20480 \\
30 & 18 & −17125520143/20480 \\
31 & 19 & −9867816803/20480 \\
32 & 20 & −2610113463/20480 \\
\end{longtable}

In particular all eight full sectors are less than2\^{}(-127446), and

\[
                 \sum_{s=25}^{32}|\text{whole sector}_s|
                  <8\,2^{-127446}=2^{-127443}.
 \tag{PI17}
\]

There is no further sum over k, r, denominator sizes or sign profiles. The next substitution s=33 has upper exponent4647589877/20480\textgreater0. That is a failure of this sufficient estimate, not evidence that a term exists or has a large contribution.

\subsection{5. Scope, chronology and unchanged geometric data}\label{scope-chronology-and-unchanged-geometric-data}

Append127443 once to the actual chronological all-top array at checked integration and preserve every earlier error, including DO099's transverse interval and DO100's distinct high-ideal rank-three sector. The companion pins the actual canonical DO10086/31 predecessor, checks both full ordered sums and the staged87/31 lists, and records the unchanged complete floors25/11. Any intervening independently checked charge is retained by an explicit canonical-order supplement. This theorem exhausts only the remaining globally nonunipotent rank-two scalar-cross branch at supports25..32. It gives neither full support25 exhaustion nor a global-unipotent or rank-three theorem.

Directly read interfaces are the full support-independent DO076(RS3--6) reciprocal-line proof; DO095(RA21--23) and K070(PT5--8) for the actual doubled G valuation and simple coefficient poles; all K096(EP1--16); all DO096(RP1--21); DO099's transverse interval; literal DO093(IV4--15) normalized inversion; the paired pencil inclusions LIFT44(PR4--7); and the complete original R059(GB7--21) count. The new content is the uniform23\textless=r\textless=s sign classification, the bound s−12 including every active lambda-unit place, and the full interval substitution with its exact union charge. No settled primary theorem is re-proved or extended beyond its checked interface.

The original Hermitian modules over Spec(O\_E), paired ideals over Spec(O\_F), polarizations, Picard data, theta line and biextensions are unchanged. Primitive integer extraction and the actual source/target inclusions retain all local denominator information. The contextual anchors remain SGA3 I4.7.3/4.7.3.1, SGA6 VI Theorem3.3, FGA232 Theorem3.1 pp148--149, FGA236 Corollaries1.7/2.3 pp229/232, SGA7 I VII2.7/2.9 pp169--174 and VIII7.4(b) pp310--311, and the already pinned cote4/cote54 locators. No new K-theory, l-adic, motivic or source-constructor statement follows.

The signed sufficient inequality, prescribed global Bessel/Mram period, selected arithmetic-image content, literal source construction, Read/STiesR/specific032 pairing and changed-family DP44 integer comparison remain open. The exact certificate supplements the complete proof by checking every sign/pole case, all actual prime powers, all eight complete endpoints and the historical error arrays; it does not enumerate rational unitary arrows or establish a global period from finite diagnostics.

\clearpage

\section{AM-LIFT-078 --- averaged content ideals over the actual real field}\label{am-lift-078-averaged-content-ideals-over-the-actual-real-field}

\textbf{Canonical checked integration --- 2026-11-23T00:03:12.610696+00:00.} Root and independent whole-proof/script/interface readings and separate4,553-assertion replays passed without correction. For every fractional ideal of the actual F and bounded nonzero real element, the entire original conjugate-content ideal sum is\textless2\^{}1065 y\^{}(21021/125). Ordered divisor ideals and full real subideal packing yield3\^{}f y\^{}(f(1+delta)) zeta\_F(1+delta)\^{}2. Atdelta1/1000 the actual imprimitive-character Euler product and all ramified factors give zeta\_F(1.001)\textless3003·5\^{}167 without GRH or a class-number estimate. Primitive v congruences also retain the exact original q affine ideal and both true eta-metric q/r packing bounds. The lemma has no sector charge, floor change or R080 dependency; the canonical87/31 error arrays remain unchanged pending a separately checked full application.

\textbf{Algebraic-geometric translation.} The count lives on fractional invertible modules and divisor ideals over Spec(O\_F), extended to the original Spec(O\_E) Hermitian modules. Actual ideal classes, the CM involution and the polarization are retained. It supplies an unconditional counting constant in the absolutely convergent range, not a motivic or critical-line L-function comparison. Exact carried contextual anchors are SGA3 I4.7.3/4.7.3.1; SGA6 VI Theorem3.3; FGA232 Theorem3.1 pp148--149; FGA236 Corollaries1.7/2.3 pp229/232; SGA7 I VII2.7/2.9 pp169--174 and VIII7.4(b) pp310--311; cote4 scan003 §§8.1.1--8.1.2, scans005--006 §§8.4.1--8.4.3 and cote54 scans005--007. No new lambda-ring/K-theory identity, l-adic sheaf or weight theorem, CM-motive identification or arithmetic source constructor follows. DP\_MULT Q3e and THETA3 G1/A7 await its separately counted application; CX\_GRH CX-ALG-024 and GCT\_PLETH T-052 retain their existing gates.

Author: /root/lift\_accounting, GPT-6 Astra Pro, ALP\_MACHINE. 22 November 2026 UTC. Controlled pending proof. Author whole-file verification, exact diagnostics and independent review are recorded separately. This proof is independent of the prospective R080 application.

\subsection{1. Actual field and complete joint count}\label{actual-field-and-complete-joint-count}

Retain the unchanged fields of K032 and DO091(TD11): \[
 E=\mathbf Q(\zeta_{28})B_7(29)\mathbf Q(\sqrt5,\sqrt{13}),
 \qquad F=E^+,\qquad [F:\mathbf Q]=f=168.                 \tag{AV1}
\] The degree-seven field \(B_7(29)\) is defined by the subgroup \(\{1,12,17,28\}\) of the 29th cyclotomic Galois group. The real field is the compositum of \(\mathbf Q(\zeta_{28})^+\), \(B_7(29)\), \(\mathbf Q(\sqrt5)\), and \(\mathbf Q(\sqrt{13})\). Thus it is an abelian subfield of \[
                \mathbf Q(\zeta_m),\qquad m=52780=28\cdot29\cdot5\cdot13.
                                                               \tag{AV2}
\] Also \(i\in\mathcal O_E^\times\). These are existing actual field data, not a new extension or coefficient field.

Let \(\Lambda\) be any nonzero fractional F-ideal, with \(N_F\Lambda=z^f\), \(z>0\). Set \(\mathfrak l=\Lambda\mathcal O_E\), \(R>0\), and \(y=R/z\). For every nonzero \(t\in\Lambda\) with \(|\tau(t)|\le R\) at all real embeddings, put \[
 \mathfrak A_F=(t)\Lambda^{-1},\qquad
 \mathfrak A=\mathfrak A_F\mathcal O_E.
\] These are integral ideals. Let \(c(t)\) count the integral E-ideals \(\mathfrak C\) satisfying \[
 \mathfrak A\subseteq\mathfrak C\subseteq\mathcal O_E,
       \qquad\mathfrak C\overline{\mathfrak C}\subseteq\mathfrak A.
                                                               \tag{AV3}
\] Then \[
 \boxed{\displaystyle
 \sum_{\substack{0\ne t\in\Lambda\\|\tau(t)|\le R\ \forall\tau}}c(t)
      <2^{1065}y^{21021/125}.}                               \tag{AV4}
\] No principality, lattice-shape, equidistribution or class-number premise is used.

In the actual rank-one pencil problem take \(t=-i\eta\), where \(\eta\ne0\) is the pure-imaginary negative entry, and take \(\Lambda\mathcal O_E=\mathfrak l\) of DO092(SN3). The actual content ideal \(\mathfrak b\) gives \(\mathfrak C=\mathfrak b\mathfrak l^{-1}\), and SN4--5 give AV3. Thus AV4 bounds the complete joint choice of \(\eta,\mathfrak b\), replacing their separate worst-case product. All other coordinate, incidence, trace and source factors remain to be counted.

\subsection{2. Content choices and ordered ideal pairs}\label{content-choices-and-ordered-ideal-pairs}

At an F-prime of exponent \(e\) in \(\mathfrak A_F\), the split case has exponents \(u,v\) in \(\mathfrak C\) with \(0\le u,v\le e\) and \(u+v\ge e\). Its exact count is \((e+1)(e+2)/2\). In the inert case the count is \(\lfloor e/2\rfloor+1\); in the ramified case it is \(e+1\), since the exponent of the extended ideal is \(2e\). All three follow by taking valuations in AV3, and all are bounded by the split count. Therefore \[
 \begin{split}
 c(t)&\le d_{3,F}(\mathfrak A_F)
 :=\prod_{\mathfrak p^e\Vert\mathfrak A_F}\frac{(e+1)(e+2)}2\\
 &=\#\{(\mathfrak K,\mathfrak L):
             \mathfrak K\mathfrak L\mid\mathfrak A_F\}.
 \end{split}                                                \tag{AV5}
\] The two ideals are ordered, nonzero and integral. The last equality counts their nonnegative local exponents with sum at most \(e\). No ideal generator is chosen.

Put \(X=y^f\). A nonzero \(t\) contributing to a pair satisfies \[
 t\in\Lambda\mathfrak K\mathfrak L,\qquad
 n:=N_F(\mathfrak K\mathfrak L)
 \le N_F\mathfrak A_F
 =|N_{F/\mathbf Q}t|/(N_F\Lambda)\le X.                       \tag{AV6}
\] Thus the original sum is empty if \(y<1\). The cutoff \(n\le X\) follows from a nonzero element; it is not inferred from counting zero in infinitely many ideals.

\subsection{3. Every subideal has a full packing bound}\label{every-subideal-has-a-full-packing-bound}

Use the real trace metric \(\|x\|^2=\sum_\tau|\tau(x)|^2\). If \(0\ne x\in\Lambda\mathfrak K\mathfrak L\), the ideal \((x)(\Lambda\mathfrak K\mathfrak L)^{-1}\) is integral. Hence \[
 |N_{F/\mathbf Q}x|\ge z^f n,\qquad
 \|x\|^2\ge f|N_{F/\mathbf Q}x|^{2/f}
               \ge fz^2n^{2/f}.                            \tag{AV7}
\] This is the real F norm and metric, with no complex half-trace factor.

The embedding box lies in the Euclidean ball of radius \(\sqrt fR\). Disjoint open balls of half the lower minimum \(\sqrt fz n^{1/f}\), centered at its lattice points, lie in the correspondingly enlarged ball. Their full volumes give \[
 \#\{t\in\Lambda\mathfrak K\mathfrak L:|\tau(t)|\le R\}
 \le(1+2y/n^{1/f})^f\le3^fy^f/n\quad(n\le y^f).             \tag{AV8}
\] Including zero only enlarges this finite count. This proof applies to every fractional subideal and every ideal class. It uses no covering-radius or covolume-density approximation.

For every \(\delta>0\), the nonnegative finite sum over AV6 is consequently bounded by \[
 \begin{split}
 \sum_t c(t)
 &\le3^fy^f\sum_{N_F(\mathfrak K\mathfrak L)\le y^f}
                 (N_F\mathfrak K\,N_F\mathfrak L)^{-1}\\
 &\le3^f y^{f(1+\delta)}\zeta_F(1+\delta)^2.
 \end{split}                                                \tag{AV9}
\] Indeed \(n^{-1}\le y^{f\delta}n^{-1-\delta}\), and the full absolutely convergent ideal-pair series is the square of the actual Dedekind zeta series. There is no uncounted class or unit factor.

\subsection{4. Actual abelian Euler factors}\label{actual-abelian-euler-factors}

Set \[
                   \delta=1/1000,\qquad \sigma=1001/1000.
                                                               \tag{AV10}
\] Inflate the 168 characters of \(\Gamma=\operatorname{Gal}(F/\mathbf Q)\) to \((\mathbf Z/m\mathbf Z)^\times\), and extend them by zero to nonunits. Use the imprimitive functions \[
 L_m(\sigma,\chi)=\sum_{n\ge1}\chi(n)n^{-\sigma}
                =\prod_{p\nmid m}(1-\chi(p)p^{-\sigma})^{-1}.
\] For \(p\nmid m\), Frobenius in \(\Gamma\) has order equal to its residue degree \(d_p\). Its character values run through every \(d_p\)-th root of unity exactly \(168/d_p\) times. The absolutely convergent Euler factors therefore give \[
 \prod_{\chi\in\widehat\Gamma}L_m(\sigma,\chi)
 =\prod_{p\nmid m}(1-p^{-d_p\sigma})^{-168/d_p}
 =\zeta_F^{(m)}(\sigma).                                    \tag{AV11}
\] This derives the required factorization directly from the actual cyclotomic Frobenius. It does not replace imprimitive functions by primitive ones and discard their omitted factors.

\subsection{5. Uniform nonprincipal bound}\label{uniform-nonprincipal-bound}

For each nonprincipal inflated character its full unit-residue sum is zero. Retain the first period. For every \(k\ge1\), subtract \((km+m)^{-\sigma}\) from all terms of that period; the total subtracted character contribution is zero. For \(1\le a\le m\), \[
 0\le(km+a)^{-\sigma}-(km+m)^{-\sigma}
       \le\sigma(m-a)(km)^{-\sigma-1}.
\] With \(\phi=\varphi(m)=16128\), this proves \[
 |L_m(\sigma,\chi)|
 \le\sum_{\substack{1\le a\le m\\(a,m)=1}}a^{-1}
            +\sigma\zeta(\sigma+1)\phi m^{-\sigma}.          \tag{AV12}
\] The first-period bound used \(a^{-\sigma}\le a^{-1}\). Also \(\zeta(\sigma+1)\le\zeta(2)<7/4\), and \((1001/1000)(7/4)<2\). The bound \(\zeta(2)<7/4\) follows by retaining its first two terms and bounding the remainder by \(\int_2^\infty x^{-2}dx\). Thus the second term of AV12 is less than \(2\phi/m\).

The exact finite rational comparisons are \[
 \sum_{\substack{1\le a\le52780\\(a,52780)=1}}a^{-1}
       <81/20,\qquad 2\cdot16128/52780<31/50.                \tag{AV13}
\] The companion verifies them directly, not by floating point. Every nonprincipal factor is therefore below \(467/100<5\) in absolute value. The principal factor satisfies \[
 0<L_m(\sigma,1)\le\zeta(\sigma)
       <1+\int_1^\infty x^{-\sigma}dx=1001.                 \tag{AV14}
\] This applies to every nonprincipal quotient character, whether complex, real, or of smaller conductor.

\subsection{6. All omitted local factors are retained}\label{all-omitted-local-factors-are-retained}

The omitted primes are exactly 2, 5, 7, 13, and 29. The actual field decomposition of K032/DO091(TD11) gives:

\begin{itemize}
\tightlist
\item
  At 2, its residue degrees in the real degree-three seventh-cyclotomic subfield and \(B_7(29)\) are 3 and 7. Thus the F residue degree is at least 21, and the number of primes is at most 8.
\item
  At 5, 7, and 13, the field \(B_7(29)\) is unramified and has residue degree 7, since these residues are outside \(\{1,12,17,28\}\). Thus each F residue degree is at least 7 and the number of primes is at most 24.
\item
  At 29 the degree-seven factor is totally ramified and the other three factors split completely. There are exactly 24 primes of norm 29.
\end{itemize}

Since their positive local factors decrease with \(\sigma>1\), their full product is less than \[
 \left(\frac{2^{21}}{2^{21}-1}\right)^8
 \prod_{p=5,7,13}\left(\frac{p^7}{p^7-1}\right)^{24}
       \left(\frac{29}{28}\right)^{24}<3.                   \tag{AV15}
\] The last inequality is exact rational arithmetic. Larger residue degrees or ramification indices only reduce this containing upper bound. Equations AV11--15 prove \[
              \boxed{\zeta_F(1001/1000)<3003\cdot5^{167}.}   \tag{AV16}
\] There is no residue-at-one, regulator, class-number or prime-distribution estimate.

Finally, \[
       3^{168}(3003\cdot5^{167})^2<2^{1065},
       \qquad168(1+1/1000)=21021/125,                       \tag{AV17}
\] by exact integers. Combining AV9--17 proves AV4.

\subsection{7. The primitive scalar keeps its original affine ideal}\label{the-primitive-scalar-keeps-its-original-affine-ideal}

In the actual reduced rank-one sector write \[
 x=(v_1q,v_2q,v_3q,r,\eta),\qquad
 v=(2,95192390527,383092671043),\qquad V=v\cdot v,
\] with \(Vq\bar q+r\bar r=\eta\bar\eta\). The integer vector \(v\) is primitive: its first entry 2 and odd second entry already generate 1. Thus the original coordinate-content ideal is precisely \(\mathfrak b=(q,r,\eta)\).

For fixed \(\eta,\mathfrak b\), support, marked double-pole set and actual labels, K077(AP7--11) puts every positive coordinate in its specified affine coset of \[
 \mathfrak c=\mathfrak b\cap\mathfrak J I_\Sigma,\qquad
                N_E\mathfrak c=N_E\mathfrak b\,T^{2d}.
\] For two allowed scalar values, \(v_j(q-q')\in\mathfrak c\) for all three \(j\). Integer Bezout gives \(q-q'\in\mathfrak c\). Conversely adding an element of \(\mathfrak c\) preserves these three affine conditions and membership in \(\mathfrak b\). Hence their nonempty containing set is an affine \(\mathfrak c\)-coset. Further content or norm constraints can reduce it. Bezout is used only for ideal membership, not for an archimedean norm estimate.

In the exact metric \(\|w\|_\eta^2=\sum_\tau|w_\tau|^2/|\eta_\tau|\), K077(AP14) gives nonzero-difference minimum squared at least \(fzT^{2d/f}\). The norm equation gives squared radii \(fR/V\) for \(q\) and \(fR\) for \(r\). Thus, putting \(y_c=R/(zT^{2d/f})\), \[
 \#q\le(1+2\sqrt{y_c/V})^{2f},\qquad
 \#r\le(1+2\sqrt{y_c})^{2f}.                                \tag{AV18}
\] These are affine-lattice bounds for the unchanged \(\mathfrak b\). If they replace the relative norm fibre count, there is no additional relative-divisor or norm-one-generator multiplier. The original norm equation may be discarded in this upper bound.

Combining AV4 and AV18 gives the full joint \(\eta,\mathfrak b,q,r\) containing-set bound at fixed actual pattern and pencil data: \[
 2^{1065}y^{f(1+1/1000)}
       (1+2\sqrt{y_c})^{2f}(1+2\sqrt{y_c/V})^{2f}.           \tag{AV19}
\] Exact content, exact marked poles and all remaining matrix equations can only reduce this count.

Multiplication by the existing Green and actual selected-trace factors leaves variable part \[
 z^\kappa(\sqrt z+A)^{2f}(\sqrt z+A/\sqrt V)^{2f},
 \quad A=2\sqrt R\,T^{-d/f},\quad
 \kappa=(a-4-1/1000)f>163f\quad(a=f=168).                    \tag{AV20}
\] It is increasing, and its ratio at \(0<z'\le z\) is at most \((z'/z)^\kappa\). This fits the established DO088/DO092 weighted-pattern comparison. It assumes neither that a cap is realized nor that a one-point criterion at the cap applies to nonmaximal patterns. Complete support/type/pair summation and a numerical sector charge remain separate work.

\subsection{8. Sources, verification and limits}\label{sources-verification-and-limits}

Direct internal inputs are K032(QH9--15), DO091(TD11), DO092(SN1--8), K075(RN11--13), K077(AP7--14), and the DO088/DO092 weighted-pattern interfaces. Their actual scopes are retained. The abelian Euler-product identity and character estimates used here are explicitly proved in sections4--6; the convergent value at1.001 is a counting constant, not a central-value or nonvanishing theorem.

The certificate checks the literal modulus, field-degree arithmetic, defining residue-degree tests, exact harmonic sum, all ramified Euler factors, final integer constant, local ideal-choice identities, primitive Bezout data and exponent identities. It does not replace the ideal-pair rearrangement, packing proof or character-block argument. Whole saved-file reading and independent mathematical review are separate from its arithmetic replay.

The original field, ideals, levels, characters, metrics, coefficient field and source data remain fixed. This lemma gives no complete support floor or sector charge by itself, no central L-value or prescribed period nonvanishing, no admissible mixed-CM source, no selected arithmetic-image content, no specific032 matching/value, no complete DP44 inequality, and no RH conclusion.

\clearpage

\section{AM-R-080 --- averaged content and both affine coordinates remove rank-one support25}\label{am-r-080-averaged-content-and-both-affine-coordinates-remove-rank-one-support25}

\textbf{Canonical checked integration --- 2026-11-23T00:07:03.395697+00:00.} Root author and independent whole-proof/script/interface readings and separate638-assertion final replays against actual DO101 passed without correction. LIFT078 has its own independent whole review and canonical integration before this application. Every original fixed-S globally rank-one-unipotent support25 term is removed with one537 charge. The primitive q coordinate retains the original affine difference ideal, so both q/r true eta-metric element counts replace the older relative-norm fibre. Independently checked LIFT078 averages every conjugate-content ideal over eta via actual subideal packing and the unconditional actual-field bound zeta\_F(1.001)\textless3003·5\^{}167. Its full constant2\^{}1065 and exponent21021/125 are retained before all Green patterns, selected traces, supports, marked subsets and original paired type/label sums. All26 double-pole cells pass; worst margin538, next1530, complete witness exponent -25183329/64. Actual88/31 errors remain below2\^{}-230 with complete floors25/11. This is a full rank-one sector result, not whole support25 or signed/global/source/selected-image/specific032/DP44/RH nonvanishing.

Author: /root, GPT-6 Astra Pro; inherited effort not externally exposed; ALP\_MACHINE. 22 November2026 UTC. Controlled full proof. Whole saved-file review, exact computation and independent review are recorded separately before promotion.

\subsection{1. Actual scope and theorem}\label{actual-scope-and-theorem-3}

Retain the unchanged original J56 all-top kernel, H=133, f=a=168, actual field, paired source/target pencils, admitted labels, finite projectors, selected Hilbert metric, favorable-point transport, central representative, phases, masses and Haar measures. Consider every globally rank-one unipotent matrix fixed pointwise on the original S, with full support s=25. Let d be the intrinsic number of actual double poles of eta=(h-I)\_nn; K077(AP6) proves that these are paired type-II places and0\textless=d\textless=25.

\textbf{Theorem.} This entire original fixed-S globally rank-one-unipotent support25 sector has total absolute normalized contribution less than2\^{}-537, including every intrinsic double-pole set, every incidence profile and all actual paired labels.

The proof combines both actual affine coordinate conditions with LIFT078's averaged count of the original content ideals over eta. It makes no full support25 or signed-nonvanishing assertion by itself. Complete floors remain25/11 until every other structural class is exhausted. Global unipotence and rank one are premises.

\subsection{2. Carried normal form and the actual affine q coordinate}\label{carried-normal-form-and-the-actual-affine-q-coordinate}

The support-independent reduction DO091(TD2--9) and DO092(SN1--3) give either an integral witness or \[
 N=h-I=x\bar x^{\mathsf t}J_0/\eta,\qquad
 x=(v q,r,\eta),\qquad
 Vq\bar q+r\bar r=\eta\bar\eta,\qquad \bar\eta=-\eta\ne0,
                                                        \tag{AQ1}
\] where v=(2,95192390527,383092671043) is the fixed primitive integer vector, \[
 V=v\cdot v=155821585821105090045582,\qquad
 q\ne0,\qquad|\eta_\tau|\le R=8844/133.
                                                        \tag{AQ2}
\] The full original negative column determines N, so no source-isometry fibre remains after its actual coordinates are counted.

At a fixed support, intrinsic double-pole set D, actual original labelled pairs and paired rank-zero G ideal J, retain \[
 \mathfrak l=J I_\Sigma I_D=\Lambda_F\mathcal O_E,\quad
 N_E\mathfrak l=z^{2f},\quad
 z=(N_EJ)^{1/(2f)}T^{-(s+d)/f},\quad N_EJ\le M^2,\quad M=Q^{167}.
                                                        \tag{AQ3}
\] The original content ideal is b=(q,r,eta). It contains (eta) and lies in l. The nonempty reduced stratum satisfies y=R/z\textgreater=V, and \[
 N_Eb\ge |N_E\eta|^{1/2}z^f.
                                                        \tag{AQ4}
\] All these statements use the actual paired original entries; no common stable lattice or principal ideal generator is posited.

Fix eta,b and the original pencil data. K077(AP7--11) puts each positive coordinate x\_j into a specified affine coset of \[
 \mathfrak c=b\cap JI_\Sigma,\qquad N_Ec=N_Eb\,T^{2d}.
                                                        \tag{AQ5}
\] For r this is already the bound used in DO092(SN9). It applies to the first three coordinates v\_jq as well. If q and q' both satisfy these actual conditions, then v\_j(q-q') lies in c for all three indices. The integer identity \[
 2\,\frac{1-95192390527}{2}+95192390527=1
                                                        \tag{AQ6}
\] therefore gives q-q' in c.~Thus all possible q lie in one affine c-coset whenever there is a solution. This is a purely algebraic difference argument; the large Bezout coefficient introduces no archimedean norm loss. Fixing r or imposing AQ1 only cuts down this containing coset.

Use the original positive eta-weighted metric \[
 \|w\|_\eta^2=\sum_\tau |w_\tau|^2/|\eta_\tau|.
                                                        \tag{AQ7}
\] AQ1--2 give squared radii fR for r and fR/V for q. For any nonzero difference in c, the arithmetic-geometric mean and AQ4--5 give \[
 \|w\|_\eta^2\ge
 f\left(N_Ec/|N_E\eta|^{1/2}\right)^{1/f}
 \ge f zT^{2d/f}.                                        \tag{AQ8}
\] Disjoint Euclidean balls in the true real2f-dimensional coordinate metric therefore bound the two actual coordinate counts by \[
 \#r\le(1+2\sqrt{y_c})^{2f},\qquad
 \#q\le(1+2\sqrt{y_c/V})^{2f},\qquad
 y_c=R/(zT^{2d/f}).                                      \tag{AQ9}
\] Both counts are uniform in the affine shifts and all actual labelled pairs. Multiplying them bounds the joint possibilities; the norm equation only reduces that product. Zero q may be counted in this upper bound, but q≠0 remains the premise used for the actual nonemptiness inequality y\textgreater=V. No relative ideal or norm-one generator multiplicity is added: AQ9 counts the actual elements q themselves.

\subsection{3. Average the actual content ideals over eta before counting coordinates}\label{average-the-actual-content-ideals-over-eta-before-counting-coordinates}

For each actual eta, let A\_F=(-i eta)Lambda\_F\^{}(-1), an integral F-ideal. DO092(SN4--6) bounds the number of its original coordinate-content ideals b locally by T\_e=(e+1)(e+2)/2 at a split F-prime, floor(e/2)+1 at an inert prime and e+1 at a ramified prime. Each is bounded by the local coefficient of d3\_F(A\_F), where \[
 d3_F(A_F)=\#\{(K,L):KL\mid A_F\}.
\] The ideals and their classes are not replaced by principal generators.

Apply the independent averaged-content theorem AM-LIFT-078 to the actual real ideal Lambda\_F of norm z\^{}f and the bound\textbar eta\_tau\textbar\textless=R. Its complete conclusion at delta=1/1000 is \[
 \sum_{\substack{0\ne\eta\in i\Lambda_F\\|\eta_\tau|\le R}}
          \#\{\hbox{eligible original content ideals }b\}
 <2^{1065}y^{1001f/1000}
 =2^{1065}y^{21021/125},\qquad y=R/z.                     \tag{AQ10}
\] Here further actual pencil or norm conditions only reduce the left side. For specificity, the proof of LIFT078 sums over K,L with n=N\_F(KL)\textless=y\^{}f.~The real lattice Lambda\_F KL has minimum at least sqrt(f)z n\^{}(1/f). Full disjoint-ball packing bounds its nonzero eta count by3\^{}f y\^{}f/n.~Summing with n\^{}delta\textless=y\^{}(f delta) gives3\^{}f y\^{}(f(1+delta)) zeta\_F(1+delta)\^{}2. This is a lattice bound for every ideal, including coarse ones, not an expected-density replacement.

The adopted real field has degree168 and lies in Q(zeta\_52780). The elementary imprimitive-character and actual ramified-Euler calculation in LIFT078 gives \[
 \zeta_F(1001/1000)<3003\cdot5^{167},\qquad
 3^{168}(3003\cdot5^{167})^2<2^{1065}.                    \tag{AQ11}
\] It uses only the finite reduced harmonic sum, nonprincipal character-period cancellation and absolutely convergent Euler products at a real point greater than one. It assumes no GRH, no L-value nonvanishing on a critical line, and no principal ideal class. The exact finite arithmetic and complete proof are part of the separate LIFT078 record. In particular AQ10 replaces the combined old eta-lattice/content estimate. It is not multiplied by that old estimate a second time.

The uniform coordinate counts AQ9 may be pulled outside this sum because their lower difference bound depends only on z,d,R,V, not the particular eta,b or affine shift. Counting all q,r in those two containing cosets still bounds every actual column and hence every rank-one matrix.

The genuine radial-density refinement DO092(SN16--19), including its exact angular integral and multiplicative near-point transport, gives the actual selected normalized trace \[
                    |c(h)|<2\,2^{-931f/2-630}z^{-f}.      \tag{AQ12}
\] Both compact factors, all word/twist parameters, the original selected metric and the favorable-point argument remain. No operator-norm substitution or stronger pointwise estimate is inferred.

\subsection{4. Complete weighted finite patterns and actual incidence sum}\label{complete-weighted-finite-patterns-and-actual-incidence-sum}

The finite bound remains2(N\_EJ)\^{}(a/2)M\^{}(-a). Multiplying it by AQ9--12 leaves, apart from positive constants depending on fixed s,d, \[
 F(z)=z^\kappa(\sqrt z+A)^{2f}(\sqrt z+A/\sqrt V)^{2f},
 \quad A=2\sqrt R\,T^{-d/f},\quad
 \kappa=(a-4-1/1000)f=3443979/125>163f.                  \tag{AQ13}
\] This is the full function after the average, trace and both coordinate packings have been included. It increases for every z\textgreater0 and satisfies F(z')/F(z)\textless=(z'/z)\^{}kappa whenever0\textless z'\textless=z.

Let z\_sd=(M/T\^{}(s+d))\^{}(1/f) be the full rank-zero cap. Omitting a paired rank-zero G prime of norm q lowers z by q\^{}(1/f). Hence the complete weighted sum is \[
 \sum_J F(z_J)\le F(z_{sd})\prod_v(1+q_v^{-\kappa/f})
 <2F(z_{sd}),                                            \tag{AQ14}
\] because every q\_v\textgreater2\^{}34 and(1+2\^{}-5542)\^{}334\textless2. All334 rank-zero choices are included; no unweighted2\^{}334 factor is silently removed. The cap need not itself satisfy the original reduced nonemptiness y\textgreater=V: the displayed positive increasing continuation still bounds every feasible actual stratum. We do not infer uniqueness of every actual q fibre by checking it only at the cap.

K077(AP18) retains every original ordered source/target pair and its true incidence weight with factor binom(f,s)binom(s,d)(T+1)\^{}(s-d). The CM partner is paired data, not an independent label. Put \[
 y_{sd}=R/z_{sd},\qquad y'_{sd}=R(T^{s-d}/M)^{1/f}.
\] At z\_sd the finite numerator cancels M\^{}a. Thus the complete d-cell bound at s25 is \[
\begin{split}
 U_{25,d}<&\;2^{3+1065-931f/2-630}
       \binom{168}{25}\binom{25}{d}(T+1)^{25-d}z_{25,d}^{-f}\\
 &\quad\cdot y_{25,d}^{21021/125}
       (1+2\sqrt{y'_{25,d}})^{2f}
       (1+2\sqrt{y'_{25,d}/V})^{2f}. 
\end{split}                                             \tag{AQ15}
\] The factor2\^{}3 consists of the finite Euler correction, complete weighted G sum and actual near-point product, each less than2. The averaged constant2\^{}1065 already counts eta and every eligible content ideal. Both affine coordinate factors count actual elements; the older28·C\_rel norm-fibre factor is replaced by this genuine count. No additional content, generator, stabilizer or class multiplier is needed.

\subsection{5. Every endpoint and the complete sector charge}\label{every-endpoint-and-the-complete-sector-charge}

The companion computes AQ15 at all26 values of d by exact rational f-th-root and square-root brackets, followed by directed rational logarithms. Each full logarithmic upper bound is strictly below the negative of its displayed margin:

\begin{longtable}[]{@{}rr@{}}
\toprule\noalign{}
d & Inverse binary margin \\
\midrule\noalign{}
\endhead
\bottomrule\noalign{}
\endlastfoot
0 & 9367 \\
1 & 10244 \\
2 & 10768 \\
3 & 10907 \\
4 & 10925 \\
5 & 10925 \\
6 & 10922 \\
7 & 10920 \\
8 & 10918 \\
9 & 10916 \\
10 & 10914 \\
11 & 10912 \\
12 & 10909 \\
13 & 10891 \\
14 & 10779 \\
15 & 10313 \\
16 & 9460 \\
17 & 8486 \\
18 & 7494 \\
19 & 6499 \\
20 & 5505 \\
21 & 4510 \\
22 & 3516 \\
23 & 2523 \\
24 & 1530 \\
25 & 538 \\
\end{longtable}

The worst cell is d25 with upper exponent strictly between−539 and−538; every other cell has margin at least1530. The complete independent integral-witness branch in DO091(TD10) has \[
 E_{\rm wit}(25)=1+168+107+5757+46af
 +(10L_T+997)25-a167L_Q=-25183329/64<-390000,              \tag{AQ16}
\] where L\_T=1019333/1024 and L\_Q=287843/4096 have their exact established directions. Its full original five-column/source count includes every type and zero label. The support-dependent formula has the same local premises here; no extension of an old numerical endpoint is assumed.

The entire support25 sector therefore has absolute contribution below \[
                  2^{-538}+25\,2^{-1530}+2^{-390000}
                  <2^{-537}.                            \tag{AQ17}
\] This includes the entire witness branch and every reduced d-cell. The intrinsic d partition, complete marked subsets, original type sums and all paired labels were included before the final union. There is no remaining coordinate, eta, content or norm-range union. The theorem follows.

The prior unpromoted q-packing draft still used a worst-case content bound and left d24/25 with positive sufficient expressions. Its files are retained as research history; those positive bounds were never actual contributions or a canonical sector claim. AQ10 is the extra proved input that closes both remaining cells.

\subsection{6. Computation, chronology and geometric scope}\label{computation-chronology-and-geometric-scope}

The logarithm certificate uses the same elementary finite atanh series with a rational positive-tail bound as K071(NB23) and DO092's executable certificate. Both f-th roots and square roots are bounded by exact rational/integer power comparisons; every reported cell has a certified unit interval. LIFT078's averaged estimate and field constant are checked independently. These calculations are arithmetic diagnostics of the written proof, not an enumeration of rational arrows.

At canonical integration append537 once to the actual ordered all-top error array and retain every earlier charge and the full independent uniform array. The companion records its predecessor and checks the full exact sums below2\^{}-230; integration rechecks the then-actual chronology. The unique231 charge and next minimum358 remain. Complete support floors stay25/11 until a separate exhaustion of all other structural classes is proved.

The affine condition lives in the original Hom module over Spec(O\_E) and its fractional ideal, with both CM partners and exact content kept. The averaging sums actual ideal divisibility lattices rather than changing the coefficient field. The true eta-dependent positive metric is used only for counting actual differences. The polarization, Picard class, biextension and selected complex realization are unchanged. Exact carried contextual anchors are SGA3 I4.7.3/4.7.3.1; SGA6 VI Theorem3.3; FGA232 Theorem3.1 pp148--149; FGA236 Corollaries1.7/2.3 pp229/232; SGA7 I VII2.7/2.9 pp169--174 and VIII7.4(b) pp310--311; cote4 scan003 §§8.1.1--8.1.2, scans005--006 §§8.4.1--8.4.3 and cote54 scans005--007.

No new lambda-ring identity, K-theory map, l-adic weight statement, motivic projector or CM-motive comparison follows. DP\_MULT Q3e and THETA3 G1/A7 gain precisely the stated smaller actual finite/real remainder. Literal mixed-CM admissible source construction, selected arithmetic-image content, prescribed global periods and the specific032 pairing/value, every actual DP44 term, CX\_GRH CX-ALG-024 and the global GCT\_PLETH T-052 comparison remain separate open gates. No Lean build or archive audit is performed.

\clearpage

\section{AM-DO-102 --- the non-all-II unipotent scalar-cross sector at support25}\label{am-do-102-the-non-all-ii-unipotent-scalar-cross-sector-at-support25}

\textbf{Canonical checked integration --- 2026-11-23T00:13:31.183876+00:00.} Root and independent whole-proof/script/interface readings and separate5,406-assertion replays passed without mathematical correction. Historical pending wording below is preserved. Every globally unipotent rank-two scalar-cross support25 term with at least one type-I or type-III place is removed with one881 charge. The actual full-positive normal form makes all original cross coordinates integral outside typeII. Its original central coordinate has affine difference lattice iJ0 I\_Sigma,F I\_D,F, retaining the shift, full real discriminant and750 transference coefficient. The complete u24 sector costs less than2\^{}-882, including50 type-position choices and both extra T factors. The padded u\textless=23 sector includes300 sets, all3\^{}25 type assignments and T\^{}25 pair cost, with explicitly charged cap30 removal and full sum below2\^{}-174292. Actual89/31 ordered errors remain below2\^{}-230 with complete floors25/11. The all-II support25 sector remains open; no signed/global/source/selected-image/specific032/DP44 or RH conclusion.

Author: \nolinkurl{/root/definite_oscillator}, ALP\_MACHINE. 23 November2026 UTC.

Controlled pending proof. Author whole-proof/script check, exact replay, independent review and canonical integration are distinct events.

\subsection{1. Actual sector and theorem}\label{actual-sector-and-theorem-3}

Keep the original J56 all-top normalized kernel, H=133, f=168, the actual E/F, all334 unmixed G reductions, the paired source/target pencils, the selected Hilbert metric and both compact factors, the central phase, Haar measures and class masses. At support s=25 use the DO099(JR24) residual: both cross lines are in the actual packet plane and the four short compressions vanish. The support-independent DO093(IV19--22)/R071(GR6--12) identities then give the full positive167 compression h\_ZZ=I. Assume, in addition, that h=I+N is globally unipotent of rank2, and that C=N\_PZ has nonzero restriction to the original S164. All unmixed G reductions remain unipotent. This is the scalar-cross branch, not the fixed-S branch.

Let i,u,t count actual active types I,II,III, respectively, so i+u+t=25. The theorem concerns every profile with u≤24. The all-II profile u=25 is excluded.

\textbf{Theorem.} The entire stated non-all-II support25 sector has total absolute contribution to the unchanged normalized kernel strictly less than

\[
                              2^{-881}.                       \tag{NT1}
\]

This is one whole-sector charge. The proof below retains every original support, type, pair, orientation, pivot and reconstruction factor. It does not claim a support-floor change or a bound for the all-II sector.

Write Q=1427908193729121522601, U=77Q²−14, R=17424/133, M=\sqrt{R/U}, and D0=7\^{}140 29\^{}144 5\^{}84 13\^{}84. The exact directed logarithms used below are

\[
 \log_2T<L_T=1019333/1024,\qquad T>2^{995}.                    \tag{NT2}
\]

\subsection{2. Original cross and central coordinates; type-dependent membership}\label{original-cross-and-central-coordinates-type-dependent-membership}

The actual globally unipotent normal form is DO097(UC3--8), equivalently DO098(CF4--7). Over E, with no finite integral-frame assertion, put

\[
 C=\ell r,\quad \ell=(1,\eta),\quad\eta\bar\eta=U,
 \quad v=J_Z^{-1}r^*,\quad q=rv>0,
\] \[
 N=\begin{pmatrix}0&-v\ell^*J_P\\ \ell r&d\ell\ell^*J_P\end{pmatrix},
 \quad J_P=\operatorname{diag}(U,-1),\quad d+\bar d=-q,
 \quad N^2Z=0,
 \quad \alpha=(N^2)_{nn}=Uq\le R.                            \tag{NT3}
\]

Every embedding satisfies \textbar\textbar r\textbar\textbar≤M. The original central coordinate w=N\_nn=−Ud has Re w=alpha/2 and \textbar Im w\textbar\textless H/2. The original164 values xi\_j=r(e\_j), the three fixed transverse values r(y1),r(y2),r(z0), one nonzero original pivot beta=eta xi\_j, and w reconstruct h uniquely. This is the proved CF4--7 reconstruction, not an inferred small isometry fibre. The first nonzero S pivot exists; we retain164\textless2\^{}8 choices. Denote the number of nonzero original xi\_j by k, 1≤k≤164.

Every displayed S/transverse input belongs separately to every actual source and target pencil at both CM partners. For type III, the exact arrows give Ny∈M\_d and (h\^{}-1−I)y=−Ny∈M\_c, because N²y=0. Thus Ny∈M\_c∩M\_d⊂Lambda at both partners. For type I, h preserves the common intersection module, which contains y and is contained in Lambda. Type II gives the usual simple target pole. The original integer Bezout coefficients of primitive u extract r(y), and the negative pivot is already an original coordinate. Consequently every counted C scalar is integral at types I and III and has at most a simple paired pole at type II. These are DO097(UC10--13)'s actual module arguments; no integrality of h itself is inferred.

Use one canonical whole-C orientation pattern at G, as in R072(PJ6) and DO098(CF8--10). There are at most4\^{}334=2\^{}668 such patterns. Let J be the corresponding E-ideal and J0 the F-ideal formed by the rank-zero G primes once. Then

\[
 |a(h)|\le2Q^{-335\cdot167}(N_EJ)^{335/2},\quad
 N_EJ\le Q^{334},\quad N_EJ\le Q^{167}N_FJ_0,
 \quad N_FJ_0\le Q^{167}.                                   \tag{NT4}
\]

All cross values and beta use this same J. No extra pattern is chosen for the transverse coordinates, the central congruence or the pivot.

Fix a padded set D of v active places containing every type-II place. Set

\[
 L=JI_D,\quad N_EI_D=T^{-2v},\quad z=(N_EL)^{1/(2f)}.
                                                                    \tag{NT5}
\]

All cross scalars lie in L. The nonzero-coordinate minimum gives kz²≤M². We will use v=24 if u=24, and v=23 if u≤23. These choices are overcounts of the actual C denominators and introduce no new local condition.

\subsection{3. The actual affine central fibre for the padded pole set}\label{the-actual-affine-central-fibre-for-the-padded-pole-set}

K071(NB5)'s entrywise pole table is I:1, II:2, III:1, at both CM partners. Thus for fixed C, one orientation pattern and actual local pairs, the differences of admissible w belong to the extended ideal J0 I\_Sigma,F I\_D,F O\_E. They are pure imaginary because Re w=alpha/2 is fixed by C. Since i is an integral unit and O\_E∩iF=iO\_F, the fibre is contained in one affine coset of

\[
 \Gamma=iJ_0 I_{\Sigma,F}I_{D,F},\qquad
 N_F(J_0I_{\Sigma,F}I_{D,F})=N_FJ_0 T^{-(s+v)}.              \tag{NT6}
\]

At padded type-I/III places the second permitted pole merely enlarges the difference lattice. The original w need not lie in Gamma; its affine shift, including the possible dyadic alpha/2, is retained. There is no division by2 in Gamma, no full-image projection assertion and no zero-target projection assumption.

The full real covolume and codifferent calculation is DO098(CF12--15), with s+v replacing48. Here s+v is49 or48. The exact original comparison

\[
                         750^{168}Q^{167}<T^{48}              \tag{NT7}
\]

therefore gives real covering radius less than sqrt(f) for Gamma. The true real dimension is f, covolume sqrt(D\_F) N\_F(Gamma), D\_F=2\^{}fD0. The translated cells lie in the radius68sqrt(f) ball. CF15's rational upper volume is below2\^{}479, so

\[
 \#w<2^{479}\frac{T^{s+v}}{N_FJ_0}
       \le2^{479}Q^{167}T^{s-v}z^{-2f}.                     \tag{NT8}
\]

The second inequality follows from NT4 and N\_EJ=T\^{}\{2v\}z\^{}\{2f\}. Alternatively enlarge to the ideal with J0 omitted, and use CF16's integral-coset packing to obtain

\[
                              \#w\le134^f T^{s+v}.           \tag{NT9}
\]

NT8 and NT9 are alternative counts of the same fibre. They are never multiplied.

\subsection{4. Complete factors transported from the checked central-fibre theorem}\label{complete-factors-transported-from-the-checked-central-fibre-theorem}

The following factors are the exact DO098(CF17--24) factors for an arbitrary ideal L; their proofs do not use that all active places are type II. The three transverse coordinates have count

\[
 4A_3(4U)^fM^{6f}z^{-6f},\qquad
 A_3=528^{3f}/((3f)!D_0^3)<2^{-4040}.                        \tag{NT10}
\]

This uses the actual Gram determinant4U and the fractional-ideal cover in that covector metric, not an original-ring bound for a fractional ideal. The pivot norm equation has at most

\[
                              2^{249}R^{f/16}z^{-f/8}         \tag{NT11}
\]

solutions, by the unchanged actual split-relative norm constant and all28 norm-one generators. The same actual selected trace obeys

\[
 |c(h)|\le\min(1,2^{-929f/2}U^{-f/2}z^{-f}).                \tag{NT12}
\]

Use theta=1 except at k=164, where the geometric interpolation exponent is7/8. The original joint S-row cap is

\[
 \binom{164}{k}47^k2^{2fk}(f!)^k
 \left(\frac{(87/32)M^2}{fkz^2}\right)^{fk}.                \tag{NT13}
\]

This includes the positions of the k original nonzero S entries. The remaining fixed multiplicity is

\[
 2^{845}=2^{1+668+168+8}:                                   \tag{NT14}
\]

Green factor2, all common orientations, the original support bound binom(168,25)≤2\^{}168, and all pivot choices. There is no residual source fibre because NT3's reconstruction is injective. The CF selected trace, mixed finite factors, central phase, source/target and Haar normalizations are unchanged.

For fixed actual type positions the full ordered-pair/absolute-incidence sum at one active place is bounded by T for type I, and1 for each of types II and III: K071(NB7--9). This uses the true globally unipotent type-I action and the exact type-III reciprocal weight. The two CM partners form one paired label choice. We first apply the uniform count for fixed actual pairs, and only then sum these weights. All profiles below are counted explicitly.

For comparison record the already whole-checked CF25--32 bounds, re-evaluated by the companion certificate. At s=v=24, the range1≤k≤64 uses NT8 and theta1, with endpoint z=M/sqrt(max(k,30)); every complete cell is \textless2\^{}-3927. The range65≤k≤163 uses NT9 and theta1, endpoint z=M/sqrt(k); every cell is \textless2\^{}-10020. The final k164 uses NT9 and theta7/8 and is \textless2\^{}-2880. Their complete union is \textless2\^{}-2879. These are arithmetic envelopes retaining all NT10--14 factors, not references to any previously charged set.

Before the S-row factor the z exponents in these three cases are respectively (2607/8)f, (2623/8)f and328f. After NT13 they are positive on the first two ranges and zero on the last. Hence every stated endpoint is valid by monotonicity.

\subsection{5. The two actual type ranges and their full sums}\label{the-two-actual-type-ranges-and-their-full-sums}

\subsubsection{5.1 Exactly24 type-II places}\label{exactly24-type-ii-places}

Here v=24 and there is exactly one other active place. Its position and its type have50 possibilities, less than2\^{}6. Their complete pair-weight cost is at most T. The cross ideal and all its caps are exactly those used in CF25--32, including

\[
                 z\le Q^{167/f}T^{-24/f}<M/\sqrt{30}.         \tag{NT15}
\]

Both alternative central counts NT8 and NT9 are T times their s=v=24 counterparts. No other count changes. Therefore, summing all k before applying the profile factor, the entire u24 sector is

\[
 <2^6T^2\,2^{-2879}<2^{-882},\qquad
                    2879-6-2L_T>882.                        \tag{NT16}
\]

The factor2\^{}6 includes all50 type-position possibilities; the original support factor remains in NT14. No second k union is added.

\subsubsection{5.2 At most23 type-II places}\label{at-most23-type-ii-places}

Pad the type-II set to a set D of23 active places. The number of padded sets is at most binom(25,23)=300\textless2\^{}9. Keep all3\^{}25\textless2\^{}50 type assignments, even those outside this subcase, as an upper count. The complete pair weight is at most T\^{}25. Thus a uniform per-k envelope may be multiplied by

\[
                              2^{59}T^{25}.                  \tag{NT17}
\]

The additional cap M/sqrt(30) is not asserted when v23. We use only z≤M/sqrt(k). For the first k range, the combined Green and NT8 T power is334v+s. It changes from335·24=8040 to334·23+25=7707, a loss of333 powers of T. After NT17 the saving is308 powers.

Removing the cap30 affects only k\textless30. Its endpoint multiplier is

\[
 (30/k)^{(54747-2fk)/2}\le30^{54411/2}<2^{136028}.            \tag{NT18}
\]

The last strict inequality follows by integer powers from30\textless2\^{}5. It is a conservative common bound, not a replacement of a negative logarithm by a positive one. With NT2, every one of the64 first-range cells, including every type/padding/pair factor, is therefore below

\[
 2^{-3927+136028+59-308\cdot995}=2^{-174300}.                 \tag{NT19}
\]

For65≤k≤163 the cap30 was inactive. The NT9 combined T power changes from337·24=8088 to336·23+25=7753, a saving335, or310 after the pair bound. Every one of these99 cells is below

\[
                       2^{-10020+59-310\cdot995}
                         =2^{-318411}.                       \tag{NT20}
\]

At k164 the z exponent remains zero and the same T change gives

\[
                       2^{-2880+59-310\cdot995}
                         =2^{-311271}.                       \tag{NT21}
\]

The entire u≤23 sector is consequently less than2\^{}-174292. Combining NT16 with NT19--21 proves NT1, since

\[
                   2^{-882}+2^{-174292}<2^{-881}.             \tag{NT22}
\]

Every active profile with u≤24 is in exactly one of the two displayed ranges. Padding is only an explicit overcount and is included once. The all-II profile is in neither range.

\subsection{6. Verification, accounting and exact scope}\label{verification-accounting-and-exact-scope}

The certificate re-evaluates the complete164 CF endpoints from exact factorials and integer power comparisons, checks all new central exponents, cap changes, type/pair/padding multiplicities, both real ideal-cover cases and the full union. Such diagnostics do not replace the original module, normal-form, selected-trace or relative-norm proofs above. The frozen DO098 and DO097 inputs are cited rather than rewritten or modified.

The actual predecessor is checks/r080\_joint\_remainder\_error.json, with88 all-top and31 uniform entries and complete support floors25/11. Its ordered tail retains DO10029550, DO101127443 and R080537. Append881 once, obtaining89/31 entries; retain both complete floors. The companion certificate reads the complete actual lists and verifies both resulting sums are less than2\^{}-230. Its explicit historical --prior option for DO101 does not change the default R080 chronology or omit R080 from the final verification. No error is dropped or reordered.

This proves only the stated non-all-II globally unipotent rank-two scalar-cross support25 reduction. It proves neither the all-II support25 bound nor a new complete floor. No signed-kernel inequality, prescribed global period, cycle value, actual HostFamP constructor, DP44 sign or RH conclusion is asserted.

\subsection{Canonical algebraic-geometric placement}\label{canonical-algebraic-geometric-placement}

The original Hom modules and determinantal conditions remain over Spec(O\_E), and the affine central fibre is a torsor under a fractional ideal in the fixed real trace realization. Its contraction and norm retain both CM partners; no integral moving frame is introduced. The polarization, theta line, Picard classes and biextensions remain unchanged. Exact carried contextual anchors: SGA3 I4.7.3/4.7.3.1; SGA6 VI Theorem3.3; FGA232 Theorem3.1 pp148--149; FGA236 Corollaries1.7/2.3 pp229/232; SGA7 I VII2.7/2.9 pp169--174 and VIII7.4(b) pp310--311; cote4 scan003 §§8.1.1--8.1.2, scans005--006 §§8.4.1--8.4.3; cote54 scans005--007. These retain the formal-group, Dieudonne and Cartier-duality context without a new comparison theorem. No new lambda-ring/K-theory equality, l-adic sheaf/weight theorem, CM-motive or motivic L-function assertion is inferred. DP\_MULT Q3e and THETA3 G1/A7 gain exactly the stated finite/real gate reduction; CX\_GRH CX-ALG-024 and GCT\_PLETH T-052 retain their separate arithmetic gates.

\clearpage

\section{AM-R-081 --- dense-row and common-ideal restrictions for rank three at support25}\label{am-r-081-dense-row-and-common-ideal-restrictions-for-rank-three-at-support25}

\textbf{Canonical checked integration --- 2026-11-23T00:18:41.677681+00:00.} Root author and independent whole-proof/script/interface readings and separate648-assertion replays passed without correction. Historical pending wording below is preserved. Every remaining rank-three scalar-cross support25 term now has at least123 nonzero original S coordinates and2\^{}33096 N\_EJ\textless T\^{}50 for the same deterministic common-C ideal. The original(k+3)-coordinate covector ball keeps the true Gramdet4U, arbitrary-ideal cover, exact beta norm fibre, actual selected trace, full paired plane-column U/T index and source fibre. At z*=2\^{}(-197/2), all122 low-ideal cells have maximum -7111295/4096; the complete high-ideal endpoint is -5721469/4096. Both complete unions cost one1395 charge, retaining the old overlapping DO10029550. Actual90/31 errors remain below2\^{}-230, complete floors25/11. Dense smaller-ideal rank3, all-II unipotent rank2 and signed/global/source/selected-image/specific032/DP44/RH gates remain open.

Author: \nolinkurl{/root}, GPT-6 Astra Pro; inherited effort is not separately exposed. ALP\_MACHINE,23 November2026 UTC. Controlled complete proof; whole saved-proof readings, independent review and exact certificates are recorded separately before canonical promotion.

\subsection{1. Actual sector and theorem}\label{actual-sector-and-theorem-4}

Keep the unchanged J56 all-top kernel, H=133,f=168,n=169, every actual paired source/target pencil, finite projector, selected Hilbert metric, both compact factors, central representative, class mass, phase and Haar measure. At support s=25 use the full positive-compression residual supplied by DO099 section6/JR24 and DO095 RA1--10:

\[
 E^{169}=Z\perp P,\quad Z=S_{164}\perp u^\perp,\quad
 \dim Z=167,\quad P=Eu\perp Ee_n,\qquad
 h=\begin{pmatrix}I_Z&B\\ C&D\end{pmatrix},\quad
 C|_S\ne0,\quad\operatorname{rank}_E(h-I)=3.                 \tag{DC1}
\]

Every unmixed G reduction remains unipotent; global unipotence is not assumed. In fact the rank-three alternative here is the non-real-trace alternative of DO095. Write Cx=(u lambda(x),eta lambda(x)), eta bar(eta)=U, using the original primitive u and U=77Q²−14, Q=1427908193729121522601. Put

\[
 k=\#\{j\in S:\lambda(e_j)\ne0\},\quad1\le k\le164,
 \qquad M=H/(2\sqrt U),\quad B_0=H/2.                       \tag{DC2}
\]

Use exactly the deterministic common-C ideal J of DO100 HI1--2: both CM primes at a rank-zero G place; at rank one the designated prime if C vanishes there and the conjugate otherwise; neither at rank at least two. The original factorization and paired inverse-adjoint relation guarantee this rule. Set

\[
 L=J I_\Sigma,\quad N_E I_\Sigma=T^{-50},\qquad
 z=(N_E L)^{1/(2f)},\qquad z_* =2^{-197/2}.                  \tag{DC3}
\]

\textbf{Theorem.} The complete absolute contribution of the following union of actual support25 rank-three terms is less than2\^{}-1395:

\[
          \{z\ge z_*\}\ \cup\ \{z<z_*,\ 1\le k\le122\}.
                                                                    \tag{DC4}
\]

Thus after removing its current-remainder intersection, every remaining rank-three support25 term satisfies both

\[
                k\ge123,\qquad 2^{33096}N_EJ<T^{50}.        \tag{DC5}
\]

The latter is exactly z\textless z\_*, since197f=33096. This strictly refines DO100's2\^{}32928 ideal threshold and adds the original-coordinate density requirement. It supplies no complete support-floor increase, signed positivity, actual arithmetic source or DP44 conclusion. Every earlier overlapping charge, including DO10029550, is retained.

\subsection{2. The complete original parameters and two Green bounds}\label{the-complete-original-parameters-and-two-green-bounds}

The fixed transverse inputs y1,y2,z0, their integral original-coordinate extraction and their genuine metric are R072 PJ4--9 and R076 CV3--10. The three vectors are

\[
 (-976627,-153032,38026,0),\quad
 (-567678,695570,-172838,0),\quad
 (256766,-7654233287,-30805128044,1).                         \tag{DC6}
\]

Their Gram G has detG=4U and G-I\_3 positive definite. The164 original S values, their three transverse values and beta=eta lambda(e\_j) at a nonzero original pivot determine C over E. Retain164\textless2\^{}8 pivot choices. Every displayed scalar belongs to the same original fractional ideal L at both CM partners. No rational ratio or moving frame is declared integral.

The same common ideal has at most2\^{}668 patterns and gives the weak pointwise Green estimate

\[
 |a(h)|\le2Q^{-335\cdot167}T^{335s}z^{335f},\qquad s=25.
                                                                    \tag{DC7}
\]

DO095 RA15--20 also gives the strong rank-three estimate, with no extra pattern or witness enumeration,

\[
 |a(h)|<2Q^{-501\cdot167}2^{9\cdot167f}T^{167t}
       \le2Q^{-501\cdot167}2^{9\cdot167f}T^{167s},
 \qquad t\le s=25.                                         \tag{DC8}
\]

The witness in DC8 is an original nonzero S entry of BC, equal to an actual three-minor, of conjugate modulus below2\^{}9 and only simple paired type-III poles. Thus both estimates concern the same actual normalized finite factor; the two mixed factors keep modulus bound one.

The genuine selected trace and norms remain

\[
 \|\lambda_\tau\|<M,\quad|\beta_\tau|<B_0,\qquad
 |c(h)|<2^{-929f/2}U^{-f/2}z^{-f}.                            \tag{DC9}
\]

For fixed covector and pivot, the exact norm equation beta bar(beta)=U xi\_j bar(xi\_j), with beta in L, has the complete bound

\[
 \#\beta\le28 C_F B_0^{f/6}z^{-f/6},\qquad
 C_F=\frac{2^{300}}{29^6 41^{14}2129^{14}}<2^{42}.            \tag{DC10}
\]

This is R076 CV11--13, using the actual relative ideal norm and all28 norm-one generators; it applies to nonprincipal and non-conjugation-invariant L. Fixing C does not justify dropping the rest of the reconstruction. Retain the complete original plane-column count and source fibre

\[
                    191^{4f}U^f T^{4s},\qquad G_2<2^{36}.    \tag{DC11}
\]

These are R072 PJ14--18/R076 CV14, with the true fixed U denominator, full paired target index and every target label including zero. Finally retain support choices2\^{}168 and the complete actual ordered-pair/absolute-incidence sum P\_T\^{}25\textless2\^{}(997*25). The Green factor2, common patterns, supports, pivots and fibre have fixed exponent

\[
                         1+668+168+8+36=881.                 \tag{DC12}
\]

No additional type-profile enumeration is needed: the full P\_T sum already contains every incidence type and its paired labels.

\subsection{3. The joint covector volume on a prescribed original support}\label{the-joint-covector-volume-on-a-prescribed-original-support}

Fix a k-element subset of the original164 S coordinates. The remaining S values vanish. Together with the three transverse values, the covector is represented by d=k+3 complex coordinates at each of the f labels, with true dual Gram diag(I\_k,G\^{}-1). Its norm is the restriction of the same positive covector norm in DC9. Hence the actual set is contained in the radius M sqrt(f) ball of this real2df-dimensional metric on L\^{}d.

The exact half-trace covolume is

\[
                  D_0^d z^{2df}/(4U)^f,
             \qquad D_0=7^{140}29^{144}5^{84}13^{84}.         \tag{DC13}
\]

This is the determinant calculation R076 CV8--10 with the original zero coordinates omitted; the three transverse coordinates and their determinant are retained. Since G-I\_3\textgreater0, transformation to this dual metric is a contraction of the product metric. The carried arbitrary-ideal cover is at most(2250/7)sqrt(df)z, less than4200 sqrt(f)z for d\textless=167. These are archimedean lattice statements, not integral changes of frame.

Disjoint fundamental cells inside the enlarged ball, with pi\textless22/7 and(22/7)f=528, therefore give the complete support-pattern count

\[
 \#\lambda_k\le {164\choose k}\,
 A_d(4U)^f(M+4200z)^{2df}z^{-2df},\qquad
 A_d=\frac{528^{df}}{(df)!D_0^d}.                            \tag{DC14}
\]

The binomial counts the original support subset. Allowing zero coordinates in its containing L\^{}d ball only enlarges the bound. The pivot still exists for actual covectors and is counted separately in DC12. No new k-dimensional integral moving subspace is introduced: the omitted coordinates are zero in the original orthonormal S coordinates.

The original integers give M\textgreater2\^{}-68. At every z\textless=z\_*\textless2\^{}-98,

\[
 4200z/M<4200/2^{30}<2^{-17},\qquad
                      (1+2^{-17})^{334f}<2.                 \tag{DC15}
\]

The final inequality is an exact integer-power comparison. Since2df\textless=334f, the full enlargement factor in DC14 is less than two. Thus throughout this lower range

\[
 \#\lambda_k<2{164\choose k}A_d(4U)^f M^{2df}z^{-2df}.       \tag{DC16}
\]

There is no omission of a covering term or substitution of an asymptotic density formula.

\subsection{4. The entire low-ideal sparse-row range}\label{the-entire-low-ideal-sparse-row-range}

Multiply DC7,DC9--12 and DC16, with every original factor retained. Its remaining z exponent is

\[
 \kappa_k=\bigl(335-2d-1-1/6\bigr)f
           =\bigl(2003/6-2d\bigr)f,\qquad d=k+3.            \tag{DC17}
\]

For every1\textless=k\textless=122 this is positive; at122 it is14084. Therefore the whole upper expression increases with z, and may be evaluated at z\_*. No lower bound on z, finite positive metric or existence at this endpoint is needed.

Use the original directed rational bounds

\[
 L_T=1019333/1024>\log_2T,\quad L_Q=287843/4096<\log_2Q,
 \quad L_H=28900/4096>\log_2H,
\] \[
 L_U=601356/4096<\log_2U,\qquad L_{191}=15519/2048>\log_2(191).
                                                                    \tag{DC18}
\]

Define a\_d as the least integer with A\_d\textless2\^{}(a\_d), and b\_k as the least integer with binom(164,k)\textless2\^{}(b\_k). These are exact strict integer comparisons of factorials and powers; equality is never rounded down. The complete k-cell has logarithm strictly less than

\[
\begin{split}
 E_k={}&881+1+5+42+a_d+b_k+4fL_{191}-929f/2+2f\\
 &+(2d+1/6)f(L_H-1)-(d-3/2)fL_U-335\cdot167L_Q\\
 &+339\cdot25L_T+997\cdot25-(197/2)\kappa_k.
\end{split}                                                  \tag{DC19}
\]

Here the1 is the full radius enlargement,5 bounds28,42 bounds C\_F, and2f is the exact4\^{}f determinant factor. The U coefficient retains both the real trace and the projected column. The T power339s contains335s from the common ideal Green numerator and4s from the actual projected column. The beta fibre is included in both the M/B exponents and kappa.

The exact companion certificate checks every one of the122 factorial/binomial comparisons and complete rational exponents. Their maximum is

\[
                  \max_{1\le k\le122}E_k
                     =E_{122}=-7111295/4096<-1736,           \tag{DC20}
\]

with a\_125=-281106 and b\_122=131. Thus the entire low-ideal sparse range is less than

\[
                         122\,2^{-1736}<2^{-1729}.           \tag{DC21}
\]

These are all coordinate-support sizes in the stated range, not an interpolation between a few numerical endpoints. The script checks each exact integer comparison and returns the full table. The first unclaimed k123 endpoint is positive; this only marks failure of this sufficient envelope and is not a claim about an actual contribution.

\subsection{5. The complete high-ideal range and the full union}\label{the-complete-high-ideal-range-and-the-full-union}

For z\textgreater=z\_*, keep the full167-dimensional covector ball of R076 CV10, with A\_167\textless2\^{}-387281, and use the strong Green estimate DC8. Its whole z dependence is

\[
                         (M/z+4200)^{334f}z^{-7f/6}.         \tag{DC22}
\]

It decreases on z\textgreater0, so this full high-ideal sector is bounded at z\_*. At that endpoint the same one-bit radius inflation DC15 applies. This branch counts all167 covector coordinates together and has no binomial or separate k union.

Define exactly the full common constant from DO100 HI13,

\[
\begin{split}
 C_{25}={}&881+1+5+42-387281+4fL_{191}-929f/2+2f\\
 &+(2005f/6)(L_H-1)-(331f/2)L_U+100L_T+24925.
\end{split}                                                  \tag{DC23}
\]

The new high endpoint is

\[
\begin{split}
 E_{\mathrm{high}}={}&C_{25}-501\cdot167L_Q+9\cdot167f
           +167\cdot25L_T+(197/2)(2011f/6)\\
          ={}&-5721469/4096<-1396.                           \tag{DC24}
\end{split}
\]

This is exactly DO100 HI14 with98 replaced by197/2; the cover remains valid and no source, target or phase is altered. Every original common-J pattern is in exactly one of the low/high ranges. We deliberately retain the complete pattern count for each sector and sum their complete upper bounds; the union is consequently

\[
                     2^{-1396}+2^{-1729}<2^{-1395}.           \tag{DC25}
\]

This proves DC4 and the remaining restrictions DC5. It does not claim removal of k\textgreater=123 with z\textless z\_*, nor does it assert that those terms exist.

\subsection{6. Chronology, exact checks and geometric scope}\label{chronology-exact-checks-and-geometric-scope}

The actual predecessor is checks/do102\_joint\_remainder\_error.json, with89 all-top and31 uniform entries. Its ordered suffix retains DO10029550,DO101127443,R080537 and DO102881. Append1395 once, retaining every previous charge. The full exact sums remain below2\^{}-230 and the complete support floors remain25/11. The companion verifies the actual ordered arrays and all122 complete cells, the precise cover and discriminant/Gram constants, both directed endpoints and the full union. This is an exact arithmetic certificate for the stated proof, not an enumeration of rational arrows.

The direct settled interfaces are DO099 section6/JR24; DO095 RA15--20; R072 PJ4--18; R076 CV4--18; and DO100 HI1--14. The new result is the original-support covector ball DC13--16 and its full composition with the sharper two-range cutoff. The preceding proofs and their exact source locators are retained; no settled theorem is re-derived under an altered normalization.

The fixed support condition defines coordinate submodules inside the original positive Hom module over Spec(O\_E). Its dual metric and determinant4U remain explicit; the common-C ideal is a paired determinantal condition, and the norm fibre retains actual ideal classes and CM norm-one units. Rational metric transformations never supply new integral frames. The original polarization, theta line, Picard classes and biextensions remain unchanged. Contextual anchors are SGA3 I4.7.3/4.7.3.1; SGA6 VI Theorem3.3; FGA232 Theorem3.1 pp148--149; FGA236 Corollaries1.7/2.3 pp229/232; SGA7 I VII2.7/2.9 pp169--174 and VIII7.4(b) pp310--311; cote4 scan003 §§8.1.1--8.1.2,scans005--006 §§8.4.1--8.4.3; cote54 scans005--007. These carry the formal-group, Dieudonne-crystal and Cartier-duality context without a new comparison theorem. No new lambda-ring/K-theory equality, l-adic sheaf/weight assertion, CM-motive or motivic L-function conclusion follows.

DP\_MULT Q3e and THETA3 G1/A7 gain exactly DC5's smaller actual finite/real gate. The signed sufficient inequality R078 J25.2, literal fully admissible mixed-CM source, selected arithmetic-image content, prescribed global Bessel/Mram value, Read/STiesR/specific032 pairing and changed-family DP44 comparison remain open. CX\_GRH CX-ALG-024 has no new RH theorem, and GCT\_PLETH T-052 retains its separate global arithmetic charge comparison. No Lean build or archive audit is performed.

\clearpage

\section{AM-R-082 --- type-conditioned rank-three dense-row and ideal cuts at support25}\label{am-r-082-type-conditioned-rank-three-dense-row-and-ideal-cuts-at-support25}

\textbf{Canonical checked integration --- 2026-11-23T00:29:38.168958+00:00.} Root and independent whole-proof/script/interface readings and separate2710-assertion final-fingerprint replays passed. Historical pending wording is preserved. R082 TP1--13 restrict every remaining rank-three support25 term, with its actual t=11..25 type-III places, to k\textgreater K\_t and2\^{}(42p\_t)N\_EJ\textless T\^{}50 for the same common-C ideal. In increasing t order, (p\_t,K\_t) are(1120,149),(1096,148),(1073,147),(1049,147),(1026,146),(1002,145),(979,144),(955,142),(932,141),(908,139),(885,137),(861,135),(838,132),(814,129),(791,125). The complete general type/pair weights and actual-t minor poles give maxima-9691203/4096 over2106 low cells and-11857677/4096 over15 high cells, one2353 charge. Actual91/31 errors stay below2\^{}-230; complete floors25/11. The all-II unipotent rank-two branch and surviving dense smaller-ideal rank-three profiles remain; signed/global/source/selected-image/specific032/DP44/RH gates remain open.

Author: \nolinkurl{/root}, GPT-6 Astra Pro; inherited effort not separately exposed. ALP\_MACHINE,23 November2026 UTC. Complete pending proof; saved whole-proof readings and exact replays are separate check events.

\subsection{1. The unchanged actual sector and the complete table}\label{the-unchanged-actual-sector-and-the-complete-table}

Retain the actual support25 rank-three scalar-cross sector DC1--3 of R081, the all-top H133 cutoff, f168, every original paired pencil, the actual finite and selected real factors, the fixed metric, central phase, class mass and Haar normalizations. The common-C ideal J is the same deterministic ideal as DO100/R081; L=J I\_Sigma and z=(N\_E L)\^{}(1/336), with N\_E I\_Sigma=T\^{}-50. The original S covector support has size k,1\textless=k\textless=164. Let t be the actual number of type-III unequal-pencil off-diagonal places.

DO074 TC1 applies with rank(h\_SS-I)=0 and at every total support. It makes every t\textless=10 term globally unipotent of rank at most two. This contradicts the present global rank three. Hence necessarily11\textless=t\textless=25; no new numerical charge is associated with that carried exclusion.

For each t use the following integers p\_t,K\_t and put z\_t=2\^{}(-p\_t/8).

\begin{longtable}[]{@{}rrr@{}}
\toprule\noalign{}
t & p\_t & K\_t \\
\midrule\noalign{}
\endhead
\bottomrule\noalign{}
\endlastfoot
11 & 1120 & 149 \\
12 & 1096 & 148 \\
13 & 1073 & 147 \\
14 & 1049 & 147 \\
15 & 1026 & 146 \\
16 & 1002 & 145 \\
17 & 979 & 144 \\
18 & 955 & 142 \\
19 & 932 & 141 \\
20 & 908 & 139 \\
21 & 885 & 137 \\
22 & 861 & 135 \\
23 & 838 & 132 \\
24 & 814 & 129 \\
25 & 791 & 125 \\
\end{longtable}

\textbf{Theorem.} For the unchanged normalized summand w(h), the complete absolute contribution of the stated union satisfies

\[
 \sum_{t=11}^{25}
 \sum_{\substack{h\text{ in the stated sector},\ t(h)=t\\
 z(h)\ge z_t\ \mathrm{or}\ [z(h)<z_t,\ k(h)\le K_t]}}
                       |w(h)|<2^{-2353}.                     \tag{TP1}
\]

Consequently every remaining rank-three support25 term satisfies, for its actual t,

\[
             k\ge K_t+1,\qquad 2^{42p_t}N_EJ<T^{50}.        \tag{TP2}
\]

All thresholds refine R081: each K\_t\textgreater=125\textgreater122 and each42p\_t\textgreater=33222\textgreater33096. The complete support floors remain25/11. This does not remove every rank-three term or the all-II globally unipotent rank-two branch, and it does not prove the signed inequality or an actual source/DP44 comparison.

\subsection{2. Exact weights for a fixed type-III count}\label{exact-weights-for-a-fixed-type-iii-count}

Use K070 PT1--6 on the unchanged full pencil and its actual subfamily of b=T-1 labels. The matrix need not be globally unipotent. In type II the exact coefficient is1/b and at most b diagonal labelled pairs are possible; their full weighted cost is at most one. In type III the exact coefficient is-1/{[}b(b-1){]} and there are at most b(b-1) off-diagonal pairs; the full absolute weighted cost is again at most one. These are paired CM choices, not independent labels at both components.

For a nonintegral type-I place, K062's general bound, retained in K070 PT6, is \textbar rho\textbar\textless=2/(T-2). Bounding all available ordered pairs by b² gives cost

\[
                  \frac{2(T-1)^2}{T-2}<3T,\qquad T>4.       \tag{TP3}
\]

Here type I does not receive the stronger globally unipotent table. Summing both possible types I and II at one of the other25-t places costs at most3T+1\textless4T. Choose the t positions with binom(25,t) and retain the original support factor2\^{}168 separately. For each prescribed pair tuple the matrix count below is uniform; only after that count may these true weights be summed. The complete type/pair factor is therefore at most

\[
 W_t={25\choose t}(4T)^{25-t}
       <2^{75-2t}T^{25-t},\qquad
 P_t:=75-2t+(25-t)L_T,\quad L_T=1019333/1024.               \tag{TP4}
\]

The binomial bound binom(25,t)\textless2\^{}25 is conservative also at t=25. TP4 includes every arrangement of the other types, every original ordered labelled pair and the exact absolute coefficient. No additional type-profile multiplier or independent CM-partner enumeration remains.

This replaces the coarse factor2\^{}(997*25) in R081. It is not multiplied by that old factor. The rank-three finite estimate also retains the actual t in its T\^{}(167t) numerator rather than substituting25.

\subsection{3. Complete unchanged matrix factors and the low-ideal expression}\label{complete-unchanged-matrix-factors-and-the-low-ideal-expression}

All counts in R081 DC6--16 are uniform for each original pair tuple and remain valid with t fixed. In particular, keep the common G patterns2\^{}668, support2\^{}168, pivots2\^{}8, Green correction2, source fibre2\^{}36; their exponent is881. The original projected negative-column count is191\^{}(4f)U\^{}f T\^{}100, retaining its fixed U denominator and full paired target index, including zero targets. The relative beta norm fibre is28C\_F(H/2)\^{}(f/6)z\^{}(-f/6), with C\_F\textless2\^{}42 and all28 generators. The actual selected trace is bounded by2\^{}(-929f/2)U\^{}(-f/2)z\^{}-f, with both compact factors. The original norms are M=H/(2sqrt U),B0=H/2, with Q=1427908193729121522601,U=77Q²-14.

For a prescribed k-element original S support, d=k+3 and the joint covector ball has exact dual Gram diag(I\_k,G\^{}-1), detG=4U, covolume D0\^{}d z\^{}(2df)/(4U)\^{}f.~Its full cover is at most4200sqrt(f)z, with \(D_0=7^{140}29^{144}5^{84}13^{84}\). The argument in R081 DC13--16 supplies

\[
 \#\lambda_k<2{164\choose k}A_d(4U)^fM^{2df}z^{-2df},
 \qquad A_d=528^{df}/((df)!D_0^d),                           \tag{TP5}
\]

whenever z\textless=z\_t: every p\_t/8\textgreater=791/8\textgreater197/2, so the previously proved one-bit full cover inflation applies. No smaller covering term, integral moving basis or asymptotic density is substituted.

Let a\_d,b\_k be the same least strict integer binary ceilings of A\_d and binom(164,k) as in R081 DC19. Put

\[
 \kappa_k=(2003/6-2d)f,
\] \[
\begin{split}
 C_k={}&881+1+5+42+a_d+b_k+4fL_{191}-929f/2+2f\\
 &+(2d+1/6)f(L_H-1)-(d-3/2)fL_U-335\cdot167L_Q
                 +339\cdot25L_T .
\end{split}                                                  \tag{TP6}
\]

The rational directed L\_Q,L\_H,L\_U,L191 are exactly R081 DC18. All K\_t\textless=149\textless164, so kappa\_k\textgreater0 in every claimed low cell. The entire weighted expression with weak335 Green therefore increases with z and is bounded at z\_t by the exact full exponent

\[
                         E_{t,k}=C_k+P_t-(p_t/8)\kappa_k.
                                                                    \tag{TP7}
\]

No lower ideal norm bound or attainment of z\_t is required. Every beta, trace, determinant, projected-column and source-fibre exponent is already present in C\_k.

The certificate checks each of the sum\_t K\_t=2106 cells by exact factorial, binomial, power and rational comparisons. Their complete maximum is

\[
               \max_{11\le t\le25,\ 1\le k\le K_t}E_{t,k}
                  =-9691203/4096<-2366,
                     \quad(t,k)=(14,147).                   \tag{TP8}
\]

Thus their full low-ideal union is less than2106*2\^{}-2366\textless2\^{}-2354. The selected K\_t are conservative choices; next-cell sufficient bounds are recorded but no conclusion of existence or nonzero contribution follows from their failure to meet the chosen margin.

\subsection{4. Complete high-ideal sectors with their actual type-III poles}\label{complete-high-ideal-sectors-with-their-actual-type-iii-poles}

For z\textgreater=z\_t retain the full167-coordinate covector ball, so there is no binomial or k union. DO095 RA15--20 gives the genuine pointwise bound

\[
       |a(h)|<2Q^{-501\cdot167}2^{9\cdot167f}T^{167t}.       \tag{TP9}
\]

The original three-minor witness and its simple paired type-III poles are unchanged. There is no extra witness or divisor-pattern count. Combined with all other original counts, the whole z expression is(M/z+4200)\^{}(334f)z\^{}(-7f/6), decreasing on z\textgreater0. Evaluate it at z\_t, where TP5's full cover control remains valid.

Define

\[
\begin{split}
 C_H={}&881+1+5+42-387281+4fL_{191}-929f/2+2f\\
 &+(2005f/6)(L_H-1)-(331f/2)L_U+100L_T\\
 &-501\cdot167L_Q+9\cdot167f.
\end{split}                                                  \tag{TP10}
\]

The complete high exponent for type count t is

\[
                H_t=C_H+167tL_T+P_t+(p_t/8)(2011f/6).      \tag{TP11}
\]

All15 exact comparisons are checked separately. Their maximum is

\[
                   \max_{11\le t\le25}H_t
                         =-11857677/4096<-2894,
                         \qquad t=11.                      \tag{TP12}
\]

Thus their full high-ideal union is less than15*2\^{}-2894\textless2\^{}-2890. Each original term has one actual t and is in one of its two norm ranges. Counting a full common-G pattern family in every containing sector is an explicit overcount. Summing both complete unions gives

\[
               2106\,2^{-2366}+15\,2^{-2894}<2^{-2353},     \tag{TP13}
\]

which proves TP1. The equivalence of the low range and TP2 follows by raising z\textless2\^{}(-p\_t/8) to336:336/8=42 and z\^{}336=N\_EJ/T\^{}50.

\subsection{5. Actual chronology and limitations}\label{actual-chronology-and-limitations}

The predecessor at author freezing is checks/r081\_joint\_remainder\_error.json with90 all-top and31 uniform entries, complete floors25/11. Append2353 once, retaining every old charge and in particular R0811395 and DO10029550. If another independently checked charge is integrated first, the explicit --prior option must use that actual complete extension; the certificate requires the entire R081 prefix and unchanged uniform array, records its exact input hash and retains every intervening entry. The canonical integration records the actual order and both full exact sums below2\^{}-230.

The proof uses the complete original R081 counts, K070 PT2--6's general incidence types and coefficient bounds, DO074 TC1's any-support exclusion for t\textless=10, and DO095 RA15--20's actual simple-pole witness. None requires global unipotence of the present rank-three term. The new steps are the full conditional type/pair sum TP3--4 and its complete low/high type-indexed endpoint table. Exact finite diagnostics supplement those arguments and do not enumerate the arithmetic arrows.

This refines only the actual rank-three support25 containing-set remainder. The all-II unipotent rank-two branch, every surviving dense smaller-ideal rank-three profile, larger supports, the signed sufficient inequality R078 J25.2, prescribed global Bessel/Mram value, literal admissible mixed-CM source, selected arithmetic-image content, Read/STiesR/specific032 pairing and DP44 comparison remain open. No RH or global oscillator comparison is claimed.

\subsection{6. Algebraic geometry and carried locators}\label{algebraic-geometry-and-carried-locators}

The type decomposition consists of incidences between the same integral pencil lattices over Spec(O\_E). The paired determinantal ideal and its original-coordinate submodule retain the same norms and fixed complex realization. Conditioning on t changes no source/target object, polarization, Picard class, theta line, biextension, metric or Haar measure. The two inequalities are counts of the actual integral correspondences with their prescribed normalized coefficients.

Exact contextual anchors remain SGA3 I4.7.3/4.7.3.1; SGA6 VI Theorem3.3; FGA232 Theorem3.1 pp148--149; FGA236 Corollaries1.7/2.3 pp229/232; SGA7 I VII2.7/2.9 pp169--174 and VIII7.4(b) pp310--311; cote4 scan003 §§8.1.1--8.1.2,scans005--006 §§8.4.1--8.4.3; cote54 scans005--007. No new lambda-ring/K-theory identity, l-adic sheaf or weight theorem, CM-motive comparison or motivic L-function conclusion is inferred. DP\_MULT Q3e and THETA3 G1/A7 gain exactly TP2's smaller finite/real gate; CX\_GRH CX-ALG-024 and GCT\_PLETH T-052 retain their distinct arithmetic gates. No Lean build or archive audit is performed.

\clearpage

\section{AM-DO-103 --- an actual line-denominator cut in the all-II support25 sector}\label{am-do-103-an-actual-line-denominator-cut-in-the-all-ii-support25-sector}

\textbf{Canonical checked integration --- 2026-11-23T00:35:51.827105+00:00.} Complete root and independent proof/script/interface reviews passed. Frozen raw proof and historical certificates remain unchanged. DO103 removes the complete all-II globally rank-two-unipotent scalar-cross support25 subset with N\_E denominator(eta)\textless=2\^{}120. Both actual CM pole lines force the same selected delta=eta-Q(-7+2sqrt(-7)) divisibility. The globally integral denominator-cleared delta ideal has norm at least7\^{}84 T\^{}p and below{[}4(U+14){]}\^{}168 N\_E D\_eta, so the exact host inequality gives p\textless=24. The shared cross/double-pole containing set, original affine central counts and all164 CF endpoints give25T*2\^{}-2879\textless2\^{}-1878. Actual92/31 ordered errors remain below2\^{}-230, complete floors25/11. Every remaining all-II support25 term has N\_E D\_eta\textgreater2\^{}120. The large-denominator branch, rank-three profiles, signed/global/source/selected-image/specific032/DP44/RH gates remain open.

Author: \nolinkurl{/root/definite_oscillator}, ALP\_MACHINE. 23 November2026 UTC.

Controlled pending proof; author whole review, independent replay/review and canonical integration are separate events.

\subsection{1. Exact hypotheses and statement}\label{exact-hypotheses-and-statement}

Keep the unchanged actual J56 all-top normalized kernel and H=133, f=168. Work at support s=25 in the DO099/R071 scalar-cross residual: h\_ZZ=I on the positive167-space Z, C=h\_PZ has nonzero restriction to S164, h=I+N is globally unipotent of rank2, and all334 unmixed G reductions are unipotent. Every active T place has actual type II. Keep all fields, CM partners, lattices, labels, finite projectors, selected real trace, packet point, phase, class masses and Haar measures.

Use the original primitive vector u, U=u*u=77Q²−14, Q=1427908193729121522601 and P=Eu⊕Ee\_n with Gram diag(U,−1). The exact DO097/DO098 normal form is

\[
 C=\ell r,\quad \ell=(1,\eta)\text{ in the }(u,e_n)\text{ basis},
 \quad \eta\bar\eta=U,\quad N_{PP}=d\ell\ell^*J_P,
 \quad N^2Z=0.                                               \tag{LD1}
\]

Thus eta is a uniquely normalized global E scalar; it is not assumed integral. Define its actual denominator ideal by

\[
 \mathfrak D_\eta=\prod_{\mathfrak p}
       \mathfrak p^{\max(0,-v_{\mathfrak p}(\eta))}.
                                                                    \tag{LD2}
\]

\textbf{Theorem.} The entire stated all-II support25 sector with

\[
                            N_E\mathfrak D_\eta\le2^{120}     \tag{LD3}
\]

has total absolute actual normalized contribution less than2\^{}-1878. Outside this charged subset, the same all-II branch satisfies the explicit strict inequality N\_E D\_eta\textgreater2\^{}120. This is a line-denominator restriction, not a complete all-II estimate or a support-floor change.

The proof identifies one common set supporting both every cross pole and every full double pole. It does not replace the paired pencil by an abstract norm sphere or assume that fixing C makes the central scalar integral.

\subsection{2. The exact small comparison scalar and the two pole lines}\label{the-exact-small-comparison-scalar-and-the-two-pole-lines}

Let j=√−7 and set

\[
 a_0=Q(-7+2j),\qquad a_0\bar a_0=77Q^2=U+14,
 \qquad \delta=\eta-a_0.                                    \tag{LD4}
\]

The actual J56 uniformizer is pi=A+Bj with A=2B+1, and K077(AP1)/LIFT44(PR1--7) give

\[
 \pi_k=\pi+Qa_0,\quad
 \pi\bar\pi=T,\quad
 \pi_k\bar\pi_k=T+Q^2U,\quad
 \pi\bar a_0+\bar\pi a_0=-14Q.                             \tag{LD5}
\]

These are exact original scalar identities. In particular a0, bar(a0), Q and U are units at every active T place. The last assertion follows also from gcd(T,77QU)=1 in the actual-host certificate.

At a selected CM prime P above an active relative place, pi is a uniformizer and bar(pi) is a unit. The actual fixed type-II lattice has c≠0 and generator

\[
 A_c=I+c\pi^{-1}zq_*,\quad z=v_0/\bar\pi,
 \quad v_0=(Qu,\pi_k),\quad q_*=-v_0^*J-\pi m,
 \quad q_*z=0.                                               \tag{LD6}
\]

The omitted/zero-label convention is the literal one of LIFT44(PR4--7) and K077(AP2--5): c=0 fixes the original standard lattice and is not an active type-II place. The pole line at P is the reduction of v0.

At the conjugate prime bar(P), the actual lattice is the inverse-adjoint dual. Its pole column is J bar(q\_*)\^{}t; LD6 gives

\[
                   J\bar q_*^t\equiv-v_0\pmod{\bar\pi}.    \tag{LD7}
\]

Thus its projective pole line is again the reduction of the same global v0, now at bar(P). This is a statement about the actual dual generator; it does not identify the two matrices or use an independent conjugate label.

In the (u,e\_n) frame the selected pole line has slope pi\_k/Q≡a0 modulo P. At bar(P), LD5 gives

\[
       \pi_k/Q\equiv a_0-14/\bar a_0=U/\bar a_0
                                    \pmod{\bar P}.           \tag{LD8}
\]

The fixed Z/P frame has determinant of modulus2U and is a unit at all selected G and T primes. The original first positive coordinate u1=2, Q and pi\_k are T-units. These facts allow the saturated line reductions below without declaring a moving frame globally integral.

\subsection{3. A common eligible set for cross poles and full double poles}\label{a-common-eligible-set-for-cross-poles-and-full-double-poles}

Call an active relative place eligible if

\[
                         v_P(\delta)\ge1,                   \tag{LD9}
\]

where P is its original selected pi-component. This implies eta is a unit at P, since a0 is a unit; eta bar(eta)=U then makes eta a unit at bar(P) as well. The denominator ideal D\_eta is therefore trivial at both components of an eligible place.

We prove that every cross pole, and every full N double pole, occurs only at an eligible relative place.

First, each fixed original S input and each of the three fixed integral transverse inputs belongs separately to the actual source and target pencils, at both components. This is the literal DO097(UC10--13) input. Because N²Z=0 and h\_ZZ=I, its image under N is its C image and lies in the target lattice. A nonintegral such image has simple pole and leading vector on that target's pole line. Therefore the saturated global line Eℓ has the same residue line as v0 at that component. If v\_P(eta)\textgreater0 the saturated line has zero negative coordinate; if v\_P(eta)\textless0 it has zero positive coordinates. Neither equals the actual pole line, which has both a positive unit coordinate and a negative unit coordinate. Hence a cross pole forces eta to be a unit and the corresponding pole-line slope equality.

Second, put K=A\_c\^{}-1 h A\_c, which is integral. Since (zq\_*)²=0,

\[
 N=K-I+c\pi^{-1}(zq_*K-Kzq_*)
             -c^2\pi^{-2}z(q_*Kz)q_*.                       \tag{LD10}
\]

If any entry has pole2, its leading image is the actual pole line. The same calculation holds in the inverse-adjoint dual generator at bar(P). On the other hand, the cross blocks have at most simple poles, N\_ZZ=0, and LD1 shows that the only possible double-pole part is the rank-one PP block with saturated image Eℓ. The same unit-coordinate argument therefore forces eta to be a unit and the same slope equality. No conclusion about individual central entries being integral is used.

At P either implication gives eta≡a0, hence LD9. At bar(P) it gives eta≡U/bar(a0). Define

\[
              \gamma=\bar a_0\eta-U=-\eta\bar\delta.         \tag{LD11}
\]

Since eta and bar(a0) are units there, this is equivalent to bar(delta) divisible by bar(P), or delta divisible by P. Thus both partners give the same eligible condition, with no second independent place count.

Consequently, if E0 is the eligible set, every original counted C scalar has at most a simple paired pole supported on E0; the full N has double poles only on E0 and at most simple poles on its complement. This applies to the three original transverse values and the negative pivot as well as to the S row. Integral primitive-u extraction is exactly the one in DO097/DO098.

\subsection{4. The denominator norm and the exact factor at seven}\label{the-denominator-norm-and-the-exact-factor-at-seven}

The element delta is nonzero, because its vanishing would contradict eta bar(eta)=U and a0 bar(a0)=U+14. The ideal

\[
                         (\delta)\mathfrak D_\eta            \tag{LD12}
\]

is integral: at each prime a0 is integral, so v(delta)≥min(v(eta),0), and D\_eta clears this possible negative valuation.

We need the actual local data at7, rather than a generic discriminant estimate. In the unchanged field, E/F splits at both F primes over7; there are four E primes, each absolute ramification index6 and residue degree14. These exact data are the actual Galois/inertia computation in the saved packet\_plane\_norm\_boundary, §§1--4, using K032(QH9--15) and K030(PH2). The companion rechecks the order84 decomposition group and its six-element inertia; complex conjugation is absent. Also v7(U)=1 and Q is prime to7. Therefore at each of the two conjugate E-prime pairs over7,

\[
 v_{\mathfrak p}(U)=6,\quad v_{\mathfrak p}(a_0)=3,
 \quad v_{\mathfrak p}(\eta)+v_{\bar{\mathfrak p}}(\eta)=6. \tag{LD13}
\]

The valuation3 follows from -7+2√−7: its two summands have valuations6 and3, with2 and Q units. Write a=v\_p(eta). The sum of the two valuations of the integral ideal LD12 is at least

\[
 \min(a,3)+\max(-a,0)+\min(6-a,3)+\max(a-6,0)\ge3.         \tag{LD14}
\]

For a≤0 or a≥6 the lower expression is3; on0≤a≤3 it is3+a, and on3≤a≤6 it is9−a. This proves the inequality for all integer valuations, including fractional eta. Each residue norm is7\^{}14 and there are two such pairs, so LD12 has norm divisible by7\^{}84.

At every eligible active place delta is divisible by the selected degree-one prime and D\_eta is a unit at both partners. These primes are distinct from those over7. If p=\textbar E0\textbar, it follows that

\[
       7^{84}T^p\le N_E((\delta)\mathfrak D_\eta)
           =|N_{E/\mathbb Q}\delta|\,N_E\mathfrak D_\eta
           <[4(U+14)]^fN_E\mathfrak D_\eta.                  \tag{LD15}
\]

The strict archimedean inequality uses \textbar eta\_tau\textbar=sqrt(U) and \textbar a0\_tau\textbar=sqrt(U+14), so \textbar delta\_tau\textbar\textless2sqrt(U+14), at all168 complex labels. This is the absolute E norm with336 embeddings, equivalently168 squared complex magnitudes.

The original integer values satisfy

\[
                 7^{84}T^{25}>2^{120}[4(U+14)]^{168}.         \tag{LD16}
\]

Thus LD3 and LD15 force p≤24. No unproved small-denominator norm solution or principality of D\_eta is assumed. The field-level global isotropy of P from the earlier norm-boundary note does not supply LD3 for an actual arrow.

\subsection{5. Complete count of the small-denominator sector}\label{complete-count-of-the-small-denominator-sector}

For each actual matrix satisfying LD3, pad E0 to a subset D of24 of its25 active places. There are at most25 choices. Fix that padded set, the actual paired labels and one canonical whole-C G pattern. Section3 gives precisely

\[
 C\text{ scalars}\in L=JI_D,\quad N_EI_D=T^{-48},\qquad
 \text{full entry poles}\le I_{\Sigma,F}I_{D,F}\mathcal O_E.
                                                                    \tag{LD17}
\]

The full cross map and the original w=N\_nn reconstruct h injectively by DO098(CF4--7). Its fixed-C differences lie in one affine coset of iJ0 I\_Sigma,F I\_D,F, not a half-lattice. Hence the two central counts are exactly DO102(NT8--9) with s25,v24:

\[
 \#w<2^{479}Q^{167}T z^{-2f},\qquad
 \#w\le134^fT^{49},\quad z=(N_EL)^{1/(2f)}.                 \tag{LD18}
\]

All original G patterns, support and pivot choices, the three-coordinate covolume, the relative norm pivot fibre, the actual selected trace and every k=1..164 endpoint are the unchanged CF8--32 factors. The actual global cap z\textless M/sqrt(30) remains valid because v24. Each alternative central count is exactly T times its s=v=24 counterpart. The complete CF32 union was less than2\^{}-2879.

Every active place is type II. After fixing and counting each actual pair, K071(NB7) gives complete pair-weight sum at most1 at each place. There is one actual type profile, no independent conjugate label, no second orientation or denominator-ideal enumeration. The only additional covering multiplicity is the25 padded choices. Thus the entire desired sector is

\[
  <25T\,2^{-2879}<2^{-1878},\qquad
  2879-5-1019333/1024=1923643/1024>1878.                      \tag{LD19}
\]

This is a complete positive upper bound for the absolute contribution, not a signed-positivity assertion. The matrix condition LD3 is not counted separately: all containing LD17 sectors are counted, so the bound is uniform over every admissible denominator ideal.

\subsection{6. Verification, chronology and remaining boundary}\label{verification-chronology-and-remaining-boundary}

The certificate checks the exact original pi/a0 identities, the actual7 decomposition data and all valuation branches, the integer comparison LD16, the paired finite pole-line equivalence, all25 padding choices and the whole CF endpoint union. The local finite examples diagnose the algebra only; actual arrow membership and normalized counting remain the proved source interfaces above.

The actual predecessor at author freeze is checks/r081\_joint\_remainder\_error.json, with90 all-top and31 uniform entries and complete support floors25/11. Append1878 once, obtaining91/31 while retaining every earlier charge, all uniform terms and both complete floors. The certificate reads the actual complete arrays, verifies the DO102881 and R0811395 ordered tail, and checks both entire sums below2\^{}-230. Any later canonical insertion must be retained separately in chronological order, without rewriting this historical predecessor.

The new restriction is on the original global eta denominator, N\_E D\_eta\textgreater2\^{}120 outside the removed all-II subset. Every actual covector value r(y) and eta r(y) lies in the same L, so the exact valuationwise identity L∩eta\^{}-1L=L D\_eta is available for a future large-denominator count. No bound for that complementary count is asserted here. The all-II large-denominator branch and any remaining rank-three branch stay open; no full support-floor, signed sufficient inequality, prescribed period, arithmetic image, DP44 or RH conclusion follows.

\subsection{Canonical algebraic-geometric context}\label{canonical-algebraic-geometric-context}

The denominator and common pole conditions are fractional ideal conditions on the original integral correspondence over Spec(O\_E). The imaginary central affine lattice retains its integral model and discriminant; the inequality changes a counted subset, not the polarization, Poincare biextension, theta line or Haar measure. The fixed complex norm is the same realization of those original modules. Exact inherited contextual anchors: SGA3 I4.7.3/4.7.3.1; SGA6 VI Theorem3.3; FGA232 Theorem3.1 pp148--149; FGA236 Corollaries1.7/2.3 pp229/232; SGA7 I VII2.7/2.9 pp169--174 and VIII7.4(b) pp310--311; cote4 scan003 §§8.1.1--8.1.2 and scans005--006 §§8.4.1--8.4.3; cote54 scans005--007. These are contextual locators, not newly established comparison theorems. DP\_MULT Q3e and THETA3 G1/A7 retain their actual-source/period/accounting obligations; CX\_GRH CX-ALG-024 and GCT\_PLETH T-052 retain their separate gates. No new lambda-ring or K-theory operation, l-adic weight theorem or motivic L-function identity is inferred.

\clearpage

\section{AM-LIFT-079 --- quadratic pure types at the actual good-CM-point interface}\label{am-lift-079-quadratic-pure-types-at-the-actual-good-cm-point-interface}

\textbf{Canonical checked integration --- 2026-11-23T00:35:51.827105+00:00.} Complete root and independent proof/script/interface reviews passed. Frozen raw proof and historical certificates remain unchanged. LIFT079 CT1--10 proves that literal CMObj and, independently, both actual LamObjGood copies force the complex target to be isogenous to a power of one CM elliptic curve, using pure complete quadratic eigenbases and the actual surjective/quasi-inverse sum. The actual LamPtClauses fibre identity transfers this constraint to the prescribed point; a point with a nonelliptic simple factor has no such witness. Thus dropping IsTen or bypassing CMObj while retaining Good, or substituting an unrelated elliptic Good object, does not supply the mixed-J point hypothesis. This is conditional on the stated point/frame premises and is not a contradiction to the accepted original source theorem. R003/DP25 simplicity is carried, not counted again. No mixed replacement constructor, finite charge, period/image/value, DP44 or RH conclusion follows.

\textbf{Status:} author whole-proof, source-interface and complete diagnostic-script review PASS, 2026-11-23. The exact completion time and frozen fingerprints are recorded in pure\_type\_good\_point\_boundary\_verification.json. This is a source-interface result, not an admissible mixed-family construction or a nonvanishing theorem. The source commit is 7cbf7a0c32fca1badf05931879a37c3f8a2b03b4. All newly retrieved source files are preserved in the adjacent source\_cm directory. No Lean build is used.

\subsection{1. The exact conclusion}\label{the-exact-conclusion}

Put \(K_0=\mathbf Q(\sqrt{-d})\), where \(d\) and the specified quadratic root belong to the literal geometric Host. Let \(C_0/\mathbf C\) be a CM elliptic curve with endomorphism algebra \(K_0\).

\textbf{CT1.} For every genuine characteristic-zero instance of the source record GlobalTheory.CM.CMObj, with parameters \(p,q\), its complex target satisfies \[
 B_{\mathbf C}\ \sim\ C_0^{\,f(p+q)},\qquad
 f=[E:\mathbf Q]/2.                                      \tag{CT1}
\] Here \(\sim\) means isogeny of abelian varieties, without an assertion that the polarizations, integral models, or \(E\)-actions are identified by that isogeny. The case \(p+q=0\) means the zero abelian variety.

\textbf{CT2.} The same conclusion is already forced directly by the unchanged LamObjGood interface, both in GlobalTheory.LamF and in its actual Formalism.L2.Pj.abelian₂ copy. It does not require constructing a CMObj or supplying IsTen first. More precisely, for an actual LamObjGood object with multiplicities \(n_\Psi,n'\), \[
 B_{\mathbf C}\ \sim\ C_0^{\,f(n_\Psi+n')}.                \tag{CT2}
\] The decisive clauses are the complete pure eigenbases and IsIsogSum.

Consequently, removing only IsTen, or bypassing only the CMObj constructor, cannot place a nonzero power of the geometrically simple three-dimensional Jacobian \(J=\operatorname{Jac}(y^2=x^7-1)\) in this unchanged good-point interface. The conclusion holds after arbitrary field enlargement or isogeny replacement of that target. It is a constraint on this particular source interface, not a claim that mixed CM theta kernels or mixed-type moduli cannot be constructed.

\subsection{2. Literal source premises}\label{literal-source-premises}

The two declarations named Host below must be kept distinct. The toroidal geometric Host specifies a quadratic root, an order element, conjugation and relative unramifiedness. The stronger HostRaw used in the parameter wrapper additionally fixes its Galois frame and conductor conditions.

\begin{longtable}[]{@{}
  >{\raggedright\arraybackslash}p{(\columnwidth - 2\tabcolsep) * \real{0.5000}}
  >{\raggedright\arraybackslash}p{(\columnwidth - 2\tabcolsep) * \real{0.5000}}@{}}
\toprule\noalign{}
\begin{minipage}[b]{\linewidth}\raggedright
Literal source
\end{minipage} & \begin{minipage}[b]{\linewidth}\raggedright
Exact fields used
\end{minipage} \\
\midrule\noalign{}
\endhead
\bottomrule\noalign{}
\endlastfoot
\nolinkurl{ReciprocityTheory/ReductiveGroups/ArithmeticGeometry/BirationalGeometry/ToroidalCompactification/TwoWeightEqOneAddRoot.lean,} 14--26 & The geometric Host has prime \(d\), \(d\bmod16=7\), \(r_t^2=-d\), \(2w_0=1+r_t\), and a complex root \(r_{t,0}\). \\
Same, 108--117 & IsLieEigen is the tangent eigencondition; LieMult gives a basis of that entire eigenspace by a unique-coefficient condition. \\
\nolinkurl{ReciprocityTheory/ReductiveGroups/ArithmeticGeometry/BirationalGeometry/ToroidalCompactification/IsGrpPullback.lean,} 17--23 & OAct is a unital additive and multiplicative action by group homomorphisms. \\
\nolinkurl{Development/AbelianVarietyTheory/Definitions/CMLattices/IsDiffOf.lean,} 43--67 & PAT is an actual polarized abelian scheme; HasType imposes the two tangent multiplicities at every injective algebraically closed characteristic-zero realization. \\
Same, 115--129 & CMObj requires pure types \((1,0)\) and \((0,1)\), target type \((p,q)\), \(E\)-linear homomorphisms, and surjectivity of their sum on the specified complex fibre. \\
\nolinkurl{Development/AbelianVarietyTheory/Definitions/ShimuraVarieties/PowObjGeometricallyConnected.lean,} 82--85, IsLinearHom & Linearity includes IsMonHom, so the maps in the preceding row are group homomorphisms. \\
\nolinkurl{Development/AbelianVarietyTheory/Definitions/CMLattices/LamLetters.lean,} 42--75, 110--142 & LamObj records complete eigenbases; LamObjGood includes those types and a positive-multiple quasi-inverse to the sum of the component homomorphisms. \\
\nolinkurl{Formalism/LocalToGlobal/ProjectiveLattices/RosatiLatticeNorms/LamObj.lean,} 38--71, 106--138 & The actual Formalism copy repeats the same eigenbasis, isogeny, polarization and metrized-line clauses. \\
\nolinkurl{Formalism/LocalToGlobal/ProjectiveLattices/RosatiLatticeNorms/LamPtClauses.lean,} 70--108 & LamPtClauses requires LamObjGood and ties its B fibre, polarization, section, action and letters to the prescribed point; LamPtSt supplies such a witness only under its stated frame hypotheses. \\
\nolinkurl{Formalism/LocalToGlobal/ProjectiveLattices/EmbKCoe.lean,} 29--33 & The coefficient embeddings are the literal restrictions of the complex embeddings. Thus all labels in the specified quadratic half have the same value on \(r_t\), also after applying any embedding of the coefficient field. \\
\nolinkurl{Formalism/SharpRestriction/RamificationFrobenius/TotientGaloisClasses/ConjEMemHostFIff.lean,} 60--66 & HostRaw.cmType is explicitly the half selected by one value of its quadratic root. \\
\end{longtable}

The polynomial \(X^2+d\) is irreducible over \(\mathbf Q\), because \(d>0\). Therefore \(K_0\) is an imaginary quadratic field. Its two embeddings are distinguished by the two values of \(r_t\). Since \(E/\mathbf Q\) is separable, exactly \(f\) embeddings of \(E\) extend each embedding of \(K_0\).

This proof concerns the actual characteristic-zero abelian varieties described by these source records. It does not claim that the finite diagnostics below mechanically verify the source proofs or the analytic theory of abelian varieties.

\subsection{3. Pure quadratic tangent action forces elliptic isotypy}\label{pure-quadratic-tangent-action-forces-elliptic-isotypy}

\textbf{CT3, elementary analytic lemma.} Let \(A/\mathbf C\) be an abelian variety of dimension \(g\), with an action of an order in an imaginary quadratic field \(K_0\). Suppose the induced action on \(\operatorname{Lie}(A)\) is scalar through a single embedding \(\iota:K_0\hookrightarrow\mathbf C\). Then \[
                         A\sim C_\iota^g,                \tag{CT3}
\] where \(C_\iota=\mathbf C/\iota(\mathcal O_{K_0})\).

\textbf{Proof.} Write the complex uniformization as \(A(\mathbf C)=V/\Lambda\). Every algebraic endomorphism lifts to its complex-linear derivative on \(V\), preserving \(\Lambda\). The rational vector space \(\Lambda_\mathbf Q\) is a \(K_0\)-vector space of dimension \(g\): the order action extends to its fraction field, and its rational dimension is \(2g\).

Choose a \(K_0\)-basis \(e_1,\ldots,e_g\) of \(\Lambda_\mathbf Q\). The \(K_0\)-linear isomorphism \(K_0^g\to\Lambda_\mathbf Q\) becomes, after tensoring over \(\mathbf Q\) with \(\mathbf R\), a complex-linear isomorphism \[
       \mathbf C^g \longrightarrow V
\] when \(K_0\otimes_\mathbf Q\mathbf R\cong\mathbf C\) is identified using \(\iota\). This is precisely where the single-embedding tangent premise is used. The lattice \(\sum_j\mathcal O_{K_0}e_j\) and \(\Lambda\) span the same rational vector space and hence are commensurable. Multiplying the former by a nonzero integer gives an inclusion in \(\Lambda\) of finite index. The resulting complex-linear map descends to a surjective holomorphic homomorphism \(C_\iota^g\to A\) with finite kernel, hence an algebraic isogeny of abelian varieties. \(\square\)

The two embeddings of \(K_0\) give the same elliptic isogeny class. Indeed \(\bar\iota(\mathcal O_{K_0})=\iota(\mathcal O_{K_0})\) as subsets of \(\mathbf C\), because the order is conjugation-stable. This assertion forgets the direction of the \(K_0\)-action; it does not identify that action with its conjugate.

The lemma is also the quadratic special case of CM-type classification: an induced CM type gives the corresponding power of the smaller CM abelian variety. Milne (2020, Proposition 3.12 and the proof of Proposition 3.13, p.~30) states this isogeny classification. Poincaré reducibility (Theorem 2.12, pp.~26--27) implies that every abelian subquotient of a power of a fixed CM elliptic curve has only that elliptic simple factor.

\subsection{4. Application to the actual CMObj}\label{application-to-the-actual-cmobj}

Let \(X\) be the source CMObj and take the complex fibre at its specified embedding \(\sigma_1\) of the coefficient field. A ring homomorphism from \(\mathcal O_E\) to \(\mathbf C\) extends uniquely to an embedding of \(E\), because we are in characteristic zero. The unital action extends to \(E\) inside the rational endomorphism algebra.

The algebra \(E\otimes_\mathbf Q\mathbf C\) is a product of copies of \(\mathbf C\), one for each embedding. Its finite-dimensional representations are direct sums of the corresponding eigenspaces. The literal LieMult basis condition therefore says exactly what is needed here; no semisimplicity of an arbitrary integral representation is being assumed.

For \(A_+=X.A(\mathrm{true})\), typeP gives one-dimensional eigenspaces at the \(f\) embeddings with the selected value of \(r_t\), and zero eigenspaces at the other embeddings. Hence \[
 \dim A_+=f,\qquad
 d\iota_{A_+}(r_t)=\sigma_1(\rho)\,\mathrm{id}.             \tag{CT4}
\] Similarly \[
 \dim A_-=f,\qquad
 d\iota_{A_-}(r_t)=-\sigma_1(\rho)\,\mathrm{id}.            \tag{CT5}
\] Since \(r_t\) generates \(K_0\), CT3 applies and both \(A_+\) and \(A_-\) are isogenous, as abelian varieties, to \(C_0^f\). These dimension and eigenspace statements are also exactly the conclusions of the source \nolinkurl{LieAlgebraEigenDecompositionCMSignatureStatement.lean,} line 10.

The target typeB gives \[
                       \dim B=f(p+q).                   \tag{CT6}
\] The maps in CMObj are group homomorphisms by IsLinearHom. Their sum defines a homomorphism \[
                  F:A_+^p\times A_-^q\longrightarrow B.
                                                               \tag{CT7}
\] The onto clause is surjectivity on all complex points of this same map. A morphism from an abelian variety is proper, so its image is closed; if the image were proper, its nonempty complement would contain a closed complex point. Thus \(F\) is surjective. Its source and target have the same dimension by CT4--CT6. Its kernel is consequently zero-dimensional and finite, and \(F\) is an isogeny. Combining with CT3 proves CT1.

This argument uses neither the definition nor an instance of IsTen. The target mixed tangent multiplicity by itself would not prove CT1: the onto condition is essential.

\subsection{5. The obstruction persists directly in LamObjGood}\label{the-obstruction-persists-directly-in-lamobjgood}

Now start with LamObjGood, without assuming that it arose from CMObj. Its xΨ and x′ carry complete Lie bases. The bijections τΨ and τ′ enumerate respectively the embeddings with \(r_t=r_{t,0}\) and their complement. The IsLieType clauses state that each basis vector has its assigned \(E\)-eigencharacter.

The coefficient map is a restriction of the original complex embedding, by EmbKCoe.lean, 29--33. Therefore all entries indexed by τΨ have one and the same value on \(r_t\) in the coefficient field; those indexed by τ′ have the opposite value. A later complex embedding of that coefficient field may interchange the two roots, but cannot split either half into different values. Thus the complete tangent spaces of AΨ and A′ have pure quadratic action. CT3 again gives \[
                         A_\Psi\sim C_0^f,\qquad A'\sim C_0^f.
                                                               \tag{CT8}
\]

The first identity in IsIsogSum gives component morphisms back from \(B\) and a positive integer \(N\) such that \[
                  F\circ G=[N]_B,                       \tag{CT9}
\] where \(F:A_\Psi^{n_\Psi}\times(A')^{n'}\to B\) is the sum of the recorded homomorphisms. Since multiplication by \(N\) on a complex abelian variety is surjective, CT9 forces \(F\) to be surjective. The recorded dimensions agree, so \(F\) is an isogeny. This proves CT2. The second IsIsogSum identity records the converse positive-multiple composition and is consistent with this conclusion, but is not needed merely for surjectivity.

There is no gap from treating a pointwise relation as a map relation: the source identity is quantified over every test object and every point, so it is a Yoneda identity. Specializing to the complex fibre is sufficient for the preceding argument.

An analogous explicit constraint is visible in the separate CMObjBundlePull interface: its Types predicate imposes the two pure eigenspace systems (143--147), and its Good predicate includes a finite surjective sumMap (156--172). It therefore has the same elliptic-isotypy consequence.

\subsection{6. Why this is an actual downstream requirement}\label{why-this-is-an-actual-downstream-requirement}

This is not an obstruction inferred merely from a similarly named unused definition.

\begin{enumerate}
\def\labelenumi{\arabic{enumi}.}
\tightlist
\item
  \nolinkurl{ExistsLamObjGoodOfCmObj.lean,} theorem at line 44, takes the literal CMObj and constructs LamObjGood. Its lines 60--63 consume typeP and typeM to produce the eigenbases. Lines 98--104 consume onto to obtain the actual inverse-isogeny clause.
\item
  \nolinkurl{ExistsGoodCmPointOfIntegralCmObjectStatement.lean,} line 16, takes CMObj and constructs a good CM point above the given moduli point.
\item
  \nolinkurl{ExistsGoodCmPointAtCmBasisFrameStatement.lean,} line 16, supplies the good CM point at the specified CM basis frame.
\item
  LamLetters.lean, 169--184, defines LamDegSt, LamIndSt and LamLenSt using LamObjGood. The actual Formalism copy in LamPtClauses.lean, 53--68, has the same premises. LamIndLenDeg.lean, 14--22, passes the global statements to that source library and identifies its scalar degree with SizeData₂.logOm under the pinned-host premise.
\item
  LamPtClauses.lean, 76--97, additionally identifies the B fibre with the actual prescribed moduli point, retaining the polarization/section, endomorphism and letter relations. CmFrameOfEq.lean, 108--125, consumes precisely this witness in lamw\_link and then applies lambda\_lower\_bound. AxHInst.lean, line 136, retains the LamPtSt, LamDegSt and LamIndSt hypotheses in the subsequent height construction.
\end{enumerate}

\textbf{CT10, the extra pointwise source gate.} Suppose the complex abelian variety at a prescribed point \((U.\mathrm{un}\ k\ x_0).x\) has a nonelliptic geometric simple factor. Then there is no witness of the unchanged LamPtClauses at that point.

Indeed, LamPtClauses includes LamObjGood. Its group pullback and displayed isomorphism identify the generic B fibre with the universal abelian fibre at \(y.x\). Its final scalar-extension equality identifies the latter fibre with the prescribed complex point. CT2 would therefore make that prescribed abelian variety isogenous to a power of \(C_0\), contradicting the nonelliptic factor. The section, frame or coefficient field cannot be chosen independently to avoid this equality.

If such a point also satisfies every frame antecedent in LamPtSt, the corresponding LamPtSt implication cannot be supplied for a proposed replacement realization. This is a conditional statement about that replacement; no existence of such a point under the original quadratic-type frame hypotheses is asserted, and no contradiction to the accepted original source theorem is claimed. In particular, the absence of a LamPtClauses witness is not inferred merely from a failure to find one.

Thus there are four distinct assertions to keep apart:

\begin{longtable}[]{@{}
  >{\raggedright\arraybackslash}p{(\columnwidth - 2\tabcolsep) * \real{0.5000}}
  >{\raggedright\arraybackslash}p{(\columnwidth - 2\tabcolsep) * \real{0.5000}}@{}}
\toprule\noalign{}
\begin{minipage}[b]{\linewidth}\raggedright
Proposed shortcut
\end{minipage} & \begin{minipage}[b]{\linewidth}\raggedright
Remaining literal requirement
\end{minipage} \\
\midrule\noalign{}
\endhead
\bottomrule\noalign{}
\endlastfoot
Drop IsTen but keep CMObj & Pure A± and onto still imply CT1. \\
Bypass CMObj but keep LamObjGood & Pure complete Lie bases and IsIsogSum still imply CT2. \\
Use an unrelated elliptic Good object for a fixed mixed point & The fibre identification in LamPtClauses contradicts CT10. \\
Supply arbitrary mixed HasType for B & That alone supplies neither the required source factors nor the isogeny comparison. \\
\end{longtable}

The uniform-signature condition Datum.dom is a separate source interface. It is not identified here as a consumer of CMObj.

\subsection{7. The Jacobian consequence and the necessary replacement}\label{the-jacobian-consequence-and-the-necessary-replacement}

The carried R003 result establishes that \(J=\operatorname{Jac}(y^2=x^7-1)\) is geometrically simple of dimension three. Its actual CM type on \(\mathbf Q(\zeta_7)\) is \(\{1,2,3\}\), whereas the two types induced from its only proper CM subfield \(\mathbf Q(\sqrt{-7})\) are \(\{1,2,4\}\) and \(\{3,5,6\}\). The primitive-type classification gives this simplicity. Any nonzero power of \(J\), and any variety isogenous to such a power, has \(J\) as its only geometric simple factor.

Such a variety cannot be isogenous to a power of an elliptic curve. Indeed a nonzero homomorphism from an elliptic curve into a simple three-dimensional abelian variety would have positive-dimensional image, forcing that image to be the whole target, which is dimensionally impossible. The same argument, combined with reducibility, excludes the proposed elliptic power as a surjective source. Applying CT1 or CT2 gives the claimed failure of the unchanged good-point interface.

For an abstract mixed-type replacement, the exact obligations exposed here are:

\begin{longtable}[]{@{}
  >{\raggedright\arraybackslash}p{(\columnwidth - 2\tabcolsep) * \real{0.5000}}
  >{\raggedright\arraybackslash}p{(\columnwidth - 2\tabcolsep) * \real{0.5000}}@{}}
\toprule\noalign{}
\begin{minipage}[b]{\linewidth}\raggedright
Item
\end{minipage} & \begin{minipage}[b]{\linewidth}\raggedright
What needs a replacement proof
\end{minipage} \\
\midrule\noalign{}
\endhead
\bottomrule\noalign{}
\endlastfoot
Labels and signature & Replace the quadratic half by the actual CM type, and prove the corresponding complete eigenbasis and signature statements. \\
Geometric comparison & Construct the actual component morphisms and their surjective or quasi-inverse sum for the chosen polarized varieties and integral lattices. \\
Metrized eigenlines & Retain the actual integral eigenline lattices and all archimedean polarization norms. \\
Degree formula & Prove the required degree identities for those lines; the old single quadratic period SizeData₂.logOm cannot be imported from its name. \\
Arithmetic evaluation & Reconcile the modified good-point data with the actual sections, letters, comparison operators and Read/032 identities. \\
\end{longtable}

The elementary independence and length mechanisms are not asserted to require elliptic curves intrinsically. The source proof LamLenOfIsHerm.lean, 14--37, uses the polarization-pullback identity for each recorded letter and its actual Hermitian norm; ExistsLamLettersIsGood.lean, 25--66, has a general finite-set block enumeration. These are potential transferable pieces, not a completed generalized source theorem. Their original hypotheses still refer to LamObjGood.

Changing the host degree, the auxiliary primes or the length does not remove CT1--CT2 while these pure-type and isogeny clauses are retained. Conversely, this proof makes no assertion that a correctly reconstructed mixed-type interface is impossible. It gives no obstruction to higher-weight tensor correspondences and no conclusion about the genuine mixed theta kernels already constructed independently.

\subsection{8. Verification and literature}\label{verification-and-literature}

The proof is characteristic-zero geometry. A finite diagnostic checks only the concrete \(\mathbf Q(\zeta_7)\) type enumeration and the literal source fingerprints/clauses. It does not prove the general isogeny statements, a Lean theorem, a moduli existence theorem or any global period nonvanishing.

The prior DP25 analysis already excluded the balanced Jacobian example using Newton slopes. The contribution here is the characteristic-zero argument for every target in these exact interfaces, and specifically CT10 at the actual pointwise hypothesis used by lamw\_link. The directly read Formalism copy and its fibre identification prevent an abstract Good object from supplying that hypothesis. The older example is not counted again as a new result.

\textbf{Primary reference.} Milne, J. S. (2020). \emph{Complex multiplication} (v0.10, July 14), Theorem 2.12, pp.~26--27; Proposition 3.12 and the proof of Proposition 3.13, p.~30. \url{https://www.jmilne.org/math/CourseNotes/CM.pdf}

The relevant primary text and the exact source files above were directly read. No large-scale archive audit, Lean build, source mutation or claim of final main-task resolution is made.

\subsection{Canonical algebraic-geometric context}\label{canonical-algebraic-geometric-context-1}

The conclusion concerns characteristic-zero fibres of the given abelian schemes with their endomorphism actions. Up to rational isogeny their degree-one CM motives have only the quadratic elliptic simple factor; this statement forgets the specified integral model and polarization. The Picard-dual and biextension structures remain those of the actual varieties and cannot be transferred to the primitive-Jacobian fibre merely by changing a label. No claim is made that higher tensor motives or reconstructed mixed moduli are impossible. Exact inherited contextual anchors: SGA3 I4.7.3/4.7.3.1; SGA6 VI Theorem3.3; FGA232 Theorem3.1 pp148--149; FGA236 Corollaries1.7/2.3 pp229/232; SGA7 I VII2.7/2.9 pp169--174 and VIII7.4(b) pp310--311; cote4 scan003 §§8.1.1--8.1.2 and scans005--006 §§8.4.1--8.4.3; cote54 scans005--007. These are contextual locators, not newly established comparison theorems. DP\_MULT Q3e and THETA3 G1/A7 retain their actual-source/period/accounting obligations; CX\_GRH CX-ALG-024 and GCT\_PLETH T-052 retain their separate gates. No new lambda-ring or K-theory operation, l-adic weight theorem or motivic L-function identity is inferred.

\clearpage

\section{AM-K097 --- the small line-denominator rank-three sector at support25}\label{am-k097-the-small-line-denominator-rank-three-sector-at-support25}

\textbf{Canonical checked integration --- 2026-11-23T00:51:48.580659+00:00.} Complete root and independent proof/script/interface reviews and separate exact replays passed. The actual canonical predecessor has92/31 entries; this integration gives93/31, complete floors25/11. Frozen author history is preserved. AM-K097 removes the complete actual global-rank-three scalar-cross support25 subset with N\_E D\_eta\textless=2\^{}120, using the cross-only pole lemma, at most24 eligible places, all25 padding choices and all15 actual type profiles. The complete normalized bound is below2\^{}-11658. No rank-three full-double-pole or unipotent central-fibre assertion is imported. One11658 charge appends to the actual DO10392/31 prefix, giving93/31; both full exact sums remain below2\^{}-230 and complete floors remain25/11. Remaining rank-three terms have N\_E D\_eta\textgreater2\^{}120 and retain R082 cuts. The signed kernel, actual selected arithmetic image, prescribed global period, mixed source construction, specific032/Read, changed-family DP44 and RH remain open.

Author: \nolinkurl{/root/kummer_extension}, 23 November2026 UTC. Controlled pending proof. Author whole saved-proof/script verification and independent review are recorded separately. No canonical file is modified.

\subsection{1. Exact unchanged sector and conclusion}\label{exact-unchanged-sector-and-conclusion}

Retain the actual J56 all-top kernel, H=133,f=168, support s=25, the original E/F, every CM partner and G/T prime, finite projectors, selected real trace and metric, compact factors, packet representative, central phase, Haar measures and class masses. Work in the DO099/R071 rank-three scalar-cross residual used in R081/R082:

\[
 E^{169}=Z\perp P,\quad Z=S_{164}\perp u^\perp,\quad
 P=Eu\perp Ee_n,\quad J_P=\operatorname{diag}(U,-1),
\] \[
 h=\begin{pmatrix}I_Z&B\\ C&D_P\end{pmatrix},\quad
 C|_S\ne0,\quad\operatorname{rank}_E(h-I)=3,
 \quad Cx=(u r(x),\eta r(x)),\quad\eta\bar\eta=U,
 \quad U=77Q^2-14.                                           \tag{RP1}
\]

Here Q=1427908193729121522601, N=h-I, and the original u is primitive. The scalar eta is uniquely normalized by the first plane coordinate, not assumed integral. All334 unmixed G reductions are unipotent; global unipotence is not assumed. Put

\[
 \mathfrak D_\eta=\prod_{\mathfrak p}\mathfrak p^{\max(0,-v_{\mathfrak p}(\eta))}.
                                                                    \tag{RP2}
\]

\textbf{Theorem.} The entire RP1 sector satisfying N\_E D\_eta≤2\^{}120 has total absolute actual normalized contribution less than

\[
                              2^{-11658}.                    \tag{RP3}
\]

After removing its intersection with the current remainder, every surviving rank-three support25 term has N\_E D\_eta\textgreater2\^{}120, in addition to the retained R082 restrictions. This removes every coordinate-support size in the stated denominator sector. It does not imply global unipotence, a full support-floor increase, positivity of the signed sum, an arithmetic image witness or a period/DP44 conclusion.

\subsection{2. The actual cross-pole condition does not require unipotence}\label{the-actual-cross-pole-condition-does-not-require-unipotence}

Use DO103 LD4--LD8's exact original scalars and paired pencil generators:

\[
 a_0=Q(-7+2\sqrt{-7}),\quad a_0\bar a_0=U+14,
 \quad\delta=\eta-a_0\ne0.                                  \tag{RP4}
\]

At the selected pi-component of an active T place the nonzero-label pencil pole line is the reduction of v0=(Qu,pi\_k), whose slope in (u,e\_n) is a0. At its conjugate component the actual inverse-adjoint dual has pole column congruent to -v0; its slope is U/bar(a0). The frame determinant has modulus2U, a unit at G and T, and both a positive coordinate of v0 and its negative coordinate are units. These are LIFT44 PR4--PR7, K077 AP1--AP3 and DO103 LD5--LD8, with the same original labels. No independent conjugate label is introduced.

Every original S input and each of the three fixed integral transverse inputs y1,y2,z0 of R072 PJ4--PJ5 belongs to every actual source and target pencil at both CM partners. If hM\_c=M\_d, then h\_ZZ=I and that membership give

\[
                    Ny=Cy=hy-y\in M_d.                     \tag{RP5}
\]

This identity uses neither N²Z=0 nor a nilpotent normal form. For target label d=0 the target is the standard lattice, so Cy is integral. For any nonzero target a nonintegral Cy has at most a simple pole with leading vector on the target pole line. The saturated image line of C is E(1,eta). If eta has positive local valuation its saturated residue line has zero negative coordinate; if its valuation is negative all positive coordinates vanish. Both contradict equality with the prescribed pole line. Thus a cross pole forces eta to be a unit and forces the appropriate slope equality.

At the selected component this is delta divisible by pi. At the opposite component use the exact identity

\[
               \bar a_0\eta-U=-\eta\bar\delta.              \tag{RP6}
\]

Since eta is a unit there, the opposite slope condition is equivalent to divisibility of delta at the same original selected component. Thus every cross pole, at either partner, is supported on the common eligible set

\[
                  E_0=\{v\in\Sigma:v_{\mathfrak P_v}(\delta)\ge1\}.
                                                                    \tag{RP7}
\]

The norm relation eta bar(eta)=U makes eta a unit at both partners of each eligible place. The denominator ideal is trivial there.

At a type-I place h preserves the full pencil and hence its intersection M, which contains each displayed input and is contained in the original lattice. Therefore Cy=hy-y is integral. This additional observation also requires no unipotence. The proof below only needs the containing pole set RP7; any extra integrality is allowed to enlarge the count.

Primitive integer Bezout coefficients of u extract r(y) from the original positive coordinates of Cy without a global U denominator. The negative pivot beta=eta r(e\_j) is already an original coordinate. Accordingly all167 counted r-values and beta have at most simple paired T poles supported on E0, and retain original integrality away the active support.

Only this cross-block conclusion is carried from the pencil. The DO103 assertion about the full double-pole part used the global rank-two unipotent PP normal form. It is not asserted here. In particular the full projected negative-column cost remains T\^{}100 below.

\subsection{3. The exact denominator cap and the padded common ideal}\label{the-exact-denominator-cap-and-the-padded-common-ideal}

The argument DO103 LD12--LD16 applies to any global eta with eta bar(eta)=U. It does not use the rank or characteristic polynomial of h. Explicitly, (delta) D\_eta is integral. At seven there are two conjugate pairs of E primes, each of ramification index6 and residue degree14. Their eta valuations are a and6-a, whereas a0 has valuation3 at each. The denominator-corrected sum in each pair is at least3. Hence the norm of that integral ideal is divisible by7\^{}84.

At every eligible active place it is divisible by the selected degree-one T prime, where the denominator is a unit. At all complex labels \textbar delta\_tau\textbar\textless2sqrt(U+14). Therefore

\[
 7^{84}T^{|E_0|}
 \le N_E((\delta)\mathfrak D_\eta)
 <[4(U+14)]^{168}N_E\mathfrak D_\eta.                       \tag{RP8}
\]

The actual exact inequality

\[
                   7^{84}T^{25}>2^{120}[4(U+14)]^{168}       \tag{RP9}
\]

proves \textbar E0\textbar≤24 under the theorem's small-denominator condition. All denominator primes, including those above seven and dividing U, are retained. The field-level existence of some norm solution for U does not supply a bound on this actual denominator.

Pad E0 to a24-element subset D of its25-place active support. There are at most25 possible such sets. Let J be the same canonical whole-C G ideal of R072/R076/DO100, with its original at-most2\^{}668 patterns, and define

\[
 L=J I_D,\quad N_E I_D=T^{-48},\quad v=24,
 \quad z=(N_E L)^{1/(2f)}\ge T^{-24/f}.                     \tag{RP10}
\]

Every counted r-value and beta belongs to this L. The common-G orientation is unchanged and remains valid for the full cross map in the nonunipotent rank-three case. No extra orientation is selected for the transverse inputs or pivot. Padding is a counted covering of the actual matrices; no denominator ideal or eta choice is separately enumerated.

\subsection{4. Complete original matrix and type factors}\label{complete-original-matrix-and-type-factors}

Write M=H/(2sqrt U), B0=H/2 and D0=7\^{}140 29\^{}144 5\^{}84 13\^{}84. R076 CV7--CV14 and R082 TP5--TP10 apply to the same L. The full167-dimensional original covector ball has dual Gram diag(I164,G\^{}-1), detG=4U, G-I3\textgreater0, and covolume D0\^{}167 z\^{}(334f)/(4U)\^{}f.~Its actual arbitrary-ideal cover is less than4200sqrt(f)z. Consequently

\[
 \#r\le A(4U)^f(M+4200z)^{334f}z^{-334f},\quad
 A=\frac{528^{167f}}{(167f)!D_0^{167}}<2^{-387281}.           \tag{RP11}
\]

This includes all zero patterns in one ball; there is no k-binomial or k union. Actual C\textbar S is nonzero, so a nonzero original S pivot exists. Its at-most164 choices are retained. The exact relative-norm fibre gives

\[
 \#\beta\le28 C_F B_0^{f/6}z^{-f/6},\quad
 C_F=\frac{2^{300}}{29^6 41^{14}2129^{14}}<2^{42}.           \tag{RP12}
\]

Nonprincipal and non-conjugation-invariant ideals are included by the actual relative ideal norm count. The selected trace and remaining reconstruction bounds stay

\[
 |c(h)|<2^{-929f/2}U^{-f/2}z^{-f},\quad
 \#\operatorname{pr}_P(he_n)<191^{4f}U^fT^{100},\quad G_2<2^{36}.
                                                                    \tag{RP13}
\]

The full paired25-place projection cost and its fixed U denominator are unchanged. No central affine-fibre or rank-two reconstruction is substituted.

The weak and strong pointwise Green bounds on the same actual finite factor are

\[
 |a(h)|\le2Q^{-335\cdot167}T^{335v}z^{335f},\qquad
 |a(h)|<2Q^{-501\cdot167}2^{9\cdot167f}T^{167t}.             \tag{RP14}
\]

The first expresses N\_EJ=T\^{}(2v)z\^{}(2f); the second is the original nonzero three-minor estimate DO095 RA15--RA20. The two mixed finite factors retain modulus bound1. No new witness enumeration is introduced.

The actual type-III count satisfies11≤t≤25: DO074 TC1 makes t≤10 globally unipotent of rank at most two, contrary to RP1. R082 TP3--TP4 gives the complete paired type/label weight

\[
 {25\choose t}(4T)^{25-t}<2^{P_t},\qquad
 P_t=75-2t+(25-t)L_T,\quad L_T=1019333/1024.                \tag{RP15}
\]

This uses the general type-I bound3T and the exact type-II and type-III weighted cost1, with no global-unipotence premise. For each fixed original pair and padded set the arithmetic count is applied first, then these weights are summed. The complete coarse factor2\^{}(997*25) is replaced by RP15, not multiplied by it. The remaining fixed factor is exactly2\^{}881: Green2, G patterns2\^{}668, support2\^{}168, pivots2\^{}8 and source fibre2\^{}36. There is no independent conjugate label count.

\subsection{5. Two full norm ranges and exact complete exponents}\label{two-full-norm-ranges-and-exact-complete-exponents}

Split the padded ideal norms at z0=2\^{}-98. The exact original integers give M\textgreater2\^{}-68 and

\[
 4200z_0/M<4200/2^{30}<2^{-17},\qquad
                  (1+2^{-17})^{334f}<2.                    \tag{RP16}
\]

Thus the entire cover enlargement in RP11 is less than a factor2 for z≤z0. In the low range, multiplying RP11--RP14 with the weak Green bound leaves exactly z\^{}(-f/6). The original integrality of J in RP10 bounds this by T\^{}(v/6)=T\^{}4. This lower-endpoint step is essential: the exponent is negative and is not maximized at z0.

In the high range use the strong Green bound. Its whole varying expression is

\[
                        (M/z+4200)^{334f}z^{-7f/6},         \tag{RP17}
\]

which decreases for z\textgreater0. It is bounded at z0, where RP16 supplies the same factor2. No ideal attaining the endpoint is assumed.

Use the directed constants LQ=287843/4096\textless log2Q, LH=28900/4096\textgreater log2H, LU=601356/4096log2(191). Define the full archimedean common constant

\[
 C_* =881+1+5+42-387281+4fL_{191}-929f/2+2f
       +(2005f/6)(L_H-1)-(331f/2)L_U.                        \tag{RP18}
\]

The1 is the complete cover inflation,5 bounds28,42 bounds CF, and2f is the determinant factor4\^{}f.~The U coefficient retains both the projected column and the genuine selected trace. The complete low and high exponents, before the final25 padding choices, are

\[
 E_t^-=C_*-335\cdot167L_Q+(335v+v/6+100)L_T+P_t,
\] \[
 E_t^+=C_*-501\cdot167L_Q+9\cdot167f
              +(167t+100)L_T+P_t+98(2011f/6),\qquad v=24.  \tag{RP19}
\]

The companion computes all30 rational cells, with exact original integer bounds for every directed constant and the full factorial volume. Their maxima are

\[
 \max_{11\le t\le25}E_t^-=-47790211/4096<-11667\quad(t=11),
\] \[
 \max_{11\le t\le25}E_t^+=-223030653/4096<-54450\quad(t=25).
                                                                    \tag{RP20}
\]

All original coordinate-support sizes were included in RP11, so no extra164 factor appears. We conservatively sum both complete norm-range bounds and all15 actual type counts, then retain all25 padding choices. The exact dyadic inequality

\[
           375\bigl(2^{-11667}+2^{-54450}\bigr)<2^{-11658}    \tag{RP21}
\]

proves RP3. Using the full pattern count on both sides of the split only enlarges this explicit bound. No failed sufficient bound is used as a lower bound or an existence argument.

\subsection{6. Chronology, interfaces and remaining obligations}\label{chronology-interfaces-and-remaining-obligations}

The intended chronological predecessor is actual DO103, checks/do103\_joint\_remainder\_error.json, with92 all-top and31 uniform errors, after R0822353 and DO1031878. The script accepts an explicit canonical --prior that retains this full prefix and every later entry. Append11658 once, preserve all uniform terms and both complete floors25/11, and check the entire ordered sums below2\^{}-230. All cuts are applied only to their intersections with the current remainder; overlapping earlier charges are retained.

The new mathematical interface is RP5--RP7: the actual cross-pole line condition survives without global unipotence. Its common eligible-set bound uses the already checked DO103 denominator-corrected norm proof. The complete count then uses R076 CV4--CV18, DO100 HI1--HI14 and R082 TP3--TP11 with v24 in the cross ideal, actual support25 in the projected column, and the original full167-dimensional metric. No full double-pole or central-entry restriction is inferred for rank three.

All modules remain inside the existing original Hom space over Spec(O\_E); only a containing fractional ideal changes in the count. The theta line, Picard/biextension data, polarization and complex realization are unaltered. Carried contextual anchors remain SGA3 I4.7.3/4.7.3.1; SGA6 VI Theorem3.3; FGA232 Theorem3.1 pp148--149; FGA236 Corollaries1.7/2.3 pp229/232; SGA7 I VII2.7/2.9 pp169--174 and VIII7.4(b) pp310--311; cote4 scan003 §§8.1.1--8.1.2, scans005--006 §§8.4.1--8.4.3 and cote54 scans005--007. No new comparison, motivic, K-theoretic, l-adic or source-constructor theorem is asserted.

The complementary large-denominator rank-three sector, the remaining global-rank-two branch, full finite signed inequality, actual restricted selected arithmetic image, prescribed periods, genuine mixed source admissibility, Read/specific032 and changed-family DP44 comparison remain separate. This finding supplies no global nonvanishing or RH conclusion and performs no Lean build or archive audit.

\clearpage

\section{AM-R-083 --- a complete large-line-denominator tail for dense rank three at support25}\label{am-r-083-a-complete-large-line-denominator-tail-for-dense-rank-three-at-support25}

\textbf{Canonical checked integration --- 2026-11-23T00:53:17.426242+00:00.} Complete root and independent proof/script/interface reviews and separate exact replays passed. The actual canonical predecessor has93/31 entries; this integration gives94/31, complete floors25/11. Frozen author history is preserved. AM-R-083 removes the entire actual global-rank-three scalar-cross support25 tail N\_E D\_eta\textgreater=2\^{}4864 with k\textgreater=126, the latter already forced on the current remainder by R082. Its full167-dimensional arbitrary-ideal covector count, exact per-denominator eta enumeration and complete ideal-zeta tail retain all original type weights, G patterns, selected traces, projected U/T column and source fibre. The complete exponent is -22330263/4096\textless-5451, charged conservatively once as5400. The actual K09793/31 prefix becomes94/31; full sums remain below2\^{}-230 and floors25/11 are unchanged. Together with K097, remaining rank-three support25 terms satisfy2\^{}120\textless N\_E D\_eta\textless2\^{}4864 and all R082 profile cuts. The all-II unipotent rank-two, signed kernel, selected arithmetic image, mixed source, prescribed global period, specific032/Read, DP44 and RH gates remain separate and open.

Author: \nolinkurl{/root}, GPT-6 Astra Pro, ALP\_MACHINE,23 November2026 UTC; inherited effort not separately exposed. Controlled pending proof. Whole saved-proof readings and separate exact replays precede promotion.

\subsection{1. Exact actual scope}\label{exact-actual-scope-2}

Keep the original J56 H133 all-top scalar-cross support25 sector of R081/R082: h\_ZZ=I on the positive167-space Z, the cross map C is nonzero on the original S164, global rank(h-I)=3, and all334 unmixed G reductions are unipotent. Keep the actual paired pencils, normalized trace, compact factors, packet, central phase, class masses, source/target and Haar. Write

\[
 Cx=(u\lambda(x),\eta\lambda(x)),\quad \eta\bar\eta=U,
 \quad U=77Q^2-14,\ Q=1427908193729121522601,
 \quad f=168,\quad M=H/(2\sqrt U).                 \tag{LD1}
\]

Let k count nonzero original S coordinates of lambda, and let D\_eta be the actual denominator ideal of eta, with d=N\_E D\_eta. The theorem concerns k\textgreater=126. Every remaining rank-three support25 term after R082 has this property, because its smallest K\_t is125. DO074 TC1 gives11\textless=t\textless=25 for its actual number t of type-III places.

\textbf{Theorem.} The entire stated sector with

\[
                  k\ge126,\qquad d\ge2^{4864}              \tag{LD2}
\]

has absolute total normalized contribution less than2\^{}-5400. Thus, after retaining R082's complete previous cuts, every remaining rank-three support25 term has d\textless2\^{}4864. No lower bound on d or full support25 exhaustion follows from this theorem.

\subsection{2. Count eta before its actual covector}\label{count-eta-before-its-actual-covector}

For any fixed integral ideal D which is the exact denominator of a solution eta*bar(eta)=U, the integral ideal A=(eta)D is coprime to D and satisfies

\[
                    A\bar A=(U)D\bar D.                   \tag{LD3}
\]

At a conjugate pair p,bar(p), put v\_p(U)=e\textgreater=0 and v\_p(eta)=a, so the other valuation is e-a. Specifying the negative parts of these two valuations specifies both valuations whenever one is negative. If neither is negative, there are at most e+1 choices, namely a=0,\ldots,e. A conjugation-fixed prime has2a=e and at most one possibility. This proves that every choice away from the prime divisors of U is forced and that the total number of numerator ideals is bounded by the product of e+1 over the relatively split F-primes dividing U. Nonprincipal candidate fractional ideals contribute no eta.

The exact factorization and actual field decomposition in the fully read packet-plane arithmetic proof, with its recursive Lucas and Frobenius certificate, are

\[
 U=3^3\,7\,12837557489173\,111395838330757\,580868899958147.
                                                               \tag{LD4}
\]

All five rational primes are relatively split in E/F. Their numbers of F-primes and valuations of U are respectively

\[
                (4,3),\ (2,6),\ (4,1),\ (12,1),\ (4,1).
                                                               \tag{LD5}
\]

For the four primes other than7, the absolute residue degrees are42,42,14,42 and ramification is1. At7 the decomposition group has order84 and inertia6, giving two F-primes of residue degree14 and valuation6. These are computations in the actual degree336 field K032 QH9--15 and the real degree7 conductor29 quotient K030 PH2, not generic splitting assumptions. The companion independently rechecks the Lucas factorization and these group computations.

For each candidate principal fractional ideal, any two actual generators with exact norm U differ by a relative-norm-one unit. The checked CM unit calculation gives28 such roots of unity. Consequently, for every fixed D,

\[
        \#\{\eta: \eta\bar\eta=U,\ D_\eta=D\}
               \le28\,4^4\,7^2\,2^{20}<2^{39}.             \tag{LD6}
\]

There is no class-number multiplier: candidate ideals are counted before testing principality. This argument needs no existence theorem for a norm solution.

\subsection{3. The exact denominator content and the original trace}\label{the-exact-denominator-content-and-the-original-trace}

Fix one canonical whole-C G pattern with integral ideal J. Put L=J I\_Sigma, N\_E I\_Sigma=T\^{}-50 and z=(N\_EL)\^{}(1/(2f)). The original primitive-u extraction and target-pencil membership in R072 PJ4--7 apply to every original S and every one of its three fixed integral transverse inputs y. They give both lambda(y) and eta lambda(y) in L. The latter is the original negative coordinate of Cy. Thus, valuation by valuation,

\[
 \lambda(y)\in L\cap\eta^{-1}L=L D_\eta,
 \qquad z'=(N_E(LD_\eta))^{1/(2f)}=z d^{1/(2f)}.           \tag{LD7}
\]

Indeed the valuation of the intersection is v(L)+max(0,-v(eta)). This holds for nonprincipal and non-conjugation-stable ideals; no moving frame is made integral. Fixing eta and all167 values determines C. We keep the earlier harmless2\^{}8 pivot factor.

The actual joint row norm is still \textbar\textbar lambda\_tau\textbar\textbar\textless M. The k nonzero original S values have total half-trace squared norm at least kf(z')², by the norm/AM--GM minimum for the ideal LD\_eta. Hence

\[
                       z'<M/\sqrt{k}\le M/\sqrt{126}.      \tag{LD8}
\]

For a nonzero original pivot xi\_j in this same ideal, R061 GS11--12's proof now gives the genuine selected trace

\[
                |c(h)|<2^{-929f/2}U^{-f/2}(z')^{-f}.       \tag{LD9}
\]

This follows by replacing its product lower bound \textbar N\_E xi\_j\textbar\^{}(1/2)\textgreater=z\^{}f by(z')\^{}f, with the identical Cartan comparison and both compact factors. It is not a change of operator or normalization.

\subsection{4. Full covector volume, cover and every remaining factor}\label{full-covector-volume-cover-and-every-remaining-factor}

R076 CV8--10 applies to the actual ideal LD\_eta. The fixed metric is diag(I164,G\^{}-1), det G=4U and G-I\textgreater0. Its covolume is D0\^{}167(z')\^{}(334f)/(4U)\^{}f and full covering radius is at most4200sqrt(f)z', where D0=7\^{}140 29\^{}144 5\^{}84 13\^{}84. Thus

\[
 \#\lambda\le A(4U)^f(M+4200z')^{334f}(z')^{-334f},
 \qquad A=528^{167f}/((167f)!D_0^{167})<2^{-387281}.         \tag{LD10}
\]

This full count includes all original coordinate supports, so there is no k union or binomial. For a nonempty actual sector LD8 gives

\[
       1+4200z'/M<1+4200/\sqrt{126}<383,
                     383^{256}<2^{2197}.                  \tag{LD11}
\]

If the actual sector is empty its contribution is zero; admitting extra covectors, including zero, in LD10 is harmless after applying this necessary bound on the fixed ideal. No small-cover approximation is used.

The original weak Green bound on the same J is

\[
 |a(h)|\le2Q^{-335\cdot167}T^{8375}(z')^{335f}d^{-335/2}.
                                                               \tag{LD12}
\]

It is exactly R072 PJ6, because N\_EJ=T\^{}50(z')\^{}(2f)/d.~The product of LD10, LD9 and LD12 has zero exponent of z':335f-334f-f=0. The remaining cover factor is bounded by LD11. Replacing the old beta norm fibre by the complete eta enumeration LD6 is the only scalar-count change.

Retain the original projected negative-column count191\^{}(4f)U\^{}f T\^{}100 and the positive rank-two source fibre2\^{}36. Retain all G patterns2\^{}668, support choices2\^{}168, pivots2\^{}8 and the Green factor2; together with the source fibre these have exponent881. No source-fibre uniqueness improvement or projected U denominator is discarded.

For fixed actual t, R082 TP3--4's complete general type/pair sum is bounded by2\^{}(75-2t)T\^{}(25-t). Type I uses the general bound\textless3T, while types II and III cost1 each. This includes every arrangement and paired label. There are15 possible t. Their directed exponents P\_t=75-2t+(25-t)L\_T decrease strictly in t, so P\_t\textless=P\_11=53+14L\_T. The two CM components are not enumerated independently.

With the exact directed logs of R081/R082, the resulting bound for a fixed D, after all t, is at most15*2\^{}C d\^{}(-335/2), where

\[
\begin{split}
C={}&881+39-387281+4fL_{191}-929f/2+2f\\
 &+334f(L_H-1)-(331f/2)L_U-335\cdot167L_Q\\
 &+8475L_T+(53+14L_T)+(2197/256)334f.             \tag{LD13}
\end{split}
\]

The U exponent includes row determinant, projected column, selected trace and M; every original factor is present. All directed-log signs in LD13 are fixed by these displayed powers.

\subsection{5. Complete ideal tail and chronology}\label{complete-ideal-tail-and-chronology}

Sum over all integral E-ideals D of norm d\textgreater=d0=2\^{}4864, allowing even ideals which cannot be denominators. Unique factorization of ideals gives the usual Euler product. At each rational prime, every residue degree is at least1 and the number of prime ideals is at most336, so

\[
 \zeta_E(2)\le\zeta_{\mathbf Q}(2)^{336}<2^{336}.
                                                               \tag{LD14}
\]

The final strict inequality also follows directly from1+sum\_(n\textgreater=2)n\^{}-2\textless1+sum\_(n\textgreater=2){[}n(n-1){]}\^{}-1=2. Therefore

\[
 \sum_{N_ED\ge d_0}(N_ED)^{-335/2}
       \le d_0^{-331/2}\zeta_E(2)<2^{336-(331/2)4864}.
                                                               \tag{LD15}
\]

Using15\textless2\^{}4, the complete directed exponent is

\[
 C+4+336-(331/2)4864
             =-22330263/4096<-5451<-5400.                   \tag{LD16}
\]

This proves the full tail theorem. It sums every denominator, every eta for that denominator, every actual type profile, every support/pattern/pivot, the full covector ball, projected column and source fibre. No unproved distribution or cancellation is used.

The author base is checks/do103\_joint\_remainder\_error.json,92/31 errors. At integration append5400 once to the actual complete extension of that base; the --prior option verifies its entire ordered prefix and retains every intervening charge. Both full exact sums remain below2\^{}-230; complete floors stay25/11. The remaining rank-three small/intermediate-denominator profiles, all-II rank-two branch, signed inequality, source/selected-image/global-period/specific032 and DP44 gates remain open. No RH or global oscillator assertion follows.

\subsection{6. Geometry and exact contextual locators}\label{geometry-and-exact-contextual-locators}

LD7 is an exact intersection of fractional rank-one O\_E modules over Spec(O\_E). Its full dual positive-module realization retains the same fixed polarization Gram matrix. The norm-one fibre is counted by integral ideals and CM units, and the tail sums these genuine ideal norms. No theta line, Picard class, biextension, integral source object, metric, class mass or Haar is rescaled.

Carried contextual locators are SGA3 I4.7.3/4.7.3.1; SGA6 VI Theorem3.3; FGA232 Theorem3.1 pp148--149; FGA236 Corollaries1.7/2.3 pp229/232; SGA7 I VII2.7/2.9 pp169--174 and VIII7.4(b) pp310--311; cote4 scan003 §§8.1.1--8.1.2 and scans005--006 §§8.4.1--8.4.3; cote54 scans005--007. They supply the inherited geometric language, not new lambda-ring/K-theory operations, l-adic weights, CM-motive comparisons or L-function identities. DP\_MULT Q3e and THETA3 G1/A7 gain precisely the smaller finite/real gate; CX\_GRH CX-ALG-024 and GCT\_PLETH T-052 retain their independent arithmetic obligations. No Lean build or archive audit is performed.

\clearpage

\section{AM-DO-104 --- the large line-denominator tail at all-II support25}\label{am-do-104-the-large-line-denominator-tail-at-all-ii-support25}

\textbf{Canonical checked integration --- 2026-11-23T00:57:18.968927+00:00.} Complete root and independent proof/script/interface reviews and separate exact replays passed. The actual canonical predecessor has94/31 entries; this integration gives95/31, complete floors25/11. Frozen author history is preserved. AM-DO-104 removes the complete actual all-II globally rank-two-unipotent scalar-cross support25 tail N\_E D\_eta\textgreater=2\^{}1820. Every original coordinate lies in LD\_eta; the full original4-by4 block and both mixed blocks give the exact affine central difference ideal iJ0 B\_eta I\_Sigma,F\^{}2 with D\_eta bar(D\_eta)=B\_eta O\_E. The full real750 cover, all167 covector values, actual selected trace, all164 row cells, original weights and the entire ideal-zeta tail are retained. The complete maximum -10346711/4096\textless-2526 yields one2518 charge. Actual R08394/31 becomes95/31, complete sums remain below2\^{}-230 and floors25/11 remain. Together with DO103 the remaining all-II rank-two branch has2\^{}120\textless N\_E D\_eta\textless2\^{}1820. No whole support25, signed-kernel, prescribed selected arithmetic image/period, mixed source, specific032/Read, DP44 or RH conclusion.

Author: \nolinkurl{/root/definite_oscillator}, ALP\_MACHINE. 23 November2026 UTC.

Controlled pending proof. Author final-file verification, exact arithmetic replay, independent review and canonical integration are separate events. No canonical file is changed here.

\subsection{1. Exact sector and conclusion}\label{exact-sector-and-conclusion-1}

Keep the actual J56 all-top normalized kernel, H=133, f=168, all original fields, CM components, pencil arrows, finite projectors, selected real trace, packet representative, central phase, Haar measures and class masses. Work in the R071 full-positive-compression residual supplied at support25 by DO099(JR24):

\[
 E^{169}=Z\perp P,\quad Z=S_{164}\perp V_3,\quad
 P=Eu\perp Ee_n,\quad h_{ZZ}=I,\quad h=I+N,
 \quad \operatorname{rank}_E N=2,\quad N^3=0.                 \tag{CD1}
\]

The cross map C:Z→P is nonzero on S. Every active T place has type II, their number is s=25, and all334 unmixed G reductions are unipotent. Global rank-two unipotence is an explicit premise. Set

\[
 Q=1427908193729121522601,\quad U=77Q^2-14,\quad
 R=17424/133,\quad M=\sqrt{R/U},\quad
 D_0=7^{140}29^{144}5^{84}13^{84},\quad D_F=2^{168}D_0.
                                                               \tag{CD2}
\]

Use the unique normalized isotropic line C(Z)=E(1,eta) in the original basis (u,e\_n), so eta bar(eta)=U. Its exact integral denominator ideal is

\[
 D=D_\eta=\prod_{\mathfrak P}\mathfrak P^{\max(0,-v_{\mathfrak P}\eta)},
 \qquad x=N_ED,\qquad X=2^{1820}.                            \tag{CD3}
\]

\textbf{Theorem.} The entire CD1 sector with x≥X has total absolute actual normalized contribution strictly less than2\^{}-2518. Together with DO103, this leaves the explicit line-denominator interval

\[
                       2^{120}<N_ED_\eta<2^{1820}             \tag{CD4}
\]

inside this all-II branch. This theorem does not remove that interval or change the complete support floors25/11. It makes no assertion about nonunipotent rank three, fixed-S matrices, the signed sum, a prescribed period or the selected arithmetic image.

\subsection{2. Exact denominator descent and eta multiplicity}\label{exact-denominator-descent-and-eta-multiplicity}

The actual E/F is quadratic and unramified at finite places, with E=F(i) and root group mu28; these are the accepted K032/K050 data. At a conjugate pair write a=v\_P eta and e=v\_P U≥0. The other eta valuation is e−a. They cannot both be negative. At an inert place the valuation is e/2≥0. Consequently D and bar(D) are coprime, supported only at relatively split primes, and valuationwise

\[
 D\bar D=B_\eta\mathcal O_E,\qquad B_\eta\subseteq\mathcal O_F,
 \qquad N_FB_\eta=N_ED=x.                                   \tag{CD5}
\]

This constructs the actual integral F-ideal B\_eta; it does not require D to be principal.

For fixed exact D, put A=(eta)D. It is integral, (A,D)=1 and

\[
                  A\bar A=(U)D\bar D.                       \tag{CD6}
\]

If D has exponent b\textgreater0 at P, A has exponent0 there and e+b at bar(P), so no choice remains at that pair. At a U-prime absent from D, the two A-exponents are j,e−j for0≤j≤e. Outside U all exponents are forced. The directly checked actual factorization and Frobenius/inertia table in \nolinkurl{packet_plane_norm_boundary.md} (NB1--6) give

\[
 U=3^3\,7\,12837557489173\,111395838330757\,580868899958147.
                                                               \tag{CD7}
\]

Every F-prime above these factors splits in E/F. The F-prime counts and valuations of U are respectively (4,3),(2,6),(4,1),(12,1),(4,1): at seven the absolute ramification/residue degrees are6/14; at the four unramified factors the residue degrees are42,42,14,42. Thus the complete number of possible A for one D is at most

\[
                     K_U=4^4\,7^2\,2^{20}=49\,2^{28}.        \tag{CD8}
\]

For each fractional principal ideal A D\^{}-1 there are at most28 generators with exact relative norm U. The quotient of two such generators is an integral norm-one unit, hence one of the accepted28 roots of unity (K075 RN17--19/K050). Nonprincipal ideals give no solution. Therefore

\[
                         \#\{\eta:D_\eta=D\}\le28K_U<2^{39}.
                                                               \tag{CD9}
\]

There is no class-number factor, independently varying archimedean phase, or omitted numerator-ideal sum. The certificate replays the exact support-prime data and checks the local exponent alternatives; the previously checked primality certificates remain explicit inputs.

\subsection{3. One enlarged common ideal for all167 covector coordinates}\label{one-enlarged-common-ideal-for-all167-covector-coordinates}

The DO097/DO098 actual normal form is

\[
 C=\ell r,\quad \ell=(1,\eta),\quad v=J_Z^{-1}r^*,\quad q=rv>0,
 \quad
 N=\begin{pmatrix}0&-v\ell^*J_P\\ \ell r&d\ell\ell^*J_P\end{pmatrix},
 \quad d+\bar d=-q,\quad J_P=\operatorname{diag}(U,-1),
 \quad \alpha=(N^2)_{nn}=Uq\le R.                            \tag{CD10}
\]

No finite integral isotropic frame is inferred. Use the fixed transverse inputs y1=(-976627,-153032,38026,0), y2=(-567678,695570,-172838,0), z0=(256766,-7654233287,-30805128044,1), whose Gram matrix G has determinant4U. Each original S input and each of these three inputs belongs separately to every actual source and target pencil at both CM components, as in DO097(UC10--13)/DO098(CF5,CF8).

Fix the original whole-C G orientation of R072(PJ6), with at most4\^{}334=2\^{}668 patterns. Its E-ideal J satisfies

\[
 |a(h)|\le2Q^{-335\cdot167}(N_EJ)^{335/2},\qquad N_EJ\le Q^{334}.
                                                               \tag{CD11}
\]

Let I\_Sigma have exponent−1 at both CM primes of every active T place, so N\_E I\_Sigma=T\^{}-50, and L=J I\_Sigma. Primitive integer Bezout coefficients of the original u extract r(y) from the positive original coordinates of C(y); its negative coordinate is eta r(y). All are in L, using the same whole-C orientation, not an extra coordinate orientation. Thus exactly

\[
 r(y)\in L\cap\eta^{-1}L=LD\quad\text{for all167 fixed inputs},
 \qquad z'=(N_E(LD))^{1/(2f)}.                              \tag{CD12}
\]

The equality is the primewise identity v(r)≥v(L)+max(0,−v(eta)). The companion CM component is included, even when eta has poles. Given eta, these167 values determine r. With d they determine the full matrix through CD10; hence no projected negative-column or lattice-isometry fibre remains. We harmlessly retain the old first-pivot multiplicity164\textless2\^{}8 although eta is now counted first.

Let k be the number of nonzero original S values xi\_j=r(e\_j), 1≤k≤164. Their actual squared norm sums to at most M² at every embedding. The nonzero ideal minimum in the complex half-trace metric and CD11 imply

\[
 k(z')^2\le M^2,\qquad
 z'\le z_G x^{1/(2f)},\qquad z_G=Q^{167/f}T^{-25/f}.
                                                               \tag{CD13}
\]

All ideal metrics, CM norms and real dimensions here are the original ones.

\subsection{4. The strengthened original central affine difference lattice}\label{the-strengthened-original-central-affine-difference-lattice}

Fix eta,r and the G pattern. If d1,d2 occur, their difference Delta d is pure imaginary by CD10. In the original four positive coordinates and the negative coordinate, the differences of the full matrix are

\[
 (\Delta d)u_i u_j,\quad
 -\bar\eta(\Delta d)u_i,\quad
 \eta(\Delta d)u_j,\quad -U\Delta d.                         \tag{CD14}
\]

The other terms depend only on the fixed C and cancel. Let J0 be the F-ideal containing each rank-zero G prime once. Both N matrices have their full entries divisible at these primes, pole order at most2 at both partners of every type-II T place, and integral entries elsewhere. Hence their differences lie in the original entry ideal K\_E=J0 O\_E I\_Sigma². Since u is a primitive integer vector, two applications of its integer Bezout functional to the4×4 block, and one to each cross block in CD14, give

\[
          \Delta d,\quad\eta\Delta d,\quad\bar\eta\Delta d\in K_E.
                                                               \tag{CD15}
\]

This extracts Delta d itself; it does not divide an integral matrix entry by U. At any paired prime, with eta valuations a,e−a and e≥0, the extra exponent required by CD15 is max(0,−a,a−e), exactly the exponent of D bar(D). At an inert prime the extra is0. Since i is an integral unit and Delta d∈iF, contraction gives the containing difference lattice

\[
          \Gamma_d=iJ_0B_\eta I_{\Sigma,F}^2,\qquad
          \operatorname{covol}(\Gamma_d)
                =\sqrt{D_F}(N_FJ_0)xT^{-50}.               \tag{CD16}
\]

The fibre is contained in one affine coset if it is nonempty. No claim places d or Im(d) in this difference ideal. The affine real-part shift −q/2 is retained, so there is no unproved half-lattice or dyadic cancellation. CD15--16 remain valid when D meets U, G or T.

The same original G table gives N\_FJ0≤N\_EJ. The accepted full-real-ideal transference estimate (DO091 TD20--22) therefore gives

\[
 \mu(\Gamma_d)<750\sqrt f\,((N_FJ_0)xT^{-50})^{1/f}
             \le750\sqrt f\,(z')^2
             \le\frac{750R}{Uk}\sqrt f.                    \tag{CD17}
\]

The actual negative entry w=−Ud satisfies \textbar Im w\textbar\textless H/2, by DO098(CF6--7). Thus \textbar Im d\textbar\textless H/(2U). Translated Voronoi cells of this affine coset lie in a ball of radius C\_k sqrt(f)/U, where

\[
 C_k=67+98300/k>H/2+750R/k,
 \qquad
 V_F(C)=\frac{((22/7)fC^2)^{f/2}}{(f/2)!\sqrt{D_F}}.
                                                               \tag{CD18}
\]

Consequently the complete central fibre obeys

\[
 \#d\le V_F(C_k)\frac{T^{50}}{U^f(N_FJ_0)x}
       \le V_F(C_k)U^{-f}T^{50}x^{-1}.                      \tag{CD19}
\]

Only the favourable N\_FJ0≥1 factor is dropped. This is a full affine lattice count with a proved covering radius, not a covolume-density assertion. The real dimension is f and the codifferent normalization is the full real one.

\subsection{5. Complete row, transverse, real and finite factors}\label{complete-row-transverse-real-and-finite-factors}

The checked CF17--19 transverse argument applies to the actual fractional ideal LD with no change of basis assumed integral. Since z'≤M and det G=4U, its complete count is

\[
 4A_3(4U)^fM^{6f}(z')^{-6f},\qquad
 A_3=\frac{528^{3f}}{(3f)!D_0^3}<2^{-4040}.                  \tag{CD20}
\]

Here the actual metric cover follows from (2250/7)\^{}2·3 tr(G\^{}-1)\textless1024\^{}-2 and (1025/1024)\^{}(6f)\textless4; it is valid for arbitrary fractional ideals, not just O\_E.

The actual selected real trace, with the same two compact factors and near-point normalization as DO098(CF22), satisfies

\[
                      |c(h)|<2^{-929f/2}U^{-f/2}(z')^{-f}.
                                                               \tag{CD21}
\]

Indeed alpha\_tau≥U\textbar xi\_j,tau\textbar² for one nonzero xi\_j∈LD, so its field norm is at least U\^{}f(z')\^{}(2f). This uses the same coefficient bound on the same selected module, even when its right side exceeds1. No interpolation is needed here.

R072(PJ10--12)/DO098(CF23) gives the complete original S-row bound

\[
 {164\choose k}47^k2^{2fk}(f!)^k
      \left(\frac{(87/32)M^2}{fk(z')^2}\right)^{fk}.         \tag{CD22}
\]

Eta is already enumerated by CD9; the former beta relative-norm fibre is not multiplied again. Conversely it is not simply omitted: eta,r,d reconstruct the full h as proved above.

For each fixed actual type-II pair, all these estimates are uniform. Summing at an active place over its at most T−1 diagonal labels against exact \textbar rho\textbar=1/(T−1) gives at most1 (K071 NB7). There is no independent conjugate label or extra type choice. Keep support2\^{}168, common patterns2\^{}668, first pivot2\^{}8 and the Green correction2, a fixed exponent845. Class/Haar factors are exactly those in the accepted normalized coefficient; the retained remaining factor has absolute value≤1.

By CD12, \[
 N_EJ=(z')^{2f}T^{50}/x.
\] The product of CD11, CD19--22 has the following dependence on x,z':

\[
                  x^{-337/2}(z')^{b_k},\qquad
                   b_k=(328-2k)f\ge0.                       \tag{CD23}
\]

Every paired T exponent is retained: the Green factor contributes335s and the central fibre2s, for total337s. After maximizing the row radius, the complete positive M power is334f. Together with the transverse, real and central U factors this becomes R\^{}(167f) U\^{}(-335f/2). This regrouping precedes all directed logarithmic substitutions.

\subsection{6. The full infinite denominator tail and its feasible endpoint}\label{the-full-infinite-denominator-tail-and-its-feasible-endpoint}

For each k set \[
 F_k(x)=x^{-337/2}
       \min\{M/\sqrt k,\ z_Gx^{1/(2f)}\}^{b_k}.             \tag{CD24}
\]

This is an upper envelope for every feasible z'; no assertion that either endpoint is attained is made. For x≥X, the minimum in CD24 grows by at most (x/X)\^{}(1/(2f)), whether or not a cap changes branch. Hence

\[
 F_k(x)\le F_k(X)(x/X)^{-\beta_k},\qquad
                    \beta_k=k+9/2>2.                       \tag{CD25}
\]

At k164 the z' exponent is0, so the same formula is exact in x. There is no union factor for the two cap branches.

Sum over all nonzero integral E-ideals, enlarging the set of exact eta denominators. Absolute convergence and CD25 give

\[
 \sum_{N_ED\ge X}F_k(N_ED)
 \le F_k(X)X^{\beta_k}\sum_{N_ED\ge X}(N_ED)^{-\beta_k}
 \le F_k(X)X^2\zeta_E(2)
 <F_k(X)X^2 2^{336}.                                        \tag{CD26}
\]

The elementary Euler bound used here is zeta\_E(2)≤zeta(2)\footnote{E:Q}\textless2\^{}336: at a rational prime each of its at most336 factors has residue norm at least p, and zeta(2)\textless1+integral\_1\^{}infinity(t\^{}-2)dt=2. This counts every ideal class without choosing generators. CD9 has already bounded eta for each exact denominator. No bounded-class or equidistribution input is used.

The exact actual-host comparison \[
                   R^{168}T^{50}>61^{168}U^{168}Q^{334}2^{1820}
                                                               \tag{CD27}
\] implies M²/{[}z\_G²X\^{}(1/f){]}\textgreater61. Therefore in evaluating CD24 at X we can use the single safe cap M/sqrt(q\_k), q\_k=max(k,61). This includes small k whose unconstrained endpoint M/sqrt(k) is infeasible.

\subsection{7. Exact164-cell certificate and total}\label{exact164-cell-certificate-and-total}

Retain the checked one-sided actual-host constants \[
 L_T=1019333/1024>\log_2T,\quad
 L_Q=287843/4096<\log_2Q,\quad
 L_U=601356/4096<\log_2U,\quad
 L_R=113/16>\log_2R.                                        \tag{CD28}
\]

After CD26 and CD27 the common logarithmic upper bound, before the two explicit k-dependent ceilings, is \[
\begin{split}
 C_*={}&845+2+39+2f-4040-929f/2
       +167fL_R-(335f/2)L_U\\
      &+337\cdot25L_T-335\cdot167L_Q
       +336-(333/2)\cdot1820
       =566578985/4096.                                     \tag{CD29}
\end{split}
\]

Here39 is the complete eta multiplicity,336 is the zeta factor, and333/2=337/2−2 includes the full infinite ideal sum. All R powers are positive when L\_R is substituted.

Let v\_k be the least strict integral binary ceiling of V\_F(C\_k), and c\_k the corresponding ceiling of the exact rational factor \[
 \frac{\binom{164}{k}47^k2^{2fk}(f!)^k87^{fk}}
 {32^{fk}f^{fk}k^{fk}q_k^{(164-k)f}}.                        \tag{CD30}
\]

The companion script uses exact factorials, powers and comparisons for all164 k. In V\_F(C\_k), sqrt(D\_F)=2\^{}84 7\^{}70 29\^{}72 5\^{}42 13\^{}42 is an integer; no floating logarithm is used. Its output gives \[
 \max_{1\le k\le164}(C_*+v_k+c_k)
  =-10346711/4096<-2526,
 \quad k=61,\quad v_{61}=1256,\quad c_{61}=-142107.           \tag{CD31}
\]

The164 original support-size cells are disjoint and exhaustive, and their complete union is \[
                 <164\,2^{-2526}<2^{-2518}.                 \tag{CD32}
\]

This proves the theorem. A positive or unproved bound in the complementary denominator interval is not interpreted as positivity, nonvanishing or existence of a surviving contribution.

\subsection{8. Exact dependencies and chronological scope}\label{exact-dependencies-and-chronological-scope}

The actual source/target and normal-form inputs are DO097(UC3--13), DO098(CF1--23), DO099(JR24), R071(GR6--12), R072(PJ4--12), K071(NB5--9), and the primitive original-coordinate extraction in R059/DO103. The actual field/support-prime input is the full \nolinkurl{packet_plane_norm_boundary.md} NB1--6 and its recursively certified primes, rooted in K032(QH9--15); norm-one generator rigidity is K075(RN17--19)/K050. Arbitrary complex-ideal cover is DO070; full real dual/covolume is DO091(TD20--22). The new steps are the exact eta denominator count, all167 LD constraints, the original4×4 central difference extraction and B\_eta descent, the full central cover, and the complete denominator tail with its feasible endpoint.

Append2518 once to the actual ordered all-top predecessor, retaining every prior error and all31 uniform entries. The default predecessor is \nolinkurl{checks/do103_joint_remainder_error.json} (92/31); the script accepts \nolinkurl{--prior} so an intervening canonical charge is retained in a separate unchanged-script replay. Floors remain25/11. A canonical integration must state the actual chronological prefix and preserve the frozen author history.

The ideals are invertible modules on the unchanged Spec(O\_E)/Spec(O\_F), with the original CM involution and polarization. Carried contextual anchors remain SGA3 I4.7.3/4.7.3.1; SGA6 VI Theorem3.3; FGA232 Theorem3.1 pp148--149; FGA236 Corollaries1.7/2.3 pp229/232; SGA7 I VII2.7/2.9 pp169--174 and VIII7.4(b) pp310--311; cote4 scan003 sections8.1.1--8.1.2, scans005--006 sections8.4.1--8.4.3, and cote54 scans005--007. No new source constructor, motivic comparison, signed-kernel conclusion, selected derivative-image primitivity, prescribed Bessel period or DP44/RH conclusion follows.

\clearpage

\section{AM-LIFT-080 --- arbitrary CM-type letters and an explicit unequal-degree Step5 bound}\label{am-lift-080-arbitrary-cm-type-letters-and-an-explicit-unequal-degree-step5-bound}

\textbf{Canonical checked integration --- 2026-11-23T00:58:26.119818+00:00.} Complete root and independent whole-proof/script/literal-interface reviews, with separate7560-assertion exact replays, passed. The frozen historical source remains unchanged. LIFT080 MT1--17 proves a generic Step5 transfer for actual mixed-type comparison lines with individually bounded normalized degrees d\_j\textgreater=d\_-. Rescaling only their archimedean norms by exp(d\_j-e) makes all auxiliary degrees e, keeps every finite ideal and local polynomial quantifier unchanged, and replaces La=L by L exp((e-d\_-)\emph{+). The exact cost is C(e)=(1/2+\textbar e\textbar+log(L)/2)(D+gn)+D log(rL)+(3D+gn)(e-d}-)\emph{+/2. Its optimum within this comparison family is C0+(-d}-)\_+ min(D+gn,(3D+gn)/2). Arbitrary-type isogeny derivatives and the literal arbitrary-set form\_letter\_le supply independence and norm mechanisms; explicit mixed products J\^{}(14h p) with complementary-action factors supply geometric inclusions/projections with N=L=1 and permutation U. The actual prescribed point/model/section, all-extension polynomial tests, effective eigenline degrees, coefficient-field/frame ties and fully revised source/DP44 room inequality remain required. The old LamPtClauses/LamObjGood pure-type boundary remains; no main-family constructor, finite charge, period/image/value, specific032/Read or RH conclusion follows.

Author: \nolinkurl{/root/lift_accounting}, 23 November 2026 UTC. \textbf{Author whole-proof/source/script review PASS, 2026-11-23 UTC; exact completion and final fingerprints are recorded in LIFT080\_AUTHOR\_VERIFICATION.json.} This result applies the unchanged generic \nolinkurl{lambda_lower_bound} to auxiliary metrized comparison lines. It does not construct the unchanged quadratic-type \nolinkurl{LamObjGood} or \nolinkurl{LamPtClauses} at a mixed point. No source declaration is edited and no Lean build is used.

\subsection{1. The exact generic source input}\label{the-exact-generic-source-input}

Fix a number field K of degree d, the genuine \nolinkurl{KeyIn K} object x, and write g=x.P.g\textgreater=1. Retain the source theta line T, determinant line A, finite line ideals, square row matrix U, nonzero polynomial G, integer r\textgreater=1, finite block dimensions m\_b\textless=r, nonnegative equivariance weights a\_bi, integer n\textgreater=0, real L\textgreater=1 and positive archimedean bounds S\_sigma. Put

\[
 D=\sum_b\sum_i a_{bi},\qquad V=gn,\qquad
 \mathcal S=\frac1d\sum_{\sigma:K\hookrightarrow\mathbf C}\log S_\sigma.
 \tag{MT1}
\]

Here D is a total polynomial weight, unrelated to the source record also called D. All logarithms in this report are natural.

Assume every premise of \nolinkurl{StatementS.Step5} except its equal-degree premise, using the \textbf{actual} good comparison lines ell\_j. In particular retain \nolinkurl{KeyCore}, \nolinkurl{LemmaHStatement}, the actual \nolinkurl{KeyLocAt} and \nolinkurl{KeyArchAt}, det U!=0, G!=0, every exact \nolinkurl{Step5Loc} condition and every \nolinkurl{Step5Arch} condition. These include the all-residue-extension polynomial integrality assertion and the prescribed unitary-frame supremum; neither follows just from the existence of integral tangent vectors.

Write

\[
 d_j=\operatorname{ndeg}(\ell_j),\qquad
 d_-\le d_j\quad(1\le j\le g).                              \tag{MT2}
\]

A proved lower bound d\_- is sufficient. The minimum of this finite set exists, but replacing an effective degree calculation by an unevaluated minimum would not give an effective DP44 estimate.

The exact source definitions are:

\[
\begin{split}
 \operatorname{ndeg}(I,n_\sigma)
  &=-\frac{\log N_{K/\mathbf Q}I}{d}
       -\frac1d\sum_\sigma\log n_\sigma,\\
 \operatorname{pDeg}(u)&=\frac{1+|u|}{2}.
\end{split}                                                   \tag{MT3}
\]

\nolinkurl{MLine.Good} requires only a nonzero fractional ideal and positive norms. The original equal-degree conclusion of Step5 is

\[
 \mathcal S\ge
 -\bigl(\operatorname{pDeg}(\ell)+(1/2)\log L\bigr)(D+V)
 -D\log(rL),                                                  \tag{MT4}
\]

under the additional premise ndeg(ell\_j)=ell/2 for all j. The symbol ell in MT4 is the scalar source \nolinkurl{logOm}, not one of the line objects.

\subsection{2. Exact normalization with all finite data fixed}\label{exact-normalization-with-all-finite-data-fixed}

\textbf{Theorem MT5.} Under the hypotheses of section1, for every real e set

\[
 \delta_e=(e-d_-)_+ =\max(0,e-d_-).
\]

Then

\[
 \boxed{\quad
 \mathcal S\ge-C(e),\qquad
 C(e)=\bigl(\tfrac12+|e|+\tfrac12\log L\bigr)(D+V)
       +D\log(rL)+\frac{3D+V}{2}\,\delta_e.
 \quad}                                                       \tag{MT5}
\]

No finite ideal, theta section, determinant line, map, polynomial, coefficient field, class representative, local level, Haar measure or S\_sigma is changed.

\textbf{Proof.} Define auxiliary metrized lines

\[
 \widetilde\ell_j.I=\ell_j.I,\qquad
 \widetilde\ell_j.n(\sigma)=\exp(d_j-e)\ell_j.n(\sigma),
 \qquad L_e=L\exp(\delta_e).                                  \tag{MT6}
\]

Their ideals are nonzero and their norms positive. There are exactly d complex embeddings, counted individually as in the literal sum, so MT3 gives

\[
 \operatorname{ndeg}(\widetilde\ell_j)=d_j-(d_j-e)=e.
 \tag{MT7}
\]

At every prime the same local generators, completions, extension maps, block matrices and polynomial values witness \nolinkurl{Step5Loc}: that definition depends on ell\_j only through ell\_j.I. In particular the determinant-ideal inequality is unchanged.

For the archimedean length clause, d\_j\textgreater=d\_- implies

\[
 L_e\widetilde\ell_j.n(\sigma)
  =L\ell_j.n(\sigma)\exp(\delta_e+d_j-e)
  \ge L\ell_j.n(\sigma).                                    \tag{MT8}
\]

The same positive Hermitian matrix H\_sigma therefore satisfies the required row-length inequalities. Its determinant, the actual \nolinkurl{gramBlock}, G's equivariance, every unitary-frame condition and S\_sigma are unchanged. This proves \nolinkurl{Step5Arch} with the auxiliary lines and L\_e. Also L\_e\textgreater=1. \nolinkurl{KeyLocAt} and \nolinkurl{KeyArchAt} involve only the original x,T,A and hence remain exactly the supplied ones.

Apply the literal theorem \nolinkurl{lambda_lower_bound : Step5} with logOm=2e, auxiliary lines MT6 and L\_e. Since pDeg(2e)=1/2+\textbar e\textbar{} and log L\_e=log L+delta\_e, its cost is

\[
 \bigl(\tfrac12+|e|+\tfrac12\log L+\tfrac12\delta_e\bigr)(D+V)
 +D\bigl(\log(rL)+\delta_e\bigr),
\]

which is exactly C(e). This proves MT5. The real exponentials are allowed in the norm field of \nolinkurl{MLine}; no algebraic scalar or valuation is attributed to them. They are auxiliary inequalities in a generic theorem, not a rescaling of the actual arithmetic operator or CM period. The normalized physical theta line T and determinant line A remain untouched. QED.

The auxiliary line need not satisfy the original \nolinkurl{IsVecLine} metric identification unless its scaling factor is one. Nor is it asserted to be the ell field of the original \nolinkurl{LamLetters.IsGood}. Those assertions are unnecessary: the conclusion is obtained by passing the function of auxiliary lines directly to generic Step5. The unchanged physical lines and their integral lattices remain the source of the hypotheses.

\subsection{3. Keeping the old parameter, and the exact best choice in this argument}\label{keeping-the-old-parameter-and-the-exact-best-choice-in-this-argument}

If one wants to retain a specified source value ell\_0=\nolinkurl{SizeData₂.logOm}, take e=ell\_0/2. Its original MT4 cost is retained, with the \textbf{explicit additional sufficient cost}

\[
 \boxed{\quad
 \Delta C=\frac{3D+gn}{2}\,(\ell_0/2-d_-)_+.
 \quad}                                                       \tag{MT9}
\]

This is an archimedean comparison penalty. It is not a finite denominator, a removed product-formula term or a new value of the actual theta multiplier. A later comparison using the old ell\_0 must add MT9 to its room inequality; it cannot silently invoke the old equality-based \nolinkurl{lamw_link} with a false \nolinkurl{LamDegSt} premise.

There is also a closed optimization of MT5. Put

\[
 u=(-d_-)_+,\qquad
 C_0=\tfrac12(1+\log L)(D+V)+D\log(rL).
 \tag{MT10}
\]

Then

\[
 \boxed{\quad
 \mathcal S\ge-C_{\rm opt},\qquad
 C_{\rm opt}=C_0+u\min\left(D+V,\frac{3D+V}{2}\right).
 \quad}                                                       \tag{MT11}
\]

This is the minimum \textbf{within the specific uniform equalization MT6 and the conservative inflation from d\_-}, using the literal pDeg. It is not a claim of optimality among all arithmetic height arguments.

If d\_-\textgreater=0, choose e=0 and delta\_e=0. If d\_-\textless0 and D\textgreater=V, choose e=d\_-; if d\_-\textless0 and D\textless=V, choose e=0. When D=V both choices, and every e between them, have the same cost. The zero cases n=0 or D=0 are included.

For completeness, when d\_-\textless0 the extra cost on e\textless=d\_- is −e(D+V), minimized at d\_-. On d\_-\textless=e\textless=0 it is affine with slope(D−V)/2. On e\textgreater=0 it is increasing with slope(D+V)+(3D+V)/2, unless both coefficients vanish. These three intervals prove MT11. If d\_-\textgreater=0, every extra term \textbar e\textbar(D+V)+delta\_e(3D+V)/2 is nonnegative and vanishes at e=0. Thus no endpoint or sign case is omitted.

All of MT5--11 are consequences of the unchanged generic theorem and exact definitions. They require neither a common quadratic CM period nor equality among the actual d\_j.

\subsection{4. Which letter constructions transfer to arbitrary CM types}\label{which-letter-constructions-transfer-to-arbitrary-cm-types}

The geometric inputs can be stated without the pure quadratic halves used by \nolinkurl{LamObjGood}. Let A\_nu and B be polarized abelian varieties over a characteristic-zero field K containing the required CM labels. Give each A\_nu a complete K-basis v\_nu,i of its tangent space by E-eigenvectors, with any finite prescribed eigenlabel set Phi\_nu. Let f\_nu,j:A\_nu→B be E-linear homomorphisms such that their sum

\[
 F:\prod_{\nu,j}A_\nu\longrightarrow B
 \tag{MT12}
\]

is an isogeny. Choose an actual K-basis of Lie(B). Form the square matrix U whose rows are the coordinates of df\_nu,j(v\_nu,i), ordered by their eigenlabels.

\textbf{MT13, independence.} This U is invertible. In characteristic zero an isogeny has invertible tangent map. Alternatively, a positive-multiple quasi-inverse G gives dG dF=N I with N nonzero, which proves injectivity; source and target dimensions agree. The rows of U are precisely the images of a complete product tangent basis. This argument does not use the shape of the CM types. It is the mathematical content of the source \nolinkurl{EqZeroOfPow.lamInd}, which uses the quasi-inverse sum and complete bases, not the quadratic root.

The label ordering is also independent of quadratic induction. The source \nolinkurl{ExistsLamLettersIsGood.blockEquiv}, lines25--44, enumerates a finite set and its complement for an arbitrary predicate. More generally one may take the disjoint union of all component label lists and enumerate each fibre. Its multiplicity at tau is the sum of the actual occurrences of tau. E-linearity ensures that these are the target eigenspaces. There is no substitution of an abstract dimension count for the actual maps.

\textbf{MT14, lengths.} Suppose f\_nu,j has the genuine polarization pullback

\[
 f_{\nu,j}^{\vee}\lambda_B f_{\nu,j}
       =\lambda_{A_\nu}\,\iota_\nu(a_{\nu,j}),
 \qquad
 |\tau(a_{\nu,j})|\le L^2
 \tag{MT14}
\]

at every required embedding. Let ell\_nu,i be the actual integral tangent eigenline with its polarization norm. The image vector then has norm at most L times that source norm. Indeed the pullback equality gives the Hermitian quadratic-form equality on the eigenvector; its scalar is the corresponding eigenvalue of a\_nu,j. Positivity and its absolute-value bound give the asserted inequality. The source theorem \nolinkurl{form_letter_le} states precisely this implication with an \textbf{arbitrary} set S of embeddings, using the exact all-torsion Weil-pairing version of MT14. It already includes the source polarization/metric normalization, so no missing pi or Gaussian factor is inserted.

If these maps extend over the specified integral models, the images of the actual tangent eigenline lattices lie in the target tangent lattice. Their wedge gives the determinant-ideal inequality in the first part of \nolinkurl{Step5Loc}. This does not supply its remaining polynomial-integrality quantifiers. Those, the Key comparisons, the prescribed point/section and the nonzero polynomial are retained inputs in section1.

Thus the actual isogeny and polarization identities suffice for the independence and length mechanisms for arbitrary CM types. The new MT5--11 shows that the later generic height argument needs only a lower bound on their actual line degrees, with the stated cost. None of these arguments declares the old quadratic-type records inhabited.

\subsection{5. Explicit mixed product letters and the remaining actual point gate}\label{explicit-mixed-product-letters-and-the-remaining-actual-point-gate}

For a concrete geometric example retain the checked principally polarized hosts

\[
 A_h=J^{14h},\quad h\in\{1,4\},\quad
 E_h=\mathbf Q(\zeta_7,i)B_7(29)R_h,
 \quad R_1=\mathbf Q,\ R_4=\mathbf Q(\sqrt5,\sqrt{13}).
 \tag{MT15}
\]

Their actual O\_Eh actions, product polarizations, CM types and period/metric coordinates are K030/K032/K035. Here {[}E\_h:Q{]}=2f and dim(A\_h)=f=42h. Let A\_h\^{}- denote the same polarized abelian variety with the action precomposed by complex conjugation of E\_h; its type is the complementary type. For nonnegative p,q with p+q\textgreater0 form

\[
 B=A_h^p\times(A_h^-)^q
 \tag{MT16}
\]

with the product polarization and the displayed component actions. The component inclusions and projections give MT12 with F an isomorphism and N=1. Their polarization pullbacks are exact with a=1. In the concatenated actual eigenbasis, U is the identity up to the specified label permutation, so det U=±1 and the length constant is L=1. The multiplicities are p on the actual type Phi\_h and q on its complement. These are genuine mixed CM-type letters, not elliptic substitutes.

The actual norms are the polarization norms. For example K035 PF9--10 identifies them in its algebraic coordinates through

\[
 H_h(\zeta,\zeta)=\sum_\tau
       \frac{\delta_\rho\pi_l^{-6}}{2G_i}|\zeta_\tau|^2,
 \qquad H_{h,\tau}|\Omega_\tau|^2=-2i\tau(\xi_h)>0.
 \tag{MT17}
\]

Any source Lie-metric convention is the one already included in \nolinkurl{form_letter_le}; MT17 is not used to replace a metrized integral tangent line by the embedding-coordinate ideal cD\_tau. K035 explicitly distinguishes that ideal from a primitive Hodge line at bad primes.

For a specified integral realization of these products, component inclusions preserve the direct-sum tangent lattice and the same argument gives MT14 on its actual lines. This does not assert that an arbitrary CM frame selected by \nolinkurl{cm_frame_of_eq} has fibre MT16, that a chosen coefficient field already contains every required model, or that it satisfies the original quadratic signature. In particular \textbf{CT10 of LIFT079 still rules out the unchanged LamPtClauses at such a nonelliptic mixed point}. MT16 is an explicit geometric replacement input, not a new literal LamPtGood instance.

To use MT5 for the actual admitted main family, the remaining requirements are therefore precise:

\begin{longtable}[]{@{}
  >{\raggedright\arraybackslash}p{(\columnwidth - 4\tabcolsep) * \real{0.3333}}
  >{\raggedright\arraybackslash}p{(\columnwidth - 4\tabcolsep) * \real{0.3333}}
  >{\raggedright\arraybackslash}p{(\columnwidth - 4\tabcolsep) * \real{0.3333}}@{}}
\toprule\noalign{}
\begin{minipage}[b]{\linewidth}\raggedright
Requirement
\end{minipage} & \begin{minipage}[b]{\linewidth}\raggedright
What this theorem supplies
\end{minipage} & \begin{minipage}[b]{\linewidth}\raggedright
What remains to be supplied for the main family
\end{minipage} \\
\midrule\noalign{}
\endhead
\bottomrule\noalign{}
\endlastfoot
Actual CM point and component maps & Explicit product maps MT15--16; arbitrary-type isogeny criterion MT12 & The prescribed point, its rational frame, integral models, polarization, section and all source identities \\
Eigenletters & Exact finite label ordering and invertibility MT13 & Compatibility with the actual fixed coefficient field and prescribed frame \\
Norms & Exact source-normalized pullback implication MT14 & Its actual map/polarization witnesses at the prescribed point \\
Degrees & Unchanged Step5 with lower bound MT2 and fully costed MT5/MT9/MT11 & An effective lower bound for the actual integral eigenline degrees; no unproved equality to logOm \\
Arithmetic polynomial & Finite ideals and all supplied local tests are preserved by MT6 & The genuine all-extension \nolinkurl{Step5Loc}, nonzero polynomial and source section/jet tie \\
DP44/source assembly & An explicit extra term if the old logOm is retained & Full Host/Sizes, Read/STiesR/032 and every changed-family comparison term \\
\end{longtable}

An arbitrary finite degree bound exists after the actual finite list of lines is constructed, but that fact alone is not an effective admissibility or DP44 estimate. No unknown bad-prime tangent-lattice index is cancelled by this argument.

\subsection{6. Exact source locators and verification scope}\label{exact-source-locators-and-verification-scope}

All sources are pinned to commit \nolinkurl{7cbf7a0c32fca1badf05931879a37c3f8a2b03b4}. Newly retrieved copies and retrieval records are in the adjacent source\_cm directory. Their hashes are recorded separately.

\begin{itemize}
\tightlist
\item
  \nolinkurl{Formalism/LocalToGlobal/CMDeterminantIdeals/DetIdeal.lean}: MLine, Good, deg and ndeg, lines10--27; the norms are arbitrary positive real functions and the ideals remain algebraic.
\item
  \nolinkurl{Formalism/LocalToGlobal/CMDeterminantIdeals/StatementC.lean}: CoeffCauchy, LamEnd, pDeg and lineNN, lines10--26.
\item
  \nolinkurl{Formalism/LocalToGlobal/CMDeterminantIdeals/StatementS.lean}: exact Step5Loc, Step5Arch and Step5, lines38--73; their full bodies are used.
\item
  \nolinkurl{Formalism/LocalToGlobal/CMDeterminantIdeals/LambdaLowerBound.lean}: literal theorem \nolinkurl{lambda_lower_bound : Step5}, line11.
\item
  \nolinkurl{Formalism/LocalToGlobal/CMDeterminantIdeals/HadamardLoc0.lean}: KeyCore and the unaltered Key, Hadamard and line-degree interfaces.
\item
  \nolinkurl{Development/AbelianVarietyTheory/Proofs/CubicalLineBundles/FormLetterLe.lean}: \nolinkurl{form_letter_le}, lines19--43, with arbitrary embedding set S; full proof read.
\item
  \nolinkurl{Development/AbelianVarietyTheory/Proofs/CubicalLineBundles/EqZeroOfPow.lean}: complete independence proof \nolinkurl{lamInd} and the characteristic-zero tangent argument.
\item
  \nolinkurl{Development/AbelianVarietyTheory/Proofs/CubicalLineBundles/ExistsLamLettersIsGood.lean}: blockEquiv25--44 and the complete eigenletter construction.
\item
  \nolinkurl{Development/AbelianVarietyTheory/Proofs/CubicalLineBundles/LamDegStLogOm.lean}: lines12--16 retain the separate elliptic/Kronecker degree theorem; it is not invoked for the mixed product.
\item
  \nolinkurl{Formalism/LocalToGlobal/ProjectiveLattices/RosatiLatticeNorms/CmFrameOfEq.lean}: actual \nolinkurl{L.lines}44--59, \nolinkurl{lamw_link}78--84 and108--125; the old equality-based degree premise and room inequality are not silently removed.
\end{itemize}

The finite certificate verifies exact degree-shift algebra, every branch of the stated cost comparison on rational model grids, arbitrary-type label/row constructions and literal source fingerprints. The proof of the theorem is the argument above, not a numerical search or a formal Lean build. No global distinguished pair, period nonvanishing, selected arithmetic-image content, uniform CM model, literal main-family existence, DP44 inequality or RH theorem is asserted.

\subsection{Canonical algebraic-geometric context}\label{canonical-algebraic-geometric-context-2}

The comparison lines are invertible fractional-ideal modules on the unchanged spectrum of the coefficient integers, equipped with positive metrics at every counted embedding. The auxiliary metric change preserves their algebraic modules, determinant functor, integral Hom maps and all finite local tests. It is an archimedean arithmetic-degree comparison, not an isomorphism of the original metrized Hodge lines and not a new CM height equality. The explicit product abelian schemes, their duals, Picard data and polarization biextensions retain their actual geometry; the isogeny derivative is a characteristic-zero tangent construction. Formal groups/Dieudonne crystals at finite places and the existing lambda-ring/K-theory and motivic realizations are unchanged by the auxiliary norms. No new l-adic weight, motivic L-function or mixed source comparison is inferred. Carried contextual locators: SGA3 I4.7.3/4.7.3.1; SGA6 VI Theorem3.3; FGA232 Theorem3.1 pp148--149; FGA236 Corollaries1.7/2.3 pp229/232; SGA7 I VII2.7/2.9 pp169--174 and VIII7.4(b) pp310--311; cote4 scan003 sections8.1.1--8.1.2, scans005--006 sections8.4.1--8.4.3; cote54 scans005--007. Programme rows: DP\_MULT Q3e/DP44/DP52/DP54--65; THETA3 G1/A7; CX\_GRH CX-ALG-024; GCT\_PLETH T-052. These locators supply context, not a claimed newly checked comparison theorem.

\clearpage

\section{AM-R-084 --- original and inverse null lines exhaust rank three at support25}\label{am-r-084-original-and-inverse-null-lines-exhaust-rank-three-at-support25}

\textbf{Canonical checked integration --- 2026-11-23T01:13:10.804962+00:00.} Complete root and independent proof/script/interface reviews and separate exact replays passed. The actual canonical predecessor has95/31 entries; this integration gives96/31, complete floors25/11. Frozen author history is preserved. AM-R-084 exhausts every remaining global-rank-three scalar-cross support25 term. The complete K097 count applies to \textless=24 eligible places without a denominator premise; the exact DO093 inversion carries this count and R083 to the inverse null line with the same original S-supportk\textgreater=126, actual weights, masses, phases and cutoff. Outside the four charged sets both slopes are fully eligible and each denominator norm is\textless2\^{}4864. Their denominator-cleared difference has the full paired T\^{}50 divisor but norm\textless(4U)\^{}168*2\^{}9728\textless T\^{}50, so the two lines coincide. DO095 then forces global rank2, a contradiction. The entire four-sector absolute union is\textless2\^{}-5398, with both inverse costs included. Actual DO10495/31 becomes96/31; every old overlapping charge remains, full sums are\textless2\^{}-230, and complete floors25/11 remain unchanged. At support25 only the separately retained all-II globally rank-two-unipotent scalar-cross branch remains, with2\^{}120\textless N\_E D\_eta\textless2\^{}1820. Higher supports and signed-kernel, selected arithmetic-image, prescribed global period, actual mixed source, specific032/Read, DP44 and RH gates remain open.

Author: \nolinkurl{/root}, GPT-6 Astra Pro, ALP\_MACHINE. 23 November2026 UTC. Inherited effort is not separately exposed.

Controlled pending proof. Whole saved-proof/script readings, exact diagnostics, independent review and canonical promotion are separate recorded events. No support-floor change is claimed here.

\subsection{1. Exact original sector and theorem}\label{exact-original-sector-and-theorem}

Retain the unchanged J56 all-top normalized kernel, its original E/F of degrees336/168, H=133, every G/T prime and CM partner, actual finite projectors, selected real module and metric, packet representative, central normalization, Haar measures and class masses. Let s=25. In the actual DO099 JR24/R071 residual write

\[
 E^{169}=Z\perp P,\qquad Z=S_{164}\perp u^\perp,\qquad
 P=Eu\perp Ee_n,\qquad J_P=\operatorname{diag}(U,-1),
\] \[
 h=\begin{pmatrix}I_Z&B\\C&D\end{pmatrix},\quad
 (h^{-1})_{ZZ}=I_Z,\quad \det h=1,\quad
 C|_S\ne0,\quad \operatorname{rank}_E(h-I)=3.                 \tag{IL1}
\]

All334 unmixed G reductions are unipotent. Normalize the isotropic image line only over E:

\[
 C=\ell r,\quad \ell=(1,\eta),\quad \eta\bar\eta=U,
 \quad k=\#\{j\in S:r(e_j)\ne0\},\quad
 U=77Q^2-14,
\] \[
 Q=1427908193729121522601,\qquad
 U=156996979348344707969383396772362042643000463.             \tag{IL2}
\]

Impose k\textgreater=126. The actual remaining support25 rank-three sector satisfies this: every R082 threshold K\_t is at least125, and the retained inequality is k\textgreater K\_t. Its type-III count is11\textless=t\textless=25, as in K097 RP15/DO074 TC1.

\textbf{Theorem.} The complete actual IL1 sector with k\textgreater=126 has total absolute normalized contribution strictly less than

\[
                              2^{-5398}.                    \tag{IL3}
\]

Consequently every remaining global-rank-three scalar-cross support25 term is exhausted, after retaining the already charged R082 and all earlier cuts. The support25 remainder is then confined to the separately retained all-II globally rank-two-unipotent branch, with DO103/DO104 interval2\^{}120\textless N\_E D\_eta\textless2\^{}1820. Higher supports and the complete floors25/11 remain unchanged. No signed-kernel positivity, selected arithmetic-image witness, prescribed global period, mixed source constructor, specific032 pairing/value, DP44 sign or RH statement follows.

\subsection{2. The eligible-place count has no denominator premise}\label{the-eligible-place-count-has-no-denominator-premise}

For any norm-U scalar eta define its exact denominator ideal and eligible subset by

\[
 D_\eta=\prod_{\mathfrak P}\mathfrak P^{\max(0,-v_{\mathfrak P}\eta)},
 \qquad
 E_\eta=\{v\in\Sigma:v_{\mathfrak P_v}(\eta-a_0)\ge1\},
 \qquad a_0=Q(-7+2\sqrt{-7}).                              \tag{IL4}
\]

Here Sigma is the actual25-place support and P\_v is its fixed selected CM component. At these places U and a0 are units; a0 bar(a0)=U+14. K097 RP4--RP7 proves, without a global-unipotence hypothesis, that every original cross coordinate can have poles only at the paired places over E\_eta. In particular eligibility makes eta a unit at both partners. The opposite pole slope is U/bar(a0), with no new independent partner label.

We require the following slightly more general form of the count inside that proof:

\[
 \sum_{h\text{ satisfying IL1},\ |E_\eta|\le24}
                |\mathcal W(h)|<2^{-11658}.                 \tag{IL5}
\]

This form does not impose any bound on N\_E D\_eta and does not require k\textgreater=126. To verify it, begin with the paragraph following K097 RP9. Its only input there is \textbar E\_eta\textbar\textless=24. Pad that subset to a24-subset of Sigma, at most25 choices. The same common whole-C G ideal J gives

\[
 L=J I_D,\quad N_E I_D=T^{-48},\quad
 z=(N_EL)^{1/336}\ge T^{-24/168}.                            \tag{IL6}
\]

All167 original covector values and the original negative pivot beta belong to L. K097 RP11--RP15 retain the entire covector ball and actual relative beta-norm fibre, the actual selected trace, the full projected-column count191\^{}(4f)U\^{}fT\^{}100, source fibre2\^{}36, all G patterns/supports/pivots, and the complete general type weights. No eta-denominator enumeration or lower/upper bound enters those factors. The low-range z exponent is -f/6 and IL6 gives the necessary T\^{}4 bound; the high-range expression is decreasing. Their full maxima in RP20 are -47790211/4096 and -223030653/4096. Thus the full25-padding and15-type union is exactly the already certified RP21 inequality

\[
          375(2^{-11667}+2^{-54450})<2^{-11658}.              \tag{IL7}
\]

This proves IL5 in its enlarged scope. The old small-denominator hypothesis occurred only in K097 RP8--RP9, where it was one sufficient way to prove \textbar E\_eta\textbar\textless=24. It is not reused here. All original zero-coordinate patterns and label-zero cases remain in the unchanged full count. No assertion about the full rank-three double-pole part is imported from the rank-two normal form.

The second complete containing-set bound is the checked R083 LD1--LD16 theorem:

\[
 \sum_{h\text{ satisfying IL1},\ k\ge126,\ N_ED_\eta\ge2^{4864}}
                  |\mathcal W(h)|<2^{-5400}.                \tag{IL8}
\]

It includes every denominator ideal and every actual eta, the full selected trace, all type weights, projected U/T column and source fibre. Its exponent is -22330263/4096\textless-5451. We use this full bound, not a density-only estimate or a bounded ideal-class count.

\subsection{3. Literal inversion and the second normalized null line}\label{literal-inversion-and-the-second-normalized-null-line}

The exact inversion calculation is DO093 IV4--IV16. Its equalities do not depend on s. The actual finite coefficient and selected real coefficient are normalized traces of unitary representations against orthogonal projections, so

\[
 a(h^{-1})=\overline{a(h)},\qquad
 c(h^{-1})=\overline{c(h)}.                                \tag{IL9}
\]

For each actual paired T pencil, reversing the ordered source/target arrow transposes the incidence matrix. Its real symmetric projector gives rho(h\^{}-1)=rho(h), including the admitted zero label. Therefore

\[
 \mathcal W(h^{-1})=\overline{\mathcal W(h)},\qquad
 |\mathcal W(h^{-1})|=|\mathcal W(h)|.                       \tag{IL10}
\]

The class mass, central factor28, finite probability Haar, real Haar, actual packet representative and coefficient-bundle norm are the same. Inversion preserves the actual cutoff, support, each of the three incidence types, all unmixed finite unipotence conditions and the determinant-one representative. The source and target lattices are reversed together at both CM partners. These are exactly the support-independent IV9--IV15 identities; no new factor two, phase, partner label or changed lattice is introduced.

We also need the precise original-coordinate support. DO095 RA3--RA4 gives detD=1 and BD\^{}-1 C=0. The block inverse formula, with identity Schur complement, is

\[
                  C^-:=(h^{-1})_{PZ}=-D^{-1}C.             \tag{IL11}
\]

The nonzero vector D\^{}-1 ell spans the inverse isotropic cross line. Its first coordinate cannot vanish, since J\_P=diag(U,-1) has no nonzero isotropic vector of the form(0,b). Write

\[
 -D^{-1}\ell=\gamma(1,\xi),\quad\gamma\in E^\times,
 \quad\xi\bar\xi=U,
 \qquad C^-=(1,\xi)(\gamma r).                            \tag{IL12}
\]

Thus the inverse row has exactly the same nonzero original S coordinates: k(h\^{}-1)=k(h). Its positive compression is I, its global rank is3, and t(h\^{}-1)=t(h). Every hypothesis of IL5 and IL8 applies to the inverse matrix, with its actual reversed source/target pair. In particular k\textgreater=126 is retained without an assertion that the chosen common-G orientation is inversion invariant.

The previously cut remainder itself need not be invariant under inversion. We only map each subset injectively into the \textbf{complete} original containing sector in IL5 or IL8. DO093 IV16 justifies this use of absolute sums and unchanged source/target-uniform reconstruction bounds. Consequently

\[
 \sum_{h\text{ satisfying IL1},\ k\ge126,\ |E_\xi|\le24}
                      |\mathcal W(h)|<2^{-11658},
\] \[
 \sum_{h\text{ satisfying IL1},\ k\ge126,\ N_ED_\xi\ge2^{4864}}
                      |\mathcal W(h)|<2^{-5400}.             \tag{IL13}
\]

\subsection{4. Uniqueness of a fully eligible bounded-denominator line}\label{uniqueness-of-a-fully-eligible-bounded-denominator-line}

Suppose a matrix lies outside all four containing sectors IL5, IL8 and IL13. Then

\[
 E_\eta=E_\xi=\Sigma,\qquad
 N_ED_\eta<2^{4864},\qquad N_ED_\xi<2^{4864}.               \tag{IL14}
\]

Put delta=eta-xi. If delta is nonzero, the ideal

\[
                      \mathfrak A=(\delta)D_\eta D_\xi     \tag{IL15}
\]

is a nonzero integral ideal. Indeed at every prime, v(delta)\textgreater=min(v(eta),v(xi)), and adding both nonnegative denominator exponents makes that lower bound nonnegative. Neither denominator is required to be principal or conjugation invariant.

At every selected prime over Sigma, full eligibility gives eta congruent to xi congruent to a0. Both scalars and both exact denominator ideals are units there and at the conjugate partner. Conjugating the selected equality and using each relative norm U gives the opposite equality as well: the paired residue for each scalar is U/bar(a0). Hence delta is divisible at both partners. All50 degree-one CM primes divide the integral ideal in IL15, so

\[
                 T^{50}\le N_E\mathfrak A
                     =|N_{E/\mathbf Q}\delta|
                         (N_ED_\eta)(N_ED_\xi).             \tag{IL16}
\]

At every complex label, \textbar eta\_tau\textbar=\textbar xi\_tau\textbar=sqrt(U). Thus \textbar delta\_tau\textbar\textless=2sqrt(U), and the full absolute E norm satisfies

\[
 |N_{E/\mathbf Q}\delta|\le(4U)^{168},\qquad
 N_E\mathfrak A<(4U)^{168}2^{9728}.                         \tag{IL17}
\]

The unchanged actual host integers obey the exact strict inequality

\[
                     T^{50}>(4U)^{168}2^{9728}.             \tag{IL18}
\]

The companion checks IL18 by integer powers using the original J56 host certificate. No floating logarithm, height estimate on a chosen generator, class-number bound or principal-denominator assertion is used. IL16--IL18 contradict delta nonzero. Therefore

\[
                                \eta=\xi.                  \tag{IL19}
\]

This uniqueness is a statement about the full two-partner residue conditions and the product of the two exact denominator norms. Replacing the paired divisor by just one component would not be the stated argument.

\subsection{5. Equal image lines contradict global rank three}\label{equal-image-lines-contradict-global-rank-three}

Equation IL19 identifies the original and inverse cross image lines. By IL11, D\^{}-1 preserves E ell; hence D ell=lambda ell for a nonzero scalar lambda. In DO095 RA5 notation b=-ell*J\_PD, so

\[
 b\ell=-\lambda\ell^*J_P\ell=0,\qquad
                         BC=0.                             \tag{IL20}
\]

DO095 RA9--RA10 then gives rank\_E(h-I)=2 in this exact residual, contradicting IL1. This uses the determinant-one, identity-positive-compression and nonzero-cross premises already displayed. It does not infer global unipotence from line equality, nor does it invoke a moving integral Witt basis. The alternative rank-two matrix may be nonunipotent; either possibility already contradicts rank3.

Thus every IL1 matrix with k\textgreater=126 belongs to at least one of the four counted sectors. Keeping the overlaps gives the complete absolute bound

\[
 \sum |\mathcal W(h)|
 <2\bigl(2^{-11658}+2^{-5400}\bigr)<2^{-5398},               \tag{IL21}
\]

which proves IL3. The two inverse-image contributions are explicitly charged; the original R083/K097 errors are not silently reused without their inverse cost. For canonical integration append5398 once to the actual ordered all-top predecessor, retaining every old overlapping charge and all31 uniform errors. The exact full sums remain below2\^{}-230. The complete support floor stays25, because the all-II rank-two branch has not yet been exhausted.

\subsection{6. Exact source, run and geometric scope}\label{exact-source-run-and-geometric-scope}

The new content is the denominator-independent extraction IL5 from the already proved K097 RP10--RP21 count, its literal inverse application, the two-line height uniqueness IL15--IL19, and their complete rank-three exhaustion. Inputs freshly read in full or at the exact displayed clauses are K097 RP1--RP21; R083 LD1--LD16; DO093 IV4--IV16; DO095 RA1--RA10 and RA24--RA26; R082 TP1--TP13; DO099 JR24/R071 GR11--GR12; and R061 GS11--GS12 for the actual selected coefficient. The underlying source identity is K062 AT13--AT16/AT24, as carried explicitly in DO093. The known finite/projector, local model and actual source normalization results are cited, not reproved or replaced.

The integral ideal IL15 is an invertible module on the unchanged Spec(O\_E). Its norm is the determinant of multiplication combined with the original ideal norm, and the nonzero section is tested at every complex realization and all selected paired finite points. Inversion is the actual transpose of the same Hermitian correspondence and its symmetric pencil coefficients. The null-line equality concerns the rational Hom module and the resulting rank stratum; it supplies no globally integral null basis. The CM abelian variety, theta line, Picard scheme, polarization and biextension remain unchanged. Formal groups/Dieudonne crystals and Cartier duality remain those already constructed; no new lambda-ring/K-theory equality, l-adic weight theorem, motivic comparison or motivic L-function result is inferred.

Carried exact contextual locators are SGA3 I4.7.3/4.7.3.1; SGA6 VI Theorem3.3; FGA232 Theorem3.1 pp148--149; FGA236 Corollaries1.7/2.3 pp229/232; SGA7 I VII2.7/2.9 pp169--174 and VIII7.4(b) pp310--311; cote4 scan003 sections8.1.1--8.1.2, scans005--006 sections8.4.1--8.4.3; cote54 scans005--007. DP\_MULT Q3e and THETA3 G1/A7 gain the exhaustive rank-three support25 reduction. CX\_GRH CX-ALG-024 and GCT\_PLETH T-052 retain their separate arithmetic gates. No accepted Alpöge theorem or archive is audited, and no Lean build is performed.

\clearpage

\section{AM-LIFT-081 --- effective degrees of the actual integral mixed-CM tangent lines}\label{am-lift-081-effective-degrees-of-the-actual-integral-mixed-cm-tangent-lines}

\textbf{Canonical checked integration --- 2026-11-23T01:18:14.298792+00:00.} Complete root and independent whole-proof/script/source-interface reviews, with separate10,946-assertion exact replays, passed. The frozen historical source remains unchanged. LIFT081 ED1--23 supplies effective intrinsic line-degree lower bounds for every actual comparison component eigenline on J14/J56, with the full integral CM action and the literal source metric n\^{}2=pi H. The good model is obtained over K0(alpha,beta), alpha\^{}7=-4 and beta\^{}2=(zeta7-1)\^{}7, of degree at most14 over K0. The actual integral vector 2q\_h p\_tau e\_i has norm below sqrt(2)q\_h at every counted embedding, with q1=406 and q4=26390. Therefore ndeg\textgreater-log(sqrt(2)q\_h), and the weaker common bounds are -10log(2) for J14 and -16log(2) for J56. These normalized degrees are invariant under further finite base change of the actual good model. Substitution into LIFT080 makes its sufficient conditional Step5 cost fully explicit: Copt=1/2(1+log L)(D+V)+D log(rL)+b\_h min(D+V,(3D+V)/2), V=gn and b\_h=10log(2) or16log(2). This resolves the effective comparison-line-degree premise only. The prescribed mixed point, model/level/coefficient-field ties, nonzero frame polynomial, all residue-extension Step5Loc conditions, source-room data, selected arithmetic image/content, global period and DP44 evaluation remain required. No claim is made about an arbitrary line in a larger multiplicity space, equal degrees, a primitive coordinate-ideal equality, an old LamPtClauses witness, a finite charge, or RH.

Author: \nolinkurl{/root/lift_accounting},23 November2026 UTC. Controlled pending proof. The final whole saved-proof/source/script review and independent reviews are separate records. No canonical file is modified.

\subsection{1. The actual objects and the effective bound}\label{the-actual-objects-and-the-effective-bound}

Let J be the Jacobian of C:y²=x⁷−1, with its actual principal polarization and the rational invariant differentials

\[
 \omega_i=x^{i-1}dx/y,\qquad 1\le i\le3.
 \tag{ED1}
\]

Write e\_i for their dual tangent vectors. Use the two existing K030/K032 product hosts

\[
 A_h=J^{14h},\quad E_1=\mathbf Q(\zeta_{28})B_7(29),\quad
 E_4=E_1(\sqrt5,\sqrt{13}),\qquad h=1,4.
 \tag{ED2}
\]

Here {[}E\_h:Q{]}=84h and dim A\_h=42h. The O\_Eh action and product principal polarization are the actual ones proved in those findings, with Rosati equal to conjugation. The notation A\_h\^{}- means the same polarized variety and integral model, with its O\_Eh action precomposed by conjugation. It does not change the complex variety or its metric. Finite repeated copies of these comparison factors are also allowed, using the eigenlines of each component separately. No bound for an arbitrary line in a higher-multiplicity isotypic sum is asserted.

Take a number field K0 containing E\_h and the coefficients already needed for the chosen generic objects. Put

\[
 \iota=\zeta_{28}^{7},\quad \lambda=\zeta_7-1,\quad
 \alpha^7=-4,\quad\beta^2=\lambda^7,\qquad K=K_0(\alpha,\beta).
 \tag{ED3}
\]

Thus {[}K:K0{]}\textless=14. The good Jacobian models of ALP-LIFT-05 give the actual product abelian scheme over O\_K. There is no requirement that this auxiliary field be CM or Galois. The host E\_h/F\_h is unchanged. The source \nolinkurl{LamObj.K₁} itself requires a number field and the indicated coefficient-field algebra structure, not a CM extension. This fact permits this auxiliary model construction; it does not identify K with an independently fixed global coefficient field or instantiate an unrelated source record.

For a holomorphic E\_h eigenline choose any nonzero algebraic tangent vector v\_tau over K and let

\[
 I_\tau=\{c\in K:c v_\tau\in\operatorname{Lie}(\mathcal A_h/\mathcal O_K)\},
 \qquad n_\tau(\sigma)=\|v_\tau\|_{\mathrm{source},\sigma}.
 \tag{ED4}
\]

This is a nonzero fractional ideal, equivalently the intrinsic saturated line intersection. The norm is exactly the one in \nolinkurl{IsVecLine}, including its factor pi. Any actual source \nolinkurl{IsVecLine} for the same vector and model has this ideal and norm: its local inverse-matrix criterion is precisely ED4, and its archimedean clause is the norm used below. For a local model basis whose rows are M in the generic Key basis, the coordinates of c v\_i are exactly c e\_i M\^{}(-1). The source requires all these coordinates to be integral exactly when c belongs to the local principal ideal generated by cl; this is the localization of ED4. No equality between this ideal and K035's embedding-coordinate ideal cD\_tau is asserted.

Define the rational clearing integers

\[
                       q_1=406,
             \qquad q_4=26390=65q_1.
 \tag{ED5}
\]

\textbf{Theorem.} Every line ED4, on A\_h or A\_h\^{}-, satisfies

\[
 \boxed{\quad
 \operatorname{ndeg}(I_\tau,n_\tau)
           >-\log(\sqrt2\,q_h).
 \quad}                                                     \tag{ED6}
\]

In particular effective common lower bounds are

\[
 d_-=-b_h,\qquad b_1=10\log2,\quad b_4=16\log2.
 \tag{ED7}
\]

All logarithms are natural. These bounds hold in every coefficient embedding and after every further finite base extension of the actual good model. They are independent of the chosen normalization of v\_tau and the number of repeated comparison copies. They do not assert that the degrees are equal, calculate the primitive eigenline ideals, or supply the prescribed mixed-point/section, polynomial integrality, Host/Sizes, Read or selected-image gates.

\subsection{2. The actual integral tangent vectors and source norm}\label{the-actual-integral-tangent-vectors-and-source-norm}

We reuse the proved good models, with their full differential bases and saturation arguments, from ALP-LIFT-05 H3--H6 and H19. At primes above2 the model and basis are

\[
 Y^2+Y=X^7,\quad x=\alpha X,\quad y=\iota(2Y+1),
 \quad\eta_i^{(2)}=X^{i-1}dX/(2Y+1),
 \quad\omega_i=\alpha^i\iota^{-1}\eta_i^{(2)}.
 \tag{ED8}
\]

The inverse tangent coordinates of e\_i have valuation -2i/7, using v\_2(2)=1. Hence2e\_i is integral for all three i, with valuations1-2i/7\textgreater0. The differentials here are dx/y, not the half-differentials; dropping that distinction would give a wrong dyadic conclusion.

At primes above7 the proved model is

\[
 Y^2=((1+\lambda X)^7-1)/\lambda^7,\quad x=1+\lambda X,
 \quad y=\beta Y,
\] \[
 \eta_j^{(7)}=X^{j-1}dX/Y
   =\beta\lambda^{-j}\sum_{a=0}^{j-1}
        {j-1\choose a}(-1)^{j-1-a}\omega_{a+1}.
 \tag{ED9}
\]

Evaluating these basis differentials on e\_i gives coefficients that are integral multiples of beta lambda\^{}(-j), whose normalized7-valuations are7/12-j/6\textgreater=1/12 for j\textless=3. Thus e\_i itself is integral there. At primes away2,7 the original omega\_i are an integral differential basis. Therefore

\[
           2e_i\in\operatorname{Lie}(\mathcal J/\mathcal O_K)
                   \quad(1\le i\le3).                      \tag{ED10}
\]

The earlier proof constructs smooth proper models, not only equations with smooth affine charts: it includes the infinity charts, geometric integrality, the saturated bases and the relative Jacobian. The2-model is smooth over O\_K{[}1/7{]}, the7-model over O\_K{[}1/2{]}, and their displayed generic coordinate changes glue them on the overlap. ED10 is consequently a statement about the actual global tangent module. Repeated products retain this direct-sum integral module.

The full integral O\_Eh action extends on this model. In the2-chart the curve automorphism x↦zeta7 x is X↦zeta7 X,Y↦Y. In the7-chart it is X↦1+zeta7 X,Y↦Y. Both maps and their inverses are integral and preserve the infinity section. The remaining actual action uses the integer matrices on the paired Jacobian factors, the multiplication matrices in K030's integral basis of the inverse-different module, and K032's integral quadratic action matrices. Such matrices act by integral group homomorphisms on the product abelian scheme. This uses the existing full-order product action; it does not invent a larger integral CM order.

Let G\_i be the actual Hodge Gram coefficients of ALP-LIFT-03/05 and K035 PF8. In the rational omega coordinates the tangent polarization form has

\[
                         H(e_i,e_i)=G_i^{-1}.
 \tag{ED11}
\]

One can read this directly from PF9--PF11: the paired coordinates w\_±=u±iv have coefficient1/(2G\_i), so their sum of squared lengths gives coefficient1/G\_i on each original factor. The later B7 and quadratic-coordinate changes merely express the same product principal metric in other coordinates.

The literal source \nolinkurl{LamLetters.IsVecLine}, lines96--106, requires

\[
                         n_\sigma(v)^2=\pi H_\sigma(v,v).
 \tag{ED12}
\]

The pi in ED12 is not inferred from PF9. K066's exact rational beta-function certificate proves G\_i/(2pi)\textgreater1 for every i. Consequently every original product basis vector satisfies

\[
                          \|e_i\|_{\mathrm{source}}<1/\sqrt2.
 \tag{ED13}
\]

The curve, product polarization and differential basis are defined over Q. In each coefficient embedding they give the same complex metric in these rational coordinates. The E-action/eigenlabels are transported; no field embedding is applied to a free transcendental period. This is the same realization convention retained in K066 PC3/PC9.

\subsection{3. Integral trace-dual projectors on the actual CM order}\label{integral-trace-dual-projectors-on-the-actual-cm-order}

We first prove the general algebraic lemma. Let E be a number field, K contain all its embeddings, and let b\_1,\ldots,b\_N be an integral basis of O\_E. Write b\_i\^{}* for its trace-dual basis. If a positive integer q satisfies

\[
                       q\mathfrak D_{E/\mathbf Q}^{-1}
                                     \subset\mathcal O_E,
 \tag{ED14}
\]

then q b\_i\^{}* is integral for every i. For an embedding tau:E→K the separability projector is

\[
                  p_\tau=\sum_i \tau(b_i^*)\,\mathrm{act}(b_i),
 \tag{ED15}
\]

where act is the actual E-action on the generic tangent space. It is the identity on the tau eigenspace and zero on every other eigenspace. Indeed the two embedding matrices are inverse under the trace pairing: the trace-dual identities imply

\[
             \sum_i\rho(b_i)\tau(b_i^*)=\delta_{\rho,\tau}.
 \tag{ED16}
\]

This also proves p\_tau²=p\_tau and sum p\_tau=I, without a choice of primitive element. Since act(b\_i) preserves the integral tangent module and q tau(b\_i\^{}*) lies in O\_K, q p\_tau preserves it. No assertion that p\_tau itself is integral is needed.

For completeness, ED14 with the actual q\_h is read directly from the full tensor integer bases of LIFT063. Their bilinear trace matrices have denominator-clearing constants as follows:

\begin{longtable}[]{@{}
  >{\raggedright\arraybackslash}p{(\columnwidth - 4\tabcolsep) * \real{0.3000}}
  >{\raggedright\arraybackslash}p{(\columnwidth - 4\tabcolsep) * \real{0.3000}}
  >{\raggedleft\arraybackslash}p{(\columnwidth - 4\tabcolsep) * \real{0.4000}}@{}}
\toprule\noalign{}
\begin{minipage}[b]{\linewidth}\raggedright
Integral constituent basis
\end{minipage} & \begin{minipage}[b]{\linewidth}\raggedright
Trace matrix or its description
\end{minipage} & \begin{minipage}[b]{\linewidth}\raggedleft
A clearing integer for its inverse
\end{minipage} \\
\midrule\noalign{}
\endhead
\bottomrule\noalign{}
\endlastfoot
1,zeta28,\ldots,zeta28\^{}11 & T\_ab=c\_28(a+b), the Ramanujan trace & 14 \\
The seven order-four Gaussian periods in B7(29) & 29I\_7−4*11\^{}t & 29 \\
1,(1+sqrt5)/2 & {[}{[}2,1{]},{[}1,3{]}{]} & 5 \\
1,(1+sqrt13)/2 & {[}{[}2,1{]},{[}1,7{]}{]} & 13 \\
\end{longtable}

These are bilinear trace matrices; the factor2 in the cyclotomic full trace has not been confused with its half-trace Hermitian matrix. The first inverse is checked exactly in the companion. The second inverse is(I+4\emph{11\^{}t)/29. The other two inverses follow from their determinants5 and13. Since the constituent discriminants are coprime, their tensor integral bases are exactly O\_Eh by LIFT063. Tensoring the inverse matrices gives the integer q\_1=14}29 or q\_4=14\emph{29}5*13 in ED14. It is irrelevant whether a single Gaussian period generates the full integer ring.

As a check, the actual field discriminants are2\^{}84 7\^{}70 29\^{}72 and2\^{}336 7\^{}280 29\^{}288 5\^{}168 13\^{}168. The normalized local exponents of the different are respectively1,5/6,6/7 and additionally1/2,1/2. These give the same clearing integers; the trace-matrix argument already proves the required inclusion.

At every complex realization Rosati equals conjugation on the actual E\_h action. Its commuting operators are normal, their distinct embedding eigenspaces are orthogonal, and p\_tau is the corresponding orthogonal projection. Thus

\[
                 \|p_\tau v\|_{\mathrm{source},\sigma}
                       \le\|v\|_{\mathrm{source},\sigma}
                       \quad\hbox{for every }\sigma.
 \tag{ED17}
\]

Under sigma the projector is the one for the transported algebraic embedding sigma∘tau. This is why the same global projector construction gives bounds in every complex realization. Precomposing the E-action by conjugation only permutes these projectors.

\subsection{4. A bounded integral vector gives the intrinsic line-degree bound}\label{a-bounded-integral-vector-gives-the-intrinsic-line-degree-bound}

The actual E\_h action on A\_h has multiplicity one for every holomorphic embedding. For a chosen such tau, p\_tau is nonzero. The original product vectors e\_i form a K-basis, so one of them, chosen once over K, has p\_tau e\_i≠0. Put

\[
                             w_\tau=2q_h p_\tau e_i.
 \tag{ED18}
\]

By ED10 and the integral-projector assertion, w\_tau is an actual nonzero integral tangent vector on that eigenline. By ED13 and ED17,

\[
                     0<\|w_\tau\|_{\mathrm{source},\sigma}
                                      <\sqrt2 q_h
                                \quad\hbox{for every }\sigma.
 \tag{ED19}
\]

The index i is not chosen separately at the complex places. A nonzero algebraic vector stays nonzero under every field embedding. Thus ED19 bounds one global vector and its actual conjugate realizations, not unrelated local choices. Moreover the action matrices, trace-dual coefficients and original e\_i are already defined over K0, so w\_tau itself is defined over K0. The radicals alpha and beta are used to exhibit the good integral model, not as new coefficients of this bounded tangent vector. This does not assert that the original object already has good reduction over K0.

Write w\_tau=c\_tau v\_tau with c\_tau∈K×. Its integrality gives c\_tau O\_K⊂I\_tau. Therefore N(I\_tau)\textless=\textbar N\_K/Q(c\_tau)\textbar, including fractional I\_tau. The exact source definition \nolinkurl{DetIdeal.ndeg} gives

\[
\begin{split}
 \operatorname{ndeg}(I_\tau,n_\tau)
  &=-\frac{\log N(I_\tau)}{[K:\mathbf Q]}
       -\frac1{[K:\mathbf Q]}\sum_\sigma\log n_\tau(\sigma)\\
  &\ge-\frac1{[K:\mathbf Q]}\sum_\sigma
                    \log\bigl(|\sigma(c_\tau)|n_\tau(\sigma)\bigr)\\
  &=-\frac1{[K:\mathbf Q]}\sum_\sigma
                       \log\|w_\tau\|_{\mathrm{source},\sigma}
   >-\log(\sqrt2 q_h).
\end{split}                                                   \tag{ED20}
\]

The sign follows from the inclusion c\_tau O\_K⊂I\_tau, not its reverse. The equality in the second line is the ordinary algebraic field-norm/product identity, and every archimedean value of the same c\_tau is included. This proves ED6. In particular no unknown bad-prime index is discarded: using only one integral vector gives a lower bound for the actual saturated line degree, irrespective of whether that vector is primitive.

Since2q\_1²\textless2\^{}20 and2q\_4²\textless2\^{}32, ED7 follows. Further finite base change tensors the smooth tangent bundle and its saturated rank-one intersection. Locally this is extension of the same principal fractional ideal; archimedean realizations repeat with their usual multiplicities. Both finite norm degree and archimedean sum scale by the extension degree, so normalized degrees and the displayed bounds are unchanged. The models needed for this assertion are the actual good models, not arbitrary unrelated lattices.

\subsection{5. Exact application to the unequal-degree Step5 theorem}\label{exact-application-to-the-unequal-degree-step5-theorem}

Use LIFT080 with any finite list of these actual comparison lines. Keep all its other hypotheses, including the prescribed Key data, target determinant and theta lines, actual row matrix, nonzero polynomial G, every all-residue-extension \nolinkurl{Step5Loc} condition, and the actual \nolinkurl{Step5Arch} bound. Put D=sum a\_bi, V=gn and retain its L\textgreater=1 and integer r\textgreater=1. ED7 supplies the previously unevaluated line-degree premise d\_j\textgreater=d\_- with d\_-=-b\_h.

For the auxiliary common degree e, the exact sufficient cost is now

\[
 C_h(e)=\bigl(\tfrac12+|e|+\tfrac12\log L\bigr)(D+V)
       +D\log(rL)+\tfrac{3D+V}{2}(e+b_h)_+.
 \tag{ED21}
\]

The optimal cost within that proved equalization method is

\[
 \boxed{\quad C_{h,\mathrm{opt}}
  =\tfrac12(1+\log L)(D+V)+D\log(rL)
       +b_h\min\{D+V,(3D+V)/2\}.\quad}                     \tag{ED22}
\]

The choices e=-b\_h when D\textgreater=V and e=0 when D\textless=V realize the two branches, including equality and zero-degree cases. If the original scalar convention logOm=ell\_0 is retained instead, the fully explicit extra cost is

\[
                  \Delta C_h=\tfrac{3D+V}{2}
                                      (\ell_0/2+b_h)_+.
 \tag{ED23}
\]

Using the sharper b\_h=log(sqrt2 q\_h) in ED21--ED23 is also valid. These are upper costs in the actual inequality, not lower bounds on an attained arithmetic value. No physical metric, finite line ideal, section, period, measure or level is rescaled; LIFT080's auxiliary lines merely apply the unchanged generic Step5 theorem. In particular the new b\_h is not silently identified with the old elliptic/Kronecker logOm.

\subsection{6. Exact remaining source and field boundary}\label{exact-remaining-source-and-field-boundary}

This finding constructs bounded integral vectors in the actual comparison eigenlines and evaluates a sufficient common degree lower bound. It does not construct the full requested source family.

\begin{longtable}[]{@{}
  >{\raggedright\arraybackslash}p{(\columnwidth - 4\tabcolsep) * \real{0.3333}}
  >{\raggedright\arraybackslash}p{(\columnwidth - 4\tabcolsep) * \real{0.3333}}
  >{\raggedright\arraybackslash}p{(\columnwidth - 4\tabcolsep) * \real{0.3333}}@{}}
\toprule\noalign{}
\begin{minipage}[b]{\linewidth}\raggedright
Interface
\end{minipage} & \begin{minipage}[b]{\linewidth}\raggedright
Supplied here
\end{minipage} & \begin{minipage}[b]{\linewidth}\raggedright
Still required
\end{minipage} \\
\midrule\noalign{}
\endhead
\bottomrule\noalign{}
\endlastfoot
Comparison good model & Existing ALP-LIFT-05 model and its products over K0(alpha,beta), degree\textless=14 & The actual prescribed target point and all required morphisms to it \\
Integral CM action & Existing full-order product action extended on those good models & Any different host/order/model must be checked separately \\
Metrized eigenlines & Their intrinsic saturated intersection and exact source piH metric; ED6 & No primitive cD identification or equal-degree assertion \\
Generic height step & ED21--ED23 with every LIFT080 premise retained & Nonzero G, complete polynomial integrality, actual Key/frame/section conditions \\
Fields & One explicit finite auxiliary number field extension; all its embeddings included & Any preassigned common coefficient field, normal-CM requirement elsewhere, descent or uniformity tie \\
Main comparison & A concrete evaluated line-degree cost component & Changed-family Host/Sizes/DP44 data, Read/specific032, selected arithmetic nonzero/content and global period \\
\end{longtable}

The auxiliary field in ED3 may be non-CM. This is permitted for the displayed \nolinkurl{LamObj.K₁} and generic Step5 number-field variable. It is not evidence that a separately constrained CM coefficient field may be enlarged in that way. Nor is the existence of these free product objects an identification with an arbitrary fixed uniformization point. LIFT079's old quadratic pure-type/fixed-point obstruction remains exactly in force, and no original \nolinkurl{LamPtGood} witness is inferred.

\subsection{7. Literal sources and verification scope}\label{literal-sources-and-verification-scope}

\begin{itemize}
\tightlist
\item
  \nolinkurl{agents/hodge_lattice_and_cm_basis.md}, ALP-LIFT-05 H1--H6, H16--H19 and the complete proof of(a): actual good models, saturated differential lattices and gluing.
\item
  \nolinkurl{agents/degree_seven_principal_host.md}, K030 PH2--PH15 and PH20--PH25; \nolinkurl{agents/degree_four_principal_host_repair.md}, K032 QH9--QH15: full integral product actions and principal polarizations.
\item
  \nolinkurl{agents/mixed_cm_period_frame_adapter.md}, K035 PF8--PF11: exact original product metric and actual coordinate transformations.
\item
  \nolinkurl{agents/actual_cm_period_cost.md} and its exact certificate, K066 PC3/PC9 and the rational inequality G\_i/(2pi)\textgreater1: all-realization convention and Hodge norm bound.
\item
  LIFT063 \nolinkurl{explicit_original_ring_cover_pending.md}, sections1--3: exact tensor integer bases and trace matrices. The bilinear trace clearing is checked explicitly here.
\item
  Recovered pinned \nolinkurl{Development/AbelianVarietyTheory/Definitions/CMLattices/LamLetters.lean}: \nolinkurl{LamObj}43--67, \nolinkurl{IsModelAt}87--94, \nolinkurl{IsVecLine}96--106. The finite inverse-matrix criterion and the factor pi are both retained.
\item
  Pinned \nolinkurl{Formalism/LocalToGlobal/CMDeterminantIdeals/DetIdeal.lean}, lines10--27: exact MLine, Good, degree and normalized degree.
\item
  Pinned \nolinkurl{ReciprocityTheory/ReductiveGroups/ArithmeticGeometry/BirationalGeometry/ToroidalCompactification/IsLieBasis.lean}, lines20--42: tangent module, complete basis and genuine KeyIn data.
\item
  LIFT080 MT1--MT17, with its unchanged literal \nolinkurl{StatementS.Step5} and \nolinkurl{LambdaLowerBound} interfaces: the cost application, not a new source axiom.
\end{itemize}

The companion verifies exact local differential-coordinate valuations, trace matrices and their denominator clearing, separability-projector identities in exact finite algebra models, the actual field constants, product-formula signs, every cost branch and source fingerprints. Such diagnostics support the mathematical proof; they do not mechanize an abelian scheme or create an admissible automorphic/source record. No Lean build, archive audit, finite remainder charge, support-floor increase, period nonvanishing or RH claim is made.

\subsection{Canonical algebraic-geometric context}\label{canonical-algebraic-geometric-context-3}

The saturated eigenline is the intersection of a characteristic-zero CM eigenspace with the actual tangent bundle of the good abelian scheme over Spec(O\_K). The trace-dual separability projector is made integral by the displayed inverse-different clearing integer; no integral splitting at ramified primes is asserted. Its Rosati-compatible complex realization is an orthogonal projection. One global integral vector controls the same metrized invertible module at all embeddings; the fractional-ideal norm and product formula give its lower arithmetic degree. This preserves the original Jacobian, dual/Picard data, principal polarization, Poincare biextension and theta line. Auxiliary metric equalization in LIFT080 changes an archimedean comparison cost only. Formal groups and their Dieudonne crystals, existing lambda-ring/K-theory operations and the motivic realizations remain the same. No new l-adic weight, motivic L-function identity or global arithmetic source comparison is inferred. Carried contextual locators: SGA3 I4.7.3/4.7.3.1; SGA6 VI Theorem3.3; FGA232 Theorem3.1 pp148--149; FGA236 Corollaries1.7/2.3 pp229/232; SGA7 I VII2.7/2.9 pp169--174 and VIII7.4(b) pp310--311; cote4 scan003 sections8.1.1--8.1.2, scans005--006 sections8.4.1--8.4.3; cote54 scans005--007. Programme rows: DP\_MULT Q3e/DP44/DP52/DP54--65; THETA3 G1/A7; CX\_GRH CX-ALG-024; GCT\_PLETH T-052. These are contextual locators, not newly proved comparison theorems.

\clearpage

\section{AM-DO-105 --- one eligible line, one local label, and the sparse middle interval}\label{am-do-105-one-eligible-line-one-local-label-and-the-sparse-middle-interval}

\textbf{Canonical checked integration --- 2026-11-23T01:19:44.326109+00:00.} Complete root and independent proof/script/interface reviews and separate exact replays passed. The actual canonical predecessor has96/31 entries; this integration gives97/31, complete floors25/11. Frozen author history is preserved. AM-DO-105 removes the entire actual all-II globally rank-two-unipotent scalar-cross support25 subset with at most24 eligible places, without any denominator premise, together with the fully eligible middle window2\^{}120\textless N\_E D\_eta\textless2\^{}1820 having1\textless=k\textless=21 nonzero original S coordinates. For each fixed support there is at most one fully eligible eta in this window and at most one actual diagonal pencil label per active place, including central-only poles. The original weighted label sum is at most(T-1)\^{}-25. The joint S/transverse beta integral, moving central cover and simultaneous x/z endpoint retain all full metric/phase/class/Haar factors. The worst claimed cell -1852891/4096\textless-452 and full fresh union give one451 charge. Actual R08496/31 becomes97/31, both complete sums remain below2\^{}-230 and complete floors25/11 remain. After R084, DO103/104 and this theorem the only support25 branch left is fully eligible all-II rank-two unipotent, with2\^{}120\textless N\_E D\_eta\textless2\^{}1820 and22\textless=k\textless=164, unique eta per support and at most one label per active place. The dense range, all higher-support cuts, signed inequality, actual mixed source, selected arithmetic image/content, prescribed global period, specific032/Read, DP44 and RH gates remain.

Author: \nolinkurl{/root/definite_oscillator}, ALP\_MACHINE. 23 November2026 UTC.

Controlled pending proof. Author saved-file verification, independent review/replay and canonical integration are separate events. Existing frozen findings and canonical records are unchanged.

\subsection{1. Actual sector and precise conclusion}\label{actual-sector-and-precise-conclusion}

Keep the actual J56 all-top normalized kernel, H=133, f=168, support s=25, the original E/F, both CM components, every actual pencil and omitted endpoint, finite projectors, class masses, selected real trace, packet point, central phase and Haar measures. Work in the DO099(JR24)/R071 remainder with

\[
 E^{169}=Z\perp P,\quad Z=S_{164}\perp V_3,\quad
 P=Eu\perp Ee_n,\quad h_{ZZ}=I,\quad h=I+N,
 \quad \operatorname{rank}_E N=2,\quad N^3=0.                 \tag{EL1}
\]

The cross map C:Z→P is nonzero on S. All25 active relative T places have actual type II; all334 unmixed G reductions are unipotent. Global rank-two unipotence is a hypothesis, not inferred here from the finite reductions.

Use the original normalization

\[
 C=\ell r,\quad\ell=(1,\eta),\quad\eta\bar\eta=U,
 \quad U=77Q^2-14,\quad Q=1427908193729121522601,
 \quad x=N_E D_\eta.                                       \tag{EL2}
\]

Here D\_eta is the exact integral denominator ideal. Write k for the number of nonzero original S coordinates of r, so1≤k≤164. Call an active place eligible exactly when eta−a0 is divisible at its original selected pi-component, where a0=Q(−7+2sqrt(−7)); this is DO103(LD4--11). Let p be the number of eligible active places.

\textbf{Theorem.} The union of the following two actual subsets has total absolute normalized contribution less than2\^{}-451:

\begin{enumerate}
\def\labelenumi{\arabic{enumi}.}
\tightlist
\item
  p≤24, with no hypothesis on x;
\item
  p=25, 2\^{}120\textless x\textless2\^{}1820, and1≤k≤21.
\end{enumerate}

The proof also establishes two exact restrictions: for each fixed25-place support, there is at most one fully eligible eta with x\textless2\^{}1820; for that eta, each active type-II place admits at most one of its actual diagonal pencil labels. The latter restriction concerns the complete weighted label sum after fixing eta. It is not a stronger bound for each previously fixed label tuple.

Together with the already charged DO103 and DO104 tails, the remaining EL1 branch therefore has

\[
 p=25,\qquad 2^{120}<x<2^{1820},\qquad 22\le k\le164,        \tag{EL3}
\]

with the stated uniqueness and local-label restrictions. No estimate for that dense complementary range, full support-floor increase or signed lower bound is asserted.

\subsection{2. The eligibility count does not require a denominator bound}\label{the-eligibility-count-does-not-require-a-denominator-bound}

DO103(LD6--11) proves in the actual paired pencils that every pole of an original C scalar and every double pole of any entry of N lies at an eligible place. The proof uses the two original pole lines and eta bar(eta)=U; the bound x≤2\^{}120 occurs only later, in LD15--16, to force p≤24. In particular the entire counting argument LD17--19 is valid whenever p≤24, independently of x.

For completeness, pad the eligible set to a24-element subset D of the25 active places, with at most25 choices. Every original C scalar lies in J I\_D, N\_E I\_D=T\^{}-48. Full entries have at most the poles I\_Sigma,F I\_D,F. The actual fixed-C central difference coset and its two counts are the unchanged DO102(NT8--9)/DO103(LD18):

\[
 \#w<2^{479}Q^{167}T z^{-2f},\qquad
 \#w\le134^fT^{49},\qquad z=(N_E(JI_D))^{1/(2f)}.           \tag{EL4}
\]

All164 original-coordinate cells, the shared G orientations, transverse metric, pivot, relative norm count and selected trace are the complete DO098(CF8--32) factors. At each active type-II place the full diagonal-label sum against \textbar rho\textbar=1/(T−1) is at most1. The two central counts differ from the support24 count by exactly T. Consequently

\[
       |\text{entire }p\le24\text{ subset}|
              <25T\,2^{-2879}<2^{-1878}.                    \tag{EL5}
\]

We freshly charge this larger containing subset in the theorem's union; the earlier DO103 charge remains in the chronological error list. There is no assertion that a previously charged set has acquired new members without an error charge.

\subsection{3. A fully eligible eta is unique in the bounded denominator window}\label{a-fully-eligible-eta-is-unique-in-the-bounded-denominator-window}

Fix the actual25-place support Sigma. At an eligible selected prime P, eta≡a0. The actual constants U and a0 are T-units. The identity eta bar(eta)=U therefore gives at the partner

\[
           \eta\equiv U/\bar a_0\pmod{\bar P}.              \tag{EL6}
\]

This is precisely the actual inverse-adjoint pole-line equality of DO103(LD7--11); it is not a second independently chosen residue condition. D\_eta is trivial at both components of every eligible place.

Suppose eta1 and eta2 are both fully eligible with denominator norms x1,x2\textless2\^{}1820. If they differ, the ideal

\[
                    (\eta_1-\eta_2)D_{\eta_1}D_{\eta_2}    \tag{EL7}
\]

is a nonzero integral ideal. Indeed each denominator clears the corresponding possible negative valuation, so their product clears the difference. By EL6 it is divisible by both degree-one CM primes at every place of Sigma, hence has norm at least T\^{}50. At every complex embedding, \textbar eta1−eta2\textbar≤2sqrt(U). Thus its norm is at most (4U)\^{}f x1 x2, whereas the exact actual-host integer comparison is

\[
                         (4U)^{168}2^{3640}<T^{50}.          \tag{EL8}
\]

This is impossible. Thus eta1=eta2. The uniqueness is per fixed relative support, before choosing any G pattern or pencil scalar label. We count neither denominator ideals, numerator ideals nor28 norm-one generators in this bounded fully eligible branch. Absence of such an eta is allowed and only decreases the count.

\subsection{4. One local label, including the central-only pole case}\label{one-local-label-including-the-central-only-pole-case}

Fix an active selected T component, with DVR O and uniformizer pi. Let Lambda be the original lattice. The literal actual pencil is

\[
 M_c=A_c\Lambda,\quad A_c=I+c\pi^{-1}zq_*,\quad
 z=v_0/\bar\pi,\quad v_0=(Qu,\pi_k),\quad
 q_*=-v_0^*J-\pi m,\quad mz=1,\quad q_*z=0.                \tag{EL9}
\]

The label set is the residue field kappa minus\{−1\}, with infinity omitted. The label0 is included but cannot be active type II. K077(AP1--6), LIFT44(PR4--7) and K070(PT1--6)'s paired construction are the retained original interfaces. We may use m=e1\^{}t/z1: K063's \nolinkurl{scalar_pencil_small_field_plane_bound} (PF12) proves equality of the actual lattices after this replacement, by an integral right change of basis. No label is changed.

Write the ambient isotropic column and row as

\[
                 w=(u,\eta),\qquad w^\dagger=(u^*,-\bar\eta).
\]

They have a unit original positive coordinate at T, and w\^{}dagger w=0. Full eligibility and EL6 give units a\_l,b\_l with

\[
 w\equiv a_lz\pmod\pi,\qquad
 w^\dagger\equiv b_lq_*\pmod\pi.                           \tag{EL10}
\]

One may take a\_l≡bar(pi)/Q and b\_l≡−1/Q. Define the integral scalars a=q\_\emph{w/pi and b=w\^{}dagger z/pi. To compare the two cancellation labels, choose integral lifts w=a\_l z+pi v and w\^{}dagger=b\_l q\_}+pi t. Expanding w\^{}dagger w=0 and using q\_*z=0 gives

\[
                            b_l a+a_l b\equiv0\pmod\pi.    \tag{EL11}
\]

Now

\[
 A_c^{-1}w=w-cza,\qquad w^\dagger A_c=w^\dagger+cbq_*.
                                                               \tag{EL12}
\]

Both are integral for every c.~Their reductions vanish respectively when a\_l−ca=0 and b\_l+cb=0. If a=0, EL11 forces b=0, and neither can vanish. Otherwise each vanishes at the same unique label c\_eta=a\_l/a=−b\_l/b. It may be an omitted label, in which case there is no admissible active arrow of the stated kind. This equality uses the actual paired row and column, not two unrelated local choices.

We show that any active type-II arrow preserving M\_c must have c=c\_eta. In ambient notation the exact EL1 normal form is

\[
 N=wR_0-V_0w^\dagger+dww^\dagger,\qquad
 R_0|_P=0,\quad V_0\in Z,\quad
 V_0=J_Z^{-1}R_0^*,\quad d+\bar d=-R_0V_0.                 \tag{EL13}
\]

Every fixed original S input and y1,y2,z0 is in each source and target lattice at both components. They form an integral Z frame at T because its determinant has only the previously retained factor2U, which is a T-unit.

If an R\_0 value has a pole at the selected component, applying N to the corresponding integral Z input gives w times that value in M\_c; thus A\_c\^{}-1 w vanishes modulo pi, forcing c\_eta. If the conjugate R\_0 row has a pole, apply the same argument in the paired Hermitian-dual lattice. Its condition, transported back by the exact inverse-adjoint pairing, is w\^{}dagger A\_c≡0, again forcing the same c\_eta. No independent partner label or sign multiplicity appears.

It remains possible that both cross rows are integral but N has a central pole. In that case R\_0 and V\_0 are integral in the fixed T-unit Z frame. Since z is in P, R\_0z=0. Since v0 is in P, v0\^{}*J V\_0=0. Therefore

\[
 R_0A_c=R_0,\qquad
 q_*V_0=-\pi mV_0,\qquad
 A_c^{-1}V_0=V_0+czmV_0.                                   \tag{EL14}
\]

Combining EL12--14, the conjugated mixed part is integral. If d were integral then the original N would also be integral, contrary to the active support. Hence v\_pi(d)\textless0. The remaining conjugated term is

\[
                  d(A_c^{-1}w)(w^\dagger A_c).              \tag{EL15}
\]

Its integrality forces at least one of the two integral vector factors to vanish modulo pi: if both have a unit entry, their outer product has a unit entry, which cannot cancel an integral mixed part. EL11--12 then force c=c\_eta. This settles the central-only case, including a genuine double pole, without assuming that a cross pole is present.

The paired condition has already been used exactly; selected integrality and its partner are linked by unitarity and the determinant-unit condition, as in DO093/R059. Thus there is at most one actual diagonal pair per active relative place. After fixing the unique eta from Section3, the complete25-place pair/weight sum is bounded by

\[
                    (T-1)^{-25}<2T^{-25}.                  \tag{EL16}
\]

This follows from the exact \textbar rho\_II\textbar=1/(T−1), not from averaging over arbitrary matrices. It retains the original two omitted points and the zero-label case.

\subsection{5. A joint-energy row/transverse count and the moving central cover}\label{a-joint-energy-rowtransverse-count-and-the-moving-central-cover}

We retain DO104(CD10--21) before its coarse endpoint substitutions. Let J be the canonical whole-C G ideal, with at most4\^{}334 patterns, N\_EJ≤Q\^{}334 and Green bound2Q\^{}(-335·167)(N\_EJ)\^{}(335/2). Put

\[
 L'=JI_\Sigma D_\eta,\quad z=(N_EL')^{1/(2f)},\quad
 R=17424/133,\quad M=\sqrt{R/U},\quad
 z_G=Q^{167/f}T^{-25/f}.                                    \tag{EL17}
\]

All167 fixed r values lie in L' by the exact intersection L∩eta\^{}-1L=L D\_eta. Their full positive squared norm is at most M² at each embedding. The original minimum and global G cap give

\[
                   z\le A_k(x):=\min(M/\sqrt k,z_Gx^{1/(2f)}).
                                                               \tag{EL18}
\]

For the k nonzero original S entries, the general R072(PJ10--12) Gaussian packing argument is valid at every total half-trace radius sqrt(f)r. Its radial counting function satisfies

\[
 B_k(r)\le K_k r^{2fk}z^{-2fk},\quad
 K_k={164\choose k}47^k2^{2fk}(f!)^k
                  \left(\frac{87/32}{fk}\right)^{fk}.       \tag{EL19}
\]

Here B\_k(0)=0 because k≥1. The exact three-coordinate Gram matrix G on y1,y2,z0 has determinant4U. Its actual L' product metric has covolume D0\^{}3 z\^{}(6f)/(4U)\^{}f and covering radius less than epsilon sqrt(f)z, with epsilon=1/1024, by DO104(CD20). Only these transverse coordinates are rounded; the S entries retain their exact positions. A point of joint radius at most sqrt(f)M, after adjoining its transverse Voronoi cell, lies in the joint ball of radius sqrt(f)(M+epsilon z).

For each fixed S point of scaled length r, the volume of its allowed transverse slice is bounded by the three-coordinate ball volume with squared radius (M+epsilon z)\^{}2−r². Sum these slices and integrate the radial counting function EL19 by parts. With a=M+epsilon z the exact beta integral is

\[
 \int_0^a(a^2-r^2)^{3f}\,d(r^{2fk})
       =\frac{(fk)!(3f)!}{((k+3)f)!}\,a^{2(k+3)f}.           \tag{EL20}
\]

The decreasing slice weight permits the polynomial upper bound for B\_k after integration by parts; boundary terms vanish at a and0. Thus the full row/transverse count is at most

\[
 K_k A_3(4U)^f\,\mathcal B_k\,
       (M+\epsilon z)^{2(k+3)f}z^{-2(k+3)f},\quad
 \mathcal B_k=\frac{(fk)!(3f)!}{((k+3)f)!},\quad
 A_3=\frac{528^{3f}}{(3f)!D_0^3}<2^{-4040}.                 \tag{EL21}
\]

This uses the single full positive norm constraint, rather than multiplying two independent M-radius bounds. It introduces no lattice density assumption or additional basis integrality.

For fixed eta and r, DO104(CD14--19)'s original4×4 Bezout extraction gives a single affine central difference coset in iJ0 B\_eta I\_Sigma,F², with covolume sqrt(DF)(N\_FJ0)xT\^{}-50. Since N\_FJ0≤N\_EJ, its full-real covering radius is less than750sqrt(f)z². Keep that z dependence, instead of replacing it by its k endpoint. Thus the complete central count is

\[
 \#d\le V_F(H/2+750Uz^2)U^{-f}T^{50}x^{-1},\quad
 V_F(C)=\frac{((22/7)fC^2)^{f/2}}{(f/2)!\sqrt{D_F}}.        \tag{EL22}
\]

We dropped only N\_FJ0≥1. The real-part shift d+bar(d)=−q is fixed after r; no assertion places d itself in the difference ideal and no half-lattice factor is inserted. Given eta,r,d, the exact normal form reconstructs h; hence no additional source-isometry or projected-column fibre is required. The harmless old first-pivot factor164\textless2\^{}8 is retained.

Finally the unchanged actual selected trace, with both compact factors and its original near-point normalization, is

\[
                  |c(h)|<2^{-929f/2}U^{-f/2}z^{-f}.          \tag{EL23}
\]

This is DO104(CD21), not an operator bound. All class/phase/Haar normalization remains exactly the accepted one, with the remaining factor of modulus at most1.

\subsection{6. One denominator endpoint and the exact21-cell union}\label{one-denominator-endpoint-and-the-exact21-cell-union}

Apart from constants independent of x,z, the product of the full counts, Green factor and real trace is

\[
 x^{-337/2}z^{b_k}
 (1+\epsilon z/M)^{2(k+3)f}
 V_F(H/2+750Uz^2),\qquad b_k=(328-2k)f\ge0.                \tag{EL24}
\]

For fixed x this is nondecreasing in z, so substitute A\_k(x). If x2≥x1\textgreater0, then A\_k(x2)/A\_k(x1)≤(x2/x1)\^{}(1/(2f)), even when the minimum changes branch. The three increasing factors in EL24 grow by at most the powers164−k, k+3 and1 of x2/x1, respectively. Including x\^{}-337/2 leaves the strictly decreasing power−1/2. This proves a single common endpoint bound on x≥X0=2\^{}120, with no union over cap branches or denominator ideals.

The actual-host exact comparisons give

\[
 R^{168}T^{50}>68036^{168}U^{168}Q^{334}2^{120},\quad
 A_k(X_0)<M/\sqrt{68036}<M/260,\quad
 H/2+750R/68036<68.                                        \tag{EL25}
\]

Since k≤164, the68036 cap is the applicable safe one for every coordinate cell. Moreover

\[
 (1+1/(1024\cdot260))^{2\cdot167\cdot168}<2,\qquad
                         V_F(68)<2^{479}.                  \tag{EL26}
\]

Keep once each Green2, support2\^{}168, common patterns2\^{}668 and pivot2\^{}8, for exponent845. Keep the additional factor2 in EL16 and factor2 in EL26. The total T exponent is now337s−s=336s. Regroup the positive total M power334f before substituting any directed logarithm; it contributes R\^{}(167f)U\^{}(-167f), and the other factors leave total U exponent−335f/2. Use

\[
 L_T=1019333/1024>\log_2T,\quad L_Q=287843/4096<\log_2Q,
 \quad L_U=601356/4096<\log_2U,\quad
                         L_R=28810/4096>\log_2R.            \tag{EL27}
\]

The last inequality is newly checked by exact integer powers; its coefficient167f is positive. The complete common logarithmic upper bound, before the central479 and rational k-factor, is

\[
\begin{split}
 C_*={}&845+1+1+2f-4040-929f/2+167fL_R-(335f/2)L_U\\
       &+336\cdot25L_T-335\cdot167L_Q-(337/2)\cdot120
          =1618188837/4096.                                \tag{EL28}
\end{split}
\]

For each1≤k≤164 let c\_k be the least integral strict binary ceiling of

\[
 \frac{\binom{164}{k}47^k2^{2fk}(f!)^k87^{fk}(fk)!(3f)!}
 {32^{fk}f^{fk}k^{fk}68036^{(164-k)f}((k+3)f)!}.             \tag{EL29}
\]

The certificate compares exact integers for every cell. The full exponents are E\_k=C\_*+479+c\_k. In the claimed sparse range the maximum is

\[
 E_{21}=-1852891/4096<-452,\qquad
 \max_{k\le20}E_k=E_{20}=-10204635/4096<-2491.              \tag{EL30}
\]

The exact dyadic comparison

\[
                    2^{-1878}+20\,2^{-2491}+2^{-452}<2^{-451}
                                                               \tag{EL31}
\]

includes the entire fresh EL5 subset and every claimed middle-window coordinate cell. This proves the theorem with one charge451. As a documented boundary, the same sufficient envelope has E22=6457893/4096\textgreater0 and E164=1007770149/4096\textgreater0. Such positive upper estimates do not assert a positive, nonzero or existing contribution; the entire k22..164 middle range remains outside this conclusion.

\subsection{7. Exact verification, chronology and retained gates}\label{exact-verification-chronology-and-retained-gates}

The companion checks the literal host comparisons, paired denominator valuation identity, exact cancellation relation and full inverse-adjoint local models. Its local models include both genuine central-only poles and poles in each cross row; they supplement the general DVR proof and do not claim globally admissible short matrices. It also checks the beta-factor formula, all164 exact factorial ceilings, the common endpoint monotonicity exponents, the complete claimed dyadic union and both full chronological error lists.

Direct source/interface dependencies are the whole DO103(LD1--19), DO104(CD10--32), R072(PJ4--12), K077(AP1--7), the actual scalar-pencil change of row in \nolinkurl{scalar_pencil_small_field_plane_bound} (PF12), and the exact paired normalization already retained in DO093/R059. Their settled norm-one, field, and selected-trace results are cited without new unproved extensions. The new steps are the denominator-independent eligibility use, bounded fully eligible uniqueness, coincident row/column label cancellation including central-only poles, the joint-energy beta integral and the simultaneous x/z endpoint.

The default chronological predecessor is \nolinkurl{checks/r084_joint_remainder_error.json}, with96 all-top and31 uniform charges and complete floors25/11. Append451 once, retaining every earlier charge including DO1031878, DO1042518 and R0845398. The script accepts \nolinkurl{--prior} for a separate unchanged-byte chronological replay if another canonical finding intervenes. Neither complete floor changes.

The original fractional ideals and affine fibres are invertible-module conditions on Spec(O\_E)/Spec(O\_F), with the same CM involution, Hermitian polarization, theta line, biextension and complex metric. Exact inherited contextual anchors remain SGA3 I4.7.3/4.7.3.1; SGA6 VI Theorem3.3; FGA232 Theorem3.1 pp148--149; FGA236 Corollaries1.7/2.3 pp229/232; SGA7 I VII2.7/2.9 pp169--174 and VIII7.4(b) pp310--311; cote4 scan003 §§8.1.1--8.1.2 and scans005--006 §§8.4.1--8.4.3; cote54 scans005--007. No new lambda-ring/K-theory equality, l-adic or motivic comparison is inferred. The signed sufficient inequality, literal source constructor, prescribed Bessel/Mram period, selected arithmetic content, specific032/Read/STiesR pairing, DP44 and RH conclusions remain open.

\clearpage

\section{AM-LIFT-082 --- arbitrary-type frame-polynomial transfer with all local extensions}\label{am-lift-082-arbitrary-type-frame-polynomial-transfer-with-all-local-extensions}

\textbf{Canonical checked integration --- 2026-11-23T01:27:26.206675+00:00.} Complete root and independent whole-proof/script/source-interface reviews, with separate23,542-assertion exact replays, passed. The frozen historical source remains unchanged. LIFT082 FT1--29 transfers one actual nonzero frame polynomial to every Step5Loc prime and complete-DVR extension, and every Step5Arch embedding, from explicit arbitrary-type geometric inputs. Generic FrPt polynomial existence and same-point nonzero evaluation produce one G before any completion. Integral eigenmodule bases commute with arbitrary compatible DVR extension by flatness; no integral eigenspace splitting is assumed. The exact relation Ub\^{}t Y=X, intrinsic theta generator unit, determinant inclusion and physical row coefficient-clearing indices give all local clauses. The original pi-normalized Hermitian metric, actual fibre/section maps and goodOf conjugation give all archimedean clauses, retaining lieFac and the dual theta norm. There is no new transfer cost beyond LIFT080/081. The actual prescribed point with Dw.mult, good model/level/Key/section/coefficient-field ties, valSet arithmetic-lattice identification, same-point nonzero F value, Host/Sizes, Read/STiesR, selected arithmetic image/content, global period, specific032 and DP44 remain required. No old CMPt.Good/LamObjGood/LamPtClauses instance, finite charge, support-floor change or RH assertion follows.

Author: \nolinkurl{/root/lift_accounting}, 23 November 2026 UTC. \textbf{Author whole-proof/source/script review PASS, 23 November 2026 UTC; exact completion and final hashes are recorded in LIFT082\_AUTHOR\_VERIFICATION.json.} This is a mathematical transfer of the local and archimedean arguments in the pinned source, with their geometric inputs displayed separately. It does not add or modify a Lean declaration. It constructs no new instance of the old quadratic-type \nolinkurl{CMPt.Good}, \nolinkurl{LamObjGood} or \nolinkurl{LamPtClauses}.

\subsection{1. Actual inputs, including the point and its signature}\label{actual-inputs-including-the-point-and-its-signature}

Fix the genuine source \nolinkurl{Host E}, \nolinkurl{Datum H}, level, moduli object \nolinkurl{Xo}, and coefficient object \nolinkurl{c}. Retain the source geometry laws used by the cited frame, section, base-change and uniformization theorems. Let K be a number field with its \textbf{specified} c.F-algebra structure. It need not be a CM field. A change of K must still satisfy every independent coefficient-field/Galois condition of the application; this theorem does not furnish such a change.

The following data and equalities are hypotheses, rather than conclusions inferred from a dimension count.

\textbf{FT1 (actual fibre and section).} There is a framed point

\[
 y\in\operatorname{FrPt}(X_o,c,K),\qquad x=y.x,
 \tag{FT1}
\]

and a genuine \nolinkurl{FibKey} at x. Its principally polarized generic abelian variety, rigidified theta line, nonzero section theta, complete tangent basis v, full-order action and map to the universal fibre are tied exactly as in \nolinkurl{FibKey}, \nolinkurl{SecTie} and \nolinkurl{SABMap}. The section of y is the image of this same theta under \nolinkurl{SecTie}. The target has the supplied good polarized integral model over O\_K, with the compatible full O\_E action and actual generic-fibre identification. In particular this is a point of the specified moduli problem with the specified level; an arbitrary abelian variety of the right dimension is insufficient.

Write g for the target dimension. We require g\textgreater=1 and the literal condition 1\textless=Dw.p+Dw.q. Let m\_tau=\nolinkurl{Dw.mult tau}. The actual tangent letters u\_tau,i, for 1\textless=i\textless=m\_tau, form the frame of y, after applying the universal-fibre map. They are actual E-eigenvectors, and the disjoint union of their label lists has \textbf{exactly} these multiplicities. Choose the source index equivalence

\[
 e:\coprod_{\tau:E\hookrightarrow\mathbf C}\{1,\ldots,m_\tau\}
       \simeq\{1,\ldots,g\}.
 \tag{FT2}
\]

Let U be the matrix whose row e(tau,i) is the coordinate row of u\_tau,i in v. It is invertible. One sufficient construction is LIFT080 MT12--14: actual E-linear component homomorphisms whose product map is an isogeny, with complete component eigenbases and precisely the required total multiplicities. The component CM types may be arbitrary. This assertion does not say that arbitrary component types have the required total multiplicities.

\textbf{FT3 (actual line lattices and metrics).} Retain the genuine theta line T and determinant line Delta, with their source \nolinkurl{MLine.Good} structures and exact \nolinkurl{KeyLocAt} and \nolinkurl{KeyArchAt} identifications for the above \nolinkurl{KeyIn}. More concretely, at each finite prime, T is the line of scalars making the fixed theta section integral, and Delta is the determinant line of the actual integral tangent lattice in the basis v. Thus a local generator c\_a of T.I makes c\_a theta a generator of the rank-one module of sections, up to a local unit; a local generator c\_D of Delta.I is the determinant of an integral tangent basis, up to a local unit. These identifications are retained, not inferred from a chosen analytic normalization. They are the identifications implemented by \nolinkurl{KeyIntegral}, \nolinkurl{KeyLocAt}, and the section-generator comparison in the pinned source.

For each letter retain an actual good metrized line ell\_j such that

\[
 c\in\ell_j.I\quad\Longrightarrow\quad c u_j
     \in\operatorname{Lie}(B_{\mathcal O_K}),
 \qquad
 \|u_j\|_\sigma\le L\,\ell_j.n(\sigma),\qquad L\ge1.
 \tag{FT3}
\]

These are geometric input implications. Actual homomorphisms extending over the stated models and the polarization pullback identities of LIFT080 MT14 prove them for arbitrary component types. The norms are the source norms: the tangent Hermitian form is pi times the Appell--Humbert form, with the source conjugate-first convention. No embedding-coordinate ideal is substituted for the intrinsic tangent line. In particular the lower degree estimates of LIFT081 apply only to its actual component lines and their permitted base changes.

\textbf{FT4 (actual form and nonzero evaluation).} Let F belong to the source space \nolinkurl{cusp Xo c Unf a n}, with \nolinkurl{Dom a n}, a\_tau,i nonnegative integers, and n\textgreater=0. Keep the actual uniformization family at every good embedding, its \nolinkurl{IsUnf} witnesses, and the common absolute value xi of its xi0 parameters, exactly as required by \nolinkurl{cusp_bound_allEmb}. Require the arithmetic condition

\[
 \forall\mathfrak l\subset\mathcal O_{c.F}\ \forall
       v\in\operatorname{valSet}(X_o,c,\mathrm{algCls},\mathfrak l),
       \quad v(F)\ge0.
 \tag{FT4}
\]

When F is supplied through \nolinkurl{recB.Mlat}, the actual \nolinkurl{LiesTotal}/arithmetic dictionary identifying that lattice with FT4 must be supplied or cited. Membership in a differently defined lattice is not automatically FT4. The exact source \nolinkurl{nonneg_iff_of_liesTotal} isolates this requirement.

Finally require one nonzero value over this \textbf{same underlying point}: for some embedding sigma\_1:K→C over its specified embedding of c.F and some complex frame w,

\[
 w.x=\sigma_1^*x,\qquad F(w)\ne0.
 \tag{FT5}
\]

It is enough to know F(y)!=0. A nonzero form somewhere on the moduli space, or a nonzero theta section, does not imply FT5 at the chosen point.

\subsection{2. The single global polynomial and its exact weight}\label{the-single-global-polynomial-and-its-exact-weight}

\textbf{Theorem FT6.} Under FT1--5 there is one nonzero polynomial G over K in the block matrix variables \nolinkurl{FrVar Dw.mult} such that:

\begin{enumerate}
\def\labelenumi{\arabic{enumi}.}
\tightlist
\item
  G is the literal frame polynomial of F at y, with all the source base-ring quantifiers;
\item
  the same G, U, line ideals and theta section satisfy \nolinkurl{Step5Loc e T.dual Delta ell U G n p} at every prime p of K;
\item
  the same G satisfies \nolinkurl{Step5Arch} at every complex embedding, with the actual source supremum
\end{enumerate}

\[
 S_\sigma=\operatorname{supC}
 (X_o,c,\mathrm{Unf},\xi,\operatorname{goodOf}(\sigma|_{c.F}),a,n,F).
 \tag{FT6}
\]

There is no new finite denominator or extra archimedean factor in these two assertions. All geometry, levels, fields, sections and input norms in section1 remain in force. If an application also needs sum a\textless=aMax or positivity of every S\_sigma, those premises are retained; the source \nolinkurl{SupDef} implies the latter for nonzero F.

\textbf{Proof of the polynomial assertion.} The directly read theorem \nolinkurl{GlobalTheory.LamT.exists_isFrPoly} applies to an arbitrary \nolinkurl{FrPt} over a number field and F in \nolinkurl{sec a n}. It does not assume \nolinkurl{CMPt.Good}. Its proof first represents the frame rule as P/(product det Y)\^{}N. The nonnegative Borel weight permits cancellation of the determinant denominator and gives G. Thus the conclusion is a polynomial, not a rational function evaluated only on invertible matrices.

In the literal definition \nolinkurl{IsFrPoly}, for \textbf{every commutative c.F-algebra S}, every c.F-algebra map phi:K→S, compatible base-changed frame z, and block matrices Y with unit determinants, one has

\[
 \operatorname{FrChg}(Y^{\mathsf t},1)(z)=z'
 \quad\Longrightarrow\quad
 F(z')=\operatorname{evalFr}(\phi G,Y).
 \tag{FT7}
\]

The literal variable convention is \nolinkurl{evalFr G Y = eval(G, v ↦ Y[v.1][v.2.2][v.2.1])}: a variable v=(tau,i,j) evaluates to Y\_tau{[}j,i{]}. Thus the first variable index is the physical matrix column, and the second is the physical row. This G is chosen before any finite-prime completion or archimedean embedding. The same-point frame-change theorem writes w in FT5 as an invertible block frame change followed by multiplication of the section by a nonzero scalar. The scalar contributes its n-th power. Hence G=0 would force F(w)=0. This proves G!=0 without choosing new polynomials separately at different primes.

Write G(Y) for evalFr(G,Y). For any invertible upper triangular block matrix R, frame covariance gives

\[
 G(YR)=\left(\prod_{\tau,i}R_{\tau,ii}^{a_{\tau i}}\right)G(Y).
 \tag{FT8}
\]

The cited \nolinkurl{isFrPoly_equiv_and_ne_zero} proves this exact \nolinkurl{IsEquiv} after every complex embedding. It also shows the column-weight convention: diagonal R scales \textbf{columns} of Y. Polynomial identity over the infinite field extends the identity from invertible Y. Consequently every nonzero monomial has degree a\_tau,i in physical matrix column i (the first \nolinkurl{FrVar} index) of block tau, and total degree

\[
 D=\sum_{\tau,i}a_{\tau i}.
 \tag{FT9}
\]

Zero weights, zero-dimensional blocks and n=0 are allowed. The target dimension g remains positive. No new source bound on D is inferred.

\subsection{3. Local bases, section units and the determinant inequality}\label{local-bases-section-units-and-the-determinant-inequality}

Fix a nonzero prime p of O\_K. Choose local generators c\_a,c\_D,c\_j of T.I, Delta.I and ell\_j.I. Let R be the completed local DVR, K\_p its fraction field, and f:K→K\_p the canonical embedding. Its residue field is finite, and

\[
 f(x)\in R\ \Longleftrightarrow\ |x|_p\le1.
 \tag{FT10}
\]

The model and every actual fibre/section/action tie base change to R. We use the integral model belonging to the actual point; no local replacement of the prescribed point or level occurs.

At a finite place, write tau\_R(a) for the source scalar \nolinkurl{f((algebraMap c.F K)(embK c tau a))}; it is the transported algebraic eigenlabel, not the complex number tau(a) inserted into R. The coefficient-field action and this transport are part of FT1. For each tau take the integral eigenspace

\[
 L_\tau=\{v\in\operatorname{Lie}(B_R):
       \iota(a)v=\tau_R(a)v\text{ for all }a\in\mathcal O_E\}.
 \tag{FT11}
\]

All tau\_R(a) belong to R because they are algebraic integers. This is a finite torsion-free R-module of rank m\_tau and is therefore free. Its rank follows by inverting the uniformizer: a finite basis of O\_E gives finitely many kernel equations, so localization recovers the characteristic-zero eigenspace. Choose its basis beta\_tau and write the original row frame as

\[
 u_\tau=U_{b,\tau}\,\beta_\tau.
 \tag{FT12}
\]

Every U\_b,tau is invertible. FT3 implies

\[
 f(c_{e(\tau,i)})\,(U_{b,\tau})_{ij}\in R.
 \tag{FT13}
\]

Equivalently, if beta\_tau=g\_tau u\_tau in the source frame-change convention, U\_b,tau=g\_tau\^{}\{-1\}. This is precisely the inverse, not the direct, frame-change matrix in \nolinkurl{Step5Loc}. The proof of \nolinkurl{letter_integral_in_eigenbasis} uses this relation after its generic \nolinkurl{image_integral_of_isLocGen} step. Its pure-type wrapper only supplies the integral maps and actual letter identities, which FT1--3 provide here.

The direct sum of the modules L\_tau need \textbf{not} equal the whole integral tangent lattice when embeddings collide modulo p.~No integral idempotents or etale splitting at ramified primes are asserted. The literal integral framed-point predicate requires each displayed eigenframe to be an R-basis of its own eigenspace, which is exactly what has been constructed.

Choose an integral full tangent basis of B\_R, whose coordinate matrix in v is M. By the actual determinant-line identification, det M=f(c\_D)epsilon\_D for epsilon\_D in R\^{}*. Every scaled letter f(c\_j)u\_j lies in this full tangent lattice. The determinant of their coordinate matrix is therefore an R-multiple of det M. Thus

\[
 \frac{f((\prod_jc_j)\det U)}{f(c_D)}\in R,
 \qquad
 |(\prod_jc_j)\det U|_p\le|c_D|_p.
 \tag{FT14}
\]

This is the exact multiplicative valuation direction in \nolinkurl{Step5Loc}, without a saturation or index equality.

For the theta line choose an integral generator theta\_0. By FT3,

\[
 \theta_0=s\theta,
 \qquad s=f(c_a)\epsilon_a,\quad\epsilon_a\in R^*.
 \tag{FT15}
\]

The directly read source \nolinkurl{integral_generator_unit} establishes this same unit relation from the \nolinkurl{KeyLocAt}/model comparisons. We have made its actual intrinsic-model identification explicit in FT3. A principal theta section may vanish at the identity; it is a nonzero generator of its module of sections, not a unit-valued theta null. Put

\[
 t=c_a^{-1};\qquad t\text{ is a local generator of }T^\vee.I.
 \tag{FT16}
\]

\subsection{4. Every complete-DVR extension in Step5Loc}\label{every-complete-dvr-extension-in-step5loc}

Now let W be \textbf{any} complete DVR with finite residue field, with fraction field K\_W, an R-algebra structure that is local, and a field map h:K\_p→K\_W satisfying

\[
 h(a)=a\text{ under the compatible maps }R\longrightarrow W
       \longrightarrow K_W.
 \tag{FT17}
\]

This is the entire extension range in \nolinkurl{Step5Loc}. No bound on residue degree or ramification is used, and W is not assumed finite as an R-module. The compatibility and injectivity of the field map make R→W injective. W is torsion-free over the DVR R and hence flat.

The eigenmodule equations in FT11 are finitely many R-linear kernel equations, using a finite Z-basis of O\_E. Flatness preserves their common kernel. The tangent module of the smooth group and the given action also commute with this base change. Hence beta\_tau base changes to a W-basis of the actual W-eigenmodule. This is the generic content of \nolinkurl{flat_of_inj}, \nolinkurl{eig_flat} and \nolinkurl{lieInt_baseChange}; it applies even when the residue eigencharacters coincide. The integral principal section theta\_0 remains a generator after the same base change. The source \nolinkurl{secInt_baseChange} carries the exact rigidification, universal-fibre and \nolinkurl{SecTie} equalities, using principal-polarization section base change.

Let Y\_tau be block matrices over K\_W and X\_tau matrices over W with unit determinants, satisfying the literal relation

\[
 h(U_{b,\tau})^{\mathsf t}Y_\tau=X_\tau.
 \tag{FT18}
\]

Taking transposes gives Y\_tau\^{}t h(U\_b,tau)=X\_tau\^{}t. Consequently the frame Y\_tau\^{}t u\_tau is X\_tau\^{}t beta\_tau. It is therefore an integral basis of the correct eigenmodule. Together with h(theta\_0) and the unchanged base-changed level/fibre data this gives an actual integral framed point y\_int. This step explains both occurrences of transpose; replacing FT18 by U\_b Y=X would be a different assertion.

The source \nolinkurl{evR_mem_range_of_valSet_nonneg} applies FT4 to this integral frame and gives

\[
 F(y_{\rm int})\in W.
 \tag{FT19}
\]

It applies at every allowed W, not only at the first completion or a finite prescribed list of extensions. The hypotheses finite residue field, completeness, actual action and actual section ties are retained exactly as in that theorem.

The frame polynomial FT7 is evaluated with section theta, whereas y\_int uses h(s)theta. The literal section weight is n, so if w\_0=F(y\_int) in W,

\[
 w_0=h(s)^n\operatorname{evalFr}(h(fG),Y).
\]

Using FT15--16 yields the required identity

\[
 \operatorname{evalFr}(h(fG),Y)
   =h(f(t^n))\,w,
 \qquad w=h(\epsilon_a)^{-n}w_0\in W.
 \tag{FT20}
\]

The image of epsilon\_a is still a unit. There is no discarded unit, theta amplitude, different, determinant denominator or scalar normalization. FT10, FT13--14, FT16 and FT20 are all the clauses of \nolinkurl{Step5Loc}. The same proof applies to every p with the G already fixed in section2.

For clarity, the ensuing coefficient assertion also retains the correct variable index. If mu is a monomial index with variable v=(tau,i,j), its associated letter is \textbf{e(tau,j)}, the second \nolinkurl{FrVar} index and the physical matrix \textbf{row} of Y. This is different from the first index carrying the right-diagonal weight in FT8. The source \nolinkurl{CoeffIntegralLoc} then gives

\[
 [G]_\mu\prod_{v=(\tau,i,j)}c_{e(\tau,j)}^{\mu_v}
        \in t^n R.
 \tag{FT21}
\]

Intersecting these local conditions gives the actual coefficient in the fractional ideal of

\[
 (T^\vee)^{\otimes n}\otimes
 \left(\bigotimes_{v=(\tau,i,j)}
          \ell_{e(\tau,j)}^{\otimes\mu_v}\right)^\vee,
 \tag{FT22}
\]

which is the literal \nolinkurl{lineNN}. At least one coefficient is nonzero because G!=0. Nothing here identifies that coefficient with a separately prescribed theta derivative, selected class sum or 032 value.

\subsection{5. The actual metric, all embeddings, and the same polynomial}\label{the-actual-metric-all-embeddings-and-the-same-polynomial}

Fix sigma:K→C. Use the actual \nolinkurl{KeyArchAt} and \nolinkurl{SecTie} data for the fibre of y. Let H\_sigma be the positive Hermitian matrix of the source tangent metric in v. Its entries are

\[
 (H_\sigma)_{ij}=\pi\,\operatorname{herm}_{H_m}(W_j,W_i),
 \tag{FT23}
\]

where W\_i are the differentials of v\_i in the actual Appell--Humbert uniformization. The source convention is conjugate-linear in the first argument of the induced coordinate form; the reversal in FT23 agrees with its \nolinkurl{hermForm} convention. The generic \nolinkurl{form_eq_pi_mul_hermForm} proves this exact identity under changes of the actual fibre and uniformization.

The principal polarization and the actual fundamental-domain volume imply

\[
 \Delta.n(\sigma)^2=(\det H_\sigma).\operatorname{Re},
 \qquad \Delta.n(\sigma)=\sqrt{(\det H_\sigma).\operatorname{Re}}.
 \tag{FT24}
\]

This is the fully generic \nolinkurl{TopLieVector.sq_eq_det_gram}: its proof uses the unimodular Riemann-form identity volume(Fd) det(H\_m)=1 and the complete tangent basis. It has no hypothesis on a quadratic CM type. Positivity chooses the displayed square root. The row estimates of FT3 give

\[
 \operatorname{hq}_{H_\sigma}(\sigma U_j)
       \le (L\ell_j.n(\sigma))^2.
 \tag{FT25}
\]

Let y\_sigma be the actual sigma base change of y. Its block metric is exactly

\[
 \operatorname{Gr}_{\tau,ik}
 =\overline{\sigma U_{e(\tau,i)}}\,
       H_\sigma\,\sigma U_{e(\tau,k)}^{\mathsf t},
 \qquad
 \operatorname{IsMetric}(y_\sigma,T.n(\sigma)^2,\operatorname{Gr}).
 \tag{FT26}
\]

The equality follows from the actual fibre/section/differential ties, not only an abstract isometry of Hermitian vector spaces. The source generic helpers \nolinkurl{exists_fibre}, \nolinkurl{exists_unif} and \nolinkurl{form_eq_gram} implement those ties. In the original \nolinkurl{exists_isMetric} wrapper, the old pure-type data are used to provide these maps and the length estimates. FT1--3 provide the exact replacements.

For any block matrices Y satisfying Y\_tau\^{}* Gr\_tau Y\_tau=I, their determinants are nonzero. Change the actual frame by Y\^{}t and leave the section unchanged. The generic \nolinkurl{isMetric_frChg} transforms the metric by

\[
 N\longmapsto |a|^2N,
 \qquad\operatorname{Gr}_\tau\longmapsto
           Y_\tau^*\operatorname{Gr}_\tau Y_\tau;
 \quad\text{here }a=1.
 \tag{FT27}
\]

The result is an orthonormal actual frame with theta norm squared T.n(sigma)\^{}2. The source \nolinkurl{cusp_bound_allEmb} therefore gives

\[
 |\operatorname{evalFr}(\sigma G,Y)|
       \le T.n(\sigma)^n S_\sigma,
 \quad\text{or equivalently}\quad
 |\operatorname{evalFr}(\sigma G,Y)|\,(T^\vee.n(\sigma))^n
       \le S_\sigma.
 \tag{FT28}
\]

For an embedding whose restriction to c.F is not good, that theorem base changes the actual framed point by complex conjugation. It conjugates the rule value once, preserves its absolute value and theta norm, and uses \nolinkurl{goodOf} exactly as in FT6. It does not replace K by a chosen CM field or relabel an arbitrary real metric. In particular the source factor \nolinkurl{lieFac=exp(-D log(2 pi xi)/2)} already present in \nolinkurl{supC} remains present. FT8 and FT24--28 prove every clause of \nolinkurl{Step5Arch} for the same G. This completes Theorem FT6.

\subsection{6. Height consequence, with every cost retained}\label{height-consequence-with-every-cost-retained}

Assume in addition the unchanged \nolinkurl{KeyCore}, \nolinkurl{LemmaHStatement}, a block bound m\_tau\textless=r with r\textgreater=1, positive S\_sigma, and a proved lower bound d\_- on the normalized degrees of these actual ell\_j. Theorem FT6 supplies the two previously retained polynomial premises of LIFT080. Set V=gn and d={[}K:Q{]}. For every real e\_0 one obtains the explicit bound

\[
 \frac1d\sum_\sigma\log S_\sigma\ge
 -\left(\frac12+|e_0|+\frac12\log L\right)(D+V)
 -D\log(rL)-\frac{3D+V}{2}(e_0-d_-)_+.
 \tag{FT29}
\]

This is exactly LIFT080 MT5, not a new cost estimate. If the actual component lines are the ones covered by LIFT081, one may take d\_-=-10 log2 for J14 or -16 log2 for J56. The optimized or old-logOm costs of LIFT080/LIFT081 then apply without an additional term from this transfer. If the old source logOm is retained, the explicit extra cost remains (3D+V)(logOm/2-d\_-)\_+/2. The maps, physical norms, ideals and polynomial have not been rescaled by this theorem.

Passing to an extension to supply a genuine model or label set is permitted here only after the actual point, section, maps, line lattices and uniformizations have been base changed with their stated ties. The generic formula averages over all embeddings of the resulting field; \nolinkurl{mean_restrict} is applicable only when its function really is the restriction of the lower-field function. No omitted orbit, field discriminant, class normalization, changed-family DP44 term or source Galois condition is cancelled by calling this a transfer.

\subsection{7. Exact reduction and remaining source gates}\label{exact-reduction-and-remaining-source-gates}

\begin{longtable}[]{@{}
  >{\raggedright\arraybackslash}p{(\columnwidth - 4\tabcolsep) * \real{0.3333}}
  >{\raggedright\arraybackslash}p{(\columnwidth - 4\tabcolsep) * \real{0.3333}}
  >{\raggedright\arraybackslash}p{(\columnwidth - 4\tabcolsep) * \real{0.3333}}@{}}
\toprule\noalign{}
\begin{minipage}[b]{\linewidth}\raggedright
Old wrapper or remaining premise
\end{minipage} & \begin{minipage}[b]{\linewidth}\raggedright
What has been proved here
\end{minipage} & \begin{minipage}[b]{\linewidth}\raggedright
What is still required in an actual application
\end{minipage} \\
\midrule\noalign{}
\endhead
\bottomrule\noalign{}
\endlastfoot
Pure-type \nolinkurl{CMPt.Good} used by the local/metric wrappers & Their frame-polynomial arguments transfer from FT1--5 for arbitrary component types & An actual point with multiplicities \nolinkurl{Dw.mult}, its exact universal fibre and level \\
Polynomial existence and nonzero value & One polynomial over K, with all base-ring identities, is nonzero from FT5 & A nonzero evaluation of the actual F on that same underlying point \\
Finite models, eigenletters and theta line & Integral letters and the actual section/Key ties give all primes and all complete-DVR extensions & The good models, full integral action, maps and intrinsic line identifications themselves \\
Arithmetic source lattice & FT4 implies integral rule values on every constructed frame & The literal \nolinkurl{recB.Mlat}/\nolinkurl{LiesTotal}/\nolinkurl{valSet} identification where needed \\
Archimedean bound & The exact source metric and uniformizations give all embeddings and \nolinkurl{goodOf} & Actual \nolinkurl{KeyArchAt}, \nolinkurl{SecTie}, \nolinkurl{IsUnf}, xi0 and cusp/weight premises \\
Unequal line degrees & LIFT080 applies with no new transfer cost; LIFT081 supplies its scoped component bounds & Compatibility of those actual component lines with the chosen point \\
Literal \nolinkurl{LamPtClauses}/\nolinkurl{lamw_link} & No new instance of the old quadratic records is asserted & An authorized proved replacement of those records and every unchanged comparison interface \\
Main source and arithmetic conclusion & A pair of generic Step5 premises has been reduced to explicit geometric inputs & Host/Sizes, Read/STiesR, selected derivative/class content, actual global period and requested 032 value \\
\end{longtable}

In particular the argument does not reverse LIFT079. It never declares a nonelliptic mixed point to satisfy the old pure-type record. It also does not construct an F-dependent point satisfying FT5, prove a global period nonzero, or identify the coefficient in FT22 with the required arithmetic detector. These are substantive remaining gates. The theorem removes the need to reprove the local and archimedean polynomial transport after a genuine arbitrary-type replacement point satisfying the displayed inputs has been constructed.

\subsection{8. Exact source addresses and verification boundary}\label{exact-source-addresses-and-verification-boundary}

All source files are pinned to commit \nolinkurl{7cbf7a0c32fca1badf05931879a37c3f8a2b03b4}; file hashes, relevant declaration locators and retrieval provenance are recorded in the companion source manifest. No source bytes are edited, no Lean build is run, and no statement below is inferred merely from a theorem's name. The complete listed proof bodies used in the transfer were read.

The source directory abbreviations in this paragraph are \nolinkurl{C=Development/AbelianVarietyTheory/Proofs/CubicalLineBundles}, \nolinkurl{D=Development/AbelianVarietyTheory/Definitions/CMLattices}, and \nolinkurl{F=Formalism/LocalToGlobal}.

\begin{itemize}
\tightlist
\item
  \nolinkurl{D/IntegralOf.lean}: \nolinkurl{SecTie}, \nolinkurl{FibKey}, the old \nolinkurl{CMPt.Good}, generic \nolinkurl{IsFrPoly}, \nolinkurl{IsMetricAt}, \nolinkurl{LieInt}, \nolinkurl{SecInt}, and \nolinkurl{integral_of}.
\item
  \nolinkurl{C/ExistsIsFramePolynomialOfRuleNumberField.lean}: \nolinkurl{exists_isFrPoly}; \nolinkurl{C/FramePolynomialEquivariantAndNeZero.lean}: \nolinkurl{evalFr_mul_tri}, \nolinkurl{isFrPoly_equiv_and_ne_zero}.
\item
  \nolinkurl{C/LieLetterIntegralOfLocalGenerator.lean}: generic \nolinkurl{image_integral_of_isLocGen}; \nolinkurl{C/LieLetterIntegralInIntegralEigenbasis.lean}: \nolinkurl{coord_eq}, \nolinkurl{isEigV_of_generic}, and the full wrapper \nolinkurl{letter_integral_in_eigenbasis}.
\item
  \nolinkurl{C/ExistsIntegralEigenbasisOfFrameAtLocalModel.lean}: \nolinkurl{lieInt_exists_and_change}; \nolinkurl{C/IntegralEigenbasisBaseChangeDvr.lean}: \nolinkurl{flat_of_inj}, \nolinkurl{eig_flat}, \nolinkurl{lieInt_baseChange}.
\item
  \nolinkurl{C/ExistsUnitIntegralGeneratorPrincipalLocalModel.lean}, \nolinkurl{C/IntegralGeneratorAssociatedLocalGeneratorCMLine.lean}, and \nolinkurl{C/IntegralGeneratorBaseChangeValuationRing.lean}: the genuine generator, exact local unit and \nolinkurl{secInt_baseChange}.
\item
  \nolinkurl{C/RuleValueIntegralOfSupLatticeValuesNonneg.lean}: \nolinkurl{evR_mem_range_of_valSet_nonneg}; \nolinkurl{C/LocalLowerBoundCmPointFramePolynomialEveryPrime.lean}: the complete old-wrapper local proof, including all extension and transpose/scalar steps.
\item
  \nolinkurl{C/TopLieVectorNormSqEqGramDet.lean}: generic \nolinkurl{sq_eq_det_gram}; \nolinkurl{C/FormEqPiMulHermForm.lean}: generic \nolinkurl{form_eq_pi_mul_hermForm}.
\item
  \nolinkurl{C/ExistsMetricDataCmPointComplexEmbedding.lean}: generic \nolinkurl{exists_fibre}, \nolinkurl{exists_unif}, \nolinkurl{form_eq_gram}, and the complete old-wrapper \nolinkurl{exists_isMetric} proof; \nolinkurl{C/MetricChangeOfFrame.lean}: \nolinkurl{isMetric_frChg}.
\item
  \nolinkurl{C/CuspFormBoundedBySupNormAtGoodEmbeddingAllEmbedd.lean}: \nolinkurl{cusp_bound_allEmb}; \nolinkurl{C/ArchimedeanHypothesisLowerBoundCmPoint.lean}: full \nolinkurl{step5Arch_cm}; \nolinkurl{F/ProjectiveLattices/L2DefiniteCuspForms.lean}: \nolinkurl{goodOf}, \nolinkurl{lieFac}, \nolinkurl{supC}, \nolinkurl{SupDef}.
\item
  \nolinkurl{F/CMDeterminantIdeals/StatementB.lean}: \nolinkurl{KeyIntegral}, \nolinkurl{KeyLocAt}, \nolinkurl{KeyArchAt}; \nolinkurl{StatementS.lean}: \nolinkurl{CoeffIntegralLoc}, \nolinkurl{Step5Loc}, \nolinkurl{Step5Arch}, \nolinkurl{Step5}; \nolinkurl{StatementC.lean}: \nolinkurl{lineNN}; \nolinkurl{StatementG.lean}: the transposed \nolinkurl{evalFr} variable convention and \nolinkurl{substFr}.
\item
  \nolinkurl{F/ProjectiveLattices/NonnegIffOfLiesTotal.lean}: exact arithmetic class, \nolinkurl{valSet}, \nolinkurl{LiesTotal}, \nolinkurl{nonneg_iff_of_liesTotal}; \nolinkurl{F/ProjectiveLattices/RosatiLatticeNorms/CmFrameOfEq.lean}: unchanged \nolinkurl{LamTieSt} and \nolinkurl{lamw_link} boundary.
\end{itemize}

The companion diagnostics check exact local matrix orientation, unit/scalar powers, determinant and coefficient-ideal directions, ramified eigenmodule examples, Gaussian-rational metric covariance, column homogeneity and the quoted source fingerprints. They are finite arithmetic diagnostics, not a substitute for the universal arguments above, a formal proof build, or construction of the still-required actual point and nonzero form value.

\subsection{Canonical algebraic-geometric context}\label{canonical-algebraic-geometric-context-4}

The actual abelian scheme, its tangent bundle and intrinsic theta/determinant invertible modules live on Spec(O\_K), with the given full CM order action. Integral eigenmodules are common kernels over each DVR and commute with flat base change, including ramified extensions; they need not split the integral tangent lattice. The section is the actual generator of the polarized theta line and is transported with its rigidification, Picard data and Poincare biextension. The same algebraic frame polynomial is evaluated under every local and complex base change. The pi-Hermitian form and determinant volume retain the original complex realization and arithmetic degree. Formal groups/Dieudonne crystals, lambda-ring/K-theory operations and motivic realizations are not altered by this conditional transfer. No new l-adic weight theorem, motivic L-function identity or global arithmetic comparison is inferred. Exact carried contextual locators: SGA3 I4.7.3/4.7.3.1; SGA6 VI Theorem3.3; FGA232 Theorem3.1 pp148--149; FGA236 Corollaries1.7/2.3 pp229/232; SGA7 I VII2.7/2.9 pp169--174 and VIII7.4(b) pp310--311; cote4 scan003 sections8.1.1--8.1.2, scans005--006 sections8.4.1--8.4.3; cote54 scans005--007. Programme rows DP\_MULT Q3e/DP44/DP52/DP54--65, THETA3 G1/A7, CX\_GRH CX-ALG-024 and GCT\_PLETH T-052 retain their stated remaining gates.

\clearpage

\section{AM-LIFT-083 --- the actual signature restriction on a Jacobian isotypic fibre}\label{am-lift-083-the-actual-signature-restriction-on-a-jacobian-isotypic-fibre}

\textbf{Canonical checked integration --- 2026-11-23T01:52:37.614492+00:00.} Complete root and independent whole-proof/script/source-interface reviews, with separate79,768-assertion exact replays, passed. The frozen historical source remains unchanged. LIFT083 JS1--25 proves the actual pinned signature obstruction for every J-isotypic complex fibre, independently of CMObj/IsTen/LamObjGood. For any quadratic Host action into End0(J\^{}N)=M\_N(Q(zeta7)), d=7 gives Lie multiplicities(2a+b,a+2b), a+b=N, and every other Host quadratic field gives balanced(3N/2,3N/2). Actual DetCond requires(f\_E p,f\_E q). Hence d=7 needs p\textless=2q and q\textless=2p, while d!=7 needs p=q. The literal pinned profiles p\_j=k+j−1,q\_j=m−j are unbalanced at every stage and terminal q\_m=0; under full SizesAt, k\textgreater m\^{}5 f\_R\textgreater=2m forces p\_j\textgreater2q\_j at every stage. Thus no solely J-isotypic point can satisfy any pinned stage under those full premises, even with enlarged field/changed primes and the old good-point wrappers bypassed. An actual enlargement A\textasciitilde J\^{}N×B must have N\textless=f\_E min(p,q) and dim B\textgreater=f\_E max(p−2q,q−2p,0); the terminal fibre contains no J factor. These are necessary geometric dimensions, not height or DP44 costs. Other CM isogeny classes, their actual full point/action/model/level/section and same-point nonzero value remain open, together with selected arithmetic content, prescribed global period, Read/STiesR/specific032 and DP44. No finite charge, support-floor change or RH conclusion follows.

Author: ALP\_MACHINE \nolinkurl{/root/lift_accounting}, 23 November 2026 UTC.

\textbf{Status.} Author whole-proof, literal-source and complete-script check PASS; final exact diagnostic replay: 79,768 assertions. Final file fingerprints and completed UTC are recorded in \nolinkurl{LIFT083_AUTHOR_VERIFICATION.json}. Independent review is separate. No Lean build or main-problem/nonvanishing conclusion is claimed.

\subsection{1. Statement and exact scope}\label{statement-and-exact-scope}

Let

\[
 C:\ y^2=x^7-1,\qquad J=\operatorname{Jac}(C),\qquad
 K_7=\mathbf Q(\zeta_7),\qquad k_7=\mathbf Q(\sqrt{-7}).
 \tag{JS1}
\]

All abelian varieties and endomorphism algebras in the geometric argument are over \(\mathbf C\). Thus the conclusions apply after any characteristic-zero coefficient-field extension and to every variety geometrically isogenous to the variety under discussion. An isogeny here need not preserve a chosen polarization, integral lattice or action; the action is transported through its rational inverse.

Let \(E\) be any CM field, \([E:\mathbf Q]=2f_E\), carrying a literal geometric \nolinkurl{Host} and \nolinkurl{Datum} from the pinned source. Write \(d\) for its positive prime with \(d\equiv7\pmod {16}\), put \(k_d=\mathbf Q(\sqrt{-d})\), and let the actual datum multiplicities be \((p,q)\). Suppose a nonzero actual complex fibre \(A\) satisfies the source \nolinkurl{DetCond} for its \(\mathcal O_E\)-action. If \(A\sim J^N\), \(N\ge1\), then:

\[
\begin{array}{ll}
d=7:&\exists a,b\in\mathbf Z_{\ge0},\ a+b=N,
 \quad(f_Ep,f_Eq)=(2a+b,a+2b);\\[2mm]
d\ne7:&N\text{ is even},\quad
 (f_Ep,f_Eq)=(3N/2,3N/2).
\end{array}
\tag{JS2}
\]

These conclusions quantify \textbf{every} unital \(\mathbf Q\)-algebra embedding of the Host quadratic field into \(M_N(K_7)\). In particular they do not assume that the chosen \(E\)-action restricts to the original central \(K_7\)-action, or even that \(E\) contains \(K_7\) as an abstract subfield. The only compulsory subfield of \(E\) used in this argument is \(k_d\), generated by its literal \nolinkurl{Host.rt}.

In the first case, necessarily

\[
 p\le2q,\qquad q\le2p,\qquad
 a=\frac{f_E(2p-q)}3,\qquad b=\frac{f_E(2q-p)}3.
 \tag{JS3}
\]

The displayed quantities must be nonnegative integers. In the second case \(p=q\).

For a source record \(D\), any \nolinkurl{Par.Pins} pins each stage \(1\le j\le m\) to

\[
 p_j=k+j-1,\qquad q_j=m-j.
 \tag{JS4}
\]

Therefore a \(J\)-isotypic actual fibre is excluded at \textbf{every} pinned stage if \(d\ne7\), and for \(d=7\) it can occur only if

\[
                         3j\le2m-k+1.
 \tag{JS5}
\]

In particular it is excluded at the terminal stage \(j=m\) under \nolinkurl{Ranks} alone. If the record also has the literal \nolinkurl{SizesAt}, it is excluded at \textbf{every} pinned stage, for every allowed Host quadratic field. This is a constraint on the proposed \(J\)-isotypic replacement point even after the old \nolinkurl{CMObj}, \nolinkurl{LamObjGood} and \nolinkurl{LamPt} wrappers are set aside.

There is also a necessary condition on an enlargement. Suppose an actual fibre decomposes up to geometric isogeny as

\[
                        A\sim J^N\times B,
 \tag{JS6}
\]

where \(B\) has no simple factor geometrically isogenous to \(J\), and \(N\ge0\). Then

\[
 N\le f_E\min(p,q),\qquad
 \dim B\ge f_E\max\{p-2q,q-2p,0\}.
 \tag{JS7}
\]

If \(d\ne7\), the stronger bounds are

\[
 N\le\frac{2f_E\min(p,q)}3,\qquad
 \dim B\ge f_E|p-q|.
 \tag{JS8}
\]

At a pinned stage \(p_j>q_j\). Thus \(N\le f_Eq_j\), and the terminal fibre \(q_m=0\) contains \textbf{no} \(J\) factor at all. These are dimension bounds; they are neither height lower bounds nor lower bounds on a DP44 cost. They leave non-\(J\)-isotypic constructions open.

\subsection{2. Literal source clauses and the field conventions}\label{literal-source-clauses-and-the-field-conventions}

The source is pinned at commit \nolinkurl{7cbf7a0c32fca1badf05931879a37c3f8a2b03b4}. The companion source manifest records exact file hashes, anchors and line locations. The following bodies were read directly.

\begin{longtable}[]{@{}
  >{\raggedright\arraybackslash}p{(\columnwidth - 2\tabcolsep) * \real{0.5000}}
  >{\raggedright\arraybackslash}p{(\columnwidth - 2\tabcolsep) * \real{0.5000}}@{}}
\toprule\noalign{}
\begin{minipage}[b]{\linewidth}\raggedright
Source
\end{minipage} & \begin{minipage}[b]{\linewidth}\raggedright
Clause used
\end{minipage} \\
\midrule\noalign{}
\endhead
\bottomrule\noalign{}
\endlastfoot
\nolinkurl{ReciprocityTheory/ReductiveGroups/ArithmeticGeometry/BirationalGeometry/ToroidalCompactification/TwoWeightEqOneAddRoot.lean}, 14--26 & \nolinkurl{Host.d_prime}, \nolinkurl{d_mod}, \nolinkurl{rt_sq} and \nolinkurl{sqrt_neg_disc_sq}; \nolinkurl{rt} generates the actual imaginary quadratic field. \\
Same file, 30--49 & \nolinkurl{Datum.p}, \nolinkurl{q}, rank and \nolinkurl{dom}; the trace expression is \(\sum_{\tau(rt)=rt_0}(p\tau(a)+q\bar\tau(a))\). \\
Same file, 113--120 & \nolinkurl{LieMult} is a complete basis of the indicated eigenmodule, not only a list of independent eigenvectors. \nolinkurl{DetCond} imposes multiplicities \(p\) and \(q\) according to the actual root of \nolinkurl{rt}. \\
Same file, 144--160 & \nolinkurl{Obj.geom} imposes this \nolinkurl{DetCond} on every algebraically closed geometric fibre, with the actual base root. \\
\nolinkurl{Formalism/LocalToGlobal/ProjectiveLattices/PartitionExtendByZeroAntitone.lean}, 12--26, 39--49 & \nolinkurl{InR}, \nolinkurl{IsDatW}, \nolinkurl{Par.PinHost}, \nolinkurl{PinDat} and \nolinkurl{Pins}; \nolinkurl{IsDatW} pins the two individual multiplicities to \nolinkurl{Ranks.sigW}. \\
\nolinkurl{Formalism/LocalRingCardinality/Params.lean}, 6--22, 28--46 & \(m\ge2\), \(k\ge m+3\), \(f_R\ge m\), \(f_R=2\,\mathrm{ord}\), and \nolinkurl{sigW j = (k+j-1,m-j)}. \\
\nolinkurl{Formalism/LocalRingCardinality/DataParams₂.lean}, 9--27 & \nolinkurl{Scale.x > 1}, \nolinkurl{eps > 0}, and \(L=x^\epsilon\). \\
\nolinkurl{Formalism/SharpRestriction/ElevenLe.lean}, 11, 35--36, 134--138 & \nolinkurl{szTh = eps/100}; the first \nolinkurl{SizeKRel} conjunct and its inclusion in \nolinkurl{SizesAt}. \\
\end{longtable}

We use \(f_E=[E:\mathbf Q]/2\) for the geometric embedding count and \(f_R\) for \nolinkurl{Ranks.f} when deriving the numerical size inequality. No equality between them is \textbf{needed} for the exclusion below: the geometric inequalities cancel \(f_E>0\), while the size argument uses only \(f_R\ge m\). This avoids inserting an unproved field-degree identification into the argument.

There are three distinct uses of a field action. The source gives an \(\mathcal O_E\)-action. After tensoring with \(\mathbf Q\), it gives an \(E\)-action on rational homology and Lie. Restriction to \(k_d\subset E\) then gives the quadratic action studied in section 4. Independently, every endomorphism of \(J^N\) is a matrix over its central endomorphism field \(K_7\). It is this \textbf{central \(K_7\)-linearity of matrices}, not an assumed linearity over the maximal real subfield of \(E\), that is used to classify the quadratic operator. No compatibility of the two embeddings of \(K_7\) is assumed.

\subsection{3. The geometric endomorphism algebra of the actual Jacobian}\label{the-geometric-endomorphism-algebra-of-the-actual-jacobian}

The curve in (JS1) has genus three. Its holomorphic differentials may be taken as

\[
                  \omega_i=x^{i-1}\,dx/y,\qquad i=1,2,3.
 \tag{JS9}
\]

The curve automorphism \((x,y)\mapsto(\zeta_7x,y)\) induces the eigenvalue \(\zeta_7^i\) on \(\omega_i\). It therefore gives the faithful field action \(K_7\hookrightarrow\operatorname{End}^0(J)\), with type \(\Phi=\{1,2,3\}\). The only proper CM subfield of the cyclic degree-six field \(K_7\) is \(k_7\). Its induced types are the quadratic-residue set \(\{1,2,4\}\) and its complement \(\{3,5,6\}\); \(\Phi\) is neither. The primitive-type conclusion and geometric simplicity were already established in R003 and used in LIFT079. Milne (2020, Proposition 3.13, p.~30) supplies the corresponding isogeny classification. Together with Proposition 3.3 and Remark 3.5, pp.~27--28, this gives

\[
                   \operatorname{End}^0_{\mathbf C}(J)=K_7.
 \tag{JS10}
\]

For completeness, the equality can also be checked directly without using simplicity as a shortcut. The faithful \(K_7\)-action makes \(H_1(J,\mathbf Q)\) a one-dimensional \(K_7\)-space. Every \(\mathbf Q\)-linear endomorphism of this space is uniquely

\[
                  T=\sum_{g\in\operatorname{Gal}(K_7/\mathbf Q)}a_g g,
                  \qquad a_g\in K_7.
 \tag{JS11}
\]

Indeed, evaluation of the six automorphisms on \(1,\zeta_7,\ldots,\zeta_7^5\) gives a Vandermonde matrix on the six distinct primitive seventh roots, so these maps form a \(K_7\)-basis of \(\operatorname{End}_{\mathbf Q}(K_7)\). On complexification, a nonzero term \(a_gg\) has a nonzero entry on every position of its embedding-permutation graph. Different \(g\) have disjoint such positions in each row, so the entries cannot cancel across \(g\). A rational endomorphism of the abelian variety preserves the Hodge decomposition. Hence a nonzero \(a_g\) requires \(g\Phi=\Phi\) (equivalently preservation of its conjugate type). The stabilizer of \(\{1,2,3\}\) in \((\mathbf Z/7\mathbf Z)^\times\) is the identity. Thus \(T\) is multiplication by an element of \(K_7\). Since that field is already present, (JS10) follows.

Consequently

\[
                 \operatorname{End}^0_{\mathbf C}(J^N)=M_N(K_7),
                 \qquad
                 \operatorname{Lie}(J^N)=
                    \bigoplus_{\tau\in\Phi}\mathbf C^N_{\tau}.
 \tag{JS12}
\]

A matrix \(M\in M_N(K_7)\) acts on the \(\tau\)-summand by \(\tau(M)\). The same description, up to conjugacy by the differential of a rational isogeny, applies to any \(J\)-isotypic fibre. Conjugating the action, changing an integral model, or enlarging the coefficient field does not change these Lie multiplicities. The standard matrix-algebra statement and invariance of isotypic summands also follow from Poincaré reducibility as stated by Milne (2020, Theorem 2.12 and its following paragraph, pp.~26--27).

\subsection{4. Every quadratic embedding into the matrix algebra}\label{every-quadratic-embedding-into-the-matrix-algebra}

Let \(\iota:k_d\hookrightarrow M_N(K_7)\) be an arbitrary unital \(\mathbf Q\)-algebra embedding, and set \(R=\iota(\sqrt{-d})\). Then

\[
                              R^2=-dI_N.
 \tag{JS13}
\]

Each matrix entry is in \(K_7\), so \(R\) is a \(K_7\)-linear operator on \(K_7^N\). This assertion does not require \(\iota\) to be a \(K_7\)-algebra map; such a requirement would not even make sense when \(k_d\not\supset K_7\).

\subsubsection{\texorpdfstring{4.1. The split case \(d=7\)}{4.1. The split case d=7}}\label{the-split-case-d7}

Put

\[
 \alpha=\zeta_7+\zeta_7^2+\zeta_7^4
          -\zeta_7^3-\zeta_7^5-\zeta_7^6.
 \tag{JS14}
\]

The exact identities \(\alpha^2=-7\) and \(\tau_u(\alpha)=(u/7)\alpha\) follow by multiplication in \(\mathbf Q[X]/\Phi_7(X)\). Because \(X^2+7=(X-\alpha)(X+\alpha)\) has distinct roots in \(K_7\), \(R\) is diagonalizable over \(K_7\), with eigenvalue multiplicities \(a,b\), \(a+b=N\). Equivalently its two projectors are \((I\pm R/\alpha)/2\), whose \(K_7\)-ranks are \(a,b\).

Among \(u=1,2,3\), two are quadratic residues and one is a nonresidue modulo seven. Applying the three actual type embeddings in (JS12), the complex Lie multiplicities of the roots \(\alpha,-\alpha\) are therefore

\[
                             (2a+b,a+2b).
 \tag{JS15}
\]

If the Host's chosen complex root is the opposite root, interchange \(a\) and \(b\). Thus the assertion is independent of that orientation and still has the existential form in (JS2). Every pair \(a,b\ge0\) does occur for a quadratic embedding into this matrix algebra, by using the diagonal matrix with \(a\) copies of \(\alpha\) and \(b\) copies of \(-\alpha\). This last assertion concerns only the rational quadratic representation; it constructs no full integral \(E\)-action or polarized moduli point.

\subsubsection{\texorpdfstring{4.2. The nonsplit case \(d\ne7\)}{4.2. The nonsplit case d\textbackslash ne7}}\label{the-nonsplit-case-dne7}

The Galois group of \(K_7/\mathbf Q\) is cyclic of order six, and it has a unique index-two subgroup. Thus \(k_7\) is its unique quadratic subfield. Since \(d\) is prime, \(k_d\ne k_7\) when \(d\ne7\); hence \(X^2+d\) is irreducible over \(K_7\). The operator \(R\) makes \(K_7^N\) a vector space over the degree-two field \(K_7[X]/(X^2+d)\). Consequently \(N\) is even, and its characteristic polynomial over \(K_7\) is

\[
                              (X^2+d)^{N/2}.
 \tag{JS16}
\]

After any of the three embeddings \(\tau\in\Phi\), each of the two complex roots has multiplicity \(N/2\). Their total Lie multiplicities are \((3N/2,3N/2)\). This proves the second case of (JS2), and it works equally for any imaginary quadratic field not isomorphic to \(k_7\), without a prime-discriminant assumption. Conversely the companion block \(\left(\begin{smallmatrix}0&-d\\1&0\end{smallmatrix}\right)\), repeated \(N/2\) times, realizes this rational quadratic representation when \(N\) is even. Again no source point is supplied by that converse.

\subsection{5. Passing from the actual determinant condition to quadratic multiplicities}\label{passing-from-the-actual-determinant-condition-to-quadratic-multiplicities}

Let \(A\) be an actual complex fibre of the source \nolinkurl{Obj}. Its unital \(\mathcal O_E\)-action extends to a unital \(E\)-action on rational homology: multiplication by a nonzero integer is invertible in \(\operatorname{End}^0(A)\). The resulting map from the field \(E\) is injective when \(A\ne0\), since its kernel is a proper ideal. Its complex Lie representation is a representation of

\[
                          E\otimes_{\mathbf Q}\mathbf C
                          \cong\prod_{\tau:E\hookrightarrow\mathbf C}\mathbf C.
 \tag{JS17}
\]

In particular it decomposes as the direct sum of the simultaneous eigenspaces indexed by all embeddings \(\tau\). This is ordinary semisimplicity in characteristic zero; it does not assert that integral eigenspaces split at ramified finite primes.

Each of the two embeddings of \(k_d\) has exactly \([E:k_d]=f_E\) extensions to \(E\). The source \nolinkurl{LieMult} gives the full dimensions of those eigenspaces, and \nolinkurl{Obj.geom} imposes \nolinkurl{DetCond} for the same physical root. Hence the two quadratic-root multiplicities are exactly

\[
                               (f_Ep,f_Eq),
                               \qquad \dim A=f_E(p+q).
 \tag{JS18}
\]

If \(A\sim J^N\), transport this \(E\)-action by a geometric rational isogeny and restrict it to \(k_d\). The isogeny's differential is an invertible complex-linear map, so it preserves the multiplicities (JS18). Section 4 applies to \textbf{that} quadratic embedding, without identifying its image with a preselected central subfield. Comparing (JS15) or (JS16) with (JS18) proves (JS2) and (JS3).

The geometric argument uses only the actual fibre and the source determinant condition. It has not used \nolinkurl{CMObj}, \nolinkurl{IsTen}, \nolinkurl{LamObjGood}, \nolinkurl{LamPtClauses}, a choice of eigenletters, a theta-section normalization or any metric estimate.

\subsection{6. The pinned profiles, including the full size predicate}\label{the-pinned-profiles-including-the-full-size-predicate}

By \nolinkurl{Par.PinDat}, \nolinkurl{IsDatW} and the literal definition of \nolinkurl{Ranks.sigW}, (JS4) holds for every \(j\) with \(1\le j\le m\). The natural-number subtractions are ordinary integer subtractions in this range. The \nolinkurl{Ranks} inequality \(k\ge m+3\) gives

\[
                  p_j-q_j=k-m+2j-1\ge2j+2>0.
 \tag{JS19}
\]

Thus the nonsplit case in section 4 is impossible at every pinned stage. In the split case \(p_j\le2q_j\) is exactly (JS5). No sufficiency for an actual \(E\)-action or moduli point is claimed when (JS5) happens to hold. At \(j=m\), \(q_j=0\) and \(p_j=k+m-1>0\), contradicting (JS3).

Now assume the literal \nolinkurl{SizesAt L}. Its first \nolinkurl{SizeKRel} conjunct, \nolinkurl{Scale.L = x^eps}, and \nolinkurl{szTh = eps/100} give

\[
  k\ge x^{\epsilon/100}m^5 f_R Lx^\epsilon
    =m^5f_R\,x^{201\epsilon/100}
    >m^5f_R.
 \tag{JS20}
\]

Strictness uses exactly \(x>1\) and \(\epsilon>0\). Since \(f_R\ge m\ge2\), one has \(m^5f_R\ge m^6\ge2m\). Therefore

\[
                    p_j\ge k>2m>2(m-j)=2q_j
                    \quad(1\le j\le m).
 \tag{JS21}
\]

This excludes every \(J\)-isotypic actual fibre at every pinned stage, including any enlarged-degree Host and any changed primes, as long as the same source signature and full \nolinkurl{SizesAt} premises are retained. The conclusion is pointwise at a record satisfying \nolinkurl{SizesAt}; it makes no assertion at an ineligible family parameter where that predicate has not been supplied. The previously checked fixed-degree exclusions in LIFT054 are not used to make this argument vacuous.

\subsection{7. What an enlargement must supply}\label{what-an-enlargement-must-supply}

Consider (JS6). Poincaré reducibility identifies \(J^N\) as the entire \(J\)-isotypic summand. There are no nonzero rational homomorphisms between it and \(B\), since their simple isogeny factors are disjoint. Therefore

\[
 \operatorname{End}^0(A)\cong M_N(K_7)\times\operatorname{End}^0(B),
 \tag{JS22}
\]

with the first factor omitted if \(N=0\). The corresponding isotypic idempotent is central, so \textbf{every} actual \(E\)-action preserves this decomposition. Its restriction to a nonzero factor remains unital and thus gives a faithful field action. No assumption that a chosen product isogeny was initially \(E\)-linear is needed; conjugating the action by that isogeny produces the stated block action.

Let \(u_+,u_-\) be the two quadratic multiplicities on \(\operatorname{Lie}(B)\), both nonnegative. If \(d=7\), then for some \(a+b=N\)

\[
        f_Ep=2a+b+u_+,
        \qquad f_Eq=a+2b+u_-.
 \tag{JS23}
\]

Both \(2a+b\) and \(a+2b\) are at least \(N\). It follows that \(N\le f_E\min(p,q)\). Also

\[
 f_E(p-2q)=-3b+u_+-2u_-\le u_+\le\dim B,
 \tag{JS24}
\]

and its sign-reversed analogue proves (JS7). For \(d\ne7\), the two contributions from \(J^N\) equal \(3N/2\); subtracting gives \(f_E(p-q)=u_+-u_-\), which proves (JS8). These arguments include \(N=0\) trivially.

For a pinned stage, (JS19) and (JS7) give

\[
 \dim B\ge f_E(k+3j-1-2m),\qquad
 \frac{\dim J^N}{\dim A}\le\frac{3(m-j)}{k+m-1},
 \tag{JS25}
\]

where the first inequality is useful only when its right side is positive. Under full \nolinkurl{SizesAt}, it is positive at every stage. At the terminal stage \(j=m\), (JS7) forces \(N=0\). These are necessary geometric restrictions on an actual enlarged source. They do not construct \(B\), its polarization, full integral action, local models, level, coefficient-field ties, section, or a nonzero value of the prescribed form at the same point.

\subsection{8. Relation to the earlier source boundary and verification scope}\label{relation-to-the-earlier-source-boundary-and-verification-scope}

R003/DP25 established the geometric simplicity of \(J\) and the incompatibility of a \(J\)-power with an unchanged elliptic-isotypic \nolinkurl{IsTen} construction. LIFT079 subsequently pinned the complete pure-type and onto clauses of \nolinkurl{CMObj}/\nolinkurl{LamObjGood}, and the fixed-fibre equality in \nolinkurl{LamPt}. Those statements remain unchanged.

The new conclusion here is (JS2) applied to the \textbf{actual \nolinkurl{IsDatW}/\nolinkurl{DetCond} profile}, even when those old good-point constructors are bypassed. It rules out making the required point solely from the \(J\)-isotypic blocks whose line metrics were bounded in LIFT081. LIFT080 and LIFT082 remain valid conditional transfer theorems. This result identifies a missing point/signature premise in a proposed application of them; it does not alter their equations or claim that a different CM isogeny class cannot satisfy their hypotheses.

The companion certificate checks exact cyclotomic multiplication, the type stabilizer and both quadratic restriction counts, finite matrix models for the two cases, integer profile inequalities and enlargement bounds, and the exact source fingerprints/anchors. Its finite models are diagnostics for the displayed arguments. They do not replace the general proof, construct an admissible point, verify a Lean development, evaluate a prescribed section or period, or establish any sign of the final DP44 comparison.

\textbf{Primary reference.} Milne, J. S. (2020, July 14). \emph{Complex multiplication} (Version 0.10), Theorem 2.12 and following paragraph (pp.~26--27), Proposition 3.3 and Remark 3.5 (pp.~27--28), and Propositions 3.12--3.13 (p.~30). \url{https://www.jmilne.org/math/CourseNotes/CM.pdf}. These passages were directly read in the author's PDF during this check. The field and matrix arguments in sections 3--7 are supplied explicitly above.

\subsection{Canonical algebraic-geometric context}\label{canonical-algebraic-geometric-context-5}

The obstruction is a statement about the characteristic-zero geometric fibre of the actual abelian scheme and its full CM action: the rational isotypic projector of h\^{}1(J) is central, and the Hodge/Lie multiplicities of every quadratic action constrain the determinant condition. It persists under geometric isogeny and coefficient-field extension but does not assert any integral eigenspace splitting or identify a new polarized lattice. The complementary weight-one factor and its Picard/polarization/Poincare-biextension data still require construction. Formal groups and Dieudonne crystals of the J factors retain the accepted DP slopes; the dimension inequality supplies no stronger divided-power gain, lambda-ring or K-theory identity, l-adic weight bound or motivic L-function identity. Exact inherited contextual locators: SGA3 I4.7.3/4.7.3.1; SGA6 VI Theorem3.3; FGA232 Theorem3.1 pp148--149; FGA236 Corollaries1.7/2.3 pp229/232; SGA7 I VII2.7/2.9 pp169--174 and VIII7.4(b) pp310--311; cote4 scan003 sections8.1.1--8.1.2 and scans005--006 sections8.4.1--8.4.3; cote54 scans005--007. DP\_MULT Q3e/DP44/DP52/DP54--65 and THETA3 G1/A7 now need a point in a sufficiently enlarged isogeny class or a changed signature interface; CX\_GRH CX-ALG-024 and GCT\_PLETH T-052 retain their separate open gates.

\clearpage

\section{AM-K099 --- total eligibility and Galois restrictions on the actual middle scalar}\label{am-k099-total-eligibility-and-galois-restrictions-on-the-actual-middle-scalar}

\textbf{Canonical checked integration --- 2026-11-23T02:02:43.558518+00:00.} Complete saved-proof and executable reviews and separate exact diagnostic replays passed. No numerical error charge is introduced. Frozen author history is preserved. AM-K099 GE1--13 restricts the actual remaining fully eligible support25 branch to total eligible count q=25 or26. The norm bounds give x\textgreater2\^{}120 for q25, x\textgreater2\^{}1116 for q26, and x\textgreater2\^{}2111 for q27, excluding q27 under x\textless2\^{}1820. Over K0=Q(sqrt(-7)), the relative Galois group of order168 acts regularly on the selected T primes. At q25 the scalar has trivial relative stabilizer and generates E over K0; at q26 the relative stabilizer has order at most2. The absolute degrees are at least168 and84 respectively. Distinct actual norm-U scalars in the bounded denominator window have eligible-set overlap at most14; set and scalar stabilizers agree. A relative stabilizer of order2 forces the denominator ideal to extend from the fixed field and x to be a square, even at ramified primes. This gives necessary structural restrictions only, no numerical sector removal, new error charge, source construction, signed positivity, period, DP44 or RH claim. The existing97/31 error arrays and complete floors25/11 remain unchanged.

Author: \nolinkurl{/root/kummer_extension}. 23 November 2026 UTC.

Controlled pending finding. The theorem is a structural necessary condition, with no numerical error charge or support-floor change. Author checks, independent reviews and canonical integration are separate events.

\subsection{1. Actual scalar and the total eligible set}\label{actual-scalar-and-the-total-eligible-set}

Keep the unchanged J56 field, original selected prime components and actual all-top support25 kernel. In the globally rank-two-unipotent scalar-cross remainder of DO105(EL1--3), the normalized image line is

\[
 C=(1,\eta)^t r,\qquad \eta\bar\eta=U,\qquad
 U=77Q^2-14,\quad Q=1427908193729121522601,
 \quad 2^{120}<x:=N_E\mathfrak D_\eta<2^{1820}.             \tag{GE1}
\]

Here \(\mathfrak D_\eta\) is the exact integral denominator ideal. Every one of the25 active relative T places is eligible. No scalar is assumed integral, and no new rational or moving integral frame is introduced.

Write \(j=\sqrt{-7}\), \(K_0=\mathbb Q(j)\), and

\[
 a_0=Q(-7+2j),\qquad a_0\bar a_0=U+14.                    \tag{GE2}
\]

Let \(\mathcal P\) be \textbf{all168} degree-one E primes above the original selected prime \((\pi_T)\) of \(K_0\), not only the25 active places. Define

\[
 \mathcal S(\eta)=\{\mathfrak p\in\mathcal P:
                  v_{\mathfrak p}(\eta-a_0)\ge1\},
 \qquad q=|\mathcal S(\eta)|.                              \tag{GE3}
\]

The actual active support is a subset of \(\mathcal S(\eta)\), hence q≥25. The comparison in GE3 is only at the selected component. At an eligible place the norm equation forces the partner value \(\eta\equiv U/\bar a_0\); it does \textbf{not} give a second divisibility of \(\eta-a_0\).

\textbf{Theorem.} Under GE1, the total eligibility count is q=25 or26. If q=26, then x\textgreater{}\(2^{1116}\). Let \(G_0=\operatorname{Gal}(E/K_0)\) and \(H=\operatorname{Stab}_{G_0}(\eta)\). Then

\begin{longtable}[]{@{}
  >{\raggedright\arraybackslash}p{(\columnwidth - 4\tabcolsep) * \real{0.3333}}
  >{\raggedright\arraybackslash}p{(\columnwidth - 4\tabcolsep) * \real{0.3333}}
  >{\raggedright\arraybackslash}p{(\columnwidth - 4\tabcolsep) * \real{0.3333}}@{}}
\toprule\noalign{}
\begin{minipage}[b]{\linewidth}\raggedright
Total eligibility
\end{minipage} & \begin{minipage}[b]{\linewidth}\raggedright
Relative stabilizer
\end{minipage} & \begin{minipage}[b]{\linewidth}\raggedright
Field consequence
\end{minipage} \\
\midrule\noalign{}
\endhead
\bottomrule\noalign{}
\endlastfoot
q=25 & H is trivial & \(K_0(\eta)=E\), so \([\mathbb Q(\eta):\mathbb Q]\ge168\) \\
q=26 & \(|H|=1\) or2 & \([E:K_0(\eta)]\le2\), so \([\mathbb Q(\eta):\mathbb Q]\ge84\) \\
\end{longtable}

The stabilizer of the set \(\mathcal S(\eta)\) in \(G_0\) is exactly H. Distinct \(G_0\)-conjugates have eligible-set overlap at most14. If H has order2, \(\mathfrak D_\eta\) extends from \(E^H\), and x is the square of an integer ideal norm. None of these assertions proves that an admissible eta exists or that the remaining kernel contribution is nonzero.

\subsection{2. The exact norm thresholds for25,26 and27 eligible places}\label{the-exact-norm-thresholds-for2526-and27-eligible-places}

The scalar \(\delta=\eta-a_0\) is nonzero, since its vanishing would contradict GE1--2. The ideal \((\delta)\mathfrak D_\eta\) is integral: valuationwise, \(v(\delta)\ge\min(v(\eta),0)\). DO103(LD13--15), using the actual four E primes over7, gives the additional norm factor \(7^{84}\), also when eta has denominators at7. Explicitly, each of the two conjugate prime pairs has ramification index6, residue degree14, and the paired denominator-cleared valuation sum is at least3. These are the carried actual-field statements; the same ideal is being used here.

At each prime in the \textbf{total} set GE3, delta is divisible by that degree-one prime. The same norm argument therefore gives

\[
 7^{84}T^q\le N_E\bigl((\delta)\mathfrak D_\eta\bigr)
       <[4(U+14)]^{168}x.                                  \tag{GE4}
\]

The strict upper bound uses \(|\eta_\tau|=\sqrt U\) and \(|(a_0)_\tau|=\sqrt{U+14}\) at every complex embedding, hence \(|\delta_\tau|<2\sqrt{U+14}\). The absolute E norm is the product of the168 squared complex magnitudes. The number of active arrows plays no role in this scalar norm argument.

Put \(B=[4(U+14)]^{168}\). The unchanged actual integer T satisfies the following strict comparisons, verified by integer powers:

\[
\begin{array}{rcl}
 2^{120}B&<&7^{84}T^{25}<2^{121}B,\\
 2^{1116}B&<&7^{84}T^{26}<2^{1117}B,\\
 2^{2111}B&<&7^{84}T^{27}<2^{2112}B.
\end{array}                                                \tag{GE5}
\]

If q≥27, GE4--5 force x\textgreater{}\(2^{2111}\), contradicting GE1. If q=26 they force x\textgreater{}\(2^{1116}\). Thus q=25 or26, and in particular every middle scalar with x≤\(2^{1116}\) has q=25. The latter is a logical restriction, not a count of those scalars.

\subsection{3. The actual relative Galois action is free on the selected primes}\label{the-actual-relative-galois-action-is-free-on-the-selected-primes}

The actual field of K032(QH9--15) is

\[
 E=\mathbb Q(\zeta_{28})B_7(29)\mathbb Q(\sqrt5,\sqrt{13}),
 \quad [E:\mathbb Q]=336,
 \quad \operatorname{Gal}(E/\mathbb Q)
       \simeq C_{21}\times C_2^4.                           \tag{GE6}
\]

The quadratic subfield \(K_0\) is the quadratic-residue character of the modulo7 factor. Consequently

\[
 |G_0|=168,\qquad G_0\simeq C_{21}\times C_2^3.             \tag{GE7}
\]

The actual large-T certificate proves T prime, T≡1 modulo52780, and the coprime factorization \(\pi_T\bar\pi_T=T\). Thus T splits completely in E, and \(\mathcal P\) is exactly the168 primes over the chosen \(K_0\) prime. A decomposition group for E/K0 is trivial. The action of \(G_0\) on \(\mathcal P\) is therefore free and transitive, not merely an abstract permutation action of the correct cardinality.

Every \(\sigma\in G_0\) fixes a0, so valuations give

\[
                 \mathcal S(\sigma\eta)=\sigma\mathcal S(\eta).
                                                               \tag{GE8}
\]

Since H fixes eta, it preserves \(\mathcal S(\eta)\). Its action on this set is free, by the preceding prime action; hence

\[
                   |H|\mid q,\qquad |H|\mid168.             \tag{GE9}
\]

For q25 their gcd is1, and for q26 it is2. Galois correspondence now gives the field assertions in Section1. Adjoining K0 to \(\mathbb Q(\eta)\) has degree at most2, which explains the absolute-degree lower bounds168 and84; the relative and absolute degrees are not conflated.

These are arithmetic consequences of a single eta. The proof does not apply a Galois automorphism to the full kernel or assume that selected real types, projectors, class representatives, phases or weights are Galois-invariant.

\subsection{4. Paired overlap, set stabilizers and denominator descent}\label{paired-overlap-set-stabilizers-and-denominator-descent}

Let eta1 and eta2 be distinct elements of E satisfying the same norm equation, each with denominator norm less than\(2^{1820}\). Suppose t primes of \(\mathcal P\) are eligible for both. At each such prime they are units and both reduce to a0. Their norm equations give the same partner residue \(U/\bar a_0\). Thus

\[
 (\eta_1-\eta_2)\mathfrak D_{\eta_1}\mathfrak D_{\eta_2}
 \quad\text{is integral and has norm at least }T^{2t}.
                                                               \tag{GE10}
\]

The product of denominator ideals clears the difference, without requiring either ideal principal or equal to the other. The archimedean bound gives

\[
 T^{2t}\le(4U)^{168}x_1x_2<(4U)^{168}2^{3640}<T^{30}.       \tag{GE11}
\]

The last comparison is exact. Therefore t≤14. For \(\sigma\in G_0\), conjugation preserves x and the norm equation. If \(\sigma\eta\ne\eta\), GE8 and GE11 imply

\[
                    |\mathcal S(\eta)\cap\sigma\mathcal S(\eta)|
                        \le14.                             \tag{GE12}
\]

Conversely, any \(\sigma\) preserving \(\mathcal S(\eta)\) gives overlap q≥25, so it must fix eta. This proves equality of the scalar and set stabilizers. In the q26, \textbar H\textbar2 case the eligible set consists of13 free H-orbits; the set orbit under \(G_0\) has size84. In the trivial-stabilizer case its orbit has size168.

For the denominator statement, let L=E\^{}H and eta∈L. For every E prime P above an L prime p, normalized valuations satisfy \(v_P(\eta)=e(P/p)v_p(\eta)\). Taking negative parts commutes with multiplication by the positive integer e. Hence

\[
 \mathfrak D_{\eta,E}=\mathfrak D_{\eta,L}\mathcal O_E,
 \qquad N_E\mathfrak D_{\eta,E}
          =(N_L\mathfrak D_{\eta,L})^{[E:L]}.               \tag{GE13}
\]

This argument allows ramification; no unramified descent is assumed. For \textbar H\textbar2 it gives the claimed integer square. More generally the absolute stabilizer of eta has order1,2 or4, since its intersection with the index-two subgroup G0 has order at most2. Its denominator norm is correspondingly a first, second or fourth power of an ideal norm in \(\mathbb Q(\eta)\). This extra restriction alone does not provide a density estimate or a smaller absolute contribution.

\subsection{5. Verification and exact scope}\label{verification-and-exact-scope}

The companion reads the actual host, checks every strict comparison GE5 and GE11 by integers, checks the actual subgroup of the modulo28 factor and its regular168-point action, and verifies the resulting stabilizer divisibilities and denominator negative-part identities. These finite diagnostics accompany the general ideal/Galois proofs; they do not enumerate actual eta solutions.

Direct carried inputs are K032(QH9--15, the group statement following QH21), K030(PH2), DO040.6--7 and its actual J56 prime certificate, the packet-plane norm-boundary §§1--4, DO103(LD4--15), and DO105(EL1--8). The certificate records their current file hashes. The exact same field, CM involution, fractional ideals and affine kernel data are retained. No new geometric comparison or source constructor is used.

There is \textbf{no new error charge}. Both complete support floors and both chronological error arrays remain unchanged. The q26 lower bound is available as an input to a separately proved count; this structural theorem alone removes no signed-kernel sector. In particular, high degree or small Galois overlap does not by itself contradict the existence of an eligible eta. The full selected arithmetic image, prescribed period, source constructor, signed sufficient inequality, DP44 and RH conclusions remain open.

\subsection{Canonical algebraic-geometric context}\label{canonical-algebraic-geometric-context-6}

The scalar and its denominator define an actual fractional ideal on Spec(O\_E); denominator descent is valuationwise and includes ramification. The Galois argument concerns this arithmetic datum only and asserts no symmetry of the selected kernel. The original fractional ideals and affine fibres are invertible-module conditions on Spec(O\_E)/Spec(O\_F), with the same CM involution, Hermitian polarization, theta line, biextension and complex metric. Exact inherited contextual anchors remain SGA3 I4.7.3/4.7.3.1; SGA6 VI Theorem3.3; FGA232 Theorem3.1 pp148--149; FGA236 Corollaries1.7/2.3 pp229/232; SGA7 I VII2.7/2.9 pp169--174 and VIII7.4(b) pp310--311; cote4 scan003 sections8.1.1--8.1.2 and scans005--006 sections8.4.1--8.4.3; cote54 scans005--007. No new lambda-ring/K-theory equality, l-adic or motivic comparison is inferred. DP\_MULT Q3e/DP44/DP52/DP54--65, THETA3 G1/A7, CX\_GRH CX-ALG-024 and GCT\_PLETH T-052 retain their actual arithmetic and global comparison gates.

\clearpage

\section{AM-K100 --- a larger-denominator middle-sector count}\label{am-k100-a-larger-denominator-middle-sector-count}

\textbf{Canonical checked integration --- 2026-11-23T02:02:45.578639+00:00.} Complete root and independent proof/script/interface reviews and separate exact replays passed. The actual canonical predecessor has97/31 entries; this integration gives98/31, complete floors25/11. Frozen author history is preserved. AM-K100 QE1--16 removes the complete actual fully eligible all-II globally rank-two-unipotent scalar-cross support25 subset with2\^{}1116\textless=x=N\_E D\_eta\textless2\^{}1820 and22\textless=k\textless=48, including either total eligible count25 or26. The exact1116 endpoint, cap1117, complete joint row/transverse beta count, central volume\textless2\^{}679 and unique-eta/label weight give worst k48 ceiling-1730011/4096\textless-422; the26 smaller cells and this cell form one complete2\^{}-421 union. All k\textless=21 contributions remain in DO105451 and are not added again here. Append421 to the actual DO10597/31 ordered totals, giving98/31; both entire sums remain below2\^{}-230 with complete floors25/11 unchanged. In particular all q26 cells up through k48 are removed. The surviving support25 possibilities have2\^{}120\textless x\textless2\^{}1116 with22\textless=k\textless=164, or2\^{}1116\textless=x\textless2\^{}1820 with49\textless=k\textless=164, in addition to K099 structural restrictions. No claim for these surviving cells, no actual source, selected arithmetic image, prescribed period, signed positivity, specific032, DP44 or RH conclusion follows.

Author: \nolinkurl{/root/kummer_extension}. 23 November 2026 UTC.

Controlled pending finding. This is a separate quantitative consequence of the complete actual DO105 count. It retains K099 as a structural input, and adds one explicitly charged subset. Frozen earlier proofs and their errors are unchanged.

\subsection{1. Actual sector and the new subset}\label{actual-sector-and-the-new-subset}

Keep exactly DO105(EL1--3): the J56 all-top support25 scalar-cross kernel, H=133, f=168, global rank-two unipotence, h\_ZZ=I, and all25 active places of actual type II. Preserve the original CM partners, paired pencils and omitted labels, finite projectors, selected real trace, central phase, class masses and Haar measures. In the original normalized line and covector, write

\[
 C=(1,\eta)^t r,\qquad\eta\bar\eta=U,
 \qquad x=N_E\mathfrak D_\eta,
 \qquad k=\#\{j\in S_{164}:r_j\ne0\}.                    \tag{QE1}
\]

All25 active places are eligible. The existing middle window is \(2^{120}<x<2^{1820}\); DO105 already charged its k≤21 portion.

\textbf{Theorem.} The entire actual subset

\[
       2^{1116}\le x<2^{1820},\qquad22\le k\le48            \tag{QE2}
\]

has total absolute normalized contribution less than\(2^{-421}\). Each subset being removed is first intersected with the current remainder; every earlier numerical error is conservatively retained.

By K099(GE4--5), total eligibility q=26 forces x\textgreater{}\(2^{1116}\). Thus QE2 includes every remaining q26 term with22≤k≤48. The theorem also covers q25 terms satisfying QE2. It needs no enumeration of total eligible sets, Galois orbits, denominator ideals or extra eta values. The k1..21 bounds below are diagnostic only; they are not added a second time to the new union. The q25 lower-denominator portion and all k≥49 remain outside this theorem.

\subsection{2. The full retained counting expression}\label{the-full-retained-counting-expression}

We use the whole DO105 proof, not just its final numerical envelope. For each fixed25-place support there is at most one fully eligible eta with x\textless{}\(2^{1820}\), by EL6--8. For that eta, at each active place there is at most one actual diagonal label, including the central-only-pole case EL9--15. The complete weighted pair sum is therefore

\[
                         (T-1)^{-25}<2T^{-25}.              \tag{QE3}
\]

No second partner label or additional factor for the possible26th eligible but inactive place appears: the original support enumeration is already retained.

Let J be the canonical whole-C G ideal and set

\[
 L'=JI_\Sigma\mathfrak D_\eta,
 \quad z=(N_EL')^{1/(2f)},\quad R=17424/133,
 \quad M=\sqrt{R/U},\quad z_G=Q^{167/f}T^{-25/f}.             \tag{QE4}
\]

The actual integral intersection \(L\cap\eta^{-1}L=L\mathfrak D_\eta\) puts all167 original r values in L'. Their full positive norm bound and the global ideal cap give

\[
                    z\le A_k(x):=\min(M/\sqrt k,z_Gx^{1/(2f)}).
                                                               \tag{QE5}
\]

DO105(EL19--23) supplies the complete joint row/transverse count, central affine count and actual selected trace:

\[
\begin{split}
 \#r&\le K_kA_3(4U)^f\mathcal B_k
              (M+\epsilon z)^{2(k+3)f}z^{-2(k+3)f},\\
 K_k&=\binom{164}{k}47^k2^{2fk}(f!)^k
                      \left(\frac{87/32}{fk}\right)^{fk},
 \quad\mathcal B_k=\frac{(fk)!(3f)!}{((k+3)f)!},\\
 A_3&=\frac{528^{3f}}{(3f)!D_0^3}<2^{-4040},\qquad\epsilon=1/1024,\\
 \#d&\le V_F(H/2+750Uz^2)U^{-f}T^{50}x^{-1},\\
 V_F(C)&=\frac{((22/7)fC^2)^{f/2}}{(f/2)!\sqrt{D_F}},\qquad
 |c(h)|<2^{-929f/2}U^{-f/2}z^{-f}.
\end{split}                                                \tag{QE6}
\]

Here \(D_0=7^{140}29^{144}5^{84}13^{84}\) and \(D_F=2^{168}D_0\). The original transverse Gram determinant is4U and its actual arbitrary-ideal covering coefficient is less than epsilon; neither a moving integral basis nor a density-only lattice estimate is used. The central factor counts the exact affine difference coset \(iJ_0B_\eta I_{\Sigma,F}^2\), with no half-lattice substitution. Given eta,r,d, the unchanged normal form reconstructs h. Thus there is no additional isometry or projected-column fibre to discard. The harmless retained pivot factor164\textless{}\(2^8\) stays present.

Keep the complete Green bound \(2Q^{-335\cdot167}(N_EJ)^{335/2}\), the single support factor\(2^{168}\), all\(4^{334}=2^{668}\) common G patterns and QE3. These are precisely the actual finite and real factors of EL24--28.

\subsection{3. The new denominator endpoint and all cap branches}\label{the-new-denominator-endpoint-and-all-cap-branches}

After the fixed factors are extracted, the x,z dependence is

\[
 x^{-337/2}z^{b_k}
 (1+\epsilon z/M)^{2(k+3)f}V_F(H/2+750Uz^2),
 \qquad b_k=(328-2k)f\ge0.                                 \tag{QE7}
\]

For fixed x this is nondecreasing in z. For x2≥x1\textgreater0 the cap ratio is at most \((x_2/x_1)^{1/(2f)}\), including a change between the two branches in QE5. The three increasing factors in QE7 grow by at most powers164−k, k+3 and1 of x2/x1. Together with x\^{}-337/2 the net exponent is−1/2. Therefore the complete bound decreases with x, and QE2 may be bounded at

\[
                              X_0=2^{1116}.                 \tag{QE8}
\]

This is a single common endpoint. There is no union over the denominator x or the two cap branches.

The exact actual-host integer comparison is

\[
 1117^{168}U^{168}Q^{334}2^{1116}
       <R^{168}T^{50}
       \le1118^{168}U^{168}Q^{334}2^{1116}.                 \tag{QE9}
\]

Consequently, at QE8,

\[
 A_k(X_0)<M/\sqrt{1117}<M/33,\qquad
 H/2+750R/1117<155.                                        \tag{QE10}
\]

Since1117\textgreater164, the displayed global cap is valid and smaller than the nonzero-coordinate cap for every1≤k≤164. Exact integer/rational volume comparisons give

\[
 (1+1/(1024\cdot33))^{2\cdot167\cdot168}<2^3,
 \qquad V_F(155)<2^{679}.                                  \tag{QE11}
\]

These replace DO105's endpoint68036, inflation2 and central479. The full root discriminant, different, actual real trace normalization and all affine shifts are unchanged.

\subsection{4. Exact factorial cells and the fresh27-cell union}\label{exact-factorial-cells-and-the-fresh27-cell-union}

Use the same directed logarithms

\[
 L_T=1019333/1024>\log_2T,\quad L_Q=287843/4096<\log_2Q,
 \quad L_U=601356/4096<\log_2U,\quad L_R=28810/4096>\log_2R.
                                                               \tag{QE12}
\]

As in EL28, regroup the total positive M power334f before using any directed bound. The single Green2, support, G patterns and pivot contribute845; QE3 contributes1 and QE11 contributes3. The complete common exponent before central679 and the rational k-factor is

\[
\begin{split}
 C_{26}={}&845+1+3+2f-4040-929f/2+167fL_R-(335f/2)L_U\\
          &+336\cdot25L_T-335\cdot167L_Q-(337/2)\cdot1116\\
        =&\;930781733/4096.                                 \tag{QE13}
\end{split}
\]

The total T exponent336·25 retains335 from Green,2 from the central count, and−1 from the unique-label weight. For each1≤k≤164 let c\_k be the least integral strict binary ceiling of

\[
 \frac{\binom{164}{k}47^k2^{2fk}(f!)^k87^{fk}(fk)!(3f)!}
 {32^{fk}f^{fk}k^{fk}1117^{(164-k)f}((k+3)f)!}.             \tag{QE14}
\]

The full bound for that coordinate count is \(E_k=C_{26}+679+c_k\). The companion compares exact integers in every factorial ceiling and every claimed cell. In the newly removed range,

\[
 \max_{22\le k\le47}E_k=E_{47}=-5236187/4096<-1278,
 \qquad E_{48}=-1730011/4096<-422.                          \tag{QE15}
\]

There are26 cells in22≤k≤47 and one at48. Their full union is strictly bounded by

\[
                      26\,2^{-1278}+2^{-422}<2^{-421}.       \tag{QE16}
\]

This proves the theorem with one charge421. The exact adjacent estimate is \(E_{49}=1755685/4096>0\), and \(E_{164}=321182245/4096>0\). These failed sufficient upper bounds imply neither a sign nor existence of a contribution. The previously charged k≤21 cells are neither included nor silently substituted in QE16.

\subsection{5. Verification, chronology and remaining gates}\label{verification-chronology-and-remaining-gates}

The full proof uses K099 only to infer that total q26 belongs to the larger-denominator interval. The count itself proves QE2 and uses the unchanged whole DO105(EL6--29) actual paired-label, common-ideal, metric, Gaussian beta-integral, central and selected-trace interfaces. No total-eligibility or Galois multiplicity is counted independently. The companion checks the exact host/covolumes, all directed powers, both sides of the cap bracket, central and transverse volume bounds, simultaneous denominator monotonicity for every k, all164 cell ceilings, and the exact new27-cell union.

The current predecessor is \nolinkurl{checks/do105_joint_remainder_error.json}, with97 all-top and31 uniform entries, complete floors25/11. Append421 once; the historical451 charge remains intact. The companion supports \nolinkurl{--prior} for a later unchanged-byte chronological replay and verifies the complete ordered prefix and both full error sums below\(2^{-230}\). Neither complete support floor changes.

The integral modules, CM action and involution, polarization, theta line, original selected module and complex metric are the accepted ones. This quantitative subset bound creates no new source family, period nonvanishing, arithmetic-image witness or signed positivity. The remaining q25 lower-denominator sector, k≥49 larger-denominator sector, actual selected derivative/class-content map, source/period/DP44 and RH conclusions remain open.

\subsection{Canonical algebraic-geometric context}\label{canonical-algebraic-geometric-context-7}

The original fractional ideals and affine fibres are invertible-module conditions on Spec(O\_E)/Spec(O\_F), with the same CM involution, Hermitian polarization, theta line, biextension and complex metric. Exact inherited contextual anchors remain SGA3 I4.7.3/4.7.3.1; SGA6 VI Theorem3.3; FGA232 Theorem3.1 pp148--149; FGA236 Corollaries1.7/2.3 pp229/232; SGA7 I VII2.7/2.9 pp169--174 and VIII7.4(b) pp310--311; cote4 scan003 sections8.1.1--8.1.2 and scans005--006 sections8.4.1--8.4.3; cote54 scans005--007. No new lambda-ring/K-theory equality, l-adic or motivic comparison is inferred. DP\_MULT Q3e/DP44/DP52/DP54--65, THETA3 G1/A7, CX\_GRH CX-ALG-024 and GCT\_PLETH T-052 retain their actual arithmetic and global comparison gates.

\clearpage

\section{AM-DO-106 --- the complete actual central fibre and its finite Fourier norm}\label{am-do-106-the-complete-actual-central-fibre-and-its-finite-fourier-norm}

\textbf{Canonical checked integration --- 2026-11-23T02:04:33.994252+00:00.} Complete saved-proof and executable reviews and separate exact diagnostic replays passed. No numerical error charge is introduced. Frozen author history is preserved. AM-DO-106 CFI1--24 proves the complete actual central fibre Gamma=B\_eta I\_Sigma,F squared, covolume sqrt(D\_F) x T\^{}-50, sufficient period Q T squared B\_eta of index Q\^{}168 T\^{}386, and the full-real codifferent dual. The original active type-II factor has Fourier l1 norm2/T for integral cross rows and1/(T-1) for a fixed cross pole. The actual mixed G projectors and every inactive-T factor are retained by the finite positive-projection trace lemma. The center-independent whole-C Green budget gives total finite Fourier l1 at most2 Q\^{}(-335*167)(N\_E J\_C)\^{}(335/2)(2/T)\^{}j(T-1)\^{}(-(25-j)), with no finite-index loss. The complete orbit lies in the original finite support; regularized Poisson retains its exact signs and covolume. Extending a disjoint collection past H133 uses at most the single complete2\^{}-2000 tail once, subject to reconciliation of earlier removals. No numerical alias estimate, sector removal, new charge or floor change follows; current98/31 errors remain. Actual mixed source, selected arithmetic content, prescribed period, signed inequality, specific032, DP44 and RH remain open.

Author: \nolinkurl{/root/definite_oscillator}, ALP\_MACHINE. 23 November 2026 UTC.

Controlled pending proof. Saved-file author verification, exact diagnostics, independent review, and canonical integration are separate events. This result adds no error charge and removes no numerical sector. DO105 and all earlier frozen files remain unchanged.

\subsection{1. Actual sector, parameter, and conclusion}\label{actual-sector-parameter-and-conclusion}

Retain the actual J56 field, every original finite level and pencil, both CM components, selected holomorphic module, packet point, central character, class masses, and Haar normalizations. Work with a matrix in the remaining DO105 sector: support Sigma of size25, all active places of type II and fully eligible, globally rank-two unipotent scalar-cross, with

\[
 E^{169}=Z\perp P,\quad Z=S_{164}\perp V_3,\quad
 P=Eu\perp Ee_n,\quad h_{ZZ}=I,\quad
 C=\ell r\ne0,\quad \ell=(1,\eta),\quad\eta\bar\eta=U.
 \tag{CFI1}
\]

Here C is nonzero on S, U=77Q²−14, f=168, and the original primitive integer vector u and fixed transverse frame are exactly DO104/105's. The middle bound2\^{}120\textless N\_E D\_eta\textless2\^{}1820 and22≤k≤164 identify the intended application, but the arithmetic lemmas below do not use their numerical endpoints. Fix eta,r and one actual matrix h0 in this sector. This is a nonempty-fibre assertion; no existence of h0 is proved.

In ambient coordinates set

\[
 w=(u,\eta),\quad w^\dagger=(u^*,-\bar\eta),\quad
 N(d)=wR_0-V_0w^\dagger+dww^\dagger,
 \quad V_0=J_Z^{-1}R_0^*,\quad R_0|_P=0,
 \quad d+\bar d=-q,\quad q=R_0V_0>0.
 \tag{CFI2}
\]

This is the actual DO104(CD10)/DO105(EL13) normal form, not an integral moving frame. Let d0 be h0's scalar. For t∈F define

\[
 n(t)=I+i t\,ww^\dagger,\qquad h(t)=h_0n(t)=n(t)h_0=h(d_0+i t).
 \tag{CFI3}
\]

The equality follows from R0w=wdagger V0=wdagger w=0. Thus n is an additive unitary group, det n=1, and the whole fibre has the same C and N²=−q w wdagger. Every h(t) is globally rank-two unipotent; no finite positivity is invoked.

Write D=D\_eta for the exact E-denominator ideal. DO104(CD3--5) gives

\[
 D\bar D=B\mathcal O_E,\quad B=B_\eta\subset\mathcal O_F,
 \quad N_FB=N_ED=x.
 \tag{CFI4}
\]

Put I=product over v∈Sigma of q\_v\^{}−1 as an F-ideal. The \textbf{complete central lattice}, without a chosen G rank pattern, is

\[
 \Gamma=BI^2\subset F,\qquad
 \operatorname{covol}(\Gamma)=\sqrt{D_F}\,xT^{-50}.
 \tag{CFI5}
\]

The theorem proves that d0+iGamma is exactly the set of central scalars preserving the chosen active pencil lattices and the original standard lattices elsewhere. Every such matrix remains in the actual finite support of K062(AT14--20). It also proves an exact finite-period/Fourier formulation, including the all-II indicator and the two mixed G projectors. The current97/31 error arrays and complete floors25/11 are unchanged. No positive-frequency suppression or signed-kernel bound is asserted.

\subsection{2. Exact central lattice at every finite place}\label{exact-central-lattice-at-every-finite-place}

The actual away-T compact is the standard lattice stabilizer, including at every ramified prime: R025(R25.1--5) defines Lambda=O\_E\^{}169 and K=Stab(hat Lambda), with principal G kernel exactly modulo Q. Thus an integral unitary matrix and its integral inverse meet the required away-T compact condition. There is no hyperspecial substitution at ramified rational primes.

Outside Sigma, h0 preserves Lambda. Hence h(t) preserves Lambda exactly when n(t) does. In the original four positive coordinates and the negative coordinate, the entries of wwdagger include u\_i u\_j, −bar(eta)u\_i, eta u\_j, and −U. Primitive integer Bezout for u, applied twice to the4×4 block and once to each mixed block, shows that the content ideal of wwdagger is

\[
 (1,\eta,\bar\eta)\mathcal O_E=B^{-1}\mathcal O_E.
 \tag{CFI6}
\]

For clarity, at a paired prime write a=v\_P(eta), e=v\_P(U)≥0; then v\_P(bar eta)=e−a, and the minimum entry valuation is min(0,a,e−a)=−max(0,−a,a−e). This is exactly −v\_P(D bar D). At an inert prime eta is integral. The accepted E/F is unramified at every finite place; the descent CFI4 applies also over the ramified rational primes. Since i is an integral unit, n(t) is integral precisely when t∈B locally. Its inverse is n(−t), so the condition is equality of lattices. This also proves necessity, with no factor2 or half-lattice at a dyadic place.

At v∈Sigma use the original selected component and the actual label c\_eta of DO105(EL9--16):

\[
 M_c=A_c\Lambda,\quad A_c=I+c\pi^{-1}zq_*,\quad q_*z=0,
 \quad w_c=A_c^{-1}w,\quad b_c=w^\dagger A_c.
 \tag{CFI7}
\]

Full eligibility makes eta a unit at both partners, and U is a T-unit, so B is a unit here. DO105 proves that w\_c and b\_c are divisible by pi at the unique admitted active label. Each has \textbf{exact} minimum valuation1. Indeed A\_c and A\_c\^{}−1 have entries of valuation at least−1; if w\_c were divisible by pi², then w=A\_c w\_c would be divisible by pi, contradicting its original positive unit coordinate u1=2. The row argument is identical using wdagger=b\_c A\_c\^{}−1. Therefore the rank-one matrix w\_c b\_c has exact content pi², and

\[
 A_c^{-1}n(t)A_c=I+i t w_c b_c\in\operatorname{Aut}(\Lambda)
 \quad\Longleftrightarrow\quad v_v(t)\ge-2.
 \tag{CFI8}
\]

The inverse has the same condition. The opposite CM lattice is the actual inverse-adjoint paired lattice; unitarity carries CFI8 to that component. The F-valuation of t is the same at both split partners. No independent conjugate label, congruence, or factor is inserted.

Since h0 preserves each chosen M\_c, CFI6--8 prove both directions of CFI5. In particular every t∈Gamma gives an actual incidence at every active place; at inactive T places it fixes Lambda, which is the included label0. At all other places it lies in the original compact. Thus every h(t) belongs to K062's full set mathscr C, even beyond the H133 cutoff. The fixed C is nonzero, and h(t)\_ZZ=I; consequently t↦h(t) is injective even in the quotient by mu28. The lattice factor in CFI5 is the full real covolume; multiplication by i has real absolute Jacobian1.

\subsection{3. Exact all-II indicator on this complete fibre}\label{exact-all-ii-indicator-on-this-complete-fibre}

At an active place call the cross data \emph{integral} when both original C blocks, at the selected and conjugate components, are integral. This is equivalent to R0 and V0 being integral in the fixed Z frame, whose determinant has only the T-unit factor2U. Let j be the number of these places,0≤j≤25.

If a cross pole occurs at either component, every h(t) retains it because C is fixed. A full-pencil-preserving unipotent has both cross blocks integral by K071(NB5). Every h(t) fixes M\_c, so the incidence classification forces unequal pencils and a unique fixed incidence, namely type II, for all t∈Gamma. Its exact normalized factor is1/(T−1).

Suppose instead that both cross rows are integral. Retain the original full-pencil pair

\[
 \mathcal M=\{x\in\Lambda:q_*x\in\pi\mathcal O\},\qquad
 \mathcal N=\Lambda+\pi^{-1}\mathcal O z.
 \tag{CFI9}
\]

DO105's same-label row replacement permits the literal q\emph{=−v0}J−pi m with mz=1 and m=e1\^{}t/z1, without changing any actual label or lattice. Since z∈P and V0∈Z,

\[
 R_0z=0,\quad q_*V_0=-\pi mV_0,\quad
 w,V_0\in\mathcal M,\quad w\in\pi\mathcal N,\quad
 R_0\mathcal N\subset\mathcal O,\quad
 w^\dagger\mathcal M\subset\pi\mathcal O,
 \quad w^\dagger\mathcal N\subset\mathcal O.
 \tag{CFI10}
\]

Hence the mixed part wR0−V0wdagger maps N into M. The same-label cancellation has a,b units in DO105(EL11--12), in particular wdagger z/pi=b is a unit. For any d∈pi\^{}−2 O, the central term d w wdagger preserves both M and N if v(d)≥−1. If v(d)=−2, its value on pi\^{}−1z has an entry of valuation−2, while the mixed part is integral; it cannot preserve N. Therefore, among the matrices already preserving M\_c,

\[
 h(d)\text{ has active type II at }v
 \quad\Longleftrightarrow\quad v_v(d)=-2.
 \tag{CFI11}
\]

For v(d)≥−1 the full pencil is preserved; when v(d)≥0 the original matrix is integral. This distinction is proved locally, rather than inferred from a cross pole.

Since h0 is active type II, d0 has valuation−2 in this integral-cross case. As t runs through the local Gamma=pi\^{}−2 O, the excluded condition d0+i t∈pi\^{}−1 O is one affine coset of the index-T subgroup pi\^{}−1 O. Write its indicator as1\_Av(t). The exact restricted local factor is thus

\[
 b_v(t)=\begin{cases}
 (T-1)^{-1},&\text{a fixed cross pole occurs},\\
 (T-1)^{-1}(1-\mathbf1_{A_v}(t)),&\text{both cross rows are integral}.
 \end{cases}                                                \tag{CFI12}
\]

The paired component has already been incorporated through the actual unitary lattice relation. CFI11--12 use one relative residue variable, not two. The global fractional-ideal residue map to this one-dimensional quotient is surjective, by the Chinese remainder theorem for the invertible ideal Gamma.

For a function on a finite abelian group, use normalized Fourier coefficients and their sum of absolute values, denoted \textbar\textbar·\textbar\textbar\_A. The second line of CFI12 has constant coefficient1/T, and T−1 nontrivial coefficients of magnitude1/{[}T(T−1){]}. Consequently

\[
 \|b_v\|_A=2/T\quad\text{in the integral-cross case},\qquad
 \|b_v\|_A=1/(T-1)\quad\text{in the cross-pole case}.
 \tag{CFI13}
\]

This is an exact Fourier norm, including the affine phase of A\_v. Its product over active places is at most (2/T)\^{}j(T−1)\^{}−(25−j), and in particular less than2\^{}25(T−1)\^{}−25. It replaces no actual diagonal-label count; eta and its unique labels were fixed before forming this central orbit.

\subsection{4. A genuine finite period, including inactive T places and mixed G factors}\label{a-genuine-finite-period-including-inactive-t-places-and-mixed-g-factors}

A sufficient common period inside Gamma is

\[
 \Gamma_{\rm per}=Q T^2B\mathcal O_F\subset\Gamma,
 \qquad [\Gamma:\Gamma_{\rm per}]=Q^f T^{2f+50}.
 \tag{CFI14}
\]

At every G place, CFI6 shows that n(t) is principal modulo the original rational Q for t∈Gamma\_per. R025(R25.2--3) and K058(PG3) therefore make its action trivial in the actual finite representation, including both mixed constituents. At every T place its original entries differ from I by multiples of pi². Every full-pencil lattice L satisfies pi Lambda⊂L⊂pi\^{}−1 Lambda, so (n(t)−I)L⊂pi Lambda⊂L, and the inverse does likewise. Thus n(t) fixes every actual source lattice, and the complete incidence matrix M\_(h0 n(t)) is unchanged. The original support/integrality conditions and CFI12 are also unchanged. There is no change at any other place.

Define on t∈Gamma the \textbf{actual restricted finite weight}

\[
 F(t)=a(h(t))\prod_{v\mid T,\ v\notin\Sigma}\rho_v(h(t))
                  \prod_{v\in\Sigma}b_v(t),
 \tag{CFI15}
\]

with a and rho exactly K062(AT8--16). It is zero when some active place ceases to be type II. The outside-Sigma matrices remain integral by Section2; their full rho factors are retained, not set to1. CFI14 makes F a function on the genuine finite quotient Gamma/Gamma\_per. No source finite character at T is missing: AT7--8 assigns value1 to each actual incidence, and the subtraction projector is exactly P=I−J/(T−1).

The following elementary Fourier fact avoids any loss equal to the large index in CFI14. If U is a unitary representation of a finite abelian group A, V is a commuting unitary, P is a finite-rank orthogonal projection, and D=rank P\textgreater0, then

\[
 f(a)=\operatorname{Tr}(PVU(a))/D\quad\Longrightarrow\quad
                           \|f\|_A\le1.                    \tag{CFI16}
\]

Indeed the character projections E\_chi commute with V, and the Fourier coefficient is Tr(E\_chi P E\_chi V)/D. The operator E\_chi P E\_chi is positive trace class, so its trace against a unitary has absolute value at most its trace. Summing over chi gives Tr P/D=1. P need not commute with U or V. This proof also works on a Hilbert space with finite-rank P and a finite-quotient action on the relevant cyclic subspace.

Apply CFI16 at each G place to the actual R, h0 and Qgood. The commutation comes from CFI3. In particular no commutation of a mixed projector with h0 or with the center is assumed. At an inactive T place realize the local group unitarily on l2 of all lattices and take P to be the mean-zero projection supported on the actual D. Then rho(h)=Tr(PU(h))/(T−2), exactly AT7--11. The compact central parameter group has finite orbit on this finite support; the kernel in CFI14 fixes that support, and by commutation also every h0-translate. It therefore acts trivially on the closed cyclic span of all those translates. CFI16 applies there, even if that span is infinite-dimensional.

Normalized Fourier coefficients of a product are convolutions, so their l1 norm is submultiplicative. This proves, without a finite-index multiplier,

\[
                         \|F\|_A\le(2/T)^j(T-1)^{-(25-j)}.
 \tag{CFI17}
\]

\subsection{5. Retaining the original Green budget before the central sum}\label{retaining-the-original-green-budget-before-the-central-sum}

A stronger form keeps the original unmixed G decay. Define a \textbf{center-independent whole-C pattern} at each of the334 unmixed relative places by the vanishing of C at its two CM components. Let J\_C be the E-ideal with both primes when both C blocks vanish, the zero side when exactly one vanishes, and neither when neither vanishes. There are at most4\^{}334 patterns, and N\_E J\_C≤Q\^{}334. This is a pattern of the fixed C, not a separate pattern for each central scalar.

Use the fixed G-unit Z/P frame, whose determinant factor2U is a G-unit. After saturating the null column and annihilating row locally, the nilpotent block has the algebraic shape

\[
 N_v=\begin{pmatrix}0&0&-v\\ r&0&d\\0&0&0\end{pmatrix}
 \quad\text{in the order }Z,w,\text{opposite line}.
 \tag{CFI18}
\]

To justify this shape, take a primitive local generator of the saturated null line and a primitive generator of its annihilating row in the rank-two P lattice. A unit-basis completion exists because the row is primitive and its kernel is generated by that column. Original integrality then makes the scalar central coefficient and the two cross vectors integral. The paired relation h\^{}−1\_ZP=−h\_ZP and h\^{}−1\_PZ=−h\_PZ in CFI2 identifies vanishing of the upper cross block with vanishing of C at the companion component. Thus the displayed r and v represent exactly the two paired C blocks. Saturation here is only over the local DVR; it asserts neither a global integral moving basis nor a positive finite form. If both are nonzero after reduction, rank N\_v=2 for every d.~If exactly one is nonzero, rank N\_v=1 for every d.~If both vanish, ranks0/1 vary with d alone. The companion component has the same rank by inverse-adjointness. These assertions follow directly from the row and column supports in CFI18, also when rv=0.

At an unmixed GL169(F\_q) place the normalized cuspidal unipotent trace is

\[
 A_2=[(q^{167}-1)(q^{168}-1)]^{-1},\qquad
 A_1=-(q^{168}-1)^{-1},\qquad A_0=1,
 \tag{CFI19}
\]

as in the directly retained Green/Gorodetsky--Hazan formula in LIFT064(GP1). If both C blocks vanish, the exact lattice CFI5 makes the central rank-one coefficient run through every F\_q residue: the central endomorphism, multiplied by a local generator of B, is primitive by CFI6. The local finite function is1 at one affine residue and −c at the other q−1 residues, c=(q168−1)\^{}−1. Its Fourier coefficients have magnitudes (1−(q−1)c)/q at zero and (1+c)/q at each nonzero character. Their sum is exactly1. Hence CFI16 loses nothing in this case.

In the other cases the coefficient is the constant A1 or A2. Multiplying the same original correction factors (bounded in total by2) gives

\[
 \|a(h(\cdot))\|_A
 \le2Q^{-335\cdot167}(N_EJ_C)^{335/2}.
 \tag{CFI20}
\]

For a one-zero place, the exponent167.5 is conservative: the true loss relative to rank2 is167. The two mixed factors remain bounded by1 by CFI16. No independent rank-zero pattern, orientation, or Fourier-cardinality factor is added. Combining CFI13, the inactive-T norm bound and CFI20 yields

\[
 \boxed{\quad\|F\|_A\le
 2Q^{-335\cdot167}(N_EJ_C)^{335/2}
       (2/T)^j(T-1)^{-(25-j)}.\quad}                         \tag{CFI21}
\]

The same original primitive-u extraction gives r(y),eta r(y)∈J\_C I\_Sigma for all167 counted inputs; their intersection is (J\_C I\_Sigma)D\_eta exactly as in DO104(CD12). Thus the Fourier norm uses the same covector congruence information as the earlier counts. It does not declare the scalar d to lie in a rank-zero ideal or retain a central-rank pattern after summing d.

\subsection{6. Exact Poisson formulation and cutoff handling}\label{exact-poisson-formulation-and-cutoff-handling}

Embed Gamma in F\_R with the full real metric Tr\_F(t²). Let

\[
 c(t)=\frac{\operatorname{Tr}_W(P_W\pi(g_p^{-1}h(t)_\infty g_p)|_W)}
                   {\dim W}.
 \tag{CFI22}
\]

This is the actual selected vector-bundle coefficient. In particular the embeddings of the chosen algebraic i, the near-point conjugation, compact factors, and central phase are left in this formula. No replacement by a scalar radial kernel or a newly selected vector occurs.

The complete orbit lies in mathscr C by Section2 and has no central-quotient repetitions. K062(AT18--20) and R048(RC48.4--6) therefore give absolute convergence of sum\_(t∈Gamma)\textbar c(t)\textbar: bounded operator norm gives only finitely many matrices, and the infinite tail is controlled. For epsilon\textgreater0 set c\_epsilon(t)=c(t)exp(−pi epsilon\textbar\textbar t\textbar\textbar²). The actual selected vectors are smooth Harish--Chandra vectors. Repeated derivatives along this commuting real center are matrix coefficients obtained by applying fixed Lie-algebra monomials to those vectors, including the fixed Ad(g\_p) conjugation. Unitarity bounds each such derivative uniformly in t by the corresponding fixed Hilbert norms. Hence c\_epsilon is Schwartz.

Use Fourier transform

\[
 \widehat\phi(\xi)=\int_{F_\mathbb R}\phi(t)
                         e^{-2\pi i\operatorname{Tr}_F(t\xi)}\,dt.
\]

The exact dual ideals are

\[
 \Gamma^*=\mathfrak D_F^{-1}B^{-1}I^{-2},\qquad
 \Gamma_{\rm per}^*=Q^{-1}T^{-2}\mathfrak D_F^{-1}B^{-1}.
 \tag{CFI23}
\]

Choose representatives theta\_chi of Gamma\_per\emph{/Gamma} for the characters of Gamma/Gamma\_per, and write F(t)=sum\_chi b\_chi exp(2pi i Tr\_F(theta\_chi t)). Periodizing a Schwartz function over Gamma and integrating its Fourier series on a fundamental domain proves the usual Poisson formula directly; rapid decay justifies both sums and the coefficient integral. Applying it to every finite character gives

\[
 \sum_{t\in\Gamma}F(t)c(t)
 =\lim_{\epsilon\downarrow0}\frac{T^{50}}{\sqrt{D_F}\,x}
   \sum_\chi b_\chi
       \sum_{\xi\in\Gamma^*}
                  \widehat c_\epsilon(\xi-\theta_\chi).
 \tag{CFI24}
\]

The limit on the left is justified by the already proved absolute convergence. CFI24 does not interchange the epsilon limit with an unbounded dual sum, nor claim a uniform bound for that sum. It makes the remaining analytic requirement exact. The bound sum\textbar b\_chi\textbar{} in CFI21 has no period-index loss.

If a disjoint collection of these central fibres is selected by its fixed support, eta and C, and each is extended from the actual H133 cut to its entire orbit while retaining F(t), the total newly included absolute tail is at most2\^{}-2000, \textbf{once for the collection}, by R048(RC48.6). This follows because all nonzero-weight terms are distinct actual-support matrices outside H133. It is not a separate2\^{}-2000 allowance for each fibre. Any change from the current set-difference remainder to a sum over complete fibres must also reconcile older removed subsets; this finding makes no new accounting claim or charge.

\subsection{7. Exact boundary, sources, and verification}\label{exact-boundary-sources-and-verification}

Positive support of a central Fourier transform, even if subsequently proved in the actual section model, does not by itself bound CFI24. The actual selected energies involve B=2\^{}955 and A=2\^{}960; a fine physical center lattice can still have a dual mesh far below those energies, with many aliases. DO080's radial fidelity estimate alone supplies no higher-power center estimate. A valid continuation must bound the actual transformed selected trace at the translated dual lattice in CFI24, or use a different exact cancellation. This theorem proves neither that analytic bound nor any numerical reduction of the dense sector.

The new proofs are CFI6--8's exact stabilizer lattice, CFI10--13's complete all-II indicator, CFI16's finite trace Fourier norm, CFI18--21's center-independent Green restoration, and CFI23--24's exact arithmetic Poisson interface. Direct carried locators are R025(R25.1--5), K058(PG1--3), K062(AT1--20), K071(NB1--7), DO104(CD3--23), and DO105(EL9--16). The unipotent finite-character normalization is Gorodetsky, O., \& Hazan, Z. (2017), \emph{On certain degenerate Whittaker models for cuspidal representations of GL\_(kn)(F\_q)}, arXiv:1707.07308, Theorem2.1/equation(2.1), as already directly checked in LIFT064(GP1); its settled formula is used without a new source-theorem proof. The exact selected section normalization is the carried K043/K046/DO080 realization from Nakahama (2016), and is not replaced here.

The companion diagnostic checks general denominator valuations, split inverse-adjoint local models, exact central lattice content, every residue in sampled central fibres, the all-II Fourier constants, finite rank cases including isotropic transverse rows, noncommuting finite projector models, actual period/covolume exponents and unchanged complete error arrays. These finite diagnostics supplement the proofs; they are not globally admissible short examples or a numerical evaluation of the actual central transform.

The integral model remains the original correspondence over Spec(O\_E); the full real codifferent and original theta/automorphic section are retained. The accepted contextual anchors remain SGA3 I4.7.3/4.7.3.1, SGA6 VI Theorem3.3, FGA232 Theorem3.1, FGA236 Corollaries1.7/2.3, and the previously recorded SGA7/Poincare-biextension locators. No new comparison theorem, lambda operation, l-adic weight statement or motivic identity is inferred. Signed inequality, prescribed period, selected arithmetic image/content, mixed-source conditions, specific032/Read, DP44 and RH gates remain open.

\subsection{Canonical algebraic-geometric context}\label{canonical-algebraic-geometric-context-8}

The exact fractional ideals and codifferent define the original arithmetic central correspondence over Spec(O\_F), including ramified and dyadic rational primes; no rational moving basis is declared globally integral. The original hermitian vector bundle and finite representation are unchanged. The original fractional ideals and affine fibres are invertible-module conditions on Spec(O\_E)/Spec(O\_F), with the same CM involution, Hermitian polarization, theta line, biextension and complex metric. Exact inherited contextual anchors remain SGA3 I4.7.3/4.7.3.1; SGA6 VI Theorem3.3; FGA232 Theorem3.1 pp148--149; FGA236 Corollaries1.7/2.3 pp229/232; SGA7 I VII2.7/2.9 pp169--174 and VIII7.4(b) pp310--311; cote4 scan003 sections8.1.1--8.1.2 and scans005--006 sections8.4.1--8.4.3; cote54 scans005--007. No new lambda-ring/K-theory equality, l-adic or motivic comparison is inferred. DP\_MULT Q3e/DP44/DP52/DP54--65, THETA3 G1/A7, CX\_GRH CX-ALG-024 and GCT\_PLETH T-052 retain their actual arithmetic and global comparison gates.

\clearpage

\section{AM-R-085 --- the exact selected coefficient on an actual central fibre}\label{am-r-085-the-exact-selected-coefficient-on-an-actual-central-fibre}

\textbf{Canonical checked integration --- 2026-11-23T02:04:34.550354+00:00.} Complete saved-proof and executable reviews and separate exact diagnostic replays passed. No numerical error charge is introduced. Frozen author history is preserved. AM-R-085 RF1--31 gives the exact genuine selected vector-valued central coefficient c(t)=P(t)/(z0-epsilon i gamma t)\^{}s, with s=kappa+K+2N and degree P\textless=s-kappa, using the original Pieri component and its exact Hilbert/PBW norm. The finite coefficient rule, partial fractions, signed gamma Fourier transform and affine Poisson phases are explicit. The positive central density for the same selected component has mass one and pointwise dominates the transformed commuting fibre, without assuming its projection commutes. Products over all168 labels lie in the actual signed orthant and vanish on its boundary faces. Ordinary absolutely convergent Poisson holds for the actual lattice and finite characters supplied separately by DO106, with no quotient-cardinality loss. This removes the unspecified analytic-model step but proves no numerical dual-lattice alias bound, finite sector removal or new error charge. Current98/31 totals and floors25/11 remain; the signed inequality, actual source, selected arithmetic image, prescribed period, specific032, DP44 and RH remain open.

Author: GPT-6 Astra Pro, \nolinkurl{/root}, ALP\_MACHINE; reasoning effort is not externally exposed. 23 November 2026 UTC.

\textbf{Pending author and independent whole-proof checks.} This is an exact analytic model and convergent summation identity, with no numerical remainder removal, finite error charge, support-floor change or nonvanishing assertion. It retains the vector-valued representation, the prescribed selected compact type and its actual Hilbert normalization.

\subsection{1. Actual representation and central parameter}\label{actual-representation-and-central-parameter}

At one real label use the actual K043 representation of U(a,1), with a=168 in the J56 application. Its minimal compact type is V\_alpha tensor z\^{}beta. The selected type W is the multiplicity-one Pieri component of homogeneous polynomial degree

\[
 N=\sum_{i=0}^{a-1}(\lambda_i-\alpha_i),\qquad d_W=\dim W.
 \tag{RF1}
\]

Use exactly the K043 Hilbert/PBW isometry. Put

\[
 k_i=\alpha_i-\alpha_{a-1},\qquad K=\sum_i k_i,
 \qquad \kappa=\alpha_{a-1}-\beta,\qquad s=\kappa+K+2N.
 \tag{RF2}
\]

The k\_i form a partition with last part zero, so V\_k is a polynomial representation of GL(a) of homogeneous degree K. The actual all-top formulas give

\[
             \kappa=2\ell+q+a>a,
 \tag{RF3}
\]

where q is the literal canonical-twist increment, not a finite-field cardinality. All these integers are positive in the stated actual family; N and K may be very large. The theorem does not substitute a scalar holomorphic discrete series.

Let h\_0 be an actual globally rank-two unipotent arrow in the retained scalar-cross branch. In the original coordinates its central fibre is

\[
 h(t)=h_0 n(t)=h_0+i_{\rm fld}t\,w w^\dagger,
 \quad n(t)=I+i_{\rm fld}t\,w w^\dagger,
 \quad w^\dagger=w^*J,\quad w^\dagger w=0,
 \tag{RF4}
\]

where J=diag(I\_a,-1), h\_0w=w and w\^{}dagger h\_0=w\^{}dagger. Here t is a real parameter at one embedding; globally it is the corresponding real component of t in F. The element i\_fld is the actual square root of -1 in E. Write

\[
             \sigma(i_{\rm fld})=\varepsilon i,
             \qquad \varepsilon\in\{1,-1\}.
 \tag{RF5}
\]

No common sign across all selected CM labels is assumed.

Conjugate by the actual prescribed near-point element g\_p and put

\[
 g(t)=g_p^{-1}h(t)g_p=g_0+\varepsilon i t\,w' w'^\dagger,
 \quad w'=g_p^{-1}w,
 \quad \gamma=|w'_{a+1}|^2>0,
 \quad z_0=(g_0)_{a+1,a+1}.
 \tag{RF6}
\]

The fixed null vector is nonzero, so its negative coordinate is nonzero and gamma is positive. The actual rank-two unipotent calculation of K085 and DO080, section7, gives

\[
     \operatorname{Re}z_0=1+\tfrac12((g_0-I)^2)_{a+1,a+1}>0.
 \tag{RF7}
\]

The central case with h\_0=I is included separately below and has z\_0=1. The argument only needs Re z\_0\textgreater0. Every g(t) has determinant one. Any extension of the representation from SU(a,1) to U(a,1) by its prescribed compact central character therefore contributes the scalar one along this fibre.

Define the genuine normalized selected trace

\[
 c(t)=d_W^{-1}\operatorname{Tr}\bigl(P_W\pi(g(t))|_W\bigr).
 \tag{RF8}
\]

This includes both compact factors of the actual transported argument.

\subsection{2. Exact statement}\label{exact-statement-2}

There is an explicitly specified finite polynomial P, with complex coefficients determined by the actual matrices and selected component, such that

\[
 \boxed{\quad c(t)=\frac{P(t)}{(z_0-\varepsilon i\gamma t)^s},
                    \qquad \deg P\le s-\kappa.\quad}
 \tag{RF9}
\]

The formula in section4 specifies P by finite polynomial operations, degree-N coefficient extraction and the trace of the actual Fischer projector. It does not require diagonalizing an infinite operator or altering the selected integral or Hilbert lattice.

Equivalently, there are uniquely determined coefficients A\_nu, with zero terms permitted, for which

\[
 c(t)=\sum_{\nu=\kappa}^{s}
          A_\nu(z_0-\varepsilon i\gamma t)^{-\nu}.
 \tag{RF10}
\]

With Fourier convention

\[
       \widehat c(\xi)=\int_{\mathbf R}c(t)e^{-2\pi i t\xi}\,dt,
 \tag{RF11}
\]

the exact answer is

\[
 \boxed{\quad
 \widehat c(\xi)=\frac{2\pi}{\gamma}
 e^{-2\pi z_0\varepsilon\xi/\gamma}
 \sum_{\nu=\kappa}^{s}
     A_\nu\frac{(2\pi\varepsilon\xi/\gamma)^{\nu-1}}{(\nu-1)!},
                   \quad\varepsilon\xi>0;\qquad
 \widehat c(\xi)=0\quad(\varepsilon\xi\le0).
 \quad}
 \tag{RF12}
\]

The value at zero is zero since kappa\textgreater1. No assertion that the individual A\_nu are nonnegative is made. The function retains the phase from Im z\_0 and the full vector-valued selection in the A\_nu.

\subsection{3. The literal Nakahama action and the denominator sign}\label{the-literal-nakahama-action-and-the-denominator-sign}

Swap the negative coordinate with the first coordinate, as in K043(PB4), to use Nakahama's SU(1,a) ball of row vectors z. Write the block matrix of g(t)\^{}(-1) as

\[
                  g(t)^{-1}=\begin{pmatrix}a(t)&b(t)\\c(t)&d(t)\end{pmatrix}.
 \tag{RF13}
\]

In this section c(t) inside a block matrix is its a-by-1 block only; the normalized trace is always the separately labelled RF8. Put

\[
 M(z,t)=c(t)z+d(t),\qquad r(z,t)=a(t)z+b(t),
 \quad D(t)=\det d(t),\quad \Delta(z,t)=\det M(z,t).
 \tag{RF14}
\]

Take the determinant shift in K043(PB4) to be -alpha\_last. Nakahama(2016), equation(5.5), then gives exactly

\[
 [\pi(g(t))F](z)=\Delta(z,t)^{-\kappa}
       \tau_k(M(z,t)^{-1})F(r(z,t)M(z,t)^{-1}).
 \tag{RF15}
\]

The convergence condition in Theorem5.1 is kappa\textgreater a, which is RF3. The action is defined on the original homogeneous Hermitian bundle; tau\_k is not discarded.

The cofactor identity for an inverse matrix of determinant one gives

\[
           D(t)=g(t)_{a+1,a+1}=z_0-\varepsilon i\gamma t.
 \tag{RF16}
\]

Indeed, the determinant of the positive a-by-a principal block of g(t)\^{}(-1) is the complementary one-by-one principal minor of g(t). In RF6 the last diagonal entry of w'w'\^{}dagger is -gamma. This proves both the sign and the scale in RF16 from the actual action convention. One cannot choose a frequency sign independently at each label to change the actual arithmetic lattice.

\subsection{4. Finite coefficient extraction and the exact degree bound}\label{finite-coefficient-extraction-and-the-exact-degree-bound}

Because w'\^{}dagger w'=0 and g\_0w'=w', the inverse g(t)\^{}(-1)=g\_0\^{}(-1)-epsilon i t w'w'\^{}dagger is an affine rank-one matrix pencil. Every minor of such a pencil is affine in t: in its multilinear determinant expansion, two columns from the rank-one increment are proportional, so all terms containing t² vanish.

Every minor of M(z,t) is also affine in t. One may see this either from the rank-one t update of M for fixed z, or by Cauchy--Binet after writing

\[
             M(z,t)=[c(t)\ d(t)]\begin{pmatrix}z\\I_a\end{pmatrix}.
 \tag{RF17}
\]

Moreover M is a rank-one update of d in z, so each of its minors has total z-degree at most one. In particular every entry of adj M is affine in t and has z-degree at most one. There are affine-in-t coefficients L\_j(t) such that

\[
                 \Delta(z,t)=D(t)+\sum_{j=1}^{a}z_jL_j(t).
 \tag{RF18}
\]

Every entry of r adj M has the same affine-in-t property and total z-degree at most one. To verify the point without multiplying two loose degree bounds, use Cramer's rule: its j-th entry is the determinant of M with its j-th row replaced by r. This is a row-selected a-by-a minor of the product g(t)\^{}(-1)(z;I\_a). Cauchy--Binet expresses it as a linear combination of minors of the affine rank-one pencil g(t)\^{}(-1), with coefficients of z-degree at most one. Thus no t² term survives in this numerator either.

For a homogeneous V\_k-valued polynomial F of degree N, the right side of RF15 can therefore be written

\[
 \frac{R_F(z,t)}{\Delta(z,t)^{L}},\qquad
 L=\kappa+K+N,
 \quad \deg_t R_F\le K+N.
 \tag{RF19}
\]

Here tau\_k(adj M) has degree K in the matrix entries, and F(r adj M) contributes N more. This uses the homogeneous polynomial representation V\_k; all its matrix entries extend polynomially to all matrices. A determinant twist has already been incorporated into kappa.

Let {[} {]}\_N mean total homogeneous z-degree N. Expanding the denominator about z=0 gives the finite coefficient rule

\[
 \left[\frac{R_F(z,t)}{\Delta(z,t)^L}\right]_N
 =\sum_{j=0}^{N}(-1)^j\binom{L+j-1}{j}
       D(t)^{-L-j}
       [R_F(z,t)(z\cdot L(t))^j]_N.
 \tag{RF20}
\]

After putting this over D(t)\^{}(L+N)=D(t)\^{}s, the j-th numerator has t-degree at most

\[
               (K+N)+j+(N-j)=K+2N=s-\kappa.
 \tag{RF21}
\]

Thus RF9 holds for every entry of the degree-N coefficient operator. The expansion is valid near z=0 since D(t) is nonzero for real t; coefficient extraction is a finite algebraic identity, so it supplies the indicated rational function for every real t.

It remains to justify taking the actual selected trace without changing its metric. Distinct homogeneous degrees are Hilbert-orthogonal by the negative compact circle, and different Pieri components within one degree are orthogonal for both the invariant Hilbert metric and the Fischer metric. On this particular W, K043(PB5--8) gives exactly

\[
       \langle\ ,\ \rangle_{\mathcal H}|_W
            =D_{\lambda,\alpha,\beta}^{-1}
                \langle\ ,\ \rangle_{\rm F}|_W.
 \tag{RF22}
\]

In RF22--23 the Hermitian pairing is linear in its first argument; with the opposite convention interchange the arguments. This specifies the intrinsic operator trace rather than its complex conjugate.

Choose a Fischer-orthonormal basis e\_1,\ldots,e\_(d\_W) of W. The Hilbert-orthonormal basis is sqrt(D\_(lambda,alpha,beta))e\_j. The two factors from this basis cancel the inverse D in RF22 when computing the matrix coefficient. Consequently the exact RF8 is

\[
 c(t)=\frac1{d_W}\sum_{j=1}^{d_W}
   \left\langle[\pi(g(t))e_j]_N,e_j\right\rangle_{\rm F}.
 \tag{RF23}
\]

This is an identity only for the selected component with its actual common norm factor. It is not an assertion that the Fischer norm agrees with the Hilbert norm on the whole polynomial space. Using RF20 in RF23 and multiplying by D(t)\^{}s is the promised explicit finite formula for P. An arbitrary finite basis of W gives the same formula by its exact Fischer Gram inverse; an artificial choice of orthonormal algebraic square roots is unnecessary. The full actual projection and its dimension are retained.

\subsection{5. Exact Fourier transform and positive spectral majorant}\label{exact-fourier-transform-and-positive-spectral-majorant}

Since D(t) is a nonconstant affine polynomial, every polynomial of degree at most s-kappa has a unique expansion in its powers D(t)\^{}j, 0\textless=j\textless=s-kappa. This proves RF10. More explicitly, with t\_* = z\_0/(epsilon i gamma),

\[
        A_\nu=
        \frac{P^{(s-\nu)}(t_*)}{(s-\nu)!}
                (-\varepsilon i\gamma)^{-(s-\nu)}.
 \tag{RF24}
\]

For Re z\_0\textgreater0 and every positive integer nu,

\[
 (z_0-\varepsilon i\gamma t)^{-\nu}
 =\frac1{(\nu-1)!}\int_0^\infty
       u^{\nu-1}e^{-z_0u}e^{\varepsilon i\gamma tu}\,du.
 \tag{RF25}
\]

The candidate Fourier density in RF12 is integrable and its inverse Fourier transform is RF25 by the change of variable xi=epsilon gamma u/(2pi). The factor is 2pi/gamma, independent of epsilon. Each original term is integrable on the real line because nu\textgreater=kappa\textgreater1. Uniqueness of the Fourier transform therefore proves RF12. Its value at zero is the continuous zero limit. There is no assertion of termwise positivity for the finite partial-fraction expansion.

There is, however, an exact positive majorant that keeps the selected type. Let p(xi) be RF12 computed for the same transported central subgroup and h\_0=I. Then

\[
          p(\xi)\ge0,\quad\int_{\mathbf R}p(\xi)d\xi=1,
          \qquad |\widehat c(\xi)|\le p(\xi).
 \tag{RF26}
\]

Here p is the actual selected central spectral density, not a scalar Bergman or Gaussian replacement. To prove RF26, let U(t)=pi(n'(t)) be the one-parameter unitary group and E(B) its spectral measure in the RF11 convention. Put V=pi(g\_0). The original central relation gives VU(t)=U(t)V, hence V commutes with E(B). For the finite-rank orthogonal projection P\_W, define

\[
 \mu(B)=d_W^{-1}\operatorname{Tr}(P_W E(B)),\qquad
 \nu(B)=d_W^{-1}\operatorname{Tr}(P_W V E(B)).
 \tag{RF27}
\]

The first measure is positive of mass one. For every B, Tr(P\_W V E(B))=Tr(E(B)P\_WE(B)V); the latter positive trace-class factor gives \textbar nu(B)\textbar\textless=mu(B). Applying this inequality to every finite measurable partition proves the total-variation domination \textbar nu\textbar\textless=mu, not merely a bound on total mass. RF9--12 for h\_0=I identify mu with the continuous density p; the same formulas identify nu with the continuous density hat c.~Thus the Radon--Nikodym domination holds everywhere by continuity, proving RF26. The projection need not commute with either U or V.

\subsection{6. All real labels and ordinary convergent Poisson summation}\label{all-real-labels-and-ordinary-convergent-poisson-summation}

For the f=168 actual labels let c\_sigma be RF8 and put

\[
 C(t)=\prod_\sigma c_\sigma(t_\sigma),\qquad
 \widehat C(\xi)=\prod_\sigma\widehat c_\sigma(\xi_\sigma).
 \tag{RF28}
\]

The Fourier transform is with Euclidean pairing sum\_sigma t\_sigma xi\_sigma, equivalently the usual real trace pairing on F. Fubini applies because each factor is integrable. Its support is the explicitly signed closed orthant epsilon\_sigma xi\_sigma\textgreater=0, and it vanishes on every boundary face. Inside that orthant it is a finite multivariate polynomial times the explicitly given decaying complex exponential. Also \textbar hat C(xi)\textbar\textless=product\_sigma p\_sigma(xi\_sigma).

Let Gamma be any full lattice in R\^{}f and t\_0 any real vector. Suppose b on t\_0+Gamma is periodic modulo a finite-index sublattice, and expand it on that finite quotient as

\[
 b(t_0+v)=\sum_\chi b_\chi e^{2\pi i\langle\theta_\chi,v\rangle},
       \qquad v\in\Gamma.
 \tag{RF29}
\]

The theta\_chi represent its actual characters; changing one by an element of Gamma* changes no result. Then the exact ordinary, absolutely convergent identity is

\[
 \boxed{\quad
 \sum_{v\in\Gamma} b(t_0+v)C(t_0+v)
 =\frac1{\operatorname{covol}\Gamma}
      \sum_\chi b_\chi
      \sum_{\xi\in\Gamma^*-\theta_\chi}
                 e^{2\pi i\langle\xi,t_0\rangle}\widehat C(\xi).
 \quad}
 \tag{RF30}
\]

No finite-quotient cardinality multiplies this formula. In particular its absolute value is at most

\[
 \frac{\sum_\chi|b_\chi|}{\operatorname{covol}\Gamma}
   \sup_\theta\sum_{\xi\in\Gamma^*-\theta}
                    \prod_\sigma p_\sigma(\xi_\sigma).
 \tag{RF31}
\]

Only the finitely many actual theta\_chi are required in that supremum; the larger supremum is a valid upper bound.

For completeness, absolute convergence and pointwise Poisson here do not need an unproved Schwartz assertion. RF9 gives c\_sigma(t)=O((1+\textbar t\textbar)\^{}(-kappa\_sigma)); the actual minimum kappa\_sigma exceeds f.~On any translate of a lattice the product C therefore has a summable radial upper bound O((1+\textbar\textbar t\textbar\textbar)\^{}(-min kappa\_sigma)), uniformly for t\_0 in a compact fundamental set. Its periodization is continuous and uniformly absolutely convergent. Each modulated function e\^{}(2pi i theta t)C(t) has the same property. RF12 gives exponential decay times a finite polynomial for its Fourier transform on its signed orthant, extended by zero; hence its values on any shifted dual lattice are absolutely summable. The Fourier coefficients of the periodization, computed by absolute Fubini, are covol(Gamma)\^{}(-1) times these transformed values. The absolutely convergent Fourier series equals the periodization by uniqueness of Fourier coefficients for continuous torus functions. Evaluating at t\_0 and removing the modulation gives RF30 with the displayed shift and phase. RF26 and the triangle inequality give RF31.

In the actual arithmetic application, Gamma, the finite periodic factor and its Fourier coefficients must be those proved for the original central orbit, with both mixed G projectors, both T partners, all support restrictions and the actual Haar normalization. The lattice/finite-factor theorem DO106 is a separate input. RF30 does not justify adding arrows outside that original support. Nor does it justify subtracting a cutoff tail repeatedly: any extension of disjoint retained fibres uses the single complete already bounded outside-H tail as specified by the arithmetic accounting.

\subsection{7. Scope, verification and exact sources}\label{scope-verification-and-exact-sources}

This removes an unspecified analytic-model step: the actual selected central coefficient, every frequency sign, its finite polynomial/gamma transform and its full convergent lattice sum are now explicitly given. It does not show that the shifted dual lattices avoid the spectral region or that their alias sum in RF31 is small. The displayed identities supply no inequality comparing the shifted dual lattice with the selected spectral range; that numerical estimate remains to be proved. No high-frequency suppression follows from support in a signed orthant alone. The complete signed kernel inequality, prescribed period, selected arithmetic content, admissible mixed source, specific032 comparison, DP44 difference and Hecke RH remain open.

Primary source: Nakahama, R. (2016). \emph{Norm computation and analytic continuation of vector valued holomorphic discrete series representations}. Journal of Lie Theory, 26(4), 927--990. Equation(5.5), printed p.959, equation(5.7), p.960, and Theorem5.1, pp.960--961; author version arXiv:1506.05919. The published text at these exact passages was read directly from \nolinkurl{research/nakahama2016.txt}, lines2280--2445, during this proof. K043(PB1--10) supplies the exact actual representation, determinant shift and selected norm; K060(RC1) supplies the homogeneous-degree orthogonality. DO080(RT1, section7) and K085(section2) are used in their already checked genuine-argument scope. DO105(EL1--3) supplies the actual unipotent central normal form. The new determinant/coefficient, Fourier and periodization arguments are proved in full above.

The companion checks finite rank-one unitary matrix pencils, exact minor degrees and signs, finite selected-degree examples with the literal action, partial fractions, inverse-Fourier normalization, and the actual integer kappa range. Such finite diagnostics do not prove the general representation-theoretic argument; that is the saved proof. No Lean build or archive audit is performed.

The algebraic-geometric setting is the same automorphic homogeneous Hermitian vector bundle and its selected Schur module: polynomial matrix coefficients and coefficient extraction are performed in that actual representation. The central fibre and its Fourier dual remain the original arithmetic group and trace pairing. Nothing here changes the CM Picard scheme, polarization or Poincare biextension, and no new l-adic weight or motivic L-function identity is asserted. Contextual carried locators are SGA3 I4.7.3/4.7.3.1, SGA6 VI Theorem3.3, FGA232 Theorem3.1 pp.148--149, FGA236 Corollaries1.7 and2.3 pp.229 and232, SGA7 I VII2.7/2.9 pp.169--174 and VIII7.4(b) pp.310--311, cote4 scan003 sections8.1.1--8.1.2 and scans005--006 sections8.4.1--8.4.3, and cote54 scans005--007. DP\_MULT Q3e and THETA3 G1/A7 receive an explicit analytic interface only; GCT\_PLETH T-052, Q-ALPOGE and ALPOGE\_GCT Q-ALPOGE-METHOD retain their separate global comparison gates.

\subsection{Canonical algebraic-geometric context}\label{canonical-algebraic-geometric-context-9}

The finite polynomial extraction takes place in the prescribed homogeneous Hermitian bundle and its selected Schur module; the central fibre and Fourier trace pairing are the arithmetic ones. The result constructs an analytic model only, not a new arithmetic or motivic comparison. The original fractional ideals and affine fibres are invertible-module conditions on Spec(O\_E)/Spec(O\_F), with the same CM involution, Hermitian polarization, theta line, biextension and complex metric. Exact inherited contextual anchors remain SGA3 I4.7.3/4.7.3.1; SGA6 VI Theorem3.3; FGA232 Theorem3.1 pp148--149; FGA236 Corollaries1.7/2.3 pp229/232; SGA7 I VII2.7/2.9 pp169--174 and VIII7.4(b) pp310--311; cote4 scan003 sections8.1.1--8.1.2 and scans005--006 sections8.4.1--8.4.3; cote54 scans005--007. No new lambda-ring/K-theory equality, l-adic or motivic comparison is inferred. DP\_MULT Q3e/DP44/DP52/DP54--65, THETA3 G1/A7, CX\_GRH CX-ALG-024 and GCT\_PLETH T-052 retain their actual arithmetic and global comparison gates.

\clearpage

\section{AM-K101 --- a uniform denominator cutoff for every remaining row support}\label{am-k101-a-uniform-denominator-cutoff-for-every-remaining-row-support}

\textbf{Canonical checked integration --- 2026-11-23T02:09:48.095975+00:00.} Complete root and independent proof/script/interface reviews and separate exact replays passed. The actual canonical predecessor has98/31 entries; this integration gives99/31, complete floors25/11. Frozen author history is preserved. AM-K101 UC1--17 removes the complete actual fully eligible all-II globally rank-two-unipotent scalar-cross support25 subset with2\^{}1585\textless=x=N\_E D\_eta\textless2\^{}1820 for every22\textless=k\textless=164. It uses the original DO105 all-factor bound with a single common denominator endpoint, exact capmax(161,k), full selected trace and moving affine central cover. The143-cell union is below2\^{}-254, with nondense maximum-1412571/4096\textless-344 and dense maximum-1048027/4096\textless-255. Append254 once to actual K10098/31, giving99/31 and both whole sums below2\^{}-230; all earlier overlaps,451 and421 are retained and complete floors25/11 remain. The only still-unremoved support25 cells therefore have2\^{}120\textless x\textless2\^{}1585,22\textless=k\textless=164, with K100 additionally excluding22\textless=k\textless=48 when x\textgreater=2\^{}1116, and all K099 structural restrictions retained. No prospective inactive-T gain, signed positivity, actual source, selected arithmetic content, prescribed period, specific032, DP44 or RH result is inferred.

Author: \nolinkurl{/root/kummer_extension}. 23 November 2026 UTC.

Controlled pending finding. This is one complete fresh subset bound from the unchanged DO105 counting expression. K099 and K100 remain frozen. No new local or archimedean input, spectral replacement, or independent denominator enumeration is introduced.

\subsection{1. Actual sector and theorem}\label{actual-sector-and-theorem-5}

Keep the complete actual DO105 middle sector: the J56 all-top signed kernel, support25, H=133, f=168, global rank-two unipotence, scalar cross block nonzero on the original S164, h\_ZZ=I, and all25 active places of type II and eligible. Retain the original paired CM components, actual labels and omitted labels, finite projectors, selected real module, central phases, class masses, Haar measures and original lattice metrics. Write

\[
 C=(1,\eta)^t r,\qquad\eta\bar\eta=U,\qquad
 x=N_E\mathfrak D_\eta,\qquad
 k=\#\{j\in S_{164}:r_j\ne0\},
 \quad U=77Q^2-14,\quad R=17424/133.                       \tag{UC1}
\]

The previously retained middle window is 2\^{}120\textless x\textless2\^{}1820. DO105 has already charged its k≤21 part, and K100 has separately charged 2\^{}1116≤x\textless2\^{}1820,22≤k≤48.

\textbf{Theorem.} The complete actual subset

\[
                   2^{1585}\le x<2^{1820},\qquad22\le k\le164
                                                               \tag{UC2}
\]

has total absolute normalized contribution less than 2\^{}-254. Each removed subset is intersected with the current remainder before removal; every earlier numerical error, including DO105's451 and K100's421, is conservatively retained. Thus after this theorem every still-unremoved middle-sector term has x\textless2\^{}1585. No complete support floor changes.

This includes both possible total-eligibility counts q25 and q26, without enumerating total eligible sets, eta denominators or Galois orbits. The already charged k≤21 cells appear only as diagnostics in the companion. The theorem is independent of any prospective inactive-T improvement.

\subsection{2. Complete original count and monotonicity}\label{complete-original-count-and-monotonicity}

For each actual25-place support, DO105 EL6--8 gives at most one eligible eta with x\textless2\^{}1820. Its EL9--15 proves at most one actual active diagonal label even when only the central part has a pole. The full weighted active pair sum is therefore

\[
                             (T-1)^{-25}<2T^{-25}.          \tag{UC3}
\]

Let J be the canonical whole-C G ideal used in the original count and put

\[
 L'=JI_\Sigma\mathfrak D_\eta,\quad
 z=(N_EL')^{1/(2f)},\quad M=\sqrt{R/U},\quad
 z_G=Q^{167/f}T^{-25/f},\quad
 A_k(x)=\min(M/\sqrt k,z_Gx^{1/(2f)}).                     \tag{UC4}
\]

All167 original covector values lie in L' by the actual integral intersection, and z≤A\_k(x). Retain the whole DO105 EL19--28 expression, displayed also in K100 QE6:

\[
\begin{split}
 \#r&\le K_kA_3(4U)^f\mathcal B_k
       (M+z/1024)^{2(k+3)f}z^{-2(k+3)f},\\
 K_k&=\binom{164}{k}47^k2^{2fk}(f!)^k
                    \left(\frac{87/32}{fk}\right)^{fk},\qquad
 \mathcal B_k=\frac{(fk)!(3f)!}{((k+3)f)!},\\
 A_3&=\frac{528^{3f}}{(3f)!D_0^3}<2^{-4040},\\
 \#d&\le V_F(133/2+750Uz^2)U^{-f}T^{50}x^{-1},\\
 V_F(C)&=\frac{((22/7)fC^2)^{f/2}}{(f/2)!\sqrt{D_F}},\qquad
 |c(h)|<2^{-929f/2}U^{-f/2}z^{-f}.
\end{split}                                                \tag{UC5}
\]

Here D0=7\^{}140 29\^{}144 5\^{}84 13\^{}84 and DF=2\^{}168 D0. In particular the original transverse Gram determinant4U, the arbitrary-ideal transverse covering coefficient, the exact pure-imaginary affine central coset, the full real codifferent, and the selected trace are unchanged. Eta,r,d reconstruct h by the actual normal form; no extra isometry or column fibre is discarded. The retained harmless pivot factor164\textless2\^{}8 remains.

Multiply UC5 by UC3, the complete Green bound

\[
                    2Q^{-335\cdot167}(N_EJ)^{335/2},       \tag{UC6}
\]

the support factor2\^{}168 and all4\^{}334 common patterns. After extracting fixed factors, the x,z dependence is exactly

\[
 x^{-337/2}z^{b_k}(1+z/(1024M))^{2(k+3)f}
                V_F(133/2+750Uz^2),\qquad
                   b_k=(328-2k)f\ge0.                     \tag{UC7}
\]

It is nondecreasing in z at fixed x. For x2≥x1, the ratio of caps in UC4 is at most (x2/x1)\^{}(1/(2f)), also when the minimum changes branch. The three increasing factors in UC7 grow by at most powers164-k, k+3 and1 of x2/x1. Including x\^{}-337/2, their net exponent is -1/2. Hence the full count decreases with x. The entire interval UC2 may therefore be bounded at the single endpoint

\[
                              X_0=2^{1585}.                \tag{UC8}
\]

There is no union over x, a denominator ideal, or the two cap branches. This argument includes k164, where b\_k=0.

\subsection{3. Exact endpoint and every finite cell}\label{exact-endpoint-and-every-finite-cell}

The actual-host integer comparison is

\[
 161^{168}U^{168}Q^{334}2^{1585}
       <R^{168}T^{50}
       \le162^{168}U^{168}Q^{334}2^{1585}.                  \tag{UC9}
\]

For each k set a\_k=max(161,k). Then at UC8,

\[
           A_k(X_0)\le M/\sqrt{a_k},\quad
 \rho_k=\left\lfloor133/2+750R/a_k\right\rfloor+1,
 \quad133/2+750UA_k(X_0)^2<\rho_k.                         \tag{UC10}
\]

Equality in the nonzero-coordinate cap causes no problem; the rounded central radius is strictly larger. All a\_k are between161 and164, so floor(sqrt(a\_k))=12. Let I\_k be the least integral strict binary ceiling of

\[
       \frac{(1024\cdot12+1)^{2(k+3)f}}
                    {(1024\cdot12)^{2(k+3)f}},             \tag{UC11}
\]

and let V\_k be the least integral strict binary ceiling of V\_F(rho\_k). The exact radii and central ceilings at the top four caps are

\begin{longtable}[]{@{}rrr@{}}
\toprule\noalign{}
a\_k & rho\_k & V\_k \\
\midrule\noalign{}
\endhead
\bottomrule\noalign{}
\endlastfoot
161 & 677 & 1036 \\
162 & 674 & 1035 \\
163 & 670 & 1034 \\
164 & 666 & 1032 \\
\end{longtable}

Every I\_k≤7, and I\_163=I\_164=7. The companion verifies these inequalities by integer powers with the exact real covolume.

Use the retained directed logarithms

\[
 L_T=1019333/1024>\log_2T,\quad
 L_Q=287843/4096<\log_2Q,\quad
 L_U=601356/4096<\log_2U,\quad
 L_R=28810/4096>\log_2R.                                  \tag{UC12}
\]

Before using these directed bounds, combine the total positive M power as 2(k+3)f+b\_k=334f. The common exponent, excluding I\_k,V\_k and the rational row factor, is

\[
\begin{split}
 C_*={}&845+1+2f-4040-929f/2+167fL_R-(335f/2)L_U\\
       &+336\cdot25L_T-335\cdot167L_Q-(337/2)\cdot1585\\
     =&\;607076901/4096.                                  \tag{UC13}
\end{split}
\]

The845 retains the single Green2, support2\^{}168, common-pattern2\^{}668 and original pivot2\^{}8. The additional1 is UC3. The total T exponent336 per support place is335 from Green,2 from the central count and minus1 from the unique-label weight.

Let B\_k be the least integral strict binary ceiling of

\[
 \frac{\binom{164}{k}47^k2^{2fk}(f!)^k87^{fk}(fk)!(3f)!}
 {32^{fk}f^{fk}k^{fk}a_k^{(164-k)f}((k+3)f)!}.            \tag{UC14}
\]

The complete cell is bounded strictly by 2\^{}E\_k, where

\[
                              E_k=C_*+I_k+V_k+B_k.         \tag{UC15}
\]

The companion checks all164 exact factorial cells; the newly removed143-cell set is exactly22≤k≤164. Its sharp maxima within this rational envelope are

\[
 \max_{22\le k\le163}E_k=E_{163}=-1412571/4096<-344,
 \qquad E_{164}=-1048027/4096<-255.                        \tag{UC16}
\]

Therefore its entire union satisfies

\[
                  142\,2^{-344}+2^{-255}<2^{-254}.         \tag{UC17}
\]

No endpoint, support, type, pattern, central or sparse-position multiplicity is omitted. In particular the dense k164 cell belongs to the same union.

For scope of the cutoff search, the same retained directed envelope at the preceding integer endpoint x=2\^{}1584 has global integer cap161 and E\_164=-357851/4096. That sufficient bound does not meet a230-compatible union target. Thus1585 is the first successful integer endpoint in the final adjacent-cutoff comparison; no optimality claim about other constants, other methods, actual contributions or noninteger cutoffs is made.

\subsection{4. Chronology, verification, and remaining gate}\label{chronology-verification-and-remaining-gate}

The actual predecessor is K100, with98 all-top and31 uniform error entries. Append254 once, retaining the full ordered prefix and the earlier421. The companion permits an unchanged-byte \nolinkurl{--prior} replay against a later extension and checks both complete sums below2\^{}-230. The complete support floors remain25/11.

The whole proof reuses DO105 EL6--29 and K100 QE3--7, including the simultaneous x/z monotonicity, the original joint Gaussian beta integral, the exact central affine lattice and the fully selected trace. Its new work is the common endpoint1585, the changing per-k cap, the full164-cell exact scan and the fresh143-cell union. Old k≤21 estimates are not recharged. K100's overlapping portion is not subtracted from its historical error.

This cuts the upper part of the surviving middle denominator window uniformly in k. It supplies no new actual eta, signed positivity, complete support-floor increase, source-point or period construction, selected derivative/class-content witness, DP44 or RH consequence. The original integral CM modules, polarization, theta correspondence and complex section metric are unchanged.

\subsection{Canonical algebraic-geometric context}\label{canonical-algebraic-geometric-context-10}

The original fractional ideals and affine fibres are invertible-module conditions on Spec(O\_E)/Spec(O\_F), with the same CM involution, Hermitian polarization, theta line, biextension and complex metric. Exact inherited contextual anchors remain SGA3 I4.7.3/4.7.3.1; SGA6 VI Theorem3.3; FGA232 Theorem3.1 pp148--149; FGA236 Corollaries1.7/2.3 pp229/232; SGA7 I VII2.7/2.9 pp169--174 and VIII7.4(b) pp310--311; cote4 scan003 sections8.1.1--8.1.2 and scans005--006 sections8.4.1--8.4.3; cote54 scans005--007. No new lambda-ring/K-theory equality, l-adic or motivic comparison is inferred. DP\_MULT Q3e/DP44/DP52/DP54--65, THETA3 G1/A7, CX\_GRH CX-ALG-024 and GCT\_PLETH T-052 retain their actual arithmetic and global comparison gates.

\clearpage

\section{AM-DO-107 --- inactive pencil weights and one paired cross congruence}\label{am-do-107-inactive-pencil-weights-and-one-paired-cross-congruence}

\textbf{Canonical checked integration --- 2026-11-23T02:14:48.815609+00:00.} Complete root and independent proof/script/interface reviews and separate exact replays passed. The actual canonical predecessor has99/31 entries; this integration gives100/31, complete floors25/11. Frozen author history is preserved. AM-DO-107 IP1--23 removes three complete subsets of the actual fully eligible all-II globally rank-two-unipotent scalar-cross support25 middle sector: any inactive paired-C-zero place for all k22..164; no paired zero with x\textgreater=2\^{}744 for all k22..164; and no paired zero with x\textless2\^{}744 and22\textless=k\textless=99. The literal included-Lambda projector gives rho=1/(T-1) at every noneligible inactive place with a nonzero paired C block. A paired zero instead supplies an actual extra normT² row ideal, giving T\^{}-335 while retaining the old central affine cover. The complete143/143/78-cell union is below2\^{}-392. Append392 to the actual K10199/31 arrays, giving100/31; all historical overlaps including451,421,254 remain, both full sums are below2\^{}-230, and complete floors25/11 are unchanged. The only support25 branch left has no inactive paired C zero, total eligibility exactly25,2\^{}120\textless N\_E D\_eta\textless2\^{}744 and100\textless=k\textless=164, with K099 field/stabilizer restrictions retained. This is a strict smaller finite gate, not support25 exhaustion, signed positivity, an admissible mixed source, selected arithmetic image, prescribed global period, specific032, DP44 or RH.

Author: /root/definite\_oscillator, ALP\_MACHINE. 23 November 2026 UTC.

Controlled pending proof. Whole saved-file author verification, unchanged exact replay, independent review and canonical integration are separate events. All earlier frozen files remain unchanged.

\subsection{1. Actual sector and theorem}\label{actual-sector-and-theorem-6}

Keep the unchanged J56 all-top kernel, H=133, f=168, original field and CM involution, all finite lattices and projectors, selected real module and packet point, central phase, class masses and probability Haar measures. Work in the remaining DO105 sector: all25 active relative T places have type II, global rank-two unipotence, and the full-positive compression from R071/DO099. In the original fixed coordinates, \[
 E^{169}=Z\perp P,\quad Z=S_{164}\perp V_3,\quad
 P=Eu\perp Ee_n,\quad h_{ZZ}=I,\quad C=(1,\eta)^t r,
 \quad\eta\bar\eta=U=77Q^2-14.                         \tag{IP1}
\] The whole C is nonzero on S. Let D=D\_eta be the exact E-denominator ideal, x=N\_E D, and k the number of nonzero original S entries of r. The current middle sector has \[
 2^{120}<x<2^{1820},\qquad22\le k\le164.                \tag{IP2}
\] All25 active places are eligible. DO105 gives at most one eta per fixed active support and at most one actual diagonal label at each active place. The original common G pattern and every real and finite factor are retained below.

An inactive relative T place is a \textbf{paired zero} when both original cross blocks C, at its two CM components, vanish after reduction. Let b count such inactive places. This is a condition in the original integral lattices, not on a divided moving frame.

\textbf{Theorem.} The total absolute normalized contribution of the union of these subsets of IP2 is less than \(2^{-392}\):

\begin{enumerate}
\def\labelenumi{\arabic{enumi}.}
\tightlist
\item
  b\textgreater0, with every22≤k≤164 and every x in IP2;
\item
  b=0 and \(x\ge2^{744}\), with every22≤k≤164;
\item
  b=0 and \(x<2^{744}\), with22≤k≤99.
\end{enumerate}

Intersect these sets with the current remainder before removal, and retain every earlier error. Any surviving member of this middle sector consequently satisfies \[
 b=0,\qquad |\mathcal S(\eta)|=25,\qquad
 2^{120}<x<2^{744},\qquad100\le k\le164.                 \tag{IP3}
\] Here S(eta) counts total eligibility among all168 selected T primes. This is a strict subset reduction, with no full support25 removal, new complete floor, signed positivity, prescribed period or source constructor.

\subsection{2. The literal inactive incidence coefficient and local implication}\label{the-literal-inactive-incidence-coefficient-and-local-implication}

At an inactive T place h and its inverse preserve the original Lambda. The actual pencil of K062(AT3--6) is determined by \[
 \mathcal M=\{v\in\Lambda:q_*v\in\pi\mathcal O\},\qquad
 \mathcal N=\Lambda+\pi^{-1}\mathcal O z.                 \tag{IP4}
\] The original pole line z lies in P. Modulo pi, q\_* annihilates Z and is the nonzero annihilator of z in P. This follows from the literal row q\emph{=−v0}J−pi m in DO105(EL9--15), which preserves the actual labels. The fixed Z/P frame has determinant supported on2U, a T-unit. Its use is a local unit-coordinate change, not a globally integral moving frame.

The standard lattice is the \textbf{included label c=0}; the omitted labels are infinity and−1. If h does not preserve the full pencil, K062(AT7--8) gives at most one actual incidence, and Lambda is already fixed. Therefore F(h)=C(h)=1. With m=T−1 the literal coefficient is \[
 \rho(h)=\frac{F(h)-C(h)/m}{T-2}
        =\frac{1-1/m}{m-1}=\frac1{T-1}.                 \tag{IP5}
\] This is the value of the original averaged projector, with no new inactive label enumeration. The opposite CM lattice is fixed by inverse-adjointness and supplies no independent incidence factor.

Suppose now that this inactive place is noneligible. In its local P lattice choose a primitive generator w0 of the saturated null line and a primitive row b0 annihilating w0. They are locally scaled versions of w=(u,eta) and wdagger=(u*,-bar(eta)). This saturation is allowed when eta has a zero or pole. The residue lines w0 and z are different: for unit eta equality is precisely the selected eligibility congruence eta≡a0; for a zero or pole the saturated line has one zero coordinate whereas z has two unit coordinates. Thus \[
                       b_0z\ne0,\qquad q_*w_0\ne0
                       \quad\pmod\pi.                  \tag{IP6}
\] The original rank-two unipotent normal form, after local saturation, is \[
 N=h-I=w_0r_0-v_0b_0+d_0w_0b_0,\qquad
 r_0|_P=0,\quad v_0\in Z.                               \tag{IP7}
\] Its coefficients are integral because N is integral and the extracting column and row are primitive. The selected C reduction vanishes exactly when r0 vanishes; the companion C reduction vanishes exactly when v0 vanishes, by the actual inverse-adjoint relation and N²Z=0. This uses a global algebraic identity of this unipotent normal form, not a positive residue-field form.

If h preserves the full pencil, its reduction preserves the pole line z and row-kernel q\_\emph{. Since that reduction is unipotent, its eigenvalues on the line and dual quotient are1. Hence Nz=0 and q\_}N=0 modulo pi. But \[
 Nz=(-v_0+d_0w_0)(b_0z)=0
       \quad\Longrightarrow\quad v_0=d_0=0,
\] because Z and P are complementary. The other equation becomes (q\_*w0)r0=0, so r0=0. Full-pencil preservation therefore forces N=0 modulo pi. In particular every noneligible inactive place that is not a paired zero contributes exactly IP5. This includes one-zero paired cross blocks. Applying the actual paired adjoint gives the same conclusion at the other CM component, not a squared coefficient.

\subsection{3. Total eligibility and simultaneous denominator-saturated content}\label{total-eligibility-and-simultaneous-denominator-saturated-content}

The total eligible set consists of selected E primes over pi\_T where eta−a0 is divisible, with a0=Q(−7+2sqrt(−7)). The denominator-cleared ideal is integral. DO103(LD13--15)'s original7-adic argument contributes7\^{}84, so the same proof over the full eligible set gives K099(GE4) \[
 7^{84}T^q\le[4(U+14)]^{168}x,\qquad q=|\mathcal S(\eta)|. \tag{IP8}
\] Exact original-host comparisons give q≤26 throughout IP2, and q26 forces x\textgreater2\^{}1116. Thus q25/26 leaves at least142 noneligible inactive places, or exactly143 when q25. In particular x\textless2\^{}744 forces q25. No Galois symmetry or new eta multiplicity is used.

At any inactive paired zero, for every one of the167 fixed original Z inputs y, both C(y)=w r(y) and its companion are divisible. Primitive integer u extracts r(y),eta r(y) in pi O. At a selected prime with a=v\_P(eta), \[
                   v_P(r(y))\ge1+\max(0,-a).             \tag{IP9}
\] The same statement holds at the partner. It retains the denominator ideal even when eta has a local pole. If J is the canonical whole-C G ideal and I\_Sigma is the inverse paired active-T ideal, every one of the167 values therefore lies in \[
 L_v=J I_\Sigma D\,(\mathfrak q_v\mathcal O_E),\qquad
 N_E(\mathfrak q_v\mathcal O_E)=T^2.                     \tag{IP10}
\] The G support and this additional ideal are disjoint. No principal ideal, inverse minor or moving global basis is chosen. Choose the first paired-zero place in any fixed ordering; bounding all choices costs at most143\textless2\^{}8.

For completeness choose one nonzero original S coordinate xi. At every embedding \textbar xi\textbar²≤R/U, R=17424/133. IP9 at all b places gives \[
                xT^{2b-50}\le|N_E(\xi)|\le(R/U)^{168}.   \tag{IP11}
\] Exact inequalities (R/U)\^{}168 T\^{}22\textless2\^{}120 and (R/U)\^{}168 T\^{}24\textless2\^{}1116 imply b≤13, and b≤12 if q26. Thus every matrix already has at least130 inactive IP5 factors. The stronger split below uses IP10 when b\textgreater0 and the full142/143 factors when b=0.

\subsection{4. Complete row, central, real and finite factors}\label{complete-row-central-real-and-finite-factors}

Put M=sqrt(R/U), epsilon=1/1024 and z=(N\_E(J I\_Sigma D))\^{}(1/336). Retain the entire DO105(EL17--23) bound: \[
\begin{split}
 \#r&\le K_k A_3(4U)^f\mathcal B_k
       (M+\epsilon z)^{2(k+3)f}z^{-2(k+3)f},\\
 K_k&={164\choose k}47^k2^{2fk}(f!)^k(87/(32fk))^{fk},\\
 \mathcal B_k&=(fk)!(3f)!/((k+3)f)!,\qquad
 A_3=528^{3f}/((3f)!D_0^3)<2^{-4040},\\
 \#d&\le V_F(H/2+750Uz^2)U^{-f}T^{50}x^{-1},\\
 V_F(C)&=((22/7)fC^2)^{f/2}/((f/2)!\sqrt{D_F}),\\
 |c(h)|&<2^{-929f/2}U^{-f/2}z^{-f}.
\end{split}                                                \tag{IP12}
\] These are the actual joint S/transverse beta integral, Gram determinant4U and arbitrary-ideal cover, one nonempty affine central coset with full real750 normalization, and genuine selected trace. Eta,r,d reconstruct h. The old harmless pivot164\textless2\^{}8 remains; no source-isometry or projected-column fibre is newly discarded.

Keep Green2Q\^{}(-335·167)(N\_EJ)\^{}(335/2), all4\^{}334 common patterns and the single support2\^{}168 factor. In zero-free cases combine the active and inactive coefficients before comparison: \[
       (T-1)^{-(25+g)}<2T^{-(25+g)},\qquad g\le143.       \tag{IP13}
\] The exact inequality2(T−1)\^{}168\textgreater T\^{}168 proves this simultaneously. The existing single active-weight factor2 still suffices; no extra factor2 or inactive label count is hidden.

\textbf{Case A: b\textgreater0.} Use z\_v=(N\_E L\_v)\^{}(1/336)=zT\^{}(1/f). The row and trace factors in IP12 apply with z\_v, since all167 values lie in L\_v. Exactly \[
 (N_EJ)^{335/2}
       =z_v^{335f}T^{335\cdot25}x^{-335/2}T^{-335}.       \tag{IP14}
\] The central affine ideal itself is unchanged. Its cover uses the old z=z\_v/T\^{}(1/f). Since the nonzero-coordinate minimum gives z\_v≤M/sqrt(k), T²\textgreater3000\^{}f and k≥22 imply \[
                 H/2+750Uz^2
        \le H/2+750R/(kT^{2/f})<68.                      \tag{IP15}
\] This is a proved affine covering estimate, not a density-only gain. The global cap z\_v≤z\_G x\^{}(1/336)T\^{}(1/f) also holds; discarding it is safe. The net z\_v power is (328−2k)f≥0, so the nonzero-coordinate cap may be substituted. The remaining x power is x\^{}(-337/2), bounded at2\^{}120. Other inactive coefficients are bounded by1 in this branch.

\textbf{Cases B and C: b=0.} Use g=142 for the complete high-denominator window; use g=143 below2\^{}744, where q25. All other factors remain IP12. The original cap z≤min(M/sqrt(k),z\_G x\^{}(1/336)) and DO105's complete monotonicity apply: after the three cap-dependent factors the net x comparison exponent is−1/2. Each whole window has one endpoint, without a denominator, cap-branch or eta union.

\subsection{5. Exact endpoints and the complete union}\label{exact-endpoints-and-the-complete-union}

Use the unchanged directed bounds \[
 L_T=1019333/1024>\log_2T>L_T^-=1019332/1024,\quad
 L_Q=287843/4096<\log_2Q,\quad
 L_U=601356/4096<\log_2U,\quad L_R=28810/4096>\log_2R.
                                                               \tag{IP16}
\] Regroup the positive total R power before substitution. Let c\_k(A) be the least integer strict binary ceiling of \[
 \frac{{164\choose k}47^k2^{2fk}(f!)^k87^{fk}(fk)!(3f)!}
 {32^{fk}f^{fk}k^{fk}A^{(164-k)f}((k+3)f)!}.             \tag{IP17}
\] Let j\_k be the strict binary ceiling of (1+1/(1024 floor(sqrt(k))))\^{}(2(k+3)f). The original complete common constant C\_*=1618188837/4096 includes Green2, support2\^{}168, patterns2\^{}668, pivot2\^{}8, active-weight2 and the original one-bit transverse inflation.

For A the full exponent, including the chosen inactive place, is \[
 E^A_k=C_*+479+c_k(k)+(j_k-1)+8-335L_T^-.                \tag{IP18}
\] All143 exact cells give \[
          \max E^A_k=E^A_{22}=-226452427/4096<-55286.      \tag{IP19}
\] For B, at X0=2\^{}744 the exact host comparison is \[
 5183^{168}U^{168}Q^{334}2^{744}
       <R^{168}T^{50}\le5184^{168}U^{168}Q^{334}2^{744}.   \tag{IP20}
\] Thus z\textless M/sqrt(5183)\textless M/71, central radius\textless86, full transverse inflation\textless2² and V\_F(86)\textless2\^{}536. The full exponent is \[
 E^B_k=C_*-(337/2)(744-120)+1+57
                      +479+c_k(5183)-142L_T^-.           \tag{IP21}
\] Its maximum is E\^{}B\_164=−1642683/4096\textless−401. For C, use the original120 endpoint,68036 cap and68 central radius: \[
 E^C_k=C_*+479+c_k(68036)-143L_T^-.                       \tag{IP22}
\] On22≤k≤99 its maximum is E\^{}C\_99=−4818891/4096\textless−1176. The adjacent E\^{}C\_100=2054197/4096\textgreater0 is only failure of this sufficient estimate.

There are143 cells in each of A and B and78 in C. They are disjoint, and further intersections with earlier removed sets only decrease the absolute bound. The deliberately conservative union is \[
       143\,2^{-55286}+143\,2^{-401}+78\,2^{-1176}<2^{-392}.
                                                               \tag{IP23}
\] Each cell already contains support, common pattern, coordinate positions, original pivot and, in A, chosen inactive-place factors. No additional pair, eta, central fibre or endpoint multiplicity is required.

\subsection{6. Verification, chronology and retained geometry}\label{verification-chronology-and-retained-geometry}

The companion verifies the original host, every directed power and cap bracket, the unique-incidence projector value, finite paired split models including one-zero cross blocks and saturated endpoint lines, exact denominator-content valuations, every complete endpoint, the union and both full ordered error lists. Finite models diagnose the general local algebra; they are not globally admissible short matrices.

Direct inputs are the complete K062(AT3--16), K071(NB1--9), DO105(EL1--29), DO104(CD10--22), DO106(CFI1--10,18--21) and K099(GE1--5). Full-Z integral inputs use the actual fixed vectors in both modules, not NB5 alone. New steps are the noneligible inactive-pencil implication, its literal rho value, denominator-saturated paired congruence, the extra-ideal central/trace distinction and complete three-branch count.

The script defaults to the current K101 predecessor and accepts --prior for unchanged-byte chronological replay. It checks the full DO105 prefix, retains K100/K101 when integrated, appends392 once and verifies both complete sums below2\^{}-230. Complete floors remain25/11. Historical subject bytes are not altered if a later charge intervenes.

All ideals are the original invertible modules over the accepted CM integer rings with the same involution and Hermitian polarization. Local unit frames change neither the global theta line and biextension nor the selected module and complex metric. Carried contextual anchors remain SGA3 I4.7.3/4.7.3.1; SGA6 VI Theorem3.3; FGA232 Theorem3.1; FGA236 Corollaries1.7/2.3; SGA7 I VII2.7/2.9 and VIII7.4(b); cote4 scans003/005--006 and cote54 scans005--007. No new lambda-ring, l-adic, motivic, selected arithmetic-image, prescribed Bessel/Mram, literal source, DP44 or RH conclusion follows. Residual IP3 and all unrelated sectors remain open.

\subsection{Canonical algebraic-geometric context}\label{canonical-algebraic-geometric-context-11}

The original fractional ideals and affine fibres are invertible-module conditions on Spec(O\_E)/Spec(O\_F), with the same CM involution, Hermitian polarization, theta line, biextension and complex metric. Exact inherited contextual anchors remain SGA3 I4.7.3/4.7.3.1; SGA6 VI Theorem3.3; FGA232 Theorem3.1 pp148--149; FGA236 Corollaries1.7/2.3 pp229/232; SGA7 I VII2.7/2.9 pp169--174 and VIII7.4(b) pp310--311; cote4 scan003 sections8.1.1--8.1.2 and scans005--006 sections8.4.1--8.4.3; cote54 scans005--007. No new lambda-ring/K-theory equality, l-adic or motivic comparison is inferred. DP\_MULT Q3e/DP44/DP52/DP54--65, THETA3 G1/A7, CX\_GRH CX-ALG-024 and GCT\_PLETH T-052 retain their actual arithmetic and global comparison gates.

\clearpage

\section{AM-DO-108 --- persistent active cross poles and the exact constant-T central fibre}\label{am-do-108-persistent-active-cross-poles-and-the-exact-constant-t-central-fibre}

\textbf{Canonical checked integration --- 2026-11-23T02:35:46.672156+00:00.} Complete root and independent proof/script/interface reviews and separate exact replays passed. The actual canonical predecessor has100/31 entries; this integration gives101/31, complete floors25/11. Frozen author history is preserved. AM-DO-108 AC1--15 removes the entire actual support25 fully eligible all-II globally rank-two-unipotent scalar-cross subset with any active paired-C-integral place, at cost less than2\^{}-55281. The exact extra normT² row ideal gives T\^{}-335; the central cover uses the old scale. Its143 complete cells append55281 to actual DO107100/31, giving101/31, all historical charges retained and both sums below2\^{}-230, floors25/11 unchanged. Remaining middle data have j=b=0, exactly25 total eligible places,2\^{}120\textless x\textless2\^{}744 and100\textless=k\textless=164. Every complete actual central orbit then keeps all168 original T factors equal to1/(T-1), with finite weight(T-1)\^{}-168 a(h(t)), sufficient periodQGamma, indexQ\^{}168, exact trace dual and ordinary Poisson formula AC15. This strictly removes active-integral profiles and removes their indicator from the exact Fourier interface; it gives no alias estimate or automatic short-sum/fibre replacement. Admissible source, selected arithmetic content, signed nonvanishing, global period, specific032, DP44 and RH remain open.

Author: /root/definite\_oscillator, ALP\_MACHINE. 23 November 2026 UTC.

Controlled pending proof. The finite subset estimate and the complete-fibre arithmetic identity below are distinct conclusions. Author saved-file checks, independent review and canonical integration remain separate events. DO105--107 and all older frozen files are unchanged.

\subsection{1. Exact sector and new removal}\label{exact-sector-and-new-removal}

Keep the actual J56 all-top kernel with f=168, H=133 and support Sigma of size25, all original fields, paired finite pencils, omitted labels, projectors, selected holomorphic module and packet point, central phase, class masses and probability Haar measures. The global rank-two unipotent scalar-cross sector is DO105(EL1--3): \[
 E^{169}=Z\perp P,\quad Z=S_{164}\perp V_3,\quad
 P=Eu\perp Ee_n,\quad h_{ZZ}=I,\quad C=(1,\eta)^t r,\quad
 \eta\bar\eta=U,
 \quad2^{120}<x=N_ED_\eta<2^{1820},\quad22\le k\le164.
                                                               \tag{AC1}
\] The C restriction to S is nonzero. All25 active places are eligible, and all have actual type II. The fixed transverse frame and original primitive integer u are those of DO104--107, with determinant supported on2U, a T-unit.

Call an active place \textbf{cross-integral} if the whole C block is integral at both CM components. This includes all167 fixed original Z inputs, not just the164 S inputs. It does not mean that h is integral: a central-only pole can still make this an active type-II place. Let j be the number of such active places, exactly as in DO106(CFI9--13).

\textbf{Finite-sector theorem.} The complete subset of AC1 with j≥1 has absolute normalized contribution less than \(2^{-55281}\). Intersect it with the current remainder and append this fresh charge once, retaining every earlier charge.

Combined with DO107, the remaining middle sector has \[
 j=0,\quad b=0,\quad|\mathcal S(\eta)|=25,\quad
 2^{120}<x<2^{744},\quad100\le k\le164.                  \tag{AC2}
\] Here b counts inactive places where both C reductions vanish. This does not remove the entire support25 sector or change either complete support floor.

\subsection{2. Exact extra ideal and complete count}\label{exact-extra-ideal-and-complete-count}

At every active eligible place, eta is a unit at both partners: the selected residue is a0, and the norm equation forces the partner residue U/bar(a0). Therefore D\_eta is a unit there. If the whole C is integral at both components, primitive integer u applied to C(y)=w r(y), w=(u,eta), gives r(y) integral at both components for every one of the167 fixed inputs. Equivalently, the original inverse active-prime denominator can be removed from all these coordinates simultaneously.

With the same canonical whole-C G ideal J and original paired active ideal I\_Sigma, this is \[
 r(y)\in L_v=J I_\Sigma D_\eta(\mathfrak q_v\mathcal O_E),
 \quad N_E(\mathfrak q_v\mathcal O_E)=T^2
 \qquad\text{for all167 inputs }y.                     \tag{AC3}
\] No new integral moving basis or division by U is used. The G and T factors have disjoint supports. This is precisely the ideal inclusion in DO107(IP10), now obtained by removing a pole at an active place rather than forcing a zero at an inactive place. Choosing the first cross-integral place in the fixed active ordering costs at most25\textless2\^{}5, instead of the143\textless2\^{}8 bound there.

This is different from DO105(EL4--5): EL5 removed at most24 eligible places and controlled full double poles as well. A cross-integral place here can remain eligible and can retain a genuine central-only double pole. We make no claim of an improved central denominator or smaller central affine lattice.

Put z=(N\_E(JI\_Sigma D\_eta))\^{}(1/336), z\_v=zT\^{}(1/f) and M=sqrt(R/U), R=17424/133. The full DO105 row and selected-trace counts apply with z\_v by AC3. The central affine difference coset remains \(iJ_0B_\eta I_{\Sigma,F}^2\), so its cover uses the \textbf{old} z, not z\_v. Exactly as in DO107(IP12--15), \[
\begin{split}
 (N_EJ)^{335/2}
   &=z_v^{335f}T^{335\cdot25}x^{-335/2}T^{-335},\\
 z_v&\le M/\sqrt{k},\\
 H/2+750Uz^2
   &\le H/2+750R/(kT^{2/f})<68\qquad(k\ge22).
\end{split}                                                \tag{AC4}
\] The last step is the exact T²\textgreater3000\^{}f comparison. The original full-real750 cover, one affine coset, covolume and shift are retained. No density-only assertion is made.

Keep all common G patterns4\^{}334, support2\^{}168, coordinate-position factors, original pivot164\textless2\^{}8, actual Green2, unique active-label weight and genuine near-point selected trace. Inactive coefficients are bounded by1 for this finite estimate. The original all-word and role bounds, central quotient, class masses and Haar factors are unchanged. Eta,r,d reconstruct h, as in DO104/105. There is no newly discarded isometry fibre.

The entire IP18 formula is therefore valid with its chosen-place exponent8 replaced by5. With C\_*=1618188837/4096, L\_T\^{}-=1019332/1024, c\_k(k) and j\_k exactly DO107(IP17--18), the full bound is \[
 E^{\rm act}_k
 =C_*+479+c_k(k)+(j_k-1)+5-335L_T^-
 =E^{A,{\rm DO107}}_k-3.                               \tag{AC5}
\] The source formula and every factorial ceiling are reread and replayed, not reconstructed from a threshold table. Thus \[
 \max_{22\le k\le164}E^{\rm act}_k
 =E^{\rm act}_{22}=-226464715/4096<-55289.
                                                               \tag{AC6}
\] There are143 coordinate cells. Their complete union, with the first-active-place choice already inside each cell, satisfies \[
                  143\,2^{-55289}<2^{-55281}.            \tag{AC7}
\] This proves the finite-sector theorem with one fresh55281 charge. A later application may restrict to AC2's narrower k-range, but the whole containing AC1 subset is what is charged here.

\subsection{3. Complete central fibres retain the original support and all-II type}\label{complete-central-fibres-retain-the-original-support-and-all-ii-type}

Now take any retained matrix satisfying AC2 and fix its eta,r. Use the original ambient normal form \[
 w=(u,\eta),\quad w^\dagger=(u^*,-\bar\eta),\quad
 N(d)=wR_0-V_0w^\dagger+dww^\dagger,\quad
 d+\bar d=-q,\quad q=R_0V_0>0.
                                                               \tag{AC8}
\] Let d0 be the scalar of the chosen matrix. DO106(CFI3--8) proves, at every finite place including all rationally ramified places, that \[
\begin{split}
 \Gamma&=B_\eta I_{\Sigma,F}^2,\qquad
       D_\eta\bar D_\eta=B_\eta\mathcal O_E,\\
 h(t)&=h_0n(t)=n(t)h_0=h(d_0+i_{\rm fld}t),\\
 n(t)&=I+i_{\rm fld}t\,ww^\dagger,\qquad t\in\Gamma,
\end{split}                                                \tag{AC9}
\] is exactly the complete central family preserving the chosen active lattices and the original standard lattices elsewhere. Its full real covolume is sqrt(D\_F)xT\^{}-50. The null line and all cross data are fixed; N²=−qwwdagger is fixed and nonzero. Every h(t) remains globally rank-two unipotent, and distinct t remain distinct modulo mu28 because h(t)\_ZZ=I.

At each active place j=0 means that at least one of the two original C blocks has a pole. It persists for every t since C is fixed. A full-pencil-preserving unipotent has both whole C blocks integral: every fixed original Z input lies in both actual modules, and h\_ZZ=I gives C(y)=hy−y in the integral module. The paired inverse-adjoint gives the companion assertion. This is the full-Z bridge already checked in DO106, not a new unproved extension of NB5's S-only wording. Thus every h(t) has unequal full pencils there. It preserves the chosen actual lattice by AC9, so it has exactly that unique fixed incidence and remains active type II, with rho=1/(T−1). This includes arbitrarily large real t; it is a finite-source assertion, not an H133 bound.

There are exactly143 inactive places because total eligibility is25 and all25 active places are eligible. They are all noneligible. At each one b=0 means at least one paired C reduction is nonzero. These reductions are fixed along AC9. Every h(t) preserves Lambda there, and DO107(IP4--7) forces unequal full pencils and the unique included Lambda incidence. Again rho=1/(T−1) for all t.

Consequently the \textbf{exact original finite T factor} is \[
               \prod_{v\mid T}\rho_v(h(t))
                       =(T-1)^{-168}\qquad(t\in\Gamma).  \tag{AC10}
\] There is no remaining active all-II indicator or inactive character phase. No omitted pencil point is inserted. Both CM components have been used through the actual paired modules, not through two independent rho factors. The same Sigma, type profile, k and eta are preserved throughout this complete orbit.

\subsection{4. Smaller exact finite period and ordinary Poisson identity}\label{smaller-exact-finite-period-and-ordinary-poisson-identity}

Write a(h(t)) for K062(AT14)'s actual normalized G trace, with both mixed projectors retained. All other away-T conditions hold by AC9. The full finite weight on this orbit is exactly \[
                         F(t)=(T-1)^{-168}a(h(t)).       \tag{AC11}
\] A sufficient period is now \textbf{QGamma}, not the larger sufficient period used before removing the indicators. At every G place the original content of wwd is B\_eta\^{}-1. Since I\_Sigma is a G-unit, t∈QGamma makes n(t) principal modulo the actual Q. Its action is therefore the identity in the entire original finite representation, including the mixed constituents. Hence a(h(t+v))=a(h(t)) for v∈QGamma.

At T no preservation of every pencil lattice by n(QGamma) is asserted or needed. The exact original coefficient AC10 is already constant on the whole Gamma orbit, and every translated point has its required chosen incidence. At all other primes the original compact condition is unchanged. Thus this is a period of the actual whole finite weight, without extending its support.

Since Gamma is an invertible F-ideal, \[
 [\Gamma:Q\Gamma]=Q^{168},\qquad
 \Gamma^*=\mathfrak D_F^{-1}B_\eta^{-1}I_{\Sigma,F}^{-2},
 \quad (Q\Gamma)^*=Q^{-1}\Gamma^*.                     \tag{AC12}
\] The dual is for the full real trace pairing. There is no factor2. The old T-dependent period index is not inserted.

Expand on this exact finite quotient \[
 a(h(t))=\sum_\chi a_\chi e^{2\pi i\operatorname{Tr}_F(\theta_\chi t)},
 \qquad\theta_\chi\in Q^{-1}\Gamma^*/\Gamma^*.
                                                               \tag{AC13}
\] DO106(CFI16--20)'s positive spectral-projection argument and original Green table apply unchanged. If J\_C is the center-independent whole-C pattern ideal, they give \[
 \sum_\chi|a_\chi|
       \le2Q^{-335\cdot167}(N_EJ_C)^{335/2}.             \tag{AC14}
\] The Chinese remainder map for the invertible ideal Gamma identifies its residues modulo Q with the actual G residue parameters. Both mixed projectors are retained and need not commute with the center. No quotient-cardinality factor accompanies AC14.

Let c(t) be the genuine selected coefficient of h0 n(t), with its actual near-point conjugation and all embedding signs. R085(RF9--31) supplies its exact rational/Fourier model and ordinary absolute Poisson identity. In the parameterization AC9 the base point is t=0, so \[
 \sum_{t\in\Gamma} F(t)c(t)
 =\frac{T^{50}}{\sqrt{D_F}\,x}(T-1)^{-168}
       \sum_\chi a_\chi
       \sum_{\xi\in\Gamma^*-\theta_\chi}\widehat c(\xi).
                                                               \tag{AC15}
\] All actual selected vectors, phases and Fourier signs remain in c and its transform. This is an exact identity, not an estimate of the shifted alias sum. No smallness follows merely from its positive-frequency support or from the smaller period.

\subsection{5. Accounting, verification and scope}\label{accounting-verification-and-scope}

The finite estimate AC7 concerns the short H133 subset and is the only new numerical charge. AC9--15 describe complete actual central orbits. They do not permit an automatic replacement of the currently retained short sum by a union of complete fibres. Such an application must partition the fixed eta,r data, reconcile previously removed inside-cutoff subsets, and use the already proved complete outside-H tail at most once across disjoint fibres. No tail is charged separately for each orbit. These statements neither reinsert earlier removed terms silently nor assert cancellation of them.

The companion replays every143 complete AC5 factorial cell, the exact first-place multiplicity and union, active paired valuation identities, original T/G unit constants and period valuation inclusions, and both full chronological error arrays. It records current DO107 as predecessor100/31, appends55281 once to give101/31, and preserves complete floors25/11. Every old charge remains, including the weaker historical subsets. The unchanged actual full error sums must remain below2\^{}-230.

Exact dependencies are DO105(EL1--29), DO104(CD10--22), DO106(CFI1--24), DO107(IP1--23), K062(AT3--16), K071(NB1--9), R025(R25.1--5) and R085(RF1--31). Their exact source/module, finite and metric conventions are retained. No new family, arithmetic image, source constructor, signed lower bound, prescribed Bessel/Mram period or DP44 conclusion is inferred.

The algebraic-geometric setting is the same CM integer rings, invertible ideals, involution, Hermitian polarization, theta line and Poincare biextension, with the same selected Schur module and complex metric. The carried anchors are SGA3 I4.7.3/4.7.3.1, SGA6 VI Theorem3.3, FGA232 Theorem3.1, FGA236 Corollaries1.7/2.3, SGA7 I VII2.7/2.9 and VIII7.4(b), cote4 scans003/005--006 and cote54 scans005--007. No additional K-theory, l-adic, motivic, period, selected arithmetic-content or RH assertion is made.

\subsection{Canonical algebraic-geometric context}\label{canonical-algebraic-geometric-context-12}

The original fractional ideals and affine fibres are invertible-module conditions on Spec(O\_E)/Spec(O\_F), with the same CM involution, Hermitian polarization, theta line, biextension and complex metric. Exact inherited contextual anchors remain SGA3 I4.7.3/4.7.3.1; SGA6 VI Theorem3.3; FGA232 Theorem3.1 pp148--149; FGA236 Corollaries1.7/2.3 pp229/232; SGA7 I VII2.7/2.9 pp169--174 and VIII7.4(b) pp310--311; cote4 scan003 sections8.1.1--8.1.2 and scans005--006 sections8.4.1--8.4.3; cote54 scans005--007. No new lambda-ring/K-theory equality, l-adic or motivic comparison is inferred. DP\_MULT Q3e/DP44/DP52/DP54--65, THETA3 G1/A7, CX\_GRH CX-ALG-024 and GCT\_PLETH T-052 retain their actual arithmetic and global comparison gates.

\clearpage

\section{AM-LIFT-084 --- Mixed signature principal points and the literal finite Tate level}\label{am-lift-084-mixed-signature-principal-points-and-the-literal-finite-tate-level}

\textbf{Canonical checked integration --- 2026-11-23T02:37:33.914026+00:00.} Complete saved-proof and executable reviews and separate exact diagnostic replays passed. No numerical error charge is introduced. Frozen author history is preserved. AM-LIFT-084 MP1--25 constructs a literal field-valued Obj and FrPt for any given actual source Host, compatible Par.PinHost/PinLevel/PinDat, specified stage Datum, existing Moduli representation and same Coeff field, after a finite extension. The actual principal ideal identity is derived from Host fields. A pure principally polarized CM block and its conjugate action supply every uniform(p,q); the explicit trace-one semicharacter supplies actual SAB prim/adjoint/cubical lines. Relative-unramified selfdual Hermitian lattice uniqueness, including dyadic places, gives the required profinite Tate isometry; the literal commutator sign and CRT preserve the exact finite level M=g0g1,A=1. All-a normalized line identities descend over one finite field; only geometric full Tate systems are asserted. Separately, on the existing E\_h fields the explicit pure block and three J-action twists give J\^{}(fq) times C\_h\^{}(p-2q) for p\textgreater=2q, attaining LIFT083 complement dimension f(p-2q). Those fixed fields still fail full Sizes. No effective field/frame/line costs, good-model/old LamPt identification, prescribed same-point F nonzero, actual full Host/Sizes constructor, selected arithmetic content, prescribed period, specific032, DP44 or RH is supplied. No numerical error charge;101/31 arrays and floors25/11 remain.

Author: /root/lift\_accounting. 23 November 2026 UTC.

\textbf{Author whole-proof, literal-interface, and complete-script check passed; independent review and canonical promotion pending.} This construction is conditional on an actual source Host, the stated Pins, and an already chosen moduli representation. The principal-ideal identity is proved below from the Host fields. The construction does not supply a full SizesAt family or a nonzero value of the prescribed section. Final frozen fingerprints and the exact diagnostic replay are recorded in the adjacent author verification and manifest.

\subsection{1. Two distinct constructions and the source target}\label{two-distinct-constructions-and-the-source-target}

There are two results to keep separate.

\begin{enumerate}
\def\labelenumi{\arabic{enumi}.}
\tightlist
\item
  Given an actual source Host and its principal fractional ideal, one can construct a principally polarized CM variety with every desired uniform signature \((p,q)\), using one pure-type block and the conjugate \textbf{action} on copies of that same block.
\item
  On the existing fields \(E_h\), \(h=1,4\), a different construction realizes the maximal permitted \(J\)-isotypic part in LIFT083: three twists of the known Jacobian block, together with pure blocks. This gives \(J^{fq}\) when \(p\ge2q\), with complement of dimension \(f(p-2q)\).
\end{enumerate}

Neither statement identifies the new point with the earlier LamPt, and neither implies that the old LamObjGood predicate holds.

Here is the exact quantifier and interface distinction. In the first result fix an actual \nolinkurl{Formalism.BasicNotions.Host}, a compatible \nolinkurl{Par} with its \nolinkurl{PinHost}, \nolinkurl{PinLevel}, and \nolinkurl{PinDat} clauses, and one stage in the literal \nolinkurl{InR} range. Keep its \nolinkurl{Dw}, \nolinkurl{Xo}, and \nolinkurl{c} unchanged. Then a finite extension of the \textbf{same} coefficient field admits an \nolinkurl{Obj} and a \nolinkurl{FrPt} at the specified stage. The proof is stagewise and also gives a common finite extension for any finite list of stages by taking their compositum. It does not construct the \nolinkurl{Par}, its moduli schemes, or the numerical Host/Sizes hypotheses from scratch.

\begin{longtable}[]{@{}
  >{\raggedright\arraybackslash}p{(\columnwidth - 2\tabcolsep) * \real{0.5000}}
  >{\raggedright\arraybackslash}p{(\columnwidth - 2\tabcolsep) * \real{0.5000}}@{}}
\toprule\noalign{}
\begin{minipage}[b]{\linewidth}\raggedright
Literal input
\end{minipage} & \begin{minipage}[b]{\linewidth}\raggedright
Clause retained in the construction
\end{minipage} \\
\midrule\noalign{}
\endhead
\bottomrule\noalign{}
\endlastfoot
\nolinkurl{DataHost.Host}, \nolinkurl{xi0}, \nolinkurl{AA}, \nolinkurl{𝔟_spec} & The actual pure fractional lattice and its integral principal trace pairing. \\
\nolinkurl{Par.PinHost} & The quadratic field element and selected root agree with the given source Host; \nolinkurl{φ₀_mem} makes the selected root \(i\sqrt d\). \\
\nolinkurl{Par.PinDat} / \nolinkurl{IsDatW} & The existing lattice \nolinkurl{Dw.La}, its exact alternating form, and its prescribed \((p,q)\) are retained. \\
\nolinkurl{Par.PinLevel} and \nolinkurl{Level} & \(M=g_0g_1\), \(A=1\), and the existing coprime conjugate factors of \(M\) are retained. \\
\nolinkurl{Moduli.rep} & A reduced field base with the constructed literal object produces a point and its cartesian cubical \nolinkurl{ObjMap}. \\
\nolinkurl{Coeff.conj} and \nolinkurl{FrPt} & All eigenlabels lie in the same coefficient field; only a finite extension for the point, finite level, frame, and section is allowed. \\
\end{longtable}

These are assumptions and preserved identities, not a construction of \nolinkurl{PinField}, an arithmetic \nolinkurl{Read}, or a full admissible source family. The separate explicit \(E_h,d=7\) result in §7 is a geometric construction in those fields and does not assume their excluded full numerical Host data.

Write \(E/F\) for a CM extension, \(K_0=\mathbb Q(\jmath)\), \(\jmath^2=-d\), and \[
 \Psi=\{\tau:E\hookrightarrow\mathbb C:\tau(\jmath)=i\sqrt d\}.
 \tag{MP1}
\] The literal HostRaw.cmType uses exactly this set. For an actual source Host, \[
 \xi=-\frac1{N(\mathfrak c_K)\jmath},\qquad
 \mathcal A=\mathfrak c_E\mathfrak b^{-1},\qquad
 (\xi)\mathcal A\bar{\mathcal A}\mathfrak D_{E/\mathbb Q}=1
 \tag{MP2}
\] is the stated FactPrincipalDatum identity. Here it follows directly from the actual Host fields: the relative different is \(\mathfrak b\bar{\mathfrak b}\), the quadratic different of \(K_0=\mathbb Q(\sqrt{-d})\) is generated by \(\jmath\), and \(\mathfrak c_K\bar{\mathfrak c_K}=(N\mathfrak c_K)\). The tower formula for the different and extension of these ideals give MP2 by cancellation. Thus we use the content of the Host record and the quadratic different formula, not the declaration of FactPrincipalDatum as an unsupported axiom.

Let \(D\) be the already specified source Datum, of rank \(r=p+q\). Its integral alternating form is denoted \(E_D\), and it is unimodular with adjoint involution conjugation. Let the source level have \(A=1\), \(M\ge3\), and conjugate coprime ideals \(\mathfrak M,\bar{\mathfrak M}\) with product \(M\mathcal O_E\). These are exactly the finite level clauses used below. In particular the pinned level \(M=g_0g_1,A=1\) is not enlarged.

\subsection{2. A pure principal block from the actual source ideal}\label{a-pure-principal-block-from-the-actual-source-ideal}

On \(V_\Psi=\mathbb C^\Psi\), let \[
 C=\mathbb C^\Psi/\Psi(\mathcal A).
 \tag{MP3}
\] The lattice is stable under the full integer ring. Put \[
 \psi(x,y)=\operatorname{Tr}_{E/\mathbb Q}(\xi x\bar y),
 \qquad
 \Omega(x,y)=-\psi(x,y)
             =\operatorname{Tr}_{E/\mathbb Q}(\xi\bar x y).
 \tag{MP4}
\] The signs are explicit: Milne's CM notes use \(\psi(x,Jx)>0\), whereas the literal source Datum.dom uses \(\Omega(Jx,x)>0\). These are the same polarization convention expressed with opposite alternating forms. Since \(\operatorname{Im}\tau(\xi)>0\) on \(\Psi\), the Hermitian form \[
 H(v,w)=\Omega(Jv,w)+i\Omega(v,w)
 \tag{MP5}
\] is positive definite and linear in its first variable. The ideal identity MP2 says exactly that MP4 is integral and unimodular. Therefore MP3 is an abelian variety with a principal polarization, not merely a polarizable torus or an isogeny class.

The integer-ring action is by holomorphic lattice maps. Its adjoint for the polarization is complex conjugation, by the trace identity. Its holomorphic \(E\)-type is \(\Psi\), one copy of every selected embedding. The type is induced from \(K_0\), so \(C\) is isogenous to a power of a CM elliptic curve with field \(K_0\); this is a structural description, not an assertion that the principal polarization is a product polarization in an arbitrary elliptic basis.

Use the same polarized variety with action \[
                 a\longmapsto \iota(\bar a)
 \tag{MP6}
\] to obtain \(C^{-}\). Its type is \(\bar\Psi\). It has the same underlying ample principal line and the same Rosati compatibility. Thus \[
                 A_{p,q}=C^{p}\times(C^{-})^{q}
 \tag{MP7}
\] has a literal product principal polarization, full integral \(\mathcal O_E\)-action, and uniform Lie multiplicities \(p\) on \(\Psi\) and \(q\) on \(\bar\Psi\). Its dimension is \(f(p+q)\), \(f=[F:\mathbb Q]\). This works for all nonnegative \(p,q\), including the definite terminal case \(q=0\).

The characteristic polynomial on the tangent space is \[
 \prod_{\tau\in\Psi}(X-\tau(a))^p
 \prod_{\tau\notin\Psi}(X-\tau(a))^q.
 \tag{MP8}
\] This is the actual DetCond condition, not an averaged signature statement. After a base embedding changes, the selected condition changes together with the image of \(\jmath\), so it remains the same algebraic determinant condition. More precisely, in characteristic zero the full \(E\)-action is separable. Its tangent representation is the direct sum of the eigenspaces just counted. Choosing a basis in each eigenspace gives the existence and uniqueness of the coefficients in every literal \nolinkurl{LieMult} clause of \nolinkurl{DetCond}; MP8 is its characteristic-polynomial expression, not a substitute for those clauses.

\subsection{3. The rigidified symmetric line and the literal prim clause}\label{the-rigidified-symmetric-line-and-the-literal-prim-clause}

Let \(w_0=(1+\jmath)/2\), so \(w_0+\bar w_0=1\). The integer-ring action makes \(w_0\) a lattice endomorphism. On any one of the polarized lattices above define \[
                \alpha(\lambda)=(-1)^{\Omega(\lambda,w_0\lambda)}.
 \tag{MP9}
\] The exponent is an integer. Adjointness and alternation give \[
 \Omega(\lambda,w_0\mu)+\Omega(\mu,w_0\lambda)
 =\Omega(\lambda,(w_0-\bar w_0)\mu)
 \equiv\Omega(\lambda,\mu)\pmod2.
 \tag{MP10}
\] Thus \(\alpha\) is an Appell--Humbert semicharacter for \(H\). It takes values in \(\{\pm1\}\), so its line is symmetric; rigidify it at the origin. The standard cubical structure on a rigidified line over an abelian variety supplies the source cubical identities. For precision, the theorem of the cube trivializes the source third-difference line on \(A^3\). Normalize that trivialization at the origin. Its restrictions to each coordinate face are the canonical cancellations of the same line factors, because both normalized trivializations differ only by a constant on \(A^2\). The two swaps and the source four-variable cocycle follow for the same reason on \(A^3\) and \(A^4\): both sides trivialize the identical tensor expression and agree at the origin. This gives the exact face, symmetry, and cocycle fields of the source CubicalStr, and uniqueness makes it compatible with descent.

The source cubeLB for two maps has signs \(+,-,-,+\) on the empty, singleton, and full subsets. Hence for the maps \(1,w_0\) its line is \[
 (1+w_0)^*L\otimes L^{-1}\otimes w_0^*L^{-1}\otimes 0^*L.
 \tag{MP11}
\] Its Hermitian form is \(H\), since \(w_0+\bar w_0=1\). Its semicharacter is \[
 \frac{\alpha((1+w_0)\lambda)}
      {\alpha(\lambda)\alpha(w_0\lambda)}
       =(-1)^{\Omega(\lambda,w_0\lambda)}
       =\alpha(\lambda).
 \tag{MP12}
\] Appell--Humbert therefore gives an isomorphism of the actual lines in MP11 and \(L\). This is the source SAB.prim statement, rather than equality only in the Néron--Severi group.

Similarly, on the product lattice of \(A\times A\), the two lines \[
 B_L(a x,y),\qquad B_L(x,\bar a y)
 \tag{MP13}
\] have equal Hermitian forms and equal semicharacters, because \(\Omega(a\lambda,\mu)=\Omega(\lambda,\bar a\mu)\). They are isomorphic as lines. This supplies the literal SAB.adj clause. The source symmetry includes its rigidification; normalize the symmetric isomorphism at the origin accordingly.

Conjugating the action replaces \(w_0\) by \(\bar w_0=1-w_0\), and MP9 is unchanged. External products multiply the semicharacters. Hence the chosen lines on MP7 satisfy every stated SAB line identity simultaneously.

Because the polarization is principal, its polarization homomorphism has trivial scheme-theoretic kernel. The source Out.KTriv tests exactly this kernel by local triviality of translation differences, including nonreduced test schemes. Thus KTriv follows from principality, not merely from triviality on complex points. Amplitude and properness give the remaining geometric line/abelian clauses.

\subsection{4. A finite local uniqueness lemma, including residue characteristic two}\label{a-finite-local-uniqueness-lemma-including-residue-characteristic-two}

\textbf{Lemma.} Suppose \(E/F\) is unramified at every finite place. Two finite projective \(\mathcal O_E\)-modules of the same rank, with unimodular alternating integral forms whose adjoint involution is conjugation, are isometric after tensoring with \(\widehat{\mathbb Z}\).

\textbf{Proof.} Decompose at each rational prime and then each finite place \(v\) of \(F\). Put \(R=\mathcal O_{F_v}\). First suppose \(S=\mathcal O_{E_v}\) is the ring of an unramified quadratic field extension. Choose an element \(u\in S\) whose residue is outside the residue field of \(R\). Then \(\nu=u-\bar u\) is a unit and \(\bar\nu=-\nu\). This remains valid in residue characteristic two: the nonzero residue \(u-\bar u\) is the finite-field trace of \(u\), and the trace map of the separable quadratic residue extension is surjective. One can choose \(u\) with that trace nonzero.

Choose a generator \(\delta\) of the inverse different of \(F_v/\mathbb Q_\ell\). Since \(E_v/F_v\) is unramified, \(\gamma=\delta\nu\) generates the absolute inverse different of \(E_v/\mathbb Q_\ell\) and is purely imaginary.

For an alternating form \(B\) with the given adjoint relation, trace duality defines a unique Hermitian form \(h\), conjugate-linear in the first variable, by \[
       B(x,a y)=\operatorname{Tr}_{E_v/\mathbb Q_\ell}
                         (\gamma a\,h(x,y)).
 \tag{MP14}
\] The adjoint relation and alternation give the Hermitian symmetry. Since the lattice is an \(S\)-module, testing \(B(x,a y)\) for every \(a\in S\) shows \[
 B(x,L)\subset\mathbb Z_\ell
 \quad\Longleftrightarrow\quad h(x,L)\subset S.
 \tag{MP15}
\] Thus unimodularity of \(B\) is exactly selfduality of the integral Hermitian lattice.

Such a Hermitian lattice has a vector of unit norm. If a diagonal Gram entry is a unit, this is immediate. Otherwise an off-diagonal entry is a unit, because the Gram determinant is a unit. For the corresponding basis vectors, the norm of \(e_i+t e_j\) reduces to a nonzero trace for a suitable residue \(t\). Surjectivity of the finite-field trace supplies that \(t\), including in characteristic two. Split off this unit-norm line by integral orthogonal projection and proceed inductively.

The norm map \(S^\times\to R^\times\) is surjective. At residue level this is the finite-field norm. On every successive principal-unit quotient it is the surjective residue trace; successive correction and completeness give a norm preimage of any unit. Rescaling each diagonal unit therefore makes the lattice the standard identity Hermitian lattice. This proves uniqueness.

At a split place \(S=R\times R\) with factor exchange, take \(\nu=(1,-1)\). Hermitian selfduality is a perfect pairing between the two free \(R\)-modules, and dual bases give the unique standard split lattice. This argument also covers dyadic split places. The local isometries combine over all finite places and all rational primes. No assertion of a global integral isometry is made. \(\square\)

The source Host.unram is precisely the needed relative unramified hypothesis. Absolute ramification is absorbed in \(\delta\); it is not replaced by an absolute unramified assertion.

\subsection{5. Exact TateIso and the finite prescribed level}\label{exact-tateiso-and-the-finite-prescribed-level}

Apply the lemma to the prescribed \((D.La,E_D)\) and the Betti lattice \((\Lambda_A,\Omega)\) of MP7. Obtain an \(\mathcal O_E\)-linear profinite isometry \[
        f:\widehat{D.La}\ \simeq\ \widehat{\Lambda_A}.
 \tag{MP16}
\] For \(n>0\), reduce it modulo \(n\) and identify \(\Lambda_A/n\Lambda_A\) with \(A[n](\mathbb C)\) by \(\lambda\mapsto\lambda/n\). This supplies additive, \(\mathcal O_E\)-linear maps \(\operatorname{ev}_n\), with kernel \(nD.La\), onto all \(n\)-torsion, and satisfying the source tower identity. Define \(\operatorname{ev}_0=0\), for which the unconditional additivity and linearity clauses also hold.

The source ConstGen.thetaLB is \(L^n\otimes t_x^*L^{-n}\). In the Appell--Humbert trivialization with \[
 j_\lambda(z)=\alpha(\lambda)
       \exp\!\left(\pi H(z,\lambda)+\frac\pi2H(\lambda,\lambda)\right),
 \tag{MP17}
\] a generator for \(x=a/n\) is \(f_x(z)=\exp(-\pi nH(z,x))\). Its lattice transformation law is correct because \(n\Omega(x,\lambda)\) is integral. The exact source CommIsOf ratio is \[
 \frac{f_x(z)f_y(z+x)}{f_y(z)f_x(z+y)}
   =\exp(-2\pi i n\Omega(x,y))
   =\exp(-2\pi i E_D(a,b)/n).
 \tag{MP18}
\] Consequently set \(\zeta_n=\exp(-2\pi i/n)\) for \(n>0\), and set \(\zeta_0=1\). The positive-index roots are primitive and satisfy \(\zeta_{nm}^m=\zeta_n\) for \(n,m>0\); these are exactly the source quantifiers. MP18 is the source TateIso.isom sign. Any two generators differ by a scalar on the proper connected abelian variety, so the ratio is independent of their choice, as required by Out.WeilIs.

Now choose \(e\in\mathcal O_E\) with \[
             e\equiv1\pmod{\mathfrak M},\qquad
             e\equiv0\pmod{\bar{\mathfrak M}}.
 \tag{MP19}
\] Define the finite level section \(\eta(a)=\operatorname{ev}_M(ea)\), and set \(\eta_A=0\) because \(A=1\). If \(a\in\mathfrak M D.La\), then \(ea\in M D.La\), so \(\eta(a)=0\). If \(a\in\bar{\mathfrak M}D.La\), then \((e-1)a\in M D.La\), so \(\eta(a)=\operatorname{ev}_M(a)\). These are the exact two source level clauses. No larger modulus or new finite prime is introduced.

\subsection{6. Algebraic descent and what is defined over one finite field}\label{algebraic-descent-and-what-is-defined-over-one-finite-field}

The CM variety \(C\) and its full endomorphism action descend to \(\overline{\mathbb Q}\), and then to some number field. The action is specified by finitely many integer-ring generators; its algebraic relations descend as well. The polarization homomorphism descends by full faithfulness for homomorphisms of abelian varieties over algebraically closed fields. To descend the particular line, first choose an algebraic line in that polarization class. The symmetric lines inducing the fixed polarization form a finite torsor under the two-torsion of the dual abelian variety. Our Appell--Humbert line is symmetric, so its class is one of these algebraic classes. Thus that line, its origin rigidification, and its uniquely normalized cubical structure are all defined over a finite extension. This argument does not claim that an arbitrary transcendental flat line descends.

The structural line isomorphisms are normalized at the origin. Once the two lines and their rigidifications are defined over \(K\), a geometric isomorphism with that normalization is unique: the only units on the proper connected abelian variety, or its square, are constants. Every Galois conjugate is therefore the same normalized isomorphism, so it descends to \(K\). Equivalently, base change for the one-dimensional space of sections of the trivial difference line, together with evaluation at the origin, gives its normalized \(K\)-generator. This applies to \textbf{every} \(a\in\mathcal O_E\) in \nolinkurl{SAB.adj} over the same \(K\), without adjoining infinitely many choices. The literal \nolinkurl{LBIso} clauses ask only for an isomorphism; the normalized one proves them. The determinant polynomial and full eigenspace statement persist under every geometric base change.

After adjoining the finitely many coordinates of the \(M\)-torsion sections in MP19 and any already prescribed coefficient field, these finite level sections descend to one finite extension \(K\). The \(A=1\) root condition is automatic. The entire profinite isometry and all torsion points are \textbf{not} asserted to be defined over \(K\). The source asks for their existence separately over every algebraically closed geometric field, which follows by transporting the construction on \(\overline{\mathbb Q}\) into such a field. Every torsion point in characteristic zero is algebraic over the chosen number field, so the transported maps remain onto all \(n\)-torsion.

The resulting abelian variety with its line, action, and finite sections is therefore an object of the literal Obj moduli problem attached to the same \(D\) and level. The exact Moduli.rep clause in PushInvertible.lean applies to every reduced base scheme and every such object. Apply it to \(\operatorname{Spec}K\), which is reduced. It supplies the point and an ObjMap to the universal family. In particular its cartesian pullback and cubical line compatibility are included; no merely set-theoretic moduli identification is being used.

The already fixed source Coeff field contains all conjugates of \(E\), by Coeff.conj. Take \(K\) to be a finite extension of this same field, rather than replacing the Coeff record. The tangent eigenspaces then have dimensions \(D.mult(\tau)\), so choose algebraic eigenbases over \(K\). A principal ample line has a one-dimensional section space over \(K\); choose a nonzero algebraic generator. Transport through the ObjMap gives every field of the literal FrPt in IsCombAt.lean. This is a qualitative field-valued frame, not an effective integral frame or an assertion of nonvanishing of a prescribed rule at that point.

The unchanged data are the Host, the pinned Datum (including its integral lattice and alternating form), the Coeff field, and the finite Level. The new data are the finite field-valued point, its eigenframe, and its principal-section generator. The extension degree, integral frame costs, and section normalization have not been bounded.

\subsection{7. An explicit pure block and maximal Jacobian content on the existing fields}\label{an-explicit-pure-block-and-maximal-jacobian-content-on-the-existing-fields}

For \(E_0=\mathbb Q(\zeta_{28})=\mathbb Q(\zeta_7,i)\), let \[
 \jmath=\zeta_7+\zeta_7^2+\zeta_7^4-\zeta_7^3-\zeta_7^5-\zeta_7^6,
 \quad b=((1+i)(1-\zeta_7))^{-1},\quad I=b\mathcal O_{E_0}.
 \tag{MP20}
\] Then \(\jmath^2=-7\) and \[
          (-1/\jmath)\,I\bar I=\mathfrak D_{E_0/\mathbb Q}^{-1}
          \quad\hbox{as fractional ideals}.
 \tag{MP21}
\] At each prime over 2, \(1+i\) has valuation one, \(1-\zeta_7\) is a unit, and the different exponent is two. At the unique prime over 7, the valuations of \(\jmath\) and \(1-\zeta_7\) are three and one, giving total exponent \(-3-2=-5\). There are no other valuations. These are ideal equalities, not merely norm equalities. The pure type is all embeddings with \(\jmath\mapsto i\sqrt7\).

Tensor \(I\) with the exact trace-selfdual \(B_7(29)\) lattice \(\pi_B^{-3}\mathcal O_B\) of K030. For \(h=4\), also tensor the weighted selfdual integral bases of \(\mathbb Q(\sqrt5)\) and \(\mathbb Q(\sqrt{13})\) from K032. This gives a principally polarized pure block \(C_h\) with full \(\mathcal O_{E_h}\)-action, dimension \(f=42h\), and lattice \[
 b\pi_B^{-3}\mathcal O_{E_h},\qquad
 \xi_{\rm pure}=(-1/\jmath)\delta_h,\quad
 \delta_1=1,\quad\delta_4=\delta_5\delta_{13}>0.
 \tag{MP22}
\] The tensor trace identities make this an actual product polarization of \(7h\) copies of the six-dimensional pure block, not merely a rational comparison.

Let \(A_{J,h}=J^{14h}\) be the actual principal block from K030/K032, and let \(\sigma\) fix the real auxiliary fields and \(i\), with \(\sigma(\zeta_7)=\zeta_7^2\). It has order three and fixes \(\jmath,w_0\). Twist its full action by \(1,\sigma,\sigma^2\). The three \(\zeta_7\)-type sets are \[
           \{1,2,3\},\qquad\{2,4,6\},\qquad\{4,1,5\}.
 \tag{MP23}
\] Every quadratic-residue label occurs twice and every nonresidue label once. Every other field embedding is carried unchanged and occurs in all its induced lifts. Therefore the product of these three twists has uniform type \((2,1)\) relative to \(K_0\).

For \(p\ge2q\), take \(q\) copies of that triple and \(p-2q\) copies of \(C_h\). Its uniform signature is \((p,q)\), its polarization is principal, and its underlying variety is \[
                J^{fq}\times C_h^{\,p-2q}.
 \tag{MP24}
\] The complementary dimension \(f(p-2q)\) attains the lower bound of LIFT083. Since the pure block is elliptic-isotypic and \(J\) is simple of dimension three, it introduces no further \(J\)-factor. At \(q=0\) the construction contains no \(J\)-factor.

The old Jacobian trace-one line is preserved on every twist because \(\sigma(w_0)=w_0\). The old convention used \((1-\jmath)/2=\bar w_0\); MP9 is unchanged by that replacement. Nevertheless no old selected polynomial, derivative, gauge-cost constant, or nonzero theta value is transported without a separate proof.

\subsection{8. Exact remaining gates and references}\label{exact-remaining-gates-and-references}

The source-wide construction is conditional on an actual Host with MP2. It does not supply its numerical SizesAt or Host existence window. The explicit \(E_h,d=7\) construction is a separate geometric model; the fixed degrees 42 and 168 still fail the already proved full source-size conditions. It cannot be relabeled as an admissible larger Host.

The point and the pinned datum have compatible \textbf{finite} integral symplectic modules; a global integral Betti isometry is neither proved nor required by TateIso. The exact source level is retained, but a finite coefficient extension may be needed for the point and its finite level sections. No infinite adelic level is descended to that finite extension.

The principal object and a field-valued FrPt do not imply the old LamObjGood/LamPt identities. LIFT082 may be applied only after its separate good-model, all-DVR, line-degree, frame, and same-point hypotheses are supplied. In particular \[
                     F(x)\ne0
 \tag{MP25}
\] for the prescribed automorphic/algebraic rule is still unevaluated. The field degree, finite denominators, new pure-block line degrees, and all DP44 terms require their own effective bounds. There is no global period, selected-content, Read/STiesR/032, DP44-sign, or RH conclusion.

Primary references directly read for this draft:

\begin{itemize}
\tightlist
\item
  Milne, J. S. (2008). \emph{Abelian varieties} (Version 2.0), Lemma 2.4 and Theorems 2.8--2.9, printed pp.~13--14; Theorem 5.1 and Corollary 5.2, printed pp.~21--22; §5.13, printed pp.~24--25: the source-sign Hermitian/Riemann dictionary, algebraization, normalized cube trivialization, and descent of a trivial line. \url{https://www.jmilne.org/math/CourseNotes/AVc.pdf}
\item
  Milne, J. S. (2020). \emph{Complex multiplication} (Version 0.10), Proposition 6.1 and its Appell--Humbert proof, printed pp.~46--47; Proposition 7.9 and Corollary 7.10, printed pp.~54--55: algebraization and descent of CM varieties and their endomorphisms. \url{https://www.jmilne.org/math/CourseNotes/CM.pdf}
\item
  K030 PH4--22 and K032 QH3--17: the already checked exact auxiliary tensor lattices, product principal polarizations, full integer-ring action, and trace-one lines.
\end{itemize}

Literal source files directly read include TwoWeightEqOneAddRoot.lean (Host, Datum, SAB, DetCond, TateIso, Obj), \nolinkurl{PartitionExtendByZeroAntitone.lean} (IsDatW and all Pins), DataHost.lean (AA, xi0, FactPrincipalDatum), ConjEMemHostFIff.lean (HostRaw.cmType), CubeSymmetricSwapFirst.lean (cubeLB and its signs), CubicalPullBaseChangeSt.lean (rigidified symmetry), NonzeroLocus.lean (RelAmple/KTriv), ConstGen.lean (thetaLB/CommIsOf/WeilIs), PushInvertible.lean (baseKA, Moduli.rep and Coeff), and IsCombAt.lean (FrPt). Their exact hashes, line counts, and literal anchors are recorded in \nolinkurl{MIXED_SIGNATURE_POINT_PENDING_MANIFEST.json}. The complete companion script was read and replayed: 37,101 exact assertions check the explicit cyclotomic principal lattice, its alternating determinant and positive principal minors, semicharacter and adjoint identities, local Hermitian residue models including dyadic lifts, source torsion sign and finite CRT level, sharp signature counts, and all 16 source/dependency fingerprints. These are consistency diagnostics accompanying the universal written proof; finite samples do not replace the local uniqueness, descent, or moduli arguments. The final author verification records the actual replay time and subject hashes. Independent whole review remains pending.

\subsection{Canonical algebraic-geometric context}\label{canonical-algebraic-geometric-context-13}

This is a principally polarized abelian scheme over a number field, with full CM integer-ring action, rigidified symmetric cubical theta line, literal adjoint biextension and finite level; a spread over an open part of Spec(O\_K) is qualitative only and no needed bad-prime control is inferred. Exact contextual anchors are FGA232 Theorem3.1 pp148--149 and FGA236 Corollaries1.7/2.3 pp229/232 for Picard representability; SGA7 I VII2.7/2.9 pp169--174 and VIII7.4(b) pp310--311 for biextensions/duality; cote4 scan003 sections8.1.1--8.1.2 and scans005--006 sections8.4.1--8.4.3, and cote54 scans005--007 for the carried formal-group/Cartier-dual language; SGA3 I4.7.3/4.7.3.1 and SGA6 VI Theorem3.3 remain contextual. The isogeny decomposition describes CM h1 motives rationally, but no new integral crystal comparison, lambda-ring/K-theory identity, l-adic weight estimate or motivic L-function consequence is asserted. DP\_MULT Q3e/DP44/DP52/DP54--65, THETA3 G1/A7, CX\_GRH CX-ALG-024 and GCT\_PLETH T-052 retain their effective arithmetic/global gates.

\clearpage

\section{AM-R-086 --- principal Jacobian blocks over a fixed general Host field}\label{am-r-086-principal-jacobian-blocks-over-a-fixed-general-host-field}

\textbf{Canonical checked integration --- 2026-11-23T02:52:52.625871+00:00.} Complete saved-proof and executable reviews and separate exact diagnostic replays passed. No numerical error charge is introduced. Frozen author history is preserved. AM-R-086 ZJ1--21 constructs principal source points with actual J content over any fixed actual Host E under the retained LIFT084 input hypotheses. The auxiliary compositum L=E Q(zeta7), c={[}L:E{]} in\{1,2,3,6\}, has an explicit integral CM trace polarization of degree\textbar D\_L\textbar7\^{}(cf/3). Zarhin Theorem1.1 gives a full integer-ring compatible principal block. Its exact conjugate-action balancing, with the even-c improvement, yields block sizes a\_c=(8,8,24,24). For p\textgreater=q the resulting exact J exponent is N=(2f/3)(q-r), r=q mod a\_c, and pure complementary dimension f(p-q+2r). If d=7 and p\textgreater=2q, a principal(16,8) block gives N=f(q-r7), r7=q mod8, complement dimension f(p-2q+3r7). Zero remainders attain the corresponding LIFT083 bounds. The unchanged Datum, finite level and Coeff give actual Obj/FrPt over a finite extension. A separate full O\_K7 lattice-approximation argument gives an underlying isogeny J\^{}N to the actual J factor prime to any prescribed finite set, and an isomorphism of p-divisible group schemes after good reduction. No old theta-line/frame compatibility is inferred. Effective field/isogeny/line/frame costs, numerical Host/Sizes existence, same-point F nonzero, actual lift identities/arithmetic image, global period, specific032 and DP44 remain open; no RH or finite numerical charge. Current101/31 arrays and complete floors25/11 are unchanged.

Author: GPT-6 Astra Pro, \nolinkurl{/root}, ALP\_MACHINE; reasoning effort not externally exposed. 23 November 2026 UTC.

\textbf{Controlled draft: whole saved-proof and independent checks pending.} The construction below keeps the actual Host field and source datum. It removes a geometric Jacobian-content premise under the same conditional source inputs as LIFT084. It does not construct a numerical Host/Sizes record, theta-lift chain, prescribed nonzero section, effective arithmetic costs, or a DP44 comparison.

\subsection{1. Statement with the fixed source inputs}\label{statement-with-the-fixed-source-inputs}

Let E be the CM field of an actual source Host, let {[}E:Q{]}=2f, and retain its quadratic subfield K\_0=Q(j), j squared equal to -d, its selected pure type Psi, and all inputs of LIFT084 section1: the actual compatible Par, PinHost/PinLevel/PinDat, stage Datum, existing Moduli representation and same Coeff field. Write (p,q) for that Datum's uniform multiplicities. Assume p\textgreater=q and p+q\textgreater0. Let C be LIFT084's pure principal block of dimension f; C\^{}- denotes the same variety with conjugated O\_E action. Set

\[
 K_7=\mathbb Q(\zeta_7),\quad
 J=\operatorname{Jac}(y^2=x^7-1),\quad
 L=E K_7,\quad c=[L:E],\quad t=[L:K_7]=cf/3.
 \tag{ZJ1}
\]

The compositum is taken inside a fixed algebraic closure compatible with the chosen complex embedding. It is auxiliary: \textbf{the source Host remains E}. Since K\_7/Q is cyclic of degree six,

\[
                       c\in\{1,2,3,6\},\qquad t\in\mathbb Z_{>0}.
 \tag{ZJ2}
\]

There exists a principally polarized abelian variety B with full integral O\_E action, Rosati conjugation, and uniform multiplicities (8c,8c), whose underlying variety satisfies

\[
                B\sim J^{16t}=J^{16cf/3}.
 \tag{ZJ3}
\]

Put b=floor(q/(8c)) and r=q-8cb, so 0\textless=r\textless8c\textless=48. Then

\[
 A_{p,q}=B^b\times C^{p-8cb}\times(C^-)^{q-8cb}
 \tag{ZJ4}
\]

has the exact source uniform signature (p,q), principal polarization and full O\_E action. It supplies a literal Obj and FrPt over a finite extension of the \textbf{same} Coeff field, at the unchanged source finite level and Datum. Its exact J-isotypic exponent and complementary dimension are

\[
 N=16tb=\frac{2f}{3}(q-r),\qquad
 \dim A_{\rm comp}=f(p-q+2r).
 \tag{ZJ5}
\]

Here the complement is the displayed pure-block product, with no additional J constituent. If d!=7 and r=0, these quantities attain LIFT083(JS8)'s upper bound on J content and lower bound on complementary dimension.

There is a useful improvement when c is even: Y alone, defined in section3, already has uniform signature (4c,4c). In that case replace the block size 8c by 4c in ZJ4--5 and take N=8tb. The common formulas N=(2f/3)(q-r) and dim A\_comp=f(p-q+2r) are unchanged with the new remainder. Consequently the construction can always use a block size a\_c given by a\_1=a\_2=8 and a\_3=a\_6=24, and a remainder strictly smaller than 24. The separate d=7 construction below improves its block size to eight.

If d=7 and p\textgreater=2q, there is a stronger construction. There exists a principal O\_E block D with uniform multiplicities (16,8) and underlying variety J\^{}(8f). For b\_7=floor(q/8), r\_7=q-8b\_7, take

\[
 A^{(7)}_{p,q}=D^{b_7}\times C^{p-16b_7}\times(C^-)^{q-8b_7}.
 \tag{ZJ6}
\]

Its exact J exponent and complementary dimension are

\[
       N_7=f(q-r_7),\qquad
       \dim A^{(7)}_{\rm comp}=f(p-2q+3r_7),\quad0\le r_7<8.
 \tag{ZJ7}
\]

When 8 divides q this attains LIFT083(JS7)'s bound. No optimality is asserted for nonzero remainders. The already checked E\_h construction of LIFT084 can do better in those special fields; it is not being superseded.

For either construction and any prescribed finite set P of rational primes, the actual J-isotypic product factor A\_J admits, after a finite coefficient extension, an isogeny

\[
                  J^N\longrightarrow A_J
              \quad\text{of degree prime to every prime in }P.
 \tag{ZJ8}
\]

At places of good reduction this induces an isomorphism of the underlying p-divisible groups for p in P. Such good reduction can be acquired over a further finite extension. Neither the extension degree nor the isogeny degree is bounded here, and the isogeny is not asserted to preserve the old selected theta line or prescribed frame.

\subsection{2. A polarized full-integer-ring block before principalization}\label{a-polarized-full-integer-ring-block-before-principalization}

The settled Jacobian calculation LIFT083(JS9--12), with DP\_MULT DP45/DP49, gives the faithful full O\_(K\_7) action and type Phi=\{1,2,3\}, relative to the automorphism (x,y) maps to (zeta\_7 x,y). Its geometric endomorphism algebra is K\_7 and J is simple. These known findings are carried, not reproved.

The field L is CM. Indeed complex conjugation preserves each constituent E and K\_7 and hence their compositum. For any embedding of L, its restrictions to the two CM constituents intertwine their intrinsic conjugations with complex conjugation. The same therefore holds on L. The fixed field is totally real, and its index is two because K\_7 is not real. Define the induced type

\[
 \Phi_L=\{\rho:L\hookrightarrow\mathbb C:\rho|_{K_7}\in\Phi\}.
 \tag{ZJ9}
\]

It contains exactly one member from each conjugate pair. Every K\_7 embedding has t extensions, so its size is 3t=cf.~The compositum degree in ZJ2 follows by restriction of the Galois group of L/E to K\_7; its fixed field there is E intersect K\_7. Thus c divides six, and {[}L:Q{]}=2cf=6t.

There is an explicit integral polarization, so no ideal-class principalization premise is inserted. Put

\[
 \xi_7=\zeta_7-\zeta_7^{-1},\qquad
 X=\mathbb C^{\Phi_L}/\Phi_L(\mathcal O_L),\qquad
 \Omega_X(x,y)=\operatorname{Tr}_{L/\mathbb Q}(\xi_7\bar x y).
 \tag{ZJ10}
\]

The element xi\_7 is integral, purely imaginary, and has positive imaginary part at the three embeddings in Phi: these are twice i times sin(2pi a/7), a=1,2,3. Therefore Omega\_X is integral and alternating on O\_L, is invariant under the complex structure, and satisfies Omega\_X(Jv,v)\textgreater0 for nonzero v. It is nondegenerate because the trace pairing is nondegenerate and xi\_7 is nonzero. The Riemann-form algebraization theorem gives X as a polarized abelian variety. Multiplication by O\_L is holomorphic and preserves the lattice, and

\[
             \Omega_X(a x,y)=\Omega_X(x,\bar a y)
                    \quad(a\in\mathcal O_L).
 \tag{ZJ11}
\]

Thus the polarization lambda\_X has the exact conjugate-adjoint property required below. It need not be principal. For clarity, its degree is explicit in the auxiliary field discriminant:

\[
              \deg\lambda_X=|D_L|\,7^t.
 \tag{ZJ12}
\]

To see this, its alternating dual lattice is xi\_7 inverse times the absolute codifferent of L, up to the harmless sign from conjugation. The index of O\_L in that dual is \textbar D\_L\textbar{} times the absolute norm of xi\_7. The norm in K\_7 is seven: xi\_7=zeta\_7 inverse times (zeta\_7 squared minus one), and the product of (1-zeta\_7\^{}a), a=1,\ldots,6, is Phi\_7(1)=7. Its norm in L is 7\^{}t. The degree of a polarization is the determinant of its integral alternating form, equivalently that dual-lattice index. ZJ12 is not an effective bound for a complete source cost because D\_L and the field/model/frame data still enter.

As a rational K\_7-vector space, L is K\_7\^{}t. Any K\_7 basis identifies its rational Hodge structure for the induced type ZJ9 with t copies of the Hodge structure of J. The corresponding lattices are commensurable, so multiplication by a common integer gives a holomorphic isogeny. Consequently

\[
                              X\sim J^t.
 \tag{ZJ13}
\]

This is an underlying isogeny statement, not an identification of polarizations or of a preselected integral O\_E action on J\^{}t. Algebraization and CM descent place X, its polarization, and finitely many O\_L action generators over a number field. Every relation then descends over that same field. We use Milne (2020), Proposition6.1 and Proposition7.9/Corollary7.10, at their stated scopes.

\subsection{3. Equivariant principalization and exact balancing}\label{equivariant-principalization-and-exact-balancing}

We use the known equivariant quaternion trick, rather than reproving it. Zarhin (2023; arXiv version5, Theorem1.1, p.2) says that an integral ring action compatible with an involution and a polarization yields a compatible principal polarization on (X times X-dual)\^{}4. The dual action is a maps to iota(bar a)-dual. Apply that theorem to O\_L and complex conjugation. It gives

\[
 Y=(X\times X^\vee)^4,
 \quad\iota^\vee(a)=\iota(\bar a)^\vee,
 \quad\mu_Y\iota_Y(a)=\iota_Y(\bar a)^\vee\mu_Y,
 \tag{ZJ14}
\]

with mu\_Y principal and all data over a number field. The action is integral and unital; restricting it to O\_E is allowed because O\_E is contained in O\_L.

The dual action in ZJ14 is crucial. The original polarization satisfies lambda\_X iota(a)=iota(bar a)-dual lambda\_X. Thus lambda\_X is an O\_L-linear isogeny to the \textbf{conjugate-dual action}, and its differential is invertible in characteristic zero. X and that dual have the same O\_L Lie type. Hence Y has type eight times Phi\_L and is isogenous to J\^{}(8t). One must not replace this by a sum of the two conjugate types at this step.

For an embedding tau:E to C let m\_tau be the number of extensions to L in Phi\_L. Each tau has c extensions, and conjugation bijects its excluded extensions with the included extensions over bar(tau). Therefore

\[
                            m_\tau+m_{\bar\tau}=c.
 \tag{ZJ15}
\]

Let Y\^{}- be the same principally polarized variety with its O\_E action precomposed by conjugation, and take B=Y times Y\^{}-. Its multiplicity at tau is 8m\_tau+8m\_bar(tau)=8c. Its product polarization is principal, its Rosati involution on O\_E is conjugation, and its underlying variety is J\^{}(16t) up to isogeny. This proves ZJ3, including uniformity at \textbf{every} E embedding, rather than only a total quadratic signature.

The exponents in ZJ4 are nonnegative because p\textgreater=q\textgreater=8cb. Adding its E-type multiplicities gives p on Psi and q on its conjugate. Its dimension is f(p+q). The pure blocks C and C\^{}- are elliptic-isotypic, whereas J is simple of dimension three. Thus the entire J-isotypic factor is exactly B\^{}b. Counting dimensions gives

\[
 3N=16cfb=2f(q-r),\qquad
 f(p+q)-3N=f(p-q+2r),
 \tag{ZJ16}
\]

which proves ZJ5 and its comparison with the already established d!=7 bound.

For the even-c improvement, let H=Gal(K\_7/(E intersect K\_7)), of order c.~Its extension sets are cosets of H. When c is even, H contains the unique order-two element, complex conjugation. Each coset is therefore a disjoint union of conjugate pairs, and Phi selects exactly one member from each pair. Thus m\_tau=c/2 for every tau. The principal Y has each E multiplicity 8(c/2)=4c and underlying J\^{}(8t). It supplies the asserted smaller balanced block without pairing with Y\^{}-. In summary the available balanced block size is

\[
 a_c=\begin{cases}4c,&c\text{ even},\\8c,&c\text{ odd},\end{cases}
 \qquad(a_1,a_2,a_3,a_6)=(8,8,24,24).
 \tag{ZJ21}
\]

With b=floor(q/a\_c), r=q-a\_cb, its J exponent is (2f/3)a\_cb. Thus the relation 3N=2f(q-r) and the complementary-dimension formula in ZJ16 hold for this smaller block as well. The integer exponents follow from 3 dividing cf; no fractional number of J factors is used.

\subsection{4. The stronger block when the Host quadratic field is Q(sqrt(-7))}\label{the-stronger-block-when-the-host-quadratic-field-is-qsqrt-7}

Assume d=7. Fix the identification of the common quadratic subfield so its selected root is the source root. If its sign is reversed in a different convention, conjugate the Jacobian action throughout; the underlying J is unchanged. Now E intersect K\_7 contains K\_0, so c is either three or one.

If c=3, each E embedding over the positive quadratic root has the three quadratic-residue extensions to K\_7. Phi contains two of those and one of the nonresidue extensions. Thus X itself has uniform O\_E multiplicities (2,1), and t=f.~The single Y in ZJ14 has uniform (16,8) and underlying J\^{}(8f). Set D=Y.

If c=1, K\_7 is contained in E and t=f/3. Use the three types

\[
 \Phi_0=\{1,2,3\},\quad
 \Phi_1=\{2,4,6\},\quad
 \Phi_2=\{4,1,5\}.
 \tag{ZJ17}
\]

For each a=0,1,2 construct X\_a from O\_E with type induced by Phi\_a. A positive integral trace coefficient is sigma\_(2\^{}(-a))(xi\_7), where sigma\_u(zeta\_7)=zeta\_7\^{}u. The construction of section2 applies verbatim with this explicit coefficient. Its underlying variety is J\^{}t: the type change is a twist of the K\_7 action on the same J. Apply ZJ14 separately to X\_a and take D=Y\_0 times Y\_1 times Y\_2. Every residue label occurs twice in ZJ17 and every nonresidue once. Therefore D has uniform multiplicities (16,8) and is isogenous to J\^{}(24t)=J\^{}(8f).

This construction does not require an automorphism of E extending sigma\_2. The three induced types and three polarized lattices are constructed separately. In both cases D has full O\_E action and a principal polarization. Since p\textgreater=2q, the exponents p-16b\_7 and q-8b\_7 in ZJ6 are nonnegative. Adding multiplicities gives exactly (p,q). The pure complement has dimension

\[
 f\{p+q-24b_7\}=f(p-2q+3r_7),
 \tag{ZJ18}
\]

and the J exponent is 8fb\_7=f(q-r\_7). This proves ZJ6--7. The terminal case q=0 has no J factor, in agreement with LIFT083; the construction does not evade that obstruction.

\subsection{5. Literal source line, Tate level, point and finite list of stages}\label{literal-source-line-tate-level-point-and-finite-list-of-stages}

Each new product is an actual principally polarized abelian variety with full O\_E action, Rosati conjugation and the exact uniform source Lie multiplicities. Its Betti lattice is a finite torsion-free O\_E module, hence projective because O\_E is Dedekind, and has rank p+q. Its integral alternating form is unimodular. These are precisely the inputs used in the universal parts of LIFT084 sections3--6; those arguments were not restricted to its particular product of pure blocks.

Choose on this new lattice the trace-one semicharacter

\[
                  \alpha(\lambda)=(-1)^{\Omega(\lambda,w_0\lambda)},
            \qquad w_0=(1+\jmath)/2.
 \tag{ZJ19}
\]

LIFT084(MP9--13) gives the actual rigidified symmetric cubical line and SAB.prim/adjoint line identities. It retains the given principal polarization, rather than asserting that a previously chosen arbitrary representative already has the required semicharacter. The same-field normalized descent argument applies to every a in O\_E. The relative-unramified local lemma, including dyadic places, matches this new selfdual lattice with the \textbf{given} Datum after profinite completion. LIFT084(MP16--19) then supplies its geometric Tate systems and the unchanged M=g\_0g\_1, A=1 level sections. No global integral Betti isometry is asserted.

The same reduced-base Moduli.rep and fixed Coeff give the literal Obj, point and FrPt over a finite extension. The exact source locators are LIFT084's manifest: TwoWeightEqOneAddRoot Host/Datum/SAB/TateIso/Obj, PartitionExtendByZeroAntitone IsDatW/Pins, PushInvertible Moduli.rep/Coeff, and IsCombAt FrPt. There is no change to the source Host E, its Datum, finite level or coefficient record.

For a finite list of pinned stages, make these constructions stage by stage and take the compositum of their fields of definition. Under the full Sizes hypotheses LIFT083 gives p\_j\textgreater2q\_j, so ZJ6 is available at every stage if d=7, and ZJ4 is available for every other d.~The exact counts are ZJ5 or ZJ7 with q\_j=m-j and p\_j=k+j-1. This produces a finite list of actual source points; it does \textbf{not} construct any lift arrow or prove the six lift identities between them.

\subsection{6. Isogenies avoiding prescribed primes and the precise local consequence}\label{isogenies-avoiding-prescribed-primes-and-the-precise-local-consequence}

We prove the limited prime-avoidance statement ZJ8, because a general isogeny by itself would not preserve the integral p-divisible group.

\textbf{Lattice lemma.} Let A be a complex abelian variety with a full O\_(K\_7) action whose rational Hodge structure is N copies of the type of J. Then for any finite set P of rational primes there is a holomorphic isogeny J\^{}N to A of degree prime to P.

\textbf{Proof.} Enlarging P by one prime is harmless, so assume it is nonempty. The zero-dimensional case is the identity map; assume N\textgreater0. Identify the rational homology spaces as K\_7\^{}N with their common complex structure. Let Lambda\_0 and Lambda\_A be the two integral lattices. Both are finite projective O\_(K\_7) modules of rank N. At each rational prime ell in P their completions are free of rank N over O\_(K\_7) tensor Z\_ell, a product of DVRs. There is consequently an invertible K\_7 tensor Q\_ell-linear map g\_ell carrying one completed lattice onto the other. The condition of doing so is open in the matrix topology: in integral bases it is GL\_N(O\_(K\_7) tensor Z\_ell).

Choose a matrix g in M\_N(K\_7) sufficiently close to all g\_ell. This simultaneous approximation follows directly by expressing every matrix entry in a fixed Q basis of K\_7 and applying rational approximation at the finitely many rational primes: multiply the local entries by a common product of those primes to make them integral, use integer CRT to the required powers, and divide back by that product. The chosen open conditions make g invertible and give g Lambda\_0 tensor Z\_ell=Lambda\_A tensor Z\_ell for every ell in P.

The quotient (g Lambda\_0+Lambda\_A)/Lambda\_A is a finite group whose order has no prime divisor in P. Choose an integer m annihilating it, also prime to P. Then mg Lambda\_0 is contained in Lambda\_A and agrees with it after completing at P. The induced complex-linear map is an isogeny and its degree is the lattice index {[}Lambda\_A:mg Lambda\_0{]}, which is prime to P. Algebraization makes it an algebraic isogeny. This proves the lemma. \(\square\)

The J-isotypic products constructed above satisfy that lemma. On X the auxiliary K\_7 type is Phi repeated t times. On its conjugate-dual in ZJ14 it is the same type, as already proved by the polarization isogeny. On Y\^{}- use the original auxiliary K\_7 action, independently of the conjugated O\_E action. In the three-type construction precompose each auxiliary K\_7 action by the inverse type twist; this again gives Phi and preserves the full integer ring. Taking their product gives a full O\_(K\_7) action of type Phi repeated N times on A\_J. It need not be the original restriction of the source O\_E action. That distinction does not alter A\_J or its source structure. The lemma proves ZJ8, and CM full faithfulness descends the resulting map to a finite extension.

All the factors are CM abelian varieties, so they acquire good reduction at any prescribed finite set of places after a finite extension (Milne, 2020, Proposition7.12, p.55). At such places both the isogeny and a quasi-inverse extend uniquely over the DVR (ibid., Corollary6.6, p.49). If n is its degree, the kernel is killed by n, so a homomorphism in the reverse direction can be chosen with both composites multiplication by n.~The extended identities hold over the DVR. For p in P, multiplication by n is an automorphism of every p-power torsion group scheme. Thus the isogeny induces an isomorphism

\[
                  J[p^\infty]^N\simeq A_J[p^\infty].
 \tag{ZJ20}
\]

The proof concerns group schemes, not only torsion points. The connected formal groups and their Dieudonne objects inherit that isomorphism. In particular one may include 11, or any specified primes congruent to 3 or 5 modulo7, in P. The previously accepted DP45/DP49 local group types therefore remain present in this J constituent after good reduction. The map has not been shown to respect the old theta line, selected harmonic directions, differential gauge, or source-normalized divided-power operation. No new slope or lift saving follows from ZJ20 alone.

\subsection{7. Effective boundary and Grothendieck interpretation}\label{effective-boundary-and-grothendieck-interpretation}

The remaining source construction problem has genuinely changed in the stated conditional scope: the fixed general Host field now admits literal principal source points containing a quantitatively specified J constituent, with a pure complementary factor. There is no requirement that E itself contain K\_7; the auxiliary compositum was used only to construct its action. Nor is a change to the existing source Datum or level hidden in the construction.

The field of definition, prime-avoiding isogeny degree, integral frames, new line degrees and all arithmetic normalizations remain uncontrolled. Potential good reduction is qualitative and does not supply the particular all-DVR/good-model predicates of LIFT082. A prescribed rule F may still vanish at every point chosen here. No Read/STiesR identification, actual arithmetic image, global period, specific032 cycle pairing, lift chain or value of any g\_j,e\_j,H has been established. DP44 remains unevaluated, and neither a shorter chain nor a universal negative cost is proved.

Geometrically the construction is an abelian scheme over a number field with O\_E action, principal polarization, Picard class, rigidified symmetric cubical theta line and Poincare biextension, with its explicit rational J-isotypic projector. It spreads over an open part of Spec(O\_K); no required bad-prime set or effective model is supplied. Its rational h1 motive decomposes into the stated J copies and the pure elliptic-isotypic complement. Isogeny-induced rational realizations and, at the selected good primes, ZJ20 are the only comparisons asserted. No new lambda-ring/K-theory identity, l-adic weight estimate or motivic L-function/RH conclusion is inferred.

The carried exact Grothendieck locators are FGA232 Theorem3.1 pp148--149; FGA236 Corollaries1.7/2.3 pp229/232; SGA7 I VII2.7/2.9 pp169--174 and VIII7.4(b) pp310--311; SGA3 I4.7.3/4.7.3.1; SGA6 VI Theorem3.3; cote4 scan003 sections8.1.1--8.1.2 and scans005--006 sections8.4.1--8.4.3; cote54 scans005--007. They identify the inherited Picard, duality and formal-group frameworks, not a proof of any new global period. Run interfaces remain DP\_MULT Q3e, DP44/DP52/DP54--65; THETA3 G1/A7; CX\_GRH CX-ALG-024; and GCT\_PLETH T-052.

\subsection{References and verification boundary}\label{references-and-verification-boundary}

\begin{itemize}
\tightlist
\item
  Zarhin, Y. G. (2023). Abelian varieties, quaternion trick and endomorphisms. In I. Cheltsov, X. Chen, L. Katzarkov, \& J. Park (Eds.), \emph{Birational geometry, Kahler--Einstein metrics and degenerations} (Springer Proceedings in Mathematics \& Statistics, Vol.409, pp.857--864). Springer. \url{https://doi.org/10.1007/978-3-031-17859-7_42}. Exact theorem used: Theorem1.1; directly read primary author version arXiv:2010.07995v5 (29 January2022), p.2, equations(3)--(6). The known theorem is cited and not reproved.
\item
  Milne, J. S. (2020). \emph{Complex multiplication} (Version0.10). \url{https://www.jmilne.org/math/CourseNotes/CM.pdf}. Proposition6.1 and Appell--Humbert discussion, pp46--47; Corollary6.6, p49; Proposition7.9/Corollary7.10, pp54--55; Proposition7.12, p55. Primary passages directly read.
\item
  ALP\_MACHINE AM-LIFT-083, JS1--25: settled endomorphism/type calculation and actual-source J-content bounds. AM-LIFT-084, MP1--25: pure principal block and the exact line/profinite-lattice/finite-level/descent/FrPt construction. Their whole-proof checks and exact replay records are retained.
\end{itemize}

The companion exact diagnostics check the cyclotomic trace polarization, degree and norm, all subgroup extension-count patterns, all displayed type sums and dimension/remainder formulas, the three-twist construction and elementary prime-avoidance lattice models. They are not a substitute for the universal CM, principalization, descent, lattice-approximation or source-interface proofs above. No finite arithmetic remainder charge is introduced.

\subsection{Canonical algebraic-geometric context}\label{canonical-algebraic-geometric-context-14}

The principal polarized abelian scheme over a number field retains the source O\_E action, actual Picard class, trace-one rigidified symmetric cubical line and Poincare biextension. The auxiliary field does not change the Host. The rational J-isotypic h1 projector gives the displayed CM-motive decomposition; after selected good reduction the prime-avoiding isogeny yields an actual p-divisible group-scheme isomorphism, hence the corresponding formal-group/Dieudonne isomorphism. It does not identify the selected theta line or source-normalized differential operation. Exact inherited contextual anchors: FGA232 Theorem3.1 pp148--149;FGA236 Corollaries1.7/2.3 pp229/232;SGA7 I VII2.7/2.9 pp169--174 and VIII7.4(b) pp310--311;SGA3 I4.7.3/4.7.3.1;SGA6 VI Theorem3.3;cote4 scan003 sections8.1.1--8.1.2 and scans005--006 sections8.4.1--8.4.3;cote54 scans005--007. No new lambda-ring/K-theory identity, l-adic weight bound or motivic L-function assertion. DP\_MULT Q3e/DP44/DP52/DP54--65 and THETA3 G1/A7 retain their effective costs/nonvanishing/arithmetic gates;CX\_GRH CX-ALG-024 and GCT\_PLETH T-052 remain open in their global scopes.

\clearpage

\section{AM-LIFT-085 --- A nonzero frame rule on a fixed-level CM isogeny orbit}\label{am-lift-085-a-nonzero-frame-rule-on-a-fixed-level-cm-isogeny-orbit}

\textbf{Canonical checked integration --- 2026-11-23T03:08:31.882138+00:00.} Complete saved-proof and executable reviews and separate exact diagnostic replays passed. No numerical error charge is introduced. Frozen author history is preserved. AM-LIFT-085 HF1--26 proves that, for an existing literal Host/Datum/Moduli/Coeff, actual CM seed Obj, original finite level A=1 and a given nonzero natural frame rule, a point in that same OE-isogeny class has a nonzero literal FrPt evaluation at the unchanged finite level. The isogeny u has degree N\^{}(2g), N prime to any prescribed finite set S containing the level primes, N=1 mod M, and u*L1=L0\^{}(N\^{}2) rigidified and cubical. Every output object and a nonzero frame descend over K1=K0(A0{[}N\^{}2{]}), with {[}K1:K0{]}\textless=N\^{}(8g\^{}2). The argument meets every analytic moduli component through an algebraic family based on its own target lattice. A weight-rule frame can be chosen by integer matrices with entries0..D+1 relative to any initial rational eigenframe. N itself, seed degree, initial-frame/theta-generator heights, models and actual line costs remain unevaluated. The given rule must itself be nonzero; no prescribed arithmetic-image nonvanishing is proved. No old pure LamObjGood or selected derivative identification, global period, specific032 pairing, DP44 or RH conclusion follows. This removes the qualitative same-rule same-point existence gate within the fixed CM isogeny class under the precise assumptions. Current101/31 error arrays and complete floors25/11 remain unchanged.

Author: /root/lift\_accounting. 23 November 2026 UTC.

\textbf{Author whole-proof/source/script check complete; exact replay and independent review are recorded separately.} This is a qualitative theorem for an already supplied source Host, Datum, level, moduli representation, coefficient field, and nonzero rule. Its isogeny denominator is not bounded. It neither constructs an admissible numerical family nor proves that the prescribed arithmetic image is nonzero.

\subsection{1. The exact assertion and the preserved data}\label{the-exact-assertion-and-the-preserved-data}

Fix the literal \nolinkurl{Tor.Host E}, \nolinkurl{Datum H}, \nolinkurl{Level H}, \nolinkurl{Moduli Dw l}, and \nolinkurl{Coeff l} of the source files listed in §9. Assume \[
 E/F\text{ is CM and unramified at every finite place},\qquad
 r=Dw.p+Dw.q\ge1,\qquad l.A=1.
 \tag{HF1}
\] The unramified condition is a field of the source Host. The source supplies an integral projective \(\mathcal O_E\)-module \(Dw.La\), a unimodular alternating form with adjoint conjugation, and its real signature domain. The level has odd \(M\ge3\) and coprime conjugate factors \(\mathfrak M\bar{\mathfrak M}=M\mathcal O_E\). In the pinned application, \nolinkurl{Par.PinLevel} says exactly \(M=g_0g_1,A=1\). We keep that level, its actual \(\mathfrak M\), \nolinkurl{Dw}, \nolinkurl{Xo}, and \nolinkurl{c} fixed throughout.

Put \(K=c.F\), \(f=[F:\mathbb Q]\), \(g=fr\), and \(X=Xo.XF(c)\). Fix a CM seed object \[
 (A_0,\iota_0,\mathcal L_0,\eta_0)
 \quad\text{over a number field }K_0\supset K,
 \tag{HF2}
\] which is an actual source \nolinkurl{Obj} of these data. Its full integer-ring action, principal line, cubical rigidification, and prescribed finite \(\eta_0\) are part of the seed. LIFT084 constructs such seeds under its stated source assumptions. The theorem below also applies to any other seed satisfying HF2; it does not require its polarization to be a product or its underlying variety to be purely Jacobian-isotypic.

Let \(F_0\ne0\) be a literal natural frame rule. In particular it may be an element of \nolinkurl{sec Xo c a n}. Let \(S\) be any finite set of rational primes containing every prime dividing \(M\). Then there exist a positive integer \(N\), a finite extension \(K_1/K_0\), an actual source object \((A_1,\mathcal L_1,\eta_1)\), an \(\mathcal O_E\)-equivariant isogeny \(u:A_0\to A_1\), and a field-valued literal frame \(y\), such that \[
 \begin{gathered}
 \gcd\!\left(N,\prod_{p\in S}p\right)=1,\qquad N\equiv1\pmod M,\\
 u^*\mathcal L_1\simeq\mathcal L_0^{\otimes N^2},\qquad
 \deg u=N^{2g},\qquad u\eta_0=\eta_1,\\
 y.x\text{ is the point of }(A_1,\mathcal L_1,\eta_1),\qquad
 \operatorname{evR}(F_0,y)\ne0.
 \end{gathered}
 \tag{HF3}
\] The line isomorphism is rigidified and cubical, not just an equality of polarization classes. One may take \[
 K_1=K_0(A_0[N^2]),\qquad
 [K_1:K_0]\le\left|\operatorname{GL}_{2g}(\mathbb Z/N^2\mathbb Z)\right|
 \le N^{8g^2}.
 \tag{HF4}
\] For \(N=1\), the group in HF4 means the trivial group. The entire conclusion, including the frame, is over this field; §7 proves the descent of all the listed data.

The same argument proves a stronger qualitative statement: the points satisfying HF3 except for the last nonvanishing condition are complex analytically dense in \textbf{every} component of the actual \(X\). The proof does not confuse this with the seed's initial integral lattice class. It starts from an arbitrary target point, keeps that target's lattice and finite level, and varies its compatible complex structure.

No bound for \(N\) in terms of the source numerical parameters is supplied. No nonzero rule is constructed here. In particular applying HF3 requires the actual arithmetic rule, rather than an unrelated nonzero section of the same weight.

\subsection{2. The trace-one element selects the actual line}\label{the-trace-one-element-selects-the-actual-line}

Write \(w=w_0\), so \(w+\bar w=1\). Use the source alternating convention \(\Omega(Jv,v)>0\). The Appell--Humbert Hermitian form is linear in the first variable and has imaginary part \(\Omega\).

For a principally polarized complex torus with full \(\mathcal O_E\)-action and Rosati conjugation, define \[
 \alpha(\lambda)=(-1)^{\Omega(\lambda,w\lambda)}.
 \tag{HF5}
\] LIFT084 MP9--12 proves that this is a semicharacter and gives the required symmetric, rigidified, cubical line. Here we also need uniqueness. If a line with this polarization has semicharacter \(\beta\), its source second-difference pullback along \((1,w)\) has semicharacter \[
 \frac{\beta((1+w)\lambda)}{\beta(\lambda)\beta(w\lambda)}
   =\exp\bigl(\pi i\Omega(\lambda,w\lambda)\bigr)=\alpha(\lambda).
 \tag{HF6}
\] Its Hermitian form is the original one, since the cross term is \(H(v,wv')+H(wv,v')=H(v,(w+\bar w)v')=H(v,v')\). Thus the literal \nolinkurl{SAB.prim} identity forces \(\beta=\alpha\). A symmetric two-torsion modification is not an extra choice compatible with \nolinkurl{prim}.

There is also a relative algebraic construction. If \(\lambda:A\to A^\vee\) is the principal polarization of an abelian scheme with Rosati-compatible action, let \(\mathcal P\) be the normalized relative Poincaré bundle and put \[
 \mathcal B_\lambda=(1\times\lambda)^*\mathcal P,
 \qquad
 \mathcal L_\lambda=(1,w)^*\mathcal B_\lambda.
 \tag{HF7}
\] The polarization of this line is \(\lambda w+w^\vee\lambda=\lambda(w+\bar w)=\lambda\). On each complex fibre HF6 identifies it with HF5, so it is relatively ample and principal. The normalized Poincaré construction is compatible with base change. The usual cubical structure of a rigidified line gives \nolinkurl{cub}; the fibrewise symmetry, adjoint, and prim identities extend as normalized line isomorphisms over a reduced base. To see the last assertion precisely, a rigidified line difference defines a section of the separated relative Picard space. If its restriction to every geometric fibre is trivial, that section equals zero on a reduced base. The line is then pulled back from the base, and its rigidification trivializes that pullback. A normalized trivialization is unique, since \(\pi_*\mathcal O_A=\mathcal O_S\). Consequently all the identities, for \textbf{every} \(a\in\mathcal O_E\), descend over the same base. This is not a choice of infinitely many field extensions.

The construction uses the relative rigidified Picard functor and Poincaré bundle. The fibre identity \((1\times\lambda_L)^*\mathcal P=\mathcal B_L\) is Milne's exact polarization convention; the rigidified relative Picard description is recorded in van der Geer--Moonen, §6.2--6.3. These are algebraic statements. No algebraicity is inferred merely from a holomorphic line on a nonproper parameter space.

\subsection{3. A literal natural rule is regular on the frame torsor}\label{a-literal-natural-rule-is-regular-on-the-frame-torsor}

The universal source object over \(X\) is an abelian scheme: it is proper and semiabelian, with the given relatively ample principal line. Its polarization is principal because \nolinkurl{KTriv} is the triviality of the translation kernel of that ample line. The coefficient field contains every image of \(E\). Thus the tangent bundle decomposes over \(K\) into eigenbundles \[
 \operatorname{Lie}(A/X)=\bigoplus_{\tau:E\hookrightarrow\mathbb C}\mathcal V_\tau,
 \qquad \operatorname{rank}\mathcal V_\tau=Dw.mult(\tau).
 \tag{HF8}
\] The ranks follow from the geometric \nolinkurl{DetCond} clauses; the idempotents of the split algebra \(E\otimes_\mathbb Q K\) give the decomposition. The sheaf \(\mathcal T=\pi_*\mathcal L\) is invertible and commutes with arbitrary base change. Indeed an ample line on a characteristic-zero abelian variety has no positive cohomology, and for a principal polarization its space of sections has dimension one. Cohomology and base change then gives the assertion. One can obtain the vanishing without an extra positivity convention: pull back the symmetric line by multiplication by a sufficiently large integer, use Serre vanishing for the resulting power, and split the pullback by the trace of the finite étale multiplication map.

Let \[
 P=\left(\prod_\tau\operatorname{Fr}(\mathcal V_\tau)\right)
       \times_X\mathcal T^\times.
 \tag{HF9}
\] This is a smooth surjective torsor over \(X\) under \(\prod_\tau\operatorname{GL}_{Dw.mult(\tau)}\times\mathbb G_m\). A zero-rank block contributes the trivial group. The literal \nolinkurl{FrPt} is precisely its functor of points: \nolinkurl{tan}, \nolinkurl{eig}, and \nolinkurl{basis} express the eigenframe, while \nolinkurl{gen} expresses a generator of \(\pi_*\mathcal L\). \nolinkurl{FrMap} is the corresponding pullback relation, including the actual pullback of the theta section. This description does not impose a CM condition on arbitrary test points.

Every affine chart of \(P\) is a reduced finite-type \(K\)-scheme. Its coordinate algebra has a presentation by a polynomial ring modulo a radical ideal and is therefore a literal \nolinkurl{TRing}. Apply the natural rule to its universal \nolinkurl{FrPt}. Naturality makes these chart functions agree after restriction to affine intersections, hence gives \[
                       f_{F_0}\in\Gamma(P,\mathcal O_P).
 \tag{HF10}
\] Conversely every rule value on a test ring is the pullback of HF10. This follows by covering the image of that test point by the affine charts and using the same naturality relation. Consequently \(F_0\ne0\) implies \(f_{F_0}\ne0\). Faithful base change to \(\mathbb C\) preserves this nonzero function. Since \(P\) is reduced, \(P_{f_{F_0}}\) is a nonempty open set. Its image \[
                   U_{F_0}=\pi(P_{f_{F_0}})\subset X
 \tag{HF11}
\] is nonempty and open because \(\pi\) is smooth. Thus it suffices to meet this \textbf{actual} open set in the base. This step proves neither that the initial seed lies in it nor that it meets a preselected connected component.

If a point \(x\in U_{F_0}\) is defined over a number field containing \(K\), its fibre of HF9 has a rational point: choose bases of the eigenspaces and a generator of the one-dimensional section space. It is then a product of general linear groups and a multiplicative group. Over an infinite field, its rational points are Zariski dense, so a frame with nonzero value exists over that same field.

This also resolves the source's \nolinkurl{evalAt} convention. A number field finite over \(K\) has a radical polynomial presentation, and any of the field-valued frames just constructed is \nolinkurl{FromTest}. Naturality identifies every permissible test presentation with HF10. Therefore \nolinkurl{evR} equals its regular-function value and does not use the default-zero branch.

\subsection{4. The rational Hermitian form and the finite component}\label{the-rational-hermitian-form-and-the-finite-component}

Take an arbitrary complex point \(x\in X(\mathbb C)\). Let \((V_x,\Lambda_x,\Omega_x,J_x)\) be its rational homology, integral homology, alternating form, and complex structure. Use the analogous notation for the seed. Both lattices are integral, projective over \(\mathcal O_E\), and selfdual for their exact alternating forms. They have the same rank and prescribed signatures.

Choose any fixed nonzero \(\xi\in E\) with \(\bar\xi=-\xi\). The trace pairing extracts a unique Hermitian form \(h\) from \(\Omega\), with one common convention in the two spaces. An \(E\)-linear map preserves \(h\) exactly when it preserves \(\Omega\): preservation of \(\Omega(av,w)\) for every \(a\in E\) and nondegeneracy of the trace pairing prove the converse.

At every finite place the two lattices, and hence the two rational forms, are isometric. This is the local selfdual theorem in LIFT084 §4, which includes the dyadic case and the absolute inverse different. Its hypotheses hold because the extension \(E/F\) is unramified. At infinity the fixed \nolinkurl{DetCond} gives equal Hermitian signatures. Therefore the torsor of exact \(E\)-linear isometries is locally trivial everywhere.

For completeness, its Hasse principle applies in both odd and even ranks. The group \[
 G=\operatorname{Res}_{F/\mathbb Q}U(h),\qquad
 G^{\mathrm{der}}=\operatorname{Res}_{F/\mathbb Q}SU(h),\qquad
 G/G^{\mathrm{der}}=\operatorname{Res}_{F/\mathbb Q}U(1)
 \tag{HF12}
\] has simply connected derived group; in rank one the derived group is trivial. The first cohomology of the last torus is \(F^\times/N_{E/F}E^\times\). The cyclic Hasse norm theorem makes its localization map injective. Milne's SVI Lemma 8.20 now makes the isometry torsor trivial. We obtain an exact rational isometry \[
                    (V_0,\Omega_0)\simeq_E(V_x,\Omega_x).
 \tag{HF13}
\] No even-rank-only consequence for a similitude group is used. We henceforth make HF13 and work in one rational space \((V,\Omega)\), while retaining \textbf{both} lattices \(\Lambda_0\) and \(\Lambda_x\).

For each rational prime in \(S\), choose a local unitary isometry \(k_p\) carrying \(\Lambda_{0,p}\) onto \(\Lambda_{x,p}\). At primes dividing \(M\), it can also carry \(\eta_0\) to \(\eta_x\). Here is the required local reason. Coprimality of \(\mathfrak M\) and \(\bar{\mathfrak M}\), with product \(M\), implies that every factor above such a prime is split in \(E/F\). The alternating pairing identifies the two integral halves as dual modules. The actual \(\eta\) is a complete basis identification on one half modulo the relevant prime power, by the \nolinkurl{TateIso.ker}, \nolinkurl{onto}, and geometric compatibility clauses. An arbitrary invertible change of that half lifts to its integral general linear group; the inverse adjoint action on the other half preserves \(\Omega\). Hence the required level matching is possible. We only use \(A=1\): there is no additional full-level root condition to match.

These conditions define nonempty open subsets of \(G(\mathbb Q_p)\). They fix the target lattice and exact finite decoration at the protected primes; they do not identify all integral genera globally.

\subsection{5. An algebraic period family through the arbitrary target}\label{an-algebraic-period-family-through-the-arbitrary-target}

Let \(\mathcal D\) be the space of positive compatible complex structures on \((V_\mathbb R,\Omega)\) that commute with \(E\) and have the prescribed multiplicities. It is a product of unitary Hermitian domains, \[
              \mathcal D\simeq\prod_{v\mid\infty}
                  U(p,q)/(U(p)\times U(q)).
 \tag{HF14}
\] The real group \(G(\mathbb R)\) acts transitively. Its standard realization as a product of matrix balls makes the weight-one Hodge filtration holomorphic. If \(pq=0\), the domain is a point, and the arguments below reduce to that case without a positive-dimensional density assertion.

Keep \(\Lambda_x\) and \(\eta_x\) fixed. Choose a neat finite-index subgroup \(\Gamma\) of their unitary stabilizer. This is an auxiliary proof cover, not a change to the output level. Such a subgroup exists by SVI Proposition 3.5. The kernel of the action on the domain is compact; its intersection with the lattice stabilizer is finite, and neatness removes it. The image is arithmetic. Baily--Borel, as stated in SVI Theorem 3.12, gives the nonsingular algebraic variety \[
                           Y=\Gamma\backslash\mathcal D.
 \tag{HF15}
\] On \(Y\), the local system defined by \(\Lambda_x\) and the tautological filtration is an integral polarized variation of Hodge structures of weight one. SVI Theorem 14.8 gives an \textbf{algebraic abelian scheme} over \(Y\). Full faithfulness transports the flat integer-ring action and the polarization; unimodularity makes it principal. Thus the algebraicity in this argument is supplied by the algebraic weight-one theorem, not inferred from density of an analytic orbit.

Apply HF7 to this algebraic family. It supplies the actual rigidified symmetric principal line, cubical structure, all adjoint identities, and prim identity. The constant finite \(\eta_x\) classes are \(\Gamma\)-invariant and give algebraic torsion sections. One way to verify algebraicity is to regard each analytic section of the finite étale torsion cover as a connected component of degree one; connected components of a complex algebraic variety agree with analytic connected components. Only finitely many generators are needed. Set \(\eta_A=0\). Each geometric fibre has the source \nolinkurl{TateIso} by the local selfdual classification and the exact level extension argument in §4. The theta commutator is precisely the one checked in LIFT084 MP19--20; a compatible primitive root tower exists over the algebraically closed residue field, and its index-one value is one. No full adelic trivialization over a number field is being asserted.

The tangent multiplicities are constant on \(\mathcal D\) and give all the literal \nolinkurl{LieMult} clauses. We have therefore constructed a genuine \nolinkurl{Obj} over the reduced algebraic base \(Y\), with its map to \(\operatorname{Spec}(KA(l))\) through the fixed complex embedding. The exact \nolinkurl{Moduli.rep} hypothesis supplies a regular map \[
                              \Phi_x:Y\longrightarrow X_\mathbb C.
 \tag{HF16}
\] Its fibre at \([J_x]\) is the original object at \(x\): HF6 identifies the actual line, and the lattice, action, and finite \(\eta_x\) are unchanged. Uniqueness in \nolinkurl{Moduli.rep} therefore gives \(\Phi_x([J_x])=x\).

This construction is repeated for an arbitrary target \(x\), with its own lattice and level. It is not a claim that the seed's single arithmetic quotient surjects onto \(X\). It is also independent of the source's \nolinkurl{IsUnf}, \nolinkurl{N1c}, or \nolinkurl{FamilyUnif} constructors. Those analytic record ties remain separate inputs in any use of LIFT082.

\subsection{6. Rational approximation with the exact finite conditions}\label{rational-approximation-with-the-exact-finite-conditions}

We use a direct weak-approximation argument for \(G\). In the \(\mathbb Q\)-vector space \[
 \mathfrak u=\{Z\in\operatorname{End}_E(V):Z^\dagger=-Z\},
 \qquad c(Z)=(1-Z)(1+Z)^{-1},
 \tag{HF17}
\] the Cayley transform is a rational isomorphism between the open set \(\det(1+Z)\ne0\) and the open unitary set \(\det(1+c(Z))\ne0\). Its inverse is the same displayed fractional expression in the unitary matrix. The identities follow by taking the adjoint and commuting polynomials in one matrix.

This open set is dense in every real or finite local group. To check density even near an element with eigenvalue \(-1\), multiply it by the central norm-one element \[
                         z(t)=\frac{1+t\jmath}{1-t\jmath},
 \qquad \bar\jmath=-\jmath,\quad \jmath\ne0.
 \tag{HF18}
\] As \(t\) tends to zero locally, this tends to one. The determinant \(\det(1+z(t)g)\) is a nonzero rational function of \(t\), so only finitely many local values are excluded. It is nonzero as a function because \(z(t)\) is nonconstant and \(\det(1+zg)\) is a nonzero polynomial in \(z\). For restriction of scalars, make this argument at each of the finitely many factors. Weak approximation in the rational vector space \(\mathfrak u\) now proves density of \[
          G(\mathbb Q)\longrightarrow
          G(\mathbb R)\times\prod_{p\in S}G(\mathbb Q_p).
 \tag{HF19}
\] This is weak approximation at a finite set. It is not strong approximation with integrality at every other prime.

Choose \(x\in U_{F_0}(\mathbb C)\). By HF16, the inverse image of that open set is a nonempty open subset of \(Y\); its inverse image \(W\) in \(\mathcal D\) contains \(J_x\). Since the real unitary group is transitive, there is a real group element sending the seed complex structure into \(W\). Approximate it using HF19 while imposing all the open finite conditions from §4. We obtain \(g_1\in G(\mathbb Q)\) such that \[
 J_1=g_1J_0g_1^{-1}\in W,\qquad
 g_1\Lambda_{0,p}=\Lambda_{x,p}\ (p\in S),\qquad
 g_1\eta_0=\eta_x\text{ modulo }M.
 \tag{HF20}
\] The torus \((V_\mathbb R,J_1)/\Lambda_x\), with its HF5 line and finite \(\eta_x\), is the fibre of HF16 at \([J_1]\), so its source point belongs to \(U_{F_0}\).

The two rational lattices are commensurable. Because HF20 is an integral isomorphism at every protected prime, a denominator \(d\) with \(d g_1\Lambda_0\subset\Lambda_x\) can be chosen prime to \(S\). Multiply it by a positive integer prime to \(S\), chosen by the Chinese remainder theorem, so that the resulting \(N\) satisfies \(N\equiv1\pmod M\). Then \[
                    u=Ng_1:A_0\longrightarrow A_1
 \tag{HF21}
\] is a holomorphic \(\mathcal O_E\)-equivariant isogeny with the exact finite \(\eta\) equality. Riemann's theorem makes it algebraic. Since \(g_1\) preserves \(\Omega\), its polarization multiplier is \(N^2\). The semicharacters satisfy \[
 \alpha_1(Ng_1\lambda)
   =(-1)^{N^2\Omega(\lambda,w\lambda)}
   =\alpha_0(\lambda)^{N^2}.
 \tag{HF22}
\] The Hermitian forms have the same \(N^2\) factor. Appell--Humbert gives the exact rigidified line isomorphism in HF3; uniqueness also preserves the cubical structure. The principal polarization degree formula gives \(\deg u=N^{2g}\). The polarized adjoint isogeny \(u^\dagger\) satisfies \[
                      u^\dagger u=[N^2],
                      \qquad \ker u\subset A_0[N^2].
 \tag{HF23}
\] In particular \(A_1\) is CM and remains in the specified seed isogeny class. All the preceding arguments work with an arbitrary nonempty open subset of \(X(\mathbb C)\), proving the density assertion.

\subsection{7. The finite field and a genuine nonzero frame}\label{the-finite-field-and-a-genuine-nonzero-frame}

Let \(K_1=K_0(A_0[N^2])\). By HF23, the kernel of \(u\) is a subgroup of a torsion group whose every point is defined over \(K_1\). In characteristic zero it is a finite étale subgroup scheme over that field. Hence the quotient \(A_1=A_0/\ker u\) and \(u\) are defined over \(K_1\). This is precisely the quotient-over-the-ground-field assertion in Milne, AV, ChapterIV, Lemma2.1, printed135, with the finite quotient construction in ChapterI, Proposition8.10 and Remark8.12, printed38--39. No finiteness theorem for the number of isogeny classes is used here. The subgroup is \(\mathcal O_E\)-stable, so the full action descends to the quotient. Its identities descend from the seed action.

The analytic principal polarization on the quotient is also defined over \(K_1\). It is the unique homomorphism \(\lambda_1\) satisfying \(u^\vee\lambda_1u=N^2\lambda_0\); equality after composition with an isogeny is injective on homomorphisms. Homomorphisms of abelian varieties over an algebraically closed characteristic-zero field do not acquire new elements after an algebraically closed extension, so this uniqueness gives the required Galois descent. Formula HF7 then supplies the actual line over \(K_1\), including its origin rigidification and all normalized structural identities. Thus there is no extra field extension to select an arbitrary symmetric line.

The finite decoration is \(\eta_1=u\eta_0\), and is therefore defined over \(K_1\). It agrees with the target decoration in HF20 because \(N\equiv1\pmod M\). All source \nolinkurl{Obj} clauses hold geometrically and hence for this descended object. The exact \nolinkurl{Moduli.rep} over \(\operatorname{Spec}K_1\) gives its point of \(X\); base change and uniqueness identify it with the complex point already in \(U_{F_0}\).

The Galois action on the free \((\mathbb Z/N^2\mathbb Z)\)-module \(A_0[N^2]\) is faithful on its field of definition. This proves HF4. It covers the quotient, action, polarization, canonical line, and finite \(\eta\), rather than only the abstract subgroup.

Because \(K\subset K_1\) already contains every embedding label of \(E\), all eigenspaces of the quotient tangent space are defined over \(K_1\). Each has a basis there. The principal theta section space is one-dimensional over \(K_1\). Thus the frame torsor fibre has a \(K_1\)-point. Its nonvanishing open has a \(K_1\)-point by §3. This proves the last clause of HF3 with a literal \nolinkurl{FrPt}, \nolinkurl{FromTest}, and nonzero \nolinkurl{evR}, without another field extension.

For a weight rule \(F_0\in\operatorname{sec}(a,n)\), there is also a bounded integer frame change relative to \textbf{any} initial rational eigenframe at this chosen point. Keep a nonzero theta section fixed. The actual \nolinkurl{exists_isFrPoly} theorem gives its frame polynomial \(G\), and \nolinkurl{isFrPoly_equiv_and_ne_zero} shows that \(G\ne0\) because the frame fibre contains a nonzero value. Put \[
                 D=\sum_\tau\sum_{i<Dw.mult(\tau)}a_{\tau,i}.
 \tag{HF24}
\] The exact equivariance implies that \(G\) is homogeneous of total degree \(D\), and has degree \(a_{\tau,i}\) in physical column \(i\). The polynomial \(G\prod_\tau\det Y_\tau\) is nonzero and has degree at most \(D+1\) in each individual matrix entry. A nonzero polynomial of individual degrees at most \(D+1\) cannot vanish on the entire Cartesian grid \(\{0,\ldots,D+1\}\) in characteristic zero; induction on the variables proves this. Consequently there are integer matrices \(B_\tau\), with \[
 0\le (B_\tau)_{ij}\le D+1,\qquad \det B_\tau\ne0,\qquad G(B)\ne0.
 \tag{HF25}
\] Apply the literal transpose in \nolinkurl{IsFrPoly}, namely \nolinkurl{FrChg (B.transpose) 1}. This gives a nonzero frame with no section rescaling. For a nonempty block of size \(m\), \[
 \|B_\tau\|_{\mathrm{op}}\le m(D+1),\qquad
 1\le|\det B_\tau|\le m!(D+1)^m,
 \tag{HF26}
\] and the entries of its inverse have denominators dividing that actual determinant. This is a relative frame-change bound. It gives no bound on the initial rational eigenframe or on the section generator.

\subsection{8. Exact cost and application boundary}\label{exact-cost-and-application-boundary}

The construction preserves the following data and changes the following objects. The distinction is necessary for a later use of LIFT080--082 or any DP44 comparison.

\begin{longtable}[]{@{}
  >{\raggedright\arraybackslash}p{(\columnwidth - 4\tabcolsep) * \real{0.3333}}
  >{\raggedright\arraybackslash}p{(\columnwidth - 4\tabcolsep) * \real{0.3333}}
  >{\raggedright\arraybackslash}p{(\columnwidth - 4\tabcolsep) * \real{0.3333}}@{}}
\toprule\noalign{}
\begin{minipage}[b]{\linewidth}\raggedright
Item
\end{minipage} & \begin{minipage}[b]{\linewidth}\raggedright
What is proved here
\end{minipage} & \begin{minipage}[b]{\linewidth}\raggedright
What remains to be supplied in an arithmetic application
\end{minipage} \\
\midrule\noalign{}
\endhead
\bottomrule\noalign{}
\endlastfoot
Host, Datum, signature, coefficient labels & The same \(E,H,Dw,c\), \(p,q\), and every exact eigenlabel & An actual admitted Host/Sizes family and the original source pins \\
Finite level & The original \(M,\mathfrak M,A=1\), with \(u\eta_0=\eta_1\) & Compatibility with any additional arithmetic structures not present in this \nolinkurl{Obj} \\
Abelian variety & A CM variety in the fixed seed's \(\mathcal O_E\)-isogeny class & It need not equal the seed, its old LamPt, or any specified product polarization \\
Principal line & \(u^*\mathcal L_1=\mathcal L_0^{N^2}\) rigidified and cubical & The source's chosen theta section is not transported as the same section of \(\mathcal L_0\) \\
Isogeny and field & \(\deg u=N^{2g}\), \([K_1:K_0]\le N^{8g^2}\) & \(N\), the initial seed-field degree, and all numerical height effects remain unevaluated \\
Protected primes & \(N\) is prime to every prime in \(S\), including 11 if desired & No full selected-jet, gauge, or content invariance follows just from this fact \\
Archimedean tangent metric & \(H_1(duv,duw)=N^2H_0(v,w)\), and the source \(\pi H\) norm scales by \(N\) & Integral eigenline degrees and actual coefficient ideals on the new models \\
Nonzero evaluation & The same nonzero natural rule has a literal field-valued frame with \(\operatorname{evR}\ne0\) & A proof that the prescribed arithmetic image/rule is itself nonzero \\
Frame change & HF25--26 relative to any rational eigenframe; no further extension & Initial frame and theta-generator heights, inverse determinants, model ties \\
Step5 & The qualitative same-point nonzero condition can be supplied for this new point & Every remaining FT1--6 geometric/model/section/letter premise, full extension quantifiers, and revised costs \\
\end{longtable}

Even if \(S\) contains 11, the equality of line powers in HF3 changes the actual theta realization under pullback. It does not identify the new selected differential with the old one. The old pure-type \nolinkurl{LamObjGood} boundary is unchanged. No central \(L\)-value nonvanishing, prescribed period, arithmetic \nolinkurl{Read}/\nolinkurl{STiesR}/032 instance, universal DP44 sign, or RH conclusion follows.

The theorem proves an existential detector for a \textbf{given nonzero rule} in a fixed isogeny class. It is not an effective hitting theorem for its nonvanishing open. The weak approximation and open-set argument leave precisely that denominator bound uncomputed.

\subsection{9. Algebraic-geometric translation, literal interfaces, and verification scope}\label{algebraic-geometric-translation-literal-interfaces-and-verification-scope}

In the run's Grothendieck language, HF7 is a construction in the rigidified relative Picard functor from the normalized Poincaré biextension, and HF21 is an actual isogeny of the original abelian schemes over a finite extension, with its line-power identity and finite level retained. The inherited contextual anchors are FGA, Bourbaki232, Theorem3.1, pp.148--149; Bourbaki236, Corollaries1.7 and2.3, pp.229 and232; and SGA7 I, ExposéVII,2.7 and2.9, pp.169--174, and ExposéVIII,7.4(b), pp.310--311. These are the previously checked run locators, not a claim of a new reading here. The precise relative Picard/base-change and Poincaré inputs used in this proof are the freshly read references below. The eigenspace bundles and their frame torsor are actual vector bundles, rather than a virtual Grothendieck-group equality; the inherited module context is SGA3 I,4.7.3--4.7.3.1.

For a protected prime \(p\in S\), HF23 also has the following strictly conditional integral interpretation. Whenever the isogeny and its adjoint extend to the relevant abelian-scheme models, their induced maps on the \(p\)-divisible groups are inverse after composing the adjoint with the automorphism \([N^2]^{-1}\), since \(p\nmid N\). This elementary consequence does not construct those models or identify their chosen theta sections, selected derivatives, or integral eigenlines. The carried formal-group/Cartier-duality context remains cote4, scan003, §§8.1.1--8.1.2 and scans005--006, §§8.4.1--8.4.3, and cote54, scans005--007. No new crystal comparison, motivic identification, weight or purity theorem, motivic \(L\)-function identity, or RH conclusion is asserted. DP\_MULT Q3e and THETA3 G1/A7 gain only the conditional same-rule nonzero point within the fixed CM isogeny class; the actual arithmetic image and all quantitative DP44 costs remain separate. CX\_GRH CX-ALG-024 and GCT\_PLETH T-052 retain their distinct arithmetic comparison gates.

All source paths are relative to the work tree and are fingerprinted in the adjacent source manifest. The principal interfaces are:

\begin{itemize}
\tightlist
\item
  \nolinkurl{TwoWeightEqOneAddRoot.lean}: Host14--27, Datum30--47, Level51--61, SAB72--84, DetCond119--122, TateIso131--142, Obj144--160, ObjMap162--166.
\item
  \nolinkurl{PartitionExtendByZeroAntitone.lean}: \nolinkurl{Par.PinLevel}43 and \nolinkurl{PinDat}47; the pinned application has \(A=1\).
\item
  \nolinkurl{PushInvertible.lean}: Moduli29--38, including representability on every reduced base; Coeff40--45 and XF49.
\item
  \nolinkurl{IsCombAt.lean}: tangent/eigen/basis39--54, FrPt67--77, exact FrMap79--85, FrBor87--95, sec100 and evR103.
\item
  \nolinkurl{AlgebraMapCAt.lean}: TRing9--12, naturality35--36, FromTest51--52, evalAt54--58, and EvalWD84--87. The default-zero evaluation branch is explicitly excluded in §3 and §7.
\item
  \nolinkurl{IntegralOf.lean}: actual \nolinkurl{IsFrPoly}97--102 and its transpose convention; \nolinkurl{ExistsIsFramePolynomialOfRuleNumberField.lean}50; \nolinkurl{FramePolynomialEquivariantAndNeZero.lean}62. These generic source theorems apply to the newly constructed rational frame, without the old pure-type \nolinkurl{CMPt.Good} record.
\item
  Canonical LIFT084 MP9--20 supplies the local selfdual classification, dyadic trace-one proof, canonical source semicharacter, and exact theta commutator. LIFT082 records the separately retained local/archimedean application hypotheses.
\end{itemize}

Primary references directly read for the argument:

\begin{enumerate}
\def\labelenumi{\arabic{enumi}.}
\tightlist
\item
  Milne, J. S. (2017 revision of 2004). \emph{Introduction to Shimura varieties}. \href{https://www.jmilne.org/math/xnotes/svi.pdf}{Author PDF}. Proposition3.5, printed34; Theorem3.12, printed38; integral refinement in §6.11, printed74--75; Theorem8.17 and Lemma8.20, printed88--89; Theorem14.8, printed122. The imported statements are neat-subgroup existence, algebraicity of arithmetic domain quotients, the rational-form qualifier, its Hasse criterion, and algebraicity/full faithfulness of polarized weight-one families. The rational form, exact level, line, and source representability checks above are additional arguments.
\item
  Milne, J. S. (2020). \emph{Class field theory} (Version4.03). \href{https://www.jmilne.org/math/CourseNotes/CFT.pdf}{Author PDF}. ChapterVIII, Theorem3.1, printed234: the cyclic Hasse norm theorem and its cohomological proof.
\item
  Milne, J. S. (2008). \emph{Abelian varieties} (Version2.00). \href{https://www.jmilne.org/math/CourseNotes/AVc.pdf}{Author PDF}. §8, printed37--40: duality, normalized Poincaré bundle, and the polarization pullback identity; Theorem11.1, printed54: polarization degree and cohomology; ChapterI, Proposition8.10 and Remarks8.11--8.12, printed38--39, and ChapterIV, Lemma2.1, printed135: the finite subgroup quotient and its ground-field descent. Riemann equivalence and isogeny compatibility are also stated in SVI6.8--6.9.
\item
  van der Geer, G., \& Moonen, B. (2011 draft). \emph{Abelian varieties}, ChapterVI, §6.2--6.3, printed88--89. \href{https://www.math.ru.nl/~bmoonen/BookAV/DualAV1.pdf}{Chapter PDF}. The rigidified relative Picard functor, its quotient by line bundles from the base, and representability/separatedness are used for HF7 and its uniqueness.
\item
  The Stacks Project Authors. \emph{Cohomology and base change: Applications}, \href{https://stacks.math.columbia.edu/tag/0BDM}{Tag0BDM}, Lemmas36.32.1--36.32.4. The rank-one direct image and arbitrary base change in HF9 follow from these statements and the fibre cohomology calculation.
\end{enumerate}

The companion executable is an exact diagnostic certificate: it checks source fingerprints, Cayley/adjoint identities on rational models, the semicharacter power and trace-one identities, local split frame changes, protected-prime denominator clearing, finite torsion-field degree bounds, and the Cartesian-grid argument on finite polynomial models. It does not numerically establish the Hasse principle, Baily--Borel, the weight-one algebraization theorem, or a uniform denominator bound. Those are proved or explicitly imported in the text above. Final author and independent whole-file verification records are separate.

\subsection{Canonical algebraic-geometric context}\label{canonical-algebraic-geometric-context-15}

The map is an actual isogeny of abelian varieties over a number field, with OE action, finite level, principal polarization and the uniquely reconstructed trace-one cubical line. Its construction uses the rigidified relative Picard functor and normalized Poincare biextension: inherited FGA232 Theorem3.1 pp148--149, FGA236 Corollaries1.7/2.3 pp229/232; SGA7 I VII2.7/2.9 pp169--174 and VIII7.4(b) pp310--311. Actual eigenbundles and frame torsors retain SGA3 I4.7.3/4.7.3.1 and SGA6 VI Theorem3.3 as context. At protected primes the p-divisible consequence assumes extension to actual abelian models; cote4 scan003 sections8.1.1--8.1.2, scans005--006 sections8.4.1--8.4.3 and cote54 scans005--007 remain contextual. No new integral crystal, lambda-ring/K-theory, l-adic weight or CM-motivic L-function comparison is claimed. DP\_MULT Q3e/DP44/DP52/DP54--65 and THETA3 G1/A7 retain the quantitative and arithmetic-image gates; CX\_GRH CX-ALG-024 and GCT\_PLETH T-052 retain their distinct global comparison gates.

\clearpage

\section{AM-K-102 --- positive-norm cone counting for the dense all-II remainder}\label{am-k-102-positive-norm-cone-counting-for-the-dense-all-ii-remainder}

\textbf{Canonical checked integration --- 2026-11-23T03:14:01.090065+00:00.} Complete root and independent proof/script/interface reviews and separate exact replays passed. The actual canonical predecessor has101/31 entries; this integration gives102/31, complete floors25/11. Frozen author history is preserved. AM-K-102 PC1--27 bounds the entire actual DO107 IP3 support25 remainder by2\^{}-3550: all-II globally rank-two-unipotent scalar-cross, exactly25 total eligible places, no inactive paired cross zero,2\^{}120\textless N\_E D\_eta\textless2\^{}744 and100\textless=k\textless=164. The relative-norm phase count, positive-cone Gaussian/Poisson majorant and true transverse metric use one joint energy constraint; the complete original complex-volume tail covers every larger ideal-norm parameter. The original25 active weighted-label factors and143 inactive pointwise factors are retained without assuming complete-orbit constancy. All195 endpoint bounds are below-3557; the65-cell union costs3550 once. Actual DO108101/31 becomes102/31, every old overlapping charge retained, both complete sums below2\^{}-230 and floors25/11 until a separate exhaustive assembly. No signed positivity, prescribed period, selected arithmetic-image witness, effective admitted source cost, specific032, DP44 or RH claim follows.

Author: GPT-6 Astra Pro, /root/kummer\_extension, ALP\_MACHINE. 23 November 2026 UTC.

\textbf{Complete author-checked proof; independent frozen-file review pending.} The companion performs exact arithmetic after the general arguments below. The author verification records the whole saved-proof/script check and final replay separately. No numerical diagnostic alone certifies the theorem.

\subsection{1. Actual sector, unchanged inputs and statement}\label{actual-sector-unchanged-inputs-and-statement}

Keep the actual J56 all-top kernel, original point, selected real module, both compact factors, central classes, canonical twists, finite projectors, source lattices, pair weights and masses. Put \[
 f=168,\quad m=164,\quad H=133,\quad R=H+H^{-1}-2=17424/133,
 \quad Q=g_0g_1,\quad U=77Q^2-14,\quad M^2=R/U.
 \tag{PC1}
\] Consider DO107's residual IP3 sector: support \(s=25\), all type II, global rank-two unipotent scalar-cross, full positive identity compression, total eligibility exactly25, no paired-C-zero inactive place, and \[
 2^{120}<x=N_ED_\eta<2^{744},\qquad100\le k\le164.          \tag{PC2}
\] Here \(k\) counts nonzero original S-coordinate covector values. DO105 fixes a unique fully eligible \(\eta\) for each support. DO105 EL13--16 proves that each retained active type-II term has at most one actual diagonal label, including a central-only pole, and has exact absolute coefficient \((T-1)^{-1}\). Its complete weighted-label sum is therefore at most \((T-1)^{-25}\). DO107 IP3--8/13 supplies the143 inactive factors \((T-1)^{-1}\) pointwise because total eligibility is25 and no inactive place has paired-C-zero. These two statements give the complete factor \((T-1)^{-168}\) used here. This is a bound on the retained support terms; DO106's possible variation on a complete central orbit includes other points and does not contradict it. No complete-orbit constancy is asserted or needed. No integral moving frame is introduced. DO108's later active-pole cut is unnecessary for this estimate but its prior error is retained.

For each fixed common G orientation set \[
 L=J I_\Sigma D_\eta,\quad N_EL=z^{2f},\quad
 N_EJ=z^{2f}T^{50}/x,\quad p=M^2/(kz^2).                  \tag{PC3}
\] The original minimum gives \(p\ge1\), and the original G cap gives \[
 z^2\le(Q^{334}T^{-50}x)^{1/f},\qquad
 kp>68036(2^{120}/x)^{1/f}.                               \tag{PC4}
\] The second comparison is checked again by exact host integer powers.

\textbf{Theorem.} The entire actual PC2 sector has absolute contribution less than \(2^{-3550}\). Remove only its intersection with the current remainder; retain every old charge. A higher full support floor requires a separate exhaustive assembly.

Use the unchanged DO104/DO105 original-coordinate reconstruction. The167 covector values comprise the \(k\) active original S coordinates and three original transverse values in their true Gram metric. The latter lattice has \[
 \operatorname{covol}\Lambda_{\rm tr}
 =D_0^3z^{6f}/(4U)^f,\quad
 \mu(\Lambda_{\rm tr})<\epsilon_0\sqrt f\,z,\quad
 \epsilon_0=1/1024,\quad D_0=7^{140}29^{144}5^{84}13^{84}.
 \tag{PC5}
\] The full joint squared norm is at most \(fM^2\). Given \(\eta,r,d\), the normal form reconstructs \(h\); there is no extra isometry fibre. The harmless old pivot factor \(164<2^8\) is retained.

\subsection{2. Relative norm phases}\label{relative-norm-phases}

Let the real fractional ideal \(\mathfrak C\) satisfy \(\mathfrak C\mathcal O_E=L\bar L\). Relative ideal norm gives \[
 N_F\mathfrak C=N_EL=z^{2f},\quad
 \operatorname{covol}\mathfrak C=\sqrt{D_F}\,z^{2f},
 \qquad D_F=2^fD_0.                                      \tag{PC6}
\] For nonzero \(r_j\in L\), \(t_j=r_j\bar r_j\in\mathfrak C\) is totally positive. The integral E-ideal \((r_j)L^{-1}\) has relative norm \((t_j)\mathfrak C^{-1}\). At a split F-prime of exponent \(e\) there are \(e+1\) choices; at nonsplit primes there is at most one. Generators of a fixed ideal with the same exact relative norm differ by a unit of absolute value one at every embedding, hence by one of the28 roots of unity. Therefore the accepted K075 argument gives \[
 \#\{r_j\in L:r_j\bar r_j=t_j\}
 \le28\,d_F((t_j)\mathfrak C^{-1}).                       \tag{PC7}
\] This uses relative F divisors, not a crude E-divisor count.

DO091 TD11 proves the complete F-prime norm list below4096:29 with multiplicity24,1681 with multiplicity84,2129 with multiplicity168. For exponent \(1/11\), norms at least2048 contribute no constant because \(e+1\le2^e\le q^{e/11}\). The local maxima for29 and1681 occur at \(e=2\) and \(e=1\), by consecutive ratios. Consequently \[
 d_F(\mathfrak a)\le C_{11}N_F(\mathfrak a)^{1/11},\qquad
 C_{11}^{11}=\frac{3^{264}2^{924}}{29^{48}41^{168}}
                         <2^{220}.                     \tag{PC8}
\] Every comparison is checked by integer powers; no new field or splitting assumption is introduced.

The joint energy and arithmetic-geometric mean over its \(fk\) positive real coordinates give \[
 \prod_j N_F((t_j)\mathfrak C^{-1})
 =z^{-2fk}\prod_{j,\tau}t_{j,\tau}\le p^{fk}.              \tag{PC9}
\] Thus all phase multiplicities together are at most \((28\,2^{20})^k p^{fk/11}\). Pull this bound out only on the original bounded set. When the lattice sums below are enlarged, no divisor assertion is made for added zero or oversized tuples.

\subsection{3. Positive-cone Gaussian majorant}\label{positive-cone-gaussian-majorant}

In the full real trace metric, \[
 \mathfrak C^*=\mathfrak D_F^{-1}\mathfrak C^{-1},\qquad
 \lambda_1(\mathfrak C^*)^2
 \ge f/(\operatorname{rd}(F)^2z^4).                      \tag{PC10}
\] This follows from the fractional-ideal norm minimum and arithmetic-geometric mean. There is no half-trace factor2. The actual field has \(\operatorname{rd}(F)<1500\).

Set \(\sigma=250z^2\). The Gaussian dual mass satisfies \[
 \Theta=\sum_{\xi\in\mathfrak C^*}
                    e^{-2\pi^2\sigma^2\|\xi\|^2}<4.      \tag{PC11}
\] For the exact constants, let \(u=(2\pi250/\operatorname{rd}(F))^2>u_0=24649/22500\). The accepted Gaussian tail (Miller and Stephens-Davidowitz,2019, equation1.1) at the nonzero minimum gives \(\Theta-1\le\eta\Theta\), with \(\eta=(eu)^{f/2}e^{-fu/2}\). It also follows directly by splitting a Gaussian at scales \(s\) and \(ts\), using the minimum for the remaining tail, and using Poisson summation to bound the scale-\(ts\) mass by \(t^{-f/2}\) times the scale-\(s\) mass; optimize \(t=f/(2s)\). The function \((eu)^{f/2}e^{-fu/2}\) decreases for \(u>1\). For \(w=u_0-1\), \(\log(1+w)\le w-w^2/2+w^3/3\), and exact rational comparison gives \(84(w^2/2-w^3/3)>1/3\). Hence \(\eta<e^{-1/3}<3/4\), proving PC11.

For any \(a>0\), take the normalized one-dimensional Gaussian \(g_\sigma\) of variance \(\sigma^2\), and put \[
 v=17/20,\quad h=\sigma v+a\sigma^2,\qquad
 F_h(y)=e^{-ay}\mathbf1_{y\ge-h}.
\] Completing the square gives, for \(y\ge0\), \[
 (F_h*g_\sigma)(y)
 =e^{-ay+a^2\sigma^2/2}\Phi((y+h)/\sigma-a\sigma)
 \ge e^{-ay+a^2\sigma^2/2}\Phi(v).                        \tag{PC12}
\] Integrating the lower Taylor bound for \(e^{-t^2/2}\), and using \(\pi<22/7\), gives \(\Phi(17/20)>4/5\). The exact rational test is \[
 I=v-v^3/6+v^5/40-v^7/336>0,\qquad
                        I^2>(3/10)^2(44/7).             \tag{PC13}
\] The convolution is Schwartz; its product therefore admits absolute Poisson summation. Its Fourier factors satisfy \(|\widehat F_h(\xi)|\le e^{ah}/a\) and \(\widehat g_\sigma(\xi)=e^{-2\pi^2\sigma^2\xi^2}\). Using PC6 and PC11--13, then setting \(a=A/z^2\), yields \[
 \sum_{t\in\mathfrak C,\ t_\tau\ge0}
 e^{-(A/z^2)\sum_\tau t_\tau}
 <\frac4{\sqrt{D_F}}
 \left\{\frac5{4A}
       \exp(250Av+250^2A^2/2)\right\}^{f}.                \tag{PC14}
\] The factor4 is once per F-lattice, not per real coordinate. Its scaling cancels against the covolume exactly. Inclusion of the origin only enlarges the original totally positive tuples.

\subsection{4. Joint S and transverse energy}\label{joint-s-and-transverse-energy}

Let \(\delta=1/32768\). For \(y\) in the Voronoi cell of \(\ell\in\Lambda_{\rm tr}\), \(\|\ell\|^2\ge(1-\delta)\|y\|^2-\mu_{\rm tr}^2/\delta\). Integrating all cells, using PC5 and \(\epsilon_0^2/\delta=1/32\), gives \[
 \sum_{\ell\in\Lambda_{\rm tr}}e^{-(A/z^2)\|\ell\|^2}
 \le\frac{(4U)^f}{D_0^3}
       \left(\frac{\pi}{A(1-\delta)}\right)^{3f}
       e^{Af/32}.                                       \tag{PC15}
\] The true Gram determinant is retained.

Use the single joint exponential majorant \[
 \mathbf1_{\sum t_{j,\tau}+\|\ell\|^2\le fM^2}
 \le e^{AfM^2/z^2}
       e^{-(A/z^2)(\sum t_{j,\tau}+\|\ell\|^2)}.
\] The complete row count, including the positions of the nonzero original S coordinates, is therefore \[
\begin{split}
 J_k(p;A)\le{}&
 {164\choose k}(28\,2^{22})^k p^{fk/11}D_F^{-k/2}
 \frac{(4U)^f}{D_0^3}
 \left(\frac{\pi}{A(1-\delta)}\right)^{3f}
 \left(\frac5{4A}\right)^{fk}\\
 &\times\exp\{Afkp+fk(250Av+250^2A^2/2)+Af/32\}.
\end{split}                                              \tag{PC16}
\] The factor \(2^{22}\) combines C11's \(2^{20}\) and the theta factor4. This uses the single actual energy constraint, rather than two independent radius bounds. Only after PC9 has bounded the original phase multiplicity are the Laplace sums enlarged.

\subsection{5. Complete inherited normalization}\label{complete-inherited-normalization}

The unchanged DO105 affine central fibre and actual selected trace give \[
 \#d\le V_F(H/2+750Uz^2)U^{-f}T^{50}x^{-1},\qquad
 |c(h)|<2^{-929f/2}U^{-f/2}z^{-f},
\] \[
 V_F(C)=\frac{((22/7)fC^2)^{f/2}}{(f/2)!\sqrt{D_F}}.
                                                               \tag{PC17}
\] The real-part shift is fixed after the row. This is a difference-module count, not a claim that \(d\) belongs to that ideal. The common Green bound is \[
              2Q^{-335(f-1)}(N_EJ)^{335/2}.                \tag{PC18}
\] Retain support \(2^{168}\), patterns \(4^{334}=2^{668}\), pivot \(2^8\), and this one Green2. All168 exact T factors together satisfy \((T-1)^{-168}<2T^{-168}\), with only this additional2. The original role, compact, mixed-projector and mass normalizations remain those of DO105/DO107; all remaining factors have modulus at most one.

For each \(k\), multiplying every factor gives \[
 2^{846-929f/2}Q^{-335(f-1)}T^{8257}x^{-337/2}U^{-3f/2}
 z^{334f}V_F(H/2+750Uz^2)J_k(p;A).                       \tag{PC19}
\] Patterns have already been counted. Inside each fixed pattern, \(p\) is fixed by its ideal norm; the alternative estimates below require no new pattern or endpoint union factor.

\subsection{6. All denominator and ideal-norm ranges}\label{all-denominator-and-ideal-norm-ranges}

Set \(p_0=68036/k\), \(P=10000\), \(\lambda=3/5\). If \(1\le p\le p_0\), choose the constant \(A=\lambda/p_0\). PC4 forces \(x\ge2^{120}(p_0/p)^f\). After substituting this in PC19, its nonconstant \(p\)-dependence is \[
 p^{(3/2+k/11)f}e^{Afkp}
                      V_F(H/2+750R/(kp)).                \tag{PC20}
\] The last factor decreases by at most \(p^{-f}\). The complete expression is strictly increasing, and is bounded at \(p=p_0,x=2^{120}\). This handles the coupled denominator cap without an additional union.

For \(p_0\le p\le P\), set \(x=2^{120}\) in its decreasing power, and choose \(A=\lambda/p\). The logarithm of PC19, apart from a constant, is \[
 f((1+1/11)k-164)\log p+b_1/p+b_2/p^2
                         +f\log(H/2+750R/(kp)),
 \qquad b_1,b_2>0.                                      \tag{PC21}
\] It is convex in \(\log p\): the linear term is affine, the reciprocal terms are positive exponentials, and the last term is log-sum-exp. The maximum is at \(p_0\) or \(P\).

For \(p\ge P\), use the original full complex \((k+3)\)-coordinate lattice, whose covolume is \(D_0^{k+3}z^{2(k+3)f}/(4U)^f\). The accepted arbitrary-ideal covering coefficient \(2250/7\) and the sharper transverse bound imply the safe total cover \((2250/7)\sqrt{f(k+3)}z\). The volume count is \[
 J_k^{\rm vol}(p)\le{164\choose k}
 \frac{((22/7)f[\sqrt{kp}+(2250/7)\sqrt{k+3}]^2)^{(k+3)f}
 (4U)^f}{((k+3)f)!\,D_0^{k+3}}.                         \tag{PC22}
\] After the remaining PC19 factors, its \(p\)-dependence is \(p^{(k-164)f}\) times a decreasing positive radius factor and a decreasing central factor. Since \(k\le164\), the bound is nonincreasing. Its maximum is at \(p=P,x=2^{120}\).

These alternatives cover every allowed \(x,z\), with no denominator-ideal enumeration or extra case multiplier.

\subsection{7. Exact certificate and one charge}\label{exact-certificate-and-one-charge}

At both cone endpoints the exponential parameter is \[
 S(k,p)=fk\left(\frac35+\frac{255}{2p}+\frac{11250}{p^2}\right)
                                      +\frac{3f}{160p}.  \tag{PC23}
\] Use \(\log_2e<13/9\), \(\pi<22/7\), \(V_F(68)<2^{479}\), and \[
 L_T=1019333/1024,\quad L_Q=287843/4096,\quad
 L_U=601356/4096,\quad L_R=28810/4096.                    \tag{PC24}
\] Group the full positive R and negative U powers before the directed substitutions. The common exponent for all three endpoint types is \[
 C_0=846+2f-929f/2+479+167fL_R-(335f/2)L_U
                 +8257L_T-335(167)L_Q-(337/2)120.         \tag{PC25}
\] Here \(C_0=1053636089/4096\). For the cone, add \(13S/9\) and the exact strict binary upper bound for \[
 {164\choose k}(28\,2^{22})^k D_F^{-k/2}D_0^{-3}
 \left(\frac{22\cdot32768}{7\cdot32767}\right)^{3f}
 A^{-(k+3)f}(5/4)^{fk}(kp)^{-167f}p^{fk/11}.              \tag{PC26}
\] The certificate handles the last fractional power by an exact integer-power ceiling divided by11. For PC22, both square roots are replaced by integer ceilings and the resulting rational value is compared by integer shifts. No floating logarithm is used.

Every one of the \(65\times3\) complete endpoint bounds is less than \(-3557\), so \[
                     65\,2^{-3557}<2^{-3550}.            \tag{PC27}
\] There is no factor3 for alternative endpoint bounds within one fixed ideal enumeration. The adjacent output records every exact rational cell. Its worst upper exponent is \(-232301590280014177/65174873927680\), at \(k=164,p=68036/164\). It also checks the phase maxima, C11, full-real metric, Gaussian/CDF rational inequalities, host cap, complete monotonicity signs and chronological arrays.

\subsection{8. Dependencies, chronology and remaining gates}\label{dependencies-chronology-and-remaining-gates}

The direct inputs are K075 RN17--19's relative-norm phase count; DO091 TD11's exact prime list; DO104 CD11--21's covector module, transverse metric, central difference and selected trace; DO105 EL16/22--27's complete reconstruction and weights; DO107 IP3--8/13's residual sector and inactive factors; and DO070 FT3--5's arbitrary-fractional-ideal cover. Their original fields, modules and hypotheses are retained. The new arguments are PC8's exponent, PC12--16's positive-cone majorant and joint count, and PC20--22's endpoint reduction.

The default predecessor is checks/do108\_joint\_remainder\_error.json, with101 all-top and31 uniform errors. Append3550 once and retain all previous charges, including451,421,254,392,55281. The script permits an unchanged-byte replay against a later complete prefix. Both dyadic sums remain below \(2^{-230}\). Floors25/11 are retained until a separate exhaustive assembly.

The ideal modules, involution, Hermitian polarization, theta line, selected Hilbert metric, finite projectors and class contractions remain the actual ones. The exact inherited contextual anchors are SGA3 I4.7.3/4.7.3.1; SGA6 VI Theorem3.3; FGA232 Theorem3.1 pp148--149; FGA236 Corollaries1.7/2.3 pp229/232; SGA7 I VII2.7/2.9 pp169--174 and VIII7.4(b) pp310--311; cote4 scan003 §§8.1.1--8.1.2 and scans005--006 §§8.4.1--8.4.3; and cote54 scans005--007. The carried run interfaces remain DP\_MULTQ3e, DP44/DP52/DP54--65, THETA3G1/A7, CX\_GRH CX-ALG-024 and GCT\_PLETH T-052. These are retained framework references, not new lambda-ring, K-theoretic, weight or motivic conclusions. No prescribed Bessel/Mram period, selected ordinary-derivative or arithmetic-content witness, effective source cost, Read/STiesR comparison, DP44 or RH conclusion follows.

\subsubsection{Primary reference}\label{primary-reference-2}

Miller, S. D., \& Stephens-Davidowitz, N. (2019). \emph{Kissing numbers and transference theorems from generalized tail bounds} (arXiv:1802.05708, version2). \url{https://arxiv.org/pdf/1802.05708}. Equation(1.1), printed pages1--2, is used with its exact Euclidean Gaussian scaling. PC11 gives its short Poisson derivation in the required specialization. All full-real covolume and dual-ideal normalizations are retained explicitly.

\subsection{Canonical algebraic-geometric context}\label{canonical-algebraic-geometric-context-16}

The positive real norms are sections of the original relative-norm ideal C with C O\_E=L bar(L), and its full trace dual is Different\_F\^{}-1 C\^{}-1. The original Hermitian metric, paired source lattices, polarization, theta line, Picard/biextension data, selected coefficient and finite projectors remain fixed. Exact inherited contextual anchors are SGA3 I4.7.3/4.7.3.1; SGA6 VI Theorem3.3; FGA232 Theorem3.1 pp148--149; FGA236 Corollaries1.7/2.3 pp229/232; SGA7 I VII2.7/2.9 pp169--174 and VIII7.4(b) pp310--311; cote4 scan003 §§8.1.1--8.1.2 and scans005--006 §§8.4.1--8.4.3; cote54 scans005--007. These are retained framework locators, not new lambda-ring/K-theory, l-adic weight, CM-motive or L-function conclusions. DP\_MULT Q3e and THETA3 G1/A7 gain this finite subset removal; CX\_GRH CX-ALG-024 and GCT\_PLETH T-052 retain their distinct gates.

\clearpage

\section{AM-R-087 --- complete support25 exhaustion and the support26 remainder}\label{am-r-087-complete-support25-exhaustion-and-the-support26-remainder}

\textbf{Canonical checked integration --- 2026-11-23T03:21:20.607320+00:00.} Complete author and independent whole-final-proof/script/interface readings and separate3140-assertion replays passed. Root read the entire independent report and verification before promotion. AM-R-087 J26.1--5 exhausts every original all-top support25 class at the unchanged H133 cutoff. The complete F/IC8 rank1..33 ranges force A=I; DO090/R080/R079 exhaust fixed S; DO099/DO095 give the precise cross rank2/3 alternatives; R082/R084 exhaust rank3; DO101 and DO102 remove nonunipotent rank2 and all non-all-II unipotent profiles; the exact DO103/104/105/107/K102 denominator, eligibility and coordinate partition exhausts the final all-II branch, including every equality boundary. The complete all-top floor is26 and the independent uniform floor remains11. Every102/31 chronological error entry is unchanged, both full sums stay below2\^{}-230 and no assembly charge is added. The actual sufficient signed gate is J26.2; global period, fully sized admissible source, selected arithmetic image, specific032/Read, actual DP44, Kummer-Hecke RH and the global oscillator comparison remain open.

Author: /root, GPT-6 Astra Pro; reasoning effort not externally exposed; ALP\_MACHINE. 23 November2026 UTC. Whole saved-proof/script reading, exact accounting and independent review are separate prerequisites for canonical promotion.

\subsection{1. The unchanged actual kernel and statement}\label{the-unchanged-actual-kernel-and-statement}

Retain the actual J56 all-top upper-endpoint kernel of K044, its already charged H=133 cutoff, f=a=168, n=169, and the original positive S of dimension164. Retain the original field, both CM components, all source and target pencils, finite projectors, central representative, packet point, selected Hilbert metric, phases, Haar measures and class masses. Write w(h) for the same normalized summand. R078 J25.1--4 supplies the complete support floor25. Subsequent containing-set removals act on that same kernel, with their complete chronological charges retained.

\textbf{Theorem.} After all checked removals through K102 there is an actual remaining subset R\_top\^{}(26) such that

\[
 \mathcal E_{\rm top}=1+\operatorname{Re}
       \sum_{h\in R_{\rm top}^{(26)}}w(h)+\epsilon,
 \qquad |\epsilon|<2^{-230},\qquad s(h)\ge26.       \tag{J26.1}
\]

No new error charge is added by this assembly. All102 chronological all-top charges and all31 uniform charges are retained exactly. The independent uniform H=8192 floor remains11. The sufficient inequality is now

\[
 \operatorname{Re}\sum_{h\in R_{\rm top}^{(26)}}w(h)
                         >-1+2^{-230}.             \tag{J26.2}
\]

J26.2 remains open. The assertion proved here is the exhaustive support25 reduction in J26.1, not its sufficient signed inequality. Every existing higher-support restriction remains available in its original scope. Nothing in this proof extends a support25-only numerical estimate to support26.

\subsection{2. All nonidentity positive compressions are already covered}\label{all-nonidentity-positive-compressions-are-already-covered}

Put A=h\_SS and r=rank\_E(A-I). The retained K068 LR3/LR21--22 residual has r\textless=33 and determinant-one representatives with all334 unmixed G reductions unipotent. K069 IC1--8 gives the exhaustive original fixed class F and exceptional isotropic-cross class IC8 when A is nonidentity. These conditions make no global-unipotence assumption. The following complete inclusive ranges each contain support25:

\begin{longtable}[]{@{}
  >{\raggedright\arraybackslash}p{(\columnwidth - 4\tabcolsep) * \real{0.3333}}
  >{\raggedright\arraybackslash}p{(\columnwidth - 4\tabcolsep) * \real{0.3333}}
  >{\raggedright\arraybackslash}p{(\columnwidth - 4\tabcolsep) * \real{0.3333}}@{}}
\toprule\noalign{}
\begin{minipage}[b]{\linewidth}\raggedright
Original class
\end{minipage} & \begin{minipage}[b]{\linewidth}\raggedright
Complete checked range needed here
\end{minipage} & \begin{minipage}[b]{\linewidth}\raggedright
Locator
\end{minipage} \\
\midrule\noalign{}
\endhead
\bottomrule\noalign{}
\endlastfoot
F, r=1 & supports18..94 & DO086 FR1--10 \\
F, r=2 & through110 & R068 FA1--9 \\
F, r=3 & through154 & R068 FA1--9 \\
F, 4\textless=r\textless=33 & through168 & R068 FA1--9 \\
IC8, r=1 or2 & supports23..28 & R069 EM1--21 \\
IC8, 3\textless=r\textless=31 & complete R075 table; its smallest ceiling is28 & R075 NR1--16 with its contiguous DO085/DO064 predecessor ranges \\
IC8, r=32 or33 & through168 & DO064 PM22--23 \\
\end{longtable}

R075 states and proves the entire interval beyond each prior ceiling, not just an endpoint value. Its rank3 ceiling is28; every larger compressed rank has a ceiling at least28. Thus the table covers every r=1..33 in each original structural class without a gap. R078 section2 already carries the same complete ranges; the present use checks their applicability at25. Consequently every support25 term remaining after those removals has A=I.

\subsection{3. Identity compression: fixed S or nonzero cross}\label{identity-compression-fixed-s-or-nonzero-cross}

Write the original block decomposition relative to S perpendicular W0, where W0 has signature(4,1), as h={[}{[}I,B{]},{[}C,D{]}{]}. The actual unitary identities R059 GB1--3 imply C\emph{J0 C=0 and B=-C}J0 D. Hence C=0 implies B=0 and h=I\_S direct-sum h0. If C is nonzero its image is a null line. These two alternatives exhaust the A=I case; no integral moving frame is introduced.

\subsubsection{3.1 The complete fixed-S sector}\label{the-complete-fixed-s-sector}

If h0 is not globally unipotent, DO090 NE1--6 removes the entire fixed-S class at supports23..26, including every global rank, characteristic degree and incidence profile. In particular it covers25.

If h0 is globally unipotent and nonidentity, the signature-one classification DO060 IB21--22 leaves exactly two possibilities: rank(h-I)=1 with (h-I)\^{}2=0, or rank(h-I)=2 with (h-I)\^{}3=0 and rank((h-I)\^{}2)=1. R080 AQ1--17 removes the entire first class at support25, including its integral-witness branch and all26 intrinsic double-pole cells. R079 PC1--25 removes the entire second class at support25, including both finite-rank ranges and every short/last original covector pivot. Its extra partial support26 conclusion is not used to enlarge this assertion. Rank zero is h=I and has support0, so cannot occur at25.

Thus the fixed-S alternatives are exhausted. The original type and label sums are inside the cited complete bounds; there is no new profile factor to multiply into the error list.

\subsubsection{3.2 The complete nonzero-cross reduction}\label{the-complete-nonzero-cross-reduction}

For C nonzero, DO099 JR1--24 removes both transverse cross directions at every support25..70 and removes the plane-C sector with any nonzero original short compression at25. At support25 its remaining matrices therefore have both cross lines in the actual packet plane and all four short compressions zero. DO099 JR24 explicitly verifies the support-independent DO093 IV19--22/R071 GR1--12 algebraic premises. It follows that, in the rational orthogonal decomposition

\[
 Z=S_{164}\perp u^\perp,\quad \dim_E Z=167,\qquad
 P=Eu\perp Ee_n,\qquad
 h=\begin{pmatrix}I_Z&B_Z\\C_Z&D_P\end{pmatrix},
 \quad (h^{-1})_{ZZ}=I_Z,
 \quad C_Z|_S\ne0,\quad\det h=1.                    \tag{J26.3}
\]

The precise DO095 RA1--10 classification now applies. Its linear-algebra proof uses J26.3 and no support count. It gives global rank3 if B\_Z C\_Z is nonzero and global rank2 if that product is zero; global ranks0 and1 are impossible with the stated nonzero S cross. At rank2, global unipotence and nonunipotence are separate exhaustive alternatives.

For the global-rank-three branch, DO074 TC1 forces the actual type-III count t to satisfy11\textless=t\textless=25. R082 TP1--13 removes every cell below its complete type-dependent thresholds. Every remaining term has k\textgreater K\_t\textgreater=125, where k is the number of nonzero original S covector coordinates. Thus k\textgreater=126. R084 IL1--21 removes the entire actual rank-three sector with k\textgreater=126, using the original and inverse null lines and charging both inverse containing sectors. Its proof does not require the previously cut remainder itself to be inversion invariant. Therefore no rank-three support25 term remains.

For the global-rank-two nonunipotent branch, DO101 PI1--17 removes the complete plane-C containing sector at supports25..32. Its hypotheses are weaker than J26.3 and so apply here. For the global-rank-two unipotent branch with at least one active type-I or type-III place, DO102 NT1--22 removes every support25 profile with u\textless=24. These two removals leave only the all-type-II globally rank-two-unipotent branch.

\subsection{4. Exhaustive denominator, eligibility and coordinate split in the last branch}\label{exhaustive-denominator-eligibility-and-coordinate-split-in-the-last-branch}

In that branch retain the exact original normalization C\_Z=(1,eta)\^{}t r, eta bar(eta)=U, the denominator ideal D\_eta, and

\[
 x=N_E D_\eta,\qquad 1\le k\le164,
 \qquad p=\#\{\text{eligible active places}\},\quad0\le p\le25.
                                                        \tag{J26.4}
\]

The cross is nonzero on S, which gives k\textgreater=1. The term p counts active eligibility, while q below counts eligibility among all168 selected places. The two are not interchanged. The following successive containing-set removals are exhaustive, with all boundary equalities assigned explicitly:

\begin{longtable}[]{@{}
  >{\raggedright\arraybackslash}p{(\columnwidth - 2\tabcolsep) * \real{0.5000}}
  >{\raggedright\arraybackslash}p{(\columnwidth - 2\tabcolsep) * \real{0.5000}}@{}}
\toprule\noalign{}
\begin{minipage}[b]{\linewidth}\raggedright
Condition on a still-retained all-II term
\end{minipage} & \begin{minipage}[b]{\linewidth}\raggedright
Complete containing-set bound
\end{minipage} \\
\midrule\noalign{}
\endhead
\bottomrule\noalign{}
\endlastfoot
x\textless=2\^{}120 & DO103 LD1--19 \\
x\textgreater=2\^{}1820 & DO104 CD1--32 \\
2\^{}120\textless x\textless2\^{}1820 and p\textless=24 & DO105 EL1--5; its theorem actually has no denominator premise in this case \\
2\^{}120\textless x\textless2\^{}1820, p=25 and1\textless=k\textless=21 & DO105 EL1--31 \\
Same open interval, p=25,22\textless=k\textless=164, and b\textgreater0 & DO107 IP1--23, first subset \\
Same open interval, p=25,22\textless=k\textless=164, b=0 and x\textgreater=2\^{}744 & DO107 IP1--23, second subset \\
2\^{}120\textless x\textless2\^{}744, p=25, b=0 and22\textless=k\textless=99 & DO107 IP1--23, third subset \\
2\^{}120\textless x\textless2\^{}744, p=25, b=0 and100\textless=k\textless=164 & K102 PC1--27, with q=25 supplied by DO107 IP3/8 \\
\end{longtable}

Here b counts inactive paired cross-zero places in the original lattices, exactly as in DO107. Outside its first three subsets, DO107 IP3 proves total eligibility q=25; its IP8 argument excludes q26 below2\^{}744. Hence every premise of K102 PC2 is present in the last row. DO105's unique eligible eta and active-label assertion, and DO107's inactive pointwise coefficients, are exactly the inputs retained in K102. No complete-central-orbit constancy or Fourier alias bound is substituted for them.

The first two rows cover both exterior denominator ranges including their endpoints. In their complement, p is either at most24 or exactly25. In the latter case k is either at most21 or at least22. Then b is either positive or zero. With b zero, x either lies at least at2\^{}744 or below it, and the remaining integer k is at most99 or at least100. These are exhaustive alternatives for every actual term, not a sampling of denominator values. Thus the last all-II branch is exhausted.

K100, K101 and DO108 supplied useful intervening containing-set removals. Their already charged errors remain in full. The larger bounds just used need not invoke their additional hypotheses: DO107 and K102 cover complete containing sectors, so applying them to their intersections with the later actual remainder is valid. No older charge is cancelled because its set is subsequently included in another bound.

\subsection{5. Complete chronological accounting and proof of the new floor}\label{complete-chronological-accounting-and-proof-of-the-new-floor}

The predecessor for this assembly is checks/k102\_joint\_remainder\_error.json, with102 all-top and31 uniform entries and the conservative preassembly floors25/11. Relative to the83-entry R078 list, the entire chronological suffix is:

\begin{longtable}[]{@{}lr@{}}
\toprule\noalign{}
Finding & Retained exponent \\
\midrule\noalign{}
\endhead
\bottomrule\noalign{}
\endlastfoot
DO099 & 12079 \\
R079 & 1120 \\
DO100 & 29550 \\
DO101 & 127443 \\
R080 & 537 \\
DO102 & 881 \\
R081 & 1395 \\
R082 & 2353 \\
DO103 & 1878 \\
K097 & 11658 \\
R083 & 5400 \\
DO104 & 2518 \\
R084 & 5398 \\
DO105 & 451 \\
K100 & 421 \\
K101 & 254 \\
DO107 & 392 \\
DO108 & 55281 \\
K102 & 3550 \\
\end{longtable}

The full arrays, not merely this suffix, are verified by exact integer shifts. There is exactly one all-top exponent231 and every other all-top exponent is at least254. Consequently

\[
 \sum_{e\in\mathcal L_{\rm top}}2^{-e}
 \le2^{-231}+101\,2^{-254}<2^{-230}.                 \tag{J26.5}
\]

The full31-entry uniform sum remains below2\^{}-230 as well. This assembly appends no entry and changes no earlier exponent. The certificate verifies the entire R078 prefix, each of the19 successive canonical error files, all retained charges, inclusive rank ceilings, the finite coordinate/type/denominator-boundary labels used above, and both exact sums. These are arithmetic and coverage diagnostics. They do not enumerate arithmetic arrows or replace the geometric partition and actual-source arguments in sections2--4.

Apply each chronological removal only to its intersection with the remainder after earlier removals. The actual sets thereby removed are disjoint even when the proved containing sets overlap. Each has absolute sum bounded by its existing complete charge. R078 already removed every support below25. Sections2--4 cover every original support25 alternative. Therefore every term left has support at least26. The original energy identity and the triangle inequality, with J26.5 and the retained earlier charges, prove J26.1.

This strictly reduces the finite/real sufficient Q1/Q2 gate from R078 J25.2 to J26.2. It does not prove J26.2, produce the prescribed Bessel/Mram value, construct a fully sized admissible mixed-CM family, evaluate the selected arithmetic-image content, identify the actual Read/STiesR or specific032 pairing, or compute the actual DP44 cost difference. Neither a shorter valid chain nor a universal strict negative DP44 difference is proved. Q3 has no Kummer-Hecke RH conclusion. Q4 keeps its separate global arithmetic oscillator-comparison gate.

\subsection{6. Algebraic-geometric placement and exact carried locators}\label{algebraic-geometric-placement-and-exact-carried-locators}

This is an exhaustion of the original determinantal, fixed-submodule and isotropic-cross strata of Hermitian correspondences over Spec(O\_E), with the original paired incidence schemes and chosen complex realizations retained. The fractional denominator and relative-norm ideals remain actual invertible modules; no principality, globally integral moving frame, or inverse on the arithmetic theta image is inferred. The polarization, theta line, Picard data and biextensions are unchanged.

The exact contextual module and determinantal anchors are SGA3, ExposeI,4.7.3/4.7.3.1, and SGA6, ExposeVI, Theorem3.3. The Picard anchors are FGA, Bourbaki232, Theorem3.1, pp148--149, and Bourbaki236, Corollaries1.7/2.3, pp229/232. Biextensions retain SGA7 I, ExposeVII,2.7/2.9, pp169--174, and ExposeVIII,7.4(b), pp310--311. Formal groups, Dieudonne crystals and Cartier duality retain cote4 scan003, §§8.1.1--8.1.2, scans005--006, §§8.4.1--8.4.3, and cote54 scans005--007. These are inherited contextual placements, not newly proved comparison theorems.

No new lambda-ring/K-theory equality, l-adic sheaf or weight statement, CM-motive comparison, or motivic L-function theorem follows. DP\_MULT Q3e, DP44/DP52/DP54--65 and THETA3 G1/A7 gain precisely J26.1--2 as a smaller finite sufficient gate. CX\_GRH CX-ALG-024 and GCT\_PLETH T-052 retain their separate arithmetic gates. The known source theorems are used at their exact displayed locators, without rederiving them. No Lean build or archive audit is performed.

\clearpage

\section{AM-R-088 --- only the identity G pattern and four double-pole sizes remain at rank one}\label{am-r-088-only-the-identity-g-pattern-and-four-double-pole-sizes-remain-at-rank-one}

\textbf{Canonical checked integration --- 2026-11-23T03:33:58.295183+00:00.} Complete root and independent proof/script/interface reviews and separate exact replays passed. The actual canonical predecessor has102/31 entries; this integration gives103/31, complete floors26/11. Frozen author history is preserved. AM-R-088 GZ1--11 removes the actual fixed-S rank-one-unipotent support26 integral-witness branch, every reduced double-pole cell d\textless=22, and every reduced pattern with at least one nonidentity unmixed G reduction. The exact averaged-content and original affine-coordinate count retains every paired type/label factor and the actual selected trace. The full334-place nonempty weighted G-pattern sum is below2\^{}-5533, replacing only the former factor2. All27 complete endpoints and the witness give one523 charge. Only the reduced all334Gidentity pattern with d23..26 remains in this branch. Actual R087102/31 becomes103/31, every old overlap retained, both whole sums below2\^{}-230, and complete floors26/11 unchanged. No full support26 exhaustion, signed positivity, actual source, selected arithmetic content, prescribed period, specific032/Read, DP44 or RH conclusion follows.

Author: GPT-6 Astra Pro; reasoning effort not externally exposed; ALP\_MACHINE /root. 23 November2026 UTC. Complete new proof; saved whole-proof/script reviews and exact replays are separate events before promotion.

\subsection{1. Actual sector and theorem}\label{actual-sector-and-theorem-7}

Keep the unchanged J56 all-top kernel, charged H=133 cutoff, f=a=168, original field E/F, selected real representation and metric, packet point, both CM partners, every paired source/target pencil and label, finite projectors, central phase, class masses and Haar measures. Consider the original fixed-S globally rank-one-unipotent sector at support s=26:

\[
 h=I_S\oplus(I+N),\qquad \operatorname{rank}_E N=1,
 \qquad N^2=0.                                      \tag{GZ1}
\]

Let d be the intrinsic number of actual double poles of eta=N\_nn.~K077 AP6 gives0\textless=d\textless=26 and every such place is typeII. At each of the334 unmixed relative G places the actual integral reduction of N has rank0 or1. The paired rank-zero ideal J includes exactly the rank-zero places. Write J\_max for the ideal containing all334 of them; N\_E J\_max=Q\^{}334, Q=g0g1.

\textbf{Theorem.} The complete actual union of the following support26 subsets has absolute normalized contribution less than2\^{}-523:

\begin{enumerate}
\def\labelenumi{\arabic{enumi}.}
\tightlist
\item
  the integral-witness branch of DO091 TD2--10;
\item
  the reduced branch with d\textless=22;
\item
  the reduced branch with at least one nonidentity unmixed G reduction, equivalently J!=J\_max.
\end{enumerate}

Consequently every still-retained fixed-S rank-one-unipotent support26 term has the reduced original form below, all334 unmixed reductions equal to the identity, and23\textless=d\textless=26. All earlier charges remain; this result introduces one new523 charge and keeps the complete floors26/11. It does not exhaust rank-one support26 or prove the signed sufficient inequality.

\subsection{2. The complete original count at variable support}\label{the-complete-original-count-at-variable-support}

The support-independent source interfaces are DO091 TD2--9, DO092 SN1--6/16--21, K077 AP6--11/18, LIFT078 AV1--9 and the primitive-coordinate argument in its section7, and R080 AQ1--15. The earlier numerical conclusion at support25 is not extended without recomputing the whole support26 expression.

DO091 splits every original column into an integral-witness branch or the reduced form

\[
 N=x\bar x^{\mathsf t}J_0/\eta,\quad
 x=(v q,r,\eta),\quad v=(2,95192390527,383092671043),
\] \[
 Vq\bar q+r\bar r=\eta\bar\eta,\quad q\ne0,
 \quad\bar\eta=-\eta\ne0,\quad |\eta_\tau|\le R=8844/133,
 \quad V=155821585821105090045582.                    \tag{GZ2}
\]

The two short integral pairings and the q=0 third-pairing witness use separate membership in the actual source and target lattices at every T place. Their algebra has no support bound. Their complete five-column/source count has only the displayed T\^{}(10s) and full pair-weight support dependence. At26 the unchanged TD10 expression gives

\[
 E_{\rm wit}(26)=1+168+107+5757+46af
 +(10L_T+997)26-a167L_Q
 =-195859503/512<-382000,                            \tag{GZ3}
\]

where L\_T=1019333/1024 and L\_Q=287843/4096 have their accepted strict directions. This bounds the entire witness branch, including all types, zero targets and possible witnesses, with no G-pattern enumeration added.

Fix s,d, an original support and intrinsic marked subset, all actual ordered pairs and the paired G pattern. Set M=Q\^{}167 and

\[
 \mathfrak l=J I_\Sigma I_D=\Lambda_F\mathcal O_E,
 \quad z=(N_EJ)^{1/(2f)}T^{-(s+d)/f},
 \quad y=R/z,\quad y_c=R/(zT^{2d/f}).                \tag{GZ4}
\]

The full entry ideal and content b=(q,r,eta) obey the same actual containments as R080 AQ3--5. The primitive first-three-coordinate Bezout identity puts both q and r in their original affine difference cosets of c=b intersect JI\_Sigma. In the true eta-weighted metric their uniform counts are

\[
 \#r\le(1+2\sqrt{y_c})^{2f},\qquad
 \#q\le(1+2\sqrt{y_c/V})^{2f}.                       \tag{GZ5}
\]

LIFT078 AV4 applies to every real fractional ideal, without a support restriction, and bounds the entire joint eta/content choice by2\^{}1065 y\^{}(1001f/1000). The actual selected trace is the support-independent DO092 SN19 bound2·2\^{}(-931f/2-630)z\^{}-f.~Both compact factors and the actual near-point transport remain. The Green coefficient is2(N\_EJ)\^{}(a/2)M\^{}-a. A full original column determines N; no isometry fibre or extra norm-generator multiplicity is added.

All these statements hold uniformly for each original labelled pair before its weight is summed. K077 AP18 gives the complete support/marked/type/pair factor binom(f,s)binom(s,d)(T+1)\^{}(s-d), including both CM components as paired data. Thus every input of R080 AQ13--15 applies with s variable; its positive extension to the cap is valid even when that artificial cap is not a nonempty actual stratum.

\subsection{3. Excluding the full rank-zero pattern has an exact large gain}\label{excluding-the-full-rank-zero-pattern-has-an-exact-large-gain}

Before G-pattern summation the complete z-dependent factor in the reduced count, up to a positive constant independent of J, is exactly

\[
 F_{s,d}(z)=z^\kappa(\sqrt z+A)^{2f}
                 (\sqrt z+A/\sqrt V)^{2f},
 \quad A=2\sqrt R\,T^{-d/f},
 \quad \kappa=(a-4-1/1000)f=3443979/125>163f.         \tag{GZ6}
\]

This is R080 AQ13 with every averaged-content, trace and both coordinate factors retained. For0\textless z'\textless=z, F(z')/F(z)\textless=(z'/z)\^{}kappa. Let z\_sd=(M/T\^{}(s+d))\^{}(1/f). Omitting one paired rank-zero place of F-norm q\_v decreases z by q\_v\^{}(1/f). Since the actual global matrix has rank one, its integral reduction has only ranks0 and1, and there is exactly one binary choice per unmixed relative place. Its other CM component is the inverse-adjoint partner, not a second choice. Therefore

\[
 \sum_{J\ne J_{\max}}F(z_J)
 \le F(z_{sd})\left\{\prod_{v=1}^{334}
                 (1+q_v^{-\kappa/f})-1\right\}.     \tag{GZ7}
\]

Every q\_v is g0 or g1 and exceeds2\^{}34. Thus q\_v\^{}(-kappa/f)\textless2\^{}-5542. The exact integer comparison

\[
 (2^{5542}+1)^{334}-2^{5542\cdot334}
                       <2^{5542\cdot334-5533}       \tag{GZ8}
\]

gives a full nonempty-pattern sum below2\^{}-5533. This is an explicit sum over all nonempty missing subsets; it does not discard their multiplicity or insert only one favorable prime.

The complete R080-style cell bound below used the unrestricted pattern bound less than2. Replacing that single factor2 by2\^{}-5533 consequently gives the relative improvement

\[
 U_{s,d}^{\rm nonfull}<2^{-5534}U_{s,d}.             \tag{GZ9}
\]

The finite Euler correction and near-point transport each still contribute their own factor2. They are not part of the replaced pattern factor.

\subsection{4. Every support26 endpoint with all original factors}\label{every-support26-endpoint-with-all-original-factors}

Put y\_sd=R/z\_sd and y'\_sd=R(T\^{}(s-d)/M)\^{}(1/f). The complete unrestricted reduced d-cell is bounded by

\[
\begin{split}
 U_{s,d}={}&2^{3+1065-931f/2-630}
 \binom fs\binom sd(T+1)^{s-d}z_{sd}^{-f}
 y_{sd}^{21021/125}\\
 &\times(1+2\sqrt{y'_{sd}})^{2f}
       (1+2\sqrt{y'_{sd}/V})^{2f}.                  \tag{GZ10}
\end{split}
\]

Here2\^{}3 consists of one finite Euler correction, one complete G-pattern bound and one genuine near-point product. The2\^{}1065 constant counts all eta and content ideals; both remaining factors count actual coordinate elements in their affine cosets. No type, norm, class, generator, phase or Haar factor is suppressed.

At s=26 the exact f-th-root and square-root brackets, followed by directed rational logarithms, give the following strict margins. A margin m means log2 U\_(26,d)\textless-m; negative margins are failed unrestricted sufficient bounds, not contributions or existence claims.

\begin{longtable}[]{@{}rr@{}}
\toprule\noalign{}
d & Margin m \\
\midrule\noalign{}
\endhead
\bottomrule\noalign{}
\endlastfoot
0 & 4402 \\
1 & 5376 \\
2 & 6254 \\
3 & 6779 \\
4 & 6918 \\
5 & 6936 \\
6 & 6936 \\
7 & 6934 \\
8 & 6932 \\
9 & 6930 \\
10 & 6928 \\
11 & 6927 \\
12 & 6925 \\
13 & 6922 \\
14 & 6904 \\
15 & 6792 \\
16 & 6326 \\
17 & 5473 \\
18 & 4500 \\
19 & 3507 \\
20 & 2513 \\
21 & 1518 \\
22 & 524 \\
23 & −470 \\
24 & −1464 \\
25 & −2456 \\
26 & −3448 \\
\end{longtable}

All27 complete cells are evaluated. Among d\textless=22 the worst margin is524 and all other22 margins are at least1518. At d23..26, GZ9 gives nonfull-pattern margins5064,4070,3078,2086, respectively. These four are all strictly positive and apply to every nonidentity-G pattern together.

It suffices to count all patterns at d\textless=22 and only nonfull patterns at d\textgreater=23; this includes the theorem's entire union without a duplicated d\textless=22 branch. Together with the whole witness branch, the complete bound is

\[
 2^{-524}+22\,2^{-1518}+2^{-5064}+2^{-4070}
       +2^{-3078}+2^{-2086}+2^{-382000}<2^{-523}.     \tag{GZ11}
\]

This proves the theorem. No extra support or marked-set union is needed after GZ10, because those factors were included before evaluating each cell. The intrinsic d partition is disjoint. The only uncovered reduced possibilities have J=J\_max and23\textless=d\textless=26; they have not been declared nonexistent or small.

\subsection{5. Chronology, verification and geometric meaning}\label{chronology-verification-and-geometric-meaning}

Append523 once to the actual complete all-top predecessor, retaining every older charge and all31 uniform entries. The default predecessor is R087's102/31 arrays with complete floors26/11; unchanged-script replay against a later complete prefix is permitted and recorded separately. Both full dyadic sums remain below2\^{}-230. Only intersection with the then-current remainder is removed. The full floor remains26 and the signed J26.2 inequality stays open.

The exact certificate checks all27 directed endpoint intervals, the full nonempty G-pattern sum GZ8, the complete claimed union, the actual witness endpoint and the entire chronological arrays. Its elementary rational log and root methods are those already checked in R080. Finite diagnostics are not arithmetic-arrow enumeration and do not replace the whole applicability argument in sections2--3.

The G condition is reduction to the identity on the original integral Hermitian modules at every unmixed finite point. The content and affine-coordinate conditions remain in the same fractional invertible modules over Spec(O\_E)/Spec(O\_F), with both CM partners, polarization, theta line, Picard data and biextension retained. The exact carried framework anchors are SGA3 I4.7.3/4.7.3.1; SGA6 VI Theorem3.3; FGA232 Theorem3.1 pp148--149; FGA236 Corollaries1.7/2.3 pp229/232; SGA7 I VII2.7/2.9 pp169--174 and VIII7.4(b) pp310--311; cote4 scan003 §§8.1.1--8.1.2 and scans005--006 §§8.4.1--8.4.3; cote54 scans005--007. These are inherited contexts, not new lambda-ring/K-theory, l-adic weight or CM-motive comparisons.

DP\_MULT Q3e and THETA3 G1/A7 gain this narrower original finite sufficient gate. The actual fully sized source, selected arithmetic image/content, prescribed period, specific032/Read and actual DP44 costs remain open. CX\_GRH CX-ALG-024 and the global GCT\_PLETH T-052 comparison keep their separate gates. No RH conclusion, Lean build or archive audit is asserted.

\subsection{Canonical algebraic-geometric context}\label{canonical-algebraic-geometric-context-17}

The reductions are the actual integral Hermitian modules at unmixed finite points; content and coordinate ideals remain fractional invertible modules over Spec(O\_E) and Spec(O\_F), with the same paired CM action, polarization, theta line, Picard scheme and biextension. Carried exact anchors: SGA3 I4.7.3/4.7.3.1; SGA6 VI Theorem3.3; FGA232 Theorem3.1 pp148--149; FGA236 Corollaries1.7/2.3 pp229/232; SGA7 I VII2.7/2.9 pp169--174 and VIII7.4(b) pp310--311; cote4 scan003 sections8.1.1--8.1.2 and scans005--006 sections8.4.1--8.4.3; cote54 scans005--007. No new lambda-ring/K-theory, l-adic weight or CM-motive comparison is asserted. DP\_MULT Q3e, THETA3 G1/A7, CX\_GRH CX-ALG-024 and GCT\_PLETH T-052 retain their stated global gates.

\clearpage

\section{AM-LIFT-086 --- actual good models and intrinsic line costs along the CM isogeny construction}\label{am-lift-086-actual-good-models-and-intrinsic-line-costs-along-the-cm-isogeny-construction}

\textbf{Canonical checked integration --- 2026-11-23T03:40:47.957010+00:00.} Complete saved-proof and executable reviews and separate exact diagnostic replays passed. No numerical error charge is introduced. Frozen author history is preserved. AM-LIFT-086 IM1--27 constructs actual everywhere-good polarized models for the LIFT085 geometric CM seed and target over K=K0(A0{[}15N\^{}2{]}), with {[}K:K0{]}\textless=(15N\^{}2)\^{}(4g\^{}2). The full integral OE action, principal polarization, canonical trace-one cubical line, isogeny and adjoint extend over the same ring. Actual intrinsic tangent-line degrees change by an amount in{[}-log N,log N{]}, the full determinant degree by at most g log N, and the theta degree by at most g log N/2 under KeyCore. Target intrinsic lines have length constant1. If the actual seed lines have degree at least-b and all genuine rule/integral/cusp/source premises hold, the sufficient Step5 cost is C\_N=(D+gn)/2+D log r+(b+log N)min(D+gn,(3D+gn)/2). No effective bound on N, initial field or arbitrary-Host b is supplied. Host/Sizes, actual rule nonzero/content, prescribed period, specific032/Read, evaluated DP44 and RH remain open. No new finite error charge; actual103/31 arrays and complete floors26/11 remain unchanged.

Author: /root/lift\_accounting, 23 November 2026 UTC. \textbf{Author-frozen theorem in the controlled pending area; the hash-specific whole-read, exact replay and independent reviews have separate records.} No canonical source or ledger is changed.

This theorem supplies the good polarized model and actual metrized tangent-line transport for the CM point constructed in LIFT085. It gives an explicit coefficient for log N, where N is that construction's integer isogeny scale. It does not bound N, the initial field, or the seed line degrees for an arbitrary new Host. It is not an effective admissible-family or DP44 conclusion.

\subsection{1. Exact objects, field and conclusions}\label{exact-objects-field-and-conclusions}

Retain the actual Host E/F, Datum, A=1 finite level, Coeff, Moduli representation and genuine geometric CM seed of LIFT085. Let the seed be \[
 (A_0,\lambda_0,\iota_0,L_0,\eta_0)/K_0,\qquad
 \dim A_0=g\ge1.                                            \tag{IM1}
\] The field K0 contains the source coefficient field and all required algebraic E-labels. The polarization is principal, the action is the full O\_E action with Rosati conjugation, and L0 is the actual origin-rigidified symmetric cubical trace-one line. Geometric CM means that A0 has sufficiently many complex multiplications after algebraic closure. The mere action of E on a variety of dimension larger than {[}E:Q{]}/2 would not establish this premise. It holds for the actual CM products in LIFT084 and R086. The isogeny target is therefore also geometrically CM.

Set \[
 K_{\rm good}=K_0(A_0[15]).                                  \tag{IM2}
\] Run LIFT085 over this field, enlarging its protected rational-prime set by3,5 if desired. Retain its actual E-linear isogeny and prescribed finite level: \[
 u:A_0\longrightarrow A_1,\quad
 u^\dagger u=[N^2],\quad uu^\dagger=[N^2],\quad
 u^\vee\lambda_1u=N^2\lambda_0,\quad
 u\eta_0=\eta_1,\qquad N\ge1.                                \tag{IM3}
\] Here \(u^\dagger=\lambda_0^{-1}u^\vee\lambda_1\) is an actual homomorphism because both polarizations are principal. These equalities concern the actual algebraic polarizations.

All field-valued target data and its integral polarized model exist over \[
 K=K_0(A_0[15N^2]),\qquad
 [K:K_0]\le|\operatorname{GL}_{2g}(\mathbf Z/(15N^2))|
             \le(15N^2)^{4g^2}.                             \tag{IM4}
\] The target field of LIFT085 over Kgood is contained in K; equality is unnecessary. If a good seed model is already supplied over K0, the factor 15 can be omitted and the bound becomes N\^{}(8g²). All normalized line degrees below are invariant under further finite base extension of these actual models. In particular, the seed degree bound may be established over its initial field and then used over Kgood and K.

\textbf{Model conclusion.} Both Ai extend to abelian schemes \(\mathcal A_i/\mathcal O_K\). Their full integer-ring actions, principal polarizations, canonical trace-one rigidified cubical lines \(\mathcal L_i\), and u,u-dagger extend over the same ring. They retain \[
 \mathcal u^*\mathcal L_1\simeq\mathcal L_0^{\otimes N^2}
       \quad\text{as origin-rigidified cubical lines}.        \tag{IM5}
\] No additional extension is needed for the infinitely quantified action, adjoint or primitive-line identities. This constructs the integral PAT/PolAV geometry used by the source; it does not assert an integral Obj over a base where the source requires every positive integer to be a unit.

Choose an algebraic seed eigenbasis v1,\ldots,vg, with its actual labels and multiplicities, and put \(w_j=du(v_j)\). Define the intrinsic ideals \[
 I_{0,j}=\{c\in K:cv_j\in\operatorname{Lie}\mathcal A_0\},
 \quad I_{1,j}=\{c\in K:cw_j\in\operatorname{Lie}\mathcal A_1\}.
                                                               \tag{IM6}
\] Equip these lines with the literal source norms \(n_{i,j,\sigma}^2=\pi H_{i,\sigma}(v_{i,j},v_{i,j})\), where v0,j=vj and v1,j=wj. Write di,j for their normalized arithmetic degrees. Then \[
 I_{0,j}\subset I_{1,j}\subset N^{-2}I_{0,j},\qquad
 n_{1,j,\sigma}=Nn_{0,j,\sigma},                              \tag{IM7}
\] \[
 d_{1,j}-d_{0,j}
 =\frac{\log|I_{1,j}/I_{0,j}|}{[K:\mathbf Q]}-\log N
       \in[-\log N,\log N].                                 \tag{IM8}
\] The quotient is a finite additive module, also for fractional ideals. The ideals have identical localizations away from N.

The target theta and determinant lines in KeyLocAt/KeyArchAt are constructed from this actual model. Under the retained source KeyCore degree identity, they satisfy \[
 |\operatorname{ndeg}\Delta_1-\operatorname{ndeg}\Delta_0|
      \le g\log N,\qquad
 |\operatorname{ndeg}T_1-\operatorname{ndeg}T_0|
      \le\tfrac g2\log N.                                   \tag{IM9}
\] This does not compare particular theta sections, arithmetic lift images or divided theta operators.

\textbf{Step5 consequence.} Suppose \(d_{0,j}\ge-b\), b\textgreater=0. Keep the actual rule, every integral-evaluation and cusp-bound premise of LIFT082, KeyCore, LemmaHStatement and all global source compatibility premises. Put \(D=\sum_{\tau,i}a_{\tau i}\), V=gn, and let r\textgreater=1 bound the actual block multiplicities. At the target point selected by LIFT085 for a nonzero rule, the nonzero frame polynomial has the sufficient bound \[
 \boxed{\quad
 \frac1{[K:\mathbf Q]}\sum_\sigma\log S_\sigma\ge-C_N,\qquad
 C_N=\tfrac12(D+V)+D\log r+
 (b+\log N)\min\{D+V,(3D+V)/2\}.
 \quad}                                                       \tag{IM10}
\] This uses target intrinsic lines and length constant1. N remains unevaluated, and b remains a premise unless supplied by a theorem for those actual seed lines.

\subsection{2. A bounded field giving good reduction everywhere}\label{a-bounded-field-giving-good-reduction-everywhere}

Fix a finite place v of K0, of residue characteristic p.~Geometric CM gives potential good reduction (Milne,2020, Proposition7.12, p.55). Thus inertia has finite image on \(T_\ell A_0\) for \(\ell\ne p\), by the Néron criterion and its converse (ibid.,6.11--6.12 and the paragraph after6.12, p.51). This does not require every geometric CM endomorphism to be defined over K0.

Choose ell=3 if p!=3, and ell=5 if p=3. Over Kgood every ell-torsion point is rational. Hence the finite inertia image lies in \[
                 1+\ell M_{2g}(\mathbf Z_\ell).              \tag{IM11}
\] This group has no nontrivial torsion for odd ell. Indeed, if \(X=1+\ell^aB\ne1\), a\textgreater=1 and B is not zero modulo ell, binomial expansion gives \[
 v_\ell(X^\ell-1)=a+1,\qquad
 v_\ell(X^m-1)=a\quad(\ell\nmid m),                           \tag{IM12}
\] where matrix valuation is the minimum entry valuation. The leading terms are \(\ell^{a+1}B\) and \(m\ell^aB\). Every other term has strictly larger valuation, including the last term for ell=3,a=1. Iterating proves \(X^n\ne1\) for every positive n.~Thus the finite inertia image is trivial, and the Néron criterion gives good reduction at every place of Kgood.

The local abelian schemes glue over its ring of integers. Start with a spread outside finitely many primes and use uniqueness of each good model and extension of generic homomorphisms at the missing DVRs (Milne6.4/6.6, p.49). Equivalently this is the everywhere proper Néron model over the Dedekind base. No good model over the smaller K0 is claimed. Good reduction survives arbitrary finite base change, so the seed has its actual model over O\_K.

Since A1 is isogenous to A0, it too has everywhere good reduction, by Milne6.14, pp51--52. That result follows from the injection of its Tate module into an unramified Tate module. The same gluing gives \(\mathcal A_1/\mathcal O_K\).

Finally \(A_0[m](\overline K_0)\simeq(\mathbf Z/m)^{2g}\) for m=15N². The torsion field is the fixed field of the kernel of this finite Galois action, so its group injects into GL\_(2g)(Z/m). Counting all square matrices gives IM4. The model extension has no unspecified factor beyond that field. K need not be CM, Galois over Q, unramified at protected primes, or equal to a previously fixed coefficient field.

\subsection{3. Full integral structure and the actual trace-one theta line}\label{full-integral-structure-and-the-actual-trace-one-theta-line}

Every generic homomorphism between two abelian schemes over a DVR extends uniquely (Milne6.4/6.6), and these extensions glue over a Dedekind base. Apply this to u,u-dagger, every action map iota\_i(a), a in O\_E, and the principal polarization and its inverse. The extended polarization is a polarization (Milne6.7), and its extended inverse makes it principal. Uniqueness preserves addition, multiplication, the unit, the adjoint relation and IM3. The infinitely quantified action entails no infinite field extension: each extension is unique over the same fixed ring, and each ring relation holds on its generic fibre.

The extended u is finite locally free of degree N\^{}(2g). Its kernel lies in \(\mathcal A_0[N^2]\). On every geometric fibre the two composites in IM3 make u an isogeny, with finite kernel and surjective image. The proper map u is therefore finite. Flatness follows from the fibre criterion, since both abelian schemes are smooth over the Dedekind base and fibre isogenies are flat. Its kernel is exactly the flat closure of the generic kernel. It realizes the fppf quotient by that closure, agreeing with the finite locally free quotient and base-change theorem (van der Geer \& Moonen,2011,4.16(iii),4.29--4.30). No reduced or étale replacement of the kernel is made at primes dividing N.

Let w0 be the literal Host element, \(w_0+\bar w_0=1\). Let \(\mathscr P_i\) be the normalized Poincaré line on \(\mathcal A_i\times\mathcal A_i^\vee\). Define the line on the square and then its pullback along a graph: \[
 \mathscr B_i=(1\times\lambda_i)^*\mathscr P_i
       \quad\text{on }\mathcal A_i\times\mathcal A_i,
 \qquad
 \mathcal L_i=(\mathrm{id},\iota_i(w_0))^*\mathscr B_i .
                                                               \tag{IM13}
\] Its polarization is \[
 \lambda_i\iota_i(w_0)+\iota_i(w_0)^\vee\lambda_i
 =\lambda_i\{\iota_i(w_0)+\iota_i(\bar w_0)\}=\lambda_i.
                                                               \tag{IM14}
\] Poincaré biextension supplies the origin rigidification, cubical structure and symmetry. Its adjoint identity and the trace-one relation give the exact SAB.adj and SAB.prim clauses, as in the universal argument of LIFT084/085. These constructions occur over O\_K, commute with base change, and restrict to the actual generic source lines. No symmetric two-torsion twist is inserted.

Functoriality of the normalized Poincaré biextension and E-linearity give \[
 u^*\mathcal L_{\lambda_1}
 \simeq\mathcal L_{u^\vee\lambda_1u}
 =\mathcal L_{N^2\lambda_0}
 \simeq\mathcal L_{\lambda_0}^{\otimes N^2}.                  \tag{IM15}
\] These are the actual rigidifications and cube maps. The associated normalized cubical Hermitian metrics have \(u^*H_1=N^2H_0\). Thus no independent complex scalar or CM period is removed in the norm comparison.

The generic level sections extend over each DVR by properness, and u eta0=eta1 extends. We retain their field-valued prescribed level. We do not assert an isomorphism from a constant étale level scheme to integral torsion at a residue characteristic dividing its order. For p not dividing N the actual maps do give \[
 \mathcal A_0[p^\infty]\simeq\mathcal A_1[p^\infty],\qquad
 \mathcal u^{-1}=[N^2]^{-1}\mathcal u^\dagger
        \text{ on finite }p\text{-power torsion}.             \tag{IM16}
\] This is an equality of group-scheme maps after the displayed coefficient extension. It supplies no identity of selected theta derivatives or arithmetic lift images.

\subsection{4. Intrinsic tangent lattices and every embedding}\label{intrinsic-tangent-lattices-and-every-embedding}

The tangent module \(M_i=\operatorname{Lie}(\mathcal A_i/\mathcal O_K)\) is finite projective of rank g. Its intersection with any nonzero generic line Kv is a saturated rank-one module, represented by IM6. This construction applies to a chosen line even if its full eigenspace has larger multiplicity. A seed degree bound must still hold for that actual line: LIFT081 does not bound every line in a larger multiplicity space.

The derivative of u sends M0 into M1, proving \(I_{0,j}\subset I_{1,j}\). If cwj belongs to M1, applying du-dagger gives \(N^2cv_j\in M_0\), proving the second inclusion in IM7. These become equalities after localizing away from N.

They also have exact base change. Locally choose finitely many integral linear equations cutting out the rational line, clearing denominators in an annihilator basis. The saturated intersection with Mi is their simultaneous kernel. Flat extension commutes with that kernel. Number-field integer-ring extensions are flat, so this is precisely the base change of the actual intrinsic line, not an unrelated lattice. For the full simultaneous O\_E eigenkernel, use a finite integral basis of O\_E and the same argument. This retains every complete-DVR extension in LIFT082.

At every complex embedding sigma, algebraic base change of the polarization identity gives \[
 \pi H_{1,\sigma}(du(v_j),du(v_j))
     =N^2\pi H_{0,\sigma}(v_j,v_j).                          \tag{IM17}
\] This is valid with the source conjugate-first convention. Positive square roots give n1=N n0. The same global vj, u and algebraic labels are transported at every embedding; no field embedding is applied to a free transcendental period.

For fractional ideals I0 contained in I1, \[
        |I_1/I_0|=\mathrm N(I_0)/\mathrm N(I_1).              \tag{IM18}
\] Multiplication by a common integer reduces this to the integral-ideal index. IM7 gives \(1\le|I_1/I_0|\le N^{2[K:\mathbf Q]}\). The literal definition is \[
 \operatorname{ndeg}(I,n)=
 -\frac{\log\mathrm N(I)}{[K:\mathbf Q]}
 -\frac1{[K:\mathbf Q]}\sum_\sigma\log n_\sigma.              \tag{IM19}
\] Substitution proves IM8 exactly. All complex embeddings are counted individually. Rescaling a vector by c in K* changes its ideal and norms by opposite product-formula terms. Thus the normalized degree and IM8 are independent of that normalization.

These ideals/norms satisfy the actual IsVecLine clauses. If rows of a local model basis have matrix M in the generic basis, cvj is integral exactly when every entry of \(c e_jM^{-1}\) is integral. This is the localized intrinsic-ideal criterion. The archimedean clause is precisely pi H. No coordinate ideal cD\_tau is substituted.

\subsection{5. Actual theta and determinant Key lines}\label{actual-theta-and-determinant-key-lines}

Choose a nonzero generic theta section theta\_i and the bases above. Define \[
 T_i.I=\{a\in K:a\theta_i\in H^0(\mathcal A_i,\mathcal L_i)\},
 \quad
 \Delta_i.I=\{a\in K:a(v_{i,1}\wedge\cdots\wedge v_{i,g})
                                      \in\det M_i\}.        \tag{IM20}
\] The section module is invertible of rank one and commutes with every base change. On each geometric fibre, an ample line on an abelian variety has zero higher cohomology: for large m invertible in the characteristic, finite étale multiplication, the projection formula and normalized trace make H\^{}j(L) a direct summand of H\^{}j({[}m{]}\emph{L); symmetry gives {[}m{]}}L=L\^{}(m²), and Serre vanishing applies. Principality and Riemann--Roch give h0=1. Proper cohomology and base change now gives the assertion over O\_K, as in LIFT085's frame-torsor proof. A chosen generic theta\_i need not generate this module at every prime.

At each prime choose a local section generator and a local tangent basis. Their scalar and determinant relative to the generic choices are exactly the generators ca,cd in KeyIntegral/KeyLocAt. The local section is nonzero on every fibre. The global tangent module need not be free: only these local KeyIn records and the complete generic eigenbasis are used. No root of the Key scalar or further extension is adjoined.

At each embedding give these lines their exact KeyArchAt norms: \[
 T_i.n_\sigma^2=
 \frac{\int_{Fd}|\theta_i(z)|^2e^{-\pi H_i(z,z)}\,dz}
      {\operatorname{vol}(Fd)},\qquad
 \Delta_i.n_\sigma^2=
 \frac{\pi^g|\det W_i|^2}{\operatorname{vol}(Fd)}.            \tag{IM21}
\] The uniformization, function, fundamental domain and Wi represent the actual Key, section and basis. The principal identity det(Hi) vol(Fd)=1 makes the second expression the determinant of the source tangent Gram matrix. A scalar change in theta\_i changes the actual ideal and all norms compatibly, leaving normalized degree unchanged.

Wedges of du and du-dagger give \[
 \Delta_0.I\subset\Delta_1.I\subset N^{-2g}\Delta_0.I,\qquad
 \Delta_1.n_\sigma=N^g\Delta_0.n_\sigma.                     \tag{IM22}
\] Thus the normalized determinant degree changes by \(\log|\Delta_1.I/\Delta_0.I|/[K:\mathbf Q]-g\log N\), lying in the stated interval. This uses the complete determinant lattice, not a false equality between it and the product of the individual line intersections.

Retain the source KeyCore and its KeyDegCore consequence. Since the source states its formula for dual lines, reversal of both degree signs and division by {[}K:Q{]} gives \[
 \operatorname{ndeg}T_i=
 \tfrac12\operatorname{ndeg}\Delta_i-\tfrac g4\log(2\pi).
                                                               \tag{IM23}
\] Subtracting proves IM9. This is an application of the stipulated Key theorem, not a new unconditional arithmetic Riemann--Roch statement. In particular no identity between u*theta1 and theta0 is inferred: the former lies in H0(L0\^{}(N²)), generally a much larger space.

\subsection{6. All local polynomial tests and the exact sufficient cost}\label{all-local-polynomial-tests-and-the-exact-sufficient-cost}

Use the actual rule supplied to LIFT085, with the polynomial/equivariance weights and all arithmetic/cusp hypotheses of LIFT082. Its chosen base point is in the nonzero frame locus. The frame polynomial relative to the reference eigenbasis wj and any nonzero theta1 is therefore nonzero, even if the value at that particular reference frame is zero. A change of basis and scalar acts transitively on the frame fibre, and a nonzero value at another frame is an evaluation of the same polynomial. The source IsFrPoly and FrChg transpose convention remains exactly the one checked in LIFT082.

The model, full action, section module, eigenkernel base change and actual Key lines above supply the geometric inputs of LIFT082. Its integral-evaluation premise on the rule remains necessary: a good model does not make every rational rule integral. With that premise, LIFT082 gives all Step5Loc quantifiers, including every compatible complete-DVR extension and its invertible frame matrices. The genuine IsUnf and cusp\_bound\_allEmb premises still supply the bounds S\_sigma in Step5Arch.

Use the target intrinsic lines as comparison lines. Their vectors have their own norms, so L=1. Their product wedge is integral in det M1, giving the determinant-ideal inequality. The images form a complete independent eigenbasis because du is invertible and E-linear. IM8 gives \[
                         d_{1,j}\ge-b-\log N.               \tag{IM24}
\] The unchanged LIFT080 height theorem now gives IM10.

Retaining the old seed lines would instead give L=N and d\_-=-b, hence \[
 C_N^{\rm seed}=\tfrac12(D+V)+D\log r+
 b\min\{D+V,(3D+V)/2\}+\tfrac{3D+V}{2}\log N.               \tag{IM25}
\] The exact improvement in this sufficient upper bound is \[
 C_N^{\rm seed}-C_N=\tfrac12(D-V)_+\log N\ge0.               \tag{IM26}
\] This is a comparison of sufficient estimates, not a lower bound on a necessary cost. If the old fixed ell0=logOm must be retained, the additional LIFT080 MT9 term is \[
             \tfrac{3D+V}{2}\,(\ell_0/2+b+\log N)_+.          \tag{IM27}
\] It is added to the old L=1 equality-based expression. The unchanged lamw\_link is not invoked with a false common-degree premise.

The constants b=10log2 or16log2 may be inserted only for the actual LIFT081 component lines in their J14/J56 field/model/action scope. They are not assigned to R086's new general-Host blocks or a pure complementary factor without a separate line theorem.

\subsection{7. Constructed interfaces and remaining gates}\label{constructed-interfaces-and-remaining-gates}

\begin{longtable}[]{@{}
  >{\raggedright\arraybackslash}p{(\columnwidth - 4\tabcolsep) * \real{0.3333}}
  >{\raggedright\arraybackslash}p{(\columnwidth - 4\tabcolsep) * \real{0.3333}}
  >{\raggedright\arraybackslash}p{(\columnwidth - 4\tabcolsep) * \real{0.3333}}@{}}
\toprule\noalign{}
\begin{minipage}[b]{\linewidth}\raggedright
Interface
\end{minipage} & \begin{minipage}[b]{\linewidth}\raggedright
Supplied for the LIFT085 point
\end{minipage} & \begin{minipage}[b]{\linewidth}\raggedright
Remaining requirement
\end{minipage} \\
\midrule\noalign{}
\endhead
\bottomrule\noalign{}
\endlastfoot
Integral model & Everywhere good model after15-torsion; no additional target-model extension & Actual seed and initial field under the Host/Datum/Moduli hypotheses \\
Principal geometry & Full integral action, principal polarization, canonical line, all adjoint/primitive identities & No old pure-type LamObjGood/LamPtClauses instance is manufactured \\
Level/formal groups & Same generic finite level; prime-to-N torsion and p-divisible isomorphisms & No old selected theta/frame/divided derivative identification \\
Field cost & {[}K:K0{]} at most (15N²)\^{}(4g²) & N, initial field degree/discriminants and all changed-field terms \\
Letters & Actual target lines, exact degree interval, determinant inclusion, all flat base changes & Numerical lower bound for actual seed lines in a new Host \\
Key lines & Actual ideals/norms and local/archimedean predicates; IM9 under KeyCore & Stipulated Key/height theorems and actual source operator ties \\
Polynomial & Nonzero frame polynomial at the chosen base point & Actual rule nonzero, integral as required, and all cusp hypotheses \\
Main comparison & Explicit sufficient penalty in log N & No effective N, DP44 room, selected content, period, Read/STiesR/032 value or RH \\
\end{longtable}

The model and target-line constructions are consequently no longer independent hypotheses once the geometric CM seed and LIFT085 isogeny are supplied. The effective problem remains: N was chosen qualitatively from a nonempty open locus, and the actual arithmetic rule and every comparison term still require proof.

The relative Picard/Poincaré and cubical biextension live over the same arithmetic base. IM16 gives precisely the formal-group/Dieudonné comparison away from N; it is not a slope gain or equality of arbitrary differential operators. Carried contextual locators are FGA232 Theorem3.1 pp148--149; FGA236 Corollaries1.7/2.3 pp229/232; SGA7 I VII2.7/2.9 pp169--174 and VIII7.4(b) pp310--311; SGA3 I4.7.3/4.7.3.1; SGA6 VI Theorem3.3; cote4 scan003 sections8.1.1--8.1.2 and scans005--006 sections8.4.1--8.4.3; cote54 scans005--007. These are inherited context, not fresh readings or new motive/weight/RH implications. DP\_MULT Q3e/DP44/DP52/DP54--65, THETA3 G1/A7, CX\_GRH CX-ALG-024 and GCT\_PLETH T-052 retain the displayed gates.

\subsection{8. Source and primary verification boundary}\label{source-and-primary-verification-boundary-1}

Literal sources are pinned to commit7cbf7a0c32fca1badf05931879a37c3f8a2b03b4. Exact hashes and read ranges are in the companion manifest.

\begin{itemize}
\tightlist
\item
  TwoWeightEqOneAddRoot.lean: Host14--25, SAB72--83, SABMap86--93, Base124--129, TateIso131--141, Obj144--160. The integral construction supplies SAB/PAT, not Obj over a base violating Base.unit.
\item
  LamLetters.lean: IsModelAt87--94 and IsVecLine96--106, including the inverse-matrix ideal condition and pi H.
\item
  IntegralOf.lean: FibKey29--39, IsFrPoly97--101, IsMetricAt103--118, LieInt120--128, SecInt130--136 and integral\_of150--158. All complete-DVR/section conventions are those of LIFT082.
\item
  DetIdeal.lean: MLine10--14, Good20, deg24--25, ndeg27 and dual33. StatementB.lean: KeyIntegral12--16, KeyLocAt18--24, KeyArchAt26--32 and KeyDeg34--37. HadamardLoc0.lean: KeyCore12--29, KeyDegCore33--36, UnimodVol38--43, LieGramDet45--52 and HadamardLoc0 54--62. StatementS.lean: full Step5Loc38--53, Step5Arch55--61 and Step5 63--73.
\item
  Canonical LIFT080 MT1--17, LIFT081 ED1--23, LIFT082 FT1--29, LIFT084 MP1--25, LIFT085 HF1--26 and R086 ZJ1--21, in their checked scopes.
\end{itemize}

Primary passages read for this report:

\begin{enumerate}
\def\labelenumi{\arabic{enumi}.}
\tightlist
\item
  Milne, J. S. (2020). \emph{Complex multiplication} (Version0.10). \url{https://www.jmilne.org/math/CourseNotes/CM.pdf}. Proposition6.4 and Corollary6.6, p.49, with displayed proofs; Proposition6.7, p.49, statement and cited-proof scope; Theorem6.12 and Corollaries6.13--6.14, pp51--52, with displayed proofs; Proposition7.12, p.55, with displayed proof. IM12 is proved here. The scanned Serre--Tate article was located but not legibly read; no direct-read claim for its proof is made.
\item
  Conrad, B. (2015, October14). \emph{Néron models, Tamagawa factors, and Tate--Shafarevich groups} (notes taken by T. Feng). \url{https://math.stanford.edu/~conrad/BSDseminar/Notes/L3.pdf}. Theorem1.1, p.2; Example2.1, p.4; Remark3.8, p.7. These confirm the Dedekind/global and local model interfaces. Semistable reduction is not needed in the15-torsion proof.
\item
  van der Geer, G., \& Moonen, B. (2011, September15). \emph{Abelian varieties}, ChapterIV, Quotients by group schemes. \url{https://www.math.ru.nl/~bmoonen/BookAV/Quotients.pdf}. Theorem4.16(iii), pp55--56; Examples4.29/4.40 and base change4.30, pp62--63 and66. This known quotient theorem is invoked, not reproved.
\item
  The proper-cohomology/base-change input is the one read and recorded in LIFT085, including Stacks36.32, tag0BDM. The canonical Poincaré/rigidified construction is the one proved in LIFT084/085.
\end{enumerate}

The exact diagnostics check matrix valuation leading terms, torsion-field bounds, fractional-ideal directions, all-embedding norm scaling, the dual Key sign, every cost case including D=0 or n=0, and literal source fingerprints. They do not replace the universal CM/model/polarization/cohomology/source arguments and are not a Lean build.

\subsection{Canonical algebraic-geometric context}\label{canonical-algebraic-geometric-context-18}

The actual abelian schemes, integer-ring action, principal Picard/Poincare construction, normalized cubical line, finite-flat isogeny and its kernel live over Spec(O\_K). Away from N the displayed adjoint inverse gives an actual p-divisible-group and formal-group/Dieudonne comparison, without an operator or slope gain. Carried exact anchors: FGA232 Theorem3.1 pp148--149; FGA236 Corollaries1.7/2.3 pp229/232; SGA7 I VII2.7/2.9 pp169--174 and VIII7.4(b) pp310--311; SGA3 I4.7.3/4.7.3.1; SGA6 VI Theorem3.3; cote4 scan003 sections8.1.1--8.1.2 and scans005--006 sections8.4.1--8.4.3; cote54 scans005--007. No new lambda-ring/K-theory, l-adic weight or CM-motive identification. DP\_MULT Q3e/DP44/DP52/DP54--65, THETA3 G1/A7, CX\_GRH CX-ALG-024 and GCT\_PLETH T-052 retain their separate arithmetic comparisons.

\clearpage

\section{AM-K-103 --- a complete denominator, eligibility and coordinate cut at all-II support26}\label{am-k-103-a-complete-denominator-eligibility-and-coordinate-cut-at-all-ii-support26}

\textbf{Canonical checked integration --- 2026-11-23T03:47:18.015018+00:00.} Complete root and independent proof/script/interface reviews and separate exact replays passed. The actual canonical predecessor has103/31 entries; this integration gives104/31, complete floors26/11. Frozen author history is preserved. AM-K-103 SC26.1--28 removes nine complete subsets of the actual all-II globally rank-two-unipotent scalar-cross support26 class. Its support-free full-Z proof, new eta window and uniqueness, exact paired ideals and inactive weights, full infinite ideal tail, joint energy and positive-norm-cone counts retain every original factor. All1,229 cells give one293 charge. The residual necessarily has v=q26, 2\^{}1116\textless N\_E D\_eta\textless2\^{}2108, b=j0 and81\textless=k\textless=164. Actual R088103/31 becomes104/31, all earlier overlapping charges retained, both entire sums below2\^{}-230 and complete floors26/11 unchanged. Other types/classes, signed positivity, actual source, selected arithmetic content, prescribed period, specific032/Read, DP44 and RH remain open.

Author: /root/kummer\_extension, ALP\_MACHINE. 23 November 2026 UTC.

\textbf{Author whole saved-proof and complete-script review completed; independent review pending.} Arithmetic diagnostics accompany the geometric and counting proofs; they do not replace them. Earlier frozen findings remain unchanged.

\subsection{1. Actual scope and the new remaining sector}\label{actual-scope-and-the-new-remaining-sector}

Keep the actual J56 all-top kernel, H=133, f=168, m=164, the original source and target pencils at both CM components, every admitted label, selected real module and near-point normalization, finite projectors, compact factors, twists, central phases, Haar measures and class masses. Work in its current finite-unipotent G remainder and assume global rank-two unipotence explicitly. Restrict to scalar-cross matrices \[
 h_{SS}=I,\quad C=h_{W_0S}\ne0,\quad h=I+N,\quad
 \operatorname{rank}_E N=2,\qquad s=26,\quad
             \text{all active places type II}.             \tag{SC26.1}
\] The already charged DO099 interval removes transverse C and transverse inverse C through support70. Thus this theorem works after that actual containing-sector removal, with C in the original packet plane \(P=Eu\oplus Ee_n\). It does not import DO099's support25 four-compression removal.

Put \[
 Q=1427908193729121522601,\quad U=77Q^2-14,\quad
 R=17424/133,\quad M^2=R/U,\quad
 D_0=7^{140}29^{144}5^{84}13^{84},\quad D_F=2^{168}D_0.
                                                               \tag{SC26.2}
\] The normalized isotropic image line has slope eta, with eta bar(eta)=U. Let \(D_\eta\) be its exact integral E-denominator ideal and \(x=N_ED_\eta\). Let v count eligible active places, q count all eligible places among the168 selected T primes, k count nonzero original S covector values, b count inactive paired-C-zero places, and j count active places where the whole paired C is integral. Eligibility means \[
 \eta\equiv a_0=Q(-7+2\sqrt{-7})\pmod{\mathfrak p}          \tag{SC26.3}
\] at the original selected component, not at both CM partners.

\textbf{Theorem.} A single new absolute error less than \(2^{-293}\) removes all members of SC26.1 except the following necessary conditions: \[
 \boxed{\quad v=q=26,\quad 2^{1116}<x<2^{2108},\quad
        b=j=0,\quad81\le k\le164.\quad}                    \tag{SC26.4}
\] For each fixed support, eta is unique in this window; each active place has at most one actual type-II label; and the retained pointwise/weighted product has all168 factors \((T-1)^{-1}\).

Every removed set is intersected with the current remainder; every old error is retained. This is an actual support26 subset reduction, not support26 exhaustion or a change to the complete floors26/11. It does not assert existence of a surviving arrow or the signed sufficient inequality.

\subsection{2. Full positive compression without a support25 count}\label{full-positive-compression-without-a-support25-count}

A nilpotent of rank two has N³=0. If N²=0, unitarity gives N†=−N and N†N=0, so im N is totally isotropic. Dimension two contradicts Witt index one at a real embedding. Hence N²≠0. There is one length-three Jordan block and all other blocks have length one, so \[
 \ell=\operatorname{im}N^2=\operatorname{im}N\cap\ker N
       \text{ is a line},\qquad N+N^\dagger=N^2.           \tag{SC26.5}
\] The line is isotropic and N² is self-adjoint. Since N\_SS=0, its S/S compression is zero. The self-adjoint rank-one expression for N² forces its image line to have zero S component: its diagonal entries on positive S coordinates are a nonzero common scalar times the squared coordinate magnitudes. Thus \(\ell\subset W_0\), N²S=0, and \[
                         0\ne N(S)=C(S)=\ell .
\] The plane-C condition gives \(\ell\subset P\). Let \(Z=S\oplus u^\perp\) be the original positive167-space orthogonal to P. Self-adjointness gives N²Z=0, and then NZ lies in im N intersect ker N. Therefore \[
                     h_{ZZ}=I,\qquad N^2Z=0.             \tag{SC26.6}
\] No support variable, positive residue-field metric, or globally integral moving frame enters this argument. It agrees with the independently checked support26 taxonomy, section5.1; the argument is explicitly proved here.

The actual DO097 UC3--13 / DO098 CF4--7 normal form now applies: \[
 C=\ell r,\quad\ell=(1,\eta),\quad
 N=\begin{pmatrix}0&-a\ell^*J_P\\
                  \ell r&d\ell\ell^*J_P\end{pmatrix},
 \quad a=J_Z^{-1}r^*,\quad
 d+\bar d=-rJ_Z^{-1}r^*,\quad J_P=\operatorname{diag}(U,-1).
                                                               \tag{SC26.7}
\] Its full positive energy is \(rJ_Z^{-1}r^*=\alpha/U\le M^2\) at every embedding. Eta, the167 original covector values and d reconstruct h uniquely. The actual negative entry satisfies \(|\operatorname{Im}(-Ud)|<H/2\).

\subsection{3. Local eligibility and the first complete containing set}\label{local-eligibility-and-the-first-complete-containing-set}

Every original S input and each of the three fixed transverse inputs belongs separately to every actual source and target pencil at both CM components. DO103 LD6--11 uses precisely this membership, SC26.6 and the literal generators \(A_c=I+c\pi^{-1}zq_*\), \(q_*z=0\). Its local proof has no support restriction. A cross pole forces the saturated line ell to equal the actual pole line. A full double pole does so as well, by the rank-one leading term of \(A_cKA_c^{-1}\). Its selected slope is a0 and its partner slope is U/bar(a0). The identity \[
                 \bar a_0\eta-U=-\eta\,\overline{\eta-a_0}
\] makes these the same selected eligibility condition. A zero or pole of eta gives one zero coordinate of the saturated line, which cannot equal the actual pole line with two unit coordinates.

Consequently all original cross scalars have simple poles only on the v eligible active places; every full matrix entry has order at most two there and at most one at the other active places.

If v≤24, pad the eligible active set to24 of the26 active places, costing at most \(\binom{26}{24}=325\). Fix the padded set. Every row, Green, transverse, relative-norm and selected-real factor is then the unchanged complete DO098 CF8--32 count with a24-place cross ideal of norm T\^{}-48 and its original cap. The central difference ideal changes from \(iJ_0I_D^2\) to \(iJ_0I_\Sigma I_D\), with norm denominator T\^{}50. Each alternative central count CF14/CF16 is multiplied by T²; its covering radius can only decrease. The support count remains bounded by2\^{}168. Every complete all-II active label-weight sum is at most one. Therefore \[
 \binom{26}{24}T^2\,2^{-2879}
 <2^{-2879+9+2L_T}<2^{-879},\quad L_T=1019333/1024.        \tag{SC26.8}
\] The164-cell union inside CF32 is already included. This bound has no eta-denominator premise.

\subsection{4. A new denominator window and uniqueness argument}\label{a-new-denominator-window-and-uniqueness-argument}

The denominator-cleared ideal \((\eta-a_0)D_\eta\) is integral. The actual paired valuations over7 give the factor7\^{}84 even with eta denominators there. Its nonzero scalar has \(|\eta-a_0|<2\sqrt{U+14}\) at every complex embedding. The proof DO103 LD12--15 / K099 GE4 thus gives, independently of support, \[
                  7^{84}T^q<[4(U+14)]^{168}x.            \tag{SC26.9}
\] The exact actual integers satisfy \[
\begin{array}{ll}
 7^{84}T^{25}>2^{120}[4(U+14)]^{168},&
 7^{84}T^{26}>2^{1116}[4(U+14)]^{168},\\
 7^{84}T^{27}>2^{2111}[4(U+14)]^{168},&
 7^{84}T^{29}>2^{4102}[4(U+14)]^{168}.
\end{array}                                               \tag{SC26.10}
\] Below x=2\^{}4096, q≤28. If v≥25 then x\textgreater2\^{}120; if q≥26 then x\textgreater2\^{}1116. At v26 and x\textless2\^{}2108, q is exactly26. The support25 q≤26 conclusion is not used.

Two norm-U scalars with denominator norms below2\^{}4096 that agree at25 selected eligible primes also agree at their partners, by their norm equations. Their difference times the product of the two denominator ideals, if nonzero, is integral of norm at least T\^{}50. But its norm is less than \[
                           (4U)^{168}2^{8192}<T^{50}.     \tag{SC26.11}
\] They coincide. This is the new bounded-denominator uniqueness proof. For v25, the eligible active subset has26\textless2\^{}5 possibilities; for v26 no such choice is required. No denominator ideal or outside-support eligible set is separately enumerated in this window.

DO105 EL9--16's local one-label calculation contains no support restriction. An eligible active place permits at most one actual diagonal label, including a central-only pole, with exact weight1/(T−1). A noneligible active place is bounded by its complete all-II label sum, at most one. This is not a replacement of its fixed-label coefficient by one.

\subsection{5. General row and central ideals}\label{general-row-and-central-ideals}

Fix eta, the support and its eligible subset, and a canonical whole-C G orientation. Its ideal J has at most4\^{}334 patterns, and \[
 |a_G(h)|\le2Q^{-335\cdot167}(N_EJ)^{335/2},\qquad
                          N_EJ\le Q^{334}.               \tag{SC26.12}
\] Let I\_v have exponent−1 at both CM primes of each eligible active place. Primitive integer u-extraction from C and its negative coordinate gives all167 values in \[
 JI_v\cap\eta^{-1}JI_v=JI_vD_\eta=:L,\quad
 z=(N_EL)^{1/(2f)},\quad N_EJ=z^{2f}T^{2v}/x.             \tag{SC26.13}
\] This is a valuationwise equality, including denominators at G or inactive T.

For fixed eta,r, central differences are imaginary. Two integer Bezout extractions from the original4×4 block and one from each mixed block give \(\Delta d,\eta\Delta d,\bar\eta\Delta d\in J_0\mathcal O_EI_\Sigma I_v\). The exact descent \(D_\eta\bar D_\eta=B_\eta\mathcal O_E\), \(N_FB_\eta=x\), yields \[
 \Gamma_d=iJ_0B_\eta I_{\Sigma,F}I_{v,F},\qquad
 \operatorname{covol}\Gamma_d
       =\sqrt{D_F}(N_FJ_0)xT^{-(26+v)}.                   \tag{SC26.14}
\] This contains differences in one nonempty affine coset; d itself is not declared integral. Its fixed real-part shift is retained.

Using \(N_FJ_0\le N_EJ\) and the full-real750 transference constant, \[
 \#d\le V_F\!\left(H/2+750Uz^2T^{(v-26)/f}\right)
                U^{-f}T^{26+v}x^{-1},\quad
 V_F(C)=\frac{((22/7)fC^2)^{f/2}}{(f/2)!\sqrt{D_F}} .
                                                               \tag{SC26.15}
\] Only \(N_FJ_0\ge1\) was discarded from the covolume. The true transverse three-coordinate covolume is \(D_0^3z^{6f}/(4U)^f\), with cover less than \(2^{-10}\sqrt f\,z\). The actual selected real trace remains \[
                       |c(h)|<2^{-929f/2}U^{-f/2}z^{-f}. \tag{SC26.16}
\] These are the original metric and selected coefficient.

A chosen extra paired C congruence multiplies L by an ideal of E-norm T². Write e=1 for this case and e=0 otherwise, and let z denote the enlarged row scale. Exactly \[
 N_EJ=z^{2f}T^{2v}x^{-1}T^{-2e},\quad kz^2\le M^2,\quad
 z^2\le(Q^{334}T^{-2(v-e)}x)^{1/f}.                       \tag{SC26.17}
\] The central ideal is unchanged: its cover uses \(z^2T^{(v-26-2e)/f}\). The real trace uses the improved row scale.

\subsection{6. Inactive factors and extra paired ideals}\label{inactive-factors-and-extra-paired-ideals}

The local DO107 IP4--7 argument is support-free. At a noneligible inactive place the saturated null line and the pencil pole line have both transverse pairings units. Full-pencil preservation forces \(Nz=0\) and \(q_*N=0\), which force every saturated unipotent coefficient to vanish. If the paired C blocks are not both zero, the full pencil is not preserved. Since the original Lambda is included label0, it is the unique incidence; the literal projector coefficient is \[
                                  |\rho|=1/(T-1).        \tag{SC26.18}
\] Inverse adjointness supplies the partner, not another coefficient.

At b inactive paired-zero places, every original r(y) gains their paired ideals beyond D\_eta, even at eta poles. A nonzero S value therefore satisfies \[
                       xT^{2b-2v}\le |N_Er(y)|\le(R/U)^f.
                                                               \tag{SC26.19}
\] The exact inequalities \((R/U)^fT^{22}<2^{120}\) and \((R/U)^fT^{24}<2^{1116}\) give b≤13 for v25 and v26, respectively. Below2\^{}4096, at least139 noneligible places are inactive when v25, and at least140 when v26. Thus b\textgreater0 still leaves at least126 or127 factors SC26.18. Choosing one paired zero costs at most142\textless2\^{}8, and all167 values gain its normT² ideal, e=1.

For b0 retain g=142 when v=q25, g=139 when v25,q≥26, g=140 when v26 in the full bounded window, and g=142 when v26,x\textless2\^{}2108, where q26.

An active place at which the whole paired C is integral supplies the same normT² row gain relative to I\_v when v26. Eta is a unit there by eligibility. The central ideal stays unchanged, since its block can still have a double pole. Choosing one such active place costs26\textless2\^{}5. This is only the support-free local portion of DO108; its numerical25-bound is not imported.

Combine active and inactive factors before comparison: \[
                        (T-1)^{-(v+g)}<2T^{-(v+g)}.      \tag{SC26.20}
\] The single exact inequality \(2(T-1)^{168}>T^{168}\) proves this for all present v+g≤168. There is no extra inactive label count or second correction2.

\subsection{7. The infinite large-denominator tail}\label{the-infinite-large-denominator-tail}

At x≥X=2\^{}4096, no uniqueness is assumed. Enlarge cross denominators to all26 active places. DO104 CD5--9's exact eta count is support-independent: each exact denominator ideal has at most \(28\cdot49\cdot2^{28}<2^{39}\) norm-U scalars, using the actual factorization of U and28 norm-one roots. The whole central ideal is \(iJ_0B_\eta I_{\Sigma,F}^2\).

For each k, the complete original-coordinate count has dependence \[
 F_k(x)=x^{-337/2}
 \min\{M/\sqrt k,\ Q^{167/f}T^{-26/f}x^{1/(2f)}\}^{(328-2k)f}.
                                                               \tag{SC26.21}
\] For x≥X, \(F_k(x)\le F_k(X)(x/X)^{-k-9/2}\). Summing over every integral E-ideal, with \(\zeta_E(2)\le\zeta(2)^{336}<2^{336}\), gives \[
                  \sum_{N_ED\ge X}F_k(N_ED)<F_k(X)X^2 2^{336}.
                                                               \tag{SC26.22}
\] All ideal classes and eta numerators are included. The actual cap at X exceeds18; use row endpoint \(M/\sqrt{\max(k,18)}\).

Retain the independent transverse factor4, A3\textless2\^{}-4040, central radius \(67+98300/k\), support2\^{}168, G patterns2\^{}668, pivot2\^{}8, one Green2, eta factor2\^{}39, and generic all-II label sums at most one. The companion recalculates all164 cells at s26; every exponent is below−12189. This is a fresh support26 tail count.

\subsection{8. One complete bounded-window formula}\label{one-complete-bounded-window-formula}

Below2\^{}4096 and v≥25, eta is fixed after its eligible subset. Define \[
 K_k=\binom{164}{k}47^k2^{2fk}(f!)^k(87/(32fk))^{fk},\quad
 A_3=\frac{528^{3f}}{(3f)!D_0^3},\quad
 \mathcal B_k=\frac{(fk)!(3f)!}{((k+3)f)!}.
\] The unchanged joint-energy bound DO105 EL19--21 is \[
 \#r\le K_kA_3(4U)^f\mathcal B_k
           (M+2^{-10}z)^{2(k+3)f}z^{-2(k+3)f}.           \tag{SC26.23}
\] It integrates transverse slice volumes over the S-ball, not two independent full-radius counts. Eta,r,d reconstruct h, while the old harmless pivot factor remains.

At lower endpoint x≥2\^{}a, set \(c(v,a,e)^f=(R/U)^fT^{2(v-e)}Q^{-334}2^{-a}\). The exact lower integer caps, with both strict power brackets verified, are:

\begin{longtable}[]{@{}rrrr@{}}
\toprule\noalign{}
v & a & e & Lower cap \\
\midrule\noalign{}
\endhead
\bottomrule\noalign{}
\endlastfoot
26 & 4096 & 0 & 18 \\
25 & 120 & 0 & 68036 \\
25 & 1116 & 0 & 1117 \\
25 & 120 & 1 & 18 \\
26 & 1116 & 1 & 1117 \\
26 & 1116 & 0 & 4124955 \\
26 & 2108 & 0 & 68850 \\
\end{longtable}

Use \(z\le M/\sqrt{\max(k,c)}\) at this endpoint. As x grows, its admissible z cap grows at most as x\^{}(1/(2f)). The z-power in the full product contributes at most x\^{}(164−k), the enlarged row radius at most x\^{}(k+3), and the central radius at most x. Against x\^{}-337/2 the net growth is at most x\^{}-1/2. Hence each joint-row window has its maximum at the stated lower x, even across a cap change.

Let d count only explicit eligible-subset and chosen-place choices, and define \[
\begin{split}
 \mathcal C(v,a,g,e,d)={}&846+2f-929f/2+479+167fL_R
 -(335f/2)L_U\\
 &+(335v+26-g-335e)L_T-335(167)L_Q-(337/2)a+d,\\
 L_T={}&1019333/1024,\quad L_Q=287843/4096,\quad
 L_U=601356/4096,\quad L_R=28810/4096 .
\end{split}                                               \tag{SC26.24}
\] The846 retains support, common patterns, pivot, one Green2 and the single factor2 in SC26.20. The full T power is positive before its upper logarithm is used. The positive R and negative U powers are already grouped.

Let c\_k(A) be the strict integral binary ceiling of \[
 \frac{\binom{164}{k}47^k2^{2fk}(f!)^k87^{fk}(fk)!(3f)!}
 {32^{fk}f^{fk}k^{fk}A^{(164-k)f}((k+3)f)!},              \tag{SC26.25}
\] and j\_k(A) the corresponding ceiling of \((1+1/(1024\lfloor\sqrt A\rfloor))^{2(k+3)f}\). The full joint endpoint is \[
                  \mathcal C-4040+c_k(\max(k,c))+j_k(\max(k,c)).
                                                               \tag{SC26.26}
\] Replace479 by the proved central-volume ceiling where needed. An extra paired ideal gives central radius below68, since \(T^{2/f}>3690\) and c≥18. For v25,q≥26, radius155 is a safe bound. These are affine-cover estimates.

\subsubsection{Cone alternatives}\label{cone-alternatives}

For three sectors below, put \(t=M^2/(kz^2)\), \(t_0=c/k>1\), \(P_k=\max(10000,t_0)\). Apply the complete K102 PC6--16 positive-norm count to this same ideal. Its relative phases are bounded by \((28\,2^{20})^k t^{fk/11}\); its full-real dual theta factor is4 per F-lattice; its transverse determinant and joint Laplace energy are unchanged. No beta-pivot or second eta count is added.

The complete factors before this row count are \[
 2^{846-929f/2+d}Q^{-335\cdot167}T^{335v+26-g}
 x^{-337/2}U^{-3f/2}z^{334f}
 V_F(H/2+750Uz^2T^{(v-26)/f}).                            \tag{SC26.27}
\] For t≤t0 the actual G cap forces \(x\ge2^a(t_0/t)^f\). Choosing the Laplace parameter A=(3/5)/t0 gives an increasing expression, by the full exponent \((3/2+k/11)f\) and the central-radius derivative from K102 PC20. On {[}t0,Pk{]}, choose A=(3/5)/t; the logarithm is affine plus positive exponentials in log t and a central log-sum-exp, hence convex. Beyond Pk the original complex-volume bound decreases, including at k164. Thus the alternative endpoints are t0,Pk,Pk, with no extra union factor. Pk may depend on k, which is already separately counted.

At every such endpoint kt is an integer and the central radius is below68. The script evaluates K102 PC23--26 and its full complex-volume alternative by exact factorials, powers and rational ceilings. The unchanged constants log2(e)\textless13/9, pi\textless22/7 and C11\textless2\^{}20 are retained.

\subsection{9. Nine containing sectors and one fresh union}\label{nine-containing-sectors-and-one-fresh-union}

All finite-window rows below have x\textless2\^{}4096. Each coordinate cell already includes support, original positions, common G patterns, central fibre, weighted labels and every displayed choice.

\begin{longtable}[]{@{}
  >{\raggedright\arraybackslash}p{(\columnwidth - 6\tabcolsep) * \real{0.2500}}
  >{\raggedright\arraybackslash}p{(\columnwidth - 6\tabcolsep) * \real{0.2500}}
  >{\raggedright\arraybackslash}p{(\columnwidth - 6\tabcolsep) * \real{0.2500}}
  >{\raggedright\arraybackslash}p{(\columnwidth - 6\tabcolsep) * \real{0.2500}}@{}}
\toprule\noalign{}
\begin{minipage}[b]{\linewidth}\raggedright
Group
\end{minipage} & \begin{minipage}[b]{\linewidth}\raggedright
Containing sector
\end{minipage} & \begin{minipage}[b]{\linewidth}\raggedright
Endpoint/count parameters
\end{minipage} & \begin{minipage}[b]{\linewidth}\raggedright
Cells and strict exponent bound
\end{minipage} \\
\midrule\noalign{}
\endhead
\bottomrule\noalign{}
\endlastfoot
A & v≤24, any x & 325T² times complete CF32 & 1 below−879 \\
B & x≥2\^{}4096, any v & Infinite ideal tail & 164 below−12189 \\
C & v25,b\textgreater0 & a120,g126,e1,d13,cap18 & 164 below−172990 \\
D & v26,b\textgreater0 & a1116,g127,e1,d8,cap1117 & 164 below−47200 \\
E & v=q25,b0 & a120,g142,e0,d5,cap68036; cone & 164 below−1567 \\
F & v25,q≥26,b0 & a1116,g139,e0,d5,cap1117,central155 & 164 below−58950 \\
G & v26,b0,x≥2\^{}2108 & a2108,g140,e0,d0,cap68850; cone & 164 below−1249 \\
H & v26,b0,x\textless2\^{}2108,k≤80 & a1116,g142,e0,d0,cap4124955; cone & 80 below−294 \\
I & v26,b0,x\textless2\^{}2108,j\textgreater0 & a1116,g142,e1,d5,cap1117 & 164 below−62130 \\
\end{longtable}

First remove A and B. Then v is25 or26 in the bounded window. Remove b\textgreater0. At v25, E and F cover every q. At v26, remove the high window, the sparse range and active paired-integrality. The remaining constraints force q26. This proves exhaustion outside SC26.4.

The1229 bounds are summed conservatively, including any overlapping containing sets; older charges are retained. For cone cells only the maximum of the alternative endpoint bounds is charged. The certificate records every rational exponent, takes a sufficient dyadic margin, and checks \[
               \sum_{\text{all1229 cells}}2^{-e_{\rm cell}}
                                      <2^{-293}.         \tag{SC26.28}
\] It neither multiplies an endpoint maximum by three nor silently deletes overlaps.

The exact largest upper exponent is \[
 -15077367706059399527/51109284870082560<-294,
\] in group H at k=80. The preliminary replay passed7,984 assertions over all1,229 complete counting cells. A separate unchanged-byte final replay pins this saved proof and the current chronological predecessor.

\subsection{10. Chronology and retained interfaces}\label{chronology-and-retained-interfaces}

The actual baseline is checks/r087\_joint\_remainder\_error.json:102 all-top and31 uniform errors, floors26/11. The script accepts a later complete prefix, preserves the entire ordered baseline and appends293 once. Both complete sums stay below2\^{}-230. No full support27 conclusion follows.

Section2 proves the support-free full-Z compression. The actual local inputs are DO097 UC3--13, DO098 CF4--16, DO103 LD6--15, DO104 CD10--22, DO105 EL9--27, DO107 IP4--15 and DO108 AC3--9, applied only as proved here. The eta multiplicity and infinite tail use DO104 CD5--9/23--26 with newly calculated support26 dependence. The norm-phase and cone estimates use the whole checked K102 PC6--26. No support25 numerical cut, denominator window, uniqueness or total eligibility hypothesis is silently extended.

The field, involution, Hermitian polarization, theta line, selected metric, projectors, class contraction and original correspondence remain fixed. Exact carried anchors are SGA3 I4.7.3/4.7.3.1; SGA6 VI Theorem3.3; FGA232 Theorem3.1 pp148--149; FGA236 Corollaries1.7/2.3 pp229/232; SGA7 I VII2.7/2.9 pp169--174 and VIII7.4(b) pp310--311; cote4 scan003 §§8.1.1--8.1.2, scans005--006 §§8.4.1--8.4.3; cote54 scans005--007. DP\_MULT Q3e/DP44/DP52/DP54--65, THETA3 G1/A7, CX\_GRH CX-ALG-024 and GCT\_PLETH T-052 retain their separate comparisons. No new lambda-ring, K-theoretic, motivic, selected ordinary-derivative/content, prescribed Bessel/Mram period, mixed source, Read/STiesR, DP44 or RH conclusion is asserted.

\subsection{Canonical algebraic-geometric context}\label{canonical-algebraic-geometric-context-19}

The reductions are the actual integral Hermitian modules at unmixed finite points; content and coordinate ideals remain fractional invertible modules over Spec(O\_E) and Spec(O\_F), with the same paired CM action, polarization, theta line, Picard scheme and biextension. Carried exact anchors: SGA3 I4.7.3/4.7.3.1; SGA6 VI Theorem3.3; FGA232 Theorem3.1 pp148--149; FGA236 Corollaries1.7/2.3 pp229/232; SGA7 I VII2.7/2.9 pp169--174 and VIII7.4(b) pp310--311; cote4 scan003 sections8.1.1--8.1.2 and scans005--006 sections8.4.1--8.4.3; cote54 scans005--007. No new lambda-ring/K-theory, l-adic weight or CM-motive comparison is asserted. DP\_MULT Q3e, THETA3 G1/A7, CX\_GRH CX-ALG-024 and GCT\_PLETH T-052 retain their stated global gates.

\clearpage

\section{AM-DO-109 --- weighted finite patterns and eligible pencils remove the short fixed-block sector at26}\label{am-do-109-weighted-finite-patterns-and-eligible-pencils-remove-the-short-fixed-block-sector-at26}

\textbf{Canonical checked integration --- 2026-11-23T03:48:34.306084+00:00.} Complete root and independent proof/script/interface reviews and separate exact replays passed. The actual canonical predecessor has104/31 entries; this integration gives105/31, complete floors26/11. Frozen author history is preserved. AM-DO-109 SP26.1--25 removes the entire actual fixed-S globally rank-two-unipotent support26 low-D\_G\textless=32 short-pivot class. The original plane reduction and paired alpha budget give2i+t\textless=5. The actual V-input eligibility and central-only label proof, true Gram metric, affine central difference lattice and full weighted four-state G sum are retained. The three complete p\textless=24,p25,p26 endpoints give one820 charge; the p26 branch proves eta uniqueness before summing actual labels. Together with already charged R079 high-D\_G and last-only pivots, this completes the whole fixed-S rank-two-unipotent support26 structural class. Actual K103104/31 becomes105/31, every overlap retained, both entire sums below2\^{}-230 and complete floors26/11 unchanged. Other support26 classes, signed positivity, actual source, selected arithmetic content, prescribed period, specific032/Read, DP44 and RH remain open.

Author: \nolinkurl{/root/definite_oscillator}, ALP\_MACHINE. 23 November2026 UTC. Controlled pending proof. Author whole-file review, exact replay, independent review and canonical integration are separate events. Existing canonical records and earlier frozen findings are unchanged.

\subsection{1. Actual sector and theorem}\label{actual-sector-and-theorem-8}

Keep the unchanged actual J56 all-top normalized kernel, its charged H=133 cutoff, f=a=168, n=169, original E/F, paired CM components, all admitted source/target pencil labels and omitted endpoints, finite projectors, selected Hilbert metric, packet point, central representative, phase, class masses and Haar measures. Work in the original K068 finite-unipotent remainder. Consider

\[
 h=I_S\oplus h_0=I+N,\quad \dim S=164,\quad
 \operatorname{rank}_E N=2,\quad N^3=0,\quad
 \operatorname{rank}_E N^2=1,\quad s=26.                 \tag{SP26.1}
\]

Let d\_v be the actual reduction rank at each of the334 unmixed relative G places and D\_G=sum\_v d\_v. Global unipotence and fixation of S are hypotheses. In the D\_G\textless=32 case let r be the original transverse covector defined below, and let a\_j=r(y\_j).

\textbf{Theorem.} The complete SP26.1 subset with D\_G\textless=32 and at least one of a1,a2 nonzero has total absolute actual normalized contribution less than2\^{}-820. Every original support, incidence type, paired label, finite pattern, covector, phase and central fibre is included.

R079 already removes the D\_G\textgreater=33 and the D\_G\textless=32,a1=a2=0 support26 subsets with its earlier1120 charge. Thus this new theorem completes the fixed-S globally rank-two-unipotent support26 class, while retaining that older charge. It makes no complete support27, signed inequality, global period, source-construction, selected arithmetic-image, specific032 or DP44 assertion.

Use the literal host constants

\[
 g_0=22285932761,\quad g_1=64072175441,\quad Q=g_0g_1,
 \quad U=77Q^2-14,\quad R=17424/133,\quad M=\sqrt{R/U},
\] \[
 L_T=1019333/1024>\log_2T,\quad
 L_Q=287843/4096<\log_2Q,\quad
 L_U=601356/4096<\log_2U,\quad
 L_R=113/16>\log_2R.                                  \tag{SP26.2}
\]

The exact host certificate and the companion below verify these directions by integer powers. No field, prime or selected real weight is changed.

\subsection{2. The original plane, metric and small type-I count}\label{the-original-plane-metric-and-small-type-i-count}

R058(HB20--22), using the separate source and target pairings of DO078(FP20--27), proves at D\_G\textless=32 that ell=im N\^{}2 lies in the actual packet plane P=Eu perpendicular Ee\_n.~Its integral short-pairing argument has no support restriction. The positive definite quotient ell\^{}perp/ell has a unipotent unitary action, hence identity. With V=u\^{}perp in the original positive E\^{}4 and Z=S perpendicular V, it follows that h\_ZZ=I and N\^{}2Z=0. The cross map C:V→P is nonzero: otherwise N would act only on P and have square zero. R077(FB8) and R079(PC6) therefore give the exact normal form

\[
 C=\ell r,\quad\ell=(1,\eta),\quad \eta\bar\eta=U,
 \quad v=J_Z^{-1}r^*,\quad q=rv>0,\quad\alpha=(N^2)_{nn}=Uq,
\] \[
 N=\begin{pmatrix}0&-v\ell^*J_P\\
                  \ell r&d\ell\ell^*J_P\end{pmatrix},
 \quad J_P=\operatorname{diag}(U,-1),\quad
 d+\bar d=-q,\quad w=N_{nn}=-Ud.                       \tag{SP26.3}
\]

Here r vanishes on S. C and w determine the entire original matrix, and Re(w)=alpha/2 is fixed by C. There is no remaining column or isometry fibre. At every embedding 0\textless alpha\_tau\textless=R and \textbar\textbar r\_tau\textbar\textbar\textless=M in the true positive dual metric.

Use the original primitive u and transverse vectors

\[
u=(2,95192390527,383092671043,12529843548438452388609),
\] \[
y_1=(-976627,-153032,38026,0),\quad
y_2=(-567678,695570,-172838,0),\quad
y_3=(256766,-7654233287,-30805128044,1).
\]

Their Gram matrix is

\[
 G=\begin{pmatrix}
 978665066829&441392856078&-303935233242\\
 441392856078&835948910828&95913020934\\
 -303935233242&95913020934&1007543201084982869062
 \end{pmatrix},\quad \det G=4U,\quad G>2^{16}I.        \tag{SP26.4}
\]

The fixed frame has determinant of modulus2U and is a unit at every selected G and T prime. No moving orthogonal frame is declared integral. Choose the first nonzero short value among a1,a2 and put beta=eta a\_j. Its actual squared length B=\textbar\textbar y\_j\textbar\textbar\^{}2 satisfies B\textless2\^{}40. We retain the harmless common pivot allowance3\textless2\^{}2 from R079 even though only two short choices occur.

The alpha ideal in R058(HB22) is unchanged:

\[
 \alpha\in U\mathfrak J_F I_{\rm II,F}^2 I_{\rm III,F},
 \quad N_F\mathfrak J_F=\prod_vq_v^{2-d_v}
                 \ge Q^{334}/g_1^{32}.                \tag{SP26.5}
\]

This uses the original entry (N\^{}2)11=−4alpha/U, its actual paired entry poles, and U's retained coprimalities; it loses no dyadic factor. Write i,u,t for the numbers of types I,II,III. The exact integer inequality

\[
 (133U)^{168}Q^{334}>17424^{168}g_1^{32}T^{46}          \tag{SP26.6}
\]

shows that mu=2u+t\textless=46 is impossible: SP26.5 would give N\_F(alpha)\textgreater R\^{}f.~Since i+u+t=26, every actual remaining profile satisfies

\[
 2u+t\ge47,\qquad 2i+t\le5,\qquad i\le2.             \tag{SP26.7}
\]

This is a fresh use of the exact inequality for the indicated profiles, not an assertion that R058's earlier support23 numerical conclusion already covered26. The twelve possible triples are i=0,t=0..5; i=1,t=0..3; and i=2,t=0..1, with u=26−i−t.

\subsection{3. Actual eligible type-II places and the common containing ideals}\label{actual-eligible-type-ii-places-and-the-common-containing-ideals}

At a selected T component use the literal pencil M\_c=A\_c Lambda, A\_c=I+c pi\^{}-1 z q\_\emph{, (zq\_})\^{}2=0, with the inverse-adjoint paired generator at the other CM component. The admitted labels omit infinity and−1 and include0; an active type-II place has c!=0. Let a0=Q(−7+2sqrt(−7)), so a0 bar(a0)=U+14. Call a type-II place eligible when eta≡a0 at its original selected pi-component. DO103(LD5--11)'s original pole-line calculation gives the paired condition eta≡U/bar(a0) at the other component. Thus both partners specify the same eligibility condition, and eta is a unit at both of them.

We need the local statements on V, although the original scalar-cross applications had C\textbar S!=0. Their proofs use the fixed integral inputs and SP26.3 only, as follows.

Every y\_j lies separately in every source and target pencil, at both CM partners. At type I, h preserves the full-pencil intersection, which contains y\_j and lies in Lambda; hence Ny\_j is integral. At type III, Ny\_j lies in M\_d, while (h\^{}-1−I)y\_j=−Ny\_j lies in M\_c because N\^{}2y\_j=0. The distinct-pencil intersection lies in Lambda, so Ny\_j is integral there too. Repeating at the partner proves paired integrality. Primitive integer Bezout extraction from u makes both a\_j and eta a\_j integral in those original scalar coordinates.

At type II, a pole of one of these C scalars gives a leading image vector on the actual target pole line. The saturated line E(u,eta) can equal that line only if eta is a local unit, because the actual line has both a positive unit coordinate and a negative unit coordinate. The slope at the selected prime is a0, and the inverse-adjoint partner has slope U/bar(a0), exactly as in DO103(LD7--11). Thus every such cross pole lies at an eligible type-II place.

The same conclusion holds for a full entry of N of pole order2. Indeed, if K=A\_c\^{}-1 h A\_c is integral, then

\[
 N=K-I+c\pi^{-1}(zq_*K-Kzq_*)
                 -c^2\pi^{-2}z(q_*Kz)q_* .             \tag{SP26.8}
\]

Any nonzero double-pole image is the actual pole line. In SP26.3 the cross blocks have at most simple poles and N\_ZZ=0, so a double pole can occur only in the rank-one PP block, with saturated image E(u,eta). This again forces the same eligibility. The argument applies in the paired generator too. It requires C\textbar V!=0 and the displayed normal form, but nowhere requires a nonzero S coordinate.

Let E\_0 be the eligible type-II set and p=\textbar E\_0\textbar. Pad E\_0 to a set D of v active places, where v will be24,25 or26 below. The original G table of R077(FB9--13)/R079(PC9,14) gives a canonical whole-C ideal J. At a relative G place of norm q\_v, its local states are: rank0 gives both CM factors (E-norm q\_v\^{}2); rank1 gives one of the two orientations (E-norm q\_v); rank2 gives neither. At each state both the three a\_j and beta have this divisibility. Consequently

\[
 L=J I_D,\quad N_EI_D=T^{-2v},\quad z=(N_EL)^{1/(2f)},
 \quad a_1,a_2,a_3,\beta\in L,
 \quad |a(h)|\le2Q^{-335\cdot167}T^{335v}z^{335f},
\] \[
 N_EJ=\prod_vq_v^{2-d_v}\le Q^{334},\qquad
 z\le z_v:=Q^{167/f}T^{-v/f}.                          \tag{SP26.9}
\]

The paired G orientation is one common orientation for the whole cross map. It is not an independent orientation for each coordinate or for each CM partner. Rank2 imposes no extra ideal condition. Both mixed finite factors retain their modulus bound one.

Let J0 be the F-product of the actual rank-zero G places once. Then N\_EJ\textless=Q\^{}167 N\_FJ0 and N\_FJ0\textless=Q\^{}167. The full original entry w has at most simple poles everywhere in the active support Sigma and one additional pole only on E\_0. After padding, the difference of any two w at fixed C lies in

\[
 \Gamma=iJ_0 I_{\Sigma,F}I_{D,F},\quad
 \operatorname{covol}(\Gamma)=\sqrt{D_F}\,N_FJ_0T^{-(s+v)},
 \quad D_F=2^{168}D_0.                                \tag{SP26.10}
\]

This is one affine difference coset, with its shift retained; w itself is not asserted to lie in the ideal. Since O\_E intersect iF=iO\_F, there is no half-lattice or factor2 loss. The original full-real codifferent and transference give cover below750sqrt(f)(N\_FJ0 T\^{}\{-(s+v)\})\^{}(1/f). The exact comparison750\^{}f Q\^{}167\textless T\^{}50 makes this less than sqrt(f) for v\textgreater=24. With \textbar Im(w)\textbar\textless133/2, the complete real central volume at radius68sqrt(f) gives

\[
 \#w<2^{479}\frac{T^{s+v}}{N_FJ_0}
       \le2^{479}Q^{167}T^{s-v}z^{-2f}.                \tag{SP26.11}
\]

All containments include every prime away from T: the original entries are integral there, the displayed G factors retain the prescribed divisibility, and the original primitive Bezout extraction introduces no denominator. The actual V/P frame is used only at its G/T unit primes and for its exact archimedean metric.

\subsection{4. Full covector, phase and trace counts; weighted G summation}\label{full-covector-phase-and-trace-counts-weighted-g-summation}

The half-trace scalar covolume is D0 z\^{}(2f), D0=7\^{}140 29\^{}144 5\^{}84 13\^{}84. In the true three-value metric G\^{}-1 the covolume is D0\^{}3 z\^{}(6f)/(4U)\^{}f.~R079(PC10--11)'s arbitrary-ideal product cover is below4sqrt(f)z, because (2250/7)\^{}2·3\textless16·2\^{}16. Thus

\[
 \#r\le A_3(4U)^f(M+4z)^{6f}z^{-6f},\qquad
 A_3=528^{3f}/((3f)!D_0^3)<2^{-4040}.                  \tag{SP26.12}
\]

This is a whole covector volume bound, valid for every ideal and affine-free row lattice, not an expected-density factor. The actual reconstruction uses a nonzero short pivot. If eta has not yet been fixed, the full relative norm equation beta bar(beta)=U a\_j bar(a\_j), with beta in L, has the R079(PC12) count

\[
 \#\beta\le28C_F(RB)^{f/12}z^{-f/6},\qquad
 28<2^5,\quad C_F<2^{42},\quad B<2^{40}.               \tag{SP26.13}
\]

This counts every actual ideal class and norm-one phase through the integral relative norm ideal at an existing solution, as in K075(RN17--19)/DO091(TD11--12). It does not assume a principal fractional ideal or a globally integral eta. Once r and beta are fixed, eta=beta/a\_j is determined.

The genuine selected trace, with both compact factors, actual metric and near-point transport, is exactly the R079(PC13) bound

\[
 |c(h)|<2^{-929f/2}(U/B)^{-f/2}z^{-f}.                 \tag{SP26.14}
\]

Its premise is the actual nonzero short pivot and the original ideal minimum. It is not an operator-norm replacement. Overcounting zero rows in SP26.12 does not assert that SP26.14 holds for such artificial rows; it remains a uniform upper bound on the actual subset being counted.

For fixed support, type assignment, padded D and labelled pair, combine SP26.9 and11--14. Apart from positive constants independent of J, the complete z dependence is

\[
 F_v(z)=z^\kappa(M+4z)^{6f},\qquad
 \kappa=(335-6-1/6-2-1)f=(1955/6)f>0.                \tag{SP26.15}
\]

For0\textless z'\textless=z, F\_v(z')/F\_v(z)\textless=(z'/z)\^{}kappa. Lowering one local E-ideal exponent from2 to1 decreases z by q\_v\^{}(-1/(2f)). Summing the two possible rank1 orientations, the rank2 state and the rank0 state therefore gives, relative to the full-rank-zero cap,

\[
 1+2q_v^{-1955/12}+q_v^{-1955/6}
                  =(1+q_v^{-1955/12})^2.              \tag{SP26.16}
\]

Across all334 places this product is below (1+2\^{}-5534)\^{}668\textless2, since each q\_v\textgreater2\^{}34 and34·1955/12\textgreater5534. Thus every finite pattern and its two genuine orientations are included for one factor2. No additional4\^{}334 or independent CM-partner factor is inserted. We may enlarge from D\_G\textless=32 to all such patterns for this upper sum after applying the low-D\_G structural restrictions to the actual subset. If eta is already fixed, SP26.13 is omitted, kappa becomes326f\textgreater1955f/6, and the same product bound remains valid.

At any fixed original labelled tuple the bounds above are uniform. K071(NB7--8) sums the complete actual paired incidence weights to at most T\^{}i for prescribed types. TypeII andIII cost at most1. Support choices cost less than2\^{}168. Padded sets, type assignments and pivots will be charged explicitly in the next section. All remaining phase, class and Haar factors stay in their original normalization, with their retained modulus bound one.

\subsection{5. Three exhaustive eligibility cases and exact endpoints}\label{three-exhaustive-eligibility-cases-and-exact-endpoints}

The integer comparisons

\[
 z_{24}<M/4,\qquad z_{25}<M/256,\qquad
 (65/64)^{6f}<2^{23},\qquad z_{26}\le z_{25}             \tag{SP26.17}
\]

are checked without floating point. For v24 use M+4z\textless2M; for v25,26 use M+4z\textless(65/64)M. There is no separate cap-region union.

For bookkeeping, when SP26.13 is retained define the following complete logarithmic upper bound with s=26:

\[
\begin{split}
 E(v,i,b_{\rm type},b_{\rm pad},b_{\rm rad})={}&
 1+1+168+b_{\rm type}+2+479+5+42-4040+2f+b_{\rm rad}
       -929f/2\\
 &+(37f/12)L_R-(5f/2)L_U+(7f/12)40
       -(49\cdot167/6)L_Q\\
 &+(49v/6+s+i)L_T+b_{\rm pad}.                         \tag{SP26.18}
\end{split}
\]

The first two1's are the Green Euler factor and the complete weighted G sum. The2f is the4\^{}f metric determinant; the radius allowance is separate. The phase/divisor allowances5+42 occur exactly once. The central T\^{}(s−v) and actual type weight T\^{}i have already been included. Every power of R,U,B is combined with its final sign before the directed logarithms are substituted.

\textbf{Case A: p\textless=24.} Pad the eligible set to24 of the26 active places, with at most binom(26,24)=325\textless2\^{}9 choices. SP26.7 gives i\textless=2. Retain the full type assignment bound3\^{}26\textless2\^{}52, even though some profiles are impossible. Use b\_rad=6f. The whole case satisfies

\[
 \log_2 U_A<E(24,2,52,9,1008)
       =-170072485/24576<-6920.                        \tag{SP26.19}
\]

\textbf{Case B: p=25.} Necessarily u\textgreater=25. The original assignments consist only of allII, or oneI and25II, or oneIII and25II. Their total number is1+2·26=53\textless2\^{}6, before any eligibility restriction. Pad to25 places, with at most26\textless2\^{}5 choices, and use i\textless=1 and b\_rad=23. The whole case satisfies

\[
 \log_2 U_B<E(25,1,6,5,23)
       =-20183369/24576<-821.                          \tag{SP26.20}
\]

\textbf{Case C: p=26.} All places are typeII and fully eligible. We first prove that eta is unique for each support within the actual D\_G\textless=32 short sector. Let D\_eta be its exact denominator ideal and x=N\_ED\_eta. The nonzero short a\_j and eta a\_j both lie in L=JI\_Sigma, so

\[
 a_j\in L\cap\eta^{-1}L=LD_\eta,
 \quad x\le(BR/U)^fT^{52}/N_EJ
        \le(2^{40}R/U)^fT^{52}g_1^{32}/Q^{334}<2^{13100}.
                                                               \tag{SP26.21}
\]

The last inequality uses only the exact host bounds R\textless2\^{}8, U\textgreater2\^{}146, Q\textgreater2\^{}70, g1\textless2\^{}36 and T\textless2\^{}996: their total exponent is13100. It is uniform in the G pattern and pivot. No eta-ideal enumeration has yet been made.

If two such eta differ, (eta1−eta2)D\_eta1 D\_eta2 is a nonzero integral ideal. Full eligibility gives the same residue a0 at every selected component, and eta bar(eta)=U gives the same residue U/bar(a0) at every partner. Both denominator ideals are units there. Consequently its norm is at least T\^{}52, while the archimedean bound gives at most(4U)\^{}f x1x2. The exact strict comparison

\[
 (4U)^f2^{26200}<T^{52}                                \tag{SP26.22}
\]

is a contradiction. Thus there is at most one eta for each actual26-place support. The assertion is before choosing any G pattern or diagonal label, and absence of such eta only lowers the count.

For completeness, the one-label result of DO105(EL9--15) applies here with C\textbar V!=0. At an eligible place write the actual column and row as w0=(u,eta) and w0\^{}dagger=(u\emph{,-bar(eta)). With integral lifts w0=a\_l z+pi v and w0\^{}dagger=b\_l q\_}+pi t, isotropy gives b\_l(q\_*w0/pi)+a\_l(w0\^{}dagger z/pi)=0 modulo pi. The two reductions A\_c\^{}-1w0 and w0\^{}dagger A\_c therefore vanish at the same unique label, if one exists. A pole in either cross row forces that label by applying N to the corresponding original Z input, at the appropriate paired component. If both rows are integral but a central pole remains, the mixed term is integral after conjugation and the central term is d(A\_c\^{}-1w0)(w0\^{}dagger A\_c); its integrality again forces the same label. This includes central-only double poles and an omitted cancellation label, which yields no actual active arrow. The argument uses the T-unit fixed V frame and does not need C\textbar S!=0.

Hence, after eta is fixed, the entire diagonal-label/weight sum is at most

\[
 (T-1)^{-26}<2T^{-26}.                                 \tag{SP26.23}
\]

This is a sum over actual labels, not a per-fixed-tuple saving used before fixing eta. The beta fibre and its28C\_F factor are now absent because beta=eta a\_j is determined. Keep the full row, central and actual trace factors. At the cap the remaining z exponent is326f, the total Q exponent is−8·167 and, after SP26.23, the T exponent is8s. The complete case is bounded by

\[
\begin{split}
 E_C={}&1+1+168+2+479-4040+2f+23-929f/2
      +3fL_R-(5f/2)L_U+(f/2)40\\
     &-8\cdot167L_Q+8sL_T+1
       =-11592835/512<-22642.                           \tag{SP26.24}
\end{split}
\]

The last1 is precisely the finite correction in SP26.23. There is no additional type, padding, eta, phase, generator or inverse-image multiplicity in this fully eligible case. The larger all-pattern sum still bounds the actual fixed-eta low-D\_G subset.

The three p ranges are exhaustive and disjoint. The complete union is therefore

\[
 U_A+U_B+U_C<2^{-6920}+2^{-821}+2^{-22642}<2^{-820}.      \tag{SP26.25}
\]

This proves the theorem. No witness or type case remains outside the charged union.

\subsection{6. Verification, chronology and scope}\label{verification-chronology-and-scope}

The companion reads the actual J56 host, checks the literal Gram and frame determinants and positive minors, all discriminant/volume constants, the alpha norm inequality, every profile and padded-choice bound, both G-pattern exponents, the exact cap and uniqueness powers, the local cancellation identity over small odd finite fields, the three rational endpoints and complete dyadic union. These diagnostics do not prove automorphic or local source interfaces by themselves; those are supplied in the written argument and whole source readings.

The default actual predecessor is \nolinkurl{checks/r088_joint_remainder_error.json}, with103 all-top and31 uniform entries and complete floors26/11. Its first102 all-top entries are the unchanged R087 base; R088's523 charge follows them. The script permits a later actual predecessor only when its all-top array extends that exact base prefix and its uniform array is unchanged. Append820 once, retain every earlier charge, and check both full dyadic sums below2\^{}-230. Remove only the intersection with the actual current remainder. The independent R088 rank-one523 charge remains before this charge. The complete floor stays26; the fixed-S rank-two class is the only entire structural class newly completed here.

The new inputs are the weighted four-state G sum, the actual V-input pole/label extension, and the low-D\_G fully eligible eta uniqueness with its complete counting application. The inherited geometry remains over the same fractional invertible modules on Spec(O\_E)/Spec(O\_F), with the same polarization, theta line, Picard and biextension data. Exact carried context: SGA3 I4.7.3/4.7.3.1; SGA6 VI Theorem3.3; FGA232 Theorem3.1 pp148--149; FGA236 Corollaries1.7/2.3 pp229/232; SGA7 I VII2.7/2.9 pp169--174 and VIII7.4(b) pp310--311; cote4 scan003 §§8.1.1--8.1.2, scans005--006 §§8.4.1--8.4.3 and cote54 scans005--007. These are inherited contexts, not new source readings or a new lambda-ring, K-theory, l-adic weight or CM-motive comparison.

DP\_MULT Q3e and THETA3 G1/A7 gain precisely this finite/real sector reduction. The remaining support26 classes, complete signed sufficient inequality, fully sized literal mixed source, selected arithmetic image/content, prescribed global period, specific032/Read pairing and every actual DP44 cost remain separate gates. CX\_GRH CX-ALG-024, global GCT\_PLETH T-052 and RH conclusions remain open. No Lean build or archive audit is asserted.

\subsection{Canonical algebraic-geometric context}\label{canonical-algebraic-geometric-context-20}

The reductions are the actual integral Hermitian modules at unmixed finite points; content and coordinate ideals remain fractional invertible modules over Spec(O\_E) and Spec(O\_F), with the same paired CM action, polarization, theta line, Picard scheme and biextension. Carried exact anchors: SGA3 I4.7.3/4.7.3.1; SGA6 VI Theorem3.3; FGA232 Theorem3.1 pp148--149; FGA236 Corollaries1.7/2.3 pp229/232; SGA7 I VII2.7/2.9 pp169--174 and VIII7.4(b) pp310--311; cote4 scan003 sections8.1.1--8.1.2 and scans005--006 sections8.4.1--8.4.3; cote54 scans005--007. No new lambda-ring/K-theory, l-adic weight or CM-motive comparison is asserted. DP\_MULT Q3e, THETA3 G1/A7, CX\_GRH CX-ALG-024 and GCT\_PLETH T-052 retain their stated global gates.

\clearpage

\section{AM-K-104 --- full-volume removal of the next all-II support26 coordinate interval}\label{am-k-104-full-volume-removal-of-the-next-all-ii-support26-coordinate-interval}

\textbf{Canonical checked integration --- 2026-11-23T03:58:30.046321+00:00.} Complete root and independent proof/script/interface reviews and separate exact replays passed. The actual canonical predecessor has105/31 entries; this integration gives106/31, complete floors26/11. Frozen author history is preserved. AM-K-104 FV1--14 removes the complete new coordinate interval81\textless=k\textless=95 from the actual K103 all-II globally rank-two-unipotent scalar-cross support26 residual. The original full joint-volume lattice, actual central affine fibre and selected trace are retained. The full334-place four-state weighted G product is below2; the exact support binomial is below2\^{}101. Complete coupled-x monotonicity gives the actual lower denominator endpoint without a missing range. All15 cells, with worst-1128627/4096, give one274 charge. The all-II residual now has v=q26, 2\^{}1116\textless x\textless2\^{}2108, b=j0 and96\textless=k\textless=164. Actual DO109105/31 becomes106/31, all prior overlaps retained, both entire error sums below2\^{}-230 and complete floors26/11 unchanged. Other classes, signed positivity, actual source, selected arithmetic content, prescribed period, specific032/Read, DP44 and RH remain open.

Author /root/kummer\_extension, 23 November2026 UTC. Controlled pending proof; final author replay and independent whole reviews are separate records.

\subsection{1. Actual scope}\label{actual-scope}

Retain every actual J56 all-top object and normalization of K103, with H133, f168, s26, all active places typeII, global rank-two unipotence, A=I and nonzero scalar cross map. Apply this result only after the current K103 removal. Its exact residual conditions are \[
 v=q=26,\quad 2^{1116}<x=N_ED_\eta<2^{2108},\quad b=j=0,
 \qquad81\le k\le164.                                      \tag{FV1}
\] The original eta is unique for each support in this window; active labels have the checked one-label bound and every inactive place supplies its retained factor. Thus the full actual weighted product is bounded by \[
 (T-1)^{-168}<2T^{-168}.                                   \tag{FV2}
\] This uses K103's new support26 proof, not the old support25 uniqueness. No statement of constancy along an untruncated central orbit is required.

\textbf{Theorem.} The complete subset of FV1 with81≤k≤95 has absolute normalized contribution less than2\^{}-274. Add this one fresh charge to the current full ordered list, intersect the removed set with the current remainder, and retain every previous charge. Within this all-II scalar-cross branch the new necessary range is96≤k≤164. Other support26 branches are unchanged; no full support27 conclusion follows.

Use \[
 Q=1427908193729121522601,\quad U=77Q^2-14,\quad
 R=17424/133,\quad M^2=R/U,\quad
 D_0=7^{140}29^{144}5^{84}13^{84}.                           \tag{FV3}
\] The original covector has k nonzero S coordinates and the three fixed transverse coordinates. Its full energy is at most fM².

\subsection{2. One actual full-volume count}\label{one-actual-full-volume-count}

Fix a support, its unique eta, and one original common-C G pattern. Put \[
 L=JI_\Sigma D_\eta,\qquad z=(N_EL)^{1/(2f)},\qquad
 N_EJ=z^{2f}T^{52}/x\le Q^{334}.                            \tag{FV4}
\] The actual transverse dual Gram determinant is4U. Therefore the real covolume of the (k+3)-coordinate containing lattice is \[
             D_0^{k+3}z^{2(k+3)f}/(4U)^f.                 \tag{FV5}
\] K102 PC22 proves the safe covering coefficient c0 sqrt((k+3)f) z, c0=2250/7, by the original fractional-ideal bound and actual transverse metric. It does not declare a moving basis integral. The volume bound for the original energy ball is consequently \[
 {\cal V}_k(z)=
 {164\choose k}
 \frac{((22/7)f[M+c_0\sqrt{k+3}z]^2)^{(k+3)f}(4U)^f}
 {((k+3)f)!\,D_0^{k+3}z^{2(k+3)f}} .                     \tag{FV6}
\] This counts the full joint ball, including all original coordinate positions. No Gaussian, positive-norm phase or beta-pivot count is multiplied into it.

Retain the actual selected real trace and affine central fibre from K103 SC26.14--16: \[
 |c(h)|<2^{-929f/2}U^{-f/2}z^{-f},\qquad
 \#d\le V_F(H/2+750Uz^2)U^{-f}T^{52}x^{-1}.               \tag{FV7}
\] The central real shift, exact denominator ideal and full-real covolume remain the same. Eta,r,d reconstruct the full matrix. We conservatively retain the old factor2\^{}8 for the reconstruction pivot.

Together with Green, the z-dependent part of the complete count for fixed eta is, up to positive factors independent of z, \[
 F_k(z)=z^{(328-2k)f}
       (M+c_0\sqrt{k+3}z)^{2(k+3)f}
       (H/2+750Uz^2)^f .                                  \tag{FV8}
\] For all k≤164 it is nondecreasing. For k≤95 its leading exponent is strictly positive, which also gives the pattern saving below.

\subsection{3. Sum the actual G patterns with their losses}\label{sum-the-actual-g-patterns-with-their-losses}

At each of the334 unmixed places, write q\_G for its actual residue size. The whole-C ideal has local norm exponents2,1,0: one choice at exponent2, at most two canonical orientations at exponent1, and one grouped state at exponent0. Every actual term is covered by this pattern set. Extra matrix ranks are already inside the containing count and do not add pattern labels.

For fixed x let \[
 z_*=(Q^{334}T^{-52}x)^{1/(2f)}
\] be the all-rank-zero cap. If a local exponent decreases by h=0,1,2, its scale decreases by q\_G\^{}(-h/(2f)). From FV8, \[
 F_k(z)/F_k(z_*)\le(z/z_*)^{(328-2k)f}.
\] The full sum over patterns is at most \[
 F_k(z_*)\prod_G
 \left(1+2q_G^{-(164-k)}+q_G^{-2(164-k)}\right)
 =F_k(z_*)\prod_G(1+q_G^{-(164-k)})^2<2F_k(z_*).           \tag{FV9}
\] Indeed k≤95 gives164−k≥69, q\_G≥g0=22285932761, and2·668\textless g0\^{}69. For0\textless t and nt\textless1, the binomial series gives(1+t)\^{}n≤1/(1−nt); apply it with n668 and t=g0\^{}-69. There is no unchanged2\^{}668 pattern factor after this sum.

The estimate uses one fixed eta for the support; there is no separate eta or denominator choice per pattern. All affine central factors were first bounded as in FV7, so their decrease with z is in the correct direction for FV9.

\subsection{4. Coupled x endpoint and exact cell formula}\label{coupled-x-endpoint-and-exact-cell-formula}

The support count satisfies the strict exact bound \[
                      {168\choose26}<2^{101}.             \tag{FV10}
\] Thus the complete constant discrete factor is2\^{}112: support101, the pattern sum1, pivot8, one global Green correction1 and the single combined T-factor correction1.

The upper cap z\_*(x) grows as x\^{}(1/(2f)). In the full product the remaining x power is x\^{}-337/2. If x is multiplied by t≥1, FV8 increases by at most \[
 t^{164-k}\,t^{k+3}\,t=t^{168};
\] the row radius ratio is at most t\^{}(1/(2f)) and the central radius ratio is at most t\^{}(1/f). Hence the full count decreases by at least t\^{}-1/2. This proves the lower endpoint x=2\^{}1116 directly for the whole volume count. It includes the range formerly handled by the cone lower-p argument; there is no unbounded-scale gap and no endpoint union.

At this endpoint define c by \[
 c^f=(R/U)^fT^{52}Q^{-334}2^{-1116}.
\] The actual exact integers give \[
                 4124955<c<4124956,\qquad
 H/2+750R/4124955<68.                                      \tag{FV11}
\] Since FV8 increases with z, replace z\_* by M/sqrt(c0'), where c0'=4124955. It also lies below the nonzero-coordinate cap for the entire range.

Use the grouped directed logarithms \[
 L_T=1019333/1024,\quad L_Q=287843/4096,\quad
 L_U=601356/4096,\quad L_R=28810/4096.
\] The complete common exponent is \[
 C_{\rm vol}=112+2f-929f/2+479+167fL_R
 -335fL_U/2+8594L_T-335(167)L_Q-(337/2)1116.               \tag{FV12}
\] Every source, measure, support, finite weight, central fibre and selected-real factor is present. The full T exponent8594 is positive before its upper bound is substituted; positive R and negative U powers are already combined.

Put S=1024 and let r\_S(n) be the least integer with r\_S(n)\^{}2≥S²n. Set \[
 a_k=7r_S(4124955)+2250r_S(k+3),\quad d_k=(k+3)f .
\] All square roots are bounded above by exact integer comparisons. Let e\_k be the least strict binary upper exponent of \[
 W_k=
 \frac{{164\choose k}(22f a_k^2)^{d_k}}
 {(343S^2)^{d_k}d_k!\,D_0^{k+3}4124955^{167f}}.           \tag{FV13}
\] Then the complete cell exponent is C\_vol+e\_k. This is FV6 at the enlarged radius, not a new metric normalization.

The accompanying exact replay checks every one of the15 new cells81..95. Their largest upper exponent is \[
                    E_{95}=-1128627/4096<-275.            \tag{FV14}
\] It also records k96 as a failed sufficient bound, not a lower estimate for any actual contribution. The exact dyadic union of the15 complete upper bounds is below2\^{}-274; in particular the last cell is below2\^{}-275 and all preceding cells are below2\^{}-2140. There is one new charge274. The earlier K103 charge293 and all other errors remain.

\subsection{5. Source, chronology and remaining scope}\label{source-chronology-and-remaining-scope}

K103 SC26.4--7 supplies the actual residual and complete positive compression; SC26.12--17 supplies the true row/central ideals and selected trace; SC26.20 supplies the combined finite correction. K102 PC22 supplies the actual fractional-lattice full-volume count. FV8--9 and the whole-volume coupled endpoint are proved here and have their complete exact arithmetic replay. No positive-definite assumption at a finite residue field is used.

The default canonical predecessor is checks/k103\_joint\_remainder\_error.json. The script permits a later full chronological prefix. If a pre-integration replay is run against K103's frozen pending certificate, that event is explicitly historical and a separate unchanged canonical-prefix replay is required. Complete all-top and uniform lists are retained, with one274 append. Floors26/11 stay unchanged.

All original fields, involution, Hermitian form, polarization, theta line, finite projectors, selected real module, central phases, class masses and Haar factors remain fixed. The carried algebraic-geometric anchors are exactly those listed in K103 section10 (SGA3, SGA6, FGA232/236, SGA7 and cote4/54), in their previously checked scopes. No selected ordinary-derivative/content, prescribed period, effective source-family, Read/STiesR, DP44 or RH conclusion follows.

\subsection{Canonical algebraic-geometric context}\label{canonical-algebraic-geometric-context-21}

The reductions are the actual integral Hermitian modules at unmixed finite points; content and coordinate ideals remain fractional invertible modules over Spec(O\_E) and Spec(O\_F), with the same paired CM action, polarization, theta line, Picard scheme and biextension. Carried exact anchors: SGA3 I4.7.3/4.7.3.1; SGA6 VI Theorem3.3; FGA232 Theorem3.1 pp148--149; FGA236 Corollaries1.7/2.3 pp229/232; SGA7 I VII2.7/2.9 pp169--174 and VIII7.4(b) pp310--311; cote4 scan003 sections8.1.1--8.1.2 and scans005--006 sections8.4.1--8.4.3; cote54 scans005--007. No new lambda-ring/K-theory, l-adic weight or CM-motive comparison is asserted. DP\_MULT Q3e, THETA3 G1/A7, CX\_GRH CX-ALG-024 and GCT\_PLETH T-052 retain their stated global gates.

\clearpage

\section{AM-R-089 --- weighted eta and content sums complete rank-one support26}\label{am-r-089-weighted-eta-and-content-sums-complete-rank-one-support26}

\textbf{Canonical checked integration --- 2026-11-23T04:03:03.951974+00:00.} Complete root and independent proof/script/interface reviews and separate exact replays passed. The actual canonical predecessor has106/31 entries; this integration gives107/31, complete floors26/11. Frozen author history is preserved. AM-R-089 WN1--18 removes the complete actual fixed-S rank-one-unipotent support26 remainder: all334Gidentity and intrinsic double-pole sizes23..26. Its universal reciprocal-norm weighted eta/content lemma uses determinant-one rectangular rescaling and full dyadic shell packing, including every unit and ideal class. The original selected trace retains its actual eta norm before summation; both genuine affine coordinates and the whole actual support/marked/type/pair sum remain. The complete rational upper exponent is-146697349/4096\textless-35000, giving one35000 charge. Together with already charged R088 this completes the whole fixed-S rank-one-unipotent support26 structural class. Actual K104106/31 becomes107/31, every old overlapping error retained, both entire sums below2\^{}-230 and complete floors26/11 unchanged. Other support26 classes, signed positivity, actual source, selected arithmetic content, prescribed period, specific032/Read, DP44 and RH remain open.

Author: GPT-6 Astra Pro; reasoning effort not externally exposed; ALP\_MACHINE /root. 23 November2026 UTC. Controlled full proof. Whole saved-file reading, exact execution and independent review are separate records before canonical promotion.

\subsection{1. Actual sector and theorem}\label{actual-sector-and-theorem-9}

Keep the unchanged J56 all-top kernel, its H=133 cutoff, f=a=168, original fields E/F, CM partners, source and target pencils, actual labels and finite projectors, selected Hilbert metric, packet point, central phases, class masses and Haar measures. The fixed-S globally rank-one-unipotent class has h=I+N, rank\_E N=1 and N squared=0. R088 GZ1--11 has already removed its entire integral-witness branch, every reduced double-pole size d\textless=22 and every nonidentity unmixed-G pattern at support s=26. Its only possible remainder has

\[
 s=26,\quad23\le d\le26,\quad J=J_{\max},\quad
 N_EJ=Q^{334}=M^2,\quad M=Q^{167}.                 \tag{WN1}
\]

Here d is the intrinsic number of actual double poles of the original negative entry eta. Every such place is type II. Global rank-one unipotence is an explicit hypothesis, not a consequence of a finite positive form.

\textbf{Theorem.} The entire actual WN1 remainder has absolute normalized contribution less than2\^{}-35000. Together with the already charged R088 removal, this completes the whole fixed-S globally rank-one-unipotent support26 class. Append35000 once to the actual chronological all-top list, retaining every older overlapping charge. This does not exhaust the other support26 structural classes or prove the signed sufficient inequality.

The new input is a reciprocal-norm-weighted version of the LIFT078 content sum. The actual selected trace is kept as a function of eta throughout that sum.

\subsection{2. A weighted fractional-ideal lemma}\label{a-weighted-fractional-ideal-lemma}

Let F be a totally real number field of degree f, Lambda a nonzero fractional ideal with N\_F Lambda=z\^{}f, and R\textgreater0. Write

\[
 Y=R^f/z^f,\quad J_Y=\lfloor\log_2Y\rfloor,\quad
 A_t=(t)\Lambda^{-1}
 \quad(0\ne t\in\Lambda).                            \tag{WN2}
\]

Suppose c(t)\textless=d\_\{3,F\}(A\_t), where the latter counts all ordered pairs of nonzero integral ideals (K,L) with KL dividing A\_t. For every delta\textgreater0 and Y\textgreater=1,

\[
 \sum_{\substack{0\ne t\in\Lambda\\|\tau t|\le R\ \forall\tau}}
 \frac{c(t)}{|N_{F/\mathbf Q}t|}
 \le6^f z^{-f}{J_Y+f\choose f}
       Y^\delta\zeta_F(1+\delta)^2.                 \tag{WN3}
\]

For Y\textless1 the sum is empty. The estimate has no regulator, principality, ideal-class, unit-count or lattice-shape hypothesis.

To prove it, first fix K,L and put n=N\_F(KL). Nonzero t contributing to this pair belongs to Lambda KL and satisfies

\[
 |N_Ft|\ge z^f n,\qquad n\le Y.                    \tag{WN4}
\]

Partition its absolute embedding sizes uniquely into dyadic shells

\[
 R2^{-j_\tau-1}<|\tau t|\le R2^{-j_\tau},
 \quad j_\tau\in\mathbf Z_{\ge0}.
\]

For a profile j set a\_tau=R2\^{}(-j\_tau) and B\_j=product\_tau a\_tau. A nonempty shell must have B\_j\textgreater=z\^{}f n.~In particular sum\_tau j\_tau\textless=J\_Y, so there are at most binom(J\_Y+f,f) possible profiles. This is the finite weak-composition identity; every sign is still present within each absolute-value shell.

The entire shell lies in the centered box with sides {[}-a\_tau,a\_tau{]}. Apply the diagonal linear map with entries B\_j\^{}(1/f)/a\_tau. Its determinant is one, and it sends this box into the Euclidean ball of radius sqrt(f) B\_j\^{}(1/f). For every nonzero difference in the transformed lattice, the product of its coordinates is unchanged. By WN4 applied to that original lattice difference and arithmetic-geometric mean, its squared Euclidean norm is at least f z\^{}2 n\^{}(2/f). Disjoint open balls of half this minimum around all points of the transformed box therefore give

\[
 \#\{t\in\Lambda KL:|\tau t|\le a_\tau\}
 \le\left(1+\frac{2B_j^{1/f}}{zn^{1/f}}\right)^f
 \le\frac{3^fB_j}{z^f n}.                           \tag{WN5}
\]

The last inequality is used only for a nonempty nonzero shell, for which B\_j\textgreater=z\^{}f n.~Counting the origin in the containing box is harmless; no origin is assigned a reciprocal norm. No algebraicity of the diagonal map is needed: the product minimum is preserved for every nonzero lattice difference. This is why arbitrarily uneven embeddings and all units are covered uniformly.

Within the shell, \textbar N\_Ft\textbar\textgreater2\^{}-f B\_j. Multiplying this reciprocal-norm bound by WN5 gives at most6\^{}f/(z\^{}f n). The centered box already counts all signs, so there is no extra2\^{}f sign factor. Summing profiles and then the ordered ideal pairs gives

\[
 \sum_t\frac{d_{3,F}(A_t)}{|N_Ft|}
 \le6^f z^{-f}{J_Y+f\choose f}
       \sum_{N_F(KL)\le Y}\frac1{N_F(KL)}.          \tag{WN6}
\]

For n\textless=Y, n\^{}-1\textless=Y\^{}delta n\^{}(-1-delta). The full absolutely convergent ordered-pair sum is zeta\_F(1+delta)\^{}2. This proves WN3. The sum rearrangement is nonnegative, and the cutoff n\textless=Y is forced by a nonzero t; it is not inferred by counting zero in infinitely many ideals. For Y\textless1, the nonzero ideal-norm minimum already rules out t.

\subsection{3. Actual content, affine coordinates and selected trace}\label{actual-content-affine-coordinates-and-selected-trace}

Use the unchanged reduced normal form of R088 GZ2 and R080 AQ1--9:

\[
 N=x\bar x^{\mathsf t}J_0/\eta,\quad x=(vq,r,\eta),
 \quad Vq\bar q+r\bar r=\eta\bar\eta,
 \quad\bar\eta=-\eta\ne0,\quad q\ne0,
\] \[
 v=(2,95192390527,383092671043),\quad
 V=155821585821105090045582,\quad
 |\eta_\tau|\le R=8844/133.                           \tag{WN7}
\]

For a fixed original support, marked set, actual labelled pairs and the single WN1 pattern,

\[
 \mathfrak l=J I_\Sigma I_D=\Lambda_F\mathcal O_E,
 \quad z_d^f=N_F\Lambda_F=M/T^{26+d},
 \quad Y_d=R^fT^{26+d}/M.                            \tag{WN8}
\]

The element t=-i eta belongs to Lambda\_F because i is the original integral unit and E=F(i). The exact content b=(q,r,eta) gives C=b l\^{}-1 and satisfies LIFT078 AV3. The actual valuation count AV5 consequently gives c(t)\textless=d3\_F((t)Lambda\_F\^{}-1). This is exactly WN3's ideal and divisor count. Also

\[
 |N_E\eta|^{1/2}=|N_F(-i\eta)|.                     \tag{WN9}
\]

No ideal generator, phase, class or extra unit multiplicity is introduced.

The genuine selected trace follows from DO088 WM4--6 with the DO092 SN16--18 radial multiplier retained before its final ideal-minimum substitution:

\[
 |c(h)|<2\,2^{-931f/2-630}|N_F(-i\eta)|^{-1}.       \tag{WN10}
\]

Indeed the actual rank-one transport has sinh(r(h\_p))=\textbar eta\_p\textbar{} and eta\_p=eta times the squared absolute pairing with the transported negative line. The same packet point gives product\_tau exp(2r\_p)\textless2. The integrated radial multiplier has eta\_a\^{}f\textless2\^{}-630 at a=f168. Both compact factors, the actual selected module and every word parameter remain those of the checked source. Formula WN10 is a bound on that actual normalized trace, including arbitrarily small individual conjugates, and is not an operator-norm replacement. We do not substitute \textbar N\_F(-i eta)\textbar\textgreater=z\_d\^{}f before summation.

For fixed eta,b and original labels, R080 AQ5--9 gives the two genuine affine-coordinate counts

\[
 C_d\le(1+2\sqrt{y_{c,d}})^{2f}
           (1+2\sqrt{y_{c,d}/V})^{2f},\qquad
 y_{c,d}=R(T^{26-d}/M)^{1/f}.                        \tag{WN11}
\]

Their exact difference ideal follows from both original paired pole congruences and the primitive integer vector v. The counts are uniform in eta, b and every affine shift, so they may be taken outside WN3. The full column determines N. No relative-norm fibre, generator count or lattice-isometry multiplicity is added.

For the actual field, LIFT078 AV16--17 gives

\[
 \zeta_F(1001/1000)<3003\,5^{167},\qquad
 3^{168}(3003\,5^{167})^2<2^{1065}.
\]

Set delta=1/1000. Since 6\^{}f=2\^{}f3\^{}f, WN3 becomes

\[
 \sum_\eta\frac{\#\{\text{actual content }b\}}{|N_F(-i\eta)|}
 <2^{1233}z_d^{-f}{J_{Y_d}+f\choose f}Y_d^{1/1000}.
                                                               \tag{WN12}
\]

This is the combined weighted eta/content sum. It replaces that combined count and trace dependence, and is not multiplied by either the old unweighted eta count or an additional z\_d\^{}-f trace bound.

\subsection{4. One complete four-cell bound}\label{one-complete-four-cell-bound}

The exact host integers satisfy, simultaneously for23\textless=d\textless=26,

\[
 1<Y_d\le Y_{26}<2^{42000},\qquad
 {42168\choose168}<2^{1600},\qquad
 Y_d^{1/1000}<2^{42}.                                \tag{WN13}
\]

Thus the entire profile binomial in WN12 is below2\^{}1600. The same integers give y\_(c,23)\textless2\^{}-40, and y\_(c,d) decreases with d.~Since V\textgreater1,

\[
 C_d<(1+2^{-19})^{4f}<2.                              \tag{WN14}
\]

The last inequality follows either by its exact integer power or from the binomial comparison (1+u)\^{}m\textless=1/(1-mu) for mu\textless1, with m=672 and u=2\^{}-19. All four actual coordinate cells are included.

At J\_max the full Green bound2(N\_EJ)\^{}(a/2)M\^{}-a equals2. There is exactly one G pattern left; no G-pattern factor2 is charged. K077 AP18 retains the entire support, marked-set, type and actual ordered-pair sum

\[
 {168\choose26}{26\choose d}(T+1)^{26-d}.             \tag{WN15}
\]

It is applied after the uniform weighted column bound for each original labelled pair. The opposite CM component remains paired data. For the four possible d,

\[
 (T+1)^{26-d}z_d^{-f}
 =\frac{T^{52}}M(1+T^{-1})^{26-d}
 <\frac{2T^{52}}M,
\] \[
 {168\choose26}<2^{101},\qquad
 \sum_{d=23}^{26}{26\choose d}=2952<2^{12}.           \tag{WN16}
\]

The ratio comparison follows from (T+1)\^{}3\textless2T\^{}3. This combines the four complete intrinsic marked-size cells; no later support or type union remains.

Multiplying WN10, WN12--16 and the actual Green bound gives the complete union bound

\[
 2^{4+1233+101+12+1600+42-931f/2-630}\,T^{52}Q^{-167}.
                                                               \tag{WN17}
\]

The exponent4 consists exactly of the Green2, near-point2, coordinate-product2 and (1+1/T) ratio2. The weighted-content constant1233, profile1600, support101, marked-set12 and Rankin42 allowances each occur once. All remaining original mixed, central phase, compact, class and probability-Haar factors retain their already normalized modulus bounds. No cancellation of distinct terms is used.

With the unchanged directed host bounds

\[
 \log_2T<1019333/1024,\qquad
 \log_2Q>287843/4096,
\]

the full logarithm in WN17 is strictly less than

\[
 4+1233+101+12+1600+42-931\cdot84-630
 +52\frac{1019333}{1024}-167\frac{287843}{4096}
 =-\frac{146697349}{4096}<-35000.                    \tag{WN18}
\]

This proves the stated single charge35000. The large spare margin is deliberate; the new argument is the weighted sum, not optimization of integer ceilings. R088's previous523 charge and every older overlap remain in the full accounting.

\subsection{5. Chronology, sources and geometric scope}\label{chronology-sources-and-geometric-scope}

The companion checks the literal host, full radial and packet constants, the actual weighted-content constant, every one of the four Y and coordinate inequalities, the full profile binomial, complete support/marked/ratio factors, exact WN18, and both complete chronological arrays. It pins the final proof and source hashes. These are exact arithmetic diagnostics of the universal written argument, not an enumeration of arithmetic arrows. The script accepts --prior for an unchanged replay against the actual canonical prefix and appends35000 exactly once. Complete support floors remain26/11 pending an exhaustive assembly of the other classes.

Direct interfaces are R088 GZ1--11 (the exact residual); R080 AQ1--9 and K077 AP6--11/18 (original normal form, content, both affine coordinates and the complete pair sum); DO088 WM3--8 (genuine norm-sensitive selected trace and packet transport); DO092 SN16--19 (the retained630-bit angular factor); and LIFT078 AV1--9/16--18 (actual eta/content transport and the accepted convergent zeta constant). WN3--6 prove the new weighted lemma for arbitrary real fractional ideals rather than citing a regulator or counting theorem with unverified constants.

The rescaling is a real counting device on the embeddings of the same fractional ideal. It does not change the CM field, coefficient field, integral Hom module, polarization, Picard class, theta line, biextension or source data. The ideal and divisor sums live on Spec(O\_F) and their original extensions to Spec(O\_E). Carried exact anchors are SGA3 I4.7.3/4.7.3.1; SGA6 VI Theorem3.3; FGA232 Theorem3.1 pp148--149; FGA236 Corollaries1.7/2.3 pp229/232; SGA7 I VII2.7/2.9 pp169--174 and VIII7.4(b) pp310--311; cote4 scan003 sections8.1.1--8.1.2 and scans005--006 sections8.4.1--8.4.3; cote54 scans005--007. These are inherited placements, not new lambda-ring, K-theory, l-adic-weight or CM-motive comparisons.

DP\_MULT Q3e and THETA3 G1/A7 obtain precisely this completed rank-one support26 class in their finite sufficient gate. The actual admissible mixed-CM source, selected arithmetic-image content, prescribed period and specific032 pairing/value, all actual DP44 terms, CX\_GRH CX-ALG-024 and the global GCT\_PLETH T-052 comparison retain their open gates. No RH conclusion, Lean build or archive audit is claimed.

\subsection{Canonical algebraic-geometric context}\label{canonical-algebraic-geometric-context-22}

The reductions are the actual integral Hermitian modules at unmixed finite points; content and coordinate ideals remain fractional invertible modules over Spec(O\_E) and Spec(O\_F), with the same paired CM action, polarization, theta line, Picard scheme and biextension. Carried exact anchors: SGA3 I4.7.3/4.7.3.1; SGA6 VI Theorem3.3; FGA232 Theorem3.1 pp148--149; FGA236 Corollaries1.7/2.3 pp229/232; SGA7 I VII2.7/2.9 pp169--174 and VIII7.4(b) pp310--311; cote4 scan003 sections8.1.1--8.1.2 and scans005--006 sections8.4.1--8.4.3; cote54 scans005--007. No new lambda-ring/K-theory, l-adic weight or CM-motive comparison is asserted. DP\_MULT Q3e, THETA3 G1/A7, CX\_GRH CX-ALG-024 and GCT\_PLETH T-052 retain their stated global gates.

\clearpage

\section{AM-DO-110 --- the complete non-all-II unipotent scalar-cross sector at support26}\label{am-do-110-the-complete-non-all-ii-unipotent-scalar-cross-sector-at-support26}

\textbf{Canonical checked integration --- 2026-11-23T04:10:31.703015+00:00.} Complete root and independent proof/script/interface reviews and separate exact replays passed. The actual canonical predecessor has107/31 entries; this integration gives108/31, complete floors26/11. Frozen author history is preserved. AM-DO-110 NT26.1--31 removes the complete actual non-all-II globally rank-two-unipotent scalar-cross support26 sector. The support-independent full-positive normal form, original paired cross/full-entry poles, 24-place eta uniqueness, actual one-label cancellation including central-only poles, full affine central lattice, common-G ideal and exact joint count are retained. Seven exhaustive branches and all1148 cells give one561 charge. Actual R089107/31 becomes108/31, every overlap retained, both complete sums below2\^{}-230 and full floors26/11 unchanged. The remaining all-II unipotent26 class is the K103/K104 dense window; other scalar-cross classes, signed positivity, actual source, arithmetic content, prescribed period, specific032/Read, DP44 and RH remain open.

Author: /root/definite\_oscillator, ALP\_MACHINE. 23 November 2026 UTC.

Controlled pending theorem. The whole saved-proof reading, complete executable reading, exact replay, independent review and canonical integration are separate recorded events. Earlier frozen files remain unchanged.

\subsection{1. Actual scope and theorem}\label{actual-scope-and-theorem-4}

Retain the actual J56 all-top normalized kernel, H=133, f=168, m=164, its original E/F and integral lattice, all334 unmixed G reductions and the two mixed factors, the actual source and target pencils at both CM partners, all admitted labels, the selected real module and both compact factors, the near-point normalization, central phases, Haar measures and class masses. Work in the current finite-unipotent G remainder. Assume global rank-two unipotence explicitly:

\[
 h=I+N,\quad h_{SS}=I,\quad C=h_{W_0S}\ne0,\quad
 \operatorname{rank}_E N=2,\quad s=26.                    \tag{NT26.1}
\]

The already charged DO099(JR1--24) interval removes transverse C and transverse inverse C through support70. Thus C lies in the original packet plane P=Eu⊕Ee\_n.~No support25 four-compression removal is applied at support26.

Let i,u,t count actual active types I,II,III, so i+u+t=26. This theorem concerns every profile with u≤25. The all-II profile is excluded.

\textbf{Theorem.} The entire stated non-all-II sector has absolute contribution less than

\[
                              2^{-561}.                    \tag{NT26.2}
\]

This is one fresh all-top charge, intersected with the actual current remainder. It does not assert full support26 exhaustion or change the complete floors26/11. In particular, the all-II remainder of K103 remains separate.

Put

\[
 Q=1427908193729121522601,\quad U=77Q^2-14,\quad
 R=17424/133,\quad M^2=R/U,\quad
 D_0=7^{140}29^{144}5^{84}13^{84},\quad D_F=2^{168}D_0.
                                                               \tag{NT26.3}
\]

Use the exact original directed bounds

\[
 995<\log_2T<L_T=1019333/1024,\quad
 \log_2Q>L_Q=287843/4096,\quad
 \log_2U>L_U=601356/4096.
                                                               \tag{NT26.4}
\]

For the bounded-window count use log2R\textless L\_R=28810/4096. Where the earlier CF count or ideal tail uses113/16 instead, retain that larger upper bound with its complete positive R power.

\subsection{2. Support-independent full positive compression and actual pole sets}\label{support-independent-full-positive-compression-and-actual-pole-sets}

The support-free argument K103(SC26.5--7) is reproduced here to specify the interface. Rank-two nilpotence gives N³=0. If N²=0, unitarity gives N†=−N and N†N=0, so im N is a two-dimensional totally isotropic space, impossible at a signature-one embedding. Thus N has one length-three Jordan block and

\[
 \ell=\operatorname{im}N^2=\operatorname{im}N\cap\ker N
 \text{ is a line},\qquad N+N^\dagger=N^2.                \tag{NT26.5}
\]

The self-adjoint rank-one expression for N² and N\_SS=0 force its S projection to vanish. Consequently N²S=0 and N(S)=C(S)=ell. Since C is plane-valued, ell⊂P. On the fixed positive complement Z=P⊥=S⊕u⊥, self-adjointness gives N²Z=0; hence NZ⊂ell and h\_ZZ=I. This uses no finite positive metric, integral moving frame or additional support cut.

The original normal form DO097(UC3--13), DO098(CF4--7), K103(SC26.7) is therefore valid:

\[
 \ell=(1,\eta),\quad \eta\bar\eta=U,\quad C=\ell r,
 \qquad
 N=\begin{pmatrix}
 0&-a\ell^*J_P\\ \ell r&d\ell\ell^*J_P
 \end{pmatrix},
 \quad a=J_Z^{-1}r^*,\quad J_P=\operatorname{diag}(U,-1),
\] \[
 d+\bar d=-rJ_Z^{-1}r^*,\qquad
 \alpha=U rJ_Z^{-1}r^*\le R,\qquad N^2Z=0.                \tag{NT26.6}
\]

Eta, the167 original covector values and d determine h injectively. The original central entry is w=−Ud, with Re w=alpha/2 and \textbar Im w\textbar\textless H/2. Let k∈\{1,\ldots,164\} count the nonzero original S values. The first nonzero original S pivot exists; the old harmless164\textless2\^{}8 pivot factor is retained throughout.

Every original S input and each of the three fixed integral transverse inputs belongs separately to every source and target pencil at both CM components. At type III, Ny∈M\_d and (h\^{}-1−I)y=−Ny∈M\_c, since N²y=0. Thus Ny lies in M\_c∩M\_d⊂Lambda. At type I, h preserves the common intersection module containing y and contained in Lambda. Primitive integer Bezout coefficients of u then extract both r(y) and eta r(y) in the original integral coordinates. Thus all167 cross values are integral at I and III.

At type II the simple cross pole is possible only if the saturated line ell is the actual pole line. The rank-one leading double-pole term of the literal shear A\_c K A\_c\^{}-1 proves the same necessity for any full matrix double pole. This is the original DO103(LD6--11) argument, valid for the full Z inputs just established. At the selected component its condition is

\[
 \eta\equiv a_0=Q(-7+2\sqrt{-7})\pmod{\mathfrak p};       \tag{NT26.7}
\]

at the partner it is eta≡U/bar(a0). The norm identity makes these one eligibility condition. A zero or pole of eta cannot give this saturated pole line, whose two displayed coordinates are units.

Let E0 be the set of eligible active type-II places, and write p=\textbar E0\textbar≤u≤25. Then cross values have simple poles only on E0; full entries have pole at most two on E0 and at most one on Sigma\E0. A noneligible type-II place is therefore treated as a simple full-entry pole even though its type remains II. This distinction is essential when p\textless u.

The exact complete ordered-pair/absolute-incidence factors are K071(NB7--8): at one active place they are at most T−1 for type I and1 for each of II and III. Both CM partners form one paired choice. These weighted sums are applied only after the actual-lattice matrix count has been bounded uniformly for fixed pairs.

\subsection{3. The p≤23 containing sector, with all type and padding factors}\label{the-p23-containing-sector-with-all-type-and-padding-factors}

Pad E0 to a set D of23 active places. There are at most binom(26,23)=2600\textless2\^{}12 padded sets. Keep all3\^{}26\textless2\^{}52 type assignments and the conservative full pair weight T\^{}26. Hence the complete extra factor is

\[
                              2^{64}T^{26}.                 \tag{NT26.8}
\]

Let J be the canonical common-C G ideal, with at most4\^{}334 patterns and

\[
 |a_G(h)|\le2Q^{-335\cdot167}(N_EJ)^{335/2},\quad
 N_EJ\le Q^{334},\quad N_EJ\le Q^{167}N_FJ_0.
                                                               \tag{NT26.9}
\]

All cross values and the original nonzero beta pivot lie in L=JI\_D, with N\_EL=z\^{}(2f), N\_EI\_D=T\^{}-46. The full-entry pole statement above gives the fixed-C central differences in the single affine coset of iJ0 I\_Sigma,F I\_D,F. Its real part is fixed; the shift, including any dyadic alpha/2, is retained. Since s+\textbar D\textbar=49 and750\^{}168 Q\^{}167\textless T\^{}48, the two complete DO102(NT6--14) central counts apply:

\[
 \#w<2^{479}Q^{167}T^{s-|D|}z^{-2f},\qquad
 \#w\le134^fT^{s+|D|}.                                   \tag{NT26.10}
\]

They are alternatives for one fibre, never multiplied. All old real, transverse, relative-norm, support, pattern and pivot factors are unchanged; their fixed multiplier is2\^{}845=2\^{}(1+668+168+8). Reconstruction is injective, so no smaller isometry fibre is inferred.

The complete CF endpoints, recomputed by the companion script, are: for1≤k≤64, the improved central count and trace exponent1 give every cell below2\^{}-3927 at s=v=24; for65≤k≤163 the other central count and exponent1 give below2\^{}-10020; k164 with exponent7/8 gives below2\^{}-2880. These are complete arithmetic envelopes, not claims that the present matrices belong to an earlier removed sector.

For the first range the combined Green/central T power is334v+s. At v23,s26 it is7708 instead of8040; after NT26.8 the net saving is306 powers of T. Removing the old cap30 costs at most2\^{}136028, by the exact DO102(NT18) bound30\^{}54411\textless2\^{}272056. Thus each of the first64 cells is below

\[
 2^{-3927+136028+64-306\cdot995}=2^{-172305}.               \tag{NT26.11}
\]

For the other ranges the power336v+s changes from8088 to7754; after NT26.8 the saving is308 powers. Their complete bounds are respectively2\^{}-316416 and2\^{}-309276. The full164-cell p≤23 union is therefore less than2\^{}-172297. This includes every profile and every possible padded pole set, without an eta-denominator condition.

\subsection{4. Remaining p≥24: denominator tail, uniqueness and complete row formula}\label{remaining-p24-denominator-tail-uniqueness-and-complete-row-formula}

Now p is24 or25. In particular i≤2. Set x=N\_E D\_eta, where D\_eta is the exact integral denominator ideal of eta. For x≥2\^{}4096, enlarge all cross and full-entry denominators to the26-place all-II budget. K103(SC26.21--22) is the complete infinite ideal tail, with the support-independent bound of fewer than2\^{}39 norm-U scalars per exact denominator ideal and zeta\_E(2)\textless2\^{}336. Every one of its164 cells is below2\^{}-12189. Here keep all3\^{}26\textless2\^{}52 types and the actual pair bound T². Hence every new tail cell is below

\[
             2^{-12189+52+2L_T}<2^{-10146},                \tag{NT26.12}
\]

and the whole tail is less than2\^{}-10138. There is no eta uniqueness or eligible-set enumeration in this branch.

For the rest assume1≤x\textless2\^{}4096. The original denominator-cleared scalar eta−a0 is nonzero, and the exact7-adic valuation argument DO103(LD12--15)/K099(GE4) gives

\[
 7^{84}T^q<[4(U+14)]^{168}x,                               \tag{NT26.13}
\]

where q counts eligibility among all168 selected places. The exact actual comparisons retained from K103 imply q≤28; q≥25 implies x\textgreater2\^{}120, and q≥26 implies x\textgreater2\^{}1116. No support25 bound q≤26 is assumed.

The new24-place uniqueness comparison is

\[
                        (4U)^{168}2^{8192}<T^{48}.          \tag{NT26.14}
\]

Indeed two norm-U scalars eligible on a fixed24-place set agree at both partners there. Their difference times their denominator ideals, if nonzero, is integral of norm at leastT\^{}48 and less than the left side of NT26.14. They coincide. Thus eta is unique after choosing E0 when p24 or25; neither its denominator ideal nor its additional eligible places are enumerated in this bounded window.

Fix eta and one common G pattern. All167 row values lie in

\[
 L=JI_p\cap\eta^{-1}JI_p=JI_pD_\eta,
 \quad z^{2f}=N_EL,
 \quad N_EJ=z^{2f}T^{2p}/x.                               \tag{NT26.15}
\]

The equality is valuationwise, including eta denominators at G and inactive T. The original four-coordinate double Bezout extraction and the mixed entries give the fixed-C central difference ideal

\[
 \Gamma_d=iJ_0B_\eta I_{\Sigma,F}I_{p,F},\qquad
 D_\eta\bar D_\eta=B_\eta\mathcal O_E,\quad N_FB_\eta=x.
                                                               \tag{NT26.16}
\]

This is one affine coset, not an assertion d∈Gamma\_d.~Its full-real covolume and its count are

\[
 \operatorname{covol}\Gamma_d=\sqrt{D_F}N_FJ_0\,xT^{-(26+p)},
\] \[
 \#d\le V_F\!\left(H/2+750Uz^2T^{(p-26)/f}\right)
                   U^{-f}T^{26+p}x^{-1},\quad
 V_F(C)=\frac{((22/7)fC^2)^{f/2}}{(f/2)!\sqrt{D_F}}.
                                                               \tag{NT26.17}
\]

Here N\_FJ0≤N\_EJ controls the covering radius; only N\_FJ0≥1 is dropped from the covolume. The half-trace normalization is not inserted into this real lattice.

A chosen extra inactive paired-C-zero place multiplies the row ideal by an E-ideal of normT². Write e=1 then and e=0 otherwise. Exactly

\[
 N_EJ=z^{2f}T^{2p}x^{-1}T^{-2e},\quad kz^2\le M^2,
 \quad z^2\le(Q^{334}T^{-2(p-e)}x)^{1/f}.                 \tag{NT26.18}
\]

The central ideal stays NT26.16: its covering radius uses z²T\^{}((p−26−2e)/f), while the trace uses the improved row scale z.

The original full joint-energy row bound, the actual trace and reconstruction are K103(SC26.16,23):

\[
 \#r\le K_k A_3(4U)^f\mathcal B_k
       (M+2^{-10}z)^{2(k+3)f}z^{-2(k+3)f},\quad
 |c(h)|<2^{-929f/2}U^{-f/2}z^{-f},                          \tag{NT26.19}
\]

where K\_k=binom(164,k)47\^{}k2\^{}(2fk)(f!)\^{}k(87/(32fk))\^{}(fk), A3=528\^{}(3f)/((3f)!D0³)\textless2\^{}-4040 and B\_k=(fk)!(3f)!/((k+3)f)!. The actual covector Gram determinant4U, the full selected trace and both compact factors remain. No scalar Bergman coefficient is substituted.

\subsection{5. Inactive factors and the common numerical envelope}\label{inactive-factors-and-the-common-numerical-envelope}

At each noneligible inactive place without paired-C-zero reduction, the support-independent DO107(IP4--7) pencil calculation gives \textbar rho\textbar=1/(T−1). The unique incidence is the admitted original Lambda. Paired inverse adjointness supplies one paired condition, not a second coefficient. Let b count all inactive paired-C-zero places. At least168−q−(26−p) inactive places are noneligible, so at least

\[
                         g=168-q-(26-p)-b                 \tag{NT26.20}
\]

such factors may be retained when this lower bound is used. Additional factors are bounded by1. Every paired-zero place contributes its extra row ideal even if eta has local zeros or poles. A nonzero original S value gives

\[
                  xT^{2b-2p}\le |N_Er(y)|\le(R/U)^f.     \tag{NT26.21}
\]

For p25, x\textgreater2\^{}120 and (R/U)\^{}f T\^{}22\textless2\^{}120 imply b≤13. For p24 the stronger exact comparison (R/U)\^{}f T\^{}22\textless1 and x≥1 imply b≤12. One chosen paired-zero place costs at most142\textless2\^{}8.

At an eligible II place the fixed-eta complete-pencil calculation DO105(EL9--16) permits at most one actual diagonal label, even if the only full pole is central. Its weight is1/(T−1). Other II places and III places have complete label-weight sum at most1, and type I contributes at mostT. Combine the negative factors before bounding:

\[
                    (T-1)^{-(p+g)}<2T^{-(p+g)}.           \tag{NT26.22}
\]

The exact2(T−1)\^{}168\textgreater T\^{}168 suffices for every p+g≤168. The fixed factor is therefore2\^{}846: the original2\^{}845 and this single2.

For a lower endpoint x≥2\^{}a, let c(p,a,e)\^{}f=(R/U)\^{}f T\^{}(2(p−e))Q\^{}-3342\^{}-a. Choose a certified integer lower cap c0≤c and put A=max(k,c0); if the global cap is discarded, put A=k. This enlarges the allowed row endpoint. If the central radius there is at most68, the complete joint-volume exponent is bounded by

\[
 \mathcal C(p,a,g,e,d,i)-4040+\mathcal R_k(A)+\mathcal I_k(A),
                                                               \tag{NT26.23}
\] \[
 \mathcal C=846+2f-929f/2+479+167fL_R-(335f/2)L_U
 +(335p+26-g-335e+i)L_T-335\cdot167L_Q-337a/2+d.
\]

Here d is the paid type/set/extra-place binary factor, not the central coordinate. R\_k(A) is the least strict binary ceiling of

\[
 \frac{\binom{164}{k}47^k2^{2fk}(f!)^k87^{fk}(fk)!(3f)!}
 {32^{fk}f^{fk}k^{fk}A^{(164-k)f}((k+3)f)!},              \tag{NT26.24}
\]

and I\_k(A) is the least strict ceiling of (1+1/(1024 floor(sqrt(A))))\^{}(2(k+3)f). The exact positive T exponent is checked before using L\_T. When comparing two envelopes with a negative T-power difference, use the lower bound995 instead.

The optimized complete dependence is proportional to

\[
 x^{-337/2}z^{(328-2k)f}(M+2^{-10}z)^{2(k+3)f}
 V_F\!\left(H/2+750Uz^2T^{(p-26-2e)/f}\right).           \tag{NT26.25}
\]

It is increasing in z. On the global-cap branch z∝x\^{}(1/(2f)), its upper logarithmic derivative in x is

\[
 -337/2+(164-k)+(k+3)+1=-1/2;
\]

on the energy-cap branch it decreases more strongly. Thus the whole expression decreases after optimizing z, and can be evaluated at x=2\^{}a. The small central-radius bounds used below are endpoint bounds after this optimization; they are not asserted uniformly in x.

The K103 cone/complex-volume alternatives may replace NT26.23 in their stated ranges. Their complete three-range proof remains unchanged: the small-t coupled-denominator branch is increasing to its endpoint, the middle logarithmic expression is convex, and the final full-volume expression decreases. A maximum of those endpoint alternatives is one cell, not three cells.

\subsection{6. Exactly25 eligible type-II places}\label{exactly25-eligible-type-ii-places}

Here u=p=25 and exactly one active place is I or III. There are26·2=52\textless2\^{}6 type/eligible-set choices; i≤1. Relative to the complete K103 p25 envelopes, the only changes are one additional binary bit and at most one power of T. The pole, central, norm and metric data are identical.

If b\textgreater0, retain g≥126, choose a paired-zero place, and use e1, a120 and cap18. K103 group C bounds every one of its164 complete cells below−172990. Every new cell is therefore below

\[
                       -172990+1+L_T<-171993.              \tag{NT26.26}
\]

If b=0 and q=25, then all142 inactive places are noneligible. Use the full K103 group E cone and volume envelope with a120 and cap68036. Its complete bound for every k is below−1567, so the new cell is below

\[
                          -1567+1+L_T<-570.                \tag{NT26.27}
\]

If b=0 and q≥26, then x\textgreater2\^{}1116 and g≥139. The complete K103 group F envelope has cap1117 and central radius155. Its bound below−58950 becomes below−57953 after the same extra bit and T. The three164-cell unions are consequently below2\^{}-171985, 2\^{}-562 and2\^{}-57945. These are fresh comparisons of complete arithmetic envelopes, not deductions from the prior removal status of an all-II set.

\subsection{7. Exactly24 eligible type-II places}\label{exactly24-eligible-type-ii-places}

Choose E0 in binom(26,24)=325 ways. At its other two active places the type assignment has at most3²−1=8 possibilities because the all-II profile is excluded. Thus all such choices cost2600\textless2\^{}12, and i≤2. Eta is unique by NT26.14. No q-subset or denominator ideal is separately counted.

If b=0, retain g≥138. Evaluate at x=1, p24,e0. The exact cap satisfies

\[
 (R/U)^fT^{48}>30^fQ^{334},\quad T^2>3690^{168},\quad
 H/2+750R/(30\cdot3690)<68.                               \tag{NT26.28}
\]

Thus use A=max(k,30). Compare with the checked K103 group C envelope p25,a120,g126,e1,d13. Its total T exponent7940 becomes7930 after including i≤2; the denominator term increases by20220 bits; the type/set factor decreases by one bit. The raw factor and covering inflation decrease because30≥18. Hence every new cell is below

\[
              -172990-10\cdot995+20220-1=-162721.          \tag{NT26.29}
\]

If b\textgreater0, NT26.21 gives b≤12, so at least126 noneligible inactive factors remain. Choose one paired zero: e1, d≤20, i≤2. The true central radius is bounded by H/2+750R/3690²\textless68 for every k. Discard the global cap and use A=k.

The new complete T exponent7607 is333 smaller than the same reference group C exponent7940; the denominator change is20220 bits and the choice change is+7 bits. For k\textless18, removing cap18 increases the raw factor by at most18\^{}(163f)\textless2\^{}136920. Its entire new inflation is less than2\^{}10, because (1025/1024)\^{}6720\textless2\^{}10; for k≥18 both factors are unchanged. Thus each cell is below

\[
 -172990-333\cdot995+20220+7+136920+10=-347168.             \tag{NT26.30}
\]

The two164-cell unions are below2\^{}-162713 and2\^{}-347160.

\subsection{8. Complete union, executable verification and chronology}\label{complete-union-executable-verification-and-chronology}

The seven disjoint branches are p≤23; p≥24 with x≥2\^{}4096; p24 below that threshold with b0 or b\textgreater0; and p25 below it with b\textgreater0, or b0/q25, or b0/q≥26. They cover every actual non-all-II profile. The complete respective union bounds are

\[
 2^{-172297},\quad2^{-10138},\quad2^{-162713},\quad
 2^{-347160},\quad2^{-171985},\quad2^{-562},\quad2^{-57945}.
                                                               \tag{NT26.31}
\]

Their sum is strictly less than2\^{}-561. Every orientation, support, original pivot, eligible/padded set, actual type, weighted ordered pair, extra-place choice, denominator tail and k-size factor has been retained once. The actual central phase and signed finite factors remain in the kernel; this proof uses their absolute bounds and asserts no sign.

The companion executable recomputes all164 original CF cells, the complete reference K103 C/E/F and infinite-tail cells, all new integer norm/cap/uniqueness comparisons, all type and pole cases, and the seven complete new unions. It verifies the coupled endpoint signs and all directed substitutions. These diagnostics supplement the original lattice, representation and module proofs; they do not enumerate global arrows or evaluate a period.

The default actual predecessor is checks/do109\_joint\_remainder\_error.json, with105 all-top and31 uniform errors and complete floors26/11. Its ordered tail includes R088523, K103293 and DO109820. Append561 once. The --prior option accepts a later complete extension of that exact prefix, retaining every intervening charge. Both complete resulting dyadic sums are checked below2\^{}-230; no term is dropped or reordered. The earlier canonical K103 charge remains even though some numerical envelopes are reused here.

This establishes only the complete non-all-II globally rank-two-unipotent scalar-cross support26 reduction. The all-II branch, other structural classes, the signed sufficient inequality, prescribed global period, selected integral image, actual source constructor, specific032 value, DP44 sign and RH remain separate.

\subsection{9. Source and algebraic-geometric placement}\label{source-and-algebraic-geometric-placement}

Exact carried interfaces are DO099 JR1--24; DO097 UC3--13; DO098 CF4--32; DO102 NT3--21; K071 NB2--9; DO103 LD6--15; DO104 CD5--32; DO105 EL9--23; DO107 IP4--15; DO108 AC3--13; K102 PC5--27; and K103 SC26.5--28. This proof uses their original field, roles, metric and normalization. Their whole reviewed mathematical interfaces are retained; no new external primary theorem is asserted from an unread citation.

The original Hom and determinantal modules remain over Spec(O\_E), and the central coordinate is counted in an affine torsor under the fixed fractional real ideal. Neither a rational normal form nor a saturated local line is declared a globally integral basis. Polarization, theta line, Picard classes and biextensions are unchanged. The inherited contextual anchors remain SGA3 I4.7.3/4.7.3.1; SGA6 VI Theorem3.3; FGA232 Theorem3.1 pp148--149; FGA236 Corollaries1.7/2.3 pp229/232; SGA7 I VII2.7/2.9 pp169--174 and VIII7.4(b) pp310--311; cote4 scan003 §§8.1.1--8.1.2, scans005--006 §§8.4.1--8.4.3; and cote54 scans005--007. No new reading or comparison theorem at these anchors is claimed. Formal-group, Dieudonne, Cartier-dual, lambda-ring, l-adic and motivic assertions keep their existing hypotheses; no such assertion follows from this finite/real count alone.

\subsection{Canonical algebraic-geometric context}\label{canonical-algebraic-geometric-context-23}

The actual fractional coordinate and content ideals remain invertible modules over Spec(O\_E) and Spec(O\_F), with both CM components, original polarizations, Picard classes, theta lines, biextensions, metric and Haar normalization unchanged. Exact carried anchors: SGA3 I4.7.3/4.7.3.1; SGA6 VI Theorem3.3; FGA232 Theorem3.1 pp148--149; FGA236 Corollaries1.7/2.3 pp229/232; SGA7 I VII2.7/2.9 pp169--174 and VIII7.4(b) pp310--311; cote4 scan003 sections8.1.1--8.1.2 and scans005--006 sections8.4.1--8.4.3; cote54 scans005--007. No new lambda-ring/K-theory, l-adic weight or CM-motive comparison follows. DP\_MULT Q3e, THETA3 G1/A7, CX\_GRH CX-ALG-024 and GCT\_PLETH T-052 retain their stated global gates.

\clearpage

\section{AM-LIFT-087 --- effective intrinsic lines on a newly chosen general-Host CM seed}\label{am-lift-087-effective-intrinsic-lines-on-a-newly-chosen-general-host-cm-seed}

\textbf{Canonical checked integration --- 2026-11-23T04:15:47.297026+00:00.} Complete saved-proof and executable reviews and separate exact diagnostic replays passed. No numerical error charge is introduced. Frozen author history is preserved. AM-LIFT-087 GS1--20 chooses an actual new general-Host Serre seed X=O\_L tensor\_\{O\_K7\} J and the actual pure block b\^{}-1 tensor\_\{O\_k\} C0, preserving the Host ideal, Datum, level, coefficient premises and R086 signature/J-content formulas. Every constructed intrinsic component eigenline has degree at least-b, with b=max(log D\_L+.5 log(2t),(sqrt(d)+1)/3+log D\_E+.5 log f). Literal integral b gives1 in b\^{}-1; full-order discriminant-cleared projectors and the exact Zarhin pullback isometry retain all integral models and source pi H norms. The new trace-polarization degree is D\_L/7\^{}(5t), distinct from R086. Under all retained actual rule, KeyCore, integrality/cusp/source premises, LIFT086 now has explicit seed b in C\_N=(D+gn)/2+D log r+(b+log N)min(D+gn,(3D+gn)/2). Effective N, initial field, actual Host/Sizes constructor, selected image, prescribed period, specific032/Read, DP44 and RH remain open. No finite charge; current108/31 arrays and complete floors26/11 unchanged.

Author: \nolinkurl{/root/lift_accounting}, 23 November 2026 UTC. \textbf{Author-frozen proof; final saved-file, source and executable verification records are adjacent. Independent review and canonical integration are separate.} This file is a new construction and does not change the seed used in the frozen R086 proof.

\subsection{1. Scope and quantitative conclusion}\label{scope-and-quantitative-conclusion}

Retain an actual source Host E, its quadratic subfield k=Q(sqrt(-d)), its actual integral ideal b with b bar(b)=Different(E/k), its selected pure type, the prescribed Datum and finite level, and all Moduli/PinDat/Coeff prerequisites of LIFT084. The integrality of b is a literal field of \nolinkurl{DataHost.Host}; the inverse module b\^{}(-1) below is fractional. Write {[}E:Q{]}=2f, D\_E=\textbar disc(E)\textbar, K7=Q(zeta7), L=E K7, c={[}L:E{]}, t={[}L:K7{]}=cf/3, and D\_L=\textbar disc(L)\textbar. The letter k here is a field, not a source rank. The actual Host has discriminant disc(k)=-d.~All logarithms below are natural.

The construction below chooses a different integral polarized induced block before applying the same equivariant principalization and signature balancing as R086. Its actual intrinsic tangent lines have explicit common lower degree bounds

\[
 b_J=\log D_L+\tfrac12\log(2t),\qquad
 b_C=\frac{\sqrt d+1}{3}+\log D_E+\tfrac12\log f,
 \qquad b=\max(b_J,b_C).
 \tag{GS1}
\]

More precisely, it supplies a complete algebraic eigenbasis on the new principal seed, and the intrinsic saturated line generated by each basis vector has normalized degree at least -b, for the literal source norm squared pi H. Repeated product blocks and conjugated actions have the same bound, independent of their multiplicity. This does not bound an arbitrary line in a higher-multiplicity eigenspace.

The source signatures and Jacobian-content formulas of R086 are retained. The seed, its polarization before principalization, and its comparison eigenlines are newly chosen. The prescribed Host, Datum and finite level are not replaced. Finite extension for the actual point, finite level and frame remains allowed exactly as in LIFT084; its total degree is not bounded here. The nonzero frame point subsequently obtained by LIFT085 still has an unevaluated isogeny scale N.

Conditional on all the actual rule, model, KeyCore, local integrality, cusp, LemmaH and source-compatibility hypotheses of LIFT082/LIFT086, substituting this b into LIFT086 gives

\[
 C_N=\frac{D+gn}{2}+D\log r
       +(b+\log N)\min\left(D+gn,\frac{3D+gn}{2}\right).
 \tag{GS2}
\]

Here g is the actual seed dimension, D is the sum of the nonnegative frame-polynomial weights, n is its theta-line power, and r\textgreater=1 is the source block-size bound. These are exactly the variables of the retained Step5 statement; N\textgreater=1 is the fixed-level isogeny scale in LIFT085/LIFT086.

This is an explicit seed-line contribution to a sufficient bound, not an effective bound for N, the full coefficient field, an admissible Host/Sizes example, or DP44. The original arithmetic selected-image, period, Read/STiesR and specific032 gates remain distinct.

\subsection{2. Projective Serre tensors with the actual polarizations}\label{projective-serre-tensors-with-the-actual-polarizations}

Let R be the maximal order of an imaginary quadratic or CM field acting on a principally polarized abelian variety A, with Rosati conjugation. For a finite projective R-module M, use the covariant tensor

\[
 M\otimes_R A:=\operatorname{Hom}_R(M^\vee,A),
 \qquad M^\vee=\operatorname{Hom}_R(M,R).
 \tag{GS3}
\]

Milne (2020, Proposition 7.33 and Remark 7.34, pp.62--63) constructs this over the base field and over an arbitrary base scheme, and proves compatibility with base change. In particular it is an abelian scheme when A is one. Each m in M defines the actual homomorphism A to M tensor A. Over C its tangent map is v to m tensor v.

For precision, a Hermitian form h:M times M to R, linear in its first variable, gives a map M to the dual of the conjugate module. Combining this map with the principal polarization of A gives a homomorphism from M tensor A to its dual. This assertion follows for a free module from the Hermitian matrix formula; for a projective module, realize it as an idempotent summand of a finite free module and restrict that formula. The free-module identities commute with the idempotents and with base change. Thus they prove the assertion for the projective module without choosing a free integral basis.

More explicitly, if Omega\_A is the actual alternating form, the induced alternating form is \[
 \Omega_M(m\otimes x,n\otimes y)=\Omega_A(x,h(n,m)y).
\] Adjointness makes this balanced over R in each argument. The reversed h arguments are required by the first-variable-linear convention; in particular this is the actual source alternating form in GS6. If h is positive at the selected complex embeddings, the corresponding Riemann form is positive and the homomorphism is a polarization. If h identifies M with its Hermitian dual, the polarization is principal. For the map defined by m its Hermitian pullback at an embedding rho of R is

\[
                         H_m=\rho(h(m,m))H_A.
 \tag{GS4}
\]

All these statements concern actual integral morphisms and polarizations. They do not identify a rational comparison basis as an integral basis, and they apply to the actual good models after finite base extension.

For the Jacobian block take

\[
 X=\mathcal O_L\otimes_{\mathcal O_{K7}}J,
 \qquad h_L(x,y)=\operatorname{Tr}_{L/K7}(x\bar y),
 \qquad J=\operatorname{Jac}(y^2=x^7-1).
 \tag{GS5}
\]

The module O\_L is finite projective over the Dedekind ring O\_K7, and h\_L is integral, positive and nondegenerate. Multiplication by O\_L is integral with adjoint conjugation. Thus X has a genuine integral polarization lambda\_X. Its type is the type of J induced from K7 to L, its dimension is 3t, and X is isogenous to J\^{}t. It is a new polarized realization of the same induced CM type used in R086.

For the pure block retain the actual Host ideal, put M=b\^{}(-1), and let C0 be the principally polarized CM elliptic curve with complex lattice c\_K from the source Host. Then

\[
 C=M\otimes_{\mathcal O_k}C_0,
 \qquad h_E(x,y)=\operatorname{Tr}_{E/k}(x\bar y).
 \tag{GS6}
\]

Its trace-Hermitian dual is

\[
 M^\#=\mathfrak D_{E/k}^{-1}\bar M^{-1}
      =(\mathfrak b\bar{\mathfrak b})^{-1}\bar{\mathfrak b}
      =\mathfrak b^{-1}=M.
 \tag{GS7}
\]

Consequently C is principal. Its analytic lattice is c\_K b\^{}(-1), and its alternating form is exactly the trace form of LIFT084 MP2--5, with the same source sign and xi=-1/(N(c\_K)sqrt(-d)). It is therefore the actual pure block, not an unrelated principally polarizable elliptic power. The nonzero integral ideal b is retained with its literal source class; it is not replaced by a more convenient ideal representative.

Choose a finite number-field extension K of the existing coefficient field on which the full actions, polarizations and the needed finite level are defined, containing all L labels, and where J and C0 have good reduction. Such a field exists by the established good models for J and CM potential good reduction for C0; LIFT086 gives an explicit 15-torsion procedure once the initial field has been chosen. The Serre constructions then have actual good models. All following degree statements are invariant under further finite base extension of these models.

\subsection{3. The source degree of the CM elliptic tangent line}\label{the-source-degree-of-the-cm-elliptic-tangent-line}

Let ell\_0 be the intrinsic tangent line of C0 over its good-model field, with source norm squared pi H. At a complex embedding choose a period basis omega\_1, omega\_1 tau for a nonzero invariant differential omega, with tau reduced in the upper half-plane and y=Im(tau). Its Hodge norm is \textbar omega\_1\textbar sqrt(y). If v is dual to omega, its source norm is

\[
                n(v)=\frac{\sqrt\pi}{|\omega_1|\sqrt y}.
 \tag{GS8}
\]

This follows because the principal Hermitian form on the complex tangent line is \textbar z\textbar{} squared divided by the area \textbar omega\_1\textbar{} squared y; the extra pi is the literal IsVecLine factor. Thus the degree of ell\_0 is the negative of the Faltings degree of the invariant-differential line, less one-half log pi.

Write Delta\_mod(tau)=q product\_\{n\textgreater=1\}(1-q\^{}n)\^{}24, q=exp(2pi i tau). The Weierstrass-discriminant normalization is Delta(tau)=(2pi)\^{}12 Delta\_mod(tau). The general discriminant formula for the Faltings height, before any averaging over ideal classes, is equation (6.2) of Barquero-Sanchez et al.~(2017, p.14). Its finite term vanishes on the good model. Therefore

\[
 \operatorname{ndeg}(\ell_0)
   =\log(2\sqrt\pi)
      +\frac1{12[K:\mathbf Q]}
          \sum_{\sigma:K\hookrightarrow\mathbf C}
               \log\bigl(y_\sigma^6
                          |\Delta_{\rm mod}(\tau_\sigma)|\bigr).
 \tag{GS9}
\]

No hypothesis that a model is defined over Q(j) is used: the general embedding sum (6.2) is the input, not the subsequent ideal-class averaging proposition. One can also obtain GS9 directly from a minimal local invariant differential, the discriminant transformation Delta(a omega)=a\^{}(-12)Delta(omega), and the product formula. This explains why arbitrary scaling of omega does not affect the degree.

Every conjugate has full CM order of discriminant -d.~Reduction of its positive binary quadratic form gives

\[
                       \sqrt3/2\le y\le\sqrt d/2.
 \tag{GS10}
\]

The upper bound is y=sqrt(d)/(2a) with positive integer a; the lower bound is the reduced fundamental-domain bound. Put r=\textbar q\textbar. Then r\textless=exp(-pi sqrt3)\textless1/200, and

\[
 24\sum_{n\ge1}\log|1-q^n|
 \ge-\frac{24r}{(1-r)^2}>-1,
 \qquad 6\log y\ge\log(27/64)>-1.
 \tag{GS11}
\]

The first inequality uses -log(1-x)\textless=x/(1-x) and the geometric series. Hence log(y\^{}6\textbar Delta\_mod(tau)\textbar)\textgreater-pi sqrt(d)-2. Since log(2sqrt(pi)) is positive and pi\textless4, GS9 proves the sufficient uniform bound

\[
             \operatorname{ndeg}(\ell_0)
                  \ge-b_\ell,\qquad
             b_\ell=\frac{\sqrt d+1}{3}.
 \tag{GS12}
\]

This bound concerns the actual good-model line and its exact source metric. It does not impose a primitive global tangent generator or an equality of differential ideals.

\subsection{4. The literal Host supplies a bounded integral inclusion}\label{the-literal-host-supplies-a-bounded-integral-inclusion}

The source is stronger here than an arbitrary fractional-ideal input. \nolinkurl{DataHost.Host}, lines64--65, supplies b as an integral ideal and identifies b bar(b) with the relative different. The different is nonzero, so b is nonzero. Consequently

\[
       \mathcal O_E\subset\mathfrak b^{-1}=M,
       \qquad 1\in M.
 \tag{GS13}
\]

Indeed 1 times b lies in O\_E, which is precisely membership in the inverse ideal. No choice of a small element or replacement of the ideal class is needed. The actual integral map j\_1:C0 to C defined by this element has

\[
       h_E(1,1)=\operatorname{Tr}_{E/k}(1)=f,
       \qquad j_1^*H_C=fH_{C_0}.
 \tag{GS14}
\]

Every holomorphic E-coordinate of 1 is nonzero. GS13--14 therefore provide a single actual integral map that detects every required pure eigenline, uniformly at every coefficient embedding. This argument uses the integral-ideal field of the literal Host; it does not assert the same bound for an arbitrary unrelated fractional module.

\subsection{5. Integral projectors and the complete eigenlines}\label{integral-projectors-and-the-complete-eigenlines}

Let A be either E or L and suppose its full order acts on the actual good model. If a\_i is an integral basis and a\_i\^{}* its trace-dual basis, the discriminant inverse-matrix formula proves D\_A a\_i\^{}* integral. Hence the separability projector

\[
        p_\tau=\sum_i\tau(a_i^*)\,\iota(a_i)
 \quad\text{satisfies}\quad
        D_Ap_\tau\operatorname{Lie}(\mathcal A/\mathcal O_K)
                      \subset\operatorname{Lie}(\mathcal A/\mathcal O_K).
 \tag{GS15}
\]

This is LIFT081 ED14--17 with a universal discriminant clearing integer instead of the special original-field constants. Rosati conjugation makes the embedding eigenspaces orthogonal. Thus p\_tau has operator norm at most one in every actual complex realization, including the source pi factor. The labels and algebraic coefficients are transported under each embedding of K.

On X, take the actual homomorphism j\_1:J to X defined by 1 in O\_L. At a K7 embedding, h\_L(1,1)=t. The algebraic vector j\_1(2e\_i) has a nonzero component in every L eigenspace extending the K7 label of e\_i: under L tensor\_\{K7,rho\} C equal to the product of its embedding copies, its coordinate is 1 in every copy. LIFT081 ED8--13 gives 2e\_i integral on the actual good Jacobian model and \textbar\textbar2e\_i\textbar\textbar\_source\textless sqrt2. Consequently

\[
       v_\tau=D_Lp_\tau j_1(2e_i)\ne0,
       \qquad v_\tau\text{ integral},
       \qquad\|v_\tau\|_\sigma<D_L\sqrt{2t}
                              \quad\text{for every }\sigma.
 \tag{GS16}
\]

There is one vector for every holomorphic L embedding. They form a complete eigenbasis because each such eigenspace has dimension one. In particular no short integral basis of O\_L and no associated basis-height factor is needed.

On C, use the actual j\_1 of GS13--14. Let v be any nonzero algebraic tangent vector of C0, and set w\_tau=D\_E p\_tau dj\_1(v). Before projection, its component at every holomorphic E label is v, hence each projected vector is nonzero. The map D\_E p\_tau dj\_1 is an integral O\_K-linear map of tangent modules: dj\_1 is induced by the actual group homomorphism, and D\_E p\_tau is the integral Lie-module operator from GS15. No group-scheme homomorphism is asserted for the projector itself. This map carries the intrinsic ideal of v into the intrinsic ideal of w\_tau. By GS14--15 its metric expansion is at most

\[
                             D_E\sqrt f.
 \tag{GS17}
\]

For a morphism of rank-one intrinsic lattices with this metric bound, inclusion of the image ideal and the product formula give

\[
 \operatorname{ndeg}(K w_\tau\cap\operatorname{Lie}\mathcal C)
  \ge\operatorname{ndeg}(\ell_0)
       -\log(D_E\sqrt f)
  \ge-b_C.
 \tag{GS18}
\]

Here the source ideal need not be principal. If I\_0 and I\_tau denote the two coefficient ideals for v and w\_tau, integrality gives I\_0 contained in I\_tau, hence N(I\_tau)\textless=N(I\_0). The literal degree is minus log of this ideal norm and minus the sum of all embedding log norms, divided by {[}K:Q{]}. Together with n\_tau(sigma)\textless=D\_E sqrt(f) n\_0(sigma), this proves the lower-degree direction in GS18. Exactly the same calculation with the integral vector GS16 and O\_K contained in its saturated coefficient ideal gives degree at least -b\_J on X. All eigenvectors are chosen algebraically once, not separately at complex embeddings. Rescaling them does not change the intrinsic normalized degree.

\subsection{6. Principalization preserves the vector norms}\label{principalization-preserves-the-vector-norms}

Apply the equivariant quaternion construction to the actual polarized X. In Zarhin (2023, arXiv v5, §4, equations(16)--(19), pp.5--6), with n\_X=deg(lambda\_X), the isogeny is

\[
 \pi:X^4\times X^4\longrightarrow(X^\vee)^4\times X^4,
 \qquad (x,y)\longmapsto(\lambda_X^4x,Ix-y),
 \qquad\pi^\vee\mu\pi=\lambda_X^8.
 \tag{GS19}
\]

The resulting mu is principal and the action on the dual is a to iota(bar(a)) dual. The displayed isogeny is fully O\_L-equivariant. The Riemann forms in GS19 imply the exact metric identity H\_mu(dpi(v),dpi(v))=H\_lambda\_X\^{}8(v,v). Therefore the source norms, with the same pi factor, are unchanged. No estimate on the entries of the quaternion matrix I is needed for this identity.

The isogeny extends between the actual good models, so it carries integral tangent vectors to integral vectors. It is invertible on the generic tangent space. Applying it to the eigenbasis in eight product copies of X therefore supplies a complete eigenbasis of the principal Y, with the same norm bounds and intrinsic degree lower bounds. The argument gives a lower bound for the saturated target line even if the image vector is not primitive.

For reference the trace dual of O\_L over O\_K7 is the relative codifferent. The integral Betti lattice of J is a rank-one projective O\_K7 module, with its perfect principal alternating pairing. Tensoring that lattice with O\_L and its trace dual shows that the polarization's dual-lattice index is the absolute norm of the relative different. This calculation is local over O\_K7 and does not assume a global free Betti generator. Thus the degree in GS5 is

\[
             n_X=N_{L/\mathbf Q}\mathfrak D_{L/K7}
                     =D_L/7^{5t},
             \qquad\deg\pi=n_X^4.
 \tag{GS20}
\]

This differs from R086's explicit initial trace polarization. Neither degree is silently identified with the other. The exact pullback identity GS19 is why principalization does not introduce an additional adverse metric factor in the present line lower bounds. No claim is made that the two differently constructed principal seeds coincide.

\subsection{7. Signature balancing, all source clauses, and remaining costs}\label{signature-balancing-all-source-clauses-and-remaining-costs}

The induced L type, dimension and conjugate-dual convention agree with the inputs to R086. Thus its exact balancing proof applies to the newly chosen Y: use Y times its conjugate O\_E action for odd c, the even-c improvement when applicable, and pure C or C minus to fill the remainder. For d=7 use the three corresponding K7 action twists of J before the same construction. Twisting the action permutes eigenlabels and leaves the underlying principal metric and integral model unchanged. Hence every component line has the common bound GS1.

The resulting principal product has the same exact multiplicity and J-content formulas of R086, with block sizes 8 or24 in the general case and the separate (16,8) block for d=7. Its actual pure block is GS6, including the terminal q=0 case. All assertions about the literal DetCond, source sign, canonical trace-one rigidified cubical line, all-a descent, finite TateIso and fixed-level Obj/FrPt are precisely the hypotheses and constructions of LIFT084/R086. This proof does not obtain a global integral Betti isometry, change the source finite level, or descend an infinite adelic level to one finite field.

For the comparison eigenlines, saturation and every flat/complete-DVR base change follow as in LIFT082/LIFT086: write the chosen rational line as the common kernel of finitely many integral linear forms, then use flatness of integer-ring extension. This supplies actual IsVecLine lattices with the source norm, not embedding-coordinate ideals. The full determinant line remains the determinant of the whole tangent lattice, not automatically the product of those saturated line lattices.

Applying LIFT085 to a nonzero rule and then LIFT086 transfers these seed lines to its selected point, with degree loss at most log N. Under all the separately retained source premises, the optimized bound is GS2. The finite extension bound relative to the chosen seed field is still the LIFT086 bound \nolinkurl{(15 N^2)^(4g^2)} (or \nolinkurl{N^(8g^2)} for an already good seed). This does not bound the initially chosen seed-field extension, N, the new rule's arithmetic coefficients, or any untransported DP44 comparison data.

The new contribution is the evaluated intrinsic seed-line lower bound for a specified general-Host construction. It is not a restatement of the special J14/J56 bounds, and those excluded fixed-degree numerical Host profiles are not used as admissible examples. All signed finite-kernel gates, selected-image content, nonzero global central value or period, actual Read/032, and RH remain open.

\subsection{8. Primary and source anchors; verification boundary}\label{primary-and-source-anchors-verification-boundary}

\begin{itemize}
\tightlist
\item
  Milne, J. S. (2020). \emph{Complex multiplication} (v0.10). Proposition7.33, Remark7.34 and Proposition7.36, pp.62--63: projective tensor construction and base change. \url{https://www.jmilne.org/math/CourseNotes/CM.pdf}
\item
  Milne, J. S. (2021). \emph{Elliptic curves} (2nd ed.). ChapterV, Theorem3.7, printedp.192: the discriminant product with its \nolinkurl{(2pi)^12} factor; ChapterII section2 gives the Weierstrass invariant normalization. \url{https://jmilne.org/math/Books/EC2.pdf}
\item
  Barquero-Sanchez, A., Cadwallader, L., Cannon, O., Genao, T., \& Masri, R. (2017). Faltings heights of CM elliptic curves and special Gamma values. \emph{Research in Number Theory,3},13. Equation(6.2), p.14 is used in its general embedding-sum form; no use of the later Q(j)-dependent averaging step is made. \url{https://doi.org/10.1007/s40993-017-0077-7}
\item
  Zarhin, Y. G. (2023). \emph{Abelian varieties, quaternion trick and endomorphisms}. Theorem1.1 and the literal isogeny/pullback/equivariance equations(16)--(19), arXiv v5 pp.5--6. \url{https://arxiv.org/abs/2010.07995}
\end{itemize}

The literal source clauses are \nolinkurl{DataHost.lean}, Host39ff, especially integral b64 and relative-different identity65, ccK86--87, xi0 at149, AA0 at153--154 and AA at156--157; \nolinkurl{LamLetters.lean}, IsModelAt87--94 and IsVecLine96--106; \nolinkurl{DetIdeal.lean}, MLine10--14, deg24--25, ndeg27 and dual33; and \nolinkurl{StatementS.lean}, full Step5Loc38--53, Step5Arch55--61 and Step5 63--73. The actual Host/Datum/SAB/TateIso/Obj clauses in \nolinkurl{TwoWeightEqOneAddRoot.lean} retain their LIFT084 scopes. Carried checked interfaces are LIFT081 ED8--20, LIFT084 MP1--19, R086 ZJ1--21, LIFT082's full all-DVR/metric source transport, and LIFT086 IM1--27. Exact byte hashes and ranges belong to the companion source manifest.

The Serre tensors, their integral integer-ring actions, the equivariant finite isogeny and the principal polarizations extend over the chosen good-model base. The canonical relative Picard/Poincare and rigidified cubical theta line are the already checked LIFT084/LIFT086 constructions. This is an evaluated intrinsic-line cost for a newly specified seed; no new motivic, slope, or divided-derivative identity follows. The exact Grothendieck/run context remains that recorded in those findings, including FGA232 Theorem3.1, FGA236 Corollaries1.7/2.3, SGA7 I VII2.7/2.9 and VIII7.4(b), and DP\_MULT Q3e/DP44/DP52/DP54--65 and THETA3 G1/A7.

The companion diagnostics check the displayed constants, trace-dual and fractional-ideal signs, exact Hermitian/quaternion matrix identities, complete type counts and every sufficient-cost branch. They do not replace the universal Serre construction, good-model, polarization, source-transfer or nonzero-locus theorems. The final whole saved-proof/source/script check, freeze and independent review are separate records.

\subsection{Canonical algebraic-geometric context}\label{canonical-algebraic-geometric-context-24}

The actual abelian schemes, integer-ring action, principal Picard/Poincare construction, normalized cubical line, finite-flat isogeny and its kernel live over Spec(O\_K). Away from N the displayed adjoint inverse gives an actual p-divisible-group and formal-group/Dieudonne comparison, without an operator or slope gain. Carried exact anchors: FGA232 Theorem3.1 pp148--149; FGA236 Corollaries1.7/2.3 pp229/232; SGA7 I VII2.7/2.9 pp169--174 and VIII7.4(b) pp310--311; SGA3 I4.7.3/4.7.3.1; SGA6 VI Theorem3.3; cote4 scan003 sections8.1.1--8.1.2 and scans005--006 sections8.4.1--8.4.3; cote54 scans005--007. No new lambda-ring/K-theory, l-adic weight or CM-motive identification. DP\_MULT Q3e/DP44/DP52/DP54--65, THETA3 G1/A7, CX\_GRH CX-ALG-024 and GCT\_PLETH T-052 retain their separate arithmetic comparisons.

\clearpage

\section{AM-R-090 --- intersecting eligible sets exhaust full-positive rank three at support26}\label{am-r-090-intersecting-eligible-sets-exhaust-full-positive-rank-three-at-support26}

\textbf{Canonical checked integration --- 2026-11-23T04:26:02.790256+00:00.} Complete root and independent proof/script/interface reviews and separate exact replays passed. The actual canonical predecessor has108/31 entries; this integration gives109/31, complete floors26/11. Frozen author history is preserved. AM-R-090 IS1--25 removes the entire actual full-positive globally rank-three scalar-cross support26 class with nonzero original S-cross, without a sparsity hypothesis. The complete eligible-set count at most24 costs6670; the complete all-k denominator tail at least2\^{}8192 costs11965. Literal inversion applies each containing-set bound once more. Outside them, the two eligible sets intersect in at least24 paired places, forcing equal null lines by the exact6395-bit denominator-cleared norm gap and contradicting global rank three. The full four-sector union is below2\^{}-6668. Actual DO110108/31 becomes109/31, every previous overlapping charge retained, both complete sums below2\^{}-230 and floors26/11 unchanged. The nonzero short-compression and dense all-II rank-two classes, signed positivity, actual source, selected arithmetic content, prescribed period, specific032/Read, DP44 and RH retain their gates.

Author: /root, GPT-6 Astra Pro; reasoning effort not externally exposed; ALP\_MACHINE. 23 November2026 UTC. Controlled proof. Whole saved-proof/script readings, exact replay, independent review and canonical integration are separate events.

\subsection{1. Actual unchanged sector and conclusion}\label{actual-unchanged-sector-and-conclusion}

Retain the actual J56 all-top kernel, H=133, f=168, the original degree336/168 fields E/F, all334 unmixed G reductions, both mixed factors, both CM partners of every original T pencil, all admitted labels, the selected real module and metric, packet representative, central normalization, class masses and Haar measures. At support s=26 assume the full-positive form

\[
 E^{169}=Z\perp P,\quad Z=S_{164}\perp u^\perp,\quad
 P=Eu\perp Ee_n,\quad J_P=\operatorname{diag}(U,-1),
\] \[
 h=\begin{pmatrix}I_Z&B\\C&D\end{pmatrix},\quad
 (h^{-1})_{ZZ}=I_Z,\quad\det h=1,\quad C|_S\ne0,
 \quad\operatorname{rank}_E(h-I)=3.                 \tag{IS1}
\]

All unmixed G reductions are unipotent. No global unipotence or sparsity is assumed. Normalize only over E:

\[
 C=\ell r,\quad\ell=(1,\eta),\quad\eta\bar\eta=U,
 \quad U=77Q^2-14,
 \quad Q=1427908193729121522601.                    \tag{IS2}
\]

The original u and fixed three integral transverse inputs are those of R072 PJ4--5/R076 CV3. All167 values of r are counted in that original fixed basis. Write x=N\_E D\_eta for the exact denominator ideal of eta. The actual number t of type-III places satisfies11\textless=t\textless=26 by the support-independent DO074 TC1: t\textless=10 would force global unipotence of rank at most two.

\textbf{Theorem.} The complete actual IS1 sector has total absolute normalized contribution

\[
                         \sum_h|\mathcal W(h)|<2^{-6668}.
                                                        \tag{IS3}
\]

This is a full rank-three theorem with no lower bound beyond k\textgreater=1 on the number of nonzero original S-row coordinates. It applies to the full-positive part of the support26 remainder whenever that form has been established. It introduces no complete support-floor increase or signed-kernel, actual-source, selected-image, prescribed-period, specific032, DP44 or RH conclusion.

\subsection{2. A complete count with at most24 eligible active places}\label{a-complete-count-with-at-most24-eligible-active-places}

Define the original eligible set

\[
 E_\eta=\{v\in\Sigma:v_{\mathfrak P_v}(\eta-a_0)\ge1\},
 \qquad a_0=Q(-7+2\sqrt{-7}),\quad a_0\bar a_0=U+14.
                                                        \tag{IS4}
\]

K097 RP4--7 is a support-independent cross-block calculation. Every original positive input belongs separately to each source and target pencil. Its image difference is Cy because h\_ZZ=I. A nonintegral cross value forces the saturated null line to be the original pole line. At the selected component the slope is a0, and at the conjugate component it is U/bar(a0). Thus a cross pole at either partner forces eligibility at the same selected place and makes eta a unit at both partners. Type-I integrality and primitive-u extraction are also unchanged. No rank-two full-double-pole or central-fibre assertion is used here.

Suppose \textbar E\_eta\textbar\textless=24. Pad E\_eta to a24-subset D of the26-place support, at most binomial(26,24)=325 choices. With the original common whole-C G ideal J put

\[
 L=JI_D,\quad N_E I_D=T^{-48},\quad
 z=(N_E L)^{1/(2f)}\ge T^{-24/f}.                 \tag{IS5}
\]

All167 original r-values and the negative pivot beta=eta r(e\_j) belong to L. The actual common G pattern count is still at most2\^{}668. With M=H/(2sqrt U), B0=H/2, the complete R076/K097 factors are

\[
 \#r\le A(4U)^f(M+4200z)^{334f}z^{-334f},\qquad
 A=\frac{528^{167f}}{(167f)!D_0^{167}}<2^{-387281},
                                                        \tag{IS6}
\] where D\_0=7\^{}140 29\^{}144 5\^{}84 13\^{}84 is the fixed scalar discriminant factor. The covector metric is diag(I164,G\^{}-1), detG=4U, G-I\textgreater0, and its full covering radius is less than4200sqrt(f)z. The entire ball includes every zero-coordinate pattern, with no binomial or k union.

\[
 \#\beta\le28C_F B_0^{f/6}z^{-f/6},\quad C_F<2^{42},
 \qquad |c(h)|<2^{-929f/2}U^{-f/2}z^{-f},
                                                        \tag{IS7}
\] \[
 \#\operatorname{pr}_P(he_n)<191^{4f}U^fT^{104},
 \qquad G_2<2^{36}.                              \tag{IS8}
\]

The exact relative-norm fibre in IS7 counts integral relative-norm ideals first and at most28 actual norm-one generators each; nonprincipal and non-conjugation-stable ideals remain included. IS8 is the full original projected-column and source-fibre count at s=26. It retains the fixed U denominator and all labels, including zero.

The two genuine pointwise Green estimates on the same original factor are

\[
 |a(h)|\le2Q^{-335\cdot167}T^{335\cdot24}z^{335f},
 \quad |a(h)|<2Q^{-501\cdot167}2^{9\cdot167f}T^{167t}.
                                                        \tag{IS9}
\]

The second is DO095 RA15--20's original nonzero three-minor bound, with only simple paired type-III poles. It needs no extra witness or pattern choice. The original two mixed finite factors still have modulus at most one. The whole fixed factor is2\^{}881: Green2, G patterns2\^{}668, support choices2\^{}168, pivots2\^{}8 and source fibre2\^{}36.

For each fixed t, the complete general type/weighted-pair sum is

\[
 {26\choose t}(4T)^{26-t}<2^{P_t},\quad
 P_t=78-2t+(26-t)L_T,\quad L_T=1019333/1024.
                                                        \tag{IS10}
\]

Indeed type I has complete weighted cost below3T, types II and III at most one, and binomial(26,t)\textless2\^{}26. This replaces the old coarse pair bound; it is not multiplied by it. Every arithmetic count is uniform for the fixed original source/target pair before weights are summed. No independent CM-partner factor appears.

\subsection{3. The entire two-range bound at support26}\label{the-entire-two-range-bound-at-support26}

Set z\_*=2\^{}-95. The actual integers give M\textgreater2\^{}-68 and

\[
 4200z_*/M<4200/2^{27}<2^{-14},\qquad
                   (1+2^{-14})^{334f}<2^5.       \tag{IS11}
\]

For z\textless=z\_* the full radius inflation in IS6 is therefore below2\^{}5. Combining IS6--9 with weak Green leaves z\^{}(-f/6). Its maximum is controlled by the lower endpoint in IS5, giving T\^{}4; the negative power is not evaluated at z\_\emph{. For z\textgreater=z\_} use strong Green. The whole remaining dependence is

\[
                   (M/z+4200)^{334f}z^{-7f/6},    \tag{IS12}
\]

which decreases on z\textgreater0. Its upper bound is therefore its value at z\_*, with the same full inflation bound. No ideal attaining this cutoff is assumed.

Use directed LQ=287843/4096, LH=28900/4096, LU=601356/4096 and L191=15519/2048, respectively lower for Q, upper for H, lower for U and upper for191. Define

\[
 C_*=881+5+5+42-387281+4fL_{191}-929f/2+2f
       +(2005f/6)(L_H-1)-(331f/2)L_U.             \tag{IS13}
\]

The two fives are the full covering inflation and28\textless2\^{}5. The complete32 cells, before325 paddings, are

\[
 E_t^-=C_*-335\cdot167L_Q+(335\cdot24+4+104)L_T+P_t,
\] \[
 E_t^+=C_*-501\cdot167L_Q+9\cdot167f
                 +(167t+104)L_T+P_t+95(2011f/6).
                                                        \tag{IS14}
\]

Their exact maxima are

\[
 \max E_t^-=-27374879/4096<-6683\quad(t=11),
 \qquad\max E_t^+=-217699105/4096<-53149\quad(t=26).
                                                        \tag{IS15}
\]

There are16 actual t values. Consequently the complete at-most24-eligible sector, with every coordinate support and denominator, has absolute contribution

\[
 325\cdot16(2^{-6683}+2^{-53149})<2^{-6670}.       \tag{IS16}
\]

No denominator condition has entered this count. It is a fresh application of the complete K097/R076 envelope at the stated support and cutoff.

\subsection{4. A complete denominator tail with no sparsity requirement}\label{a-complete-denominator-tail-with-no-sparsity-requirement}

Now count the complete IS1 sector with x\textgreater=2\^{}8192, without any eligibility premise. The fixed-denominator enumeration R083 LD3--6 gives fewer than2\^{}39 norm-U scalars per exact integral denominator ideal. Its actual U factorization and decomposition counts remain those already checked there; no class-group factor is omitted.

For fixed eta use the common original G ideal and all26 active pole places:

\[
 L=JI_\Sigma,\quad N_E I_\Sigma=T^{-52},\quad
 L\cap\eta^{-1}L=LD_\eta,\quad
 z'=(N_E(LD_\eta))^{1/(2f)}.                       \tag{IS17}
\]

Every original r-value lies in this exact intersection, including at G and inactive primes. Since the original S-cross is nonzero, at least one such S value is nonzero. Norm/AM--GM and the actual joint row bound give z'\textless M, without the old k\textgreater=126 premise. Hence the full167-coordinate covering inflation satisfies

\[
              1+4200z'/M<4201<2^{49/4}.            \tag{IS18}
\]

Use the full covector volume IS6 with z', the same selected trace with z', and the same full projected-column/source-fibre count. The weak Green factor is

\[
 2Q^{-335\cdot167}T^{335\cdot26}(z')^{335f}x^{-335/2}.
                                                        \tag{IS19}
\]

The powers of z' cancel exactly:335f-334f-f=0. Fixing eta removes the beta norm fibre; it is replaced by the entire eta count2\^{}39. All patterns/supports/pivots and the original selected metric remain. The16 type factors are each bounded by P11=56+15L\_T. Thus for fixed denominator D the whole contribution is bounded by16 times2\^{}C x\^{}(-335/2), where

\[
 C=881+39-387281+4fL_{191}-929f/2+2f
     +334f(L_H-1)-(331f/2)L_U-335\cdot167L_Q
     +(335\cdot26+104)L_T+(56+15L_T)+(49/4)334f.
                                                        \tag{IS20}
\]

The complete ideal tail satisfies, as in R083 LD14--15,

\[
 \sum_{N_ED\ge2^{8192}}(N_ED)^{-335/2}
       \le2^{-(331/2)8192}\zeta_E(2)
       <2^{336-(331/2)8192}.                       \tag{IS21}
\]

All integral ideals, including those that are not denominators, are included. With16=2\^{}4 the full exponent is

\[
 C+4+336-(331/2)8192=-49009671/4096<-11965.          \tag{IS22}
\]

This proves an absolute bound2\^{}-11965 for the entire large-denominator sector. No sum over k or lower sparsity threshold is needed, because the full volume already counts every original coordinate pattern and z'\textless M follows from a single nonzero S value.

\subsection{5. Inversion and the intersecting eligible sets}\label{inversion-and-the-intersecting-eligible-sets}

DO093 IV4--16 gives the literal identity W(h\^{}-1)=conjugate(W(h)). Inversion transposes the same original paired pencil incidence and retains every real symmetric weight, actual trace, central factor, class mass, Haar measure, cutoff, support and incidence type. The complete containing sectors in IS16 and IS22 may therefore be used for the inverse even if the already cut remainder is not inverse-closed.

DO095 RA3--5 gives detD=1 and BD\^{}-1 C=0. Hence

\[
 C^-=(h^{-1})_{PZ}=-D^{-1}C=\gamma(1,\xi)r,
 \quad\gamma\ne0,\quad\xi\bar\xi=U.              \tag{IS23}
\]

The normalized first component is nonzero because a nonzero isotropic vector in diag(U,-1) cannot have first component zero. In particular the inverse S-cross remains nonzero; no sparsity assumption is needed. Full positive compression, rank3 and the finite unipotence premises are preserved. Applying each entire bound to the inverse costs exactly one additional copy of that bound.

Outside these four counted sectors, x\_eta and x\_xi are both less than2\^{}8192, and each eligible set has at least25 of the same26 active places. Their intersection therefore has at least24 places. At each common place both scalars equal the selected slope a0 and their conjugate-partner slopes both equal U/bar(a0). Both denominators are units at both partners.

If delta=eta-xi were nonzero, (delta)D\_eta D\_xi would be an integral ideal. The valuation inequality min(a,b)+max(0,-a)+max(0,-b)\textgreater=0 proves this even for nonprincipal and non-conjugation-stable denominator ideals. All48 paired degree-one primes above the24 common places divide this ideal. Its archimedean norm satisfies

\[
 T^{48}\le |N_E\delta|x_\eta x_\xi
           <(4U)^{168}2^{16384}<T^{48},             \tag{IS24}
\]

where the final strict inequality is checked with the original exact integers. The bound on delta uses \textbar eta\_tau\textbar=\textbar xi\_tau\textbar=sqrt U at every complex label. Thus eta=xi. There is no enumeration of the two eligible sets or their intersection: this is a deterministic implication for each matrix after both containing-set bounds have been paid.

Equality of the null lines means D preserves E ell. With b=-ell*J\_P D, this gives b ell=0 and BC=0. The exact DO095 RA9 alternative then forces global rank(h-I)=2, contradicting IS1. This is the same support-independent rank argument as R084 IL19--20. It does not assert global unipotence from line equality.

Every IS1 matrix is consequently in at least one of the four counted sectors, and their full union is

\[
                  2(2^{-6670}+2^{-11965})<2^{-6668}.
                                                        \tag{IS25}
\]

Overlaps with earlier charges and between containing sectors are retained. Both inverse costs are explicitly included.

\subsection{6. Exact diagnostics, chronology and open gates}\label{exact-diagnostics-chronology-and-open-gates}

The companion independently recomputes the actual host, all directed integer logarithms, the complete factorial volume, the full covering inequalities, all32 type/range cells, the entire padded union, the new all-k ideal-tail endpoint, the48-prime uniqueness comparison and both complete ordered error sums. Finite valuation/overlap diagnostics corroborate the displayed algebra; they do not enumerate arithmetic arrows or prove a global period.

The author base is checks/do110\_joint\_remainder\_error.json,108 all-top and31 uniform entries, complete floors26/11. Append6668 once. Its explicit --prior option accepts only a complete later canonical extension of that base, preserving every intervening entry and the entire uniform array. Both full sums stay below2\^{}-230. No full support27 floor is claimed: the other structural classes remain under their separate results.

The original normalized signed inequality, actual fully admissible mixed-CM family and all six lift identities, controlled arithmetic selected-image content, prescribed Bessel/Mram and global period, specific032/Read/STiesR and the actual evaluated DP44 comparison remain open. No Kummer-Hecke RH or global oscillator comparison follows.

\subsection{7. Run, source and algebraic-geometric placement}\label{run-source-and-algebraic-geometric-placement}

The used checked interfaces are K097 RP4--7 and RP10--21 (cross-only poles and complete relative-norm/full-volume envelope); R076 CV4--18 (actual full covector metric and fibre); R083 LD3--15 (exact denominator enumeration, intersection, trace and ideal tail); R084 IL9--20 (literal inverse line and paired difference); DO093 IV4--16 (actual coefficients/source-target inversion); DO095 RA1--20 (full-positive rank classification and original minor); R072 PJ4--18 (original scalar data, projected column and source fibre); R082 TP3--4 (general weighted type count); DO074 TC1 (support-independent t\textless=10 exclusion). Known results are invoked at those exact scopes; the new support26 endpoints, no-sparsity tail and25-of26 intersection are proved above.

The ideals remain invertible modules on the original Spec(O\_E). IS17 is their actual intersection, while IS24 is the divisor and field norm of a nonzero section in that unchanged module. The full covector volume uses the fixed Hermitian polarization, with its determinant and finite discriminant retained. Inversion is the original adjoint Hermitian correspondence with both source/target CM partners. No new integral null frame, theta line, polarization, Picard class, biextension or Haar normalization is introduced.

Exact carried contextual locators: SGA3 I4.7.3/4.7.3.1; SGA6 VI Theorem3.3; FGA232 Theorem3.1 pp148--149; FGA236 Corollaries1.7/2.3 pp229/232; SGA7 I VII2.7/2.9 pp169--174 and VIII7.4(b) pp310--311; cote4 scan003 sections8.1.1--8.1.2 and scans005--006 sections8.4.1--8.4.3; cote54 scans005--007. DP\_MULT Q3e and THETA3 G1/A7 gain precisely the smaller finite remainder. No new formal-group slope, lambda-ring/K-theory identity, l-adic weight, CM-motive comparison or motivic L-function result is asserted; CX\_GRH CX-ALG-024 and GCT\_PLETH T-052 retain their separate arithmetic gates. No Lean build or archive audit is performed.

\subsection{Canonical algebraic-geometric context}\label{canonical-algebraic-geometric-context-25}

The reductions are the actual integral Hermitian modules at unmixed finite points; content and coordinate ideals remain fractional invertible modules over Spec(O\_E) and Spec(O\_F), with the same paired CM action, polarization, theta line, Picard scheme and biextension. Carried exact anchors: SGA3 I4.7.3/4.7.3.1; SGA6 VI Theorem3.3; FGA232 Theorem3.1 pp148--149; FGA236 Corollaries1.7/2.3 pp229/232; SGA7 I VII2.7/2.9 pp169--174 and VIII7.4(b) pp310--311; cote4 scan003 sections8.1.1--8.1.2 and scans005--006 sections8.4.1--8.4.3; cote54 scans005--007. No new lambda-ring/K-theory, l-adic weight or CM-motive comparison is asserted. DP\_MULT Q3e, THETA3 G1/A7, CX\_GRH CX-ALG-024 and GCT\_PLETH T-052 retain their stated global gates.

\clearpage

\section{AM-DO-111 --- unitary positive compression removes the nonzero-short-compression sector at support26}\label{am-do-111-unitary-positive-compression-removes-the-nonzero-short-compression-sector-at-support26}

\textbf{Canonical checked integration --- 2026-11-23T04:28:11.459550+00:00.} Complete root and independent proof/script/interface reviews and separate exact replays passed. The actual canonical predecessor has109/31 entries; this integration gives110/31, complete floors26/11. Frozen author history is preserved. AM-DO-111 VC1--25 proves the actual positive transverse compression is unitary when both original scalar-cross directions lie in the retained plane. Its four original integral short witnesses have the sharper all-embedding bound2\^{}41. The complete simultaneous finite restoration, arbitrary-ideal row volume, relative norm fibre, original selected trace, full W0 column, source fibre and all paired type weights give exact exponent-87517717/4096\textless-21366 for the entire nonzero-short-compression support26 sector. The complementary integral norm argument gives full positive167 compressionI and a positive166 fixed space, hence global rank at most3; global unipotence is not inferred. Actual R090109/31 becomes110/31 with one21366 append, every old overlapping charge retained, both entire error sums below2\^{}-230 and complete floors26/11 unchanged. Whole support26 assembly remains separate; dense all-II rank-two, signed positivity, actual admissible source, selected arithmetic content, prescribed period, specific032/Read, DP44 and RH remain open.

Author: /root/definite\_oscillator, ALP\_MACHINE. 23 November2026 UTC.

Controlled pending proof. Author whole-file verification, exact arithmetic, independent review and canonical integration are recorded separately. No earlier frozen finding is changed.

\subsection{1. Actual scope and conclusion}\label{actual-scope-and-conclusion-1}

Keep the original J56 all-top normalized kernel, H=133, f=168, n=169, S of dimension164 and the original W0 of signature(4,1). Every field, lattice, finite projector, source/target pencil and CM partner, admitted label, class mass, Haar measure, central phase, selected Hilbert metric, compact factor and actual near-point representative remains unchanged. Work in the existing K068 remainder, where all334 unmixed G reductions are unipotent. Global unipotence is not assumed.

Write h\_SS=I, C=h\_(W0,S)≠0, and N=h−I. At support26 the already charged DO099(JR1--12) interval removes transverse C and transverse lower cross block of h\^{}-1. Thus both are in the original packet plane

\[
 P=Eu\oplus Ee_n,\qquad
 u=(2,95192390527,383092671043,12529843548438452388609),
\] \[
 U=\langle u,u\rangle=77Q^2-14,
 \qquad Q=1427908193729121522601.                          \tag{VC1}
\]

Let V=u⊥ in the four positive W0 coordinates, so dim V=3 and the original orthogonal decomposition over E is S⊕V⊕P. Use the original integral short vectors

\[
 y_1=(-976627,-153032,38026,0),\qquad
 y_2=(-567678,695570,-172838,0),
\] \[
 \|y_1\|^2=978665066829,\qquad \|y_2\|^2=835948910828.     \tag{VC2}
\]

Their negative coordinates are zero. The unchanged Hermitian pairing convention is that used in R070 and DO099. Define the same four original entries

\[
                 \Delta_{ab}=\langle y_a,Ny_b\rangle,
                          \qquad1\le a,b\le2.             \tag{VC3}
\]

\textbf{Theorem.} At s=26, the entire stated doubly-plane scalar-cross sector with at least one nonzero Delta\_ab has absolute normalized contribution less than

\[
                              2^{-21366}.                  \tag{VC4}
\]

Every actual incidence profile is included. No additional type-III threshold or global-rank restriction is required in the count. In particular it covers the nonunipotent N1 branch of the source-pinned support26 taxonomy. There is one new all-top charge and no uniform charge.

After this containing-sector removal, all four short compressions vanish in the scalar-cross remainder at support26. The support-independent original integral norm argument then gives full positive167 compression equal to the identity and a positive166-dimensional pointwise fixed space. Hence global rank(h−I)≤3. This is a structural consequence within that class; it is not a complete support27 assertion.

\subsection{2. Exact positive-compression lemma}\label{exact-positive-compression-lemma}

The following algebraic argument applies over the original E and at every real embedding. No finite-field positivity or integral moving basis is used.

Because C is plane-valued, the V/S block of h is zero. Because the lower cross block of h\^{}-1 is also plane-valued, its V/S block is zero. Taking the actual Hermitian adjoint shows that the S/V block of h is zero. Thus, over the original orthogonal decomposition, h has the form

\[
 h=\begin{pmatrix}
 I&0&B_S\\
 0&A&B_V\\
 C_S&C_V&D
 \end{pmatrix}.                                           \tag{VC5}
\]

Here C\_S:S→P is nonzero. From unitarity on the S columns,

\[
                        C_S^\dagger J_P C_S=0.
\]

Its image is therefore the isotropic line ell in the nondegenerate plane P; equivalently this is the original rank-one C assertion R059(GB3--6). The S/V column pairing gives

\[
                        C_S^\dagger J_P C_V=0.             \tag{VC6}
\]

Inside a nondegenerate Hermitian plane the orthogonal complement of a nonzero isotropic line is that same line. Hence im C\_V⊂ell and C\_V†J\_P C\_V=0. Applying unitarity to the V columns proves

\[
                       A^\dagger G_V A=G_V.                \tag{VC7}
\]

Thus A=h\_VV is genuinely unitary for the actual positive form on V. The terms of hy in P can be nonzero; VC7 is a compression identity, not pointwise fixation of V. In particular, at every complex embedding,

\[
          \|A\|_{\rm op}=1,\qquad\|A-I\|_{\rm op}\le2.   \tag{VC8}
\]

The inverse-plane hypothesis matters. For example, in the proof-only real model with ordered coordinates S,V,P\_+,P\_- and form diag(1,1,1,−1), let

\[
 n=\begin{pmatrix}
 1&0&-1&1\\0&1&0&0\\1&0&1/2&1/2\\1&0&-1/2&3/2
 \end{pmatrix},\quad
 k=\begin{pmatrix}
 1&0&0&0\\0&3/5&-4/5&0\\0&4/5&3/5&0\\0&0&0&1
 \end{pmatrix}.
                                                               \tag{VC9}
\]

Both are unitary and h=nk has h\_SS=1 and plane-valued C, but h\_SV=−4/5 and h\_VV=3/5. Its inverse cross block is not plane-valued. This example is only a boundary for the algebraic hypothesis; it is not an actual global pencil arrow or a different adopted family.

\subsection{3. Integral short witness and unchanged finite restoration}\label{integral-short-witness-and-unchanged-finite-restoration}

Each y\_a belongs separately to every original source and target pencil at both CM components. If h carries the actual source to the target, hy\_b and y\_a belong to that paired selfdual target. Therefore

\[
                         \Delta_{ab}\in\mathcal O_E.       \tag{VC10}
\]

This is the original arbitrary-arrow argument R070(SC4), not an assertion that h stabilizes one lattice. Since y\_a∈V and P⊥V, VC7--8 now give the stronger exact bound

\[
 |\tau\Delta_{ab}|=|\langle y_a,(A-I)y_b\rangle|
 \le2\|y_a\|\|y_b\|
 \le2\cdot978665066829<2^{41}.                            \tag{VC11}
\]

The old general cutoff bound2\^{}47 is unnecessary in this doubly-plane sector. This change is a bound on the original integral witness; the actual selected trace is retained separately below.

Use the original whole-C ideal J\_C and simultaneous compression ideal J\_Delta from the fully proved paired finite lemma R070(SC7--11). At finite rank d, the deterministic table is:

\begin{longtable}[]{@{}
  >{\raggedleft\arraybackslash}p{(\columnwidth - 6\tabcolsep) * \real{0.2667}}
  >{\raggedright\arraybackslash}p{(\columnwidth - 6\tabcolsep) * \real{0.2000}}
  >{\raggedleft\arraybackslash}p{(\columnwidth - 6\tabcolsep) * \real{0.2667}}
  >{\raggedleft\arraybackslash}p{(\columnwidth - 6\tabcolsep) * \real{0.2667}}@{}}
\toprule\noalign{}
\begin{minipage}[b]{\linewidth}\raggedleft
Finite rank
\end{minipage} & \begin{minipage}[b]{\linewidth}\raggedright
Zero lower C blocks at the two CM components
\end{minipage} & \begin{minipage}[b]{\linewidth}\raggedleft
E-norm exponent in J\_C
\end{minipage} & \begin{minipage}[b]{\linewidth}\raggedleft
E-norm exponent in J\_Delta
\end{minipage} \\
\midrule\noalign{}
\endhead
\bottomrule\noalign{}
\endlastfoot
0 & Both & 2 & 2 \\
1 & Exactly one & 1 at that component & 2 \\
1 & Both & 2 & 0 \\
2 & Neither & 0 & 2 \\
2 & Exactly one & 1 at that component & 0 \\
2 & Both & 2 & 0 \\
At least3 & Any & 0 & 0 \\
\end{longtable}

The relevant normalized Green losses are0,168,335,501 at ranks0,1,2,3. Restoration by167 times the J\_C exponent and167/2 times the J\_Delta exponent dominates each deficit from501. Higher unipotent ranks already have at least that loss. Every assigned J\_C factor divides the complete original C block; every assigned J\_Delta factor divides all four compressions. These are simultaneous conditions proved with the actual inverse-adjoint pairing, not independent choices for each entry.

One nonzero actual Delta\_ab therefore gives N\_EJ\_Delta\textless2\^{}(41·336). There is no additional four-witness factor. The complete finite bound is

\[
 |a_G(h)|\le2Q^{-501\cdot167}(N_EJ_C)^{167}
                                      2^{41\cdot167f}.       \tag{VC12}
\]

There are at most4\^{}334=2\^{}668 choices of J\_C, and N\_EJ\_C≤Q\^{}334. The varying J\_Delta has been bounded uniformly; there is no second pattern count. Both mixed finite factors retain modulus at most one. Global nilpotence is never inferred from finite nilpotence.

\subsection{4. Complete original-coordinate count}\label{complete-original-coordinate-count}

Let L=J\_C I\_Sigma, with N\_EI\_Sigma=T\^{}-52 and z=(N\_EL)\^{}(1/(2f)). The original primitive u extraction at both CM partners gives

\[
 Ce_i=(u\xi_i,\beta_i),\quad \xi_i,\beta_i\in L,\quad
 \|\xi_\tau\|<M:=H/(2\sqrt U),\quad |\beta_\tau|<b:=H/2.
                                                               \tag{VC13}
\]

The joint norm follows from the actual negative row of h and isotropy of C, exactly DO099(JR14). It is not a coordinatewise replacement for that norm. The rational quotient of a nonzero beta pivot by xi reconstructs the whole cross line over E only.

The common-pattern cap and original arbitrary-ideal metric cover give

\[
 z\le Q^{167/f}T^{-26/f}<2^{-84},\quad M<2^{-67},\quad
 \mu(L^{164})<4117\sqrt f\,z.                              \tag{VC14}
\]

The exact integer comparisons are Q\^{}1672\^{}(84f)\textless T\^{}26, H²2\^{}132\textless U and(2250/7)²·164\textless4117². This is the true complex half-trace metric on the entire original row, with covolume D0\^{}164 z\^{}(328f), where D0=7\^{}14029\^{}1445\^{}8413\^{}84. Its full count is

\[
 \#\xi\le A_{164}(M+4117z)^{328f}z^{-328f},\quad
 A_{164}=\frac{528^{164f}}{(164f)!D_0^{164}}.                 \tag{VC15}
\]

The single ball includes all coordinate supports; there is no k union or additional position count. The zero row may be included in the upper count, while reconstruction is asserted only for actual C≠0.

Retain the exact relative norm fibre DO099(JR18), using the actual split-prime constant of DO092(SN10--11) and all28 norm-one generators:

\[
                             \#\beta<2^{249}b^{f/8}z^{-f/8}.
                                                               \tag{VC16}
\]

Its integral norm ideal is formed only for actual solutions; nonprincipal ideals contribute no generator. No integrality of eta=beta/xi is asserted. There remain at most164\textless2\^{}8 pivot choices.

All other factors are exactly R070(SC14--16)/DO099(JR19):

\[
 |c(h)|<2^{-929f/2}U^{-f/2}z^{-f},\quad
 \#\{D e_n\}<2^{5757}T^{10s},\quad G_5<2^{107}.             \tag{VC17}
\]

The first is the actual selected normalized trace, with both compact factors and the near-point representative. The second is the full original five-complex-coordinate target-overlattice ball, uniform in every admitted label including zero. It makes no real-coordinate coarea assertion. The third is the actual projective source-lattice fibre: C fixes the positive163-space ker C, and C together with D e\_n reconstructs the complete negative column before the same rank-five finite automorphism bound is applied. No moving basis is declared integral.

Finally retain the entire original ordered-pair/absolute-incidence sum

\[
                 P_T^{26}<2^{997\cdot26},\qquad
                 P_T=2(T-1)^2/(T-2).                       \tag{VC18}
\]

It already includes every type profile, paired label and local weight. No separate3\^{}26 factor is added or omitted. The fixed binary exponent is

\[
             1+668+168+8+107+5757+249=6958.                 \tag{VC19}
\]

These are respectively the Green correction, common patterns, support count, pivot, source fibre, full target column and relative norm fibre.

\subsection{5. One exact endpoint and the complete charge}\label{one-exact-endpoint-and-the-complete-charge}

Since N\_EJ\_C=T\^{}(2s)z\^{}(2f), combining VC12 and VC15--17 leaves

\[
 z^{819}(M+4117z)^{328f},\qquad
 819=334f-328f-f/8-f=39f/8>0.                               \tag{VC20}
\]

It increases for z\textgreater0. The single endpoint VC14 gives

\[
                 M+4117z<135189/2^{84},
\] \[
 2^{-4069114}\le
 A_{164}(135189/2^{84})^{328f}<2^{-4069113}.                 \tag{VC21}
\]

The interval in VC21 concerns the explicit rational22/7 volume envelope. The true volume is bounded above by it; no equality with the pi-volume is asserted.

Use the actual directed logarithms

\[
 L_T=1019333/1024>\log_2T,\quad L_Q=287843/4096<\log_2Q,
\quad L_U=601356/4096<\log_2U,\quad L_H=28900/4096>\log_2H.
\]

The full logarithmic exponent, with every factor above retained, is

\[
\begin{split}
 E_{26}={}&6958-929f/2+41\cdot167f+(f/8)(L_H-1)
          -(f/2)L_U-501\cdot167L_Q\\
         &+(344L_T+997)26-4069113-819\cdot84\\
        ={}&-87517717/4096<-21366.                           \tag{VC22}
\end{split}
\]

The T power344s is334s from the whole-C numerator and10s from the original target column. Both the beta norm fibre and selected trace have already been included in the z exponent819. There is one complete sector, one row endpoint and no extra union over supports, inverse directions, witness choices, rank distributions or k. This proves VC4.

\subsection{6. The resulting support26 full-positive identity}\label{the-resulting-support26-full-positive-identity}

For completeness, the structural consequence is established with the actual original integral vector, rather than importing a prior numerical support range. Outside VC4, all four Delta\_ab vanish. By VC7, A is unitary on positive V. The vanishing2×2 compression means that for y∈Y=span(y1,y2), the Y projection of Ay is y. Equality of norms forces its component perpendicular to Y to vanish. Thus A fixes Y pointwise. This says only that the V compression fixes Y; hy can still differ from y by a nonzero vector of ell in P.

Use the original integers and rational vector of DO094/R071:

\[
 M_0=u_1^2+u_2^2+u_3^2=155821585821105090045582,
\] \[
 z_0=(256766,-7654233287,-30805128044,1),\quad
 t_0=(-u_4(u_1,u_2,u_3)/M_0,1).                             \tag{VC23}
\]

Here z0∈V is integral, z0−t0∈Y, t0⊥Y and \textbar\textbar t0\textbar\textbar²=U/M0. Since A is unitary and fixes Y, At0=alpha t0 with alpha bar(alpha)=1. The original two-lattice pairing gives

\[
 \delta_z=\langle z_0,(h-I)z_0\rangle
          =(U/M_0)(\alpha-1)\in\mathcal O_E.               \tag{VC24}
\]

The vector z0 belongs separately to every actual pencil: at T it is (y3−y1+y2)/2 with2 a unit, and away T its original integer coordinates suffice. No rational t0 is declared integral.

The exact comparisons are gcd(U,M0)=9 and M0²\textgreater36U. Substituting alpha=1+(M0/U)delta\_z into its norm-one identity gives

\[
 U(\delta_z+\bar\delta_z)=-M_0\delta_z\bar\delta_z.
\]

Coprimality of U/9 and M0/9, by integer Bezout, makes nu=delta\_z bar(delta\_z)/(U/9) an algebraic integer in F. Every real conjugate is nonnegative and at most36U/M0²\textless1. If nonzero, its positive integral field norm would be less than1. Thus nu=0, delta\_z=0 and alpha=1, exactly the support-independent R071(GR1--10) argument.

It follows that A=I\_V. By VC6, N maps S⊕V into the single nonzero line ell. Consequently

\[
 h_{S\oplus V,S\oplus V}=I,\qquad
 \dim_E(\ker(h-I)\cap(S\oplus V))=166,\qquad
                         \operatorname{rank}_E(h-I)\le3. \tag{VC25}
\]

No individual y\_a or z0 is asserted fixed by h. The new compression identity, the actual integral scalar and the positive166-space all concern the original arrow. No finite positive metric or global nilpotence was inferred.

\subsection{7. Verification, accounting and source boundaries}\label{verification-accounting-and-source-boundaries}

The companion certificate checks the exact original short norms and orthogonality, the paired finite restoration deficits, all retained dimension and multiplicity exponents, the actual T/Q/U/H directed bounds, the fractional-ideal cap, the full target and scalar-row volume ceilings, the single exact exponent, and the original scalar-rigidity integers. Reduced exact unitary models check the block lemma and the necessity of the inverse-plane premise; these are algebraic diagnostics, not asserted source-admissible global arrows.

The current canonical predecessor is checks/do110\_joint\_remainder\_error.json, with108 all-top and31 uniform entries. The executable accepts a later complete extension of the actual R089 prefix and records the supplied predecessor. Append21366 once, retaining all earlier errors and both complete floors26/11. The final author and peer records identify the exact actual chronology used; an historical replay is not silently relabeled as a later ledger check. Both entire dyadic sums are checked below2\^{}-230.

Direct mathematical inputs are the full R070(SC1--25), DO099(JR1--24), R061(GS11--15), DO070's original fractional-ideal cover, DO092(SN10--11), DO098(CF20--21), and R071(GR1--12). The new step is the genuine positive-compression unitary identity VC5--8, its sharper original integral witness VC11, and the fully retained support26 evaluation VC22. The actual selected trace is unchanged. There is no new representation, field, source constructor, projector, lattice normalization or Haar convention.

The original Hermitian Hom modules and integral pairing sections remain over Spec(O\_E); field norms remain determinants of multiplication and the fixed space is the kernel of the actual original linear map. Polarization, Picard, theta-line and biextension data are unchanged. The inherited contextual anchors remain SGA3 I4.7.3/4.7.3.1; SGA6 VI Theorem3.3; FGA232 Theorem3.1 pp148--149; FGA236 Corollaries1.7/2.3 pp229/232; SGA7 I VII2.7/2.9 pp169--174 and VIII7.4(b) pp310--311; cote4 scan003 §§8.1.1--8.1.2, scans005--006 §§8.4.1--8.4.3; and cote54 scans005--007. These are carried context, with no new reading or comparison theorem asserted here. The signed sufficient inequality, prescribed period, selected arithmetic image, source admissibility, specific032 value, DP44 comparison and RH remain separate.

\subsection{Canonical algebraic-geometric context}\label{canonical-algebraic-geometric-context-26}

The reductions are the actual integral Hermitian modules at unmixed finite points; content and coordinate ideals remain fractional invertible modules over Spec(O\_E) and Spec(O\_F), with the same paired CM action, polarization, theta line, Picard scheme and biextension. Carried exact anchors: SGA3 I4.7.3/4.7.3.1; SGA6 VI Theorem3.3; FGA232 Theorem3.1 pp148--149; FGA236 Corollaries1.7/2.3 pp229/232; SGA7 I VII2.7/2.9 pp169--174 and VIII7.4(b) pp310--311; cote4 scan003 sections8.1.1--8.1.2 and scans005--006 sections8.4.1--8.4.3; cote54 scans005--007. No new lambda-ring/K-theory, l-adic weight or CM-motive comparison is asserted. DP\_MULT Q3e, THETA3 G1/A7, CX\_GRH CX-ALG-024 and GCT\_PLETH T-052 retain their stated global gates.

\clearpage

\section{AM-LIFT-088 --- an explicit field for the complete general-Host CM seed}\label{am-lift-088-an-explicit-field-for-the-complete-general-host-cm-seed}

\textbf{Canonical checked integration --- 2026-11-23T04:38:48.463885+00:00.} Complete saved-proof and executable reviews and separate exact diagnostic replays passed. No numerical error charge is introduced. Frozen author history is preserved. AM-LIFT-088 SF1--16 supplies an explicit field for the entire newly chosen LIFT087 general-Host seed over the same given Coeff field B: {[}K\_seed:B{]}\textless=168\emph{23040}h(-d)\emph{M\^{}(4g\^{}2)\textless=3870720}d*M\^{}(4g\^{}2), g=f(p+q). The actual prescribed elliptic CM action descends over B(j) by scalar differential injectivity. The explicit J good-model radicals, elliptic15-torsion, actual projective Serre modules, equivariant principalization, canonical trace-one cubical line, original finite M-level and literal Obj/FrPt all lie over the stated field. Under every retained LIFT085 nonzero-rule/source hypothesis, {[}K\_out:B{]} is at most that bound timesN\^{}(8g\^{}2), with the existing intrinsic Step5 cost unchanged. EffectiveN, absolute degree of the givenB, field discriminants, coefficient/section heights, numerical Host/Sizes existence, selected arithmetic content, prescribed period, specific032/Read, evaluated DP44 and RH remain open. No finite error charge; current110/31 arrays and complete floors26/11 are unchanged.

Author: \nolinkurl{/root/lift_accounting}, 23 November 2026 UTC. \textbf{Author-frozen proof; the final whole-file, literal-source and complete-executable checks are recorded in the adjacent verification packet. Independent review and canonical promotion are separate.} This is a new consequence for the seed chosen in LIFT087. No frozen earlier proof, canonical record or error ledger is changed.

\subsection{1. Actual data and the relative field bound}\label{actual-data-and-the-relative-field-bound}

Retain the actual Host, its quadratic field \(k=\mathbf Q(\sqrt{-d})\), its specified ideal \(\mathfrak c_K\), the stage Datum, the compatible Pins, the existing Moduli representation and the same Coeff record used in LIFT084--087. Thus d is a prime congruent to7 modulo16, \(\operatorname{disc}k=-d\), and \(B=c.F\subset\mathbf C\) is the given number field. The source Coeff.conj puts every conjugate of E in B. In particular the prescribed copy of k lies in B. The actual finite level is

\[
                  M=g_0g_1\ge3,\qquad A=1.                 \tag{SF1}
\]

Write \([E:\mathbf Q]=2f\), and let \(A_0\) be the newly chosen principal seed of LIFT087 for this stage. Retain that construction's signature hypotheses, in particular p\textgreater=q and p+q\textgreater0. Its abelian dimension is \(g=f(p+q)\); this is different from the moduli dimension \(fpq\) appearing in \nolinkurl{Datum.dim}. Let \(h=h(-d)\) be the class number of the maximal order of k.

\textbf{Theorem.} The complete seed, with its actual good polarized model, full integral action, canonical trace-one rigidified symmetric cubical line, prescribed finite level and a literal field-valued FrPt, can be defined over a finite extension \(K_{\rm seed}/B\) satisfying

\[
 [K_{\rm seed}:B]\le168\cdot23040\,h\,M^{4g^2}
       \le3\,870\,720\,d\,M^{4g^2}.                         \tag{SF2}
\]

The source Host, Datum, Coeff and Level records are unchanged. This is a choice of a seed and its auxiliary field, not a bound for every model of the same geometric variety or for an arbitrary prescribed earlier descent field. It retains the new polarized seed and intrinsic comparison lines of LIFT087, rather than identifying that seed with the earlier R086 polarization.

Under every additional hypothesis of the fixed-level nonzero-rule construction LIFT085 and the cost transfer LIFT086, the target can subsequently be defined over

\[
 K_{\rm out}=K_{\rm seed}(A_0[N^2]),\qquad
 [K_{\rm out}:B]\le3\,870\,720\,d\,M^{4g^2}N^{8g^2}.        \tag{SF3}
\]

N is the same unevaluated integer isogeny scale in those theorems. Equations SF2--3 bound relative degrees. Absolute degrees require multiplication by \([B:\mathbf Q]\). They give no bound for the discriminant of either field, the heights of its defining coefficients, N, the prescribed arithmetic content, a period or the DP44 expression.

\subsection{2. The prescribed elliptic factor and its full CM action}\label{the-prescribed-elliptic-factor-and-its-full-cm-action}

Let \(j_* = j(\mathbf C/\mathfrak c_K)\). The Hilbert class polynomial of the maximal order has integer coefficients, degree h, and has this j-invariant as a root. Hence

\[
                         [B(j_*):B]\le h.                  \tag{SF4}
\]

The input is Sutherland's Theorem20.12 and the preceding ideal-class torsor description (2023, pp.5--6). We do not need irreducibility, an equality with the Hilbert class field, or a choice of a different ideal class.

Choose an elliptic curve C0 over \(B(j_*)\) with this j-invariant. One explicit model, when \(j_*\notin\{0,1728\}\), is

\[
 y^2=x^3-\frac{3j_*}{j_*-1728}x-\frac{2j_*}{j_*-1728}.
                                                               \tag{SF5}
\]

Its discriminant is nonzero and the formula \(j=1728(4a^3)/(4a^3+27b^2)\) gives exactly \(j_*\). If either exceptional j-invariant were used, the usual models \(y^2=x^3+1\) and \(y^2=x^3-x\) supply it over the same field; no extension is required. Choose a complex origin-preserving isomorphism with \(\mathbf C/\mathfrak c_K\), thereby fixing the prescribed geometric \(\mathcal O_k\)-action and the principal polarization. The latter is the canonical elliptic principal polarization, so it is defined over the model field without a theta-line choice.

Here the \textbf{entire} \(\mathcal O_k\)-action descends over \(B(j_*)\). First, every geometric endomorphism of C0 is defined over \(\overline{\mathbf Q}\). For a nonzero CM endomorphism a, its finite kernel lies in C0{[}n{]} for the positive integer \(n=a\bar a\). The kernel is algebraic. Its quotient and an isomorphism from that quotient to C0 are defined over \(\overline{\mathbf Q}\), so the endomorphism is algebraic there. Equivalently one may use full faithfulness for homomorphisms of abelian varieties over algebraically closed extensions.

Differentiation is injective on the geometric endomorphism ring of an elliptic curve in characteristic zero: a nonzero endomorphism is an isogeny, is separable, and therefore has nonzero differential. For the fixed action, the derivative of a is the prescribed scalar \(\rho(a)\) in the one-dimensional Lie space. That scalar belongs to k and thus to \(B(j_*)\); this assertion is independent of a basis of the Lie line. If \(\sigma\) fixes \(B(j_*)\), then

\[
           d(a^\sigma)=\sigma(\rho(a))=\rho(a)=da,
                   \qquad a^\sigma=a.                     \tag{SF6}
\]

Galois descent now defines every a over that same field. All ring and adjoint identities hold there because they hold after faithfully flat base change. This is simultaneous descent of the full action, not an extension for a finite list followed by an unproved assertion about all a. Upon applying any embedding of the resulting field into C, the derivatives become the corresponding conjugate embedding of \(\rho(a)\); no independent reassignment of the CM labels or transcendental period is made.

For completeness an elementary sufficient class-number bound is

\[
                              h(-d)\le d.                  \tag{SF7}
\]

Indeed the quadratic Minkowski theorem gives an integral representative in every ideal class of norm at most \(2\sqrt d/\pi\), hence at most \(r=\lfloor\sqrt d\rfloor\). In a quadratic field the number of ideals of norm n is at most the ordinary divisor number \(d_2(n)\): at a split prime the coefficient for exponent e is e+1, at an inert prime it is zero or one, and at a ramified prime it is one. Thus it is also at most n.~Consequently \(h\le\sum_{n=1}^r n=r(r+1)/2\le r^2\le d\), since r\textgreater=1. This uses Milne's Theorem4.3 and its proof (2020b, pp.70,80--81), with elementary unique ideal factorization. No analytic class-number estimate is needed.

\subsection{3. One explicitly bounded field of good reduction}\label{one-explicitly-bounded-field-of-good-reduction}

Put \(K_j=B(j_*)\), take \(\zeta_7\) primitive, and choose

\[
 i^2=-1,\quad \alpha^7=-4,\quad
 \beta^2=(\zeta_7-1)^7,\qquad
 K_J=K_j(\zeta_7,i,\alpha,\beta).
                                                               \tag{SF8}
\]

The successive defining polynomials have degrees at most6,2,7,2. Therefore \([K_J:K_j]\le168\), irrespective of intersections or elements already present in B. No degree of E or of its normal closure is added here: all conjugates of E are already in B by Coeff.conj, and adjoining \(\zeta_7\) supplies every conjugate label of \(L=E\mathbf Q(\zeta_7)\).

LIFT081 ED3/ED8--10 gives the actual everywhere-good model of \(J=\operatorname{Jac}(y^2=x^7-1)\) over this field, with full \(\mathcal O_{\mathbf Q(\zeta_7)}\)-action. In particular the characteristic2 chart is \(Y^2+Y=X^7\), with \(x=\alpha X\), \(y=i(2Y+1)\), while the characteristic7 chart uses \(x=1+(\zeta_7-1)X\), \(y=\beta Y\). The infinity charts and saturated differential bases are part of the carried theorem. Thus SF8 covers the complete good model and its action, not merely the displayed affine equations.

Now define

\[
                         K_{\rm good}=K_J(C_0[15]).         \tag{SF9}
\]

The finite Galois action embeds into \(\operatorname{GL}_2(\mathbf Z/15)\), whose order is

\[
 |\operatorname{GL}_2(\mathbf F_3)|\,
 |\operatorname{GL}_2(\mathbf F_5)|
       =(9-1)(9-3)(25-1)(25-5)=23040.
                                                               \tag{SF10}
\]

C0 has potential good reduction because it has geometric CM. Over SF9 it has good reduction everywhere by the precise finite-inertia argument of LIFT086 IM11--12: at residue characteristic different from3 use its rational3-torsion, and at characteristic3 use its rational5-torsion. Finite inertia is then contained in the torsion-free odd-prime principal congruence subgroup, hence is trivial. This proves

\[
                       [K_{\rm good}:B]\le168h\cdot23040.
                                                               \tag{SF11}
\]

The full elliptic action extends uniquely to the good model. The construction does not require the field to be CM, Galois over Q, or unramified at the original bad primes. Good reduction persists under the subsequent finite extensions.

\subsection{4. No additional field for the projective modules or principalization}\label{no-additional-field-for-the-projective-modules-or-principalization}

All the projective modules here are the actual modules in LIFT087:

\[
 X=\mathcal O_L\otimes_{\mathcal O_{K7}}J,
       \qquad C=\mathfrak b^{-1}\otimes_{\mathcal O_k}C_0.
                                                               \tag{SF12}
\]

They are finite projective over their indicated maximal orders. Covariantly, \(M\otimes_R A\) means \(\operatorname{Hom}_R(M^\vee,A)\). Milne's Proposition7.33 and Remark7.34 construct this over the \textbf{same} base and commute with arbitrary base change. One may realize M by an idempotent over a finite free R-module; the entries act by already defined integral endomorphisms, and its summand is an abelian scheme. A chosen presentation or a Steinitz ideal class therefore entails no extension of the coefficient field and no assumption that M is free.

Multiplication by \(\mathcal O_L\) on the first module, or by \(\mathcal O_E\) on the second, consists of module maps over R. They induce the full corresponding action over \(K_{\rm good}\), and over its ring of integers. The trace-Hermitian forms used in LIFT087 are maps of those modules; their polarization homomorphisms are defined over the same base. In the pure case the actual identity \((\mathfrak b^{-1})^\#=\mathfrak b^{-1}\) makes the polarization principal. The ideal \(\mathfrak b\), the class \(\mathfrak c_K\), and their actual trace form have not been replaced.

Apply the equivariant Zarhin construction to X. Its target is \((X^\vee)^4\times X^4\); its isogeny and principal polarization are defined over the starting field by Theorem1.1 and the explicit equations16--19. The integral quaternion coefficients are integers. Its dual action is the specified conjugate-dual action, and \(\pi^\vee\mu\pi=\lambda_X^8\). Consequently there is no extension to choose an isotropic subgroup, no extra root needed for the polarization, and no new metric scalar. The actual variety is already a product of X and its dual, so it has a good model over the same ring; the polarization and full action extend uniquely. Action twists, conjugation of the E-action and the finite products giving the R086/LIFT087 signatures are also performed over this base.

This constructs the chosen seed \((A_0,\lambda_0,\iota_0)\) and its good model over \(K_{\rm good}\). Every conjugate of every L label lies there. In particular the projectors and algebraic component eigenvectors used for the intrinsic bounds in LIFT087 may be chosen over this field. This does not assert that an arbitrary integral basis or a theta section of small height has been found.

\subsection{5. The actual line, level and complete source object}\label{the-actual-line-level-and-complete-source-object}

A polarization homomorphism alone need not specify a line over its base field. Here use exactly the trace-one line of LIFT085 HF5--7 and LIFT086 IM13--15. For the normalized Poincare line \(\mathscr P\), put

\[
 \mathscr B=(1\times\lambda_0)^*\mathscr P,
       \qquad \mathcal L_0=(1,\iota_0(w_0))^*\mathscr B.
                                                               \tag{SF13}
\]

The source has \(w_0+\bar w_0=1\). Thus SF13 has polarization \(\lambda_0\) and is the actual origin-rigidified symmetric cubical line. The Poincare construction is over \(K_{\rm good}\), including its good model. The normalized all-a adjoint and primitive-line identities descend uniquely over that same field. No two-torsion twist, square root, or infinite sequence of additional fields is involved. This supplies the literal SAB line, symmetry, cube, adj and prim fields.

Adjoin only the actual M-torsion:

\[
                         K_{\rm seed}=K_{\rm good}(A_0[M]).
                                                               \tag{SF14}
\]

Its Galois group embeds into \(\operatorname{GL}_{2g}(\mathbf Z/M)\), so its degree is at most \(M^{4g^2}\). The profinite isometry proved in LIFT084 matches the fixed integral Datum lattice to the actual polarized seed lattice, with the same \(\mathcal O_E\)-action. Choose \(e\in\mathcal O_E\) satisfying \(e=1\bmod\mathfrak M\) and \(e=0\bmod\bar{\mathfrak M}\). Then

\[
               \eta(a)=\operatorname{ev}_M(ea),\qquad
               \eta_A=0.                                   \tag{SF15}
\]

Every value of eta belongs to A0{[}M{]}, so SF14 defines the entire map over one field. The infinitely indexed map has a finite quotient domain. LIFT084 MP19 proves both exact kernel and complementary-ideal agreement clauses; M is not replaced by M squared or by a new prime. The root condition at A=1 is automatic. The source TateIso tower at all positive indices is required separately over every algebraically closed geometric field. That geometric tower is supplied by the retained lattice isometry and its actual theta commutator sign; it is not asserted to descend as an infinite full level over SF14.

The same seed satisfies the exact DetCond with the prescribed multiplicities, and its principal ample line has the required KTriv/section property. Hence this is an object of the literal Obj moduli problem over the reduced base \(\operatorname{Spec}K_{\rm seed}\). Its good integral model is a separate abelian-scheme statement; we do not assert an integral Obj where Base.unit requires all positive integers to be units.

\subsection{6. A field-valued FrPt without another extension}\label{a-field-valued-frpt-without-another-extension}

The actual \nolinkurl{Moduli.rep} clause applies to this \textbf{object over this field}, not merely its geometric isomorphism class. It supplies a morphism \(\operatorname{Spec}K_{\rm seed}\to X_o\) and the cubical cartesian ObjMap to the universal family over the same base. There is no new field-of-moduli versus field-of-definition step.

Because Coeff.conj splits E over B, the tangent space is a direct sum of K\_seed-vector eigenspaces of the actual multiplicities. Choose a basis in each eigenspace over K\_seed; equivalently retain the constructed LIFT087 component eigenvectors. These give the precise tan, eig and unique-combination basis clauses of FrPt. The actual principal ample line has a one-dimensional space of sections over K\_seed: proper base change identifies its scalar extension with the one-dimensional geometric space. Choose a nonzero vector in that K\_seed-line. It supplies the FrPt.s/gen clauses without adjoining its geometric coordinates from an arbitrary prechosen generator.

All these statements concern the actual selected field-valued frame. They do not assert small coefficients for an arbitrary frame or a specified scalar normalization of the theta section. Applying any embedding of K\_seed transports the full object, its algebraic actions, frame and generator together. In particular no separate field extension per CM label or per complex embedding has been counted or assumed.

Combining SF4, SF8, SF10 and SF14 proves SF2. For a finite list of stages one may choose the same elliptic and Jacobian ground field. Adjoining their torsion fields gives the sufficient common bound \(168h\cdot23040\,M^{4\sum_j g_j^2}\); this is just the product of the individual Galois-image bounds, not an assertion that those torsion fields are disjoint.

\subsection{7. Subsequent nonzero point and retained arithmetic costs}\label{subsequent-nonzero-point-and-retained-arithmetic-costs}

If the actual natural rule is nonzero and every geometric/source hypothesis of LIFT085 holds, run its fixed-level orbit argument from this seed. Because the seed already has good reduction, its complete target and model are defined over SF3, with no second factor15. The scale N remains selected qualitatively by that argument. Neither SF2 nor the finite degree of an algebraic point bounds the isogeny denominator required to reach the rule's nonvanishing locus.

The LIFT087 intrinsic lower bound b is unchanged by these finite base changes. Under the retained LIFT082/LIFT086 hypotheses, the sufficient Step5 cost remains

\[
 \frac{D+gn}{2}+D\log r+
  (b+\log N)\min\{D+gn,(3D+gn)/2\}.                         \tag{SF16}
\]

The table records the exact change in the pending obligations.

\begin{longtable}[]{@{}
  >{\raggedright\arraybackslash}p{(\columnwidth - 4\tabcolsep) * \real{0.3333}}
  >{\raggedright\arraybackslash}p{(\columnwidth - 4\tabcolsep) * \real{0.3333}}
  >{\raggedright\arraybackslash}p{(\columnwidth - 4\tabcolsep) * \real{0.3333}}@{}}
\toprule\noalign{}
\begin{minipage}[b]{\linewidth}\raggedright
Item
\end{minipage} & \begin{minipage}[b]{\linewidth}\raggedright
Supplied here
\end{minipage} & \begin{minipage}[b]{\linewidth}\raggedright
Retained separate obligation
\end{minipage} \\
\midrule\noalign{}
\endhead
\bottomrule\noalign{}
\endlastfoot
Initial seed field & SF2, including full actions, principalization, canonical line, finite level and frame & Absolute degree of the given B; all field discriminants and height bounds \\
Good model & Same-field construction from the explicit J model and C0{[}15{]} & Compatibility with any independently specified integral arithmetic input \\
Prescribed nonzero rule & Conditional target field SF3 through LIFT085 & Actual rule nonzero, component/source hypotheses, and effective N \\
Intrinsic lines & The same actual LIFT087 bounds and LIFT086 transport & All KeyCore, local integrality, cusp and comparison hypotheses \\
Arithmetic/global chain & No selected value or error charge is introduced & Host/Sizes constructor, selected content, period, Read/STiesR/032 and evaluated DP44 \\
\end{longtable}

\subsection{8. Sources and verification boundary}\label{sources-and-verification-boundary-4}

Fresh literal reads: \nolinkurl{TwoWeightEqOneAddRoot.lean} Host14--27, Datum30--46, Level50--65, SAB72--84, Base124--129, TateIso131--142 and Obj144--160; \nolinkurl{PushInvertible.lean} Moduli29--38 and Coeff40--46; \nolinkurl{IsCombAt.lean} FrPt67--76 and its surrounding exact tangent/eigenvector definitions. The source manifest records their actual bytes and the other carried interfaces. LIFT081 ED3/ED8--10, LIFT084 MP16--19, LIFT085 HF5--9, LIFT086 IM2/IM11--16 and the complete LIFT087 construction are retained with their checked scopes.

Primary references: Sutherland, A. V. (2023, November21), \emph{18.783 Elliptic Curves, Lecture20}, pp.5--6, Theorem20.12 and its displayed proof, \url{https://math.mit.edu/classes/18.783/2023/LectureNotes20.pdf}; Milne, J. S. (2020a), \emph{Complex multiplication}, v0.10, Proposition7.33/Remark7.34/Proposition7.36, pp.62--63, and the carried good-reduction locators, \url{https://www.jmilne.org/math/CourseNotes/CM.pdf}; Milne, J. S. (2020b), \emph{Algebraic number theory}, v3.08, Theorem4.3 p.70 and proof pp.80--81, \url{https://www.jmilne.org/math/CourseNotes/ANT.pdf}; Zarhin, Y. G. (2023), \emph{Abelian varieties, quaternion trick and endomorphisms}, Theorem1.1 and equations16--19 in the author v5, pp.4--6, \url{https://arxiv.org/pdf/2010.07995}.

The companion exact diagnostics check the explicit elliptic model, class-number majorants, finite torsion-group cardinalities, field-factor arithmetic, all-conjugate scalar descent models, projective-module/idempotent constructions, finite-level CRT clauses and all frozen source fingerprints. These diagnostics supplement the universal written proof. They neither enumerate all CM fields nor prove a new moduli representability theorem. Final whole-file review, actual replay time, subject hashes and independent reviews are separate records. No canonical files, numerical remainder arrays or main-goal statuses are changed here.

\subsection{Canonical algebraic-geometric context}\label{canonical-algebraic-geometric-context-27}

The actual abelian schemes, integer-ring action, principal Picard/Poincare construction, normalized cubical line, finite-flat isogeny and its kernel live over Spec(O\_K). Away from N the displayed adjoint inverse gives an actual p-divisible-group and formal-group/Dieudonne comparison, without an operator or slope gain. Carried exact anchors: FGA232 Theorem3.1 pp148--149; FGA236 Corollaries1.7/2.3 pp229/232; SGA7 I VII2.7/2.9 pp169--174 and VIII7.4(b) pp310--311; SGA3 I4.7.3/4.7.3.1; SGA6 VI Theorem3.3; cote4 scan003 sections8.1.1--8.1.2 and scans005--006 sections8.4.1--8.4.3; cote54 scans005--007. No new lambda-ring/K-theory, l-adic weight or CM-motive identification. DP\_MULT Q3e/DP44/DP52/DP54--65, THETA3 G1/A7, CX\_GRH CX-ALG-024 and GCT\_PLETH T-052 retain their separate arithmetic comparisons. Precise primary locator clarification: Zarhin authorv5 Theorem1.1 is printedp2; section4 construction is pp4--6 and equations16--19 with equivariance are pp5--6. The compressed page range in the frozen paragraph is read with this exact split. The same-field Serre construction is Milne CM2020v0.10 Proposition7.33/Remark7.34/Proposition7.36 pp62--63. The initial elliptic degree uses Sutherland2023 Lecture20 Theorem20.12 pp5--6; its elementary class-number majorant specializes Milne ANT2020v3.08 Theorem4.3 printedp70 and its proof pp80--81.

\clearpage

\section{AM-R-091 --- the exact remaining support26 branch}\label{am-r-091-the-exact-remaining-support26-branch}

\textbf{Canonical checked integration --- 2026-11-23T04:54:10.929679+00:00.} Complete saved-proof and executable reviews and separate exact diagnostic replays passed. No numerical error charge is introduced. Frozen author history is preserved. AM-R-091 DR1--7 consolidates every checked support26 cut through DO111 into one actual-remainder theorem. The only possible support26 terms are all-type-II globally rank-two-unipotent scalar-cross matrices with full positive167 compression, nonzero original S cross, v=q26, b=j0, strict denominator window2\^{}1116\textless N\_E D\_eta\textless2\^{}2108 and96\textless=k\textless=164. The actual retained support26 set and actual higher-support remainder split the unchanged energy exactly, with the complete110/31 errors below2\^{}-230 and floors26/11. Every earlier cut and overlap charge is retained. No charge or independent new movement is assigned to this consolidation. DR5, actual fully sized source, selected arithmetic content, prescribed period, specific032/Read, evaluated DP44, Kummer-Hecke RH and global oscillator comparison remain open.

Author: /root, GPT-6 Astra Pro; reasoning effort not externally exposed; ALP\_MACHINE. 23 November2026 UTC. Whole saved-proof and executable readings, exact diagnostics and independent review are recorded separately before promotion. Earlier frozen findings are unchanged.

\subsection{1. Actual remainder and theorem}\label{actual-remainder-and-theorem}

Keep the actual J56 all-top upper-endpoint kernel of K044, its charged H=133 cutoff, f=a=168, n=169, original positive S of dimension164, original field and integral Hermitian lattice, both CM partners, every admitted source and target pencil and label, the selected real coefficient and Hilbert metric, the packet point, finite projectors, compact factors, central representative and phases, Haar measures and class masses. Let w(h) be that same normalized summand. R087 J26.1 already proves a complete support floor26. Apply the subsequent checked containing-set removals chronologically through DO111, always to their intersections with the actual remainder. Denote the resulting actual set by R. This definition keeps every earlier cut at every support; it does not replace R by a convenient larger set.

Define its disjoint pieces \[
 D_{26}=\{h\in R:s(h)=26\},\qquad
 R_{\ge27}=\{h\in R:s(h)\ge27\}.                         \tag{DR1}
\] In the original S perpendicular W0 decomposition write h={[}{[}A,B{]},{[}C,D{]}{]}, with W0 of signature(4,1). Let P=Eu perpendicular Ee\_n be the original packet plane and Z=P-perpendicular=S perpendicular u-perpendicular, of dimension167. For the normal form below, let k count the nonzero original S covector values; D\_eta is the exact integral E-denominator ideal of eta and x=N\_E D\_eta. Let v count eligible active places, q eligible places among all168 selected T primes, b inactive paired-C-zero places and j active places at which the entire paired C is integral. These are the definitions of K103 SC26.2--4, not interchangeable rank or incidence variables.

\textbf{Theorem.} Every h in D26 necessarily satisfies all of the following: \[
\begin{gathered}
 A=I,\quad C|_S\ne0,\quad h_{ZZ}=I,\quad(h^{-1})_{ZZ}=I,
 \quad\det h=1,\\
 h=I+N,\quad\operatorname{rank}_E N=2,\quad
 N^3=0,\quad\operatorname{rank}_E N^2=1,\\
 \text{all26 active places have type II},\qquad
 v=q=26,\quad b=j=0,\\
 2^{1116}<x<2^{2108},\qquad96\le k\le164.               \tag{DR2}
\end{gathered}
\] The original and inverse cross directions lie in P. In the Z perpendicular P coordinates the unchanged actual normal form is \[
 N=\begin{pmatrix}0&-a\ell^*J_P\\\ell r&d\ell\ell^*J_P\end{pmatrix},
 \quad\ell=(1,\eta)^t,\quad\eta\bar\eta=U,\quad
 J_P=\operatorname{diag}(U,-1),\quad a=J_Z^{-1}r^*,
\] \[
 d+\bar d=-rJ_Z^{-1}r^*,\quad
 U rJ_Z^{-1}r^*\le17424/133\text{ at every embedding},
 \quad |\operatorname{Im}(-Ud)|<133/2.                    \tag{DR3}
\] The167 original covector values, eta and d reconstruct h injectively. At each fixed support eta is unique in the stated window, each active place has at most one actual type-II label, and all168 factors (T-1)\^{}(-1) of the retained pointwise/weighted product remain. Eligibility is eta congruent to Q(-7+2sqrt(-7)) at the selected CM prime, with the true partner condition from its norm equation. Neither a second independent partner condition nor a sum over additional eta is inserted.

The actual energy identity therefore has the form \[
 \mathcal E_{\rm top}=1+\operatorname{Re}
  \left(\sum_{h\in D_{26}}w(h)+\sum_{h\in R_{\ge27}}w(h)\right)
  +\epsilon,\qquad |\epsilon|<2^{-230}.                  \tag{DR4}
\] The sufficient finite/real inequality is \[
 \operatorname{Re}\left(\sum_{h\in D_{26}}w(h)
                 +\sum_{h\in R_{\ge27}}w(h)\right)
                              >-1+2^{-230}.             \tag{DR5}
\] DR5 is open. DR2 is a necessary containing description of D26, not a claim that every matrix satisfying DR2 occurs, that D26 is nonempty, or that the complex coefficients have one sign. R\textgreater=27 retains all earlier restrictions in their proved scopes. This assembly adds no error charge, keeps the complete floors26/11 and consolidates already accepted reductions; the act of consolidation is not new independent movement.

\subsection{2. Exhaustion of every nonidentity positive compression}\label{exhaustion-of-every-nonidentity-positive-compression}

The unchanged K068 LR3/LR21--22 finite-unipotent remainder has rank\_E(A-I)\textless=33 and determinant-one representatives. K069 IC1--8, as used with its full scope in R087 section2, partitions A!=I into the original fixed class F and exceptional isotropic-cross class IC8. All334 unmixed G reductions are unipotent; global unipotence is not assumed here. The accepted complete ranges are:

\begin{longtable}[]{@{}
  >{\raggedright\arraybackslash}p{(\columnwidth - 4\tabcolsep) * \real{0.3333}}
  >{\raggedright\arraybackslash}p{(\columnwidth - 4\tabcolsep) * \real{0.3333}}
  >{\raggedright\arraybackslash}p{(\columnwidth - 4\tabcolsep) * \real{0.3333}}@{}}
\toprule\noalign{}
\begin{minipage}[b]{\linewidth}\raggedright
Original class
\end{minipage} & \begin{minipage}[b]{\linewidth}\raggedright
Inclusive support coverage
\end{minipage} & \begin{minipage}[b]{\linewidth}\raggedright
Exact run locator
\end{minipage} \\
\midrule\noalign{}
\endhead
\bottomrule\noalign{}
\endlastfoot
F, rank1 & 18..94 & DO086 FR1--10 \\
F, rank2 & through110 & R068 FA1--9 \\
F, rank3 & through154 & R068 FA1--9 \\
F, ranks4..33 & through168 & R068 FA1--9 \\
IC8, ranks1 or2 & 23..28 & R069 EM1--21 \\
IC8, ranks3..31 & R075 complete contiguous table, minimum final ceiling28 & R075 NR1--16 with its stated DO085/DO064 predecessor intervals \\
IC8, ranks32 or33 & through168 & DO064 PM22--23 \\
\end{longtable}

Every interval includes26. In particular R075 supplies the entire interval above its respective predecessor ceiling, not just a checked endpoint, so it leaves no intervening support gap. The table covers all compressed ranks1..33 in both classes. Every retained support26 term consequently has A=I. We use the cited complete theorems at their stated ranges and do not extend any support25-only estimate.

\subsection{3. Exhaustion of the fixed-S alternatives}\label{exhaustion-of-the-fixed-s-alternatives}

For A=I, the original unitary identities R059 GB1--3, retained in R087 section3, give C\emph{J0 C=0 and B=-C}J0 D. If C=0, then B=0 and h=I\_S direct-sum h0 on W0. Otherwise C has nonzero isotropic image. These alternatives are exhaustive and keep the original integral coordinates.

In the fixed-S case, DO090 NE1--6 removes every globally nonunipotent h0 at support26, with no global-rank, characteristic-degree or incidence-profile restriction. If h0 is nonidentity and globally unipotent, DO060 IB21--22 gives exactly rank1 with N²=0 or rank2 with N³=0 and rank N²=1. Rank0 is the identity with support0.

For rank1, R088 GZ1--9 removes the complete integral-witness branch, every intrinsic double-pole count d\textless=22, and every nonidentity unmixed-G pattern. Its only possible remaining branch has no such witness, all G reductions the identity and23\textless=d\textless=26. R089 WN1 and its full weighted-eta proof remove exactly that remaining class. No profile or label factor has been omitted: each cited theorem bounds its entire containing sector with those factors already present.

For rank2, R079 PC1--25 removes at support26 every D\_G\textgreater=33 pattern, and every D\_G\textless=32 pattern whose two short original covector values r(y1),r(y2) both vanish. DO109 SP26.1--4 and its complete count remove the complementary low-D\_G branch where at least one of those two values is nonzero. These are exhaustive alternatives for the integer D\_G and the two actual values. The last covector direction is included in R079, rather than silently discarded. Thus no nonidentity fixed-S support26 term remains.

\subsection{4. Exhaustion of all cross classes except rank-two unipotence}\label{exhaustion-of-all-cross-classes-except-rank-two-unipotence}

We now have A=I and C!=0. DO099 JR1--3/JR1--24 removes both the forward and inverse transverse cross sectors at supports25..70. In particular each surviving cross direction at26 lies in the original plane P. No support25 four-compression estimate is reused at26.

DO111 VC1--4/VC23--25 supplies that missing support26 step: it removes the entire both-plane class with at least one nonzero short compression Delta\_ab, a,b in\{1,2\}. Its proof assumes neither global unipotence nor a type profile. In its complement, the unitary transverse compression fixes the two short directions, and its original integral last-direction identity forces that compression to be I. Hence h\_ZZ=I and N(S perpendicular u-perpendicular) has one-dimensional isotropic image; there is a positive166-dimensional fixed subspace. The original positive167-dimensional compression is the identity. These statements do not assert that each short vector is fixed by the full h.

The unitary inverse is the Hermitian adjoint, so (h\^{}-1)\_ZZ=(h\_ZZ)\^{}dagger=I in the same orthogonal decomposition. The determinant-one central representative and C\textbar S!=0 were retained throughout. Thus every algebraic premise of DO095 RA1--10 is met at support26. That support-independent proof gives exactly two global ranks: rank3 when B\_Z C\_Z!=0 and rank2 when B\_Z C\_Z=0. Rank0 contradicts C!=0; rank1 would put the entire image in its plane cross line and force B\_Z=0, contrary to the same unitary identities and invertibility of the plane block. There is no omitted rank1 cross class.

R090 IS1--3 and its whole proof remove every resulting global-rank-three support26 term. It uses both original and inverse eligible lines and covers all k\textgreater=1; no sparse-coordinate restriction is needed. The actual type-III range11..26 is already included. Inversion is used on complete containing sectors with their charges retained, not as an unproved invariance of a previously reduced subset.

At global rank2, global unipotence and global nonunipotence are exhaustive alternatives. DO101 PI1--3/PI1--17 removes the entire plane-C nonunipotent class at every support25..32, hence at26. The only remaining cross class is globally rank-two unipotent. DO110 NT26.1--2/NT26.1--31 then removes every profile with at least one active type I orIII, equivalently with at most25 active type-II places. The only possible profile is all26 type II. All of these are complete containing-set removals and impose no new arithmetic source construction.

\subsection{5. Exact final denominator and coordinate conditions}\label{exact-final-denominator-and-coordinate-conditions}

The hypotheses of K103 SC26.1 now hold. Its SC26.5--7 proof supplies the support-independent full-positive normal form, nilpotence statements and exact reconstruction in DR3. Its SC26.2--4 theorem, established by nine exhaustive containing-set removals SC26.8--28, leaves only \[
 v=q=26,\quad b=j=0,\quad2^{1116}<x<2^{2108},
                         81\le k\le164.                 \tag{DR6}
\] Its exact lower norm comparison excludes x=2\^{}1116. Its complete group G removes x\textgreater=2\^{}2108, including equality. Its group H removes k\textless=80. Its other groups cover every failure of the stated eligibility and paired-integrality conditions. Thus the endpoints in DR6 are strict for the actual remaining set, not merely open ranges chosen to simplify notation.

K104 FV1 and its full-volume theorem FV1--11 remove every term in this same window with81\textless=k\textless=95. The residual integer range is therefore exactly the necessary96\textless=k\textless=164. K103's bounded-denominator uniqueness and one-label statement give the remaining assertions after DR3. The true norm equation supplies the partner slope U/bar(a0), and every one of the168 factors(T-1)\^{}(-1) remains in the containing bound. No complete central orbit, real-part-free central lattice, or Fourier constancy is asserted.

Sections2--5 establish every part of DR2 for actual h in D26. They establish a necessary description and preserve every previously imposed restriction; they do not identify D26 with all matrices in that description. In particular an estimate for that larger containing class may be useful later, but its actual sign or arithmetic realizability has not been proved here.

\subsection{6. Chronological accounting, diagnostic scope and the energy identity}\label{chronological-accounting-diagnostic-scope-and-the-energy-identity}

The accepted complete-floor file checks/r087\_joint\_remainder\_error.json has102 all-top and31 uniform charges, with floors26/11. The eight subsequent chronological additions are:

\begin{longtable}[]{@{}lr@{}}
\toprule\noalign{}
Finding & Retained exponent \\
\midrule\noalign{}
\endhead
\bottomrule\noalign{}
\endlastfoot
R088 & 523 \\
K103 & 293 \\
DO109 & 820 \\
K104 & 274 \\
R089 & 35000 \\
DO110 & 561 \\
R090 & 6668 \\
DO111 & 21366 \\
\end{longtable}

The current file is checks/do111\_joint\_remainder\_error.json. Its entire110/31 arrays are retained verbatim. There is exactly one all-top exponent231 and every other all-top exponent is at least254. Therefore \[
 \sum_{e\in\mathcal L_{\rm top}}2^{-e}
 \le2^{-231}+109\,2^{-254}<2^{-230}.                    \tag{DR7}
\] The entire31-entry uniform sum is separately below2\^{}-230. The companion executable verifies the complete chronological prefixes, each historical single charge, exact support floors, every source fingerprint, inclusive compressed-rank coverage, finite profile/coordinate splits and both complete dyadic sums. Its finite labels are diagnostics of the stated partition; they are not an enumeration of arithmetic arrows or a substitute for the source applicability and universal proof above.

Each containing-set estimate is applied to its intersection with the actual remainder at its chronological stage. Those removed intersections are disjoint even when the original containing sets overlap. The triangle inequality bounds their total contribution by the full retained charge sum. R087 already removed every support below26. Section1's D26 and R\textgreater=27 partition exactly what remains. Splitting that actual sum proves DR4 from the retained original energy identity and DR7. DR5 is sufficient for positivity because its strict lower bound dominates the worst permitted negative epsilon. Neither DR5 nor positivity is proved.

This paper records the consolidation of existing smaller sectors; it is not an extra charge, a floor27 theorem, or an independently new strict gate reduction beyond those accepted sector theorems. The uniform floor remains11. Fully sized admissible mixed-CM data, the actual selected arithmetic image and its content, prescribed period and Bessel/Mram values, specific032/Read/STiesR, all six actual lift identities, controlled nonvanishing and actual DP44 entries remain separate open gates. In particular neither a shorter valid chain nor a universal strict negative actual DP44 cost difference follows. Kummer-Hecke RH and the global oscillator arithmetic comparison remain open.

\subsection{7. Algebraic-geometric placement and exact contextual locators}\label{algebraic-geometric-placement-and-exact-contextual-locators}

The case split concerns actual fixed-submodule, determinantal and isotropic-cross strata of the original Hermitian correspondences over Spec(O\_E). The paired incidence modules, fractional denominator ideals, polarization, theta line and chosen complex realization stay fixed. Splitting a finite/real kernel remainder does not construct an inverse on an arithmetic theta image or a motivic comparison. No fractional ideal is declared principal and no rational moving frame is promoted to a global integral frame.

The carried module/determinantal context is SGA3, ExposeI,4.7.3/4.7.3.1, and SGA6, ExposeVI, Theorem3.3. Picard context remains FGA, Bourbaki232, Theorem3.1, pp148--149, and Bourbaki236, Corollaries1.7/2.3, pp229/232. Biextension context remains SGA7 I, ExposeVII,2.7/2.9, pp169--174, and ExposeVIII,7.4(b), pp310--311. Formal groups, Dieudonne crystals and Cartier duality retain cote4 scan003, sections8.1.1--8.1.2, scans005--006, sections8.4.1--8.4.3, and cote54 scans005--007. These are inherited exact framework locators, not new identifications.

No new lambda-ring/K-theory identity, l-adic sheaf or weight theorem, CM-motive identification or motivic L-function assertion is obtained. DP\_MULT Q3e with DP44/DP52/DP54--65 and THETA3 G1/A7 receive DR1--5 as the consolidated finite gate. CX\_GRH CX-ALG-024 and GCT\_PLETH T-052 keep their separate arithmetic gates. The source theorems are cited at their displayed run IDs and internal locators and are not re-proved. Alpoge's theorem and Lean archive remain stipulated; no build or archive audit is performed.

\subsection{Canonical algebraic-geometric context}\label{canonical-algebraic-geometric-context-28}

This is the original fixed-submodule/determinantal/isotropic-cross stratification over Spec(O\_E), with actual fractional ideals and paired incidence modules. Carried anchors: SGA3 I4.7.3/4.7.3.1; SGA6 VI Theorem3.3; FGA232 Theorem3.1 pp148--149; FGA236 Corollaries1.7/2.3 pp229/232; SGA7 I VII2.7/2.9 pp169--174 and VIII7.4(b) pp310--311; cote4 scan003 sections8.1.1--8.1.2, scans005--006 sections8.4.1--8.4.3, cote54 scans005--007. No new lambda-ring/K-theory, l-adic weight, formal-group or CM-motive identification. DP\_MULT Q3e/DP44/DP52/DP54--65 and THETA3 G1/A7 receive the consolidated finite gate; CX\_GRH CX-ALG-024 and GCT\_PLETH T-052 retain their arithmetic comparisons.

\clearpage

\section{AM-K-105 --- complex Gaussian counting for the next all-II support26 interval}\label{am-k-105-complex-gaussian-counting-for-the-next-all-ii-support26-interval}

\textbf{Canonical checked integration --- 2026-11-23T04:57:58.240734+00:00.} Complete root and independent proof/script/interface reviews and separate exact replays passed. The actual canonical predecessor has110/31 entries; this integration gives111/31, complete floors26/11. Frozen author history is preserved. AM-K-105 CG1--25 removes all39 complete coordinate sectors96\textless=k\textless=134 from the actual all-II globally rank-two-unipotent support26 scalar-cross remainder. The scalar real-dimension336 Poisson bound, full coupled Gram determinant, original selected trace, full affine central covolume and every finite/source factor remain. Restoring(N\_FJ0)\^{}(-1) and its actual common-G correlation gives the full weighted pattern factor\textless2Q\^{}(-167); complete coupled-x monotonicity gives the exact lower denominator endpoint. The worst cell is-9009839/6144\textless-1466 and the preceding38 are below-2831, so one1465 charge covers the union. Actual DO111110/31 becomes111/31, every overlap retained and both complete sums below2\^{}-230; floors26/11 stay fixed. With R091, the only possible support26 branch now has135\textless=k\textless=164, the same v=q26,b=j0 and strict denominator window2\^{}1116\textless x\textless2\^{}2108. Signed positivity, actual source, selected arithmetic content, period, specific032/Read, DP44 and RH remain open.

Author: /root/kummer\_extension, GPT-6 Astra Pro, ALP\_MACHINE. 23 November 2026 UTC. Controlled pending proof. Author saved-proof/script checks, exact replay and independent whole review are separate records.

\subsection{1. Actual sector and theorem}\label{actual-sector-and-theorem-10}

Retain every original J56 all-top object and normalization of K103/K104: H=133, f=168, original fields E/F, paired source and target pencils, finite projectors, selected real module, compact factors, twists, central phases, class masses and Haar measures. Apply this theorem after the actual K103 and K104 cuts, in their all-II support26 globally rank-two-unipotent scalar-cross branch. Its remaining conditions are \[
 s=26,\qquad v=q_{\rm elig}=26,\qquad
 2^{1116}<x=N_ED_\eta<2^{2108},\qquad b=j=0,\qquad96\le k\le164.
 \tag{CG1}
\] Here k counts the nonzero original S covector coordinates. The support-free full positive compression, N²Z=0 and exact unipotent normal form are K103 SC26.5--7. Eta is unique for each fixed support in this denominator window by SC26.11. This uniqueness is retained before summing any G patterns.

The actual 26 active weighted-label factors and 142 inactive pointwise factors give \[
 (T-1)^{-168}<2T^{-168}. \tag{CG2}
\] The active one-label result includes central-only poles; the additional j=0 condition is already part of CG1. No statement about a larger untruncated central orbit is substituted for this retained-set bound.

\textbf{Theorem.} The complete subset of CG1 with \[
                         96\le k\le134
\] has absolute normalized contribution less than \(2^{-1465}\). Charge this fresh subset once and retain every older, possibly overlapping charge. Within the all-II branch the new necessary coordinate range is \(135\le k\le164\). The denominator window and other conditions in CG1 remain unchanged. This is not a full support27 assertion.

Put \[
 Q=1427908193729121522601,\quad U=77Q^2-14,\quad
 R=17424/133,\quad M^2=R/U,\quad
 D_0=7^{140}29^{144}5^{84}13^{84},\quad D_F=2^fD_0.
 \tag{CG3}
\] For one canonical common-C G ideal J and the fixed eta, the original row ideal is \[
 L=JI_\Sigma D_\eta,\qquad N_EL=z^{2f},\qquad
 N_EJ=z^{2f}T^{52}/x\le Q^{334}.
 \tag{CG4}
\] All k original S values and the three fixed transverse values lie in L, including at the denominator primes of eta. This is K103 SC26.12--17's exact intersection, not an integral null-frame assertion. Eta, these row values and the central scalar d reconstruct the original matrix.

\subsection{2. The exact scalar dual minimum in real dimension336}\label{the-exact-scalar-dual-minimum-in-real-dimension336}

Use the original half-trace norm \(\|a\|^2=\sum_{\sigma\ {\rm CM}}|\sigma(a)|^2\). Its Euclidean dual ideal is \[
 L^\#=2\operatorname{Different}_E^{-1}\bar L^{-1}.
\] The actual discriminant identity is \(|D_E|=2^{2f}D_0^2\). Consequently \[
 N_EL^\#=D_0^{-2}z^{-2f},\qquad
 \lambda_1(L^\#)^2\ge fD_0^{-2/f}z^{-2}
                  >\frac{f}{732^2z^2}.                 \tag{CG5}
\] The inequality follows from the fractional-ideal norm minimum and arithmetic-geometric mean over the f complex labels. The exact integer comparison \(D_0<732^f\) is checked. Thus732 bounds the \textbf{half} root discriminant; the full root discriminant is twice this quantity. No extra division by2 is introduced. These are the same dual and metric conventions as DO070 FT3.

For completeness, let a real lattice Lambda have dimension n and minimum squared norm at least m², and write \[
 \Theta_\Lambda(s)=\sum_{\lambda\in\Lambda}e^{-s\|\lambda\|^2}.
\] For \(0<t<1\), the nonzero minimum and Poisson summation give \[
 \Theta_\Lambda(s)-1
 \le e^{-(1-t)sm^2}\Theta_\Lambda(ts)
 \le e^{-(1-t)sm^2}t^{-n/2}\Theta_\Lambda(s).             \tag{CG6}
\] Indeed the Gaussian Poisson formula is \[
 \Theta_\Lambda(ts)=\frac{(\pi/(ts))^{n/2}}{\operatorname{covol}\Lambda}
       \Theta_{\Lambda^*}\!\left(\frac{\pi^2}{ts}\right),
\] and its dual sum decreases as its parameter increases. This proves the second inequality of CG6, with its full dimension-dependent factor. All sums converge absolutely. This is the Gaussian specialization of Miller and Stephens-Davidowitz (2019, equations1.1 and2.1); the needed specialization is proved here.

Set \(A_0=2^{-16}\). For any \(0<A\le A_0\), use Lambda=L\#, n=2f=336, \[
 s=\pi^2z^2/A,\qquad m^2=f/(732^2z^2),\qquad
 u_0=\frac{(157/50)^2\,65536}{732^2}>1.
\] Then \(sm^2>fu_0\). Taking \(t=1/u_0\) in CG6 gives the exact scalar theta loss \[
 \eta_0=(e u_0)^f e^{-fu_0}.
\] With \(w=u_0-1\), the elementary inequality \(\log(1+w)\le w-w^2/2+w^3/3\) implies \[
 -\log\eta_0\ge f(w^2/2-w^3/3)
 =\frac{28175870985287512680844}{9169519753215087890625}>3.
 \tag{CG7}
\] Also \(e^3>1+3+3^2/2+3^3/6+3^4/24>16\). Hence eta0\textless1/16, and CG6 yields \[
 \sum_{\xi\in L^\#}e^{-\pi^2z^2\|\xi\|^2/A}<16/15
                     \quad(0<A\le2^{-16}).             \tag{CG8}
\] The factor \(t^{-168}\), not a one-dimensional theta estimate, is what produces the exponent f in CG7.

\subsection{3. Full coupled Gram matrix and one joint row count}\label{full-coupled-gram-matrix-and-one-joint-row-count}

Fix k original S positions and put \(d_k=k+3\). The true covector metric is \[
 H_k=\operatorname{diag}(I_k,G^{-1}),\qquad
 \det G=4U,\qquad G-I_3>0,
\] where G is the Gram matrix of the three fixed integral transverse inputs from R072/R076. In its Euclidean realization the row lattice is \(H_k^{1/2}L^{d_k}\), of exact real covolume \[
 \frac{D_0^{d_k}z^{2d_kf}}{(4U)^f}.                     \tag{CG9}
\] Its dual is \(H_k^{-1/2}(L^\#)^{d_k}\). The dual norm has Gram diag(I\_k,G), so it is at least the ordinary product norm on \((L^\#)^{d_k}\). Therefore its dual Gaussian mass is at most the product of d\_k scalar masses in CG8: \[
 \Theta_{\rm row}^*(\pi^2z^2/A)<(16/15)^{d_k}.           \tag{CG10}
\] This product domination is essential. CG8 is applied in real dimension336 separately to each scalar ideal; it is not applied in dimension336d\_k with an unchanged scalar minimum. The full transverse determinant remains in CG9.

The original joint energy is \(\|r\|^2\le fM^2\). The single Gaussian majorant and Poisson formula consequently give, for \(w=M^2/z^2\), \[
 \#r\le {164\choose k}(16/15)^{d_k}
       \frac{(4U)^f}{D_0^{d_k}}
       \left(\frac{\pi}{A}\right)^{d_kf}e^{Afw},
           \qquad0<A\le2^{-16}.                         \tag{CG11}
\] All actual complex values, including their phases, are counted directly. Eta has already been fixed by the support's uniqueness theorem; no relative-norm, beta-pivot or second eta enumeration is added. The old harmless reconstruction-pivot factor \(164<2^8\) is retained below.

The enlarged lattice sum includes tuples with zero or excessive coordinates. The actual nonzero-pivot trace bound is first applied uniformly to the actual k-sector, and only its cardinality is then enlarged. No trace statement is made for added zero tuples.

\subsection{4. Genuine central, finite and selected-real factors}\label{genuine-central-finite-and-selected-real-factors}

Retain the full covolume in K103 SC26.14 before its SC26.15 enlargement. Let J0 be the original F-ideal of rank-zero G places, as in DO098 CF10. The precise bound is \[
 \#d\le \frac{V_F(H/2+750Uz^2)\,U^{-f}T^{52}x^{-1}}{N_FJ_0},\qquad
 |c(h)|<2^{-929f/2}U^{-f/2}z^{-f},
\] \[
 V_F(C)=\frac{((22/7)fC^2)^{f/2}}{(f/2)!\sqrt{D_F}}.      \tag{CG12}
\] The central real-part shift is fixed by the row, and the count is in its actual difference module. Its covolume is \(\sqrt{D_F}(N_FJ_0)xT^{-52}\) before passing from d to the original negative entry −Ud. Its full-real cover750 is still enlarged using \(N_FJ_0\le N_EJ\), exactly as in K103. This enlargement of the numerator does not require discarding \(N_FJ_0\) from the denominator. No new local lemma is used to restore that positive divisor. There is no extra source isometry fibre because eta,r,d reconstruct h.

The common whole-C Green estimate is \[
 |a_G(h)|\le2Q^{-335\cdot167}(N_EJ)^{335/2}.              \tag{CG13}
\] Both mixed finite projectors remain with their modulus-one bounds. Combining CG2--4 and CG9--13 leaves, apart from positive factors independent of w, \[
 x^{-337/2}(N_FJ_0)^{-1}\,
 F_k(w),\quad
 F_k(w)=w^{-167f}
       \left(\frac{\pi}{A(w)}\right)^{d_kf}
       e^{A(w)fw}V_F(H/2+750R/w),
\] \[
                      A(w)=\min\{2^{-16},d_k/w\}.        \tag{CG14}
\] The full T power is \(335\cdot26+52-168=8594\), and the grouped U and R powers are \(U^{-335f/2}R^{167f}\). All original measure and coefficient normalizations remain unchanged.

\subsection{5. Weighted sum of every actual G pattern}\label{weighted-sum-of-every-actual-g-pattern}

The function CG14 is continuous at \(w=d_k2^{16}\). In the constant-A branch, differentiation with respect to log w gives \[
 \frac{d\log F_k}{d\log w}
 =-167f+A_0fw
 -f\,\frac{750R/w}{H/2+750R/w}
 \le-(164-k)f,
\] because \(A_0w\le d_k\). In the other branch \(A=d_k/w\), the first two factors have combined w power \((d_k-167)f\); the central factor decreases. Thus the same bound holds there. Consequently, for all \(w_2\ge w_1>0\), \[
 F_k(w_2)\le F_k(w_1)(w_2/w_1)^{-(164-k)f}.              \tag{CG15}
\] This is a bound on the complete count, including its moving central radius. It compares uniform norm-dependent upper bounds for the different original ideals; it does not require those ideals to be homothetic.

For fixed eta and x, the all-rank-zero G cap is \[
 z_*=(Q^{334}T^{-52}x)^{1/(2f)},\qquad
 w_*=M^2/z_*^2 .
\] The actual whole-C ideal has local norm exponents2,1,0 at each of the334 places, with multiplicities at most1,2,1. Exponent2 occurs only at a rank-zero place, contributing exponent1 to J0; the other states contribute exponent0 to J0. Thus the source-pinned DO098 CF10 inequality is, place by place and globally, \[
 N_EJ\le Q^{167}N_FJ_0,\qquad
 \frac1{N_FJ_0}\le Q^{-167}(w/w_*)^f.                    \tag{CG15a}
\] The second inequality uses \(N_EJ=Q^{334}(w_*/w)^f\). It remains valid when denominator primes of eta overlap G, because J and J0 are the original finite-rank ideals and the row ideal already includes D\_eta separately.

A decrease by h in a local J exponent multiplies w by \(q_G^{h/f}\). By CG15, the complete enlarged factor \((w/w_*)^fF_k(w)\) decreases at least with exponent \(-(163-k)f\). Summing the complete pattern set therefore gives \[
 \sum_{\rm patterns}\frac{F_k(w)}{N_FJ_0}
 \le Q^{-167}F_k(w_*)\prod_G(1+q_G^{-(163-k)})^2
 <2Q^{-167}F_k(w_*) \quad(96\le k\le134).                \tag{CG16}
\] For the last strict inequality, \(163-k\ge29\), every \(q_G\ge g_0=22285932761\), and the exact comparison \(2\cdot668<g_0^{29}\) gives \((1+g_0^{-29})^{668}\le(1-668g_0^{-29})^{-1}<2\). The exponent states and the two orientations are the same original ones as K104 FV9; no additional finite rank labels occur.

The pattern sum replaces2\^{}668 and leaves the extra factor \(Q^{-167}\) at the all-rank-zero endpoint. At that endpoint \(N_FJ_0=Q^{167}\), so this is its exact recovered divisor. Eta and its denominator do not vary independently with the pattern. The exact support count is \({168\choose26}<2^{101}\). Together with the pattern factor2, pivot2\^{}8, one Green2 and the single correction2 in CG2, the complete discrete exponent is \[
                           101+1+8+1+1=112.             \tag{CG17}
\]

\subsection{6. The coupled denominator endpoint}\label{the-coupled-denominator-endpoint}

After CG16 the recovered factor \(Q^{-167}\) is independent of x; no factor \(w_*^{-f}\) remains. For the fixed support's actual eta, \(w_*(x)\) is proportional to \(x^{-1/f}\). Let \[
 \epsilon(w)=\frac{750R/w}{H/2+750R/w},\qquad0<\epsilon(w)<1.
\] In the constant-A branch, the full product after CG16 has logarithmic x derivative \[
 \frac{d}{d\log x}\log\{x^{-337/2}F_k(w_*(x))\}
 =-\frac32-A_0w_*+\epsilon(w_*)<-\frac12-A_0w_*<0.
\] In the branch \(A=d_k/w_*\), it is \[
                     -\frac32-d_k+\epsilon(w_*)<-\frac12-d_k<0.
 \tag{CG18}
\] The expressions agree at the switching point. Hence the complete bound decreases with x throughout the whole actual window, and its maximum is at the lower numerical endpoint \(x=2^{1116}\). This includes all patterns and the central radius, not merely a selected variable power.

At that endpoint put \(c=w_*\). The unchanged exact host integers give \[
 c^f=(R/U)^fT^{52}Q^{-334}2^{-1116},\qquad
                  c_-=4124955<c<c_+=4124956.             \tag{CG19}
\] For all claimed k, \(c_+<(k+3)2^{16}\), so the endpoint uses A0. Also \[
 H/2+750R/c_-<68,\qquad V_F(68)<2^{479}.                 \tag{CG20}
\] The second comparison is checked after squaring, with the exact \(D_F\); no root-discriminant or real-lattice factor is suppressed.

There is one whole endpoint per k. Neither the piecewise choice of A nor the weighted G sum creates a new endpoint or range union.

\subsection{7. Exact rational endpoint and complete fresh union}\label{exact-rational-endpoint-and-complete-fresh-union}

Use the previously directed bounds \[
 L_T=1019333/1024,\quad L_Q=287843/4096,\quad
 L_U=601356/4096,\quad L_R=28810/4096,
\] upper for T,R and lower for Q,U. The complete common exponent is \[
\begin{split}
 C_G={}&112+2f-929f/2+479+167fL_R-(335f/2)L_U\\
       &+8594L_T-336\cdot167L_Q-(337/2)1116\\
     ={}&211150679/512 .
\end{split}                                             \tag{CG21}
\] The2f is the retained determinant factor4\^{}f.~All positive T and R powers and all negative U powers have been combined before substituting directed bounds.

For \(d_k=k+3\), let e\_k be the least strict integral binary upper exponent of the rational number \[
 W_k=
 \frac{{164\choose k}16^{d_k}(22\cdot65536)^{d_kf}}
 {15^{d_k}7^{d_kf}D_0^{d_k}c_-^{167f}}.                 \tag{CG22}
\] The theta product is exactly the factor \((16/15)^{d_k}\). Since \(\log_2e<13/9\), the positive energy exponential is bounded by \[
 \frac{13}{9}\frac{fc_+}{65536}
                 =93842749/6144.                       \tag{CG23}
\] The complete k-cell exponent is \(C_G+93842749/6144+e_k\). Every term in this formula comes from CG11--17 and CG20; there is no additional Gaussian or phase multiplier.

The companion checks all39 fresh cells96 through134 by exact integer powers and rational arithmetic. The largest exponent is \[
 E_{134}=-9009839/6144<-1466.                            \tag{CG24}
\] All preceding38 cells are below−2831; in particular \[
 E_{133}=-17396399/6144<-2831.
\] Thus the full new union satisfies \[
 \sum_{k=96}^{134}|\text{complete k-cell}|
       <2^{-1466}+38\,2^{-2831}<2^{-1465}.               \tag{CG25}
\] The next-cell diagnostic gives \(-617135/6144>-101\), too large to append under the required complete error budget below \(2^{-230}\). It supplies no removal, lower bound or existence statement. No removal for k135..164 is claimed here.

\subsection{8. Exact verification, chronology and remaining obligations}\label{exact-verification-chronology-and-remaining-obligations}

The certificate recomputes the actual host, half-trace dual normalization constants, rational dimension-dependent theta inequality, full Gram minors/determinant, central rank-zero correlation, weighted pattern inequality, support count, coupled endpoint bracket, central real volume, all directed logs and all39 complete cells. Rational samples of the two derivative identities corroborate the universal piecewise calculation, not replace it. Both full ordered error arrays are checked.

The base is canonical checks/do111\_joint\_remainder\_error.json with110 all-top and31 uniform entries, complete floors26/11. The script accepts a later complete canonical extension of that base without reordering it. Append1465 exactly once; retain274,293 and every earlier charge, as well as the entire uniform array. Both full sums remain below2\^{}-230. No automatic support-floor change is associated with this coordinate interval.

The literal interfaces are K103 SC26.1--7/11--20 (actual sector, unique eta, full positive normal form, original row and central ideals, selected trace and all168 factors); K104 FV1--10 (the preceding coordinate cut, actual G-pattern states and support count); DO098 CF10 and K103 SC26.14--15 (the restored central divisor); R076 CV8--10 (the fixed coupled Gram and covolume); DO070 FT3 (the half-trace dual); and K102 PC11/17--19 (Poisson convention and unchanged source factors). The dimension-aware scalar estimate, coupled Gaussian row count and both monotonicities are proved anew above. Only the new controlled endpoint is counted here.

The underlying original ideals, field, involution, Hermitian metric, polarization, theta line, Picard/biextension data, finite projectors and selected coefficient are unchanged. Exact inherited context remains SGA3 I4.7.3/4.7.3.1; SGA6 VI Theorem3.3; FGA232 Theorem3.1 pp148--149; FGA236 Corollaries1.7/2.3 pp229/232; SGA7 I VII2.7/2.9 pp169--174 and VIII7.4(b) pp310--311; cote4 scan003 sections8.1.1--8.1.2/scans005--006 sections8.4.1--8.4.3 and cote54 scans005--007. These are contextual anchors, not new comparison or motivic claims.

The dense all-II coordinates135..164, signed sufficient inequality, actual selected arithmetic-image content, prescribed global Bessel/Mram value, effective fully admissible source construction, Read/STiesR/specific032, changed-family DP44 and RH remain separate. This proof makes no Lean-build or archive-audit claim.

\subsubsection{Primary reference}\label{primary-reference-3}

Miller, S. D., \& Stephens-Davidowitz, N. (2019). \emph{Kissing numbers and transference theorems from generalized tail bounds} (arXiv:1802.05708v2). \url{https://arxiv.org/pdf/1802.05708}. Equation1.1, printed pages1--2, and the Poisson convention in equation2.1, printed page4, were freshly read. CG6--8 supplies the complete specialization in real dimension336.

\subsection{Canonical algebraic-geometric context}\label{canonical-algebraic-geometric-context-29}

The original fractional invertible modules, Hermitian form, involution, polarization, theta line, paired incidences and selected complex realization are unchanged. Carried exact anchors: SGA3 I4.7.3/4.7.3.1; SGA6 VI Theorem3.3; FGA232 Theorem3.1 pp148--149; FGA236 Corollaries1.7/2.3 pp229/232; SGA7 I VII2.7/2.9 pp169--174 and VIII7.4(b) pp310--311; cote4 scan003 sections8.1.1--8.1.2, scans005--006 sections8.4.1--8.4.3 and cote54 scans005--007. No new lambda-ring/K-theory, l-adic weight, formal-group or CM-motive comparison. DP\_MULT Q3e/DP44/DP52/DP54--65 and THETA3 G1/A7 gain the stated smaller finite gate; CX\_GRH CX-ALG-024 and GCT\_PLETH T-052 retain their arithmetic gates. Primary bibliographic precision: Stephen D. Miller and Noah Stephens-Davidowitz (2019), arXiv1802.05708v2, equations1.1/2.1, printedpp1--2/4. The frozen source manifest misspells the first given name as Steven; its S. D. proof citation and mathematics are unchanged.

\clearpage

\section{AM-LIFT-089 --- Full-action and half-level bounds for the same seed field}\label{am-lift-089-full-action-and-half-level-bounds-for-the-same-seed-field}

\textbf{Canonical checked integration --- 2026-11-23T04:59:16.985229+00:00.} Complete saved-proof and executable reviews and separate exact diagnostic replays passed. No numerical error charge is introduced. Frozen author history is preserved. AM-LIFT-089 HF89.1--18 sharpens the degree of exactly the same LIFT088 seed and output fields. For r=p+q,g=fr, the actual projective OE-lattice is free modulo each finite ideal, and the full defined action bounds its torsion Galois image. The literal m/bar(m) level halves are perfect Weil duals; Coeff.zeta\_mem already gives mu\_M inB. Thus Kgood(A{[}M{]})=Kgood(A{[}m{]}) and {[}Kseed:Kgood{]}\textless=M\^{}(f r\^{}2). The same target field Kseed(A{[}N\^{}2{]}) has relative degree\textless=N\^{}(4f r\^{}2), including ramifiedN andN1. Consequently {[}Kout:B{]}\textless=3870720\emph{d}M\^{}(f r\^{}2)*N\^{}(4f r\^{}2), under every actual LIFT085 nonzero-rule/source hypothesis. All actions, lines, level, Obj/FrPt and intrinsic costs are unchanged. EffectiveN, absoluteBdegree, discriminants, coefficient/frame heights, actual Host/Sizes existence, selected arithmetic content, period, specific032/Read, DP44 and RH remain open. No error charge; current111/31 arrays and complete floors26/11 are unchanged.

Author: /root/lift\_accounting, ALP\_MACHINE. 23 November 2026 UTC.

\textbf{Author-frozen proof.} The final whole-file, source-interface and executable checks are recorded in the adjacent verification packet; independent review and canonical integration are separate. This refines the relative field-degree estimates in checked LIFT088. The seed, polarization and canonical line, finite level, chosen object and frames are unchanged. No effective isogeny scale or prescribed nonzero arithmetic rule is constructed.

\subsection{1. Statement in the literal source data}\label{statement-in-the-literal-source-data}

Retain all hypotheses and the newly chosen seed of LIFT087--088. There are actual source Tor.Host, Datum, Level, Moduli and Coeff records; the complete integer-ring action is defined over the stated ground field; the level has \(A=1\); and every compatibility used in LIFT084--088 remains an input. Put

\[
 \mathcal O=\mathcal O_E,\quad [E:\mathbb Q]=2f,\quad
 r=Dw.p+Dw.q\ge1,\quad g=fr,\quad B=c.F.
 \tag{HF89.1}
\]

Write \(\mathfrak m=l.Mb\). The source gives

\[
 \mathfrak m\overline{\mathfrak m}=M\mathcal O,\qquad
 \mathfrak m+\overline{\mathfrak m}=\mathcal O,\qquad M\ge3.
 \tag{HF89.2}
\]

Let \(K_{\rm good}\supset B\) be the actual field in LIFT088 before adding level torsion. It carries the same everywhere-good polarized CM seed \(A_0\), its full action and canonical rigidified cubical line. LIFT088 proves

\[
 [K_{\rm good}:B]\le168\cdot23040\,h(-d).
 \tag{HF89.3}
\]

The source Coeff.zeta\_mem puts \(\zeta_{M'}\) in \(B\), where \(M'=\operatorname{lcm}(M,A)\). Thus \(A=1\) gives \textbf{\(\mu_M\subset B\)}. This is an existing coefficient-field condition, not a new root extension.

\textbf{Theorem.} The very same seed field chosen in LIFT088 satisfies

\[
\begin{split}
 K_{\rm seed}
 &=K_{\rm good}(A_0[M])
   =K_{\rm good}(A_0[\mathfrak m]),\\
 [K_{\rm seed}:K_{\rm good}]
 &\le|\operatorname{GL}_r(\mathcal O/\mathfrak m)|
 \le M^{fr^2}.
\end{split}
\tag{HF89.4}
\]

Consequently

\[
 [K_{\rm seed}:B]\le
 3\,870\,720\,h(-d)M^{fr^2}
 \le3\,870\,720\,dM^{fr^2}.
 \tag{HF89.5}
\]

Under the complete nonzero-rule and orbit hypotheses of LIFT085--086, their same target field

\[
 K_{\rm out}=K_{\rm seed}(A_0[N^2])
 \tag{HF89.6}
\]

satisfies

\[
 [K_{\rm out}:K_{\rm seed}]
 \le|\operatorname{GL}_r(\mathcal O/N^2\mathcal O)|
 \le N^{4fr^2}.
 \tag{HF89.7}
\]

For \(N=1\), interpret the extension as trivial and the upper bound as one. Therefore

\[
 [K_{\rm out}:B]\le3\,870\,720\,dM^{fr^2}N^{4fr^2}.
 \tag{HF89.8}
\]

The estimates include every object, action, line, finite level, field-valued moduli point and frame previously descended over those same fields. They do not change the absolute degree of \(B\), any discriminant, the denominator \(N\), intrinsic line degrees or the coefficient-field realization.

\subsection{2. The projective lattice modulo an ideal}\label{the-projective-lattice-modulo-an-ideal}

The source Datum.La is a finite projective \(\mathcal O\)-module of constant rank \(r\). For every nonzero integral ideal \(\mathfrak a\),

\[
 \Lambda/\mathfrak a\Lambda
 \simeq(\mathcal O/\mathfrak a)^r
 \quad\text{as }\mathcal O/\mathfrak a\text{-modules}.
 \tag{HF89.9}
\]

A globally free integer lattice is not assumed. Factor \(\mathfrak a=\prod_{\mathfrak p}\mathfrak p^{e_{\mathfrak p}}\). Chinese remaindering decomposes \(\mathcal O/\mathfrak a\) into the local Artinian factors \(\mathcal O/\mathfrak p^{e_{\mathfrak p}}\). At each \(\mathfrak p\), the localization of the projective lattice is free of rank \(r\), since \(\mathcal O_{\mathfrak p}\) is a discrete valuation ring. Reducing a local basis gives a free rank-\(r\) module over its Artinian factor. The finitely many bases combine under the central idempotents to a basis over the product ring. This proves HF89.9 without a class-group field or a basis-height factor. In particular,

\[
\begin{split}
 |\Lambda/\mathfrak a\Lambda|&=(N_E\mathfrak a)^r,\\
 |\operatorname{Aut}_{\mathcal O}
       (\Lambda/\mathfrak a\Lambda)|
 &=|\operatorname{GL}_r(\mathcal O/\mathfrak a)|
 \le(N_E\mathfrak a)^{r^2}.
\end{split}
\tag{HF89.10}
\]

The last inequality counts all matrices. It holds for arbitrary prime powers, including rational primes ramified in \(E\); the finite quotient need not be a field.

\subsection{3. Literal geometric torsion and the level half}\label{literal-geometric-torsion-and-the-level-half}

Over an algebraic closure of \(K_{\rm good}\), TateIso.ker, onto, add and lin give an \(\mathcal O\)-linear isomorphism

\[
 \Lambda/n\Lambda\xrightarrow{\;\sim\;}A_0[n]\qquad(n\ge1).
 \tag{HF89.11}
\]

This uses the actual geometric Tate isomorphism. It is not an infinite tower trivialization over a number field.

Choose \(e\in\mathcal O\) with

\[
 e\equiv1\pmod{\mathfrak m},\qquad
 e\equiv0\pmod{\overline{\mathfrak m}}.
 \tag{HF89.12}
\]

Then \(e^2=e\) and \(\bar e=1-e\) modulo \(M\). The resulting decomposition is

\[
 A_0[M]=P\oplus Q,\qquad
 P=eA_0[M]=A_0[\mathfrak m],\quad
 Q=(1-e)A_0[M]=A_0[\overline{\mathfrak m}].
 \tag{HF89.13}
\]

Under HF89.11, \(P=e\Lambda/M\Lambda\simeq\Lambda/\mathfrak m\Lambda\). The actual level map is

\[
 \eta(a)=\operatorname{ev}_M(ea).
 \tag{HF89.14}
\]

It kills \(\mathfrak m\Lambda\), and on \(\overline{\mathfrak m}\Lambda\) it equals \(\operatorname{ev}_M(a)\), exactly as Obj.eta\_ker and the final clause of Obj.geom require. Its image is all of \(P\). No independently varied decoration or theta frame is introduced.

Conjugation preserves ideal norms. Hence

\[
 (N_E\mathfrak m)^2=M^{2f},\qquad N_E\mathfrak m=M^f.
 \tag{HF89.15}
\]

The field generated by all points of \(P\) is finite Galois: characteristic-zero torsion is finite étale and \(P\) is defined over \(K_{\rm good}\). Its Galois action commutes with the full \(\mathcal O\)-action, which is already defined over that field. The faithful torsion representation embeds its Galois group in \(\operatorname{Aut}_{\mathcal O}(P)\). HF89.10 and HF89.15 now give the degree bound in HF89.4.

\subsection{4. The opposite half adds no field}\label{the-opposite-half-adds-no-field}

The principal polarization and Rosati-compatible action are defined over \(K_{\rm good}\). Their Weil pairing is Galois equivariant and perfect, with

\[
 e_M(ax,y)=e_M(x,\bar a y).
 \tag{HF89.16}
\]

In the literal Tate model, this follows from Datum.adj and TateIso.isom applied to the unimodular alternating pairing. Both \(P\) and \(Q\) are isotropic. Indeed, for \(x,y\in P\), insert \(a=e\) in HF89.16 and use \(\bar e\,y=(1-e)y=0\); the other half is identical. Perfection gives

\[
 P\times Q\longrightarrow\mu_M,\qquad
 Q\simeq\operatorname{Hom}_{\mathbb Z}(P,\mu_M).
 \tag{HF89.17}
\]

Explicitly, an element of \(Q\) orthogonal to \(P\) is also orthogonal to \(Q\), and hence to the full torsion, so it is zero. Both sides of the last map have order \(M^g\), which proves surjectivity. The dual cardinality follows by decomposing the finite \(M\)-killed abelian group into cyclic groups: a cyclic factor of order \(d\mid M\) has exactly \(d\) homomorphisms to \(\mu_M\).

Because \(\mu_M\subset B\subset K_{\rm good}\), a Galois element fixing \(P\) pointwise also fixes \(Q\): its pairing with each \(p\in P\) is unchanged, and the cross pairing is nondegenerate. Thus adjoining \(P\) makes all \(M\)-torsion rational. This proves the exact equality of fields in HF89.4.

The proof uses the existing roots supplied by Coeff.zeta\_mem. It does not attach such roots for free to a smaller field, and it does not assert trivial cyclotomic character at arbitrary torsion orders.

\subsection{5. The same target torsion field}\label{the-same-target-torsion-field}

Applying HF89.11 and HF89.9 at \(n=N^2\) gives

\[
 A_0[N^2]\simeq(\mathcal O/N^2\mathcal O)^r,\qquad
 |\mathcal O/N^2\mathcal O|=N^{4f}.
 \tag{HF89.18}
\]

Galois again acts faithfully and \(\mathcal O\)-linearly. Counting all matrices proves HF89.7. No primeness, splitting or absolute unramifiedness condition on \(N\) is needed for this inequality. The prime-avoidance conditions in LIFT085 may still be imposed.

The quotient, full action, principal polarization, canonical trace-one line, finite \(\eta\), Moduli.rep point and eigenframe are exactly those of LIFT085--088. Their all-object descent remains the earlier theorem; this corollary sharpens the size of their existing torsion field.

\subsection{6. Cost and quantifier table}\label{cost-and-quantifier-table}

\begin{longtable}[]{@{}
  >{\raggedright\arraybackslash}p{(\columnwidth - 6\tabcolsep) * \real{0.2143}}
  >{\raggedleft\arraybackslash}p{(\columnwidth - 6\tabcolsep) * \real{0.2857}}
  >{\raggedleft\arraybackslash}p{(\columnwidth - 6\tabcolsep) * \real{0.2857}}
  >{\raggedright\arraybackslash}p{(\columnwidth - 6\tabcolsep) * \real{0.2143}}@{}}
\toprule\noalign{}
\begin{minipage}[b]{\linewidth}\raggedright
Quantity
\end{minipage} & \begin{minipage}[b]{\linewidth}\raggedleft
LIFT088 bound
\end{minipage} & \begin{minipage}[b]{\linewidth}\raggedleft
Bound here
\end{minipage} & \begin{minipage}[b]{\linewidth}\raggedright
Scope
\end{minipage} \\
\midrule\noalign{}
\endhead
\bottomrule\noalign{}
\endlastfoot
Same initial level field over \(K_{\rm good}\) & \(M^{4g^2}\) & \(M^{fr^2}\) & Full action, the literal half level, and \(\mu_M\subset B\) \\
Same target torsion field over \(K_{\rm seed}\) & \(N^{8g^2}\) & \(N^{4fr^2}\) & Full action alone, including absolute ramification \\
Relative output degree over \(B\) & \(3\,870\,720\,dM^{4g^2}N^{8g^2}\) & \(3\,870\,720\,dM^{fr^2}N^{4fr^2}\) & Same seed, output and coefficient field \\
Absolute degree, discriminants and \(N\) & Unevaluated & Unevaluated & The actual \([B:\mathbb Q]\) remains; no new height estimate \\
\end{longtable}

For finitely many stages over the common pre-level field, the level factor becomes \(M^{f\sum_jr_j^2}\). If stage \(j\) uses scale \(N_j\), its target factor is \(N_j^{4fr_j^2}\). The common \(168\cdot23040\,h(-d)\) ground-field factor is used once when the simultaneous LIFT088 construction uses its common good field.

Every field embedding transports the same full action and projective quotient. Their ranks and finite-ring cardinalities are unchanged. This does not introduce a new coefficient field.

The main effective gate remains \(N\). The input to LIFT085 must still be the actual nonzero rule. Host, Sizes, Pins, Moduli, Unf and Step5 prerequisites retain their separate scopes. No prescribed arithmetic image, global period, central L-value, 032 value, DP44 sign or RH is established.

\subsection{7. Exact source addresses}\label{exact-source-addresses}

The directly relevant source interfaces are:

\begin{itemize}
\tightlist
\item
  TwoWeightEqOneAddRoot.lean: Datum 30--49; Level 51--61 and \(M'\) at63; TateIso 131--142; Obj 144--160.
\item
  PushInvertible.lean: Moduli 29--38; Coeff 40--46, especially conj at43 and zeta\_mem at44.
\item
  \nolinkurl{PartitionExtendByZeroAntitone.lean:} PinLevel at43, retaining \(A=1\) and the actual \(M=g_0g_1\).
\item
  IsCombAt.lean: actual FrPt 67--76.
\item
  Canonical LIFT088 SF1--16: the same good, seed and target fields and all-object descent.
\item
  Canonical LIFT085 HF3--4 and §§6--7; canonical LIFT086 IM11--18: actual quotient and model descent.
\end{itemize}

The new proof is the elementary projective Artinian quotient, full-action Galois centralizer and source-normalized Weil-dual-half calculation. No Mumford--Tate or open-image theorem is used, and equality of the Galois image with the full general linear group is not asserted.

\subsection{Canonical algebraic-geometric context}\label{canonical-algebraic-geometric-context-30}

The finite-etale torsion modules carry the actual OE-action and perfect polarization pairing. The two literal idempotent halves are Cartier dual through the existing roots of unity; all descended abelian schemes and cubical Picard/Poincare data are those already constructed. Carried exact anchors: FGA232 Theorem3.1 pp148--149; FGA236 Corollaries1.7/2.3 pp229/232; SGA7 I VII2.7/2.9 pp169--174 and VIII7.4(b) pp310--311; SGA3 I4.7.3/4.7.3.1; SGA6 VI Theorem3.3; cote4 scan003 sections8.1.1--8.1.2 and scans005--006 sections8.4.1--8.4.3; cote54 scans005--007. No new formal-group slope, lambda-ring/K-theory identity, l-adic weight theorem or CM-motive comparison follows. DP\_MULT Q3e/DP44/DP52/DP54--65, THETA3 G1/A7, CX\_GRH CX-ALG-024 and GCT\_PLETH T-052 retain their separate arithmetic comparisons.

\clearpage

\section{AM-DO-112 --- an enlarged Gaussian cap for two further all-II support26 cells}\label{am-do-112-an-enlarged-gaussian-cap-for-two-further-all-ii-support26-cells}

\textbf{Canonical checked integration --- 2026-11-23T05:07:09.474210+00:00.} Complete root and independent proof/script/interface reviews and separate exact replays passed. The actual canonical predecessor has111/31 entries; this integration gives112/31, complete floors26/11. Frozen author history is preserved. AM-DO-112 EC1--22 removes both complete coordinate sectors k135 and k136 from the actual all-II globally rank-two-unipotent support26 scalar-cross remainder. The larger Gaussian cap59/(50*65536), with the full real dimension336 tail and scalar theta bound26, preserves the original row metric, affine central divisor, selected trace, every finite factor and complete coupled pattern/denominator monotonicities. Exact whole-sector exponents are-691859209/307200\textless-2252 and-283283209/307200\textless-922; their union is below2\^{}-921. Actual K105111/31 becomes112/31, every historical overlap retained and both complete sums below2\^{}-230; floors26/11 stay fixed. The only possible support26 branch now has137\textless=k\textless=164, v=q26,b=j0 and strict denominator window2\^{}1116\textless x\textless2\^{}2108. The positive k137 diagnostic is unclaimed. Signed positivity, actual source, selected arithmetic content, period, specific032/Read, actual DP44 and RH remain open.

Author: /root/definite\_oscillator, ALP\_MACHINE. 23 November 2026 UTC. Controlled pending proof. Whole saved-proof/script review, exact replay and independent review are separate records. Earlier frozen subjects are unchanged.

\subsection{1. The actual sector and the complete new cut}\label{the-actual-sector-and-the-complete-new-cut}

Retain the original J56 all-top kernel, H=133, f=168, the original fields, actual paired source and target pencils and admitted labels, selected real representation and Hilbert metric, both compact factors, two mixed finite projectors, twists, central phases, Haar measures and class masses. Work in the all-II, globally rank-two-unipotent scalar-cross support26 remainder after K103, K104 and K105. Its necessary conditions include \[
 s=26,\quad v=q_{\rm elig}=26,\quad
 2^{1116}<x=N_E D_\eta<2^{2108},\quad b=j=0,\quad135\le k\le164 .
 \tag{EC1}
\] Here k is the number of nonzero original S covector coordinates; eta has norm U and is unique for a fixed support in this window. The support-free positive compression, exact normal form and uniqueness are K103 SC26.5--7/11, not a transplanted support25 count. K105 CG1--25 is the separately checked predecessor; its frozen proof and certificate remain unchanged.

\textbf{Theorem.} The complete subset of EC1 with k=135 or136 has absolute normalized contribution less than \(2^{-921}\). Remove its intersection with the current remainder once and retain every previous charge. The all-II coordinate range is then necessarily \(137\le k\le164\). This is a strict subset reduction, not a support27 theorem or an existence assertion.

The sole new analytic choice is the larger uniform Gaussian cap \[
 A_*=\frac{59}{50\cdot65536}.
 \tag{EC2}
\] The proof pays its larger scalar dual-theta loss explicitly. No numerical optimization is presumed in the argument.

\subsection{2. Exact carried arithmetic and metric interfaces}\label{exact-carried-arithmetic-and-metric-interfaces}

Put \[
 Q=1427908193729121522601,\quad U=77Q^2-14,\quad
 R=17424/133,\quad M^2=R/U,\quad
 D_0=7^{140}29^{144}5^{84}13^{84},\quad D_F=2^fD_0 .
 \tag{EC3}
\] For a canonical common whole-C finite G ideal J, let \[
 L=JI_\Sigma D_\eta,\qquad N_EL=z^{2f},\qquad
 N_EJ=z^{2f}T^{52}/x\le Q^{334}.
 \tag{EC4}
\] Every counted original S value and each of the three fixed integral transverse values belongs to this same ideal, including at denominator primes of eta. K103 SC26.13 proves the equality \(JI_\Sigma\cap\eta^{-1}JI_\Sigma=JI_\Sigma D_\eta\) valuationwise. No integral moving null frame is used.

The true row metric is \(H_k=\operatorname{diag}(I_k,G^{-1})\). The actual fixed three-vector Gram satisfies \(\det G=4U\) and \(G-I_3>0\); hence the real row covolume is \[
 \operatorname{covol}(H_k^{1/2}L^{k+3})
       =D_0^{k+3}z^{2(k+3)f}/(4U)^f .
 \tag{EC5}
\] This is R076 CV8--10 and K105 CG9--10. The actual joint energy is \(\|r\|^2\le fM^2\). Eta, all row values and the central scalar d reconstruct the original matrix, so no uncounted isometry fibre is discarded.

Let J0 be the F-ideal of the rank-zero G places. For fixed eta and row, all admissible central scalars lie in at most one affine coset of \(iJ_0B_\eta I_{\Sigma,F}^2\), with \(N_FB_\eta=x\), because their differences lie in this ideal. The ideal covolume is \(\sqrt{D_F}(N_FJ_0)xT^{-52}\). K103 SC26.14--16 and DO098 CF10 give \[
 \#d\le
 \frac{V_F(H/2+750Uz^2)\,U^{-f}T^{52}x^{-1}}{N_FJ_0},
 \qquad
 V_F(C)=\frac{((22/7)fC^2)^{f/2}}{(f/2)!\sqrt{D_F}},
 \tag{EC6}
\] \[
 |c(h)|<2^{-929f/2}U^{-f/2}z^{-f},\qquad
 |a_G(h)|\le2Q^{-335\cdot167}(N_EJ)^{335/2},
 \qquad N_EJ\le Q^{167}N_FJ_0 .
 \tag{EC7}
\] The full-real cover750 in EC6 is the original conservative enlargement using \(N_FJ_0\le N_EJ\); retaining J0 in the covolume denominator is compatible with that enlargement. The real-part shift is fixed after the row. This is not a claim that d itself lies in the difference ideal. The actual selected trace in EC7 retains both compact factors and the near-point transport.

Each active place has at most one actual type-II label, including a central-only pole, and each inactive place contributes its original pointwise factor. Thus all168 actual T factors together satisfy \[
                     (T-1)^{-168}<2T^{-168}.
 \tag{EC8}
\] K103 SC26.18--20 and the residual v=q26,b=j0 supply this retained-set statement. Nothing here replaces it by a coefficient on a larger extended central orbit.

\subsection{3. A full dimension336 theta bound at the larger cap}\label{a-full-dimension336-theta-bound-at-the-larger-cap}

In the original half-trace complex metric, \[
 L^\#=2\mathfrak D_E^{-1}\bar L^{-1},\qquad
 \lambda_1(L^\#)^2\ge fD_0^{-2/f}z^{-2}>
                  \frac{f}{732^2z^2}.
 \tag{EC9}
\] The exact integer inequality \(D_0<732^f\) is retained. The number732 bounds the half root discriminant, as in DO070 FT3 and K105 CG5.

For any real lattice Lambda of dimension n and minimum squared norm at least m², write \(\Theta_\Lambda(s)=\sum_{\lambda\in\Lambda}e^{-s\|\lambda\|^2}\). For \(0<t<1\), the nonzero minimum and the full Gaussian Poisson formula imply \[
 \Theta_\Lambda(s)-1
 \le e^{-(1-t)sm^2}\Theta_\Lambda(ts)
 \le e^{-(1-t)sm^2}t^{-n/2}\Theta_\Lambda(s).
 \tag{EC10}
\] The second inequality follows because the dual Gaussian in \(\Theta_\Lambda(ts)=(\pi/(ts))^{n/2}
(\operatorname{covol}\Lambda)^{-1}
\Theta_{\Lambda^*}(\pi^2/(ts))\) decreases as its argument increases. All sums converge absolutely. This directly proves the Gaussian specialization of Miller and Stephens-Davidowitz (2019), equations1.1 and2.1.

For \(0<A\le A_*\), take \[
 n=336,\quad s=\pi^2z^2/A,\quad m^2=f/(732^2z^2),\quad
 u_*=\frac{(157/50)^2}{A_*732^2}>1,\quad t=1/u_* .
\] Then \(sm^2>fu_*\). Put \(w=u_*-1\), so \(0<w<1\). The logarithm inequality \(\log(1+w)\le w-w^2/2+w^3/3\) gives \[
 -\log\bigl((eu_*)^f e^{-fu_*}\bigr)
 \ge \delta:=f(w^2/2-w^3/3)
 =\frac{4812033215271993624452}
        {120526515033315938296875}>0.
 \tag{EC11}
\] Exact rational arithmetic verifies \[
                        \delta+\delta^2/2>1/25 .
 \tag{EC12}
\] Consequently \(e^\delta>1+\delta+\delta^2/2>26/25\), so the multiplier in EC10 is less than25/26. Rearranging EC10 yields the complete scalar estimate \[
              \Theta_{L^\#}(\pi^2z^2/A)<26
                         \qquad(0<A\le A_*).
 \tag{EC13}
\] The exponent in EC10 is \(t^{-168}\). This is a dimension336 result, not a one-dimensional estimate.

The dual row metric is \(\operatorname{diag}(I_k,G)\), which dominates the ordinary product metric. Applying EC13 separately to k+3 scalar ideals gives a dual row mass at most \(26^{k+3}\). The joint Gaussian count is therefore \[
 \#r\le {164\choose k}26^{k+3}\frac{(4U)^f}{D_0^{k+3}}
            \left(\frac{\pi}{A}\right)^{(k+3)f}
            e^{AfM^2/z^2},\quad0<A\le A_* .
 \tag{EC14}
\] This counts the original complex phases directly. Apply the actual nonzero-pivot trace bound before enlarging the row set; no trace assertion is made about added zero tuples.

\subsection{4. Every finite pattern and the coupled x endpoint}\label{every-finite-pattern-and-the-coupled-x-endpoint}

Write \(d_k=k+3\), \(q=M^2/z^2\) and choose the admissible piecewise parameter \[
 A(q)=\min\{A_*,d_k/q\},\qquad
 F_k(q)=q^{-167f}(\pi/A(q))^{d_kf}e^{A(q)fq}
                    V_F(H/2+750R/q).
 \tag{EC15}
\] Combining EC4--8/14 gives a positive fixed factor times \(x^{-337/2}(N_FJ_0)^{-1}F_k(q)\). The full T power is8594 and the grouped U,R powers are \(U^{-335f/2}R^{167f}\), exactly as in K105 CG14.

Set \(\epsilon(q)=(750R/q)/(H/2+750R/q)\), with \(0<\epsilon<1\). In the constant branch, \[
 f^{-1}\frac{d\log F_k}{d\log q}=-167+A_*q-\epsilon
                                              \le-(164-k).
\] In the branch \(A=d_k/q\) the normalized derivative is \(-167+d_k-\epsilon\). The function is continuous at their common boundary.

The all-rank-zero cap is \[
 z_*=(Q^{334}T^{-52}x)^{1/(2f)},\qquad q_*=M^2/z_*^2.
\] From the exact central correlation in EC7, \[
 (N_FJ_0)^{-1}\le Q^{-167}(q/q_*)^f .
 \tag{EC16}
\] Thus the complete enlarged q factor decreases at least with exponent \(-(163-k)f\). The original local states of J have exponents2,1,0 and multiplicities1,2,1; exponent2 occurs only at rank zero. Summing all of them gives \[
 \sum_{\rm patterns}(N_FJ_0)^{-1}F_k(q)
 \le Q^{-167}F_k(q_*)\prod_G(1+q_G^{-(163-k)})^2
 <2Q^{-167}F_k(q_*),\quad k=135,136 .
 \tag{EC17}
\] Indeed \(163-k\ge27\), \(q_G\ge g_0=22285932761\), and \(2\cdot668<g_0^{27}\) proves the last product bound. No additional rank, orientation, eta or witness factor is introduced.

The recovered factor \(Q^{-167}\) at this endpoint is independent of x. Since \(q_*\) is proportional to \(x^{-1/f}\), the complete logarithmic x derivative is \[
 -3/2-A_*q_*+\epsilon(q_*)<-1/2-A_*q_*<0
\] in the constant branch and \(-3/2-d_k+\epsilon(q_*)<-1/2-d_k<0\) in the other. Therefore the whole bound, including every pattern and the moving central radius, is largest at \(x=2^{1116}\). No homothety of different actual ideals is assumed; these are uniform norm-dependent bounds.

\subsection{5. Exact endpoint and complete two-cell union}\label{exact-endpoint-and-complete-two-cell-union}

The unchanged exact host values give \[
 c_-=4124955<q_*(2^{1116})<4124956=c_+,\qquad
 (c_+)A_*<138,\qquad H/2+750R/c_-<68,\qquad V_F(68)<2^{479}.
 \tag{EC18}
\] Both claimed endpoints use the constant branch. The harmless original pivot factor \(164<2^8\) is retained. With support choice \({168\choose26}<2^{101}\), one pattern factor2, one Green2 and the single correction2 in EC8, the discrete exponent is112.

Use the directed logarithmic bounds \[
 L_T=1019333/1024,\quad L_Q=287843/4096,\quad
 L_U=601356/4096,\quad L_R=28810/4096
\] upper for T,R and lower for Q,U. Combining all powers before these substitutions gives the common exponent \[
 C=112+2f-929f/2+479+167fL_R-(335f/2)L_U
       +8594L_T-336\cdot167L_Q-(337/2)1116
   =211150679/512 .
 \tag{EC19}
\] For \(d_k=k+3\), let \(e_k\) be the least strict integral binary upper exponent of \[
 \frac{{164\choose k}26^{d_k}(22\cdot3276800)^{d_kf}}
      {(7\cdot59)^{d_kf}D_0^{d_k}c_-^{167f}} .
 \tag{EC20}
\] Using \(\log_2 e<13/9\), the complete endpoint is \(E_k=C+(13/9)fA_*c_++e_k\). The companion proves all ceilings with exact integers: \[
\begin{array}{c|r|c}
 k&e_k&E_k\\ \hline
 135&-432679&-691859209/307200<-2252\\
 136&-431349&-283283209/307200<-922\\
 137&-430020&124985591/307200>0 .
\end{array}
 \tag{EC21}
\] The last row is only an adjacent failed sufficient-bound diagnostic. For the two claimed cells, \[
             2^{-2252}+2^{-922}<2^{-921}.
 \tag{EC22}
\] There is no extra union over Gaussian branches, finite patterns or denominator parameters.

\subsection{6. Verification, exact predecessor and scope}\label{verification-exact-predecessor-and-scope}

The final certificate checks the actual host, full Gram determinant and dual domination, half-trace dual normalization, EC11--13's dimension-sensitive theta constant, both monotonicity identities on rational diagnostic samples, exact endpoint bracket, complete central volume, every directed integer-power bound, all three displayed fractions and the two-cell union. The general lattice, local, selected-trace and monotonicity arguments are the mathematical proof above; the diagnostic samples do not replace them.

The actual predecessor is canonical checks/k105\_joint\_remainder\_error.json, with111 all-top and31 uniform entries, complete floors26/11. The script accepts that predecessor or a later complete canonical extension, preserving the entire ordered arrays. Append921 once, obtaining112/31 at the K105 prefix, and keep every old overlapping charge. The exact full sums remain below2\^{}-230.

The explicit source interfaces are K105 CG1--25; K103 SC26.5--7/11--20; K104 FV1--10; DO098 CF8--10; DO070 FT3; and R076 CV8--10. The new cap, theta26 bound and EC21 cells are proved here. The actual representation, arithmetic fields, ideal classes, involution, Hermitian metric, polarization, CM theta line, canonical traces and original finite projectors are unchanged. The inherited context in K105 remains SGA3 I4.7.3/4.7.3.1; SGA6 VI Theorem3.3; FGA232 Theorem3.1 pp148--149; FGA236 Corollaries1.7/2.3 pp229/232; SGA7 I VII2.7/2.9 pp169--174 and VIII7.4(b) pp310--311; cote4 scan003 sections8.1.1--8.1.2/scans005--006 sections8.4.1--8.4.3 and cote54 scans005--007. These are carried contextual locators, not new readings or comparison claims.

Coordinates137..164, the full signed sufficient inequality, selected arithmetic-image content, prescribed Bessel/Mram value, effective admitted source construction, Read/STiesR/specific032, changed-family DP44 and RH remain separate. No Lean-build or archive-audit claim is made.

\subsubsection{Primary reference}\label{primary-reference-4}

Miller, S. D., \& Stephens-Davidowitz, N. (2019). \emph{Kissing numbers and transference theorems from generalized tail bounds} (arXiv:1802.05708v2). \url{https://arxiv.org/pdf/1802.05708}. Equations1.1 and2.1, printed pages1--2 and4, were freshly read during the preceding K105 review. EC10--13 supplies the complete required specialization at the new cap.

\subsection{Canonical algebraic-geometric context}\label{canonical-algebraic-geometric-context-31}

The original fractional invertible modules, Hermitian form, involution, polarization, theta line, paired incidences and selected complex realization are unchanged. Carried exact anchors: SGA3 I4.7.3/4.7.3.1; SGA6 VI Theorem3.3; FGA232 Theorem3.1 pp148--149; FGA236 Corollaries1.7/2.3 pp229/232; SGA7 I VII2.7/2.9 pp169--174 and VIII7.4(b) pp310--311; cote4 scan003 sections8.1.1--8.1.2, scans005--006 sections8.4.1--8.4.3 and cote54 scans005--007. No new lambda-ring/K-theory, l-adic weight, formal-group or CM-motive comparison. DP\_MULT Q3e/DP44/DP52/DP54--65 and THETA3 G1/A7 gain the stated smaller finite gate; CX\_GRH CX-ALG-024 and GCT\_PLETH T-052 retain their arithmetic gates. Primary Gaussian/Poisson context remains Miller, S. D., and Stephens-Davidowitz, N. (2019), arXiv1802.05708v2, equations1.1/2.1, printedpp1--2/4; the new-cap specialization is proved in full here.

\clearpage

\section{AM-LIFT-090 --- Two generating CM varieties control the same seed fields}\label{am-lift-090-two-generating-cm-varieties-control-the-same-seed-fields}

\textbf{Canonical checked integration --- 2026-11-23T05:12:59.068171+00:00.} Complete saved-proof and executable reviews and separate exact diagnostic replays passed. No numerical error charge is introduced. Frozen author history is preserved. AM-LIFT-090 SG1--21 sharply bounds the same actual LIFT087 seeds and LIFT088 fields through the two original generating CM varieties J and prescribed C0. Integral projective split maps and their duals put every seed n-torsion in Kgood(J{[}n{]},C0{[}n{]}) at every n, with no prime-to-index restriction; full maximal-order rank-one actions bound that field by n\^{}8. Thus the same seed and output fields cost M\^{}8 and N\^{}16, and a finite list of stages sharing J,C0,Kgood costs one M\^{}8 and one lcm(N\_j)\^{}16. The same C0{[}15{]} field has degree\textless=225 before good reduction, sharpening the same ground factor to37800\emph{h(-d)\textless=37800}d.~The complete relative bound is37800\emph{d}M\^{}8*lcm(N\_j)\^{}16. Retain the smaller LIFT089 bounds where applicable. The actual models, actions, polarizations, canonical line, finite level, Obj/FrPt and intrinsic cost remain unchanged. EffectiveN, absoluteBdegree, discriminants, coefficient heights, all nonzero-rule/source premises, selected arithmetic content, prescribed period, specific032/Read, actual DP44 and RH remain open. No error charge; current112/31 arrays and complete floors26/11 remain.

Author: /root/lift\_accounting, ALP\_MACHINE. 23 November 2026 UTC.

\textbf{Author-frozen proof.} Final whole-file, source and executable verification records are adjacent; independent review and canonical integration are separate. This is a refinement for the actual seed chosen in LIFT087 and the actual fields chosen in LIFT088. It does not replace that seed by an isogenous product, alter its polarization, or choose a new coefficient field. The isogeny scale \(N\) remains unevaluated.

\subsection{1. Statement and unchanged data}\label{statement-and-unchanged-data}

Retain the complete source Host, Datum, Level, Pins, Moduli and Coeff hypotheses of LIFT084--088. Write \(B=c.F\), \(k=\mathbb Q(\sqrt{-d})\), \(K_7=\mathbb Q(\zeta_7)\), \([E:\mathbb Q]=2f\), \(r=p+q\), and \(g=fr\). The actual level has \(M=g_0g_1\ge3\) and \(A=1\). In the same fixed algebraic closure used in LIFT088 put

\[
 K=K_{\rm good},\qquad
 J=\operatorname{Jac}(y^2=x^7-1),\qquad
 C_0(\mathbb C)\simeq\mathbb C/\mathfrak c_K .
 \tag{SG1}
\]

Here \(C_0\) is the prescribed elliptic model, with its prescribed full \(\mathcal O_k\)-action. Both \(J\) and \(C_0\), their canonical principal polarizations, their full \(\mathcal O_{K_7}\) and \(\mathcal O_k\) actions, and their everywhere-good models are defined over this \(K\). LIFT088 gives

\[
 [K:B]\le168\cdot23040\,h(-d)
          =3\,870\,720\,h(-d)\le3\,870\,720\,d .
 \tag{SG2}
\]

For \(n\ge1\) define the auxiliary containing field

\[
                         T_n=K(J[n],C_0[n]).
 \tag{SG3}
\]

This notation does not replace an actual seed field by \(T_n\). Let \(A_0\) be precisely the constructed principal seed of LIFT087.

\textbf{Theorem.} For every positive integer \(n\),

\[
       K(A_0[n])\subset T_n,\qquad [T_n:K]\le n^8.
 \tag{SG4}
\]

The inclusion remains true simultaneously for any finite list of the LIFT087 stage seeds built over the same \(K\) from the same \(J,C_0\), including their specified action twists and signature-balancing products. Consequently the very same fields of LIFT088 satisfy

\[
\begin{aligned}
 K_{\rm seed}&=K(A_0[M]),&
 [K_{\rm seed}:K]&\le M^8,\\
 K_{\rm out}&=K_{\rm seed}(A_0[N^2]),&
 [K_{\rm out}:K_{\rm seed}]&\le N^{16}.
\end{aligned}
\tag{SG5}
\]

The second line is used only under all the additional nonzero-rule, fixed-component, isogeny and source hypotheses of LIFT085--086. The resulting relative bound is

\[
            [K_{\rm out}:B]\le
                3\,870\,720\,d\,M^8N^{16}.
 \tag{SG6}
\]

There is no condition that \(n\) be prime to a Serre-module index, a polarization degree, a discriminant or the characteristic of a residue field. The torsion fields here are characteristic-zero fields; good integral models are retained from LIFT088. This is not a claim about the order of a specific Galois image or equality in SG4.

\subsection{2. Actual projective summands, with split maps}\label{actual-projective-summands-with-split-maps}

Let \(R\) be a commutative ring acting on an abelian variety \(A/K\), with every element of \(R\) acting over \(K\). For a finite projective \(R\)-module \(P\), use the same covariant Serre tensor as LIFT087:

\[
        \mathcal T_A(P)=P\otimes_R A
              :=\operatorname{Hom}_R(P^\vee,A),\qquad
        P^\vee=\operatorname{Hom}_R(P,R).
 \tag{SG7}
\]

Choose a finite projective complement \(P'\) and an actual module isomorphism \(P\oplus P'\simeq R^s\). Denote the inclusion and retraction by

\[
             i:P\longrightarrow R^s,\qquad
             p:R^s\longrightarrow P,\qquad pi=1_P.
 \tag{SG8}
\]

Functoriality of SG7 gives homomorphisms over \(K\)

\[
 I:\mathcal T_A(P)\longrightarrow A^s,\qquad
 Q:A^s\longrightarrow\mathcal T_A(P),\qquad
 QI=1_{\mathcal T_A(P)}.
 \tag{SG9}
\]

Equivalently, \(e=ip\in M_s(R)\) satisfies \(e^2=e\); the induced endomorphism \(IQ\) of \(A^s\) is its actual integral matrix action. Its image is the indicated abelian direct factor. The complementary factor is the image of \(1-e\), and addition and the two projections are inverse maps between the product of the factors and \(A^s\).

These are integral module maps and honest abelian-variety homomorphisms. Entries in \(R\) act over \(K\); no generator of a Steinitz class, rational change of basis, isogeny denominator, or extension field has been introduced. Milne (2020, Proposition 7.33 and its proof, Remark 7.34, pp.62--63) supplies precisely this projective direct-factor construction and base-change compatibility. The same maps exist on the actual good models.

Let \(F/K\) be any extension over which all points of \(A[n]\) are rational. If \(x\in\mathcal T_A(P)[n]\), then \(I(x)\in A[n]^s\) is \(F\)-rational and \[
                               x=Q(I(x)).
 \tag{SG10}
\] Thus \(x\) is \(F\)-rational. In fact restriction of \(I\) identifies this torsion with \(eA[n]^s\): the reverse inclusion follows by applying \(Q\) to an \(n\)-torsion point of \(A^s\). In particular, \[
                  K(\mathcal T_A(P)[n])\subset K(A[n]).
 \tag{SG11}
\]

This proof does not lift an \(n\)-torsion point through an arbitrary isogeny. The split identity \(QI=1\) is the reason no larger torsion order is needed.

\subsection{3. Dual factors and the actual Zarhin target}\label{dual-factors-and-the-actual-zarhin-target}

Assume \(A\) has a principal polarization \(\lambda_A:A\xrightarrow{\sim}A^\vee\) defined over \(K\). Dualizing SG9 gives

\[
 Q^\vee:\mathcal T_A(P)^\vee\longrightarrow(A^\vee)^s,\qquad
 I^\vee:(A^\vee)^s\longrightarrow\mathcal T_A(P)^\vee,
 \qquad I^\vee Q^\vee=1.
 \tag{SG12}
\]

These maps are again defined over \(K\). The isomorphism \(\lambda_A^s:A^s\xrightarrow{\sim}(A^\vee)^s\) implies that every \(n\)-torsion point of \((A^\vee)^s\) is rational over \(K(A[n])\). Applying the same split-map argument as SG10 yields

\[
           K(\mathcal T_A(P)^\vee[n])\subset K(A[n]).
 \tag{SG13}
\]

The principal polarization used here is that of the original \(J\) or \(C_0\). It is not the generally nonprincipal trace polarization on a larger tensor. Its inverse is an actual morphism, so SG13 is valid also at primes dividing that trace-polarization degree. In particular this argument does not add roots of unity to account for a dual: it uses a ground-field principal-polarization isomorphism directly. Roots that are already forced by rational torsion and the Weil pairing cause no extra field factor.

Now use the exact two modules of LIFT087:

\[
 X=\mathcal O_L\otimes_{\mathcal O_{K_7}}J,\qquad
 C=\mathfrak b^{-1}\otimes_{\mathcal O_k}C_0,\qquad
 L=E K_7.
 \tag{SG14}
\]

Both modules are finite projective over the indicated Dedekind maximal orders. The fractional nature of \(\mathfrak b^{-1}\) does not affect finite projectivity. SG11--13 show that \(X[n]\) and \(X^\vee[n]\) are rational over \(K(J[n])\), while \(C[n]\) and \(C^\vee[n]\) are rational over \(K(C_0[n])\).

The equivariant quaternion construction used in LIFT087 gives a principal polarization on the \textbf{actual variety}

\[
                            Y=(X^\vee)^4\times X^4.
 \tag{SG15}
\]

This is Zarhin's Theorem1.1 and its equation(5); equations(16)--(19) supply the isogeny, pullback and equivariance. We use the displayed target itself, not an unspecified quotient isogenous to it. The construction changes its polarization and specifies the conjugate-dual action. It does not change its \(n\)-torsion as an underlying variety. Thus \(Y[n]\) is rational over \(K(J[n])\), at every \(n\), without a prime-to-degree restriction.

The finite products, conjugations of the \(E\)-action and action twists used in LIFT087's signature balancing leave the underlying varieties in this class. In the \(d=7\) construction, a twist of the \(K_7\)-action on \(J\) is the same ground-field variety with its action precomposed by a field automorphism; all resulting endomorphisms are already defined over \(K\). A conjugated action in the general construction is likewise an action change, not an unaccounted geometric conjugate or descent twist. Products of such factors therefore have all \(n\)-torsion over \(T_n\). This proves the inclusion in SG4 for the exact seed and for every indicated stage product.

\subsection{4. The two rank-one torsion representations}\label{the-two-rank-one-torsion-representations}

Here is the required integral rank-one argument, including ramified primes. Suppose a degree-\(2a\) CM field \(F_0\) acts faithfully on an abelian variety \(A/K\) of dimension \(a\), its full ring \(R_0=\mathcal O_{F_0}\) preserves the integral Betti lattice, and every element of \(R_0\) acts over \(K\). For a fixed complex embedding of \(K\), let \(\Lambda_A=H_1(A(\mathbb C),\mathbb Z)\).

The rational vector space \(\Lambda_A\otimes\mathbb Q\) is a vector space over \(F_0\). Its rational dimension and \([F_0:\mathbb Q]\) are both \(2a\), so its \(F_0\)-dimension is one. The lattice is finitely generated and torsion-free over \(R_0\), hence finite projective of rank one over the Dedekind domain \(R_0\). There is no assertion that this rank-one projective module is globally free.

For each \(n\ge2\), factor \(nR_0\) into prime powers. At each prime, the localized rank-one projective module is free over the discrete valuation ring. Reduction gives a rank-one free module over the corresponding local Artinian quotient. Chinese remaindering combines these finitely many bases. Thus

\[
             \Lambda_A/n\Lambda_A\simeq R_0/nR_0,\qquad
             A[n](\overline K)\simeq R_0/nR_0
                         \quad\text{as }R_0\text{-modules}.
 \tag{SG16}
\]

The second isomorphism is the analytic \(n^{-1}\Lambda_A/\Lambda_A\) torsion identification, compatible with the algebraic action. These are geometric module identifications; they do not assert a rational basis of torsion over \(K\). This argument agrees with Milne's Remark7.6 for the maximal-order case, and proves the finite quotient statement directly. Absolute ramification of a prime of \(F_0\) makes the finite quotient nonreduced as a ring; it does not affect the rank-one freeness just proved.

The torsion field \(K(A[n])\) is a finite Galois extension, with faithful action on the geometric torsion. Since the full \(R_0\)-action is already defined over \(K\), that action commutes with Galois. Consequently

\[
 \operatorname{Gal}(K(A[n])/K)\hookrightarrow
    \operatorname{Aut}_{R_0}(R_0/nR_0)
       =(R_0/nR_0)^\times,\qquad
 [K(A[n]):K]\le n^{2a}.
 \tag{SG17}
\]

The right side counts all ring elements, so no surjectivity, openness, or exact CM reciprocity image is needed. In SG1, the faithful full actions have \(a=3,F_0=K_7\) and \(a=1,F_0=k\), respectively. Their degrees are at most \(n^6\) and \(n^2\). The degree of their compositum is at most the product, even when the fields intersect. Hence \([T_n:K]\le n^8\), completing SG4. For \(n=1\) all torsion fields are \(K\) and every asserted bound is one.

The same argument is valid after replacing \(K\) by any extension containing it: the full actions remain defined, and the torsion image is a subgroup of the preceding image. This base-change form is needed for the second line of SG5.

\subsection{5. Same fields, simultaneous stages, and all descended objects}\label{same-fields-simultaneous-stages-and-all-descended-objects}

For the same \(K_{\rm seed}=K(A_0[M])\), SG4 gives its first bound in SG5. No enlarged containing field is substituted for it. Over that \(K_{\rm seed}\), apply SG4 and SG17 to \(n=N^2\). Then

\[
 K_{\rm out}\subset
 K_{\rm seed}(J[N^2],C_0[N^2]),\qquad
 [K_{\rm out}:K_{\rm seed}]\le (N^2)^8=N^{16}.
 \tag{SG18}
\]

This proves SG5--6 for the exact fields already used. The complete all-object descent of LIFT088 and LIFT085 §§6--7 is unchanged: the good quotient and isogeny, full integral action, principal polarization, canonical trace-one rigidified cubical line and its normalized identities, finite \(\eta\), Moduli.rep point and rational eigenframes are all defined over those same fields. This corollary changes their degree estimates only. It does not derive algebraicity from an analytic orbit or silently choose a line or theta generator over a different field.

For a finite list of stages using the same \(J,C_0,K\) and common \(M\), define the actual composita

\[
\begin{split}
 K_\Sigma&=K\bigl(A_{0,j}[M]:j\in I\bigr),\\
 K_{\Sigma,\rm out}
   &=K_\Sigma\bigl(A_{0,j}[N_j^2]:j\in I\bigr),\qquad
 L_N=\operatorname{lcm}_{j\in I}N_j.
\end{split}
\tag{SG19}
\]

All level fields lie in the single \(T_M\), rather than in independently chosen copies. Since \(N_j^2\mid L_N^2\), every target torsion group is rational over \(K_\Sigma(J[L_N^2],C_0[L_N^2])\). Therefore

\[
\begin{split}
 [K_\Sigma:K]&\le M^8,\\
 [K_{\Sigma,\rm out}:K_\Sigma]&\le L_N^{16},\\
 [K_{\Sigma,\rm out}:B]&\le3\,870\,720\,d\,M^8L_N^{16}.
\end{split}
\tag{SG20}
\]

These are inclusions of the actual torsion fields in one fixed algebraic closure. Transporting a field embedding carries the varieties, module maps, full actions and torsion fields together; no separate extension per CM label is used. The actual eigenframes and theta generator remain those covered by the earlier all-object descent.

The common ground-field factor is counted once. For a common \(N_j=N\), use \(L_N=N\). These statements require the same prescribed elliptic curve and Jacobian and the same common good field. They do not claim this bound for arbitrary unrelated CM seeds merely sharing signatures, or for different prescribed elliptic ideal classes without proving their own common progenitor construction. The proof never replaces a seed by an arbitrary isogenous product.

\subsection{6. Exact cost change and remaining gates}\label{exact-cost-change-and-remaining-gates}

\begin{longtable}[]{@{}
  >{\raggedright\arraybackslash}p{(\columnwidth - 6\tabcolsep) * \real{0.2143}}
  >{\raggedleft\arraybackslash}p{(\columnwidth - 6\tabcolsep) * \real{0.2857}}
  >{\raggedleft\arraybackslash}p{(\columnwidth - 6\tabcolsep) * \real{0.2857}}
  >{\raggedright\arraybackslash}p{(\columnwidth - 6\tabcolsep) * \real{0.2143}}@{}}
\toprule\noalign{}
\begin{minipage}[b]{\linewidth}\raggedright
Item
\end{minipage} & \begin{minipage}[b]{\linewidth}\raggedleft
Prior sufficient bound
\end{minipage} & \begin{minipage}[b]{\linewidth}\raggedleft
Bound for the same fields here
\end{minipage} & \begin{minipage}[b]{\linewidth}\raggedright
Retained scope
\end{minipage} \\
\midrule\noalign{}
\endhead
\bottomrule\noalign{}
\endlastfoot
One stage's level field over \(K\) & \(M^{4g^2}\) in LIFT088; \(M^{fr^2}\) in LIFT089 & \(M^8\) & Exact projective Serre/Zarhin seed \\
One target field over its seed field & \(N^{8g^2}\) in LIFT088; \(N^{4fr^2}\) in LIFT089 & \(N^{16}\) & All additional LIFT085--086 hypotheses \\
A finite simultaneous list's common level field & A product of stage bounds & One \(M^8\) & Same \(J,C_0,K\) and common level \\
A finite simultaneous list's target fields & A product of target bounds & \(L_N^{16}\) & \(L_N\) is the actual lcm; it remains unevaluated \\
\end{longtable}

There is also an immediate refinement of the same ground field in SG2. LIFT088 defines \(K=K_J(C_0[15])\), where the full \(\mathcal O_k\)-action on \(C_0\) is already defined over \(B(j_*)\subset K_J\). The proof of SG17 needs no good-reduction hypothesis, so it applies before adjoining this torsion. Keeping the existing \(K_J\) and \(K\) gives \[
 [K:K_J]\le15^2=225,\qquad
 [K:B]\le168\cdot225\,h(-d)
           =37\,800\,h(-d)\le37\,800\,d .
 \tag{SG21}
\] Thus every occurrence of the conservative factor \(3\,870\,720\,d\) in SG6 and SG20 may also be replaced by \(37\,800\,d\). This changes no good-model construction or field. The original SG2 remains a valid carried bound.

When a previous exponent is smaller, take the minimum of the two valid bounds. In particular, combining with LIFT089 gives \(M^{\min(8,fr^2)}\) and \(N^{\min(16,4fr^2)}\) for one stage. No prior field or descent statement is weakened.

In the single-stage form, the logarithm of the sufficient relative degree is now at most \[
       \log(3\,870\,720)+\log d+8\log M+16\log N.
\] The intrinsic eigenline lower bound \(b\) from LIFT087 and the exact optimized Step5 cost from LIFT086 are unchanged. The latter still contains its \(\log N\) term. Absolute degrees still require \([B:\mathbb Q]\). No discriminant, coefficient-height, or normalized scalar bound is obtained merely from these relative degrees.

The main effective gate remains \(N\). The actual prescribed rule must still be nonzero, and all Host/Sizes, component, Pins, Moduli, model, polynomial, cusp, KeyCore and comparison hypotheses retain their original quantifiers. No prescribed arithmetic selected-image witness, global central value or period, Read/STiesR/032 instance, DP44 sign or RH statement follows.

\subsection{7. Sources and verification boundary}\label{sources-and-verification-boundary-5}

The exact actual seed and all its underlying products are LIFT087 GS3--7, GS19 and §7. The same-field construction and the precise definitions of \(K,K_{\rm seed},K_{\rm out}\) are LIFT088 SF8--15 and §§4--7. The quotient/action/line/level/FrPt descent is LIFT085 HF21--26 and §7, with LIFT086's finite-flat good-model transport. The literal interfaces retained are TwoWeightEqOneAddRoot.lean, Datum30--49, Level51--61, TateIso131--142 and Obj144--160; PushInvertible.lean, Moduli29--38 and Coeff40--46; \nolinkurl{PartitionExtendByZeroAntitone.lean,} PinLevel43; IsCombAt.lean, FrPt67--76. The source manifest records exact byte hashes and the actual read scopes.

Milne, J. S. (2020). \emph{Complex multiplication} (version0.10). Proposition7.33 and its complete proof, Remark7.34, pp.62--63; Remark7.6, p.53. \url{https://www.jmilne.org/math/CourseNotes/CM.pdf}

Zarhin, Y. G. (2023). \emph{Abelian varieties, quaternion trick and endomorphisms}. Author arXiv version5, Theorem1.1 and equation(5), printedp.2; §4, equations(16)--(19), printedpp.5--6. \url{https://arxiv.org/pdf/2010.07995}

The fresh ingredients are the explicit split-map/dual-map torsion inclusion and the use of the two original rank-one CM torsion representations for all stages at once. The relative Picard/Poincare, cubical line, finite-flat quotient and actual good models are carried with their previously checked scopes. No new motive, weight, slope or operator identity is claimed.

The companion exact diagnostics check integral idempotent and dual splitting models at composite torsion orders, finite rank-one ring automorphisms, nonprincipal ideal quotients, common-lcm containment and field-exponent arithmetic, including SG21, together with all source fingerprints. They supplement the universal proof; they do not compute the actual Galois groups or evaluate \(N\). Final whole-file review, exact executable replay, frozen hashes and independent review are separate records.

\subsection{Canonical algebraic-geometric context}\label{canonical-algebraic-geometric-context-32}

The finite-etale torsion modules carry the actual OE-action and perfect polarization pairing. The two literal idempotent halves are Cartier dual through the existing roots of unity; all descended abelian schemes and cubical Picard/Poincare data are those already constructed. Carried exact anchors: FGA232 Theorem3.1 pp148--149; FGA236 Corollaries1.7/2.3 pp229/232; SGA7 I VII2.7/2.9 pp169--174 and VIII7.4(b) pp310--311; SGA3 I4.7.3/4.7.3.1; SGA6 VI Theorem3.3; cote4 scan003 sections8.1.1--8.1.2 and scans005--006 sections8.4.1--8.4.3; cote54 scans005--007. No new formal-group slope, lambda-ring/K-theory identity, l-adic weight theorem or CM-motive comparison follows. DP\_MULT Q3e/DP44/DP52/DP54--65, THETA3 G1/A7, CX\_GRH CX-ALG-024 and GCT\_PLETH T-052 retain their separate arithmetic comparisons. The new split inclusions and their duals are integral homomorphisms of the existing abelian schemes; finite-etale geometric torsion fields are bounded through the two original rank-one maximal-order Galois representations. No new CM reciprocity or motive theorem is required.

\clearpage

\section{AM-K-106 --- the actual radial profile and every nonzero cross coordinate}\label{am-k-106-the-actual-radial-profile-and-every-nonzero-cross-coordinate}

\textbf{Canonical checked integration --- 2026-11-23T05:13:37.724136+00:00.} Complete root and independent proof/script/interface reviews and separate exact replays passed. The actual canonical predecessor has112/31 entries; this integration gives113/31, complete floors26/11. Frozen author history is preserved. AM-K-106 NT1--14 removes the complete coordinate sector k137 from the actual all-II globally rank-two-unipotent support26 scalar-cross remainder. The exact universal radial-profile factor carries through the original rank-two near-point trace and yields630 bits. Every nonzero original S coordinate lies in the same fractional ideal; the proved product-of-sums inequality gives k\^{}-84 and596 further bits at137. Applied before the full DO112 Gaussian cardinality enlargement, these constant factors preserve every actual metric, finite weight, central affine covolume and complete pattern/denominator count. The exact whole-sector exponent is-251641609/307200\textless-819. Actual DO112112/31 becomes113/31, all previous overlapping charges retained, both complete sums below2\^{}-230 and floors26/11 unchanged. The only possible support26 branch now has138\textless=k\textless=164 with the same v=q26,b=j0 and2\^{}1116\textless x\textless2\^{}2108. Signed positivity, selected arithmetic content, prescribed period, full effective source, specific032/Read, actual DP44 and RH remain open.

Author: /root/kummer\_extension, GPT-6 Astra Pro, ALP\_MACHINE. 23 November 2026 UTC. Controlled pending proof. Final author verification, exact replay, independent review and canonical integration are distinct records. No earlier frozen proof is amended.

\subsection{1. Exact retained sector and new assertion}\label{exact-retained-sector-and-new-assertion}

Keep the actual J56 all-top family at H=133, the original fields E/F with f=168, the actual selected compact module and Hilbert norm, every allowed word and canonical twist, both compact factors, the prescribed near-point argument, paired source/target pencils, finite projectors, central phases, Haar measures and class masses. Apply the new count in the all-II support26 global rank-two-unipotent scalar-cross remainder after K103, K104, K105 and DO112. Thus \[
 s=26,\quad v=q_{\rm elig}=26,\quad
 2^{1116}<x=N_ED_\eta<2^{2108},\quad b=j=0,\quad137\le k\le164.
 \tag{NT1}
\] The integer k counts nonzero original S coordinates of the actual common covector. The full positive compression, N²Z=0 and original normal form are proved in K103 SC26.5--7, and its SC26.11 proves uniqueness of eta for each support in this denominator window.

\textbf{Theorem.} The complete subset of NT1 with k=137 has absolute normalized contribution less than \(2^{-819}\). Append this one fresh charge while preserving every older charge and its order. The remaining coordinate range is \(138\le k\le164\), with the other restrictions in NT1 unchanged. There is no support27 conclusion or existence claim for the complement.

The improvement has two separately proved ingredients: a carried radial-profile factor that applies before the rank-two near-point comparison, and a new field-norm inequality using all k nonzero original coordinates. Neither ingredient changes the Gaussian lattice count.

\subsection{2. The carried radial profile is valid for the rank-two argument}\label{the-carried-radial-profile-is-valid-for-the-rank-two-argument}

Write \(a=168\) for the positive real-group rank and \(f=168\) for the number of CM labels; the equal numerical values play different roles. DO080 bounds the genuine normalized selected trace \[
 c_W(g)=\frac{\operatorname{Tr}_W(P_W\pi(g)|_W)}{\dim W},
\] with its full vector-valued metric. In DO080 RT9, at radial coordinate Y=y, the conditional density of the first complex coordinate \(z_1=u+it\), with \(R_y=\sqrt{1-y}\), is \[
 \frac{a-1}{\pi R_y^2}
 \left(1-\frac{u^2+t^2}{R_y^2}\right)^{a-2}
 {\bf1}_{u^2+t^2<R_y^2}.
 \tag{NT2}
\] Its Cartan change of variable satisfies the same positive Jacobian estimate \[
 \left|\frac{du}{dy'}\right|
 \le \frac{\sqrt y}{2\sinh(r)y'^2}.
\] Retain the pointwise bound on the density factor by \((1-t^2/R_y^2)^{a-2}\), and enlarge the admissible t interval to \([-R_y,R_y]\). Compared with the uniform density used in DO080 RT12, its exact multiplier is \[
 \vartheta_a=\frac12\int_{-1}^1(1-t^2)^{a-2}\,dt
 =\frac{4^{a-1}}{2(a-1)\binom{2a-2}{a-1}}.
 \tag{NT3}
\] This is the canonical DO092 SN16--18 argument. The elementary integral follows by integration by parts from \(I_0=2\) and \(I_j=2j I_{j-1}/(2j+1)\).

The multiplier is independent of y,y' and of the two compact factors of g. The forward and reverse transition kernels each gain \(\vartheta_a\), so their geometric mean in DO080 RT14 gains exactly \(\vartheta_a\), not its square root. Detailed balance and the same variance bound therefore give \[
 |c_W(g)|<
 \frac{\vartheta_a\,2^{-931/2}}{\sinh r}\quad(r>0).
 \tag{NT4}
\] At r=0 retain the bound one. This is an estimate for the actual selected trace, not an operator-norm estimate. No global rank-one assumption enters NT2--4: DO092 uses rank one only in the subsequent application SN19, which is not used here.

Let \(N=h-I\) for the actual globally rank-two-unipotent matrix, and put \(\alpha=(N^2)_{nn}\). The original DO080 section7 calculation gives, at each embedding, \[
 \alpha>0,\qquad
 \alpha_p=(g_p^{-1}N^2g_p)_{nn}>\alpha/3,\qquad
 \sinh^2 r(g_p^{-1}hg_p)\ge\alpha_p .
 \tag{NT5}
\] It uses the actual near-point radius \(r_p<1/2\) and an archimedean normalized null vector only. No such vector is made integral at finite primes. Applying NT4 at this same argument yields \[
 |c_\tau(g_p^{-1}hg_p)|
 <\vartheta_a\,2^{-929/2}\alpha_\tau^{-1/2}.
\] Consequently the genuine product trace satisfies \[
 |c(h)|<
 \vartheta_{168}^{168}\,2^{-929f/2}
 \operatorname N_{F/\mathbb Q}(\alpha)^{-1/2},
 \qquad
 \vartheta_{168}^{168}<2^{-630}.
 \tag{NT6}
\] The last comparison is exact: \(\vartheta_{168}^4<2^{-15}\), followed by the power42. No compact intertwiner, word-dependent metric, source Haar factor or near-point loss has been discarded.

\subsection{3. The norm of the sum retains every nonzero coordinate}\label{the-norm-of-the-sum-retains-every-nonzero-coordinate}

Retain the original covector ideal and scale \[
 L=JI_\Sigma D_\eta,\qquad N_EL=z^{2f},
\] as in K103 SC26.13. Its k nonzero S coordinates are \(\xi_1,\ldots,\xi_k\in L\). In the fixed original metric, \[
 \alpha/U=rJ_Z^{-1}r^*
       \ge \sum_{i=1}^k|\xi_i|^2
 \quad\hbox{at every CM label},\qquad U=77Q^2-14.
 \tag{NT7}
\] This is K103 SC26.7 / DO104 CD10. The transverse contribution has the actual positive metric G\^{}-1 and is merely discarded from this lower bound; no rational moving basis is declared integral.

Every nonzero \(\xi_i\in L\) satisfies \[
 \prod_{\tau\ {\rm CM}}|\xi_{i,\tau}|^2
       =|N_E\xi_i|\ge N_EL=z^{2f}.
 \tag{NT8}
\] Indeed \((\xi_i)L^{-1}\) is an integral ideal. This is a global norm statement; it does not claim a separate lower bound on each complex absolute value.

Here is the required direction of the sum inequality. For positive \(a_{i,\tau}\), put \(s_\tau=\sum_i a_{i,\tau}\). By arithmetic-geometric mean, \[
 \sum_i\left(\prod_\tau\frac{a_{i,\tau}}{s_\tau}\right)^{1/f}
 \le\frac1f\sum_\tau\sum_i\frac{a_{i,\tau}}{s_\tau}=1.
\] Multiplying by \((\prod_\tau s_\tau)^{1/f}\) proves \[
 \left(\prod_\tau\sum_i a_{i,\tau}\right)^{1/f}
       \ge \sum_i\left(\prod_\tau a_{i,\tau}\right)^{1/f}.
 \tag{NT9}
\] Apply this with \(a_{i,\tau}=|\xi_{i,\tau}|^2\), which is positive because each chosen coordinate is a nonzero field element. Equations NT7--9 imply \[
 \operatorname N_F(\alpha)\ge
 U^f\prod_\tau\sum_i|\xi_{i,\tau}|^2
 \ge U^f(kz^2)^f .
\] Substitution in NT6 gives the new uniform bound \[
 \boxed{\quad
 |c(h)|<
 2^{-630}\,k^{-f/2}\,
 2^{-929f/2}U^{-f/2}z^{-f}.
 \quad} \tag{NT10}
\] This holds for every actual k-sector satisfying the stated normal form and ideal membership. It is applied before enlarging the set of row tuples in a Gaussian majorant. No assertion about the trace of the added zero tuples is needed. The new factor is independent of x, z and the finite G pattern, so both complete monotonicity arguments remain unchanged.

\subsection{4. The full k=137 Gaussian count}\label{the-full-k137-gaussian-count}

Use the larger dimension336 Gaussian cap of DO112, \[
 A_*=\frac{59}{50\cdot65536},\qquad
 \Theta_{L^\#}(\pi^2z^2/A)<26\quad(0<A\le A_*).
 \tag{NT11}
\] The half-trace dual is \(2\mathfrak D_E^{-1}\bar L^{-1}\). Its minimum, the full factor \(t^{-168}\), and the rational theta inequality are the universal EC9--13 proof, not a one-dimensional estimate. For d\_k=k+3=140, the true coupled primal metric diag(I\_k,G\^{}-1) has covolume \(D_0^{d_k}z^{2d_kf}/(4U)^f\); the dual Gram diag(I\_k,G) dominates the scalar product metric. Thus the full row mass costs \(26^{d_k}\) and the Gaussian row bound is EC14 with k=137.

All other factors are retained individually: the central covolume divisor \(N_FJ_0\), central radius \(H/2+750Uz^2\), Green estimate \(2Q^{-335\cdot167}(N_EJ)^{335/2}\), original support count \(\binom{168}{26}<2^{101}\), original row-position count \(\binom{164}{137}\), conservative pivot factor \(164<2^8\), both mixed projectors, and all168 actual T factors \((T-1)^{-168}<2T^{-168}\). Eta is fixed once by the support's uniqueness theorem; no beta or separate eta enumeration is added.

For clarity, the weighted G argument is checked anew at this k. The exact correlation \(N_EJ\le Q^{167}N_FJ_0\) and the complete q derivative from EC15--17 give exponent \(163-k=26>0\). Hence \[
 \prod_G(1+q_G^{-26})^2<2,\qquad
                 2\cdot668<g_0^{26}.
 \tag{NT12}
\] The all-rank-zero endpoint retains the factor \(Q^{-167}\); the old2\^{}668 pattern count is not also included. The complete x derivative is still strictly negative on both branches of \(A(q)=\min(A_*,d_k/q)\). The fixed multiplier in NT10 does not affect it. There is one whole endpoint at \(x=2^{1116}\).

At that endpoint, \(c_-=4124955<q_*<4124956=c_+\), \(A_*c_+<140\), the central radius is below68 and \(V_F(68)<2^{479}\). These are the same rational inequalities, explicitly rechecked for the new cell. The complete discrete exponent remains112.

\subsection{5. One exact fresh charge and chronology}\label{one-exact-fresh-charge-and-chronology}

Use the unchanged directed constants and grouped common exponent from K105 CG21 / DO112 EC19, \[
 C_G=211150679/512.
\] The k=137 rational factor is \[
 W_{137}=
 \frac{\binom{164}{137}26^{140}
                  (22\cdot50\cdot65536)^{140f}}
 {7^{140f}59^{140f}D_0^{140}c_-^{167f}}.
\] Its least strict integral binary upper exponent is -430020. With the unchanged positive-exponential bound \(\log_2e<13/9\), the complete old-trace cell has exponent \[
 C_G+\frac{13}{9}fA_*c_+-430020
                         =124985591/307200.
 \tag{NT13}
\] This calculation includes all patterns, supports, labels, phases, central fibres and the whole denominator window. It is independently recomputed; a failed next-cell diagnostic is not itself being treated as a theorem.

The exact integer inequalities \[
 \vartheta_{168}^{168}<2^{-630},\qquad
                      137^{84}>2^{596}
\] and NT10 now give \[
 E_{137}<\frac{124985591}{307200}-630-596
       =-\frac{251641609}{307200}<-819.
 \tag{NT14}
\] There is one k-cell, and its support, row-position, G-pattern and label sums are already inside NT13. No additional multiplicity or hidden endpoint union is required. This proves the complete charge819.

The final chronological replay must use canonical DO112, preserving its complete112-entry all-top prefix and the unchanged31-entry uniform array. Append819 once to obtain113/31, with the complete floors26/11 unchanged. All old charges, including K1051465 and DO112921, remain. No bound for k138..164 is claimed here.

\subsection{6. Source and verification boundaries}\label{source-and-verification-boundaries}

The actual trace and near-point interface is canonical DO080 RT1--3 and sections3--7. The radial sharpening is canonical DO092 SN16--18, whose universal coefficient statement is distinguished from its later rank-one application SN19. The original energy and row ideals are K103 SC26.7/13 and DO104 CD10--13/21; the exact coupled Gram and Gaussian conventions are K105 CG5--18 and DO112 EC9--18. This proof supplies the all-coordinate norm inequality, the explicit rank-two transfer of the radial factor and the complete fresh k137 count.

The underlying representation-theoretic primary input remains Nakahama, R. (2016), \emph{Norm computation and analytic continuation of vector valued holomorphic discrete series representations}, Journal of Lie Theory,26(4),927--990, equation5.7 and Theorem5.1, pp960--961, \url{https://arxiv.org/abs/1506.05919}, through the checked exact K043/K046/DO080 dictionary. The Gaussian primary convention remains Miller, S. D., and Stephens-Davidowitz, N. (2019), \emph{Kissing numbers and transference theorems from generalized tail bounds}, arXiv:1802.05708v2, equations1.1/2.1, \url{https://arxiv.org/pdf/1802.05708}. No scalar-fibre substitute or new representation-theoretic comparison is used.

The companion recomputes every new rational and integer inequality and preserves the full ordered error arrays. Finite positive-matrix examples supplement, but do not replace, the universal AM-GM proof NT9. The universal source and metric statements are mathematical proof obligations, not consequences of a finite replay.

The same original field, involution, polarization, theta line, Picard/biextension data, finite projectors and selected coefficient are retained. Exact carried context remains SGA3 I4.7.3/4.7.3.1; SGA6 VI Theorem3.3; FGA232 Theorem3.1 pp148--149; FGA236 Corollaries1.7/2.3 pp229/232; SGA7 I VII2.7/2.9 pp169--174 and VIII7.4(b) pp310--311; cote4 scan003 sections8.1.1--8.1.2/scans005--006 sections8.4.1--8.4.3 and cote54 scans005--007. No new motivic or arithmetic comparison is asserted. Dense k138..164, signed positivity, prescribed selected arithmetic-image content, global Bessel/Mram value, effective source construction, Read/STiesR/032, evaluated DP44 and RH remain separate.

\subsection{Canonical algebraic-geometric context}\label{canonical-algebraic-geometric-context-33}

The original fractional invertible modules, Hermitian form, involution, polarization, theta line, paired incidences and selected complex realization are unchanged. Carried exact anchors: SGA3 I4.7.3/4.7.3.1; SGA6 VI Theorem3.3; FGA232 Theorem3.1 pp148--149; FGA236 Corollaries1.7/2.3 pp229/232; SGA7 I VII2.7/2.9 pp169--174 and VIII7.4(b) pp310--311; cote4 scan003 sections8.1.1--8.1.2, scans005--006 sections8.4.1--8.4.3 and cote54 scans005--007. No new lambda-ring/K-theory, l-adic weight, formal-group or CM-motive comparison. DP\_MULT Q3e/DP44/DP52/DP54--65 and THETA3 G1/A7 gain the stated smaller finite gate; CX\_GRH CX-ALG-024 and GCT\_PLETH T-052 retain their arithmetic gates. Primary Gaussian/Poisson context remains Miller, S. D., and Stephens-Davidowitz, N. (2019), arXiv1802.05708v2, equations1.1/2.1, printedpp1--2/4; the new-cap specialization is proved in full here. The sharper trace is the same genuine vector-valued selected normalized trace and original near-point argument; no operator-norm or rank-one transport substitute is used. Nakahama(2016), Journal of Lie Theory26(4),927--990, equation5.7 and Theorem5.1 printed960--961, supplies the carried representation convention.

\clearpage

\section{AM-LIFT-091 --- relative root discriminants of the same CM torsion fields}\label{am-lift-091-relative-root-discriminants-of-the-same-cm-torsion-fields}

\textbf{Canonical checked integration --- 2026-11-23T05:28:47.119486+00:00.} Complete saved-proof and executable reviews and separate exact diagnostic replays passed. No numerical error charge is introduced. Frozen author history is preserved. AM-LIFT-091 RD1--18 bounds root discriminants of the same actual LIFT088--090 finite torsion fields. For a finite Galois extension of actual degree D, unramified outside an integral ideal S, rd(L)/rd(K)\textless=D\emph{Norm(rad S)\^{}(1/{[}K:Q{]}); the complete base-normalized local different and actual-degree wild contribution are proved. Thus the simultaneous seed satisfies rd(KSigma)\textless=rd(Kgood)}rad(M)\emph{M\^{}8, and the same actual output satisfies rd(KSigma,out)\textless=rd(Kgood)}rad(M L\_N)\emph{lcm(M,L\_N²)\^{}8\textless=rd(Kgood)}M\^{}9\emph{L\_N\^{}17. Its relative increase over KSigma is at most rad(L\_N)}L\_N\^{}16. The literal N\_j=1 mod M congruence is retained, so no false level/output-overlap saving is claimed. Full all-object descent and the37800\emph{d}M\^{}8*L\_N\^{}16 relative degree bound are unchanged. Initial rd(Kgood), absolute B degree/discriminant, effective N\_j, heights, all nonzero-rule/source conditions, selected content, prescribed period, specific032/Read and evaluated DP44 remain open. No numerical error charge;113/31 full arrays and floors26/11 remain.

Author: /root/lift\_accounting, ALP\_MACHINE. 23 November 2026 UTC.

\textbf{Author-frozen proof.} The final whole-file, source and executable verification is recorded in the adjacent author packet. This is a consequence for the exact fields of canonical LIFT088--090. It does not replace the given coefficient field, the chosen polarized seed, its good model, or the target selected by LIFT085. The isogeny scale remains unevaluated. The initial root discriminant is retained as an explicit input.

\subsection{1. Statement for the actual fields}\label{statement-for-the-actual-fields}

Retain every hypothesis and object of LIFT087--090, including the given Coeff field B, actual Host and Datum, prescribed elliptic factor C0, fixed J, all projective Serre summands and their duals, equivariant principalization, finite level M, actual good models and full integral actions. Set

\[
 K=K_{\mathrm{good}},\qquad
 T_n=K(J[n],C_0[n]),\qquad
 \operatorname{rd}(F)=|\operatorname{disc}(F/\mathbf Q)|^{1/[F:\mathbf Q]} .
 \tag{RD1}
\]

For a positive integer n, write rad(n) for the product of its distinct rational prime divisors, with rad(1)=1. All fields are subfields of the same fixed algebraic closure.

For one stage the exact fields remain

\[
 K_{\mathrm{seed}}=K(A_0[M]),\qquad
 K_{\mathrm{out}}=K_{\mathrm{seed}}(A_0[N^2]).
 \tag{RD2}
\]

The second field is used only under all the nonzero-rule, component, isogeny and source hypotheses of LIFT085--086. For a finite list of stages built from the same J, C0 and K, with common level M, use the actual fields

\[
\begin{split}
 K_\Sigma&=K(A_{0,j}[M]:j\in I),\\
 K_{\Sigma,\mathrm{out}}
   &=K_\Sigma(A_{0,j}[N_j^2]:j\in I),\\
 L_N&=\operatorname{lcm}_{j\in I}N_j,\qquad
 \ell=\operatorname{lcm}(M,L_N^2).
\end{split}
\tag{RD3}
\]

The lcm of an empty list is one. The assertions below are bounds for these fields, not a change to the containing fields T\_n.

The literal LIFT085 HF20--23 construction also gives N\_j=1 modulo M and makes every N\_j prime to its prescribed protected primes. In particular gcd(M,L\_N)=1 for the actual target list, so ell=M L\_N\^{}2. We retain the general lcm expression in the formulas, but do not claim a level/target overlap saving that this exact source congruence excludes.

\textbf{Theorem.} All these extensions over K are finite Galois. The same actual fields satisfy

\[
\begin{split}
 \operatorname{rd}(K_\Sigma)
   &\le \operatorname{rd}(K)\,\operatorname{rad}(M)M^8,\\
 \operatorname{rd}(K_{\Sigma,\mathrm{out}})
   &\le \operatorname{rd}(K)\,
       \operatorname{rad}(M L_N)\,[K_{\Sigma,\mathrm{out}}:K]\\
   &\le \operatorname{rd}(K)\,
       \operatorname{rad}(M L_N)\ell^8\\
   &\le \operatorname{rd}(K)\,M^9L_N^{17}.
\end{split}
\tag{RD4}
\]

The first line also applies to the one-stage seed. For a single target, take L\_N=N. A useful separate same-field relative bound is

\[
 \frac{\operatorname{rd}(K_{\Sigma,\mathrm{out}})}
      {\operatorname{rd}(K_\Sigma)}
       \le \operatorname{rad}(L_N)L_N^{16}.
 \tag{RD5}
\]

The direct second line of RD4 retains the actual support and degree in one inequality. For these literal source scales the level and target supports are disjoint; shared-stage target primes are still counted once through L\_N. None of these estimates requires K or B to be Galois over Q, unramified over Q, or CM. The local ramification above primes dividing M or L\_N is fully included.

\subsection{2. Exact local different inequality}\label{exact-local-different-inequality}

Let L\_w/K\_v be a finite extension of finite extensions of Q\_p.~Let ord\_v and ord\_w be the integer valuations normalized by their respective uniformizers, and let e=e(w/v). Thus ord\_w(a)=e ord\_v(a) for a in K\_v. Write

\[
 a_w=\operatorname{ord}_w(\mathfrak D_{L_w/K_v}).
\]

Then

\[
                 a_w\le e-1+e\,\operatorname{ord}_v(e).
 \tag{RD6}
\]

Here ord\_v(e) is the valuation of the integer e in K\_v; it equals e(K\_v/Q\_p) ord\_p(e), not generally ord\_p(e).

To prove RD6, let K'\_v be the maximal unramified intermediate field. Its normalized valuation restricts to ord\_v and its different is trivial; the different of L\_w/K'\_v equals that of L\_w/K\_v. The extension L\_w/K'\emph{v is totally ramified of degree e. Choose a uniformizer pi of L\_w. The integral ring is O}\{K'\_v\}{[}pi{]}, and pi has an Eisenstein minimal polynomial

\[
 f(X)=X^e+a_{e-1}X^{e-1}+\cdots+a_0.
 \tag{RD7}
\]

The different is generated by f'(pi). Put a\_e=1 and ignore zero derivative terms. The valuation of its i-th term, for 1\textless=i\textless=e, is

\[
 \operatorname{ord}_w(i a_i\pi^{i-1})
       =e\,\operatorname{ord}_{K'_v}(i a_i)+i-1.
 \tag{RD8}
\]

These valuations have distinct residues modulo e. Consequently the nonzero terms have a unique minimum, and there is no cancellation at that minimum. The valuation of their sum is the minimum and is at most the valuation of the leading derivative term e pi\^{}\{e-1\}. This is e ord\_v(e)+e-1, proving RD6. The case e=1 is included: the extension is unramified and the different exponent is zero.

The exact primary inputs are Crew (2019), Proposition1.2.5.1, printed p.30 (integral Eisenstein generator); Proposition2.1.3.2, p.54 (different generated by the derivative); Corollary2.1.4.3, p.55 (removing the unramified intermediate field); and the different tower in Proposition2.1.4.5, pp.55--56. The upper bound RD6 is proved here. The nearby lower-bound ramification lemma is not being used as an upper estimate.

An instructive normalization check is K\_v=Q\_2(sqrt(2)) and L\_w=Q\_2(2\^{}(1/4)). The latter is a degree-two Eisenstein extension of K\_v. Its derivative 2pi has ord\_w equal to 5. RD6 gives 2-1+2*ord\_v(2)=5, whereas incorrectly substituting ord\_2(2)=1 would give 3.

\subsection{3. A finite-prime ideal-support bound}\label{a-finite-prime-ideal-support-bound}

Let L/K be a finite Galois extension of number fields of degree D. Let S be a nonzero integral ideal of O\_K such that L/K is unramified at every finite prime not dividing S. Define rad(S) to have exponent one at every prime ideal dividing S, and put d\_K={[}K:Q{]}. Then

\[
 \boxed{\quad
 \log\operatorname{rd}(L)-\log\operatorname{rd}(K)
 \le \frac{\log N_{K/\mathbf Q}\operatorname{rad}(S)}{d_K}
        +\log D .
 \quad}
 \tag{RD9}
\]

The finite primes themselves, not all primes above their rational residue characteristics, occur in the first term.

Here are the complete normalizations. For a prime v of K, Galoisness makes the local ramification index e\_v, residue degree f\_v and different exponent a\_v independent of w above v. If there are g\_v such primes, then g\_v e\_v f\_v=D. Localizing the trace pairing and using the different as inverse trace dual gives the relative discriminant exponent

\[
 \operatorname{ord}_v(\mathfrak d_{L/K})
     =\sum_{w\mid v} f(w/v)a_w=g_v f_v a_v.
\]

The different tower and ideal norm give

\[
 |\operatorname{disc}(L/\mathbf Q)|
   =|\operatorname{disc}(K/\mathbf Q)|^D
       N_{K/\mathbf Q}\mathfrak d_{L/K}.
\]

Taking normalized logarithms proves the exact identity

\[
 \log\operatorname{rd}(L)-\log\operatorname{rd}(K)
   =\frac1{d_K}\sum_v\frac{a_v}{e_v}\log N(v).
 \tag{RD10}
\]

This can also be derived by applying the local different tower over Q\_p at every w and summing its valuations with absolute residue degrees. The trace-dual identification and local discriminant norm are Crew §2.1.5, equations(2.1.5.1)--(2.1.5.2), p.56, and §2.1.6, equation(2.1.6.3), p.58. These local identities recover the global discriminant formula by localization; no globally free relative integral basis is assumed.

Only v dividing S contributes to RD10. By RD6 its contribution is at most

\[
 \frac1{d_K}\sum_{v\mid S}
     \left(1-\frac1{e_v}+\operatorname{ord}_v(e_v)\right)\log N(v).
 \tag{RD11}
\]

Since e\_v divides the Galois degree D, the integer D/e\_v is integral and ord\_v(e\_v)\textless=ord\_v(D). The tame term in RD11 is at most log N(rad S)/d\_K. The entire wild term is bounded by

\[
 \frac1{d_K}\sum_{v\mid S}\operatorname{ord}_v(D)\log N(v)
 \le \frac1{d_K}\log N(D O_K)=\log D.
 \tag{RD12}
\]

This proves RD9. It also proves the sharper version retaining the actual local tame terms and the restricted principal ideal (D)\_S in RD11--12.

In rational-prime notation, ord\_v(D)=e(K\_v/Q\_p) ord\_p(D), log N(v)=f(K\_v/Q\_p)log p, and sum\_\{v\textbar p\}e(K\_v/Q\_p)f(K\_v/Q\_p)=d\_K. Thus the wild contribution above a rational p is at most ord\_p(D)log p.~The factor d\_K cancels exactly. There is no factor equal to the number of primes over p and no assumption of absolute unramifiedness.

For S=n O\_K, or for a subset of that support,

\[
 \frac{\log N(\operatorname{rad} S)}{d_K}
       \le\log\operatorname{rad}(n).
 \tag{RD13}
\]

Indeed at p dividing n, sum\_\{v\textbar p\}f(K\_v/Q\_p)\textless=d\_K. This is only a convenient weaker specialization; RD9 retains the actual ideal support whenever it is known.

\subsection{4. Ramification of the actual finite torsion fields}\label{ramification-of-the-actual-finite-torsion-fields}

The exact seed A0 and every stage seed A0,j have everywhere-good abelian-scheme models over O\_K by LIFT088 SF8--13. This is the same K=K\_good as in LIFT090, not merely a potentially good field.

At a finite prime v of residue characteristic p not dividing n, the good model's n-torsion is finite etale. Its geometric points are fixed by inertia. Equivalently, Milne (2020), Corollary6.11 and the paragraph immediately following it, printed p.51, says precisely that good reduction makes prime-to-p torsion unramified. Hence

\[
 K(A_{0,j}[n])/K
       \text{ is unramified at every }v\nmid n.
 \tag{RD14}
\]

This implication alone is needed; we do not apply a converse reduction criterion here. It remains true after every finite base extension because the good model base changes to an abelian scheme. It applies to every finite torsion order n, including prime powers, at primes whose residue characteristic does not divide n; arbitrary absolute ramification of K is allowed.

The field K(A0,j{[}n{]}) is the fixed field of the kernel of the finite Galois action on the entire n-torsion group. It is therefore finite Galois. Finite composita are also Galois. The actual fields in RD2--3 are such composita over K. Away from their indicated torsion orders, inertia acts trivially on every constituent; it acts trivially on the compositum. Thus K\_Sigma/K is unramified outside M O\_K and K\_\{Sigma,out\}/K is unramified outside M L\_N O\_K. Over K\_Sigma the target extension is unramified outside L\_N O\_\{K\_Sigma\}.

This argument does not claim an infinite adelic level over a finite field. It concerns only the finite torsion groups actually used in LIFT088--090. It does not claim unramifiedness at a protected prime dividing M or L\_N.

\subsection{5. Degree and support are combined for the same fields}\label{degree-and-support-are-combined-for-the-same-fields}

The exact projective split-map, dual-map and maximal-order quotient arguments of LIFT090 SG4/SG18--20 give

\[
 [K_\Sigma:K]\le M^8,\qquad
 [K_{\Sigma,\mathrm{out}}:K_\Sigma]\le L_N^{16}.
 \tag{RD15}
\]

They concern the actual fields and remain valid for every same-progenitor stage simultaneously. Applying RD9/13/14 to each step proves the first line of RD4 and RD5.

There is a direct bound on the complete extension that is valid even without using the source congruence. Every stage's M-torsion and N\_j\^{}2-torsion is contained in its ell-torsion, because both orders divide ell=lcm(M,L\_N\^{}2). The literal split inclusions of LIFT090 therefore give

\[
 K_{\Sigma,\mathrm{out}}\subset
       T_\ell=K(J[\ell],C_0[\ell]),\qquad
 [K_{\Sigma,\mathrm{out}}:K]\le\ell^8
                         \le M^8L_N^{16}.
 \tag{RD16}
\]

Only the degree is bounded by this containing field; no larger field replaces the actual output. RD14 already proves the ramification support of the actual output itself. Applying RD9 once, over K, proves

\[
 \operatorname{rd}(K_{\Sigma,\mathrm{out}})
 \le \operatorname{rd}(K)\operatorname{rad}(M L_N)
       [K_{\Sigma,\mathrm{out}}:K]
 \le \operatorname{rd}(K)\operatorname{rad}(M L_N)\ell^8.
 \tag{RD17}
\]

Finally rad(M L\_N)\textless=M L\_N and ell\textless=M L\_N\^{}2 give the last line of RD4. For the actual HF21 scales, ell=M L\_N\^{}2 and rad(M L\_N)=rad(M)rad(L\_N), by their congruence; the bound remains exactly the one claimed. If N\_j=1 for all j, this reduces to the seed-field bound. No disjointness of torsion fields or product of stage counts is assumed. If the complete actual ramification support is smaller than M L\_N O\_K, retain RD9 with that smaller ideal instead.

\subsection{6. Exact source and cost interface}\label{exact-source-and-cost-interface}

This finding changes a sufficient bound on the normalized discriminant of already constructed fields. The absolute field degree is still the actual one; LIFT090's separate bound 37800\emph{d}M\^{}8*L\_N\^{}16 is retained, with {[}B:Q{]} needed for absolute degrees. For the same fields,

\[
 \log\operatorname{rd}(K_{\Sigma,\mathrm{out}})
 \le\log\operatorname{rd}(K_{\mathrm{good}})
       +\log\operatorname{rad}(M L_N)+8\log\ell
 \le\log\operatorname{rd}(K_{\mathrm{good}})
       +9\log M+17\log L_N.
 \tag{RD18}
\]

\begin{longtable}[]{@{}
  >{\raggedright\arraybackslash}p{(\columnwidth - 4\tabcolsep) * \real{0.3333}}
  >{\raggedright\arraybackslash}p{(\columnwidth - 4\tabcolsep) * \real{0.3333}}
  >{\raggedright\arraybackslash}p{(\columnwidth - 4\tabcolsep) * \real{0.3333}}@{}}
\toprule\noalign{}
\begin{minipage}[b]{\linewidth}\raggedright
Item
\end{minipage} & \begin{minipage}[b]{\linewidth}\raggedright
Supplied here
\end{minipage} & \begin{minipage}[b]{\linewidth}\raggedright
Unchanged prerequisite or remaining cost
\end{minipage} \\
\midrule\noalign{}
\endhead
\bottomrule\noalign{}
\endlastfoot
Actual seed and output fields & RD2--3, without replacement & The same J, prescribed C0, K\_good and all source records \\
Finite ramification support & RD14, including full Galois composita & Everywhere-good model; only prime-to-order torsion is unramified \\
Normalized discriminant increase & RD9 and RD17--18 & Initial rd(K\_good), retained explicitly \\
Finite simultaneous stages & One ell and one radical support & Same progenitors; actual L\_N remains unevaluated \\
Models, action, polarization and theta line & No additional extension or choice & LIFT085--090's full all-object descent \\
Intrinsic Step5 degree/cost & No change to b or its log N transport & All IntegralOf, model, frame, polynomial, cusp and KeyCore hypotheses \\
Arithmetic main chain & No numerical error charge or new witness & Host/Sizes, Pins/what2, nonzero rule, selected content, period, Read/STiesR/032 and evaluated DP44 \\
\end{longtable}

The line, cubical identities, finite eta, Moduli.rep point and FrPt are already defined over the fields in RD2--3 by LIFT085 HF21--26, LIFT086 IM11--16 and LIFT088 SF12--16. No new generator or point is chosen to obtain the discriminant bound. Field conjugation transports the whole data and leaves absolute discriminants invariant; no independent choice per CM label is added.

The literal source clauses remain TwoWeightEqOneAddRoot.lean, Host14--27, Datum30--49, Level51--61, SAB72--84, Base124--129, TateIso131--142 and Obj144--160; PushInvertible.lean, Moduli29--38 and Coeff40--46; \nolinkurl{PartitionExtendByZeroAntitone.lean,} Par.eO22, PinLevel43 and PinField45; and IsCombAt.lean, FrPt67--76. The good integral model is not asserted to be an integral Obj over a base failing Base.unit. A smaller abstract coefficient field does not supply the fixed source's eO/PinField identifications.

The discriminant of K\_good and the given B, the actual effective N, coefficient and section heights, specified scalar normalizations and all source-construction gates are not bounded here. In particular a relative root-discriminant estimate is not a proof of a changed-family DP44 sign. No signed finite-kernel, global period or RH conclusion follows.

\subsection{7. Primary references and verification}\label{primary-references-and-verification}

Crew, R. (2019, April6). \emph{Local class field theory}. Proposition1.2.5.1, p.30; Proposition2.1.3.2, p.54; Corollary2.1.4.3 and Proposition2.1.4.5, pp.55--56; §2.1.5, equations(2.1.5.1)--(2.1.5.2), p.56; §2.1.6, equation(2.1.6.3), p.58. \url{https://people.clas.ufl.edu/rcrew/files/LCFT.pdf}

Milne, J. S. (2020). \emph{Complex multiplication} (version0.10). Corollary6.11 and the following paragraph, p.51. \url{https://www.jmilne.org/math/CourseNotes/CM.pdf}

The displayed primary passages were freshly read. The local upper bound and the finite-prime ideal-support summation are proved in full above. The same actual seed, projective torsion and all-object descent inputs are the checked LIFT085--090 results, with their exact scopes retained. The companion checks normalized valuation models, the distinct-residue argument, local-to-global degree factors, rational-support and lcm inequalities, same-field exponent arithmetic and source fingerprints. Such finite diagnostics supplement the universal proof and do not compute the actual ramification or an effective N.

The relative Picard/Poincare, origin-rigidified cubical line and finite-flat good-model data stay exactly those of the carried construction. Carried contextual anchors remain FGA232 Theorem3.1 pp148--149; FGA236 Corollaries1.7/2.3 pp229/232; SGA7 I VII2.7/2.9 pp169--174 and VIII7.4(b) pp310--311; SGA3 I4.7.3/4.7.3.1; SGA6 VI Theorem3.3; cote4 scan003 §§8.1.1--8.1.2 and scans005--006 §§8.4.1--8.4.3; cote54 scans005--007. No new lambda-ring, formal-group slope, l-adic weight or CM-motive comparison is asserted. Only the finite-etale torsion ramification and discriminant of the actual descended fields are added.

\subsection{Canonical algebraic-geometric context}\label{canonical-algebraic-geometric-context-34}

The actual everywhere-good abelian schemes and finite-etale prime-to-order torsion give finite Galois covers outside the stated torsion primes. The different and trace pairing quantify the same finite base changes. Picard/Poincare, normalized cubical line, full actions and finite-flat quotients are unchanged. Carried anchors: FGA232 Theorem3.1 pp148--149; FGA236 Corollaries1.7/2.3 pp229/232; SGA7 I VII2.7/2.9 pp169--174 and VIII7.4(b) pp310--311; SGA3 I4.7.3/4.7.3.1; SGA6 VI Theorem3.3; cote4 scan003 sections8.1.1--8.1.2 and scans005--006 sections8.4.1--8.4.3; cote54 scans005--007. No new formal-group slope, lambda-ring/K-theory, l-adic weight or CM-motive comparison. DP\_MULT Q3e/DP44/DP52/DP54--65 receives a relative field-discriminant reduction; THETA3 G1/A7, CX\_GRH CX-ALG-024 and GCT\_PLETH T-052 retain their arithmetic comparison gates.

\clearpage

\section{AM-R-092 --- root discriminant of the explicit general-Host seed field}\label{am-r-092-root-discriminant-of-the-explicit-general-host-seed-field}

\textbf{Canonical checked integration --- 2026-11-23T05:35:43.233179+00:00.} Complete saved-proof and executable reviews and separate exact diagnostic replays passed. No numerical error charge is introduced. Frozen author history is preserved. AM-R-092 ID1--18 bounds the initial root discriminant for the exact SF5 seed C0 and radical/torsion tower of LIFT088, over the actual source coefficient field B. The maximal-order class-field step B(j)/B is unramified; the actual radical tower has relative root-discriminant factor at most76832. Every CM conjugate satisfies \textbar sigma(j)\textbar\textless32 exp(pi sqrt(d)), and the actual C0{[}15{]} extension is treated at its true degree before applying the bound225. Consequently rd(Kgood)\textless2\^{}50 exp(2 pi sqrt(d)) rd(B). Composed with LIFT091, rd(KSigma,out)\textless2\^{}50 exp(2 pi sqrt(d)) rd(B) rad(M L\_N) lcm(M,L\_N\^{}2)\^{}8\textless=2\^{}50 exp(2 pi sqrt(d)) rd(B) M\^{}9 L\_N\^{}17. The stated log absolute discriminant bound retains the exact relative degree37800\emph{d}M\^{}8*L\_N\^{}16 and actual {[}B:Q{]},rd(B). This removes the separate initial-good-field discriminant premise for this fixed seed, not for arbitrary twists. Absolute B degree/discriminant, effective N\_j, heights, Host/Sizes existence, selected nonzero arithmetic rule/content, period, specific032/Read and evaluated DP44 remain open. No new finite-error charge;113/31 full arrays and complete floors26/11 unchanged.

Author: /root, GPT-6 Astra Pro; reasoning effort not externally exposed; ALP\_MACHINE. 23 November 2026 UTC.

\textbf{Author proof, awaiting the separately recorded whole saved-file and executable checks and independent review.} This finding uses the exact explicit elliptic model SF5 and field tower SF8--9 of LIFT088, the same-progenitor torsion bound of LIFT090, and the ideal-support different lemma LIFT091 RD9--13. The chosen model is part of the assertion: no bound uniform over arbitrary twists having the same j-invariant is asserted.

\subsection{1. Actual fields and statement}\label{actual-fields-and-statement}

Retain the literal Host, Datum, Moduli, Coeff and Pins hypotheses of LIFT085--090. In particular, d is a prime congruent to 7 modulo 16, k=Q(sqrt(-d)) has maximal-order discriminant -d, the prescribed ideal class is c\_K, and the given coefficient field B=c.F contains the specified k. Let

\[
 j=j(\mathbf C/\mathfrak c_K),\qquad K_j=B(j).
 \tag{ID1}
\]

Use the explicit LIFT088 SF5 elliptic factor

\[
 C_0:\quad y^2=x^3-\frac{3j}{j-1728}x-\frac{2j}{j-1728},
 \tag{ID2}
\]

with the prescribed complex CM action and canonical principal polarization, and set, in the fixed algebraic closure,

\[
\begin{split}
 K_J&=K_j(\zeta_7,i,\alpha,\beta),\\
 i^2&=-1,\qquad \alpha^7=-4,\qquad \beta^2=(\zeta_7-1)^7,\\
 K_{\rm good}&=K_J(C_0[15]).
\end{split}
\tag{ID3}
\]

These are precisely the displayed seed fields, without an additional splitting field, twist, frame generator, class representative or level change. The complete polarized general-Host seed and its good model are the ones constructed from this C0 and the original J=Jac(y\^{}2=x\^{}7-1) in LIFT087--090. Write rd(F)=\textbar disc(F/Q)\textbar\^{}(1/{[}F:Q{]}).

\textbf{Theorem.} For this explicit choice,

\[
 \operatorname{rd}(K_j)=\operatorname{rd}(B),\qquad
 \operatorname{rd}(K_J)\le76832\,\operatorname{rd}(B),
 \tag{ID4}
\]

and

\[
 \boxed{\quad
 \operatorname{rd}(K_{\rm good})
 <2^{50}\exp(2\pi\sqrt d)\operatorname{rd}(B).
 \quad}
 \tag{ID5}
\]

For a finite list of stages built from the same J, C0 and K\_good, let K\_Sigma, K\_\{Sigma,out\}, L\_N=lcm\_j N\_j, and ell=lcm(M,L\_N\^{}2) be the actual fields and integers of LIFT091 RD2--3. Under every retained nonzero-rule and source hypothesis for the output stages,

\[
\begin{split}
 \operatorname{rd}(K_\Sigma)
 &<2^{50}e^{2\pi\sqrt d}\operatorname{rd}(B)
      \operatorname{rad}(M)M^8,\\
 \operatorname{rd}(K_{\Sigma,\mathrm{out}})
 &<2^{50}e^{2\pi\sqrt d}\operatorname{rd}(B)
      \operatorname{rad}(M L_N)\ell^8\\
 &\le2^{50}e^{2\pi\sqrt d}\operatorname{rd}(B)M^9L_N^{17}.
\end{split}
\tag{ID6}
\]

The last weak inequality is a comparison of the explicit right sides; the resulting conservative bound for the actual root discriminant is also strict. In the literal LIFT085 construction N\_j=1 modulo M, so gcd(M,L\_N)=1 and ell=M L\_N\^{}2. We do not claim a level/output prime-overlap saving in that particular case. Shared-stage primes are still counted only once.

The remaining initial arithmetic inputs are the absolute degree and root discriminant of the given B. The integers N\_j, coefficient/section heights and all actual source and arithmetic nonvanishing conditions remain separate. This supplies no evaluated DP44 cost or RH assertion.

\subsection{2. The maximal-order class-field step is unramified}\label{the-maximal-order-class-field-step-is-unramified}

The singular modulus j is an algebraic integer by the Hilbert class polynomial theorem (Sutherland, 2023a, Theorem20.12, pp.5--6), already used in LIFT088 SF4. Because its CM order is the full maximal order, k(j)/k is the Hilbert class field and is unramified at every finite prime (Sutherland, 2023b, §21.4, p.7, and Corollary21.10). This is an additional class-field input beyond LIFT088, whose degree argument alone did not use the Hilbert class-field identification.

Since k is already a subfield of B, K\_j=B k(j). Unramified local extensions remain unramified after arbitrary finite base extension (Crew, 2019, Proposition1.2.4.6 and Corollary1.2.4.7(ii), p.29). Applying this at every finite prime of B shows that K\_j/B is unramified. Its relative different and discriminant are the unit ideals. The discriminant tower formula therefore yields

\[
 |\operatorname{disc}(K_j/\mathbf Q)|
   =|\operatorname{disc}(B/\mathbf Q)|^{[K_j:B]},
\]

which proves the first equality in ID4. This equality does not require B/Q to be Galois or B to be a CM field. It does not replace B by k(j).

For ID2, j cannot be 0 or 1728. Those elliptic curves have an origin-preserving automorphism of order 3 or 4, respectively, so their characteristic-zero CM endomorphism field is Q(sqrt(-3)) or Q(i). This is incompatible with the prescribed full order of discriminant -d, d\textgreater=7. The exact automorphism input is Milne (2021), Proposition IV.7.13 and Remark IV.7.14(a), p.140. Thus all denominators and nonzero principal ideals used below are legitimate.

\subsection{3. Ramification in the explicit Jacobian radical tower}\label{ramification-in-the-explicit-jacobian-radical-tower}

Write F0=K\_j, F1=F0(zeta7), F2=F1(i), F3=F2(alpha), and F4=F3(beta)=K\_J. Each successive extension is Galois. The first is a base change of Q(zeta7)/Q and has degree dividing 6. The second and fourth have degree 1 or 2. Since zeta7 is already in F2, the third is a cyclic Kummer extension of degree 1 or 7.

The supports are read directly from the defining polynomials. The cyclotomic polynomial Phi7 has discriminant -7\^{}5 up to the irrelevant sign convention for absolute discriminants, so its algebra is etale away from 7. The polynomial X\^{}2+1 has discriminant -4. The polynomial X\^{}7+4 has discriminant supported at 2 and 7, as follows either from its derivative 7X\^{}6 or from the binomial discriminant formula. Finally X\^{}2-(zeta7-1)\^{}7 has discriminant 4(zeta7-1)\^{}7, and zeta7-1 is a unit away from 7 because its norm from Q(zeta7) is 7. Consequently the relative extensions are unramified outside the indicated rational supports even if the base field is already ramified. Taking a field factor of a finite etale polynomial algebra preserves this statement.

Use LIFT091 RD9--13 with the actual degree of each step. In the first step every relative ramification index divides 6; since the only possible residue characteristic is 7, the step is tame. Its normalized local different contribution is at most the sum of log N(v)/{[}F0:Q{]} over v above 7, hence at most log 7. For the other steps RD9 gives the product of the actual degree and the radical rational support. It is only after applying that lemma that the degree is replaced by 2, 7 or 2.

\begin{longtable}[]{@{}
  >{\raggedright\arraybackslash}p{(\columnwidth - 6\tabcolsep) * \real{0.2000}}
  >{\raggedleft\arraybackslash}p{(\columnwidth - 6\tabcolsep) * \real{0.2667}}
  >{\raggedleft\arraybackslash}p{(\columnwidth - 6\tabcolsep) * \real{0.2667}}
  >{\raggedleft\arraybackslash}p{(\columnwidth - 6\tabcolsep) * \real{0.2667}}@{}}
\toprule\noalign{}
\begin{minipage}[b]{\linewidth}\raggedright
Step
\end{minipage} & \begin{minipage}[b]{\linewidth}\raggedleft
Relative degree
\end{minipage} & \begin{minipage}[b]{\linewidth}\raggedleft
Possible rational ramification
\end{minipage} & \begin{minipage}[b]{\linewidth}\raggedleft
Sufficient root-discriminant ratio
\end{minipage} \\
\midrule\noalign{}
\endhead
\bottomrule\noalign{}
\endlastfoot
F1/F0 & divides 6 & 7, tame & 7 \\
F2/F1 & 1 or 2 & 2 & 4 \\
F3/F2 & 1 or 7 & 2, 7 & 98 \\
F4/F3 & 1 or 2 & 2, 7 & 28 \\
\end{longtable}

Multiplication gives

\[
 \frac{\operatorname{rd}(K_J)}{\operatorname{rd}(K_j)}
 \le7\cdot4\cdot98\cdot28=76832.
 \tag{ID7}
\]

This proves the second assertion of ID4. No claim that the four extensions are disjoint is used. Their displayed degrees are bounds for each successive extension and their actual local ramification is included.

\subsection{4. Every conjugate of the actual singular modulus}\label{every-conjugate-of-the-actual-singular-modulus}

Each complex conjugate of j is the j-invariant of an elliptic curve with the same maximal CM order. Choose a reduced positive primitive binary quadratic form (a,b,c) of discriminant -d for its ideal class. Its associated

\[
 \tau=\frac{-b+i\sqrt d}{2a}
\]

is in the closed standard modular fundamental domain. In particular a\textgreater=1, \textbar Re(tau)\textbar\textless=1/2, \textbar tau\textbar\textgreater=1, and therefore

\[
 \frac{\sqrt3}{2}\le y=\operatorname{Im}\tau
       \le\frac{\sqrt d}{2},\qquad
 r=|q|=e^{-2\pi y}\le e^{-\pi\sqrt3}<\frac1{200}.
 \tag{ID8}
\]

The ideal-class description is the one preceding Sutherland's Theorem20.12; the closed fundamental-domain reduction is Milne (2021), Proposition V.1.3 and proof, p.177. One can see the upper bound without an estimate for a: integrality gives a\textgreater=1. The elementary final numerical inequality is certified exactly by pi\textgreater157/50, sqrt3\textgreater12/7 and sum\_\{n=0\}\^{}\{21\} ((157/50)(12/7))\^{}n/n! \textgreater{} 200.

Milne's displayed modular j formula and equation(42), p.183, together with Theorem V.3.7, p.192, give the normalized identity

\[
 j(\tau)=q^{-1}\frac{E_4(\tau)^3}
                       {\prod_{n\ge1}(1-q^n)^{24}},
 \qquad
 E_4(\tau)=1+240\sum_{n\ge1}\sigma_3(n)q^n.
 \tag{ID9}
\]

The normalization is exact: g4=60G4, G4=2 zeta(4) E4=(pi\^{}4/45)E4, and 1728(60 pi\^{}4/45)\^{}3=(2 pi)\^{}12. Thus no residual scalar in the modular discriminant is omitted.

For 0\textless r\textless=1/200, absolute convergence allows the elementary divisor-sum rearrangement

\[
 \sum_{n\ge1}\sigma_3(n)r^n
 =\sum_{m\ge1}\frac{m^3r^m}{1-r^m}
 \le\frac1{1-r}\sum_{m\ge1}m^3r^m
 =\frac{r(1+4r+r^2)}{(1-r)^5}.
\]

The rational function on the right is increasing on (0,1), for example by its nonnegative power series. Hence

\[
 |E_4(\tau)|
 \le1+\frac{240r(1+4r+r^2)}{(1-r)^5}
 \le\frac{703769200999}{312079600999}<\frac52.
 \tag{ID10}
\]

Likewise, using -log(1-t)\textless=t/(1-t) and absolute convergence,

\[
\begin{split}
 \log\prod_{n\ge1}|1-q^n|^{-24}
 &\le24\sum_{n\ge1}-\log(1-r^n)\\
 &\le\frac{24r}{(1-r)^2}
 \le\frac{4800}{39601}<\frac12.
\end{split}
\tag{ID11}
\]

Since e\textless4, the product in ID11 is less than 2. Combining ID8--11 gives, for every complex embedding sigma of any field containing j,

\[
 |\sigma(j)|
 <\frac{125}{4} e^{\pi\sqrt d}
 <32e^{\pi\sqrt d}.
 \tag{ID12}
\]

This bounds every conjugate, not only the prescribed complex embedding. Consequently it controls the normalized absolute norm from K\_J even when B has larger degree or is not Galois over Q.

\subsection{5. The actual elliptic model and its 15-torsion field}\label{the-actual-elliptic-model-and-its-15-torsion-field}

For the coefficients a0=-3j/(j-1728), b0=-2j/(j-1728) in ID2, direct algebra gives

\[
 \Delta(C_0)=-16(4a_0^3+27b_0^2)
       =2^{12}3^6\,\frac{j^2}{(j-1728)^3}.
 \tag{ID13}
\]

At a prime of K\_J outside the support of 6j(j-1728), both Weierstrass coefficients are integral and this discriminant is a unit. The displayed cubic and its point at infinity therefore form a smooth proper elliptic model at that prime. At a good prime whose residue characteristic does not divide 15, the entire 15-torsion is unramified. This is the finite-etale torsion implication in Milne (2020), Corollary6.11 and its following paragraph, p.51, used in LIFT091 RD14. It follows that the actual finite Galois extension

\[
 K_{\rm good}/K_J=K_J(C_0[15])/K_J
\]

is unramified outside the prime-ideal support of

\[
 S=(30j(j-1728))\mathcal O_{K_J}.
 \tag{ID14}
\]

The full O\_k action is already defined over K\_j by LIFT088 SF6. The maximal-order rank-one torsion argument of LIFT090 SG21 thus gives the actual degree m={[}K\_good:K\_J{]}\textless=15\^{}2=225. Applying LIFT091 RD9 to the actual degree m, and only then using this bound, yields

\[
 \frac{\operatorname{rd}(K_{\rm good})}{\operatorname{rd}(K_J)}
 \le m\,N_{K_J/\mathbf Q}(\operatorname{rad}S)^{1/[K_J:\mathbf Q]}
 \le225\,N_{K_J/\mathbf Q}(\operatorname{rad}S)^{1/[K_J:\mathbf Q]}.
 \tag{ID15}
\]

Local ramification indices divide the actual degree m. No divisibility by the numerical upper bound 225 is assumed.

The generator of S is a nonzero algebraic integer. Therefore N(rad S)\textless=N(S). Put C\_d=32 exp(pi sqrt d)\textgreater1. ID12 and the triangle inequality give

\[
 |\sigma(j(j-1728))|
 <C_d(C_d+1728)\le1729C_d^2
\]

for every embedding of K\_J into C. Multiplying all {[}K\_J:Q{]} bounds and taking that root proves

\[
 N(\operatorname{rad}S)^{1/[K_J:\mathbf Q]}
 <30\cdot1729\cdot32^2\exp(2\pi\sqrt d).
 \tag{ID16}
\]

This uses the actual prime-ideal support; it does not first replace S by every prime over a rational norm and then incur an additional field-degree factor.

Combining ID7, ID15 and ID16 proves

\[
\begin{split}
 \operatorname{rd}(K_{\rm good})
 &<C\,e^{2\pi\sqrt d}\operatorname{rd}(B),\\
 C&=76832\cdot225\cdot30\cdot1729\cdot32^2\\
  &=918207553536000<2^{50}.
\end{split}
\tag{ID17}
\]

This establishes ID5. The good model itself, full action, trace-one cubical line, finite level, Moduli.rep object and FrPt are already the complete LIFT088--090 constructions over these fields; the discriminant estimate makes no new choice of such objects.

\subsection{6. Composition and exact remaining arithmetic obligations}\label{composition-and-exact-remaining-arithmetic-obligations}

LIFT091 RD4 and RD17 apply to the same K\_good, same projective Serre summands, same duals and principalization, and same finite level/target fields. Substitution of ID5 proves ID6.

For reference, LIFT090 SG20--21 also gives {[}K\_\{Sigma,out\}:B{]}\textless=37800 d M\^{}8 L\_N\^{}16. Since a number-field root discriminant is at least one, ID6 yields the explicit absolute-discriminant majorant

\[
\begin{split}
 \log|\operatorname{disc}(K_{\Sigma,\mathrm{out}}/\mathbf Q)|
 <{}&37800\,d\,[B:\mathbf Q]M^8L_N^{16}\\
 &{}\cdot\bigl(\log\operatorname{rd}(B)+50\log2
       +2\pi\sqrt d+9\log M+17\log L_N\bigr).
\end{split}
\tag{ID18}
\]

This is a bound conditional on the displayed actual inputs, not their evaluation. In particular it does not bound the given B from an abstract smaller coefficient field, does not construct the literal fixed eO/PinField data over such a smaller field, and does not bound the nonvanishing isogeny scale N\_j.

\begin{longtable}[]{@{}
  >{\raggedright\arraybackslash}p{(\columnwidth - 4\tabcolsep) * \real{0.3333}}
  >{\raggedright\arraybackslash}p{(\columnwidth - 4\tabcolsep) * \real{0.3333}}
  >{\raggedright\arraybackslash}p{(\columnwidth - 4\tabcolsep) * \real{0.3333}}@{}}
\toprule\noalign{}
\begin{minipage}[b]{\linewidth}\raggedright
Gate component
\end{minipage} & \begin{minipage}[b]{\linewidth}\raggedright
Change proved here
\end{minipage} & \begin{minipage}[b]{\linewidth}\raggedright
Remaining input
\end{minipage} \\
\midrule\noalign{}
\endhead
\bottomrule\noalign{}
\endlastfoot
Initial class-field discriminant & ID4 removes any increase from adjoining j & The same maximal-order CM class and B containing k \\
Explicit J good-model radical field & Fixed factor 76832 & The actual SF8 radical choices \\
Complete initial good field & ID5, via the SF5 model and full action & rd(B) and the actual Host discriminant d \\
Complete simultaneous output fields & ID6 and ID18 & {[}B:Q{]}, actual N\_j and all nonzero-rule/source hypotheses \\
Main lift cost and selected pairing & No cost sign or value supplied & Host/Sizes, actual selected arithmetic content, period, Read/STiesR/032 and DP44 \\
\end{longtable}

This strictly reduces the field-discriminant subgate for the explicit general-Host seed: its additional root-discriminant cost is bounded by d and the known fixed radical tower, rather than retained as a separate unknown rd(K\_good). It leaves Q1 and Q2 open and does not change the finite support remainder, its signs, or any numerical error array.

The algebraic-geometric objects remain abelian schemes over Spec(O\_Kgood), their relative Picard/Poincare biextension, symmetric rigidified cubical line, finite-flat isogenies and finite-etale prime-to-order torsion. The different measures the finite base change; it does not alter a formal-group slope or provide a comparison of arithmetic theta operators. The carried exact contextual locators are FGA232 Theorem3.1 pp.148--149; FGA236 Corollaries1.7/2.3 pp.229/232; SGA7 I VII2.7/2.9 pp.169--174 and VIII7.4(b) pp.310--311; SGA3 I4.7.3/4.7.3.1; SGA6 VI Theorem3.3; cote4 scan003 §§8.1.1--8.1.2 and scans005--006 §§8.4.1--8.4.3; cote54 scans005--007. These are carried contexts, not newly read or newly identified motives. No new lambda-ring/K-theory operation, l-adic weight identity, CM-motive equivalence or L-function comparison is asserted. DP\_MULT Q3e/DP44/DP52/DP54--65 receives this field-cost reduction; THETA3 G1/A7, CX\_GRH CX-ALG-024 and GCT\_PLETH T-052 retain their respective arithmetic gates.

\subsection{7. Sources and check boundary}\label{sources-and-check-boundary}

Crew, R. (2019, April 6). \emph{Local class field theory}. Proposition1.2.4.6 and Corollary1.2.4.7, p.29; Proposition1.2.5.1, p.30; Proposition2.1.3.2, p.54; Corollary2.1.4.3 and Proposition2.1.4.5, pp.55--56. \url{https://people.clas.ufl.edu/rcrew/files/LCFT.pdf}

Milne, J. S. (2020). \emph{Complex multiplication} (version0.10). Corollary6.11 and following paragraph, p.51; the full-action inputs are used through LIFT088--091. \url{https://www.jmilne.org/math/CourseNotes/CM.pdf}

Milne, J. S. (2021). \emph{Elliptic curves} (2nd ed.). Proposition IV.7.13 and Remark IV.7.14(a), p.140; Proposition V.1.3 and proof, p.177; displayed j formula and equation(42), p.183; Theorem V.3.7 and proof, p.192. \url{https://www.jmilne.org/math/Books/EC2.pdf}

Sutherland, A. V. (2023a, November 21). \emph{18.783 Elliptic Curves, Lecture20}. Theorem20.12 and the preceding CM ideal-class description, pp.5--6. \url{https://math.mit.edu/classes/18.783/2023/LectureNotes20.pdf}

Sutherland, A. V. (2023b, November 28). \emph{18.783 Elliptic Curves, Lecture21}. §21.4 and Corollary21.10, p.7. \url{https://math.mit.edu/classes/18.783/2023/LectureNotes21.pdf}

The local proof dependencies are AM-LIFT-085 HF20--26; AM-LIFT-087 GS3--20; AM-LIFT-088 SF1--16; AM-LIFT-090 SG1--21; AM-LIFT-091 RD6--17. The literal clauses remain TwoWeightEqOneAddRoot.lean Host14--27, Datum30--49, Level51--61, SAB72--84, Base124--129, TateIso131--142, Obj144--160; PushInvertible.lean Moduli29--38 and Coeff40--46; \nolinkurl{PartitionExtendByZeroAntitone.lean} Par.eO22 and PinField45; IsCombAt.lean FrPt67--76. Exact file hashes and scoped reads are recorded separately. No Lean build or archive audit is needed or claimed.

The companion exact calculations check the normalized Weierstrass and modular identities, rational analytic majorants, radical tower arithmetic, degree/divisibility distinctions and source fingerprints. They do not prove the CM class-field theorem by sampling, compute the actual ramification of these fields, or evaluate B or any N\_j. The saved proof and entire executable must pass their recorded whole-file reviews before canonical promotion.

\subsection{Canonical algebraic-geometric context}\label{canonical-algebraic-geometric-context-35}

The exact CM elliptic curve, its maximal-order action and the fixed Jacobian radical model give the same actual abelian schemes after the established good-reduction base change. Unramified class-field base change, integral Weierstrass smoothness, finite-etale prime-to-order torsion and the different/trace pairing control the finite covers and root discriminants. Picard/Poincare data, normalized cubical line, quotient/action/level/FrPt descent and integral maps remain those already checked. Carried anchors: FGA232 Theorem3.1 pp148--149; FGA236 Corollaries1.7/2.3 pp229/232; SGA7 I VII2.7/2.9 pp169--174 and VIII7.4(b) pp310--311; SGA3 I4.7.3/4.7.3.1; SGA6 VI Theorem3.3; cote4 scan003 sections8.1.1--8.1.2 and scans005--006 sections8.4.1--8.4.3; cote54 scans005--007. No new lambda-ring/K-theory identity, crystalline slope, l-adic weight theorem, motivic comparison or Hecke RH statement. DP\_MULT Q3e/DP44/DP52/DP54--65 receives this effective initial-field cost reduction; THETA3 G1/A7, CX\_GRH CX-ALG-024 and GCT\_PLETH T-052 retain their actual arithmetic comparison gates.

\clearpage

\section{AM-LIFT-092 --- the conductor of the actual good-CM torsion fields}\label{am-lift-092-the-conductor-of-the-actual-good-cm-torsion-fields}

\textbf{Canonical checked integration --- 2026-11-23T05:40:30.780735+00:00.} Complete saved-proof and executable reviews and separate exact diagnostic replays passed. No numerical error charge is introduced. Frozen author history is preserved. AM-LIFT-092 CD1--19 strengthens the root-discriminant cost of the same actual LIFT088--090 fields. For the two full-maximal-order CM progenitors J,C0 over the actual everywhere-good Kgood, every actual F inside Kgood(J{[}n{]},C0{[}n{]}) is abelian and has conductor dividing n O\_Kgood and discriminant dividing (n O\_Kgood)\^{}(D-1), D={[}F:Kgood{]}. Thus rd(F)/rd(Kgood)\textless=n\^{}(1-1/D)\textless=n, including all dyadic/CM-ramified primes and D=1. The same simultaneous output satisfies rd(KSigma,out)\textless=rd(Kgood)\emph{lcm(M,L\_N\^{}2)=rd(Kgood)}M\emph{L\_N\^{}2 under the actual N\_j=1 mod M; its relative output increase over KSigma is at most L\_N\^{}2. Combining the separate canonical R092 bound gives rd(KSigma,out)\textless2\^{}50 exp(2 pi sqrt(d))}rd(B)\emph{M}L\_N\^{}2. LIFT090 relative degree37800\emph{d}M\^{}8*L\_N\^{}16 is unchanged; the sufficient absolute-discriminant logarithm is this degree times {[}B:Q{]} times (log rd(B)+50 log2+2 pi sqrt(d)+log M+2 log L\_N). All original fields, varieties, polarizations, line, level, quotient, action, source points and frames are retained. Actual B degree/discriminant, effective N\_j, metric/height costs, Host/Sizes, selected nonzero arithmetic content, period, specific032/Read and evaluated DP44 remain open. No finite-error charge;113/31 full arrays and complete floors26/11 unchanged.

Author: /root/lift\_accounting, ALP\_MACHINE. 23 November 2026 UTC.

\textbf{Author-frozen proof.} Complete saved-file, source and executable verification is recorded in the adjacent author packet; independent review and canonical integration are separate. This strengthens the relative discriminant estimate of LIFT091 for its same actual fields. It uses the full maximal-order CM actions on the fixed progenitors J and C0 over the actual everywhere-good field. It does not change that field, any polarized seed, the finite level, the target point, or an isogeny scale.

\subsection{1. Actual objects and theorem}\label{actual-objects-and-theorem}

Retain the complete source and construction hypotheses of canonical LIFT085--090. Write K=K\_good, E\_J=Q(zeta7), and E\_C=k for the prescribed imaginary quadratic Host field acting on C0. The abelian varieties J and C0 are defined over K, have dimensions three and one, have their full O\_\{E\_J\} and O\_\{E\_C\} actions over K, and have everywhere-good abelian-scheme models over O\_K. Their specified principal polarizations are already over K. For n\textgreater=1 put

\[
 T_n=K(J[n],C_0[n]).
 \tag{CD1}
\]

These are exactly the two-progenitor torsion fields of LIFT090. Let K subset F subset T\_n be any intermediate field and put D={[}F:K{]}. Then F/K is abelian and Galois, and its finite conductor and relative discriminant satisfy

\[
 \mathfrak f(F/K)\mid nO_K,\qquad
 \mathfrak d_{F/K}\mid(nO_K)^{D-1}.
 \tag{CD2}
\]

In particular,

\[
 \boxed{\quad
 \frac{\operatorname{rd}(F)}{\operatorname{rd}(K)}
       \le n^{1-1/D}\le n .\quad}
 \tag{CD3}
\]

The conductor has no infinite part. All residue characteristics, including two and primes ramifying in either CM field or K, are included. No absolute-unramified assumption is made. For D=1 the first bound is one, as required; n=1 gives F=K.

For a finite same-progenitor list of stages, retain the literal fields

\[
\begin{split}
 K_\Sigma&=K(A_{0,j}[M]:j\in I),\\
 K_{\Sigma,\mathrm{out}}
   &=K_\Sigma(A_{0,j}[N_j^2]:j\in I),\\
 L_N&=\operatorname{lcm}_{j\in I}N_j,\qquad
 \ell=\operatorname{lcm}(M,L_N^2).
\end{split}
\tag{CD4}
\]

Use lcm(empty)=1. Denote their actual degrees over K by D\_seed and D\_out, and the relative output degree over K\_Sigma by D\_rel. The same fields obey

\[
\begin{split}
 \operatorname{rd}(K_\Sigma)
   &\le\operatorname{rd}(K)M^{1-1/D_{\mathrm{seed}}}
    \le\operatorname{rd}(K)M,\\
 \operatorname{rd}(K_{\Sigma,\mathrm{out}})
   &\le\operatorname{rd}(K)\ell^{1-1/D_{\mathrm{out}}}
    \le\operatorname{rd}(K)\ell
    \le\operatorname{rd}(K)M L_N^2,\\
 \frac{\operatorname{rd}(K_{\Sigma,\mathrm{out}})}
      {\operatorname{rd}(K_\Sigma)}
   &\le L_N^{2(1-1/D_{\mathrm{rel}})}\le L_N^2.
\end{split}
\tag{CD5}
\]

For the actual scales of LIFT085 HF20--23, every N\_j is congruent to one modulo M. Hence gcd(M,L\_N)=1 and ell=M L\_N\^{}2. The general lcm notation does not assert an extra level/target overlap saving in this source case. Shared-stage target factors are counted once by L\_N. A single stage is the special case L\_N=N.

\subsection{2. The exact CM reciprocity formula}\label{the-exact-cm-reciprocity-formula}

Let A be either J or C0 and let E be its full CM field. Temporarily write Phi for its actual CM type and E\^{}* for its reflex field, using a fixed embedding of K into an algebraic closure of Q. The full E-action is defined over K, so the reflex field embeds in K. The adelic E-linear Tate module is rank one over E tensor\_Q A\_f.~The action therefore factors through the abelianized Galois group.

Use the arithmetic-Frobenius normalization of the Artin map, as in Conrad's \emph{Main Theorem of Complex Multiplication}, notation on p.2. Theorem3.2 and Remark3.3, pp.8--9, give the exact formula

\[
 \rho_A(\operatorname{Art}_K(s))
       =\lambda_A(s)\,r_A(s_f)^{-1},\qquad
 r_A=N_\Phi\circ N_{K/E^*},
 \tag{CD6}
\]

where lambda\_A is the discrete E-valued CM character and r\_A is the actual reflex-type norm. This describes the Galois action on the whole rational Tate module, not just an isogeny class or a semisimplification. The formula is transported to the integral lattice in section4 below.

At every finite place v of K, the actual good reduction of A implies

\[
                 \lambda_A(O_{K_v}^{\times})=1.
 \tag{CD7}
\]

This is exactly Conrad's Theorem3.5(1), p.10. Its proof chooses a Tate prime different from the residue characteristic and uses the good-reduction inertia criterion to show that the E-valued lambda character is trivial on the entire local unit group. Thus CD7 also holds at places dividing n.~It is stronger than choosing a sufficiently small subgroup on which lambda happens to be trivial. We use neither the subsequent L-function assertion nor its purity discussion.

For a local unit u, viewed as an idele supported at v, CD6--7 give

\[
       \rho_A(\operatorname{Art}_K(u))=r_A(u)^{-1}.
 \tag{CD8}
\]

The inverse is retained. Changing the Artin convention would invert this formula, which leaves the congruence kernel unchanged; no such convention change is made here. This new reciprocity calculation does not alter a theta, Bessel or Haar normalization elsewhere in the source.

\subsection{3. Integral principal congruences under the type norm}\label{integral-principal-congruences-under-the-type-norm}

Fix a rational prime p and m\textgreater=1. The morphism of Q-tori r\_A is a map

\[
  (K\otimes_\mathbf Q\mathbf Q_p)^{\times}
       \longrightarrow(E\otimes_\mathbf Q\mathbf Q_p)^{\times}.
 \tag{CD9}
\]

For every embedding tau:E-\textgreater overline Q, Conrad's Remark3.6, pp.10--11, writes tau composed with r\_A as a product of embeddings of K, indexed by a subset Psi\_tau of Hom(K,overline Q). The displayed formula there is first expressed on global elements and identifies the algebraic characters. Even starting only with its principal elements congruent to one modulo the displayed auxiliary ideal, these elements form a full additive lattice translate and are Zariski dense in Res\_\{K/Q\}(G\_m); hence they determine that character identity. It therefore remains the same identity of tori after base change to Q\_p and an algebraic closure. Equivalently, this follows from the determinant definition of the reflex norm in the paragraph preceding Remark3.3. There are integral, nonnegative exponents; no roots or division by a degree are taken.

Let v lie above p and let

\[
                    u\in1+p^mO_{K_v}.
 \tag{CD10}
\]

Put u in the v component of K tensor Q\_p and one in every other component. Each p-adic embedding of K belonging to v extends continuously to K\_v and sends u to an element of 1+p\^{}m O\_\{C\_p\}; all other component factors are one. Thus the embedded type-norm identity gives

\[
 \tau(r_A(u))\in1+p^mO_{\mathbf C_p}
       \quad\hbox{for every }\tau:E\hookrightarrow\overline{\mathbf Q}_p.
 \tag{CD11}
\]

Here the normalized p-adic valuation on C\_p has v\_p(p)=1. Every finite extension of Q\_p has its unique extended absolute valuation. If y belongs to E tensor Q\_p and every embedding has v\_p(tau(y))\textgreater=m, then y/p\^{}m belongs to the product of the rings of integers of the completions of E above p.~Consequently

\[
 r_A(u)-1\in p^m(O_E\otimes\mathbf Z_p),\qquad
 r_A(u)^{-1}-1\in p^m(O_E\otimes\mathbf Z_p).
 \tag{CD12}
\]

The inverse assertion follows because 1+p\^{}m times an integral element is a unit in each local factor. This argument uses valuation rings, not an etale integral splitting of O\_E at p.~It applies unchanged if p=2 or p ramifies in E, E\^{}*, K or their common closure. It does not require choosing an integral eigenbasis for a ramified tensor product.

In the integer-normalized valuation at v, CD10 is

\[
  u\in U_v^{e_v^{\mathrm{abs}}m},\qquad
  e_v^{\mathrm{abs}}=e(K_v/\mathbf Q_p),\quad
  U_v^a=1+\mathfrak p_v^a\ (a\ge1).
 \tag{CD13}
\]

This exponent is exactly ord\_v(p\^{}m). No absolute ramification factor is dropped. The use of the rational integer p\^{}m is what makes the contraction to every maximal CM order in CD12 exact.

\subsection{4. From the rational action to the actual torsion lattice}\label{from-the-rational-action-to-the-actual-torsion-lattice}

The integral Tate module T\_p(A) is stable under the full O\_E tensor Z\_p action. For J and C0 it is locally a rank-one projective module over that maximal order, as checked in LIFT090 section4, SG16--17. No chosen rational eigenbasis replaces this integral lattice. If b belongs to p\^{}m(O\_E tensor Z\_p), then

\[
                    bT_p(A)\subset p^mT_p(A).
 \tag{CD14}
\]

By CD8 and CD12, Art\_K(u) is therefore the identity on T\_p(A)/p\^{}m T\_p(A)=A{[}p\^{}m{]}. The equality is the actual compatible torsion quotient. It is unaffected by a nonprincipal global projective ideal because the full order acts integrally before any local trivialization.

The maximal-order hypothesis is substantive. For illustration only, in E=Q(sqrt5) let O\_E=Z{[}(1+sqrt5)/2{]} and R=Z{[}sqrt5{]}. The unit a=2+sqrt5 equals 1+2*(1+sqrt5)/2. It is one modulo 2O\_E but not one modulo 2R. Multiplication by a on the basis (1,sqrt5) of R has matrix {[}{[}2,5{]},{[}1,2{]}{]}, which is not the identity modulo two. Thus a congruence in the maximal order could not silently be transferred to an arbitrary nonmaximal-order lattice. The actual progenitors have the full order, so CD14 is available.

For a local unit at v of residue characteristic p, the type norm in CD6 has no nontrivial components at rational primes different from p.~CD7 then makes its action exactly the identity on every prime-to-p torsion subgroup. Combining this with CD14, if p\^{}m is the p-part of n, the kernel of the n-torsion action contains U\_v\^{}\{ord\_v(n)\} when p divides n, and it contains O\_\{K\_v\}\^{}\{times\} when p does not divide n.~The same conclusion holds simultaneously for J and C0.

\subsection{5. Conductor and discriminant of the actual subfields}\label{conductor-and-discriminant-of-the-actual-subfields}

The full CM actions imply that the torsion image for each progenitor is commutative: it consists of scalar automorphisms of the rank-one CM Tate module. Thus T\_n/K is finite abelian, and every intermediate F/K is Galois and abelian. This is not an assertion that every subfield of an arbitrary Galois extension is Galois.

The two simultaneous local kernels of section4 say that the Artin map to Gal(T\_n/K) kills the idele subgroup

\[
 W_{nO_K}=K_\infty^{\times}
   \prod_{v\nmid n}O_{K_v}^{\times}
   \prod_{v\mid n}(1+nO_{K_v}).
 \tag{CD15}
\]

The infinite component is killed because K contains a CM reflex field and has no real embeddings; its infinite multiplicative group is a product of connected copies of C\^{}\{times\}. More concretely, this actual K already contains zeta7. The finite quotient action also kills the diagonal K\^{}\{times\} by reciprocity. Milne's \emph{Class Field Theory}, V, Proposition4.6 and Proposition4.7, pp.172--174, identify the resulting quotient with a quotient of the ray class group of modulus nO\_K. Therefore the conductor divides nO\_K. Each actual intermediate F gives a further quotient of that same map, so its conductor also divides nO\_K. No ray class field replaces F.

Let G=Gal(F/K), with \textbar G\textbar=D. Every finite character of G has conductor dividing nO\_K; the trivial character has conductor O\_K. The conductor--discriminant formula (Milne, V, Theorem3.27, p.165) gives

\[
 \mathfrak d_{F/K}
       =\prod_{\chi\in\widehat G}\mathfrak f_0(\chi)
       \mid(nO_K)^{D-1}.
 \tag{CD16}
\]

Taking norms and using the exact discriminant tower yields

\[
 \begin{split}
 \log\operatorname{rd}(F)-\log\operatorname{rd}(K)
 &=\frac{\log N_{K/\mathbf Q}\mathfrak d_{F/K}}
          {[K:\mathbf Q]D}\\
 &\le\frac{D-1}{D}\log n .
 \end{split}
 \tag{CD17}
\]

This proves CD2--3. In the trivial extension there are no nontrivial characters, the empty discriminant product is O\_K, and CD17 is equality zero. Absolute ramification is already incorporated through ord\_v(n) in CD13; no additional factor depending on {[}K:Q{]} or the number of primes over p appears.

\subsection{6. The same seeds, targets and complete source data}\label{the-same-seeds-targets-and-complete-source-data}

LIFT090's actual integral idempotent split maps and their duals put every stage's n-torsion field inside T\_n.~These are maps over K, not lifts through an isogeny of unspecified index. Zarhin's chosen target is the literal product of X and X\^{}dual used in that construction; its polarization changes do not change this field inclusion. Hence

\[
  K_\Sigma\subset T_M,\qquad
  K_{\Sigma,\mathrm{out}}\subset T_\ell.
 \tag{CD18}
\]

These are the unchanged actual fields of CD4. Applying CD3 to them gives the first two inequalities of CD5. Over K\_Sigma the same J and C0 still have their full actions and everywhere-good models. Their base-changed torsion fields contain the actual target extension for the common order L\_N\^{}2. Repeating the same theorem over this base proves the final relative inequality in CD5. No disjointness of fields or multiplication of stage degrees is assumed.

The level and quotient field definitions remain those of LIFT085 HF20--26, LIFT088 SF12--16 and LIFT090 SG18--21. The quotient, integral O\_E action, principal polarization, canonical trace-one line, rigidification/cubical identities, finite eta, represented Moduli point and FrPt have the same all-object descent as before. Frames need no new extension once their eigenspaces and theta line are defined. The good integral model remains distinct from an integral Obj whose Base.unit field-level predicate might fail.

\begin{longtable}[]{@{}
  >{\raggedright\arraybackslash}p{(\columnwidth - 4\tabcolsep) * \real{0.3333}}
  >{\raggedright\arraybackslash}p{(\columnwidth - 4\tabcolsep) * \real{0.3333}}
  >{\raggedright\arraybackslash}p{(\columnwidth - 4\tabcolsep) * \real{0.3333}}@{}}
\toprule\noalign{}
\begin{minipage}[b]{\linewidth}\raggedright
Quantity
\end{minipage} & \begin{minipage}[b]{\linewidth}\raggedright
New bound for the same object
\end{minipage} & \begin{minipage}[b]{\linewidth}\raggedright
Retained prerequisite or open cost
\end{minipage} \\
\midrule\noalign{}
\endhead
\bottomrule\noalign{}
\endlastfoot
Initial good field & No change & rd(K\_good), supplied separately if proved \\
Relative seed discriminant & rd(K\_seed)/rd(K\_good)\textless=M & Actual good J,C0 and full maximal-order actions \\
Relative complete output discriminant & \textless=M L\_N\^{}2 & Same finite target list and actual scales \\
Target over the level field & \textless=L\_N\^{}2 & Same descended fields and models \\
Absolute or relative field degrees & LIFT090's previous M\^{}8 and L\_N\^{}16 bounds unchanged & Existing coefficient field and {[}B:Q{]} \\
Intrinsic line and Step5 metric cost & Unchanged LIFT086--087 formulas & Actual log N, models, sections, frames and polynomial hypotheses \\
Main arithmetic conclusions & No new witness or finite charge & Nonzero rule, effective N, selected image/content, period, Read/STiesR/032 and DP44 \\
\end{longtable}

Thus the relative logarithmic discriminant cost is at most

\[
 \log\operatorname{rd}(K_{\Sigma,\mathrm{out}})
 \le\log\operatorname{rd}(K_{\mathrm{good}})
          +\log M+2\log L_N.
 \tag{CD19}
\]

This improves LIFT091's sufficient 9logM+17logL\_N estimate for its same actual CM fields. LIFT091's general ideal-support lemma is unchanged and still applies to extensions lacking these CM-good-reduction hypotheses. A choice of a smaller abstract Coeff field does not supply the fixed eO/PinField identifications. No effective scale, new point, nonzero evaluation, selected arithmetic value or DP44 sign is proved by the conductor bound.

\subsection{7. Literal and primary verification scope}\label{literal-and-primary-verification-scope}

The literal Host/Datum/Level/SAB/Base/TateIso/Obj clauses are those in TwoWeightEqOneAddRoot.lean; Moduli and Coeff are in PushInvertible.lean; the fixed eO/PinLevel/PinField clauses are in \nolinkurl{PartitionExtendByZeroAntitone.lean;} FrPt is in IsCombAt.lean. They are retained in exactly the source scopes already read for LIFT091, with unchanged fingerprints. The actual full-order progenitor and same-field descent proof is canonical LIFT090. Its integral projective modules, dual maps and finite simultaneous list are essential; a rational isotypy alone is insufficient.

Fresh primary inputs are Conrad's \emph{Main Theorem of Complex Multiplication}, notation p.2, Theorem3.2 and its proof pp.8--9, the reflex-norm paragraph and Remark3.3 p.9, Theorem3.5(1) and proof p.10, and the product-of-embeddings identity in Remark3.6 pp.10--11. The new conductor consequence is proved here. Milne's \emph{Class Field Theory}, version4.03, V, Propositions4.6--4.7 pp.172--174 and Theorem3.27 p.165, supply the exact idelic/conductor and discriminant normalizations. No claim that their entire books or papers were newly read is made.

The companion diagnostics test products and inverses of integral principal congruences, ramified ring models, integral versus nonmaximal-order lattice behavior, exact character/discriminant exponents, absolute-ramification normalization, same-field finite lists and source fingerprints. They supplement the universal proof rather than compute actual conductors or N.

Conrad, B. (n.d.). \emph{Main theorem of complex multiplication}. \url{https://math.stanford.edu/~conrad/papers/mainthm.pdf}

Milne, J. S. (2020). \emph{Class field theory} (version4.03). \url{https://www.jmilne.org/math/CourseNotes/CFT.pdf}

The relative Picard/Poincare, origin-rigidified cubical line and finite-flat good-model language remains the carried LIFT085--090 construction, with its established FGA232/FGA236, SGA7 I VII/VIII, SGA3 I, SGA6 VI and cote4/cote54 anchors. This finding adds finite CM class-field ramification information only. It makes no new comparison of lambda-rings, formal-group slopes, weights, motives or an RH implication.

\subsection{Canonical algebraic-geometric context}\label{canonical-algebraic-geometric-context-36}

The actual full CM Tate lattices carry integral maximal-order actions at every prime. CM reciprocity and good reduction make their same torsion covers factor through the finite ray-class modulus of the torsion order; the character-conductor product measures the different of each actual abelian subcover. The Picard/Poincare and origin-rigidified cubical line, finite-flat quotients, full actions, finite level and framed points are unchanged. Carried anchors: FGA232 Theorem3.1 pp148--149; FGA236 Corollaries1.7/2.3 pp229/232; SGA7 I VII2.7/2.9 pp169--174 and VIII7.4(b) pp310--311; SGA3 I4.7.3/4.7.3.1; SGA6 VI Theorem3.3; cote4 scan003 sections8.1.1--8.1.2 and scans005--006 sections8.4.1--8.4.3; cote54 scans005--007. No new crystalline slope, lambda-ring/K-theory identity, l-adic weight theorem, CM-motive comparison or Hecke RH result. DP\_MULT Q3e/DP44/DP52/DP54--65 receives a sharper finite-field ramification cost; THETA3 G1/A7, CX\_GRH CX-ALG-024 and GCT\_PLETH T-052 retain their separate arithmetic comparisons.

\clearpage

\section{AM-K-107 --- a uniform selected-trace refinement and the next complete coordinate cell}\label{am-k-107-a-uniform-selected-trace-refinement-and-the-next-complete-coordinate-cell}

\textbf{Canonical checked integration --- 2026-11-23T05:59:10.807934+00:00.} Complete root and independent proof/script/interface reviews and separate exact replays passed. The actual canonical predecessor has113/31 entries; this integration gives114/31, complete floors26/11. Frozen author history is preserved. AM-K-107 SV1--36 removes the entire k138 cell of the actual all-II globally rank-two-unipotent scalar-cross support26 remainder. The genuine selected density yields a full tail-inclusive squared coefficient below2\^{}-933 and a single actual near-point product below2. Together with the exact radial factor, all-coordinate norm and original saturated transverse lattice, the complete count is bounded by2\^{}(-50428781/204800)\textless2\^{}-246. The original scalar theta is in real dimension336, with full21\^{}(k+3) and all original central, support, row, finite and pattern factors. The remaining branch has139\textless=k\textless=164 and2\^{}(71443/64)\textless x\textless2\^{}2108, v=q26,b=j0. Exactly one246 append gives114/31 ordered errors with all old overlaps retained, both full sums below2\^{}-230, unique minimum231 and next minimum246; complete floors26/11 remain. Mandatory corrected boundary conventions: SV6 is strict for h\textgreater=1 and equality at h0; SV11 quotient is for r\textgreater0 and the bound at r0 is1. The frozen proof is read with the explicit boundary addendum reproduced below. Signed positivity, selected arithmetic content, period, effective source, specific032/Read, actual DP44 and RH remain open.

Author: /root/kummer\_extension, with independently supplied saturated-basis and exact-endpoint checks by /root/definite\_oscillator. 23 November 2026 UTC. Controlled pending proof. Final author read, exact replay, independent review and canonical integration are separate records.

\subsection{1. Actual scope and the complete new count}\label{actual-scope-and-the-complete-new-count}

Keep the actual J56 all-top family, H=133, f=a=168, every admitted word and canonical twist, the selected compact module with its genuine vector-valued Hilbert metric, both compact factors, the original fields and projective modules, paired source/target pencils, finite projectors, central phases, Haar measures and class masses.

After K106 the possible support26 branch has global rank-two unipotence, scalar cross, all-II type, and \[
 v=q_{\rm elig}=26,\qquad b=j=0,\qquad
 2^{1116}<x=N_ED_\eta<2^{2108},\qquad138\le k\le164.
 \tag{SV1}
\] Here k is the number of nonzero original S covector coordinates. K103 SC26.5--7 proves the support-free full positive compression and normal form; SC26.11 proves uniqueness of eta for each support in this window.

\textbf{Theorem.} The genuine selected trace obeys the uniform local bound SV11--12 and the actual rank-two near-point bound SV16. With every original arithmetic and finite factor retained, the complete subset of SV1 with k=138 has absolute normalized contribution less than \(2^{-246}\). Append246 once to the current remainder budget, preserving every earlier charge and both ordered arrays. The only possible coordinate range in this all-II branch becomes139 through164. No assertion for support27 is made.

The proof retains the entire small-probability term before deriving a uniform trace estimate; no such term is dropped from a later lattice count. It then uses a saturated integral basis in the original transverse lattice, a dimension336 scalar theta estimate, and one coupled endpoint for the complete cell. The finite arithmetic proves a sufficient upper bound, not a nonzero matrix or a lower bound for the remaining cells.

\subsection{2. Exact positive moments control the lower radial tail}\label{exact-positive-moments-control-the-lower-radial-tail}

At one CM label let Y be the radial variable under the genuine selected probability density of DO080 RT6--8, and write \[
 Z_r=Y^{-1},\quad M_1={\bf E}Z_r,\quad m=M_1-1,\quad
 v=\operatorname{Var}(Z_r),\quad B=2^{955}.
 \tag{SV2}
\] The notation \(Z_r\) is distinct from the rational positive space in the arithmetic normal form. The exact K046 BL8--10 moments apply to each vector in the selected irreducible module and hence to its averaged density. For every nonnegative integer h, \[
 M_1=\prod_i\frac{r_i+m_i-1}{r_i-2},\qquad
 {\bf E}Y^h=\prod_i\frac{(r_i-1)_h}{(r_i+m_i)_h}.
 \tag{SV3}
\] There is no upper restriction on h for the positive moment. The actual all-word estimates DO080 RT21--23 give \[
 32<M_1<34,\qquad
 S_\varepsilon:=\sum_i
 \frac{m_i+1}{(r_i-3)(r_i+m_i-1)}<B^{-1},
 \qquad
 \frac vm< V_0:=\frac{34^2}{31(B-1)}.
 \tag{SV4}
\] All \(r_i>3\) and \(m_i\ge0\). In particular \(31<m<33\). The interior word changes \(m_i\), not \(r_i\), as corrected in DO079-C1.

For each \(j\ge0\), direct subtraction gives \[
 \frac{(r_i+m_i-1)(r_i-1+j)}
      {(r_i-2)(r_i+m_i+j)}
 =1+\frac{(j+1)(m_i+1)}
          {(r_i-2)(r_i+m_i+j)}
 \le1+(j+1)\varepsilon_i .
 \tag{SV5}
\] Multiplying SV5 and using \(\log(1+t)\le t\) proves the full moment estimate \[
 M_1^h{\bf E}Y^h
 \le\exp\!\left(\frac{h(h+1)}2S_\varepsilon\right)
 <\exp\!\left(\frac{h(h+1)}{2B}\right).
 \tag{SV6}
\] This is a uniform analytic consequence of the literal moments, not an assumption that Y is Gaussian.

Take \(\epsilon=1/256\) and \(h=2^{19}\). On the lower-tail event \(\mathcal E=\{Z_r-1<(1-\epsilon)m\}\), one has \[
 Z_r<M_1\left(1-\frac{31}{8192}\right).
\] Markov's inequality applied to the positive variable \(Y^h\), followed by SV6, yields \[
 \Pr(\mathcal E)
 <\left(1-\frac{31}{8192}\right)^h
       \exp\!\left(\frac{h(h+1)}{2B}\right)
 <e^{-1983}<2^{-2800}.
 \tag{SV7}
\] The last constants are exact: \(31h/8192=1984\), \(h(h+1)<2B\), \(e>27/10\), \((27/10)^7>2^{10}\), and \(1983(10/7)>2800\).

\subsection{3. Retain the entire denominator in the spread estimate}\label{retain-the-entire-denominator-in-the-spread-estimate}

Set \(S=2\sqrt{Z_r-1}\) and \(s_0=2\sqrt m\), exactly as in DO080 RT16--19. The identity is \[
 {\bf E}(S-s_0)^2
 =4{\bf E}\frac{(Z_r-M_1)^2}
   {(\sqrt{Z_r-1}+\sqrt m)^2}.
 \tag{SV8}
\] On \(\mathcal E^c\) the denominator is at least \(C m\), where \[
 C=(1+\sqrt{255}/16)^2>\frac{399}{100}.
\] On \(\mathcal E\), \(0<Z_r-1<m\), so the denominator is at least m and the squared numerator is at most \(m^2\). Thus the complete second moment, including the lower tail, satisfies \[
 {\bf E}(S-s_0)^2
 \le \frac{4v}{Cm}+4m\Pr(\mathcal E)
 <\frac{4V_0}{C}+132\,2^{-2800}
 =:\sigma_*^2.
 \tag{SV9}
\] The elementary comparison \(C>399/100\) follows by squaring the positive rational number \(12761/800<\sqrt{255}\).

The same positive-weight Cauchy--Schwarz argument as DO080 RT19 now gives \[
 \left(\int_0^\infty\sqrt{q_S(t)}\,dt\right)^2
 \le 2\pi\sigma_* .
 \tag{SV10}
\] This uses \({\bf E}(S-s_0)^2\le\sigma_*^2\); it does not require \(s_0={\bf E}S\). The exact radial-density multiplier of DO092 SN16--18 and K106 NT2--4 is \[
 \vartheta=\frac{4^{167}}{334\binom{334}{167}}.
\] The genuine selected trace therefore obeys \[
 |c_W(g)|\le
 \min\left\{1,\frac{\vartheta\,2(a-1)\sigma_*}{\sinh r(g)}\right\}.
 \tag{SV11}
\] Both compact factors and the selected Hilbert metric remain present. This is a trace estimate, not an operator estimate.

For a=168 the exact squared coefficient, including its entire tail term, is \[
 [2(a-1)\sigma_*]^2
 <
 \frac{515834944}{31(B-1)}\frac{100}{399}
       +528\cdot167^2\,2^{-2800}
 <2^{-933}.
 \tag{SV12}
\] Every comparison in SV12 is a rational integer comparison. Thus the tail is absorbed uniformly before any row, central-fibre, support or finite-pattern enumeration. There is no uncharged additive tail sector.

\subsection{4. The actual near point and all nonzero original coordinates}\label{the-actual-near-point-and-all-nonzero-original-coordinates}

The literal actual point of LIFT41 NP12, also used in DO088 WM5, satisfies at every CM label \[
 \sinh^2r_{p,\tau}=\frac{77Q^4-14Q^2}{T}<2^{-704}.
 \tag{SV13}
\] Its compact directions do not affect this Cartan radius. This point is the same one used by the rank-two near-point argument; no rank-one matrix hypothesis is needed to identify its radius.

For \(N=h-I\) and \(\alpha=(N^2)_{nn}>0\), DO080 section7 proves \[
 \alpha_{p,\tau}\ge e^{-2r_{p,\tau}}\alpha_\tau,\qquad
 \sinh^2r(g_p^{-1}hg_p)\ge\alpha_{p,\tau}.
 \tag{SV14}
\] Keep this multiplicative estimate, without replacing it by the old local loss2. Since \(r_{p,\tau}<2^{-352}\) and \(168\,2^{-352}<1/2\), the geometric-series bound for the exponential gives \[
 \prod_\tau e^{r_{p,\tau}}<2.
\] Consequently SV11--14 give \[
 |c(h)|<
 2\,\vartheta^{168}\,2^{-933f/2}
       N_F(\alpha)^{-1/2}.
 \tag{SV15}
\]

The original ideal is \(L=JI_\Sigma D_\eta\), \(N_EL=z^{2f}\). The actual k nonzero S coordinates \(\xi_i\in L\) satisfy \[
 \alpha_\tau\ge U\sum_{i=1}^k|\xi_{i,\tau}|^2,\qquad
 \prod_\tau|\xi_{i,\tau}|^2\ge z^{2f}.
\] The universal product-of-sums inequality proved in K106 NT9 gives \[
 N_F(\alpha)\ge U^f(kz^2)^f.
\] Applying it to SV15 before any Gaussian cardinality enlargement yields \[
 |c(h)|<
 2\,\vartheta^{168}k^{-84}\,
      2^{-933f/2}U^{-f/2}z^{-f}.
 \tag{SV16}
\] For the numerical cell we retain the exact inequalities \[
 \vartheta^{168}<2^{-649},\qquad 138^{84}>2^{597}.
\] All trace factors are uniform in x, z and the finite G pattern. The former complete normalization \(2^{-929f/2}\) is improved by \(2^{1-2f}=2^{-335}\), before the additional radial precision. No ambient zero tuple is assigned a trace estimate.

\subsection{5. A saturated original transverse basis}\label{a-saturated-original-transverse-basis}

The following geometric improvement was independently supplied and source-checked by /root/definite\_oscillator in saturated\_transverse\_basis\_pending.md. We include its complete required proof here.

Use the original packet vector \[
 u=(2,95192390527,383092671043,12529843548438452388609),
 \quad u^tu=U=77Q^2-14
\] and integral columns \[
 y_1=(-976627,-153032,38026,0),\quad
 y'_2=(-283839,347785,-86419,0),\quad
 z_0=(256766,-7654233287,-30805128044,1).
 \tag{SV17}
\] All are perpendicular to u. Their four maximal minors are \((u_4,-u_3,u_2,-u_1)\), with gcd1. Their image is therefore the saturated kernel of \(u^t:\mathbf Z^4\to\mathbf Z\). Tensoring the split exact sequence with \(\mathcal O_E\) gives an integral basis of the original transverse intersection. Appending u has determinant of modulus U, not1; no projection denominator U is removed.

K088 IP5--6 applies to every original integral y perpendicular to u. At the selected T component the actual shear has \(q_*y=-\pi my\), so \(A_c^{-1}y=y+czmy\) is integral. At the CM partner, y pairs integrally with the standard lattice and as zero with the pole direction z, proving membership in the exact dual lattice. Thus \(y'_2\) itself belongs to each actual source and target, including at dyadic places. This is not division of a previously known lattice member by2.

The normal form and the original primitive-u Bezout extraction put both \(r(y)\) and \(\eta r(y)\) in \(JI_\Sigma\). At G the new frame has determinant U, a G-unit; at the other finite places use the original integral inputs and actual modules. Applying the exact denominator intersection gives all three new values in the same \(L=JI_\Sigma D_\eta\), including at poles of eta.

Their true Gram is \[
 G'=\begin{pmatrix}
 978665066829&220696428039&-303935233242\\
 220696428039&208987227707&47956510467\\
 -303935233242&47956510467&1007543201084982869062
 \end{pmatrix},\qquad\det G'=U.
 \tag{SV18}
\] The leading principal minors of \(G'-I\) are \[
 978665066828,\quad155821585819917437751047,\quad
 156996979347148096974710532833471513864461899,
\] all positive. For d=k+3 original value coordinates the primal metric is exactly \(\operatorname{diag}(I_k,(G')^{-1})\), with real covolume \[
             D_0^d z^{2df}/U^f.
 \tag{SV19}
\] Its dual Gram \(\operatorname{diag}(I_k,G')\) still dominates the scalar product metric. Relative to the old enlarged value lattice, the Gaussian prefactor improves by exactly \(4^{-f}=2^{-336}\). The central difference module, its conservative cover and all finite multiplicities are unchanged.

\subsection{6. Directed scalar theta and the exact eligibility endpoint}\label{directed-scalar-theta-and-the-exact-eligibility-endpoint}

Put \(D_0=7^{140}29^{144}5^{84}13^{84}\), \(D_F=2^fD_0\), \(R=17424/133\), \(M^2=R/U\), and retain the DO112 cap \(A_*=59/3276800\). In the original half-trace metric the scalar dual is \[
 L^\#=2\mathfrak D_E^{-1}\bar L^{-1},\qquad
 \lambda_1(L^\#)^2\ge fD_0^{-2/f}z^{-2}.
 \tag{SV20}
\] This is DO070 FT3 / K105 CG5 / DO112 EC9. The exact integer comparison \(2^{168}D_0<1463^{168}\) gives \(r_0=D_0^{1/f}<1463/2\).

Here is the full dimension-dependent theta estimate used below. For any real lattice \(\Lambda\) of dimension n with minimum squared norm at least \(\mu^2\), write \(\Theta_\Lambda(s)=\sum_{\lambda\in\Lambda}e^{-s\|\lambda\|^2}\). For \(0<t<1\), the minimum and Gaussian Poisson summation give \[
 \Theta_\Lambda(s)-1
 \le e^{-(1-t)s\mu^2}\Theta_\Lambda(ts)
 \le e^{-(1-t)s\mu^2}t^{-n/2}\Theta_\Lambda(s).
 \tag{SV21}
\] Indeed the dual Gaussian in \(\Theta_\Lambda(ts)=(\pi/(ts))^{n/2}
(\operatorname{covol}\Lambda)^{-1}
\Theta_{\Lambda^*}(\pi^2/(ts))\) decreases when its argument increases. This proves the required specialization directly; all sums converge absolutely.

Take \(n=336\), \(s=\pi^2z^2/A\), \(0<A\le A_*\), and use the rational lower bound \[
 u_0=\frac{(314159/100000)^2}{A_*(1463/2)^2}
     =\frac{50532289167872}{49328816796875}>1,\qquad t=u_0^{-1}.
 \tag{SV22}
\] The bound \(\pi>314159/100000\) follows without decimal evaluation: Machin's identity is certified by the exact tangents5/12,120/119 and1, and the positive angle \(4\arctan(1/5)-\arctan(1/239)\) is less than 4/5 and hence less than \(\pi/2\). The alternating arctangent series gives \[
 \pi>
 16\sum_{j=0}^{5}\frac{(-1)^j}{(2j+1)5^{2j+1}}-\frac4{239}
 >\frac{314159}{100000}.
\] For \(w=u_0-1\), with \(0<w<1\), \[
 \delta=f(w^2/2-w^3/3)
 =\frac{843397831864796206225313406896708926716}
 {17147628093021546628449738025665283203125},\qquad
 \delta+\delta^2/2>\frac1{20}.
 \tag{SV23}
\] The inequality \(\log(1+w)\le w-w^2/2+w^3/3\) implies that the multiplier in SV21 is strictly less than \(e^{-\delta}<20/21\). Rearranging therefore yields \[
          \Theta_{L^\#}(\pi^2z^2/A)<21\qquad(0<A\le A_*).
 \tag{SV24}
\] The Poisson loss is \(t^{-168}\), so this is a real dimension336 result. The true dual metric from SV18--19 dominates the scalar product metric; its entire mass is thus bounded by \(21^d\), not by21. Poisson summation in the primal metric and the actual joint energy \(\|r\|^2\le fM^2\) give \[
 \#r\le {164\choose k}\,21^d\,\frac{U^f}{D_0^d}
       \left(\frac{\pi}{A}\right)^{df}
       e^{AfM^2/z^2},\qquad d=k+3.
 \tag{SV25}
\] The complex phases are counted in this very lattice sum. There is no additional relative-norm, phase, eta or source-isometry enumeration. Eta, all original row values and the central scalar reconstruct the matrix. SV16 is applied to actual nonzero k-tuples before their cardinality is enlarged.

The denominator floor is the exact selected-eligibility norm inequality K103 SC26.9 / K099 GE4: \[
 x>\frac{7^{84}T^{26}}{[4(U+14)]^{168}}.
 \tag{SV26}
\] It includes the actual denominator and valuations above7. It uses one selected prime per eligible place; it does not double that T exponent. The strict integer comparison \[
 (7^{84}T^{26})^{64}>
       2^{71443}[4(U+14)]^{10752}
\] proves \[
                         x>2^{X_0},\qquad X_0=71443/64.
 \tag{SV27}
\] Eta remains unique per support by the earlier bounded-denominator theorem. No support, label or scalar is newly enumerated by improving this endpoint.

\subsection{7. Complete count and chronological replay}\label{complete-count-and-chronological-replay}

Retain the original central affine bound of K103 SC26.14--16, before discarding its rank-zero G divisor, and the exact DO098 CF10 correlation: \[
 \#d_{\rm cent}\le
 \frac{V_F(H/2+750Uz^2)\,U^{-f}T^{52}x^{-1}}{N_FJ_0},
 \quad
 V_F(C)=\frac{((22/7)fC^2)^{f/2}}{(f/2)!\sqrt{D_F}},
\] \[
 |a_G(h)|\le2Q^{-335\cdot167}(N_EJ)^{335/2},\qquad
 N_EJ=z^{2f}T^{52}/x\le Q^{167}N_FJ_0 .
 \tag{SV28}
\] The central real-part shift is fixed by the row. The covolume is that of the actual difference ideal, not an assertion that the scalar itself is in that ideal. Its full-real cover750 and the U projection denominator remain unchanged. All26 active weighted-label factors and142 inactive pointwise factors give \((T-1)^{-168}<2T^{-168}\), by K103 SC26.18--20 / DO105 EL16. This bound is on the retained set; it is not a claim about a larger untruncated central orbit. Both mixed projectors keep their modulus-one bounds, with the original source, phase, class-mass and Haar normalization.

For \(q=M^2/z^2\), choose \(A(q)=\min(A_*,d/q)\) and put \[
 F_k(q)=q^{-167f}\left(\frac{\pi}{A(q)}\right)^{df}
             e^{A(q)fq}V_F(H/2+750R/q).
 \tag{SV29}
\] Combining SV16/25/28 and the T factors gives, apart from positive factors independent of q, \(x^{-337/2}(N_FJ_0)^{-1}F_k(q)\). The new trace and basis factors are independent of q and x, so the complete monotonicity can be checked directly. Write \(\epsilon(q)=(750R/q)/(H/2+750R/q)\in(0,1)\). The normalized logarithmic q derivative of \(F_k\) is \[
 -167+A_*q-\epsilon(q)
 \quad\hbox{or}\quad -167+d-\epsilon(q),
 \tag{SV30}
\] in its two branches, and the values agree at the boundary.

For fixed eta, the all-rank-zero endpoint is \[
 z_*=(Q^{334}T^{-52}x)^{1/(2f)},\qquad q_*=M^2/z_*^2.
\] The exact correlation in SV28 gives \((N_FJ_0)^{-1}\le Q^{-167}(q/q_*)^f\). At k=138 the enlarged q factor consequently decreases at least with exponent \(-(163-k)f=-25f\). The original local J exponents2,1,0 have multiplicities1,2,1, with exponent2 only at rank zero. Summing all such states and their two orientations gives \[
 \sum_{\rm patterns}\frac{F_{138}(q)}{N_FJ_0}
 \le Q^{-167}F_{138}(q_*)\prod_G(1+q_G^{-25})^2
 <2Q^{-167}F_{138}(q_*).
 \tag{SV31}
\] The last strict comparison follows by the exact integer inequality \((1+g_0^{-25})^{668}<2\), \(g_0=22285932761\). It is valid even if denominator primes of eta overlap G: J,J0 retain their original finite-rank definitions and D\_eta is separately in L. No ideal homothety or independently variable eta is assumed.

Since \(q_*\) is proportional to \(x^{-1/f}\), the full logarithmic x derivative of \(x^{-337/2}F_k(q_*)\), after SV31, is \[
 -3/2-A_*q_*+\epsilon(q_*)< -1/2-A_*q_*<0
 \quad\hbox{or}\quad
 -3/2-d+\epsilon(q_*)<-1/2-d<0 .
 \tag{SV32}
\] The recovered \(Q^{-167}\) is independent of x. Thus the entire bound, including the central radius and all patterns, is largest at \(x=2^{X_0}\). The piecewise Gaussian choice introduces no endpoint union.

At that endpoint the exact host integers give \[
 q_*^f=\frac{R^fT^{52}}{U^fQ^{334}2^{X_0}},\qquad
 c_-=4119905<q_*<4119906=c_+,
\] \[
 A_*c_+<141,\qquad H/2+750R/c_-<68,\qquad V_F(68)<2^{479}.
 \tag{SV33}
\] The root brackets are checked by integer powers64f; the volume bound is checked after squaring with the exact \(D_F\). The endpoint is in the constant-A branch. Retain the conservative reconstruction-pivot factor \(164<2^8\), the full support count \({168\choose26}<2^{101}\), one weighted-pattern factor2, one Green2, and the single combined T correction2. Their discrete exponent is \(101+8+1+1+1=112\).

The directed bounds \[
 L_T=1019333/1024,\quad L_Q=287843/4096,\quad
 L_U=601356/4096,\quad L_R=28810/4096
\] are upper for T,R and lower for Q,U. All signed powers are combined before substitution. For transparent comparison with the old determinant and trace, write \[
 \begin{split}
 C(X_0)={}&112+2f-929f/2+479+167fL_R-(335f/2)L_U\\
          &+8594L_T-336\cdot167L_Q-(337/2)X_0
          =211125067/512 .
 \end{split}
 \tag{SV34}
\] Here \(8594=335\cdot26+52-168\); the additional Q exponent167 is precisely the restored central divisor in SV31. The old determinant factor \(4^f\) is the displayed2f and will be removed exactly once below. Likewise the old selected factor \(-929f/2\) is replaced using the exact335-bit gain proved in SV16.

For d=141 define the positive rational number \[
 W_{138}=
 \frac{{164\choose138}21^{141}(22\cdot3276800)^{141f}}
 {(7\cdot59)^{141f}D_0^{141}c_-^{167f}}.
 \tag{SV35}
\] Its least strict integral binary upper exponent is \(-428684\). The inherited bounds \(\pi<22/7\), \(\log_2e<13/9\) are unchanged, and \((13/9)fA_*c_+=3686629219/204800\). The complete new cell exponent, with saturation336, full tail-inclusive trace335, radial649 and all-coordinate597 each used once, is \[
 \begin{split}
 E_{138}
 &\le\frac{211125067}{512}+
       \frac{3686629219}{204800}
       -428684-336-335-649-597\\
 &=-\frac{50428781}{204800}<-246.
 \end{split}
 \tag{SV36}
\] All support, label, row-position, pivot, central and G-pattern counts are already inside this one bound. There is no further endpoint or range union.

The exact certificate independently checks the universal moment-factor identity, actual word/twist parameter chambers, the full probability-tail constant, actual near point, saturated integer lattice and true Gram, dimension336 theta arithmetic, the eligibility root, both complete monotonicity formulas on diagnostic samples, all signed powers, exact rational ceilings and the full ordered budgets. The universal probability, Hilbert, local-module and derivative arguments are the proof above; finite samples do not replace them.

The actual base is canonical checks/k106\_joint\_remainder\_error.json, 113 all-top and31 uniform entries, floors26/11. The companion accepts that prefix or a later complete canonical ordered extension. Append246 once, giving114/31 at this base, with all old overlapping charges unchanged. Both full sums remain below \(2^{-230}\); the next-smallest charge after231 is246. This cell removal does not increase the support floor. Coordinates 139 through164, signed positivity, selected arithmetic-image content, global period, effective source, Read/specific032, actual DP44 and RH remain separate obligations.

\subsection{8. Primary and carried source scope}\label{primary-and-carried-source-scope}

The actual Hilbert model is Nakahama, R. (2016), \emph{Norm computation and analytic continuation of vector valued holomorphic discrete series representations}, Journal of Lie Theory,26(4),927--990, equation5.7 and Theorem5.1, pp960--961, through K043/K046's literal dictionary. K046 BL8--10 gives the actual positive moments; DO080 RT20--24 supplies their uniform all-word bounds. DO080 RT8--19 gives the exact density, detailed balance and Cauchy argument; DO092 SN16--18 supplies the universal radial profile. LIFT41 NP12 and DO088 WM5 identify the actual point's radius; DO080 section7 is the same rank-two transport.

The original paired modules are K088 IP4--6, DO098 CF4--10 and K103 SC26.5--20. The Gaussian convention is Miller, S. D., and Stephens-Davidowitz, N. (2019), \emph{Kissing numbers and transference theorems from generalized tail bounds}, arXiv:1802.05708v2, equations1.1/2.1, \url{https://arxiv.org/pdf/1802.05708}, with the complete required specialization proved in DO112 EC9--17. The saturated-basis calculation is incorporated above with its full integral-input argument.

The same field, involution, polarization, theta line, canonical traces, projective modules and finite projectors remain fixed. Carried context remains SGA3 I4.7.3/4.7.3.1; SGA6 VI Theorem3.3; FGA232 Theorem3.1 pp148--149; FGA236 Corollaries1.7/2.3 pp229/232; SGA7 I VII2.7/2.9 pp169--174 and VIII7.4(b) pp310--311; cote4 scan003 sections8.1.1--8.1.2/scans005--006 sections8.4.1--8.4.3 and cote54 scans005--007. These are inherited contextual locators, not new primary readings or new arithmetic comparison assertions.

\subsection{Canonical algebraic-geometric context}\label{canonical-algebraic-geometric-context-37}

The original fractional invertible modules, Hermitian form, involution, polarization, theta line, paired incidences and selected complex realization are unchanged. Carried exact anchors: SGA3 I4.7.3/4.7.3.1; SGA6 VI Theorem3.3; FGA232 Theorem3.1 pp148--149; FGA236 Corollaries1.7/2.3 pp229/232; SGA7 I VII2.7/2.9 pp169--174 and VIII7.4(b) pp310--311; cote4 scan003 sections8.1.1--8.1.2, scans005--006 sections8.4.1--8.4.3 and cote54 scans005--007. No new lambda-ring/K-theory, l-adic weight, formal-group or CM-motive comparison. DP\_MULT Q3e/DP44/DP52/DP54--65 and THETA3 G1/A7 gain the stated smaller finite gate; CX\_GRH CX-ALG-024 and GCT\_PLETH T-052 retain their arithmetic gates. Primary Gaussian/Poisson context remains Miller, S. D., and Stephens-Davidowitz, N. (2019), arXiv1802.05708v2, equations1.1/2.1, printedpp1--2/4; the new-cap specialization is proved in full here. The sharper trace is the same genuine vector-valued selected normalized trace and original near-point argument; no operator-norm or rank-one transport substitute is used. Nakahama(2016), Journal of Lie Theory26(4),927--990, equation5.7 and Theorem5.1 printed960--961, supplies the carried representation convention. The new bound uses the same selected Hilbert metric and integral saturated kernel of the original primitive packet row; the projection norm U and central affine difference module remain. The exact tail is absorbed before counting.

\subsubsection{Mandatory checked boundary qualification}\label{mandatory-checked-boundary-qualification}

\section{AM-K-107 --- explicit boundary precision for the preserved frozen proof}\label{am-k-107-explicit-boundary-precision-for-the-preserved-frozen-proof}

Author: /root/kummer\_extension. Independent observations: /root/lift\_accounting.

The preserved subject is sharpened\_radial\_variance\_support\_twenty\_six.md, SHA256 3ad12d483d85b2f0e20917d5f78c0b5a31f4e9a726ea1fab18c746c4788dda2a. The parent requested that these frozen bytes remain unchanged. Read that subject with the following two precise qualifications.

\begin{enumerate}
\def\labelenumi{\arabic{enumi}.}
\tightlist
\item
  In SV6 the final strict inequality holds for integer h\textgreater=1. At h=0 both sides are exactly1. Equivalently the final displayed sign is \textless= for h\textgreater=0, and is strict for h\textgreater=1. Every use in SV7 has h=2\^{}19, so its strict probability bound and all subsequent constants are unchanged.
\item
  The quotient in SV11 is stated for r(g)\textgreater0. At r(g)=0 retain the actual normalized bound \textbar c\_W(g)\textbar\textless=1. The nonidentity rank-two application SV14--16 has positive alpha and hence positive Cartan radius; no division by zero occurs in that application.
\end{enumerate}

These are boundary qualifications, not a new estimate or sector charge. The full lower-tail absorption, saturated original lattice, theta21, complete exponent -50428781/204800 and charge246 are unchanged. The mathematical subject for the whole author/peer/root verdict is the frozen proof together with this explicit qualification. The script remains unchanged, and its final6,118 exact checks already include the h=0 factor identity; its tail calculation uses h=2\^{}19. This note does not amend a canonical file or error array.

\clearpage

\section{Active type-II local saturation: a boundary for additional congruence savings}\label{active-type-ii-local-saturation-a-boundary-for-additional-congruence-savings}

\textbf{Canonical checked integration --- 2026-11-23T06:00:48.470982+00:00.} Complete saved-proof and executable reviews and separate exact diagnostic replays passed. No numerical error charge is introduced. Frozen author history is preserved. AM-DO-113 LS1--4 proves local sufficiency of the retained simple cross and double central pole allowances in the original paired unipotent normal form, conditional on the actually admitted cancellation label and the unitary relation. Both cancellation factors are divisible by the uniformizer, all three conjugated terms are integral separately, and nilpotence gives an integral inverse. Thus reimposing this same active-pencil preservation alone supplies no additional local linear or quadratic congruence beyond those allowances and the unitary relation. The existing pending row is closed in this local scope. No arbitrary local parameter is asserted to globalize, satisfy shortness/G/total support, or attain the enlarged count. No numerical charge or main finite-sector removal; complete114/31 arrays and floors26/11 are unchanged.

\textbf{Pending scoped algebraic note. No error charge or new sector removal.} Author: /root/definite\_oscillator. 23 November 2026 UTC.

\subsection{1. Actual local hypotheses}\label{actual-local-hypotheses}

Fix one original selected split T place and its companion component. Retain the actual normal form and the literal same-pencil representative from DO105 EL9--15 and K103 SC26.6--7/12--16: \[
 M_c=A_c\Lambda,\qquad A_c=I+c\pi^{-1}zq_*,
 \quad q_*z=0,\quad mz=1,\quad q_*=-v_0^\dagger-\pi m .
\] Here z and v0 lie in the packet plane P, the integral functional m is the allowed original representative, and the fixed original positive Z frame is integral and invertible at T. The paired lattice is the exact inverse-adjoint companion. We make no assertion about a positive residue-field metric.

For the actual isotropic column w=(u,eta) and row w†, let \[
 N=wR-Vw^\dagger+dww^\dagger,\qquad
 R|_P=0,\quad V\in Z,\quad V=J_Z^{-1}R^*,\quad
 d+\bar d=-RV .
 \tag{LS1}
\] Thus N³=0 and h=I+N is the original paired-unitary unipotent normal form. Eligibility makes eta and its conjugate T-units. Assume the unique cancellation label c\_eta is actually admitted. DO105 EL10--12 then gives \[
 w_c=A_c^{-1}w=\pi W,\qquad
 t_c=w^\dagger A_c=\pi T_0,\qquad W,T_0\text{ integral}.
 \tag{LS2}
\] An eligible eta for which this label is omitted is not being assigned an arrow.

\subsection{2. The apparent additional constraints are already satisfied}\label{the-apparent-additional-constraints-are-already-satisfied}

The two exact identities DO105 EL14 are \[
 RA_c=R,\qquad
 q_*V=-\pi mV,\qquad A_c^{-1}V=V+cz\,mV .
 \tag{LS3}
\] They remain valid when R and V have simple poles. They use Rz=0 and v0†V=0, not integrality of those two vectors.

Suppose R and V have at most a simple pole in the original T-unit frame and d has at most a double pole. Since c,z,m are integral, \(A_c^{-1}V\) still has at most a simple pole. Conjugating LS1 and using LS2--3 gives \[
 A_c^{-1}NA_c
 =\pi WR-\pi(A_c^{-1}V)T_0+\pi^2dWT_0 .
 \tag{LS4}
\] Every term is integral separately. As N³=0, the inverse of I+A\_c\^{}−1NA\_c is the integral polynomial I−A\_c\^{}−1NA\_c+(A\_c\^{}−1NA\_c)². Hence h stabilizes M\_c, not merely maps it into itself. Its prescribed inverse-adjoint companion stabilizes the paired lattice.

This proves sufficiency of the very pole orders used in the current row and central count, once the actual cancellation label exists. Necessity of these orders is already in K103: all fixed Z inputs belong to the original source and target, giving simple cross poles; the original full-entry target bound gives order at most two for the central scalar in the T-unit plane frame. Thus there is no extra local quadratic or linear congruence from this active same-pencil condition beyond these pole allowances and the unitary normal-form relation.

The statement concerns local preservation. It does not say that arbitrary locally allowed parameters occur globally, meet the shortness cutoff, satisfy the chosen G conditions, or have the exact desired total support. If an original S row value genuinely has a simple pole, h is nonintegral there; the existing type-I integrality lemma rules out type I in that subcase. The one-label argument is unchanged. This is not a construction of a global short arrow.

\subsection{3. Paired rational-DVR diagnostic with a transverse component}\label{paired-rational-dvr-diagnostic-with-a-transverse-component}

The companion uses exact rational matrices at p=5,7,11 in coordinates (S,T,P+,P−), J=diag(1,1,1,−1). Put \[
 z=(0,0,1,1+2p)^t,\quad
 m=(0,1,1,0),\quad
 b=(1+p)/(1+2p),\quad
 q_*=(0,-p,-1-p,b).
\] Then mz=1 and q\_\emph{z=0. The selected original v0† is (0,0,1,−b), so q\_}=−v0†−pm exactly. Take w=(0,0,1,1)\^{}t, w†=(0,0,1,−1), and c=(p−1)/2. This is a nonzero admitted label, distinct from−1. Direct calculation gives the two cancellation divisibilities in LS2.

The row is R=(a/p,b1/p,0,0), its paired adjoint vector is V=(2/p,3/p,0,0)\^{}t, and d=epsilon/p² for epsilon=±1. The companion component uses the exchanged row/vector and dbar=−RV−d.~This is the split-field involution; no real positivity is inferred from these rational component matrices. The chosen m has mV=3/p≠0, so LS3's transverse correction is actually exercised.

The diagnostic checks both component matrices, N³=0, exact paired unitarity, exact inverse-adjoint companion A, genuine original row and central poles, and integrality of both conjugated matrices and their inverses at p.~It also checks rank(N)=2 over the rational component field. These examples are a finite local scope test, not a proof of a global family. The universal calculation is LS1--4.

\subsection{4. Consequence for the dense count}\label{consequence-for-the-dense-count}

The present Gaussian row/central enlargement does not leave an unpaid active-T congruence that can be recovered solely by reimposing hM\_c=M\_c in this normal form. A further saving must use a genuinely additional global condition, a relation not represented by these pole orders, a stronger actual selected-trace estimate, or cancellation. This note claims no impossibility theorem for other counting arguments and no numerical lower bound for the contribution.

Exact carried sources: canonical agents/eligible\_line\_uniqueness\_and\_middle\_support.md, EL9--15; agents/support\_twenty\_six\_all\_ii\_reduction.md, SC26.6--7/12--16; agents/type\_iii\_cross\_integrality.md for the separate type-I integrality scope. The original packet, involution, lattice pairing and full selected representation remain unchanged.

\subsection{Canonical algebraic-geometric context}\label{canonical-algebraic-geometric-context-38}

The proof computes the integral stabilizer of the original paired lattices over the selected split local ring, retaining the Hermitian form, involution, admitted pencil label, unipotent group coordinates and polarization. No positive residue-field metric is assumed. It gives no new lambda-ring/K-theory identity, l-adic weight theorem, CM-motive or oscillator comparison. Carried exact anchors: SGA3 I4.7.3/4.7.3.1; SGA6 VI Theorem3.3; FGA232 Theorem3.1 pp148--149; FGA236 Corollaries1.7/2.3 pp229/232; SGA7 I VII2.7/2.9 pp169--174 and VIII7.4(b) pp310--311; cote4 scan003 sections8.1.1--8.1.2 and scans005--006 sections8.4.1--8.4.3; cote54 scans005--007. DP\_MULT Q3e/DP44 and THETA3 G1/A7 gain a checked boundary for this specific congruence-count route, not a completed source/period construction; CX\_GRH CX-ALG-024 and GCT\_PLETH T-052 retain their separate arithmetic gates.

\clearpage

\section{AM-DO-114 --- spherical packing and three further complete coordinate sectors}\label{am-do-114-spherical-packing-and-three-further-complete-coordinate-sectors}

\textbf{Canonical checked integration --- 2026-11-23T06:22:26.870952+00:00.} Complete root and independent proof/script/interface reviews and separate exact replays passed. The actual canonical predecessor has114/31 entries; this integration gives115/31, complete floors26/11. Frozen author history is preserved. AM-DO-114 SP1-27 proves a uniform scalar theta below3/2 in real dimension336 for A\textless=1/43500 by a complete finite spherical-cap and nonzero-shell argument. Applied to the original saturated transverse Gram, actual selected trace, full affine central fibre and every finite pattern, the whole139..141 union is below2\^{}-354. The necessary actual support26 all-II rank-two-unipotent range becomes142\textless=k\textless=164 with2\^{}(71443/64)\textless x\textless2\^{}2108,v=q26,b=j0. Exactly one354 append gives115/31 complete ordered errors, both full sums below2\^{}-230, with every historical overlap retained and complete floors26/11 unchanged. No support27, signed positivity, selected arithmetic content, period, effective source, specific032/Read, evaluated DP44 or RH conclusion.

Author: /root/definite\_oscillator. 23 November 2026 UTC. Controlled pending proof. The final author read, executable replay, independent whole review and canonical integration are separate events.

\subsection{1. Exact retained sector and statement}\label{exact-retained-sector-and-statement}

Keep the actual J56 all-top family, H=133, f=168, the original fields E/F, involution, Hermitian metric, paired source and target pencils, admitted labels, selected vector-valued real module and its Hilbert metric, both compact factors, canonical twists, finite projectors, phases, class masses and Haar normalization. After canonical K107 the only possible support26 branch has global rank-two unipotence, scalar cross, all-II type and

\[
 s=26,\quad v=q_{\rm elig}=26,\quad b=j=0,\quad
 2^{X_0}<x=N_ED_\eta<2^{2108},\quad
 X_0=71443/64,\quad139\le k\le164.                 \tag{SP1}
\]

Here k counts the nonzero original S covector coordinates. K103 SC26.5--7 gives the support-free full positive compression and exact unipotent normal form. Its SC26.11 fixes eta uniquely for each support in this denominator window. K107 SV26--27 proves the sharper lower endpoint, with the actual valuations over7. No eta or denominator ideal is independently enumerated below.

\textbf{Theorem.} The complete subset of SP1 with \(139\le k\le141\) has absolute normalized contribution less than \(2^{-354}\). Append354 once, retaining every prior charge and both complete ordered arrays. The new necessary coordinate range in this all-II branch is \(142\le k\le164\), with every other SP1 condition unchanged. The complete support floors remain26/11. No support27 or signed-positivity assertion follows.

The new ingredient is an explicit scalar-lattice theta bound in real dimension336, proved from a finite spherical-packing estimate. It permits \(A\le1/43500\) with scalar theta below3/2. The proof retains the complete central affine count and checks both Gaussian-parameter branches and the coupled denominator endpoint. It uses no new decay estimate beyond the checked K107 selected trace.

Set

\[
 Q=1427908193729121522601,\quad U=77Q^2-14,
 \quad R=17424/133,\quad M^2=R/U,
\] \[
 D_0=7^{140}29^{144}5^{84}13^{84},\qquad D_F=2^fD_0.
                                                        \tag{SP2}
\]

\subsection{2. A dimension-explicit finite packing estimate}\label{a-dimension-explicit-finite-packing-estimate}

First let \(v_1,\ldots,v_m\) be unit vectors in \(\mathbb R^n\) with pairwise inner products at most zero. Their Gram matrix G is positive semidefinite of rank at most n, with diagonal entries1 and off-diagonal entries in{[}-1,0{]}. Therefore

\[
 \sum_{i\ne j}G_{ij}\ge-m,\qquad
 \operatorname{tr}(G^2)
 =m+\sum_{i\ne j}G_{ij}^2
 \le m-\sum_{i\ne j}G_{ij}\le2m.
\]

Applying Cauchy--Schwarz to its nonzero eigenvalues gives

\[
 m^2=(\operatorname{tr}G)^2\le n\operatorname{tr}(G^2)
 \le2nm,\qquad m\le2n.                                \tag{SP3}
\]

This is the orthoplex bound; the argument above is complete. Rankin(1955), Theorem1(iv), is its classical spherical-packing reference. No asymptotic packing estimate or unspecified dimensional constant is used.

Let a finite set X in \(B_n(\rho)\) have mutual distances at least1, where \(\rho\ge1\). The spherical lift

\[
 x\longmapsto \widetilde x
 =\rho^{-1}(x,\sqrt{\rho^2-\|x\|^2})\in S^n
 \subset\mathbb R^{n+1}
\]

does not decrease distances before division by rho. Thus distinct lifted points have inner product at most

\[
 s_\rho=1-\frac1{2\rho^2}\in[1/2,1).
                                                        \tag{SP4}
\]

Fix a spherical cap of angular radius theta with

\[
 \cos\theta=\sqrt{s_\rho},\qquad
 \sin\theta=\frac1{\sqrt2\rho}.
\]

For each lifted point in the cap, write its pole coordinate as a and its orthogonal projection as p.~Then \(a\ge\sqrt{s_\rho}\). For two such points,

\[
 \langle p_i,p_j\rangle
 =\langle\widetilde x_i,\widetilde x_j\rangle-a_i a_j
 \le s_\rho-a_i a_j\le0.                              \tag{SP5}
\]

A point with zero projection is the cap pole. It cannot coexist in this cap with another point, because the other point's inner product with the pole would be at least \(\sqrt{s_\rho}>s_\rho\). Otherwise normalize all projections, which lie in \(\mathbb R^n\), and use SP3. Thus every cap contains at most2n lifted points. The sphere has ambient dimension n+1; the projected vectors have dimension n.~They are not conflated.

The normalized area of this cap is

\[
 \mu_n(\theta)=c_n\int_0^\theta\sin^{n-1}\phi\,d\phi,
 \qquad c_n=\frac{\Gamma((n+1)/2)}{\sqrt\pi\Gamma(n/2)}.
\]

Since \(0\le\cos\phi\le1\) on this interval, integrating the derivative of \(\sin^n\phi\) gives

\[
 \mu_n(\theta)\ge\frac{c_n}{n}\sin^n\theta.
\]

Average the cap occupancy over its pole using rotationally invariant probability measure. Each point has probability \(\mu_n(\theta)\) of belonging to the cap, so

\[
 \#X\le C_n(\sqrt2\rho)^n,
 \qquad C_n=\frac{2n^2}{c_n},\qquad \rho\ge1.           \tag{SP6}
\]

This proof works for arbitrary finite separated sets. It assumes neither a positive metric over a finite field nor any special shape of an ideal lattice. For n\textgreater=2, the weaker \(c_n\ge1/2\) follows from \(\int_0^\pi\sin^{n-1}\phi\,d\phi\le2\); the application instead retains the exact c\_n.

\subsection{3. From nonzero shells to a uniform scalar theta bound}\label{from-nonzero-shells-to-a-uniform-scalar-theta-bound}

Let Lambda be a real n-dimensional lattice with minimum at least lambda\textgreater0. Write N(rho) for the number of its \textbf{nonzero} vectors of length at most rho lambda. It is zero for rho\textless1; for rho\textgreater=1, apply SP6 to the scaled lattice points. With \(a=t\lambda^2>0\), Tonelli's theorem for the nonnegative summands gives

\[
 \Theta_\Lambda(t)-1
 =\sum_{v\ne0}e^{-t\|v\|^2}
 =\int_1^\infty 2a\rho e^{-a\rho^2}N(\rho)\,d\rho.
                                                        \tag{SP7}
\]

The bound SP6 makes the integral finite. Substituting it and enlarging the resulting integral to zero yields

\[
 \Theta_\Lambda(t)-1
 \le C_n2^{n/2}a^{-n/2}\Gamma(n/2+1).
                                                        \tag{SP8}
\]

The zero lattice vector has been excluded from N and restored by the1 on the left. This is a nonzero-shell estimate with rho\textgreater=1, even though the final integral is enlarged to rho\textgreater=0.

For n=336=2f, the half-integer gamma identity gives

\[
 c_{336}=168\binom{336}{168}/4^{168},\qquad
 C_{336}=\frac{2\cdot336^2\cdot4^{168}}
                 {168\binom{336}{168}}<32000.            \tag{SP9}
\]

The final strict comparison is an exact integer inequality in the companion.

For each actual scalar fractional ideal L with \(N_EL=z^{2f}\), the original half-trace metric and K107 SV20 / DO070 FT3 give

\[
 L^\#=2\mathfrak D_E^{-1}\bar L^{-1},\qquad
 \lambda_1(L^\#)^2\ge fD_0^{-2/f}z^{-2},\qquad
 r_0=D_0^{1/f}<1463/2.                                 \tag{SP10}
\]

The last strict bound follows from \(2^{168}D_0<1463^{168}\). The factor2 in the dual has already been incorporated in this normalization; no additional half-trace conversion is made. The estimate is uniform in the ideal class and denominators.

Take

\[
 A_1=1/43500,\quad 0<A\le A_1,\quad
 t=\pi^2z^2/A,\quad
 a_1=\frac{168\cdot43500\cdot(314159/100000)^2}{(1463/2)^2}.
                                                        \tag{SP11}
\]

The exact bound \(\pi>314159/100000\), proved by the alternating Machin series in K107 SV22, implies \(t\lambda_1(L^\#)^2>a_1\). Applying SP8--10 and the exact rational comparison

\[
 32000\cdot2^{168}\cdot168!\,a_1^{-168}<1/2
\]

proves the new uniform scalar estimate

\[
 \boxed{\quad\Theta_{L^\#}(\pi^2z^2/A)<3/2
                    \quad(0<A\le1/43500).\quad}         \tag{SP12}
\]

The factorial168! comes from \(\Gamma(f+1)\) in SP8. This is a336-dimensional estimate applied separately to scalar ideals, not to the entire row with the same scalar minimum.

\subsection{4. Complete actual row, central and selected-trace factors}\label{complete-actual-row-central-and-selected-trace-factors}

Fix the actual support, its unique eta and one canonical common whole-C G orientation. With the same original local ideal as K103 SC26.12--14,

\[
 L=JI_\Sigma D_\eta,\quad N_EL=z^{2f},\quad
 N_EJ=z^{2f}T^{52}/x\le Q^{334}.                         \tag{SP13}
\]

Every counted S value and each of the three saturated transverse values lies in L, including at eta denominator primes. The saturated original transverse basis is K107 SV17--19; its Gram G' has determinant U and \(G'-I>0\). Consequently, for d=k+3, the primal row metric is \(\operatorname{diag}(I_k,(G')^{-1})\), its exact real covolume is \(D_0^d z^{2df}/U^f\), and its dual Gram dominates the ordinary product metric. Using SP12 separately in each scalar ideal gives dual row mass less than \((3/2)^d\).

Gaussian Poisson summation, with Fourier convention \(\widehat{e^{-t\|\cdot\|^2}}(\xi)=(\pi/t)^{m/2}e^{-\pi^2\|\xi\|^2/t}\), and the original joint energy \(\|r\|^2\le fM^2\), now give

\[
 \#r\le\binom{164}{k}(3/2)^d\frac{U^f}{D_0^d}
            (\pi/A)^{df}e^{AfM^2/z^2},\qquad0<A\le A_1.
                                                        \tag{SP14}
\]

All complex phases are inside this lattice count. The actual nonzero k-sector trace is bounded before its cardinality is enlarged to include zero tuples. No coefficient estimate is assigned to the added tuples.

Retain the entire genuine selected bound K107 SV16, including its uniformly absorbed lower-probability tail and the actual multiplicative near-point product:

\[
 |c(h)|<2\vartheta^{168}k^{-84}\,
          2^{-933f/2}U^{-f/2}z^{-f},\qquad
 \vartheta=\frac{4^{167}}{334\binom{334}{167}},\quad
 \vartheta^{168}<2^{-649}.                              \tag{SP15}
\]

The all-coordinate norm inequality proves this for every nonzero k, not only K107's counted cell138. Its local quotient is used at positive radius, which follows from alpha\textgreater0 in this nonidentity rank-two sector. K107's explicit h=0/r=0 boundary qualification is retained. Neither an operator estimate nor a scalar Bergman model is substituted.

The central difference ideal, fixed real-part shift and conservative full-real cover750 are unchanged. If J0 is the original rank-zero F-ideal, K103 SC26.14--16 and DO098 CF10 give

\[
 \#d_{\rm cent}\le
 \frac{V_F(H/2+750Uz^2)U^{-f}T^{52}x^{-1}}{N_FJ_0},
 \quad V_F(C)=\frac{((22/7)fC^2)^{f/2}}{(f/2)!\sqrt{D_F}},
\] \[
 |a_G(h)|\le2Q^{-335\cdot167}(N_EJ)^{335/2},\qquad
 N_EJ\le Q^{167}N_FJ_0.                                \tag{SP16}
\]

This counts one affine coset if nonempty, rather than declaring its scalar integral. Eta, the row and the central scalar reconstruct the matrix exactly. No source-isometry fibre or phase factor has been dropped. The two mixed finite projectors retain their modulus-one bounds.

On SP1 the26 active weighted label factors and142 inactive pointwise factors give

\[
                     (T-1)^{-168}<2T^{-168}.             \tag{SP17}
\]

This uses actual eligibility, the one admitted active label including central-only poles, and b=j=0. It is a statement on the retained set and does not assert a bound on an enlarged central orbit.

\subsection{5. Every finite pattern and both coupled parameter branches}\label{every-finite-pattern-and-both-coupled-parameter-branches}

Write \(q=M^2/z^2\), \(d=k+3\), and choose

\[
 A(q)=\min(A_1,d/q),\qquad
 F_k(q)=q^{-167f}(\pi/A(q))^{df}e^{A(q)fq}
                   V_F(H/2+750R/q).                     \tag{SP18}
\]

Combining SP13--17 leaves a fixed positive factor times

\[
 x^{-337/2}(N_FJ_0)^{-1}F_k(q).
\]

The extra selected factors in SP15 are independent of q and x. The full T exponent is \(335\cdot26+52-168=8594\), and the grouped U and R powers are \(U^{-335f/2}R^{167f}\).

Set \(\epsilon(q)=(750R/q)/(H/2+750R/q)\in(0,1)\). In the constant-A branch, the normalized logarithmic derivative of F is \(-167+A_1q-\epsilon(q)\). In the branch A=d/q it is \(-167+d-\epsilon(q)\). The function and derivatives agree at the switching point, and in both branches

\[
 \frac1f\frac{d\log F_k}{d\log q}\le-(164-k).            \tag{SP19}
\]

The scalar theta bound3/2 is uniform throughout both branches. No derivative of an unknown lattice theta function has been taken.

For fixed eta and x the all-rank-zero endpoint is

\[
 z_*=(Q^{334}T^{-52}x)^{1/(2f)},\qquad q_*=M^2/z_*^2.
\]

The exact correlation in SP16 implies

\[
 (N_FJ_0)^{-1}\le Q^{-167}(q/q_*)^f.
\]

Thus the complete enlarged q factor decreases with exponent at least \((163-k)f\). The actual local J exponents2,1,0 have multiplicities1,2,1; exponent2 occurs only at rank zero. A decrease by h multiplies q by \(q_G^{h/f}\). Summing every original state and its two orientations gives

\[
 \sum_{\rm patterns}(N_FJ_0)^{-1}F_k(q)
 \le Q^{-167}F_k(q_*)\prod_G(1+q_G^{-(163-k)})^2
 <2Q^{-167}F_k(q_*),\qquad139\le k\le141.              \tag{SP20}
\]

Indeed \(163-k\ge22\), \(q_G\ge g_0=22285932761\), and the exact integer comparison \((1+g_0^{-22})^{668}<2\) suffices. No additional pattern or ideal-class enumeration appears. This remains valid if eta denominator primes meet G: J and J0 retain their original finite-rank definitions, while D\_eta is separately present in L. No homothety of different ideals is assumed.

After SP20 the recovered \(Q^{-167}\) is independent of x. Since \(q_*\) is proportional to \(x^{-1/f}\), the full logarithmic x derivative is

\[
 -3/2-A_1q_*+\epsilon(q_*)<-1/2-A_1q_*<0,
 \quad\hbox{or}\quad
 -3/2-d+\epsilon(q_*)<-1/2-d<0.                          \tag{SP21}
\]

Therefore the entire bound, including the moving central radius and every pattern, is largest at \(x=2^{X_0}\). The actual lower bound is strict. The piecewise parameter produces no extra endpoint union.

At this endpoint the exact host integers give

\[
 q_*^f=R^fT^{52}/(U^fQ^{334}2^{X_0}),\qquad
 4119905=c_-<q_*<c_+=4119906,
\] \[
 A_1c_+<142,\qquad H/2+750R/c_-<68,\qquad
 V_F(68)<2^{479}.                                      \tag{SP22}
\]

All three claimed cells are in the constant-A branch. The root brackets use exact integer powers64f; the full-real volume bound is checked after squaring with the actual D\_F.

\subsection{6. Exact endpoints, complete union and chronology}\label{exact-endpoints-complete-union-and-chronology}

Retain support \(\binom{168}{26}<2^{101}\), the harmless original reconstruction pivot164\textless2\^{}8, the one Green2, weighted-pattern2 and combined T-correction2. Their exponent is112. Positions \(\binom{164}{k}\) remain inside the Gaussian rational factor. The selected near-point factor2 is already inside SP15 and must not be added a second time. There is no new eta, inverse, type, witness or endpoint multiplier.

Use directed bounds

\[
 L_T=1019333/1024,\quad L_Q=287843/4096,\quad
 L_U=601356/4096,\quad L_R=28810/4096,
\]

upper for T,R and lower for Q,U. As a transparent comparison with K107, define

\[
 C(X_0)=112+2f-929f/2+479+167fL_R-(335f/2)L_U
 +8594L_T-336\cdot167L_Q-(337/2)X_0
 =211125067/512.                                       \tag{SP23}
\]

The displayed old determinant factor2f is removed exactly once by the saturated-basis saving336. The old selected normalization is improved exactly once by335, since \(2^{-929f/2}2^{-335}=2^{1-933f/2}\). This is simply SP14--16 regrouped; no extra model change is made.

For d=k+3 let e\_k be the least strict integral binary upper exponent of

\[
 W_k=\frac{\binom{164}{k}3^d(22\cdot43500)^{df}}
 {2^d7^{df}D_0^d c_-^{167f}}.                           \tag{SP24}
\]

The theta product is exactly \((3/2)^d\). Retain \(\pi<22/7\) and \(\log_2 e<13/9\), so the positive energy exponent is at most

\[
 (13/9)fA_1c_+=249940964/10875.                         \tag{SP25}
\]

Let g\_k be the integer in the following table, with \(k^{84}>2^{g_k}\). The complete cell exponent is

\[
 E_k=C(X_0)+249940964/10875+e_k-336-335-649-g_k.
                                                        \tag{SP26}
\]

Every comparison is by integer powers and rational arithmetic:

\begin{longtable}[]{@{}lrrl@{}}
\toprule\noalign{}
k & e\_k & g\_k & E\_k \\
\midrule\noalign{}
\endhead
\bottomrule\noalign{}
\endlastfoot
139 & −436305 & 597 & −16065218807/5568000 \textless{} −2885 \\
140 & −435039 & 598 & −9021698807/5568000 \textless{} −1620 \\
141 & −433773 & 599 & −1978178807/5568000 \textless{} −355 \\
\end{longtable}

Thus the complete new union is bounded by

\[
 \sum_{k=139}^{141}|\text{complete k-sector}|
 <2^{-2885}+2^{-1620}+2^{-355}<2^{-354}.                  \tag{SP27}
\]

The optional next-cell diagnostic is \(E_{142}=5059773193/5568000>0\), using e\_142=−432508 and g\_142=600. It is not a removal, lower bound or existence assertion. No cell beyond141 is claimed.

The exact canonical predecessor is checks/k107\_joint\_remainder\_error.json, with114 all-top and31 uniform entries and complete floors26/11. The companion accepts a later complete ordered extension of that prefix. Append354 once, giving115/31 at the stated base, with all old overlapping charges unchanged. Both complete sums stay below2\^{}-230; the minimum231 and the existing246 remain. No support-floor change or signed sufficient inequality is inferred.

\subsection{7. Verification and source scope}\label{verification-and-source-scope}

The executable checks the exact cap normalization and32000 comparison, scalar theta factorial inequality, half-trace/root-discriminant constants, original saturated Gram and determinant, exact eligibility and root bracket, original central volume, full signed powers and discrete factors, all four displayed rational endpoints, both derivative branches, the whole finite-pattern bound and both chronological arrays. Finite Gram/cap models check normalization and boundary cases; they do not replace SP3--8's universal argument. No asymptotic numerical sphere-packing table is an input.

Fresh complete canonical K107 SV1--36 and its explicit boundary addendum were read, as was the complete canonical K105 CG1--25 proof. Fresh exact local inputs include K103 SC26.5--20, DO098 CF4--13 and DO105 EL9--16. DO112 EC3--17 and its carried complete author/peer proof fix the unchanged Gaussian convention and original factor accounting. The source manifest identifies the actual files, fingerprints and exact reading scopes. K107 supplies the genuinely selected trace, integral saturated basis and all-word bounds; no fresh whole Nakahama or SGA/FGA reading is claimed here.

The primary spherical reference is Rankin, R. A. (1955). The closest packing of spherical caps in n dimensions. \emph{Proceedings of the Glasgow Mathematical Association, 2}(3),139--144. \url{https://doi.org/10.1017/S2040618500033219}. Theorem1(iv), printed139, and its section2 proof, printed142, were directly checked. The Gram argument, n+1 lift, cap averaging, exact constant and theta integration required here are all proved explicitly above. The Gaussian Fourier convention is the carried Miller, S. D., and Stephens-Davidowitz, N. (2019), \emph{Kissing numbers and transference theorems from generalized tail bounds}, arXiv:1802.05708v2, equations1.1/2.1; no stronger kissing estimate from that paper is used.

The field, involution, fractional projective modules, polarization, theta line, selected complex realization and paired incidences remain unchanged. Inherited exact context is SGA3 I4.7.3/4.7.3.1; SGA6 VI Theorem3.3; FGA232 Theorem3.1 pp148--149; FGA236 Corollaries1.7/2.3 pp229/232; SGA7 I VII2.7/2.9 pp169--174 and VIII7.4(b) pp310--311; cote4 scan003 sections8.1.1--8.1.2/scans005--006 sections8.4.1--8.4.3 and cote54 scans005--007. These are contextual anchors only. No new lambda-ring, K-theory, l-adic weight, formal-group or CM-motive comparison is made. The remaining finite cells, signed kernel, actual selected arithmetic-image content, prescribed global period, effective source and auxiliary scales, Read/STiesR/specific032, DP44 and RH remain distinct obligations.

\subsection{Canonical algebraic-geometric context}\label{canonical-algebraic-geometric-context-39}

The original invertible fractional modules, Hermitian form, involution, paired incidences, polarization, theta line and genuine selected complex realization remain fixed. The exact dimension336 half-trace dual is2 D\_E\^{}-1 conjugate(L)\^{}-1; the saturated original transverse Gram and central affine difference module are retained. The finite cap proof is given in full, with classical context Rankin(1955), Theorem1(iv), printed139 and proof printed142, doi10.1017/S2040618500033219. Gaussian convention: Miller and Stephens-Davidowitz(2019), arXiv1802.05708v2,equations1.1/2.1, carried directly checked scope. Contextual exact anchors: SGA3 I4.7.3/4.7.3.1; SGA6 VI Theorem3.3; FGA232 Theorem3.1pp148-149; FGA236 Corollaries1.7/2.3pp229/232; SGA7 I VII2.7/2.9pp169-174 and VIII7.4(b)pp310-311; cote4 scan003 sections8.1.1-8.1.2/scans005-006 sections8.4.1-8.4.3, cote54 scans005-007. No new crystal, lambda-ring/K-theory, l-adic weight or CM-motive comparison. DP\_MULT Q3e/DP44/DP52/DP54-65 and THETA3 G1/A7 gain only this finite reduction; CX\_GRH CX-ALG-024 and GCT\_PLETH T-052 retain their global comparison gates.

\clearpage

\section{AM-K-108 --- the actual angular profile and a complete energy-weighted interval}\label{am-k-108-the-actual-angular-profile-and-a-complete-energy-weighted-interval}

\textbf{Canonical checked integration --- 2026-11-23T06:30:15.160648+00:00.} Complete root and independent proof/script/interface reviews and separate exact replays passed. The actual canonical predecessor has115/31 entries; this integration gives116/31, complete floors26/11. Frozen author history is preserved. AM-K-108 AE1-30 retains the full actual Cartesian conditional angular density, cancels its D ratio in detailed balance and inserts exp(-2U\textbar\textbar r\textbar\textbar\^{}2) in the original lattice sum. With the separately checked DO114 scalar theta, original transverse Gram, central affine fibre, common ideal correlation and all finite factors, the complete142..159 union is below2\^{}-585. The necessary support26 all-II rank-two-unipotent remainder becomes160\textless=k\textless=164, with2\^{}(71443/64)\textless x\textless2\^{}2108,v=q26,b=j0. Exactly one585 charge gives116/31 ordered errors, both full sums below2\^{}-230; floors26/11 remain. No160+, signed positivity, arithmetic-content nonvanishing, period, effective source, specific032/Read, evaluated DP44 or RH conclusion.

Author: /root/kummer\_extension. 23 November 2026 UTC. Controlled pending proof. Final author read, exact replay, independent whole review and canonical integration are separate records. DO114's uniform scalar-theta lemma is a required separately checked input; no frozen predecessor is edited.

\subsection{1. Exact scope and result}\label{exact-scope-and-result}

Retain the actual J56 all-top family, H=133, f=a=168, original fields, involution, Hermitian metric, paired source and target pencils, admitted labels, finite projectors, selected vector-valued representation and its Hilbert metric, both compact factors, words, twists, phases, Haar measures and class masses.

After K107 and the separate DO114 interval the possible all-II support26 rank-two-unipotent scalar-cross terms satisfy \[
 v=q_{\rm elig}=26,\quad b=j=0,\quad
 2^{X_0}<x=N_ED_\eta<2^{2108},\quad
 X_0=71443/64,\quad142\le k\le164.
 \tag{AE1}
\] Here k counts nonzero original S covector values. K103 SC26.5--7/11 supplies the full positive normal form and unique eta for each support. None of these facts is inferred from a different support size.

\textbf{Theorem.} The actual selected trace has the additional angular factor proved in AE9--10. In the original rank-two normal form that factor gives the energy estimate AE15. Combining it with the uniform dimension336 scalar-theta bound of DO114 SP3--12 removes the complete interval \(142\le k\le159\) at absolute normalized cost less than \(2^{-585}\). The separate139..141 charge is retained; it is not replaced or repeated. Coordinates160..164 and a full support27 conclusion are not claimed.

Put \[
 \nu=a-\tfrac32=\tfrac{333}{2},\quad
 \vartheta_a=\tfrac12\int_{-1}^{1}(1-t^2)^{a-2}\,dt
 =\frac{4^{a-1}}{2(a-1)\binom{2a-2}{a-1}},
\] \[
 U=77Q^2-14,\quad Q=1427908193729121522601,\quad
 R_0=17424/133,\quad M^2=R_0/U,\quad
 D_0=7^{140}29^{144}5^{84}13^{84},\quad D_F=2^fD_0.
 \tag{AE2}
\] The radius R used temporarily in the kernel proof below is different from the fixed energy bound \(R_0\).

\subsection{2. Retaining the full Cartesian density}\label{retaining-the-full-cartesian-density}

DO080 RT6--14 derives the classical fidelity upper bound from the exact finite trace on the actual selected compact module. Its normalized radial probability density is p(y), and its conditional radial transition is \(K_r(y,y')\). For \(r>0\), write \[
 c=\cosh r,\quad s=\sinh r,\quad R=\sqrt{1-y},\quad D=\sqrt{y/y'}.
\] The first complex sphere coordinate is \(z_1=u+it\), and \[
 f_y(u,t)=\frac{a-1}{\pi R^2}
 \left(1-\frac{u^2+t^2}{R^2}\right)^{a-2}
 {\bf1}_{u^2+t^2<R^2}.
 \tag{AE3}
\] This is the actual uniform-sphere marginal because the selected density is K-invariant, with the full vector-valued metric already included. The original radial change and its real Jacobian are \[
 y'=\frac{y}{(c+su)^2+s^2t^2},\qquad
 u=\frac{\sqrt{D^2-s^2t^2}-c}{s},\qquad
 \left|\frac{du}{dy'}\right|\le\frac{\sqrt y}{2s y'^2}.
 \tag{AE4}
\] The positive square root is forced by \(c+su\ge c-sR>0\). The Jacobian inequality follows from \((c-sR)^2-y=(s-cR)^2\ge0\), exactly as in DO080 RT10--11. Infeasible pairs are extended by zero.

Suppose first that c\textgreater D. Then \(u\le(D-c)/s<0\). Put \(K=1-(c-D)^2/(s^2R^2)\). A feasible pair requires K\textgreater0; otherwise the conditional density is zero almost everywhere. The entire density power in AE3 is bounded by \((K-t^2/R^2)_+^{a-2}\), with its support retained. Thus \[
 \int (K-t^2/R^2)_+^{a-2}\,dt
 =R K^{a-3/2}\int_{-1}^1(1-v^2)^{a-2}\,dv.
 \tag{AE5}
\] The change is \(t=R\sqrt K\,v\). At a=2 interpret the zero power with its indicator; AE5 is then the interval-length identity.

Since \(R^2\le1\), \[
 K\le1-\frac{(c-D)^2}{s^2}
 =\frac{2cD-D^2-1}{c^2-1}
 \le\frac{2D}{c+1}.
 \tag{AE6}
\] After putting the terms over the denominator \(c^2-1\), the final difference is exactly \((D-1)^2/(c^2-1)\). If instead c\textless=D, retain the unchanged DO092 beta integral, and observe that \(2D/(c+1)\ge1\). Consequently both cases prove the uniform directional estimate \[
 K_r(y,y')\le
 \vartheta_a\frac{a-1}{\pi s}\frac{\sqrt y}{R y'^2}
 \left(\frac{2\sqrt{y/y'}}{c+1}\right)^\nu .
 \tag{AE7}
\] No band in y,y', Gaussian approximation to the selected density, or discarded probability tail occurs in this estimate.

\subsection{3. Detailed balance and the genuine selected trace}\label{detailed-balance-and-the-genuine-selected-trace}

The full invariant ball measure gives DO080 RT13: \(w(y)K_r(y,y')=w(y')K_r(y',y)\) almost everywhere. The inverse Cartan element has the same radial transition kernel, since it is K-conjugate to the original one. Applying AE7 also in the reverse direction replaces D by \(D^{-1}\). The geometric mean cancels these two factors exactly, yielding \[
 \sqrt{K_r(y,y')K_r(y',y)}
 \le \vartheta_a\frac{a-1}{\pi\sinh r}
 (yy')^{-3/4}((1-y)(1-y'))^{-1/4}
 \left(\frac2{\cosh r+1}\right)^\nu .
 \tag{AE8}
\] This cancellation concerns a nonnegative upper kernel, not the signed selected coefficient. Both actual compact factors remain in the latter.

The exact change \(S=2\sqrt{Y^{-1}-1}\), together with K107 SV8--12's tail-inclusive spread estimate, gives \[
 |c_W(g)|\le
 \min\left\{1,\frac{\vartheta_a\,2^{-933/2}}{\sinh r}
                  \left(\frac2{\cosh r+1}\right)^\nu\right\},
 \qquad r>0 .
 \tag{AE9}
\] At r=0 use \(|c_W(g)|\le1\). The full lower tail in K107 was absorbed into the coefficient before this step; no additive error is now left over. Equivalently the extra multiplier is \[
                  \operatorname{sech}^{\,2a-3}(r/2).
 \tag{AE10}
\] AE9 is a normalized selected-trace bound, not an operator estimate and not a substitution of a scalar lowest-type matrix coefficient.

\subsection{4. Exact near-point transport and the original energy}\label{exact-near-point-transport-and-the-original-energy}

For the nonidentity globally rank-two-unipotent matrix put \(\alpha=(N^2)_{nn}>0\). The original normal form has \[
 \alpha_\tau=U\,r_\tau J_Z^{-1}r_\tau^*,\qquad
 0<\alpha_\tau\le R_0<132 .
 \tag{AE11}
\] At the actual near point, DO080 section7 and LIFT41 NP12 give \[
 \alpha_{p,\tau}\ge e^{-2r_{p,\tau}}\alpha_\tau,\quad
 \sinh^2r(h_{p,\tau})\ge\alpha_{p,\tau},\quad r_{p,\tau}<2^{-352}.
\] Define \[
 \Psi(x)=\left(\frac2{\sqrt{1+x}+1}\right)^\nu,\qquad x\ge0.
 \tag{AE12}
\] For \(t=\sqrt{1+x}\), one has \(x\,d\log((\sqrt{1+x}+1)/2)/dx=x/[2t(t+1)]=(t-1)/(2t)\in[0,1/2]\). Thus \(\Psi(e^{-2r_p}\alpha)\le e^{\nu r_p}\Psi(\alpha)\). Combining this with the inverse-sinh factor gives only \[
 \exp\left((\nu+1)\sum_\tau r_{p,\tau}\right)<2,
\] because \((\nu+1)f\,2^{-352}<1/2\). This replaces, rather than multiplies, the old global near-point factor2. Consequently \[
 |c(h)|<
 2\vartheta_a^f2^{-933f/2}N_F(\alpha)^{-1/2}
                   \prod_\tau\Psi(\alpha_\tau).
 \tag{AE13}
\]

Let \(g(x)=\log((\sqrt{1+x}+1)/2)\). It is concave, nonnegative and has g(0)=0; its derivative \(1/[2\sqrt{1+x}(1+\sqrt{1+x})]\) decreases. For \(0\le x\le132\), \[
 g(x)\ge \frac{g(132)}{132}x,\qquad
 g(132)>\log6>\frac74,\qquad \frac{\nu}{132}\frac74>2 .
 \tag{AE14}
\] Indeed \(\sqrt{133}>11\), while \(e<11/4\) and \((11/4)^7<6^4\) prove the logarithm comparison by exact rational arithmetic. Therefore \(\Psi(x)\le e^{-2x}\) throughout the actual range.

The original nonzero S coordinates belong to \(L=JI_\Sigma D_\eta\), \(N_EL=z^{2f}\). K106 NT7--10 gives \(N_F(\alpha)\ge U^f(kz^2)^f\). Applying that bound in AE13 and then AE14 produces the whole energy-sensitive estimate \[
 |c(h)|<
 2\vartheta_a^f k^{-84}2^{-933f/2}U^{-f/2}z^{-f}
                            e^{-2U\|r\|^2},
 \qquad \|r\|^2=\sum_\tau r_\tau J_Z^{-1}r_\tau^* .
 \tag{AE15}
\] This does not replace a field norm by an arithmetic mean. The norm inequality and the exact sum of energies are two separate valid consequences of the same original alpha.

\subsection{5. Weighted Gaussian sum in the actual coupled lattice}\label{weighted-gaussian-sum-in-the-actual-coupled-lattice}

For d=k+3, retain the saturated original transverse basis from K107 SV17--19. Its Gram \(G'\) has determinant U and \(G'-I>0\). The true primal metric is \(\operatorname{diag}(I_k,(G')^{-1})\), with covolume \[
                  D_0^d z^{2df}/U^f .
 \tag{AE16}
\] Both CM components and all denominator primes remain in the original ideal L; no moving rational basis is declared integral.

The uniformly valid DO114 SP3--12 scalar estimate is \[
 \Theta_{L^\#}(\pi^2z^2/A)<3/2,\qquad 0<A\le A_1=1/43500.
 \tag{AE17}
\] It follows from its full dimension336 spherical-cap bound and nonzero-shell integral, not from a one-dimensional minimum estimate. The true dual Gram dominates the product metric; hence its complete theta loss is \((3/2)^d\).

Write \(q=M^2/z^2\) and choose \[
                  A(q)=\min(A_1,2R_0/q).
 \tag{AE18}
\] This satisfies \(A(q)\le2Uz^2\). For every nonnegative row energy E, \[
                         e^{-2UE}\le e^{-A(q)E/z^2}.
 \tag{AE19}
\] Apply AE15 to the actual nonzero k-sector first, and then enlarge its weighted row sum to the full coupled lattice. Poisson summation gives \[
 \sum_{\rm actual\ rows}e^{-2U\|r\|^2}
 \le {164\choose k}(3/2)^d\frac{U^f}{D_0^d}
                            \left(\frac{\pi}{A(q)}\right)^{df}.
 \tag{AE20}
\] There is no positive Laplace-energy exponential. The added zero or out-of-ball tuples are used only in this nonnegative weighted sum; no selected-trace statement is assigned to them. Every complex phase is counted by this lattice sum, with no additional norm-one or pivot phase enumeration. Eta, the row and central scalar still reconstruct h.

\subsection{6. All original finite and central factors; both monotonicities}\label{all-original-finite-and-central-factors-both-monotonicities}

Retain the exact common-C G ideal and central rank-zero correlation: \[
 N_EJ=z^{2f}T^{52}/x,\quad
 |a_G(h)|\le2Q^{-335\cdot167}(N_EJ)^{335/2},\quad
 N_EJ\le Q^{167}N_FJ_0 .
\] For the actual central affine coset, retain \[
 \#d_{\rm cent}\le
 \frac{V_F(H/2+750Uz^2)U^{-f}T^{52}x^{-1}}{N_FJ_0},
 \quad V_F(C)=\frac{((22/7)fC^2)^{f/2}}{(f/2)!\sqrt{D_F}}.
 \tag{AE21}
\] These are K103 SC26.12--16 and DO098 CF10. The full-real cover750, fixed real-part shift and original U denominator are unchanged. The difference module is not confused with the affine scalar itself. All26 active weighted-label and142 inactive pointwise factors are \((T-1)^{-168}<2T^{-168}\). Both mixed projectors and the original source, phase, mass and Haar normalizations remain exactly as in K107.

Combining AE15/20/21 leaves, apart from fixed positive factors, \[
 x^{-337/2}(N_FJ_0)^{-1}F_k(q),\qquad
 F_k(q)=q^{-167f}\left(\frac{\pi}{A(q)}\right)^{df}
                              V_F(H/2+750R_0/q).
 \tag{AE22}
\] Let \(\epsilon(q)=(750R_0/q)/(H/2+750R_0/q)\in(0,1)\). The normalized logarithmic derivatives of F in its two branches are \[
 -167-\epsilon(q)\quad(q<2R_0/A_1),\qquad
 -167+d-\epsilon(q)\quad(q>2R_0/A_1).
 \tag{AE23}
\] The function is continuous at the switch. Its one-sided derivatives need not agree, but both satisfy the stated bounds.

For fixed eta and x, the all-rank-zero endpoint is \[
 z_*=(Q^{334}T^{-52}x)^{1/(2f)},\quad q_*=M^2/z_*^2,\quad
 (N_FJ_0)^{-1}\le Q^{-167}(q/q_*)^f .
\] For \(142\le k\le159\), the complete enlarged q factor decreases at least with exponent \(-(163-k)f\le-4f\). Summing the original local J states2,1,0 with multiplicities1,2,1 gives \[
 \sum_{\rm patterns}\frac{F_k(q)}{N_FJ_0}
 \le Q^{-167}F_k(q_*)\prod_G(1+q_G^{-(163-k)})^2
 <2Q^{-167}F_k(q_*).
 \tag{AE24}
\] The last comparison uses \((1+g_0^{-4})^{668}<2\). Denominator overlap with G causes no change: D\_eta is already in L, and J,J0 retain their original finite-rank definitions. This is a uniform ideal-norm bound, not an assumed homothety.

The recovered \(Q^{-167}\) is independent of x, while q\_* is proportional to \(x^{-1/f}\). The full logarithmic x derivatives after AE24 are \[
 -3/2+\epsilon(q_*)<-1/2,\qquad
 -3/2-d+\epsilon(q_*)<-1/2-d.
 \tag{AE25}
\] Thus the entire bound, including central radius and every pattern, is largest at \(x=2^{X_0}\). No new eta enumeration or Gaussian-branch union is introduced.

\subsection{7. Exact endpoints and the complete fresh union}\label{exact-endpoints-and-the-complete-fresh-union}

The same original integers as K107 give \[
 q_*^f=\frac{R_0^fT^{52}}{U^fQ^{334}2^{X_0}},\qquad
 c_-=4119905<q_*<4119906=c_+,\quad A_1c_+<2R_0 .
 \tag{AE26}
\] Hence the endpoint uses the constant-A branch, and \(H/2+750R_0/c_-<68\), \(V_F(68)<2^{479}\). Retain support \({168\choose26}<2^{101}\), the harmless original pivot \(164<2^8\), one Green2, one weighted-pattern2, and the combined T correction2. Their discrete exponent is112. The complete row-position count remains inside the next rational factor.

Use the unchanged directed constants \[
 L_T=1019333/1024,\quad L_Q=287843/4096,\quad
 L_U=601356/4096,\quad L_R=28810/4096,
\] upper for T,R0 and lower for Q,U. Regroup all signed powers first. As a transparent comparison with the old determinant and selected coefficient, \[
 C(X_0)=112+2f-929f/2+479+167fL_R-(335f/2)L_U
       +8594L_T-336\cdot167L_Q-(337/2)X_0
       =211125067/512 .
 \tag{AE27}
\] Here \(8594=335\cdot26+52-168\). The old determinant2f is removed once by336, and the selected coefficient is replaced once by335; the latter already contains the complete near-point factor2 including the new angular term. The exact radial bound is \(\vartheta_a^f<2^{-649}\).

For d=k+3 let \(e_k\) be the least strict integer binary upper exponent of \[
 W_k=\frac{{164\choose k}3^d(22\cdot43500)^{df}}
           {2^d7^{df}D_0^d c_-^{167f}},
 \qquad g_k=\lfloor\log_2(k^{84})\rfloor .
 \tag{AE28}
\] All k in the claimed interval are not powers of2, so \(k^{84}>2^{g_k}\). The complete exponent is \[
                         E_k=C(X_0)+e_k-336-335-649-g_k .
 \tag{AE29}
\] There is no \(\log_2e\) term because AE20 has no positive energy exponential. Exact rational arithmetic gives the following complete cells:

\begin{longtable}[]{@{}lrrr@{}}
\toprule\noalign{}
k & e\_k & g\_k & E\_k \\
\midrule\noalign{}
\endhead
\bottomrule\noalign{}
\endlastfoot
142 & −432508 & 600 & −11302069/512 \\
143 & −431242 & 601 & −10654389/512 \\
144 & −429976 & 602 & −10006709/512 \\
145 & −428711 & 603 & −9359541/512 \\
146 & −427446 & 603 & −8711861/512 \\
147 & −426180 & 604 & −8064181/512 \\
148 & −424915 & 605 & −7417013/512 \\
149 & −423650 & 606 & −6769845/512 \\
150 & −422385 & 607 & −6122677/512 \\
151 & −421120 & 608 & −5475509/512 \\
152 & −419855 & 608 & −4827829/512 \\
153 & −418591 & 609 & −4181173/512 \\
154 & −417326 & 610 & −3534005/512 \\
155 & −416062 & 611 & −2887349/512 \\
156 & −414797 & 611 & −2239669/512 \\
157 & −413533 & 612 & −1593013/512 \\
158 & −412269 & 613 & −946357/512 \\
159 & −411006 & 614 & −300213/512 \\
\end{longtable}

The largest is \(E_{159}<-586\); all preceding17 are below−1848. Thus \[
 \sum_{k=142}^{159}|\text{complete k-sector}|
 <2^{-586}+17\,2^{-1848}<2^{-585}.
 \tag{AE30}
\] The adjacent sufficient-bound diagnostic is \(E_{160}=345931/512>0\). It is not a lower bound or an existence claim. No cell beyond159 is removed by this theorem.

The frozen base is canonical K107114/31. The script accepts a complete later canonical ordered extension and appends585 once. At the planned actual DO114115/31 predecessor this gives116/31. It retains354 for the earlier139..141 interval,246, and every other old entry, including all overlaps. Both complete error sums remain below \(2^{-230}\), and the complete support floors remain26/11.

\subsection{8. Source and verification scope}\label{source-and-verification-scope-2}

The new analytic argument is AE3--15. The direct source is the complete genuine Hilbert-density/coarea/detailed-balance proof DO080 RT6--19 and the exact conditional beta integral DO092 SN16--18. K107 SV4--16, together with its mandatory boundary clarification, supplies the uniformly absorbed lower tail and actual near-point geometry. Its SV17--19 and K103 SC26.5--20 supply the original integral lattice and full arithmetic factors. DO114 SP3--12 supplies the new uniformly valid scalar theta lemma; its separate139..141 charge is not a substitute for this proof's count.

The primary selected Hilbert model remains Nakahama, R. (2016). Norm computation and analytic continuation of vector valued holomorphic discrete series representations. Journal of Lie Theory,26(4),927--990, equation5.7 and Theorem5.1, printedpp960--961, through the actual K043/K046 dictionary. The Gaussian Fourier convention remains Miller, S. D., and Stephens-Davidowitz, N. (2019), Kissing numbers and transference theorems from generalized tail bounds, arXiv:1802.05708v2, equations1.1/2.1. The dimension-explicit packing input is wholly proved in DO114, with Rankin, R. A. (1955), The closest packing of spherical caps in n dimensions, Proceedings of the Glasgow Mathematical Association,2(3),139--144, as its classical reference. No stronger packing theorem is assumed.

The certificate corroborates the universal angular identities, full beta scaling, actual near-point loss, energy constant, exact scalar theta inequality, original integral Gram, both one-sided monotonicities, every complete endpoint and both ordered arrays. Finite models do not replace the proofs of the kernel bound, local modules or continuum derivative inequalities.

All original fields, involution, projective modules, polarization, theta line and selected complex realization remain fixed. Exact inherited context is SGA3 I4.7.3/4.7.3.1; SGA6 VI Theorem3.3; FGA232 Theorem3.1 pp148--149; FGA236 Corollaries1.7/2.3 pp229/232; SGA7 I VII2.7/2.9 pp169--174 and VIII7.4(b) pp310--311; cote4 scan003 sections8.1.1--8.1.2, scans005--006 sections8.4.1--8.4.3 and cote54 scans005--007. These are inherited contextual anchors, not new primary readings or comparison claims. Coordinates160..164, full signed positivity, actual selected arithmetic content, prescribed global period, effective source/auxiliary-scale construction, Read/specific032, actual DP44 and RH remain separate obligations.

\subsection{Canonical algebraic-geometric context}\label{canonical-algebraic-geometric-context-40}

The original paired invertible modules, Hermitian metric, saturated transverse basis, theta line and actual selected complex realization remain fixed. The new result is an analytic density and energy estimate on that realization, not a new algebraic correspondence. Nakahama(2016), Journal of Lie Theory26(4),927-990, equation5.7 and Theorem5.1, printed960-961, supplies the carried actual invariant selected norm; the named passage was directly reread. The finite cap proof of DO114 is retained with its Rankin(1955), Theorem1(iv), pp139,142 context. Exact contextual anchors remain SGA3 I4.7.3/4.7.3.1; SGA6 VI Theorem3.3; FGA232 Theorem3.1pp148-149; FGA236 Corollaries1.7/2.3pp229/232; SGA7 I VII2.7/2.9pp169-174 and VIII7.4(b)pp310-311; cote4 scan003 sections8.1.1-8.1.2/scans005-006 sections8.4.1-8.4.3 and cote54 scans005-007. No new crystal, lambda-ring/K-theory, l-adic weight, CM-motive or global period comparison follows. DP\_MULT Q3e/DP44/DP52/DP54-65 and THETA3 G1/A7 gain only this finite gate reduction; CX\_GRH CX-ALG-024 and GCT\_PLETH T-052 retain their global gates.

\clearpage

\section{AM-R-093 --- One determinant character controls the prescribed d=7 elliptic torsion}\label{am-r-093-one-determinant-character-controls-the-prescribed-d7-elliptic-torsion}

\textbf{Canonical checked integration --- 2026-11-23T06:31:34.501119+00:00.} Complete saved-proof and executable reviews and separate exact diagnostic replays passed. No numerical error charge is introduced. Frozen author history is preserved. AM-R-093 DC1-19, on the actual d=7 branch, proves one finite-prime-unramified quadratic character epsilon with N\_E7/k(rho\_J)=epsilon\emph{chi\_cyc}rho\_C\_plus at every Tate prime. The integral rank-three norm at all primes, including2/7, gives K(J{[}n{]},C0{[}n{]}) contained in K(J{[}n{]})K\_epsilon, with relative degree at most2, simultaneously for every n.~Same actual projective/dual seeds and source congruences then give {[}Kgood:B{]}\textless=16128 and {[}K\_Sigma,out:B{]}\textless=32256\emph{M\^{}3}L\_N\^{}(129/16). The norm/type computation and arithmetic Artin convention are explicit; epsilon is neither evaluated nor assumed trivial. Root-discriminant bounds remain R092/LIFT092 on the same fields. Absolute coefficient-field degree/discriminant, effectiveN, Host/Sizes, selected arithmetic nonzero content, period, Read/specific032, evaluatedDP44 and RH remain open. No finite-error charge;116/31 and floors26/11 unchanged.

Author: /root, GPT-6 Astra Pro; effort not externally exposed; run ALP\_MACHINE. 23 November 2026 UTC. Draft until the adjacent whole-proof, executable and source verification records are completed. No canonical gate or field has been changed.

\subsection{1. The same actual fields and the theorem}\label{the-same-actual-fields-and-the-theorem}

Specialize the existing LIFT087--090 construction to the actual Host parameter \(d=7\). Put \(k=\mathbb Q(\sqrt{-7})\), \(E_7=\mathbb Q(\zeta_7)\), \(J=\operatorname{Jac}(y^2=x^7-1)\), and retain the prescribed elliptic curve \(C_0\) of LIFT088 SF5 and the same field \(K=K_{\rm good}\) of SF8--9. Thus \(E_7\subset K\), both varieties have everywhere good reduction over K, and their full maximal-order actions and canonical principal polarizations are defined over K. The coefficient field B, ideal class, polarizations, source level, lines, frames and source hypotheses are unchanged. Write \[
 T_n=K(J[n],C_0[n]),\qquad n\ge1.                         \tag{DC1}
\] All the exact seeds of LIFT087 and their action-twisted stage products have n-torsion over this same T\_n, by the integral split maps of LIFT090 SG7--17.

There is one continuous quadratic character, possibly trivial, \[
 \epsilon:G_K\longrightarrow\{1,-1\},                   \tag{DC2}
\] unramified at every finite prime, such that the full adelic characters satisfy \[
 N_{E_7/k}(\rho_J)=\epsilon\,\chi_{\rm cyc}\,\rho_C^+.
                                                               \tag{DC3}
\] Here \(\rho_C^+\) is the character of the existing C0 action oriented so that its Lie embedding of k is the specified inclusion in K. It is either the existing character or its coefficient conjugate. This notation does not replace C0, twist it, or change its prescribed action. If \(K_\epsilon\) is the fixed field of the kernel of epsilon, then, simultaneously for every n, \[
 T_n\subset K(J[n])K_\epsilon,\qquad
 [T_n:K(J[n])]\le2.                                      \tag{DC4}
\] For n=1 the extension is trivial, and for n=2 the right-hand extra extension is unnecessary. Only one K\_epsilon is used for all stages and all n.

For the literal level M and the unevaluated scale \(N=L_N=\operatorname{lcm}_j N_j\) of LIFT085/LIFT090, the same common output field satisfies the sufficient bound \[
 [K_{\Sigma,\rm out}:B]\le
 32256\,M^3N^8
 \prod_{\substack{p\mid N\\p\equiv6\pmod7}}(1+1/p)^2
 \le32256\,M^3N^{129/16}.                                \tag{DC5}
\] The single-stage statement is obtained by N=N\_1. In particular the same seed field costs at most \(32256M^3\) over B. DC5 does not bound N, the absolute degree or discriminant of B, or any DP44 term.

\subsection{2. The exact reflex-norm identity}\label{the-exact-reflex-norm-identity}

Identify \(G=\operatorname{Gal}(E_7/\mathbb Q)\) with \((\mathbb Z/7)^\times\), using \(a(\zeta_7)=\zeta_7^a\). The retained Jacobian calculation LIFT083 JS9--12 and R086 §3 gives its actual CM type \(\Phi=\{1,2,3\}\) for the specified action. Its stabilizer under G is trivial: multiplying this three-element set by each of 2,3,4,5,6 changes it. Its reflex field is therefore E7 and its reflex type is \(\Phi^{-1}=\{1,4,5\}\). These are precisely the inverse-type and norm conventions of Lombardo (2017, Definitions 2.7--2.9, p.476). Let \(H=\{1,2,4\}=\operatorname{Gal}(E_7/k)\). The multiplicities of the products \(h r\), for \(h\in H,r\in\Phi^{-1}\), are

\begin{longtable}[]{@{}lrrrrrr@{}}
\toprule\noalign{}
automorphism a & 1 & 2 & 3 & 4 & 5 & 6 \\
\midrule\noalign{}
\endhead
\bottomrule\noalign{}
\endlastfoot
multiplicity & 2 & 2 & 1 & 2 & 1 & 1 \\
\end{longtable}

The three products in each row h are (1,4,5), (2,1,3), and (4,2,6). Consequently, on the restriction-of-scalars tori, and hence on ideles, \[
 N_{E_7/k}\circ r_J=N_{K/\mathbb Q}\,r_C,
 \quad r_J=N_{\Phi}\circ N_{K/E_7},\quad r_C=N_{K/k}.
                                                               \tag{DC6}
\] Indeed each embedding of E7 in H occurs twice and every other embedding once; the first factor on the right contributes one of every embedding and the second contributes one of H. Applying N\_(K/E7) proves the statement for arbitrary K containing E7; K need not be CM or Galois over Q. The identity holds at both embeddings of k, since conjugation interchanges H and its complement. This is a calculation of the actual norm maps, with no assumption on the size of a Mumford--Tate or Galois image.

\subsection{3. A finite-order quotient of the two CM characters}\label{a-finite-order-quotient-of-the-two-cm-characters}

Use arithmetic Artin reciprocity and the covariant Tate representations. Conrad's Theorem 3.2 and Remark 3.3 (printed pp.8--9) give continuous characters \(\lambda_J:\mathbb A_K^\times\to E_7^\times\) and \(\lambda_C:\mathbb A_K^\times\to k^\times\), with their targets discrete, trivial on the infinite factors, satisfying \[
 \rho_J(\operatorname{Art}_K(s))=\lambda_J(s)r_J(s_f)^{-1},
 \quad \rho_C^+(\operatorname{Art}_K(s))=\lambda_C(s)r_C(s_f)^{-1}.
                                                               \tag{DC7}
\] For a diagonal \(c\in K^\times\), the same result gives \(\lambda_J(c)=r_J(c)\), \(\lambda_C(c)=r_C(c)\). Let \(\nu(s)\in\mathbb Q_{>0}^\times\) be the positive absolute norm of the fractional ideal determined by s\_f.~Define \[
 e(s)=\frac{N_{E_7/k}\lambda_J(s)}{\nu(s)\lambda_C(s)}\in k^\times.
                                                               \tag{DC8}
\] This is a continuous homomorphism with discrete target. Since K contains E7, it is totally imaginary. Thus for every nonzero c in K, \(\nu(c)=|N_{K/\mathbb Q}c|=N_{K/\mathbb Q}c\). DC6--7 imply e(c)=1. Both lambda characters are trivial on every finite local unit group: this is exactly Conrad's good-reduction criterion, Theorem 3.5(1), printed p.10, applied to the actual everywhere-good models. The same is true of nu. All three characters are trivial on the infinite factors as well. It follows that e factors through \[
 \mathbb A_K^\times/
 \bigl(K^\times K_\infty^\times\prod_{v<\infty}\mathcal O_{K_v}^\times\bigr)
 \simeq\operatorname{Cl}(K).                              \tag{DC9}
\] This is the modulus-one case of Milne, CFT, Proposition V.4.6; its finiteness is the modulus-one case of Theorem V.1.7 (pp.150--151,173). Therefore the image of e is a finite subgroup of k\^{}times. The only roots of unity in Q(sqrt(-7)) are ±1: a root of unity in a quadratic field has cyclotomic degree at most two, so its order belongs to \{1,2,3,4,6\}; the nonreal cases generate Q(sqrt(-3)) or Q(i). Hence e has values in \{±1\}. Global class field theory now identifies e with the unique finite character epsilon in DC2; its kernel contains every local unit group, so it is unramified at every finite prime. This step does not assume that Cl(K) is trivial or that either algebraic Hecke character is determined by its infinity type alone.

For completeness, in this arithmetic Artin convention the cyclotomic character is \[
 \chi_{\rm cyc}(\operatorname{Art}_K(s))=
          \nu(s)N_{K/\mathbb Q}(s_f)^{-1}.                 \tag{DC10}
\] Both sides lie in the integral finite rational idelic units. This is the norm-compatible arithmetic Artin formula over Q, with a positive rational generator of the finite ideal factored out; it is the normalization used in Conrad's proof immediately before Theorem 3.2, printed p.8. In particular a local uniformizer at a prime of residue cardinality q acts by q on prime-to-q roots of unity. Taking norms in DC7, then using DC6 and DC8--10, proves DC3 on the Artin image. Continuity extends the identity to G\_K.

\subsection{4. Integral norm and torsion at every prime}\label{integral-norm-and-torsion-at-every-prime}

The norm in DC3 preserves integral congruences, including at 2 and 7. Here is an explicit integral verification, avoiding a rational eigenbasis. Put \(\omega=(1+\sqrt{-7})/2\), so \(\omega^2-\omega+2=0\). The Gaussian period \(\zeta_7+\zeta_7^2+\zeta_7^4\) is \(\omega-1\), with the choice of square root matched to H. Consequently the minimal polynomial over k is \[
 P(X)=X^3-(\omega-1)X^2-\omega X-1.                      \tag{DC11}
\] Its conjugate product is \(\Phi_7(X)=1+X+\cdots+X^6\). Since \(\mathcal O_{E_7}=\mathbb Z[\zeta_7]\), one has \(\mathcal O_{E_7}=\mathcal O_k[\zeta_7]\), freely of rank three on 1,zeta7,zeta7². Multiplication by zeta7 has matrix \[
 Z=\begin{pmatrix}0&0&1\\1&0&\omega\\0&1&\omega-1\end{pmatrix}.
                                                               \tag{DC12}
\] Thus N\_(E7/k) is the determinant of multiplication in this integral free module. If a is congruent to 1 modulo n O\_E7, its multiplication matrix is I+nB, so its norm is 1 modulo n O\_k. This remains valid after completing at any rational prime, even where the finite quotient is nonreduced.

By LIFT090 SG16--17, the actual maximal-order Tate lattices of J and C0 are locally free of rank one. A Galois element acting trivially on J{[}n{]} therefore has \(\rho_J\equiv1\pmod n\) in the full integral order. The principal Weil pairing of J is perfect and defined over K, so the same element has \(\chi_{\rm cyc}\equiv1\pmod n\). Reducing DC3 by the integral norm just proved gives \(\rho_C^+\equiv\epsilon^{-1}\pmod n\). In particular an element fixing J{[}n{]} and K\_epsilon fixes C0{[}n{]}. Coefficient conjugation preserves the congruence ideal, so the same conclusion holds with the prescribed, unoriented C0 action. This proves DC4 by fixed fields. At n=2 both signs are one, proving the stated improvement. The argument applies to the same epsilon for every n; it does not introduce one quadratic extension at each level or each stage.

\subsection{5. Exact Jacobian multiplier counts}\label{exact-jacobian-multiplier-counts}

Write c for complex conjugation on E7. Its principal pairing and Rosati action imply \(a c(a)=\chi\) for the scalar a giving the Galois action on J{[}p\^{}t{]}, with \(\chi\in(\mathbb Z/p^t)^\times\). This follows directly from Galois equivariance and perfection of the pairing: for all x,y, e(ax,ay)=e(x,c(a)ay)=e(x,chi y). The rank-one full-order action is faithful. No equality of the actual image with all such scalars is assumed.

For p different from 7 the order of p modulo 7 is 1,2,3 or6, and the kernel of a -\textgreater{} a c(a) on \((\mathcal O_{E_7}/p^t)^\times\) has size \[
 K_J(p^t)=p^{3t}F_J(p),\quad
 F_J(p)=\begin{cases}
 (1-p^{-1})^3,&p\equiv1\ (7),\\
 (1+p^{-1})^3,&p\equiv6\ (7),\\
 1-p^{-3},&\operatorname{ord}_7p=3,\\
 1+p^{-3},&\operatorname{ord}_7p=6.
 \end{cases}                                                \tag{DC13}
\] To prove this, factor the cyclotomic algebra into its unramified local factors. In order1 there are three pairs of degree-one factors exchanged by c; the kernel consists of three arbitrary units and their inverses. In order2 there are three unramified quadratic factors, each with its norm to Z\_p.~In order3 there are two degree-three factors exchanged by c; the kernel is one arbitrary unit. In order6 there is one unramified degree-six factor with its norm to the unramified degree-three fixed ring. The unramified quadratic norm on residue fields is surjective, with kernel q+1, and at each further p-adic step it is induced by the surjective residue trace. Indeed N(1+p\^{}j x)=1+p\^{}j Tr(x) modulo p\^{}(j+1) for j\textgreater=1; this also holds for p=2. The kernel sizes at t=1 and the successive q factors give DC13. This argument explicitly includes the dyadic order3 case.

At p=7, E7/Q7 is totally ramified of degree6 and its c-fixed field has degree3. The unit group modulo 7\^{}t in E7 has size 6\emph{7\^{}(6t-1). Its norm image consists exactly of the units of the fixed ring whose residue is a square in F7. Necessity follows by reduction, because conjugation acts trivially on that residue field. For sufficiency, such a unit has a square root in the fixed ring by odd-prime Hensel lifting, and its norm as an element of E7 is its square. The fixed-ring image thus has size 3}7\^{}(3t-1), whence \[
                    K_J(7^t)=2\,7^{3t}.                  \tag{DC14}
\] The same argument works modulo 7\^{}t by first lifting a unit of the fixed quotient to a local unit. Because zeta7 already belongs to K, the scalar cyclotomic multiplier at 7\^{}t is congruent to1 modulo7, giving at most 7\^{}(t-1) choices. Thus the allowed Jacobian scalar set at 7\^{}t has size at most \((2/7)7^{4t}<7^{4t}\).

At an unramified p, allowing every scalar multiplier gives the upper bound \(\varphi(p^t)K_J(p^t)=p^{4t}(1-p^{-1})F_J(p)\). The last factor is at most1 in orders1,3,6. For order6 this follows from \((1-x)(1+x^3)<1\) for 0\textless x\textless1. In order2 it is \((1-x)(1+x)^3=(1-x^2)(1+x)^2\le(1+x)^2\). Chinese remaindering therefore proves, for every n\textgreater=1, \[
 [K(J[n]):K]\le n^4
     \prod_{\substack{p\mid n\\p\equiv6\pmod7}}(1+1/p)^2.
                                                               \tag{DC15}
\] It is an upper bound on an image, not an independence assertion.

\subsection{6. The actual d=7 level, one quadratic factor, and the ground field}\label{the-actual-d7-level-one-quadratic-factor-and-the-ground-field}

The literal NTypes.g\_cond is g\_i=1 modulo QS\emph{q}d.~With d=7 it forces g\_i=1 modulo7. This uses DataN.lean line916, and the actual level is M=g\_0g\_1 by \nolinkurl{ExistsParameterConditions.lean} lines55--57. For a general Host with d different from7 this congruence must not be inferred. Here every prime dividing M splits in E7. Coeff.zeta\_mem already puts mu\_M in B (LIFT089 HF89.1--4). The multiplier is consequently one at level M, and the allowed Jacobian scalar set has size \[
 V_J(M)=M^3\prod_{p\mid M}(1-p^{-1})^3\le M^3.            \tag{DC16}
\] The formula covers repeated prime powers as well. All actual N\_j are congruent to1 modulo M (LIFT085 HF21), so gcd(M,N)=1. For ell=M N², the actual scalar image injects into the product of the allowed sets at M and N². DC4 uses a single epsilon at this one combined level. Thus \[
 [T_\ell:K]\le2V_J(M)N^8
       \prod_{\substack{p\mid N\\p\equiv6\pmod7}}(1+1/p)^2.
                                                               \tag{DC17}
\] Each exact stage seed and output field is a subfield of T\_ell by LIFT090 SG18--20. No factor2 is charged separately to the seed and target steps. For every exceptional prime p\textgreater=13, \((1+1/p)^{32}\le(14/13)^{32}<13\le p\); the strict integer comparison is \(14^{32}<13^{33}\). Therefore \[
 \prod_{\substack{p\mid N\\p\equiv6\pmod7}}(1+1/p)^2
       \le\operatorname{rad}(N)^{1/16}\le N^{1/16}.       \tag{DC18}
\]

The same ground field also has the d=7 bound \[
                         [K:B]\le16128.                 \tag{DC19}
\] Indeed the quadratic Minkowski bound gives an integral representative of each class of k of norm at most 2sqrt(7)/pi\textless2, so h(-7)=1. Thus the existing j\_* is already in B by LIFT088 SF4. Because B contains k, adjoining zeta7 costs at most3, rather than6; the existing radicals i, alpha\^{}7=-4 and beta²=(zeta7-1)\^{}7 cost at most2,7,2. Hence \([K_J:B]\le84\). Before good reduction, the full O\_k action on C0 already descends over B by SF6. Its 15-torsion field over K\_J embeds into \((\mathcal O_k/15)^\times\) by LIFT090 SG16--17, which requires no good reduction. The polynomial X²-X+2 is irreducible modulo3 and modulo5, so this group has order (9-1)(25-1)=192. Consequently the very same K=K\_J(C0{[}15{]}) has degree at most84*192=16128. This does not replace any of the fields or models. Combining DC17--19 proves DC5.

\subsection{7. Arithmetic and geometric scope}\label{arithmetic-and-geometric-scope}

DC3 is an equality of the actual rank-one adelic characters. Equivalently, on every rational Tate module it is the k-linear determinant of the three-dimensional k-module T(J) tensored with Q, with the elliptic character, the cyclotomic character and one finite unramified character. DC11--12 additionally give the integral congruence statement needed for actual finite étale torsion schemes. On the good models the maximal-order actions and polarizations remain morphisms of abelian schemes over O\_K. The Picard/Poincare construction, finite level and canonical cubical line are precisely those of LIFT087--090. No new motivic isomorphism is asserted from an equality of Galois characters; no theta comparison follows from it.

The inherited exact Grothendieck anchors remain FGA232 Theorem3.1 pp.148--149; FGA236 Corollaries1.7/2.3 pp.229/232; SGA7 I, VII2.7/2.9 pp.169--174 and VIII7.4(b) pp.310--311; SGA3 I4.7.3/4.7.3.1 and SGA6 VI Theorem3.3. Formal-group and Cartier-duality context is cote4 scan003 §§8.1.1--8.1.2, scans005--006 §§8.4.1--8.4.3, and cote54 scans005--007. These are carried contextual locators, not a claim of a new reading or a new integral crystal, lambda-ring or K-theory comparison.

The same root-discriminant estimates of R092 and LIFT092 remain valid: rd(K\_good)\textless2\^{}50 exp(2pi sqrt(7)) rd(B), and rd(K\_out)\textless=rd(K\_good) M N². A smaller sufficient degree does not by itself replace those conductor estimates. DP\_MULT Q3e/DP44/DP52/DP54--65 and THETA3 G1/A7 gain the d=7 same-field degree reduction only. Absolute B degree/discriminant, effective N, actual numerical Host/Sizes, selected arithmetic nonzero-rule content, prescribed period, specific032/Read/STiesR and the evaluated DP44 comparison are still open. CX\_GRH CX-ALG-024 and GCT\_PLETH T-052 retain their own global comparison gates. There is no RH claim and no finite analytical error charge.

\subsection{8. Sources and verification boundary}\label{sources-and-verification-boundary-6}

Primary sources used in their stated scopes: Conrad, B. (n.d.). \emph{Main theorem of complex multiplication}, Theorem3.2 and Remarks3.3--3.4, printed pp.8--10; Theorem3.5(1), p.10; Remark3.6, pp.10--11, and the arithmetic Artin normalization immediately before Theorem3.2, p.8. \url{https://math.stanford.edu/~conrad/papers/mainthm.pdf} . Lombardo, D. (2017). Galois representations attached to abelian varieties of CM type. \emph{Bulletin de la Société Mathématique de France, 145}(3), 469--501. Definitions2.7--2.9, p.476. \url{https://doi.org/10.24033/bsmf.2745} . Milne, J. S. (2020). \emph{Class field theory} (v4.03), TheoremV.1.7 pp.150--151, PropositionV.4.6 pp.172--173 and PropositionV.4.7 pp.173--174. \url{https://www.jmilne.org/math/CourseNotes/CFT.pdf} . The quadratic Minkowski theorem is the retained Milne ANT v3.08, Theorem4.3 p.70, proof pp.80--81, already used in LIFT088 SF7 and R092.

The companion exact diagnostic checks the finite CM-type table, polynomial and integral determinant identities, local norm-size arithmetic, the exceptional-prime exponent, the ground constants and frozen source hashes. It supplements the universal proof, and is not an exhaustive computation of a global Galois image, an evaluation of epsilon, or a proof of effective N. Whole-file and independent checks, their actual timestamps and subject hashes, are recorded separately after the saved files are read.

\subsection{Canonical algebraic-geometric context}\label{canonical-algebraic-geometric-context-41}

The actual full-order l-adic realizations of J and C0 and their principal Weil pairings yield an integral determinant-character identity in the k-component of the CM torus, up to one unramified quadratic Artin character. Coefficient conjugation only orients the existing C0 character; no curve, model, action, ideal class, seed or source coefficient field is replaced. The identity is not asserted to arise from an algebraic correspondence or isomorphism of CM motives. Conrad, Main theorem of complex multiplication, Theorem3.2, Remarks3.3-3.4 and Theorem3.5(1), printedpp8-10, with Remark3.6pp10-11; Lombardo(2017), Definitions2.7-2.9,p476; Milne CFTv4.03,V Theorem1.7pp150-151 and Propositions4.6-4.7pp172-174 are directly checked in their indicated scopes. Carried geometric anchors: FGA232 Theorem3.1pp148-149; FGA236 Corollaries1.7/2.3pp229/232; SGA7 I VII2.7/2.9pp169-174 and VIII7.4(b)pp310-311; SGA3 I4.7.3/4.7.3.1; SGA6 VI Theorem3.3; cote4 scan003 sections8.1.1-8.1.2/scans005-006 sections8.4.1-8.4.3, cote54 scans005-007. No new crystalline slope, lambda-ring/K-theory or weight theorem follows. DP\_MULT Q3e and DP44/DP52/DP54-65 gain the d7 same-field cost reduction; THETA3 G1/A7, CX\_GRH CX-ALG-024 and GCT\_PLETH T-052 retain their global comparison gates.

\clearpage

\section{AM-LIFT-093 --- Principal pairings bound the same two-generator torsion fields}\label{am-lift-093-principal-pairings-bound-the-same-two-generator-torsion-fields}

\textbf{Canonical checked integration --- 2026-11-23T06:37:31.608745+00:00.} Complete saved-proof and executable reviews and separate exact diagnostic replays passed. No numerical error charge is introduced. Frozen author history is preserved. AM-LIFT-093 PA1-27 refines the general-d degree for the exact same Kgood, J,C0 and original source seeds. Principal Weil pairings impose one shared scalar multiplier. Exact all-prime norm counts, including2/7/d and the overlappingd7 case, give {[}F(J{[}n{]},C0{[}n{]}):F{]}\textless=U(n)=n\^{}5\emph{product\_(p\textbar n,p=6 mod7)(1+1/p)\^{}3 for every finiteF/Kgood, and U(N\^{}2)\textless=N\^{}(81/8). The actual source level has multiplier1 and boundV(M)=M\^{}4}product\_(p\textbar M)F\_J(p)(1-1/p)\textless2M\^{}4; no universal g mod7 condition is inserted. The existing C0{[}15{]} extension costs at most192, so {[}Kgood:B{]}\textless=32256h(-d)\textless=32256d and the actual common output satisfies {[}Kout\_Sigma:B{]}\textless=32256d V(M)U(L\_N\^{}2)\textless64512d M\^{}4 L\_N\^{}(81/8). R093 retains its independent strongerd7 bound. Same fields, all objects and intrinsic Step5 costs remain; effectiveN, absoluteB, Host/Sizes, selected nonzero arithmetic content, period, Read/specific032, evaluatedDP44 and RH are open. No finite-error charge;116/31 and floors26/11 unchanged.

Author: /root/lift\_accounting, ALP\_MACHINE, 23 November 2026 UTC.

\textbf{Author-frozen after whole saved-proof, source-interface and executable review.} This report changes degree estimates for the actual fields and the newly chosen seeds of LIFT087--090. It changes no canonical record, polarization, integral action, level, frame, or error charge. The adjacent author verification records the final subject hashes, exact diagnostic replay and completed review time. No independent verdict is asserted here.

\subsection{1. Statement for the actual fields}\label{statement-for-the-actual-fields-1}

Retain every Host, Datum, Pins, Moduli, Coeff and seed-construction hypothesis of LIFT084--090. In particular, \(d\) is a prime with \(d\equiv7\pmod {16}\), \(k=\mathbb Q(\sqrt{-d})\), and its order is the full maximal order. Put

\[
 K_7=\mathbb Q(\zeta _7),\quad
 J=\operatorname{Jac}(y^2=x^7-1),\quad
 K=K_{\rm good},\quad T_n(F)=F(J[n],C_0[n])\quad(F\supset K).
 \tag{PA1}
\]

Here \(C_0\), its prescribed ideal class and full \(\mathcal O_k\)-action are exactly those of LIFT088. Its ground field is exactly

\[
 K_j=B(j_*),\quad K_J=K_j(\zeta _7,i,\alpha,\beta),\quad
 \alpha^7=-4,\quad\beta^2=(\zeta _7-1)^7,\quad
 K=K_J(C_0[15]).                                      \tag{PA2}
\]

Both canonical principal polarizations and both full maximal-order actions are defined over \(K\). In particular \textbf{\(\zeta _7\in K\)}. The actual level is \(M=g_0g_1\), with distinct primes, and \(A=1\); the existing Coeff.zeta\_mem gives \(\mu_M\subset B\subset K\). No new root extension is made.

For \(p\ne7\), let \(t_p=\operatorname{ord}_{7}(p)\), and define

\[
 F_J(p)=
 \begin{cases}
 (1-p^{-1})^3,&t_p=1,\\
 (1+p^{-1})^3,&t_p=2,\\
 1-p^{-3},&t_p=3,\\
 1+p^{-3},&t_p=6,
 \end{cases}
 \qquad
 \mathcal U(n)=n^5
   \prod_{\substack{p\mid n\\p\equiv6\ (7)}}(1+p^{-1})^3.
 \tag{PA3}
\]

Empty products are one. For \(n=1\), every torsion extension below is trivial. For \(n>1\), \(\mathcal U(n)\) is an integer: an exceptional prime power contributes \(p^{5a-3}(p+1)^3\).

\textbf{Theorem.} For every finite extension \(F/K\) and every \(n\ge1\),

\[
 [T_n(F):F]\le\mathcal U(n),\qquad
 [T_{N^2}(F):F]\le\mathcal U(N^2)\le N^{81/8}\le N^{11}.
 \tag{PA4}
\]

The exponents include all prime powers, the primes \(2,7,d\), and \(N=1\). No prime-to-discriminant or prime-to-polarization-degree condition is imposed. Good reduction is retained from the construction; the pairing count itself concerns characteristic-zero torsion and does not require a reduction argument.

For the source level, set

\[
 \mathcal S(M)=\prod_{p\mid M}F_J(p)(1-p^{-1}),\qquad
 \mathcal V(M)=M^4\mathcal S(M).
 \tag{PA5}
\]

Both level primes split in \(k\), are odd, and differ from \(7,d\), as verified in §6. Consequently the same actual seed field satisfies

\[
 [K_{\rm seed}:K]\le\mathcal V(M)<2M^4,
 \qquad K_{\rm seed}=K(A_0[M]).                         \tag{PA6}
\]

If neither level prime is \(6\pmod7\), then \(\mathcal S(M)<1\), so the sharper bound is \(\mathcal V(M)<M^4\). This applies in particular when both are \(1\pmod7\), as the literal congruence implies whenever \(7\mid QS\,q\,d\), including \(d=7\). We do not silently impose this additional divisibility on every Host.

The same ground field in PA2 satisfies

\[
 [K:B]\le168\cdot192\,h(-d)=32256h(-d)\le32256d.
 \tag{PA7}
\]

For a finite list of the same-seed stages of LIFT090, let \(L_N=\operatorname{lcm}_jN_j\). Under all the nonzero-rule and target-construction hypotheses of LIFT085--086, the same actual composita obey

\[
 \begin{split}
 [K_\Sigma:K]&\le\mathcal V(M),\\
 [K_{\rm out,\Sigma}:K_\Sigma]&\le\mathcal U(L_N^2)
                                  \le L_N^{81/8},\\
 [K_{\rm out,\Sigma}:B]&\le32256d\,\mathcal V(M)
                                      \mathcal U(L_N^2)
                            <64512d\,M^4L_N^{81/8}.
 \end{split}                                             \tag{PA8}
\]

These are degree estimates for the original fields, not replacements by the containing torsion fields. Their source objects, integral models and intrinsic costs are unchanged. \(N_j\), the absolute degree of the given \(B\), the exact coefficient-ring pin, the nonzero rule, selected arithmetic content, prescribed period, Read/032 and the evaluated DP44 expression remain separate obligations.

\subsection{2. The integral scalar actions and one shared multiplier}\label{the-integral-scalar-actions-and-one-shared-multiplier}

LIFT090 SG7--17 proves, for every \(n\), the actual integral direct-summand inclusions

\[
 K(A_{0,j}[n])\subset K(J[n],C_0[n]),\qquad
 J[n]\simeq\mathcal O_{K_7}/n\mathcal O_{K_7},\quad
 C_0[n]\simeq\mathcal O_k/n\mathcal O_k                  \tag{PA9}
\]

as geometric modules for their indicated orders. The rank-one identifications use the actual projective Betti lattices localized at every prime and Chinese remaindering, not global principality. The split maps and their duals are defined over \(K\), and the Zarhin target is the actual product of a projective Serre factor and its dual. Thus PA9 holds also at primes dividing a module index, and after every finite extension \(F/K\).

Every Galois automorphism over \(F\) acts by units \(a,b\) in these two rings. The Rosati involutions of the actual principal polarizations restrict to complex conjugation. For \(J\), the curve automorphism \((x,y)\mapsto(\zeta _7x,y)\) fixes infinity and preserves the effective-divisor theta locus \(W_2\subset J\). Hence it preserves the canonical polarization, its Rosati adjoint is its inverse, and this gives conjugation on \(\mathbb Z[\zeta _7]=\mathcal O_{K_7}\). For \(C_0\), the canonical elliptic polarization induces the positive involution on its imaginary quadratic endomorphism field, which is conjugation. The theta construction is Milne (2008, III §6, Theorem6.6, pp.105--108); the Rosati identity and positive-involution description are Milne (2008, I §14, p.61) and Milne (2020b, Corollary1.40 and Example2.9, pp.21,25--26).

The finite Weil pairings are perfect because the polarizations are principal. Their Galois equivariance and adjoint identity give, for a Galois element \(\sigma\),

\[
 \begin{split}
 e_J(ax,ay)&=e_J(x,\bar a a y)=e_J(x,y)^{\chi_n(\sigma)},\\
 e_C(bu,bv)&=e_C(u,\bar b b v)=e_C(u,v)^{\chi_n(\sigma)}.
 \end{split}                                             \tag{PA10}
\]

Here \(\chi_n\) is the action on the \textbf{same} group \(\mu_n\), so there is one common multiplier. Perfectness implies that \(a\bar a-\chi_n(\sigma)\) and \(b\bar b-\chi_n(\sigma)\) act as zero. Faithfulness of the rank-one modules in PA9 then proves the equalities in their finite rings,

\[
                       a\bar a=b\bar b=\chi_n(\sigma).
 \tag{PA11}
\]

Milne (2008, I §13, pp.57--58, Proposition13.2) supplies exactly the finite-pairing compatibility, including all \(n\) in characteristic zero. The choice of geometric rank-one bases changes neither scalar multiplication nor PA11.

Here is the exact local counting interface. Put \(R_J=\mathcal O_{K_7}\otimes\mathbb Z_p\), \(R_C=\mathcal O_k\otimes\mathbb Z_p\), and work modulo \(p^a\), \(a\ge1\). Let \(I_J,I_C\subset(\mathbb Z/p^a)^\times\) consist of the rational scalar units occurring as \(u\bar u\) in the respective rings. Let \(k_J,k_C\) be the orders of their norm-one kernels. If \(S\subset(\mathbb Z/p^a)^\times\) is an allowed multiplier subgroup, then the ambient pair group has exact order

\[
 \left|\{(u,v):u\bar u=v\bar v\in S\}\right|
                  =|I_J\cap I_C\cap S|\,k_J k_C.         \tag{PA12}
\]

Indeed each nonempty norm fibre is a coset of its norm-one kernel. This is an equality for the displayed finite group, followed by an upper bound for the actual Galois image. In particular, two ramified norm images do not produce two independent multiplier choices or two independent half-factors. The scalar map \(\mathbb Z/p^a\to R_J/p^aR_J\) is injective: if a rational scalar lies in \(p^aR_J\), its quotient by \(p^a\) is integral over \(\mathbb Z_p\) and belongs to \(\mathbb Q_p\), hence lies in \(\mathbb Z_p\). The same argument applies to \(R_C\).

\subsection{3. Exact norm counts away from the odd ramified primes}\label{exact-norm-counts-away-from-the-odd-ramified-primes}

For \(p\ne7\), cyclotomic factorization gives the following local algebras with involution. The decomposition degree is \(t_p\); complex conjugation is the element \(-1\) in \((\mathbb Z/7)^\times\). The four cases are:

\begin{longtable}[]{@{}
  >{\raggedleft\arraybackslash}p{(\columnwidth - 4\tabcolsep) * \real{0.3636}}
  >{\raggedright\arraybackslash}p{(\columnwidth - 4\tabcolsep) * \real{0.2727}}
  >{\raggedleft\arraybackslash}p{(\columnwidth - 4\tabcolsep) * \real{0.3636}}@{}}
\toprule\noalign{}
\begin{minipage}[b]{\linewidth}\raggedleft
\(t_p\)
\end{minipage} & \begin{minipage}[b]{\linewidth}\raggedright
Algebra with conjugation
\end{minipage} & \begin{minipage}[b]{\linewidth}\raggedleft
Norm-one kernel modulo \(p^a\)
\end{minipage} \\
\midrule\noalign{}
\endhead
\bottomrule\noalign{}
\endlastfoot
1 & Three split pairs \(\mathbb Z_p\times\mathbb Z_p\), with swaps & \(p^{3a}(1-p^{-1})^3\) \\
2 & Three unramified quadratic extensions of \(\mathbb Q_p\) & \(p^{3a}(1+p^{-1})^3\) \\
3 & Two unramified cubic factors interchanged & \(p^{3a}(1-p^{-3})\) \\
6 & An unramified quadratic extension of the unramified cubic field & \(p^{3a}(1+p^{-3})\) \\
\end{longtable}

Thus

\[
                 k_J(p^a)=p^{3a}F_J(p),\qquad
                 I_J(p^a)=(\mathbb Z/p^a)^\times
                       \quad(p\ne7).                     \tag{PA13}
\]

The cyclotomic maximal order and ramification assertions are Milne (2020a, Proposition6.2 and Theorem6.4, pp.96--99); Example8.18 (p.143) identifies the Frobenius degree with the order of \(p\bmod7\).

For clarity, all norm-surjectivity and prime-power steps used in this table can be proved directly. A split pair has norm \((x,y)\mapsto xy\), with one freely chosen unit. If the fixed unramified local field has residue field \(\mathbb F_q\), an unramified quadratic extension has residue norm \(z\mapsto z^{q+1}\), surjective onto \(\mathbb F_q^\times\), with kernel \(q+1\). On each subsequent principal-unit layer,

\[
 N(1+p^r x)\equiv1+p^r\operatorname{Tr}(x)\pmod {p^{r+1}}
                           \quad(r\ge1).                 \tag{PA14}
\]

The residue trace is surjective; choosing each new coefficient corrects the next digit of a target norm. Induction proves norm surjectivity modulo every \(p^a\). The kernel size is consequently the quotient of the unit-group orders,

\[
 \frac{q^{2(a-1)}(q^2-1)}{q^{a-1}(q-1)}
                    =q^{a-1}(q+1).
\]

This proof includes residue characteristic2. In that characteristic the finite-field trace is not the scalar map \(x\mapsto2x\): for a quadratic finite-field extension, \(x\mapsto x+x^q\) has kernel the \(q\)-element base field, so its image has \(q\) elements and is the entire base field. PA14 remains valid for \(p=2,r=1\), since the quadratic cross term is divisible by \(p^{r+1}\).

For \(p\ne d\), let \(\epsilon_p=\left(\frac{-d}{p}\right)\), with its Kronecker value at2. The quadratic split or unramified calculation gives

\[
 k_C(p^a)=p^a(1-\epsilon_p/p),\qquad
 I_C(p^a)=(\mathbb Z/p^a)^\times.                         \tag{PA15}
\]

In particular \(\epsilon_2=1\), since \(d\equiv7\pmod {16}\). Directly, the maximal-order polynomial \(X^2-X+(d+1)/4\) reduces to \(X(X-1)\) modulo2 with simple roots, so the algebra splits over \(\mathbb Z_2\). For \(J\), \(t_2=3\). No dyadic ramified norm formula is being imported.

\subsection{4. Odd ramification and the already defined seventh roots}\label{odd-ramification-and-the-already-defined-seventh-roots}

At7, \(K_7/\mathbb Q\) is totally ramified of degree6, with residue field \(\mathbb F_7\). Its fixed field \(K_7^+\) has degree3 and the relative quadratic extension is tamely ramified. The finite unit-group sizes are

\[
 |(\mathcal O_{K_7}/7^a)^\times|=6\,7^{6a-1},\qquad
 |(\mathcal O_{K_7^+}/7^a)^\times|=6\,7^{3a-1}.           \tag{PA16}
\]

The norm image consists exactly of the units with square residue in \(\mathbb F_7^\times\). Necessity holds because conjugation acts trivially on the residue field. Conversely, any unit in the fixed local field with square residue is a square there, by Hensel lifting at an odd prime. If \(v^2=u\) in the fixed field, then its relative norm is \(N(v)=v^2=u\). This proves the full image statement, also modulo \(7^a\), and hence

\[
 k_J(7^a)=2\,7^{3a},\qquad
 I_J(7^a)=\{c\in(\mathbb Z/7^a)^\times:
                           c\bmod7\text{ is a square}\}.
 \tag{PA17}
\]

Likewise the quadratic extension \(k/\mathbb Q\) is ramified only at the odd prime \(d\). Its unit-group order modulo \(d^a\) is \((d-1)d^{2a-1}\), while its scalar norm image has order \(\varphi(d^a)/2\). The same square-residue/Hensel argument gives

\[
 k_C(d^a)=2d^a,\qquad
 I_C(d^a)=\{c\in(\mathbb Z/d^a)^\times:
                           c\bmod d\text{ is a square}\}.
 \tag{PA18}
\]

These formulae concern quotienting by the rational integer \(p^a\), not by a single local uniformizer power \(a\). In PA16 the corresponding uniformizer exponents are \(6a\) and \(3a\), respectively. This distinction is essential for the powers of \(7\).

Use PA12 first with the full scalar unit group. At \(d=7\), the common image is the same square-residue subgroup of index2, so only one half-factor appears. The actual field gives a further restriction: PA2 contains \(\zeta _7\), and \(F\supset K\) fixes it. Therefore every actual cyclotomic multiplier at \(7^a\) is \(1\pmod7\). Set

\[
 S_{7^a}=\ker((\mathbb Z/7^a)^\times\to
                         (\mathbb Z/7)^\times),\qquad
 |S_{7^a}|=7^{a-1}.                                    \tag{PA19}
\]

This subgroup is contained in \(I_J\) and also in \(I_C\), whether or not \(d=7\). At every other prime we may conservatively take all scalar multipliers. Consequently PA12 yields the following exact orders for the chosen ambient local groups, normalized by \(p^{5a}\):

\begin{longtable}[]{@{}lr@{}}
\toprule\noalign{}
Prime case & Ambient order divided by \(p^{5a}\) \\
\midrule\noalign{}
\endhead
\bottomrule\noalign{}
\endlastfoot
\(p\ne7,d\) & \((1-p^{-1})(1-\epsilon_p/p)F_J(p)\) \\
\(p=d\ne7\) & \((1-p^{-1})F_J(p)\) \\
\(p=7,d\ne7\) & \(\frac{2}{7}(1-\epsilon_7/7)\) \\
\(p=d=7\) & \(4/7\) \\
\end{longtable}

The third row is at most \(16/49\), and the fourth is less than1. In the second row the index2 norm-image restriction cancels the factor2 in \(k_C\). The last row uses the shared multiplier set of size \(7^{a-1}\), not its square and not two unrelated choices.

\subsection{\texorpdfstring{5. Uniform degree and the exponent \(81/8\)}{5. Uniform degree and the exponent 81/8}}\label{uniform-degree-and-the-exponent-818}

Let \(0<x<1\). The elementary inequalities

\[
 (1-x)(1\pm x)\le1,\qquad
 (1-x^2)(1+x^3)\le1                                  \tag{PA20}
\]

show that the first two rows in the preceding table are at most1 when \(t_p=1,3,6\). For example, in the first row the worst quadratic sign is \(\epsilon_p=-1\), producing \((1-p^{-2})F_J(p)\); in the \(t_p=6\) case the second inequality in PA20 applies. For the ramified second row, \((1-x)(1+x^3)<1\) as well. If \(t_p=2\), discard the remaining factors, which are at most1, and retain \((1+p^{-1})^3\). The condition \(t_p=2\) is exactly \(p\equiv6\pmod7\).

Chinese remaindering gives an injection of the actual Galois image into the product of the local ambient groups. Their dependencies across primes can only reduce the image. Multiplying their orders proves the first assertion of PA4, including \(2,7,d\). No independence of local characters or exact Galois-image theorem is asserted.

Every exceptional prime is at least13. The exact integer inequality

\[
                            14^{24}<13^{25}              \tag{PA21}
\]

implies \((1+1/13)^{24}<13\). Since the left side \((1+1/p)^{24}\) decreases with \(p\), and the right side increases, it follows that

\[
 (1+p^{-1})^3<p^{1/8}\quad(p\ge13).
\]

The primes dividing \(N^2\) are exactly those dividing \(N\). Therefore

\[
 \mathcal U(N^2)
 \le N^{10}\operatorname{rad}(N)^{1/8}
 \le N^{81/8}\le N^{11}.                              \tag{PA22}
\]

For \(N=1\), read every inequality as equality1. The use of \(\operatorname{rad}(N)\), rather than the larger bound \(\operatorname{rad}(N^2)\le N^2\), is what gives \(81/8\).

\subsection{6. The exact source level and its separate factor}\label{the-exact-source-level-and-its-separate-factor}

The following are literal source facts, not congruences inferred from earlier concrete prime examples:

\begin{itemize}
\tightlist
\item
  DataHost.lean, Host fields g\_prime and g\_ne, gives two distinct rational primes.
\item
  DataN.lean, NTypes.g\_cond at line916, gives \(g_i\equiv1\pmod {QS\,q\,d}\).
\item
  ConjEMemHostFIff.lean, HostRaw.two\_le\_e and tower\_dvd, gives \(16\mid QS\). In particular \(g_i\equiv1\pmod4\), so \(g_i\ne2,7\).
\item
  \nolinkurl{ExistsParameterConditions.lean,} g\_mod12--21, g\_jacobi29--39 and splitG44--49, proves the congruence \(g_i\equiv1\pmod d\), hence \(g_i\ne d\), and that \(k\) splits at each \(g_i\). Its levelG\_M55--57 identifies \(M=g_0g_1,A=1\).
\item
  \nolinkurl{PartitionExtendByZeroAntitone.lean,} Par.PinLevel at line43, and PushInvertible.lean, Coeff.zeta\_mem at line45, retain exactly this level and \(\mu_M\subset B\).
\end{itemize}

At a level prime \(p\), the cyclotomic multiplier is1 because \(F\supset K\supset\mu_M\). PA12 therefore has multiplier factor1. Since \(k\) splits there, PA13 and PA15 give the exact ambient order

\[
             k_J(p)k_C(p)=p^4F_J(p)(1-p^{-1}).            \tag{PA23}
\]

Multiplying over the two distinct primes gives \(\mathcal V(M)\). Each factor \(F_J(p)(1-p^{-1})\) is less than1 for \(t_p=1,3,6\). In the remaining case it is less than \((1+p^{-1})^3\le(14/13)^3\). Thus

\[
 \mathcal S(M)\le(14/13)^6<2,\qquad
 14^6=7529536<2\cdot4826809=2\cdot13^6.                \tag{PA24}
\]

If there is no exceptional prime, the product is less than1. A general Host need not be assigned a congruence modulo7 merely from g\_mod; the exact factor PA5 is valid in every case. When \(7\mid QS\,q\,d\), the literal g\_cond does supply the sharper split-\(K_7\) branch, where \(\mathcal V(M)=(g_0-1)^4(g_1-1)^4<M^4\). For \(d=7\) this branch is automatic. No change of primes is used in either branch.

The inclusions PA9 now prove PA6 for exactly \(K(A_0[M])\). This is a smaller containing group for its existing Galois image, not a new choice of a level vector or of the prescribed finite level.

\subsection{7. Ground field, simultaneous stages, and preserved costs}\label{ground-field-simultaneous-stages-and-preserved-costs}

Before the good-model extension, the entire \(\mathcal O_k\)-action on \(C_0\) is already defined over \(K_J\), by LIFT088 SF4--8. Its integral rank-one torsion module therefore gives

\[
 [K_J(C_0[15]):K_J]\le
 |(\mathcal O_k/15)^\times|
 \le(3^2-1)(5^2-1)=192.                               \tag{PA25}
\]

Here \(d\) is neither3 nor5. At an unramified quadratic prime the unit count is \((p-1)^2\) in the split case and \(p^2-1\) in the inert case, so the displayed upper bound includes both. This is the \textbf{same} field \(K_J(C_0[15])\), and the existing finite-inertia proof of everywhere-good reduction remains unchanged. With \([K_J:B(j_*)]\le168\) and \([B(j_*):B]\le h(-d)\le d\), PA25 proves PA7.

For a finite list of stages, LIFT090's actual projective maps and conjugate-dual maps place every \(A_{0,j}[n]\) in the one \(T_n(K)\). Let

\[
 K_\Sigma=\prod_jK(A_{0,j}[M]),\qquad
 K_{\rm out,\Sigma}
       =K_\Sigma\bigl(A_{0,j}[N_j^2]\text{ for all }j\bigr).
 \tag{PA26}
\]

The product denotes a compositum in the fixed algebraic closure. Then \(K_\Sigma\subset T_M(K)\), and the second field is contained in \(T_{L_N^2}(K_\Sigma)\). Apply PA6 over \(K\) and PA4 over \(K_\Sigma\), and then PA7. This proves PA8. A product of one torsion-field bound per stage is unnecessary. No finite collection of new roots, level trivialisations or dual torsion fields is being omitted: the all-object descent is exactly LIFT088 and the same-field inclusions are exactly LIFT090.

The intrinsic comparison of LIFT086--087 depends on the same isogeny scale \(N\) and the same seed line bound \(b\). The sufficient Step5 quantity stays

\[
 \frac{D+gn}{2}+D\log r+
 (b+\log N)\min\{D+gn,(3D+gn)/2\}.                    \tag{PA27}
\]

Here the section degree \(n\) in PA27 has its earlier source meaning; it is unrelated to the torsion integer \(n\) in PA1--4. Finite field base change does not change the normalized intrinsic degree. The new result neither evaluates \(N\) nor changes this cost. Retain the minimum with LIFT089's \(M^{fr^2}\), \(N^{4fr^2}\), or LIFT090's earlier bounds whenever that is smaller.

\begin{longtable}[]{@{}
  >{\raggedright\arraybackslash}p{(\columnwidth - 4\tabcolsep) * \real{0.3333}}
  >{\raggedright\arraybackslash}p{(\columnwidth - 4\tabcolsep) * \real{0.3333}}
  >{\raggedright\arraybackslash}p{(\columnwidth - 4\tabcolsep) * \real{0.3333}}@{}}
\toprule\noalign{}
\begin{minipage}[b]{\linewidth}\raggedright
Item
\end{minipage} & \begin{minipage}[b]{\linewidth}\raggedright
Supplied by this theorem
\end{minipage} & \begin{minipage}[b]{\linewidth}\raggedright
Retained requirement
\end{minipage} \\
\midrule\noalign{}
\endhead
\bottomrule\noalign{}
\endlastfoot
Ground field \(K\) & Same field, degree at most \(32256h(-d)\) over \(B\) & Actual given \(B\), its absolute degree and coefficient-ring pin \\
Actual seed level & Exact factor \(\mathcal V(M)<2M^4\), and the stated sharper branch & Every original Pins, lattice and moduli hypothesis \\
One target or a finite list & \(\mathcal U(N^2)\le N^{81/8}\), or one \(\mathcal U(L_N^2)\) & Effective values of the isogeny scales and all LIFT085 nonzero-rule hypotheses \\
Models, line, action, frame & Same descended objects as LIFT088--090 & No new small-height frame or normalized coefficient assertion \\
Root discriminants & No change to the proved LIFT091--092 conductor bounds & Initial-field bound and the same unevaluated \(N\) \\
Main arithmetic problem & A strictly smaller conditional field-degree factor & Prescribed arithmetic image, global period, Read/032, DP44 and RH \\
\end{longtable}

No finite-kernel error charge is introduced. This is not the separate \(d=7\) determinant-character comparison pursued in R093, nor does it rely on that comparison.

\subsection{8. Primary and literal verification boundary}\label{primary-and-literal-verification-boundary}

The local finite-ring calculations above are proved explicitly, including the principal-unit lifting at2 and the uniformizer exponents at7. The companion executable checks exact small finite rings, prime-power formulae, both ramified cases and their shared multiplier images, the \(\zeta_7\) restriction, the source-level branches, exact exponent inequalities and same-field arithmetic constants. Finite diagnostics do not prove the universal theorem or calculate an actual torsion Galois group. The proof and the named primary inputs supply that argument.

The adjacent source manifest records the actual frozen bytes and precise read scopes for LIFT085/088/089/090, the literal Host/Level/Coeff/Pins fields, NTypes.g\_cond, the derived splitG theorem and the primary locators below. It distinguishes a proved scalar-action inclusion from any unproved surjectivity of the Galois image. It also records that only a finite source level descends over the given number field; the full Tate tower remains a geometric condition.

\subsubsection{References}\label{references-6}

Milne, J. S. (2008). \emph{Abelian varieties} (Version2.00). ChapterI, §§13--14, pp.57--61; ChapterIII, §6, Theorem6.6 and Summary6.11, pp.105--108. \url{https://www.jmilne.org/math/CourseNotes/AV.pdf}

Milne, J. S. (2020a). \emph{Algebraic number theory} (Version3.08). Proposition6.2, Theorem6.4, pp.96--99; Example8.18, p.143. \url{https://www.jmilne.org/math/CourseNotes/ANT.pdf}

Milne, J. S. (2020b). \emph{Complex multiplication} (Version0.10). Corollary1.40, p.21; Example2.9, pp.25--26. \url{https://www.jmilne.org/math/CourseNotes/CM.pdf}

ALP\_MACHINE. (2026). AM-LIFT-088, SF1--16; AM-LIFT-089, HF89.1--18; AM-LIFT-090, SG1--21; AM-LIFT-085/086 fixed-level and all-object descent interfaces. Their canonical frozen paths and hashes are identified in the accompanying source manifest.

\subsection{Canonical algebraic-geometric context}\label{canonical-algebraic-geometric-context-42}

The actual principal Picard/Poincare pairings and full maximal-order scalar actions constrain the l-adic and finite-etale torsion realizations of the same two CM progenitors. The new count is an image-container estimate, not an image-surjectivity or motivic-comparison theorem. Literal source precision: in the frozen PushInvertible.lean, Coeff.zeta\_mem is line44, not the isolated45 printed in the frozen author body; the fullCoeff40-46 locator was already correct. Read that occurrence with the immutable locator-precision addendum. No source or mathematical premise changes. Fresh primary anchors: Milne AVv2.00(2008),I13 Proposition13.2pp57-58,I14p61,III6.6p105 and Summary6.11pp107-108; Milne ANT v3.08(2020),Proposition6.2pp96-98,Theorem6.4p99,Example8.18p143; Milne CMv0.10(2020),Corollary1.40p21,Example2.9pp25-26. Carried geometric anchors: FGA232 Theorem3.1pp148-149; FGA236 Corollaries1.7/2.3pp229/232; SGA7 I VII2.7/2.9pp169-174 and VIII7.4(b)pp310-311; SGA3 I4.7.3/4.7.3.1; SGA6 VI Theorem3.3; cote4 scan003 sections8.1.1-8.1.2/scans005-006 sections8.4.1-8.4.3 and cote54 scans005-007. There is no new crystalline slope, lambda-ring/K-theory identity, weight or CM-motive map. DP\_MULT Q3e/DP44/DP52/DP54-65 gain a smaller same-field cost; THETA3 G1/A7, CX\_GRH CX-ALG-024 and GCT\_PLETH T-052 retain their global gates.

\clearpage

\section{AM-K-109 --- full transverse metric and the next complete energy sector}\label{am-k-109-full-transverse-metric-and-the-next-complete-energy-sector}

\textbf{Canonical checked integration --- 2026-11-23T06:40:48.560679+00:00.} Complete root and independent proof/script/interface reviews and separate exact replays passed. The actual canonical predecessor has116/31 entries; this integration gives117/31, complete floors26/11. Frozen author history is preserved. AM-K-109 TG1-19 retains the full all-A scalar theta polynomial and the original transverse Gram G\textgreater2\^{}36 I, with true determinantU. Using the entire actual K108 energy, it removes the complete k160 support26 cell with exact exponent-163509/512\textless-319. Original module membership, central affine fibre, common ideal correlation, all finite factors and both q/x monotonicities are retained. The necessary actual all-II rank-two-unipotent remainder becomes161\textless=k\textless=164 with2\^{}(71443/64)\textless x\textless2\^{}2108,v=q26,b=j0. Exactly one319 append gives117/31 complete ordered errors, both entire sums below2\^{}-230; complete floors26/11 remain. No161+, signed positivity, selected arithmetic content, effective source, period, specific032/Read, evaluatedDP44 or RH conclusion.

Author: /root/kummer\_extension. 23 November 2026 UTC. Controlled pending proof. Author whole read, exact replay, independent whole review and canonical integration are separate records. The frozen K108 and DO114 proofs are not modified.

\subsection{1. Original sector and statement}\label{original-sector-and-statement}

Retain the actual all-top J56 family, H=133 and f=a=168, every allowed word and canonical twist, the original fields and involution, projective modules, Hermitian polarization, selected vector-valued Hilbert metric, both compact factors, original paired pencils, phases, finite projectors, Haar measures and class masses.

After K108 the still possible all-II support26 globally rank-two-unipotent scalar-cross branch has \[
 v=q_{\rm elig}=26,\quad b=j=0,\quad
 2^{X_0}<x=N_ED_\eta<2^{2108},\quad
 X_0=71443/64,\quad160\le k\le164.
 \tag{TG1}
\] Here k is the number of nonzero original S covector coordinates. The normal form, unique eta for each support in this exact window and original paired membership are K103 SC26.5--13, retained in K108 AE1/11/16.

\textbf{Theorem.} The complete k=160 subset of TG1 has absolute normalized contribution less than \(2^{-319}\). Append319 once to the actual predecessor, retaining K108585, DO114354 and every prior ordered charge. No claim for161..164, complete support27, signed positivity or arithmetic nonvanishing follows.

Use the unchanged constants \[
 Q=1427908193729121522601,\quad U=77Q^2-14,\quad
 R_0=17424/133,\quad M^2=R_0/U,
\] \[
 D_0=7^{140}29^{144}5^{84}13^{84},\quad D_F=2^fD_0,
 \quad \pi_-=314159/100000,\quad r_+=1463/2 .
 \tag{TG2}
\] The proof retains the full parameter dependence of DO114's scalar theta estimate, and uses a larger lower bound for the actual transverse Gram. It does not replace the underlying ideal lattice by a rectangular lattice.

\subsection{2. Scalar theta for every parameter and the true transverse Gram}\label{scalar-theta-for-every-parameter-and-the-true-transverse-gram}

For a scalar fractional ideal L with \(N_EL=z^{2f}\), the original half-trace metric gives \[
 L^\#=2\mathfrak D_E^{-1}\bar L^{-1},\qquad
 \lambda_1(L^\#)^2\ge fD_0^{-2/f}z^{-2},\qquad
 D_0^{1/f}<r_+ .
\] DO114 SP3--10 proves, for every real n-dimensional lattice of minimum lambda and every t\textgreater0, \[
 \Theta_\Lambda(t)-1
 \le C_n2^{n/2}\Gamma(n/2+1)(t\lambda^2)^{-n/2},
 \qquad C_{336}<32000.
\] The proof retains dimension336 and its factorial168!. It comes from the orthoplex bound, spherical lift, cap averaging and nonzero-shell integral; there is no restriction to a small Gaussian parameter in this inequality.

Define the fixed positive rational number \[
 \mathcal K=32000\,2^f f!
             \left(\frac{r_+^2}{f\pi_-^2}\right)^f .
 \tag{TG3}
\] Using the exact strict \(\pi>\pi_-\), the whole scalar estimate is therefore \[
 \Theta_{L^\#}(\pi^2z^2/A)<1+\mathcal K A^f
                         \qquad(A>0).
 \tag{TG4}
\] DO114's boxed3/2 estimate is only one specialization of TG4; no restriction \(A\le1/43500\) is imposed here.

The original saturated transverse inputs from K107 SV17--19 have the actual Gram \[
 G'=\begin{pmatrix}
 978665066829&220696428039&-303935233242\\
 220696428039&208987227707&47956510467\\
 -303935233242&47956510467&1007543201084982869062
 \end{pmatrix},\qquad\det G'=U.
 \tag{TG5}
\] Put \(c_0=2^{36}\). The three leading principal minors of \(G'-c_0 I\) are exactly \[
 909945590093,\qquad78929108079151063344782,\qquad
 79524486185942701840832982385541548780235919.
\] All are positive. Sylvester's criterion proves \(G'>c_0I\). This is an estimate for the fixed real energy metric; it is not a positive metric assertion over a finite field.

The primal row metric for d=k+3 values is \(\operatorname{diag}(I_k,(G')^{-1})\), with exact real covolume \[
                         D_0^d z^{2df}/U^f.
 \tag{TG6}
\] The dual metric is \(\operatorname{diag}(I_k,G')\). Hence its first k scalar theta factors are bounded by TG4, and its transverse three factors are bounded by \[
                  [1+\mathcal K(A/c_0)^f]^3.
 \tag{TG7}
\] All three coordinates still range over the same original ideal L. This domination changes a quadratic-form upper estimate only; it does not remove the projection denominator U or impose stronger membership.

\subsection{3. Retain the entire actual energy in the Gaussian sum}\label{retain-the-entire-actual-energy-in-the-gaussian-sum}

K108 AE13--15 proves, before any lattice enlargement, the actual selected-trace inequality \[
 |c(h)|<
 2\vartheta_a^f k^{-84}2^{-933f/2}U^{-f/2}z^{-f}
                       e^{-2U\|r\|^2},
 \qquad
 \vartheta_a=\frac{4^{167}}{334\binom{334}{167}} .
 \tag{TG8}
\] Its full angular profile, lower-tail absorption, both compact factors, all-coordinate field-norm inequality and combined actual near-point factor2 are retained. The local quotient is for positive Cartan radius; at radius zero the underlying local trace is bounded by one. Nonidentity rank-two terms have positive alpha.

For a nonempty original k-sector, each nonzero S coordinate has squared Euclidean norm at least \(fz^2\), whereas the actual row has energy at most \(fM^2\). Therefore \[
 q=M^2/z^2\ge k,\qquad
 A=2Uz^2=2R_0/q.
 \tag{TG9}
\] At k=160, \(0<A\le2R_0/160<2\). The exact rational comparison \[
 \mathcal K(2/c_0)^f<1/16,\qquad (17/16)^3<2
 \tag{TG10}
\] makes the entire transverse dual theta in TG7 strictly below2. This bound holds for every admitted ideal and every q in the whole sector, not only at the final endpoint.

The energy weight in TG8 is now exactly \(e^{-A\|r\|^2/z^2}\). Poisson summation with the carried Fourier convention and TG4--7/10 gives \[
 \sum_{\rm actual\ rows}e^{-2U\|r\|^2}
 \le {164\choose k}\frac{2U^f}{D_0^d}\,
       \pi^{df}(A^{-f}+\mathcal K)^k A^{-3f}.
 \tag{TG11}
\] The two in this formula is the transverse theta bound, independent of the earlier selected near-point2. No positive Laplace-energy exponential appears.

The true selected bound is applied to actual nonzero k-tuples first. Only the resulting nonnegative weighted row sum is enlarged to the entire coupled lattice. Added zero or out-of-ball tuples are not assigned a trace estimate. The parameter inequality q\textgreater=k was already established for the original nonempty sector before this enlargement. All complex phases remain counted by that same lattice sum, with no additional phase or eta enumeration.

\subsection{4. Complete finite factors and both monotonicities}\label{complete-finite-factors-and-both-monotonicities}

The original local ideal and complete affine central count are unchanged: \[
 L=JI_\Sigma D_\eta,\quad N_EJ=z^{2f}T^{52}/x,\quad
 |a_G(h)|\le2Q^{-335\cdot167}(N_EJ)^{335/2},
\] \[
 \#d_{\rm cent}\le
 \frac{V_F(H/2+750Uz^2)U^{-f}T^{52}x^{-1}}{N_FJ_0},
 \quad N_EJ\le Q^{167}N_FJ_0,
\] \[
 V_F(C)=\frac{((22/7)fC^2)^{f/2}}{(f/2)!\sqrt{D_F}} .
 \tag{TG12}
\] These are K103 SC26.12--16, DO098 CF10 and K108 AE21. They retain the actual affine shift, full-real750 cover, original U denominator, common whole-C orientation and injective reconstruction from eta, row and central scalar. At both CM components the same integral-input proof, including eta denominator primes, remains in force. The two mixed finite projectors keep their accepted bounds.

The actual26 active weighted-label and142 inactive pointwise factors give \((T-1)^{-168}<2T^{-168}\). This is a statement on the retained actual set, not a constant coefficient claim on an enlarged central orbit.

Apart from fixed positive factors, TG8/11/12 give \[
 x^{-337/2}(N_FJ_0)^{-1}F_k(q),
\] \[
 F_k(q)=q^{-167f}\pi^{df}
 \left[(q/(2R_0))^f+\mathcal K\right]^k
 (q/(2R_0))^{3f}V_F(H/2+750R_0/q).
 \tag{TG13}
\] The factor2 in TG11 is kept outside this function and will remain inside the exact rational endpoint.

Put \[
 w(q)=\frac{q^f}{q^f+(2R_0)^f\mathcal K}\in(0,1),\qquad
 \epsilon(q)=\frac{750R_0/q}{H/2+750R_0/q}\in(0,1).
\] A direct derivative gives \[
 \frac1f\frac{d\log F_k}{d\log q}
                         =-164+k\,w(q)-\epsilon(q).
 \tag{TG14}
\] No derivative of an unknown lattice theta function is taken.

For fixed eta and x, the all-rank-zero endpoint is \[
 z_*=(Q^{334}T^{-52}x)^{1/(2f)},\quad q_*=M^2/z_*^2,
 \quad (N_FJ_0)^{-1}\le Q^{-167}(q/q_*)^f.
\] At k=160, the complete enlarged q factor decreases at least with exponent−3f. Summing every original local state2,1,0 with multiplicities1,2,1 therefore gives \[
 \sum_{\rm patterns}\frac{F_{160}(q)}{N_FJ_0}
 \le Q^{-167}F_{160}(q_*)\prod_G(1+q_G^{-3})^2
 <2Q^{-167}F_{160}(q_*).
 \tag{TG15}
\] The exact sufficient comparison is \((1+g_0^{-3})^{668}<2\), where \(g_0=22285932761\). Denominator overlap with G does not alter the ideal table or introduce an ideal-class factor. The parameter-dependent estimate is uniform in the ideal, without an assumed homothety.

After TG15 the recovered \(Q^{-167}\) is independent of x, while \(q_*\propto x^{-1/f}\). The full logarithmic x derivative is \[
 -337/2-[-164+160w(q_*)-\epsilon(q_*)]
 =-9/2-160w(q_*)+\epsilon(q_*)<-7/2.
 \tag{TG16}
\] Thus the complete bound, including every finite pattern and the moving central radius, is largest at \(x=2^{X_0}\). There is no Gaussian branch, additional endpoint union or independent eta count.

\subsection{5. One exact complete cell and chronology}\label{one-exact-complete-cell-and-chronology}

Use the unchanged exact coupled bracket and central bound: \[
 q_*^f=\frac{R_0^fT^{52}}{U^fQ^{334}2^{X_0}},\quad
 c_-=4119905<q_*<4119906=c_+,
\] \[
 H/2+750R_0/c_-<68,\qquad V_F(68)<2^{479}.
 \tag{TG17}
\] Retain support101, original harmless pivot8, Green1, weighted-pattern1 and the combined T correction1. Their exponent remains112. The additional transverse theta2 is retained inside the rational factor below. The selected near-point2 is already in the trace saving335 and is not counted again.

Group the full signed U/R0/T/Q/x powers before using the same directed bounds as K108. The unchanged comparison constant is \[
 C=112+2f-929f/2+479+167fL_R-(335f/2)L_U+
       8594L_T-336\cdot167L_Q-(337/2)X_0
   =211125067/512,
\] where \(L_T=1019333/1024,\ L_Q=287843/4096,\ L_U=601356/4096,\ L_R=28810/4096\), upper for T,R0 and lower for Q,U. R0 in the following rational expression remains its exact fixed value.

With k=160,d=163, form \[
 W=\frac{2\binom{164}{160}(22/7)^{163f}
    [(c_+/(2R_0))^f+\mathcal K]^{160}
    (c_+/(2R_0))^{3f}}
             {D_0^{163}c_-^{167f}}.
\] Exact integer numerator/denominator comparison gives \[
 W<2^{-410738},\qquad160^{84}>2^{615},\qquad
 \vartheta_a^f<2^{-649}.
 \tag{TG18}
\] Both appearances of the cap use the safe directed endpoint: c\_+ in its increasing numerator and c\_- in its decreasing denominator. The full parameter monotonicity was proved before this separate rational enclosure.

The complete sector exponent is therefore \[
 E_{160}
 =C-410738-336-335-649-615
 =-163509/512<-319 .
 \tag{TG19}
\] The saturated-basis saving336, selected coefficient improvement335 and radial649 occur exactly once. There is one complete coordinate cell and no further coordinate union. Append319 once; no other charge is replaced.

The executable accepts the actual canonical K108 predecessor, or a later complete ordered extension, and retains both full arrays. At K108's116/31 base the result is117/31, with complete floors26/11 unchanged and both full sums below \(2^{-230}\). The already charged142..159 interval,139..141 interval and k138 cell remain in their original chronological positions. Cells161..164, signed sufficient positivity and a full support27 assembly remain open.

\subsection{6. Verification and source scope}\label{verification-and-source-scope-1}

The new proof consists of the all-A specialization TG3--4, exact stronger Gram inequality TG5, original energy-weighted Poisson sum TG9--11 and complete monotonic/count argument TG13--19. The exact Gram check was also independently requested from /root/definite\_oscillator; it is a targeted algebra check, not a substitute for the full theorem review.

Direct dependencies are K108 AE1--30, DO114 SP3--12, K107 SV17--19, K103 SC26.5--20, DO098 CF4--13 and DO070 FT3. K108 retains the genuine DO080/Nakahama model and complete lower-tail/near-point normalization; there is no new selected module. The universal scalar packing argument is completely written in DO114 and independently whole-reviewed. Finite diagnostics corroborate matrix arithmetic, rational inequalities, count exponents and chronological arrays; they do not prove the original local membership or universal Poisson and differentiation steps.

Primary normalization remains Nakahama, R. (2016), \emph{Norm computation and analytic continuation of vector valued holomorphic discrete series representations}, Journal of Lie Theory,26(4),927--990, equation5.7 and Theorem5.1, printed960--961; Miller, S. D., and Stephens-Davidowitz, N. (2019), \emph{Kissing numbers and transference theorems from generalized tail bounds}, arXiv:1802.05708v2, equations1.1/2.1; and Rankin, R. A. (1955), \emph{The closest packing of spherical caps in n dimensions}, Proceedings of the Glasgow Mathematical Association,2(3),139--144. These are carried with the exact fresh/carry scopes recorded in K108 and DO114, not claimed as new whole-paper readings here.

The unchanged contextual anchors are SGA3 I4.7.3/4.7.3.1; SGA6 VI Theorem3.3; FGA232 Theorem3.1 pp148--149; FGA236 Corollaries1.7/2.3 pp229/232; SGA7 I VII2.7/2.9 pp169--174 and VIII7.4(b) pp310--311; cote4 scan003 sections8.1.1--8.1.2/scans005--006 sections8.4.1--8.4.3 and cote54 scans005--007. No new motivic, lambda-ring, K-theory, weight, formal-group or CM comparison is made. The actual selected arithmetic image, prescribed period, effective source/auxiliary construction, Read/specific032, actual DP44 and RH remain separate obligations.

\subsection{Canonical algebraic-geometric context}\label{canonical-algebraic-geometric-context-43}

The original paired fractional modules and exact integral saturated transverse lattice retain their true Hermitian Gram and projection denominator. The stronger real dual-metric inequality and Gaussian count act on the same actual selected complex realization and theta line. No rational orthogonal frame is declared integral. The previous exact Nakahama(2016)equation5.7/Theorem5.1pp960-961, Gaussian Fourier convention of Miller and Stephens-Davidowitz(2019)equations1.1/2.1, and Rankin(1955)Theorem1(iv) context remain in the directly checked scopes of K108/DO114. Carried geometric anchors: SGA3 I4.7.3/4.7.3.1; SGA6 VI Theorem3.3; FGA232 Theorem3.1pp148-149; FGA236 Corollaries1.7/2.3pp229/232; SGA7 I VII2.7/2.9pp169-174 and VIII7.4(b)pp310-311; cote4 scan003 sections8.1.1-8.1.2/scans005-006 sections8.4.1-8.4.3 and cote54 scans005-007. No new formal-group slope, lambda-ring/K-theory identity, l-adic weight or CM-motive comparison. DP\_MULT Q3e/DP44/DP52/DP54-65 and THETA3 G1/A7 gain this finite reduction; CX\_GRH CX-ALG-024 and GCT\_PLETH T-052 retain their global gates.

\clearpage

\section{AM-DO-115 --- a cubic spherical-cap bound and the complete k161 sector}\label{am-do-115-a-cubic-spherical-cap-bound-and-the-complete-k161-sector}

\textbf{Canonical checked integration --- 2026-11-23T06:53:45.453886+00:00.} Complete root and independent proof/script/interface reviews and separate exact replays passed. The actual canonical predecessor has117/31 entries; this integration gives118/31, complete floors26/11. Frozen author history is preserved. AM-DO-115 DC1-24 proves an all-parameter scalar theta estimate using the exact positive cubic kernel, finite cap occupancy and nonzero-shell integral in real dimension336. The original selected angular profile gives beta777/352, and the unchanged saturated transverse Gram gives a complete257/256 theta factor. With every original paired ideal, affine central fibre, finite factor and both continuum monotonicities retained, the complete k161 support26 cell is below2\^{}-233, with exact exponent-119477/512. The necessary actual all-II rank-two-unipotent remainder becomes162\textless=k\textless=164 with2\^{}(71443/64)\textless x\textless2\^{}2108,v=q26,b=j0. Exactly one233 append gives118/31 complete ordered errors, both entire sums below2\^{}-230; unique minimum231,next233 and complete floors26/11 remain. No162+, signed positivity, selected arithmetic content, effective source, period, specific032/Read, evaluatedDP44 or RH conclusion.

Author: /root/definite\_oscillator. 23 November 2026 UTC. Controlled pending proof. Author verification, independent review and canonical integration are separate events. This document does not alter any frozen predecessor.

\subsection{1. Retained actual sector and theorem}\label{retained-actual-sector-and-theorem}

Retain the actual all-top J56 family, H=133 and f=a=168, its original E/F, involution, Hermitian polarization, fractional projective modules, paired source/target pencils and admitted labels, selected vector-valued Hilbert metric, both compact factors, canonical words and twists, finite projectors, phases, Haar normalization and class masses. After K109 the remaining all-II support26 scalar-cross globally rank-two-unipotent set has \[
161\le k\le164,\quad v=q_{\rm elig}=26,\quad b=j=0,
\quad 2^{X_0}<x=N_ED_\eta<2^{2108},\qquad X_0=71443/64.
\tag{DC1}
\] Here k counts the nonzero original S coordinates. K103 SC26.5--13 supplies the normal form, original membership and unique eta for each support in this window; K107 SV26--27 supplies the exact strict x floor. No eta or denominator ideal is independently enumerated.

\textbf{Theorem.} The complete k=161 subset of DC1 has absolute normalized contribution less than \(2^{-233}\). Append233 once to the actual K109 predecessor, retaining all earlier charges and both complete ordered arrays. The full support floors remain26/11. The new coordinate condition is k\textgreater=162 in this branch only. No complete support27, signed positivity, arithmetic content or period assertion follows.

Use the unchanged constants \[
Q=1427908193729121522601,\quad U=77Q^2-14,
\quad R=17424/133,\quad M^2=R/U,
\] \[
D_0=7^{140}29^{144}5^{84}13^{84},\quad D_F=2^fD_0,
\quad n=2f=336,\quad t=49/1000.
\tag{DC2}
\] The new scalar estimate is valid for every Gaussian parameter. The actual energy improvement and three-coordinate estimate below are uniform throughout the original sector.

\subsection{2. A self-contained cubic positive-kernel bound}\label{a-self-contained-cubic-positive-kernel-bound}

For unit vectors v,w in real n-space put x=\textless v,w\textgreater{} and \[
P_0(x)=1,\quad P_1(x)=x,\quad
P_2(x)=\frac{nx^2-1}{n-1},\quad
P_3(x)=\frac{(n+2)x^3-3x}{n-1}.
\tag{DC3}
\] Each is a positive kernel: P0 is constant and P1 is the ordinary Gram kernel. For P2 use the tensors \(H_2(v)=v\otimes v-I/n\), whose inner product is \(x^2-1/n\). For P3 define \[
S(v)_{ijk}=v_i\delta_{jk}+v_j\delta_{ik}+v_k\delta_{ij},
\qquad H_3(v)=v^{\otimes3}-S(v)/(n+2).
\] Direct contraction gives \[
\langle v^{\otimes3},S(w)\rangle=3x,\quad
\langle S(v),S(w)\rangle=3(n+2)x,
\quad\langle H_3(v),H_3(w)\rangle=x^3-3x/(n+2).
\] The displayed kernels are positive multiples of these Gram kernels. Consequently, for every finite unit-vector list, the double sum of each Pj is nonnegative; the P0 sum is the square of its cardinality. No assumed positive finite-field form is involved.

Set \[
a=\frac{1+t}{1+nt}=1049/17464,\qquad F(X)=(X-t)(X+a)^2.
\tag{DC4}
\] Its expansion \(F=\sum_{j=0}^3 f_jP_j\) has \[
f_0=(2a-t)/n-ta^2,\quad f_1=3/(n+2)+a^2-2at,
\quad f_2=(2a-t)(n-1)/n,\quad f_3=(n-1)/(n+2).
\] At the fixed exact n,t,a these are, respectively, \[
\frac{74539509}{2134939072000},\quad
\frac{42505405671}{6442941128000},\quad
\frac{3467987}{48899200},\quad\frac{335}{338},
\tag{DC5}
\] all positive. If m unit vectors have pairwise inner products \textless=t, each off-diagonal F term is nonpositive. Positivity of the four kernels therefore gives \[
f_0m^2\le\sum_{i,j}F(\langle v_i,v_j\rangle)
\le mF(1),\qquad
m\le\mathcal M:=\frac{F(1)}{f_0}
=\frac{n(1-t)(2+(n+1)t)}{1-nt^2}
=\frac{369723123}{12079}.
\tag{DC6}
\] The argument includes coincident-point exclusion, since the permitted off-diagonal bound t\textless1 already makes all vectors distinct. This is the degree-three linear-programming estimate of Delsarte, Goethals, and Seidel (1977), Theorem4.3 and Example4.6. The tensor calculation above proves all positive-kernel input needed here; no general harmonic-analysis theorem is left unstated.

\subsection{3. Finite cap occupancy, exact dimension and scalar theta}\label{finite-cap-occupancy-exact-dimension-and-scalar-theta}

Let X be a finite 1-separated subset of the Euclidean n-ball of radius rho, where rho\textgreater=1. Lift it to the unit sphere in real n+1-space by \[
x\longmapsto\widetilde x=\rho^{-1}(x,\sqrt{\rho^2-\|x\|^2}).
\] Distances do not decrease before the division by rho, so distinct lifted points have inner product at most \(s=1-1/(2\rho^2)\in[1/2,1)\). Choose a cap with pole height h and angle theta satisfying \[
h^2=\frac{s-t}{1-t},\qquad
\sin^2\theta=\frac1{2\rho^2(1-t)}.
\tag{DC7}
\] Since t\textless1/3, \[
h^2-s^2=\frac{(1-s)[s-t(1+s)]}{1-t}>0.
\] Thus h\textgreater s. A point at the cap pole cannot coexist with another cap point, because their inner product would be at least h\textgreater s. Otherwise normalize the projections of all cap points onto the n-dimensional pole complement. If their pole coordinates are u,v\textgreater=h, their normalized projected inner product is at most \[
\frac{s-uv}{\sqrt{(1-u^2)(1-v^2)}}.
\] Its derivative with respect to u has the sign of us-v, which is negative since v\textgreater=h\textgreater s\textgreater=us; the analogous statement holds in v. The maximum is therefore at u=v=h and equals t. Applying DC6 in the projected dimension n bounds every cap occupancy by M. The lifted sphere has ambient dimension n+1; no extra dimension is inserted into DC6.

The normalized cap area is \[
\mu_n(\theta)=c_n\int_0^\theta\sin^{n-1}\phi\,d\phi
\ge(c_n/n)\sin^n\theta,
\quad c_n=\frac{\Gamma((n+1)/2)}{\sqrt\pi\Gamma(n/2)}.
\] The inequality follows by inserting cos(phi)\textless=1 and integrating the derivative of sin(phi)\^{}n.~Averaging cap occupancy over the pole gives \[
\#X\le\frac{n\mathcal M}{c_n}[2(1-t)]^{n/2}\rho^n.
\tag{DC8}
\] For a real lattice of minimum lambda, count only nonzero vectors in the radius-rho*lambda ball. That count vanishes for rho\textless1. With b=t\_G lambda\^{}2\textgreater0, Tonelli and DC8 give \[
\Theta_\Lambda(t_G)-1
=\int_1^\infty2b\rho e^{-b\rho^2}N(\rho)\,d\rho
\le\frac{n\mathcal M}{c_n}[2(1-t)]^{n/2}
\Gamma(n/2+1)b^{-n/2}.
\tag{DC9}
\] The last integral was enlarged to rho\textgreater=0 only after applying a bound valid for nonzero shells rho\textgreater=1. The zero vector is restored exactly once by the1 on the left.

For n=336, \[
c_{336}=168\binom{336}{168}/4^{168},\qquad
D_0^{1/f}<r_+:=365729/500,
\quad\pi_-:=314159/100000<\pi<\pi_+:=3141593/10^6.
\tag{DC10}
\] The root-discriminant comparison is the exact integer inequality \(500^{168}D_0<365729^{168}\). The two pi bounds follow from the alternating Machin series: \(\pi=16\arctan(1/5)-4\arctan(1/239)\). Its tangent and angle range are checked exactly; six terms of the first series and the first term of the second give the lower bound, while seven terms and the two-term lower second series give the upper bound.

The actual half-trace scalar dual, without a change of normalization, is \[
L^\#=2\mathfrak D_E^{-1}\bar L^{-1},\quad
N_EL=z^{2f},\quad
\lambda_1(L^\#)^2\ge fD_0^{-2/f}z^{-2}.
\] This is DO070 FT3, retained in DO114 SP10 and K109 TG3--4. Define the explicit positive rational constant \[
\mathcal K_c=
\frac{336\mathcal M}{c_{336}}(951/500)^{168}168!
\left(\frac{(365729/500)^2}{168(314159/100000)^2}\right)^{168}.
\tag{DC11}
\] DC9--10 prove, uniformly in the fractional ideal, its class, denominators and all A\textgreater0, \[
\boxed{\ \Theta_{L^\#}(\pi^2z^2/A)<1+\mathcal K_c A^{168}.\ }
\tag{DC12}
\] No rectangular model for L, small-A hypothesis, or assumed homothety of different ideals is used.

\subsection{4. Exact actual energy and transverse theta}\label{exact-actual-energy-and-transverse-theta}

K108 AE3--15 proves the actual selected angular multiplier \(\Psi(\alpha)=\exp[-\nu g(\alpha)]\), where \[
\nu=333/2,\qquad g(y)=\log((\sqrt{1+y}+1)/2),\qquad0\le\alpha\le R<132.
\] Here g(0)=0 and g is concave. Since sqrt133\textgreater11, g(132)\textgreater log6\textgreater7/4. Concavity along {[}0,132{]} yields \[
\nu g(\alpha)\ge\frac{333}{2}\frac{7}{4\cdot132}\alpha
=\beta\alpha,\qquad\beta=777/352.
\tag{DC13}
\] The inequality log6\textgreater7/4 is certified, for example by the positive series for log6 with parameter5/7 and a finite partial sum. Thus the same K108 proof retains the stronger factor \(e^{-\beta U\|r\|^2}\), since alpha\textgreater=U times the row energy. This does not reapply the angular profile after counting. Its all-word variance bound, exact radial factor, actual multiplicative near-point product and both compact factors remain exactly those in K108; the combined near-point loss is still the already charged factor2, not a newly bounded perturbation.

The actual saturated transverse Gram is \[
G'=\begin{pmatrix}978665066829&220696428039&-303935233242\\
220696428039&208987227707&47956510467\\
-303935233242&47956510467&1007543201084982869062\end{pmatrix},
\quad\det G'=U.
\] Its leading minors after subtracting2\^{}37 I are \[
841226113357,\quad11481363302936327071374,\quad
11567969518504471682008247001894572054490255.
\tag{DC14}
\] All are positive, so G'\textgreater2\^{}37 I. The original saturated integral inputs and determinant are unchanged from K107 SV17--19 and K109 TG5--7. This real inequality asserts no finite-field positivity.

Write q=M²/z². For any nonempty original k=161 sector, the norm of each nonzero S coordinate gives q\textgreater=k. Set \[
A=\beta Uz^2=\beta R/q<2.
\] For the transverse three scalar factors use the already proved DO114/K109 coefficient \(\mathcal K_o=32000\,2^{168}168![(1463/2)^2/(168\pi_-^2)]^{168}\). Its exact comparison \(\mathcal K_o(1/43500)^{168}<1/2\), together with \(87000\cdot2^{20}<2^{37}\), gives \[
\mathcal K_o(2/2^{37})^{168}<2^{-3361},\qquad
(1+2^{-3361})^3<257/256.
\tag{DC15}
\] Thus the complete transverse dual theta is below257/256 uniformly for every admitted q. This is a replacement of the transverse2 in K109 only; its exact determinantU and original projection denominators remain.

The true selected trace is bounded first on the actual nonzero sector by \[
2\vartheta^f k^{-84}2^{-933f/2}U^{-f/2}z^{-f}
e^{-\beta U\|r\|^2},\qquad
\vartheta=4^{167}/[334\binom{334}{167}].
\] Only its resulting positive row sum is enlarged. Poisson summation for the actual primal metric diag(I\_k,(G')\^{}\{-1\}), whose real covolume is \(D_0^{k+3}z^{2(k+3)f}/U^f\), gives \[
\sum_{\rm actual\ rows}e^{-\beta U\|r\|^2}
\le\binom{164}{k}\frac{257U^f}{256D_0^{k+3}}\pi^{(k+3)f}
(A^{-f}+\mathcal K_c)^k A^{-3f}.
\tag{DC16}
\] Added zero or out-of-ball lattice tuples are not assigned a selected trace estimate. There is no positive Gaussian compensating exponential. Complex phases remain in the same lattice sum.

\subsection{5. Original finite factors, affine central fibre and all parameters}\label{original-finite-factors-affine-central-fibre-and-all-parameters}

Retain K109 TG12--17 in full: \[
L=JI_\Sigma D_\eta,\quad N_EJ=z^{2f}T^{52}/x,\quad
|a_G(h)|\le2Q^{-335\cdot167}(N_EJ)^{335/2},
\] \[
\#d_{\rm cent}\le
\frac{V_F(H/2+750Uz^2)U^{-f}T^{52}x^{-1}}{N_FJ_0},
\quad N_EJ\le Q^{167}N_FJ_0,
\] \[
V_F(C)=\frac{((22/7)fC^2)^{f/2}}{(f/2)!\sqrt{D_F}},
\qquad(T-1)^{-168}<2T^{-168}.
\tag{DC17}
\] These preserve one affine central coset, its fixed real-part shift, full750 cover, exact original U denominator, the shared whole-C orientation at both CM partners, eta denominator intersections and original reconstruction from eta,row,central scalar. No isometry fibre is dropped. Both mixed projectors keep their modulus-one bounds. The last inequality combines26 active weighted labels and142 inactive pointwise factors on the retained set; it is not an extension to arbitrary central-orbit terms.

Apart from fixed positive factors, the whole bound is \(x^{-337/2}(N_FJ_0)^{-1}F_k(q)\), where \[
F_k(q)=q^{-167f}\pi^{(k+3)f}
[(q/(\beta R))^f+\mathcal K_c]^k(q/(\beta R))^{3f}
V_F(H/2+750R/q).
\] Put \(w=q^f/[q^f+(\beta R)^f\mathcal K_c]\) and \(\epsilon=(750R/q)/(H/2+750R/q)\), both in(0,1). Then \[
f^{-1}d\log F_k/d\log q=-164+kw-\epsilon.
\tag{DC18}
\] At fixed x the all-rank-zero endpoint is \[
z_*=(Q^{334}T^{-52}x)^{1/(2f)},\quad q_*=M^2/z_*^2,
\quad(N_FJ_0)^{-1}\le Q^{-167}(q/q_*)^f.
\] For k161 the complete enlarged q factor decreases at least as q\^{}\{-2f\}. The actual local states2,1,0 have multiplicities1,2,1. Summing them, including the paired orientations exactly once, gives \[
\sum_{\rm patterns}(N_FJ_0)^{-1}F_{161}(q)
\le Q^{-167}F_{161}(q_*)\prod_G(1+q_G^{-2})^2
<2Q^{-167}F_{161}(q_*).
\tag{DC19}
\] The sufficient integer inequality is \((1+22285932761^{-2})^{668}<2\). Overlap of eta denominators with G does not change the carried J/J0 table and adds no ideal-class count. The bound is uniform in L.

Since q\_* is proportional to x\^{}\{-1/f\}, the coupled full x derivative is \[
-337/2-[-164+161w-\epsilon]=-9/2-161w+\epsilon<-7/2.
\tag{DC20}
\] Hence the complete bound is maximal at x=2\^{}\{X0\}. No endpoint or Gaussian-parameter union is introduced. At that point the exact coupled bracket and central volume are \[
q_*^f=R^fT^{52}/(U^fQ^{334}2^{X_0}),\qquad
c_-=4119905<q_*<4119906=c_+,
\] \[
H/2+750R/c_-<68,\qquad V_F(68)<2^{479}.
\tag{DC21}
\]

\subsection{6. Complete exact endpoint and ledger}\label{complete-exact-endpoint-and-ledger}

Retain the support bound binom(168,26)\textless2\^{}101, harmless original pivot164\textless2\^{}8, Green2, pattern2 and combined T correction2. Their exponent is112. The additional257/256 is kept inside the rational endpoint. The near-point2 is already in the trace saving335 and is not added again. There is no extra type, eta, phase, witness, inverse or coordinate union.

With the same directed host bounds as K109, \[
L_T=1019333/1024,\quad L_Q=287843/4096,\quad
L_U=601356/4096,\quad L_R=28810/4096,
\] upper for T,R and lower for Q,U, the complete grouped constant is \[
C=112+2f-929f/2+479+167fL_R-(335f/2)L_U+
8594L_T-336\cdot167L_Q-(337/2)X_0=211125067/512.
\tag{DC22}
\] The T exponent8594 is335*26+52-168; the Q exponent includes the recovered Q\^{}\{-167\}. Keep exact R in the rational factor \[
W=\frac{(257/256)\binom{164}{161}\pi_+^{164f}
[(c_+/(\beta R))^f+\mathcal K_c]^{161}
(c_+/(\beta R))^{3f}}{D_0^{164}c_-^{167f}}.
\tag{DC23}
\] Exact integer numerator/denominator comparisons give \[
W<2^{-410652},\qquad161^{84}>2^{615},\qquad
\vartheta^{168}<2^{-649}.
\] Both directions of the cap enclosure are safe: increasing numerator at c\_+, decreasing denominator at c\_-. DC18--20 establish the full continuous comparison before this enclosure.

The complete cell exponent is \[
E_{161}=C-410652-336-335-649-615
=-119477/512<-233.
\tag{DC24}
\] The saturated determinant improvement336, coefficient/near-point improvement335, exact radial649 and all-coordinate615 are each used once. There is exactly one complete coordinate cell. The executable retains every old charge, appends233 once and verifies both full ordered totals below2\^{}-230. At the actual K109117/31 base this gives118/31. The existing231 minimum remains unique; the new233 becomes the next-smallest charge. The complete support floors remain26/11. Cells162..164 and all separate global/arithmetic gates remain open.

\subsection{7. Sources and scope of verification}\label{sources-and-scope-of-verification}

The new universal mathematical ingredients are DC3--12; DC13 retains a sharper rational consequence of the already proved actual angular profile; DC14--24 retain the complete original metric/count and strengthen its exact endpoint. Finite polynomial/cap and derivative models are diagnostics, not substitutes for the universal proofs or original module/source interfaces.

The freshly read primary source is Delsarte, P., Goethals, J. M., \& Seidel, J. J. (1977). Spherical codes and designs. \emph{Geometriae Dedicata, 6}(3),363--388. \url{https://doi.org/10.1007/BF03187604}. Exact fresh scope: Theorem3.3 and its addition formula; Theorem3.6, Corollaries3.7--3.8 and proofs; Theorem4.3 and proof; Example4.6 and the following Example4.7, printed367--370. Its Gegenbauer polynomials use a different positive normalization; DC3 fixes the normalization used here and proves positivity directly. No whole-paper reading or later degree bound is claimed.

Direct carried dependencies are DO114 SP3--12 (cap normalization, nonzero shells and actual scalar dual), K108 AE3--15/21--30 (genuine angular profile, all-word trace, actual near-point normalization and complete finite factors), K109 TG1--19 (full original Gram, weighted Poisson count and all factors), K107 SV17--19/26--27 (saturated input basis and eligibility floor), K103 SC26.5--20 and DO098 CF4--13 (original paired ideal, affine central fibre, correlation and reconstruction), and DO070 FT3. The exact source manifest fingerprints these files. Nakahama(2016)equation5.7/Theorem5.1pp960--961 and the Miller--Stephens-Davidowitz Gaussian Fourier convention remain in their already checked scopes; no fresh whole-source reading is asserted.

The unchanged algebraic-geometric context is SGA3 I4.7.3/4.7.3.1; SGA6 VI Theorem3.3; FGA232 Theorem3.1pp148--149; FGA236 Corollaries1.7/2.3pp229/232; SGA7 I VII2.7/2.9pp169--174 and VIII7.4(b)pp310--311; cote4 scan003 sections8.1.1--8.1.2/scans005--006 sections8.4.1--8.4.3 and cote54 scans005--007. These are inherited context only. The field, polarization, theta line and selected arithmetic realization are unchanged. No new lambda-ring, K-theory, l-adic weight, formal-group or CM-motive comparison is made. Actual selected arithmetic content, prescribed period, effective source/auxiliary scales, Read/STiesR/specific032, DP44 and RH remain distinct obligations.

\subsection{Canonical algebraic-geometric context}\label{canonical-algebraic-geometric-context-44}

The universal positive-kernel and Gaussian estimates act on the unchanged actual selected complex realization. Original projective integral modules, paired CM lattices, saturated transverse basis, Hermitian metric, projection denominator and theta line are retained. Delsarte, Goethals, and Seidel(1977),Theorems3.3/3.6/4.3,Corollaries3.7-3.8 and Examples4.6-4.7,printed367-370,provide the exact primary context; the needed degree-three kernels and full cap/count form are proved explicitly. The previous exact Nakahama(2016)equation5.7/Theorem5.1pp960-961, Gaussian Fourier convention of Miller and Stephens-Davidowitz(2019)equations1.1/2.1, and Rankin(1955)Theorem1(iv) context remain in the directly checked scopes of K108/DO114. Carried geometric anchors: SGA3 I4.7.3/4.7.3.1; SGA6 VI Theorem3.3; FGA232 Theorem3.1pp148-149; FGA236 Corollaries1.7/2.3pp229/232; SGA7 I VII2.7/2.9pp169-174 and VIII7.4(b)pp310-311; cote4 scan003 sections8.1.1-8.1.2/scans005-006 sections8.4.1-8.4.3 and cote54 scans005-007. No new formal-group slope, lambda-ring/K-theory identity, l-adic weight or CM-motive comparison. DP\_MULT Q3e/DP44/DP52/DP54-65 and THETA3 G1/A7 gain this finite reduction; CX\_GRH CX-ALG-024 and GCT\_PLETH T-052 retain their global gates.

\clearpage

\section{AM-K-110 --- retain the exact angular profile in the actual transverse Gaussian}\label{am-k-110-retain-the-exact-angular-profile-in-the-actual-transverse-gaussian}

\textbf{Canonical checked integration --- 2026-11-23T07:01:17.694071+00:00.} Complete root and independent proof/script/interface reviews and separate exact replays passed. The actual canonical predecessor has118/31 entries; this integration gives119/31, complete floors26/11. Frozen author history is preserved. AM-K-110 PM1-22 retains the full actual angular profile as a positive Laplace mixture, includes the entire low-parameter tail in an explicit atom, and computes the exact inverse third moment. Weighted AM-GM and the all-minimum/all-parameter DO115 scalar theorem control arbitrary embedding weights in the unchanged original ideal and transverse Gram. Every transverse correction is retained, giving a uniform2180-bit row improvement. With all original affine central, paired ideal, finite and source factors and both continuum monotonicities, the complete k162 support26 cell is below2\^{}-1071, with exact exponent-548533/512. The necessary actual all-II rank-two-unipotent remainder becomes163\textless=k\textless=164 with2\^{}(71443/64)\textless x\textless2\^{}2108,v=q26,b=j0. Exactly one1071 append gives119/31 complete ordered errors, both entire sums below2\^{}-230; minimum231,next233 and complete floors26/11 remain. No k163/164, signed positivity, selected arithmetic content, effective source, period, specific032/Read, evaluatedDP44 or RH conclusion.

Author: /root/kummer\_extension. 23 November 2026 UTC. Controlled pending theorem. Author verification, independent review and canonical integration are separate events. No frozen predecessor is changed.

\subsection{1. Exact original trace and purpose}\label{exact-original-trace-and-purpose}

Keep the actual all-top J56 support26 globally rank-two-unipotent scalar-cross branch, H=133, f=a=168, v=q\_elig=26, b=j=0, unique eta per support and \(2^{71443/64}<x=N_ED_\eta<2^{2108}\). All original fields, involution, projective modules, paired source and target pencils, finite projectors, phases, selected vector-valued Hilbert metric, both compact factors, words, twists, class masses and Haar measures remain fixed.

\textbf{Theorem.} After DO115's complete k161 reduction, the complete k=162 subset of this actual branch has absolute normalized contribution less than \(2^{-1071}\). Append1071 once to the actual DO115 predecessor, retaining both full ordered error arrays and every earlier charge. At the118/31 predecessor the result is119/31. The complete support floors remain26/11; the new remaining coordinate condition is k\textgreater=163 in this branch only. The positive-mixture and anisotropic-lattice inequalities below are universal in their explicitly stated ranges.

The genuinely selected estimate K108 AE13, before its uniform exponential simplification, is \[
 |c(h)|<
 2\vartheta_a^f2^{-933f/2}N_F(\alpha)^{-1/2}
                  \prod_\tau\Psi(\alpha_\tau),
 \quad
 \Psi(t)=\left(\frac2{\sqrt{1+t}+1}\right)^\nu,
 \quad \nu=333/2.
 \tag{PM1}
\] Here \(\alpha_\tau=U r_\tau J_Z^{-1}r_\tau^*>0\), \(\alpha_\tau\le R_0=17424/133<132\), and \(\vartheta_a^f<2^{-649}\). The full K107 lower tail is already uniformly absorbed; the leading2 is the single actual near-point factor, including the angular transport. It is not multiplied again below.

The original k nonzero S coordinates in \(L=JI_\Sigma D_\eta\), \(N_EL=z^{2f}\), imply \[
 N_F(\alpha)\ge U^f(kz^2)^f,\qquad
 q=M^2/z^2\ge k,\quad M^2=R_0/U .
 \tag{PM2}
\] Use PM2 on the actual nonzero sector before any Gaussian enlargement. The purpose of this proof is to retain PM1's full profile and obtain a uniform multiplicative improvement over the complete K109 row bound, without an uncharged small-probability or mixture tail.

\subsection{2. An explicit positive Laplace measure}\label{an-explicit-positive-laplace-measure}

For t,s\textgreater0 the elementary subordination identity is \[
 e^{-t\sqrt s}
 =\frac{t}{2\sqrt\pi}\int_0^\infty
       b^{-3/2}e^{-t^2/(4b)}e^{-sb}\,db .
 \tag{PM3}
\] For completeness, put \(F(c)=\int_0^\infty e^{-u^2-c^2/u^2}\,du\). For c\textgreater0, differentiation on compact positive c-ranges is justified by an integrable majorant. Substituting v=c/u gives \(F'(c)=-2F(c)\); dominated convergence at c=0 gives \(F(0)=\sqrt\pi/2\). Thus \(F(c)=(\sqrt\pi/2)e^{-2c}\). Substitution \(b=t^2/(4u^2)\) proves PM3. This is also the positive-real specialization of NIST DLMF10.32.10 and10.39.2, directly checked; the argument above fixes the normalization independently.

The gamma integral and PM3 give a positive measure on b\textgreater0: \[
 d\mu_\nu(b)=
 \frac{2^\nu}{2\sqrt\pi\Gamma(\nu)}
 b^{-3/2}e^{-b}
 \left(\int_0^\infty t^\nu e^{-t-t^2/(4b)}\,dt\right)db,
\] \[
                   \Psi(\alpha)=\int_0^\infty e^{-b\alpha}\,d\mu_\nu(b)
                     \qquad(\alpha\ge0).
 \tag{PM4}
\] Every integrand is nonnegative, so Tonelli justifies both orders of integration. Evaluation at alpha=0 gives total mass1. The measure is an analytic counting device; no new finite projector, coefficient field or freely chosen source vector is introduced.

For b\textless2 and \(0\le\alpha\le132\), \[
 e^{2\alpha}\int_{b<2}e^{-b\alpha}\,d\mu_\nu(b)
 \le e^{264}\int_{b<2}e^{-132b}\,d\mu_\nu(b)
 \le e^{264}\Psi(132)<2^{-39}.
 \tag{PM5}
\] The last inequality is exact: K108 AE14 proves \(\log((\sqrt{133}+1)/2)>7/4\), hence the logarithm of its preceding expression is below \(264-(333/2)(7/4)=-219/8\). The elementary bounds \(e>27/10\), \((27/10)^7>2^{10}\) and \(2190/56>39\) suffice.

Define the finite positive majorizing measure \[
 \widetilde\mu=\mu_\nu|_{[2,\infty)}+2^{-39}\delta_2 .
\] It is supported on b\textgreater=2, has mass below2, and on the whole actual alpha range \[
                         \Psi(\alpha)\le
 \int_{[2,\infty)}e^{-b\alpha}\,d\widetilde\mu(b).
 \tag{PM6}
\] The discarded mixing interval is retained in the explicit atom. There is no additive error left for the arithmetic sum.

\subsection{3. The exact inverse third moment}\label{the-exact-inverse-third-moment}

By Tonelli and the elementary gamma integral, \[
 m_3:=\int_0^\infty b^{-3}\,d\mu_\nu(b)
 =\frac12\int_0^\infty \alpha^2\Psi(\alpha)\,d\alpha .
\] Let \(t=(\sqrt{1+\alpha}-1)/(\sqrt{1+\alpha}+1)\). Then \(\alpha=4t/(1-t)^2\), \(d\alpha=4(1+t)(1-t)^{-3}dt\), and \(\Psi(\alpha)=(1-t)^\nu\). Since nu\textgreater6, \[
 m_3=32\{\mathrm B(3,\nu-6)+\mathrm B(4,\nu-6)\}
 =\frac{64\nu}{(\nu-6)(\nu-5)(\nu-4)(\nu-3)}
 =\frac{18944}{1224323425}.
 \tag{PM7}
\] The beta reductions use ordinary convergent integrals and positive rational gamma recurrences; no analytic continuation is used.

Write \(M_j=\int b^{-j}\,d\widetilde\mu(b)\) for j=0,1,2,3, and put \[
 \overline M_3=m_3+2^{-42}
 =\frac{83316594330677601}{5384631367784149811200}.
 \tag{PM8}
\] Then \(M_3\le\overline M_3\), \(M_j<2\) for all four j, and \[
 \overline M_3>2^{-17},\qquad
                         (8\overline M_3)^{168}<2^{-2180}.
 \tag{PM9}
\] The last two statements are exact integer comparisons. The factor8 in PM9 will compare with K109's parameter2; it is not an additional loss in PM1.

\subsection{4. Anisotropic embeddings without changing the original ideal}\label{anisotropic-embeddings-without-changing-the-original-ideal}

Fix \(b_\tau\ge2\), let \(\mathcal B=(\prod_\tau b_\tau)^{1/f}\), and put \(t_0=Uz^2\). On one scalar ideal L use the real metric \(\sum_\tau b_\tau|v_\tau|^2\). Multiplication by the positive diagonal square-root matrix changes its real covolume by \(\prod_\tau b_\tau=\mathcal B^f\), and its dual is the inverse diagonal transform of the original half-trace dual \(L^\#=2\mathfrak D_E^{-1}\bar L^{-1}\).

For a nonzero vector in that dual, the original algebraic norm minimum and arithmetic-geometric mean give \[
 \sum_\tau b_\tau^{-1}|\xi_\tau|^2
 \ge \frac{f}{D_0^{2/f}z^2\mathcal B}.
 \tag{PM10}
\] This inequality holds for arbitrary positive embedding weights, without asserting that the diagonally transformed lattice is an ideal or that different ideals are homothetic.

Let \(\mathcal K\) be any proved scalar coefficient no larger than K109 TG3 such that \[
 \Theta_{L^\#}(\pi^2z^2/A)\le1+\mathcal K A^f
                         \quad\hbox{for every }A>0
\] by the same minimum-based finite packing bound. Repeating that bound with PM10 and applying Poisson gives the scalar S Gaussian sum \[
 \sum_{v\in L}e^{-U\sum b_\tau|v_\tau|^2}
 \le\frac{\pi^f}{D_0}
       \{(t_0\mathcal B)^{-f}+\mathcal K\}
 \le\frac{\pi^f}{D_0}
       \{(2t_0)^{-f}+\mathcal K\}.
 \tag{PM11}
\] All embedding factors and the dimension336 scalar minimum remain explicit. The second inequality uses only \(\mathcal B\ge2\).

For the three true transverse values, the original Gram satisfies \(G'>c_0I\), \(c_0=2^{36}\), and det G'=U, as proved in K109 TG5. The same diagonal scaling in each embedding retains this inequality for its dual. Thus its Gaussian sum is bounded by \[
 \frac{\pi^{3f}U^f}{D_0^3t_0^{3f}}
 (\prod_\tau b_\tau)^{-3}
 \left[1+\mathcal Kc_0^{-f}t_0^f\prod_\tau b_\tau\right]^3.
 \tag{PM12}
\] The actual primal covolume and projection denominator U are unchanged. In particular no arithmetic variable has been divided by a real mixing parameter.

\subsection{5. Absorb the entire mixture, including every transverse correction}\label{absorb-the-entire-mixture-including-every-transverse-correction}

Put \(\rho=\mathcal K c_0^{-f}t_0^f\). For the present dense range k\textgreater=161, the original energy ceiling gives \(t_0\le R_0/k<1\). The exact K109 coefficient satisfies \[
                 0<\rho\le\mathcal K_{\rm old}c_0^{-f}<2^{-3400}.
 \tag{PM13}
\] Integrating PM12 over the product measure \(\widetilde\mu^{\otimes f}\) gives exactly, apart from its displayed constant prefactor, \[
 \sum_{j=0}^3\binom3j\rho^j M_{3-j}^{\,f}.
\] Using PM8--9 and \(M_j<2\), \[
 \begin{split}
 \sum_{j=0}^3\binom3j\rho^j M_{3-j}^{\,f}
 &\le\overline M_3^{\,f}
       +2^f(3\rho+3\rho^2+\rho^3)\\
 &<\overline M_3^{\,f}+7\,2^f2^{-3400}
 <2\,\overline M_3^{\,f}.
 \end{split}
 \tag{PM14}
\] For the last comparison, \(\overline M_3^f>2^{-2856}\), whereas the added term is below \(2^{-3229}\). The zero-vector terms of the transverse dual theta, every positive theta correction and the mass of the added atom are all included.

Apply PM1--2 and PM6 to actual nonzero rows before enlarging any lattice sum. The right side is nonnegative, so Tonelli permits summation and integration in either order. Pull the uniform S bound PM11 outside that integral, retain all original coordinate positions, and then use PM14. The resulting full weighted row bound is \[
 {164\choose k}\frac{2U^f\pi^{(k+3)f}}{D_0^{k+3}}\,
       \{(2Uz^2)^{-f}+\mathcal K\}^{k}
       (Uz^2)^{-3f}\overline M_3^{\,f}.
 \tag{PM15}
\] This replaces the K109 full weighted row factor by at most \((8\overline M_3)^f<2^{-2180}\), with the same admissible scalar coefficient \(\mathcal K\). The remaining selected prefactor is unchanged. Neither the field-norm lower bound nor the actual pointwise alpha ceiling is asserted for the extra zero or out-of-ball lattice tuples; they were used only on the original set before this positive enlargement.

\subsection{6. Full-count interface}\label{full-count-interface}

For any specified original coordinate cell with k\textless=162, the full K109 finite, central and source factors remain unchanged. The scalar coefficient \(\mathcal K\) is fixed and positive. Hence, with \[
 w(q)=\frac{q^f}{q^f+(2R_0)^f\mathcal K},\qquad
 \epsilon(q)=\frac{750R_0/q}{H/2+750R_0/q},
\] the normalized q derivative remains \[
                  -164+kw(q)-\epsilon(q).
 \tag{PM16}
\] Restoring the exact central rank-zero correlation gives uniform decrease at least163-k. For k161 or162 this is at least2 or1, respectively; the entire original finite-state product is bounded by the corresponding checked product \(\prod_G(1+q_G^{-(163-k)})^2\). The full coupled x derivative is still \[
                    -9/2-kw(q_*)+\epsilon(q_*)<-7/2.
 \tag{PM17}
\] The multiplier PM9 is independent of every ideal, q and x, so no mixture, endpoint, pattern or denominator-ideal enumeration is added.

For clarity, the complete carried arithmetic factors being used are \[
 N_EJ=z^{2f}T^{52}/x,\qquad
 |a_G(h)|\le2Q^{-335\cdot167}(N_EJ)^{335/2},
\] \[
 \#d_{\rm cent}\le
 \frac{V_F(H/2+750Uz^2)U^{-f}T^{52}x^{-1}}{N_FJ_0},
 \quad N_EJ\le Q^{167}N_FJ_0,\quad (T-1)^{-168}<2T^{-168},
\] where \(V_F(C)=((22/7)fC^2)^{f/2}/((f/2)!\sqrt{D_F})\). The central count is one actual affine coset with its prescribed real-part shift; the full750 cover and original U denominator survive. The common whole-C paired orientation, denominator intersections and reconstruction from eta,row and central scalar are unchanged. The two mixed projectors retain their modulus-one bounds. The T expression comprises26 active weighted labels and142 inactive pointwise factors on the retained set.

For fixed x, \(z_*=(Q^{334}T^{-52}x)^{1/(2f)}\), \(q_*=M^2/z_*^2\) is the all-rank-zero endpoint, and \((N_FJ_0)^{-1}\le Q^{-167}(q/q_*)^f\). Applying PM16 after this correlation sums all four local states2,1,0 with multiplicities1,2,1. This proves the complete product bound stated above, rather than choosing a separate endpoint for each pattern. PM17 then differentiates the complete endpoint including x\^{}\{-337/2\}. The retained unique eta for each support makes this a supremum over x, with no hidden denominator-ideal sum.

The actual support, original pivot, all168 T coefficients, mixed projectors, central affine shift, full-real cover750, complete Q/T/U/R powers and exact chronological arrays are retained below. No coefficient is treated as constant on an enlarged central orbit.

\subsection{7. Complete fresh k162 endpoint and chronology}\label{complete-fresh-k162-endpoint-and-chronology}

The application uses only k=162, disjoint from DO115's k161. Its scalar input is the frozen DO115 DC3--12 minimum-based cubic-cap bound, with n=336 and \[
 c_n=168{336\choose168}/4^{168},\quad t=49/1000,\quad
 M=369723123/12079,
\] \[
 \mathcal K_{\rm cub}=
 \frac{336M}{c_n}\left(\frac{951}{500}\right)^{168}168!
 \left(\frac{(365729/500)^2}{168(314159/100000)^2}\right)^{168}.
 \tag{PM18}
\] DO115 DC9 proves this for every Euclidean lattice with the stated minimum, at every positive Gaussian parameter; DC10--12 then gives the exact coefficient. This generality applies directly after PM10's real diagonal transformation. An ideal-only fixed-A bound would not suffice. The exact arithmetic also checks K\_cub\textless K\_old. Neither DO115's larger uniform energy coefficient nor its sharper upper bound for pi is imported.

At k162 the complete four-state G sum costs less than2, since its exponent is163-k=1 and \((1+g_0^{-1})^{668}<2\), with g0=22285932761. Keep the actual support bound2\^{}101, original pivot2\^{}8, Green2, this pattern2 and the full168-place T correction2. Their combined exponent is112. The exact old central correlation N\_EJ\textless=Q\^{}167 N\_FJ0 and its moving affine cover are retained, so the central real-volume ceiling remains479 at the endpoint. The full physical constants are \[
 f=168,\quad D_0=7^{140}29^{144}5^{84}13^{84},\quad
 Q=1427908193729121522601,\quad U=77Q^2-14,
 \quad R_0=17424/133.
\] The genuine eligible denominator floor is X0=71443/64. Its coupled cap obeys \[
 q_*^f=\frac{R_0^fT^{52}}{U^fQ^{334}2^{X_0}},\qquad
 4119905<q_*<4119906,
 \quad H/2+750R_0/q_*<68.
 \tag{PM19}
\] The exact packet T is read from the original J56 host certificate. Every positive and negative power is retained before directed logarithms. With \(L_T=1019333/1024\), \(L_Q=287843/4096\), \(L_U=601356/4096\), \(L_R=28810/4096\), the full common exponent is \[
 C=112+2f-929f/2+479+167fL_R-(335f/2)L_U
     +8594L_T-336\cdot167L_Q-(337/2)X_0
   =211125067/512.
 \tag{PM20}
\] This includes the original source, central, relative-phase and original support factors from K109. The T exponent is335\emph{26+52-168=8594, the Q exponent is335}167+167, and the denominator exponent is337/2. Nothing is removed by replacing a finite signed coefficient with a constant on a full central orbit.

Put d=165, c\_-=4119905, c\_+=4119906. The rational Gaussian endpoint is \[
 W=\frac{2{164\choose162}(22/7)^{165f}}
 {D_0^{165}c_-^{167f}}
 \left[(c_+/(2R_0))^f+\mathcal K_{\rm cub}\right]^{162}
 (c_+/(2R_0))^{3f}.
\] Exact integer arithmetic gives \[
 W<2^{-409309},\qquad
 2^{616}<162^{84},\qquad (8\overline M_3)^f<2^{-2180}.
 \tag{PM21}
\] Subtract the saturated-basis336, complete improved trace335, actual radial649, all-coordinate616 and mixture2180 exactly once. The complete endpoint is \[
 E=C-409309-336-335-649-616-2180
       =-548533/512<-1071.
 \tag{PM22}
\] Both continuum reductions PM16--17, all patterns, all supports and the unique-eta/one-label accounting were already included. There is one coordinate cell, so its fresh absolute charge is1071, without another coordinate union. The executable accepts the actual canonical DO115 predecessor or a later complete ordered extension, retains both arrays and appends1071 once. The full two error sums remain below \(2^{-230}\); the existing minimum231 and next-smallest233 remain unchanged. No estimate for k163 or164 is asserted: at k163 the uniform163-k pattern exponent is zero, so the product2 used here cannot simply be carried over. Canonical integration remains a separate event after whole reviews.

\subsection{8. Primary and carried source boundary}\label{primary-and-carried-source-boundary}

PM3 was directly checked against the NIST Digital Library of Mathematical Functions, equations10.32.10 and10.39.2, on the positive real domain. Exact primary URLs are \url{https://dlmf.nist.gov/10.32.E10} and \url{https://dlmf.nist.gov/10.39.E2}. The independent elementary Gaussian proof is included above; no complex branch or unstated convergence continuation enters the result. PM4--9 are new positive-integral deductions, not claims quoted from a theta theorem.

The genuine selected density and compact-factor normalization remain DO080/K107/K108 through the actual Nakahama equation5.7/Theorem5.1 model. The scalar dimension, finite packing and Fourier convention remain DO114/K109 and the new DO115 DC3--12 minimum-based cubic bound. The original integral saturated basis, paired denominator intersection, common G orientation, exact affine central lattice and reconstruction remain K103 SC26.5--20, DO098 CF4--13 and K107 SV17--19/26--27. K109 TG1--19 supplies the complete original count. These carried inputs are read as full mathematical interfaces, not inferred from finite diagnostic examples.

The unchanged algebraic-geometric context is SGA3 I4.7.3/4.7.3.1; SGA6 VI Theorem3.3; FGA232 Theorem3.1pp148--149; FGA236 Corollaries1.7/2.3pp229/232; SGA7 I VII2.7/2.9pp169--174 and VIII7.4(b)pp310--311; cote4 scan003 sections8.1.1--8.1.2/scans005--006 sections8.4.1--8.4.3 and cote54 scans005--007. These are inherited context only. The field, polarization, theta line and selected arithmetic realization are unchanged. No new lambda-ring, K-theory, l-adic weight, formal-group or CM-motive comparison is made.

The full signed inequality, remaining coordinate cells, actual selected arithmetic image, prescribed period, effective source and auxiliary construction, Read/specific032, actual DP44 and RH remain separate obligations.

\subsection{Canonical algebraic-geometric context}\label{canonical-algebraic-geometric-context-45}

The positive real mixture is an estimate in the existing selected complex realization; it introduces no arithmetic scalar, field extension, new projector or freely chosen source vector. Original integral projective modules and paired CM ideals retain their exact half-trace dual and true transverse covolume. NIST DLMF equations10.32.10/10.39.2 supply the exact positive-real integral context, and the full normalized elementary proof is included in PM3. Delsarte, Goethals, and Seidel(1977),Theorems3.3/3.6/4.3,Corollaries3.7-3.8 and Examples4.6-4.7,printed367-370,provide the exact primary context; the needed degree-three kernels and full cap/count form are proved explicitly. The previous exact Nakahama(2016)equation5.7/Theorem5.1pp960-961, Gaussian Fourier convention of Miller and Stephens-Davidowitz(2019)equations1.1/2.1, and Rankin(1955)Theorem1(iv) context remain in the directly checked scopes of K108/DO114. Carried geometric anchors: SGA3 I4.7.3/4.7.3.1; SGA6 VI Theorem3.3; FGA232 Theorem3.1pp148-149; FGA236 Corollaries1.7/2.3pp229/232; SGA7 I VII2.7/2.9pp169-174 and VIII7.4(b)pp310-311; cote4 scan003 sections8.1.1-8.1.2/scans005-006 sections8.4.1-8.4.3 and cote54 scans005-007. No new formal-group slope, lambda-ring/K-theory identity, l-adic weight or CM-motive comparison. DP\_MULT Q3e/DP44/DP52/DP54-65 and THETA3 G1/A7 gain this finite reduction; CX\_GRH CX-ALG-024 and GCT\_PLETH T-052 retain their global gates.

\clearpage

\section{AM-DO-116 --- a degree-five cap bound for every scalar Gaussian parameter}\label{am-do-116-a-degree-five-cap-bound-for-every-scalar-gaussian-parameter}

\textbf{Canonical checked integration --- 2026-11-23T07:07:39.176956+00:00.} Complete saved-proof and executable reviews and separate exact diagnostic replays passed. No numerical error charge is introduced. Frozen author history is preserved. AM-DO-116 QC1--18 proves a universal real336-dimensional lattice theta bound for every positive Gaussian parameter from its actual nonzero minimum, using an exact positive degree-five spherical polynomial and the full cap/nonzero-shell argument. The same theorem applies to arbitrary positive embedding weights with the original half-trace minimum. Its exact coefficient K5 satisfies128K5²\textless Kcub², a scalar improvement of more than3.5 bits. No complete arithmetic sector, numerical error charge, support-floor change or independent main-gate movement is asserted. Current119/31 errors, floors26/11 and remainingk163..164 are unchanged.

Author: /root/definite\_oscillator. 23 November 2026 UTC. Controlled pending scalar theorem. Author, independent and root reviews are separate events. No numerical sector or error charge is asserted here.

\subsection{1. Statement and separation from arithmetic applications}\label{statement-and-separation-from-arithmetic-applications}

Fix n=336 and f=168. Put \[
t_5=883/10000,\qquad a_5=53/500,\qquad b_5=59/100000,
\quad\mathcal M_5=396215477657563473/122893087073,
\] \[
c_{336}=168\binom{336}{168}/4^{168},\qquad
C_5=\frac{336\mathcal M_5}{c_{336}}[2(1-t_5)]^{168}168!.
\tag{QC1}
\] \textbf{Universal theorem.} For every Euclidean lattice Lambda of real dimension336, every lower bound lambda\textgreater0 for its nonzero minimum and every s\textgreater0, \[
\Theta_\Lambda(s):=\sum_{v\in\Lambda}e^{-s\|v\|^2}
\le1+C_5(s\lambda^2)^{-168}.
\tag{QC2}
\] The constant does not require Lambda to be an ideal, rectangular, rationally orthogonal or homothetic to any fixed lattice. There is no upper restriction on s or its reciprocal.

For the actual original half-trace scalar ideals with \(N_EL=z^{2f}\), let \[
D_0=7^{140}29^{144}5^{84}13^{84},\quad
r_+=365729/500,\quad\pi_-=314159/100000,
\] \[
\mathcal K_5=C_5\left(\frac{r_+^2}{168\pi_-^2}\right)^{168}.
\tag{QC3}
\] Then \[
\Theta_{L^\#}(\pi^2z^2/A)<1+\mathcal K_5A^{168}
\qquad\hbox{for every }A>0.
\tag{QC4}
\] Writing \(\mathcal K_c\) for the exact DO115 DC11 cubic coefficient, the exact comparison is \[
128\mathcal K_5^2<\mathcal K_c^2,
\qquad0<\mathcal K_5<\mathcal K_c<\mathcal K_{\rm old},
\tag{QC5}
\] where \(\mathcal K_{\rm old}\) is K109 TG3. Thus the scalar coefficient improves by strictly more than3.5 binary bits. This is a comparison of positive coefficients, not by itself a sector-count saving. No arithmetic pattern, affine central fibre, trace multiplier or error ledger is changed. Any application must retain those separate factors and establish its own complete endpoint.

\subsection{2. The positive kernels through degree five}\label{the-positive-kernels-through-degree-five}

For n\textgreater=3 define normalized polynomials P\_l by P0=1,P1=X and \[
(l+n-2)P_{l+1}(X)=(2l+n-2)XP_l(X)-lP_{l-1}(X).
\tag{QC6}
\] The first six have value1 at X=1. In particular \[
P_2=\frac{nX^2-1}{n-1},\quad
P_3=\frac{(n+2)X^3-3X}{n-1},
\] \[
P_4=\frac{(n+2)(n+4)X^4-6(n+2)X^2+3}{(n-1)(n+1)},
\] \[
P_5=\frac{(n+4)(n+6)X^5-10(n+4)X^3+15X}{(n-1)(n+1)}.
\tag{QC7}
\] For each l=0,\ldots,5, the kernel \(P_l(\langle v,w\rangle)\) on the unit sphere in real n-space is positive semidefinite. Here is its normalization and proof.

Let H\_l be the finite-dimensional real space of homogeneous harmonic polynomials of degree l, restricted to the unit sphere and equipped with the probability-area inner product. The restriction is injective, since a homogeneous polynomial vanishing on the sphere vanishes everywhere. H\_l is nonzero: the real part of \((x_1+ix_2)^l\) supplies an example. It is invariant under the orthogonal group. For an orthonormal basis Y\_j set \[
K_l(v,w)=\sum_{j=1}^{h_l}Y_j(v)Y_j(w),\qquad h_l=\dim H_l.
\] This is a Gram kernel. It is independent of the orthonormal basis and invariant under simultaneous orthogonal transformations. Hence its diagonal is constant; integrating the diagonal gives \(K_l(v,v)=h_l\).

For fixed unit v, K\_l(v,w) is a homogeneous harmonic polynomial in w, invariant under the stabilizer of v. On the unit sphere it therefore depends polynomially only on \(u=\langle v,w\rangle\), with degree at most l. The elementary Laplacian calculation for \(r^lp(u)\), where r=\textbar\textbar w\textbar\textbar{} and u=\(\langle v,w\rangle/r\), gives \[
(1-u^2)p''(u)-(n-1)up'(u)+l(l+n-2)p(u)=0.
\tag{QC8}
\] Indeed the stabilizer-invariant polynomial is a sum of powers of the axial coordinate and the squared orthogonal radius; the same formula follows by differentiating these terms. This equation has a unique polynomial solution of degree at most l up to scalar: if a nonzero solution had degree j\textless l, its leading coefficient would be multiplied by \(l(l+n-2)-j(j+n-2)>0\), a contradiction; the remaining coefficients are then determined recursively by its leading one. Direct substitution shows that QC6--7 supply the solution with value1 at u=1 for l\textless=5. It follows that \[
K_l(v,w)=h_lP_l(\langle v,w\rangle).
\tag{QC9}
\] The positive factor h\_l proves the asserted positive semidefiniteness. This is the addition-formula normalization of Delsarte, Goethals, and Seidel(1977), Theorem3.3, with \(P_l=Q_l/Q_l(1)\). The preceding argument supplies the required positivity directly; no sign of an orthogonal-polynomial coefficient is presumed from a plot.

\subsection{3. Exact degree-five polynomial and finite spherical occupancy}\label{exact-degree-five-polynomial-and-finite-spherical-occupancy}

At n336 use the exact polynomial \[
F_5(X)=(X-t_5)(X^2+a_5X+b_5)^2.
\tag{QC10}
\] Its expansion \(F_5=\sum_{l=0}^5 f_lP_l\) has the six coefficients

\begin{longtable}[]{@{}rl@{}}
\toprule\noalign{}
l & f\_l \\
\midrule\noalign{}
\endhead
\bottomrule\noalign{}
\endlastfoot
0 & 122893087073/354900000000000000 \\
1 & 67872507/1062500000000 \\
2 & 862570663/714000000000 \\
3 & 3887273/171000000 \\
4 & 27930223/229840000 \\
5 & 22579/23256 \\
\end{longtable}

Every coefficient is strictly positive. The exact companion checks the polynomial identity coefficient by coefficient; it does not merely compare sample values. For comparison the power-basis coefficients from degree0 through5 are \[
-3073723/10^{14},\ -668529/62500000000,\
-607033/625000000,\ -15759/2500000,\ 1237/10000,\ 1.
\tag{QC11}
\] If m unit vectors have pairwise inner products at most t5, the off-diagonal F5 terms are nonpositive, because the remaining factor in QC10 is a real square. QC9 and f0\textgreater0 give \[
f_0m^2\le\sum_{i,j}F_5(\langle v_i,v_j\rangle)
\le mF_5(1).
\] Therefore \[
m\le\frac{F_5(1)}{f_0}
=\frac{(1-t_5)(1+a_5+b_5)^2}{f_0}
=\mathcal M_5.
\tag{QC12}
\] This is the elementary linear-programming argument of Delsarte et al., Theorem4.3, applied to the specific polynomial above. No optimality of the chosen rational parameters is claimed or needed.

\subsection{4. Euclidean-ball lift and the universal theta inequality}\label{euclidean-ball-lift-and-the-universal-theta-inequality}

Let X be any finite 1-separated subset of the n-dimensional ball of radius rho\textgreater=1. Its lift \[
x\longmapsto\rho^{-1}(x,\sqrt{\rho^2-\|x\|^2})\in S^n
\] has pairwise inner products at most \(s_\rho=1-1/(2\rho^2)\ge1/2\). Choose the spherical cap with pole height h and angle theta defined by \[
h^2=\frac{s_\rho-t_5}{1-t_5},\qquad
\sin^2\theta=\frac1{2\rho^2(1-t_5)}.
\tag{QC13}
\] The fixed t5 satisfies0\textless t5\textless1/3, so \[
h^2-s_\rho^2
=\frac{(1-s_\rho)[s_\rho-t_5(1+s_\rho)]}{1-t_5}>0.
\] Consequently h\textgreater s\_rho. A point at the cap pole cannot coexist with another cap point. For any two other cap points, with pole coordinates u,v\textgreater=h, the normalized projected inner product in the n-dimensional pole complement is at most \[
\frac{s_\rho-uv}{\sqrt{(1-u^2)(1-v^2)}}.
\] Its derivatives in u,v are negative, because their respective signs are those of \(us_\rho-v\) and \(vs_\rho-u\), both negative. Its maximum is at u=v=h and equals t5. QC12 therefore bounds every cap occupancy by M5. The lift is in n+1 dimensions; the code bound is applied to projections in dimension n.

The normalized cap area satisfies \[
\mu_n(\theta)=c_n\int_0^\theta\sin^{n-1}\phi\,d\phi
\ge\frac{c_n}{n}\sin^n\theta,
\qquad c_n=\frac{\Gamma((n+1)/2)}{\sqrt\pi\Gamma(n/2)}.
\] Averaging the occupancy over the cap pole gives \[
\#X\le\frac{n\mathcal M_5}{c_n}[2(1-t_5)]^{n/2}\rho^n.
\tag{QC14}
\] Apply QC14 to nonzero lattice points after scaling by lambda. Their number N(rho) vanishes for rho\textless1. With b=s lambda², Tonelli gives \[
\Theta_\Lambda(s)-1
=\int_1^\infty2b\rho e^{-b\rho^2}N(\rho)\,d\rho
\le\frac{n\mathcal M_5}{c_n}[2(1-t_5)]^{n/2}
\Gamma(n/2+1)b^{-n/2}.
\tag{QC15}
\] Only the final majorant integral is extended to zero. This proves QC2 with the exact factorial168!, and explicitly excludes and restores the zero vector once. It is valid for every lattice shape and every s\textgreater0.

\subsection{5. Original scalar dual, arbitrary real metric and coefficient comparison}\label{original-scalar-dual-arbitrary-real-metric-and-coefficient-comparison}

The original half-trace dual and norm minimum, retained in DO070 FT3, DO114 SP10 and DO115 DC10--12, are \[
L^\#=2\mathfrak D_E^{-1}\bar L^{-1},\quad
N_EL=z^{2f},\quad
\lambda_1(L^\#)^2\ge\frac{f}{D_0^{2/f}z^2}.
\tag{QC16}
\] The exact integer comparisons \(500^{168}D_0<365729^{168}\) and \(\pi>314159/100000\), with the same Machin-series proof as DO115, imply QC4 from QC2. The original half-trace factor2 is already in QC16; no new conversion is made.

More generally, if positive real embedding weights b\_tau are inserted and \(B=(\prod b_\tau)^{1/f}\), the corresponding inverse-diagonal dual has minimum at least \[
\lambda_b^2\ge\frac{f}{D_0^{2/f}z^2B}.
\tag{QC17}
\] Indeed the original algebraic norm bound for a nonzero dual element is \(\prod_\tau|\xi_\tau|^2\ge D_0^{-2}z^{-2f}\). Applying arithmetic-geometric mean to the positive numbers \(|\xi_\tau|^2/b_\tau\) gives QC17. No nonzero conjugate vanishes. This is also the explicitly normalized argument in K110 PM10. QC2 applies to that transformed Euclidean lattice whether or not it is itself an ideal. It gives the same coefficient K5 with the parameter multiplied by B. This establishes precisely the scalar input needed for a later anisotropic Gaussian mixture; it does not perform that mixture or an arithmetic count.

The exact cubic coefficient of DO115 is \[
\mathcal K_c=\frac{336\mathcal M_c}{c_{336}}[2(1-t_c)]^{168}168!
\left(\frac{r_+^2}{168\pi_-^2}\right)^{168},
\quad t_c=49/1000,\quad\mathcal M_c=369723123/12079.
\] All common normalizations cancel in the quotient: \[
\frac{\mathcal K_5}{\mathcal K_c}
=\frac{\mathcal M_5}{\mathcal M_c}
\left(\frac{1-t_5}{1-t_c}\right)^{168}.
\tag{QC18}
\] The companion proves the strict rational inequality128 times the square of this quotient\textless1 by integer powers. Thus QC5 follows, without numerical logarithms or a floating optimization. The same companion verifies Kc\textless Kold. Positivity and finiteness are explicit.

\subsection{6. Exact verification and primary source scope}\label{exact-verification-and-primary-source-scope}

The executable verifies the normalized recurrence and value at1, the harmonic differential equation for all six polynomials, exact coefficientwise expansion of F5, positivity, its diagonal/cardinality constant, rational cap identities, the dimension336 area and168! normalization, original field/root/pi constants, the complete strict coefficient ratio and the anisotropic minimum direction. Universal positivity and cap counting follow from the written proof, not from finite model checks. The default output lies only in the author's pending directory; an explicit --output supports independent unchanged replay.

Fresh primary scope: Delsarte, P., Goethals, J. M., \& Seidel, J. J. (1977). Spherical codes and designs. \emph{Geometriae Dedicata, 6}(3),363--388. \url{https://doi.org/10.1007/BF03187604}. Definition2.1 and the first five polynomials, Theorem2.4, Theorem3.3/addition formula and Theorem4.3/proof were read directly, printed364--369. Their Q\_l satisfy Q\_l(1)\textgreater0; QC6--9 use P\_l=Q\_l/Q\_l(1) and include the required harmonic normalization proof. No whole-paper reading or assertion of an optimal code bound is made.

The exact cap and half-trace interfaces retain DO115 DC7--12 and DO114 SP3--10; original field constants retain the literal J56 record and DO070 FT3. K110's mixture is cited only for the explicit interpretation of arbitrary embedding weights, and QC17 has its elementary minimum proof stated. This theorem does not require K110's numerical charge. The true integral saturated Gram, paired incidence, central affine lattice, selected trace and full finite-pattern/label sum remain obligations of each separate application.

No canonical ledger, support floor or finite sector is changed by this scalar theorem. In particular it does not by itself remove k163 or164. The original E/F, involution, polarization, theta line, selected complex/arithmetic realization and source data remain unchanged. Carried context is SGA3 I4.7.3/4.7.3.1; SGA6 VI Theorem3.3; FGA232 Theorem3.1pp148--149; FGA236 Corollaries1.7/2.3pp229/232; SGA7 I VII2.7/2.9pp169--174 and VIII7.4(b)pp310--311; cote4 scan003 sections8.1.1--8.1.2/scans005--006 sections8.4.1--8.4.3 and cote54 scans005--007. These are inherited anchors only; no new lambda-ring, K-theory, l-adic weight, formal-group or CM-motive comparison is asserted. The signed kernel, selected arithmetic content, prescribed period, effective source/auxiliary scales, Read/specific032, actual DP44 and RH remain separate.

\subsection{Canonical algebraic-geometric context}\label{canonical-algebraic-geometric-context-46}

The positive real metric estimate leaves original integral Hermitian modules, paired fractional ideals, polarization, theta line, selected complex/arithmetic realizations and source data unchanged. Delsarte, Goethals, and Seidel(1977),Definition2.1,Theorems2.4/3.3/4.3 and Corollary3.8,printed364--369,are the exact primary context; QC6--15 supplies the normalized proof and new explicit polynomial. Carried anchors: SGA3 I4.7.3/4.7.3.1; SGA6 VI Theorem3.3; FGA232 Theorem3.1pp148--149; FGA236 Corollaries1.7/2.3pp229/232; SGA7 I VII2.7/2.9pp169--174 and VIII7.4(b)pp310--311; cote4 scan003 sections8.1.1--8.1.2/scans005--006 sections8.4.1--8.4.3 and cote54 scans005--007. No new formal-group slope, lambda/K-theory identity, l-adic weight or CM-motive comparison. DP\_MULT Q3e,THETA3 G1/A7,CX\_GRH CX-ALG-024 andGCT\_PLETH T-052 retain their global gates.

\clearpage

\section{AM-K-111 --- a degree-fifteen positive kernel and uniform scalar theta}\label{am-k-111-a-degree-fifteen-positive-kernel-and-uniform-scalar-theta}

\textbf{Canonical checked integration --- 2026-11-23T07:14:13.149211+00:00.} Complete saved-proof and executable reviews and separate exact diagnostic replays passed. No numerical error charge is introduced. Frozen author history is preserved. AM-K-111 HP1--18 proves the universal real336-dimensional lattice-minimum/all-positive-parameter scalar theorem using a fixed exact degree-fifteen polynomial with16 strictly positive normalized Gegenbauer coefficients. Its full harmonic positivity, cap projection, nonzero-shell integral and original half-trace/anisotropic normalization give K15²\textless2\^{}-21 Kcub². This is only a scalar coefficient improvement; the positive zero-dual term, all arithmetic patterns and entire row endpoint remain required. No error charge or standalone main-gate movement. Current119/31 arrays,floors26/11 and actual residualk163..164 remain unchanged.

Author: /root/kummer\_extension. 23 November 2026 UTC. Controlled pending scalar theorem. Author, independent and root verification are separate events. No frozen predecessor is changed.

\subsection{1. The exact scalar statement and its boundary}\label{the-exact-scalar-statement-and-its-boundary}

Put n=336, f=168 and \[
 t=2159/10000,\qquad
 c_n=168{336\choose168}/4^{168}.
 \tag{HP1}
\] Let \[
\begin{split}
 p(X)=10^{-12}(&-32766+383921X+47138047X^2+34060017X^3\\
 &-7736258949X^4-23081334653X^5
       +261518505090X^6+10^{12}X^7),
\end{split}
\] \[
                         F(X)=(X-t)p(X)^2 .
 \tag{HP2}
\] The polynomial is fixed and rational. A floating search suggested the parameters; no optimization claim or floating inequality is used in the theorem.

Define \[
 \mathcal M_{15}
 =\frac{6015684926187167549292201834533348561577984}
        {2618017117694612446410293519},
\quad
 C_{15}=\frac{336\mathcal M_{15}}{c_n}
                    (7841/5000)^{168}168!.
 \tag{HP3}
\]

\textbf{Universal theorem.} For every full-rank Euclidean lattice Lambda of real dimension336 and every positive lower bound lambda for its nonzero minimum, at every s\textgreater0, \[
       \Theta_\Lambda(s):=\sum_{v\in\Lambda}e^{-s\|v\|^2}
       \le 1+C_{15}(s\lambda^2)^{-168}.
 \tag{HP4}
\] For the actual J56 half-trace scalar ideal, put \[
 D_0=7^{140}29^{144}5^{84}13^{84},\quad
 r_+=365729/500,\quad \pi_-=314159/100000,
\] \[
 \mathcal K_{15}
 =C_{15}\left(\frac{r_+^2}{168\pi_-^2}\right)^{168}.
 \tag{HP5}
\] For every fractional ideal L with \(N_EL=z^{336}\), its original dual satisfies, uniformly in its ideal class and denominators, \[
 \Theta_{L^\#}(\pi^2z^2/A)<1+\mathcal K_{15}A^{168}
                         \qquad(A>0).
 \tag{HP6}
\] The same universal minimum-based theorem applies to any positive real diagonal transform; no transformed lattice is asserted to remain an ideal.

If K\_c is the complete DO115 DC11 cubic coefficient, then \[
                    (\mathcal K_{15}/\mathcal K_c)^2<2^{-21}.
 \tag{HP7}
\] This is a comparison of positive scalar coefficients. It is not a claim that an entire row polynomial or a sector contribution gains10.5 bits per coordinate: the positive zero-dual term, original q range, full pattern sum, selected trace, central fibre and exact numerical endpoint must still be retained. This theorem appends no error charge and does not itself remove a coordinate cell.

\subsection{2. Positive normalized kernels through degree fifteen}\label{positive-normalized-kernels-through-degree-fifteen}

For n=336 define P0=1, P1=X and \[
 (\ell+n-2)P_{\ell+1}
 =(2\ell+n-2)XP_\ell-\ell P_{\ell-1}\qquad(\ell\ge1).
 \tag{HP8}
\] For every \(0\le\ell\le15\), these have degree ell, value1 at1, and satisfy \[
 (1-X^2)P_\ell''-(n-1)XP_\ell'
               +\ell(\ell+n-2)P_\ell=0.
 \tag{HP9}
\] The exact executable verifies HP8--9 coefficient by coefficient, including degrees zero and one. These are finite rational polynomial identities, not numerical tests at sample points.

Here is the complete positive-kernel argument. Let H\_ell be the real homogeneous harmonic polynomials of degree ell, restricted to the unit sphere, with its probability-area inner product. Restriction is injective by homogeneity. The space is nonzero, since the real part of (x1+i x2)\^{}ell is harmonic, including the constant1 at ell0. It is finite-dimensional and orthogonally invariant. Choose an orthonormal basis Y\_j and form \[
                       K_\ell(v,w)=\sum_jY_j(v)Y_j(w).
\] This is a Gram kernel, independent of basis and invariant under simultaneous rotations. Its diagonal is constant; integrating it gives \(K_\ell(v,v)=h_\ell=\dim H_\ell>0\).

For fixed unit v, the kernel is a homogeneous harmonic polynomial in w invariant under the stabilizer of v. Such an invariant polynomial is a sum of powers of the axial coordinate and squared orthogonal radius. On the sphere it is therefore a polynomial in \(u=\langle v,w\rangle\) of degree at most ell. Computing the Laplacian of \(r^\ell q(u)\) gives the differential equation HP9. Its polynomial solution of degree at most ell is unique up to scalar: a leading degree j\textless ell would have nonzero coefficient multiplier \(\ell(\ell+n-2)-j(j+n-2)>0\); at degree ell the lower coefficients are determined recursively. HP8--9 and P\_ell(1)=1 consequently identify \[
                  K_\ell(v,w)=h_\ell P_\ell(\langle v,w\rangle).
 \tag{HP10}
\] Thus every displayed normalized P\_ell is a positive semidefinite kernel. In particular, for any finite unit-vector list, the full double sum of P\_ell is nonnegative. The ell0 sum is the square of the cardinality.

This is the normalization of the addition formula of Delsarte, Goethals, and Seidel (1977), Theorem3.3, with \(P_\ell=Q_\ell/Q_\ell(1)\). The direct harmonic argument is included. The same general argument is also recorded in the pending DO116 QC6--9 proof; the present theorem supplies the independent degree0--15 identities and does not require a degree-five numerical extrapolation.

\subsection{3. Exact polynomial expansion and finite occupancy}\label{exact-polynomial-expansion-and-finite-occupancy}

The exact expansion \(F=\sum_{\ell=0}^{15}f_\ell P_\ell\) has the following coefficients.

\begin{longtable}[]{@{}
  >{\raggedleft\arraybackslash}p{(\columnwidth - 2\tabcolsep) * \real{0.5714}}
  >{\raggedright\arraybackslash}p{(\columnwidth - 2\tabcolsep) * \real{0.4286}}@{}}
\toprule\noalign{}
\begin{minipage}[b]{\linewidth}\raggedleft
ell
\end{minipage} & \begin{minipage}[b]{\linewidth}\raggedright
f\_ell
\end{minipage} \\
\midrule\noalign{}
\endhead
\bottomrule\noalign{}
\endlastfoot
0 & 2618017117694612446410293519/5064668570357760000000000000000000000000000 \\
1 & 25642602737559917343694363/235522161940000000000000000000000000000 \\
2 & 49270993895458521006571454987/9989484359680000000000000000000000000000 \\
3 & 27264576833619334674831717023/159600900326400000000000000000000000000 \\
4 & 1319180964541869111389213974237/396668904320000000000000000000000000000 \\
5 & 32422895056018777837235070320113/550707661440000000000000000000000000000 \\
6 & 2453081830638655671834891939307919/3650162534400000000000000000000000000000 \\
7 & 504734173631194480428057501773284769/66499470553600000000000000000000000000000 \\
8 & 29659547606497059902226044525712687247/528955788632064000000000000000000000000000 \\
9 & 376732329958280660264997999744761310193/851641657015800000000000000000000000000000 \\
10 & 1514931689182620170634314702329194069897/703858802892800000000000000000000000000000 \\
11 & 1678258246987540890198121713891102687/135640414720000000000000000000000000000 \\
12 & 3633660155162743492753565397042642811/97115505868800000000000000000000000000 \\
13 & 32165265258571248320217928094677439/202083832668160000000000000000000000 \\
14 & 475007414177162440250677/2009673474048000000000000 \\
15 & 30931304169353/41810448138240 \\
\end{longtable}

Every numerator and denominator is positive. HP8, HP2 and this table give an exactly checkable polynomial identity; the executable reconstructs all power-basis coefficients and checks every table entry. This is not an inference from positivity at a finite set of points.

If m unit vectors have pairwise inner products at most t, every off-diagonal F term is nonpositive, since the second factor in HP2 is a real square. HP10 and f0\textgreater0 therefore give \[
 f_0m^2\le\sum_{i,j}F(\langle v_i,v_j\rangle)\le mF(1).
\] Thus \[
                             m\le F(1)/f_0=\mathcal M_{15}.
 \tag{HP11}
\] The exact numerator/denominator in HP3 is obtained from \(F(1)=(1-t)p(1)^2\). No claim that the degree-fifteen polynomial is optimal is made or needed.

\subsection{4. The cap argument for all Euclidean lattices}\label{the-cap-argument-for-all-euclidean-lattices}

Let X be a finite1-separated subset of the n-ball of radius rho\textgreater=1. Its lifted unit vectors in real n+1-space are \[
                    x\longmapsto\rho^{-1}
                           (x,\sqrt{\rho^2-\|x\|^2}).
\] Before division by rho, distances do not decrease. Their pairwise inner products are therefore at most \(s_\rho=1-1/(2\rho^2)\in[1/2,1)\). Choose the cap with pole height h and angle theta satisfying \[
 h^2=\frac{s_\rho-t}{1-t},\qquad
 \sin^2\theta=\frac1{2\rho^2(1-t)}.
 \tag{HP12}
\] Our exact t satisfies0\textless t\textless1/3. Hence \[
 h^2-s_\rho^2
 =\frac{(1-s_\rho)[s_\rho-t(1+s_\rho)]}{1-t}>0,
\] so h\textgreater s\_rho. A cap point at its pole cannot coexist with another cap point. Otherwise, normalize the projections onto the n-dimensional pole complement. For two pole coordinates u,v\textgreater=h the projected inner product is bounded above by \[
 \frac{s_\rho-uv}{\sqrt{(1-u^2)(1-v^2)}}.
\] Its derivative in u has the sign of u s\_rho-v\textless0, and likewise in v. Its maximum at u=v=h is exactly t. HP11 bounds every cap occupancy by M15. The code dimension remains336, although the spherical lift has ambient dimension337.

The normalized cap area is \[
 \mu_n(\theta)
 =c_n\int_0^\theta\sin^{n-1}\phi\,d\phi
 \ge(c_n/n)\sin^n\theta.
\] The inequality follows by inserting cos(phi)\textless=1 in the nonnegative integral. Averaging occupancy over all cap poles gives the full finite bound \[
                \#X\le\frac{n\mathcal M_{15}}{c_n}
                          [2(1-t)]^{n/2}\rho^n .
 \tag{HP13}
\] The constant c\_n in HP1 is the exact gamma quotient \(\Gamma((n+1)/2)/[\sqrt\pi\Gamma(n/2)]\) for n336.

For a full-rank lattice of minimum lambda, the nonzero counting function N(rho) in the radius-rho*lambda ball vanishes below1. With b=s lambda², nonnegative integration gives \[
\begin{split}
 \Theta_\Lambda(s)-1
 &=\int_1^\infty2b\rho e^{-b\rho^2}N(\rho)\,d\rho\\
 &\le\frac{n\mathcal M_{15}}{c_n}[2(1-t)]^{n/2}
        \Gamma(n/2+1)b^{-n/2}.
\end{split}
 \tag{HP14}
\] Only the final majorant integral was enlarged to radius0. The zero vector is excluded and restored exactly once. The exact gamma factor is168!, proving HP4 at every s\textgreater0. No lattice-shape or small-Gaussian-parameter assumption appears.

\subsection{5. Actual half-trace and anisotropic scope}\label{actual-half-trace-and-anisotropic-scope}

For the original field and half-trace metric, DO070 FT3 and K109 TG3--4 give \[
 L^\#=2\mathfrak D_E^{-1}\bar L^{-1},\qquad
 \lambda_1(L^\#)^2\ge fD_0^{-2/f}z^{-2}.
 \tag{HP15}
\] The exact integer comparison \(500^{168}D_0<365729^{168}\) gives \(D_0^{1/f}<r_+\); DO115 DC10's directly verified Machin bound gives pi\textgreater pi\_-. Applying HP4 with s=pi²z²/A proves HP6 with the precise coefficient HP5. The inequalities in r\_+ and pi\_- give the strict upper envelope.

More generally, for any positive real embedding weights b\_tau and geometric mean B, the diagonal transform of the original dual satisfies \[
 \sum_\tau b_\tau^{-1}|\xi_\tau|^2
             \ge fD_0^{-2/f}z^{-2}B^{-1}\quad(\xi\ne0).
 \tag{HP16}
\] This follows from the algebraic norm minimum and AM--GM, as in K110 PM10. HP4 applies directly to this real Euclidean lattice with its true changed covolume. There is no additional arithmetic ideal, coordinate, phase or norm-one multiplicity introduced by this scalar lemma.

The true coupled transverse Gram, original integral saturated input lattice, full S support and all finite projectors remain outside HP4 itself. An application must use the correct primal/dual metrics and all original membership arguments; scalar positivity over the reals implies nothing about a residue-field Hermitian form.

\subsection{6. Exact comparison and no standalone charge}\label{exact-comparison-and-no-standalone-charge}

The common factors in HP5 and DO115 DC11 cancel. With \(\mathcal M_c=369723123/12079\), \[
 \frac{\mathcal K_{15}}{\mathcal K_c}
 =\frac{\mathcal M_{15}}{\mathcal M_c}
                    \left(\frac{7841}{9510}\right)^{168}.
 \tag{HP17}
\] Writing M15=N15/D15 and Mc=Nc/Dc in lowest terms, the exact integer inequality is \[
            2^{21}(N_{15}D_c)^2\,7841^{336}
                         <(D_{15}N_c)^2\,9510^{336}.
 \tag{HP18}
\] This proves HP7 without a logarithm approximation. The complete rational coefficient is retained in the certificate.

A row factor such as \((A^{-168}+\mathcal K)^k\) contains a positive term unaffected by replacing K. Thus HP7 alone cannot justify a uniform2\^{}(-10.5k) row gain. Nor does it justify extending a previous finite-pattern product at k163 or164: those require their separate complete pattern sums. No error ledger is modified here. The support floors, signed sufficient inequality, actual selected arithmetic content, period, effective source, Read/specific032, DP44 and RH remain unchanged.

\subsection{7. Primary and carried sources}\label{primary-and-carried-sources}

The primary reference is Delsarte, P., Goethals, J. M., \& Seidel, J. J. (1977). Spherical codes and designs. \emph{Geometriae Dedicata,6}(3),363--388. \url{https://doi.org/10.1007/BF03187604}. Fresh scope: Theorem3.3/addition formula, Theorem3.6 and Corollaries3.7--3.8 with proofs, Theorem4.3/proof and Examples4.6--4.7, printed367--370. Its Q\_l normalization is related to HP8 by a positive factor. The present proof establishes the required normalization and positivity through degree15 directly.

The exact cap and layer-cake interfaces were freshly whole-reviewed in DO115 DC3--12. The general harmonic construction was also independently discussed with /root/definite\_oscillator in the pending DO116 QC6--9 proof; this theorem proves all degree0--15 identities and its own explicit rational polynomial. The DO116 degree-five coefficient and its count consequences are not dependencies.

The original scalar metric is DO070 FT3; the entire all-parameter Fourier and true-Gram conventions remain K109 TG3--11. The genuine selected trace and original arithmetic factors remain K108/K110, with their exact inherited Nakahama5.7/Theorem5.1 scopes. Carried algebraic-geometric context is SGA3 I4.7.3/4.7.3.1; SGA6 VI Theorem3.3; FGA232 Theorem3.1pp148--149; FGA236 Corollaries1.7/2.3pp229/232; SGA7 I VII2.7/2.9pp169--174 and VIII7.4(b)pp310--311; cote4 scan003 sections8.1.1--8.1.2/scans005--006 sections8.4.1--8.4.3 and cote54 scans005--007. These are inherited anchors, not fresh whole-source readings or new comparison claims. This scalar theorem constructs no new source model or algebraic-geometric comparison. Finite exact diagnostics verify rational identities and constants; the universal lattice and positive-kernel proofs are HP8--16.

\subsection{Canonical algebraic-geometric context}\label{canonical-algebraic-geometric-context-47}

The polynomial and cap proof estimate the existing selected real metric and leave every original integral module, paired ideal, polarization, theta line and arithmetic source fixed. Delsarte,Goethals,andSeidel(1977),Definition2.1,Theorems2.4/3.3/4.3 andCorollary3.8,printed364--369,are the exact primary context; HP8--16 proves the normalized degree0--15 and all-minimum/all-parameter form explicitly. Carried anchors: SGA3 I4.7.3/4.7.3.1; SGA6 VI Theorem3.3; FGA232 Theorem3.1pp148--149; FGA236 Corollaries1.7/2.3pp229/232; SGA7 I VII2.7/2.9pp169--174 and VIII7.4(b)pp310--311; cote4 scan003 sections8.1.1--8.1.2/scans005--006 sections8.4.1--8.4.3 and cote54 scans005--007. No new formal-group slope, lambda/K-theory identity, l-adic weight or CM-motive comparison. DP\_MULT Q3e,THETA3 G1/A7,CX\_GRH CX-ALG-024 andGCT\_PLETH T-052 retain their global gates.

\clearpage

\section{AM-LIFT-094 --- Effective algebraic Hecke escape on a fixed indefinite component}\label{am-lift-094-effective-algebraic-hecke-escape-on-a-fixed-indefinite-component}

\textbf{Canonical checked integration --- 2026-11-23T07:20:56.098184+00:00.} Complete saved-proof and executable reviews and separate exact diagnostic replays passed. No numerical error charge is introduced. Frozen author history is preserved. AM-LIFT-094 HE0-22 proves conditional effective algebraic Hecke escape on a fixed indefinite arithmetic family. For a reduced algebraic bad locus of total degreeD and dimension\textless=s, finite equal-ample-pullback correspondences of degree\textless=delta with no nonempty proper closed forward-invariant locus admit escape in at mostB\_delta,s(D) steps, where A\_delta(D)=2 delta\^{}D D\^{}2, B\_delta,0(D)=D and B\_delta,s(D)=D+B\_delta,s-1(A\_delta(D)). For actual projective O\_E lattice rankr\textgreater=3 indefinite at every real place, f={[}F:Q{]}, an explicitly cleared rational Witt frame gives t=1+3McQ\_S, two physical original-M-level Hecke correspondences with delta=t\^{}(4fr\^{}2) and multiplier t\^{}2. Every forward orbit is dense; the conditional accumulatedN\textless=t\^{}B is obtained. The actual regular LIFT085 period-family map and nonempty inverse image of the prescribed nonvanishing open remain hypotheses; H,D,c,N0 and admitted Host/Sizes costs are not evaluated. Same original fields, quotient/action/line/level/frame conditions and independent entry multiplierN0 remain. No selected arithmetic nonzero content, global period, specific032/Read/STiesR, DP44 or RH conclusion. No finite-error charge;119/31 and floors26/11 unchanged.

Author: /root/lift\_accounting. 23 November 2026 UTC.

\textbf{Conditional theorem; final author verification is recorded separately.} This theorem has two precise conclusions. First, finite algebraic correspondences with equal ample pullbacks and no nonempty proper closed invariant subset admit the explicit degree/dimension escape recursion HE4. Second, on any fixed unitary arithmetic domain of rank at least three that is indefinite at every real place, two explicitly constructed correspondences satisfy that invariant-subset condition. Their common physical isogeny multiplier and covering degree are

\[
 t=1+3Mc\prod_{p\in S}p,\qquad \delta=t^{4fr^2},
 \qquad N\le t^{B_{\delta,s}(D)}.\tag{HE0}
\]

Here c clears the stated finite list of rational Witt-frame matrices in the \textbf{actual} lattice; D is the degree of the actual pulled-back bad locus in a chosen very ample Baily--Borel power; and s bounds its dimension. These are explicit inputs, not values supplied by the literal Sizes clauses. The source application uses the algebraic period-family map already proved in LIFT085, not an assumed isomorphism between the literal moduli scheme and a PEL quotient. It requires the same rule to be nonzero on that family. An entry isogeny from a different integral component contributes its separate multiplier N\_0. No unconditional effective source N, new nonzero arithmetic rule, finite error charge, or DP44 conclusion is claimed. An analytic Cayley period map is never treated as a bounded-degree polynomial.

\subsection{1. The precise algebraic setting}\label{the-precise-algebraic-setting}

Let X be a nonempty reduced quasi-projective complex variety, with a projective closure Xbar and a very ample line bundle H on Xbar. Its irreducible components need not have the same dimension. For a reduced closed Z contained in X, define deg(Z) as the sum of the H-degrees of the projective closures of all its irreducible components. Thus every component has positive integral degree and their number is at most deg(Z).

Fix finitely many diagrams

\[
 X\xleftarrow{p_i}C_i\xrightarrow{q_i}X.\tag{HE1}
\]

Both maps are finite étale and surjective. Suppose each diagram has a projective compactification with finite maps pbar\_i,qbar\_i to Xbar, whose inverse images of X are exactly C\_i, and with

\[
 \bar p_i^*H\simeq\bar q_i^*H.\tag{HE2}
\]

Each finite étale map p\_i has degree at most an integer delta\textgreater=1 on every connected component of its target, counting all points in its generic fibre. We use all branches of each correspondence. The forward orbit of x consists of points reached by finitely many choices y in q\_i(p\_i\^{}\{-1\}(x)); the empty word is allowed. Assume

\[
 \text{there is no nonempty proper closed }W\subset X
 \text{ with }q_i(p_i^{-1}W)\subset W\text{ for every }i.
 \tag{HE3}
\]

For example, density of every forward orbit implies HE3. Density of just one seed orbit is not silently substituted for HE3.

For delta,D\textgreater=1 set

\[
 A_\delta(D)=2\delta^D D^2,\qquad
 B_{\delta,0}(D)=D,\qquad
 B_{\delta,s}(D)=D+B_{\delta,s-1}(A_\delta(D))\quad(s\ge1).
 \tag{HE4}
\]

These are explicit integers. They are monotone in delta,D,s. Set the escape bound for the empty set to zero.

\textbf{Abstract assertion.} If Z is a proper closed subset of X, deg(Z)\textless=D and dim(Z)\textless=s, every x in X has a forward path of length at most B\_\{delta,s\}(D) ending outside Z.

The bound concerns existence of a branch and a word. It does not assert that every branch escapes, or that a traced Hecke operator has nonzero value.

\subsection{2. A projective separating hypersurface}\label{a-projective-separating-hypersurface}

We use the following elementary degree fact, with a proof to make the only degree input explicit.

\textbf{Separation lemma.} If Zbar is a reduced closed subset of projective space of total component degree D and w is outside Zbar, there is a homogeneous polynomial f of positive degree at most D, vanishing on Zbar and nonzero at w. If Zbar is empty, take f=1.

First let V be an irreducible component of dimension r and degree d.~Choose a linear subspace L through w, of projective dimension n-r-1, disjoint from V. Such an L exists: intersect r+1 general hyperplanes through w, successively avoiding the finitely many components that remain. Since w is not in V, the hyperplanes have no base point on V, and r+1 such hyperplanes have empty intersection with V. If r=n-1, V itself is a hypersurface of degree d and already gives the assertion.

Otherwise choose a hyperplane C in L not containing w. Projection with center C maps V to a closed subset of P\^{}\{r+1\}. The projection is finite on V: each fibre is the intersection of V with an affine space, is also proper, and is therefore finite. Its image has dimension r, hence is a hypersurface. Its degree is at most d, by the finite projection degree formula. The projected point of w lies outside this image, because its fibre is L and L misses V. Pulling back the equation of that image gives a polynomial of degree at most d that vanishes on V but not at w.

Perform this construction for every irreducible component of Zbar and multiply the resulting polynomials. The sum of their degrees is at most D. This proves the separation lemma, including reducible and non-equidimensional Zbar.

It follows in particular that, for any irreducible V' not contained in Zbar, a polynomial of degree at most D can be chosen vanishing on Zbar but not identically on V': choose w in V' minus Zbar.

The basic finiteness of projection is also the argument in Stacks Project, Section 43.23 (Tag 0B1N), Lemma 43.23.1. Its proof uses that a proper affine fibre is finite. The degree comparison used here is the proper pushforward/line-bundle projection formula, which does not require a nonsingular compactification; see Lemma 42.26.4, Tag 02SU.

\subsection{3. The finite graph of largest-dimensional components}\label{the-finite-graph-of-largest-dimensional-components}

Let r=dim(Z), and let Z\_1,\ldots,Z\_m be its components of dimension r. Thus m\textless=D. Add one exit vertex, denoted \emph{. For each i and j, every irreducible component A of p\_i\^{}\{-1\}(Z\_j) maps surjectively to Z\_j under p\_i, since p\_i is finite étale. Its q\_i-image is irreducible, closed and of dimension r. If that image is contained in Z, it equals one of Z\_1,\ldots,Z\_m; add the corresponding directed edge. Otherwise add an edge to }.

Every vertex Z\_j reaches * in at most m edges. Indeed, if the reachable vertices from Z\_j contained no exit, the union of their components would be a nonempty proper closed subset of X stable under every branch of every correspondence. This contradicts HE3. A shortest path to * repeats no component vertex, so its length is at most m.

Fix one such path for each Z\_j. Fibre products of its correspondence diagrams, together with the chosen component branches, give an irreducible variety A\_j and finite maps

\[
 p^{(j)}:A_j\longrightarrow Z_j,\qquad
 q^{(j)}:A_j\longrightarrow X,\qquad \ell_j\le m,
 \tag{HE5}
\]

where the first map is surjective, of degree e\_j\textless=delta\^{}\{ell\_j\}, and the second image is not contained in Z. To justify the composition, the selected branch maps are finite and dominant onto the intermediate component. Their generic fibres are finite field extensions of its function field. Their tensor product is nonzero; choosing a prime and taking its closure in the fibre product gives an irreducible component dominating both selected branches. Iterate this construction. The generic degree over the initial component is at most the product of the branch degrees. The final image remains the prescribed exit image because all selected maps are dominant. The individual irreducible component of a finite étale cover of a possibly singular component is not assumed itself to be étale; only the original covers are used as finite étale maps.

The same fibre products of projective compactifications give projective closures of these chosen branches with finite endpoint maps. The line equality HE2 telescopes to

\[
 (\bar p^{(j)})^*H\simeq(\bar q^{(j)})^*H.\tag{HE6}
\]

No open image is replaced by a closed image without the finiteness hypothesis.

\subsection{4. One explicit dimension-lowering cut}\label{one-explicit-dimension-lowering-cut}

Choose, by Section 2, a homogeneous polynomial f\_j of degree d\_j\textless=D vanishing on Zbar but not identically on the closure of q\^{}\{(j)\}(A\_j). Its pullback is a nonzero section of ((qbar\^{}\{(j)\})\^{}*H)\^{}\{d\_j\} on the irreducible projective closure Abar\_j. Let Y\_j be the intersection with X of the image under pbar\^{}\{(j)\} of its zero set.

Then Y\_j is closed in X, contained in Z\_j, and has dimension at most r-1. For r=0 it is empty. Projective properness makes the image closed, and finiteness preserves the dimension of every nonempty component. Since the pulled-back section is not identically zero, its zero divisor on Abar\_j has no component of dimension r.

For r\textgreater=1 its effective divisor degree, the line identity HE6 and the projection formula give

\[
 \begin{aligned}
 \deg_H(Y_j)
 &\le d_j\deg_{(\bar q^{(j)})^*H}(\bar A_j)\\
 &=d_j\deg_{(\bar p^{(j)})^*H}(\bar A_j)\\
 &=d_j e_j\deg_H(\overline{Z_j})
 \le D\delta^D\deg_H(\overline{Z_j}).
 \end{aligned}\tag{HE7}
\]

Boundary components, multiplicities and finite image degrees only enlarge the positive effective-cycle degree on the right; none is discarded in a direction that could underestimate Y\_j. The first-Chern-class projection formula is valid for these possibly singular projective closures, by Tag 02SU.

Let Z\_low be the union of the components of Z of dimension strictly less than r, and put

\[
 Y=Z_{\rm low}\cup\bigcup_{j=1}^mY_j.\tag{HE8}
\]

Then Y is closed, dim(Y)\textless=r-1, and

\[
 \deg(Y)\le D+\delta^D D^2\le A_\delta(D).\tag{HE9}
\]

Define Bad\_m(Z) to consist of points for which every path of length at most m remains in Z. We claim

\[
                      Bad_m(Z)\subset Y.\tag{HE10}
\]

The empty path first places such a point in Z. If it is in Z\_low there is nothing to prove. Otherwise choose a largest-dimensional component Z\_j containing it. The finite surjection p\^{}\{(j)\} has a point above it. Every endpoint of that selected branch is reached by a path of length ell\_j\textless=m; hence its q\^{}\{(j)\}-image lies in Z, and f\_j vanishes there. The original point therefore belongs to Y\_j. This proves HE10. We have not asserted that Y equals Bad\_m(Z), or used an existential branch condition where a universal one is needed.

\subsection{5. Escape induction and the actual multiplier consequence}\label{escape-induction-and-the-actual-multiplier-consequence}

For r=0, Section 3 already proves escape within m\textless=D steps. For r\textgreater=1 apply the induction hypothesis to Y, with dimension at most r-1 and degree at most A\_delta(D). This gives a path from any point to a point outside Y in at most B\_\{delta,r-1\}(A\_delta(D)) steps. By HE10, a point outside Y has at least one path of length at most m\textless=D ending outside Z. Concatenate the paths. This proves HE4 and the abstract assertion.

Here is the precise conditional isogeny consequence. Suppose every branch in the chosen finite family is an actual OE-linear polarized isogeny of the same source moduli problem, with a positive integer n\_i\textless=L satisfying

\[
 u_i^*\mathcal L_{i+1}\simeq\mathcal L_i^{n_i^2},\quad
 n_i\equiv1\pmod M,\quad (n_i,\prod_{p\in S}p)=1,
 \quad u_i\eta_i=\eta_{i+1},\tag{HE11}
\]

including the rigidified cubical line and exact finite level. Assume the bound is uniform over the branch, not merely over one representative. For a path of length ell, composition gives the identical conditions with

\[
 N=\prod_{i=1}^{\ell}n_i\le L^\ell
       \le L^{B_{\delta,s}(D)}.\tag{HE12}
\]

The same source LIFT085 polarization-adjoint argument gives ker(u) contained in A\_0{[}N\^{}2{]}. Its all-object descent, and LIFT090--093's same-field degree estimates, then apply to this N under their original hypotheses. The quantitative assertion is conditional on HE1--3, HE11 and the actual degree D; it does not require a polynomial pullback through an analytic period map.

\subsection{6. Exact interface with the run and the unresolved inputs}\label{exact-interface-with-the-run-and-the-unresolved-inputs}

Write c\_src for the fixed source coefficient record, to distinguish it from the integer frame-clearing input c in Sections 7--10. For LIFT085 the actual bad locus is the closed complement of

\[
 U_{F_0}=\pi(P_{f_{F_0}})\subset X= X_o.XF(c_{\rm src}),\tag{HE13}
\]

with F\_0 the same given nonzero natural rule. The frame polynomial degree in HF24--26 bounds changes of frame over an already chosen point; it is not a degree bound for HE13 in an unspecified projective embedding of X.

The literal \nolinkurl{Moduli} declaration in \nolinkurl{PushInvertible.lean}, lines 29--38, supplies a quasi-projective smooth scheme, its universal object, and the exact representation property. It supplies no numerical embedding degree, no specified very ample line or power, and no finite Hecke generating set. The literal \nolinkurl{sec} in \nolinkurl{IsCombAt.lean}, lines 98--101, supplies naturality and Borel weight conditions; those clauses alone do not identify the numerical degree of HE13 with the physical matrix degree D of HF24.

The primary Hecke framework is compatible with HE1--2 in a suitably identified Shimura model: Daw, Javanpeykar and Kühne (2020), Sections 3.1--3.2, give the finite Hecke diagrams and equality of pullbacks of the Baily--Borel line bundle. Their Proposition 10.1 uses the same projection/degree cutting mechanism. Its hypotheses involve a positive-dimensional special subvariety and are not asserted for our CM seed. The present abstract proof is provided separately and makes no appeal to that theorem to force an orbit or a point to escape.

The remaining source application requires the following concrete data, with proofs.

\begin{longtable}[]{@{}
  >{\raggedright\arraybackslash}p{(\columnwidth - 2\tabcolsep) * \real{0.5000}}
  >{\raggedright\arraybackslash}p{(\columnwidth - 2\tabcolsep) * \real{0.5000}}@{}}
\toprule\noalign{}
\begin{minipage}[b]{\linewidth}\raggedright
Required input
\end{minipage} & \begin{minipage}[b]{\linewidth}\raggedright
Current status
\end{minipage} \\
\midrule\noalign{}
\endhead
\bottomrule\noalign{}
\endlastfoot
An actual period-family map Phi:Y→X and a numerical very ample power H on the Baily--Borel compactification of Y & HF12--16 supplies the regular family map; no isomorphism with X is needed. A numerical very-ampleness bound is not supplied here. \\
A finite family satisfying HE3 on a chosen arithmetic family & Sections 7--9 prove this for two correspondences on a fixed indefinite quotient. They do not enter other integral components. \\
Finite endpoint extensions and exact HE2 for that H & Section 9 proves them on the auxiliary arithmetic quotient, with a chosen unevaluated very ample power. \\
A numerical D for the inverse image of the actual same-rule bad locus HE13 & Unevaluated. The frame-only degree HF24 is insufficient. \\
A numerical delta and a branch-uniform L satisfying all of HE11 & Sections 7--10 give delta=t\^{}(4fr²) and L=t once the actual frame denominator c is supplied; component entry is separate. \\
Nonzero F\_0 & Still an explicit hypothesis. This proof constructs no selected arithmetic nonzero rule. \\
\end{longtable}

Density in one connected component cannot justify HE3 on a disconnected source space. Either a finite correspondence family connects all relevant components with explicit cost, or the theorem must be applied on a component on which the actual rule is proved nonzero, with the cost of an actual seed in that component included. Similarly, the q=0 compact terminal situation is not given a positive-dimensional domain argument.

Thus the theorem supplies a fully explicit conditional recursion. Sections 7--10 make its finite-family constants explicit from an actual Witt frame on one indefinite family. The actual degree and entry data are still needed before it is an effective source N or a proof of any admissible numerical window.

\subsection{7. A fixed indefinite rational component}\label{a-fixed-indefinite-rational-component}

We now prove the finite-family and density inputs on one fixed arithmetic domain, with its original rational Hermitian form and lattice. The argument does not supply the degree of the actual pulled-back rule locus or the cost of entering this family from a different integral component.

Let E/F be a CM extension, {[}F:Q{]}=f, and let (V,h) have E-dimension r\textgreater=3 and signature (p\_tau,q\_tau) with p\_tau,q\_tau\textgreater=1 at every real place. Let Lambda be an actual full projective OE-lattice of rank r in V. Write \(\mathcal D\) for the product of the corresponding unitary Hermitian symmetric domains. The theorem first concerns its own congruence quotient Y. In the source application, Y has the regular map \(\Phi:Y\to X_\mathbb C\) constructed by LIFT085 HF15--16. We use the closed inverse-image bad locus on Y. No isomorphism, surjectivity onto a component, or finite degree is assumed for \(\Phi\).

The rational Hermitian space contains a hyperbolic plane. Indeed, if Delta\_h is its determinant class in F\^{}\emph{/N(E\^{}}), choose a Hermitian space of dimension r-2, signatures (p\_tau-1,q\_tau-1), and determinant class -Delta\_h. The required sign condition is exactly (-1)\^{}(q\_tau-1), since sign(-Delta\_h)=(-1)\^{}(q\_tau+1). Existence and uniqueness in the Hermitian classification then identify V with H orthogonal-sum A. Deligne (1982), Proposition 4.1, printed page 28 of the official TeX transcription, states precisely this existence/uniqueness by dimension, signatures and determinant class. The chosen isometry need not be integral, and its denominators are retained below. No rank-two or definite assertion is made by this argument.

Choose this E-rational Witt frame. The Gram matrix is

\[
 J=\begin{pmatrix}0&0&1\\0&B&0\\1&0&0\end{pmatrix},
 \qquad B=B^*,\quad\det B\ne0.\tag{HE14}
\]

Let P\_+,P\_- be the E-linear projections to its first and last lines. For v in A and z in E with bar(z)=-z put

\[
 n(v,z)=\begin{pmatrix}
 1&-v^*B&z-\tfrac12v^*Bv\\
 0&I&v\\0&0&1
 \end{pmatrix}.\tag{HE15}
\]

Direct multiplication gives n(v,z)\^{}*J n(v,z)=J and determinant one. For rational scalar b, n(bv,0)=exp(bX(v)), where X(v)\^{}3=0, and n(0,bz)=exp(bZ(z)). The opposite matrices are obtained by interchanging the first and last coordinates. Fix Q-bases of A and E\^{}-=\{z:bar(z)=-z\}. Choose a positive integer c such that cP\_+,cP\_- are in End\_OE(Lambda), cX(v) and its opposite lie in 2 End\_OE(Lambda) for every selected basis vector, and cZ(z) and its opposite lie in End\_OE(Lambda). A common integer exists because the list is finite and all maps are rational. This c is an explicit input, not a bounded output of the source Sizes hypotheses.

\subsection{8. Congruence, all elementary root directions, and density}\label{congruence-all-elementary-root-directions-and-density}

Let M be the prescribed finite level and S a finite set of protected rational primes; in the source application take S to include the primes dividing M, as in LIFT085. Put L\_0=3M, Q\_S=product\_\{p in S\}p (empty product one), and

\[
 t=1+L_0cQ_S,\qquad
 a=tP_++(1-P_+-P_-)+t^{-1}P_-.\tag{HE16}
\]

Thus t\textgreater1, t=1 mod L\_0, gcd(t,Q\_S)=1, and a lies in SU(h)(F). Let Gamma be the subgroup of SU(h)(F) preserving Lambda and congruent to one modulo L\_0 Lambda. The matrices n(L\_0cv,0), n(0,L\_0cz), and their opposites belong to Gamma. For example the horizontal exponential is I+L\_0cX+(L\_0cX)\^{}2/2; the choice cX in 2 End makes its square term integral and its difference from I divisible by L\_0. The inverse has the same property. This explicitly supplies every basis direction of both opposite unipotent radicals, including their central directions.

Both ta and ta\^{}\{-1\} preserve Lambda. Indeed

\[
 ta=tI+(t-1)(tP_+-P_-),\qquad
 ta^{-1}=tI+(t-1)(tP_--P_+).\tag{HE17}
\]

At every prime not dividing t, a and a\^{}\{-1\} preserve the local lattice and are congruent to one modulo L\_0: use (t-1)P\_\pm=L\_0Q\_S(cP\_\pm). In particular all protected local lattices and the prescribed M-level remain fixed. Any additional finite congruence conditions would have to be included in L\_0 and in the definition of Gamma; they are not silently supplied by membership in S.

The arithmetic group Gamma is neat. On the underlying integral Z-lattice every gamma has the form I+3T. In an algebraic closure of Q\_3 each eigenvalue is 1+3b with b integral, because b is an eigenvalue of T. Every product of integer powers of these eigenvalues therefore lies in 1+3 O\_bar(Q\_3). Its only root of unity is one: a root of unity of order prime to 3 has nontrivial residue unless it is one; a primitive 3\^{}e-th root has valuation v\_3(zeta-1)=1/(2*3\^{}(e-1))\textless1. This is the neatness condition, not merely absence of finite-order matrices. Passing to this auxiliary level is for the proof; the final objects are forgotten back to their original M-level.

Conjugation satisfies

\[
 a n(v,z)a^{-1}=n(tv,t^2z).\tag{HE18}
\]

Consequently the closure of the group generated by Gamma,a,a\^{}\{-1\} contains the full real one-parameter subgroups in every chosen horizontal and central direction: powers and conjugates have parameters L\_0c Z{[}1/t{]} or L\_0c Z{[}1/t\^{}2{]}, both dense in R. A Q-basis of A, respectively E\^{}-, is a real basis after restriction of scalars and tensoring with R. Thus the closure contains both entire real opposite unipotent radicals, independently at every real place.

These radicals generate the full real Lie algebra su(p\_tau,q\_tau). Here is a direct check that avoids omitting any root direction. In the complexification, X(v) as v varies, together with X(iv), separates the elementary matrices E\_\{1j\} and E\_\{jr\} for every middle index j, since B is invertible. The opposite radical similarly supplies E\_\{j1\} and E\_\{rj\}. Their commutators supply every middle E\_\{ij\}, E\_\{1r\}, E\_\{r1\}, and all diagonal differences. They therefore span sl\_r(C). The real generated algebra is contained in su(p\_tau,q\_tau), whose complexification is sl\_r(C), so it has the full real dimension and equals su(p\_tau,q\_tau). The independent real-place directions prove the same statement for the product. Since SU(p\_tau,q\_tau) is connected, the closure contains their full product, which acts transitively on \(\mathcal D\). Connectedness follows, for example, from the Cartan decomposition: its maximal compact S(U(p\_tau) times U(q\_tau)) is connected, and the noncompact factor is a vector space.

Every orbit of the group generated by Gamma and \(a^{\pm1}\) is consequently dense in \(\mathcal D\). Each a lies in the identity real group, so this is a statement on the same connected domain and the same rational Hermitian space. It does not move a seed to a different adelic component.

\subsection{9. Exact finite indices, algebraicity, and the invariant-subset bridge}\label{exact-finite-indices-algebraicity-and-the-invariant-subset-bridge}

For \(a\), put \(\Gamma_a=\Gamma\cap a^{-1}\Gamma a\). The two maps are \(p([I])=[I]\) and \(q([I])=[aIa^{-1}]\) from \(\Gamma_a\backslash\mathcal D\) to \(Y=\Gamma\backslash\mathcal D\). The second correspondence is defined with a inverted. These formulas also fix the branch order used below. If gamma is congruent to one modulo t\^{}2 Lambda, write gamma=I+t\^{}2T. Then

\[
 a\gamma a^{-1}=I+(ta)T(ta^{-1})\tag{HE19}
\]

is integral, as is its inverse. At the primes dividing L\_0, conjugation preserves the L\_0 congruence by HE17 and gcd(t,L\_0)=1. Hence the principal t\^{}2-congruence subgroup of Gamma lies in Gamma intersect a\^{}\{-1\}Gamma a; the identical argument works with a inverted. Both covering degrees are at most

\[
 \delta=t^{4fr^2}.\tag{HE20}
\]

Indeed Lambda/t\^{}2 Lambda is projective of rank r over OE/t\^{}2 OE and hence free over this finite Artinian ring. Its endomorphism ring has cardinality \textbar OE/t\^{}2 OE\textbar\^{}\{r\^{}2\}=t\^{}\{4fr\^{}2\}. The index of a reduction-image subgroup is bounded by that entire cardinality. No global free OE-basis of Lambda is assumed.

For this neat arithmetic Hermitian quotient the diagrams are algebraic finite étale maps: finite-index inclusions give finite topological covers of the nonsingular quotient, hence their algebraic maps have isomorphic tangent maps. The exact compactification input is Klingler--Yafaev, Proposition 5.3.2(1), Remark 5.3.4, and the arithmetic-lattice extension immediately after (5.1), printed pages 24--25 of the author version. It gives extensions to the Baily--Borel compactifications and equality of the pullbacks of their canonical ample automorphic lines. It applies to arithmetic lattices in a semisimple group of Hermitian type, so SU need not itself be presented as a full Shimura datum. The group here has no real compact factors, by the indefinite hypothesis.

These extended maps are finite. They are proper, and their pullback of an ample line is ample. A positive-dimensional fibre would contain a curve on which that pullback has degree zero, contradicting ampleness. Proper quasi-finite morphisms are finite. Their inverse images of the open quotient are exactly the open sources. Indeed, for either extended map, its open source C is already finite, hence proper, over Y. Its inclusion into the inverse image of Y is therefore proper (the latter is separated over Y). This inclusion is also open and dense; properness makes its image closed, so it is the whole inverse image. Apply this on every connected component. This argument retains the entire open cover and requires no assertion that a boundary point may map into the interior.

Choose a positive power H of the target's ample automorphic line that is very ample. Its two pullbacks agree. The same choice works for the correspondences of a and its inverse and for every iterated fibre product in Sections 3--4. Thus HE1--2 hold on Y with HE20, without any claim about an unknown compactification of the literal source X. The chosen power and the degree of the bad locus with respect to it remain unevaluated.

The inverse image under p of {[}I{]} consists of the classes of gamma I gamma\^{}\{-1\} for gamma in Gamma, modulo Gamma\_a. Its q-images are the classes of a gamma I gamma\^{}\{-1\} a\^{}\{-1\}. Thus forward paths using the two correspondences are exactly the images of words in Gamma and \(a^{\pm1}\); elements of Gamma do not cost a step on the quotient. Choosing all branches is essential. Their density just proved implies HE3: a nonempty algebraically closed forward-invariant subset contains the entire forward orbit of any of its points, and is analytically closed, so density forces it to be the whole component. This supplies the missing bridge rather than assuming HE3 from qualitative density of a larger, unspecified family.

\subsection{10. Branch-uniform physical isogeny multiplier}\label{branch-uniform-physical-isogeny-multiplier}

Use the exact algebraic period family over Y supplied by LIFT085 HF15--16. Its fixed lattice is Lambda and a domain point is a compatible complex structure I. The integral linear map ta, or ta\^{}\{-1\}, is complex-linear from I to aIa\^{}\{-1\}, respectively a\^{}\{-1\}Ia, and scales the Riemann form by t\^{}2. The same is true for every branch representative gamma\_1 a\^{}\{+/-1\} gamma\_2, since Gamma preserves the original lattice and polarization. All such maps are OE-linear. Their action on the prescribed M-level is the identity because t=1 mod M and every gamma and a are congruent to one there. Thus their positive physical multiplier is the same t, uniformly over all branches. The positive Hermitian metrics also satisfy \(H_{\rm target}(duv,duw)=t^2H_{\rm source}(v,w)\): combine the t\^{}2 alternating-form identity with the transported complex structure. Therefore no extra archimedean operator norm of an integral branch representative is inserted into this exact relative metric identity.

For the actual canonical trace-one line used in LIFT085, functoriality of the rigidified Poincare bundle gives the exact rigidified cubical isomorphism u\^{}*L\_target=L\_source\^{}\{t\^{}2\}; this is its HF5--7 identity, not an identification of numerical polarization classes alone. This step requires the literal period-family/trace-one-line identification from LIFT085. For a different source line its corresponding exact pullback condition remains a hypothesis.

Let F\_0 be the same nonzero literal natural rule and let U\_\{F\_0\} be its actual nonvanishing open from HF11. Assume precisely that

\[
 W=\Phi^{-1}(U_{F_0})\ne\varnothing,\qquad
 Z=Y\setminus W,\qquad \deg_H Z\le D,\quad\dim Z\le s.\tag{HE21}
\]

Because Phi is regular, Z is a proper closed subset of Y. The theorem uses its degree on this auxiliary quotient, including all irreducible components. For the uniform source signature its ambient dimension is f p q, so one may use s=f p q-1 when Z is nonempty. The recursion gives

\[
 N\le t^{B_{t^{4fr^2},s}(D)}.\tag{HE22}
\]

For a seed represented by a point of Y, every step is an actual OE-isogeny between fibres of that same period family. Composition therefore produces an actual object in U\_\{F\_0\}, and the source nonzero-frame argument in HF9--11 and Section 7 of LIFT085 gives a nonzero literal evR over the same torsion field. Its generic field-degree bound is at most N\^{}\{8g\^{}2\}, g=fr; the stronger same-seed field results apply only when their original J/C\_0 hypotheses hold. The auxiliary 3M marking is used only to prove the statement on Y and is forgotten. It does not force adjoining extra 3-torsion or roots of unity: descent is applied to the final original-level object and its kernel in A\_0{[}N\^{}2{]}, exactly as in HF23 and Section 7.

If a controlled seed first enters this family through an actual isogeny with multiplier N\_0, the total multiplier is at most N\_0 times HE22. That entry isogeny must separately satisfy OE-linearity, exact finite level, the rigidified line-power identity, and the protected-prime conditions. Its multiplier is not bounded by density inside the family. Taking an arbitrary target x in U\_\{F\_0\} and its HF16 family does ensure HE21, but it does not bound the controlled seed's entry into that family.

The remaining numerical/application inputs are now more precise: the actual Hermitian frame c, the very ample automorphic power H, the true pulled-back same-rule bad-locus degree D, the nonzero-rule/family condition HE21, and the entry multiplier N\_0 when needed. The finite-family size is two and its degree/multiplier bounds are HE20/HE16 once these exact frame data are given. No actual Sizes window, effective admitted N, terminal definite case, arithmetic selected nonzero rule, or DP44 conclusion is asserted.

\subsection{11. Literal-source, algebraic-geometric and verification boundary}\label{literal-source-algebraic-geometric-and-verification-boundary}

The exact source fields are unchanged. \nolinkurl{TwoWeightEqOneAddRoot.lean} has Host14--27 (including finite unramifiedness and the trace-one w\_0), Datum30--47 (actual projective lattice, alternating form and signatures), Level53--64, SAB72--84 (including every adj identity and prim), TateIso131--142, and Obj144--160. The finite level is the literal PinLevel43 in \nolinkurl{PartitionExtendByZeroAntitone.lean}, with A=1; PinDat47 and PinField45 are not rebuilt by this theorem. \nolinkurl{PushInvertible.lean} gives Moduli29--38, Coeff40--45, and XF49. Its representation clause is what supplies the regular Phi from the algebraic family. No field replacement or equality of a new coefficient ring with what2.O is inferred.

The natural-rule and frame fields are \nolinkurl{IsCombAt.lean}, FrPt67--77, FrMap79--85, sec100 and evR103, with the test/naturality/default-zero conventions in \nolinkurl{AlgebraMapCAt.lean}, TRing9--12, IsNatural35--36, FromTest51--52, evalAt54--58 and EvalWD84--87. They are used through the checked HF9--11 proof. The same actual rule is evaluated; a convenient nonzero section of the same weight is not substituted.

In the inherited Grothendieck language, the physical line identity is an identity of rigidified line bundles constructed from the normalized relative Poincare biextension; the quotient is an actual abelian variety over its finite torsion field; the tangent eigenframes are bases of actual vector bundles. No virtual K-group equality, motivic comparison, weight/purity statement or new Hodge-cycle algebraicity result is used. At protected primes, any p-divisible-group isomorphism still has the original conditional good-model/extension premise of LIFT085 Section 9. Changes of theta section, selected derivatives, integral eigenlines, invariant costs and Step5 hypotheses are not removed by this existence statement.

The diagnostic script checks the finite graph lemma in exhaustive small models, the recursion's arithmetic and boundary cases, exact rational Hermitian/unipotent/conjugation/lattice identities, the elementary Lie-generator span, and the stated literal fingerprints. Such finite diagnostics do not prove the Hermitian classification, Baily--Borel algebraicity or real density; their complete arguments/imports are given above. The certificate is not a numerical evaluation of H, D, c or N\_0. The final saved-file author and independent reviews are separate records.

\subsection{References and reading scope}\label{references-and-reading-scope}

Daw, C., Javanpeykar, A., \& Kühne, L. (2020). Effective estimates for the degrees of maximal special subvarieties. \emph{Selecta Mathematica, 26}, Article 2. \url{https://doi.org/10.1007/s00029-019-0528-1}. Freshly read Sections 3.1--3.2, 9.1--9.2 and Proposition 10.1 with its proof; other sections are not asserted as a fresh whole reading.

The Stacks Project Authors. (2026). \emph{The Stacks Project}, Section 43.23 (Tag 0B1N), Lemma 43.23.1, \url{https://stacks.math.columbia.edu/tag/0B1N}; Lemma 42.26.4 (Tag 02SU), \url{https://stacks.math.columbia.edu/tag/02SU}. Both displayed statements and proofs were freshly read. The latter is the general first-Chern-class/proper-pushforward formula, not only the nonsingular-variety projection formula.

Deligne, P. (1982). Hodge cycles on abelian varieties. In \emph{Hodge cycles, motives, and Shimura varieties}. Official TeX transcription, Proposition 4.1 and Corollary 4.2, printed pages 28--29, \url{https://www.jmilne.org/math/Documents/Deligne82.pdf}. The Hermitian classification statement and proof were freshly read; no motivic consequence is used here.

Klingler, B., \& Yafaev, A. (2014). The André--Oort conjecture. \emph{Annals of Mathematics, 180}(3), 867--925. Author version, \url{https://www2.mathematik.hu-berlin.de/~klingleb/papiers/KY1309.pdf}. Proposition 5.3.2(1), Remark 5.3.4, and the arithmetic-lattice extension after (5.1), printed pages 24--25, were freshly read. Only these unconditional compactification/line-bundle statements are imported; the paper's conditional André--Oort results are not used.

\subsection{Canonical algebraic-geometric context}\label{canonical-algebraic-geometric-context-48}

The new argument is algebraic degree descent on the actual regular period family, with finite Hecke correspondences, equal ample Chern classes and proper pushforward. It preserves the physical CM Picard/Poincare cubical line and its multiplier under original-level isogenies. Mandatory exact locator: Datum in TwoWeightEqOneAddRoot.lean is lines30-49, replacing the isolated30-47 in the frozen proof; the final condition at49 was directly read. The complete immutable locator addendum accompanies this theorem. Fresh primary anchors: Deligne(1982) Proposition4.1/Corollary4.2pp28-29; Klingler-Yafaev(2014) Proposition5.3.2(1),Remark5.3.4 and extension after(5.1)pp24-25; Daw-Javanpeykar-Kuhne(2020) sections3.1-3.2,Proposition10.1 with its distinct positive-dimensional scope; Stacks0B1N Lemma43.23.1 and02SU Lemma42.26.4. Carried geometric anchors from LIFT085: FGA232 Theorem3.1pp148-149, FGA236 Corollaries1.7/2.3pp229/232; SGA7 I VII2.7/2.9pp169-174,VIII7.4(b)pp310-311; cote4 scan003 sections8.1.1-8.1.2/scans005-006 sections8.4.1-8.4.3 and cote54 scans005-007. No new crystalline slope, lambda-ring/K-theory identity, weight or CM-motive comparison is established. DP\_MULT Q3e/DP44/DP52/DP54-65 gain a conditional finite algebraic escape input; THETA3 G1/A7,CX\_GRH CX-ALG-024 and GCT\_PLETH T-052 retain their global gates.

\clearpage

\section{AM-LIFT-094 mandatory locator precision}\label{am-lift-094-mandatory-locator-precision}

Status: locator-only clarification to the frozen final proof; no mathematical change.

In Section 11 of \nolinkurl{effective_algebraic_hecke_escape.md}, read the literal \nolinkurl{Datum} declaration in \nolinkurl{TwoWeightEqOneAddRoot.lean} as \textbf{lines 30--49}, replacing the displayed range 30--47. The declaration starts at line 30, and its final signature, commutation, and trace condition ends at line 49. This exact final condition was included in the complete 168-line source read and in the frozen source-file fingerprint. No hypothesis, formula, proof, executable diagnostic, or conclusion changes.

The source path is \nolinkurl{research/alpoge_source/ReciprocityTheory/ReductiveGroups/ArithmeticGeometry/BirationalGeometry/ToroidalCompactification/TwoWeightEqOneAddRoot.lean}, SHA-256 \nolinkurl{3646386e9cd25287a50eb1ad82456c5b26316207741b8ed7298d46aa4f9d7826}. The exact declaration was directly reread at lines 30--49 for this addendum.

The original frozen proof, script, source manifest, author verification, and independent replay remain unchanged. This addendum is to accompany the delivered theorem and the canonical source-locator record.

Frozen subject fingerprints:

\begin{itemize}
\tightlist
\item
  Proof: \nolinkurl{988caa0818b1089f28687a8036e0599e67b7dc00b1859dc8855d0cf89fc0ec6c}.
\item
  Script: \nolinkurl{7b6d714a46437336222e2d2d43bf729da0c8937eeaec79d91d55c53b020c1dda}.
\item
  Source manifest: \nolinkurl{0333aa171953179de90b49f2ac651aa0a47e7eca1ee74499bfac7ee1b699b445}.
\end{itemize}

Completed UTC: 2026-11-23T07:15:26.830640+00:00

\clearpage

\section{AM-R-094 --- complete polynomial pattern summation and the support27 floor}\label{am-r-094-complete-polynomial-pattern-summation-and-the-support27-floor}

\textbf{Canonical checked integration --- 2026-11-23T07:22:21.994934+00:00.} Complete root and independent proof/script/interface reviews and separate exact replays passed. The actual canonical predecessor has119/31 entries; this integration gives120/31, complete floors27/11. Frozen author history is preserved. AM-R-094 PP1-24 removes the entire remaining conditional J56 all-top support26 class. R091 plus K109/DO115/K110 leaves only actual all-II global-rank-two-unipotent scalar-cross cells k163/164; the full correlated finite-pattern containing polynomial and low-parameter angular mixture give complete exponents -199349/512\textless-389 and -159413/512\textless-311, using DO116 and K111 respectively. Their union is below2\^{}-310. Append310 once to the exact K110119/31 arrays:120/31, both entire sums\textless2\^{}-230, complete floors27/11. The actual remainder now has support at least27 with every earlier higher cut and original finite/real/arithmetic convention retained. PP3, the sufficient signed higher-support inequality, remains open, as do actual Host/Sizes, selected nonzero arithmetic content, global period, specific032/Read/STiesR and evaluatedDP44/RH. No actual shorter mixed-CM chain or universal strict-negative cost theorem is asserted.

Author: /root, GPT-6 Astra Pro; inherited effort not externally exposed; ALP\_MACHINE. 23 November2026 UTC. Controlled pending theorem. Whole saved-proof/script reviews, exact replays and canonical integration are separate events. No earlier frozen theorem is edited.

\subsection{1. Actual scope and theorem}\label{actual-scope-and-theorem-5}

Keep the original conditional J56 all-top kernel, H=133, f=168, n=169, its original positive S of dimension164, all fields and integral Hermitian modules, both CM partners, every admitted source/target pencil and label, the selected real coefficient and metric, central phases, Haar measures and class masses. Apply every checked containing-set removal chronologically through K110 to its intersection with the actual remainder. Let R be what remains. R091 DR1--7 applies to this smaller actual set. Every support26 term therefore has the original and inverse positive167 compressions equal to I, determinant one, nonzero original S cross, globally rank-two-unipotent normal form and all26 active places of type II. The later checked cuts give \[
 v=q_{\rm elig}=26,\quad b=j_{\rm int}=0,\quad
 2^{71443/64}<x=N_ED_\eta<2^{2108},\quad k\in\{163,164\}.
 \tag{PP1}
\] The symbols q\_elig and j\_int here are eligibility and active-integrality counts. Below q denotes a positive real scale and j a polynomial index; these are different variables.

\textbf{Theorem.} The entire actual support26 portion of R has absolute normalized contribution strictly below2\^{}-310. More precisely the full k163 and k164 cells are below2\^{}-389 and2\^{}-311, respectively. Appending310 once to the complete K110119/31 error arrays gives120/31 errors, with both complete sums below2\^{}-230. The all-top support floor becomes27; the uniform floor remains11.

The remaining actual energy has the form \[
 \mathcal E_{\rm top}=1+\operatorname{Re}\sum_{h\in R_{\ge27}}w(h)
                  +\epsilon,\qquad |\epsilon|<2^{-230},
 \tag{PP2}
\] where R\_ge27 retains every previous cut in its proved scope. Its sufficient signed inequality \[
       \operatorname{Re}\sum_{h\in R_{\ge27}}w(h)>-1+2^{-230}
 \tag{PP3}
\] is still open. This is an absolute finite/real remainder theorem for the fixed conditional kernel, not an admitted Host/Sizes constructor, arithmetic nonvanishing theorem, mixed lift chain or evaluated DP44 saving.

\subsection{2. Two all-minimum scalar inputs}\label{two-all-minimum-scalar-inputs}

Use DO116 QC1--18 at k163 and K111 HP1--18 at k164. Both prove for every real336-dimensional Euclidean lattice with nonzero minimum at least lambda and every s\textgreater0 a theta bound1+C(s lambda\^{}2)\^{}-168. This includes arbitrary positive real diagonal changes of metric. Their exact original half-trace scalar coefficients are \[
 \mathcal K_d=\frac{336\mathcal M_d}{c_{336}}
       [2(1-t_d)]^{168}168!
       \left(\frac{(365729/500)^2}{168(314159/100000)^2}\right)^{168},
 \quad c_{336}=168\binom{336}{168}/4^{168},
 \tag{PP4}
\] with \[
 t_5=883/10000,\quad \mathcal M_5=396215477657563473/122893087073,
\] \[
 t_{15}=2159/10000,\quad
 \mathcal M_{15}=
 6015684926187167549292201834533348561577984/
 2618017117694612446410293519.
 \tag{PP5}
\] The full positive-kernel, spherical-cap, nonzero-shell and normalization proofs are in those scalar results. No optimality assumption is used. Both coefficients are positive and smaller than the cubic coefficient of DO115 and K109's original coefficient. The new proof uses the entire positive quantity A\^{}-168+K\_d; it does not infer a coordinatewise saving from the coefficient ratio alone.

Put R0=17424/133, U=77Q\^{}2-14, Q=1427908193729121522601 and D0=7\^{}140 29\^{}144 5\^{}84 13\^{}84. K110 PM3--15 gives a positive majorizing mixture with all mixing parameters at least2, exact inverse third-moment majorant \[
 \overline M_3=83316594330677601/5384631367784149811200,
 \qquad(8\overline M_3)^{168}<2^{-2180}.
 \tag{PP6}
\] Its entire removed low-parameter interval is included in its explicit atom. For either coefficient PP4, its unchanged all-minimum proof applies after the inverse-diagonal transform of K110 PM10. Its complete original transverse correction remains bounded using K\_old and the true Gram comparison G'\textgreater2\^{}36 I. K110 PM13--15 requires k\textgreater=161, which both cells satisfy; it does not require k\textless=162. Only the later pattern argument PM16 had that restriction, and that is replaced below.

In particular the full weighted original row bound, for N\_E L=z\^{}(2f), is \[
 \binom{164}{k}\frac{2U^f\pi^{(k+3)f}}{D_0^{k+3}}
   \{(2Uz^2)^{-f}+\mathcal K_d\}^{k}
   (Uz^2)^{-3f}\overline M_3^{\,f}.
 \tag{PP7}
\] All actual nonzero-coordinate norm and pointwise alpha conditions are used before the nonnegative lattice enlargement. Extra zero or out-of-ball tuples receive no selected trace claim. Every transverse dual correction and zero-vector contribution is included in PP7. The rest of the original selected prefactor is unchanged.

\subsection{3. Original finite states and central correlation}\label{original-finite-states-and-central-correlation}

The complete source/finite interfaces are K103 SC26.5--20, DO098 CF8--15, K107 SV17--20/26--27 and K110 PM15--22. For one common whole-C orientation, \[
 L=JI_{26}D_\eta,\quad N_EJ=z^{2f}T^{52}/x,\quad
 |a_G(h)|\le2Q^{-335\cdot167}(N_EJ)^{335/2}.
 \tag{PP8}
\] There are334 unmixed G places, with residue norms q\_v in\{g0,g1\}, \[
 g0=22285932761,\quad g1=64072175441,\quad
 \prod_v q_v=Q^{167}.
\] At each place the norm exponent e\_v of J is0,1,2 for reduction rank at least2,1,0, with multiplicities1,2,1. The rank-zero central ideal J0 has exponent0,0,1 in those three cases. Thus \[
 N_EJ\le Q^{334},\qquad N_EJ\le Q^{167}N_FJ_0.
 \tag{PP9}
\] This is one common pattern for the entire cross map; neither transverse coordinates nor the pivot select independent patterns.

The actual central fibre is one affine coset with its prescribed real-part shift. Its full-real750 cover and original U denominator give \[
 \#d_{\rm cent}\le
 V_F(H/2+750Uz^2)U^{-f}T^{52}x^{-1}(N_FJ_0)^{-1},
\quad V_F(C)=\frac{((22/7)fC^2)^{f/2}}{(f/2)!\sqrt{D_F}},
\quad D_F=2^fD_0.
 \tag{PP10}
\] No assertion of integrality of the affine shift is made. The167 original covector values, eta and d reconstruct h injectively. For each support eta is unique in PP1's window. All26 active one-label weights and142 inactive pointwise factors remain, with(T-1)\^{}-168\textless2T\^{}-168; both mixed projectors retain their modulus-one estimates.

Define \[
 z_*=(Q^{334}T^{-52}x)^{1/(2f)},\quad
 q=R_0/(Uz^2),\quad q_*=R_0/(Uz_*^2),\quad
 \lambda=q/q_*\ge1.
 \tag{PP11}
\] Then \[
 \lambda^f=Q^{334}/N_EJ=\prod_v q_v^{2-e_v},\qquad
 (N_FJ_0)^{-1}\le Q^{-167}\lambda^f.
 \tag{PP12}
\] The original row energy implies q\textgreater=k. The positive formulas below are evaluated algebraically at q\_* when forming a majorant; no new selected trace is assigned at an artificial endpoint.

\subsection{4. Full polynomial sum, including its two exceptional degrees}\label{full-polynomial-sum-including-its-two-exceptional-degrees}

After the unchanged factors independent of q and the correlation Q\^{}-167 are removed, the relevant positive expression is \[
 F_k(q)=q^{-167f}\{(q/(2R_0))^f+\mathcal K_d\}^{k}
       (q/(2R_0))^{3f}V_F(H/2+750R_0/q).
 \tag{PP13}
\] Put u\_\emph{=(q\_}/(2R0))\^{}f and w\_\emph{=u\_}/(u\_\emph{+K\_d), so0\textless w\_}\textless1. The central volume decreases with q. Therefore PP12 and the binomial theorem give, for each complete pattern, \[
 \lambda^f\frac{F_k(q)}{F_k(q_*)}
 \le\sum_{j=0}^{k}\binom kj(1-w_*)^{k-j}w_*^j
                         \lambda^{(j-163)f}.
 \tag{PP14}
\] All terms are nonnegative. Summing over all four states at every place gives the exact containing-pattern expression \[
 \sum_{j=0}^{k}\binom kj(1-w_*)^{k-j}w_*^j P_j,
 \qquad P_j=\prod_v(1+q_v^{j-163})^2.
 \tag{PP15}
\] It follows directly that \[
 \begin{array}{ll}
 P_j<(1+g0^{-1})^{668}<2,&j\le162,\\
 P_{163}=2^{668},&\\
 P_{164}=Q^{334}\prod_v(1+q_v^{-1})^2<2Q^{334}.&
 \end{array}
 \tag{PP16}
\] This explicitly includes the degree163 term where the old uniform negative power vanished, and the degree164 term where it changed sign.

At the lowest denominator endpoint X0=71443/64, the exact host comparison gives q\_\emph{\textless4119906. The companion proves for both scalar constants \[
 \frac{(4119906/(2R_0))^{168}}{\mathcal K_d}<2^{-200}.
 \tag{PP17}
\] Since q\_} decreases as x increases, w\_\emph{\textless2\^{}-200 throughout PP1. For k163 or164 the degree163 binomial weight is at most164w\_}\^{}163; the degree164 weight, if present, is at most w\_*\^{}164. Also Q\textless2\^{}71. Hence the two exceptional contributions together are less than \[
 164\,2^{668-200\cdot163}
       +2^{1+71\cdot334-200\cdot164}<1.
 \tag{PP18}
\] The remaining binomial weights sum to at most one, so PP15 is strictly below3, and in particular below4. This is a uniform complete finite-pattern bound, not a numerical sample at the all-rank-zero state. We conservatively charge4 rather than optimizing3. It costs one additional binary bit compared with the preceding product2.

\subsection{5. Complete coupled denominator monotonicity}\label{complete-coupled-denominator-monotonicity}

Set \[
 w(q)=\frac{(q/(2R_0))^f}{(q/(2R_0))^f+\mathcal K_d},\qquad
 \varepsilon(q)=\frac{750R_0/q}{H/2+750R_0/q}.
\] Both lie strictly between zero and one. Differentiating the entire expression PP13, including its moving affine cover, gives \[
 \frac1f\frac{d\log F_k}{d\log q}=-164+kw(q)-\varepsilon(q).
\] The complete outside denominator factor is x\^{}-337/2, and q\_* is proportional to x\^{}-1/f.~Therefore \[
 \frac{d\log[x^{-337/2}F_k(q_*)]}{d\log x}
       =-9/2-kw(q_*)+\varepsilon(q_*)<-7/2.
 \tag{PP19}
\] No uniform monotonicity in q after central correlation is required. We first summed all patterns by PP14--18 and only then used this complete x derivative. The pattern majorant4 is independent of x. The full count consequently has its supremum at x=2\^{}X0. Unique eta for each fixed support justifies taking this supremum; no denominator-ideal or eta enumeration is inserted.

\subsection{6. Entire original endpoint, with both new scalar constants}\label{entire-original-endpoint-with-both-new-scalar-constants}

The exact packet T is read from the original J56 host certificate. Direct integer powers establish \[
 4119905<q_*<4119906\quad(x=2^{X_0}),\qquad
 H/2+750R_0/4119905<68,\qquad V_F(68)<2^{479}.
 \tag{PP20}
\] Keep the actual support bound2\^{}101, original pivot2\^{}8, Green2, complete pattern4 and full168-place T correction2. Their combined exponent is113. The saturated-basis gain336, improved selected trace335, actual radial gain649 and coordinate gain floor(log2(k\^{}84)) are precisely the complete K110 PM20--22 factors; PP7 gives its same additional mixture gain2180. These are each retained once.

With L\_T=1019333/1024, L\_Q=287843/4096, L\_U=601356/4096 and L\_R=28810/4096, the directed common exponent is \[
 \begin{split}
 C={}&113+2f-929f/2+479+167fL_R-(335f/2)L_U\\
     &+8594L_T-336\cdot167L_Q-(337/2)X_0
       =211125579/512.
 \end{split}
 \tag{PP21}
\] The T power is335\emph{26+52-168=8594, and the Q power is335}167+167. The true original selected, central, relative-phase, U, R and denominator powers are unchanged. The earlier factor2 in the full Gaussian row bound is retained below, separately from the enumerated finite factors.

For d=k+3 and c\_-=4119905,c\_+=4119906 define the entirely rational majorant \[
 W_k(\mathcal K)=
 \frac{2\binom{164}{k}(22/7)^{df}}{D_0^d c_-^{167f}}
       \{(c_+/(2R_0))^f+\mathcal K\}^{k}
       (c_+/(2R_0))^{3f}.
 \tag{PP22}
\] Exact integer arithmetic gives

\begin{longtable}[]{@{}
  >{\raggedright\arraybackslash}p{(\columnwidth - 8\tabcolsep) * \real{0.1765}}
  >{\raggedright\arraybackslash}p{(\columnwidth - 8\tabcolsep) * \real{0.1765}}
  >{\raggedleft\arraybackslash}p{(\columnwidth - 8\tabcolsep) * \real{0.2353}}
  >{\raggedleft\arraybackslash}p{(\columnwidth - 8\tabcolsep) * \real{0.2353}}
  >{\raggedright\arraybackslash}p{(\columnwidth - 8\tabcolsep) * \real{0.1765}}@{}}
\toprule\noalign{}
\begin{minipage}[b]{\linewidth}\raggedright
Cell
\end{minipage} & \begin{minipage}[b]{\linewidth}\raggedright
Scalar coefficient
\end{minipage} & \begin{minipage}[b]{\linewidth}\raggedleft
Strict ceiling for log2 W
\end{minipage} & \begin{minipage}[b]{\linewidth}\raggedleft
Coordinate gain
\end{minipage} & \begin{minipage}[b]{\linewidth}\raggedright
Complete exponent
\end{minipage} \\
\midrule\noalign{}
\endhead
\bottomrule\noalign{}
\endlastfoot
k163 & K5 & −408627 & 617 & −199349/512\textless−389 \\
k164 & K15 & −408548 & 618 & −159413/512\textless−311 \\
\end{longtable}

Here each complete exponent equals \[
       C+\operatorname{ceil}_{\rm strict}(\log_2 W)
                      -336-335-649-\operatorname{gain}_k-2180.
 \tag{PP23}
\] The companion verifies the least strict ceilings, rather than relying on floating-point logarithms. Both entire continuum reductions, all four-state patterns, every support and coordinate position, the central affine fibre and actual label weights are already included. Summing these two complete cells gives \[
       \sum_{h\in R,\,s(h)=26}|w(h)|
         <2^{-389}+2^{-311}<2^{-310}.
 \tag{PP24}
\] No extra union over patterns, eta or endpoint choices is needed or omitted.

\subsection{7. Exhaustion, chronological errors and remaining gate}\label{exhaustion-chronological-errors-and-remaining-gate}

R091 DR1--7 proves that all earlier support26 classes have already been removed; the later canonical cuts through K110 give exactly the necessary containing conditions PP1. PP24 therefore removes the whole support26 portion of the actual current remainder. It is not just a theorem about a selected sparse subset. Every prior containing set was applied to its intersection with the then-current remainder; overlaps retain their full earlier charges. The complete K110 all-top list119 and uniform list31 are copied in their original order. Exactly310 is appended once to the first list. Both full dyadic sums remain strictly below2\^{}-230, with the unique231 minimum and next233 unchanged.

Splitting off the actual removed support26 set and using the triangle inequality proves PP2. R091's actual higher-support set, with every later restriction in its proved scope, is exactly the new remainder. This proves the complete all-top support floor27 and leaves the uniform floor11. PP3 remains the signed finite/real gate. No whole central orbit is substituted for the selected subset, no Fourier constancy is assumed, and no new complex coefficient sign is inferred from the absolute count.

The complete proof is PP1--24. Exact rational polynomial-product identities on small finite state spaces, derivative diagnostics and endpoint replay supplement it; they are not an enumeration of arithmetic arrows. Every source interface and complete saved proof is reviewed separately. The two scalar theorems are independently checked before this result is promoted. Their coefficient ratios alone are explicitly insufficient for the application.

\subsection{8. Sources and algebraic-geometric placement}\label{sources-and-algebraic-geometric-placement}

The original strata, fractional ideals and common paired incidence modules live over Spec(O\_E) and Spec(O\_F). The actual polarization, theta line, selected realization, Picard and biextension data are fixed. PP14--18 is a finite positive majorization on those original strata, not a replacement arithmetic correspondence. The scalar positivity context is Delsarte, P., Goethals, J. M., and Seidel, J. J.(1977), \emph{Spherical codes and designs}, Geometriae Dedicata6,363--388, Definition2.1, Theorems2.4/3.3/4.3 and Corollary3.8; the exact new polynomials and full cap/layer-cake proofs are DO116 QC1--18 and K111 HP1--18. The mixture input is K110 PM3--15, with NIST DLMF10.32.10/10.39.2 as its exact positive-real normalization context. Their known statements are cited rather than re-proved here.

Carried exact framework locators are SGA3 I4.7.3/4.7.3.1; SGA6 VI Theorem3.3; FGA232 Theorem3.1pp148--149; FGA236 Corollaries1.7/2.3pp229/232; SGA7 I VII2.7/2.9pp169--174 and VIII7.4(b)pp310--311; cote4 scan003 sections8.1.1--8.1.2 and scans005--006 sections8.4.1--8.4.3; cote54 scans005--007. No new formal-group slope, lambda-ring/K-theory identity, l-adic weight or CM-motive identification is obtained. These are carried anchors, not new whole-text readings.

DP\_MULT Q3e/DP44/DP52/DP54--65 and THETA3 G1/A7 receive the strictly smaller finite gate PP3. An actually admitted fully sized mixed source, selected arithmetic content, prescribed period, Read/specific032/STiesR, all six lift identities and evaluated DP44 entries remain open. CX\_GRH CX-ALG-024 still requires the Kummer-Hecke RH input; GCT\_PLETH T-052 and Q-ALPOGE retain the global arithmetic oscillator map. Neither RH nor a shorter valid lift chain or a universal strict negative actual DP44 cost difference is claimed.

\subsection{Canonical algebraic-geometric context}\label{canonical-algebraic-geometric-context-49}

This is a bound on the same selected real theta coefficient and finite correspondences for the original integral Hermitian modules. The local paired fractional ideals and affine central difference ideal are retained before nonnegative enlargement; the containing polynomial is not an equality for the original signed coefficients. Picard/Poincare line, polarization, finite-etale torsion realizations, projectors and arithmetic source are unchanged. Primary scalar context carried through fully checked DO116/K111: Delsarte,Goethals,andSeidel(1977),Definition2.1,Theorems2.4/3.3/4.3 and Corollary3.8,pp364-369. Geometric anchors carried through R091 and LIFT085: FGA232 Theorem3.1pp148-149,FGA236 Corollaries1.7/2.3pp229/232; SGA7 I VII2.7/2.9pp169-174,VIII7.4(b)pp310-311; SGA3 I4.7.3/4.7.3.1; SGA6 VI Theorem3.3; cote4 scan003 sections8.1.1-8.1.2/scans005-006 sections8.4.1-8.4.3 and cote54 scans005-007. No new Dieudonne slope, lambda/K-theory relation, l-adic weight or CM-motive comparison is asserted. DP\_MULT Q3e/DP44/DP52/DP54-65 and THETA3 G1/A7 have a strictly smaller conditional finite remainder; their actual arithmetic chain/pairing gates, CX\_GRH CX-ALG-024 and GCT\_PLETH T-052 remain open globally.

\clearpage

\section{Scope precision for AM-R-093: the literal full Host excludes d = 7}\label{scope-precision-for-am-r-093-the-literal-full-host-excludes-d-7}

\textbf{Status:} immutable source-interface addendum; no change to AM-R-093 DC1--19, no new finding, numerical error charge, or main-gate movement. Prepared by \nolinkurl{/root/lift_accounting} on 23 November 2026 UTC. This note distinguishes a conditional CM-field theorem from the additional requirements for its application to the full pinned source.

\subsection{1. Exact declarations and the elementary lower bound}\label{exact-declarations-and-the-elementary-lower-bound}

The literal full arithmetic record is \nolinkurl{Formalism.BasicNotions.Host (R : Ranks)} in

\nolinkurl{research/alpoge_source/Formalism/LocalRingCardinality/DataHost.lean}, lines 39--54.

Its field \nolinkurl{qd_big}, exactly at line 54, is

\begin{verbatim}
qd_big : 4 * R.t * R.r + 1 < q ∧ 4 * R.t * R.r + 1 < d
\end{verbatim}

The rank record is in

\nolinkurl{research/alpoge_source/Formalism/LocalRingCardinality/Params.lean}, lines 6--22, with \nolinkurl{Ranks.r} at line 32. The following are the sufficient lower hypotheses actually needed here:

\begin{longtable}[]{@{}
  >{\raggedright\arraybackslash}p{(\columnwidth - 4\tabcolsep) * \real{0.3000}}
  >{\raggedleft\arraybackslash}p{(\columnwidth - 4\tabcolsep) * \real{0.4000}}
  >{\raggedright\arraybackslash}p{(\columnwidth - 4\tabcolsep) * \real{0.3000}}@{}}
\toprule\noalign{}
\begin{minipage}[b]{\linewidth}\raggedright
Literal source field
\end{minipage} & \begin{minipage}[b]{\linewidth}\raggedleft
Exact locator
\end{minipage} & \begin{minipage}[b]{\linewidth}\raggedright
Mathematical consequence
\end{minipage} \\
\midrule\noalign{}
\endhead
\bottomrule\noalign{}
\endlastfoot
\nolinkurl{two_le_m : 2 ≤ m} & line 13 & \(m\ge2\) \\
\nolinkurl{m_add_three_le_k : m + 3 ≤ k} & line 14 & \(k\ge m+3\ge5\) \\
\nolinkurl{def r : ℕ := R.k + R.m - 1} & line 32 & \(r=k+m-1\ge6\) \\
\nolinkurl{rank_add_five_le : k + m - 1 + 5 ≤ t} & line 16 & \(t\ge r+5\ge11\) \\
the second conjunct of \nolinkurl{Host.qd_big} & \nolinkurl{DataHost.lean}, line 54 & \(d>4tr+1\) \\
\end{longtable}

All subtraction is the literal natural-number subtraction. Since \(k+m\ge7\), it equals ordinary integer subtraction in the displayed expression. Therefore

\[
  4R.t\,R.r+1\ge4\cdot11\cdot6+1=265,
  \qquad d>265.
\]

These fields alone exclude \(d=7\). Neither the large \nolinkurl{SizesAt} conditions nor a numerical estimate for the actual large ranks is needed.

The actual parity fields make the conservative estimate stronger: \nolinkurl{k_even} at line 10 upgrades \(k\ge5\) to \(k\ge6\), hence \(r\ge7\); \nolinkurl{rank_add_five_le} then gives \(t\ge12\), and \nolinkurl{t_odd} at line 12 gives \(t\ge13\). Thus every full record in fact satisfies

\[
  d>4\cdot13\cdot7+1=365.
\]

The exclusion uses \textbf{\nolinkurl{Ranks.t} and \nolinkurl{Ranks.r} from these exact declarations}. It does not identify \nolinkurl{Ranks.t} with any separately introduced \nolinkurl{TPar.t}, and it does not need \nolinkurl{three_m_lt_t}, \nolinkurl{m_even}, a primality hypothesis, or a congruence condition on \(d\).

\subsection{2. The geometric Host and the actual pin are distinct}\label{the-geometric-host-and-the-actual-pin-are-distinct}

The geometric record

\nolinkurl{ReciprocityTheory.ChartRestriction.G.Geo.Tor.Host E}

is defined in

\nolinkurl{research/alpoge_source/ReciprocityTheory/ReductiveGroups/ArithmeticGeometry/BirationalGeometry/ToroidalCompactification/TwoWeightEqOneAddRoot.lean}, lines 14--26.

It requires that \(d\) be prime and congruent to 7 modulo 16, the specified square root and half-root, conjugation, and relative unramifiedness. It contains neither a \nolinkurl{Ranks} record nor \nolinkurl{qd_big}. Consequently this note does \textbf{not} exclude the \(d=7\) CM/geometric objects considered at that interface, and it does not change any of the conditional constructions or determinant-character identities for those objects.

The passage to the full source is not free: in

\nolinkurl{research/alpoge_source/Formalism/LocalToGlobal/ProjectiveLattices/PartitionExtendByZeroAntitone.lean}, lines 39--41,

\nolinkurl{Par.PinHost} begins with the exact equality

\begin{verbatim}
Q.H.d = D.host.d
\end{verbatim}

and lines 47--49 retain \nolinkurl{PinDat} and the full conjunction \nolinkurl{Pins}. Hence a geometric host with \nolinkurl{Q.H.d = 7} cannot satisfy \nolinkurl{PinHost} for an actual full \nolinkurl{D.host : Formalism.BasicNotions.Host D.R}. Substitution into the equality would give \nolinkurl{D.host.d = 7}, contradicting the previous section. No conclusion is asserted about other, non-pinned geometric data.

\subsection{3. Required reading of AM-R-093}\label{required-reading-of-am-r-093}

The unchanged canonical subject is

\nolinkurl{research/d7_determinant_character_torsion_fields.md}, AM-R-093 DC1--19.

Its sections 1--5 prove a conditional theorem for the specified actual CM varieties \(J\) and \(C_0\), their full actions and pairings, and the everywhere-good field containing \(\mathbb Q(\zeta_7)\). The single unramified quadratic character, the integral determinant congruence, and the joint torsion-field containment remain precisely in that stated CM scope. Section 6 adds the indicated conditional source-level congruences to obtain its \(M^3\) field estimate.

Phrases such as ``the actual Host parameter \(d=7\)'' in the opening and ``actual \(d=7\) level'' in section 6 must therefore be read as a \textbf{conditional specialization of those CM/level hypotheses}, not as an existence theorem for a full \nolinkurl{Formalism.BasicNotions.Host} or a \nolinkurl{D7} package satisfying all pins. No nonempty full-Host \(d=7\) branch was intended or proved. If the full \nolinkurl{Host.qd_big} and rank fields are additionally imposed, that specialization has no instance, as shown above.

This addendum does not weaken or rewrite DC1--19, alter their constants, or replace any CM variety, action, field, level, or proof. It records the precise application boundary: \textbf{AM-R-093 does not provide a field-cost bound for an admitted full-source \(d=7\) family.} The general-\(d\) pairing-aware theorem AM-LIFT-093, canonical at

\nolinkurl{agents/pairing_aware_same_seed_torsion_fields.md},

is the relevant admitted-source field-cost route, conditional on its own explicit construction and level hypotheses. The actual Host/Sizes construction, effective isogeny scale, selected nonzero evaluation, and final comparison remain separate unresolved inputs.

\subsection{4. Verification boundary}\label{verification-boundary}

The four literal declarations above were freshly read in their displayed ranges, together with all eight sections and the canonical context of AM-R-093. The inequalities in section 1 are a universal elementary deduction from the exact source fields; no finite numerical experiment substitutes for that deduction. The adjacent verification record fingerprints the exact source and canonical subjects and this immutable note. No canonical file was modified.

\clearpage

\section{AM-DO-117 --- the actual selected trace removes nonunipotent fixed blocks at support twenty-seven}\label{am-do-117-the-actual-selected-trace-removes-nonunipotent-fixed-blocks-at-support-twenty-seven}

\textbf{Canonical checked integration --- 2026-11-23T07:43:26.263878+00:00.} Complete root and independent proof/script/interface reviews and separate exact replays passed. The actual canonical predecessor has120/31 entries; this integration gives121/31, complete floors27/11. Frozen author history is preserved. AM-DO-117 FC1-25 removes the entire original conditional J56 all-top H133 fixed-S nonunipotent support27 sector, at every global rank and incidence profile. The five complete containing-set exponents are -5531348963/4096, -3078643005/4096, -2767418421/1024, -9563431915/2048 and -120294901/4096; their union is below2\^{}-29365. The new small scalar remains in the full paired I\_Sigma, while its exact original near-point transport uses only the universal RT2 selected Hilbert trace bound. Append29365 once to the complete R094120/31 predecessor:121/31, both full sums\textless2\^{}-230, complete floors27/11 unchanged. This is a strictly smaller conditional support27 remainder, with no full support27 exhaustion, actual admitted Host/Sizes, selected arithmetic content, prescribed global period, specific032 pairing, evaluatedDP44 or RH conclusion.

Author: \nolinkurl{/root/definite_oscillator}, 23 November 2026 UTC. Controlled pending theorem. Author, independent and root whole checks are separate events. Earlier frozen subjects and canonical arrays are unchanged.

\subsection{1. Complete sector and result}\label{complete-sector-and-result}

Retain the original J56 all-top family, every allowed interior word and canonical twist, the original finite projectors, paired source and target pencils, class masses, phases and actual near point. Use

\[
 f=a=168,\quad n=169,\quad H=133,\quad
 h=I_S\oplus h_0,\quad \dim_E S=164,\quad
 h_0\in U(4,1)(F),\quad \det h_0=1.                 \tag{FC1}
\]

The entire sector considered here consists of nonunipotent matrices in FC1 with nonzero original finite coefficient, all 334 unmixed G reductions unipotent, and exactly s=27 nonintegral relative T places. No global-rank or incidence-type restriction is imposed. The original cutoff gives \(\|h_\tau\|_{\rm op}<133\) at every label.

\textbf{Theorem.} The total absolute contribution of this complete sector is less than

\[
                              2^{-29365}.              \tag{FC2}
\]

Only its intersection with the current remainder is removed. This is one fresh all-top charge, not an exhaustive support-floor update.

The small scalar used below retains its original full-support denominator ideal \(I_\Sigma\). In particular, this proof does \textbf{not} replace s by the number of type-III places in DO089 NC18. Characteristic coefficients have the narrower pole statement of K070 PT5; the small scalar is not a characteristic coefficient.

\subsection{2. Original quadratic alternatives and all other cases}\label{original-quadratic-alternatives-and-all-other-cases}

Put

\[
 N=h_0-I,\qquad H_h=h_0+h_0^{-1}-2I
       =-N^\dagger N=N^2h_0^{-1}.                       \tag{FC3}
\]

Use the original integers and transverse inputs of DO089 NC4--5:

\[
\begin{split}
 Q&=1427908193729121522601,\qquad U=77Q^2-14,\\
 u&=(2,95192390527,383092671043,12529843548438452388609),\\
 y_1&=(-976627,-153032,38026,0),\\
 y_2&=(-567678,695570,-172838,0),\\
 y_3&=(104583,-15309315176,-61610045224,2).
\end{split}                                             \tag{FC4}
\]

Their negative coordinates are zero. The positive vectors y\_j span \(T_0=u^\perp\); write \(P=T_0^\perp=Eu\oplus Ee_n\). This fixed plane has signature (1,1). Their squared Euclidean norms are

\[
 L_1=978665066829,\quad L_2=835948910828,\quad
 L_3=4030172803672366725045.                              \tag{FC5}
\]

The complete DO089 proof NC3--19 is independent of the value s=22. Its use at arbitrary s, including its precise denominator and counting dependence, was recorded in DO090 NE3--6. We retain its following alternatives without changing the original lattice frame.

\begin{enumerate}
\def\labelenumi{\arabic{enumi}.}
\tightlist
\item
  Some integral compression \(\langle y_i,H_hy_j\rangle\) is nonzero. Its absolute value at each label is below \(2^{81}\), and it has no T poles.
\item
  The nontrivial characteristic degree is k=3,4 or5. The corresponding nonzero coefficient \(e_k(N)\) has absolute value below \(2^{b_k}\), where \((b_3,b_4,b_5)=(13,13,12)\), and at most simple poles at the actual type-III partners.
\item
  The characteristic degree is two and the entire transverse compression vanishes. DO089 NC13--15 proves, including all possible generalized-primary multiplicities, that the global rank is two and
\end{enumerate}

\[
 y=\operatorname{tr}N\in F^\times,\quad |\tau y|<136,
 \quad H_h=yP_V,\quad H_h^2=yH_h,
 \quad H_h(T_0)\subset P.                               \tag{FC6}
\]

Here V is the nontrivial primary plane and P\_V its Hermitian orthogonal projection. Every vector H\_h y\_j is isotropic. If at least one is nonzero, it has the unique form

\[
                         H_hy_j=(u b,\beta),\qquad b\ne0,
                |\beta_\tau|=\sqrt U\,|b_\tau|.         \tag{FC7}
\]

The original paired source and target membership, followed by the integer Bezout relation \(\gcd(u_i)=1\), proves exactly

\[
               b\in I_\Sigma,\qquad N_E I_\Sigma=T^{-2s}.
                                                               \tag{FC8}
\]

Both CM partners are included. At a G reduction of rank d, the paired scalar valuation is at least \(2(2-d)_+\), as in NC7. The original rank-two Green comparison therefore gives the pointwise bound

\[
 |a(h)|\le 2 B^{335f/2}T^{335s/2}Q^{-335(f-1)}
 \quad\text{if } |b_\tau|<B\text{ at every label}.       \tag{FC9}
\]

This bound includes the two mixed contractions and the whole Euler correction. One may choose a nonzero j for each matrix: all choices give the same uniform bound, so no extra witness-choice count is introduced.

If instead H\_h(T\_0)=0, then V=P and DO090 NE3--4 gives

\[
             y\in U J_{G,F} I_{\mathrm{III},F},\qquad
                   N_FJ_{G,F}=Q^{167}.                 \tag{FC10}
\]

Since t\textless=s=27, the exact unchanged integer data satisfy

\[
 U^{168}Q^{167}>136^{168}T^{27},                         \tag{FC11}
\]

so this last alternative is empty. Its conservative binary margin is \(33504601/4096>0\), using the directed bounds in section 5 below and \(136<2^8\). This retains the actual F-ideal norm, not its square; the U and G ideals are coprime to the T denominator. No anisotropy or norm-solvability assertion is used.

\subsection{3. A stronger bound for the same small scalar}\label{a-stronger-bound-for-the-same-small-scalar}

For a unitary matrix of signature (a,1), write r for its Cartan radius. Its singular values are \(e^r,e^{-r}\) and ones, so \(\|h\|_{\rm op}=e^r\). Its negative diagonal entry has modulus \(\cosh r\). Unitarity of the last row implies that the Euclidean norm of the remaining part of that row is \(\sinh r\). The same statements hold for the inverse, whose Cartan radius is also r. These are statements about the original real Hermitian form, not a residue-field positive metric.

The vector y\_j has zero negative coordinate. Taking the last coordinate in FC7 and using FC3 consequently gives

\[
 \sqrt U\,|b_\tau|
 \le 2\sinh r_\tau\,\|y_j\|
 <\|h_\tau\|_{\rm op}\|y_j\|
 <133\sqrt{L_3}.                                      \tag{FC12}
\]

The middle strict inequality is \(2\sinh r=e^r-e^{-r}<e^r\). This is stronger than applying the full operator norm \(\|H_h\|\le268\) to y\_j. The exact integer comparison

\[
                    133^2 L_3\,2^{60}<U                \tag{FC13}
\]

therefore proves the uniform bound \(|b_\tau|<2^{-30}\). There is no square-root extension of E; square roots occur only in real inequalities. All arithmetic remains in the original integral coordinates and fractional ideal FC8.

\subsection{4. Actual near-point transport and selected trace}\label{actual-near-point-transport-and-selected-trace}

Use the literal LIFT41 NP5--7 representative

\[
 z_{\rm pkt}=(Qu,\pi_k)/\bar\pi,\quad c=(z_{\rm pkt})_n,\qquad
 b_{\rm pkt}=z_{\rm pkt}-e_n,\qquad
 g_p=I+\frac{b_{\rm pkt}b_{\rm pkt}^*J}{1-\bar c}.
                                                               \tag{FC14}
\]

Its positive component is proportional to u. Thus b\_pkt lies in the fixed plane P, and \(b_{\rm pkt}^*J y_j=0\). In particular

\[
                       g_p y_j=y_j,\qquad g_p(P)=P.     \tag{FC15}
\]

This is an exact identity at every label, not merely a statement that the point is close to the origin. A different representative g\_p k of the same point, with k compact, gives compact conjugation of the argument; the normalized selected trace is unchanged because the actual selected type is compact-stable.

Let \(h'=g_p^{-1}hg_p\), with Cartan radius r', and let r\_p be the Cartan radius of g\_p.~Because FC15 fixes y\_j,

\[
                   H_{h'}y_j=g_p^{-1}H_hy_j.            \tag{FC16}
\]

For any nonzero isotropic vector w with negative coordinate beta, the positive part has Euclidean norm \textbar beta\textbar. The positive and negative parts of g\_p e\_n have norms \(\sinh r_p\) and \(\cosh r_p\). The reverse triangle inequality in their Hermitian pairing proves

\[
        |(g_p^{-1}w)_n|\ge e^{-r_p}|\beta|.             \tag{FC17}
\]

Applying this to FC7 and bounding the last rows of h' and its inverse as in FC12 yields

\[
        2\sinh r'\,\|y_j\|
          \ge e^{-r_p}\sqrt U\,|b_\tau|>0.             \tag{FC18}
\]

The original normalized selected trace is exactly DO080 RT1 at this argument. Its universal RT2 estimate, retaining both compact factors and the genuine vector-valued selected Hilbert metric, gives

\[
 |c_{W,\tau}(h')|
 <\frac{2^{-465}}{\sinh r'}
 \le\frac{2^{-464}e^{r_p}\|y_j\|}{\sqrt U\,|b_\tau|}
 <\frac{2^{-501}e^{r_p}}{|b_\tau|}.                    \tag{FC19}
\]

The last strict comparison uses the exact inequalities \(L_3<2^{72}\) and \(U>2^{146}\). It is valid even if its right side exceeds one; no lower cutoff on an individual conjugate of b is imposed. This is a direct transport of the isotropic H\_h image. It does not apply DO080's unipotent-only RT3 to a nonunipotent matrix, and it does not replace a small additive displacement error by an unproved multiplicative one.

LIFT41 NP12 and the unchanged integers give

\[
 \sinh^2 r_p=(77Q^4-14Q^2)/T<2^{-704},\qquad
             \prod_\tau e^{r_{p,\tau}}<2.              \tag{FC20}
\]

For the second inequality, \(r_p\le\sinh r_p<2^{-352}\), so \(\sum_\tau r_{p,\tau}<168\,2^{-352}<1/2\); use \(e^x\le(1-x)^{-1}<2\) for \(0\le x<1/2\). It costs one factor two for the full archimedean product, not one per label.

For the same nonzero algebraic b chosen in FC7, FC8 implies

\[
              \prod_\tau|b_\tau|=|N_Eb|^{1/2}\ge T^{-s}.
\]

Hence the complete actual selected real factor obeys

\[
                       |c(h)|<2^{1-501f}T^s.           \tag{FC21}
\]

The new trace bound and FC9 apply simultaneously to each matrix before any containing lattice set is counted. Enlarging that matrix count afterward does not assume that every enlarged lattice point has the same geometric properties.

\subsection{5. Full original columns, pair weights and exact exponents}\label{full-original-columns-pair-weights-and-exact-exponents}

Keep the full five-complex-coordinate negative-column count of DO089 NC22, with all target labels including zero:

\[
 \#h\le G_5V_5(223)T^{10s},\qquad G_5<2^{107},
 \quad V_5(223)=\frac{[((22/7)f223^2)]^{5f}}{(5f)!D_0^5}<2^{5757},
 \quad D_0=7^{140}29^{144}5^{84}13^{84}.                 \tag{FC22}
\]

The source contains \(\pi e_n\), so the column lies in the original target overlattice \(\pi^{-1}M_d\). Its paired covolume is \(D_0^5T^{-10s}\); it contains the original integral five-coordinate lattice, including at target zero. The original radius below \(95\sqrt f\) and the accepted integral-ring cover give \(223\sqrt f\). The same-pair source fibre is bounded by the actual GL5 gcd at the two G primes. No real-part slice, determinant enumeration or new class-mass factor is introduced.

Support choices are bounded by \(2^{168}\). The complete ordered-pair bound is

\[
                 P_T=\frac{2(T-1)^2}{T-2}<2^{997}.      \tag{FC23}
\]

For the integral-compression and quadratic-small-scalar alternatives use \(P_T^s\), which already sums all actual types and labels. For characteristic degree k, the exact type-III ordered off-diagonal pairs cancel the reciprocal absolute local weight. With \(L_k=169k-k(k+1)/2\), the full type/position sum is the unchanged

\[
 \sum_{t=0}^s {s\choose t}T^{L_kt/k}P_T^{s-t}
 =(T^{L_k/k}+P_T)^s<2^sT^{L_ks/k}.                     \tag{FC24}
\]

These are precisely NC22--23, including the two mixed contractions, every class and central phase, and the real bound one for the four alternatives where FC21 is not used. There is no extra G-rank pattern count or witness multiplicity.

Put \(L_T=1019333/1024\), \(L_Q=287843/4096\), strict upper/lower bounds for \(\log_2T\) and \(\log_2Q\), and \(C=1+168+107+5757=6033\). The five complete support-27 exponents are

\[
\begin{split}
 E_A&=C+81(335f/2)+(10L_T+997)27-335\cdot167L_Q
      =-5531348963/4096,\\
 E_3&=C+27+13(501f/3)+(10+501/3)27L_T-501\cdot167L_Q
      =-3078643005/4096,\\
 E_4&=C+27+13(666f/4)+(10+666/4)27L_T-666\cdot167L_Q
      =-2767418421/1024,\\
 E_5&=C+27+12(830f/5)+(10+830/5)27L_T-830\cdot167L_Q
      =-9563431915/2048,\\
 E_C&=C-30(335f/2)+(11+335/2)27L_T+997\cdot27
                  -335\cdot167L_Q-501f+1
      =-120294901/4096.
\end{split}                                             \tag{FC25}
\]

In the final line, the extra T exponent one and the extra global factor two are exactly FC21; the paired Green pole exponent remains \(335s/2\) from the full ideal FC8. Each exponent is strictly below -29368. Their complete absolute union is therefore less than \(5\,2^{-29368}<2^{-29365}\). The pure-plane case was proved empty in FC11. This proves FC2 for all global ranks and all incidence profiles.

\subsection{6. Chronology, sources and limits}\label{chronology-sources-and-limits}

The certificate accepts the actual canonical predecessor through \nolinkurl{--prior}. It verifies its entire K110 119-term all-top prefix and unchanged 31-term uniform array, retains any later intervening canonical entries, and appends 29365 exactly once. The historical prefix and any later chronology replay are stored separately; no frozen author output is overwritten. The complete dyadic sums remain below \(2^{-230}\). Existing support floors are retained: a full support-27 exhaustion requires separate theorems for the other structural classes.

Direct carried inputs are DO089 NC1--24 and DO090 NE1--6, including their full original negative-column and paired Green proof; K070 PT1--6 for the exact distinction between coefficient and entry poles; LIFT41 NP1--12 for the unchanged point and its exact representative; and the whole DO080 RT1--24 selected density theorem. The new steps FC12--21 are proved above for the nonunipotent quadratic alternative and use only RT2's universal bound. The original Hilbert model remains Nakahama, R. (2016), \emph{Norm computation and analytic continuation of vector valued holomorphic discrete series representations}, \emph{Journal of Lie Theory, 26}(4), 927--990, equation (5.7) and Theorem 5.1, through the checked K043/K046 dictionary. That source is carried at its established scope; no new whole-paper reading is claimed here.

The same original integral-module and determinantal context is retained: SGA3 I4.7.3/4.7.3.1; SGA6 VI Theorem3.3; FGA232 Theorem3.1 pp148--149; FGA236 Corollaries1.7/2.3 pp229/232; SGA7 I VII2.7/2.9 pp169--174 and VIII7.4(b) pp310--311; cote4 scan003 sections8.1.1--8.1.2, scans005--006 sections8.4.1--8.4.3 and cote54 scans005--007. These are inherited anchors, not fresh comparison constructions or reading claims. No finite projector, field, prime, pencil, real representation, metric or arithmetic image has been changed.

The companion verifies exact constants, all five full bounds, original fibre/volume factors, paired Green losses and rational unitary/transport models. Those finite diagnostics do not replace the universal matrix, Hermitian-primary or lattice arguments. This theorem makes no assertion about a signed sufficient sum, actual selected arithmetic content, prescribed Bessel/Mram period, effective source, Read/specific032, DP44 or RH.

\subsection{Canonical algebraic-geometric context}\label{canonical-algebraic-geometric-context-50}

This is an estimate for the same selected real theta coefficient on the original integral Hermitian modules and finite correspondences. Exact null-image transport preserves the original fractional ideal, packet, compact factors, vector-valued metric and Hilbert/PBW normalization. The carried primary analytic source is Nakahama(2016), equation(5.7) and Theorem5.1, through the checked K043/K046 dictionary and DO080 RT1-24. The unchanged geometric anchors are SGA3 I4.7.3/4.7.3.1; SGA6 VI Theorem3.3; FGA232 Theorem3.1 pp148-149; FGA236 Corollaries1.7/2.3 pp229/232; SGA7 I VII2.7/2.9 pp169-174 and VIII7.4(b) pp310-311; cote4 scan003 sections8.1.1-8.1.2 and scans005-006 sections8.4.1-8.4.3; cote54 scans005-007. These are inherited locators, not new whole-source readings or constructions. No new Dieudonne slope, lambda/K-theory relation, l-adic weight, CM-motive map or global arithmetic seesaw is asserted. DP\_MULT Q3e/DP44 and THETA3 G1 retain their actual arithmetic gates; CX\_GRH CX-ALG-024 and GCT\_PLETH T-052 retain their established scope.

\clearpage

\section{AM-K-112 --- every active type-II place is eligible in the full-positive unipotent form}\label{am-k-112-every-active-type-ii-place-is-eligible-in-the-full-positive-unipotent-form}

\textbf{Canonical checked integration --- 2026-11-23T07:45:26.546837+00:00.} Complete saved-proof and executable reviews and separate exact diagnostic replays passed. No numerical error charge is introduced. Frozen author history is preserved. AM-K-112 UE1-14 proves that every active type-II place in the exact full-positive globally rank-two-unipotent form is eligible, without an eta-denominator bound. The noneligible active subset is empty, with zero error charge. In the actual all-II scalar-cross support27 remainder, v=s=27, x\textgreater2\^{}2111; x\textless2\^{}3107 forces total eligibility27; x\textless2\^{}12000 gives total eligibility at most36 and at most one eta per fixed support. Fixed-eta active labels contribute at most(T-1)\^{}-27. The full168-factor bound additionally requires total eligibility27 and no inactive paired-C-zero place. Original paired row and central difference ideals retain their support27 denominator power54. Current DO117121/31 exact errors and complete floors27/11 remain unchanged; no full support27 estimate, actual source/content/period, specific032, DP44 or RH conclusion is asserted.

Author: \nolinkurl{/root/kummer_extension}. 23 November 2026 UTC. Controlled pending proof. This is a support-independent local exclusion and its explicitly bounded arithmetic consequences. Author, independent and root verification are separate events. No canonical file or error ledger is changed.

\subsection{1. Actual scope and conclusion}\label{actual-scope-and-conclusion-2}

Keep the actual J56 field, involution, original Hermitian module, source and target pencils at both CM components, omitted labels, class masses, finite projectors, selected real representation, central phases and Haar measures. Suppose the original determinant-one representative has the exact full-positive globally rank-two-unipotent form \[
 E^{169}=Z\perp P,\qquad P=Eu\perp Ee_n,\qquad
 h=I+N,\quad h_{ZZ}=I,\quad N^2Z=0,\quad
 \operatorname{rank}_E N=2,\quad N^3=0.
 \tag{UE1}
\] Its isotropic line im(N²) in P has normalized slope eta with eta bar(eta)=U. The original constants are \[
 Q=1427908193729121522601,\quad U=77Q^2-14,
 \quad a_0=Q(-7+2\sqrt{-7}).
 \tag{UE2}
\] An original selected T place is eligible when eta is integral there and eta≡a0 modulo its selected prime. This is the actual selected pole-line condition; its paired counterpart follows from the norm equation and is not independently prescribed.

\textbf{Theorem.} Every active type-II place of UE1 is eligible. Thus, in the all-II support-s branch, its active eligibility count is exactly v=s, without any bound on the eta denominator. The noneligible active type-II containing subset is empty and has exactly zero contribution. There is no error charge for this exclusion.

After fixing eta, the actual complete active-label/weight sum is at most (T−1)\^{}(-s). This does not assert that eta is globally unique without a separately proved denominator bound, and it supplies no inactive factor by itself.

The full-positive hypothesis is essential to this application. For the scalar-cross globally rank-two-unipotent branch, it follows after the accepted DO099 plane-C removal through support70 from the support-independent proof K103 SC26.5--7: the unique image line of N² lies in P, and N²Z=0 forces NZ into that line. Consequently the theorem applies at support27 in that actual containing remainder. It is not a statement about arbitrary globally nonunipotent matrices or merely finite unipotent reductions.

\subsection{2. An exact rational-DVR lemma}\label{an-exact-rational-dvr-lemma}

Let O be a DVR with fraction field K and uniformizer pi. Let \[
 \Lambda=Z_O\oplus P_O,\qquad \operatorname{rank}P_O=2.
\] Take a primitive z∈P\_O and an integral row q on Lambda satisfying qz=0, q(Z\_O)⊂pi O and with q restricted to P primitive. Let c∈O\^{}× and \[
 A_c=I+c\pi^{-1}zq,\qquad A_c^{-1}=I-c\pi^{-1}zq.
 \tag{UE3}
\] Let w∈P\_O and b∈P\_O\^{}vee be primitive, with bw=0, and extend b by zero on Z. Assume the residue line of w differs from that of z. Since q\textbar P is a primitive annihilator of z in the rank-two residue plane, this is equivalent to \[
                       qw\in O^\times,\qquad bz\in O^\times.
 \tag{UE4}
\] For an arbitrary K-row R vanishing on P, a K-column V∈Z\_K and d∈K, put \[
                    N=wR-Vb+dwb.
 \tag{UE5}
\] \textbf{Local lemma.} The following conditions are equivalent: \[
 \begin{split}
 &A_c^{-1}NA_c\in\operatorname{End}_O(\Lambda);\\
 &R(Z_O)\subset\pi O,\quad V\in\pi Z_O,\quad d\in\pi^2O.
 \end{split}
 \tag{UE6}
\] This is a valuation statement over the original local integral frame. No positivity over the residue field, global integral isotropic basis, Hermitian adjoint identity or height bound is needed for the lemma.

\subsection{3. Necessity, including every possible cancellation}\label{necessity-including-every-possible-cancellation}

The two transformed factors are \[
 A_c^{-1}w=w-c\pi^{-1}z(qw),\qquad
 bA_c=b+c\pi^{-1}(bz)q.
 \tag{UE7}
\] By UE4 each has exact minimum coordinate valuation −1. The restriction of bA\_c to P has the same exact minimum −1. Indeed its leading polar coefficient is a unit times the primitive row q\textbar P; the integral term cannot cancel that pole. Likewise the leading polar column in A\_c\^{}(-1)w is a unit times primitive z.

For any y∈Z\_O, A\_c\^{}(-1)y is integral because qy∈pi O. Also Ny=wR(y), since b vanishes on Z. If the first condition of UE6 holds, then \[
 A_c^{-1}Ny=(A_c^{-1}NA_c)(A_c^{-1}y)
                         =(A_c^{-1}w)R(y)
                         \in\Lambda.
 \tag{UE8}
\] The exact pole in the column therefore forces R(y)∈pi O for every y, giving the first divisibility in UE6.

The difference A\_c\^{}(-1)−I has image in P, so the Z-row projection of left multiplication by A\_c\^{}(-1) is unchanged. Moreover A\_c(P)⊂P\_K. The Z-by-P block of the conjugate is consequently exactly \[
                  (A_c^{-1}NA_c)_{ZP}=-V(bA_c)|_P.
 \tag{UE9}
\] No wR or dwb term has a Z-row component. The exact pole of the restricted row forces every coordinate of V to lie in pi O. This selected-component argument avoids any inference from a division of a paired lattice element.

Since Rz=0, RA\_c=R. Since q(Z\_O)⊂pi O and V∈pi Z\_O, qV∈pi²O and hence A\_c\^{}(-1)V∈pi Lambda. Thus the two conjugated mixed terms \[
                 (A_c^{-1}w)R,\qquad -(A_c^{-1}V)(bA_c)
 \tag{UE10}
\] are integral separately. Subtracting them from the integral conjugate leaves d(A\_c\^{}(-1)w)(bA\_c) integral. Its rank-one outer product has exact minimum coordinate valuation −2: choose one coordinate of valuation −1 in each factor. Therefore d∈pi²O. This excludes both a central pole and any purported cancellation of it by mixed terms. The argument also treats d=0, whose valuation is understood as infinite.

\subsection{4. Sufficiency and inactivity in the original standard module}\label{sufficiency-and-inactivity-in-the-original-standard-module}

Conversely, assume all three divisibilities in UE6. The two mixed terms in UE10 are integral, by the preceding identities, and d times the outer product in UE7 is integral because its pole order is at most two. Their sum is A\_c\^{}(-1)NA\_c, proving sufficiency.

The original N in UE5 is itself integral, and in fact N Lambda⊂pi Lambda. Independently of these valuations, \[
 N^2=-(RV)wb,\qquad N^3=0,
 \qquad (I+N)^{-1}=I-N+N^2.
 \tag{UE11}
\] Thus both h and its inverse preserve the original standard lattice. The same divisibilities also make A\_cprime\^{}(-1)NA\_cprime integral for every finite cprime∈O, including cprime=0; this follows from the identical calculation with cprime allowed to be nonunit. This last statement concerns the finite shear family and does not require an additional assertion about an omitted infinity label.

\subsection{5. Verification of the actual source hypotheses}\label{verification-of-the-actual-source-hypotheses}

At the original selected T component, the actual fixed Z/P frame is integral and invertible: the original frame determinant has factors only 2U, which are T-units. This is K103 SC26.5--7 and DO098 CF4--7. It is a local unit change of original coordinates, not a global integrality assertion for a moving null basis.

The actual pencil is exactly \[
 M_c=A_c\Lambda,\quad A_c=I+c\pi^{-1}zq_*,\quad
 z=v_0/\bar\pi\in P,\quad
 q_*=-v_0^\dagger-\pi m,\quad mz=1,\quad q_*z=0.
 \tag{UE12}
\] K063 PF12 proves that the allowed choice m=e1\^{}t/z1 leaves each actual lattice and each label unchanged, by an integral right change of basis. Both omitted labels remain exactly as prescribed. Since v0∈P and Z is orthogonal to P, q\_*(Z\_O)⊂pi O. Its P restriction is primitive and annihilates z modulo pi; the actual pole line has two nonzero plane coordinates. These are DO105 EL9--14 and DO107 IP4--7. No new pencil is substituted.

Saturate the actual image line locally, choosing primitive w∈P\_O and a primitive row b annihilating it, extended by zero on Z. The exact normal form K103 SC26.7 / DO098 CF4 becomes UE5 after absorbing the two local scalar rescalings into R,V,d.~The properties R\textbar P=0 and V∈Z are unaffected. This construction permits zeros and poles of eta and does not assume that R,V or d are initially integral.

If the place is noneligible, the saturated residue line differs from z. For unit eta equality is exactly eta≡a0. If eta has a zero or pole, one coordinate of its saturated line is zero, whereas both coordinates of the actual pole line are units. The norm equation makes the paired line condition consistent, but UE8--10 already proves the necessary divisibilities using the selected component alone. Thus UE4 holds.

An active type-II place has a fixed actual pencil M\_c.~Its label cannot be zero, since M0 is the standard lattice and activity means nonintegrality there; hence c is a residue unit. The original type-II fixed incidence gives A\_c\^{}(-1)hA\_c∈GL(Lambda), and therefore the first condition of UE6. That lemma forces h and h\^{}(-1) integral in Lambda. The exact inverse-adjoint companion is then integral at the other CM component as well. This contradicts the original active support. Theorem UE1--2 follows.

This uses the literal distinction in K071 NB3--7: type II is a fixed incidence, whereas type III is off-diagonal. The proof does not replace an off-diagonal source/target pair by a fixed one. DO113 LS1--4 treats the complementary eligible situation, where the transformed column and row vanish to order at least one at the cancellation label and maximal poles are locally allowed. The present noneligible exclusion is consistent with that result and makes no claim about global realizability of eligible local parameters.

\subsection{6. Exact arithmetic consequences at support27, without a count}\label{exact-arithmetic-consequences-at-support27-without-a-count}

Write x=N\_E D\_eta for the exact denominator norm and q\_tot for total eligibility among all168 selected T places. The theorem gives v=s=27 in the all-II27 branch, hence q\_tot≥27. The support-independent denominator-cleared norm argument DO103 LD12--15 / K099 GE4 remains \[
                  7^{84}T^{q_{\rm tot}}<[4(U+14)]^{168}x.
 \tag{UE13}
\] The strict actual-host integer comparisons in the certificate give x\textgreater2\^{}2111 and show q\_tot≤36 whenever x\textless2\^{}12000. They also show that x\textless2\^{}3107 forces q\_tot=27. These are necessary conditions, not bounds on the contribution of the omitted denominator ranges.

If two norm-U slopes are eligible at the same27 active places and both have denominator norm below2\^{}12000, their denominator-cleared difference is divisible at all54 CM primes. Its norm is at least T\^{}54 if nonzero and at most (4U)\^{}168 times the product of the denominator norms. The exact inequality \[
                         (4U)^{168}2^{24000}<T^{54}
 \tag{UE14}
\] therefore proves uniqueness of eta per fixed27-place support in this bounded window. Neither a denominator-ideal count nor a multiplicity28 is inserted there. No uniqueness is asserted outside the stated window.

For fixed eta, the support-free actual one-label argument DO105 EL9--16 now applies at every active place. The complete active label-weight sum is at most (T−1)\^{}(-27), including central-only poles and the possibility that the required label is omitted. This statement also applies when eta is enumerated first in a separate large-denominator count. It does not turn an eta sum into one term without UE14.

Inactive factors remain separate: DO107 supplies (T−1)\^{}(-1) at each noneligible inactive place which is not a paired-C-zero place. Together with the active factors, this proves the full168-factor bound when q\_tot=27 and there is no inactive paired zero. Those two conditions are not consequences of eligibility alone, and no necessity claim about other possible inactive coefficients is made. For the original globally normalized central scalar d, the difference ideal remains i J0 B\_eta I\_Sigma,F², with real norm denominator T\^{}54, and the full row ideal remains J I\_Sigma D\_eta, with N\_E I\_Sigma=T\^{}-54. No support26 denominator window, numerical charge, or T exponent is carried into a support27 count by this theorem.

\subsection{7. Verification and source boundary}\label{verification-and-source-boundary}

The companion verifies the local equivalence by exact rational matrices, both transformed valuations, original-input membership, the Z-by-P block identity, nilpotence and the integral inverse. It includes nontrivial q\textbar Z and multidimensional Z models. It also reads the actual host and verifies UE13--14's explicit integer comparisons. These diagnostics support but do not replace the universal DVR proof UE3--11.

Exact source interfaces are K063 PF12; K071 NB3--7; DO098 CF4--7; DO105 EL9--16; DO107 IP4--10; K103 SC26.5--7 and its local eligibility discussion; DO113 LS1--4; DO103 LD12--15 and K099 GE4. The field and full source/target Hermitian pairing, projectors, compact factors, phases and Haar measures are unchanged. The original FGA/SGA/Picard/biextension context is inherited; no new arithmetic realization or model comparison is asserted. Exact carried contextual locators are SGA3 I4.7.3/4.7.3.1; SGA6 VI Theorem3.3; FGA232 Theorem3.1 pp148--149; FGA236 Corollaries1.7/2.3 pp229/232; SGA7 I VII2.7/2.9 pp169--174 and VIII7.4(b) pp310--311; cote4 scan003 sections8.1.1--8.1.2 and scans005--006 sections8.4.1--8.4.3; cote54 scans005--007. These are inherited context, not new primary comparison claims. A complete support27 count, its signed sufficient inequality, actual selected arithmetic content, prescribed period, effective source costs, Read/specific032, DP44 and RH remain separate.

\subsection{Canonical algebraic-geometric context}\label{canonical-algebraic-geometric-context-51}

The theorem computes the integral stabilizer of the unchanged paired Hermitian lattices over the original selected split DVR. The original finite correspondences, omitted labels, inverse-adjoint companion, Picard/Poincare polarization data and selected real coefficient are retained. It constructs no new arithmetic realization. Carried contextual anchors: SGA3 I4.7.3/4.7.3.1; SGA6 VI Theorem3.3; FGA232 Theorem3.1 pp148-149; FGA236 Corollaries1.7/2.3 pp229/232; SGA7 I VII2.7/2.9 pp169-174 and VIII7.4(b) pp310-311; cote4 scan003 sections8.1.1-8.1.2 and scans005-006 sections8.4.1-8.4.3; cote54 scans005-007. These are inherited exact locators, not fresh primary readings or new lambda/K-theory, l-adic weight, Dieudonne or CM-motive theorems. DP\_MULT Q3e/DP44 and THETA3 G1 have a smaller conditional finite remainder; their actual arithmetic constructions, CX\_GRH CX-ALG-024 and GCT\_PLETH T-052 retain their separate global gates.

\clearpage

\section{AM-LIFT-095 --- Integral hyperbolic frames and the actual Hecke clearing input}\label{am-lift-095-integral-hyperbolic-frames-and-the-actual-hecke-clearing-input}

\textbf{Canonical checked integration --- 2026-11-23T07:54:50.914760+00:00.} Complete saved-proof and executable reviews and separate exact diagnostic replays passed. No numerical error charge is introduced. Frozen author history is preserved. AM-LIFT-095 IH1-6 proves an integral hyperbolic pair inside every stated selfdual relative-unramified indefinite Hermitian lattice of rank at least3, by local integral pairs and strong approximation on SU. The original nonterminal source W\_j.L0 satisfies these hypotheses conditionally on the admitted full record, and the same E-linear operators preserve its actual AA multiple. Choosing auxiliary doubled horizontal and integral imaginary root bases gives c=1 in LIFT094: t=1+3M Q\_S, delta=t\^{}(4fr\^{}2), L=t; any separately required larger congruence modulus is retained. This removes the unevaluated auxiliary clearing integer without a replacement lattice, seed translation or extra isogeny. The actual very ample line H, same-rule bad-locus degree D, nonempty family intersection, component-entry multiplier N0, selected nonzero rule and full Host/Sizes remain unresolved inputs. No evaluated total N, period, specific032, DP44 or RH conclusion; no finite charge. Current DO117121/31 arrays and floors27/11 remain unchanged.

\textbf{Complete author subject; final whole-proof/source/script verification and exact replay are recorded separately. Independent whole review is pending.} This is the scoped source application AM-LIFT-095 of AM-LIFT-094. It leaves the actual very ample line, bad-locus degree, component-entry cost, nonzero rule and Host/Sizes construction unresolved. Written 23 November 2026 UTC by \nolinkurl{/root/lift_accounting}.

\subsection{1. Statement and exact scope}\label{statement-and-exact-scope-1}

Let E/F be a CM extension unramified at every finite place (split factors are allowed). Let (V,h) be a nondegenerate Hermitian E-space of rank r≥3, conjugate-linear in its first variable, indefinite at every real place of F. Let L⊂V be a full projective O\_E-lattice with L=L\^{}sharp. Then there exist e,f∈L with

\[h(e,e)=h(f,f)=0,\qquad h(e,f)=1.\tag{IH1}\]

and an integral hyperbolic splitting

\[L=(\mathcal O_E e\oplus\mathcal O_E f)\perp A.\tag{IH2}\]

written in the ordered coordinates O\_E e ⊕ A ⊕ O\_E f,

where A=L∩(Ee+Ef)\^{}perp is a full projective lattice in the complement. For any nonzero fractional ideal a of E, the same rational hyperbolic frame has integral root operators and projections on aL. In the notation of LIFT094 HE14--16, one may choose the auxiliary horizontal Q-basis in 2A and the central Q-basis in O\_E\^{}-; the sufficient clearing integer is then c=1. Without the doubling, c=2 suffices.

This chooses auxiliary root directions inside the original rational Hermitian space. It does not replace a prescribed source basis, lattice, polarization, level, finite type, or CM point. The statement is only for the indefinite stages; the definite terminal stage is excluded.

\subsection{2. Integral local hyperbolic pairs, including residue characteristic two}\label{integral-local-hyperbolic-pairs-including-residue-characteristic-two}

Fix a finite place v of F. First suppose E\_v/F\_v is the unramified quadratic field extension, with integer rings S/R. The finite projective dual commutes with localization and completion, so L\_v is again selfdual. A unimodular Hermitian S-lattice has nondegenerate reduction over the quadratic extension of residue fields. Such a reduction has a vector of nonzero norm. Indeed, if all diagonal values vanish, choose x,y with h(x,y)≠0. For a residue scalar a, the norm of x+ay is the nonzero-trace expression a h(x,y)+bar(a h(x,y)); surjectivity of the quadratic finite-field trace makes one such norm nonzero. This proof includes characteristic two. Lifting that vector gives a unit norm vector, which splits orthogonally. Repeating gives an orthogonal S-basis with all diagonal coefficients in R\^{}×.

The norm S\^{}×→R\^{}× is surjective. The residue norm is surjective; successive lifts use

\[N(1+\pi^j x)\equiv1+\pi^j\operatorname{Tr}(x)\pmod{\pi^{j+1}},\qquad j\ge1.\tag{IH3}\]

and the surjective residue trace. Completeness gives a lift at every precision, also in residue characteristic two. Thus the first two diagonal units u1,u2 admit s∈S\^{}× with N(s)=−u1/u2. The vector e=e1+s e2 is primitive and isotropic. Perfection of the lattice pairing gives f0∈L\_v with h(e,f0)=1. The trace S→R is surjective: it is nonzero modulo π, so its R-linear image contains a unit and hence equals R. Choose a∈S with a+bar a=h(f0,f0). Then f=f0−ae has norm zero and still satisfies h(e,f)=1. This gives an integral hyperbolic pair and a selfdual orthogonal complement.

If E\_v=F\_v×F\_v, its involution exchanges the factors. A selfdual Hermitian lattice is a pair of free R-modules in perfect duality; choosing a basis of one and its dual identifies it with the standard diagonal Hermitian lattice. The same construction works with the split norm (x,y)↦xy and trace (x,y)↦x+y, both surjective. Thus no dyadic exception or assumption that 2 is invertible enters. Only relative unramifiedness of E\_v/F\_v is used; F\_v may be absolutely ramified over Q\_2.

\subsection{3. One integral global pair by strong approximation}\label{one-integral-global-pair-by-strong-approximation}

The rational Hermitian space contains a hyperbolic plane by the dimension/signature/determinant classification already checked in LIFT094 §7: Deligne (1982), Proposition 4.1, printed p.28 in the official transcription. Choose one rational hyperbolic pair e0,f0.

At every finite v, section 2 gives an integral hyperbolic pair ev,fv in L\_v. Witt extension gives a local unitary isometry taking (e0,f0) to (ev,fv). It can be made special unitary without changing that pair: its target orthogonal complement has dimension r−2≥1 and a nonisotropic line. Multiplication on that line by the inverse determinant is a unitary map because the determinant has norm one; this changes the determinant to one. In the split case the same assertion follows from the perfect dual decomposition. No integrality of this determinant-correction map is required; it fixes ev,fv exactly.

For all but finitely many finite places, e0,f0 already lie in L\_v, so the identity is sufficient. At each remaining place, the condition that g e0 and g f0 lie in L\_v defines a nonempty open subset of SU(V)(F\_v). At the other places use the compact open stabilizer of L\_v. Their restricted product is a nonempty open subset of SU(V)(A\_F,f).

Over an algebraic closure, SU(V) becomes SL\_r. Thus SU(V) is connected, simply connected, and absolutely almost simple of type A\_(r−1); its product at the infinite places is noncompact, since the Hermitian signatures are indefinite. Rapinchuk (2013), Theorem 2.3, printed p.281, therefore gives density of SU(V)(F) in SU(V)(A\_F,f). Equivalently apply Milne (2017), Theorem 4.16, printed p.48, to Res\_(F/Q)SU(V), with the noncompact-type definition at p.40. Choose g in the indicated open subset. Then e=g e0 and f=g f0 belong to every completed L\_v and hence to L itself, because a full Dedekind lattice is the intersection of its finite localizations inside V. Their rational pairing is unchanged.

This argument uses strong approximation on the group and local transitivity on ordered hyperbolic pairs. It does not assert strong approximation for an arbitrary homogeneous space, nor does it assert that every unitary lattice genus has one class. The construction is inside the fixed lattice L and the same rational Hermitian space.

\subsection{4. Integral projections and all root directions}\label{integral-projections-and-all-root-directions}

With the source convention that h is conjugate-linear in its first variable, define

\[P_+(x)=e\,h(f,x),\qquad P_-(x)=f\,h(e,x).\tag{IH4}\]

Selfduality implies h(L,L)⊂O\_E, so these maps preserve L. They are complementary idempotent line projections with P\_+P\_-=P\_-P\_+=0 and P\_+\^{}dagger=P\_-. Each isotropic line projection is not asserted to be self-adjoint. Their sum is the h-orthogonal projection onto the hyperbolic plane. The residual projection 1−P\_+−P\_- identifies the complement A, giving the decomposition in section 1. The complement is projective and full; it need not have a free O\_E-basis.

For v∈A, the horizontal root operator in coordinates (a,x,b)∈Ee⊕A\_E⊕Ef is

\[X(v)(a,x,b)=(-h(v,x),bv,0).\tag{IH5}\]

Its cube is zero, and X(v)\^{}dagger=−X(v). It preserves L, as does its opposite obtained by exchanging e,f.~For z∈O\_E with bar z=−z, Z(z)(a,x,b)=(zb,0,0) and its opposite preserve L. Choose a Q-basis of the rational vector space A\_E from the underlying integral lattice A, and double it. This is still a Q-basis, and every resulting X(v) and opposite lies in 2 End\_OE(L). A Q-basis of E\^{}- can likewise be chosen in O\_E\^{}-: differences a−bar a of integral elements span E\^{}- over Q.

Every map here is E-linear. Consequently any such map preserving L also preserves the O\_E-span aL for every fractional ideal a: it commutes with each scalar in a. Thus all of LIFT094's exact requirements on cP\_±, cX(v), cX\^{}op(v), cZ(z), cZ\^{}op(z) hold on aL with c=1. Only auxiliary Q-basis choices are being made; no change to the physical source frame is inferred.

\subsection{5. Literal full-source application}\label{literal-full-source-application}

The actual source has hT:V→ₗ⋆{[}E{]}V→ₗ{[}E{]}E (\nolinkurl{DataSpaces.lean}, line34), with L0 a full lattice and \nolinkurl{baseLattice_isUnimodular} at lines44--46. The defining \nolinkurl{IsUnimodular} is \nolinkurl{dual h L=L} in the retrieved exact \nolinkurl{agents/lift_sources/Formalism/HermitianLattices/StatementLat.lean}, lines31--38. The imported \nolinkurl{Formalism.E.hdual} is the ordinary integral dual, namely the intersection of the inverse images of O\_E under h(y,−) for y∈L; this is the freshly retrieved pinned \nolinkurl{Formalism/FockGram/StatementHerm.lean}, lines9--12, stored with the present source packet. \nolinkurl{Spaces.cm} supplies the CM/unramified input (\nolinkurl{DataSpaces.lean}, lines14--19,79). The \nolinkurl{deg2} and \nolinkurl{fixed} clauses identify F with the fixed field of the source conjugation. Its literal \nolinkurl{IsCMUnram} in \nolinkurl{StatementMass.lean}, lines44--48, therefore rules out every ramified finite prime of this actual quadratic extension: conjugation would otherwise be trivial on the residue field, contradicting the stipulated integral θ−bar θ outside that prime.

For every 1≤j\textless m, \nolinkurl{Spaces.W_rank} and \nolinkurl{W_sigval} (\nolinkurl{DataSpaces.lean}, lines121--124), together with \nolinkurl{Ranks.sigW} (\nolinkurl{Params.lean}, line46), give rank r=k+m−1 and signatures (k+j−1,m−j), both positive. The exact rank hypotheses m≥2, k≥m+3 and k even give k≥6 and r≥7 (\nolinkurl{Params.lean}, lines10,13--16,32). Thus all preceding hypotheses apply to the original (W\_j,hT,L0).

The exact \nolinkurl{IsDatW} condition (\nolinkurl{PartitionExtendByZeroAntitone.lean}, lines24--31) identifies the pinned datum lattice with the O\_E-span of AA·L0 by the original injective O\_E-linear map, and the alternating form with ±Tr\_(E/Q)(xi0 hT). The unitary exponentials built from the root operators, and the rational unitary element a of HE16, transport through this fixed identification and preserve that alternating form. The individual projections and nilpotent operators are not asserted to be isometries. Multiplication by the fixed scalar xi0 and the sign does not change the unitary group; the original finite level is preserved by the HE17 congruences, not by an additional assumption on the residues of e or f.~No replacement of AA by a principal ideal is made.

Accordingly LIFT094 may use, on every such indefinite source family,

\[t=1+3M Q_S,\qquad\delta=t^{4fr^2},\qquad L=t.\tag{IH6}\]

where Q\_S is the product of the protected rational primes. Its complete density, branch-isogeny, polarization, exact M-level and all-object field arguments then apply without an unevaluated clearing integer. If the protected set consists only of the primes dividing M, t=1+3M rad(M). If additional exact finite congruences require LIFT094 to use a larger auxiliary modulus L\_0, the same proof gives c=1 and t=1+L\_0 Q\_S, with no removal of those congruences. The strong-approximation element was used only to choose a rational hyperbolic pair inside the unchanged lattice; it is not an extra isogeny step or a translation of the original seed, so it contributes no component-entry degree. Its rational-coordinate height is not bounded or required for this denominator assertion.

The numerical very ample line, the degree D of the same-rule bad locus, its nonempty intersection with this component, and any entry cost N0 remain exactly as in LIFT094. No effective total N, admitted Host/Sizes construction, nonzero selected image, or DP44 inequality is claimed.

\subsection{6. Fresh primary locators and verification boundary}\label{fresh-primary-locators-and-verification-boundary}

Milne, J. S. (2017). \emph{Introduction to Shimura varieties} (revised September16,2017), Definition3.18 p.40, Theorem4.16 p.48. \url{https://www.jmilne.org/math/xnotes/svi.pdf} . These statements were freshly read on23November2026; web references turn905677825view0 and turn613764452view1/2.

Rapinchuk, A. S. (2013). Strong approximation for algebraic groups. In \emph{Thin groups and superstrong approximation}, MSRI Publications61, pp.269--298. Theorem2.3, p.281. \url{https://library.slmath.org/books/Book61/files/70rapi.pdf} . The exact theorem and hypotheses were freshly read, web references turn769642170view0 and turn905677825view3.

Deligne's rational Hermitian classification is the carried checked LIFT094 input, not a newly downloaded source. The adjacent source manifest records every exact local source fingerprint and its actual reading scope. The companion certificate checks finite local norm/trace and hyperbolic-pair models, all root/projection identities in oblique integral frames, the exact c=1 congruences, and the preserved polarization multiplier. These diagnostics are not a computational proof of strong approximation or a construction of the original source. Final saved-file author and independent checks have separate timestamped records. No finite analytical error charge or proof of the full source application is asserted.

\subsection{Canonical algebraic-geometric context}\label{canonical-algebraic-geometric-context-52}

The construction is an integral orthogonal direct-summand decomposition of the same finite projective Hermitian module over Spec(O\_E), followed by integral E-linear root endomorphisms and the same fractional-ideal multiple. It retains the source alternating form signed Tr(xi0 h), the original polarization and exact M-level. The resulting isogenies use the unchanged rigidified Poincare/cubical-line transport of LIFT085/LIFT094. Fresh primary group input: Rapinchuk(2013),Theorem2.3,p281; Milne(2017),Definition3.18,p40 and Theorem4.16,p48. Carried rational classification: Deligne(1982),Proposition4.1,p28. Carried projective-module/exterior context: SGA3 I4.7.3/4.7.3.1 and SGA6 VI Theorem3.3; Picard/biextension context: FGA232 Theorem3.1 pp148-149,FGA236 Corollaries1.7/2.3 pp229/232,SGA7 I VII2.7/2.9 pp169-174 and VIII7.4(b) pp310-311. Formal-crystal/Cartier-duality anchors remain cote4 scan003 sections8.1.1-8.1.2,scans005-006 sections8.4.1-8.4.3 and cote54 scans005-007. No new formal slope, motivic comparison, l-adic weight or global oscillator map is inferred. DP\_MULT Q3e/DP44/DP52/DP54-65 and THETA3 G1/A7 receive a sharper conditional actual-source cost interface; CX\_GRH CX-ALG-024 and GCT\_PLETH T-052 keep their separate global gates.

\clearpage

\section{AM-DO-118 --- the complete fixed-block rank-two unipotent sector at support27}\label{am-do-118-the-complete-fixed-block-rank-two-unipotent-sector-at-support27}

\textbf{Canonical checked integration --- 2026-11-23T07:55:49.779104+00:00.} Complete root and independent proof/script/interface reviews and separate exact replays passed. The actual canonical predecessor has121/31 entries; this integration gives122/31, complete floors27/11. Frozen author history is preserved. AM-DO-118 UB1-29 removes the entire actual conditional J56 fixed-S globally rank-two-unipotent support27 sector, all original profiles and labels. The rank35 high envelope retains the full five-column and same-pair isometry fibre. In the rank34 low branch, the true row metric, full beta norm fibre, affine central cosets and actual selected trace are retained; fixing row and beta before eta-dependent labels permits the one-label bound without global eta uniqueness. The six complete strict dyadic bounds have exponents8792,4491,3928,23649,17510,11371; their union is below2\^{}-3927. Append3927 once to DO117121/31, giving122/31 and both full sums below2\^{}-230; complete floors27/11 remain unchanged. This does not remove the other support27 classes, establish a support28 floor, or solve the signed inequality, actual admitted source, selected arithmetic content, prescribed period, specific032 pairing, evaluatedDP44 or RH.

Author: \nolinkurl{/root/definite_oscillator}, ALP\_MACHINE. 23 November2026 UTC. Controlled pending proof. Saved-file author verification, exact replay, independent whole review and canonical integration are separate events. Existing canonical records and earlier frozen findings are unchanged.

\subsection{1. Actual scope and theorem}\label{actual-scope-and-theorem-6}

Retain the actual J56 all-top normalized kernel at the charged cutoff H=133, f=a=168, n=169, its original E/F, paired CM components, admitted source/target pencils and omitted endpoints, finite projectors, selected Hilbert metric, favorable packet point, central representative, phases, class masses and Haar measures. Work in the original K068 finite-unipotent remainder. The sector considered is

\[
 h=I_S\oplus h_0=I+N,\quad \dim S=164,\quad
 \operatorname{rank}_E N=2,\quad N^3=0,\quad
 \operatorname{rank}_E N^2=1,\qquad s=27.                 \tag{UB1}
\]

These are the signature-one globally rank-two-unipotent fixed-S terms. Global unipotence and fixation of S are explicit hypotheses, not inferred from the finite reductions. At the334 unmixed relative G places put d\_v=rank(N mod v) in\{0,1,2\} and D\_G=sum\_v d\_v. Write i,u,t for the numbers of actual types I,II,III, so i+u+t=27.

\textbf{Theorem.} The total absolute actual normalized contribution of the entire UB1 sector, including every actual support, incidence profile, paired label, finite pattern, covector, phase and central fibre, is strictly less than2\^{}-3927. Its intersection with the current remainder can therefore be removed with one fresh charge3927, retaining all earlier charges.

The proof separates D\_G\textgreater=35 from D\_G\textless=34. In the latter case it obtains the original packet-plane normal form, then separates the last-only transverse pivot from short pivots with at most24, exactly25, exactly26 or exactly27 eligible type-II places. The new count order retains the complete relative norm fibre before summing local labels. It requires no global uniqueness theorem for eta, no denominator bound on eta, and no new local active-II eligibility theorem.

No complete support28 floor, scalar-cross removal, signed lower bound, actual selected arithmetic image, prescribed period, literal source constructor, specific032/Read or DP44 conclusion is asserted.

Use the unchanged host constants

\[
 g_0=22285932761,\quad g_1=64072175441,\quad Q=g_0g_1,
 \quad U=77Q^2-14,\quad R=17424/133,\quad M=\sqrt{R/U},
\] \[
 L_T=1019333/1024>\log_2T,\quad L_Q=287843/4096<\log_2Q,
 \quad L_U=601356/4096<\log_2U,\quad L_R=113/16>\log_2R.
                                                               \tag{UB2}
\]

Every direction is checked against the literal J56 prime T by exact integer powers. In particular2\^{}995\textless T\textless2\^{}996 and2\^{}34\textless g0\textless g1\textless2\^{}36.

\subsection{2. High finite rank, the low-rank plane and the exact type budget}\label{high-finite-rank-the-low-rank-plane-and-the-exact-type-budget}

Put alpha=(N\^{}2)\_nn.~The unchanged unitary identity and original paired ideals give

\[
 0<\alpha_\tau\le R,\qquad
 \mathfrak J_F=\prod_v\mathfrak p_v^{2-d_v},\qquad
 J=N_F\mathfrak J_F=Q^{334}/\prod_vq_v^{d_v}.
                                                               \tag{UB3}
\]

The complete R058(HB3--14) envelope is support-independent in its displayed formulas. With z0=(JT\^{}-2s)\^{}(1/(2f)), it retains the full original negative-column count2\^{}5379 T\^{}(8s+r0), fixed-column fibre G5\textless2\^{}107, gamma count2\^{}1062((375+R/(2z0\^{}2))/376)\^{}f, actual selected trace, all support and type assignments, original zero-target coarea, both mixed factors and3\^{}334\textless2\^{}530 G patterns. Its increasing z0 part is

\[
 z_0^{334f}\left(\frac{375+R/(2z_0^2)}{376}\right)^f
 =z_0^{332f}(375z_0^2+R/2)^f/376^f
 <2^{42}z_0^{55776}\quad(0<z_0\le1).                    \tag{UB4}
\]

The actual profile comparison is at the same J: the square-pole exponent is mu=2u+t, the gamma factor decreases when its norm floor increases, and the remaining pair/coarea and real factor is at most2\^{}(i+t)T\^{}(-t/2)\textless=2\^{}s. This is the original R058(HB9--11) comparison of upper bounds, without positing an all-II matrix.

If D\_G\textgreater=35, then J\textless=Q\^{}334/g0\^{}35. The strict integer comparison

\[
                 Q^{334}2^{21\cdot1499}<g_0^{35}T^{54}   \tag{UB5}
\]

gives z0\textless2\^{}(-1499/16). Thus the complete high-rank contribution has logarithmic upper bound

\[
\begin{split}
 E_H={}&1+5379+1062+107+168+3\cdot27+530-464f\\
       &-335\cdot167L_Q+343\cdot27L_T+42-55776\cdot1499/16\\
     ={}&-36014199/4096<-8792.                            \tag{UB6}
\end{split}
\]

All its type and pattern unions are already included; no later multiplicity is attached to UB6.

Suppose D\_G\textless=34. The original integral source/target pairings of DO078(FP20--27), retained in R058(HB20--22), give b\_j=\textless y\_j,N\^{}2 y\_j\textgreater{} in O\_F intersect J\_F. They use separate membership in every actual source and target module at both CM partners, not a common stable lattice. The old short-vector maximum is L\_old=4030172803672366725045. The exact inequality

\[
 Q^{334}>g_1^{69}(R L_{\rm old})^{168}                   \tag{UB7}
\]

and J\textgreater=Q\^{}334/g1\^{}34\textgreater=Q\^{}334/g1\^{}69 force all three pairings to vanish. Hence the rank-one image of N\^{}2 is contained in the original packet plane P=Eu perpendicular Ee\_n.~The original entry (N\^{}2)11=−4alpha/U and its exact dyadic/G/T coprimalities give

\[
 \operatorname{im}N^2\subset P,\qquad
 \alpha\in U\mathfrak J_F I_{\rm II,F}^2I_{\rm III,F}.   \tag{UB8}
\]

This is the support-free plane implication; none of the old numerical emptiness ranges is silently extended. The fresh integer inequality

\[
 (133U)^{168}Q^{334}>17424^{168}g_1^{34}T^{45}            \tag{UB9}
\]

excludes mu\textless=45 by N\_F(alpha)\textless=R\^{}f.~Consequently every remaining low-rank profile satisfies

\[
 2u+t\ge46,\qquad 2i+t\le8,\qquad i\le4.               \tag{UB10}
\]

There are25 possible type triples: i=0,..,4 and t=0,..,8−2i, with u=27−i−t. We retain the larger full assignment count3\^{}27\textless2\^{}54 throughout.

\subsection{3. Original normal form, transverse metric and containing ideals}\label{original-normal-form-transverse-metric-and-containing-ideals}

Let V=u\^{}perp in the original positive E\^{}4 and Z=S perpendicular V. Since the quotient ell\^{}perp/ell for ell=im N\^{}2 is positive definite, its unipotent unitary action is identity. Thus h\_ZZ=I and N(V) is contained in ell. The cross map C:V→P is nonzero; otherwise N acts on P alone and has square zero. R077(FB8)/R079(PC6) give

\[
 C=\ell r,\quad\ell=(1,\eta),\quad\eta\bar\eta=U,
 \quad v=J_Z^{-1}r^*,\quad q=rv>0,\quad\alpha=Uq,
\] \[
 N=\begin{pmatrix}0&-v\ell^*J_P\\
                 \ell r&d\ell\ell^*J_P\end{pmatrix},
 \quad J_P=\operatorname{diag}(U,-1),\quad
 d+\bar d=-q,\quad w=N_{nn}=-Ud.                         \tag{UB11}
\]

Here r vanishes on S. C and w reconstruct the whole original matrix, with no remaining column or isometry fibre; Re(w)=alpha/2 is fixed by C. In particular N\^{}2Z=0.

Use the actual primitive integer vector and fixed transverse vectors

\[
 u=(2,95192390527,383092671043,12529843548438452388609),
\] \[
 y_1=(-976627,-153032,38026,0),\quad
 y_2=(-567678,695570,-172838,0),\quad
 y_3=(256766,-7654233287,-30805128044,1).
\]

For a\_j=r(y\_j), the true covector norm is q=aG\^{}-1 a*, with

\[
 G=\begin{pmatrix}
 978665066829&441392856078&-303935233242\\
 441392856078&835948910828&95913020934\\
 -303935233242&95913020934&1007543201084982869062
 \end{pmatrix},\quad\det G=4U,\quad G>2^{16}I.           \tag{UB12}
\]

The fixed frame has determinant of modulus2U, a unit at G and T. Its metric is not replaced by an integral orthogonal frame. When a1 or a2 is nonzero, choose the first such pivot, put beta=eta a\_j and B=\textbar\textbar y\_j\textbar\textbar\^{}2\textless2\^{}40. Retain the harmless common pivot allowance3\textless2\^{}2.

Call an active type-II place eligible if eta is congruent to a0=Q(−7+2sqrt(−7)) at its original selected component. Its partner condition is eta congruent to U/bar(a0), using eta bar(eta)=U and the actual inverse-adjoint line. Let p be the number of such places. The support-free original-input proof in DO109(SP26.8--10), which explicitly covers V rather than requiring C\textbar S!=0, gives the following facts.

Each y\_j lies in every original source and target module at both components. TypeI preserves their intersection, so Ny\_j is integral. At typeIII, Ny\_j lies in the target, while (h\^{}-1−I)y\_j=−Ny\_j lies in the distinct source because N\^{}2y\_j=0; the intersection is contained in the original integral lattice. Thus the cross entries are integral at I and III. At II a simple pole of a cross scalar forces its saturated image line to be the actual pencil pole line, hence forces eligibility at both partners. A full-entry double pole has the actual pencil pole line as its leading image; UB11 shows it can occur only in the rank-one PP block, so it likewise forces eligibility. Primitive integer Bezout extraction from u introduces no denominator. These are containing-ideal statements; they do not assert that every type-II place is eligible.

Pad the eligible set to v active places D. At G use the original canonical whole-C ideal J: rank0 gives both CM factors of E-norm q\_v\^{}2; rank1 gives one of its two orientations of E-norm q\_v; rank2 gives neither. One common orientation controls all a\_j and beta. Therefore

\[
 L=JI_D,\quad N_EI_D=T^{-2v},\quad z=(N_EL)^{1/(2f)},
 \quad a_1,a_2,a_3,\beta\in L,\quad
 z\le z_v=Q^{167/f}T^{-v/f},
\] \[
 |a(h)|\le2Q^{-335\cdot167}T^{335v}z^{335f}.             \tag{UB13}
\]

Both mixed finite factors retain their modulus bounds one. There is no independent orientation choice for each coordinate or each partner.

Let J0 be the F-product of rank-zero G primes once. Its original table gives N\_EJ\textless=Q\^{}167 N\_FJ0 and N\_FJ0\textless=Q\^{}167. Full entries have one pole at most everywhere in the active support Sigma and a second only on the eligible set. For fixed C and an actual labelled tuple the central differences therefore lie in

\[
 \Gamma=iJ_0I_{\Sigma,F}I_{D,F},\quad
 \operatorname{covol}(\Gamma)=\sqrt{D_F}N_FJ_0T^{-(s+v)},
 \quad D_F=2^{168}D_0,\quad D_0=7^{140}29^{144}5^{84}13^{84}.
                                                               \tag{UB14}
\]

This is one containing affine coset after C and the labels are fixed. The shift is retained; w itself is not asserted to lie in the difference ideal. Since O\_E intersect iF=iO\_F there is no half-lattice. For v\textgreater=24 the actual full-real codifferent cover is below750sqrt(f)(N\_FJ0 T\^{}-(s+v))\^{}(1/f)\textless sqrt(f), by750\^{}fQ\^{}167\textless T\^{}50 and s+v\textgreater=51. With \textbar Im(w)\textbar\textless133/2 and the full real volume V\_F(68)\textless2\^{}479,

\[
 \#w<2^{479}\frac{T^{s+v}}{N_FJ_0}
       \le2^{479}Q^{167}T^{s-v}z^{-2f}.                 \tag{UB15}
\]

All ramified and unselected primes retain the original full-entry integrality; the displayed fractional ideals contain no other denominators.

\subsection{4. Full row and norm fibre before the paired label sum}\label{full-row-and-norm-fibre-before-the-paired-label-sum}

The arbitrary-ideal three-coordinate cover in the true G\^{}-1 metric is below4sqrt(f)z, because(2250/7)\^{}2·3\textless16·2\^{}16. The exact covolume is D0\^{}3 z\^{}(6f)/(4U)\^{}f.~Since q\_tau\textless=M\^{}2, the full original row count is

\[
 \#r\le A_3(4U)^f(M+4z)^{6f}z^{-6f},\quad
 A_3=528^{3f}/((3f)!D_0^3)<2^{-4040}.                    \tag{UB16}
\]

For every fixed actual nonzero pivot, beta satisfies beta bar(beta)=U a\_j bar(a\_j) and beta in L. The integral relative norm ideal exists whenever a solution exists. K075(RN17--19) and DO091(TD11--12) give

\[
 \#\beta\le28C_F(RB)^{f/12}z^{-f/6},\quad
 C_F=\frac{2^{300}}{29^6 41^{14}2129^{14}}<2^{42},
 \quad28<2^5.                                         \tag{UB17}
\]

Indeed the relative ideal has F-norm at most(RB)\^{}f z\^{}-2f; each ideal has at most28 norm-one generators. The bound allows nonprincipal and non-conjugation-stable L, counting absence of a generator as zero. The full beta fibre is retained in every short-pivot case below.

The exact selected trace DO080(RT3), with both compact factors, original metric and original near-point transport, gives

\[
 N_F\alpha\ge(U/B)^f z^{2f},\qquad
 |c(h)|<2^{-929f/2}(U/B)^{-f/2}z^{-f}.                   \tag{UB18}
\]

This bound is applied only to actual nonzero pivots. Including artificial zero rows in UB16 is merely an enlargement of their number, not a trace claim for those rows.

We now change the order of the finite upper sum. Fix the support, its type assignment, padded eligible set, canonical G pattern, a nonzero original row r and one of the beta values counted in UB17. Then eta=beta/a\_j and hence the complete C are determined. No global uniqueness of eta across distinct rows or beta values is needed.

For clarity, the one-label argument of DO105(EL9--15) is algebraic at fixed eta. In the actual pencil A\_c=I+c pi\^{}-1 z q\_\emph{, an eligible saturated column w0=(u,eta) and paired row w0\^{}dagger have reductions a\_l z and b\_l q\_}, with a\_l,b\_l units. Set a=q\_\emph{w0/pi and b=w0\^{}dagger z/pi. Isotropy and q\_}z=0 give

\[
 b_l a+a_l b=0\pmod\pi,\qquad
 A_c^{-1}w_0=w_0-cza,\quad
 w_0^\dagger A_c=w_0^\dagger+cbq_* .                    \tag{UB19}
\]

The column and row reductions therefore vanish at the same unique possible label c\_eta. If a=0 neither vanishes; an omitted label gives no admissible active arrow. A pole in either cross row forces that label by applying N to an original Z input at the appropriate CM component. If both rows are integral but a central pole remains, the actual mixed terms are integral after conjugation: R0 A\_c=R0 and A\_c\^{}-1 V0=V0+czmV0. The residual term is d(A\_c\^{}-1w0)(w0\^{}dagger A\_c); its integrality forces one factor to vanish modulo pi and hence the same label. This covers central-only poles. The fixed V frame is a T-unit frame, so the argument has its literal original-input meaning here. No independent paired label, sign or partner multiplicity is introduced.

Thus every eligible active type-II place admits at most one of its original diagonal pairs, independently of the central parameter. Its true absolute weight is1/(T−1). At the other places K071(NB7--8) retains the full weighted ordered-pair bounds T for I and1 for II or III. Consequently, at fixed r,beta and prescribed type/eligible sets,

\[
 \sum_{\text{admissible labelled tuples}}|\rho_T|
                      \le T^i(T-1)^{-p}.                \tag{UB20}
\]

For each such tuple UB15 bounds its possibly different affine central coset uniformly. Therefore the central-and-label sum is at most that bound times UB20. Summing this uniform inequality over every row and its complete beta fibre proves the claimed order rigorously. It does not apply a label saving to a fixed-tuple matrix count before eta is determined, remove the beta fibre, or assume that different labels share a coset. All sums are over nonnegative absolute upper bounds. The original phases, class masses and Haar normalization retain their existing modulus bounds.

For the p\textless=24 branch we may discard the new factor(T−1)\^{}-p and use the original upper weight T\^{}i. For p=25,26,27 we retain it. The exact integer comparisons give

\[
             (T-1)^{-p}<2T^{-p}<2^{1-995p}
                         \quad(p=25,26,27).              \tag{UB21}
\]

\subsection{5. The complete weighted G sum and the short-pivot endpoints}\label{the-complete-weighted-g-sum-and-the-short-pivot-endpoints}

Multiply UB13 and UB15--18. Apart from constants independent of J, the complete z dependence is

\[
 F_v(z)=z^{\kappa}(M+4z)^{6f},\qquad
 \kappa=(335-6-1/6-2-1)f=(1955/6)f>0.                   \tag{UB22}
\]

For0\textless z'\textless=z, F\_v(z')/F\_v(z)\textless=(z'/z)\^{}kappa. Relative to the all-rank-zero cap, lowering a local ideal exponent from2 to1 multiplies z by q\_v\^{}-1/(2f). Both genuine rank-one orientations and the rank-two state give the exact local majorant

\[
 1+2q_v^{-1955/12}+q_v^{-1955/6}
                =(1+q_v^{-1955/12})^2.
\] \[
 \prod_{v=1}^{334}(1+q_v^{-1955/12})^2
           <(1+2^{-5534})^{668}<2.                       \tag{UB23}
\]

This includes every original finite pattern for one factor2. Enlarging the pattern sum beyond D\_G\textless=34 is valid after imposing the structural restrictions on the actual subset. There is no additional2\^{}668 orientation count in these short branches. The label factor UB20 is uniform in the fixed data and independent of this z comparison.

At the common caps the exact host comparisons give

\[
 z_{24}<M/4,\qquad z_{25}<M/256,\qquad z_{26},z_{27}\le z_{25},
 \qquad(65/64)^{6f}<2^{23}.                             \tag{UB24}
\]

Hence use radius2M for v24 and(65/64)M for v\textgreater=25. No cap-region union is needed. Define, retaining the full beta fibre,

\[
\begin{split}
 E(v,i,b_{\rm type},b_{\rm pad},b_{\rm rad})={}&
 1+1+168+b_{\rm type}+2+479+5+42-4040+2f+b_{\rm rad}-929f/2\\
 &+(37f/12)L_R-(5f/2)L_U+(7f/12)40\\
 &-(49\cdot167/6)L_Q+(49v/6+27+i)L_T+b_{\rm pad}.
\end{split}                                             \tag{UB25}
\]

The two initial1's are the Green Euler factor and weighted G sum. The2f is the4\^{}f determinant, the radius allowance is separate, and5+42 retains the entire relative norm/phase fibre. The central T\^{}(27−v) and type-I T\^{}i have already been included. Powers of R,U,B are combined with their final signs before using directed logarithms.

There are fewer than2\^{}168 supports and3\^{}27\textless2\^{}54 type assignments. For p\textless=24, pad to24 of27 places: binom(27,24)=2925\textless2\^{}12. For p25,26,27 use the actual eligible set, with binom(27,p)\textless=2925\textless2\^{}12 choices. These allowances are retained separately from the one-label saving. The four complete short branches satisfy

\begin{longtable}[]{@{}
  >{\raggedright\arraybackslash}p{(\columnwidth - 6\tabcolsep) * \real{0.2143}}
  >{\raggedright\arraybackslash}p{(\columnwidth - 6\tabcolsep) * \real{0.2143}}
  >{\raggedleft\arraybackslash}p{(\columnwidth - 6\tabcolsep) * \real{0.2857}}
  >{\raggedleft\arraybackslash}p{(\columnwidth - 6\tabcolsep) * \real{0.2857}}@{}}
\toprule\noalign{}
\begin{minipage}[b]{\linewidth}\raggedright
Eligible II count
\end{minipage} & \begin{minipage}[b]{\linewidth}\raggedright
Parameters before the label gain
\end{minipage} & \begin{minipage}[b]{\linewidth}\raggedleft
Complete exponent
\end{minipage} & \begin{minipage}[b]{\linewidth}\raggedleft
Strict bound
\end{minipage} \\
\midrule\noalign{}
\endhead
\bottomrule\noalign{}
\endlastfoot
p\textless=24 & E(24,4,54,12,1008) & −96557629/24576 & \textless−3928 \\
p=25 & E(25,2,54,12,23)+1−995·25 & −581207129/24576 & \textless−23649 \\
p=26 & E(26,1,54,12,23)+1−995·26 & −143444991/8192 & \textless−17510 \\
p=27 & E(27,0,54,12,23)+1−995·27 & −279462817/24576 & \textless−11371 \\
\end{longtable}

\[
          \log_2 U_p<E_p\quad(p\le24,25,26,27).          \tag{UB26}
\]

The first line uses UB10. In the other lines u\textgreater=p implies i\textless=27−p.~The larger full type assignment factor is kept even when the actual profile is unique. All original pivot choices remain in2\^{}2. There is no additional union over p0 through24 in the first containing padded branch.

\subsection{6. The last-only pivot and the complete union}\label{the-last-only-pivot-and-the-complete-union}

If a1=a2=0, then a3!=0. Its exact Schur-complement metric is

\[
 q=|a_3|^2/W,\quad W=(G^{-1}_{33})^{-1}
      =4U/\det G_{\{1,2\},\{1,2\}},\quad2^{16}<W<2^{70}.
                                                               \tag{UB27}
\]

Use R079(PC22--24) with the original all-support ideal L=JI\_Sigma and its original central ideal iJ0 I\_I,F I\_II,F\^{}2 I\_III,F. The central factor is2\^{}479 Q\^{}167 z\^{}-2f T\^{}(u−s), and the full paired incidence bound T\^{}i leaves T\^{}-t\textless=1. No eligibility restriction is used in this branch.

The exact one-coordinate lattice has covolume D0 z\^{}(2f)/W\^{}f, cover below4sqrt(f)z and count A1 W\^{}f(M+4z)\^{}(2f)z\^{}-2f, where A1=528\^{}f/(f!D0)\textless2\^{}-1083. Its scalar norm fibre and actual trace are UB17--18 with B=W. The exact host inequality R\^{}fT\^{}54\textgreater25\^{}fU\^{}fQ\^{}334 gives z\_max\textless M/5 and radius below2M. Retain all2\^{}668 original G patterns, full type2\^{}54, support2\^{}168 and the harmless common pivot2\^{}2. The resulting complete endpoint is

\[
\begin{split}
 E_L={}&1+668+168+54+2+479+5+42-1083+2f-929f/2\\
 &+(13f/12)L_R-(3f/2)L_U+(19f/12)70\\
 &-(25\cdot167/6)L_Q+(31\cdot27/6)L_T\\
 ={}&-110371081/24576<-4491.                            \tag{UB28}
\end{split}
\]

This is the same support-free complete formula with a new exact support27 evaluation, not a reuse of an old numerical range. The last-only coordinate is counted in its own genuine lattice; no zero-coordinate density multiplier is inserted.

The high branch, the low last-only branch and the four low short branches exhaust UB1. Combining UB6, UB26 and UB28 gives the exact strict dyadic inequality

\[
 2^{-8792}+2^{-4491}+2^{-3928}+2^{-23649}
              +2^{-17510}+2^{-11371}<2^{-3927}.           \tag{UB29}
\]

Every branch is a complete bound on its indicated subset. Every support/type/label/phase/central/fibre factor has been counted above, and no additional branch multiplier is needed. This proves the theorem.

\subsection{7. Verification, chronology and retained geometric scope}\label{verification-chronology-and-retained-geometric-scope}

The companion verifies the literal host and directed powers, high-rank cap, budget69 plane comparison, low-D\_G34 pole-budget45 contradiction, all25 type triples, the true Gram matrix and Schur complement, original frame units, real/complex covolumes, all radius and weighted-pattern inequalities, actual paired cancellation models and a finite fibre-order diagnostic. It checks all six exact endpoints, the strict union and both full chronological error arrays. These diagnostics supplement the complete source/metric/normal-form proof; they do not enumerate actual global arrows or evaluate a period.

Direct source inputs are R058(HB2--24), R077(FB6--17), R079(PC6--24), DO109(SP26.3--18), DO105(EL9--16), K071(NB5,7--9), K075(RN17--19), DO091(TD11--12) and DO080(RT3). The full original pencil and paired inverse-adjoint interfaces in those proofs are retained. The new steps are the rank35/34 split at27, the actual mu\textgreater=46 restriction, the row-plus-beta-before-label sum UB19--21, and the complete six-branch count. K112's newer local eligibility lemma is not a premise.

At integration append3927 once to the then-actual ordered all-top error list, retaining every prior charge and the unchanged uniform list. The script accepts \nolinkurl{--prior} for an unchanged-byte chronological replay. Its historical base is the canonical R094120/31 record, complete floors27/11. The default actual predecessor is \nolinkurl{checks/do117_joint_remainder_error.json}, with121/31 entries after the distinct fixed-S nonunipotent29365 charge. Every later predecessor must extend that complete121-entry prefix. The present staged arrays have122/31 entries and retain both complete floors27/11. This theorem removes only the stated structural class at27, so neither complete floor changes.

The exact covector and affine ideal lattices are modules over the original Spec(O\_E) and Spec(O\_F), with the same CM involution, Hermitian polarization, theta line, biextension and complex metric. The carried module/exterior anchors remain SGA3 I4.7.3/4.7.3.1 and SGA6 VI Theorem3.3. The Picard/biextension anchors remain FGA232 Theorem3.1 pp148--149, FGA236 Corollaries1.7/2.3 pp229/232, and SGA7 I VII2.7/2.9 pp169--174 and VIII7.4(b) pp310--311. The inherited formal-crystal and Cartier-duality context is cote4 scan003 §§8.1.1--8.1.2, scans005--006 §§8.4.1--8.4.3 and cote54 scans005--007. These contextual sources are not claimed as newly read here.

No new lambda-ring equality, K-theory comparison, l-adic or motivic statement is inferred. The signed remainder, actual mixed source and prescribed global period, selected arithmetic image/content, specific032 matching/value, all actual DP44 terms and the separate RH-dependent conclusions remain as independently recorded. No Lean build, archive audit or change of selected projector, weights, fields or primes is performed.

\subsection{Canonical algebraic-geometric context}\label{canonical-algebraic-geometric-context-53}

The count uses the same integral Hermitian modules and original finite correspondences. The row and central difference ideals remain fractional projective modules on Spec(O\_E) and Spec(O\_F); the exact relative norm fibre includes all ideal classes and norm-one phases. The metric is the true fixed Gram/Schur metric, and the actual selected theta trace retains its Hilbert normalization and packet transport. Carried exact anchors are SGA3 I4.7.3/4.7.3.1; SGA6 VI Theorem3.3; FGA232 Theorem3.1 pp148-149; FGA236 Corollaries1.7/2.3 pp229/232; SGA7 I VII2.7/2.9 pp169-174 and VIII7.4(b) pp310-311; cote4 scan003 sections8.1.1-8.1.2 and scans005-006 sections8.4.1-8.4.3; cote54 scans005-007. No new primary reading, formal slope, lambda/K-theory relation, l-adic weight or CM-motive comparison is inferred. DP\_MULT Q3e/DP44 and THETA3 G1/A7 receive only this smaller conditional finite remainder; CX\_GRH CX-ALG-024 and GCT\_PLETH T-052 retain their separate global gates.

\clearpage

\section{AM-K-113 --- a complete all-II support27 containing-sector reduction}\label{am-k-113-a-complete-all-ii-support27-containing-sector-reduction}

\textbf{Canonical checked integration --- 2026-11-23T08:05:02.778598+00:00.} Complete root and independent proof/script/interface reviews and separate exact replays passed. The actual canonical predecessor has122/31 entries; this integration gives123/31, complete floors27/11. Frozen author history is preserved. AM-K-113 SC27.1-27 removes seven complete containing groups in the actual conditional J56 all-II global-rank-two-unipotent scalar-cross support27 sector. The33 complete endpoints retain all four transverse mixture terms, every G pattern, all inactive-zero subsets, chosen active-integral positions, coordinate sizes and a separately enumerated infinite eta-ideal tail. Their full union is below2\^{}-304. The necessary residual is b=j0,2\^{}2111\textless x\textless2\^{}3328, either total eligibility27 and100\textless=k\textless=164, or total eligibility28,k164 with x\textgreater2\^{}3107. Append304 once to actual DO118122/31, giving123/31 with both full sums\textless2\^{}-230; complete floors27/11 remain unchanged. No whole support27 exhaustion, actual admitted Host/Sizes, signed inequality, selected arithmetic content, prescribed period, specific032 pairing, evaluatedDP44 or RH conclusion.

Author: \nolinkurl{/root/kummer_extension}, 23 November2026 UTC. Controlled pending proof. Whole author, independent and root verification are separate events. No canonical artifact is changed by this file.

\subsection{1. Original sector and precise conclusion}\label{original-sector-and-precise-conclusion}

Retain the actual J56 all-top H133 kernel, original field and CM involution, paired source/target pencils and omitted labels, class masses, finite projectors, central phases, Haar measures, selected vector-valued real metric, every allowed word and canonical twist. Consider the globally rank-two-unipotent scalar-cross branch at exactly27 active type-II places, with nonzero cross map on the original164-dimensional S block. The accepted DO099 plane-C reduction through70 and the support-independent full-positive argument of K112/ K103 give \[
 E^{169}=Z\perp P,\quad Z=S\perp V_3,\quad P=Eu\perp Ee_n,
 \quad h=I+N,\quad h_{ZZ}=I,\quad N^2Z=0,\quad N^3=0.
 \tag{SC27.1}
\] All334 unmixed G reductions are unipotent. Write the original isotropic line as E(1,eta) in the fixed plane coordinates, with eta bar(eta)=U. Let x=N\_E D\_eta be its exact denominator norm; r be total eligibility among168 selected T places; k the number of nonzero original S covector values; b the number of inactive places at which both paired C blocks vanish; and j the number of active places at which the whole paired C is integral.

\textbf{Theorem.} A union of complete containing subsets of this sector has total absolute contribution less than2\^{}-304. Outside that union, necessarily \[
 b=j=0,\quad 2^{2111}<x<2^{3328},\quad
 \begin{cases}
 r=27,\quad100\le k\le164;\quad\text{or}\\
 r=28,\quad k=164,\quad x>2^{3107}.
 \end{cases}
 \tag{SC27.2}
\] There is one fresh combined304 charge. This is not a full support27 exhaustion: fixed-S and other structural sectors are separate, and the displayed all-II remainder is not shown empty. No support26 numerical charge or denominator window is imported.

The physical constants remain \[
 f=a=168,\ H=133,\ R=17424/133,\ M^2=R/U,\
 Q=1427908193729121522601,\ U=77Q^2-14,
 \quad D_0=7^{140}29^{144}5^{84}13^{84},\quad D_F=2^{168}D_0.
 \tag{SC27.3}
\]

\subsection{2. Eligibility, bounded uniqueness, and exact original ideals}\label{eligibility-bounded-uniqueness-and-exact-original-ideals}

K112 UE1--12 proves that every active type-II place is eligible, without a denominator restriction. Thus all27 active places are eligible. Its proof uses the actual selected-component Z-by-P block and is independent of a numerical support count; it is applied here at its canonical checked scope.

Put B={[}4(U+14){]}\^{}168. The exact original denominator-cleared norm argument gives \[
 7^{84}T^r<Bx,\qquad
 (4U)^{168}2^{24000}<T^{54},\qquad7^{84}T^{37}>B2^{12000}.
 \tag{SC27.4}
\] Consequently, below x=2\^{}12000 there is at most one eta for each fixed27-place active support, and27\textless=r\textless=36. The uniqueness uses the difference of two eta values, hence divisibility at all54 paired CM primes. The first inequality uses eta−a0 only at selected primes and is not doubled. Fresh exact lower bounds x\textgreater2\^{}(a\_r) are

\begin{longtable}[]{@{}rrrrrrrrrrr@{}}
\toprule\noalign{}
r & 27 & 28 & 29 & 30 & 31 & 32 & 33 & 34 & 35 & 36 \\
\midrule\noalign{}
\endhead
\bottomrule\noalign{}
\endlastfoot
a\_r & 2111 & 3107 & 4102 & 5098 & 6093 & 7088 & 8084 & 9079 & 10075 & 11070 \\
\end{longtable}

Every table entry is checked by the integer inequality7\^{}84 T\^{}r\textgreater B2\^{}(a\_r). No denominator-ideal or root-of-unity multiplicity is counted in this bounded unique-eta window.

Fix the original canonical whole-C G pattern J and rank-zero real ideal J0. As in DO098 CF10 and DO104 CD11--12, \[
 |a_G(h)|\le2Q^{-335\cdot167}(N_EJ)^{335/2},\qquad
 N_EJ\le Q^{334},\qquad N_EJ\le Q^{167}N_FJ_0.
 \tag{SC27.5}
\] The four local states have paired exponent2,1,0 and multiplicities1,2,1. They are the original common orientation states, not separate coordinate patterns.

Every original integral Z input, including the integral saturated transverse basis, satisfies r(y),eta r(y) in J I\_Sigma, where N\_E I\_Sigma=T\^{}-54. Valuationwise intersection therefore gives membership in J I\_Sigma D\_eta. Each inactive paired-zero place adds its full paired ideal, of E-norm T², even when eta has a local pole. If one active integral-C place is selected, it adds one more such ideal, since eta is a unit there. These are the support-independent local arguments DO107 IP4--10 and K103 SC26.17--19, not their old numerical applications.

Let e denote the number of added paired ideals. For the b\textgreater0 group take all b inactive ideals, so e=b. For the b=0,j\textgreater0 group select one active place, so e=1. Otherwise e=0. For the actual row ideal L, \[
 N_EL=z^{2f},\quad N_EJ=z^{2f}T^{54-2e}/x,\quad
 kz^2\le M^2,\quad q=M^2/z^2\ge k.
 \tag{SC27.6}
\] Every one of the167 fixed covector values belongs to this same L. There is no extra beta norm fibre: eta, this complete covector and the central scalar reconstruct the original h.

For fixed eta,row and pattern, the original globally normalized central scalar has one affine difference coset in \[
 \Gamma_d=iJ_0B_\eta I_{\Sigma,F}^2,\quad N_FB_\eta=x,
 \quad\operatorname{covol}\Gamma_d=\sqrt{D_F}(N_FJ_0)xT^{-54}.
 \tag{SC27.7}
\] The unchanged full real750 cover, with the actual imaginary radius H/(2U), gives \[
 \#d\le \frac{V_F(C_e(q))U^{-f}T^{54}x^{-1}}{N_FJ_0},\quad
 C_e(q)=H/2+750R T^{-2e/f}/q,
 \quad V_F(C)=\frac{((22/7)fC^2)^{f/2}}{(f/2)!\sqrt{D_F}}.
 \tag{SC27.8}
\] Adding a row ideal does not change this central difference ideal. The real-part shift and all dyadic conditions are retained; no half-lattice is substituted.

\subsection{3. Complete finite weights and subset multiplicities}\label{complete-finite-weights-and-subset-multiplicities}

Once eta is fixed, the actual one-label proof DO105 EL9--16 gives the entire active label-weight sum at most(T−1)\^{}-27. It includes central-only poles and omitted cancellation labels. Below2\^{}12000, eta is unique per active support by SC27.4; no extra eta sum is omitted.

At a noneligible inactive place which is not a paired-C-zero place, the exact actual coefficient has magnitude1/(T−1). There are at least \[
 g=168-r-b
 \quad\text{such places, and}\quad
 (T-1)^{-(27+g)}<2T^{-(27+g)}.
 \tag{SC27.9}
\] Other inactive coefficients retain their contraction bound1. There is no new inactive-label count. The single correction2 follows from2(T−1)\^{}168\textgreater T\^{}168.

For clarity, b\textless=14 throughout this bounded window: a nonzero S value gives x T\^{}(2b−54)\textless=(R/U)\^{}168, while x\textgreater2\^{}2111 and the exact inequality(R/U)\^{}168 T\^{}24\textless2\^{}2111 exclude b\textgreater=15. Thus g\textgreater=118 and the exponent27+g is between145 and168. The proof nevertheless bounds all nonempty inactive subsets by2\^{}141, including subsets which cannot occur; this is a safe full containing sum.

The number of active supports is below2\^{}104 and the retained original first-pivot multiplicity is below2\^{}8. For b\textgreater0 sum every subset of the141 inactive positions, costing less than2\^{}141. For b=0,j\textgreater0 choose one of the27 active positions, costing less than2\^{}5; multiple eligible choices are allowed to overcount. The two mixed contractions remain bounded by one. No class, determinant or independently varying phase multiplicity is added.

Combining the complete T powers of Green, central and actual finite factors gives, before any directed logarithms, \[
 335(27-e)+54-(27+g)=
 \begin{cases}8904+r-334b,&e=b,\\8904+r-335,&b=0,e=1,\\8904+r,&b=e=0.
 \end{cases}
 \tag{SC27.10}
\] In particular the gain from b is not obtained by retaining one extra ideal while discarding the other b−1 ideals or their coefficient costs.

\subsection{4. The full positive-mixture row majorant at every coordinate size}\label{the-full-positive-mixture-row-majorant-at-every-coordinate-size}

Use the genuine selected trace K108 AE13, with the complete K107 lower-tail absorption and actual near-point transport. Before any positive Gaussian enlargement, \[
 |c(h)|<2\vartheta_a^f2^{-933f/2}N_F(\alpha)^{-1/2}
             \prod_\tau\Psi(\alpha_\tau),\quad
 \Psi(t)=\left(\frac2{\sqrt{1+t}+1}\right)^{333/2},
 \quad 0<\alpha_\tau\le R<132,
\] \[
 N_F\alpha\ge U^f(kz^2)^f,\qquad\vartheta_a^f<2^{-649}.
 \tag{SC27.11}
\] The last field-norm inequality uses the k original nonzero S coordinates. These restrictions are not asserted for the zero or out-of-ball tuples introduced by the ensuing positive enlargement.

K110 PM3--10 supplies the positive measure mu supported on b\textgreater=2, including its explicit low-parameter atom2\^{}-39 at2. Its moments satisfy \[
 M_3\le m:=\frac{83316594330677601}{5384631367784149811200},
 \qquad M_0,M_1,M_2<2.
 \tag{SC27.12}
\] This is the whole majorizing measure; no additive analytic tail remains outside the count.

Use the canonical K111 degree15 scalar coefficient \[
 c_{336}=168{336\choose168}/4^{168},\quad
 M_{15}=\frac{6015684926187167549292201834533348561577984}
 {2618017117694612446410293519},
\] \[
 K=\frac{336M_{15}}{c_{336}}
       (7841/5000)^{168}168!
       \left(\frac{(365729/500)^2}{168(314159/100000)^2}\right)^{168}.
 \tag{SC27.13}
\] The theorem is valid for every positive Gaussian parameter and every real336-dimensional lattice with the required minimum. Hence it applies after the product-one/positive diagonal embedding changes in K110 PM10, without asserting that the changed real lattice is an ideal.

The actual saturated transverse basis has det G'=U and G'\textgreater2\^{}36 I. Its membership is obtained for each original integral input, not by dividing an old lattice value by2. Its primal covolume and dual metric are exactly those in K107/K109. Put \[
 S(q)=(q/(2R))^f+K,\qquad
 \rho(q)=K\left(\frac{R}{2^{36}q}\right)^f,\qquad
 A(q)=m^f+2^f\{3\rho(q)+3\rho(q)^2+\rho(q)^3\}.
 \tag{SC27.14}
\] The complete weighted row sum is at most \[
 {164\choose k}\frac{U^f\pi^{(k+3)f}}{D_0^{k+3}}
               S(q)^k(q/R)^{3f}A(q).
 \tag{SC27.15}
\] Indeed the uniform scalar S bound is K110 PM11, and its transverse formula PM12 integrates to exactly the sum of four nonnegative terms binom(3,h)rho\^{}h M\_(3−h)\^{}f.~SC27.12 bounds it by A(q). This proof does not impose the older dense-only condition Uz²\textless1. In particular none of the last three terms is dropped when q becomes small. All compact factors, the genuine Hilbert metric and the full positive tail are retained in SC27.11 before summation.

\subsection{5. Coupled endpoint monotonicity and the entire finite pattern sum}\label{coupled-endpoint-monotonicity-and-the-entire-finite-pattern-sum}

Expand A(q) into terms indexed h=0,1,2,3. Apart from their positive constant coefficients, the complete q dependence before the central correlation is \[
 F_{k,h}(q)=q^{-(164+h)f}S(q)^k V_F(C_e(q)).
 \tag{SC27.16}
\] Let w=(q/(2R))\^{}f/S(q), and epsilon=(C\_e(q)−H/2)/C\_e(q). Both lie strictly between0 and1. The normalized logarithmic q derivative and the full coupled x derivative, at fixed J, are \[
 f^{-1}\frac{d\log F_{k,h}}{d\log q}
       =-164-h+kw-\epsilon<0,
\] \[
 \frac{d\log\{x^{-337/2}F_{k,h}(q(x))\}}{d\log x}
       =-9/2+h-kw+\epsilon<-1/2.
 \tag{SC27.17}
\] Here q(x) is proportional to x\^{}-1/f.~Thus every one of the four correction terms decreases with x, at every1\textless=k\textless=164. It is legitimate to evaluate the containing positive majorant at a lower endpoint even if that endpoint itself would violate the original q\textgreater=k feasibility condition. The latter condition was used only on actual rows in SC27.11.

For fixed x and e define the all-rank-zero endpoint and pattern ratio by \[
 q_*^f=\frac{R^fT^{54-2e}}{U^fQ^{334}x},\qquad
 D=Q^{334}/N_EJ,\quad q=q_*D^{1/f},
 \quad (N_FJ_0)^{-1}\le Q^{-167}D.
 \tag{SC27.18}
\] After expanding S\^{}k, the scalar degree l in a transverse term h has pattern exponent l−163−h. The full containing product is exactly \[
 \prod_{\mathfrak q\mid G}\left(1+q_{\mathfrak q}^{l-163-h}\right)^2.
 \tag{SC27.19}
\] It includes the multiplicities1,2,1 of every original local state. This is the same deliberate central-correlated majorant used in R094; the original pattern polynomial is no larger.

For l\textless=162+h, the product is below2, by(1+g0\^{}-1)\^{}668\textless2. The only possible exceptional degrees have exponent0 or1, with products at most2\^{}668 or2Q\^{}334. At every endpoint used below with k\textgreater162, the exact certificate verifies \[
             (q_* /(2R))^f/K<2^{-200}.
 \tag{SC27.20}
\] The normalized positive binomial masses then make the sum of both exceptional tails below1, even after their complete products: use164·2\^{}(668−200·163)+2\^{}(1+71·334−200·164)\textless1. Therefore the entire pattern sum for each of the four correction terms is below3\textless4. For the r27 sparse group k\textless=99, every degree is nonexceptional, so the same4 bound is valid without SC27.20. This is a full pattern sum, not a maximum over one preferred G pattern.

For the b\textgreater0 group, first carry out SC27.17 and the pattern sum at the lower x endpoint. The central radius at the all-rank-zero endpoint is \[
 C_e(q_*)=H/2+750R/q_*(e=0),
 \tag{SC27.21}
\] so it is independent of b. The logarithmic derivative with respect to b of a transverse term, divided by log T, is−6+2h−2kw\textless0. Thus every nonempty subset's endpoint is bounded by the e=b=1 endpoint, after which all its possible positions are summed by2\^{}141. For one active integral-C place the corresponding derivative is−7+2h−2kw\textless0; it suffices to keep one such place and all27 choices. No false uniform derivative in r is used: every r=27,\ldots,36 endpoint is evaluated separately.

Finally the complete endpoint's k-dependent majorant, including the original-coordinate field-norm factor, is proportional to \[
                {164\choose k}\,\mathcal A^k k^{-84},
 \qquad \mathcal A=(22/7)^f S(q_*)/D_0>2^{92}.
 \tag{SC27.22}
\] Its adjacent ratio is greater than \(\mathcal A/(164\cdot2^{84})>1\). Thus all k up to K0 are bounded by the endpoint K0, and their full union costs K0 (bounded by2\^{}8 for164/163 or2\^{}7 for99). This accounts for every coordinate size and every subset of S; no unproved endpoint sampling replaces the continuum or discrete argument.

\subsection{6. A separate complete infinite denominator tail}\label{a-separate-complete-infinite-denominator-tail}

For x\textgreater=2\^{}12000, eta is enumerated first, with no bounded-uniqueness assertion. DO104 CD5--9 applies at arbitrary support: for each exact denominator ideal there are fewer than2\^{}39 norm-U slopes, including the full numerator-ideal choices and all28 norm-one roots. The original reconstruction from eta,row,d is injective.

Use only the older support-independent actual trace bound2\^{}-929f/2 U\^{}-f/2 z\^{}-f, full original transverse count4 A3(4U)\^{}f M\^{}(6f)z\^{}-6f with A3\textless2\^{}-4040, and the complete spherical/Gaussian S count DO104 CD20--22. All active type-II labels are summed generically with total weight at most1 per place. Retain support2\^{}168, all patterns2\^{}668, pivot2\^{}8 and Green2, giving the original845-bit factor. No K112 one-label saving or new Gaussian scalar improvement is needed for this tail.

The original actual norm cap z\textless=M/sqrt(k) gives the complete dependence x\^{}-337/2 z\^{}((328−2k)f), with nonnegative z exponent. It can be bounded by its norm endpoint at every x\textgreater=X. Summing over all integral E ideals and using zeta\_E(2)\textless2\^{}336 yields the factor X²2\^{}336. This is a convergent whole ideal sum, without ideal generators or a class-number restriction. Set X=2\^{}12000.

The central radius is at most67+98300/k\textless=98367. The complete k-dependent S factor is at most2\^{}111356, since binom(164,k)\textless=2\^{}164,47\^{}k\textless2\^{}984, f!\textless=f\^{}f,87/32\textless4 and both negative k powers are at most1. Sum all164 coordinate cells by2\^{}8. Therefore the entire tail is bounded by2 raised to \[
\begin{split}
 E_{\rm tail}={}&845+2+39+2f-4040-929f/2
       +167f(113/16)-(335f/2)L_U\\
 &+337\cdot27L_T-335\cdot167L_Q+336-(333/2)12000\\
 &+\lceil\log_2 V_F(98367)\rceil_{\rm strict}+111356+8
 =-3162562095/4096.
\end{split}
 \tag{SC27.23}
\] All positive R powers were grouped before the safe upper bound log2 R\textless113/16. This contains the entire unbounded denominator tail. No fixed-window multiplicity or support26 tail charge is borrowed.

\subsection{7. Complete exact endpoints and one union}\label{complete-exact-endpoints-and-one-union}

Use the directed original logarithms \[
 L_T=1019333/1024,\quad L_Q=287843/4096,\quad
 L_U=601356/4096,\quad L_R=28810/4096,
 \tag{SC27.24}
\] with upper signs for T,R and lower signs for Q,U. All are verified by integer powers. In a bounded endpoint x=2\^{}a, the certificate encloses q\_* between positive rationals q\_- and q\_+ by exact168th-power comparisons. Its central radius uses the e=0 lower bracket as in SC27.21. Put v\_a equal to the strict integral binary ceiling of that full central volume.

The common exponent is \[
 C_{r,a,b_*,j_*}=117-933f/2-649+v_a+167fL_R-(335f/2)L_U
 +(8904+r-334b_*-335j_*)L_T-336\cdot167L_Q-(337/2)a.
 \tag{SC27.25}
\] Here (b\_\emph{,j\_})=(1,0) for group B after its all-b monotonicity, (0,1) for group C after choosing one active place, and (0,0) otherwise. In particular j\_* is a chosen-place indicator, not the original number j of active integral-C places. The117 contains support104, pivot8, Green1, complete T correction1, pattern4 costing2, and the single selected near-point factor1. The saturated determinant U, not4U, is already used in SC27.15; there is no additional saturation subtraction. The actual radial factor649 is charged once. The exact rational row expression is \[
 W_{k,a,e}=\frac{{164\choose k}(22/7)^{(k+3)f}
 S(q_+)^k(q_+/R)^{3f}A(q_-)}
 {D_0^{k+3}q_-^{167f}k^{84}}.
 \tag{SC27.26}
\] Every numerator and denominator is integral after clearing explicit rational factors. The certificate computes the least strict binary ceiling of W and adds the complete padding stated below. No floating logarithm enters any certified endpoint.

\begin{longtable}[]{@{}
  >{\raggedright\arraybackslash}p{(\columnwidth - 8\tabcolsep) * \real{0.1667}}
  >{\raggedleft\arraybackslash}p{(\columnwidth - 8\tabcolsep) * \real{0.2222}}
  >{\raggedleft\arraybackslash}p{(\columnwidth - 8\tabcolsep) * \real{0.2222}}
  >{\raggedright\arraybackslash}p{(\columnwidth - 8\tabcolsep) * \real{0.1667}}
  >{\raggedleft\arraybackslash}p{(\columnwidth - 8\tabcolsep) * \real{0.2222}}@{}}
\toprule\noalign{}
\begin{minipage}[b]{\linewidth}\raggedright
Containing group
\end{minipage} & \begin{minipage}[b]{\linewidth}\raggedleft
r endpoints
\end{minipage} & \begin{minipage}[b]{\linewidth}\raggedleft
largest k
\end{minipage} & \begin{minipage}[b]{\linewidth}\raggedright
additional complete padding
\end{minipage} & \begin{minipage}[b]{\linewidth}\raggedleft
maximum full exponent
\end{minipage} \\
\midrule\noalign{}
\endhead
\bottomrule\noalign{}
\endlastfoot
A: x\textgreater=2\^{}12000 & unbounded & 164 & already in SC27.23 & −3162562095/4096 \\
B: b\textgreater0, x\textless2\^{}12000 & 27..36 & 164 & 2\^{}141 all subsets and2\^{}8 all k & −1695235/1024 \\
C: b=0,j\textgreater0, x\textless2\^{}12000 & 27..36 & 164 & 2\^{}5 active choices and2\^{}8 all k & −356729/128 \\
D: b=j=0,r\textgreater=29, x\textless2\^{}12000 & 29..36 & 164 & 2\^{}8 all k & −2855283/1024 \\
E: b=j=0,r27 or28, x\textgreater=2\^{}3328 below the tail & 27,28 & 164 & 2\^{}8 all k & −39143/128 \\
F: b=j=0,r27,k\textless=99 below the tail & 27 & 99 & 2\^{}7 all k & −809725/1024 \\
G: b=j=0,r28,k\textless=163 below the tail & 28 & 163 & 2\^{}8 all k & −70695/128 \\
\end{longtable}

For B,C,D,F,G use a=a\_r from section2, so their containing sets may extend beyond an earlier removal; the union is deliberately conservative. E uses a=3328 at both r values. There are exactly1+10+10+8+2+1+1=33 full endpoints. They include all overlaps; within one endpoint the two-dimensional ideal/denominator continuum was bounded before the count.

The worst endpoint is E with r28: q\_- =6785380593/4096, q\_+=3392690297/2048, v\_a=474, C=47976985/128, and the rational ceiling of W is−375134. Including the8-bit coordinate union gives−39143/128. The exact sum of all33 dyadic upper bounds is strictly less than \[
                              2^{-304}.
 \tag{SC27.27}
\] Removing their union leaves exactly the necessary alternatives SC27.2, within the stated original all-II sector. Their possible existence, sign or nonvanishing is not asserted.

\subsection{8. Verification, chronological prefix and remaining source gates}\label{verification-chronological-prefix-and-remaining-source-gates}

The executable checks every fresh eligibility floor and uniqueness norm comparison, the whole active/inactive correction, all position and coordinate sums, the all-parameter scalar coefficient, every root enclosure, each complete endpoint, the infinite ideal tail and the full33-term union. Exact finite rational derivative and pattern checks accompany, but do not replace, the universal positive-integral and lattice arguments above. The entire preliminary replay is preserved separately.

Its chronological baseline is R094120/31 with complete floors27/11, and \nolinkurl{--prior} accepts a later complete canonical ordered extension. The final author replay uses actual DO118122/31, retaining DO11729365 followed by DO1183927. The whole all-top prefix and all31 uniform entries are retained;304 is appended exactly once, giving123/31. Both full dyadic error totals remain below2\^{}-230. The current all-top floor stays27, and the uniform floor stays11. Chronology updates are separate replays of unchanged subjects.

Direct inputs are canonical K112 UE1--14; K103 SC26.5--20 only at its support-independent local/module scope; DO098 CF4--13; DO104 CD5--26; DO107 IP4--10 and DO108's local active-integral statement; K107's actual saturated membership, metric and complete trace; K108 AE13--21; K110 PM3--14; K111 HP1--18; and R094 PP5--18's exact positive polynomial pattern argument. Every application preserves the original all168 selected places and both CM components. No freely varying source coefficient or replacement selected module is used.

Primary context remains Nakahama(2016), equation5.7/Theorem5.1, for the genuine vector-valued Hilbert normalization; NIST DLMF10.32.10/10.39.2 for the positive-real subordination identity, proved in full in K110; and Delsarte, Goethals, and Seidel(1977), Theorems3.3/3.6/4.3, for the harmonic-code context fully normalized in K111. The half-trace Fourier convention and ideal transference are the checked DO070/K105 interfaces. This theorem introduces no new primary comparison claim.

Carried exact geometric anchors are SGA3 I4.7.3/4.7.3.1; SGA6 VI Theorem3.3; FGA232 Theorem3.1 pp148--149; FGA236 Corollaries1.7/2.3 pp229/232; SGA7 I VII2.7/2.9 pp169--174 and VIII7.4(b) pp310--311; cote4 scan003 sections8.1.1--8.1.2, scans005--006 sections8.4.1--8.4.3; and cote54 scans005--007. These are inherited context. The full signed sufficient inequality, remaining structural sectors, actual selected arithmetic content, prescribed period, effective source/field costs, Read/specific032, DP44 and RH remain separate.

\subsection{Canonical algebraic-geometric context}\label{canonical-algebraic-geometric-context-54}

All local row and central difference ideals are the original fractional projective modules over Spec(O\_E)/Spec(O\_F), with the same CM involution, polarization, Picard/Poincare line, finite correspondences and selected real metric. Positive mixture, scalar packing and finite-pattern polynomials estimate these original objects without replacing their arithmetic source. Primary analytic context is carried through checked K107/K108/K110: Nakahama(2016),equation5.7 and Theorem5.1 pp960-961; NIST DLMF10.32.10/10.39.2. Scalar context is Delsarte,Goethals,andSeidel(1977),Theorems3.3/3.6/4.3 through fully checked K111 HP1-18. Carried geometric anchors: SGA3 I4.7.3/4.7.3.1; SGA6 VI Theorem3.3; FGA232 Theorem3.1 pp148-149; FGA236 Corollaries1.7/2.3 pp229/232; SGA7 I VII2.7/2.9 pp169-174 and VIII7.4(b) pp310-311; cote4 scan003 sections8.1.1-8.1.2/scans005-006 sections8.4.1-8.4.3 and cote54 scans005-007. No new formal slope, lambda/K-theory relation, l-adic weight or CM-motive comparison is inferred. DP\_MULT Q3e/DP44 and THETA3 G1/A7 receive only the stated smaller conditional finite gate; CX\_GRH CX-ALG-024 and GCT\_PLETH T-052 remain globally open at their established gates.

\clearpage

\section{AM-R-095 --- the complete rank-one fixed-block sector at support27}\label{am-r-095-the-complete-rank-one-fixed-block-sector-at-support27}

\textbf{Canonical checked integration --- 2026-11-23T08:11:53.450261+00:00.} Complete root and independent proof/script/interface reviews and separate exact replays passed. The actual canonical predecessor has123/31 entries; this integration gives124/31, complete floors27/11. Frozen author history is preserved. AM-R-095 FW1-21 removes the entire actual conditional J56 fixed-S globally rank-one-unipotent support27 sector, all original profiles, labelled pairs, double-pole sizes and unmixed G patterns. A uniform weighted eta/content cutoff precedes the complete334-place binary pattern product; the remaining z exponent is165f, both genuine affine coordinates and the whole support/marked/type sum remain. The reduced exponent is-54242909/4096 and the complete witness exponent is-95126187/256; their union is below2\^{}-13000. Append13000 once to actual K113123/31, giving124/31 with both full sums\textless2\^{}-230; complete floors27/11 remain. No whole support27 exhaustion, signed inequality, actual admitted Host/Sizes, selected arithmetic content, prescribed period, specific032 pairing, evaluatedDP44 or RH conclusion.

Author: GPT-6 Astra Pro; inherited effort not externally exposed; ALP\_MACHINE \nolinkurl{/root}. 23 November 2026 UTC. New composition proof. The saved whole-proof reading, exact execution and independent review are separate dated events; this file alone does not assert canonical promotion.

\subsection{1. Actual sector and claim}\label{actual-sector-and-claim}

Keep the conditional J56 all-top kernel, f=a=168, H=133, the same fields E/F, original integral Hermitian modules, both CM partners, actual source and target pencils and labels, finite projectors, selected real representation and Hilbert metric, packet point, central phases, class masses and probability Haar measures. The actual sector considered here has

\[
h=I_S\oplus(I+N),\qquad \operatorname{rank}_E N=1,
\qquad N^2=0,\qquad s=27.                         \tag{FW1}
\]

Global rank and unipotence are premises. They are not inferred from a finite reduction. The entire sector means every original support, type profile, labelled incidence and unmixed G pattern satisfying FW1.

\textbf{Theorem.} Its full absolute normalized contribution is less than2\^{}-13000. This is one fresh charge on its intersection with the current remainder. It does not establish a complete support28 floor or the sufficient signed inequality. The argument freshly evaluates the complete support27 count; it does not extend an earlier numerical support range by assertion.

The direct accepted inputs are R088 GZ2--5 and its support-independent source explanation in section2, and R089 WN2--12. R088 supplies the original witness/reduced dichotomy, all coordinate and incidence factors and the full binary pattern structure. R089 WN3 is the already proved reciprocal-norm weighted fractional-ideal lemma, whose statement allows every nonzero real fractional ideal. Its proof is cited, not repeated. Its application here is new: the dyadic-profile cutoff is made uniform over \textbf{all} G patterns and all intrinsic double-pole sizes before a new monotone pattern sum.

\subsection{2. Original witness and reduced alternatives}\label{original-witness-and-reduced-alternatives}

Set

\[
Q=22285932761\cdot64072175441=1427908193729121522601,
\quad M=Q^{167},\quad R=8844/133,
\] \[
L_T=1019333/1024,\qquad L_Q=287843/4096,
\qquad \log_2T<L_T,\quad \log_2Q>L_Q.              \tag{FW2}
\]

The same literal T is read from the J56 host certificate, with2\^{}995\textless T\textless2\^{}996. The support-free witness expression from R088 GZ3 gives, freshly at27,

\[
E_{\rm wit}=1+168+107+5757+46af
 +(10L_T+997)27-a167L_Q
 =-95126187/256<-371000.                            \tag{FW3}
\]

This includes the whole original five-column/source-isometry count, all witnesses, support and type assignments, zero targets and finite factors. No later pattern multiplicity is added to FW3.

In the other alternative the same original column determines N and has

\[
N=x\bar x^{\mathsf t}J_0/\eta,\quad x=(vq,r,\eta),
\quad v=(2,95192390527,383092671043),
\] \[
Vq\bar q+r\bar r=\eta\bar\eta,\quad q\ne0,
\quad \bar\eta=-\eta\ne0,\quad |\eta_\tau|\le R,
\quad V=155821585821105090045582>2^{77}.             \tag{FW4}
\]

Let0\textless=d\textless=27 be the intrinsic number of original double poles of eta; every such place is typeII. Let J be the paired rank-zero ideal at the334 unmixed relative G places. The integral rank-one reduction has rank0 or1, giving exactly one binary choice at each such place. The inverse-adjoint CM partner is paired data and is not another independent choice. For fixed support, marked set, original labelled pairs and J, write

\[
\mathfrak l=J I_\Sigma I_D=\Lambda_F\mathcal O_E,
\qquad z^f=N_F\Lambda_F=(N_EJ)^{1/2}T^{-(27+d)},
\] \[
Y=R^f/z^f,\qquad y_c=R/(zT^{2d/f}).               \tag{FW5}
\]

These are the actual paired entry and content ideals of R088 GZ4, not enlarged denominator guesses. The two original affine difference cosets satisfy

\[
C_d(z)=(1+2\sqrt{y_c})^{2f}
          (1+2\sqrt{y_c/V})^{2f}                  \tag{FW6}
\]

as a uniform upper bound for their joint count at fixed eta, actual content and original labels. The primitive-coordinate extraction and true metric are unchanged. A complete column determines N, so no isometry fibre, additional relative-norm fibre, phase multiplicity or norm-generator count is appended to FW6.

\subsection{3. Weighted eta and content at every pattern}\label{weighted-eta-and-content-at-every-pattern}

Put t=-i eta in the actual real fractional ideal Lambda\_F; the original integral unit i identifies these elements bijectively. The actual content valuation count of R089 WN7--9 gives

\[
c(t)\le d_{3,F}((t)\Lambda_F^{-1}),\qquad
|N_E\eta|^{1/2}=|N_Ft|.                            \tag{FW7}
\]

The genuine selected normalized trace, before replacing its eta norm by an ideal minimum, satisfies R089 WN10:

\[
|c(h)|<2\,2^{-931f/2-630}|N_Ft|^{-1}.              \tag{FW8}
\]

The630-bit integrated angular factor, both actual compact factors, original packet transport and selected Hilbert normalization remain. This estimate is uniform in the stated sector and every pattern; its derivation uses no assumption J=J\_max. It is a selected trace estimate, not a matrix-operator norm substitute.

If Y\textless1 there is no nonzero t. Otherwise R089 WN3, applied with delta=1/1000, and its exact actual-field constant WN12 imply

\[
\sum_{0\ne t\in\Lambda_F,\ |\tau t|\le R}
       \frac{c(t)}{|N_Ft|}
 <2^{1233}z^{-f}
    {\lfloor\log_2Y\rfloor+f\choose f}Y^{1/1000}. \tag{FW9}
\]

All real signs, units, ideal classes and content ideals occur within this sum. The two affine counts are uniform in each actual content and shift, so FW6 can be pulled outside it. Neither an unweighted eta count nor an additional z\^{}-f selected-trace bound is multiplied into FW9.

Since N\_EJ\textgreater=1 and27+d\textless=54, the literal host integers give the uniform bounds

\[
Y\le R^fT^{54}<2^{55000},\qquad
{55168\choose168}<2^{1650},\qquad
Y^{1/1000}<2^{55}.                                 \tag{FW10}
\]

Consequently FW9 is at most2\^{}(1233+1650+55)z\^{}-f, for every d and every G pattern. The zero-contribution Y\textless1 cases also obey this positive upper bound. FW10 counts all dyadic embedding profiles; no regulator or lattice-shape estimate is assumed.

\subsection{4. Complete weighted pattern sum}\label{complete-weighted-pattern-sum}

The original finite Green bound is2(N\_EJ)\^{}(a/2)M\^{}-a. Fix d and the original labels. After FW10, the remaining z-dependent part of the product of this Green bound, the weighted sum and both affine-coordinate bounds is, up to a positive constant independent of J,

\[
F_d(z)=z^{(a-3)f}(\sqrt z+A_d)^{2f}
                  (\sqrt z+A_d/\sqrt V)^{2f},
\qquad A_d=2\sqrt R\,T^{-d/f}.                    \tag{FW11}
\]

Indeed the Green numerator contributes z\^{}(af), the weighted eta/content sum contributes z\^{}-f, and the two rewritings of FW6 contribute z\^{}-2f. The profile and Rankin factors have already been replaced by uniform constants. Thus no remaining J dependence has been omitted. Since a-3=165\textgreater0, for0\textless z'\textless=z,

\[
F_d(z')/F_d(z)\le(z'/z)^{165f}.                    \tag{FW12}
\]

At the maximal rank-zero pattern, N\_EJ\_max=M\^{}2 and

\[
z_d=(M/T^{27+d})^{1/f}.                             \tag{FW13}
\]

Removing one relative rank-zero prime of norm q\_v decreases z by q\_v\^{}(1/f). Summing every binary missing subset therefore gives the containing bound

\[
\sum_JF_d(z_J)
\le F_d(z_d)\prod_{v=1}^{334}(1+q_v^{-165})
<2F_d(z_d).                                        \tag{FW14}
\]

Every q\_v is g0 or g1, exceeds2\^{}34, and hence q\_v\^{}-165\textless2\^{}-5610. The last inequality follows from the exact integer comparison (2\^{}5610+1)\^{}334\textless2\^{}(5610*334+1). The complete334-choice product includes the full identity pattern. A positive containing value at a cap that is not arithmetically realized is legitimate by FW12. There is no later2\^{}334 multiplicity.

At that cap, y\_c=R(T\^{}(27-d)/M)\^{}(1/f) decreases with d.~A fresh comparison at d=0 gives

\[
R^fT^{27}<2^{97f}M,
\qquad y_c<2^{97},\quad y_c/V<2^{20},
\] \[
C_d(z_d)<(1+2^{50})^{336}(1+2^{11})^{336}
          <2^{20500}\qquad(0\le d\le27).          \tag{FW15}
\]

Both coordinates, including their actual V metric, remain in this bound. In particular no small-coordinate premise from support26 is imported.

\subsection{5. Entire support, marked-set and original-pair union}\label{entire-support-marked-set-and-original-pair-union}

At J\_max the finite numerator divided by M\^{}a equals1, and its Euler factor2 remains. The full original support/marked/type/label sum is R088 GZ5's source interface, K077 AP18:

\[
{168\choose27}{27\choose d}(T+1)^{27-d}.           \tag{FW16}
\]

It is applied after the uniform bound for each original labelled pair. Every type and both paired components remain in it. At the cap,

\[
(T+1)^{27-d}z_d^{-f}
 =T^{54}M^{-1}(1+T^{-1})^{27-d}
 <2T^{54}M^{-1},                                   \tag{FW17}
\]

because the actual integer T obeys(T+1)\^{}27\textless2T\^{}27. Also

\[
{168\choose27}<2^{104},\qquad
\sum_{d=0}^{27}{27\choose d}=2^{27}.               \tag{FW18}
\]

Combining FW8--18 bounds the \textbf{entire reduced sector}, with all patterns and all d, by

\[
2^{4+1233+1650+55+20500+104+27-931f/2-630}
          T^{54}Q^{-167}.                          \tag{FW19}
\]

The exponent4 consists of exactly the Green Euler factor2, the actual near-point factor2, the full G-pattern factor2 and the incidence-ratio factor2. The20500 bound already counts both affine coordinates. The weighted eta/content constant1233, all embedding profiles1650, Rankin factor55, support104 and full marked-set union27 occur once each. All other original mixed, central phase, class and normalized Haar factors have their inherited modulus bounds. No cancellation of different arrows is used.

The full logarithm in FW19 is strictly less than

\[
4+1233+1650+55+20500+104+27-931\cdot84-630
 +54\frac{1019333}{1024}-167\frac{287843}{4096}
 =-\frac{54242909}{4096}<-13242.                    \tag{FW20}
\]

Together with FW3, the entire sector is bounded by

\[
2^{-13242}+2^{-371000}<2^{-13000}.                  \tag{FW21}
\]

This proves the theorem, with deliberately unused margin. The two branches exhaust the original witness/reduced dichotomy. All old errors are retained even where their underlying sets overlap the new containing bound.

\subsection{6. Accounting, verification and geometric scope}\label{accounting-verification-and-geometric-scope}

The unchanged companion checks the literal host and radius, actual weighted-content/radial constants, all28 double-pole sizes, every uniform profile/coordinate/ratio comparison, the full334-factor pattern inequality, original type polynomial, exact FW3/FW20/FW21 and the complete chronological ordered arrays. The default predecessor is DO118122/31 with complete floors27/11; unchanged replay on a later canonical prefix is permitted and recorded separately. It appends13000 once and changes neither complete floor. Exact diagnostics supplement the universal composition proof; they do not enumerate globally admissible arrows.

Direct accepted sources are R088 GZ2--5 and its section2 explanation of the support-free original dichotomy, GZ6--7 for the paired binary pattern mechanism; R089 WN2--12 for the universal weighted lemma, actual content transport and norm-sensitive selected trace. Their underlying precise interfaces remain DO091 TD2--9; DO092 SN1--6/16--21; K077 AP6--11/18; R080 AQ1--15; LIFT078 AV1--9/16--18; DO088 WM3--8. The earlier endpoint ranges are not reused. No new external theorem is introduced.

The fractional entry, content and coordinate ideals are the same invertible modules over Spec(O\_F) and Spec(O\_E); their ideal norms and paired finite reductions carry the same CM action. The counting rescaling belongs to R089's real embedding argument and changes no integral module, polarization, Picard class, theta line or biextension. Carried exact geometric anchors are SGA3 I4.7.3/4.7.3.1; SGA6 VI Theorem3.3; FGA232 Theorem3.1 pp148--149; FGA236 Corollaries1.7/2.3 pp229/232; SGA7 I VII2.7/2.9 pp169--174 and VIII7.4(b) pp310--311; cote4 scan003 sections8.1.1--8.1.2 and scans005--006 sections8.4.1--8.4.3; cote54 scans005--007. These are inherited placements, not new lambda-ring, K-theory, l-adic-weight, Dieudonne-crystal or CM-motive comparisons.

DP\_MULT Q3e/DP44/DP52/DP54--65 and THETA3 G1/A7 receive precisely the smaller conditional finite remainder proved here. The actual admitted mixed source and full Host/Sizes, selected arithmetic content, prescribed global period, specific032 pairing and value, actual DP44 cost, CX\_GRH CX-ALG-024 RH gate and global GCT\_PLETH T-052/Q-ALPOGE comparison remain unresolved. No Lean build or archive-validity audit is performed.

\subsection{Canonical algebraic-geometric context}\label{canonical-algebraic-geometric-context-55}

The entry, content and coordinate ideals remain the same fractional invertible modules over Spec(O\_F)/Spec(O\_E), with the original paired CM action. R089 real-embedding packing bounds the reciprocal norm on these ideals without changing the polarization, Picard class, theta line, Poincare biextension or selected realization. Carried exact anchors: SGA3 I4.7.3/4.7.3.1; SGA6 VI Theorem3.3; FGA232 Theorem3.1 pp148-149; FGA236 Corollaries1.7/2.3 pp229/232; SGA7 I VII2.7/2.9 pp169-174 and VIII7.4(b) pp310-311; cote4 scan003 sections8.1.1-8.1.2/scans005-006 sections8.4.1-8.4.3 and cote54 scans005-007. No new external theorem, Dieudonne slope, lambda/K-theory relation, l-adic weight or CM-motive comparison. DP\_MULT Q3e/DP44/DP52/DP54-65 and THETA3 G1/A7 receive only this smaller conditional finite remainder; CX\_GRH CX-ALG-024 and GCT\_PLETH T-052 retain their global gates.

\clearpage

\section{AM-DO-119 --- every active place is eligible, and the complete non-all-II unipotent sector at support27}\label{am-do-119-every-active-place-is-eligible-and-the-complete-non-all-ii-unipotent-sector-at-support27}

\textbf{Canonical checked integration --- 2026-11-23T08:24:41.985104+00:00.} Complete root and independent proof/script/interface reviews and separate exact replays passed. The actual canonical predecessor has124/31 entries; this integration gives125/31, complete floors27/11. Frozen author history is preserved. AM-DO-119 NA1-29 proves every active place of the actual full-positive globally rank-two-unipotent form is eligible, including distinct source/target labels and every type. Its complete non-all-II scalar-cross support27 bound retains all original paired rows, full-entry poles, affine central cosets, finite weights, relative norm fibres and the entire denominator-ideal tail. Six exhaustive branches give one412 charge. Append412 once to actual R095124/31, yielding125/31 with both complete sums below2\^{}-230 and floors27/11 unchanged. The dense all-II27 and separate nonzero-short-compression sector remain open; no full support28, actual admitted source/Host/Sizes, selected arithmetic content, prescribed period, specific032 pairing, evaluatedDP44 or RH conclusion.

Author: \nolinkurl{/root/definite_oscillator}. 23 November2026 UTC. Controlled pending proof. Whole author reading, unchanged executable replay, independent review and canonical integration are separate events. No canonical file is changed here.

\subsection{1. Original sector and conclusion}\label{original-sector-and-conclusion}

Keep the actual J56 all-top normalized kernel, H=133, f=168, original E/F and integral Hermitian module, both CM partners, the original source and target pencils and omitted labels, all334 unmixed G factors and both mixed projectors, selected real module and compact factors, packet point, central phases, class masses and Haar measures. Work in the finite-unipotent G remainder and assume global rank-two unipotence explicitly. In the scalar-cross branch, \[
 h=I+N,\qquad h_{SS}=I,\qquad C=h_{W_0S}\ne0,
 \qquad \operatorname{rank}_E N=2,\qquad s=27.             \tag{NA1}
\] The already charged DO099 JR1--24 interval removes transverse C and inverse C through support70. Retain its plane-C remainder. The support-independent K103 SC26.5--7 / DO110 NT26.5--6 proof gives \[
 E^{169}=Z\perp P,\quad Z=S_{164}\perp V_3,\quad
 P=Eu\perp Ee_n,\quad N^3=0,\quad N^2Z=0,\quad h_{ZZ}=I.
                                                               \tag{NA2}
\] Indeed N² is nonzero, since otherwise the unitary identity would make the rank-two image totally isotropic in signature one. Its unique image line is the image of C and lies in P; self-adjointness of N² gives N²Z=0, and the single length-three Jordan block gives NZ in that line. No support26 numerical cut is used.

Write i,u,t for the active types I,II,III, so i+u+t=27. This theorem concerns \textbf{u≤26}, excluding the all-II branch owned separately by K113.

\textbf{Theorem.} Every active place in NA2, of any of the three types, is eligible for the actual isotropic line. The entire stated non-all-II support27 sector has absolute normalized contribution less than \[
                              2^{-412}.                    \tag{NA3}
\] The eligibility statement is support-independent. The numerical assertion is only the complete NA1 sector at support27. It supplies one fresh charge and does not assert a full support28 floor, a signed lower bound or an arithmetic nonzero period.

Use the unchanged constants \[
 Q=1427908193729121522601,\quad U=77Q^2-14,\quad
 R=17424/133,\quad M^2=R/U,\quad
 D_0=7^{140}29^{144}5^{84}13^{84},\quad D_F=2^{168}D_0.
                                                               \tag{NA4}
\] Put LT=1019333/1024, LQ=287843/4096, LU=601356/4096 and LR=28810/4096. These are upper bounds for log2T,log2R and lower bounds for log2Q,log2U, respectively; log2T\textgreater995. The coarser113/16 upper bound for log2R is retained only in the old CF and infinite-tail formulas in which its complete positive power is unchanged.

\subsection{2. A local lemma for distinct original source and target labels}\label{a-local-lemma-for-distinct-original-source-and-target-labels}

Let O be a DVR with uniformizer pi, fraction field K, and original local decomposition Lambda=Z\_O⊕P\_O with rank(P\_O)=2. Fix primitive z in P\_O and an integral row q with qz=0, q(Z\_O) contained in pi O and q\textbar P primitive. For each finite residue label c choose a lift which is zero or a unit, and put \[
 A_c=I+c\pi^{-1}zq,\qquad A_c^{-1}=I-c\pi^{-1}zq.
                                                               \tag{NA5}
\] Let w in P\_O and b in P\_O dual be primitive with bw=0, extended by zero on Z. Suppose the saturated residue line of w differs from the residue line of z. In a plane this implies qw and bz are units. Thus A\_d\^{}\{-1\}w has minimum coordinate valuation minus1 when d is a unit and0 when d=0; bA\_c restricted to P has the analogous valuation minus1 or0. These assertions use the primitive leading polar coefficient, which an integral term cannot cancel.

For any K-row r vanishing on P, any K-column v in Z\_K and any scalar a in K, put \[
                       N=wr-vb+awb.                         \tag{NA6}
\] The local assertion is \[
 A_d^{-1}(I+N)A_c\in\operatorname{End}_O(\Lambda)
       \quad\Longrightarrow\quad N\in\operatorname{End}_O(\Lambda)
 \qquad(c,d\text{ arbitrary finite labels}).              \tag{NA7}
\] This extends the diagonal K112 UE3--11 implication to the actual separate source and target. It is not an assertion that the two labels coincide.

Here is the complete proof. Write epsilon\_c=1 for a unit label and0 for label0, and similarly epsilon\_d.~For any y in Z\_O, both A\_c\^{}\{-1\}y and A\_d\^{}\{-1\}y are integral, because qy lies in pi O. If K\_cd denotes the integral matrix in NA7, then \[
 K_{cd}A_c^{-1}y-A_d^{-1}y=(A_d^{-1}w)r(y)\in\Lambda.
                                                               \tag{NA8}
\] Hence r(Z\_O) is contained in pi\^{}\{epsilon\_d\}O. The Z-by-P block of K\_cd is exactly −v(bA\_c)\textbar P: the identity correction has no Z row, and neither other term has a Z row. Consequently \[
              v\in\pi^{\epsilon_c}Z_O.                    \tag{NA9}
\] As rz=0, rA\_c=r. Also qv belongs to pi\^{}\{epsilon\_c+1\}O, so A\_d\^{}\{-1\}v belongs to pi\^{}\{epsilon\_c\}Lambda. Both transformed mixed terms \[
                 (A_d^{-1}w)r,
                 \qquad -(A_d^{-1}v)(bA_c)                 \tag{NA10}
\] are therefore integral. The exact remaining equality is \[
 K_{cd}=I+(c-d)\pi^{-1}zq
            +(A_d^{-1}w)r-(A_d^{-1}v)(bA_c)
            +a(A_d^{-1}w)(bA_c).                           \tag{NA11}
\] The identity correction has pole order at most1. The outer product in the last term has exact minimum valuation −epsilon\_c−epsilon\_d: choose independently a minimum coordinate in its column and row. If epsilon\_c+epsilon\_d≥1, integrality of K\_cd forces \[
                    v_\pi(a)\ge\epsilon_c+\epsilon_d-1\ge0.
                                                               \tag{NA12}
\] Otherwise that term would have a pole of order at least2, which cannot cancel the identity correction or the two integral mixed terms. If c=d=0, NA7 already says I+N is integral. Thus N is integral in all cases. In the two-unit case the argument gives a in pi O, but it does not claim the stronger diagonal pi² condition when c and d differ. In the one-zero-label case integrality of a suffices; no additional independence argument or cancellation assumption is required.

The identities bw=0, r\textbar P=0 and v in Z also give N²=−(rv)wb and N³=0. Thus (I+N)\^{}\{-1\}=I−N+N² is integral. This universal DVR argument uses neither a finite positive metric nor an arithmetic realization of the local parameters.

\subsection{3. Every active original place is eligible}\label{every-active-original-place-is-eligible}

The actual full-positive normal form, in the plane coordinates (u,e\_n), is \[
 \ell=(1,\eta),\quad \eta\bar\eta=U,\quad
 N=\begin{pmatrix}0&-a\ell^*J_P\\ \ell r&d\ell\ell^*J_P\end{pmatrix},
 \quad a=J_Z^{-1}r^*,\quad J_P=\operatorname{diag}(U,-1),
 \quad d+\bar d=-rJ_Z^{-1}r^*.
                                                               \tag{NA13}
\] Here alpha=UrJ\_Z\^{}\{-1\}r\^{}*≤R and eta,r,d determine the original matrix injectively. The original S row is nonzero, so its first nonzero pivot exists. At T the fixed Z/P frame is integral and invertible: its determinant factors are only2U, both T-units. This does not assert global integrality of a moving null basis.

Use the actual pencil M\_c=A\_cLambda, with z=v0/bar(pi), q=−v0\^{}dagger−pi m and mz=1. The allowed replacement m=e1\^{}t/z1 from K063 PF12 leaves every actual labelled lattice unchanged by an integral right change. The same literal identity gives qz=0 exactly. Since v0,z lie in P and Z is orthogonal to P, q(Z\_O) is contained in pi O; its P restriction is primitive. All these are the literal K063 PF12 / DO105 EL9--14 / K112 UE12 interfaces.

Saturate the line im(N²) locally, taking primitive w and its primitive annihilator b in P. Absorb their local rescalings in r,v,a; NA13 then has exactly NA6 with r\textbar P=0 and v in Z. No initial integrality of those scalars is assumed. Noneligibility is precisely that the saturated line differs from z: for unit eta the equality is eta≡a0 with a0=Q(−7+2sqrt(−7)); a zero or pole of eta gives one zero coordinate, whereas both actual pole-line coordinates are units.

Every contributing original local factor has an actual incidence hM\_c=M\_d, by K070 PT1--3's C(h)\textgreater0 support. Both admitted labels are finite because infinity is omitted. This applies to full-pencil type I as well as unique-incidence types II and III; it does not replace a type-III arrow by a fixed pencil. Noneligibility would invoke NA7 and make h and h\^{}\{-1\} integral in the selected original lattice. The unchanged inverse-adjoint companion is then integral at the paired CM component, whose original Gram is a T-unit matrix. This contradicts activity. Therefore all active places are eligible.

Let q\_tot denote total eligibility among all168 selected places and x=N\_E D\_eta the exact denominator norm. At support27, \[
 q_{\rm tot}\ge27,\qquad x>2^{2111},\qquad
 x<2^{12000}\Longrightarrow q_{\rm tot}\le36.             \tag{NA14}
\] These follow from the original denominator-cleared norm inequality 7\^{}84 T\^{}\{q\_tot\}\textless{[}4(U+14){]}\^{}168 x and the exact host comparisons in the certificate. If q\_tot≥28, then x\textgreater2\^{}3107. In the bounded window x\textless2\^{}12000, eta is unique per full27-place support: two such scalars have a difference divisible at both CM partners of all27 places, contradicting \[
                           (4U)^{168}2^{24000}<T^{54}.       \tag{NA15}
\] The partner congruence is eta≡U/bar(a0), not an independently selected residue. No eta, denominator-ideal or norm-one-phase multiplicity is counted in this bounded window. Large denominators are counted separately below.

\subsection{4. Original row, central coset and complete weights}\label{original-row-central-coset-and-complete-weights}

The original paired module inclusions with N²Z=0 imply that every one of the167 fixed row values and its eta multiple is integral at types I and III, and has at most a simple pole at type II. Full entries have pole at most1,2,1 at I,II,III. The proof is DO110 NT26.6--7: every fixed input lies in both original modules; type III uses Ny in M\_c intersection M\_d, and type I preserves the common intersection. No pole is removed from a full entry merely because a cross entry is integral.

Let J be the original common-C G ideal, J0 the F-ideal of rank-zero G places. Retain all4\^{}334 patterns and \[
 |a_G(h)|\le2Q^{-335\cdot167}(N_EJ)^{335/2},\quad
 N_EJ\le Q^{334},\quad N_EJ\le Q^{167}N_FJ_0.
                                                               \tag{NA16}
\] For a padded cross-pole set D of v places containing the u type-II places, put L=JI\_D D\_eta, z\^{}\{2f\}=N\_EL. If one chosen inactive paired-C-zero place is additionally retained, write e=1 and multiply L by its paired ideal of normT²; otherwise e=0. Exactly \[
 N_EJ=z^{2f}T^{2v}x^{-1}T^{-2e},\quad
 kz^2\le M^2,\quad z^2\le(Q^{334}T^{-2(v-e)}x)^{1/f}.
                                                               \tag{NA17}
\] The denominator intersection is valuationwise even when eta has poles at G or inactive T. The extra paired-zero factor applies simultaneously to r(y) and eta r(y); it is not inferred from only one CM component.

After eta,r and the actual labelled tuple are fixed, the central coordinate lies in one affine difference coset of iJ0 B\_eta I\_Sigma,F I\_D,F, where D\_eta bar(D\_eta)=B\_eta O\_E and N\_F B\_eta=x. The original4-by4 Bezout extraction is DO104 CD14--19. Its complete real covolume and cover give \[
 \#d\le V_F\!\left(H/2+750Uz^2T^{(v-27-2e)/f}\right)
              U^{-f}T^{27+v}x^{-1},
 \quad V_F(C)=\frac{((22/7)fC^2)^{f/2}}{(f/2)!\sqrt{D_F}}.
                                                               \tag{NA18}
\] Only N\_FJ0≥1 is dropped from the denominator. The covering radius uses the old ideal before the extra paired-zero factor; this explains its T\^{}\{-2e/f\}. The affine real-part shift is retained and no half-lattice factor is inserted.

At every type-II place eta is fixed, so DO105 EL9--16 gives at most one actual diagonal label, including central-only poles and omitted labels. Its complete weight is at most(T−1)\^{}\{-1\}. Type I has complete ordered-pair/weight sum at mostT; type III at most1. After actual matrix bounds are taken uniformly for each tuple, the total active factor is at most T\^{}i(T−1)\^{}\{-u\}.

Let b be the number of inactive paired-C-zero places. Since every active place is eligible, there are exactly168−q\_tot noneligible places and all are inactive. DO107 IP4--10 gives one factor(T−1)\^{}\{-1\} at each which is not paired-zero. Thus any lower bound g≤168−q\_tot−b may be retained. Combine the factors before using \[
                    (T-1)^{-(u+g)}<2T^{-(u+g)},\qquad
                    2(T-1)^{168}>T^{168}.                  \tag{NA19}
\] All other factors are bounded by1. In particular this is a count on the actual retained set; it does not extend a truncated central orbit for free.

The original nonzero S value satisfies xT\^{}\{2b-2u\}≤(R/U)\^{}f.~The exact comparison (R/U)\^{}f T\^{}24\textless2\^{}2111 and NA14 therefore imply b≤12 when u≤25 and b≤13 when u=26. No subset of all inactive places is separately enumerated unless a chosen paired-zero place is used.

\subsection{5. The complete arithmetic envelopes, without numerical transplant}\label{the-complete-arithmetic-envelopes-without-numerical-transplant}

The low-u and infinite-tail branches use the original full CF factors. In the bounded window the same joint-energy row count as DO110 NT26.19--25 is \[
 \#r\le K_kA_3(4U)^f\mathcal B_k
 (M+2^{-10}z)^{2(k+3)f}z^{-2(k+3)f},\quad
 |c(h)|<2^{-929f/2}U^{-f/2}z^{-f},                         \tag{NA20}
\] where \[
 K_k=\binom{164}{k}47^k2^{2fk}(f!)^k(87/(32fk))^{fk},\quad
 A_3=528^{3f}/((3f)!D_0^3)<2^{-4040},\quad
 \mathcal B_k=(fk)!(3f)!/((k+3)f)!.
\] The actual transverse Gram determinant4U, full selected trace, phases and reconstruction remain. The old harmless pivot factor2\^{}8 is kept. The fixed factor2\^{}846 consists of support2\^{}168, patterns2\^{}668, pivot2\^{}8, one Green2 and the single correction2 in NA19. No additional isometry or negative-column fibre is inserted or deleted.

For a lower bound x≥2\^{}a define c(v,a,e)\^{}f=(R/U)\^{}f T\^{}\{2(v-e)\}Q\^{}\{-334\}2\^{}\{-a\}. At its lower endpoint use z≤M/sqrt(max(k,c0)) for any certified c0≤c.~The full joint expression is proportional to \[
 x^{-337/2}z^{(328-2k)f}(M+2^{-10}z)^{2(k+3)f}
 V_F(H/2+750Uz^2T^{(v-27-2e)/f}).                         \tag{NA21}
\] It is increasing in z. Along the global cap, z grows at most as x\^{}\{1/(2f)\}, so the net upper x-power is −337/2+(164−k)+(k+3)+1=−1/2. This proves the same single lower-x endpoint, including the moving central radius and the cap change.

For clarity define the full directed common exponent, with u0 a lower bound on u, i0 an upper bound on i, and d\_num the paid type/set/extra-place binary factor, \[
\begin{split}
 \mathcal C(s,v,u_0,a,g,e,i_0,d_{\rm num};V)={}&846+2f-929f/2+V
 +167fL_R-(335f/2)L_U\\
 &+(336v+s-u_0-g-335e+i_0)L_T
 -335\cdot167L_Q-(337/2)a+d_{\rm num}.
\end{split}                                                \tag{NA22}
\] Its T coefficient is positive in every application before LT is used. V is the strict real-volume ceiling479 for radius68 or679 for radius155. The exact row and inflation ceilings are those of DO110 NT26.24: \[
 \mathcal R_k(A)=\left\lceil\log_2
 \frac{\binom{164}{k}47^k2^{2fk}(f!)^k87^{fk}(fk)!(3f)!}
 {32^{fk}f^{fk}k^{fk}A^{(164-k)f}((k+3)f)!}\right\rceil_{\rm strict},
\] and the strict ceiling of (1+1/(1024 floor(sqrt(A))))\^{}\{2(k+3)f\}. The joint endpoint is C−4040 plus these two ceilings.

The companion freshly recomputes all164 complete old reference envelopes, as arithmetic expressions with every factor, not as assertions that present matrices were previously removed:

\begin{longtable}[]{@{}
  >{\raggedright\arraybackslash}p{(\columnwidth - 4\tabcolsep) * \real{0.3000}}
  >{\raggedright\arraybackslash}p{(\columnwidth - 4\tabcolsep) * \real{0.3000}}
  >{\raggedleft\arraybackslash}p{(\columnwidth - 4\tabcolsep) * \real{0.4000}}@{}}
\toprule\noalign{}
\begin{minipage}[b]{\linewidth}\raggedright
Reference
\end{minipage} & \begin{minipage}[b]{\linewidth}\raggedright
Complete parameters and row alternative
\end{minipage} & \begin{minipage}[b]{\linewidth}\raggedleft
Every k-cell is strictly below
\end{minipage} \\
\midrule\noalign{}
\endhead
\bottomrule\noalign{}
\endlastfoot
C & C(26,25,25,120,126,1,0,13;479),cap18,joint row & −172990 \\
E & C(26,25,25,120,142,0,0,5;479),cap68036,full cone/volume & −1567 \\
F & C(26,25,25,1116,139,0,0,5;679),cap1117,joint row & −58950 \\
\end{longtable}

Reference E uses the full K102 PC6--26 / K103 SC26.27 positive-norm cone. Put t0=68036/k and Pk=max(10000,t0). Its three alternatives are the two cone endpoints t0,Pk and the complex-volume endpoint Pk. For t≤t0 the coupled denominator constraint and full central-radius derivative give an increasing expression; on{[}t0,Pk{]} the logarithm is convex; beyond Pk the complete volume expression decreases. The maximum is one k-cell, with no factor3. The companion recomputes the exact phase constant, real theta4, joint Laplace energy and both endpoint formulas; no beta-pivot or second eta sum is added. These statements depend on the displayed cap, pole powers and original metric, and so apply to the new actual lattice count with the explicitly changed parameters below.

\subsection{6. Six exhaustive branches and one fresh charge}\label{six-exhaustive-branches-and-one-fresh-charge}

\textbf{A. u≤23, every x.} Pad the cross-pole set to23 places. Pay binom(27,23)=17550\textless2\^{}15, all3\^{}27\textless2\^{}54 type choices, and the conservative pair factorT\^{}27. The original full-C/central CF24 formulas have combined T powers334v+s and336v+s. At v23,s27, after the pair factor these save respectively304 and306 T powers against the complete support24 formulas. Retain the same cap30 removal cost2\^{}136028. The three complete coordinate ranges1..64,65..163,164 are therefore below \[
 2^{-170310},\qquad2^{-314421},\qquad2^{-307281},            \tag{NA23}
\] per cell. The full164-cell branch is below2\^{}-170302. This uses the old full relative beta norm fibre and all G orientations; it needs no eta uniqueness.

\textbf{B. u≥24, x≥2\^{}12000.} Enlarge all poles to the full27-place all-II budget. Enumerate eta first, keeping the complete fewer-than2\^{}39 norm-U scalars per exact denominator ideal and zeta\_E(2)\textless2\^{}336 from DO104 CD5--9/23--26. The tail bound is F\_k(x)≤F\_k(X)(x/X)\^{}\{-k-9/2\}, followed by the full ideal sum ≤F\_k(X)X²2\^{}336. Retain generic all-II label sums≤1 before adding all3\^{}27\textless2\^{}54 actual types and T\^{}3, since i≤3 here. Keep the original2\^{}845, transverse inflation4, determinant4\^{}f and A3; no fixed denominator or phase is omitted.

For an explicit coarse certificate, drop the endpoint cap to k, bound every central radius by98367, and bound every original S-coordinate factor by2\^{}111356 using binom(164,k)≤2\^{}164,47\textless2\^{}6,(f!)\^{}k≤f\^{}\{fk\},87/32\textless4 and all denominator factors≥1. The complete exponent, including the164-cell union, is \[
\begin{split}
 E_B={}&845+2+39+2f-4040-929f/2+167f(113/16)
 -(335f/2)L_U\\
 &+337\cdot27L_T-335\cdot167L_Q+336-(333/2)12000\\
 &+\lceil\log_2V_F(98367)\rceil_{\rm strict}+111356+8+54+3L_T
 <-700000 .
\end{split}                                                 \tag{NA24}
\] This is a complete infinite tail, not a numerical window assumption.

All remaining branches have24≤u≤26 and x\textless2\^{}12000. Eta is unique by NA15, and q\_tot≤36. We deliberately weaken x\textgreater2\^{}2111 to x≥2\^{}2110, and x\textgreater2\^{}3107 to x≥2\^{}3106 where convenient. Since T\textgreater2\^{}995, \[
                       T^2/2^{1990}>1.                    \tag{NA25}
\] Hence each fresh cap below is at least its displayed old reference cap. This direction avoids any reversed logarithm comparison.

\textbf{C. u24 or25, every b.} Pad to v25, paying351\textless2\^{}9 sets and all3\^{}27\textless2\^{}54 types, hence d\_num=63. Here i≤3, b≤12 and g≥120; retain at least24 active one-label factors. Use a2110,e0,cap18,V479. The central radius is below68 because T\^{}\{2/f\}\textgreater3690 and H/2+750R/(18·3690)\textless68. Relative to reference C the common exponent changes by \[
       \Delta_C=346L_T-(337/2)1990+50=4688929/512.           \tag{NA26}
\] Both row factors decrease under the cap enlargement. Every cell is below−163831, and the whole branch below2\^{}-163823.

\textbf{D. u26, b\textgreater0.} There are27·2=54\textless2\^{}6 choices of the sole type-I/III position and type, i≤1 and b≤13. Retain g≥119 and choose one inactive paired-zero place in at most141\textless2\^{}8 ways. Use v26,e1,a2110,cap18,V479,d\_num14. The central radius retains the old z scale, with T\^{}\{(26−27−2)/f\}; it is below68. Relative to reference C the common change is \[
       \Delta_D=344L_T-(337/2)1990+1=911127/128.             \tag{NA27}
\] Every cell is below−165871; the164-cell branch is below2\^{}-165863. No sum over all paired-zero subsets is needed, since a chosen one suffices for this containing count.

\textbf{E. u26,b0,q\_tot27.} Every one of the141 inactive places is noneligible, so g=141. Use v26,e0,a2110,d\_num6,i≤1 and the full cone/volume cap68036. The central radius below68 follows even after discarding its extra T\^{}\{-1/f\}. Relative to reference E, \[
       \Delta_E=338L_T-(337/2)1990+1=586509/512.             \tag{NA28}
\] All cone and volume factors, including their full parameter-range proof, are unchanged after this cap enlargement. Every complete cell is below −1567+586509/512=−215795/512\textless−421. The164-cell branch is below2\^{}-413. The exact diagnostic maximum is stronger; no stronger bound is needed here.

\textbf{F. u26,b0,q\_tot≥28.} Now x\textgreater2\^{}3107 and g≥132. Use v26,e0,a3106,d\_num6,i≤1,cap1117 and retain the conservative old central radius155. Relative to reference F the change is again Delta\_D, since3106−1116=1990. Every cell is below−51831, and the whole branch below2\^{}-51823. This one containing branch covers every q\_tot from28 to36, without a q-subset enumeration or a claim of monotonicity in q\_tot.

These six branches are disjoint and exhaustive: first u≤23, then the infinite tail, then u24/25, then the three u26 cases. Their complete union satisfies \[
 2^{-170302}+2^{-700000}+2^{-163823}+2^{-165863}
           +2^{-413}+2^{-51823}<2^{-412}.                  \tag{NA29}
\] Every original support, type, label, padded set, common G orientation, pivot, selected trace, real/transverse metric, affine central fibre and eta/ideal multiplicity is retained once. Older overlapping numerical charges remain in the chronological array.

\subsection{7. Exact verification, chronology and retained scope}\label{exact-verification-chronology-and-retained-scope}

The executable verifies the new general-arrow local implication in exact rational models with nonzero q\textbar Z, both zero and nonzero source/target labels, negative candidate valuations, nilpotence and the inverse-adjoint partner. These models are diagnostics, not global arrows. It recomputes every original CF cell and every complete reference C/E/F cone/volume cell, all new profile/denominator/cap comparisons, all six full branch unions and both entire ordered error arrays.

The actual default predecessor is checks/k113\_joint\_remainder\_error.json,123 all-top/31 uniform, complete floors27/11; it retains the entire DO118122/31 prefix and appends K113's304. The --prior option accepts a later complete extension without reordering or dropping any charge. Append412 once. The unique231 term, distinct233 term and all other preceding charges remain; both full sums are checked below2\^{}-230. No complete support28 floor is asserted here.

Direct interfaces are DO099 JR1--24; K070 PT1--6; K071 NB2--9; DO098 CF4--32; DO103 LD6--15; DO104 CD5--26; DO105 EL9--27; DO107 IP4--15; K102 PC6--26; K103 SC26.5--28; DO110 NT26.5--31; and K112 UE3--14. The new lemma is NA5--12 for distinct labels; K112's diagonal result is not silently applied to type III. All primary inputs are the already checked theorems in those exact original models; no new external representation or lattice theorem is asserted.

The actual fractional ideals remain invertible modules over Spec(O\_E) and Spec(O\_F), and the central coordinate is counted in a torsor under its real difference ideal. Polarization, theta line, Picard and Poincare/biextension identities, selected metric and Haar measures remain unchanged. Carried contextual locators are SGA3 I4.7.3/4.7.3.1; SGA6 VI Theorem3.3; FGA232 Theorem3.1pp148--149; FGA236 Corollaries1.7/2.3pp229/232; SGA7 I VII2.7/2.9pp169--174 and VIII7.4(b)pp310--311; cote4 scan003 sections8.1.1--8.1.2,scans005--006 sections8.4.1--8.4.3, and cote54 scans005--007. These are inherited context, not new primary readings or comparison claims. The all-II27 branch, other support27 structures, signed sufficient inequality, actual source, selected arithmetic content, prescribed period, Read/specific032, DP44 and RH gates remain separate.

\subsection{Canonical algebraic-geometric context}\label{canonical-algebraic-geometric-context-56}

The new lemma concerns the original labelled lattice correspondence over each selected split DVR. Its fractional row and central ideals remain the same invertible modules over Spec(O\_E)/Spec(O\_F); the real central fibre is an affine torsor with the original shift, metric and codifferent normalization. Actual paired CM involution, polarizations, Picard and theta lines, Poincare biextensions, selected realization, classes and Haar measures are unchanged. Exact carried anchors: SGA3 I4.7.3/4.7.3.1; SGA6 VI Theorem3.3; FGA232 Theorem3.1 pp148-149; FGA236 Corollaries1.7/2.3 pp229/232; SGA7 I VII2.7/2.9 pp169-174 and VIII7.4(b) pp310-311; cote4 scan003 sections8.1.1-8.1.2 and scans005-006 sections8.4.1-8.4.3; cote54 scans005-007. No new lambda/K-theory, l-adic weight, Dieudonne or CM-motive theorem is inferred. DP\_MULT Q3e/DP44 and THETA3 G1/A7 gain only this smaller conditional finite remainder; CX\_GRH CX-ALG-024 and GCT\_PLETH T-052 retain their global gates.

\clearpage

\section{AM-K-114 --- full-positive rank three at support27}\label{am-k-114-full-positive-rank-three-at-support27}

\textbf{Canonical checked integration --- 2026-11-23T08:31:21.646287+00:00.} Complete root and independent proof/script/interface reviews and separate exact replays passed. The actual canonical predecessor has125/31 entries; this integration gives126/31, complete floors27/11. Frozen author history is preserved. AM-K-114 FP27.1-29 removes the entire full-positive-compression global-rank-three scalar-cross support27 sector with nonzero original S-cross. All34 complete ideal/type cells,2925 padded sets, original full167-covector metric, relative norm fibre, projected U denominator/T108 column and G2 source fibre are retained. The complete eta-first all-ideal tail requires no coordinate sparsity. Literal inversion and23 common paired eligible positions force equal null lines outside four containing sectors, contradicting rank3. Their full union costs1650 once. Append1650 to actual DO119125/31, yielding126/31, both entire sums below2\^{}-230 and complete floors27/11 unchanged. Full positive compression remains a premise; the separate nonzero-short-compression27 and dense all-II27 sectors remain open. No full support28, actual admitted Host/Sizes, selected content, prescribed period, specific032, evaluatedDP44 or RH conclusion.

Author: \nolinkurl{/root/kummer_extension},23 November2026 UTC. Controlled pending proof. Author, independent and root whole-file checks and exact replays are separate events. No canonical artifact is changed.

\subsection{1. Exact original sector and conclusion}\label{exact-original-sector-and-conclusion}

Retain the actual J56 all-top H133 normalized kernel, f=168, original E/F of degrees336/168, the determinant-one central representative, all334 unmixed finite-unipotent G reductions, both mixed projectors, every actual paired T pencil and admitted label, selected real Hilbert metric, two compact factors, central phases, class masses and Haar measures. At exactly s=27 active relative T places assume \[
 E^{169}=Z\perp P,\quad Z=S_{164}\perp u^\perp,
 \quad P=Eu\perp Ee_n,\quad J_P=\operatorname{diag}(U,-1),
\] \[
 h=\begin{pmatrix}I_Z&B\\ C&D\end{pmatrix},\qquad
 (h^{-1})_{ZZ}=I_Z,\qquad\det h=1,\qquad C|_S\ne0,
 \qquad\operatorname{rank}_E(h-I)=3.
 \tag{FP27.1}
\] The full-positive identity is an explicit hypothesis. DO111's support26 removal of nonzero short compression does not establish it at support27. The latter nonzero-compression sector is separate. This theorem neither removes that sector nor gives a complete support28 floor.

\textbf{Theorem.} The complete absolute normalized contribution of FP27.1 is strictly less than \[
                              2^{-1650}.
 \tag{FP27.2}
\] There is no coordinate-sparsity or denominator hypothesis. There is one fresh combined1650 charge, preserving every previous overlapping charge. The complete support floors stay27/11.

Use the unchanged original constants \[
 H=133,\quad Q=1427908193729121522601,\quad U=77Q^2-14,
 \quad D_0=7^{140}29^{144}5^{84}13^{84},
 \quad M=H/(2\sqrt U),\quad B_0=H/2.
 \tag{FP27.3}
\] Write C=ell r with ell=(1,eta), eta bar(eta)=U. The actual negative-row norm bounds the full167-dimensional covector r by M at each embedding. The primitive original u, the164 original S inputs, and the three integral transverse inputs of R072 PJ4--5 are retained. Their positive Gram is diag(I164,G), det G=4U and G-I\textgreater0. No moving rational null frame is declared integral.

DO074 TC1 is support-independent: because the original S compression is I, a type-III count t\textless=10 would force global unipotence and global rank at most2. Thus every FP27.1 term has \[
                              11\le t\le27.
 \tag{FP27.4}
\] The global-rank-three classification and pointwise minor bound used below are DO095 RA1--20. Neither global unipotence nor an extra determinant phase is assumed.

\subsection{2. Cross-only eligibility and a complete padded count}\label{cross-only-eligibility-and-a-complete-padded-count}

For the actual active set Sigma define \[
 E_\eta=\{v\in\Sigma:v_{\mathfrak P_v}(\eta-a_0)\ge1\},
 \quad a_0=Q(-7+2\sqrt{-7}),\quad a_0\bar a_0=U+14.
 \tag{FP27.5}
\] K097 RP4--7 proves the required local statement without global unipotence. Each original positive input y belongs separately to every actual source and target at both CM components; h\_ZZ=I gives Ny=Cy=hy-y in the target. A nonintegral Cy has its leading pole on the actual target pole line. Saturation excludes an eta zero or pole there, so eta is a unit and has the selected residue a0. At the companion component the slope is U/bar(a0); the exact norm relation makes this the same selected eligibility condition. At type I the full-pencil intersection makes Cy integral. Primitive integer Bezout extraction from u yields r(y), and the negative pivot beta=eta r(e\_j) is an original coordinate.

In particular all167 original r-values and beta have only simple paired poles supported on E\_eta. This is a cross-only statement. No rank-two unipotent assertion about full central double poles is used, and the full projected negative-column cost will remain.

Suppose \textbar E\_eta\textbar\textless=24. Pad E\_eta to a24-element subset Delta of the27 active positions. The complete number of containing choices is \[
                       \binom{27}{24}=2925.
 \tag{FP27.6}
\] For the original common whole-C orientation ideal J put \[
 L=JI_\Delta,\quad N_EI_\Delta=T^{-48},\quad
 z=(N_EL)^{1/(2f)}\ge T^{-24/f}.
 \tag{FP27.7}
\] All167 covector values and beta lie in this same fractional ideal. The at-most2\^{}668 G patterns are common to the entire cross map, with no independent transverse or pivot orientations. There is no separate eta count in this part: fixing r and beta reconstructs eta over E whenever the chosen original S pivot is nonzero.

The complete original-coordinate estimates are \[
 \#r\le A(4U)^f(M+4200z)^{334f}z^{-334f},\quad
 A=\frac{528^{167f}}{(167f)!D_0^{167}}<2^{-387281},
 \tag{FP27.8}
\] \[
 \#\beta\le28C_F B_0^{f/6}z^{-f/6},\quad
 C_F=\frac{2^{300}}{29^6 41^{14}2129^{14}}<2^{42},\qquad
 |c(h)|<2^{-929f/2}U^{-f/2}z^{-f}.
 \tag{FP27.9}
\] The full covector metric is diag(I164,G\^{}-1), with exact real covolume D0\^{}167 z\^{}(334f)/(4U)\^{}f and an arbitrary-ideal covering radius below4200 sqrt(f)z, by R076 CV8--10. Its single ball includes every coordinate support, so no binomial or k union is needed. The scalar fibre counts integral relative-norm ideals first and at most28 actual norm-one generators each, as in R076 CV11--13; nonprincipal and non-conjugation-stable ideals remain included. The trace is the actual selected normalized trace with both original compact factors, as in R061 GS11--12/R076 CV7.

Once C is known, retain the complete original projected negative-column and source-fibre estimates \[
       \#\operatorname{pr}_P(he_n)<191^{4f}U^fT^{108},
                      \qquad G_2<2^{36}.
 \tag{FP27.10}
\] R072 PJ14--18 proves the exact projection lattice (U\^{}-1 O\_E)u plus O\_E e\_n and its paired target index T\^{}(4s), including label zero. Fixing C and the plane projection reconstructs the entire negative column by unitarity. The remaining source-lattice isometry fibre fixes the actual positive166-space ker C and is bounded by G2. Neither the fixed U denominator nor that fibre is omitted.

The two genuine pointwise Green estimates on the same original G coefficient are \[
 |a(h)|\le2Q^{-335\cdot167}T^{335\cdot24}z^{335f},
 \qquad |a(h)|<2Q^{-501\cdot167}2^{9\cdot167f}T^{167t}.
 \tag{FP27.11}
\] The second is DO095 RA15--20: a nonzero original S entry of BC is an exact three-minor, has conjugates below2\^{}9 and only simple paired type-III poles. Its uniform norm bound requires no witness or additional G-pattern enumeration. Both mixed projectors retain modulus at most1.

The fixed factor is2\^{}881, comprising one Green2, G patterns2\^{}668, support choices2\^{}168, pivot2\^{}8 and the actual G2 fibre2\^{}36. For fixed t the complete general type and weighted paired-label sum is \[
 \binom{27}{t}(4T)^{27-t}<2^{P_t},\qquad
 P_t=81-2t+(27-t)L_T,\quad L_T=1019333/1024.
 \tag{FP27.12}
\] This is the exact R082 TP3--4 argument: type I costs below3T and type II costs at most1 at a non-type-III position; type III costs at most1. Counts are uniform for each original source/target pair before these weights are summed. There is no independent conjugate label and no extra3\^{}27 or coarse P\_T\^{}27 multiplier.

\subsection{3. Both complete ideal-norm ranges}\label{both-complete-ideal-norm-ranges}

Set z\_\emph{=2\^{}-92. The actual integer comparisons give M\textgreater2\^{}-68 and \[
 4200z_*/M<4200/2^{24}<2^{-11},\qquad
                    (1+2^{-11})^{334f}<2^{40}.
 \tag{FP27.13}
\] For z\textless=z\_} this bounds the full radius inflation in FP27.8 by2\^{}40. Combining its volume, the relative norm fibre and selected trace with weak Green leaves exactly z\^{}(-f/6), whose upper bound is T\^{}4 from FP27.7. The negative exponent is evaluated at the original lower ideal endpoint, not at z\_*.

For z\textgreater=z\_* use strong Green. The entire remaining dependence is \[
                 (M/z+4200)^{334f}z^{-7f/6},
 \tag{FP27.14}
\] strictly decreasing for z\textgreater0. Its supremum is therefore bounded at z\_*, with the same full2\^{}40 inflation. No ideal is assumed to attain the artificial cutoff. Both containing bounds are paid below, so possible overlap at the cutoff causes no omission.

Use the directed original integer logarithms \[
 L_Q=287843/4096,\quad L_H=28900/4096,\quad
 L_U=601356/4096,\quad L_{191}=15519/2048,
 \tag{FP27.15}
\] lower for Q,U and upper for H,191. Define \[
 C_*=881+40+5+42-387281+4fL_{191}-929f/2+2f
             +(2005f/6)(L_H-1)-(331f/2)L_U.
 \tag{FP27.16}
\] The40 is the entire radius inflation,5 bounds28,2f retains the determinant4U, and all projected-column/selected-metric U factors are included. The complete34 cells, before the2925 paddings, are \[
 E_t^-=C_*-335\cdot167L_Q+(335\cdot24+4+108)L_T+P_t,
\] \[
 E_t^+=C_*-501\cdot167L_Q+9\cdot167f
                  +(167t+108)L_T+P_t+92(2011f/6).
 \tag{FP27.17}
\] Exact rational arithmetic gives \[
 \max_{11\le t\le27}E_t^-=-6832571/4096<-1668\quad(t=11),
\] \[
 \max_{11\le t\le27}E_t^+=-212240581/4096<-51816\quad(t=27).
 \tag{FP27.18}
\] Thus the entire at-most24-eligible sector, at every denominator and every coordinate support, satisfies \[
      2925\cdot17(2^{-1668}+2^{-51816})<2^{-1652}.
 \tag{FP27.19}
\] Every support, type, padding, original label and source fibre has been counted. This is a fresh support27 count, not an extension of an old numerical charge by analogy.

\subsection{4. The whole denominator tail, without sparsity}\label{the-whole-denominator-tail-without-sparsity}

Let D\_eta be the exact denominator ideal and x=N\_E D\_eta. For each fixed integral ideal D, R083 LD3--6 proves \[
 \#\{\eta:\eta\bar\eta=U,\ D_\eta=D\}
             \le28\cdot4^4\cdot7^2\cdot2^{20}<2^{39}.
 \tag{FP27.20}
\] It counts all numerator ideals using the actual factorization of U and every norm-one root; nonprincipal candidates contribute no generator. There is no class-number factor or existence assertion.

Fix eta and the common G pattern J, and now permit all27 active paired pole places. Every original r-value lies in \[
 L'=JI_\Sigma\cap\eta^{-1}JI_\Sigma
       =JI_\Sigma D_\eta,\quad N_EI_\Sigma=T^{-54},
 \quad z'=(N_EL')^{1/(2f)}.
 \tag{FP27.21}
\] The valuation equality is v(L)+max(0,-v(eta)), valid for arbitrary fractional ideals. Because one original S value is nonzero, its ideal norm/AM--GM minimum and the joint real norm imply z'\textless M. Consequently \[
                 1+4200z'/M<4201<2^{49/4}.
 \tag{FP27.22}
\] Use the full167-dimensional volume and selected trace with z', together with the same entire projected column and G2 fibre. Weak Green is exactly \[
       2Q^{-335\cdot167}T^{335\cdot27}(z')^{335f}x^{-335/2}.
 \tag{FP27.23}
\] The z' powers cancel:335f-334f-f=0. Eta enumeration replaces the beta norm fibre; it is not multiplied by that old fibre. There is no sparsity restriction or k union. Every original type count is bounded by P\_11=59+16L\_T, and17\textless2\^{}5.

Define the complete fixed-denominator exponent \[
\begin{split}
 C_{\rm tail}={}&881+39-387281+4fL_{191}-929f/2+2f
       +334f(L_H-1)-(331f/2)L_U-335\cdot167L_Q\\
 &+(335\cdot27+108)L_T+(59+16L_T)+(49/4)334f.
\end{split}
 \tag{FP27.24}
\] For X=2\^{}10240, the complete sum over all integral E ideals is bounded by \[
 \sum_{N_ED\ge X}(N_ED)^{-335/2}
       \le X^{-331/2}\zeta_E(2)<X^{-331/2}2^{336}.
 \tag{FP27.25}
\] The elementary Euler bound is zeta\_E(2)\textless=zeta\_Q(2)\^{}336\textless2\^{}336; it includes every ideal class and even ideals which cannot be denominators. The full tail exponent, including all17 types, is \[
 C_{\rm tail}+5+336-(331/2)10240
                 =-51015031/4096<-12454.
 \tag{FP27.26}
\] Thus every denominator x\textgreater=2\^{}10240 is covered. No finite denominator window or support26 tail exponent has been imported.

\subsection{5. Literal inversion and23 common paired places}\label{literal-inversion-and23-common-paired-places}

DO093 IV4--16 proves the exact original identity W(h\^{}-1)=conjugate(W(h)). Inversion transposes the actual paired incidence matrix, retaining admitted zero labels and omitted labels, every absolute local weight and incidence type, both mixed projectors, the actual real trace and compact factors, cutoff H, central representative, class mass and finite/real Haar measures. The complete containing sectors FP27.19 and FP27.26 may each be applied to the inverse once, whether or not an already cut remainder is inverse-closed.

DO095 RA3--5 gives det D=1 and BD\^{}-1 C=0. Thus \[
 (h^{-1})_{PZ}=-D^{-1}C=\gamma(1,\xi)r,
 \quad\gamma\ne0,\quad\xi\bar\xi=U.
 \tag{FP27.27}
\] The first component of a nonzero null vector in diag(U,-1) cannot vanish. The original inverse S-cross is therefore nonzero; the full-positive identities, global rank and finite-unipotent premises are preserved. No independent inverse choice is made at individual embeddings.

Outside these four complete direct/inverse containing sectors, both denominator norms are less than2\^{}10240. Each eligible set has at least25 of the same27 active positions, so their intersection has at least23 positions. This is a deterministic implication, not an additional enumeration of eligible sets.

At a common position eta and xi both equal a0 in the selected residue field, and both partner residues equal U/bar(a0). Their denominators are units there. If delta=eta-xi were nonzero, the ideal (delta)D\_eta D\_xi would be integral and divisible by all46 paired degree-one primes at those23 positions. Its archimedean norm would satisfy \[
 T^{46}\le|N_E\delta|x_\eta x_\xi
       <(4U)^{168}2^{20480}<T^{46}.
 \tag{FP27.28}
\] The last strict inequality is an actual integer comparison, with a conservative308-bit gap. The first integrality statement follows primewise from min(a,b)+max(0,-a)+max(0,-b)\textgreater=0; no principality or conjugation stability of either denominator is assumed.

Hence eta=xi. By FP27.27, D preserves the original null line. With b=-ell*J\_P D this implies b ell=0, hence BC=0. DO095 RA9 then gives global rank(h-I)=2, contradicting FP27.1. This line-equality implication is support-independent and does not infer global unipotence.

Every FP27.1 term is consequently covered, and the whole direct/inverse union is \[
                   2(2^{-1652}+2^{-12454})<2^{-1650}.
 \tag{FP27.29}
\] Both inverse costs and all overlaps are retained. This proves FP27.2.

\subsection{6. Exact verification and chronological scope}\label{exact-verification-and-chronological-scope}

The executable checks the exact original host and directed integer logarithms, the full factorial volume and relative-norm constant, complete radius inflation, every one of the34 type/range cells, the2925-by17 padded union, the entire ideal tail and the46-prime denominator contradiction. Exact finite valuation and eligible-set models supplement the universal arguments. They do not enumerate arithmetic matrices or supply an arithmetic image witness.

The chronological baseline is R094120/31 with complete floors27/11. The default final author predecessor is canonical R095124/31, retaining K113304 followed by R09513000; \nolinkurl{--prior} permits a later full ordered canonical extension. Append1650 once, giving125/31 while preserving the full all-top and uniform arrays and every overlapping earlier charge. Both full error sums must remain below2\^{}-230. Complete floors remain27/11: the nonzero-short-compression sector and other support27 structural sectors are outside this theorem.

\subsection{7. Source and geometric boundaries}\label{source-and-geometric-boundaries}

Direct inherited interfaces are K097 RP4--7/RP10--21, R076 CV4--18, R083 LD3--15, R084 IL9--20, R090 IS1--25, DO093 IV4--16, DO095 RA1--20, R072 PJ4--18, R082 TP3--4 and DO074 TC1. The new work is the explicit support27 padded/t-profile count, the full all-k denominator tail and the25-of27/23-paired-place contradiction. No support26 numerical charge or missing compression premise is imported.

All ideals remain the original fractional invertible modules over Spec(O\_E) or Spec(O\_F). The relative norm fibre, denominator intersection and divisor-field norm use these exact modules. The fixed positive Hermitian Gram, projected U denominator, paired source/target lattices, polarization, theta line, Picard/biextension data and selected realization are unchanged.

Exact carried contextual locators are SGA3 I4.7.3/4.7.3.1; SGA6 VI Theorem3.3; FGA232 Theorem3.1 pp148--149; FGA236 Corollaries1.7/2.3 pp229/232; SGA7 I VII2.7/2.9 pp169--174 and VIII7.4(b) pp310--311; cote4 scan003 sections8.1.1--8.1.2, scans005--006 sections8.4.1--8.4.3; and cote54 scans005--007. These are inherited context, not new primary comparison claims. No new source constructor, selected arithmetic content, signed sufficient inequality, prescribed period, specific032/Read, evaluated DP44 or RH conclusion is asserted.

\subsection{Canonical algebraic-geometric context}\label{canonical-algebraic-geometric-context-57}

The actual fractional coordinate and denominator ideals remain invertible modules over Spec(O\_E)/Spec(O\_F). Relative ideal norm and the paired divisor of the line difference use those same modules; the projected U denominator, original Hermitian Gram, source/target lattice pair, selected metric, finite projectors and class/Haar factors are retained. Inversion is the original adjoint correspondence with its theta/Picard/Poincare-biextension data, not a new integral model. Exact inherited anchors: SGA3 I4.7.3/4.7.3.1; SGA6 VI Theorem3.3; FGA232 Theorem3.1 pp148-149; FGA236 Corollaries1.7/2.3 pp229/232; SGA7 I VII2.7/2.9 pp169-174 and VIII7.4(b) pp310-311; cote4 scan003 sections8.1.1-8.1.2 and scans005-006 sections8.4.1-8.4.3; cote54 scans005-007. No new lambda/K-theory, l-adic, Dieudonne or CM-motive comparison is inferred. DP\_MULT Q3e/DP44 and THETA3 G1/A7 receive only this conditional finite reduction; CX\_GRH CX-ALG-024 and GCT\_PLETH T-052 retain their arithmetic gates.

\clearpage

\section{AM-LIFT-096 --- Theta jets and the degree of the same-rule bad locus}\label{am-lift-096-theta-jets-and-the-degree-of-the-same-rule-bad-locus}

\textbf{Canonical checked integration --- 2026-11-23T08:45:49.089474+00:00.} Complete saved-proof and executable reviews and separate exact diagnostic replays passed. No numerical error charge is introduced. Frozen author history is preserved. AM-LIFT-096 TJ1-18 proves the same-rule bad-locus degree bound on the actual principal-3M special-unitary family of LIFT085/LIFT094, under its stated nonempty same-rule open-locus hypothesis. For g=fr,d=fpq, total frame weight w and theta weight n, take k=gw+n/2 for the literal even-n Dom; squaring the same rule gives k=2gw+n in general without changing its reduced zero set. A finite proof cover of degree at most e\_g=2\^{}(8g\^{}2+2g+4), first jets of the fourth power of the original symmetric principal line, and exact theta multiplier identities produce scalar cuts. For H=(psi\_bar\^{}\emph{L\_S)\^{}h very ample and V=deg\_H(Ybar), proper pushforward gives D\textless=D\_}=e\_g V sum\_\{j=0\}\^{}d k\^{}j, including mixed-dimensional components. When k=0 and the bad locus is proper it is empty. With LIFT095 c=1, LIFT094 then gives the conditional within-family multiplier N\textless=t\^{}\{B\_\{t\^{}\{4fr\^{}2\},fpq-1\}(D\_\emph{)\}, t=1+3M Q\_S; any independent component-entry multiplier N0 is retained. The proof cover contributes only to D\_}, not to the physical isogeny or output field. The numerical h,V, nonempty actual-family intersection, controlled N0, selected arithmetic nonzero rule and admitted full Host/Sizes remain open. No evaluated total N, six-identity chain, actual period, specific032 pairing, DP44 sign or RH conclusion; no finite error charge. Current K114126/31 exact arrays and complete floors27/11 are unchanged.

Author: /root/lift\_accounting. 23 November 2026 UTC.

\textbf{Complete author subject; exact diagnostic, saved-file verification and independent-review statuses are recorded separately. Not canonically integrated.} The argument below replaces an independent bad-locus degree input by an explicit expression involving the weight and the degree of a specified ambient arithmetic compactification. It does not evaluate the latter degree, a very-ampleness exponent, the component-entry multiplier, the actual same-rule nonempty-family condition, or the admitted Host/Sizes window. No effective actual source N, selected arithmetic nonzero rule, global period, DP44 or RH conclusion is asserted. The pre-CP46 research snapshot is retained separately and is not relabelled as this final subject.

\subsection{1. Exact family and proposed conclusion}\label{exact-family-and-proposed-conclusion}

Retain a literal source Host, Datum, Level, Moduli and Coeff, with A=1 as in LIFT085/094. Use exactly the LIFT094 arithmetic family Y=Gamma\D, where Gamma consists of the elements of SU(V)(F) preserving the actual lattice Lambda and acting as the identity on Lambda/(3M)Lambda. This group is neat by HE17--18. Assume the uniform signature (p,q), with p,q\textgreater=1 and r=p+q\textgreater=3. The actual nonterminal source applications have r\textgreater=7. Put f={[}F:Q{]}, g=fr and d=dim(Y)=fpq. Let Phi:Y-\textgreater X\_C be the regular period-family map of LIFT085 HF15--16 on this family. It need not be an isomorphism, surjective, or finite onto the literal moduli scheme. The original finite eta is preserved by Gamma; the auxiliary 3M marking is forgotten on the physical source objects.

The family has an algebraic principally polarized abelian scheme pi:A-\textgreater Y, its actual full OE action, finite marking, and the rigidified symmetric principal line L selected by the source trace-one element. Write

\[
 E_Y=e^*\Omega^1_{A/Y}=\bigoplus_\tau E_\tau,
 \quad \lambda=\det E_Y,
 \quad T=\pi_*L,
 \quad r_\tau=\operatorname{rank} E_\tau.                 \tag{TJ1}
\]

Here E\_tau is dual to the literal tau-eigen Lie bundle. Its rank is Dw.mult(tau), and sum\_tau r\_tau=g. The physical line L is not an auxiliary automorphic ample line.

Fix the \textbf{same} natural rule F0 in source sec(a,n), with all a(tau,i),n nonnegative integers. Put

\[
 w_\tau=\sum_{i<r_\tau}a(\tau,i),\qquad w=\sum_\tau w_\tau.
                                                               \tag{TJ2}
\]

Let Z be the reduced closed locus on Y where every actual eigenframe and theta-section generator gives value zero. Thus Z=Y minus Phi\^{}\{-1\}(U\_F0), with U\_F0 exactly LIFT085 HF11. For the escape application assume Z is proper. Nonzero F0 on the entire literal moduli space does not by itself prove this for a preselected Y.

Define

\[
 e_g=2^{8g^2+2g+4},\qquad
 k=gw+n/2\quad(n\text{ even}).                         \tag{TJ3}
\]

For arbitrary n replace F0 by F0\^{}2 only for constructing the equations; its zero locus is identical, and use k=2gw+n.~This is sanctioned by the literal sec\_mul and evalAt\_mul identities. In the actual Dom(a,n) case, Even n is literally a field of Dom, so the smaller expression applies.

Sections 2--5 construct a finite etale proof cover f':Y'-\textgreater Y of degree at most e\_g and finitely many algebraic sections of lambda\^{}k on Y' whose common zero locus is f'\^{}\{-1\}Z. Section 6 specifies an auxiliary very ample line H on the normal Baily--Borel compactification Ybar, with

\[
             H|_Y=\lambda^{a_H},\qquad a_H=(g+1)h,\quad h\ge1.
                                                               \tag{TJ4}
\]

For V=deg\_H(Ybar), the proposed conclusion is

\[
 \deg_H(\overline Z)\le
 e_g V\sum_{j=0}^{d} k^j.                              \tag{TJ5}
\]

Degree means the sum of degrees of all reduced irreducible components, including components of different dimensions, exactly as in LIFT094. If k=0 and Z is proper, the argument proves Z empty. Otherwise the integer on the right is a permissible D for LIFT094. The expression is uniform in the coefficients of F0, but V and h remain unevaluated.

\subsection{2. The actual coefficient bundle, including the theta sign}\label{the-actual-coefficient-bundle-including-the-theta-sign}

LIFT085 HF8--11 identifies the literal FrPt functor with the torsor of eigen Lie frames and nonzero generators of T. The natural rule is an algebraic regular function on this torsor. The exact source polynomial theorem gives, at every number-field point and fixed theta generator, a polynomial in the physical Lie-frame matrix coordinates. Diagonal frame changes give multihomogeneous degree w\_tau in the tau block. The physical transpose in IsFrPoly affects the coordinate display but not this degree or the underlying bundle of polynomial functions on ordered frames.

This fibrewise polynomial description is algebraic over the base. Here is the needed argument, rather than an assumption of regular coefficients. Work first on the literal X over the number field Coeff.F, on an affine open trivializing all eigen Lie bundles and T. The exact Moduli.smooth and qproj fields make X smooth and of finite type over this field, so it and its frame torsor are reduced. The restriction of the rule is a regular function on that open times the product of general linear groups and G\_m. Theta scaling is exactly s-\textgreater c s and F0-\textgreater c\^{}n F0; hence the last factor is t\^{}n times a regular function on the matrix factors.

The space of polynomials of the stated block degrees is finite-dimensional over Coeff.F. Its evaluations at matrices in the product of GL\_r(Coeff.F) span its dual, since these rational points are Zariski dense in the affine matrix space. Choose finitely many such evaluations giving an injective evaluation map and a linear left inverse. Applying that inverse to the regular values of the rule produces regular coefficient functions on the affine base. At each closed point the residue field is a number field, so the literal exists\_isFrPoly theorem identifies the rule with this polynomial. Equality on all closed points, and reducedness of the base and frame torsor, proves equality as regular functions. This construction glues uniquely. Only then pull it back along the regular map Phi, which need not be defined by polynomial period coordinates.

Consequently the rule determines an algebraic section of

\[
 W_F=\left(\bigotimes_\tau
     \operatorname{Sym}^{w_\tau}(E_\tau^{\oplus r_\tau})\right)
          \otimes T^{-n}.                              \tag{TJ6}
\]

The direct-sum multiplicity counts the ordered frame columns. The actual polynomial may lie in a smaller highest-weight subbundle, but no such identification is needed. The sign T\^{}\{-n\} follows from the literal positive scaling F0(c s)=c\^{}n F0(s): a degree-n function of a vector in T is an element of T\^{}\{-n\}. Similarly, a polynomial in Lie vectors is a tensor of cotangent vectors, so E\_tau, not E\_tau\^{}*, occurs in TJ6.

The coefficient section vanishes at y exactly when the polynomial is zero. A polynomial over C vanishing on all invertible frame matrices is zero, since those matrices form a dense open subset. Thus its zero locus is exactly Z. No preferred vacuum, new rule, or selected arithmetic derivative is substituted.

\subsection{3. A bounded finite theta cover}\label{a-bounded-finite-theta-cover}

The actual L is symmetric and rigidified. Its fourth power L\^{}4 is totally symmetric, has polarization type delta=(4,\ldots,4), and is very ample on every geometric fibre. The last assertion is the classical Lefschetz very-ampleness result for a power at least four; Lubicz--Robert (2022), Introduction p.1, explicitly records the n\textgreater3 form used here.

Over C, the symmetric-theta-structure torsor for L\^{}4 is finite etale. Candelori (2016), Lemma 2.5.3 and Corollary 2.5.4, give its group extension

\[
 0\longrightarrow K(\delta)[2]\longrightarrow
 \operatorname{Aut}_{\rm sym}G(\delta)
 \longrightarrow\operatorname{Sp}(K(\delta))\longrightarrow0.
                                                               \tag{TJ7}
\]

Here K(delta)=(Z/4)\^{}\{2g\}, so \textbar K(delta){[}2{]}\textbar=2\^{}\{2g\}. The elementary containing-matrix count gives \textbar Sp\_\{2g\}(Z/4)\textbar\textless=4\^{}\{(2g)\^{}2\}=2\^{}\{8g\^{}2\}. Hence the torsor degree is at most 2\^{}\{8g\^{}2+2g\}. This is a cover of the same family for a geometric degree proof; it is not a replacement of the physical source level, coefficient record, or output object.

A symmetric theta structure chooses a maximal symmetric splitting H of the theta group. The quotient phi:A-\textgreater B=A/H has a normalized symmetric principal line Theta with phi\^{}*Theta=L\^{}4. The standard weight-one representation gives, algebraically,

\[
 \pi_*L^4\simeq W\otimes K,\qquad
 W=\mathcal O^{4^g},\qquad K=\pi_{B*}\Theta.             \tag{TJ8}
\]

This is Candelori Theorem 4.1.3 and the complete proof of Theorem 4.2.1: the invariant line is the pushforward of the descended principal theta line. No analytic scalar normalization of a theta basis is required. Since phi is an isogeny in characteristic zero, its differential identifies the two Hodge bundles and their determinants.

For each normalized symmetric principal theta line, Candelori's key formula in the proof of Theorem 3.3.1 says that the inverse square of its pushforward is the Hodge determinant times a theta multiplier line. The latter is associated to a mu\_4 torsor (Definition 3.2.2), including both parities of theta characteristic. Apply this separately to the original (A,L) and the descended (B,Theta). Pull back both mu\_4 torsors, at a degree cost at most 16. On the resulting finite etale cover Y', of degree e\textless=e\_g,

\[
                    T^{-2}\simeq\lambda,
                    K^{-2}\simeq\lambda.               \tag{TJ9}
\]

Every connected component of this nonempty finite etale cover dominates connected Y; alternatively retain the whole cover and its total degree e. All subsequent estimates allow the whole cover. Trivializing these finite multiplier lines adds no claim about an arithmetic field of a chosen source point.

\subsection{4. First jets generate an explicit twist of the Hodge bundle}\label{first-jets-generate-an-explicit-twist-of-the-hodge-bundle}

Let V4=pi\_*L\^{}4 on Y'. Along the origin, the first principal-parts bundle has exact sequence

\[
 0\longrightarrow E_Y\longrightarrow J^1_e(L^4)
     \longrightarrow\mathcal O_{Y'}\longrightarrow0,    \tag{TJ10}
\]

using the given origin rigidification. The quotient of its second exterior power by the second exterior power of E\_Y is canonically E\_Y. Composing with the first-jet map gives

\[
                     \bigwedge^2 V4\longrightarrow E_Y.\tag{TJ11}
\]

In a local trivialization it sends s\_i wedge s\_j to s\_i(0) d s\_j(0)-s\_j(0) d s\_i(0). Changing the local trivialization adds the same multiple of the values to the differentials, and the two extra terms cancel. Therefore TJ11 requires no chosen connection or global splitting of TJ10.

The map is fibrewise surjective. Very ampleness gives a section s\_0 nonzero at the origin and an immersion by the complete linear system. The differentials at the origin of the ratios s\_j/s\_0 span the cotangent space; multiplying them by s\_0(0)\^{}2 gives the images in TJ11. Nakayama then proves surjectivity of the bundle map everywhere.

Using V4=W tensor K and K\^{}\{-2\}=lambda, TJ11 becomes a surjection

\[
           (\bigwedge^2W)\otimes\mathcal O_{Y'}
                         \twoheadrightarrow E_Y\otimes\lambda.
                                                               \tag{TJ12}
\]

Thus E\_Y tensor lambda is globally generated by finitely many actual algebraic jet minors. Taking the (g-1)-st exterior power yields global generation of

\[
 \bigwedge^{g-1}(E_Y\otimes\lambda)
                =E_Y^*\otimes\lambda^g.                \tag{TJ13}
\]

Each eigenbundle E\_tau\^{}* is a direct-summand quotient of E\_Y\^{}\emph{, so E\_tau\^{}} tensor lambda\^{}g is globally generated. This step uses the actual split OE action and the full Hodge determinant; it does not confuse a rank-r eigenblock determinant with lambda of rank g.

\subsection{5. Scalar equations for the same rule}\label{scalar-equations-for-the-same-rule}

Symmetric powers, direct sums and tensor products of globally generated locally free sheaves are globally generated. In characteristic zero, dualizing a symmetric power gives the corresponding symmetric power of the dual. For even n, TJ6, TJ9 and TJ13 therefore show

\[
                         W_F^*\otimes\lambda^k
                 \quad\hbox{is globally generated},\qquad k=gw+n/2.
                                                               \tag{TJ14}
\]

The generators may be chosen as finitely many tensor products of the jet-minor generators from TJ12. Pair them with the actual section s\_F of W\_F. This gives finitely many algebraic scalar sections sigma\_j of lambda\^{}k. At a point these pairings all vanish if and only if s\_F vanishes, because the generator values span the complete dual fibre. Hence their common zero locus is precisely f'\^{}\{-1\}Z.

If n is odd, perform the same construction with the literal square rule F0\^{}2 in sec(2a,2n), justified by IsCombAt.sec\_mul. Since the base is reduced and of characteristic zero, its pointwise vanishing locus is the same Z. The resulting k is 2gw+n.~This is only a geometric device for bounding that locus; the eventual nonzero value is still a value of F0.

The exponent k is an explicit twist obtained from global generation. It is not asserted to equal the projective degree of Z.

\subsection{6. The Siegel canonical line gives a permitted auxiliary H}\label{the-siegel-canonical-line-gives-a-permitted-auxiliary-h}

The trace-symplectic representation of Res\_\{F/Q\}SU(V) maps the domain D holomorphically into Siegel space H\_g. It is a totally geodesic embedding: at a base complex structure its Cartan involution is conjugation by that complex structure, and the inclusion into the symplectic group intertwines this involution. The tangent map is injective because the action on V is faithful modulo its finite central domain kernel. It is compatible with the tautological weight-one family of LIFT085.

The group Gamma from LIFT094 is contained in the principal 3M congruence subgroup of the integral symplectic group for the actual unimodular alternating lattice. Choose a symplectic Z-basis of that lattice. The integral 3M level is neat, and the family thus has a regular map to the corresponding Siegel arithmetic quotient. This is an auxiliary proof marking. The original finite eta, OE action and trace-one physical line are unchanged when this marking is forgotten.

Let L\_S be the canonical ample Baily--Borel line of that Siegel quotient, and let psi\_bar be the extension of the period-domain inclusion to Baily--Borel compactifications. Klingler--Yafaev (2014), Proposition 5.3.6 and its arithmetic-lattice setup, gives

\[
                \psi_{\rm bar}^*L_S\otimes L_Y^{-1}
                     \quad\hbox{nef},                  \tag{TJ15}
\]

where L\_Y is the canonical ample line on Ybar. A tensor product of an ample line and a nef line is ample on a projective variety. Therefore A\_H=psi\_bar\^{}*L\_S is ample. Choose h\textgreater=1 such that H=A\_H\^{}h is very ample. This is a specified permitted choice of the auxiliary line in LIFT094, rather than an identification of the physical theta line with that auxiliary line. Its exponent h is not numerically bounded here.

On the Siegel open, the algebraic Kodaira--Spencer isomorphism is Omega\^{}1=Sym\^{}2(E\_S). The determinant identity det(Sym\^{}2 E\_S)=(det E\_S)\^{}\{g+1\} gives L\_S\textbar=lambda\_S\^{}\{g+1\}. Pulling back proves TJ4. Kieffer--Page--Robert (2025), Section 4.3 and Proposition 4.8, printed pp.24--25 of the 45-page final author version, explicitly states and describes this algebraic all-g Kodaira--Spencer isomorphism before specializing later applications to genus two; the determinant computation here is elementary. The earlier 70-page full version has the same input in Section 4.3.1 and Proposition 4.22, printed pp.40--41. One can also derive it from the usual universal polarized deformation calculation.

The two Hecke correspondences retain equal pullbacks of H. Their rational unitary elements map to rational symplectic elements; apply the functorial canonical-line identity of Klingler--Yafaev Proposition 5.3.2(1) on the Siegel Hecke diagrams, then pull back. Thus choosing this H preserves the exact HE2 needed by LIFT094. No unitary determinant-character trivialization is assumed.

\subsection{7. Extension across the boundary}\label{extension-across-the-boundary}

The Q-group Res\_\{F/Q\}SU(V) has absolute simple factors of type A\_\{r-1\}; for r\textgreater=3 it has no PGL\_2 quotient. Milne SVI Remark 3.13(b), equivalently the setup of Klingler--Yafaev Section 5.3, gives codim(Ybar minus Y)\textgreater=2. If the group is anisotropic, the boundary is empty and the same argument is simpler. No rank-two modular-curve assertion is used.

Normalize Ybar in the function fields of the finitely many connected components of Y'. The result Ybar' is normal and projective, finite over Ybar; its inverse image of Y is exactly Y', since Y' is already finite etale and normal there. The boundary still has codimension at least two under a finite map. Set H'=f\_bar\^{}*H. It is ample and globally generated, though it need not be very ample; that suffices for the degree argument below.

Raise each scalar section sigma\_j to the positive integer a\_H. On Y' it is a section of lambda\^{}\{k a\_H\}=H'\^{}k. Its zero locus is unchanged. A line bundle on a normal Noetherian variety equals the pushforward of its restriction across a codimension-at-least-two complement. Therefore every sigma\_j\^{}\{a\_H\} extends uniquely to an algebraic section of H'\^{}k on Ybar'. This is the normal Hartogs/reflexive extension lemma, Stacks Tag 0EBJ (Lemma 31.13.12 in the presently read version), applied to an invertible sheaf. No arbitrary holomorphic function is asserted to be algebraic, and no unproved growth condition for a meromorphic automorphic form is substituted.

Let B be the reduced common zero set of these extended sections. Then B intersect Y'=f'\^{}\{-1\}Z. The closure of f'\^{}\{-1\}Z is the union of those irreducible components of B meeting Y'. Extra components contained entirely in the boundary can only enlarge the degree bound.

\subsection{8. Mixed-dimensional intersection bound}\label{mixed-dimensional-intersection-bound}

Here is a direct bound independent of the number of scalar equations. Let Xbar be a projective reduced pure-d-dimensional variety, L an ample globally generated line, and B the reduced common zero set of finitely many sections of L\^{}k, with k\textgreater=1 and B containing no component of Xbar. Write V\_X=deg\_L(Xbar).

Start with the effective cycle {[}Xbar{]}. Choose a generic linear combination of the given sections, avoiding every component not contained in B. Its divisor intersects those components properly. Whenever an irreducible component of the resulting cycle is contained in B, freeze it and retain its positive multiplicity; apply the next generic combination only to the remaining components. There are finitely many components at every stage, and the ground field is infinite, so one combination can avoid all of them. At codimension j the total degree of the active plus newly frozen effective cycle is at most k\^{}j V\_X, by the divisor-degree formula. After d stages only points remain; one more generic combination avoids all unfrozen points and removes them.

Every irreducible component of B occurs among the frozen components. At its generic point all chosen sections vanish. The sequence of active supports continues to contain that point until it becomes a component of the current support, which happens no later than its codimension; it is then frozen. Distinct frozen cycles and their multiplicities can overcount, which is harmless. Consequently

\[
                   \deg_L(B)\le V_X\sum_{j=0}^{d}k^j.   \tag{TJ16}
\]

This proof uses effective cycles on possibly singular varieties; proper intersections with a nonzero section of a line on an integral component give effective Cartier divisors, and the first-Chern-class degree formula is valid without nonsingularity. Embedded components do not reduce the reduced-component upper bound. The projection formula used here and below is Stacks Tag 02SU, already checked in LIFT094.

Apply this to the whole finite normal compactification Ybar'. Since f'\^{}\{-1\}Z is proper on every component, the extended scalar sections are not all identically zero on any component. Its total H'-degree is eV by finite pushforward. The closure of f'\^{}\{-1\}Z has degree at most eV sum\_j k\^{}j. Every component of Z has a component of its finite inverse image mapping onto it; the finite degree and the projection formula show that the degree of Z cannot exceed the summed degree upstairs. This proves TJ5 using e\textless=e\_g.

For k=0, the extended scalar sections are regular functions on each connected normal projective component of Ybar', hence constants. Properness of f'\^{}\{-1\}Z means that at least one such constant is nonzero on each component. Thus Z is empty. This boundary is separate from the k\textgreater=1 intersection induction.

\subsection{9. Consequence for LIFT094 and retained costs}\label{consequence-for-lift094-and-retained-costs}

For k\textgreater=1 define the integer

\[
 D_* =2^{8g^2+2g+4}\,\deg_H(Ybar)\sum_{j=0}^{fpq}k^j.   \tag{TJ17}
\]

It is a genuine degree upper bound for the same-rule inverse-image bad locus, assuming the family meets U\_F0. Combine it with the unchanged escape recurrence of LIFT094, and with c=1 from LIFT095 when its actual relative-unramified/selfdual hypotheses apply. On that family, if the controlled seed is already represented,

\[
 N\le t^{B_{t^{4fr^2},fpq-1}(D_*)},\qquad
 t=1+3M\prod_{p\in S}p.                               \tag{TJ18}
\]

If the seed enters through an additional admitted isogeny, its independent multiplier N0 multiplies the right side. When Z is empty the path length is zero and the within-family multiplier is one. The same original-field/torsion/model/line-cost theorems apply only under their unchanged hypotheses.

The theta cover is not part of the physical Hecke path. It is used solely to bound a degree of a closed subset of Y; after pushing that bound down, LIFT094 acts on the original Y and the original finite source level. In particular its degree e\_g is a factor in D\_*, not an extra output-field or physical-isogeny factor.

\begin{longtable}[]{@{}
  >{\raggedright\arraybackslash}p{(\columnwidth - 4\tabcolsep) * \real{0.3333}}
  >{\raggedright\arraybackslash}p{(\columnwidth - 4\tabcolsep) * \real{0.3333}}
  >{\raggedright\arraybackslash}p{(\columnwidth - 4\tabcolsep) * \real{0.3333}}@{}}
\toprule\noalign{}
\begin{minipage}[b]{\linewidth}\raggedright
Quantity or condition
\end{minipage} & \begin{minipage}[b]{\linewidth}\raggedright
What this route supplies
\end{minipage} & \begin{minipage}[b]{\linewidth}\raggedright
What remains
\end{minipage} \\
\midrule\noalign{}
\endhead
\bottomrule\noalign{}
\endlastfoot
Actual physical rule and theta line & The same F0 and the HF7 trace-one line throughout & A concrete nonzero selected arithmetic rule is not constructed \\
Finite cover degree & e\textless=2\^{}(8g\^{}2+2g+4) & No output-field interpretation is asserted \\
Coefficient dual twist & k=gw+n/2 for Even n; general k=2gw+n by the same zero locus & No claim that k equals deg Z \\
Compactification line & H=(psi\_bar\^{}*L\_S)\^{}h with H very ample and H & =lambda\^{}((g+1)h) \\
Same-rule bad-locus degree & D\_* once V and h are supplied & Its numerical value is not evaluated \\
Family nonvanishing & Conditional exact Phi\^{}\{-1\}(U\_F0) nonempty & No transfer from a different component without entry data \\
Actual Host/Sizes and DP44 & All old conditions and physical cost terms retained & No admissible effective window or DP44 sign \\
\end{longtable}

This is potentially a strict reduction of the geometric degree gate: D is no longer an independent input once the chosen compactification's total degree is bounded. It is not a completion of the remaining effectiveness problem.

\subsection{10. Fresh literal and primary reading record}\label{fresh-literal-and-primary-reading-record}

Literal source bodies read during this route:

\begin{itemize}
\tightlist
\item
  \nolinkurl{research/alpoge_source/Formalism/LocalToGlobal/ProjectiveLattices/EmbKCoe.lean:} complete 35 lines; Ix11, ex13--15, Dom21, embF27--28 and embK30--34. Dom includes Even n; sec alone does not.
\item
  \nolinkurl{research/alpoge_source/Formalism/LocalToGlobal/ProjectiveLattices/IsCombAt.lean:} complete 127 lines; FrPt67--77, FrMap79--85, FrBor87--96, sec100--101, evR103--104, sec\_mul106--125.
\item
  \nolinkurl{research/alpoge_source/Formalism/LocalToGlobal/ProjectiveLattices/AlgebraMapCAt.lean:} complete 147 lines; naturality35--36, weight38--39, secW41--48, FromTest51--52, evalAt54--58, evalAt\_mul72--74, EvalWD84--87.
\item
  \nolinkurl{recovery_pending/lift/source_cm/Development/AbelianVarietyTheory/Proofs/CubicalLineBundles/ExistsIsFramePolynomialOfRuleNumberField.lean:} complete saved theorem/proof, exists\_isFrPoly50; explicit number-field and every-base-change polynomial scope. This does not supply a degree for the moduli variable by itself.
\item
  \nolinkurl{recovery_pending/lift/source_cm/Development/AbelianVarietyTheory/Definitions/CMLattices/IntegralOf.lean:} complete 163 lines freshly read, including IsFrPoly97--102 and its physical transpose convention.
\item
  \nolinkurl{research/alpoge_source/ReciprocityTheory/ReductiveGroups/ArithmeticGeometry/BirationalGeometry/ToroidalCompactification/TwoWeightEqOneAddRoot.lean:} complete 168 lines; actual Host14--27, Datum30--49, SAB72--84, TateIso131--142, Obj144--160.
\item
  \nolinkurl{recovery_pending/lift/source_cm/ReciprocityTheory/ReductiveGroups/ArithmeticGeometry/BirationalGeometry/ToroidalCompactification/PushInvertible.lean:} complete 124 lines freshly read; Moduli29--38, including qproj, smooth and rep; Coeff40--45 and XF49--50. These supply the reduced finite-type literal base, not a numerical projective embedding degree.
\item
  Accepted canonical LIFT085 HF7--16 and HF24--26; LIFT094 HE0--22 and mandatory locator addendum; LIFT095 integral hyperbolic-frame theorem. LIFT094 and LIFT095 were now wholly reread in bounded, untruncated chunks; LIFT085's stated interface sections were reread. A pre-CP46 truncated LIFT094 output is superseded by this complete reading.
\end{itemize}

Primary passages freshly read, with their limited imported roles:

\begin{enumerate}
\def\labelenumi{\arabic{enumi}.}
\tightlist
\item
  Candelori, L. (2016). The transformation laws of algebraic theta functions. arXiv:1609.04486v1. \url{https://arxiv.org/pdf/1609.04486}. Sections 2.3--2.6, especially Lemma2.5.3/Corollary2.5.4 printed9--10; Section3.2 and Definition3.2.2 printed20; key formula in Theorem3.3.1 proof printed21; Theorems4.1.3/4.2.1 and complete proofs printed28--29. These give the finite symmetric theta cover, its finite multiplier lines and the algebraic weight-one bundle decomposition. The jet, coefficient and degree arguments above are supplied here.
\item
  Lubicz, D., \& Robert, D. (2022, June29). Fast change of level and applications to isogenies. \url{https://people.maths.bris.ac.uk/~jb12407/ANTS-XV/papers/ANTS-XV_lubicz-robert.pdf}. Introduction p.1, the n\textgreater3 Lefschetz very-ampleness statement. No isogeny-algorithm complexity is used.
\item
  Klingler, B., \& Yafaev, A. (2014). The Andre--Oort conjecture. Annals of Mathematics,180(3),867--925. Author PDF \url{https://www2.mathematik.hu-berlin.de/~klingleb/papiers/KY1309.pdf}. Section5.3, Proposition5.3.2(1), Remark5.3.4, Proposition5.3.6 and its arithmetic-lattice reduction, printed23--26. The compactification and line-bundle assertions are unconditional; no GRH/Andre--Oort conclusion is used.
\item
  Milne, J. S. (2017 revision). Introduction to Shimura varieties. \url{https://www.jmilne.org/math/xnotes/svi.pdf}. Theorem3.12 and Remark3.13 printed38--39, including the no-PGL2 boundary-codimension condition. The actual family uses the already checked Theorem14.8 and HF15 construction.
\item
  Kieffer, J., Page, A., \& Robert, D. (2025). Computing isogenies from modular equations in genus two. Journal of Algebra,666,331--386. \url{https://doi.org/10.1016/j.jalgebra.2024.11.029}. Final author PDF \url{https://www.math.u-bordeaux.fr/~damienrobert/pro/publications/articles/modular_isogenies_g2.pdf}, Section4.3 and Proposition4.8, printed24--25 (45-page version), freshly read. The earlier author full version \url{https://www.math.u-bordeaux.fr/~damienrobert/pro/publications/articles/modular_isogenies_g2_full.pdf}, Section4.3.1 and Proposition4.22, printed40--41 (70-page version), was also read. Only the all-g algebraic Omega=Sym2(Hodge) identity and its determinant are used, not a genus-two algorithm or an analytic coordinate-degree estimate.
\item
  The Stacks Project Authors. Reflexive modules, Tag0EBJ, \url{https://stacks.math.columbia.edu/tag/0EBJ}, presently Lemma31.13.12; exact statement and proof read. The first-Chern-class proper pushforward formula is Tag02SU, already read for LIFT094.
\end{enumerate}

\subsection{11. Verification boundary and the remaining application}\label{verification-boundary-and-the-remaining-application}

The companion executable checks complete first-jet evaluation and its gauge/basis covariance, signed exterior identities, the all-rank symmetric-square determinant, finite recovery of homogeneous coefficient functions on invertible matrices, both theta-weight conventions, exact small theta-structure counts, mixed-dimensional degree bookkeeping, and every frozen literal/dependency fingerprint. These are finite exact diagnostics. They do not computationally prove the imported theta, very-ampleness, Baily--Borel, Kodaira--Spencer or normal-extension theorems; those inputs and the universal algebraic arguments are stated above. Complete final saved-file author and independent readings have separate hash-specific, timestamped records.

In the inherited Grothendieck formulation, the new construction uses the relative principal-parts sequence of the actual rigidified cubical line, the finite etale theta-structure torsor, descent of its invariant section line, the representable frame torsor over the smooth source moduli scheme, and proper Chern-class pushforward on a normal compactification. It neither replaces the source line by a numerical polarization class nor changes the prescribed rule. No virtual K-theory identity, slope/purity assertion, CM-motive comparison, selected derivative invariance or global oscillator map is inferred. The FGA Picard/Poincare and SGA7 cubical-line inputs are precisely the carried LIFT085 context; no additional unverified archive theorem is invoked here.

The actual total compactification degree is the next numerical gate. Neither the natural-rule matrix degree nor a polynomial Cayley grid bounds it. The current construction supplies scalar algebraic cuts on a finite cover of the true family and pushes their geometric degree down; it does not turn the analytic period map into a polynomial map.

\subsection{Canonical algebraic-geometric context}\label{canonical-algebraic-geometric-context-58}

The construction uses the relative first-principal-parts sequence of the original rigidified cubical line, finite etale theta-structure and multiplier torsors, the descended invariant theta line, the representable frame torsor on the smooth literal moduli scheme, and proper first-Chern-class pushforward on the normal Baily-Borel compactification. Fresh primary locators: Candelori(2016) Lemma2.5.3/Corollary2.5.4 pp9-10, Definition3.2.2 p20, Theorem3.3.1 proof p21, Theorems4.1.3/4.2.1 pp28-29; Lubicz-Robert(2022) Introduction1.1 p1; Klingler-Yafaev(2014) Proposition5.3.2(1), Remark5.3.4, Proposition5.3.6 pp24-26; Milne(2017) Theorem3.12/Remark3.13 pp38-39; Kieffer-Page-Robert(2025) Section4.3/Proposition4.8 pp24-25 of the45-page author version; Stacks Tags0EBJ/02SU. The standard ample-plus-nef step has the additional exact root/peer locator Lazarsfeld(2004), Positivity in Algebraic Geometry I, Corollary1.4.10 pp53-54; this is a locator addition, not a change to the frozen proof. Carried Picard/biextension anchors: FGA232 Theorem3.1 pp148-149; FGA236 Corollaries1.7/2.3 pp229/232; SGA7 I VII2.7/2.9 pp169-174 and VIII7.4(b) pp310-311. Projective-module/exterior context: SGA3 I4.7.3/4.7.3.1 and SGA6 VI Theorem3.3. Formal-crystal and Cartier-duality context remains cote4 scan003 sections8.1.1-8.1.2,scans005-006 sections8.4.1-8.4.3, and cote54 scans005-007. These carried anchors do not turn the construction into a new slope, lambda-ring, K-theory, l-adic weight, CM-motive or global oscillator comparison. DP\_MULT Q3e/DP44/DP52/DP54-65 and THETA3 G1/A7 receive a more explicit conditional geometric-cost interface; CX\_GRH CX-ALG-024 and GCT\_PLETH T-052 retain their separate global gates.

\clearpage

\section{AM-DO-120 --- positive short energy and rank-four restoration at supports27--31}\label{am-do-120-positive-short-energy-and-rank-four-restoration-at-supports2731}

\textbf{Canonical checked integration --- 2026-11-23T08:54:03.101648+00:00.} Complete root and independent proof/script/interface reviews and separate exact replays passed. The actual canonical predecessor has126/31 entries; this integration gives127/31, complete floors27/11. Frozen author history is preserved. AM-DO-120 QW1-29 removes the entire original doubly-plane scalar-cross sector with nonzero four-entry short compression at every support27..31, including all global characteristic types, actual coordinate supports and paired labels. The original integral positive F witness q\_j=\textbar\textbar(A-I)y\_j\textbar\textbar\^{}2 has every conjugate below2\^{}42. The complete paired finite low-rank proof restores the rank-four Green baseline666 with coefficients331/2 and335/2; both original full source/target costs, selected trace, arbitrary-ideal row, relative norm fibre and all incidence weights remain. All five full endpoint bounds give one270871 charge. Appending it once to K114126/31 yields127/31, with both entire sums below2\^{}-230 and complete floors27/11 unchanged. The complementary scalar-cross remainder at supports27..31 has full positive compressionI, a positive166 fixed space and global rank at most3 by the support-free original integral scalar proof; global unipotence is not inferred. The dense all-II27 window remains open; no complete support28, actual admitted Host/Sizes, selected arithmetic content, prescribed period, specific032 value, evaluatedDP44 or RH conclusion.

Author: /root/definite\_oscillator,23 November2026 UTC. Controlled pending proof. Complete saved-proof, executable, source and independent review events are recorded separately. No canonical file is changed by this document.

\subsection{1. Original sector and complete interval}\label{original-sector-and-complete-interval}

Retain the actual J56 all-top H133 normalized kernel, f=168, n=169, original CM field E/F of degrees336/168, all334 unmixed finite-unipotent G reductions, both mixed projectors, every admitted paired T-pencil label, actual selected Hilbert metric and both compact factors, original near-point representative, central phase, class mass and Haar measures. Put S=E\^{}164. After the already charged DO099 interval, both the original and inverse S-crosses lie in the original packet plane \[
 P=Eu\oplus Ee_n,\qquad
 u=(2,95192390527,383092671043,12529843548438452388609),
\] \[
 U=u^*u=77Q^2-14,\quad Q=1427908193729121522601,
 \quad V=u^\perp\subset E^4,\quad Z=S\perp V.
 \tag{QW1}
\] Here P has form diag(U,-1), V is positive of dimension3, and the original S-cross is nonzero. DO099 supplies the two-plane premise throughout25--70; no later short-compression removal is imported.

The exact algebraic argument DO111 VC5--8 gives \[
 h=\begin{pmatrix}K&B\\C&D\end{pmatrix}_{Z\oplus P},\quad
 K=\operatorname{diag}(I_S,A),\quad A^\dagger G_V A=G_V,
 \quad (h^{-1})_{ZZ}=K^{-1},
 \quad \operatorname{rank}B,\operatorname{rank}C\le1.
 \tag{QW2}
\] In particular the S/V and V/S blocks of h and its inverse are zero. The full P-by-Z cross C has one common null image, by pairing the S and V columns. This is a rational identity and a complex positive-compression statement; it does not make a finite form positive and does not imply A=I.

Use the original integral positive vectors \[
 y_1=(-976627,-153032,38026,0),\quad
 y_2=(-567678,695570,-172838,0),
\] \[
 L_1=\|y_1\|^2=978665066829,\quad
 L_2=\|y_2\|^2=835948910828<L_1,
 \quad\Delta_{ij}=\langle y_i,(h-I)y_j\rangle.
 \tag{QW3}
\] \textbf{Theorem.} The union of the entire original doubly-plane scalar-cross sectors with \[
 27\le s\le31,\qquad (\Delta_{ij})_{1\le i,j\le2}\ne0
 \tag{QW4}
\] has absolute normalized contribution strictly below2\^{}-270871. Every incidence profile, original source/target pair, coordinate support and global characteristic type is included. This adds one fresh all-top charge270871, retaining all previous overlapping charges. The complete support floors remain27/11.

Outside this union at these supports, the support-independent DO111 VC23--25/R071 integral scalar argument gives K=I\_Z, a positive166-dimensional fixed space, and global rank(h-I)\textless=3. This consequence is only inside the stated doubly-plane scalar-cross class. No full support28 exhaustion or global unipotence is claimed.

\subsection{2. A positive integral F-field witness}\label{a-positive-integral-f-field-witness}

Each y\_j belongs separately to every actual source and target pencil at both CM components. The actual paired selfdual target therefore gives Delta\_jj in O\_E, as in R070 SC4 and DO111 VC10. Since y\_j is perpendicular to P, its diagonal compression is \textless y\_j,(A-I)y\_j\textgreater. Unitarity of A gives the exact identity \[
 q_j:=\langle(A-I)y_j,(A-I)y_j\rangle
       =-\Delta_{jj}-\overline{\Delta_{jj}}\in\mathcal O_F.
 \tag{QW5}
\] At every real embedding this is nonnegative and \[
              0\le\tau(q_j)\le4L_j\le4L_1<2^{42}.
 \tag{QW6}
\] If all four Delta entries do not vanish, at least one of (A-I)y\_1,(A-I)y\_2 is nonzero. The corresponding q\_j is a nonzero totally positive element of F. Indeed its positive norm cannot vanish at any embedding. We use whichever of these two original elements is nonzero only after assigning a divisor which divides both. There is no two-witness union.

The degree in its norm bound is f=168, not2f. Thus any integral F-ideal J\_q dividing that nonzero q\_j satisfies \[
                          N_FJ_q<2^{42f}.
 \tag{QW7}
\] No denominator, division by2, moving frame or inferred global quotient lattice occurs in QW5--7.

\subsection{3. Universal finite low-rank compression lemma}\label{universal-finite-low-rank-compression-lemma}

Work over either residue field at an actual split G place, in the fixed G-unit S/V/P frame. The transverse frame (y\_1,y\_2,z\_0), with \[
 z_0=(256766,-7654233287,-30805128044,1),
\] has Gram determinant4U, and2U is a unit at every actual G prime. All displayed matrices and their inverses are integral there. Let N=h-I, d=rank N, E\_0=A-I. The reduction of h is unipotent by the original K068 cut. Thus N\^{}(d+1)=0. The paired component retains the exact inverse-adjoint relation and A\_c is the corresponding inverse adjoint of A. No finite positivity is used.

Denote the original lower S-cross by C\_S and upper S-cross by B\_S. The zero/nonzero status of B\_S equals that of the lower S-cross at the companion component: block inversion with the identity S block gives an invertible Schur factor, exactly R070 SC7. Consequently the labels ``both'', ``one'' and ``neither'' zero C blocks below have their original paired meaning.

Projection of im N onto S has rank one if B\_S is nonzero. Its kernel then has dimension d-1. All N(V) lie in this kernel. If C\_S is also nonzero, the kernel contains the nonzero line N(S) in P, which projects to zero on V. It follows that \[
 \operatorname{rank}E_0\le d-2\quad(B_S,C_S\ne0),
 \qquad
 \operatorname{rank}E_0\le d-1\quad\text{exactly one nonzero}.
 \tag{QW8}
\] For the second inequality when B\_S=0,C\_S!=0, apply the same dimension argument to the transpose. In particular E\_0=0 at d=2 with both S-crosses nonzero, and at d=1 with exactly one nonzero. At d=0 it is zero as well.

We next prove the needed square-zero conclusion in the remaining low-rank cells. Jacobi's complementary-minor identity and QW2 give det D=det h/det K!=0. Block inversion gives \[
                             BD^{-1}C=0.
 \tag{QW9}
\] If either entire cross B or C is zero, h is block triangular, so K and A are unipotent. Whenever QW8 bounds rank E\_0 by1, this implies E\_0\^{}2=0. This whole-cross-zero case is handled before choosing rank-one factors.

If both entire crosses are nonzero, factor \[
              B=Kv b,\qquad C=\ell r,\qquad
                   bD^{-1}\ell=0,\quad\kappa=b\ell,
 \tag{QW10}
\] with nonzero column v, row b, column ell and row r. All this is linear algebra over the current residue field. The S parts are B\_S=v\_S b and C\_S=ell r\_S. No globally integral null basis is asserted.

If kappa=0, ker b is the line spanned by ell. The two-dimensional Cayley identity D\^{}-1=((tr D)I-D)/det D and QW10 give bD ell=0. Hence D preserves ell. Both the line ell and Z plus ell are invariant under h, and the induced map on (Z plus ell)/ell is K. Thus A is unipotent. Combined with rank E\_0\textless=1, this gives E\_0\^{}2=0.

Suppose d=3 and both S-crosses are nonzero. If kappa!=0, choose an S row i with v\_i!=0 and an S column j with r\_j!=0; these indices need not coincide. In rows \{i,P\} and columns \{j,P\}, N has the3-by3 submatrix \[
 M_{ij}=\begin{pmatrix}0&v_i b\\\ell r_j&D-I\end{pmatrix}.
 \tag{QW11}
\] The identity b adj(D)ell=0 and adj(D-I)=adj(D)-I show \[
                      \det M_{ij}=v_i r_j\kappa\ne0.
 \tag{QW12}
\] Its inverse sends (0,ell)\^{}t to (1/r\_j,0,0)\^{}t. Since rank N=3, the Schur complement relative to this minor is zero. Its V-by-V entry is E\_0 minus the product of (0,A v\_V b), M\_ij\^{}-1 and (0,ell r\_V)\^{}t. That product vanishes because its first factor has zero S-column entry. Therefore E\_0=0. Together with the kappa=0 case and QW8 this proves \[
        d=3,\ B_S,C_S\ne0\quad\Longrightarrow\quad E_0^2=0.
 \tag{QW13}
\]

Now suppose d=2 and exactly one S-cross is nonzero. For C\_S=0,B\_S!=0, N\^{}3=0 and the zero S/V block of h\^{}-1=I-N+N\^{}2 give \[
                  0=(N^2)_{SV}=B_S C_V=v_S\kappa r_V.
 \tag{QW14}
\] If kappa!=0 then r\_V=0; also r\_S=0 because C\_S=0. Hence the entire C is zero, and the triangular case applies. If kappa=0 the invariant-flag argument above applies. In either case A is unipotent and QW8 gives E\_0\^{}2=0. For the opposite S-cross orientation, transpose the entire block argument. QW2 and QW9 are preserved with K\^{}t,D\^{}t and the exchanged crosses, so (E\_0\^{}t)\^{}2=0 and again E\_0\^{}2=0.

Finally, if d=1 and both S-crosses are zero, N\^{}2=0. Its inverse V/V block is I-E\_0, while QW2 makes that block A\^{}-1. Thus (I+E\_0)(I-E\_0)=I and E\_0\^{}2=0. These arguments cover every cell used below, without a claim about higher finite ranks.

\subsection{4. Exact one- and two-power F-valuations and Green restoration}\label{exact-one--and-two-power-f-valuations-and-green-restoration}

Over the original G-adic paired lattices, \[
 q_j=\langle y_j,(2I-A-A^{-1})y_j\rangle
     =\langle E_0y_j,E_0y_j\rangle.
 \tag{QW15}
\] If E\_0\^{}2=0 modulo the prime, then A\^{}-1=I-E\_0 there and the first formula gives q\_j=0 modulo that F prime. If E\_0=0 modulo the prime, A\_c=G\_c\^{}-1 A\^{}-t G\_c also equals I at the companion. The second formula is then a product of two integral vectors each divisible by the local uniformizer, so v\_F(q\_j)\textgreater=2. The G-unit Gram creates no denominator. These statements hold for both j=1,2, including when a reduced y\_j vanishes.

Define the original S-cross ideal J\_C by the indicated vanishing CM components, and an F-ideal J\_q with the indicated exponent. All choices are deterministic. The complete table is

\begin{longtable}[]{@{}
  >{\raggedleft\arraybackslash}p{(\columnwidth - 8\tabcolsep) * \real{0.2222}}
  >{\raggedright\arraybackslash}p{(\columnwidth - 8\tabcolsep) * \real{0.1667}}
  >{\raggedleft\arraybackslash}p{(\columnwidth - 8\tabcolsep) * \real{0.2222}}
  >{\raggedleft\arraybackslash}p{(\columnwidth - 8\tabcolsep) * \real{0.2222}}
  >{\raggedright\arraybackslash}p{(\columnwidth - 8\tabcolsep) * \real{0.1667}}@{}}
\toprule\noalign{}
\begin{minipage}[b]{\linewidth}\raggedleft
Finite rank d
\end{minipage} & \begin{minipage}[b]{\linewidth}\raggedright
Zero lower S-crosses at the CM pair
\end{minipage} & \begin{minipage}[b]{\linewidth}\raggedleft
E-norm exponent j in J\_C
\end{minipage} & \begin{minipage}[b]{\linewidth}\raggedleft
F-prime exponent w in J\_q
\end{minipage} & \begin{minipage}[b]{\linewidth}\raggedright
Reason for w
\end{minipage} \\
\midrule\noalign{}
\endhead
\bottomrule\noalign{}
\endlastfoot
0 & Both & 2 & 2 & E\_0=0 \\
1 & Exactly one & 1 & 2 & E\_0=0 by QW8 \\
1 & Both & 2 & 1 & E\_0\^{}2=0 by the inverse identity \\
2 & Neither & 0 & 2 & E\_0=0 by QW8 \\
2 & Exactly one & 1 & 1 & QW14 and its transpose \\
2 & Both & 2 & 0 & No compression condition required \\
3 & Neither & 0 & 1 & QW13 \\
3 & Exactly one & 1 & 0 & No compression condition required \\
3 & Both & 2 & 0 & No compression condition required \\
At least4 & Any & 0 & 0 & Baseline rank-four decay \\
\end{longtable}

Rank one with neither lower cross zero is impossible: a rank-one N with N\_SS=0 cannot have both B\_S and C\_S nonzero. At every assigned factor, J\_C divides all original C\_S entries, while J\_q divides both q\_1 and q\_2. There are at most4\^{}334=2\^{}668 J\_C patterns, including the empty choice. The varying J\_q will be norm-bounded uniformly, so there is no second pattern enumeration.

The actual normalized unipotent Green exponents are \[
 L_d=d(337-d)/2,\quad(L_0,L_1,L_2,L_3,L_4)=(0,168,335,501,666).
 \tag{QW16}
\] The accepted original Euler correction for the first four factors is below2 over all334 places; both mixed factors have modulus at most1. Every table row satisfies \[
                \frac{331}{2}j+\frac{335}{2}w\ge666-L_d.
 \tag{QW17}
\] For d\textgreater=4 the rank-four baseline suffices. Therefore the original finite coefficient satisfies \[
 |a_G(h)|\le
 2Q^{-666\cdot167}(N_EJ_C)^{331/2}(N_FJ_q)^{335/2}
 <2Q^{-666\cdot167}(N_EJ_C)^{331/2}2^{42\cdot335f/2}.
 \tag{QW18}
\] The final inequality uses the one nonzero actual q\_j and QW7. This is a pointwise bound with no witness-choice factor. It changes no finite projector or measure.

\subsubsection{Sharp scope boundary for the finite lemma}\label{sharp-scope-boundary-for-the-finite-lemma}

The generic conclusion that A must be unipotent is false. Over any odd residue characteristic, in coordinates S1,V2,P2, take \[
 h=\begin{pmatrix}
 1&0&0&1&0\\
 0&1&0&1&0\\
 0&0&2&1&0\\
 1&-1&-1/2&1&1\\
 0&0&0&-1/2&0
 \end{pmatrix}.
 \tag{QW19}
\] It has characteristic polynomial (X-1)\^{}5 and rank(h-I)=4, while K=diag(1,1,2) and (h\^{}-1)\_ZZ=diag(1,1,1/2). Both S-crosses are nonzero; C has image E e\_1 in P and the inverse C has image E e\_2. Pair the other split component by J\^{}-1 h\^{}-t J with hyperbolic J\_P. Then all required split Hermitian and inverse-compression identities hold, but A=diag(1,2) is not unipotent. One may pad V by an identity coordinate. This is a local scope counterexample, not a globally admissible short pencil arrow. It does not contradict the low-rank table.

\subsection{5. Entire unchanged original matrix count}\label{entire-unchanged-original-matrix-count}

Put L=J\_C I\_Sigma, N\_E I\_Sigma=T\^{}(-2s) and z=(N\_E L)\^{}(1/(2f)). The original primitive u extraction, separately at both CM partners and at every actual source/target label, gives \[
 C_S e_i=(u\xi_i,\beta_i),\quad \xi_i,\beta_i\in L,
 \quad\|\xi_\tau\|<M=H/(2\sqrt U),\quad|\beta_\tau|<b=H/2.
 \tag{QW20}
\] The whole-C ideal here means the complete original W0-by-S cross, not an unpaid extension to the three V inputs. No separate A enumeration or projected-column reconstruction is used.

The actual half-trace row lattice has covolume D0\^{}164 z\^{}(328f), with \[
 D_0=7^{140}29^{144}5^{84}13^{84},\qquad
 A_{164}=\frac{528^{164f}}{(164f)!D_0^{164}}.
\] The original arbitrary-ideal covering radius is below4117sqrt(f)z. Hence the single full row ball gives \[
                   \#\xi\le A_{164}(M+4117z)^{328f}z^{-328f}.
 \tag{QW21}
\] This includes every original zero pattern in one ball; no k-binomial or k-union is needed. The actual row is nonzero, so retain at most164\textless2\^{}8 pivot choices for reconstruction. The full relative norm fibre, including all28 possible norm-one generators, remains DO099 JR18/DO111 VC16: \[
                         \#\beta<2^{249}b^{f/8}z^{-f/8}.
 \tag{QW22}
\] The quotient beta/xi reconstructs the line only over E. No integral moving basis or global denominator hypothesis is assumed.

Retain every other original factor exactly as in DO111 VC17--19: \[
 |c(h)|<2^{-929f/2}U^{-f/2}z^{-f},\quad
 \#(D_{W0}e_n)<2^{5757}T^{10s},\quad G_5<2^{107},
\] \[
                  P_T^s<2^{997s},\qquad P_T=2(T-1)^2/(T-2).
 \tag{QW23}
\] The selected trace retains the original near-point and both compact factors. The target column is the full five-complex-coordinate original target-overlattice count, uniform at every label including zero. Given C\_S and that W0 column, unitarity reconstructs its S component. The unchanged source-lattice fibre fixes the original positive163-space ker C\_S and the negative column and has the full G5 bound. No forward-line stabilization, compact-compression enumeration, smaller source fibre or two-coordinate target replacement is inferred.

P\_T\^{}s includes all original types, paired labels and absolute incidence weights. No additional3\^{}s, partner-label or rank-pattern count is inserted. Including the Green correction,2\^{}668 common patterns,2\^{}168 supports,2\^{}8 pivots, full G5, target constant and beta fibre, the fixed exponent remains \[
                      1+668+168+8+107+5757+249=6958.
 \tag{QW24}
\]

\subsection{6. All five exact endpoints and one union}\label{all-five-exact-endpoints-and-one-union}

The original ideal cap is N\_EJ\_C\textless=Q\^{}334. Since N\_EJ\_C=T\^{}(2s)z\^{}(2f), multiplication of QW18 and QW21--23 leaves the entire varying factor \[
             z^{315}(M+4117z)^{328f},\qquad
       315=331f-328f-f/8-f=15f/8>0.
 \tag{QW25}
\] It increases for every z\textgreater0. The count is thus bounded at the genuine common-ideal upper endpoint, without assuming an ideal attains it. The exact original integers imply M\textless2\^{}-67 and the following strict caps and least strict rational-volume ceilings:

\begin{longtable}[]{@{}
  >{\raggedleft\arraybackslash}p{(\columnwidth - 4\tabcolsep) * \real{0.3333}}
  >{\raggedleft\arraybackslash}p{(\columnwidth - 4\tabcolsep) * \real{0.3333}}
  >{\raggedleft\arraybackslash}p{(\columnwidth - 4\tabcolsep) * \real{0.3333}}@{}}
\toprule\noalign{}
\begin{minipage}[b]{\linewidth}\raggedleft
s
\end{minipage} & \begin{minipage}[b]{\linewidth}\raggedleft
b\_s with z\textless2\^{}-b\_s
\end{minipage} & \begin{minipage}[b]{\linewidth}\raggedleft
v\_s, the strict binary ceiling of A164(2\^{}-67+4117*2\^{}-b\_s)\^{}(328f)
\end{minipage} \\
\midrule\noalign{}
\endhead
\bottomrule\noalign{}
\endlastfoot
27 & 90 & -4071532 \\
28 & 96 & -4071571 \\
29 & 101 & -4071571 \\
30 & 107 & -4071571 \\
31 & 113 & -4071571 \\
\end{longtable}

Each cap is checked by Q\^{}167*2\^{}(b\_s f)\textless T\^{}s, and each ceiling by exact integer numerator/denominator powers. The row volume uses the explicit rational upper bound22/7 for pi, not an asserted equality with the pi-volume.

Use the same directed logarithms \[
 L_T=1019333/1024>\log_2T,\quad L_Q=287843/4096<\log_2Q,
\] \[
 L_U=601356/4096<\log_2U,\quad L_H=28900/4096>\log_2H.
 \tag{QW26}
\] All original factors give the complete exponent \[
\begin{split}
 E_s={}&6958-929f/2+42\cdot335f/2+(f/8)(L_H-1)
       -(f/2)L_U-666\cdot167L_Q\\
      &+(341L_T+997)s+v_s-315b_s.
\end{split}
 \tag{QW27}
\] The T power341s is331s from the original C ideal plus10s from the full target column. Both the beta fibre and genuine trace have already entered QW25. Exact rational arithmetic gives

\begin{longtable}[]{@{}rr@{}}
\toprule\noalign{}
s & Complete E\_s \\
\midrule\noalign{}
\endhead
\bottomrule\noalign{}
\endlastfoot
27 & -3328741449/2048 \\
28 & -2635465079/2048 \\
29 & -1941463717/2048 \\
30 & -1248107475/2048 \\
31 & -554751233/2048 \\
\end{longtable}

All five are strictly below-270874. Therefore \[
       \sum_{s=27}^{31}|\text{complete N1 contribution at }s|
                         <5\,2^{-270874}<2^{-270871}.
 \tag{QW28}
\] There is no additional endpoint, coordinate, inverse or witness union. The same displayed sufficient count at s32 is positive138605009/2048 and proves nothing about that sector. No existence or lower-bound conclusion is drawn from this failure.

\subsection{7. Structural consequence, verification and exact scope}\label{structural-consequence-verification-and-exact-scope}

When all four short compressions vanish, unitary A fixes span(y1,y2) pointwise. DO111 VC23--24 then uses the original integral vector z0, M0=u1\^{}2+u2\^{}2+u3\^{}2, gcd(U,M0)=9 and M0\^{}2\textgreater36U to force its remaining unitary scalar to equal1. This proof uses the actual original integral pairing and the algebraic-integer norm; it has no support parameter. Thus A=I and QW2 implies \[
 h_{ZZ}=I_Z,\qquad\dim(\ker(h-I)\cap Z)=166,
                         \operatorname{rank}_E(h-I)\le3.
 \tag{QW29}
\] Individual y\_j need not be fixed by h, since their P components may be nonzero. This statement does not assume global unipotence and is not applied before the short-energy sector has been bounded.

The companion checks the exact fixed frame, field degree and witness sizes; every local table inequality; the universal low-rank identities through exact symbolic and finite paired models; the explicit rank-four boundary model; both one- and two-power paired valuations; the complete original target/source/pair constants; all five full endpoints; the positive s32 diagnostic; and both complete ordered error sums. Models supplement the universal proof and assert no admissible global matrix exists.

The default current predecessor is checks/k114\_joint\_remainder\_error.json,126 all-top and31 uniform entries, complete floors27/11. The executable accepts a later complete ordered canonical extension. Append270871 once, preserving every earlier overlapping charge and every uniform entry. Both complete sums must remain below2\^{}-230. This finding alone makes no complete support-floor increase.

Direct inherited mathematical interfaces are DO099 JR1--24, DO111 VC1--25, R070 SC1--19, R061 GS11--15, DO070's arbitrary fractional-ideal cover, DO092 SN10--11, K068's actual finite-unipotence cut and K079 ER7--9's normalized Green exponents/Euler correction. No new external representation theorem is invoked. The new proof is QW5--18 and its complete original-count interval evaluation QW20--28.

The original Hermitian modules, invertible ideals and integral pairings remain over Spec(O\_E) and Spec(O\_F); their field norms are the same determinants of multiplication. Polarization, theta line, Picard/biextension data, complex realization, finite projectors and Haar/class normalizations are unchanged. Carried contextual anchors remain SGA3 I4.7.3/4.7.3.1; SGA6 VI Theorem3.3; FGA232 Theorem3.1 pp148--149; FGA236 Corollaries1.7/2.3 pp229/232; SGA7 I VII2.7/2.9 pp169--174 and VIII7.4(b) pp310--311; cote4 scan003 sections8.1.1--8.1.2, scans005--006 sections8.4.1--8.4.3; and cote54 scans005--007. These are inherited context, not a fresh primary-reading claim. DP\_MULT Q3e/DP44, THETA3 G1/A7, CX\_GRH CX-ALG-024 and GCT\_PLETH T-052 retain their separate signed, source, arithmetic-image, prescribed-period and global comparison gates. No Lean build or archive audit is performed.

\subsection{Canonical algebraic-geometric context}\label{canonical-algebraic-geometric-context-59}

The new q\_j are integral trace/norm pairing sections over Spec(O\_F) from the original Hermitian modules over Spec(O\_E). Their assigned F-divisor is simultaneous for both original short vectors; its norm is the same determinant of multiplication. The finite low-rank flags and Schur minors live in the actual G-unit reductions and assert no integral global quotient or moving integral frame. The original fractional row ideals, CM involution, metric, finite projectors, classes, Haar measures, polarization, theta line, Picard data and Poincare biextensions remain unchanged. Carried exact anchors: SGA3 I4.7.3/4.7.3.1; SGA6 VI Theorem3.3; FGA232 Theorem3.1 pp148-149; FGA236 Corollaries1.7/2.3 pp229/232; SGA7 I VII2.7/2.9 pp169-174 and VIII7.4(b) pp310-311; cote4 scan003 sections8.1.1-8.1.2 and scans005-006 sections8.4.1-8.4.3; cote54 scans005-007. No new lambda-ring/K-theory, l-adic weight, Dieudonne or CM-motive comparison is inferred. DP\_MULT Q3e/DP44 and THETA3 G1/A7 receive this strictly smaller conditional finite remainder; CX\_GRH CX-ALG-024 and GCT\_PLETH T-052 retain their global gates.

\clearpage

\section{AM-LIFT-097 --- An explicit very-ampleness exponent on the same arithmetic compactification}\label{am-lift-097-an-explicit-very-ampleness-exponent-on-the-same-arithmetic-compactification}

\textbf{Canonical checked integration --- 2026-11-23T08:55:55.887761+00:00.} Complete saved-proof and executable reviews and separate exact diagnostic replays passed. No numerical error charge is introduced. Frozen author history is preserved. AM-LIFT-097 VA1-15 evaluates the previously unspecified very-ampleness exponent on the same actual principal-3M arithmetic compactification X. For d=fpq,g=fr and A=psi\_bar\^{}\emph{L\_S, the exact integer h\_d=2\^{}(d+1)(d+1)(d+2)! makes H=A\^{}h\_d very ample. The no-PGL2 boundary condition and canonical-line extension identify K\_X as Cartier ample; the neat Baily-Borel pair is log canonical, A-K\_X is nef, and both applications of effective freeness/very ampleness retain log bigness on every lc/slc stratum. The original Hodge restriction and Hecke equal-pullback identities are unchanged. Writing v\_A=int\_X c1(A)\^{}d, the full ambient degree is h\_d\^{}d v\_A, so LIFT096 gives D\_}=2\^{}(8g\^{}2+2g+4)h\_d\^{}d v\_A sum\_\{j=0\}\^{}d k\^{}j, with k=gw+n/2 for even n or2gw+n after the same-zero square. The independent quantity v\_A is a positive unevaluated integer. LIFT094/095 then retain N\_within\textless=t\^{}\{B\_\{t\^{}(4fr\^{}2),d-1\}(D\_*)\},t=1+3M Q\_S, with every protected congruence condition and any separate component-entry N0. The same-family nonempty open, selected arithmetic nonzero rule and admitted Host/Sizes remain open. No evaluated actual N, new physical field/level/isogeny, six-identity chain, period, specific032 pairing, DP44 or RH conclusion; no finite error charge. Current DO120127/31 arrays and complete floors27/11 are unchanged.

Author: \nolinkurl{/root/lift_accounting}. 23 November 2026 UTC.

\textbf{Complete author subject; final saved-file, source and diagnostic verification are recorded separately. Not canonically integrated.} This is a separate consequence for the family and auxiliary line of AM-LIFT-096. Its frozen proof is unchanged. The result evaluates the previously unspecified exponent h. The remaining total intersection number, nonempty same-rule family intersection, component-entry multiplier and actual Host/Sizes construction are retained explicitly. There is no numerical finite-kernel charge or actual global-period, DP44 or RH conclusion.

\subsection{1. Exact statement and source family}\label{exact-statement-and-source-family}

Retain the precise family of LIFT094 and LIFT096. Thus E/F is the admitted CM extension, f={[}F:Q{]}, (V,h\_V) is the actual Hermitian space of uniform signature (p,q), p,q\textgreater=1, r=p+q\textgreater=3, and Lambda is its original lattice. Set

\[
 \Gamma=\{\gamma\in SU(V)(F):\gamma\Lambda=\Lambda,
              \ (\gamma-1)\Lambda\subseteq3M\Lambda\},
 \qquad Y=\Gamma\backslash\mathcal D,
 \qquad d=fpq,\quad g=fr.                              \tag{VA1}
\]

Use one connected domain and its quotient, with the same actual finite eta and the regular family map to the literal source moduli scheme. The auxiliary 3M marking is forgotten on the physical objects exactly as in LIFT096. The argument is conditional on that admitted family; it does not construct a Host, change a coefficient field, or establish the nonvanishing of a selected rule on this component.

Let X=Ybar be the normal Baily--Borel compactification. Let psi\_bar:X-\textgreater Sbar be the extension of the trace-symplectic domain map to the Siegel Baily--Borel compactification at the actual integral 3M level. Write L\_S for the latter's canonical ample automorphic line and set

\[
             A=\psi_{\rm bar}^{*}L_S,
 \qquad h_d=2^{d+1}(d+1)(d+2)!,
 \qquad H=A^{\otimes h_d}.                              \tag{VA2}
\]

Then H is very ample on this exact X. On Y it is the line

\[
                    H|_Y\simeq\lambda^{(g+1)h_d},       \tag{VA3}
\]

where lambda is the determinant of the full rank-g Hodge bundle of the actual abelian family. The equal-pullback identities required for both LIFT094 Hecke correspondences hold for H. No new cover, source level or physical isogeny is introduced by this choice.

Let

\[
                  v_A=\int_X c_1(A)^d.                 \tag{VA4}
\]

This is a positive integer, not evaluated here. The exact total H-degree is

\[
                  \deg_H X=h_d^d v_A.                 \tag{VA5}
\]

Consequently the same-rule degree bound of LIFT096 becomes explicit in h:

\[
 D_* =2^{8g^2+2g+4}\,h_d^d v_A\sum_{j=0}^{d}k^j,
 \qquad
 k=gw+n/2\quad(n\ {
m even}).                         \tag{VA6}
\]

For arbitrary n the same-zero-locus square construction uses k=2gw+n.~Here w is the actual sum of frame weights of the same F0, not a new rule or a replacement arithmetic derivative. When k=0 and the same-family bad locus is proper it is empty, as already proved in LIFT096; no escape is required within that component.

\subsection{2. Normality, boundary codimension and the actual Cartier canonical divisor}\label{normality-boundary-codimension-and-the-actual-cartier-canonical-divisor}

The group Gamma is neat by the integral 3M congruence argument in LIFT094 HE17--18. Hence Y is smooth. The connected domain is irreducible as a complex manifold, so its quotient and its normal projective Baily--Borel compactification X are irreducible. These are the connected-component versions of Baily--Borel algebraicity, as recorded in Klingler--Yafaev (2014, Proposition 5.3.2 and Remark 5.3.4). The construction is not performed on a possibly disconnected union with unstated component degrees.

The Q-group Res\_\{F/Q\}SU(V) has absolute factors of type A\_(r-1). Since r\textgreater=3 it has no simple quotient of type PGL\_2 over Q. Milne's boundary statement (2017, Theorem 3.12 and Remark 3.13(b)) therefore gives

\[
                         \operatorname{codim}_X(X\setminus Y)\ge2.
                                                               \tag{VA7}
\]

If the arithmetic quotient is already compact, the boundary is empty and the same reasoning applies. No assertion for a rank-two modular-curve factor is imported.

Let L\_Y be the ample line on X in Klingler--Yafaev Proposition 5.3.2(1), with its arithmetic-lattice extension in Remark 5.3.4. Its restriction is the line of top-degree forms K\_Y. All codimension-one points of X lie in Y by VA7. On a normal variety a rank-one reflexive sheaf is determined by its restriction to an open containing all codimension-one points. Thus

\[
                   \mathcal O_X(K_X)\simeq L_Y.        \tag{VA8}
\]

In particular K\_X is a Cartier divisor and is ample. This identifies the actual canonical divisor, not merely a positive rational power or a numerically equivalent class. The statement uses normality and reflexive extension; it does not require, or infer, that X is Cohen--Macaulay or Gorenstein in every codimension.

The normal pair (X,0) is log canonical. Indeed, Alexeev (1996, Section 3, Theorem 3.2 and Corollary 3.3, printed pp.~9--10) gives log canonicity of the Baily--Borel pair with its divisorial boundary for a neat arithmetic quotient. That boundary divisor is zero by VA7. This is the unconditional locally symmetric-space statement in Section 3, not a use of the conditional moduli discussion in Section 2. Its proof is the toroidal/log-canonical model construction from Mumford; no conjectural minimal model program is needed for this application. Normality also makes (X,0) a semi-log canonical pair: the normalization is X itself, its conductor is empty, and the requisite S2 and codimension-one regularity hold.

\subsection{3. Positivity on every log-canonical center}\label{positivity-on-every-log-canonical-center}

The trace-symplectic representation gives the holomorphic totally geodesic embedding of domains checked in LIFT096 Section 6. Its image preserves the integral symplectic lattice and has level 3M, so the exact two neat arithmetic groups satisfy the hypotheses of Klingler--Yafaev Proposition 5.3.6. That proposition yields nefness on the whole compactification, including curves contained in the boundary, of

\[
                         Q=A-K_X.                     \tag{VA9}
\]

We use additive notation for Cartier divisors in VA9 and below. The underlying line statement is A tensor L\_Y\^{}\{-1\} nef. Since K\_X is ample and Q is nef, A=K\_X+Q is ample. Thus

\[
                         2A-K_X=A+Q                  \tag{VA10}
\]

is ample as well. The sum of an ample Cartier class and a nef class is ample on a projective variety.

An ample divisor restricts to an ample divisor on every positive-dimensional closed subvariety, hence is big there. On a zero-dimensional stratum its restriction is big in the usual dimension-zero convention. It follows that VA10 is nef and log big with respect to (X,0): it is big on X and on \textbf{every} log-canonical center, with no assumption that the centers have a particular boundary shape. This is stronger than the nef-and-big condition on X alone and is exactly the condition required by the effective theorem used next.

\subsection{4. Effective freeness followed by corrected very ampleness}\label{effective-freeness-followed-by-corrected-very-ampleness}

Fujino (2009, Theorem 1.1 and Remark 1.2, p.~475) applies to the projective log-canonical pair (X,0), the nef Cartier divisor A, and a=2, using VA10. It gives the following explicit basepoint-free multiple:

\[
                m_d A,\qquad
                m_d=2^{d+1}(d+1)!(d+2).               \tag{VA11}
\]

The theorem's formula is m(d,a)=2\^{}(d+1)(d+1)!(ceil(a)+d). There is no unrecorded rounding or dimension shift in VA11.

Put D=m\_d A. Then D is an ample Cartier divisor with basepoint-free complete linear system. Moreover

\[
              D-K_X=(m_d-1)A+Q                       \tag{VA12}
\]

is ample because m\_d\textgreater1. As in Section 3, it is nef and big on every semi-log canonical stratum of (X,0), including X. All hypotheses of Fujino (2017, Lemma 7.1, printed pp.~367--368) therefore hold. Although that article's title concerns surfaces, Section 7 is explicitly independent of its surface results, and Lemma 7.1 is stated and proved for arbitrary dimension d.~It gives very ampleness of (d+1)D. Hence

\[
 (d+1)m_d
  =2^{d+1}(d+1)(d+1)!(d+2)
  =2^{d+1}(d+1)(d+2)!
  =h_d.                                                \tag{VA13}
\]

This proves VA2. The corrected published Lemma 7.1 is the imported very-ampleness input. We do not substitute the older unmodified Kollár lemma that Fujino explicitly discusses before this statement. The cited proof uses log-canonical vanishing, Castelnuovo--Mumford regularity and the section-ring multiplication maps to prove an actual closed immersion, not merely a finite morphism or separation of generic points.

For clarity, the first values for dimensions d=1,2,3,4 are respectively 48, 576, 7680 and 115200. The actual VA1 family has d\textgreater=2, but the abstract effective calculation also applies in dimension one when its geometric hypotheses hold. A useful conservative size bound is

\[
 h_d\le2^{d+1}(d+1)(d+2)^{d+2}.                        \tag{VA14}
\]

This is only a size estimate; the exact factorial expression VA13 is retained in every subsequent cost.

\subsection{5. Compatibility and the full degree cost}\label{compatibility-and-the-full-degree-cost}

The choice of h\_d changes only the auxiliary projective embedding. It leaves the original Host, Datum, coefficient record, abelian scheme, full OE action, physical trace-one theta line, finite eta, weights and rule F0 untouched. LIFT096's algebraic Siegel Kodaira--Spencer determinant identity gives A\textbar\_Y=lambda\^{}(g+1), hence VA3. The equal-pullback identity for A on both Hecke diagrams follows from the canonical automorphic line functoriality already checked there; taking the same integer tensor power preserves it exactly.

For an ample Cartier line on a projective integral d-fold, top intersection is the positive integer v\_A in VA4. Multilinearity of first-Chern-class intersection gives c1(A\^{}h)\^{}d=h\^{}d c1(A)\^{}d, proving VA5. Thus the large factor h\_d\^{}d is part of the bound and cannot be discarded when changing projective embeddings.

Insert VA5 into the unchanged LIFT096 bound. That argument constructs equations on a finite proof cover, raises their powers to match H, extends them across the codimension-two boundary, and pushes the degree down. It gives VA6 on the original family. The integer h\_d does not change the effective equation exponent k relative to H: the common power (g+1)h\_d changes lambda\^{}k into H\^{}k. The theta cover's degree remains a factor of D\_* only; it is not an output field degree or an isogeny multiplier.

With all the unchanged LIFT094/LIFT095 hypotheses, the resulting conditional within-family escape estimate is therefore

\[
 N_{\rm within}\le
 t^{B_{\delta,d-1}\!\left(2^{8g^2+2g+4}h_d^d v_A
                          \sum_{j=0}^{d}k^j\right)},
 \quad t=1+3M\prod_{\ell\in S}\ell,
 \quad\delta=t^{4fr^2}.                               \tag{VA15}
\]

Any required component-entry multiplier N0 is separate and multiplies this bound. The parameter t must still satisfy every previously imposed congruence and protected-prime condition. If a numerical upper bound v\_A\textless=V0 is subsequently proved, an integer ceiling of V0 can be substituted into VA15 by monotonicity. No value for V0 is asserted here.

\begin{longtable}[]{@{}
  >{\raggedright\arraybackslash}p{(\columnwidth - 4\tabcolsep) * \real{0.3333}}
  >{\raggedright\arraybackslash}p{(\columnwidth - 4\tabcolsep) * \real{0.3333}}
  >{\raggedright\arraybackslash}p{(\columnwidth - 4\tabcolsep) * \real{0.3333}}@{}}
\toprule\noalign{}
\begin{minipage}[b]{\linewidth}\raggedright
Interface
\end{minipage} & \begin{minipage}[b]{\linewidth}\raggedright
Supplied here
\end{minipage} & \begin{minipage}[b]{\linewidth}\raggedright
Still required for the actual application
\end{minipage} \\
\midrule\noalign{}
\endhead
\bottomrule\noalign{}
\endlastfoot
Same component X and auxiliary line A & Normal projective X, Cartier ample K\_X, and A-K\_X nef & The admitted family and literal source identification \\
Very ample embedding & H=A\^{}h\_d with exact h\_d=2\^{}(d+1)(d+1)(d+2)! & No additional unknown exponent h \\
Ambient intersection degree & deg\_H(X)=h\_d\^{}d v\_A & A numerical bound on the same v\_A, with field, rank and level factors \\
Same-rule bad locus & VA6 under LIFT096's exact same-family nonempty condition & Proof that this component meets U\_F0 \\
Entry and arithmetic costs & The conditional expression VA15 and its exact change of line & N0, actual Host/Sizes window and selected source construction \\
\end{longtable}

This is an explicit reduction of the auxiliary very-ampleness input. It does not turn an analytic period map into a bounded-degree polynomial, prove a central L-value nonzero, or identify the arithmetic selected rule with another section.

\subsection{6. Primary sources, literal interfaces and verification scope}\label{primary-sources-literal-interfaces-and-verification-scope}

The exact inherited source family and same-rule map are those of frozen LIFT096 TJ1--4, TJ15 and TJ17--18, and canonical LIFT094 HE0--3, HE17--22. The source Coeff and Moduli clauses still have their original field, smoothness, quasi-projectivity and representability scope; no strengthening is made. The accompanying source manifest records their precise current files and fingerprints, with targeted versus whole fresh reading distinguished.

Primary passages freshly read for this proof:

\begin{enumerate}
\def\labelenumi{\arabic{enumi}.}
\tightlist
\item
  Alexeev, V. (1996). \emph{Log canonical singularities and complete moduli of stable pairs} (preliminary version, August 15; arXiv:alg-geom/9608013v1). \url{https://arxiv.org/pdf/alg-geom/9608013}. Section 1 definitions, and Section 3, Theorem 3.2/Corollary 3.3 with their full neat arithmetic-quotient discussion, printed pp.~9--10. Only that unconditional Section 3 application is used.
\item
  Fujino, O. (2009). Effective base point free theorem for log canonical pairs---Kollár type theorem. \emph{Tohoku Mathematical Journal, 61}, 475--481. \url{https://doi.org/10.2748/tmj/1264084495}. Final author PDF \url{https://www.math.kyoto-u.ac.jp/~fujino/effective-final.pdf}. Theorem 1.1/Remark 1.2, p.~475; definition of log bigness, p.~476; the effective proof in Section 2.2. The coefficient in VA11 uses this published formula.
\item
  Fujino, O. (2017). Effective basepoint-free theorem for semi-log canonical surfaces. \emph{Publications of the Research Institute for Mathematical Sciences, 53}, 349--370. \url{https://doi.org/10.4171/PRIMS/53-3-1}. Published PDF \url{https://ems.press/content/serial-article-files/41296?nt=1}. Section 7, Lemma 7.1 and its complete proof, pp.~367--368; the independent, all-dimensional scope is stated at the start of Section 7 and in the introduction. Remark 7.2, pp.~368--369, corroborates the factorial scale, but the argument above uses the lemma directly.
\item
  Klingler, B., \& Yafaev, A. (2014). The André--Oort conjecture. \emph{Annals of Mathematics, 180}(3), 867--925. \url{https://www2.mathematik.hu-berlin.de/~klingleb/papiers/KY1309.pdf}. Proposition 5.3.2(1)/Remark 5.3.4, author printed p.~24; Proposition 5.3.6 and its arithmetic-lattice reduction, pp.~25--26. Only the unconditional compactification and line comparison are imported.
\item
  Milne, J. S. (2017 revision). \emph{Introduction to Shimura varieties}. \url{https://www.jmilne.org/math/xnotes/svi.pdf}. Theorem 3.12/Remark 3.13(b), printed pp.~38--39, as directly read for LIFT096, give the normal Baily--Borel setting and the no-PGL2 boundary-codimension condition.
\end{enumerate}

The companion executable checks the exact two-stage factorial constant, positivity identities as divisor polynomials, all-dimensional degree scaling, the unchanged theta/Hodge weights and same-rule substitutions, protected-prime and level exponents, and every recorded subject/dependency fingerprint. These are arithmetic and interface diagnostics, not a computational proof of log canonicity, the effective freeness theorem or very ampleness. The universal arguments and imported theorems above are the proof. Independent mathematical/source reviews and whole saved-file author verification have separate actual timestamps and hashes.

The present result remains in the algebraic compactification and coherent-sheaf part of the inherited Grothendieck formulation: canonical reflexive extension, Cartier positivity, log-canonical vanishing, algebraic section-ring multiplication and Chern-class intersection. It adds no motive, weight, Frobenius, CM-type admissibility or Riemann-hypothesis implication. The next quantitative input is the actual v\_A, not another frame-polynomial degree.

\subsection{Canonical algebraic-geometric context}\label{canonical-algebraic-geometric-context-60}

This is a coherent-sheaf and algebraic compactification result for the same original family: rank-one reflexive canonical extension over a normal scheme, actual Cartier positivity, log-canonical vanishing, algebraic section-ring generation and top first-Chern-class intersection. Exact fresh primary inputs are Alexeev(1996), Section3 Theorem3.2/Corollary3.3 pp9-10, independent of the conditional general-moduli discussion; Fujino(2009), Theorem1.1/Remark1.2 p475 and log-bigness definition p476; Fujino(2017), independent all-dimensional Section7 Lemma7.1 and full proof pp367-368; Klingler-Yafaev(2014), Proposition5.3.2(1), Remark5.3.4, Proposition5.3.6 authorpp24-26; Milne(2017), Theorem3.12/Remark3.13 pp38-39. The standard ample-plus-nef implication is carried with Lazarsfeld(2004), Corollary1.4.10 pp53-54. The prescribed theta line, OE action, Picard data and Poincare/cubical-line transports are unchanged. Carried contextual anchors: FGA232 Theorem3.1 pp148-149; FGA236 Corollaries1.7/2.3 pp229/232; SGA7 I VII2.7/2.9 pp169-174 and VIII7.4(b) pp310-311; SGA3 I4.7.3/4.7.3.1; SGA6 VI Theorem3.3; cote4 scan003 sections8.1.1-8.1.2 and scans005-006 sections8.4.1-8.4.3, cote54 scans005-007. No new lambda-ring/K-theory equality, l-adic sheaf/weight, Dieudonne-crystal or CM-motive comparison follows. DP\_MULT Q3e/DP44/DP52/DP54-65 and THETA3 G1/A7 receive a conditional cost bound with one fewer geometric input; CX\_GRH CX-ALG-024 and GCT\_PLETH T-052 retain their arithmetic/global gates.

\clearpage

\section{AM-R-096 --- resolution of the fixed sufficient-envelope test}\label{am-r-096-resolution-of-the-fixed-sufficient-envelope-test}

Recorded 2026-11-23T08:57:46.719641+00:00. \textbf{The proposed certification by the unchanged DO111 VC12--22 envelope at support27 fails.} This is a statement about that fixed sufficient estimate, not a lower bound on the actual sector or a refutation of its smallness. The original08:10:50 probe and08:17:20 exploratory note are preserved; the standalone reproduction was completed later at08:41:35.513194 UTC.

The exact actual inputs and directed logarithms are those of DO111 VC1--3 and VC13--22, with f=168,H=133 and s=27. The unchanged common-ideal cap gives z\textless2\^{}-90. The complete original row radius is8392725/2\^{}90, and the least strict binary ceiling of the rational row-volume envelope is−4071532. The entire old z exponent is819. With L\_T=1019333/1024,L\_Q=287843/4096,L\_U=601356/4096,L\_H=28900/4096, its complete exponent is

\[
\begin{split}
E={}&6958-929f/2+41(167)f+(f/8)(L_H-1)-(f/2)L_U\\
 &-501(167)L_Q+(344L_T+997)27-4071532-819(90)\\
 ={}&1289132235/4096>0.
\end{split}
\]

Every equality and strict directed comparison is reproduced by \nolinkurl{recovery_pending/root/r096_existing_envelope_reproducer.py}, whose saved08:41:35 output is \nolinkurl{r096_existing_envelope_replay.json}. Its integer numerator/denominator comparisons check the exact row-volume ceiling; it retains the full original target/source fibre, selected trace, paired finite weights and all positions. Thus this displayed upper bound is larger than1 and cannot certify the required subunit upper bound. This proves the negative answer to the fixed-estimate question. It proves no existence, sign, contribution lower bound or impossibility of a sharper count.

The research sector is separately resolved by the new and whole-reviewed AM-DO-120 QW1--29, canonically integrated at08:54:03.101648 UTC. Its positive O\_F witness, new finite rank-four restoration and complete original count prove a different bound, below2\^{}-270871 for all nonzero-short-compression sectors at supports27--31. The new theorem does not change the value or historical failure of the old formula. R096 therefore has no remaining mathematical subtask and no new error charge; the actual current lists are DO120127/31, complete floors27/11. The failed sufficient estimate is not credited as additional review movement beyond DO120's substantive gate contraction.

This resolution changes no original Hermitian module, theta line, Picard/biextension datum, projector or normalization. It has no separate lambda/K-theory, l-adic, crystalline or motivic consequence. Exact carried frameworks and global application gates are those recorded in DO111 and DO120; no new Grothendieck source is invoked. The actual signed, source-admissibility, selected arithmetic content, prescribed period, specific032, DP44 and RH gates remain open.

\clearpage

\section{AM-R-097 --- the complete support27 necessary-window assembly}\label{am-r-097-the-complete-support27-necessary-window-assembly}

\textbf{Canonical checked integration --- 2026-11-23T09:16:12.287665+00:00.} Complete saved-proof and executable reviews and separate exact diagnostic replays passed. No numerical error charge is introduced. Frozen author history is preserved. AM-R-097 NW1-9 exhaustively composes the checked structural cuts at support27. Every remaining actual support27 arrow is full-positive, scalar-cross, globally rank-two-unipotent with all27 active places typeII, b=j0 and one of two necessary alternatives: total eligibility27,100\textless=k\textless=164,2\^{}2111\textless x\textless2\^{}3328; or total eligibility28,k164,2\^{}3107\textless x\textless2\^{}3328. DO120 supplies full positive compression before K114 is applied. The signed energy retains the actual R27 subset of this containing window and R\_ge28, with \textbar epsilon\textbar\textless2\^{}-230. No set equality, existence, emptiness, sign or floor28 conclusion. No new charge; actual DO120127/31 ordered arrays and complete floors27/11 are unchanged. Actual mixed-CM source, selected arithmetic content, six identities, prescribed period, specific032 pairing, evaluatedDP44, Kummer-Hecke RH and global oscillator comparison remain open.

Author: \nolinkurl{/root/definite_oscillator}, under the parent's R097 reservation. 23 November 2026 UTC. Pending whole-proof composition; canonical promotion and independent review are separate events. This file makes no canonical edit and adds no numerical charge.

\subsection{1. The unchanged actual kernel and assertion}\label{the-unchanged-actual-kernel-and-assertion}

Retain the actual conditional J56 all-top kernel at its already charged cutoff H=133, f=a=168 and n=169. Keep the original E/F, integral Hermitian modules, both CM partners, admitted source and target pencils and omitted labels, finite projectors, determinant-one central representative, selected real representation and Hilbert metric, packet point, compact factors, phases, Haar measures and class masses. All subsequent containing-set removals are applied to their intersections with the actual remainder after earlier removals. Every previously charged containing estimate remains in its chronological error list.

R094 PP1--24 gives the complete support floor27. The later canonical removals, in their actual order, are DO117, DO118, K113, R095, DO119, K114 and DO120. Let R be the actual remainder after all of them. Its error arrays have127 all-top entries and31 uniform entries. This assembly leaves those arrays unchanged.

\textbf{Theorem.} Every term of R at support27 belongs to the following necessary containing class W27. In the original rational orthogonal decomposition

\[
 E^{169}=Z\perp P,\qquad Z=S_{164}\perp u^\perp,
 \qquad P=Eu\perp Ee_n,\qquad J_P=\operatorname{diag}(U,-1),
 \tag{NW1}
\]

it has determinant one, all334 unmixed G reductions unipotent, nonzero original S-cross, and

\[
 h=I+N,\quad \operatorname{rank}_E N=2,\quad N^3=0,
 \quad\operatorname{rank}_E N^2=1,\quad N^2Z=0,
 \quad h_{ZZ}=(h^{-1})_{ZZ}=I_Z.
 \tag{NW2}
\]

Every one of the27 active relative T places is of type II and is eligible. Let eta be the original packet-plane slope, eta bar(eta)=U; let D\_eta be its exact denominator ideal and x=N\_E D\_eta; let r\_elig be total eligibility among all168 selected T places; let k count the nonzero original S covector values; let b count inactive places where both paired C blocks vanish; and let j count active places where the whole paired C is integral. Then

\[
 b=j=0,\qquad 2^{2111}<x<2^{3328},\qquad
 \begin{cases}
 r_{\rm elig}=27,&100\le k\le164;\\
 r_{\rm elig}=28,&k=164,\quad x>2^{3107}.
 \end{cases}
 \tag{NW3}
\]

These are exactly the two K113 SC27.2 windows. The second line retains the common upper bound x\textless2\^{}3328 and its stronger lower bound x\textgreater2\^{}3107. No equality endpoint is added. Total eligibility is distinct from the27 active places and from any matrix rank.

Let R27 be the actual set R intersect support27, so R27 is a subset of W27, and let R\_ge28 retain every original higher-support restriction already proved. The exact remaining energy has the form

\[
 \mathcal E_{\rm top}
  =1+\operatorname{Re}\!\left(\sum_{h\in R_{27}}w(h)
                    +\sum_{h\in R_{\ge28}}w(h)\right)+\epsilon,
 \qquad |\epsilon|<2^{-230}.
 \tag{NW4}
\]

The same sufficient signed inequality, with the two actual sums in parentheses greater than -1+2\^{}-230, remains open. A containing class is not substituted for the actual signed set. No absolute estimate, emptiness or sign for either window NW3 is proved here. The complete floors remain27/11, not28/11.

\subsection{2. All nonidentity original positive compressions are already excluded at27}\label{all-nonidentity-original-positive-compressions-are-already-excluded-at27}

Put A\_S=h\_SS and r\_S=rank\_E(A\_S-I). K068 LR3/LR20--22 gives the actual determinant-one representative, global integrality of A\_S, finite unipotence at all334 unmixed G places, and r\_S\textless=33. These are all-support statements proved at a larger cutoff; restriction to the actual H133 set preserves them. Finite unipotence is not global unipotence.

For r\_S\textgreater=1, K069 IC1--8 is exhaustive: either h fixes K\_A=ker(A\_S-I) intersect ker((A\_S-I)\^{}*) pointwise (the original class F), or it is in the precise exceptional isotropic-cross class IC8. These original classes are not replaced by arbitrary moving-space conditions. The following complete inclusive ranges contain27:

\begin{longtable}[]{@{}
  >{\raggedright\arraybackslash}p{(\columnwidth - 6\tabcolsep) * \real{0.2308}}
  >{\raggedleft\arraybackslash}p{(\columnwidth - 6\tabcolsep) * \real{0.3077}}
  >{\raggedright\arraybackslash}p{(\columnwidth - 6\tabcolsep) * \real{0.2308}}
  >{\raggedright\arraybackslash}p{(\columnwidth - 6\tabcolsep) * \real{0.2308}}@{}}
\toprule\noalign{}
\begin{minipage}[b]{\linewidth}\raggedright
Original class
\end{minipage} & \begin{minipage}[b]{\linewidth}\raggedleft
Compressed ranks
\end{minipage} & \begin{minipage}[b]{\linewidth}\raggedright
Complete range used at27
\end{minipage} & \begin{minipage}[b]{\linewidth}\raggedright
Exact locator
\end{minipage} \\
\midrule\noalign{}
\endhead
\bottomrule\noalign{}
\endlastfoot
F & 1 & 18 through94 & DO086 FR1--10 \\
F & 2 & through110 & R068 FA1--9 \\
F & 3 & through154 & R068 FA1--9 \\
F & 4 through33 & through168 & R068 FA1--9 \\
IC8 & 1 or2 & 23 through28 & R069 EM1--21 \\
IC8 & 3 through31 & complete contiguous ceilings, every one at least28 & R075 NR1--16, with its DO085/DO064 predecessor ranges \\
IC8 & 32 or33 & 5 through168, with earlier small supports already removed & DO064 PM22--23 \\
\end{longtable}

For clarity the R075 ceiling sequence for r\_S=3,\ldots,31 is

\[
 (28,34,40,45,51,56,62,67,73,78,84,89,94,99,105,
 110,115,120,125,130,135,140,145,150,155,160,164,168,168).
 \tag{NW5}
\]

R075 removes all fresh cells above each previous ceiling, not merely the displayed final endpoint. Its union with those predecessors is therefore contiguous. In particular the minimal rank3 ceiling28 covers27. Rank31 is already covered by its predecessor through168, even though its fresh R075 interval is empty. The two classes cover every possible r\_S=1,\ldots,33; no global characteristic or type restriction is inserted. It follows that every support27 term of R has

\[
                              h_{SS}=I_S.
 \tag{NW6}
\]

No new multiple of an old error charge is incurred by applying these whole-sector bounds at a single support.

\subsection{3. The complete fixed-S alternative}\label{the-complete-fixed-s-alternative}

Write h in the original S perpendicular W0 blocks, and C\_S=h\_W0,S. If C\_S=0, NW6 says that h fixes every vector of S. Unitarity preserves the orthogonal complement of the nondegenerate fixed positive space, hence h=I\_S perpendicular h0 with h0 in U(4,1) and determinant one. Conversely that fixed-S form has C\_S=0. This also rules out a separate upper-cross-only class: the upper cross vanishes by unitarity when the lower cross vanishes.

If h0 is not globally unipotent, DO117 FC1--25 removes the entire original fixed-S support27 sector, at every global rank, characteristic degree and incidence profile. It retains the full paired I\_Sigma denominator for its small scalar; this application does not import a type-III-only pole improvement.

If h0 is globally unipotent and not the identity, the actual signature-one classification DO060 IB21--22 gives precisely rank(h-I)=1 with N\^{}2=0, or rank(h-I)=2 with N\^{}3=0 and rank(N\^{}2)=1. R095 FW1--21 removes the complete first class at27. DO118 UB1--29 removes the complete second class at27, including its high finite-rank sector and every low-rank short or last-only pivot and original type/label assignment. Rank zero means h=I and has support zero, so is incompatible with support27.

Thus no fixed-S term remains. Every support27 term of R satisfies C\_S!=0. This is a global algebraic dichotomy; it does not use a positive form on a finite residue field.

\subsection{4. From the nonzero cross to the actual three-dimensional action}\label{from-the-nonzero-cross-to-the-actual-three-dimensional-action}

DO099 JR1--24 removes the entire original and inverse transverse S-cross containing sectors at every support25 through70. The absolute inverse transport in that theorem uses the original swapped source/target pair and selected normalization. Its containing-set conclusion applies to the present remainder without requiring that the remainder itself be inversion invariant. At support27 both remaining cross directions therefore lie in the same fixed packet plane P of NW1; their individual null image lines need not coincide.

The support-independent unitary block identity DO111 VC5--8 then gives

\[
 h=\begin{pmatrix}K&B\\ C&D\end{pmatrix}_{Z\oplus P},
 \quad K=\operatorname{diag}(I_S,A_V),\quad A_V^*G_VA_V=G_V,
 \quad(h^{-1})_{ZZ}=K^{-1},\quad\operatorname{rank}B,C\le1.
 \tag{NW7}
\]

Here C restricted to S is nonzero. This identity alone does not force A\_V=I, finite unipotence of A\_V, or coincidence of the two crossing lines. The exact rank-four paired local counterexample in DO120 QW19 explicitly prevents that shortcut.

Let Delta\_ij=\textless y\_i,(h-I)y\_j\textgreater{} for the original two short integral positive inputs y1,y2. DO120 QW1--29 removes the complete doubly-plane class with this four-entry compression nonzero at every support27 through31. Its proof retains the full original column and source fibre and all types; no further hypothesis is needed at27. Consequently every remaining Delta\_ij is zero.

The support-independent structural consequence QW29 now applies. Unitarity of A\_V and the two zero diagonal energy identities make A\_V fix span(y1,y2). The original integral third-vector compression and exact R071 scalar norm inequality force its final unitary scalar to equal one. Thus K=I\_Z and

\[
 h_{ZZ}=(h^{-1})_{ZZ}=I_Z,\qquad
 \dim(\ker(h-I)\cap Z)=166,\qquad\operatorname{rank}_E(h-I)\le3.
 \tag{NW8}
\]

Only this point establishes full positive compression at27. The earlier support26 numerical DO111 removal is not used at27. The y\_i need not individually be fixed by h, since their P components may remain nonzero.

There is no global rank-zero or rank-one scalar-cross escape. Indeed unitarity with upper-left I\_Z gives C\^{}\emph{J\_PC=0, so C has rank one since C\_S!=0. Jacobi's identity, determinant one and the inverse positive compression give detD=1. The off-diagonal identity gives B=-J\_Z\^{}\{-1\}C\^{}}J\_PD, so B!=0. If N=h-I had rank one, writing N=ab\^{}* and using its zero Z-by-Z block would force either a\_Z=0 or b\_Z=0, hence respectively B=0 or C=0. Both contradict the established nonzero crosses. Rank zero contradicts C\_S!=0 directly. Therefore

\[
                         \operatorname{rank}_E(h-I)\in\{2,3\}.
 \tag{NW9}
\]

The full DO095 RA1--10 algebra applies with these precise premises, despite the historical support23/24 introduction of that proof: its equations are support-independent. In particular BC nonzero is exactly global rank three; BC zero is exactly global rank two with the real reciprocal eigenvalue described in RA10. This imports no global nilpotence.

\subsection{5. Every remaining characteristic/type alternative except all-II is removed}\label{every-remaining-characteristictype-alternative-except-all-ii-is-removed}

If the global rank is three, NW8 and the original nonzero S-cross give every premise FP27.1 of K114. Its whole theorem FP27.1--29 removes the full positive-compression rank-three support27 sector with no coordinate-sparsity or denominator assumption. Its type-III range11 through27 follows from the support-independent DO074 TC1, so no t\textless=10 class is omitted. The full original projected-column and source fibre are already inside K114's bound. K114 is applied after NW8, never to the uncut N1 sector.

At global rank two, global unipotence and nonunipotence are exhaustive alternatives. DO101 PI1--17 removes the entire nonunipotent scalar-cross plane-C interval25 through32. Its support-independent reciprocal-line argument requires only the original S identity, nonzero C\_S, determinant one, rank two and nonunipotence. These hypotheses are weaker than NW8. Its actual trace-pole threshold is not applied to a zero trace witness or to a unipotent matrix.

Hence the global rank-two remainder is unipotent. DO060 IB21--22 gives N\^{}3=0 and rankN\^{}2=1, and the support-independent normal-form argument DO119 NA1--2 gives N\^{}2Z=0. For i,u,t the actual active type counts, i+u+t=27. DO119 NA1--29 removes every such sector with u\textless=26, retaining separate source/target labels, all original paired weights, row and central fibres and the infinite denominator tail. Its all-active eligibility lemma is valid for all types, but its numerical theorem is used only for the non-all-II class.

The only remaining possibility is u=27 and i=t=0. All active places are type II. K113 SC27.1--27 applies exactly to this globally rank-two-unipotent scalar-cross class, with nonzero original S row and NW2. Its complete seven-group reduction yields precisely NW3. In particular b=j=0 is a necessary result of that charged theorem, not an assumed absence of paired zero/integral positions. K112's local active-II result supplies eligibility without any support26 numerical transplant.

This proves the claimed exhaustive necessary-window assertion. Every support27 term has passed one of the two K069 classes, the fixed/cross split, the two transverse cuts, the short-compression split, the rank-two/three split, global unipotence, and the original type partition. No existence claim for an arrow in either final window follows.

\subsection{6. Exact chronology, no new charge, and verification scope}\label{exact-chronology-no-new-charge-and-verification-scope}

The complete chronological all-top prefix through R094 has120 entries. The successive additions are

\begin{longtable}[]{@{}lrr@{}}
\toprule\noalign{}
Canonical finding & Single appended exponent & All-top length \\
\midrule\noalign{}
\endhead
\bottomrule\noalign{}
\endlastfoot
DO117 & 29365 & 121 \\
DO118 & 3927 & 122 \\
K113 & 304 & 123 \\
R095 & 13000 & 124 \\
DO119 & 412 & 125 \\
K114 & 1650 & 126 \\
DO120 & 270871 & 127 \\
\end{longtable}

All31 uniform entries remain unchanged. The companion reads every original canonical predecessor in this chain and checks exact prefix equality, each appended exponent, all complete floors and both full dyadic sums below2\^{}-230. This assembly appends nothing; in particular it does not recharge the interval DO120 or any support27 specialization of an older containing theorem. The unique231 minimum and next233 are retained.

R094 already excludes supports below27. Sections2--5 cover all support27 structural alternatives except the exact necessary windows NW3. Intersecting each old containing removal with the then-current remainder makes the actual removed sets disjoint while retaining every earlier bound. The unchanged energy identity and triangle inequality prove NW4. All earlier restrictions at supports28 and above remain in their original scope, including DO120's27--31 structural consequence where its two-plane premise holds.

The diagnostic certificate verifies source fingerprints and the displayed inclusive tables, the entire chronological arrays, finite structural/type partitions and the strict numerical boundary conventions. These checks do not enumerate arithmetic arrows, reprove the imported source theorems, or prove that an open window is populated. Whole saved-proof and source-interface readings are separate from the diagnostic replay. The source manifest records exact locators and the fresh targeted reading scope; earlier whole checks of those findings remain their own historical events.

\subsection{7. Retained arithmetic and geometric meaning}\label{retained-arithmetic-and-geometric-meaning}

The original integral Hermitian modules, CM involution, paired fractional ideals, finite incidence correspondences, theta line, Picard/Poincare data, polarization, biextensions and actual selected metric are unchanged. The only new conclusion is their complete support27 taxonomy. No moving rational frame is assigned global integrality; no finite field is given a positive metric; no central orbit is substituted for a truncated actual sum.

Carried framework anchors, at the scopes already recorded in R094 and the cited findings, are SGA3 I4.7.3/4.7.3.1; SGA6 VI Theorem3.3; FGA232 Theorem3.1 pp148--149; FGA236 Corollaries1.7/2.3 pp229/232; SGA7 I VII2.7/2.9 pp169--174 and VIII7.4(b) pp310--311; cote4 scan003 sections8.1.1--8.1.2 and scans005--006 sections8.4.1--8.4.3; and cote54 scans005--007. No fresh whole-source reading, Dieudonne slope, lambda/K-theory identity, l-adic weight or CM-motive comparison is asserted here.

The signed finite/real inequality in NW4, admitted full Host/Sizes source, same selected arithmetic nonzero content, prescribed global Bessel/Mram period, Read/specific032/STiesR, all six lift identities, evaluated DP44, Kummer-Hecke RH and the global oscillator comparison remain open at their established scopes. This result supplies neither an unconditional period nor a shorter admitted arithmetic chain.

\subsection{Canonical algebraic-geometric context}\label{canonical-algebraic-geometric-context-61}

This assembly concerns the same integral projective OE modules, paired CM ideals, polarization and finite correspondences over Spec(OE)/Spec(OF), with the same selected Hilbert realization. Positive compression is proved over the original global field; it is not transferred to residue-field positivity. Original Picard/theta/Poincare data and actual signed coefficients are retained. Carried exact anchors: SGA3 I4.7.3/4.7.3.1; SGA6 VI Theorem3.3; FGA232 Theorem3.1 pp148-149; FGA236 Corollaries1.7/2.3 pp229/232; SGA7 I VII2.7/2.9 pp169-174 and VIII7.4(b) pp310-311; cote4 scan003 sections8.1.1-8.1.2 and scans005-006 sections8.4.1-8.4.3; cote54 scans005-007. These are unchanged contextual interfaces, not fresh primary-source readings or new formal-group/Dieudonne, lambda-ring/K-theory, l-adic-weight or CM-motive comparison claims. DP\_MULT Q3e/DP44 and THETA3 G1/A7 retain the actual-family gates; CX\_GRH CX-ALG-024 and GCT\_PLETH T-052 remain open at their recorded global gates.

\clearpage

\section{AM-K-115 --- shifted Gamma moments in the actual all-II support27 row count}\label{am-k-115-shifted-gamma-moments-in-the-actual-all-ii-support27-row-count}

\textbf{Canonical checked integration --- 2026-11-23T09:17:53.346404+00:00.} Complete root and independent proof/script/interface reviews and separate exact replays passed. The actual canonical predecessor has127/31 entries; this integration gives128/31, complete floors27/11. Frozen author history is preserved. AM-K-115 GM1-23 proves the cap-dependent shape83/rate7/2 positive Gamma majorant, retains its entire low-weight interval by a paid positive shift, and controls every inverse moment0..k+3. All scalar and transverse terms, original saturated metric, central affine coset, source reconstruction,168 T factors and four-state finite patterns remain. The complete ten cells with total eligibility27,b=j0,2\^{}2111\textless x\textless2\^{}3328,k100..109 have total absolute contribution below2\^{}-945. The worst full exponent is-2907383/3072\textless-946, every other is below-3166. One945 append changes DO120127/31 to128/31, both entire sums remain below2\^{}-230 and complete floors27/11 stay. Combined with R097, the remaining actual support27 set is contained in the full-positive scalar-cross rank-two-unipotent all-II class with either total eligibility27,110\textless=k\textless=164,2\^{}2111\textless x\textless2\^{}3328, or eligibility28,k164,2\^{}3107\textless x\textless2\^{}3328. The next positive sufficient envelope is not an actual lower bound. No floor28, signed nonvanishing, admitted source, selected arithmetic content, prescribed period, specific032 value, evaluatedDP44 or RH conclusion.

Author: \nolinkurl{/root/kummer_extension}, 23 November2026 UTC. Controlled pending proof. Exact arithmetic, complete saved-proof/script/source author checks and independent review are separate events. No canonical file is changed here.

\subsection{1. Actual remaining sector and theorem}\label{actual-remaining-sector-and-theorem}

Retain the original J56 all-top H133 normalized kernel, field E/F of degrees336/168, CM involution, all paired source/target pencils and admitted labels, finite projectors, actual vector-valued selected Hilbert metric and both compact factors, near-point representative, class masses, central phases, Haar measures, words and canonical twists. Work only in the actual globally rank-two-unipotent scalar-cross all-II support27 branch remaining after canonical K113. Write r for total eligibility, k for the number of nonzero original S covector values, b for inactive paired-C-zero places, and j for active integral-C places. In the present sector \[
 r=27,\quad b=j=0,\quad100\le k\le164,\quad
 2^{2111}<x=N_E\mathfrak D_\eta<2^{3328}.
 \tag{GM1}
\] The separate K113 alternative r=28,k=164 is not covered here.

\textbf{Theorem.} The complete containing union of the ten coordinate cells100\textless=k\textless=109 inside GM1 has absolute normalized contribution less than2\^{}-945. Append945 once to the actual canonical error list, retaining every earlier overlapping charge. Consequently the r=27 alternative may be restricted to110\textless=k\textless=164. This theorem gives no r=28 reduction, no complete support28 floor and no signed nonvanishing conclusion.

The physical constants and actual energy are \[
 f=168,\quad H=133,\quad R=17424/133<132,\quad
 Q=1427908193729121522601,\quad U=77Q^2-14,
\] \[
 D_0=7^{140}29^{144}5^{84}13^{84},\quad D_F=2^{168}D_0,
 \quad\nu=333/2,\quad m=83.
 \tag{GM2}
\] K113 SC27.1 gives the original full-positive rank-two normal form, with N\^{}2Z=0 and N\^{}3=0. Its actual selected estimate, directly carried from K108 AE13 and the full K107 near-point/tail argument, is \[
 |c(h)|<2\vartheta_a^f2^{-933f/2}(N_F\alpha)^{-1/2}
                    \prod_\tau\Psi(\alpha_\tau),
 \quad\Psi(t)=\left(\frac2{\sqrt{1+t}+1}\right)^\nu,
\] \[
 0<\alpha_\tau\le R,\quad\vartheta_a^f<2^{-649},\quad
 N_F\alpha\ge U^f(kz^2)^f.
 \tag{GM3}
\] Here L=J I\_Sigma D\_eta is the actual row ideal, N\_EL=z\^{}(2f), and the last inequality uses all k nonzero original S coordinates. GM3 is applied to the original rows before any positive Gaussian enlargement. Its energy cap and nonzero-coordinate condition are not asserted for the extra lattice tuples admitted below.

\subsection{2. A positive Gamma majorant on the entire actual energy interval}\label{a-positive-gamma-majorant-on-the-entire-actual-energy-interval}

We prove the explicit inequality \[
 \Psi(t)\le(1+2t/7)^{-83}\qquad(0\le t\le132).
 \tag{GM4}
\] Put v=sqrt(1+t) and \[
 F(v)=\frac{333}{2}\log\frac{v+1}{2}
                  -83\log\frac{2v^2+5}{7}.
\] Its derivative has the sign of P(v)=2v²-664v+1665. This polynomial is strictly decreasing on1\textless=v\textless=12, has P(1)\textgreater0 and P(12)\textless0, and thus F has at most one interior critical point, a maximum, on{[}1,sqrt133{]}. Its minimum on that interval is attained at an endpoint. At1, F=0. At sqrt133, use the exact integer comparisons \[
 1153^2<133\cdot10000,\qquad
                (1253/200)^{333}>(271/7)^{166}.
 \tag{GM5}
\] They give F(sqrt133)\textgreater0. This proves GM4, including equality at t=0. The assertion is cap-dependent; no extension of GM4 beyond132 is used.

Set beta=7/2. The probability measure \[
 d\gamma(b)=\frac{\beta^{83}}{82!}b^{82}e^{-\beta b}\,db,
                    \qquad b>0,
 \tag{GM6}
\] has Laplace transform(1+t/beta)\^{}-83. The normalization and transform follow by83 elementary integrations by parts, or by the same integer gamma integral. Every integrand is positive.

For any sigma\textgreater0, on the original0\textless=t\textless=R, \[
 \Psi(t)\le e^{\sigma R}
                \int_0^\infty e^{-(\sigma+b)t}\,d\gamma(b).
 \tag{GM7}
\] Indeed the right side is e\^{}\{sigma(R-t)\}(1+t/beta)\^{}-83. Thus the entire low-b region is retained after a positive shift. There is no discarded mixing interval, divergent inverse moment or separately unpaid tail. In a product over f embeddings the one paid factor is e\^{}(sigma R f).

\subsection{3. Every moment needed by the full positive expansion}\label{every-moment-needed-by-the-full-positive-expansion}

For a fixed k define \[
 d=k+3>83,\quad\sigma_k=(d-83)/R,\quad
 \Lambda_j=\int_0^\infty(\sigma_k+b)^{-j}\,d\gamma(b)
                     \quad(0\le j\le d).
 \tag{GM8}
\] All these integrals converge because sigma\_k\textgreater0; Lambda0=1. Dropping e\^{}-beta b, and substituting b=sigma\_k u/(1-u), gives the ordinary convergent beta integral \[
 \Lambda_d< B_d:=(7/2)^{83}\sigma_k^{83-d}
                         \frac{(d-84)!}{(d-1)!}.
 \tag{GM9}
\] Both beta arguments83 and d-83 are positive integers. This is not analytic continuation of an unshifted high inverse moment.

Choose the positive rational beta\_k in the table below so that beta\_k\^{}d B\_d\textless1. Holder's inequality on the probability space gamma gives \[
 \Lambda_j\le\Lambda_d^{j/d}\Lambda_0^{1-j/d}
                 \le B_d^{j/d}\le\beta_k^{-j},\qquad0\le j\le d.
 \tag{GM10}
\] At j=0 both sides are1; for positive j the last inequality is strict. Hence every degree of the scalar and transverse polynomial, including degree zero, is controlled by the same explicitly certified effective parameter. The gamma rate beta=7/2, shift sigma\_k and effective beta\_k are distinct quantities.

The paid shift is e\^{}((d-83)f). The elementary rational comparison e\textless87/32 and(87/32)\^{}9\textless2\^{}13 proves \[
                      e^{(d-83)f}<2^{13(d-83)f/9}.
 \tag{GM11}
\] The executable bounds the exponential series by its first nine terms plus a geometric remainder before making these integer-power comparisons.

\subsection{4. Anisotropic actual lattices and all mixture terms}\label{anisotropic-actual-lattices-and-all-mixture-terms}

Use the canonical all-parameter/all-minimum scalar coefficient of K111: \[
 c_{336}=168{336\choose168}/4^{168},\quad
 M_{15}=\frac{6015684926187167549292201834533348561577984}
 {2618017117694612446410293519},
\] \[
 K=\frac{336M_{15}}{c_{336}}(7841/5000)^{168}168!
       \left(\frac{(365729/500)^2}{168(314159/100000)^2}\right)^{168}.
 \tag{GM12}
\] The original saturated transverse Gram has det G'=U and G'\textgreater c0 I with c0=2\^{}36. Its values belong to the actual common ideal for every original integral Z input; this is the K107/K109 interface retained in K113, not division of a previous lattice value by two. The primal covolume, actual projection denominator and half-trace dual normalization remain unchanged.

For each positive embedding weight c\_tau=sigma\_k+b\_tau, the AM-GM dual minimum of K110 PM10 applies. It is valid for arbitrary positive weights, not only c\_tau\textgreater=2. The diagonally transformed real lattice need not be an ideal. Let t0=Uz²=R/q and B=(product c\_tau)\^{}(1/f). The scalar Poisson majorant is \[
 \frac{\pi^f}{D_0}\{(t_0B)^{-f}+K\}.
\] The true three-dimensional transverse factor is \[
 \frac{\pi^{3f}U^f}{D_0^3t_0^{3f}}
        (\prod c_\tau)^{-3}
           [1+Kc_0^{-f}t_0^f\prod c_\tau]^3.
 \tag{GM13}
\] These are the first, untruncated inequalities in K110 PM11--12. Their source theorem is universal in the positive Gaussian parameter and real scalar minimum.

Expand the k scalar factors and all three transverse factors before integrating. A term with l scalar volume factors and h transverse theta corrections has the inverse moment degree l+3-h, which ranges from0 to d.~All coefficients are nonnegative. Tonelli and GM10 therefore control every term, with no additive remainder. The full weighted row sum is bounded by \[
 e^{(d-83)f}{164\choose k}\frac{U^f\pi^{df}}{D_0^d}
 \left\{\left(\frac q{\beta_k R}\right)^f+K\right\}^{k}
 \left(\frac q{\beta_k R}\right)^{3f}
 \left[1+K\left(\frac{\beta_k R}{c_0q}\right)^f\right]^3.
 \tag{GM14}
\] Only after GM3 and GM7 have been applied to the original rows is this positive lattice sum enlarged. The actual selected prefactor, compact factors, energy cap and norm-product restriction have not been reassigned to the extra zero or out-of-ball tuples.

\subsection{5. Complete arithmetic weights, central fibre and continuum}\label{complete-arithmetic-weights-central-fibre-and-continuum}

All arithmetic interfaces are precisely K113 SC27.4--10. In GM1 the slope eta is unique for each fixed27-place active support because x\textless2\^{}3328\textless2\^{}12000 and the paired difference bound gives (4U)\^{}168·2\^{}24000\textless T\^{}54. There is no denominator-ideal sum or separate beta norm fibre. Eta, the complete167 original covector values and the original central scalar reconstruct h.

There are no inactive paired-zero places or active integral-C places. All27 active weighted label factors and all141 inactive pointwise factors contribute at most(T-1)\^{}-168\textless2T\^{}-168. The actual Green and central bounds are \[
 |a_G(h)|\le2Q^{-335\cdot167}(N_EJ)^{335/2},\quad
 N_EJ=z^{2f}T^{54}/x,\quad N_EJ\le Q^{334},
\] \[
 \#d_{\rm cent}\le
 \frac{V_F(C(q))U^{-f}T^{54}x^{-1}}{N_FJ_0},\quad
 C(q)=H/2+750R/q,\quad
 V_F(C)=\frac{((22/7)fC^2)^{84}}{84!\sqrt{D_F}},
\] \[
                  N_EJ\le Q^{167}N_FJ_0.
 \tag{GM15}
\] This is one actual affine central coset, with its original real-part shift, full750 covering radius and all dyadic conditions. The shift in the analytic measure changes none of these arithmetic ideals.

The complete multiplicities are support\textless2\^{}104, original pivot\textless2\^{}8, Green correction2, full T correction2, selected near-point factor2, and a conservative full pattern bound4. Their fixed binary exponent is117. Both mixed projectors keep their contraction bounds. No omitted target, source stabilizer or eta enumeration is replaced by a smaller unproved object: the reconstruction is the same full rank-two normal form used in K113.

For fixed k let S\_k(q)=(q/(beta\_k R))\^{}f+K. In a transverse correction h, the q-dependent expression before the central correlation is proportional to \[
 F_{k,h}(q)=q^{-(164+h)f}S_k(q)^k V_F(C(q)).
 \tag{GM16}
\] Put w=(q/(beta\_k R))\^{}f/S\_k(q), epsilon=(C(q)-H/2)/C(q), so0\textless w,epsilon\textless1. Direct differentiation gives \[
 \frac1f\frac{d\log F_{k,h}}{d\log q}=-164-h+kw-\epsilon<0,
\] \[
 \frac{d\log\{x^{-337/2}F_{k,h}(q(x))\}}{d\log x}
            =-9/2+h-kw+\epsilon<-1/2.
 \tag{GM17}
\] The effective beta\_k and paid shift depend only on k, not on x, q or the pattern. Thus the complete continuum is bounded at the lower x endpoint2111, before any pattern sum.

For fixed x write \[
 q_*^f=\frac{R^fT^{54}}{U^fQ^{334}x},\quad
 D=Q^{334}/N_EJ,\quad q=q_*D^{1/f},\quad
 (N_FJ_0)^{-1}\le Q^{-167}D.
 \tag{GM18}
\] Expanding S\_k(q)\^{}k gives, for scalar degree l and transverse correction h, the exact containing pattern product \[
             \prod_{\mathfrak q\mid G}(1+q_{\mathfrak q}^{l-163-h})^2.
 \tag{GM19}
\] Since l\textless=k\textless=109, its exponent is at most-54. The actual primes satisfy(1+g0\^{}-1)\^{}668\textless2, which bounds every product; the retained4 bound is conservative. It includes all four local states with multiplicities1,2,1. No exceptional high-degree pattern or separate rank distribution has been omitted.

\subsection{6. Ten exact complete endpoints and one fresh union}\label{ten-exact-complete-endpoints-and-one-fresh-union}

At x=2\^{}2111 the true cap satisfies \[
 251124964<q_*<251124965,
 \qquad \lceil\log_2V_F(H/2+750R/251124964)\rceil_{\rm strict}=474.
 \tag{GM20}
\] Both root comparisons are exact168th-power integer comparisons. Use the unchanged directed logarithms LT=1019333/1024, LQ=287843/4096, LU=601356/4096 and LR=28810/4096, with upper signs for T,R and lower for Q,U. Let \[
 C=117-933f/2-649+474+167fL_R-(335f/2)L_U
          +8931L_T-336\cdot167L_Q-(337/2)2111.
 \tag{GM21}
\] The complete T power8931 is335·27+54-168. The Q power includes the full central correlation. The original U,R and x powers are unchanged. No scalar-code coefficient is being substituted for the physical selected module.

For q\_-=251124964 and q\_+=251124965 compute the strict binary ceiling g\_k of the complete rational expression \[
 W_k=\frac{{164\choose k}(22/7)^{df}S_k(q_+)^k
       (q_+/(\beta_kR))^{3f}
       [1+K(\beta_kR/(c_0q_-))^f]^3}
           {D_0^d q_-^{167f}k^{84}}.
 \tag{GM22}
\] Then E\_k=C+13(d-83)f/9+g\_k is a full upper exponent, including the paid shift. The exact rational values are:

\begin{longtable}[]{@{}rrrr@{}}
\toprule\noalign{}
k & Effective beta\_k & Strict g\_k & Complete E\_k \\
\midrule\noalign{}
\endhead
\bottomrule\noalign{}
\endlastfoot
100 & 26437/4096 & -604739 & -64500983/3072 \\
101 & 26163/4096 & -602745 & -19209981/1024 \\
102 & 6475/1024 & -600754 & -50768119/3072 \\
103 & 12823/2048 & -598764 & -43909367/3072 \\
104 & 12701/2048 & -596778 & -12354301/1024 \\
105 & 12583/2048 & -594792 & -30216439/3072 \\
106 & 24939/4096 & -592810 & -23382263/3072 \\
107 & 1545/256 & -590829 & -5517053/1024 \\
108 & 6127/1024 & -588850 & -9726199/3072 \\
109 & 1519/256 & -586873 & -2907383/3072 \\
\end{longtable}

The worst bound is strictly below-946, and every other bound is below-3166. Thus the full ten-cell union satisfies \[
             \sum_{k=100}^{109}2^{E_k}
              <2^{-946}+9\,2^{-3166}<2^{-945}.
 \tag{GM23}
\] Each cell already includes its full coordinate-position binomial, active supports, slope uniqueness, central coset, all G patterns and all168 T factors. The ten-cell sum is the only remaining union. The same sufficient envelope at k110 is positive1300739/1024 and makes no claim for that cell. It is not a lower bound on the actual sum.

\subsection{7. Chronology, checks and remaining scope}\label{chronology-checks-and-remaining-scope}

The actual current predecessor is checks/do120\_joint\_remainder\_error.json,127 all-top and31 uniform entries, complete floors27/11. The executable accepts a later complete ordered extension and preserves the historical author output. Append945 once, giving128/31 on that predecessor. Both complete dyadic sums must remain below2\^{}-230; every older overlapping charge and uniform entry remains.

The certificate checks the cap-dependent profile endpoint, gamma/beta factorial normalization, every effective moment parameter, degree-zero and all positive expansion-degree coverage, original field constants and directed logarithms, all continuum derivative signs, the full finite-pattern envelope, both cap roots, the real central-volume ceiling, every complete rational row exponent, the full ten-cell union and complete prior/new error arrays. These diagnostics supplement the universal proofs above. They do not assert the existence of a rational short arrow or replace the actual selected module.

Direct mathematical interfaces are K113 SC27.1--22, K110 PM3--17, K111 HP1--18, K109's actual saturated Gram and all-parameter dual estimate, and K108 AE13 with the complete K107 near-point/tail correction. All their original source/class/lattice and metric normalizations remain. The new proof is the bounded-energy Gamma majorant, fully paid shift, all-degree Holder envelope and its complete support27 count. No new external representation or algebraic-geometric comparison theorem is invoked.

Carried exact contextual anchors remain SGA3 I4.7.3/4.7.3.1; SGA6 VI Theorem3.3; FGA232 Theorem3.1 pp148--149; FGA236 Corollaries1.7/2.3 pp229/232; SGA7 I VII2.7/2.9 pp169--174 and VIII7.4(b) pp310--311; cote4 scan003 sections8.1.1--8.1.2 and scans005--006 sections8.4.1--8.4.3; and cote54 scans005--007. These are inherited context. DP\_MULT Q3e/DP44, THETA3 G1/A7, CX\_GRH CX-ALG-024 and GCT\_PLETH T-052 keep their separate signed, source, selected arithmetic-content, prescribed-period and global comparison gates. No complete support28 assertion, Lean build or archive audit is made.

\subsection{Canonical algebraic-geometric context}\label{canonical-algebraic-geometric-context-62}

The analytic probability measure is a bound on the same selected real realization; it introduces no arithmetic parameter, field extension, physical source vector, lattice rescaling or projector. All original projective OE modules, fractional paired ideals, CM involution, polarization, theta line, Picard and Poincare data remain. The new Gamma/beta/Holder and positive-integral arguments are proved directly; original representation and harmonic-code primary scopes are inherited through K107/K108/K110/K111/K113. Carried exact anchors: SGA3 I4.7.3/4.7.3.1; SGA6 VI Theorem3.3; FGA232 Theorem3.1 pp148-149; FGA236 Corollaries1.7/2.3 pp229/232; SGA7 I VII2.7/2.9 pp169-174 and VIII7.4(b) pp310-311; cote4 scan003 sections8.1.1-8.1.2 and scans005-006 sections8.4.1-8.4.3; cote54 scans005-007. No new formal-group/Dieudonne, lambda-ring/K-theory, l-adic-weight or CM-motive comparison. DP\_MULT Q3e/DP44 and THETA3 G1/A7 receive a strictly smaller conditional finite gate; CX\_GRH CX-ALG-024 and GCT\_PLETH T-052 retain their global gates.

\clearpage

\section{AM-LIFT-098 --- Explicit volume of the same arithmetic period family}\label{am-lift-098-explicit-volume-of-the-same-arithmetic-period-family}

\textbf{Canonical checked integration --- 2026-11-23T09:27:02.631222+00:00.} Complete saved-proof and executable reviews and separate exact diagnostic replays passed. No numerical error charge is introduced. Frozen author history is preserved. AM-LIFT-098 AV1-20 evaluates the previously unevaluated automorphic volume for the same actual principal-3M special-unitary arithmetic family. Put n=3M, r=p+q\textgreater=3, f={[}F:Q{]}, g=fr, d=fpq, a=r\^{}2-1, a\_r=product\_\{j=1\}\^{}\{r-1\} j!/(2pi)\^{}(j+1), b\_pq=product\_\{i=0\}\^{}\{p-1\}i!/(q+i)!, and e\_v=product\_\{j=2\}\^{}r(1-epsilon\_v\^{}j q\_v\^{}(-j))\^{}(-1), epsilon\_v=1 at split and -1 at inert finite places. The literal relative-unramified CM and ordinary selfdual-lattice assumptions give hyperspecial models even at dyadic places. The original invertible ideal multiple changes no stabilizer or congruence index. Exact strong approximation and finite orders give mu\_Gamma=D\_F\^{}(a/2)a\_r\^{}f n\^{}(fa) product\_\{v not dividing n\}e\_v with compact SU mass one. The physical Hodge representation S+Q\^{}dual gives lambda\^{}r=K\_Y\^{}2 and A\^{}r=K\_X\^{}(2(g+1)); Mumford proportionality yields v\_A={[}2(g+1){]}\^{}d d! b\_pq\^{}f mu\_Gamma. An explicit integer upper bound is Vcal=n\^{}(fa)D\_F\^{}ceil(a/2)2\^{}(f(r-1)){[}2(g+1){]}\^{}d d!(product\_\{j=1\}\^{}\{r-1\}j!)\^{}f.~Thus D\_arith=2\^{}(8g\^{}2+2g+4)h\_d\^{}d Vcal sum\_\{j=0\}\^{}d k\^{}j, h\_d=2\^{}(d+1)(d+1)(d+2)!, retains the full same-rule cost, and N\_within\textless=t\^{}\{B\_\{t\^{}(4fr\^{}2),d-1\}(D\_arith)\}, t=1+3M Q\_S. The independent volume input has been replaced by actual field, rank, signature and level data. Any component-entry N0, same-family nonempty open, selected arithmetic nonzero rule and admitted Host/Sizes remain open. No evaluated actual N, six-identity mixed-CM chain, period, specific032 pairing, DP44 difference or RH conclusion. No finite error charge; current K115128/31 arrays and complete floors27/11 unchanged.

\textbf{Complete author subject; final saved-file, source and diagnostic verification are recorded separately. No canonical or independent-review claim.} This is a source-conditional volume theorem for the exact family of AM-LIFT-094--097. It does not construct an admitted source, an entry point in the same-rule nonzero locus, or a bound for the entry denominator. All constants below refer to the original Hermitian lattice and the auxiliary proof level, not to a replacement finite kernel.

\subsection{1. Exact family and statement}\label{exact-family-and-statement}

\textbf{AV1.} Let \(F\) be totally real of degree \(f\), let \(E/F\) be a CM quadratic extension unramified at every finite place, and let \((V,h)\) have \(E\)-dimension \(r=p+q\ge3\), with \(p,q\ge1\) and signature \((p,q)\) at every selected real place. Let \(L_0\subset V\) be a full selfdual \(\mathcal O_E\)-lattice for the ordinary Hermitian pairing. Let \(\mathfrak a\) be an invertible fractional \(\mathcal O_E\)-ideal and put \(\Lambda=\mathfrak aL_0\). This includes the actual lattice transported by \nolinkurl{IsDatW}, which is \nolinkurl{AA · L0}; its trace alternating form and its sign are retained. Let \(M\ge1\) be the given positive integer physical level. Put

\[
n=3M,\quad G=\operatorname{SU}(V,h)/F,\quad a=r^2-1,\quad
g=fr,\quad d=fpq,
\] \[
\Gamma_0=\{\gamma\in G(F):\gamma\Lambda=\Lambda\},\qquad
\Gamma=\{\gamma\in\Gamma_0:(\gamma-1)\Lambda\subset n\Lambda\}.
\tag{AV1}
\]

Use precisely the connected symmetric domain \(\mathcal D=\prod_{\tau:F\hookrightarrow\mathbb R}\mathcal D_{p,q}\), the family \(Y=\Gamma\backslash\mathcal D\), and its normal Baily--Borel compactification \(X\) from LIFT096--097. The extra marking at \(3\) is a proof-level refinement subsequently forgotten; the physical level \(M\), finite eta, rational Hermitian space and source period-family map are unchanged. The group is special unitary, not the full unitary group or an adjoint normalizer.

Let \(\lambda=\det(e^*\Omega^1_{\mathcal A/Y})\) be the determinant of the full rank-\(g\) physical Hodge bundle. Let \(A\) be the specific Cartier ample line of LIFT096, with \(A|_Y=\lambda^{g+1}\). Its degree is \(v_A=\int_Xc_1(A)^d\). These are the already proved geometric hypotheses of that same family; they are not deduced here from an arbitrary abstract rule.

\textbf{AV2.} Write \(D_F=|\operatorname{disc}(F)|\). For a finite place \(v\), let \(q_v\) be its residue cardinality and let \(\epsilon_v=1\) if \(v\) splits in \(E\), and \(\epsilon_v=-1\) otherwise. Define

\[
 e_v=\prod_{j=2}^{r}(1-\epsilon_v^{j}q_v^{-j})^{-1},\qquad
 a_r=\prod_{j=1}^{r-1}\frac{j!}{(2\pi)^{j+1}},\qquad
 b_{p,q}=\frac{\prod_{i=0}^{p-1}i!}{\prod_{i=0}^{p-1}(q+i)!}.
\tag{AV2}
\]

Then the exact volume/intersection identity is

\[
\boxed{
v_A=[2(g+1)]^d d!\,b_{p,q}^{f}\,
D_F^{a/2}a_r^{f}n^{fa}\prod_{v\nmid n}e_v .}
\tag{AV3}
\]

In particular the following explicit integer bounds it:

\[
\boxed{
\mathcal V=n^{fa}D_F^{\lceil a/2\rceil}
2^{f(r-1)}[2(g+1)]^d d!
\left(\prod_{j=1}^{r-1}j!\right)^f,
\qquad v_A\le\mathcal V .}
\tag{AV4}
\]

The Euler product in AV3 converges absolutely. The proof below identifies the Haar normalization with the proportionality constant; AV3 is not obtained by merely renaming a group covolume as a projective degree.

\subsection{2. The original lattice, all finite places, and the exact index}\label{the-original-lattice-all-finite-places-and-the-exact-index}

\textbf{AV3a.} Multiplication by the invertible ideal \(\mathfrak a\) preserves both stabilizer and principal-congruence conditions. Indeed, \(\gamma(\mathfrak aL_0)=\mathfrak a\gamma L_0\), and multiplication by \(\mathfrak a^{-1}\) reverses each inclusion. Similarly,

\[
(\gamma-1)\mathfrak aL_0\subset n\mathfrak aL_0
\Longleftrightarrow (\gamma-1)L_0\subset nL_0.
\tag{AV5}
\]

No generator of \(\mathfrak a\), moving integral basis or index between \(\Lambda\) and \(L_0\) enters. This is especially relevant when \nolinkurl{AA} is not principal.

\textbf{AV4a.} At a finite place put \(R=\mathcal O_{F_v}\), \(S=\mathcal O_E\otimes_{\mathcal O_F}R\). The hypothesis says \(S/R\) is quadratic finite étale (including \(R\times R\)). The localized lattice is free over each local factor and its Hermitian pairing is perfect. The determinant-one unitary group scheme of this lattice becomes \(\operatorname{SL}_r\) after the finite étale base change splitting \(S/R\): the two factors of the module are dual through the perfect pairing, so an automorphism on one factor determines the other. Consequently it is a smooth connected reductive \(R\)-model with generic fibre \(G_{F_v}\). This argument uses finite étale descent and never divides by \(2\). It therefore includes relatively unramified dyadic extensions even when \(F_v/\mathbb Q_2\) is ramified. By Tits (1979, §3.8.1, p.55), its \(R\)-points \(P_v\) are hyperspecial. In particular the generic group is the unramified quasi-split form at \(v\).

For comparison, the direct dyadic hyperbolic splitting in LIFT095 used integral trace surjectivity; the reductive-model argument here gives the full group scheme and its reduction, which are what the volume calculation needs. A ramified quadratic dyadic extension is not being covered by either assertion.

\textbf{AV5a.} The special fibre has order

\[
|\overline G_v(\mathbb F_{q_v})|=q_v^{r(r-1)/2}
\prod_{j=2}^{r}(q_v^j-\epsilon_v^{j})
=q_v^{a}/e_v.
\tag{AV6}
\]

Here is an elementary verification of the unitary case, valid also in characteristic \(2\). In a nondegenerate Hermitian \(m\)-space over \(\mathbb F_{q^2}\), let \(A_m\) count norm-zero vectors and \(B_m\) count vectors of any one fixed nonzero norm. The norm map onto \(\mathbb F_q^\times\) makes the latter counts equal, and

\[
A_m+(q-1)B_m=q^{2m},\qquad
A_{m+1}=A_m+(q^2-1)B_m.
\]

Starting from \(A_0=1,B_0=0\), induction gives \(B_m=q^{m-1}(q^m-(-1)^m)\). Counting ordered orthonormal bases gives the product from \(j=1\) to \(r\) for \(U_r(q)\). The determinant is onto the norm-one group of order \(q+1\) (diagonal matrices suffice); division by \(q+1\) gives AV6. The split case is the usual ordered-basis count for \(\operatorname{GL}_r(q)\), divided by \(q-1\). Cameron (§5.1, Proposition 5.1, p.58) gives the same finite-unitary formula and proof.

Write \(a_v=v_v(n)\ge1\) at \(v\mid n\). Smoothness over the complete DVR gives surjective successive reduction maps; their kernels have the cardinality of the additive Lie algebra, \(q_v^a\). Thus

\[
|P_v/P_v(a_v)|=q_v^{(a_v-1)a}|\overline G_v(\mathbb F_{q_v})|
=q_v^{a_va}/e_v.
\tag{AV7}
\]

There is no assumption that \(L_0\) is globally free. The reductions and the smooth integral models are local, and they glue as the actual coherent collection of stabilizers.

\textbf{AV6a.} Strong approximation applies to the simply connected absolutely almost simple group \(G\), away from all archimedean places: \(G(F\otimes\mathbb R)\) is noncompact since \(p,q\ge1\). Consequently \(G(F)\) is dense in \(G(\mathbb A_{F,f})\), and \(\Gamma_0=G(F)\cap\prod_vP_v\) maps onto every finite quotient of the open compact product. Applying this to AV7 proves the exact index

\[
[\Gamma_0:\Gamma]=\prod_{v\mid n}q_v^{a_va}e_v^{-1}
=n^{fa}\prod_{v\mid n}e_v^{-1}.
\tag{AV8}
\]

This uses the same strong-approximation theorem and hypotheses as LIFT095; it introduces no extra adelic class-number factor. The level-\(3M\) group is neat: in the integral rational representation on \(\Lambda\), every eigenvalue of an element congruent to \(1\bmod3\) is \(1\bmod3\) in a \(3\)-adic algebraic closure; so is every multiplicative combination of these eigenvalues, and no nontrivial root of unity has that congruence. Hence no torsion or ineffective central stabilizer remains in \(\Gamma\backslash\mathcal D\).

\subsection{3. Prasad's normalized covolume and the level cancellation}\label{prasads-normalized-covolume-and-the-level-cancellation}

\textbf{AV7a.} Normalize the Haar measure at each real place as in Prasad (1989, §§3.5--3.6, pp.110--111): the invariant top exterior form is chosen so that its absolute measure on the compact real form \(\operatorname{SU}(r)\) has total mass \(1\). Let \(\mu_\infty\) be the product measure on \(G_\infty=\prod_\tau\operatorname{SU}(p,q)\). This is not the Euler--Poincaré measure.

Prasad's Theorem 3.7, its §§3.3--3.4 hypotheses, and Remark 3.11 apply to the coherent hyperspecial collection above. The Tamagawa number of this simply connected group is \(1\). Its absolute type is \(A_{r-1}\), whose exponents are \(1,\ldots,r-1\). Its quasi-split splitting extension is \(E/F\). Since that extension is unramified at every finite place,

\[
D_E/D_F^2=N_{F/\mathbb Q}(\mathfrak D_{E/F})=1;
\]

thus the relative-discriminant factor of Theorem 3.7 is exactly \(1\), regardless of the parity-dependent exponent of type \(^{2}A\). At a hyperspecial place both reductive dimensions in Prasad's local factor equal \(a\), so the factor is exactly \(q_v^a/|\overline G_v(\mathbb F_{q_v})|=e_v\). It follows that

\[
\mu_\infty(\Gamma_0\backslash G_\infty)
=D_F^{a/2}a_r^f\prod_v e_v.
\tag{AV9}
\]

The quotient for \(\Gamma_0\) is a group quotient even if that larger arithmetic group has torsion. Multiplication by the finite index AV8 is therefore valid without an orbifold convention. We obtain

\[
\boxed{\mu_\infty(\Gamma\backslash G_\infty)
=D_F^{a/2}a_r^f n^{fa}\prod_{v\nmid n}e_v.}
\tag{AV10}
\]

In particular every level Euler factor has cancelled. Strong approximation also identifies the relevant principal arithmetic quotient without an extra genus, ideal-class, connected-component or normalizer index. No assertion about all components of the original source moduli is made: AV10 belongs to this particular fixed rational Hermitian family.

\subsection{4. Haar measure and the proportionality constant}\label{haar-measure-and-the-proportionality-constant}

\textbf{AV8a.} Set

\[
K_\infty=\prod_\tau S(U(p)\times U(q)),\qquad
G^c=\prod_\tau SU(r),\qquad
\mathcal D^c=G^c/K_\infty=\prod_\tau\operatorname{Gr}(p,r).
\]

Give \(K_\infty\) Haar mass \(1\). Use this same compact group and same measure for the quotient of \(G_\infty\) and of \(G^c\). By Weil disintegration the invariant quotient density on \(\mathcal D^c\) has total mass \(1\), because the Prasad-normalized compact group has mass \(1\). At the tangent-space level the two real forms are \(\mathfrak k\oplus\mathfrak p\) and \(\mathfrak k\oplus i\mathfrak p\). Analytic continuation preserves absolute exterior volume, so these are precisely the paired invariant densities used in Hirzebruch--Mumford proportionality. The \(K_\infty\)-fibres over \(Y\) have mass \(1\); neatness makes their action free. Therefore the proportionality constant is

\[
c(\Gamma)=\frac{\operatorname{vol}(Y)}{\operatorname{vol}(\mathcal D^c)}
=\mu_\infty(\Gamma\backslash G_\infty).
\tag{AV11}
\]

Mumford (1977, Theorem 3.2 and its proof, pp.262--263) explicitly identifies the constant with this ratio and supplies the sign \((-1)^d\) on degree-\(d\) characteristic numbers of canonically extended automorphic bundles.

For an independent normalization check, Borel and Prasad (1989, §4.2, p.137) identify the absolute Euler--Poincaré measure with \(c_v\mu_v\), where \(c_v\) is the Euler characteristic of the compact dual. Here that is \(\binom{r}{p}\) at each real place. Thus

\[
|\chi(\Gamma)|=\binom{r}{p}^{f}\mu_\infty(\Gamma\backslash G_\infty).
\tag{AV12}
\]

If one starts instead with an Euler characteristic, division by this factor is necessary. AV10 starts with Prasad's compact-mass-one Haar, so no such factor is inserted in AV11.

\textbf{AV9a.} The real group \(SU(p,q)\) and its maximal compact \(S(U(p)\times U(q))\) are connected. We use one connected domain and one arithmetic group, not a union of determinant classes. There is no continuous unitary central torus because the group is \(SU\). Its finite centre is contained in the same \(K_\infty\) on both sides. We have not replaced \(SU\) by its adjoint quotient; the compact normalization and quotient integration are already compatible with that centre. Neatness gives \(\Gamma\cap Z(SU)=1\). If both groups were passed to the adjoint quotient, the finite covering degree and normalization of the compact form would cancel in AV11; appending an isolated factor \(|Z|\) would be incorrect.

\subsection{5. The physical Hodge line and its exact compact-dual degree}\label{the-physical-hodge-line-and-its-exact-compact-dual-degree}

\textbf{AV10a.} On one compact-dual factor let \(S\) be the tautological rank-\(p\) bundle and \(Q\) the rank-\(q\) quotient of the constant \(\mathbb C^r\). The complexification of the underlying trace symplectic representation is \(W\oplus W^*\), and the Hodge Lagrangian at a \(p\)-plane \(S\subset W\) is \(S\oplus\operatorname{Ann}(S)\). Since \(\operatorname{Ann}(S)=Q^*\), the full physical Hodge bundle on the compact dual is \(S\oplus Q^*\), with the roles of \(p,q\) exchanged if the opposite CM representative is used. Its determinant and the canonical line are

\[
\lambda_c=(\det S)^2(\det\mathbb C^r)^{-1},\qquad
K_{\operatorname{Gr}(p,r)}=(\det S)^r(\det\mathbb C^r)^{-p}.
\tag{AV13}
\]

These identities follow respectively from \(\det S\det Q=\det\mathbb C^r\) and \(T=\operatorname{Hom}(S,Q)\). On the actual special unitary group the determinant character of \(\mathbb C^r\) is trivial. Therefore, on the product domain and its quotient,

\[
\lambda^r\simeq K_Y^2,\qquad A^r|_Y\simeq K_Y^{2(g+1)}.
\tag{AV14}
\]

There is no untracked flat determinant twist. The canonical Cartier line \(K_X\) and the line \(A\) are the actual lines established in LIFT097 and LIFT096. The Baily--Borel boundary has codimension at least two for these \(r\ge3\) type-\(A\) factors. Since \(X\) is normal, the isomorphism of line bundles AV14 extends uniquely across that boundary:

\[
A^r\simeq K_X^{2(g+1)}.
\tag{AV15}
\]

This also identifies the intersection normalization without assuming an arbitrary extension of the physical Hodge line. On a smooth toroidal resolution the canonical automorphic extension of the cotangent is the logarithmic cotangent, and its determinant is the pullback of the Baily--Borel canonical automorphic line: Mumford, Proposition 3.4(a,b) and proof, pp.264--265. Thus its Chern numbers compute \(K_X^d\), and AV15 computes \(A^d\). There is no additional toroidal boundary intersection term.

\textbf{AV11a.} Write \(H_\tau=c_1(\det S_\tau^*)\) on the compact dual. Then \(c_1(K_{\mathcal D^c})=-r\sum_\tau H_\tau\). The Plücker degree of one Grassmannian is

\[
\int_{\operatorname{Gr}(p,r)}H^{pq}=(pq)!
\frac{\prod_{i=0}^{p-1}i!}{\prod_{i=0}^{p-1}(q+i)!}
=(pq)!b_{p,q}.
\tag{AV16}
\]

For completeness, repeated Pieri multiplication by the hyperplane class counts standard tableaux of the rectangular \(p\times q\) diagram. The hook-length formula divides \((pq)!\) by \(\prod_{i=0}^{p-1}\prod_{j=0}^{q-1}(i+j+1)=\prod_{i=0}^{p-1}(q+i)!/i!\), proving AV16. Smirnov (2015, §2.9, Theorems 2.31--2.32, pp.18--19) gives these two steps. This proof makes the symmetric \(p,q\) normalization explicit.

Only the term \(H_\tau^{pq}\) at every factor survives in the product intersection. Hence

\[
(-1)^d\int_{\mathcal D^c}c_1(K_{\mathcal D^c})^d
=r^d d! b_{p,q}^{f}.
\tag{AV17}
\]

Mumford proportionality, AV11, and AV15 now give

\[
v_A=\left(\frac{2(g+1)}r\right)^d K_X^d
=[2(g+1)]^d d!b_{p,q}^{f}\mu_\infty(\Gamma\backslash G_\infty).
\tag{AV18}
\]

Combining AV18 with AV10 proves AV3. Positivity agrees with \(A\) ample; the two \((-1)^d\) signs have cancelled.

\subsection{6. An explicit integer majorant and the complete escape cost}\label{an-explicit-integer-majorant-and-the-complete-escape-cost}

\textbf{AV12a.} For each \(v\), all factors of \(e_v\) are positive, and

\[
e_v\le\prod_{j=2}^{r}(1-q_v^{-j})^{-1}.
\]

The absolutely convergent product is bounded by \(\prod_{j=2}^r\zeta_F(j)\). Primewise, a rational prime has at most \(f\) primes above it, all residue degrees at least one; therefore \(\zeta_F(j)\le\zeta(j)^f\), including rational primes ramified in \(F\). For every integer \(j\ge2\),

\[
\zeta(j)\le\zeta(2)=1+\sum_{m=2}^{\infty}m^{-2}
<1+\int_1^{\infty}x^{-2}\,dx=2.
\]

Thus \(\prod_{v\nmid n}e_v<2^{f(r-1)}\). Also \(a_r\le\prod_{j=1}^{r-1}j!\), \(b_{p,q}\le1\), and \(D_F^{a/2}\le D_F^{\lceil a/2\rceil}\). Substitution proves the explicit integer AV4. The deliberately coarse removal of \(\pi\), the Plücker denominator and the unused level Euler factors only enlarges the bound.

\textbf{AV13a.} Retain the already checked quantities

\[
h_d=2^{d+1}(d+1)(d+2)!,\qquad e_g=2^{8g^2+2g+4},
\]

and the same-rule scalar exponent \(k=gw+n_F/2\) in the source even-weight case, or \(k=2gw+n_F\) after squaring the same rule, as in LIFT096. Here \(n_F\) is the rule's theta weight and is unrelated to the proof level \(n=3M\). Then

\[
\deg_{A^{h_d}}X=h_d^dv_A\le h_d^d\mathcal V,
\qquad
D_*\le D_{\rm arith}:=e_g h_d^d\mathcal V\sum_{j=0}^{d}k^j.
\tag{AV19}
\]

Every number on the right is explicit in the given field discriminant, \(f,r,p,q,M\), and the fixed rule weights. The same conditional recursion of LIFT094--096 therefore gives

\[
N\le t^{B_{\delta,d-1}(D_{\rm arith})},\qquad
t=1+3M Q_S,\quad\delta=t^{4fr^2},
\tag{AV20}
\]

once the identical family, projective correspondence extension, same-rule nonempty intersection and entry-point hypotheses of those findings hold. The monotonicity of that recursion permits the larger integer degree. The auxiliary \(c=1\) comes from LIFT095, not from a change of the original lattice.

\subsection{7. Exact source applicability and limits}\label{exact-source-applicability-and-limits}

The source clauses read for this theorem are \nolinkurl{CMData.cmUnram/totReal/deg2/fixed} (DataSpaces14--19), the actual pairing and full selfdual lattice (26--46), actual W-rank/signature (121--124), \nolinkurl{IsCMUnram} (StatementMass44--48), ordinary \nolinkurl{hdual} (StatementHerm9--12) and \nolinkurl{IsUnimodular} (StatementLat31--38). At a ramified quadratic prime conjugation is trivial on the residue field, contradicting the source's required integral \(\theta-\bar\theta\notin\mathfrak p\); hence those literal fields imply the finite relative-unramified premise. No \nolinkurl{SiegelMassRatio} proposition is used as a theorem.

\nolinkurl{IsDatW} in PartitionExtendByZeroAntitone24--31 identifies the exact \nolinkurl{AA · L0} lattice and trace alternating form; \nolinkurl{PinLevel}43 fixes \(M=g_0g_1\) and the source level-record parameter \(A_{\rm level}=1\) (unrelated to the ample line \(A\) above), while \nolinkurl{PinDat}47 supplies the actual signature. The actual interior \(W_j\) has \(p=k+j-1\), \(q=m-j>0\), \(r=k+m-1\) by Params32,46 and DataSpaces121--124. The terminal definite \(q=0\) case is not covered. All conclusions are conditional on the literal source instance and geometric family required by LIFT096--097.

\begin{longtable}[]{@{}
  >{\raggedright\arraybackslash}p{(\columnwidth - 2\tabcolsep) * \real{0.5000}}
  >{\raggedright\arraybackslash}p{(\columnwidth - 2\tabcolsep) * \real{0.5000}}@{}}
\toprule\noalign{}
\begin{minipage}[b]{\linewidth}\raggedright
Quantity or condition
\end{minipage} & \begin{minipage}[b]{\linewidth}\raggedright
Treatment here
\end{minipage} \\
\midrule\noalign{}
\endhead
\bottomrule\noalign{}
\endlastfoot
Original \(L_0\) and \nolinkurl{AA · L0} & Same stabilizer and same principal congruence kernel, AV5. \\
Finite places above \(2\) & Included by finite étale reductive descent; no unramified-over-\(\mathbb Q_2\) assumption. \\
Proof level \(3M\) & Exact smooth-group index AV8 and Euler cancellation AV10. \\
Haar, compact centre, connected component & Fixed \(SU\), compact-form mass 1, \(K\)-mass 1, same connected domain; AV11--12. \\
Hodge and canonical lines & Physical trace representation and exact \(SU\) determinant identity AV13--15. \\
\(v_A\), very-ample degree and bad-locus degree & Explicit AV3--4 and all multipliers retained in AV19. \\
Same-rule nonempty family intersection, entry \(N_0\) & Remain hypotheses; no source existence from a volume. \\
Full Host/Sizes, fixed coefficient pins, arithmetic Read/032, DP44/RH & Remain open; no sign or nonvanishing inference. \\
\end{longtable}

In particular this is not a transfer of old fixed-\(T\) conditional-kernel constants to a newly chosen admissible Host. It evaluates a geometric input of the same-family conditional escape theorem. It does not establish that an admitted family and a nonzero prescribed rule meet in the required component.

\subsection{8. Primary references and verification boundary}\label{primary-references-and-verification-boundary}

Borel, A., \& Prasad, G. (1989). Finiteness theorems for discrete subgroups of bounded covolume in semi-simple groups. \emph{Publications Mathématiques de l'IHÉS, 69},119--171. \url{https://doi.org/10.1007/BF02698843}. Exact normalization: §4.2,p.137.

Cameron, P. J. (n.d.). \emph{Notes on classical groups}, Chapter5,§5.1,Proposition 5.1 and proof,p.58. \url{https://webspace.maths.qmul.ac.uk/p.j.cameron/class_gps/ch5.pdf}.

Mumford, D. (1977). Hirzebruch's proportionality theorem in the non-compact case. \emph{Inventiones Mathematicae, 42},239--272. \url{https://doi.org/10.1007/BF01389790}. Canonical extensions: Theorem 3.1,pp.256--257; exact volume ratio and sign: Theorem 3.2 and proof,pp.262--263; logarithmic cotangent/canonical pullback: Proposition 3.4 and proof,pp.264--265. Author-hosted primary: \url{https://www.dam.brown.edu/people/mumford/alg_geom/papers/1977b--HirzProp-GT.pdf}.

Prasad, G. (1989). Volumes of S-arithmetic quotients of semi-simple groups. \emph{Publications Mathématiques de l'IHÉS, 69},91--114. \url{https://doi.org/10.1007/BF02698841}. §§3.3--3.6,pp.109--111; Theorem 3.7,p.111; Remark 3.11,pp.112--113; type-A exponents§1.5,p.96.

Smirnov, E. (2015). \emph{Grassmannians, flag varieties, and Gelfand--Zetlin polytopes} (arXiv:1508.03010v1). \url{https://arxiv.org/abs/1508.03010}. §2.9,Theorems 2.31--2.32,pp.18--19.

Tits, J. (1979). Reductive groups over local fields. In \emph{Automorphic forms, representations and L-functions},Part1,Proceedings of Symposia in Pure Mathematics33,29--69. Hyperspecial integral model criterion: §3.8.1,p.55.

The primary statements and their precise normalizations have been read. The companion executable checks actual finite-field and unramified dyadic-ring matrix counts, smooth congruence indices, every ordered local splitting/ramification pattern through degree six, Euler-factor cancellation, independent rectangular-tableau counts, compact-dual signs and determinant characters, the complete integer degree cost, and exact source fingerprints. These finite diagnostics do not replace the universal arguments, Prasad's theorem, proportionality or strong approximation. The final complete saved-file author check, exact computations and independent reviews have separate actual timestamps and subject hashes.

\subsection{9. Algebraic and analytic comparison scope}\label{algebraic-and-analytic-comparison-scope}

The intersection \(v_A\) is the degree of the algebraic first-Chern class of the actual Cartier line on the normal projective compactification. Pullback to a smooth toroidal resolution preserves its top degree by the projection formula. Mumford's good-metric theorem represents the Chern class of the canonical extension by an integrable Chern form; Proposition 3.4 identifies that extension with the algebraic pullback used here. Thus AV18 compares the same algebraic intersection with its compact-dual Chern number, rather than replacing a line by a merely numerical analytic normalization. The finite-place construction uses the reductive model of the unchanged selfdual projective module, and the congruence index is the degree of its finite level quotient. The existing physical trace-one line, rigidification, cubical/Poincare transport and complete source field remain those of LIFT085--096. No new motive, weight, formal slope, crystalline comparison or global nonvanishing statement follows from this volume identity.

\subsection{Canonical algebraic-geometric context}\label{canonical-algebraic-geometric-context-63}

The calculation concerns the same smooth reductive unitary group scheme over the integer ring, obtained by finite-etale descent, its smooth congruence reductions, the actual Picard/Hodge bundle and the top intersection of its canonical extension on the normal arithmetic compactification. Exact fresh primary locators: Prasad(1989), Sections1.5 and3.3-3.7 pp96,109-111 and Remark3.11 pp112-113; Borel-Prasad(1989), Section4.2 p137; Mumford(1977), Theorem3.2 and proof pp262-263 and Proposition3.4(a,b) and proof pp264-265; Tits(1979), Section3.8.1 p55; Cameron, Section5.1 Proposition5.1 and proof p58; Smirnov(2015), Section2.9 Theorems2.31-2.32 pp18-19. Rapinchuk(2013), Theorem2.3 p281 and Milne(2017), Theorem4.16 p48 retain the accepted LIFT095 checks. The original OE action, determinant line, Picard data and Poincare/cubical-line transports are unchanged. Carried contextual anchors: FGA232 Theorem3.1 pp148-149; FGA236 Corollaries1.7/2.3 pp229/232; SGA7 I VII2.7/2.9 pp169-174 and VIII7.4(b) pp310-311; SGA3 I4.7.3/4.7.3.1; SGA6 VI Theorem3.3; cote4 scan003 sections8.1.1-8.1.2 and scans005-006 sections8.4.1-8.4.3, cote54 scans005-007. These are contextual uses, not newly reread Grothendieck primary proofs. No new lambda-ring/K-theory equality, l-adic weight, Dieudonne-crystal or CM-motive comparison follows. DP\_MULT Q3e/DP44/DP52/DP54-65 and THETA3 G1/A7 receive a conditional explicit arithmetic bound with the volume parameter eliminated; CX\_GRH CX-ALG-024 and GCT\_PLETH T-052 retain their arithmetic/global gates.

\clearpage

\section{AM-K-116 --- reciprocal-energy shifted mixture in the actual support27 count}\label{am-k-116-reciprocal-energy-shifted-mixture-in-the-actual-support27-count}

\textbf{Canonical checked integration --- 2026-11-23T09:35:39.990812+00:00.} Complete root and independent proof/script/interface reviews and separate exact replays passed. The actual canonical predecessor has128/31 entries; this integration gives129/31, complete floors27/11. Frozen author history is preserved. AM-K-116 RE1-28 retains the actual reciprocal-energy factor in a shifted positive half-Gamma convolution. The infinite total mass is explicit: every positive inverse moment is bounded by a convergent Gamma/Beta formula, and the unique polynomial degree0 term is made integrable by reserving theta=2\^{}-20 of the actual product-energy exponential before lattice enlargement. Every scalar/transverse term, original saturated metric and covolume, full167-entry reconstruction, central affine750 cover, all168 T factors and all paired four-state G patterns remain. The complete total-eligibility27,b=j\_int=0,2\^{}2111\textless x\textless2\^{}3328,k110 cell has full exponent-1539841/3072\textless-501. One501 append changes actual K115128/31 to129/31, retaining every old overlapping charge; both full sums remain below2\^{}-230 and complete floors27/11 stay. With R097/K115 the remaining actual support27 set is contained in the full-positive scalar-cross globally rank-two-unipotent all-II class with either total eligibility27,111\textless=k\textless=164,2\^{}2111\textless x\textless2\^{}3328 or total eligibility28,k164,2\^{}3107\textless x\textless2\^{}3328. The k111 positive sufficient envelope5269759/3072 is not a lower bound or impossibility theorem. No complete support28, signed nonvanishing, admitted source, selected arithmetic content, prescribed period, specific032 value, evaluated DP44 or RH conclusion.

Author: \nolinkurl{/root/kummer_extension}, 23 November2026 UTC. Controlled pending proof. Whole saved-proof/script/source checks, exact replay, independent review and canonical promotion are separate events. No canonical file is changed.

\subsection{1. Exact remaining sector and claim}\label{exact-remaining-sector-and-claim}

Retain the original conditional J56 all-top H133 kernel, E/F of degrees336/168, all paired source/target pencils and admitted labels, central representative, finite projectors, actual selected vector-valued Hilbert metric and both compact factors, near point, canonical twists, phases, class masses and Haar measures. Work in the canonical K113/K115 remaining all-II globally rank-two-unipotent scalar-cross support27 branch. Its total eligibility is27, its inactive paired-C-zero count b and active integral-C count j\_int are zero, and \[
 2^{2111}<x=N_E D_\eta<2^{3328},\qquad110\le k\le164.
 \tag{RE1}
\] Here k counts the nonzero original S covector values. The total-eligibility28 alternative is separate.

\textbf{Theorem.} The complete k=110 cell in RE1 has absolute normalized contribution less than2\^{}-501. Append501 once to the actual canonical error list, retaining every older overlapping bound. Consequently the total-eligibility27 branch can be restricted to111\textless=k\textless=164. This is not a complete support28 theorem and gives no conclusion for the other eligibility alternative or for the actual signed sum.

Use the unchanged constants \[
 f=168,\quad H=133,\quad R=17424/133<132,\quad
 Q=1427908193729121522601,\quad U=77Q^2-14,
\] \[
 D_0=7^{140}29^{144}5^{84}13^{84},\quad D_F=2^{168}D_0,
 \quad m=83,\quad\beta=7/2,\quad\theta=2^{-20}.
 \tag{RE2}
\] For the application set d=k+3=113 and sigma=(d-m-1/2)/R=7847/34848. The following positive-integral argument is stated with d\textgreater m+1/2, before its explicit k110 use.

The genuine selected coefficient is bounded by \[
 |c(h)|<2\vartheta_a^f2^{-933f/2}
                  \prod_\tau\alpha_\tau^{-1/2}\Psi(\alpha_\tau),
 \quad\Psi(t)=\left(\frac2{\sqrt{1+t}+1}\right)^{333/2},
\] \[
 0<\alpha_\tau\le R,\quad\vartheta_a^f<2^{-649},\quad
 \prod_\tau\alpha_\tau\ge\alpha_0^f,\qquad
 \alpha_0=Ukz^2=Rk/q,
 \quad q=R/(Uz^2).
 \tag{RE3}
\] Here alpha=(N\^{}2)\_\{nn\} and alpha\_tau=U r\_tau J\_Z\^{}\{-1\}r\_tau\^{}* is the true positive row energy at that embedding. These are the original K107/K108/K113 norms and all-k product floor. They apply to actual nonzero rows, before any Gaussian enlargement. Unlike K115, the reciprocal square root in RE3 will remain inside the integral. It is not also replaced by the old uniform alpha0\^{}(-f/2) prefactor.

\subsection{2. A positive convolution with finite positive-degree moments}\label{a-positive-convolution-with-finite-positive-degree-moments}

Canonical K115 GM4--7 proves on0\textless=t\textless=132 that \[
 \Psi(t)\le(1+t/\beta)^{-m},\qquad
 d\gamma(b)=\frac{\beta^m}{\Gamma(m)}b^{m-1}e^{-\beta b}\,db,
 \quad\int d\gamma=1.
 \tag{RE4}
\] For t\textgreater0, the positive half-Gamma integral is \[
 t^{-1/2}=\frac1{\sqrt\pi}\int_0^\infty s^{-1/2}e^{-st}\,ds.
 \tag{RE5}
\] Let mu\_sigma be the positive pushforward of dgamma(b) ds/sqrt(pi s) under c=sigma+b+s. It is locally finite and sigma-finite, supported on c\textgreater sigma, and has infinite total mass. On the original0\textless t\textless=R, \[
 t^{-1/2}\Psi(t)
 \le e^{\sigma R}\int e^{-ct}\,d\mu_\sigma(c).
 \tag{RE6}
\] Indeed the right side equals e\^{}(sigma(R-t)) t\^{}(-1/2)(1+t/beta)\^{}(-m). The whole half-Gamma tail is retained. The paid product factor is e\^{}(sigma R f).

For each integer j\textgreater=1, Tonelli and the ordinary convergent beta integral give \[
 M_j:=\int c^{-j}\,d\mu_\sigma(c)
 =\frac{\Gamma(j-1/2)}{\Gamma(j)}
                 \int(\sigma+b)^{1/2-j}\,d\gamma(b)<\infty.
 \tag{RE7}
\] The inner s integral uses beta arguments1/2 and j-1/2, both positive. This formula does not assert a finite M0. That single divergent polynomial degree is treated separately in section3.

For1\textless=j\textless=m, dropping the positive shift gives \[
 M_j\le A_j:=
 \frac{\Gamma(j-1/2)\Gamma(m-j+1/2)}{\Gamma(j)\Gamma(m)}
                          \beta^{j-1/2}.
 \tag{RE8}
\] The gamma argument m-j+1/2 is positive, including j=m. For j\textgreater=m+1=84, instead drop e\^{}(-beta b) in RE7 and use a beta integral with arguments m and j-m-1/2. The result is \[
 M_j<A_j:=\beta^m\sigma^{m+1/2-j}
                         \frac{\Gamma(j-m-1/2)}{\Gamma(j)}.
 \tag{RE9}
\] Every argument is positive. No analytic continuation, divergent unshifted high moment, or unproved probability normalization for mu\_sigma is used.

For exact arithmetic put H\_n=(2n)!/(4\^{}n n!) and pi\_+=22/7. From Gamma(n+1/2)=sqrt(pi) H\_n, define rational bounds P\_j\textgreater=M\_j\^{}f by \[
 P_j=(\pi_+^2\beta)^{84}
 \left[\frac{\beta^{j-1}H_{j-1}H_{m-j}}
 {(j-1)!(m-1)!}\right]^{168}\quad(1\le j\le83),
\] \[
 P_j=(\pi_+\sigma)^{84}
 \left[\frac{\beta^{83}\sigma^{83-j}H_{j-84}}{(j-1)!}
 \right]^{168}\quad(j\ge84).
 \tag{RE10}
\] The square roots disappear because f=168 is even. These are exactly the two branches used by the executable; it does not evaluate an approximate special function.

\subsection{3. The zero-degree term is made integrable before enlargement}\label{the-zero-degree-term-is-made-integrable-before-enlargement}

For positive mixing weights c\_tau put B=(product c\_tau)\^{}(1/f). On the actual rows RE3 gives, by AM-GM, \[
 \sum_\tau c_\tau\alpha_\tau
       \ge f\left(\prod_\tau c_\tau\alpha_\tau\right)^{1/f}
       \ge f\alpha_0 B.
 \tag{RE11}
\] Thus reserve a fraction theta of the actual exponential: \[
 e^{-\sum c_\tau\alpha_\tau}
 \le e^{-\theta f\alpha_0B}
                   e^{-(1-\theta)\sum c_\tau\alpha_\tau}.
 \tag{RE12}
\] Only the second factor is subsequently summed over the enlarged full row lattice. This reservation needs a product-energy floor, not a separate lower bound on each alpha\_tau. No energy bound is assigned to the extra tuples admitted by that enlargement.

For integer j\textgreater=1 the damping in RE12 can simply be bounded by1, giving the finite integral M\_j\^{}f.~For degree zero use y e\^{}(-y)\textless=e\^{}(-1)\textless1 for every y\textgreater0, hence \[
 e^{-\theta f\alpha_0B}\le(\theta\alpha_0B)^{-f}.
\] Tonelli now gives the finite bound \[
 \int e^{-\theta f\alpha_0B}\prod_\tau d\mu_\sigma(c_\tau)
       \le\left(\frac{M_1}{\theta\alpha_0}\right)^f
       \le P_1\left(\frac{q}{\theta Rk}\right)^f.
 \tag{RE13}
\] In particular the infinite mass of mu\_sigma never appears as an unpaid constant. The last bound retains the whole mixing region, including every anisotropic tail.

\subsection{4. Full original scalar and transverse polynomial}\label{full-original-scalar-and-transverse-polynomial}

Use the actual common row ideal L=J I\_Sigma D\_eta, NE L=z\^{}(2f). The saturated original transverse Gram G' has det G'=U and G'\textgreater c0 I for c0=2\^{}36. Its three values belong to the actual ideal because the input basis is integral; no lattice value has been divided by2. The scalar half-trace real dimension is336. K111 HP12--18 supplies, for every positive Gaussian parameter and every lattice minimum, the unchanged scalar coefficient \[
 c_{336}=168{336\choose168}/4^{168},\qquad
 M_{15}=\frac{6015684926187167549292201834533348561577984}
 {2618017117694612446410293519},
\] \[
 K=\frac{336M_{15}}{c_{336}}(7841/5000)^{168}168!
       \left[\frac{(365729/500)^2}{168(314159/100000)^2}\right]^{168}.
 \tag{RE14}
\] The original K110 PM10--12 anisotropic minimum and true transverse covolume apply to the weights (1-theta)c\_tau. Write t0=Uz²=R/q. The complete lattice Gaussian is bounded by \[
 {164\choose k}\frac{U^f\pi^{df}}{D_0^d}
 \left\{((1-\theta)t_0B)^{-f}+K\right\}^{k}
 ((1-\theta)t_0)^{-3f}B^{-3f}
 [1+Kc_0^{-f}((1-\theta)t_0)^fB^f]^3.
 \tag{RE15}
\] All original coordinate-position choices are retained. The product dual bound uses the true G', while the primal covolume keeps its determinant U. The universal scalar theorem is applied to each anisotropically transformed real lattice, not to a fictitious one-dimensional theta series.

Expand with l scalar volume factors and h transverse corrections, where0\textless=l\textless=k and0\textless=h\textless=3. The inverse moment degree is j=l+3-h. Except at the unique pair(l,h)=(0,3), it is an integer in1..d. Combining RE10 and RE13 gives the following full weighted row majorant: \[
 e^{\sigma Rf}{164\choose k}\frac{U^f\pi^{df}}{D_0^d}
                          \mathcal P_k(q),
 \tag{RE16}
\] where \[
 \begin{split}
 \mathcal P_k(q)={}&
 \sum_{\substack{0\le l\le k,\ 0\le h\le3\\j=l+3-h\ge1}}
 {k\choose l}{3\choose h}K^{k-l+h}c_0^{-hf}
 \left[\frac{q}{(1-\theta)R}\right]^{jf}P_j\\
 &+K^{k+3}c_0^{-3f}P_1\left[\frac{q}{\theta Rk}\right]^f.
 \end{split}
 \tag{RE17}
\] The second line is the full zero-degree contribution, not a discarded correction. Every term is positive. No unweighted Holder claim for the infinite measure is made. The original trace was used in RE3/RE6/RE11 before the only enlargement, RE15.

\subsection{5. All arithmetic factors and both continuum directions}\label{all-arithmetic-factors-and-both-continuum-directions}

The exact arithmetic inputs remain K113 SC27.4--10 and canonical K115 GM15--19. For the bounded window in RE1, eta is unique per fixed27-place support. Its original covector and central scalar reconstruct the whole matrix. All27 active weighted and141 inactive pointwise factors give(T-1)\^{}(-168)\textless2T\^{}(-168). The actual Green and central bounds are \[
 |a_G(h)|\le2Q^{-335\cdot167}(N_EJ)^{335/2},\quad
 N_EJ=z^{2f}T^{54}/x\le Q^{334},
\] \[
 \#d_{\rm cent}\le
 \frac{V_F(C(q))U^{-f}T^{54}x^{-1}}{N_FJ_0},\quad
 C(q)=H/2+750R/q,
\] \[
 V_F(C)=\frac{((22/7)fC^2)^{84}}{84!\sqrt{D_F}},\quad
                 N_EJ\le Q^{167}N_FJ_0.
 \tag{RE18}
\] The original affine shift, full750 covering radius, dyadic conditions, phases, source reconstruction and both mixed contraction factors are unchanged. There is no extra eta or norm-generator enumeration, and no missing target/source fibre. Support\textless2\^{}104, pivot\textless2\^{}8, the Green correction2, T correction2, near-point factor2 and conservative pattern factor4 give fixed exponent117.

There is no longer an external alpha0\^{}(-f/2), or its k\^{}(-84) factor. The z power outside RE17 is335f, so before central correlation each positive-degree term has q dependence \[
 F_j(q)=q^{(j-335/2)f}V_F(C(q)).
 \tag{RE19}
\] The repaired zero-degree term has this same form with effective degree j\_eff=1. Put epsilon=(C(q)-H/2)/C(q) in(0,1). For1\textless=j\textless=113, \[
 \frac1f\frac{d\log F_j}{d\log q}=j-335/2-\epsilon<0,
\] \[
 \frac{d\log\{x^{-337/2}F_j(q(x))\}}{d\log x}
                              =-1-j+\epsilon<0.
 \tag{RE20}
\] All mixing parameters and moment bounds depend only on k, not on x or J. This checks the entire denominator continuum term by term, including the repaired term, before the finite-pattern sum.

For fixed x let \[
 q_*^f=\frac{R^fT^{54}}{U^fQ^{334}x},\quad
 D=Q^{334}/N_EJ,\quad q=q_*D^{1/f},\quad
 (N_FJ_0)^{-1}\le Q^{-167}D.
 \tag{RE21}
\] The four local states give the complete product \[
 \prod_{\mathfrak q\mid G}(1+q_{\mathfrak q}^{j-333/2})^2.
 \tag{RE22}
\] At every j\textless=113 its exponent is at most-107/2. The repaired term uses j\_eff=1. The actual inequality(1+g0\^{}(-1))\^{}668\textless2 therefore bounds every pattern product; the retained4 remains conservative. This includes all1,2,1 local multiplicities and every common orientation. No per-coordinate pattern or preferred identity pattern is substituted.

\subsection{6. Exact complete endpoint and every polynomial term}\label{exact-complete-endpoint-and-every-polynomial-term}

By RE20 it suffices to use x=2\^{}2111. The unchanged exact cap bounds and full central volume are \[
 q_-=251124964<q_*<q_+=251124965,\qquad
           V_F(H/2+750R/q_-)<2^{474}.
 \tag{RE23}
\] The full rational directed common exponent is \[
 \begin{split}
 C={}&117-933f/2-649+474+(335f/2)L_R-(335f/2)L_U\\
    &+8931L_T-336\cdot167L_Q-(337/2)2111
       =593387605/1024,
 \end{split}
 \tag{RE24}
\] where LT=1019333/1024, LQ=287843/4096, LU=601356/4096 and LR=28810/4096 have the same upper/lower directions as K115. The T exponent is335·27+54−168. The new R exponent335f/2 reflects the reciprocal norm's movement inside the integral; the U exponent remains−335f/2.

At k110,d113 use sigma=7847/34848 in every P\_j. Define \[
 V=K[((1-\theta)R)/q_+]^f,\qquad V<2^{-935},
\] \[
 \mathcal M=[q_+/((1-\theta)R)]^{df}P_d.
 \tag{RE25}
\] For a nonmain positive-degree term the exact ratio to M is \[
 {k\choose l}{3\choose h}V^{d-j}c_0^{-hf}P_j/P_d.
\] For the repaired zero-degree term the ratio is \[
 V^d c_0^{-3f}[q_+/(\theta Rk)]^f P_1/P_d.
 \tag{RE26}
\] The executable uses least strict binary ceilings of the exact rational factors. It checks every443 nonmain term; their largest certified upper exponent is−957. In particular the complete nonmain sum is less than443·2\^{}-100\textless1 times M. Keeping the conservative factor2 for the entire polynomial gives \[
 W=\frac{{164\choose110}(22/7)^{113f}}{D_0^{113}q_-^{335f/2}}
                  2\mathcal M,
       \qquad W<2^{-587140}.
 \tag{RE27}
\] Both half-integer moments and all polynomial corrections are inside this rational comparison.

The entire paid shift is \[
 e^{\sigma Rf}=e^{(59/2)f}<2^{13(59/2)f/9}=2^{21476/3}.
\] Consequently the full absolute exponent of the k110 cell is \[
 C+21476/3-587140=-1539841/3072<-501.
 \tag{RE28}
\] This proves the theorem with one501 charge. There is no further coordinate-size union; the complete binomial and every original support/pattern/fibre factor already occur in the cell. The same sufficient envelope at k111 has positive exponent5269759/3072. That failed upper bound supplies no lower bound, existence claim or obstruction to a different estimate.

\subsection{7. Exact checks, source boundaries and remaining gate}\label{exact-checks-source-boundaries-and-remaining-gate}

The actual current predecessor is checks/k115\_joint\_remainder\_error.json,128 all-top and31 uniform entries with floors27/11. Append501 once and retain the entire chronological prefix and all uniform entries. The executable accepts a later full ordered extension without overwriting historical outputs. Both full dyadic sums remain below2\^{}-230; the complete floors stay27/11.

The certificate checks the original host and exact logarithm directions, the physical cap, bounded eta uniqueness, original support/pivot/all168T factors, both cap roots, central volume, all half-integer gamma recurrences and moment degrees, all positive polynomial pairs including the unique zero degree, both continuum signs, the full four-state pattern inequality, all443 relative terms, the complete endpoint and exact ledger sums. The universal proofs of Tonelli, the original product-energy application and finite moment identities are explicit above; finite diagnostics do not replace them.

Primary normalization: National Institute of Standards and Technology. (n.d.). \emph{Digital Library of Mathematical Functions}, equations5.9.1,5.12.1/5.12.3 and5.4.6, \url{https://dlmf.nist.gov/5.9.E1}, \url{https://dlmf.nist.gov/5.12.E1}, \url{https://dlmf.nist.gov/5.4.E6}. These positive-real Gamma/Beta domains and Gamma(1/2)=sqrt(pi) were directly read for this proof. All beta arguments used above are positive. The half-integer recurrence also follows immediately by integration by parts. No analytic continuation is needed.

Direct inherited mathematical interfaces are canonical K115 GM1--23; K113 SC27.1--10; K110 PM10--12 at its first untruncated all-positive-weight scope; K111 HP12--18; K109 TG3--12; K108 AE11--21; and K107's actual saturated membership, all-k norm floor and full near-point/tail factor. No new representation, selected module, arithmetic correspondence or source family is substituted.

Carried geometric locators remain SGA3 I4.7.3/4.7.3.1; SGA6 VI Theorem3.3; FGA232 Theorem3.1 pp148--149; FGA236 Corollaries1.7/2.3 pp229/232; SGA7 I VII2.7/2.9 pp169--174 and VIII7.4(b) pp310--311; cote4 scan003 sections8.1.1--8.1.2 and scans005--006 sections8.4.1--8.4.3; and cote54 scans005--007. These are inherited context. DP\_MULT Q3e/DP44, THETA3 G1/A7, CX\_GRH CX-ALG-024 and GCT\_PLETH T-052 keep their signed, source, selected arithmetic-content, prescribed-period and global comparison gates. No full support28 assertion, Lean build or archive audit is made.

\subsection{Canonical algebraic-geometric context}\label{canonical-algebraic-geometric-context-64}

The positive infinite measure is an analytic majorant of the same selected real realization. Its reserved damping is imposed on the actual original rows before enlargement; it introduces no arithmetic parameter, field extension, physical source vector, lattice rescaling or projector. All original projective OE modules, paired fractional ideals, CM involution, polarization, theta line, Picard and Poincare data remain. Fresh primary normalization: NIST DLMF equations5.9.1,5.12.1/5.12.3 and5.4.6, with their positive-real domains; exact URLs and separate root/peer reading scopes are preserved. The universal moment and zero-degree arguments are proved fully. Original representation and harmonic-code inputs retain their checked K107/K108/K110/K111/K113/K115 scopes. Carried exact anchors: SGA3 I4.7.3/4.7.3.1; SGA6 VI Theorem3.3; FGA232 Theorem3.1 pp148-149; FGA236 Corollaries1.7/2.3 pp229/232; SGA7 I VII2.7/2.9 pp169-174 and VIII7.4(b) pp310-311; cote4 scan003 sections8.1.1-8.1.2 and scans005-006 sections8.4.1-8.4.3; cote54 scans005-007. No new formal-group/Dieudonne, lambda-ring/K-theory, l-adic-weight or CM-motive comparison. DP\_MULT Q3e/DP44 and THETA3 G1/A7 receive the smaller conditional finite gate; CX\_GRH CX-ALG-024 and GCT\_PLETH T-052 retain their arithmetic/global gates.

\clearpage

\section{AM-LIFT-099 --- Bounded ray-class entry into the same-rule nonzero family}\label{am-lift-099-bounded-ray-class-entry-into-the-same-rule-nonzero-family}

\textbf{Canonical checked integration --- 2026-11-23T09:43:36.940264+00:00.} Complete saved-proof and executable reviews and separate exact diagnostic replays passed. No numerical error charge is introduced. Frozen author history is preserved. AM-LIFT-099 RE1-26 supplies bounded component entry and effective same-rule escape for a supplied actual literal admitted Host/Datum/Level/Moduli/Coeff with A=1, fixed CM seed, relative-unramified E/F, ordinary selfdual IsDatW ideal multiple, and uniform indefinite rank r=p+q\textgreater=3. Let f={[}F:Q{]}, g=fr, d=fpq, n=3M, m=3M Q\_S, D\_E=D\_F\^{}2. Exact determinant congruence images and SU strong approximation identify the unitary finite double quotient with the norm-one determinant quotient; Cl\_E(m) surjects onto it. Half-trace duality and all-shape transference give a prime-to-m representative B of every ray class with N\_E B\textless=(2m)\^{}(2f)D\_E. The actual lattice B e+bar(B)\^{}(-1)f+Lprime, multiplied by the original AA and retaining J0, admits an original-level entry isogeny with N0=1 mod m and N0\textless=C\_ent=m(2m)\^{}(2f)D\_F\^{}2. All such bounded families cover the literal moduli complex points, so one meets the same supplied nonzero-rule open. With t=1+m, delta=t\^{}(4fr\^{}2), k=gw+n\_F/2 for even n\_F or k=2gw+n\_F after the same-zero square, Vcal=(3M)\^{}(f(r\^{}2-1))D\_F\^{}ceil((r\^{}2-1)/2)2\^{}(f(r-1)){[}2(g+1){]}\^{}d d!(product\_(j=1)\^{}(r-1)j!)\^{}f, h\_d=2\^{}(d+1)(d+1)(d+2)!, e\_g=2\^{}(8g\^{}2+2g+4), and D\_\emph{=e\_g h\_d\^{}d Vcal sum\_(j=0)\^{}d k\^{}j, the entire scale satisfies N=N0 t\^{}ell\textless=C\_ent t\^{}\{B\_delta,d-1(D\_})\}, B\_delta,0(D)=D and B\_delta,s(D)=D+B\_delta,s-1(2 delta\^{}D D\^{}2). The exact action, principal trace-one cubical line, finite eta and nonzero literal frame descend over the original K0(A0{[}N\^{}2{]}); generic relative degree\textless=N\^{}(8g\^{}2). Under precisely the original J,C0,Kgood same-seed hypotheses, {[}Kout:Bcoef{]}\textless64512 d\_CM M\^{}4 N\^{}(81/8) and rd(Kout)\textless2\^{}50 exp(2pi sqrt(d\_CM))rd(Bcoef) M N\^{}2, with the full composed N and actual coefficient-field degree/discriminant retained. No ray-class or proof-cover field is adjoined. Selected arithmetic-rule nonzero, actual Host/Sizes admission, remaining analytic/model/section/fibre Step5 inputs, Read/specific032, evaluated DP44 and RH remain open. Definite terminal signatures and rank below3 are excluded. No finite error charge; complete129/31 arrays and floors27/11 unchanged.

Author: /root/lift\_accounting. 23 November 2026 UTC.

\textbf{Complete author subject; hash-specific saved-file, source and executable verification are recorded separately. Independent review and canonical integration are separate.} This theorem is conditional on an existing literal admitted source and a supplied nonzero rule. It gives a bounded entry isogeny into an arithmetic family on which that same rule is nonzero. It does not construct the selected arithmetic rule or the admitted Host/Sizes instance. No canonical change or finite-kernel charge is made.

\subsection{1. The fixed source and the assertion}\label{the-fixed-source-and-the-assertion}

Retain the literal Host, Datum, Level, Moduli and Coeff of LIFT085, with source level-record parameter A=1 and physical level M. Let E/F be the actual CM quadratic extension, unramified at every finite place, with {[}F:Q{]}=f.~Assume the nonterminal uniform signature (p,q), p,q\textgreater=1, r=p+q\textgreater=3, and put g=fr and d=fpq. Retain the actual IsDatW lattice description AA·L0, where AA is an invertible fractional O\_E-ideal and L0 is full and ordinarily Hermitian-selfdual. All coefficient labels, the trace alternating form and its sign, and the finite eta are fixed.

Fix an actual CM seed Obj A0 over K0 containing the original coefficient field, with full action, principal trace-one line \(\mathcal L_0\), cubical rigidification and original finite eta0. Let F0 be the same given nonzero natural rule in sec(a,n\_F). Let S be a finite set of protected rational primes containing those dividing M. Write

\[
 Q_S=\prod_{p\in S}p,\qquad n=3M,\qquad m=3M Q_S,\qquad
 D_E=|\operatorname{disc}(E)|=D_F^2.
 \tag{RE1}
\]

There is an O\_E-equivariant entry isogeny u0:A0-\textgreater AB to another actual Obj of precisely these data such that

\[
\begin{gathered}
 u_0^*\mathcal L_B\simeq\mathcal L_0^{N_0^2},\qquad
 \deg u_0=N_0^{2g},\qquad u_0\eta_0=\eta_B,\\
 (N_0,Q_S)=1,\qquad N_0\equiv1\pmod m,\qquad
 1\le N_0\le C_{\rm ent}:=m(2m)^{2f}D_E,
\end{gathered}                                                    \tag{RE2}
\]

with the line isomorphism rigidified and cubical. The CM seed AB belongs to an arithmetic principal-n special-unitary family YB whose regular literal period map PhiB satisfies

\[
                 \Phi_B^{-1}(U_{F_0})\ne\varnothing.       \tag{RE3}
\]

Here UF0 is exactly the nonzero-frame open of LIFT085 HF11. The entry point itself is not asserted to have nonzero rule value. It is the controlled seed in a family where the same-rule escape argument can be applied. The entry object, full action, line and eta descend over K0(A0{[}N0\^{}2{]}); no auxiliary ideal-class field is adjoined.

Combining this entry theorem with LIFT094--098 yields the following explicit bound. Put \(w=\sum_\tau\sum_{i<Dw.mult(\tau)}a(\tau,i)\), the same total physical frame weight as in LIFT096. Use the actual rule exponent

\[
 k=gw+n_F/2\quad(n_F\text{ even}),\qquad
 k=2gw+n_F\quad\text{after squaring in general},
\] \[
\begin{split}
 a_G&=r^2-1,\\
 \mathcal V&=(3M)^{fa_G}D_F^{\lceil a_G/2\rceil}
 2^{f(r-1)}[2(g+1)]^d d!\left(\prod_{j=1}^{r-1}j!\right)^f,\\
 h_d&=2^{d+1}(d+1)(d+2)!,\qquad e_g=2^{8g^2+2g+4},\\
 D_*&=e_gh_d^d\mathcal V\sum_{j=0}^{d}k^j,\qquad
 t=1+m,\qquad\delta=t^{4fr^2}.
\end{split}                                                       \tag{RE4}
\]

For the already proved integer recursion

\[
 B_{\delta,0}(D)=D,\qquad
 B_{\delta,s}(D)=D+B_{\delta,s-1}(2\delta^D D^2),
 \tag{RE5}
\]

there is a final original-level O\_E-isogeny u:A0-\textgreater A1 and a literal field-valued frame y with

\[
\begin{gathered}
 \operatorname{evR}(F_0,y)\ne0,\qquad
 u^*\mathcal L_1\simeq\mathcal L_0^{N^2},\qquad u\eta_0=\eta_1,\\
 (N,Q_S)=1,\quad N\equiv1\pmod M,\quad\deg u=N^{2g},\\
 N\le C_{\rm ent}\,t^{B_{\delta,d-1}(D_*)}.
\end{gathered}                                                    \tag{RE6}
\]

All final data and the nonzero frame are over K0(A0{[}N\^{}2{]}), so the generic field bound is N\^{}(8g\^{}2). The stronger same-J/C0 torsion bounds apply only with their original seed hypotheses. The integer bound RE6 is explicit in the given field, signature, level, rule weights and protected set. It is not a numerical verification of the full Host/Sizes inequalities or of the nonzero selected arithmetic image.

\subsection{2. The actual marked genus and an integral plane in the seed}\label{the-actual-marked-genus-and-an-integral-plane-in-the-seed}

Let Lambda0 be the integral homology of the supplied seed with its exact Riemann form Omega and O\_E action. The argument of LIFT085 HF12--13 applies to every complex point x of the literal moduli space: local selfdual classification, the prescribed signatures and the unitary Hasse principle give an exact E-linear rational isometry from its rational homology to that of the seed. This works in odd and even ranks, and is not the even-rank-only similitude assertion. We use the same fixed extraction of a Hermitian form from Omega as in that proof.

This identification also relates the seed to the given IsDatW rational space. Its lattice is locally isometric to AA·L0 at every finite place. Therefore

\[
 L=\mathrm{AA}^{-1}\Lambda_0
 \quad\text{is a full ordinary selfdual Hermitian lattice},\qquad
 \Lambda_0=\mathrm{AA}\,L.                                  \tag{RE7}
\]

Indeed each completion of L is isometric to the completion of L0, because the local isometry is E-linear and commutes with the ideal AA. Duality commutes with completion, and equality of full Dedekind lattices can be tested at every finite place. This does not assert that the seed lattice and the prescribed datum lattice are globally integrally isometric.

LIFT095 gives, inside this same L, an integral hyperbolic decomposition

\[
 L=(\mathcal O_E e\oplus\mathcal O_E f)\perp L',\qquad
 h(e,e)=h(f,f)=0,\quad h(e,f)=1.                              \tag{RE8}
\]

The rank r\textgreater=3, finite relative unramifiedness, and all real indefinite signatures are exactly its hypotheses. No height bound for the rational coordinates of e,f is required below: all denominator bounds will be ideal inclusions in this unchanged lattice, not bounds on a matrix in an unrelated basis. The orthogonal complement L' may be projective and need not be free.

Write U=U(h), T=ker(Res\_(E/F)G\_m-\textgreater G\_m), with determinant U-\textgreater T. Let K(n) be the compact open subgroup of U(A\_F,f) preserving Lambda0 locally and acting as identity modulo n Lambda0. Multiplication by the invertible ideal AA cancels in both conditions, so K(n) is also the corresponding subgroup for L.

\subsection{3. The exact determinant quotient}\label{the-exact-determinant-quotient}

The determinant maps U(F)-\textgreater T(F) and U(F\_v)-\textgreater T(F\_v) are surjective. To see the rational assertion, choose a nonisotropic E-line in V and act by z in T(F) on that line and by identity on its orthogonal complement. The determinant is z. The same proof works locally.

There is also the exact integral statement

\[
 \det K(n)=T(n):=
 \prod_{v\nmid n}\mathcal O_{E_v}^{\,1}\times
 \prod_{v\mid n}\{z\in\mathcal O_{E_v}^{\,1}:z\equiv1\bmod n\mathcal O_{E_v}\}.
                                                               \tag{RE9}
\]

For a relatively unramified nonsplit place, the local selfdual proof in LIFT095 §2 splits off a vector of unit norm, including residue characteristic two. Its orthogonal projection is integral. Acting by z on that line fixes the lattice; when z-1 belongs to n O\_E,v its difference from identity belongs to n End(L\_v). This proves surjectivity at every congruence depth. At a split place choose a line in one free factor and its dual in the other; the same construction has determinant z and its inverse-adjoint partner. The reverse inclusion in RE9 follows directly by taking the determinant of a matrix congruent to identity. No residue-characteristic exception is omitted.

Strong approximation for SU(h), in the exact scope already checked in LIFT095, proves a bijection

\[
 U(F)\backslash U(\mathbb A_{F,f})/K(n)
 \ \xrightarrow{\ \det\ }\
 T(F)\backslash T(\mathbb A_{F,f})/T(n).                       \tag{RE10}
\]

Here is the injectivity argument, including the nonnormal compact subgroup issue. If two determinants differ by an element of T(F) and an element of T(n), lift those two elements to U(F) and K(n) using the preceding surjectivities. After adjusting the two representatives, their determinants agree. The resulting difference lies in SU(A\_F,f). The desired equality of double cosets asks for an element of SU(F) in a translate of a conjugate of K(n) intersect SU(A\_F,f), a nonempty open subset. Strong approximation supplies it. This does not require K(n) to be normal in U(A\_F,f). Surjectivity follows from local determinant surjectivity and the integral version at almost all places.

All factors of SU(h) at infinity are noncompact, and U(p,q) is connected. Thus a rational unitary isometry preserves the same positive compatible-complex-structure domain. No additional real sign class is lost in RE10.

\subsection{4. A ray-class surjection with the literal level retained}\label{a-ray-class-surjection-with-the-literal-level-retained}

Use the finite ray modulus m O\_E, with no infinite conditions since E is CM. The map

\[
 b\longmapsto b/\bar b:
 \mathbb A_{E,f}^{\times}\longrightarrow T(\mathbb A_{F,f})
                                                               \tag{RE11}
\]

is surjective. Locally this is Hilbert 90; in the split case it is the explicit map (b1,b2)-\textgreater(b1/b2,b2/b1). At almost all places the target is a norm-one unit and b can be chosen to be a unit: in an unramified field extension, divide a Hilbert-90 solution by a power of a base-field uniformizer, leaving b/bar b unchanged. Consequently the local solutions form an actual restricted-product idele. This includes dyadic unramified places.

Principal ideles map into T(F). At places not dividing m, integral units map into integral norm-one units. At places dividing m, units congruent to one modulo m map into units congruent to one modulo m and hence modulo n.~Thus RE11 induces a surjection

\[
 \operatorname{Cl}_E(m) woheadrightarrow
 T(F)\backslash T(\mathbb A_{F,f})/T(n).                       \tag{RE12}
\]

The idelic/ideal identification is precisely Milne, CFT V, Proposition4.6 and proof, pp.172--173. Every ray class has an integral ideal B prime to m as representative. One can first take a fractional ideal prime to m and then multiply by an integer congruent to one modulo m that clears its denominators. Such an integer exists by the ordinary Chinese remainder theorem, since these denominators are prime to m. The next section replaces this initially unbounded representative by a bounded one in exactly the same ray class.

There is no assertion that the ray-class map is injective, or that the unitary genus has one class. Its surjectivity is sufficient. Its modulus covers the full proof level n=3M and every protected prime. In particular it also controls the original half-level eta, whose two split halves are related by the actual perfect pairing.

\subsection{5. A small ideal in every ray class, with all ideal shapes allowed}\label{a-small-ideal-in-every-ray-class-with-all-ideal-shapes-allowed}

The following estimate is valid for any CM field E of degree 2f and any integer m\textgreater=2; finite relative unramifiedness is not needed in this section. Equip E\_R with the half-trace Euclidean metric

\[
 \|x\|^2=\tfrac12\operatorname{Tr}_{E/\mathbb Q}(x\bar x)
       =\sum_{\sigma\in\Phi}|\sigma x|^2.
                                                               \tag{RE13}
\]

For an arbitrary nonzero fractional ideal I, its Euclidean dual is

\[
 I^*=2\mathfrak D_E^{-1}\bar I^{-1},\qquad
 N_E(I^*)=2^{2f}D_E^{-1}N_E(I)^{-1}.                           \tag{RE14}
\]

Indeed the inner-product integrality condition is exactly the full trace-dual condition after multiplication by 1/2 and conjugation. If x is a nonzero element of I\emph{, the inclusion (x) subset I} gives \textbar N\_E x\textbar\textgreater=N\_E(I*). Arithmetic--geometric mean on the f complex absolute squares therefore gives

\[
 \lambda_1(I^*)\ge
 2\sqrt f\,D_E^{-1/(2f)}N_E(I)^{-1/(2f)}.
                                                               \tag{RE15}
\]

Miller and Stephens-Davidowitz (2019), Theorem3.2 and equation(3.3), printed p.8, give for every real dimension s

\[
 \mu(\mathcal L)\lambda_1(\mathcal L^*)
 \le s/(2\pi)+3\sqrt{s}/\pi<2s.
\]

Apply this with s=2f and RE15. For every ideal shape,

\[
 \mu(I)<2\sqrt f\,D_E^{1/(2f)}N_E(I)^{1/(2f)}.                \tag{RE16}
\]

The coarse strict inequality uses only pi\textgreater3 and sqrt(s)\textless=s. There is no class-number, regulator, basis-height or principality hypothesis.

Start with an integral representative B prime to m of a ray class and put I=B\^{}-1, so O\_E subset I. The coset 1+mI contains an element a of norm at most mu(mI)=m mu(I): choose a closest lattice point to -1. At every prime dividing m, I\_v=O\_E,v, hence a is congruent to one modulo m there. In particular a is nonzero; zero cannot lie in this coset. Moreover a belongs to I, so Bnew=aB is an integral ideal prime to m in precisely the same ray class. By RE16 and arithmetic--geometric mean,

\[
\begin{split}
 N_E(B_{\rm new})
 &=|N_Ea|N_EB
 \le (\|a\|/\sqrt f)^{2f}N_EB\\
 &<(2m)^{2f}D_E N_E(I)N_EB
 =(2m)^{2f}D_E.
\end{split}                                                       \tag{RE17}
\]

Thus every ray class has a representative B with

\[
            (B,m)=1,\qquad N_E B\le (2m)^{2f}D_E.           \tag{RE18}
\]

This is a bound for a representative of each class, not a density estimate for ideals in a coset. The zero exclusion occurs before the norm estimate. The same proof applies to arbitrary inverse fractional ideals I and does not extrapolate a ring-only covering estimate.

\subsection{6. The actual bounded entry lattice and physical isogeny}\label{the-actual-bounded-entry-lattice-and-physical-isogeny}

For each bounded ideal B in RE18 define, in the same rational space and hyperbolic plane RE8,

\[
 L_B=(B e\oplus\bar B^{-1}f)\perp L',\qquad
 \Lambda_B=\mathrm{AA}\,L_B.                                 \tag{RE19}
\]

This lattice is ordinarily selfdual: the two ideal coefficients are exact Hermitian duals because h(e,f)=1. At every place dividing m it agrees with L and hence LambdaB agrees with Lambda0. At any finite place choose b\_v generating B\_v, taking b\_v=1 at v\textbar m. The unitary map diag(b\_v,bar(b\_v)\^{}-1,1) on the hyperbolic plane and its complement sends L\_v to LB,v and has determinant b\_v/bar b\_v. Its finite adelic class is exactly the image of the ideal class under RE11--12. Unit changes of the local generators lie in the integral stabilizer and do not alter this class.

Give V\_R/LambdaB the original seed complex structure J0. The same positive Riemann form is integral and unimodular, since LambdaB is locally isometric to Lambda0. It defines a principal polarization. The trace-one construction of LIFT085 HF5--7 supplies the exact symmetric rigidified cubical line. Full O\_E action and all signature clauses remain unchanged. The finite etaB is the canonical transport of eta0 through the equality of the local lattices at M; the two finite M-torsion modules are canonically identified by those completions. The single finite adelic isometry \(g_B\) transports the complete profinite lattice and its pairing, hence the entire coherent family of all-n TateIso maps and the same primitive-root tower; it is not a separate unrelated choice at each n.~At M it agrees with the finite eta identification. This proves the geometric TateIso clauses without a full infinite-level trivialization over a number field. LIFT085 HF15--16 supplies the genuine regular period family for LambdaB and its actual Obj structures, rather than only an analytic torus construction.

Let b=N\_EB. Since O\_E/B has cardinality b, multiplication by b annihilates it; therefore b belongs to B. The same holds for bar B. Directly from RE19,

\[
 b\Lambda_0\subset\Lambda_B,\qquad
 b\Lambda_B\subset\Lambda_0.                                \tag{RE20}
\]

These are inclusions of the actual projective ideal multiples, with no coordinate-height factor. Since b is coprime to m, choose the integer c with 1\textless=c\textless=m and cb congruent to one modulo m. Put N0=cb. The map N0 times the identity of V\_R is complex-linear for J0, integral from Lambda0 to LambdaB, O\_E-linear, and has polarization multiplier N0\^{}2. Its exact eta action is the prescribed one because N0=1 mod M. The trace-one semicharacter or normalized Poincare formula gives its rigidified cubical line-power identity. Principality then gives degree N0\^{}(2g).

The second inclusion in RE20, or the polarized adjoint identity, gives

\[
                       \ker u_0\subset A_0[N_0^2].         \tag{RE21}
\]

The standard quotient and the complete all-object descent of LIFT085 §7 place AB, u0, its O\_E action, principal polarization, canonical line and etaB over K0(A0{[}N0\^{}2{]}). No ideal-class field, arbitrary theta-characteristic extension or infinite level field is required. Finally N0\textless=m b and RE18 prove RE2.

\subsection{7. Why these bounded families meet every literal target}\label{why-these-bounded-families-meet-every-literal-target}

Take an arbitrary x in the complex points of the literal moduli scheme. By LIFT085 HF12--13 identify its rational homology exactly with the seed space. Choose local unitary isometries carrying Lambda0 to its integral lattice Lambdax. At primes dividing M choose them to carry the actual eta0 to etax; this is possible by the split half-level argument in HF13--14. Choose arbitrary additional compatible markings at the primes of n=3M. These local isometries form g\_x in U(A\_F,f).

By RE10--12 and the bounded ray representative RE18, there exist a bounded ideal B, gamma in U(F), and k0 in K(n) with

\[
                         g_x=\gamma g_B k_0.                \tag{RE22}
\]

Consequently Lambdax=gamma LambdaB, with the same equality of the actual finite eta at M. Let Jx be its compatible complex structure in this rational identification. The point gamma\^{}-1 Jx gamma of the domain for LambdaB gives an object isomorphic to the original target, including its canonical trace-one line and all source structures. The exact Moduli.rep uniqueness identifies its literal period-map value with x. Thus

\[
 X(\mathbb C)=\bigcup_{B\text{ as in RE18}}\Phi_B(Y_B(\mathbb C)).
                                                               \tag{RE23}
\]

The union is finite, since only finitely many integral ideals have bounded norm. Neither an isomorphism of each YB with a moduli component nor finiteness of PhiB is needed. An unbounded rational isometry gamma is used only to identify a target with a point of the same family; it is not an entry isogeny, and its denominator does not enter RE20--21 or the arithmetic cost.

Because F0 is the same nonzero natural rule, LIFT085 HF9--11 gives a nonempty open UF0 in the literal moduli space. Choose x in that open and apply RE23. For at least one bounded B, the regular pullback PhiB\^{}-1(UF0) is nonempty. This proves RE3. The entry seed AB is a point of that family because it uses the fixed lattice LambdaB and the original J0. This proves the needed controlled seed-to-family entry without claiming that the rule is nonzero on the original seed component.

\subsection{8. Composition with the checked within-family escape and exact scope}\label{composition-with-the-checked-within-family-escape-and-exact-scope}

The family YB is formed with the special-unitary principal-n subgroup for LambdaB. It has the same dimension d, original M-level, rational Hermitian space, signatures and field discriminant. The ordinary lattice LB is again selfdual, so LIFT095 supplies c=1 within it. The exact compactification and branch proofs of LIFT094 apply with t=1+m and delta=t\^{}(4fr\^{}2). Its branch maps have the same polarization multiplier t\^{}2 and preserve the exact original eta, regardless of the chosen integral hyperbolic frame.

LIFT096 applies to the same rule F0 on this family; RE3 makes its bad locus proper. LIFT097 supplies h\_d.~LIFT098's volume estimate applies uniformly to every LB, because it uses the selfdual hyperspecial collection and principal-n index, not a particular global lattice class. Hence RE4 bounds the same-rule bad-locus degree. The LIFT094 recursion gives an escape path from AB of length ell\textless=B\_(delta,d-1)(D\_*), or length zero if its bad locus is empty. Every step is an actual original-level isogeny with positive multiplier t. Composing with u0 gives N=N0 t\^{}ell and proves RE6.

The full composition has kernel in A0{[}N\^{}2{]} by its polarized adjoint. Its final quotient, action, canonical line and finite level descend over K0(A0{[}N\^{}2{]}) by LIFT085. A nonzero literal frame exists over that same field because the eigenframe/theta-generator torsor is a product of general linear groups and G\_m with a nonempty regular-function open. The source FromTest clause and evalAt convention therefore give nonzero evR rather than its default-zero branch. The finite proof markings and theta covers are not adjoined to this output field.

Here are the complete field and discriminant costs when the supplied seed is exactly the same-seed construction of LIFT087--093. This is an additional hypothesis for this corollary, not an assertion about an arbitrary CM seed in section1. Write Bcoef=c.F, let d\_CM be the actual imaginary-quadratic Host prime (distinct notation from the dimension d=fpq), and retain precisely the original J,C0,Kgood, full maximal-order actions, everywhere-good models and K0=Kgood(A0{[}M{]}). The split integral Serre and dual maps of LIFT090 give

\[
 K_{\rm out}=K_0(A_0[N^2])
 \subset K_{\rm good}(J[MN^2],C_0[MN^2]),
 \qquad (M,N)=1.                                      \tag{RE24}
\]

This uses the entire composed scale N=N0 t\^{}ell, including the component entry. It does not apply a same-seed statement to the changed lattice AB in isolation. LIFT093 PA7--8 and LIFT092 CD2--5, together with canonical R092's initial-field bound, yield

\[
\begin{split}
 [K_{\rm out}:B_{\rm coef}]&<64512\,d_{\rm CM} M^4N^{81/8},\\
 \operatorname{rd}(K_{\rm out})
 &<2^{50}e^{2\pi\sqrt{d_{\rm CM}}}
       \operatorname{rd}(B_{\rm coef})\,M N^2.
\end{split}                                             \tag{RE25}
\]

Both inequalities concern the original actual output field, not a replacement by the containing torsion field. Conductor divisibility and the subfield discriminant estimate are exactly those of LIFT092; the finite proof covers do not enter. In particular a sufficient absolute-discriminant bound is

\[
\begin{split}
 \log|\operatorname{disc}K_{\rm out}|
 <{}&[B_{\rm coef}:\mathbb Q]\,64512d_{\rm CM}M^4N^{81/8}\\
 &\cdot\bigl(\log\operatorname{rd}(B_{\rm coef})+50\log2
       +2\pi\sqrt{d_{\rm CM}}+\log M+2\log N\bigr).
\end{split}                                             \tag{RE26}
\]

The explicit upper bound in RE6 may be substituted for N everywhere in RE25--26, since all factors increase for N\textgreater=1. The original coefficient field's degree and discriminant remain inputs. There is no ray-class field factor: the ray quotient was used to select a lattice ideal, and the actual entry quotient descends in the original seed's finite torsion field. The bounded ideal selection adds the full multiplier C\_ent in RE6, not an unrecorded class-number multiplier or a field-degree saving. A list of stages may use the existing lcm bounds only with the same original two progenitors and all their shared-field hypotheses.

\begin{longtable}[]{@{}
  >{\raggedright\arraybackslash}p{(\columnwidth - 2\tabcolsep) * \real{0.5000}}
  >{\raggedright\arraybackslash}p{(\columnwidth - 2\tabcolsep) * \real{0.5000}}@{}}
\toprule\noalign{}
\begin{minipage}[b]{\linewidth}\raggedright
Gate or datum
\end{minipage} & \begin{minipage}[b]{\linewidth}\raggedright
Conclusion of this theorem
\end{minipage} \\
\midrule\noalign{}
\endhead
\bottomrule\noalign{}
\endlastfoot
Nonzero same-rule family and entry N0 & Established with RE2--3 for any supplied nonzero rule under the stated actual source/seed hypotheses. \\
Auxiliary c, very-ample h, ambient volume V and bad-locus D & The checked LIFT095--098 inputs apply uniformly after entry, with every factor in RE4 retained. \\
Original level, coefficient record, eigenlabels and trace-one line & Preserved exactly; no replacement what2/PinField or source ring is inferred. \\
Full source analytic-family records and Step5 inputs & This constructs the regular geometric family and field-valued FrPt used in LIFT085. It does not identify an arbitrary IsUnf/N1c/FamilyUnif record or remove the separately listed LIFT082 model, section, polynomial and fibre hypotheses. \\
N and all-object output field & Explicit upper bound RE6 and field K0(A0{[}N\^{}2{]}); stronger field-cost results retain their original seed hypotheses. \\
Chosen seed-line/field/discriminant inputs & The existing LIFT087--093 results may be composed only in their stated actual seed scopes. \\
Actual selected arithmetic rule nonzero, admitted full Host/Sizes, Read/032, DP44/RH & Remain open. No numerical admissibility or arithmetic-image conclusion is inferred from the conditional N bound. \\
Definite terminal q=0 or rank below three & Not covered; neither strong approximation nor the indefinite Hecke escape is asserted there. \\
\end{longtable}

The improvement over LIFT085 is a uniform bound on the finite-component entry, followed by the actual algebraic same-rule escape. It does not reuse the invalid idea that an analytic Cayley period map is a polynomial of bounded degree.

\subsection{9. Primary inputs, literal interfaces and verification boundary}\label{primary-inputs-literal-interfaces-and-verification-boundary}

Miller, S. D., \& Stephens-Davidowitz, N. (2019). \emph{Kissing numbers and transference theorems from generalized tail bounds} (arXiv:1802.05708v2). Theorem3.2, equation(3.3), and full proof, printed p.8. \url{https://arxiv.org/abs/1802.05708}. The finite-dimensional bound is used for every fractional ideal, not an asymptotic version or a ring-only estimate.

Milne, J. S. (2020). \emph{Class field theory} (v4.03). II Proposition1.22 and Corollary1.23, p.67 (Hilbert90); III Proposition1.1 and its proof, pp.97--98 (unramified unit solutions); V Proposition4.6 and its proof, pp.172--173 (the exact idelic ray-class quotient). \url{https://www.jmilne.org/math/CourseNotes/CFT.pdf}. These passages were freshly read for this proof. No ray-class field extension is used in the output construction.

Milne, J. S. (2017). \emph{Introduction to Shimura varieties}, Lemma8.20 and proof, p.89. \url{https://www.jmilne.org/math/xnotes/svi.pdf}. The exact unitary Hasse-principle scope was freshly reread, while LIFT085's already checked application with cyclic Hasse norm remains the carried source-to-rational-space bridge. Strong approximation is the complete already checked LIFT095 input (Rapinchuk2013, Theorem2.3,p281; Milne2017, Theorem4.16,p48).

The freshly read literal definitions are Host/Datum/Level/SAB/DetCond/TateIso/Obj in TwoWeightEqOneAddRoot.lean, including the complete Datum30--49 and Obj144--160; Moduli.rep, Coeff and XF in PushInvertible.lean29--49; and the exact IsDatW/PinLevel/PinDat clauses. FrPt, naturality, sec, FromTest and evalAt are those of LIFT085's full torsor argument and are fingerprinted with their exact source bodies in the final packet. The literal exists\_isFrPoly theorem \nolinkurl{(ExistsIsFramePolynomialOfRuleNumberField.lean50–112)} and isFrPoly\_equiv\_and\_ne\_zero \nolinkurl{(FramePolynomialEquivariantAndNeZero.lean62–115)} quantify arbitrary nonnegative a,n and every F in sec; they impose no additional Dom premise. The source even-n Dom permits the first value of k in RE4; otherwise squaring preserves the same reduced zero set. Source hypotheses are not replaced by names of records.

The argument takes place in actual projective modules, their finite congruence quotients and the finite ideal class quotient. The entry line is the normalized relative Poincare/trace-one line, and the quotient is an actual abelian variety over its finite torsion field. The existing cubical and all-action descent proofs are retained. No new motivic, lambda-ring, weight, crystalline-slope or global period assertion follows.

The companion executable checks finite unramified quadratic unit and norm-one images including dyadic prime powers, integral unitary determinant lifts at every tested congruence depth, nonnormal finite double quotients, nonprincipal maximal-order ideal cosets and their exact determinant covolumes, paired-prime lattice inclusions, full rational polarization matrices and degrees, protected-level inverses, the integer escape recursion and exact source fingerprints. Its finite models do not prove strong approximation, universal transference, the Hasse principle, algebraic period-family existence or literal source admission; these are explicit universal proof inputs above. The hash-specific final author and independent whole-review records are separate and must not be inferred from an executable PASS alone.

\subsection{Canonical algebraic-geometric context}\label{canonical-algebraic-geometric-context-65}

The new reduction is a finite ray-class quotient controlling integral projective Hermitian lattices in the same rational genus over Spec(O\_E), with exact congruence level and a bounded ordinary ideal norm. The chosen lattice becomes an actual principally polarized CM abelian variety over the original seed torsion field; its line-power identity is the normalized rigidified Poincare/cubical identity, and the nonzero rule is evaluated on an actual eigenframe torsor. New primary locators: Miller and Stephens-Davidowitz(2019),Theorem3.2/equation3.3 and proof,p8; Milne(2020),CFT II1.22-1.23 p67,III1.1 pp97-98,V4.6 pp172-173; Milne(2017),SVI Lemma8.20 p89. Strong approximation is carried Rapinchuk(2013),Theorem2.3 p281 and Milne(2017),Theorem4.16 p48 through LIFT095. Carried Grothendieck anchors: projective modules/exterior operations SGA3 I4.7.3/4.7.3.1 and SGA6 VI Theorem3.3; Picard/Poincare FGA232 Theorem3.1 pp148-149,FGA236 Corollaries1.7/2.3 pp229/232,SGA7 I VII2.7/2.9 pp169-174 and VIII7.4(b) pp310-311; formal crystals and Cartier duality cote4 scan003 sections8.1.1-8.1.2,scans005-006 sections8.4.1-8.4.3 and cote54 scans005-007. These are the existing contextual anchors, not new proofs of crystalline slopes, motivic realization comparisons, lambda-ring positivity or l-adic weights. DP\_MULT Q3e/DP44/DP52/DP54-65 and THETA3 G1/A7 receive an explicit conditional source-point cost, with the actual selected rule and six-identity/period matching still separate. CX\_GRH CX-ALG-024 and GCT\_PLETH T-052 retain their own global gates.

\clearpage

\section{AM-DO-121 --- exact unipotent Cartan profile and a complete support27 interval}\label{am-do-121-exact-unipotent-cartan-profile-and-a-complete-support27-interval}

\textbf{Canonical checked integration --- 2026-11-23T10:07:26.832081+00:00.} Complete root and independent proof/script/interface reviews and separate exact replays passed. The actual canonical predecessor has129/31 entries; this integration gives130/31, complete floors27/11. Frozen author history is preserved. AM-DO-121 EC1-24 retains the exact original unipotent identity cosh r\textgreater=1+alpha/2, hence sinh squared r\textgreater=alpha(1+alpha/4), in the genuine selected coefficient. Its positive infinite-mass Gamma167/rate4 convolution has every required convergent inverse moment; original product-energy reservation repairs the sole degree0 term before enlargement. The complete actual total-eligibility27/support27 all-II globally rank-two-unipotent scalar-cross union k111..139 has absolute normalized contribution below2\^{}-993. Every scalar/transverse polynomial term, original saturated metric and paired ideal, full167-entry reconstruction, unique eta, central affine750 cover, all168 finite factors and correlated four-state patterns is paid. One993 append changes actual K116129/31 to130/31, with both entire sums\textless2\^{}-230, all old overlapping errors retained and floors27/11 unchanged. The remaining necessary support27 alternatives are total eligibility27,140\textless=k\textless=164,2\^{}2111\textless x\textless2\^{}3328, or total eligibility28,k164,2\^{}3107\textless x\textless2\^{}3328; b=j\_int=0 throughout. The next k140 sufficient envelope767061/1024 is positive and is not a lower bound or impossibility theorem. No complete support28, signed nonvanishing, admitted source, selected arithmetic content, prescribed period, specific032 pairing value, evaluated DP44 saving or RH consequence.

Author: \nolinkurl{/root/definite_oscillator}. 23 November2026 UTC. \textbf{Controlled pending proof; no canonical promotion.} Whole-author verification, final exact replay, independent whole review and root review are recorded as separate events. Historical exploratory and preliminary arithmetic are retained separately.

\subsection{1. Actual sector and theorem}\label{actual-sector-and-theorem-11}

Keep the original J56 all-top/H133 kernel, the actual CM fields and involution, every source/target pencil and label, full selected vector-valued module and Hilbert metric, both compact factors, canonical twists, phases, Haar measures, class masses and all finite projectors. Put \[
 f=a=168,\quad H=133,\quad R=17424/133,
 \quad Q=1427908193729121522601,\quad U=77Q^2-14,
\] \[
 D_0=7^{140}29^{144}5^{84}13^{84},\qquad D_F=2^{168}D_0.
 \tag{EC1}
\]

The current predecessor is canonical K116,129 all-top/31 uniform entries and complete floors27/11. Its remaining total-eligibility27 branch is the actual globally rank-two-unipotent scalar-cross support27 sector with all active typesII, \[
 b=j_{\rm int}=0,\qquad 2^{2111}<x=N_ED_\eta<2^{3328},
 \qquad111\le k\le164.
 \tag{EC2}
\] Here k counts nonzero original S covector values. The separate total-eligibility28 alternative has k164 and log2x\textgreater3107; no cut for that alternative is claimed here.

\textbf{Theorem.} On EC2 the complete interval111..139 has absolute normalized contribution below2\^{}-993. One new993 charge is appended, retaining every previous overlapping charge and both ordered arrays. The sufficient estimate at k140 is positive and makes no removal claim. This theorem does not imply complete support28 or signed nonvanishing.

\subsection{2. The original unipotent identity at the actual near point}\label{the-original-unipotent-identity-at-the-actual-near-point}

Let h=I+N be an actual matrix in EC2. The established full-positive normal form gives rank\_E N=2 and N\^{}3=0. At the original negative coordinate, \[
 \alpha=(N^2)_{nn}=U r J_Z^{-1}r^*>0,
 \qquad0<\alpha_\tau\le R.
 \tag{EC3}
\] These are literal K113 SC27.5--7 / DO098 CF4--6 identities, not a finite-field positivity claim. Let h\_p=g\_p\^{}\{-1\}hg\_p=I+N\_p at the original actual near point. Nilpotence and unitarity are preserved. In the same diagonal Hermitian convention, \[
 (h_p^{-1})_{nn}=\overline{(h_p)_{nn}},\qquad
 h_p^{-1}=I-N_p+N_p^2.
\] Consequently \[
 \alpha_p=(N_p^2)_{nn}=N_{p,nn}+\overline{N_{p,nn}},
 \qquad\Re(h_{p,nn})=1+\alpha_p/2.
 \tag{EC4}
\] DO080's actual selected radial fidelity, section7, already identifies this same alpha\_p and proves alpha\_p\textgreater=exp(-2r\_p)alpha. It also identifies \textbar h\_p,nn\textbar=cosh r(h\_p); no replacement finite null frame is needed. Thus \[
 \cosh r(h_p)\ge1+\alpha_p/2,
 \qquad\sinh^2r(h_p)\ge\alpha_p+\alpha_p^2/4.
 \tag{EC5}
\] The old estimate retained only the first summand in EC5. The quadratic term is now retained.

The accepted universal actual coefficient estimate K108 AE9 is \[
 |c_W(h_p)|\le
 \vartheta_a2^{-933/2}\frac1{\sinh r(h_p)}
 \left(\frac2{\cosh r(h_p)+1}\right)^{a-3/2}.
\] For nonidentity rank-two unipotence EC3--5 give \[
 |c_W(h_p)|\le\vartheta_a2^{-933/2}
 \alpha_p^{-1/2}(1+\alpha_p/4)^{-167}.
 \tag{EC6}
\] This is the same selected trace, not a scalar representation substitute. The angular exponent166.5 and the additional reciprocal-sinh half power combine to167.

Write F(t)=t\^{}-1/2(1+t/4)\^{}-167. Its absolute logarithmic elasticity satisfies \[
 -\frac{d\log F(t)}{d\log t}
 =\frac12+\frac{167t}{t+4}<\frac{335}{2}.
\] Therefore F(exp(-2r\_p)alpha)\textless=exp(335r\_p)F(alpha). The literal near-point bound r\_p\textless2\^{}-352 implies exp(335 sum\_tau r\_p)\textless2. This \textbf{replaces} the previous global near-point2 and does not multiply another2. Combining all embeddings, with vartheta\_a\^{}168\textless2\^{}-649, \[
 |c(h)|<2\,2^{-649}\,2^{-933f/2}
 \prod_\tau\alpha_\tau^{-1/2}(1+\alpha_\tau/4)^{-167}.
 \tag{EC7}
\] The whole lower-probability tail was already absorbed into933 by K107. None is discarded here.

\subsection{3. An unshifted positive measure with every required inverse moment}\label{an-unshifted-positive-measure-with-every-required-inverse-moment}

Let gamma be the Gamma probability measure of shape m=167 and rate beta=4. For all t\textgreater0, \[
 (1+t/4)^{-167}=\int_0^\infty e^{-bt}\,d\gamma(b),
 \qquad t^{-1/2}=\int_0^\infty e^{-st}\frac{ds}{\sqrt{\pi s}}.
\] Let mu be their convolution under c=b+s. It is a locally finite positive measure with infinite total mass. Indeed its density is proportional to the integral from0 to c of b\^{}166(c-b)\^{}(-1/2)exp(-4b)db, which is locally integrable. Its total mass is infinite because the half-Gamma factor has infinite mass and the Gamma factor has mass1. One has \[
 F(t)=\int_0^\infty e^{-ct}\,d\mu(c).
 \tag{EC8}
\] There is no paid shift or cap enlargement loss. Positivity permits Tonelli. For integer1\textless=j\textless=167, \[
 M_j=\int c^{-j}\,d\mu(c)
 =4^{j-1/2}\frac{\Gamma(j-1/2)\Gamma(167+1/2-j)}
                  {\Gamma(j)\Gamma(167)}<\infty.
 \tag{EC9}
\] The convolution Beta integral has j-1/2\textgreater0, and the remaining Gamma argument167.5-j is positive throughout the stated range. No M0 is assigned. FormulaEC9 uses ordinary convergent integrals, with NIST DLMF5.9.1,5.12.1/3 and5.4.1/6 as the normalization anchors.

Let H\_n=(2n)!/(4\^{}n n!) and pi\_+=22/7. Exact rational powers satisfying P\_j\textgreater=M\_j\^{}f are \[
 P_j=(\pi_+^2\,4)^{84}
 \left[\frac{4^{j-1}H_{j-1}H_{167-j}}
 {(j-1)!166!}\right]^{168}.
 \tag{EC10}
\] The independent recurrence is M\_(j+1)/M\_j=4(j-1/2)/{[}j(166.5-j){]}. Every actual row in the interval uses only j\textless=142; the universal moment identity holds through167.

\subsection{4. Repair the sole zero-degree term before enlargement}\label{repair-the-sole-zero-degree-term-before-enlargement}

The original row ideal is L=J I\_Sigma D\_eta, with NE L=z\^{}(2f), and q=R/(Uz²). Every one of the k nonzero original S values belongs to L at both CM partners. The genuine-row product inequality gives \[
 \prod_\tau\alpha_\tau\ge\alpha_0^f,
 \qquad\alpha_0=Ukz^2=Rk/q.
 \tag{EC11}
\] For positive c\_tau let B=(prod c\_tau)\^{}(1/f). Before enlarging any row set, \[
 \sum c_\tau\alpha_\tau\ge f\alpha_0B.
\] Reserve theta=2\^{}-20 of this actual exponent and enlarge only the remainder: \[
 e^{-\sum c\alpha}\le
 e^{-\theta f\alpha_0B}e^{-(1-\theta)\sum c\alpha}.
 \tag{EC12}
\] For j\textgreater=1, dropping the reserved damping leaves the finite product moment M\_j\^{}f.~For the only degree-zero term use \[
 e^{-\theta f\alpha_0B}\le(\theta\alpha_0B)^{-f},
 \qquad\int e^{-\theta f\alpha_0B}\,d\mu^{\otimes f}
 \le P_1\left(\frac{q}{\theta Rk}\right)^f.
 \tag{EC13}
\] This is the complete K116 RE11--13 repair, now with the unshifted measureEC8. It does not assert a per-embedding lower bound for alpha, finite total mass, or a nonexistent zeroth moment. The old separate k\^{}-84 / z\^{}-f trace bound is not also retained.

\subsection{5. Entire original lattice polynomial and arithmetic factors}\label{entire-original-lattice-polynomial-and-arithmetic-factors}

Use the actual saturated integral transverse basis. Its Gram G' has detG'=U and G'\textgreater2\^{}36 I; both CM membership and the original half-trace primal/dual conventions are unchanged. Put c0=2\^{}36 and let K be the exact K111 HP12--18 scalar coefficient \[
 c_{336}=168\binom{336}{168}/4^{168},\quad
 M_{15}=\frac{6015684926187167549292201834533348561577984}
 {2618017117694612446410293519},
\] \[
 K=\frac{336M_{15}}{c_{336}}(7841/5000)^{168}168!
 \left[\frac{(365729/500)^2}{168(314159/100000)^2}\right]^{168}.
 \tag{EC14}
\] K110 PM10--12's first, untruncated Poisson inequalities hold for arbitrary positive embedding weights. They yield the full scalar polynomial and all four transverse terms; no simplification requiring B\textgreater=2 is imported. With d=k+3, the integrated row factor is at most \[
 \binom{164}{k}\frac{U^f\pi_+^{df}}{D_0^d}\mathcal P_k(q),
 \tag{EC15}
\] where \[
 \mathcal P_k(q)=
 \sum_{\substack{0\le l\le k,\ 0\le h\le3\\j=l+3-h\ge1}}
 \binom kl\binom3h K^{k-l+h}c_0^{-hf}
 \left(\frac{q}{(1-\theta)R}\right)^{jf}P_j
 +K^{k+3}c_0^{-3f}P_1
 \left(\frac{q}{\theta Rk}\right)^f.
 \tag{EC16}
\] The last summand is exactly(l,h)=(0,3), not an omitted remainder. Real diagonal metric changes do not change L or its paired denominator orientation.

Keep the full original affine central coset, including its real-part shift and dyadic possibility. Its count is \[
 V_F(H/2+750R/q)\frac{T^{54}}{U^f x N_FJ_0},
 \quad V_F(C)=\frac{(\pi_+f C^2)^{84}}{84!\sqrt{D_F}}.
 \tag{EC17}
\] The750 full-real cover is retained. No half-lattice or inferred smaller fibre is used. Fixing eta, all167 original covector values and the central parameter determines the full matrix. Bounded eta uniqueness per fixed support follows from the original27 eligibility congruences and the norm gap throughout x\textless2\^{}3328; it introduces no unpaid label or beta fibre.

At these actual total27,b=j\_int0 terms, all27 active weighted factors and141 inactive pointwise factors are present: \[
 (T-1)^{-168}<2T^{-168}.
\] The unchanged Green factor is at most2Q\^{}(-335*167)(NEJ)\^{}(335/2), with \[
 N_EJ=z^{2f}T^{54}/x\le Q^{334},
 \qquad N_EJ\le Q^{167}N_FJ_0.
 \tag{EC18}
\] Support choices\textless2\^{}104, the harmless original pivot choices\textless2\^{}8, Green2, T2, actual near-point2 and conservative pattern4 give fixed117 bits. Every old source/target, phase, class, Haar and selected-trace factor remains.

\subsection{6. Uniform ideal, denominator and finite-pattern bounds}\label{uniform-ideal-denominator-and-finite-pattern-bounds}

For a positive degree j, after all common powers are collected the q factor is \[
 F_j(q)=q^{(j-335/2)f}V_F(H/2+750R/q).
 \tag{EC19}
\] The repaired degree-zero term has effective j=1. If epsilon=750R/{[}q(H/2+750R/q){]}, then0\textless epsilon\textless1, and \[
 \frac1f\frac{d\log F_j}{d\log q}=j-335/2-\epsilon<0.
\] At q\_\emph{\^{}f=R\^{}fT\^{}54/(U\^{}fQ\^{}334x), the coupled x derivative of x\^{}-337/2 F\_j(q\_}) is -1-j+epsilon\textless0. Thus every term is bounded at x=2\^{}2111 with the exact q cap, not an independently inconsistent pair of endpoints.

Writing D=Q\^{}334/NEJ gives q=q\_*D\^{}(1/f) and (NFJ0)\^{}-1\textless=Q\^{}-167D. For each degree the original four-state multiplicities1,2,1 produce \[
 \prod_{v=1}^{334}(1+q_v^{j-333/2})^2.
 \tag{EC20}
\] Here j\textless=142, hence exponent\textless=-24.5\textless-1, and the product is below2 by the original g0 comparison; the complete count conservatively retains4. There is no extra orientation or independent pattern multiplier. This uniform bound is not asserted for a hypothetical j167 application.

\subsection{7. Complete exact endpoint and remaining scope}\label{complete-exact-endpoint-and-remaining-scope}

The original directed rational logarithms and exact cap checks are \[
 L_T=1019333/1024,\quad L_Q=287843/4096,
 \quad L_U=601356/4096,\quad L_R=28810/4096,
\] \[
 251124964<q_*<251124965,\quad
 V_F(H/2+750R/251124964)<2^{474}.
 \tag{EC21}
\] With all preceding factors, the common exponent is unchanged from K116: \[
 C=117-933f/2-649+474+(335f/2)(L_R-L_U)
       +8931L_T-336\cdot167L_Q-(337/2)2111
   =593387605/1024.
 \tag{EC22}
\] Unlike K116, there is no positive paid-shift exponent. Every cell retains all4(k+1) terms ofEC16. Exact moment ceilings and the same scalar ratio ceiling-935 bound the full nonmain polynomial by the main term, giving a conservative factor2. The full cell exponent is \[
 E_k=C+\left\lceil\log_2\left[
 2\binom{164}k\frac{\pi_+^{(k+3)f}}{D_0^{k+3}}
 q_-^{-335f/2}
 \left(\frac{q_+}{(1-\theta)R}\right)^{(k+3)f}P_{k+3}
 \right]\right\rceil_{\rm strict}.
 \tag{EC23}
\] The following exact full-cell table verifies the entire interval111..139, with exact worst E139=-1017771/1024\textless-993.5, penultimate E138=-2792363/1024\textless-2726, and every preceding cell below-2726. The exact failed next endpoint is E140=767061/1024\textgreater0. Hence \[
 \sum_{k=111}^{139}2^{E_k}
 <3\,2^{-995}+28\,2^{-2726}<2^{-993}.
 \tag{EC24}
\]

\begin{longtable}[]{@{}rrr@{}}
\toprule\noalign{}
k & Complete rational factor ceiling & Full exponent E\_k \\
\midrule\noalign{}
\endhead
\bottomrule\noalign{}
\endlastfoot
111 & -626285 & -47928235/1024 \\
112 & -624726 & -46331819/1024 \\
113 & -623162 & -44730283/1024 \\
114 & -621593 & -43123627/1024 \\
115 & -620020 & -41512875/1024 \\
116 & -618442 & -39897003/1024 \\
117 & -616859 & -38276011/1024 \\
118 & -615270 & -36648875/1024 \\
119 & -613677 & -35017643/1024 \\
120 & -612078 & -33380267/1024 \\
121 & -610473 & -31736747/1024 \\
122 & -608863 & -30088107/1024 \\
123 & -607247 & -28433323/1024 \\
124 & -605626 & -26773419/1024 \\
125 & -603998 & -25106347/1024 \\
126 & -602364 & -23433131/1024 \\
127 & -600724 & -21753771/1024 \\
128 & -599077 & -20067243/1024 \\
129 & -597424 & -18374571/1024 \\
130 & -595764 & -16674731/1024 \\
131 & -594096 & -14966699/1024 \\
132 & -592422 & -13252523/1024 \\
133 & -590739 & -11529131/1024 \\
134 & -589050 & -9799595/1024 \\
135 & -587352 & -8060843/1024 \\
136 & -585645 & -6312875/1024 \\
137 & -583930 & -4556715/1024 \\
138 & -582207 & -2792363/1024 \\
139 & -580474 & -1017771/1024 \\
\end{longtable}

The coordinate-position binomial, supports, pivots, unique slope, full central coset, all finite patterns and local factors are already in each cell. In EC24, 2\^{}(-1/2)\textless3/4 supplies the first dyadic comparison. This29-cell sum is the only remaining union. The k140 sufficient envelope is positive and is retained solely as a diagnostic. The total-eligibility28,k164 branch remains unclaimed; its full scalar correction and high-degree pattern boundary cannot be discarded.

The complete executable checks all exact host constants, EC4--7 algebraic identities and elasticity, all positive Gamma arguments and moment recurrences, every scalar/transverse term includingEC13, full continuum derivative signs and pattern multiplicities, cap roots, all29 endpoints and the complete union. It retains the actual K116129/31 ordered prefix, appends993 once, preserves all uniform entries and verifies both full dyadic sums\textless2\^{}-230, with floors27/11 unchanged. The author event, independent whole check and root whole check remain distinct from these diagnostic assertions and are required before canonical acceptance.

\subsection{8. Source and scope record}\label{source-and-scope-record}

Direct interfaces are DO080 actual\_selected\_radial\_fidelity §7; K108 AE8--13; K107 SV8--19 including its mandatory boundary corrections; K113 SC27.1--15; K116 RE11--28; K110 PM10--12; K111 HP12--18; K109 TG3--12; DO098 CF4--16 and DO104 CD10--19. NIST DLMF5.9.1,5.12.1/3 and5.4.1/6 provide the positive Gamma/Beta normalizations; all new instances have explicitly positive arguments. K's separate targeted identity check is not a whole proof/count review.

All original projective OE modules, involution, polarization, theta/Picard/Poincare data and chosen selected realization remain. Inherited context only: SGA3 I4.7.3/4.7.3.1; SGA6 VI Theorem3.3; FGA232 Theorem3.1 pp148--149; FGA236 Corollaries1.7/2.3 pp229/232; SGA7 I VII2.7/2.9 pp169--174 and VIII7.4(b) pp310--311; cote4 scan003 sections8.1.1--8.1.2 and scans005--006 sections8.4.1--8.4.3; cote54 scans005--007. DP\_MULT Q3e/DP44, THETA3 G1/A7, CX\_GRH CX-ALG-024 and GCT\_PLETH T-052 retain their signed/source/selected-content/prescribed-period/global comparison gates. No Lean build, archive audit, global short-arrow existence or support28 theorem is asserted.

\subsection{Canonical algebraic-geometric context}\label{canonical-algebraic-geometric-context-66}

The strengthened majorant applies to the same actual selected real realization and original unipotent endomorphism. The positive infinite measure is only an analytic integration device; it adds no arithmetic parameter, field, physical vector, projector or lattice rescaling. All original projective OE modules, CM involution, polarization, cubical theta line, Picard and Poincare data remain. The new identities and ordinary convergent inverse moments are proved in full, using NIST DLMF5.9.1,5.12.1/5.12.3 and5.4.6 with explicit positive arguments. Genuine vector-valued action/metric remains the established Nakahama2016 section5.1, equations5.1/5.3/5.5/5.7 and Theorem5.1 pp960-961; independent fresh versus carried source scopes are recorded. Inherited context only: SGA3 I4.7.3/4.7.3.1; SGA6 VI Theorem3.3; FGA232 Theorem3.1 pp148-149; FGA236 Corollaries1.7/2.3 pp229/232; SGA7 I VII2.7/2.9 pp169-174 and VIII7.4(b) pp310-311; cote4 scan003 sections8.1.1-8.1.2 and scans005-006 sections8.4.1-8.4.3; cote54 scans005-007. No new formal-group/Dieudonne, lambda-ring/K-theory, l-adic-weight or CM-motive comparison. DP\_MULT Q3e/DP44 and THETA3 G1/A7 receive the strictly smaller conditional finite gate; CX\_GRH CX-ALG-024 and GCT\_PLETH T-052 retain their arithmetic/global gates.

\clearpage

\section{AM-R-098 --- a weighted radial Jacobian and a clamped positive profile}\label{am-r-098-a-weighted-radial-jacobian-and-a-clamped-positive-profile}

\textbf{Canonical checked integration --- 2026-11-23T10:18:52.302786+00:00.} Complete saved-proof and executable reviews and separate exact diagnostic replays passed. No numerical error charge is introduced. Frozen author history is preserved. AM-R-098 WJ1-28 proves a new bound for the same actual selected trace with factor(cosh r/sinh r){[}2/(cosh r+1){]}\^{}(333/2)/sqrt30, retaining both compact factors and the full probability tail. The original unipotent energy yields shape166 with whole-field coefficient(2/15)\^{}84 and the single old near-point2. A positive clamp at delta=1/(1000R), imposed on actual rows before enlargement, pays less than125/104 and makes every required inverse moment through167 finite. The original un-clamped167th moment and total mass diverge; the highest moment includes both clamp atom and full upper integral, and the sole polynomial degree0 is repaired by original product-energy reservation. Every term of the all-positive-weight scalar/transverse polynomial is retained. This is an analytic interface only: it introduces no charge, source admission, central-fibre gain, high-degree finite-pattern shortcut, signed nonvanishing, selected arithmetic-image detection, specific032 value, evaluated DP44 or RH result. Actual130/31 error arrays and floors27/11 remain unchanged.

Author: \nolinkurl{/root}, GPT-6 Astra Pro; inherited reasoning effort is not externally exposed. 23 November 2026 UTC. \textbf{Controlled pending proof.} This is a new analytic estimate for the existing actual selected trace. It gives no complete arithmetic sector charge by itself. Whole-author, independent whole review and exact diagnostics are separate events.

\subsection{1. Scope and imported exact interfaces}\label{scope-and-imported-exact-interfaces}

Keep the actual J56 selected vector-valued holomorphic representation, its Hilbert metric, both compact factors, source/target labels, original fields and involution, all phases, projectors, Haar measures and class masses. In the notation of DO080 RT4--14, the invariant ball measure is disintegrated as \nolinkurl{w(y) dy d omega}, the actual selected probability has density \nolinkurl{p(y)}, and its conditional radial transition is \nolinkurl{K_r(y,y')}. The established finite-trace fidelity inequality is \[
 |c_W(g)|\le\iint\sqrt{p(y)p(y')}
       \sqrt{K_r(y,y')K_r(y',y)}\,dy\,dy'.                 \tag{WJ1}
\] This is the original finite selected trace, not an infinite trace or operator norm. The inverse Cartan element is compact-conjugate to the original, so detailed balance applies with both actual compact factors retained.

Put \[
 a=f=168,\quad \nu=a-\tfrac32=\tfrac{333}{2},\quad
 \vartheta=\frac{4^{167}}{334\binom{334}{167}},\quad
 R=\frac{17424}{133}<132.                                 \tag{WJ2}
\] The unchanged original arithmetic constants used below are \nolinkurl{Q=1427908193729121522601}, \nolinkurl{U=77Q^2-14} and \nolinkurl{D_0=7^140 29^144 5^84 13^84}.

K107 SV4 and SV8--12 provide, for the exact selected radial variable, \[
 S=2\sqrt{Y^{-1}-1},\quad s_0=2\sqrt m,\quad31<m<33,
 \quad {\bf E}(S-s_0)^2\le\sigma_*^2,
 \quad(334\sigma_*)^2<2^{-933}.                            \tag{WJ3}
\] This includes the entire lower-probability tail. The variance center need not be the mean of S. In particular \nolinkurl{0<sigma_*<2^-12}. The established original near-point bounds are \nolinkurl{r_p<2^-352} at every embedding and \nolinkurl{alpha_p>=exp(-2r_p)alpha}, with \nolinkurl{0<alpha<=R}, for the actual nonidentity globally rank-two-unipotent class. All these are imported at their checked scope; no support-specific numerical conclusion is imported.

\textbf{Theorem.} For every actual selected trace with these interfaces and Cartan radius \nolinkurl{r>0}, \[
 |c_W(g)|\le\min\left\{1,
 \frac{\vartheta\,2^{-933/2}}{\sqrt{30}}
 \frac{\cosh r}{\sinh r}
 \left(\frac2{\cosh r+1}\right)^\nu\right\}.               \tag{WJ4}
\] For the original rank-two-unipotent matrices this implies the positive profile bound WJ17 below. A clamp of its representing positive measure has every inverse moment needed through167, with explicit rational majorants WJ23--25 and a global cost less than125/104. Neither the unclamped167th moment nor a zeroth moment is asserted to be finite. Any arithmetic count using the theorem must still pay every original source, central-fibre and finite-pattern factor.

\subsection{2. An exact feasible-pair Jacobian}\label{an-exact-feasible-pair-jacobian}

Write \nolinkurl{c=cosh r}, \nolinkurl{s=sinh r>0}, \nolinkurl{R_y=sqrt(1-y)}, \nolinkurl{D=sqrt(y/y')}. In the original first complex coordinate \nolinkurl{z_1=u+it}, put \[
 v=c+su=\sqrt{D^2-s^2t^2}>0,
 \qquad u^2+t^2\le1-y.                                   \tag{WJ5}
\] The positive square root follows from \nolinkurl{c-s R_y>0}. The disk constraint is equivalent to \[
 D^2-2cv+c^2\le s^2(1-y),
 \quad\text{hence}\quad
 2cv\ge D^2+1+s^2y\ge2D.                                 \tag{WJ6}
\] Here \nolinkurl{c^2-s^2=1} and \nolinkurl{D^2+1-2D=(D-1)^2>=0}. Thus \nolinkurl{v>=D/c} for every feasible pair. The exact real Jacobian of the same original map is therefore \[
 \left|\frac{du}{dy'}\right|
   =\frac{y}{2s y'^2v}
   \le\frac{c\sqrt y}{2s y'^{3/2}}.                       \tag{WJ7}
\] This inequality is valid on the entire feasible region, including both \nolinkurl{c>D} and \nolinkurl{c<=D}. There is no radial band or probability conditioning.

The original normalized angular density is \[
 \frac{a-1}{\pi R_y^2}
 \left(1-\frac{u^2+t^2}{R_y^2}\right)^{a-2}
 {\bf1}_{u^2+t^2<R_y^2}.                                 \tag{WJ8}
\] Use the accepted K108 AE5--6 angular-height calculation at its support-independent scope: if \nolinkurl{c>D}, the integral in t of the density power, without its prefactor, is bounded by \nolinkurl{2 vartheta R_y [2D/(c+1)]^nu}; if \nolinkurl{c<=D}, the unchanged full beta integral gives the same bound because \nolinkurl{2D/(c+1)>=1}. Multiplying by WJ7 yields, with zero extension for infeasible pairs, \[
 K_r(y,y')\le
 \vartheta\frac{(a-1)c}{\pi s}
 \frac{\sqrt y}{\sqrt{1-y}\,y'^{3/2}}
 \left(\frac{2\sqrt{y/y'}}{c+1}\right)^\nu.               \tag{WJ9}
\] Apply WJ9 in both directions. Their D factors cancel exactly, and the geometric mean is \[
 \sqrt{K_r(y,y')K_r(y',y)}\le
 \vartheta\frac{(a-1)c}{\pi s}
 (yy')^{-1/2}[(1-y)(1-y')]^{-1/4}
 \left(\frac2{c+1}\right)^\nu.                            \tag{WJ10}
\] This cancellation is for a nonnegative upper kernel; the signed selected coefficient is not identified with a radial function.

\subsection{3. A weighted whole-probability Cauchy estimate}\label{a-weighted-whole-probability-cauchy-estimate}

By WJ1 and WJ10, the remaining integral is the square of \[
 I=\int_0^1\sqrt{p(y)}\,y^{-1/2}(1-y)^{-1/4}\,dy
   =\int_0^\infty\sqrt{q_S(S)}
                  \left(\frac4{S^2+4}\right)^{1/4}\,dS.  \tag{WJ11}
\] The equality is exact: \nolinkurl{y=4/(S^2+4)} and \nolinkurl{|dy/dS|=8S/(S^2+4)^2}. In particular the new multiplier is part of the same selected probability density; no scalar-module replacement occurs.

Put \nolinkurl{sigma=sigma_*}. Weighted Cauchy--Schwarz, first on finite intervals and then by monotone convergence, gives \[
 I^2\le2\sigma^2 J,\qquad
 J=\int_0^\infty
 \frac{2/\sqrt{S^2+4}}{(S-s_0)^2+\sigma^2}\,dS.            \tag{WJ12}
\] The first Cauchy factor is at most \nolinkurl{E(S-s_0)^2+sigma^2<=2sigma^2}. To bound J, split the integral only for estimation, retaining both parts. If \nolinkurl{|S-s_0|<=1/8}, then \[
 S\ge2\sqrt{31}-1/8>0,\quad
 S^2+4\ge128+1/64-\sqrt{31}/2>124,
\] since \nolinkurl{sqrt(31)<6}. Thus its numerator is less than \nolinkurl{1/sqrt(31)}. On the complement it is at most1 and the denominator is at least \nolinkurl{(S-s_0)^2}. Extending these nonnegative integrals to the appropriate subsets of the real line gives \[
 J<\frac\pi{\sigma\sqrt{31}}+
       2\int_{1/8}^\infty t^{-2}\,dt
   =\frac\pi{\sigma\sqrt{31}}+16.                        \tag{WJ13}
\] No probability tail has been removed. Consequently \[
 I^2<\frac{2\pi\sigma}{\sqrt{31}}+32\sigma^2
     <\frac{2\pi\sigma}{\sqrt{30}}.                      \tag{WJ14}
\] For completeness, the last comparison is exact. One has \[
 \frac1{\sqrt{30}}-\frac1{\sqrt{31}}
 =\frac1{\sqrt{930}(\sqrt{31}+\sqrt{30})}>\frac1{372},
\] using \nolinkurl{sqrt(930)<31} and \nolinkurl{sqrt(31)+sqrt(30)<12}. On the other hand, \nolinkurl{sigma<2^-12} and \nolinkurl{pi>3} give \nolinkurl{16 sigma/pi<1/768<1/372}. This is precisely the inequality needed after dividing by \nolinkurl{2 pi sigma}.

Combining WJ10, WJ14 and \nolinkurl{2(a-1)sigma=334sigma<2^-933/2} proves WJ4. At \nolinkurl{r=0} retain the ordinary bound \nolinkurl{|c_W(g)|<=1}. The proof has preserved the full lower-probability tail inside sigma and has introduced no additive tail charge.

\subsection{4. Original unipotent energy and one near-point factor}\label{original-unipotent-energy-and-one-near-point-factor}

The checked actual identity DO080 section7 / DO121 EC4--5 gives \nolinkurl{cosh r(h_p)>=1+alpha_p/2}. The function \[
 G(c)=\frac{c}{\sqrt{c^2-1}}
              \left(\frac2{c+1}\right)^\nu,\qquad c>1,
\] is decreasing because \[
 \frac{d\log G}{dc}
 =-\frac1{c(c^2-1)}-\frac\nu{c+1}<0.                     \tag{WJ15}
\] Thus WJ4 implies \[
 |c_W(h_p)|\le\frac{\vartheta2^{-933/2}}{\sqrt{30}}
 (1+\alpha_p/2)\alpha_p^{-1/2}(1+\alpha_p/4)^{-167}
 \le\frac2{\sqrt{30}}\vartheta2^{-933/2}F(\alpha_p),
 \quad F(t)=t^{-1/2}(1+t/4)^{-166}.                       \tag{WJ16}
\] The last inequality uses \nolinkurl{1+t/2<=2(1+t/4)}. It is a valid new bound, not a claim that this weakened profile is pointwise smaller than DO121's profile for every t.

Its absolute logarithmic elasticity is \nolinkurl{1/2+166t/(t+4)<333/2}. The actual multiplicative near-point transport therefore costs at most \nolinkurl{exp(333 sum r_p)<2}; indeed \nolinkurl{333*168*2^-352<1/2<log 2}. This is the one existing global near-point2. The original energy has \nolinkurl{0<alpha_tau<=R} before the clamp below. Since \nolinkurl{vartheta^f<2^-649}, multiplication over all f embeddings gives \[
 |c(h)|<2\,2^{-649}\,2^{-933f/2}
            (2/15)^{84}\prod_\tau F(\alpha_\tau).         \tag{WJ17}
\] All original selected normalization, words and compact factors remain. No external \nolinkurl{k^-84 z^-f} estimate is simultaneously multiplied into WJ17.

\subsection{5. Positive convolution, the divergent moment, and the paid clamp}\label{positive-convolution-the-divergent-moment-and-the-paid-clamp}

Let gamma be the Gamma probability measure with shape166 and rate4 and let \nolinkurl{d eta(s)=ds/sqrt(pi s)} on the positive half-line. Their convolution under \nolinkurl{c=b+s} is a positive locally finite measure mu with infinite total mass, and \[
 F(t)=\int_0^\infty e^{-ct}\,d\mu(c),\qquad t>0.           \tag{WJ18}
\] The ordinary positive Gamma and Beta integrals give, for integer \nolinkurl{1<=j<=166}, \[
 M_j:=\int c^{-j}\,d\mu(c)
 =4^{j-1/2}
 \frac{\Gamma(j-1/2)\Gamma(166+1/2-j)}
      {\Gamma(j)\Gamma(166)}<\infty.                     \tag{WJ19}
\] Here both Beta/Gamma arguments are positive. This is the same general positive convolution identity as DO121 EC8--10, now applied with shape166. Its range ends at166. In fact the density is \[
 \rho_\mu(c)=\frac{4^{166}}{\Gamma(166)\sqrt\pi}
       \int_0^c b^{165}(c-b)^{-1/2}e^{-4b}\,db
 \le A c^{331/2},\qquad
 A=\frac{4^{166}}{\Gamma(333/2)}.                         \tag{WJ20}
\] The beta integral proves the inequality. For \nolinkurl{0<c<=delta}, the reverse bound with an additional factor \nolinkurl{exp(-4delta)} also holds. Hence the unclamped167th inverse moment diverges at0 as the integral of \nolinkurl{c^-3/2}. No analytic continuation assigns it a finite value.

Choose \[
 \delta=\frac1{1000R}>0,
 \qquad\mu_\delta=(c\mapsto\max\{c,\delta\})_*\mu.        \tag{WJ21}
\] This is an explicit positive pushforward, including an atom at delta of finite mass. Its total mass remains infinite. For each actual \nolinkurl{0<alpha<=R}, the pointwise inequality \nolinkurl{exp(-c alpha)<=exp(delta R) exp(-max(c,delta)alpha)} gives \[
 \prod_\tau F(\alpha_\tau)
 \le e^{f/1000}
 \int e^{-\sum_\tau c_\tau\alpha_\tau}
                       \,d\mu_\delta^{\otimes f}(c)
 <\frac{125}{104}
 \int e^{-\sum_\tau c_\tau\alpha_\tau}
                       \,d\mu_\delta^{\otimes f}(c).     \tag{WJ22}
\] Indeed \nolinkurl{f/1000=21/125<1} and the positive series gives \nolinkurl{e^u<1/(1-u)=125/104}. The clamp is imposed on actual rows after near-point transport and before lattice enlargement; the original upper energy bound R is used at exactly that point.

For \nolinkurl{1<=j<=166}, monotonicity of the inverse power gives \nolinkurl{M_j(delta)<=M_j}. At the highest needed degree, WJ20 gives the fully convergent estimate \[
 \begin{split}
 M_{167}(\delta)
 &=\delta^{-167}\int_0^\delta d\mu(c)
                       +\int_\delta^\infty c^{-167}d\mu(c)\\
 &\le A\delta^{-1/2}\left(\frac2{333}+2\right).
 \end{split}                                             \tag{WJ23}
\] Both pieces are required. The first is the clamp atom, and the second is bounded by \nolinkurl{A integral_delta^infty c^-3/2 dc}. In particular this is not a truncated integral with an unpaid discarded interval.

For exact rational use define \nolinkurl{H_n=(2n)!/(4^n n!)}, \nolinkurl{pi_+=22/7} and \nolinkurl{pi_-=314159/100000<pi}. Since f=168 is even, explicit rational upper bounds for all product moments are \[
 P_j=(4\pi_+^2)^{84}
 \left[\frac{4^{j-1}H_{j-1}H_{166-j}}
 {(j-1)!165!}\right]^{168},\quad1\le j\le166,             \tag{WJ24}
\] \[
 P_{167}=
 \left[\frac{4^{166}}{H_{166}}\left(2+\frac2{333}\right)\right]^{168}
       \delta^{-84}\pi_-^{-84}.                          \tag{WJ25}
\] They satisfy \nolinkurl{M_j(delta)^f<=P_j}. WJ25 uses \nolinkurl{Gamma(333/2)=sqrt(pi)H_166}; its lower bound for pi occurs because pi is in the denominator. The standard exact pi brackets are the already checked arithmetic brackets used in K111/DO121, not newly measured decimals. Every rational power can be evaluated by integer arithmetic.

\subsection{6. Original product-energy reservation and the full polynomial interface}\label{original-product-energy-reservation-and-the-full-polynomial-interface}

Let the actual common paired row ideal be \nolinkurl{L=J I_Sigma D_eta} with \nolinkurl{NE L=z^(2f)} and \nolinkurl{q=R/(Uz^2)}. For \nolinkurl{1<=k<=164} nonzero original S values, the accepted actual product-energy inequality is \[
 \prod_\tau\alpha_\tau\ge(Rk/q)^f.                      \tag{WJ26}
\] This is imposed on original rows, not on the enlarged Gaussian lattice. For positive clamped weights put \nolinkurl{B_c=(prod c_tau)^(1/f)}, \nolinkurl{alpha_0=Rk/q} and \nolinkurl{theta=2^-20}. The arithmetic-geometric mean gives \nolinkurl{sum c_tau alpha_tau>=f alpha_0 B_c}. Thus, before enlargement, \[
 e^{-\sum c\alpha}\le
 e^{-\theta f\alpha_0B_c}e^{-(1-\theta)\sum c\alpha}.      \tag{WJ27}
\] For the unique degree-zero scalar/transverse term use \nolinkurl{e^-theta f alpha_0 B_c <= (theta alpha_0 B_c)^-f}, because \nolinkurl{e^-u<=1/u} for every positive u. Integrating yields at most \nolinkurl{P_1(q/(theta Rk))^f}. For positive degrees one may drop the reserved damping and use WJ24--25. This proves convergence without a zeroth inverse moment; clamping alone did not make the total mass finite.

For any independently valid universal scalar-theta coefficient K in the original real dimension336, with \nolinkurl{c_0=2^36}, \nolinkurl{d=k+3}, the complete integrated row factor is consequently bounded by \[
 \binom{164}{k}\frac{U^f\pi_+^{df}}{D_0^d}
 \left\{
 \sum_{\substack{0\le l\le k,\ 0\le h\le3\\j=l+3-h\ge1}}
 \binom kl\binom3h K^{k-l+h}c_0^{-hf}
 \left(\frac{q}{(1-\theta)R}\right)^{jf}P_j
 +K^{k+3}c_0^{-3f}P_1\left(\frac{q}{\theta Rk}\right)^f
 \right\}.                                               \tag{WJ28}
\] This follows from K110 PM10--12's first untruncated Poisson inequalities, valid for every positive embedding weight, and the original saturated primal/dual Gram and covolume of K109. For \nolinkurl{k<=164}, every degree is at most167 and hence all displayed moments are finite. The final term is exactly \nolinkurl{(l,h)=(0,3)}; it is not optional. Multiply WJ28 by the coefficient in WJ17 and by125/104 from WJ22, and by all original arithmetic, finite, central-coset and support costs of the sector under consideration.

In particular WJ28 does not justify replacing the original correlated finite-pattern sum by a constant at degree167: its pattern exponent differs from the negative exponents used in DO121. Any such high-degree application must supply its own full four-state bound. It also supplies no new bound on the central affine fibre and does not remove its original750 cover or denominator correlation. Distinct pointwise trace bounds are not multiplied together without a new common proof.

\subsection{7. Verification, sources and algebraic-geometric scope}\label{verification-sources-and-algebraic-geometric-scope}

The companion executable checks the symbolic feasible-pair and Jacobian identities, exact weighted-tail constants, derivative signs, near-point and clamp losses, all166 finite inverse-moment recurrences and positive arguments, the exact highest-moment rational bound, and the complete polynomial degree bookkeeping through k164. These are finite exact diagnostics of the universal proof, not a substitute for reading the proof and not an actual global-arrow enumeration. It produces no error charge or canonical prefix mutation.

Direct inherited theorem interfaces: DO080 RT4--19 and section7; K107 SV4 and SV8--12, with its canonical boundary qualifications; K108 AE3--8; DO121 EC4--5 and EC8--13; K109 TG3--12; K110 PM10--12; and K113 SC27.5--15 only when an arithmetic sector is subsequently composed. Ordinary positive integral normalization: National Institute of Standards and Technology (n.d.), \emph{Digital Library of Mathematical Functions}, equations5.9.1,5.12.1/5.12.3,5.4.1/5.4.6, \url{https://dlmf.nist.gov/5.9.E1} , \url{https://dlmf.nist.gov/5.12} and \url{https://dlmf.nist.gov/5.4} . The new uses have their convergence domains proved explicitly above. Genuine selected action/metric is carried from Nakahama, R. (2016), \emph{Norm computation and analytic continuation of vector valued holomorphic discrete series representations}, Journal of Lie Theory26(4),927--990, section5.1 equations5.1/5.3/5.5/5.7 and Theorem5.1 pp960--961, in the established K043/DO080 normalization; no new module identification is claimed.

This is an estimate at the archimedean realization of the same projective OE modules and polarized CM abelian data. The positive measure, analytic weights and clamp do not alter the scheme over its original arithmetic base, the involution, Hermitian lattice, theta line, Picard scheme or Poincare biextension. Carried contextual anchors are SGA3 I4.7.3/4.7.3.1; SGA6 VI Theorem3.3; FGA232 Theorem3.1 pp148--149; FGA236 Corollaries1.7/2.3 pp229/232; SGA7 I VII2.7/2.9 pp169--174 and VIII7.4(b) pp310--311; cote4 scan003 sections8.1.1--8.1.2 and scans005--006 sections8.4.1--8.4.3; cote54 scans005--007. These locators are inherited context, not new primary readings or newly proved comparisons. No new formal-group/Dieudonne, lambda-ring/K-theory, l-adic weight, CM-motive or motivic L-function result follows from this real analytic inequality alone. DP\_MULT Q3e/DP44, THETA3 G1/A7, CX\_GRH CX-ALG-024 and GCT\_PLETH T-052 retain their current arithmetic and global gates until a separately checked composition changes them. No Lean build or archive theorem-validity audit is performed.

\subsection{Canonical algebraic-geometric context}\label{canonical-algebraic-geometric-context-67}

This is an estimate at the archimedean realization of the same projective OE modules and polarized CM abelian data. The positive measure, analytic weights and clamp do not alter the scheme over its original arithmetic base, the involution, Hermitian lattice, theta line, Picard scheme or Poincare biextension. Carried contextual anchors are SGA3 I4.7.3/4.7.3.1; SGA6 VI Theorem3.3; FGA232 Theorem3.1 pp148--149; FGA236 Corollaries1.7/2.3 pp229/232; SGA7 I VII2.7/2.9 pp169--174 and VIII7.4(b) pp310--311; cote4 scan003 sections8.1.1--8.1.2 and scans005--006 sections8.4.1--8.4.3; cote54 scans005--007. These locators are inherited context, not new primary readings or newly proved comparisons. No new formal-group/Dieudonne, lambda-ring/K-theory, l-adic weight, CM-motive or motivic L-function result follows from this real analytic inequality alone. DP\_MULT Q3e/DP44, THETA3 G1/A7, CX\_GRH CX-ALG-024 and GCT\_PLETH T-052 retain their current arithmetic and global gates until a separately checked composition changes them. No Lean build or archive theorem-validity audit is performed.

\clearpage

\section{AM-K-117 --- degree29 scalar theta and the total-eligibility28 remainder}\label{am-k-117-degree29-scalar-theta-and-the-total-eligibility28-remainder}

\textbf{Canonical checked integration --- 2026-11-23T10:30:52.475955+00:00.} Complete root and independent proof/script/interface reviews and separate exact replays passed. The actual canonical predecessor has130/31 entries; this integration gives131/31, complete floors27/11. Frozen author history is preserved. AM-K-117 D29.1-27 proves the universal degree29 scalar Gaussian lattice bound from the fixed exact normalized Gegenbauer polynomial and, with AM-R-098 WJ1-28, removes the complete actual support27 total-eligibility28 branch k164,b=j\_int0,2\^{}3107\textless x\textless2\^{}3328 with absolute normalized contribution below2\^{}-278. The entire660-term polynomial includes the repaired zero degree and the unique degree167 term with its full Q\^{}167 pattern cost. Every original167-entry reconstruction, paired ideal, selected metric, central affine750 cover, unique eta, support and all finite factors is paid; the one eligible inactive place keeps contraction1. One278 append preserves all130/31 predecessor entries, giving131/31 with both entire sums below2\^{}-230 and floors27/11 unchanged. Only the eligibility27 alternative140\textless=k\textless=164,2\^{}2111\textless x\textless2\^{}3328,b=j\_int0 remains at actual support27. No complete support28, signed nonvanishing, actual admitted source, selected arithmetic nonzero rule, prescribed period, evaluated DP44 cost or RH conclusion.

Author: \nolinkurl{/root/kummer_extension}. 23 November 2026 UTC. \textbf{Controlled pending complete proof.} R098 is the separately proved analytic input pinned in section4. No canonical promotion or charge is made by this file. Exact author, independent and root reviews remain separate events.

\subsection{1. Two statements and their exact scope}\label{two-statements-and-their-exact-scope}

The first statement is a universal scalar theorem in real dimension n=336, with f=168. It applies to every full-rank Euclidean lattice, every positive lower bound for its nonzero minimum and every positive Gaussian parameter. It is not an assertion about a finite-field metric or a whole arithmetic sector by itself.

The second statement concerns exactly the original J56 all-top, H133, globally rank-two-unipotent scalar-cross support27 branch with all active typesII. Retain every original source/target pencil, selected module, Hilbert metric, compact factor, word, twist, phase, projector, class mass and Haar normalization. The accepted K113 necessary window here is \[
 r_{\mathrm{elig}}=28,\quad k=164,\quad b=j_{\mathrm{int}}=0,
 \qquad2^{3107}<x=N_ED_\eta<2^{3328}.
 \tag{D29.1}
\] There is exactly one eligible inactive place; it is not assigned a noneligible coefficient. No total-eligibility27 window or error charge is imported.

\textbf{Theorem.} Using the separately proved R098 WJ17--28 analytic lemma, the complete containing sector D29.1 has absolute normalized contribution below2\^{}-278. The full polynomial has660 positive terms, including its repaired zero-degree term and the exceptional degree167 finite-pattern term. One278 charge suffices. This theorem makes no claim about the distinct eligibility27 k140..164 branch or complete support28.

Keep \[
 H=133,\quad R=17424/133,\quad Q=1427908193729121522601,
 \quad U=77Q^2-14,
\] \[
 D_0=7^{140}29^{144}5^{84}13^{84},\quad D_F=2^{168}D_0,
 \quad c_0=2^{36},\quad\theta=2^{-20}.
 \tag{D29.2}
\]

\subsection{2. A fixed rational polynomial and30 positive coefficients}\label{a-fixed-rational-polynomial-and30-positive-coefficients}

Put t=3261/10000 and p(X)=10\^{}-40 sum\_(i=0)\^{}14 a\_i X\^{}i, with the following fixed integers. They are mathematical data, not runtime optimization variables.

\begin{longtable}[]{@{}rr@{}}
\toprule\noalign{}
i & a\_i \\
\midrule\noalign{}
\endhead
\bottomrule\noalign{}
\endlastfoot
0 & 467035257309642805123331562 \\
1 & 80706128497300306343343393509 \\
2 & -896747485906017370732648699802 \\
3 & -68253127749041731200047829468619 \\
4 & 193469986188896577033338662619971 \\
5 & 14588671108851652483552680805364850 \\
6 & -2790200368251377807361018527495220 \\
7 & -1201746595827807685617021847335365024 \\
8 & -1268121273479902686867904738314760747 \\
9 & 43845687011100523654367843700241562010 \\
10 & 76766905895335943301603420494006647710 \\
11 & -707202963916244613740763137836577177204 \\
12 & -1539347306526975564077555218127409229230 \\
13 & 4077376859640910257041936285673548439273 \\
14 & 10000000000000000000000000000000000000000 \\
\end{longtable}

Define \[
 F(X)=(X-t)p(X)^2,\quad P_0=1,\quad P_1=X,
\] \[
 (\ell+n-2)P_{\ell+1}=(2\ell+n-2)XP_\ell-\ell P_{\ell-1}.
 \tag{D29.3}
\] The exact expansion F=sum\_(ell=0)\^{}29 f\_ell P\_ell has strictly positive coefficients. A compact exact sign certificate is2\^{}(e\_ell)\textless=f\_ell\textless2\^{}(e\_ell+1), with \[
\begin{split}
(e_0,\ldots,e_{29})={}&(-86,-78,-72,-66,-61,-57,-52,-48,-44,-41,\\
&-37,-34,-31,-28,-25,-23,-20,-18,-16,-14,\\
&-12,-10,-8,-7,-6,-4,-4,-3,-3,-2).
\end{split}
 \tag{D29.4}
\] Every comparison is an exact rational inequality. The executable reconstructs all30 rational coefficients from D29.3, checks all30 bounds, and recombines the complete power polynomial. No positivity inference from sampled X values is used. In particular \[
 f_0=\frac{17314247081621745624513323162146089564059753484541557086550620198717365927118078952293}
 {806689762279303805877092352000000000000000000000000000000000000000000000000000000000000000000000000000000000000}.
 \tag{D29.5}
\] Set M29=F(1)/f0. Explicitly, \[
 M_{29}=\frac{33747058477132956267553486841592409194296848334887458558564511205312018487139081178119693408828488034189049856}
 {752793351374858505413622746180264763654771890632241612458722617335537649005133867491}.
 \tag{D29.6}
\] A floating search suggested these fixed rational parameters. No optimality, floating inequality or numerical optimization is used after D29.3 is fixed.

\subsection{3. Universal harmonic, cap and Gaussian proof}\label{universal-harmonic-cap-and-gaussian-proof}

The normalized recurrence D29.3 gives P\_ell(1)=1 and \[
 (1-X^2)P_\ell''-(n-1)XP_\ell'
       +\ell(\ell+n-2)P_\ell=0.
 \tag{D29.7}
\] The executable checks these identities coefficient by coefficient through degree29. Here is the complete positive-kernel justification. Restrict the real homogeneous harmonic polynomials of degree ell to the sphere, use probability-area inner product, and choose an orthonormal basis Y\_i. This space is nonzero because the real part of (x1+i x2)\^{}ell is harmonic, including the constant at ell0. Its kernel sum\_i Y\_i(v)Y\_i(w) is a rotation-invariant Gram kernel. Its constant diagonal is h\_ell=dim H\_ell\textgreater0. For fixed v, invariance under the stabilizer of v makes the restriction a polynomial of degree at most ell in the inner product: every invariant homogeneous polynomial is a sum of powers of the axial coordinate and squared orthogonal radius. Computing its homogeneous Laplacian gives D29.7. The polynomial solution of degree at most ell is unique up to scale: the leading coefficient at a smaller degree j would have nonzero multiplier ell(ell+n-2)-j(j+n-2), and the remaining coefficients are recursively determined. Normalization at1 identifies the kernel with h\_ell P\_ell. Thus every P\_ell is positive semidefinite on unit-vector lists, including ell0.

For an m-point code with off-diagonal inner products at most t, F is nonpositive on those inner products because its second factor is a square. D29.4 and positive kernels give f0 m²\textless=sum\_(i,j)F(\textless v\_i,v\_j\textgreater)\textless=mF(1). Hence m\textless=M29. This is the normalized Delsarte--Goethals--Seidel addition-formula bound, with its proof included here.

Let X be1-separated in the Euclidean n-ball of radius rho\textgreater=1. Lift x to rho\^{}-1(x,sqrt(rho²-\textbar\textbar x\textbar\textbar²)) on the sphere in dimension n+1. Pairwise inner products are at most s\_rho=1-1/(2rho²). Choose a cap of pole height h with \[
 h^2=(s_\rho-t)/(1-t),\qquad
 \sin^2\vartheta=[2rho^2(1-t)]^{-1}.
 \tag{D29.8}
\] Since0\textless t\textless1/3 and s\_rho\textgreater=1/2, h\textgreater s\_rho. A pole point cannot share the cap with another code point. Otherwise normalize projection to the n-dimensional pole complement. The projected pairwise bound is(s\_rho-uv)/sqrt((1-u²)(1-v²)); each derivative is negative for u,v\textgreater=h, so its maximum is t at u=v=h. The cap occupancy is at most M29.

With c\_n=168 binom(336,168)/4\^{}168, the normalized cap area is c\_n integral\_0\^{}vartheta sin\^{}(n-1)(phi)dphi \textgreater=(c\_n/n)sin\^{}n(vartheta). Averaging over all poles gives \[
 \#X\le\frac{336M_{29}}{c_n}[2(1-t)]^{168}rho^{336}.
 \tag{D29.9}
\] The code dimension is336; the lift dimension337 does not replace it in the harmonic bound. For a lattice of nonzero minimum at least lambda, its nonzero count vanishes below rho1. Positive layer-cake integration and enlargement of the final majorant integral to0 give \[
 \Theta_\Lambda(s)\le1+C_{29}(s\lambda^2)^{-168},
\quad C_{29}=\frac{336M_{29}}{c_n}(6739/5000)^{168}168!,
\quad s>0.
 \tag{D29.10}
\] The zero vector is restored exactly once. There is no small-parameter or lattice-shape restriction.

For the actual half-trace ideal dual, lambda1²\textgreater=f D0\^{}(-2/f) z\^{}-2, with D0\^{}(1/f)\textless365729/500 and pi\textgreater314159/100000. Thus the all-parameter scalar coefficient is \[
 K_{29}=C_{29}
 \left[\frac{(365729/500)^2}{168(314159/100000)^2}\right]^{168}.
 \tag{D29.11}
\] The same theorem applies after any positive real diagonal embedding change, using the original algebraic norm minimum and AM--GM. The changed real lattice need not be an ideal. The full dimension336, half-trace factor2 and all covolumes remain.

With the accepted exact K15 coefficient, direct integer arithmetic gives \[
 32\left[\frac{M_{29}}{M_{15}}(6739/7841)^{168}\right]^2<1,
 \qquad K_{29}/K_{15}<2^{-5/2}.
 \tag{D29.12}
\] This scalar comparison is not substituted for the complete row calculation below.

\subsection{4. Required R098 trace and the clamped positive measure}\label{required-r098-trace-and-the-clamped-positive-measure}

This section records the precise external analytic input. Its complete proof is R098 WJ1--28, preserved at \nolinkurl{recovery_pending/root/weighted_radial_jacobian_clamped_profile.md}, SHA256 \nolinkurl{e08b11673224309bf263b91c059006c58b67f3b9f4303cf2d1a63a5d884bb11f}. The author and my separate whole proof/script/source review are recorded in their distinct packets; my unchanged56,291 replay passed10:15:33 UTC and the complete peer report was saved/read back at10:17:52 UTC. R098 retains the exact all-feasible Cartesian Jacobian, the complete selected radial probability measure, its whole-tail weighted Cauchy integral and both compact factors. For the actual original positive energies0\textless alpha\_tau\textless=R, WJ17/WJ22 supply the product trace bound \[
 |c(h)|<2\,2^{-649}\,2^{-933f/2}
             (2/15)^{84}(125/104)
 \int \exp(-\sum c_\tau\alpha_\tau)
                      \prod_\tau d\nu_\delta(c_\tau),
 \quad\delta=1/(1000R).
 \tag{D29.13}
\] Here mu166 is the positive convolution of Gamma(shape166,rate4) with ds/sqrt(pi s), and nu\_delta is its pushforward by c -\textgreater{} max(c,delta). The cap loss125/104 bounds exp(f/1000). The single near-point2 uses the exact elasticity bound and replaces the earlier factor. The actual alpha cap is used before enlargement; it is not asserted on newly added lattice tuples.

For1\textless=j\textless=166, clamping only decreases the inverse moments, and ordinary convergent integration gives \[
 N_j:=\int c^{-j}d\nu_\delta(c)
 \le4^{j-1/2}
 \frac{\Gamma(j-1/2)\Gamma(166.5-j)}{\Gamma(j)\Gamma(166)}.
 \tag{D29.14}
\] For j167, the unclamped moment would diverge. Its density satisfies \[
 \frac{d\mu_{166}}{dc}\le
       \frac{4^{166}}{\Gamma(166.5)}c^{165.5}.
\] Integrating below and above delta separately proves the finite bound \[
 N_{167}\le\frac{4^{166}}{\Gamma(166.5)}
                  \delta^{-1/2}\left(2+\frac2{333}\right).
 \tag{D29.15}
\] No finite zeroth moment is asserted. Let H\_j=(2j)!/(4\^{}j j!), pi\_+=22/7 and pi\_-=314159/100000. Rational upper bounds for N\_j\^{}f are \[
 P_j=(4\pi_+^2)^{84}
 \left[\frac{4^{j-1}H_{j-1}H_{166-j}}{(j-1)!165!}\right]^{168}
 \quad(1<=j<=166),
\] \[
 P_{167}=(1/(\pi_-\delta))^{84}
 \left[\frac{4^{166}(668/333)}{H_{166}}\right]^{168}.
 \tag{D29.16}
\] The inverse pi direction in the last formula is essential. All mixing tails remain inside these positive integrals.

\subsection{5. The original rows, central coset and all167 finite factors}\label{the-original-rows-central-coset-and-all167-finite-factors}

The support-independent original normal form and K113 give h\_ZZ=I,N²Z=0,N³=0, alpha=U r J\_Z\^{}-1 r*\textgreater0. The original row ideal is \[
 L=JI_\Sigma D_\eta,\quad N_EL=z^{2f},\quad
 N_EJ=z^{2f}T^{54}/x,\quad q=R/(Uz^2).
 \tag{D29.17}
\] All164 S values are nonzero; all167 original values belong to this same ideal at both CM partners. The actual saturated transverse Gram has detG'=U and G'\textgreater2\^{}36 I, with primal covolume D0\^{}167 z\^{}(334f)/U\^{}f.~No real change of weights alters this arithmetic membership.

On actual rows, product(alpha\_tau)\textgreater=(Rk/q)\^{}f.~For positive mixing coordinates put B=(product c\_tau)\^{}(1/f) and alpha0=Rk/q. Then AM--GM gives sum c\_tau alpha\_tau\textgreater=f alpha0 B. Before positive Gaussian enlargement reserve theta=2\^{}-20 of the exponent: \[
 e^{-\sum c_\tau\alpha_\tau}
 \le e^{-\theta f\alpha_0B}
       e^{-(1-\theta)\sum c_\tau\alpha_\tau}.
\] Only the second factor is enlarged to the full row lattice. For its positive degrees the first factor is bounded by1. For degree zero, y exp(-y)\textless=1 gives exp(-theta f alpha0 B)\textless=(theta alpha0 B)\^{}(-f). Tonelli therefore bounds its entire mixing integral by (N1/(theta alpha0))\^{}f\textless=P1{[}q/(theta Rk){]}\^{}f.~This treats the infinite total mass without replacing it by a constant.

Writing t0=Uz²=R/q, the complete anisotropic Gaussian factor before expansion is \[
 \{((1-\theta)t_0B)^{-f}+K_{29}\}^{k}
 ((1-\theta)t_0B)^{-3f}
 [1+K_{29}c_0^{-f}((1-\theta)t_0B)^f]^3,
\] times the true covolume prefactor displayed below. It follows by applying D29.10 to each changed scalar real lattice, using the original algebraic dual minimum and G'\textgreater c0 I. Expanding all scalar and transverse choices gives degree j=l+3-h. For every positive degree use D29.16, and for the unique zero degree use the preceding damping estimate. This yields the complete row polynomial \[
 {164\choose k}\frac{U^f\pi_+^{(k+3)f}}{D_0^{k+3}}
 \left\{\sum_{\substack{0<=l<=k,\ 0<=h<=3\\j=l+3-h>=1}}
 {k\choose l}{3\choose h}K_{29}^{k-l+h}c_0^{-hf}
 \left[\frac q{(1-\theta)R}\right]^{jf}P_j
 +K_{29}^{k+3}c_0^{-3f}P_1\left[\frac q{\theta Rk}\right]^f\right\}.
 \tag{D29.18}
\] Here k164, so the coordinate binomial equals1. No external k\^{}-84 or z\^{}-f reciprocal-trace factor is counted again. The complete actual phase space is contained in this lattice sum; no extra phase or beta norm-fibre multiplicity is omitted.

For each fixed27-place active support, the exact slope is unique throughout x\textless2\^{}3328 by the paired norm gap(4U)\^{}168 2\^{}24000\textless T\^{}54. Thus eta is fixed before counting its row and the full central affine coset. The extra eligible inactive place is determined by eta, not an independent variable or label. The original support count is binom(168,27)\textless2\^{}104, and we conservatively retain the old pivot164\textless2\^{}8.

The27 active weighted factors and140 noneligible inactive pointwise factors give \[
 (T-1)^{-167}<2T^{-167}.
 \tag{D29.19}
\] The remaining eligible inactive factor and both mixed contractions are bounded by1. The active one-label argument and the noneligible inactive argument have their original hypotheses b=j\_int0. No168th coefficient is inserted.

Retain the complete original Green and affine central bounds \[
 |a_G(h)|\le2Q^{-335\cdot167}(N_EJ)^{335/2},\quad
 N_EJ\le Q^{334},\quad N_EJ\le Q^{167}N_FJ_0,
\] \[
 \#d\le\frac{V_F(C(q))U^{-f}T^{54}x^{-1}}{N_FJ_0},\quad
 C(q)=H/2+750R/q,\quad
 V_F(C)=\frac{(\pi_+ f C^2)^{84}}{84!\sqrt{D_F}}.
 \tag{D29.20}
\] The real-part shift, dyadic conditions, full-real750 cover and primitive-u extraction remain. Eta, all167 row values and d reconstruct the entire original matrix; no source/target fibre is discarded. The resulting T power is335*27+54-167=8932.

\subsection{6. Both continuum directions and the full exceptional pattern}\label{both-continuum-directions-and-the-full-exceptional-pattern}

For positive degree j the complete q dependence before central correlation is \[
 F_j(q)=q^{(j-335/2)f}V_F(C(q)).
\] Its normalized logarithmic q derivative is j-335/2-epsilon\textless0 for every j\textless=167, where epsilon=(C-H/2)/C in(0,1). The repaired term has effective degree1. At \[
 q_*^f=\frac{R^fT^{54}}{U^fQ^{334}x},\qquad
 D=Q^{334}/N_EJ,\quad q=q_*D^{1/f},
 \tag{D29.21}
\] the derivative of x\^{}-337/2 F\_j(q\_*(x)) with respect to log x is-1-j+epsilon\textless0. All moment bounds, including the clamped167 bound, are independent of x and J. Therefore the entire containing count decreases with x and may be bounded at x=2\^{}3107.

The correlation(N\_FJ0)\^{}-1\textless=Q\^{}-167D gives pattern exponent j-333/2. For j\textless=166 this is at most-1/2. If s0=floor(sqrt(g0)), the exact inequality(s0+1)\^{}668\textless2s0\^{}668 proves \[
 \prod_{\mathfrak q|G}(1+q_{\mathfrak q}^{j-333/2})^2<2.
\] For the exceptional j167, only(l,h)=(164,0) occurs. Its full containing product is \[
 \prod_{\mathfrak q|G}(1+\sqrt{q_{\mathfrak q}})^2
 =Q^{167}\prod_{\mathfrak q|G}(1+q_{\mathfrak q}^{-1/2})^2
 <2Q^{167}.
 \tag{D29.22}
\] All334 original places and multiplicities1,2,1 are included. We conservatively retain a common factor4, and multiply that exceptional polynomial term by Q\^{}167. The repaired zero term uses the ordinary degree1 pattern. No unsupported uniform pattern estimate is used at degree167.

\subsection{7. Exact entire endpoint and single charge}\label{exact-entire-endpoint-and-single-charge}

The original integer inequalities give the fresh coupled cap \[
 4122987<q_*<4122988\quad\hbox{at }x=2^{3107},
\quad7^{84}T^{28}>[4(U+14)]^{168}2^{3107}.
 \tag{D29.23}
\] Write q\_-,q\_+ for those two integer bounds. The full moving central-volume ceiling is474. With the original directed logs LT=1019333/1024, LQ=287843/4096, LU=601356/4096 and LR=28810/4096, all common costs are \[
 C=117-933f/2-649+474+(335f/2)(L_R-L_U)
       +8932L_T-336\cdot167L_Q-(337/2)3107
   =211276557/512.
 \tag{D29.24}
\] The fixed117 retains support104,pivot8,Green1,T1,near-point1 and common-pattern2. The R098 trace gain and clamping cost remain explicit rational multipliers below; neither is hidden in C.

Put A=q\_+/((1-theta)R), V=K29/A\^{}f.~Exact strict ceilings are log2V\textless58 and log2(1/V)\textless-57. The dominant polynomial term is(l,h)=(0,0),j3, namely K29\^{}164 A\^{}(3f)P3. For any other positive degree the exact ratio to this term is \[
 {164\choose l}{3\choose h}
       V^{3-j}c_0^{-hf}P_j/P_3,
 \tag{D29.25}
\] with an additional Q\^{}167 at j167. The repaired j0 ratio is V³ c0\^{}-3f {[}q\_+/(theta Rk){]}\^{}f P1/P3. All659 nonmain terms have certified relative upper exponent at most-991. In particular their full sum is below659*2\^{}-100\textless1 times the dominant term. These are exact rational comparisons of the complete coefficients, not an asymptotic scalar plateau approximation.

Consequently the complete rational factor satisfies \[
 \frac{2\pi_+^{167f}}{D_0^{167}q_-^{335f/2}}
 K_{29}^{164}A^{3f}P_3\,(2/15)^{84}(125/104)
 <2^{-412928}.
 \tag{D29.26}
\] Thus the entire sector has exponent \[
 C-412928=-142579/512<-278.
 \tag{D29.27}
\] There is only this one k and total-eligibility branch. All supports, labels, original patterns and affine fibres have already been counted. No further coordinate or inactive-place union is omitted. One278 charge is therefore sufficient after the required whole reviews.

The final executable must preserve the actual DO121130/31 prefix, or a later complete ordered extension, append278 once, retain all31 uniform entries, and check both whole dyadic sums below2\^{}-230. Floors remain27/11. Historical author and peer outputs are not overwritten when chronology advances.

\subsection{8. Primary and literal source boundaries}\label{primary-and-literal-source-boundaries}

The scalar primary input is Delsarte, Goethals, and Seidel(1977), Definition2.1, Theorems2.4 and3.2--3.3, Corollary3.8 and Theorem4.3 with their proofs, printed364--369. The normalized harmonic proof, cap projection and shell integration are included above. The exact actual half-trace minimum and anisotropic interpretation are the accepted DO070 FT3, K109 TG3--7, K110 PM10--12 and K111 HP12--18 interfaces. No numerical sector conclusion is inferred merely from a better scalar coefficient.

R098 is the separately required whole analytic source for D29.13. The moment bounds D29.14--16 use ordinary positive Gamma/Beta integrals, DLMF5.9.1,5.12.1/3 and5.4.6. The true original lattice, finite factors and central source are K113 SC27.1--21, K116 RE11--22, DO080 section7, K107 SV2--19, K108 AE2--15, DO098 CF4--16 and DO104 CD10--19. The DO121 eligibility27 numerical charge is neither replaced nor reused for this eligibility28 count.

The carried projective modules, involution, polarization, theta/Picard/Poincare data and selected realization remain fixed. SGA3 I4.7.3/4.7.3.1; SGA6 VI Theorem3.3; FGA232 Theorem3.1 pp148--149; FGA236 Corollaries1.7/2.3 pp229/232; SGA7 I VII2.7/2.9 pp169--174 and VIII7.4(b) pp310--311; cote4 scan003 sections8.1.1--8.1.2/scans005--006 sections8.4.1--8.4.3; and cote54 scans005--007 are inherited context. DP\_MULT Q3e/DP44, THETA3 G1/A7, CX\_GRH CX-ALG-024 and GCT\_PLETH T-052 keep their signed/source/selected-image/period/global-comparison gates. No archive theorem-validity audit or Lean build is undertaken.

Delsarte, P., Goethals, J. M., \& Seidel, J. J.(1977). Spherical codes and designs. \emph{Geometriae Dedicata,6}(3),363--388. \url{https://doi.org/10.1007/BF03187604}

National Institute of Standards and Technology.(n.d.). \emph{Digital Library of Mathematical Functions}, §§5.4,5.9,5.12. \url{https://dlmf.nist.gov/5.9}

\subsection{Canonical algebraic-geometric context}\label{canonical-algebraic-geometric-context-68}

The fixed harmonic positive-kernel polynomial controls the same real scalar dual lattice after arbitrary positive diagonal rescaling; no transformed lattice is asserted to be an arithmetic ideal. The actual projective OE modules, half-trace pairing, CM involution, polarized theta line, Picard and Poincare data stay unchanged. Delsarte, Goethals, and Seidel (1977), Definition2.1, Theorem2.4, Theorems3.2-3.3, Corollary3.8 and Theorem4.3, printed364-368, provide the exact primary normalization context. NIST DLMF5.9.1,5.12.1/3,5.4.6 are used only with positive real integral arguments, including the full clamped degree167 atom and tail. Inherited context only: SGA3 I4.7.3/4.7.3.1; SGA6 VI Theorem3.3; FGA232 Theorem3.1 pp148-149; FGA236 Corollaries1.7/2.3 pp229/232; SGA7 I VII2.7/2.9 pp169-174 and VIII7.4(b) pp310-311; cote4 scan003 sections8.1.1-8.1.2 and scans005-006 sections8.4.1-8.4.3; cote54 scans005-007. No new formal-group/Dieudonne, lambda-ring/K-theory, l-adic-weight or CM-motive comparison. DP\_MULT Q3e/DP44 and THETA3 G1/A7 receive the strictly smaller conditional finite gate; CX\_GRH CX-ALG-024 and GCT\_PLETH T-052 retain their arithmetic/global gates.

\clearpage

\section{AM-DO-122 --- the original central-energy coset at support27}\label{am-do-122-the-original-central-energy-coset-at-support27}

\textbf{Canonical checked integration --- 2026-11-23T10:32:01.596526+00:00.} Complete root and independent proof/script/interface reviews and separate exact replays passed. The actual canonical predecessor has131/31 entries; this integration gives132/31, complete floors27/11. Frozen author history is preserved. AM-DO-122 CE1-29 removes the complete actual support27 total-eligibility27 all-II globally rank-two-unipotent scalar-cross k140 sector with absolute contribution below2\^{}-359. MANDATORY CUTOFF QUALIFICATION: the prose after CE8 is read as norm(h\_tau)\textless=H=133, as proved by R048-RC RC48.4/RC48.6 and the paragraph after RC48.7; CE8 already has the correct non-strict sign. The frozen proof is read together with recovery\_pending/definite/DO122\_CUTOFF\_SOURCE\_ADDENDUM.md, whose full text is included in the checked paper. The joint physical central variable w=-U Im d, original affine ideal, full750 cover, weighted Voronoi integral and U\^{}f physical Jacobian are retained. The central integral consumes one profile power, giving shape166; all moments used and every positive polynomial term are paid, with original-row zero repair, all finite factors and correlated patterns. One359 append to actual K117131/31 gives132/31; both entire sums remain below2\^{}-230 and floors27/11 stay unchanged. The sole support27 necessary alternative is now total eligibility27,141\textless=k\textless=164,2\^{}2111\textless x\textless2\^{}3328,b=j\_int0. The positive k141 envelope is only a failed sufficient bound. No full support28, signed nonvanishing, actual admitted source, selected arithmetic content, prescribed period, evaluated DP44 or RH claim.

Author: \nolinkurl{/root/definite_oscillator}. 23 November2026 UTC. \textbf{Controlled pending proof.} Author, independent and root whole-proof/source/script checks are distinct from exact diagnostic replays. No canonical artifact is changed here.

\subsection{1. Exact sector and complete theorem}\label{exact-sector-and-complete-theorem}

Retain the original J56 all-top/H133 kernel, actual CM fields and involution, source/target pencils and labels, full selected vector-valued module and Hilbert metric, both compact factors, phases, canonical twists, Haar measures, class masses and every finite projector. Set \[
 f=a=168,\quad H=133,\quad R=17424/133,\quad
 Q=1427908193729121522601,\quad U=77Q^2-14,
\] \[
 D_0=7^{140}29^{144}5^{84}13^{84},\qquad D_F=2^{168}D_0.
 \tag{CE1}
\] The actual predecessor is canonical DO121,130 all-top/31 uniform entries and complete floors27/11. Work in its remaining globally rank-two-unipotent scalar-cross support27 sector, with all active typesII and total eligibility27: \[
 b_{\rm inactive}=j_{\rm int}=0,\qquad
 2^{2111}<x=N_ED_\eta<2^{3328},\qquad k=140.
 \tag{CE2}
\] Here k is the number of nonzero original S covector values. All334 unmixed G reductions are unipotent. The separate total-eligibility28,k164 branch is not included.

\textbf{Theorem.} The complete actual contribution of CE2 has absolute normalized value below2\^{}-359. Append one359 charge, retaining all previous ordered charges and both complete sums. Neither a k141 removal nor a complete support28 assertion follows. The k141 positive sufficient envelope is retained as a diagnostic only.

The new step sums a decaying majorant over the original affine central coset, using its full physical scaling and covering radius. The row, central and source factors are combined only after fixing the actual slope and whole-C pattern.

\subsection{2. A joint actual trace bound in energy and the original central entry}\label{a-joint-actual-trace-bound-in-energy-and-the-original-central-entry}

The accepted full-positive normal form, K113 SC27.1--7 and DO098 CF4--6, gives \[
 h=I+N,\quad N^3=0,\quad \operatorname{rank}_E N=2,
 \quad\alpha=(N^2)_{nn}=U rJ_Z^{-1}r^*>0,
 \quad0<\alpha_\tau\le R.
 \tag{CE3}
\] In that exact original normal form the central scalar d satisfies \[
 N_{nn}=-Ud,\qquad d+\bar d=-rJ_Z^{-1}r^*,
 \qquad w:=\Im N_{nn}=-U\Im d.
 \tag{CE4}
\] Thus, at each real label, h\_nn=1+alpha/2+iw. The real variable w here is the imaginary part of the original negative entry; it is not an independently normalized Heisenberg coordinate.

The unipotent inverse-adjoint identity also holds at the actual near point h\_p=g\_p\^{}-1 h g\_p.~Indeed h\_p\^{}-1=I−N\_p+N\_p² and (h\_p\^{}-1)\_nn=overline(h\_p,nn) imply \[
 \alpha_p=(N_p^2)_{nn}=2\Re N_{p,nn},\quad
 c_p=\cosh r(h_p)=|h_{p,nn}|,
 \quad c_p^2=(1+\alpha_p/2)^2+w_p^2.
 \tag{CE5}
\] These are the exact DO080 section7 identities retained in DO121 EC3--5. Put b(t)=1+t/4 and p=333/4. For t\textgreater0 and real v, let c=sqrt((1+t/2)\^{}2+v²). Then \[
 c^2-1=t\,b(t)+v^2\ge t\,b(t),\qquad
 (c+1)^2\ge4b(t)^2+v^2,
 \tag{CE6}
\] since the difference in the second inequality is2(c−1−t/2)\textgreater=0. The accepted selected coefficient K108 AE9, including the entire K107 lower-tail absorption, therefore gives \[
 |c_W(h_p)|\le\vartheta_a2^{-933/2}F(\alpha_p,w_p),
\] \[
 F(t,v)=t^{-1/2}b(t)^{-167}
       \left(1+\frac{v^2}{4b(t)^2}\right)^{-p}.
 \tag{CE7}
\] The exponent167 combines the angular166.5 with the reciprocal-sinh half power. This is a bound for the original selected trace with arbitrary compact factors, not a replacement scalar model.

We require the original w, rather than w\_p, for the arithmetic lattice. At an archimedean place choose the original real orthonormal frame fixing the negative vector e\_n.~Let v=g\_p e\_n and multiply v by a unit scalar so its negative component is the positive number cosh r\_p; the quadratic entry is unchanged. Its positive part has norm sinh r\_p.~For0\textless=r\_p\textless1/2, \[
 \|v-e_n\|\le e^{r_p}-1\le2r_p,\quad
 \|v\|\le e^{r_p}<2,\quad
 |h_{p,nn}-h_{nn}|\le6Hr_p.
 \tag{CE8}
\] The last inequality follows by expanding the quadratic J-pairing and using the actual operator cutoff \textbar\textbar h\textbar\textbar\textless H. It needs no near-identity estimate for the full g\_p or its compact factors. DO080 section7 separately gives alpha\_p\textgreater=exp(−2r\_p)alpha.

For u=v²/(4b(t)²+v²) in{[}0,1{]}, direct differentiation yields \[
 -\partial_{\log t}\log F
 =\frac12+(167-2pu)\frac{t}{t+4}\in(0,335/2),
 \qquad |\partial_v\log F|\le p/(2b(t))\le p/2.
 \tag{CE9}
\] Consequently \[
 F(\alpha_p,w_p)
 \le\exp((335+3Hp)r_p)F(\alpha,w),\quad
 335+3Hp<2^{16}.
 \tag{CE10}
\] The literal r\_p\textless2\^{}-352 makes exp(2\^{}16 f 2\^{}-352)\textless2. This is the single old global near-point2, replaced by a proved joint transport bound; no second near-point factor is multiplied. With vartheta\_a\^{}168\textless2\^{}-649, \[
 |c(h)|<2\,2^{-649}\,2^{-933f/2}
 \prod_\tau F(\alpha_\tau,w_\tau).
 \tag{CE11}
\]

\subsection{3. Original paired difference ideal and physical covolume}\label{original-paired-difference-ideal-and-physical-covolume}

Fix the actual eta, all167 original covector values r, and the common whole-C G pattern J. Eta,r,d determine the entire original matrix injectively through DO098 CF4. There is no projected-column or isometry fibre left after these data are fixed.

Let J0 be the F-ideal containing each rank-zero G prime once. The original finite table, with local norm-exponent states2,1,0 and multiplicities1,2,1, gives \[
 N_EJ\le Q^{334},\qquad N_EJ\le Q^{167}N_FJ_0,
 \qquad N_FJ_0\le Q^{167}.
 \tag{CE12}
\] Every original row value lies in L=J I\_Sigma D\_eta, with N\_E I\_Sigma=T\^{}-54. This includes both CM partners and every integral saturated transverse input. Write \[
 N_EL=z^{2f}=N_EJ\,xT^{-54},\quad
 q=R/(Uz^2)\ge k.
 \tag{CE13}
\]

For two admitted d values with the fixed preceding data, Delta d is pure imaginary. In the original four-positive-plus-negative block the differences are (Delta d)u\_i u\_j, −bar(eta)(Delta d)u\_i, eta(Delta d)u\_j and −U Delta d.~Primitive integer Bezout extraction gives Delta d,eta Delta d,bar(eta)Delta d in the original entry ideal J0 O\_E I\_Sigma². This is the support-independent DO104 CD14--16 argument, explicitly applied at27 through K113 SC27.7. At each paired prime, if the eta valuations are A,e−A with e\textgreater=0, the extra exponent is max(0,−A,A−e), exactly that of D\_eta bar(D\_eta). Its contraction is B\_eta with N\_F B\_eta=x. Thus the d fibre is contained in one affine coset of \[
 \Gamma_d=iJ_0B_\eta I_{\Sigma,F}^2,\qquad
 \operatorname{covol}\Gamma_d
 =\sqrt{D_F}(N_FJ_0)xT^{-54}.
 \tag{CE14}
\] At an inert prime there is no extra denominator exponent. This remains valid if D\_eta meets U,G or T. The real-part shift is retained: neither d nor Im d is asserted to belong to the difference ideal. No division by2 or byU is used in the primitive extraction.

Passing to the physical variable w=−U Im d multiplies real covolume by U\^{}f.~The embeddings of i contribute only signs, so the containing affine lattice has \[
 \operatorname{covol}\Gamma_w
 =U^f\sqrt{D_F}(N_FJ_0)xT^{-54},\qquad
 \mu(\Gamma_w)<750\sqrt f\,U z_c,
\] \[
 z_c^f=(N_FJ_0)xT^{-54},\qquad
 U z_c=\frac Rq\left(\frac{N_FJ_0}{N_EJ}\right)^{1/f}.
 \tag{CE15}
\] The full-real750 cover, including the full real codifferent/discriminant normalization, is the accepted DO098 CF12--13 / DO104 CD17 bound. The last identity displays the exact q and J0 correlation; a larger row ideal does not alter the central lattice.

Here the literal global norm bound is much stronger than the coarse estimate750R/q. By CE12, x\textless2\^{}3328 and the exact actual-host inequality \[
 U^fQ^{167}2^{3328+83f}<T^{54},
 \qquad U z_c<2^{-83}.
 \tag{CE16}
\] This holds simultaneously for every actual q,J,J0 in CE2. It is obtained directly from the central ideal, rather than combining incompatible extrema of q and a freely varying divisor ratio.

\subsection{4. A complete weighted affine-coset bound}\label{a-complete-weighted-affine-coset-bound}

For a fixed original row put b\_tau=1+alpha\_tau/4 and \[
 \phi(w)=\prod_\tau(1+w_\tau^2/(4b_\tau^2))^{-p}.
\] EquationCE9 gives \textbar\textbar grad log phi\textbar\textbar\textless=p sqrt(f)/2 everywhere in R\^{}f.~Each translated Voronoi cell of Gamma\_w has the same physical covolume and lies within the full covering radius of its center. Hence for every y in the cell centered at w, phi(w)\textless=exp(375pf U z\_c)phi(y). Summing the disjoint cell integrals gives the complete affine bound \[
 \sum_{w\in w_0+\Gamma_w}\phi(w)
 \le\frac{\exp(375pf U z_c)}{\operatorname{covol}\Gamma_w}
 \prod_\tau\int_{\mathbb R}(1+t^2/(4b_\tau^2))^{-p}\,dt.
 \tag{CE17}
\] All integrals converge because p\textgreater1/2; monotone convergence justifies the count over the whole infinite coset. Actual admitted central points form a subset. No assertion makes every extended lattice point an original arrow or assigns it a finite coefficient. We sum the positive analytic upper envelope, while keeping the fixed actual pattern's finite coefficient bound.

Substitution followed by the ordinary positive Beta integral gives \[
 \int_{\mathbb R}(1+t^2/(4b^2))^{-p}\,dt
 =2b\sqrt\pi\,\frac{\Gamma(p-1/2)}{\Gamma(p)}
 \le\mathcal G b,\qquad
 \mathcal G=4\sqrt{\pi/329}.
 \tag{CE18}
\] For completeness, Cauchy--Schwarz in the ordinary Gamma integral proves Gamma(p−1/2)²\textless=Gamma(p−1)Gamma(p); recurrence gives the ratio\textless=1/sqrt(p−1), with p−1=329/4\textgreater0. No asymptotic Gamma estimate is used. NIST DLMF5.9.1 and5.12.1/3 provide the positive integral normalization.

The exact check375pf2\^{}-83\textless1/2 and exp(1/2)\textless2 bound the entire Voronoi inflation by a single additional2. Combining CE11,CE14--18 yields \[
 \sum_{\text{actual central fibre}}|c(h)|
 <4\,2^{-649}\,2^{-933f/2}
 \frac{\mathcal G^f T^{54}}{U^f\sqrt{D_F}(N_FJ_0)x}
 \prod_\tau\alpha_\tau^{-1/2}(1+\alpha_\tau/4)^{-166}.
 \tag{CE19}
\] The two factors2 in4 are the genuine near-point transport and the weighted Voronoi comparison. The central integral consumes one power of1+alpha/4. It is therefore shape166, rather than the shape167 before this summation. The U\^{}f physical Jacobian remains explicit.

\subsection{5. All row terms, the zero-degree repair, and finite costs}\label{all-row-terms-the-zero-degree-repair-and-finite-costs}

Let mu be the convolution of the Gamma probability measure of shape m=166,rate4 and ds/sqrt(pi s) on the positive half-line. It has infinite total mass, and \[
 t^{-1/2}(1+t/4)^{-166}=\int e^{-ct}\,d\mu(c),\qquad
 M_j=4^{j-1/2}\frac{\Gamma(j-1/2)\Gamma(166.5-j)}
 {\Gamma(j)\Gamma(166)},\quad1\le j\le166.
 \tag{CE20}
\] Both Gamma arguments are positive. There is no M0 or finite167th moment claimed. For H\_n=(2n)!/(4\^{}n n!), pi\_+=22/7, exact rational bounds are \[
 P_j=(4\pi_+^2)^{84}
 \left[\frac{4^{j-1}H_{j-1}H_{166-j}}
 {(j-1)!165!}\right]^{168}\ge M_j^f.
 \tag{CE21}
\]

The actual genuine-row norm floor is prod alpha\_tau\textgreater=(Ukz²)\^{}f=(Rk/q)\^{}f.~Reserve theta=2\^{}-20 of sum c\_tau alpha\_tau before enlarging the row. With B=(prod c\_tau)\^{}(1/f), the reserved factor is exp(−theta fRkB/q). For the sole zero-degree term use e\^{}-u\textless=u\^{}-1 to obtain \[
 \int e^{-\theta fRkB/q}\,d\mu^{\otimes f}
 \le P_1\left(\frac q{\theta Rk}\right)^f.
 \tag{CE22}
\] For positive degrees, drop this reserved factor and integrate their finite moments. This is the actual K116/DO121 repair; neither a per-embedding alpha floor nor a finite mass of mu is inferred. The original cutoff0\textless alpha\textless=R was used for actual matrices before enlargement. No zero or out-of-cutoff row introduced later receives an unproved arithmetic premise.

Retain the exact K111 scalar coefficient and actual saturated transverse metric: \[
 c_{336}=168\binom{336}{168}/4^{168},\quad
 M_{15}=\frac{6015684926187167549292201834533348561577984}
 {2618017117694612446410293519},
\] \[
 K=\frac{336M_{15}}{c_{336}}(7841/5000)^{168}168!
 \left[\frac{(365729/500)^2}{168(314159/100000)^2}\right]^{168},
 \qquad\det G'=U,\quad G'>c_0I,\quad c_0=2^{36}.
 \tag{CE23}
\] The first untruncated Poisson inequalities K110 PM10--12 apply to every positive embedding weight. With d\_k=k+3 they give, after the preceding actual-row restriction, the full row factor binom(164,k)U\^{}f pi\_+\^{}(d\_k f)/D0\^{}(d\_k) times \[
 \mathcal P_k(q)=
 \sum_{\substack{0\le l\le k,\ 0\le h\le3\\j=l+3-h\ge1}}
 \binom kl\binom3h K^{k-l+h}c_0^{-hf}
 \left(\frac q{(1-\theta)R}\right)^{jf}P_j
 +K^{k+3}c_0^{-3f}P_1\left(\frac q{\theta Rk}\right)^f.
 \tag{CE24}
\] Every scalar term and all four transverse terms are retained, including the unique(l,h)=(0,3) repaired term. At k140 the largest moment is143; the diagnostic141 uses144. The row U\^{}f cancels exactly the physical central Jacobian U\^{}-f in CE19.

The actual bounded uniqueness K113 SC27.4 supplies at most one eta per original27-place support throughout x\textless2\^{}3328. Fixing it before the row and central sum is legitimate. DO105 EL10--16 fixes the only possible active label from eta at both CM partners, including central-only poles; its support-independent local proof is already applied at27 in K113 SC27.9. In CE2, j\_int=0 additionally makes each active cross pole persist when the central coordinate varies. Thus the weighted central bound is for one containing affine coset for the prescribed original pair, not a union over uncounted central-dependent labels. The141 inactive pointwise factors are fixed by the same cross data, as in DO108 AC8--10's local argument. The complete T factor is(T−1)\^{}-168\textless2T\^{}-168. The finite coefficient is2Q\^{}(-335·167)(N\_EJ)\^{}(335/2). Support choices\textless2\^{}104, original harmless first-pivot choices\textless2\^{}8, Green2,T2,near-point2,weighted-cover2 and conservative pattern4 give118 fixed bits. No beta norm fibre is independently counted after eta is fixed, and no source or negative-column fibre is silently omitted: the original reconstruction above is injective.

\subsection{6. Uniform q,x and four-state control}\label{uniform-qx-and-four-state-control}

Collecting all factors except the fixed central constant gives, for each positive degree j, and for effective degree1 of the repaired term, \[
 x^{-337/2}q^{(j-335/2)f},\qquad
 q_*^f=\frac{R^fT^{54}}{U^fQ^{334}x}.
 \tag{CE25}
\] Its normalized q derivative is j−335/2\textless0. After substituting q\_*(x), its x logarithmic derivative is exactly−1−j\textless0. The uniform cover bound CE16 has already been paid, so no untracked moving-radius or q-dependent exponential remains.

For D=Q\^{}334/N\_EJ, one has q=q\_*D\^{}(1/f) and (N\_FJ0)\^{}-1\textless=Q\^{}-167 D. For degree j the whole original four-state sum is bounded by \[
 \prod_{v=1}^{334}(1+q_v^{j-333/2})^2<2,
 \qquad j\le144,
 \tag{CE26}
\] using j−333/2\textless=−22.5\textless−1 and the exact old g0 product inequality. We retain4 conservatively. There is no additional orientation or independent pattern factor. This statement does not extend the present argument to the divergent shape166 degree167 or the exceptional high-degree pattern branch.

\subsection{7. Complete exact endpoint and chronology}\label{complete-exact-endpoint-and-chronology}

The unchanged actual-host directed logarithms are L\_T=1019333/1024,L\_Q=287843/4096,L\_U=601356/4096,L\_R=28810/4096. At the coupled endpoint x=2\^{}2111, \[
 251124964<q_*<251124965,\qquad
 \frac{\mathcal G^f}{\sqrt{D_F}}
 \le\frac{(16\pi_+/329)^{84}}{\sqrt{D_F}}<2^{-1110}.
 \tag{CE27}
\] The T power is335·27+54−168=8931. All finite, selected and central factors yield the complete common exponent \[
 C=118-933f/2-649-1110+(335f/2)(L_R-L_U)
       +8931L_T-336\cdot167L_Q-(337/2)2111
   =591766613/1024.
 \tag{CE28}
\] There is no shift loss and no previous central ball volume474. The sole new cover2 is already included in118. Exact rational moment ceilings and the scalar-ratio ceiling−935 bound all nonmain terms of CE24 by the main one. For k140 the563 nonmain terms each have relative upper exponent at most−496; their entire sum is below the main term. The complete cell includes the conservative factor2 multiplying that main term: \[
 E_k=C+\left\lceil\log_2\left[
 2\binom{164}k\frac{\pi_+^{(k+3)f}}{D_0^{k+3}}
 q_-^{-335f/2}
 \left(\frac{q_+}{(1-\theta)R}\right)^{(k+3)f}P_{k+3}
 \right]\right\rceil_{\rm strict}.
 \tag{CE29}
\]

\begin{longtable}[]{@{}
  >{\raggedleft\arraybackslash}p{(\columnwidth - 6\tabcolsep) * \real{0.2667}}
  >{\raggedleft\arraybackslash}p{(\columnwidth - 6\tabcolsep) * \real{0.2667}}
  >{\raggedleft\arraybackslash}p{(\columnwidth - 6\tabcolsep) * \real{0.2667}}
  >{\raggedright\arraybackslash}p{(\columnwidth - 6\tabcolsep) * \real{0.2000}}@{}}
\toprule\noalign{}
\begin{minipage}[b]{\linewidth}\raggedleft
k
\end{minipage} & \begin{minipage}[b]{\linewidth}\raggedleft
Full positive polynomial factor ceiling
\end{minipage} & \begin{minipage}[b]{\linewidth}\raggedleft
Complete exponent
\end{minipage} & \begin{minipage}[b]{\linewidth}\raggedright
Status
\end{minipage} \\
\midrule\noalign{}
\endhead
\bottomrule\noalign{}
\endlastfoot
140 & −578257 & −368555/1024 & Strictly below−359; one complete sector \\
141 & −576494 & 1436757/1024 & Positive sufficient envelope; no removal \\
\end{longtable}

Every coordinate-position binomial, active support, retained pivot, unique eta, actual affine central fibre, source reconstruction, G pattern and local T factor is already in the first row. There is no remaining endpoint or position union. Append359 exactly once to the actual DO121130/31 prefix, giving131/31 with complete floors27/11 unchanged and both full dyadic sums\textless2\^{}-230. If an accepted later extension precedes this result, the executable preserves its entire ordered prefix before appending359.

\subsection{8. Sources and retained boundaries}\label{sources-and-retained-boundaries}

Direct inputs are DO080 section7; K108 AE8--13; K107 SV8--19 with its mandatory boundary qualifications; DO098 CF4--16; DO104 CD10--19; DO105 EL10--16; DO108 AC8--10's local fixed-cross argument; K113 SC27.1--15; K110 PM10--12's first untruncated inequalities; K111 HP12--18; K109 TG3--12; and canonical DO121 EC11--24. The new joint transport, weighted affine comparison and exact central integral are proved in CE5--19. NIST DLMF5.9.1 and5.12.1/3 are used only with positive ordinary-integral arguments; the Gamma-ratio bound is proved by Cauchy--Schwarz above.

National Institute of Standards and Technology. (n.d.). \emph{Digital Library of Mathematical Functions}, §§5.9 and5.12. \url{https://dlmf.nist.gov/5.9.E1} and \url{https://dlmf.nist.gov/5.12.E3}.

Nakahama, R. (2016). Norm computation and analytic continuation of vector valued holomorphic discrete series representations. \emph{Journal of Lie Theory,26}(4),927--990. \url{https://arxiv.org/abs/1506.05919}. The literal vector-valued action and metric enter through the accepted DO080/K107/K108 interfaces; this proof does not replace their selected module.

All original projective O\_E modules, involution, polarization, theta/Picard/Poincare data and selected realization remain. Inherited context only: SGA3 I4.7.3/4.7.3.1; SGA6 VI Theorem3.3; FGA232 Theorem3.1 pp148--149; FGA236 Corollaries1.7/2.3 pp229/232; SGA7 I VII2.7/2.9 pp169--174 and VIII7.4(b) pp310--311; cote4 scan003 sections8.1.1--8.1.2 and scans005--006 sections8.4.1--8.4.3; cote54 scans005--007. DP\_MULT Q3e/DP44, THETA3 G1/A7, CX\_GRH CX-ALG-024 and GCT\_PLETH T-052 retain their signed/source/selected-content/prescribed-period/global comparison gates. No fresh reading of those inherited context sources, actual arrow existence, numerical q28 result, complete support28 theorem or signed nonvanishing is asserted.

\subsection{Canonical algebraic-geometric context}\label{canonical-algebraic-geometric-context-69}

The physical coordinate change is a real analytic treatment of the same original central affine module, preserving its shift, CM partner conditions and polarization. It adds no arithmetic parameter or field and changes no selected metric. NIST DLMF5.9.1,5.12.1/3 and5.4.6 are applied with explicit positive integral arguments. Inherited context only: SGA3 I4.7.3/4.7.3.1; SGA6 VI Theorem3.3; FGA232 Theorem3.1 pp148-149; FGA236 Corollaries1.7/2.3 pp229/232; SGA7 I VII2.7/2.9 pp169-174 and VIII7.4(b) pp310-311; cote4 scan003 sections8.1.1-8.1.2 and scans005-006 sections8.4.1-8.4.3; cote54 scans005-007. No new formal-group/Dieudonne, lambda-ring/K-theory, l-adic-weight or CM-motive comparison. DP\_MULT Q3e/DP44 and THETA3 G1/A7 receive the strictly smaller conditional finite gate; CX\_GRH CX-ALG-024 and GCT\_PLETH T-052 retain their arithmetic/global gates.

\clearpage

\section{AM-DO-122 --- exact operator-cutoff locator and harmless boundary precision}\label{am-do-122-exact-operator-cutoff-locator-and-harmless-boundary-precision}

Author: \nolinkurl{/root/definite_oscillator}. Saved 2026-11-23T10:18:13.850289+00:00 after the parent's source-precision request. Frozen proof/script bytes remain unchanged.

The retained original cutoff is \textbf{R048-RC}, \nolinkurl{research/reconstructed_r048_all_top_small_cutoff.md}, §RC48.4, equationRC48.6 and the paragraph immediately following RC48.7 (local lines98--110). The complete RC48.1--4 source was freshly read for this clarification. Its equation sums the original rational support representatives g with max\_tau log\textbar\textbar g\_tau\textbar\textbar\textgreater2R+1/100. Since2R+1/100=489/100 and exp(489/100)\textless133, every original matrix with max\_tau\textbar\textbar g\_tau\textbar\textbar\textgreater133 lies in the charged annulus. These are the rational h representatives before actual near-point conjugation. Their unitary central representatives have the same operator norm.

The literal retained bound is therefore \textbf{\textbar\textbar h\_tau\textbar\textbar\textless=H=133}. In the sentence immediately after CE8, read the prose phrase \nolinkurl{actual operator cutoff ||h||<H} as \nolinkurl{actual operator cutoff ||h||<=H}. The displayed CE8 is already non-strict:

\[
 |h_{p,nn}-h_{nn}|\le6Hr_p.
\]

Its quadratic-pairing proof uses only\textbar\textbar h\textbar\textbar\textless=H, \textbar\textbar v−e\_n\textbar\textbar\textless=2r\_p and\textbar\textbar v\textbar\textbar+1\textless3. At r\_p=0 both sides vanish. Every later inequality, the joint transport coefficient335+3Hp, the single nearpoint2, the exact central-cover/row count and the359 charge are unchanged. No tail is charged again and no cutoff is asserted for the independent uniform H8192 family.

This is an explicit source locator and endpoint wording qualification. It introduces no new numerical estimate. The final mathematical subject is the original frozen proof read with this clarification. Its complete7571-check final replay is unchanged and does not need repetition for this textual precision; the script already retains non-strict exact model identities.

Source SHA256: \nolinkurl{a36916dd92cfcd4a13d9f7dd8906222dff73a80b5aee433280818becaba0adc0}. Frozen proof SHA256: \nolinkurl{8b05dd12db9c2633b889c0cdc845858e9dce6ce88ad571a67ee0ad466d6b6969}. Frozen executable SHA256: \nolinkurl{4f7b410625c9cae7e505ec517ab8ff272e5f2a6867a27dae3de6eb8f16230b32}.

\clearpage

\section{AM-LIFT-100 --- the literal half-level theta nonvanishing law}\label{am-lift-100-the-literal-half-level-theta-nonvanishing-law}

\textbf{Canonical checked integration --- 2026-11-23T10:40:33.872893+00:00.} Complete saved-proof and executable reviews and separate exact diagnostic replays passed. No numerical error charge is introduced. Frozen author history is preserved. AM-LIFT-100 HN1-19 proves the literal ThNe0 law given an actual algebraic FrPt over a finite number-field extensionK of the coefficient field, the original symmetric principal rigidified line and odd ideal half-level, literal IsNab with its prescribed conjugation, and literal ThValTot. For every unchanged nonzero P in harmA of total degreeD=DF with D+2\textless=M, some individual permitted n-tuple produces a nonzero actual thGen rule. The original ideal b satisfies b bar(b)=(M); C=A{[}b{]} has orderM\^{}g, H=A{[}b\^{}2{]} has orderM\^{}(2g), and the actual eta identifies La/bLa withC. Unique normalized odd lifts of H/C form exactlyM\^{}g permitted sections spanning the pullback of H0(A/C,L\_B\^{}M), with pi\^{}*L\_B=L\^{}M and L\_B principal. Bauer-Szemberg interpolation, the original common IsNab chart and ordinary highest symbols prove detection, retaining the literal M\^{}(-D) and factorials. The construction occurs on a reduced test-ring presentation ofK(A{[}M\^{}2{]}); no additional scalar-root field is required. At mostM\^{}(gn) tuples suffice existentially; no particular tuple or its norm/content is evaluated. Dom, n\textgreater0 and r\textless=n remain as in ThNe0; cusp-valuedness uses a separately supplied ThInCusp. This is a conditional named source-law constructor, not nonvanishing of the restricted original coupled selected image. Actual Moduli/FrPt existence, IsNab, ThValTot, other LawsP/LawsSh, full Host/Sizes, prescribed period, specific032, actual g\_j,e\_j,H and evaluatedDP44 remain open. No charge; current132/31 arrays and complete floors27/11 unchanged.

\textbf{Author:} /root/lift\_accounting, GPT-6 Astra Pro. \textbf{Date:} 23 November 2026 UTC.

\textbf{Status:} complete mathematical text, pending independent whole review. The separate author verification records the completed whole-file/source reading, exact diagnostics and frozen hashes. This proves a conditional constructor for the actual ThNe0 law. It does not identify the original coupled selected image with this entire theta family. No canonical file or arithmetic error list is changed.

\subsection{1. Statement and exact hypotheses}\label{statement-and-exact-hypotheses-1}

Retain an actual geometric Host E, Datum H, odd Level H, Coeff l and Moduli Dw l in the source. Put \[
 [E:\mathbb Q]=2f,\qquad r=Dw.p+Dw.q,\qquad g=fr,
 \qquad M=l.M,\qquad \mathfrak b=l.Mb.
 \tag{HN1}
\] Thus \(M\ge3\) is odd, \(\mathfrak b\overline{\mathfrak b}=M\mathcal O_E\), and \(\mathfrak b+\overline{\mathfrak b}=\mathcal O_E\). No generator of \(\mathfrak b\) is assumed or chosen.

Assume the following actual data, rather than their names alone:

\begin{enumerate}
\def\labelenumi{\arabic{enumi}.}
\tightlist
\item
  An algebraic frame point \(y\in\operatorname{FrPt}(X_o,c;K)\), where \(K/c.F\) is a finite number-field extension. Its abelian variety, full \(\mathcal O_E\)-action, original rigidified symmetric principal line \(L\), actual half-level \(\eta\), eigenframe and theta generator \(s\) are those of that point.
\item
  A source operator nab satisfying the literal IsNab Xo c nab, with the prescribed conjugation. It is not replaced by an independently normalized derivative.
\item
  The literal ThValTot Xo c nab, on all the source's rings and frame points.
\end{enumerate}

An actual supplied point from the checked mixed-principal-point/seed constructions may be used when all their hypotheses match. This theorem does not construct the ambient D7, all of LawsP, a missing Moduli.rep, or the remaining full Host/Sizes records.

\subsubsection{Theorem HN2}\label{theorem-hn2}

For every \(n>0\), weight \(a\), and nonzero \(P\in\operatorname{harmA}(c,a,n)\), put \[
 D=DF(Dw,a)=\sum_\tau\sum_{i<Dw.mult(\tau)}a(\tau,i).
\] If \(D+2\le M\), there is an actual permitted tuple \(x:\operatorname{Fin}(n)\to Dw.La\) and a source rule \[
 F\in\operatorname{thGen}(X_o,c,\mathrm{nab},a,n),\qquad F\ne0,
 \tag{HN2}
\] using that same polynomial \(P\). A nonzero value occurs on a test-ring presentation of a finite extension of the given frame point's field. More precisely, it is enough to extend \(K\) to \[
 K_1=K(A[M^2]).
 \tag{HN3}
\] No further field is required for the normalized symmetric theta lifts. There are exactly \(M^g\) distinct permitted sections per column and at most \(M^{gn}\) tuples to consider. This is a finite existence assertion, not an evaluated witness or a lower bound for its value.

In particular these hypotheses imply the source proposition ThNe0 Xo c nab. Its additional Dom a n and \(r\le n\) hypotheses may be retained verbatim; the argument does not need to weaken or change them. If the separately supplied ThInCusp applies, the resulting rule belongs to its prescribed cusp subspace. No cuspidality statement is inferred without that source law.

\subsection{2. The actual ideal halves and the actual level map}\label{the-actual-ideal-halves-and-the-actual-level-map}

The source's geometric TateIso identifies all torsion with the quotients of the original projective \(\mathcal O_E\)-lattice Dw.La, compatibly with the action and the alternating form. Its local modules are free of rank \(r\). Consequently for every nonzero integral ideal \(\mathfrak a\), \[
 |A[\mathfrak a]|=(N_{E/\mathbb Q}\mathfrak a)^r.
 \tag{HN4}
\] One checks this after localizing at the finitely many primes dividing \(\mathfrak a\); the index of a prime-power ideal on a free rank-\(r\) module is its ideal norm to the \(r\)-th power. This proves (HN4) for the actual projective lattice and introduces no class-number factor.

Conjugation preserves ideal norms. Taking norms in \(\mathfrak b\overline{\mathfrak b}=(M)\) gives \(N\mathfrak b=M^f\). Define \[
 C=A[\mathfrak b],\qquad H=A[\mathfrak b^2],\qquad D_0=A[M].
\] The coprime-ideal decomposition, or its local free-module calculation, gives \[
 |C|=M^g,\quad |H|=M^{2g},\quad |D_0|=M^{2g},\quad
 H\cap D_0=C,\quad [M]:H/C\overset\sim\longrightarrow C.
 \tag{HN5}
\]

This is also the literal half-level already present in the source, not an additional freely chosen torsion subgroup. In Obj.geom, the actual map \(\eta\) agrees with TateIso.ev M on \(\overline{\mathfrak b}Dw.La\), while eta\_ker kills \(\mathfrak b Dw.La\). Since these submodules sum to Dw.La, the Chinese remainder decomposition proves \[
 Dw.La/\mathfrak b Dw.La\overset{\eta}{\simeq}C.
 \tag{HN6}
\] Indeed, if \(u=v+w\) with \(v\in\overline{\mathfrak b}Dw.La\) and \(w\in\mathfrak b Dw.La\), then \(\eta(u)=\mathrm{ev}_M(v)\). Its kernel is exactly \(\mathfrak b Dw.La\), by the local free-module intersection \(\overline{\mathfrak b}Dw.La\cap\mathfrak b Dw.La=M Dw.La\). This also proves surjectivity onto the \(\mathfrak b\)-torsion factor. Thus every point in (HN5)'s final group is represented by a permitted label \(x\in Dw.La\).

\subsection{3. Odd symmetric lifts and principal descent}\label{odd-symmetric-lifts-and-principal-descent}

Use the original normalized symmetry and rigidification of \(L\). The cubical multiplication isomorphism is \[
 \alpha_M:[M]^*L\overset\sim\longrightarrow L^{M^2},
 \tag{HN7}
\] with its normalization at zero. In the literal convention this is exactly IsNormMul; its inverse is not silently substituted in IsTwistPull.

The generator required as \(\alpha_0\) in IsVtOver is on the universal group, not merely on the chosen fibre. Construct (HN7) there first. The given rigidified CubicalStr and normalized symmetry yield, by pulling the cubical identity to integer multiples and induction, \[
 [m]^*L\simeq
 L^{m(m+1)/2}\otimes([-1]^*L)^{m(m-1)/2}\simeq L^{m^2}.
\] The identity at \(m=0\) uses the original rigidification, and at \(m=1\) is the identity; the cube's face identities preserve that normalization at each induction step. This supplies the global generator \(\alpha_0\), whose pullback is the generator \(\alpha\) in IsVtOver. The normalized symmetry and this construction commute with base change. Subsequent uses of \(L,A\) denote their actual fibre.

For an odd exponent \(q\) and a \(q\)-torsion point of the theta group quotient, there is a unique lift \(U\) satisfying \(U^q=1\) and \(\iota(U)=U^{-1}\), where \(\iota\) is induced by the line's symmetry. Start with any lift. The product \(\iota(U)U\) is central. A square-root rescaling makes it one; then \(U^q\) is central and its own inverse. Changing the sign of \(U\), if necessary, makes \(U^q=1\), because \(q\) is odd. Two such lifts differ by a scalar satisfying both \(z^2=1\) and \(z^q=1\), hence by one. The theta-group involution is independent of rescaling the symmetry. On an isotropic subgroup these lifts form a subgroup: products commute and satisfy the same two identities, so uniqueness applies.

The actual \(\mathcal O_E\)-adjoint law on the cubical line gives the balanced commutator identity \[
 e_{L^q}(a x,y)=e_{L^q}(x,\bar a y).
 \tag{HN8}
\] It follows by applying the source SAB.adj isomorphism to the two translation restrictions defining the commutator. The pairing is perfect on \(A[q]\), since \(L\) is principal. We use the commutator convention induced by the source IsActOf operators. Replacing it by the inverse convention in a literature statement changes neither this identity nor its annihilator calculations; no source operator is replaced.

The subgroup \(C\) is isotropic for \(e_{L^M}\). Choose \(u\in\mathfrak b\), \(v\in\overline{\mathfrak b}\) with \(u+v=1\). For \(x,y\in C\), one has \(y=vy\), hence \(e_{L^M}(x,y)=e_{L^M}(\bar v x,y)=1\). The same argument with \(\mathfrak b^2\) proves that \(H\) is isotropic for \(e_{L^{M^2}}\). Their orders in (HN5) make them maximal isotropic at their respective levels.

The unique symmetric lifts of \(C\) in the theta group of \(L^M\) give descent data for the finite étale quotient \(\pi:A\to B=A/C\). Effective descent supplies an ample line \(L_B\) with the actual isomorphism \[
 \pi^*L_B\simeq L^M.
 \tag{HN9}
\] The Euler-characteristic degree formula gives \(\chi(L_B)=\chi(L^M)/\deg\pi=M^g/|C|=1\), so \(L_B\) is principal. Here \(\chi(L_B)\) is the line's Riemann--Roch degree, while the degree of its polarization homomorphism is \(\chi(L_B)^2=1\). It inherits symmetry and a rigidification. This quotient is used for interpolation; it is not identified with a separately prescribed moduli point and does not change the original \(L\). No principal generator of \(\mathfrak b\), self-isogeny of \(A\), or unchecked equality with a different chosen line is needed.

The theta-group facts used here are Mumford (1966), §1, Proposition 1 and its degree computation, p.~291, Theorem 1, pp.~293--294, and the invariant-section interpretation in Theorem 2's proof, p.~297. The specific ideal halves, normalized odd lifts and comparisons above are proved explicitly.

\subsection{4. The exact permitted span and its scalar convention}\label{the-exact-permitted-span-and-its-scalar-convention}

Let \(s_M=\alpha_M([M]^*s)\). For \(t'\in H\) let \[
 S_{t'}=U_{t'}s_M\in H^0(A,L^{M^2}),
 \tag{HN10}
\] where \(U_{t'}\) is the unique symmetric lift with exponent \(M^2\), acting by the literal IsActOf convention. The deck group \(D_0\) of \([M]\) has its canonical lift on \([M]^*L\), fixes \([M]^*s\), and is compatible with symmetry. Transport by (HN7) identifies it with the unique symmetric lift. Thus it fixes \(s_M\) exactly.

The descent lifts of \(C\) on \(L^M\), raised to the \(M\)-th tensor power, and the restriction of the \(D_0\) deck lifts are both symmetric finite-order lifts on \(L^{M^2}\); uniqueness identifies them. Since \(C\subset H\) and \(H\) is isotropic, every \(S_{t'}\) is \(C\)-invariant. Effective descent therefore places the orbit in \[
 W_M=\pi^*H^0(B,L_B^M)\subset H^0(A,L^{M^2}).
 \tag{HN11}
\]

The group \(D_0\) is maximal isotropic for \(L^{M^2}\), by its deck lift and order \(M^{2g}\). The \(D_0\)-character of \(S_{t'}\) is the actual operator commutator \[
 d\longmapsto U_dU_{t'}U_d^{-1}U_{t'}^{-1}.
 \tag{HN12}
\] These characters agree for two roots exactly when their difference lies in \(D_0^\perp\cap H=D_0\cap H=C\). Distinct characters of a finite abelian group in characteristic zero have linearly independent eigenspaces. Every orbit vector is nonzero because its operator is invertible. One representative per coset in \(H/C\) therefore gives \(M^g\) independent vectors in \(W_M\). Riemann--Roch gives \(h^0(B,L_B^M)=M^g\). Hence \[
 \{S_{t'}:t'\in H/C\}\text{ is a basis of }W_M.
 \tag{HN13}
\] Canonical lifts in \(H\) form a subgroup, and those in \(C\) fix \(s_M\), so changing a representative by an element of \(C\) does not change the section. This exact statement uses the normalized odd lifts; unnormalized lifts would instead give scalar multiples.

Now \(t=[M]t'\in C\), and (HN6) supplies \(x\in Dw.La\) with \(\eta(x)=t\). The root \(t'\) is killed by \(\mathfrak b^2\). Use \(\alpha_M\) and its required base pullback for the two generators in IsVtOver, \(s_1=s_M\), and \(U=U_{t'}\). Its root, annihilation, normalization, twist-pull and action clauses are exactly (HN7) and (HN10). Thus (HN13) comes from the \textbf{literal permitted family}, not the full \(M^2\)-torsion orbit or arbitrary Schwartz data.

All data in (HN10) are defined over \(K_1\) in (HN3). The torsion points are rational there; the normalized symmetric lift over each is geometrically unique. Galois invariance and descent define it over that same field. The square-root argument above proves geometric existence, whereas uniqueness proves descent; no extra root extension is asserted. The normalized multiplication isomorphism is already over \(K\). The auxiliary quotient adds no field extension.

\subsection{\texorpdfstring{5. The exact \(D+2\) jet threshold}{5. The exact D+2 jet threshold}}\label{the-exact-d2-jet-threshold}

Bauer and Szemberg (1997), Theorem 1 and Corollary 2.2, prove that a tensor product of \(D+2\) ample line bundles on a complex abelian variety is \(D\)-jet ample. Thus for \(M\ge D+2\), \[
 H^0(B,L_B^M)\twoheadrightarrow
 H^0(B,L_B^M\otimes\mathcal O_B/\mathfrak m_0^{D+1}).
 \tag{HN14}
\] If \(M>D+2\), use \(D+1\) factors \(L_B\) and the final factor \(L_B^{M-D-1}\). For \(D=0\) this is global generation for two ample factors. The number-field statement follows by extending to \(\mathbb C\) and faithfully descending the cokernel of this finite-dimensional map. Only characteristic zero is asserted; K022's stronger integral threshold is unchanged.

The étale quotient \(\pi\) identifies the completed local rings at zero. Since \(\pi^*L_B^M=L^{M^2}\) by the \(M\)-th power of (HN9), pulling back (HN14) yields \[
 W_M\twoheadrightarrow
 H^0(A,L^{M^2}\otimes\mathcal O_A/\mathfrak m_0^{D+1}).
 \tag{HN15}
\] The actual line comparison is an invertible local frame change and preserves jet surjectivity. When \(g=0\), the assertion follows directly from the section space of a point, with empty tangent indices.

This is the new ideal-quotient adapter. DP59 and K022 used a principal \(\alpha\)-isogeny on their stated Jacobian products and an integral interpolation threshold. No such special isogeny or reduction-at-11 hypothesis is imported here.

\subsection{6. Literal derivatives, leading symbols, and an individual tuple}\label{literal-derivatives-leading-symbols-and-an-individual-tuple}

RowDeg implies that every supported monomial of \(P\) has total degree \(D\). Thus P.totalDegree in thExpr is \(D\). Choose a supported monomial \(\beta\) with coefficient \(c_\beta\ne0\), and let \(\beta_b\) be its exponents in column \(b\), \(D_b=|\beta_b|\). Then \(\sum_bD_b=D\).

At the given frame over \(K_1\), choose algebraic formal parameters \(z_\ell\) dual to its actual eigenvectors. IsNab supplies a common formal group chart and a rigidified line frame. Its Dof, DmonOf and DpolyOf are invariant differential operators. Because the formal group law has identity linear term, their highest symbols at zero are the ordinary differential monomials in the specified tangent vectors. The coefficient-field clause Coeff.conj splits \(E\otimes_{\mathbb Q}K_1\) into its embedding factors, and the characteristic-zero \(E\)-action decomposes the tangent space accordingly. Thus FrPt.basis in the individual eigenspaces gives a basis of the full \(g\)-dimensional tangent space. IsChart.basis and IsChart.tan then make the tangent matrix invertible and force the chart dimension to be \(g\). Its compatible maps on all jet rings induce an isomorphism of completed local rings. Explicitly, the inverse formal map is constructed successively in each total degree by solving with that same invertible linear term. In particular a germ vanishing to order \(d\) is killed at zero by every operator of order less than \(d\).

The gauge in nabVal multiplies the section germ by \(\exp(M^2h_L)\). Its constant is one: IsCanonGauge.ax0/ax1 and the zero-preserving action give \(h_L(0)=0\). ChartAt.rig0 identifies the degree-zero line frame with the original rigidification. The leading-symbol calculation therefore keeps the actual connection and is unaffected by its higher coefficients.

Using (HN15), choose \(T_b\in W_M\) for each column, whose germ in an algebraic local frame equal to the rigidification at zero satisfies \[
 T_b(z)\equiv M^{D_b}z^{\beta_b}\pmod{\mathfrak m_0^{D+1}}.
 \tag{HN16}
\] For any differential monomial of \(P\), a column order below \(D_b\) gives zero. If no column order is below \(D_b\), total degree \(D\) forces all column orders to equal \(D_b\); the highest-symbol formula then kills the monomial unless its exponents are \(\beta\). Consequently the literal multilinear expression is \[
 M^{-D}\sum_{\gamma\in\operatorname{supp}P}c_\gamma
       \prod_b\mathrm{nab}_y(M^2,T_b,\gamma_b)
 =c_\beta\prod_{\ell,b}\beta_{\ell b}!\ne0.
 \tag{HN17}
\] This is the ordinary derivative factorial. It is neither divided out nor claimed to be a unit at a finite prime.

IsNab also supplies the needed linearity in each section. At a fixed frame its common chart assigns a germ in a fixed invertible line frame; sums and scalar multiples of sections have the corresponding germs. The nabVal formula is linear in that germ. Thus the left side of (HN17) is multilinear in the \(n\) sections.

Expand each \(T_b\) in (HN13). If every tuple of basis sections had zero value, multilinearity would make (HN17) zero. Some tuple \((S_{t'_b})_b\) therefore has nonzero value. By (HN5)--(HN6) it supplies actual labels \(x_b\in Dw.La\), and each section satisfies IsVtOver. The interpolating linear combination proves existence; the final assertion concerns one individual permitted tuple.

\subsection{7. From the algebraic value to the literal rule}\label{from-the-algebraic-value-to-the-literal-rule}

The source Rules are functions on reduced finitely generated test rings TRing c.F, not on arbitrary complex algebras. This distinction is retained.

Since \(K_1/c.F\) is finite separable, choose a primitive element with irreducible polynomial \(q\). The field \(R=c.F[T]/(q)\), with its radical defining ideal, is a literal test ring. Transport the actual point, eigenframe, generator and torsion data through \(R\simeq K_1\). The preceding argument can be performed over \(R\): jet surjectivity descends to it and IsNab is quantified over it.

For the selected tuple and the unchanged \(P\), ThValTot and existsUnique\_thRule supply a rule \(F\) satisfying IsThRule. At this test-ring frame use the identity faithfully flat map and the actual sections in ThValAt. Its thExpr is the nonzero value just found. Uniqueness in ThValTot forces \(F\)'s value to be that scalar, and therefore \[
 F\ne0\quad\text{as an actual source rule}.
 \tag{HN18}
\] This does not rely on choice-based evalAt at an arbitrary complex point, on an unproved FromTest assertion, or on transporting an unassumed naturality law for nab. It works in a test ring where all stipulated source laws apply.

The polynomial is still the supplied \(P\in\operatorname{harmA}\), so this proves HN2. Retaining Dom, \(n>0\), \(r\le n\), \(D+2\le M\), and the existence of \(P\) proves exactly ThNe0.

\subsection{8. Scope, costs, and the selected-image boundary}\label{scope-costs-and-the-selected-image-boundary}

The quantifiers are \[
 \text{given an actual algebraic frame }y,\quad
 \forall\,0\ne P\in\operatorname{harmA},\quad
 \exists\text{ an individual permitted tuple }x\text{ detected at }y
 \text{ after (HN3)}.
 \tag{HN19}
\] They do not say that every prescribed tuple, coupled class sum or iterated-image rule is nonzero.

\begin{itemize}
\tightlist
\item
  The ideal, level, original line, physical eigenframe, theta generator, coefficient field and polynomial remain fixed. No arbitrary-section family replaces thGen.
\item
  The field extension is (HN3); the finite family has size \(M^{gn}\). For a checked same-\(J,C_0\) seed, the existing torsion-field bounds apply to this inclusion under their original hypotheses. No new degree or discriminant constant is asserted.
\item
  There is no positive finite-prime content estimate, archimedean norm estimate or evaluated tuple. The value (HN17) belongs to an interpolating linear combination; it is not asserted to be the chosen basis tuple's value.
\item
  The quotient \(B\) only proves interpolation and does not change the original point, level or tuple realization.
\item
  IsNab, ThValTot and an actual algebraic point remain hypotheses. The other LawsP/LawsSh fields, complete Host/Sizes admission and specified arithmetic source realization are not inferred.
\end{itemize}

The original coupled family has a restricted shared-covector sum followed by a selected differential/branching map. This theorem ranges over the full literal thGen tuple family. It establishes no equality between those families. LIFT27's raw injectivity, K095's class sum, terminal brC/PtilFF, literal Read/032 and DP44 keep their earlier status. The result constructs one named source law; it does not replace the original selected-rule gate.

\subsection{9. Source and literature addresses}\label{source-and-literature-addresses}

\begin{enumerate}
\def\labelenumi{\arabic{enumi}.}
\tightlist
\item
  TwoWeightEqOneAddRoot.lean: geometric Host14--26, Datum30--49, Level53--62, SAB72--82, TateIso131--142 and Obj144--160. The projective lattice, conjugation, line and half-level are retained. PushInvertible.lean: Moduli29--38 and Coeff40--45. CubeSymmetricSwapFirst.lean: RigLB144--148 and CubicalStr167--186, including its actual face normalization.
\item
  IsCombAt.lean: FrPt67--77 and FrMap79--85.
\item
  IsThRule.lean: IsVtOver32--41, KilledBy63--64, thExpr66--70, ThValAt72--78, IsThRule80--82 and harmA/thGen86--91.
\item
  ThProps.lean: DF28, ThNe030--32 and existsUnique\_thRule44--52. HasDataCondition.lean: nabPin47, thTot51, thNe53. RecPJThAllNeBot.lean: consumer80--87.
\item
  ThGerm.lean: IsChart58--71, IsGerm172--180, ChartAt189--206 and IsNab208--217. CocycleRightZeroThd.lean: Dof26--28, DmonOf30--32, DpolyOf34--36, IsCanonGauge174--183 and nabVal187--189. RowDeg.lean: RowDeg40--41 and HHa/HaO43--48.
\item
  AlgebraMapCAt.lean: TRing9--12, Rules30, FromTest51--52 and evalAt54--57.
\item
  The two additionally recovered files are in recovery\_pending/lift/half\_level\_sources: ThetaPt.lean, SymLift44--51, IsTwistPull73--81 and IsActOf83--84; ThetaVV.lean, IsNormMul103--108. Both match the original DP59 member hashes. Retrieval and reading are not a build or validity audit.
\end{enumerate}

\textbf{Primary references}

Bauer, T., \& Szemberg, T. (1997). Higher order embeddings of abelian varieties. \emph{Mathematische Zeitschrift, 224}, 449--455. Theorem 1 and Corollary 2.2. Author PDF, pages 2 and 6: \url{https://www.mathematik.uni-marburg.de/~tbauer/1997-Higher-order-embeddings.pdf} . The theorem, proof and corollary, including the characteristic-zero convention, were read directly.

Mumford, D. (1966). On the equations defining abelian varieties. I. \emph{Inventiones Mathematicae, 1}, 287--354. §1(II), p.289; Proposition1 and degree computation, p.291; Theorem1, pp.293--294; Proposition3, pp.295--296; Theorem2, p.297. \url{https://pazuki.perso.math.cnrs.fr/index_fichiers/Mumford66.pdf} . The named statements and the relevant descent, commutator and invariant-section proofs were read directly.

\subsection{10. Verification boundary}\label{verification-boundary-1}

The companion diagnostics check finite odd-level Heisenberg and quotient models, ordinary highest symbols with nonlinear formal-group/gauge models, and source anchors and hashes. They do not formally prove the general geometric theorem or evaluate the original arithmetic image. The separate author packet records the final whole-file reading, exact diagnostic completion time, complete source manifest and frozen hashes; it preserves preliminary computations separately. No independent review is claimed by the author.

\subsection{Canonical algebraic-geometric context}\label{canonical-algebraic-geometric-context-70}

This is a characteristic-zero generic-fibre statement for the original polarized CM abelian scheme and its projective OE lattice; no extension over all of Spec(Z), integral content or new reduction-at11 theorem is asserted. The finite etale ideal subgroup C gives an actual quotient isogeny and descent of the cubical theta line, with exact Picard identity pi\^{}*L\_B=L\^{}M. The deck-character decomposition is in the finite-dimensional section representation; it does not replace the physical selected vector. Primary exact locators: Mumford(1966),section1(II)-(V) p289,Proposition1 and degree computation p291,Theorem1 pp293-294,Proposition3 pp295-296,Theorem2/invariant-section proof p297; Bauer and Szemberg(1997),Theorem1 PDFp2 and Corollary2.2 PDFp6,with full Theorem2.1 proof PDFpp3-5 and complex-ground-field convention. Actual IsVtOver32-41,ThNe030-32,DF28,IsNab208-217,IsNormMul103-107 and SymLift44-51 are pinned in the exact source record. Carried contextual anchors: SGA3 I4.7.3/4.7.3.1; SGA6 VI Theorem3.3; FGA232 Theorem3.1 pp148-149; FGA236 Corollaries1.7/2.3 pp229/232; SGA7 I VII2.7/2.9 pp169-174 and VIII7.4(b) pp310-311; cote4 scan003 sections8.1.1-8.1.2/scans005-006 sections8.4.1-8.4.3 and cote54 scans005-007. Formal derivatives retain the actual formal group and gauge, but no new Dieudonne slope or Cartier-dual divisibility follows. Section dimensions and exact line identities are compatible with the existing K-theory/exterior-power framework; no lambda-positivity, l-adic weight or new CM-motive L-function comparison is proved. The auxiliary quotient is isogenous to the original CM variety, without proving the required arithmetic oscillator or period matching. DP\_MULT Q3e/DP44/DP52/DP54-65 and THETA3 G1/A7 receive this named conditional source-law construction; CX\_GRH CX-ALG-024 and GCT\_PLETH T-052 retain their global gates.

\clearpage

\section{AM-K-118 --- full selected radial-band probability and its complete arithmetic tail}\label{am-k-118-full-selected-radial-band-probability-and-its-complete-arithmetic-tail}

\textbf{Canonical checked integration --- 2026-11-23T10:51:04.474372+00:00.} Complete saved-proof and executable reviews and separate exact diagnostic replays passed. No numerical error charge is introduced. Frozen author history is preserved. AM-K-118 BT1-22 proves that the entire actual all-top selected probability outsideI={[}1/35,1/31{]} is below2\^{}-40000000, uniformly in every original word and twist. Both high moments are in their true convergence domain. Fidelity restriction costs at most2\^{}9 times its square root over all168 embeddings, including radius0 and both compact factors. The full bounded original all-II globally rank-two-unipotent scalar-cross support27 family with nonzeroS cross and x\textless2\^{}3328 has complete nonarchimedean normalized mass below2\^{}11394877; original167 inputs, full affine central coset, supports, labels, G patterns and source/Haar factors are counted. The entire discarded arithmetic contribution is therefore below2\^{}-8605114. The retained-band main term remains to be estimated and this input appends no standalone charge or cell removal. Current132/31 arrays and complete floors27/11 remain unchanged; no unbounded-denominator, support28, signed nonvanishing, source admission, selected-image content, period, specific032, evaluatedDP44 or RH conclusion.

Author: \nolinkurl{/root/kummer_extension}, 23 November 2026 UTC. Controlled pending proof. Whole author, independent and root verification are distinct events. No canonical file or error array is changed.

\subsection{1. The actual measure and the scope of the conclusion}\label{the-actual-measure-and-the-scope-of-the-conclusion}

Use the actual J56 all-top upper endpoint, every admitted binary word and canonical twist, and the genuine vector-valued selected Hilbert metric. At one CM label, DO080 RT4--8 defines the compact-invariant probability density \nolinkurl{rho} on the ball against the invariant measure. Put \[
 Y=1-|z|^2,\qquad Z=Y^{-1},\qquad B=2^{955},\quad f=a=168,
 \qquad I=[1/35,1/31],\quad h=2^{30}.                 \tag{BT1}
\] The exact statement proved below is \[
             \Pr\{Y\notin I\}<2^{-40,000,000}.         \tag{BT2}
\] It is uniform in every actual selected word and allowed twist. The probability is the complete normalized selected probability; there is no conditioning on a smaller set, scalar-fibre replacement or uncharged lower tail.

For an arbitrary real group element, retain both compact factors and let \nolinkurl{F(g)} be DO080 RT7's nonnegative fidelity integral. Let \nolinkurl{F_I(g)} be the same integral restricted to the simultaneous conditions \nolinkurl{Y(z) in I} and \texttt{Y(g\^{}\{-1\}z)\ in\ I}. Then \[
 0\le F_I(g)\le F(g)\le1,\qquad
 |c_W(g)|\le F(g)\le F_I(g)+2\,2^{-20,000,000}.        \tag{BT3}
\] This includes Cartan radius zero. The product over all168 embeddings therefore obeys \[
 |c(h_0)|\le\prod_\tau F_{I,\tau}(g_{p,\tau}^{-1}h_0g_{p,\tau})
                      +2^{9-20,000,000}.              \tag{BT4}
\] The arithmetic matrix is denoted \nolinkurl{h_0} to distinguish it from the integer moment order \nolinkurl{h} in BT1.

For the entire bounded all-II globally rank-two-unipotent scalar-cross support27 sector of K113, with the original nonzero S-cross and \nolinkurl{x=N_E D_eta<2^3328}, replacing the full product fidelity by the restricted product adds total absolute normalized contribution less than \[
                         2^{-8,605,114}.              \tag{BT5}
\] The assertion covers every original coordinate size, both total-eligibility alternatives and every finite G pattern. It introduces no standalone error charge or removal of a coordinate cell. A subsequent retained-band estimate must retain BT5 in its complete union. The infinite denominator tail, nonunipotent sectors and support28 are outside BT5.

\subsection{2. Both signs of the exact high moment}\label{both-signs-of-the-exact-high-moment}

Write the original K046 BL2 parameters as \nolinkurl{r_i,m_i}. K046 BL8--10, in the actual Nakahama normalization, gives \[
 M_1=\mathbb E Z=\prod_i\frac{r_i+m_i-1}{r_i-2},\qquad
 \mathbb E Y^h=\prod_i\prod_{j=0}^{h-1}
                     \frac{r_i-1+j}{r_i+m_i+j},         \tag{BT6}
\] and, for integer \nolinkurl{h<min_i(r_i-1)}, \[
 \mathbb E Z^h=\prod_i\prod_{j=0}^{h-1}
                     \frac{r_i+m_i-1-j}{r_i-2-j}.       \tag{BT7}
\] The formula holds for every vector of the selected compact type, hence for its normalized average. Joint scalar norm formulas also hold for the actual tensor selected type. No finite-dimensional scalar surrogate is introduced.

DO080 RT21--23 proves uniformly in the actual word and twist that \[
 32<M_1<34,\qquad \sum_i\epsilon_i<B^{-1},\qquad
 \epsilon_i=\frac{m_i+1}{(r_i-3)(r_i+m_i-1)}.             \tag{BT8}
\] Here the literal endpoint has \nolinkurl{2^484<ell<2^485}, \[
 \begin{split}
 E&=B+\ell+96,\\
 r_i&=E+(4a-3-4i)\ell+q+1\quad(0\le i<a-1),\\
 r_{a-1}&=2\ell+q+1,\quad0\le q\le\ell,\\
 m_0&=2^{960}+83-668\ell,\quad
 m_i=4\ell-1-\ell\omega_i\quad(1\le i<a-1),\quad
 m_{a-1}=12.
 \end{split}                                           \tag{BT9}
\] In particular \nolinkurl{2h<min_i(r_i-2)} for BT1. The new negative moment is strictly inside the convergent chamber; its validity is not inferred from the much smaller second-moment case.

For the positive moment K107 SV5 gives, for every integer \nolinkurl{j>=0}, \[
 \frac{(r_i+m_i-1)(r_i-1+j)}{(r_i-2)(r_i+m_i+j)}
 \le1+(j+1)\epsilon_i.
\] Taking products and using \nolinkurl{log(1+t)<=t} gives, for \nolinkurl{h>=1}, \[
 M_1^h\mathbb E Y^h
 <\exp\!\left(\frac{h(h+1)}{2B}\right).                \tag{BT10}
\] At order zero the corresponding two sides are exactly1.

For the negative moment put \nolinkurl{A_i=r_i+m_i-1} and \nolinkurl{b_i=r_i-2}. Direct subtraction gives \[
 \frac{b_i(A_i-j)}{A_i(b_i-j)}
 =1+\frac{j(m_i+1)}{(r_i+m_i-1)(r_i-2-j)}
 =1+j\epsilon_i\frac{r_i-3}{r_i-2-j}.                   \tag{BT11}
\] For \nolinkurl{0<=j<h} and \nolinkurl{2h<r_i-2}, all denominators are positive and the last quotient is less than2. Multiplication, followed by the same logarithm inequality and BT8, proves \[
 M_1^{-h}\mathbb E Z^h
 <\exp\!\left(\frac{h(h-1)}{B}\right)\quad(h\ge2).    \tag{BT12}
\] At orders zero and one use equality instead of the strict sign. The application uses \nolinkurl{h=2^30}.

Markov's inequality now gives the two complete tails \[
 \Pr(Y>1/31)<(31/32)^h\exp\!\left(\frac{h(h+1)}{2B}\right),
\] \[
 \Pr(Y<1/35)<(34/35)^h\exp\!\left(\frac{h(h-1)}B\right). \tag{BT13}
\] Both exponential correction exponents are less than1. Since \nolinkurl{1-t<=exp(-t)}, their sum is less than \nolinkurl{2 exp(1-h/35)}. The exact inequalities \[
 e>\frac{65}{24}>\frac{27}{10},\qquad
 (27/10)^7>2^{10},\qquad
 \frac{17}{7}-\frac{2h}{49}
       =-\frac{2147483529}{49}<-40,000,000               \tag{BT14}
\] prove BT2. This uses a symbolic moment bound, not numerical summation of \nolinkurl{h} factors or an asymptotic concentration theorem.

\subsection{3. A restricted fidelity with the full probability tail retained}\label{a-restricted-fidelity-with-the-full-probability-tail-retained}

For the actual selected density, write \[
 F(g)=\int\sqrt{\rho(z)\rho(g^{-1}z)}\,d\mu(z).
\] The invariant measure and probability normalization give \nolinkurl{F(g)<=1} by Cauchy--Schwarz. On the subset where \nolinkurl{Y(z) notin I}, the same inequality bounds the integral by \nolinkurl{sqrt(epsilon)}, where \nolinkurl{epsilon=Pr(Y notin I)}. On the subset where \texttt{Y(g\^{}\{-1\}z)\ notin\ I}, invariance and the substitution \nolinkurl{z=gw} give the same bound. These two subsets contain the complement of the simultaneous band; overlapping points may be counted twice. Thus \[
                     F(g)-F_I(g)\le2\sqrt\epsilon.     \tag{BT15}
\] The right side is uniform in \nolinkurl{g}, including both compact factors and the genuine near point. No positivity of the actual signed coefficient is asserted.

For numbers \nolinkurl{0<=u_tau<=v_tau<=1}, telescoping the product gives \[
       0\le\prod v_\tau-\prod u_\tau
                                  \le\sum(v_\tau-u_\tau).
\] Apply this with \nolinkurl{u=F_I}, \nolinkurl{v=F}, and use \nolinkurl{2f=336<2^9}. The actual tensor selected normalized trace is bounded by the product fidelities, as in DO080 RT1--7 and the product application RT3. This proves BT3--4. The band is imposed on the original measure before an arithmetic or Gaussian enlargement; no independent freedom in the selected vector or terminal functional is introduced.

\subsection{4. A complete coarse arithmetic mass for the bounded support27 sector}\label{a-complete-coarse-arithmetic-mass-for-the-bounded-support27-sector}

Use exactly K113 SC27.1--10's original normal form, source/target pencils, phases, class and Haar normalization, finite projectors and reconstruction. Keep \nolinkurl{x<2^3328} and all27 active places of type II. K112 supplies active eligibility; K113 SC27.4 supplies at most one \nolinkurl{eta} per fixed27-place active support throughout the larger range \nolinkurl{x<2^12000}. Thus no norm-one generator, ideal-class or free phase multiplicity is omitted here. Its exact paired gap also directly implies the present smaller-window uniqueness.

Let \nolinkurl{I_Sigma} have exponent−1 at both CM components over each active T place, so \[
 N_E I_\Sigma=T^{-54},\qquad
 N_F I_{\Sigma,F}^{2}=T^{-54},\qquad T<2^{996}.         \tag{BT16}
\] All167 original covector values belong to \nolinkurl{J I_Sigma D_eta}, and hence to the larger \nolinkurl{I_Sigma}, because both J and the exact denominator ideal are integral. This enlargement is performed only for the tail's finite cardinality bound. It makes no claim that these larger tuples satisfy the original normal form.

In the original positive covector metric, DO098 CF4 and DO104 CD10 give \[
 \alpha=U\,rJ_Z^{-1}r^*\le R<U.
\] Consequently every original input y satisfies \nolinkurl{|r(y)|<=sqrt(R/U)||y||} at each embedding. Use the164 original unit S inputs and the actual saturated transverse inputs of K107 SV17--19, \[
 y_1=(-976627,-153032,38026,0),\quad
 y'_2=(-283839,347785,-86419,0),\quad
 z_0=(256766,-7654233287,-30805128044,1).
\] Each of their squared lengths is less than \nolinkurl{2^80}. Therefore all167 values have modulus less than \nolinkurl{2^40} at every embedding. This is a bound on the values of the original integral inputs, including \nolinkurl{y'_2} itself; no dyadic division of an old lattice value is used.

For a nonzero value \nolinkurl{v in I_Sigma}, the integral ideal \texttt{(v)I\_Sigma\^{}\{-1\}} has norm at least1. The product formula and AM--GM give the half-trace Euclidean minimum \[
 \|v\|\ge\sqrt f\,T^{-27/f}.                           \tag{BT17}
\] In real dimension \nolinkurl{2f=336}, disjoint balls of half this minimum around lattice points in the radius \nolinkurl{2^40 sqrt(f)} ball yield \[
 \#\{v\in I_\Sigma: |v_\tau|<2^{40}\ \forall\tau\}
 \le(1+2^{41}T^{27/f})^{336}
 <2^{42\cdot336}T^{54}<2^{67896}.                       \tag{BT18}
\] The zero value is included. Counting all167 values independently enlarges all coordinate-size and S-position subsets at once; no further \nolinkurl{binom(164,k)} or union over k is needed.

For fixed eta, row and original G pattern, the complete central fibre is in one affine coset with difference lattice \[
 iJ_0B_\eta I_{\Sigma,F}^{2}\subseteq iI_{\Sigma,F}^{2},
\] by DO104 CD14--19 and K113 SC27.7. Its imaginary scalar has absolute value less than \nolinkurl{H/(2U)<2^8} at every embedding. The affine real-part shift and all dyadic conditions remain; the enlarged affine imaginary lattice has minimum at least \texttt{sqrt(f)T\^{}\{-54/f\}}. The same packing argument, now in real dimension168, gives \[
 \#d\le(1+2^9T^{54/f})^{168}
       <2^{10\cdot168}T^{54}<2^{55464}.                 \tag{BT19}
\] This is an affine packing bound, not an unproved covolume-density estimate. Eta, the167 values and d reconstruct the full original matrix injectively. There is no discarded negative-column or source-isometry fibre.

Keep all original finite multiplicities. The number of active supports is less than \nolinkurl{2^104}; retain the harmless pivot factor \nolinkurl{164<2^8} and every one of the \nolinkurl{4^334=2^668} original common G patterns. For each such pattern, \[
 |a_G(h_0)|\le2Q^{-335\cdot167}(N_EJ)^{335/2}\le2,
\] since \nolinkurl{N_EJ<=Q^334}. At each active type-II place the complete diagonal-label sum is at most1: there are at most \nolinkurl{T-1} actual labels, each with absolute normalized weight \nolinkurl{1/(T-1)}. This weaker generic sum suffices; no new active-label uniqueness or inactive factor is claimed. All inactive coefficients and both mixed contractions retain their bound1. The class/Haar factors are exactly those in the accepted normalized count, as in DO104 section5 and K113 SC27.8--10; their retained remaining factor has absolute value at most1. No new source fibre or independent phase is inserted or suppressed.

Thus the entire nonarchimedean normalized absolute mass multiplying a uniform real-trace bound is less than \[
 2^{104+8+668+1+167\cdot67896+55464}
                         =2^{11394877}.               \tag{BT20}
\] This is deliberately very coarse. It covers both surviving eligibility alternatives and all coordinate sizes without borrowing any small sector charge.

\subsection{5. Complete tail composition and its limits}\label{complete-tail-composition-and-its-limits}

Multiply the uniform additive product-fidelity error BT4 by BT20. Every original arrow, support, pattern, finite label, row and central affine fibre has already been included. The result is \[
 2^{11394877}\,2^{9-20000000}
                         =2^{-8605114},                \tag{BT21}
\] which proves BT5. There is no summation over selected words: the bound is pointwise and uniform for each allowed word and twist, as is the original requested kernel estimate. Likewise the168 embedding loss has already been paid once in BT4.

More precisely, if \nolinkurl{w(h_0)} denotes the original normalized nonarchimedean weight in this bounded sector, then \[
 \sum_{h_0}|w(h_0)c(h_0)|
 \le\sum_{h_0}|w(h_0)|\prod_\tau F_{I,\tau}(g_{p,\tau}^{-1}h_0g_{p,\tau})
                                      +2^{-8605114}.     \tag{BT22}
\] The first term is still the actual restricted fidelity and must be estimated separately. Neither its smallness nor signed nonvanishing is asserted. This theorem does not append a charge by itself; a subsequent complete arithmetic theorem may pay BT21 in its union. No bound for an unbounded denominator family is inferred from this finite containing count.

\subsection{6. Source and verification boundary}\label{source-and-verification-boundary-1}

The literal all-word/twist probability law and its convergence chamber are K046 BL1--14 and DO080 RT4--8/RT20--24. K107 SV3--6 is used with its canonical boundary amendment at order0. The new negative-moment estimate is proved explicitly in BT11--12. The saturated original inputs and memberships are K107 SV17--19; original energy, injective reconstruction and source/fibre normalization are DO098 CF4--10 and DO104 CD10--19/section5. Support27 eligibility, bounded uniqueness and the precise finite containing set are K112 UE1--12 and K113 SC27.1--10. None of their numerical removal charges is reused.

Primary metric source, freshly read in the retained original: Nakahama, R. (2016). \emph{Norm computation and analytic continuation of vector valued holomorphic discrete series representations}. Journal of Lie Theory,26(4),927--990, equation(5.7) and Theorem5.1, printed960--961. The exact source dictionary and all shifts remain K043/K046's checked interface. No theorem-validity audit or Lean build of the Alpöge archive is performed.

The actual field, involution, Hermitian form, projective modules, polarization, theta line and finite projectors remain unchanged. Carried context is SGA3 I4.7.3/4.7.3.1; SGA6 VI Theorem3.3; FGA232 Theorem3.1 pp148--149; FGA236 Corollaries1.7/2.3 pp229/232; SGA7 I VII2.7/2.9 pp169--174 and VIII7.4(b) pp310--311; cote4 scan003 sections8.1.1--8.1.2/scans005--006 sections8.4.1--8.4.3 and cote54 scans005--007. These are inherited contextual anchors. There is no new source admission, selected arithmetic image, prescribed period, DP44, specific032 or RH conclusion.

\subsection{Canonical algebraic-geometric context}\label{canonical-algebraic-geometric-context-71}

This is an estimate at the archimedean realization of the same projective OE modules and polarized CM abelian data. The positive measure and simultaneous radial restriction do not alter the scheme over its original arithmetic base, the involution, Hermitian lattice, theta line, Picard scheme or Poincare biextension. Carried contextual anchors are SGA3 I4.7.3/4.7.3.1; SGA6 VI Theorem3.3; FGA232 Theorem3.1 pp148--149; FGA236 Corollaries1.7/2.3 pp229/232; SGA7 I VII2.7/2.9 pp169--174 and VIII7.4(b) pp310--311; cote4 scan003 sections8.1.1--8.1.2 and scans005--006 sections8.4.1--8.4.3; cote54 scans005--007. These locators are inherited context, not new primary readings or newly proved comparisons. No new formal-group/Dieudonne, lambda-ring/K-theory, l-adic weight, CM-motive or motivic L-function result follows from this real analytic inequality alone. DP\_MULT Q3e/DP44, THETA3 G1/A7, CX\_GRH CX-ALG-024 and GCT\_PLETH T-052 retain their current arithmetic and global gates until a separately checked composition changes them. No Lean build or archive theorem-validity audit is performed.

\clearpage

\section{AM-R-099 --- retained angular height on the full paid radial band}\label{am-r-099-retained-angular-height-on-the-full-paid-radial-band}

\textbf{Canonical checked integration --- 2026-11-23T10:55:51.699048+00:00.} Complete saved-proof and executable reviews and separate exact diagnostic replays passed. No numerical error charge is introduced. Frozen author history is preserved. AM-R-099 AH1-25 derives an extra exp{[}-(111/50)(cosh r-6/5){]} from the same actual selected angular kernel restricted toI={[}1/35,1/31{]} at both radial arguments. The old and weighted Jacobians give separate alternative bounds, retaining the exact detailed-balance cancellation and whole radial spread. The original unipotent Cartan identity yields joint J\_lambda(t,w)=t\^{}(-1/2)(1+t/4)\^{}(-167)(1+w\^{}2/{[}4(1+t/4)\^{}2{]})\^{}(-333/4)exp(-111t/100), with the original physical w=-U Im d.~The non-strict original norm\textless=133 cutoff gives a single full near-point2, and the entire intercept costs less than2\^{}108. Original central integration may consume exactly one radial power, leaving shape166; its arithmetic application remains separate. The weighted energy-only alternative retains its own(2/15)\^{}84 coefficient and is not multiplied into the joint bound. The exact shifted measure has infinite mass but finite every positive integer inverse moment; the zero-degree term needs original-row damping before enlargement. The actual trace is bounded only after adding the complete K118 tail2\^{}-8605114 in its precise bounded support27 scope. No standalone error charge or coordinate removal; current132/31 arrays and complete floors27/11 unchanged.

Author: \nolinkurl{/root}, GPT-6 Astra Pro; inherited effort is not externally exposed. 23 November 2026 UTC. \textbf{Controlled pending proof.} Whole author and independent reviews, exact diagnostics, and any arithmetic application are separate events. No canonical error is appended here.

\subsection{1. Actual selected measure and statement}\label{actual-selected-measure-and-statement}

Retain the original J56 selected vector-valued representation, its actual Hilbert metric, both compact factors, binary word and canonical twist, all fields, polarizations, labels, finite projectors, phases, Haar measures and class masses. Let rho be the compact-invariant probability density of DO080 RT4-8 against the invariant ball measure mu. Put \[
 a=f=168,\quad\nu=333/2,\quad
 \vartheta=\frac{4^{167}}{334\binom{334}{167}},\quad
 I=[1/35,1/31],\quad Y(z)=1-|z|^2.
 \tag{AH1}
\] For an actual real group element g, define the restricted nonnegative fidelity \[
 F_I(g)=\int_{Y(z)\in I,\,Y(g^{-1}z)\in I}
                \sqrt{\rho(z)\rho(g^{-1}z)}\,d\mu(z).
 \tag{AH2}
\] The untruncated fidelity F(g) is DO080 RT7's upper bound for the normalized selected trace. In particular 0\textless=F\_I(g)\textless=F(g)\textless=1. \textbf{The selected trace is not asserted to be bounded by F\_I alone.} K118 BT1-22 supplies its separately proved complete tail comparison, stated precisely in section7 below.

Write c=cosh r and s=sinh r for the Cartan radius r\textgreater0. Put \[
 \gamma=111/50,\qquad\lambda=111/100,\qquad
 E(c)=\exp[-\gamma(c-6/5)].
 \tag{AH3}
\] The two alternative bounds proved from the same restricted angular kernel are \[
 F_I(g)\le\frac{\vartheta2^{-933/2}}s
                \left(\frac2{c+1}\right)^\nu E(c),
 \tag{AH4}
\] \[
 F_I(g)\le\frac{\vartheta2^{-933/2}}{\sqrt{30}}
              \frac cs\left(\frac2{c+1}\right)^\nu E(c).
 \tag{AH5}
\] Each can also be bounded by1. At r=0 only F\_I\textless=1 is needed; the quotient formulas are not evaluated. These are alternative kernel bounds, and they are never multiplied together. AH4 permits the original joint central-energy treatment without consuming an additional radial power. AH5 permits the weighted energy-only profile.

\subsection{2. An exact angular-height loss}\label{an-exact-angular-height-loss}

Use the actual radial transition K\_r(y,y') of DO080 RT9-14. For y,y' inI, write \[
 D=\sqrt{y/y'},\qquad R_y=\sqrt{1-y},\qquad
 D,D^{-1}\le\sqrt{35/31}<11/10.
 \tag{AH6}
\] The last strict inequality follows from3500\textless3751. In the first complex coordinate z\_1=u+it, \[
 y'=\frac{y}{(c+su)^2+s^2t^2},\quad
 u=\frac{\sqrt{D^2-s^2t^2}-c}s,
 \quad u^2+t^2<1-y.
 \tag{AH7}
\] The positive square root is forced by c-sR\_y\textgreater0, as in the exact source Jacobian proof. If c\textgreater=6/5 then c\textgreater D and c\textgreater D\^{}\{-1\}. Set \[
 K=1-\frac{(c-D)^2}{(c^2-1)(1-y)}.
\] The first-case angular argument of K108 AE5 bounds the integral of the density power by2 vartheta R\_y K\^{}nu whenK\textgreater0. IfK\textless=0 the directional density is zero almost everywhere. No restriction on the integration variable t beyond its actual feasible interval is imposed.

A direct common-denominator calculation gives the exact identity \[
 1-\frac{K}{2D/(c+1)}=
 \frac{(D-1)^2+y(c-1)(c+1-2D)}
      {2D(c-1)(1-y)}.
 \tag{AH8}
\] All denominators are positive. Since y/(1-y)\textgreater=1/34, D\textless=11/10 and c+1-2D\textgreater=c-6/5\textgreater=0, it follows that \[
 1-\frac{K}{2D/(c+1)}
 \ge\frac5{374}(c-6/5)\ge\frac{c-6/5}{75}.
 \tag{AH9}
\] The second comparison is375\textgreater=374. If a feasible pair hasK\textgreater0, AH9 forces x=(c-6/5)/75\textless1. The elementary inequality1-x\textless=exp(-x) then proves \[
 K^\nu\le\left(\frac{2D}{c+1}\right)^\nu
             \exp[-\nu(c-6/5)/75]
 =\left(\frac{2D}{c+1}\right)^\nu E(c).
 \tag{AH10}
\] WhenK\textless=0 the zero extension satisfies the same upper bound. This also covers arbitrarily largec: an impossible positive-height case contributes zero, not a negative density.

If1\textless c\textless6/5, the established all-feasible K108 AE5-7 estimate applies in its two original cases c\textgreater D and c\textless=D. Its old height factor is{[}2D/(c+1){]}\^{}nu. SinceE(c)\textgreater1 in this range, the same AH10 upper bound follows without claiming c\textgreater D. Thus the directional bound is uniform for everyr\textgreater0 onI timesI. The reverse direction has the same proof withD replaced byD\^{}\{-1\}.

\subsection{3. The two real Jacobians and detailed balance}\label{the-two-real-jacobians-and-detailed-balance}

The exact old Cartesian bound is \[
 \left|\frac{du}{dy'}\right|
          \le\frac{\sqrt y}{2s y'^2}.
\] Together with the normalized sphere marginal and AH10 this gives \[
 K_r(y,y')\le\vartheta\frac{a-1}{\pi s}
       \frac{\sqrt y}{\sqrt{1-y}\,y'^2}
       \left(\frac{2\sqrt{y/y'}}{c+1}\right)^\nu E(c).
 \tag{AH11}
\] Alternatively R098 WJ5-7 gives the all-feasible exact bound \[
 \left|\frac{du}{dy'}\right|
          \le\frac{c\sqrt y}{2s y'^{3/2}},
\] so the same height argument gives \[
 K_r(y,y')\le\vartheta\frac{(a-1)c}{\pi s}
       \frac{\sqrt y}{\sqrt{1-y}\,y'^{3/2}}
       \left(\frac{2\sqrt{y/y'}}{c+1}\right)^\nu E(c).
 \tag{AH12}
\] Both formulas hold only on the indicated band; outside it this proof makes no improved assertion.

Disintegration and detailed balance on the symmetric setI timesI give exactly \[
 F_I(g)=\iint_{I\times I}\sqrt{p(y)p(y')}
             \sqrt{K_r(y,y')K_r(y',y)}\,dy\,dy'.
 \tag{AH13}
\] Indeed DO080 RT13 cancels the invariant radial weights before restriction, and restricting both radial variables preserves that identity. Compact factors preserveY andrho; they remain present in the actual selected coefficient. Applying AH11 in both directions cancelsD andD\^{}\{-1\}, while the geometric mean retains one entire factorE(c). The remaining radial integral is the square of \[
 \int_I\sqrt{p(y)}\,y^{-3/4}(1-y)^{-1/4}\,dy.
\] It is bounded by its full positive integral. K107 SV8-12's full-tail spread estimate, in the exact coordinateS=2sqrt(Y\^{}\{-1\}-1), bounds this square by2pi sigma, with334 sigma\textless2\^{}\{-933/2\}. Multiplication by(a-1)/pi gives AH4.

Applying AH12 in both directions instead gives the square of \[
 \int_I\sqrt{p(y)}\,y^{-1/2}(1-y)^{-1/4}\,dy.
\] The full positive integral has square less than2pi sigma/sqrt30 by R098 WJ11-14. This proves AH5. Those full radial integrals are used only to bound the retained-band integral; the trace outside the band still requires section7's explicit additive payment. No probability mass is silently conditioned away.

\subsection{4. A joint original central-energy profile}\label{a-joint-original-central-energy-profile}

For an actual nonidentity globally rank-two-unipotent original matrix, keep \[
 \alpha=(N^2)_{nn}>0,\quad0<\alpha_\tau\le R=17424/133<132,
 \quad w=\Im N_{nn}=-U\Im d,\quad H=133.
 \tag{AH14}
\] DO122 CE3-6 and its mandatory cutoff addendum apply at their support-independent normal-form scope. At the actual near point, \[
 c_p^2=(1+\alpha_p/2)^2+w_p^2,\quad
 \alpha_p\ge e^{-2r_p}\alpha,\quad
 |w_p-w|\le6Hr_p,\quad r_p<2^{-352}.
 \tag{AH15}
\] Here the original representative has norm(h\_tau)\textless=H. This is the exact non-strict R048-RC cutoff; no strict version or bound on the full compact part ofg\_p is assumed.

Putb(t)=1+t/4, p\_0=333/4, and \[
 J_\lambda(t,v)=t^{-1/2}b(t)^{-167}
       \left(1+\frac{v^2}{4b(t)^2}\right)^{-p_0}e^{-\lambda t}.
 \tag{AH16}
\] The exact inequalitiesc\^{}2-1\textgreater=t b(t) and(c+1)\^{}2\textgreater=4b(t)\^{}2+v\^{}2 are the original CE6. Also c\textgreater=1+t/2, so \[
 E(c)\le\exp(111/250)e^{-\lambda t}.
\] Applying these inequalities to the single AH4 kernel, rather than multiplying separate trace estimates, proves \[
 F_I(h_p)\le\vartheta2^{-933/2}\exp(111/250)
                            J_\lambda(\alpha_p,w_p).
 \tag{AH17}
\] Foru=v\^{}2/(4b(t)\^{}2+v\^{}2) in{[}0,1{]}, direct differentiation gives \[
 -\partial_{\log t}\log J_\lambda
 =\tfrac12+(167-2p_0u)\frac{t}{t+4}+\lambda t,
 \quad |\partial_v\log J_\lambda|\le p_0/(2b(t))\le p_0/2.
 \tag{AH18}
\] The first expression is positive and at most335/2+lambda t. Monotonicity in t permits replacement ofalpha\_p byexp(-2r\_p)alpha even ifalpha\_p exceedsalpha. Integrating the logarithmic derivative only over the interval{[}exp(-2r\_p)alpha,alpha{]}, where t\textless=R, and then transportingw gives \[
 J_\lambda(\alpha_p,w_p)
 \le\exp[(335+2\lambda R+3Hp_0)r_p]J_\lambda(\alpha,w).
\] The exact coefficient335+2lambdaR+3Hp\_0 is less than2\^{}16. Hence the entire product transport costs less than2. It replaces the single old global near-point factor. Sincevartheta\^{}f\textless2\^{}-649, \[
 \prod_\tau F_{I,\tau}(h_{p,\tau})
 <2\,2^{-649}\,2^{-933f/2}\exp(111f/250)
                         \prod_\tau J_\lambda(\alpha_\tau,w_\tau).
 \tag{AH19}
\] This is a bound for the restricted product; the actual trace still carries its explicit tail. The exponential factor is independent of the central variable. Thus the original full affine central-coset argument may integrate AH16 in that variable, consuming exactly one power ofb and leaving the shape166 profile. Such an arithmetic application must still prove its complete row count and error union; none is asserted here.

\subsection{5. Weighted energy-only alternative and the exact cost}\label{weighted-energy-only-alternative-and-the-exact-cost}

The weighted Jacobian gives a separate consequence. The function \[
 G(c)=\frac c{\sqrt{c^2-1}}\left(\frac2{c+1}\right)^\nu
\] is decreasing for c\textgreater1, since its logarithmic derivative is -1/{[}c(c\^{}2-1){]}-nu/(c+1). BothG(c) andE(c) decrease. Therefore AH5, c\_p\textgreater=1+alpha\_p/2 and1+t/2\textless=2(1+t/4) give \[
 F_I(h_p)\le\frac2{\sqrt{30}}\vartheta2^{-933/2}
       \exp(111/250)\Phi_{166,\lambda}(\alpha_p),
 \quad\Phi_{m,\lambda}(t)=t^{-1/2}(1+t/4)^{-m}e^{-\lambda t}.
 \tag{AH20}
\] For m166, its logarithmic elasticity is at most333/2+lambda t. The same original-energy argument bounds the whole transport byexp{[}(333+2lambdaR)sum r\_p{]}\textless2. Hence the energy-only product has the additional factor(2/15)\^{}84 relative to the coefficient inAH19, withproductPhi166 replacingproductJ. This is an alternative to the joint central bound, not an additional multiplier for it.

The entire exponential cost satisfies \[
                 \exp(111f/250)<2^{108}.
 \tag{AH21}
\] An exact rational proof is supplied in the companion certificate: forx=111/250, the positive Taylor sum through degree12 plus{[}x\^{}13/13!{]}/{[}1-x/14{]} boundsexp(x) above, because subsequent term ratios are at mostx/14. That rational bound is less than39/25, and39\^{}168\textless2\^{}108 25\^{}168. No floating logarithm is used. Keeping the exact exponential or this explicit power-of-two majorant are both valid.

\subsection{6. Positive shifted measure and the infinite constant term}\label{positive-shifted-measure-and-the-infinite-constant-term}

For any positive integer m, letb have Gamma probability distribution of shapem and rate4, and letds/sqrt(pi s) be the positive half-line measure. Push their product forward byc=lambda+b+s. The resulting positive locally finite measuremu\_m,lambda has infinite total mass and represents \[
 \Phi_{m,\lambda}(t)=\int e^{-ct}\,d\mu_{m,\lambda}(c),\quad t>0.
 \tag{AH22}
\] All integer positive inverse moments are finite. Tonelli followed by the ordinary positive Beta integral in s gives \[
 M_j=\int c^{-j}\,d\mu_{m,\lambda}(c)
 =\frac{\Gamma(j-1/2)}{\Gamma(j)}
       {\bf E}(\lambda+b)^{1/2-j}
 \le\frac{\Gamma(j-1/2)}{\Gamma(j)}\lambda^{1/2-j},
 \quad j\ge1.
 \tag{AH23}
\] Every argument is positive. The coarser final bound establishes finiteness; an arithmetic application may need a sharper independently certified moment. No analytic continuation, unproved quadrature, finite total mass or zero-th inverse moment is used. NIST DLMF5.9.1 and5.12.1/3 provide the ordinary integral normalizations; the translated convolution and bounds here are proved directly.

The original K116/DO121 zero-degree repair continues to apply. On an actual row withk\textgreater=1 original nonzeroS inputs, the exact norm floor isN\_F(alpha)\textgreater=(Rk/q)\^{}f.~For positivec\_tau andtheta in(0,1), AM-GM first gives \[
 e^{-\theta\sum c_\tau\alpha_\tau}
 \le e^{-\theta f(\prod c_\tau)^{1/f}Rk/q}
 \le(\theta Rk/q)^{-f}\prod c_\tau^{-1}.
 \tag{AH24}
\] Integrating the final product usesM\_1\^{}f.~This reservation is made on actual rows before Gaussian enlargement. It does not assign a finite mass to the constant polynomial term. All other expanded polynomial degrees must still be retained with their appropriate moments and finite-pattern costs.

\subsection{7. Full arithmetic tail and conclusion boundary}\label{full-arithmetic-tail-and-conclusion-boundary}

Apply the independently proved K118 BT3-5/BT20-22 to precisely its bounded original all-II globally rank-two-unipotent scalar-cross support27 sector with nonzeroS cross and x\textless2\^{}3328. Its full selected probability outsideI is less than2\^{}-40000000; the complete original arithmetic tail, after the168-factor product loss, is less than2\^{}-8605114. Consequently for any actual subsetS of that sector and its original normalized finite weightsw(h), \[
 \sum_{h\in S}|w(h)c(h)|
 \le\sum_{h\in S}|w(h)|\prod_\tau F_{I,\tau}(h_{p,\tau})
                              +2^{-8605114}.
 \tag{AH25}
\] AH19 orAH20 can bound the retained product on the right. The tail must be paid exactly as an explicit term in a later complete union; it is not appended as a standalone charge by this analytic theorem. No bound for an unbounded denominator family or different support follows from this containing count. No actual arrow existence, sign, full support floor, selected arithmetic-image content, source admission, prescribed period, evaluatedDP44, specific032 pairing or Hecke RH conclusion is asserted.

\subsection{8. Exact source and geometric context}\label{exact-source-and-geometric-context}

Actual selected density, invariant measure, finite-trace inequality, both compact factors, disintegration and detailed balance: DO080 RT4-19. Whole-tail spread: K107 SV3-12 with its canonical zero-radius/order-zero qualifications. Full conditional angular density and two-case height: K108 AE3-9. Weighted feasible Jacobian and radial Cauchy: R098 WJ5-14. Original joint Cartan identity and physical central variable: DO122 CE3-10, with its mandatory non-strict-cutoff addendum and R048-RC RC48.4/RC48.6, paragraph afterRC48.7. Full probability and original arithmetic tail: K118 BT1-22. Each is used at its stated exact scope; no old numerical interval charge is imported.

The underlying primary selected metric is Nakahama (2016), equation5.7 and Theorem5.1, printed960-961, through the already checked K043/K046/DO080 dictionary. Positive integral identities are NIST DLMF5.9.1,5.12.1/3 and5.4.6. This new proof is an estimate in the same real realization and not a new arithmetic source or representation. The projective OE modules, involution, polarization, cubical line and Picard/Poincare data remain unchanged. Inherited contextual anchors are SGA3 I4.7.3/4.7.3.1; SGA6 VI Theorem3.3; FGA232 Theorem3.1 pp148-149; FGA236 Corollaries1.7/2.3 pp229/232; SGA7 I VII2.7/2.9 pp169-174 and VIII7.4(b) pp310-311; cote4 scan003 sections8.1.1-8.1.2, scans005-006 sections8.4.1-8.4.3, and cote54 scans005-007. These are carried anchors, not new formal-group slopes, Cartier-dual divisibility, lambda-ring positivity, K-theoretic or l-adic weight comparisons, or CM-motive L-function identities. DP\_MULT Q3e/DP44 and THETA3 G1/A7 retain the original arithmetic/period gates; CX\_GRH CX-ALG-024 and GCT\_PLETH T-052 retain their global gates.

\subsection{Canonical algebraic-geometric context}\label{canonical-algebraic-geometric-context-72}

This is an estimate at the archimedean realization of the same projective OE modules and polarized CM abelian data. The positive shifted measure and simultaneous radial restriction do not alter the scheme over its original arithmetic base, the involution, Hermitian lattice, theta line, Picard scheme or Poincare biextension. Carried contextual anchors are SGA3 I4.7.3/4.7.3.1; SGA6 VI Theorem3.3; FGA232 Theorem3.1 pp148--149; FGA236 Corollaries1.7/2.3 pp229/232; SGA7 I VII2.7/2.9 pp169--174 and VIII7.4(b) pp310--311; cote4 scan003 sections8.1.1--8.1.2 and scans005--006 sections8.4.1--8.4.3; cote54 scans005--007. These locators are inherited context, not new primary readings or newly proved comparisons. No new formal-group/Dieudonne, lambda-ring/K-theory, l-adic weight, CM-motive or motivic L-function result follows from this real analytic inequality alone. DP\_MULT Q3e/DP44, THETA3 G1/A7, CX\_GRH CX-ALG-024 and GCT\_PLETH T-052 retain their current arithmetic and global gates until a separately checked composition changes them. No Lean build or archive theorem-validity audit is performed.

\clearpage

\section{AM-K-119 --- a positive rational majorant for the shifted Gamma166 profile}\label{am-k-119-a-positive-rational-majorant-for-the-shifted-gamma166-profile}

\textbf{Canonical checked integration --- 2026-11-23T10:55:52.116162+00:00.} Complete saved-proof and executable reviews and separate exact diagnostic replays passed. No numerical error charge is introduced. Frozen author history is preserved. AM-K-119 RM1-19 proves the global positive majorant t\^{}(-1/2)(1+t/4)\^{}(-166)exp(-111t/100)\textless=t\^{}(-1/2)(1+t/4)\^{}(-166)(1+t/400)\^{}(-444). Its Gamma166/rate4 plus Gamma444/rate400 plus half-Gamma convolution has infinite mass and inverse moments exactly for1/2\textless j\textless610+1/2. Every required integer degree1..167 is evaluated as pi times an explicit positive rational, with all full coefficients and strict product168 bounds recorded. Finite principal parts plus integration-by-parts recurrence agree exactly with the independent73704-cell rectangular-integral/binomial calculation at every degree. The complete scalar/transverse row polynomial retains all terms and repairs its sole zero degree by original-row energy reservation before enlargement. Actual selected-kernel, intercept, near-point, centralU scaling, all finite/source/class/Haar/support/pattern and complete band-tail costs remain separate. This analytic input appends no charge and removes no sector; current132/31 arrays and complete floors27/11 unchanged.

Author: /root/kummer\_extension. 23 November 2026 UTC. Controlled pending theorem. Whole author, independent and root reviews are separate events. No canonical file, error array or numerical sector charge is changed.

\subsection{1. The exact shifted profile and a positive upper measure}\label{the-exact-shifted-profile-and-a-positive-upper-measure}

Fix \[
 p=166,\quad \lambda=111/100,\quad N=444,\quad
 a_0=4,\quad b_0=400,\quad N/b_0=\lambda.                 \tag{RM1}
\] For \(t>0\), the shifted profile of interest is \[
 F_\lambda(t)=t^{-1/2}(1+t/4)^{-166}e^{-\lambda t}.
\] The elementary inequality \(\log(1+u)\le u\) gives the global positive majorant \[
 F_\lambda(t)\le G(t):=
 t^{-1/2}(1+t/4)^{-166}(1+t/400)^{-444}.                  \tag{RM2}
\] No finite energy cutoff or extra multiplicative constant is required for RM2. It is an upper inequality, not an equality with the shifted profile.

Let \(\gamma_{m,b}\) denote the Gamma probability measure of positive integer shape \(m\) and rate \(b\), with density \[
       \frac{b^m}{(m-1)!}u^{m-1}e^{-bu}\,du.
\] Let \(\xi(du)=u^{-1/2}du/\sqrt{\pi}\), and form the positive convolution \[
       \nu=\gamma_{166,4}*\gamma_{444,400}*\xi.           \tag{RM3}
\] The positive Gamma integral and Tonelli's theorem prove \[
                       G(t)=\int_0^\infty e^{-ct}\nu(dc).
                                                               \tag{RM4}
\] The analogous exact shifted measure is the translate by \(\lambda\) of \(\gamma_{166,4}*\xi\). Both measures have infinite total mass. Indeed their Laplace transforms tend to infinity like \(t^{-1/2}\) as \(t\downarrow0\), and monotone convergence applies. They are not probability measures, and no zeroth inverse moment is assigned a finite value.

\textbf{Theorem.} For every integer \(1\le j\le167\), the exact shifted inverse moment is bounded above by \[
 L_j^\lambda
 :=\frac1{\Gamma(j)}\int_0^\infty
 t^{j-3/2}(1+t/4)^{-166}e^{-\lambda t}\,dt
 \le L_j:=\int c^{-j}\nu(dc)=\pi r_j,                    \tag{RM5}
\] where every \(r_j\) is the explicitly defined positive rational in RM13--16 below. The companion exact table gives every rational \(r_j\) and every strict dyadic upper ceiling for \[
                       P_j=\bigl((22/7)r_j\bigr)^{168}. \tag{RM6}
\] In particular \(L_j^{168}<P_j\). The true shifted moments converge for every real \(j>1/2\); the rational-majorant moments converge precisely for \[
                    1/2<j<610+1/2.                     \tag{RM7}
\] Thus all integer moments through610 exist for the majorant, including the167 moments required by the original row polynomial. Nothing is asserted for its zeroth moment.

The elementary strict bracket \(\pi<22/7\) can be certified without decimal arithmetic: polynomial division gives \(\int_0^1 x^4(1-x)^4/(1+x^2)\,dx=22/7-\pi>0\).

This is a reusable analytic and exact-arithmetic theorem. Applying it to an actual selected coefficient requires a separately checked bound by \(F_\lambda\), with its original coefficient, near-point, band-tail, central-fibre and finite factors retained. It does not multiply distinct coefficient estimates or remove an arithmetic sector by itself.

\subsection{2. Moment normalization and convergence}\label{moment-normalization-and-convergence}

For \(j>0\) and \(c>0\), the ordinary positive Gamma integral gives \[
 c^{-j}=\frac1{\Gamma(j)}\int_0^\infty t^{j-1}e^{-ct}\,dt.
\] Applying Tonelli to RM4 therefore gives \[
 L_j=\frac1{\Gamma(j)}\int_0^\infty
       t^{j-3/2}(1+t/a_0)^{-p}(1+t/b_0)^{-N}\,dt.        \tag{RM8}
\] Near zero this integrand is comparable to \(t^{j-3/2}\). Near infinity it is comparable to \(t^{j-3/2-p-N}\). These two comparisons prove RM7 for \(j>0\), including divergence at its boundary exponents. Both measures are locally finite and have infinite mass, so their mass on \([1,\infty)\) is infinite; for \(j\le0\), \(c^{-j}\ge1\) there proves divergence as well. For the exact shifted integral the extra exponential supplies integrability at infinity for every \(j\), so only the lower condition \(j>1/2\) remains. RM2 inside these nonnegative integrals proves RM5's inequality without an exchange of a conditionally convergent integral.

For positive integer \(j\) in the convergent range, put \(t=s^2\). Then \[
 L_j=\frac{2a_0^p b_0^N}{(j-1)!}
       \int_0^\infty
       \frac{s^{2j-2}\,ds}{(s^2+a_0)^p(s^2+b_0)^N}.     \tag{RM9}
\] The factor2 is the full substitution Jacobian. The half-Gamma normalization in RM3 has already been accounted for through its Laplace transform \(t^{-1/2}\); no extra \(1/\sqrt\pi\) or \(\Gamma(j-1/2)\) is inserted into RM9.

\subsection{3. Two exact base integrals by finite partial fractions}\label{two-exact-base-integrals-by-finite-partial-fractions}

For positive integers \(u,v\), let \(\Delta=b_0-a_0=396\). The following rational-function identity holds: \[
 \frac1{(x+a_0)^u(x+b_0)^v}
 =\sum_{\ell=1}^{u}\frac{A_\ell(u,v)}{(x+a_0)^\ell}
  +\sum_{\ell=1}^{v}\frac{B_\ell(u,v)}{(x+b_0)^\ell},    \tag{RM10}
\] where \[
 A_\ell(u,v)=
 \frac{(-1)^{u-\ell}\binom{u+v-\ell-1}{u-\ell}}
 {\Delta^{u+v-\ell}},\qquad
 B_\ell(u,v)=
 \frac{(-1)^u\binom{u+v-\ell-1}{v-\ell}}
 {\Delta^{u+v-\ell}}.                                  \tag{RM11}
\] To prove the identity, expand the factor regular at \(x=-a_0\) in powers of \(x+a_0\), and retain the terms of orders0 through \(u-1\). This gives precisely the displayed \(A_\ell\). Expanding at \(x=-b_0\) gives \(B_\ell\). The difference in RM10 has no finite pole and tends to zero at infinity; as a rational function it is identically zero. This proves the whole identity, not only a sample of its values.

For a positive integer \(r\) and integer \(\ell\ge1\), \[
 J_\ell(r):=\frac1\pi\int_0^\infty(s^2+r^2)^{-\ell}\,ds
 =\frac{\binom{2\ell-2}{\ell-1}}
        {2^{\,2\ell-1}r^{\,2\ell-1}}.                  \tag{RM12}
\] The positive beta integral proves this, or one may start from \(J_1(r)=1/(2r)\) and integrate by parts to obtain \(J_{\ell+1}(r)=(2\ell-1)J_\ell(r)/(2\ell r^2)\). All these integrals converge; each finite partial-fraction term is individually integrable.

Define the positive rational \[
 \mathcal I(u,v)=
       \sum_{\ell=1}^{u}A_\ell(u,v)J_\ell(2)
      +\sum_{\ell=1}^{v}B_\ell(u,v)J_\ell(20).
\] Although this finite expression has signs, RM10--12 identify it exactly as \[
 \mathcal I(u,v)=\frac1\pi\int_0^\infty
         (s^2+a_0)^{-u}(s^2+b_0)^{-v}\,ds>0.             \tag{RM13}
\] All arithmetic in the signed sums is rational and exact; a cancellation of decimal approximations is never used.

Let \(C=2a_0^p b_0^N\). The first two rational coefficients are \[
 r_1=C\mathcal I(p,N),\qquad
 r_2=C\{\mathcal I(p-1,N)-a_0\mathcal I(p,N)\}.           \tag{RM14}
\] The second identity is the numerator identity \(s^2=(s^2+a_0)-a_0\) inside RM9, and hence its right side is positive.

\subsection{4. An exact recurrence, an independent finite formula and the complete table}\label{an-exact-recurrence-an-independent-finite-formula-and-the-complete-table}

For integer \(j\ge1\), differentiate \[
 \frac{s^{2j-1}}{(s^2+a_0)^{p-1}(s^2+b_0)^{N-1}}.
\] At zero its boundary value is zero. At infinity it is zero whenever \(2j+3<2p+2N\), which includes every recurrence use through moment167. Put \[
 D_j=2p+2N-2j-3,\qquad
 B_j=(2j-1)(a_0+b_0)-2(p-1)b_0-2(N-1)a_0.
\] Integrating the derivative and collecting the three powers of \(s^2\) gives the exact recurrence \[
 r_{j+2}=
 \frac{jB_j r_{j+1}+(2j-1)a_0b_0 r_j}
      {j(j+1)D_j}\quad(1\le j\le165).                 \tag{RM15}
\] Every denominator is positive. Together with RM14 this defines all167 rational coefficients. Induction using the actual integral identity RM9 proves that every computed coefficient is exactly \(L_j/\pi\), and hence positive. Positivity is not inferred merely from a floating recurrence.

A second finite expression independently checks the whole table. Extend \(\mathcal I\) by \(\mathcal I(0,v)=J_v(20)\) and \(\mathcal I(u,0)=J_u(2)\), for positive remaining argument. The elementary identity \[
 \mathcal I(u,v)=
 \frac{\mathcal I(u,v-1)-\mathcal I(u-1,v)}{\Delta}
 \quad(u,v\ge1)
\] follows by taking the difference of the two integrands. A rectangular array \(0\le u\le166,\ 0\le v\le444\), excluding \((0,0)\), gives all \(\mathcal I(u,444)\) independently of the partial-fraction sums. For \(1\le j\le167\), expansion of \(s^{2j-2}=((s^2+a_0)-a_0)^{j-1}\) now yields \[
 r_j=\frac{C}{(j-1)!}\sum_{r=0}^{j-1}
          \binom{j-1}{r}(-a_0)^{j-1-r}\mathcal I(p-r,N).
                                                               \tag{RM16}
\] All required \(p-r\) are nonnegative, including zero in the last formula. Every term is a convergent ordinary integral; there is no analytic continuation of a divergent individual term. The exact executable compares RM16 with RM14--15 for all167 degrees.

The companion certificate stores every full rational \(r_j\), not only a decimal logarithm, and every strict dyadic ceiling of RM6. Selected entries of that complete table are:

\begin{longtable}[]{@{}
  >{\raggedleft\arraybackslash}p{(\columnwidth - 24\tabcolsep) * \real{0.0769}}
  >{\raggedleft\arraybackslash}p{(\columnwidth - 24\tabcolsep) * \real{0.0769}}
  >{\raggedleft\arraybackslash}p{(\columnwidth - 24\tabcolsep) * \real{0.0769}}
  >{\raggedleft\arraybackslash}p{(\columnwidth - 24\tabcolsep) * \real{0.0769}}
  >{\raggedleft\arraybackslash}p{(\columnwidth - 24\tabcolsep) * \real{0.0769}}
  >{\raggedleft\arraybackslash}p{(\columnwidth - 24\tabcolsep) * \real{0.0769}}
  >{\raggedleft\arraybackslash}p{(\columnwidth - 24\tabcolsep) * \real{0.0769}}
  >{\raggedleft\arraybackslash}p{(\columnwidth - 24\tabcolsep) * \real{0.0769}}
  >{\raggedleft\arraybackslash}p{(\columnwidth - 24\tabcolsep) * \real{0.0769}}
  >{\raggedleft\arraybackslash}p{(\columnwidth - 24\tabcolsep) * \real{0.0769}}
  >{\raggedleft\arraybackslash}p{(\columnwidth - 24\tabcolsep) * \real{0.0769}}
  >{\raggedleft\arraybackslash}p{(\columnwidth - 24\tabcolsep) * \real{0.0769}}
  >{\raggedleft\arraybackslash}p{(\columnwidth - 24\tabcolsep) * \real{0.0769}}@{}}
\toprule\noalign{}
\begin{minipage}[b]{\linewidth}\raggedleft
j
\end{minipage} & \begin{minipage}[b]{\linewidth}\raggedleft
1
\end{minipage} & \begin{minipage}[b]{\linewidth}\raggedleft
2
\end{minipage} & \begin{minipage}[b]{\linewidth}\raggedleft
3
\end{minipage} & \begin{minipage}[b]{\linewidth}\raggedleft
100
\end{minipage} & \begin{minipage}[b]{\linewidth}\raggedleft
140
\end{minipage} & \begin{minipage}[b]{\linewidth}\raggedleft
141
\end{minipage} & \begin{minipage}[b]{\linewidth}\raggedleft
142
\end{minipage} & \begin{minipage}[b]{\linewidth}\raggedleft
143
\end{minipage} & \begin{minipage}[b]{\linewidth}\raggedleft
150
\end{minipage} & \begin{minipage}[b]{\linewidth}\raggedleft
160
\end{minipage} & \begin{minipage}[b]{\linewidth}\raggedleft
166
\end{minipage} & \begin{minipage}[b]{\linewidth}\raggedleft
167
\end{minipage} \\
\midrule\noalign{}
\endhead
\bottomrule\noalign{}
\endlastfoot
strict upper exponent for \(P_j\) & −315 & −1390 & −2366 & −82428 & −108534 & −109115 & −109693 & −110266 & −114169 & −119404 & −122351 & −122828 \\
\end{longtable}

The value in this row is the least integer \(b\) checked to satisfy \(P_j<2^b\). The rational \(P_j\) itself is available whenever a complete count benefits from retaining its precision. Floating hypergeometric comparisons in the historical probe are not used anywhere in this theorem or certificate.

\subsection{5. The infinite-mass zero term and the original positive row polynomial}\label{the-infinite-mass-zero-term-and-the-original-positive-row-polynomial}

The analytic measure alone does not repair a zero-degree term. Let the actual nonzero original row energies satisfy \[
 \alpha_\tau>0,\qquad \prod_\tau\alpha_\tau\ge\alpha_0^{\,f},
 \quad f=168,
\] and fix \(0<\theta<1\). For positive parameters \(c_\tau\) in the product measure, put \(\mathcal B=(\prod_\tau c_\tau)^{1/f}\). On those actual rows, AM--GM gives \[
 \sum_\tau c_\tau\alpha_\tau\ge f\mathcal B\alpha_0,
 \qquad
 e^{-\theta\sum c_\tau\alpha_\tau}
 \le(\theta\alpha_0\mathcal B)^{-f}.                     \tag{RM17}
\] The last inequality follows from \(e^{-u}\le1/u\) for \(u>0\). Consequently its full positive integral is at most \[
             (\theta\alpha_0)^{-f}L_1^f
                       <(\theta\alpha_0)^{-f}P_1.       \tag{RM18}
\] This estimate is applied before enlargement to an ambient lattice or the zero tuple. It uses a finite first inverse moment and never integrates a constant against the infinite-mass measure.

For clarity, the exact algebraic interface to the original coupled Gaussian count is the following. Use K109/K110's original covolume and dual metric, an independently proved all-parameter scalar coefficient \(K\), the original bound \(G'>c_0I\) with \(c_0=2^{36}\), \(d=k+3\), \(1\le k\le164\), and \(q=R/(Uz^2)\). The actual product-energy floor is \(\alpha_0=Rk/q\). After reserving RM17 and applying the original all-positive-weight Poisson inequalities, the complete row factor is bounded by \[
 \binom{164}{k}\frac{U^f\pi^{df}}{D_0^d}
 \left[
 \sum_{\substack{0\le l\le k,\ 0\le h\le3\\j=l+3-h\ge1}}
 \binom{k}{l}\binom3h K^{k-l+h}c_0^{-hf}
       \left(\frac{q}{(1-\theta)R}\right)^{jf}P_j
 +K^{k+3}c_0^{-3f}P_1
       \left(\frac{q}{\theta Rk}\right)^f
 \right].                                                \tag{RM19}
\] For each positive degree the reserved damping can be dropped and its moment is bounded by RM6. The last term is exactly the unique pair \((l,h)=(0,3)\), repaired by RM18. Every degree is at most167. All terms are nonnegative and the sum is finite, so Tonelli applies with the stated finite bounds.

RM19 is the same complete scalar/transverse algebra as R098 WJ28, with the new measure and its exact moment table. It introduces no independent beta norm fibre, row restriction, terminal-functional freedom or finite-pattern shortcut. At degree167 any exceptional positive finite-pattern exponent must still be counted in a subsequent arithmetic theorem. The central coset, its \(U^f\) Jacobian, all source/target/Haar/class factors, finite supports and labels, any external exponential coefficient, near-point factor and K118's complete band-tail term all remain outside RM19 and must be paid by the caller. In particular no independent central-profile and band-profile inequalities are multiplied without a new joint proof.

\subsection{6. Scope and source boundaries}\label{scope-and-source-boundaries}

The new ingredients are the positive Gamma444 approximation of the exact exponential shift, the finite rational evaluation of every required inverse moment, and its explicit compatibility with the existing zero-degree reservation. The theorem makes no sector contribution, error-charge, full support28, signed-sum or nonvanishing assertion.

Primary normalization, freshly read with its positive-real domains: National Institute of Standards and Technology. (n.d.). \emph{Digital Library of Mathematical Functions}, equations5.9.1 and5.12.1--5.12.3, \url{https://dlmf.nist.gov/5.9} and \url{https://dlmf.nist.gov/5.12}. All new convergence conditions, finite partial fractions and integration-by-parts identities are proved above; neither analytic continuation nor a special-function numerical evaluator is required.

Exact carried interfaces are K109 TG3--12, K110 PM10--12, K116 RE11--20 and R098 WJ18--28, with their original actual-row and infinite-mass qualifications. K118 supplies a separately checked full-probability arithmetic tail only when used at its stated bounded support27 scope. An application to the actual trace additionally requires the separate genuine selected-kernel band estimate; this theorem does not assert that estimate before its review. The original scheme, projective modules, involution, polarization, theta line, Poincare pairing, selected vector and finite projectors are unchanged. Carried contextual anchors remain SGA3 I4.7.3/4.7.3.1; SGA6 VI Theorem3.3; FGA232 Theorem3.1 pp148--149; FGA236 Corollaries1.7/2.3 pp229/232; SGA7 I VII2.7/2.9 pp169--174 and VIII7.4(b) pp310--311; cote4 scan003 sections8.1.1--8.1.2/scans005--006 sections8.4.1--8.4.3 and cote54 scans005--007. These are inherited contextual locators, not new arithmetic comparison results. No original-archive theorem-validity audit or Lean build is performed.

\subsection{Canonical algebraic-geometric context}\label{canonical-algebraic-geometric-context-73}

This is an estimate at the archimedean realization of the same projective OE modules and polarized CM abelian data. The positive Gamma convolution and exact inverse-moment evaluation do not alter the scheme over its original arithmetic base, the involution, Hermitian lattice, theta line, Picard scheme or Poincare biextension. Carried contextual anchors are SGA3 I4.7.3/4.7.3.1; SGA6 VI Theorem3.3; FGA232 Theorem3.1 pp148--149; FGA236 Corollaries1.7/2.3 pp229/232; SGA7 I VII2.7/2.9 pp169--174 and VIII7.4(b) pp310--311; cote4 scan003 sections8.1.1--8.1.2 and scans005--006 sections8.4.1--8.4.3; cote54 scans005--007. These locators are inherited context, not new primary readings or newly proved comparisons. No new formal-group/Dieudonne, lambda-ring/K-theory, l-adic weight, CM-motive or motivic L-function result follows from this real analytic inequality alone. DP\_MULT Q3e/DP44, THETA3 G1/A7, CX\_GRH CX-ALG-024 and GCT\_PLETH T-052 retain their current arithmetic and global gates until a separately checked composition changes them. No Lean build or archive theorem-validity audit is performed.

\clearpage

\section{AM-R-100 --- a narrow-band joint profile with a three-bit field cost}\label{am-r-100-a-narrow-band-joint-profile-with-a-three-bit-field-cost}

\textbf{Canonical checked integration --- 2026-11-23T11:05:35.837469+00:00.} Complete saved-proof and executable reviews and separate exact diagnostic replays passed. No numerical error charge is introduced. Frozen author history is preserved. AM-R-100 NB1-16 proves the actual selected fidelity restricted to I*={[}16/529,16/527{]} has the same original joint physical-central kernel with lambda32/25. The universal angular loss coefficient527/33858 exceeds1/65 and both conditional-density branches, zero height, full radial spread and detailed balance are retained. Original non-strict norm\textless=133 yields precisely one full near-point2; the entire intercept exp(128f/13175) costs less than2\^{}3. Central integration may remove one radial power to shape166 while retaining its full physical U\^{}f covariance and original finite costs. Passage from restricted fidelity to actual selected traces explicitly requires the new narrow probability NB15; under it K118 band-independent fidelity/mass reasoning gives the full additive tail2\^{}-8605114. No wide-band probability implication, K120 moment premise, standalone charge or sector removal is inferred. Existing132/31 arrays and complete floors27/11 remain unchanged.

Author: \nolinkurl{/root}, GPT-6 Astra Pro; inherited effort is not externally exposed. 23 November 2026 UTC. \textbf{Controlled pending proof.} Whole author and independent checks are separate. This analytic theorem introduces no arithmetic error charge.

\subsection{1. Original object and precise conclusion}\label{original-object-and-precise-conclusion}

Keep the actual selected J56 representation and its vector-valued Hilbert metric, both compact factors, fields, involution, original pencils, finite projectors, phases, class masses and Haar normalization. Use DO080 RT4--14's invariant measure and selected probability density. For \[
 I_*=[16/529,16/527],\qquad
 F_*(g)=\int_{Y(z)\in I_*,\,Y(g^{-1}z)\in I_*}
       \sqrt{\rho(z)\rho(g^{-1}z)}\,d\mu(z),
 \tag{NB1}
\] only the restricted nonnegative fidelity is estimated here. Write \[
 a=f=168,\quad \nu=333/2,\quad
 \vartheta=\frac{4^{167}}{334\binom{334}{167}},\quad
 c_* =529/527,\quad \gamma_*=64/25,\quad \lambda_*=32/25.
 \tag{NB2}
\] For every Cartan radius \(r>0\), with \(c=\cosh r\), \(s=\sinh r\), \[
 F_*(g)\le \frac{\vartheta2^{-933/2}}s
       \left(\frac2{c+1}\right)^\nu
       \exp[-\gamma_*(c-c_*)].
 \tag{NB3}
\] Also \(0\le F_*\le1\); at \(r=0\) only that bound is used.

For the actual nonidentity globally rank-two-unipotent matrices in the original norm cutoff \(\|h_\tau\|\le133\), define \[
 \alpha=(N^2)_{nn}>0,\quad \alpha_\tau\le R=17424/133,
 \quad w=\Im N_{nn}=-U\Im d,\quad b(t)=1+t/4,
\] \[
 J_*(t,v)=t^{-1/2}b(t)^{-167}
       \left(1+\frac{v^2}{4b(t)^2}\right)^{-333/4}
       e^{-\lambda_*t}.
 \tag{NB4}
\] At the unchanged actual near points, the entire field product satisfies \[
 \prod_\tau F_*(g_{p,\tau}^{-1}h_\tau g_{p,\tau})
 <2\,2^{-649}\,2^{-933f/2}\,2^3
                      \prod_\tau J_*(\alpha_\tau,w_\tau).
 \tag{NB5}
\] The first2 is the single complete near-point cost. The factor \(2^3\) pays the entire field intercept. The selected trace is not bounded by the restricted product alone; section4 states its required full-tail comparison.

\subsection{2. The new angular constants in the same kernel}\label{the-new-angular-constants-in-the-same-kernel}

For \(y,y'\in I_*\), put \(D=\sqrt{y/y'}\). Then \[
 D,D^{-1}\le\sqrt{529/527}<528/527<c_*,
 \qquad \frac y{1-y}\ge\frac{16}{513}.
 \tag{NB6}
\] The strict squared comparison is \(529\cdot527=528^2-1\).

When \(c\ge c_*\), both forward and reverse directions lie in the first angular case. The exact R099 AH8 identity for \(K=1-(c-D)^2/[(c^2-1)(1-y)]\) gives \[
 1-\frac K{2D/(c+1)}
 =\frac{(D-1)^2+y(c-1)(c+1-2D)}
        {2D(c-1)(1-y)}
 \ge\frac{527}{33858}(c-c_*)
 \ge\frac{c-c_*}{65}.
 \tag{NB7}
\] Indeed \(c+1-2D\ge c-c_*\ge0\), and \(65\cdot527=34255>33858\). Every denominator is positive. For feasible positive height, the last inequality forces \(0\le(c-c_*)/65<1\). Thus \[
 K^\nu\le\left(\frac{2D}{c+1}\right)^\nu
             e^{-(\nu/65)(c-c_*)}
 \le\left(\frac{2D}{c+1}\right)^\nu
             e^{-\gamma_*(c-c_*)},
 \tag{NB8}
\] since \(\nu/65=333/130>64/25\). If \(K\le0\), the directional density is zero almost everywhere, and its zero extension obeys the same bound. For \(1<c<c_*\), use K108 AE5--7's two original cases, while the new exponential exceeds1. No first-case formula is imposed when \(c\le D\).

The remaining kernel argument is exactly the accepted R099 AH11--13: multiply the old real Cartesian Jacobian into this one angular density, apply it also in the reverse direction, and use detailed balance on the symmetric set \(I_*\times I_*\). The geometric mean cancels \(D\) and retains one complete exponential factor. The restricted radial integral is bounded by the full positive K107 SV8--12 spread integral. Its full lower tail is already included in the coefficient \(2^{-933/2}\). This proves NB3 with unchanged normalization and both compact factors. It does not remove the selected trace outside \(I_*\).

\subsection{3. Original joint energy, physical central variable and total field cost}\label{original-joint-energy-physical-central-variable-and-total-field-cost}

The exact original Cartan identity and joint inequalities are DO122 CE3--6, read with its mandatory non-strict-cutoff addendum. They hold at the actual near point with \(t=\alpha_p\), \(v=w_p\): \[
 c_p^2=(1+t/2)^2+v^2,\quad
 c_p^2-1\ge t b(t),\quad
 (c_p+1)^2\ge4b(t)^2+v^2,\quad c_p\ge1+t/2.
 \tag{NB9}
\] Consequently the single NB3 kernel yields \[
 F_*(h_p)\le\vartheta2^{-933/2}e^{128/13175}
                         J_*(\alpha_p,w_p).
 \tag{NB10}
\] Here \(\gamma_*(c_*-1)=128/13175\), and \(\gamma_*/2=\lambda_*\). The factor in NB10 comes from that same angular height, not from multiplication of two independent trace bounds.

The accepted physical near-point estimates are \[
 \alpha_p\ge e^{-2r_p}\alpha,\qquad
 |w_p-w|\le6Hr_p,\qquad r_p<2^{-352},\quad H=133.
 \tag{NB11}
\] They use the original rational representative's bound \(\|h_\tau\|\le H\), not a bound on the full compact part of \(g_p\). R099 AH18's derivative calculation applies with the new positive \(\lambda_*\): \[
 0<-\partial_{\log t}\log J_*
 \le335/2+\lambda_*t,\qquad
 |\partial_v\log J_*|\le(333/4)/2.
 \tag{NB12}
\] The energy comparison is made only between \(e^{-2r_p}\alpha\) and \(\alpha\), where \(t\le R\). Monotonicity permits this even when \(\alpha_p>\alpha\). Hence \[
 J_*(\alpha_p,w_p)
 \le e^{C_*r_p}J_*(\alpha,w),\quad
 C_*=335+2\lambda_*R+3H(333/4)<2^{16}.
 \tag{NB13}
\] Since \(2^{16}f2^{-352}<1/2\), the product transport costs less than2. This replaces the earlier single near-point2. The exact original \(\vartheta^f<2^{-649}\) remains unchanged.

The whole intercept costs less than three bits: \[
 \exp\left(\frac{128f}{13175}\right)<e^2
                 <(11/4)^2=121/16<8.
 \tag{NB14}
\] The first inequality is \(128\cdot168=21504<26350\). For a self-contained rational upper bracket, the positive exponential series through degree4 and its geometric remainder give \(e\le65/24+1/100<11/4\): the tail starts at \(1/5!\), and later term ratios are at most \(1/6\). Combining NB10--14 proves NB5.

The exponential in \(J_*\) is independent of the physical central variable. Accordingly the original DO122 CE14--20 weighted affine central integral can consume precisely one factor \(b\) while retaining the shift \(\lambda_*\), leaving \(t^{-1/2}(1+t/4)^{-166}e^{-\lambda_*t}\). The complete original affine lattice, \(U^f\) covolume, full750 cover, finite patterns and row count must still be proved in any arithmetic application. No separate weighted-Jacobian coefficient is multiplied into NB5.

\subsection{4. Explicit probability premise for an actual-trace application}\label{explicit-probability-premise-for-an-actual-trace-application}

The theorem NB3--5 concerns the restricted fidelity and needs no concentration premise. To pass to actual traces, suppose a separately proved uniform statement gives \[
                         \Pr(Y\notin I_*)<2^{-40000000}.
 \tag{NB15}
\] K118 BT15--22 then applies with this new measurable band: its two Cauchy--Schwarz complement estimates and product telescoping depend only on that probability, not on the old endpoints. Its complete original bounded support27 containing mass also has no dependence on the band. Thus for every actual subset \(\mathscr S\) of its all-II globally rank-two-unipotent scalar-cross support27 family with nonzero original S-cross and \(x<2^{3328}\), \[
 \sum_{h\in\mathscr S}|w(h)c(h)|
 \le\sum_{h\in\mathscr S}|w(h)|\prod_\tau F_*(h_{p,\tau})
                              +2^{-8605114}.
 \tag{NB16}
\] This is a conditional application with the explicit premise NB15; it does not infer the narrower probability from the old wider band. A separate exact moment calculation may establish NB15. The new Gamma512/rate400 majorant of the remaining shape166 profile is also a separate input, obtained by the K119 positive-integral method with \(512/400=32/25\). No unreviewed table is part of this theorem.

The exact shifted representing measure itself has infinite total mass and finite all positive integer inverse moments, by the general R099 AH22--24 statement. A row-polynomial degree0 still requires the actual-row product-energy reserve before Gaussian enlargement. No arithmetic cell or signed sum is estimated by NB3--16 alone.

\subsection{5. Sources and geometric scope}\label{sources-and-geometric-scope}

The actual selected object and invariant disintegration are DO080 RT4--19; the two angular cases are K108 AE3--9; the complete radial spread is K107 SV8--12; the exact angular identity, same-kernel composition and positive shifted-measure scope are R099 AH8--24. Original joint Cartan and physical central transport are DO122 CE3--20 with \nolinkurl{DO122_CUTOFF_SOURCE_ADDENDUM.md}; the original cutoff is R048-RC RC48.4/RC48.6 and the paragraph afterRC48.7. K118 BT15--22 supplies only the explicitly conditional full arithmetic-tail comparison. None of their numerical removal charges is imported.

Primary metric normalization remains Nakahama, R. (2016), \emph{Norm computation and analytic continuation of vector valued holomorphic discrete series representations}, Journal of Lie Theory,26(4),927--990, equation5.7 and Theorem5.1, printed960--961, through the checked selected-module dictionary. Ordinary positive integral normalization is National Institute of Standards and Technology (n.d.), \emph{Digital Library of Mathematical Functions},5.9.1 and5.12.1/3. The new rational constants and same-kernel inequality are proved above.

This is an archimedean estimate on the same original projective OE module and polarized CM data. It changes no scheme, involution, theta line, Picard scheme, Poincare biextension or arithmetic selected image. Inherited exact contextual anchors remain SGA3 I4.7.3/4.7.3.1; SGA6 VI Theorem3.3; FGA232 Theorem3.1 pp148--149; FGA236 Corollaries1.7/2.3 pp229/232; SGA7 I VII2.7/2.9 pp169--174 and VIII7.4(b) pp310--311; cote4 scan003 sections8.1.1--8.1.2/scans005--006 sections8.4.1--8.4.3 and cote54 scans005--007. No new formal-group/Dieudonne slope, lambda-ring/K-theory positivity, l-adic weight or CM-motive L-function identity is asserted. DP\_MULT Q3e/DP44, THETA3 G1/A7, CX\_GRH CX-ALG-024 and GCT\_PLETH T-052 retain their arithmetic/global gates. There is no original-theorem validity audit or Lean build.

\subsection{Canonical algebraic-geometric context}\label{canonical-algebraic-geometric-context-74}

This is an estimate at the archimedean realization of the same projective OE modules and polarized CM abelian data. The positive shifted measure and simultaneous radial restriction do not alter the scheme over its original arithmetic base, the involution, Hermitian lattice, theta line, Picard scheme or Poincare biextension. Carried contextual anchors are SGA3 I4.7.3/4.7.3.1; SGA6 VI Theorem3.3; FGA232 Theorem3.1 pp148--149; FGA236 Corollaries1.7/2.3 pp229/232; SGA7 I VII2.7/2.9 pp169--174 and VIII7.4(b) pp310--311; cote4 scan003 sections8.1.1--8.1.2 and scans005--006 sections8.4.1--8.4.3; cote54 scans005--007. These locators are inherited context, not new primary readings or newly proved comparisons. No new formal-group/Dieudonne, lambda-ring/K-theory, l-adic weight, CM-motive or motivic L-function result follows from this real analytic inequality alone. DP\_MULT Q3e/DP44, THETA3 G1/A7, CX\_GRH CX-ALG-024 and GCT\_PLETH T-052 retain their current arithmetic and global gates until a separately checked composition changes them. No Lean build or archive theorem-validity audit is performed.

\clearpage

\section{AM-DO-123 --- retained-band integration over the original central coset}\label{am-do-123-retained-band-integration-over-the-original-central-coset}

\textbf{Canonical checked integration --- 2026-11-23T11:07:57.846500+00:00.} Complete root and independent proof/script/interface reviews and separate exact replays passed. The actual canonical predecessor has132/31 entries; this integration gives133/31, complete floors27/11. Frozen author history is preserved. AM-DO-123 BC1-27 removes the entire original bounded support27/total-eligibility27 all-II full-positive scalar-cross globally rank-two-unipotent k141..142 union, with b=j\_int0 and2\^{}2111\textless x\textless2\^{}3328. The same R099 selected angular kernel, original physical central coordinate w=-U Im d, full affine covolume/U scaling/750 cover, shape166/shift111/100 moments and actual-row zero-degree repair retain every polynomial, finite, label, support, pivot, class/Haar and pattern factor. Full exponents-3344299/1024 and-1688491/1024 plus the entire K118 tail2\^{}-8605114 yield one1647 charge. Appending it once to actual DO122132/31 gives133/31 with both entire sums below2\^{}-230 and complete floors27/11 unchanged. The sole remaining support27 necessary alternative is eligibility27,143\textless=k\textless=164,2\^{}2111\textless x\textless2\^{}3328,b=j\_int0. The k143 envelope-28587/1024 is insufficient and unclaimed. No full support28, signed nonvanishing, source admission, selected arithmetic-image content, prescribed period, specific032 value, evaluatedDP44 or RH conclusion.

Author: \nolinkurl{/root/definite_oscillator}, 23 November 2026 UTC. \textbf{Controlled pending proof.} Whole author, independent and root reviews are distinct from exact diagnostic replays. R099 and K119 completed their separate whole reviews and were canonically accepted at10:55:51 and10:55:52 UTC respectively; this application still requires its own whole independent and root reviews. No canonical file or error array is changed here.

\subsection{1. Actual sector and complete statement}\label{actual-sector-and-complete-statement}

Retain the original J56 all-top H133 kernel, CM fields and involution, source/target pencils and omitted labels, all selected vector-valued Hilbert spaces, compact factors, binary words, canonical twists, phases, Haar measures, class masses and finite projectors. Put \[
 f=a=168,\quad H=133,\quad R=17424/133,\quad
 Q=1427908193729121522601,\quad U=77Q^2-14,
\] \[
 D_0=7^{140}29^{144}5^{84}13^{84},\qquad D_F=2^{168}D_0.
 \tag{BC1}
\] The actual baseline is DO122,132 all-top and31 uniform entries, with complete floors27/11. Consider its remaining globally rank-two-unipotent scalar-cross support27 family with all active typesII, all334 unmixed G reductions unipotent, and \[
 q_{\rm total}=27,\quad b_{\rm inactive}=j_{\rm int}=0,
 \quad 2^{2111}<x=N_ED_\eta<2^{3328},\quad k\in\{141,142\}.
 \tag{BC2}
\] Here k is the number of nonzero original S covector values. The radial row parameter q below is different from total eligibility.

\textbf{Theorem.} The complete actual absolute normalized contribution of BC2 is less than2\^{}-1647. This includes the entire K118 discarded-band tail once, after both coordinate sectors are combined. Append one1647 charge to the full actual chronological predecessor; retain every previous charge and both complete error sums. The floors remain27/11. The k143 sufficient envelope fails the required230-bit budget and is explicitly unclaimed.

The improvement integrates the same retained-band angular kernel in the original physical central coordinate. It does not multiply independently proved trace bounds or reuse an old numerical sector charge.

\subsection{2. One restricted kernel with its full external tail}\label{one-restricted-kernel-with-its-full-external-tail}

For the actual compact-invariant selected density rho against invariant measure mu, let F\_I(g) be its fidelity integral restricted to both radial variables Y(z),Y(g\^{}-1 z) in I={[}1/35,1/31{]}. It is nonnegative and at most1, but the actual selected coefficient is not bounded by F\_I alone. K118 BT1--22 proves, for every actual subset S of the bounded all-II family in BC2 and its original finite weights, \[
 \sum_{h\in S}|a(h)c(h)|
 \le\sum_{h\in S}|a(h)|\prod_\tau F_{I,\tau}(h_{p,\tau})
       +2^{-8605114}.
 \tag{BC3}
\] Its tail already includes the full168-label probability loss and a containing count of all supports, rows, labels, patterns and central affine cosets. It is valid simultaneously for the union of the two k sectors here. No per-cell multiplication of that tail is needed.

Write h=I+N in the accepted full-positive normal form. The original covector r and isotropic slope eta satisfy \[
 N^3=0,\quad \operatorname{rank}_E N=2,\quad
 \eta\bar\eta=U,\quad
 \alpha=(N^2)_{nn}=U rJ_Z^{-1}r^*>0,\quad\alpha_\tau\le R,
\] \[
 N_{nn}=-Ud,\qquad \Re N_{nn}=\alpha/2,\qquad
 w=\Im N_{nn}=-U\Im d.
 \tag{BC4}
\] These are the original DO098 CF4--6 identities, applied at support27 through K113 SC27.1--7. The original arithmetic representative satisfies \textbar\textbar h\_tau\textbar\textbar≤133, the non-strict R048-RC cutoff stated in the mandatory DO122 cutoff addendum.

Let p=333/4, lambda=111/100 and b(t)=1+t/4. R099 AH4 and AH16--19 prove from a single restricted angular kernel that \[
 \prod_\tau F_{I,\tau}(h_{p,\tau})
 <2\,2^{-649}\,2^{-933f/2}\exp(111f/250)
   \prod_\tau J_\lambda(\alpha_\tau,w_\tau),
\] \[
 J_\lambda(t,v)=t^{-1/2}b(t)^{-167}
       \left(1+\frac{v^2}{4b(t)^2}\right)^{-p}e^{-\lambda t}.
 \tag{BC5}
\] Specifically the same original Jacobian kernel receives the height factor exp{[}-(111/50)(cosh(r)-6/5){]}. The exact identities cosh(r)\^{}2=(1+t/2)\^{}2+v\^{}2, sinh(r)\^{}2≥t b(t) and (cosh(r)+1)\^{}2≥4b(t)\^{}2+v\^{}2 give BC5 directly. The alternative weighted-Jacobian energy profile of R099 AH20 is not multiplied into this formula.

The actual near-point identities are alpha\_p≥e\^{}-2r\_p alpha and \textbar w\_p-w\textbar≤6Hr\_p.~For u=v²/(4b(t)²+v²), \[
 -\partial_{\log t}\log J_\lambda
 =\frac12+(167-2pu)\frac{t}{t+4}+\lambda t>0,
 \qquad |\partial_v\log J_\lambda|\le p/2.
 \tag{BC6}
\] Only the original interval e\^{}-2r\_p alpha≤t≤alpha≤R is used for the transport derivative. Its full coefficient335+2lambdaR+3Hp is less than2\^{}16. Since r\_p\textless2\^{}-352, the product transport is below2, exactly the single near-point factor already shown in BC5. R099 AH21's rational Taylor bound gives \[
                       \exp(111f/250)<2^{108}.
 \tag{BC7}
\] The additional band cost is therefore fully paid before the row enlargement.

\subsection{3. Original row and central difference lattice}\label{original-row-and-central-difference-lattice}

Fix eta and the canonical common whole-C G pattern J. The original four local states have E-norm exponents2,1,0 with multiplicities1,2,1. Let J0 be the real ideal consisting of each rank-zero G prime once. The exact original table gives \[
 |a_G(h)|\le2Q^{-335\cdot167}(N_EJ)^{335/2},\quad
 N_EJ\le Q^{334},\quad N_EJ\le Q^{167}N_FJ_0,
 \quad N_FJ_0\le Q^{167}.
 \tag{BC8}
\] At both original CM partners, every integral Z input has r(y),eta r(y) in J I\_Sigma. Their exact valuation intersection gives all167 values in \[
 L=J I_\Sigma D_\eta,\quad N_EI_\Sigma=T^{-54},\quad
 N_EL=z^{2f}=N_EJ\,xT^{-54},\quad q=R/(Uz^2)\ge k.
 \tag{BC9}
\] The three saturated transverse inputs are actual integral inputs, not values divided by2. Their Gram matrix G' has determinant U and G'\textgreater2\^{}36 I. The true covector metric is diag(I\_k,(G')\^{}-1), and the real covolume is D0\^{}(k+3) z\^{}(2(k+3)f)/U\^{}f.

For fixed eta,r and pattern, all central d values lie in one affine coset with difference lattice \[
 \Gamma_d=iJ_0B_\eta I_{\Sigma,F}^{2},\quad N_FB_\eta=x,
 \quad \operatorname{covol}\Gamma_d
       =\sqrt{D_F}(N_FJ_0)xT^{-54}.
 \tag{BC10}
\] This is the original four-positive-plus-negative entry calculation DO104 CD14--19, transported to27 in K113 SC27.7. The primitive Bezout extraction gives Delta d,eta Delta d,bar(eta)Delta d in the full-entry ideal, with no division byU or2. The local extra exponent is max(0,-v(eta),-v(bar eta)), precisely that of D\_eta bar(D\_eta); its real contraction is B\_eta. Zeros or poles meeting U,G,T cause no exception. The real-part shift is retained.

In the physical variable w=-U Im d the full real lattice therefore satisfies \[
 \operatorname{covol}\Gamma_w
 =U^f\sqrt{D_F}(N_FJ_0)xT^{-54},\quad
 \mu(\Gamma_w)<750\sqrt f\,U z_c,
\] \[
 z_c^f=(N_FJ_0)xT^{-54},\qquad
 U z_c=\frac Rq\left(\frac{N_FJ_0}{N_EJ}\right)^{1/f}<2^{-83}.
 \tag{BC11}
\] The last global inequality follows from the exact actual-host comparison U\^{}f Q\^{}167 2\^{}(3328+83f)\textless T\^{}54. It holds for every simultaneous actual q,J,J0, rather than treating a row norm and a divisor ratio as independent extrema. The coefficient750 is the full real codifferent/trace covering bound.

\subsection{4. Complete weighted integration of the physical central coset}\label{complete-weighted-integration-of-the-physical-central-coset}

For the fixed row set b\_tau=1+alpha\_tau/4 and phi(w)=prod\_tau(1+w\_tau²/(4b\_tau²))\^{}-p.~The exponential exp{[}-lambda sum alpha\_tau{]} is constant on that central fibre. The global bound \textbar\textbar grad log phi\textbar\textbar≤p sqrt(f)/2 and the true covering radius imply, by integration over translated Voronoi cells, \[
 \sum_{w\in w_0+\Gamma_w}\phi(w)
 \le\frac{\exp(375pf U z_c)}{\operatorname{covol}\Gamma_w}
          \prod_\tau\int_{\mathbb R}
                 (1+t^2/(4b_\tau^2))^{-p}\,dt.
 \tag{BC12}
\] This sums a positive upper envelope on a containing affine lattice; no enlarged point is declared to be an actual arrow. Monotone convergence is legitimate, since p\textgreater1/2 and all displayed one-dimensional integrals converge. The original pattern's finite coefficient bound remains fixed throughout this enlargement.

The ordinary positive Beta integral and Cauchy--Schwarz for Gamma give \[
 \int_{\mathbb R}(1+t^2/(4b^2))^{-p}\,dt
 =2b\sqrt\pi\frac{\Gamma(p-1/2)}{\Gamma(p)}
 \le\mathcal G b,\qquad \mathcal G=4\sqrt{\pi/329}.
 \tag{BC13}
\] Indeed Gamma(p-1/2)\^{}2≤Gamma(p-1)Gamma(p) and p-1=329/4. The exact bound375pf2\^{}-83\textless1/2 makes the complete Voronoi inflation less than2. Combining BC5 and BC10--13 gives \[
 \sum_{\text{actual central fibre}}\prod_\tau F_{I,\tau}(h_{p,\tau})
 <4\,2^{-649}\,2^{-933f/2}\exp(111f/250)
 \frac{\mathcal G^fT^{54}}{U^f\sqrt{D_F}(N_FJ_0)x}
 \prod_\tau\Phi(\alpha_\tau),
\] \[
                   \Phi(t)=t^{-1/2}(1+t/4)^{-166}e^{-\lambda t}.
 \tag{BC14}
\] The two factors2 in4 are exactly near-point transport and the weighted Voronoi comparison. The central integral consumes one power of1+t/4, leaving shape166. The physical U\^{}f Jacobian has not been dropped.

\subsection{5. All positive moments, including the repaired zero degree}\label{all-positive-moments-including-the-repaired-zero-degree}

Let mu be the pushforward of Gamma(shape166,rate4) times ds/sqrt(pi s) by c=lambda+b+s. It is a positive locally finite measure of infinite total mass, and Phi(t)=int exp(-ct)dmu(c). Its positive inverse moments are \[
 M_j=\frac1{\Gamma(j)}\int_0^\infty
          t^{j-3/2}(1+t/4)^{-166}e^{-\lambda t}\,dt,
 \qquad j\ge1.
 \tag{BC15}
\] All are finite because lambda\textgreater0 and j\textgreater1/2. The K119 rational moment theorem bounds them by pi r\_j. Its explicit N=444 positive integral majorant uses \[
 e^{-\lambda t}\le(1+t/400)^{-444},\quad
 r_j=\frac{2\,4^{166}400^{444}}{\pi(j-1)!}
       \int_0^\infty
       \frac{s^{2j-2}}{(s^2+4)^{166}(s^2+400)^{444}}\,ds.
 \tag{BC16}
\] Two exact partial-fraction base integrals and an integration-by-parts recurrence give rational r\_j for every1≤j≤167. The theorem proves positivity and normalization of these constants; finite recurrence diagnostics are not substituted for that positive integral proof. Define the exact strict dyadic ceilings \[
 b_j=\left\lceil\log_2[((22/7)r_j)^{168}]\right\rceil_{\rm strict},
 \qquad P_j=2^{b_j}>M_j^f.
 \tag{BC17}
\] Every required value is independently checked from the final frozen K119 table. No zeroth moment or finite total mass is asserted.

On a genuine row the product energy floor is prod alpha\_tau≥(Ukz²)\^{}f=(Rk/q)\^{}f.~Reserve theta=2\^{}-20 of the Gaussian exponent before enlarging the row. With B=(prod c\_tau)\^{}(1/f), AM-GM and e\^{}-u≤u\^{}-1 imply \[
 e^{-\theta\sum c_\tau\alpha_\tau}
 \le e^{-\theta fBRk/q}
 \le(\theta Rk/q)^{-f}\prod c_\tau^{-1}.
 \tag{BC18}
\] Thus the sole zero-degree term integrates to at most P1(q/(theta Rk))\^{}f.~Every positive degree integrates with its corresponding P\_j. No per-embedding energy floor is assigned to a zero or out-of-cutoff tuple introduced by the later positive Gaussian enlargement.

Keep the complete universal degree15 scalar coefficient K of K111 HP1--18, with its original half-trace minimum, and c0=2\^{}36. Explicitly \[
 c_{336}=168\binom{336}{168}/4^{168},\quad
 M_{15}=\frac{6015684926187167549292201834533348561577984}
 {2618017117694612446410293519},
\] \[
 K=\frac{336M_{15}}{c_{336}}(7841/5000)^{168}168!
 \left[\frac{(365729/500)^2}{168(314159/100000)^2}\right]^{168}.
 \tag{BC19}
\] The exact K110 PM10--12 product-one change and anisotropic Poisson bound apply for every positive c\_tau, without declaring the changed real lattice to be an ideal. For d\_k=k+3 the complete row factor is binom(164,k)U\^{}f pi\_+\^{}(d\_k f)/D0\^{}d\_k times \[
 \mathcal P_k(q)=
 \sum_{\substack{0\le l\le k,\ 0\le h\le3\\ j=l+3-h\ge1}}
 \binom kl\binom3h K^{k-l+h}c_0^{-hf}
       \left(\frac q{(1-\theta)R}\right)^{jf}P_j
 +K^{k+3}c_0^{-3f}P_1\left(\frac q{\theta Rk}\right)^f.
 \tag{BC20}
\] All scalar degrees and all four transverse terms remain, including(l,h)=(0,3). The row U\^{}f cancels the physical central U\^{}-f in BC14 exactly. The actual energy and band restrictions were applied before this enlargement; no new actual-arrow assertion is made about its points.

\subsection{6. Every original finite factor and both endpoint derivatives}\label{every-original-finite-factor-and-both-endpoint-derivatives}

K112 active eligibility and K113 bounded paired uniqueness give at most one eta for each original27-place support throughout x\textless2\^{}3328. Eta is therefore fixed before the row and central count. DO105 EL9--16's actual same-label argument fixes the only possible active label, at both partners and including central-only poles. In BC2 the active cross pole persists when the central coordinate varies. The complete central bound is for one containing affine coset for that actual pair; no central-dependent label union is suppressed.

The141 inactive factors are fixed by the same cross data, as in DO108 AC8--10. The full168-factor bound is(T−1)\^{}-168\textless2T\^{}-168. Keep supports\textless2\^{}104, the retained original first-pivot choices\textless2\^{}8, Green2, the T correction2, near-point2, weighted-cover2 and conservative full-pattern4. These give118 fixed bits. The full T power is335·27+54−168=8931. Class/Haar masses and mixed contractions remain exactly those of the original normalized coefficient. Eta,r,d reconstruct the original matrix, so there is no unpaid beta, negative-column or source-isometry fibre.

For each positive polynomial degree j, and effective degree1 for the repaired term, all q and x dependence is \[
                 x^{-337/2}q^{(j-335/2)f}.
 \tag{BC21}
\] Its q derivative divided byf is j−335/2\textless0 for all degrees here. At the all-rank-zero endpoint \[
 q_*^f=\frac{R^f T^{54}}{U^fQ^{334}x},
 \qquad q=q_*D^{1/f},\quad D=Q^{334}/N_EJ,
 \quad (N_FJ_0)^{-1}\le Q^{-167}D,
 \tag{BC22}
\] the coupled x logarithmic derivative is exactly−1−j\textless0. The cover inflation was already bounded uniformly in BC11--14; no untracked moving central radius remains.

For each degree the entire original four-state containing sum is \[
                 \prod_{v=1}^{334}(1+q_v^{j-333/2})^2<2.
 \tag{BC23}
\] The two target sectors use j≤145; the failed diagnostic uses j≤146. Thus j−333/2≤−20.5\textless−1 throughout, and the exact old inequality(1+1/g0)\^{}668\textless2 suffices. We retain4 conservatively in the118 bits. There is no additional pattern or orientation factor and no extension to an exceptional degree167 is needed.

\subsection{7. Exact complete endpoints and the paid union}\label{exact-complete-endpoints-and-the-paid-union}

Use the actual-host directed logarithms L\_T=1019333/1024, L\_Q=287843/4096, L\_U=601356/4096 and L\_R=28810/4096, with the same upper/lower directions as DO122. At x=2\^{}2111 the exact bracket is251124964\textless q\_*\textless251124965. Also \[
        \frac{\mathcal G^f}{\sqrt{D_F}}
        \le\frac{(16(22/7)/329)^{84}}{\sqrt{D_F}}<2^{-1110}.
 \tag{BC24}
\] After BC7's108 band bits are included, the common exponent is \[
 C=118-933f/2-649-1110+108
     +(335f/2)(L_R-L_U)+8931L_T-336\cdot167L_Q
     -(337/2)2111
   =591877205/1024.
 \tag{BC25}
\] All factors precede the directed logarithmic substitution. In particular no old central ball volume or independent weighted-Jacobian improvement appears.

The strict scalar ratio ceiling is ceil(log2{[}K((1-theta)R/q\_-)\^{}f{]})=-935. Exact moment ceilings and the full binomial/transverse coefficients bound every nonmain term of BC20 relative to its(l,h)=(k,0) term. For k141 the567 nonmain terms are each at most2\^{}-358 times that term; for k142 the571 nonmain terms are each at most2\^{}-361 times it. Their sums are strictly less than the main term. The entire positive polynomial is therefore bounded by twice the main term. The dyadic moment P\_(k+3) is included in that complete rational ceiling.

\begin{longtable}[]{@{}
  >{\raggedleft\arraybackslash}p{(\columnwidth - 10\tabcolsep) * \real{0.1739}}
  >{\raggedleft\arraybackslash}p{(\columnwidth - 10\tabcolsep) * \real{0.1739}}
  >{\raggedleft\arraybackslash}p{(\columnwidth - 10\tabcolsep) * \real{0.1739}}
  >{\raggedleft\arraybackslash}p{(\columnwidth - 10\tabcolsep) * \real{0.1739}}
  >{\raggedleft\arraybackslash}p{(\columnwidth - 10\tabcolsep) * \real{0.1739}}
  >{\raggedright\arraybackslash}p{(\columnwidth - 10\tabcolsep) * \real{0.1304}}@{}}
\toprule\noalign{}
\begin{minipage}[b]{\linewidth}\raggedleft
k
\end{minipage} & \begin{minipage}[b]{\linewidth}\raggedleft
Moment degree
\end{minipage} & \begin{minipage}[b]{\linewidth}\raggedleft
b\_(k+3)
\end{minipage} & \begin{minipage}[b]{\linewidth}\raggedleft
Full positive factor ceiling
\end{minipage} & \begin{minipage}[b]{\linewidth}\raggedleft
Complete exponent
\end{minipage} & \begin{minipage}[b]{\linewidth}\raggedright
Scope
\end{minipage} \\
\midrule\noalign{}
\endhead
\bottomrule\noalign{}
\endlastfoot
141 & 144 & −110835 & −581271 & −3344299/1024 & Complete target sector \\
142 & 145 & −111401 & −579654 & −1688491/1024 & Complete target sector \\
143 & 146 & −111962 & −578033 & −28587/1024 & Failed230-bit sufficient envelope only \\
\end{longtable}

For clarity, the full factor ceiling is \[
 \left\lceil\log_2\!\left[
 2\binom{164}k\frac{\pi_+^{(k+3)f}}{D_0^{k+3}}
 q_-^{-335f/2}
 \left(\frac{q_+}{(1-\theta)R}\right)^{(k+3)f}
 \right]\right\rceil_{\rm strict}+b_{k+3}.
 \tag{BC26}
\] The same arithmetic retains every correction for the failed next cell, whose maximum relative exponent is−366. A negative upper estimate of about−27.9 bits is insufficient for the requested error ledger; it gives no k143 removal.

The first two complete bounds are respectively below2\^{}-3265 and2\^{}-1648. Paying the whole K118 tail once gives the exact dyadic union \[
            2^{-3265}+2^{-1648}+2^{-8605114}<2^{-1647}.
 \tag{BC27}
\] There is no extra multiplicity for two endpoint branches or separate witness choices. The actual chronological prefix, all132 baseline all-top charges and all31 uniform charges are retained. The executable accepts a later complete prefix without changing the mathematical subjects, rejects a duplicate1647 and checks both full sums below2\^{}-230.

\subsection{8. Verification, source and conclusion boundary}\label{verification-source-and-conclusion-boundary}

The final source packet fingerprints R099's same-kernel AH4/AH14--24 input and its actual metric sources, the K119 positive rational-integral moment theorem and exact output table, K118's complete probability/arithmetic tail, and DO122's entire original central geometry with its mandatory non-strict-cutoff addendum. The continuous inequalities are proved above and in those stated inputs; finite models supplement them without claiming to enumerate global arrows.

Additional literal interfaces are DO098 CF1--16, DO104 CD10--19 and section5, K113 SC27.1--10, DO105 EL9--16, DO108 AC8--10, K107 SV17--19, K109's actual transverse metric, K110 PM10--12 and K111's all-minimum/all-positive-parameter scalar theorem. The ordinary Gamma/Beta normalizations are NIST DLMF5.9.1 and5.12.1/3. The selected norm/metric source remains Nakahama(2016), equation5.7 and Theorem5.1, printed960--961, through the original checked K043/K046/DO080 dictionary. No scalar representation is substituted.

The same projective modules, involution, Hermitian form, polarization, cubical theta line and finite projectors are retained. Carried FGA/SGA context is unchanged; no fresh reading or new formal-group, K-theoretic, l-adic-weight or CM-motive theorem is inferred from it. There is no full support28, actual selected arithmetic-image nonvanishing, prescribed period, source admission, specific032 or evaluated DP44/RH conclusion. The only proposed main-gate change is the complete two-cell union BC2, after all independent and root reviews pass.

\subsection{Canonical algebraic-geometric context}\label{canonical-algebraic-geometric-context-75}

The physical coordinate change is a real analytic treatment of the same original central affine module, preserving its shift, CM partner conditions and polarization. It adds no arithmetic parameter or field and changes no selected metric. NIST DLMF5.9.1,5.12.1/3 and5.4.6 are applied with explicit positive integral arguments. Inherited context only: SGA3 I4.7.3/4.7.3.1; SGA6 VI Theorem3.3; FGA232 Theorem3.1 pp148-149; FGA236 Corollaries1.7/2.3 pp229/232; SGA7 I VII2.7/2.9 pp169-174 and VIII7.4(b) pp310-311; cote4 scan003 sections8.1.1-8.1.2 and scans005-006 sections8.4.1-8.4.3; cote54 scans005-007. No new formal-group/Dieudonne, lambda-ring/K-theory, l-adic-weight or CM-motive comparison. DP\_MULT Q3e/DP44 and THETA3 G1/A7 receive the strictly smaller conditional finite gate; CX\_GRH CX-ALG-024 and GCT\_PLETH T-052 retain their arithmetic/global gates.

\clearpage

\section{AM-K-120 --- a narrower actual selected band and its positive shifted moments}\label{am-k-120-a-narrower-actual-selected-band-and-its-positive-shifted-moments}

\textbf{Canonical checked integration --- 2026-11-23T11:08:21.210996+00:00.} Complete saved-proof and executable reviews and separate exact diagnostic replays passed. No numerical error charge is introduced. Frozen author history is preserved. AM-K-120 NB1-24 proves the actual all-word/twist mean2111/64\textless M1\textless2113/64 and full probability outside I\emph{={[}16/529,16/527{]} below2\^{}-40000000, using h=2\^{}35 in the true negative-moment chamber. Both event complements and all168 embeddings, multiplied by the complete original bounded mass2\^{}11394877, give the full additive arithmetic tail2\^{}-8605114. Separately the shifted shape166 lambda32/25 profile is dominated by the positive Gamma166/rate4, Gamma512/rate400 and half-Gamma convolution. All167 exact rational inverse moments pi}r\_j are evaluated by two independent rational integral methods; the real chamber is1/2\textless j\textless678.5. Infinite mass and the original-row first-moment repair for the sole zero degree are retained, with every scalar/transverse polynomial term and true primal covolume. This theorem supplies neither the same-kernel angular implication nor the complete sector count; all intercept, central, finite and high-degree pattern costs remain separately payable. No standalone charge or main-gate change; current133/31 arrays and complete floors27/11 unchanged.

Author: \nolinkurl{/root/kummer_extension}. 23 November 2026 UTC. Controlled pending theorem. Whole author, independent and root checks are separate events. This proof changes no canonical file, ordered error array or sector charge. K118 and K119 remain frozen historical subjects.

\subsection{1. Actual probability, the new band and exact statement}\label{actual-probability-the-new-band-and-exact-statement}

Use the original J56 all-top upper endpoint, every admitted binary word and canonical twist, the genuine selected vector-valued Hilbert metric and the actual compact-invariant probability density of DO080 RT4--8. Put \[
 Y=1-|z|^2,\quad Z=Y^{-1},\quad M_1=\mathbb E Z,
 \quad a=f=168,\quad B=2^{955},\quad A=2^{960},
\] \[
 I_*=[16/529,16/527],\qquad h=2^{35}.                 \tag{NB1}
\] The first conclusion is the uniform actual mean and full-probability bound \[
 \frac{2111}{64}<M_1<\frac{2113}{64},\qquad
                 \Pr(Y\notin I_*)<2^{-40,000,000}.    \tag{NB2}
\] The probability in NB2 is not conditioned on another event. For the original bounded all-II globally rank-two-unipotent scalar-cross support27 family, with nonzero original S-cross and \(x=N_ED_\eta<2^{3328}\), replacing all full fidelities by their simultaneous \(I_*\)-restricted fidelities has total absolute normalized cost less than \[
                            2^{-8,605,114}.             \tag{NB3}
\] This is the same complete containing family as K118 BT5. Its entire arithmetic mass is retained. No restricted main-term estimate or numerical sector removal is asserted by NB3.

The second conclusion is a separate reusable positive-moment statement. Define \[
 \lambda_*=32/25,\qquad
 \Phi_*(t)=t^{-1/2}(1+t/4)^{-166}e^{-\lambda_*t},
 \quad t>0.                                             \tag{NB4}
\] There is an explicit positive rational majorant of NB4 whose first167 inverse moments are exactly \(\pi r_j^*\), with \(r_j^*>0\) rational. The companion table stores every rational and a strict dyadic ceiling for \[
                      P_j^*=((22/7)r_j^*)^{168}.         \tag{NB5}
\] The actual selected trace is not asserted to obey NB4 by this theorem. A separate same-kernel angular proof must supply that implication, its intercept and near-point factors. The probability statement and the analytic moment statement are proved independently below; neither substitutes for that missing application step.

\subsection{2. The literal selected first moment is uniformly close to33}\label{the-literal-selected-first-moment-is-uniformly-close-to33}

Let \(\ell\) be the actual endpoint, \(2^{484}<\ell<2^{485}\), and let the original twist parameter be \(0\le q\le\ell\). DO080 RT22, with the corrected word convention, gives \[
 E=B+\ell+96,\quad
 r_i=E+(4a-3-4i)\ell+q+1\quad(0\le i<a-1),
 \quad r_{a-1}=2\ell+q+1,
\] \[
 m_0=A+83-668\ell,\quad
 m_i=4\ell-1-\ell\omega_i\ (1\le i<a-1),\quad
 m_{a-1}=12.                                            \tag{NB6}
\] Every \(m_i\ge0\), and \[
                 M_1=\prod_i\frac{r_i+m_i-1}{r_i-2}.   \tag{NB7}
\] Each factor increases with \(m_i\) and decreases with \(r_i\). Thus its maximum over all actual words and twists occurs at the zero word and \(q=0\). The exact telescoping product from RT22 is \[
 M_1\le M_+:=
 \frac{(33B+2\ell+179)(2\ell+12)}
      {(B+6\ell+95)(2\ell-1)}.
\] All other factors exceed1, so the first factor alone at \(q=\ell\) gives \[
 M_1>M_-:=\frac{33B+3\ell+179}{B+671\ell+95}.           \tag{NB8}
\] No monotonicity in a fictitious word-dependent \(r_i\) is used.

The lower comparison \(M_->2111/64\) is exactly \[
              B>1,416,289\ell+189,089.                 \tag{NB9}
\] It follows, for example, from \(\ell<2^{485}\) and \(B=2^{955}\). For the upper comparison, clearing the positive denominator yields \[
\begin{split}
 &2113(B+6\ell+95)(2\ell-1)
       -64(33B+2\ell+179)(2\ell+12)\\
 &=B(2\ell-27457)+25100\ell^2+364344\ell-338207>0.
\end{split}                                             \tag{NB10}
\] Both terms are positive at the actual \(\ell>2^{484}\). This proves the mean interval in NB2 by integer inequalities, uniformly in all actual words and twists.

\subsection{3. Both complete high-moment tails}\label{both-complete-high-moment-tails}

Keep the exact normalized positive and negative moments of K046 BL8--10 and K118 BT6--12. With \[
 \epsilon_i=\frac{m_i+1}{(r_i-3)(r_i+m_i-1)},\qquad
                    \sum_i\epsilon_i<B^{-1},
\] their ordinary convergence chamber and exact factor inequalities give, for the present \(h\ge2\), \[
 M_1^h\mathbb E Y^h
   <\exp\!\left(\frac{h(h+1)}{2B}\right),\qquad
 M_1^{-h}\mathbb E Z^h
   <\exp\!\left(\frac{h(h-1)}B\right).                  \tag{NB11}
\] The actual minimum satisfies \(\min_i(r_i-2)>2h\), since \(2\ell-1>2^{36}\). In particular the negative moment is strictly within \(h<\min_i(r_i-1)\); no analytic continuation or inference from just a second moment occurs. Both correction exponents in NB11 are less than1.

For completeness, the negative factor used in NB11 is exactly \[
 \frac{(r_i-2)(r_i+m_i-1-j)}
      {(r_i+m_i-1)(r_i-2-j)}
 =1+j\epsilon_i\frac{r_i-3}{r_i-2-j}
 \le1+2j\epsilon_i\quad(0\le j<h).
\] The positive normalized factor is at most \(1+(j+1)\epsilon_i\). Products and \(\log(1+u)\le u\) give NB11 with every factor included.

Markov's inequality and NB2's mean interval now give \[
 \Pr(Y>16/527)
 <\left(\frac{2108}{2111}\right)^h
       \exp\!\left(\frac{h(h+1)}{2B}\right),
\] \[
 \Pr(Y<16/529)
 <\left(\frac{2113}{2116}\right)^h
       \exp\!\left(\frac{h(h-1)}B\right).               \tag{NB12}
\] Both ratios are at most \(1-3/2116\). Their full sum is therefore less than \[
                  2\exp(1-3h/2116).
\] The exact inequalities \(e>65/24>27/10\) and \((27/10)^7>2^{10}\) give, since the exponent is negative, \[
 \log_2\Pr(Y\notin I_*)
 <\frac{17}{7}-\frac{30h}{7\cdot2116}
 =-\frac{257698028767}{3703}<-40,000,000.                \tag{NB13}
\] This proves the complete probability statement. It is a symbolic high-moment argument, not numerical multiplication of \(h\) factors or a Gaussian approximation to the selected density.

\subsection{4. Exact fidelity comparison and the complete original arithmetic payment}\label{exact-fidelity-comparison-and-the-complete-original-arithmetic-payment}

Let \(F_{I_*}(g)\) denote DO080's fidelity integral restricted to the two conditions \(Y(z),Y(g^{-1}z)\in I_*\). Invariance of the original measure and Cauchy--Schwarz, applied separately to each complementary event, give \[
 0\le F_{I_*}(g)\le F(g)\le1,\quad
 |c_W(g)|\le F(g)\le F_{I_*}(g)+2\,2^{-20,000,000}.     \tag{NB14}
\] This is uniform in the real argument, both compact factors and the genuine near point, including radius0. The product of168 fidelities telescopes with loss at most their sum, giving \[
 |c(h_0)|\le\prod_\tau F_{I_*,\tau}(h_{p,\tau})
                         +2^{9-20,000,000}.             \tag{NB15}
\] The notation \(h_0\) distinguishes the original arithmetic matrix from the moment order.

The same unchanged arithmetic containing family has complete nonarchimedean normalized absolute mass less than \(2^{11,394,877}\), by K118 BT16--20. This count is independent of the chosen band. It retains the unique actual slope per27-place support below \(x=2^{3328}\); every167 original row value in the containing ideal; the full real168-dimensional affine central fibre; every one of the original \(4^{334}\) common patterns; all support and retained pivot choices; the complete active weighted-label sum; both mixed contractions and the original class/Haar normalization. It does not require the uncharged restricted main term to be small.

Multiplying the full mass by NB15 gives, for any actual subset \(\mathcal S\) of this original bounded family, \[
 \sum_{h_0\in\mathcal S}|w(h_0)c(h_0)|
 \le\sum_{h_0\in\mathcal S}|w(h_0)|
        \prod_\tau F_{I_*,\tau}(h_{p,\tau})
                      +2^{-8,605,114}.                 \tag{NB16}
\] The new band is a different restriction. NB16 is proved directly for it; it is not inferred by inclusion from the wider K118 band. The same numerical tail size is paid once in a subsequent complete union. No extra charge is appended here, and no unbounded-denominator or different-support count is inferred.

\subsection{5. A separate exact positive measure for the stronger exponential shift}\label{a-separate-exact-positive-measure-for-the-stronger-exponential-shift}

For \(t>0\), the inequality \(\log(1+u)\le u\) gives \[
 \Phi_*(t)\le G_*(t):=
 t^{-1/2}(1+t/4)^{-166}(1+t/400)^{-512}.                \tag{NB17}
\] It holds globally, with no cutoff loss, because \(512/400=32/25\). Let \[
 \nu_*=\gamma_{166,4}*\gamma_{512,400}*
                   [u^{-1/2}du/\sqrt\pi],              \tag{NB18}
\] where both Gamma factors are probability normalized by shape and rate. Its Laplace transform is exactly \(G_*\). Its total mass is infinite. Its inverse moments converge precisely for \[
                          1/2<j<678+1/2.               \tag{NB19}
\] Indeed Tonelli and the ordinary Gamma integral give \[
 L_j^*=\int c^{-j}\nu_*(dc)
 =\frac1{\Gamma(j)}\int_0^\infty
 t^{j-3/2}(1+t/4)^{-166}(1+t/400)^{-512}dt              \tag{NB20}
\] for \(j>0\). Its endpoint powers prove NB19; local finiteness and infinite tail mass prove divergence for \(j\le0\). The exact shifted measure, translated by \(\lambda_*\), has finite moments for every real \(j>1/2\), and its moment is at most NB20. The zeroth inverse moment remains infinite for both measures.

Here is the full exact rational evaluation, specialized rather than numerically interpolated from K119. Set \(p=166,N=512,A_0=4,B_0=400,\Delta=396\), and let \(\mathcal I(u,v)\) be the positive integral \[
 \mathcal I(u,v)=\pi^{-1}\int_0^\infty
              (s^2+A_0)^{-u}(s^2+B_0)^{-v}ds.
\] The K119 RM10--13 finite partial-fraction identity, valid for all positive integers \(u,v\), evaluates this as a rational number using the one-pole coefficients \[
 J_\ell(r)=\frac{\binom{2\ell-2}{\ell-1}}
                    {2^{2\ell-1}r^{2\ell-1}},\quad r=2,20.
\] With \(C=2A_0^pB_0^N\), the first two normalized moments are \[
 r_1^*=C\mathcal I(p,N),\qquad
 r_2^*=C[\mathcal I(p-1,N)-A_0\mathcal I(p,N)].          \tag{NB21}
\] For \(1\le j\le165\) define \[
 D_j=2p+2N-2j-3>0,\qquad
 B_j=(2j-1)(A_0+B_0)-2(p-1)B_0-2(N-1)A_0,
\] \[
 r_{j+2}^*=
 \frac{jB_jr_{j+1}^*+(2j-1)A_0B_0r_j^*}
      {j(j+1)D_j}.                                     \tag{NB22}
\] Exactly as in K119 RM9/RM15, the substitution \(t=s^2\) and the derivative of \(s^{2j-1}/[(s^2+A_0)^{p-1}(s^2+B_0)^{N-1}]\) prove that \(L_j^*=\pi r_j^*\). Both boundary values vanish in every stated use. Each rational is positive by its integral identity.

The complete table is independently checked using \(\mathcal I(0,v)=J_v(20)\), \(\mathcal I(u,0)=J_u(2)\) and \[
 \mathcal I(u,v)=
 [\mathcal I(u,v-1)-\mathcal I(u-1,v)]/\Delta.
\] A rectangular166 by512 array, followed by \[
 r_j^*=\frac C{(j-1)!}\sum_{r=0}^{j-1}
 \binom{j-1}{r}(-A_0)^{j-1-r}\mathcal I(p-r,N),         \tag{NB23}
\] checks every one of the167 values without using the partial-fraction bases or the moment recurrence. All terms are convergent ordinary integrals, including \(p-r=0\) at degree167. No floating special function is used. The certificate also checks \(r_j^*<r_j\) for the K119 Gamma444 majorant; analytically this is immediate from its strictly smaller positive integrand. This comparison does not replace any required external intercept or tail cost.

\subsection{6. Infinite mass, every polynomial degree and source boundaries}\label{infinite-mass-every-polynomial-degree-and-source-boundaries}

For the original nonzero rows, retain the exact product floor \(\prod_\tau\alpha_\tau\ge(Rk/q)^f\) and reserve \(0<\theta<1\) before enlargement. The AM--GM argument K119 RM17--19 gives the whole integrated degree-zero bound \[
                  P_1^*(q/(\theta Rk))^f.              \tag{NB24}
\] Every positive polynomial degree uses its finite \(P_j^*\). The entire original row polynomial is exactly K119 RM19 with \(P_j\) replaced by \(P_j^*\): all scalar terms, all four transverse terms, original primal covolume \(D_0^{k+3}z^{2(k+3)f}/U^f\), real dimension336 and actual saturated Gram remain. The sole zero-degree term is still \((l,h_{\rm tr})=(0,3)\); no moment of infinite mass is assigned a finite value. At \(k=164\), the degree167 finite-pattern cost must still be supplied by a separate complete count.

The new results are NB2/NB16's actual probability and full containing arithmetic tail, and NB17--24's independent positive moment table. A new selected-kernel angular theorem and a new complete arithmetic count remain distinct requirements. Their external exponential coefficient, actual near-point transport, physical central-coset scaling/cover, original finite factors and high-degree pattern corrections cannot be omitted or multiplied from unrelated bounds. No numerical error charge, new complete support floor, signed nonvanishing, admitted source, selected arithmetic image, prescribed period or evaluated DP44 conclusion is asserted here.

Direct inherited interfaces are DO080 RT4--8/RT20--24; K046 BL8--10's actual convergent moments; K118 BT6--22's factor, fidelity and complete containing-mass arguments; K119 RM1--19's ordinary positive integral identities and exact row-polynomial interface; K109 TG3--12, K110 PM10--12, K116 RE11--20 and R098 WJ18--28 at their unchanged source scopes. Primary normalization: National Institute of Standards and Technology. (n.d.). \emph{Digital Library of Mathematical Functions}, equations5.9.1 and5.12.1--5.12.3, \url{https://dlmf.nist.gov/5.9} and \url{https://dlmf.nist.gov/5.12}. Genuine selected metric and norm normalization remain Nakahama (2016), equation5.7 and Theorem5.1, printed960--961, through the accepted K043/K046 dictionary. No theorem-validity audit or Lean build of the original archive is performed.

The original fields, involution, projective modules, polarization, cubical theta line, Picard/Poincare pairing, actual selected vector and finite projectors are unchanged. Inherited contextual anchors remain SGA3 I4.7.3/4.7.3.1; SGA6 VI Theorem3.3; FGA232 Theorem3.1 pp148--149; FGA236 Corollaries1.7/2.3 pp229/232; SGA7 I VII2.7/2.9 pp169--174 and VIII7.4(b) pp310--311; cote4 scan003 sections8.1.1--8.1.2/scans005--006 sections8.4.1--8.4.3 and cote54 scans005--007. These are carried context, not new source constructions or arithmetic comparison theorems.

\subsection{Canonical algebraic-geometric context}\label{canonical-algebraic-geometric-context-76}

This is an estimate at the archimedean realization of the same projective OE modules and polarized CM abelian data. The positive Gamma convolution and exact inverse-moment evaluation do not alter the scheme over its original arithmetic base, the involution, Hermitian lattice, theta line, Picard scheme or Poincare biextension. Carried contextual anchors are SGA3 I4.7.3/4.7.3.1; SGA6 VI Theorem3.3; FGA232 Theorem3.1 pp148--149; FGA236 Corollaries1.7/2.3 pp229/232; SGA7 I VII2.7/2.9 pp169--174 and VIII7.4(b) pp310--311; cote4 scan003 sections8.1.1--8.1.2 and scans005--006 sections8.4.1--8.4.3; cote54 scans005--007. These locators are inherited context, not new primary readings or newly proved comparisons. No new formal-group/Dieudonne, lambda-ring/K-theory, l-adic weight, CM-motive or motivic L-function result follows from this real analytic inequality alone. DP\_MULT Q3e/DP44, THETA3 G1/A7, CX\_GRH CX-ALG-024 and GCT\_PLETH T-052 retain their current arithmetic and global gates until a separately checked composition changes them. No Lean build or archive theorem-validity audit is performed.

\clearpage

\section{AM-LIFT-101 --- All-ring harmonic theta values and the literal connection constructor}\label{am-lift-101-all-ring-harmonic-theta-values-and-the-literal-connection-constructor}

\textbf{Canonical checked integration --- 2026-11-23T11:17:37.956514+00:00.} Complete saved-proof and executable reviews and separate exact diagnostic replays passed. No numerical error charge is introduced. Frozen author history is preserved. AM-LIFT-101 AT1-25 proves two precise conditional constructors. Given actual Host, Datum, odd Level, Coeff and Moduli records, every literal IsNab yields all-ring ThValTot, actual FrMap naturality and harmonic independence. With the accepted CanonGaugeSt theorem, those same records permit a choice of NabV satisfying IsNab, including NabZero at all integer powers. The proof supplies coherent formal charts and frames over arbitrary coefficient algebras, global primitive/adjunction normalization, finite-order corrected additive frames, ordinary derivatives, exact harmonic cancellation and faithfully-flat descent of the original odd ideal-half-level theta sections. Where the operator is free to be chosen, independent nabPin, nabUniq and thTot premises are removed. Combined with AM-LIFT-100 and an actual FrPt over a finite number-field extension, unchanged nonzero harmonic P and DF+2\textless=M, the literal ThNe0 law follows without separate chosen-IsNab/ThValTot assumptions. This does not identify a separately prescribed operator, construct the actual arithmetic source or Host/Sizes, prove the other source laws, or establish nonvanishing of the original coupled selected image. No error charge; complete133/31 arrays and floors27/11 unchanged.

\textbf{Status:} complete author-reviewed mathematical text; final frozen-hash diagnostic replay and its author record are supplied alongside this file. Independent review and canonical adoption are separate. This constructs named source laws conditional on the actual geometric records; it does not construct the original arithmetic source, complete Host/Sizes package, or coupled selected image. All author files remain in the pending directory.

\subsection{1. Statement and quantifiers}\label{statement-and-quantifiers}

Retain an actual geometric \nolinkurl{Host E}, \nolinkurl{Datum H}, odd \nolinkurl{Level H}, \nolinkurl{Coeff l}, and \nolinkurl{Moduli Dw l}. Write \(F=c.F\), \(M=l.M\), \(\mathfrak b=l.Mb\), \(r=Dw.p+Dw.q\), and \(g=[E:\mathbf Q]r/2\). The universal group \(A/X_o\), its full \(\mathcal O_E\)-action, rigidified symmetric line \(L\), cubical structure, global adjunction isomorphisms, primitive isomorphism, and actual half-level \(\eta\) are those in the given source records. In particular

\[
 M\ge3\text{ is odd},\qquad
 \mathfrak b\overline{\mathfrak b}=M\mathcal O_E,\qquad
 \mathfrak b+\overline{\mathfrak b}=\mathcal O_E. \tag{AT1}
\]

Use the CM involution on \(\mathcal O_E\). Its equality with \nolinkurl{H.bar} is exactly the first clause of \nolinkurl{IsNab}. We use the accepted formal statement \nolinkurl{CanonGaugeSt}, supplied by \nolinkurl{CanonB.canonGauge}, with its full commutative-\(\mathbf Q\)-algebra scope; no assertion about a build or the validity of the source archive is made. The accepted \nolinkurl{CanonGaugeIndependenceAll} statements provide the corresponding formal chart/gauge identities. Below we prove the geometric applicability and the finite-order comparison needed by the literal records.

\textbf{Theorem A --- the law from a pinned connection.} For every actual \nolinkurl{nab : NabV Xo c} satisfying the literal \nolinkurl{IsNab Xo c nab}, the literal \nolinkurl{ThValTot Xo c nab} holds. More precisely, take any commutative \(F\)-algebra \(S\), any \(y\in\operatorname{FrPt}(X_o,c;S)\), any row degree \(a:(E\hookrightarrow\mathbf C)\times\mathbf N\to\mathbf N\), any \(n\in\mathbf N\), any \(x:\operatorname{Fin}(n)\to Dw.La\), and any \(P\in\operatorname{harmA}(c,a,n)\). There is a unique

\[
 v_{y,x,P}\in S\quad\text{such that}\quad
 \operatorname{ThValAt}(X_o,c,\mathrm{nab},y,x,P,v_{y,x,P}). \tag{AT2}
\]

Here \(S\) may have nilpotents, may be disconnected, and need not be finitely generated or Noetherian. For every actual \nolinkurl{FrMap} along an \(F\)-algebra map \(S\to S'\), the values commute with that map. Two pinned connections give the same values for the same original data. The faithfully flat witness in \nolinkurl{ThValAt} can be taken finite étale of rank \(M^{gn}\); for \(n=0\) it is the identity map. This rank is the simultaneous half-level-root torsor rank, not a bound for an arbitrary global arithmetic field.

\textbf{Theorem B --- existence of a pinned connection.} Under these same actual geometric records and the accepted \nolinkurl{CanonGaugeSt}, one can choose a \nolinkurl{NabV Xo c} satisfying the literal \nolinkurl{IsNab Xo c nab}. The construction is over all the rings quantified in \nolinkurl{IsNab}, and uses global normalized primitive/adjunction generators on the universal group. Its per-point auxiliary chart/frame choices need not be natural. Theorem A proves naturality of the resulting harmonic theta expression.

Consequently \nolinkurl{NabUniq Xo c} holds as well. For a family of already supplied actual stage records, these constructions provide the connection entries and the named \nolinkurl{nabPin}, \nolinkurl{nabUniq}, and \nolinkurl{thTot} clauses. They do not supply the other clauses of \nolinkurl{LawsP}, \nolinkurl{LawsSh}, or \nolinkurl{TopExists}. Theorem A applies to an already prescribed operator only when that operator satisfies its original \nolinkurl{IsNab}; Theorem B chooses an operator where such a choice is permitted.

\subsection{2. Properness, global normalizations, and the actual tangent frame}\label{properness-global-normalizations-and-the-actual-tangent-frame}

The universal \nolinkurl{SAB.semi} supplies smoothness and geometrically connected fibres, while \nolinkurl{Obj.proper} supplies properness. Thus \(A/X_o\) is an abelian scheme; this is also the exact source lemma \nolinkurl{PrimUniqPf.isAbelianScheme_univ}. Every power \(A^j/X_o\) has the same properties. For such a morphism,

\[
 (\pi_j)_*\mathcal O_{A^j}=\mathcal O_{X_o},
 \quad\text{universally under base change}. \tag{AT3}
\]

The Stacks Project, Lemma 36.32.6, proves (AT3) for proper flat finitely presented morphisms with geometrically reduced connected fibres. In particular any automorphism of an invertible sheaf on \(A^j\) is multiplication by a base unit. Two global generators of the same invertible difference sheaf which agree along the zero section are identical. This holds over arbitrary base changes, including nonreduced ones.

Apply this to \nolinkurl{SAB.prim}. An existing global generator of \(L\otimes P(X,[w_0]X)^{-1}\), where \(P\) denotes the cubical two-variable line, can be multiplied by the inverse of its value at zero. There is therefore a unique global \nolinkurl{wP} satisfying the literal \nolinkurl{IsPrim}, including the specified rigidification. It is global on the universal \(A\), as required in \nolinkurl{ChartAt}, not just on the selected fibre.

Likewise each \nolinkurl{SAB.adj a} has a unique normalized global generator on \(A^2\). Normalization at \((0,0)\) gives the canonical normalization on both axes: its restriction to either axis is a unit on \(A\), hence comes from the base by (AT3). All addition, multiplication, symmetry, and biextension compatibility diagrams for these normalized maps commute. Their two sides are global isomorphisms of the same line on a power of \(A\); their ratio is a base unit and is one at the origin. This establishes the compatibilities simultaneously for every \(a\in\mathcal O_E\), without adjoining or choosing infinitely many extra scalars.

For a given actual frame point \(y\) over \(S\), its eigenframe is a full tangent frame. Indeed \nolinkurl{Coeff.conj} and separability give

\[
 E\otimes_{\mathbf Q}F\simeq\prod_{\tau:E\hookrightarrow\mathbf C} F,
 \qquad
 \operatorname{Lie}(A_S/S)=\bigoplus_\tau\operatorname{Lie}(A_S/S)_\tau. \tag{AT4}
\]

After tensoring with any \(F\)-algebra the idempotents remain orthogonal and sum to one. The source \nolinkurl{FrPt.basis} is a basis on each summand, so the vectors indexed by \nolinkurl{Lt1 Dw} form a basis of the entire Lie module. This argument does not require reducedness of \(S\), a pure CM type, or a separately chosen minimal real representation. Nonzero elements of \(F\) remain units in every nonzero \(F\)-algebra.

The Host relations imply

\[
 w_0+\overline{w_0}=1. \tag{AT5}
\]

To see this in the actual source, every complex embedding sends \nolinkurl{rt} to a nonzero square root of the negative rational number \(-d\); conjugation negates it. Add \(2w_0=1+\mathrm{rt}\) to its conjugate and divide by two. The same observation proves that \nolinkurl{InPsi} contains exactly one embedding in each conjugate pair.

\subsection{3. The complete formal chart and frame construction}\label{the-complete-formal-chart-and-frame-construction}

Fix a commutative \(F\)-algebra \(S\) and \(y\in\operatorname{FrPt}(S)\). Put \(d=g=|\mathrm{Lt1}(Dw)|\), using the fixed finite indexing equivalence. Smoothness constructs compatible formal coordinates at the identity with the tangent basis (AT4). Here is an all-affine-base argument rather than a field-only chart assertion.

At order one, the specified basis gives the morphism from the first infinitesimal neighbourhood of zero in \(\mathbf A^d_S\). Each subsequent truncated polynomial algebra surjects onto the preceding one with square-zero kernel. Smoothness permits a lift of the morphism on these affine thickenings. The global affine lifting statement, including the vanishing of the quasi-coherent obstruction torsor, is the Stacks Project, Definition 37.11.1 and Lemma 37.11.7. Each lift preserves the previous one. Its derivative is unchanged. Locally for a smooth chart the associated graded map is the symmetric algebra of the cotangent isomorphism; induction gives an isomorphism at every finite order. These local statements glue because they concern the same morphism. Thus the resulting compatible sequence is a formal-coordinate isomorphism.

Pull group multiplication and each actual \(\mathcal O_E\)-action through these coordinates. They define a commutative formal group law \(G\) and an action \(A_a\), with every \nolinkurl{IsChart.zero}, \nolinkurl{trunc}, \nolinkurl{law}, \nolinkurl{act}, \nolinkurl{tan}, and \nolinkurl{basis} clause. This construction works for arbitrary affine \(S\). No use of the source \nolinkurl{Moduli.rep} is made here; that field has a reduced-scheme hypothesis, whereas our point \(y\) is already given.

The two-variable cubical line is trivial on each infinitesimal coordinate scheme. Explicitly, let \(R_m=S[Z]/(Z)^{m+1}\), and let \(V_m\) be its corresponding invertible module. Its rigidified reduction modulo the nilpotent augmentation ideal has a generator. A lift of that generator generates \(V_m\): the map \(R_m\to V_m\) is an isomorphism modulo the nilpotent ideal, its cokernel vanishes by nilpotent Nakayama, and its kernel vanishes after splitting the surjection to the projective module \(V_m\). Generators can be lifted compatibly through all \(m\).

Normalize these biframe generators on both axes. If their axis ratios to the canonical generators are \(u(X)\) and \(v(Y)\), multiply by \(u(X)^{-1}v(Y)^{-1}u(0)\). The two ratios agree at the origin, and the formula makes both restrictions canonical. These adjustments respect truncation. We have obtained actual coherent generators \(e_m\) of the exact line \nolinkurl{FrTy}, with both axis normalizations.

In these generators, the given cubical maps, symmetry, and the normalized global adjunctions of §2 are units. Their compatible truncations define \(\phi_1,\phi_2,s,\operatorname{ad}_a\) as formal series. The cube's face identities give the six zero-variable conditions; its symmetry and cocycle identities give left/right associativity, commutativity and interchange. The symmetry square and its compatibility with the cube give \nolinkurl{sy_inv} and \nolinkurl{sy_law}. The normalized global adjunction diagrams give \nolinkurl{adjoint_compat_left/right}, \nolinkurl{ad_one}, \nolinkurl{ad_mul}, \nolinkurl{ad_add}, and \nolinkurl{ad_sy}. All diagrams are the same original maps, expressed in these generators; the normalization argument (AT3) fixes their scalars. Thus these series and \(G,A_a\) form the full literal \nolinkurl{FB}, and the defining equalities of \nolinkurl{IsFrame}, including the global \nolinkurl{wL} for every \(a\), hold.

Apply \nolinkurl{CanonB.canonGauge} with \(R=k=S\), \(O=\mathcal O_E\). This is legitimate: \(\mathcal O_E\) is finite free over \(\mathbf Z\), reduced, and has its actual CM involution, while \(S\) is a \(\mathbf Q\)-algebra. Obtain \(h\) satisfying every clause of \nolinkurl{IsCanonGauge}. The source definition has no field or reducedness requirement on \(k\).

With the unique global \nolinkurl{wP} of §2, the line generators are exactly

\[
 e_{L,m}=w_P\,e_m(X,A_{w_0}X). \tag{AT6}
\]

Their value at the origin is the original rigidification, supplying \nolinkurl{ChartAt.rig0}. Every section of \(L^k\), \(k\ge0\), has a unique compatible power-series germ \(f\) in the generators \(e_{L,m}^k\). Equality of all finite truncations gives uniqueness. This proves the entire germ-existence clause of \nolinkurl{IsNab}, using one chart/frame/gauge for all \(k,F\) at this \(S,y\).

Choose these auxiliary data for each \(S,y\), and define the nonnegative part of \nolinkurl{NabV} by the exact source formula

\[
 \mathrm{nab}_y(k,F,\alpha)
 =\left.D_G^{\alpha}\!\left(f\exp\bigl(kh(X,A_{w_0}X)\bigr)\right)\right|_{X=0}.
 \tag{AT7}
\]

The vectors in \(D_G^\alpha\) are the prescribed physical eigenframe through \nolinkurl{uvec}; the order and ordinary derivative normalization are those of \nolinkurl{DmonOf}. For negative integer \(k\), which \nolinkurl{IsNab} does not constrain, set the coefficient at \(\alpha=0\) to the rigidified value of \(F\) at zero, and the other coefficients to zero. The nonnegative construction supplies every clause of \nolinkurl{IsNab}, thereby proving Theorem B. The value at \(\alpha=0\) in (AT7) is also the rigidified value because \(h\) vanishes on the axes. Thus this particular extension additionally satisfies \nolinkurl{NabZero} for all integers \(k\).

The existence just established is conditional on the actual geometric records and the accepted formal gauge theorem. It is not conditional on an unproved existence of a full formal chart over arbitrary \(S\): the affine lifting and invertible-module argument above supplies that step.

\subsection{4. Arbitrary pinned witnesses: do not assume biframe truncation compatibility}\label{arbitrary-pinned-witnesses-do-not-assume-biframe-truncation-compatibility}

For Theorem A we must compare any witnesses admitted by the literal \nolinkurl{IsNab}, not only the coherent ones constructed above. An actual \nolinkurl{ChartAt} has a compatible \nolinkurl{x} through \nolinkurl{IsChart.trunc}, but it contains \textbf{no truncation-compatibility axiom for its generators \(e_m\)}. We consequently work at one finite order \(D\ge\deg P\). We do not infer a single formal biframe from that record.

First, its chart is a formal-coordinate isomorphism. The \nolinkurl{IsChart.tan/basis} clauses and (AT4) make its derivative invertible. Smoothness and the formal inverse construction then give the inverse at every order. Over a nonzero ring its dimension is \(g\); over the zero ring all scalar equalities below are automatic.

There is a canonical additive coordinate system on the formal commutative group once its physical tangent basis is fixed. For completeness, over a \(\mathbf Q\)-algebra the commuting invariant derivations \(D_1,\ldots,D_g\) give the Taylor map

\[
 f\longmapsto\sum_{\alpha\in\mathbf N^g}
       \frac{e^*(D^\alpha f)}{\alpha!}Z^\alpha. \tag{AT8}
\]

This map is defined on the sheaf of functions near the identity; Leibniz proves multiplicativity, and invariance proves that it intertwines the group law with addition. Its derivative is the prescribed isomorphism. Solving degree by degree gives its formal inverse. Equivalently, one may construct it in local smooth coordinates and glue: any homomorphism between additive formal groups over a \(\mathbf Q\)-algebra is linear, because its homogeneous terms of degree \(j\ne1\) are killed by the invertible rational number \(2^j-2\). Thus the homomorphism with this tangent map is unique. The construction and uniqueness commute with arbitrary maps of coefficient rings. This is the multivariable commutative \(\mathbf Q\)-theorem; (AT8) supplies the particular normalization used here.

In these additive coordinates, \(\mathcal O_E\) acts by its tangent action and the actual frame vectors are the coordinate vectors. The invariant derivatives become ordinary partial derivatives. This is also the content of \nolinkurl{TransPolySt}/\nolinkurl{NabTransSt}, with the derivative matrix applied to \nolinkurl{uvec}; no unrecorded second coordinate conjugation or factorial renormalization occurs.

At the chosen order \(D\), gauge-correct the actual biframe by

\[
 \widetilde e_D=e_D\exp(-h)\pmod{(X,Y)^{D+1}}. \tag{AT9}
\]

This sign is forced by the source equations. For example \nolinkurl{IsFrame.ph1} expresses the cubic generator as

\[
 \phi_1\,\frac{e_D(X+_G Y,Z)}{e_D(X,Z)e_D(Y,Z)},
 \qquad
 \phi_1=\exp\{h(X,Z)+h(Y,Z)-h(X+_G Y,Z)\}.
 \tag{AT10}
\]

Substituting (AT9) makes its new multiplier one. The second group law, symmetry and every adjunction similarly have multiplier one, by the literal gauge equations. The axes of \(e_D\) are already normalized: restrict \nolinkurl{IsFrame.ph1} or \nolinkurl{ph2} to a zero variable and use the cube's corresponding face and the FB zero-variable condition. Hence (AT9) is axis-normalized as well.

The global adjunction generator used by an arbitrary \nolinkurl{IsFrame.ad} agrees with the normalized one of §2 \textbf{after pullback to the selected point}. Indeed their ratio is a global base unit by (AT3); restricting the frame equation to its origin, using the normalized axes and \nolinkurl{adjoint_zero}, shows that this unit restricts to one in \(S\). This does not claim that a single point forces equality on the whole universal base. Equality on its pulled-back group, and consequently after every subsequent base change, is exactly what (AT11) requires. For \nolinkurl{IsPrim} the normalization condition is itself global, so any two such primitive generators already agree universally. The comparison therefore uses the original pulled-back structural maps with no residual scalar ambiguity.

\subsection{5. The finite-order balanced quadratic comparison}\label{the-finite-order-balanced-quadratic-comparison}

Take two arbitrary pinned witnesses for the same \(S,y\) and actual sections. In the common additive coordinates, the ratio of their corrected biframe generators at order \(D\) is a unit \(R_D(X,Y)\) with both axis restrictions one. Its logarithm \(b_D=\log R_D\) satisfies, modulo total degree \(D+1\),

\[
\begin{split}
 b_D(X+X',Y)&=b_D(X,Y)+b_D(X',Y),\\
 b_D(X,Y+Y')&=b_D(X,Y)+b_D(X,Y'),\\
 b_D(Y,X)&=b_D(X,Y),\\
 b_D(aX,Y)&=b_D(X,\overline aY)\qquad(a\in\mathcal O_E).
\end{split} \tag{AT11}
\]

All logarithms and exponentials here are finite polynomials in the nilpotent augmentation ideal. These identities follow by comparing the same multiplier-one maps in (AT9). The first two equations force \(b_D\) to be bilinear: axis vanishing excludes degree zero in either block, and substituting \(X'=X\) kills every \(X\)-degree other than one by \(2^j-2\in\mathbf Q^\times\); do the same in \(Y\). For \(D<2\), \(b_D=0\). The conclusion is valid over rings with nilpotents.

Write \(i\) for a vector in the \(\tau\)-eigenspace and \(j\) for a vector in the \(\tau'\)-eigenspace. The last identity in (AT11) gives

\[
 [\tau(a)-\tau'(\overline a)]\,b_D(i,j)=0\quad\text{for all }a\in\mathcal O_E.
 \tag{AT12}
\]

If \(\tau'\ne\overline\tau\), choose an integral element separating the embeddings. The bracket is a nonzero element of \(F\), hence a unit in \(S\). Thus \(b_D(i,j)=0\) except between conjugate eigenspaces. The diagonal comparison in the primitive line is

\[
 Q_D(Z)=b_D(Z,w_0Z)=\tfrac12 b_D(Z,Z), \tag{AT13}
\]

by symmetry, balance and (AT5). The coefficients of \(Q_D\) therefore involve only conjugate pairs. If \(\widetilde e'_D=\widetilde e_D\exp(b_D)\), the corrected germs of the same section of \(L^k\) satisfy

\[
 g'_j(Z)=g_j(Z)\exp(-kQ_D(Z))\pmod{(Z)^{D+1}}. \tag{AT14}
\]

For the \(n\) columns, set \(Q_{D,n}=\sum_{j=1}^nQ_D(Z_j)\). The equal power \(k\) in every column is essential. The literal harmonic condition in \nolinkurl{RowDeg.harmK}, with inverse matrix the identity, says

\[
 \sum_{j=1}^n
 \partial_{\tau,i,j}\partial_{\overline\tau,i',j}P=0
 \quad(\tau\in\mathrm{InPsi},\ i,i'). \tag{AT15}
\]

It follows that \(Q_{D,n}(\partial)P=0\). There are no cross-column terms. The canonical bijection from the source \nolinkurl{LtN} to \nolinkurl{Lt1 × Fin n} identifies this exactly with the column operator in \nolinkurl{GaugeIndependenceColumns}; it is not the false claim that each separate column is harmonic.

The ordinary derivative pairing obeys

\[
 \langle P,\exp(-kQ_{D,n})G\rangle
 =\langle\exp(-kQ_{D,n}(\partial))P,G\rangle
 =\langle P,G\rangle,
 \qquad \langle P,G\rangle=P(\partial)G(0). \tag{AT16}
\]

This is a finite identity since \(P\) has finite degree; it follows monomial by monomial from the Leibniz rule and ordinary factorials. It proves equality of the full harmonic column expression for the two witnesses. It does not assert equality of their individual connection coefficients. Nor is compatibility of \(e_D\) for different \(D\) needed: for each polynomial choose its one finite order, and carry out (AT9)--(AT16) at that order.

\subsection{\texorpdfstring{6. All-ring naturality and independence of \nolinkurl{nab}}{6. All-ring naturality and independence of }}\label{all-ring-naturality-and-independence-of-nab}

Every formal operation used at order \(D\)---coefficient extraction, the finite logarithm/exponential, multiplication, substitution, ordinary differentiation and the rigidified pullback of a section---commutes with a map of commutative \(F\)-algebras \(S\to S'\). The universal primitive and adjunction maps commute with it by their uniqueness. Push the chosen witness at \(y\) through a given \nolinkurl{FrMap} to \(y'\), and compare it with the witness selected by \nolinkurl{IsNab} at \(y'\) using §5. At this finite order its axes, group/action laws, gauge identities, germs and physical tangent vectors are the required ones.

Hence for any actual tuple \(F_j\in H^0(A_y,L^k)\), and its transported tuple \(F'_j\),

\[
 \phi\!\left(\sum_{\alpha}P_\alpha
        \prod_j\mathrm{nab}_y(k,F_j,\alpha_j)\right)
 =\sum_{\alpha}\phi(P_\alpha)
        \prod_j\mathrm{nab}_{y'}(k,F'_j,\alpha_j). \tag{AT17}
\]

Both \(y,y'\) here are actual \nolinkurl{FrPt} and the map is an actual \nolinkurl{FrMap}. We do not use \nolinkurl{Moduli.rep} on a nonreduced algebra, or assume that a point with all frames exists over every map. The normalized factor \(M^{-\deg P}\) lies in \(F\) and is preserved, so (AT17) is the base-change identity for the source \nolinkurl{thExpr}.

The same proof at the identity coefficient map compares any two \nolinkurl{IsNab} operators. It proves \nolinkurl{ThExprUniqSt} on harmonic inputs, and later identifies their \nolinkurl{ThValAt} relations and \nolinkurl{thGen} sets. The named \nolinkurl{NabUniq} is therefore a consequence of pointwise harmonic equality, stronger than merely assuming its set equality in advance.

\subsection{7. Canonical odd lifts over arbitrary bases}\label{canonical-odd-lifts-over-arbitrary-bases}

Set \(q=M^2\). Work over any affine base change of the actual abelian scheme. The normalized symmetric structure of \(L^q\) gives an involution \(\iota\) of its theta group which fixes the central \(\mathbf G_m\) and covers \(x\mapsto-x\). The involution is independent of rescaling the symmetry: the rescaling comes from a base unit by (AT3), and cancels in the source \nolinkurl{symGen} formula.

Let \(x\) be a \(q\)-torsion point. The theta-group torsor over \(x\) has a lift \(V\) locally for the faithfully flat topology. Write

\[
 V^q=\alpha,\qquad \iota(V)=\beta V^{-1}.
 \quad\text{Then}\quad\beta^q=\alpha^2. \tag{AT18}
\]

These are units in the local base ring. Define, with \(q\) odd,

\[
 U=\alpha^{-1}\beta^{(q-1)/2}V. \tag{AT19}
\]

Direct substitution gives \(U^q=1\) and \(\iota(U)=U^{-1}\). If \(V\) is replaced by \(uV\), then \(\alpha\) and \(\beta\) become \(u^q\alpha\) and \(u^2\beta\), and the prefactor in (AT19) changes by \(u^{-1}\). Therefore \(U\) is independent of the chosen local lift and descends. Any other such lift differs by a unit \(z\) satisfying \(z^q=z^2=1\); Bézout for \(q,2\) gives \(z=1\) over every ring. This proves the exact source \nolinkurl{SymLift} conditions and uniqueness, functorially under all base changes. No square-root extension is necessary for this all-ring proof.

On an isotropic subgroup these lifts are multiplicative: the lifts commute, and their product satisfies the two characterizing identities for the sum point. All assertions are scheme-theoretic and hold over nonreduced bases; geometric-point uniqueness alone would not suffice for this step.

The cubical multiplication identity provides on the universal group the unique normalized isomorphism

\[
 \alpha_0:[M]^*L\xrightarrow{\sim}L^{M^2}. \tag{AT20}
\]

It is the original \nolinkurl{IsNormMul}, as in LIFT100 HN7: pull the cubical identity to integer multiples, start with the rigidification at zero and the identity at one, and induct. Uniqueness follows from (AT3). The symmetry commutes with this isomorphism because their two composites are normalized global maps of the same line. Thus the deck lifts of \(A[M]\) on \([M]^*L\), transported through (AT20), are exactly its unique symmetric finite-order lifts. They fix the normalized section \(s_M=\alpha_0([M]^*s)\) exactly.

\subsection{8. The actual half-level root torsor and root-independent sections}\label{the-actual-half-level-root-torsor-and-root-independent-sections}

Put

\[
 C=A[\mathfrak b],\qquad H_b=A[\mathfrak b^2].
 \quad [M]:H_b\longrightarrow C
 \text{ is finite étale surjective with kernel }C. \tag{AT21}
\]

Indeed the torsion is finite étale because all integers are invertible in the characteristic-zero base. The coprime decomposition (AT1), or the literal geometric \nolinkurl{TateIso} and the original projective \(\mathcal O_E\)-lattice modulo \(M^2\), gives (AT21) on geometric fibres. Its degree is \(|C|=M^g\); finite étale morphisms and their equality can be checked on those fibres.

Every original level point \(\eta(x)\) lies in \(C\). This uses \nolinkurl{Obj.geom} rather than an unstated global \(\mathcal O_E\)-linearity of \(\eta\). First \(M\eta(x)=0\) by \nolinkurl{eta_add/eta_ker}, since \(M Dw.La\subset\mathfrak bDw.La\). Write \(x=x_b+x_{\bar b}\) using the two coprime submodules. At every geometric point, \nolinkurl{eta_ker} kills \(x_b\), and \nolinkurl{Obj.geom} identifies \(\eta(x_{\bar b})\) with \nolinkurl{TateIso.ev M}. Multiplication by an element of \(\mathfrak b\) then gives a multiple of \(M\), hence zero. Equality of sections of the finite étale group \(A[M]\) can be checked on geometric points: its diagonal is open and closed. This proves the equality over arbitrary bases, including nilpotents.

The subgroup \(H_b\) is isotropic in the theta group of \(L^{M^2}\). The adjunction gives the balanced commutator identity \(e_q(ax,y)=e_q(x,\bar a y)\). Choose \(a\in\mathfrak b^2\), \(b\in\overline{\mathfrak b}^2\) with \(a+b=1\). For \(x,y\in H_b\), the first factor in \(e_q((a+b)x,y)\) is one because \(ax=0\), and the second is one by balance because \(\bar b y=0\). This is an identity of the original theta-group operators, with its actual convention.

For \(t=\eta(x)\in C(S)\), choose a root \(t'\in H_b\) after the finite étale root torsor of (AT21) and define

\[
 S_t=U_{t'}\,s_M. \tag{AT22}
\]

Two roots differ by \(c\in C\subset A[M]\cap H_b\). The canonical \(H_b\)-lifts are multiplicative, and the canonical lift of \(c\) is the deck lift fixing \(s_M\). Hence (AT22) is independent of the root, with no scalar ambiguity. Descent of sections therefore gives an actual global section \(S_t\in H^0(A_S,L^{M^2})\). It commutes with every base change. Each \nolinkurl{IsVtOver} witness is this section after its stipulated pullback: its global normalized \(\alpha_0\), its symmetry lift, its twisted pullback and its action are unique by the preceding arguments and the explicit source equations.

The section \(S_t\) need not have an \(S\)-rational root \(t'\), so we do not assert \nolinkurl{IsVtOver} over \(S\) without a witness. The finite étale cover will supply that witness in \nolinkurl{ThValAt}.

\subsection{\texorpdfstring{9. Proof of \nolinkurl{ThValTot} and its uniqueness clause}{9. Proof of  and its uniqueness clause}}\label{proof-of-thvaltot-and-its-uniqueness-clause}

For the literal tuple \(x=(x_1,\ldots,x_n)\) and harmonic \(P\), define in \(S\)

\[
 v_{y,x,P}=M^{-\deg P}
   \sum_{\alpha\in\operatorname{supp}P} P_\alpha
   \prod_{j=1}^n\mathrm{nab}_y\bigl(M^2,S_{\eta(x_j)},\alpha_j\bigr).
 \tag{AT23}
\]

Take the fibre product of the \(n\) root torsors (AT21). It is affine, finite étale, faithfully flat, and has rank \(M^{gn}\) over \(S\). The pulled-back frame point exists with exactly the source \nolinkurl{FrMap}: tangent frames pull back, and its theta-section generator remains a generator by flat base change for sections of the proper abelian scheme (Stacks Project, Lemma 30.5.2, with \(i=0\)). The morphism is quasi-compact and quasi-separated, and the sheaf is invertible, so every hypothesis of that lemma holds over arbitrary \(S\). Only this flat base change is needed to produce the witness; we are not using an unproved representability claim over all nonreduced algebras. Its universal roots and (AT19)--(AT22) give the literal \nolinkurl{IsVtOver} tuple. Equation (AT17), including \(M^{-\deg P}\), now proves that the image of (AT23) equals the required \nolinkurl{thExpr}. This is existence in \nolinkurl{ThValAt}.

Conversely suppose \(v'\in S\) has any literal \nolinkurl{ThValAt} witness over a faithfully flat \(S\to S'\). Its \nolinkurl{FrMap} identifies the actual point, eigenframe and section. Each of its permitted translated sections is the corresponding base change of \(S_{\eta(x_j)}\), by §8. Equation (AT17) therefore gives

\[
 \phi(v')=\operatorname{thExpr}_{y'}=\phi(v_{y,x,P}). \tag{AT24}
\]

A faithfully flat ring map is injective, so \(v'=v_{y,x,P}\). This proves Theorem A, for all the source rings. It covers \(P=0\), \(n=0\), disconnected or nonreduced \(S\), and the zero ring without extra nonvanishing hypotheses.

Two pinned operators yield the same (AT23) by §5. Their \nolinkurl{ThValAt} relations, hence every fixed-\(x,P\) \nolinkurl{IsThRule}, agree. This implies \nolinkurl{NabUniq}, retaining the literal \nolinkurl{Dom} restriction if it is requested even though the comparison itself only needs the harmonic condition. The source \nolinkurl{existsUnique_thRule} now constructs the unique rule on its actual reduced finitely generated \nolinkurl{TRing} test objects; no enlargement of that definition is used.

\subsection{10. Named-law consequences and exact boundary}\label{named-law-consequences-and-exact-boundary}

For any already supplied family of actual \nolinkurl{Xo j}, coefficient and level records, choose the above operator at each permitted index. Then

\[
 \mathrm{nabPin},\qquad \mathrm{nabUniq},\qquad \mathrm{thTot}
 \quad\text{are supplied in their literal quantifiers}. \tag{AT25}
\]

For a separately prescribed operator, only the implications from its existing \nolinkurl{IsNab} are asserted. Filling (AT25) is not a replacement for compatibility with separately prescribed \nolinkurl{dNab}, short-theta, cusp, integral-model, comparison, or automorphic operators. It does not prove the remaining \nolinkurl{LawsCore}, \nolinkurl{LawsP} or \nolinkurl{LawsSh} fields.

Combined with LIFT100, an actual finite-number-field \nolinkurl{FrPt}, a nonzero harmonic polynomial, and the original \(DF+2\le M\) bound give that theorem's named \nolinkurl{ThNe0} law without separately assuming \nolinkurl{ThValTot} or a freely chosen pinned connection. The point, polynomial and source records remain actual premises. No norm, denominator, arithmetic content, DP44 sign or global period value is evaluated here.

In particular the complete \nolinkurl{thGen} tuple family remains distinct from the original coupled class sum followed by its fixed derivative/branching functional. Raw theta injectivity and the selected-image gate retain their previous scopes. This construction proves formal source laws; it does not assert that the actual selected arithmetic map is nonzero or that a complete admitted Host/Sizes instance exists.

\subsection{11. Literal interface map and verification boundary}\label{literal-interface-map-and-verification-boundary}

The final source manifest records paths, hashes, exact read extents and anchors. The main interfaces are:

\begin{longtable}[]{@{}
  >{\raggedright\arraybackslash}p{(\columnwidth - 2\tabcolsep) * \real{0.5000}}
  >{\raggedright\arraybackslash}p{(\columnwidth - 2\tabcolsep) * \real{0.5000}}@{}}
\toprule\noalign{}
\begin{minipage}[b]{\linewidth}\raggedright
Interface
\end{minipage} & \begin{minipage}[b]{\linewidth}\raggedright
Exact use
\end{minipage} \\
\midrule\noalign{}
\endhead
\bottomrule\noalign{}
\endlastfoot
\nolinkurl{TwoWeightEqOneAddRoot} --- Host, Level, SAB, Obj, TateIso & Actual involution, trace-one element, ideal halves, global cubical/adjunction/primitive maps, properness and geometric level identification \\
\nolinkurl{PushInvertible} --- Moduli and Coeff & Given universal proper family and splitting coefficient field; \nolinkurl{rep} is never invoked on nonreduced rings \\
\nolinkurl{IsCombAt} --- FrPt, FrMap & Full physical tangent eigenframe, actual theta generator, exact base change \\
\nolinkurl{CocycleRightZeroThd} --- FB, IsCanonGauge, nabVal & Every biframe coherence and the sign \(f\exp(kh_L)\) \\
\nolinkurl{ConnectionKillsSeries} --- CanonGaugeSt; \nolinkurl{CanonB.canonGauge} & Accepted existence over arbitrary commutative \(\mathbf Q\)-algebras \\
\nolinkurl{ChartChg}; \nolinkurl{CanonGaugeIndependenceAll} & Exact physical chart transformation and full equal-column gauge operator; §5 verifies the geometric applicability at finite order \\
\nolinkurl{ThGerm} --- IsChart, IsFrame, IsPrim, IsGerm, ChartAt, IsNab & Full constructor, especially global \nolinkurl{wP/wL}, all-ring quantifiers and the absence of an \nolinkurl{e}-truncation axiom \\
\nolinkurl{RowDeg} --- lap, harmK, HaO & Conjugate-pair harmonic identities summed over equal columns \\
\nolinkurl{ThetaPt} --- SymLift, IsTwistPull, IsActOf; \nolinkurl{ThetaVV} --- IsNormMul & Actual odd-lift, normalized universal multiplication and action conventions \\
\nolinkurl{IsThRule}; \nolinkurl{ThProps}; \nolinkurl{DatIndep} & Exact faithfully flat existence/uniqueness, actual test-ring rules and \nolinkurl{NabUniq} \\
\nolinkurl{PrimUniqPf} & The exact proper-to-abelian and global-normalization source bridge; the needed global-function input is independently supplied by (AT3) \\
\nolinkurl{HasDataCondition} & Only the named clauses (AT25), without inferring the rest of the source record \\
\end{longtable}

The two newly retrieved formal interfaces, \nolinkurl{ChartChg.lean} and \nolinkurl{WickVanishing/NabTrans.lean}, were read and saved from the pinned public source commit; their retrieval record states only that targeted scope. The exact earlier \nolinkurl{ThetaPt/ThetaVV} copies remain the hash-matched DP59 interfaces recorded in LIFT100. No source compilation, archive-wide validation or audit is claimed.

The companion diagnostics check finite odd-lift descent algebra (including nilpotent coefficients), truncated biadditivity, ordinary columnwise harmonic cancellation, and finite formal coordinate/gauge signs, together with frozen source fingerprints. They do not formally certify the general geometric theorem or evaluate the original arithmetic image. The author verification separately records whole saved-file reading, exact executable completion, source/primary reads and frozen hashes; it does not assert independent review.

\subsection{References and primary scopes}\label{references-and-primary-scopes}

Mumford, D. (1966). On the equations defining abelian varieties. I. \emph{Inventiones Mathematicae, 1}, 287--354. §1, pp.~289--297, for the theta group, commutator and descent interpretation. The particular all-ring normalization (AT18)--(AT22) is proved above. \url{https://pazuki.perso.math.cnrs.fr/index_fichiers/Mumford66.pdf}

The Stacks Project Authors. (2026). \emph{The Stacks Project}. Definition 37.11.1 and Lemma 37.11.7 (affine infinitesimal lifting, including the quasi-coherent torsor proof), \url{https://stacks.math.columbia.edu/tag/02GZ} ; Lemma 36.32.6 (universal global functions), \url{https://stacks.math.columbia.edu/tag/0BDM} ; Lemma 30.5.2 (flat base change), \url{https://stacks.math.columbia.edu/tag/02KH} . These exact statements and proofs were read directly for this argument. The invertible-module lifting and formal Taylor comparisons are proved explicitly above.

Zink, T. (1984). \emph{Cartiertheorie kommutativer formaler Gruppen}. Teubner-Texte zur Mathematik, 68. English translation, Chapter I, §5, Theorem 1.27, gives the commutative \(\mathbf Q\)-theorem used in (AT8); here the invariant-derivation Taylor construction fixes its normalization explicitly. The indexed primary excerpt supplies the statement and initial proof; direct PDF retrieval was unavailable during this reading, so no whole-PDF review is claimed. \url{https://perso.univ-rennes1.fr/matthieu.romagny/articles/Zink_Cartier_theory_of_commutative_formal_groups.pdf}

\subsection{Canonical algebraic-geometric context}\label{canonical-algebraic-geometric-context-77}

The construction lives over arbitrary commutative coefficient algebras of characteristic zero, including nonreduced rings. Proper SAB records give abelian schemes; universal degree-zero cohomology normalizes global line automorphisms, formal smoothness supplies compatible jets, and nilpotent lifting supplies line generators. The cubical line and its biextension induce an additive corrected frame; the OE eigenspace decomposition and trace-one CM weight give the exact quadratic harmonic cancellation. Odd normalized theta lifts and finite etale ideal-half root torsors yield actual sections and flat descent, without principalizing the ideal or identifying global units away from the chosen point. Primary locators: Stacks Lemmas36.32.5-6(tag0BDM),Definition37.11.1/Lemma37.11.7(tag02GZ),Lemma30.5.2(tag02KH); Mumford(1966),section1II-V p289,Proposition1 pp290-291,Proposition2 pp291-292,Theorem1 pp293-294,Proposition3 pp295-296,Theorem2 p297. Exact Alpoge literal interfaces and scoped reads are recorded in the root source record and mandatory locator addendum. Carried contextual locators: SGA7 I VII2.7/2.9 pp169-174 and VIII7.4(b) pp310-311; FGA232 Theorem3.1 pp148-149 and FGA236 Corollaries1.7/2.3 pp229/232; cote4 scan003 sections8.1.1-8.1.2/scans005-006 sections8.4.1-8.4.3 and cote54 scans005-007; SGA3 I4.7.3/4.7.3.1 and SGA6 VI Theorem3.3. These are geometric/source-law results, with no new integral Dieudonne slope, lambda-positivity, l-adic weight theorem or CM-motive L-function comparison. DP\_MULT Q3e/DP44/DP52/DP54-65 and THETA3 G1/A7 gain the stated source-law reduction; CX\_GRH CX-ALG-024 and GCT\_PLETH T-052 keep their arithmetic gates.

\clearpage

\section{AM-LIFT-101: source-locator precision addendum}\label{am-lift-101-source-locator-precision-addendum}

Recorded UTC: 2026-11-23T11:10:05.159986+00:00.

This mandatory locator addendum accompanies the frozen AM-LIFT-101 proof, script and source manifest. It corrects four stale line totals in the source manifest's prose \nolinkurl{read_scope} fields. Its numerical \nolinkurl{lines} fields and exact source hashes already record the correct values. No mathematical statement, source applicability, diagnostic result, frozen proof, script or manifest changes.

\begin{longtable}[]{@{}
  >{\raggedright\arraybackslash}p{(\columnwidth - 6\tabcolsep) * \real{0.2143}}
  >{\raggedleft\arraybackslash}p{(\columnwidth - 6\tabcolsep) * \real{0.2857}}
  >{\raggedleft\arraybackslash}p{(\columnwidth - 6\tabcolsep) * \real{0.2857}}
  >{\raggedright\arraybackslash}p{(\columnwidth - 6\tabcolsep) * \real{0.2143}}@{}}
\toprule\noalign{}
\begin{minipage}[b]{\linewidth}\raggedright
Literal file
\end{minipage} & \begin{minipage}[b]{\linewidth}\raggedleft
Stale prose total
\end{minipage} & \begin{minipage}[b]{\linewidth}\raggedleft
Actual complete extent
\end{minipage} & \begin{minipage}[b]{\linewidth}\raggedright
Exact SHA-256
\end{minipage} \\
\midrule\noalign{}
\endhead
\bottomrule\noalign{}
\endlastfoot
\nolinkurl{research/alpoge_source/Formalism/LocalToGlobal/ProjectiveLattices/DatIndep.lean} & 24 & 1--25 & \nolinkurl{05d03155235c56bfa33ed85c9730b03bf816a9825c8d274628e5f6dcf2d23520} \\
\nolinkurl{research/alpoge_source/Formalism/SharpRestriction/Divisibility/GeneratorPowers/PrimUniqPf.lean} & 176 & 1--173 & \nolinkurl{1db5758dc87df3cc148a1a96072170b1a260d0a61688dc85b47f210b14e85b1d} \\
\nolinkurl{research/alpoge_source/Formalism/SharpRestriction/Divisibility/GeneratorPowers/ChartFrameTrA.lean} & 85 & 1--77 & \nolinkurl{e0f9bb37b8f6eb1dd240df306bed1d52fb7fec2e0d322e65bcd37d78a2601e1b} \\
\nolinkurl{research/alpoge_source/Formalism/SharpRestriction/Divisibility/GeneratorPowers/ThExprUniqOfChartIndep2.lean} & 50 & 1--41 & \nolinkurl{9632492053ce5efdf8ef8c352f1b32e24da89a81d108bfc48747086252797cc9} \\
\end{longtable}

Each of these four complete actual files has been freshly read in full for this addendum. ``Whole file'' means precisely the actual complete extent displayed above; the earlier prose is not a claim to have read nonexistent extra lines. The frozen original source manifest is retained unchanged for provenance.

Frozen proof SHA-256: \nolinkurl{bd64500a0cea9421502ba6a1dcad1ec15cb1a674371b0a4a3cf78bab774a17b2}.

Frozen script SHA-256: \nolinkurl{d87dd2b5cdb371d07650d91ca253a206f5dadfdf964823accb72d7914a3adfcd}.

Frozen source-manifest SHA-256: \nolinkurl{2df68c8fc7fd147b4d564b427d52c70e4b36ca724e479f3a88b098b54a5be1dd}.

\clearpage

\section{AM-DO-124 --- the narrow-band original central-coset bound at k143}\label{am-do-124-the-narrow-band-original-central-coset-bound-at-k143}

\textbf{Canonical checked integration --- 2026-11-23T11:24:11.746929+00:00.} Complete root and independent proof/script/interface reviews and separate exact replays passed. The actual canonical predecessor has133/31 entries; this integration gives134/31, complete floors27/11. Frozen author history is preserved. AM-DO-124 CB1-27 removes the entire original bounded support27/total-eligibility27 all-II full-positive scalar-cross globally rank-two-unipotent k143 sector, with b=j\_int0 and2\^{}2111\textless x\textless2\^{}3328. The same accepted R100 narrow-band angular kernel and K120 actual narrow probability/moments retain the original physical central coordinate w=-U Im d, complete affine covolume/U scaling/750 cover, shape166/lambda32/25 profile, all167 inverse moments and actual-row zero-degree repair. Every original polynomial, finite, label, support, pivot, class/Haar and pattern factor is paid. The full exponent-725931/1024\textless-708 plus the entire new narrow-band tail2\^{}-8605114 gives one707 charge. Appending exactly once to actual DO123133/31 gives134/31 with both entire sums below2\^{}-230 and complete floors27/11 unchanged. The sole remaining support27 necessary alternative is total eligibility27,144\textless=k\textless=164,2\^{}2111\textless x\textless2\^{}3328,b=j\_int0. The k144 sufficient envelope924757/1024\textgreater0 is explicitly unclaimed. No full support28, signed nonvanishing, source admission, selected arithmetic-image content, prescribed period, specific032 value, evaluatedDP44 or RH conclusion.

Author: /root/definite\_oscillator, 23 November 2026 UTC. \textbf{Controlled pending proof.} R100, DO123 and K120 were canonically accepted at11:05:35,11:07:57 and11:08:21 UTC. This new arithmetic application requires its own whole author, independent and root reviews. Exact executable diagnostics supplement the continuous proof. No canonical file is changed here.

\subsection{1. Actual sector and complete statement}\label{actual-sector-and-complete-statement-1}

Retain the original J56 all-top H133 kernel, CM fields and involution, source/target pencils and omitted labels, all selected vector-valued Hilbert spaces, compact factors, binary words, canonical twists, phases, Haar measures, class masses and finite projectors. Put \[
 f=a=168,\quad H=133,\quad R=17424/133,\quad
 Q=1427908193729121522601,\quad U=77Q^2-14,
\] \[
 D_0=7^{140}29^{144}5^{84}13^{84},\qquad D_F=2^{168}D_0.
 \tag{CB1}
\] The actual baseline is DO123,133 all-top and31 uniform entries, with complete floors27/11. Consider its remaining globally rank-two-unipotent scalar-cross support27 family with all active typesII, all334 unmixed G reductions unipotent, and \[
 q_{\rm total}=27,\quad b_{\rm inactive}=j_{\rm int}=0,
 \quad 2^{2111}<x=N_ED_\eta<2^{3328},\quad k=143.
 \tag{CB2}
\] Here k is the number of nonzero original S covector values. The radial row parameter q below is different from total eligibility.

\textbf{Theorem.} The complete actual absolute normalized contribution of CB2 is less than2\^{}-707. This includes the entire K120 narrow-band discarded-fidelity tail once. Append one707 charge to the full actual chronological predecessor; retain every previous charge and both complete error sums. The floors remain27/11. The k144 sufficient envelope is positive and is explicitly unclaimed.

The improvement integrates the same retained-band angular kernel in the original physical central coordinate. It does not multiply independently proved trace bounds or reuse an old numerical sector charge.

\subsection{2. One restricted kernel with its full external tail}\label{one-restricted-kernel-with-its-full-external-tail-1}

For the actual compact-invariant selected density rho against invariant measure mu, let F\_I(g) be its fidelity integral restricted to both radial variables Y(z),Y(g\^{}-1 z) in I={[}16/529,16/527{]}. It is nonnegative and at most1, but the actual selected coefficient is not bounded by F\_I alone. K120 NB1--13 proves the new concentration premise, and its NB14--16 together with R100 NB15--16 and the band-independent K118 BT15--22 comparison prove, for every actual subset S of the bounded all-II family in CB2 and its original finite weights, \[
 \sum_{h\in S}|a(h)c(h)|
 \le\sum_{h\in S}|a(h)|\prod_\tau F_{I,\tau}(h_{p,\tau})
       +2^{-8605114}.
 \tag{CB3}
\] For every actual word and canonical twist, K120 derives2111/64\textless M1=E(Y\^{}-1)\textless2113/64 directly from the original RT22 factors. Positive and negative integer moments of order h=2\^{}35 give Markov ratios at most2108/2111 and2113/2116, each at most1−3/2116. The full positive moment corrections cost less than e. Thus the two-sided probability is less than2\^{}-40000000; this is proved for the narrower interval itself, not inferred from an older band.

Its tail already includes the full168-label probability loss and a containing count of all supports, rows, labels, patterns and central affine cosets. It is valid simultaneously for the single complete k143 sector here. No per-cell multiplication of that tail is needed.

Write h=I+N in the accepted full-positive normal form. The original covector r and isotropic slope eta satisfy \[
 N^3=0,\quad \operatorname{rank}_E N=2,\quad
 \eta\bar\eta=U,\quad
 \alpha=(N^2)_{nn}=U rJ_Z^{-1}r^*>0,\quad\alpha_\tau\le R,
\] \[
 N_{nn}=-Ud,\qquad \Re N_{nn}=\alpha/2,\qquad
 w=\Im N_{nn}=-U\Im d.
 \tag{CB4}
\] These are the original DO098 CF4--6 identities, applied at support27 through K113 SC27.1--7. The original arithmetic representative satisfies \textbar\textbar h\_tau\textbar\textbar≤133, the non-strict R048-RC cutoff stated in the mandatory DO122 cutoff addendum.

Let p=333/4, lambda=32/25 and b(t)=1+t/4. R100 NB3--14 prove from a single restricted angular kernel that \[
 \prod_\tau F_{I,\tau}(h_{p,\tau})
 <2\,2^{-649}\,2^{-933f/2}\exp(128f/13175)
   \prod_\tau J_\lambda(\alpha_\tau,w_\tau),
\] \[
 J_\lambda(t,v)=t^{-1/2}b(t)^{-167}
       \left(1+\frac{v^2}{4b(t)^2}\right)^{-p}e^{-\lambda t}.
 \tag{CB5}
\] Specifically the same original Jacobian kernel receives the height factor exp{[}-(64/25)(cosh(r)-529/527){]}. The exact identities cosh(r)\^{}2=(1+t/2)\^{}2+v\^{}2, sinh(r)\^{}2≥t b(t) and (cosh(r)+1)\^{}2≥4b(t)\^{}2+v\^{}2 give CB5 directly. The forward and reverse height estimates use D,D\^{}-1\textless528/527 and y/(1−y)≥16/513. The exact height loss is at least(c−529/527)/65 when c≥529/527; the other Cartan case uses the old positive bound. Since333/130\textgreater64/25, the same detailed-balanced kernel gives this exponential. Its local intercept is128/13175. No weighted-Jacobian alternative is multiplied into this formula.

The actual near-point identities are alpha\_p≥e\^{}-2r\_p alpha and \textbar w\_p-w\textbar≤6Hr\_p.~For u=v²/(4b(t)²+v²), \[
 -\partial_{\log t}\log J_\lambda
 =\frac12+(167-2pu)\frac{t}{t+4}+\lambda t>0,
 \qquad |\partial_v\log J_\lambda|\le p/2.
 \tag{CB6}
\] Only the original interval e\^{}-2r\_p alpha≤t≤alpha≤R is used for the transport derivative. Its full coefficient335+2lambdaR+3Hp is less than2\^{}16. Since r\_p\textless2\^{}-352, the product transport is below2, exactly the single near-point factor already shown in CB5. R100 NB14's positive Taylor bound e\textless11/4 and21504\textless26350 give \[
                       \exp(128f/13175)<e^2<(11/4)^2<2^3.
 \tag{CB7}
\] The additional band cost is therefore fully paid before the row enlargement.

\subsection{3. Original row and central difference lattice}\label{original-row-and-central-difference-lattice-1}

Fix eta and the canonical common whole-C G pattern J. The original four local states have E-norm exponents2,1,0 with multiplicities1,2,1. Let J0 be the real ideal consisting of each rank-zero G prime once. The exact original table gives \[
 |a_G(h)|\le2Q^{-335\cdot167}(N_EJ)^{335/2},\quad
 N_EJ\le Q^{334},\quad N_EJ\le Q^{167}N_FJ_0,
 \quad N_FJ_0\le Q^{167}.
 \tag{CB8}
\] At both original CM partners, every integral Z input has r(y),eta r(y) in J I\_Sigma. Their exact valuation intersection gives all167 values in \[
 L=J I_\Sigma D_\eta,\quad N_EI_\Sigma=T^{-54},\quad
 N_EL=z^{2f}=N_EJ\,xT^{-54},\quad q=R/(Uz^2)\ge k.
 \tag{CB9}
\] The three saturated transverse inputs are actual integral inputs, not values divided by2. Their Gram matrix G' has determinant U and G'\textgreater2\^{}36 I. The true covector metric is diag(I\_k,(G')\^{}-1), and the real covolume is D0\^{}(k+3) z\^{}(2(k+3)f)/U\^{}f.

For fixed eta,r and pattern, all central d values lie in one affine coset with difference lattice \[
 \Gamma_d=iJ_0B_\eta I_{\Sigma,F}^{2},\quad N_FB_\eta=x,
 \quad \operatorname{covol}\Gamma_d
       =\sqrt{D_F}(N_FJ_0)xT^{-54}.
 \tag{CB10}
\] This is the original four-positive-plus-negative entry calculation DO104 CD14--19, transported to27 in K113 SC27.7. The primitive Bezout extraction gives Delta d,eta Delta d,bar(eta)Delta d in the full-entry ideal, with no division byU or2. The local extra exponent is max(0,-v(eta),-v(bar eta)), precisely that of D\_eta bar(D\_eta); its real contraction is B\_eta. Zeros or poles meeting U,G,T cause no exception. The real-part shift is retained.

In the physical variable w=-U Im d the full real lattice therefore satisfies \[
 \operatorname{covol}\Gamma_w
 =U^f\sqrt{D_F}(N_FJ_0)xT^{-54},\quad
 \mu(\Gamma_w)<750\sqrt f\,U z_c,
\] \[
 z_c^f=(N_FJ_0)xT^{-54},\qquad
 U z_c=\frac Rq\left(\frac{N_FJ_0}{N_EJ}\right)^{1/f}<2^{-83}.
 \tag{CB11}
\] The last global inequality follows from the exact actual-host comparison U\^{}f Q\^{}167 2\^{}(3328+83f)\textless T\^{}54. It holds for every simultaneous actual q,J,J0, rather than treating a row norm and a divisor ratio as independent extrema. The coefficient750 is the full real codifferent/trace covering bound.

\subsection{4. Complete weighted integration of the physical central coset}\label{complete-weighted-integration-of-the-physical-central-coset-1}

For the fixed row set b\_tau=1+alpha\_tau/4 and phi(w)=prod\_tau(1+w\_tau²/(4b\_tau²))\^{}-p.~The exponential exp{[}-lambda sum alpha\_tau{]} is constant on that central fibre. The global bound \textbar\textbar grad log phi\textbar\textbar≤p sqrt(f)/2 and the true covering radius imply, by integration over translated Voronoi cells, \[
 \sum_{w\in w_0+\Gamma_w}\phi(w)
 \le\frac{\exp(375pf U z_c)}{\operatorname{covol}\Gamma_w}
          \prod_\tau\int_{\mathbb R}
                 (1+t^2/(4b_\tau^2))^{-p}\,dt.
 \tag{CB12}
\] This sums a positive upper envelope on a containing affine lattice; no enlarged point is declared to be an actual arrow. Monotone convergence is legitimate, since p\textgreater1/2 and all displayed one-dimensional integrals converge. The original pattern's finite coefficient bound remains fixed throughout this enlargement.

The ordinary positive Beta integral and Cauchy--Schwarz for Gamma give \[
 \int_{\mathbb R}(1+t^2/(4b^2))^{-p}\,dt
 =2b\sqrt\pi\frac{\Gamma(p-1/2)}{\Gamma(p)}
 \le\mathcal G b,\qquad \mathcal G=4\sqrt{\pi/329}.
 \tag{CB13}
\] Indeed Gamma(p-1/2)\^{}2≤Gamma(p-1)Gamma(p) and p-1=329/4. The exact bound375pf2\^{}-83\textless1/2 makes the complete Voronoi inflation less than2. Combining CB5 and CB10--13 gives \[
 \sum_{\text{actual central fibre}}\prod_\tau F_{I,\tau}(h_{p,\tau})
 <4\,2^{-649}\,2^{-933f/2}\exp(128f/13175)
 \frac{\mathcal G^fT^{54}}{U^f\sqrt{D_F}(N_FJ_0)x}
 \prod_\tau\Phi(\alpha_\tau),
\] \[
                   \Phi(t)=t^{-1/2}(1+t/4)^{-166}e^{-\lambda t}.
 \tag{CB14}
\] The two factors2 in4 are exactly near-point transport and the weighted Voronoi comparison. The central integral consumes one power of1+t/4, leaving shape166. The physical U\^{}f Jacobian has not been dropped.

\subsection{5. All positive moments, including the repaired zero degree}\label{all-positive-moments-including-the-repaired-zero-degree-1}

Let mu be the pushforward of Gamma(shape166,rate4) times ds/sqrt(pi s) by c=lambda+b+s. It is a positive locally finite measure of infinite total mass, and Phi(t)=int exp(-ct)dmu(c). Its positive inverse moments are \[
 M_j=\frac1{\Gamma(j)}\int_0^\infty
          t^{j-3/2}(1+t/4)^{-166}e^{-\lambda t}\,dt,
 \qquad j\ge1.
 \tag{CB15}
\] All are finite because lambda\textgreater0 and j\textgreater1/2. The accepted K120 NB17--24 positive rational moment theorem bounds them by pi r\_j. Its explicit N=512 positive integral majorant uses \[
 e^{-\lambda t}\le(1+t/400)^{-512},\quad
 r_j=\frac{2\,4^{166}400^{512}}{\pi(j-1)!}
       \int_0^\infty
       \frac{s^{2j-2}}{(s^2+4)^{166}(s^2+400)^{512}}\,ds.
 \tag{CB16}
\] Two exact partial-fraction base integrals and an integration-by-parts recurrence give rational r\_j for every1≤j≤167. The theorem proves positivity and normalization of these constants; finite recurrence diagnostics are not substituted for that positive integral proof. Define the exact strict dyadic ceilings \[
 b_j=\left\lceil\log_2[((22/7)r_j)^{168}]\right\rceil_{\rm strict},
 \qquad P_j=2^{b_j}>M_j^f.
 \tag{CB17}
\] Every required value is independently checked from the final frozen K120 table. No zeroth moment or finite total mass is asserted.

On a genuine row the product energy floor is prod alpha\_tau≥(Ukz²)\^{}f=(Rk/q)\^{}f.~Reserve theta=2\^{}-20 of the Gaussian exponent before enlarging the row. With B=(prod c\_tau)\^{}(1/f), AM-GM and e\^{}-u≤u\^{}-1 imply \[
 e^{-\theta\sum c_\tau\alpha_\tau}
 \le e^{-\theta fBRk/q}
 \le(\theta Rk/q)^{-f}\prod c_\tau^{-1}.
 \tag{CB18}
\] Thus the sole zero-degree term integrates to at most P1(q/(theta Rk))\^{}f.~Every positive degree integrates with its corresponding P\_j. No per-embedding energy floor is assigned to a zero or out-of-cutoff tuple introduced by the later positive Gaussian enlargement.

Keep the complete universal degree15 scalar coefficient K of K111 HP1--18, with its original half-trace minimum, and c0=2\^{}36. Explicitly \[
 c_{336}=168\binom{336}{168}/4^{168},\quad
 M_{15}=\frac{6015684926187167549292201834533348561577984}
 {2618017117694612446410293519},
\] \[
 K=\frac{336M_{15}}{c_{336}}(7841/5000)^{168}168!
 \left[\frac{(365729/500)^2}{168(314159/100000)^2}\right]^{168}.
 \tag{CB19}
\] The exact K110 PM10--12 product-one change and anisotropic Poisson bound apply for every positive c\_tau, without declaring the changed real lattice to be an ideal. For d\_k=k+3 the complete row factor is binom(164,k)U\^{}f pi\_+\^{}(d\_k f)/D0\^{}d\_k times \[
 \mathcal P_k(q)=
 \sum_{\substack{0\le l\le k,\ 0\le h\le3\\ j=l+3-h\ge1}}
 \binom kl\binom3h K^{k-l+h}c_0^{-hf}
       \left(\frac q{(1-\theta)R}\right)^{jf}P_j
 +K^{k+3}c_0^{-3f}P_1\left(\frac q{\theta Rk}\right)^f.
 \tag{CB20}
\] All scalar degrees and all four transverse terms remain, including(l,h)=(0,3). The row U\^{}f cancels the physical central U\^{}-f in CB14 exactly. The actual energy and band restrictions were applied before this enlargement; no new actual-arrow assertion is made about its points.

\subsection{6. Every original finite factor and both endpoint derivatives}\label{every-original-finite-factor-and-both-endpoint-derivatives-1}

K112 active eligibility and K113 bounded paired uniqueness give at most one eta for each original27-place support throughout x\textless2\^{}3328. Eta is therefore fixed before the row and central count. DO105 EL9--16's actual same-label argument fixes the only possible active label, at both partners and including central-only poles. In CB2 the active cross pole persists when the central coordinate varies. The complete central bound is for one containing affine coset for that actual pair; no central-dependent label union is suppressed.

The141 inactive factors are fixed by the same cross data, as in DO108 AC8--10. The full168-factor bound is(T−1)\^{}-168\textless2T\^{}-168. Keep supports\textless2\^{}104, the retained original first-pivot choices\textless2\^{}8, Green2, the T correction2, near-point2, weighted-cover2 and conservative full-pattern4. These give118 fixed bits. The full T power is335·27+54−168=8931. Class/Haar masses and mixed contractions remain exactly those of the original normalized coefficient. Eta,r,d reconstruct the original matrix, so there is no unpaid beta, negative-column or source-isometry fibre.

For each positive polynomial degree j, and effective degree1 for the repaired term, all q and x dependence is \[
                 x^{-337/2}q^{(j-335/2)f}.
 \tag{CB21}
\] Its q derivative divided byf is j−335/2\textless0 for all degrees here. At the all-rank-zero endpoint \[
 q_*^f=\frac{R^f T^{54}}{U^fQ^{334}x},
 \qquad q=q_*D^{1/f},\quad D=Q^{334}/N_EJ,
 \quad (N_FJ_0)^{-1}\le Q^{-167}D,
 \tag{CB22}
\] the coupled x logarithmic derivative is exactly−1−j\textless0. The cover inflation was already bounded uniformly in CB11--14; no untracked moving central radius remains.

For each degree the entire original four-state containing sum is \[
                 \prod_{v=1}^{334}(1+q_v^{j-333/2})^2<2.
 \tag{CB23}
\] The target sector uses j≤146; the failed diagnostic uses j≤147. Thus j−333/2≤−19.5\textless−1 throughout, and the exact old inequality(1+1/g0)\^{}668\textless2 suffices. We retain4 conservatively in the118 bits. There is no additional pattern or orientation factor and no extension to an exceptional degree167 is needed.

\subsection{7. Exact complete endpoints and the paid union}\label{exact-complete-endpoints-and-the-paid-union-1}

Use the actual-host directed logarithms L\_T=1019333/1024, L\_Q=287843/4096, L\_U=601356/4096 and L\_R=28810/4096, with the same upper/lower directions as DO122. At x=2\^{}2111 the exact bracket is251124964\textless q\_*\textless251124965. Also \[
        \frac{\mathcal G^f}{\sqrt{D_F}}
        \le\frac{(16(22/7)/329)^{84}}{\sqrt{D_F}}<2^{-1110}.
 \tag{CB24}
\] After CB7's3 band bits are included, the common exponent is \[
 C=118-933f/2-649-1110+3
     +(335f/2)(L_R-L_U)+8931L_T-336\cdot167L_Q
     -(337/2)2111
   =591769685/1024.
 \tag{CB25}
\] All factors precede the directed logarithmic substitution. In particular no old central ball volume or independent weighted-Jacobian improvement appears.

The strict scalar ratio ceiling is ceil(log2{[}K((1-theta)R/q\_-)\^{}f{]})=-935. Exact moment ceilings and the full binomial/transverse coefficients bound every nonmain term of CB20 relative to its(l,h)=(k,0) term. For k143 the575 nonmain terms are each at most2\^{}-353 times that term; their sum is strictly less than the main term. The entire positive polynomial is therefore bounded by twice the main term. The dyadic moment P146 is included in that complete rational ceiling. Every one of the four transverse terms and every scalar degree is present.

\begin{longtable}[]{@{}
  >{\raggedleft\arraybackslash}p{(\columnwidth - 10\tabcolsep) * \real{0.1739}}
  >{\raggedleft\arraybackslash}p{(\columnwidth - 10\tabcolsep) * \real{0.1739}}
  >{\raggedleft\arraybackslash}p{(\columnwidth - 10\tabcolsep) * \real{0.1739}}
  >{\raggedleft\arraybackslash}p{(\columnwidth - 10\tabcolsep) * \real{0.1739}}
  >{\raggedleft\arraybackslash}p{(\columnwidth - 10\tabcolsep) * \real{0.1739}}
  >{\raggedright\arraybackslash}p{(\columnwidth - 10\tabcolsep) * \real{0.1304}}@{}}
\toprule\noalign{}
\begin{minipage}[b]{\linewidth}\raggedleft
k
\end{minipage} & \begin{minipage}[b]{\linewidth}\raggedleft
Moment degree
\end{minipage} & \begin{minipage}[b]{\linewidth}\raggedleft
b\_(k+3)
\end{minipage} & \begin{minipage}[b]{\linewidth}\raggedleft
Full positive factor ceiling
\end{minipage} & \begin{minipage}[b]{\linewidth}\raggedleft
Complete exponent
\end{minipage} & \begin{minipage}[b]{\linewidth}\raggedright
Scope
\end{minipage} \\
\midrule\noalign{}
\endhead
\bottomrule\noalign{}
\endlastfoot
143 & 146 & −112538 & −578609 & −725931/1024 & Complete target sector \\
144 & 147 & −113109 & −576997 & 924757/1024 & Failed positive sufficient envelope only \\
\end{longtable}

For clarity, the full factor ceiling is \[
 \left\lceil\log_2\!\left[
 2\binom{164}k\frac{\pi_+^{(k+3)f}}{D_0^{k+3}}
 q_-^{-335f/2}
 \left(\frac{q_+}{(1-\theta)R}\right)^{(k+3)f}
 \right]\right\rceil_{\rm strict}+b_{k+3}.
 \tag{CB26}
\] The same arithmetic retains all579 nonmain corrections for the failed next cell, whose maximum relative exponent is−356. Its complete positive upper estimate gives no k144 removal. It is an honest boundary of this sufficient envelope, not a lower bound for that sector.

The complete k143 main bound is strictly below2\^{}-708. Paying the entire narrow-band arithmetic tail once gives \[
                     2^{-708}+2^{-8605114}<2^{-707}.
 \tag{CB27}
\] There is no extra endpoint or witness union. The actual chronological prefix, all133 baseline all-top charges and all31 uniform charges are retained. The executable accepts a later complete prefix without changing these mathematical subjects, rejects a duplicate707 and checks both full sums below2\^{}-230. The original floors27/11 remain unchanged.

\subsection{8. Verification, source and conclusion boundary}\label{verification-source-and-conclusion-boundary-1}

The final source packet fingerprints R100 NB1--16's same-kernel narrow-band profile, K120 NB1--24's actual narrower probability and exact positive moment table, the band-independent K118 probability-to-arithmetic comparison, and accepted DO123's complete original affine central/row/pattern geometry, reproduced above. DO122's mandatory non-strict-cutoff addendum is retained. The continuous inequalities are proved above and in those stated inputs; finite models supplement them without claiming to enumerate global arrows. The unchanged DO123141/142 charge is retained only as a chronological predecessor, not reused for this new143 cell.

Additional literal interfaces are DO098 CF1--16, DO104 CD10--19 and section5, K113 SC27.1--10, DO105 EL9--16, DO108 AC8--10, K107 SV17--19, K109's actual transverse metric, K110 PM10--12 and K111's all-minimum/all-positive-parameter scalar theorem. The ordinary Gamma/Beta normalizations are NIST DLMF5.9.1 and5.12.1/3. The selected norm/metric source remains Nakahama(2016), equation5.7 and Theorem5.1, printed960--961, through the original checked K043/K046/DO080 dictionary. No scalar representation is substituted.

The same projective modules, involution, Hermitian form, polarization, cubical theta line and finite projectors are retained. Carried FGA/SGA context is unchanged; no fresh reading or new formal-group, K-theoretic, l-adic-weight or CM-motive theorem is inferred from it. There is no full support28, actual selected arithmetic-image nonvanishing, prescribed period, source admission, specific032 or evaluated DP44/RH conclusion. The only proposed main-gate change is the complete single-cell sector CB2, after all independent and root reviews pass.

\subsection{Canonical algebraic-geometric context}\label{canonical-algebraic-geometric-context-78}

The physical coordinate change is a real analytic treatment of the same original central affine module, preserving its shift, CM partner conditions and polarization. It adds no arithmetic parameter or field and changes no selected metric. NIST DLMF5.9.1,5.12.1/3 and5.4.6 are applied with explicit positive integral arguments. Inherited context only: SGA3 I4.7.3/4.7.3.1; SGA6 VI Theorem3.3; FGA232 Theorem3.1 pp148-149; FGA236 Corollaries1.7/2.3 pp229/232; SGA7 I VII2.7/2.9 pp169-174 and VIII7.4(b) pp310-311; cote4 scan003 sections8.1.1-8.1.2 and scans005-006 sections8.4.1-8.4.3; cote54 scans005-007. No new formal-group/Dieudonne, lambda-ring/K-theory, l-adic-weight or CM-motive comparison. DP\_MULT Q3e/DP44 and THETA3 G1/A7 receive the strictly smaller conditional finite gate; CX\_GRH CX-ALG-024 and GCT\_PLETH T-052 retain their arithmetic/global gates.

\clearpage

\section{AM-LIFT-102 --- Prescribed exponential frames and the literal connection pin}\label{am-lift-102-prescribed-exponential-frames-and-the-literal-connection-pin}

\textbf{Canonical checked integration --- 2026-11-23T11:47:40.333149+00:00.} Complete saved-proof and executable reviews and separate exact diagnostic replays passed. No numerical error charge is introduced. Frozen author history is preserved. AM-LIFT-102 TheoremsA-E/EC1-16 proves the prescribed-frame converse and the actual-package harmonic comparison. Given the actual geometric records, every coherent IsExpChart/IsExpGauge is realized by the original Poincare biframe with h0, so its nonnegative ordinary Taylor operator satisfies literal IsNab with the exact product factorial. The normalized Appell-Humbert reduced frame at every actual complex point gives such data, for mixed signatures as well as pure types, and admits a pinned all-ring extension. The entire ShortValAt and class relations are unchanged between two pinned operators on every harmK input. In an actually supplied OMPins/OneMapTieP package, CovPin already pins the operative DeltaDat.nab2 and Pins.io supplies range(io)=span(HHa) contained in harmK. Thus no additional equality of individual ordinary coefficients is needed to compare its original thLaValAt, IsThClsD and OneMapRow relation in the reduced gauge. All original witnesses, roots, theta-group sections, beta/index/cardinality/N-degree normalizations, coefficient eF map, vol/gam/ep/cdelta weights, phases and selected maps remain fixed. The package itself, absent short witnesses and sec/cdelta laws, mixed integral coefficient descent and original coupled selected-image nonvanishing remain unconstructed. Unknown unpinned operators still require their own witness. No error charge; complete134/31 arrays and floors27/11 unchanged.

\textbf{Status: completed mathematical proof, 2026-11-23; hash-specific author verification and exact diagnostics are recorded in the companion packet.} Independent review and canonical adoption are separate. AM-LIFT-101 is the accepted canonical input, including its mandatory locator addendum. This result pins prescribed exponential-frame operators in the literal all-ring source and gives an actual complex reduced-frame realization; it does not identify an arbitrary unknown stored operator, construct a complete admitted source, or prove nonvanishing of the original coupled selected image.

\subsection{1. Theorems and exact quantifiers}\label{theorems-and-exact-quantifiers}

Retain the actual geometric Host/Datum/odd Level/Coeff/Moduli records and the accepted formal gauge input of AM-LIFT-101. At an actual frame point y over an arbitrary commutative coefficient-field algebra S, suppose the \emph{prescribed} data are the literal coherent \nolinkurl{IsExpChart y x} and \nolinkurl{IsExpGauge y x e}. These include the physical tangent frame, diagonal integral CM action, all truncation compatibilities, origin rigidification, and the globally normalized beta generators specified in \nolinkurl{IsExpGauge.coc}.

\textbf{Theorem A --- exact prescribed-frame converse.} For every such prescribed frame e, there is an actual \nolinkurl{ChartAt} with additive formal group law, diagonal integral CM action, trivial formal biextension multipliers, and canonical gauge h=0, whose primitive line frame is \textbf{exactly e} at every order. Its ordinary source connection coefficients equal \nolinkurl{expVal} of the prescribed germ for every monomial, not merely after harmonic contraction. A family of such prescribed data at every S,y therefore pins its prescribed nonnegative Taylor operator by literal \nolinkurl{IsNab}.

LIFT101 constructed some pinned operator and proved harmonic independence of pinned operators. The present converse supplies a sufficient geometric witness for an operator already prescribed by its exponential frame. It does not assert that any unknown preselected operator has that witness.

The source already proves the balanced-quadratic ratio theorem in \nolinkurl{GaugeUniqPf.gaugeUniq_of_global}; that classification is an existing accepted interface, not a new source theorem credited here. The new step is realization of its quadratic ambiguity by the actual Poincare biframe, with exact ordinary-derivative equality.

\subsection{2. Reference biframe and global beta normalization}\label{reference-biframe-and-global-beta-normalization}

Use LIFT101 AT3--10 to construct a coherent corrected biframe E0 over S, in the unique additive formal coordinates determined by the physical eigenframe. Every addition, symmetry and integral CM adjunction multiplier is one. These are the pulled-back original structural maps. The unique global normalized primitive wP gives a coherent one-variable line frame

\begin{verbatim}
e0(Z) = wP E0(Z,w0 Z).
\end{verbatim}

The trace-one identity is w0 + bar(w0) = 1. Global normalized beta\_a exists from the primitive and normalized adjunction generators: the actual \nolinkurl{ESecToUnit.exists_beta} supplies precisely the line combination with maps (1+a bar(a), a, 1, a bar(a)) and exponents (1,-2,-1,-1), normalized to rig\^{}(-3) on the universal zero section. The argument displayed in \nolinkurl{ExpGauge.expGaugeOf_of} provides its exponential cocycle identity for this corrected coherent chart: specialize phi to the identity in the additive coordinates. The required action-equivariance equality, which that source argument obtains from its general \nolinkurl{ExpActOfSt} input, here follows directly from the already diagonal action. Thus no separate universal \nolinkurl{LinEigSt} or \nolinkurl{ExpActOfSt} premise is silently introduced. The remaining inputs are the actual original cubical/adjunction identities and the normalized beta just constructed. The truncation premise is supplied by this particular construction, not inferred for an arbitrary \nolinkurl{ChartAt}.

Any two globally normalized beta\_a generators are equal: their ratio is a unit on the proper universal abelian scheme, hence a base unit, and zero-section normalization makes that base unit one. This is LIFT101 AT3 and the literal \nolinkurl{GaugeBeta.betaUniq_of_global}. It is stronger than equality at a single unrelated point, and uses the actual universal normalization.

The prescribed exponential chart equals the reference additive chart. To justify comparison, the actual \nolinkurl{FrPt.basis} and \nolinkurl{Coeff.conj} split the full tangent module; the \nolinkurl{IsExpChart.tan} clause therefore makes its derivative an isomorphism. Compose with the reference formal-coordinate isomorphism of LIFT101 AT4--8. The resulting formal map has invertible derivative, so the formal inverse is constructed degree by degree over S; it intertwines the two additive laws. An additive formal homomorphism over a Q-algebra is linear: compare its value at 2Z with twice its value at Z, and use the invertible rational factors 2\^{}j-2 for j != 1 (the zero constant term is already fixed). The physical derivative is the identity. This proves equality at every finite order, over nonreduced and non-Noetherian S as well as fields.

\subsection{3. Ratio classification, including arbitrary nonreduced bases}\label{ratio-classification-including-arbitrary-nonreduced-bases}

Write e = e0 R. The coherent unit R is a formal series over S with R(0)=1. Comparison using the common global beta\_a gives, for every integral a,

\begin{verbatim}
R((1+a bar(a))Z) = R(aZ)^2 R(Z) R(a bar(a)Z).                 (EC1)
\end{verbatim}

Let q=log R. Formal logarithms and exponentials are defined coefficientwise over every Q-algebra; this uses augmentation-adic nilpotence at each finite order, not reducedness or convergence. Taking a=2 gives

\begin{verbatim}
q(5Z) = 2q(2Z) + q(Z) + q(4Z).                            (EC2)
\end{verbatim}

For a homogeneous degree j coefficient the multiplier is c\_j=5\^{}j-2*2\^{}j-1-4\^{}j. Its values at j=0,1,2 are -3,-4,0. At j=3 it is44, and

\begin{verbatim}
c_(j+1)=5c_j+4^j+6*2^j+4
\end{verbatim}

is positive for every j\textgreater=3. Thus q is homogeneous quadratic. This elementary calculation matches the degree conclusion of the accepted gauge-ratio statement and is included to make the all-ring use explicit.

Let B be its symmetric polarization, so q(Z)=B(Z,Z)/2. Substituting q(Z)=B(Z,Z)/2 into the logarithm of the full identity (EC1) yields

\begin{verbatim}
B(Z,(a bar(a))Z) = B(aZ,aZ).                              (EC3)
\end{verbatim}

Polarizing (EC3) then shows that, for vectors in the tau and tau' eigenspaces, its coefficient is killed by

\begin{verbatim}
tau(a bar(a)) + tau'(a bar(a)) - 2 tau(a) tau'(a).          (EC4)
\end{verbatim}

If tau'=bar(tau), this polynomial vanishes identically. Otherwise it is a nonzero polynomial in the independent split embedding coordinates of E tensor\_Q F. The integral lattice O\_E is Zariski dense in that affine space: choose an integral basis and apply induction to a polynomial vanishing on all integer coordinate tuples. Hence some integral a makes (EC4) a nonzero element of F, therefore a unit of any nonzero F-algebra S. The corresponding B coefficient is zero. The zero ring is automatic.

Consequently B is symmetric, bilinear, zero on nonconjugate pairs, and satisfies

\begin{verbatim}
B(aX,Y)=B(X,bar(a)Y),       B(Z,w0 Z)=B(Z,Z)/2=q(Z).       (EC5)
\end{verbatim}

This permits arbitrary mixed signatures and multiplicities. It does not assert q=0 in that situation. K018's additional pure J-type uniqueness is a different statement.

\subsection{4. Exact biframe realization and ordinary derivative pin}\label{exact-biframe-realization-and-ordinary-derivative-pin}

Define at every finite order

\begin{verbatim}
E(X,Y)=E0(X,Y) exp(B(X,Y)).                               (EC6)
\end{verbatim}

Bilinearity, symmetry and the adjoint balance in (EC5) show that both addition multipliers, the symmetry multiplier and every integral CM adjunction multiplier remain one. The axes remain canonically normalized. These identities use the same universal cubical/adjunction maps as E0; no new global map is substituted. Compatibility in the jet order follows from a single formal exponential.

The same globally normalized primitive now gives

\begin{verbatim}
wP E(Z,w0 Z)=e0(Z) exp(q(Z))=e(Z).                        (EC7)
\end{verbatim}

Take the literal \nolinkurl{FB.triv} for the additive group and its diagonal O\_E-action, uvec equal to the physical coordinate vectors, and h=0. All \nolinkurl{ChartAt} fields follow from the prescribed chart, (EC6), the original global primitive, and (EC7), including \nolinkurl{rig0}. Every section of L\^{}k has a unique germ in e\^{}k, since e is a coherent generator. Thus the source \nolinkurl{IsGerm} and prescribed \nolinkurl{IsExpGerm} are the same relation in these coordinates.

The source formula \nolinkurl{nabVal} is D\^{}alpha applied to f exp(k h\_L) at zero. Here h=0 and the formal group is additive. Therefore for all k\textgreater=0, all actual sections F and all monomials alpha,

\begin{verbatim}
nab_y(k,F,alpha) = alpha! [Z^alpha] f = expVal(f,alpha).  (EC8)
\end{verbatim}

The factorial is the product of the source ordinary factorials; it is not a divided derivative or an independently rescaled analytic differential. The source \nolinkurl{IsNab} requires one chart witness per S,y for all k,F, and (EC6) supplies exactly that. Negative k can be extended by the LIFT101 rigidified zero-order convention; it is outside the \nolinkurl{IsNab} requirement.

For an already supplied \nolinkurl{NabV}, equality with (EC8) at every S,y,k,F,alpha is the exact remaining operator hypothesis. The actual CM involution supplies the required \nolinkurl{StarRing} and star=H.bar clause. The coordinate vectors indexed by \nolinkurl{ltIdx} give the \nolinkurl{IsChart.basis} field directly. The one coherent biframe and one global primitive at each S,y work for every k and every section F, not a new chart depending on the derivative order. For a family whose nonnegative operators are defined by prescribed exponential gauges, (EC8) proves its \nolinkurl{IsNab}. Unconstrained ringwise formulas do not acquire this property automatically.

\subsection{5. Actual complex reduced frame, without pure-type uniqueness}\label{actual-complex-reduced-frame-without-pure-type-uniqueness}

\textbf{Theorem B --- the actual complex reduced-frame application.} Let sigma:F-\textgreater C be any embedding and let y be an actual \nolinkurl{FrPt} over the literal \nolinkurl{CAt sigma} coefficient algebra. Identify its underlying complex field and use the physical eigenframe. The following argument applies to any signature allowed by the given records; it does not require a pure CM type or that the abelian variety be a simple CM variety. Retain its actual abelian variety, full integral CM action, rigidified symmetric L and polarization. Write A(C)=V/Lambda and use the physical tangent vectors for V. The Appell--Humbert presentation of L has multiplier

\begin{verbatim}
j_lambda(z)=alpha(lambda) exp(pi H(z,lambda)+pi H(lambda,lambda)/2),  (EC9)
\end{verbatim}

where H is Hermitian, linear in its first argument, and E=Im H is integral on Lambda. Origin rigidification fixes the otherwise constant ambiguity of the line isomorphism: two isomorphisms of these same analytic line bundles differ by a nowhere-zero global function on the compact connected abelian variety, hence by a constant. This is a statement about the specified reduced Appell--Humbert presentation, not about arbitrary frames on the universal cover. Symmetry of L implies alpha(lambda) in \{+1,-1\}. The Rosati condition gives H(ax,y)=H(x,bar(a)y).

Put T=a bar(a). The Hermitian form of the beta\_a line combination is zero, because

\begin{verbatim}
H((1+T)x,(1+T)y)-2H(ax,ay)-H(x,y)-H(Tx,Ty)=0.             (EC10)
\end{verbatim}

Indeed T is selfadjoint and H(ax,ay)=H(x,Ty). The semicharacter also cancels:

\begin{verbatim}
alpha((1+T)lambda)/(alpha(a lambda)^2 alpha(lambda) alpha(Tlambda))=1, (EC11)
\end{verbatim}

since E(lambda,Tlambda)=0 and alpha(a lambda)\^{}2=1. Therefore the constant-one frame on the universal covering vector space obeys exactly the beta\_a exponential identity. The original globally normalized beta\_a restricts to this frame: the two fibrewise global generators have constant ratio on the proper complex abelian variety, fixed to one by the origin normalization. This comparison does not infer equality of unrelated universal sections from a single point.

Analytic exponential coordinates induce algebraic finite jets: choose a smooth algebraic etale coordinate system at zero; its analytification is a local analytic coordinate system, and both completed local rings are C{[}{[}Z{]}{]}. The Taylor series of the covering exponential consequently defines compatible morphisms from every algebraic truncated polynomial ring. The same local identification carries the holomorphic line generator to generators of all algebraic finite jets. Thus the reduced frame gives the literal \nolinkurl{IsExpChart} and \nolinkurl{IsExpGauge}, with all exact physical and global-normalization clauses. In particular the holomorphic one-variable frame is not asserted to be a global nonvanishing section of the ample line on A: it is a frame on the covering vector space and on the formal neighbourhood. The beta difference line, whose multiplier has just been proved trivial, does have the normalized global generator required in \nolinkurl{IsExpGauge.coc}.

The converse construction therefore pins the ordinary reduced-theta derivatives for all complex frame points, including mixed types. It makes no claim that the reduced frame's individual coefficients descend to the original number field or are integral. Those stronger pure-J conclusions in K018 used an additional zero-ambiguity argument.

Primary anchor freshly read: Bannai and Kobayashi, author arXiv:math/0610163v4, section1.2 beginning printedp9; Theorem1.7 and Definition1.8 are on printedp10 (PDFindex9). The published article is Bannai and Kobayashi (2010). The exact multiplier in that primary text is (EC9). The new direct retrieval is recorded in the source manifest; only the Appell--Humbert classification and reduced-frame convention are used, not a p-adic integrality theorem. The beta cancellation and exact algebraic-jet application are proved here.

\subsection{6. Simultaneous all-ring extension and the prescribed-family boundary}\label{simultaneous-all-ring-extension-and-the-prescribed-family-boundary}

\textbf{Theorem C --- an all-ring extension of the prescribed complex operator.} For each complex coefficient-field embedding and each actual complex frame point, take the just-verified prescribed reduced frame. If the same actual record is described twice, use its one normalized reduced presentation; uniqueness in section5 makes the prescriptions agree. For all other coefficient algebras/points, use a coherent corrected frame from LIFT101. The source \nolinkurl{NabV}/\nolinkurl{IsNab} does not demand naturality of every ordinary coefficient or auxiliary choice. Classical selection of the actual chart witnesses at each quantified record therefore defines an all-ring pinned family matching the prescribed reduced derivatives at the complex points. Choose the negative-power entries as in LIFT101 AT7; the resulting family also satisfies \nolinkurl{NabZero} at every integer power. Every nonnegative ordinary coefficient of an actual section at a complex point equals its reduced-theta ordinary derivative in the physical tangent directions, with exactly the factorial in \nolinkurl{expVal}. LIFT101 then supplies naturality and independence of its harmonic theta expressions.

This is a constructor with prescribed complex values, not proof of equality to an arbitrary previously stored \nolinkurl{Top.nab} or \nolinkurl{DeltaDat.nab2}. For those records, the needed equality with their actual specified exponential frames must be supplied or checked. If a fixed operator is already given by such frames on every quantified algebra, the theorem certifies it without replacing the operator. If only its complex description is given, the theorem constructs a compatible all-ring extension; it does not identify its previously unspecified noncomplex entries.

\subsection{7. Exact short-expression and class-sum transport}\label{exact-short-expression-and-class-sum-transport}

\textbf{Theorem D --- short-expression independence between actual pins.} For two pinned operators on the same stage, the LIFT101 AT15--16 finite-order comparison applies to every polynomial in \nolinkurl{harmK} with the identity inverse matrix: its proof uses only the displayed conjugate-pair Laplacian equations, without a row-invariance, row-degree or coefficient-integrality premise. It therefore applies to HHa, \nolinkurl{harmL} in the literal \nolinkurl{ThShort}, and the full coefficient-field span of HHa. This comparison permits all arbitrary witnesses allowed by \nolinkurl{IsNab}; in particular it does not add a biframe-truncation axiom to an existing \nolinkurl{CovPin}. The source \nolinkurl{IsNab} at \nolinkurl{ThGerm.lean}209 forces every existentially supplied \nolinkurl{StarRing} to have star=H.bar, so the two pins have the same actual CM involution. The coherent converse of Theorem A is needed for the new prescribed extension, not imposed retrospectively on the stored operator.

For any fixed witness in \nolinkurl{ShortValAt}, keep all original data: the faithfully flat ring map, frame, complete z-enumeration and \nolinkurl{GeoInj}, each tp root, each original \nolinkurl{IsUUOver} section, the lattice Lambda, N, beta, w, finite-class masses and every phase. Each q,i summand has the same power k=N\^{}2 in all columns, so the harmonic comparison proves equality term by term after polynomial contraction. It leaves the prefactor

\begin{verbatim}
beta^-1 * mIdx^-1 * card(LevQ^n)^-1 * N^(-deg P)
\end{verbatim}

and the entire q/i/class sum unchanged. Consequently \nolinkurl{ShortValAt} is equivalent for the two operators on such inputs. Replacing only \nolinkurl{DeltaDat.nab2} by a different already-pinned operator preserves the corresponding \nolinkurl{IsShortRule} relation and the same \nolinkurl{shortGen} subset. This is conditional on both pins; it cannot justify replacing an unknown operator without its exact witness.

The actual \nolinkurl{shortGen} first argument \nolinkurl{_nab} is unused; the operative connection is \nolinkurl{DeltaDat.nab2}. \nolinkurl{CovPin} requires \nolinkurl{IsLev}, l2.A=1 and the latter connection's \nolinkurl{IsNab}. The new converse can supply that last conjunct from the prescribed exponential frame, while leaving the other two requirements untouched. \nolinkurl{ShortValTot} also requires existence of a section in \nolinkurl{sec}, which is not established merely by this termwise comparison. Nevertheless, if that law or a \nolinkurl{ShortInCusp}/\nolinkurl{ShortNe0} law is already supplied for one pinned operator, it transfers to the other: the complete defining relations and sets are equal. This transports an existing assertion; it does not prove any of these existence or nonzero assertions from scratch. Root/character laws and the original selected coupled-image nonvanishing remain separate.

\subsection{8. A prescribed complex short witness and the exact remaining admission}\label{a-prescribed-complex-short-witness-and-the-exact-remaining-admission}

For a literal \nolinkurl{ShortValAt} witness already given over \nolinkurl{CAt sigma}, each actual section F(q,i,b) has an analytic reduced theta germ theta(q,i,b;Z) in the frame of section5. The operator from Theorem C gives, without an extra scalar,

\[
 \sum_{m\in\operatorname{supp}P}P_m\prod_b
 \operatorname{nab}_y(N^2,F_{qib},m_b)
 =\left.P(\partial_{Z_1},\ldots,\partial_{Z_n})
       \prod_b\theta_{qib}(Z_b)\right|_{Z_1=\cdots=Z_n=0}. \tag{EC12}
\]

The source prefactor and weights are exactly

\[
 \beta^{-1}\,\operatorname{mIdx}^{-1}\,
 |\operatorname{LevQ}^n|^{-1}N^{-\deg P}
 \sum_q w(q)\sum_i (\text{EC12}). \tag{EC13}
\]

The actual theta-group actions defining each F, every permitted root, the complete z-enumeration with its \nolinkurl{GeoInj} condition, the coarse-class measure, and the original selected polynomial remain those of that witness. Formula (EC12) holds for arbitrary P; comparison with a different already-pinned connection in Theorem D additionally uses the actual harmonic condition. Neither formula changes the coefficient coevaluation or moves a fixed terminal branching map through a physical eigenframe change.

This is an evaluation of an existing witness in the genuine reduced gauge. It does not create the finite z/tp witness, descend a newly invented one, or identify it with a particular originally prescribed coupled theta construction without that construction's own literal dictionary. A general faithfully flat witness over a C-algebra need not admit a C-point; no such point is assumed to reduce it. Theorem D works directly over that actual algebra, while (EC12) is asserted for a complex-valued witness when one is supplied.

The strongest operator consequence is therefore precise. If the nonnegative part of a preselected \nolinkurl{Top.nab} or \nolinkurl{DeltaDat.nab2} is supplied at \textbf{every} quantified ring and point by the prescribed exponential gauges of Theorem A, its literal \nolinkurl{IsNab} is proved without replacing that operator. If only the genuine complex reduced operators are prescribed, Theorem C supplies an all-ring extension with exactly those complex values. It does not assert that an arbitrary previously stored extension equals this one on the other coefficient algebras. If an old extension was already pinned, Theorem D identifies its full harmonic short expressions with those of the constructed extension.

\subsection{9. The genuine coupled package already supplies the operative pin}\label{the-genuine-coupled-package-already-supplies-the-operative-pin}

\textbf{Theorem E --- exact comparison inside a supplied class package.} The preceding unknown-operator qualification is not an additional premise when the original package already supplies its literal pin. There are three directly read suppliers:

\begin{enumerate}
\def\labelenumi{\arabic{enumi}.}
\tightlist
\item
  \nolinkurl{HasDataCondition.lean}, \nolinkurl{LawsSh.dNab}, line60, gives \nolinkurl{IsNab} for each operative \nolinkurl{(T.dδ j n).nab₂} at every stated stage.
\item
  \nolinkurl{ThClassPins.lean}, \nolinkurl{OneMapTieP}, lines176--201, has the literal \nolinkurl{CovPin Xo c' dδ} conjunct at line199.
\item
  \nolinkurl{OneMapT.lean}, \nolinkurl{OMDat.Pins.cov}, line104, is exactly that \nolinkurl{CovPin} for the stored \nolinkurl{O.dδ}. Its \nolinkurl{OMPins} definition, lines143--144, includes this \nolinkurl{Pins} record along with the coefficient-field equality.
\end{enumerate}

These are fields of the supplied genuine records, not existence theorems constructing those records. Expanding \nolinkurl{ThClassPinsRec.lean}62--64 gives, at every tuple,

\[
 \operatorname{CovPin}(d_\delta)
 \Longrightarrow
 \operatorname{IsLev}(d_\delta(t)),\quad l_2.A=1,\quad
 \operatorname{IsNab}((d_\delta(t)).\mathrm{nab}_2). \tag{EC14}
\]

The corresponding harmonic premise is also literal. \nolinkurl{OMDat.Pins.io} at line102 supplies \nolinkurl{IoLFRec}; \nolinkurl{ThClassPinsRec.lean}45--47 unfolds this to \nolinkurl{IoLF}. The range clause in \nolinkurl{ThClass.lean}123--131 is

\[
 \operatorname{range}(io)
 =\operatorname{span}_F\bigl(\operatorname{HHa}(Dw,n,a)\bigr)
 \subseteq\operatorname{harmK}(Dw,n). \tag{EC15}
\]

The inclusion follows because \nolinkurl{harmK} is a submodule, and every HHa element belongs to it by \nolinkurl{RowDeg.lean}. It does not require an unproved claim that an arbitrary coefficient vector is integral or belongs to a particular lattice \nolinkurl{harmL}. Thus \textbf{every actual coefficient} io(x\_cl) used by the original class expression is harmonic. The field isomorphism \nolinkurl{dδ.eF} preserves the derivative equations coefficientwise; the labels and the identity inverse matrix in those equations are unchanged. It also carries the source \nolinkurl{embF} scalars through the literal equality of coefficient fields.

Now compare the stored operator, whose pin follows from (EC14), with the reduced-frame extension of Theorem C on that same stage. Apply Theorem D to \nolinkurl{thLaValAt}, which is exactly \nolinkurl{ShortValAt} with beta=1 (\nolinkurl{ThClass.lean}17--19). For every class, preserve the same lattice La, root power Nc, weight w, coefficient io(x\_cl), ring, point and scalar value. This gives equality of the class relations before the outer sum. The definition \nolinkurl{ThClsValAt} at lines37--41 then keeps exactly

\[
 \operatorname{evR}(\xi,y)
 \sum_{\mathrm{cl}}\operatorname{vol}\,
      \operatorname{gam}(\mathrm{cl})^{-1}
      \operatorname{ep}(\mathrm{cl})\,v_{\mathrm{cl}}. \tag{EC16}
\]

For the actual covering-stage record, \nolinkurl{IsThClsD} at lines184--193 uses these same individual values, transports coefficients through \nolinkurl{dδ.eF}, and multiplies by the original \nolinkurl{evR(dδ.cδ,y₂)}. Each quantifier, each original \nolinkurl{FrOver}, and the entire class sum is therefore unchanged. No choice of a new section or a new coefficient tensor is involved.

Finally the definition \nolinkurl{OneMapRow} at lines195--208 uses \nolinkurl{IsThClsD} in its existence and universal comparison clauses. Replacing only the operative nab₂ entries by the already-pinned comparison extensions leaves these relations identical; it preserves the same \nolinkurl{Cu}, \nolinkurl{thAd}, admissibility predicate Pl, central group element z, root-of-unity ze, phases ph and cc, \nolinkurl{lieFac}, actual uniformization and Haar-dependent dictionaries. The other conjuncts of \nolinkurl{OneMapTieP} involve unchanged data, and their \nolinkurl{CovPin} remains true because the new connections are pinned. At the level of this literal arithmetic class relation, no separate equality of the individual ordinary coefficients of the stored and reduced operators is required.

This is the main actual-package reduction. It does \textbf{not} assert existence of \nolinkurl{OMPins}, \nolinkurl{OneMapTieP}, \nolinkurl{LawsSh}, or the desired arithmetic source. Given such a package, its already-present pin plus range clause suffice to compare the original full harmonic class expression in the genuine reduced gauge. Any further fixed branching or selected derivative is still the original map applied to that same expression; this theorem neither enlarges its image nor proves a nonzero selected output. Where no source package supplies (EC14), the prescribed-operator qualification in section8 remains necessary.

\subsection{10. Unchanged data, costs, and source-law table}\label{unchanged-data-costs-and-source-law-table}

\begin{longtable}[]{@{}
  >{\raggedright\arraybackslash}p{(\columnwidth - 2\tabcolsep) * \real{0.5000}}
  >{\raggedright\arraybackslash}p{(\columnwidth - 2\tabcolsep) * \real{0.5000}}@{}}
\toprule\noalign{}
\begin{minipage}[b]{\linewidth}\raggedright
Item
\end{minipage} & \begin{minipage}[b]{\linewidth}\raggedright
Exact consequence
\end{minipage} \\
\midrule\noalign{}
\endhead
\bottomrule\noalign{}
\endlastfoot
Host, Datum, full integral CM action, Moduli, Level, Coeff & Kept literally; their existence is a premise. No new Host/Sizes instance is inferred. \\
Actual physical eigenframe and theta generator & Kept at the same actual FrPt; no move to a different point, finite component or selected family. \\
Polarization, cubical maps, global adjunctions and primitive & Same original maps, normalized by their prescribed rigidifications. \\
Coefficient field and any arithmetic model & No extension or model change is made by the formal converse. Complex coefficients are not asserted to descend or to be integral. \\
Ordinary coefficients & Exactly expVal, including the full product factorial; no added formal-degree or Haar factor. \\
Levels, theta-group lifts, N, mIdx, class masses, characters & Kept in (EC12)--(EC13) and in the termwise comparison. \\
Polynomial and terminal branching/derivative & Kept unchanged; harmonicity is used only for comparison of two pins. \\
\nolinkurl{nabPin} and \nolinkurl{dNab} & Supplied where the operator is defined by the verified frame family, or freely filled with the constructed extension; unknown preselected entries still require their actual witness. \\
\nolinkurl{CovPin} & Only its \nolinkurl{IsNab} conjunct is supplied by the converse. The original \nolinkurl{IsLev} and l2.A=1 conjuncts remain required. \\
\nolinkurl{ThValTot}, \nolinkurl{NabUniq}, \nolinkurl{thTot} & Follow from canonical LIFT101 for the verified pin. \\
\nolinkurl{ShortValAt}, \nolinkurl{IsShortRule}, \nolinkurl{shortGen} & Exactly preserved between two already-pinned connections with all other records fixed. \\
\nolinkurl{ShortValTot}, cusp and nonzero laws & An existing law transfers under that equality; existence and nonzero output are not obtained solely from pinning. \\
\end{longtable}

There is no numerical error charge and no change to any established finite remainder. The construction incurs no new field degree, level, isogeny degree, metric, period, denominator or cardinality factor. This last statement records unchanged data, not an effective bound for the still-unconstructed admitted source. There is no consequence for the DP44 sign, global central value, selected arithmetic content or RH without the remaining actual hypotheses.

\subsection{11. Literal source interface map and verification scope}\label{literal-source-interface-map-and-verification-scope}

The companion source record fingerprints each exact file and states whether its complete body or only the required locators were freshly read. The principal interfaces are:

\begin{longtable}[]{@{}
  >{\raggedright\arraybackslash}p{(\columnwidth - 2\tabcolsep) * \real{0.5000}}
  >{\raggedright\arraybackslash}p{(\columnwidth - 2\tabcolsep) * \real{0.5000}}@{}}
\toprule\noalign{}
\begin{minipage}[b]{\linewidth}\raggedright
File and declaration
\end{minipage} & \begin{minipage}[b]{\linewidth}\raggedright
Use
\end{minipage} \\
\midrule\noalign{}
\endhead
\bottomrule\noalign{}
\endlastfoot
\nolinkurl{IsExpChart.lean}: IsExpChart, IsExpGauge, IsExpGerm, expVal & Coherent all-order physical chart/frame, globally normalized beta, and exact ordinary factorial. \\
\nolinkurl{ThGerm.lean}: IsChart, IsFrame, IsPrim, ChartAt, IsNab & Every actual chart and biframe field, universal primitive and adjunction generator, full ring/point/section quantifiers. \\
\nolinkurl{CocycleRightZeroThd.lean}: FB.triv, IsCanonGauge, hL, nabVal & Multiplier-one biframe, h=0 and the literal positive gauge-correction sign. \\
\nolinkurl{ExpTie.lean}: FrTruncC, IsExpGaugeOf, ExpGaugeOfSt, TieSt & Coherence is supplied for the reference construction; no coherence is inferred for an arbitrary ChartAt. \\
\nolinkurl{ExpGauge.lean}: expGaugeOf\_of/holds; \nolinkurl{ESecToUnit.lean}: exists\_beta & Reference gauge and global normalized beta from actual primitive/adjunction data. \\
\nolinkurl{GaugeBeta.lean}: betaUniq\_of\_global & Global normalization equality, not equality inferred universally from a single point. \\
\nolinkurl{GaugeUniq.lean}, \nolinkurl{GaugeUniqPf.lean} & Existing accepted balanced-quadratic classification, reproduced in the exact all-ring scope needed for the converse. \\
\nolinkurl{IsThRule.lean}: NabV and actual theta expression & No ordinary-coefficient naturality axiom; the actual powers, colOf indexing and normalization are retained. \\
\nolinkurl{RowDeg.lean}: lap, harmK, HHa & Conjugate-pair Laplacians summed over equal-power columns. \\
\nolinkurl{ThShort.lean}: DeltaDat, harmL, ShortValAt, IsShortRule, shortGen & Fixed operative nab2, actual original witness data and complete weighted expression. \\
\nolinkurl{ThClassPinsRec.lean}: CovPin; \nolinkurl{HasDataCondition.lean}: LawsP/LawsSh & Exactly which named source-law conjunct is supplied and which remains. \\
\nolinkurl{ThClass.lean}: IoLF, thLaValAt, ThClsValAt, IsThClsD, OneMapRow & Actual harmonic coefficient range and complete class-sum transport in Theorem E. \\
\nolinkurl{ThClassPins.lean}: OneMapTieP; \nolinkurl{OneMapT.lean}: OMDat.Pins, OMPins & The literal supplying records for the operative pin and io range; their existence is not inferred. \\
\end{longtable}

Canonical LIFT101 AT3--17 is used with its full all-affine-base construction, finite-order arbitrary-pin comparison and mandatory locator addendum. Canonical K018's general Appell--Humbert calculation is compatible with section5. Its additional pure-J uniqueness and integral slope assertions are not extended to mixed types.

The diagnostics test exact coefficient recurrences, split-embedding separation, balanced-biframe identities, truncated logarithm/exponential over a ring with nilpotents, ordinary derivative factorials, equal-column harmonic contraction, and the Hermitian/semicharacter cancellation. They also fingerprint the exact source files. These finite computations do not replace the general geometric proof, certify source compilation, or evaluate the original coupled selected image.

\subsection{12. References and precise primary scope}\label{references-and-precise-primary-scope}

Bannai, K., \& Kobayashi, S. (2010). Algebraic theta functions and the p-adic interpolation of Eisenstein--Kronecker numbers. \emph{Duke Mathematical Journal, 153}(2), 229--295. Author version arXiv:math/0610163v4 (December11,2007), section1.2, Theorem1.7 and Definition1.8, printedp10, with the setup on printedp9. \url{https://arxiv.org/pdf/math/0610163} . The Appell--Humbert classification and displayed reduced multiplier were freshly read. No arithmetic integrality or ordinary-reduction theorem from that article is used here.

The Stacks Project Authors. (2026). \emph{The Stacks Project}. Lemma36.32.6, universal global functions of proper flat families with geometrically reduced connected fibres; Definition37.11.1 and Lemma37.11.7, affine formal smoothness; Lemma30.5.2, flat base change. \url{https://stacks.math.columbia.edu/tag/0BDM} ; \url{https://stacks.math.columbia.edu/tag/02GZ} ; \url{https://stacks.math.columbia.edu/tag/02KH} . These are carried through the checked LIFT101 applicability proof; this finding does not claim a new complete reading of those chapters. The complex-to-algebraic finite-jet identification is proved in section5 using smooth local coordinates.

The relative cubical/Poincare interpretation is the same one as canonical LIFT101. No new motive, l-adic weight, Dieudonne slope or global automorphic nonvanishing implication is attached to this prescribed-frame comparison.

\subsection{Canonical algebraic-geometric context}\label{canonical-algebraic-geometric-context-79}

The construction is over arbitrary characteristic-zero coefficient algebras and retains the given integral CM action, polarization, cubical maps, global adjunctions and primitive. A balanced quadratic change is realized by twisting the same formal Poincare biframe; it changes no abelian scheme, Picard point, source field, level or chosen arithmetic model. The ordinary factorial and the harmonic contraction are exact. Bannai and Kobayashi(2010), Algebraic theta functions and the p-adic interpolation of Eisenstein-Kronecker numbers, Duke Mathematical Journal153(2),229-295, doi10.1215/00127094-2010-024; author arXiv math/0610163v4 section1.2, Theorem1.7/Definition1.8, printedpp9-10, supplies the freshly read Appell-Humbert classification/multiplier statement. Universal global functions and formal smoothness are the checked LIFT101 inputs, Stacks Lemmas36.32.5-6(tag0BDM), Definition37.11.1/Lemma37.11.7(tag02GZ), Lemma30.5.2(tag02KH). Carried exact geometric context: SGA7 I, VII2.7/2.9 pp169-174 and VIII7.4(b) pp310-311; FGA232 Theorem3.1 pp148-149 and FGA236 Corollaries1.7/2.3 pp229/232; SGA3 I4.7.3/4.7.3.1; SGA6 VI Theorem3.3; cote4 scan003 sections8.1.1-8.1.2 and scans005-006 sections8.4.1-8.4.3; cote54 scans005-007. These contextual locators are carried from the documented earlier readings, not claimed as new whole-text checks. No new integral extension over Z, Dieudonne slope or Cartier-duality formula, lambda-ring positivity, K-theory filtration, l-adic weight theorem or CM-motive L-function comparison is obtained. DP\_MULT Q3e/DP44/DP52/DP54-65 and THETA3 G1/A7 gain only the stated source comparison; CX\_GRH CX-ALG-024 and GCT\_PLETH T-052 retain their global arithmetic gates.

\clearpage

\section{AM-K-121 --- exact coarea and paired narrow-band height}\label{am-k-121-exact-coarea-and-paired-narrow-band-height}

\textbf{Canonical checked integration --- 2026-11-23T11:48:20.355741+00:00.} Complete saved-proof and executable reviews and separate exact diagnostic replays passed. No numerical error charge is introduced. Frozen author history is preserved. AM-K-121 KP1-34 proves the exact complete coarea kernel and all-c paired-height bound for the original J56 selected density restricted to16/529\textless=Y\textless=16/527. Both real branches and both directions retain one common Beta expectation bounded by23479/23450; the paired factor is(yyprime)\^{}1/4*c\^{}-1/2, and the same kernel has height{[}1-(8/513)(c-1){]}\_+\^{}(333/2), supported at c\textless521/8. The original near-point and physical central affine lattice are transported by explicit additive cutoff broadenings2\^{}-340 and2\^{}-74, retaining the complete U\^{}f covolume, affine shift and750 cover. The full central integral has a genuine positive containing u-mixture after a globally paid exponential tangent; the compact profile itself is not asserted to be a positive Laplace transform. Every positive inverse moment is finite; the infinite zero degree still requires the actual-row reserve before enlargement. The same selected-spread coefficient is refined using every actual word/twist and S\_epsilon\textless97/(100B), with the entire old lower tail absorbed and the separate complete band-complement2\^{}-8605114 tail retained. The conservative coarea and spread field ceilings are-423 and-15, respectively; they are analytic factors and supply no sector charge alone. A complete original polynomial/moment/pattern/endpoint count remains required. No error charge; complete134/31 arrays, floors27/11 and k144..164 finite gate remain unchanged.

Author: \nolinkurl{/root/kummer_extension}. 23 November 2026 UTC. \textbf{Controlled pending mathematical proof.} This is a reusable analytic and original-affine-coset input, not a sector count. Independent reviews, executable diagnostics and any later arithmetic application are separate events. No canonical file or error array is modified.

\subsection{1. Actual selected object and statement}\label{actual-selected-object-and-statement}

Retain the original J56 selected vector-valued representation and its actual Hilbert metric, every binary word and canonical twist, both compact factors, the original fields, polarized projective modules, finite pencils and projectors, phases, class masses and Haar normalization. Let the selected radial probability density and the invariant ball measure be those of DO080 RT4--14. Set \[
 a=f=168,\quad\nu=a-\tfrac32=333/2,\quad p=333/4,
 \quad I=[16/529,16/527],\quad\delta=8/513,
\] \[
 \vartheta=\frac{4^{167}}{334\binom{334}{167}},\quad
 C_\beta=\frac{23479}{23450},\quad
 \kappa_I=\frac4{\sqrt{527}}C_\beta,\quad
 H(c)=1-\delta(c-1).
 \tag{KP1}
\] For the restricted nonnegative fidelity \[
 F_I(g)=\int_{Y(z),Y(g^{-1}z)\in I}
      \sqrt{\rho(z)\rho(g^{-1}z)}\,d\mu(z),
 \qquad Y(z)=1-|z|^2,
\] the theorem below proves, for every Cartan radius r\textgreater0, c=cosh r and s=sinh r, \[
 F_I(g)\le\kappa_I\,\vartheta2^{-933/2}
 \frac1{s\sqrt c}\left(\frac2{c+1}\right)^\nu
                      H(c)_+^\nu.
 \tag{KP2}
\] As always F\_I\textless=1, and at r=0 only that bound is used. The exact cutoff is c\textless521/8 on the positive retained kernel. The coarea factor and the compact height in KP2 are obtained simultaneously from one conditional kernel. They are not a product of unrelated coefficient estimates.

Sections 5--6 transport this factor through the actual near point and original physical central affine lattice by explicit additive broadenings of H. No relative logarithmic bound is asserted for H at its zero. Section 7 gives a containing family of positive shifted moment profiles and preserves the infinite constant-term obligation. The selected trace is never bounded by F\_I alone: the full K120/K118 arithmetic tail is retained in section 8.

\subsection{2. Exact coarea integral, including both directions}\label{exact-coarea-integral-including-both-directions}

Fix y,y' in (0,1), and write \[
 R_y=\sqrt{1-y},\quad D=\sqrt{y/y'},\quad z_1=u+it,
 \quad y'=\frac y{(c+su)^2+s^2t^2}.
\] The actual first-complex-coordinate marginal is \[
 \frac{a-1}{\pi R_y^2}
 \left(1-\frac{u^2+t^2}{R_y^2}\right)^{a-2}
                         {\bf1}_{u^2+t^2<R_y^2}.
 \tag{KP3}
\] On the original disk c+su\textgreater=c-sR\_y\textgreater0. Put \[
 v=c+su=\sqrt{D^2-s^2t^2},\qquad
 v_0=\frac{1+D^2+s^2y}{2c}>0,\qquad
 \Delta=D-v_0,
\] \[
 K=1-\frac{(c-D)^2}{s^2R_y^2},\qquad
 \Delta=\frac{s^2R_y^2K}{2c}.
 \tag{KP4}
\] The original disk condition is exactly v\textgreater v\_0; indeed \[
 R_y^2-u^2-t^2=\frac{2c(v-v_0)}{s^2}.
\] Thus if Delta\textless=0 the conditional density is zero almost everywhere. If Delta\textgreater0, its complete support has v\_0\textless v\textless=D. In particular 0\textless Delta\textless D. These identities hold whether c\textgreater D or c\textless=D. The calculation does not impose a first-case formula on an omitted second case.

The exact real Jacobian and the change from each sign of t are \[
 \left|\frac{\partial u}{\partial y'}\right|
          =\frac y{2s\,y'^2v},\qquad
 \left|\frac{dt}{dv}\right|
          =\frac v{s\sqrt{D^2-v^2}}.
\] The two signs cancel the first factor 1/2. The complete conditional transition is therefore \[
 K_r(y,y')=\frac{(a-1)y}{\pi R_y^2s^2y'^2}
 \left(\frac{2c}{s^2R_y^2}\right)^{a-2}
 \int_{v_0}^{D}\frac{(v-v_0)^{a-2}}{\sqrt{D^2-v^2}}\,dv.
 \tag{KP5}
\] There is no singularity omitted at t=0: this change has an integrable square-root endpoint.

Set tau=D-v, tau=Delta V and kappa=Delta/(2D). Then 0\textless kappa\textless1/2. Euler's positive Beta integral gives \[
 K_r(y,y')=\vartheta\frac{a-1}{\pi s}
       \frac y{R_y y'^2\sqrt{cD}}K^\nu
       E_a(\kappa),
\] \[
 E_a(\kappa)={\bf E}(1-\kappa V)^{-1/2},\qquad
                     V\sim\operatorname{Beta}(1/2,a-1).
 \tag{KP6}
\] Here B(1/2,a-1)=2 vartheta and sqrt(Delta/(2D))=s R\_y sqrt(K)/(2sqrt(cD)). These two factors are both present. Compared with K108 AE7's directional prefactor, the new factor is sqrt(y/(cD)) E\_a(kappa).

The function v-\textgreater(1-kappa v)\^{}(-1/2) is convex on {[}0,1{]}. Its chord majorizes it, and hence \[
 E_a(\kappa)\le1+\frac{(1-\kappa)^{-1/2}-1}{2a-1}
 <1+\frac{\sqrt2-1}{335}<\frac{23479}{23450}=C_\beta.
 \tag{KP7}
\] The final rational comparison uses sqrt2\textless99/70. This is a probability expectation, not a one-dimensional lattice theta estimate.

In the reverse direction D'=D\^{}(-1), y'=y/D\^{}2. Direct substitution gives v\_0'=v\_0/D\^{}2 and Delta'=Delta/D\^{}2. Thus kappa'=kappa, so the two directional expectations are exactly equal. Their geometric mean is one E\_a(kappa), not its square. The paired prefactor from KP6 relative to the old AE8 prefactor is \[
                      (yy')^{1/4}c^{-1/2}E_a(\kappa).
 \tag{KP8}
\]

\subsection{3. The exact paired height has a diagonal maximum on this band}\label{the-exact-paired-height-has-a-diagonal-maximum-on-this-band}

Where the two conditional heights are positive, define the normalized paired height \[
 \mathcal H=
 \sqrt{\frac K{2D/(c+1)}
             \frac {K'}{2D^{-1}/(c+1)}}.
\] Put u=sqrt(y/(1-y)), v=sqrt(y'/(1-y')), w=uv, and e=u/v+v/u-2. This w is a radial-odds variable, distinct from the physical central variable used below. A common-denominator calculation gives \[
 \mathcal H=
 \frac{2c\sqrt{(1+w)^2+we}-(2+e)-(c^2+1)w}{2(c-1)}.
 \tag{KP9}
\] The positive sign of the numerator follows from feasibility. For the actual band, \[
 m:=16/513\le w\le16/511=:M,\quad e\ge0,
 \quad e\le\frac{(w-m)^2}{wm},
 \quad\frac{w-m}{m(1+w)}\le\frac2{527}.
 \tag{KP10}
\] To prove the e inequality, write u=sqrt(m) U, v=sqrt(m) V with U,V\textgreater=1. The identity (U\^{}2-1)(V\^{}2-1)\textgreater=0 implies U/V+V/U\textless=UV+(UV)\^{}(-1). The final quotient in KP10 increases with w and equals2/527 at M.

Concavity of the square root gives \[
 \mathcal H\le1-\frac{(c-1)w}{2}
       +\frac e{2(c-1)}\left(\frac{cw}{1+w}-1\right).
 \tag{KP11}
\] For1\textless c\textless=2, the last bracket is negative because cw/(1+w)\textless=32/527\textless1. Thus mathcal H\textless=1-(c-1)m/2. For c\textgreater=2, discard the negative -1 in that bracket and apply KP10. The excess over1-(c-1)m/2 is at most \[
 \frac{w-m}{2}
 \left[-(c-1)+\frac{c(w-m)}{(c-1)m(1+w)}\right]\le0,
\] because (c-1)\^{}2/c\textgreater=1/2\textgreater2/527. This also includes w=m by continuity. Consequently the uniform, all-c estimate is \[
       0<\mathcal H\le1-\frac{8}{513}(c-1)=H(c)
 \tag{KP12}
\] on every positive feasible pair. If the right side is nonpositive, no positive feasible pair exists; zero extension gives the corresponding compact-support bound. This proves c\textless521/8 for positive retained density. There is no separate small-c exponential intercept.

\subsection{4. Detailed balance and the full radial spread}\label{detailed-balance-and-the-full-radial-spread}

On the symmetric restriction I times I the original disintegration and detailed balance give \[
 F_I(g)=\iint_{I\times I}\sqrt{p_Y(y)p_Y(y')}
                   \sqrt{K_r(y,y')K_r(y',y)}\,dy\,dy'.
 \tag{KP13}
\] Compact factors preserve Y and the original selected density. Combining KP6, KP8 and KP12 cancels D and D\^{}(-1), retains exactly H(c)\^{}nu, and multiplies the old AE8 radial kernel by (yy')\^{}(1/4)c\^{}(-1/2)E\_a(kappa). On the band (yy')\^{}(1/4)\textless=4/sqrt527. The remaining positive radial integral is at most the entire K107 SV8--12 spread integral, including its previously absorbed lower tail. Its exact coefficient is vartheta2\^{}(-933/2). This proves KP2. No replacement of the vector-valued metric, compact factors, or selected coefficient by a scalar lowest-type coefficient occurs.

\subsection{5. Near-point transport without a false cutoff elasticity}\label{near-point-transport-without-a-false-cutoff-elasticity}

For the actual nonidentity globally rank-two-unipotent original matrices write \[
 t=\alpha=(N^2)_{nn}>0,\quad a_t=1+t/2,\quad
 w=\Im N_{nn},\quad c(t,w)=\sqrt{a_t^2+w^2},\quad b(t)=1+t/4.
 \tag{KP14}
\] The original cutoff is non-strict \textbar\textbar h\_tau\textbar\textbar\textless=H0=133, so c\textless=H0 and t\textless=17424/133\textless132. The actual near-point Cartan radius satisfies r\_p\textless2\^{}(-352). DO080 and DO122, with its cutoff addendum, supply \[
 t_p\ge e^{-2r_p}t,\qquad |w_p-w|\le6H_0r_p,
 \qquad c_p^2=(1+t_p/2)^2+w_p^2.
 \tag{KP15}
\] The Cartan operator norm identity \textbar\textbar g\textbar\textbar=exp r and submultiplicativity give r(h\_p)\textgreater=r(h)-2r\_p.~If r(h)\textgreater=2r\_p, cosh(r(h)-2r\_p)\textgreater=exp(-2r\_p)cosh r(h). If r(h)\textless2r\_p, c\_p\textgreater=1\textgreater=exp(-2r\_p)c.~Hence in both cases \[
 c_p\ge e^{-2r_p}c,\quad c_p^{-1/2}\le e^{r_p}c^{-1/2},
 \quad H(c_p)\le H(c)+2\delta H_0r_p<H(c)+\eta_0,
 \quad\eta_0=2^{-340}.
 \tag{KP16}
\] This additive broadening is essential: log H has no bounded elasticity at its zero. No relative estimate for H is used.

Retain the old joint base profile \[
 J_0(t,w)=t^{-1/2}b(t)^{-167}
                (1+w^2/(4b(t)^2))^{-p}.
\] The exact Cartan inequalities give s\_p\^{}(-1){[}2/(c\_p+1){]}\^{}nu\textless=J\_0(t\_p,w\_p). DO122 CE9--10 yields \[
 J_0(t_p,w_p)c_p^{-1/2}
 \le e^{(336+3H_0p)r_p}J_0(t,w)c^{-1/2}.
 \tag{KP17}
\] The coefficient is below2\^{}16; therefore the entire168-embedding transport costs less than2. Combining the single-kernel KP2 with KP16--17 and vartheta\^{}f\textless2\^{}-649 gives \[
 \prod_\tau F_I(h_{p,\tau})
 <2\,\kappa_I^f2^{-649}2^{-933f/2}
 \prod_\tau J_0(t_\tau,w_\tau)c(t_\tau,w_\tau)^{-1/2}
                       [H(c(t_\tau,w_\tau))+\eta_0]_+^\nu.
 \tag{KP18}
\] The two source factors in the right side arose by factoring the same KP2 bound before transport. No second near-point factor or independent trace estimate is multiplied into KP18.

\subsection{6. The original physical central coset, with its height retained}\label{the-original-physical-central-coset-with-its-height-retained}

For the bounded original support27 all-II scalar-cross family x=N\_E D\_eta\textless2\^{}3328, fix eta, all167 original covector values and the common four-state G pattern J. The actual central values lie in one affine coset of the physical lattice Gamma\_w, with \[
 \operatorname{covol}\Gamma_w
   =U^f\sqrt{D_F}(N_FJ_0)xT^{-54},\quad
 \mu(\Gamma_w)<750\sqrt f\,U z_c,\quad
 z_c^f=(N_FJ_0)xT^{-54},\quad U z_c<2^{-83}.
 \tag{KP19}
\] Here the ideal J0 is distinct from the profile J\_0. These are exactly DO122 CE12--16 and the original primitive paired-entry extraction, including the full physical U\^{}f factor and the full-real750 cover. No larger row ideal is substituted for the central ideal.

For fixed t, c(t,w) is1-Lipschitz in the physical w. Every coordinate displacement inside a translated Voronoi cell is less than750sqrt(f)2\^{}(-83)\textless2\^{}(-69). Since delta\textless1/64, \[
 [H(c(t,w_{\rm centre}))+\eta_0]_+
 \le[H(c(t,w_{\rm cell}))+\eta_1]_+,
 \qquad\eta_1=2^{-74}.
 \tag{KP20}
\] Indeed the required broadening is less than2\^{}(-75)+2\^{}(-340)\textless2\^{}(-74). This comparison crosses the cutoff safely. For the other positive weight, \[
 |\partial_w\log(J_0(t,w)c(t,w)^{-1/2})|
 \le(p+1)/2.
\] Thus its whole-field Voronoi inflation is at most exp(375(p+1)f U z\_c)\textless2. Integrating over all disjoint translated cells proves the complete containing affine sum \[
 \sum_{\text{actual central fibre}}\prod_\tau F_I(h_{p,\tau})
 <4\kappa_I^f2^{-649}2^{-933f/2}
 \frac{T^{54}}{U^f\sqrt{D_F}(N_FJ_0)x}
                       \prod_\tau\mathcal I_{\eta_1}(t_\tau),
 \tag{KP21}
\] where \[
 \mathcal I_\eta(t)=t^{-1/2}b(t)^{-167}
 \int_{\mathbb R}(1+w^2/(4b(t)^2))^{-p}
       c(t,w)^{-1/2}[H(c(t,w))+\eta]_+^\nu\,dw.
 \tag{KP22}
\] The exact physical c-dependence remains in this integral. It is finite for every t\textgreater0 and has finite central support. The two factors2 in KP21 are precisely the complete near-point and complete Voronoi factors. The original real-part shift of d is retained in the affine coset. No new point of the enlarged analytic lattice is declared to be an original arrow.

The finite coefficient bound for a fixed actual pattern may multiply this inequality. Its complete pattern correlation, type and support sum, original row count and error union still have to be supplied by an application. KP21 alone is not an arithmetic sector removal.

\subsection{7. Positive containing moment profiles and the zero degree}\label{positive-containing-moment-profiles-and-the-zero-degree}

A simpler reusable majorant follows without discarding the original central covolume. Since c(t,w)\textgreater=a\_t, KP22 and the original positive Beta central integral give \[
 \mathcal I_\eta(t)\le\mathcal G\,Phi_\eta(t),\quad
 \mathcal G=4\sqrt{\pi/329},
\] \[
 \Phi_\eta(t)=t^{-1/2}(1+t/4)^{-166}(1+t/2)^{-1/2}
                       [1+\eta-\beta t]_+^\nu,
 \quad\beta=4/513.
 \tag{KP23}
\] This is a containing energy profile; KP22 is stronger and remains available. Put T\_eta=(1+eta)/beta. For any lambda\textgreater=nu/T\_eta, concavity of the logarithm at t0=T\_eta-nu/lambda gives, globally for t\textgreater0, \[
 [1+\eta-\beta t]_+^\nu
 \le\exp(C_{\lambda,\eta}-\lambda t),\qquad
 C_{\lambda,\eta}=\nu\log(\nu\beta/\lambda)
                         +\lambda T_\eta-\nu.
 \tag{KP24}
\] For t\textgreater=T\_eta the left side is zero. At lambda=nu/T\_eta the tangent is at zero. The complete field intercept is exp(f C\_lambda,eta); it is not free and need not be at most one.

Let B4 and B2 be independent Gamma probability variables of shape166,rate4 and shape1/2,rate2. Push forward their product with ds/sqrt(pi s) by z=lambda+B4+B2+s. The positive measure mu\_lambda is locally finite, has infinite total mass, and represents \[
 t^{-1/2}(1+t/4)^{-166}(1+t/2)^{-1/2}e^{-\lambda t}
                      =\int_0^\infty e^{-zt}\,d\mu_\lambda(z).
 \tag{KP25}
\] Tonelli and positive Gamma/Beta integrals give every integer positive inverse moment \[
 M_j(\lambda)=\frac1{\Gamma(j)}
 \int_0^\infty t^{j-3/2}(1+t/4)^{-166}(1+t/2)^{-1/2}
                                    e^{-\lambda t}\,dt<\infty,
 \quad j\ge1.
 \tag{KP26}
\] The integrability at zero requires j\textgreater1/2, and the positive shift handles infinity. If a future exact arithmetic certificate replaces exp(-lambda t) by a positive Gamma\_N/rate(N/lambda) majorant, it must prove all relevant moments and pay every intercept. No such finite table is asserted here.

The full physical central dependence also has a positive containing mixture; a compactly supported profile is \textbf{not} itself represented as a positive Laplace transform. Choose gamma\textgreater=nu*delta/(1+eta) and put \[
 C^c_{\gamma,\eta}=\nu\log(\nu\delta/\gamma)
          +\gamma\left(1+\frac{1+\eta}{\delta}\right)-\nu.
 \tag{KP27}
\] The exponential tangent majorizes {[}H(c)+eta{]}\_+\^{}nu by exp(C\^{}c\_gamma,eta-gamma\emph{c) for every c\textgreater=1. In the original physical central integral set w=2b(t)u, A=1+t/2, and L=sqrt(1+u\^{}2). The exact norm support inequality is \[
 c=\sqrt{A^2+(A+1)^2u^2}\ge A L+u^2/L,
 \quad c^2-(A L+u^2/L)^2=\frac{u^2}{1+u^2}.
 \tag{KP28}
\] It follows that KP22 is bounded by the positive shifted mixture \[
 2e^{C^c_{\gamma,\eta}}t^{-1/2}(1+t/4)^{-166}
 (1+t/2)^{-1/2}
 \int_{\mathbb R}(1+u^2)^{-p-1/4}
 e^{-\gamma(L+u^2/L)}e^{-(\gamma L/2)t}\,du.
 \tag{KP29}
\] Thus KP25--26 apply at each positive shift lambda(u)=gamma}L/2, with the entire positive u-weight and the complete tangent intercept retained. The u-measure is finite, every positive inverse moment is finite by lambda(u)\textgreater=gamma/2, and the full representing measure still has infinite total mass. Tonelli is applicable throughout. This is a genuine globally positive Gaussian mixture containing the central integral, not an illicit use of the compact profile's Mellin moments in a Poisson estimate. Its numerical moments and a complete sector application remain separate work.

The degree-zero term remains infinite. On each actual row, before Gaussian enlargement, keep the original norm floor N\_F(alpha)\textgreater=(Rk/q)\^{}f.~For theta in(0,1) and positive z\_tau, \[
 e^{-\theta\sum_\tau z_\tau\alpha_\tau}
 \le(\theta Rk/q)^{-f}\prod_\tau z_\tau^{-1}.
 \tag{KP30}
\] Integration uses M1\^{}f and converts precisely the zero-degree term to effective degree1. The remaining positive polynomial terms retain their individual inverse moments and all four transverse contributions. A new application must prove its full q/x monotonicities and finite-pattern sum; none is inferred from a dominant-term diagnostic.

\subsection{8. Probability, actual weights, and result boundary}\label{probability-actual-weights-and-result-boundary}

K120 proves the uniform actual selected probability outside I is below2\^{}(-40000000). Its complete K118 original containing mass and168-factor product argument apply unchanged to every actual subset S of the bounded support27 all-II globally rank-two-unipotent scalar-cross sector with nonzero original S-cross and x\textless2\^{}3328. Consequently \[
 \sum_{h\in S}|w(h)c(h)|
 \le\sum_{h\in S}|w(h)|\prod_\tau F_I(h_{p,\tau})
                                      +2^{-8605114}.
 \tag{KP31}
\] KP18 or KP21 may bound the retained product at their indicated scopes. The tail is an explicit term in a later complete union, not a standalone charge here. All original local weights and finite enumeration costs remain present. There is no new k144 removal, support-floor increment, arithmetic source construction, selected-image nonvanishing or signed positivity assertion.

The preliminary floating files \nolinkurl{exact_height_curvature_initial_probe.json} and the separate DO application diagnostic are not part of this theorem's proof or any numerical charge. They do not turn a failed sufficient bound into a lower bound or impossibility statement.

\subsection{9. A compatible full selected-spread refinement}\label{a-compatible-full-selected-spread-refinement}

This section refines the same radial coefficient; it does not multiply a second independent trace bound into KP2. Retain B=2\^{}955 and the literal ell=ell\_max from the actual J56 record in \nolinkurl{enlarged_gap_admissible_interval_certificate.json}. The actual parameters and simultaneous word/twist monotonicity are DO080 RT20--23, freshly used in K120 NB6--10. At the zero-word,q=0 maximum, \[
 \epsilon_0<\frac{32}{33B},\qquad
 \sum_{i=1}^{166}\epsilon_i<\frac{664\ell}{B^2},\qquad
 \epsilon_{167}=\frac{13}{(2\ell-2)(2\ell+12)}.
\] The two exact integer comparisons for the actual ell are \[
 664\cdot6600\ell<B,\qquad
 (2\ell-2)(2\ell+12)>85800 B.
\] Consequently every actual word and twist satisfies \[
 S_\epsilon<\frac1B\left(\frac{32}{33}
                     +\frac1{6600}+\frac1{6600}\right)
             =\frac{97}{100B}.
 \tag{KP32}
\] K120's exact mean interval gives 2111/64\textless M1\textless2113/64. The function x\^{}2/(x-1) increases for x\textgreater2, so the exact variance identity in RT21 yields \[
 \frac{\operatorname{Var}(Y^{-1})}{M_1-1}
 <\frac{97}{100B-97}\frac{2113^2}{64\cdot2049}
 =:V_*.
 \tag{KP33}
\] The lower radial tail proof K107 SV5--7 remains valid with its original full moment bound and tail2\^{}(-2800), because KP32 also implies the old S\_epsilon\textless1/B and the old mean interval remains true. Therefore its entire spread argument SV8--12 gives a squared coefficient \[
 C_{\rm rad}^2<\mathcal D_*:=
 334^2\frac{400}{399}V_*+528\cdot167^2\,2^{-2800}
                         <2^{-933}.
 \tag{KP34}
\] Every occurrence of2\^{}(-933/2) in the local bound can consequently be replaced by sqrt(mathcal D\_\emph{), and every full-field2\^{}(-933f/2) by mathcal D\_}\^{}(f/2). The exact rational relative factor is (2\^{}933 mathcal D\_*)\^{}84, with no floating ceiling required. This includes the full old lower tail inside the coefficient. It leaves the separately paid K120 band-complement tail in KP31 unchanged. A numerical application must use this replacement once only.

\subsection{10. Sources and inherited geometric context}\label{sources-and-inherited-geometric-context}

The actual selected density, invariant measure, radial conditional kernel, both compact factors and detailed balance are DO080 RT4--19. The old sphere marginal, exact coordinate conventions and comparison kernel are K108 AE3--9. The full spread constant, including its original absorbed tail, is K107 SV3--12 with the mandatory zero-radius/order-zero boundary addendum. R099 AH7--13 and R100 NB6--8 supply the exact narrow-band algebraic starting conventions; KP4--12 above rederive and strengthen the full coarea/paired-height step directly. The original Cartan identity and joint transport are DO121 EC3--8 and DO122 CE3--10 with its non-strict H133 cutoff addendum. The full physical affine lattice and covering radius are DO122 CE12--18 through the original DO104 CD14--17 and DO098 CF12--13. K120 NB1--16 and K118 BT15--22 provide precisely the separately paid band probability and complete containing tail. The actual-row degree-zero reserve is K116 RE9--16 / K119 RM19.

The ordinary positive Beta and Gamma integral normalizations are NIST DLMF5.12.1/3 and5.9.1, with strictly positive real arguments. The selected Hilbert metric remains Nakahama (2016), equation5.7 and Theorem5.1, printed960--961, through the accepted K043/K046/DO080 dictionary. No theorem-validity audit or Lean build of the source archive is performed.

All projective order modules, CM involution, polarization, cubical line, Picard/Poincare data and arithmetic periods are unchanged. Carried contextual anchors are SGA3 I4.7.3/4.7.3.1; SGA6 VI Theorem3.3; FGA232 Theorem3.1 pp148--149; FGA236 Corollaries1.7/2.3 pp229/232; SGA7 I VII2.7/2.9 pp169--174 and VIII7.4(b) pp310--311; cote4 scan003 sections8.1.1--8.1.2/scans005--006 sections8.4.1--8.4.3 and cote54 scans005--007. These are unchanged context, not new slope, duality, lambda-ring, weight or L-function claims. DP\_MULT Q3e/DP44, THETA3 G1/A7, CX\_GRH CX-ALG-024 and GCT\_PLETH T-052 retain their original arithmetic and global gates.

\subsection{Canonical algebraic-geometric context}\label{canonical-algebraic-geometric-context-80}

The theorem is a real-place estimate on the unchanged selected theta coefficient and its original arithmetic affine fibre. It retains the same polarized projective order modules, CM involution, saturated lattices, cubical/Poincare data, finite projectors, phases, Haar measures and arithmetic periods. NIST DLMF5.9.1 and5.12.1/3 are used with positive real Gamma/Beta arguments; Nakahama(2016), Norm computation and analytic continuation of vector valued holomorphic discrete series representations, Journal of Lie Theory26(4),927-990, equation5.7 and Theorem5.1, printed960-961, remains the actual Hilbert-metric input through K043/K046/DO080. Carried exact geometric context: SGA7 I, VII2.7/2.9 pp169-174 and VIII7.4(b) pp310-311; FGA232 Theorem3.1 pp148-149 and FGA236 Corollaries1.7/2.3 pp229/232; SGA3 I4.7.3/4.7.3.1; SGA6 VI Theorem3.3; cote4 scan003 sections8.1.1-8.1.2 and scans005-006 sections8.4.1-8.4.3; cote54 scans005-007. These contextual locators are carried from the documented earlier readings, not claimed as new whole-text checks. No scheme over Z or Picard/motive is reconstructed, and no new Dieudonne slope, Cartier-duality formula, lambda-ring/K-theory positivity, l-adic weight or CM-motive L-function theorem follows. DP\_MULT Q3e/DP44, THETA3 G1/A7, CX\_GRH CX-ALG-024 and GCT\_PLETH T-052 retain the same actual source, arithmetic image, period and global comparison gates.

\clearpage

\section{AM-DO-125 --- a closed positive central mixture at k144}\label{am-do-125-a-closed-positive-central-mixture-at-k144}

\textbf{Canonical checked integration --- 2026-11-23T11:58:50.948644+00:00.} Complete root and independent proof/script/interface reviews and separate exact replays passed. The actual canonical predecessor has134/31 entries; this integration gives135/31, complete floors27/11. Frozen author history is preserved. AM-DO-125 CM1--28 removes the complete original bounded support27/total-eligibility27 all-II full-positive scalar-cross globally rank-two-unipotent k144 cell, b=j\_int0 and2\^{}2111\textless x\textless2\^{}3328. Accepted K121 same-kernel coarea/height/variance is globally bounded by a paid exponential tangent before transport, then the exact physical central integral is bounded by a positive Gamma166/rate4, half-Gamma/rate2, half-Gamma/rate3514/29, infinite half-Gamma and shift29/20 convolution. Both complete quadrature tails and8064 intervals for degree147 yield a directed rational moment enclosure; finite-endpoint Hölder covers every lower positive degree and the actual-row reserve repairs degree0. All579 corrections, original supports/pivots/Green/T factors, full750 cover/U scaling and four finite small inflations remain. The full exponent is below−7622655/32768\textless−232.5, and the entire2\^{}-8605114 tail is paid once, giving one232 charge. The exact135/31 full error sums stay below2\^{}-230 with floors27/11. The sole remaining support27 necessary alternative is145\textless=k\textless=164 with the same denominator interval; the positive k145 envelope removes nothing. Signed PP3, higher-support remainder, actual source and selected arithmetic image, prescribed period, specific032 value, actual lift data/evaluatedDP44 and Hecke RH remain open.

Author: \nolinkurl{/root/definite_oscillator}, 23 November 2026 UTC. \textbf{Controlled pending complete proof.} K121 is canonically accepted; this application requires its own whole author, independent and root checks. No canonical file is changed. Numerical development files are preserved separately and are not substituted for the exact certificate.

\subsection{1. Original sector and statement}\label{original-sector-and-statement-1}

Retain the actual all-top J56 H133 normalized kernel, fields, CM involution, original polarized projective modules and source/target pencils, all finite projectors, phases, class masses, Haar normalization, binary words and canonical twists. The selected vector-valued Hilbert metric and both compact factors are unchanged. Put \[
 f=a=168,\quad H=133,\quad R=17424/133,\quad
 Q=1427908193729121522601,\quad U=77Q^2-14,
\] \[
 D_0=7^{140}29^{144}5^{84}13^{84},\quad D_F=2^{168}D_0,
 \quad p=333/4,\quad\nu=333/2,\quad\delta=8/513.
 \tag{CM1}
\] The exact predecessor is DO124, with all134 all-top and31 uniform charges and complete floors27/11. Consider its actual globally rank-two-unipotent scalar-cross support27 sector \[
 q_{\mathrm{total}}=27,\quad b_{\mathrm{inactive}}=j_{\mathrm{int}}=0,
 \quad 2^{2111}<x=N_ED_\eta<2^{3328},\quad k=144,
 \tag{CM2}
\] with all active typesII and every334 unmixed G reduction unipotent. The radial parameter q below is distinct from total eligibility.

\textbf{Theorem.} The complete original absolute normalized contribution of CM2 is less than \(2^{-232}\), including the entire narrow-band arithmetic tail once. Append one232 charge to the complete actual predecessor. Both full error sums remain below2\^{}-230; floors27/11 remain unchanged. The following k145 endpoint is a failed sufficient envelope and is not removed.

This theorem uses a globally positive Laplace mixture containing the exact physical central integral. It does not use the Mellin moments of a compactly supported height in a Gaussian or Poisson argument.

\subsection{2. Same actual kernel, full tail and exact variance}\label{same-actual-kernel-full-tail-and-exact-variance}

Let I={[}16/529,16/527{]} and let F\_I be the actual selected fidelity restricted to both radial variables in I. K120 NB1--16 and K118 BT15--22 give, for every actual subset S of the bounded sector, \[
 \sum_{h\in S}|a(h)c(h)|
 \le\sum_{h\in S}|a(h)|\prod_\tau F_I(h_{p,\tau})
                          +2^{-8605114}.
 \tag{CM3}
\] All168 factors and the complete original finite containing mass are already paid in this tail. The original selected coefficient is not bounded by F\_I alone.

Write \(h=I+N\), \(N^3=0\), rank N=2. In the accepted original normal form, \[
 t=\alpha=(N^2)_{nn}=UrJ_Z^{-1}r^*>0,\quad t_\tau\le R,
 \quad N_{nn}=-Ud,\quad\Re N_{nn}=t/2,
\] \[
 w=\Im N_{nn}=-U\Im d,\quad A(t)=1+t/2,
 \quad b(t)=1+t/4,\quad c(t,w)=\sqrt{A(t)^2+w^2}.
 \tag{CM4}
\] The actual original representative satisfies the non-strict cutoff \textbar\textbar h\_tau\textbar\textbar\textless=133. The mandatory DO122 cutoff addendum is retained.

Define the exact rational constants \[
 \vartheta=\frac{4^{167}}{334\binom{334}{167}},\quad
 C_\beta=23479/23450,\quad\kappa_I=4C_\beta/\sqrt{527},
\] \[
 V_* =\frac{97}{100B-97}\frac{2113^2}{64\cdot2049},\quad
 \mathcal D_*=334^2\frac{400}{399}V_*
             +528\cdot167^2\,2^{-2800},\qquad B=2^{955}.
 \tag{CM5}
\] K121 KP3--13 derives from the full two-real-coordinate coarea kernel and detailed balance the simultaneous factor \(\kappa_I c^{-1/2}[1-\delta(c-1)]_+^\nu\). Its KP32--34 derives the entire spread coefficient \(\sqrt{\mathcal D_*}\) from the actual all-word/twist epsilon sum and mean, including its absorbed lower radial tail. Accordingly, at the near point, that same kernel is bounded by \[
 \kappa_I\vartheta\sqrt{\mathcal D_*}\,
 J(t_p,w_p)c_p^{-1/2}[1-\delta(c_p-1)]_+^\nu,
\] \[
 J(t,w)=t^{-1/2}b(t)^{-167}
             (1+w^2/(4b(t)^2))^{-p}.
 \tag{CM6}
\] No two unrelated trace bounds are multiplied. Both actual compact factors are retained. K121's normalization includes both signs of the real coarea variable, one paired Beta expectation, and no missing factor2.

\subsection{3. A global exponential tangent before either transport}\label{a-global-exponential-tangent-before-either-transport}

Choose \(\gamma=29/10\). Since \(\gamma>\nu\delta\), concavity of the logarithm proves globally for every c\textgreater=1 \[
 [1-\delta(c-1)]_+^\nu\le e^{C-\gamma c},\quad
 C=\nu\log(\nu\delta/\gamma)+\gamma(1+\delta)/\delta-\nu.
 \tag{CM7}
\] The tangent point is \((1+\delta)/\delta-\nu/\gamma>1\). Outside the compact support the left side is zero. The whole intercept is retained. We apply CM7 at c\_p before transporting or extending any lattice; no relative logarithmic bound on the compact cutoff is needed.

K121 KP15--17 and DO122 CE8--10 give r\_p\textless2\^{}-352, \(t_p\ge e^{-2r_p}t\), \(|w_p-w|\le6Hr_p\), and \(c_p\ge e^{-2r_p}c\). Since original c\textless=H, \[
 J(t_p,w_p)c_p^{-1/2}e^{-\gamma c_p}
 \le e^{(335+3Hp+1+2\gamma H)r_p}
                       J(t,w)c^{-1/2}e^{-\gamma c}.
 \tag{CM8}
\] This uses the original monotonicity of J in t and its bounded w-gradient, the Cartan inequality for c, and \(1-e^{-2r_p}\le2r_p\). The complete product inflation has natural logarithm less than1/2048 and hence is less than \(2^{1/1024}\), because log2\textgreater1/2. This is the entire near-point cost, replacing the old2.

For each fixed row and actual G pattern, the original primitive paired-entry extraction gives one affine central coset with \[
 \operatorname{covol}\Gamma_w
 =U^f\sqrt{D_F}(N_FJ_0)xT^{-54},\quad
 \mu(\Gamma_w)<750\sqrt f\,U z_c,
\] \[
 z_c^f=(N_FJ_0)xT^{-54},\quad
 U z_c=\frac Rq\left(\frac{N_FJ_0}{N_EJ}\right)^{1/f}<2^{-83}.
 \tag{CM9}
\] These are the exact DO124 CB8--12 / DO122 CE12--16 interfaces. In particular the physical U\^{}f Jacobian, affine real-part shift, full-real750 covering constant and original central ideal remain present. The inequality follows from the literal host comparison \(U^fQ^{167}2^{3328+83f}<T^{54}\), simultaneously for every actual coupled row and pattern.

The global physical w-gradient of the transported positive weight satisfies \[
 |\partial_w\log[J(t,w)c(t,w)^{-1/2}e^{-\gamma c(t,w)}]|
 \le p/2+1/4+\gamma.
 \tag{CM10}
\] Indeed \textbar w\textbar/c\textless=1 and \textbar w\textbar/(2c²)\textless=1/(4A)\textless=1/4. Integrating over disjoint translated Voronoi cells bounds the actual central sum by its full real integral divided by CM9's covolume, with whole-field inflation at most \(\exp[750f(p/2+1/4+\gamma)2^{-83}]<2^{1/1024}\). The positive integral converges. Enlarged points are not asserted to satisfy any original arithmetic condition.

\subsection{4. Exact central normalization and the closed positive profile}\label{exact-central-normalization-and-the-closed-positive-profile}

Set \(w=2b(t)u\), \(L=\sqrt{1+u^2}\). The support-plane inequality is exact: \[
 c=\sqrt{A^2+(A+1)^2u^2}\ge AL+u^2/L,
 \quad c^2-(AL+u^2/L)^2=u^2/(1+u^2).
 \tag{CM11}
\] Also \(c^{-1/2}\le A^{-1/2}L^{-1/2}\). Consequently the central integral of CM8, including e\^{}C, is at most \[
 \mathcal G_1 t^{-1/2}b(t)^{-166}A(t)^{-1/2}
 \mathbb E\exp[C-\gamma(L+u^2/L)-\gamma Lt/2],
\] \[
 \mathcal G_1=2\int_{\mathbb R}(1+u^2)^{-167/2}du
 =\frac{2\cdot4^{83}82!83!}{166!}.
 \tag{CM12}
\] The expectation uses the probability normalization of this exact u-density. In particular the outside2 in CM12 is retained; the Beta normalization alone is half of \(\mathcal G_1\).

Under \(Z=u^2/(1+u^2)\), the probability law is Beta(1/2,83). The elementary inequalities \[
 L\ge1+Z/2,\qquad 2L-L^{-1}\ge1+3Z/2
 \tag{CM13}
\] follow by differentiating on0\textless=Z\textless1 or by their positive binomial series. Let S be an independent Gamma(167/2,rate1) variable. The standard Beta--Gamma change of variables gives SZ\textasciitilde Gamma(1/2,rate1). Conditional Jensen, applied to the convex exponential, gives for every beta\textgreater=0 \[
 \mathbb E e^{-\beta Z}
 \le\mathbb E e^{-\beta SZ/(167/2)}
 =(1+2\beta/167)^{-1/2}.
 \tag{CM14}
\] Thus the full expectation in CM12 is bounded by \[
 e^{C-\gamma}e^{-\gamma t/2}
                  [1+\gamma(t+6)/334]^{-1/2}.
\] For the exact chosen gamma, put \[
 \lambda=29/20,\quad r_*=3514/29,\quad
 D_*=(1480/1653)^{333/2}e^{1557/80}\sqrt{1670/1757},
\] \[
 \Phi(t)=t^{-1/2}(1+t/4)^{-166}(1+t/2)^{-1/2}
                  (1+t/r_*)^{-1/2}e^{-\lambda t}.
 \tag{CM15}
\] The original central integral is therefore at most \(\mathcal G_1D_*\Phi(t)\). Both the complete tangent intercept and the probability-normalization factor \(\sqrt{1670/1757}\) are explicit.

The function Phi is the Laplace transform of the positive convolution of Gamma(166,rate4), Gamma(1/2,rate2), Gamma(1/2,rate r\_*), the measure ds/sqrt(pi s), and the fixed shift lambda. The Gamma factors are probability normalized. The convolution is locally finite, has infinite total mass and finite inverse moments \[
 M_j=\frac1{\Gamma(j)}\int_0^\infty
 t^{j-3/2}(1+t/4)^{-166}(1+t/2)^{-1/2}
                 (1+t/r_*)^{-1/2}e^{-\lambda t}dt,
 \qquad j>1/2.
 \tag{CM16}
\] Positivity justifies Tonelli. This is a global containing profile, not the compact height itself. No finite M0 is asserted.

\subsection{5. Exact moment enclosure and all lower degrees}\label{exact-moment-enclosure-and-all-lower-degrees}

The certificate encloses M147 by ordinary positive integration with no floating arithmetic. For j147 (and separately j148 as the failed-next diagnostic), write the integrand in CM16 including1/(j−1)! as F\_j(t). On{[}1,64{]} use mesh h=1/128, giving8064 intervals. Its logarithmic derivatives are \[
 D_j(t)=\frac{j-3/2}{t}-\frac{166}{t+4}
       -\frac1{2(t+2)}-\frac1{2(t+r_*)}-\lambda,
\] \[
 D_j'(t)=-\frac{j-3/2}{t^2}+\frac{166}{(t+4)^2}
       +\frac1{2(t+2)^2}+\frac1{2(t+r_*)^2}.
 \tag{CM17}
\] For each mesh interval{[}a,b{]}, endpoint term bounds enclose D\_j; the sum of the absolute displayed derivative terms at a bounds \textbar D\_j'\textbar. An explicit positive Fmax is obtained by using b\^{}(j−1) in the numerator and a in every decreasing denominator and exponential. Therefore \[
 |F_j''(t)|\le F_{\max}\{
 \max(D_{\min}^2,D_{\max}^2)+|D_j'|_{\max}\}.
 \tag{CM18}
\] The absolute trapezoid error on that interval is at most h³ times CM18 divided by12, proved by integrating the linear-interpolation remainder. Every interval error is added with outward rounding.

Both complete tails are included. On{[}0,1{]} discard the factors bounded by1; on{[}64,infinity) use the powers of t in all three denominators. Since j\textless337/2, the resulting rigorous bounds are \[
 \int_0^1F_j(t)dt\le\frac1{(j-1/2)(j-1)!},
\] \[
 \int_{64}^\infty F_j(t)dt
 \le\frac{4^{166}\sqrt{2r_*}}{\lambda(j-1)!}
            64^{j-337/2}e^{-64\lambda}.
 \tag{CM19}
\] All quantities are computed with integers and rational numbers. Values are scaled by2\^{}700 and rounded outwards on a2\^{}-256 grid. Square roots use integer square-root inequalities. For e\^{}-x, reduce x/2\^{}m\textless=1/16, bracket the alternating series by degrees25 and24, round outwards and square m times with directed integer divisions. The output stores the complete rational lower and upper bounds, every aggregate trapezoid/tail/error component, precision and mesh. Finite models do not replace this continuum error proof.

Let \(\overline M_{147}\) be this exact upper bound. A deliberately loose \(\overline M_1\) is the accepted K120 first rational moment times355/113: pointwise, Phi is smaller than its shape166 shift32/25 profile, which is bounded by K120's positive Gamma512 majorant. Hölder on the positive representing measure now proves for every1\textless=j\textless=147 \[
 M_j\le\overline M_1^{(147-j)/146}
                 \overline M_{147}^{(j-1)/146}.
 \tag{CM20}
\] This interpolation is valid although the total measure is infinite, because both endpoint moments are finite. The same argument with148 is used solely in the failed-next diagnostic.

On a genuine row \(\prod\alpha_\tau\ge(Rk/q)^f\). Reserve theta=2\^{}-20 of the actual Gaussian exponent before enlargement. For the sole degree-zero term AM--GM and e\^{}-u\textless=u\^{}-1 give the full repair \[
 \int e^{-\theta\sum c_\tau\alpha_\tau}d\mu^{\otimes f}
 \le \overline M_1^f(q/(\theta Rk))^f.
 \tag{CM21}
\] No per-embedding floor is applied to subsequently enlarged row tuples.

\subsection{6. Complete original row, source and finite-pattern factors}\label{complete-original-row-source-and-finite-pattern-factors}

All geometric and reconstruction interfaces are exactly DO124 CB8--23, retained here explicitly. Every167 integral row value is in L=J I\_Sigma D\_eta at both CM partners. The three saturated transverse inputs have Gram determinant U and Gram\textgreater2\^{}36I. The actual covector metric has real covolume \(D_0^{k+3}z^{2(k+3)f}/U^f\), where \[
 z^{2f}=N_EJ\,xT^{-54},\quad q=R/(Uz^2)\ge k,
 \quad |a_G(h)|\le2Q^{-335\cdot167}(N_EJ)^{335/2},
\] \[
 N_EJ\le Q^{334},\qquad N_EJ\le Q^{167}N_FJ_0.
 \tag{CM22}
\] K113 bounded paired uniqueness fixes at most one eta per27-place support before row counting. The original same-label argument at both partners fixes the active label, including the central-only boundary; b=j\_int=0 makes all168 T factors constant on the full central fibre. Eta,r,d reconstruct the original matrix injectively. No beta, negative-column or source-isometry fibre is dropped.

Retain the K111 all-positive-parameter scalar coefficient \[
 K=\frac{336M_{15}}{c_{336}}(7841/5000)^{168}168!
 \left[\frac{(365729/500)^2}{168(314159/100000)^2}\right]^{168},
\] \[
 c_{336}=168\binom{336}{168}/4^{168},\qquad
 M_{15}=6015684926187167549292201834533348561577984/
                         2618017117694612446410293519.
 \tag{CM23}
\] The exact product-one anisotropic Gaussian argument retains all scalar and four transverse terms. Apart from \(\binom{164}{k}U^f\pi_+^{(k+3)f}/D_0^{k+3}\), the complete row polynomial is \[
 \sum_{\substack{0\le l\le k,\ 0\le h\le3\\j=l+3-h\ge1}}
 \binom kl\binom3h K^{k-l+h}2^{-36hf}
       [q/((1-\theta)R)]^{jf}\overline M_j^f
 +K^{k+3}2^{-108f}\overline M_1^f[q/(\theta Rk)]^f.
 \tag{CM24}
\] The row U\^{}f cancels the physical central U\^{}-f exactly. The original support count is the exact binomial168 choose27; retain all164 first-pivot choices and Green2. The complete T product is(T−1)\^{}-168, and the resulting T power is335·27+54−168=8931.

For each positive or effective degree j, after original central correlation, the q,x dependence is \(x^{-337/2}q^{(j-335/2)f}\). Its q derivative is negative, and at \[
 q_*^f=R^fT^{54}/(U^fQ^{334}x)
 \tag{CM25}
\] the coupled x derivative is exactly−1−j. The four-state pattern sum for each degree is \(\prod_{v=1}^{334}(1+q_v^{j-333/2})^2\). Here j\textless=147, so every exponent is below−1; the complete product is at most(1+1/g0)\^{}668. Its natural logarithm is below668/g0\textless1/2048. The entire T correction has logarithm at most168/(T−1)\textless1/2048. Together with CM8 and CM10, these are \textbf{four} original inflations, each less than2\^{}(1/1024). There is no hidden orientation or second pattern factor. Replacing their former conservative2,2,2,4 bounds is justified by these original formulas.

At the coupled lower endpoint x=2\^{}2111, the exact bracket is q\_-=251124964\textless q\_*\textless251124965=q\_+. The full579 nonmain terms in CM24 are checked using CM20. The scalar ratio \(K[(1-\theta)R/q_-]^f\) and the finite binomial/transverse coefficients give maximum relative dyadic exponent−150. The repaired zero term is included. Their entire sum is below the main upper term; the complete polynomial is bounded by \textbf{twice} that term. This conservative factor2 is retained.

\subsection{7. Directed full endpoint and entire tail payment}\label{directed-full-endpoint-and-entire-tail-payment}

Use \(\pi_+=355/113\), proved in the certificate by rational alternating arctangent intervals and Machin's identity. The identity follows from the tangent addition law for4 arctan(1/5)−arctan(1/239), whose angle is in(0,pi/2) and whose tangent is1. The exact intervals give pi\textless355/113; no decimal value of pi is used.

The full retained endpoint is bounded by \[
\begin{split}
 \mathcal B={}&4\,2^{4/1024}\binom{168}{27}164\,
 \mathcal D_*^{84}\vartheta^{168}
 (\kappa_I\mathcal G_1D_*)^{168}D_F^{-1/2}
 (R/U)^{335f/2}T^{8931}Q^{-336\cdot167}
 2^{-337\cdot2111/2}\\
 &\times\binom{164}{144}\pi_+^{147f}D_0^{-147}
 q_-^{-335f/2}[q_+/((1-\theta)R)]^{147f}
                         \overline M_{147}^{168}.
\end{split}
 \tag{CM26}
\] The initial4 consists exactly of Green2 and the complete-polynomial2. Every source, support, row, central and local factor precedes the exact evaluation. There is no old central ball volume or independent weighted-Jacobian estimate multiplied into this formula.

The executable uses 160-bit directed fixed-point logarithms, not floating point. For a positive rational, first extract its exact integral binary exponent; the remaining logarithm is evaluated using \(\log v=2\sum_{n\ge0}[(v-1)/(v+1)]^{2n+1}/(2n+1)\) on1\textless=v\textless=2. It retains100 terms and the full geometric tail, with outward divisions. The same certified log2 interval supplies1/log2 for the sole exponential intercept. Every sign is propagated. It proves \[
 \log_2\mathcal B<-7622655/32768<-232-1/2.
 \tag{CM27}
\] All numerical moment and logarithm intervals are stored as exact rational values. The historical floating probes are explicitly not proof inputs.

The entire original narrow-band tail is then paid once: \[
 \mathcal B+2^{-8605114}<2^{-232-1/2}+2^{-8605114}<2^{-232}.
 \tag{CM28}
\] For example2\^{}-1/2\textless3/4 and2\^{}-8605114\textless2\^{}-234 prove the last strict inequality. There is no additional endpoint or position union. The script independently preserves all134/31 ordered predecessor arrays, appends232 once and checks both complete sums below2\^{}-230. The previous arrays plus231 would fail; no231 charge is asserted.

The exact directed endpoint ceilings, including all nonmain terms, are:

\begin{longtable}[]{@{}
  >{\raggedleft\arraybackslash}p{(\columnwidth - 6\tabcolsep) * \real{0.2667}}
  >{\raggedleft\arraybackslash}p{(\columnwidth - 6\tabcolsep) * \real{0.2667}}
  >{\raggedleft\arraybackslash}p{(\columnwidth - 6\tabcolsep) * \real{0.2667}}
  >{\raggedright\arraybackslash}p{(\columnwidth - 6\tabcolsep) * \real{0.2000}}@{}}
\toprule\noalign{}
\begin{minipage}[b]{\linewidth}\raggedleft
Original nonzero S count k
\end{minipage} & \begin{minipage}[b]{\linewidth}\raggedleft
Positive degree
\end{minipage} & \begin{minipage}[b]{\linewidth}\raggedleft
Complete upper exponent
\end{minipage} & \begin{minipage}[b]{\linewidth}\raggedright
Scope
\end{minipage} \\
\midrule\noalign{}
\endhead
\bottomrule\noalign{}
\endlastfoot
144 & 147 & −7622655/32768 & Complete target cell with579 corrections \\
145 & 148 & 44737731/32768 & Failed next sufficient envelope with583 corrections \\
\end{longtable}

The separately evaluated k145, degree148 envelope is positive. It is an honest failure of this sufficient estimate, not a lower bound on its sector and not an impossibility statement. No k145+ or complete support28 claim follows.

\subsection{8. Source, verification and geometric boundaries}\label{source-verification-and-geometric-boundaries}

The primary new analytic input is accepted K121 KP1--34, especially the exact same-kernel coarea/height KP3--13, original Cartan transport KP15--17, physical support-plane positive mixture KP27--30 and variance KP32--34. The closed Beta--Gamma Jensen step and exact quadrature proof are supplied above. NIST DLMF5.9.1 and5.12.1/3 give the ordinary positive Gamma/Beta normalizations; the Beta--Gamma product law also follows directly by the Jacobian change from two independent Gamma variables. No asymptotic or analytically continued integral is used.

All original arithmetic factors are DO124 CB8--23, with DO122 CE3--19 and its mandatory non-strict cutoff addendum, DO098 CF4--16, DO104 CD10--19/section5, K113 SC27.1--10, DO105 EL9--16, DO108 AC8--10, K107 SV17--19, K110 PM10--12 and K111 HP1--18. K120 NB1--16 supplies the actual narrower probability and entire arithmetic tail, and NB17--24 supplies only the loose first moment used here. The source packet fingerprints each applied interface.

Nakahama (2016), equation5.7 and Theorem5.1, printed960--961, remains the actual selected Hilbert-metric dictionary through K043/K046/DO080. No scalar representation substitutes for it. Projective order modules, polarization, cubical theta line, Picard/Poincare data, CM involution and finite phases remain unchanged. The carried FGA/SGA/cote contextual anchors are inherited, not new primary readings or new formal-group, K-theoretic, weight or motivic comparisons. There is no new source admission, selected arithmetic-image nonvanishing, prescribed period, specific032 or evaluated DP44/RH conclusion. No Lean build or archive audit was performed.

\textbf{Verification boundary.} Exact executable diagnostics supplement the complete mathematical argument; neither finite coarea models nor quadrature samples alone establish it. Author whole read, independent whole review, root review and unchanged replays are distinct recorded events. Only the single complete CM2 sector is proposed for removal.

\subsection{Canonical algebraic-geometric context}\label{canonical-algebraic-geometric-context-81}

The exact positive profile treats the same original central affine module, real shift, CM partner incidence, polarization and selected Hilbert metric. The full central Jacobian and original projective row covolume are retained. Ordinary positive Gamma/Beta formulas are NIST DLMF5.9.1 and5.12.1/3; the Beta--Gamma/Jensen reduction and continuous exact quadrature are proved in CM11--20. Inherited context only: SGA3 I4.7.3/4.7.3.1; SGA6 VI Theorem3.3; FGA232 Theorem3.1 pp148-149; FGA236 Corollaries1.7/2.3 pp229/232; SGA7 I VII2.7/2.9 pp169-174 and VIII7.4(b) pp310-311; cote4 scan003 sections8.1.1-8.1.2 and scans005-006 sections8.4.1-8.4.3; cote54 scans005-007. No new formal-group/Dieudonne, lambda-ring/K-theory, l-adic-weight or CM-motive comparison. DP\_MULT Q3e/DP44 and THETA3 G1/A7 receive the strictly smaller conditional finite gate; CX\_GRH CX-ALG-024 and GCT\_PLETH T-052 retain their arithmetic/global gates.

\clearpage

\section{AM-LIFT-103 --- The all-ring short-value trace constructor}\label{am-lift-103-the-all-ring-short-value-trace-constructor}

\textbf{Canonical checked integration --- 2026-11-23T12:02:46.901586+00:00.} Complete saved-proof and executable reviews and separate exact diagnostic replays passed. No numerical error charge is introduced. Frozen author history is preserved. AM-LIFT-103 ST1--19 constructs the entire first all-ring existence/uniqueness conjunct of literal ShortValTot from the actual geometric records, Lambda.Unimod/TrivAt(M)/LevOK(N,M), an operative IsNab pin and P in identity harmK. The canonical finite-etale dual-row kernel has exactly the source zCard, and every GeoInj witness is a full enumeration. Allowed half-primary root differences lie in that subgroup by TrivAt(M); unique normalized odd theta lifts give the exact complete section-trace identity with the original U\_t U\_z order. Harmonic contraction, a finite-etale construction cover, flat H0/eigenbasis base change and faithfully flat element descent produce the unique value over every coefficient algebra, retaining nilpotents and all original beta/mIdx/cardinality/N-degree/weight factors. Every original faithfully flat witness agrees. Arbitrary-map naturality compares supplied actual FrMap points and does not assume FrPt.gen under nonflat pullback. Applying the theorem to the actual cover-stage data removes the short-value scalar constructor from the conditional source gate. The second original-base sec/cdelta/IsShortRule conjunct remains open, as do source/package existence, selected arithmetic content/nonzero, cusp/spanning laws, period, specific032 value and evaluatedDP44/RH. No numerical error charge;135/31 complete arrays and floors27/11 unchanged.

\textbf{Status: final author subject, 2026-11-23; independent and root reviews pending.} The complete author/source review and 131,640 exact diagnostics are recorded in the adjacent final verification packet, which identifies the frozen bytes and actual completion times. This concerns precisely the first conjunct of the literal \nolinkurl{ShortValTot}; the existence of a rule in \nolinkurl{sec} on the original base, with its prescribed \nolinkurl{DeltaDat.cδ} descent, is a separate requirement.

\subsection{1. Statement with the actual source quantifiers}\label{statement-with-the-actual-source-quantifiers}

Retain an actual Host, Datum, odd Level, Coeff and Moduli from the supplied source. Write O=O\_E, L₀=Dw.La, r=Dw.p+Dw.q, {[}E:Q{]}=2f and g=fr. The original alternating form EF on L₀ is perfect and satisfies EF(ax,y)=EF(x,bar(a)y). The actual object has its original symmetric rigidified principal line, cubical maps, integral O-action and geometric \nolinkurl{TateIso}. No new moduli object or universal Tate trivialization is assumed.

Let Λ be the literal \nolinkurl{GenLat E n} and suppose

\[
 \Lambda=\Lambda^{\#},\qquad
 \Lambda.\operatorname{TrivAt}(M),\qquad
 \Lambda.\operatorname{LevOK}(N,M),\qquad M=l.M. \tag{ST1}
\]

Here \# is the O-valued Hermitian dual for the original identity form on Eⁿ; \nolinkurl{LevOK} means N is odd, M divides N, and NOⁿ lies in Λ∩Oⁿ. Let \nolinkurl{nab} satisfy literal \nolinkurl{IsNab}, with its actual conjugation. Let w be any original coefficient-field-valued weight on \nolinkurl{(Fin n → LevQ Dw l)}, let β be any coefficient-field scalar, and let P belong to the identity-matrix \nolinkurl{harmK} for these rows and n columns. In particular P may be any actual \nolinkurl{harmL} or any coefficient in the \nolinkurl{IoLF} image.

\textbf{Theorem ST2.} For every commutative c.F-algebra S, including nonreduced and non-Noetherian algebras, and every actual y∈FrPt(Xo,c;S),

\[
 \exists!v\in S:\quad
 \operatorname{ShortValAt}(X_o,c,\mathrm{nab},w,\Lambda,\beta,N,P,y,v).
 \tag{ST2}
\]

The value commutes with arbitrary coefficient-algebra maps between supplied actual points related by \nolinkurl{FrMap}. Every witness permitted by the original existential definition computes this same value. No particular choice of enumeration, allowed roots, normalized multiplication generator or symmetric lift is imposed on an existing witness. This naturality statement does not assert existence of a base-changed \nolinkurl{FrPt.gen} under every nonflat map.

Thus (ST1), an actual pin, and the original harmonic condition construct the entire all-ring existence-and-uniqueness conjunct in \nolinkurl{ThShort.lean}287--293. Applying this at the cover stage uses exactly dδ.Xo₂,c₂,nab₂,w(x₀),eF(β),l₂ and the mapped P. The extra \nolinkurl{TrivAt bad} and \nolinkurl{AdmX} premises of \nolinkurl{ShortValTot} can be retained; they are not needed for this first conjunct. This theorem does not construct the subsequent f∈sec or its \nolinkurl{IsShortRule} relation.

The cases n=0, r=0 and S=0 are included. Empty products and the rank-zero torsion group have their literal singleton meanings. For the zero ring every scalar equals zero; all the scheme statements below are statements over its empty spectrum.

\subsection{2. The original Hermitian module is a genuine lattice}\label{the-original-hermitian-module-is-a-genuine-lattice}

First, the definition of \nolinkurl{GenLat.Unimod} does imply that Λ is a full finitely generated projective O-lattice; this is not added silently to the definition.

Let W=EΛ. If W were proper, its orthogonal complement for the identity Hermitian form would contain nonzero v. This form is positive definite at every complex place. Every cv, c∈E, pairs trivially with Λ, hence belongs to Λ by selfduality. In particular h(v,v)/m² would be integral for every positive integer m. Since h(v,v) is a nonzero algebraic number, this is impossible. Hence W=Eⁿ. Choose an E-basis from Λ. Its O-span L is a lattice, and Λ⊂L\#; since O is Noetherian, Λ is finitely generated. It is torsion-free and therefore projective over the Dedekind domain O.

The inclusion NOⁿ⊂Λ and selfduality give

\[
 NO^n\subseteq\Lambda\subseteq N^{-1}O^n. \tag{ST3}
\]

For example, h(v,Ne\_j)=Nbar(v\_j) is integral for v∈Λ, proving the second inclusion. Therefore all quotients and additive indices appearing in \nolinkurl{mIdx} and \nolinkurl{zCard} are finite and positive. There is no default infinite-index interpretation in this application.

Define the ordinary row dual

\[
 \Lambda^\vee=\{\mu\in(E^n)^*: \mu\Lambda\subseteq O\}.
 \tag{ST4}
\]

It is a finite projective lattice. Every μ∈Λ∨ has Nμ∈Oⁿ, by NOⁿ⊂Λ. The source's quantification over all matrices in \nolinkurl{InZh} is equivalent to imposing the single-row equations for all μ∈Λ∨: a row can be placed in a matrix with all other rows zero, and every matrix in the source condition consists of such rows. A finite O-generating set of Λ∨ suffices.

Selfduality also gives norm(det Λ)=1. Indeed det(Λ\#)=bar(det Λ)⁻¹ for the identity Hermitian form; take absolute ideal norms. Consequently

\[
 |\Lambda/NO^n|=N^{2fn},\qquad
 zCard(\Lambda,N,r)=|\Lambda/NO^n|^r=N^{2gn}. \tag{ST5}
\]

One may verify the index formula locally. At conjugate pairs the determinant valuations are opposite and the residue cardinalities coincide; at a conjugation-fixed prime the determinant valuation is zero. This uses no choice of a free basis of Λ and no class-number factor. The exact source index formula, rather than its simplified value, is retained in every normalization below.

\subsection{3. A canonical finite étale kernel, not a chosen Tate tensor}\label{a-canonical-finite-uxe9tale-kernel-not-a-chosen-tate-tensor}

Pull the original universal object to y and denote its abelian scheme by A/S. Properness and the semiabelian structure make it an abelian scheme. For every positive integer k, {[}k{]} is finite étale over this characteristic-zero base. This also follows directly: its differential is k times the identity, so it is étale; it is proper and has finite geometric fibres, so it is finite. Its fibres have k\^{}(2g) points. The étale and proper statements are stable under arbitrary base change; no reducedness of S is used.

Define a subgroup scheme of A{[}N²{]}ⁿ by the actual integral O-endomorphisms

\[
 \mathcal H_\Lambda
 =\bigcap_{\mu\in\Lambda^\vee}
       \ker\bigl(N\mu:A[N^2]^n\longrightarrow A[N^2]\bigr).
 \tag{ST6}
\]

Only finitely many μ are needed by section2. Kernels and finite intersections in the category of finite étale group schemes are again finite étale. Equivalently these are kernels of finite locally constant sheaves. This definition is canonical over S and commutes with arbitrary base change. It is precisely the predicate \nolinkurl{InZh}: the matching integral cm entries in that definition are the uniquely determined coefficients Nμ.

Use the actual \nolinkurl{Obj.geom} TateIso only at a geometric point to compute the fibre, not to define (ST6) over S. It identifies A{[}N²{]} with L₀/N²L₀, with the original O-action and EF pairing. Finite projectivity of L₀ gives

\[
 (\mathcal H_\Lambda)_{\bar s}
 \simeq
 \frac{N\Lambda\otimes_O L_0}{N^2O^n\otimes_O L_0}
 \simeq(\Lambda/NO^n)\otimes_O L_0. \tag{ST7}
\]

For the first equality, the equations Nμ(x)=0 modulo N²L₀ are the double-dual characterization of NΛ. More explicitly, at an O-prime write N as a unit times π\^{}a and write Λ=U diag(π\^{}\{e\_i\})Oⁿ, with U integral invertible. Inclusions (ST3) give −a≤e\_i≤a. The dual-row equations say π\^{}\{a−e\_i\}(U⁻¹x)\_i=0 modulo π\^{}\{2a\}, equivalently (U⁻¹x)\_i∈π\^{}\{a+e\_i\}O modulo π\^{}\{2a\}. This is exactly NΛ/N²Oⁿ. The local module L₀ is free of rank r, so the calculation repeats in each of its r coordinates. This local basis is used only for the rank/kernel verification; no basis is chosen in the definition over S. The inclusion NΛ⊂Oⁿ in (ST3) makes (ST7) an actual subgroup of the displayed torsion quotient.

Every localization of L₀ at a maximal ideal is free of rank r. Hence a finite O-module C satisfies \textbar C⊗L₀\textbar=\textbar C\textbar\^{}r, by its primary decomposition. Thus the rank of (ST6) is exactly the source \nolinkurl{zCard}, constant on S. There is no canonical global identification (ST7) asserted: different Tate trivializations give the same canonical kernel and the same rank.

The subgroup is isotropic for the commutator pairing of the product line L\^{}\{N²\} on Aⁿ. In (ST7), the pairing of Nλ⊗u and Nμ⊗v is trivial because the numerator is N² times

\[
 \sum_j EF_E(\lambda_j u,\mu_j v)
   =EF_E\!\left(u,\left(\sum_j\bar\lambda_j\mu_j\right)v\right)\in\mathbb Z.
 \tag{ST8}
\]

Here EF\_E is the rational extension of the actual alternating form; the identity is its actual adjoint law. Hermitian selfduality makes the scalar sum an element of O. Since (ST5) is the square root of the order of A{[}N²{]}ⁿ, the subgroup is maximal isotropic. Equality of this pairing to one over S follows from equality of morphisms of finite étale sheaves, checked on their geometric stalks. This geometric-stalk argument is used for discrete torsion data only; it will not be used to compare scalar values in a nonreduced ring.

\subsection{4. The actual half-level and all its allowed roots}\label{the-actual-half-level-and-all-its-allowed-roots}

Put b=l.Mb. The literal level conditions give b bar(b)=(M) and b+bar(b)=O. The original η induces

\[
 L_0/bL_0\ \simeq\ A[b], \tag{ST9}
\]

as finite étale sheaves. The proof is the same exact Chinese-remainder argument as LIFT100 HN4--6: \nolinkurl{Obj.geom} specifies η on bar(b)L₀, \nolinkurl{eta_ker} kills bL₀, and their sum is L₀. Geometric verification of the resulting finite étale isomorphism suffices over every S. Therefore the actual \nolinkurl{levPt(levRep(q))} is killed by M and belongs to the b-primary half.

For each source label q and column, consider the finite étale scheme of points t satisfying

\[
 Nt=\eta(\operatorname{levRep}(q)),\qquad
 t\text{ is killed by a power of }b. \tag{ST10}
\]

It is nonempty on every geometric fibre and thus is a finite étale surjection onto S. To see this without a principal generator for b, use the primary decomposition of A{[}N²{]}. Since M divides N, multiplication by N on its b-primary part maps onto the b-primary part of A{[}N{]}. The specified η point belongs to this image. A point in (ST10) automatically has N²t=0. The seemingly unbounded power of b in the source condition is therefore equivalent here to membership in this fixed finite primary summand.

Any two allowed roots t,t' differ by

\[
 \delta=t'-t\in A[N]_{b\text{-primary}}. \tag{ST11}
\]

\nolinkurl{TrivAt M} is used at exactly this point. Its actual denominator k is coprime to M and gives Λ\_p=O\_pⁿ for every p dividing M. Consequently the primary part of (ST6) at every prime over b is A{[}N{]}ⁿ. Thus each tuple of differences (ST11) belongs to H\_Λ.

The entire b-primary summand of A{[}N²{]} is isotropic. Choose an ideal I supported on the primes of b that kills that summand. I and bar(I) are coprime. With u∈I, v∈bar(I), u+v=1, one has uy=0 and bar(v)x=0. The actual balanced commutator identity then gives e(x,y)=e(x,vy)=e(bar(v)x,y)=1 for x,y in this summand. In particular each pair t,δ in (ST10)--(ST11) commutes in the theta group. No phase is discarded.

\subsection{\texorpdfstring{5. Every permitted \nolinkurl{IsUUOver} section has one normalization}{5. Every permitted  section has one normalization}}\label{every-permitted-isuuover-section-has-one-normalization}

Construct the normalized multiplication isomorphism on the original universal abelian scheme before pulling it to S:

\[
 \alpha_N:[N]^*L\overset\sim\longrightarrow L^{N^2}. \tag{ST12}
\]

The original cubical identity and normalized symmetry give it by the same induction as LIFT100 HN7. Its origin restriction is the original rigidification to the power N²−1, exactly \nolinkurl{ThetaVV.IsNormMul}103--108. Thus α₀ is a global generator on the source's universal group, as required in \nolinkurl{IsUUOver}; it is not a substitute defined only on one fibre. Any two such normalized generators have a global base-unit ratio, and origin normalization makes it one. This uses proper universal global functions, as in LIFT101. The required pulled generator α is then forced by the literal \nolinkurl{LB.Gen.Corr} clause.

For q₀=N² odd, every point of A{[}q₀{]} has a unique theta lift U\_x with U\_x\^{}\{q₀\}=1 and symmetry ι(U\_x)=U\_x⁻¹. The original principal polarization identifies the kernel of the polarization of L\^{}\{q₀\} with A{[}q₀{]}, so its theta-group fibre over each such point is a G\_m-torsor. This is exactly the situation of \nolinkurl{ThetaPt.SymLift}44--50. The assertion holds over arbitrary S: locally for the fppf topology choose any theta lift, rescale by a square root to arrange the symmetry relation, and then rescale by its q₀-th power, which lies in μ₂, to arrange order q₀. Two normalized lifts differ by a unit z with z²=z\^{}\{q₀\}=1; Bezout for 2 and q₀ gives z=1 over every ring. Uniqueness descends the lift and makes its formation commute with base change. The theta-group involution is unchanged by a scalar rescaling of the symmetric line isomorphism.

For clarity, the initial local lift is a G\_m-torsor and the square-root cover is faithfully flat because 2 is invertible. These are construction covers; they introduce no ambiguity after unique descent. The assertion is not deduced by ignoring nilpotents in S. On an isotropic subgroup the normalized lifts form a group, since commuting products satisfy both defining identities and uniqueness applies.

Let s be the actual y.s and set s\_N=α\_N({[}N{]}\^{}*s). The literal \nolinkurl{IsTwistPull} fixes this section, including its orientation. For any permitted z,t, every witness to \nolinkurl{IsUUOver} gives precisely

\[
 F(z,t)=U_t U_z s_N. \tag{ST13}
\]

Indeed α₀, its pullback and both symmetric lifts have just been shown unique; each \nolinkurl{IsTwistPull}/\nolinkurl{IsActOf} relation determines its section by equalities on an open cover. These identities apply to every existentially allowed witness, even when each summand initially supplies its own α₀ or symmetry generator. There is no free scalar, chosen character or additional sign left in (ST13).

\subsection{6. The original complete trace is independent of all root choices}\label{the-original-complete-trace-is-independent-of-all-root-choices}

Use the product line and external product sections on Aⁿ. For a tuple t=(t\_b), define the section

\[
 \mathcal T_t
 =U_t\sum_{z\in\mathcal H_\Lambda}U_z\,(s_N\boxtimes\cdots\boxtimes s_N).
 \tag{ST14}
\]

The sum is the finite étale trace of this family of sections. It may be constructed on a finite étale cover enumerating H\_Λ; different enumerations give exactly the same sum, so effective descent of sections defines it independently of that cover.

For an arbitrary allowed \nolinkurl{z} in \nolinkurl{ShortValAt}, \nolinkurl{InZh} places its entries in H\_Λ. \nolinkurl{GeoInj} makes the induced map from the constant scheme with zCard labels into H\_Λ injective on every geometric fibre. Both finite étale schemes have the same rank. This is therefore an isomorphism. In particular \textbf{every allowed z is the full enumeration}, not just the enumeration constructed for existence. Over a disconnected base its permutation relative to another enumeration can vary by an open-and-closed partition; equality of the complete sum still holds.

If t' is another allowed tuple, δ=t'−t belongs to H\_Λ by section4. Both t and δ lie in the same isotropic b-primary half, so U\_(t+δ)=U\_t U\_δ. The lifts on H\_Λ form a group, and translation by δ permutes it. Consequently, as an exact section identity over S,

\[
 U_{t'}\sum_{z\in\mathcal H_\Lambda}U_z
 =U_t\sum_{z\in\mathcal H_\Lambda}U_z. \tag{ST15}
\]

This retains the source ordering U\_t U\_z. It proves root independence before differentiating or specializing to fields. No unexplained commutator factor is absorbed into β or w. The actual \nolinkurl{TrivAt M} and half-primary conditions are necessary for this argument; arbitrary N-th roots are not substituted for the permitted ones.

\subsection{7. Harmonic contraction and the exact all-ring scalar}\label{harmonic-contraction-and-the-exact-all-ring-scalar}

For a fixed actual point y and a fixed positive k, \nolinkurl{IsNab} makes every operator on a section linear over its coefficient algebra: its common chart assigns the section's germ in one line frame and applies a fixed linear differential operator. It follows that the source expression

\[
 \mathcal D_{P,y}\Bigl(\boxtimes_b F_b\Bigr)
 =\sum_{m\in\operatorname{supp}P}P_m\prod_b
       \mathrm{nab}_y(N^2,F_b,\operatorname{colOf}(m,b)) \tag{ST16}
\]

is multilinear. It extends linearly to the finite sum of external product sections in (ST14).

There is also a coordinate-free way to see that the extension depends on the actual product section, rather than on a presentation as a sum of pure tensors: choose the common \nolinkurl{IsNab} chart, multiply the germs on the n factors and apply the resulting finite-order differential functional at the origin of Aⁿ. Equation (ST16) is exactly its expansion. Thus every equality of sections in (ST14)--(ST15) gives the same scalar.

Although ordinary coefficients of a ringwise chosen pin need not be natural, this harmonic functional is natural under arbitrary ring maps and actual \nolinkurl{FrMap}. Pull the finite-order chart of the first point to the new base, and compare it with the chart of the new pin. The LIFT101 AT15--16 argument, made explicit for all \nolinkurl{harmK} in LIFT102 Theorem D, annihilates the balanced quadratic gauge change because every column uses the same k=N². It works for arbitrary actual sections of L\^{}\{N²\}, not only the special sections used in \nolinkurl{ThValAt}. Only finitely many jets, up to deg P, are used. No new coherence axiom is imposed on the arbitrary stored pin.

On a finite étale cover R of S where the roots and enumeration exist, compute exactly the source expression

\[
 v_R=
 \beta^{-1}\,mIdx(\Lambda,N,r)^{-1}
 |\operatorname{LevQ}^n|^{-1}N^{-\deg P}
 \sum_q w(q)\,\mathcal D_{P,y_R}(\mathcal T_{t(q)}).
 \tag{ST17}
\]

All coefficient-field scalars are mapped into R. Every cardinality and normalization here is literally the original one; in particular mIdx is not replaced by zCard. The β inverse is the source field inverse. If β=0, its inverse is zero and the scalar is zero; no division assumption is silently introduced. The integer indices are positive by section2, and N≥1.

There is one finite étale surjective affine cover R/S adequate for this computation. The scheme of enumerations of the finite étale constant-rank H\_Λ is finite étale and surjective. The finitely many root schemes in (ST10) are also finite étale and surjective. Their fibre product over the affine base is affine and gives the required faithfully flat algebra. Flat base change for H⁰ transports the actual theta generator y.s together with the literal \nolinkurl{FrPt.gen} property to R. The full physical tangent basis also pulls back: the coefficient field splits the O-action into its embedding summands, and the finite projective Lie module and its specified bases commute with this flat base change. Thus y\_R is an actual \nolinkurl{FrPt}, and the source \nolinkurl{FrMap} is supplied, not merely a point with a chosen torsion tuple. These are the same flat H⁰ and eigenbasis interfaces used in LIFT101. The normalized lifts and α\_N descend uniquely and require no additional arbitrary choices.

On R⊗\_S R the two pullbacks of v\_R agree: harmonic naturality compares the two pins, while (ST15) and the full enumeration comparison identify the sections. Faithfully flat descent of elements gives a unique v∈S mapping to v\_R. This proves existence in (ST2), using the actual enumeration, roots and sections just constructed as a literal \nolinkurl{ShortValAt} witness.

\subsection{8. Uniqueness against every original witness}\label{uniqueness-against-every-original-witness}

Let an arbitrary original witness use S→S', an arbitrary faithfully flat map, together with its own FrMap, complete z tuple, allowed roots, α₀, lifts and sections. S' is not assumed reduced, finite, étale or of finite type. Compare it with the constructed cover over the common faithfully flat algebra

\[
 R'=R\otimes_S S'. \tag{ST18}
\]

Sections5--6 force both witnesses to use the same α\_N and symmetric lifts, a full enumeration of the same canonical H\_Λ, and trace sections differing only by an allowed root change. Equation (ST15) makes those trace sections equal. Harmonic naturality then makes their scalar contractions equal over R'. Faithful flatness implies equality over S'. If v'∈S was the scalar supplied by that witness, injectivity of S→S' gives v'=v.

The same comparison after any coefficient-algebra map S→T and any supplied actual point related by \nolinkurl{FrMap} shows that the image of v is its unique value. The comparison does not require a theorem asserting \nolinkurl{FrPt.gen} under every nonflat map. All newly constructed frame-point pullbacks in the existence and common-cover arguments are along flat maps, where the preceding H⁰ argument applies. No scalar equality is inferred merely from equality at geometric points. The only geometric-point arguments in the proof identify finite étale torsion schemes and their discrete pairings; all trace, derivative and descent equalities take place in the full rings and retain their nilpotents.

This establishes existence, uniqueness, arbitrary base change and equivalence with every permitted witness. It does not merely define a preferred smaller subrelation of \nolinkurl{ShortValAt}.

\subsection{9. Exact application and the unsupplied second conjunct}\label{exact-application-and-the-unsupplied-second-conjunct}

In the literal \nolinkurl{ShortValTot}, use this theorem with the actual cover-stage data dδ.Xo₂, dδ.c₂, dδ.nab₂, dδ.w(x₀), dδ.eF(β\_H), and the actual coefficient map applied to P. The hypotheses \nolinkurl{Λ.Unimod}, \nolinkurl{Λ.TrivAt dδ.l₂.M} and \nolinkurl{Λ.LevOK N dδ.l₂.M} are already present. \nolinkurl{harmL} supplies the identity Laplacians, and the coefficient-field equality preserves those equations. An actual \nolinkurl{CovPin} or \nolinkurl{LawsSh.dNab} supplies the operative pin, as verified in LIFT102. No change of finite level, scalar, original class sum, selected polynomial or source object occurs.

The remaining \nolinkurl{ShortValTot} clause is

\[
 \exists f\in\operatorname{sec}(X_o,c,a,n):
 \operatorname{IsShortRule}(d_\delta,\Lambda,N,x_0,P,f). \tag{ST19}
\]

It asks that the cover-stage value, multiplied by the actual cδ, descend to a rule on the original moduli object with the specified covariance. Neither the all-ring scalar construction nor the existing \nolinkurl{IsLev} statement alone proves that invariance. The literal \nolinkurl{PhiPin}, \nolinkurl{PsiTie}, covering geometry and cδ clauses must be matched to (ST19) separately. This proof does not assert (ST19), cusp-valuedness, spanning, or a nonzero selected output.

The auxiliary finite étale covers are used solely to construct and descend the actual scalar over its original ring S. They do not change S, the coefficient field, the original source records or any DP44 cost. No numerical error charge is added.

\subsection{10. Sources, primary scope, and verification boundary}\label{sources-primary-scope-and-verification-boundary}

Literal sources freshly read for this construction are \nolinkurl{ThShort.lean} in full, especially GenLat39--64, InZh87--92, DeltaDat/FrOver/IsLev114--188, harmonic input235--237, GeoInj239--242 and ShortValAt244--259; \nolinkurl{TwoWeightEqOneAddRoot.lean}30--49, Level54--62, SAB72--84 and TateIso/Obj131--160; the complete \nolinkurl{IsCombAt.lean}, especially FrPt67--75 and FrMap79--85; and the exact recovered \nolinkurl{ThetaPt.lean} and \nolinkurl{ThetaVV.lean} definitions of SymLift, IsTwistPull, IsActOf and IsNormMul. The companion manifest records actual file hashes and exact reading extents.

Canonical LIFT100 HN4--8 supplies the already-checked ideal-half and normalized odd-lift algebra. Canonical LIFT101 supplies the all-affine formal comparison and universal normalized global-functions input. LIFT102, accepted at11:47:40 UTC on23November2026, supplies the actual-package pin and all-harmK comparison; its canonical proof fingerprint is recorded in the companion manifest.

Fresh primary readings are:

\begin{itemize}
\tightlist
\item
  Milne, J. S. (2008). \emph{Abelian varieties} (Version2.00), ChapterI, Proposition7.1 and Theorem7.2 with proof, printedpp32--33. \url{https://www.jmilne.org/math/CourseNotes/AV.pdf} . The degree and étaleness statements concern fibres; the relative finite-étale deduction is given in section3.
\item
  The Stacks Project Authors. (2026). Lemma59.64.4 (Tag03RV), with its proof, and Lemma41.18.3 (Tag04HN), with its proof: finite étale schemes and finite locally constant sheaves, and étale local trivialization. \url{https://stacks.math.columbia.edu/tag/03RV} ; \url{https://stacks.math.columbia.edu/tag/04HN} . The enumeration cover and canonical finite-kernel application are proved above.
\item
  The Stacks Project Authors. (2026). Lemma35.3.6 (Tag023M), with its proof, and Lemma59.16.3 (Tag03O9): faithfully flat descent of elements without Noetherian or reduced hypotheses. \url{https://stacks.math.columbia.edu/tag/023M} ; \url{https://stacks.math.columbia.edu/tag/03O9} . The actual equality to which descent is applied is (ST15)--(ST18).
\item
  The Stacks Project Authors. (2026). Lemma30.5.2 (Tag02KH), complete statement and proof: flat base change for a quasi-coherent module on a quasi-compact, quasi-separated morphism. \url{https://stacks.math.columbia.edu/tag/02KH} . Applied to the original proper abelian scheme and its line bundle, this preserves the specified free rank-one H⁰ module and its actual generator along the finite étale construction covers.
\end{itemize}

The exact companion diagnostics verify 131,640 assertions. They include all original-coordinate dual-kernel, cardinality and pairing tests in fourteen local diagonal maximal-order models; a non-diagonal rational unitary lattice; actual half-primary root changes and full trace permutations; multi-column traces; nilpotent coefficient transport; and explicit boundary failures when either the half-primary or TrivAt condition is dropped. They also check the eleven source fingerprints and thirty-five literal anchors. These finite tests illustrate and check the displayed algebra; they are not a substitute for the general arbitrary-base proof.

The original cubical/Poincare interpretation is retained. The proof establishes a section-valued finite trace and its faithfully flat scalar descent; it makes no claim about a new motive, Frobenius weight, slope, or global automorphic nonvanishing. No Lean compilation or source/archive validity claim is made. The first preliminary run and the final hash-specific run are preserved separately in the author packet.

\subsection{Canonical algebraic-geometric context}\label{canonical-algebraic-geometric-context-82}

The kernel is a canonical finite-etale subgroup of the original abelian scheme, defined by integral dual-row endomorphisms. It uses the actual projective CM module, alternating adjoint form, principal polarization and cubical line; no global Tate basis or changed field/source is introduced. The original theta-group trace is an equality of sections and its scalar descends faithfully flatly over arbitrary characteristic-zero coefficient rings. Exact primary applications: Milne(2008), Abelian varieties Version2.00, ChapterI Proposition7.1/Theorem7.2, printedpp32--33; Stacks Lemma59.64.4(tag03RV), Lemma41.18.3(tag04HN), Lemma35.3.6(tag023M), Lemma59.16.3(tag03O9), Lemma30.5.2(tag02KH). Relative and all-witness applications are proved in ST3--18. Carried exact geometric context: SGA7 I, VII2.7/2.9 pp169-174 and VIII7.4(b) pp310-311; FGA232 Theorem3.1 pp148-149 and FGA236 Corollaries1.7/2.3 pp229/232; SGA3 I4.7.3/4.7.3.1; SGA6 VI Theorem3.3; cote4 scan003 sections8.1.1-8.1.2 and scans005-006 sections8.4.1-8.4.3; cote54 scans005-007. These contextual locators are carried from the documented earlier readings, not claimed as new whole-text checks. No new integral extension over Z, Dieudonne slope or Cartier-duality formula, lambda-ring positivity, K-theory filtration, l-adic weight theorem or CM-motive L-function comparison is obtained. DP\_MULT Q3e/DP44/DP52/DP54-65 and THETA3 G1/A7 gain only the stated source comparison; CX\_GRH CX-ALG-024 and GCT\_PLETH T-052 retain their global arithmetic gates.

\clearpage

\section{AM-LIFT-104 --- The cover short-value rule and its exact evaluation}\label{am-lift-104-the-cover-short-value-rule-and-its-exact-evaluation}

\textbf{Canonical checked integration --- 2026-11-23T12:24:26.860928+00:00.} Complete saved-proof and executable reviews and separate exact diagnostic replays passed. No numerical error charge is introduced. Frozen author history is preserved. AM-LIFT-104 CV1--21 constructs the exact cover-stage Rules/sec(a,n) object from LIFT103 for the original HHa polynomial and operative pin. Infinite-field RowInv identities extend to every coefficient algebra; commuting invariant derivations and the literal physical frame change give the positive row character. The original theta generator contributes t\_none\^{}n, with every full-trace normalization retained. On FromTest the actual evaluator equals the unique scalar for every presentation. Actual IsLev nonzero evR(cdelta) supplies FromTest at every required algebraically closed cover frame, so H2=cdelta*G2 belongs to cover sec(a,n) and has precisely the required field evaluations. IsShortRule is equivalent to original-base descent of this particular H2 under all original FrOver quantifiers. The actual CovPin/LawsFull suppliers force both auxiliary levels A=A2=1, removing that potential mismatch in those scopes. Original-base descent itself, source/package existence, integral content, cusp/spanning, selected arithmetic nonzero, period, specific032 value and evaluatedDP44/RH remain open. No numerical charge;135/31 complete arrays and floors27/11 unchanged.

\textbf{Status: complete author-verified mathematical text, 2026-11-23; independent review and canonical acceptance remain pending.} This constructs the actual cover-stage rule from the all-ring scalar of LIFT103. It does not descend that rule across the original DeltaDat.r.

\subsection{1. Statement and the exact source boundary}\label{statement-and-the-exact-source-boundary}

Fix the actual Host, Datum, odd Level, Coeff and Moduli records, an operative \nolinkurl{IsNab} pin, and the data Λ,N,w,β of LIFT103. Thus Λ is Hermitian selfdual, \nolinkurl{TrivAt(l.M)} and \nolinkurl{LevOK(N,l.M)} hold, w is the original coefficient-field weight, and β is the original coefficient-field scalar. Let

\[
 P\in\operatorname{HHa}(Dw,n,a),
 \qquad V_S(y)=\text{the unique original ShortValAt value at }y.
 \tag{CV1}
\]

All inputs, including the original finite trace, permitted roots and normalization factors, are retained. No nonzero hypothesis on β or P is added. The HHa hypothesis includes the original identity-matrix harmonic equations, RowInv and RowDeg. It does not impose an extra integral-polynomial condition.

\textbf{Theorem CV2.} There is a unique rule \(G\in\operatorname{Rul}(X_o,c)\) whose value on every literal TRing point is \(V\). This rule belongs to the actual

\[
 G\in\operatorname{sec}(X_o,c,a,n).
 \tag{CV2}
\]

For every coefficient algebra S and every actual y with the source \nolinkurl{FromTest} property,

\[
 \operatorname{evR}(G,y)=V_S(y).
 \tag{CV3}
\]

No equality (CV3) is claimed at an arbitrary point outside FromTest; the actual source evaluator is defined to be zero there.

Apply the theorem to the original cover-stage objects of an actual DeltaDat d.~Assume the actual \nolinkurl{IsLev}, an operative pin for d.nab₂, and the lattice/polynomial hypotheses already present in \nolinkurl{ShortValTot}. Then the cover rule \(G_2\) exists and

\[
 H_2=d.c_\delta G_2\in\operatorname{sec}(d.X_{o,2},d.c_2,a,n).
 \tag{CV4}
\]

At every algebraically closed cover point \(y_2\), the actual \nolinkurl{IsLev} forces FromTest, and therefore

\[
 \operatorname{evR}(H_2,y_2)
 =\operatorname{evR}(d.c_\delta,y_2)V(y_2).
 \tag{CV5}
\]

Consequently the remaining second conjunct of ShortValTot is precisely the original-base section descent of this particular H₂ along every supplied FrOver relation. It is not an additional scalar, ordinary-coefficient or cover-rule existence problem. We do not prove that original-base descent here.

\subsection{2. Extending the literal RowInv identity to arbitrary algebras}\label{extending-the-literal-rowinv-identity-to-arbitrary-algebras}

Write F=c.F. For each embedding τ let r\_τ=Dw.mult τ, and let X\_\{τ,i,b\} be the actual row/column polynomial variables. Source RowInv says that for every i'\textgreater=0 with i'\textless i and every t∈F,

\[
 P(\ldots,X_{\tau,i,b}+tX_{\tau,i',b},\ldots)=P(X)
 \quad\text{for every column }b. \tag{CV6}
\]

The field F has characteristic zero and is infinite. Regard the left side minus the right side as a polynomial in an indeterminate T, with coefficients in F{[}X{]}. Every coefficient in X is a one-variable polynomial over F that vanishes for all t∈F. Each coefficient is therefore zero. Thus (CV6) is an identity in F{[}T,X{]}, not merely an assertion at its field-valued points.

After any F-algebra map F→S and any specialization T→s∈S, the identity remains valid. This step retains every nilpotent in S and does not assume S is reduced, Noetherian, an integral domain, or a field.

Let C\_τ be any lower-triangular matrix over S with unit diagonal entries t\_\{τ,i\}. Write C\_τ=D\_τ U\_τ, with D\_τ diagonal and U\_τ lower unitriangular. A lower unitriangular matrix over any commutative ring is a product of elementary lower transvections: successively eliminate each strictly lower entry by subtracting its multiple of the pivot row, whose pivot is one. This uses no division except by the already specified diagonal units in forming U\_τ.

Applying (CV6) to these transvections, and then the source RowDeg identity monomial by monomial, yields

\[
 P(C_\tau X_\tau\text{ for every }\tau)
 =\left(\prod_{\tau,i}t_{\tau,i}^{a(\tau,i)}\right)P(X)
 \quad\text{in }S[X]. \tag{CV7}
\]

The substitution acts simultaneously on the same row in every column. There is no inverse matrix and no inverse character. The statement includes P=0 and empty row/column sets.

\subsection{3. Exact directional-operator substitution}\label{exact-directional-operator-substitution}

We first verify the orientation directly from the literal connection definition, without asserting an unproved change of the chosen gauge.

In \nolinkurl{CocycleRightZeroThd.lean}25--36, Dof(F,u) is the translation-invariant vector field associated to u for the formal commutative group law F. Its displayed coefficients are S-linear in u. These vector fields commute. Indeed, the formal translations by two independent parameters commute by the group-law commutativity and associativity; taking the coefficient linear in each parameter in either order gives

\[
 D_uD_v=D_vD_u. \tag{CV8}
\]

Explicitly, use two square-zero indeterminates ε and δ, and expand the two equal formal expressions f(F(F(X,εu),δv)) and f(F(F(X,δv),εu)). Their εδ coefficients are the two sides of (CV8). This is a coefficient identity over S. It remains valid with nilpotents; it does not use an analytic exponential or field-valued specialization. The literal DmonOf fold order is consequently immaterial for these operators.

Let \(u'_i=\sum_j C_{ij}u_j\). S-linearity, (CV8), and expansion of each monomial give

\[
 D_{u'}(P)=D_u\bigl(P(X_i\mapsto\sum_j C_{ij}X_j)\bigr).
 \tag{CV9}
\]

Here D(P) denotes the source polynomial in commuting directional operators, independently in each column. The fixed corrected germ \(f\exp(kh_L)\) in \nolinkurl{FB.nabVal}187--190 is not altered by this substitution. Taking its constant coefficient preserves (CV9).

For comparison with the usual coordinate description, if u'=C u then old coordinates are Z=C\^{}t Z'. Hence \(\partial_{Z'_i}=\sum_j C_{ij}\partial_{Z_j}\), exactly (CV9). The transpose occurs in the coordinate change, not in the substitution on the row labels of P.

\subsection{4. Reusing an actual chart at the changed physical frame}\label{reusing-an-actual-chart-at-the-changed-physical-frame}

Suppose the actual source relation FrBor(t,y,y') holds over S. \nolinkurl{IsCombAt.lean}87--96 says that the base points agree and the physical tangent vectors satisfy

\[
 u'_{\tau,i}=\sum_{j\leq i}C_{\tau,ij}u_{\tau,j},
 \qquad C_{\tau,ii}=t_{\tau,i}\in S^\times.
 \tag{CV10}
\]

The asserted linear combinations are exactly those in the Lie module of the abelian scheme: the group operation on dual-number tangent points is addition, and epsMul gives its scalar multiplication. The old and new bases therefore have the matrices in section 2, over the actual ring S.

Choose the ChartAt supplied by the old pin at y. Keep its underlying formal group coordinates, formal group law, action, biframe, primitive generator and canonical gauge unchanged. Replace only its tangent-coordinate map by

\[
 \operatorname{uvec}'(\tau,i)=
 \sum_j C_{\tau,ij}\operatorname{uvec}(\tau,j). \tag{CV11}
\]

This is admissible chart data at y'. The source IsChart.tan follows from (CV10); its basis property follows because C\_τ is invertible. All other ChartAt clauses concern the common geometric point and the unchanged geometric/gauge data. In particular ChartAt does not contain the selected theta-section generator y.s. The existence of germs for every section carries over, because the underlying point and formal chart have not changed.

We do not identify this auxiliary chart with the chart arbitrarily selected by the stored pin at y'. Instead apply the pointwise finite-order comparison proved in LIFT101 AT11--17 and LIFT102 Theorem D. Its proof compares any two admissible chart/gauge data for the same physical point, with actual germs; it does not require them to be selected by one global choice. It says their full harmonic column contractions agree. Every column here has the same line power k=N², and P obeys the original identity-matrix harmonic equations. Thus the actual stored connection at y' can be evaluated in the auxiliary chart (CV11) for this full expression.

In that auxiliary chart, (CV9) and (CV7) give precisely the factor

\[
 \chi_a(t)=\prod_{\tau,i}t_{\tau,i}^{a(\tau,i)}.
 \tag{CV12}
\]

The argument compares whole harmonic contractions. It neither asserts equality of ordinary connection coefficients at two frames nor chooses a new operative nab₂ in the source record.

\subsection{5. The theta-generator factor is exactly t\_none\^{}n}\label{the-theta-generator-factor-is-exactly-t_nonen}

Let λ=t\_none. The remaining FrBor clause supplies an actual map between the two pullbacks of the same universal group. Its two projection equalities force the canonical identity identification of these pullbacks. Its section equality therefore says exactly

\[
 s'=\lambda s. \tag{CV13}
\]

The normalized multiplication generator α\_N is geometric and fixed. Pullback by {[}N{]} and its action on sections are linear over S. Hence \(s_N'=\lambda s_N\). The normalized odd theta lifts are also geometric and fixed. Each original IsUUOver section therefore changes by

\[
 F'(z,t_b)=\lambda F(z,t_b). \tag{CV14}
\]

The canonical kernel, complete enumeration, allowed roots and weights depend on the common underlying point, not on its tangent frame or the scaling of its theta-section generator. Choose one of the finite étale construction covers from LIFT103 for that common point. Both actual frame points pull back along it, and the FrBor relation pulls back as well. Formula (CV14) holds for every summand there, and hence for the full trace section.

There are n columns. On the fixed chart, a germ f for F becomes λf for λF, directly from the source IsGerm section identity; DpolyOf and constant-coefficient extraction are S-linear in that germ. Thus each actual pinned column is linear in its section. Multilinearity gives λ\^{}n, rather than λ\^{}(N²n): the section s is first pulled back by {[}N{]} and then sent through a line isomorphism; it is not raised to the N²-th tensor power. Combined with (CV12), this yields

\[
 V_S(y')=
 \lambda^n\chi_a(t)V_S(y).
 \tag{CV15}
\]

One may verify (CV15) on the common faithfully flat cover and descend equality by injectivity. Every β inverse, mIdx, zCard enumeration, LevQ cardinality, N-degree and w factor remains exactly the original one. The same proof applies if β=0 or P=0, when the value is zero.

\subsection{6. The actual rule and its evaluation boundary}\label{the-actual-rule-and-its-evaluation-boundary}

For a literal test record p:TRing F and y:FrPt(p.R), define

\[
 G(p,y)=V_{p.R}(y). \tag{CV16}
\]

This is a rule of exactly the source type \nolinkurl{Rules}, which is a family on the reduced finite-type test algebras p.R. Its uniqueness with property (CV16) is pointwise uniqueness; it does not mean sec of this weight is one-dimensional.

LIFT103 arbitrary-map naturality between supplied FrMap points proves IsNatural. Equation (CV15) gives IsWeight for the exact exponent function

\[
 \operatorname{ex}(a,n)(\mathrm{none})=n,
 \qquad\operatorname{ex}(a,n)(\mathrm{some}(\tau,i))=a(\tau,i),
 \tag{CV17}
\]

which is \nolinkurl{EmbKCoe.lean}13--15. By the literal definition of secW and sec, this proves (CV2).

Suppose now that y over any F-algebra S is FromTest. Every witnessing presentation consists of p, φ:p.R→S and y₀ with FrMap(φ,y₀,y). LIFT103 naturality gives

\[
 \phi(G(p,y_0))=\phi(V_{p.R}(y_0))=V_S(y). \tag{CV18}
\]

This holds for every presentation. In particular it holds for the one chosen by the actual evalAt definition, proving (CV3). No additional EvalWD axiom is needed for this particular rule: (CV18) proves its independence of the chosen presentation directly.

Outside FromTest, evalAt is literally zero. We neither replace that branch nor claim every point over an arbitrary nonreduced algebra is FromTest. The all-ring scalar V of LIFT103 continues to exist independently of this evaluator boundary.

\subsection{7. The actual IsLev clause supplies all required cover FromTest points}\label{the-actual-islev-clause-supplies-all-required-cover-fromtest-points}

Let d be the actual DeltaDat and assume its literal IsLev. That condition contains

\[
 d.c_\delta\in\operatorname{sec}(d.X_{o,2},d.c_2,0,0)
 \quad\text{and}\quad
 \operatorname{evR}(d.c_\delta,y_2)\ne0
 \tag{CV19}
\]

for every algebraically closed field and every cover frame point y₂. By \nolinkurl{AlgebraMapCAt.lean}54--58, if y₂ were not FromTest, that evaluator would be zero. Therefore every such y₂ is FromTest. This is a direct implication of the actual IsLev clause; it does not assume a separate AllTest theorem for the cover moduli.

Apply sections 1--6 with d.Xo₂,d.c₂,d.nab₂, the exact mapped P, weight d.w(x₀), scalar d.eF(betaH) and the original Λ,N. The field identification preserves RowInv, RowDeg and harmonicity, so the HHa premise is supplied by the original harmL condition. The operative pin is the actual dNab/CovPin supplier already identified in LIFT102.

Define H₂=d.cδ·G₂ as the pointwise product of the original Rules objects. The source sec\_mul theorem gives (CV4). The source evaluator is multiplicative, and (CV19) supplies FromTest for every field point used by IsShortRule, so (CV5) follows. Nothing is divided by cδ in this construction.

\subsection{8. Exact remaining original-base condition}\label{exact-remaining-original-base-condition}

For any f∈sec(Xo,c,a,n), the actual relation IsShortRule is now equivalent to

\[
 \operatorname{evR}(f,y)=\operatorname{evR}(H_2,y_2)
 \tag{CV20}
\]

for every field, coefficient-compatibility datum and pair y,y₂ with FrOver quantified in its original definition. In one direction use the unique V from LIFT103 and (CV5). In the other direction use that unique scalar as the literal ShortValAt witness and then (CV5). This retains every original FrOver condition and every cδ factor.

Equation (CV20) is the residual original-base descent problem; it is not proved by this theorem. The cover-stage rule exists, and its evaluation at every required cover field point is determined. To conclude the second ShortValTot conjunct one must still construct an original-base f satisfying (CV20).

For the actual coupled packages, the previously observed auxiliary-A mismatch has a literal supplier and is not an additional premise: \nolinkurl{ThClassPinsRec.lean}62--63 defines CovPin to include d.IsLev, d.l₂.A=1 and the operative IsNab. Since IsLev contains Level.Dvd, the original positive l.A divides one and hence also equals one (\nolinkurl{TwoWeightEqOneAddRoot.lean}53--66). The actual OMPins/OneMapTieP route already supplies this CovPin, as read in LIFT102. Likewise \nolinkurl{FullExists.lean}40 supplies the new-A equality through LawsFull.lvA when that full record is given. Thus, in those actual scopes,

\[
 l.A=d.l_2.A=1. \tag{CV21}
\]

The extra new-A identity in PsiTie is then automatic, and the root-of-unity/base-field clause at A is unchanged. This removes that particular potential mismatch for these genuine packages; it does not claim that bare IsLev alone forces (CV21).

We do not assume that the declared ToroidalLevelCompatibility interface is automatically supplied for an arbitrary stored d.r. The independent original-base investigation must still retain the actual finite cover, its level action and all required fibre invariance. This proof uses neither PhiPin nor PsiTie to bypass that issue.

\subsection{9. Scope, sources and diagnostic boundary}\label{scope-sources-and-diagnostic-boundary}

No additional finite extension, choice of source object, new polynomial, level change, norm bound or numerical error charge is introduced. The construction retains the actual package, and remains conditional on its geometric records and operative pin. It does not prove source admission, original-base descent, integral coefficient descent, cusp/spanning, selected arithmetic nonzero or the final DP44 inequality.

Literal fresh readings include the complete IsCombAt, RowDeg, EmbKCoe and ThGerm files; AlgebraMapCAt 1--115; CocycleRightZeroThd 1--76 and 160--224; ThShort 114--235 and 244--355; complete HasDataCondition, ThClassPinsRec and FullExists; and TwoWeightEqOneAddRoot 1--100. The companion manifest records precise fresh and carried extents and source hashes. The new deductions are the elementary polynomial identity (CV6)--(CV7), commuting invariant-operator calculation (CV8)--(CV9), exact FrBor trace scaling, construction of the actual Rules object and the FromTest implication (CV19). The all-ring scalar and pointwise harmonic comparison are the accepted LIFT103 and LIFT101/LIFT102 inputs.

This is a mathematical source-interface argument, not a Lean compilation or source/archive validity claim. The companion executable checks finite exact models over Q{[}ε{]}/(ε³), including full lower-triangular row transformations, identity-matrix harmonic cancellation across columns, a nonadditive formal-group chart, actual generator scaling and evaluator branch controls. These diagnostics accompany the general proof and do not substitute for arbitrary-ring reasoning.

\subsection{Canonical algebraic-geometric context}\label{canonical-algebraic-geometric-context-83}

The construction uses the original CM abelian scheme, projective OE lattice, rigidified cubical line and actual physical tangent eigenframes. Formal invariant differentiation and identity-matrix harmonic contraction commute with the stated coefficient maps, retaining nilpotents. The cover rule is a natural function on the literal reduced finite-type frame tests with exactly the source Borel character; its evaluation boundary is derived from the source FromTest branch. The original-base section descent remains explicit. Exact Alpoge interfaces: RowDeg.lean33--48; IsCombAt.lean67--127; EmbKCoe.lean13--15; \nolinkurl{ThGerm.lean56–74/189–217;} \nolinkurl{CocycleRightZeroThd.lean25–36/166–190;} AlgebraMapCAt.lean9--98; \nolinkurl{ThShort.lean155–181/263–293;} ThClassPinsRec.lean62--63; FullExists.lean40; \nolinkurl{TwoWeightEqOneAddRoot.lean53–66.} The non-elementary geometric inputs are accepted LIFT101--103 at their documented primary scopes. Carried exact context: SGA7 I VII2.7/2.9 pp169--174 and VIII7.4(b) pp310--311; FGA232 Theorem3.1 pp148--149 and FGA236 Corollaries1.7/2.3 pp229/232; SGA3 I4.7.3/4.7.3.1; SGA6 VI Theorem3.3; cote4 scan003 sections8.1.1--8.1.2 and scans005--006 sections8.4.1--8.4.3; cote54 scans005--007. These are carried contextual readings, not new whole-text checks. No new integral extension over Z, Dieudonne slope, lambda-ring positivity, K-theory filtration, l-adic weight theorem or CM-motive L-function comparison follows. DP\_MULT Q3e/DP44/DP52/DP54--65 and THETA3 G1/A7 gain this conditional source-rule reduction; CX\_GRH CX-ALG-024 and GCT\_PLETH T-052 retain their global gates.

\clearpage

\section{AM-K-122 --- an actual selected evaluation coherence cut}\label{am-k-122-an-actual-selected-evaluation-coherence-cut}

\textbf{Canonical checked integration --- 2026-11-23T12:27:23.492137+00:00.} Complete saved-proof and executable reviews and separate exact diagnostic replays passed. No numerical error charge is introduced. Frozen author history is preserved. AM-K-122 HC1--18, with the mandatory explicit nonzero-original-S-cross scope addendum, bounds the actual selected trace by its simultaneous narrow-band high-coherence integral plus the complete band error and2\^{}(-B/4) locally. Algebraic evaluation has plus-N circle weight; its genuine metric adjoint range lies in the selected high-weight spectral subspace. Two-projector exponential factorization, exact two-singular-value calculation and actual all-word/twist weights give overlap below2\^{}(-B/4) whenever projective angular coherence is at most9/10. Insertion into the original finite trace before absolute values retains both actual fibre metrics, isometric transport and every compact factor. The full original bounded support27 all-II globally rank-two-unipotent scalar-cross family with nonzero original S-cross and x\textless2\^{}3328 has paid global errors2\^{}-8605114 and2\^{}(11394885-B/4), the latter below2\^{}-40000000. The actual high-coherence main integral remains unbounded. This analytic input adds no sector charge, removes no k145..164 cell and changes no complete support floor;135/31 arrays and floors27/11 remain. Source admission, selected arithmetic image/content/nonzero, signedPP3, period, specific032 and evaluatedDP44/RH remain open.

Author: \nolinkurl{/root/kummer_extension}. 23 November 2026 UTC. \textbf{Research-pending complete analytic theorem; author and independent checks are separate.} No canonical file, numerical sector charge or support floor is changed. The theorem removes a precisely defined part of the original coefficient integral with a fully paid arithmetic error. Estimating the remaining high-coherence integral is a separate problem.

\subsection{1. Actual selected representation and conclusion}\label{actual-selected-representation-and-conclusion}

Retain the all-top J56 upper endpoint, every admitted binary word and canonical twist, the actual selected irreducible compact type \[
 W=V_\lambda\subset V_\alpha\otimes\operatorname{Sym}^N(\mathbf C^a),
 \qquad a=f=168,\quad B=2^{955},\quad 2^{484}<\ell<2^{485}.
 \tag{HC1}
\] The selected Hilbert metric, the homogeneous bundle metric at every ball point, the two compact factors, and the original action are exactly K043 PB1--8 and DO080 RT1--8. The polynomial degree and compact weights have their original convention; the interior word changes the m\_i, not the r\_i.

For two nonzero ball points define their projective angular coherence by \[
 \chi(z,z')=\frac{|\langle z,z'\rangle|}{|z|\,|z'|}\in[0,1].
 \tag{HC2}
\] At a zero point assign chi any value: all selected sections have positive homogeneous degree, so their evaluation and the original density vanish there. Let \[
 I_*=[16/529,16/527],\quad
 F_{I_*,>} (g)=\int_{Y(z),Y(g^{-1}z)\in I_*\atop
                         \chi(z,g^{-1}z)>9/10}
       \sqrt{\rho(z)\rho(g^{-1}z)}\,d\mu(z).
 \tag{HC3}
\] This is a restriction of the original nonnegative fidelity, not a new scalar representation or a signed coefficient formula.

\textbf{Theorem.} At every real label, for every original allowed word/twist and every group element g, \[
 |c_W(g)|\le F_{I_*,>}(g)+2\sqrt{\Pr(Y\notin I_*)}
                                  +2^{-B/4}.
 \tag{HC4}
\] The low-coherence part of the pointwise selected trace integral gains the factor 2\^{}(-B/4) relative to its original Cauchy majorant. The statement includes arbitrary compact factors and r=0; it divides by neither a Cartan radius nor a coordinate.

On the original bounded all-II globally rank-two-unipotent scalar-cross support27 family with x=N\_E D\_eta\textless2\^{}3328, the full arithmetic composition is \[
 \sum_{h_0}|w(h_0)c(h_0)|
 \le\sum_{h_0}|w(h_0)|\prod_\tau
 F_{I_*,>,\tau}(g_{p,\tau}^{-1}h_0g_{p,\tau})
 +2^{-8605114}+2^{11394885-B/4}.
 \tag{HC5}
\] All supports, finite patterns and labels, original rows and affine central fibres, class/Haar factors and source normalization are included. The last term is less than 2\^{}(-40000000). No charge is appended: the remaining main integral is not estimated by this theorem.

\subsection{2. The genuine evaluation range has high compact weight}\label{the-genuine-evaluation-range-has-high-compact-weight}

For a unit row direction u let P\_u=u\^{}*u be the rank-one Hermitian projection on C\^{}a. Let T\_z:W→E\_z be evaluation into the actual Hermitian bundle fibre. Here T\_u denotes algebraic evaluation of a homogeneous polynomial at the unit vector u into the fixed compact fibre V\_alpha; no boundary value or boundary Hermitian metric is asserted. Every metric adjoint below is taken for an interior point z=r u with 0\textless r\textless1. Homogeneity gives the same adjoint range at every nonzero point z=r u on that ray, after the usual coordinate identification. Changing the positive fibre metric multiplies evaluation on the left by an invertible operator; hence it does not change the range of its adjoint as a subspace of W. The selected Hilbert metric is a positive scalar multiple of the Fischer metric on this one irreducible W, by PB5. No such assertion is made for the whole ambient degree-N space.

Consider the compact one-parameter group k(t)=exp(it P\_u) in the positive U(a) block. The literal compact action satisfies \[
 T_u\tau_\lambda(k(t))=e^{iNt}\tau_\alpha(k(t))T_u.
 \tag{HC6}
\] For u=e\_1 this is direct substitution in homogeneous degree-N polynomials: the evaluation coordinate receives the phase e\^{}(it), and the fibre retains its own compact action. Conjugation gives every unit u. This is the plus-N convention already fixed by lambda=alpha+m in PB2--4; the opposite coordinate convention would require conjugating both sides, not silently changing one weight.

Every eigenvalue of d tau\_alpha(P\_u) lies in {[}alpha\_last,alpha\_0{]}. For example, after subtracting alpha\_last from each highest weight, this is a polynomial representation with each coordinate weight between0 and alpha\_0−alpha\_last; compact conjugation handles general u. Thus the adjoint evaluation range lies in the spectral subspace \[
 \operatorname{ran}T_z^*\subseteq
 {f1}_{[N+\alpha_{a-1},\infty)}(d\tau_\lambda(P_u))W.
 \tag{HC7}
\] This follows by taking adjoints of the finite-dimensional intertwining identity HC6, or by matching its circle characters. It does not identify this range with the whole spectral subspace and does not use the unproved low-ladder Gram estimate.

Twist only this finite compact representation by det\^{}(-alpha\_last), and denote its polynomial highest weight by \[
 \mu_i=\lambda_i-\alpha_{a-1}\ge0,\qquad
 X_u=d\tau_\mu(P_u),\qquad
 Q_u={\bf1}_{[N,\infty)}(X_u).
 \tag{HC8}
\] This auxiliary finite-dimensional determinant twist is used only to bound the same spectral subspaces. It changes neither the original group action in the trace nor the source representation. HC7 says T\_z=T\_z Q\_u and T\_z\^{}\emph{=Q\_u T\_z\^{}} in the actual selected Hilbert metric.

\subsection{3. An exact two-projection exponential bound}\label{an-exact-two-projection-exponential-bound}

Let X,Y be finite-dimensional Hermitian operators and P,Q their spectral projections on eigenvalues at least N. For any s\textgreater0, \[
 PQ=P e^{-sX/2}\,[e^{sX/2}e^{sY/2}]\,e^{-sY/2}Q,
\] so \[
 \|PQ\|\le e^{-sN}\|e^{sX/2}e^{sY/2}\|.
 \tag{HC9}
\] There is no commutation assumption between X and Y. Only each spectral projection commutes with its own operator.

Apply HC9 to X\_u,X\_v. Put b=e\^{}(s/2)\textgreater1 and \[
 A_u=I+(b-1)P_u,\qquad A_v=I+(b-1)P_v.
\] The middle operator is the genuine finite representation tau\_mu(A\_u A\_v). If x=\textbar\textless u,v\textgreater\textbar, the only two nontrivial singular values of A\_u A\_v are sigma\_+\textgreater=sigma\_-\textgreater=1 and satisfy \[
 \sigma_+\sigma_-=b^2,\qquad
 \sigma_+-\sigma_-=(b^2-1)x,
\] \[
 \sigma_+=\frac{\sqrt{4b^2+(b^2-1)^2x^2}+(b^2-1)x}{2}.
 \tag{HC10}
\] For proof, restrict to span(u,v). The determinant is b\^{}2, and the trace of the squared singular values is 2b\^{}2+(b\^{}2−1)\^{}2x\^{}2. Outside that span the product is the identity. The limiting cases x=0 and x=1 follow by continuity. The inequality sigma\_-\textgreater=1 also follows from sigma\_+\textless=b\^{}2.

Singular-value decomposition and the actual unitary compact norm imply \[
 \|\tau_\mu(A_u A_v)\|
                  =\sigma_+^{\mu_0}\sigma_-^{\mu_1}.
 \tag{HC11}
\] Indeed the two unitary factors act isometrically. On a positive diagonal matrix the weight spaces are orthogonal. Every weight is majorized by the highest weight mu, and the logarithms of the diagonal singular values are ordered; their largest pairing is attained by mu. This proves HC11 directly, without a Hilbert/Fischer substitution for an unrelated subspace. The remaining singular values equal1 and contribute no omitted factors.

Choose b=4/3. When x\textless=9/10, the exact formula HC10 gives \[
 \sigma_+<173/100,\qquad\sigma_-=(16/9)/\sigma_+.
 \tag{HC12}
\] The strict comparison is checked by squaring positive rationals: with v=(b\^{}2−1)(9/10)=7/10, one has (2·173/100−v)\^{}2\textgreater4b\^{}2+v\^{}2.

\subsection{4. The actual all-word inequalities give an enormous uniform gain}\label{the-actual-all-word-inequalities-give-an-enormous-uniform-gain}

Use the literal DO080 RT22 and K120 NB6 formulas. If s\_omega is the number of one-bits among the166 interior positions, then \[
 N=32B-(4+s_\omega)\ell-71,
 \quad0\le s_\omega\le166.
 \tag{HC13}
\] Since r\_i=alpha\_i−beta−i, the fibre width and the first two selected weights after HC8's determinant shift are \[
 \alpha_0-\alpha_{167}=B+668\ell-71,
 \quad \mu_0=33B+12,
\] \[
 \mu_1=B+668\ell-71-\ell\omega_1,
 \qquad \mu_{167}=12.
 \tag{HC14}
\] The canonical twist cancels from these differences. All other mu\_i are nonnegative by dominance. The exact endpoint inequalities ell\textless2\^{}485 and B=2\^{}955 imply, uniformly in every word/twist, \[
 N>\tfrac{319}{10}B,\qquad
 \mu_0<\tfrac{331}{10}B,\qquad
 \mu_1<\tfrac{11}{10}B.
 \tag{HC15}
\] Because both singular values are at least1, HC9--12 and HC15 give \[
 \|Q_uQ_v\|
 <(4/3)^{-(319/5)B}
        \sigma_+^{(331/10)B}\sigma_-^{(11/10)B}
 \le\left[(173/100)^{320}(3/4)^{616}\right]^{B/10}
 <2^{-B/4}.
 \tag{HC16}
\] The last strict inequality is the exact integer-power comparison \[
                 2^5(173/100)^{640}(3/4)^{1232}<1.
 \tag{HC17}
\] This is not an asymptotic estimate in B or an empirical test over words. HC13--15 prove the whole parameter range; HC17 uses only modest integer powers.

\subsection{5. Insert the overlap into the original trace before taking absolute values}\label{insert-the-overlap-into-the-original-trace-before-taking-absolute-values}

At z and z'=g\^{}(-1)z let U\_g:E\_z'→E\_z be the original isometric bundle transport. In the actual orthonormal basis of W the integrand in DO080 RT5--7 is the finite trace of T\_z\^{}* U\_g T\_z'. Cyclicity and HC8 give \[
 \operatorname{Tr}(T_z^*U_gT_{z'})
 =\operatorname{Tr}(Q_vQ_uT_z^*U_gT_{z'}).
\] The trace-norm inequality and Hilbert--Schmidt Hölder inequality imply \[
 |\operatorname{Tr}(T_z^*U_gT_{z'})|
 \le\|Q_vQ_u\|\,\|T_z\|_{\rm HS}\|T_{z'}\|_{\rm HS}.
 \tag{HC18}
\] This retains the actual fibre metrics: U\_g is an isometry between them, and the ranges of the adjoint evaluation maps are the HC7 subspaces. No fibre condition number, unnormalized polynomial quotient, or scalar Bergman replacement is present. Multiplication by c\_tau/d\_W turns the product of the two Hilbert--Schmidt norms into exactly sqrt(rho(z)rho(z')).

On chi\textless=9/10 use HC16; on its complement use the ordinary bound1 for the overlap. The total original fidelity is at most1 by Cauchy and invariant measure. The part outside the simultaneous radial band costs at most2sqrt(Pr(Y notin I\_*)), exactly K118 BT15. This proves HC4. It is uniform in g and its compact factors; the integral is the ordinary convergent selected Hilbert realization, in Nakahama's actual convergence chamber. No analytic continuation of divergent norms is used.

For completeness, each local quantity is bounded by1 even if the right side of HC4 exceeds1. Set v\_tau=min(1,F\_high,tau+epsilon\_tau), u\_tau=F\_high,tau and telescope products. This gives the sum of the local additive errors while retaining u\_tau in {[}0,1{]}. The term involving coherence is at most f·2\^{}(-B/4)\textless2\^{}(8−B/4).

K118 BT16--20 gives the full bounded support27 normalized nonarchimedean absolute mass \textless2\^{}11394877. Its count includes all167 original values, the affine central parameter, the full4\^{}334 common patterns, active supports and labels, both mixed contractions and original source/class/Haar normalization. Eta remains unique per support by the original bounded reconstruction. Using the actual x\textless2\^{}3328 cutoff, this same containing mass applies to every surviving k and both original eligibility alternatives. Therefore the coherence term costs \textless2\^{}(11394885−B/4). The independently proved narrow-band arithmetic tail remains exactly2\^{}(-8605114). This proves HC5 and pays every additive error before any proposed row or Gaussian enlargement.

\subsection{6. Scope, source record and remaining problem}\label{scope-source-record-and-remaining-problem}

HC5 is a reusable strict cut on the actual coefficient integral. It does not estimate the remaining high-coherence integral, claim a better finite-sector exponent, remove k145, or change the135/31 ledger. In particular, the restriction chi\textgreater9/10 depends on the actual two ball points and cannot be replaced by an independent compact rotation, a fictitious scalar coefficient or the old radial height alone. A later application must prove its angular/row/central comparison and pay both HC5 errors.

The carried selected action and ordinary convergence are Nakahama(2016), equation5.5, equation5.7 and Theorem5.1, through the exact K043 PB2--8 dictionary. These specific primary passages were directly reread in the retained published text for this theorem. DO080 RT1--8 supplies the actual normalized finite trace, evaluation density and isometric bundle transport; RT22, with the mandatory corrected word convention, supplies HC13--14. K120 NB1--13 gives the actual narrow band, complete tail and exact ell bounds. K118 BT15--22 supplies the full arithmetic containing mass and product composition. None of their numerical removal charges is reused.

The spectral projection, singular-value and trace-norm steps above are proved explicitly. A separate exploratory harmonic-divisibility lead motivated this argument but is not required for HC4--5. Its normalized quotient is not used. No low-ladder Gram matrix, finite Fischer overlap or minimum-weight dimension is substituted for the genuine evaluation-range projection.

The actual base field, involution, Hermitian/projective modules, theta data and prescribed selected family remain fixed. No source constructor, selected arithmetic image nonvanishing, prescribed period, DP44, specific032 or RH conclusion follows. Accepted Alpöge files are used only at their statement/interface boundaries; there is no theorem-validity audit or Lean build.

\subsubsection{References}\label{references-7}

Nakahama, R. (2016). Norm computation and analytic continuation of vector valued holomorphic discrete series representations. \emph{Journal of Lie Theory, 26}(4),927--990. Equations5.5 and5.7, Theorem5.1, printed959--961. \url{https://jolt.centre-mersenne.org/item/10.5802/jolt.919.pdf} . Fresh scoped primary read, with the exact source convention K043.

\subsection{Mandatory preserved scope addendum}\label{mandatory-preserved-scope-addendum}

\section{AM-K-122 --- mandatory containing-family scope precision}\label{am-k-122-mandatory-containing-family-scope-precision}

The preserved proof \nolinkurl{selected_evaluation_high_weight_coherence.md}, SHA-256 \nolinkurl{8dfc27c46fc86b1266c66947ce7fb1e02d423a2121484927966df8c36d6a2350}, and its executable remain unchanged.

In HC5 and every use of the K118 BT16--20 containing mass in section5, the original bounded all-II globally rank-two-unipotent scalar-cross support27 family has \textbf{nonzero original S-cross}, in addition to x=N\_E D\_eta\textless2\^{}3328 and the other stated original hypotheses. This is the precise K118 BT5/BT16--20 family, where the actual slope/row/central reconstruction applies. In particular it contains every presently surviving k145..164 cell. No arithmetic mass statement for a zero-S-cross family is made.

The local selected-coefficient statements HC4 and HC6--18 do not use this arithmetic premise and retain their full stated group-element and all-word scope. No constant, executable comparison, ordered error array or numerical sector charge changes.

\subsection{Canonical algebraic-geometric context}\label{canonical-algebraic-geometric-context-84}

This is a real-place spectral/evaluation estimate on the unchanged selected theta coefficient, retaining the original polarized projective order modules, CM involution, cubical/Poincare data, finite projectors, phases, Haar measures and arithmetic periods. The range of the true evaluation adjoint and the actual invariant Hermitian fibre metric determine the angular cut. Nakahama(2016), Norm computation and analytic continuation of vector valued holomorphic discrete series representations, Journal of Lie Theory26(4),927--990, equations5.5/5.7 and Theorem5.1, printed959--961, is freshly read in the retained published primary with the exact K043 PB2--8 convention; DO080 RT1--8/RT22 supplies the source-normalized trace and corrected word formula. K118 BT5/BT15--22 and K120 NB1--16 supply the precise bounded nonzero-S-cross family and full arithmetic mass/tail. Carried exact geometric context: SGA7 I, VII2.7/2.9 pp169-174 and VIII7.4(b) pp310-311; FGA232 Theorem3.1 pp148-149 and FGA236 Corollaries1.7/2.3 pp229/232; SGA3 I4.7.3/4.7.3.1; SGA6 VI Theorem3.3; cote4 scan003 sections8.1.1-8.1.2 and scans005-006 sections8.4.1-8.4.3; cote54 scans005-007. These contextual locators are carried from the documented earlier readings, not claimed as new whole-text checks. No scheme over Z or Picard/motive is reconstructed, and no new Dieudonne slope, Cartier-duality formula, lambda-ring/K-theory positivity, l-adic weight or CM-motive L-function theorem follows. DP\_MULT Q3e/DP44, THETA3 G1/A7, CX\_GRH CX-ALG-024 and GCT\_PLETH T-052 retain the same actual source, arithmetic image, period and global comparison gates.

\clearpage

\section{AM-DO-126 --- The integral isotropic quotient and the support27 row count}\label{am-do-126-the-integral-isotropic-quotient-and-the-support27-row-count}

\textbf{Canonical checked integration --- 2026-11-23T12:42:47.616444+00:00.} Complete root and independent proof/script/interface reviews and separate exact replays passed. The actual canonical predecessor has135/31 entries; this integration gives136/31, complete floors27/11. Frozen author history is preserved. AM-DO-126 IQ1--36 removes the complete original bounded support27/total-eligibility27 all-II full-positive scalar-cross globally rank-two-unipotent union145\textless=k\textless=159, b=j\_int0 and2\^{}2111\textless x\textless2\^{}3328. The saturated original isotropic quotient is a projective selfdual O\_E-module, including at2/U; every actual row lies in J I\_Sigma D\_eta times its ordinary row dual. Its standard S164 summand and selfdual rank3 complement give exact transverse covolume D0\^{}3 z\^{}(6f), while the independent physical central U\^{}-f remains. A genuine real-dimension1008 Gaussian bound retains both transverse terms. The unchanged positive central mixture, exact continuous moments148/162, finite-endpoint Holder, repaired zero term, all4590 positive terms, all original finite/source factors and the entire band tail yield one582 charge. Worst directed term is−19510525/32768\textless−595; the entire union is\textless2\^{}-582. Exact136/31 full error sums remain\textless2\^{}-230 with floors27/11. The sole remaining support27 necessary alternative is160\textless=k\textless=164 with the same denominator interval. The positive k160 sufficient envelope removes nothing. Signed PP3, higher-support remainder, actual source and selected arithmetic image, prescribed period, specific032 value, actual lift data/evaluatedDP44 and Hecke RH remain open.

Author: \nolinkurl{/root/definite_oscillator}, 23 November 2026 UTC. Controlled pending proof. Complete author verification, independent review and canonical acceptance are separate events. No canonical file is changed here.

\subsection{1. Actual sector and proposed theorem}\label{actual-sector-and-proposed-theorem}

Retain the original J56 all-top data, the standard finite compact, every original source and target pencil, both CM components, selected module and Hilbert metric, compact factors, class masses, phases and Haar normalizations. The actual predecessor is DO125, with135 all-top and31 uniform charges and complete support floors27/11. Its remaining all-II globally rank-two-unipotent scalar-cross support27 sector has

\[
 s=q_{\rm elig}=27,\quad b=j_{\rm int}=0,\quad
 2^{2111}<x=N_ED_\eta<2^{3328},\quad145\le k\le164.
 \tag{IQ1}
\]

Here k is the number of nonzero original S coordinates. All334 unmixed G reductions are unipotent. The scalar eta is unique for each fixed active support in this window by K113 SC27.4--10, before any G pattern or row is counted.

\textbf{Theorem.} The complete original absolute normalized contribution of the subset

\[
                         145\le k\le159
 \tag{IQ2}
\]

is less than \(2^{-582}\), including the entire original narrow-band tail once. Append one582 charge to the actual predecessor, retaining every old charge. Both full error sums remain below2\^{}-230; the complete floors27/11 remain unchanged. The k160 calculation below is only a failed sufficient envelope.

Use

\[
 f=168,\ H=133,\ R=17424/133,\ Q=1427908193729121522601,
 \quad U=77Q^2-14,
\] \[
 D_0=7^{140}29^{144}5^{84}13^{84},\quad
 D_F=2^fD_0,\quad |D_E|=D_F^2=2^{2f}D_0^2.
 \tag{IQ3}
\]

The new step is an original integrality condition on the row at the primes dividing U. It replaces an enlarged transverse value lattice by its actual containing selfdual quotient lattice. It does not remove the independent physical central Jacobian.

\subsection{2. The exact ambient lattice and original normal form}\label{the-exact-ambient-lattice-and-original-normal-form}

R025 R25.1--3 defines the actual compact by

\[
 \Lambda=O_E^{169},\qquad J=\operatorname{diag}(I_{168},-1),
 \qquad K=\operatorname{Stab}(\widehat\Lambda).
 \tag{IQ4}
\]

Thus Lambda is unimodular for the displayed O\_E-valued Hermitian form, including at primes over2 and all rationally ramified primes. The actual original matrix preserves Lambda away from the selected T support; no hyperspecial replacement at a ramified place is made. This is also the precise compact interface in DO106 CFI6--8.

Use the actual rational orthogonal decomposition

\[
 E^{169}=Z\perp P,\quad Z=S\perp V,\quad
 S=E^{164},\quad V=u^\perp\subset E^4,\quad
 P=Eu\perp Ee_n,
\] \[
 u=(2,95192390527,383092671043,12529843548438452388609),
 \qquad u^*u=U.
 \tag{IQ5}
\]

The original integer vector u is primitive. The globally unipotent normal form DO098 CF4 / DO125 CM4, in the original ambient coordinates, is

\[
 w=(0_S,u,\eta),\quad\eta\bar\eta=U,\quad
 N=h-I=w\mathcal R-vw^\dagger+dww^\dagger,
\] \[
 \mathcal R|_P=0,\quad r=\mathcal R|_Z,\quad
 v=J_Z^{-1}r^*,\quad d+\bar d=-rv,\quad N^3=0,
 \quad\operatorname{rank}N=2.
 \tag{IQ6}
\]

Here \(w^\dagger=w^*J\), so \(w^\dagger w=0\). In particular

\[
 Nw=0,\qquad Ny=w\mathcal R(y)\ (y\in\ell^\perp),
 \quad\ell=Ew,\qquad
 \alpha=(N^2)_{nn}=U\|r\|^2>0.
 \tag{IQ7}
\]

All real conjugates of alpha are at most R. The original representative has the non-strict cutoff \(\|h\|\le H\). We retain the mandatory DO122 cutoff addendum.

\subsection{3. The primitive isotropic quotient, with no free-basis assumption}\label{the-primitive-isotropic-quotient-with-no-free-basis-assumption}

Put \(D=D_\eta=\{a\in O_E:a\eta\in O_E\}\). Primitive integer Bezout for u gives the exact fractional-line equality

\[
 \ell\cap\Lambda=Dw.
 \tag{IQ8}
\]

Indeed aw is integral if and only if a and a eta are integral. This applies at every finite prime, even when eta has a pole or a zero and even over2. The line module is saturated because it is the intersection of a vector-space line with Lambda.

Define

\[
 \mathscr K_\eta=(\ell^\perp\cap\Lambda)/(\ell\cap\Lambda).
 \tag{IQ9}
\]

This is a finite projective O\_E-module of rank167, with a perfect integral Hermitian pairing. Here is a local proof that includes dyadic places. Localize over a finite F-place, so the coefficient algebra is either the quadratic local integer ring or the product of the two split integer rings. The saturated line has a primitive local generator e. Since the ambient Hermitian module is unimodular, choose f with \(\langle e,f\rangle=1\). Its plane has Gram matrix

\[
 \begin{pmatrix}0&1\\1&a\end{pmatrix},\qquad a\in O_F,
 \quad\det=-1.
 \tag{IQ10}
\]

The inverse Gram matrix is integral. Thus the plane is an integral orthogonal direct summand, its complement is unimodular, and its contribution to e-perpendicular modulo e is zero. The quotient is exactly that unimodular complement. This proof divides neither by2 nor by the norm of a primitive vector. In the split case it is the same argument over the product algebra, with conjugation exchanging its factors; both CM pairings are retained. Consequently the local quotients are free and selfdual, and they glue to the projective selfdual module IQ9. No global generator of D, global partner f, or global integral orthogonal basis is selected.

Over E, \(\ell^\perp=Z\oplus\ell\). Projection to Z therefore identifies the quotient in IQ9 isometrically with an O\_E-lattice in the original positive Hermitian space Z. Its metric is positive definite at every complex place. Since w has zero S coordinates, the original standard S module splits orthogonally:

\[
 \mathscr K_\eta=O_E^{164}\perp\mathscr K_{\eta,V},
 \qquad\operatorname{rank}_{O_E}\mathscr K_{\eta,V}=3,
 \qquad\mathscr K_{\eta,V}=\mathscr K_{\eta,V}^{\dagger}.
 \tag{IQ11}
\]

The last dual is the O\_E-valued Hermitian dual. The ordinary E-linear row dual \(\mathscr K_\eta^\vee=\operatorname{Hom}_{O_E}(\mathscr K_\eta,O_E)\) carries the corresponding positive dual metric and is likewise unimodular. Under Hermitian Riesz identification it is conjugate-linearly isometric to \(\mathscr K_\eta\); this convention does not identify I with bar(I).

\subsection{4. The actual row belongs to the new lattice at every prime}\label{the-actual-row-belongs-to-the-new-lattice-at-every-prime}

Let J be the original whole-C ideal, with the deterministic paired orientation of R072 PJ6. Let \(I_\Sigma\) be the original paired simple-pole ideal, \(N_EI_\Sigma=T^{-54}\), and set

\[
 I=JI_\Sigma D,\qquad N_EI=z^{2f}=N_EJ\,xT^{-54}.
 \tag{IQ12}
\]

We claim the stronger global membership

\[
                       r\in I\mathscr K_\eta^\vee.
 \tag{IQ13}
\]

First work away from G and T. The original matrix and N are integral on Lambda there. For every \(y\in\ell^\perp\cap\Lambda\), IQ7 gives \(Ny=w\mathcal R(y)\in\Lambda\). By IQ8 this says \(\mathcal R(y)\in D\). The row vanishes on ell, so it descends to IQ9 and proves \(r\in D\mathscr K_\eta^\vee\) locally. This is a condition from \textbf{all} original integral vectors in ell-perpendicular, not an inference from the three old fixed transverse inputs. It includes every prime dividing U, every possible eta denominator prime, all dyadic primes and all rationally ramified primes outside G/T.

At every G or T place, \(\gcd(2U,QT)=1\) by the actual host arithmetic. The standard lattice consequently splits integrally as

\[
 \Lambda=Z_O\perp P_O,\quad
 P_O=O_Eu\oplus O_Ee_n,\quad
 Z_O=O_E^{164}\perp(u^\perp\cap O_E^4).
 \tag{IQ14}
\]

The projector along u uses U\^{}-1, which is a unit here. Since the perpendicular of ell in the Hermitian plane P is ell, projecting IQ9 gives exactly \(\mathscr K_\eta=Z_O\) at these primes. This remains true when eta has opposite nonzero valuations at a split G pair: its primitive line module in P may vary, but the positive quotient is still the same Z\_O.

The actual G/T frame from R072 PJ4--6 has determinant2U. The already proved whole-C conditions therefore hold on \textbf{every} input of Z\_O in the same orientation. They give \(r(y),\eta r(y)\in JI_\Sigma\) for every such y. The valuationwise intersection is exactly

\[
 JI_\Sigma\cap\eta^{-1}JI_\Sigma=JI_\Sigma D.
 \tag{IQ15}
\]

This proves IQ13 at G/T, including both CM partners. At an active T place eta is a unit by eligibility; at inactive places or G the intersection formulation also covers zeros and poles. No new orientation or congruence condition is assigned. Combining the local statements proves IQ13 globally. In particular primes over2 are handled by the away-G/T proof and never by declaring a determinant2U frame invertible there.

For fixed eta and J, the admissible row is thus contained in

\[
             I^{164}\perp I\mathscr K_{\eta,V}^{\vee}.
 \tag{IQ16}
\]

The selfdual quotient depends on eta, but eta is already unique for the fixed active support. No isomorphism-class or basis multiplicity is added. We use only a containing lattice; sufficiency of IQ13 for a complete arrow is not asserted.

\subsection{5. Exact metric covolume and dual minimum for a projective Hermitian lattice}\label{exact-metric-covolume-and-dual-minimum-for-a-projective-hermitian-lattice}

Use the same real metric as the source count, \(\|v\|^2=\sum_\tau h_\tau(v_\tau,v_\tau)\), the half trace over the f complex labels. Let K be any positive selfdual rank-m projective O\_E-lattice and I any fractional ideal of norm z\^{}(2f).

The real Euclidean dual of K is \(2\mathfrak D_E^{-1}K\). Multiplication of a rank-m projective lattice by a fractional ideal A changes its real covolume by \((N_EA)^m\), as follows locally from its determinant ideal or finite indices; principality is not required. The exact discriminant IQ3 gives

\[
 N_E(2\mathfrak D_E^{-1})=D_0^{-2}.
\]

Comparing the covolume of the real dual with the reciprocal covolume proves

\[
 \operatorname{covol}K=D_0^m,\qquad
 \operatorname{covol}(IK)=D_0^m z^{2mf},\qquad
 (IK)^\#=2\mathfrak D_E^{-1}\bar I^{-1}K.
 \tag{IQ17}
\]

The same formulas apply to the ordinary row dual with its dual Hermitian metric. This is an exact field-and-ideal formula, with no globally free lattice assumption.

If \(v\ne0\) lies in AK, then \(h(v,v)\) is totally positive and belongs to \(A\bar A\cap F\). The latter ideal has F-norm \(N_EA\); equivalently this follows place by place, at split, inert and ramified relative places. The principal norm is therefore at least \(N_EA\). AM--GM gives

\[
 \|v\|^2\ge f(N_EA)^{1/f}.
\]

For the Euclidean dual in IQ17 this becomes

\[
 \lambda_1((IK)^\#)^2\ge \frac{f}{r_0^2z^2},
 \qquad r_0=D_0^{1/f}<365729/500.
 \tag{IQ18}
\]

For positive embedding weights c\_tau and \(B=(\prod c_\tau)^{1/f}\), the same product norm gives the weighted dual minimum

\[
 \sum_\tau c_\tau^{-1}\|v_\tau\|^2
 \ge\frac{f}{r_0^2z^2B}.
 \tag{IQ19}
\]

The minimum has f, not mf, on the right. The larger real dimension will be paid explicitly in the next theta bound. IQ17 for m=3 is precisely where the old covolume \(D_0^3z^{6f}/U^f\) is replaced by \(D_0^3z^{6f}\).

\subsection{6. A genuine dimension1008 transverse Gaussian bound}\label{a-genuine-dimension1008-transverse-gaussian-bound}

DO114 SP3--8 proves for every real n-dimensional lattice with minimum at least lambda, and every t\textgreater0,

\[
 \Theta(t)-1\le\frac{2n^2}{c_n}2^{n/2}\Gamma(n/2+1)
                           (t\lambda^2)^{-n/2},\quad
 c_n=\frac{\Gamma((n+1)/2)}{\sqrt\pi\Gamma(n/2)}.
 \tag{IQ20}
\]

Its finite separated-set proof lifts to S\^{}n, projects cap points into dimension n, applies the2n orthoplex bound, and integrates only nonzero lattice shells of radius at least the minimum before enlarging the integral. Thus IQ20 is valid for the nonrectangular projective quotient lattice; no scalar-dimension constant is reused in dimension1008.

Put n=6f=1008, \(c_n=504\binom{1008}{504}/4^{504}\), \(\pi_-=314159/100000\), and define the explicit rational constant

\[
 K_3=\frac{2\cdot1008^2}{c_{1008}}2^{504}504!
       \left[\frac{(365729/500)^2}{168\pi_-^2}\right]^{504}.
 \tag{IQ21}
\]

Let \(\theta=2^{-20}\). Gaussian Poisson summation and IQ17--20 give, for the full transverse row lattice and arbitrary positive embedding weights,

\[
 \sum_{r_V\in I\mathscr K_{\eta,V}^{\vee}}
 e^{-(1-\theta)U\sum c_\tau\|r_{V,\tau}\|^2}
 \le\frac{\pi_+^{3f}}{D_0^3}
 \left\{\left[\frac{q}{(1-\theta)RB}\right]^{3f}+K_3\right\},
 \quad q=\frac{R}{Uz^2},\quad\pi_+=355/113.
 \tag{IQ22}
\]

The Fourier convention is \(\widehat{e^{-t\|x\|^2}}(\xi)=(\pi/t)^{n/2}e^{-\pi^2\|\xi\|^2/t}\). The weighted dual minimum gives the term \(K_3[(1-\theta)Uz^2B]^{3f}\) before multiplication by the Poisson volume. Both terms in IQ22 are retained for every B and q; no small-parameter or dense-only assumption is used.

For each of the k standard S coordinates use the unchanged K111/DO125 scalar coefficient

\[
 K=\frac{336M_{15}}{c_{336}}(7841/5000)^{168}168!
 \left[\frac{(365729/500)^2}{168\pi_-^2}\right]^{168},
\] \[
 c_{336}=168\binom{336}{168}/4^{168},\quad
 M_{15}=\frac{6015684926187167549292201834533348561577984}
                 {2618017117694612446410293519}.
 \tag{IQ23}
\]

Its Gaussian factor is \((\pi_+^f/D_0)\{[q/((1-\theta)RB)]^f+K\}\). Orthogonality IQ11 permits the product with IQ22; no rectangular model is imposed on the rank3 factor.

\subsection{7. The actual central integral and positive mixture are unchanged}\label{the-actual-central-integral-and-positive-mixture-are-unchanged}

We retain the entire accepted DO125 CM3--16 analytic argument, with its original selected module and both compact factors. Put

\[
 \vartheta=\frac{4^{167}}{334\binom{334}{167}},\quad
 \kappa_I=\frac{4}{\sqrt{527}}\frac{23479}{23450},\quad
 \mathcal G_1=\frac{2\cdot4^{83}82!83!}{166!},
\] \[
 V_* =\frac{97}{100B_0-97}\frac{2113^2}{64\cdot2049},\quad
 \mathcal D_*=334^2\frac{400}{399}V_*+528\cdot167^2\,2^{-2800},
 \quad B_0=2^{955},
\] \[
 D_*=(1480/1653)^{333/2}e^{1557/80}\sqrt{1670/1757},
\] \[
 \Phi(t)=t^{-1/2}(1+t/4)^{-166}(1+t/2)^{-1/2}
       (1+t/(3514/29))^{-1/2}e^{-29t/20}.
 \tag{IQ24}
\]

The physical central coordinate is still \(w_{\rm phys}=\Im N_{nn}=-U\Im d\). For fixed eta,row and actual G pattern, its admissible scalars lie in the original affine difference coset with

\[
 \operatorname{covol}\Gamma_w
 =U^f\sqrt{D_F}(N_FJ_0)xT^{-54}.
 \tag{IQ25}
\]

This follows from the original primitive4-by4 entry extraction, not from the quotient-lattice lemma. In particular the actual central factor U\^{}-f is independent of IQ17 and remains in the complete count. The affine real-part shift, possible dyadic shift, full-real750 covering bound and J0 divisor are retained.

The same original coarea kernel, exponential tangent before transport, support-plane inequality in the physical central variable and normalized Beta--Gamma Jensen bound give central integral at most \(\mathcal G_1D_*\Phi(\alpha)\), multiplied by the actual spread factor \(\vartheta\kappa_I\sqrt{\mathcal D_*}\) at each embedding. This is one containing coefficient bound, not a product of independent trace estimates. The near-point and central Voronoi inflations remain below2\^{}(1/1024) each, by the original non-strict H133 cutoff and \(Uz_c<2^{-83}\). K120/K118 gives the entire omitted narrow-band contribution less than2\^{}-8605114, to be charged once for the final interval.

Phi is a positive Laplace convolution with infinite total mass: the probability Gamma factors have shapes166,1/2,1/2 and rates4,2,3514/29, convolved with ds/sqrt(pi s) and the fixed shift29/20. For every j\textgreater1/2 its inverse moment is

\[
 M_j=\frac1{\Gamma(j)}\int_0^\infty
 t^{j-3/2}(1+t/4)^{-166}(1+t/2)^{-1/2}
 (1+t/(3514/29))^{-1/2}e^{-29t/20}\,dt.
 \tag{IQ26}
\]

There is no zeroth inverse moment. On the original genuine row, the unchanged k nonzero S coordinates give \(\prod\alpha_\tau\ge(Rk/q)^f\). Reserve theta of the genuine Gaussian exponent before any row enlargement. The exact DO125 CM21 AM--GM repair gives, for the unique degree-zero term,

\[
 \int e^{-\theta\sum c_\tau\alpha_\tau}d\mu^{\otimes f}
 \le\overline M_1^f[q/(\theta Rk)]^f.
 \tag{IQ27}
\]

No per-embedding alpha floor is imposed on subsequently enlarged tuples.

\subsection{8. Exact moments and the complete new row polynomial}\label{exact-moments-and-the-complete-new-row-polynomial}

The executable reproduces the DO125 positive quadrature at j=148,162 and163: interval{[}1,64{]}, mesh1/128, exact second-derivative trapezoid error on every subinterval, complete{[}0,1{]} and{[}64,infinity) tails, outward integer exponentials and square roots. The proof of the right tail only requires j\textless337/2 and therefore applies to all three degrees. The exact rational upper bound for M1 is the accepted K120 j1 positive rational moment times355/113; dropping the two extra half factors and replacing29/20 by32/25 proves its containing scope. No moment is inferred from a floating integral.

For j between1 and148, and then between148 and162, Holder on the same positive measure gives

\[
 M_j\le M_a^{(b-j)/(b-a)}M_b^{(j-a)/(b-a)},
 \quad(a,b)=(1,148)\ \text{or}\ (148,162).
 \tag{IQ28}
\]

Both endpoints are finite; no interpolation against an infinite M0 is used. These define the positive upper bounds \(\overline M_j\). M163 is used only in the failed-next diagnostic.

Set d=k+3. Combining the k scalar factors with IQ22 and integrating the positive convolution retains the complete polynomial

\[
 \begin{split}
 \mathcal P_k(q)={}&\sum_{l=0}^k\binom kl K^{k-l}
   [q/((1-\theta)R)]^{(l+3)f}\overline M_{l+3}^f\\
 &+K_3\sum_{l=1}^k\binom kl K^{k-l}
   [q/((1-\theta)R)]^{lf}\overline M_l^f\\
 &+K_3K^k\overline M_1^f[q/(\theta Rk)]^f.
 \end{split}
 \tag{IQ29}
\]

Apart from \(\binom{164}{k}\pi_+^{df}/D_0^d\), this is the full row sum. There are exactly2(k+1) positive terms, including the repaired zero term. Unlike the former enlarged transverse value count, no U\^{}f precedes this polynomial.

\subsection{9. Every finite, source, central and coupled-endpoint factor}\label{every-finite-source-central-and-coupled-endpoint-factor}

The full original Green estimate and common ideals are unchanged:

\[
 |a_G(h)|\le2Q^{-335\cdot167}(N_EJ)^{335/2},\quad
 N_EJ\le Q^{334},\quad N_EJ\le Q^{167}N_FJ_0.
 \tag{IQ30}
\]

The support count is exactly \(\binom{168}{27}\), the retained first-pivot multiplicity is164, and both mixed finite factors keep their modulus-one bounds. Eta,r,d reconstruct the actual matrix injectively, so no original beta, negative-column or source-isometry fibre is discarded. The complete active one-label and inactive factors are \((T-1)^{-168}\), constant on the actual central fibre because b=j\_int=0. Their T exponent after combining Green and IQ25 is335·27+54−168=8931.

For positive degree j in IQ29, or effective j=1 for its repaired zero term, the varying factor before the explicit J0 comparison is

\[
 x^{-337/2}q^{(j-335/2)f},\qquad
 q_*^f=\frac{R^fT^{54}}{U^fQ^{334}x}.
 \tag{IQ31}
\]

All j here are at most162. The displayed q derivative is negative. More explicitly, put rho=(N\_EJ)/Q\^{}334=(q\emph{/q)\^{}f.~The second comparison in IQ30 gives (N\_FJ0)\^{}-1\textless=Q\^{}-167 rho\^{}-1. Relative to the all-zero G cap, the full Green, row and central term therefore gains rho\^{}(333/2−j), whose exponent is positive. This proves the required endpoint comparison after central correlation as well. At q} the coupled x derivative is exactly−1−j. Thus the full window is bounded at x=2\^{}2111 and the all-zero G cap. The complete four-state pattern sum for each term is

\[
 \prod_{v=1}^{334}(1+q_v^{j-333/2})^2
 \le(1+1/g_0)^{668}<2^{1/1024}.
 \tag{IQ32}
\]

This retains both rank-one orientations and all rank0/rank1/rank-at-least2 states. There is no extra2\^{}668 pattern multiplier. The full T correction is also below2\^{}(1/1024). Together with near-point and central cover, exactly four inflations of this size occur. The quotient lattice may vary with eta but not with the G pattern; its universal bound is uniform in that lattice.

The exact endpoint brackets are

\[
 q_-=251124964<q_*<251124965=q_+,
 \tag{IQ33}
\]

verified by168th powers. For each k define the entire common prefactor

\[
 \begin{split}
 \mathcal A_k={}&2\,2^{4/1024}\binom{168}{27}164\,
 \mathcal D_*^{84}\vartheta^{168}
 (\kappa_I\mathcal G_1D_*)^{168}D_F^{-1/2}U^{-f}\\
 &\times(R/U)^{335f/2}T^{8931}Q^{-336\cdot167}
 2^{-337\cdot2111/2}\binom{164}{k}\pi_+^{df}D_0^{-d}
 q_-^{-335f/2}.
 \end{split}
 \tag{IQ34}
\]

The initial2 is exactly Green2. For every positive-degree summand use q+ in its positive power in IQ29; use q+ also in IQ27. The independent central U\^{}-f in IQ34 is the full arithmetic saving from IQ13 and IQ17. No physical-coordinate Jacobian is absorbed into a new row coordinate.

\subsection{10. Exact interval union and the next failed cell}\label{exact-interval-union-and-the-next-failed-cell}

The complete certificate uses160-bit directed logarithms with100-term atanh series and full geometric remainders, as in DO125. It proves the actual Machin upper bound pi\textless355/113 and lower bound pi\textgreater314159/100000, the exact discriminant/root bound, all host coprimalities, both q brackets, every moment enclosure and both monotonicities. All coefficients and signs in IQ34 are retained before directed evaluation.

For every one of the

\[
                 \sum_{k=145}^{159}2(k+1)=4590
 \tag{IQ35}
\]

positive terms of IQ29 times IQ34, the directed logarithmic upper bound is strictly below−595. The worst term is the main l159, degree162 term, with directed upper ceiling −19510525/32768. This is a complete containing sum, including all transverse theta corrections and every repaired zero term; no main-term approximation is used.

The whole original narrow-band tail is paid once. Since it is less than2\^{}-595, the full interval is bounded by

\[
 4590\,2^{-595}+2^{-8605114}
 <4591\,2^{-595}<2^{-582}.
 \tag{IQ36}
\]

The unchanged predecessor arrays are preserved in order,582 is appended once, and both complete error sums remain below2\^{}-230. The k160 main term with the exact degree163 upper moment has directed lower bound68793421/65536\textgreater0 for this chosen sufficient envelope; no k160 or higher removal is asserted.

\subsection{11. Source, geometric and verification boundary}\label{source-geometric-and-verification-boundary}

The new algebraic content is IQ8--19: saturation in the \textbf{actual} standard ambient lattice, the nonprincipal selfdual quotient, every-prime row containment, and its exact covolume/minimum in the unchanged positive metric. The new counting content is the dimension1008 theta factor IQ21--22 and its complete two-term replacement IQ29. All fixed-source costs are inherited with exact scope from R025 R25.1--3; R072 PJ4--6; DO098 CF4--13; DO104 CD10--19; DO106 CFI3--8 and its section2 opening compact statement; K113 SC27.4--10; DO125 CM3--28; and the universal finite-dimensional packing proof DO114 SP3--8. K105 CG5 fixes the exact half-trace/codifferent convention. The copied quadrature proof is verified anew at its three stated degrees.

The isotropic quotient is a projective Hermitian module over Spec(O\_E). It changes no original matrix, arithmetic source, labelled pencil, polarization, theta line, Picard/Poincare data, CM involution or selected real realization. No new motive, weight, Dieudonne, lambda-ring or K-theory comparison is claimed. The carried contextual anchors remain SGA3 I4.7.3/4.7.3.1; SGA6 VI Theorem3.3; FGA232 Theorem3.1 pp148--149 and FGA236 Corollaries1.7/2.3 pp229/232; SGA7 I VII2.7/2.9 and VIII7.4(b); cote4 sections8.1.1--8.1.2 and8.4.1--8.4.3, and cote54. They are inherited context, not a new primary reading.

The actual selected representation dictionary remains Nakahama(2016), equation5.7 and Theorem5.1, pp960--961, through the accepted complete analytic interfaces. Rankin(1955), Theorem1(iv), is the classical orthoplex reference; the required finite-dimensional estimate is proved in the accepted SP3--8 argument and applied here with its correct new dimension. No scalar coefficient substitutes for the selected module. No source admission, signed nonzero, prescribed period, specific032 value, evaluated DP44/RH conclusion, complete floor28, Lean build or archive audit follows.

Development floating screens, the initial D0-to-float overflow, exact moment-only probes, preliminary full replay and final author replay are recorded separately. Finite models and arithmetic diagnostics supplement the complete proof and do not establish its projective-lattice statements by enumeration.

\subsection{Canonical algebraic-geometric context}\label{canonical-algebraic-geometric-context-85}

The new object is the saturated isotropic quotient (ell-perpendicular intersect Lambda)/(ell intersect Lambda) over Spec(O\_E), with its descended perfect Hermitian pairing and nonprincipal projective module. Integral local plane splitting uses determinant−1 and no division by2. Its positive metric and exact codifferent covolume retain the original polarization, CM involution, cubical line, Picard/Poincare data and all finite phases. The actual arithmetic arrow is only placed in a necessary containing module. Inherited context only: SGA3 I4.7.3/4.7.3.1; SGA6 VI Theorem3.3; FGA232 Theorem3.1 pp148--149; FGA236 Corollaries1.7/2.3 pp229/232; SGA7 I VII2.7/2.9 pp169--174 and VIII7.4(b) pp310--311; cote4 scan003 sections8.1.1--8.1.2 and scans005--006 sections8.4.1--8.4.3; cote54 scans005--007. No new Dieudonne, lambda-ring/K-theory, l-adic-weight or CM-motive comparison. DP\_MULT Q3e/DP44 and THETA3 G1/A7 gain the strictly smaller conditional finite gate; CX\_GRH CX-ALG-024 and GCT\_PLETH T-052 retain their arithmetic/global gates.

\clearpage

\section{AM-LIFT-105 --- The standard-lattice short theta unit}\label{am-lift-105-the-standard-lattice-short-theta-unit}

\textbf{Canonical checked integration --- 2026-11-23T12:58:50.923946+00:00.} Complete saved-proof and executable reviews and separate exact diagnostic replays passed. No numerical error charge is introduced. Frozen author history is preserved. AM-LIFT-105 SU1-22 proves the literal ShortThetaUnit for the standard auxiliary lattice O\^{}n and any actual operative IsNab on the original geometric data. The same unit is the inverse of the actual selected beta0, independent of N, x0, polynomial and every original ring/witness. Global normalized multiplication/primitive identities give same-point ordinary derivative scaling without added truncation coherence. The complete standard A{[}N{]}\^{}n deck trace cancels exactly mIdx in the source action order, every allowed N-root gives its allowed M-root, and RowDeg cancels the full N/M degree factor. Harmonic all-ring comparison and faithful-flat injectivity retain nilpotents and every original witness. Nonstandard lattice beta laws, original-base descent, source/package, cusp/spanning, selected arithmetic nonzero/content, period, specific032 value and evaluatedDP44/RH remain separate. No numerical charge;136/31 complete arrays and floors27/11 unchanged.

\textbf{Status: final author subject, 2026-11-23; independent and root reviews pending.} The complete author proof/source/script review and exact diagnostic records are in the adjacent final packet, which identifies the frozen bytes and actual completion times. The claim concerns the literal \nolinkurl{ShortThetaUnit} source law for the original standard lattice and an actual pinned operator. It does not construct an arithmetic source, a covering-stage descent datum, or a nonzero selected coupled image.

\subsection{1. Exact statement and quantifiers}\label{exact-statement-and-quantifiers}

Retain actual Host, Datum, odd Level, Coeff and Moduli records from the supplied source, their original integral CM action and symmetric rigidified principal line, and the accepted formal gauge theorem used in LIFT101. Write

\[
 O=\mathcal O_E,\quad F=c.F,\quad M=l.M,\quad
 \mathfrak b=l.Mb,\quad r=Dw.p+Dw.q,\quad g=[E:\mathbf Q]r/2.
 \tag{SU1}
\]

Let \nolinkurl{nab} satisfy the literal \nolinkurl{IsNab Xo c nab}, with its required involution equal to H.bar. For every n, the conclusion is the original proposition

\[
 \operatorname{ShortThetaUnit}(X_o,c,\mathrm{nab},n).
 \tag{SU2}
\]

More explicitly, put \(\Lambda_0=O^n\), exactly \nolinkurl{GenLat ⟨stdLat E n⟩}, and retain the actual chosen scalar

\[
 \beta_0=\operatorname{betaH}(c,\Lambda_0).
 \tag{SU3}
\]

For every row weight a, positive odd N with M dividing N, x0 in Dw.La\^{}n, and P in the literal \nolinkurl{harmA c a n}, the unit in \nolinkurl{ShortThetaUnit} may be taken to be \textbf{the same}

\[
 u=\beta_0^{-1}\in\mathcal O_F^\times.
 \tag{SU4}
\]

It is independent of N, x0, P, the coefficient algebra, the frame point and all witnesses. Write Q=(Fin n → LevQ Dw l). For every commutative F-algebra S, every actual y∈FrPt(S), every v∈S and every function w:Q→S satisfying the two original \nolinkurl{ShortValAt} and \nolinkurl{ThValAt} hypotheses, the exact equality is

\[
 v=\beta_0^{-1}|Q|^{-1}
   \sum_{q\in Q}\zeta_M^{\operatorname{pairL}(x_0,\operatorname{levRep}(q))}w(q).
 \tag{SU5}
\]

All coefficient-field scalars in this display are mapped to S. The phase is the positive exponent in \nolinkurl{x0Dat}; the cardinality and representatives are the original source definitions. Nilpotents, disconnected algebras and non-Noetherian algebras are included. The zero ring, n=0, P=0, and g=0 are included with their literal empty-product meanings. No assertion that the selected beta equals 1 is needed.

\subsection{2. The actual beta and index factors}\label{the-actual-beta-and-index-factors}

For Λ0, its intersection Ln with the standard lattice is Λ0 itself. Every determinant in \nolinkurl{detSet Λ0} is integral over Z, and 1 is in that set, including the empty determinant for n=0. Each embedding embF therefore gives

\[
 \operatorname{idealAt}(c,\Lambda_0,\tau)=\mathcal O_F.
 \tag{SU6}
\]

The scalar 1 satisfies both \nolinkurl{BetaIs} and \nolinkurl{BetaAbs}: their two lattice products coincide and the remaining additive index is 1. Thus the successful branch of \nolinkurl{betaH} applies, and its chosen β0 has

\[
 \beta_0\ne0,\qquad \mathcal O_F\beta_0=\mathcal O_F.
 \tag{SU7}
\]

This equality implies β0 belongs to O\_F and has an inverse in O\_F. It proves (SU4) without replacing the source's noncomputably selected beta. The additional actual BetaAbs identity says every complex absolute value of β0 is 1; that fact is not needed for (SU5).

The standard lattice is unimodular for the original identity Hermitian form and is \nolinkurl{TrivAt bad} for every bad, using k=1. Its literal LevOK is precisely odd N, M\textbar N and the automatic inclusion NO\^{}n⊂O\^{}n.~As O has Z-rank 2f,

\[
 zCard(\Lambda_0,N,r)=mIdx(\Lambda_0,N,r)
   =|O^n/NO^n|^r=N^{2gn}>0.
 \tag{SU8}
\]

In particular the index used in the denominator is the full kernel cardinality, not its square root. Q is finite and nonempty, since L0 is a finite projective O-lattice and O/b is finite; the zero tuple supplies an element. Thus all denominator scalars in (SU5) and \nolinkurl{ShortValAt} are nonzero rational units in F.

\subsection{3. The normalized multiplication maps and the primitive frame}\label{the-normalized-multiplication-maps-and-the-primitive-frame}

Let A be the original universal abelian scheme and L its original symmetric rigidified line. Use the global normalized multiplication maps of LIFT100 HN7/LIFT103 ST12,

\[
 \alpha_d:[d]^*L\longrightarrow L^{d^2}
 \qquad(d\ge1).
 \tag{SU9}
\]

Their value on the zero section is precisely \nolinkurl{IsNormMul}, namely the original rigidification raised to d²−1. These are maps on the universal group before any pullback to y. Any two maps with the same source, target and normalization agree: their ratio is a global unit on A, hence a base unit, and its zero-section value is 1. The same holds on every power A\^{}j and after arbitrary base change. This is the actual properness/connected-fibre argument of LIFT101 AT3; its primary input is Stacks Lemma36.32.6.

Let B(X,Y) be the source's two-variable cubical line, with its normalized biextension laws. Iterating those laws gives a global map

\[
 \psi_d:([d]\times[d])^*B\longrightarrow B^{d^2}.
 \tag{SU10}
\]

The axes and zero values are canonical. Commutativity and associativity follow from the original cubical maps; alternatively their competing global maps have the same normalization and the preceding uniqueness fixes their scalar. The integral endomorphism w0 commutes with {[}d{]}. The original normalized primitive is an isomorphism

\[
 w_P:B(X,[w_0]X)\longrightarrow L_X.
 \tag{SU11}
\]

This is the orientation in \nolinkurl{ThGerm.eL}: the generator of L is wP times the generator of B. The two global maps from {[}d{]}\^{}*B(X,{[}w0{]}X) to L\^{}(d²) have the same zero normalization, hence

\[
 \alpha_d\circ[d]^*w_P
   =w_P^{\otimes d^2}\circ\psi_d.
 \tag{SU12}
\]

No equality at an arbitrary single fibre is used to conclude a universal identity. Both sides of (SU12) are universal global isomorphisms and are normalized along the universal zero section. This also proves compatibility with the normalized symmetry. Applying the same argument directly to the maps {[}ab{]}\^{}*L→L\^{}(a²b²) yields

\[
 \alpha_{ab}=\alpha_a^{\otimes b^2}\circ[a]^*\alpha_b.
 \tag{SU13}
\]

There is no free phase or coefficient-field scalar in these identities.

Here is their exact implication for an arbitrary literal pinned witness; no extra biframe-truncation axiom is assumed. Fix one finite order D sufficient for the derivatives in question. LIFT101 AT8--11 puts its chart in the additive coordinates determined by the physical tangent frame and gauge-corrects its biframe by E\_D=e\_D exp(−h), modulo order D+1. The original addition, symmetry and adjunction multipliers then equal one. Iterating each addition law d times, at this same finite order, gives

\[
 \psi_d([d,d]^*E_D)=E_D^{d^2}.
 \tag{SU14}
\]

The corrected primitive line frame is e\_D\^{}L=wP E\_D(Z,w0Z). Equations (SU12) and (SU14) prove

\[
 \alpha_d[d]^*e_D^L=(e_D^L)^{d^2}\pmod{(Z)^{D+1}}.
 \tag{SU15}
\]

The same actual ChartAt is available for every nonnegative line power and every section by \nolinkurl{IsNab}. If F is a section of L\^{}k and F\_d=α\_d\^{}k{[}d{]}\^{}*F, its corrected germ in that witness is therefore f(dZ) when F has corrected germ f(Z). Invariant physical derivatives become ordinary derivatives in these additive coordinates, with derivative of {[}d{]} exactly d times the identity. Thus, for every multiindex η,

\[
 \mathrm{nab}_y(kd^2,F_d,\eta)
    =d^{|\eta|}\mathrm{nab}_y(k,F,\eta).
 \tag{SU16}
\]

Choose D≥\textbar η\textbar{} to justify this assertion. It is a pointwise identity within the same supplied pin and does not assert arbitrary-base-change naturality of individual ordinary coefficients. The finite-order argument covers the noncoherent biframe sequences permitted by the actual source. Ordinary factorials occur identically on both sides; no divided derivative has been introduced.

\subsection{4. The full standard kernel fixes the pulled-back section}\label{the-full-standard-kernel-fixes-the-pulled-back-section}

For an n-tuple z of torsion points, the exact source predicate specializes to

\[
 \operatorname{InZh}(\Lambda_0,N,z)
       \quad\Longleftrightarrow\quad z\in A[N]^n.
 \tag{SU17}
\]

For the forward direction take mu=I and cm=N I in \nolinkurl{InZh}. The resulting row equation is Nz\_b=0. Conversely a matrix mu carrying O\^{}n into O\^{}n has all entries in O, by applying it to the standard vectors. Every permitted cm is Nmu, so each defining row sum vanishes on N-torsion. The N²-torsion condition also follows.

By LIFT103 ST5--7 and ST14, every permitted GeoInj witness in \nolinkurl{ShortValAt} enumerates the \textbf{entire} finite étale subgroup A{[}N{]}\^{}n, of rank (SU8). This includes a disconnected base, where the comparison permutation may vary between open-and-closed pieces.

Write s\_N=α\_N{[}N{]}\^{}\emph{s for the actual source section. On {[}N{]}\^{}}L there is a canonical deck lift of each z∈A{[}N{]}, since {[}N{]}∘t\_z={[}N{]}. It fixes {[}N{]}\^{}*s. Its Nth power is the identity, and inversion sends the lift at z to its inverse. Transport it through the normalized symmetric α\_N. The resulting lift on L\^{}(N²) has point z, N²th power 1 and the source symmetry. It is consequently the unique normalized odd \nolinkurl{SymLift} of LIFT103 ST12--13. Therefore

\[
 U_z^{(N)}s_N=s_N,
 \qquad
 U_t^{(N)}U_z^{(N)}s_N=U_t^{(N)}s_N.
 \tag{SU18}
\]

This uses the literal order z first, t second in \nolinkurl{IsUUOver}. We do not commute U\_z and U\_t or discard a commutator phase. Every summand section allowed by the original relation is thus independent of its z label. The sum over labels contributes exactly mIdx, and cancels precisely its inverse in \nolinkurl{ShortValAt}.

\subsection{5. Every allowed N-root gives an allowed M-root}\label{every-allowed-n-root-gives-an-allowed-m-root}

Write N=dM, where d≥1 is odd. For every source label q and column, take an arbitrary permitted short-value root t with Nt=η(levRep q), killed by some power of b. Put t\_M=dt. Then Mt\_M=η(levRep q). The latter point is killed by b. Because b and bar(b) are coprime and b bar(b)=(M), it follows that

\[
 \mathfrak b^2t_M=0.
 \tag{SU19}
\]

For clarity, at a prime dividing b the bar(b) factor is a unit, so bMt\_M=0 gives b²t\_M=0. At the other primes t\_M has no component, since it is killed by a power of b. Equivalently this is the primary decomposition of the same finite étale torsion group, not a choice of a principal generator for b. Thus t\_M satisfies exactly \nolinkurl{KilledBy(l.Mb^2)} in \nolinkurl{IsVtOver}.

Let U\_M be its unique symmetric lift on L\^{}(M²) with M²th power one, and set F\_M=U\_M s\_M. Pull this lift back by {[}d{]} along the pair (t,U\_M), using {[}d{]}t=t\_M, and transport by α\_d\^{}(M²). This produces a theta lift at t on L\^{}(N²). Its N²th power is identity: t has N²t=0, while U\_M\^{}(N²)=1 because M² divides N². Its symmetry follows from the normalized symmetry of α\_d.~These are the required N² conditions, not merely the M² conditions of the original lift. Odd-lift uniqueness therefore identifies it with U\_N at t.

This construction can be described without an informal lift choice: pullback gives a homomorphism from the fibre product of the theta group of L\^{}(M²) with A under {[}d{]}, whose elements are (t,U) with {[}d{]}t=pt(U), into the theta group of {[}d{]}\^{}*L\^{}(M²). It preserves multiplication, the identity and inversion. The N²-power check is made in that fibre product. Transport through α\_d\^{}(M²) is an isomorphism of theta groups preserving the normalized symmetry. It therefore works over every coefficient ring.

Together with (SU13), it proves the exact section equality

\[
 F_N:=U_Ns_N=\alpha_d^{M^2}[d]^*F_M.
 \tag{SU20}
\]

Both sections satisfy the original literal relations, with every allowed normalized multiplication generator and lift agreeing by LIFT103 uniqueness. No assumption that {[}d{]} is an isomorphism of abelian schemes, or that d is prime to M, is used.

\subsection{6. The total-degree normalization and arbitrary witnesses}\label{the-total-degree-normalization-and-arbitrary-witnesses}

For P≠0, its RowDeg condition makes every support monomial have the same total degree

\[
 D=P.\operatorname{totalDegree}
    =\sum_{\tau}\sum_{i<\operatorname{mult}(\tau)}a(\tau,i).
 \tag{SU21}
\]

This follows by summing the source row degrees over all rows and columns. There is no assertion that each column has the same degree. If m is any support monomial, the sum of the degrees of its actual \nolinkurl{colOf m b} is D. Apply (SU16) to (SU20) in each column and multiply. The unnormalized entire polynomial contraction is multiplied by d\^{}D. Since N=dM,

\[
 N^{-D}\sum_m P_m\prod_b\mathrm{nab}(N^2,F_{N,b},\mathrm{col}_b m)
 =M^{-D}\sum_m P_m\prod_b\mathrm{nab}(M^2,F_{M,b},\mathrm{col}_b m).
 \tag{SU22}
\]

The right side is precisely \nolinkurl{thExpr} at the M-root constructed in section5. For P=0 both support sums are zero, so (SU22) holds without identifying a nonempty support degree. For n=0 the same calculation concerns the scalar constant polynomial, and all section products are empty.

Now take any original \nolinkurl{ShortValAt} witness S→S′ with its supplied FrMap point. The root and section calculations above apply to its actual roots and every one of its allowed sections. Its full z sum cancels mIdx by section4. For each q, (SU22) is a literal \nolinkurl{ThValAt} witness over S′, using the identity faithfully flat extension there and the M-root t\_M. LIFT101's harmonic all-ring uniqueness and FrMap naturality apply, since \nolinkurl{harmA} is contained in identity harmK. They identify this value with the image of the separately supplied w(q) over S. This use of harmonic naturality is distinct from the pointwise ordinary scaling in (SU16).

Substitution into the defining equality for the short witness gives the image of (SU5) in S′. A faithfully flat map of rings is injective, so (SU5) holds in S. This proves agreement for \textbf{every original} short witness and every original family of theta-value witnesses, not only for witnesses built in this proof. It is an equality in S itself; comparison of geometric points is not substituted for scalar descent over a ring with nilpotents.

\subsection{7. The precise source-law consequence}\label{the-precise-source-law-consequence}

The unit (SU4) now supplies the original existential in \nolinkurl{ThShort.ShortThetaUnit} with its correct order of quantifiers. Thus (SU2) follows from an actual IsNab pin and the original geometric records. In scopes where LIFT101 is allowed to choose the operator, its constructor also supplies this law for that chosen operator. For a preselected operator, its actual IsNab witness remains required; the law is not inferred for arbitrary NabV.

The standard-lattice statement does not require BetaExists or KPsiIn for arbitrary lattices. It does not replace the actual Fourier character, alter the old level, make a new covering-stage choice, or assume an integral standard basis for the original projective lattice Dw.La. Only the separate auxiliary Λ0=O\^{}n is standard. The original Dw.La remains projective throughout the finite torsion/rank computation.

No claim is made here about original-base sec/cdelta descent, nonstandard lattice short laws, cusp membership, spanning, ShortNe0, or nonvanishing of the original coupled selected map. The actual source/Host/Sizes/package existence, period and evaluated DP44 gates remain. No finite-sector count or numerical error charge changes.

\subsection{8. Exact sources, verification scope, and references}\label{exact-sources-verification-scope-and-references}

The adjacent source manifest records exact full hashes, line totals and actual reading scopes. The essential literal interfaces are:

\begin{longtable}[]{@{}
  >{\raggedright\arraybackslash}p{(\columnwidth - 2\tabcolsep) * \real{0.5000}}
  >{\raggedright\arraybackslash}p{(\columnwidth - 2\tabcolsep) * \real{0.5000}}@{}}
\toprule\noalign{}
\begin{minipage}[b]{\linewidth}\raggedright
Source
\end{minipage} & \begin{minipage}[b]{\linewidth}\raggedright
Exact role
\end{minipage} \\
\midrule\noalign{}
\endhead
\bottomrule\noalign{}
\endlastfoot
\nolinkurl{ThShort.lean}20--32,37--106,244--261,313--323 & Original IsUUOver order, standard lattice and indices, beta choice, matrix kernel, full ShortValAt normalization and exact ShortThetaUnit quantifiers. \\
\nolinkurl{IsThRule.lean}32--43,57--80,85--88 & Original M-root condition, normalized multiplication/pullback, colOf, thExpr, ThValAt and harmA. \\
\nolinkurl{RowDeg.lean}40--48 & Constant row multidegrees, identity harmonicity, and the actual integral polynomial subset. \\
\nolinkurl{ThGerm.lean}108--221 & Actual biframe/cubical/adjunction equations, normalized primitive orientation, corrected germ and one pin witness for all powers/sections. \\
Recovered hash-matched \nolinkurl{ThetaPt.lean} and \nolinkurl{ThetaVV.lean} & SymLift, IsTwistPull, theta action, and the exact origin normalization IsNormMul. \\
Canonical LIFT101 AT3, AT8--17, AT18--25 & Universal global-unit uniqueness, finite-order arbitrary-witness gauge correction, physical additive coordinates, and all-ring harmonic value uniqueness/naturality. \\
Canonical LIFT102 Theorem D & Compatible carried arbitrary-pin harmonic scope; no extra ordinary-coefficient comparison premise. \\
Canonical LIFT103 ST5--15 & Exact finite kernel rank/full enumeration, actual half-primary roots, universal normalized multiplication maps, odd lifts and every-witness section identities. \\
\end{longtable}

The direct new argument is the globally normalized multiplication/primitive identity (SU12), its finite-order actual-pin consequence (SU16), and its complete composition with the standard deck sum and the N/M normalization. No build, proof-assistant execution or archive-validity claim is made. The finite executable diagnostics check arithmetic kernel/rank models, odd-lift normalization, nilpotent coefficient/gauge models, ordinary multicolumn derivative scaling, index cancellation and the original phase/cardinality direction. They do not replace the universal proof.

\textbf{Primary and carried context.} The Stacks Project Authors. (2026). \emph{The Stacks Project}, Section36.32, Lemmas36.32.5--6 (tag0BDM), \url{https://stacks.math.columbia.edu/tag/0BDM} . The full displayed statements and proofs were freshly read for (SU9)--(SU13), including arbitrary base change. The accepted geometric inputs from LIFT101--103 retain their documented Milne abelian-variety and SGA7 cubical/biextension readings; those are carried readings, not a newly asserted full-source check. A fresh attempt to open Milne's CourseNotes/AV.pdf during this author turn timed out and is not counted as a successful new primary read.

The proof uses the actual Picard/Poincare line, rigidification, cubical law and finite étale torsion sheaves. The functorial multiplication and odd-lift arguments are identities of those line bundles and group schemes before differentiation. They are valid over arbitrary characteristic-zero affine coefficient bases, so no reduced-point shortcut is used. No new integral model, Dieudonne slope, motivic realization, cohomological weight, K-theory positivity, or L-value statement is inferred from this source-law identity.

\subsection{Canonical algebraic-geometric context}\label{canonical-algebraic-geometric-context-86}

The new identity is formed on the actual rigidified Picard/Poincare line and cubical biextension over the original characteristic-zero moduli base. The original projective CM lattice is retained; only the separate auxiliary standard lattice is used. Normalized multiplication and odd lifts are universal line/group-scheme identities before the full harmonic contraction. Exact Alpoge interfaces: \nolinkurl{ThShort.lean20-32,37-106,244-261,313-323;} \nolinkurl{IsThRule.lean32-43,57-80,85-88;} RowDeg.lean40-48; ThGerm.lean108-221; recovered \nolinkurl{ThetaPt.lean44-70,73-95,152-161} and ThetaVV.lean103-108. Primary normalization input: Stacks36.32.5-6, tag0BDM. Accepted LIFT101 AT3,AT8-21, LIFT102 TheoremD and LIFT103 ST5-17 retain their complete prior scope. Carried exact context: SGA7 I VII2.7/2.9 pp169--174 and VIII7.4(b) pp310--311; FGA232 Theorem3.1 pp148--149 and FGA236 Corollaries1.7/2.3 pp229/232; SGA3 I4.7.3/4.7.3.1; SGA6 VI Theorem3.3; cote4 scan003 sections8.1.1--8.1.2 and scans005--006 sections8.4.1--8.4.3; cote54 scans005--007. These are carried contextual readings, not new whole-text checks. No new integral extension over Z, Dieudonne slope, lambda-ring positivity, K-theory filtration, l-adic weight theorem or CM-motive L-function comparison follows. DP\_MULT Q3e/DP44/DP52/DP54--65 and THETA3 G1/A7 gain this conditional source-rule reduction; CX\_GRH CX-ALG-024 and GCT\_PLETH T-052 retain their global gates.

\clearpage

\section{AM-R-101 --- Canonical half-level descent of the original short rule}\label{am-r-101-canonical-half-level-descent-of-the-original-short-rule}

\textbf{Canonical checked integration --- 2026-11-23T13:04:52.861846+00:00.} Complete saved-proof and executable reviews and separate exact diagnostic replays passed. No numerical error charge is introduced. Frozen author history is preserved. AM-R-101 FD1-23 constructs both literal ShortValTot conjuncts for every original admissible lattice, polynomial and frame witness when actual IsLev, new A=1, operative IsNab and PhiPin/PsiTie are supplied. The principal theta section module is an invertible line with arbitrary base change; the exact physical frame scheme represents the literal Rules and field evaluator. Half-level rigidity and reduced-base Moduli.rep construct a canonical finite-etale half-level torsor independently of stored d.r. Its exact Sim/LevAut action gives the positive psi Fourier character, inverse psi full trace and cancelling PsiTie scalar. The full harmonic cover rule descends effectively to an original-base sec(a,n), and every original FrOver witness is identified with its canonical framed pullback. No global d.r equality or circular full LawsSh premise is assumed. Actual CovPin supplies the first three hypotheses; dPhi supplies the character package only for n\textgreater=r. The separate small-n supplier, source/package existence, cusp/spanning, selected arithmetic nonzero/content, period, specific032 value and evaluatedDP44/RH remain open here. No numerical charge;136/31 arrays and floors27/11 unchanged.

Author: \nolinkurl{/root}, ALP\_MACHINE. \textbf{Proposed complete argument; not yet frozen, author-verified, independently reviewed or canonically accepted.} This file is research work. The theorem below is conditional on the exact named data; it does not construct an admitted source or assume a missing global compatibility for the stored map.

\subsection{1. Exact hypothesis and conclusion}\label{exact-hypothesis-and-conclusion}

Retain an actual \nolinkurl{Host}, \nolinkurl{Datum}, old odd \nolinkurl{Level}, \nolinkurl{Coeff} and \nolinkurl{Moduli}, and an actual \nolinkurl{DeltaDat d}. Put

\[
 O=\mathcal O_E,\quad L_0=Dw.La,\quad b=l.Mb,\quad
 b_2=d.l_2.Mb,\quad M=l.M,\quad M_2=d.l_2.M,
 \quad m=d.pd,\quad I=d.I. \tag{FD1}
\]

Assume \nolinkurl{d.IsLev}, \nolinkurl{d.l₂.A=1}, an operative \nolinkurl{IsNab} for the stored \nolinkurl{d.nab₂}, and actual \nolinkurl{ph,psi,x₀} satisfying \nolinkurl{d.PhiPin ph psi x₀} and \nolinkurl{d.PsiTie psi}. The coefficient fields are identified by the original \nolinkurl{d.eF}, not by an independently chosen embedding. For the desired number of columns, retain exactly the lattice, polynomial and admissibility hypotheses in \nolinkurl{ShortValTot}.

\textbf{Proposed theorem FD2.} These inputs construct both conjuncts of the literal \nolinkurl{ShortValTot}. In particular there is an original-base rule

\[
 f\in\operatorname{sec}(X_o,c,a,n)
 \quad\text{with}\quad
 \operatorname{IsShortRule}(d,\Lambda,N,x_0,P,f). \tag{FD2}
\]

The all-ring first conjunct is accepted LIFT103. The new assertion is the second conjunct for the particular cover rule of LIFT104. No uniqueness of a section beyond its determined evaluations is asserted here. No full \nolinkurl{LawsSh}, \nolinkurl{LawsFull} or \nolinkurl{ShortValTot} record is used as a premise in deriving the second conjunct.

Actual \nolinkurl{CovPin} supplies \nolinkurl{IsLev}, the new-A equality and the operative pin. The actual unbundled \nolinkurl{LawsSh.dPhi} clause supplies PhiPin/PsiTie only in its stated range \nolinkurl{D.R.r≤n}; it is not a supplier for every smaller n.~Thus this result removes the section-descent obligation exactly where those particular fields are supplied. It does not construct dPhi, handle an arbitrary absent character package, or manufacture the full source-law record.

\subsection{2. The theta section module commutes with every base change}\label{the-theta-section-module-commutes-with-every-base-change}

Write \(\pi:A\to X\) for either actual universal abelian scheme, after the given coefficient-field base change, and L for its original line. Properness and the semiabelian condition make A an abelian scheme. The literal \nolinkurl{RelAmple} gives a relatively ample line; \nolinkurl{KTriv} says that its translation kernel is trivial. On a geometric fibre this is a principal polarization. The line is symmetric and rigidified with the original cubical structure.

For completeness, the needed cohomological form follows from a short argument. On any characteristic-zero geometric fibre, the cubical multiplication identity gives \([q]^*L\simeq L^{q^2}\). Multiplication by a positive integer q is finite flat of degree \(q^{2g}\). The finite-flat trace of1 is that degree, so division by it splits \(\mathcal O_A\to[q]_*\mathcal O_A\). Tensor by L and use the finite-morphism projection formula. Thus \(H^i(A,L)\) is a direct summand of \(H^i(A,[q]^*L)\). For q sufficiently large the latter vanishes when i\textgreater0 by Serre vanishing. Consequently all higher cohomology of L vanishes. The polarization degree formula \(\deg\phi_L=\chi(L)^2\), together with principality, gives

\[
 h^0(A_s,L_s)=1,\qquad H^i(A_s,L_s)=0\quad(i>0). \tag{FD3}
\]

The multiplication identity is already established with the source rigidification in LIFT100/LIFT103. The finite-flat degree and the polarization formula are the cited standard abelian-variety results; no new version of them is asserted. Serre vanishing is Stacks Lemma30.17.1(1)⇒(2). Milne, Theorem13.3(a), supplies the degree formula.

The morphism π is proper, flat and of finite presentation, and L is invertible. By Stacks Lemma36.30.4, \(R\pi_*L\) is perfect and its formation commutes with arbitrary base change. Lemma36.31.3, applied to FD3 at every point, makes it an invertible module in degree zero. Hence

\[
 \mathcal E=\pi_*L\text{ is invertible},\qquad
 \mathcal E\otimes\mathcal O_T\simeq\pi_{T,*}L_T
 \quad\text{for every }T\to X. \tag{FD4}
\]

This proves the nonflat base-change assertion needed below; it is not inferred from the merely flat existence-cover argument in LIFT103. The proof even retains nonreduced T. If the dimension is zero, the same statements reduce to the unit section module.

\subsection{3. The actual physical frame functor and reduced tests}\label{the-actual-physical-frame-functor-and-reduced-tests}

The finite projective Lie module of A splits over the actual coefficient field F into its O-eigenspaces, since \nolinkurl{Coeff.conj} supplies every embedding of E. The geometric \nolinkurl{DetCond} identifies their constant ranks with \nolinkurl{Dw.mult}. Let \(\mathcal L_\tau\) denote these vector bundles and use FD4 for the theta-section line. The scheme

\[
 \mathcal F_X=
 \left(\prod_\tau\operatorname{Isom}(\mathcal O_X^{Dw.mult\,\tau},
                                  \mathcal L_\tau)\right)
 \times_X\operatorname{Isom}(\mathcal O_X,\mathcal E) \tag{FD5}
\]

represents exactly \nolinkurl{FrPt} over arbitrary F-algebras. To verify representability rather than assume it, trivialize each finite locally free module on an open of X. The first factor there is the matrix affine space with determinant inverted; the last factor is an invertible scalar coordinate. Their usual change-of-basis transition maps glue, and their S-points are precisely bases and a generator. The actual dual-number tangent definitions give the same Lie vectors and scalar combinations. FD4 identifies the global section generator in \nolinkurl{FrPt.gen} with a basis of the last line. Thus \nolinkurl{FrMap} is exactly base change of these represented data. In particular any FrPt pulls back under any coefficient-algebra map.

The moduli X is smooth and of finite type over F. Therefore \(\mathcal F_X\) is smooth, reduced, separated and of finite type over F. It has an affine-open cover by rings that are reduced finitely generated F-algebras. Each is isomorphic to a literal \nolinkurl{TRing F}, because it is a quotient of a finite-variable polynomial ring by a radical ideal.

A natural \nolinkurl{Rules} object is consequently a regular function on \(\mathcal F_X\). Indeed, evaluate it at the universal framed point on each such affine open. On affine opens of pairwise overlaps, naturality identifies the restrictions. The resulting regular functions glue. Conversely a regular function pulls back to every literal test point and is natural. Independence of the chosen TRing presentations follows from naturality under the corresponding algebra isomorphisms. These two operations are inverse by restricting an arbitrary test point to inverse images of the affine opens. No assertion about a universal moduli object over nonreduced bases is used; only the already supplied universal abelian scheme and its vector bundles are pulled back.

The source \nolinkurl{FrBor} is the lower-triangular change of each Lie basis and scalar change of the theta generator. Accordingly \nolinkurl{sec(a,n)} consists precisely of these regular functions with the displayed positive character of LIFT104. Every field-valued frame, in particular every algebraically closed one, factors through an affine open of \(\mathcal F_X\). Hence it is \nolinkurl{FromTest}, and the literal evaluator of any natural rule equals the value of its regular function at that point:

\[
 \operatorname{evR}(G,y)=G_{\rm reg}(y)
 \quad\text{for every field-valued FrPt }y. \tag{FD6}
\]

This supplies the original-base evaluator as well as the cover evaluator. It leaves untouched the source's zero branch at arbitrary non-FromTest points over general rings.

\subsection{4. The old half level has no automorphisms}\label{the-old-half-level-has-no-automorphisms}

First, \nolinkurl{Level.Dvd}, the new A=1 and old positivity give

\[
                 l.A=d.l_2.A=1. \tag{FD7}
\]

The two base fields \nolinkurl{KA} are therefore literally the same subfield \(\mathbf Q(H.rt0)\) of \(\mathbf C\). The actual level identities and cancellation of nonzero ideals in O give

\[
 M_2=Mm,\quad(m,M)=1,\quad b_2=bI,\quad
 I\bar I=(m),\quad I+\bar I=O. \tag{FD8}
\]

Here m is positive and odd, since both levels have positive odd M. The ideal cancellation uses \(b\bar b=(M)\), \(b_2\bar b_2=(M_2)\), and b₂=bI; it does not require b or I to be principal.

On a geometric fibre, the original eta gives an O-linear isomorphism \(L_0/bL_0\to A[b]\). This is an immediate CRT consequence of its specified values on \(\bar bL_0\), its b-kernel and the actual \nolinkurl{TateIso}. The apparent absence of an explicit O-linearity field in the eta record is harmless: these equations imply O-linearity at geometric points, and therefore on the reduced universal base as an equality of sections of a finite étale torsion scheme. Its pullbacks then retain that equality.

An O-linear polarized automorphism fixing eta fixes A{[}b{]}. It preserves both ideal summands \(A[M]=A[b]\oplus A[\bar b]\). The original Weil pairing is perfect and balanced, and both halves are isotropic. Its cross pairing is therefore perfect. If v lies in A{[}bar b{]} and α fixes A{[}b{]}, then \(e(u,\alpha v-v)=1\) for every u in A{[}b{]}; the difference lies in both its annihilator A{[}b{]} and A{[}bar b{]}, so it is zero. Thus α fixes all A{[}M{]}. The group of automorphisms preserving the polarization is finite by Milne Proposition17.5(a), so α has finite order. Since M≥3, Proposition17.5(b) then gives α=1.

Over a reduced base, two such group maps equal at every geometric fibre are equal: the source abelian scheme is smooth over a reduced base and hence reduced, and their separated equalizer has the whole underlying space. Two line isomorphisms over the same group map differ by a base unit, by universal \(\pi_*\mathcal O_A=\mathcal O_X\); rigidification makes that unit1. Consequently there is at most one normalized old-level object isomorphism. This uniqueness is the only rigidity needed for the following use of \nolinkurl{Moduli.rep}.

\subsection{5. Construct the canonical half-level cover independently of d.r}\label{construct-the-canonical-half-level-cover-independently-of-d.r}

Over the old moduli scheme, form the sheaf

\[
 Y=\operatorname{Isom}_{O/I}(L_0/IL_0,A[I]). \tag{FD9}
\]

The indicated source is the constant finite O/I-module scheme. Torsion in characteristic zero is finite étale, and the ideal-kernel is a union of its finite étale components. Thus both sides are finite locally constant modules. Their isomorphism sheaf is finite étale: on an étale cover trivializing the modules it is a finite constant set, and the transition permutations glue it. The geometric \nolinkurl{TateIso} makes it nonempty on every fibre. It is a right torsor under the constant finite group

\[
                     G=\operatorname{Aut}_{O/I}(L_0/IL_0). \tag{FD10}
\]

No free global O-basis is chosen. Projectivity of L₀ is sufficient; its reduction at I is free locally on the finite product of Artinian local rings O/I.

A point of Y specifies precisely the extra half marking. The CRT decomposition for b and I, which are coprime by (m,M)=1, gives a marking eta₂ with old component \([m]^{-1}\eta\) in A{[}b{]} and its specified new component in A{[}I{]}. Conversely eta₂ maps to \([m]\eta_2=\eta\), and its I-component gives that point of Y. All sections, O-actions, cubical line and determinant condition are the original ones.

We verify the geometric Tate condition for every such new marking. Start with any of the old fibre's full geometric Tate isomorphisms. At each prime of O dividing I, an automorphism of L₀/IL₀ lifts to the corresponding completed free lattice: lift its matrix and note its determinant is a unit. On the conjugate prime choose the inverse adjoint for the perfect alternating pairing, and choose the identity at all other primes. Primes fixed by conjugation cannot divide I, because I and bar I are coprime. The resulting full profinite O-linear automorphism preserves EF and is identity at all primes dividing M. Composing the initial Tate isomorphism with it realizes the specified I-half marking, while the conjugate half is forced by the pairing. The harmless units M and m in the CRT identifications are incorporated in the chosen I-matrix; they do not restrict which isomorphisms occur. The old half marking is unchanged. This proves the actual new Obj.geom clause, including compatible primitive roots; the multiplier is1.

The old object and Y are reduced, and Y is finite étale over the smooth old moduli. Therefore the given \nolinkurl{Moduli₂.rep} applies to the new object on Y. Conversely the universal high-level object, with eta multiplied by m and trivial A-marking, defines an old-level object on the reduced high-level moduli and hence an old classifying map. Its I-component defines a map to Y. Uniqueness in both representability clauses and the old-level rigidity of section4 make these two maps inverse, with their normalized family isomorphisms. We obtain a canonical finite étale surjection

\[
              \pi:X_{o,2}\longrightarrow X_o,
 \qquad X_{o,2}\simeq Y, \tag{FD11}
\]

over their common KA, together with the canonical cartesian O-linear cubical family map. Base change by the actual equal coefficient fields gives the same torsor over F. This construction does not assert \(d.r=\pi\) as a global map and does not use the declared \nolinkurl{ToroidalLevelCompatibility} proposition as a supplied theorem about d.r.

\subsection{6. The deck transformations satisfy the literal Sim and LevAut clauses}\label{the-deck-transformations-satisfy-the-literal-sim-and-levaut-clauses}

For k∈G, the completed isometry constructed in section5 defines compatible isometries of every \(L_0/nL_0\), identity at all primes dividing M. Finite projectivity gives

\[
 \operatorname{End}_O(L_0)/n\operatorname{End}_O(L_0)
 \simeq\operatorname{End}_{O/n}(L_0/nL_0). \tag{FD12}
\]

Choose an integral representative f(n) of each reduction and put ν(n)=1; choose any representative at n=0. Compatibility of the fixed completed isometry gives the exact \nolinkurl{Sim.tower}. Its inverses modulo n lift by the same argument, giving \nolinkurl{Sim.inv}. Pairing preservation modulo n gives \nolinkurl{Sim.mult}. The resulting literal Sim is identity modulo b at the old level, and the two A-conditions are automatic by FD7. It belongs to \nolinkurl{InGam l l₂} and acts by k on LaI.

Modify the high-level eta by a↦eta(g.f(M₂)a), keeping the underlying polarized object. The reduced \nolinkurl{Moduli₂.rep} constructs an automorphism ρ of X₂ with the exact orientation

\[
 \nu_\rho\bigl(\eta_2(g a)\bigr)=\eta_2'(a), \tag{FD13}
\]

where νρ is its cartesian O-linear cubical family map. The inverse completed isometry gives the inverse map. This is exactly the source \nolinkurl{LevAut}, including its base-KA condition. It fixes the old marking, so it is the corresponding deck action of the torsor FD11. Every pair of geometric points in one π-fibre is related by such a ρ. Their physical frames and theta generator transport uniquely through the canonical old-object identification. Thus FD11 also induces a finite étale G-torsor

\[
                  \mathcal F_{X_2}\longrightarrow\mathcal F_X. \tag{FD14}
\]

The fibres here are full framed fibres, not merely unframed geometric points.

\subsection{7. The original finite Fourier character has the required sign}\label{the-original-finite-fourier-character-has-the-required-sign}

Let c=\nolinkurl{cHalf H I}. The Chinese remainder theorem gives \(c\in\bar I\), \(c-1\in I\). The pairing

\[
 B_I(\bar u,\bar v)=EF(u,cv)\pmod m,
 \quad L_0/\bar I L_0\times L_0/IL_0\longrightarrow\mathbf Z/m \tag{FD15}
\]

is well defined and perfect. For well-definedness, changing v by IL₀ makes cv change by mL₀; changing u by bar I L₀ moves that factor to the second variable by the original adjoint identity and has the same effect. Modulo m, c is the projector onto the I-quotient summand. Unimodularity of EF and the complementary isotropic CRT summands make the resulting cross pairing perfect. This proof also covers nonprincipal ideals.

For k∈G there is therefore a unique adjoint kst on the conjugate quotient such that

\[
                  B_I(kst\,u,v)=B_I(u,kv). \tag{FD16}
\]

It is O-linear with the appropriate conjugate adjunction. For the completed unitary g above it is the inverse of g's action on the bar-I quotient. Equation FD16 is exactly the congruence hypothesis in PhiPin, including its arbitrary qRep representatives. PhiPin gives \(ph(kst x)=\psi(k)^{-1}ph(x)\).

The Fourier formula in PhiPin now yields, by the finite change of variable y=kst x,

\[
 \phi_\delta(gq)
 =\sum_x ph(x)\zeta_m^{\sum_b B_I(x_b,kq_b)}
 =\psi(k)\sum_y ph(y)\zeta_m^{\sum_b B_I(y_b,q_b)}
 =\psi(k)\phi_\delta(q). \tag{FD17}
\]

All identities are in the original coefficient field via eF. The actual root \(\zeta_{M_2'}^{M_2'/m}\) has order m, since A₂=1 and M₂=Mm. Representatives change the exponents by multiples of m, so there is no chosen-representative ambiguity.

The old factor \(\zeta_M^{\operatorname{pairL}(x_0,q)}\) is unchanged by g. Indeed gq−q lies in bL₀ modulo M₂, while \nolinkurl{cSep l} lies in bar b; their product lies in ML₀ and EF is integral. The complete original weight therefore has transformation

\[
                       w(gq)=\psi(k)w(q). \tag{FD18}
\]

No assumption on the value or support of x₀ is needed for this invariance. The other PhiPin conditions remain available for their original source uses but are not substituted for the Fourier identity.

\subsection{8. Transport the actual trace and harmonic functional}\label{transport-the-actual-trace-and-harmonic-functional}

We first spell out the geometric comparison needed here, since \nolinkurl{IsNab} alone does not assert equivariance of its arbitrary chosen ordinary coefficients at different moduli points. A cartesian O-linear cubical family isomorphism preserves the original normalized multiplication generator, normalized odd theta lifts, dual-row kernel, all root torsors and the source action order. Each follows from its defining pullback identity and the uniqueness already proved in LIFT101/LIFT103. The original normalized primitive generator is preserved as well: the pulled primitive and the intrinsic primitive differ by a base unit, and their rigidifications force that unit to1.

Transport an actual finite-order chart and biframe through the family isomorphism, together with its physical tangent frame and section generator. Its formal group law, action, corrected germs and ordinary directional contraction are unchanged after that transport. The global primitive and adjunction maps match by the preceding normalization; the cubical map gives the other frame equations. Compare this admissible chart with the one chosen by the stored IsNab at the target. LIFT101 AT11--17/LIFT102 TheoremD identifies their complete equal-line-power harmonic contractions. This proves functoriality of the actual full scalar trace under such family isomorphisms, without demanding equality of individual ordinary coefficients.

For y₂ and its deck translate y₂′ with the transported physical data, FD13 identifies the q-labelled section trace at y₂′ with the (gq)-labelled section trace at y₂. The canonical kernel is the same under that identification, so no new permutation multiplicity occurs. If T\_y(q) denotes its full original harmonic contraction, the unique scalar from LIFT103 satisfies

\[
 V(y_2')=C\sum_q w(q)T_{y_2}(gq)
       =\psi(k)^{-1}C\sum_q w(q)T_{y_2}(q)
       =\psi(k)^{-1}V(y_2). \tag{FD19}
\]

The change of variable is a bijection of the original LevQ tuple space. The constant C is exactly \(\beta^{-1}mIdx^{-1}|LevQ^n|^{-1}N^{-\deg P}\). It is identical on the two points. If β=0 the entire equality remains valid with value0. No phase is dropped in transporting U\_t U\_z, and no reordering of those operators is required.

The literal PsiTie applies to this k,g,ρ: all its \nolinkurl{InGam}, new-A identity, \nolinkurl{LevAut} and LaI conditions were checked in section6. It gives

\[
 c_\delta(y_2')=\psi(k)c_\delta(y_2). \tag{FD20}
\]

Together with FD19, this proves exact invariance of the LIFT104 cover rule:

\[
               H_2(y_2')=H_2(y_2),\qquad H_2=c_\delta G_2. \tag{FD21}
\]

These are actual field evaluations, not arbitrary selections of a geometric scalar. LIFT104 or FD6 identifies them with the regular function attached to the literal cover rule.

\subsection{9. Effective descent and the exact source evaluator}\label{effective-descent-and-the-exact-source-evaluator}

View H₂ as a regular function on \(\mathcal F_{X_2}\), using section3. For each g∈G, the difference g*H₂−H₂ vanishes at all algebraically closed points by FD21. The frame scheme is reduced and of finite type over F; such a regular function is zero. For example on each affine chart its value lies in every maximal ideal after extension to an algebraic closure, so it lies in the nilradical, which is zero. Thus invariance holds as an identity of regular functions, rather than only of point values.

For a finite étale G-torsor, its two pullbacks on the fibre product agree exactly when all deck pullbacks agree. Affine-locally this is the ordinary faithfully flat equalizer for elements. Hence H₂ descends to a unique regular function F on the original frame scheme. The right Borel/frame action commutes with deck transport. Since H₂ has the original positive weight (a,n), faithful-flat injectivity gives the same character for F. Section3 turns F into a literal

\[
                        f\in\operatorname{sec}(X_o,c,a,n). \tag{FD22}
\]

Now take any original field, coefficient-compatibility data and actual \nolinkurl{FrOver d y₂ y} from \nolinkurl{IsShortRule}. We must compare it with the canonical π, not assume the stored r has that global description. The supplied d.m is cartesian and O-linear, while \nolinkurl{EtaOK} gives exactly the old marking \([m]\eta_2\). Coefficient compatibility makes both old objects live over the same field and KA-map. Therefore the object at y and the canonical old object at π(y₂) are isomorphic as old-level objects. \nolinkurl{Moduli.rep} over this reduced field forces their classifying points to agree. The old-level rigidity of section4 makes the isomorphism the canonical one. The two tangent equalities in FrOver identify the physical frames. Its line-generator witness is normalized at zero, so proper universal functions force it to equal the canonical line isomorphism; its SecMaps and pullback clauses then identify the theta generators as well. Consequently the image of y₂ in the frame torsor FD14 is exactly y.

By descent and FD6,

\[
 \operatorname{evR}(f,y)=H_2(y_2)
 =\operatorname{evR}(d.c_\delta,y_2)V(y_2). \tag{FD23}
\]

The scalar V is the literal unique ShortValAt from LIFT103. Equation FD23 is all of the second original \nolinkurl{IsShortRule} conjunct, for every pair quantified there. It neither replaces FrOver by a smaller chosen collection nor asserts a global identity d.r=π. LIFT104 CV20 now supplies the required second ShortValTot conjunct; together with LIFT103 it proves FD2 in the stated hypothesis scope.

\subsection{10. Boundary, primary locators and pending verification}\label{boundary-primary-locators-and-pending-verification}

If m=1, then I=O and the torsor and Fourier group are trivial; the same formulas apply. Zero polynomial, zero beta, n=0 and rank-zero modules use their usual singleton/empty-index conventions. The comparison uses characteristic-zero generic fibres and the actual coefficient field. It claims no integral extension or alteration of the supplied level, polynomial, lattice, theta source, arithmetic image, periods or DP44 data.

The argument is intended to discharge a named section-law obligation from d.IsLev, A₂=1, the operative pin and the actual PhiPin/PsiTie fields. It does not supply those records. In particular the source dPhi's n≥r condition is retained; no all-n construction is inferred from that restricted supplier. ShortInCusp, spanning, selected arithmetic nonzero/content, the global period, the specific032 pairing and evaluatedDP44/RH remain separate.

Primary citations used at their stated scopes:

\begin{itemize}
\tightlist
\item
  Milne, J. S. (1986/2022). \emph{Abelian varieties}. Author-corrected version of the chapter in \emph{Arithmetic geometry}, Theorem13.3(a) and Proposition17.5(a,b), with Lemma17.6; author PDF pp25 and37. \url{https://www.jmilne.org/math/xnotes/AVs.pdf} . The theorem statement and complete displayed proposition/lemma proofs were freshly read; no fresh whole-chapter check is claimed.
\item
  The Stacks Project Authors. (2026). \emph{The Stacks Project}. Lemma30.17.1, tag01XO (Serre vanishing criterion); Lemma36.30.4, tag0B91 (proper perfect cohomology and arbitrary base change); Lemma36.31.3, tag0BDK (constant one-degree perfect complex); Lemmas36.32.5--6, tag0BDM (universal global functions); Lemma59.64.4, tag03RV and Lemma41.18.3, tag04HN (finite étale local systems); Lemma35.3.6, tag023M and Lemma59.16.3, tag03O9 (faithfully flat descent). The new cohomological statements were directly read; finite-étale and descent inputs retain the full readings recorded in LIFT103.
\end{itemize}

The original interface scopes are \nolinkurl{TwoWeightEqOneAddRoot} Host/Datum/Level/TateIso/Obj/ObjMap, \nolinkurl{PushInvertible} Moduli/Coeff, \nolinkurl{NonzeroLocus} RelAmple/KTriv, \nolinkurl{HeckeKer} Sim/InGam/LevAut, \nolinkurl{ThShort} DeltaDat/EtaOK/FrOver/IsLev/PhiPin/PsiTie/ShortValTot, and \nolinkurl{AlgebraMapCAt}/\nolinkurl{IsCombAt} Rules/FrPt/sec/evaluator. They are accepted interfaces, not objects of an archive validity audit. There is no Lean build.

The canonical geometric dictionary, if this proof is accepted, remains the original finite étale level sheaf, rigidified cubical/Poincare line, Picard section line and its faithfully flat descent. The arithmetic source and all integral, Dieudonné, K-theory, weight and CM-motive L-function gates remain unchanged. Exact inherited FGA/SGA/cote locators will be attached in the reviewed integration record rather than invented here.

\subsection{Canonical algebraic-geometric context}\label{canonical-algebraic-geometric-context-87}

This is effective descent of a Borel-equivariant regular function on the genuine physical frame scheme over the original characteristic-zero moduli base. The Picard theta section line commutes with arbitrary base change, while its canonical finite-etale half-level torsor is built from the original projective OE-module and rigidified cubical/Poincare line. The Tate and inverse-adjoint data preserve the original polarization. Exact primary locators: Milne Abelian varieties Theorem13.3(a), Proposition17.5(a,b) and Lemma17.6, printed25/37; Stacks30.17.1 (01XO),36.30.4 (0B91),36.31.3 (0BDK),36.32.5-6 (0BDM),59.64.4 (03RV),41.18.3 (04HN),35.3.6 (023M),59.16.3 (03O9). Exact Alpoge interfaces and root fresh/carried extents are the eighteen-file canonical\_half\_level\_short\_rule\_descent\_sources.json; in particular ThShort DeltaDat.hF/EtaOK/FrOver/IsLev/PhiPin/PsiTie/ShortValTot, reduced-base PushInvertible.Moduli.rep, HeckeKer.Sim/InGam/LevAut, and actual CubeMap.IsPullCubical. Carried exact context: SGA7 I VII2.7/2.9 pp169--174 and VIII7.4(b) pp310--311; FGA232 Theorem3.1 pp148--149 and FGA236 Corollaries1.7/2.3 pp229/232; SGA3 I4.7.3/4.7.3.1; SGA6 VI Theorem3.3; cote4 scan003 sections8.1.1--8.1.2 and scans005--006 sections8.4.1--8.4.3; cote54 scans005--007. These are carried contextual readings, not new whole-text checks. No new integral extension over Z, Dieudonne slope, lambda-ring positivity, K-theory filtration, l-adic weight theorem or CM-motive L-function comparison follows. DP\_MULT Q3e/DP44/DP52/DP54--65 and THETA3 G1/A7 gain this conditional source-rule reduction; CX\_GRH CX-ALG-024 and GCT\_PLETH T-052 retain their global gates.

\clearpage

\section{AM-DO-127 --- exact Green--central correlation and five complete denominator tails}\label{am-do-127-exact-greencentral-correlation-and-five-complete-denominator-tails}

\textbf{Canonical checked integration --- 2026-11-23T13:08:27.305123+00:00.} Complete root and independent proof/script/interface reviews and separate exact replays passed. The actual canonical predecessor has136/31 entries; this integration gives137/31, complete floors27/11. Frozen author history is preserved. AM-DO-127 GC1-16 removes five complete containing subsets of the original support27/eligibility27 full-positive scalar-cross globally rank-two-unipotent all-II branch, with b=j\_int0: k160/161/162/163/164 and log2x\textgreater=2120/2130/2140/2150/2165 respectively, inside the original x\textless2\^{}3328 window. The exact normalized finite Green formula gives2 Q\^{}(-335*167)(NEJ)\^{}167 NFJ0, whose NFJ0 cancels the independent physical central divisor. The complete pattern is product(1+q\_v\^{}(j-167))\^{}2, including exactly2\^{}668 at degree167. The actual projective selfdual quotient, full1008 transverse theta, physical U\^{}-f, same positive profile, all1630 terms and one entire band tail remain. Worst directed term -13968857/32768\textless-426 gives the full union\textless2\^{}-415. Exactly one415 append gives137/31 full sums\textless2\^{}-230 and unchanged floors27/11. The sole support27 necessary branch has2\^{}2111\textless x\textless2\^{}(X\_k) with the five respective upper exponents. Low-x cells, signed PP3, higher supports, actual source/selected arithmetic image, period, specific032 value, evaluatedDP44 and Hecke RH remain open.

\textbf{Controlled pending theorem.} Author whole verification, independent whole review and canonical integration are separate events. The earlier CP52 snapshot and its floating screens are preserved unchanged under \nolinkurl{historical/do127_cp52_pending} and in their original diagnostic files. This final mathematical text uses the complete exact certificate, not those screens. No canonical artifact is edited here.

\subsection{1. Exact current scope and inherited factors}\label{exact-current-scope-and-inherited-factors}

The actual predecessor is \nolinkurl{checks/do126_joint_remainder_error.json}, SHA256 \nolinkurl{c80115b13d58a6e8a1e309d995c67e72fe641a19368c0224dd3e6ef658a84b85}, with all136 all-top and31 uniform charges, complete floors27/11. Its sole support27 necessary branch, outside the already charged subsets, is globally rank-two-unipotent, full-positive scalar-cross, all type II, total eligibility27, b=j\_int=0, with160\textless=k\textless=164 and2\^{}2111\textless x=N\_E D\_eta\textless2\^{}3328. Here k counts nonzero original S covector coordinates; b counts inactive paired-C-zero places; j\_int counts active places where the whole paired C is integral.

\textbf{Theorem.} The union of the following five complete subsets of that original branch has total absolute actual normalized contribution strictly below2\^{}-415:

\begin{longtable}[]{@{}rr@{}}
\toprule\noalign{}
k & newly covered denominator range \\
\midrule\noalign{}
\endhead
\bottomrule\noalign{}
\endlastfoot
160 & 2\^{}2120\textless=x\textless2\^{}3328 \\
161 & 2\^{}2130\textless=x\textless2\^{}3328 \\
162 & 2\^{}2140\textless=x\textless2\^{}3328 \\
163 & 2\^{}2150\textless=x\textless2\^{}3328 \\
164 & 2\^{}2165\textless=x\textless2\^{}3328 \\
\end{longtable}

Thus outside one freshly charged union, the upper denominator bounds become respectively2\^{}2120,2\^{}2130,2\^{}2140,2\^{}2150,2\^{}2165. The lower bound2\^{}2111 and every other literal premise remain. This is a necessary-window reduction, not emptiness or a complete support28 result.

Retain the actual J56 field, standard compact, selected vector-valued representation, allowed words and twists, both original CM partners, finite projectors, source and target pencils, omitted labels, packet representative, central phase, Haar normalization and class mass. The entire accepted DO126 IQ1--36 quotient theorem is used, including its projective selfdual quotient at every prime, genuine1008-dimensional transverse theta factor K3, independent physical-central U\^{}-f, full750 central cover and exact eta/row/central reconstruction. In particular no global free basis, division by2 or unproved unit property at primes dividing U is introduced.

Set f=168, H=133, R=17424/133, Q=1427908193729121522601, U=77Q\^{}2−14, D0=7\^{}140\emph{29\^{}144}5\^{}84*13\^{}84 and D\_F=2\^{}168 D0. T and g0\textless g1 are the exact J56 host integers, Q=g0g1. The accepted scalar theta coefficient is K from DO126 IQ23; the full transverse coefficient is K3 from IQ21. Both are restated in section4. Only the finite character/central-ideal correlation is sharpened; all other analytic and reconstruction constants are retained.

\subsection{2. The exact finite identity}\label{the-exact-finite-identity}

At every one of the334 full cuspidal G places, h is integral and N=h−I satisfies N\^{}3=0 and rank N\textless=2 after reduction. Thus the reduction is unipotent of rank d=0,1,2. Green's formula, as explicitly stated in Gorodetsky and Hazan (2017, Theorem2.1, equation2.1, and Corollary2.2, printed p.5), and the original specialization AM-LIFT-064 GP1 give normalized absolute character

\[
 |a_v|=\left[\prod_{i=169-d}^{168}(q_v^i-1)\right]^{-1}.
 \tag{GC1}
\]

The same original common-row ideal J and central rank-zero ideal J0 have the DO098 CF9--10 exponents

\begin{longtable}[]{@{}
  >{\raggedleft\arraybackslash}p{(\columnwidth - 6\tabcolsep) * \real{0.2500}}
  >{\raggedleft\arraybackslash}p{(\columnwidth - 6\tabcolsep) * \real{0.2500}}
  >{\raggedleft\arraybackslash}p{(\columnwidth - 6\tabcolsep) * \real{0.2500}}
  >{\raggedleft\arraybackslash}p{(\columnwidth - 6\tabcolsep) * \real{0.2500}}@{}}
\toprule\noalign{}
\begin{minipage}[b]{\linewidth}\raggedleft
finite rank d
\end{minipage} & \begin{minipage}[b]{\linewidth}\raggedleft
exponent of N\_EJ
\end{minipage} & \begin{minipage}[b]{\linewidth}\raggedleft
exponent of N\_FJ0
\end{minipage} & \begin{minipage}[b]{\linewidth}\raggedleft
leading character exponent
\end{minipage} \\
\midrule\noalign{}
\endhead
\bottomrule\noalign{}
\endlastfoot
0 & 2 & 1 & 0 \\
1 & 1 & 0 & 168 \\
2 & 0 & 0 & 335 \\
\end{longtable}

Both rank-one orientations are retained as distinct local states. In each row of the table the leading finite character weight is exactly

\[
 q_v^{-335}(q_v^{v(N_EJ)})^{167}q_v^{v(N_FJ_0)}.
 \tag{GC2}
\]

For completeness, when d=1 the denominator is q\_v\^{}168−1, and when d=2 it is (q\_v\^{}167−1)(q\_v\^{}168−1). The value depends on rank, including both possible rank-two unipotent Jordan types. The parameter character is evaluated at1. The product of all334 residue norms is Q\^{}167. There are at most668 factors of the form (1−q\_v\^{}-i)\^{}-1, all with i\textgreater=167. Bernoulli's inequality and the exact host comparison g0\^{}167\textgreater1336 give

\[
 \prod_{v,i}(1-q_v^{-i})^{-1}
 \le(1-g_0^{-167})^{-668}<2.
\]

The two actual mixed factors retain their modulus-one bounds. Consequently

\[
 |a_G(h)|\le2Q^{-335\cdot167}(N_EJ)^{167}N_FJ_0.
 \tag{GC3}
\]

This is an absolute containing bound. It does not assert positivity of the finite product or impose an unproved Jordan-type condition beyond actual global rank two and unipotence.

The physical central affine lattice still has covolume

\[
 U^f\sqrt{D_F}(N_FJ_0)xT^{-54}.
 \tag{GC4}
\]

Thus its factor (N\_FJ0)\^{}-1 cancels the final factor in GC3 exactly, before summing patterns. The original central lattice, its shift, its covering radius, and the independent U\^{}-f do not change.

\subsection{3. The complete pattern and coupled variables}\label{the-complete-pattern-and-coupled-variables}

Put q=R/(Uz\^{}2), z\^{}(2f)=(N\_EJ)xT\^{}-54, and

\[
 q_*^f=R^fT^{54}/(U^fQ^{334}x),\qquad
 \rho=(N_EJ)/Q^{334}=(q_*/q)^f.
 \tag{GC5}
\]

For a positive row-polynomial degree j, the Green--central--row product relative to the all-identity G cap has factor rho\^{}(167-j). The repaired zero-degree term has effective j=1, exactly as in IQ27. All genuine degrees now satisfy1\textless=j\textless=167. Lowering a local J exponent from2 to1 gives q\_v\^{}(j-167), and lowering it to0 gives its square. Both rank-one orientations are paid once. The complete pattern product is therefore

\[
 \mathcal B_j=\prod_{v=1}^{334}(1+q_v^{j-167})^2.
 \tag{GC6}
\]

For1\textless=j\textless=166, log B\_j\textless=668/g0\textless1/2048\textless(log2)/1024 by exact rational comparisons. Thus B\_j\textless2\^{}(1/1024). For j167, the product is exactly4\^{}334=2\^{}668. This includes all four states at every relative place; no Q\^{}167 exception remains. At fixed x, q\^{}((j-167)f) is nonincreasing, including equality for j167. At q=q*, the coupled x power is exactly x\^{}(-1-j), hence every term decreases strictly throughout its claimed high-denominator interval.

The fixed denominator eta is unique for each27-place support throughout the whole window, by K113 SC27.4. Taking an endpoint uniformly over that unique value does not add a denominator-ideal count, a discretization in x, or another choice of eta. The positive quotient lattice may vary with eta, but the accepted theta estimate is uniform over the actual selfdual projective lattice.

The complete prefactor before the row polynomial can equivalently be written

\[
\begin{split}
\mathcal A_k(x,q)={}&2\,2^{3/1024}\binom{168}{27}164\,
 \mathcal D_*^{84}\vartheta^{168}(\kappa_I\mathcal G_1D_*)^{168}
 D_F^{-1/2}U^{-f}\binom{164}{k}\pi_+^{(k+3)f}D_0^{-(k+3)}\\
 &\times(R/U)^{167f}T^{8904}Q^{-335\cdot167}x^{-168}q^{-167f}.
\end{split}
\tag{GC7}
\]

The three small inflations are the actual near-point transport, central Voronoi cover and complete168-place T correction. The finite pattern factor GC6 is attached separately to each term. Green2, all support and164 first-pivot choices, both original mixed finite factors, and the full physical normalization remain. The T exponent is334*27+54-168=8904 after substitution of (N\_EJ)\^{}167; this equals the accepted DO126 all-identity endpoint after GC5 is inserted.

\subsection{4. The unchanged actual positive profile and full quotient theta}\label{the-unchanged-actual-positive-profile-and-full-quotient-theta}

The original same-kernel coarea, paid tangent and physical central integral are exactly DO125 CM3--16 / DO126 IQ24--27. In particular this proof never substitutes a compact-support Mellin integral into a positive Laplace argument. With B0=2\^{}955, the constants in GC7 are

\[
\begin{gathered}
 \vartheta=\frac{4^{167}}{334\binom{334}{167}},\quad
 \kappa_I=\frac4{\sqrt{527}}\frac{23479}{23450},\quad
 \mathcal G_1=\frac{2\cdot4^{83}82!83!}{166!},\\
 V_* =\frac{97}{100B_0-97}\frac{2113^2}{64\cdot2049},\quad
 \mathcal D_*=334^2\frac{400}{399}V_*+528\cdot167^2\,2^{-2800},\\
 D_*=(1480/1653)^{333/2}e^{1557/80}\sqrt{1670/1757},\\
 \Phi(t)=t^{-1/2}(1+t/4)^{-166}(1+t/2)^{-1/2}
 (1+t/(3514/29))^{-1/2}e^{-29t/20}.
\end{gathered}
\tag{GC8}
\]

Its positive Laplace measure is the convolution of independent probability Gamma factors of shapes166,1/2,1/2 and rates4,2,3514/29, the infinite half-Gamma measure ds/sqrt(pi s), and shift29/20. It has infinite total mass; only its positive inverse moments are used. For j\textgreater1/2,

\[
 M_j=\frac1{\Gamma(j)}\int_0^\infty
 t^{j-3/2}(1+t/4)^{-166}(1+t/2)^{-1/2}
 (1+t/(3514/29))^{-1/2}e^{-29t/20}\,dt.
 \tag{GC9}
\]

Put theta=2\^{}-20. On the original genuine row the k nonzero S coordinates give product(alpha\_tau)\textgreater=(Rk/q)\^{}f.~Reserve theta of the actual Gaussian exponent before any tuple enlargement. The AM--GM argument of DO125 CM21 then controls the sole degree-zero term by

\[
 \int e^{-\theta\sum c_\tau\alpha_\tau}d\mu^{\otimes f}
 \le\overline M_1^f[q/(\theta Rk)]^f.
 \tag{GC10}
\]

There is no per-embedding alpha floor on the enlarged row set and no invocation of an infinite M0. The remaining Gaussian uses(1−theta)c\_tau.

The actual projective rank3 quotient from DO126 IQ8--22 has covolume D0\^{}3 z\^{}(6f), with the physical positive metric isometric to the original transverse space. Its dual lower bound is f/(r0\^{}2 z\^{}2 B), where B=(product c\_tau)\^{}(1/f); the numerator is f, not3f. In real dimension1008 the full theta factor is

\[
 K_3=\frac{2\cdot1008^2}{c_{1008}}2^{504}504!
 \left[\frac{(365729/500)^2}{168\pi_-^2}\right]^{504},
 \quad c_{1008}=504\binom{1008}{504}/4^{504}.
\]

For each standard S coordinate retain the accepted K111 scalar coefficient

\[
 \begin{gathered}
 K=\frac{336M_{15}}{c_{336}}(7841/5000)^{168}168!
 \left[\frac{(365729/500)^2}{168\pi_-^2}\right]^{168},\quad
 c_{336}=168\binom{336}{168}/4^{168},\\
 M_{15}=\frac{6015684926187167549292201834533348561577984}
 {2618017117694612446410293519},\quad
 \pi_-={314159}/{100000},\quad \pi_+={355}/{113}.
 \end{gathered}
 \tag{GC11}
\]

The executable verifies both directions of these rational pi bounds by Machin's identity and alternating arctangent intervals. No new scalar-only optimization is used.

Combining the k scalar theta factors with the genuine rank3 theta factor, integrating the same positive measure, and retaining the repair GC10 gives all2(k+1) terms

\[
\begin{split}
 \mathcal P_k(q)={}&\sum_{l=0}^k\binom kl K^{k-l}
 [q/((1-\theta)R)]^{(l+3)f}\overline M_{l+3}^{f}\\
 &+K_3\sum_{l=1}^k\binom kl K^{k-l}
 [q/((1-\theta)R)]^{lf}\overline M_l^{f}\\
 &+K_3K^k\overline M_1^{f}[q/(\theta Rk)]^f.
\end{split}
\tag{GC12}
\]

The first summand has effective degree j=l+3; the second has j=l; the last has effective j=1. Multiply each individual term by GC7 and its own GC6, before taking any union. The factor binom(164,k) pi\_+\^{}((k+3)f)/D0\^{}(k+3) is already in GC7. There is no compensating U\^{}f row factor, because the actual quotient is selfdual; the independent physical-central U\^{}-f stays present.

Fixing eta, the complete quotient covector, and the original central scalar reconstructs h injectively by the original normal form. The selected active labels are fixed by that same eta and actual nonintegral whole C; the actual b=j\_int=0 condition makes all168 T weights constant on each complete central fibre. Their combined weight is(T−1)\^{}-168. The first-pivot choice164 and support count binom(168,27) remain explicitly paid. Fixing a finite rank pattern before enlarging its central coset allows that coset to contain other residue patterns; this is a containing count with the original pattern's fixed weight. Summing all four states afterward is therefore valid without claiming independent geometric realizability of those states.

\subsection{5. All moment degrees, quadrature tails and directed arithmetic}\label{all-moment-degrees-quadrature-tails-and-directed-arithmetic}

The accepted DO126 exact anchor certificate supplies the actual rational upper constants for M148,M162,M163; its complete SHA256 is checked by the executable. The M1 upper constant is the accepted K120 j1 positive rational integral times355/113. Dropping the two extra half factors and replacing29/20 by32/25 proves this is a containing upper bound for GC9.

The new executable recomputes M164,M165,M166,M167. It copies the accepted integer/Fraction quadrature routine: mesh1/128 on{[}1,64{]}, a separate exact second-derivative trapezoid error on each of8064 intervals, both entire tails{[}0,1{]} and{[}64,infinity), 256-bit outward exponentials and square roots, and scaling2\^{}700. On the right tail the power exponent is j−337/2\textless0 for every new degree, so the same decreasing-power/exponential estimate applies. The left tail is bounded by1/{[}(j−1)!(j−1/2){]}. No part of either tail is discarded.

For every intermediate integer degree, Holder on this same positive measure uses its nearest finite certified endpoints a\textless b:

\[
 M_j\le\overline M_a^{(b-j)/(b-a)}
          \overline M_b^{(j-a)/(b-a)},
 \quad a<j<b,\qquad
 a,b\in\{1,148,162,163,164,165,166,167\}.
 \tag{GC13}
\]

No probability normalization of the infinite convolution and no zeroth moment are used. The exact finite endpoint constants themselves, not lower bounds for true moments, are evaluated for any failed sufficient envelope.

All logarithms use160-bit directed rational bounds: normalize the argument into{[}1,2{]}, expand log(v)=2 sum z\^{}(2m+1)/(2m+1), z=(v−1)/(v+1), retain100 terms and an explicit positive geometric tail. Negative multipliers reverse the selected endpoints. Thus no floating log enters the certificate.

The three nonpattern natural-log inflation bounds used in GC7 are respectively

\[
 \frac{f(335+3Hp+1+2\gamma H)}{2^{352}},\quad
 \frac{750f(p/2+1/4+\gamma)}{2^{83}},\quad
 \frac{168}{T-1},\qquad p=333/4,\quad\gamma=29/10.
 \tag{GC14}
\]

Each is strictly below1/2048. The physical central estimate U z\_c\textless2\^{}-83 is uniform over the whole original x\textless2\^{}3328 window by the exact host comparison U\^{}f Q\^{}167 2\^{}(3328+83f)\textless T\^{}54. The pattern inflation is separate: below1/1024 bit for j\textless=166 and exactly668 bits for j167. There is no extra generic polynomial factor2, because every positive term is explicitly enumerated.

\subsection{6. The complete endpoint and error union}\label{the-complete-endpoint-and-error-union}

For each k, write X\_k for the lower exponent in the theorem's table. GC5 and168th powers prove the following strict integer brackets, while the full exact term list gives the displayed rational upper maxima:

\begin{longtable}[]{@{}rrrr@{}}
\toprule\noalign{}
k & X\_k & q\_-\textless q\_*\textless q\_+ & largest directed log2 upper bound \\
\midrule\noalign{}
\endhead
\bottomrule\noalign{}
\endlastfoot
160 & 2120 & 241970978,241970979 & −13968857/32768 \\
161 & 2130 & 232190704,232190705 & −28656345/65536 \\
162 & 2140 & 222805740,222805741 & −15250679/32768 \\
163 & 2150 & 213800110,213800111 & −2092789/4096 \\
164 & 2165 & 200969543,200969544 & −6113299/8192 \\
\end{longtable}

Negative q powers use q\_-, positive q powers use q\_+. Every one of the1630 terms is strictly below2\^{}-426. In particular the k164 main-volume term of degree167 pays the entire2\^{}668 pattern factor. The global maximum is the k160,l160 main-transverse term of degree163:

\[
 E_{\max}\le-13968857/32768<-426.
 \tag{GC15}
\]

The entire original omitted narrow-band contribution is less than2\^{}-8605114 by accepted K120 NB14--16 / K118 BT15--22. This is already a full containing mass bound and is charged once for the union, not once per k. Therefore

\[
 \sum_{k=160}^{164}\!\text{claimed-sector contribution}
 <1630\,2^{-426}+2^{-8605114}
 <1631\,2^{-426}<2^{-415}.
 \tag{GC16}
\]

This proves the theorem. The executable checks the complete actual136/31 arrays, appends415 once, leaves the31 uniform exponents unchanged, and verifies all full error sums below2\^{}-230 by exact integer arithmetic. The complete support floors remain27/11.

The failed common cutoff2160 remains explicit: the chosen complete k164 main-term upper constant has directed log2 lower endpoint6143847/65536\textgreater0. At k160 and X2111, the analogous chosen upper constant has lower endpoint68793421/65536\textgreater0. These are failures of those sufficient envelopes, not lower bounds for the original sector or counterexamples to a stronger theorem.

\subsection{7. Preserved history, source scope and remaining main gates}\label{preserved-history-source-scope-and-remaining-main-gates}

The original CP52 draft, moment-only258100-check probe, exact-input/floating-output screens, and separate paid-tangent scan remain unchanged. The latter found only a floating sampled gain of about33.41 field bits near gamma3.02; it is not a rigorous optimum and it is not used in this theorem. The first complete268600-check replay is preserved separately from the final proof-hash replay. The accepted DO126 proof and all its original records remain immutable.

The principal newly refreshed input is Green's exact rank-dependent cuspidal character formula. The original J/J0 table and central coset are DO098 CF8--14; their support-free membership content is used with the current support27 transfer K113 SC27.1--10 and the full accepted DO126 IQ1--36. The physical positive-mixture/zero-degree/cover argument is DO125 CM3--27, retained in DO126 IQ24--34; every scope is restated above where the new finite correlation interacts with it. The new arithmetic diagnostic proves its exact inequalities but does not substitute for the projective quotient or original-arrow geometry.

\textbf{Primary reference.} Gorodetsky, O., \& Hazan, Z. (2017). \emph{On certain degenerate Whittaker models for cuspidal representations of GL\_(kn)(F\_q)} (arXiv:1707.07308v2). Theorem2.1, equation2.1 and Corollary2.2, printed p.5. \url{https://arxiv.org/pdf/1707.07308}. The full stated character theorem and its rank consequence were freshly opened and read for this proof; its original Green attribution is retained.

No complete support28 estimate, signed PP3 inequality, actual selected arithmetic-image statement, prescribed period, specific032 value, evaluated DP44 bound or RH conclusion follows. The original global/source/period and Grothendieck-context gates remain as in the canonical predecessor; this packet claims no new primary FGA/SGA reading or descent theorem.

\subsection{Canonical algebraic-geometric context}\label{canonical-algebraic-geometric-context-88}

The refined finite formula is the character of the original cuspidal finite representation at its unipotent reduction, correlated with the actual row ideal and central congruence ideal. It is applied to the same saturated isotropic quotient over Spec(O\_E), projective Hermitian dual and physical real central coset. No lattice, polarization, cubical/Poincare line, source arrow, compact factor or selected CM phase is replaced. Primary formula: Gorodetsky and Hazan(2017), arXiv1707.07308v2, Theorem2.1/equation2.1 and Corollary2.2, printed5, retaining Green attribution. Inherited context only: SGA3 I4.7.3/4.7.3.1; SGA6 VI Theorem3.3; FGA232 Theorem3.1 pp148--149; FGA236 Corollaries1.7/2.3 pp229/232; SGA7 I VII2.7/2.9 pp169--174 and VIII7.4(b) pp310--311; cote4 scan003 sections8.1.1--8.1.2 and scans005--006 sections8.4.1--8.4.3; cote54 scans005--007. No new Dieudonne, lambda-ring/K-theory, l-adic-weight or CM-motive comparison. DP\_MULT Q3e/DP44 and THETA3 G1/A7 gain the strictly smaller conditional finite gate; CX\_GRH CX-ALG-024 and GCT\_PLETH T-052 retain their arithmetic/global gates.

\clearpage

\section{AM-LIFT-106 --- Admissible rank and the actual all-n short-value supplier}\label{am-lift-106-admissible-rank-and-the-actual-all-n-short-value-supplier}

\textbf{Canonical checked integration --- 2026-11-23T13:18:23.409985+00:00.} Complete saved-proof and executable reviews and separate exact diagnostic replays passed. No numerical error charge is introduced. Frozen author history is preserved. AM-LIFT-106 AR1-18 proves that actual AdmX forces n\textgreater=p+q for both Boolean branches on the original finite projective OE lattice. Below that rank shortGen is empty and both original ShortValTot conjuncts are vacuous for every already selected DeltaDat. Actual PinDat gives p+q=R.r, including the terminal definite datum. Above that rank the actual dPhi range supplies exactly the character hypotheses of accepted R101; unbundled dLev/dNab/newA1/PinDat therefore fill all-n shTot with the original data unchanged. Carrying the other eight supplied fields completes LawsSh. Cusp and uniqueness below r are vacuous; upper-range cusp/spanning, selected arithmetic nonzero/content, source/package existence, period,032 and evaluatedDP44/RH remain separate. No new charge;137/31 arrays and27/11 floors unchanged.

Author: \nolinkurl{/root/lift_accounting}. Author whole-proof/source/script verification is complete; the exact author diagnostic and freeze records are adjacent. Independent review and canonical acceptance remain separate. It treats the supplied interfaces as the original assumptions; it is not a Lean build or an archive-validity audit. It changes no selected \nolinkurl{DeltaDat}, coefficient field, lattice, level, finite function, character, operator or class sum.

\subsection{1. The precise improvement}\label{the-precise-improvement}

Write \(O=\mathcal O_E\), \(L=\mathrm{Dw.La}\), \(r_0=\mathrm{Dw.p}+\mathrm{Dw.q}\), \(b=l.\mathrm{Mb}\), and

\[
 J=\overline b=\mathrm{Mfr}(l),\qquad
 b\overline b=(M),\qquad M=l.M\ge3. \tag{AR1}
\]

The source \nolinkurl{Datum} makes \(L\) finite projective over \(O\) and gives its \(O\)-module rank as \(r_0\). Its actual admissibility condition is

\[
 \operatorname{AdmX}(\mathrm{rk},x_0)
 \Longleftrightarrow
 \operatorname{span}_O\{x_0(i):i\in\operatorname{Fin}(n),\ 
                  \mathrm{rk}=\mathrm{true}\ \text{or}\ i<r_0\}
       +JL=L. \tag{AR2}
\]

\textbf{Theorem A.} For either Boolean value and every tuple \(x_0:\operatorname{Fin}(n)\to L\),

\[
                \operatorname{AdmX}(\mathrm{rk},x_0)
                    \Longrightarrow r_0\le n. \tag{AR3}
\]

Consequently \nolinkurl{shortGen} is empty and both conjuncts of \nolinkurl{ShortValTot} are vacuous when \(n<r_0\), for every already selected \nolinkurl{DeltaDat}. This is an impossibility of an antecedent, not a claim that a chosen scalar or section is zero.

For the actual pinned family, \(r_0=D.R.r\). AM-R-101 constructs both short-value conjuncts from the supplied \nolinkurl{IsLev}, new \(A=1\), operative \nolinkurl{IsNab}, \nolinkurl{PhiPin} and \nolinkurl{PsiTie}. The literal \nolinkurl{dPhi} provides the last two predicates only for \(n\ge D.R.r\), exactly the range remaining after AR3. Thus these unbundled source fields supply the \textbf{all-\(n\)} \nolinkurl{shTot} field for the unchanged family. Sections 5--7 state the quantifiers and the record-completion consequence explicitly.

\subsection{2. The ideal is proper and every residue fibre has the stated rank}\label{the-ideal-is-proper-and-every-residue-fibre-has-the-stated-rank}

The ring \(O\) is the full integer ring of a number field and hence an integral domain with fraction field \(E\). The field has characteristic zero. In particular \(M\ne0\), so the principal ideal \((M)\) is nonzero.

It is also proper. If \(M\) were a unit in \(O\), then \(1/M\in O\cap\mathbf Q\) would be a rational algebraic integer. More explicitly, a monic integral equation for \(1/M\), multiplied by its degree-th power of \(M\), would give an equation \(1+Mz=0\) with \(z\in\mathbf Z\), impossible for \(M\ge3\). This argument uses the full integer ring, not the field \(E\) or a localization in which \(M\) has been inverted.

The source \nolinkurl{H.bar} is a ring automorphism and its induced map on ideals preserves and reflects the top ideal. If \(J=O\), then \(b=O\), contradicting \(bJ=(M)\ne O\). If \(J=0\), its product with \(b\) is zero, contradicting \((M)\ne0\). Hence

\[
             0\ne J\ne O. \tag{AR4}
\]

Choose a maximal ideal \(\mathfrak p\supseteq J\), and write \(k_{\mathfrak p}=O/\mathfrak p\). Such an ideal exists because \(J\) is proper. In fact the following rank assertion holds at every maximal ideal of \(O\):

\[
 \dim_{k_{\mathfrak p}}(L/\mathfrak pL)=r_0. \tag{AR5}
\]

Here is the complete rank comparison. Finite projectivity makes \(L_{\mathfrak p}\) finite free over \(O_{\mathfrak p}\); this is the implication (2) to (5) in Stacks Lemma 10.78.2, tag 00NX. If its rank is \(s\), extension to \(E\) gives

\[
 E\otimes_O L\simeq E^s,\qquad
 L/\mathfrak pL\simeq k_{\mathfrak p}^{s}. \tag{AR6}
\]

To match \(s\) with the literal \nolinkurl{Module.finrank O L}, finite projectivity embeds \(L\) into a finite free module, so \(L\) is torsion free. A finite family in \(L\) is \(O\)-linearly independent if and only if its images in \(E\otimes_O L\) are \(E\)-linearly independent: clear all denominators in any proposed relation. Conversely choose an \(E\)-basis of \(E\otimes_O L\). Its finitely many entries are finite linear combinations of images of elements of \(L\); multiplying each basis vector by a nonzero common denominator places it in the image of \(L\), preserving independence. Thus the maximum independent-family size is the finite dimension of \(E\otimes_O L\). The source rank equality identifies that dimension with \(r_0\). This proves AR5 without a global free \(O\)-basis or a class-number assumption.

Nothing in AR4--AR6 uses a principal generator of \(b\), splitting of any chosen rational prime, or the relative-unramified field hypothesis. Those stronger conditions remain in the actual source but are unnecessary for this lower bound.

\subsection{3. The admissible rank bound, including the Boolean exception}\label{the-admissible-rank-bound-including-the-boolean-exception}

Let

\[
 I_{\mathrm{rk}}=
 \{i\in\operatorname{Fin}(n):\mathrm{rk}=\mathrm{true}\text{ or }i<r_0\}.
 \tag{AR7}
\]

Assume AR2 and reduce the sum of submodules modulo a maximal ideal \(\mathfrak p\supseteq J\). Every element of \(JL\) maps to zero. Therefore the images of the displayed tuple entries span \(L/\mathfrak pL\). By AR5, ordinary vector-space dimension gives

\[
 r_0\le |I_{\mathrm{rk}}|=
 \begin{cases}
 n,&\mathrm{rk}=\mathrm{true},\\
 \min(n,r_0),&\mathrm{rk}=\mathrm{false}.
 \end{cases} \tag{AR8}
\]

In both cases \(r_0\le n\), proving Theorem A. In particular setting \nolinkurl{rk := decide (n = nV)} cannot bypass the inequality, even if \(nV\) were an arbitrary external natural number. No hypothesis about \(nV\), and no selection of a different tuple, is needed.

There is a useful exact refinement in the false branch. If \(n\ge r_0\), the map defined by its first \(r_0\) entries is

\[
 (O/J)^{r_0}\longrightarrow L/JL. \tag{AR9}
\]

For \nolinkurl{rk=false}, AR2 is precisely its surjectivity. The target is finite projective of constant rank \(r_0\) after base change from \(O\). Localizing at every maximal ideal of \(O/J\), the map is a surjective square matrix over a local ring, hence invertible: reduction gives nonzero determinant in the residue field, so its determinant is a unit. Thus AR9 is an isomorphism. Conversely an isomorphism AR9 gives AR2. Therefore, in the false branch, admissibility says that these first \(r_0\) classes are a genuine \(O/J\)-basis, not merely a rationally independent family. This is a consequence about the tuple already supplied; no admissible tuple is constructed here.

The case \(r_0=0\) causes no exception. AR6 and torsion freeness imply \(L=0\); the empty basis statements have their ordinary meaning. The actual pinned family below has positive rank, but the general theorem does not need that restriction.

\subsection{4. Exact small-n consequences for the original predicates}\label{exact-small-n-consequences-for-the-original-predicates}

Fix any old \nolinkurl{Moduli} and \nolinkurl{Coeff}, any \nolinkurl{bad,nV,a}, any \(n<r_0\), and any existing \nolinkurl{DeltaDat d} with that number of columns. No level, connection or character hypothesis on \(d\) is needed in this section.

First, the literal definition \nolinkurl{shortGen}, \nolinkurl{ThShort.lean271–275}, requires an existential tuple \(x_0\) satisfying AR2. Theorem A excludes every such tuple. Hence

\[
 \operatorname{shortGen}(\mathrm{nab},\mathrm{bad},nV,a,n,d)
                         =\varnothing. \tag{AR10}
\]

The set is empty, rather than the singleton containing the zero rule. A rule cannot enter it without the impossible admissibility witness, regardless of its value. The externally displayed \nolinkurl{_nab} does not affect this argument, and no condition on the operative \nolinkurl{d.nab₂} is used either.

Second, \nolinkurl{ShortValTot}, lines287--294, quantifies over \(\Lambda,N,x_0,P\) and then assumes the four lattice conditions, \nolinkurl{AdmX}, and harmonic membership before asserting a conjunction. The first member of that conjunction is the all-ring unique scalar statement:

\[
 \forall S,y_2,\ \exists!v:\ S,\quad
        \operatorname{ShortValAt}(d.\mathrm{nab}_2,\ldots,y_2,v).
 \tag{AR11}
\]

The second is the original-base section statement:

\[
 \exists f\in\operatorname{sec}(X_o,c,a,n),\quad
             \operatorname{IsShortRule}(d,\Lambda,N,x_0,P,f).
 \tag{AR12}
\]

For \(n<r_0\), both AR11 and AR12 remain inside the implication whose \nolinkurl{AdmX} premise is false. Logical elimination of that premise proves \nolinkurl{ShortValTot}. It does \textbf{not} assert AR11 for arbitrary inadmissible input, produce a zero scalar, claim an arbitrary-ring base-change theorem, or construct a section satisfying AR12 for an inadmissible tuple. This distinguishes the two quantified conclusions from unconditional existence claims.

Third, AR10 makes \nolinkurl{ShortInCusp} true for every proposed cusp submodule. Likewise the short-generator sets for any two choices of the operative connection are both empty, so the equality required by \nolinkurl{LawsFull.shUniq} holds below \(r_0\), without needing connection independence there. These statements say nothing about cusp membership or connection comparison in the nonempty range.

\subsection{5. Exact identification with the actual pinned rank}\label{exact-identification-with-the-actual-pinned-rank}

Let \(D\) and \(T:\mathrm{Top}(D)\) be actual source data. Assume the unbundled \nolinkurl{T.Q.PinDat} and fix \(j\) with \nolinkurl{InR D j}, so \(1\le j\le D.R.m\). The literal \nolinkurl{IsDatW} gives

\[
 (T.Q.Dw(j)).p=D.R.k+j-1,\qquad
 (T.Q.Dw(j)).q=D.R.m-j. \tag{AR13}
\]

Both subtractions are the displayed natural-number subtractions. Since \(j\ge1\) and \(j\le D.R.m\), they do not truncate in this sum. Thus

\[
 (T.Q.Dw(j)).p+(T.Q.Dw(j)).q
 =D.R.k+D.R.m-1=D.R.r. \tag{AR14}
\]

This includes \(j=1\) and the terminal \(j=D.R.m\), where the second signature is zero. The same deduction is explicitly written in \nolinkurl{DepRun.lean87–92}; the proof here uses its hypotheses directly, not an unrecorded signature convention. \nolinkurl{LawsM.pins} supplies \nolinkurl{Par.Pins}, whose fourth member is \nolinkurl{PinDat}; \nolinkurl{LawsCore} extends \nolinkurl{LawsM}, and \nolinkurl{LawsP} extends \nolinkurl{LawsCore}. These are legitimate rank suppliers that do not involve \nolinkurl{LawsSh.shTot}.

For orientation, the actual \nolinkurl{Ranks} constraints also imply \(r>0\), and \nolinkurl{nV=r+t>r}. Neither fact is needed to treat the two Boolean branches in AR8. The lower-bound argument remains valid before imposing the complete size windows.

\subsection{6. The unchanged-family all-n supplier theorem}\label{the-unchanged-family-all-n-supplier-theorem}

\textbf{Theorem B.} Retain the same actual \(D,T\) and \nolinkurl{PinDat}. For every \(j\) in \nolinkurl{InR}, suppose that the following predicates are supplied for every \(n\ge D.R.r\):

\begin{enumerate}
\def\labelenumi{\arabic{enumi}.}
\tightlist
\item
  \nolinkurl{(T.dδ j n).IsLev};
\item
  \nolinkurl{(T.dδ j n).l₂.A = 1};
\item
  an actual operative \nolinkurl{IsNab} for \nolinkurl{(T.dδ j n).nab₂}, with its required \nolinkurl{StarRing};
\item
  actual \nolinkurl{ph,psi,x₀} satisfying \nolinkurl{(T.dδ j n).PhiPin ph psi x₀} and \nolinkurl{(T.dδ j n).PsiTie psi}.
\end{enumerate}

Then for every \(j\) in \nolinkurl{InR}, every \(a,n\) in the original \nolinkurl{Dom} range, the literal conclusion is

\[
 \operatorname{ShortValTot}(T.Q.Xo(j),T.Q.c,T.\mathrm{nab}(j),
       \mathrm{badN}(D),D.R.nV,a,n,T.d\delta(j,n)). \tag{AR15}
\]

\textbf{Proof.} Split the natural-number alternatives \(n<r\) and \(r\le n\). In the first, AR14 and section4 prove both conditional conjuncts of AR15 for the already selected record. In the second, apply AM-R-101 FD1--23 to that exact record and the four supplied predicates. Its all-ring value is the accepted LIFT103 construction, and its original-base section is the canonical half-level/Fourier descent. It retains the original coefficient identification, all normalizing factors, \nolinkurl{cδ}, \nolinkurl{φδ}, stored operative connection, supplied \nolinkurl{FrOver} witnesses and class/level data. This proves AR15 for all \(n\). No choice or equality for an arbitrary preselected absent character package is inferred. \(\square\)

The theorem only needs the four inputs in the upper range. In the actual source, \nolinkurl{LawsSh.dLev} and \nolinkurl{dNab} have the stronger all-\(n\) quantifier; the separate \nolinkurl{LawsFull.lvA} clause has the stronger all-\(n\) equality. The actual \nolinkurl{dPhi} is exactly

\[
 \forall j\in\mathrm{InR},\ \forall n\ge r,\quad
 \exists\mathrm{ph},\psi,x_0:\quad
 \psi\ne1\ \wedge\ \mathrm{PhiPin}\ \wedge\ \mathrm{PsiTie}.
 \tag{AR16}
\]

It therefore supplies every character datum needed in Theorem B. The nontriviality condition \(\psi\ne1\) remains an actual hypothesis of that source clause, even though R101's transformation/descent calculation does not use it. For an actual indexed tuple package, \nolinkurl{CovPin} supplies the first three predicates together; the character predicates must still be supplied for the same operative record. Merely having some unrelated nonzero finite function is not a supplier.

There is no circular use of a completed law record. The theorem takes the displayed clauses as \textbf{unbundled hypotheses}. The occurrences of their names inside \nolinkurl{LawsSh} and \nolinkurl{LawsFull} identify exact source interfaces, not a proof that those entire records exist. In particular it never assumes \nolinkurl{LawsSh.shTot}, or \nolinkurl{LawsFull.sh}, while constructing AR15.

\subsection{7. Record completion and the remaining ranges}\label{record-completion-and-the-remaining-ranges}

Theorem B gives a direct record-completion result. Suppose that the same actual \(T\) has \nolinkurl{PinDat}, the separate all-\(n\) new-\(A=1\) clause, and the following supplied fields with their literal quantifiers:

\[
 \mathrm{beta},\ \mathrm{dLev},\ \mathrm{dNab},\ \mathrm{dPd},\quad
 \mathrm{dPhi},\ \mathrm{dRoot},\ \mathrm{shCusp},\ \mathrm{shNe}.
 \tag{AR17}
\]

Then inserting AR15 as \nolinkurl{shTot} completes \nolinkurl{LawsSh T}; all eight other fields are carried verbatim. This removes \nolinkurl{shTot} as an independent law obligation \textbf{conditional on those actual remaining fields}. It does not construct the finite root-class data, prescribed finite operator, cusp space, or nonzero selected image.

The input \nolinkurl{shCusp} in AR17 may itself be restricted to \(n\ge r\), because AR10 proves its lower range. The same reasoning extends a supplied upper-range \nolinkurl{shUniq} equality to the original all-\(n\) field of \nolinkurl{LawsFull}; both generator sets are empty below \(r\), including after the literal update of \nolinkurl{nab₂}. This logical extension does not prove the upper-range comparison or construct the other \nolinkurl{LawsFull} fields.

There is no conflict with \nolinkurl{shNe} or \nolinkurl{SpanningAtDominantWeight}. Their displayed threshold is

\[
 n\ge 8\left(\frac{[E:\mathbf Q]}2\right)r_0+4\ge r_0. \tag{AR18}
\]

The half-degree is a positive integer. One can see positivity already from the source \nolinkurl{Host.rt_sq}: \(E\) contains a root of \(X^2+d\), with \(d\) a positive prime, so it contains an imaginary quadratic field and has degree at least2. The existing CM assumption gives the same conclusion. Thus the nonzero and spanning laws never ask for a member in the empty range. AR18 verifies compatibility of quantifiers; it proves neither law in its substantive range, and it does not prove that the cusp submodule is nonzero.

All data and costs are literally retained:

\begin{longtable}[]{@{}
  >{\raggedright\arraybackslash}p{(\columnwidth - 2\tabcolsep) * \real{0.5000}}
  >{\raggedright\arraybackslash}p{(\columnwidth - 2\tabcolsep) * \real{0.5000}}@{}}
\toprule\noalign{}
\begin{minipage}[b]{\linewidth}\raggedright
Item
\end{minipage} & \begin{minipage}[b]{\linewidth}\raggedright
Effect of this theorem
\end{minipage} \\
\midrule\noalign{}
\endhead
\bottomrule\noalign{}
\endlastfoot
\nolinkurl{Host}, \nolinkurl{Sizes}, \nolinkurl{Par}, \nolinkurl{PinHost}, \nolinkurl{PinLevel}, \nolinkurl{PinField}, \nolinkurl{PinDat} & Existing inputs retained; only \nolinkurl{PinDat} is used for AR14. \\
Every selected \nolinkurl{T.dδ j n}, including \(n<r\) & Unchanged. No zero completion or replacement is performed. \\
\nolinkurl{pd}, \nolinkurl{I}, old/new levels, moduli, \nolinkurl{r}, \nolinkurl{m}, \nolinkurl{hF}, \nolinkurl{eF} & Unchanged. The only newly filled field is the proposition \nolinkurl{shTot}. \\
\nolinkurl{nab₂}, \nolinkurl{φδ}, \nolinkurl{cδ}, finite phases and normalizations & Unchanged; R101 is applied to the same operative data in its upper range. \\
\nolinkurl{dPd}, \nolinkurl{dRoot}, \nolinkurl{halfI}, \nolinkurl{XiVal}, fixed \nolinkurl{xi}, coupled coefficients and class sums & No field is changed, so no transport or replacement identity is required. Their existence remains separate. \\
Integer/metric/field/period costs and the finite error arrays & No new cost or charge; current DO127 137/31 arrays and floors 27/11 are unchanged. \\
Integral selected image, upper-content witness, global period, exact032 package and DP44 & No nonvanishing, equality, admission or sign is inferred. \\
\end{longtable}

\subsection{8. Literal and primary source addresses}\label{literal-and-primary-source-addresses}

All paths in this table are relative to \nolinkurl{/tmp/alpa66/work}. The companion source JSON records the exact file hashes and the actual read extents.

\begin{longtable}[]{@{}
  >{\raggedright\arraybackslash}p{(\columnwidth - 2\tabcolsep) * \real{0.5000}}
  >{\raggedright\arraybackslash}p{(\columnwidth - 2\tabcolsep) * \real{0.5000}}@{}}
\toprule\noalign{}
\begin{minipage}[b]{\linewidth}\raggedright
File
\end{minipage} & \begin{minipage}[b]{\linewidth}\raggedright
Exact used interface
\end{minipage} \\
\midrule\noalign{}
\endhead
\bottomrule\noalign{}
\endlastfoot
\nolinkurl{research/alpoge_source/ReciprocityTheory/ReductiveGroups/ArithmeticGeometry/BirationalGeometry/ToroidalCompactification/TwoWeightEqOneAddRoot.lean} & \nolinkurl{Host}14--26; finite/projective/rank \nolinkurl{Datum}30--42; \nolinkurl{Level}53--62, especially \nolinkurl{three_le}55 and \nolinkurl{prod}59. \\
\nolinkurl{research/alpoge_source/Formalism/LocalToGlobal/ProjectiveLattices/ThShort.lean} & \nolinkurl{Mfr}93, \nolinkurl{AdmX}107--109; \nolinkurl{DeltaDat}114--136; \nolinkurl{shortGen}271--275; both \nolinkurl{ShortValTot} conjuncts287--294; \nolinkurl{ShortInCusp}296--298; spanning/nonzero322--333. \\
\nolinkurl{research/alpoge_source/Formalism/LocalToGlobal/ProjectiveLattices/PartitionExtendByZeroAntitone.lean} & \nolinkurl{InR}12, \nolinkurl{IsDatW}24--31, \nolinkurl{PinDat}47, \nolinkurl{Pins}49. \\
\nolinkurl{research/alpoge_source/Formalism/LocalRingCardinality/Params.lean} & \nolinkurl{Ranks}6--22, \nolinkurl{r}32, \nolinkurl{nV}34, \nolinkurl{sigW}46; \nolinkurl{FactSigW}81 is a declared proposition, not used as an assumed proof. \\
\nolinkurl{research/alpoge_source/Formalism/LocalToGlobal/ProjectiveLattices/DepRun.lean} & Explicit source deduction of the same rank equality87--92 inside \nolinkurl{nuFinAt_of}. \\
\nolinkurl{research/alpoge_source/Formalism/LocalToGlobal/ProjectiveLattices/CuspidalL2Subspace.lean} & \nolinkurl{LawsM.pins}93--95 and \nolinkurl{LawsCore}131. \\
\nolinkurl{research/alpoge_source/Formalism/LocalToGlobal/ProjectiveLattices/HasDataCondition.lean} & \nolinkurl{Top}20--23; \nolinkurl{LawsP}32; exact \nolinkurl{LawsSh}57--73, especially restricted \nolinkurl{dPhi}64 and all-\(n\) \nolinkurl{shTot}69--70. \\
\nolinkurl{research/alpoge_source/Formalism/LocalToGlobal/ProjectiveLattices/FullExists.lean} & Fixed \nolinkurl{xi}20; \nolinkurl{XiVal}28--31; \nolinkurl{lvA}40; \nolinkurl{halfI}42; all-\(n\) \nolinkurl{shUniq}44--47; \nolinkurl{xiVal}49. \\
\nolinkurl{research/alpoge_source/Formalism/LocalToGlobal/ProjectiveLattices/ThClassPinsRec.lean} & Exact \nolinkurl{CovPin}62--63, including new \(A=1\) and the operative \nolinkurl{nab₂}. \\
\nolinkurl{agents/canonical_half_level_short_rule_descent.md} & Accepted AM-R-101 FD1--23; exact character-supplied upper-range theorem, promoted13:04:52.861846 UTC after complete independent review. \\
\end{longtable}

The independent mathematical input for local freeness was freshly read in the full displayed statement and proof of Stacks Lemma10.78.2. The remaining rank, quotient-generation and logical arguments are proved above. The companion exact diagnostics test residue matrix ranks, both Boolean selectors, quotient bases, natural-number endpoints and quantified range completion. They are not substitutes for the universal projective-module argument or R101's geometric descent proof.

\subsubsection{References}\label{references-8}

The Stacks Project Authors. (2026). \emph{The Stacks Project}, Lemma10.78.2, tag00NX, ``Finite projective modules.'' \url{https://stacks.math.columbia.edu/tag/00NX}

ALP\_MACHINE. (2026). \emph{AM-R-101: Canonical half-level descent of the original short rule} (frozen source-level research proof, FD1--23). Accepted canonical SHA256 \nolinkurl{f76cdf6160daafd0e6df128df8b303e2e09910a8244bdba0d4203bb116cc218c}; source and independent-review records are retained alongside the proof.

\subsection{Canonical algebraic-geometric context}\label{canonical-algebraic-geometric-context-89}

Over Spec OE the physical CM lattice is a vector bundle of constant rank p+q. The proper original half-level ideal has a residue point, and generation of its restriction requires at least that many sections. The false-Boolean branch is an actual basis of the restricted vector bundle. This finite-flat rank statement closes the missing lower-column range of the characteristic-zero framed theta descent without changing any Picard/Poincare or class-character data. Primary locator: Stacks Lemma10.78.2 (00NX). Exact source interfaces: TwoWeightEqOneAddRoot14-62; ThShort93,107-109,271-298,322-333; PartitionExtendByZeroAntitone12,24-49; Params32,34,46; DepRun87-92; HasDataCondition57-73; FullExists40-47. Carried exact context: SGA7 I VII2.7/2.9 pp169--174 and VIII7.4(b) pp310--311; FGA232 Theorem3.1 pp148--149 and FGA236 Corollaries1.7/2.3 pp229/232; SGA3 I4.7.3/4.7.3.1; SGA6 VI Theorem3.3; cote4 scan003 sections8.1.1--8.1.2 and scans005--006 sections8.4.1--8.4.3; cote54 scans005--007. These are carried contextual readings, not new whole-text checks. No new integral extension over Z, Dieudonne slope, lambda-ring positivity, K-theory filtration, l-adic weight theorem or CM-motive L-function comparison follows. DP\_MULT Q3e/DP44/DP52/DP54--65 and THETA3 G1/A7 gain this conditional source-rule reduction; CX\_GRH CX-ALG-024 and GCT\_PLETH T-052 retain their global gates.

\clearpage

\section{AM-R-102 --- Cited beta law and the prescribed coefficient-field supplier}\label{am-r-102-cited-beta-law-and-the-prescribed-coefficient-field-supplier}

\textbf{Canonical checked integration --- 2026-11-23T13:30:52.108107+00:00.} Complete saved-proof and executable reviews and separate exact diagnostic replays passed. No numerical error charge is introduced. Frozen author history is preserved. AM-R-102 CB1-8 applies accepted GlobalTheory.BetaEx.betaExists directly to every original T.Q.Dw(j), T.Q.c, n and badN(D), filling LawsSh.beta. KPsiIn remains the internal antecedent; the original selected betaH has its successful-choice specification only when that antecedent is supplied. Together with accepted LIFT106, the seven other supplied clauses and separate PinDat/newA1 assemble LawsSh without any DeltaDat change. Accepted kPsiIn\_iff/reflexRoots\_iff/reflexLe\_of\_basis identify the exact same-field character supplier; the actual class/root data are not constructed. The accepted coefficient constructor supplies an eventual extension-field law, not the prescribed eO/PinField identity. No duplicate determinant theorem or archive audit is credited. No charge;137/31 arrays and27/11 floors unchanged; actual source, selected nonzero/content, cusp/spanning, period,032 and evaluatedDP44/RH remain separate.

Author: \nolinkurl{/root}. This result records applicability of accepted archive theorems and the resulting conditional source-record reduction. It does not re-prove those theorems, validate the archive, or claim a new determinant-index theorem. Author and independent checks are recorded separately.

\subsection{1. The exact accepted beta supplier}\label{the-exact-accepted-beta-supplier}

At the fixed recovered source commit \nolinkurl{7cbf7a0c32fca1badf05931879a37c3f8a2b03b4}, Alpöge's archive declares

\[
 \texttt{GlobalTheory.BetaEx.betaExists}:
 \quad\forall E,H,\mathrm{Dw},l,c,n,\mathrm{bad},\quad
 \mathrm{BetaExists}\,(\mathrm{Dw}:=\mathrm{Dw})\,c\,n\,\mathrm{bad}.
 \tag{CB1}
\]

The exact declaration is \nolinkurl{Development/AbelianVarietyTheory/Proofs/UnitaryExistence/ExistsIndexIdeal.lean322–324}; its source body is325--341. The field, number-field and CM-field instances and the original \nolinkurl{Host}, \nolinkurl{Datum}, \nolinkurl{Level} and \nolinkurl{Coeff} arguments are all explicit. CB1 is cited as an accepted theorem, in accordance with the task. It is not assigned new theorem credit here.

The target predicate is the literal \nolinkurl{ThShort.lean283–285} implication

\[
 \mathrm{KPsiIn}(\mathrm{Dw},c)\Longrightarrow
 \forall\Lambda:\mathrm{GenLat}(E,n),\quad
 \Lambda.\mathrm{Unimod}\Longrightarrow\Lambda.\mathrm{TrivAt}(\mathrm{bad})
 \Longrightarrow
 \exists\beta\ne0,\quad\mathrm{BetaIs}(\Lambda,\beta)
                         \wedge\mathrm{BetaAbs}(\Lambda,\beta).
 \tag{CB2}
\]

Thus CB1 supplies the proposition \nolinkurl{BetaExists} for every \(n\) and \nolinkurl{bad}, with \nolinkurl{KPsiIn} still inside its antecedent. It neither assumes a positive \nolinkurl{TrivAt} witness nor replaces the source's unrestricted \nolinkurl{GenLat} by a globally free lattice. These are statements of the accepted theorem's scope, not claims that its proof has been audited.

\textbf{Application A.} For every already existing original \nolinkurl{Top T}, CB1 fills the field

\[
 \forall j\in\mathrm{InR},\ \forall n,\quad
 \mathrm{BetaExists}\,(\mathrm{Dw}:=T.Q.\mathrm{Dw}(j))\,T.Q.c\,n\,
                         (\mathrm{badN}(D)). \tag{CB3}
\]

\textbf{Proof of the application.} Instantiate CB1 with the actual \(E=D.\mathrm{host}.E\), \(H=T.Q.H\), \(\mathrm{Dw}=T.Q.\mathrm{Dw}(j)\), \(l=T.Q.l\), \(c=T.Q.c\), the given \(n\), and \nolinkurl{badN D}. These are exactly the arguments in \nolinkurl{HasDataCondition.lean58}. No equality of different coefficient records, lattice replacement, or additional source existence is needed. This proves the literal \nolinkurl{LawsSh.beta} field. \(\square\)

\subsection{2. The actual selected scalar and all-n record completion}\label{the-actual-selected-scalar-and-all-n-record-completion}

For a particular self-dual and \nolinkurl{TrivAt} lattice, an actual \nolinkurl{KPsiIn} supplier discharges the antecedent in CB2. The existential condition in the original \nolinkurl{betaH} definition is then true, so its selected element satisfies

\[
 \mathrm{betaH}(\Lambda)\ne0,\qquad
 \mathrm{BetaIs}(\Lambda,\mathrm{betaH}(\Lambda)),\qquad
 \mathrm{BetaAbs}(\Lambda,\mathrm{betaH}(\Lambda)). \tag{CB4}
\]

This is the specification of the successful choice in \nolinkurl{ThShort.lean84–85}. The unsuccessful branch is one. No equality between the selected element and an independently produced witness is asserted, and its phase remains the original one. The norm in \nolinkurl{BetaAbs} is the full additive index of \nolinkurl{Ln} in \nolinkurl{La}, raised to the physical rank \(p+q\); it is not \nolinkurl{mIdx} or \nolinkurl{zCard}. All complex embeddings of the actual coefficient field remain in that predicate.

\textbf{Application B.} Assume actual \nolinkurl{PinDat}, the separate new-\(A=1\) clause, and the seven fields \nolinkurl{dLev}, \nolinkurl{dNab}, \nolinkurl{dPd}, \nolinkurl{dPhi}, \nolinkurl{dRoot}, \nolinkurl{shCusp}, \nolinkurl{shNe} for the same original family, with their literal quantifiers. Then \nolinkurl{LawsSh T} follows without changing any \nolinkurl{DeltaDat}.

\textbf{Proof.} CB3 supplies \nolinkurl{beta}. Accepted AM-LIFT-106 AR1--18 supplies \nolinkurl{shTot}: below the physical rank the admissibility antecedent is impossible; at and above it the original \nolinkurl{dPhi} range supplies accepted AM-R-101's character hypotheses. Insert these two propositions alongside the seven supplied fields into the nine-field record \nolinkurl{HasDataCondition.lean57–73}. This is an assembly of unbundled assumptions and never assumes a complete \nolinkurl{LawsSh} or \nolinkurl{LawsFull}. \(\square\)

The separate \nolinkurl{LawsP.kPsi} field is not needed merely to prove that the implication CB3 holds. It is needed when the actual witness conclusion CB4 is used. This distinguishes a completed implication field from existence of the required character data.

\subsection{3. Equivalent root data in the same prescribed field}\label{equivalent-root-data-in-the-same-prescribed-field}

The accepted \nolinkurl{GlobalTheory.LatticeBetaCompare.kPsiIn_iff} and \nolinkurl{reflexRoots_iff} give

\[
 \mathrm{KPsiIn}(\mathrm{Dw},c)
   \Longleftrightarrow\mathrm{ReflexLe}(\mathrm{Dw},c)
   \Longleftrightarrow\mathrm{ReflexRoots}(\mathrm{Dw},c).
 \tag{CB5}
\]

Exact declarations: \nolinkurl{MultAddMultConjugate.lean269–272} and460--474; the consolidated \nolinkurl{typings_agree} declaration557--582 includes both equivalences and the finite-basis implication below. These are accepted input theorems. Their algebraic proofs are not repeated.

The definition \nolinkurl{ReflexNormIndexIdeal.lean64–66} makes the final predicate completely explicit: for each nonzero integral ideal \(\mathfrak a\subset\mathcal O_E\), it asks for

\[
 h>0,\quad\gamma\in\mathcal O_E,\quad\psi\in c.F,\qquad
 \mathfrak a^h=(\gamma),\qquad
 \psi^h=\prod_{\tau:E\hookrightarrow\mathbf C}
                   \mathrm{embF}(c,\tau)(\gamma)^{\mathrm{Dw.mult}(\tau)}.
 \tag{CB6}
\]

All entries, including \(\psi\), lie in the same already prescribed field. The equivalence does not say that arbitrary previously fixed choices of \(h\) and \(\gamma\) admit such a root: the existence quantifiers in CB6 are retained. In particular it does not discard the unit or root-of-unity adjustment implicit in changing a principal generator.

The accepted \nolinkurl{reflexLe_of_basis}, declaration447--458 of the same file, also gives

\[
       \mathrm{ReflexBasis}(\mathrm{Dw},c)
                    \Longrightarrow\mathrm{KPsiIn}(\mathrm{Dw},c). \tag{CB7}
\]

The precise finite table, defined in \nolinkurl{ReflexNormIndexIdeal.lean68–72}, consists of nonzero ideals \(\mathfrak c_i\), exponents \(h_i\), generators \(\gamma_i\) and elements \(\psi_i\in c.F\). Every class must have a unique expression \(\prod_i[\mathfrak c_i]^{a_i}\) with \(0\le a_i<h_i\), and the table must satisfy \(\mathfrak c_i^{h_i}=(\gamma_i)\) and the CB6 root equality for each \(i\). Thus an actual table satisfying the source definition is sufficient. No such table in the user's prescribed field has been produced by this result, and CB7 is not claimed to be an equivalence for a fixed proposed table.

\textbf{Application C.} Given these same-field root data for every original \(j\) in \nolinkurl{InR}, use CB5 or CB7 to fill the exact \nolinkurl{LawsP.kPsi} clause at \nolinkurl{HasDataCondition.lean55}. Applications A and B then remove \nolinkurl{beta} and \nolinkurl{shTot} as independent obligations in the unchanged family, while their original character supplier remains explicit. This is direct specialization of the cited interfaces, not a claim that the rest of \nolinkurl{LawsP} or the actual family exists.

\subsection{4. What the general coefficient constructor supplies}\label{what-the-general-coefficient-constructor-supplies}

The accepted archive theorem \nolinkurl{GlobalTheory.exists_coeff_typeNorms_principal} has the exact scope

\[
 \forall E,H,\mathrm{Dw},l,c,\quad
 \exists c':\mathrm{Coeff}(l),\quad c.F\subseteq c'.F\ \wedge
 \forall c'':\mathrm{Coeff}(l),\quad
        c'.F\subseteq c''.F\Longrightarrow\mathrm{KPsiIn}(\mathrm{Dw},c'').
 \tag{CB8}
\]

Its declaration is \nolinkurl{ExistsCoeffTypeNormsPrincipal.lean23–25}, with complete source through74. It is an actual theorem rather than the declaration of a desired proposition. It gives an eventual supplier in coefficient fields containing an extension. No numerical bound on that extension's degree is included in this theorem's statement.

For the fixed original \nolinkurl{Par}, however, the coefficient record is \nolinkurl{Q.c}; \nolinkurl{Q.eO} is a specified integer-ring isomorphism to \nolinkurl{D.what₂.O}, and \nolinkurl{PinField} fixes its complex realization. The constructor CB8 does not return these original equalities or an identified original field. Therefore applying it and replacing \nolinkurl{Q.c} by the output would require separate transport of the pinned source package and its field costs. CB8 alone does not discharge the same-field condition in Applications B/C for the already selected package. If its output happened to have the required identified field, that identification would itself have to be supplied.

The source also defines a proposition called \nolinkurl{ReflexFieldExists} with a class-number degree bound. A definition of that proposition is not used here as a theorem that it holds. Likewise \nolinkurl{RecCond.OKPsi} is the universally quantified desired fixed-package law, not a constructor. These are distinctions between literal source interfaces; they are not negative conclusions about whether some other accepted theorem can supply them.

\subsection{5. Geometry, programme consequences and the preserved exploration}\label{geometry-programme-consequences-and-the-preserved-exploration}

The accepted beta theorem identifies a determinant-index line of the inclusion of number-ring lattices with the corresponding weighted type-norm line in the coefficient ring. The source norm is the cardinality of the finite torsion quotient over \(\operatorname{Spec}\mathcal O_E\); the required scalar trivializes the weighted invertible ideal over \(\operatorname{Spec}\mathcal O_{c.F}\) with its stated archimedean norms. CB5/CB7 keep that trivialization in the prescribed field; CB8 permits passage to a larger coefficient field.

Applications A--C give a precise reduction of the original conditional source-construction gate. The beta implication is supplied by the cited archive theorem; all-\(n\) short-value existence is supplied by the accepted R101/LIFT106 composition; the remaining character obligation can be tested against the exact same-field root or finite class-table interfaces. No numerical error charge, selected integral theta image, cusp/spanning theorem, global period, specific032 pairing, actual \(g_j,e_j,H\), DP44 sign or RH claim is supplied. The current137/31 exact error arrays and27/11 support floors are unchanged.

For DP\_MULT Q3e/DP44/DP52/DP54--65 and THETA3 G1/A7 this is a source-law applicability reduction, conditional on original admitted data and the remaining laws. It creates no new Dieudonne slope, integral extension across bad fibres, lambda-ring positivity, l-adic weight theorem or motivic L-function comparison. The existing FGA/SGA Picard and biextension context is carried at its previously recorded exact locators. CX\_GRH CX-ALG-024 and GCT\_PLETH T-052 retain their global gates.

The preliminary nine-section determinant derivation and its unexecuted finite arithmetic script are retained under \nolinkurl{R102_EXPLORATORY_REDIRECT.json}. Once CB1 was recovered, that derivation was withdrawn from promotion because the task requires citation of accepted Alpöge results. It receives no new theorem or movement credit and is not represented as an archive audit. The present result is solely the exact applicability and record-completion argument just proved.

\subsection{6. References and verification scope}\label{references-and-verification-scope-1}

Alpöge, L. (n.d.). \emph{Remark}, step8 and the accompanying Lean archive, fixed source commit \nolinkurl{7cbf7a0c32fca1badf05931879a37c3f8a2b03b4}. \url{https://alpo.ge/remark.pdf}

\begin{longtable}[]{@{}
  >{\raggedright\arraybackslash}p{(\columnwidth - 2\tabcolsep) * \real{0.5000}}
  >{\raggedright\arraybackslash}p{(\columnwidth - 2\tabcolsep) * \real{0.5000}}@{}}
\toprule\noalign{}
\begin{minipage}[b]{\linewidth}\raggedright
Accepted source
\end{minipage} & \begin{minipage}[b]{\linewidth}\raggedright
Exact locator and use
\end{minipage} \\
\midrule\noalign{}
\endhead
\bottomrule\noalign{}
\endlastfoot
\nolinkurl{Development/AbelianVarietyTheory/Proofs/UnitaryExistence/ExistsIndexIdeal.lean} & \nolinkurl{GlobalTheory.BetaEx.betaExists}, declaration322--324: universal original \nolinkurl{BetaExists}. \\
\nolinkurl{Development/AbelianVarietyTheory/Proofs/UnitaryExistence/MultAddMultConjugate.lean} & \nolinkurl{GlobalTheory.LatticeBetaCompare.kPsiIn_iff}269--272; \nolinkurl{reflexLe_of_basis}447--458; \nolinkurl{reflexRoots_iff}460--474; \nolinkurl{typings_agree}557--589. \\
\nolinkurl{Development/AbelianVarietyTheory/BlockTheory/ExistenceTori/ReflexNormIndexIdeal.lean} & \nolinkurl{reflexNorm/reflexIdeal/IsBe/ReflexLe}51--62; exact \nolinkurl{ReflexRoots/ReflexBasis}64--72; \nolinkurl{ReflexFieldExists}110--111 is a proposition definition only. \\
\nolinkurl{Development/AbelianVarietyTheory/Proofs/UnitaryExistence/ExistsCoeffTypeNormsPrincipal.lean} & \nolinkurl{GlobalTheory.exists_coeff_typeNorms_principal}, declaration23--25: extension-field eventual supplier. \\
\nolinkurl{Formalism/LocalToGlobal/ProjectiveLattices/ThShort.lean} & \nolinkurl{BetaIs/BetaAbs/betaH}74--85; \nolinkurl{KPsiIn/BetaExists}277--285. \\
\nolinkurl{Formalism/LocalToGlobal/ProjectiveLattices/HasDataCondition.lean} & Original \nolinkurl{Top}20--23; \nolinkurl{LawsP.kPsi}55; full \nolinkurl{LawsSh}57--73. \\
\nolinkurl{Formalism/LocalToGlobal/ProjectiveLattices/PartitionExtendByZeroAntitone.lean} & \nolinkurl{Par.eO}20 and \nolinkurl{PinField}45: original coefficient-field identification. \\
\nolinkurl{Formalism/LocalToGlobal/ProjectiveLattices/FullExists.lean} & \nolinkurl{lvA}40 and unchanged full-package fields. \\
\nolinkurl{Formalism/LocalToGlobal/ProjectiveLattices/RecCond.lean} & \nolinkurl{OKPsi}7, definition only. \\
\nolinkurl{agents/admissible_rank_all_n_short_value.md} & Accepted AM-LIFT-106 AR1--18, exact all-\(n\) supplier and noncircular record completion. \\
\end{longtable}

The accompanying manifest records the recovered bytes, source URLs and exact root read scopes. Verification checks the complete application proof and the literal supplied/required signatures. No Lean build or proof-validity audit of the accepted archive is performed.

\subsection{Canonical algebraic-geometric context}\label{canonical-algebraic-geometric-context-90}

This is applicability of the accepted determinant-index/type-norm trivialization over the original number-ring schemes, followed by exact assembly of the existing characteristic-zero theta source record. The prescribed integer-ring isomorphism and complex embedding are retained. Accepted source commit7cbf7a0c32fca1badf05931879a37c3f8a2b03b4: ExistsIndexIdeal.betaExists322-324; MultAddMultConjugate.kPsiIn\_iff269-272,reflexLe\_of\_basis447-458,reflexRoots\_iff460-474,typings\_agree557-589; ReflexNormIndexIdeal64-72; ExistsCoeffTypeNormsPrincipal23-25. Literal consumers: ThShort74-85,277-285; HasDataCondition55,57-73; PartitionExtendByZeroAntitone20,45; FullExists40. Carried exact context: SGA7 I VII2.7/2.9 pp169--174 and VIII7.4(b) pp310--311; FGA232 Theorem3.1 pp148--149 and FGA236 Corollaries1.7/2.3 pp229/232; SGA3 I4.7.3/4.7.3.1; SGA6 VI Theorem3.3; cote4 scan003 sections8.1.1--8.1.2 and scans005--006 sections8.4.1--8.4.3; cote54 scans005--007. These are carried contextual readings, not new whole-text checks. No new integral extension over Z, Dieudonne slope, lambda-ring positivity, K-theory filtration, l-adic weight theorem or CM-motive L-function comparison follows. DP\_MULT Q3e/DP44/DP52/DP54--65 and THETA3 G1/A7 gain this conditional source-rule reduction; CX\_GRH CX-ALG-024 and GCT\_PLETH T-052 retain their global gates.

\clearpage

\section{AM-K-123 --- a complete positive-mixture high-coherence cut}\label{am-k-123-a-complete-positive-mixture-high-coherence-cut}

\textbf{Canonical checked integration --- 2026-11-23T13:44:41.209164+00:00.} Complete root and independent proof/script/interface reviews and separate exact replays passed. The actual canonical predecessor has137/31 entries; this integration gives138/31, complete floors27/11. Frozen author history is preserved. AM-K-123 HCX1-26 removes the complete original bounded support27/eligibility27 full-positive scalar-cross globally rank-two-unipotent all-II k160 cell, b=j\_int0 and2\^{}2111\textless x\textless2\^{}3328. The actual all-word selected evaluation gives a chi\textgreater937/1000 restriction with both full arithmetic errors paid. A525772-test outward interval cover proves the complete physical-coherence cap; exact three-real-dimensional Beta disintegration gives a fraction below2\^{}-10 with a paid margin. The same-kernel positive four-component central mixture retains the original affine U scaling, full750 cover, Young regularization and extra transport2. Nine full quadratures, finite Holder anchors, actual-row zero repair, the true1008-dimensional transverse theta, exact Green/J0 correlation, all322 terms and both full tails give one600 charge. Exact138/31 full sums remain below2\^{}-230 and floors27/11 remain. The sole support27 necessary alternative is161\textless=k\textless=164,2\^{}2111\textless x, with upper exponents2130/2140/2150/2165 respectively. Signed PP3, higher supports, actual source/selected arithmetic image/content/nonzero, prescribed period, specific032 value, evaluatedDP44 and RH remain open.

Author: \nolinkurl{/root/kummer_extension}. 23 November 2026 UTC. \textbf{Research-pending complete theorem; author verification, independent review and canonical acceptance remain separate.} This document does not alter a canonical file. The full exact certificate and all original interfaces are part of the result.

\subsection{1. Original sector and conclusion}\label{original-sector-and-conclusion-1}

Retain the original J56 all-top selected kernel, f=a=168, every binary word and canonical twist, the actual polarized projective order modules, CM involution, finite pencils/projectors, two mixed G factors, selected Hilbert and fibre metrics, class/Haar normalization and source fibre. We work only in the remaining support27, globally rank-two-unipotent, full-positive, all-II scalar-cross branch with nonzero original S-cross. Thus total eligibility is27, b=j\_int=0, and \[
  2^{2111}<x=N_ED_\eta<2^{3328},\quad 1\le k\le164,
  \quad t_\tau=\alpha_\tau>0,\quad t_\tau\le R=17424/133.
  \tag{HCX1}
\] The integral isotropic quotient, all-prime row ideal, standard S164 summand and actual selfdual rank3 transverse complement are exactly DO126 IQ1--36. No fixed transverse basis or division by U at its prime divisors is introduced. The original physical central coordinate is w\_tau=Im(h\_nn)\_tau; its affine lattice and U\^{}f Jacobian remain those of DO122/DO125/DO126. The stronger exact Green/J0 identity is the accepted DO127 GC1--16.

\textbf{Theorem.} The entire HCX1 k=160 containing sector has absolute normalized contribution less than 2\^{}-600. In particular this removes its still-open interval 2\^{}2111\textless x\textless2\^{}2120 after DO127. All old ordered errors remain; one600 is appended to the actual DO127137/31 predecessor, with complete floors27/11 unchanged. No k161--164 removal or support28 conclusion is asserted.

The proof first restricts the actual coefficient to high angular coherence, then proves a complete conditional cap on one real-parameter rectangle. Its complement is covered by a genuine positive Laplace mixture. The full scalar/transverse polynomial is subsequently counted with the original finite, ideal, central, source and metric factors.

\subsection{2. A sharper actual selected-evaluation cut}\label{a-sharper-actual-selected-evaluation-cut}

Put B=2\^{}955 and let ell be the original endpoint with 2\^{}484\textless ell\textless2\^{}485. K122 HC6--11 applies to the actual adjoint evaluation ranges in the selected Hilbert metric, not a substitute Fischer overlap. The auxiliary polynomial highest weight satisfies, for s\_omega in{[}0,166{]}, \[
 N=32B-(4+s_\omega)\ell-71,\quad
 \mu_0=33B+12,\quad
 \mu_1=B+(668-\omega_1)\ell-71.
 \tag{HCX2}
\] Consequently mu0−mu1\textless32B and N−mu1\textgreater(154999/5000)B, uniformly in all words and canonical twists. The two-projection factorization is unchanged. Choose b=101/100 and chi0=937/1000. The exact two-singular-value formula HC10 gives sigma\_plus\textless637163/625000: the positive-rational square comparison is included in the certificate. Therefore \[
 \|Q_uQ_v\|\le\sigma_+^{\mu_0-\mu_1}b^{-2(N-\mu_1)}
 <2^{-B/5000}\quad\text{if }|\langle u,v\rangle|\le\chi_0,
 \tag{HCX3}
\] by the exact integer inequality \[
  2(637163/625000)^{160000}(100/101)^{309998}<1.
 \tag{HCX4}
\] Insertion into the literal finite selected trace is precisely K122 HC18, with the actual isometric fibre transport and both compact factors. The common compact change used below acts on both points simultaneously and preserves this projective coherence.

Let I={[}16/529,16/527{]}. The original K120 narrow-band tail, K118 full containing arithmetic mass and K122 product telescoping give \[
 \sum_h|a(h)c(h)|\le\sum_h|a(h)|\prod_\tau F_{I,\chi>\chi_0}(h_{p,\tau})
       +2^{-8605114}+2^{11394885-B/5000}.
 \tag{HCX5}
\] Both errors are fully arithmetic, not merely local constants. The second is below2\^{}-40000000 by an integer comparison. Nonzero original S-cross is explicitly retained in the containing family. No new eta multiplicity, independent rotation or signed-coefficient premise occurs.

\subsection{3. The exact physical normal form and conditional cap}\label{the-exact-physical-normal-form-and-conditional-cap}

At one real place, a simultaneous compact conjugation sends the fixed null line to C(e1+en) and the nonzero horizontal vector to the second positive coordinate. This follows directly from DO098 CF4: first normalize the null vector's negative coordinate by a phase, then apply a unitary positive change and a unitary change on its orthogonal complement. The latter sends the horizontal vector to a real vector of length sqrt(t); its sign can be chosen negative. These are compact conjugations of the same matrix, not independently chosen left/right rotations. They preserve h\_nn, t=(N²)\_nn, the radial measure, the selected compact type, and chi. In this basis the matrix acting on the first two positive and the negative coordinate is \[
 g=\begin{pmatrix}1-A-iw&-v&A+iw\\v&1&-v\\-A-iw&-v&1+A+iw\end{pmatrix},
 \qquad A=t/2,\quad v=\sqrt t,
 \tag{HCX6}
\] and it is the identity on the other positive coordinates. The band-restricted high-coherence fidelity is invariant under inversion, by the change of variable z↔gz and the symmetry of chi. Thus we may compute the forward action of this g; inversion preserves t and replaces w by−w, which is also covered below. Direct multiplication verifies the Hermitian identity g\^{}* diag(1,1,−1) g=diag(1,1,−1), (g−I)\^{}3=0, and ((g−I)\^{}2)\_nn=t. The same construction applies to each original near-point conjugate. Its t is positive because this is a noncentral rank-two unipotent.

Write d=1+A+iw, c=\textbar d\textbar{} and s²=c²−1. For a ball point z put rho=d−(A+iw)z1−v z2. The exact action is \[
 z'_1=1-(1-z_1)/\rho,\quad z'_2=(z_2-v(1-z_1))/\rho,
 \quad z'_\perp=z_\perp/\rho.
 \tag{HCX7}
\] Fix y=1−\textbar z\textbar² and y'=1−\textbar z'\textbar² in I, D=sqrt(y/y'), R\_y²=1−y, and rho=D(d/c)e\^{}\{i theta\}. The two-coordinate denominator fibre is \[
 z_P=z_0+\xi e,\quad
 z_0=(\rho-d)(-A+iw,-v)/s²,\quad e=(v,-A-iw)/s.
 \tag{HCX8}
\] Here z0 is orthogonal to the unit vector e. Put \[
 S=R_y²-|\rho-d|²/s²,\quad
 S_*={2D\over c+1}-y-{(D-1)^2\over s²},\quad
 L²={4Dc\over s²}.
 \tag{HCX9}
\] The remaining positive-coordinate energy is m=S−\textbar xi\textbar². The original coherence satisfies \[
 \chi^2={|n|²\over R_y²(D²-y)},\quad
 n=R_y²+\bar z_1(\rho-1)-v(1-z_1)\bar z_2.
 \tag{HCX10}
\] Expand n=n0+L0 xi+M0 bar(xi)+Q0\textbar xi\textbar². For exact interval evaluation the following identities avoid cancellation. Set \[
 C={A(A²+w²)+iw(s²+2A)\over s^4}.
\] Then \[
 n_0=-y+{D\over c}e^{-i\theta}
       +C\left[cD e^{i\theta}+{D\over c}e^{-i\theta}-D²-1\right],
\] \[
 |L_0|+|M_0|={2A\sqrt t\over s^3}
  \left(|\rho-d|+|\bar d\rho-1|\right),\qquad
 |Q_0|={2A\sqrt{A²+w²}\over s²}.
 \tag{HCX11}
\] These follow by substitution in HCX8--10 and the identity \textbar rho\textbar²=D². They are polynomial/rational identities, not numerical model assumptions.

\textbf{Complete rectangle cap.} For 6\textless=t\textless=48, \textbar w\textbar\textless=4 and every y,y' in I, chi\textgreater chi0 forces \[
                         m<(19/20)S_*.
 \tag{HCX12}
\] If m\textgreater=(19/20)S\emph{, then S\textgreater=(19/20)S}. Choose the representative theta in{[}−pi,pi{]} modulo2pi. Parameterize all such phases by u in{[}−1,1{]} using \[
 h=\sin(\theta/2)=u\sqrt{S_*s²/(80Dc)},\quad
 \cos\theta=1-2h²,\quad\sin\theta=2h\sqrt{1-h²}.
\] The relevant radius for xi is r²=(S*/20)(1−u²). HCX11 bounds \textbar n\textbar{} by \textbar n0\textbar+(\textbar L0\textbar+\textbar M0\textbar)r+\textbar Q0\textbar r². The certificate proves this is strictly less than chi0 R\_y sqrt(D²−y) over the entire rectangle, using outward 72-bit rational intervals and integer square roots. It starts from672 rational boxes covering t6..48,w0..4,u−1..1, bisects only into two exact children, and accepts a box only after the strict upper comparison. Both radial variables remain their full independent intervals throughout. Every denominator is checked positive; square-root lower truncation at0 only enlarges an enclosure. The final cover has263,222 accepted boxes,525,772 tested boxes and no unprocessed box. Negative w is complex conjugation together with theta→−theta, already covered by u's two signs. Thus this is a finite continuum certificate, not a sampled optimization. The subdivision heuristic has no logical role beyond choosing which exact cover to test next.

\subsection{4. Exact conditional measure and the positive complement}\label{exact-conditional-measure-and-the-positive-complement}

The two-coordinate marginal of uniform spherical measure is \[
 {(a-1)(a-2)\over\pi²R_y^4}
       (1-|z_P|²/R_y²)^{a-3}\,d²z_1d²z_2.
\] The complex denominator change contributes s\^{}-2. Its polar radial change at fixed y is D dD=y dy'/(2y'\^{}2), in absolute value. Hence the restricted transition density is \[
 {(a-1)(a-2)y\over2\pi²R_y^{2a-2}s²y'^2}
 \int d\theta\int_{|\xi|²<S}(S-|\xi|²)^{a-3}
                    {\bf1}_{\chi>\chi_0}\,d²\xi.
 \tag{HCX13}
\] This derives the normalization from the original sphere marginal; there is no additional factor2. Set q=L sin(theta/2). Then m=S\emph{−q²−\textbar xi\textbar² and the phase density is 2/{[}L sqrt(1−q²/L²){]}. Feasibility gives 0\textless S}/L²\textless1/2, because its difference from1/2 is −(D²+1+ys²)/(4Dc). Its extra weight lies between1 and sqrt2. Radial integration on the three-real-dimensional ball (q,Re xi,Im xi) now proves the conditional fraction bound \[
 \varepsilon_*=\sqrt2\Pr\{X\ge1/20\},\qquad
                  X\sim\operatorname{Beta}(3/2,166).
 \tag{HCX14}
\] Concavity of sqrt(1−z), integrated on0\textless=z\textless=19/20, gives \[
 \varepsilon_*\le {353\sqrt{10}\over3340\,\vartheta}
                       (19/20)^{166}
                  <2^{-10}-2^{-24},\quad
 \vartheta={4^{167}\over334\binom{334}{167}}.
 \tag{HCX15}
\] The final inequality is squared between positive rationals and checked exactly.

The same event is symmetric under interchanging the two original points. The inverse normal form has the same t and opposite w. Thus the identical conditional fraction bounds both detailed-balance directions, and their geometric mean retains that fraction once. For every t\textgreater0,w in R the fraction is bounded by \[
 \varepsilon_*+e^{9(6-t)}+e^{(5/4)(t-48)}+(w/4)^{16}.
 \tag{HCX16}
\] Inside the rectangle this follows from HCX14; outside at least one complement term is at least1. This is an inequality for the conditional kernel, not for each individual high-coherence point. It multiplies the SAME genuine kernel that produced DO125 CM6--15; no independent coefficient estimates are multiplied.

\subsection{5. Physical central integration, transport and a genuine positive measure}\label{physical-central-integration-transport-and-a-genuine-positive-measure}

Let p=333/4, gamma=29/10, lambda=29/20, b(t)=1+t/4, A(t)=1+t/2 and R\emph{=3514/29. The original DO125 joint kernel is proportional to \[
 t^{-1/2}b(t)^{-167}\left(1+{w²\over4b(t)²}\right)^{-p}
                          c^{-1/2}e^{C_\gamma-\gamma c},
 \quad c=\sqrt{A(t)²+w²}.
 \tag{HCX17}
\] All coarea, selected-variance, spread, tangent and physical-central constants are unchanged. The support-plane inequality used there is c\textgreater=A sqrt(1+u²)+u²/sqrt(1+u²), with w=2bu. Also c\^{}-1/2\textless=A\^{}-1/2(1+u²)\^{}-1/4. Therefore insertion of(w/W)\^{}(2r) tilts the exact central Beta law to Beta(r+1/2,83−r), for0\textless=r\textless83. Its total shape stays167/2. The same conditional Gamma Jensen inequality yields, relative to the unchanged G1 D} normalization, \[
 C_r\Phi_r(t),\quad
 \Phi_r(t)=t^{-1/2}b(t)^{-(166-2r)}A(t)^{-1/2}
    (1+t/R_*)^{-(r+1/2)}e^{-\lambda t},
\] \[
 C_r=(4/W²)^r{\Gamma(r+1/2)\Gamma(83-r)\over\Gamma(1/2)\Gamma83}
                    (1670/1757)^r.
 \tag{HCX18}
\] The exact central normalization is G1=2·4\^{}83·82!83!/166!, with the same D*=(1480/1653)\^{}(333/2)e\^{}(1557/80)sqrt(1670/1757). The factors U\^{}-f, sqrt(D\_F)\^{}-1, N\_FJ0\^{}-1 and x\^{}-1T\^{}54 have not changed.

There is no logarithmic-derivative argument at the zero of w\^{}16. The original covering radius gives each physical coordinate displacement at most750 sqrt(f) U z\_c; the accepted uniform inequality U z\_c\textless2\^{}-83 and sqrt(f)\textless13 make this less than2\^{}-69. Original near-point transport gives \textbar w\_p−w\textbar\textless=6H r\_p with r\_p\textless2\^{}-352. The combined displacement is below delta=2\^{}-68. For e=2\^{}-20, convexity gives \[
 (|w|+\delta)^{16}\le(1+e)^{15}|w|^{16}
                         +(1+1/e)^{15}\delta^{16}.
 \tag{HCX19}
\] After division by4\^{}16 the extra constant is below2\^{}-800 and is absorbed by the strict margin in HCX15. The first coefficient is retained exactly in the new moment coefficient. This includes dyadic and affine-coset shifts; it does not center the lattice or discard a nonzero coset.

Both radial transport directions follow directly from N²=−t omega omega*J with omega=(u,1), \textbar u\textbar=1: t\_p/t=\textbar\textless omega,g\_p e\_n\textgreater\_J\textbar², whose square root lies between cosh(r\_p)−sinh(r\_p) and cosh(r\_p)+sinh(r\_p). Thus t\_p/t lies in{[}e\^{}-2r\_p,e\^{}2r\_p{]}. The new rates9 and5/4 cost a finite extra factor: t\textless=R and all displayed logarithmic powers are bounded explicitly. A conservative total natural-log bound f·10000(H+1)·2\^{}-352 is below1/2048. The original baseline near-point and central Lipschitz inflations remain exactly those in GC7. We pay an additional complete factor2 in the final count for these new scalar transports, beyond the old three1/1024-bit allowances. HCX19's multiplicative factor is additionally retained in C8, so no gain is credited by suppressing it.

Define the global positive profile \[
 \Psi(t)=2^{-10}\Phi_0(t)
       +e^{54}\Phi_0(t)e^{-9t}
       +e^{-60}\Phi_0(t)e^{5t/4}
       +(1+2^{-20})^{15}C_8\Phi_8(t),\quad W=4.
 \tag{HCX20}
\] Every shape in its normalized Gamma factors is positive, including166−16=150 and8+1/2; every exponential shift is positive, with the smallest1/5. Thus Psi is the Laplace transform of an explicit positive measure nu, the sum of four positive convolutions. Its total mass is infinite because of the unnormalized half-Gamma factor t\^{}-1/2. Positive inverse moments are nevertheless finite: \[
 M_j=\int c^{-j}d\nu(c)={1\over\Gamma(j)}
                  \int_0^\infty t^{j-1}\Psi(t)\,dt,
                 \qquad j=1,\ldots,163.
 \tag{HCX21}
\] Tonelli applies to nonnegative integrands. No M0, signed inverse transform or compact-support Laplace representation is used.

\subsection{6. Exact complete moments and the original zero-degree repair}\label{exact-complete-moments-and-the-original-zero-degree-repair}

For j148,160,163 the executable integrates the low, high and even-central components separately. For each positive integrand it uses8064 directed trapezoid intervals on{[}1,64{]}, exact second-derivative error, and both{[}0,1{]} and{[}64,infinity) tails. The general integrand, before its positive coefficient, is \[
 {t^{j-3/2}\over\Gamma(j)}(1+t/4)^{-(166-2r)}
       (1+t/2)^{-1/2}(1+t/R_*)^{-(r+1/2)}e^{-\lambda't}.
 \tag{HCX22}
\] Its full logarithmic derivatives are used. At infinity the remaining power after bounding the three denominators is beta=j−337/2+r. When beta\textgreater0, t\^{}beta\textless=64\^{}beta exp(beta(t−64)/64); hence the remaining exponential rate is lambda'−max(beta,0)/64\textgreater0 in every new enclosure. This explicitly handles the r8,j163 positive-power tail. The high-radial component may have a deliberately loose but finite tail; it is included with its e\^{}-60 coefficient, not dropped as small.

The baseline moment anchors148/162/163 are exact accepted DO126 intervals. Baseline160 uses finite-endpoint Holder and a directed rational14th root. The nine new integrals give72,576 intervals. Their coefficients e\^{}54 and e\^{}-60 have directed rational bounds from alternating exponential series and outward squaring. The full four-component sum is performed before a field power is taken.

A safe first moment is M1\textless2\^{}79. Indeed each probability Gamma factor is at most1, C8(1+2\^{}-20)\^{}15\textless1, and integrating t\^{}-1/2 exp(−lambda't) gives sqrt(pi/lambda'). In particular pi\textless16/5 makes the smallest-rate1/5 integral less than4. The low-t complement term is less than e\^{}54, the other three total less than7, and e\^{}54\textless2\^{}78 follows from the proved log2(e)\textless13/9. All intermediate j use Holder between adjacent finite anchors1,148,160,163 on the same possibly infinite positive measure. Neither probability normalization nor a zeroth moment is assumed.

The sole zero-degree polynomial term is repaired BEFORE any row enlargement, using the original k nonzero S entries and their product-norm floor, exactly K116 RE7--14 and DO126 IQ22--28. Reserve theta=2\^{}-20 of the damping, then the repaired contribution is M1\^{}f(q/(theta R k))\^{}f.~All positive degrees use(q/((1−theta)R))\^{}(jf)M\_j\^{}f.~The floor is never applied to newly admitted zero coordinates. The enlarged Gaussian count keeps the original projective selfdual rank3 quotient and its1008-real-dimensional theta; no scalar estimate is substituted for it.

\subsection{7. Every arithmetic factor, pattern and continuum endpoint}\label{every-arithmetic-factor-pattern-and-continuum-endpoint}

Write q=R/(Uz²), z\^{}(2f)=(N\_EJ)xT\^{}-54, and \[
 (q_*)^f={R^fT^{54}\over U^fQ^{334}x}.
 \tag{HCX23}
\] DO127's literal Green identity bounds the complete G product by2 Q\^{}(-335·167)(N\_EJ)\^{}167 N\_FJ0. It cancels the ORIGINAL central N\_FJ0 denominator. For every effective degree j, all four local states sum to \[
 B_j=\prod_v(1+q_v^{j-167})².
 \tag{HCX24}
\] Here1\textless=j\textless=163, so B\_j\textless2\^{}(1/1024); every rank-one orientation is counted. The true k-dependent polynomial has2(k+1) terms, \[
 \sum_{l=0}^k\binom kl K^{k-l}
       (q/((1-\theta)R))^{(l+3)f}M_{l+3}^f
 +K_3\sum_{l=1}^k\binom kl K^{k-l}
       (q/((1-\theta)R))^{lf}M_l^f
 +K_3K^kM_1^f(q/(\theta Rk))^f.
 \tag{HCX25}
\] K is the checked degree15,336-real-dimensional scalar theta constant; K3 is the checked1008-real-dimensional quotient theta constant with its correct minimum numerator f.~Both the full transverse volume and its theta correction remain. There is no extraneous4\^{}f mixture count: M\_j is already the sum of every positive component at that embedding, and M\_j\^{}f retains their complete product expansion.

The prefactor is exactly DO127 GC7, with ONE additional factor2 from HCX19--20 transport. It therefore retains binomial(168,27), the164 pivot, binomial(164,k), the original Dvar\^{}84, vartheta\^{}168, coarea factor, G1D* to the168th power, sqrt(D\_F)\^{}-1, physical U\^{}-168, (R/U)\^{}(167f), T\^{}8904, Q\^{}(-335·167), x\^{}-168q\^{}(-167f), pi\_+\^{}((k+3)f) and D0\^{}(-(k+3)). In particular \[
                       8904=54\cdot168-168.
 \tag{HCX26}
\] No source fibre, finite label, Haar/class factor or U divisor is changed.

At fixed x every term has q power(j−167)f\textless0. At q=q*, its complete x power is−1−j\textless0, also for the repaired effective degree1. The whole continuum is therefore bounded at x=2\^{}2111 and the all-identity q cap. Original eta uniqueness per27-place support holds throughout x\textless2\^{}3328 by K113; there is no extra denominator-ideal sum or continuum multiplicity. Every affine central fibre and finite pattern is retained before enlargement and positive summation.

The exact certificate recomputes all322 k160 terms with the new finite moments and actual DO127 predecessor. The strict q bracket is251124964\textless q*\textless251124965. The maximal directed logarithmic upper is−40383083/65536, at l160 in the main transverse term, degree163. Every term satisfies log2(term)\textless−610, including the additional full factor2. It also verifies the original q bracket by168th powers. The positive errors in HCX5 are each below2\^{}-610. Since324\textless2\^{}(610−600), the complete union is strictly below2\^{}-600. All earlier errors are preserved in order and600 is appended once. Numerical endpoint details and the complete322-term array belong to the final exact certificate; no floating diagnostic supplies a theorem premise.

\subsection{8. Source use, verification and remaining boundary}\label{source-use-verification-and-remaining-boundary}

The actual selected action/evaluation normalization and ordinary convergence are K043 PB1--8, DO080 RT1--8/RT22 and K122 HC1--18 plus its mandatory nonzero-S-cross addendum. K120 NB1--16 and K118 BT15--22 provide the exact narrow-band/full-arithmetic errors. The original unipotent normal form and physical w are DO098 CF4--7; all current all-prime ideals and quotient metrics are DO126 IQ1--36. K121 KP1--34 and DO125 CM3--21 supply the same-kernel coarea/tangent/central positive-mixture normalization. DO127 GC1--16 supplies the exact Green/J0 cancellation, all-four-state pattern, complete prefactor and endpoint directions. No earlier numerical removal charge is reused as a factor.

All new geometric, probability, positivity, moment and counting steps are proved here. The finite diagnostic models do not prove a missing source constructor or replace the continuum argument. The actual selected arithmetic image/nonzero, period, signedPP3, original source admission, specific032, evaluatedDP44 and RH gates remain separate. The Alpöge archive is used only at its statement/interface boundaries; no theorem-validity audit or Lean build occurs. Carried FGA/SGA context remains precisely that recorded in the accepted dependencies; no new whole-text reading or algebraic-geometric construction is claimed.

The initial sampled cap and moment searches, two interrupted subdivision heuristics, and other diagnostic files remain dated pending history. They are not final certificates. The final author packet records the whole saved-proof/script read, all source hashes and full exact replay. Independent whole review remains required before canonical promotion.

\subsubsection{Primary literature carried through the accepted interfaces}\label{primary-literature-carried-through-the-accepted-interfaces}

Nakahama, R. (2016). Norm computation and analytic continuation of vector valued holomorphic discrete series representations. \emph{Journal of Lie Theory, 26}(4),927--990. Equations5.5/5.7 and Theorem5.1, printed959--961, with the actual K043 convention as recorded in K122. \url{https://jolt.centre-mersenne.org/item/10.5802/jolt.919.pdf} . The selected realization and ordinary convergence are carried from that checked interface; no new full-paper reading is claimed here.

Gorodetsky, O., \& Hazan, Z. (2017). \emph{On certain degenerate Whittaker models for cuspidal representations of GL\_(kn)(F\_q).} arXiv:1707.07308v2. Theorem2.1, equation2.1 and Corollary2.2, printed5. \url{https://arxiv.org/pdf/1707.07308} . The rank-dependent normalized Green formula and its exact all-pattern consequence are used through the separately checked DO127, whose primary page was freshly read in this agent's independent review.

\subsection{Canonical algebraic-geometric context}\label{canonical-algebraic-geometric-context-91}

The new estimate is on the original real-place selected theta coefficient and its actual homogeneous Hermitian bundle. Simultaneous compact conjugation and conditional sphere disintegration preserve the original polarized projective OE modules, selected CM phase, finite projector, cubical/Poincare line, Haar/class normalization and prescribed arithmetic arrow. The same projective selfdual isotropic quotient over Spec(OE) and original affine central ideal remain in the count. The positive real Laplace mixture is an analytic majorant, not a new algebraic theta correspondence. Primary normalization: Nakahama(2016), equations5.5/5.7 and Theorem5.1, printed959-961; exact finite Green specialization: Gorodetsky and Hazan(2017), Theorem2.1/equation2.1 and Corollary2.2, printed5, through the accepted inputs. Inherited context only: SGA3 I4.7.3/4.7.3.1; SGA6 VI Theorem3.3; FGA232 Theorem3.1 pp148--149; FGA236 Corollaries1.7/2.3 pp229/232; SGA7 I VII2.7/2.9 pp169--174 and VIII7.4(b) pp310--311; cote4 scan003 sections8.1.1--8.1.2 and scans005--006 sections8.4.1--8.4.3; cote54 scans005--007. No new Dieudonne, lambda-ring/K-theory, l-adic-weight or CM-motive comparison. DP\_MULT Q3e/DP44 and THETA3 G1/A7 gain the strictly smaller conditional finite gate; CX\_GRH CX-ALG-024 and GCT\_PLETH T-052 retain their arithmetic/global gates.

\clearpage

\section{AM-LIFT-107 --- Adjacent signatures and the same-field CM-type generator law}\label{am-lift-107-adjacent-signatures-and-the-same-field-cm-type-generator-law}

\textbf{Canonical checked integration --- 2026-11-23T13:50:09.162185+00:00.} Complete saved-proof and executable reviews and separate exact diagnostic replays passed. No numerical error charge is introduced. Frozen author history is preserved. AM-LIFT-107 AG1-19 proves that, for the original Par/PinDat and coefficient field F, every signature KPsiIn is equivalent to the two actual stages j1/j2 and to one normalized rank-one CM-type generator law. Explicit adjacent Bezout extraction preserves F and every embedded absolute norm, with no relation imposed between independently chosen witness phases. A finite normalized class-generator table is necessary and sufficient; a fixed principal-power root table is sufficient only, retaining its unit-phase distinction. Accepted R102 then supplies the unchanged beta implication. Actual table existence, original source/Host/Sizes and selected arithmetic nonzero/content, period,032 and evaluatedDP44/RH remain open. No new charge; current138/31 arrays and27/11 floors unchanged. Frozen137/31 author context is historical.

Author: \nolinkurl{/root/lift_accounting}. \textbf{Research proof pending independent whole review.} The accompanying author-verification and pending manifest record the exact frozen bytes and completed checks. Accepted Alpöge theorems are cited and applied as supplied inputs; this is not a Lean build or an audit of those accepted theorems.

\subsection{1. The original family and the reduced arithmetic obligation}\label{the-original-family-and-the-reduced-arithmetic-obligation}

Fix an actual source \nolinkurl{D} and \nolinkurl{Q : Par D} with \nolinkurl{Q.PinDat}. Write \(E=D.\mathrm{host}.E\), \(H=Q.H\), \(l=Q.l\), \(c=Q.c\), and \(F=c.F\). This \(F\) always means the \textbf{original coefficient field}, not the maximal real subfield of \(E\). The source \nolinkurl{Coeff.conj} places every embedding of \(E\) into this same subfield of \(\mathbf C\); \nolinkurl{embF} and \nolinkurl{embK} are the literal resulting maps.

Put \(k=D.R.k\), \(m=D.R.m\), and, for \(1\le j\le m\),

\[
 p_j=k+j-1,\qquad q_j=m-j,\qquad
 r=k+m-1,\qquad \delta_j=p_j-q_j.
 \tag{AG1}
\]

These are the actual signatures of \(Q.Dw(j)\) by \nolinkurl{IsDatW/PinDat}. The \nolinkurl{Ranks} record makes \(k,m\) even, \(m\ge2\), and \(k\ge m+3\). Consequently \(k-m\ge4\), \(r\) is odd, and

\[
 u=(k-m)/2\in\mathbf N,\quad u\ge2,\qquad
 h=(r-1)/2\in\mathbf N,\qquad
 \delta_j=2u+2j-1>0.
 \tag{AG2}
\]

All natural-number subtractions in AG1 are exact on this range. In particular stages \(j=1,2\) both belong to the original \nolinkurl{InR}.

Let

\[
 \Sigma^+=\{\tau:E\hookrightarrow\mathbf C:
                    \tau(H.rt)=H.rt0\},\qquad
 \Sigma^-=\operatorname{Hom}(E,\mathbf C)\setminus\Sigma^+.
 \tag{AG3}
\]

For each nonzero integral ideal \(I\) of \(\mathcal O_E\), define ideals of \(\mathcal O_F\) by

\[
 \Phi(I)=\prod_{\tau\in\Sigma^+}\tau(I)\mathcal O_F,\qquad
 \Psi(I)=\prod_{\tau\in\Sigma^-}\tau(I)\mathcal O_F,\qquad
 n_I=N_{E/\mathbf Q}I.
 \tag{AG4}
\]

Here \(\tau(I)\mathcal O_F\) denotes extension through the actual \nolinkurl{embF/embK}. It is independent of an ideal generator.

Define the \textbf{rank-one CM-type generator law} \(\mathcal C(H,c)\):

\[
 \forall\,0\ne I\subseteq\mathcal O_E,\quad
 \exists\chi_I\in F^\times,\quad
 (\chi_I)=\Phi(I),\qquad
 |\sigma(\chi_I)|^2=n_I
       \quad\text{for every }\sigma:F\hookrightarrow\mathbf C.
 \tag{AG5}
\]

This is an arithmetic predicate on the existing field and CM type. It does not require constructing a rank-one moduli datum.

\textbf{Theorem A.} With the original data and \nolinkurl{PinDat} as above, the following are equivalent:

\begin{enumerate}
\def\labelenumi{\arabic{enumi}.}
\tightlist
\item
  the entire original \nolinkurl{LawsP.kPsi} obligation \(\forall j\in[1,m],\ \mathrm{KPsiIn}(Q.Dw(j),c)\);
\item
  only the two original obligations at \(j=1\) and \(j=2\);
\item
  the single same-field law \(\mathcal C(H,c)\).
\end{enumerate}

The proof supplies explicit formulas in the same \(F\). It does not alter the source's selected \nolinkurl{betaH} values.

\subsection{2. Exact type, ideal and principal-element identities}\label{exact-type-ideal-and-principal-element-identities}

The host has \(H.rt^2=-d\), with \(d\) a positive prime, and \(H.rt0^2=-d\). Every \(\tau(H.rt)\) is therefore one of the two distinct roots \(\pm H.rt0\). Complex conjugation exchanges them. Since \(E\) is CM, \(\Sigma^+\) contains exactly one embedding from each conjugate pair. The literal definition \nolinkurl{Datum.mult} is \(p_j\) on \(\Sigma^+\) and \(q_j\) on \(\Sigma^-\).

In the group of nonzero fractional ideals of \(\mathcal O_F\),

\[
             \Phi(I)\Psi(I)=(n_I).
 \tag{AG6}
\]

For completeness, this use of the usual ideal norm is independent of any unramified or split-place convention. Choose \(a>0\) with \(I^a=(x)\) using finiteness of the ideal class group of \(E\). Then

\[
 [\Phi(I)\Psi(I)]^a
   =\left(\prod_{\tau:E\hookrightarrow\mathbf C}\tau(x)\right)
   =(N_{E/\mathbf Q}x)=(n_I^a).
 \tag{AG7}
\]

The middle element is a rational number, so it has the same image in the given \(F\). The last ideal equality follows from the ordinary norm of a principal ideal. The fractional-ideal group of a number ring is torsion free, by unique prime-ideal factorization; cancel the positive exponent \(a\). This proves AG6 at all primes, with all ramification and residue degrees retained. The standard facts used here are the ideal-group theorem, principal-ideal norm compatibility and class-group finiteness (Sutherland, 2021a, Theorem 3.11; 2021c, Propositions 6.6--6.7; 2021d, Corollary 14.22).

For \(x\in E^\times\) write

\[
             N_\Phi(x)=\prod_{\tau\in\Sigma^+}\operatorname{embF}(c,\tau)(x)
                   \in F^\times.
 \tag{AG8}
\]

Then \(\Phi((x))=(N_\Phi(x))\), and for every embedding \(\sigma\) of \(F\),

\[
             |\sigma N_\Phi(x)|^2=|N_{E/\mathbf Q}x|.
 \tag{AG9}
\]

Indeed the embeddings \(\sigma\circ\operatorname{embF}(c,\tau)\), as \(\tau\) varies through all embeddings of \(E\), are distinct and therefore permute that complete set. Composition preserves the conjugate pairing: an embedding of a CM field carries its canonical involution to complex conjugation. The image of \(\Sigma^+\) still contains one element of every pair. Multiplying a half-product by its complex conjugate gives the product over all embeddings, namely the positive rational norm of \(x\). This argument does not assume that \(F\) is itself CM, normal, or Galois. The field-norm product is the usual separable-embedding formula (Sutherland, 2021b, Theorem 4.50).

The weighted ideal occurring in the literal \nolinkurl{KPsiIn} at stage \(j\) is consequently

\[
       B_j(I)=\Phi(I)^{p_j}\Psi(I)^{q_j}
              =\Phi(I)^{\delta_j}(n_I)^{q_j}.
 \tag{AG10}
\]

The definition writes this as an \(\mathcal O_F\)-submodule of \(F\) generated by the embedded ideal elements. Such a submodule is exactly the corresponding extended ideal viewed inside \(F\). The accepted \nolinkurl{LatticeBetaCompare.isBe_iff/kPsiIn_iff} interfaces identify these two literal presentations. We use that accepted identification, not an independently assumed free ideal basis.

\subsection{3. A rank-one witness supplies every original stage}\label{a-rank-one-witness-supplies-every-original-stage}

Assume AG5. For any fixed nonzero \(I\), choose its witness \(\chi_I\). For every original stage put

\[
              \beta_j(I)=\chi_I^{\delta_j}n_I^{q_j}\in F^\times.
 \tag{AG11}
\]

Its principal ideal is AG10. Since \(B_j(I)\) is integral, the element \(\beta_j(I)\) is integral as well; indeed \(\beta_j(I)\mathcal O_F\subseteq\mathcal O_F\) implies \(\beta_j(I)\in\mathcal O_F\). The literal predicate only requires an element of \(F\), but its integral ideal equality is thus exact.

At every embedding of the original coefficient field,

\[
 |\sigma\beta_j(I)|^2
   =n_I^{\delta_j}n_I^{2q_j}
   =n_I^{p_j+q_j}=n_I^r.
 \tag{AG12}
\]

Thus AG11 supplies the actual \nolinkurl{KPsiIn} witness at every \(j\), including the terminal definite stage \(j=m\), where \(q_m=0\) and \(\beta_m=\chi_I^r\). No change of \(H,l,c,Q.Dw\), finite phases, moduli, class sums or source pin occurs.

The implication from every stage to the two stages is immediate because \(m\ge2\). The converse reduction is proved next.

\subsection{4. Two adjacent original stages recover the rank-one law}\label{two-adjacent-original-stages-recover-the-rank-one-law}

Assume the two actual \nolinkurl{KPsiIn} predicates at \(j=1,2\). For each fixed nonzero ideal \(I\), choose any of their respective witnesses \(\beta_1,\beta_2\in F^\times\). They need not be compatible choices as \(I\) varies, and no phase relation between them is assumed. Define

\[
          \chi_I=\beta_1^{\,u+1}\beta_2^{-u}n_I^{-h},
          \qquad h=(r-1)/2.
 \tag{AG13}
\]

All inverses exist. The actual adjacent signatures give

\[
 \begin{aligned}
 (u+1)p_1-up_2&=k-u=(r+1)/2=h+1,\\
 (u+1)q_1-uq_2&=m-1+u=(r-1)/2=h.
 \end{aligned}
 \tag{AG14}
\]

Taking fractional ideals, using AG6, and then cancelling yields

\[
 (\chi_I)=\Phi(I)^{h+1}\Psi(I)^h(n_I)^{-h}=\Phi(I).
 \tag{AG15}
\]

In particular the quotient in AG13 is actually integral. At every \(\sigma:F\hookrightarrow\mathbf C\), the two KPsiIn norm equalities give

\[
 |\sigma\chi_I|^2
   =(n_I^r)^{u+1}(n_I^r)^{-u}n_I^{-2h}
   =n_I^{r-2h}=n_I.
 \tag{AG16}
\]

This is precisely AG5 and completes Theorem A. Equivalently, \((u+1)\delta_1-u\delta_2=1\); the positive odd differences at the two adjacent stages remove the possible power ambiguity. The argument does not assert that a single terminal predicate permits the same extraction.

Theorem A is about existence of the witnesses in the original predicates. In particular AG11 does not redefine \nolinkurl{betaH}. Once the predicate is supplied, the original selected \nolinkurl{betaH} retains its own defining specification and any original phase. No formula for that selected phase follows from AG11 or AG13.

\subsection{5. A finite table in the same field}\label{a-finite-table-in-the-same-field}

Fix any finite system of nonzero integral ideals \(I_1,\ldots,I_s\) whose classes generate \(\operatorname{Cl}(\mathcal O_E)\). The following finite table is equivalent to AG5:

\[
 \chi_i\in F^\times,\qquad
 (\chi_i)=\Phi(I_i),\qquad
 |\sigma\chi_i|^2=NI_i\quad\text{for every }i,\sigma.
 \tag{AG17}
\]

Necessity follows by applying AG5 to these ideals. For sufficiency, for an arbitrary nonzero integral \(I\) write \(I=(x)\prod_i I_i^{a_i}\) as fractional ideals, with \(x\in E^\times\) and integer exponents \(a_i\). Define

\[
                \chi_I=N_\Phi(x)\prod_i\chi_i^{a_i}.
 \tag{AG18}
\]

Its ideal is \(\Phi(I)\) by multiplicativity. By AG9 and the table, the square of every embedded absolute value is \(|N_{E/\mathbf Q}x|\prod_i(NI_i)^{a_i}=NI\). The integral ideal equality again implies \(\chi_I\in\mathcal O_F\). This establishes the equivalence without class-dependent field extensions.

One may choose a nontrivial invariant-factor system with orders \(a_1',\ldots,a_s'\ge2\), so every class has unique coordinates \(0\le a_i<a_i'\) and \(\prod_i a_i'=h_E\). Then \(s\le\log_2h_E\). For class number one the table is empty, and AG8 itself supplies AG5 for every ideal. These are exact finite-table sizes, not a numerical evaluation of the actual class group.

For comparison, a stronger prescribed table

\[
 I_i^{a_i'}=(\gamma_i),\qquad
 \psi_i^{a_i'}=N_\Phi(\gamma_i),\qquad \psi_i\in F
 \tag{AG19}
\]

is sufficient for AG17. The principal-ideal equality follows by raising to the positive power \(a_i'\) and using torsion freeness; the all-embedding norm equality follows by taking the unique positive \(a_i'\)-th root of the corresponding positive real equality. No choice of a complex root outside \(F\) is permitted in AG19.

We do \textbf{not} claim that AG17 implies AG19 for arbitrary already fixed \(\gamma_i\). Its discrepancy is the unit \(\chi_i^{a_i'}/N_\Phi(\gamma_i)\), whose absolute value is one at every embedding. It is a root of unity by the usual Kronecker theorem, but it need not be one or have an \(a_i'\)-th root in \(F\). Increasing the principal exponent by its order can remove that discrepancy; this changes the fixed exponent/table and is not used in Theorem A. The accepted archive \nolinkurl{ReflexRoots} equivalence has existential principal exponents and generators, and is consistent with this precise distinction. The unit-phase assertion uses the finite kernel of the all-place divisor map (Sutherland, 2021e, Corollary 15.8): the unit has absolute value one at finite places as well as at infinity.

\subsection{6. Exact original-law and cost consequences}\label{exact-original-law-and-cost-consequences}

Theorem A fills the original \nolinkurl{LawsP.kPsi} field from either the two adjacent obligations or the single finite table AG17. It does not assume or construct a completed \nolinkurl{LawsP} record. The actual supplied \nolinkurl{PinDat} alone supplies the signatures used in the reduction.

The accepted universal \nolinkurl{GlobalTheory.BetaEx.betaExists} then supplies the original beta implication, as recorded by canonical AM-R-102, \nolinkurl{agents/cited_beta_law_and_pinned_coefficient_supplier.md} (SHA-256 \nolinkurl{c9873b1a7d5b40e8718cbae3cc6b1d2c35cae7f210124102131a06653b71da6b}). A same-field witness discharges that implication for every required lattice; the originally chosen \nolinkurl{betaH} is left unchanged. The other source laws, actual selected theta nonzero/content, prescribed global period,032 value and DP44 inequality remain separate.

\begin{longtable}[]{@{}
  >{\raggedright\arraybackslash}p{(\columnwidth - 2\tabcolsep) * \real{0.5000}}
  >{\raggedright\arraybackslash}p{(\columnwidth - 2\tabcolsep) * \real{0.5000}}@{}}
\toprule\noalign{}
\begin{minipage}[b]{\linewidth}\raggedright
Original data or cost
\end{minipage} & \begin{minipage}[b]{\linewidth}\raggedright
Effect
\end{minipage} \\
\midrule\noalign{}
\endhead
\bottomrule\noalign{}
\endlastfoot
\(E,H,l,F\), \nolinkurl{what₂/eO/PinField} & Literally unchanged; no extension degree or discriminant factor. \\
All actual \nolinkurl{Dw(j)} and \nolinkurl{PinDat} & Unchanged; their exact signatures drive the equivalence. \\
\nolinkurl{betaH}, \nolinkurl{DeltaDat}, operative connections and finite phases & No replacement or equality of independently chosen witnesses. \\
Class sums, lattices, periods, metrics and DP44 terms & Unchanged; no cancellation or vanishing conclusion. \\
Number of scalar generator tests & One CM-type law, or one finite class-generator table, replaces the family of signature laws. \\
Actual table existence, field pin and Host/Sizes & Still required; the theorem supplies no numerical or algebraic table in the prescribed field. \\
Error arrays/support floors & No new charge; current137/31 and27/11 are unchanged. \\
\end{longtable}

The extension theorem \nolinkurl{exists_coeff_typeNorms_principal} is a separate accepted supplier after enlarging a coefficient record. It is not used to change \(F\) here. The exact \nolinkurl{Par.F_eq} pin boundary remains in force.

The interpretation on \(\operatorname{Spec}\mathcal O_F\) is a normalized trivialization of the invertible ideal \(\Phi(I)\), including the specified archimedean norms. The finite table concerns the image of the ideal class group with this metric condition. It does not replace the condition by ordinary principalization, or assert a new motivic, weight, positivity, or global nonvanishing statement.

\subsection{7. Exact sources, references and verification scope}\label{exact-sources-references-and-verification-scope}

The literal interfaces are:

\begin{itemize}
\tightlist
\item
  \nolinkurl{Formalism/LocalRingCardinality/Params.lean}6--22,32,46: even k,m, m\textgreater=2, k\textgreater=m+3 and actual signatures.
\item
  \nolinkurl{Formalism/LocalToGlobal/ProjectiveLattices/PartitionExtendByZeroAntitone.lean}12,24--31,47: InR, exact IsDatW and PinDat.
\item
  \nolinkurl{ReciprocityTheory/ReductiveGroups/ArithmeticGeometry/BirationalGeometry/ToroidalCompactification/TwoWeightEqOneAddRoot.lean}14--49: actual Host and Datum.
\item
  Recovered \nolinkurl{.../ToroidalCompactification/FlagBasis.lean}76: literal Datum.mult; whole119-line file read.
\item
  \nolinkurl{PushInvertible.lean}40--46: the actual finite coefficient field, every conjugate embedding and existing root/ground-field inclusions.
\item
  \nolinkurl{Formalism/LocalToGlobal/ProjectiveLattices/EmbKCoe.lean}25--33: original field and embedding maps.
\item
  \nolinkurl{ThShort.lean}74--85,277--285 and \nolinkurl{HasDataCondition.lean}55: exact KPsiIn and its original all-j consumer.
\item
  Accepted \nolinkurl{MultAddMultConjugate.lean}262--272,460--474 and \nolinkurl{ReflexNormIndexIdeal.lean}51--72: actual ideal/submodule presentation and existing root interfaces.
\item
  Accepted \nolinkurl{ExistsIndexIdeal.lean}322--324 and \nolinkurl{ExistsCoeffTypeNormsPrincipal.lean}23--24: respectively the universal beta implication and the separate coefficient-extension supplier.
\item
  \nolinkurl{MapRingEquivTransApply.lean}52--69: the actual pinned-field equality.
\end{itemize}

The companion source JSON records full local paths, exact hashes, actual line totals and genuinely completed reading scopes. The companion executable tests signature exponents, fractional-ideal valuation identities, exact Gaussian-rational phase models, finite class-coordinate reconstruction and the fixed-root boundary. It also checks the recorded literal interfaces. These finite diagnostics are not a proof of an arbitrary number-field assertion; that proof is AG1--19 above. No original class group, scalar table or selected arithmetic value is evaluated by the executable.

\subsubsection{References}\label{references-9}

Sutherland, A. V. (2021a, September 15). \emph{Properties of Dedekind domains} (18.785 Number Theory I, Lecture 3, Theorem 3.11 and Corollary 3.13, pp.~3--4). Massachusetts Institute of Technology. \url{https://ocw.mit.edu/courses/18-785-number-theory-i-fall-2021/mit18_785f21_lec3.pdf}

Sutherland, A. V. (2021b, September 20). \emph{Étale algebras, norm and trace} (18.785 Number Theory I, Lecture 4, Lemma 4.49 and Theorem 4.50, p.~12). Massachusetts Institute of Technology. \url{https://ocw.mit.edu/courses/18-785-number-theory-i-fall-2021/mit18_785f21_lec4.pdf}

Sutherland, A. V. (2021c, September 27). \emph{Ideal norms and the Dedekind--Kummer theorem} (18.785 Number Theory I, Lecture 6, Definition 6.5 and Propositions 6.6--6.7, p.~4). Massachusetts Institute of Technology. \url{https://ocw.mit.edu/courses/18-785-number-theory-i-fall-2021/mit18_785f21_lec6.pdf}

Sutherland, A. V. (2021d, October 27). \emph{Geometry of numbers} (18.785 Number Theory I, Lecture 14, Theorems 14.18 and 14.20, Lemma 14.21 and Corollary 14.22, pp.~6--8). Massachusetts Institute of Technology. \url{https://ocw.mit.edu/courses/18-785-number-theory-i-fall-2021/mit18_785f21_lec14.pdf}

Sutherland, A. V. (2021e, November 1). \emph{Dirichlet's unit theorem} (18.785 Number Theory I, Lecture 15, Lemma 15.7 and Corollary 15.8, p.~4). Massachusetts Institute of Technology. \url{https://ocw.mit.edu/courses/18-785-number-theory-i-fall-2021/mit18_785f21_lec15.pdf}

Alpöge source. Accepted supplied statements at repository commit \nolinkurl{7cbf7a0c32fca1badf05931879a37c3f8a2b03b4} and the literal source copies identified in the companion source record. Only their stated interfaces are applied; no validity claim beyond their accepted status is made here.

\subsection{Canonical algebraic-geometric context}\label{canonical-algebraic-geometric-context-92}

Over Spec(O\_F), the two signature-weighted invertible ideals determine the CM half-type ideal by an explicit integral Bezout combination and the rational norm line. The finite table is a normalized trivialization of these ideal classes with all prescribed archimedean norms, not ordinary principalization alone. The original polarized CM projective lattices, Poincare/cubical lines, finite phases and coefficient-field pin remain unchanged. Exact literal locators: Params6-22/32/46; PartitionExtendByZeroAntitone12/24-31/47; Datum.mult in FlagBasis76; Coeff in PushInvertible40-46; EmbKCoe25-33; ThShort277-285; HasDataCondition55; accepted MultAddMultConjugate262-272/460-474; MapRingEquivTransApply52-69. Primary ideal/norm/class/unit locators are Sutherland2021 Lectures3/4/6/14/15 as listed in the proof. Carried exact context: SGA7 I VII2.7/2.9 pp169--174 and VIII7.4(b) pp310--311; FGA232 Theorem3.1 pp148--149 and FGA236 Corollaries1.7/2.3 pp229/232; SGA3 I4.7.3/4.7.3.1; SGA6 VI Theorem3.3; cote4 scan003 sections8.1.1--8.1.2 and scans005--006 sections8.4.1--8.4.3; cote54 scans005--007. These are carried contextual readings, not new whole-text checks. No new integral extension over Z, Dieudonne slope, lambda-ring positivity, K-theory filtration, l-adic weight theorem or CM-motive L-function comparison follows. DP\_MULT Q3e/DP44/DP52/DP54--65 and THETA3 G1/A7 gain this conditional source-rule reduction; CX\_GRH CX-ALG-024 and GCT\_PLETH T-052 retain their global gates.

\clearpage

\section{AM-DO-128 --- actual finite Heisenberg zero mode and positive-definite completion}\label{am-do-128-actual-finite-heisenberg-zero-mode-and-positive-definite-completion}

\textbf{Canonical checked integration --- 2026-11-23T14:03:40.841641+00:00.} Complete saved-proof and executable reviews and separate exact diagnostic replays passed. No numerical error charge is introduced. Frozen author history is preserved. AM-DO-128 HZ1-24 gives the exact signed full-Heisenberg zero mode and additive Fourier coefficients, qualitative paired CRT completion preserving every original T incidence, full original-kernel positive definiteness and a convergent full-group nonnegative sum, and exact physical-frequency covering bounds. Identity repetition, cutoff, coset and actual selected-transform limitations are explicit. This structural theorem alone makes no numerical main-gate reduction or new charge. Current K123138/31 arrays, absolute floors27/11 and support27 k161-164 remain unchanged; frozen137/31 and k160 statements are historical. Original source/Host/Sizes, selected arithmetic image/nonzero, period,032 and DP44/RH remain open.

Author: \nolinkurl{/root/definite_oscillator}, 23 November 2026 UTC.

\textbf{Pending scoped result; no numerical sector removal or error charge.} The actual predecessor is DO127, with137 all-top and31 uniform charges and complete floors27/11. This note preserves the original finite representation, two mixed projectors, actual T pencils, selected Hilbert trace, central character and all normalizations. It identifies an exact finite cancellation, proves the applicable full-group positivity statement, and records precisely why neither statement yet estimates the remaining truncated sum.

\subsection{1. Original setting and the two different completions}\label{original-setting-and-the-two-different-completions}

Let E/F, f=168, Lambda=O\_E\^{}169, J=diag(I\_168,-1), Q=g0\emph{g1 and U=77Q\^{}2-14 be the actual J56 data. The remaining DO127 sector has h=I+N, N\^{}3=0, rank\_E N=2, full positive compression, null line ell=Ew with w=(0\_S,u,eta), eta}bar(eta)=U, and k=160,\ldots,164. Its exact denominator x=N\_E D\_eta satisfies2111\textless log2(x)\textless2165, with the narrower DO127 staircase retained. No existence in these necessary windows is asserted.

For a fixed ell write U\_ell for the rational unipotent radical of its stabilizer. Its Levi quotient is the positive Hermitian space ell\^{}perp/ell of E-dimension167. The saturated original lattice quotient is projective and selfdual by DO126(IQ8-19); no global O\_E basis is assumed. Every matrix considered here belongs to this same actual U\_ell(F). There are two distinct constructions below:

\begin{enumerate}
\def\labelenumi{\arabic{enumi}.}
\tightlist
\item
  An arithmetic congruence subgroup Gamma0, used to complete the finite residue coordinates of an actual affine coset h0 Gamma0 while fixing all actual T incidences.
\item
  An arithmetic lattice Gamma in U\_ell(F), used for a full-group positive-definite sum. It need not preserve each chosen T incidence.
\end{enumerate}

They must not be substituted for one another. Neither completion retains the original H=133 cutoff automatically. (HZ1)

\subsection{2. Exact finite signed Green calculation}\label{exact-finite-signed-green-calculation}

At an unmixed split G place with residue field k of order q, the actual unitary group identifies with GL\_169(k) on the selected CM component; the opposite component is its inverse adjoint. In a Witt chart the full unipotent radical is

\[
 U_m(k)=\left\{
 \begin{pmatrix}1&a&z\\0&I_m&b\\0&0&1\end{pmatrix}:
 a\in k^{1\times m},\ b\in k^{m\times1},\ z\in k\right\},
 \qquad m=167.                                                     \tag{HZ2}
\]

The three possible ranks of u-I are0,1,2. Their counts, also valid for an embedded U\_m with m\textless=167, are

\[
 n_0=1,\quad n_1=(q-1)+2q(q^m-1),\quad n_2=q(q^m-1)^2.
                                                                    \tag{HZ3}
\]

Indeed both nonzero a,b give independent image directions and rank2; exactly one nonzero gives rank1 for every z; with both zero the central coordinate decides rank0 or1. The normalized actual unmixed cuspidal character has the Green values

\[
 A_0=1,\quad A_1=-\frac1{q^{168}-1},\quad
 A_2=\frac1{(q^{167}-1)(q^{168}-1)}.                                \tag{HZ4}
\]

These are the literal unipotent specialization of Green's character formula, as stated in Gorodetsky and Hazan (2017, Theorem2.1 and Corollary2.2). Their absolute values are retained in DO127(GC1-3); the rank-one minus sign here comes from that literal primary formula. Thus

\[
 \sum_{u\in U_m(k)} A(u)
 =\frac{q(q^{167}-q^m)^2}{(q^{167}-1)(q^{168}-1)}.                  \tag{HZ5}
\]

The full m=167 average is exactly zero. For m\textless167 it is positive; there is no zero-mode conclusion for a smaller row subspace merely because it is unipotent. In a unitary realization of the finite representation, \textbar U\textbar\^{}\{-1\}sum\_U R(u) is the orthogonal projection onto U-invariants. Its normalized trace is zero by HZ5, so the projection itself is zero. This is the proper-parabolic Jacquet vanishing relevant to this exact radical, with no replacement of the actual representation.

There is also an exact additive-coordinate Fourier statement. Put n=q\^{}167 and Den=(n-1)(qn-1), and use normalized Fourier coefficients on the additive space k\^{}167 x k\^{}167 x k. For row frequencies xi,chi and central frequency zeta, the coefficient is

\[
 \widehat A(\xi,\chi,\zeta)=
 \begin{cases}
 [n(qn-1)]^{-1},&\zeta\ne0,\\
 \mathrm{Den}^{-1},&\zeta=0,\ \xi\ne0,\ \chi\ne0,\\
 0,&\zeta=0\text{ and }(\xi=0\text{ or }\chi=0).
 \end{cases}                                                       \tag{HZ6}
\]

Here a frequency is nonzero as a complete row functional. To prove this, write the original rank function exactly as

\[
 A=A_2+(A_1-A_2)(1_{a=0}+1_{b=0})
       +(A_2-A_1)1_{a=b=0}+(1-A_1)1_{a=b=z=0}.
                                                                    \tag{HZ7}
\]

Additive orthogonality gives HZ6. All coefficients are nonnegative; their sum is1, and the trivial coefficient is0. This is additive Fourier analysis of the displayed coordinates, not an assertion that the nonabelian group U is abelian. Translation by a general h0 can add phases and need not preserve this nonnegativity.

\subsection{3. The actual paired CM coordinates supply HZ2}\label{the-actual-paired-cm-coordinates-supply-hz2}

At every G place, 2U is a unit. A primitive generator e of the saturated isotropic line extends locally to a unimodular plane: choose f with \textless e,f\textgreater=1 and replace f by f-\textless f,f\textgreater e/2. The quotient is the original selfdual projective quotient. Over the semilocal ring obtained by retaining just the finitely many G primes, the projective modules in question have constant rank and are free: lift a basis at each maximal ideal simultaneously by CRT; Nakayama gives a surjection from a free module of that rank, and localization makes it an isomorphism. Consequently one may choose a single E-rational Witt chart integral and invertible at all G places, with quotient Gram matrix H\_K integral and invertible there. This is only a semilocal chart; it neither trivializes the global projective module nor sets its Gram matrix equal to the identity. (HZ8)

In this chart a rational point has coordinates y in E\^{}167 and t in F:

\[
 u(y,t)=\begin{pmatrix}
 1&-y^*H_K&i t-\tfrac12 y^*H_Ky\\
 0&I&y\\0&0&1
 \end{pmatrix}.                                                     \tag{HZ9}
\]

At a selected split component P, the residues of y\_P and y\_bar(P) are independent by the Chinese remainder theorem in O\_E. Therefore b=y\_P and a=-(bar y)\emph{P\^{}t H}(K,P) range independently over the two complete row spaces in HZ2. The F variable t has an independent residue and i is a unit, so z is arbitrary as well. The other matrix component is still the prescribed inverse adjoint; no extra label, finite vector or independent unitary component has been inserted.

The passage from the pure-imaginary central coordinate it to z changes z by the quadratic term ab/2. If both rows are nonzero the rank2 value does not depend on z; if either row is zero then ab=0. Hence the rank function HZ7, and therefore HZ6, remain valid in the actual paired row and pure-imaginary central coordinates after the invertible local linear coordinate changes. (HZ10)

\subsection{4. CRT completion and its precise cost boundary}\label{crt-completion-and-its-precise-cost-boundary}

Choose a positive integer M coprime to Q, sufficiently divisible at all denominators of the chart and H\_K, at2, and at T. Then points with y in M O\_E\^{}167 and t in M O\_F are integral with integral inverse outside G and are principal to order2 in the original matrix frame at every T place. Quadratic terms are harmless after increasing the fixed M. The choice is possible because all chart denominators at G are units. M is not assigned a height or index bound here. (HZ11)

More intrinsically, let Gamma0 be the intersection of U\_ell(F) with the original standard stabilizers away T, the principal matrix congruence subgroups of order2 at T, and the G-integral subgroups. It is a group. The points just constructed and CRT show that its reduction is onto the product of the full HZ2 groups at G. One may require identity at the two mixed G places while keeping surjectivity at all334 unmixed places. This uses CRT on the paired row coordinates of HZ9 and on its single F central coordinate, not an assumed product of unrelated a,b.

For an actual h0 in the original finite support, right multiplication by delta in Gamma0 preserves every T incidence matrix. In fact each actual pencil lattice L obeys pi Lambda subset L subset pi\^{}-1 Lambda; delta-I and delta\^{}-1-I in pi\^{}2 End(Lambda) imply delta L=L. Thus

\[
 M_{h_0\delta}=M_{h_0}\quad\text{at every T place}.                 \tag{HZ12}
\]

Away T both h0 and delta are in the original compact, so their product is also. The full actual finite coefficient along h0 Gamma0 has zero average over its finite G quotient: at any unmixed G place, left multiplication by the reduction of h0 permutes the full U\_167, and HZ5 applies. Mixed factors and all T factors remain exactly present.

This is a completion theorem, not a short-sum estimate. The CRT points need not satisfy \textbar\textbar h0 delta\textbar\textbar\textless=133. No cost for M, the congruence index, or the required representatives has been proved. In particular the finite average cannot be inserted as the value of the original archimedean-weighted sum. No assertion that every point of this affine coset retains the exact k or the old removal status is needed or made.

\subsection{5. The full actual kernel is positive definite}\label{the-full-actual-kernel-is-positive-definite}

Retain literally K062(AT7-16) and define on the rational group

\[
 F(g)=a(g)c(g)\prod_{v\mid T}\rho_v(g_v),\qquad F(1)=1.            \tag{HZ13}
\]

The archimedean factor is the normalized trace against the positive orthogonal projector onto the actual selected W, after the one fixed conjugation by g\_p.~It is a positive-definite function. At G the same argument applies to the two actual mixed projectors and every unmixed representation. Extending such a function from a compact subgroup by zero is positive definite: partition any finite Gram matrix by cosets of that subgroup, obtaining positive semidefinite blocks. The remaining away-T stabilizer indicators have the same property. No hyperspecial claim is made at ramified rational primes.

At T use the permutation representation on l\^{}2 of all local lattices. The finite-rank orthogonal projection P onto the mean-zero vectors supported on the actual D=k\{-1\} has rank T-2. The actual incidence coefficient is Tr(P U(g))/(T-2), exactly K062(AT7-12), hence is positive definite even though individual values may be negative. Products of these positive-definite functions are positive definite by tensor products. There is no commutation requirement on the projectors. This proves HZ13 on the same original group with all finite and selected factors retained. (HZ14)

For a fixed ell one can choose a finitely generated arithmetic lattice Gamma in U\_ell(F) containing every original finite-support element of U\_ell(F). To see this without a global O\_E frame, use the rational Witt chart and its F-Lie algebra. K062(AT18) gives g,g\^{}-1 in T\^{}-2 Mat(O\_E). Since log(I+N)=N-N\^{}2/2, all relevant logarithms lie in a fixed denominator-bounded Z-lattice. Choose a Z-lattice containing their horizontal projections and a central Z-lattice containing their central projections and half the brackets of the chosen horizontal generators. Their BCH exponential is a discrete, finitely generated, two-step nilpotent subgroup containing every required element. (HZ15)

The restriction F\textbar Gamma is absolutely summable. Its nonzero support is inside K062's original mathscr C; bounded archimedean parts are finite by AT18, and AT20 supplies an absolutely summable full tail. A rational central root of unity other than1 is not unipotent, so this restriction introduces no multiplicity from the original mu\_28 quotient.

Take Følner boxes with horizontal widths L and central widths L\^{}2. A fixed translation changes horizontal coordinates by O(1) and central coordinates by O(L), giving a boundary proportion tending to zero. Positive definiteness implies

\[
 0\le {1\over |E_L|}\sum_{x,y\in E_L}F(x^{-1}y)
 =\sum_{g\in\Gamma}{|E_L\cap E_Lg^{-1}|\over |E_L|}F(g).
                                                                    \tag{HZ16}
\]

The overlap coefficients lie in{[}0,1{]} and tend to1 for every fixed g. Absolute summability and dominated convergence therefore prove

\[
                   \sum_{g\in\Gamma}F(g)\ge0.                    \tag{HZ17}
\]

This proof requires neither an abelian Fourier transform of Gamma nor a replacement of the actual selected trace by a scalar K-type model.

\subsection{6. Why HZ17 does not remove the remaining sector}\label{why-hz17-does-not-remove-the-remaining-sector}

The exact support27/k164 remainder omits the identity. Multiplying F by its indicator gives a function f with f(1)=0. A positive-definite function with f(1)=0 is identically zero, since its two-by-two Gram minors imply \textbar f(g)\textbar\textless=f(1). Thus positivity of HZ13 supplies no positivity of that restricted remainder. Nor is g↦F(h0 g) generally positive definite on a coset subgroup. (HZ18)

Completing a null-line sector to Gamma restores the identity with weight1 and other previously removed subsets. A nonidentity unipotent in U\_ell has intrinsic line ell: the quotient is positive at every archimedean embedding, so a nonzero horizontal vector has nonzero Hermitian norm and N\^{}2 is nonzero. In rank2 the line is im(N\^{}2), and in central rank1 it is im(N). Different line sectors therefore intersect at the identity only, but completing each separately repeats this weight1. HZ17 bounds the nonidentity sum only below by-1 per line. This bound can be attained in the finite example HZ5, whose identity value is1 and whose full sum is0. No small-error bound follows from it.

All previous absolute-removal charges and the single full H133 tail must be reconciled by an actual disjoint set difference. They cannot be paid independently for every coset or every null line. The original positive energy statement K062(AT24) is consistent with HZ17 and does not provide the strict PP3 inequality merely from nonnegativity. (HZ19)

\subsection{7. Exact central-frequency scale in the remaining window}\label{exact-central-frequency-scale-in-the-remaining-window}

DO106(CFI4-8) and DO108(AC8-12) are support-independent arithmetic arguments. At support27, for fixed original row and eta, the complete central lattice is Gamma\_c=B\_eta I\_Sigma,F\^{}2 with norm x T\^{}-54. The physical central coordinate has lattice W=U Gamma\_c and full real covolume U\^{}f sqrt(D\_F) x T\^{}-54. The sufficient G period is QW. Thus the geometric covolume scale of its dual is

\[
 \delta_c=\left({T^{54}\over Q^f U^f\sqrt{D_F}\,x}\right)^{1/f}.
                                                                    \tag{HZ20}
\]

This is a covolume scale, not a minimum-length or coordinate mesh claim. The accompanying integer comparisons prove, throughout the containing window2111\textless log2 x\textless2165,

\[
                    2^{84}<\delta_c<2^{86}.                       \tag{HZ21}
\]

For reference, the normalized containing scalar complex row ideal sqrt(U) D\_eta I\_Sigma has half-trace covolume D0 U\^{}168 x T\^{}-54. Its Q-period dual covolume scale lies strictly between2\^{}5 and2\^{}6. This reference does not identify the full original quotient/central parameter lattice with a product of these scalar ideals.

A stronger rigorous observation avoids interpreting covolume as a mesh. Let L=W\^{}*=U\^{}-1 Diff\_F\^{}-1 B\_eta\^{}-1 I\_Sigma,F\^{}-2. Then N\_F L=T\^{}54/(U\^{}f D\_F x). The real-ideal dual/transference proof of DO091(TD20-22), valid for every fractional F-ideal, gives

\[
 \mu(L)<750\sqrt f\,(N_FL)^{1/f}<2^{164}.                         \tag{HZ22}
\]

The last inequality is an exact integer check at x=2\^{}2111. Every finite character class of W/QW, equivalently every coset in Q\^{}-1 L/L, has a representative of norm below2\^{}164. More generally, for any prescribed sign vector epsilon in\{+1,-1\}\^{}f, approximate the point2\^{}165 epsilon by that coset within mu(L). The resulting representative has all epsilon\_i xi\_i\textgreater2\^{}164 and

\[
               \|\xi\|<(2\sqrt f+1)2^{164}<2^{169}.              \tag{HZ23}
\]

Thus even restricting to the required archimedean sign orthant does not force finite nontrivial character classes to very large frequencies. The selected energies B=2\^{}955 and A=2\^{}960 are retained. HZ23 is not a claim that the actual Fourier density is nonzero, large or small at these representatives; it establishes the missing lower-frequency obstruction to an automatic zero-mode argument. An estimate of the actual selected transform and all its aliases remains necessary.

\subsection{8. Checked scope and remaining obligation}\label{checked-scope-and-remaining-obligation}

HZ2-12 gives exact finite Jacquet cancellation with the original paired-coordinate map and a genuine CRT completion. HZ13-17 gives full-group positivity with every original finite, selected and mass normalization unchanged. HZ18-23 identifies the unresolved cutoff, identity, coset and frequency costs. None of these statements changes the current137/31 arrays, removes k164 or any other cell, proves the signed PP3 gate, supplies a nonzero selected arithmetic image, or proves any prescribed period, source existence or DP44 cost. (HZ24)

The diagnostic script verifies HZ3-7 by exact finite models, the actual Q,T,U/discriminant inequalities in HZ21-23, source hashes and the unchanged actual predecessor. It does not certify the infinite mathematical arguments by finite examples; those are the written proofs above. All future quantitative use must retain the complete archimedean transform or a separately proved full-group lower bound.

\subsection{References and exact inherited interfaces}\label{references-and-exact-inherited-interfaces}

Gorodetsky, O., \& Hazan, Z. (2017). \emph{On certain degenerate Whittaker models for cuspidal representations of GL\_kn(F\_q)}. arXiv:1707.07308. \url{https://arxiv.org/abs/1707.07308}. Freshly read introduction, projection formula(1.5), Theorem2.1 and Corollary2.2. The Green character formula is used only in its displayed unipotent specialization.

K062, AT1-25: original finite orbit, partial incidence, positive projector, normalized kernel, full denominator and absolute tail. DO106, CFI1-24: exact original central lattice and finite Fourier scope. DO108, AC8-15: cross-pole full central fibre and Q-period. DO126, IQ8-19: original selfdual projective quotient and all-prime saturation. DO127, GC1-16: current necessary staircase and exact Green normalization. DO091, TD20-22: full real discriminant/dual covering bound. These inherited findings are kept in their accepted scopes.

\subsection{Canonical algebraic-geometric context}\label{canonical-algebraic-geometric-context-93}

The arithmetic group completion retains the original paired Hermitian projective lattice, finite flags, induced CM field, positive definite quotient and finite projector lines. The real central trace dual is the inverse different times the inverse fractional ideal; its covolume and covering radius have distinct roles. Full positive-definite matrix coefficients and their arithmetic sums do not supply positivity of a selected support mask or a global arithmetic theta cycle. Exact internal locators are K062 AT1-25; DO106 CFI1-24; DO108 AC8-15; DO126 IQ1-36; DO127 GC1-16; DO107 TD20-22, together with Gorodetsky and Hazan(2017), Theorem2.1/equation2.1/Corollary2.2, printed5. Carried exact context: SGA7 I VII2.7/2.9 pp169--174 and VIII7.4(b) pp310--311; FGA232 Theorem3.1 pp148--149 and FGA236 Corollaries1.7/2.3 pp229/232; SGA3 I4.7.3/4.7.3.1; SGA6 VI Theorem3.3; cote4 scan003 sections8.1.1--8.1.2 and scans005--006 sections8.4.1--8.4.3; cote54 scans005--007. These are carried contextual readings, not new whole-text checks. No new integral extension over Z, Dieudonne slope, lambda-ring positivity, K-theory filtration, l-adic weight theorem or CM-motive L-function comparison follows. DP\_MULT Q3e/DP44/DP52/DP54--65 and THETA3 G1/A7 gain this conditional source-rule reduction; CX\_GRH CX-ALG-024 and GCT\_PLETH T-052 retain their global gates.

\clearpage

\section{AM-LIFT-108 --- Quadratic norm image and the actual coefficient law}\label{am-lift-108-quadratic-norm-image-and-the-actual-coefficient-law}

\textbf{Canonical checked integration --- 2026-11-23T14:10:03.940004+00:00.} Complete saved-proof and executable reviews and separate exact diagnostic replays passed. No numerical error charge is introduced. Frozen author history is preserved. AM-LIFT-108 NI1-21 identifies the actual induced quadratic type ideal with the extension of N\_E/K(I), for K=Q(H.rt) embedded as specified in the original F. The complete original kPsi law is equivalent to a finite normalized table on G=image Cl(E)-\textgreater Cl(K), with arbitrary integral representatives of generator classes. Every nonnegative supplied signature is covered in the forward direction. Small quadratic representatives have norm≤(2/pi)sqrt(d), scalar height≤half log of that bound, and explicit class-power roots supply a new coefficient record of relative degree≤\textbar G\textbar≤h\_K, persistent over further upper fields. The original-field table remains unevaluated; the new field does not transport eO/PinField or changed DP44 costs. The geometric d7 class-number-one specialization is excluded by the full source d\textgreater365 and cannot relabel J14/J56 Fields.Kq. No actual proper-image claim or new charge; current138/31 and27/11 floors remain.

Author: \nolinkurl{/root/lift_accounting}. \textbf{Research proof pending independent whole review.} The accompanying author-verification and pending manifest record the exact frozen subjects and completed checks. Accepted Alpöge statements are cited as supplied inputs. This is a source application and quantitative refinement, not a validity audit of those inputs.

\subsection{1. The exact quadratic subfield and the unchanged coefficient field}\label{the-exact-quadratic-subfield-and-the-unchanged-coefficient-field}

Let \(E\) be the actual CM field, let \(H\) be the geometric source Host, and let \(c:\mathrm{Coeff}(l)\) be the original coefficient record. Write \(F=c.F\). The following argument uses the literal

\[
 H.rt^2=-d,\qquad H.rt0^2=-d,\qquad
 d\text{ prime},\quad d\equiv7\pmod {16}.
 \tag{NI1}
\]

It does not require a new moduli datum or a change to the original coefficient field. Put

\[
 K=\mathbf Q(H.rt)\subseteq E,\qquad
 \iota:K\hookrightarrow F,\quad\iota(H.rt)=H.rt0.
 \tag{NI2}
\]

The polynomial \(X^2+d\) is irreducible over \(\mathbf Q\), so \(K\) is imaginary quadratic. The embedding \(\iota\) exists because \(H.rt0\in KA(l)\subseteq F\); alternatively it is the restriction of any of the embeddings in \(\Sigma^+\) from LIFT107 AG3. Its image is the specified embedded quadratic field, not a free quadratic subfield. Since \(-d\equiv1\pmod4\), the standard quadratic integer ring formula gives

\[
 \mathcal O_K=\mathbf Z[(1+H.rt)/2],\qquad D_K=-d.
 \tag{NI3}
\]

The element is the literal \(H.w0\) in \(E\), by the Host half-root identity. The ring and discriminant formula is the accepted quadratic formula (Milne, 2020, Example 2.41, p.~39).

Let \(\Phi(I)\) denote LIFT107 AG4's product of the extended ideals over embeddings \(\tau\) satisfying \(\tau(H.rt)=H.rt0\). Define

\[
 G=\operatorname{im}\left(
 N_{E/K}:\operatorname{Cl}(\mathcal O_E)
                 \longrightarrow\operatorname{Cl}(\mathcal O_K)
                         \right),\qquad g=|G|.
 \tag{NI4}
\]

The ideal norm induces this class-group map because it is multiplicative and sends \((x)\) to \((N_{E/K}x)\). Neither surjectivity onto \(\operatorname{Cl}(\mathcal O_K)\) nor triviality of its image is assumed. In particular no assertion that this image is proper in the actual prime-discriminant CM setting is made.

\subsection{2. The type ideal is the relative norm ideal}\label{the-type-ideal-is-the-relative-norm-ideal}

For every nonzero fractional ideal \(I\) of \(E\),

\[
 \boxed{\quad
 \Phi(I)=\iota(N_{E/K}I)\mathcal O_F,
 \qquad N_{K/\mathbf Q}(N_{E/K}I)=N_{E/\mathbf Q}I.
 \quad}
 \tag{NI5}
\]

\textbf{Proof.} The embeddings in \(\Sigma^+\) are exactly all embeddings of \(E\) extending the fixed \(\iota:K\to\mathbf C\). The separable field-norm formula therefore gives, for every \(x\in E^\times\),

\[
 \prod_{\tau\in\Sigma^+}\tau(x)=\iota(N_{E/K}x).
 \tag{NI6}
\]

For a fractional ideal \(I\), choose \(a>0\) and \(x\in E^\times\) with \(I^a=(x)\). Multiplicativity and principal-ideal norm compatibility give

\[
 \Phi(I)^a=(\iota N_{E/K}x)
      =\bigl(\iota(N_{E/K}I)\mathcal O_F\bigr)^a.
 \tag{NI7}
\]

The fractional-ideal group of \(\mathcal O_F\) is torsion free, so NI7 gives the first equality. The second equality follows on a prime \(\mathfrak P\) above \(\mathfrak p\) above \(p\) from \(N_{E/K}\mathfrak P=\mathfrak p^{f(\mathfrak P/\mathfrak p)}\) and \(f(\mathfrak P/p)=f(\mathfrak P/\mathfrak p)f(\mathfrak p/p)\); multiplicativity then gives all fractional ideals. These are the standard norm formulas, with no omitted ramification hypothesis (Sutherland, 2021a, Theorem 3.11; 2021b, Theorem 4.50; 2021c, Propositions 6.6--6.7 and Theorem 6.10). In particular the argument does not replace a relative ideal norm by a selected set of split primes. \(\square\)

Every \(z\in K^\times\) satisfies, at every embedding of any number field containing \(\iota(K)\),

\[
              |\sigma\iota(z)|^2=|N_{K/\mathbf Q}z|.
 \tag{NI8}
\]

Indeed the restriction to \(K\) is one of its two conjugate complex embeddings. This fact is the reason that a correction by a principal ideal of \(K\) preserves the required metric normalization.

\subsection{3. The smaller normalized table is equivalent to the original law}\label{the-smaller-normalized-table-is-equivalent-to-the-original-law}

Choose nonzero integral \(K\)-ideals \(J_1,\ldots,J_s\) whose classes generate \(G\). Consider the finite condition

\[
 \exists\xi_i\in F^\times:\quad
  (\xi_i)=\iota(J_i)\mathcal O_F,\qquad
  |\sigma\xi_i|^2=N_{K/\mathbf Q}J_i
       \quad\text{for every }i,\sigma:F\hookrightarrow\mathbf C.
 \tag{NI9}
\]

\textbf{Theorem A.} NI9 is equivalent to the rank-one CM-type law \(\mathcal C(H,c)\) of LIFT107 AG5. For an actual original family with \nolinkurl{PinDat}, it is therefore equivalent to the original entire \nolinkurl{LawsP.kPsi} obligation, and to its two adjacent stages.

\textbf{Necessity.} For each \(J_i\), choose an integral \(E\)-ideal \(I_i\) whose norm class is \([J_i]\). Such an integral representative exists by multiplying a fractional representative by a nonzero rational integer. There is \(z_i\in K^\times\) such that

\[
       N_{E/K}I_i=(z_i)J_i.
 \tag{NI10}
\]

If \(\chi_{I_i}\) is supplied by the same-field law, set \(\xi_i=\chi_{I_i}/\iota(z_i)\). NI5 gives its exact ideal in NI9, and NI8 and norm transitivity give its squared absolute norm. The integral ideal equality implies \(\xi_i\in\mathcal O_F\) even when \(z_i\) is not integral. Thus \(J_i\) itself need not be the literal norm of an integral \(E\)-ideal; its class is sufficient.

\textbf{Sufficiency.} Given an arbitrary nonzero integral \(I\), write

\[
 N_{E/K}I=(z)\prod_iJ_i^{a_i},\quad z\in K^\times,
 \ a_i\in\mathbf Z,
 \qquad
 \chi_I=\iota(z)\prod_i\xi_i^{a_i}.
 \tag{NI11}
\]

The fractional ideal of \(\chi_I\) is \(\Phi(I)\) by NI5. Using NI8 and NI9 at every embedding gives

\[
 |\sigma\chi_I|^2
 =|N_{K/\mathbf Q}z|\prod_i(NJ_i)^{a_i}
 =N_{K/\mathbf Q}(N_{E/K}I)=N_{E/\mathbf Q}I.
 \tag{NI12}
\]

Again the integral ideal equality proves integrality. This proves both directions with no change to \(F\).

The forward implication to \nolinkurl{KPsiIn} works for \textbf{any} supplied \(\mathrm{Datum}(H)\), not only the original positive-gap family. For its nonnegative ranks \(p,q\), put \(n=NI\) and

\[
       \beta_I=\chi_I^p(n/\chi_I)^q.
 \tag{NI13}
\]

LIFT107 AG6 gives \((n/\chi_I)=\Psi(I)\); both factors are integral and have squared embedded absolute norm \(n\). Hence NI13 has the literal weighted ideal \(\Phi(I)^p\Psi(I)^q\) and squared norm \(n^{p+q}\). Negative powers of \(\chi_I\) do not impose a new denominator condition. This includes \(p=0\), \(q=0\) and \(p=q=0\). The converse from an arbitrary single pair \((p,q)\) is not asserted; the original two-stage converse is exactly LIFT107.

\subsection{4. Table size, small representatives and a bounded new coefficient record}\label{table-size-small-representatives-and-a-bounded-new-coefficient-record}

Choose an invariant-factor basis for \(G\) with nontrivial orders \(h_1,\ldots,h_s\ge2\). Then

\[
       \prod_i h_i=g\le h_K,
       \qquad s\le\log_2g\le\log_2h_K.
 \tag{NI14}
\]

If \(G\) is trivial, the table is empty and NI11 uses only a principal generator in \(K\). Thus the same-field law holds already in the original \(F\). This is a statement about the norm image; it does not require that \(\operatorname{Cl}(\mathcal O_E)\) be trivial.

Minkowski's theorem permits each generator class to be represented by an integral \(J_i\) with

\[
            1\le NJ_i\le B_d:=\frac{2}{\pi}\sqrt d.
 \tag{NI15}
\]

This is the degree-two, one-complex-place specialization of the standard Minkowski constant and the small-ideal representative theorem (Sutherland, 2021d, Theorems 14.18 and 14.20). In particular the table has a bound depending only on the quadratic discriminant, not on the much larger degree of \(E\). An entirely explicit coarse class-number bound, if needed without computing the group, is

\[
 h_K\le (2B_d)^{\log_2B_d},\qquad
 s\le (\log_2B_d)\log_2(2B_d),
 \tag{NI16}
\]

using Sutherland's ideal-count bound (2021d, Lemma 14.21). Since \(d\ge7\), \(B_d>1\) and these expressions are defined. NI16 is only a bound; the theorem does not claim that \(G\) or its table has been evaluated for a prescribed Host.

\textbf{Theorem B (explicit new coefficient record).} There exists an upper coefficient field \(F_+\supseteq F\), and the corresponding record \(c_+\) on the same \(H,l\), such that NI9 and hence \nolinkurl{KPsiIn} hold for every supplied \(\mathrm{Datum}(H)\), with

\[
                       [F_+:F]\le g\le h_K.
 \tag{NI17}
\]

Here \(F_+\) is a \textbf{new coefficient record}. NI17 makes no claim that it satisfies a previously prescribed \nolinkurl{eO/PinField}.

\textbf{Proof.} For each basis ideal choose a generator \(J_i^{h_i}=(\gamma_i)\) with \(\gamma_i\in\mathcal O_K\setminus\{0\}\). Choose one complex root \(\alpha_i\) of

\[
                   \alpha_i^{h_i}=\iota(\gamma_i),
                   \qquad F_+=F(\alpha_1,\ldots,\alpha_s).
 \tag{NI18}
\]

Successive extension degrees are at most \(h_i\), proving NI17; no irreducibility, normality or roots-of-unity extension is needed. The coefficient axioms persist literally: \(F_+\) is finite over \(\mathbf Q\), contains every old conjugate of \(E\), the old \(\zeta_{l.M'}\), and \(KA(l)\). These are exactly \nolinkurl{Coeff.fin/conj/ zeta_mem/le}. No new Moduli or pinned parameter record is inferred.

The principal ideal equality raised to \(h_i\) is \((\alpha_i)^{h_i}=(\iota\gamma_i)=
 (\iota J_i\mathcal O_{F_+})^{h_i}\), so torsion freeness gives \((\alpha_i)=\iota J_i\mathcal O_{F_+}\). Every embedding \(\sigma\) of \(F_+\) restricts to one of the two embeddings of \(K\), so

\[
 |\sigma\alpha_i|^{2h_i}
   =|\sigma\iota\gamma_i|^2
   =N_{K/\mathbf Q}(\gamma_i)=(NJ_i)^{h_i}.
 \tag{NI19}
\]

Taking the positive real \(h_i\)-th root proves NI9. This is a new-field sufficiency construction; it does not assert a converse for fixed principal generators in the old \(F\). Theorem A and NI13 now give the claimed predicates. The empty-table case has \(F_+=F\) and degree one. \(\square\)

Moreover the very same witnesses supply the laws over \textbf{every} upper coefficient field \(F''\supseteq F_+\). The ideal equalities extend to \(\mathcal O_{F''}\) and each embedding of \(F''\) restricts to an embedding of \(F_+\) with the already proved norm equality. Thus this is the actual \(c_+\le c''\) persistence required by the accepted upper-coefficient supplier, with NI17 bounding the first chosen upper field. No bound is asserted for an arbitrarily larger \(c''\) chosen afterward.

For the small representatives in NI15, any normalized table entry \(\xi_i\), and in particular every chosen \(\alpha_i\), has the usual absolute logarithmic algebraic height

\[
 h_{\mathrm{abs}}(\xi_i)=\tfrac12\log NJ_i
       \le\tfrac12\log B_d<\tfrac14\log d.
 \tag{NI20}
\]

Indeed it is an algebraic integer and every archimedean conjugate has the same modulus \(\sqrt{NJ_i}\ge1\). The defining formula in this integral case is \(h_{\mathrm{abs}}(\xi)=[F:\mathbf Q]^{-1}
\sum_{\sigma:F\hookrightarrow\mathbf C}\log\max(1,|\sigma\xi|)\), with \(F_+\) replacing \(F\) for a new root, so it is exactly the stated value. This is the height of a scalar, not a replacement for the source's Hodge-line, period or DP44 cost. The degree estimate NI17 does not by itself bound the discriminant of the changed coefficient field or certify any original field pin.

The accepted \nolinkurl{exists_coeff_typeNorms_principal} supplies qualitative upper-coefficient existence. Theorem B records the additional bound NI17 and scalar normalization coming from the \textbf{actual induced quadratic type} and the smaller norm image. It is not a re-proof or validity check of that accepted theorem.

\subsection{5. The class-number-one specialization and its exact source boundary}\label{the-class-number-one-specialization-and-its-exact-source-boundary}

If \(K\) has class number one, \(G\) is trivial. For every \(I\), choose \(\chi\in K^\times\) generating \(N_{E/K}I\). Then \(|\sigma\iota\chi|^2=NI\), so NI13 supplies \nolinkurl{KPsiIn} for every supplied datum in the original \(F\), without any extension.

For \(d=7\), the accepted Minkowski bound gives

\[
                   \frac{2}{\pi}\sqrt7<2.
 \tag{NI21}
\]

Every ideal class of \(\mathbf Q(\sqrt{-7})\) therefore has an integral representative of norm one; such an ideal is the unit ideal. Hence \(h_K=1\). This is a direct application of the standard small-ideal theorem, not a new class-number theorem or an appeal to a full class-number-one classification.

\textbf{The \(d=7\) conclusion is conditional at the geometric Host and coefficient interface. It is not an admitted full-source family.} The full arithmetic \nolinkurl{Host.qd_big} requires \(d>4R.t\,R.r+1\). The literal ranks already force \(d>265\), and their actual parity strengthens this to \(d>365\). \nolinkurl{Par.PinHost} requires equality of the geometric and full arithmetic parameters. Thus no full pinned Host in this source has \(d=7\). The immutable R093 scope addendum records the precise declarations; it is retained unchanged.

The existing \(J^{14}\) and \(J^{56}\) polarized product constructions also do not provide a contrary pin. K030 uses the induced type from \(\Phi_0=\{1,3,9,15,17,23\}\) on \(\mathbf Q(\zeta_{28})\). K032's opening and §6, QH37, explicitly retain \(\mathrm{Fields.Kq}=\mathbf Q\). The external \(\mathbf Q(\sqrt{-7})\) containing chosen prime generators cannot be relabeled as that source-defined common quadratic field. Accordingly this theorem does not identify their prescribed CM type with a missing admitted \nolinkurl{PinHost}. Its general-\(d\) norm-image reduction remains the relevant literal-source result.

\subsection{6. Exact application and cost table}\label{exact-application-and-cost-table}

\begin{longtable}[]{@{}
  >{\raggedright\arraybackslash}p{(\columnwidth - 4\tabcolsep) * \real{0.3333}}
  >{\raggedright\arraybackslash}p{(\columnwidth - 4\tabcolsep) * \real{0.3333}}
  >{\raggedright\arraybackslash}p{(\columnwidth - 4\tabcolsep) * \real{0.3333}}@{}}
\toprule\noalign{}
\begin{minipage}[b]{\linewidth}\raggedright
Item
\end{minipage} & \begin{minipage}[b]{\linewidth}\raggedright
Same-field Theorem A
\end{minipage} & \begin{minipage}[b]{\linewidth}\raggedright
New-record Theorem B
\end{minipage} \\
\midrule\noalign{}
\endhead
\bottomrule\noalign{}
\endlastfoot
Original \(E,H,l\) and quadratic embedding & Fixed & Fixed \\
Coefficient field & Original \(F\) & Explicit \(F_+\supseteq F\), relative degree at most \(g\) \\
Arithmetic condition & Normalized table on \(G\) is necessary and sufficient & Chosen class-power roots supply the table \\
Number of entries & \(s\le\log_2g\le\log_2h_K\) & Same \(s\) roots; product of degrees at most \(g\) \\
Scalar height with small ideals & At most \(\frac12\log B_d\) & Same bound \\
\nolinkurl{betaH} and prescribed \nolinkurl{DeltaDat} & Unchanged & No transport or redefinition of an old selected record is asserted \\
\nolinkurl{what₂/eO/PinField} & Retained exactly & Joint reconstruction remains required \\
Coefficient discriminant and all other DP44 terms & Unchanged & Not identified with their old values; numerical changed-field comparison remains required \\
Actual source, selected theta image, global period and032 & Still required & Still required \\
Numerical finite-sector accounting & No new charge;current138/31,floors27/11 & No new charge;current138/31,floors27/11 \\
\end{longtable}

On \(\operatorname{Spec}\mathcal O_K\), the relative ideal norm is the usual norm on invertible ideals; extension along \(\iota\) gives the ideal line over \(\operatorname{Spec}\mathcal O_F\). The reduction uses its image in the finite Picard group together with the full archimedean normalization. Ordinary principalization without that normalization is not substituted. No motivic, weight, positivity, global nonvanishing or DP44 sign conclusion is drawn.

\subsection{7. Exact sources and verification scope}\label{exact-sources-and-verification-scope-1}

The literal inputs are \nolinkurl{TwoWeightEqOneAddRoot.lean}14--26 and30--49, \nolinkurl{FlagBasis.lean}76, \nolinkurl{PushInvertible.lean}40--46, \nolinkurl{EmbKCoe.lean}25--33, \nolinkurl{ThShort.lean}277--285, \nolinkurl{HasDataCondition.lean}55, \nolinkurl{PartitionExtendByZeroAntitone.lean}24--49 and \nolinkurl{MapRingEquivTransApply.lean}52--69. The source uses the already accepted ideal/submodule comparison \nolinkurl{isBe_iff/kPsiIn_iff} and the accepted coefficient supplier; no accepted proof is audited.

AM-LIFT-107 AG1--19 supplies the original two-signature equivalence. Its accepted canonical source is \nolinkurl{agents/adjacent_signature_same_field_type_norm.md}, SHA-256 \nolinkurl{8b2ff2ea66e277a11600c2f5f2166ad36633b3c2238e0b9800fd376a4dab2694}. The original137/31 accounting in its frozen source is historical; the current K123 prefix has138/31. AM-R-102 provides the accepted beta-law application. The existing full-Host exclusion is \nolinkurl{recovery_pending/lift/r093_full_host_qd_big_scope_precision.md}; K030's opening and K032's opening/§6 are read for the precise \(J^{14}/J^{56}\) scope, not as an actual type/pin constructor.

\subsubsection{Primary references}\label{primary-references-3}

Milne, J. S. (2020). \emph{Algebraic number theory} (Version 3.08, Example 2.41, p.~39). \url{https://www.jmilne.org/math/CourseNotes/ANT.pdf}

Sutherland, A. V. (2021a, September 15). \emph{Properties of Dedekind domains} (MIT18.785 Lecture3, Theorem3.11, pp.3--4). \url{https://ocw.mit.edu/courses/18-785-number-theory-i-fall-2021/mit18_785f21_lec3.pdf}

Sutherland, A. V. (2021b, September 20). \emph{Étale algebras, norm and trace} (MIT18.785 Lecture4, Theorem4.50, p.12). \url{https://ocw.mit.edu/courses/18-785-number-theory-i-fall-2021/mit18_785f21_lec4.pdf}

Sutherland, A. V. (2021c, September 27). \emph{Ideal norms and the Dedekind--Kummer theorem} (MIT18.785 Lecture6, Definition6.5, Propositions6.6--6.7, Theorem6.10 and Proposition6.12/Remark6.13, pp.4--5,7--8). \url{https://ocw.mit.edu/courses/18-785-number-theory-i-fall-2021/mit18_785f21_lec6.pdf}

Sutherland, A. V. (2021d, October 27). \emph{The geometry of numbers} (MIT18.785 Lecture14, Theorems14.18/14.20 and Lemma14.21, pp.6--8). \url{https://ocw.mit.edu/courses/18-785-number-theory-i-fall-2021/mit18_785f21_lec14.pdf}

The exact source manifest records the local paths, hashes and actual reading scopes, including the whole frozen LIFT107 proof and the targeted \(J\)-product scope. The executable gives finite checks of relative-norm valuation formulas, finite norm-image/table transport, quadratic integral norm models, all nonnegative \((p,q)\) endpoints, the coefficient-degree bound and the strict \(d=7\) source boundary. It is not a proof of an arbitrary number-field assertion or a computation of the actual norm-image table in the original pinned field. The universal arguments are NI1--21 above; independent whole review remains required before canonical use.

\subsection{Canonical algebraic-geometric context}\label{canonical-algebraic-geometric-context-94}

On Spec(O\_K), relative ideal norm pushes invertible ideal classes from E to the finite Picard group of the specified imaginary quadratic field; extension along its original embedding gives the exact CM-type line over Spec(O\_F). The table is a trivialization with every archimedean modulus, not principalization alone. The new coefficient field is an explicit root tower with bounded relative degree and normalized scalar height, while its arithmetic discriminant and prescribed field identity remain separate. Exact source locators are TwoWeightEqOneAddRoot14-26/30-49; FlagBasis76; PushInvertible40-46; EmbKCoe25-33; ThShort277-285; HasDataCondition55; Partition24-49; MapRing52-69; DataHost39-54 and Params6-22/32. Primary norm/class/Minkowski locators are Sutherland2021 Lectures3/4/6/14 and Milne2020 ANT Example2.41, as stated in the proof. Carried exact context: SGA7 I VII2.7/2.9 pp169--174 and VIII7.4(b) pp310--311; FGA232 Theorem3.1 pp148--149 and FGA236 Corollaries1.7/2.3 pp229/232; SGA3 I4.7.3/4.7.3.1; SGA6 VI Theorem3.3; cote4 scan003 sections8.1.1--8.1.2 and scans005--006 sections8.4.1--8.4.3; cote54 scans005--007. These are carried contextual readings, not new whole-text checks. No new integral extension over Z, Dieudonne slope, lambda-ring positivity, K-theory filtration, l-adic weight theorem or CM-motive L-function comparison follows. DP\_MULT Q3e/DP44/DP52/DP54--65 and THETA3 G1/A7 gain this conditional source-rule reduction; CX\_GRH CX-ALG-024 and GCT\_PLETH T-052 retain their global gates.

\clearpage
\part{Proofs awaiting final whole-proof, source or recovery-file reconciliation}

The following contributions are retained during recovery and final reconciliation. Some have whole mathematical or fresh saved-file checks already recorded; their exact check histories and remaining source/file conditions govern their status. Inclusion does not upgrade a pending condition.

\clearpage

\section{Exact norm boundary for the residual packet plane}\label{exact-norm-boundary-for-the-residual-packet-plane}

Author: \nolinkurl{/root/definite_oscillator}, 22 November 2026 UTC. This is a bounded negative route result, not a new sector removal. It changes no canonical file, error charge or support floor.

\subsection{1. Actual field and conclusion}\label{actual-field-and-conclusion}

Use the unchanged J56 field of K032(QH9--QH11, QH15):

\[
 E=\mathbb Q(\zeta_{28})B\mathbb Q(\sqrt5,\sqrt{13}),\qquad
 [B:\mathbb Q]=7,\quad B\subset\mathbb Q(\zeta_{29}),\qquad F=E^+.
\]

The four factors have their actual coprime conductors and the stated linear disjointness. Thus

\[
 \operatorname{Gal}(E/\mathbb Q)
 = (\mathbb Z/28\mathbb Z)^\times
   \times ((\mathbb Z/29\mathbb Z)^\times/H)
   \times\{\pm1\}^2,
 \quad H=\{1,12,17,28\}.
 \tag{NB1}
\]

Complex conjugation is exactly \nolinkurl{(-1,H,1,1)}. The quadratic extension E/F is unramified at every finite place, as proved for this actual field in K032. The original packet plane has Gram matrix \nolinkurl{diag(U,-1)}, with

\[
\begin{split}
 Q&=1427908193729121522601,\\
 U&=77Q^2-14=156996979348344707969383396772362042643000463.
\end{split}\tag{NB2}
\]

\textbf{Conclusion.} U is a local norm from E/F at every place and hence a global norm. In particular this packet plane is isotropic over E. Its anisotropy cannot be used to eliminate the doubly-plane scalar-cross residual.

This conclusion provides neither a bounded-height nor an integral norm solution. It does not produce an actual arrow in the prescribed short sum, and does not contradict any of its retained finite conditions.

\subsection{2. Complete integer factorization and certified primes}\label{complete-integer-factorization-and-certified-primes}

The exact factorization is

\[
 U=3^3\cdot7\cdot12837557489173
      \cdot111395838330757\cdot580868899958147.\tag{NB3}
\]

The companion certificate verifies the full product and exact multiplicities. Each displayed factor is prime by a recursive Lucas primality certificate, with completely factored \nolinkurl{p-1}, a witness of full order, and recursively certified prime divisors. Its verification uses only integer products, modular powers and gcds. It does not rely on a probable-prime assertion or the behavior of an unfactored composite.

\subsection{3. Actual Frobenius computation at the unramified support}\label{actual-frobenius-computation-at-the-unramified-support}

For a rational prime p outside 2,5,7,13,29, represent its Frobenius in NB1 by

\[
  (p\bmod28,\ p^4\bmod29,\ (p/5),\ (p/13)).\tag{NB4}
\]

The second coordinate identifies the quotient by H with the subgroup of seventh roots in the residue multiplicative group. Let d\_E be the order of this tuple, and d\_F its order modulo the order-two subgroup generated by complex conjugation. A prime of F above p splits in E precisely when \nolinkurl{d_E=d_F}. Since the extension is abelian, the result holds for every prime above p.

The exact table is:

\begin{longtable}[]{@{}
  >{\raggedleft\arraybackslash}p{(\columnwidth - 14\tabcolsep) * \real{0.1250}}
  >{\raggedleft\arraybackslash}p{(\columnwidth - 14\tabcolsep) * \real{0.1250}}
  >{\raggedleft\arraybackslash}p{(\columnwidth - 14\tabcolsep) * \real{0.1250}}
  >{\raggedleft\arraybackslash}p{(\columnwidth - 14\tabcolsep) * \real{0.1250}}
  >{\raggedleft\arraybackslash}p{(\columnwidth - 14\tabcolsep) * \real{0.1250}}
  >{\raggedleft\arraybackslash}p{(\columnwidth - 14\tabcolsep) * \real{0.1250}}
  >{\raggedleft\arraybackslash}p{(\columnwidth - 14\tabcolsep) * \real{0.1250}}
  >{\raggedleft\arraybackslash}p{(\columnwidth - 14\tabcolsep) * \real{0.1250}}@{}}
\toprule\noalign{}
\begin{minipage}[b]{\linewidth}\raggedleft
p
\end{minipage} & \begin{minipage}[b]{\linewidth}\raggedleft
p mod28
\end{minipage} & \begin{minipage}[b]{\linewidth}\raggedleft
Order in the 28 factor
\end{minipage} & \begin{minipage}[b]{\linewidth}\raggedleft
Order in the B factor
\end{minipage} & \begin{minipage}[b]{\linewidth}\raggedleft
(p/5)
\end{minipage} & \begin{minipage}[b]{\linewidth}\raggedleft
(p/13)
\end{minipage} & \begin{minipage}[b]{\linewidth}\raggedleft
d\_E
\end{minipage} & \begin{minipage}[b]{\linewidth}\raggedleft
d\_F
\end{minipage} \\
\midrule\noalign{}
\endhead
\bottomrule\noalign{}
\endlastfoot
3 & 3 & 6 & 7 & -1 & 1 & 42 & 42 \\
12837557489173 & 25 & 3 & 7 & -1 & 1 & 42 & 42 \\
111395838330757 & 1 & 1 & 7 & -1 & -1 & 14 & 14 \\
580868899958147 & 3 & 6 & 7 & -1 & -1 & 42 & 42 \\
\end{longtable}

The certificate computes every tuple power and the quotient order directly. In particular, these are not splitting conclusions inferred only from one Legendre symbol. All four rational primes in the table split relatively in E/F.

\subsection{4. The ramified rational prime seven}\label{the-ramified-rational-prime-seven}

At seven, inertia in NB1 is

\[
 I_7=\{(a,1,1,1):a\in(\mathbb Z/28\mathbb Z)^\times,
                         a\equiv1\pmod4\}.
\tag{NB5}
\]

It has order six. A lift of the unramified Frobenius on the other factors is

\[
       \varphi_7=(15,7^4\bmod29,-1,-1).
\tag{NB6}
\]

The first entry is one modulo seven and three modulo four; changing this lift by inertia changes no decomposition group. The decomposition group is generated by \nolinkurl{I_7} and \nolinkurl{varphi_7}. Every element of that group has the same sign on the modulo-four factor and on the \nolinkurl{sqrt5} factor. Complex conjugation has modulo-four sign -1 and \nolinkurl{sqrt5} sign +1, so it is not in the decomposition group. Hence every F prime above seven also splits in E/F. The certificate additionally enumerates the group: it has order84, does not contain complex conjugation, and its inertia has order6. Thus the absolute ramification index in F is six, consistently with the relative split conclusion.

\subsection{5. The norm consequence and its limit}\label{the-norm-consequence-and-its-limit}

At every prime dividing U the relative algebra is split, so every local element is a norm. At the other finite primes U is a unit; because E/F is finite-unramified, the unit norm map is surjective. At each real place U is positive, so it is a norm from C to R. The local statement is therefore complete.

Milne (2020), Chapter III, Proposition1.2, printed p.98, supplies the unramified unit-norm statement. The cyclic Hasse norm theorem, Chapter VIII, Theorem3.1, printed p.234, then gives an element \nolinkurl{b} in E with \nolinkurl{b*bar(b)=U}. These primary statements were directly read in the author's PDF. The vector \nolinkurl{(1,b)} is isotropic for \nolinkurl{diag(U,-1)}.

This records the boundary of the anisotropy strategy. The existence of b comes with no assertion that it has the actual integral content, level conditions, finite incidences, archimedean bound or prescribed real coefficient of a surviving term. Those remain separate. In particular, no period, arithmetic image, signed inequality or DP44 conclusion follows.

\subsubsection{Reference}\label{reference}

Milne, J. S. (2020). \emph{Class field theory} (Version4.03). Chapter III, Proposition1.2, p.98; Chapter VIII, Theorem3.1, p.234. \url{https://www.jmilne.org/math/CourseNotes/CFTc.pdf}

The actual field and finite-unramified input are the directly read K032(QH9--QH15), with K030(PH2) for the exact subgroup H. No host or coefficient field is changed.

\clearpage

\section{Auxiliary note --- the integral common target in K093 is not of bounded height}\label{auxiliary-note-the-integral-common-target-in-k093-is-not-of-bounded-height}

Author: \nolinkurl{/root/kummer_extension}, 23 November 2026 UTC.

\textbf{Scope.} This note checks one proposed interpolation adapter. It gives no new finite-sector charge, selected-image nonvanishing, upper-content bound, or universal obstruction to torsion interpolation. In particular it gives no lower bound growing with T for the norm of the rational Weyl vector itself.

\subsection{1. The exact input from the accepted construction}\label{the-exact-input-from-the-accepted-construction}

K093, section 2, equation RT4 and its construction, chooses a common target factor \nolinkurl{w_f} in the actual finite projective middle lattice \nolinkurl{L1}. The source dual row carries the varying multiplier. At each position prime \nolinkurl{p} above the completely split rational prime T, that construction imposes

\[
w_f\in\mathfrak p^2L_1.
\]

Its opposite-component residues realize the prescribed primitive middle covector. Thus the common target is nonzero for the nontrivial family. Further divisibility away T is retained, including all bad frame places and the exact joint Heisenberg stabilizer. None of those further conditions weakens the displayed inclusion.

Let \nolinkurl{a} be the product of these selected position primes, one from each complex-conjugate pair. Since T is completely split,

\[
\mathfrak a\overline{\mathfrak a}=T O_E,
\qquad w_f\in\mathfrak a^2L_1.
\tag{LH1}
\]

The actual middle lattice is integral for its positive Hermitian pairing \nolinkurl{h1}; it is not replaced by a rational coordinate lattice. The integrality and positive trace metric are the actual selfdual-class interfaces of LIFT48, section 3, OB8--OB9. The proof below uses this Hermitian pairing, before multiplying by the codifferent trace coefficient.

\subsection{2. A norm statement with the exact trace metric}\label{a-norm-statement-with-the-exact-trace-metric}

Put \nolinkurl{alpha=h1(w_f,w_f)}. Sesquilinearity and integrality give

\[
\alpha\in\mathfrak a^2\overline{\mathfrak a}^{\,2}O_E
       =T^2O_E.
\]

The scalar is fixed by complex conjugation and positive at every real embedding. Therefore

\[
\alpha=T^2\beta,\qquad
0\ne\beta\in O_F,\qquad \tau(\beta)>0
\quad\text{for every }\tau:F\hookrightarrow\mathbb R.
\tag{LH2}
\]

The contraction to \nolinkurl{O_F} is exact: \nolinkurl{alpha/T^2} is integral in E and lies in F, hence belongs to the integral closure \nolinkurl{O_F}. There is no class-generator assumption on \nolinkurl{a} or \nolinkurl{L1}. Consequently

\[
N_{F/\mathbb Q}(\alpha)
 =T^{2f}N_{F/\mathbb Q}(\beta)\ge T^{2f}.
\tag{LH3}
\]

This is a field-norm bound, not the assertion \nolinkurl{tau(alpha)>=T^2} separately at every embedding. With the actual positive coefficients \nolinkurl{q_tau=-2i tau(xi)} retained from LIFT48, the exact metric satisfies

\[
\begin{aligned}
H_1(w_f,w_f)
 &=T^2\sum_\tau q_\tau\tau(\beta)\\
 &\ge fT^2\left(\prod_\tau q_\tau\right)^{1/f}
       N_{F/\mathbb Q}(\beta)^{1/f}
 \ge {2f\over \operatorname{rd}(E)}T^2.
\end{aligned}
\tag{LH4}
\]

The last equality of constants is exactly OB8--OB9. No finite valuation is assigned to the complex coefficients \nolinkurl{q_tau} or to a CM period.

\subsection{3. The narrow consequence for a line-pullback adapter}\label{the-narrow-consequence-for-a-line-pullback-adapter}

Suppose an interpolation argument uses the integral map on the standard rank-one CM module \nolinkurl{O_E -> L1}, \nolinkurl{x -> w_f x}, and the actual Hermitian polarization. Its pulled-back Hermitian form is \nolinkurl{alpha x conjugate(y)}, hence \nolinkurl{T^2 beta x conjugate(y)}. In that adapter the integral Riemann form, relative to the same codifferent convention, is divisible by \nolinkurl{T^2}. It is therefore not a polarization of degree bounded independently of T.

More explicitly, if the unit rank-one CM module is given its carried principal trace polarization, multiplication of its Hermitian form by \nolinkurl{alpha} changes the determinant of its alternating form by

\[
N_{E/\mathbb Q}(\alpha)
 =N_{F/\mathbb Q}(\alpha)^2\ge T^{4f}.
\tag{LH5}
\]

This determinant computation is conditional on using that integral module map and that polarization adapter. It does not identify the full prescribed selected two-edge evaluation with a section on this rank-one subvariety. The same initial coefficient vector, varying middle classes, terminal functional, ordinary residual derivatives, and the shared finite sum of K093/K095 have not been removed.

A zero estimate based on the largeness of T must therefore calculate the actual pulled-back polarization and derivative divisor. It cannot treat the particular common integral target in RT4 as a fixed-height direction. This is the only excluded shortcut.

\subsection{4. Why the conclusion is not a height obstruction for the rational translate}\label{why-the-conclusion-is-not-a-height-obstruction-for-the-rational-translate}

K093 represents the finite multiplier by \nolinkurl{q=w_f v_lambda/T}, with the exact source row, Heisenberg phase, and class frame. The preceding proof concerns \nolinkurl{w_f}, not \nolinkurl{w_f/T} or \nolinkurl{q}. Division by T removes the forced \nolinkurl{T^2} factor from a Hermitian self-pairing. The divided target need not be an integral map from the unchanged unit CM module; replacing its domain lattice would require its own polarization and torsion-grid calculation.

The elementary split example already shows the distinction. Over \nolinkurl{E=Q(i)}, let \nolinkurl{T=A^2+B^2} be an odd split prime and \nolinkurl{pi=A+Bi}. In the unit Hermitian lattice choose \nolinkurl{w=pi^2}. Then \nolinkurl{w} is divisible by the square of the position prime and is a unit at its conjugate, while

\[
h(w,w)=T^2,
\qquad h(w/T,w/T)=1.
\tag{LH6}
\]

For T=5 one can take \nolinkurl{w=3+4i}; the rational vector \nolinkurl{(3+4i)/5} has norm one and is not in \nolinkurl{Z[i]}. This example is only a check of the distinction. It is not a construction of the actual mixed-CM finite tensor or a witness for its selected derivative.

The accepted K093 conclusion remains unchanged: RT4 is an exact operator identity without a height upper bound; RT12 retains positive ordinary residual derivative orders; RT13 retains the full class sum. This auxiliary observation supplies neither their nonvanishing nor a general impossibility theorem.

\subsection{Source and verification record}\label{source-and-verification-record-3}

The exact accepted inputs freshly read for this note are \nolinkurl{agents/rank_one_translate_residual_derivative.md}, section 2 (RT4 and the module-CRT construction), and \nolinkurl{agents/actual_normalized_theta_operator_bound.md}, section 3 (OB8--OB9). The norm and determinant conclusions above are direct deductions, not quotations from a torsion zero-estimate theorem. The companion record pins both source hashes and distinguishes this auxiliary scope from a new main-gate finding.

\clearpage

\section{AM-K098 --- the Gaussian row reduces the rank-three denominator tail}\label{am-k098-the-gaussian-row-reduces-the-rank-three-denominator-tail}

Author: \nolinkurl{/root/kummer_extension}, ALP\_MACHINE, 23 November2026 UTC.

\textbf{Unpromoted fallback draft.} The author has not completed the whole saved-file/script verification or final certificate replay. Root's proposed R084 paired-line argument now has priority and may remove the entire rank-three sector. The exact endpoint probe behind this draft is retained as a separate calculation; it is not a final proof verdict. No canonical file is modified here.

\subsection{1. Exact sector and conclusion}\label{exact-sector-and-conclusion-2}

Keep the actual J56 all-top normalized kernel at H=133 and f=168. Retain all original fields, CM components, paired pencils, finite projectors, selected real trace, compact factors, packet, central phase, class masses, source fibre and Haar. Consider the support25 scalar-cross sector after the existing R082 cuts, with

\[
 h_{ZZ}=I,\quad \dim_E Z=167,\quad
 \operatorname{rank}_E(h-I)=3,\quad C|_{S_{164}}\ne0,
 \qquad k=\#\{i:r(e_i)\ne0\}\ge126.                 \tag{GT1}
\]

Every unmixed G reduction is unipotent, as in the carried kernel. Global unipotence is not assumed. In fact DO095 RA9 makes the global rank-three alternative nonunipotent, and DO074 gives its actual type-III count 11\textless=t\textless=25. The lower bound on k follows from the complete R082 table, whose smallest retained threshold is125.

Normalize the cross line over E as

\[
 Cx=(u r(x),\eta r(x)),\quad \eta\bar\eta=U,\quad
 Q=1427908193729121522601,\quad U=77Q^2-14,
 \quad M=\frac{H}{2\sqrt U}.                              \tag{GT2}
\]

Let D=D\_eta be the exact integral denominator ideal of eta, and put d=N\_E D. This rational normalization does not make an isotropic frame integral.

\textbf{Theorem.} The entire GT1 sector with

\[
                         d\ge 2^{3358}                     \tag{GT3}
\]

has total absolute actual normalized contribution strictly less than2\^{}-383. Thus the rank-three remainder, after K097 and this theorem, has

\[
                          2^{120}<d<2^{3358}.                \tag{GT4}
\]

This improves the upper endpoint2\^{}4864 in R083. Remove only the intersection with the current remainder and retain every earlier numerical error, including R083's5400 charge. Neither the rank-two branch nor the complete support floors25/11 change.

\subsection{2. Exact denominator enumeration and the common cross ideal}\label{exact-denominator-enumeration-and-the-common-cross-ideal}

The complete eta enumeration is R083 LD3--6, also checked independently in DO104 CD5--9. For one exact integral denominator D, the cleared numerator A=(eta)D is integral and coprime to D. At a relatively split pair, eta valuations are a,e-a with e=v(U)\textgreater=0. A negative valuation specifies both numerator exponents. When neither is negative, only the e+1 allocations of e remain. Relatively inert places have no denominator and one possible valuation.

In the actual field the complete factorization is

\[
 U=3^3\,7\,12837557489173\,111395838330757\,580868899958147.
                                                               \tag{GT5}
\]

All F-primes above these factors split in E/F. Their numbers and U-valuations are (4,3),(2,6),(4,1),(12,1),(4,1), respectively. The exact Frobenius/inertia calculation and recursive prime certificates are those in the fully read packet-plane norm boundary NB1--6. There are therefore at most 4\^{}4 7\^{}2 2\^{}20 numerator ideals for one D. Each candidate principal fractional ideal has at most28 norm-U generators, by the accepted CM norm-one unit calculation. Nonprincipal candidates give no solution. Hence

\[
              \#\{\eta:D_\eta=D,\ \eta\bar\eta=U\}<2^{39}.
                                                               \tag{GT6}
\]

This is uniform in the actual ideal D, not merely in its norm. There is no class-number factor or uncounted numerator-ideal sum.

Fix one canonical whole-C G pattern, with ideal J and at most2\^{}668 patterns. Put

\[
 L=J I_\Sigma,\quad N_E I_\Sigma=T^{-50},\quad
 z'=(N_E(LD))^{1/(2f)}.                                    \tag{GT7}
\]

R072 PJ4--8 and R083 LD7 give both r(y) and eta r(y) in L for each original S input and each of the three fixed integral transverse inputs y1,y2,z0. The source and target memberships hold at both CM partners, and primitive integer Bezout coefficients of the original u extract r(y) without dividing by U. Therefore

\[
                r(y)\in L\cap\eta^{-1}L=LD                 \tag{GT8}
\]

for all167 counted values, even if D meets G,T or U. The last equality is the exact primewise maximum of the two valuation requirements. Given eta and these values, the complete cross map is determined; the original projected-column and source-fibre counts will still be retained below.

At each archimedean label the original negative row gives \textbar\textbar r\_tau\textbar\textbar\textless M, as in R072 PJ7; this bound applies without global unipotence. A nonzero element of LD has half-trace length at least sqrt(f)z'. Thus the k nonzero original S values give the necessary condition

\[
                         k(z')^2<M^2.                       \tag{GT9}
\]

No moving positive basis or finite-field positivity is used.

\subsection{3. Joint S-row count and the true three-dimensional transverse metric}\label{joint-s-row-count-and-the-true-three-dimensional-transverse-metric}

Apply the exact general spherical-cap/Gaussian lemma R072 PJ10--12 to the original k nonzero coordinates in LD. Their complete joint count is bounded by

\[
 {164\choose k}47^k2^{2fk}(f!)^k
       \left(\frac{(87/32)M^2}{fk(z')^2}\right)^{fk}.
                                                               \tag{GT10}
\]

The binomial already counts the original positions. The factor87/32 bounds e, and the47 bound is for the actual real dimension336. This is a bound for every positive k, not an extension of the numerical support range of an earlier theorem.

For the three transverse inputs retain the original value-coordinate Gram matrix G of R072 PJ4--5, with determinant4U. The covector metric is aG\^{}-1a*, so the lattice (LD)\^{}3 in this metric has covolume D\_0\^{}3(z')\^{}\{6f\}/(4U)\^{}f, where

\[
 D_0=7^{140}29^{144}5^{84}13^{84},\qquad
 A_3=\frac{528^{3f}}{(3f)!D_0^3}<2^{-4040}.                  \tag{GT11}
\]

The arbitrary-complex-ideal cover used in DO098 CF17--19 gives covering radius at most (2250/7)sqrt(3 tr(G\^{}-1))sqrt(f)z'. Its exact Gram computation shows that the coefficient before sqrt(f)z' is less than1/1024. Since z'\textless M by GT9, translated Voronoi cells fit in a ball of radius(1+1/1024)M sqrt(f), and (1025/1024)\^{}\{6f\}\textless4. Consequently the full transverse count is

\[
                      4A_3(4U)^f M^{6f}(z')^{-6f}.          \tag{GT12}
\]

The derivation uses only the actual covector norm, the fixed Gram matrix and the arbitrary-ideal cover. In particular DO098's special unipotent value M=sqrt(R/U) is not imported: GT12 here uses exactly the larger general M in GT2. The transverse coordinates are counted in their true metric; no product of three unrelated coordinate-radius estimates or original-ring-only cover is used. Multiplying GT10 and GT12 enlarges the joint covector ball, which is a valid upper count.

\subsection{4. Every remaining actual factor and the monotone cap}\label{every-remaining-actual-factor-and-the-monotone-cap}

The same whole-C Green estimate and the actual selected real coefficient are R083 LD9/LD12:

\[
 |a(h)|\le2Q^{-335\cdot167}T^{8375}(z')^{335f}d^{-335/2},
 \qquad
 |c(h)|<2^{-929f/2}U^{-f/2}(z')^{-f}.                       \tag{GT13}
\]

The real bound is the original selected trace with both compact factors and near-point normalization. Its product lower bound uses a nonzero original xi in LD, and so has exactly z' rather than the old z. It is not an operator norm or a new vector.

Retain the entire R076/R083 projected negative-column count

\[
                         191^{4f}U^fT^{100}                 \tag{GT14}
\]

and the positive rank-two source fibre below2\^{}36. Also retain support choices2\^{}168, all2\^{}668 G patterns, the old pivot factor2\^{}8 and the one Green correction2. These last factors, including the source fibre, give exponent881. Eta is counted by GT6 instead of an additional beta norm fibre; it is not left uncounted.

For fixed actual t, the complete general type/pair sum of R082 TP3--4 is

\[
       2^{75-2t}T^{25-t},\qquad 11\le t\le25.               \tag{GT15}
\]

It contains every arrangement and paired label. It is valid for the nonunipotent rank-three case, including the admitted zero target. Its logarithm is at most P\_11=53+14L\_T. Summing the15 possibilities costs less than2\^{}4. No additional 3\^{}25 or independent conjugate-label sum is multiplied.

Before maximizing the ideal-dependent part, the product of GT10--14 has

\[
             d^{-335/2}(z')^{(328-2k)f},\qquad126\le k\le164.
                                                               \tag{GT16}
\]

The z' exponent is nonnegative, including the zero exponent at k164. Use GT9 as an upper cap. Its combination with the k\^{}\{-fk\} already in GT10 yields k\^{}\{-164f\}, while the complete remaining M power is334f. Define the exact rational row factor

\[
 F_k=\frac{\binom{164}{k}47^k2^{2fk}(f!)^k87^{fk}}
           {32^{fk}f^{fk}k^{164f}}.                         \tag{GT17}
\]

The resulting M and U powers, before any directed estimates, are

\[
          (H/2)^{334f}U^{-331f/2}.                          \tag{GT18}
\]

This includes the transverse determinant U\^{}f, the projected column U\^{}f, and the selected trace U\^{}-f/2. Every paired T power is retained: 335s from Green and4s from the projected column, for8475 at s25, followed by the actual type factor GT15.

\subsection{5. The entire ideal tail and exact39-cell arithmetic}\label{the-entire-ideal-tail-and-exact39-cell-arithmetic}

For X=2\^{}3358, sum over all integral E-ideals of norm at leastX, enlarging the set of possible eta denominators. The elementary ideal Euler bound and the inequality d\^{}-335/2\textless=X\^{}-331/2 d\^{}-2 give

\[
 \sum_{N_ED\ge X}(N_ED)^{-335/2}
       \le X^{-331/2}\zeta_E(2)<X^{-331/2}2^{336}.           \tag{GT19}
\]

This includes every ideal class and the full unbounded tail. The bound zeta\_E(2)\textless=zeta\_Q(2)\^{}336\textless2\^{}336 follows prime by prime, with no analytic continuation, equidistribution or class-number estimate. The per-ideal eta multiplicity GT6 has already been retained. There is no extra union over attained versus unattained cap endpoints.

Use the actual strict rational logarithmic bounds

\[
 L_T=1019333/1024,\quad L_Q=287843/4096,\quad
 L_U=601356/4096,\quad L_H=28900/4096,\quad
 L_{191}=15519/2048,                                        \tag{GT20}
\]

where T,H,191 are bounded above and Q,U below. All their signs follow from the explicit positive products before substitution. Let c\_k be the least strict integral binary ceiling of F\_k. Combining GT6 and GT10--19 gives, for each k, the strict logarithmic upper bound

\[
                         E_k=C_*+c_k-(331/2)3358,           \tag{GT21}
\]

with

\[
\begin{split}
 C_*={}&881+39+2-4040+4fL_{191}-929f/2+2f\\
       &+334f(L_H-1)-(331f/2)L_U-335\cdot167L_Q\\
       &+8475L_T+(53+14L_T)+336+4
        =2872231273/4096.                                  \tag{GT22}
\end{split}
\]

The +2 is the full transverse-cover inflation, +2f is its4\^{}f determinant factor, +39 counts every eta for a denominator, +336 is the infinite ideal sum, and +4 already sums all15 actual t values.

The companion script computes F\_k by exact binomials, factorials and integer powers for every k126..164 and proves both sides of every strict binary bracket. Its complete maximum is

\[
 \max E_k=E_{164}=-1596055/4096<-389,
                  \qquad c_{164}=-145869.                  \tag{GT23}
\]

There are39 support-size cells. Their union is therefore strictly less than39*2\^{}-389\textless2\^{}-383. This proves the theorem. No positive failed-bound diagnostic is used as a lower bound or an existence statement.

\subsection{6. Chronology, exact interfaces and remaining gates}\label{chronology-exact-interfaces-and-remaining-gates}

The frozen author predecessor is canonical DO104's complete95/31 arrays. The certificate accepts \nolinkurl{--prior} for an unchanged-script replay against any later complete extension, checking the whole frozen prefix. Append383 once to that actual prefix and preserve all31 uniform terms. Both full sums stay below2\^{}-230; support floors remain25/11. Prior R083's5400 charge and every earlier exclusion remain conservatively retained; the newly removed sector is intersected with the current remainder.

The proof uses R083 LD1--15 for the exact actual rank-three scope, eta multiplicity, LD content, selected trace and projected/source/type factors; R072 PJ4--12 for the original coordinates and general Gaussian bound; DO098 CF17--19 for the arbitrary-ideal transverse metric computation; R082 TP3--4 for the complete nonunipotent type/pair sum; and DO095 RA9 with DO074 for the actual t range. The new calculation combines the joint S-row count with the true transverse metric and sums the entire denominator tail. It does not change the initial or terminal operator, level, class input, metric or coefficient lattice.

Carried geometric anchors remain SGA3 I4.7.3/4.7.3.1; SGA6 VI Theorem3.3; FGA232 Theorem3.1 pp148--149; FGA236 Corollaries1.7/2.3 pp229/232; SGA7 I VII2.7/2.9 pp169--174 and VIII7.4(b) pp310--311; cote4 scan003 sections8.1.1--8.1.2, scans005--006 sections8.4.1--8.4.3, and cote54 scans005--007. These are the inherited module and polarization interfaces, not new motivic or lambda-ring assertions. The selected arithmetic image/content, same-input coarse-class cancellation, prescribed global period, literal family construction, signed inequality, DP44 and RH gates remain open.

\clearpage

\section{Pending arithmetic interface for the actual central Poisson route}\label{pending-arithmetic-interface-for-the-actual-central-poisson-route}

Author: /root/definite\_oscillator. 2026-11-23, before CP37. This is an exploration note, not a frozen finding, sector removal, or complete source/model review. DO105 is frozen and unchanged. Parent owns the actual holomorphic center/Cayley spectral model.

\subsection{Current actual scope}\label{current-actual-scope}

After canonical DO105 the support25 remainder is globally rank-two unipotent scalar-cross, all25 places type II and eligible, one eta per support in 2\^{}120\textless N\_E D\_eta\textless2\^{}1820, and22\textless=k\textless=164. The actual selected trace is the K058/K062/DO080 trace with its vector-valued bundle metric and both compact factors, not a scalar Bergman coefficient.

Write w=(u,eta), wdagger=(u\emph{,-bar(eta)), wdagger}w=0, and the DO105 normal form

N=w R0-V0 wdagger+d w wdagger, R0\emph{w=0, wdagger}V0=0, d+bar(d)=-R0*V0.

For real-field t, n(t)=I+i t w wdagger is unitary and n(t)h(d)=h(d+i t)=h(d)n(t). In particular the rational central fibre has a true additive group action; the real coefficient along it retains the original near-point conjugation.

\subsection{Exact finite Fourier norm lemma (proved algebraically here)}\label{exact-finite-fourier-norm-lemma-proved-algebraically-here}

Let A be a finite abelian group, U its unitary representation, V a commuting unitary operator, and P any finite-rank orthogonal projection. Put D=rank(P)\textgreater0 and f(a)=Tr(P V U(a))/D. The spectral projectors E\_chi for U commute with V, and

b\_chi=Tr(P V E\_chi)/D=Tr(E\_chi P E\_chi V)/D.

Since E\_chi P E\_chi is positive and V is unitary, \textbar b\_chi\textbar\textless=Tr(E\_chi P E\_chi)/D. Consequently sum\_chi \textbar b\_chi\textbar\textless=1. P need not commute with U or V. This applies directly to both mixed G projectors as well as the334 unmixed finite representations, once the actual arithmetic central period is fixed. It introduces no cardinality-of-period or new G-pattern factor.

The same argument is available for the full T incidence coefficient rho: realize the local group on l2 of all lattices, put P equal to the projection onto mean-zero functions supported on the actual D of sizeT-1, and use rho(h)=Tr(P U(h))/(T-2), exactly K062 AT7-11. A compact central parameter subgroup has finite orbit on this finite support, so a finite quotient suffices. The action of h0 commutes with the central subgroup. This argument controls the FULL incidence trace, not its multiplication by an arbitrary remaining-sector indicator. A proof using all-II restriction must retain that distinction.

\subsection{Arithmetic lattice and period to finish checking}\label{arithmetic-lattice-and-period-to-finish-checking}

DO104 CD14-16 gives the containing fixed-C difference lattice i J0 B\_eta I\_Sigma,F\^{}2. Without fixing a G rank pattern, the expected exact central stabilizer lattice is i B\_eta I\_Sigma,F\^{}2: outside Sigma the primitive-u Bezout extraction gives t,eta\emph{t,bar(eta)}t integral, equivalently t in B\_eta; at a fully eligible chosen label both A\_c\^{}-1\emph{w and wdagger}A\_c vanish mod pi. Their minimal valuations should each be exactly1 because A\_c and A\_c\^{}-1 have pole at most1 and w, wdagger are primitive. Then the conjugated central rank-one endomorphism has exact content pi\^{}2, allowing t in pi\^{}-2 O. Both CM components must remain explicit. Equality with the entire actual away-T maximal compact, including ramified places, is still to be reread rather than assumed.

A conservative prospective period is i Q T\^{}2 B\_eta O\_F: at G its central transvection is principal, and at every T place its entries are divisible by pi\^{}2, which fixes every actual pencil lattice. This lies in the prospective affine stabilizer lattice. Its large index should not be counted if the full finite Fourier l1 lemma is used. Exact all-II/support indicators are periodic as well, but their Fourier l1 norm needs a separate explicit bound.

\subsection{Cutoff and analytic boundary}\label{cutoff-and-analytic-boundary}

K062 AT14-20 was freshly read: the actual finite support mathscr C consists of the original away-T maximal condition and at least one actual incidence at each T place. R048 RC48.4-6 was freshly read: the entire infinite absolute actual-support tail outsideH133 is \textless2\^{}-2000. Thus a complete central extension may be justified by this tail only after proving that every added term remains in mathscr C and retains its actual finite weight. Any double charge/repartition of older removed sectors belongs to a separate accounting proof.

No spectral suppression follows merely from positive central spectrum. At the actual selected degree the spectral energy can be around2\^{}955--2\^{}960, whereas the physical central dual mesh from a representative middle denominator is only around2\^{}150. Many aliases can occur. Parent is checking the exact Nakahama action/Cayley center. No scalar-density replacement, higher-power fidelity estimate, signed positivity, or numerical saving is asserted.

\subsection{Fresh source reads in this continuation}\label{fresh-source-reads-in-this-continuation}

\begin{itemize}
\tightlist
\item
  K062 averaged\_scalar\_pencil\_trace.md, AT1-20, especially exact actual D, partial permutation, AT14 finite normalization and full infinite support.
\item
  K058 projected\_selected\_packet\_gram.md, PG1-8, actual Rgood and two mixed projectors.
\item
  DO105 EL1-3 and EL9-16, same actual row/column label cancellation, including central-only poles.
\item
  DO104 CD10-23, normal form, original primitive-u difference extraction, B\_eta descent and containing affine lattice.
\item
  K071 NB1-7, original three-type classification and paired unipotent pole table.
\item
  LIFT064 GP1-8, actual Green rank signs and the fact that mixed identity factors are1.
\end{itemize}

The proposed arithmetic lattice equality, restricted all-II Fourier weight, actual positive-frequency model and complete signed count remain to be proved. This note carries no author whole-finding PASS and no error-array mutation.

\clearpage

\textbf{Checkpoint 39 immutable pending-paper snapshot.} This contribution has not been promoted by this snapshot. Its original author status and all subsequent verification records govern its claims. The original working file may continue to change; this copy is fixed solely for faithful paper assembly.

\section{Actual low-ladder overlap: exact metric Gram reduction}\label{actual-low-ladder-overlap-exact-metric-gram-reduction}

\textbf{Research route, 2026-11-23. No new sector estimate or error charge.}

Author: ALP\_MACHINE /root/definite\_oscillator. This note retains the actual all-top J56 representation and the corrected DO079-C1 word convention. The identities below reduce the overlap problem to explicit positive Gram matrices. They do not claim that those matrices have been quantitatively controlled at the actual 168-dimensional parameters.

\subsection{1. The actual low-ladder projector}\label{the-actual-low-ladder-projector}

At one real label retain K043(PB1--8), K060(RC1--9), and DO079(SO1--3, SO18--19). Write

\[
 H_N=V_\alpha\otimes\operatorname{Sym}^N(\mathbf C^a),
 \qquad W=V_\lambda\subset H_N,
 \qquad C=\partial_a,
 \qquad B=C^{*_{H}}.
 \tag{LG1}
\]

The star is for the genuine Nakahama Hilbert norm. Coordinate a can replace the first coordinate of K060 by an actual compact permutation. Every orthogonal rank-one ladder meeting H\_N has a line of index j and lowest energy kappa, with C acting by the exact coefficient sqrt(j(j+kappa−1)). Define

\[
 M_J:=\ker(C^{J+1}|_{H_N})
      =\{F\in H_N:\deg_{z_a}F\le J\},
 \qquad 0\le J\le N.
 \tag{LG2}
\]

The equality on the right is an algebraic statement about vector-valued polynomials, before choosing an inner product. The \textbf{Hilbert orthogonal} projector Q\_J onto M\_J is exactly the sum of the orthogonal ladder-line projectors with j at most J. Indeed, C\^{}(J+1) kills precisely those lines; on every other line its coefficient is nonzero. Images from different invariant orthogonal ladders remain orthogonal, so they cannot cancel. Multiplicity spaces are retained. In particular the cumulative overlap in SO19 is

\[
 w_{\le J}=\frac{1}{d_\lambda}\operatorname{Tr}(P_\lambda Q_J),
 \qquad d_\lambda=\dim V_\lambda.
 \tag{LG3}
\]

This is an overlap trace. It is generally different from dim(W intersect M\_J)/d\_lambda. It is also generally different from the trace obtained by replacing Q\_J with the Fischer orthogonal projector onto the same polynomial subspace.

\subsection{2. An exact positive Gram matrix for the genuine projector}\label{an-exact-positive-gram-matrix-for-the-genuine-projector}

Use the Fischer norm gamma! on monomials and the original unitary norm on V\_alpha, exactly as in K043. The multiplicity-free Pieri decomposition of H\_N is orthogonal both for that norm and for the actual Hilbert norm. For every occurring compact highest weight mu, set

\[
 D_\mu=\prod_{i=0}^{a-1}
       (\alpha_i-\beta-i)_{\mu_i-\alpha_i}>0,
 \qquad
 G_N=\sum_\mu D_\mu^{-1}P_\mu^F.
 \tag{LG4}
\]

Nakahama's Theorem 5.1, in the literal K043(PB4--5) translation, says that G\_N is the positive metric operator: the Hilbert inner product is the Fischer inner product with G\_N inserted. This statement uses \textbf{all ambient Pieri components}. The scalar D\_lambda cancels between input and output for an operator already compressed to W, but it does not cancel from the metric on an unrelated subspace M\_J that meets several components.

Let V\_J be any full-column-rank matrix whose image is M\_J, written in a Fischer-orthonormal basis of H\_N. Its columns need not be orthonormal. The positive Gram matrix

\[
 A_J=V_J^*G_NV_J
 \tag{LG5}
\]

is invertible, and the genuine orthogonal projector is

\[
 Q_J=V_J A_J^{-1}V_J^*G_N.
 \tag{LG6}
\]

For proof, this operator fixes every V\_Jv, has image M\_J, and satisfies Q\_J\^{}*G\_N=G\_NQ\_J. These three identities characterize the orthogonal projection in the G\_N metric. Equivalently, the error of projecting any vector is G\_N-orthogonal to every column of V\_J. Since P\_lambda commutes with G\_N and G\_N P\_lambda=D\_lambda\^{}−1 P\_lambda, cyclicity of the finite-dimensional trace gives the exact answer

\[
 \boxed{\quad
 w_{\le J}
 =\frac{1}{D_\lambda d_\lambda}
   \operatorname{Tr}\!\left(
      A_J^{-1}V_J^*P_\lambda^F V_J\right).
 \quad}
 \tag{LG7}
\]

Every matrix in LG7 is finite, and every positive scalar D\_mu is given by the actual integer Pochhammer product. This is an exact reduction, not an asymptotic statement or an efficient numerical algorithm for the present enormous N.

\subsection{3. Reduction by the spectator compact group}\label{reduction-by-the-spectator-compact-group}

Let K' be the compact U(a−1) acting on the other positive coordinates and fixing both the distinguished positive coordinate and the negative coordinate. It commutes with the actual SU(1,1). We work in U(a,1), so this subgroup requires no determinant-one substitution. The compact restriction

\[
 V_\mu|_{K'}=\bigoplus_{\nu\prec\mu}V_\nu
 \tag{LG8}
\]

is multiplicity-free. This is the ordinary unitary branching rule; the determinant shifts in K043 do not alter the statement. The U(1) coordinate complementary to K' acts on the nu summand of V\_mu by weight \textbar mu\textbar−\textbar nu\textbar. Every ambient Pieri mu in H\_N has the same total \textbar mu\textbar=\textbar alpha\textbar+N=\textbar lambda\textbar. Therefore one can group H\_N into K' isotypic blocks

\[
 H_N=\bigoplus_\nu V_\nu\otimes\mathcal M_\nu,
 \qquad
 \mathcal M_\nu=\bigoplus_{\mu:\nu\prec\mu}\mathbf C e_\mu.
 \tag{LG9}
\]

Choose the e\_mu to be Fischer-orthonormal. On this multiplicity space G\_N is the diagonal matrix G\_nu=diag(D\_mu\^{}−1). The polynomial space M\_J is K'-invariant, so choose a full-rank matrix V\_(nu,J) for its multiplicity subspace. Put A\_(nu,J)=V\_(nu,J)\^{}\emph{G\_nu V\_(nu,J), and v\_(lambda,nu,J)=V\_(nu,J)\^{}}e\_lambda whenever nu interlaces lambda. LG7 becomes

\[
 \boxed{\quad
 p_\nu(J)=D_\lambda^{-1}
 v_{\lambda,\nu,J}^*A_{\nu,J}^{-1}v_{\lambda,\nu,J},
 \qquad
 w_{\le J}=\sum_{\nu\prec\lambda}
       \frac{\dim V_\nu}{d_\lambda}\,p_\nu(J).
 \quad}
 \tag{LG10}
\]

The value is zero if the relevant kernel multiplicity subspace is zero. The exact inequalities 0≤p\_nu(J)≤1 and p\_nu(J)≤p\_nu(J+1) follow because they are traces against nested Hilbert projections, not from an assumed monotonicity of inverse Gram entries.

There is a concrete compact-branching description of the columns in V\_(nu,J). Restrict V\_alpha to K', with summands V\_tau for tau interlacing alpha, and write a monomial as z\_a\^{}p times a polynomial of degree N−p in the spectator coordinates. For each 0≤p≤min(J,N), the relevant columns are the ordinary Pieri summands V\_nu inside V\_tau tensor Sym\^{}(N−p)(C\^{}(a−1)), when they occur. Their coefficients in the row basis e\_mu are ordinary reduced compact Clebsch--Gordan coefficients. Defining these coefficients by the Fischer-unitary inclusions fixes their normalization. No quantum-group or non-unit-normalized coefficient formula is silently substituted.

This blocks the desired calculation by compact branching types. The matrices can still be large, and no uniform bound on their inverse has yet been proved for the actual family.

\subsection{4. Retain cancellations inside each compact branching block}\label{retain-cancellations-inside-each-compact-branching-block}

On the entire K' isotypic block indexed by nu the actual rank-one compact energy is the scalar

\[
 h_\nu=N+|\lambda|-|\nu|-\beta.
 \tag{LG11}
\]

This follows from RC3: polynomial coordinate degree plus fibre coordinate weight is \textbar lambda\textbar−\textbar nu\textbar. Set p\_nu(−1)=0 and q\_(nu,j)=p\_nu(j)−p\_nu(j−1). Then q\_(nu,j)≥0, sum\_j q\_(nu,j)=1, and the lowest energy on an occurring j line is kappa=h\_nu−2j\textgreater1. Thus the scalar by which the genuine selected radial compression acts on V\_nu is exactly

\[
 d_\nu(r)=\sum_j q_{\nu,j}
  \cosh(r)^{-(h_\nu-2j)}
  P_j^{(0,h_\nu-2j-1)}(1-2\tanh^2r).
 \tag{LG12}
\]

To justify scalarity, the radial element commutes with K', as do both degree and selected projections. Schur's lemma applies because W restricted to K' is multiplicity-free. To justify the formula, insert the orthogonal ladder decomposition in that block. The squared overlap of the one selected multiplicity line with the j eigenspace is exactly q\_(nu,j) from LG10.

For an actual Cartan decomposition g=k\_1 a\_r k\_2, retain U=pi\_W(k\_2k\_1). The safe bound is

\[
 \boxed{\quad
 |c_W(g)|\le \frac{\|P_W\pi(a_r)P_W\|_1}{d_\lambda}
 =\sum_{\nu\prec\lambda}
   \frac{\dim V_\nu}{d_\lambda}|d_\nu(r)|
 \le\sum_{\nu,j}\frac{\dim V_\nu}{d_\lambda}
       q_{\nu,j}|d_{h_\nu-2j,j}(r)|.
 \quad}
 \tag{LG13}
\]

The first inequality is the trace-norm inequality \textbar Tr(UA)\textbar≤\textbar\textbar A\textbar\textbar\_1 for the actual unitary U. It does not replace U by the identity. The middle expression retains cancellations between ladders inside one K' type; the last expression is the grouped SO19 bound. The first expression remains valid even when U mixes K' types.

The original Haagerup--Schlichtkrull bound used in K060 depends on 2j+kappa=h\_nu. Merely deleting a small-j part does not improve that same bound on the rest. A quantitative gain from LG10 requires either a separately justified j-sensitive Jacobi estimate in the actual varying-parameter range, or an estimate of the signed combinations LG12 before applying the trace norm in LG13. Neither is asserted here.

\subsection{5. A small exact model detects the metric issue}\label{a-small-exact-model-detects-the-metric-issue}

This example is a scope diagnostic, not a replacement of the actual a=168 representation. Take a=2, fibre V\_(1,0), degree N=1, and scalar parameter kappa\textgreater2. In the compact weight-one plane use the Fischer-orthonormal vectors

\[
 e_1=z_1v_2,\qquad e_2=z_2v_1.
\]

The symmetric and alternating Pieri lines are (e\_1+e\_2)/sqrt(2) and (e\_1−e\_2)/sqrt(2). Their exact PB1 factors are D\_+=kappa+1 and D\_−=kappa−1. Consequently

\[
 G=\frac{1}{\kappa^2-1}
       \begin{pmatrix}\kappa&-1\\-1&\kappa\end{pmatrix},
 \qquad M_0=\mathbf C e_2.
 \tag{LG14}
\]

The squared genuine overlap of this kernel line with the selected symmetric line is (kappa−1)/(2kappa); for the alternating line it is (kappa+1)/(2kappa). Both Fischer overlaps equal 1/2. Therefore there is no general equality of the two overlap notions, and there is not even a type-independent domination in one direction. At kappa=3 the two genuine values are 1/3 and 2/3. This is an exact consequence of the same PB1 metric and LG6.

On the entire symmetric compact type of dimension three, the lowest-coordinate weight line belongs to M\_0 and the middle weight contributes the preceding overlap. The full cumulative j=0 weight is (1+(kappa−1)/(2kappa))/3, which at kappa=3 is 4/9. Its minimum-coordinate subspace ratio is only 1/3. Thus the existence and normalized dimension of the minimum-coordinate subspace do not determine the cumulative overlap, even in this elementary genuine holomorphic model.

\subsection{6. What is proved, and what remains}\label{what-is-proved-and-what-remains}

LG2--13 are exact linear-algebra and representation-theoretic consequences of the retained actual model and ordinary compact branching. LG14 supplies an exact small-model check of why its metric data are necessary. The companion diagnostic verifies these projection and block identities in explicit finite models and checks the lowest-ladder kernel identity. It does not compute the actual huge Gram matrices, prove a new selected-trace decay rate, remove a support sector, or change the 101/31 error list after canonical DO108.

The remaining quantitative problem is to estimate the explicit inverse positive Gram matrices in LG10, or to estimate the corresponding signed radial combinations LG12. Replacing those matrices by coordinate-weight dimensions or by a scalar Bergman model is not justified. The old bounded-j obstruction and the corrected word formulas in DO079-C1 remain unchanged.

\subsection{Sources and source-check boundary}\label{sources-and-source-check-boundary}

Nakahama, R. (2016). Norm computation and analytic continuation of vector valued holomorphic discrete series representations. \emph{Journal of Lie Theory, 26}(4), 927--990. Theorem 5.1, printed pages 960--961, and equation (5.5), printed page 959. The complete relevant local primary text was directly reread in research/nakahama2016.txt, together with the exact K043(PB1--8) dictionary.

The compact branching rule in LG8 is the same ordinary rule already used and checked in DO079(SO15). Its application here is only to finite-dimensional unitary modules, with all multiplicities in H\_N retained. The actual orthogonal ladders are K060(RC1--9), and the exact trace-norm convention is DO079(SO18--19), both freshly read in full at the relevant passages. DO079-C1 was freshly read in full and remains mandatory for actual all-word weights.

A new primary lead was directly opened: Sun, Y. (2016). A representation-theoretic proof of the branching rule for Macdonald polynomials. \emph{Mathematical Research Letters, 23}(3), 887--927, section 6.1 and Theorem 6.1, printed pages 918--919. It records explicit multiplicity-free reduced Clebsch--Gordan coefficients. No numerical coefficient from that quantum-group formula is used here: its classical limit and Fischer-unitary normalization have not been reconciled for this task. The exact coefficients in section 3 are defined by the unitary compact inclusions instead.

\clearpage

\section{Dense all-II support27: higher-moment route checkpoint}\label{dense-all-ii-support27-higher-moment-route-checkpoint}

Status: exploratory route analysis, 23 November2026 UTC. This is not a containing-sector theorem, exact endpoint certificate, impossibility proof or new error charge. Current canonical predecessor is K114126/31 with complete floors27/11.

The actual K113 remaining all-II branch has b=j=0 and either total eligibility27,100\textless=k\textless=164,2\^{}2111\textless x\textless2\^{}3328, or total eligibility28,k164,2\^{}3107\textless x\textless2\^{}3328. All source, phase, selected-trace, paired ideal and central-fibre hypotheses remain those of canonical K113. No support26 window is imported.

The previously saved \nolinkurl{support_twenty_seven_dense_envelope_failed_probe} remains unchanged. At the lower eligibility27 denominator endpoint, q* is about2.51e8 and the scalar volume term divided by K15 is about2\^{}768. Thus improving only the universal spherical-code coefficient does not substantially improve that volume-dominated endpoint. Its large positive upper bound is a failure of that sufficient estimate, not a lower bound on the actual sum.

A new analytic possibility is to retain the whole scalar polynomial inside K110's positive mixture, together with the genuine per-embedding cap alpha\_tau\textless=R=17424/133. Pulling its S-factor out at b=2 and retaining only M3 loses higher inverse moments in this range. The original positive measure in K110 PM4 has, for an integer j with0\textless2j\textless nu=333/2, \[
\int b^{-j}\,d\mu_\nu(b)
=4^j\nu\,\frac{\Gamma(\nu-2j)}{\Gamma(\nu-j+1)}
=\frac{4^j\nu}{\prod_{l=0}^{j}(\nu-2j+l)}.
\] This follows by the same positive Tonelli integral and substitution alpha=4t/(1-t)\^{}2 used in PM7. No analytic continuation is permitted: for j\textgreater=84 the original inverse moment diverges. The actual full row degrees reach167, so those moments cannot be inserted without retaining the original energy cap or a fully paid lower-b majorant.

For the existing modified measure, its low-b atom2\^{}-39 at b2 dominates high moments and limits the improvement. An exploratory exact-rational scan of the convergent moments has its smallest value of log2(2\^{}j M\_j) near j69 (about-198.27), but this is not a full count bound. The original low-b atom is a bound after multiplication by exp(2R), so a small raw low-b probability is not sufficient to delete it.

One possible complete adaptation is to choose a new lower atom b0\textgreater0 and pay exp(b0 R)Psi(R) on the original alpha range; scalar powers would then use the full finite expansion and higher truncated moments. Every resulting positive term would need its own exact anisotropic scalar minimum, actual transverse determinant, central correlation, finite-pattern product and q/x monotonicity. The bounded actual nonzero-row norm floor must be applied before any Gaussian enlargement. No such combined theorem or complete endpoint scan has yet been proved here.

The finite-G rank-two coefficient cannot simply be replaced by a rank-three baseline: a globally rank-two matrix has finite rank at most2. Likewise the exact central rank-zero correlation is already present in K113. Any further G or central saving needs a new proved algebraic/congruence statement. Work is checkpointed here before the parent-assigned independent LIFT096 review.

\clearpage

\section{Finite split-component boundary for positive compression}\label{finite-split-component-boundary-for-positive-compression}

Status: exact scoped counterexample, recorded 23 November 2026 UTC by \nolinkurl{/root/kummer_extension}. No numerical error charge and no global-arrow realization are asserted.

Use coordinates S1, V2, P2, with split Hermitian Gram J=diag(I3,J\_H), where J\_H has off-diagonal entries1 and diagonal entries0. The selected component is \[
h=\begin{pmatrix}
1&0&0&1&0\\
0&1&0&1&0\\
0&0&2&1&0\\
1&-1&-1/2&1&1\\
0&0&0&-1/2&0
\end{pmatrix},\qquad
h^{-1}=\begin{pmatrix}
1&0&0&0&2\\
0&1&0&0&2\\
0&0&1/2&0&1\\
0&0&0&0&-2\\
-1&1&1/4&1&5/2
\end{pmatrix}.
\]

Direct multiplication gives these inverse identities. The matrix N=h-I satisfies N\^{}5=0 and N\^{}4 nonzero; its rank is4 at every tested odd prime, including both actual G characteristics. Thus h is unipotent, with characteristic polynomial (X-1)\^{}5, while its positive compression is diag(1,1,2). The corresponding inverse compression is diag(1,1,1/2). In particular A=diag(1,2) is not unipotent. Both S/V off-diagonal blocks vanish for h and its inverse.

The entire lower cross blocks are respectively \[
C=\begin{pmatrix}1&-1&-1/2\\0&0&0\end{pmatrix},\qquad
C_{\rm inv}=\begin{pmatrix}0&0&0\\-1&1&1/4\end{pmatrix}.
\] Their original S restrictions are nonzero, and their full images are the two fixed null lines of J\_H. Define the companion CM component to be J\^{}\{-1\}h\^{}\{-t\}J. Then the pair is genuinely unitary over the split quadratic algebra with the component-swapping involution. The companion has the reciprocal positive compression and the same stated null-line properties. No positive-definite form over the residue field is used.

One may pad V by an identity coordinate and S by163 identity coordinates to match the S164/V3/P2 block sizes. A split Hermitian change of P basis may identify its nondegenerate plane pairing with the original local plane pairing; this does not assert preservation of the actual pencils, residue labels or global integral coordinates. The example establishes only that full finite unipotence, paired inverse-positive compression and the displayed null-cross structural hypotheses do not by themselves force A unipotent. It does not construct a global source arrow, satisfy all actual finite incidence conditions or refute a stronger low-rank lemma. Its full finite rank is4; ranks at most3 are separate.

The companion exact script verifies the inverse, nilpotence, both split-Hermitian equations, all positive and cross blocks, and finite ranks at3,5,7,11,13 and the two actual G characteristics. It passed39 checks at2026-11-23T08:33:40.024224+00:00. These finite checks support the displayed rational identities, which give the example independently. No canonical file was changed.

\clearpage

\section{Scoped assessment: theta constants do not supply h=1 on the required compactification}\label{scoped-assessment-theta-constants-do-not-supply-h1-on-the-required-compactification}

23 November 2026 UTC. Author: \nolinkurl{/root/lift_accounting}.

This is an approach note, not a new theorem, finite charge or main-gate movement. It explains why the explicit exponent in AM-LIFT-097 is obtained from log-canonical very ampleness rather than from an unverified theta-coordinate embedding at the boundary. Frozen LIFT096 is unchanged.

The relevant line in LIFT096 is the pullback of the Siegel canonical Baily--Borel line. The escape argument needs a genuinely very ample line on the entire normal projective compactification, including boundary points. Immersion of the open moduli space, or a finite normalization map defined by theta constants, does not imply that condition.

Freitag and Salvati Manni (2016), Introduction, printed pp.~1--2, distinguish these assertions for the map defined by even theta constants on the Satake compactification at level (4,8). They state injectivity and the open embedding property, but also the failure of the corresponding global isomorphism onto its image in genera at least six. Their result for genus three does not cover the present family of abelian dimension g=fr. The associated ring normalization and point separation therefore do not justify replacing the auxiliary very-ampleness exponent by one. No claim is made that every possible theta-based projective embedding fails.

The complete AM-LIFT-097 route instead uses the actual canonical Cartier divisor on the neat Baily--Borel component, positivity of 2A-K\_X on every log-canonical center, and the published effective freeness and very-ampleness results of Fujino. Its explicit exponent and the resulting h\^{}d factor are retained in the degree cost. The total intersection v\_A, the actual same-rule nonempty-family condition and component-entry multiplier remain separate inputs.

Reference: Freitag, E., \& Salvati Manni, R. (2016). \emph{On the variety associated to the ring of theta constants in genus 3} (arXiv:1606.04468). \url{https://arxiv.org/pdf/1606.04468}. Fresh direct reading: abstract and full introduction, printed pp.~1--2; no claim of a whole-paper review.

\clearpage

\section{Terminal frames, branching, and the original selected map}\label{terminal-frames-branching-and-the-original-selected-map}

\textbf{Status:} source-only research clarification, following the immutable 09:48:10 UTC snapshot. This is not a new nonvanishing theorem, finding, arithmetic charge, or claim of source admission. Canonical files and the original snapshot are unchanged.

\subsection{1. The abstract orbit statement and its missing application}\label{the-abstract-orbit-statement-and-its-missing-application}

For a finite-dimensional irreducible representation \(V\) of a group \(G\), every nonzero linear functional \(b\in V^\vee\) has \[
\bigcap_{g\in G}\ker(b\rho(g))=0.
\] Indeed this intersection is a \(G\)-stable subspace contained in the proper subspace \(\ker b\). The assertion applies equally to a nonzero linear map by composing with a suitable output functional. Applying it to the present source would require \textbf{actual permitted maps} obtained by inserting \(\rho(g)\) before the fixed terminal branch, while keeping the original coefficient representation, finite tensors, shared middle covector contraction, class masses, terminal extension and source pins unchanged. Neither the irreducibility argument nor freedom to choose coordinates establishes that requirement.

If instead a change of coordinates is represented by an isomorphism \(J_g\), the same intrinsic branch has coordinate expression \(bJ_g^{-1}\), while the same vector becomes \(J_gv\). Thus \((bJ_g^{-1})(J_gv)=b(v)\). Taking all such coordinate descriptions does not enlarge the detecting family. This elementary identity is the relevant distinction; it is not a theorem forbidding every possible new geometric construction.

\subsection{2. The literal frame conversion preserves the selected map}\label{the-literal-frame-conversion-preserves-the-selected-map}

The file \nolinkurl{research/alpoge_source/Development/AbelianVarietyTheory/Definitions/LatticeFrames/OnFramesRdApply.lean} supplies the following exact interfaces.

\begin{itemize}
\tightlist
\item
  \nolinkurl{RealD}, lines 24--45, includes the chosen frame identifications and a separate terminal \nolinkurl{low} map. \nolinkurl{RealD.Is.low}, lines 62--65, requires the prescribed interlacing and subgroup intertwining and that \nolinkurl{low} is nonzero on the ambient coefficient space. It does not assert nonzero restriction to the actual theta image.
\item
  \nolinkurl{ClsK.Kval}, lines 115--121, already quantifies over \textbf{all} output evaluation frames. The class-integrated theta output is identified as a complete frame function; introducing another evaluation frame is not a new theta operator.
\item
  \nolinkurl{ClsT.Is.tail}, lines 147--151, applies \nolinkurl{low} to that entire frame function, then applies the terminal functional and the original class sum with its stabilizer and volume factors. The ordering is explicit.
\item
  \nolinkurl{Tup.onFrames}, lines 223--249, copies \nolinkurl{om}, \nolinkurl{pt}, \nolinkurl{CC}, \nolinkurl{CT}, \nolinkurl{kom}, \nolinkurl{Ch}, \nolinkurl{dd}, \nolinkurl{e}, \nolinkurl{Col}, \nolinkurl{Ram}, \nolinkurl{S}, \nolinkurl{K}, \nolinkurl{tail}, and all displayed scalar constants unchanged. Only \nolinkurl{rd} is pulled back along the inverse frame identification.
\item
  \nolinkurl{Tup.onFrames_rd_apply}, lines 261--263, is the exact equality of evaluations at corresponding coordinates. \nolinkurl{KerX.of_onFrames}, \nolinkurl{PerX.of_onFrames}, and \nolinkurl{RdX.of_onFrames}, lines 274--298, transport the existing predicates. \nolinkurl{RealD.onFrames}, lines 303--304, retains the same \nolinkurl{low}.
\end{itemize}

Consequently the literal \nolinkurl{onFrames} conversion proves a coordinate identification of the same source family. It supplies no operation that replaces the actual terminal map by \nolinkurl{low} composed with an arbitrary independent full-rank group action while holding the remaining source data fixed.

\subsection{3. Physical eigenframes and finite decoration frames are distinct}\label{physical-eigenframes-and-finite-decoration-frames-are-distinct}

\nolinkurl{research/alpoge_source/Formalism/LocalToGlobal/ProjectiveLattices/IsCombAt.lean}, \nolinkurl{FrPt} at lines 67--77, permits a complete physical eigentangent frame and a nonzero generator of the theta line on a moduli object. These choices are genuine variables of a frame rule. They are not the finite integral \nolinkurl{Frs} or \nolinkurl{liftFr} choices in the arithmetic class construction. The latter change coordinate descriptions of fixed finite decorations subject to the exact level and character constraints.

This note does \textbf{not} infer that every physical eigenframe variation is disallowed. Rather, to use it to detect the original branched image one still needs an exact source equality identifying the resulting arithmetic rule with \nolinkurl{low} after an independently varied action. The already checked frame covariance is insufficient for that equality. The full evaluation-frame dual span in K037 CR16 is already used in the unbranched theta identification; it does not move an action past \nolinkurl{low}.

For comparison, K037 CR38--39 explicitly transports both the coefficient realization and the oscillator insertion, with their determinant factors. K038 CA24--30 transports the full finite tensor and preserves the complete class-sum pairing. Those formulas cannot be used to rotate only the terminal branch and silently leave their other factors fixed.

\subsection{4. The fixed arithmetic terminal construction}\label{the-fixed-arithmetic-terminal-construction}

The file \nolinkurl{research/alpoge_source/Formalism/LocalToGlobal/Columns/FiniteSumNonzeroLargeDepth.lean} defines \nolinkurl{brC} at lines 197--198 from the actual interlacing branch and the conjugate of the prescribed \nolinkurl{piT}. It defines \nolinkurl{tilQ} at lines 204--205 from \nolinkurl{DecoSet.tilH}, \nolinkurl{evzC} at lines 212--213 by restriction along that extension, and \nolinkurl{PtilFF} at lines 215--218 as the original class functional applied to the coefficient pairing after \nolinkurl{brC}. \nolinkurl{PtilC}, lines 220--222, composes precisely those maps with the original column and terminal modules. \nolinkurl{StdLatUnimodular.lean}, lines 165--176 and 700 onward, fixes the extended lattice by the actual added norm-one line and its level covector; it does not supply arbitrary rational lines in that slot.

The allowed finite level-group normalizers and terminal block-equivariance formulas therefore provide covariance of the prescribed map. A full general linear orbit before \nolinkurl{brC} is an additional construction, not a consequence of these definitions. The actual full middle coevaluation and shared covector/class sum must remain in any proposed construction.

DO039 gives a separate, already checked limitation for the standard integral lattice in its stated old J14/J56 realization: its finitely many integral norm-one line branches have a nonzero common kernel on the full ambient coefficient space. This does not show that the actual global theta image lies in that kernel, does not cover arbitrary lattice classes, and does not prove vanishing of the original selected map. It only prevents replacing the missing source argument by a universal ambient separation assertion for those integral lines.

\subsection{5. What remains open and the next distinct interface}\label{what-remains-open-and-the-next-distinct-interface}

The raw two-edge global theta injectivity of LIFT27 MR25--31 remains upstream of the selected terminal restriction; MR32 states that boundary. K059 IS24 already gives an exact finite-generator test for the original selected map. Rephrasing that test as a finite evaluation matrix is not a new nonvanishing proof.

A separate literal source family is \nolinkurl{thGen} in \nolinkurl{ProjectiveLattices/IsThRule.lean}, lines 89--91. It allows permitted level tuples and harmonic polynomials; \nolinkurl{ThNe0} in \nolinkurl{ThProps.lean}, lines 30--32, asks for a nonzero rule when \nolinkurl{DF + 2 ≤ M}, together with its other hypotheses. \nolinkurl{HasDataCondition.lean}, line 53, carries this as an explicit \nolinkurl{LawsP.thNe} field. A geometric proof for that permitted family might fill that specific source law, but it must not be identified with the original coupled selected image without a further equality. The theorem \nolinkurl{recPJ_thAll_ne_bot}, lines 80--87, consumes that field rather than constructing it.

The present conclusion is therefore limited and exact: \textbf{the verified frame-coordinate conversion supplies no new independent branch orbit.} The physical-frame-to-original-selected-rule equality, actual nonzero selected rule, admitted complete source, literal Read/032 identification and DP44 remain open. No impossibility claim for a different source-sanctioned detector is made.

\subsection{Verification scope}\label{verification-scope-5}

The cited source bodies, the complete \nolinkurl{IsThRule} and \nolinkurl{ThProps} files, and the named canonical report portions were read directly during this investigation. No Lean build was run. The companion metadata records exact source and report fingerprints and the completion timestamp; these are read/provenance records, not an executable proof certificate.

\clearpage

\section{AM-K-121 --- exact narrow-band height curvature (pending snapshot)}\label{am-k-121-exact-narrow-band-height-curvature-pending-snapshot}

Owner: \nolinkurl{/root/kummer_extension}. 23 November 2026 UTC. \textbf{Exploratory mathematical snapshot; not author-frozen, independently reviewed, or an arithmetic charge.} The frozen R099/R100/K120 and DO124 subjects are unchanged.

\subsection{1. Exact retained factor from the same angular kernel}\label{exact-retained-factor-from-the-same-angular-kernel}

Use the actual selected fidelity and both compact factors of DO080, the narrow band I={[}16/529,16/527{]}, and R099 AH8. Write nu=333/2, c=cosh(r)\textgreater1, s=sinh(r), D=sqrt(y/y'). The exact directional height is

K=1-(c-D)\^{}2/{[}(c\^{}2-1)(1-y){]}.

For c\textgreater=c0=529/527 the full R100 NB7 proof, before its last relaxation, gives

1-K/{[}2D/(c+1){]} \textgreater= a(c-c0), a=527/33858.

On feasible positive-height points, 0\textless=a(c-c0)\textless1. Therefore the directional density retains the full factor

L(c)={[}1-a(c-c0){]}\_+\^{}nu,

rather than replacing it by an exponential. If the height is nonpositive, the density is zero. For 1\textless c\textless c0 the original two-case height estimate suffices because L(c)\textgreater1. Thus the same L(c) bounds both forward and reverse directions for every positive radius, including the cases c\textless=D or c\textless=D\^{}(-1) through the original two-case estimate. Detailed balance cancels D and retains one entire L(c), not its square.

Using the old Cartesian Jacobian and the same full radial spread integral gives the candidate exact retained-band estimate

F\_I(g) \textless= vartheta 2\^{}(-933/2) s\^{}(-1) {[}2/(c+1){]}\^{}nu L(c).

This is derived within one angular kernel. It is not a product of distinct trace estimates. The actual trace outside the band still requires K120's independently paid full arithmetic tail.

\subsection{2. Actual energy factor and finite support}\label{actual-energy-factor-and-finite-support}

At an actual rank-two-unipotent near point, c=sqrt((1+t/2)\^{}2+w\^{}2)\textgreater=1+t/2. Since L decreases,

L(c) \textless= A(t)={[}A0-bt{]}\_+\^{}nu, A0=16930/16929, b=527/67716, T0=A0/b=67720/527, lambda0=nu/T0=175491/135440.

Consequently the same original joint central profile can be bounded by

t\^{}(-1/2)(1+t/4)\^{}(-167) {[}1+w\^{}2/(4(1+t/4)\^{}2){]}\^{}(-333/4) A(t).

Original central integration would leave shape166 times A(t). Retaining c itself rather than this energy-only relaxation could improve that integral further; this is not yet quantified. The finite support T0 is below the original energy cap R=17424/133, so the vanishing part must be handled as genuine zero support.

Direct relative near-point transport of A(t) is NOT yet justified: its logarithmic derivative diverges at T0. A later proof must transport the feasible support itself or use a globally smooth containing majorant before applying the existing near-point argument. No uniform factor2 for A(t) is asserted here.

\subsection{3. Reusable globally exponential tangent majorants}\label{reusable-globally-exponential-tangent-majorants}

For any lambda\textgreater=lambda0, concavity of log(1-t/T0) on its positive domain gives

A(t) \textless= exp(C\_lambda-lambda\emph{t), t\textgreater0, C\_lambda=nu}log(nu\emph{b/lambda)+lambda}T0-nu.

For t\textgreater=T0 the left side is zero, so the same inequality remains valid. This is a global positive exponential majorant; its intercept must be paid across all168 embeddings. For a fixed rational lambda, its joint profile has the original R099 derivative bound with lambda replacing32/25. Thus the same original-row near-point transport can be re-established with the complete intercept retained.

The shifted positive half-Gamma/Gamma166 representation has infinite total mass. Every positive inverse moment is finite, but the sole row-polynomial degree-zero term still needs the original theta=2\^{}-20 energy reservation before Gaussian enlargement. All finite patterns, physical U\^{}f central covolume, full750 cover, q/x derivatives, and original arithmetic-tail payment remain obligations of any application.

\subsection{4. Exact paired-height algebra worth retaining}\label{exact-paired-height-algebra-worth-retaining}

Put x=sqrt(y), z=sqrt(y'). Where both directional first-case heights are positive, their normalized geometric mean is exactly

sqrt\{ {[}K\_f/(2D/(c+1)){]} {[}K\_r/(2D\^{}(-1)/(c+1)){]} \} = {[}2c\emph{x}z-(c\^{}2-1)\emph{x\^{}2}z\^{}2-x\^{}2-z\^{}2{]} /{[}2\emph{x}z\emph{(c-1)}sqrt((1-x\^{}2)(1-z\^{}2)){]}.

With p=sqrt(y/(1-y)), q=sqrt(y'/(1-y')), w0=p*q and e=p/q+q/p-2\textgreater=0, it is

{[}2c\emph{sqrt((1+w0)\^{}2+w0}e)-(2+e)-(c\^{}2+1)*w0{]}/{[}2(c-1){]}.

For c\emph{w0\textless=1+w0 this is at most 1-(c-1)}w0/2, because the difference from its e=0 value has numerator

e*{[}2c\emph{w0/(sqrt((1+w0)\^{}2+w0}e)+1+w0)-1{]} \textless=0.

In particular this diagonal comparison holds uniformly on the narrow band when c\textless=527/16. The remaining larger-c region requires a separate estimate; no all-c diagonal-maximizer claim is made.

\subsection{5. Floating scale check, not a sector conclusion}\label{floating-scale-check-not-a-sector-conclusion}

\nolinkurl{exact_height_curvature_initial_probe.json} records an explicitly floating quadrature at degree j147 after shape166 central integration. It compares with the exact exp(-32t/25) integral, not the larger certified K120 rational table. The full A(t) factor improves that moment-product comparison by about300.11 bits. Tangent lambda29/20 gives about294.47 bits after paying its entire approximately267.63-bit field intercept. These figures are only diagnostics and do not establish a certified moment or complete count.

DO124's complete k144 sufficient envelope is +924757/1024, retaining all579 nonmain terms. Thus this first curvature estimate alone does not visibly close the required230-bit budget. That failed sufficient comparison is not a lower bound for the actual arithmetic sum and does not exclude a stronger same-kernel central or paired-height argument.

The checkpoint snapshot preserves the exact derivation and its unresolved transport/count obligations. No standalone charge, support-floor change, or k144 removal is claimed.

\clearpage

\section{Pending research: the cover rule and the remaining original-base descent}\label{pending-research-the-cover-rule-and-the-remaining-original-base-descent}

Saved \nolinkurl{2026-11-23T12:00:26.833067+00:00}. \textbf{Research snapshot only: no new finding ID, no final author/script/peer check, no canonical claim.} LIFT103 remains frozen independently.

\subsection{Concrete next deduction under consideration}\label{concrete-next-deduction-under-consideration}

For the actual cover stage, restrict the unique all-ring LIFT103 value to the literal TRing points. This gives a Rules object. Its arbitrary-map naturality supplies IsNatural. For P in HHa, the expected exact FrBor law is

v(y') = (t\_none)\^{}n * product\_(tau,i) t\_(tau,i)\^{}(a(tau,i)) * v(y).

The physical frame in IsCombAt.FrBor is lower triangular: u'\emph{i=sum}(j\textless=i) C\_ij u\_j, diagonal t\_i. Thus derivative substitution is P(X\_i -\textgreater{} sum\_j C\_ij X\_j), not its inverse. RowInv, RowDeg and the decomposition C=D U predict exactly the stated positive character. RowInv is initially quantified over the infinite coefficient field; its polynomial identity extends to arbitrary algebra coefficients before applying a triangular matrix with unit diagonal. Theta-generator scaling is linear in each of the n factors, hence exponent n, independently of N\^{}2. Balanced quadratic gauge changes remain killed by the original all-column harmonic condition. The exact proof and finite diagnostics are not yet written.

The resulting cover rule should therefore belong to sec(Xo2,c2,a,n). Multiplication by actual cdelta in sec(0,0), provided by IsLev, preserves this weight. This constructs a cover-stage rule only; it does not identify a rule on the original base.

\subsection{Exact FromTest supplier}\label{exact-fromtest-supplier}

AlgebraMapCAt.lean51--58 defines FromTest and evalAt: outside FromTest, evalAt is zero. ThShort.lean181--189's IsLev requires evR(cdelta)(y2) nonzero at every algebraically closed cover point. Hence each such y2 is FromTest. This gives the exact supplier needed to identify the cover rule's evR with the unique all-ring scalar by LIFT103 naturality, without assuming an additional AllTest clause for the cover moduli. This is an elementary deduction from the actual definitions; it does not prove original-base descent.

\subsection{Remaining descent issue}\label{remaining-descent-issue}

IsShortRule quantifies over every original algebraically closed y and every y2 with FrOver. To produce f in original sec, one still needs cdelta(y2)*v(y2) to be independent of the compatible lift and to descend as a rule with the exact original character. IsLev supplies existence of geometric lifts and the cover-stage scalar rule, but no direct all-ring descent equality.

PhiPin gives a full finite Fourier formula for phi\_delta and a dual-character covariance. PsiTie gives cdelta covariance under specified LevAut automorphisms, restricted by g.InGam and trivial action modulo l2.A. The fresh reading has not yet established that these automorphisms cover every pair of FrOver lifts. In particular IsLev permits l2.A to enlarge l.A, while PsiTie explicitly imposes triviality modulo l2.A; that potential scope mismatch must be resolved from the actual Sim/LevAut definitions and covering geometry, not assumed away. No counterexample to the full genuine package is asserted.

The source Moduli.rep is quantified over reduced bases. It may help prove the requisite finite level cover and its actual deck action, but it is not a supplied arbitrary-nonreduced representability theorem. The original Rules objects use reduced finite-type TRing points, while arbitrary-ring scalar existence in LIFT103 has the stronger separate scope. Keep these roles distinct.

No actual package, selected arithmetic nonzero, integral coefficient descent, cusp or spanning theorem is constructed by this snapshot. No numerical error charge or canonical changes.

\subsection{Fresh source readings}\label{fresh-source-readings}

ThShort112--235 and299--355; complete EmbKCoe35, RowDeg50, IsCombAt127 and HasDataCondition83; AlgebraMapCAt1--115; complete available PushInvertible containing Moduli.rep. The exact current file hashes below identify the read bytes. Sim/InGam/LevAut definitions were not found in the currently searched source directories and remain a targeted retrieval task, not a presumed theorem.

\begin{itemize}
\tightlist
\item
  \nolinkurl{research/alpoge_source/Formalism/LocalToGlobal/ProjectiveLattices/ThShort.lean} --- SHA \nolinkurl{910822e9773ea37f97728e494db9fffac2e9707478637452fa38b1c5e88df3b2}.
\item
  \nolinkurl{research/alpoge_source/Formalism/LocalToGlobal/ProjectiveLattices/IsCombAt.lean} --- SHA \nolinkurl{96c9bfeb8c6e48d0847d85ce6bfb3c864dc70f894f5a8b637959c3f9accacbd9}.
\item
  \nolinkurl{research/alpoge_source/Formalism/LocalToGlobal/ProjectiveLattices/AlgebraMapCAt.lean} --- SHA \nolinkurl{9bd5d736d09e07e1ae9f6f40433f7825208226cd2e1f455ce1f49c6ff02c9ac4}.
\item
  \nolinkurl{research/alpoge_source/Formalism/LocalToGlobal/ProjectiveLattices/EmbKCoe.lean} --- SHA \nolinkurl{e986226fd1bfca169529b42b3faf4231baf37eb0e9ea00592aa4f63f12d93212}.
\item
  \nolinkurl{research/alpoge_source/Formalism/LocalToGlobal/ProjectiveLattices/RowDeg.lean} --- SHA \nolinkurl{5bfd910715a5fb61ac23df068f6e10106fa14c042cb3810fc1fc04003150ba9f}.
\item
  \nolinkurl{research/alpoge_source/Formalism/LocalToGlobal/ProjectiveLattices/HasDataCondition.lean} --- SHA \nolinkurl{87ecf2483c1a5437dc46bc9b363a67494f19533b227898e21b6c0124115526fd}.
\item
  \nolinkurl{recovery_pending/lift/source_cm/ReciprocityTheory/ReductiveGroups/ArithmeticGeometry/BirationalGeometry/ToroidalCompactification/PushInvertible.lean} --- SHA \nolinkurl{f3400474c76aa3eeaa0cb85edf2e9c549811616ee2d43fd74685088a42662554}.
\end{itemize}

\subsection{Targeted source retrieval after the initial snapshot}\label{targeted-source-retrieval-after-the-initial-snapshot}

The pinned HeckeKer.lean has now been read in full and saved in short\_descent\_sources. It supplies Sim, Sim.InGam, Moduli.LevAut and the explicit ToroidalLevelCompatibility statement. That statement gives a finite etale surjective coefficient-field cover, geometric-fibre transitivity by g.InGam, and exact stabilizer g.TrivAt. It is a Prop-valued source interface, not by itself a proved supplier or a theorem applying to an arbitrary stored DeltaDat.r. The exact N1 supplier and comparison with the stored r/m still require tracing.

Sim.InGam imposes the old full-A identity only, whereas PsiTie adds the new full-A identity. The discrepancy is literal. If l2.A=l.A, the additional condition follows at once. If the full-A level grows, no complete deck-invariance theorem has been obtained. This source distinction is not a counterexample to a specially constructed actual DeltaDat, and choosing the unchanged full-A level is a potential constructive route.

\clearpage

\section{High-coherence phase geometry --- pending K123 route}\label{high-coherence-phase-geometry-pending-k123-route}

Owner: \nolinkurl{/root/kummer_extension}. This is a working derivation and explicitly not a frozen theorem, numerical sector removal or canonical edit. K122 is accepted; the stronger threshold and geometric comparison below still need their complete exact certificate and full source/normalization review. The floating values are diagnostics only.

\subsection{Exact candidate reduction}\label{exact-candidate-reduction}

For a noncentral globally rank-two unipotent, simultaneous compact conjugation sends its fixed null line and nonzero horizontal row to the first two positive coordinates. With t\textgreater0, v=sqrt(t), A=t/2 and physical central coordinate w, the inverse can be written on those coordinates and the negative coordinate as

\[
g=\begin{pmatrix}1-A-iw&-v&A+iw\\v&1&-v\\-A-iw&-v&1+A+iw\end{pmatrix}.
\]

All other positive coordinates are fixed. This canonical-form/source bridge must be proved from the actual normal form before final use. Write d=1+A+iw, c=\textbar d\textbar{} and s²=c²−1. At z in the ball, rho=d−(A+iw)z1−v z2 and

\[
z'_1=1-(1-z_1)/\rho,\quad z'_2=(z_2-v(1-z_1))/\rho,\quad z'_\perp=z_\perp/\rho.
\]

Fix y=1−\textbar z\textbar², y'=1−\textbar z'\textbar², D=sqrt(y/y'), R²=1−y and rho=D exp(i(arg d+theta)). The original denominator constraint gives

\[
z_0=(\rho-d)(-A+iw,-v)/s^2,\quad e=(v,-A-iw)/s,\quad
(z_1,z_2)=z_0+\xi e,
\]

where z0 is orthogonal to the unit vector e. Put

\[
S(\theta)=R^2-|\rho-d|^2/s^2,
\quad S_{max}=R^2-(c-D)^2/s^2.
\]

The remaining positive energy is m=S(theta)−\textbar xi\textbar². The exact original positive-coordinate coherence has numerator

\[
n=R^2+\bar z_1(\rho-1)-v(1-z_1)\bar z_2,
\qquad \chi^2=|n|^2/[R^2(D^2-y)].
\]

Thus all physical central values and denominator phases remain explicit. Expanding n=n0+L xi+M bar(xi)+Q\textbar xi\textbar² gives

\[
L=v e_1\bar z_{0,2},\quad
M=\bar e_1(\rho-1)-v(1-z_{0,1})\bar e_2,\quad
Q=v e_1\bar e_2.
\]

For an exact trial eta in (0,1), if at every feasible phase

\[
|n_0|+(|L|+|M|)\sqrt{[S-\eta S_{max}]_+}
+|Q|[S-\eta S_{max}]_+
<\chi_0 R\sqrt{D^2-y}
\]

whenever S\textgreater=eta Smax, then high coherence chi\textgreater=chi0 forces m\textless eta Smax. This is a sufficient scalar inequality; no floating scan proves it over a continuum.

\subsection{Conditional integral and three-real-dimensional tail}\label{conditional-integral-and-three-real-dimensional-tail}

Uniform-sphere disintegration gives the candidate exact restricted transition density

\[
K_>(y,y')=\frac{(a-1)(a-2)y}{2\pi^2 R^{2a-2}s^2 y'^2}
\int d\theta\int_{|\xi|^2<S(\theta)}
[S(\theta)-|\xi|^2]^{a-3}{\bf1}_{\chi>\chi_0}\,d^2\xi.
\]

The complex marginal on the first two coordinates is (a−1)(a−2)/(pi²R⁴)*(1−\textbar zP\textbar²/R²)\^{}(a−3); the change of the complex denominator contributes s\^{}−2 and its polar radial Jacobian is y/(2y'\^{}2). The unreduced formula must be matched directly with K121 KP5 before a final normalization claim.

Set L²=4Dc/s² and u=L sin(theta/2). Then

\[
m=S_{max}-u^2-|\xi|^2,\qquad
d\theta=2\,du/[L\sqrt{1-u^2/L^2}].
\]

Feasibility gives 0\textless Smax/L²\textless1/2, because D²+c²−R²s²\textgreater0. Hence the additional phase density lies between1 andsqrt2. If the preceding cap m\textless eta Smax is proved, the ratio to the unrestricted conditional transition is at most

\[
\sqrt2\,\Pr\{X\ge1-\eta\},\qquad X\sim\operatorname{Beta}(3/2,a-2).
\]

This follows by radial integration on the three-real-dimensional ball (u,Re xi,Im xi), whose density is (Smax−radius²)\^{}(a−3). It is not a one-dimensional theta estimate, and its use still needs a full all-phase cap, both directions, original near-point transport, a valid positive row majorant and arithmetic composition. A compact-support Mellin profile must not silently be treated as a positive Laplace mixture.

\subsection{Stronger selected overlap cutoff candidate}\label{stronger-selected-overlap-cutoff-candidate}

K122's actual finite-weight argument permits chi0=937/1000. With b=101/100 the literal weights give N\textgreater31.9999B, mu0\textless33.0001B and mu1\textless1.0001B. The two-singular-value formula gives sigma\_plus\textless637163/625000. The sufficient exact integer-power comparison is

\[
2(637163/625000)^{160000}(100/101)^{309998}<1,
\]

which would give overlap\textless2\^{}(−B/5000). The full K122 arithmetic mass still needs to be explicitly carried when freezing this refinement. No comparison has been accepted merely from its decimal logarithm.

\subsection{Floating diagnostics, not bounds}\label{floating-diagnostics-not-bounds}

At y=y'=1/33 and coherence0.9, SLSQP searches over every feasible denominator phase and one complex xi found: t=4,w=0 maximal remaining ratio0.666627; t=13,w=0 ratio0.734547; t=40,w=0 ratio0.917825; t=13,w=3 ratio0.780987. These are local numerical optima, not certified maxima or sector estimates.

A separate sampled triangle bound with eta0.95 over6\textless=t\textless=32 and0\textless=w\textless=4 had maximum coherence upper expression about0.89737, below the proposed0.937 cutoff. This is only a sampled diagnostic; the complete band y,y' and parameter continuum have not been certified. No error charge or support floor follows.

\subsection{Post-CP52 positive-composition route (still pending)}\label{post-cp52-positive-composition-route-still-pending}

The exact integer comparison displayed above passed in \nolinkurl{coherence_positive_majorant_probe.py}; that single comparison is exact, while its parameter screen is explicitly floating only. The screen samples all phase signs and both endpoints/midpoint of each narrow radial variable. At eta=19/20 it finds a maximum about0.914 on 6\textless=t\textless=48, \textbar w\textbar\textless=4. This does not establish the required continuum inequality.

If that complete rectangle cap is certified, the conditional Beta bound has the explicit concavity estimate

\[
\epsilon_\eta\le\frac{\sqrt2\,\eta^{166}}{166B(3/2,166)}
\left(\sqrt{1-\eta}+\frac{\eta}{334\sqrt{1-\eta}}\right).
\]

At eta19/20 its floating value is0.000976416121\ldots, slightly below2\^{}-10. That last comparison also needs an exact rational/root certificate; a looser1/1000 is available if useful. The reversed kernel has the same t and opposite physical w, and the same statement with y,y' exchanged, so a cap valid for both variables over the entire band applies to both detailed-balance directions. This uses the same original positive-coordinate coherence; there is no independently chosen rotation.

A proposed global composition uses, for any a\textgreater0, 0\textless b\textless29/20 and integer1\textless=r\textless83,

\[
1_{\text{high coherence}}\text{ (conditional fraction)}
\le\epsilon_\eta+e^{a(6-t)}+e^{b(t-48)}+(w/4)^{2r}.
\]

Indeed inside the rectangle the fraction is at most epsilon; outside, at least one displayed complement term is at least one. It is an inequality for the conditional kernel, not a pointwise inequality for the event. The resulting restricted fidelity is bounded by multiplying this nonnegative factor into the SAME CM6/CM7 kernel, then pairing directions and integrating the original central coset. Near-point and Voronoi perturbations of the extra physical moment still require explicit payment. A safe route is to enlarge the even moment by an explicit tiny positive regularizer and expand it; no logarithmic derivative of w\^{}(2r) at zero may be used.

The central even moment has a closed positive-mixture interface. Write gamma=29/10, lambda=29/20, Rstar=3514/29 and b(t)=1+t/4, A(t)=1+t/2. In CM12's exact central integral, insertion of (w/W)\^{}(2r) changes its normalized law to

\[
Z\sim\operatorname{Beta}(r+1/2,83-r),\qquad0\le r<83.
\]

The total shape remains167/2. The same support-plane and conditional Jensen argument therefore gives, relative to the unchanged G1 Dstar factor, the positive row profile

\[
C_r\,t^{-1/2}b(t)^{-(166-2r)}A(t)^{-1/2}
 e^{-29t/20}(1+t/Rstar)^{-(r+1/2)},
\]

\[
C_r=(4/W^2)^r
\frac{\Gamma(r+1/2)\Gamma(83-r)}{\Gamma(1/2)\Gamma83}
(1670/1757)^r.
\]

The central U\^{}f covolume, original ideal and G1 normalization remain unchanged. For integral r\textless=82 every Gamma shape is positive. The scalar t\^{}-1/2 still gives an infinite positive measure, so M0 is not introduced; each positive inverse moment is obtained by its original integral. The sole zero-degree polynomial term must retain the actual nonzero-row repair, using the new measure's finite M1. In particular the potentially large e\^{}(6a) factor cannot be discarded from that M1 or from lower-degree interpolation.

The two radial complement terms are also positive profiles: replace lambda by lambda+a and lambda-b, with coefficients e\^{}(6a) and e\^{}(-48b), respectively. Thus the sum is a genuine global positive Laplace majorant once the geometric and transport steps are proved. It does not treat a compact profile's Mellin moments as a Laplace representation.

A separate explicitly floating moment diagnostic, \nolinkurl{coherence_mixture_moment_probe.py/.json}, uses every profile term before integrating and finds at j163 about1667 field bits of gain for eta.95, W4, r8 and optimized a\textasciitilde9,b\textasciitilde1.259. The rational choices a9,b5/4 appear to retain this margin. The diagnostic is conditional on the unproved continuum cap and does not include a complete sector polynomial, source-cost replay or error charge. Its approximate main-volume gain suggests a possible k160 application after DO126, not a claim for k161 or beyond.

\clearpage

\section{Existing LIFT102 consequence: the literal LawsFull.shUniq clause}\label{existing-lift102-consequence-the-literal-lawsfull.shuniq-clause}

Status: source-context mapping only, recorded 2026-11-23. This note is not a new finding, does not claim new movement, and changes no canonical file. It identifies a named record field already supplied by the accepted LIFT102 Theorem D. No new executable certificate is needed for this direct application.

\subsection{1. Exact source clause}\label{exact-source-clause}

The complete 55-line \nolinkurl{FullExists.lean} was freshly read. Its \nolinkurl{LawsFull.shUniq} field, lines 44--47, quantifies over every actual stage \nolinkurl{j} with \nolinkurl{InR D j}, every \nolinkurl{n}, every supplied \nolinkurl{StarRing (O_E)}, and every alternative operator \nolinkurl{nab'} on the same stored second-stage object and coefficient field. Given its \nolinkurl{IsNab} witness, it asks for

\[
 \operatorname{shortGen}\bigl(\ldots,
       \{T.d_\delta(j,n)\text{ with nab}_2:=\mathrm{nab}'\}\bigr)
 =
 \operatorname{shortGen}\bigl(\ldots,T.d_\delta(j,n)\bigr)
\]

for all \nolinkurl{a} with \nolinkurl{Dom a n}. The notation suppresses only arguments that are literally identical on both sides: the original first object, coefficient field, first connection, \nolinkurl{badN D}, \nolinkurl{D.R.nV}, \nolinkurl{a}, and \nolinkurl{n}. The update changes only the operative \nolinkurl{nab₂} entry.

The complete 83-line \nolinkurl{HasDataCondition.lean} was freshly read. \nolinkurl{LawsSh.dNab}, line 60, supplies \nolinkurl{IsNab} for that original operative \nolinkurl{(T.dδ j n).nab₂} for every \nolinkurl{j}, \nolinkurl{InR D j}, and \nolinkurl{n}. This implication needs the unbundled \nolinkurl{dNab} clause; it need not assume the existence of an entire \nolinkurl{LawsFull} record to construct the \nolinkurl{shUniq} field.

\subsection{2. Application of the already accepted theorem}\label{application-of-the-already-accepted-theorem}

Accepted \nolinkurl{agents/prescribed_exponential_frame_connection_bridge.md}, section 7, Theorem D, compares two actual pins on the same object and coefficient field. It proves equality of each contracted \nolinkurl{ShortValAt} summand and hence equivalence of \nolinkurl{IsShortRule} and equality of the same \nolinkurl{shortGen} subsets. Its proof uses the harmonic polynomial condition present in those relations and the common nonnegative power \nolinkurl{N²} in every column. It retains every permitted faithfully flat witness, theta-group lift, enumeration, root, section, finite weight, and scalar.

The existentially presented star operations cause no mismatch: the source \nolinkurl{IsNab} pins each of them to the actual CM involution, as checked in LIFT102 Theorem D. The replacement connection in the literal \nolinkurl{shUniq} field has exactly the second \nolinkurl{IsNab} witness needed by that theorem. The first witness is supplied by the unbundled \nolinkurl{dNab} clause.

Consequently, the accepted theorem directly constructs the literal \nolinkurl{LawsFull.shUniq} proposition for all \nolinkurl{n}. There is no extra equality of ordinary coefficients, reduced-gauge choice, admissible-lattice existence, scalar descent, or selected-image nonvanishing hypothesis in this application. There is also no new low-rank argument here: Theorem D already compares the defining relations for every \nolinkurl{n}.

\subsection{3. Fixed data and remaining source obligations}\label{fixed-data-and-remaining-source-obligations}

\begin{longtable}[]{@{}
  >{\raggedright\arraybackslash}p{(\columnwidth - 2\tabcolsep) * \real{0.5000}}
  >{\raggedright\arraybackslash}p{(\columnwidth - 2\tabcolsep) * \real{0.5000}}@{}}
\toprule\noalign{}
\begin{minipage}[b]{\linewidth}\raggedright
Record item
\end{minipage} & \begin{minipage}[b]{\linewidth}\raggedright
Exact status in this mapping
\end{minipage} \\
\midrule\noalign{}
\endhead
\bottomrule\noalign{}
\endlastfoot
Actual \nolinkurl{Top}, \nolinkurl{Full}, \nolinkurl{Q}, \nolinkurl{Dw}, \nolinkurl{dδ} family & Given and retained. No existence theorem for these records is claimed. \\
\nolinkurl{DeltaDat.nab₂} & The two supplied pinned entries are compared in the exact update appearing in \nolinkurl{shUniq}. \\
All other \nolinkurl{DeltaDat} fields & Retained literally, including levels, maps, characters, roots, normalization sections, and \nolinkurl{cδ}. \\
Original class sums and selected harmonic polynomial & Retained; no orbit or span enlargement is used. \\
\nolinkurl{LawsFull.shUniq} & Already follows from LIFT102 Theorem D and the actual \nolinkurl{dNab} clause. \\
Other \nolinkurl{LawsFull} fields & Not constructed by this mapping. In particular \nolinkurl{p}, \nolinkurl{j}, \nolinkurl{sh}, \nolinkurl{canon}, \nolinkurl{lvA}, \nolinkurl{halfI}, and \nolinkurl{xiVal} remain separate requirements. \\
Costs and finite remainder arrays & No change and no charge. \\
Actual admitted source and selected arithmetic image & Not supplied or proved nonzero by this mapping. \\
\end{longtable}

This note does not strengthen or amend the frozen LIFT102 theorem. It records the exact named downstream consumer so that the already proved connection-independence consequence is not counted as a new source gate or a new finding.

\subsection{4. Sources read}\label{sources-read}

\begin{itemize}
\tightlist
\item
  \nolinkurl{research/alpoge_source/Formalism/LocalToGlobal/ProjectiveLattices/FullExists.lean}: all 55 lines, especially \nolinkurl{LawsFull} and \nolinkurl{shUniq}.
\item
  \nolinkurl{research/alpoge_source/Formalism/LocalToGlobal/ProjectiveLattices/HasDataCondition.lean}: all 83 lines, especially the actual \nolinkurl{Top}, \nolinkurl{LawsSh.dNab}, and other named obligations.
\item
  \nolinkurl{agents/prescribed_exponential_frame_connection_bridge.md}: fresh sections 5--10, including the complete Theorem D and the genuine-package supplier analysis. Its full original author and independent reviews are already recorded by the accepted result.
\end{itemize}

The adjacent source record stores the current exact hashes and read scopes. No additional primary theorem is used beyond the accepted LIFT102 proof.

\clearpage

\section{KPsiIn: accepted upstream suppliers and the fixed-field application gate}\label{kpsiin-accepted-upstream-suppliers-and-the-fixed-field-application-gate}

Status: source-interface research snapshot, 2026-11-23. No new theorem, movement credit, source-validity audit, build, or canonical change is claimed. The statements below are cited as accepted supplied-source theorems. The purpose is to identify their exact applicability to the original records.

\subsection{1. The beta law already has an exact upstream supplier}\label{the-beta-law-already-has-an-exact-upstream-supplier}

Recovered source: \nolinkurl{Development/AbelianVarietyTheory/Proofs/UnitaryExistence/ExistsIndexIdeal.lean}, namespace \nolinkurl{GlobalTheory.BetaEx}, declaration 322--324:

\begin{verbatim}
theorem betaExists :
  ∀ (E : Type) [Field E] [NumberField E] [IsCMField E]
    (H : Host E) (Dw : Datum H) (l : Level H)
    (c : Coeff l) (n bad : ℕ),
    Formalism.L2.Pj.BetaExists (Dw := Dw) c n bad
\end{verbatim}

Its statement has no \nolinkurl{bad ≠ 1} premise. The shorter \nolinkurl{Proofs/CuspFormEvaluations/BetaExistsOfNeOne.lean}13--15 invokes this theorem and states a weaker interface with that additional premise. The universal supplier, rather than the weaker wrapper, applies to \nolinkurl{LawsSh.beta} at \nolinkurl{HasDataCondition.lean}58.

This does not construct \nolinkurl{KPsiIn}. The definition of \nolinkurl{BetaExists} at \nolinkurl{ThShort.lean}283--285 already has \nolinkurl{KPsiIn} as an implication antecedent. The latter remains the separate \nolinkurl{LawsP.kPsi} clause at line 55 of \nolinkurl{HasDataCondition.lean}. Re-proving the universal beta theorem is not needed to use this accepted supplier.

\subsection{2. The actual coefficient-extension constructor}\label{the-actual-coefficient-extension-constructor}

Recovered source: \nolinkurl{Development/AbelianVarietyTheory/Proofs/UnitaryExistence/ExistsCoeffTypeNormsPrincipal.lean}, namespace \nolinkurl{GlobalTheory}, declaration 23--24, full file 74 lines:

\[
 \forall E,H,Dw,l,c,\quad
 \exists c':\mathrm{Coeff}(l),\quad c.F\subseteq c'.F
 \ \land\
 \forall c'':\mathrm{Coeff}(l),\
 c'.F\subseteq c''.F\Longrightarrow
       \mathrm{KPsiIn}(Dw,c'').
\]

Here \nolinkurl{E} has the original field, number-field and CM instances. \nolinkurl{H} and \nolinkurl{Dw} are arbitrary actual \nolinkurl{Host} and \nolinkurl{Datum}, and the level \nolinkurl{l} is unchanged. The theorem supplies a new coefficient record; it does not state that the starting \nolinkurl{c} already satisfies \nolinkurl{KPsiIn}. Its proof's displayed construction uses class representatives \nolinkurl{I(x)}, principal powers \nolinkurl{I(x)^h=(gamma(x))}, and roots

\[
 \psi(x)^h
  =\prod_{\tau:E\hookrightarrow\mathbf C}
          \tau(\gamma(x))^{\mathrm{Dw.mult}(\tau)}.
\]

The finite coefficient field contains these roots and a basis of the starting field. The accepted \nolinkurl{exists_coeff_containing} declaration in \nolinkurl{IsIntegralRatZetaOfCoeff.lean}25--26 supplies a coefficient record containing any specified finite set of algebraic complex numbers, including all required field-embedding and level data.

The displayed theorem does not give a numerical degree or discriminant bound. No bound is inferred merely from its name or from the finite construction.

\subsection{3. Exact criterion in the same coefficient field}\label{exact-criterion-in-the-same-coefficient-field}

Recovered \nolinkurl{BlockTheory/ExistenceTori/ReflexNormIndexIdeal.lean} defines \nolinkurl{reflexNorm}, \nolinkurl{reflexIdeal}, \nolinkurl{IsBe}, \nolinkurl{ReflexLe}, \nolinkurl{ReflexRoots} and \nolinkurl{ReflexBasis} on the same actual \nolinkurl{Dw,c}.

\nolinkurl{ReflexRoots Dw c} asks, for every nonzero integral ideal \nolinkurl{a} of \nolinkurl{O_E}, for \nolinkurl{h>0}, \nolinkurl{gamma∈O_E} and \nolinkurl{psi∈c.F} satisfying

\[
 a^h=(\gamma),\qquad
 \psi^h=\prod_\tau\operatorname{embF}(c,\tau)(\gamma)^{
                   \mathrm{Dw.mult}(\tau)}.
\]

\nolinkurl{ReflexBasis Dw c} instead records such principal powers and roots for a finite invariant-factor system of ideal classes, together with the exact unique class-coordinate representation.

The accepted \nolinkurl{GlobalTheory.LatticeBetaCompare} interfaces in \nolinkurl{Proofs/UnitaryExistence/MultAddMultConjugate.lean} are:

\begin{longtable}[]{@{}
  >{\raggedright\arraybackslash}p{(\columnwidth - 2\tabcolsep) * \real{0.5000}}
  >{\raggedright\arraybackslash}p{(\columnwidth - 2\tabcolsep) * \real{0.5000}}@{}}
\toprule\noalign{}
\begin{minipage}[b]{\linewidth}\raggedright
Source declaration
\end{minipage} & \begin{minipage}[b]{\linewidth}\raggedright
Exact consequence
\end{minipage} \\
\midrule\noalign{}
\endhead
\bottomrule\noalign{}
\endlastfoot
\nolinkurl{kPsiIn_iff}269--272 & \nolinkurl{KPsiIn Dw c ↔ ReflexLe Dw c}. \\
\nolinkurl{reflexLe_of_basis}447--458 & \nolinkurl{ReflexBasis Dw c → ReflexLe Dw c}. \\
\nolinkurl{reflexRoots_iff}460--474 & \nolinkurl{ReflexRoots Dw c ↔ ReflexLe Dw c}. \\
\nolinkurl{typings_agree}557--589 & Includes both \nolinkurl{KPsiIn Dw c ↔ ReflexRoots Dw c} and \nolinkurl{ReflexBasis Dw c → KPsiIn Dw c} explicitly. \\
\end{longtable}

Thus a verified root table in the original \nolinkurl{c.F} supplies \nolinkurl{KPsiIn} without changing that field. The stated criteria do not themselves assert that the root table exists in the prescribed field.

The separate definition \nolinkurl{ReflexFieldExists} in \nolinkurl{ReflexNormIndexIdeal.lean}110--111 specifies a desired extension bound involving \nolinkurl{classNumber E}. That definition alone is not an existence theorem or a numerical bound available for the actual application; no such inference is made here.

\subsection{4. The original pin is stronger than being a Coeff record}\label{the-original-pin-is-stronger-than-being-a-coeff-record}

The original \nolinkurl{Par} at \nolinkurl{Formalism/LocalToGlobal/ProjectiveLattices/PartitionExtendByZeroAntitone.lean}14--20 contains \nolinkurl{eO : O(c.F) ≃ D.what₂.O}. Its \nolinkurl{PinField} at line 45 also identifies the actual complex embedding through the prescribed \nolinkurl{D.what₂} data. \nolinkurl{Pins}49 includes this condition.

The accepted \nolinkurl{Par.F_eq} in \nolinkurl{MapRingEquivTransApply.lean}68--69 states that two parameter records for the same \nolinkurl{D} satisfying their actual \nolinkurl{PinField} have equal coefficient fields as subfields of \nolinkurl{C}. The preceding \nolinkurl{coe_mem_of_pinField}52--58 and \nolinkurl{F_le}60--66 show its exact interface.

Therefore the extension constructor in section 2 cannot simply be substituted into an already pinned \nolinkurl{Par} while keeping \nolinkurl{D.what₂}, \nolinkurl{eO} and \nolinkurl{PinField} unchanged. It applies directly to an unpinned \nolinkurl{Coeff} choice, or to a new jointly constructed source whose coefficient pin and all dependent data are supplied compatibly. This is an application boundary, not a claim that such a new source cannot be constructed.

\subsection{5. Remaining work and provenance}\label{remaining-work-and-provenance}

At fixed original data, the useful next supplier is an actual \nolinkurl{ReflexBasis}/\nolinkurl{ReflexRoots} witness in the already pinned field. For a newly chosen source, the extension theorem must be accompanied by a compatible \nolinkurl{what₂/eO/PinField} construction and the resulting arithmetic cost bounds. Neither task is completed by this snapshot. No selected theta image, cusp/spanning assertion, global period, specific032 value, or DP44 sign is inferred.

All recovered files come from the fixed source commit \nolinkurl{7cbf7a0c32fca1badf05931879a37c3f8a2b03b4}. Their exact URLs, retrieval times and SHA-256 hashes are in \nolinkurl{kpsi_sources/retrieval_01.json} through \nolinkurl{retrieval_03.json}. The adjacent snapshot source record distinguishes whole-file readings from targeted declaration/interface readings.

\clearpage

\section{AM-K-124 --- Nested coherence caps and a complete positive mixture at support twenty-seven}\label{am-k-124-nested-coherence-caps-and-a-complete-positive-mixture-at-support-twenty-seven}

Owner: \nolinkurl{/root/kummer_extension}. \textbf{Pending research draft.} The predecessor K123 is frozen and is not modified here. The new continuum diagnostics and complete arithmetic count are being checked; no numerical charge is asserted by this draft until the final certificate and whole reviews are complete.

\subsection{1. Actual containing family and the retained selected projection}\label{actual-containing-family-and-the-retained-selected-projection}

We retain the original bounded all-II, globally rank-two-unipotent scalar-cross family at support 27, with nonzero original S-cross, full positive compression, total eligibility 27, no inactive paired-cross zero and no active integral-cross place. Let k be the number of nonzero original S-row coordinates, x=N\_E D\_eta and f=168. The containing denominator window is

\[
2^{2111}<x<2^{3328}.\tag{NC1}
\]

The proposed new application is k=161 or162 throughout this window. Earlier overlapping bounds are retained; no assertion is made for k163, k164 or a full support-28 floor.

All finite labels, original projected columns, compact K-factors, Haar normalization, central affine cosets, original primitive null line and the actual selected Hilbert metric remain those of the literal K113/DO126/DO127 interfaces. K123 HCX1--5 supplies the genuine selected high-coherence restriction with chi0=937/1000, the band I={[}16/529,16/527{]}, and the complete containing-family comparison

\[
\sum_h |a(h)c_W(h)|\le
\sum_h |a(h)|\prod_\sigma F_{I,\chi>\chi_0}(h_\sigma)
+2^{-8605114}+2^{11394885-B/5000},\qquad B=2^{955}.\tag{NC2}
\]

This is the actual selected projection estimate. It includes the nonzero-S-cross premise needed in the arithmetic reconstruction. It is applied before any lattice enlargement. We do not replace the selected projection by a Fischer projection, and we do not introduce independent rotations in the two kernel directions.

\subsection{2. Exact coherence numerator and a stronger disk bound}\label{exact-coherence-numerator-and-a-stronger-disk-bound}

Use the same simultaneous compact normal form as K123 HCX6--11. At one embedding set A=t/2, v=sqrt(t), d=1+A+iw, c=\textbar d\textbar{} and s²=c²−1, with t\textgreater0 and physical w real. On the two moving positive coordinates and the negative coordinate,

\[
g=\begin{pmatrix}1-A-iw&-v&A+iw\\v&1&-v\\-A-iw&-v&1+A+iw\end{pmatrix}.
\tag{NC3}
\]

The original unitary matrix is the identity on all other positive coordinates. This simultaneous conjugation preserves the selected K-type, h\_nn, t, radial measure and the positive-coordinate coherence.

For y,y' in I, put D=sqrt(y/y'), R\_y²=1−y and rho=D(d/c)e\^{}\{i theta\}. On the denominator fibre,

\[
z_P=z_0+\xi e,\quad
z_0=(\rho-d)(-A+iw,-v)/s^2,\quad e=(v,-A-iw)/s,
\]

\[
S_*={2D\over c+1}-y-{(D-1)^2\over s^2},\quad
m=S_*-{4Dc\over s^2}\sin^2(\theta/2)-|\xi|^2.
\tag{NC4}
\]

The exact coherence numerator is

\[
n=R_y^2+\bar z_1(\rho-1)-v(1-z_1)\bar z_2
=n_0+L\xi+M\bar\xi+Q|\xi|^2,
\quad \chi^2={|n|^2\over R_y^2(D^2-y)}.\tag{NC5}
\]

Here L=v e\_1 bar(z0,2), M=bar(e1)(rho−1)−v(1−z0,1)bar(e2), and Q=v e1 bar(e2)=2A(−A+iw)/s². These are identities in the original normal form, not asymptotic expansions.

The new elementary bound retains the quadratic term's direction. If \textbar xi\textbar≤r, define a=\textbar n0\textbar, b=\textbar n0+Qr²\textbar{} and z=(\textbar L\textbar+\textbar M\textbar)r. Convexity of the complex norm in the real argument \textbar xi\textbar² gives

\[
|n_0+Q|\xi|^2|\le (1-u^2)a+u^2b,\qquad u=|\xi|/r\in[0,1].
\]

Consequently

\[
|n|\le \mathcal C(a,b,z):=
\begin{cases}
b+z,&b\ge a\ \text{or}\ z\ge2(a-b),\\
a+z^2/[4(a-b)],&b<a\ \text{and}\ z<2(a-b).
\end{cases}\tag{NC6}
\]

At r=0 the value is \textbar n0\textbar. The formula follows by maximizing a+(b−a)u²+zu over {[}0,1{]}; it includes z=0 and a=b. More importantly, the maximum is increasing separately in a,b,z, since each coefficient of these three inputs is nonnegative for every u. Thus outward upper intervals for all three inputs can safely be substituted in NC6. This step preserves the negative quadratic cancellation that the earlier sum \textbar n0\textbar+(\textbar L\textbar+\textbar M\textbar)r+\textbar Q\textbar r² discarded.

\subsection{3. Nested physical-central shells and exact continuum obligation}\label{nested-physical-central-shells-and-exact-continuum-obligation}

The proposed shell data are the following rational pairs. Each row concerns the interval from the preceding W to its displayed W, starting at zero.

\begin{longtable}[]{@{}rrr@{}}
\toprule\noalign{}
W & eta & central moment r \\
\midrule\noalign{}
\endhead
\bottomrule\noalign{}
\endlastfoot
2 & 845/1000 & 3 \\
5/2 & 871/1000 & 4 \\
3 & 889/1000 & 5 \\
7/2 & 902/1000 & 7 \\
4 & 911/1000 & 8 \\
9/2 & 918/1000 & 9 \\
5 & 926/1000 & 11 \\
6 & 939/1000 & 13 \\
7 & 951/1000 & 16 \\
8 & 962/1000 & 18 \\
10 & 981/1000 & 22 \\
\end{longtable}

The exact continuum obligation, independently checked by integer intervals, is

\[
6\le t\le48,\quad W_{i-1}\le |w|\le W_i,\quad y,y'\in I,
\quad m\ge\eta_i S_*\quad\Longrightarrow\quad\chi<937/1000.
\tag{NC7}
\]

The final executable must cover the entire continuum; the sampled precursor is explicitly insufficient. The following exact normalization removes artificial variation from the ratio D without identifying y and y'. Set

\[
 p=\sqrt{yy'},\qquad e=D+D^{-1}-2,\qquad
 \bar S=S_*/D={2\over c+1}-p-{e\over s^2}.
 \tag{NC7a}
\]

Both original radial variables remain independent. Their images satisfy p in I and 0≤e\textless1/527²: after possibly replacing D by its inverse, 1≤D≤sqrt(529/527) and D−1\textless1/527. We enlarge these two normalized parameters independently to the displayed containing intervals; no equality y=y' is imposed.

Write h=sin(theta/2)=u sqrt((1−eta)bar(S)s²/(4c)), with −1≤u≤1, and let bar(r)²=r²/D=(1−eta)bar(S)(1−u²). Direct substitution into NC5 gives

\[
 {n_0\over D}=-p+{e^{-i\theta}\over c}
 +C\bigl[(c+c^{-1})\cos\theta-2-e
                    +i(c-c^{-1})\sin\theta\bigr],
\quad C={A(A^2+w^2)+iw(s^2+2A)\over s^4},
 \tag{NC7b}
\]

and Q=2A(−A+iw)/s² as before. The normalized coherence denominator is exactly

\[
 {R_y\sqrt{D^2-y}\over D}
   =\sqrt{(1-y)(1-y')}=\sqrt{(1-p)^2-pe}.\tag{NC7c}
\]

The linear coefficients satisfy \textbar L\textbar=t sqrt(t)\textbar rho−d\textbar/s³ and \textbar M\textbar=t sqrt(t)\textbar bar(d)rho−1\textbar/s³. In particular

\[
 { |\rho-d|^2+|\bar d\rho-1|^2\over D}
 =2(c-1)^2+(c^2+1)e+8ch^2.
\]

Cauchy--Schwarz therefore gives the containing linear bound

\[
 {(|L|+|M|)r\over D}
 \le {4A\sqrt t\over s^3}
 \sqrt{(c-1)^2+(c^2+1)e/2+4ch^2}\,\sqrt{\bar r^2}.
 \tag{NC7d}
\]

Apply NC6 to the upper bounds for \textbar n0/D\textbar, \textbar n0/D+Q bar(r)²\textbar{} and the right-hand side of NC7d. This is an exact all-parameter bound, not an expansion near D=1. Conjugation handles negative w; inversion exchanges y,y', sends w to−w, and preserves t and chi. Hence the same cap applies in both detailed-balance directions.

The diagnostic uses 72-bit integer outward intervals, exact integer square roots and rational subdivision of t,w,u. The full p and e intervals are retained in every box. A box is accepted only when its complete normalized NC6 upper bound divided by the positive denominator NC7c is strictly below937/1000. Every denominator must have a positive lower bound. The final coverage counts and all unresolved-box checks will be inserted only after the exact run completes. The earlier slower unnormalized diagnostic and its first two complete shells remain preserved as history; they are not substituted for the final eleven-shell check.

As eta\_i increases with i, the preceding stronger shells imply the corresponding cap throughout \textbar w\textbar≤W\_i. Thus no uncounted containing-subset or central-sign factor is introduced.

\subsection{4. Conditional beta probabilities and a global positive majorant}\label{conditional-beta-probabilities-and-a-global-positive-majorant}

K123 HCX13's exact three-real-dimensional coarea integral applies unchanged. It includes the original complex marginal, s\^{}-2 denominator Jacobian and y/(2y'\^{}2) radial-polar Jacobian. After the substitution q=L sin(theta/2), the extra phase density lies between1 andsqrt2. Hence NC7 gives the forward and reverse conditional kernel fraction

\[
\epsilon_i\le {\sqrt2\,\eta_i^{166}\over\vartheta}
\left(\sqrt{1-\eta_i}+{\eta_i\over334\sqrt{1-\eta_i}}\right),
\quad \vartheta={4^{167}\over334\binom{334}{167}}.
\tag{NC8}
\]

This uses Beta(3/2,166), not a one-dimensional theta estimate. Its square is the exact rational number

\[
E_i^2={2\eta_i^{332}\bigl(1-\eta_i+\eta_i/334\bigr)^2
\over\vartheta^2(1-\eta_i)}.\tag{NC9}
\]

Choose increasing rational e\_i≥E\_i by an outward integer square root on a fixed dyadic grid. Put delta\_i=e\_(i+1)−e\_i for i\textless11 and delta\_11=1−e\_11. The final certificate checks all delta\_i\textgreater0. The nested step majorant telescopes, and therefore the conditional high-coherence fraction is bounded for all t\textgreater0,w real by

\[
e_1+\sum_{i=1}^{11}\delta_i(w/W_i)^{2r_i}
+e^{9(6-t)}+e^{(5/4)(t-48)}.\tag{NC10}
\]

For a central shell, all preceding powers are≥1 and their coefficients sum to the necessary e\_i; outside the last shell the sum is≥1. Outside6≤t≤48 a radial exponential is≥1. This is a bound on the original conditional kernel fraction. It is not a false pointwise indicator bound and does not treat a compact function as a Laplace transform. The same factor bounds both detailed-balance directions and enters their geometric mean once.

\subsection{5. Physical central integration, transport and the positive measure}\label{physical-central-integration-transport-and-the-positive-measure}

Multiply NC10 into the SAME positive kernel of DO125/K123 before physical central integration. Let gamma=29/10, lambda=29/20, b(t)=1+t/4, A(t)=1+t/2, Rstar=3514/29. For any integer0≤r\textless83 the exact insertion (w/W)\^{}(2r) yields, relative to the unchanged G1Dstar physical central normalization,

\[
C(r,W)\Phi_r(t),\qquad
\Phi_r(t)=t^{-1/2}b(t)^{-(166-2r)}A(t)^{-1/2}
(1+t/Rstar)^{-(r+1/2)}e^{-29t/20},
\tag{NC11}
\]

\[
C(r,W)=(4/W^2)^r
{\Gamma(r+1/2)\Gamma(83-r)\over\Gamma(1/2)\Gamma83}
(1670/1757)^r.
\tag{NC12}
\]

The derivation is the exact w=2bu substitution, the same physical c=sqrt(A²+w²) support plane, and the tilted Beta(r+1/2,83−r) conditional Jensen argument; the total beta shape stays167/2. All Gamma shapes in the table are positive. The physical U\^{}-f, sqrt(D\_F)\^{}-1, N\_F(J0)\^{}-1, x\^{}-1 and T54 factors are unchanged.

Original near-point transport gives e\^{}-2rp≤tp/t≤e\^{}2rp and \textbar wp−w\textbar≤6Hrp. The actual central750 cover, including sqrt(f), gives total coordinate displacement\textless2\^{}-68. Apply the elementary Young inequality separately to every even power, with e=2\^{}-20:

\[
(|w|+\delta)^{2r}\le(1+e)^{2r-1}|w|^{2r}
 +(1+e^{-1})^{2r-1}\delta^{2r}.\tag{NC13}
\]

No logarithmic derivative of w at zero is used. The full small atom sum must be checked and absorbed into an explicit added base coefficient. Each nonconstant coefficient retains its factor(1+2\^{}-20)\^{}(2r−1). The scalar radial transports and old three nonpattern inflations are all retained; one additional full factor2 beyond the original DO127 prefactor pays the new scalar transport. The final certificate must verify this with the actual H=133 and rp\textless2\^{}-352.

Thus the proposed global row profile is the sum of14 positive profiles:

\[
\Psi(t)=(e_1+\epsilon_{\rm atom})\Phi_0(t)
+e^{54}\Phi_0(t)e^{-9t}+e^{-60}\Phi_0(t)e^{5t/4}
+\sum_{i=1}^{11}\delta_i(1+2^{-20})^{2r_i-1}C(r_i,W_i)\Phi_{r_i}(t).
\tag{NC14}
\]

The smallest exponential shift is1/5. Therefore Psi is the Laplace transform of one positive, possibly infinite measure, and every positive inverse moment used below is finite. No moment of degree zero is introduced. All terms are combined before taking the f-th power; there is no omitted multiplicity factor.

\subsection{6. Exact moments and the original zero-degree repair}\label{exact-moments-and-the-original-zero-degree-repair}

For j≥1 define

\[
M_j={1\over\Gamma(j)}\int_0^\infty t^{j-1}\Psi(t)\,dt.\tag{NC15}
\]

The final executable encloses the required positive component moments on {[}0,1{]}, {[}1,T{]} and {[}T,infinity), retaining both complete tails. Here T=256 for the slow exponential rate1/5, T=128 for central powers r≥16, and T=64 otherwise. It uses the inherited integer/Fraction quadrature with step1/128, local second-derivative bounds and outward exponential/square-root arithmetic. If the residual right-tail power beta=j−337/2+r is positive, then t\^{}beta≤T\^{}beta exp(beta(t−T)/T) for t≥T. The complete tail has exponential rate lambda−max(beta,0)/T, and its positivity is checked for every component. The first complete diagnostic with uniform T=64 was a failed sufficient upper bound and is preserved; no negative inference about the actual sector follows from it.

The finite interpolation anchors are1,148,160,164,165. Between two finite anchors, Holder is applied directly to the SAME positive measure; no probability normalization or finite total mass is assumed. The accepted baseline at160 is itself bounded by directed Holder interpolation from its accepted148 and162 anchors. All fourteen new components are added before taking any field power or interpolation.

A safe first-moment bound is M1\textless2\^{}79. Indeed the scalar factors other than t\^{}−1/2 and the exponential are at most1. Their complete integrals are less than2 at the ordinary rate29/20, less than4 at rate1/5, and less than1 at rate209/20, using the certified bounds on pi. All central coefficients are less than1; their sum and the base/slow terms give less than30, while exp54\textless2\^{}78. This retains the large low-t coefficient rather than discarding it. The unique nominal zero-degree term first reserves theta=2\^{}−20 of the original actual nonzero-row energy, giving its M1\^{}f(q/(theta Rk))\^{}f factor before Gaussian enlargement.

\subsection{7. Complete original arithmetic polynomial and endpoint obligations}\label{complete-original-arithmetic-polynomial-and-endpoint-obligations}

The genuine rank167 isotropic quotient, its selfdual rank3 transverse part and the exact Green/J0 cancellation are those of DO126/DO127. Let K be the accepted scalar336-dimensional degree-15 coefficient and K3 the genuine1008-dimensional transverse coefficient. For each k the full polynomial has2(k+1) terms:

\[
\sum_{l=0}^k\binom{k}{l}K^{k-l}
\left({q\over(1-\theta)R}\right)^{(l+3)f}M_{l+3}^f
+K_3\sum_{l=1}^k\binom{k}{l}K^{k-l}
\left({q\over(1-\theta)R}\right)^{lf}M_l^f
+K_3K^kM_1^f\left({q\over\theta Rk}\right)^f.
\tag{NC16}
\]

Every local Green rank and both rank-one orientations are retained. At effective degree j the exact four-state pattern is product\_v(1+q\_v\^{}(j−167))². Here j≤165, so the full product is bounded by2\^{}(1/1024). The complete q exponent is(j−167)f\textless0 and the coupled x derivative is−1−j\textless0, allowing the all-zero pattern cap and the endpoint X=2111. Neither the pattern product nor the denominator endpoint is substituted before this monotonicity proof.

The original full prefactor is the DO127 GC7 expression with exactly one additional factor2 for the new transports. Explicitly it retains support binom(168,27), the164 pivot choice, binom(164,k), Dvar\^{}(f/2), vartheta\^{}f, G1Dstar\^{}f, sqrt(D\_F)\^{}−1, U\^{}−f, (R/U)\^{}(167f), T\^{}8904, Q\^{}(−335·167), x\^{}−168 q\^{}(−167f), coarea\^{}(f/2), pi\_+\^{}((k+3)f), D0\^{}(−k−3), Green2 and the three old nonpattern inflations. At X2111 the unchanged exact cap is251124964\textless q*\textless251124965, proved by168th powers.

The final full theorem requires all324+326=650 positive terms for k161 and162, together with both entire arithmetic errors in NC2, to fit one new complete charge. The endpoint and charge will be stated only after the exact replay. No floating dominant-term screen is a sector bound.

\subsection{8. Source and verification boundary}\label{source-and-verification-boundary-2}

The new analytic input is NC6 and its full all-parameter shell certificate, followed by the telescoping positive majorant and all moment coefficients. The actual selected projection/coarea/normal form and complete arithmetic source interfaces are carried from the frozen whole-reviewed K122/K123, DO125, DO126, DO127, K113, K118 and K120 inputs with explicit final hashes. The scalar and transverse theta bounds are the genuine dimension336 and1008 statements. No unproved all-ring theta construction, selected-image nonzero claim, new eta uniqueness, or source validity audit is used. Canonical files and frozen author subjects remain unchanged.

\subsection{Pending exact computation status (not an author freeze)}\label{pending-exact-computation-status-not-an-author-freeze}

The complete revised arithmetic diagnostic passed 4,936,976 exact assertions and615,936 quadrature intervals at2026-11-23T14:09:11.597321+00:00. It retains all650 terms; its largest directed exponent is−14798105/32768, at k162 and effective degree165. Conditional on every NC7 continuum cap, all terms and both full errors fit one proposed charge350. The output is explicitly PENDING\_CAP\_ARITHMETIC\_ONLY. The eleven-shell normalized continuum run has no observed complete result, so this is not a complete theorem verdict or an author freeze. The earlier safe but insufficient right-tail-at64 envelope is retained separately.

\clearpage

\section{AM-DO-129 --- Single-lattice positive completion of the support27 remainder}\label{am-do-129-single-lattice-positive-completion-of-the-support27-remainder}

Author: \nolinkurl{/root/definite_oscillator}, 23 November 2026 UTC.

\textbf{Controlled pending theorem.} This is a one-sided allocation of the actual signed kernel. It is not an absolute estimate for its support27 remainder. The actual predecessor is K123, with138 all-top and31 uniform charges and complete absolute support floors27/11. Those ordered arrays and floors are unchanged. Whole author verification, independent review and canonical acceptance are separate events.

\subsection{1. Original kernel and the precise new inequality}\label{original-kernel-and-the-precise-new-inequality}

Keep every original J56 datum, the actual H133 cutoff, CM involution and paired lattices, both omitted pencil labels, the selected vector-valued Hilbert projector and its two compact factors, words, twists, phases, finite mixed projectors, Haar probabilities and class masses. Write

\[
 f=168,\quad m=T-1,\quad H=133,\quad R=17424/133,
 \quad Q=1427908193729121522601,\quad U=77Q^2-14.
 \tag{PA1}
\]

Let \(F(g)=a(g)c(g)\prod_{v\mid T}\rho_v(g)\) be exactly K062 AT14--16, divided by its original positive normalization as in AT24. Thus the complete normalized energy is \(\mathcal E=\sum_{g\in\mu_{28}\backslash\mathscr C}F(g)\), and the identity class contributes1. The entire series is absolutely convergent. All representatives in a unipotent group below are unique modulo \(\mu_{28}\); no factor28 is inserted or removed.

Use the actual predecessor's successive set differences to partition the nonidentity original support into the already removed set \(E\), the current support27 set \(R_{27}\), and the remaining set \(R_{\ge28}\). The last set is the current unremoved higher-support set, not a claim that all matrices of support at least28 survive. Include the original outside-H tail and removed integral terms in E. With

\[
 \epsilon_{\rm old}=\sum_{e\in\text{K123 all-top array}}2^{-e},
 \qquad \sum_{g\in E}|F(g)|<\epsilon_{\rm old}<2^{-230},
 \tag{PA2}
\]

we take real parts of the current signed sums; their chosen successive set differences need not be separately inversion invariant. R097 and the subsequent checked cuts, finally K123, give the exact necessary conditions on \(R_{27}\): full positive compression, global rank-two unipotence, all27 active places of type II, exactly27 total eligible places, \(b=j_{\rm int}=0\), and

\[
 161\le k\le164,\qquad 2^{2111}<x=N_ED_\eta<2^{X_k},
 \quad(X_{161},X_{162},X_{163},X_{164})=(2130,2140,2150,2165).
 \tag{PA3}
\]

\textbf{Theorem.} There is a finite sum \(P\ge0\) of complete, absolutely convergent positive-definite Heisenberg sums, using actual lattice tuples from the original scalar orbit, and a real number \(\delta\), such that

\[
 \mathcal E=1+\mathop{\rm Re}\sum_{g\in R_{\ge28}}F(g)+P+\delta,
 \qquad |\delta|<\epsilon_{\rm old}+2^{-1266}<2^{-230}.
 \tag{PA4}
\]

Consequently the remaining sufficient signed-energy gate is

\[
 \mathop{\rm Re}\sum_{g\in R_{\ge28}}F(g)>-1+2^{-230}.
 \tag{PA5}
\]

It would imply \(\mathcal E>0\). This theorem does not prove PA5. Separately, \[
 \mathop{\rm Re}\sum_{g\in R_{27}}F(g)
       >-\epsilon_{\rm old}-2^{-1266}.
 \tag{PA6}
\] Neither statement says \(\sum_{R_{27}}|F|<2^{-1266}\), or that R27 is empty. In particular there is no absolute support-floor28 assertion, and1266 is recorded as a separate one-sided completion debit, not appended to either absolute-error array.

\subsection{2. A finite family of original null lines and their fixed label tuples}\label{a-finite-family-of-original-null-lines-and-their-fixed-label-tuples}

Let \(\Xi\) be the set of normalized parameters eta actually occurring in R27. It may be empty. Its matrices have intrinsic isotropic line

\[
 \ell_\eta=Ew,\quad w=(0_S,u,\eta),\quad
 u=(2,95192390527,383092671043,12529843548438452388609),
 \quad\eta\bar\eta=U.
 \tag{PA7}
\]

For such eta, its full eligible set Sigma has size27 and is exactly the active support of every R27 matrix with this eta. The paired bounded uniqueness K113 SC27.4 applies before any completion: throughout \(x<2^{12000}\) there is at most one eta per fixed27-place support. Thus

\[
                  |\Xi|\le\binom{168}{27}<2^{104}.
 \tag{PA8}
\]

This counts the original bounded lines only. No assertion about the number of eta values in an unrestricted completed family is used.

At each v in Sigma, DO105 EL9--15 fixes one admissible nonzero residue label \(c_\eta(v)\) depending only on eta and the original pencil. In the literal notation there, \(A_c=I+c\pi^{-1}zq_*\), and the two cancellation equations \(A_c^{-1}w\equiv0\) and \(w^\dagger A_c\equiv0\) give the same unique label. This uses both actual CM partners. It includes central-only poles. The label is admissible and nonzero because eta here comes from an existing R27 matrix; no existence of a new arrow is asserted. At all other places take the included standard label0. Denote this complete actual tuple by \(l_\eta\in D^{168}\), where \(D=\mathbb F_T\setminus\{-1\}\) and infinity is also omitted. Every R27 matrix with this line has exactly this diagonal incidence at every T place. At active places this is type II; at inactive places DO107 IP4--7 and \(b=0\) give the unique included-standard incidence. The tuple does not depend on the row or central scalar. (PA9)

Let \(U_\eta\) be the rational unipotent radical of the stabilizer of \(\ell_\eta\), and define

\[
 \Gamma_\eta=U_\eta(F)\cap
 \bigl(\prod_{v\nmid T}\operatorname{Stab}\Lambda_v\bigr)
 \cap\bigl(\prod_{v\mid T}\operatorname{Stab}L_{l_\eta(v)}\bigr).
 \tag{PA10}
\]

All stabilizers are the original ones. Their inverse-adjoint CM partners are part of the condition, not extra freely chosen labels. Every element of Gamma\_eta lies in the original finite support mathscr C: it has the displayed actual incidence at each T place. It is integral outside Sigma, so its actual nonintegral support is a subset of Sigma. In particular completion introduces no support greater than27. It does not preserve H. In particular the inactive label is exactly0, so its stabilizer is the original standard compact there, at both paired components. At the27 other places the fixed nonzero labels may be off-standard; this is why only containment of support in Sigma is initially claimed. Neither of the two omitted labels is ever used.

In the original decomposition \(E^{169}=Z\perp P\), \(Z=S_{164}\perp V_3\), every element of this radical has

\[
 h=I+N,\quad N=w\mathcal R-vw^\dagger+dww^\dagger,\quad
 \mathcal R|_P=0,\quad v=J_Z^{-1}r^*,\quad r=\mathcal R|_Z,
 \quad d+\bar d=-rv.
 \tag{PA11}
\]

If r is nonzero, positive definiteness of the quotient at an archimedean place gives \(N^2\ne0\), rank N=2 and intrinsic line im N²=ell\_eta. If r=0 and h is not the identity, rank N=1 and im N=ell\_eta. Therefore two Gamma\_eta groups intersect only at the identity. This uses archimedean positivity, not positivity of a residue-field Hermitian form.

\subsection{3. The positive comparison kernel, with its exact identity cost}\label{the-positive-comparison-kernel-with-its-exact-identity-cost}

For a local lattice L, its stabilizer characteristic function is the unitary permutation coefficient \(\langle U(g)\delta_L,\delta_L\rangle\) on the space of all local lattices. Define the auxiliary kernel

\[
 F_\eta(g)=m^{-168}a(g)c(g)
       \prod_{v\mid T}\mathbf1_{gL_{l_\eta(v)}=L_{l_\eta(v)}}.
 \tag{PA12}
\]

The away-T compact indicator is already included in a(g). K062 AT14--15 shows that these are precisely the diagonal finite-orbit vectors of the original family, at its one common real center. Thus PA12 uses no new source point, frame, finite representation or archimedean coefficient. It is a comparison kernel, not a replacement assertion about the original all-T projection P. No operator inequality \(P_\eta\le P\) is claimed.

Both a and c are normalized traces against their original positive finite-rank projectors in unitary representations; at the two mixed places the actual projectors are retained. Extension by zero from the original compact remains positive definite, by splitting Gram matrices into compact cosets. The lattice stabilizer coefficients are positive definite, and products of such coefficients are positive definite. Consequently PA12 is positive definite and

\[
                         F_\eta(1)=m^{-168}.
 \tag{PA13}
\]

At a unique diagonal incidence, the original local partial permutation matrix is \(M_g=E_{ll}\). The actual mean-zero projector is \(P=I-J/m\), of rank m−1, hence

\[
 \rho(g)=\frac{\operatorname{Tr}(PM_g)}{m-1}
        =\frac{1-1/m}{m-1}=\frac1m.
 \tag{PA14}
\]

Thus F\_eta equals F on every completed matrix with unique diagonal incidences at all168 places, in particular on R27. At a nonunique pencil this equality is not asserted: an included-fixed-line unipotent action can have \(\rho=1/[m(m-1)]\) while PA12 has1/m locally. Old absolute charges will not be transferred to those terms.

The series \(P_\eta=\sum_{g\in\Gamma_\eta}F_\eta(g)\) is absolutely convergent and nonnegative. Here is the full applicable group argument. The original incidence condition gives \(g,g^{-1}\in T^{-2}M_{169}(O_E)\) by K062 AT18. In a fixed rational Witt chart the logarithms \(N-N^2/2\) have bounded denominators. Their horizontal parts lie in a fixed finitely generated Z-lattice. Enlarge the central Z-lattice to contain their central parts and all half-brackets of horizontal generators. BCH produces a discrete finitely generated two-step nilpotent group containing Gamma\_eta. On this group retain PA12, zero outside its compact conditions. It vanishes outside Gamma\_eta, so its sum is the sum over that group even if further coefficients vanish. Bounded real pieces are finite, and the complete original absolute tail of c proves summability.

Use Følner boxes of horizontal width L and central width L². For every fixed g the relative overlap tends to1. The normalized Gram sum is nonnegative, since PA12 is positive definite, and equals \(\sum_g (|E_L\cap E_Lg^{-1}|/|E_L|)F_\eta(g)\). Dominated convergence proves \(P_\eta\ge0\). The finitely many groups give a finite real sum

\[
                         P=\sum_{\eta\in\Xi}P_\eta\ge0.
 \tag{PA15}
\]

No signed estimate is inferred merely from positivity of the restricted support mask. That mask still vanishes at the identity and is not positive definite in general.

\subsection{4. Exact partition of the completion and reconciliation of old charges}\label{exact-partition-of-the-completion-and-reconciliation-of-old-charges}

Within the disjoint nonidentity parts of the Gamma\_eta groups define:

\begin{itemize}
\tightlist
\item
  G: the noncentral elements with norm at most H at every embedding and unique diagonal incidences at all T places;
\item
  B: the noncentral elements with this norm cutoff and at least one nonunique pencil;
\item
  C: the central elements with this cutoff, \textbf{including the identity once for each eta};
\item
  L: all elements outside the cutoff, with their unique nonidentity line.
\end{itemize}

The letter L here denotes a set, not an ideal or local lattice. Set \(C_\Sigma=\sum_\eta\sum_{C\cap\Gamma_\eta}F_\eta\), and define \(B_\Sigma,L_\Sigma\) similarly. Then

\[
              P=\sum_{g\in G}F(g)+B_\Sigma+C_\Sigma+L_\Sigma.
 \tag{PA16}
\]

R27 is contained in G. Conversely every g in G has actual support exactly Sigma: at v in Sigma its admitted label is nonzero, so if g were integral it would also fix the included standard lattice, contradicting uniqueness. At all other places it fixes the standard lattice. Therefore G has support27 and all active incidences are actual type II. Its positive compression and global rank-two unipotence follow from PA11, also when its S row is zero. The complete fixed-S/scalar-cross taxonomy and current cuts then imply

\[
                      G\setminus R_{27}\subset E.
 \tag{PA17}
\]

This uses the full support27 coverage, not just a numerical k threshold. In particular zero S row is the fixed-S rank-two branch already removed by DO118. There is no unexamined rank-zero/rank-one escape in G because its row is nonzero and the positive quotient gives rank two.

All sums being absolutely convergent, PA16 and the original disjoint partition yield the exact identity

\[
 \mathcal E=1+\mathop{\rm Re}\sum_{R_{\ge28}}F+P+
       \underbrace{\mathop{\rm Re}\left(\sum_{E\setminus G}F
                      -B_\Sigma-C_\Sigma-L_\Sigma\right)}_{\delta}.
 \tag{PA18}
\]

The restored previously removed terms G\R27 cancel algebraically. They are not charged a second time. The full old absolute array remains a valid containing bound for E\G and is retained without deleting any entry. This identity differs from applying the separate inequality PA6 and then paying all old errors again.

\subsection{5. A nonunique pencil gives an extra paired ideal in the full row}\label{a-nonunique-pencil-gives-an-extra-paired-ideal-in-the-full-row}

Fix eta and a member of B. At a place with a nonunique actual incidence, the two complete pencils coincide: two shared lattices recover their intersection M and sum N, as in K062 AT7--9. If v is in Sigma, preservation of M puts every original integral Z input and its image in M subset Lambda. Since h\_ZZ=I, the whole cross block is integral. The inverse-adjoint paired modules give integrality at the other CM component as well. This is the full-Z argument of DO108 AC3 and DO106, not an S-only inference.

If v is outside Sigma it is noneligible, since total eligibility is27. Here g fixes Lambda. Saturate its line locally to a primitive w0 and its annihilating row to b0. The two transverse scalars q\_\emph{w0 and b0z are units, including eta zeros and poles. Full-pencil preservation and unipotence give Nz=0 and q\_}N=0 modulo pi. In the original decomposition these equations force the transverse vector, row and central coefficient all to vanish, exactly DO107 IP6--7. Thus both cross blocks are divisible by pi, after using the actual paired inverse adjoint. There is no independent CM choice or finite positivity assumption. (PA19)

Let J be the canonical whole-C G ideal, J0 the original central rank-zero F ideal, \(I_\Sigma\) the paired simple-pole ideal of norm T\^{}-54, and \(D=D_\eta\). DO126's actual selfdual projective quotient is \(\mathscr K_\eta=(\ell^\perp\cap\Lambda)/(\ell\cap\Lambda)\), with \(\ell\cap\Lambda=Dw\). Its integral standard summand is O\_E\^{}164 and its remaining selfdual summand has rank3. The same all-original-input argument at primes outside G/T, including every prime over2 or U, and the determinant-2U unit frame only at G/T, give

\[
 r\in I\mathscr K_\eta^\vee,\quad I=JI_\Sigma D,
 \qquad r\in I_v\mathscr K_\eta^\vee,
 \quad I_v=I(\mathfrak q_vO_E),\quad N_EI_v=T^2N_EI.
 \tag{PA20}
\]

At an active place the second inclusion removes its allowed simple pole; at an inactive place it adds a paired zero. At each component it follows from \(r(y),\eta r(y)\) in the improved entry ideal and the exact intersection with eta\^{}-1 of that ideal. At all other primes the old proof is unchanged. This is the entire quotient-row lattice, including its nonprincipal rank3 component, not just167 arbitrarily divided coordinates. Choose one such place; there are at most168\textless2\^{}8 choices.

Write \(z^{2f}=N_EI\), \(z_v=zT^{1/f}\), \(q=R/(Uz^2)\), \(q_v=q/T^{2/f}\). For a nonzero r, the positive quotient norm lies in \(I_v\bar I_v\cap F\), whose F-norm is N\_EI\_v. The product formula consequently gives

\[
                N_F\alpha\ge(Uz_v^2)^f=(R/q_v)^f.
 \tag{PA21}
\]

This is valid even when every S coordinate is zero. No per-embedding energy floor is imposed on the enlarged lattice points.

For each fixed row the original complete central difference coset is \(iJ_0B_\eta I_{\Sigma,F}^2\). It is unchanged by PA20. At active labels \(A_c^{-1}w\) and \(w^\dagger A_c\) each have exact valuation1; elsewhere primitive original-coordinate extraction gives the exact B\_eta condition. These are DO106 CFI4--8, with no numerical support restriction. At G one further retains J0 exactly as in DO127 GC3--4. The physical covolume remains

\[
             U^f\sqrt{D_F}(N_FJ_0)xT^{-54}.
 \tag{PA22}
\]

In particular the row improvement does not improve or alter the central cover, its real-part/possible dyadic shift, or the independent U\^{}-f.

\subsection{6. Complete noncentral complement count, including its zero term}\label{complete-noncentral-complement-count-including-its-zero-term}

Use the whole actual coarea/positive-mixture bound DO125 CM3--27 and DO126 IQ24--29, with the exact DO127 GC3--14 Green/J0 cancellation. These analytic bounds hold for every nonzero row in PA11 at H133; they do not require a nonzero S coordinate. Their dependence on k was in the lattice count and the genuine-row energy floor, which we specify anew. The coefficient is still the actual selected trace, not a scalar model.

Count all164 standard S coordinates freely, including zero, and retain the full selfdual rank3 Gaussian factor in real dimension1008. Thus use the entire DO127 k164 polynomial, with its330 positive terms, but replace q in the row polynomial by q\_v. In the repaired zero term replace the old164-coordinate energy floor by PA21's floor1. This is legitimate because the theta=2\^{}-20 Gaussian reserve is applied before enlarging the genuine nonzero row set. The same finite M1 is used; no infinite M0 is invoked. The inherited harmless factor164 in the complete prefactor is retained. Reconstruction from eta, the complete quotient row and d is injective, so there is no uncounted beta, column or source-isometry fibre.

The central profile and all its original constants, affine750 cover, physical Jacobian, near-point transport and moment enclosures are unchanged. The finite weight is now exactly m\^{}-168 on Gamma\_eta, so the same complete T correction applies even at the nonunique pencils; we do not compare it with their original rho. The exact full Green factor \(2Q^{-335\cdot167}(N_EJ)^{167}N_FJ_0\) still cancels J0 in PA22. For every effective degree1\textless=j\textless=167 the four-state product is still \(\prod(1+q_G^{j-167})^2\), including2\^{}668 at j167. (PA23)

Let E\_j denote each individual certified directed upper exponent in the complete DO127 k164, X2165 table. All330 entries satisfy E\_j\textless−746; the largest is −6113299/8192. Their exact all-identity-cap x dependence is x\^{}(-1-j), and their G-pattern comparison is nonincreasing in q, with equality allowed at degree167. Thus moving the containing x endpoint from2165 to2111 contributes exactly54(j+1) bits. Replacing q by q\_v in the row factor contributes −2j log2T. The zero term alone acquires 164\^{}f from the weaker energy floor. All these are multipliers of the same full positive term, not a comparison of partial envelopes.

The exact host satisfies T\textgreater2\^{}995, while164\^{}168\textless2\^{}1344. Including the chosen-place factor168\textless2\^{}8 gives

\[
 E_j^{\rm bad}<E_j+54(j+1)-1990j+8
        +1344\,\mathbf1_{\rm repaired\ zero}<-1276.
 \tag{PA24}
\]

Indeed for every ordinary positive term the additive change is at most −1874, and for the repaired zero term it is at most−530. The executable checks all330 individual entries and their index pattern. The actual largest transferred upper exponent is −1185269675/65536, at the repaired zero term; the simpler uniform−1276 is sufficient and is what is charged.

The omitted narrow-band tail must be checked for this auxiliary count. Its original uniform selected-probability error is less than2\^{}(9−20000000). The coarse BT16--20 counting proof applies directly to the present finite family: \textbar Xi\textbar\textless2\^{}104; all167 original fixed-input values lie in the larger I\_Sigma and have modulus\textless2\^{}40; their separate packing counts are\textless2\^{}67896; for a fixed row the original central affine coset is contained in a translate of i I\_Sigma,F² and its count is\textless2\^{}55464. The one label tuple per eta replaces, and cannot exceed, the old generic label-weight bound1. Retaining every harmless2\^{}668 pattern and164 pivot and Green2 gives the same containing mass\textless2\^{}11394877. This argument includes zero S entries and does not borrow a support25/26 numerical charge. It gives one full auxiliary band error\textless2\^{}-8605114, not one per witness place. Therefore

\[
 \sum_B|F_\eta(g)|<330\,2^{-1276}+2^{-8605114}
                    <331\,2^{-1276}<2^{-1267}.
 \tag{PA25}
\]

\subsection{7. Central terms, repeated identities and the single global outer tail}\label{central-terms-repeated-identities-and-the-single-global-outer-tail}

For a central element PA11 has r=0, d purely imaginary. Original primitive extraction and the fixed cancellation labels give the complete central lattice \(iB_\eta I_{\Sigma,F}^2\), contained in \(iI_{\Sigma,F}^2\). The identity lies in it. The non-strict H cutoff implies \textbar Im d\textbar\textless=H/U\textless2\^{}8 at every embedding; this weaker bound suffices. The larger real ideal has Euclidean minimum at least \(\sqrt f T^{-54/f}\). Packing disjoint half-minimum balls in dimension168 therefore gives, including zero,

\[
 \#C_\eta\le(1+2^9T^{54/f})^{168}
           <2^{1680}T^{54}<2^{55464},
 \tag{PA26}
\]

using T\textless2\^{}996. Since \textbar a\textbar,\textbar c\textbar\textless=1 and the actual comparison weight is m\^{}-168, the exact correction 2(T−1)\^{}168\textgreater T\^{}168 and PA8 give

\[
 \sum_C|F_\eta(g)|<2^{104+1+55464-168\cdot995}
                          =2^{-111591}.
 \tag{PA27}
\]

This includes the repeated identity debit; there is no additional unit identity cost per line. No rank-one Green decay is needed or silently used.

Every nonidentity element outside H occurs for at most one eta. The full R048-RC §§RC48.3--4, equations (RC48.6)--(RC48.7) and the following H133 paragraph give an estimate on \textbar c(g)\textbar{} itself over the entire original mathscr C, namely \(\sum_{\|g\|>133}|c(g)|<2^{-2000}\). All the completed points lie in that original support. Hence one use of that single global tail gives

\[
 \sum_L|F_\eta(g)|\le m^{-168}\sum_{\|g\|>133}|c(g)|
                   <2^{-169159}.
 \tag{PA28}
\]

There is no eta multiplier. The old original tail remains within PA2; PA28 is the additional, explicitly m\^{}-168-weighted subtraction required by the auxiliary completion. Equivalently both tail occurrences can be bounded by one full original \textbar c\textbar{} tail with multiplier1+m\^{}-168. This is not a separate unweighted tail payment for each line or each fibre. Here the unchanged a(g), rather than the archimedean c(g), carries every original away-T standard-compact indicator and the two mixed projectors. The only auxiliary finite modification is the displayed T stabilizer coefficient; its nonzero support supplies the same original incidence condition required for the global \textbar c\textbar{} tail.

Combining PA25,27,28 gives

\[
       |B_\Sigma|+|C_\Sigma|+|L_\Sigma|<2^{-1266}.
 \tag{PA29}
\]

PA18 proves PA4. Applying PA16 directly and bounding its additional G\R27 set by PA2 proves PA6. These are different statements with the different error accounting explicitly displayed.

\subsection{8. Verification, chronology and boundaries}\label{verification-chronology-and-boundaries}

The complete diagnostic verifies the actual host inequalities, all330 old fully certified positive terms and their new multipliers, the full pattern indices including j167, zero-degree floor change, both coarse packing exponents, line count and union. Finite original-PGL2 examples verify the exact diagonal normalization and demonstrate failure of a pointwise old-rho bound at a nonunique pencil. Valuation checks supplement PA19--22; they do not replace their all-prime written proof.

Its default predecessor is checks/k123\_joint\_remainder\_error.json, with all138/31 exponents and floors27/11 retained. No absolute charge is added. The separate1266 one-sided debit and the exact inequality epsilon\_old+2\^{}-1266\textless2\^{}-230 are checked without relabeling that debit as an absolute estimate. The script accepts a later complete extension of the same ordered prefix, provided this exact combined inequality still holds. The frozen DO128 no-charge statement and all earlier author histories remain unchanged.

This theorem advances a signed higher-support gate. It proves neither that the current support27 sum is nonnegative nor that its absolute value is small. It establishes no global period, nonzero selected arithmetic content, source admission, effective DP44 cost, specific032 value or RH conclusion. The actual projective modules, Hermitian polarization, theta line, finite orbit and complex selected metric remain the same throughout.

\subsubsection{Exact dependencies}\label{exact-dependencies}

K062 AT1--25: actual diagonal orbit, projector, kernel and masses. DO105 EL9--16: eta-dependent paired cancellation label. DO106 CFI2--10: complete original central ideal at all primes. DO107 IP4--10 and DO108 AC3/AC8--10: full-Z/paired-zero implications. K113 SC27.1--10: bounded eta uniqueness and support27 ideal scope. DO126 IQ4--34: original selfdual quotient, all-prime row membership, full1008 theta and the genuine-row zero-term repair. DO127 GC1--16 and its accepted complete certificate: exact Green/J0 correlation, all330 k164 terms, continuum and pattern factors. K118 BT15--22 / K120 NB14--16: uniform narrow-band error and the complete coarse containing mass. R048-RC §§RC48.3--4, equations (RC48.6)--(RC48.7) and the following H133 paragraph: the original global \textbar c\textbar{} tail at the non-strict H133 cutoff. R097 with all subsequent cuts through K123: the complete current support27 taxonomy, not merely an endpoint table. DO128 HZ13--19: separately checked full-group positivity mechanism and the obstruction that PA12--29 explicitly resolve for this finite family.

The original finite character input remains Gorodetsky, O., \& Hazan, Z. (2017), \emph{On certain degenerate Whittaker models for cuspidal representations of GL\_kn(F\_q)}, Theorem2.1/Corollary2.2, printed p.5, \url{https://arxiv.org/abs/1707.07308}, in the accepted DO127 scope. No new whole-paper, FGA/SGA, Lean or archive-validity audit is claimed here.

\clearpage

\section{AM-LIFT-109 --- CM norm surjectivity and the necessary pinned-field degree}\label{am-lift-109-cm-norm-surjectivity-and-the-necessary-pinned-field-degree}

Status: author development; final status is supplied by the separate hash-specific author packet.

\subsection{1. Literal hypotheses and result}\label{literal-hypotheses-and-result}

Retain the geometric Host, Datum, Level and Coeff interfaces used in AM-LIFT-107 and AM-LIFT-108. In particular, E is a CM number field, d is a rational prime with d congruent to 7 modulo 16, and K = Q(rt) is the imaginary quadratic subfield of E with rt squared equal to -d.~The chosen embedding iota: K -\textgreater{} F sends rt to rt0, where F = c.F is the original coefficient field. Coeff.fin makes F a number field and Coeff.le/Coeff.conj retain this embedding. Neither F/K nor E/K is assumed Galois. In particular, F is not assumed CM. We identify K with iota(K) when discussing F. For intersections inside C, choose an embedding of E extending iota and identify E with its image. Such an extension exists because C is algebraically closed; Coeff.conj puts this image inside F. Transport along this embedding does not change the ideal norm or the class-group statements.

Write

\[
 A_K=\operatorname{Cl}(\mathcal O_K),\qquad
 h_K=|A_K|,\qquad e_K=\exp(A_K),\qquad D_F=[F:K].
 \tag{NS1}
\]

The class group here is the ordinary ideal class group of the maximal order. Since K is imaginary quadratic, no narrow-class qualification is needed.

For a nonzero integral E-ideal I, let Phi(I) be the product of its images under the embeddings tau satisfying tau(rt)=rt0, extended to O\_F, exactly as in LIFT107 AG3 and LIFT108 NI5. The normalized rank-one condition is

\[
 \begin{split}
 \mathcal C(H,c):\quad \forall I\ne0\ \exists\chi_I\in F^\times:
 \quad &(\chi_I)=\Phi(I),\\
 &|\sigma(\chi_I)|^2=N_{E/\mathbf Q}I
       \quad\text{for every }\sigma:F\hookrightarrow\mathbf C.
 \end{split}
 \tag{NS2}
\]

For the actual original signatures and PinDat, LIFT107 proves that NS2 is equivalent to the two original stage-1 and stage-2 KPsiIn laws, and to all original stage laws. It does not require any compatibility between the phases of their chosen witnesses.

\textbf{Theorem.} Under these literal field hypotheses:

\[
 N_{E/K}: \operatorname{Cl}(\mathcal O_E)\longrightarrow A_K
       \quad\text{is surjective}.
 \tag{NS3}
\]

Consequently the group G in LIFT108 is the full group A\_K, not a possibly smaller group in this actual CM, prime-discriminant setting. If NS2 holds in the original F, then

\[
                         e_K\mid D_F.
 \tag{NS4}
\]

More generally, if a finite extension F\_+/F supplies NS2 with the same Host and quadratic embedding, then

\[
       \frac{e_K}{\gcd(e_K,D_F)}\ \bigm|\ [F_+:F].
 \tag{NS5}
\]

If A\_K is cyclic and gcd(h\_K,D\_F)=1, LIFT108's normalized new-coefficient construction of relative degree at most h\_K has relative degree exactly h\_K. This is a conditional optimality statement for that changed-field problem. It neither constructs the original source record nor transports its PinField.

For a full pinned package satisfying its original kPsi law, a further elementary consequence is {[}F:K{]} at least seven and {[}F:Q{]} at least fourteen. More precisely, the order of either ideal class above two is an odd divisor of {[}F:K{]} and its two-power is at least (d+1)/4. Section 6 proves this without a class-number classification or an actual class-group computation.

\subsection{2. Prime discriminant and the Hilbert class field}\label{prime-discriminant-and-the-hilbert-class-field}

The actual congruence and primality give

\[
  \mathcal O_K=\mathbf Z[(1+\sqrt{-d})/2],
  \qquad \operatorname{disc}(K)=-d .
 \tag{NS6}
\]

This is the standard quadratic-integer calculation in Milne (2020a, Example 2.41), already applied in LIFT108. There is exactly one prime divisor of this fundamental discriminant. Gauss's genus theorem, in Cox (1989, Theorem 6.1(iii)-(iv), pp.~121-127), therefore gives

\[
          A_K/A_K^2=1,\qquad h_K\text{ is odd}.
 \tag{NS7}
\]

Indeed, the quotient has order 2 to the power 1-1; squaring is surjective on the finite abelian group A\_K and hence bijective, so the group has no element of order two. This applies the cited genus theorem; it does not present a new proof of genus theory.

Let H\_K be the Hilbert class field of K, in the chosen algebraic closure in C. The Artin isomorphism and the standard quadratic conjugation calculation give

\[
 \operatorname{Gal}(H_K/K)\simeq A_K,\qquad
 \operatorname{Gal}(H_K/\mathbf Q)
       \simeq A_K\rtimes\langle c\rangle,\qquad
 cac^{-1}=a^{-1}.
 \tag{NS8}
\]

Here c is actual complex conjugation, of order two, not an arbitrary lift of the nontrivial quadratic automorphism. Cox Lemma 5.28, p.~110, proves H\_K/Q is Galois; the proof of Theorem 6.1, Eq. (6.4) and its following paragraph, pp.~122-123, identifies the conjugation action with inversion. These are the exact standard results used in NS8.

Set L = E intersection H\_K. It contains K. As a subfield of H\_K containing K, it corresponds to a subgroup S of A\_K. Every such subgroup is stable under inversion and hence normal in the group in NS8. Thus L/Q is Galois, and

\[
 \operatorname{Gal}(L/K)\simeq A_K/S,\qquad
 cbc^{-1}=b^{-1}\quad(b\in A_K/S).
 \tag{NS9}
\]

The CM hypothesis on E is essential at this point. To see explicitly what it supplies, take a Galois closure E' of E. Milne (2020b, Proposition 1.4 and Corollary 1.5, p.~10) show that E' is CM and has a unique complex conjugation. This conjugation is central in Gal(E'/Q): its conjugate by any automorphism is again the unique complex conjugation. It therefore preserves every subfield of E', including L. For every complex embedding of L, its restriction is the usual complex conjugation. It is nontrivial on K, so L is CM by the same Proposition 1.4. Since L/Q is Galois, its complex conjugation is central.

NS9 now says that every b in A\_K/S equals its inverse. This quotient has odd order by NS7, and an odd-order group of exponent at most two is trivial. Therefore

\[
                              E\cap H_K=K.
 \tag{NS10}
\]

This argument does not assume E/Q or E/K is Galois. The Galois closure is used only to justify the inherited CM conjugation.

\subsection{3. The full Frobenius and residue-degree argument}\label{the-full-frobenius-and-residue-degree-argument}

Let C=E H\_K. Because H\_K/K is finite Galois, C/E is finite Galois, and restriction gives an injection

\[
 \operatorname{Gal}(C/E)\hookrightarrow
                       \operatorname{Gal}(H_K/K).
 \tag{NS11}
\]

The fixed field in H\_K of its image is H\_K intersection E: an element of H\_K fixed by all automorphisms of C/E lies in E, and the converse is immediate. By finite Galois correspondence and NS10, the image in NS11 is the full group. Thus restriction is an isomorphism. In particular C/E is abelian, including when E/K is not Galois.

Fix any a in A\_K and let sigma in Gal(C/E) correspond under NS11 to its Artin image in Gal(H\_K/K). Chebotarev in the finite abelian extension C/E supplies a prime ideal P of E, unramified in C, whose Frobenius is sigma. It may be chosen outside any prescribed finite set of E-primes. This is the abelian case of Milne (2020c, V, Theorem 3.23, p.~164). No effective bound for the norm of P is being used or claimed.

Put p=P intersection O\_K, q=\textbar O\_K/p\textbar{} and f={[}O\_E/P:O\_K/p{]}. Choose a prime Q of C above P, and write q\_H=Q intersection O\_H\_K. On the residue field of Q, Frobenius for C/E raises elements to

\[
                   |O_E/P|=q^f.
 \tag{NS12}
\]

Its restriction to H\_K preserves q\_H and therefore induces the q-to-the-f power map on its residue field. Because H\_K/K is unramified, uniqueness of the Frobenius element gives

\[
 \begin{split}
 \sigma|_{H_K}
   &=\operatorname{Frob}_{p}(H_K/K)^f\\
   &=\operatorname{Art}_{H_K/K}(p^f)
     =\operatorname{Art}_{H_K/K}(N_{E/K}P).
 \end{split}
 \tag{NS13}
\]

The last equality is the ideal-norm formula N\_\{E/K\}P=p\^{}f, with the residue degree, not the ramification index. This formula and its compatibility with principal ideals are Sutherland (2021, Theorem 6.10 and Proposition 6.7) and Milne (2020c, V, p.~157). The identical Frobenius calculation is written for intermediate abelian fields in Milne V, Proposition 3.3, p.~157; NS12-NS13 explicitly perform it in the compositum C for the present possibly non-Galois E/K.

Under the Artin isomorphism A\_K -\textgreater{} Gal(H\_K/K), NS13 says {[}N\_\{E/K\}P{]}=a. Norms of principal ideals are principal, so the ideal norm descends to a map of ideal class groups. Since a was arbitrary, this proves NS3. The argument also proves surjectivity using ideals supported away from any fixed finite set; it supplies no height or prime-size bound. Ramified primes in E/K are permitted throughout the theorem. They can simply be included in the finite set avoided in this existence argument.

\subsection{4. Full-class normalized table and the original field}\label{full-class-normalized-table-and-the-original-field}

LIFT108 NI5 identifies the actual CM-type ideal with the relative norm:

\[
                \Phi(I)=\iota(N_{E/K}I)\mathcal O_F.
 \tag{NS14}
\]

For completeness, the identification uses all embeddings of E extending iota, hence the ordinary field norm on principal ideals. Passing to a principal power of I and cancelling that power in the free fractional-ideal group gives NS14 for all I, including those above ramified primes. It is the same identity already proved in LIFT108, not a substitute type.

If NS2 holds, let J be any nonzero integral K-ideal. By NS3 choose an integral E-ideal I whose norm has class {[}J{]}. Then N\_\{E/K\}I=(z)J for some z in K multiplicative. Set xi=chi\_I/iota(z). NS14 and the quadratic norm formula give

\[
   (\xi)=\iota(J)\mathcal O_F,\qquad
   |\sigma(\xi)|^2=N_{K/\mathbf Q}J
           \quad\text{for every }\sigma:F\hookrightarrow\mathbf C.
 \tag{NS15}
\]

The normalization follows because every embedding of K is one of its two complex conjugate embeddings, so \textbar sigma(iota(z))\textbar{} squared equals \textbar N\_K/Q(z)\textbar, while N\_E/Q(I)=\textbar N\_K/Q(z)\textbar{} N\_K/Q(J).

Thus LIFT108's normalized table is a table for the entire quadratic class group under the present source hypotheses. Conversely such a normalized full-class table supplies NS2 by LIFT108 NI8-NI11. The same-field equivalence remains archimedean as well as ideal-theoretic.

Now apply the ideal norm from F to iota(K) to the first identity in NS15. Standard extension followed by norm gives

\[
 \iota(J)^{D_F}
   =N_{F/\iota(K)}\bigl(\iota(J)\mathcal O_F\bigr)
   =\bigl(N_{F/\iota(K)}\xi\bigr).
 \tag{NS16}
\]

Sutherland (2021, Corollary 6.11, p.~5) proves this with all ramification present: at a prime the exponent is the sum of e\_P f\_P, equal to {[}F:K{]}. No unramified or Galois hypothesis on F/K enters. Therefore {[}J{]} raised to D\_F is trivial for every class {[}J{]}, proving NS4.

Equivalently, for the actual original signature family,

\[
 \bigl(\mathrm{KPsiIn}_1\ \&\ \mathrm{KPsiIn}_2\bigr)
 \quad\Longrightarrow\quad
 e_K\mid[F:K].
 \tag{NS17}
\]

The premise can be replaced by all original stage laws or the actual LawsP.kPsi field. It cannot be replaced here by a KPsiIn law of an arbitrary single Datum: a balanced signature, for example, can have a rational norm witness without supplying NS2. The adjacent original-signature reduction is precisely why NS17 applies to the named law.

\subsection{5. Changed-field lower bound and the cyclic optimum}\label{changed-field-lower-bound-and-the-cyclic-optimum}

Suppose F\_+/F is finite and supplies NS2 with the same Host and quadratic embedding. Applying NS4 to F\_+ and the tower formula gives

\[
                     e_K\mid D_F[F_+:F].
 \tag{NS18}
\]

Let g\_0=gcd(e\_K,D\_F). After writing e\_K=g\_0 e\_0 and D\_F=g\_0 D\_0 with gcd(e\_0,D\_0)=1, Euclid's lemma gives e\_0 dividing {[}F\_+:F{]}. This is exactly NS5. No condition on the sign, discriminant, normality or CM nature of F\_+ has been added.

If A\_K is cyclic, then e\_K=h\_K. If also gcd(h\_K,D\_F)=1, every field supplying NS2 over F has degree at least h\_K and in fact degree divisible by h\_K. LIFT108's explicit class-power-root construction supplies a new Coeff field F\_+ with degree at most \textbar G\textbar=h\_K. Hence that field satisfies

\[
                              [F_+:F]=h_K.
 \tag{NS19}
\]

The construction supplies the normalized generators, so NS19 is not inferred from principalization alone. The upper bound comes from that construction and the lower bound from NS18. For a noncyclic class group this proof claims only the exponent lower bound, which can be strictly smaller than h\_K.

When h\_K=1, e\_K=1, NS4-NS5 are vacuous and LIFT108 already supplies NS2 in the original F. A degree-one extension is then allowed. None of these boundary cases turns d=7 into an admitted full arithmetic Host: the previously checked qd\_big and PinHost exclusion remains.

\subsection{6. The split prime above two gives an explicit order bound}\label{the-split-prime-above-two-gives-an-explicit-order-bound}

Put w=(1+sqrt(-d))/2. Since d is congruent to 7 modulo 16, its exact integral polynomial has reduction

\[
  X^2-X+\frac{d+1}{4}\ \equiv\ X(X-1)\pmod2.
 \tag{NS20}
\]

Using the integral basis in NS6, the two evaluation maps O\_K -\textgreater{} F\_2, with w sent to zero and one respectively, have distinct kernels

\[
  \mathfrak p=(2,w),\qquad
  \overline{\mathfrak p}=(2,w-1),\qquad
  N\mathfrak p=N\overline{\mathfrak p}=2,\qquad
  \mathfrak p\overline{\mathfrak p}=(2).
 \tag{NS21}
\]

The last equality can also be checked without a splitting criterion: all products of the displayed generators are in (2), while 2w-(2w-2)=2 is in their product. Conjugation sends w to 1-w and exchanges the two ideals.

Let h\_2 be the order of {[}p{]} in A\_K. It is a positive odd integer by NS7. Since p to the power h\_2 is principal and integral, write

\[
   \mathfrak p^{h_2}=(\alpha),\quad
   \alpha=\frac{a+b\sqrt{-d}}2\in\mathcal O_K,\quad
   a,b\in\mathbf Z,\quad a\equiv b\pmod2,\quad
   \frac{a^2+d b^2}{4}=2^{h_2}.
 \tag{NS22}
\]

Here the norm equality is the ordinary positive quadratic norm, using N(p)=2. One has b not equal to zero. Otherwise alpha is rational and integral, so (alpha) is preserved by conjugation. This would give p to the power h\_2 equal to pbar to the power h\_2, contrary to unique factorization of ideals and the distinct primes in NS21. Equivalently, a rational principal ideal has equal valuations at p and pbar, while the ideal in NS22 does not.

If b is odd, then a is odd, so a squared is at least one and b squared is at least one. If b is even and nonzero, b squared is at least four. Both cases imply

\[
                    2^{h_2}\ge\frac{d+1}{4}.
 \tag{NS23}
\]

In the even case the stronger lower bound is d, which is greater than (d+1)/4 because d is at least seven. No least-prime or class-number theorem enters this step.

For an original field satisfying NS2, NS4 yields

\[
       h_2\mid e_K\mid D_F,\qquad
       h_2\text{ odd},\qquad
       d+1\le4\,2^{D_F}.
 \tag{NS24}
\]

Thus any proposed upper bound for the original degree D\_F gives a finite necessary upper bound for d.~The integer h\_2 itself has not been evaluated here.

Finally the literal full ranks have m at least two and k at least m+3, with k and m even. Hence k is at least six and r=k+m-1 is at least seven. The condition r+5 at most t, together with odd t, gives t at least thirteen. Full Host.qd\_big and PinHost therefore give

\[
 d>4tr+1\ge365,\qquad
 2^{h_2}\ge(d+1)/4>91.5>2^6,
 \qquad h_2\ge7,\quad [F:K]\ge7,\quad [F:\mathbf Q]\ge14.
 \tag{NS25}
\]

These are necessary conditions on a supplied full pinned source satisfying its original kPsi laws. They do not assert that a d in the remaining range admits all Host, Sizes, lattice, level or selected-image data.

\subsection{7. Exact source consequences and unchanged costs}\label{exact-source-consequences-and-unchanged-costs}

\begin{longtable}[]{@{}
  >{\raggedright\arraybackslash}p{(\columnwidth - 4\tabcolsep) * \real{0.3333}}
  >{\raggedright\arraybackslash}p{(\columnwidth - 4\tabcolsep) * \real{0.3333}}
  >{\raggedright\arraybackslash}p{(\columnwidth - 4\tabcolsep) * \real{0.3333}}@{}}
\toprule\noalign{}
\begin{minipage}[b]{\linewidth}\raggedright
Item
\end{minipage} & \begin{minipage}[b]{\linewidth}\raggedright
Consequence proved here
\end{minipage} & \begin{minipage}[b]{\linewidth}\raggedright
Remaining requirement
\end{minipage} \\
\midrule\noalign{}
\endhead
\bottomrule\noalign{}
\endlastfoot
Actual CM field E and prime discriminant -d & G=Cl(O\_K) & A literal admitted Host and all its ties \\
Original coefficient field F & e\_K divides {[}F:K{]} is necessary for the original kPsi law & Normalized generators in that very field \\
Split-prime order h\_2 & h\_2 is odd, divides {[}F:K{]}, and 2 to the h\_2 is at least (d+1)/4 & Its actual value is not computed \\
Full pinned Host with original kPsi & {[}F:K{]} at least 7, {[}F:Q{]} at least 14, and d+1 at most 4 times 2 to the {[}F:K{]} & No sufficiency or nonempty admitted family \\
A proposed extension F\_+/F & e\_K/gcd(e\_K,{[}F:K{]}) divides its degree & Actual new Coeff and reconstruction of any changed pinned records \\
Cyclic class group, coprime old degree & LIFT108's relative degree h\_K is optimal & The hypotheses must hold for the proposed actual field \\
Ideal principalization & Necessary consequence of normalized witnesses & It is not substituted for their absolute-value conditions \\
Level and bad primes & Norm-class representatives can avoid any fixed finite set & No effective norm bound for those E-prime representatives \\
Selected betaH, DeltaDat, what2/eO/PinField & No alteration & Existing selected records remain prescribed \\
Quantitative comparison & No DP44 term is removed or given a new sign & All real, finite, period, line and discriminant costs remain \\
Numerical sector accounting & No new charge; K123 accounting reference: 138/31, floors27/11 & Subsequent numerical prefixes remain separate \\
Main arithmetic conclusion & No selected-image, global-period, Read or 032 nonzero assertion & Those main gates remain open \\
\end{longtable}

The algebraic-geometric content of NS16 is the composite of pullback and norm on invertible ideal lines over the arithmetic curves Spec(O\_K) and Spec(O\_F): on their Picard groups it is multiplication by the extension degree. The degree test is an obstruction to a prescribed-field normalized trivialization; it is not a construction of that trivialization. No motive, weight or RH consequence is asserted.

\subsection{8. Exact literature and verification scope}\label{exact-literature-and-verification-scope}

Literal sources are the actual Host declaration, including its IsCMField parameter, in \nolinkurl{TwoWeightEqOneAddRoot.lean14-26;} Coeff40-46 in the recovered PushInvertible.lean; EmbKCoe.lean 25-33; ThShort.lean277-281; HasDataCondition.lean55; and IsDatW/PinHost/PinField/PinDat in \nolinkurl{PartitionExtendByZeroAntitone.lean24-49.} The fixed-field equality consumer is \nolinkurl{MapRingEquivTransApply.lean52-69.} The explicit full-Host bound uses Params.lean6-22 and32, and DataHost.lean39-54, including qd\_big54. These are accepted source interfaces applied as stated. No accepted Alpoge proof is re-derived or validity-audited.

LIFT107 is the canonical accepted adjacent-signature equivalence at agents/adjacent\_signature\_same\_field\_type\_norm.md, SHA256 8b2ff2ea66e277a11600c2f5f2166ad36633b3c2238e0b9800fd376a4dab2694. LIFT108 is the accepted normalized norm-image and new-coefficient theorem at agents/quadratic\_norm\_image\_coefficient\_law.md, SHA256 ff062862f62e2e1bd79a48eb4b4f7d866450e068a52fa9ba469d481b7e145de6, accepted at14:10:03.940004 UTC after independent whole reviews. The original frozen author proof remains unchanged.

\subsubsection{Primary references}\label{primary-references-4}

Cox, D. A. (1989). \emph{Primes of the form \(x^2+ny^2\): Fermat, class field theory, and complex multiplication}. Wiley. Exact inputs: Theorem 5.23, p.~109; Lemma 5.28 and its proof, p.~110; Theorem 6.1 statement and genus quotient, pp.~121-127; Eq. (6.4) and the generalized-dihedral calculation, pp.~122-123. The retained scanned first-edition PDF is the University of Toronto copy; the source record distinguishes the retrieved PDF from the web OCR passages actually read.

Milne, J. S. (2020a). \emph{Algebraic number theory} (Version 3.08). Example 2.41, p.~39, for the exact maximal quadratic order and discriminant.

Milne, J. S. (2020b). \emph{Complex multiplication} (Version 0.10). Proposition 1.4 with its proof, Corollary 1.5 with its proof and Remark 1.6, p.~10, for unique conjugation and CM Galois closure.

Milne, J. S. (2020c). \emph{Class field theory} (Version 4.03). V, Artin map and Proposition 3.3 with proof, pp.~156-157; Theorems 3.5-3.6 and Example 3.9, pp.~158-159; Theorem 3.23 and its displayed abelian-case proof, p.~164. The stated standard class-field and genus theorems are used as literature inputs; their entire books are not claimed read.

Sutherland, A. V. (2021). \emph{18.785 Number Theory I, Lecture 6: Ideal norms and the Dedekind-Kummer theorem}. MIT. Proposition 6.7, Theorem 6.10 and Corollary 6.11, pp.~4-5; Proposition 6.12 and Remark 6.13, pp.~7-8.

The companion exact diagnostic checks finite generalized-dihedral models, residue-degree powers, norm/extension exponents, the necessary divisibility quotient, cyclic optimality arithmetic and the split-two quadratic-norm estimate. Even-order countermodels retain the prime-discriminant boundary. Diagnostics do not compute the class group, exponent, normalized table or coefficient degree of an admitted source, and do not replace the universal proof above.

\clearpage
\part{Source audit and geometric dictionary}

\clearpage

\section{Source audit: why CMData plus generic thetaG does not instantiate the mixed family}\label{source-audit-why-cmdata-plus-generic-thetag-does-not-instantiate-the-mixed-family}

Run ALP\_MACHINE, \nolinkurl{/root/lift_accounting}, 21 November 2026 UTC. This is an exact source audit and refinement of the current gate, not a new theorem replacing accepted DP52 or DP57. The generic-interface distinction is already stated immediately after the proof of DP52 in the intake. No Lean execution is claimed.

\subsection{The field repair is real but has a deliberately narrower target}\label{the-field-repair-is-real-but-has-a-deliberately-narrower-target}

Accepted DP57 proves that \(E'=\mathbb Q(\zeta_{28})\), \(F'=(E')^+\), and \(A=J^2\) furnish the indicated integral CM action, principal product polarization, unramified CM field condition, and an actual generic adelic endpoint. In particular it really meets \nolinkurl{CMData}, whose four fields are \nolinkurl{cmUnram}, \nolinkurl{totReal}, \nolinkurl{deg2}, and \nolinkurl{fixed} in \nolinkurl{Formalism/LocalRingCardinality/DataSpaces.lean:14–19}. The geometric type is induced from no imaginary quadratic subfield. DP57 explicitly does not construct the old \nolinkurl{Datum.dom}/\nolinkurl{CMObj.HasType} input; its separate old \nolinkurl{8 ∣ f} wrapper also fails since \(f=6\). Those are carried findings, not new obstructions derived here.

\nolinkurl{CMData} is sufficient for the generic \nolinkurl{Spaces} layer, subject to its other geometric inputs. It is not the type of the input of the source \nolinkurl{thetaG_exists} theorem. That theorem begins with \nolinkurl{variable (D : D7)}.

\subsection{The actual input to thetaG is substantially larger}\label{the-actual-input-to-thetag-is-substantially-larger}

\nolinkurl{Formalism/LocalRingCardinality/RhoDetPrimeEqMuTildeMulChiPow.lean} defines \nolinkurl{D7} from line 76. Besides \nolinkurl{Ranks} and \nolinkurl{Spaces}, it includes:

\begin{itemize}
\tightlist
\item
  The full old \nolinkurl{Host}, \nolinkurl{Rows.Data}, all member column chains, their periods, fixed coefficient fields and integral row/column lattices.
\item
  Exact ties between \nolinkurl{perCol} and the member \nolinkurl{PtilCol} maps, both at the base tuple and at twisted tuples; actual lattice embeddings and compatibility with archimedean norms.
\item
  The local theory, local comparison systems, theta-jet data, finite places, stage characters, their arithmetic and Galois compatibilities, and the remaining global data.
\end{itemize}

This is not an inference from the name \nolinkurl{D7}: these are visible fields in its declaration. For example \nolinkurl{perCol_base_eq_PtilCol} appears at lines 169--171, \nolinkurl{perCol_tie} at 181--183, and the theta-jet block starts at line 311 in the recovered copy.

The old \nolinkurl{Host} itself extends \nolinkurl{Formalism.HC.HostAx}, rather than \nolinkurl{CMData} (\nolinkurl{DataHost.lean:39}). Its predecessor \nolinkurl{HostRaw} includes the special Galois/cyclotomic field package and \(\sqrt{-d}\). Its CM type is literally the set of embeddings taking that element to \(i\sqrt d\): \nolinkurl{Formalism/SharpRestriction/RamificationFrobenius/TotientGaloisClasses/ConjEMemHostFIff.lean:66}. This is an induced imaginary-quadratic type. Supplying DP57's generic field and geometric CM action therefore does not, by itself, supply this old arithmetic \nolinkurl{D7} package.

\subsection{What the generic constructor actually proves}\label{what-the-generic-constructor-actually-proves}

\nolinkurl{Formalism/LocalRingCardinality/SpanPowerMembership/ThetaFamilyAdmissibleThree.lean} defines one auxiliary Hilbert space \(HG=HR\oplus HC\oplus\mathbb C\) at line 16. Its row embeddings land in the first factor, its column embeddings in the second, and \nolinkurl{v0} in the last (\nolinkurl{rowJ}, \nolinkurl{colJ}, \nolinkurl{v0}, lines 18--33). The \nolinkurl{thetaG} structure is at lines 71--153: its automorphic spaces all equal \nolinkurl{HG}, the cusp/type spaces are the full space, group actions are identity, Schwartz spaces are all endomorphisms, and periods are inner products. The specified lowering operator is \nolinkurl{rowJ ∘ io ∘ rowQ}.

For an already supplied \nolinkurl{D7}, the final theorem at lines 277--279 gives exactly

\begin{verbatim}
∃ X : ThetaX D,
  TieRowIsom X ∧ TieIo X ∧ AdmS D X ∧
  TieRowTyp X ∧ TypStable X.
\end{verbatim}

It proves all six components of \nolinkurl{AdmS} at lines 215--274. The Rallis representation used for this construction is the line spanned by the extra vector \nolinkurl{v0} (line 216). In particular the construction's existence does not prove that the actual abelian theta germs occur in its prescribed row/column images. This is an auxiliary realization of the interface, already recognized as such by accepted DP52.

There is a further explicit omission: the theorem does \textbf{not} supply \nolinkurl{TieBm}. In \nolinkurl{DataTheta.lean:747–752}, that predicate identifies the actual terminal bilinear period \nolinkurl{D.what₂.perCol} with the specific \nolinkurl{Btop} expression on the conjugated column and embedded terminal row, including its \nolinkurl{mY} factor. \nolinkurl{TieRowIsom} and \nolinkurl{TieIo} are distinct predicates at lines 739--745. Their proofs do not establish this terminal identification. No claim that \nolinkurl{TieBm} is identically zero or impossible for the generic construction is made: its noncanonical conjugation is not silently assumed to preserve the three factors.

\subsection{Exact conclusion for the proposed shortcut}\label{exact-conclusion-for-the-proposed-shortcut}

The proposed implication ``DP57 supplies CMData for J², and thetaG supplies the six identities, so only one finite coefficient matrix remains'' does not follow from these sources. The theorem providing thetaG requires the still-unconstructed arithmetic \nolinkurl{D7} first and its conclusion omits \nolinkurl{TieBm}; more importantly its theta maps have not been identified with the classical mixed-CM correspondence or the abelian theta kernels. The numerical source \nolinkurl{SizeData₂} also remains tied to the selected arithmetic derivatives and lattice errors, rather than to the bare CM orbit certificate.

There are useful concrete pieces: accepted DP52/DP57 provide the generic endpoint, AM-R-004 supplies the actual analytic CM torus and theta germ, ALP-LIFT-05 gives its full stable Hodge normalization, ALP-LIFT-06 computes the directional bare-certificate/Cauchy equality, and ALP-LIFT-07 fixes its exact characteristic and theta-null finite normalization. They reduce the unknown analytic normalization and identify a genuine source-line mismatch. They do not yet fill the mixed arithmetic constructor, the actual six identities on its tied theta family, or the final quantitative selected-value comparison.

The smallest honest next local target is to construct the prescribed trace-one section and allowed odd-level translates in explicit CM coordinates, compare their reduced gauges to the integral connection lattices at the selected primes, and identify the resulting finite coefficient matrix with the actual theta-lift row/column map. Even that local matrix identification is not a substitute for all global constructor and norm hypotheses. This source audit makes no universal nonexistence assertion about mixed-type theta correspondences.

\clearpage

\section{Working audit: the remaining fixed G/Q seesaw tensor and its frame scalar}\label{working-audit-the-remaining-fixed-gq-seesaw-tensor-and-its-frame-scalar}

Author: GPT-6 Astra Pro, ALP\_MACHINE /root; effort not separately exposed. 22 November 2026 UTC. \textbf{Historical working audit, not a promoted finding.}

\textbf{Resolution after full checks:} R034 resolves the passive Q scalar and entire local Q tensor. DO052 resolves the distinction between the uncut G kernel and the generic function, the two normalized averages, and the required diagonal source frame with its honest Weil scalar. The old conditional and unresolved wording below is retained as a record of the approach and its correction; it is not the current theorem or an unresolved proof promoted to the checked part. Full global row/type and arithmetic normalization pins remain separate in the current register.

\textbf{Review correction, 22 November 2026, 06:57--07:00 UTC.} The G identification in section2 is conditional on a source-vector adaptation that the earlier draft had not supplied. The unchanged uncut vertical kernel in LIFT31(VA11) and R032(R32.5) is the integral residue identity cell. LIFT22(CG7--CG9) is the sharp horizontal kernel with a generic U\_a sum and the source \nolinkurl{entB} last-column scaling by q. A frame congruent to1 cannot create either operation. The exact source \nolinkurl{GAt.frame} and \nolinkurl{GAt.Vec} require a moved-frame equation with \nolinkurl{bas_0=diag(I_a,q)} and the generic \nolinkurl{lvP} vector. The team is deriving the normalized generic projection and honest Weil multiplier, with their actions on the unchanged P-character vertical map. No G common-tensor claim has been promoted from this audit.

The Q passive-model scalar is being written separately with the unchanged raw Q vector and normalized vacuum. It must not be used to hide or settle this distinct G adaptation.

This records the exact interface readings and the remaining scalar question before checkpoint15. It neither reopens the settled local GGP tests nor identifies them with the finite Schwartz identities considered here.

\subsection{1. Exact common local Fourier interface}\label{exact-common-local-fourier-interface}

The accepted source \nolinkurl{Definitions/LatticeFrames/IsCombAt.lean}, \nolinkurl{LastLoc} (lines49--66), identifies the rank-zero lastBi tensor in compatible split models with

\[
 \Xi(X,(Z,\eta))
 =\Phi_1(Z)\int_{F^b}\Phi_0([X,y])\psi_k(-\eta y)\,dy.
\]

The local additive character is not the bare split character. In \nolinkurl{ExistsUnitarySupportedAtPlaceOfMulEqIdem.lean}, lines43--44, \nolinkurl{Place.IsChar} is the equality \(\psi_k(x)=\psi_E(\xi_0\iota_v(x))\). The same file defines the core frame, lines75--76, by a positive block from the actual complement isometry and a last column \(s\,i_L(\zeta)\). The required line condition in LastLoc is \(\tau z_0=-\xi_0s\zeta\).

For the actual LIFT30 negative line, put \(c=\bar\pi_T\), \(\zeta=c\), \(z_0=\pi_T\), and \(\tau=-\xi_0c/\pi_T\). Then \(s=1\), the actual first core frame is \(\alpha=g\operatorname{diag}(I_a,c)\), and the terminal full frame is \(\gamma=\operatorname{diag}(I_N,\pi_T)\). The terminal coordinate is consequently the literal source row \(\eta=t/\pi_T\) in K038(CA2--CA3), not an unscaled last row.

The source theorem \nolinkurl{SplitPlaceDoubleStepKernelLocalModelsStatement} is accepted at this interface. Its use does not assert that the arithmetic D7 host record exists on the mixed host.

\subsection{2. The already checked finite functions give the Fourier sum}\label{the-already-checked-finite-functions-give-the-fourier-sum}

At G, LIFT22(CG7--CG9) and K038(CA15--CA18) give

\[
 \Phi_0([X,y])=f_U(X)\mathbf1_{q^{-1}e_a+O^b}(y),
 \qquad \Phi_1(Z)=H_G(Z).
\]

With conductor-O \(\psi_k\) and additive volume \(\operatorname{vol}(O)=1\), the displayed integral is

\[
 f_U(X)H_G(Z)\mathbf1_{O^b}(\eta)
                    \psi_k(-q^{-1}\eta e_a).
\]

This is exactly the sole surviving G summand of the prescribed horizontal Y tensor. It has no projective index or additional Fourier volume. LIFT30(NL2,NL8--NL9) gives \(g\equiv1\) and \(c\equiv1\pmod G\), so both the raw finite functions and the finite character factors are unchanged by this particular G frame transport. This is consistent with \nolinkurl{SeeSawLastTwoStepsValueFixedLevelPlaces.lean}, \nolinkurl{last_value} line254 onward, which explicitly chooses \(c_G=1\).

At Q, use the complete LIFT23 functions Q5--Q7, including all supports, the determinant residue character and the full allowed quotient \(q^{-M}O^b/q^MO^b\). The first function is periodic by \(q^MO^b\). Its Fourier integral is therefore exactly

\[
 q^{-Mb}\mathbf1_{q^{-M}O^b}(\eta)
     \sum_{y\in q^{-M}O^b/q^MO^b}
                     F_X(y)\psi_k(-\eta y),
\]

because each contributing coset has volume \(q^{-Mb}\). The same orthogonality proves vanishing outside the displayed dual lattice. This is the source \nolinkurl{ftY} coefficient, not a normalization invented for the integral. LIFT23's proved cofactor/Gauss formula evaluates this same full sum; it need not be rederived.

\subsection{3. The Q unit phase must be placed in the correct normalization}\label{the-q-unit-phase-must-be-placed-in-the-correct-normalization}

Away from T, \(g\) is integral with determinant one. At Q the raw first vertical determinant function is invariant under right \(g^{-1}\), because its source rank is \(r=a+1\). For bare matrix coordinates, multiplication of its last source column by \(c^{-1}\) gives

\[
 \Phi_Q([X,c^{-1}y])=\chi_Q(c)^{-1}\Phi_Q([X,y]).
\]

All support conditions are preserved, since c is a Q-unit. Thus in a normalized Heisenberg model obtained by the passive substitution \(W\mapsto W\alpha\), the unchanged raw function has the displayed factor \(\chi_Q(c)^{-1}\). This is not automatically the scalar of the honest framed theta section: K038(CA4--CA5) retains the determinant factor of the split Weil action, and its source-frame factor here is \(\chi_Q(c)\). It can cancel the bare-coordinate factor when both are part of the same frame transport.

\textbf{Unresolved in this audit:} place these two factors in the literal \nolinkurl{VForm.vf}/\nolinkurl{w.sc 0} convention, using the same basic vacuum and the same fixed arithmetic row as K038. It is not valid to count the character twice, to set it to one merely because c is a unit, or to switch from passive coordinates to a Weil group action without its character. The factors are explicitly computable roots of unity and are units at11, but this observation alone does not identify the fixed source scalar. No completed G/Q common-tensor finding is promoted yet.

\subsection{4. Exact Q orbit weights and the source period scalar}\label{exact-q-orbit-weights-and-the-source-period-scalar}

\nolinkurl{Definitions/LatticeFrames/CosInj.lean}, QAt lines113--118 and QFam.Pins lines145--150, fixes the numerical factors:

\[
 w_1(\mathcal O)=I_1|\mathcal O|,\qquad w_2=I_2,\qquad
 c_{Q,v}=(I_1I_2)^{-1}.
\]

The actual vertical middle and terminal compacts are the parahorics of block sizes \((r,t)\) and \((r,2t)\), respectively, by LIFT31 section2 and DO038(19a--19b). Thus, with \(b=r+t\), \(N=r+2t\), and \(\kappa_j(q)=\prod_{i=1}^j(1-q^{-i})\),

\[
 I_1={b\brack r}_q=q^{rt}\frac{\kappa_b}{\kappa_r\kappa_t},
 \quad I_2={N\brack r}_q=q^{2rt}\frac{\kappa_N}{\kappa_r\kappa_{2t}},
\] \[
 c_{Q,v}=q^{-3rt}
       \frac{\kappa_r^2\kappa_t\kappa_{2t}}{\kappa_b\kappa_N}.
\]

For an orbit-invariant integrated summand, multiplication by c\_Q cancels the two indices and leaves the exact orbit cardinality, which converts representatives to the full Y sum. K038(CA9--CA11) is the already checked honest-character covariance needed for that invariance. These source Q period factors remain distinct from the canonical endpoint GGP factor1 of K056 and from the raw vertical Hermitian factors in LIFT38. A complete proof of the fixed global period law must retain the initial trace weights and its same input function as well.

\subsection{5. What remains to finish this approach}\label{what-remains-to-finish-this-approach}

Resolve section3's active/passive frame scalar, record the exact source \nolinkurl{VForm} realization, and then combine the G/Q full-Schwartz identities with DO051's fixed T operator and DO050's selected real recoupling. The intended common tensor is always \(C=\mathrm{op}_T(\mathrm{lastBi}(\Phi_0,\Phi_1))\), chosen before the auxiliary tuple varies. Its real compact averaging is allowed inside the middle period integral, not as a pointwise replacement by two highest vectors. Even a completed identity will leave the permitted nonzero terminal arithmetic row, its controlled detector, global integral preimage and actual DP44 budget open.

\clearpage

\section{Current research status --- after AM-LIFT-108}\label{current-research-status-after-am-lift-108}

Updated 2026-11-23T14:10:03.940004+00:00.

\subsection{Q1 --- OPEN}\label{q1-open}

The actual J14/J56 two-edge mixed-CM kernels and all six m2 analytic predicates are constructed (K031/K032/K038--044). K035--039/K041, R032, DO038 and DO050--055 give the full-Y arithmetic kernel, genuine finite projector, selected recoupling, terminal realization and original-vacuum last-two law. DO055 proves that a nonzero permitted moved all-T Mram bottom period gives a nonzero pairing on the prescribed arithmetic image without a prior-image or integral-vertical-inverse existence premise. LITERAL SOURCE ADMISSIBILITY: LIFT080 supplies a conditional generic mixed-line height transfer with its explicit unequal-degree cost and arbitrary-type geometric independence/length mechanisms; the effective line-degree, prescribed point, model, polynomial-integrality and source-room premises remain. LIFT079 CT10 also requires replacing the actual pointwise LamPtClauses/Good pure-factor interface for every prescribed mixed fibre with a nonelliptic simple factor; this is conditional on its frame antecedents, and dropping IsTen alone is insufficient. R003 excludes the unchanged elliptic IsTen promotion for powers of J; R023 excludes HostRaw on the fixed J14/J56 fields. LIFT54 excludes full SizesAt with m2 at every eligible FamP point and both fixed degrees at every Scale. LIFT55 excludes the entire unchanged L343 profile even after rebuilding degree, ranks and primes: full SizeG forces L\textgreater2\^{}300 and full sizes force m\textgreater2\^{}7778, f\textgreater2\^{}33598, k\textgreater2\^{}73197. A retained literal Host containing K14/K56 also requires m\textgreater2\^{}11388 by R023/o\_tie/SizeM. These necessary interface obstructions do not prove theta-kernel nonexistence or a DP44 sign. A candidate needs changed field/scale/length or a proved replacement mixed interface. LIFT51--53 retain the full-space Gram/image-rank incompatibility, necessary gap departures and code capacity bounds; they are not a full constructor. LONGER VERTICAL TOWER: LIFT56--58 prove forward cusp-valuedness, every prescribed finite G/Q and basic prior-image injection, real nonzero tests, the first positive-Witt convergent formula with factor2 and every-embedding Rallis nonvanishing. On the genuine smooth multiplicity spaces the unchanged maps satisfy U=D\emph{, ker U=0 and D(C\_m)=C\_0. LIFT59 now proves e0 C\_0!=0 on the actual standard bottom Hermitian model, conditional on its genuine finite/splitting data, actual discrete-series lists and the explicit BC3 numerical bounds implied by literal SizesAt/NCGaps. It constructs an arbitrary-m literal goodTuple with the required aggregate scalar character and an actual Poincare form satisfying \textbar P(1)-\textbar mu(E)\textbar\textbar\textless1/2 at the unchanged principal-G/maximal-away-G level. The prescribed vertical image is consequently nonzero. This is a bottom cusp-space value, not the prescribed Bessel/Mram functional. DO068 now constructs every genuine simultaneous finite Loc/common-vacuum package for the actual BF/FB mixed coordinates, including almost-all basic support and compact-finite pure-tensor witnesses. Its subgroup and bad finite factors are free data. The prescribed Bessel/Mram period, full source/arithmetic ties, arithmetic norm control and a rebuilt comparison baseline remain; the genuine positive-rank free-data implication is now supplied by DO069 in its stated conditional scope. CURRENT FINITE NONVANISHING GATE (AM-R-094 PP1-24; AM-R-097 NW1-9; AM-K-115 GM1-23; AM-K-116 RE1-28; AM-DO-121 EC1-24; AM-K-117 D29.1-27; AM-DO-122 CE1-29; AM-DO-123 BC1-27; AM-DO-124 CB1-27; AM-DO-125 CM1-28; AM-DO-126 IQ1-36; AM-DO-127 GC1-16; AM-K-123 HCX1-26): At the unchanged charged J56 all-top H133 cutoff, the actual energy is1+Re(sum over actual R27 plus actual R\_ge28)+epsilon, \textbar epsilon\textbar\textless2\^{}-230. R087 exhausts all supports through25 and R091/K109/DO115/K110/R094 exhaust support26. R097 followed by K115/K116/DO121 exhaustively proves that actual R27 is contained in the full-positive scalar-cross globally rank-two-unipotent all-II class with b=j0 and the sole remaining necessary alternative: total eligibility27,161\textless=k\textless=164, with2\^{}2111\textless x and the respective strict upper bounds x\textless2\^{}2130 for k161, x\textless2\^{}2140 for k162, x\textless2\^{}2150 for k163, and x\textless2\^{}2165 for k164. K117 removes the entire former total-eligibility28 branch with one278 charge. No equality with the containing class, existence, sign or absolute-smallness assertion is made. Full positive compression is supplied by DO120 before K114 is applied. All138 chronological all-top and independent31 uniform errors have exact full sum\textless2\^{}-230; current exact arrays are checks/k123\_joint\_remainder\_error.json. The complete all-top floor is27 and the independent uniform H8192 complete floor remains11. Earlier higher-support cuts remain: F,r1 through94; F,r2 through110; F,r3 through154; F ranks4..33 through168; IC8 ranks1/2 through28 and ranks3..31 through the full R075 table; IC8 ranks32/33 through168; both transverse-cross directions through70; rank-two nonunipotent plane-cross through32. DO120 removes all nonzero short compressions in the scalar-cross two-plane sector throughout27..31, with complete positive compression and global rank\textless=3 on its complement. R055 retains its scalar-cross rank-two-unipotent cuts(s13-48,u\textless=10),(s49-84,u\textless=9),(s85-120,u\textless=8),(s121-156,u\textless=7),(s157-168,u\textless=6). All other higher-support restrictions retain their accepted scope. Every paired CM incidence, class mass, phase, finite projector, packet and selected Hilbert normalization is unchanged. K115 removes the complete r27,k100..109 union with one945 charge and K116 the full k110 cell with one501 charge, and DO121 the complete k111..139 union with one993 charge; all smaller and higher-support historical restrictions remain in scope. The sufficient signed inequality PP3 is Re(sum over actual R\_ge27)\textgreater-1+2\^{}-230 and remains unproved. Actual admitted mixed-CM source, selected arithmetic-image content, prescribed global period, specific032 pairing/normalization/value and evaluatedDP44 remain open. Historical support26 windows and earlier structural possibilities below are superseded by the checked complete reductions above. DO122 removes the complete k140 cell with one359 charge using the original weighted central coset and mandatory non-strict-cutoff qualification. DO123 removes complete k141..142 with one1647 charge, including the entire retained-band tail once; its failed k143 envelope gives no further removal. DO124 removes the complete k143 cell with one707 charge, using the actual narrower band, full physical central integral, all167 Gamma512 moments and the entire new band tail once; its positive k144 envelope gives no further removal. DO125 removes the complete k144 cell with one232 charge, using the closed positive physical-central mixture, exact continuous moment enclosure, all lower-degree corrections and the entire tail once; its positive k145 sufficient envelope gives no further removal. DO126 removes the complete k145..159 union with one582 charge using the actual nonprincipal selfdual isotropic quotient, all-prime row containment, full dimension1008 transverse theta factor and independently retained central U\^{}-f; every4590 term and the entire tail are paid. Its positive k160 sufficient envelope removes no further cell. DO127 removes the five complete higher-denominator tails at X2120/2130/2140/2150/2165 with one415 charge, using exact rank-dependent Green/central-J0 cancellation, every finite pattern including degree167 cost2\^{}668, all1630 positive terms and the entire narrow-band tail once. Those higher-denominator cuts alone leave the low-denominator cells; the complete support floor remains27. K123 removes the complete k160 cell with one600 charge, retaining the actual selected high-coherence restriction, full continuum cap, same-kernel positive central mixture, all322 terms and both complete arithmetic tails. All four remaining k161..164 denominator windows and higher-support cuts stay in force; the complete support floor remains27. J14 retains R042: H8192 outside the degree12 plane, primary thresholds only for s\textless=17 and error\textless2\^{}-793; its sign is open. BUDGET: R037--039, DO057 and LIFT47 give integral input/detector spans and signed T cost. DO056 and LIFT48--50 retain the actual sufficient operator majorant, finite volumes/projectors, genus masses, hooks, coevaluation, CM periods and corrected full algebraic scalar. The allowed Weyl detector certifies D/1221\textless g\_Delta\textless=D/1220, while the intrinsic coevaluation/period norm factor exceeds that certificate by\textgreater220D in mean binary logarithm. This is neither an operator lower bound nor an actual restricted-image content upper bound. Sharp actual image content, other signed finite contents, full all-prime LinkG/LinkDIV/LinkMu and every changed-field/endpoint/height term in DP44 remain unresolved. No shorter admissible chain or universal strict negative DP44 difference is proved. LIFT60 constructs the exact arbitrary-rho geometric/fractional-ideal Hida flags and finite data on principal-extended-different standard hosts. It gives both-component K2=unipSym, the exact root-character signs and compact diagonal correction, and all finite RowLevelMatchA/B/C under the genuine grp\_levels/grp\_record ties. Its finite index is product\_w q\_w\^{}k product\_i(q\_w\^{}i-1), with original Haar factors retained. The genuine real relative model and full projected analytic row reconstruction are now supplied by LIFT061/062; concrete arithmetic row ties and the actual Bessel/Mram value remain separate. DO069 proves the positive-rank free-data period implication for the genuine R014/R015 realization, with R016/DO068 inputs, including both parity cases and the full uncontracted-to-pointwise comparison. It leaves the initial upper period, arithmetic Frames.Good and norm budget open. LIFT061 constructs the actual positive-rank real Wu/reproduction/decay data, with positive seed energy and its inverse square-root cost unevaluated; LIFT061-U proves local uniqueness up to phase. LIFT062 proves full projected-row reconstruction and cuspidality with actual source continuity, finite projector and Haar. LIFT062-T proves the translated EvalTieT under the literal genuine comparison fields, retaining tS, complex coefficient conjugation, coset volumes and Haar ratio. The concrete J/032 STiesR.rowTie, any required unshifted RowsJoinPinned transport, and actual Bessel/Mram period/value remain open. LIFT065 derives the exact range, nP-scaled norm, level containment and normalized row map from a supplied actual Read plus models\_HE and the genuine translated evaluation. The exact row scale is(c}nP)\^{}(-1/2). With the retained actual HCFinite input and source DataS.hhF, a freely chosen full-HH isometry can satisfy the row identity. LIFT066 now proves complete SIso/Pu/Gterm/column compatibility after the literal reS refill, from a genuine SPack, AutTie, retained LocTied, actual Haar tie, full dimension and normalized row map. Its SIso and row/column pins are identities on the installed genuine data. This is not an isomorphism with arbitrary old operators. STiesR still needs a compatible RecT and the exact forms, and every old comparison link or changed-family DP44 term remains to be supplied for the refilled record. The actual Read/Hida-column/compatible-record construction, prescribed period and032 value remain open. LIFT067 identifies the literal positive-rank T character coefficients(bar(pi),pi), the prescribed standardized D\emph{tau test and exact N-Haar modulus q\footnote{-rho(k+rho+1)}. LIFT068 now evaluates the actual fixed D}tau test conditionally on the stated generic tempered representations and both conductor-one hypotheses, retaining all phases and the Whittaker/Haar factors. LIFT069 now proves the larger generic tempered hypothesis at split T for every constituent of the chosen LIFT059 bottom subspace under its genuine full source package, via a full-degree finite cuspidal place and exact large-theta parameter cancellation. LIFT070 now supplies the smaller generic tempered conductor-one T line for every existing initial Mram constituent under the genuine full SizeData2/SizesAt package, using the actual common non-top label and honest eta parity correction. Thus the prescribed physical D\emph{tau test is locally nonzero against the chosen LIFT059/LIFT069 bottom constituents. LIFT069--070 alone did not extend the chosen-bottom assertion to all e0; LIFT076 below now proves that extension. A fully sized mixed-CM source remains unconstructed. LIFT071 now calibrates the canonical positive-rank normalized physical T factor alpha\^{}natural exactly as q/{[}(q-1)(q\^{}k-1){]} under the specified probability/product Haar, with inverse square-root cost\textless q\^{}(k/2); the literal source normalized factor is c\_H}c\_N times this value, and the unnormalized alpha is Lscript times alpha\^{}natural. Its scalar-one comparison is proved by compact calibration and canonical regularization. LIFT072 now constructs the actual positive-rank split G test inside the same admitted generic finite representation, including mixed types, for every prescribed smaller vector: alpha\_G=1/d\_varsigma, or c\_H\emph{c\_N/d\_varsigma in literal source Haar. Its normalized factor retains Lscript, and its product is\textgreater Q\^{}(-f}k\emph{(k+1)/2). LIFT073 now gives every actual away-G/T spherical factor alpha\_v=c\_H,v}c\_N,v\emph{Lscript\_v and normalized factor c\_H,v}c\_N,v, including all actual dyadic places by a direct split calculation. Thus the prescribed finite local factors are determined under the retained genuine representation/source hypotheses. The global Bessel/Mram period remains open; no source Haar parameter is set to1 and no integral bound for the Hilbert vector is inferred. K089 now proves unit11-content of each actual full-Y auxiliary section before the selected differential and class contractions, with the exact vacuum scalar and finite phases retained. Its finite all-Hasse-jet unit minor does not prove survival of the prescribed selected ordinary derivative/class matrix. The remaining sharp-content question is that exact selected map on these primitive inputs, with the established lower gain removed; no freely chosen auxiliary jet or added T/Y orbit is substituted. K090 now gives the exact integral Mram tensor-augmentation basis on the unchanged free strict/projective T cover, with source class norm2\^{}(f/2) times the coefficient norm and unchanged frame/Haar costs. It identifies the exact selected matrix as the alternating2\^{}f actual columns and proves saturated auxiliary-section injection of that source before subD/class/projector contractions. Actual Sad preserves the full local source lattice at11. The selected matrix/content/nonzero gate remains MA10--11, with every fixed derivative, class weight and genuine projector retained. K091 now supplies a literal integral retraction of both symbolic selected Weyl substitutions and their complete two-edge composition, including the fixed-triple tensor-augmentation auxiliary-section image. Thus universal selected-polynomial substitution adds no Smith factor before the prescribed ordinary-derivative functional. Its abstract integral-lattice reduction is distinguished from nonsaturated ambient polynomial reduction. K092 now supplies a common-unit full-stabilizer cutoff with a primitive symbolic selected theta family and exact contracted full-middle-mass division. No uncontracted section division or selected-derivative survival is asserted. The evaluated selected matrix MA10--11 still requires the actual derivative/gauge, weighted geometric class sums, terminal brC/projectors and their exact content. LIFT074 proves the exact conditional global detector transfer on a specified distinguished embedded pair: \textless F,Pu(T)\textgreater=c\_Per\emph{conjugate(b),c\_Per=c1}mu\_f(K2)\textgreater0,with all prescribed finite tests and the complete unit-Wu contraction. The pairwise global distinction premise and actual arithmetic source/Read/032 matching remain open; b is unevaluated, so no numerical DP44 cost follows. K093 identifies the exact scalar-decoration Fourier/Heisenberg matrix RT13, with class-dependent actual rank-one rational translates and all original phases/masses. Its selected ordinary evaluation necessarily retains at least D\_tau-nu\_tau,1\textgreater0 first-edge derivatives; finite Fourier units do not determine the content of those derivatives or their class sum. The exact evaluated matrix remains open. LIFT075 now excludes the direct compact-source theta cuspidal construction for every permitted positive rank gap: the added GL1 string forces one repeated linear Green factor when d\textgreater=2, and the literal goodTuple degree-one witness kills d=1. This does not kill the coupled map or the transported-leading-r intermediate projector and does not settle the prescribed global Bessel period. K094 now gives an integral full-covector Fourier basis of the actual localized Mram domain and cancels the complete Fourier sum exactly into the signed2\^{}f corner columns FC13. Their span is the actual image/content ideal FC14, retaining full middle mass, a shared P, central phases and positive residual ordinary derivatives. The evaluated corner and middle-class sum remains open; a further source restriction retains its Fourier inverse image. LIFT076 extends the simple generic cohomological tempered-base-change theorem to every existing bottom constituent of the genuine full e0, using squarefree full inertial multiplicities, actual unitarity and the original all-degree divisor witnesses. The actual T spherical representation therefore satisfies the larger-representation tempered/generic premise without selecting LIFT059's special tuple. This remains a classification theorem; source existence, global distinction, arithmetic image/value and DP44 costs are open. DO094 adds the exact finite transverse-scalar implications (rank\textless=3 with both C nonzero, rank\textless=2 with at least one, rank\textless=1 without a C condition), and the original integral delta\_z witness. Its sharp excluded finite cells remain distinct; no numerical charge follows. K095 removes the whole middle T-covector sum and fine-class mass by the exact composite incidence and legitimate contracted coarse-stabilizer cut. The actual initial Fourier column is MF12: a sum over coarse middle classes with the same P and the original ordinary selected evaluation of the quadratic phase psi(u2ZXmu). Auxiliary primitivity precedes that evaluation; selected content/nonzero and coarse-class cancellation remain open. R071 now proves the remaining compact transverse scalar is exactly1, using integral trace/norm and the strict exact bound36U/M²\textless1. Thus the actual support23/24 scalar-cross residual fixes a positive166-dimensional E-space and has rank(h-I)\textless=3; its remaining invariant complement has signature(2,1). This supersedes the weaker positive165/rank\textless=4 structural gate from DO093. No integral moving basis, global nilpotence, full support removal or new numerical charge is inferred. LIFT077 now calibrates the exact prescribed coupling through BPC2025 Theorem1.3.6.1: GC17 retains aS²/(aH\emph{aG}a0\emph{d\_sigma}2\^{}(t\_pi+1)), embedded nT/nF, the partial central Rankin--Selberg ratio, raw G dimensions and once-normalized T factors. It applies also when the period is zero. For the specified existing genuine pair, a nonzero central Rankin--Selberg value is equivalent to nonzero b; neither its nonvanishing nor quantitative size is proved. GC19 retains the literal c\_Per,c0,nT in the source Gram lower bound. Source construction, selected arithmetic content, Read/STiesR/specific032 and DP44 remain open. R072 removes the R071 support23 scalar-cross terms with6\textless=k\textless=164, where k counts nonzero original S scalar coordinates, with one3600 charge. Its fixed packet-plane projection retains the U denominator, exact paired T\^{}(4s) target index and rank-two source fibre; the first row uses the actual joint norm and spherical-cap/Gaussian packing. The remaining support23 row has k\textless=5. No support24 or complete-floor conclusion is inferred. In the actual R071 signature(2,1) residual, BC is skew-adjoint rank one exactly at global rank3; BC=0 is exactly global rank2 with real reciprocal eigenvalue lambda. Every original S entry of BC is an exact G-unit three-minor, has conjugates\textless2\^{}9 and only simple paired typeIII poles. The pointwise Green envelope is2 Q\^{}(-501\emph{167)2\^{}(9}167f)T\^{}(167t), with no C-pattern/witness factor. In rank2, doubled G divisibility of (lambda-1)\^{}2/lambda forces global unipotence for t\textless=22. The existing DO082 removal therefore leaves rank-two support23 only nonunipotent with all23 active places typeIII; at24 its nonunipotent branch needs t\textgreater=23. No new charge or complete-floor change. Actual central value, arithmetic/source and DP44 gates remain open. The actual global rank3 R071 terms at s23,k\textless=63 and s24,k\textless=51 have combined absolute contribution\textless2\^{}-21500. The bare support-ideal count retains all original sparse cross coordinates, three transverse values, pivot, fixed packet-plane U denominator, paired T\^{}4 column index, positive rank-two source fibre and complete pair weights. The actual DO095 witness gives the uniform no-pattern rank-three Green bound. All114 cells are checked. Combined with R072 this removes every rank3 support23 term; DO095 then leaves the support23 scalar-cross remainder rank2 nonunipotent, k\textless=5, t23. At support24 remaining rank3 has k\textgreater=52. Complete floors23/11 and signed/global/source/arithmetic/DP44 gates remain open. K096 proves exact emptiness of the remaining nonunipotent rank-two support23 scalar-cross subcase, without a sparsity condition or a new charge: actual-pole parity leaves eleven possible paired cross poles, incompatible with the original row norm. Together with DO082 and R073 this exhausts the R071 support23 scalar-cross residual. R074 now supplies the complete structural support23 assembly, without a new charge. R074 supersedes all earlier historical complete-floor23 descriptions: the current complete all-top floor is24 and the exact current sufficient inequality is J24.2, with unchanged77/31 errors. The entire globally nonunipotent rank-two R071 support24 sector has absolute contribution\textless2\^{}-2962487. The nonzero reciprocal trace has23 or24 actual poles; principal-norm parity and one exactly normalized inversion leave at most twelve paired cross-pole places, including the possible active eigenvalue-unit place. Primitive original extraction gives normT\^{}-24 without a U denominator. The full original R059 count retains both five-coordinate columns, source fibre, all G patterns, padded subset choices and pair weights; the complete exponent is -60671740183/20480. Exactly one2962487 charge is added, leaving78/31 entries with both full sums below2\^{}-230. Complete floors remain24/11. Globally unipotent rank2, rank3, signed/global/source/arithmetic and DP44 gates remain open. In the original IC8 class at compressed ranks3..31, the full symmetric compression has rank r at every actual G factor. Finite unipotence therefore forces full rank\textgreater=r+1; at equality the unique nontrivial Jordan block fixes every untouched original column at both CM partners. The resulting paired ideal and baseline rank r+2 retain every original DO085 coordinate, volume, target, source fibre and pair factor. All145 stated fresh cells contribute\textless2\^{}-8866, with exact maximum -36351867/4096. The complete rank3 ceiling becomes28, and with the earlier F/IC8 cuts every remaining term at supports24..28 has A=I. One8866 charge gives79/31 entries with both full sums\textless2\^{}-230. Complete floors24/11 remain; signed/global/source/selected-image/032/DP44 gates are open. The entire globally unipotent rank-two R071 support24 sector with at least one type-I or type-III place has absolute contribution\textless2\^{}-71310. The exact block normal form gives N²Z=0 and alpha=Uq\textless=H+H\^{}-1-2, so all original S and fixed transverse inputs have only type-II poles and the full covector radius is sqrt(R/U). At most24 padded23-place ideals cover every non-all-II profile, with all164 coordinate sectors, common G patterns, genuine selected trace, true projected-column U denominator, paired T\^{}4 index and source fibre retained. The complete maximum is -292137655/4096; one71310 charge gives80/31 entries, both full sums\textless2\^{}-230. Complete floors24/11 remain. The all-II unipotent case, rank3 and fixed-block residuals, signed/global/source/arithmetic/032/DP44 gates remain open. The entire global-rank-three R071 support24 cross sector has absolute contribution\textless2\^{}-5710, with no sparsity premise. The full positive167 covector is counted in its actual Gram metric, with exact determinant4U and arbitrary-ideal covering radius; its last scalar is counted by the exact relative norm fibre, including nonprincipal and non-conjugation-stable ideals. Two pointwise Green bounds on the same original G patterns meet at z=2\^{}-100; full low/high exponents are -23390187/4096 and -361072601/4096. All original trace, paired targets, projected U denominator, source fibre, support, pivots and incidence weights are retained. One5710 append gives81/31 complete errors below2\^{}-230 and unchanged floors24/11. Rank-two all-II and fixed-block residuals, signed/global/source/selected-image/032/DP44 gates remain open. The entire all-II globally rank-two-unipotent R071 scalar-cross support24 sector has absolute contribution\textless2\^{}-2879. The actual cross map and original central entry w reconstruct the full matrix. Its affine difference lattice iJ0 I\_F² retains the rank-zero G congruence, real codifferent and full discriminant, giving the complete central count2\^{}479 Q\^{}167 z\^{}-2f. The true transverse Gram volume, exact split-relative norm fibre, every original all-II paired weight and the actual selected trace are retained. All164 coordinate sectors are exhausted by64/99/1 complete monotone ranges with strict maxima -2010899/512,-41042253/4096,-11800269/4096. One2879 append gives82/31 errors below2\^{}-230 and unchanged floors24/11. Together with DO097 and R076/DO096 this removes every R071 scalar-cross support24 residual. The fixed-S rank-two sector, signed/global/source/selected-image/032/DP44 gates remain open. The entire original fixed-S globally rank-two-unipotent support24 sector has absolute contribution\textless2\^{}-6845. The complete R058 high finite-rank envelope extends with exponent -119372763/4096. At low finite rank the original square pairings force the null image into the packet plane; a type-I place is impossible by the unchanged pole-budget46 norm contradiction. All remaining II/III profiles are counted through the true transverse3 dual Gram metric, exact pivot relative norm fibre and original affine central-entry lattice, retaining both CM partners, full discriminant and actual incidence weights. The complete low exponent is -168257413/24576; the high/low union contributes one6845 charge. Every old charge, including DO098, is retained, giving83/31 errors below2\^{}-230 and unchanged complete floors24/11 until a separate assembly. Signed/global/source/selected-image/032/DP44 gates remain open. Every actual all-top support24 term is exhausted by the complete original F/IC8 partition at compressed ranks1..33, every fixed-S nonunipotent/rank-one/rank-two class, and every scalar-cross global-rank/type/coordinate alternative. R075 covers all nonidentity positive compressions; DO090/DO092/R077 exhaust fixed S; R070/DO093/R071 reduce cross terms to the precise DO095 rank2/3 alternatives, exhausted by R076,DO096,DO097,DO098. The complete all-top support floor is25, uniform floor11. All83/31 chronological errors are unchanged and both full sums remain below2\^{}-230. No new charge. The actual remaining signed gate is J25.2; prescribed global/source/selected-image/specific032/DP44 and RH/global oscillator gates remain open. Both complete C-transverse directions, with every incidence profile and no global-unipotence restriction, are removed at supports25..70. At support25 the plane-C/nonzero-original-short-compression sector is also removed. The full joint row, both original columns and complete source fibre give transverse maximum -49537435/4096; the complete plane row, exact relative norm fibre and actual selected trace give -435793141/4096. Their entire union costs one12079 charge. The remaining scalar-cross support25 matrices have both cross blocks in the original packet plane, all four short compressions zero, positive166-dimensional fixed space and global rank at most3 by the support-free norm-one lemma. At26..70 only the two plane conditions are concluded. Actual84/31 ordered errors remain below2\^{}-230 with complete floors25/11. No signed/global/source/selected-image/specific032/DP44 or RH conclusion. The entire original fixed-S globally rank-two-unipotent sector at support25 is removed for every actual incidence profile. Also removed at26 are D\_G\textgreater=33 and the D\_G\textless=32 covector sector r(y1)=r(y2)=0. The original central affine lattice has poles I1/II2/III1; its T\^{}(u-s) factor cancels actual pair weights to T\^{}(-t)\textless=1. The true three-covector short pivots and one-coordinate Schur-complement last pivot retain all source/norm/trace/metric factors. Complete low exponent at25 is -27553337/24576; high endpoints are -77307919/4096 and -35243075/4096; the last-only26 endpoint is -78939175/8192. One1120 charge covers the complete union, giving85/31 errors below2\^{}-230 with unchanged complete floors25/11. Short-pivot low-rank26, fixed-S rank-one25, remaining scalar-cross25, signed/global/source/selected-image/specific032/DP44 and RH gates remain open. The complete global-rank-three scalar-cross support25 sector with the deterministic original common-C ideal satisfying2\^{}32928 N\_EJ\textgreater=T\^{}50 is removed with one29550 charge. The whole original167-coordinate dual-metric ball, relative norm fibre, selected trace, paired projected column, source fibre, all patterns and complete actual incidence weights remain. The full monotone high-ideal endpoint is -121040253/4096. Every remaining rank-three support25 term satisfies2\^{}32928 N\_EJ\textless T\^{}50. The weaker low-ideal extension is a failed sufficient bound4090765885/12288\textgreater0; no actual contribution is inferred. A complete local typeII shear model shows why C-integrality cannot be substituted. Actual86/31 errors remain below2\^{}-230 with complete floors25/11. No whole support25, signed/global/source/selected-image/specific032/DP44 or RH conclusion. The entire remaining globally nonunipotent rank-two scalar-cross interval25..32 is removed after the retained DO099 transverse cut. The support-independent DO076 reciprocal-line theorem gives lambda in F and inverse-C=-lambda\^{}-1 C without requiring full-Z at26+. At least23 actual lambda poles and one global inversion give at least12 positive valuations; all active lambda-unit places are retained in the complete paired C-pole cap s-12. Every original column, source fibre, pattern, pivot, support and labelled-pair weight remains. The complete worst exponent is -2610113463/20480 at32, and eight supports cost one127443 charge. Actual87/31 ordered errors remain below2\^{}-230 with complete floors25/11. Support33 is a failed sufficient bound only. Global-unipotent rank2, rank3, signed/global/source/selected-image/specific032/DP44 and RH gates remain separately open. Every original fixed-S globally rank-one-unipotent support25 term is removed with one537 charge. The primitive q coordinate retains the original affine difference ideal, so both q/r true eta-metric element counts replace the older relative-norm fibre. Independently checked LIFT078 averages every conjugate-content ideal over eta via actual subideal packing and the unconditional actual-field bound zeta\_F(1.001)\textless3003·5\^{}167. Its full constant2\^{}1065 and exponent21021/125 are retained before all Green patterns, selected traces, supports, marked subsets and original paired type/label sums. All26 double-pole cells pass; worst margin538, next1530, complete witness exponent -25183329/64. Actual88/31 errors remain below2\^{}-230 with complete floors25/11. This is a full rank-one sector result, not whole support25 or signed/global/source/selected-image/specific032/DP44/RH nonvanishing. Every globally unipotent rank-two scalar-cross support25 term with at least one type-I or type-III place is removed with one881 charge. The actual full-positive normal form makes all original cross coordinates integral outside typeII. Its original central coordinate has affine difference lattice iJ0 I\_Sigma,F I\_D,F, retaining the shift, full real discriminant and750 transference coefficient. The complete u24 sector costs less than2\^{}-882, including50 type-position choices and both extra T factors. The padded u\textless=23 sector includes300 sets, all3\^{}25 type assignments and T\^{}25 pair cost, with explicitly charged cap30 removal and full sum below2\^{}-174292. Actual89/31 ordered errors remain below2\^{}-230 with complete floors25/11. The all-II support25 sector remains open; no signed/global/source/selected-image/specific032/DP44 or RH conclusion. Every remaining rank-three scalar-cross support25 term now has at least123 nonzero original S coordinates and2\^{}33096 N\_EJ\textless T\^{}50 for the same deterministic common-C ideal. The original(k+3)-coordinate covector ball keeps the true Gramdet4U, arbitrary-ideal cover, exact beta norm fibre, actual selected trace, full paired plane-column U/T index and source fibre. At z\emph{=2\^{}(-197/2), all122 low-ideal cells have maximum -7111295/4096; the complete high-ideal endpoint is -5721469/4096. Both complete unions cost one1395 charge, retaining the old overlapping DO10029550. Actual90/31 errors remain below2\^{}-230, complete floors25/11. Dense smaller-ideal rank3, all-II unipotent rank2 and signed/global/source/selected-image/specific032/DP44/RH gates remain open. R082 TP1--13 restrict every remaining rank-three support25 term, with its actual t=11..25 type-III places, to k\textgreater K\_t and2\^{}(42p\_t)N\_EJ\textless T\^{}50 for the same common-C ideal. In increasing t order, (p\_t,K\_t) are(1120,149),(1096,148),(1073,147),(1049,147),(1026,146),(1002,145),(979,144),(955,142),(932,141),(908,139),(885,137),(861,135),(838,132),(814,129),(791,125). The complete general type/pair weights and actual-t minor poles give maxima-9691203/4096 over2106 low cells and-11857677/4096 over15 high cells, one2353 charge. Actual91/31 errors stay below2\^{}-230; complete floors25/11. The all-II unipotent rank-two branch and surviving dense smaller-ideal rank-three profiles remain; signed/global/source/selected-image/specific032/DP44/RH gates remain open. DO103 removes the complete all-II globally rank-two-unipotent scalar-cross support25 subset with N\_E denominator(eta)\textless=2\^{}120. Both actual CM pole lines force the same selected delta=eta-Q(-7+2sqrt(-7)) divisibility. The globally integral denominator-cleared delta ideal has norm at least7\^{}84 T\^{}p and below{[}4(U+14){]}\^{}168 N\_E D\_eta, so the exact host inequality gives p\textless=24. The shared cross/double-pole containing set, original affine central counts and all164 CF endpoints give25T}2\^{}-2879\textless2\^{}-1878. Actual92/31 ordered errors remain below2\^{}-230, complete floors25/11. Every remaining all-II support25 term has N\_E D\_eta\textgreater2\^{}120. The large-denominator branch, rank-three profiles, signed/global/source/selected-image/specific032/DP44/RH gates remain open. LIFT079 CT1--10 proves that literal CMObj and, independently, both actual LamObjGood copies force the complex target to be isogenous to a power of one CM elliptic curve, using pure complete quadratic eigenbases and the actual surjective/quasi-inverse sum. The actual LamPtClauses fibre identity transfers this constraint to the prescribed point; a point with a nonelliptic simple factor has no such witness. Thus dropping IsTen or bypassing CMObj while retaining Good, or substituting an unrelated elliptic Good object, does not supply the mixed-J point hypothesis. This is conditional on the stated point/frame premises and is not a contradiction to the accepted original source theorem. R003/DP25 simplicity is carried, not counted again. No mixed replacement constructor, finite charge, period/image/value, DP44 or RH conclusion follows. AM-K097 removes the complete actual global-rank-three scalar-cross support25 subset with N\_E D\_eta\textless=2\^{}120, using the cross-only pole lemma, at most24 eligible places, all25 padding choices and all15 actual type profiles. The complete normalized bound is below2\^{}-11658. No rank-three full-double-pole or unipotent central-fibre assertion is imported. One11658 charge appends to the actual DO10392/31 prefix, giving93/31; both full exact sums remain below2\^{}-230 and complete floors remain25/11. Remaining rank-three terms have N\_E D\_eta\textgreater2\^{}120 and retain R082 cuts. The signed kernel, actual selected arithmetic image, prescribed global period, mixed source construction, specific032/Read, changed-family DP44 and RH remain open. AM-R-083 removes the entire actual global-rank-three scalar-cross support25 tail N\_E D\_eta\textgreater=2\^{}4864 with k\textgreater=126, the latter already forced on the current remainder by R082. Its full167-dimensional arbitrary-ideal covector count, exact per-denominator eta enumeration and complete ideal-zeta tail retain all original type weights, G patterns, selected traces, projected U/T column and source fibre. The complete exponent is -22330263/4096\textless-5451, charged conservatively once as5400. The actual K09793/31 prefix becomes94/31; full sums remain below2\^{}-230 and floors25/11 are unchanged. Together with K097, remaining rank-three support25 terms satisfy2\^{}120\textless N\_E D\_eta\textless2\^{}4864 and all R082 profile cuts. The all-II unipotent rank-two, signed kernel, selected arithmetic image, mixed source, prescribed global period, specific032/Read, DP44 and RH gates remain separate and open. AM-DO-104 removes the complete actual all-II globally rank-two-unipotent scalar-cross support25 tail N\_E D\_eta\textgreater=2\^{}1820. Every original coordinate lies in LD\_eta; the full original4-by4 block and both mixed blocks give the exact affine central difference ideal iJ0 B\_eta I\_Sigma,F\^{}2 with D\_eta bar(D\_eta)=B\_eta O\_E. The full real750 cover, all167 covector values, actual selected trace, all164 row cells, original weights and the entire ideal-zeta tail are retained. The complete maximum -10346711/4096\textless-2526 yields one2518 charge. Actual R08394/31 becomes95/31, complete sums remain below2\^{}-230 and floors25/11 remain. Together with DO103 the remaining all-II rank-two branch has2\^{}120\textless N\_E D\_eta\textless2\^{}1820. No whole support25, signed-kernel, prescribed selected arithmetic image/period, mixed source, specific032/Read, DP44 or RH conclusion. LIFT080 MT1--17 proves a generic Step5 transfer for actual mixed-type comparison lines with individually bounded normalized degrees d\_j\textgreater=d\_-. Rescaling only their archimedean norms by exp(d\_j-e) makes all auxiliary degrees e, keeps every finite ideal and local polynomial quantifier unchanged, and replaces La=L by L exp((e-d\_-)\emph{+). The exact cost is C(e)=(1/2+\textbar e\textbar+log(L)/2)(D+gn)+D log(rL)+(3D+gn)(e-d}-)\emph{+/2. Its optimum within this comparison family is C0+(-d}-)\emph{+ min(D+gn,(3D+gn)/2). Arbitrary-type isogeny derivatives and the literal arbitrary-set form\_letter\_le supply independence and norm mechanisms; explicit mixed products J\^{}(14h p) with complementary-action factors supply geometric inclusions/projections with N=L=1 and permutation U. The actual prescribed point/model/section, all-extension polynomial tests, effective eigenline degrees, coefficient-field/frame ties and fully revised source/DP44 room inequality remain required. The old LamPtClauses/LamObjGood pure-type boundary remains; no main-family constructor, finite charge, period/image/value, specific032/Read or RH conclusion follows. AM-R-084 exhausts every remaining global-rank-three scalar-cross support25 term. The complete K097 count applies to \textless=24 eligible places without a denominator premise; the exact DO093 inversion carries this count and R083 to the inverse null line with the same original S-supportk\textgreater=126, actual weights, masses, phases and cutoff. Outside the four charged sets both slopes are fully eligible and each denominator norm is\textless2\^{}4864. Their denominator-cleared difference has the full paired T\^{}50 divisor but norm\textless(4U)\^{}168\emph{2\^{}9728\textless T\^{}50, so the two lines coincide. DO095 then forces global rank2, a contradiction. The entire four-sector absolute union is\textless2\^{}-5398, with both inverse costs included. Actual DO10495/31 becomes96/31; every old overlapping charge remains, full sums are\textless2\^{}-230, and complete floors25/11 remain unchanged. At support25 only the separately retained all-II globally rank-two-unipotent scalar-cross branch remains, with2\^{}120\textless N\_E D\_eta\textless2\^{}1820. Higher supports and signed-kernel, selected arithmetic-image, prescribed global period, actual mixed source, specific032/Read, DP44 and RH gates remain open. LATEST EFFECTIVE DEGREE REFINEMENT: LIFT081 ED1--23 supplies effective intrinsic line-degree lower bounds for every actual comparison component eigenline on J14/J56, with the full integral CM action and the literal source metric n\^{}2=pi H. The good model is obtained over K0(alpha,beta), alpha\^{}7=-4 and beta\^{}2=(zeta7-1)\^{}7, of degree at most14 over K0. The actual integral vector 2q\_h p\_tau e\_i has norm below sqrt(2)q\_h at every counted embedding, with q1=406 and q4=26390. Therefore ndeg\textgreater-log(sqrt(2)q\_h), and the weaker common bounds are -10log(2) for J14 and -16log(2) for J56. These normalized degrees are invariant under further finite base change of the actual good model. Substitution into LIFT080 makes its sufficient conditional Step5 cost fully explicit: Copt=1/2(1+log L)(D+V)+D log(rL)+b\_h min(D+V,(3D+V)/2), V=gn and b\_h=10log(2) or16log(2). This resolves the effective comparison-line-degree premise only. The prescribed mixed point, model/level/coefficient-field ties, nonzero frame polynomial, all residue-extension Step5Loc conditions, source-room data, selected arithmetic image/content, global period and DP44 evaluation remain required. No claim is made about an arbitrary line in a larger multiplicity space, equal degrees, a primitive coordinate-ideal equality, an old LamPtClauses witness, a finite charge, or RH. AM-DO-105 removes the entire actual all-II globally rank-two-unipotent scalar-cross support25 subset with at most24 eligible places, without any denominator premise, together with the fully eligible middle window2\^{}120\textless N\_E D\_eta\textless2\^{}1820 having1\textless=k\textless=21 nonzero original S coordinates. For each fixed support there is at most one fully eligible eta in this window and at most one actual diagonal pencil label per active place, including central-only poles. The original weighted label sum is at most(T-1)\^{}-25. The joint S/transverse beta integral, moving central cover and simultaneous x/z endpoint retain all full metric/phase/class/Haar factors. The worst claimed cell -1852891/4096\textless-452 and full fresh union give one451 charge. Actual R08496/31 becomes97/31, both complete sums remain below2\^{}-230 and complete floors25/11 remain. After R084, DO103/104 and this theorem the only support25 branch left is fully eligible all-II rank-two unipotent, with2\^{}120\textless N\_E D\_eta\textless2\^{}1820 and22\textless=k\textless=164, unique eta per support and at most one label per active place. The dense range, all higher-support cuts, signed inequality, actual mixed source, selected arithmetic image/content, prescribed global period, specific032/Read, DP44 and RH gates remain. LATEST ARBITRARY-TYPE POLYNOMIAL REFINEMENT: LIFT082 FT1--29 transfers one actual nonzero frame polynomial to every Step5Loc prime and complete-DVR extension, and every Step5Arch embedding, from explicit arbitrary-type geometric inputs. Generic FrPt polynomial existence and same-point nonzero evaluation produce one G before any completion. Integral eigenmodule bases commute with arbitrary compatible DVR extension by flatness; no integral eigenspace splitting is assumed. The exact relation Ub\^{}t Y=X, intrinsic theta generator unit, determinant inclusion and physical row coefficient-clearing indices give all local clauses. The original pi-normalized Hermitian metric, actual fibre/section maps and goodOf conjugation give all archimedean clauses, retaining lieFac and the dual theta norm. There is no new transfer cost beyond LIFT080/081. The actual prescribed point with Dw.mult, good model/level/Key/section/coefficient-field ties, valSet arithmetic-lattice identification, same-point nonzero F value, Host/Sizes, Read/STiesR, selected arithmetic image/content, global period, specific032 and DP44 remain required. No old CMPt.Good/LamObjGood/LamPtClauses instance, finite charge, support-floor change or RH assertion follows. LATEST ACTUAL J-ISOTYPIC SIGNATURE RESTRICTION: LIFT083 JS1--25 proves the actual pinned signature obstruction for every J-isotypic complex fibre, independently of CMObj/IsTen/LamObjGood. For any quadratic Host action into End0(J\^{}N)=M\_N(Q(zeta7)), d=7 gives Lie multiplicities(2a+b,a+2b), a+b=N, and every other Host quadratic field gives balanced(3N/2,3N/2). Actual DetCond requires(f\_E p,f\_E q). Hence d=7 needs p\textless=2q and q\textless=2p, while d!=7 needs p=q. The literal pinned profiles p\_j=k+j−1,q\_j=m−j are unbalanced at every stage and terminal q\_m=0; under full SizesAt, k\textgreater m\^{}5 f\_R\textgreater=2m forces p\_j\textgreater2q\_j at every stage. Thus no solely J-isotypic point can satisfy any pinned stage under those full premises, even with enlarged field/changed primes and the old good-point wrappers bypassed. An actual enlargement A\textasciitilde J\^{}N×B must have N\textless=f\_E min(p,q) and dim B\textgreater=f\_E max(p−2q,q−2p,0); the terminal fibre contains no J factor. These are necessary geometric dimensions, not height or DP44 costs. Other CM isogeny classes, their actual full point/action/model/level/section and same-point nonzero value remain open, together with selected arithmetic content, prescribed global period, Read/STiesR/specific032 and DP44. No finite charge, support-floor change or RH conclusion follows. AM-K099 GE1--13 restricts the actual remaining fully eligible support25 branch to total eligible count q=25 or26. The norm bounds give x\textgreater2\^{}120 for q25, x\textgreater2\^{}1116 for q26, and x\textgreater2\^{}2111 for q27, excluding q27 under x\textless2\^{}1820. Over K0=Q(sqrt(-7)), the relative Galois group of order168 acts regularly on the selected T primes. At q25 the scalar has trivial relative stabilizer and generates E over K0; at q26 the relative stabilizer has order at most2. The absolute degrees are at least168 and84 respectively. Distinct actual norm-U scalars in the bounded denominator window have eligible-set overlap at most14; set and scalar stabilizers agree. A relative stabilizer of order2 forces the denominator ideal to extend from the fixed field and x to be a square, even at ramified primes. This gives necessary structural restrictions only, no numerical sector removal, new error charge, source construction, signed positivity, period, DP44 or RH claim. The existing97/31 error arrays and complete floors25/11 remain unchanged. AM-K100 QE1--16 removes the complete actual fully eligible all-II globally rank-two-unipotent scalar-cross support25 subset with2\^{}1116\textless=x=N\_E D\_eta\textless2\^{}1820 and22\textless=k\textless=48, including either total eligible count25 or26. The exact1116 endpoint, cap1117, complete joint row/transverse beta count, central volume\textless2\^{}679 and unique-eta/label weight give worst k48 ceiling-1730011/4096\textless-422; the26 smaller cells and this cell form one complete2\^{}-421 union. All k\textless=21 contributions remain in DO105451 and are not added again here. Append421 to the actual DO10597/31 ordered totals, giving98/31; both entire sums remain below2\^{}-230 with complete floors25/11 unchanged. In particular all q26 cells up through k48 are removed. The surviving support25 possibilities have2\^{}120\textless x\textless2\^{}1116 with22\textless=k\textless=164, or2\^{}1116\textless=x\textless2\^{}1820 with49\textless=k\textless=164, in addition to K099 structural restrictions. No claim for these surviving cells, no actual source, selected arithmetic image, prescribed period, signed positivity, specific032, DP44 or RH conclusion follows. AM-DO-106 CFI1--24 proves the complete actual central fibre Gamma=B\_eta I\_Sigma,F squared, covolume sqrt(D\_F) x T\^{}-50, sufficient period Q T squared B\_eta of index Q\^{}168 T\^{}386, and the full-real codifferent dual. The original active type-II factor has Fourier l1 norm2/T for integral cross rows and1/(T-1) for a fixed cross pole. The actual mixed G projectors and every inactive-T factor are retained by the finite positive-projection trace lemma. The center-independent whole-C Green budget gives total finite Fourier l1 at most2 Q\^{}(-335}167)(N\_E J\_C)\^{}(335/2)(2/T)\^{}j(T-1)\^{}(-(25-j)), with no finite-index loss. The complete orbit lies in the original finite support; regularized Poisson retains its exact signs and covolume. Extending a disjoint collection past H133 uses at most the single complete2\^{}-2000 tail once, subject to reconciliation of earlier removals. No numerical alias estimate, sector removal, new charge or floor change follows; current98/31 errors remain. Actual mixed source, selected arithmetic content, prescribed period, signed inequality, specific032, DP44 and RH remain open. AM-R-085 RF1--31 gives the exact genuine selected vector-valued central coefficient c(t)=P(t)/(z0-epsilon i gamma t)\^{}s, with s=kappa+K+2N and degree P\textless=s-kappa, using the original Pieri component and its exact Hilbert/PBW norm. The finite coefficient rule, partial fractions, signed gamma Fourier transform and affine Poisson phases are explicit. The positive central density for the same selected component has mass one and pointwise dominates the transformed commuting fibre, without assuming its projection commutes. Products over all168 labels lie in the actual signed orthant and vanish on its boundary faces. Ordinary absolutely convergent Poisson holds for the actual lattice and finite characters supplied separately by DO106, with no quotient-cardinality loss. This removes the unspecified analytic-model step but proves no numerical dual-lattice alias bound, finite sector removal or new error charge. Current98/31 totals and floors25/11 remain; the signed inequality, actual source, selected arithmetic image, prescribed period, specific032, DP44 and RH remain open. AM-K101 UC1--17 removes the complete actual fully eligible all-II globally rank-two-unipotent scalar-cross support25 subset with2\^{}1585\textless=x=N\_E D\_eta\textless2\^{}1820 for every22\textless=k\textless=164. It uses the original DO105 all-factor bound with a single common denominator endpoint, exact capmax(161,k), full selected trace and moving affine central cover. The143-cell union is below2\^{}-254, with nondense maximum-1412571/4096\textless-344 and dense maximum-1048027/4096\textless-255. Append254 once to actual K10098/31, giving99/31 and both whole sums below2\^{}-230; all earlier overlaps,451 and421 are retained and complete floors25/11 remain. The only still-unremoved support25 cells therefore have2\^{}120\textless x\textless2\^{}1585,22\textless=k\textless=164, with K100 additionally excluding22\textless=k\textless=48 when x\textgreater=2\^{}1116, and all K099 structural restrictions retained. No prospective inactive-T gain, signed positivity, actual source, selected arithmetic content, prescribed period, specific032, DP44 or RH result is inferred. AM-DO-107 IP1--23 removes three complete subsets of the actual fully eligible all-II globally rank-two-unipotent scalar-cross support25 middle sector: any inactive paired-C-zero place for all k22..164; no paired zero with x\textgreater=2\^{}744 for all k22..164; and no paired zero with x\textless2\^{}744 and22\textless=k\textless=99. The literal included-Lambda projector gives rho=1/(T-1) at every noneligible inactive place with a nonzero paired C block. A paired zero instead supplies an actual extra normT² row ideal, giving T\^{}-335 while retaining the old central affine cover. The complete143/143/78-cell union is below2\^{}-392. Append392 to the actual K10199/31 arrays, giving100/31; all historical overlaps including451,421,254 remain, both full sums are below2\^{}-230, and complete floors25/11 are unchanged. The only support25 branch left has no inactive paired C zero, total eligibility exactly25,2\^{}120\textless N\_E D\_eta\textless2\^{}744 and100\textless=k\textless=164, with K099 field/stabilizer restrictions retained. This is a strict smaller finite gate, not support25 exhaustion, signed positivity, an admissible mixed source, selected arithmetic image, prescribed global period, specific032, DP44 or RH. AM-DO-108 AC1--15 removes the entire actual support25 fully eligible all-II globally rank-two-unipotent scalar-cross subset with any active paired-C-integral place, at cost less than2\^{}-55281. The exact extra normT² row ideal gives T\^{}-335; the central cover uses the old scale. Its143 complete cells append55281 to actual DO107100/31, giving101/31, all historical charges retained and both sums below2\^{}-230, floors25/11 unchanged. Remaining middle data have j=b=0, exactly25 total eligible places,2\^{}120\textless x\textless2\^{}744 and100\textless=k\textless=164. Every complete actual central orbit then keeps all168 original T factors equal to1/(T-1), with finite weight(T-1)\^{}-168 a(h(t)), sufficient periodQGamma, indexQ\^{}168, exact trace dual and ordinary Poisson formula AC15. This strictly removes active-integral profiles and removes their indicator from the exact Fourier interface; it gives no alias estimate or automatic short-sum/fibre replacement. Admissible source, selected arithmetic content, signed nonvanishing, global period, specific032, DP44 and RH remain open. AM-LIFT-084 MP1--25 constructs a literal field-valued Obj and FrPt for any given actual source Host, compatible Par.PinHost/PinLevel/PinDat, specified stage Datum, existing Moduli representation and same Coeff field, after a finite extension. The actual principal ideal identity is derived from Host fields. A pure principally polarized CM block and its conjugate action supply every uniform(p,q); the explicit trace-one semicharacter supplies actual SAB prim/adjoint/cubical lines. Relative-unramified selfdual Hermitian lattice uniqueness, including dyadic places, gives the required profinite Tate isometry; the literal commutator sign and CRT preserve the exact finite level M=g0g1,A=1. All-a normalized line identities descend over one finite field; only geometric full Tate systems are asserted. Separately, on the existing E\_h fields the explicit pure block and three J-action twists give J\^{}(fq) times C\_h\^{}(p-2q) for p\textgreater=2q, attaining LIFT083 complement dimension f(p-2q). Those fixed fields still fail full Sizes. No effective field/frame/line costs, good-model/old LamPt identification, prescribed same-point F nonzero, actual full Host/Sizes constructor, selected arithmetic content, prescribed period, specific032, DP44 or RH is supplied. No numerical error charge;101/31 arrays and floors25/11 remain. AM-R-086 ZJ1--21 constructs principal source points with actual J content over any fixed actual Host E under the retained LIFT084 input hypotheses. The auxiliary compositum L=E Q(zeta7), c={[}L:E{]} in\{1,2,3,6\}, has an explicit integral CM trace polarization of degree\textbar D\_L\textbar7\^{}(cf/3). Zarhin Theorem1.1 gives a full integer-ring compatible principal block. Its exact conjugate-action balancing, with the even-c improvement, yields block sizes a\_c=(8,8,24,24). For p\textgreater=q the resulting exact J exponent is N=(2f/3)(q-r), r=q mod a\_c, and pure complementary dimension f(p-q+2r). If d=7 and p\textgreater=2q, a principal(16,8) block gives N=f(q-r7), r7=q mod8, complement dimension f(p-2q+3r7). Zero remainders attain the corresponding LIFT083 bounds. The unchanged Datum, finite level and Coeff give actual Obj/FrPt over a finite extension. A separate full O\_K7 lattice-approximation argument gives an underlying isogeny J\^{}N to the actual J factor prime to any prescribed finite set, and an isomorphism of p-divisible group schemes after good reduction. No old theta-line/frame compatibility is inferred. Effective field/isogeny/line/frame costs, numerical Host/Sizes existence, same-point F nonzero, actual lift identities/arithmetic image, global period, specific032 and DP44 remain open; no RH or finite numerical charge. Current101/31 arrays and complete floors25/11 are unchanged. AM-LIFT-085 HF1--26 proves that, for an existing literal Host/Datum/Moduli/Coeff, actual CM seed Obj, original finite level A=1 and a given nonzero natural frame rule, a point in that same OE-isogeny class has a nonzero literal FrPt evaluation at the unchanged finite level. The isogeny u has degree N\^{}(2g), N prime to any prescribed finite set S containing the level primes, N=1 mod M, and u\emph{L1=L0\^{}(N\^{}2) rigidified and cubical. Every output object and a nonzero frame descend over K1=K0(A0{[}N\^{}2{]}), with {[}K1:K0{]}\textless=N\^{}(8g\^{}2). The argument meets every analytic moduli component through an algebraic family based on its own target lattice. A weight-rule frame can be chosen by integer matrices with entries0..D+1 relative to any initial rational eigenframe. N itself, seed degree, initial-frame/theta-generator heights, models and actual line costs remain unevaluated. The given rule must itself be nonzero; no prescribed arithmetic-image nonvanishing is proved. No old pure LamObjGood or selected derivative identification, global period, specific032 pairing, DP44 or RH conclusion follows. This removes the qualitative same-rule same-point existence gate within the fixed CM isogeny class under the precise assumptions. Current101/31 error arrays and complete floors25/11 remain unchanged. AM-K-102 PC1--27 bounds the entire actual DO107 IP3 support25 remainder by2\^{}-3550: all-II globally rank-two-unipotent scalar-cross, exactly25 total eligible places, no inactive paired cross zero,2\^{}120\textless N\_E D\_eta\textless2\^{}744 and100\textless=k\textless=164. The relative-norm phase count, positive-cone Gaussian/Poisson majorant and true transverse metric use one joint energy constraint; the complete original complex-volume tail covers every larger ideal-norm parameter. The original25 active weighted-label factors and143 inactive pointwise factors are retained without assuming complete-orbit constancy. All195 endpoint bounds are below-3557; the65-cell union costs3550 once. Actual DO108101/31 becomes102/31, every old overlapping charge retained, both complete sums below2\^{}-230 and floors25/11 until a separate exhaustive assembly. No signed positivity, prescribed period, selected arithmetic-image witness, effective admitted source cost, specific032, DP44 or RH claim follows. AM-R-087 J26.1--5 exhausts every original all-top support25 class at the unchanged H133 cutoff. The complete F/IC8 rank1..33 ranges force A=I; DO090/R080/R079 exhaust fixed S; DO099/DO095 give the precise cross rank2/3 alternatives; R082/R084 exhaust rank3; DO101 and DO102 remove nonunipotent rank2 and all non-all-II unipotent profiles; the exact DO103/104/105/107/K102 denominator, eligibility and coordinate partition exhausts the final all-II branch, including every equality boundary. The complete all-top floor is26 and the independent uniform floor remains11. Every102/31 chronological error entry is unchanged, both full sums stay below2\^{}-230 and no assembly charge is added. The actual sufficient signed gate is J26.2; global period, fully sized admissible source, selected arithmetic image, specific032/Read, actual DP44, Kummer-Hecke RH and the global oscillator comparison remain open. AM-R-088 GZ1--11 removes the actual fixed-S rank-one-unipotent support26 integral-witness branch, every reduced double-pole cell d\textless=22, and every reduced pattern with at least one nonidentity unmixed G reduction. The exact averaged-content and original affine-coordinate count retains every paired type/label factor and the actual selected trace. The full334-place nonempty weighted G-pattern sum is below2\^{}-5533, replacing only the former factor2. All27 complete endpoints and the witness give one523 charge. Only the reduced all334Gidentity pattern with d23..26 remains in this branch. Actual R087102/31 becomes103/31, every old overlap retained, both whole sums below2\^{}-230, and complete floors26/11 unchanged. No full support26 exhaustion, signed positivity, actual source, selected arithmetic content, prescribed period, specific032/Read, DP44 or RH conclusion follows. AM-LIFT-086 IM1--27 constructs actual everywhere-good polarized models for the LIFT085 geometric CM seed and target over K=K0(A0{[}15N\^{}2{]}), with {[}K:K0{]}\textless=(15N\^{}2)\^{}(4g\^{}2). The full integral OE action, principal polarization, canonical trace-one cubical line, isogeny and adjoint extend over the same ring. Actual intrinsic tangent-line degrees change by an amount in{[}-log N,log N{]}, the full determinant degree by at most g log N, and the theta degree by at most g log N/2 under KeyCore. Target intrinsic lines have length constant1. If the actual seed lines have degree at least-b and all genuine rule/integral/cusp/source premises hold, the sufficient Step5 cost is C\_N=(D+gn)/2+D log r+(b+log N)min(D+gn,(3D+gn)/2). No effective bound on N, initial field or arbitrary-Host b is supplied. Host/Sizes, actual rule nonzero/content, prescribed period, specific032/Read, evaluated DP44 and RH remain open. No new finite error charge; actual103/31 arrays and complete floors26/11 remain unchanged. AM-K-103 SC26.1--28 removes nine complete subsets of the actual all-II globally rank-two-unipotent scalar-cross support26 class. Its support-free full-Z proof, new eta window and uniqueness, exact paired ideals and inactive weights, full infinite ideal tail, joint energy and positive-norm-cone counts retain every original factor. All1,229 cells give one293 charge. The residual necessarily has v=q26, 2\^{}1116\textless N\_E D\_eta\textless2\^{}2108, b=j0 and81\textless=k\textless=164. Actual R088103/31 becomes104/31, all earlier overlapping charges retained, both entire sums below2\^{}-230 and complete floors26/11 unchanged. Other types/classes, signed positivity, actual source, selected arithmetic content, prescribed period, specific032/Read, DP44 and RH remain open. AM-DO-109 SP26.1--25 removes the entire actual fixed-S globally rank-two-unipotent support26 low-D\_G\textless=32 short-pivot class. The original plane reduction and paired alpha budget give2i+t\textless=5. The actual V-input eligibility and central-only label proof, true Gram metric, affine central difference lattice and full weighted four-state G sum are retained. The three complete p\textless=24,p25,p26 endpoints give one820 charge; the p26 branch proves eta uniqueness before summing actual labels. Together with already charged R079 high-D\_G and last-only pivots, this completes the whole fixed-S rank-two-unipotent support26 structural class. Actual K103104/31 becomes105/31, every overlap retained, both entire sums below2\^{}-230 and complete floors26/11 unchanged. Other support26 classes, signed positivity, actual source, selected arithmetic content, prescribed period, specific032/Read, DP44 and RH remain open. AM-K-104 FV1--14 removes the complete new coordinate interval81\textless=k\textless=95 from the actual K103 all-II globally rank-two-unipotent scalar-cross support26 residual. The original full joint-volume lattice, actual central affine fibre and selected trace are retained. The full334-place four-state weighted G product is below2; the exact support binomial is below2\^{}101. Complete coupled-x monotonicity gives the actual lower denominator endpoint without a missing range. All15 cells, with worst-1128627/4096, give one274 charge. The all-II residual now has v=q26, 2\^{}1116\textless x\textless2\^{}2108, b=j0 and96\textless=k\textless=164. Actual DO109105/31 becomes106/31, all prior overlaps retained, both entire error sums below2\^{}-230 and complete floors26/11 unchanged. Other classes, signed positivity, actual source, selected arithmetic content, prescribed period, specific032/Read, DP44 and RH remain open. AM-R-089 WN1--18 removes the complete actual fixed-S rank-one-unipotent support26 remainder: all334Gidentity and intrinsic double-pole sizes23..26. Its universal reciprocal-norm weighted eta/content lemma uses determinant-one rectangular rescaling and full dyadic shell packing, including every unit and ideal class. The original selected trace retains its actual eta norm before summation; both genuine affine coordinates and the whole actual support/marked/type/pair sum remain. The complete rational upper exponent is-146697349/4096\textless-35000, giving one35000 charge. Together with already charged R088 this completes the whole fixed-S rank-one-unipotent support26 structural class. Actual K104106/31 becomes107/31, every old overlapping error retained, both entire sums below2\^{}-230 and complete floors26/11 unchanged. Other support26 classes, signed positivity, actual source, selected arithmetic content, prescribed period, specific032/Read, DP44 and RH remain open. AM-DO-110 NT26.1--31 removes the complete actual non-all-II globally rank-two-unipotent scalar-cross support26 sector. The support-independent full-positive normal form, original paired cross/full-entry poles, 24-place eta uniqueness, actual one-label cancellation including central-only poles, full affine central lattice, common-G ideal and exact joint count are retained. Seven exhaustive branches and all1148 cells give one561 charge. Actual R089107/31 becomes108/31, every overlap retained, both complete sums below2\^{}-230 and full floors26/11 unchanged. The remaining all-II unipotent26 class is the K103/K104 dense window; other scalar-cross classes, signed positivity, actual source, arithmetic content, prescribed period, specific032/Read, DP44 and RH remain open. AM-LIFT-087 GS1--20 chooses an actual new general-Host Serre seed X=O\_L tensor\_\{O\_K7\} J and the actual pure block b\^{}-1 tensor\_\{O\_k\} C0, preserving the Host ideal, Datum, level, coefficient premises and R086 signature/J-content formulas. Every constructed intrinsic component eigenline has degree at least-b, with b=max(log D\_L+.5 log(2t),(sqrt(d)+1)/3+log D\_E+.5 log f). Literal integral b gives1 in b\^{}-1; full-order discriminant-cleared projectors and the exact Zarhin pullback isometry retain all integral models and source pi H norms. The new trace-polarization degree is D\_L/7\^{}(5t), distinct from R086. Under all retained actual rule, KeyCore, integrality/cusp/source premises, LIFT086 now has explicit seed b in C\_N=(D+gn)/2+D log r+(b+log N)min(D+gn,(3D+gn)/2). Effective N, initial field, actual Host/Sizes constructor, selected image, prescribed period, specific032/Read, DP44 and RH remain open. No finite charge; current108/31 arrays and complete floors26/11 unchanged. AM-R-090 IS1--25 removes the entire actual full-positive globally rank-three scalar-cross support26 class with nonzero original S-cross, without a sparsity hypothesis. The complete eligible-set count at most24 costs6670; the complete all-k denominator tail at least2\^{}8192 costs11965. Literal inversion applies each containing-set bound once more. Outside them, the two eligible sets intersect in at least24 paired places, forcing equal null lines by the exact6395-bit denominator-cleared norm gap and contradicting global rank three. The full four-sector union is below2\^{}-6668. Actual DO110108/31 becomes109/31, every previous overlapping charge retained, both complete sums below2\^{}-230 and floors26/11 unchanged. The nonzero short-compression and dense all-II rank-two classes, signed positivity, actual source, selected arithmetic content, prescribed period, specific032/Read, DP44 and RH retain their gates. AM-DO-111 VC1--25 proves the actual positive transverse compression is unitary when both original scalar-cross directions lie in the retained plane. Its four original integral short witnesses have the sharper all-embedding bound2\^{}41. The complete simultaneous finite restoration, arbitrary-ideal row volume, relative norm fibre, original selected trace, full W0 column, source fibre and all paired type weights give exact exponent-87517717/4096\textless-21366 for the entire nonzero-short-compression support26 sector. The complementary integral norm argument gives full positive167 compressionI and a positive166 fixed space, hence global rank at most3; global unipotence is not inferred. Actual R090109/31 becomes110/31 with one21366 append, every old overlapping charge retained, both entire error sums below2\^{}-230 and complete floors26/11 unchanged. Whole support26 assembly remains separate; dense all-II rank-two, signed positivity, actual admissible source, selected arithmetic content, prescribed period, specific032/Read, DP44 and RH remain open. AM-LIFT-088 SF1--16 supplies an explicit field for the entire newly chosen LIFT087 general-Host seed over the same given Coeff field B: {[}K\_seed:B{]}\textless=168}23040\emph{h(-d)}M\^{}(4g\^{}2)\textless=3870720\emph{d}M\^{}(4g\^{}2), g=f(p+q). The actual prescribed elliptic CM action descends over B(j) by scalar differential injectivity. The explicit J good-model radicals, elliptic15-torsion, actual projective Serre modules, equivariant principalization, canonical trace-one cubical line, original finite M-level and literal Obj/FrPt all lie over the stated field. Under every retained LIFT085 nonzero-rule/source hypothesis, {[}K\_out:B{]} is at most that bound timesN\^{}(8g\^{}2), with the existing intrinsic Step5 cost unchanged. EffectiveN, absolute degree of the givenB, field discriminants, coefficient/section heights, numerical Host/Sizes existence, selected arithmetic content, prescribed period, specific032/Read, evaluated DP44 and RH remain open. No finite error charge; current110/31 arrays and complete floors26/11 are unchanged. AM-R-091 DR1--7 consolidates every checked support26 cut through DO111 into one actual-remainder theorem. The only possible support26 terms are all-type-II globally rank-two-unipotent scalar-cross matrices with full positive167 compression, nonzero original S cross, v=q26, b=j0, strict denominator window2\^{}1116\textless N\_E D\_eta\textless2\^{}2108 and96\textless=k\textless=164. The actual retained support26 set and actual higher-support remainder split the unchanged energy exactly, with the complete110/31 errors below2\^{}-230 and floors26/11. Every earlier cut and overlap charge is retained. No charge or independent new movement is assigned to this consolidation. DR5, actual fully sized source, selected arithmetic content, prescribed period, specific032/Read, evaluated DP44, Kummer-Hecke RH and global oscillator comparison remain open. AM-K-105 CG1--25 removes all39 complete coordinate sectors96\textless=k\textless=134 from the actual all-II globally rank-two-unipotent support26 scalar-cross remainder. The scalar real-dimension336 Poisson bound, full coupled Gram determinant, original selected trace, full affine central covolume and every finite/source factor remain. Restoring(N\_FJ0)\^{}(-1) and its actual common-G correlation gives the full weighted pattern factor\textless2Q\^{}(-167); complete coupled-x monotonicity gives the exact lower denominator endpoint. The worst cell is-9009839/6144\textless-1466 and the preceding38 are below-2831, so one1465 charge covers the union. Actual DO111110/31 becomes111/31, every overlap retained and both complete sums below2\^{}-230; floors26/11 stay fixed. With R091, the only possible support26 branch now has135\textless=k\textless=164, the same v=q26,b=j0 and strict denominator window2\^{}1116\textless x\textless2\^{}2108. Signed positivity, actual source, selected arithmetic content, period, specific032/Read, DP44 and RH remain open. AM-LIFT-089 HF89.1--18 sharpens the degree of exactly the same LIFT088 seed and output fields. For r=p+q,g=fr, the actual projective OE-lattice is free modulo each finite ideal, and the full defined action bounds its torsion Galois image. The literal m/bar(m) level halves are perfect Weil duals; Coeff.zeta\_mem already gives mu\_M inB. Thus Kgood(A{[}M{]})=Kgood(A{[}m{]}) and {[}Kseed:Kgood{]}\textless=M\^{}(f r\^{}2). The same target field Kseed(A{[}N\^{}2{]}) has relative degree\textless=N\^{}(4f r\^{}2), including ramifiedN andN1. Consequently {[}Kout:B{]}\textless=3870720\emph{d}M\^{}(f r\^{}2)\emph{N\^{}(4f r\^{}2), under every actual LIFT085 nonzero-rule/source hypothesis. All actions, lines, level, Obj/FrPt and intrinsic costs are unchanged. EffectiveN, absoluteBdegree, discriminants, coefficient/frame heights, actual Host/Sizes existence, selected arithmetic content, period, specific032/Read, DP44 and RH remain open. No error charge; current111/31 arrays and complete floors26/11 are unchanged. AM-DO-112 EC1--22 removes both complete coordinate sectors k135 and k136 from the actual all-II globally rank-two-unipotent support26 scalar-cross remainder. The larger Gaussian cap59/(50}65536), with the full real dimension336 tail and scalar theta bound26, preserves the original row metric, affine central divisor, selected trace, every finite factor and complete coupled pattern/denominator monotonicities. Exact whole-sector exponents are-691859209/307200\textless-2252 and-283283209/307200\textless-922; their union is below2\^{}-921. Actual K105111/31 becomes112/31, every historical overlap retained and both complete sums below2\^{}-230; floors26/11 stay fixed. The only possible support26 branch now has137\textless=k\textless=164, v=q26,b=j0 and strict denominator window2\^{}1116\textless x\textless2\^{}2108. The positive k137 diagnostic is unclaimed. Signed positivity, actual source, selected arithmetic content, period, specific032/Read, actual DP44 and RH remain open. AM-LIFT-090 SG1--21 sharply bounds the same actual LIFT087 seeds and LIFT088 fields through the two original generating CM varieties J and prescribed C0. Integral projective split maps and their duals put every seed n-torsion in Kgood(J{[}n{]},C0{[}n{]}) at every n, with no prime-to-index restriction; full maximal-order rank-one actions bound that field by n\^{}8. Thus the same seed and output fields cost M\^{}8 and N\^{}16, and a finite list of stages sharing J,C0,Kgood costs one M\^{}8 and one lcm(N\_j)\^{}16. The same C0{[}15{]} field has degree\textless=225 before good reduction, sharpening the same ground factor to37800\emph{h(-d)\textless=37800}d.~The complete relative bound is37800\emph{d}M\^{}8\emph{lcm(N\_j)\^{}16. Retain the smaller LIFT089 bounds where applicable. The actual models, actions, polarizations, canonical line, finite level, Obj/FrPt and intrinsic cost remain unchanged. EffectiveN, absoluteBdegree, discriminants, coefficient heights, all nonzero-rule/source premises, selected arithmetic content, prescribed period, specific032/Read, actual DP44 and RH remain open. No error charge; current112/31 arrays and complete floors26/11 remain. AM-K-106 NT1--14 removes the complete coordinate sector k137 from the actual all-II globally rank-two-unipotent support26 scalar-cross remainder. The exact universal radial-profile factor carries through the original rank-two near-point trace and yields630 bits. Every nonzero original S coordinate lies in the same fractional ideal; the proved product-of-sums inequality gives k\^{}-84 and596 further bits at137. Applied before the full DO112 Gaussian cardinality enlargement, these constant factors preserve every actual metric, finite weight, central affine covolume and complete pattern/denominator count. The exact whole-sector exponent is-251641609/307200\textless-819. Actual DO112112/31 becomes113/31, all previous overlapping charges retained, both complete sums below2\^{}-230 and floors26/11 unchanged. The only possible support26 branch now has138\textless=k\textless=164 with the same v=q26,b=j0 and2\^{}1116\textless x\textless2\^{}2108. Signed positivity, selected arithmetic content, prescribed period, full effective source, specific032/Read, actual DP44 and RH remain open. AM-LIFT-091 RD1--18 bounds root discriminants of the same actual LIFT088--090 finite torsion fields. For a finite Galois extension of actual degree D, unramified outside an integral ideal S, rd(L)/rd(K)\textless=D}Norm(rad S)\^{}(1/{[}K:Q{]}); the complete base-normalized local different and actual-degree wild contribution are proved. Thus the simultaneous seed satisfies rd(KSigma)\textless=rd(Kgood)\emph{rad(M)}M\^{}8, and the same actual output satisfies rd(KSigma,out)\textless=rd(Kgood)\emph{rad(M L\_N)}lcm(M,L\_N²)\^{}8\textless=rd(Kgood)\emph{M\^{}9}L\_N\^{}17. Its relative increase over KSigma is at most rad(L\_N)\emph{L\_N\^{}16. The literal N\_j=1 mod M congruence is retained, so no false level/output-overlap saving is claimed. Full all-object descent and the37800}d\emph{M\^{}8}L\_N\^{}16 relative degree bound are unchanged. Initial rd(Kgood), absolute B degree/discriminant, effective N\_j, heights, all nonzero-rule/source conditions, selected content, prescribed period, specific032/Read and evaluated DP44 remain open. No numerical error charge;113/31 full arrays and floors26/11 remain. AM-R-092 ID1--18 bounds the initial root discriminant for the exact SF5 seed C0 and radical/torsion tower of LIFT088, over the actual source coefficient field B. The maximal-order class-field step B(j)/B is unramified; the actual radical tower has relative root-discriminant factor at most76832. Every CM conjugate satisfies \textbar sigma(j)\textbar\textless32 exp(pi sqrt(d)), and the actual C0{[}15{]} extension is treated at its true degree before applying the bound225. Consequently rd(Kgood)\textless2\^{}50 exp(2 pi sqrt(d)) rd(B). Composed with LIFT091, rd(KSigma,out)\textless2\^{}50 exp(2 pi sqrt(d)) rd(B) rad(M L\_N) lcm(M,L\_N\^{}2)\^{}8\textless=2\^{}50 exp(2 pi sqrt(d)) rd(B) M\^{}9 L\_N\^{}17. The stated log absolute discriminant bound retains the exact relative degree37800\emph{d}M\^{}8\emph{L\_N\^{}16 and actual {[}B:Q{]},rd(B). This removes the separate initial-good-field discriminant premise for this fixed seed, not for arbitrary twists. Absolute B degree/discriminant, effective N\_j, heights, Host/Sizes existence, selected nonzero arithmetic rule/content, period, specific032/Read and evaluated DP44 remain open. No new finite-error charge;113/31 full arrays and complete floors26/11 unchanged. AM-LIFT-092 CD1--19 strengthens the root-discriminant cost of the same actual LIFT088--090 fields. For the two full-maximal-order CM progenitors J,C0 over the actual everywhere-good Kgood, every actual F inside Kgood(J{[}n{]},C0{[}n{]}) is abelian and has conductor dividing n O\_Kgood and discriminant dividing (n O\_Kgood)\^{}(D-1), D={[}F:Kgood{]}. Thus rd(F)/rd(Kgood)\textless=n\^{}(1-1/D)\textless=n, including all dyadic/CM-ramified primes and D=1. The same simultaneous output satisfies rd(KSigma,out)\textless=rd(Kgood)}lcm(M,L\_N\^{}2)=rd(Kgood)\emph{M}L\_N\^{}2 under the actual N\_j=1 mod M; its relative output increase over KSigma is at most L\_N\^{}2. Combining the separate canonical R092 bound gives rd(KSigma,out)\textless2\^{}50 exp(2 pi sqrt(d))\emph{rd(B)}M\emph{L\_N\^{}2. LIFT090 relative degree37800}d\emph{M\^{}8}L\_N\^{}16 is unchanged; the sufficient absolute-discriminant logarithm is this degree times {[}B:Q{]} times (log rd(B)+50 log2+2 pi sqrt(d)+log M+2 log L\_N). All original fields, varieties, polarizations, line, level, quotient, action, source points and frames are retained. Actual B degree/discriminant, effective N\_j, metric/height costs, Host/Sizes, selected nonzero arithmetic content, period, specific032/Read and evaluated DP44 remain open. No finite-error charge;113/31 full arrays and complete floors26/11 unchanged. AM-K-107 SV1--36 removes the entire k138 cell of the actual all-II globally rank-two-unipotent scalar-cross support26 remainder. The genuine selected density yields a full tail-inclusive squared coefficient below2\^{}-933 and a single actual near-point product below2. Together with the exact radial factor, all-coordinate norm and original saturated transverse lattice, the complete count is bounded by2\^{}(-50428781/204800)\textless2\^{}-246. The original scalar theta is in real dimension336, with full21\^{}(k+3) and all original central, support, row, finite and pattern factors. The remaining branch has139\textless=k\textless=164 and2\^{}(71443/64)\textless x\textless2\^{}2108, v=q26,b=j0. Exactly one246 append gives114/31 ordered errors with all old overlaps retained, both full sums below2\^{}-230, unique minimum231 and next minimum246; complete floors26/11 remain. Mandatory corrected boundary conventions: SV6 is strict for h\textgreater=1 and equality at h0; SV11 quotient is for r\textgreater0 and the bound at r0 is1. The frozen proof is read with the explicit boundary addendum reproduced below. Signed positivity, selected arithmetic content, period, effective source, specific032/Read, actual DP44 and RH remain open. AM-DO-114 SP1-27 proves a uniform scalar theta below3/2 in real dimension336 for A\textless=1/43500 by a complete finite spherical-cap and nonzero-shell argument. Applied to the original saturated transverse Gram, actual selected trace, full affine central fibre and every finite pattern, the whole139..141 union is below2\^{}-354. The necessary actual support26 all-II rank-two-unipotent range becomes142\textless=k\textless=164 with2\^{}(71443/64)\textless x\textless2\^{}2108,v=q26,b=j0. Exactly one354 append gives115/31 complete ordered errors, both full sums below2\^{}-230, with every historical overlap retained and complete floors26/11 unchanged. No support27, signed positivity, selected arithmetic content, period, effective source, specific032/Read, evaluated DP44 or RH conclusion. AM-K-108 AE1-30 retains the full actual Cartesian conditional angular density, cancels its D ratio in detailed balance and inserts exp(-2U\textbar\textbar r\textbar\textbar\^{}2) in the original lattice sum. With the separately checked DO114 scalar theta, original transverse Gram, central affine fibre, common ideal correlation and all finite factors, the complete142..159 union is below2\^{}-585. The necessary support26 all-II rank-two-unipotent remainder becomes160\textless=k\textless=164, with2\^{}(71443/64)\textless x\textless2\^{}2108,v=q26,b=j0. Exactly one585 charge gives116/31 ordered errors, both full sums below2\^{}-230; floors26/11 remain. No160+, signed positivity, arithmetic-content nonvanishing, period, effective source, specific032/Read, evaluated DP44 or RH conclusion. AM-R-093 DC1-19, on the actual d=7 branch, proves one finite-prime-unramified quadratic character epsilon with N\_E7/k(rho\_J)=epsilon\emph{chi\_cyc}rho\_C\_plus at every Tate prime. The integral rank-three norm at all primes, including2/7, gives K(J{[}n{]},C0{[}n{]}) contained in K(J{[}n{]})K\_epsilon, with relative degree at most2, simultaneously for every n.~Same actual projective/dual seeds and source congruences then give {[}Kgood:B{]}\textless=16128 and {[}K\_Sigma,out:B{]}\textless=32256\emph{M\^{}3}L\_N\^{}(129/16). The norm/type computation and arithmetic Artin convention are explicit; epsilon is neither evaluated nor assumed trivial. Root-discriminant bounds remain R092/LIFT092 on the same fields. Absolute coefficient-field degree/discriminant, effectiveN, Host/Sizes, selected arithmetic nonzero content, period, Read/specific032, evaluatedDP44 and RH remain open. No finite-error charge;116/31 and floors26/11 unchanged. AM-LIFT-093 PA1-27 refines the general-d degree for the exact same Kgood, J,C0 and original source seeds. Principal Weil pairings impose one shared scalar multiplier. Exact all-prime norm counts, including2/7/d and the overlappingd7 case, give {[}F(J{[}n{]},C0{[}n{]}):F{]}\textless=U(n)=n\^{}5\emph{product\_(p\textbar n,p=6 mod7)(1+1/p)\^{}3 for every finiteF/Kgood, and U(N\^{}2)\textless=N\^{}(81/8). The actual source level has multiplier1 and boundV(M)=M\^{}4}product}(p\textbar M)F\_J(p)(1-1/p)\textless2M\^{}4; no universal g mod7 condition is inserted. The existing C0{[}15{]} extension costs at most192, so {[}Kgood:B{]}\textless=32256h(-d)\textless=32256d and the actual common output satisfies {[}Kout\_Sigma:B{]}\textless=32256d V(M)U(L\_N\^{}2)\textless64512d M\^{}4 L\_N\^{}(81/8). R093 retains its independent strongerd7 bound. Same fields, all objects and intrinsic Step5 costs remain; effectiveN, absoluteB, Host/Sizes, selected nonzero arithmetic content, period, Read/specific032, evaluatedDP44 and RH are open. No finite-error charge;116/31 and floors26/11 unchanged. AM-K-109 TG1-19 retains the full all-A scalar theta polynomial and the original transverse Gram G\textgreater2\^{}36 I, with true determinantU. Using the entire actual K108 energy, it removes the complete k160 support26 cell with exact exponent-163509/512\textless-319. Original module membership, central affine fibre, common ideal correlation, all finite factors and both q/x monotonicities are retained. The necessary actual all-II rank-two-unipotent remainder becomes161\textless=k\textless=164 with2\^{}(71443/64)\textless x\textless2\^{}2108,v=q26,b=j0. Exactly one319 append gives117/31 complete ordered errors, both entire sums below2\^{}-230; complete floors26/11 remain. No161+, signed positivity, selected arithmetic content, effective source, period, specific032/Read, evaluatedDP44 or RH conclusion. AM-DO-115 DC1-24 proves an all-parameter scalar theta estimate using the exact positive cubic kernel, finite cap occupancy and nonzero-shell integral in real dimension336. The original selected angular profile gives beta777/352, and the unchanged saturated transverse Gram gives a complete257/256 theta factor. With every original paired ideal, affine central fibre, finite factor and both continuum monotonicities retained, the complete k161 support26 cell is below2\^{}-233, with exact exponent-119477/512. The necessary actual all-II rank-two-unipotent remainder becomes162\textless=k\textless=164 with2\^{}(71443/64)\textless x\textless2\^{}2108,v=q26,b=j0. Exactly one233 append gives118/31 complete ordered errors, both entire sums below2\^{}-230; unique minimum231,next233 and complete floors26/11 remain. No162+, signed positivity, selected arithmetic content, effective source, period, specific032/Read, evaluatedDP44 or RH conclusion. AM-K-110 PM1-22 retains the full actual angular profile as a positive Laplace mixture, includes the entire low-parameter tail in an explicit atom, and computes the exact inverse third moment. Weighted AM-GM and the all-minimum/all-parameter DO115 scalar theorem control arbitrary embedding weights in the unchanged original ideal and transverse Gram. Every transverse correction is retained, giving a uniform2180-bit row improvement. With all original affine central, paired ideal, finite and source factors and both continuum monotonicities, the complete k162 support26 cell is below2\^{}-1071, with exact exponent-548533/512. The necessary actual all-II rank-two-unipotent remainder becomes163\textless=k\textless=164 with2\^{}(71443/64)\textless x\textless2\^{}2108,v=q26,b=j0. Exactly one1071 append gives119/31 complete ordered errors, both entire sums below2\^{}-230; minimum231,next233 and complete floors26/11 remain. No k163/164, signed positivity, selected arithmetic content, effective source, period, specific032/Read, evaluatedDP44 or RH conclusion. DO116 provides an exact stronger all-minimum/all-parameter scalar coefficient for future complete applications; it changes no actual finite sector or main gate by itself. K111 supplies an exact stronger all-minimum/all-parameter scalar coefficient; it introduces no separate arithmetic charge or main-gate reduction before a complete application. AM-LIFT-094 HE0-22 proves conditional effective algebraic Hecke escape on a fixed indefinite arithmetic family. For a reduced algebraic bad locus of total degreeD and dimension\textless=s, finite equal-ample-pullback correspondences of degree\textless=delta with no nonempty proper closed forward-invariant locus admit escape in at mostB\_delta,s(D) steps, where A\_delta(D)=2 delta\^{}D D\^{}2, B\_delta,0(D)=D and B\_delta,s(D)=D+B\_delta,s-1(A\_delta(D)). For actual projective O\_E lattice rankr\textgreater=3 indefinite at every real place, f={[}F:Q{]}, an explicitly cleared rational Witt frame gives t=1+3McQ\_S, two physical original-M-level Hecke correspondences with delta=t\^{}(4fr\^{}2) and multiplier t\^{}2. Every forward orbit is dense; the conditional accumulatedN\textless=t\^{}B is obtained. The actual regular LIFT085 period-family map and nonempty inverse image of the prescribed nonvanishing open remain hypotheses; H,D,c,N0 and admitted Host/Sizes costs are not evaluated. Same original fields, quotient/action/line/level/frame conditions and independent entry multiplierN0 remain. No selected arithmetic nonzero content, global period, specific032/Read/STiesR, DP44 or RH conclusion. No finite-error charge;119/31 and floors26/11 unchanged. AM-R-094 PP1-24 removes the entire remaining conditional J56 all-top support26 class. R091 plus K109/DO115/K110 leaves only actual all-II global-rank-two-unipotent scalar-cross cells k163/164; the full correlated finite-pattern containing polynomial and low-parameter angular mixture give complete exponents -199349/512\textless-389 and -159413/512\textless-311, using DO116 and K111 respectively. Their union is below2\^{}-310. Append310 once to the exact K110119/31 arrays:120/31, both entire sums\textless2\^{}-230, complete floors27/11. The actual remainder now has support at least27 with every earlier higher cut and original finite/real/arithmetic convention retained. PP3, the sufficient signed higher-support inequality, remains open, as do actual Host/Sizes, selected nonzero arithmetic content, global period, specific032/Read/STiesR and evaluatedDP44/RH. No actual shorter mixed-CM chain or universal strict-negative cost theorem is asserted. Mandatory R093 source-scope precision: its d7 CM/level determinant and torsion theorem is conditional and supplies no nonempty full pinned Host. Literal Params6-22/32 and DataHost54 give r\textgreater=7,t\textgreater=13,d\textgreater365; PinHost39-41 fixes the same d.~Thus the d7 M\^{}3 specialization cannot serve an admitted full-source family. General-d LIFT093 is the admitted-source field-cost route; actual Host/Sizes, N and selected arithmetic content remain open. Full checked addendum: recovery\_pending/lift/r093\_full\_host\_qd\_big\_scope\_precision.md. This clarifies scope and adds no error or movement credit. AM-DO-117 FC1-25 removes the entire original conditional J56 all-top H133 fixed-S nonunipotent support27 sector, at every global rank and incidence profile. The five complete containing-set exponents are -5531348963/4096, -3078643005/4096, -2767418421/1024, -9563431915/2048 and -120294901/4096; their union is below2\^{}-29365. The new small scalar remains in the full paired I\_Sigma, while its exact original near-point transport uses only the universal RT2 selected Hilbert trace bound. Append29365 once to the complete R094120/31 predecessor:121/31, both full sums\textless2\^{}-230, complete floors27/11 unchanged. This is a strictly smaller conditional support27 remainder, with no full support27 exhaustion, actual admitted Host/Sizes, selected arithmetic content, prescribed global period, specific032 pairing, evaluatedDP44 or RH conclusion. AM-K-112 UE1-14 proves that every active type-II place in the exact full-positive globally rank-two-unipotent form is eligible, without an eta-denominator bound. The noneligible active subset is empty, with zero error charge. In the actual all-II scalar-cross support27 remainder, v=s=27, x\textgreater2\^{}2111; x\textless2\^{}3107 forces total eligibility27; x\textless2\^{}12000 gives total eligibility at most36 and at most one eta per fixed support. Fixed-eta active labels contribute at most(T-1)\^{}-27. The full168-factor bound additionally requires total eligibility27 and no inactive paired-C-zero place. Original paired row and central difference ideals retain their support27 denominator power54. Current DO117121/31 exact errors and complete floors27/11 remain unchanged; no full support27 estimate, actual source/content/period, specific032, DP44 or RH conclusion is asserted. AM-LIFT-095 IH1-6 proves an integral hyperbolic pair inside every stated selfdual relative-unramified indefinite Hermitian lattice of rank at least3, by local integral pairs and strong approximation on SU. The original nonterminal source W\_j.L0 satisfies these hypotheses conditionally on the admitted full record, and the same E-linear operators preserve its actual AA multiple. Choosing auxiliary doubled horizontal and integral imaginary root bases gives c=1 in LIFT094: t=1+3M Q\_S, delta=t\^{}(4fr\^{}2), L=t; any separately required larger congruence modulus is retained. This removes the unevaluated auxiliary clearing integer without a replacement lattice, seed translation or extra isogeny. The actual very ample line H, same-rule bad-locus degree D, nonempty family intersection, component-entry multiplier N0, selected nonzero rule and full Host/Sizes remain unresolved inputs. No evaluated total N, period, specific032, DP44 or RH conclusion; no finite charge. Current DO117121/31 arrays and floors27/11 remain unchanged. AM-DO-118 UB1-29 removes the entire actual conditional J56 fixed-S globally rank-two-unipotent support27 sector, all original profiles and labels. The rank35 high envelope retains the full five-column and same-pair isometry fibre. In the rank34 low branch, the true row metric, full beta norm fibre, affine central cosets and actual selected trace are retained; fixing row and beta before eta-dependent labels permits the one-label bound without global eta uniqueness. The six complete strict dyadic bounds have exponents8792,4491,3928,23649,17510,11371; their union is below2\^{}-3927. Append3927 once to DO117121/31, giving122/31 and both full sums below2\^{}-230; complete floors27/11 remain unchanged. This does not remove the other support27 classes, establish a support28 floor, or solve the signed inequality, actual admitted source, selected arithmetic content, prescribed period, specific032 pairing, evaluatedDP44 or RH. AM-K-113 SC27.1-27 removes seven complete containing groups in the actual conditional J56 all-II global-rank-two-unipotent scalar-cross support27 sector. The33 complete endpoints retain all four transverse mixture terms, every G pattern, all inactive-zero subsets, chosen active-integral positions, coordinate sizes and a separately enumerated infinite eta-ideal tail. Their full union is below2\^{}-304. The necessary residual is b=j0,2\^{}2111\textless x\textless2\^{}3328, either total eligibility27 and100\textless=k\textless=164, or total eligibility28,k164 with x\textgreater2\^{}3107. Append304 once to actual DO118122/31, giving123/31 with both full sums\textless2\^{}-230; complete floors27/11 remain unchanged. No whole support27 exhaustion, actual admitted Host/Sizes, signed inequality, selected arithmetic content, prescribed period, specific032 pairing, evaluatedDP44 or RH conclusion. AM-R-095 FW1-21 removes the entire actual conditional J56 fixed-S globally rank-one-unipotent support27 sector, all original profiles, labelled pairs, double-pole sizes and unmixed G patterns. A uniform weighted eta/content cutoff precedes the complete334-place binary pattern product; the remaining z exponent is165f, both genuine affine coordinates and the whole support/marked/type sum remain. The reduced exponent is-54242909/4096 and the complete witness exponent is-95126187/256; their union is below2\^{}-13000. Append13000 once to actual K113123/31, giving124/31 with both full sums\textless2\^{}-230; complete floors27/11 remain. No whole support27 exhaustion, signed inequality, actual admitted Host/Sizes, selected arithmetic content, prescribed period, specific032 pairing, evaluatedDP44 or RH conclusion. AM-DO-119 NA1-29 proves every active place of the actual full-positive globally rank-two-unipotent form is eligible, including distinct source/target labels and every type. Its complete non-all-II scalar-cross support27 bound retains all original paired rows, full-entry poles, affine central cosets, finite weights, relative norm fibres and the entire denominator-ideal tail. Six exhaustive branches give one412 charge. Append412 once to actual R095124/31, yielding125/31 with both complete sums below2\^{}-230 and floors27/11 unchanged. The dense all-II27 and separate nonzero-short-compression sector remain open; no full support28, actual admitted source/Host/Sizes, selected arithmetic content, prescribed period, specific032 pairing, evaluatedDP44 or RH conclusion. AM-K-114 FP27.1-29 removes the entire full-positive-compression global-rank-three scalar-cross support27 sector with nonzero original S-cross. All34 complete ideal/type cells,2925 padded sets, original full167-covector metric, relative norm fibre, projected U denominator/T108 column and G2 source fibre are retained. The complete eta-first all-ideal tail requires no coordinate sparsity. Literal inversion and23 common paired eligible positions force equal null lines outside four containing sectors, contradicting rank3. Their full union costs1650 once. Append1650 to actual DO119125/31, yielding126/31, both entire sums below2\^{}-230 and complete floors27/11 unchanged. Full positive compression remains a premise; the separate nonzero-short-compression27 and dense all-II27 sectors remain open. No full support28, actual admitted Host/Sizes, selected content, prescribed period, specific032, evaluatedDP44 or RH conclusion. AM-LIFT-096 TJ1-18 proves the same-rule bad-locus degree bound on the actual principal-3M special-unitary family of LIFT085/LIFT094, under its stated nonempty same-rule open-locus hypothesis. For g=fr,d=fpq, total frame weight w and theta weight n, take k=gw+n/2 for the literal even-n Dom; squaring the same rule gives k=2gw+n in general without changing its reduced zero set. A finite proof cover of degree at most e\_g=2\^{}(8g\^{}2+2g+4), first jets of the fourth power of the original symmetric principal line, and exact theta multiplier identities produce scalar cuts. For H=(psi\_bar\^{}\emph{L\_S)\^{}h very ample and V=deg\_H(Ybar), proper pushforward gives D\textless=D\_}=e\_g V sum\_\{j=0\}\^{}d k\^{}j, including mixed-dimensional components. When k=0 and the bad locus is proper it is empty. With LIFT095 c=1, LIFT094 then gives the conditional within-family multiplier N\textless=t\^{}\{B\_\{t\^{}\{4fr\^{}2\},fpq-1\}(D\_\emph{)\}, t=1+3M Q\_S; any independent component-entry multiplier N0 is retained. The proof cover contributes only to D\_}, not to the physical isogeny or output field. The numerical h,V, nonempty actual-family intersection, controlled N0, selected arithmetic nonzero rule and admitted full Host/Sizes remain open. No evaluated total N, six-identity chain, actual period, specific032 pairing, DP44 sign or RH conclusion; no finite error charge. Current K114126/31 exact arrays and complete floors27/11 are unchanged. AM-DO-120 QW1-29 removes the entire original doubly-plane scalar-cross sector with nonzero four-entry short compression at every support27..31, including all global characteristic types, actual coordinate supports and paired labels. The original integral positive F witness q\_j=\textbar\textbar(A-I)y\_j\textbar\textbar\^{}2 has every conjugate below2\^{}42. The complete paired finite low-rank proof restores the rank-four Green baseline666 with coefficients331/2 and335/2; both original full source/target costs, selected trace, arbitrary-ideal row, relative norm fibre and all incidence weights remain. All five full endpoint bounds give one270871 charge. Appending it once to K114126/31 yields127/31, with both entire sums below2\^{}-230 and complete floors27/11 unchanged. The complementary scalar-cross remainder at supports27..31 has full positive compressionI, a positive166 fixed space and global rank at most3 by the support-free original integral scalar proof; global unipotence is not inferred. The dense all-II27 window remains open; no complete support28, actual admitted Host/Sizes, selected arithmetic content, prescribed period, specific032 value, evaluatedDP44 or RH conclusion. AM-LIFT-097 VA1-15 evaluates the previously unspecified very-ampleness exponent on the same actual principal-3M arithmetic compactification X. For d=fpq,g=fr and A=psi\_bar\^{}\emph{L\_S, the exact integer h\_d=2\^{}(d+1)(d+1)(d+2)! makes H=A\^{}h\_d very ample. The no-PGL2 boundary condition and canonical-line extension identify K\_X as Cartier ample; the neat Baily-Borel pair is log canonical, A-K\_X is nef, and both applications of effective freeness/very ampleness retain log bigness on every lc/slc stratum. The original Hodge restriction and Hecke equal-pullback identities are unchanged. Writing v\_A=int\_X c1(A)\^{}d, the full ambient degree is h\_d\^{}d v\_A, so LIFT096 gives D\_}=2\^{}(8g\^{}2+2g+4)h\_d\^{}d v\_A sum\_\{j=0\}\^{}d k\^{}j, with k=gw+n/2 for even n or2gw+n after the same-zero square. The independent quantity v\_A is a positive unevaluated integer. LIFT094/095 then retain N\_within\textless=t\^{}\{B\_\{t\^{}(4fr\^{}2),d-1\}(D\_\emph{)\},t=1+3M Q\_S, with every protected congruence condition and any separate component-entry N0. The same-family nonempty open, selected arithmetic nonzero rule and admitted Host/Sizes remain open. No evaluated actual N, new physical field/level/isogeny, six-identity chain, period, specific032 pairing, DP44 or RH conclusion; no finite error charge. Current DO120127/31 arrays and complete floors27/11 are unchanged. AM-R-097 NW1-9 exhaustively composes the checked structural cuts at support27. Every remaining actual support27 arrow is full-positive, scalar-cross, globally rank-two-unipotent with all27 active places typeII, b=j0 and one of two necessary alternatives: total eligibility27,100\textless=k\textless=164,2\^{}2111\textless x\textless2\^{}3328; or total eligibility28,k164,2\^{}3107\textless x\textless2\^{}3328. DO120 supplies full positive compression before K114 is applied. The signed energy retains the actual R27 subset of this containing window and R\_ge28, with \textbar epsilon\textbar\textless2\^{}-230. No set equality, existence, emptiness, sign or floor28 conclusion. No new charge; actual DO120127/31 ordered arrays and complete floors27/11 are unchanged. Actual mixed-CM source, selected arithmetic content, six identities, prescribed period, specific032 pairing, evaluatedDP44, Kummer-Hecke RH and global oscillator comparison remain open. AM-K-115 GM1-23 proves the cap-dependent shape83/rate7/2 positive Gamma majorant, retains its entire low-weight interval by a paid positive shift, and controls every inverse moment0..k+3. All scalar and transverse terms, original saturated metric, central affine coset, source reconstruction,168 T factors and four-state finite patterns remain. The complete ten cells with total eligibility27,b=j0,2\^{}2111\textless x\textless2\^{}3328,k100..109 have total absolute contribution below2\^{}-945. The worst full exponent is-2907383/3072\textless-946, every other is below-3166. One945 append changes DO120127/31 to128/31, both entire sums remain below2\^{}-230 and complete floors27/11 stay. Combined with R097, the remaining actual support27 set is contained in the full-positive scalar-cross rank-two-unipotent all-II class with either total eligibility27,110\textless=k\textless=164,2\^{}2111\textless x\textless2\^{}3328, or eligibility28,k164,2\^{}3107\textless x\textless2\^{}3328. The next positive sufficient envelope is not an actual lower bound. No floor28, signed nonvanishing, admitted source, selected arithmetic content, prescribed period, specific032 value, evaluatedDP44 or RH conclusion. AM-LIFT-098 AV1-20 evaluates the previously unevaluated automorphic volume for the same actual principal-3M special-unitary arithmetic family. Put n=3M, r=p+q\textgreater=3, f={[}F:Q{]}, g=fr, d=fpq, a=r\^{}2-1, a\_r=product\_\{j=1\}\^{}\{r-1\} j!/(2pi)\^{}(j+1), b\_pq=product\_\{i=0\}\^{}\{p-1\}i!/(q+i)!, and e\_v=product\_\{j=2\}\^{}r(1-epsilon\_v\^{}j q\_v\^{}(-j))\^{}(-1), epsilon\_v=1 at split and -1 at inert finite places. The literal relative-unramified CM and ordinary selfdual-lattice assumptions give hyperspecial models even at dyadic places. The original invertible ideal multiple changes no stabilizer or congruence index. Exact strong approximation and finite orders give mu\_Gamma=D\_F\^{}(a/2)a\_r\^{}f n\^{}(fa) product\_\{v not dividing n\}e\_v with compact SU mass one. The physical Hodge representation S+Q\^{}dual gives lambda\^{}r=K\_Y\^{}2 and A\^{}r=K\_X\^{}(2(g+1)); Mumford proportionality yields v\_A={[}2(g+1){]}\^{}d d! b\_pq\^{}f mu\_Gamma. An explicit integer upper bound is Vcal=n\^{}(fa)D\_F\^{}ceil(a/2)2\^{}(f(r-1)){[}2(g+1){]}\^{}d d!(product\_\{j=1\}\^{}\{r-1\}j!)\^{}f.~Thus D\_arith=2\^{}(8g\^{}2+2g+4)h\_d\^{}d Vcal sum\_\{j=0\}\^{}d k\^{}j, h\_d=2\^{}(d+1)(d+1)(d+2)!, retains the full same-rule cost, and N\_within\textless=t\^{}\{B\_\{t\^{}(4fr\^{}2),d-1\}(D\_arith)\}, t=1+3M Q\_S. The independent volume input has been replaced by actual field, rank, signature and level data. Any component-entry N0, same-family nonempty open, selected arithmetic nonzero rule and admitted Host/Sizes remain open. No evaluated actual N, six-identity mixed-CM chain, period, specific032 pairing, DP44 difference or RH conclusion. No finite error charge; current K115128/31 arrays and complete floors27/11 unchanged. AM-K-116 RE1-28 retains the actual reciprocal-energy factor in a shifted positive half-Gamma convolution. The infinite total mass is explicit: every positive inverse moment is bounded by a convergent Gamma/Beta formula, and the unique polynomial degree0 term is made integrable by reserving theta=2\^{}-20 of the actual product-energy exponential before lattice enlargement. Every scalar/transverse term, original saturated metric and covolume, full167-entry reconstruction, central affine750 cover, all168 T factors and all paired four-state G patterns remain. The complete total-eligibility27,b=j\_int=0,2\^{}2111\textless x\textless2\^{}3328,k110 cell has full exponent-1539841/3072\textless-501. One501 append changes actual K115128/31 to129/31, retaining every old overlapping charge; both full sums remain below2\^{}-230 and complete floors27/11 stay. With R097/K115 the remaining actual support27 set is contained in the full-positive scalar-cross globally rank-two-unipotent all-II class with either total eligibility27,111\textless=k\textless=164,2\^{}2111\textless x\textless2\^{}3328 or total eligibility28,k164,2\^{}3107\textless x\textless2\^{}3328. The k111 positive sufficient envelope5269759/3072 is not a lower bound or impossibility theorem. No complete support28, signed nonvanishing, admitted source, selected arithmetic content, prescribed period, specific032 value, evaluated DP44 or RH conclusion. LATEST SOURCE-POINT EFFECTIVITY REDUCTION AM-LIFT-099: for a supplied nonzero literal rule and admitted indefinite source/CM seed, component entry is now bounded by C\_ent=m(2m)\^{}(2f)D\_F\^{}2, m=3M Q\_S, and at least one of its bounded-lattice regular families meets that same-rule open. With all LIFT095-098 factors retained, N\textless=C\_ent(1+m)\^{}\{B\_\{(1+m)\^{}(4fr\^{}2),fpq-1\}(D\_})\}, D\_\emph{=e\_g h\_(fpq)\^{}(fpq) Vcal sum\_(j=0)\^{}(fpq)k\^{}j, k=gw+n\_F/2 or2gw+n\_F after squaring. The original K0(A0{[}N\^{}2{]}) carries all final objects and nonzero frame; the actual same-J/C0 corollary has degree\textless64512d\_CM M\^{}4N\^{}(81/8), rd\textless2\^{}50exp(2pi sqrt(d\_CM))rd(Bcoef)MN\^{}2 and retains {[}Bcoef:Q{]} in the absolute-discriminant cost. The component-entry multiplier is included in this full N, without a ray-class field or proof-cover extension. Independent geometric c,h,volume,degree,nonempty-family and entry costs are removed under these stated source/rule hypotheses. Actual selected arithmetic nonzero, full Host/Sizes, remaining source-family/model/section/polynomial/fibre Step5 inputs, Read/specific032, actual g\_j,e\_j,H and evaluated DP44 remain open. No charge or finite-floor change. AM-DO-121 EC1-24 retains the exact original unipotent identity cosh r\textgreater=1+alpha/2, hence sinh squared r\textgreater=alpha(1+alpha/4), in the genuine selected coefficient. Its positive infinite-mass Gamma167/rate4 convolution has every required convergent inverse moment; original product-energy reservation repairs the sole degree0 term before enlargement. The complete actual total-eligibility27/support27 all-II globally rank-two-unipotent scalar-cross union k111..139 has absolute normalized contribution below2\^{}-993. Every scalar/transverse polynomial term, original saturated metric and paired ideal, full167-entry reconstruction, unique eta, central affine750 cover, all168 finite factors and correlated four-state patterns is paid. One993 append changes actual K116129/31 to130/31, with both entire sums\textless2\^{}-230, all old overlapping errors retained and floors27/11 unchanged. The remaining necessary support27 alternatives are total eligibility27,140\textless=k\textless=164,2\^{}2111\textless x\textless2\^{}3328, or total eligibility28,k164,2\^{}3107\textless x\textless2\^{}3328; b=j\_int=0 throughout. The next k140 sufficient envelope767061/1024 is positive and is not a lower bound or impossibility theorem. No complete support28, signed nonvanishing, admitted source, selected arithmetic content, prescribed period, specific032 pairing value, evaluated DP44 saving or RH consequence. AM-K-117 D29.1-27 proves the universal degree29 scalar Gaussian lattice bound from the fixed exact normalized Gegenbauer polynomial and, with AM-R-098 WJ1-28, removes the complete actual support27 total-eligibility28 branch k164,b=j\_int0,2\^{}3107\textless x\textless2\^{}3328 with absolute normalized contribution below2\^{}-278. The entire660-term polynomial includes the repaired zero degree and the unique degree167 term with its full Q\^{}167 pattern cost. Every original167-entry reconstruction, paired ideal, selected metric, central affine750 cover, unique eta, support and all finite factors is paid; the one eligible inactive place keeps contraction1. One278 append preserves all130/31 predecessor entries, giving131/31 with both entire sums below2\^{}-230 and floors27/11 unchanged. Only the eligibility27 alternative140\textless=k\textless=164,2\^{}2111\textless x\textless2\^{}3328,b=j\_int0 remains at actual support27. No complete support28, signed nonvanishing, actual admitted source, selected arithmetic nonzero rule, prescribed period, evaluated DP44 cost or RH conclusion. AM-DO-122 CE1-29 removes the complete actual support27 total-eligibility27 all-II globally rank-two-unipotent scalar-cross k140 sector with absolute contribution below2\^{}-359. MANDATORY CUTOFF QUALIFICATION: the prose after CE8 is read as norm(h\_tau)\textless=H=133, as proved by R048-RC RC48.4/RC48.6 and the paragraph after RC48.7; CE8 already has the correct non-strict sign. The frozen proof is read together with recovery\_pending/definite/DO122\_CUTOFF\_SOURCE\_ADDENDUM.md, whose full text is included in the checked paper. The joint physical central variable w=-U Im d, original affine ideal, full750 cover, weighted Voronoi integral and U\^{}f physical Jacobian are retained. The central integral consumes one profile power, giving shape166; all moments used and every positive polynomial term are paid, with original-row zero repair, all finite factors and correlated patterns. One359 append to actual K117131/31 gives132/31; both entire sums remain below2\^{}-230 and floors27/11 stay unchanged. The sole support27 necessary alternative is now total eligibility27,141\textless=k\textless=164,2\^{}2111\textless x\textless2\^{}3328,b=j\_int0. The positive k141 envelope is only a failed sufficient bound. No full support28, signed nonvanishing, actual admitted source, selected arithmetic content, prescribed period, evaluated DP44 or RH claim. CONDITIONAL LITERAL SOURCE-LAW REDUCTION AM-LIFT-100: given an actual algebraic frame, the original odd ideal half-level and principal rigidified line, literal IsNab and ThValTot, the independent ThNe0 premise follows at the exact source thresholdDF+2\textless=M. It supplies one individual permitted tuple for every unchanged nonzero harmonic polynomial overK(A{[}M\^{}2{]}), via the actual ideal-half orbit and quotient jet interpolation, with no principalization or further scalar-root field. Actual source/frame existence, IsNab/ThValTot, every other required source law and full Host/Sizes remain; the restricted coupled selected image is not identified with this full thGen family and its nonzero/content/period gate remains. No new arithmetic charge or evaluatedDP44. AM-DO-123 BC1-27 removes the entire original bounded support27/total-eligibility27 all-II full-positive scalar-cross globally rank-two-unipotent k141..142 union, with b=j\_int0 and2\^{}2111\textless x\textless2\^{}3328. The same R099 selected angular kernel, original physical central coordinate w=-U Im d, full affine covolume/U scaling/750 cover, shape166/shift111/100 moments and actual-row zero-degree repair retain every polynomial, finite, label, support, pivot, class/Haar and pattern factor. Full exponents-3344299/1024 and-1688491/1024 plus the entire K118 tail2\^{}-8605114 yield one1647 charge. Appending it once to actual DO122132/31 gives133/31 with both entire sums below2\^{}-230 and complete floors27/11 unchanged. The sole remaining support27 necessary alternative is eligibility27,143\textless=k\textless=164,2\^{}2111\textless x\textless2\^{}3328,b=j\_int0. The k143 envelope-28587/1024 is insufficient and unclaimed. No full support28, signed nonvanishing, source admission, selected arithmetic-image content, prescribed period, specific032 value, evaluatedDP44 or RH conclusion. CONDITIONAL ALL-RING SOURCE-LAW REDUCTION AM-LIFT-101: actual Host/Datum/oddLevel/Coeff/Moduli records and accepted CanonGaugeSt suffice to choose NabV with IsNab, all-ring ThValTot and harmonic NabUniq; thus independent connection/totality premises are removed where the operator may be chosen. With the actual FrPt over a finite number-field extension and DF+2\textless=M, AM-LIFT-100 then supplies literal ThNe0. For a separately prescribed operator its compatibility remains required. Actual source/Host/Sizes, other source laws, original coupled selected-image nonvanishing, period and evaluatedDP44 remain open; no new charge. AM-DO-124 CB1-27 removes the entire original bounded support27/total-eligibility27 all-II full-positive scalar-cross globally rank-two-unipotent k143 sector, with b=j\_int0 and2\^{}2111\textless x\textless2\^{}3328. The same accepted R100 narrow-band angular kernel and K120 actual narrow probability/moments retain the original physical central coordinate w=-U Im d, complete affine covolume/U scaling/750 cover, shape166/lambda32/25 profile, all167 inverse moments and actual-row zero-degree repair. Every original polynomial, finite, label, support, pivot, class/Haar and pattern factor is paid. The full exponent-725931/1024\textless-708 plus the entire new narrow-band tail2\^{}-8605114 gives one707 charge. Appending exactly once to actual DO123133/31 gives134/31 with both entire sums below2\^{}-230 and complete floors27/11 unchanged. The sole remaining support27 necessary alternative is total eligibility27,144\textless=k\textless=164,2\^{}2111\textless x\textless2\^{}3328,b=j\_int0. The k144 sufficient envelope924757/1024\textgreater0 is explicitly unclaimed. No full support28, signed nonvanishing, source admission, selected arithmetic-image content, prescribed period, specific032 value, evaluatedDP44 or RH conclusion. PRESCRIBED-FRAME AND ACTUAL-PACKAGE REDUCTION AM-LIFT-102: coherent exponential gauges supply a literal IsNab with exact ordinary coefficients; actual complex reduced frames have an all-ring pinned extension. For a supplied OMPins/OneMapTieP class package, its existing CovPin and io range already give the operative pin and harmonicity, so its full unchanged harmonic class relation agrees with the reduced-frame evaluation without a separate individual-coefficient equality premise. Actual package/source existence, short-value section laws not supplied, selected arithmetic content/nonzero output, period and DP44 remain open. No numerical charge. CHECKED ANALYTIC INPUT AM-K-121: the exact same-kernel coarea and paired height, cutoff-safe original physical central sum, positive exponential-tangent u-mixture and full actual selected-spread refinement are available. All arithmetic polynomial/moment/pattern/source obligations remain; these analytic factors alone remove no k cell and receive no independent movement credit. The full band tail is retained; no numerical charge. AM-DO-125 CM1--28 removes the complete original bounded support27/total-eligibility27 all-II full-positive scalar-cross globally rank-two-unipotent k144 cell, b=j\_int0 and2\^{}2111\textless x\textless2\^{}3328. Accepted K121 same-kernel coarea/height/variance is globally bounded by a paid exponential tangent before transport, then the exact physical central integral is bounded by a positive Gamma166/rate4, half-Gamma/rate2, half-Gamma/rate3514/29, infinite half-Gamma and shift29/20 convolution. Both complete quadrature tails and8064 intervals for degree147 yield a directed rational moment enclosure; finite-endpoint Hölder covers every lower positive degree and the actual-row reserve repairs degree0. All579 corrections, original supports/pivots/Green/T factors, full750 cover/U scaling and four finite small inflations remain. The full exponent is below−7622655/32768\textless−232.5, and the entire2\^{}-8605114 tail is paid once, giving one232 charge. The exact135/31 full error sums stay below2\^{}-230 with floors27/11. The sole remaining support27 necessary alternative is145\textless=k\textless=164 with the same denominator interval; the positive k145 envelope removes nothing. Signed PP3, higher-support remainder, actual source and selected arithmetic image, prescribed period, specific032 value, actual lift data/evaluatedDP44 and Hecke RH remain open. ALL-RING SHORT-VALUE REDUCTION AM-LIFT-103: with the actual geometric records, selfdual level-trivial Lambda and operative pinned nab2, the full first ShortValTot conjunct is constructed for every coefficient algebra and every original faithfully flat witness, with exact full trace and source normalization. The remaining ShortValTot obligation is the original-base sec/cdelta/IsShortRule descent conjunct; existence of the actual source/package, cusp/spanning laws, selected-image nonzero/content, prescribed period and DP44 remain separate. No numerical charge. COVER-RULE REDUCTION AM-LIFT-104: the cover sec(a,n) rule H2=cdelta}G2 is constructed with its exact positive FrBor character and t\_none\^{}n generator factor. IsLev supplies every field FromTest point needed for its actual evaluation. The remaining second ShortValTot conjunct is precisely an original-base f in sec(a,n) whose evaluation equals H2 for every original FrOver pair. Actual CovPin/LawsFull supplies A=A2=1; bare IsLev is not asserted to do so. Original-base descent, actual source/package, cusp/spanning, selected nonzero/content and DP44 remain separate. No numerical charge. GENUINE SELECTED COHERENCE INPUT AM-K-122: the original coefficient integral may be restricted to simultaneous I* and projective angular coherence\textgreater9/10, with full bounded nonzero-S-cross arithmetic errors2\^{}-8605114+2\^{}(11394885-B/4). The projection bound uses the actual selected metric and all words/twists/compact factors. The remaining high-coherence integral still needs a complete angular/row/central estimate; no cell removal, charge or support-floor change is inferred. AM-DO-126 IQ1--36 removes the complete original bounded support27/total-eligibility27 all-II full-positive scalar-cross globally rank-two-unipotent union145\textless=k\textless=159, b=j\_int0 and2\^{}2111\textless x\textless2\^{}3328. The saturated original isotropic quotient is a projective selfdual O\_E-module, including at2/U; every actual row lies in J I\_Sigma D\_eta times its ordinary row dual. Its standard S164 summand and selfdual rank3 complement give exact transverse covolume D0\^{}3 z\^{}(6f), while the independent physical central U\^{}-f remains. A genuine real-dimension1008 Gaussian bound retains both transverse terms. The unchanged positive central mixture, exact continuous moments148/162, finite-endpoint Holder, repaired zero term, all4590 positive terms, all original finite/source factors and the entire band tail yield one582 charge. Worst directed term is−19510525/32768\textless−595; the entire union is\textless2\^{}-582. Exact136/31 full error sums remain\textless2\^{}-230 with floors27/11. The sole remaining support27 necessary alternative is160\textless=k\textless=164 with the same denominator interval. The positive k160 sufficient envelope removes nothing. Signed PP3, higher-support remainder, actual source and selected arithmetic image, prescribed period, specific032 value, actual lift data/evaluatedDP44 and Hecke RH remain open. STANDARD-LATTICE SOURCE LAW AM-LIFT-105: the actual ShortThetaUnit obligation follows from the original geometric data and operative IsNab, with unit beta0 inverse and the full original normalization for every witness. This removes that named standard-lattice law from the remaining conditional source-construction gate. Nonstandard beta laws, original-base descent, cusp/spanning and selected arithmetic nonzero remain separate. No numerical charge. ORIGINAL-BASE DESCENT AM-R-101: the LIFT104 cover sec rule now descends to an original-base sec(a,n) with every literal IsShortRule evaluation, whenever actual IsLev/newA1/IsNab/PhiPin/PsiTie hold. Both ShortValTot conjuncts follow; no global identification of stored d.r is required. CovPin supplies the first three and dPhi supplies the character predicates for n\textgreater=r. This removes the named original-base descent obligation in that actual supplied scope. Small-n supply, source existence, cusp/spanning, selected arithmetic nonzero and evaluatedDP44 remain separate. AM-DO-127 GC1-16 removes five complete containing subsets of the original support27/eligibility27 full-positive scalar-cross globally rank-two-unipotent all-II branch, with b=j\_int0: k160/161/162/163/164 and log2x\textgreater=2120/2130/2140/2150/2165 respectively, inside the original x\textless2\^{}3328 window. The exact normalized finite Green formula gives2 Q\^{}(-335*167)(NEJ)\^{}167 NFJ0, whose NFJ0 cancels the independent physical central divisor. The complete pattern is product(1+q\_v\^{}(j-167))\^{}2, including exactly2\^{}668 at degree167. The actual projective selfdual quotient, full1008 transverse theta, physical U\^{}-f, same positive profile, all1630 terms and one entire band tail remain. Worst directed term -13968857/32768\textless-426 gives the full union\textless2\^{}-415. Exactly one415 append gives137/31 full sums\textless2\^{}-230 and unchanged floors27/11. The sole support27 necessary branch has2\^{}2111\textless x\textless2\^{}(X\_k) with the five respective upper exponents. Low-x cells, signed PP3, higher supports, actual source/selected arithmetic image, period, specific032 value, evaluatedDP44 and Hecke RH remain open. ALL-N SOURCE LAW AM-LIFT-106: actual AdmX forces n\textgreater=p+q=R.r; below r both ShortValTot conjuncts are vacuous for the unchanged selected DeltaDat. The original restricted dPhi supplies R101 in every remaining n, so unbundled PinDat/dLev/dNab/newA1/dPhi remove shTot as a separate all-n obligation. Upper-range cusp/spanning, actual source/package and selected arithmetic nonzero remain separate. CITED SOURCE SUPPLIER AM-R-102: accepted GlobalTheory.BetaEx.betaExists supplies the literal beta implication for every original datum/n/bad; together with LIFT106, beta and shTot need no separate proof under the seven remaining supplied LawsSh clauses. Actual KPsiIn is still needed for the selected beta witness and is equivalent to accepted same-field ReflexRoots, with finite ReflexBasis sufficient. The eventual coefficient-extension constructor does not by itself identify the prescribed Q.c.F/eO/PinField. No new determinant theorem or source existence is claimed. AM-K-123 HCX1-26 removes the complete original bounded support27/eligibility27 full-positive scalar-cross globally rank-two-unipotent all-II k160 cell, b=j\_int0 and2\^{}2111\textless x\textless2\^{}3328. The actual all-word selected evaluation gives a chi\textgreater937/1000 restriction with both full arithmetic errors paid. A525772-test outward interval cover proves the complete physical-coherence cap; exact three-real-dimensional Beta disintegration gives a fraction below2\^{}-10 with a paid margin. The same-kernel positive four-component central mixture retains the original affine U scaling, full750 cover, Young regularization and extra transport2. Nine full quadratures, finite Holder anchors, actual-row zero repair, the true1008-dimensional transverse theta, exact Green/J0 correlation, all322 terms and both full tails give one600 charge. Exact138/31 full sums remain below2\^{}-230 and floors27/11 remain. The sole support27 necessary alternative is161\textless=k\textless=164,2\^{}2111\textless x, with upper exponents2130/2140/2150/2165 respectively. Signed PP3, higher supports, actual source/selected arithmetic image/content/nonzero, prescribed period, specific032 value, evaluatedDP44 and RH remain open. SAME-FIELD SIGNATURE REDUCTION AM-LIFT-107: under the actual PinDat, all original KPsiIn stages are equivalent to stages1/2 and to the rank-one CM-type ideal-generator law. The exact reduced obligation is one finite normalized Cl(O\_E)-generator table inside the original F; adjacent-stage Bezout formulas retain every norm, original selected betaH and field pin. The table is not yet supplied, and ordinary principalization or a freely enlarged field does not discharge it. QUADRATIC NORM-IMAGE REFINEMENT AM-LIFT-108: the original-field normalized table may be taken on G=image Cl(E)-\textgreater Cl(K), K=Q(H.rt), via Phi(I)=iota(N\_E/K I)O\_F. The exact table is still required inside the original F. A separate explicit root tower supplies a new coefficient record of degree≤\textbar G\textbar≤h\_K with small normalized scalar heights and upper-field persistence, but does not discharge old eO/PinField or changed-field DP44 costs. No claim that actual G is proper; full-source d\textgreater365 excludes the geometric d7 shortcut.

\subsection{Q2 --- OPEN}\label{q2-open}

The actual J14/J56 two-edge mixed-CM kernels and all six m2 analytic predicates are constructed (K031/K032/K038--044). K035--039/K041, R032, DO038 and DO050--055 give the full-Y arithmetic kernel, genuine finite projector, selected recoupling, terminal realization and original-vacuum last-two law. DO055 proves that a nonzero permitted moved all-T Mram bottom period gives a nonzero pairing on the prescribed arithmetic image without a prior-image or integral-vertical-inverse existence premise. The conditional MainX kernel/covolume and terminal-mass scalar is DO055(CI18--20), cP=vT, with the literal inverse-branch functional and original unequal class masses. CURRENT FINITE NONVANISHING GATE (AM-R-094 PP1-24; AM-R-097 NW1-9; AM-K-115 GM1-23; AM-K-116 RE1-28; AM-DO-121 EC1-24; AM-K-117 D29.1-27; AM-DO-122 CE1-29; AM-DO-123 BC1-27; AM-DO-124 CB1-27; AM-DO-125 CM1-28; AM-DO-126 IQ1-36; AM-DO-127 GC1-16; AM-K-123 HCX1-26): At the unchanged charged J56 all-top H133 cutoff, the actual energy is1+Re(sum over actual R27 plus actual R\_ge28)+epsilon, \textbar epsilon\textbar\textless2\^{}-230. R087 exhausts all supports through25 and R091/K109/DO115/K110/R094 exhaust support26. R097 followed by K115/K116/DO121 exhaustively proves that actual R27 is contained in the full-positive scalar-cross globally rank-two-unipotent all-II class with b=j0 and the sole remaining necessary alternative: total eligibility27,161\textless=k\textless=164, with2\^{}2111\textless x and the respective strict upper bounds x\textless2\^{}2130 for k161, x\textless2\^{}2140 for k162, x\textless2\^{}2150 for k163, and x\textless2\^{}2165 for k164. K117 removes the entire former total-eligibility28 branch with one278 charge. No equality with the containing class, existence, sign or absolute-smallness assertion is made. Full positive compression is supplied by DO120 before K114 is applied. All138 chronological all-top and independent31 uniform errors have exact full sum\textless2\^{}-230; current exact arrays are checks/k123\_joint\_remainder\_error.json. The complete all-top floor is27 and the independent uniform H8192 complete floor remains11. Earlier higher-support cuts remain: F,r1 through94; F,r2 through110; F,r3 through154; F ranks4..33 through168; IC8 ranks1/2 through28 and ranks3..31 through the full R075 table; IC8 ranks32/33 through168; both transverse-cross directions through70; rank-two nonunipotent plane-cross through32. DO120 removes all nonzero short compressions in the scalar-cross two-plane sector throughout27..31, with complete positive compression and global rank\textless=3 on its complement. R055 retains its scalar-cross rank-two-unipotent cuts(s13-48,u\textless=10),(s49-84,u\textless=9),(s85-120,u\textless=8),(s121-156,u\textless=7),(s157-168,u\textless=6). All other higher-support restrictions retain their accepted scope. Every paired CM incidence, class mass, phase, finite projector, packet and selected Hilbert normalization is unchanged. K115 removes the complete r27,k100..109 union with one945 charge and K116 the full k110 cell with one501 charge, and DO121 the complete k111..139 union with one993 charge; all smaller and higher-support historical restrictions remain in scope. The sufficient signed inequality PP3 is Re(sum over actual R\_ge27)\textgreater-1+2\^{}-230 and remains unproved. Actual admitted mixed-CM source, selected arithmetic-image content, prescribed global period, specific032 pairing/normalization/value and evaluatedDP44 remain open. Historical support26 windows and earlier structural possibilities below are superseded by the checked complete reductions above. DO122 removes the complete k140 cell with one359 charge using the original weighted central coset and mandatory non-strict-cutoff qualification. DO123 removes complete k141..142 with one1647 charge, including the entire retained-band tail once; its failed k143 envelope gives no further removal. DO124 removes the complete k143 cell with one707 charge, using the actual narrower band, full physical central integral, all167 Gamma512 moments and the entire new band tail once; its positive k144 envelope gives no further removal. DO125 removes the complete k144 cell with one232 charge, using the closed positive physical-central mixture, exact continuous moment enclosure, all lower-degree corrections and the entire tail once; its positive k145 sufficient envelope gives no further removal. DO126 removes the complete k145..159 union with one582 charge using the actual nonprincipal selfdual isotropic quotient, all-prime row containment, full dimension1008 transverse theta factor and independently retained central U\^{}-f; every4590 term and the entire tail are paid. Its positive k160 sufficient envelope removes no further cell. DO127 removes the five complete higher-denominator tails at X2120/2130/2140/2150/2165 with one415 charge, using exact rank-dependent Green/central-J0 cancellation, every finite pattern including degree167 cost2\^{}668, all1630 positive terms and the entire narrow-band tail once. Those higher-denominator cuts alone leave the low-denominator cells; the complete support floor remains27. K123 removes the complete k160 cell with one600 charge, retaining the actual selected high-coherence restriction, full continuum cap, same-kernel positive central mixture, all322 terms and both complete arithmetic tails. All four remaining k161..164 denominator windows and higher-support cuts stay in force; the complete support floor remains27. J14 retains R042: H8192 outside the degree12 plane, primary thresholds only for s\textless=17 and error\textless2\^{}-793; its sign is open. LITERAL SOURCE ADMISSIBILITY: LIFT080 supplies a conditional generic mixed-line height transfer with its explicit unequal-degree cost and arbitrary-type geometric independence/length mechanisms; the effective line-degree, prescribed point, model, polynomial-integrality and source-room premises remain. LIFT079 CT10 also requires replacing the actual pointwise LamPtClauses/Good pure-factor interface for every prescribed mixed fibre with a nonelliptic simple factor; this is conditional on its frame antecedents, and dropping IsTen alone is insufficient. R003 excludes the unchanged elliptic IsTen promotion for powers of J; R023 excludes HostRaw on the fixed J14/J56 fields. LIFT54 excludes full SizesAt with m2 at every eligible FamP point and both fixed degrees at every Scale. LIFT55 excludes the entire unchanged L343 profile even after rebuilding degree, ranks and primes: full SizeG forces L\textgreater2\^{}300 and full sizes force m\textgreater2\^{}7778, f\textgreater2\^{}33598, k\textgreater2\^{}73197. A retained literal Host containing K14/K56 also requires m\textgreater2\^{}11388 by R023/o\_tie/SizeM. These necessary interface obstructions do not prove theta-kernel nonexistence or a DP44 sign. A candidate needs changed field/scale/length or a proved replacement mixed interface. LIFT51--53 retain the full-space Gram/image-rank incompatibility, necessary gap departures and code capacity bounds; they are not a full constructor. LONGER VERTICAL TOWER: LIFT56--58 prove forward cusp-valuedness, every prescribed finite G/Q and basic prior-image injection, real nonzero tests, the first positive-Witt convergent formula with factor2 and every-embedding Rallis nonvanishing. On the genuine smooth multiplicity spaces the unchanged maps satisfy U=D\emph{, ker U=0 and D(C\_m)=C\_0. LIFT59 now proves e0 C\_0!=0 on the actual standard bottom Hermitian model, conditional on its genuine finite/splitting data, actual discrete-series lists and the explicit BC3 numerical bounds implied by literal SizesAt/NCGaps. It constructs an arbitrary-m literal goodTuple with the required aggregate scalar character and an actual Poincare form satisfying \textbar P(1)-\textbar mu(E)\textbar\textbar\textless1/2 at the unchanged principal-G/maximal-away-G level. The prescribed vertical image is consequently nonzero. This is a bottom cusp-space value, not the prescribed Bessel/Mram functional. DO068 now constructs every genuine simultaneous finite Loc/common-vacuum package for the actual BF/FB mixed coordinates, including almost-all basic support and compact-finite pure-tensor witnesses. Its subgroup and bad finite factors are free data. The prescribed Bessel/Mram period, full source/arithmetic ties, arithmetic norm control and a rebuilt comparison baseline remain; the genuine positive-rank free-data implication is now supplied by DO069 in its stated conditional scope. Remaining: a nonzero permitted moved Mram period, quantitative arithmetic compatibility/value, and the specific032 ball-cycle arithmetic identification. The conditional pairing is not the complete requested definite032/MainX result. LIFT60 constructs the exact arbitrary-rho geometric/fractional-ideal Hida flags and finite data on principal-extended-different standard hosts. It gives both-component K2=unipSym, the exact root-character signs and compact diagonal correction, and all finite RowLevelMatchA/B/C under the genuine grp\_levels/grp\_record ties. Its finite index is product\_w q\_w\^{}k product\_i(q\_w\^{}i-1), with original Haar factors retained. The genuine real relative model and full projected analytic row reconstruction are now supplied by LIFT061/062; concrete arithmetic row ties and the actual Bessel/Mram value remain separate. DO069 proves the positive-rank free-data period implication for the genuine R014/R015 realization, with R016/DO068 inputs, including both parity cases and the full uncontracted-to-pointwise comparison. It leaves the initial upper period, arithmetic Frames.Good and norm budget open. LIFT061 constructs the actual positive-rank real Wu/reproduction/decay data, with positive seed energy and its inverse square-root cost unevaluated; LIFT061-U proves local uniqueness up to phase. LIFT062 proves full projected-row reconstruction and cuspidality with actual source continuity, finite projector and Haar. LIFT062-T proves the translated EvalTieT under the literal genuine comparison fields, retaining tS, complex coefficient conjugation, coset volumes and Haar ratio. The concrete J/032 STiesR.rowTie, any required unshifted RowsJoinPinned transport, and actual Bessel/Mram period/value remain open. LIFT065 derives the exact range, nP-scaled norm, level containment and normalized row map from a supplied actual Read plus models\_HE and the genuine translated evaluation. The exact row scale is(c}nP)\^{}(-1/2). With the retained actual HCFinite input and source DataS.hhF, a freely chosen full-HH isometry can satisfy the row identity. LIFT066 now proves complete SIso/Pu/Gterm/column compatibility after the literal reS refill, from a genuine SPack, AutTie, retained LocTied, actual Haar tie, full dimension and normalized row map. Its SIso and row/column pins are identities on the installed genuine data. This is not an isomorphism with arbitrary old operators. STiesR still needs a compatible RecT and the exact forms, and every old comparison link or changed-family DP44 term remains to be supplied for the refilled record. The actual Read/Hida-column/compatible-record construction, prescribed period and032 value remain open. LIFT067 identifies the literal positive-rank T character coefficients(bar(pi),pi), the prescribed standardized D\emph{tau test and exact N-Haar modulus q\footnote{-rho(k+rho+1)}. LIFT068 now evaluates the actual fixed D}tau test conditionally on the stated generic tempered representations and both conductor-one hypotheses, retaining all phases and the Whittaker/Haar factors. LIFT069 now proves the larger generic tempered hypothesis at split T for every constituent of the chosen LIFT059 bottom subspace under its genuine full source package, via a full-degree finite cuspidal place and exact large-theta parameter cancellation. LIFT070 now supplies the smaller generic tempered conductor-one T line for every existing initial Mram constituent under the genuine full SizeData2/SizesAt package, using the actual common non-top label and honest eta parity correction. Thus the prescribed physical D\emph{tau test is locally nonzero against the chosen LIFT059/LIFT069 bottom constituents. LIFT069--070 alone did not extend the chosen-bottom assertion to all e0; LIFT076 below now proves that extension. A fully sized mixed-CM source remains unconstructed. LIFT071 now calibrates the canonical positive-rank normalized physical T factor alpha\^{}natural exactly as q/{[}(q-1)(q\^{}k-1){]} under the specified probability/product Haar, with inverse square-root cost\textless q\^{}(k/2); the literal source normalized factor is c\_H}c\_N times this value, and the unnormalized alpha is Lscript times alpha\^{}natural. Its scalar-one comparison is proved by compact calibration and canonical regularization. LIFT072 now constructs the actual positive-rank split G test inside the same admitted generic finite representation, including mixed types, for every prescribed smaller vector: alpha\_G=1/d\_varsigma, or c\_H\emph{c\_N/d\_varsigma in literal source Haar. Its normalized factor retains Lscript, and its product is\textgreater Q\^{}(-f}k\emph{(k+1)/2). LIFT073 now gives every actual away-G/T spherical factor alpha\_v=c\_H,v}c\_N,v\emph{Lscript\_v and normalized factor c\_H,v}c\_N,v, including all actual dyadic places by a direct split calculation. Thus the prescribed finite local factors are determined under the retained genuine representation/source hypotheses. The global Bessel/Mram period remains open; no source Haar parameter is set to1 and no integral bound for the Hilbert vector is inferred. K089 now proves unit11-content of each actual full-Y auxiliary section before the selected differential and class contractions, with the exact vacuum scalar and finite phases retained. Its finite all-Hasse-jet unit minor does not prove survival of the prescribed selected ordinary derivative/class matrix. The remaining sharp-content question is that exact selected map on these primitive inputs, with the established lower gain removed; no freely chosen auxiliary jet or added T/Y orbit is substituted. K090 now gives the exact integral Mram tensor-augmentation basis on the unchanged free strict/projective T cover, with source class norm2\^{}(f/2) times the coefficient norm and unchanged frame/Haar costs. It identifies the exact selected matrix as the alternating2\^{}f actual columns and proves saturated auxiliary-section injection of that source before subD/class/projector contractions. Actual Sad preserves the full local source lattice at11. The selected matrix/content/nonzero gate remains MA10--11, with every fixed derivative, class weight and genuine projector retained. K091 now supplies a literal integral retraction of both symbolic selected Weyl substitutions and their complete two-edge composition, including the fixed-triple tensor-augmentation auxiliary-section image. Thus universal selected-polynomial substitution adds no Smith factor before the prescribed ordinary-derivative functional. Its abstract integral-lattice reduction is distinguished from nonsaturated ambient polynomial reduction. K092 now supplies a common-unit full-stabilizer cutoff with a primitive symbolic selected theta family and exact contracted full-middle-mass division. No uncontracted section division or selected-derivative survival is asserted. The evaluated selected matrix MA10--11 still requires the actual derivative/gauge, weighted geometric class sums, terminal brC/projectors and their exact content. LIFT074 proves the exact conditional global detector transfer on a specified distinguished embedded pair: \textless F,Pu(T)\textgreater=c\_Per\emph{conjugate(b),c\_Per=c1}mu\_f(K2)\textgreater0,with all prescribed finite tests and the complete unit-Wu contraction. The pairwise global distinction premise and actual arithmetic source/Read/032 matching remain open; b is unevaluated, so no numerical DP44 cost follows. K093 identifies the exact scalar-decoration Fourier/Heisenberg matrix RT13, with class-dependent actual rank-one rational translates and all original phases/masses. Its selected ordinary evaluation necessarily retains at least D\_tau-nu\_tau,1\textgreater0 first-edge derivatives; finite Fourier units do not determine the content of those derivatives or their class sum. The exact evaluated matrix remains open. LIFT075 now excludes the direct compact-source theta cuspidal construction for every permitted positive rank gap: the added GL1 string forces one repeated linear Green factor when d\textgreater=2, and the literal goodTuple degree-one witness kills d=1. This does not kill the coupled map or the transported-leading-r intermediate projector and does not settle the prescribed global Bessel period. K094 now gives an integral full-covector Fourier basis of the actual localized Mram domain and cancels the complete Fourier sum exactly into the signed2\^{}f corner columns FC13. Their span is the actual image/content ideal FC14, retaining full middle mass, a shared P, central phases and positive residual ordinary derivatives. The evaluated corner and middle-class sum remains open; a further source restriction retains its Fourier inverse image. LIFT076 extends the simple generic cohomological tempered-base-change theorem to every existing bottom constituent of the genuine full e0, using squarefree full inertial multiplicities, actual unitarity and the original all-degree divisor witnesses. The actual T spherical representation therefore satisfies the larger-representation tempered/generic premise without selecting LIFT059's special tuple. This remains a classification theorem; source existence, global distinction, arithmetic image/value and DP44 costs are open. DO094 adds the exact finite transverse-scalar implications (rank\textless=3 with both C nonzero, rank\textless=2 with at least one, rank\textless=1 without a C condition), and the original integral delta\_z witness. Its sharp excluded finite cells remain distinct; no numerical charge follows. K095 removes the whole middle T-covector sum and fine-class mass by the exact composite incidence and legitimate contracted coarse-stabilizer cut. The actual initial Fourier column is MF12: a sum over coarse middle classes with the same P and the original ordinary selected evaluation of the quadratic phase psi(u2ZXmu). Auxiliary primitivity precedes that evaluation; selected content/nonzero and coarse-class cancellation remain open. R071 now proves the remaining compact transverse scalar is exactly1, using integral trace/norm and the strict exact bound36U/M²\textless1. Thus the actual support23/24 scalar-cross residual fixes a positive166-dimensional E-space and has rank(h-I)\textless=3; its remaining invariant complement has signature(2,1). This supersedes the weaker positive165/rank\textless=4 structural gate from DO093. No integral moving basis, global nilpotence, full support removal or new numerical charge is inferred. LIFT077 now calibrates the exact prescribed coupling through BPC2025 Theorem1.3.6.1: GC17 retains aS²/(aH\emph{aG}a0\emph{d\_sigma}2\^{}(t\_pi+1)), embedded nT/nF, the partial central Rankin--Selberg ratio, raw G dimensions and once-normalized T factors. It applies also when the period is zero. For the specified existing genuine pair, a nonzero central Rankin--Selberg value is equivalent to nonzero b; neither its nonvanishing nor quantitative size is proved. GC19 retains the literal c\_Per,c0,nT in the source Gram lower bound. Source construction, selected arithmetic content, Read/STiesR/specific032 and DP44 remain open. R072 removes the R071 support23 scalar-cross terms with6\textless=k\textless=164, where k counts nonzero original S scalar coordinates, with one3600 charge. Its fixed packet-plane projection retains the U denominator, exact paired T\^{}(4s) target index and rank-two source fibre; the first row uses the actual joint norm and spherical-cap/Gaussian packing. The remaining support23 row has k\textless=5. No support24 or complete-floor conclusion is inferred. In the actual R071 signature(2,1) residual, BC is skew-adjoint rank one exactly at global rank3; BC=0 is exactly global rank2 with real reciprocal eigenvalue lambda. Every original S entry of BC is an exact G-unit three-minor, has conjugates\textless2\^{}9 and only simple paired typeIII poles. The pointwise Green envelope is2 Q\^{}(-501\emph{167)2\^{}(9}167f)T\^{}(167t), with no C-pattern/witness factor. In rank2, doubled G divisibility of (lambda-1)\^{}2/lambda forces global unipotence for t\textless=22. The existing DO082 removal therefore leaves rank-two support23 only nonunipotent with all23 active places typeIII; at24 its nonunipotent branch needs t\textgreater=23. No new charge or complete-floor change. Actual central value, arithmetic/source and DP44 gates remain open. The actual global rank3 R071 terms at s23,k\textless=63 and s24,k\textless=51 have combined absolute contribution\textless2\^{}-21500. The bare support-ideal count retains all original sparse cross coordinates, three transverse values, pivot, fixed packet-plane U denominator, paired T\^{}4 column index, positive rank-two source fibre and complete pair weights. The actual DO095 witness gives the uniform no-pattern rank-three Green bound. All114 cells are checked. Combined with R072 this removes every rank3 support23 term; DO095 then leaves the support23 scalar-cross remainder rank2 nonunipotent, k\textless=5, t23. At support24 remaining rank3 has k\textgreater=52. Complete floors23/11 and signed/global/source/arithmetic/DP44 gates remain open. K096 proves exact emptiness of the remaining nonunipotent rank-two support23 scalar-cross subcase, without a sparsity condition or a new charge: actual-pole parity leaves eleven possible paired cross poles, incompatible with the original row norm. Together with DO082 and R073 this exhausts the R071 support23 scalar-cross residual. R074 now supplies the complete structural support23 assembly, without a new charge. R074 supersedes all earlier historical complete-floor23 descriptions: the current complete all-top floor is24 and the exact current sufficient inequality is J24.2, with unchanged77/31 errors. The entire globally nonunipotent rank-two R071 support24 sector has absolute contribution\textless2\^{}-2962487. The nonzero reciprocal trace has23 or24 actual poles; principal-norm parity and one exactly normalized inversion leave at most twelve paired cross-pole places, including the possible active eigenvalue-unit place. Primitive original extraction gives normT\^{}-24 without a U denominator. The full original R059 count retains both five-coordinate columns, source fibre, all G patterns, padded subset choices and pair weights; the complete exponent is -60671740183/20480. Exactly one2962487 charge is added, leaving78/31 entries with both full sums below2\^{}-230. Complete floors remain24/11. Globally unipotent rank2, rank3, signed/global/source/arithmetic and DP44 gates remain open. In the original IC8 class at compressed ranks3..31, the full symmetric compression has rank r at every actual G factor. Finite unipotence therefore forces full rank\textgreater=r+1; at equality the unique nontrivial Jordan block fixes every untouched original column at both CM partners. The resulting paired ideal and baseline rank r+2 retain every original DO085 coordinate, volume, target, source fibre and pair factor. All145 stated fresh cells contribute\textless2\^{}-8866, with exact maximum -36351867/4096. The complete rank3 ceiling becomes28, and with the earlier F/IC8 cuts every remaining term at supports24..28 has A=I. One8866 charge gives79/31 entries with both full sums\textless2\^{}-230. Complete floors24/11 remain; signed/global/source/selected-image/032/DP44 gates are open. The entire globally unipotent rank-two R071 support24 sector with at least one type-I or type-III place has absolute contribution\textless2\^{}-71310. The exact block normal form gives N²Z=0 and alpha=Uq\textless=H+H\^{}-1-2, so all original S and fixed transverse inputs have only type-II poles and the full covector radius is sqrt(R/U). At most24 padded23-place ideals cover every non-all-II profile, with all164 coordinate sectors, common G patterns, genuine selected trace, true projected-column U denominator, paired T\^{}4 index and source fibre retained. The complete maximum is -292137655/4096; one71310 charge gives80/31 entries, both full sums\textless2\^{}-230. Complete floors24/11 remain. The all-II unipotent case, rank3 and fixed-block residuals, signed/global/source/arithmetic/032/DP44 gates remain open. The entire global-rank-three R071 support24 cross sector has absolute contribution\textless2\^{}-5710, with no sparsity premise. The full positive167 covector is counted in its actual Gram metric, with exact determinant4U and arbitrary-ideal covering radius; its last scalar is counted by the exact relative norm fibre, including nonprincipal and non-conjugation-stable ideals. Two pointwise Green bounds on the same original G patterns meet at z=2\^{}-100; full low/high exponents are -23390187/4096 and -361072601/4096. All original trace, paired targets, projected U denominator, source fibre, support, pivots and incidence weights are retained. One5710 append gives81/31 complete errors below2\^{}-230 and unchanged floors24/11. Rank-two all-II and fixed-block residuals, signed/global/source/selected-image/032/DP44 gates remain open. The entire all-II globally rank-two-unipotent R071 scalar-cross support24 sector has absolute contribution\textless2\^{}-2879. The actual cross map and original central entry w reconstruct the full matrix. Its affine difference lattice iJ0 I\_F² retains the rank-zero G congruence, real codifferent and full discriminant, giving the complete central count2\^{}479 Q\^{}167 z\^{}-2f. The true transverse Gram volume, exact split-relative norm fibre, every original all-II paired weight and the actual selected trace are retained. All164 coordinate sectors are exhausted by64/99/1 complete monotone ranges with strict maxima -2010899/512,-41042253/4096,-11800269/4096. One2879 append gives82/31 errors below2\^{}-230 and unchanged floors24/11. Together with DO097 and R076/DO096 this removes every R071 scalar-cross support24 residual. The fixed-S rank-two sector, signed/global/source/selected-image/032/DP44 gates remain open. The entire original fixed-S globally rank-two-unipotent support24 sector has absolute contribution\textless2\^{}-6845. The complete R058 high finite-rank envelope extends with exponent -119372763/4096. At low finite rank the original square pairings force the null image into the packet plane; a type-I place is impossible by the unchanged pole-budget46 norm contradiction. All remaining II/III profiles are counted through the true transverse3 dual Gram metric, exact pivot relative norm fibre and original affine central-entry lattice, retaining both CM partners, full discriminant and actual incidence weights. The complete low exponent is -168257413/24576; the high/low union contributes one6845 charge. Every old charge, including DO098, is retained, giving83/31 errors below2\^{}-230 and unchanged complete floors24/11 until a separate assembly. Signed/global/source/selected-image/032/DP44 gates remain open. Every actual all-top support24 term is exhausted by the complete original F/IC8 partition at compressed ranks1..33, every fixed-S nonunipotent/rank-one/rank-two class, and every scalar-cross global-rank/type/coordinate alternative. R075 covers all nonidentity positive compressions; DO090/DO092/R077 exhaust fixed S; R070/DO093/R071 reduce cross terms to the precise DO095 rank2/3 alternatives, exhausted by R076,DO096,DO097,DO098. The complete all-top support floor is25, uniform floor11. All83/31 chronological errors are unchanged and both full sums remain below2\^{}-230. No new charge. The actual remaining signed gate is J25.2; prescribed global/source/selected-image/specific032/DP44 and RH/global oscillator gates remain open. Both complete C-transverse directions, with every incidence profile and no global-unipotence restriction, are removed at supports25..70. At support25 the plane-C/nonzero-original-short-compression sector is also removed. The full joint row, both original columns and complete source fibre give transverse maximum -49537435/4096; the complete plane row, exact relative norm fibre and actual selected trace give -435793141/4096. Their entire union costs one12079 charge. The remaining scalar-cross support25 matrices have both cross blocks in the original packet plane, all four short compressions zero, positive166-dimensional fixed space and global rank at most3 by the support-free norm-one lemma. At26..70 only the two plane conditions are concluded. Actual84/31 ordered errors remain below2\^{}-230 with complete floors25/11. No signed/global/source/selected-image/specific032/DP44 or RH conclusion. The entire original fixed-S globally rank-two-unipotent sector at support25 is removed for every actual incidence profile. Also removed at26 are D\_G\textgreater=33 and the D\_G\textless=32 covector sector r(y1)=r(y2)=0. The original central affine lattice has poles I1/II2/III1; its T\^{}(u-s) factor cancels actual pair weights to T\^{}(-t)\textless=1. The true three-covector short pivots and one-coordinate Schur-complement last pivot retain all source/norm/trace/metric factors. Complete low exponent at25 is -27553337/24576; high endpoints are -77307919/4096 and -35243075/4096; the last-only26 endpoint is -78939175/8192. One1120 charge covers the complete union, giving85/31 errors below2\^{}-230 with unchanged complete floors25/11. Short-pivot low-rank26, fixed-S rank-one25, remaining scalar-cross25, signed/global/source/selected-image/specific032/DP44 and RH gates remain open. The complete global-rank-three scalar-cross support25 sector with the deterministic original common-C ideal satisfying2\^{}32928 N\_EJ\textgreater=T\^{}50 is removed with one29550 charge. The whole original167-coordinate dual-metric ball, relative norm fibre, selected trace, paired projected column, source fibre, all patterns and complete actual incidence weights remain. The full monotone high-ideal endpoint is -121040253/4096. Every remaining rank-three support25 term satisfies2\^{}32928 N\_EJ\textless T\^{}50. The weaker low-ideal extension is a failed sufficient bound4090765885/12288\textgreater0; no actual contribution is inferred. A complete local typeII shear model shows why C-integrality cannot be substituted. Actual86/31 errors remain below2\^{}-230 with complete floors25/11. No whole support25, signed/global/source/selected-image/specific032/DP44 or RH conclusion. The entire remaining globally nonunipotent rank-two scalar-cross interval25..32 is removed after the retained DO099 transverse cut. The support-independent DO076 reciprocal-line theorem gives lambda in F and inverse-C=-lambda\^{}-1 C without requiring full-Z at26+. At least23 actual lambda poles and one global inversion give at least12 positive valuations; all active lambda-unit places are retained in the complete paired C-pole cap s-12. Every original column, source fibre, pattern, pivot, support and labelled-pair weight remains. The complete worst exponent is -2610113463/20480 at32, and eight supports cost one127443 charge. Actual87/31 ordered errors remain below2\^{}-230 with complete floors25/11. Support33 is a failed sufficient bound only. Global-unipotent rank2, rank3, signed/global/source/selected-image/specific032/DP44 and RH gates remain separately open. Every original fixed-S globally rank-one-unipotent support25 term is removed with one537 charge. The primitive q coordinate retains the original affine difference ideal, so both q/r true eta-metric element counts replace the older relative-norm fibre. Independently checked LIFT078 averages every conjugate-content ideal over eta via actual subideal packing and the unconditional actual-field bound zeta\_F(1.001)\textless3003·5\^{}167. Its full constant2\^{}1065 and exponent21021/125 are retained before all Green patterns, selected traces, supports, marked subsets and original paired type/label sums. All26 double-pole cells pass; worst margin538, next1530, complete witness exponent -25183329/64. Actual88/31 errors remain below2\^{}-230 with complete floors25/11. This is a full rank-one sector result, not whole support25 or signed/global/source/selected-image/specific032/DP44/RH nonvanishing. Every globally unipotent rank-two scalar-cross support25 term with at least one type-I or type-III place is removed with one881 charge. The actual full-positive normal form makes all original cross coordinates integral outside typeII. Its original central coordinate has affine difference lattice iJ0 I\_Sigma,F I\_D,F, retaining the shift, full real discriminant and750 transference coefficient. The complete u24 sector costs less than2\^{}-882, including50 type-position choices and both extra T factors. The padded u\textless=23 sector includes300 sets, all3\^{}25 type assignments and T\^{}25 pair cost, with explicitly charged cap30 removal and full sum below2\^{}-174292. Actual89/31 ordered errors remain below2\^{}-230 with complete floors25/11. The all-II support25 sector remains open; no signed/global/source/selected-image/specific032/DP44 or RH conclusion. Every remaining rank-three scalar-cross support25 term now has at least123 nonzero original S coordinates and2\^{}33096 N\_EJ\textless T\^{}50 for the same deterministic common-C ideal. The original(k+3)-coordinate covector ball keeps the true Gramdet4U, arbitrary-ideal cover, exact beta norm fibre, actual selected trace, full paired plane-column U/T index and source fibre. At z\emph{=2\^{}(-197/2), all122 low-ideal cells have maximum -7111295/4096; the complete high-ideal endpoint is -5721469/4096. Both complete unions cost one1395 charge, retaining the old overlapping DO10029550. Actual90/31 errors remain below2\^{}-230, complete floors25/11. Dense smaller-ideal rank3, all-II unipotent rank2 and signed/global/source/selected-image/specific032/DP44/RH gates remain open. R082 TP1--13 restrict every remaining rank-three support25 term, with its actual t=11..25 type-III places, to k\textgreater K\_t and2\^{}(42p\_t)N\_EJ\textless T\^{}50 for the same common-C ideal. In increasing t order, (p\_t,K\_t) are(1120,149),(1096,148),(1073,147),(1049,147),(1026,146),(1002,145),(979,144),(955,142),(932,141),(908,139),(885,137),(861,135),(838,132),(814,129),(791,125). The complete general type/pair weights and actual-t minor poles give maxima-9691203/4096 over2106 low cells and-11857677/4096 over15 high cells, one2353 charge. Actual91/31 errors stay below2\^{}-230; complete floors25/11. The all-II unipotent rank-two branch and surviving dense smaller-ideal rank-three profiles remain; signed/global/source/selected-image/specific032/DP44/RH gates remain open. DO103 removes the complete all-II globally rank-two-unipotent scalar-cross support25 subset with N\_E denominator(eta)\textless=2\^{}120. Both actual CM pole lines force the same selected delta=eta-Q(-7+2sqrt(-7)) divisibility. The globally integral denominator-cleared delta ideal has norm at least7\^{}84 T\^{}p and below{[}4(U+14){]}\^{}168 N\_E D\_eta, so the exact host inequality gives p\textless=24. The shared cross/double-pole containing set, original affine central counts and all164 CF endpoints give25T}2\^{}-2879\textless2\^{}-1878. Actual92/31 ordered errors remain below2\^{}-230, complete floors25/11. Every remaining all-II support25 term has N\_E D\_eta\textgreater2\^{}120. The large-denominator branch, rank-three profiles, signed/global/source/selected-image/specific032/DP44/RH gates remain open. LIFT079 CT1--10 proves that literal CMObj and, independently, both actual LamObjGood copies force the complex target to be isogenous to a power of one CM elliptic curve, using pure complete quadratic eigenbases and the actual surjective/quasi-inverse sum. The actual LamPtClauses fibre identity transfers this constraint to the prescribed point; a point with a nonelliptic simple factor has no such witness. Thus dropping IsTen or bypassing CMObj while retaining Good, or substituting an unrelated elliptic Good object, does not supply the mixed-J point hypothesis. This is conditional on the stated point/frame premises and is not a contradiction to the accepted original source theorem. R003/DP25 simplicity is carried, not counted again. No mixed replacement constructor, finite charge, period/image/value, DP44 or RH conclusion follows. AM-K097 removes the complete actual global-rank-three scalar-cross support25 subset with N\_E D\_eta\textless=2\^{}120, using the cross-only pole lemma, at most24 eligible places, all25 padding choices and all15 actual type profiles. The complete normalized bound is below2\^{}-11658. No rank-three full-double-pole or unipotent central-fibre assertion is imported. One11658 charge appends to the actual DO10392/31 prefix, giving93/31; both full exact sums remain below2\^{}-230 and complete floors remain25/11. Remaining rank-three terms have N\_E D\_eta\textgreater2\^{}120 and retain R082 cuts. The signed kernel, actual selected arithmetic image, prescribed global period, mixed source construction, specific032/Read, changed-family DP44 and RH remain open. AM-R-083 removes the entire actual global-rank-three scalar-cross support25 tail N\_E D\_eta\textgreater=2\^{}4864 with k\textgreater=126, the latter already forced on the current remainder by R082. Its full167-dimensional arbitrary-ideal covector count, exact per-denominator eta enumeration and complete ideal-zeta tail retain all original type weights, G patterns, selected traces, projected U/T column and source fibre. The complete exponent is -22330263/4096\textless-5451, charged conservatively once as5400. The actual K09793/31 prefix becomes94/31; full sums remain below2\^{}-230 and floors25/11 are unchanged. Together with K097, remaining rank-three support25 terms satisfy2\^{}120\textless N\_E D\_eta\textless2\^{}4864 and all R082 profile cuts. The all-II unipotent rank-two, signed kernel, selected arithmetic image, mixed source, prescribed global period, specific032/Read, DP44 and RH gates remain separate and open. AM-DO-104 removes the complete actual all-II globally rank-two-unipotent scalar-cross support25 tail N\_E D\_eta\textgreater=2\^{}1820. Every original coordinate lies in LD\_eta; the full original4-by4 block and both mixed blocks give the exact affine central difference ideal iJ0 B\_eta I\_Sigma,F\^{}2 with D\_eta bar(D\_eta)=B\_eta O\_E. The full real750 cover, all167 covector values, actual selected trace, all164 row cells, original weights and the entire ideal-zeta tail are retained. The complete maximum -10346711/4096\textless-2526 yields one2518 charge. Actual R08394/31 becomes95/31, complete sums remain below2\^{}-230 and floors25/11 remain. Together with DO103 the remaining all-II rank-two branch has2\^{}120\textless N\_E D\_eta\textless2\^{}1820. No whole support25, signed-kernel, prescribed selected arithmetic image/period, mixed source, specific032/Read, DP44 or RH conclusion. LIFT080 MT1--17 proves a generic Step5 transfer for actual mixed-type comparison lines with individually bounded normalized degrees d\_j\textgreater=d\_-. Rescaling only their archimedean norms by exp(d\_j-e) makes all auxiliary degrees e, keeps every finite ideal and local polynomial quantifier unchanged, and replaces La=L by L exp((e-d\_-)\emph{+). The exact cost is C(e)=(1/2+\textbar e\textbar+log(L)/2)(D+gn)+D log(rL)+(3D+gn)(e-d}-)\emph{+/2. Its optimum within this comparison family is C0+(-d}-)\emph{+ min(D+gn,(3D+gn)/2). Arbitrary-type isogeny derivatives and the literal arbitrary-set form\_letter\_le supply independence and norm mechanisms; explicit mixed products J\^{}(14h p) with complementary-action factors supply geometric inclusions/projections with N=L=1 and permutation U. The actual prescribed point/model/section, all-extension polynomial tests, effective eigenline degrees, coefficient-field/frame ties and fully revised source/DP44 room inequality remain required. The old LamPtClauses/LamObjGood pure-type boundary remains; no main-family constructor, finite charge, period/image/value, specific032/Read or RH conclusion follows. AM-R-084 exhausts every remaining global-rank-three scalar-cross support25 term. The complete K097 count applies to \textless=24 eligible places without a denominator premise; the exact DO093 inversion carries this count and R083 to the inverse null line with the same original S-supportk\textgreater=126, actual weights, masses, phases and cutoff. Outside the four charged sets both slopes are fully eligible and each denominator norm is\textless2\^{}4864. Their denominator-cleared difference has the full paired T\^{}50 divisor but norm\textless(4U)\^{}168\emph{2\^{}9728\textless T\^{}50, so the two lines coincide. DO095 then forces global rank2, a contradiction. The entire four-sector absolute union is\textless2\^{}-5398, with both inverse costs included. Actual DO10495/31 becomes96/31; every old overlapping charge remains, full sums are\textless2\^{}-230, and complete floors25/11 remain unchanged. At support25 only the separately retained all-II globally rank-two-unipotent scalar-cross branch remains, with2\^{}120\textless N\_E D\_eta\textless2\^{}1820. Higher supports and signed-kernel, selected arithmetic-image, prescribed global period, actual mixed source, specific032/Read, DP44 and RH gates remain open. LATEST EFFECTIVE DEGREE REFINEMENT: LIFT081 ED1--23 supplies effective intrinsic line-degree lower bounds for every actual comparison component eigenline on J14/J56, with the full integral CM action and the literal source metric n\^{}2=pi H. The good model is obtained over K0(alpha,beta), alpha\^{}7=-4 and beta\^{}2=(zeta7-1)\^{}7, of degree at most14 over K0. The actual integral vector 2q\_h p\_tau e\_i has norm below sqrt(2)q\_h at every counted embedding, with q1=406 and q4=26390. Therefore ndeg\textgreater-log(sqrt(2)q\_h), and the weaker common bounds are -10log(2) for J14 and -16log(2) for J56. These normalized degrees are invariant under further finite base change of the actual good model. Substitution into LIFT080 makes its sufficient conditional Step5 cost fully explicit: Copt=1/2(1+log L)(D+V)+D log(rL)+b\_h min(D+V,(3D+V)/2), V=gn and b\_h=10log(2) or16log(2). This resolves the effective comparison-line-degree premise only. The prescribed mixed point, model/level/coefficient-field ties, nonzero frame polynomial, all residue-extension Step5Loc conditions, source-room data, selected arithmetic image/content, global period and DP44 evaluation remain required. No claim is made about an arbitrary line in a larger multiplicity space, equal degrees, a primitive coordinate-ideal equality, an old LamPtClauses witness, a finite charge, or RH. AM-DO-105 removes the entire actual all-II globally rank-two-unipotent scalar-cross support25 subset with at most24 eligible places, without any denominator premise, together with the fully eligible middle window2\^{}120\textless N\_E D\_eta\textless2\^{}1820 having1\textless=k\textless=21 nonzero original S coordinates. For each fixed support there is at most one fully eligible eta in this window and at most one actual diagonal pencil label per active place, including central-only poles. The original weighted label sum is at most(T-1)\^{}-25. The joint S/transverse beta integral, moving central cover and simultaneous x/z endpoint retain all full metric/phase/class/Haar factors. The worst claimed cell -1852891/4096\textless-452 and full fresh union give one451 charge. Actual R08496/31 becomes97/31, both complete sums remain below2\^{}-230 and complete floors25/11 remain. After R084, DO103/104 and this theorem the only support25 branch left is fully eligible all-II rank-two unipotent, with2\^{}120\textless N\_E D\_eta\textless2\^{}1820 and22\textless=k\textless=164, unique eta per support and at most one label per active place. The dense range, all higher-support cuts, signed inequality, actual mixed source, selected arithmetic image/content, prescribed global period, specific032/Read, DP44 and RH gates remain. LATEST ARBITRARY-TYPE POLYNOMIAL REFINEMENT: LIFT082 FT1--29 transfers one actual nonzero frame polynomial to every Step5Loc prime and complete-DVR extension, and every Step5Arch embedding, from explicit arbitrary-type geometric inputs. Generic FrPt polynomial existence and same-point nonzero evaluation produce one G before any completion. Integral eigenmodule bases commute with arbitrary compatible DVR extension by flatness; no integral eigenspace splitting is assumed. The exact relation Ub\^{}t Y=X, intrinsic theta generator unit, determinant inclusion and physical row coefficient-clearing indices give all local clauses. The original pi-normalized Hermitian metric, actual fibre/section maps and goodOf conjugation give all archimedean clauses, retaining lieFac and the dual theta norm. There is no new transfer cost beyond LIFT080/081. The actual prescribed point with Dw.mult, good model/level/Key/section/coefficient-field ties, valSet arithmetic-lattice identification, same-point nonzero F value, Host/Sizes, Read/STiesR, selected arithmetic image/content, global period, specific032 and DP44 remain required. No old CMPt.Good/LamObjGood/LamPtClauses instance, finite charge, support-floor change or RH assertion follows. LATEST ACTUAL J-ISOTYPIC SIGNATURE RESTRICTION: LIFT083 JS1--25 proves the actual pinned signature obstruction for every J-isotypic complex fibre, independently of CMObj/IsTen/LamObjGood. For any quadratic Host action into End0(J\^{}N)=M\_N(Q(zeta7)), d=7 gives Lie multiplicities(2a+b,a+2b), a+b=N, and every other Host quadratic field gives balanced(3N/2,3N/2). Actual DetCond requires(f\_E p,f\_E q). Hence d=7 needs p\textless=2q and q\textless=2p, while d!=7 needs p=q. The literal pinned profiles p\_j=k+j−1,q\_j=m−j are unbalanced at every stage and terminal q\_m=0; under full SizesAt, k\textgreater m\^{}5 f\_R\textgreater=2m forces p\_j\textgreater2q\_j at every stage. Thus no solely J-isotypic point can satisfy any pinned stage under those full premises, even with enlarged field/changed primes and the old good-point wrappers bypassed. An actual enlargement A\textasciitilde J\^{}N×B must have N\textless=f\_E min(p,q) and dim B\textgreater=f\_E max(p−2q,q−2p,0); the terminal fibre contains no J factor. These are necessary geometric dimensions, not height or DP44 costs. Other CM isogeny classes, their actual full point/action/model/level/section and same-point nonzero value remain open, together with selected arithmetic content, prescribed global period, Read/STiesR/specific032 and DP44. No finite charge, support-floor change or RH conclusion follows. AM-K099 GE1--13 restricts the actual remaining fully eligible support25 branch to total eligible count q=25 or26. The norm bounds give x\textgreater2\^{}120 for q25, x\textgreater2\^{}1116 for q26, and x\textgreater2\^{}2111 for q27, excluding q27 under x\textless2\^{}1820. Over K0=Q(sqrt(-7)), the relative Galois group of order168 acts regularly on the selected T primes. At q25 the scalar has trivial relative stabilizer and generates E over K0; at q26 the relative stabilizer has order at most2. The absolute degrees are at least168 and84 respectively. Distinct actual norm-U scalars in the bounded denominator window have eligible-set overlap at most14; set and scalar stabilizers agree. A relative stabilizer of order2 forces the denominator ideal to extend from the fixed field and x to be a square, even at ramified primes. This gives necessary structural restrictions only, no numerical sector removal, new error charge, source construction, signed positivity, period, DP44 or RH claim. The existing97/31 error arrays and complete floors25/11 remain unchanged. AM-K100 QE1--16 removes the complete actual fully eligible all-II globally rank-two-unipotent scalar-cross support25 subset with2\^{}1116\textless=x=N\_E D\_eta\textless2\^{}1820 and22\textless=k\textless=48, including either total eligible count25 or26. The exact1116 endpoint, cap1117, complete joint row/transverse beta count, central volume\textless2\^{}679 and unique-eta/label weight give worst k48 ceiling-1730011/4096\textless-422; the26 smaller cells and this cell form one complete2\^{}-421 union. All k\textless=21 contributions remain in DO105451 and are not added again here. Append421 to the actual DO10597/31 ordered totals, giving98/31; both entire sums remain below2\^{}-230 with complete floors25/11 unchanged. In particular all q26 cells up through k48 are removed. The surviving support25 possibilities have2\^{}120\textless x\textless2\^{}1116 with22\textless=k\textless=164, or2\^{}1116\textless=x\textless2\^{}1820 with49\textless=k\textless=164, in addition to K099 structural restrictions. No claim for these surviving cells, no actual source, selected arithmetic image, prescribed period, signed positivity, specific032, DP44 or RH conclusion follows. AM-DO-106 CFI1--24 proves the complete actual central fibre Gamma=B\_eta I\_Sigma,F squared, covolume sqrt(D\_F) x T\^{}-50, sufficient period Q T squared B\_eta of index Q\^{}168 T\^{}386, and the full-real codifferent dual. The original active type-II factor has Fourier l1 norm2/T for integral cross rows and1/(T-1) for a fixed cross pole. The actual mixed G projectors and every inactive-T factor are retained by the finite positive-projection trace lemma. The center-independent whole-C Green budget gives total finite Fourier l1 at most2 Q\^{}(-335}167)(N\_E J\_C)\^{}(335/2)(2/T)\^{}j(T-1)\^{}(-(25-j)), with no finite-index loss. The complete orbit lies in the original finite support; regularized Poisson retains its exact signs and covolume. Extending a disjoint collection past H133 uses at most the single complete2\^{}-2000 tail once, subject to reconciliation of earlier removals. No numerical alias estimate, sector removal, new charge or floor change follows; current98/31 errors remain. Actual mixed source, selected arithmetic content, prescribed period, signed inequality, specific032, DP44 and RH remain open. AM-R-085 RF1--31 gives the exact genuine selected vector-valued central coefficient c(t)=P(t)/(z0-epsilon i gamma t)\^{}s, with s=kappa+K+2N and degree P\textless=s-kappa, using the original Pieri component and its exact Hilbert/PBW norm. The finite coefficient rule, partial fractions, signed gamma Fourier transform and affine Poisson phases are explicit. The positive central density for the same selected component has mass one and pointwise dominates the transformed commuting fibre, without assuming its projection commutes. Products over all168 labels lie in the actual signed orthant and vanish on its boundary faces. Ordinary absolutely convergent Poisson holds for the actual lattice and finite characters supplied separately by DO106, with no quotient-cardinality loss. This removes the unspecified analytic-model step but proves no numerical dual-lattice alias bound, finite sector removal or new error charge. Current98/31 totals and floors25/11 remain; the signed inequality, actual source, selected arithmetic image, prescribed period, specific032, DP44 and RH remain open. AM-K101 UC1--17 removes the complete actual fully eligible all-II globally rank-two-unipotent scalar-cross support25 subset with2\^{}1585\textless=x=N\_E D\_eta\textless2\^{}1820 for every22\textless=k\textless=164. It uses the original DO105 all-factor bound with a single common denominator endpoint, exact capmax(161,k), full selected trace and moving affine central cover. The143-cell union is below2\^{}-254, with nondense maximum-1412571/4096\textless-344 and dense maximum-1048027/4096\textless-255. Append254 once to actual K10098/31, giving99/31 and both whole sums below2\^{}-230; all earlier overlaps,451 and421 are retained and complete floors25/11 remain. The only still-unremoved support25 cells therefore have2\^{}120\textless x\textless2\^{}1585,22\textless=k\textless=164, with K100 additionally excluding22\textless=k\textless=48 when x\textgreater=2\^{}1116, and all K099 structural restrictions retained. No prospective inactive-T gain, signed positivity, actual source, selected arithmetic content, prescribed period, specific032, DP44 or RH result is inferred. AM-DO-107 IP1--23 removes three complete subsets of the actual fully eligible all-II globally rank-two-unipotent scalar-cross support25 middle sector: any inactive paired-C-zero place for all k22..164; no paired zero with x\textgreater=2\^{}744 for all k22..164; and no paired zero with x\textless2\^{}744 and22\textless=k\textless=99. The literal included-Lambda projector gives rho=1/(T-1) at every noneligible inactive place with a nonzero paired C block. A paired zero instead supplies an actual extra normT² row ideal, giving T\^{}-335 while retaining the old central affine cover. The complete143/143/78-cell union is below2\^{}-392. Append392 to the actual K10199/31 arrays, giving100/31; all historical overlaps including451,421,254 remain, both full sums are below2\^{}-230, and complete floors25/11 are unchanged. The only support25 branch left has no inactive paired C zero, total eligibility exactly25,2\^{}120\textless N\_E D\_eta\textless2\^{}744 and100\textless=k\textless=164, with K099 field/stabilizer restrictions retained. This is a strict smaller finite gate, not support25 exhaustion, signed positivity, an admissible mixed source, selected arithmetic image, prescribed global period, specific032, DP44 or RH. AM-DO-108 AC1--15 removes the entire actual support25 fully eligible all-II globally rank-two-unipotent scalar-cross subset with any active paired-C-integral place, at cost less than2\^{}-55281. The exact extra normT² row ideal gives T\^{}-335; the central cover uses the old scale. Its143 complete cells append55281 to actual DO107100/31, giving101/31, all historical charges retained and both sums below2\^{}-230, floors25/11 unchanged. Remaining middle data have j=b=0, exactly25 total eligible places,2\^{}120\textless x\textless2\^{}744 and100\textless=k\textless=164. Every complete actual central orbit then keeps all168 original T factors equal to1/(T-1), with finite weight(T-1)\^{}-168 a(h(t)), sufficient periodQGamma, indexQ\^{}168, exact trace dual and ordinary Poisson formula AC15. This strictly removes active-integral profiles and removes their indicator from the exact Fourier interface; it gives no alias estimate or automatic short-sum/fibre replacement. Admissible source, selected arithmetic content, signed nonvanishing, global period, specific032, DP44 and RH remain open. AM-LIFT-084 MP1--25 constructs a literal field-valued Obj and FrPt for any given actual source Host, compatible Par.PinHost/PinLevel/PinDat, specified stage Datum, existing Moduli representation and same Coeff field, after a finite extension. The actual principal ideal identity is derived from Host fields. A pure principally polarized CM block and its conjugate action supply every uniform(p,q); the explicit trace-one semicharacter supplies actual SAB prim/adjoint/cubical lines. Relative-unramified selfdual Hermitian lattice uniqueness, including dyadic places, gives the required profinite Tate isometry; the literal commutator sign and CRT preserve the exact finite level M=g0g1,A=1. All-a normalized line identities descend over one finite field; only geometric full Tate systems are asserted. Separately, on the existing E\_h fields the explicit pure block and three J-action twists give J\^{}(fq) times C\_h\^{}(p-2q) for p\textgreater=2q, attaining LIFT083 complement dimension f(p-2q). Those fixed fields still fail full Sizes. No effective field/frame/line costs, good-model/old LamPt identification, prescribed same-point F nonzero, actual full Host/Sizes constructor, selected arithmetic content, prescribed period, specific032, DP44 or RH is supplied. No numerical error charge;101/31 arrays and floors25/11 remain. AM-R-086 ZJ1--21 constructs principal source points with actual J content over any fixed actual Host E under the retained LIFT084 input hypotheses. The auxiliary compositum L=E Q(zeta7), c={[}L:E{]} in\{1,2,3,6\}, has an explicit integral CM trace polarization of degree\textbar D\_L\textbar7\^{}(cf/3). Zarhin Theorem1.1 gives a full integer-ring compatible principal block. Its exact conjugate-action balancing, with the even-c improvement, yields block sizes a\_c=(8,8,24,24). For p\textgreater=q the resulting exact J exponent is N=(2f/3)(q-r), r=q mod a\_c, and pure complementary dimension f(p-q+2r). If d=7 and p\textgreater=2q, a principal(16,8) block gives N=f(q-r7), r7=q mod8, complement dimension f(p-2q+3r7). Zero remainders attain the corresponding LIFT083 bounds. The unchanged Datum, finite level and Coeff give actual Obj/FrPt over a finite extension. A separate full O\_K7 lattice-approximation argument gives an underlying isogeny J\^{}N to the actual J factor prime to any prescribed finite set, and an isomorphism of p-divisible group schemes after good reduction. No old theta-line/frame compatibility is inferred. Effective field/isogeny/line/frame costs, numerical Host/Sizes existence, same-point F nonzero, actual lift identities/arithmetic image, global period, specific032 and DP44 remain open; no RH or finite numerical charge. Current101/31 arrays and complete floors25/11 are unchanged. AM-LIFT-085 HF1--26 proves that, for an existing literal Host/Datum/Moduli/Coeff, actual CM seed Obj, original finite level A=1 and a given nonzero natural frame rule, a point in that same OE-isogeny class has a nonzero literal FrPt evaluation at the unchanged finite level. The isogeny u has degree N\^{}(2g), N prime to any prescribed finite set S containing the level primes, N=1 mod M, and u\emph{L1=L0\^{}(N\^{}2) rigidified and cubical. Every output object and a nonzero frame descend over K1=K0(A0{[}N\^{}2{]}), with {[}K1:K0{]}\textless=N\^{}(8g\^{}2). The argument meets every analytic moduli component through an algebraic family based on its own target lattice. A weight-rule frame can be chosen by integer matrices with entries0..D+1 relative to any initial rational eigenframe. N itself, seed degree, initial-frame/theta-generator heights, models and actual line costs remain unevaluated. The given rule must itself be nonzero; no prescribed arithmetic-image nonvanishing is proved. No old pure LamObjGood or selected derivative identification, global period, specific032 pairing, DP44 or RH conclusion follows. This removes the qualitative same-rule same-point existence gate within the fixed CM isogeny class under the precise assumptions. Current101/31 error arrays and complete floors25/11 remain unchanged. AM-K-102 PC1--27 bounds the entire actual DO107 IP3 support25 remainder by2\^{}-3550: all-II globally rank-two-unipotent scalar-cross, exactly25 total eligible places, no inactive paired cross zero,2\^{}120\textless N\_E D\_eta\textless2\^{}744 and100\textless=k\textless=164. The relative-norm phase count, positive-cone Gaussian/Poisson majorant and true transverse metric use one joint energy constraint; the complete original complex-volume tail covers every larger ideal-norm parameter. The original25 active weighted-label factors and143 inactive pointwise factors are retained without assuming complete-orbit constancy. All195 endpoint bounds are below-3557; the65-cell union costs3550 once. Actual DO108101/31 becomes102/31, every old overlapping charge retained, both complete sums below2\^{}-230 and floors25/11 until a separate exhaustive assembly. No signed positivity, prescribed period, selected arithmetic-image witness, effective admitted source cost, specific032, DP44 or RH claim follows. AM-R-087 J26.1--5 exhausts every original all-top support25 class at the unchanged H133 cutoff. The complete F/IC8 rank1..33 ranges force A=I; DO090/R080/R079 exhaust fixed S; DO099/DO095 give the precise cross rank2/3 alternatives; R082/R084 exhaust rank3; DO101 and DO102 remove nonunipotent rank2 and all non-all-II unipotent profiles; the exact DO103/104/105/107/K102 denominator, eligibility and coordinate partition exhausts the final all-II branch, including every equality boundary. The complete all-top floor is26 and the independent uniform floor remains11. Every102/31 chronological error entry is unchanged, both full sums stay below2\^{}-230 and no assembly charge is added. The actual sufficient signed gate is J26.2; global period, fully sized admissible source, selected arithmetic image, specific032/Read, actual DP44, Kummer-Hecke RH and the global oscillator comparison remain open. AM-R-088 GZ1--11 removes the actual fixed-S rank-one-unipotent support26 integral-witness branch, every reduced double-pole cell d\textless=22, and every reduced pattern with at least one nonidentity unmixed G reduction. The exact averaged-content and original affine-coordinate count retains every paired type/label factor and the actual selected trace. The full334-place nonempty weighted G-pattern sum is below2\^{}-5533, replacing only the former factor2. All27 complete endpoints and the witness give one523 charge. Only the reduced all334Gidentity pattern with d23..26 remains in this branch. Actual R087102/31 becomes103/31, every old overlap retained, both whole sums below2\^{}-230, and complete floors26/11 unchanged. No full support26 exhaustion, signed positivity, actual source, selected arithmetic content, prescribed period, specific032/Read, DP44 or RH conclusion follows. AM-LIFT-086 IM1--27 constructs actual everywhere-good polarized models for the LIFT085 geometric CM seed and target over K=K0(A0{[}15N\^{}2{]}), with {[}K:K0{]}\textless=(15N\^{}2)\^{}(4g\^{}2). The full integral OE action, principal polarization, canonical trace-one cubical line, isogeny and adjoint extend over the same ring. Actual intrinsic tangent-line degrees change by an amount in{[}-log N,log N{]}, the full determinant degree by at most g log N, and the theta degree by at most g log N/2 under KeyCore. Target intrinsic lines have length constant1. If the actual seed lines have degree at least-b and all genuine rule/integral/cusp/source premises hold, the sufficient Step5 cost is C\_N=(D+gn)/2+D log r+(b+log N)min(D+gn,(3D+gn)/2). No effective bound on N, initial field or arbitrary-Host b is supplied. Host/Sizes, actual rule nonzero/content, prescribed period, specific032/Read, evaluated DP44 and RH remain open. No new finite error charge; actual103/31 arrays and complete floors26/11 remain unchanged. AM-K-103 SC26.1--28 removes nine complete subsets of the actual all-II globally rank-two-unipotent scalar-cross support26 class. Its support-free full-Z proof, new eta window and uniqueness, exact paired ideals and inactive weights, full infinite ideal tail, joint energy and positive-norm-cone counts retain every original factor. All1,229 cells give one293 charge. The residual necessarily has v=q26, 2\^{}1116\textless N\_E D\_eta\textless2\^{}2108, b=j0 and81\textless=k\textless=164. Actual R088103/31 becomes104/31, all earlier overlapping charges retained, both entire sums below2\^{}-230 and complete floors26/11 unchanged. Other types/classes, signed positivity, actual source, selected arithmetic content, prescribed period, specific032/Read, DP44 and RH remain open. AM-DO-109 SP26.1--25 removes the entire actual fixed-S globally rank-two-unipotent support26 low-D\_G\textless=32 short-pivot class. The original plane reduction and paired alpha budget give2i+t\textless=5. The actual V-input eligibility and central-only label proof, true Gram metric, affine central difference lattice and full weighted four-state G sum are retained. The three complete p\textless=24,p25,p26 endpoints give one820 charge; the p26 branch proves eta uniqueness before summing actual labels. Together with already charged R079 high-D\_G and last-only pivots, this completes the whole fixed-S rank-two-unipotent support26 structural class. Actual K103104/31 becomes105/31, every overlap retained, both entire sums below2\^{}-230 and complete floors26/11 unchanged. Other support26 classes, signed positivity, actual source, selected arithmetic content, prescribed period, specific032/Read, DP44 and RH remain open. AM-K-104 FV1--14 removes the complete new coordinate interval81\textless=k\textless=95 from the actual K103 all-II globally rank-two-unipotent scalar-cross support26 residual. The original full joint-volume lattice, actual central affine fibre and selected trace are retained. The full334-place four-state weighted G product is below2; the exact support binomial is below2\^{}101. Complete coupled-x monotonicity gives the actual lower denominator endpoint without a missing range. All15 cells, with worst-1128627/4096, give one274 charge. The all-II residual now has v=q26, 2\^{}1116\textless x\textless2\^{}2108, b=j0 and96\textless=k\textless=164. Actual DO109105/31 becomes106/31, all prior overlaps retained, both entire error sums below2\^{}-230 and complete floors26/11 unchanged. Other classes, signed positivity, actual source, selected arithmetic content, prescribed period, specific032/Read, DP44 and RH remain open. AM-R-089 WN1--18 removes the complete actual fixed-S rank-one-unipotent support26 remainder: all334Gidentity and intrinsic double-pole sizes23..26. Its universal reciprocal-norm weighted eta/content lemma uses determinant-one rectangular rescaling and full dyadic shell packing, including every unit and ideal class. The original selected trace retains its actual eta norm before summation; both genuine affine coordinates and the whole actual support/marked/type/pair sum remain. The complete rational upper exponent is-146697349/4096\textless-35000, giving one35000 charge. Together with already charged R088 this completes the whole fixed-S rank-one-unipotent support26 structural class. Actual K104106/31 becomes107/31, every old overlapping error retained, both entire sums below2\^{}-230 and complete floors26/11 unchanged. Other support26 classes, signed positivity, actual source, selected arithmetic content, prescribed period, specific032/Read, DP44 and RH remain open. AM-DO-110 NT26.1--31 removes the complete actual non-all-II globally rank-two-unipotent scalar-cross support26 sector. The support-independent full-positive normal form, original paired cross/full-entry poles, 24-place eta uniqueness, actual one-label cancellation including central-only poles, full affine central lattice, common-G ideal and exact joint count are retained. Seven exhaustive branches and all1148 cells give one561 charge. Actual R089107/31 becomes108/31, every overlap retained, both complete sums below2\^{}-230 and full floors26/11 unchanged. The remaining all-II unipotent26 class is the K103/K104 dense window; other scalar-cross classes, signed positivity, actual source, arithmetic content, prescribed period, specific032/Read, DP44 and RH remain open. AM-LIFT-087 GS1--20 chooses an actual new general-Host Serre seed X=O\_L tensor\_\{O\_K7\} J and the actual pure block b\^{}-1 tensor\_\{O\_k\} C0, preserving the Host ideal, Datum, level, coefficient premises and R086 signature/J-content formulas. Every constructed intrinsic component eigenline has degree at least-b, with b=max(log D\_L+.5 log(2t),(sqrt(d)+1)/3+log D\_E+.5 log f). Literal integral b gives1 in b\^{}-1; full-order discriminant-cleared projectors and the exact Zarhin pullback isometry retain all integral models and source pi H norms. The new trace-polarization degree is D\_L/7\^{}(5t), distinct from R086. Under all retained actual rule, KeyCore, integrality/cusp/source premises, LIFT086 now has explicit seed b in C\_N=(D+gn)/2+D log r+(b+log N)min(D+gn,(3D+gn)/2). Effective N, initial field, actual Host/Sizes constructor, selected image, prescribed period, specific032/Read, DP44 and RH remain open. No finite charge; current108/31 arrays and complete floors26/11 unchanged. AM-R-090 IS1--25 removes the entire actual full-positive globally rank-three scalar-cross support26 class with nonzero original S-cross, without a sparsity hypothesis. The complete eligible-set count at most24 costs6670; the complete all-k denominator tail at least2\^{}8192 costs11965. Literal inversion applies each containing-set bound once more. Outside them, the two eligible sets intersect in at least24 paired places, forcing equal null lines by the exact6395-bit denominator-cleared norm gap and contradicting global rank three. The full four-sector union is below2\^{}-6668. Actual DO110108/31 becomes109/31, every previous overlapping charge retained, both complete sums below2\^{}-230 and floors26/11 unchanged. The nonzero short-compression and dense all-II rank-two classes, signed positivity, actual source, selected arithmetic content, prescribed period, specific032/Read, DP44 and RH retain their gates. AM-DO-111 VC1--25 proves the actual positive transverse compression is unitary when both original scalar-cross directions lie in the retained plane. Its four original integral short witnesses have the sharper all-embedding bound2\^{}41. The complete simultaneous finite restoration, arbitrary-ideal row volume, relative norm fibre, original selected trace, full W0 column, source fibre and all paired type weights give exact exponent-87517717/4096\textless-21366 for the entire nonzero-short-compression support26 sector. The complementary integral norm argument gives full positive167 compressionI and a positive166 fixed space, hence global rank at most3; global unipotence is not inferred. Actual R090109/31 becomes110/31 with one21366 append, every old overlapping charge retained, both entire error sums below2\^{}-230 and complete floors26/11 unchanged. Whole support26 assembly remains separate; dense all-II rank-two, signed positivity, actual admissible source, selected arithmetic content, prescribed period, specific032/Read, DP44 and RH remain open. AM-LIFT-088 SF1--16 supplies an explicit field for the entire newly chosen LIFT087 general-Host seed over the same given Coeff field B: {[}K\_seed:B{]}\textless=168}23040\emph{h(-d)}M\^{}(4g\^{}2)\textless=3870720\emph{d}M\^{}(4g\^{}2), g=f(p+q). The actual prescribed elliptic CM action descends over B(j) by scalar differential injectivity. The explicit J good-model radicals, elliptic15-torsion, actual projective Serre modules, equivariant principalization, canonical trace-one cubical line, original finite M-level and literal Obj/FrPt all lie over the stated field. Under every retained LIFT085 nonzero-rule/source hypothesis, {[}K\_out:B{]} is at most that bound timesN\^{}(8g\^{}2), with the existing intrinsic Step5 cost unchanged. EffectiveN, absolute degree of the givenB, field discriminants, coefficient/section heights, numerical Host/Sizes existence, selected arithmetic content, prescribed period, specific032/Read, evaluated DP44 and RH remain open. No finite error charge; current110/31 arrays and complete floors26/11 are unchanged. AM-R-091 DR1--7 consolidates every checked support26 cut through DO111 into one actual-remainder theorem. The only possible support26 terms are all-type-II globally rank-two-unipotent scalar-cross matrices with full positive167 compression, nonzero original S cross, v=q26, b=j0, strict denominator window2\^{}1116\textless N\_E D\_eta\textless2\^{}2108 and96\textless=k\textless=164. The actual retained support26 set and actual higher-support remainder split the unchanged energy exactly, with the complete110/31 errors below2\^{}-230 and floors26/11. Every earlier cut and overlap charge is retained. No charge or independent new movement is assigned to this consolidation. DR5, actual fully sized source, selected arithmetic content, prescribed period, specific032/Read, evaluated DP44, Kummer-Hecke RH and global oscillator comparison remain open. AM-K-105 CG1--25 removes all39 complete coordinate sectors96\textless=k\textless=134 from the actual all-II globally rank-two-unipotent support26 scalar-cross remainder. The scalar real-dimension336 Poisson bound, full coupled Gram determinant, original selected trace, full affine central covolume and every finite/source factor remain. Restoring(N\_FJ0)\^{}(-1) and its actual common-G correlation gives the full weighted pattern factor\textless2Q\^{}(-167); complete coupled-x monotonicity gives the exact lower denominator endpoint. The worst cell is-9009839/6144\textless-1466 and the preceding38 are below-2831, so one1465 charge covers the union. Actual DO111110/31 becomes111/31, every overlap retained and both complete sums below2\^{}-230; floors26/11 stay fixed. With R091, the only possible support26 branch now has135\textless=k\textless=164, the same v=q26,b=j0 and strict denominator window2\^{}1116\textless x\textless2\^{}2108. Signed positivity, actual source, selected arithmetic content, period, specific032/Read, DP44 and RH remain open. AM-LIFT-089 HF89.1--18 sharpens the degree of exactly the same LIFT088 seed and output fields. For r=p+q,g=fr, the actual projective OE-lattice is free modulo each finite ideal, and the full defined action bounds its torsion Galois image. The literal m/bar(m) level halves are perfect Weil duals; Coeff.zeta\_mem already gives mu\_M inB. Thus Kgood(A{[}M{]})=Kgood(A{[}m{]}) and {[}Kseed:Kgood{]}\textless=M\^{}(f r\^{}2). The same target field Kseed(A{[}N\^{}2{]}) has relative degree\textless=N\^{}(4f r\^{}2), including ramifiedN andN1. Consequently {[}Kout:B{]}\textless=3870720\emph{d}M\^{}(f r\^{}2)\emph{N\^{}(4f r\^{}2), under every actual LIFT085 nonzero-rule/source hypothesis. All actions, lines, level, Obj/FrPt and intrinsic costs are unchanged. EffectiveN, absoluteBdegree, discriminants, coefficient/frame heights, actual Host/Sizes existence, selected arithmetic content, period, specific032/Read, DP44 and RH remain open. No error charge; current111/31 arrays and complete floors26/11 are unchanged. AM-DO-112 EC1--22 removes both complete coordinate sectors k135 and k136 from the actual all-II globally rank-two-unipotent support26 scalar-cross remainder. The larger Gaussian cap59/(50}65536), with the full real dimension336 tail and scalar theta bound26, preserves the original row metric, affine central divisor, selected trace, every finite factor and complete coupled pattern/denominator monotonicities. Exact whole-sector exponents are-691859209/307200\textless-2252 and-283283209/307200\textless-922; their union is below2\^{}-921. Actual K105111/31 becomes112/31, every historical overlap retained and both complete sums below2\^{}-230; floors26/11 stay fixed. The only possible support26 branch now has137\textless=k\textless=164, v=q26,b=j0 and strict denominator window2\^{}1116\textless x\textless2\^{}2108. The positive k137 diagnostic is unclaimed. Signed positivity, actual source, selected arithmetic content, period, specific032/Read, actual DP44 and RH remain open. AM-LIFT-090 SG1--21 sharply bounds the same actual LIFT087 seeds and LIFT088 fields through the two original generating CM varieties J and prescribed C0. Integral projective split maps and their duals put every seed n-torsion in Kgood(J{[}n{]},C0{[}n{]}) at every n, with no prime-to-index restriction; full maximal-order rank-one actions bound that field by n\^{}8. Thus the same seed and output fields cost M\^{}8 and N\^{}16, and a finite list of stages sharing J,C0,Kgood costs one M\^{}8 and one lcm(N\_j)\^{}16. The same C0{[}15{]} field has degree\textless=225 before good reduction, sharpening the same ground factor to37800\emph{h(-d)\textless=37800}d.~The complete relative bound is37800\emph{d}M\^{}8\emph{lcm(N\_j)\^{}16. Retain the smaller LIFT089 bounds where applicable. The actual models, actions, polarizations, canonical line, finite level, Obj/FrPt and intrinsic cost remain unchanged. EffectiveN, absoluteBdegree, discriminants, coefficient heights, all nonzero-rule/source premises, selected arithmetic content, prescribed period, specific032/Read, actual DP44 and RH remain open. No error charge; current112/31 arrays and complete floors26/11 remain. AM-K-106 NT1--14 removes the complete coordinate sector k137 from the actual all-II globally rank-two-unipotent support26 scalar-cross remainder. The exact universal radial-profile factor carries through the original rank-two near-point trace and yields630 bits. Every nonzero original S coordinate lies in the same fractional ideal; the proved product-of-sums inequality gives k\^{}-84 and596 further bits at137. Applied before the full DO112 Gaussian cardinality enlargement, these constant factors preserve every actual metric, finite weight, central affine covolume and complete pattern/denominator count. The exact whole-sector exponent is-251641609/307200\textless-819. Actual DO112112/31 becomes113/31, all previous overlapping charges retained, both complete sums below2\^{}-230 and floors26/11 unchanged. The only possible support26 branch now has138\textless=k\textless=164 with the same v=q26,b=j0 and2\^{}1116\textless x\textless2\^{}2108. Signed positivity, selected arithmetic content, prescribed period, full effective source, specific032/Read, actual DP44 and RH remain open. AM-LIFT-091 RD1--18 bounds root discriminants of the same actual LIFT088--090 finite torsion fields. For a finite Galois extension of actual degree D, unramified outside an integral ideal S, rd(L)/rd(K)\textless=D}Norm(rad S)\^{}(1/{[}K:Q{]}); the complete base-normalized local different and actual-degree wild contribution are proved. Thus the simultaneous seed satisfies rd(KSigma)\textless=rd(Kgood)\emph{rad(M)}M\^{}8, and the same actual output satisfies rd(KSigma,out)\textless=rd(Kgood)\emph{rad(M L\_N)}lcm(M,L\_N²)\^{}8\textless=rd(Kgood)\emph{M\^{}9}L\_N\^{}17. Its relative increase over KSigma is at most rad(L\_N)\emph{L\_N\^{}16. The literal N\_j=1 mod M congruence is retained, so no false level/output-overlap saving is claimed. Full all-object descent and the37800}d\emph{M\^{}8}L\_N\^{}16 relative degree bound are unchanged. Initial rd(Kgood), absolute B degree/discriminant, effective N\_j, heights, all nonzero-rule/source conditions, selected content, prescribed period, specific032/Read and evaluated DP44 remain open. No numerical error charge;113/31 full arrays and floors26/11 remain. AM-R-092 ID1--18 bounds the initial root discriminant for the exact SF5 seed C0 and radical/torsion tower of LIFT088, over the actual source coefficient field B. The maximal-order class-field step B(j)/B is unramified; the actual radical tower has relative root-discriminant factor at most76832. Every CM conjugate satisfies \textbar sigma(j)\textbar\textless32 exp(pi sqrt(d)), and the actual C0{[}15{]} extension is treated at its true degree before applying the bound225. Consequently rd(Kgood)\textless2\^{}50 exp(2 pi sqrt(d)) rd(B). Composed with LIFT091, rd(KSigma,out)\textless2\^{}50 exp(2 pi sqrt(d)) rd(B) rad(M L\_N) lcm(M,L\_N\^{}2)\^{}8\textless=2\^{}50 exp(2 pi sqrt(d)) rd(B) M\^{}9 L\_N\^{}17. The stated log absolute discriminant bound retains the exact relative degree37800\emph{d}M\^{}8\emph{L\_N\^{}16 and actual {[}B:Q{]},rd(B). This removes the separate initial-good-field discriminant premise for this fixed seed, not for arbitrary twists. Absolute B degree/discriminant, effective N\_j, heights, Host/Sizes existence, selected nonzero arithmetic rule/content, period, specific032/Read and evaluated DP44 remain open. No new finite-error charge;113/31 full arrays and complete floors26/11 unchanged. AM-LIFT-092 CD1--19 strengthens the root-discriminant cost of the same actual LIFT088--090 fields. For the two full-maximal-order CM progenitors J,C0 over the actual everywhere-good Kgood, every actual F inside Kgood(J{[}n{]},C0{[}n{]}) is abelian and has conductor dividing n O\_Kgood and discriminant dividing (n O\_Kgood)\^{}(D-1), D={[}F:Kgood{]}. Thus rd(F)/rd(Kgood)\textless=n\^{}(1-1/D)\textless=n, including all dyadic/CM-ramified primes and D=1. The same simultaneous output satisfies rd(KSigma,out)\textless=rd(Kgood)}lcm(M,L\_N\^{}2)=rd(Kgood)\emph{M}L\_N\^{}2 under the actual N\_j=1 mod M; its relative output increase over KSigma is at most L\_N\^{}2. Combining the separate canonical R092 bound gives rd(KSigma,out)\textless2\^{}50 exp(2 pi sqrt(d))\emph{rd(B)}M\emph{L\_N\^{}2. LIFT090 relative degree37800}d\emph{M\^{}8}L\_N\^{}16 is unchanged; the sufficient absolute-discriminant logarithm is this degree times {[}B:Q{]} times (log rd(B)+50 log2+2 pi sqrt(d)+log M+2 log L\_N). All original fields, varieties, polarizations, line, level, quotient, action, source points and frames are retained. Actual B degree/discriminant, effective N\_j, metric/height costs, Host/Sizes, selected nonzero arithmetic content, period, specific032/Read and evaluated DP44 remain open. No finite-error charge;113/31 full arrays and complete floors26/11 unchanged. AM-K-107 SV1--36 removes the entire k138 cell of the actual all-II globally rank-two-unipotent scalar-cross support26 remainder. The genuine selected density yields a full tail-inclusive squared coefficient below2\^{}-933 and a single actual near-point product below2. Together with the exact radial factor, all-coordinate norm and original saturated transverse lattice, the complete count is bounded by2\^{}(-50428781/204800)\textless2\^{}-246. The original scalar theta is in real dimension336, with full21\^{}(k+3) and all original central, support, row, finite and pattern factors. The remaining branch has139\textless=k\textless=164 and2\^{}(71443/64)\textless x\textless2\^{}2108, v=q26,b=j0. Exactly one246 append gives114/31 ordered errors with all old overlaps retained, both full sums below2\^{}-230, unique minimum231 and next minimum246; complete floors26/11 remain. Mandatory corrected boundary conventions: SV6 is strict for h\textgreater=1 and equality at h0; SV11 quotient is for r\textgreater0 and the bound at r0 is1. The frozen proof is read with the explicit boundary addendum reproduced below. Signed positivity, selected arithmetic content, period, effective source, specific032/Read, actual DP44 and RH remain open. AM-DO-114 SP1-27 proves a uniform scalar theta below3/2 in real dimension336 for A\textless=1/43500 by a complete finite spherical-cap and nonzero-shell argument. Applied to the original saturated transverse Gram, actual selected trace, full affine central fibre and every finite pattern, the whole139..141 union is below2\^{}-354. The necessary actual support26 all-II rank-two-unipotent range becomes142\textless=k\textless=164 with2\^{}(71443/64)\textless x\textless2\^{}2108,v=q26,b=j0. Exactly one354 append gives115/31 complete ordered errors, both full sums below2\^{}-230, with every historical overlap retained and complete floors26/11 unchanged. No support27, signed positivity, selected arithmetic content, period, effective source, specific032/Read, evaluated DP44 or RH conclusion. AM-K-108 AE1-30 retains the full actual Cartesian conditional angular density, cancels its D ratio in detailed balance and inserts exp(-2U\textbar\textbar r\textbar\textbar\^{}2) in the original lattice sum. With the separately checked DO114 scalar theta, original transverse Gram, central affine fibre, common ideal correlation and all finite factors, the complete142..159 union is below2\^{}-585. The necessary support26 all-II rank-two-unipotent remainder becomes160\textless=k\textless=164, with2\^{}(71443/64)\textless x\textless2\^{}2108,v=q26,b=j0. Exactly one585 charge gives116/31 ordered errors, both full sums below2\^{}-230; floors26/11 remain. No160+, signed positivity, arithmetic-content nonvanishing, period, effective source, specific032/Read, evaluated DP44 or RH conclusion. AM-R-093 DC1-19, on the actual d=7 branch, proves one finite-prime-unramified quadratic character epsilon with N\_E7/k(rho\_J)=epsilon\emph{chi\_cyc}rho\_C\_plus at every Tate prime. The integral rank-three norm at all primes, including2/7, gives K(J{[}n{]},C0{[}n{]}) contained in K(J{[}n{]})K\_epsilon, with relative degree at most2, simultaneously for every n.~Same actual projective/dual seeds and source congruences then give {[}Kgood:B{]}\textless=16128 and {[}K\_Sigma,out:B{]}\textless=32256\emph{M\^{}3}L\_N\^{}(129/16). The norm/type computation and arithmetic Artin convention are explicit; epsilon is neither evaluated nor assumed trivial. Root-discriminant bounds remain R092/LIFT092 on the same fields. Absolute coefficient-field degree/discriminant, effectiveN, Host/Sizes, selected arithmetic nonzero content, period, Read/specific032, evaluatedDP44 and RH remain open. No finite-error charge;116/31 and floors26/11 unchanged. AM-LIFT-093 PA1-27 refines the general-d degree for the exact same Kgood, J,C0 and original source seeds. Principal Weil pairings impose one shared scalar multiplier. Exact all-prime norm counts, including2/7/d and the overlappingd7 case, give {[}F(J{[}n{]},C0{[}n{]}):F{]}\textless=U(n)=n\^{}5\emph{product\_(p\textbar n,p=6 mod7)(1+1/p)\^{}3 for every finiteF/Kgood, and U(N\^{}2)\textless=N\^{}(81/8). The actual source level has multiplier1 and boundV(M)=M\^{}4}product}(p\textbar M)F\_J(p)(1-1/p)\textless2M\^{}4; no universal g mod7 condition is inserted. The existing C0{[}15{]} extension costs at most192, so {[}Kgood:B{]}\textless=32256h(-d)\textless=32256d and the actual common output satisfies {[}Kout\_Sigma:B{]}\textless=32256d V(M)U(L\_N\^{}2)\textless64512d M\^{}4 L\_N\^{}(81/8). R093 retains its independent strongerd7 bound. Same fields, all objects and intrinsic Step5 costs remain; effectiveN, absoluteB, Host/Sizes, selected nonzero arithmetic content, period, Read/specific032, evaluatedDP44 and RH are open. No finite-error charge;116/31 and floors26/11 unchanged. AM-K-109 TG1-19 retains the full all-A scalar theta polynomial and the original transverse Gram G\textgreater2\^{}36 I, with true determinantU. Using the entire actual K108 energy, it removes the complete k160 support26 cell with exact exponent-163509/512\textless-319. Original module membership, central affine fibre, common ideal correlation, all finite factors and both q/x monotonicities are retained. The necessary actual all-II rank-two-unipotent remainder becomes161\textless=k\textless=164 with2\^{}(71443/64)\textless x\textless2\^{}2108,v=q26,b=j0. Exactly one319 append gives117/31 complete ordered errors, both entire sums below2\^{}-230; complete floors26/11 remain. No161+, signed positivity, selected arithmetic content, effective source, period, specific032/Read, evaluatedDP44 or RH conclusion. AM-DO-115 DC1-24 proves an all-parameter scalar theta estimate using the exact positive cubic kernel, finite cap occupancy and nonzero-shell integral in real dimension336. The original selected angular profile gives beta777/352, and the unchanged saturated transverse Gram gives a complete257/256 theta factor. With every original paired ideal, affine central fibre, finite factor and both continuum monotonicities retained, the complete k161 support26 cell is below2\^{}-233, with exact exponent-119477/512. The necessary actual all-II rank-two-unipotent remainder becomes162\textless=k\textless=164 with2\^{}(71443/64)\textless x\textless2\^{}2108,v=q26,b=j0. Exactly one233 append gives118/31 complete ordered errors, both entire sums below2\^{}-230; unique minimum231,next233 and complete floors26/11 remain. No162+, signed positivity, selected arithmetic content, effective source, period, specific032/Read, evaluatedDP44 or RH conclusion. AM-K-110 PM1-22 retains the full actual angular profile as a positive Laplace mixture, includes the entire low-parameter tail in an explicit atom, and computes the exact inverse third moment. Weighted AM-GM and the all-minimum/all-parameter DO115 scalar theorem control arbitrary embedding weights in the unchanged original ideal and transverse Gram. Every transverse correction is retained, giving a uniform2180-bit row improvement. With all original affine central, paired ideal, finite and source factors and both continuum monotonicities, the complete k162 support26 cell is below2\^{}-1071, with exact exponent-548533/512. The necessary actual all-II rank-two-unipotent remainder becomes163\textless=k\textless=164 with2\^{}(71443/64)\textless x\textless2\^{}2108,v=q26,b=j0. Exactly one1071 append gives119/31 complete ordered errors, both entire sums below2\^{}-230; minimum231,next233 and complete floors26/11 remain. No k163/164, signed positivity, selected arithmetic content, effective source, period, specific032/Read, evaluatedDP44 or RH conclusion. DO116 provides an exact stronger all-minimum/all-parameter scalar coefficient for future complete applications; it changes no actual finite sector or main gate by itself. K111 supplies an exact stronger all-minimum/all-parameter scalar coefficient; it introduces no separate arithmetic charge or main-gate reduction before a complete application. AM-LIFT-094 HE0-22 proves conditional effective algebraic Hecke escape on a fixed indefinite arithmetic family. For a reduced algebraic bad locus of total degreeD and dimension\textless=s, finite equal-ample-pullback correspondences of degree\textless=delta with no nonempty proper closed forward-invariant locus admit escape in at mostB\_delta,s(D) steps, where A\_delta(D)=2 delta\^{}D D\^{}2, B\_delta,0(D)=D and B\_delta,s(D)=D+B\_delta,s-1(A\_delta(D)). For actual projective O\_E lattice rankr\textgreater=3 indefinite at every real place, f={[}F:Q{]}, an explicitly cleared rational Witt frame gives t=1+3McQ\_S, two physical original-M-level Hecke correspondences with delta=t\^{}(4fr\^{}2) and multiplier t\^{}2. Every forward orbit is dense; the conditional accumulatedN\textless=t\^{}B is obtained. The actual regular LIFT085 period-family map and nonempty inverse image of the prescribed nonvanishing open remain hypotheses; H,D,c,N0 and admitted Host/Sizes costs are not evaluated. Same original fields, quotient/action/line/level/frame conditions and independent entry multiplierN0 remain. No selected arithmetic nonzero content, global period, specific032/Read/STiesR, DP44 or RH conclusion. No finite-error charge;119/31 and floors26/11 unchanged. AM-R-094 PP1-24 removes the entire remaining conditional J56 all-top support26 class. R091 plus K109/DO115/K110 leaves only actual all-II global-rank-two-unipotent scalar-cross cells k163/164; the full correlated finite-pattern containing polynomial and low-parameter angular mixture give complete exponents -199349/512\textless-389 and -159413/512\textless-311, using DO116 and K111 respectively. Their union is below2\^{}-310. Append310 once to the exact K110119/31 arrays:120/31, both entire sums\textless2\^{}-230, complete floors27/11. The actual remainder now has support at least27 with every earlier higher cut and original finite/real/arithmetic convention retained. PP3, the sufficient signed higher-support inequality, remains open, as do actual Host/Sizes, selected nonzero arithmetic content, global period, specific032/Read/STiesR and evaluatedDP44/RH. No actual shorter mixed-CM chain or universal strict-negative cost theorem is asserted. Mandatory R093 source-scope precision: its d7 CM/level determinant and torsion theorem is conditional and supplies no nonempty full pinned Host. Literal Params6-22/32 and DataHost54 give r\textgreater=7,t\textgreater=13,d\textgreater365; PinHost39-41 fixes the same d.~Thus the d7 M\^{}3 specialization cannot serve an admitted full-source family. General-d LIFT093 is the admitted-source field-cost route; actual Host/Sizes, N and selected arithmetic content remain open. Full checked addendum: recovery\_pending/lift/r093\_full\_host\_qd\_big\_scope\_precision.md. This clarifies scope and adds no error or movement credit. AM-DO-117 FC1-25 removes the entire original conditional J56 all-top H133 fixed-S nonunipotent support27 sector, at every global rank and incidence profile. The five complete containing-set exponents are -5531348963/4096, -3078643005/4096, -2767418421/1024, -9563431915/2048 and -120294901/4096; their union is below2\^{}-29365. The new small scalar remains in the full paired I\_Sigma, while its exact original near-point transport uses only the universal RT2 selected Hilbert trace bound. Append29365 once to the complete R094120/31 predecessor:121/31, both full sums\textless2\^{}-230, complete floors27/11 unchanged. This is a strictly smaller conditional support27 remainder, with no full support27 exhaustion, actual admitted Host/Sizes, selected arithmetic content, prescribed global period, specific032 pairing, evaluatedDP44 or RH conclusion. AM-K-112 UE1-14 proves that every active type-II place in the exact full-positive globally rank-two-unipotent form is eligible, without an eta-denominator bound. The noneligible active subset is empty, with zero error charge. In the actual all-II scalar-cross support27 remainder, v=s=27, x\textgreater2\^{}2111; x\textless2\^{}3107 forces total eligibility27; x\textless2\^{}12000 gives total eligibility at most36 and at most one eta per fixed support. Fixed-eta active labels contribute at most(T-1)\^{}-27. The full168-factor bound additionally requires total eligibility27 and no inactive paired-C-zero place. Original paired row and central difference ideals retain their support27 denominator power54. Current DO117121/31 exact errors and complete floors27/11 remain unchanged; no full support27 estimate, actual source/content/period, specific032, DP44 or RH conclusion is asserted. AM-LIFT-095 IH1-6 proves an integral hyperbolic pair inside every stated selfdual relative-unramified indefinite Hermitian lattice of rank at least3, by local integral pairs and strong approximation on SU. The original nonterminal source W\_j.L0 satisfies these hypotheses conditionally on the admitted full record, and the same E-linear operators preserve its actual AA multiple. Choosing auxiliary doubled horizontal and integral imaginary root bases gives c=1 in LIFT094: t=1+3M Q\_S, delta=t\^{}(4fr\^{}2), L=t; any separately required larger congruence modulus is retained. This removes the unevaluated auxiliary clearing integer without a replacement lattice, seed translation or extra isogeny. The actual very ample line H, same-rule bad-locus degree D, nonempty family intersection, component-entry multiplier N0, selected nonzero rule and full Host/Sizes remain unresolved inputs. No evaluated total N, period, specific032, DP44 or RH conclusion; no finite charge. Current DO117121/31 arrays and floors27/11 remain unchanged. AM-DO-118 UB1-29 removes the entire actual conditional J56 fixed-S globally rank-two-unipotent support27 sector, all original profiles and labels. The rank35 high envelope retains the full five-column and same-pair isometry fibre. In the rank34 low branch, the true row metric, full beta norm fibre, affine central cosets and actual selected trace are retained; fixing row and beta before eta-dependent labels permits the one-label bound without global eta uniqueness. The six complete strict dyadic bounds have exponents8792,4491,3928,23649,17510,11371; their union is below2\^{}-3927. Append3927 once to DO117121/31, giving122/31 and both full sums below2\^{}-230; complete floors27/11 remain unchanged. This does not remove the other support27 classes, establish a support28 floor, or solve the signed inequality, actual admitted source, selected arithmetic content, prescribed period, specific032 pairing, evaluatedDP44 or RH. AM-K-113 SC27.1-27 removes seven complete containing groups in the actual conditional J56 all-II global-rank-two-unipotent scalar-cross support27 sector. The33 complete endpoints retain all four transverse mixture terms, every G pattern, all inactive-zero subsets, chosen active-integral positions, coordinate sizes and a separately enumerated infinite eta-ideal tail. Their full union is below2\^{}-304. The necessary residual is b=j0,2\^{}2111\textless x\textless2\^{}3328, either total eligibility27 and100\textless=k\textless=164, or total eligibility28,k164 with x\textgreater2\^{}3107. Append304 once to actual DO118122/31, giving123/31 with both full sums\textless2\^{}-230; complete floors27/11 remain unchanged. No whole support27 exhaustion, actual admitted Host/Sizes, signed inequality, selected arithmetic content, prescribed period, specific032 pairing, evaluatedDP44 or RH conclusion. AM-R-095 FW1-21 removes the entire actual conditional J56 fixed-S globally rank-one-unipotent support27 sector, all original profiles, labelled pairs, double-pole sizes and unmixed G patterns. A uniform weighted eta/content cutoff precedes the complete334-place binary pattern product; the remaining z exponent is165f, both genuine affine coordinates and the whole support/marked/type sum remain. The reduced exponent is-54242909/4096 and the complete witness exponent is-95126187/256; their union is below2\^{}-13000. Append13000 once to actual K113123/31, giving124/31 with both full sums\textless2\^{}-230; complete floors27/11 remain. No whole support27 exhaustion, signed inequality, actual admitted Host/Sizes, selected arithmetic content, prescribed period, specific032 pairing, evaluatedDP44 or RH conclusion. AM-DO-119 NA1-29 proves every active place of the actual full-positive globally rank-two-unipotent form is eligible, including distinct source/target labels and every type. Its complete non-all-II scalar-cross support27 bound retains all original paired rows, full-entry poles, affine central cosets, finite weights, relative norm fibres and the entire denominator-ideal tail. Six exhaustive branches give one412 charge. Append412 once to actual R095124/31, yielding125/31 with both complete sums below2\^{}-230 and floors27/11 unchanged. The dense all-II27 and separate nonzero-short-compression sector remain open; no full support28, actual admitted source/Host/Sizes, selected arithmetic content, prescribed period, specific032 pairing, evaluatedDP44 or RH conclusion. AM-K-114 FP27.1-29 removes the entire full-positive-compression global-rank-three scalar-cross support27 sector with nonzero original S-cross. All34 complete ideal/type cells,2925 padded sets, original full167-covector metric, relative norm fibre, projected U denominator/T108 column and G2 source fibre are retained. The complete eta-first all-ideal tail requires no coordinate sparsity. Literal inversion and23 common paired eligible positions force equal null lines outside four containing sectors, contradicting rank3. Their full union costs1650 once. Append1650 to actual DO119125/31, yielding126/31, both entire sums below2\^{}-230 and complete floors27/11 unchanged. Full positive compression remains a premise; the separate nonzero-short-compression27 and dense all-II27 sectors remain open. No full support28, actual admitted Host/Sizes, selected content, prescribed period, specific032, evaluatedDP44 or RH conclusion. AM-LIFT-096 TJ1-18 proves the same-rule bad-locus degree bound on the actual principal-3M special-unitary family of LIFT085/LIFT094, under its stated nonempty same-rule open-locus hypothesis. For g=fr,d=fpq, total frame weight w and theta weight n, take k=gw+n/2 for the literal even-n Dom; squaring the same rule gives k=2gw+n in general without changing its reduced zero set. A finite proof cover of degree at most e\_g=2\^{}(8g\^{}2+2g+4), first jets of the fourth power of the original symmetric principal line, and exact theta multiplier identities produce scalar cuts. For H=(psi\_bar\^{}\emph{L\_S)\^{}h very ample and V=deg\_H(Ybar), proper pushforward gives D\textless=D\_}=e\_g V sum\_\{j=0\}\^{}d k\^{}j, including mixed-dimensional components. When k=0 and the bad locus is proper it is empty. With LIFT095 c=1, LIFT094 then gives the conditional within-family multiplier N\textless=t\^{}\{B\_\{t\^{}\{4fr\^{}2\},fpq-1\}(D\_\emph{)\}, t=1+3M Q\_S; any independent component-entry multiplier N0 is retained. The proof cover contributes only to D\_}, not to the physical isogeny or output field. The numerical h,V, nonempty actual-family intersection, controlled N0, selected arithmetic nonzero rule and admitted full Host/Sizes remain open. No evaluated total N, six-identity chain, actual period, specific032 pairing, DP44 sign or RH conclusion; no finite error charge. Current K114126/31 exact arrays and complete floors27/11 are unchanged. AM-DO-120 QW1-29 removes the entire original doubly-plane scalar-cross sector with nonzero four-entry short compression at every support27..31, including all global characteristic types, actual coordinate supports and paired labels. The original integral positive F witness q\_j=\textbar\textbar(A-I)y\_j\textbar\textbar\^{}2 has every conjugate below2\^{}42. The complete paired finite low-rank proof restores the rank-four Green baseline666 with coefficients331/2 and335/2; both original full source/target costs, selected trace, arbitrary-ideal row, relative norm fibre and all incidence weights remain. All five full endpoint bounds give one270871 charge. Appending it once to K114126/31 yields127/31, with both entire sums below2\^{}-230 and complete floors27/11 unchanged. The complementary scalar-cross remainder at supports27..31 has full positive compressionI, a positive166 fixed space and global rank at most3 by the support-free original integral scalar proof; global unipotence is not inferred. The dense all-II27 window remains open; no complete support28, actual admitted Host/Sizes, selected arithmetic content, prescribed period, specific032 value, evaluatedDP44 or RH conclusion. AM-LIFT-097 VA1-15 evaluates the previously unspecified very-ampleness exponent on the same actual principal-3M arithmetic compactification X. For d=fpq,g=fr and A=psi\_bar\^{}\emph{L\_S, the exact integer h\_d=2\^{}(d+1)(d+1)(d+2)! makes H=A\^{}h\_d very ample. The no-PGL2 boundary condition and canonical-line extension identify K\_X as Cartier ample; the neat Baily-Borel pair is log canonical, A-K\_X is nef, and both applications of effective freeness/very ampleness retain log bigness on every lc/slc stratum. The original Hodge restriction and Hecke equal-pullback identities are unchanged. Writing v\_A=int\_X c1(A)\^{}d, the full ambient degree is h\_d\^{}d v\_A, so LIFT096 gives D\_}=2\^{}(8g\^{}2+2g+4)h\_d\^{}d v\_A sum\_\{j=0\}\^{}d k\^{}j, with k=gw+n/2 for even n or2gw+n after the same-zero square. The independent quantity v\_A is a positive unevaluated integer. LIFT094/095 then retain N\_within\textless=t\^{}\{B\_\{t\^{}(4fr\^{}2),d-1\}(D\_\emph{)\},t=1+3M Q\_S, with every protected congruence condition and any separate component-entry N0. The same-family nonempty open, selected arithmetic nonzero rule and admitted Host/Sizes remain open. No evaluated actual N, new physical field/level/isogeny, six-identity chain, period, specific032 pairing, DP44 or RH conclusion; no finite error charge. Current DO120127/31 arrays and complete floors27/11 are unchanged. AM-R-097 NW1-9 exhaustively composes the checked structural cuts at support27. Every remaining actual support27 arrow is full-positive, scalar-cross, globally rank-two-unipotent with all27 active places typeII, b=j0 and one of two necessary alternatives: total eligibility27,100\textless=k\textless=164,2\^{}2111\textless x\textless2\^{}3328; or total eligibility28,k164,2\^{}3107\textless x\textless2\^{}3328. DO120 supplies full positive compression before K114 is applied. The signed energy retains the actual R27 subset of this containing window and R\_ge28, with \textbar epsilon\textbar\textless2\^{}-230. No set equality, existence, emptiness, sign or floor28 conclusion. No new charge; actual DO120127/31 ordered arrays and complete floors27/11 are unchanged. Actual mixed-CM source, selected arithmetic content, six identities, prescribed period, specific032 pairing, evaluatedDP44, Kummer-Hecke RH and global oscillator comparison remain open. AM-K-115 GM1-23 proves the cap-dependent shape83/rate7/2 positive Gamma majorant, retains its entire low-weight interval by a paid positive shift, and controls every inverse moment0..k+3. All scalar and transverse terms, original saturated metric, central affine coset, source reconstruction,168 T factors and four-state finite patterns remain. The complete ten cells with total eligibility27,b=j0,2\^{}2111\textless x\textless2\^{}3328,k100..109 have total absolute contribution below2\^{}-945. The worst full exponent is-2907383/3072\textless-946, every other is below-3166. One945 append changes DO120127/31 to128/31, both entire sums remain below2\^{}-230 and complete floors27/11 stay. Combined with R097, the remaining actual support27 set is contained in the full-positive scalar-cross rank-two-unipotent all-II class with either total eligibility27,110\textless=k\textless=164,2\^{}2111\textless x\textless2\^{}3328, or eligibility28,k164,2\^{}3107\textless x\textless2\^{}3328. The next positive sufficient envelope is not an actual lower bound. No floor28, signed nonvanishing, admitted source, selected arithmetic content, prescribed period, specific032 value, evaluatedDP44 or RH conclusion. AM-LIFT-098 AV1-20 evaluates the previously unevaluated automorphic volume for the same actual principal-3M special-unitary arithmetic family. Put n=3M, r=p+q\textgreater=3, f={[}F:Q{]}, g=fr, d=fpq, a=r\^{}2-1, a\_r=product\_\{j=1\}\^{}\{r-1\} j!/(2pi)\^{}(j+1), b\_pq=product\_\{i=0\}\^{}\{p-1\}i!/(q+i)!, and e\_v=product\_\{j=2\}\^{}r(1-epsilon\_v\^{}j q\_v\^{}(-j))\^{}(-1), epsilon\_v=1 at split and -1 at inert finite places. The literal relative-unramified CM and ordinary selfdual-lattice assumptions give hyperspecial models even at dyadic places. The original invertible ideal multiple changes no stabilizer or congruence index. Exact strong approximation and finite orders give mu\_Gamma=D\_F\^{}(a/2)a\_r\^{}f n\^{}(fa) product\_\{v not dividing n\}e\_v with compact SU mass one. The physical Hodge representation S+Q\^{}dual gives lambda\^{}r=K\_Y\^{}2 and A\^{}r=K\_X\^{}(2(g+1)); Mumford proportionality yields v\_A={[}2(g+1){]}\^{}d d! b\_pq\^{}f mu\_Gamma. An explicit integer upper bound is Vcal=n\^{}(fa)D\_F\^{}ceil(a/2)2\^{}(f(r-1)){[}2(g+1){]}\^{}d d!(product\_\{j=1\}\^{}\{r-1\}j!)\^{}f.~Thus D\_arith=2\^{}(8g\^{}2+2g+4)h\_d\^{}d Vcal sum\_\{j=0\}\^{}d k\^{}j, h\_d=2\^{}(d+1)(d+1)(d+2)!, retains the full same-rule cost, and N\_within\textless=t\^{}\{B\_\{t\^{}(4fr\^{}2),d-1\}(D\_arith)\}, t=1+3M Q\_S. The independent volume input has been replaced by actual field, rank, signature and level data. Any component-entry N0, same-family nonempty open, selected arithmetic nonzero rule and admitted Host/Sizes remain open. No evaluated actual N, six-identity mixed-CM chain, period, specific032 pairing, DP44 difference or RH conclusion. No finite error charge; current K115128/31 arrays and complete floors27/11 unchanged. AM-K-116 RE1-28 retains the actual reciprocal-energy factor in a shifted positive half-Gamma convolution. The infinite total mass is explicit: every positive inverse moment is bounded by a convergent Gamma/Beta formula, and the unique polynomial degree0 term is made integrable by reserving theta=2\^{}-20 of the actual product-energy exponential before lattice enlargement. Every scalar/transverse term, original saturated metric and covolume, full167-entry reconstruction, central affine750 cover, all168 T factors and all paired four-state G patterns remain. The complete total-eligibility27,b=j\_int=0,2\^{}2111\textless x\textless2\^{}3328,k110 cell has full exponent-1539841/3072\textless-501. One501 append changes actual K115128/31 to129/31, retaining every old overlapping charge; both full sums remain below2\^{}-230 and complete floors27/11 stay. With R097/K115 the remaining actual support27 set is contained in the full-positive scalar-cross globally rank-two-unipotent all-II class with either total eligibility27,111\textless=k\textless=164,2\^{}2111\textless x\textless2\^{}3328 or total eligibility28,k164,2\^{}3107\textless x\textless2\^{}3328. The k111 positive sufficient envelope5269759/3072 is not a lower bound or impossibility theorem. No complete support28, signed nonvanishing, admitted source, selected arithmetic content, prescribed period, specific032 value, evaluated DP44 or RH conclusion. LATEST SOURCE-POINT EFFECTIVITY REDUCTION AM-LIFT-099: for a supplied nonzero literal rule and admitted indefinite source/CM seed, component entry is now bounded by C\_ent=m(2m)\^{}(2f)D\_F\^{}2, m=3M Q\_S, and at least one of its bounded-lattice regular families meets that same-rule open. With all LIFT095-098 factors retained, N\textless=C\_ent(1+m)\^{}\{B\_\{(1+m)\^{}(4fr\^{}2),fpq-1\}(D\_})\}, D\_\emph{=e\_g h\_(fpq)\^{}(fpq) Vcal sum\_(j=0)\^{}(fpq)k\^{}j, k=gw+n\_F/2 or2gw+n\_F after squaring. The original K0(A0{[}N\^{}2{]}) carries all final objects and nonzero frame; the actual same-J/C0 corollary has degree\textless64512d\_CM M\^{}4N\^{}(81/8), rd\textless2\^{}50exp(2pi sqrt(d\_CM))rd(Bcoef)MN\^{}2 and retains {[}Bcoef:Q{]} in the absolute-discriminant cost. The component-entry multiplier is included in this full N, without a ray-class field or proof-cover extension. Independent geometric c,h,volume,degree,nonempty-family and entry costs are removed under these stated source/rule hypotheses. Actual selected arithmetic nonzero, full Host/Sizes, remaining source-family/model/section/polynomial/fibre Step5 inputs, Read/specific032, actual g\_j,e\_j,H and evaluated DP44 remain open. No charge or finite-floor change. AM-DO-121 EC1-24 retains the exact original unipotent identity cosh r\textgreater=1+alpha/2, hence sinh squared r\textgreater=alpha(1+alpha/4), in the genuine selected coefficient. Its positive infinite-mass Gamma167/rate4 convolution has every required convergent inverse moment; original product-energy reservation repairs the sole degree0 term before enlargement. The complete actual total-eligibility27/support27 all-II globally rank-two-unipotent scalar-cross union k111..139 has absolute normalized contribution below2\^{}-993. Every scalar/transverse polynomial term, original saturated metric and paired ideal, full167-entry reconstruction, unique eta, central affine750 cover, all168 finite factors and correlated four-state patterns is paid. One993 append changes actual K116129/31 to130/31, with both entire sums\textless2\^{}-230, all old overlapping errors retained and floors27/11 unchanged. The remaining necessary support27 alternatives are total eligibility27,140\textless=k\textless=164,2\^{}2111\textless x\textless2\^{}3328, or total eligibility28,k164,2\^{}3107\textless x\textless2\^{}3328; b=j\_int=0 throughout. The next k140 sufficient envelope767061/1024 is positive and is not a lower bound or impossibility theorem. No complete support28, signed nonvanishing, admitted source, selected arithmetic content, prescribed period, specific032 pairing value, evaluated DP44 saving or RH consequence. AM-K-117 D29.1-27 proves the universal degree29 scalar Gaussian lattice bound from the fixed exact normalized Gegenbauer polynomial and, with AM-R-098 WJ1-28, removes the complete actual support27 total-eligibility28 branch k164,b=j\_int0,2\^{}3107\textless x\textless2\^{}3328 with absolute normalized contribution below2\^{}-278. The entire660-term polynomial includes the repaired zero degree and the unique degree167 term with its full Q\^{}167 pattern cost. Every original167-entry reconstruction, paired ideal, selected metric, central affine750 cover, unique eta, support and all finite factors is paid; the one eligible inactive place keeps contraction1. One278 append preserves all130/31 predecessor entries, giving131/31 with both entire sums below2\^{}-230 and floors27/11 unchanged. Only the eligibility27 alternative140\textless=k\textless=164,2\^{}2111\textless x\textless2\^{}3328,b=j\_int0 remains at actual support27. No complete support28, signed nonvanishing, actual admitted source, selected arithmetic nonzero rule, prescribed period, evaluated DP44 cost or RH conclusion. AM-DO-122 CE1-29 removes the complete actual support27 total-eligibility27 all-II globally rank-two-unipotent scalar-cross k140 sector with absolute contribution below2\^{}-359. MANDATORY CUTOFF QUALIFICATION: the prose after CE8 is read as norm(h\_tau)\textless=H=133, as proved by R048-RC RC48.4/RC48.6 and the paragraph after RC48.7; CE8 already has the correct non-strict sign. The frozen proof is read together with recovery\_pending/definite/DO122\_CUTOFF\_SOURCE\_ADDENDUM.md, whose full text is included in the checked paper. The joint physical central variable w=-U Im d, original affine ideal, full750 cover, weighted Voronoi integral and U\^{}f physical Jacobian are retained. The central integral consumes one profile power, giving shape166; all moments used and every positive polynomial term are paid, with original-row zero repair, all finite factors and correlated patterns. One359 append to actual K117131/31 gives132/31; both entire sums remain below2\^{}-230 and floors27/11 stay unchanged. The sole support27 necessary alternative is now total eligibility27,141\textless=k\textless=164,2\^{}2111\textless x\textless2\^{}3328,b=j\_int0. The positive k141 envelope is only a failed sufficient bound. No full support28, signed nonvanishing, actual admitted source, selected arithmetic content, prescribed period, evaluated DP44 or RH claim. CONDITIONAL LITERAL SOURCE-LAW REDUCTION AM-LIFT-100: given an actual algebraic frame, the original odd ideal half-level and principal rigidified line, literal IsNab and ThValTot, the independent ThNe0 premise follows at the exact source thresholdDF+2\textless=M. It supplies one individual permitted tuple for every unchanged nonzero harmonic polynomial overK(A{[}M\^{}2{]}), via the actual ideal-half orbit and quotient jet interpolation, with no principalization or further scalar-root field. Actual source/frame existence, IsNab/ThValTot, every other required source law and full Host/Sizes remain; the restricted coupled selected image is not identified with this full thGen family and its nonzero/content/period gate remains. No new arithmetic charge or evaluatedDP44. AM-DO-123 BC1-27 removes the entire original bounded support27/total-eligibility27 all-II full-positive scalar-cross globally rank-two-unipotent k141..142 union, with b=j\_int0 and2\^{}2111\textless x\textless2\^{}3328. The same R099 selected angular kernel, original physical central coordinate w=-U Im d, full affine covolume/U scaling/750 cover, shape166/shift111/100 moments and actual-row zero-degree repair retain every polynomial, finite, label, support, pivot, class/Haar and pattern factor. Full exponents-3344299/1024 and-1688491/1024 plus the entire K118 tail2\^{}-8605114 yield one1647 charge. Appending it once to actual DO122132/31 gives133/31 with both entire sums below2\^{}-230 and complete floors27/11 unchanged. The sole remaining support27 necessary alternative is eligibility27,143\textless=k\textless=164,2\^{}2111\textless x\textless2\^{}3328,b=j\_int0. The k143 envelope-28587/1024 is insufficient and unclaimed. No full support28, signed nonvanishing, source admission, selected arithmetic-image content, prescribed period, specific032 value, evaluatedDP44 or RH conclusion. CONDITIONAL ALL-RING SOURCE-LAW REDUCTION AM-LIFT-101: actual Host/Datum/oddLevel/Coeff/Moduli records and accepted CanonGaugeSt suffice to choose NabV with IsNab, all-ring ThValTot and harmonic NabUniq; thus independent connection/totality premises are removed where the operator may be chosen. With the actual FrPt over a finite number-field extension and DF+2\textless=M, AM-LIFT-100 then supplies literal ThNe0. For a separately prescribed operator its compatibility remains required. Actual source/Host/Sizes, other source laws, original coupled selected-image nonvanishing, period and evaluatedDP44 remain open; no new charge. AM-DO-124 CB1-27 removes the entire original bounded support27/total-eligibility27 all-II full-positive scalar-cross globally rank-two-unipotent k143 sector, with b=j\_int0 and2\^{}2111\textless x\textless2\^{}3328. The same accepted R100 narrow-band angular kernel and K120 actual narrow probability/moments retain the original physical central coordinate w=-U Im d, complete affine covolume/U scaling/750 cover, shape166/lambda32/25 profile, all167 inverse moments and actual-row zero-degree repair. Every original polynomial, finite, label, support, pivot, class/Haar and pattern factor is paid. The full exponent-725931/1024\textless-708 plus the entire new narrow-band tail2\^{}-8605114 gives one707 charge. Appending exactly once to actual DO123133/31 gives134/31 with both entire sums below2\^{}-230 and complete floors27/11 unchanged. The sole remaining support27 necessary alternative is total eligibility27,144\textless=k\textless=164,2\^{}2111\textless x\textless2\^{}3328,b=j\_int0. The k144 sufficient envelope924757/1024\textgreater0 is explicitly unclaimed. No full support28, signed nonvanishing, source admission, selected arithmetic-image content, prescribed period, specific032 value, evaluatedDP44 or RH conclusion. PRESCRIBED-FRAME AND ACTUAL-PACKAGE REDUCTION AM-LIFT-102: coherent exponential gauges supply a literal IsNab with exact ordinary coefficients; actual complex reduced frames have an all-ring pinned extension. For a supplied OMPins/OneMapTieP class package, its existing CovPin and io range already give the operative pin and harmonicity, so its full unchanged harmonic class relation agrees with the reduced-frame evaluation without a separate individual-coefficient equality premise. Actual package/source existence, short-value section laws not supplied, selected arithmetic content/nonzero output, period and DP44 remain open. No numerical charge. CHECKED ANALYTIC INPUT AM-K-121: the exact same-kernel coarea and paired height, cutoff-safe original physical central sum, positive exponential-tangent u-mixture and full actual selected-spread refinement are available. All arithmetic polynomial/moment/pattern/source obligations remain; these analytic factors alone remove no k cell and receive no independent movement credit. The full band tail is retained; no numerical charge. AM-DO-125 CM1--28 removes the complete original bounded support27/total-eligibility27 all-II full-positive scalar-cross globally rank-two-unipotent k144 cell, b=j\_int0 and2\^{}2111\textless x\textless2\^{}3328. Accepted K121 same-kernel coarea/height/variance is globally bounded by a paid exponential tangent before transport, then the exact physical central integral is bounded by a positive Gamma166/rate4, half-Gamma/rate2, half-Gamma/rate3514/29, infinite half-Gamma and shift29/20 convolution. Both complete quadrature tails and8064 intervals for degree147 yield a directed rational moment enclosure; finite-endpoint Hölder covers every lower positive degree and the actual-row reserve repairs degree0. All579 corrections, original supports/pivots/Green/T factors, full750 cover/U scaling and four finite small inflations remain. The full exponent is below−7622655/32768\textless−232.5, and the entire2\^{}-8605114 tail is paid once, giving one232 charge. The exact135/31 full error sums stay below2\^{}-230 with floors27/11. The sole remaining support27 necessary alternative is145\textless=k\textless=164 with the same denominator interval; the positive k145 envelope removes nothing. Signed PP3, higher-support remainder, actual source and selected arithmetic image, prescribed period, specific032 value, actual lift data/evaluatedDP44 and Hecke RH remain open. ALL-RING SHORT-VALUE REDUCTION AM-LIFT-103: with the actual geometric records, selfdual level-trivial Lambda and operative pinned nab2, the full first ShortValTot conjunct is constructed for every coefficient algebra and every original faithfully flat witness, with exact full trace and source normalization. The remaining ShortValTot obligation is the original-base sec/cdelta/IsShortRule descent conjunct; existence of the actual source/package, cusp/spanning laws, selected-image nonzero/content, prescribed period and DP44 remain separate. No numerical charge. COVER-RULE REDUCTION AM-LIFT-104: the cover sec(a,n) rule H2=cdelta}G2 is constructed with its exact positive FrBor character and t\_none\^{}n generator factor. IsLev supplies every field FromTest point needed for its actual evaluation. The remaining second ShortValTot conjunct is precisely an original-base f in sec(a,n) whose evaluation equals H2 for every original FrOver pair. Actual CovPin/LawsFull supplies A=A2=1; bare IsLev is not asserted to do so. Original-base descent, actual source/package, cusp/spanning, selected nonzero/content and DP44 remain separate. No numerical charge. GENUINE SELECTED COHERENCE INPUT AM-K-122: the original coefficient integral may be restricted to simultaneous I* and projective angular coherence\textgreater9/10, with full bounded nonzero-S-cross arithmetic errors2\^{}-8605114+2\^{}(11394885-B/4). The projection bound uses the actual selected metric and all words/twists/compact factors. The remaining high-coherence integral still needs a complete angular/row/central estimate; no cell removal, charge or support-floor change is inferred. AM-DO-126 IQ1--36 removes the complete original bounded support27/total-eligibility27 all-II full-positive scalar-cross globally rank-two-unipotent union145\textless=k\textless=159, b=j\_int0 and2\^{}2111\textless x\textless2\^{}3328. The saturated original isotropic quotient is a projective selfdual O\_E-module, including at2/U; every actual row lies in J I\_Sigma D\_eta times its ordinary row dual. Its standard S164 summand and selfdual rank3 complement give exact transverse covolume D0\^{}3 z\^{}(6f), while the independent physical central U\^{}-f remains. A genuine real-dimension1008 Gaussian bound retains both transverse terms. The unchanged positive central mixture, exact continuous moments148/162, finite-endpoint Holder, repaired zero term, all4590 positive terms, all original finite/source factors and the entire band tail yield one582 charge. Worst directed term is−19510525/32768\textless−595; the entire union is\textless2\^{}-582. Exact136/31 full error sums remain\textless2\^{}-230 with floors27/11. The sole remaining support27 necessary alternative is160\textless=k\textless=164 with the same denominator interval. The positive k160 sufficient envelope removes nothing. Signed PP3, higher-support remainder, actual source and selected arithmetic image, prescribed period, specific032 value, actual lift data/evaluatedDP44 and Hecke RH remain open. STANDARD-LATTICE SOURCE LAW AM-LIFT-105: the actual ShortThetaUnit obligation follows from the original geometric data and operative IsNab, with unit beta0 inverse and the full original normalization for every witness. This removes that named standard-lattice law from the remaining conditional source-construction gate. Nonstandard beta laws, original-base descent, cusp/spanning and selected arithmetic nonzero remain separate. No numerical charge. ORIGINAL-BASE DESCENT AM-R-101: the LIFT104 cover sec rule now descends to an original-base sec(a,n) with every literal IsShortRule evaluation, whenever actual IsLev/newA1/IsNab/PhiPin/PsiTie hold. Both ShortValTot conjuncts follow; no global identification of stored d.r is required. CovPin supplies the first three and dPhi supplies the character predicates for n\textgreater=r. This removes the named original-base descent obligation in that actual supplied scope. Small-n supply, source existence, cusp/spanning, selected arithmetic nonzero and evaluatedDP44 remain separate. AM-DO-127 GC1-16 removes five complete containing subsets of the original support27/eligibility27 full-positive scalar-cross globally rank-two-unipotent all-II branch, with b=j\_int0: k160/161/162/163/164 and log2x\textgreater=2120/2130/2140/2150/2165 respectively, inside the original x\textless2\^{}3328 window. The exact normalized finite Green formula gives2 Q\^{}(-335*167)(NEJ)\^{}167 NFJ0, whose NFJ0 cancels the independent physical central divisor. The complete pattern is product(1+q\_v\^{}(j-167))\^{}2, including exactly2\^{}668 at degree167. The actual projective selfdual quotient, full1008 transverse theta, physical U\^{}-f, same positive profile, all1630 terms and one entire band tail remain. Worst directed term -13968857/32768\textless-426 gives the full union\textless2\^{}-415. Exactly one415 append gives137/31 full sums\textless2\^{}-230 and unchanged floors27/11. The sole support27 necessary branch has2\^{}2111\textless x\textless2\^{}(X\_k) with the five respective upper exponents. Low-x cells, signed PP3, higher supports, actual source/selected arithmetic image, period, specific032 value, evaluatedDP44 and Hecke RH remain open. ALL-N SOURCE LAW AM-LIFT-106: actual AdmX forces n\textgreater=p+q=R.r; below r both ShortValTot conjuncts are vacuous for the unchanged selected DeltaDat. The original restricted dPhi supplies R101 in every remaining n, so unbundled PinDat/dLev/dNab/newA1/dPhi remove shTot as a separate all-n obligation. Upper-range cusp/spanning, actual source/package and selected arithmetic nonzero remain separate. CITED SOURCE SUPPLIER AM-R-102: accepted GlobalTheory.BetaEx.betaExists supplies the literal beta implication for every original datum/n/bad; together with LIFT106, beta and shTot need no separate proof under the seven remaining supplied LawsSh clauses. Actual KPsiIn is still needed for the selected beta witness and is equivalent to accepted same-field ReflexRoots, with finite ReflexBasis sufficient. The eventual coefficient-extension constructor does not by itself identify the prescribed Q.c.F/eO/PinField. No new determinant theorem or source existence is claimed. AM-K-123 HCX1-26 removes the complete original bounded support27/eligibility27 full-positive scalar-cross globally rank-two-unipotent all-II k160 cell, b=j\_int0 and2\^{}2111\textless x\textless2\^{}3328. The actual all-word selected evaluation gives a chi\textgreater937/1000 restriction with both full arithmetic errors paid. A525772-test outward interval cover proves the complete physical-coherence cap; exact three-real-dimensional Beta disintegration gives a fraction below2\^{}-10 with a paid margin. The same-kernel positive four-component central mixture retains the original affine U scaling, full750 cover, Young regularization and extra transport2. Nine full quadratures, finite Holder anchors, actual-row zero repair, the true1008-dimensional transverse theta, exact Green/J0 correlation, all322 terms and both full tails give one600 charge. Exact138/31 full sums remain below2\^{}-230 and floors27/11 remain. The sole support27 necessary alternative is161\textless=k\textless=164,2\^{}2111\textless x, with upper exponents2130/2140/2150/2165 respectively. Signed PP3, higher supports, actual source/selected arithmetic image/content/nonzero, prescribed period, specific032 value, evaluatedDP44 and RH remain open. SAME-FIELD SIGNATURE REDUCTION AM-LIFT-107: under the actual PinDat, all original KPsiIn stages are equivalent to stages1/2 and to the rank-one CM-type ideal-generator law. The exact reduced obligation is one finite normalized Cl(O\_E)-generator table inside the original F; adjacent-stage Bezout formulas retain every norm, original selected betaH and field pin. The table is not yet supplied, and ordinary principalization or a freely enlarged field does not discharge it. QUADRATIC NORM-IMAGE REFINEMENT AM-LIFT-108: the original-field normalized table may be taken on G=image Cl(E)-\textgreater Cl(K), K=Q(H.rt), via Phi(I)=iota(N\_E/K I)O\_F. The exact table is still required inside the original F. A separate explicit root tower supplies a new coefficient record of degree≤\textbar G\textbar≤h\_K with small normalized scalar heights and upper-field persistence, but does not discharge old eO/PinField or changed-field DP44 costs. No claim that actual G is proper; full-source d\textgreater365 excludes the geometric d7 shortcut.

\subsection{Q3 --- OPEN}\label{q3-open}

K021 identifies the exact source Lprime/L height slot, proves direct CM-type height-class substitution cannot supply the cubic S3 standard component, and supplies the explicit convergent ray-theta Mellin period carrying that slot. Remaining: a new arithmetic comparison to that period, prime-norm/conductor costs, and the every-epsilon window. No Hecke RH proved.

\subsection{Q4 --- PARTLY\_PROVED}\label{q4-partly_proved}

AM-DO-002 canonical obstruction; AM-DO-003 Fourier/delta-module map; AM-DO-005 bicanonical local comparison; global arithmetic compatibility remains open.

Review 40; next review 2026-11-23T14:35:32+00:00. Latest durable checkpoint 50; next checkpoint due 2026-11-23T14:20:32+00:00. Current 138/31 exact errors, floors27/11.

\subsection{Checkpoint54 pending and save boundary}\label{checkpoint54-pending-and-save-boundary}

AM-DO-129 is author-frozen, with full root final proof/script/source precheck PASS and independent3291 replay at14:14:52.717535. Its proposed energy identity uses a nonnegative sum of original single-lattice Heisenberg kernels of identity mass(T-1)\^{}-168; all330 inherited complete terms, central identities and one global outer tail give a separate1266 one-sided debit. It proposes reducing the signed gate to the actual remaining support≥28 set. Whole peer review and a distinct one-sided accounting integration remain pending. It makes no absolute R27 estimate and cannot append1266 to the absolute arrays or claim an absolute support floor28. The exact frozen author packet and root precheck are retained.

AM-LIFT-109 is author-frozen: complete122879 diagnostic at14:14:45.834073, whole author assessment14:15:42.093433. It proves norm-class surjectivity using odd quadratic class number, central CM conjugation and the full non-Galois Frobenius argument, then a necessary exponent-divisibility test for the original coefficient degree. Its elementary split-two ideal class bound would give{[}F:K{]}≥7 and{[}F:Q{]}≥14 for a supplied full pinned source satisfying kPsi. Root has read the development and independently checked the split-two derivation; final frozen full proof/source review remains pending. DO's own122879 replay at14:16:47 passed, but its complete source review/report is still pending. No canonical conclusion or charge is taken from it yet.

AM-K-124 remains a coherent pending proof. Its full650-term arithmetic diagnostic completed14:09:11.597321 with4936976 assertions and615936 quadrature intervals. The full eleven-shell continuum cap is not established; proposed350 charge is conditional and is not applied. The former cap launch has no observed completion or usable live receipt. The exact progress record, failed64-tail envelope, partial earlier shells, changed proof/script history and a reversible owned read-only output guard are retained. No cap PASS or live execution is claimed. All agent writes are held for capture, with no observed active writer; the guard prevents the only possible late output write from the unobserved launch.

\subsection{Main task status and durable-save limitation}\label{main-task-status-and-durable-save-limitation}

Q1,Q2,Q3 remain open and Q4 partly proved. No shorter admitted lift chain, evaluated g\_j/e\_j/H or DP44 sign, specific032 value, prescribed global period, selected arithmetic nonzero image, or Kummer-Hecke RH has been supplied. The current paper keeps accepted and pending findings in separate parts.

Latest confirmed durable checkpoint is50. Exact immutable checkpoints51,52,53 await durable saving because the13:41:44 ordered retry failed at prepare\_uploads with library\_storage\_limit\_exceeded, before transfer or finalization. Their ZIP hashes and prior checkpoint50 receipt remain preserved; no earlier ZIP was overwritten, renamed, deleted or nested. Checkpoint54 will be captured at its scheduled slot and the exact51/52/53/54 files retried in order. Capture alone is not a successful durable save.

\clearpage

\section{DO069 recovery: exact source-interface reconciliation}\label{do069-recovery-exact-source-interface-reconciliation}

Author: \nolinkurl{/root/definite_oscillator}. New read-only source checks after the 22 November 2026 recovery. All new artifacts remain in the authorized \nolinkurl{recovery_pending/definite} directory. Existing proof files, indices, ledgers, and checkpoints have not been modified.

\subsection{What this record establishes}\label{what-this-record-establishes}

The reconstructed \nolinkurl{DO069_final_five_part_proof.md} is a new file. The outage mathematical-text reviews at 14:42:58 and 14:46:53 UTC did not certify its bytes. After recovery I read the restored source definitions and comparison interfaces listed below, then whole-read the reconstructed proof and its fresh exact-diagnostic script. The new source reads resolve the specific ordering, parity, phase, and pointwise-equality questions that those earlier reviews carried explicitly.

This is an interface check of the stipulated pinned source, not a Lean build or a new audit of the proofs of stipulated theorems. It does not establish an arbitrary abstract \nolinkurl{ThetaX} instance or arithmetic \nolinkurl{Frames.Good}.

Pinned commit: \nolinkurl{7cbf7a0c32fca1badf05931879a37c3f8a2b03b4}. Source fingerprints are recorded in the companion \nolinkurl{source_fingerprints.json}.

\subsection{Exact clauses read}\label{exact-clauses-read}

\begin{longtable}[]{@{}
  >{\raggedright\arraybackslash}p{(\columnwidth - 4\tabcolsep) * \real{0.3333}}
  >{\raggedright\arraybackslash}p{(\columnwidth - 4\tabcolsep) * \real{0.3333}}
  >{\raggedright\arraybackslash}p{(\columnwidth - 4\tabcolsep) * \real{0.3333}}@{}}
\toprule\noalign{}
\begin{minipage}[b]{\linewidth}\raggedright
Source
\end{minipage} & \begin{minipage}[b]{\linewidth}\raggedright
Directly checked clauses
\end{minipage} & \begin{minipage}[b]{\linewidth}\raggedright
Consequence used in DO069
\end{minipage} \\
\midrule\noalign{}
\endhead
\bottomrule\noalign{}
\endlastfoot
\nolinkurl{Formalism/SharpRestriction/CuspSpanNormBounds/CayleySignature.lean} & \nolinkurl{car}, \nolinkurl{vth}, \nolinkurl{bB}, \nolinkurl{PerNZ}, and the complete \nolinkurl{ThetaUnfoldingFreeData} definition, especially lines 168--175 & The horizontal input and upper cusp vector precede the existential new vertical input. The conclusion asks for \nolinkurl{fact} and \nolinkurl{kfin}, followed by the implication for every cusp vector equal to that lift. It does not ask that the final tensor remain in \nolinkurl{Loc}. \\
\nolinkurl{Formalism/LocalRingCardinality/DataTheta.lean} & \nolinkurl{TPar} at lines 501--504: actual \nolinkurl{t}, \nolinkurl{t_odd}, and \nolinkurl{t_ge}; the separate \nolinkurl{ofWord} specialization & Oddness and the lower bound apply to the actual parameter \nolinkurl{P.t}. No equality between general \nolinkurl{P.t} and \nolinkurl{Ranks.t} is inferred. \\
\nolinkurl{Development/AbelianVarietyTheory/Definitions/LatticeFrames/IsThetaFam.lean} & \nolinkurl{VolM}, \nolinkurl{VolU}, all \nolinkurl{Step} fields, and \nolinkurl{Step.lift} at lines 99--100 & The output is literally \nolinkurl{liftS s.om s.vW.μ s.vW.𝓕 G Φ}. The actual \nolinkurl{vT} ordering is \nolinkurl{VolM D σ V₀.rank}. \\
\nolinkurl{Development/AbelianVarietyTheory/Definitions/TwistedInnerProducts/TowerEx.lean} & \nolinkurl{Lnk}, \nolinkurl{Loc}, its annihilator/vacuum conditions, \nolinkurl{UpBF} at lines 129--138, \nolinkurl{UpFB} at lines 140--148, and the separate \nolinkurl{Frames.Good} clauses & Both unfolding conclusions are nonzero scalar periods of the uncontracted \nolinkurl{Step.lift}. The earlier \nolinkurl{σ=0} FB boundary still creates a rank-one lower Fourier--Jacobi period. Arithmetic frame conditions are additional. \\
\nolinkurl{Formalism/LocalToGlobal/BesselHermitianTheta/TiedWeil.lean} & \nolinkurl{TwRule}, \nolinkurl{om_swap}, and the complete \nolinkurl{WeilTied.th} field & The labels exchange under the transpose/dual-Hom transport. The full-Schwartz theta comparison is pointwise, with its literal nonzero \nolinkurl{petC} factor. This comparison, rather than a proposed extension of an algebraic intertwiner, is the link to the abstract theta output. \\
\nolinkurl{Formalism/LocalToGlobal/IndefiniteReal/OrthogonalComplementCharacterisation.lean} & \nolinkurl{IsIo.integ} and \nolinkurl{IsIo.eq}, including \nolinkurl{∀ F : CCs, ∀ φ v g} & The fixed coefficient-model equality is pointwise on its allowed polynomial and real input modules. It can be invoked after extracting an actual compact-finite input. \\
\nolinkurl{Formalism/LocalToGlobal/IndefiniteReal/LinksChar.lean} & The complete \nolinkurl{IsItW} fields and polynomial domains & These fields do not justify a Schwartz-continuous extension of an arbitrary algebraic intertwiner. DO069 makes no such inference. \\
\nolinkurl{Formalism/LocalToGlobal/BesselHermitianTheta/TiedKfin.lean} & \nolinkurl{IsFockCx}, including compact commutation and definiteness, \nolinkurl{IsPolFock}, and the complete \nolinkurl{KfinTied.kfin} equivalence & Compatible actual source/target compacts admit the Fock/Hermite finite-vector extraction used in the proof. The conclusion is reconciled with the literal polynomial-annihilator formulation of \nolinkurl{kfin}. \\
\nolinkurl{Development/AbelianVarietyTheory/Definitions/SemilocalFrames/IsBesPosOfBesselTied.lean} & \nolinkurl{tPt}, \nolinkurl{cenCo}, \nolinkurl{jacOp}, nonzero \nolinkurl{β} and \nolinkurl{c} in \nolinkurl{FJDat}, \nolinkurl{IsFJPos}, \nolinkurl{FJTied}, \nolinkurl{PerTied}, and \nolinkurl{IsTranslSum} & The arbitrary-position normalization uses actual nonzero rational coefficients. The finite translate is transferred to the new vertical tensor. The FJ phase and the WA/WO transport are typed explicitly. \\
\nolinkurl{Development/AbelianVarietyTheory/Definitions/TwistedInnerProducts/WittSmallGroupEmbeddingMemPointwiseStabilizerIso.lean} & \nolinkurl{psiFlag}, \nolinkurl{ptS}, \nolinkurl{cenS}, \nolinkurl{jacS}, \nolinkurl{thetaJS}, \nolinkurl{liftS}, \nolinkurl{liftT}, and \nolinkurl{fjSrcFun} & The source integrand contains the conjugate of \nolinkurl{thetaJS} and of the generic flag character. The corrected dagger/central-phase calculation matches this once, with no second coordinate conjugation. \\
\nolinkurl{Formalism/LocalToGlobal/IndefiniteReal/LocalNewformFunction.lean} & Oriented \nolinkurl{WO}, symplectic/Heisenberg definitions, and \nolinkurl{toTWsmall} at lines 80--85 & The map \nolinkurl{J(w)=wᵗℓ} is the inverse rank-ordered WA/WO conversion used in the reconstructed proof. \\
\nolinkurl{Development/AbelianVarietyTheory/Proofs/DecoratedLatticeProofs/WeilRepConjSourceLineEqTargetLineStatement.lean} & The exact statement and all line, norm, character, and conjugate-adjoint hypotheses & The BF choice \nolinkurl{ζ′=bar(c)}, \nolinkurl{ζ=ξ₀ζ′}, \nolinkurl{τ=-ξ₀} gives the actual \nolinkurl{κ=-D_h}; \nolinkurl{Ξ₂(z)=f(-z)} compensates that coordinate sign exactly. The stipulated proof body was not re-audited. \\
\end{longtable}

The central versions were used where available. The semilocal FJ file and the line-comparison statement were read from the already retained \nolinkurl{definite_oscillator_seesaw_source} cache, whose exact bytes are fingerprinted below.

\subsection{Analytic order retained after the reread}\label{analytic-order-retained-after-the-reread}

The nonzero conclusion is first the uncontracted scalar theta-period functional on the genuine Schwartz realization. Fixing the genuine smooth moderate-growth cusp input, weighted theta summation and Miller--Schmid, Theorem 2.24, give a Schwartz-seminorm majorant. Compact-finite Hermite density then supplies one pure tensor finite under the actual compatible compacts. Only at that input is the fixed coefficient/Harish--Chandra equality invoked. The pointwise \nolinkurl{WeilTied.th} and \nolinkurl{IsIo.eq} fields now make the a.e.-to-smooth fallback unnecessary for those exact records; it remains a valid conditional observation for any other realization that is initially only a.e.

The nonzero normalization factors remain explicit. Nothing in the reread sets them to one, supplies a new arithmetic level, or changes the fixed horizontal tensor. The output need not remain in the auxiliary \nolinkurl{Loc} space.

\subsection{Fresh checks and their limits}\label{fresh-checks-and-their-limits}

\nolinkurl{recovery_exact_diagnostics.py} is a fresh reconstruction, not the lost pre-outage script. Its recorded replay passed 17,952 assertions at 15:27:18.382724 UTC, including 2,880 rank/parity cases, 144 exact rational Gaussian phase models, all stated R053 profile comparisons, K085 constants, and exact string-presence checks for the principal source-order clauses. The phase tests include the nonzero-coordinate half-pairing, translated argument, and the BF coordinate parity.

The mathematical proof, source interfaces, and code were read separately. The script is not an automated proof of analytic continuity, a Lean proof checker, or a complete previous-error-ledger replay. Historical reviews remain historical; this new recovery record does not retroactively change their scope.

\subsection{Remaining scope}\label{remaining-scope}

Within the named genuine R014/R015 realization and its actual comparison records, the previously carried source-order/parity/pointwise-equality prerequisites are now reconciled. The broader construction still requires the actual genuine comparison records themselves, the accepted DO068 local construction and R016 terminal cases, and the stated upper nonzero period premise. No assertion is made for an arbitrary abstract \nolinkurl{Hyp3} realization. Literal Host/SizeData, arithmetic \nolinkurl{Frames.Good}, prescribed finite row levels, norm budgets, a global period value, DP44, and RH remain outside this result.

\clearpage

\section{LIFT62 --- restored source reconciliation and remaining concrete pin}\label{lift62-restored-source-reconciliation-and-remaining-concrete-pin}

\textbf{Status:} controlled recovery note, 2026-11-22. This note records fresh reads of restored/pinned source bytes and corrects the source-interface wording of \nolinkurl{lift62_outage_draft.md}. It is not a reconstruction of historical final bytes, a Lean build, or a construction of the requested J/032 package. Existing promoted files are unchanged.

\subsection{1. Smooth vectors and the actual continuity field}\label{smooth-vectors-and-the-actual-continuity-field}

The freshly fetched pinned \nolinkurl{Formalism/LocalToGlobal/BesselHermitianTheta/TiedType.lean}, lines 18--26, defines \nolinkurl{SmoothVec} by smooth dependence on every finite product of the actual exponentials \nolinkurl{exp(t_i X_i)}, with each \nolinkurl{X_i} in the actual Lie algebra. It defines \nolinkurl{Sobolev₂ ell} by one finite list of Lie words and a constant controlling \nolinkurl{|ell(w)|} by the sum of the Hilbert norms of those differentiated vectors. Thus, on the smooth-vector Fréchet space of the genuine strongly continuous unitary representation, this is precisely a continuous linear functional of finite Sobolev order. The smooth vectors are a linear subspace: sums and scalar multiples preserve the displayed smooth orbit maps. Consequently their linear span in \nolinkurl{DSModel.sm} is the same vector space, not an enlargement by nonsmooth vectors.

\nolinkurl{IndefiniteReal/ScalarOnPartsFull.lean}, lines 82--90, defines \nolinkurl{IsMultMap} with the following actual fields: smooth automorphy of every value; cuspidality of every value; the fixed finite-level character; equivariance for the full real unitary group in each coordinate; and, separately in each coordinate and at each adelic point, an extension to a linear functional satisfying \nolinkurl{Sobolev₂}. The space \nolinkurl{Mult} is the linear span of those maps. All five properties persist under finite linear combinations.

The source supplies separate continuity, not a joint product-Sobolev inequality as a field. The required implication is elementary. A separately continuous finite multilinear scalar map on Fréchet spaces is jointly continuous. Induct on the number of factors. By induction, after fixing the first input the map in the remaining factors is jointly continuous. Choose decreasing countable neighborhood bases in those factors. For each common index n let A\_n be the closed, balanced, convex set of first inputs for which the absolute value is at most one on the corresponding product neighborhood. Every first input lies in some A\_n.~Baire's theorem gives one A\_n with nonempty interior. Subtraction of two points in this interior, followed by scaling, yields a neighborhood of zero in the first factor and a uniform bound on its product with the remaining neighborhoods. Multilinearity then gives a product of continuous seminorms. Each smooth-vector seminorm can be increased to a finite sum of Lie-word Hilbert norms. This proves the joint product-Sobolev bound used for simultaneous K-finite approximation in ER20; it is a deduction, not a silently imported source field. In the forward direction the source asks for an algebraic linear functional on the whole local Hilbert space whose restriction has this bound. Extend the continuous functional on the smooth subspace algebraically to the whole vector space; the Sobolev inequality is required only on smooth vectors and is unchanged. This is not an assertion that the extension is Hilbert-continuous.

\nolinkurl{BesselHermitianTheta/TiedSmooth.lean}, lines 44--51, defines \nolinkurl{IsAutFunSm} to include a single polynomial growth exponent for all derivative lists (with constants allowed to depend on the list and the finite component), in addition to smoothness and central-finiteness. The source therefore already supplies uniform moderate growth in the converse. \nolinkurl{IsAutFunDef.lean} gives actual rational left invariance, continuity and an open finite level; its cusp condition quantifies over every rational isotropic flag and its normalized unipotent quotient integral. The LIFT62 compact-centre/Wallach argument addresses exactly that constant-term condition.

\subsection{2. Ordering of the Hilbert reconstruction}\label{ordering-of-the-hilbert-reconstruction}

In the converse ER20, the invariant form first yields a \textbf{bounded Hilbert operator} A extending the Harish--Chandra restriction of T. Calling it a full-group intertwiner at that point is premature. For a smooth input v, choose K-finite v\_n tending to v in the smooth-vector Fréchet topology. The preceding continuity deduction gives pointwise convergence T(v\_n) to T(v); boundedness gives A(v\_n) to A(v) in L2; an almost-everywhere convergent subsequence identifies A(v)=T(v) almost everywhere. Now the already stipulated full-group equivariance of T implies equivariance of A on smooth inputs, and hence on the Hilbert space by density. The fixed Sobolev estimate then supplies the canonical smooth representative. Because this representative and T(v) are continuous and agree almost everywhere for a full-support Haar measure, they agree everywhere.

This correction only concerns the order of the proof. It neither assumes a bounded extension of an arbitrary algebraic Harish--Chandra map nor requires integrability of arbitrary smooth matrix coefficients. The forward rank-one reconstruction remains on the algebraic cyclic span of the chosen unit vector, where every coefficient is a finite linear combination of translates of the one stated L1 coefficient.

\subsection{3. Exact source paths and analytic projection}\label{exact-source-paths-and-analytic-projection}

The correct recovered paths are:

\begin{itemize}
\tightlist
\item
  \nolinkurl{Formalism/LocalToGlobal/IndefiniteReal/ScalarOnPartsFull.lean} --- actual contragredient parameter, DS model, smooth multipliers and inner product.
\item
  \nolinkurl{Formalism/LocalToGlobal/IndefiniteReal/AutomorphicEvaluation.lean} --- lowest vector, evaluation and norm identity.
\item
  \nolinkurl{Formalism/LocalToGlobal/SiegelLatticeTheory/HaarTransversals/ProjRows.lean} --- coefficient projection, Schur convolution and conditional Hom bijection.
\item
  \nolinkurl{Formalism/LocalToGlobal/SiegelLatticeTheory/HaarTransversals/RowsDef.lean} --- Hilbert row space and conditional row bijection.
\item
  \nolinkurl{Formalism/LocalToGlobal/SiegelLatticeTheory/FlagStabilizerTensors/RowsJoinPinned.lean} --- actual instantiation, Haar ratio and the final arithmetic pin.
\item
  \nolinkurl{Formalism/LocalToGlobal/SiegelLatticeTheory/FlagStabilizerTensors/RowDict.lean} and \nolinkurl{STiesR.lean} --- row conjugation and its exact supplied compatibility field.
\end{itemize}

These correct the two misplaced \nolinkurl{HaarTransversals/RowsJoinPinned} and \nolinkurl{HaarTransversals/AutomorphicEvaluation} locators in the retained outage draft. The analytic equality is between the actual smooth multiplier space and the actual projected automorphic space, with the stated finite level and finite good projector. It does not identify an unrelated record named \nolinkurl{what₂.EE} merely by notation.

\subsection{4. What STiesR supplies, conditionally}\label{what-stiesr-supplies-conditionally}

The complete freshly read \nolinkurl{STiesR.lean} extends \nolinkurl{STies} and includes an actual instance equality, a \nolinkurl{GoodOp X}, and

\nolinkurl{rowTie : ∀ hF : FinX Y, RowDictTie X G (hI hF).eHH}.

\nolinkurl{RowDictTie} is \nolinkurl{IsConjOnto} for \nolinkurl{iotaH ∘ D.sEE} and the canonical conjugation map \nolinkurl{conjHH}. Its witness is a linear isometry from the actual \nolinkurl{D.what₂.EE} onto the corresponding \nolinkurl{EEC}; it is not merely injectivity of \nolinkurl{sEE}.

Let j\_eval be the analytic evaluation equivalence \nolinkurl{EP ≃ EEC}, with exactly the Haar ratio c from \nolinkurl{RowsJoinPinned}, and let j\_row be the rowTie witness \nolinkurl{what₂.EE ≃ EEC}. Then

\nolinkurl{k = j_eval.symm ∘ j_row}

is the required linear equivalence \nolinkurl{what₂.EE ≃ EP}, and \nolinkurl{j_eval(k F)=j_row F} proves the final conjugated \nolinkurl{sEE} equality. The norm identity retains c; no isometry convention sets it to one. This deduction is exact \textbf{conditional on an actual STiesR instance and the matching evaluation data}.

\subsection{5. The translation and package gates remain}\label{the-translation-and-package-gates-remain}

\nolinkurl{IndefiniteReal/EvalTieTOne.lean} defines the translated evaluation using \texttt{T(v0)(iota\^{}\{-1\}(x)\ *\ t)}, together with the literal conjugated-level/character relation. \nolinkurl{TrISepWith.lean} uses that construction at the supplied \nolinkurl{tS}, and its \nolinkurl{seeDef} retains the scalar \nolinkurl{uS * kappaS}. In contrast, \nolinkurl{RowsJoinPinned.lean} uses unshifted evaluation \texttt{T(v0)(iotaX\^{}\{-1\}(x))}. One must supply the actual transported RealDat/level/frame equality before substituting one for the other. Neither tS nor its character is omitted by the analytic theorem.

The complete restored \nolinkurl{Development/AbelianVarietyTheory/Definitions/LatticeFrames/CxTiesOf.lean} defines \nolinkurl{RealPack} with an \nolinkurl{S : STiesR R} as a field. It therefore packages, rather than proves, the rowTie needed above. The related \nolinkurl{BotHyp}, \nolinkurl{HypAt}, \nolinkurl{LevHyp} and finite compatibility results retain their actual level, frame, period and row pins as assumptions. No concrete constructor of this STiesR for the requested J/032 instance has been recovered in this audit. A bare abstract D7 only supplies its own anti-isometric \nolinkurl{sEE}; it does not supply rowTie surjectivity.

Accordingly the exact analytic bijection and the conditional k construction are proved at their stated interfaces. The concrete J/032 record equality and an instantiated arithmetic STiesR remain open. This note does not infer a prescribed Bessel/Mram period, an effective nonzero determinant, or a literal full family from that analytic bijection.

\subsection{6. Verification boundary}\label{verification-boundary-2}

All source declarations quoted above were directly read after recovery. Fetched files use the pinned commit \nolinkurl{7cbf7a0c32fca1badf05931879a37c3f8a2b03b4}; fresh SHA metadata are in \nolinkurl{source_recovery/retrieval_*.json}. The historical intake manifest was unavailable, so these fetches are not asserted to reproduce a separately hashed old local copy. The entire retained LIFT62 mathematical draft was read after export; this note supplies its explicit continuity and ordering corrections. Final integrated proof text, independent frozen-file review and any new full certificate remain pending. No Lean compilation was run.

\clearpage

\section{Checkpoint54 research snapshot}\label{checkpoint54-research-snapshot}

Snapshot UTC: 2026-11-23T14:18:50.336436+00:00. Capture remains scheduled for2026-11-23 14:20:32 UTC. Review40 recorded mathematical movement; the next hourly review is14:35:32. The stop counter is zero. No final package is authorized by that cadence yet.

\subsection{Canonical results after checkpoint53}\label{canonical-results-after-checkpoint53}

AM-K-123 was accepted13:44:41: the full k160 support27 sector is removed by a new600 absolute charge. The exact current budget is checks/k123\_joint\_remainder\_error.json, with138 all-top and31 uniform entries and complete absolute floors27/11. Both full sums remain below2\^{}-230.

AM-LIFT-107 was accepted13:50:09: all original kPsi signatures are equivalent to stages1/2 and one normalized CM-type ideal-generator law in the same field, with a finite class-generator table. AM-DO-128 was accepted14:03:40: exact finite Heisenberg zero mode, paired CRT, full-kernel positivity and frequency bounds; it alone makes no numerical main-gate reduction. AM-LIFT-108 was accepted14:10:03: the actual type ideal is the quadratic relative norm ideal, so the normalized table may be taken on the norm-image class group; explicit roots give a new coefficient record of relative degree at most its order, with small scalar heights. This construction does not identify the new field with the original PinField.

The canonical support27 necessary branch remains total eligibility27, all active II, full positive compression, global rank-two unipotence, b=j\_int0,161≤k≤164 and2\^{}2111\textless x\textless2\^{}Xk, with X161=2130,X162=2140,X163=2150,X164=2165. The actual signed sum remains unproved; the untouched higher-support restrictions are retained. No one-sided support28 gate is installed yet.

\subsection{Pending findings and exact review boundaries}\label{pending-findings-and-exact-review-boundaries}

AM-DO-129 is author-frozen, with full root final proof/script/source precheck PASS and independent3291 replay at14:14:52.717535. Its proposed energy identity uses a nonnegative sum of original single-lattice Heisenberg kernels of identity mass(T-1)\^{}-168; all330 inherited complete terms, central identities and one global outer tail give a separate1266 one-sided debit. It proposes reducing the signed gate to the actual remaining support≥28 set. Whole peer review and a distinct one-sided accounting integration remain pending. It makes no absolute R27 estimate and cannot append1266 to the absolute arrays or claim an absolute support floor28. The exact frozen author packet and root precheck are retained.

AM-LIFT-109 is author-frozen: complete122879 diagnostic at14:14:45.834073, whole author assessment14:15:42.093433. It proves norm-class surjectivity using odd quadratic class number, central CM conjugation and the full non-Galois Frobenius argument, then a necessary exponent-divisibility test for the original coefficient degree. Its elementary split-two ideal class bound would give{[}F:K{]}≥7 and{[}F:Q{]}≥14 for a supplied full pinned source satisfying kPsi. Root has read the development and independently checked the split-two derivation; final frozen full proof/source review remains pending. DO's own122879 replay at14:16:47 passed, but its complete source review/report is still pending. No canonical conclusion or charge is taken from it yet.

AM-K-124 remains a coherent pending proof. Its full650-term arithmetic diagnostic completed14:09:11.597321 with4936976 assertions and615936 quadrature intervals. The full eleven-shell continuum cap is not established; proposed350 charge is conditional and is not applied. The former cap launch has no observed completion or usable live receipt. The exact progress record, failed64-tail envelope, partial earlier shells, changed proof/script history and a reversible owned read-only output guard are retained. No cap PASS or live execution is claimed. All agent writes are held for capture, with no observed active writer; the guard prevents the only possible late output write from the unobserved launch.

\subsection{Main task status and durable-save limitation}\label{main-task-status-and-durable-save-limitation-1}

Q1,Q2,Q3 remain open and Q4 partly proved. No shorter admitted lift chain, evaluated g\_j/e\_j/H or DP44 sign, specific032 value, prescribed global period, selected arithmetic nonzero image, or Kummer-Hecke RH has been supplied. The current paper keeps accepted and pending findings in separate parts.

Latest confirmed durable checkpoint is50. Exact immutable checkpoints51,52,53 await durable saving because the13:41:44 ordered retry failed at prepare\_uploads with library\_storage\_limit\_exceeded, before transfer or finalization. Their ZIP hashes and prior checkpoint50 receipt remain preserved; no earlier ZIP was overwritten, renamed, deleted or nested. Checkpoint54 will be captured at its scheduled slot and the exact51/52/53/54 files retried in order. Capture alone is not a successful durable save.

\clearpage

\section{Geometric dictionary and exact Grothendieck locators}\label{geometric-dictionary-and-exact-grothendieck-locators}

Author: \nolinkurl{/root}, GPT-6 Astra Pro; ALP\_MACHINE, 21 November 2026 UTC. This is a dictionary of the proved constructions and their remaining comparison problems. It does not attribute the new numerical formulas or theta identities to Grothendieck. Statements carried from the supplied runs keep their recorded intake status. The source-reading events below distinguish fresh primary inspection from carried references.

\subsection{Current additions through the selected even chain and raw finite datum}\label{current-additions-through-the-selected-even-chain-and-raw-finite-datum}

The following table extends the coverage below. Its locators refer to the exact primary passages already identified and qualified in this dictionary; it does not attribute the new calculations to those texts.

\begin{longtable}[]{@{}
  >{\raggedright\arraybackslash}p{(\columnwidth - 4\tabcolsep) * \real{0.3333}}
  >{\raggedright\arraybackslash}p{(\columnwidth - 4\tabcolsep) * \real{0.3333}}
  >{\raggedright\arraybackslash}p{(\columnwidth - 4\tabcolsep) * \real{0.3333}}@{}}
\toprule\noalign{}
\begin{minipage}[b]{\linewidth}\raggedright
Findings
\end{minipage} & \begin{minipage}[b]{\linewidth}\raggedright
Geometric object and exact comparison
\end{minipage} & \begin{minipage}[b]{\linewidth}\raggedright
Grothendieck framework and locator
\end{minipage} \\
\midrule\noalign{}
\endhead
\bottomrule\noalign{}
\endlastfoot
ALP-LIFT-23 & Exact Fourier contraction of the entire coupled Q theta-group family, supported on the cofactor line, with an explicit Gauss line and integral content & SGA7 I VII2.7/2.9, pp.169--174, retains the finite theta group and phase; SGA3 I4.7.3 organizes its character eigenspace. The finite character line has its actual arithmetic scalar, whose11 content and Weil height are computed. The finite Fourier contraction is a map of coefficient modules, not the subsequent global automorphic pushforward or its terminal mass. \\
ALP-LIFT-22 & The exact composition of two finite theta-section correspondences through their prescribed intermediate character module & SGA7 I VII2.7/2.9 retains the theta-group line and phase; SGA3 I4.7.3 and SGA6 VI3.3 organize the actual induced flag coefficients. The composition is a saturated integral injection after the explicit Haar scalar, a stronger assertion than an equality in K0. The global arithmetic row lattice and period still need comparison. \\
AM-K034 & A rigidified prime-to11 torsion Picard-zero line with its unique integral finite-holonomy connection, and equality of complete translated principal-parts modules & FGA236 Corollaries1.7/2.3, pp.229/232, and SGA7 I VII2.7/2.9 give the Picard/biextension setting; cote54 scans003/005/007 distinguishes universal-extension data from a flat structure on the ample line. EGA IV4 16.11.2/16.11.3(ii) retains ordinary differential factorials. The actual auxiliary ideal equality is proved by the integral torsion frame and independent jets, not by the virtual line class alone. \\
AM-R-018 & Exact dominant-weight extremes and finite numerical bounds on the actual selected coefficient bundles & SGA3 I4.7.3 organizes weights, SGA6 VI3.3 the flag/Schur bundle. The source spread includes its rank; this is explicit numerical representation data, not a consequence of the motivic weight. The same crystalline selected-ideal computation of K033 is compatible with the displayed finite source Sizes intervals, without a global height-window conclusion. \\
AM-R-017 & Explicit split prime divisors of a norm element, a projective residue flag with exact central rational stabilizer, and descent of an actual automorphic coefficient section & SGA3 I4.7.3/4.7.3.1 gives character gradings and SGA6 VI3.3 the flag coefficient module. The integral polarized lattice is in the FGA232 Theorem3.1, pp.148--149, setting. Its exact class mass retains the stabilizer28; the projective cover degree is an11-unit but has large real size. A finite-etale realization of the characters records their monodromy, not theta nonvanishing or a height-window theorem. \\
ALP-LIFT-21 & The actual finite theta-group section mapped through the open matrix-rank correspondence to a parabolically induced coefficient module & SGA7 I VII2.7/2.9, pp.169--174, retains the theta line and phase; SGA3 I4.7.3 and SGA6 VI3.3 organize the character and flag-module action. The normalized map is an actual saturated integral injection with the computed content and metric. Minguez's analytic boundary theorem is a separate input; the geometric coefficient map does not calibrate the global terminal mass. \\
AM-DO-035 & Exact local character cancellation, a nontrivial dilation eigensheaf on the projective residue orbit, and its integral finite Fourier theta section & SGA3 I4.7.3 retains the character grading; SGA7 I VII2.7/2.9 retains the Heisenberg line and phase. FGA236 Corollary2.3, p.232, and EGA IV4 Theorem16.11.2/Remark16.11.3(ii), pp.54--55, organize the Picard section and its ordinary jets. The raw and filtered norms and class indices are actual data. Linear lattice motion and affine theta translation are distinct operations on the section. \\
AM-K033 & Frobenius-coupled CM coordinates on the actual product principal polarization, a primitive Schur differential, and the exact valuation of its permitted section image ideal & Cote4 scan003 sections8.1.1--8.1.2/scans005--006 sections8.4.1--8.4.3 and cote54 scans005--007 give the crystalline/duality framework. EGA IV4 16.11.2/16.11.3(ii), pp.54--55, retains ordinary derivative factorials; SGA6 VI3.3 supplies the selected coefficient module; SGA7 I VII2.7/2.9 retains the primitive biextension. The sharp slope1/1220 is proved for this actual image ideal with fixed-baseline constants; the Newton polygon or virtual motive alone would not establish it. \\
AM-K032 & Positive weighted trace modules over two real quadratic integer rings, their actual product Picard polarization and a full integral CM action on J56 & FGA232 Theorem3.1, pp.148--149, and FGA236 Corollary2.3, p.232, give the polarized/Picard setting; SGA7 I VII2.7/2.9 retains the actual symmetric theta line. The weighted trace identity constructs the ppav itself. Its8-divisible real degree does not change the primitive Hodge-type common-field obstruction; the5/13 coordinate denominators and arch metric remain explicit arithmetic data. \\
ALP-LIFT-20 & The literal initial coupled arithmetic section as a finite theta-group element, its primitive affine coefficient vector and an exact finite Steinberg contraction & SGA7 I VII2.7/2.9, pp.169--174, retains the theta-group phases and line; SGA6 VI3.3 supplies the flag coefficient module and SGA3 I4.7.3 its character grading. The exact projector scalar is d\_St²/GL-order and can have11-adic denominators even though the raw section is primitive. This is a nonzero first contraction for the actual G factor, not a statement that later arithmetic pushforwards are saturated. \\
AM-K031 & Actual norm-pulled rank-one Hecke characters, their exact tame conductor divisor and integral Hodge grading, and a preserved polarized alpha-isogeny family & SGA3 I4.7.3 supplies character gradings; SGA7 I VII2.7/2.9 and FGA236 Corollary2.3, p.232, organize the symmetric theta line and allowed isogeny pullback. The numerical ranks and finite residue fields are explicit arithmetic data. The common CM field obstruction remains a property of the actual Hodge representation, not a defect removed by a virtual class. \\
AM-DO-034 & A new primitive Schur differential on the actual J14 Picard line, integral principal parts at11, and an explicitly nonzero automorphic image with its metric & EGA IV4 Theorem16.11.2/Remark16.11.3(ii), pp.54--55, distinguish ordinary derivative factorials; FGA236 Corollary1.7, p.229, gives the Picard setting; SGA6 VI3.3 supplies the flag/Schur module; SGA7 I VII2.7/2.9 retains the transported theta line. The actual42-block selected section survives the first map. Its hook metric, pi\_B coordinate scalar and ramified29 index remain additional data, and do not determine later integral images. \\
AM-R-016 & Orthogonal factorization of the polarized Hom space, conjugate and dual theta lines, and actual compact period contractions & SGA7 I VII2.7/2.9, pp.169--174, organizes the biextension and theta-group tensor products; SGA3 I4.7.3 retains the splitting-character grading. The final two period transfers are exact integrals with the same Haar measures, and at m2 they complete the six analytic predicates. No integral arithmetic row identity or bounded free-data witness follows from this tensor identity alone. \\
ALP-LIFT-19 & Distinct tame character lines, a finite induced coefficient module and an explicit primitive column-type vector with nontrivial metric & SGA3 I4.7.3 gives the grading and SGA6 VI3.3 the flag/representation-module setting. The Mackey summand is an actual submodule over the displayed coefficient ring; its Hermitian denominator is the calculated finite flag index. Its virtual K-class and its type occurrence do not identify the global arithmetic image lattice or the raw theta operator. \\
AM-K030 & An explicit trace-self-dual fractional ideal in the real degree-seven field, integral symmetric endomorphisms of the product ppav J\^{}14, and the actual polarized CM lattice & FGA232 Theorem3.1, pp.148--149, and FGA236 Corollary2.3, p.232, supply the Picard/abelian-scheme setting; SGA7 I VII2.7/2.9, pp.169--174, records the transported theta line and semicharacter. The inverse different, xi\_K=xi\_0*pi\_B\^{}(-6), trace-one diagonal, and integral11 eigenbasis are calculated explicitly. The crystalline direct sum survives this prime-to11 integral transport, while a selected highest-weight upper witness and the source's stronger host record require separate proofs. \\
AM-DO-033 & A primitive integral finite-Steinberg coefficient vector, its actual finite theta-group realization, and an explicit automorphic matrix-coefficient section with exact metric & SGA3 I4.7.3 organizes the finite character and grading; SGA6 VI3.3 supplies the representation-bundle setting; SGA7 I VII2.7/2.9 retains the exact theta-group line and finite phase. The class mass, Steinberg dimension and non-unit Hermitian norm are actual numerical data. A nonzero coefficient section is shown to survive this first theta operator, but no equality of virtual classes identifies the later prescribed image lattice. \\
AM-K029 & Descent of the induced angular CM character through every integral unit, extension from principal ideals, and exact correction of its quadratic restriction through an exponent-two capitulation kernel & SGA3 I Proposition4.7.3/Corollary4.7.3.1 supplies the character grading; FGA232 Theorem3.1, pp.148--149, and FGA236 Corollary2.3, p.232, supply the CM/Picard setting. The resulting normalized algebraic Hecke character has rank-one realizations; the unit, class-group and coefficient-field arguments are explicit arithmetic input. No effective height or number-field RH follows from the corresponding motivic or l-adic class. \\
ALP-LIFT-18 & Actual tame character lines on the norm-one torus with an exact integral-unit descent condition and fixed conductor divisor; a nonempty selected-weight locus & SGA3 I4.7.3 identifies the torus character lattice and gradings; SGA6 VI3.3 records the induced coefficient-bundle class. Its finite character is ramified on precisely the displayed G divisor. Cote4 scan012 section8.6 supplies a monodromy question, not a theorem converting this explicit tame character into an analytic nonvanishing bound. \\
AM-DO-032 & The source-defined terminal Fourier family, its finite character quotient/radical, a saturated integral coefficient submodule and a non-unimodular restricted pairing & SGA7 I VII2.7/2.9, pp.169--174, organizes the finite theta-group action and line; SGA6 VI3.3 supplies the representation-module language. The orthogonal projector denominator is additional integral metric data. Saturation in the coefficient lattice does not make the restricted pairing perfect, and its finite Gram is not the global arithmetic terminal mass. \\
AM-R-015 & The exact signed real weights, opposing character shifts, compact U(t) Levi and first-occurrence block of the actual vertical correspondence & SGA3 I4.7.3 and SGA6 VI3.3 organize weights and flag/representation classes; SGA7 I VII2.7/2.9 retains the polarized kernel line and its orientation. Li's regular real theta theorem, Atobe's first-occurrence theorem and Howe duality are separate analytic inputs. No equality in K-theory or of crystalline slopes substitutes for those representation-level correspondences. \\
AM-K028 & A prescribed infinity character fails to descend through the integral units; the actual CM Hodge representation of an ordinary J-power imposes the divisibility7 condition on multiplicity & SGA3 I Proposition4.7.3 and Corollary4.7.3.1 organize character gradings, while FGA232 Theorem3.1, pp.148--149, and FGA236 Corollary2.3, p.232, give the CM/Picard setting. The global idele-unit descent obstruction is proved explicitly here; it is not removed by passing to a virtual motive or by an integral change of theta frame. \\
ALP-LIFT-17 & An absent finite representation type in a characteristic-zero coefficient bundle, proved using flag stabilizers; exact tame conductor divisor & SGA3 I4.7.3 describes the character decomposition and SGA6 VI Theorem3.3 the flag/representation setting. The projector is zero on the actual finite flag modules, hence on all compact-generated spherical local modules. Its conductor records ramification of the corresponding rank-one character sheaf; the dictionary does not identify this divisor with the complete arithmetic height cost. \\
ALP-LIFT-16 & Actual unramified local tests, a full theta-group translate supported at29, exact principal-level degrees and arithmetic class masses & SGA7 I VII2.7 and2.9, pp.169--174, describe the theta-group/biextension lines; SGA3 I4.7.3 describes the character actions. The class index is an actual finite cover degree. Neither its local zeta test nor its unknown depth bound is a consequence of a virtual K-class. \\
AM-DO-030 & Simultaneous ordinary jets at finitely many distinct torsion sections, translated Picard-zero lines, and a finite theta-group operator with its exact norm & FGA236 Corollary1.7, p.229, and Corollary2.3, p.232, give the Picard setting; EGA IV4 Theorem16.11.2 and Remark16.11.3(ii), pp.54--55, distinguish these ordinary derivatives from divided differential operators; SGA7 I VII2.7/2.9 organizes the transported line and its finite phase. A new29 torsion choice is an actual change of section. \\
AM-DO-031 & Complete terminal finite Fourier datum, explicit lattice windows, a primitive coefficient vector away11 and a rational point with exact raw value & SGA7 I VII2.7/2.9 and VIII Proposition7.4(b), pp.310--311, in its stated extension scope. The finite group-algebra element, its character phase, and its integral content are retained; geometric nonzero sections need not survive a specified theta projection. \\
AM-K026--027 & Explicit full-Hyp3 Schur weights, exact hook metrics, and the common-weight crystalline lower bound on the actual formal theta line & SGA6 V and VI Theorem3.3 give the Schur and lambda/K-theory language; SGA3 I4.7.3 its weight grading. Cote4 scan003, sections8.1.1--8.1.2, scans005--006, sections8.4.1--8.4.3, and cote54 scans005--007 give the crystalline/duality setting. The selected slope1/1220 and retained factorial cost are calculations in the actual line and integral differential module, not just its Newton polygon. \\
AM-R-014 & Genuine rational Witt-lower Hermitian spaces, parabolic constant terms and stability of full smooth real isotypic realizations & Reductive groups and their character gradings belong to SGA3 I4.7.3; coefficient-bundle classes to SGA6 VI3.3; theta lines to SGA7 I VII2.7/2.9. Wu's Proposition3.1 is the separate analytic tower theorem. None of these framework identifications supplies the remaining source identities or arithmetic host pins. \\
AM-K019 & Quasi-periodic theta logarithmic derivatives and the polynomial auxiliary-variable connection; a full ordinary differential lattice has unbounded creation powers & EGA IV4, Theorem16.11.2 and Remark16.11.3(ii), pp.54--55, distinguish divided differential operators from the ordinary Weyl action. Cote54 scans003,005,007 supply the universal extension and duality setting, not a nilpotence theorem for this entire connection. \\
AM-K020; ALP-LIFT-13 & Exact evaluated theta derivatives and the primitive direction at the height-three primes3 and5 & Cote4 scan003, section8.1.1--8.1.2; scans005--006, section8.4.1--8.4.3; cote54 scans005--007. The computed scalar slopes concern the evaluated section, not just the isocrystal Newton slopes. \\
AM-R-009--010 & A principal theta line on the actual product A=J squared, its E-prime action, norm-one torus characters, and permitted odd-isogeny translates & FGA232 Theorem3.1, pp.148--149; FGA236 Corollaries1.7 and2.3, pp.229/232; SGA7 I VII2.7 and2.9.3, pp.169--174. The explicit period and its Gaussian metric are additional analytic data. \\
ALP-LIFT-11--12 & Tensoring the automorphic coefficient bundle by an actual unitary determinant line, then a nonzero second theta integral & SGA3 I Proposition4.7.3 and Corollary4.7.3.1 describe the character grading; SGA6 VI Theorem3.3 describes the corresponding representation/flag class. Rallis nonvanishing and its measures are separate automorphic input, not consequences of K-theory. \\
AM-DO-025--026 & Selected integral highest weights, all-polynomial theta-group kernel, real covariance, adjoint, and weighted class matrix & SGA3 I Proposition4.7.3; SGA6 V lambda-operations and VI Theorem3.3; SGA7 I VII2.7/2.9 for the biextension line. The arithmetic terminal matrix is an actual map with a metric and cannot be replaced by its virtual class. \\
AM-DO-027 & An allowed scalar normalization of the literal arithmetic datum; equal-volume chart phases are integral roots of unity & Rank-one character bundles in SGA3 I4.7.3 and SGA6 VI3.3. This is an actual integral isomorphism of each existing input line. It does not construct a missing raw arithmetic section or terminal functional. \\
AM-R-011 & Integral Schur modules, determinant insertion, hook-normalized Fischer metrics and a rational central idempotent with an explicit denominator & SGA6 V and VI3.3, with the carried-source status stated below. Over Z, the denominator primes matter; characteristic-zero decomposition alone does not provide an integral splitting. \\
AM-K022; AM-R-012 & Ordinary principal parts at the zero section, integral jet surjectivity at explicit level, determinantal ideals of the allowed section map, and a concrete first automorphic class function & EGA IV4 section16, especially Theorem16.11.2 and Remark16.11.3(ii); FGA236 Corollary1.7; SGA7 I VII2.7/2.9. The finite matrix uses ordinary jets; evaluating a derivative introduces the displayed factorial. The class-function mass is additional adelic data. \\
AM-K021 & A finite Artin character component, CM-type height classes factoring through a central-conjugation quotient, and a ray-theta Mellin period carrying L-prime/L & The Artin projector is a finite-etale, weight-zero motivic summand and can also occur inside a CM Hom motive. SGA4--5 supply its sheaf/trace language; the finite-field weight results cited below do not give the required arithmetic height comparison or number-field RH. \\
AM-DO-028; ALP-LIFT-14 & The polarized Hom space as a tensor product with its dual skew-Hermitian factor; determinant character transport; actual selected nonzero theta iteration & The character and dual-representation operations lie in SGA3 I4.7.3 and SGA6 VI3.3. The tensor/Hom symplectic map retains an actual biextension normalization in SGA7 I VII2.7/2.9. The Rallis formula, its analytic L-values and Haar factors are additional automorphic theorems, not conclusions of these geometric frameworks. \\
ALP-LIFT-15 & Full finite theta-group eigenspaces, their orthonormal Weyl basis, and one nonzero translated section at each analytic edge & SGA7 I VII2.7, pp.169--170, and VII2.9.1/Proposition2.9.3, pp.172--174, organize the biextension and theta-group lines; FGA236 Corollary2.3, p.232, supplies the Picard-functor setting. The varying finite level, norm and orbit index must be retained in any arithmetic comparison. \\
AM-R-013 & Algebraic tensor-factorized theta images, branching of the compact coefficient bundle to an actual codimension-one Hermitian subspace, and its positive restriction pairing & SGA3 I4.7.3 and SGA6 VI3.3 give the representation-bundle language. FGA232 Theorem3.1, pp.148--149, and SGA7 I VIII Theorem7.1, pp.299--300, describe the polarized/Picard and extension settings with their stated hypotheses. They do not identify the analytic restriction norm with the source's arithmetic terminal functional. \\
AM-DO-029 & A finite theta-group Fourier expansion of the literal nonterminal raw arithmetic Xi, its support/period lattices, exact phases and norm & SGA7 I VII2.7/2.9, with VIII Proposition7.4(b), pp.310--311, in its stated fully faithful extension scope. The explicit integral group-algebra element and its denominators are stronger data than the class of its line. Its selected jet evaluation can still vanish. \\
AM-K023; AM-K025 & The full integral formal section module, restricted differential operators, the finite mod11 evaluation quotient and adjacent-gap factorial content & Cote54 scans003,005,007 and cote4 scan003 section8.1.1--8.1.2, scans005--006 section8.4.1--8.4.3 supply the universal-extension/crystalline framework. EGA IV4 Theorem16.11.2 and Remark16.11.3(ii), pp.54--55, distinguish the ordinary derivative factorials from divided operators. The finite quotient and its unit Smith rank are computed here; no higher divided-lattice stability is inferred. \\
AM-K024 & A primitive integral polynomial already in its selected source/target Schur summand, with exact ordinary-CM gauge and permitted jet realization & SGA6 V and VI3.3 describe its Schur class; SGA3 I4.7.3 its highest-weight grading. EGA IV4 section16 and FGA236 Corollary1.7, p.229, organize the actual finite principal-parts map. Its 1/1220 upper local slope is a computed evaluation property, not a consequence of a virtual representation identity. \\
\end{longtable}

For every entry involving a metric, the period, archimedean norm, finite lattice and actual map remain explicit data. A statement in a representation lambda-ring, a Dieudonne crystal, or an l-adic Grothendieck group does not determine those missing normalizations.

\subsection{1. What the archive pages actually supply}\label{what-the-archive-pages-actually-supply}

\textbf{Cote 54.} The inspected manuscript concerns the universal vector extension of an abelian variety, its de Rham realization, Cartier duality, and the restricted Lie structure in characteristic \(p\). In the official PDF, scan 003 (manuscript page 2) displays the Hodge exact sequence and the vector extension; scan 005 (manuscript page 4) discusses the Cartier dual construction; scan 007 (manuscript page 6) treats Frobenius, Verschiebung, and the restricted operation on the vector kernel. The zero restricted operation there belongs to that vector kernel. It is not a statement that every invariant derivation of an abelian variety has zero \(p\)-th power. AM-K013 and AM-K017 compute precisely the different restricted powers that occur on the actual Jacobian.

\textbf{Cote 4.} Scan 003 (printed page 2), §8.1.1, relates crystalline and de Rham cohomology through the stated derived comparison; §8.1.2 records the associated Tor issue. Scan 005 (printed page 4), §8.4.1--2, singles out the Frobenius-bijective stable part. The manuscript's terminology ``semi-simple'' in this passage refers to that Frobenius-stable part; it must not be substituted for ordinary linear diagonalizability. Scan 006 (printed page 5), Lemma 8.4.3, compares the stable parts of de Rham and coherent cohomology and uses the Frobenius action on the filtration. The subsequent §8.5 contains conditional ingredients. Scan 012 (printed page 11), §8.6, asks a monodromy question for the indicated finite-field local systems and explains a possible trace/congruence consequence. It is a proposed route in the manuscript, not a theorem about all number-field Hecke \(L\)-functions.

These are the exact archive loci used here. Their content organizes the realization and monodromy questions; no newly computed theta slope or automorphic nonvanishing is inferred from an archival question.

\subsection{2. Formal groups, crystals, and integral differential operators}\label{formal-groups-crystals-and-integral-differential-operators}

Over a finite extension where the Jacobian has good reduction, its \(p\)-divisible group and Cartier dual determine dual Dieudonné objects. The accepted DP\_MULT DP30--34 and DP49 distinguish the one-dimensional height-three factor from its two-dimensional dual. The three labelled tangent directions are not three independent height-three formal lines. Frobenius permutes the CM labels and Verschiebung controls the restricted derivations. This is why the restricted nilpotence exponents in AM-K013 are direction dependent.

The exact finite-level framework is SGA 3, Exposé VII\_A, §§6.2--6.4, 7.2--7.4, and 8.1, as carried in DP\_MULT. For the de Rham/finite-group comparison, the carried primary locators are Oda (1969), Theorem 3.19, p.~89, and Corollary 5.11, p.~131. DP\_MULT's source audit additionally fixes Tate (1967), §2.2, Proposition 1, p.~162, equation (3), p.~163, equation (4), p.~164, and §2.3, Proposition 3, p.~166; Zink's dual display operators are Definition 19, equations (22)--(24), pp.~143--144, and the formal-biextension injection is Proposition 121 and equation (243), pp.~241 and 246. Those established results are cited, not proved again in this run.

The new theta statements concern the completed local rings and actual sections of a polarized abelian scheme. In AM-K012 the explicit ordered divisor chart is étale on a dense open and supplies a ring stable under the invariant derivations. AM-K017 extends a logarithmic Hessian across the theta divisor by checking every horizontal and vertical divisor and using normality. It then uses properness to turn a global regular function on the special fibre into a constant. This is a scheme-theoretic argument on the actual integral model. It does not replace a missing integral reduced theta frame by a complex analytic one.

The resulting Hurwitz series use the ordinary invariant operators (D\_i\^{}n), with factorials in the multiplication law. Their reductions and restricted-power identities are explicit. A Chern connection on an ample theta line is not thereby made into a global flat algebraic connection or a Dieudonné crystal. To make that additional assertion one would have to construct the relevant stratification, integrable connection, and Frobenius structure. The present proof does not do so.

For the oscillator, the same distinction has an exact EGA locator: EGA IV\_4, Theorem 16.11.2, p.~54, and Remarks 16.11.3(ii), p.~55. Over an integral base the full sheaf of differential operators contains divided derivatives and is generally larger than the algebra generated by coordinates and ordinary derivations. AM-DO-003 constructs an integral Weyl-algebra module, becoming the familiar delta \(D\)-module after rationalization. It does not claim an integral action of the full EGA differential-operator sheaf.

\subsection{3. Picard schemes, the theta line, and finite models}\label{picard-schemes-the-theta-line-and-finite-models}

The Jacobian is a Picard variety. The precise carried FGA references are Bourbaki exposé 232, Theorem 3.1, pp.~148--149; exposé 236, Corollary 1.7, p.~229, Corollary 2.3, p.~232, and Proposition 2.10(ii), p.~236. They give the representability and base-change framework used in DP\_MULT DP22 and DP49. The actual divisor and section calculations in the present run are additional calculations on that Picard scheme.

ALP-LIFT-07 distinguishes the non-effective characteristic (D\_\{135\}) from the invariant trace-one theta characteristic. They differ by a nonzero element of (J{[}2{]}). This is a distinction in the relative Picard scheme, so an analytic normalization or a scalar change cannot remove it. Pullback by an odd integer preserves the nontrivial two-torsion difference. AM-DO-015 instead uses the specified Mumford translation; AM-R-006 computes its actual scalar.

ALP-LIFT-08 computes the leading determinant section on an integral Picard chart, including residue characteristic two. Its first nonzero jet is the determinant of the cup-product matrix, (u\_1u\_3-u\_2\^{}2). The Hodge-basis transformations in ALP-LIFT-05 then determine the vertical contents at 2 and 7. The objects are line bundles, their primitive sections, and their integral jet images over specified discrete valuation rings. An equality after tensoring with the fraction field would not determine these contents.

AM-R-004--006 and ALP-LIFT-02--06 give the corresponding Betti lattice, de Rham basis, complex metric, determinant metric, and exact scalar comparison. Thus the height identities live on the metrized Hodge determinant line, with the finite integral lattices retained. The archimedean Cauchy estimate and the finite derivative estimate are different realizations of the same normalized section; neither may be changed independently in the cost calculation.

\subsection{4. Biextensions and the phase problem}\label{biextensions-and-the-phase-problem}

The polarized Poincaré line on \(J\times J^\vee\) is a \(\mathbf G_m\)-biextension, rigidified along both zero sections. Pullbacks, duality, and transposition organize the theta-group and trace-one constructions. The exact inspected SGA 7 I loci are:

\begin{itemize}
\tightlist
\item
  Exposé VII, §2.7, printed pp.~169--170: transposition of a biextension. These are PDF pages 174--175 in the original scan.
\item
  Exposé VII, equation (2.9.1) and Proposition 2.9.3, printed pp.~172--174: the forgetful functor to birigidified torsors and its fully faithful form under the displayed hypotheses. Example 2.9.5, p.~175, gives the abelian/Picard setting.
\item
  Exposé VIII, Theorem 7.1(a--b), printed pp.~299--300: extension over a trait. In part (b), the component groups enter a canonical obstruction pairing into \(\mathbf Q/\mathbf Z\); under the stated bounded-torsion hypothesis its vanishing is equivalent to extension. The connected-fibre case is explicitly separated.
\item
  Exposé VIII, Proposition 7.4(b), printed pp.~310--311: full faithfulness over a reduced base for the indicated smooth group schemes with connected generic fibres; over a normal base, with one group having connected fibres, the stated generic-fibre criterion describes the essential image.
\end{itemize}

The local finite Gauss factors of AM-DO-017 are additional data in a specified analytic/Heisenberg realization. The uniqueness of a birigidified biextension structure does not say that every chosen intertwiner of its realizations has scalar one. AM-DO-020 proves that scalar assertion at relative split places by evaluating positive normalized lattice indicators. AM-DO-021 exhibits the actual residual compact sign at the ramified place 7. This is a normalization issue inside an existing realization, not a counterexample to the SGA 7 existence or uniqueness theorems.

The common period in AM-DO-018 is retained for every moved lattice. Its coefficient denominators are intersection indices, and its norm is computed with specified probability Haar measure. In geometric language the Heisenberg eigenspaces are the one-dimensional section spaces of principal theta lines on isogenous complex tori, compared over common finite covers. Their phase and arithmetic integral structure are separate from the isomorphism class of the line bundle.

\subsection{5. Group schemes, weights, and the oscillator}\label{group-schemes-weights-and-the-oscillator}

The exact fresh SGA 3 locator for a diagonalizable-group action is Exposé I, Proposition 4.7.3; Corollary 4.7.3.1 gives the corresponding graded-algebra/equivariant-affine-scheme statement. Accordingly the charges in T-052 are actual characters of a torus, and its degree pieces are equivariant locally free modules. The nonnegative equal-charge construction of the accepted T-053 is a specified graded coordinate subalgebra. It is not merely a dimension count.

AM-DO-002 and AM-DO-007 compare the actual character weights of the canonical ball representation with the required polynomial degrees. AM-DO-006 proves an obstruction for the entire creation-operator weight set. AM-DO-003 changes polarization by a Fourier/Weyl algebra map and records why the resulting delta module is not the positive Fock vacuum. AM-DO-005 supplies a selected bicanonical comparison with an exact local Haar pairing, so the obstruction must not be generalized to every representation or determinant twist.

AM-R-007 belongs to the reductive-group side: the exact determinant sequence separates (SU\(V\)) from the norm-one torus in (U\(V\)), and strong approximation gives representatives for a Hermitian lattice genus with level structure. SGA 3 supplies the group-scheme framework; the approximation theorem actually used is Rapinchuk (2013), Theorem 2.3, p.~281. The source locator is not relabelled as a theorem of SGA 3. ALP-LIFT-09 then constructs an actual automorphic character and theta integral on the norm-one torus. That analytic integral requires its Haar measure and splitting, beyond the algebraic group scheme and its representation class.

\subsection{6. Lambda rings, K-theory, and what a class does not determine}\label{lambda-rings-k-theory-and-what-a-class-does-not-determine}

For a finite locally free representation \(V\), its exterior powers define the lambda operations on its exact Grothendieck group. The carried exact locators are SGA 6, Exposé V, §2 (lambda operations), §3 (gamma operations and filtration), Proposition 4.2, Definition 4.4, Proposition 4.8, and Corollary 4.9 (the universal representation-ring description); Exposé VI, Theorem 3.3, is the associated Schur-functor framework used by GCT\_PLETH. Cote 133, scan 006, records the archival pre-lambda/classifying-topos framework, as carried in that run's source study. The SGA 6 download attempt in this run failed; these numbered references remain carried references, not fresh primary inspections.

Each finite degree of the oscillator polynomial algebra has a class formed by symmetric and Schur operations. A tensor-to-Hom map such as AM-DO-009 induces the corresponding equality of representation classes after dualizing the source factor. The map itself retains the contragredient action, the alternating form and the determinant roots. The representation class alone forgets the positive metric, the integral section lattice, finite Weil phases and the actual map on automorphic forms.

Likewise AM-R-002's content ideal is data of an actual integral linear map applied to an actual vector. Its least denominator and dual metric can be expressed through determinant lines and metrized lattices; they are not recovered from a virtual rank. AM-R-001's cost subtraction and ALP-LIFT-06's directional cancellation compare those retained numerical realizations. No lambda-ring identity is used to erase a height, a covolume or an endpoint denominator.

\subsection{7. CM motives, Artin motives, and the RH boundary}\label{cm-motives-artin-motives-and-the-rh-boundary}

AM-R-003 uses the primitive CM type to separate the simple rational weight-one Hodge structure of (J) from an elliptic induced type. The corresponding CM realizations are those of (h\^{}1(J)), with its endomorphism action and polarization. Products and field extensions do not turn that simple isotypy class into the source's elliptic Serre-tensor class. AM-DO-010 instead identifies specified CM eigensummands with actual Galois twists up to isogeny. It does not identify their integral models without the separate lattice data.

For AM-K001--004 the order-(r) Kummer character is a rank-one Artin representation over the cyclotomic base. Over the open arithmetic base excluding its ramified primes it is a lisse rank-one sheaf with finite monodromy. Its induced representation is the stated augmentation representation over \(\mathbf Q\). The exact representation identity therefore carries the ramified Euler factors as recorded in the proof. For (r=3), the explicit isomorphism of two CM elliptic curves over the Kummer extension identifies the character with a line in their Hom realization. This gives an actual descent cocycle, rather than an analogy based on dimension.

The proper/smooth realization framework is SGA 4, Exposé XII, Theorems 5.1--5.2; Exposé XV, Theorem 2.1; Exposé XVI, Theorem 2.1, and SGA 5, Exposé VI, §§1.2.1, 1.4.2--1.4.3, pp.~265--266, as carried in CX\_GRH/GRH\_NIGHT. For finite-field weights the exact carried locators are Deligne's \emph{Weil I}, Theorem 1.6, p.~276, and \emph{Weil II}, Corollary 3.3.9, p.~207. They control Frobenius eigenvalues in their stated finite-field cohomological settings. In particular an Artin line has weight zero. These statements do not put the zeros of its number-field Hecke \(L\)-function on the critical line.

AM-K005 isolates the remaining analytic implication through an explicit every-epsilon arithmetic window. AM-K007--011 supply ray characters, splitting twists and CM hosts while retaining ramified conductors, prime norms and host costs. Cote 4 §8.6's monodromy question is relevant background for such realization questions; it does not supply the missing arithmetic window. The accepted Dirichlet theorem of Alpöge remains an input with its original scope. No statement of Masullo's draft is used as an unconditional theorem.

\subsection{8. Coverage of findings and the programme's rows}\label{coverage-of-findings-and-the-programmes-rows}

\begin{longtable}[]{@{}
  >{\raggedright\arraybackslash}p{(\columnwidth - 4\tabcolsep) * \real{0.3333}}
  >{\raggedright\arraybackslash}p{(\columnwidth - 4\tabcolsep) * \real{0.3333}}
  >{\raggedright\arraybackslash}p{(\columnwidth - 4\tabcolsep) * \real{0.3333}}@{}}
\toprule\noalign{}
\begin{minipage}[b]{\linewidth}\raggedright
Findings in this run
\end{minipage} & \begin{minipage}[b]{\linewidth}\raggedright
Concrete geometric content
\end{minipage} & \begin{minipage}[b]{\linewidth}\raggedright
Original row and remaining boundary
\end{minipage} \\
\midrule\noalign{}
\endhead
\bottomrule\noalign{}
\endlastfoot
AM-R-001--002; ALP-LIFT-01 & Metrized integral row/endpoint maps, their content ideals, and exact changed-family accounting & DP\_MULT DP44, DP52, Q3e; the actual terminal vector and all costs remain required. \\
AM-R-003; ALP-LIFT-02--06 & Primitive CM isotypy; polarized Hodge lattices and determinant metrics; finite normalization and directional balance & DP\_MULT DP49, DP57, DP64, Q3e; a bare-certificate equality is not a selected-gain upper bound. \\
AM-R-004--006; ALP-LIFT-07--08; AM-DO-011--016 & Actual Picard/theta lines, characteristic translation, integral leading determinant, periods, scalar contents and archimedean bounds & DP\_MULT DP22, DP49, DP54--65; source-line identification and normalization are explicit. \\
AM-K012--018 & Integral Abel charts, algebraic logarithmic Hessians, restricted powers, explicit Hurwitz jets, growth, mixed-jet counterexamples, and exact full-CM canonical-gauge matching & DP\_MULT Q3e and the local step-2 input; the two fixed mixed-line slopes now apply to the actual section, while conjugate-action products retain a separate gauge condition. \\
AM-DO-017--022; AM-R-007 & Principal theta eigenspaces on commensurable lattices, exact normalized averages, split-place phase, ramified compact character, and all-finite phase repair over the unramified CM host & THETA3 G1 and DP\_MULT Q3e; full arithmetic rows, later images and terminal mass remain distinct equations. \\
ALP-LIFT-09--10 & Genuine first theta integrals from an anisotropic torus, with actual splitting, permitted half-torsion orbit, nonvanishing, exact finite-index input norms and character projections & DP\_MULT Q3e; later definite lifts and the prescribed chain are not supplied by this first image. The first-lift host E and the repaired all-finite host E-prime still require a common-tower comparison. \\
AM-DO-001, 004, 009--010 & Signed Hermitian spaces, positive multiplicity space, exact tensor-to-Hom map and CM isogeny realization & THETA3 A7 and G1; the common positive arithmetic pin and global MainX value need their own comparison. \\
AM-DO-002--003, 005--008 & Equivariant oscillator modules, integral Weyl/Fourier map, degree and charge obstructions, selected bicanonical pairing, finite scalar equations & GCT\_PLETH T-052/T-053, Q-ALPOGE; ALPOGE\_GCT Q-ALPOGE-METHOD; scoped obstructions coexist with a selected local comparison. \\
AM-K000--006 & Artin--Schreier finite-field construction, Kummer Artin motive, cubic conductor, CM Hom line, analytic pole/window reduction & CX\_GRH CX-ALG-024; the general Kummer RH and factoring implication remain conditional on the missing window. \\
AM-K007--011 & Explicit finite etale/ray characters, determinant-root twists and arithmetic CM-host lattices & CX\_GRH and GRH\_NIGHT Kummer gate; relative prime/conductor and quantitative tower adapters remain. \\
\end{longtable}

New findings added after this version must be added to this table before the final package. The table is a coverage record, not a new proof of the findings it indexes.

\subsection{References and source status}\label{references-and-source-status}

Grothendieck, A. (n.d.). \emph{Manuscript cote 4} {[}Digitized manuscript{]}. Archives Grothendieck, Université de Montpellier. \url{https://grothendieck.umontpellier.fr/4.pdf} . Scans and paragraph locators are listed in §1; primary PDF and the relevant page images were inspected.

Grothendieck, A. (n.d.). \emph{Manuscript cote 54} {[}Digitized manuscript{]}. Archives Grothendieck, Université de Montpellier. \url{https://grothendieck.umontpellier.fr/54.pdf} . Exact used scans 003, 005, 007; primary inspection in this run.

Grothendieck, A. (1962). \emph{Techniques de construction et théorèmes d'existence en géométrie algébrique V--VI: Les schémas de Picard}. Séminaire Bourbaki, exposés 232 and 236. Exact theorem/page locators in §3, carried through the supplied DP\_MULT DP22 and its intake.

Grothendieck, A., with Dieudonné, J. (1967). Éléments de géométrie algébrique IV, quatrième partie. \emph{Publications Mathématiques de l'IHÉS, 32}, 5--361. Exact used locus IV\_4, 16.11.2--16.11.3(ii), pp.~54--55; freshly read by the oscillator agent.

Demazure, M., \& Grothendieck, A. (Eds.). (1970). \emph{Schémas en groupes (SGA 3), Tome I} (Lecture Notes in Mathematics, Vol. 151). Springer. Exposé I, 4.7.3 and 4.7.3.1 were freshly read in the recomposed French text; VII\_A locators in §2 are carried.

Berthelot, P., Grothendieck, A., \& Illusie, L. (Eds.). (1971). \emph{Théorie des intersections et théorème de Riemann--Roch (SGA 6)} (Lecture Notes in Mathematics, Vol. 225). Springer. Exact carried lambda/K-theory locators in §6; no successful fresh PDF retrieval claimed.

Grothendieck, A., Raynaud, M., \& Rim, D. S. (1972). \emph{Groupes de monodromie en géométrie algébrique (SGA 7 I)} (Lecture Notes in Mathematics, Vol. 288). Springer. Exact freshly inspected VII--VIII locators in §4. The original-scan PDF has five front-matter pages: PDF page number equals printed page plus five at these loci.

Oda, T. (1969). The first de Rham cohomology group and Dieudonné modules. \emph{Annales Scientifiques de l'École Normale Supérieure, 2}(1), 63--135. Theorem 3.19 and Corollary 5.11 are used through the accepted DP\_MULT realization results.

Deligne, P. (1974). La conjecture de Weil. I. \emph{Publications Mathématiques de l'IHÉS, 43}, 273--307. Theorem 1.6, p.~276.

Deligne, P. (1980). La conjecture de Weil. II. \emph{Publications Mathématiques de l'IHÉS, 52}, 137--252. Corollary 3.3.9, p.~207. Both weight references are carried framework inputs, with their finite-field scope retained.

\subsection{10. Additional checked pairing and twist refinements}\label{additional-checked-pairing-and-twist-refinements}

AM-DO-024 supplies a positive comparison between the actual finite Weil translate and the direct Heisenberg average on the E-prime host, together with the exact probability-Haar pairing and adjoints. These are identities of metrized section lines in the biextension/Heisenberg realization of section4, stronger data than the isomorphism class of a Picard line. They do not determine an automorphic quotient mass. AM-R-008 supplies actual integral diagonal maps for tensoring compatible coefficient modules by a unitary determinant character, within the SGA3 character and SGA6 representation K-theory frameworks of sections5--6. Its all-prime content preservation comes from root-of-unity scalars, not from equality in K0. The first-lift target-level membership and its refinement cost remain additional obligations. AM-DO-023 is held pending reconciliation of the general character convention; its unchanged-character local vanishing has a separate narrower scope.

\subsection{11. Unchanged first raw kernel and its mass}\label{unchanged-first-raw-kernel-and-its-mass}

AM-DO-036 proves positivity for the actual unchanged first raw theta tensor on the principally polarized J14/J56 hosts. Its finite lattice realization is the section/Heisenberg construction organized by the rigidified Poincare biextension of SGA7 I, VII2.7/2.9 (pp169--174) and VIII7.4(b) (pp310--311), and the Picard representability interfaces FGA232 Theorem3.1 (pp148--149) and FGA236 Corollary1.7 (p229),2.3 (p232). Those scheme results provide the object and functoriality, not this analytic positivity. The latter is a convergent doubled theta norm calculation with anisotropic target, using GQT7.1 (p33); its target quotient mass remains an explicit scalar. It therefore removes the first fixed-kernel vanishing gate of DP\_MULT Q3e and narrows THETA3 G1, without supplying an integral later image lattice or terminal mass identity. The common CM weight-one motive is retained; its Dieudonne and Cartier-dual realization remains the cote4/cote54 correspondence recorded in sections2--3. No finite-field weight statement or equality in representation K-theory supplies the missing global mass or number-field RH.

\subsection{12. Mixed periods, embedding ideals and local quotient scalars}\label{mixed-periods-embedding-ideals-and-local-quotient-scalars}

AM-K035 expresses the polarized Picard variety through its Betti/de Rham comparison with an actual per-embedding period vector. Its coordinate ideals are integral module intersections over Spec(O\_M), and its reduced theta jets are formal sections of the rigidified line in the Picard/biextension setting of FGA232 Theorem3.1 (pp148--149), FGA236 Corollary1.7 (p229),2.3 (p232), SGA7 I VII2.7/2.9 (pp169--174) and VIII7.4(b) (pp310--311). The all-order algebraic formal coordinates use the EGA IV4 16.11.2--16.11.3(ii) interface (pp54--55) already recorded above. The explicit analytic period and unit-Fock constants are new comparisons, not consequences of representability alone. At11 its cD is a unit, so it preserves the previously proved Dieudonne/Cartier-dual selected-jet content; cote4 and cote54 scans5--7 remain the exact historical framework locators of sections2--3.

ALP-LIFT-24 and AM-DO-037 calculate actual local quotient maps and their Haar/L/Gauss scalars. Their finite inputs are section representations of the same Heisenberg extension; the literal DegOne condition is the statement e=f=1 for the relevant finite localization. An identity in the representation lambda-ring or K-group (SGA6 VI3.3) records types but does not determine these periods, integral content or the global terminal mass. Likewise finite-field weights of the CM motive do not supply number-field RH for the Kummer character. These scope distinctions carry forward DP\_MULT Q3e and THETA3 G1 without changing CX\_GRH CX-ALG-024.

\subsection{13. Fixed-edge descent and compact contractions}\label{fixed-edge-descent-and-compact-contractions}

K036 is a descent statement for complete formal sections and their torsion translations on the fixed Picard/biextension family. It uses the exact algebraic rational-isometry scalar and finite-class reduction, not representability alone; its FGA/SGA7/EGA locators are those in sections4 and12. The explicit nc measures a finite isogeny/Heisenberg induction coefficient, while Cartier/Dieudonne slopes (cote4; cote54 scans5--7) concern the separate formal section content. LIFT25 computes an actual contraction in the compact representation category and its Fock metric, refining the SGA3 I4.7.3/4.7.3.1 and SGA6 VI3.3 class information. The dimension denominator and invariant-tensor norm are distinct. No virtual representation or motivic equality erases these scalar costs.

\subsection{14. Integral branching with its metric}\label{integral-branching-with-its-metric}

R019 makes the SGA3/SGA6 representation/flag restriction of section13 quantitative: the actual integral minor vector, spectral integer and metrized branch have exact factorial expressions. SGA3 I4.7.3/4.7.3.1 and SGA6 VI3.3 organize the representations, while Molev Proposition2.4 supplies the classical contravariant norm input and the new path calculation fixes its source normalization. The integer scalar is not an integral splitting theorem for every chosen lattice. The Picard/biextension and cote4/cote54 formal realizations remain those already identified; no new motivic or number-field RH consequence follows from this metric refinement.

\subsection{15. ALP-LIFT-26: fixed coupled analytic image}\label{alp-lift-26-fixed-coupled-analytic-image}

The two selected correspondences have a nonzero composition on their actual complex automorphic realization, with prescribed local sections and every multiplicity copy retained. In the Picard/biextension interpretation of sections12--14, the tensor identity uses the same transported line-bundle frames; its character and volume factors are part of the realization, not an identification of integral lattices. The exact Grothendieck locators remain FGA232 Theorem3.1 pp148--149 and236 Corollary1.7 p229/2.3 p232 for Picard geometry, SGA7I VII2.7/2.9 pp169--174 and VIII7.4(b) pp310--311 for biextensions, and EGAIV4 16.11.2--3(ii) pp54--55 for jets. None of these citations identifies the image of the arithmetic class map or its terminal branching projection. In lambda/K-theory the selected Schur factors are the same factors already recorded in section14; in motivic or ell-adic language no new algebraic correspondence comparison or RH statement follows from this analytic nonvanishing theorem.

\subsection{16. AM-K037: integral polynomial functor and class-map realization}\label{am-k037-integral-polynomial-functor-and-class-map-realization}

The map CR7 is a morphism of the specified integral Weyl polynomial functors over Z; the monic-determinant descent proves it extends from GL to the full matrix scheme on the relevant image. Its evaluation against the independently fixed algebraic theta jet is the actual arithmetic class map. This places the new fact directly in the polynomial-functor/lambda-ring language of the earlier SGA6 VI3.3 comparison, while its jet realization uses EGAIV4 16.11.2--3(ii), pp54--55. The fixed line bundle and reduced theta section retain the Picard and biextension locators of sections12--15 (FGA232/236 and SGA7I VII--VIII). Coevaluation is the usual duality morphism for the selected coefficient representation; its norm and finite stabilizers are part of the realization and are not removed by the K-theory description. No new identification of an integral motivic image, ell-adic weights or Kummer RH is claimed.

\subsection{17. AM-DO-038: terminal contraction and mass}\label{am-do-038-terminal-contraction-and-mass}

The fixed branch is a morphism between the selected coefficient realizations; the finite class pairing contracts it with the dual row and weights each class by its actual finite stabilizer. The top automorphic realization realizes this same pairing and introduces the exact adelic level volume. This is a comparison at the complex/arithmetic coefficient level, not a new identity of schemes or motives. Its polynomial/jet part has the same SGA6 VI3.3 and EGAIV4 16.11.2--3(ii) pp54--55 locators; its fixed line/frame transports retain FGA232 Theorem3.1 pp148--149, FGA236 Corollary1.7 p229 and2.3 p232, and SGA7I VII2.7/2.9 pp169--174 andVIII7.4(b) pp310--311. Stabilizer cancellation and compact Schur integration do not imply nonzero algebraic cycle classes or ell-adic purity. The latter realizations require the independently stated correspondence comparison, still open in THETA3 G1.

\subsection{18. AM-R-020: integral Weyl-image sandwich}\label{am-r-020-integral-weyl-image-sandwich}

The new inclusions are statements about finite integral representation modules over Spec(Z) and their localizations, and they survive extension to Spec(O\_M). They refine the rational restriction/Schur-functor information organized by SGA3I4.7.3/4.7.3.1 and SGA6VI3.3: equality of rational representation classes does not determine an integral image or its elementary divisors. The exact finite-place bounds use the specified source divided-power lattice and normalized duality. Their application to the Picard/biextension theta section retains the existing FGA232 Theorem3.1 pp148--149, FGA236 Corollary1.7 p229/2.3 p232, SGA7I VII2.7/2.9 pp169--174 andVIII7.4(b) pp310--311 locators; R020 adds no claim about that earlier image. The Dieudonne/Cartier-dual framework remains cote4 and cote54 scans5--7 as recorded in sections2--3. No new motivic correspondence, ell-adic weight or Kummer RH implication is asserted.

\subsection{19. AM-K038 and ALP-LIFT-27: coupled section descent and faithful analytic image}\label{am-k038-and-alp-lift-27-coupled-section-descent-and-faithful-analytic-image}

K038 is an explicit descent construction of the two composed theta-section differential operators on the finite decorated lattice quotient, with coefficient modules over O\_M and actual stabilizer weights. Its polynomial scheme maps use the existing SGA3I4.7.3/4.7.3.1, SGA6VI3.3 and EGAIV4 16.11.2--3(ii) pp54--55 interfaces. The section/translation geometry retains FGA232 Theorem3.1 pp148--149, FGA236 Corollary1.7 p229/2.3 p232, SGA7I VII2.7/2.9 pp169--174 andVIII7.4(b) pp310--311. The strict/projective refinement is an actual quotient mass identity, not a change in the underlying Picard or motivic object. LIFT27 proves faithfulness of the specified full analytic image on its initial ramified space; it does not identify a faithful algebraic correspondence after terminal restriction. The formal selected content remains within the cote4/cote54 Cartier--Dieudonne framework already recorded. Neither analytic injectivity nor equality of lambda-ring classes proves a nonzero terminal motive, an ell-adic weight assertion or Kummer RH.

\subsection{20. AM-R-021: character twists and numerical weight windows}\label{am-r-021-character-twists-and-numerical-weight-windows}

Tensoring by a determinant character shifts both u and v equally; the v-width and first u−v gap are invariant. The full arithmetic budget requires an inequality between these invariants, violated by the current equally spaced family. This is a numerical restriction on coefficient representations within the SGA3I4.7.3/4.7.3.1 and SGA6VI3.3 framework, not a new geometric obstruction to all theta correspondences. It neither changes the existing Picard/biextension/formal crystal constructions nor supplies an ell-adic or motivic RH conclusion. The exact source and prior-run locators are R021§1--3.

\subsection{21. AM-DO-039: finite integral tests and relative representation models}\label{am-do-039-finite-integral-tests-and-relative-representation-models}

The common split-line kernel is a kernel of morphisms of the specified finite polynomial representation modules over Z and their localizations. Its characteristic-zero dimension is not determined by a single highest line. This refines the coefficient-restriction picture of SGA3I4.7.3/4.7.3.1 and SGA6VI3.3; it does not assert a kernel in the actual theta-correspondence image. The local Q open/closed orbits are relative representation models of one fixed analytic vector and are not different motivic realizations. The Picard, biextension and jet comparisons retain FGA232 Theorem3.1 pp148--149, FGA236 Corollary1.7 p229/2.3 p232, SGA7I VII2.7/2.9 pp169--174 andVIII7.4(b) pp310--311, and EGAIV4 16.11.2--3(ii) pp54--55. No new cycle nonvanishing, ell-adic weight assertion or L-function theorem follows.

\subsection{22. AM-K039: actual section integrality and class descent}\label{am-k039-actual-section-integrality-and-class-descent}

The new divisibility theorem is about the integral local-section module of the fixed principal theta line and its actual class descent over a DVR above11. Prime-to11 Heisenberg translation is identified through the normalized torsion difference connection, so the reduced exponential is retained. Its Picard/biextension/jet setting has the exact FGA232 Theorem3.1 pp148--149, FGA236 Corollary1.7 p229/2.3 p232, SGA7I VII2.7/2.9 pp169--174 andVIII7.4(b) pp310--311, EGAIV4 16.11.2--3(ii) pp54--55 locators already recorded. The selected differential functor acts on integral Weyl modules over Z (SGA3I4.7.3/4.7.3.1; SGA6VI3.3), while the numerical11 bound uses the established Cartier--Dieudonne comparison (cote4; cote54 scans5--7). Finite class stabilizers are explicit descent denominators; a K-theory class or a motivic description does not remove them. These inclusions give no nonzero terminal cycle, ell-adic weight theorem or Kummer RH.

\subsection{23. AM-R-022: a feasible coefficient-weight cloud}\label{am-r-022-a-feasible-coefficient-weight-cloud}

The new profile supplies explicit highest weights and binary interlacing grids in the same integral representation framework (SGA3I4.7.3/4.7.3.1; SGA6VI3.3), with exact real Scale parameters and rational coefficient weights. It changes the coefficient representations used by the jets, not the underlying J-power, Picard scheme or biextension line; the exact FGA/SGA7/EGA jet locators are those in R022§5 and sections12--22 above. Numerical window feasibility does not identify the integral theta image or a motivic correspondence and implies no ell-adic weight or RH assertion.

\subsection{24. ALP-LIFT-28: coherent jets and a bottom detector}\label{alp-lift-28-coherent-jets-and-a-bottom-detector}

The new bottom detector is an application of coherent finite-jet interpolation on a projective minimal compactification. Its exact Grothendieck input is EGAIII.1 Theorem2.2.1(ii), p100: sufficiently ample twists kill higher cohomology of the coherent kernel of finite-jet restriction. The finite jet space retains EGAIV4 16.11.2--3(ii), pp54--55. Mumford1977 Theorem3.1 and Proposition3.4(a),(b) supply the canonical extension and ample logarithmic-canonical line. These results construct an actual analytic/coherent detector at the explicit level, with stabilizer-equivariant coefficients; they do not identify it with a prescribed theta correspondence or bound its norm. The Picard/biextension locators remain FGA232/236 and SGA7I VII--VIII as above. No ell-adic cycle or motivic L-function conclusion follows.

\subsection{25. AM-K040: coherent CM character types and arithmetic coefficient fields}\label{am-k040-coherent-cm-character-types-and-arithmetic-coefficient-fields}

The principal/base-ideal compatibility builds actual algebraic angular characters for every weight in a fixed lattice, with one finite value field and preserved finite conductors. In the CM-motive language this is the infinity-type/finite-realization datum of rank-one character realizations; it is not a proof of RH or an automatic comparison with the mixed theta motive. The underlying CM Picard scheme and biextension line retain FGA232 Theorem3.1 pp148--149, FGA236 Corollary1.7 p229/2.3 p232 and SGA7I VII2.7/2.9 pp169--174, VIII7.4(b) pp310--311. The new theorem is an ideal/character construction, not a new theorem from those geometric locators. Its finite radical-field bound controls degree only; no unrecorded discriminant or height bound is supplied.

\subsection{26. AM-DO-040: certified prime and decorated lattice data}\label{am-do-040-certified-prime-and-decorated-lattice-data}

The explicit algebraic integers piT give finite etale split residue fibres and the actual decorated integral lattice support for the replacement coefficient weights. Their deterministic primality, norm and stabilizer assertions are arithmetic inputs to the existing scheme/jet comparison, not new theorems about Picard schemes. The exact geometric locators remain FGA232/236, SGA7I VII--VIII and EGAIV4 16.11.2--3(ii) stated above. Strict/projective level masses and finite character-field enlargements remain numerical costs; the lambda-ring or motivic description does not cancel them. No new nonvanishing or L-function conclusion follows.

\subsection{27. ALP-LIFT-29: translated finite packet and the same coherent jets}\label{alp-lift-29-translated-finite-packet-and-the-same-coherent-jets}

The finite level and translation define an actual collection of special points with exactly the original isotropy groups; the displayed local matrix calculation fixes the intersection of their integral models at T. The coherent detector therefore uses the same EGAIII.1 Theorem2.2.1(ii), p100, and EGAIV4 16.11.2--3(ii), pp54--55, finite-jet interfaces as LIFT28. The translation is an isometry of the complex automorphic realization; it is not asserted to preserve an integral section lattice or have zero arithmetic height. The Picard/biextension locators remain FGA232 Theorem3.1 pp148--149, FGA236 Corollary1.7 p229/2.3 p232 and SGA7I VII2.7/2.9 pp169--174, VIII7.4(b) pp310--311. No identification of the prescribed correspondence, motive, ell-adic cycle or L-function follows from the packet comparison alone.

\subsection{28. AM-R-023: finite quotient and scale constraints}\label{am-r-023-finite-quotient-and-scale-constraints}

The fixed-field obstruction is a finite Galois quotient statement over Spec(Q), independent of the valid mixed Picard/section construction. The conditional rank bounds concern the size of a full CM realization of h1(J\^{}N); they do not remove R003's separate simple-isogeny/type obstruction. The exact FGA/SGA7 Picard and biextension locators, SGA6VI3.3 representation framework and cote4/cote54 crystal realizations are retained in R023§5. No change in a virtual representation, motive or ell-adic realization relaxes the displayed integer and Scale inequalities. No L-function theorem is obtained.

\subsection{29. AM-K041: actual new-profile integral image}\label{am-k041-actual-new-profile-integral-image}

The all-word bound is an inclusion of the actual finite-class coherent Weyl coefficient modules over the DVR above11, with both stabilizer denominators. The primitive theta section is transported through the normalized torsion/Poincare connection on its formal Picard completion, retaining the Chern exponential. Schur operations retain SGA6VI3.3; the exact Picard/biextension and jet locators are FGA232 Theorem3.1 pp148--149, FGA236 Corollary1.7 p229/2.3 p232, SGA7I VII2.7/2.9 pp169--174 andVIII7.4(b) pp310--311, EGAIV4 16.11.2--3(ii) pp54--55. The separate unrestricted analytic injectivity does not imply a nonzero restricted motivic or ell-adic correspondence. No Kummer RH result follows.

\subsection{30. AM-R-024 and ALP-LIFT-30: terminal integral rows and negative-line transport}\label{am-r-024-and-alp-lift-30-terminal-integral-rows-and-negative-line-transport}

R024 gives inclusions of the actual image in Weyl/dual-Weyl coefficient lattices over the DVR above11, retaining normalized pairing, finite stabilizers and every EqFun constraint. The exact SGA3I4.7.3/4.7.3.1 and SGA6VI3.3 frameworks do not remove these indices. LIFT30 gives a primitive integral negative line and rational isometry of its positive complement, followed by a separately required real-vector transport; a unitary Hilbert norm does not imply bounded algebraic coordinate height. The fixed Picard/biextension/formal jet locators remain FGA232 Theorem3.1 pp148--149, FGA236 Corollary1.7 p229/2.3 p232, SGA7I VII2.7/2.9 pp169--174 andVIII7.4(b) pp310--311, EGAIV4 16.11.2--3(ii) pp54--55. No algebraic-cycle nonvanishing, ell-adic weight or Kummer RH statement is inferred.

R024.18--21 applies R002's metrized image-line balance to the actual class-independent branching scalar. Its principal finite divisor and all-embedding real logarithmic norm agree, so the corresponding scalar has degree zero. This is the product formula on the number-field base, not a cancellation of the remaining integral correspondence, restricted image, or detector norm.

\subsection{31. AM-K042: neat quotients and bounded coherent interpolation}\label{am-k042-neat-quotients-and-bounded-coherent-interpolation}

The existing arithmetic quotient is neat, so the exact packet has no point-isotropy correction. The improved initial denominator is a property of the actual integral class sum, not its K-theory class. The bounded-twist question is about principal-parts sheaves and a finite-point ideal kernel on the bottom Shimura variety. Exact framework locators remain EGAIV.4 16.11.2--3(ii), EGAIII.1 Theorem2.2.1(ii) p100, SGA6VI3.3, FGA232 Theorem3.1 pp148--149, FGA236 Corollary1.7 p229/2.3 p232 and SGA7I VII2.7/2.9 pp169--174, VIII7.4(b) pp310--311. Qualitative Serre vanishing gives no bound inside NB11's finite progression; no motivic, ell-adic-weight or Kummer RH conclusion is asserted.

\subsection{32. AM-DO-041: local quotient and its actual measure}\label{am-do-041-local-quotient-and-its-actual-measure}

The result is a characteristic-zero local representation quotient with a computed integral support and Haar measure. Its scalar is not an equality of integral lattices; its finite divisor and Hermitian norm must be retained when returning to the coefficient correspondence. The inherited Schur framework remains SGA3I4.7.3/4.7.3.1 and SGA6VI3.3, with the same FGA232/236, SGA7I VII--VIII and EGAIV.4 principal-parts locators. No motive, ell-adic weight, algebraic cycle or Kummer RH conclusion follows from the local quotient computation.

\subsection{33. AM-R-025: integral frames and the complete residue action}\label{am-r-025-integral-frames-and-the-complete-residue-action}

The comparison is of actual integral standard lattice stabilizers and their finite etale split residue quotients, acting on the principal-level automorphic space. It is stronger than a matching of subgroup intersections or virtual coefficient classes. SGA3I4.7.3/4.7.3.1 and SGA6VI3.3 remain the group/Schur frameworks, while the fixed CM Picard/biextension and jet locators are FGA232 Theorem3.1 pp148--149, FGA236 Corollary1.7 p229/2.3 p232, SGA7I VII2.7/2.9 pp169--174/VIII7.4(b) pp310--311 and EGAIV.4 16.11.2--3(ii) pp54--55. A concrete finite action does not by itself identify a motivic projector or prove a nonzero algebraic cycle, ell-adic weight, or Kummer RH statement.

\subsection{34. ALP-LIFT-31: an actual adjoint image equality}\label{alp-lift-31-an-actual-adjoint-image-equality}

The result identifies the image of the actual analytic theta correspondence at fixed arithmetic level and real type; it is not merely an equality in a representation ring. Its integral coefficient realization still retains the actual finite indices, Hermitian metrics and potential failure of integral saturation. The retained frameworks are SGA3I4.7.3/4.7.3.1, SGA6VI3.3, FGA232 Theorem3.1 pp148--149, FGA236 Corollary1.7 p229/2.3 p232, SGA7I VII2.7/2.9 pp169--174 andVIII7.4(b) pp310--311, and EGAIV.4 16.11.2--3(ii) pp54--55. Analytic surjectivity supplies no algebraic-cycle nonvanishing, effective lattice inverse, ell-adic weight or Kummer RH conclusion.

\subsection{35. AM-K043 and AM-R-026: metrized jets and finite integral correspondences}\label{am-k043-and-am-r-026-metrized-jets-and-finite-integral-correspondences}

K043 computes the scalar between the PBW and Fischer models of the selected principal-parts component, retaining the exact metric and integral valuations. R026 gives a genuine finite etale residue-group projector and an explicit integral multiple, with real coefficients after conjugation. Neither calculation supplies a global inverse or algebraic-cycle nonvanishing. The exact inherited locators remain EGAIV.4 16.11.2--3(ii) pp54--55, EGAIII.1 Theorem2.2.1(ii) p100, SGA3I4.7.3/4.7.3.1, SGA6VI3.3, FGA232 Theorem3.1 pp148--149, FGA236 Corollary1.7 p229/2.3 p232, SGA7I VII2.7/2.9 pp169--174 andVIII7.4(b) pp310--311. No new motive, ell-adic weight or RH conclusion is inferred.

\subsection{36. AM-DO-042: finite etale spectral labels}\label{am-do-042-finite-etale-spectral-labels}

The squarefree polynomials define finite etale algebras over the actual prime residue fields, with explicitly controlled Frobenius orbit degrees and prime-to-order projections. Their use as labels of genuine characteristic-zero finite-group projectors retains the integral denominator in R026. These finite spectra do not identify a motivic summand or prove the automorphic period. The SGA3I4.7.3/4.7.3.1 and SGA6VI3.3 group/representation frameworks, the established FGA232/236 and SGA7I VII--VIII Picard/biextension locators, and EGAIV.4 16.11.2--3(ii) principal-parts locators remain unchanged. No Kummer RH consequence is asserted.

\subsection{37. AM-DO-043: a prescribed scalar character in the finite projector}\label{am-do-043-a-prescribed-scalar-character-in-the-finite-projector}

The norm-compatible finite torus duality and squarefree spectral polynomials select an actual characteristic-zero residue-group summand with prescribed CM torsion character. The finite28average is integral at11; other projector denominators remain. This is a finite etale/group-algebra calculation in the SGA3I4.7.3/4.7.3.1 and SGA6VI3.3 framework, not a motivic decomposition or a nonzero CM cycle. The established FGA232/236, SGA7I VII--VIII and EGAIV.4 principal-parts locators remain the geometric comparison context.

\subsection{38. AM-K044: a bounded family of coefficient bundles}\label{am-k044-a-bounded-family-of-coefficient-bundles}

The enlarged interval is a family of actual dominant coefficient representations and canonical twists, with a different coherent automorphic sheaf for each ell. It does not turn qualitative EGAIII.1 Theorem2.2.1(ii), p100, into an effective uniform principal-parts bound; the relevant jet objects remain EGAIV.4 16.11.2--3(ii), pp54--55. Schur and residue actions retain SGA3I4.7.3/4.7.3.1 and SGA6VI3.3; the established FGA232/236 and SGA7I VII--VIII Picard/biextension locators remain unchanged. No motivic, ell-adic-weight, cycle nonvanishing or Kummer RH statement follows from these numerical inequalities.

\subsection{39. AM-K045: exact finite residue exceptional sets}\label{am-k045-exact-finite-residue-exceptional-sets}

The gcd certificate concerns the prime divisors and prime-to-part Frobenius spectra in the actual finite residue fields. It supplies a hypothesis of the cited characteristic-zero/finite-characteristic representation criterion; it does not construct an ell-adic sheaf or prove its purity. Finite etale and residue-group interpretations retain SGA3I4.7.3/4.7.3.1 and SGA6VI3.3; the established Picard/biextension and principal-parts locators remain unchanged. No motivic or Kummer RH consequence follows.

\subsection{40. ALP-LIFT-32 and AM-R-027: equivariant principal parts and a projected analytic period}\label{alp-lift-32-and-am-r-027-equivariant-principal-parts-and-a-projected-analytic-period}

The finite residue correspondence acts on the actual equivariant jet fibres, including nontrivial coarse stabilizers and their scalar linearization. Characteristic-zero exactness of finite projectors combines with coherent Serre vanishing (EGAIII.1 Theorem2.2.1(ii), p100) and principal parts (EGAIV.4 16.11.2--3(ii), pp54--55) to realize the compatible datum at sufficiently large twist. SGA3I4.7.3/4.7.3.1 and SGA6VI3.3 give the inherited group/Schur framework; FGA232 Theorem3.1 pp148--149, FGA236 Corollary1.7 p229/2.3 p232 and SGA7I VII2.7/2.9 pp169--174, VIII7.4(b) pp310--311 retain the fixed Picard/biextension comparison. The resulting positive complex period is not an algebraic-cycle nonvanishing theorem, a motivic projector or an ell-adic weight/RH statement.

\subsection{41. AM-R-028: a split map of localized finite value lattices}\label{am-r-028-a-split-map-of-localized-finite-value-lattices}

The Whittaker construction gives an explicit integral splitting over O\_L localized at11, compatible with the actual finite stabilizer action even when its order is not an11-unit. It therefore refines a representation-ring equality to a lattice statement. The residue-group/Schur framework remains SGA3I4.7.3/4.7.3.1 and SGA6VI3.3; principal-parts application retains EGAIV.4 16.11.2--3(ii), pp54--55. The fixed FGA232/236 and SGA7I VII--VIII Picard/biextension comparisons remain required. Global integral coherent cohomology and motivic or ell-adic nonvanishing do not follow.

\subsection{42. AM-R-029: an integral finite spectral summand}\label{am-r-029-an-integral-finite-spectral-summand}

The squarefree finite etale spectrum p+3+3 selects a genuine characteristic-zero residue-group idempotent integral over O\_L at11. It gives a direct summand of each supplied stable lattice, including equivariant principal-parts modules, without proving that global integral cusp sections surject onto them. Exact inherited frameworks are SGA3I4.7.3/4.7.3.1, SGA6VI3.3 and EGAIV.4 16.11.2--3(ii), pp54--55; the fixed FGA232/236 and SGA7I VII--VIII Picard/biextension locators remain. No motivic decomposition, ell-adic purity, algebraic-cycle nonvanishing or RH theorem is asserted.

\subsection{43. AM-K046: a metrized analytic principal-parts splitting at one point}\label{am-k046-a-metrized-analytic-principal-parts-splitting-at-one-point}

The Poincare construction realizes the selected principal-parts fibre at the standard arithmetic point with an explicit complex Hilbert inverse. Its finite-volume metric and PBW normalization are retained; it is not yet a splitting of the global integral coherent evaluation map. EGAIV.4 16.11.2--3(ii), pp54--55, and EGAIII.1 Theorem2.2.1(ii), p100, remain the jet/coherent frameworks, with SGA3I4.7.3/4.7.3.1 and SGA6VI3.3 for coefficient actions. The fixed FGA/SGA7 Picard and biextension bridge remains required. No algebraic-cycle nonvanishing, motivic or ell-adic purity/RH result is inferred.

\subsection{44. ALP-LIFT-33: integral invariant pairings on actual class fibres}\label{alp-lift-33-integral-invariant-pairings-on-actual-class-fibres}

The exact invariant Weyl/dual-Weyl lattices exhibit non-scalar elementary divisors after the actual stabilizer denominator. This refines characteristic-zero representation and K-theory data to a localized integral module calculation; it does not identify the image of the global theta correspondence. The inherited exact locators remain SGA3I4.7.3/4.7.3.1, SGA6VI3.3, EGAIV.4 16.11.2--3(ii), pp54--55, and the FGA232/236, SGA7I VII--VIII Picard/biextension locators above. No new algebraic cycle, motivic projector or ell-adic purity claim follows.

\subsection{45. AM-K047: an arithmetic translate with small real displacement}\label{am-k047-an-arithmetic-translate-with-small-real-displacement}

The constructed rational unitary unipotent belongs to an actual finite-adelic translated packet group. It obstructs one analytic principal-parts construction at the prescribed point; it does not obstruct the coherent sheaf itself or every evaluation map. The principal-parts framework remains EGAIV.4 16.11.2--3(ii), pp54--55, and the qualitative coherence framework EGAIII.1 Theorem2.2.1(ii), p100. The fixed Picard/biextension and motive comparisons remain open and no RH consequence is asserted.

\subsection{46. AM-DO-044: an equivariant cohomology obstruction to full principal-parts interpolation}\label{am-do-044-an-equivariant-cohomology-obstruction-to-full-principal-parts-interpolation}

The actual free finite cover gives a representation-ring Euler characteristic equal to the regular residue-group class times the local-system rank and base Euler characteristic. Algebraic-coefficient boundary vanishing and regular-weight middle cohomology turn it into a dimension bound; exact determinant fibres identify the compact-dual Euler measure with the definite genus mass. This obstructs surjectivity of the full selected evaluation map in the stated weight interval, while preserving the distinction between coherent large-twist existence and an effective finite window (EGAIII.1 Theorem2.2.1(ii),p100; EGAIV.4 16.11.2--3(ii),pp54--55). The inherited SGA3I4.7.3/4.7.3.1 and SGA6VI3.3 action/K-theory frameworks, and the established FGA232/236 and SGA7I VII--VIII Picard/biextension locators remain. No particular algebraic cycle is shown zero and no motivic, ell-adic purity or RH assertion follows.

\subsection{47. AM-R-030: matched ball volume and definite mass}\label{am-r-030-matched-ball-volume-and-definite-mass}

The exact ball-Haar total volume is m\_n{[}Kmax:K{]}/(a!)\^{}f, with orbifold and ineffective central stabilizers retained. The new density calculation matches the compact dual CP\^{}a at the tangent origin and gives its volume 1/a!, hence Euler density (n!)\^{}f.~This is a metric and arithmetic-class normalization, not an additional algebraic-cycle or motivic identification. K046 normalized selected principal parts retain EGA III.1 Theorem 2.2.1(ii), p.100, and EGA IV.4 16.11.2--3(ii), pp.54--55; the run's previously verified Picard/biextension and K-theory locators remain unchanged. Theorem provenance: Kionke (2014), Theorems 3.1 and 3.3, pp.383,386--387; DO044.13--15 for the exact SU determinant-fibre comparison. Whole-proof checks: root69 and independent DO review.

\subsection{48. ALP-LIFT-34: actual middle trace and integral cutoff}\label{alp-lift-34-actual-middle-trace-and-integral-cutoff}

The finite free residue action yields a 0/1 transversal section whose finite-group trace is the original full-Y Schwartz tensor. Complete integral Weyl/dual-Weyl coevaluation and rational theta automorphy turn that trace into an exact factor\textbar Gamma\textbar. Prime-to11 Heisenberg expansion constructs the divided section independently in the integral theta-line module. The result belongs to the existing biextension/Picard and Dieudonne-crystal framework (SGA7I ExposesVII--VIII and the cote4/cote54 scan locators already verified above), with the actual integral Schur functors of K037 retained. It supplies no new motivic projector or purity theorem. Middle mass is cancelled on this image; the terminal mass remains384/357.

\subsection{49. AM-DO-045: direct oscillator occurrence obstruction}\label{am-do-045-direct-oscillator-occurrence-obstruction}

Przebinda1996 Theorems1.13/1.19,(1.16--18),p97, gives the infinitesimal-character inclusion; the actual highest-weight widths refute the prescribed direct real occurrence at top labels. At split finite places, Minguez2008 Theorem1(2),p718, forces a GL1 block. Under the finite Green characteristic-ring isomorphism, Gow--Vinroot2007 section3,Theorem3.2,pp8--9, induction is the product in the Schur alphabets; a positive degree-one alphabet remains. This is a representation-ring/K-theoretic support obstruction consistent with the previously cited lambda-ring and SGA6 locators, not a construction of an algebraic or ell-adic projector. Literal IsGood clause5 kills that support in every actual e0 block. The branching restriction period and the existing two-step tower are different maps and remain outside this obstruction.

\subsection{50. ALP-LIFT-35: a terminal finite-action fixed locus}\label{alp-lift-35-a-terminal-finite-action-fixed-locus}

The actual terminal row is a decorated integral Hermitian lattice with empty T set. Its standard flag stabilizer contains an order11 permutation of the complementary t rows; those rows vanish in the firstG residue of every raw kernel input. The quotient therefore has a fixed locus on the entire prescribed residue support, and the middle free-orbit trace argument does not extend. This is an exact finite-group action statement, consistent with the finite-etale quotient framework of SGA3I4.7.3/4.7.3.1 already located above. It proves no assertion about the integral theta value on that fixed locus and gives no motivic, ell-adic or period-vanishing conclusion.

\subsection{51. AM-K048: normalized local spectral sampling}\label{am-k048-normalized-local-spectral-sampling}

The actual CM sign pattern identifies the central-frequency lattice with a signed inverse-different image. K043 PBW normalization and the entire selected compact module are retained in the unitary Laguerre transform (Groenevelt2003 section3.1 eq3.2,section3.2 eq3.7,section3.3 Prop3.2,PDFpp9--12). Exact norm and packing estimates give the stated local sampling tail only under top-count\textless=f/2. This is an analytic representation estimate, not an automorphic Fourier coefficient bound or a motivic comparison. Existing Dieudonne/Picard/jet dictionary locators are unchanged.

\subsection{52. AM-R-031: explicit global standard evaluation}\label{am-r-031-explicit-global-standard-evaluation}

A product-exponential chart embedded by K046 gives an explicit Hilbert-valued Sobolev evaluation inequality. Its Haar normalization is proved directly from the Frobenius quotient metric and a compact Givens volume bound, and the raw/average/adelic norm conversions use R030 exactly. This concerns metrized differential evaluation of automorphic sections. Algebraic principal-parts surjectivity still uses EGAIII.1 Theorem2.2.1(ii),p100 and EGAIV.4 16.11.2--3(ii),pp54--55; the inequality supplies neither surjectivity nor an algebraic-cycle/ell-adic identification.

\subsection{53. AM-DO-046: actual automorphic parameters and normalized periods}\label{am-do-046-actual-automorphic-parameters-and-normalized-periods}

The actual cohomological coefficient systems have wide infinitesimal-character gaps, forcing generic Arthur parameters. After an honest odd Hecke-character parity correction, each cuspidal GL block falls within Caraiani2012 Theorem1.2,p2312. This uses a known cohomological/purity-based temperedness theorem; it neither constructs a new motive nor proves RH for its L-function. The actual finite goodTuple acts through the Green representation ring, whose degree grading rules out proper global blocks at the bottom. The normalized GGP identity is analytic and retains the chosen cycle/test-vector question. The established SGA3I4.7.3/4.7.3.1 and SGA6VI3.3 group/K-theory, EGAIV.4 16.11.2--3(ii),pp54--55 principal-parts, FGA232/236 and SGA7I VII--VIII Picard/biextension locators remain the comparison framework. A motivic or algebraic-cycle identification of the actual restricted period remains unproved.

\subsection{54. ALP-LIFT-36: an effective Hecke boundary factor}\label{alp-lift-36-an-effective-hecke-boundary-factor}

The actual stage ratio is an honest unitary Hecke character with nonzero angular weights, constructed from the CM character lattice. Its quantitative boundary estimate is an analytic inequality for its finite Euler product, with every local omission retained. It does not assert an algebraic period formula, a motivic critical-value theorem or any statement about zeros in the critical strip. The fixed Picard/biextension identifications of FGA232/236 and SGA7I VII--VIII and the cote4/cote54 Dieudonne-crystal comparisons remain the arithmetic context, not consequences of the estimate. The new bound removes an ineffective finite factor from the actual theta-preimage calculation; Q3 Kummer RH is untouched.

\subsection{55. AM-K049: arithmetic-invariant cutoff and a metrized selected jet}\label{am-k049-arithmetic-invariant-cutoff-and-a-metrized-selected-jet}

The rational central lattice supplies exact automorphic invariance of a periodic functional-calculus cutoff. Its spectral support and Lie derivatives control the actual analytic principal-parts evaluation at one translated point, with the canonical coefficient metric, PBW scalar and mass-volume conversion retained. The exact algebraic principal-parts locators remain EGAIV.4 16.11.2--3(ii),pp54--55; qualitative coherent realization remains EGAIII.1 Theorem2.2.1(ii),p100. This analytic small-norm estimate does not identify a rational or integral jet image, nor prove an algebraic cycle zero. The established SGA3/SGA6, FGA232/236 and SGA7I VII--VIII framework remains; no additional motivic projector, ell-adic weight or RH consequence is inferred.

\subsection{56. ALP-LIFT-37 and AM-DO-047: exact local Hermitian correspondences}\label{alp-lift-37-and-am-do-047-exact-local-hermitian-correspondences}

The prescribed finite Schwartz kernels determine exact positive Hermitian forms on the actual preceding theta image. Integral residue-frame support, finite parabolic indices and local matrix disintegration fix the G scalar; an absolute induced-spherical majorant and finite-dimensional continuation fix the Q scalar. These are metric normalizations of the existing local oscillator correspondence, with raw additive lattices and character ramification retained. They do not assert a split integral global correspondence or determine its Smith factors. The fixed SGA3/SGA6 representation and K-theory, FGA232/236 Picard and SGA7I VII--VIII biextension locators remain applicable; no new motivic or ell-adic realization follows from a local positive form.

\subsection{57. AM-DO-048: exact real Hermitian normalization}\label{am-do-048-exact-real-hermitian-normalization}

The formal degree calibrates the analytic Hilbert metric on the actual cohomological coefficient representation: d(tau)=a!dim V\_(h-rho) in the fixed ball Haar. It distinguishes scalar matrix coefficients from vector-valued evaluation. This changes a metric in the existing principal-parts comparison, not an integral lattice or motivic period identification. Exact inherited locators remain EGAIV.4 16.11.2--3(ii),pp54--55; SGA3I4.7.3/4.7.3.1; SGA6VI3.3, with the FGA232/236 and SGA7I VII--VIII Picard/biextension locators above. No ell-adic weight or RH conclusion follows.

\subsection{58. AM-K050: bounded genuine automorphic sections}\label{am-k050-bounded-genuine-automorphic-sections}

The fixed finite residue-group summand and central character admit a nonzero holomorphic cusp section at the unchanged arithmetic level throughout the stated bounded gap interval. Its finite stabilizer calculation occurs in the representation ring, using the Deligne--Lusztig cohomological character formula (1976,Theorem4.2,p124; Proposition7.4,p141; Theorem8.3,p147). Analytic Poincare nonvanishing does not establish an integral coherent principal-parts splitting, an algebraic-cycle value, or a motivic projector. The exact comparison locators remain EGAIII.1 Theorem2.2.1(ii),p100; EGAIV.4 16.11.2--3(ii),pp54--55; SGA3I4.7.3/4.7.3.1 and SGA6VI3.3, and the fixed FGA/SGA7 Picard-biextension locators. Q3 Hecke RH is unchanged.

\subsection{59. AM-K051: full-interval analytic section existence}\label{am-k051-full-interval-analytic-section-existence}

The radius and negative-moment refinement extends K050 to both complete admitted weight intervals without altering the finite representation or arithmetic level. It removes the remaining bounded-existence subgate, while preserving the difference between analytic nonzero cusp sections and integral principal-parts surjectivity or algebraic-cycle nonvanishing. The exact Grothendieck locators are those of section58; no new motivic, ell-adic or RH statement is involved.

\subsection{60. ALP-LIFT-38: an explicit Hilbert inverse of the analytic correspondence}\label{alp-lift-38-an-explicit-hilbert-inverse-of-the-analytic-correspondence}

The actual two-step vertical correspondence has an explicitly bounded complex Petersson right inverse. Its exact scalar includes definite masses, both successive local Hermitian forms, Hecke Euler values, fixed real formal degrees and prescribed source-intertwiner scales. A Hilbert inverse does not give a splitting of the localized integral theta lattice or determine its Smith factors. The fixed FGA232/236 and SGA7I VII--VIII Picard/biextension bridge, SGA3I4.7.3/4.7.3.1 and SGA6VI3.3 coefficient action/K-theory framework, and EGAIV.4 16.11.2--3(ii),pp54--55 principal-parts comparison remain required. No new motive, algebraic cycle, ell-adic weight or Hecke RH follows from the analytic bound.

\subsection{61. AM-R-032: the same finite correspondence on classes and analytic sections}\label{am-r-032-the-same-finite-correspondence-on-classes-and-analytic-sections}

The genuine residue-group projector is identified coefficient by coefficient with the literal arithmetic character idempotent. Rational lattice isomorphisms transport the integral Weyl and dual-Weyl coefficient modules, and the terminal class realization intertwines this action with analytic right translation. Its complex least-norm vertical inverse respects the same projector. These are finite correspondence and coefficient-functor statements within the already located SGA3I4.7.3/4.7.3.1 and SGA6VI3.3 representation/K-theory framework. The FGA232/236 and SGA7I VII--VIII Picard/biextension and cote4/cote54 crystal comparisons are carried as previously verified exact locators. An analytic inverse is not an integral splitting, and this finite projector equality constructs no new motivic or ell-adic projector, weight theorem or algebraic cycle. The fixed real theta/CM comparison remains a separate gate.

\subsection{62. AM-DO-049: metric normalization and transport of a correspondence}\label{am-do-049-metric-normalization-and-transport-of-a-correspondence}

An isometric analytic Howe contraction can be chosen in the constructed family; changing its scalar changes the multiplicity correspondence, its adjoint, least inverse, spectral scalar and period comparison by the proved factors. The literal arithmetic Hilbert ties and fixed unit-extremal auxiliary embedding constrain such changes. This is a metric statement over the complex realization and does not modify an integral correspondence or its motive for free. The exact principal-parts locators remain EGAIV.4 16.11.2--3(ii),pp54--55 and EGAIII.1 Theorem2.2.1(ii),p100; the Picard/biextension and K-theory locators in section61 remain unchanged. No RH conclusion follows.

\subsection{63. AM-K052: actual residue types and the spherical Hecke metric}\label{am-k052-actual-residue-types-and-the-spherical-hecke-metric}

The literal T norm-kernel condition forces a nontrivial tame central character and conductor-one newform. Its finite compact representation is an induced residue-group module, and the unchanged lower translate has a positive exact Gram factor on every such type. The minuscule spherical Hecke scalar is computed in the representation ring of the dual general linear group through Gross1998(3.14),p8 and Proposition6.4,p11. This representation-ring calculation is compatible with the already located SGA3I4.7.3/4.7.3.1 and SGA6VI3.3 dictionary; it does not supply an algebraic/motivic projector, a global cycle value, or a new ell-adic purity assertion. Temperedness is the carried DO046 theorem. Exact principal-parts and FGA/SGA7 Picard-biextension locators are unchanged. A nonzero compact projection is not the nonvanishing of the prescribed full local period or of its global central L-value.

\subsection{64. ALP-LIFT-39: primitive polynomial modules and the CM metric period}\label{alp-lift-39-primitive-polynomial-modules-and-the-cm-metric-period}

Primitive highest polynomials in the two actual joint-harmonic modules define rational Chevalley/Fock intertwiners, whose exact norms are hook products. The construction is compatible with the coefficient Schur functors in SGA3I4.7.3/4.7.3.1 and SGA6VI3.3. Rationality does not establish an integral inverse or its Smith factors. The theta realization is an isometry for the specified additive measures, so the literal unit vacuum introduces an explicit positive CM metric period; a negative Hermitian column conjugates its CM line and coordinate factors. The existing FGA232 Theorem3.1 pp148--149, FGA236 Corollary1.7 p229/2.3 p232, and SGA7I VII2.7/2.9 pp169--174, VIII7.4(b) pp310--311 locate the Picard/biextension framework. The cote4/cote54 Cartier--Dieudonne comparisons remain as previously verified, not new consequences of a hook norm. No complex metric factor is assigned an11-adic valuation, and no new motive, ell-adic weight or RH assertion is inferred.

\subsection{65. AM-K053: the actual tame test and the Weil--Deligne conductor}\label{am-k053-the-actual-tame-test-and-the-weildeligne-conductor}

The actual conductor-one determinant condition forces precisely one tame rank-one character and an unramified complement, with zero monodromy. Deligne1973 formula8.12.1 printed572 supplies the exact conductor of the full Weil--Deligne object; it must precede passage to a semisimplified Weil parameter. The normalized local test is a character-projector calculation, an epsilon/Gauss factor and a Schur-functor evaluation, compatible with SGA3I4.7.3/4.7.3.1 and SGA6VI3.3 in the previously located representation and lambda-ring framework. This is a local analytic realization calculation; it constructs no new global motive, no integral theta-image saturation and no Hecke RH statement. The existing FGA/Picard, SGA7/biextension and cote4/cote54 dictionary remains the geometric location of the underlying CM family, rather than a source of unproved global nonvanishing.

\subsection{66. AM-R-033: reductive special-fibre types and the full local period}\label{am-r-033-reductive-special-fibre-types-and-the-full-local-period}

The same smooth GL group scheme over the local integer ring and its reductive special fibre carry the actual generic finite type. Morris1999 Proposition4.1 turns the cuspidal special-fibre type into a compact-induction realization whose complete period is a finite invariant projection. The character/Schur operations sit in the existing SGA3I4.7.3/4.7.3.1 and SGA6VI3.3 framework. This is a local analytic realization result, not an integral-motive construction or a new crystal slope. The established FGA232/236 Picard and SGA7I VII--VIII biextension locators, and cote4/cote54 Cartier--Dieudonne dictionary, continue to describe the unchanged CM object. The two mixed-factor places and global arithmetic image are not inferred from this realization.

\subsection{67. AM-K054: exact local period measure and the adjoint central direction}\label{am-k054-exact-local-period-measure-and-the-adjoint-central-direction}

The canonical T factor is rational after the explicit local Rankin--Selberg and full split-unitary adjoint factors are divided out. Zhang2014 section1.1/Remark7 and Proposition4.10 give the precise local realization dictionary; the scalar direction in the unitary adjoint must be retained. This is a compatibility of analytic realizations and finite character normalizations, with Schur/Cauchy operations in the previously located lambda-ring language of SGA6VI3.3. Its rational eleven-adic valuation is not the valuation of an integral CM theta row, a new motive, or a consequence for RH. The Picard/biextension and cote4/cote54 locators remain unchanged and must not be credited with global nonvanishing.

\subsection{68. AM-DO-050: coevaluation, branching and the polarized real realization}\label{am-do-050-coevaluation-branching-and-the-polarized-real-realization}

The source pairBK duality supplies the canonical coefficient coevaluation. Cauchy multiplicity one and codimension-one branching identify its selected horizontal component, with its exact compact dimensions and hook/PBW metric. These are tensor/Schur operations in SGA3I4.7.3/4.7.3.1 and the lambda/K-theory framework of SGA6VI3.3, followed by a polarized real oscillator realization. The negative CM line conjugates its period; the positive amplitudes balance by the actual ranks. The unchanged Picard and biextension object is located at FGA232 Theorem3.1 pp148--149, FGA236 Corollary1.7 p229 and2.3 p232, SGA7I VII2.7/2.9 pp169--174 andVIII7.4(b) pp310--311. The cote4/cote54 comparison supplies no additional finite frame phase or global theta nonvanishing.

\subsection{69. ALP-LIFT-40: integral Schur modules, divided powers and explicit dual lattices}\label{alp-lift-40-integral-schur-modules-divided-powers-and-explicit-dual-lattices}

The generated modules are honest lattices over Z, transported under flat localization to the actual coefficient DVR. Compact divided powers preserve their integral Schur actions, while the ordinary noncompact creators define the intrinsic PBW lattice. Its dual is explicitly the minimal dual-Weyl lattice tensored with divided monomials. These distinctions belong to the integral representation/group-scheme framework already located in SGA3I4.7.3/4.7.3.1 and the Schur operations of SGA6VI3.3. They are not a Cartier--Dieudonne formal-group slope: cote4/cote54 and the earlier DP results remain distinct. The explicit integral section is a coefficient-module section; it does not produce a section of the global arithmetic theta-image lattice, a motivic projector over Z, or a new L-function statement.

\subsection{70. AM-K055: finite generic modules and the full local period}\label{am-k055-finite-generic-modules-and-the-full-local-period}

The finite GL generic summand is an exact representation over a cyclotomic coefficient field, with integral-group-scheme and Schur-operation language supplied by the already located SGA3 I4.7.3/4.7.3.1 and SGA6 VI3.3. Its passage to a smooth local Whittaker model is proved by a support-and-uniqueness argument; this analytic realization is not an additional assertion in those SGA texts. The prescribed inverse-generic vector and its conjugate finite function give an exact full local period1/dim(sigma). The probability and selfdual Haar comparison is explicit. This settles the original mixed G factors in the MainX local pairing without changing the Picard/biextension objects located in FGA232 Theorem3.1 pp148--149, FGA236 Corollary1.7 p229 and2.3 p232, SGA7 I VII2.7/2.9 pp169--174 andVIII7.4(b) pp310--311. It neither constructs a motivic projector over Z nor proves a global motivic L-value nonzero. Cote4/cote54 formal-crystal slope results remain unchanged.

\subsection{71. AM-K056: finite-etale unitary integral models and unramified local factors}\label{am-k056-finite-etale-unitary-integral-models-and-unramified-local-factors}

The endpoint stabilizers are the A-points of explicit reductive group schemes over each valuation ring A. A finite-etale quadratic base change splits them to GL\_n, and faithfully flat descent supplies the smooth reductive form. This is the group-scheme/descent use of the framework already located in SGA3 I4.7.3/4.7.3.1; it is independent of the absolute different of the CM trace polarization. Their hyperspecial-fixed local representations have the unramified canonical GGP factor1. The associated local L-functions use their actual unitary normalization and Gross motive, as in Zhang2014(1.1); this proves no global motivic central-value assertion. The earlier Picard, biextension and cote4/cote54 crystal dictionary is unchanged, and no new divided-power slope is extracted from this local group-scheme fact.

\subsection{72. AM-K057: exact real realizations and complete L-normalization}\label{am-k057-exact-real-realizations-and-complete-l-normalization}

The compact Schur coefficient is an invariant pairing of the actual coefficient representation; its PBW/Fischer scale is fixed before taking the analytic period. This uses the already located coefficient-functor and lambda/K-theory framework of SGA3 I4.7.3/4.7.3.1 and SGA6 VI3.3. The real Weil-group parameter and its full unitary adjoint determine the archimedean factors through the exact primary locators Knapp1994(3.6),p404/(4.6),p406 and Mok2015(2.1.1),pp6--7/(7.1.2--3),p190. These are analytic realization data; no new comparison isomorphism of motives or ell-adic sheaves is constructed. The actual central Rankin value remains unproved nonzero. The established FGA232/236 Picard, SGA7 I VII--VIII biextension and cote4/cote54 crystal locators remain the geometric context; neither an exact gamma product nor positivity of the local tests yields a new crystal slope, algebraic-cycle value or Hecke RH theorem.

\subsection{73. AM-DO-051: a fixed finite Heisenberg correspondence and its exact metric}\label{am-do-051-a-fixed-finite-heisenberg-correspondence-and-its-exact-metric}

The explicit operator is a finite combination of the translations and characters of the same theta/Heisenberg representation. Its coefficients lie in Z{[}zeta\_T,1/T{]}, hence are integral at11, while the actual source-model phase remains a separate realization scalar. This gives a concrete correspondence on the fixed local theta module; it is not a new integral splitting of the global theta-image lattice. The existing Picard/biextension locators FGA232 Theorem3.1 pp148--149, FGA236 Corollary1.7 p229/2.3 p232, and SGA7 I VII2.7/2.9 pp169--174, VIII7.4(b) pp310--311 remain the geometric interpretation. The finite Fourier norm and T-power denominators do not change the cote4/cote54 Cartier--Dieudonne slopes, or yield a motivic projector, ell-adic weight theorem or RH statement. Schur/PBW real factors remain those already located in SGA3 I4.7.3/4.7.3.1 and SGA6 VI3.3 and are not absorbed into this finite calculation.

\subsection{74. ALP-LIFT-41: a rational negative line and an actual analytic jet}\label{alp-lift-41-a-rational-negative-line-and-an-actual-analytic-jet}

The vector z is an explicit rational point of the norm-minus-one Hermitian quadric, primitive in the stated integral local lattices and with both split congruence components fixed. Its algebraic coordinates and local isometry class are scheme-level data; the near-origin bound concerns their complex analytic realization. The selected one-point inverse is an analytic section/jet theorem on the existing ball quotient, not a surjection of integral Picard or theta-image lattices. The relevant coefficient-functor interpretation remains SGA3 I4.7.3/4.7.3.1 and SGA6 VI3.3; the already located FGA232 Theorem3.1 pp148--149 and FGA236 Corollary1.7 p229/2.3 p232 give the Picard setting, while SGA7 I VII2.7/2.9 pp169--174 and VIII7.4(b) pp310--311 give its biextension context. No new motive, ell-adic weight statement or comparison map is claimed. The cote4/cote54 scans5--7 formal-group and Cartier-dual results retain their previous scopes; an archimedean jet inverse does not improve their mixed-line slopes.

\subsection{75. AM-R-034: the actual Q theta-module coordinate change}\label{am-r-034-the-actual-q-theta-module-coordinate-change}

The comparison is an explicit change of the same finite Heisenberg realization, preserving the basic lattice vector and carrying a residue-character root of unity. It is a concrete theta-module identity over the existing cyclotomic coefficient field, with its complex realization fixed; it is not a new motivic comparison theorem. Its parahoric indices describe the finite flag quotients of the integral general linear group schemes, and their exact rational factors remain part of the period trace. The Picard and biextension interpretation uses the previously checked FGA232 Theorem3.1 pp148--149, FGA236 Corollary1.7 p229/2.3 p232, and SGA7 I VII2.7/2.9 pp169--174, VIII7.4(b) pp310--311. The source-frame unit is not an ell-adic weight or a formal-group slope; cote4/cote54 scans5--7 remain unchanged. The selected coefficient functors retain their SGA3 I4.7.3/4.7.3.1 and SGA6 VI3.3 locators. No arithmetic cycle value for032 or Kummer RH follows from this local identity.

\subsection{76. AM-DO-052: generic character summands and the scaled integral frame}\label{am-do-052-generic-character-summands-and-the-scaled-integral-frame}

The finite generic averaging operator is a character idempotent over the coefficient ring with the G residue primes inverted; these denominators are units at11. It acts on the same finite theta module and commutes with the genuine central projector, while the exact source basis determines a separate rational modulus and tame-character scalar. This is an explicit correspondence on the existing coefficient realization, with equality of analytic maps on the character summand; no integral inverse or new motivic projector is supplied. Its coefficient-functor context remains SGA3 I4.7.3/4.7.3.1 and the lambda/K-theory setting SGA6 VI3.3. The fixed Picard and biextension locators remain FGA232 Theorem3.1 pp148--149, FGA236 Corollary1.7 p229/2.3 p232, SGA7 I VII2.7/2.9 pp169--174 and VIII7.4(b) pp310--311. The source-frame scaling is not a change in the cote4/cote54 scans5--7 Dieudonne slopes, an ell-adic weight result or a value of an algebraic cycle.

\subsection{77. ALP-LIFT-42: finite representation characters and arithmetic separation}\label{alp-lift-42-finite-representation-characters-and-arithmetic-separation}

The finite reductions are points of the actual integral GL group schemes, and their ordinary cuspidal characters come from regular finite-field multiplicative characters in Green's classification. The scalar-reduction estimate uses norms of nonzero integral coordinates at every embedding; the packing estimate is in the complex analytic realization and retains every finite automorphism of the base point. These inputs do not themselves construct a motivic or ell-adic projector, establish a weight theorem, or give an integral theta-image bound. Their coefficient representation setting retains SGA3 I4.7.3/4.7.3.1 and SGA6 VI3.3. The previously fixed FGA232 Theorem3.1 pp148--149, FGA236 Corollary1.7 p229/2.3 p232 and SGA7 I VII2.7/2.9 pp169--174, VIII7.4(b) pp310--311 remain the Picard/biextension context; cote4/cote54 scans5--7 formal-group slopes are unchanged. The primary finite character locators are Green1955Section5Example(ii)pp430--431 and Gorodetsky--Hazan2017Theorem2.1(2.1)p5, explicitly read in this run.

\subsection{78. AM-DO-053: the identity initial Q correspondence and its exact finite orbit}\label{am-do-053-the-identity-initial-q-correspondence-and-its-exact-finite-orbit}

At m2 the initial local correspondence is literally the identity, and its compact fibre has probability mass1. The middle primitive-vector orbit is a finite quotient of the same parahoric integral group scheme; its cardinality and both flag indices enter the trace with the specified weights. These are explicit theta-module/finite-correspondence statements, not an assertion that every global moved correspondence is identity. The Picard/biextension interpretation retains FGA232 Theorem3.1 pp148--149, FGA236 Corollary1.7 p229/2.3 p232 and SGA7 I VII2.7/2.9 pp169--174, VIII7.4(b) pp310--311. The coefficient action retains SGA3 I4.7.3/4.7.3.1 and SGA6 VI3.3. Cote4/cote54 scans5--7 crystal slopes, ell-adic weights and motive comparison questions are unchanged; the integrated orbit trace is not a new032 algebraic-cycle value or Kummer RH theorem.

\subsection{79. AM-R-035: one fixed finite tensor and its translated correspondences}\label{am-r-035-one-fixed-finite-tensor-and-its-translated-correspondences}

R035 assembles the finite oscillator representations and their actual Heisenberg correspondences before varying the auxiliary tuple. The common tensor is fixed in the original global model; the local passive characters and phases are part of its comparison, not a new choice of motivic period. The geometric interpretation retains the Poincare biextension of SGA7 I VII2.7/2.9 pp169--174 and VIII7.4(b) pp310--311, and FGA232 Theorem3.1 pp148--149 / FGA236 Corollary1.7 p229 and2.3 p232. It does not construct the specific032 algebraic-cycle class or a new motive comparison.

\subsection{80. AM-K058: a nonzero equivariant selected analytic fibre}\label{am-k058-a-nonzero-equivariant-selected-analytic-fibre}

The genuine finite projector and selected archimedean coefficient commute with the retained finite group actions; their Poincare Gram has a positive eigenvalue on the actual invariant fibre. In representation-scheme language this is a nonzero complex equivariant section with the stated finite-character cut (SGA3 I4.7.3/4.7.3.1); the associated Schur/charge operations retain SGA6 VI3.3. It is an analytic nonzero restriction, not an integral theta-image saturation, algebraic cycle, ell-adic weight assertion or RH theorem. The Dieudonne/crystal statements at cote4 and cote54 scans5--7 remain the previously cited local slopes.

\subsection{81. ALP-LIFT-43: joint character projection and the actual finite correspondence}\label{alp-lift-43-joint-character-projection-and-the-actual-finite-correspondence}

The exact ramified projector is the product of augmentation-kernel projectors of the selected residue-torus actions. Its mass-weighted descended fibre calculation retains every stabilizer, as required for an equivariant section rather than a free-orbit model. The action and invariant-module interpretation use SGA3 I4.7.3/4.7.3.1; the theta-module/biextension interpretation retains SGA7 I VII2.7/2.9 pp169--174 and VIII7.4(b) pp310--311. The field-norm exclusion uses the existing integral Hermitian lattice and all CM embeddings. No new integral projector, theta-image saturation, ell-adic weight, cycle class, crystal slope or motivic L-function theorem is asserted.

\subsection{82. AM-DO-054: one fixed theta tensor and the original period law}\label{am-do-054-one-fixed-theta-tensor-and-the-original-period-law}

The tensor is a section in the already fixed theta/Heisenberg model, transported by a finite correspondence before the auxiliary tuple varies. The original-vacuum period law retains its ramified augmentation kernels. The biextension interpretation uses the already verified SGA7 I, Exp.VII2.7/2.9, pp169--174, and Exp.VIII7.4(b), pp310--311; the equivariant-module language uses SGA3 I4.7.3/4.7.3.1. The result proves an analytic comparison in that fixed model, without asserting a new crystalline, ell-adic weight or motivic L-function theorem.

\subsection{83. AM-K059: integral sections, finite descent and spanning}\label{am-k059-integral-sections-finite-descent-and-spanning}

Local clearing is membership in a full lattice after localization at every finite coefficient-field place. The stabilizer and genuine-character sums give integral sections whose fraction-field span is the required invariant summand; they do not assert saturation. This is the finite equivariant-module framework of SGA3 I4.7.3/4.7.3.1 already located in this run. The literal Weyl/dual-Weyl pairing and theta-line interpretation retain the preceding SGA7 locators. CharT forces muT in the coefficient field, without changing the CM host. No new formal-group slope, lambda-ring identity, motive or weight purity claim follows from this clearing argument.

\subsection{84. ALP-LIFT-44: the actual lattice modification and fixed complement}\label{alp-lift-44-the-actual-lattice-modification-and-fixed-complement}

The selected T correspondence is the explicit pencil of full lattices between two fixed lattices, with its inverse-adjoint CM partner. The positive orthogonal complement transports integrally. This is a concrete modification of the same Hermitian bundle; its finite character action is an equivariant descent datum in the preceding SGA3 module language. The small semisimple correspondence changes the finite/coefficient decoration even when its points agree in a coarse quotient. No identification of a coarse point with its decorated section, no new cycle class, and no change to the existing CM Dieudonne crystal is asserted.

\subsection{85. AM-DO-055: faithful coefficient reading and the prescribed image}\label{am-do-055-faithful-coefficient-reading-and-the-prescribed-image}

The finite orbit identity with reciprocal stabilizer masses is descent of an associated finite coefficient bundle. The full theta section is read faithfully in the fixed Cauchy modules; its terminal contraction is the already fixed bilinear pairing. The SGA7 I VII2.7/2.9 and VIII7.4(b) biextension locators and SGA3 I4.7.3/4.7.3.1 equivariance locators remain the precise framework anchors. CI23 is a conditional nonzero implication on the actual arithmetic image, not a new theorem of motivic nonvanishing or an identification with the specific032 algebraic cycle. The formal-group slopes, Picard construction and CM motive remain those already recorded for DP\_MULT and the Jacobian.

\subsection{86. AM-R-036: an explicit section of the fixed lattice correspondence}\label{am-r-036-an-explicit-section-of-the-fixed-lattice-correspondence}

The fifth-power arithmetic produces a rational unitary element satisfying the actual selected lattice modification and its conjugate dual. It stays in the same decorated scalar orbit; its complex points are separated by the proved J14 estimates. The integral Hermitian-bundle and equivariant descent interpretation uses the preceding lattice dictionary and SGA3 I4.7.3/4.7.3.1. This does not identify a new algebraic theta cycle, imply an ell-adic weight estimate, alter the CM motive or change a Dieudonne slope.

\subsection{87. AM-K060: the real realization and its fixed homogeneous degree}\label{am-k060-the-real-realization-and-its-fixed-homogeneous-degree}

The commuting compact-circle grading and positive rank-one ladders concern the real unitary realization of the existing oscillator/theta module. The selected homogeneous component and its PBW normalization are the same finite coefficient representation used in the earlier symmetric-power/Schur dictionary. The Jacobi estimate is an analytic Hilbert-norm assertion, not a statement of purity for an ell-adic realization. The theta/biextension interpretation still uses SGA7 I VII2.7/2.9 pp169--174 and VIII7.4(b) pp310--311; no new motivic L-function or formal-group consequence is asserted without the remaining arithmetic comparison.

\subsection{88. AM-R-037: standard Schur sections and the real quotient metric}\label{am-r-037-standard-schur-sections-and-the-real-quotient-metric}

The literal integral bideterminant section is the polynomial coefficient of a unit tensor in exterior powers after orthogonal projection to its Cartan summand. Exterior powers and their Schur constituents are the same lambda-operation framework used earlier in this dictionary; the representation-module interpretation uses SGA3 I4.7.3/4.7.3.1. The normalized bound is a metric statement in the complex realization, not integral orthogonal projection, saturation or an ell-adic purity theorem. Finite class-frame lattices and the theta/biextension interpretation remain unchanged, with the existing SGA7 I VII2.7/2.9 and VIII7.4(b) locators. No new crystalline slope or motivic nonvanishing is asserted.

\subsection{89. LIFT45: finite character projectors and arithmetic point separation}\label{lift45-finite-character-projectors-and-arithmetic-point-separation}

The all-place projector is the product of the complements of the finite dilation invariants on the actual transported coefficient sheaf. Its mass-weighted energy inequality is a complex Hermitian realization statement. The changed-place lattice denominator ideal and its CM norm convert the mixed pole of the actual Hecke correspondence into geometric separation. This uses the existing SGA3 I4.7.3/4.7.3.1 coefficient-module and SGA7 I VII2.7/2.9, VIII7.4(b) theta-correspondence dictionary; it proves neither a new weight-purity theorem nor motivic period nonvanishing.

\subsection{90. R038: bounded integral highest-weight sections in the split realization}\label{r038-bounded-integral-highest-weight-sections-in-the-split-realization}

The standard Weyl module is an integral highest-weight representation of the split general linear group scheme. Integral unipotent points preserve that module; after the existing coefficient-field extension the finite cyclotomic grid spans its characteristic-zero realization. Exterior powers provide the lambda-operation tensor model and the normalized archimedean bound. The field-level Fourier extraction is not integral saturation. The representation-module locator remains SGA3 I4.7.3/4.7.3.1; the existing CM/Picard and theta-biextension locators are unchanged. No new Dieudonne slope or ell-adic weight statement follows.

\subsection{91. K061: an actual two-point complex realization}\label{k061-an-actual-two-point-complex-realization}

The genuine-good selected evaluation kernel is summed with its true arithmetic quotient and central Haar factor. Its normalized two-point Gram is positive definite. This is an archimedean automorphic realization statement, not a purity or motivic nonvanishing theorem. In the existing coefficient-sheaf dictionary (SGA3 I4.7.3/4.7.3.1), it proves a nonconstant restriction under one joint finite dilation, while the intersection of every local augmentation factor remains open. The previous CM/abelian Picard and theta-biextension locators remain the same.

\subsection{92. DO056: correspondence adjoints, determinant norms and finite lattices}\label{do056-correspondence-adjoints-determinant-norms-and-finite-lattices}

The punctured Fourier correspondence has its exact adjoint after inverse-automorphism masses are included. Its complex scaled isometry and away-T integral automorphism are distinct realizations of that same correspondence. The normalized T cancellation is a determinant/metric normalization statement; the remaining source period is not assigned a finite valuation. The terminal determinant-line formula keeps the full admitted class dimension. Exterior powers and the product formula transport the already proved scalar content into the determinant line. Existing Picard/biextension references are SGA7 I VII2.7/2.9 and VIII7.4(b); no additional Dieudonne or ell-adic weight conclusion is asserted.

\subsection{93. R039: elementary modifications and the exact integral weight range}\label{r039-elementary-modifications-and-the-exact-integral-weight-range}

The one-covector kernel is an elementary modification of the Hermitian locally free module on Spec(O\_E). Its Schur/Weyl functors have the exact signed scalar conductor given by the maximal highest weight, including their determinant-line twists. This is an integral representation statement in the already located SGA3 I4.7.3/4.7.3.1 and exterior/Schur-operation framework of SGA6 VI3.3. It does not identify the conductor of the theta image. The determinant twist is a line-bundle operation, and the fixed Picard/biextension setting remains FGA232 Theorem3.1 pp148--149, FGA236 Corollary1.7 p229/2.3 p232, SGA7 I VII2.7/2.9 pp169--174 and VIII7.4(b) pp310--311. The beta divisor matches its archimedean norm by the product formula; it gives no new cote4/cote54 scans5--7 crystalline slope, motivic comparison or ell-adic purity theorem.

\subsection{94. DO057: bounded sections of Hermitian coefficient bundles}\label{do057-bounded-sections-of-hermitian-coefficient-bundles}

The actual integral coefficient bundles admit explicit small sections that span after extension to the coefficient field. Independent sections of the determinant line and its dual preserve the signed twists without raising a frame index to the determinant power. The residue frame at11 remains invertible, while global saturation is not asserted. This is a quantitative metric and integral-module construction on the same bundles, with the preceding SGA3/SGA6 coefficient-functor locators and the same FGA/SGA7 Picard and theta-biextension anchors. It removes an unspecified frame-height premise from the conditional MainX image witness, but supplies neither the specific032 cycle nor a nonzero motivic period. No change is made to the established CM motive, formal-group slopes or ell-adic weights.

\subsection{95. LIFT46: character divisors, finite normalization and the different}\label{lift46-character-divisors-finite-normalization-and-the-different}

The algebraic character values define explicit principal divisors on the normalized finite coefficient cover of Spec(O\_K). Their Weil heights use those divisors and the fixed archimedean type. The different estimate is a trace-dual statement, located exactly at Stacks49.8.2--3 and49.12.4--5 (tags0BW0,0BWB,0BWI); it controls ramification without assuming the monogenic order is normal. These coefficient fields and divisor operations refine the existing CM realization, rather than constructing a new CM motive or proving its L-function RH. The fixed Picard and biextension context remains FGA232 Theorem3.1 pp148--149, FGA236 Corollary1.7 p229/2.3 p232, SGA7 I VII2.7/2.9 pp169--174 and VIII7.4(b) pp310--311. The cote4/cote54 scans5--7 slopes and SGA6 VI3.3 Schur operations are unchanged.

\subsection{96. K062: finite incidence correspondences and augmentation traces}\label{k062-finite-incidence-correspondences-and-augmentation-traces}

The actual scalar orbit is an open subset of a projective line of elementary lattice modifications, with precisely two omitted points. Rational isometries induce partial incidence correspondences. Applying the augmentation projector before summation gives the exact trace F−C/(T−1), including negative terms and nonfree orbit masses. This is equivariant module descent in the preceding SGA3 I4.7.3/4.7.3.1 framework, followed by an analytic realization and a proved infinite-tail bound. It does not construct a new algebraic cycle, prove an ell-adic weight estimate or identify the specific032 cycle. The remaining finite signed arithmetic sum is not supplied by the existing FGA/Picard, SGA7/biextension or cote4/cote54 crystal statements.

\subsection{97. LIFT47: trace duality and arithmetic Hermitian frames}\label{lift47-trace-duality-and-arithmetic-hermitian-frames}

The integral vector bundle on Spec(O\_E) retains its actual Hermitian and coefficient structures. The codifferent describes its underlying trace-dual Z-lattice, and the norm of one totally positive element gives a dimension-independent dual minimum. Euclidean transference then produces bounded sections with prescribed residues; it changes no bundle, theta correspondence or polarization. Its exact new external input is Banaszczyk1993, definitions pp625--626 and Theorems2.1--2.2 pp631--632. The arithmetic trace-dual background is the already checked Stacks49.8.2--3/49.12.4--5, and the unchanged determinant/Schur operations are located in SGA6 VI3.3. The FGA232/236 Picard, SGA7 I VII2.7/2.9 and VIII7.4(b) biextension, and cote4/cote54 scans5--7 crystal entries remain those fixed earlier. No archimedean lattice bound is recast as an ell-adic purity or RH assertion.

\subsection{98. DO058: denominator ideals and finite isometry fibres}\label{do058-denominator-ideals-and-finite-isometry-fibres}

The denominator ideal is an effective divisor supported on the selected split T fibres of Spec(O\_E). Hermitian duality pairs the two CM components, giving2s\textless=e\textless=4s. The fractional-ideal product formula bounds scalar reductions; fixing a negative section makes each integral isometry fibre a finite group with faithful reduction at both actual G primes. These are integral-lattice and finite-group statements within the already located SGA3 I4.7.3/4.7.3.1 descent framework. They neither modify the Picard/biextension data of FGA232/236 and SGA7 I VII--VIII nor supply an ell-adic weight theorem. The exact source and new proof locators are DO058 DS3--15; the original Green input remains LIFT42 FC5 and Gorodetsky--Hazan2017 Theorem2.1 p5.

\subsection{99. K063: descent of labelled modification pencils}\label{k063-descent-of-labelled-modification-pencils}

The actual elementary-modification projective line, including its two omitted points, descends to Q(sqrt(-7)). Its augmentation projector subtracts exactly a rank-two endpoint subspace. Finite coefficient descent repeats the same incidence trace at all split primes above a subfield prime; fractional-ideal packing then bounds the complete specified unitary-plane subfamily. This is an explicit module/correspondence realization, not a new cycle identification with032. The source locators are K063 PF6--21, the same SGA3 descent entries above, and the unchanged SGA6 VI3.3 representation operations. Cote4/cote54 scans5--7, FGA Picard and SGA7 biextension interpretations remain in their earlier proved scopes.

\subsection{100. K064: norm congruences on finite local thickenings}\label{k064-norm-congruences-on-finite-local-thickenings}

The fractional coefficient classes are paired by complex conjugation and reduced on the length-two local thickenings O\_L/P². The unitary norm equation becomes an exact determinant-zero condition there, with T\^{}6+T\^{}5−T\^{}3 solutions. The subsequent integral-coset count uses the original arithmetic vector lattice and archimedean metric; it changes no correspondence. Exact locators are K064 NR6--12 and the unchanged SGA3 descent/SGA6 VI3.3 representation entries. This elementary local count supplies no identification with032, no new Picard or biextension theorem, and no global L-function zero statement.

\subsection{101. DO059: a bounded quotient of the actual target modification lattice}\label{do059-a-bounded-quotient-of-the-actual-target-modification-lattice}

Multiplication of the standard negative section by the uniformizer places it in every actual source modification lattice. Its image therefore lies in the uniformizer inverse of the target lattice, a finite over-lattice with exactly T\^{}n residue classes at each CM component. Counting those actual classes gives the improved arithmetic fibre bound, after the analytic selected-moment estimate is sharpened. The exact new locators are DO059 SC24--32; the underlying modification/descent context is the preceding SGA3 entry. The cotangent/crystalline, Picard, biextension and motive dictionaries remain as previously located and are not used to assert the unresolved sign of the remaining sum.

\subsection{102. LIFT48: the metric on the actual biextension theta section}\label{lift48-the-metric-on-the-actual-biextension-theta-section}

The principal theta line and its normalized Hermitian metric give the Zak expansion of its squared norm; the zero Fourier coefficient is exactly the Fock norm. The finite modifications and complete coevaluation then give a bound on the actual arithmetic theta correspondence, with its CM period scalar and stack masses retained. Exact new locators are LIFT48 OB1--28; the underlying biextension/Picard and representation contexts retain the already recorded SGA7 I VII--VIII, FGA232/236 and SGA6 VI3.3 locators. This norm estimate proves no new comparison of motives, no theta-cycle pin and no L-function zero statement.

\subsection{103. K065: uniqueness within an arithmetic coefficient coset}\label{k065-uniqueness-within-an-arithmetic-coefficient-coset}

The transverse section of the fixed rank-two plane is so small at every archimedean embedding that two lifts of one integral coefficient class would contradict the nonzero integral norm. This removes a lift multiplicity while preserving the same local modification scheme and its determinant-zero equation. Exact new locators are K065 TL1--7; the algebraic descent and finite local thickening framework is the preceding SGA3 entry. It is not a new crystalline or motivic comparison.

\subsection{104. R040: combining disjoint strata of the same correspondence}\label{r040-combining-disjoint-strata-of-the-same-correspondence}

The larger-field fixed-plane stratum repeats its nonintegral support seven times, so is disjoint from the entire one-place stratum. R040 R40.1--6 combines their already checked bounds in the original Hermitian kernel and stack mass. The remainder is still a signed arithmetic correspondence. Existing Picard, biextension and motive locators are unchanged; no Grothendieck text is cited as settling the remaining nonvanishing inequality.

\subsection{105. LIFT49: integral Hermitian modifications and their finite discriminant modules}\label{lift49-integral-hermitian-modifications-and-their-finite-discriminant-modules}

An integral short frame identifies each positive selfdual lattice with an overlattice of a bounded Gram lattice. Such overlattices inject into the submodules of its finite discriminant module; local DVR bases and CRT show each submodule has at most the original rank many generators. LIFT49 GM1--9 therefore bounds the actual standard genus mass by a broader count. GM10--21 use the accepted standard genus mass ratios and original strict finite volumes. The earlier SGA3 descent and SGA6 VI3.3 coefficient-lattice context remains; no new representability, crystalline or motivic theorem is claimed. The branch erratum applies the already proved metrized scalar invariance R024.18--21 with the full algebraic scalar, retaining its signed T contribution.

\subsection{106. K066: rational theta section and its metrized realizations}\label{k066-rational-theta-section-and-its-metrized-realizations}

The Q-defined q-normalized section is a section of the rigidified Picard theta line; its canonical L2 norm identifies the common vacuum modulus. Algebraic coordinate generators and their finite norms remain separate from the transcendental complex creation normalization. K066 PC3--21 gives the exact distinction and every coefficient realization. The carried exact locators remain FGA232 Theorem3.1, pp148--149; FGA236 Corollary1.7, p229, and2.3, p232; SGA7 I VII2.7/2.9, pp169--174, andVIII7.4(b), pp310--311, as already recorded. No new comparison of crystalline and Betti periods or motivic L-functions is asserted.

\subsection{107. DO060: an integral subbundle block and finite modifications}\label{do060-an-integral-subbundle-block-and-finite-modifications}

The actual local lattice modifications are supported on five fixed coordinates. Their complementary positive block is integral on both conjugate components. DO060 IB3--19 then passes from integral minors and characteristic coefficients to an arithmetic restriction on the residual unitary action; IB20--32 uses signature one and finite Green characters to bound the full two-place stratum. The underlying coefficient representations and descent have the existing SGA3 I4.7.3/4.7.3.1 and SGA6 VI3.3 locators. The new quantitative argument does not establish a new theta-cycle or motive comparison.

\subsection{108. R041: partition of the same metrized correspondence}\label{r041-partition-of-the-same-metrized-correspondence}

R041 R41.1--6 removes three disjoint arithmetic index sets from the original Hermitian kernel, preserving the same positive stack-mass normalizer. The remaining signed correspondence is explicit but its nonvanishing inequality is open. This is a refinement within the already located Picard/biextension setting, not an added result of Grothendieck's texts or an L-function zero theorem.

\subsection{109. K067: primary residue-field types of the unchanged coefficient representation}\label{k067-primary-residue-field-types-of-the-unchanged-coefficient-representation}

Green's exact finite cuspidal character formula supplies primary-support and degree-dependent bounds at the unmixed finite places. K067 PS4--21 applies it to the actual lattice correspondence, leaving both mixed projectors unchanged. The distinction between a finite residue-field eigenvalue and a global algebraic root is explicit. These are finite representation calculations in the previously recorded SGA3/SGA6 representation setting, not new conclusions about l-adic weights, motives or their L-functions.

\subsection{110. R042: complementary arithmetic strata with their original weights}\label{r042-complementary-arithmetic-strata-with-their-original-weights}

R042 R42.1--6 intersects the new threshold sets with the preceding plane complement. The exact metrized kernel and its stack mass remain fixed, with every removed subset charged by an absolute bound. The existing Picard/biextension locators remain unchanged; the residual arithmetic sign is not inferred from the geometric dictionary.

\subsection{111. LIFT50: integral coefficient lattices and a single algebraic normalization}\label{lift50-integral-coefficient-lattices-and-a-single-algebraic-normalization}

LIFT50 HG1--37 distinguishes Weyl from dual-Weyl lattices, their hook comparison, the algebraic scalar s\_B*C\_D, and the transcendental creation/period norm. The finite contents and archimedean mean use the same algebraic normalization. The actual coefficient lattice definitions are RowForms.lean281--286 and the jet convention DatumVec.lean130--134. The carried geometric locators are SGA6 VI3.3 for the representation/K-theory setting; SGA7 I VII2.7/2.9, pp169--174, and VIII7.4(b), pp310--311 for the rigidified biextension; FGA232 Theorem3.1, pp148--149, and FGA236 Corollary1.7, p229, and2.3, p232 for the Picard setting. No new motivic period comparison or L-function zero statement follows from this norm budget.

\subsection{112. K068: integral exterior minors and positive fixed subspaces}\label{k068-integral-exterior-minors-and-positive-fixed-subspaces}

K068 LR4--22 uses the actual integral block, its exterior-power matrix coefficients and their paired finite valuations to restrict the arithmetic correspondence. The fixed subspace is E-defined and positive, but need not have a bounded integral frame. Its finite reductions are unipotent; this is separate from global unipotence. The exterior/representation operations live in the already cited SGA3 I4.7.3/4.7.3.1 and SGA6 VI3.3 setting. The new quantitative minor and Green-character argument is proved here, not attributed to an unproved weight or motive comparison.

\subsection{113. DO061: traces of exterior powers of the actual lattice modification}\label{do061-traces-of-exterior-powers-of-the-actual-lattice-modification}

DO061 TR4--7 gives pole one for characteristic coefficients of the actual rank-one lattice correspondence, even when individual exterior entries have pole two. The two-trace theorem and paired arithmetic norms then force global unipotence for the remaining three-place terms; the resulting complete sum bound is quantitative arithmetic, not a weight-purity assertion. The exterior and representation constructions retain SGA3 I4.7.3/4.7.3.1 and SGA6 VI3.3. The kernel remains on the same Picard/biextension objects located in FGA232 Theorem3.1, pp148--149, FGA236 Corollary1.7, p229, and2.3, p232, and SGA7 I VII2.7/2.9, pp169--174, andVIII7.4(b), pp310--311. No new formal-crystal or motivic comparison is inferred.

\subsection{114. R043: the same correspondence on a smaller arithmetic index set}\label{r043-the-same-correspondence-on-a-smaller-arithmetic-index-set}

R043 R43.1--9 preserves the Hermitian kernel, finite projectors and stack mass while removing bounded subsets by intersections. The remaining arithmetic correspondence has a positive fixed subspace of dimension at least126 and unipotent finite reductions, with no uniform integral frame for its moving complement. The existing exact Grothendieck locators describe its geometric objects; they do not prove its remaining signed inequality, the specific032 cycle identification or a Kummer L-function zero theorem.

\subsection{115. LIFT51: rank in the coefficient category obstructs the literal Gram links}\label{lift51-rank-in-the-coefficient-category-obstructs-the-literal-gram-links}

LIFT51 LI5--18 keeps full HH separate from projected EE. The source dimension identity and KD=cPu*Pu force a necessary dimension ratio for a small full diagonal defect. This is elementary rank; its representation/K-theory setting remains SGA3 I4.7.3/4.7.3.1 and SGA6 VI3.3. Exact source locators are DataRows.lean62--71; \nolinkurl{RhoDetPrimeEqMuTildeMulChiPow.lean231–233,303–307,345–354;} \nolinkurl{RelevantZeroGramTerms.lean149–151,324–343,509–522;} \nolinkurl{GlobalTheoremInputs.lean55–57,142–156;} and \nolinkurl{Chosen.lean14,53–54,75,79–84.} No new motivic/weight comparison follows. The result refutes preservation of the quantitative links for this cohort, not every mixed-type construction.

\subsection{116. DO062: translated exterior traces and the remaining unipotent correspondence}\label{do062-translated-exterior-traces-and-the-remaining-unipotent-correspondence}

DO062 RU3--9 uses characteristic-polynomial translation in the same exterior representation setting (SGA3 I4.7.3/4.7.3.1; SGA6 VI3.3), with integral finite modifications and all paired valuations retained. The moving positive fixed space is used only at archimedean realizations, where its Hermitian and Euclidean perpendiculars coincide; no integral trivialization is inferred. RU12--13 separates unequal from equal local pencil labels and exhibits the latter's surviving pole2. The existing Picard/biextension locators remain unchanged. This gives no new theorem about motives, l-adic weights, a specific032 cycle or an L-function zero.

\subsection{117. LIFT52: dimensions of the fixed arithmetic coefficient spaces under parameter changes}\label{lift52-dimensions-of-the-fixed-arithmetic-coefficient-spaces-under-parameter-changes}

The accepted cohomological/mass dimension comparison DO044.17--24 remains on the same genuine finite projector and strict initial/maximal bottom levels. LIFT52 PC3--11 varies the weights or a hypothetical split T prime in its exact rational upper bound, then uses convexity. It does not identify the resulting Hilbert-space rank with an unproved algebraic-cycle map. The coefficient representation and rank/K-theory setting retains SGA3 I4.7.3/4.7.3.1 and SGA6 VI3.3. The geometric Picard/biextension locators are unchanged. Changing EE or Mu0 is a change of the actual objects and source ties, not a scalar equivalence, and no new motive or L-function conclusion follows.

\subsection{118. K069: kernels, exterior coordinates and an arithmetic count}\label{k069-kernels-exterior-coordinates-and-an-arithmetic-count}

K069 IC4--8 describes the cross map on the kernel of the integral block and its exact isotropic alternative. IC10--15 uses determinantal charts only for uniqueness and original-coordinate projections only for injection; no inverse minor is declared integral and no moving lattice is trivialized. These are kernel/rank operations in the coefficient representation setting already located at SGA3 I4.7.3/4.7.3.1 and SGA6 VI3.3. The arithmetic count uses the actual finite lattice modification, not a new descent or comparison theorem. It leaves the same Picard/biextension objects, specific032 arithmetic pin and motivic/L-function questions unresolved.

\subsection{119. LIFT53: a finite code constraint on the coefficient-row family}\label{lift53-a-finite-code-constraint-on-the-coefficient-row-family}

LIFT53 CB2--11 compares an injective binary labeling with the literal size required of the row family. This is a cardinality and parameter-quantifier argument, independent of Picard periods or theta-cycle nonvanishing. The coefficient objects retain the existing SGA3/SGA6 representation and rank setting; the precise new source locators are DataS.lean25--27, \nolinkurl{Chosen.lean14,56–57,75–87,} \nolinkurl{ParamsHeightRankConstant.lean111–112,126–140,187–190} and FamCore.lean113--114. No additional geometric or motivic comparison is inferred from code growth.

\subsection{120. LIFT54: the eligible size predicate is stronger than the rank record}\label{lift54-the-eligible-size-predicate-is-stronger-than-the-rank-record}

LIFT54 SL1--16 is an exact numerical and quantifier obstruction inside the displayed source parameter interface, with its complete Lean file/line locators. It does not obstruct the existing schemes, Dieudonne crystals, theta biextensions or mixed CM kernels and does not identify a numerical parameter contradiction with motivic vanishing. The prior FGA232/236 and SGA7 I VII2.7/2.9 interpretations of those objects are unchanged. The new requirement belongs to their attempted quantitative source realization, rather than a new theorem in those cited texts.

\subsection{121. DO063: actual adelic lattices and positive Euclidean counting}\label{do063-actual-adelic-lattices-and-positive-euclidean-counting}

DO063 FV5--19 retains the same finitely modified module over the ring of integers and its paired Hermitian dual. Its covolume and trace-dual statements are arithmetic lattice calculations with the explicit different ideal; no indefinite form is treated as a positive metric. The existing vector-bundle and coefficient-representation interpretation at SGA3 I4.7.3/4.7.3.1 is retained. FV22--34 is the archimedean realization and exact signed correspondence count for the original CM theta object, not a new etale comparison or proof of the specific032 algebraic-cycle pin.

\subsection{122. K070: local incidence and characteristic coefficients}\label{k070-local-incidence-and-characteristic-coefficients}

K070 PT1--10 uses the original two-dimensional residue quotient and its projective line of lattices. An invariant lattice gives integral characteristic coefficients, the same exterior-power determinant data already organized by the SGA6 VI3.3 lambda operations. The argument carefully distinguishes these coefficients from matrix entries in a different lattice frame. Its signed local weights and global norm test are arithmetic statements about this fixed correspondence; no new scheme descent, purity, or motivic L-function result is asserted.

\subsection{123. LIFT55: source scale elimination}\label{lift55-source-scale-elimination}

LIFT55 US1--10 eliminates literal inequalities in ElevenLe.lean and Scale. It excludes the unchanged L343 parameter realization, including after changing ranks and field degree, but does not refute abstract rank objects, the existing CM Picard/theta constructions, or their realizations described in FGA and SGA7. A replacement source interface is a separate theorem to prove. No numerical floor is represented as a geometric or motivic nonexistence theorem.

\subsection{124. K071: nilpotence on a residue quotient and the actual lattice modification}\label{k071-nilpotence-on-a-residue-quotient-and-the-actual-lattice-modification}

K071 NB2--5 uses the actual two-dimensional quotient of paired lattices: nilpotence there makes the square carry the larger lattice into the smaller. It is a statement about integral matrix entries in the specified frame, stronger than the characteristic-coefficient integrality of K070, and does not change the underlying CM theta object. NB12--18 then uses the existing vector-module realization, its different-ideal Euclidean dual and original selected representation; the SGA3 I4.7.3/4.7.3.1 and SGA6 VI3.3 identifications carried in the preceding dictionary entries remain unchanged. The weighted incidence count is a calculation in this particular arithmetic correspondence, not a new descent, purity or motivic L-function theorem.

\subsection{125. DO064: integral rank and specialization of the actual positive block}\label{do064-integral-rank-and-specialization-of-the-actual-positive-block}

DO064 PM6--12 describes the actual endomorphism A of the positive CM lattice. After inverting28 its weighted path/cycle presentation preserves the rank of A−I under every actual G specialization. This is a calculation with this integral module and its fibres; the vector-module/group-scheme dictionary already anchored at SGA3 I4.7.3/4.7.3.1 is unchanged. The fractional ideal count in PM15--23 retains both CM prime partners and the different-ideal Euclidean dual from DO063. It does not construct an arithmetic cycle, a new motivic correspondence, or a pure sheaf whose trace supplies the remaining signed period.

\subsection{126. LIFT56: analytic first occurrence and the limits of the geometric dictionary}\label{lift56-analytic-first-occurrence-and-the-limits-of-the-geometric-dictionary}

LIFT56 VF4--10 uses the exact signed real parameters of the existing CM correspondence. Proper forward Witt-lower theta quotients vanish; the tower constant-term map therefore gives cusp-valued upward correspondences with zero allowed. This is an analytic statement about the realized oscillator representation and its integral transforms, not an identification with a morphism of the Picard schemes or biextensions. The actual CM Jacobian, formal-group/Cartier-dual data and the earlier FGA/SGA7 I VII--VIII locators remain those of the carried dictionary. No purity, motivic L-function, or RH conclusion follows from this real first-occurrence calculation. The positive-rank arithmetic/free-data comparison remains the next interface.

\subsection{127. K072: selected grading and the actual fractional coordinate module}\label{k072-selected-grading-and-the-actual-fractional-coordinate-module}

The stronger energy is the existing homogeneous polynomial degree in the actual K043 PBW realization; it is not a new motivic weight. The active-coordinate module is the original direct sum of fractional CM ideals, with both conjugate primes and its different-ideal Euclidean dual retained. The module and symmetric-power dictionary already anchored in SGA3 I4.7.3/4.7.3.1 and SGA6 VI3.3 is unchanged. Its analytic compression bound does not assert l-adic purity or a new realization comparison.

\subsection{128. DO065: an exact linear-algebraic exception}\label{do065-an-exact-linear-algebraic-exception}

The exceptional direction is controlled by im(N), ker(N) and im(N²) in the existing Hermitian vector bundle over Spec(E). In signature one, the nonzero square has a single isotropic image line. Its position relative to the fixed positive summand forces the exceptional compression A to be the identity. This is a calculation in that vector-module realization and its unitary group, using the carried SGA3 dictionary; it does not construct the arithmetic theta-cycle correspondence demanded by THETA3 G1 or the specific032 pin.

\subsection{129. R044: determinantal rank and the paired CM norm}\label{r044-determinantal-rank-and-the-paired-cm-norm}

R044 PC5 is an equation of the rank-one determinantal locus in the existing Hermitian matrix module. Its two split fibres are exchanged by the CM involution. Counting both valuations is essential: one vanishes at rank one, both at rank zero. The norm of the genuine integral ideal(x)d\_U supplies the global bound without a principal-ideal assumption. This remains within the carried SGA3 I4.7.3/4.7.3.1 vector-module dictionary and is not a new arithmetic cycle, motivic realization, or purity theorem.

\subsection{130. LIFT57: local integral transforms and retained normalizations}\label{lift57-local-integral-transforms-and-retained-normalizations}

The unchanged matrix test functions realize the same local oscillator correspondence in each rank. The G map is a compact finite-type projector; the Q map is an actual previous-image transform whose Schur complement exposes the old determinant-character strings. Their Gaussian-binomial and Haar constants are local measure/comparison factors, not an integral isomorphism or a primitive theta coefficient. The earlier schemes/vector-bundle and lambda-ring descriptions remain unchanged; no new global Picard, biextension or CM-motive map is supplied by these local injections.

\subsection{131. DO066: simultaneous CM-prime congruences in the actual vector module}\label{do066-simultaneous-cm-prime-congruences-in-the-actual-vector-module}

The rank-one determinantal equation on the original S block forces a common orientation at each split fibre. Its CM-conjugate fibre is related by inverse transpose, giving one actual ideal for all cross coordinates. This refines the existing SGA3 I4.7.3/4.7.3.1 vector-module dictionary by arithmetic conditions on its integral sections; it does not construct a new motivic or cycle correspondence. The norm and trace metrics are retained separately.

\subsection{132. K073: the difference-column lattice and paired rank-zero fibres}\label{k073-the-difference-column-lattice-and-paired-rank-zero-fibres}

The zero reduction is imposed at both CM partners. The resulting ideal multiplies the actual target lattice B\_Sigma, and the essential inclusion of the standard lattice permits subtraction of e\_n.~The count uses this genuine integral-module containment, not residue-density division. Its oscillator trace is an analytic metric factor on the same representation, not an l-adic weight or new purity statement.

\subsection{133. LIFT58: global adjunction on actual cusp multiplicity spaces}\label{lift58-global-adjunction-on-actual-cusp-multiplicity-spaces}

The analytic integral transforms are adjoint in their actual Petersson/Haar metrics. Local oscillator quotients and global cuspidal multiplicity maps preserve the prescribed previous images, giving an analytic surjection with exact scalar conventions. This does not provide a primitive integral map of Picard schemes, a biextension trivialization, or an integral inverse on CM motives. The earlier FGA and SGA7 locators describe the unchanged geometric objects; the new assertion lies at their automorphic realization and leaves the arithmetic comparison open.

\subsection{134. DO067 and R045: coefficient ideals and the reduced finite period gate}\label{do067-and-r045-coefficient-ideals-and-the-reduced-finite-period-gate}

Characteristic coefficients stay in the original number field even when a positive complement is chosen only over that field or its complex embeddings. Their paired ideal norms force the first two coefficients to vanish in the stated small-block range. The complete s5 and conditional s6 removals then concern the existing rational-point sum and its original class masses. They reduce DP\_MULT Q3e and THETA3 G1's analytic period gate, without changing the formal-group/crystal, Picard, biextension or CM-motive interfaces and without a CX-ALG-024 Hecke-RH consequence.

\subsection{135. K074: the square of the actual rank-zero ideal}\label{k074-the-square-of-the-actual-rank-zero-ideal}

At a rank-zero split fibre the nilpotent matrix has entries in the paired prime ideal, so every square entry lies in its square. This is multiplication in the existing endomorphism module, viewed through the previously recorded SGA3 I4.7.3/4.7.3.1 vector-module dictionary. The global ideal norm and the exact T pole divisor then constrain the same rational-point count. It is not a new infinitesimal deformation theorem, divided-power gain, biextension comparison or purity result. DP\_MULT Q3e and THETA3 G1 lose the rank-two part of the first remaining uniform sector; the actual rank-one period, source construction and net cost remain.

\subsection{136. K075: content ideals and a relative norm fibre}\label{k075-content-ideals-and-a-relative-norm-fibre}

The exact rank-one nilpotent in the existing endomorphism module is determined by one negative column. Its coordinate-content fractional ideal is retained on Spec(O\_E); the final coordinate lies in a relative-norm fibre over Spec(O\_F). These elementary arithmetic data refine the rational-point count of the unchanged Hermitian correspondence. The recorded SGA3 I4.7.3/4.7.3.1 module interpretation remains applicable. No norm-counting inequality is promoted to a scheme isomorphism, new biextension, Picard comparison or motivic L-function theorem. DP\_MULT Q3e and THETA3 G1 lose further fixed-block terms.

\subsection{137. R046: the whole small-support sector}\label{r046-the-whole-small-support-sector}

The F/IC8 decomposition and the actual nilpotent rank split exhaust every term at supports6 through9. After the new counts the uniform remainder starts at7 and the allowed all-top remainder at9, with the original signed normalization and error\textless2\^{}-230. This is a strict reduction of the conditional arithmetic-period gate, with the same SGA3 vector module and CM-motive objects already indexed. It supplies no algebraic cycle identification, nonzero terminal mass, full source constructor or DP44 sign.

\subsection{138. K076/R047: the metric covolume of the actual content module}\label{k076r047-the-metric-covolume-of-the-actual-content-module}

The arithmetic module on Spec(O\_E) is unchanged, with the recorded SGA3 I4.7.3/4.7.3.1 interpretation. Its positive real metric has covolumeD0 zb\^{}f; the Euclidean dual uses the absolute different and the exact half-trace factor2. Transference and actual ball volume sharpen only its rational-point count. After the checked F/IC8 case split, uniform supportseven and all-top supportnine disappear; the remainders start at8 and10, respectively. The actual signed period, source ties, specific032 cycle and DP44 cost remain. There is no new biextension, Picard, CM-motive or Hecke-RH assertion.

\subsection{139. LIFT59: a nonzero actual automorphic bottom type}\label{lift59-a-nonzero-actual-automorphic-bottom-type}

On the specified standard Hermitian module over Spec(O\_E), the finite regular-series projector is constructed with its full scalar character. The Poincare function supplies a nonzero analytic automorphic realization at the actual finite level; its origin value is controlled in the original normalization. The same SGA3 I4.7.3/4.7.3.1 module and previously indexed CM-motive dictionary describe the underlying spaces. No rational algebraic cycle or integral motivic realization is extracted merely from analytic cusp existence. In particular the nonzero value is not the specific032 Bessel period, a MainX terminal mass or a Hecke L-function zero statement. LIFT58 transports this nonzero type along its prescribed analytic tower, strictly reducing that conditional main gate; arithmetic comparison and DP44 costs remain.

\subsection{140. K077/K078: affine ideal congruences and the ten-place remainder}\label{k077k078-affine-ideal-congruences-and-the-ten-place-remainder}

The local pencil prescribes actual affine cosets in the coordinate module over Spec(O\_E); their difference ideals retain the paired CM places and have the exact additional norm T\^{}(2d). Relative norm fibres are arithmetic data for Spec(O\_E)-\textgreater Spec(O\_F), within the already indexed SGA3 I4.7.3/4.7.3.1 vector-module interpretation. Their weighted count removes actual rational-isometry classes from the unchanged oscillator sum. It creates no new algebraic cycle, biextension map, motive or weight theorem. DP\_MULT Q3e and THETA3 G1 retain the period/arithmetic comparison, now with complete uniform support floor10. CX-ALG-024 is unaffected.

\subsection{141. DO068: local lattice support and analytic oscillator realization}\label{do068-local-lattice-support-and-analytic-oscillator-realization}

Unimodular Hermitian modules and primitive direct summands are integral objects over Spec(O\_E); First2015 Corollary4.9 extends their isometries at the good local rings. The actual rational symplectic polarization gives a forward theta realization, compatible with the earlier biextension/formal-group dictionary but not itself an algebraic correspondence of the requested CM motives. The constructed compact subgroup and finite Schwartz sections are free analytic data. The exact common-vacuum and Fourier quotient normalizations remove a Loc-construction premise from the longer conditional theta passage; they do not identify 032 cycles, integral theta-image lattices, lambda/K-theoretic charges or a new l-adic purity statement. Those global geometric comparisons remain at their recorded gates.

\subsection{142. LIFT60: different-scaled flags and the exact finite row interface}\label{lift60-different-scaled-flags-and-the-exact-finite-row-interface}

The ideals d\^{}\{-(rho-a)\} and their inverses are explicit invertible modules over Spec(O\_E), realized inside the original standard Hermitian lattice when the extended different is principal. This belongs to the existing SGA3 I4.7.3/4.7.3.1 module-and-stabilizer dictionary; the flag and primitive negative line are integral geometric data, not only their real signatures. The finite quotient representation, character and normalized projector are retained exactly, including both CM components. The result supplies no new biextension comparison, l-adic weight assertion or motivic cycle. Its impact on DP\_MULT Q3e/THETA3 G1 is to remove the specified finite Hida/row adapter premise on these hosts; the real relative model and global arithmetic pairing remain.

\subsection{143. Recovered finite reductions K079--K085 and R048--R054: exact arithmetic modules}\label{recovered-finite-reductions-k079k085-and-r048r054-exact-arithmetic-modules}

These findings concern the original Hermitian module over \nolinkurl{Spec(O_E)}, its rational isometries, integral coordinate ideals and the already specified signed oscillator kernel. The module/stabilizer interpretation is the carried SGA3 Exposé I, Proposition4.7.3 and Corollary4.7.3.1 dictionary. Exterior coefficients and the coefficient/flag class retain the carried SGA6 Exposé VI, Theorem3.3 locator. These are framework locators, not attributions of the present numerical bounds to Grothendieck. The full statements and proofs are the individually indexed canonical files in the paper.

\begin{longtable}[]{@{}
  >{\raggedright\arraybackslash}p{(\columnwidth - 2\tabcolsep) * \real{0.5000}}
  >{\raggedright\arraybackslash}p{(\columnwidth - 2\tabcolsep) * \real{0.5000}}@{}}
\toprule\noalign{}
\begin{minipage}[b]{\linewidth}\raggedright
Finding
\end{minipage} & \begin{minipage}[b]{\linewidth}\raggedright
Arithmetic-geometric content and exact current role
\end{minipage} \\
\midrule\noalign{}
\endhead
\bottomrule\noalign{}
\endlastfoot
AM-K079 & Exterior minors of the original fixed-block endomorphism have paired integral divisibility. Eighteen exterior-rank slices are removed, with32739 charged in each family. No rank-one conclusion is inferred. \\
AM-K080 & The small-prime factorization and norm29 restriction give exact one-eighth divisor bounds for the actual relative norm fibre. The complete uniform fixed-W0 support10 contribution has charge1367. The reconstruction has a new identity. \\
AM-R-048 & The all-top selected metric gives H133 after a separately charged annulus. The six original structural cases then give complete all-top support10 removal. It changes no underlying lattice or projector. \\
AM-R-049 & The uniform H8192 six-case exhaustion removes every support10 class, with the additional4093 charge; it does not import the all-top cutoff. \\
AM-K081 & Characteristic coefficients have the exact pole-one type-III divisor and paired trace divisibility. The824-cell arithmetic calculation removes the stated sectors, with4086 in each family. \\
AM-K082 & The same actual coefficient module at H133 removes151 additional all-top cells, charge5001. It is a metric refinement within the original correspondence. \\
AM-R-050 & The nonunipotent rank-one positive block supplies the actual trace/norm witness in the indicated type-III sectors. Charges4094 and4088 apply respectively to the uniform and all-top families. \\
AM-K083 & For global rank at most3, nonunipotence yields a nonzero trace witness in the original number field. Its paired ideal and finite pole table remove78/75 additional cells with797/1166 charges. \\
AM-R-051 & Exhaustion of the original F/IC8, rank and column classes leaves only globally unipotent rank-two matrices with A=I at all-top support11; the extra195718 charge is retained. \\
AM-K084 & The original diagonal coefficient of N squared is a totally positive F-element, with the stated type-I/II/III pole divisor. Its arithmetic norm bound removes76/77 profiles, with92023/93271 charges. Positivity of this coefficient does not assert positivity of a character or period. \\
AM-R-052 & Original-coordinate type-I integrality and the actual scalar-cross radius reduce all-top support11 to i=0 and t=0,1,2; charge34539. \\
AM-R-053 & The corresponding uniform estimate retains H8192 and reduces its globally unipotent rank-two scalar-cross support11 profiles to i=0 and t=0,1,2,3; charge32559. Its other structural sectors are not silently removed. \\
AM-K085 & The actual all-type-II fixed incidence gives an affine real-coordinate slice with its original half-trace covolume and selected homogeneous energy. It removes the last all-top fixed-W0 support11 profile, charge685. \\
AM-R-054 & The exact structural exhaustion combines the previous module counts with C1-corrected DO070--072. Every all-top support11 class is removed; the remaining all-top support is at least12. It adds no error charge, and both full sums of42/31 retained errors remain below2\^{}-230. \\
\end{longtable}

This is a reduction of the rational-point contribution to the existing arithmetic theta pairing, not a construction of a new algebraic correspondence. It strictly reduces the conditional DP\_MULT Q3e and THETA3 G1 gate. The signed sum, literal full family, arithmetic image content, MainX/032 value and DP44 comparison remain open. In particular the class of a vector bundle in K-theory and its exterior operations do not determine the sign of this numerical pairing.

\subsection{144. DO070--DO072, C1 and LIFT063: original lattice metrics and the admitted zero label}\label{do070do072-c1-and-lift063-original-lattice-metrics-and-the-admitted-zero-label}

AM-LIFT-063 proves maximality of the stated tensor integer rings, their trace Gram determinants, original-ring covering radii and \nolinkurl{P0=(1/2)O_F}. These are actual integral modules on \nolinkurl{Spec(O_E)} and \nolinkurl{Spec(O_F)}, with their archimedean metrics. They do not give the same cover to an arbitrary fractional ideal. The exact primary inputs are Milne, \emph{Algebraic Number Theory}, Theorems4.9/6.4; Stacks Project Tag0AP2, Lemma15.45.4; and Thaine (2005), section2, printed209. Their separately recorded source checks are not a new proof of those literature results.

AM-DO-070 uses the actual codifferent/half-trace dual and a finite-dimensional transference theorem for every relevant fractional ideal. It removes the scalar-cross support11 profile i=0,t=2. AM-DO-071 uses the original-coordinate support and the CM norm/trace inequality to remove the t=1 cells with k\textless=107. AM-DO-072 combines two valid bounds for the same arithmetic lattice points: a nonzero Gaussian mass estimate and a Voronoi-volume estimate. It removes t=0 for all k and the remaining t=1 cells k\textgreater=108. Its auxiliary Gaussian is a Euclidean counting device and is not a replacement for the arithmetic theta kernel.

The exact correction AM-DO-070-C1 is part of all three conclusions. The actual pencil parameter scheme has the admitted point0. At a nonzero target its full negative-coordinate image has two poles; at target0 it has one. The unchanged full target determinant therefore produces an extra inverse pole in the real-part-kernel covolume. Complete pair weighting contributes2 per type-III relative place. An active type-II fixed incidence cannot have target0. The revised charges are43330 for DO070 and838 for DO071; the joint DO072 charge1482 is unchanged because its largest envelope occurs at t=0. Historical unqualified projection assertions are explicitly refuted and repaired; earlier event records remain visible.

The covolumes are metric data attached to the same integral modules, not new Picard functor or biextension identities. The carried formal/crystal/Cartier-duality locators remain Grothendieck cote4, scan003 §§8.1.1--8.1.2 and scans005--006 §§8.4.1--8.4.3; cote54 scans005--007. The polarized Jacobian and theta/Picard line retain FGA, Bourbaki232 Theorem3.1, pp.148--149; Bourbaki236 Corollaries1.7/2.3, pp.229/232; and SGA7 I, Exposé VII §§2.7/2.9, pp.169--174. These previously qualified framework locators are carried unchanged. No new assertion about their existence theorems is proved by the point counts.

\subsection{145. DO069 and LIFT061/061-U/062/062-T: analytic realizations with exact source ties}\label{do069-and-lift061061-u062062-t-analytic-realizations-with-exact-source-ties}

AM-DO-069 proves the positive-rank free-data implication for the genuine R014/R015 oscillator realization, using R016 and DO068. Both directions retain the actual splitting characters, polynomial/parity convention, full flag coordinates, uncontracted continuity, compact-finite pure tensors and the pointwise supplied realization identity. This is an analytic implication on a specified realization. It does not create arithmetic \nolinkurl{Frames.Good}, certify a norm budget or prove the initial global period nonzero.

AM-LIFT-061 constructs the actual relative real Wu model under the stated discrete-series/NCGaps hypotheses. The induced Hilbert projection, compact reproducer, decay estimate and tensor assembly keep the original Haar and coefficient conventions. Its nonzero seed has energy E\textgreater0, but E and the resulting inverse-square-root cost are not evaluated. The genuine discrete-series minimal K-type has multiplicity one, so the literal source lowKSub equals that full K-module; no highest-weight condition is inserted into IsLowUnit. AM-LIFT-061-U identifies the local projected line and proves uniqueness up to phase without estimating its costs.

AM-LIFT-062 gives the matrix-unit reconstruction of the entire finite-projected bottom multiplier space and proves its cuspidality under the actual arithmetic hypotheses. Its converse uses the source Sobolev continuity, a continuous Hilbert extension and smooth representatives in the stated order. It retains the full projector and the actual quotient-Haar ratio. The last comparison with an independently supplied arithmetic space is conditional on its actual rowTie witness.

AM-LIFT-062-T transports evaluation by the literal finite tS, with \nolinkurl{ConjLevel(tS^-1)}, before taking the fixed spaces. The finite projector is transported in a common principal-level ambient space, preserving complex coefficient conjugation and the exact coset-volume cancellation. It proves the translated EvalTieT under the genuine comparison fields. It neither discards tS nor constructs the requested concrete STiesR/032 instance; any unshifted RowsJoinPinned comparison remains to be proved where required.

These results live in the automorphic/archimedean realization of the same Hermitian modules and polarized CM data. The carried geometric locators are SGA3 I4.7.3/4.7.3.1 for the modules/stabilizers, FGA232 Theorem3.1 and FGA236 Corollaries1.7/2.3 for the Picard framework, and SGA7 I VII2.7/2.9 for the biextension framework; the cote4/cote54 crystal/duality locators in §144 are unchanged. No primitive integral morphism of Picard schemes, new biextension trivialization, rational algebraic cycle or motivic realization map is extracted from Hilbert-space reconstruction alone. Likewise no equality in the SGA6 VI3.3 coefficient/flag dictionary, no l-adic weight assertion, and no motivic L-function zero statement follows. DP\_MULT Q3e and THETA3 A7/G1 lose the stated analytic construction premises; GCT\_PLETH T-052 retains its global arithmetic charge comparison, and CX\_GRH CX-ALG-024 retains the Hecke-RH gate.

\subsection{146. LIFT064: a negative full finite factor with the same global unitary invariants}\label{lift064-a-negative-full-finite-factor-with-the-same-global-unitary-invariants}

The formula h(2y) in GP3 is a polynomial family of integral unitary matrices in the original coordinates. Thus it gives points of the existing unitary group scheme over the integer base, with its rank-two unipotent endomorphism and positive coefficient alpha. CRT prescribes its fibres at the actual split G primes. One relative place has rank one, the others have rank two or the specified mixed identity, and Green's exact character formula makes the complete normalized finite factor negative.

This is a counterexample to an implication based only on the listed global invariants and paired ideal divisibility. It is not a counterexample in the actual nonintegral T-pencil short sector. Algebraically, the invariant F-element alpha has a conjugation-invariant E-ideal and square absolute E-norm without forcing even relative valuations. The finite character sign is not determined by those ideal properties. The exact module/stabilizer framework remains SGA3 I4.7.3/4.7.3.1; no stronger statement about the Picard/biextension/crystal objects at the carried FGA/SGA7/cote4/cote54 locators follows. There is no new lambda-ring, K-theory, sheaf-weight, CM-motive or L-function theorem. The conditional Q1/Q2 signed main gate is unchanged; any successful restricted sign argument must use more of its literal T-incidence or archimedean conditions.

\subsection{147. DO073: the actual positive rank-five stabilizer module}\label{do073-the-actual-positive-rank-five-stabilizer-module}

The fixed original-coordinate column determines the rational subspace U\_x. Its orthogonal complement after adjoining the negative line is a positive five-dimensional E-space. Intersecting that space with the actual source lattice gives a full projective O\_E-module on Spec(O\_E). The finite automorphism group preserves this module; its localizations at the actual completely split G primes are free rank five and reduce faithfully. This is the exact arithmetic module and stabilizer interpretation in the carried SGA3 I4.7.3/4.7.3.1 framework. The construction supplies no globally integral orthogonal moving frame or height estimate for such a frame.

The resulting finite group-order bound removes the specified globally unipotent scalar-cross support12 profiles, reducing the existing DP\_MULT Q3e/THETA3 G1 finite-period gate without changing its sign normalization. It does not give full support12 exhaustion or an actual DP44 conclusion. The carried Picard and biextension locators remain FGA232 Theorem3.1, FGA236 Corollaries1.7/2.3 and SGA7 I VII2.7/2.9; the formal/crystal/duality locators remain cote4 scan003 §§8.1.1--8.1.2, scans005--006 §§8.4.1--8.4.3 and cote54 scans005--007. No new lambda/K-theory identity at SGA6 VI3.3, l-adic weight theorem, algebraic cycle identification, CM-motive comparison or Hecke L-function result is inferred from this finite stabilizer count.

\subsection{148. K086, DO074 and LIFT065: smaller integral and realization gates}\label{k086-do074-and-lift065-smaller-integral-and-realization-gates}

AM-K086 uses the actual intersection of two unequal pencil lattices. Because N squared annihilates S, both N and the inverse displacement send each original S-vector into that intersection; the paired adjoint makes the opposite cross block integral too. Thus the S-column denominator divisor is supported only at typeII. This is an inclusion of the original modules on Spec(O\_E), within the carried SGA3 I4.7.3/4.7.3.1 module/stabilizer framework. The corrected pair-weighted coarea keeps2\^{}(i+t)T\^{}i; a Gaussian/Voronoi count then removes the90 non-allII support12 globally rank-two-unipotent scalar-cross profiles, with one885 charge. No new module is substituted for the actual fractional ideals, and no positive sign is inferred.

AM-DO-074 concerns characteristic coefficients in that same original field/module, with the same exact typeIII pole divisor. It first forces the trace coefficient to vanish. The signature-one reciprocal-pair geometry then lowers the second coefficient bound before a second ideal-norm contradiction. The resulting r\textless=22,t\textless=10 rigidity is an actual rational-isometry classification, at the stated all-top metric cutoff. The carried exterior-coefficient/flag framework is SGA6 VI Theorem3.3, together with SGA3 I4.7.3/4.7.3.1. This is not an assertion that a K-theory class or lambda-operation alone determines analytic unipotence or period nonvanishing.

AM-LIFT-065 identifies the actual Read range and metric scale under its stated genuine analytic comparison. The row isometry has scale(c*nP)\^{}(-1/2). Its optional extension is a finite-dimensional Hilbert-space construction with both old and target finiteness supplied; it is not an integral isometry of arithmetic theta-image lattices. An actual Read and simultaneous spectral/column compatibility still have to be constructed for the requested032 package.

For all three findings the polarized CM Jacobian, Picard and biextension objects retain the previously qualified exact locators: FGA232 Theorem3.1, pp.148--149; FGA236 Corollaries1.7/2.3, pp.229/232; SGA7 I VII2.7/2.9, pp.169--174. Formal-group/crystal/Cartier-duality locators remain cote4 scan003 §§8.1.1--8.1.2, scans005--006 §§8.4.1--8.4.3 and cote54 scans005--007. No new algebraic cycle, Picard or biextension comparison, l-adic weight theorem, CM-motive identity or motivic L-function conclusion follows. DP\_MULT Q3e and THETA3 G1 gain the stated conditional finite/source reductions; CX\_GRH CX-ALG-024 and GCT\_PLETH T-052 retain their distinct unresolved comparisons.

\subsection{149. DO075 and R055: paired congruence modules and restricted denominator divisors}\label{do075-and-r055-paired-congruence-modules-and-restricted-denominator-divisors}

AM-DO-075 identifies160 fixed original coordinates at every rank-two unmixed G reduction of an IC8 block of rank1/2. It uses the image line of N squared and the paired inverse-adjoint relation, not positivity in a finite field. The resulting paired ideal enters the original fractional-ideal module count; convexity in its norm removes the all-top support12 sector with one1859541 charge. AM-R-055 combines K086's actual pencil-intersection integrality with complete support-dependent covolumes and removes118062 higher-support scalar-cross unipotent profiles with one1263 charge. Both are module/congruence statements over Spec(O\_E) and estimates on their genuine archimedean realizations. The exact carried module/stabilizer locator is SGA3 I4.7.3/4.7.3.1; the separate exterior-coefficient context is SGA6 VI Theorem3.3. Neither count substitutes a K-theory class for actual lattice coordinates, metrics or point counts.

\subsection{150. LIFT066: a genuine spectral refill is not an integral realization theorem}\label{lift066-a-genuine-spectral-refill-is-not-an-integral-realization-theorem}

AM-LIFT-066 constructs identity spectral intertwiners only after the source's reS operation installs the genuine Poincare/term/pair operators. The retained Haar scalar is identified by uniqueness against one positive finite-measure set, and the exact local-model isomorphisms remain inputs. The normalized row map retains(c*nP)\^{}(-1/2), every translate and the actual columns. The unchanged/replaced table in SR18 preserves all arithmetic comparison obligations. This is a theorem about the analytic realization of a supplied arithmetic package, not an isomorphism of integral theta-image modules or an automatic cycle-class comparison.

For these three findings the polarized CM Jacobian, its Picard schemes and biextensions retain FGA232 Theorem3.1 pp.148--149; FGA236 Corollaries1.7/2.3 pp.229/232; SGA7 I VII2.7/2.9 pp.169--174. The crystal/Cartier-duality anchors remain cote4 scan003 §§8.1.1--8.1.2, scans005--006 §§8.4.1--8.4.3, and cote54 scans005--007. These are carried contextual locators, not newly proved motivic consequences. No new lambda-ring identity, l-adic weight theorem, CM-motive isomorphism or motivic L-function identity follows. DP\_MULT Q3e and THETA3 G1 have the stated smaller finite/spectral gates; CX\_GRH CX-ALG-024 and GCT\_PLETH T-052 retain their distinct open arithmetic comparisons.

\subsection{151. K087 and DO076: finite determinantal ideals and a reciprocal metric}\label{k087-and-do076-finite-determinantal-ideals-and-a-reciprocal-metric}

AM-K087 reads the actual finite rank of h−I into the norm of the original paired/oriented congruence ideal: each relative G factor contributes q\^{}(2−d). A uniform upper cap on that ideal norm and a monotone full-coordinate envelope remove the sector with rank sum at least70. The remaining265 identity reductions are a statement about the original integral model at closed points, not a sheaf-weight or motivic positivity theorem. The module/stabilizer context is the carried SGA3 I4.7.3/4.7.3.1, and the exterior/determinantal context is SGA6 VI Theorem3.3.

AM-DO-076 uses the rational nontrivial primary plane of a global-rank-two nonunipotent isometry. The original scalar compression produces a rational isotropic eigenline and lambda in F; inversion makes its norm at least one for an upper count. The resulting real-linear metric change has its determinant retained exactly. It does not replace the arithmetic Hermitian lattice or assert that the square roots of absolute eigenvalues are algebraic units. This removes a specified nonunipotent sector with the original pairing and source lattice unchanged.

For both findings the geometric anchors remain the carried polarized-Jacobian/Picard/biextension locators FGA232 Theorem3.1 pp.148--149; FGA236 Corollaries1.7/2.3 pp.229/232; SGA7 I VII2.7/2.9 pp.169--174, and the crystal/duality references cote4 scan003 §§8.1.1--8.1.2, scans005--006 §§8.4.1--8.4.3 and cote54 scans005--007. No new lambda-ring identity, integral K-theory identification, l-adic weight theorem, CM-motive comparison or motivic L-function consequence is proved. The finite gates for DP\_MULT Q3e and THETA3 G1 shrink in the stated sectors; the distinct CX\_GRH CX-ALG-024 and GCT\_PLETH T-052 arithmetic gates remain.

\subsection{152. DO077, R056 and K088: determinantal reductions, real metrics and integral Hermitian pairings}\label{do077-r056-and-k088-determinantal-reductions-real-metrics-and-integral-hermitian-pairings}

DO077 reads the rank of the original finite reduction of h−I into its actual normalized Green trace; twenty-two ranks at least three supply a uniform extra factor without creating a new subset enumeration. R056 then removes the full remaining nonunipotent scalar-cross sector by covering the archimedean realization with finitely many product-one real metrics. Its dual-minimum argument preserves the arithmetic module and its covolume; these metric changes are not integral automorphisms or algebraic units.

K088 uses an actual h-stable projective selfdual Hermitian module patched from the type-II local lattices. The original integral test vectors belong to it at both CM partners. Their square pairings are integral F-elements in a determinantal G-ideal. The exact norm contradiction confines the nilpotent square image, after which one original matrix entry forces the additional numerator divisor U. The low-rank all-II sector is empty, with no new numerical error.

These arguments concern the original integral-model closed fibres and real realization. Their carried module/stabilizer context is SGA3 I4.7.3/4.7.3.1; the exterior/determinantal context is SGA6 VI Theorem3.3. The unchanged polarized Jacobian, Picard and biextension locators are FGA232 Theorem3.1 pp.148--149; FGA236 Corollaries1.7/2.3 pp.229/232; SGA7 I VII2.7/2.9 pp.169--174. The formal-crystal/Cartier-dual background remains cote4 scan003 §§8.1.1--8.1.2, scans005--006 §§8.4.1--8.4.3, and cote54 scans005--007. None is invoked as a new source for analytic positivity. No new lambda-ring identity, K-theory comparison, l-adic weight theorem, CM-motive isomorphism or motivic L-function assertion follows. The actual DP\_MULT Q3e and THETA3 G1 finite/real gates shrink; CX\_GRH CX-ALG-024 and GCT\_PLETH T-052 retain their separate arithmetic gates.

\subsection{153. LIFT067: the character on a flagged unitary group and its modulus}\label{lift067-the-character-on-a-flagged-unitary-group-and-its-modulus}

The physical norm-minus-T vector determines the final two root coefficients of the literal additive character. Their values(bar(pi),pi) are sections of the actual flag-root coordinate lines; conjugating by the rational flag diagonal D identifies their character with the unit-marked convention. The representation vector becomes D*tau and the unipotent Haar measure acquires exactly q\footnote{-rho(k+rho+1)}. This is a rational-group/analytic-realization identity preserving the integral level data, not an equality of integral polarizations or arithmetic row forms. The compact-stabilizer obstruction concerns an altered character/vector pair and does not make the actual period vanish.

The source is HidaSetting hJ/rootCoord/hA/Nfin/psiFin at24,39--41,81--82,105--115, together with LIFT60 HF3--6/HF12 and LIFT30 NL3--12. The geometric background uses the same carried Picard/biextension and crystal locators listed in §152. No new l-adic sheaf weight, lambda/K-theory equality, CM-motivic period comparison or integer DP44 norm budget is supplied by character standardization alone. The missing fixed local value and global Bessel/Mram period are explicitly retained in DP\_MULT Q3e and THETA3 G1; the specific032 cycle pin remains distinct.

\subsection{154. DO078: two actual selfdual modules and an integral square pairing}\label{do078-two-actual-selfdual-modules-and-an-integral-square-pairing}

The arbitrary-arrow argument uses two distinct projective Hermitian modules over Spec(O\_E), the actual source and target of the labelled correspondence. An original integral vector in both gives an integral pairing with hy in the target and with h\^{}-1y in the source. Their sum yields an integral F-valued square coefficient. This does not assert one stable lattice or identify different labelled lattices. The finite-rank determinantal ideals and original numerator-U divisibility then give the exact empty low-rank sector; the complementary high-rank estimate retains the real coarea and every zero-target incidence.

The carried module/stabilizer and exterior-power contexts are SGA3 I4.7.3/4.7.3.1 and SGA6 VI Theorem3.3. The underlying polarized CM Jacobian, Picard and biextension objects remain those at FGA232 Theorem3.1 pp.148--149; FGA236 Corollaries1.7/2.3 pp.229/232; SGA7 I VII2.7/2.9 pp.169--174. Crystal/Cartier-dual context remains cote4 scan003 §§8.1.1--8.1.2, scans005--006 §§8.4.1--8.4.3, and cote54 scans005--007. No new lambda-ring, K-theory, l-adic weight, CM-motive or motivic L-function comparison is inferred. The actual DP\_MULT Q3e/THETA3 G1 finite signed gate is smaller; CX\_GRH CX-ALG-024 and GCT\_PLETH T-052 retain their unresolved arithmetic interfaces.

\subsection{155. R057: complete support-twelve exhaustion in the original realization}\label{r057-complete-support-twelve-exhaustion-in-the-original-realization}

The joint theorem combines the exact original F/IC8 partition, fixed-block and scalar-cross alternatives, global-unipotence rank classification, integral Hermitian square-pairing reductions, and real metric counts. Its fresh F,r1 application uses the complete conditional unipotent count, explicitly beyond the old t\textless=8 removed union. The arithmetic modules, labelled correspondences, selected representations, class masses and CM phases remain unchanged; the actual normalized energy now has only support\textgreater=13 terms and a complete52-charge error below2\^{}-230.

The module/stabilizer and determinantal contexts remain SGA3 I4.7.3/4.7.3.1 and SGA6 VI Theorem3.3. Picard/polarization/biextension locators remain FGA232 Theorem3.1 pp.148--149; FGA236 Corollaries1.7/2.3 pp.229/232; SGA7 I VII2.7/2.9 pp.169--174. The formal-crystal and Cartier-dual background is cote4 scan003 §§8.1.1--8.1.2, scans005--006 §§8.4.1--8.4.3, and cote54 scans005--007. These are carried contextual anchors, not sources of the new numerical theorem. No new lambda/K-theory identity, l-adic weight result, CM-motive identification or motivic L-function implication is asserted. DP\_MULT Q3e and THETA3 G1/A7 have a strictly smaller finite/real gate; their prescribed period and arithmetic cost remain open. CX\_GRH CX-ALG-024 and GCT\_PLETH T-052 keep their distinct missing arithmetic comparisons.

\subsection{156. LIFT068: the conditional fixed local Bessel realization}\label{lift068-the-conditional-fixed-local-bessel-realization}

The universal matrix recurrence is an identity in a split local representation realization: row and column compact idempotents, integral flag pivots, their determinant characters, and the Haar character of the physical unipotent radical. It fixes a vector and character rather than defining a new algebraic cycle. The exact integral root transport is compatible with the split group/flag viewpoint of SGA3 I4.7.3/4.7.3.1. The class-map and theta-line comparison still use the carried Picard/biextension anchors FGA232 Theorem3.1 pp.148--149; FGA236 Corollaries1.7/2.3 pp.229/232; SGA7 I VII2.7/2.9 pp.169--174. These are contextual locators, not sources of the new JPSS evaluation.

The formal-crystal background remains cote4 scan003 sections8.1.1--8.1.2, scans005--006 sections8.4.1--8.4.3, and cote54 scans005--007 for Cartier duality. LIFT068 changes none of its divisibility slopes. It gives no new K-theory or lambda-ring identity, l-adic sheaf/weight theorem, motive, or motivic L-function conclusion. For DP\_MULT Q3e and THETA3 G1/A7, the prescribed finite local test is now evaluated under explicit generic tempered premises; global nonvanishing, the canonical local/global normalization and the actual arithmetic Read/032 comparison remain. CX\_GRH CX-ALG-024 and GCT\_PLETH T-052 retain their distinct unproved comparisons.

\subsection{157. K089: primitive actual theta sections and their full principal parts}\label{k089-primitive-actual-theta-sections-and-their-full-principal-parts}

On the good product abelian scheme over the coefficient DVR at11, the prime-to11 multiplication cover supplies an integral theta-group representation. Its matrix units and the actual primitive trace-one section pin an integral isomorphism, not just a projective comparison. Faithfully flat splitting tests the content of the original section. The finite principal-parts map is a split injection with a unit maximal minor, whereas the fixed selected ordinary differential/class map is a later morphism whose reduction kernel is still unknown.

The precise carried principal-parts locator is EGA IV4 Theorem16.11.2 and Remark16.11.3(ii), pp54--55; the factorial distinction remains. The Picard/theta-line context is FGA232 Theorem3.1 pp148--149, FGA236 Corollaries1.7/2.3 pp229/232, and SGA7 I VII2.7/2.9 pp169--174, with VIII7.4(b) pp310--311 retained only in its stated scope. Mumford1966 section1 Proposition3/Theorem2 pp295--298 and section4 Theorem1/corollary pp338--340 provide the actual finite theta-group and very-ampleness inputs. Formal crystals and Cartier duality retain cote4 scan003 sections8.1.1--8.1.2, scans005--006 sections8.4.1--8.4.3, and cote54 scans005--007. No change in the height-three divided-power slope is inferred from unit algebraic section content.

The same selected Schur classes remain in the SGA6 VI3.3 lambda/K-theory setting. The statement is an integral realization and special-fibre section theorem; it supplies no new l-adic weight, motivic projector or motivic L-function conclusion. It removes an intermediate content uncertainty in DP\_MULT Q3e and the arithmetic theta realization of THETA3 G1/A7, while their selected image, actual pairing value and cost remain open. CX\_GRH CX-ALG-024 and GCT\_PLETH T-052 retain their separate unproved maps.

\subsection{158. LIFT069: a chosen automorphic realization with simple tempered transfer}\label{lift069-a-chosen-automorphic-realization-with-simple-tempered-transfer}

The chosen finite cuspidal representation is a genuine summand of the finite reductive representation ring; its occurrence at a split closed point excludes proper parabolic induction. The actual auxiliary theta realization identifies the same local parameter by two compatible classifications. Exact cancellation of its common character string proves a simple generic source parameter. The integral odd-rank HC grid supplies a dominant algebraic coefficient representation, and the cited cohomological temperedness theorem then applies. This is a representation and realization statement, not a construction of a new CM motive or an identification of the automorphic parameter with the motive of J.

The established group/weight and coefficient-class locators remain SGA3 I4.7.3/4.7.3.1 and SGA6 VI3.3. The same FGA232/236 and SGA7 I VII--VIII Picard/biextension locators in section157, and the cote4/cote54 formal-crystal locators there, remain unchanged contextual frameworks. Caraiani2012 Theorem1.2 p2312 is the cited finite temperedness input; no new purity theorem or RH statement is claimed. DP\_MULT Q3e and THETA3 G1/A7 gain the larger finite generic tempered premise of the prescribed local test on the chosen LIFT059 subspace. Their initial representation, global period, arithmetic Read and032 value remain distinct. No consequence for CX\_GRH CX-ALG-024 or GCT\_PLETH T-052 is obtained without their required arithmetic comparisons.

\subsection{159. DO079: an extremal Schur weight and a realization-level obstruction}\label{do079-an-extremal-schur-weight-and-a-realization-level-obstruction}

The selected polynomial representation is the same Schur functor of the fixed module. After removing its determinant character, the minimum-coordinate weight space is the Schur functor on the remaining coordinates; its exact Weyl dimension ratio is an algebraic representation statement. The compact-weight and exterior/Schur constructions retain SGA3 I4.7.3/4.7.3.1 and SGA6 VI Theorem3.3 as their established group and lambda/K-theory contexts. The actual complex Hilbert realization then gives the quantified compression obstruction and the weighted-trace inequality. This neither changes the underlying integral lattice nor proves a new statement about its special-fibre selected image.

The CM/Picard and biextension context remains FGA232 Theorem3.1 pp148--149; FGA236 Corollaries1.7/2.3 pp229/232; SGA7 I VII2.7/2.9 pp169--174 and VIII7.4(b) pp310--311. The formal-crystal/Cartier-dual context is still cote4 scan003 sections8.1.1--8.1.2, scans005--006 sections8.4.1--8.4.3, and cote54 scans005--007. The proof yields no new l-adic sheaf, weight, motivic decomposition or L-function conclusion. For DP\_MULT Q3e and THETA3 G1/A7, the proposed uniform operator shortcut is excluded but the normalized-trace route stays open. The finding is relevant to the complex oscillator realization of GCT\_PLETH T-052 only in this explicit selected-model sense; its arithmetic comparison and CX\_GRH CX-ALG-024 remain unresolved.

\subsection{160. DO079-C1: corrected selected Schur weight}\label{do079-c1-corrected-selected-schur-weight}

The fibre representation is fixed across words, while the selected highest weight changes by the stated letter. The corrected coordinate difference preserves the extremal Schur dimension bound. The exact SGA3 I4.7.3/4.7.3.1 and SGA6 VI3.3 locators, and all FGA/SGA7/cote4/cote54 locators in section159, retain precisely the same contextual scope. This correction adds no new arithmetic, l-adic, K-theoretic or motivic conclusion.

\subsection{161. LIFT070: an actual coefficient grid and its tempered automorphic realization}\label{lift070-an-actual-coefficient-grid-and-its-tempered-automorphic-realization}

The literal source labels provide a common non-top embedding, and the coefficient-weight inequalities hold for every original word. The actual odd eta character shifts each GL block onto its own integral coefficient grid. The representation and Schur-weight framework remains SGA3 I4.7.3/4.7.3.1 and SGA6 VI3.3. The finite strict covector stabilizer and its unweighted augmentation condition remain the original integral module maps; conductor-one genericity is proved after temperedness, not read from a characteristic-zero idempotent alone.

The unchanged Picard/biextension context is FGA232 Theorem3.1 pp148--149; FGA236 Corollaries1.7/2.3 pp229/232; SGA7 I VII2.7/2.9 pp169--174 and VIII7.4(b) pp310--311. Crystal/Cartier duality remains cote4 scan003 sections8.1.1--8.1.2, scans005--006 sections8.4.1--8.4.3, and cote54 scans005--007. The result applies cited cohomological temperedness to an actual automorphic realization; it constructs no new CM motive, l-adic comparison or purity theorem. DP\_MULT Q3e and THETA3 G1/A7 gain the initial local premise for the fixed physical T period in this genuine fully sized branch. Their source existence, global period, actual arithmetic value and DP44 cost remain open. CX\_GRH CX-ALG-024 and GCT\_PLETH T-052 gain no unconditional arithmetic conclusion from this local premise alone.

\subsection{162. DO080: the true Hermitian realization and its radial trace bound}\label{do080-the-true-hermitian-realization-and-its-radial-trace-bound}

The selected algebraic compact type retains its original Schur realization and metric, in the established SGA3 I4.7.3/4.7.3.1 and SGA6 VI3.3 framework. Its complex homogeneous Hermitian bundle gives a positive probability density from the full finite-rank orthogonal projector. The new trace inequality is an analytic realization theorem using invariant measure, not an algebraic positivity theorem on the correspondence or a new lambda-ring identity. Both compact factors and every PBW metric scalar remain in the actual representation.

The original theta line and polarization retain FGA232 Theorem3.1 pp148--149; FGA236 Corollaries1.7/2.3 pp229/232; SGA7 I VII2.7/2.9 pp169--174 and VIII7.4(b) pp310--311. The formal-crystal/Cartier-dual layer remains cote4 scan003 sections8.1.1--8.1.2, scans005--006 sections8.4.1--8.4.3, and cote54 scans005--007. No new integral comparison, l-adic weight, motive or L-function theorem follows. This supplies a stronger analytic input for the actual DP\_MULT Q3e/THETA3 G1 finite kernel; a complete arithmetic count is still needed for any smaller remainder. GCT\_PLETH T-052 keeps its separate arithmetic comparison, and CX\_GRH CX-ALG-024 has no RH consequence here.

\subsection{163. R058: source/target integral pairings and actual selected-trace removal}\label{r058-sourcetarget-integral-pairings-and-actual-selected-trace-removal}

The new fixed-W0 sector theorem combines integral pairings on distinct original Hermitian source and target modules with the actual selected complex trace. Its square numerator ideal includes the exact original-coordinate U factor after the low finite-rank norm contradiction. No single h-stable lattice or artificial integral moving frame is introduced. The full all-profile count removes supports13..23 in the stated rank-two-unipotent block and retains the same normalized kernel, phases and class weights.

The carried module and determinantal anchors are SGA3 I4.7.3/4.7.3.1 and SGA6 VI Theorem3.3. The CM Picard/theta/biextension context is FGA232 Theorem3.1 pp148--149; FGA236 Corollaries1.7/2.3 pp229/232; SGA7 I VII2.7/2.9 pp169--174 and VIII7.4(b) pp310--311. Formal crystals and Cartier duality retain cote4 scan003 §§8.1.1--8.1.2, scans005--006 §§8.4.1--8.4.3 and cote54 scans005--007. These are contextual locators, not sources of the new numerical theorem. No new lambda-ring/K-theory equality, l-adic purity or weight result, CM-motive identification or motivic L-function assertion follows. DP\_MULT Q3e and THETA3 G1/A7 have a smaller finite gate; every final arithmetic period/cost remains open. CX\_GRH CX-ALG-024 and GCT\_PLETH T-052 retain their separate missing comparisons.

\subsection{164. K090: the actual augmentation module and saturated auxiliary theta image}\label{k090-the-actual-augmentation-module-and-saturated-auxiliary-theta-image}

The strict scalar-decoration cover gives a free finite group torsor on the existing class set. Over the coefficient DVR, its unweighted norm kernels are the tensor product of integral augmentation ideals. The split difference basis is defined without inverting the group order, and the literal diamond action preserves the coefficient lattice and underlying Picard/Hom data. Actual finite support evaluations give a split restricted injection, transported through the fixed prime-to11 theta reading and descended by faithful flatness. All assertions precede selected polynomial derivatives and class/projector contractions.

The Picard and theta-group context is FGA232 Theorem3.1 pp148--149; FGA236 Corollaries1.7/2.3 pp229/232; SGA7 I VII2.7/2.9 pp169--174 and VIII7.4(b) pp310--311; Mumford(1966) §1 Proposition3, Theorem2, pp295--298, carried through K089. Module/stabilizer and exterior-power contexts remain SGA3 I4.7.3/4.7.3.1 and SGA6 VI Theorem3.3. Formal crystals and Cartier duality retain cote4 scan003 §§8.1.1--8.1.2, scans005--006 §§8.4.1--8.4.3 and cote54 scans005--007. No new lambda-ring or K-theory equality, l-adic purity/weight theorem, CM-motive identification or motivic L-function conclusion is inferred. The remaining DP\_MULT Q3e and THETA3 G1/A7 arithmetic image is the literal selected alternating matrix; CX\_GRH CX-ALG-024 and GCT\_PLETH T-052 still require their separate arithmetic comparisons.

\subsection{165. DO081: integral scalar pairings on actual source and target modules}\label{do081-integral-scalar-pairings-on-actual-source-and-target-modules}

The new scalar pairings belong to the integral Hermitian form on each actual target module, while the inputs belong to the corresponding source. Their separate membership at both CM components removes T denominators without asserting one stable lattice. The original isotropic-plane alternative and integer Bezout row then give an exact complementary count. All original theta, Picard, coefficient and realization data are retained.

The carried contextual locators are SGA3 I4.7.3/4.7.3.1; SGA6 VI Theorem3.3; FGA232 Theorem3.1 pp148--149; FGA236 Corollaries1.7/2.3 pp229/232; SGA7 I VII2.7/2.9 pp169--174 and VIII7.4(b) pp310--311; cote4 scan003 §§8.1.1--8.1.2, scans005--006 §§8.4.1--8.4.3 and cote54 scans005--007. These are not attributions of the new numeric bound. No new lambda/K-theory equality, l-adic purity/weight theorem, motive identification or L-function conclusion follows. DP\_MULT Q3e and THETA3 G1/A7 have the stated smaller finite gate, while prescribed arithmetic period/cost, CX\_GRH CX-ALG-024 and GCT\_PLETH T-052 retain their missing comparisons.

\subsection{166. R059: rank-one cross blocks and integral source/target reconstruction}\label{r059-rank-one-cross-blocks-and-integral-sourcetarget-reconstruction}

On the original Hermitian modules, identity S compression forces a rank-one isotropic cross block C. Integral transverse pairings, or the primitive original-plane scalar row, together with one original pivot reconstruct C over E. Unitarity and the original W0 negative column then reconstruct the full negative column. No divided coordinate is declared integral, no global nilpotence is imposed, and no lattice or CM realization changes. The finite orientation count uses only the actual unipotent G reductions.

Context remains SGA3 I4.7.3/4.7.3.1; SGA6 VI Theorem3.3; FGA232 Theorem3.1 pp148--149; FGA236 Corollaries1.7/2.3 pp229/232; SGA7 I VII2.7/2.9 pp169--174 and VIII7.4(b) pp310--311; cote4 scan003 §§8.1.1--8.1.2, scans005--006 §§8.4.1--8.4.3 and cote54 scans005--007. No new lambda/K-theory equality, l-adic weight theorem, CM-motive identification or L-function assertion follows. The actual finite gate toward DP\_MULT Q3e and THETA3 G1/A7 is smaller; every final arithmetic period/cost and the separate CX\_GRH CX-ALG-024/GCT\_PLETH T-052 comparisons remain open.

\subsection{167. R060: integral exhaustion and the support15 remainder}\label{r060-integral-exhaustion-and-the-support15-remainder}

The unchanged CM Hermitian realization admits a finer finite-incidence decomposition: every support13/14 stratum is now charged using actual paired ideal conditions, source/target integral lattices and exact Green factors. The fresh exceptional-block estimate is a norm-ideal count; its endpoint retains Q\^{}(-655(f−1)). No scheme, CM field, polarization, theta map or arithmetic normalization changes. The resulting smaller analytic remainder is a gate toward DP\_MULT Q3e and THETA3 G1/A7, not a period value or a DP44 sign.

Carried context: SGA3 I4.7.3/4.7.3.1; SGA6 VI Theorem3.3; FGA232 Theorem3.1 pp148--149; FGA236 Corollaries1.7/2.3 pp229/232; SGA7 I VII2.7/2.9 pp169--174 and VIII7.4(b) pp310--311; cote4 scan003 §§8.1.1--8.1.2, scans005--006 §§8.4.1--8.4.3 and cote54 scans005--007. These are contextual locators previously carried, not a new primary-source audit. No new lambda-ring/K-theory identity, l-adic weight theorem, motivic comparison or L-function consequence is inferred.

\subsection{168. DO082: a numerator ideal and the actual selected trace}\label{do082-a-numerator-ideal-and-the-actual-selected-trace}

On the actual all-II plane, the squared column ideal is (J\_F O\_E)I\_Sigma², with norm J²T\^{}(-4s). The original integral entry further places alpha in UJ\_F I\_Sigma,F². Its norm, together with the unique totally positive alpha recovered from B, supplies a joint ideal-lattice and selected-Hilbert-trace bound. This is a statement on the fixed integral Hermitian modules and CM embeddings; it does not promote a product-norm inequality to separate embedding inequalities.

The preceding exact carried SGA/FGA/cote locators provide the same algebraic context. The charge895 excludes the complete global-rank2-unipotent scalar-cross support23 sector in the actual analytic kernel. It proves no surviving selected arithmetic value, concrete032 cycle identity, DP44 saving, CX\_GRH CX-ALG-024 RH premise, or GCT\_PLETH T-052 charge comparison.

\subsection{169. AM-DO-083 --- unchanged integral realizations}\label{am-do-083-unchanged-integral-realizations}

RO1--18 remove all globally rank-one fixed-W0 terms at supports15..20 with one1953 charge. Full type/pair/double-pole counts and actual selected trace retained. This theorem alone does not exhaust any full support.

The same original Hermitian modules, CM realizations and integral Weyl lattices remain. Carried contextual locators: SGA3 I4.7.3/4.7.3.1; SGA6 VI Theorem3.3; FGA232 Theorem3.1 pp148--149; FGA236 Corollaries1.7/2.3 pp229/232; SGA7 I VII2.7/2.9 pp169--174 and VIII7.4(b) pp310--311; cote4 scan003 §§8.1.1--8.1.2, scans005--006 §§8.4.1--8.4.3 and cote54 scans005--007. These are carried framework anchors, not new primary audits or new lambda/K-theory, l-adic weight, motivic comparison or L-function theorems. DP\_MULT Q3e and THETA3 G1/A7 retain the stated actual arithmetic gates; CX\_GRH CX-ALG-024 and GCT\_PLETH T-052 remain distinct.

\subsection{170. AM-R-061 --- unchanged integral realizations}\label{am-r-061-unchanged-integral-realizations}

GS1--23 remove the entire scalar-cross A=I class at support22, without global unipotence, with one78919 charge. Same-J scalar pivot, actual selected trace, fixed kerC rank5 fibre and full target count retained.

The same original Hermitian modules, CM realizations and integral Weyl lattices remain. Carried contextual locators: SGA3 I4.7.3/4.7.3.1; SGA6 VI Theorem3.3; FGA232 Theorem3.1 pp148--149; FGA236 Corollaries1.7/2.3 pp229/232; SGA7 I VII2.7/2.9 pp169--174 and VIII7.4(b) pp310--311; cote4 scan003 §§8.1.1--8.1.2, scans005--006 §§8.4.1--8.4.3 and cote54 scans005--007. These are carried framework anchors, not new primary audits or new lambda/K-theory, l-adic weight, motivic comparison or L-function theorems. DP\_MULT Q3e and THETA3 G1/A7 retain the stated actual arithmetic gates; CX\_GRH CX-ALG-024 and GCT\_PLETH T-052 remain distinct.

\subsection{171. AM-K091 --- unchanged integral realizations}\label{am-k091-unchanged-integral-realizations}

SR1--13 give an integral left inverse for the literal selected Weyl substitution and its full symbolic two-edge composition. No retraction of the actual derivative evaluation, geometric class sum or terminal operations is asserted.

The same original Hermitian modules, CM realizations and integral Weyl lattices remain. Carried contextual locators: SGA3 I4.7.3/4.7.3.1; SGA6 VI Theorem3.3; FGA232 Theorem3.1 pp148--149; FGA236 Corollaries1.7/2.3 pp229/232; SGA7 I VII2.7/2.9 pp169--174 and VIII7.4(b) pp310--311; cote4 scan003 §§8.1.1--8.1.2, scans005--006 §§8.4.1--8.4.3 and cote54 scans005--007. These are carried framework anchors, not new primary audits or new lambda/K-theory, l-adic weight, motivic comparison or L-function theorems. DP\_MULT Q3e and THETA3 G1/A7 retain the stated actual arithmetic gates; CX\_GRH CX-ALG-024 and GCT\_PLETH T-052 remain distinct.

\subsection{172. LIFT071: the exact local realization metric of the positive-rank T correspondence}\label{lift071-the-exact-local-realization-metric-of-the-positive-rank-t-correspondence}

The result calibrates a local Hermitian pairing on the already specified automorphic realization of the integral correspondence. It keeps the coefficient modules, scalar character and finite-level lattice unchanged. The determinant weights come from invariant quotient measures, and are not new arithmetic-image divisibility. The original module/exterior-power framework is SGA3 I4.7.3/4.7.3.1 and SGA6 VI3.3. Picard and biextension anchors remain FGA232 Theorem3.1 pp148--149; FGA236 Corollaries1.7/2.3 pp229/232; SGA7 I VII2.7/2.9 pp169--174 and VIII7.4(b) pp310--311. Formal groups/crystals/Cartier duality retain cote4 scan003 §§8.1.1--8.1.2, scans005--006 §§8.4.1--8.4.3, and cote54 scans005--007. These carried contextual locators are not sources for the new analytic scalar. No new scheme identification, lambda-ring or K-theory identity, l-adic sheaf/weight statement or CM motive/L-function theorem is inferred. DP\_MULT Q3e and THETA3 G1/A7 lose this local normalization uncertainty only under the genuine source premises; their global period, arithmetic pin and cost remain open. CX\_GRH CX-ALG-024 and GCT\_PLETH T-052 receive no unconditional global consequence.

\subsection{173. DO084: one original integral rank-one chart ball}\label{do084-one-original-integral-rank-one-chart-ball}

The original integral rank-one327-entry pivot chart lies in one Frobenius ball. Its ring-volume count replaces only the old independent-entry chart count, retaining every original seven-column, Green, source-fibre, support and pencil factor. The241 fresh F,r1 nonunipotent cells of FV1--18 contribute\textless2\^{}-4255. Together with the old R050/DO074 cuts the entire nonunipotent F,r1 sector through support18 is removed, with the listed higher-support tails. This finding alone does not raise the complete support floor.

In algebraic-geometric terms this counts original integral Hermitian endomorphism charts and labelled correspondences without changing their coefficient modules, polarizations or Picard/biextension data. The carried module/determinantal anchors are SGA3 I4.7.3/4.7.3.1 and SGA6 VI3.3. Picard/biextension anchors are FGA232 Theorem3.1 pp148--149; FGA236 Corollaries1.7/2.3 pp229/232; SGA7 I VII2.7/2.9 pp169--174 and VIII7.4(b) pp310--311. Crystal/Cartier duality retains cote4 scan003 §§8.1.1--8.1.2, scans005--006 §§8.4.1--8.4.3 and cote54 scans005--007. These are contextual locators; no new scheme, lambda-ring/K-theory identity, l-adic sheaf/weight theorem, CM motive or motivic L-function claim follows. DP\_MULT Q3e and THETA3 G1/A7 obtain precisely the stated smaller analytic remainder. CX\_GRH CX-ALG-024 and GCT\_PLETH T-052 retain their distinct missing arithmetic comparisons.

\subsection{174. R062: full-column congruences and the support18 signed remainder}\label{r062-full-column-congruences-and-the-support18-signed-remainder}

Every actual all-top support15,16,17 class is removed by the unchanged K069 F/IC8 and fixed-S/scalar-cross partition. A new full-column count for all support17 terms with every unmixed finite rank\textgreater=3 contributes\textless2\^{}-177258. It retains original ideals, full target-overlattices, all pairs and G\_n, with no A or pattern multiplicity. Together with DO084 and carried complete classes the all-top remainder has support\textgreater=18. All61 chronological all-top and31 uniform errors have exact sum\textless2\^{}-230. Uniform floor11 and every global arithmetic gate remain open.

In algebraic-geometric terms this counts original integral Hermitian endomorphism charts and labelled correspondences without changing their coefficient modules, polarizations or Picard/biextension data. The carried module/determinantal anchors are SGA3 I4.7.3/4.7.3.1 and SGA6 VI3.3. Picard/biextension anchors are FGA232 Theorem3.1 pp148--149; FGA236 Corollaries1.7/2.3 pp229/232; SGA7 I VII2.7/2.9 pp169--174 and VIII7.4(b) pp310--311. Crystal/Cartier duality retains cote4 scan003 §§8.1.1--8.1.2, scans005--006 §§8.4.1--8.4.3 and cote54 scans005--007. These are contextual locators; no new scheme, lambda-ring/K-theory identity, l-adic sheaf/weight theorem, CM motive or motivic L-function claim follows. DP\_MULT Q3e and THETA3 G1/A7 obtain precisely the stated smaller analytic remainder. CX\_GRH CX-ALG-024 and GCT\_PLETH T-052 retain their distinct missing arithmetic comparisons.

\subsection{175. DO085: minimal finite rank and original paired column congruences}\label{do085-minimal-finite-rank-and-original-paired-column-congruences}

In the actual IC8 class at block ranks3..31, equality of full and compressed finite rank forces every untouched original column and the paired row to be fixed. The resulting paired ideal and exact Green jump jointly improve both endpoint counts. All182 listed fresh rank/support cells contribute\textless2\^{}-7548, including the full rank3 class through support23. Every original incidence, source fibre and error remains; the complete all-top support floor stays18.

This is a determinantal rank statement on the reductions of the original integral Hermitian module, coupled to counting its original-coordinate columns. The carried module/determinantal anchors are SGA3 I4.7.3/4.7.3.1 and SGA6 VI3.3. The Picard/biextension context is FGA232 Theorem3.1 pp148--149; FGA236 Corollaries1.7/2.3 pp229/232; SGA7 I VII2.7/2.9 pp169--174 and VIII7.4(b) pp310--311. Formal groups/crystals/Cartier duality retain cote4 scan003 §§8.1.1--8.1.2, scans005--006 §§8.4.1--8.4.3 and cote54 scans005--007. No new scheme, lambda-ring/K-theory identity, l-adic sheaf or weight theorem, CM motive comparison or motivic L-function statement is inferred. DP\_MULT Q3e and THETA3 G1/A7 gain exactly these smaller remaining sets; CX\_GRH CX-ALG-024 and GCT\_PLETH T-052 retain their distinct global arithmetic gaps.

\subsection{176. AM-DO-086: original integral charts and untouched cross images}\label{am-do-086-original-integral-charts-and-untouched-cross-images}

The full actual F class with block rank1 at every support18..94 has total absolute contribution\textless2\^{}-8622, without global-unipotence or type-III restrictions. One original nonzero integral entry controls paired finite rank-zero loss; DO084 original chart geometry and every genuine seven-coordinate/fibre/pair factor are retained. This removes the entire F,r1 class through94 together with the old complete floor18, but supplies no full support-floor theorem by itself.

The original integral Hermitian modules, coefficient lattices and labelled correspondences remain unchanged. One original integral endomorphism entry or an untouched rank-one isotropic image is counted; no moving integral basis or arithmetic theta-image inverse is introduced. The carried module/determinantal anchors are SGA3 I4.7.3/4.7.3.1 and SGA6 VI3.3. Picard/biextension anchors remain FGA232 Theorem3.1 pp148--149; FGA236 Corollaries1.7/2.3 pp229/232; SGA7 I VII2.7/2.9 pp169--174 and VIII7.4(b) pp310--311. Crystals/Cartier duality retain cote4 scan003 §§8.1.1--8.1.2, scans005--006 §§8.4.1--8.4.3 and cote54 scans005--007. These are contextual locators; no new scheme identification, lambda-ring/K-theory identity, l-adic sheaf or weight theorem, CM motive or motivic L-function assertion follows. DP\_MULT Q3e and THETA3 G1/A7 gain precisely the stated smaller sufficient remaining sets; all actual source, selected arithmetic image, prescribed period,032 value and DP44 costs remain open. CX\_GRH CX-ALG-024 and GCT\_PLETH T-052 retain their distinct arithmetic gaps.

\subsection{177. AM-R-063: original integral charts and untouched cross images}\label{am-r-063-original-integral-charts-and-untouched-cross-images}

The entire actual IC8 block-rank1/2 class at supports18..21 has total absolute contribution\textless2\^{}-52215. The cross map on160 untouched original coordinates has rank0/1; its zero/transverse/plane alternatives and nine retained negative-column coordinates reconstruct the whole column. The original source fibre counts every A without extra block enumeration. All original targets, paired primes, incidence weights and finite Green factors remain. This is an entire-class bound without global unipotence, not a standalone full support-floor theorem.

The original integral Hermitian modules, coefficient lattices and labelled correspondences remain unchanged. One original integral endomorphism entry or an untouched rank-one isotropic image is counted; no moving integral basis or arithmetic theta-image inverse is introduced. The carried module/determinantal anchors are SGA3 I4.7.3/4.7.3.1 and SGA6 VI3.3. Picard/biextension anchors remain FGA232 Theorem3.1 pp148--149; FGA236 Corollaries1.7/2.3 pp229/232; SGA7 I VII2.7/2.9 pp169--174 and VIII7.4(b) pp310--311. Crystals/Cartier duality retain cote4 scan003 §§8.1.1--8.1.2, scans005--006 §§8.4.1--8.4.3 and cote54 scans005--007. These are contextual locators; no new scheme identification, lambda-ring/K-theory identity, l-adic sheaf or weight theorem, CM motive or motivic L-function assertion follows. DP\_MULT Q3e and THETA3 G1/A7 gain precisely the stated smaller sufficient remaining sets; all actual source, selected arithmetic image, prescribed period,032 value and DP44 costs remain open. CX\_GRH CX-ALG-024 and GCT\_PLETH T-052 retain their distinct arithmetic gaps.

\subsection{178. AM-R-064: exhaustive original-module support partition}\label{am-r-064-exhaustive-original-module-support-partition}

The original full all-top support18,19,20 classes are exhausted by the two complete structural tables J21.1--4. The actual all-top remainder now has support at least21, with the unchanged64/31 chronological error arrays and both sums below2\^{}-230. No new error charge is introduced. The independent uniform support floor remains11. The sufficient signed inequality, prescribed global period, arithmetic image and source family, specific032 value and DP44 cost comparison remain open.

The assembly combines rank conditions on the reductions of original integral Hermitian modules, the fixed-space/isotropic-image partition, and their original labelled theta correspondences. It changes no lattice, coefficient module, polarization, finite level, Picard scheme, biextension or CM realization. The root-coordinate reconstructions in the cited counts are over E and are not integral inverses of the arithmetic theta map.

The carried module/stabilizer context is SGA3 I4.7.3/4.7.3.1 and the exterior/determinantal context is SGA6 VI Theorem3.3. The Picard/biextension anchors remain FGA232 Theorem3.1 pp148--149; FGA236 Corollaries1.7/2.3 pp229/232; SGA7 I VII2.7/2.9 pp169--174 and VIII7.4(b) pp310--311. Formal groups/crystals/Cartier duality retain cote4 scan003 §§8.1.1--8.1.2, scans005--006 §§8.4.1--8.4.3 and cote54 scans005--007. These are contextual locators, not sources of this new support exhaustion. No new scheme identification, lambda-ring/K-theory identity, l-adic sheaf/weight theorem, CM motive comparison or motivic L-function result is asserted. The change to DP\_MULT Q3e and THETA3 G1/A7 is precisely J21.1. CX\_GRH CX-ALG-024 and GCT\_PLETH T-052 keep their distinct open arithmetic maps.

All new quantitative inputs have separately checked full proofs: DO085, DO086 and R063. The remaining cited inputs are the already checked K069 partition/table, DO064 ranks32/33 table, DO067 nonunipotence criterion, K081/K082 tables, DO083, R058, R059/R061 and R062. The present certificate checks every finite table index and complete historical error list, not occurrence of a rational matrix. Whole human-language proof reading is separately recorded. No Lean build or archive audit is performed.

\subsection{179. AM-DO-087: actual divisor and selected-trace composition}\label{am-do-087-actual-divisor-and-selected-trace-composition}

The entire actual all-top H133 fixed-W0 globally rank-one unipotent sector at support21 has absolute contribution less than2\^{}-1649. The exact K080 divisor constants and relative-F norm count are composed with DO083's actual selected Hilbert trace, all intrinsic double-pole sets, all weighted incidence types and paired finite factors. Global unipotence is an explicit hypothesis. No standalone complete-floor, signed nonvanishing, source-family or DP44 conclusion is made.

This is an original-coordinate count of the rank-one endomorphism and its negative column, including exact content ideals, their affine cosets, the relative-F norm equation and the labelled theta incidences. No moving integral basis, arithmetic theta-image inverse or integral norm bound for a selected Hilbert vector is introduced. Contextual module/determinantal anchors remain SGA3 I4.7.3/4.7.3.1 and SGA6 VI Theorem3.3; Picard/biextension anchors remain FGA232 Theorem3.1 pp148--149, FGA236 Corollaries1.7/2.3 pp229/232, SGA7 I VII2.7/2.9 pp169--174 and VIII7.4(b) pp310--311. Formal groups, crystals and Cartier duality retain cote4 scan003 §§8.1.1--8.1.2, scans005--006 §§8.4.1--8.4.3 and cote54 scans005--007. These locators are contextual; no new scheme identification, lambda-ring/K-theory identity, l-adic weight/sheaf theorem, CM-motive comparison or motivic L-function result is asserted. DP\_MULT Q3e and THETA3 G1/A7 gain precisely this additional containing-set bound. CX\_GRH CX-ALG-024 and GCT\_PLETH T-052 retain their separate open arithmetic maps.

\subsection{180. LIFT072: prescribed local positive-rank G realization}\label{lift072-prescribed-local-positive-rank-g-realization}

For the actual positive-rank split G model, every prescribed smaller inverse-generic K\^{}1-fixed vector admits a bottom vector inside the same admitted generic finite Pi\_0 with canonical Hilbert Bessel period1/d\_varsigma. This includes the actual mixed regular-series types. With literal source Haar the value is c\_H\emph{c\_N/d\_varsigma. The normalized local factor remains1/(d\_varsigma}Lscript), and the product bound is greater than Q\^{}(-f\emph{k}(k+1)/2). The genuine LIFT069 chosen-bottom and LIFT070 full-Sizes hypotheses remain. No global period, integral bound for the Hilbert vector, actual Read/032 value, source family or DP44 cost is supplied.

The finite extension is an intertwiner between generic representation spaces on the specified residue-group quotient of the original Hermitian module. Its exact compact character and selfdual Haar comparison concern the analytic realization of that same labelled correspondence. They give no integral inverse or saturation of the arithmetic theta image. The carried module/exterior-power context is SGA3 I4.7.3/4.7.3.1 and SGA6 VI Theorem3.3; Picard and biextension anchors remain FGA232 Theorem3.1 pp148--149, FGA236 Corollaries1.7/2.3 pp229/232, SGA7 I VII2.7/2.9 pp169--174 and VIII7.4(b) pp310--311. Formal groups/crystals/Cartier duality retain cote4 scan003 §§8.1.1--8.1.2, scans005--006 §§8.4.1--8.4.3 and cote54 scans005--007. These contextual locators imply no new scheme identification, lambda-ring/K-theory identity, l-adic sheaf or weight theorem, CM-motive comparison or motivic L-function statement. DP\_MULT Q3e and THETA3 G1/A7 lose the stated local G test obstruction under the retained source premises; all global arithmetic maps and costs remain open. CX\_GRH CX-ALG-024 and GCT\_PLETH T-052 gain no unconditional global conclusion.

\subsection{181. AM-R-065: complete original-module support21 partition}\label{am-r-065-complete-original-module-support21-partition}

The complete actual all-top support21 class is exhausted by all F/IC8 ranks1..33 and every A=I alternative. The fixed-W0 nonunipotent case uses K083's explicit t21 low-rank cell and K079's rank\textgreater=4 sector; the new DO087 bound supplies the full globally rank-one unipotent case. The actual all-top support floor is22, with unchanged65/31 chronological error arrays, both sums\textless2\^{}-230, and no new charge. The uniform floor remains11. The signed inequality and all source/arithmetic/global-period/DP44 gates remain open.

The argument partitions rank conditions and fixed/isotropic submodules of the original integral Hermitian modules and their labelled correspondences. It changes no lattice, polarization, coefficient module, finite level, Picard scheme, biextension or CM realization. No rational-coordinate reconstruction is promoted to an integral arithmetic-image inverse.

The carried module/stabilizer context is SGA3 I4.7.3/4.7.3.1, and the exterior/determinantal context is SGA6 VI Theorem3.3. The Picard/biextension anchors remain FGA232 Theorem3.1 pp148--149; FGA236 Corollaries1.7/2.3 pp229/232; SGA7 I VII2.7/2.9 pp169--174 and VIII7.4(b) pp310--311. Formal groups/crystals/Cartier duality retain cote4 scan003 §§8.1.1--8.1.2, scans005--006 §§8.4.1--8.4.3 and cote54 scans005--007. These are contextual locators, not proofs of this assembly. No new scheme identification, lambda-ring/K-theory identity, l-adic sheaf/weight result, CM motive comparison or motivic L-function theorem is asserted. DP\_MULT Q3e and THETA3 G1/A7 gain precisely J22.1; CX\_GRH CX-ALG-024 and GCT\_PLETH T-052 retain their distinct arithmetic gates.

Every quantitative containing-set theorem used here has its own accepted whole proof. The new input DO087 has root and independent whole-proof/script checks and separate358-assertion replays. The present whole proof checks exhaustiveness and exact scope; its certificate is a finite table/error diagnostic and does not enumerate rational support or compute a period. No Lean build or archive audit is performed.

\subsection{182. AM-DO-088: actual affine metric and weighted finite patterns}\label{am-do-088-actual-affine-metric-and-weighted-finite-patterns}

The entire actual all-top H133 globally rank-one unipotent fixed-W0 sector at support22 has absolute contribution\textless2\^{}-358. The exact affine difference-ideal metric and its half-trace dual give a new ball count; both weighted envelopes retain a common pre-cap subset ratio, summing every paired G rank-zero pattern with a product\textless2. The actual same packet distance makes the total multiplicative near-point loss\textless2. All23 integer marked sizes and original type/pair costs are included. Global unipotence is a hypothesis; other structural sectors and the signed/global/arithmetic/DP44 gates remain open.

The new metric is the explicit positive realization of an original fractional-ideal module; its dual uses the absolute different and half-trace factor2. The affine cosets, integral endomorphism column and labelled incidences are unchanged. No moving integral basis, arithmetic theta-image inverse or selected integral-unit vector is introduced. The carried module/exterior context is SGA3 I4.7.3/4.7.3.1 and SGA6 VI Theorem3.3. Picard/biextension anchors remain FGA232 Theorem3.1 pp148--149; FGA236 Corollaries1.7/2.3 pp229/232; SGA7 I VII2.7/2.9 pp169--174 and VIII7.4(b) pp310--311. Crystals and Cartier duality retain cote4 scan003 §§8.1.1--8.1.2, scans005--006 §§8.4.1--8.4.3 and cote54 scans005--007. These are contextual locators; no new scheme, lambda-ring/K-theory, l-adic weight/sheaf or CM-motive/L-function assertion follows. DP\_MULT Q3e and THETA3 G1/A7 gain this precise containing-set removal; CX\_GRH CX-ALG-024 and GCT\_PLETH T-052 keep their separate arithmetic gates.

\subsection{183. AM-K092: primitive full-mass symbolic cutoff}\label{am-k092-primitive-full-mass-symbolic-cutoff}

The actual full middle-stabilizer cutoff can preserve the common11-unit witness and gives a primitive symbolic selected theta family on the integral tensor-augmentation source. Its contracted selected value equals the literal full-mass-divided contribution, with every coordinate pullback, ordinary derivative, Chern gauge, creation/period scalar and terminal projector retained. Full mass division is asserted only after the invariant middle coevaluation. The same input occurs in every middle-class summand; derivative annihilation and class cancellation remain possible.

This is an integral finite-module splitting on the unchanged theta-group and Picard/Hom data. The common transversal is conjugated by the actual class/polarization operator, and inverse Weil transport also transports evaluation functionals. No division by a possibly nonunit stabilizer order occurs in the integral section construction. The actual Weyl module and tensor augmentation are split before the selected derivative; faithful-flat descent concerns the symbolic cokernel only. Mumford(1966), §1 Proposition3 and Theorem2 pp295--298, remains carried through K089. The carried exact contextual anchors are FGA232 Theorem3.1 pp148--149; FGA236 Corollaries1.7/2.3 pp229/232; SGA7 I VII2.7/2.9 pp169--174 and VIII7.4(b) pp310--311; SGA3 I4.7.3/4.7.3.1; SGA6 VI Theorem3.3. Formal crystals and Cartier duality retain cote4 scan003 §§8.1.1--8.1.2, scans005--006 §§8.4.1--8.4.3 and cote54 scans005--007. No new lambda-ring/K-theory identity, l-adic purity or weight theorem, or CM-motive/L-function assertion is inferred. DP\_MULT Q3e and THETA3 G1/A7 retain the actual arithmetic comparison gates; CX\_GRH CX-ALG-024 and GCT\_PLETH T-052 retain their separate gates.

\subsection{184. AM-R-066: joint original fractional-ideal ball}\label{am-r-066-joint-original-fractional-ideal-ball}

The entire actual IC8 rank-one/rank-two exceptional class at support22 has absolute contribution\textless2\^{}-160208, with no global-unipotence hypothesis. The original untouched160-coordinate block has rank one, so its Hilbert--Schmidt and operator norms agree and place all coordinates in one53760-dimensional ball. The exact ideal-product dual, covering radius, volume and original column/type/pair multiplicities are retained. The complete all-top support floor remains22; this alone does not exhaust support22.

This is a bound on the original product of160 fractional-ideal modules in its fixed positive Euclidean realization. The absolute different, half-trace dual and discriminant are unchanged; the rank-one cross map is a restriction/projection of the original labelled integral endomorphism. Banaszczyk(1993), Theorem2.2 p632, is carried through K076 with no new claim of an integral image inverse. The carried exact contextual anchors are FGA232 Theorem3.1 pp148--149; FGA236 Corollaries1.7/2.3 pp229/232; SGA7 I VII2.7/2.9 pp169--174 and VIII7.4(b) pp310--311; SGA3 I4.7.3/4.7.3.1; SGA6 VI Theorem3.3. Formal crystals and Cartier duality retain cote4 scan003 §§8.1.1--8.1.2, scans005--006 §§8.4.1--8.4.3 and cote54 scans005--007. No new lambda-ring/K-theory identity, l-adic purity or weight theorem, or CM-motive/L-function assertion is inferred. DP\_MULT Q3e and THETA3 G1/A7 retain the actual arithmetic comparison gates; CX\_GRH CX-ALG-024 and GCT\_PLETH T-052 retain their separate gates.

\subsection{185. LIFT073: actual finite spherical normalization}\label{lift073-actual-finite-spherical-normalization}

Under the genuine actual representation/lattice/character hypotheses, every prescribed finite away-G/T spherical unit-vector Bessel factor is alpha\_v=c\_H,v\emph{c\_N,v}Lscript\_v, and alpha\_v\^{}natural=c\_H,v*c\_N,v\textgreater0. The actual fractional different-scaled flag gives primitive coordinates even at absolute ramification. Every actual dyadic place is CM-split; a new full spherical descending-row support argument supplies the split calculation in all residue characteristics. The odd nonsplit calculation uses the precise primitive-coordinate literature formula. Together with LIFT071/072 this determines the finite local tests conditionally; global period/nonzero, source existence, integral image, Read/STiesR,032 value and DP44 remain open.

The unchanged selfdual Hermitian module over the finite etale quadratic algebra defines a smooth reductive unitary model, descending from GL\_n after that base change; its integral points are the stated hyperspecial compact. The different-scaled fractional flag and actual unit marked-line isometry preserve the prescribed Picard/Hom correspondence. This is a local analytic normalization and integral-model statement, not an integral theta-image inverse or a comparison of global motives. Exact contextual anchors are SGA3 I4.7.3/4.7.3.1 and SGA6 VI Theorem3.3; FGA232 Theorem3.1 pp148--149; FGA236 Corollaries1.7/2.3 pp229/232; SGA7 I VII2.7/2.9 pp169--174 and VIII7.4(b) pp310--311. Formal crystals/Cartier duality retain cote4 scan003 §§8.1.1--8.1.2, scans005--006 §§8.4.1--8.4.3 and cote54 scans005--007. No new lambda-ring/K-theory equality, l-adic sheaf/weight assertion, CM-motive identification or motivic L-function theorem is obtained. DP\_MULT Q3e and THETA3 G1/A7 lose the stated finite local-test gate under their retained genuine inputs; CX\_GRH CX-ALG-024 and GCT\_PLETH T-052 keep their distinct global arithmetic gates.

\subsection{186. DO089: integral compression and the complete quadratic primary alternatives}\label{do089-integral-compression-and-the-complete-quadratic-primary-alternatives}

The entire actual all-top H133 nonunipotent fixed-W0 sector at support22 has absolute contribution\textless2\^{}-776095. The proof uses integral transverse compression of h+h\^{}-1-2I without global nilpotence, simple-pole characteristic coefficients, full generalized-primary decomposition, and the quadratic isotropic-image or empty exact-plane alternatives. Every global rank, incidence type, original column and source fibre is included. This sector removal alone leaves the complete support floor22 and all signed/global/arithmetic/DP44 gates open.

The operator H and its compression are endomorphism pairings of the original paired integral Hermitian modules; two distinct source/target lattices are retained. Generalized-primary subspaces and the fixed signature-(1,1) plane are used over E or at its real realizations, with no moving integral frame. The scalar recovered by the integer Bezout relation lies in the original fractional ideal. Exact contextual anchors remain SGA3 I4.7.3/4.7.3.1; SGA6 VI Theorem3.3; FGA232 Theorem3.1 pp148--149; FGA236 Corollaries1.7/2.3 pp229/232; SGA7 I VII2.7/2.9 pp169--174 and VIII7.4(b) pp310--311; cote4 scan003 §§8.1.1--8.1.2, scans005--006 §§8.4.1--8.4.3 and cote54 scans005--007. No new scheme identification, lambda-ring/K-theory identity, l-adic sheaf/weight theorem, CM-motive comparison or motivic L-function assertion follows. DP\_MULT Q3e and THETA3 G1/A7 gain the stated containing-set removal; CX\_GRH CX-ALG-024 and GCT\_PLETH T-052 retain their separate arithmetic gates.

\subsection{187. R067: complete original-module support22 partition}\label{r067-complete-original-module-support22-partition}

Every actual all-top support22 term is exhausted by both original F/IC8 classes at every rank1..33 and all identity-compression alternatives. DO089 supplies the whole nonunipotent fixed-W0 sector, DO088 the whole globally rank-one sector, R066 the whole IC8 rank1/2 sector, and every earlier complete range is retained. The actual complete all-top support floor is23. All68/31 chronological error entries are unchanged and both sums remain\textless2\^{}-230; no new charge is added. The uniform floor remains11. J23.2 and every global/arithmetic/source/DP44 gate remain open.

The assembly partitions rank conditions and fixed/isotropic submodules in the original integral Hermitian modules and labelled correspondences. Every lattice, polarization, finite level, coefficient module, Picard datum and biextension remains the original one. No rational-coordinate count gives an inverse on the arithmetic theta image.

The exact carried contextual anchors are SGA3 I4.7.3/4.7.3.1 and SGA6 VI Theorem3.3; FGA232 Theorem3.1 pp148--149; FGA236 Corollaries1.7/2.3 pp229/232; SGA7 I VII2.7/2.9 pp169--174 and VIII7.4(b) pp310--311. Formal crystals and Cartier duality retain cote4 scan003 §§8.1.1--8.1.2, scans005--006 §§8.4.1--8.4.3 and cote54 scans005--007. These locators provide context, not a substitute proof of the exhaustion. No new scheme identification, lambda-ring/K-theory equality, l-adic sheaf/weight result, CM-motive comparison or motivic L-function theorem follows. DP\_MULT Q3e and THETA3 G1/A7 gain exactly J23.1; CX\_GRH CX-ALG-024 and GCT\_PLETH T-052 retain their distinct gates.

Every quantitative input retains its own whole-proof checks and source restrictions. This new proof concerns complete partition coverage and exact error accounting. The certificate is a finite diagnostic; it neither enumerates all rational matrices nor evaluates a global period. No Lean build or archive audit is performed.

\subsection{188. DO090: the full real-field coefficient ideal and supports23 through26}\label{do090-the-full-real-field-coefficient-ideal-and-supports23-through26}

The entire actual all-top H133 nonunipotent fixed-W0 sector at supports23..26, with all global ranks and incidence profiles, has absolute contribution\textless2\^{}-65341. The original full-G coefficient ideal makes the exact quadratic plane empty through26, and all five complete DO089 bounds retain their original columns, paired modules, fibres and finite factors. The complete all-top floor remains23 and the uniform floor11; other structural classes and all signed/global/arithmetic/DP44 gates remain open.

The real-field trace is a section of the original fractional ideal U J\_G,F I\_III,F. Its F-ideal norm, Q\^{}167, is not doubled by the two CM components. The original integral Hermitian modules and their determinantal characteristic coefficients are retained. Contextual exact anchors are SGA3 I4.7.3/4.7.3.1; SGA6 VI Theorem3.3; FGA232 Theorem3.1 pp148--149; FGA236 Corollaries1.7/2.3 pp229/232; SGA7 I VII2.7/2.9 pp169--174 and VIII7.4(b) pp310--311; cote4 scan003 §§8.1.1--8.1.2, scans005--006 §§8.4.1--8.4.3 and cote54 scans005--007. These are the carried module, Picard and biextension frameworks, not newly constructed comparison maps. No new lambda-ring/K-theory identity, l-adic sheaf or weight theorem, CM motive or motivic L-function assertion follows. DP\_MULT Q3e and THETA3 G1/A7 gain the stated containing-set removal; CX\_GRH CX-ALG-024 and GCT\_PLETH T-052 retain their arithmetic gates.

\subsection{189. R068: original integral fixed-class counts at the actual cutoff}\label{r068-original-integral-fixed-class-counts-at-the-actual-cutoff}

At the unchanged actual all-top H133 cutoff, the entire original fixed class F is covered through support110 for rank(A-I)=2,154 for rank3,and all168 places for every rank4..33. Exactly1627 fresh rank/support cells beyond K069 receive one4085 charge. This uses the original integral block chart, r+6 original column coordinates, full paired target-overlattice and source fibre, and complete Green numerator; it assumes neither global unipotence nor the IC8 partial-monomial classification. The resulting remainder has F rank2 only at support\textgreater=111,F rank3 only at support\textgreater=155,and no F rank4..33. The complete all-top floor remains23 and uniform floor11; all signed/global/arithmetic/DP44 gates remain open.

The argument specializes a proved count on the original integral positive block and its fixed submodule. It keeps the original labelled correspondences, finite fields, Hilbert metric, Picard/Hom data and coefficient realization. Both stored positive-block entries and selected column coordinates belong to the original modules. Their injective rational encoding is not an integral arithmetic-image inverse.

Exact carried contextual anchors are SGA3 I4.7.3/4.7.3.1 and SGA6 VI Theorem3.3; FGA232 Theorem3.1 pp148--149; FGA236 Corollaries1.7/2.3 pp229/232; SGA7 I VII2.7/2.9 pp169--174 and VIII7.4(b) pp310--311; cote4 scan003 §§8.1.1--8.1.2, scans005--006 §§8.4.1--8.4.3 and cote54 scans005--007. No new scheme identification, lambda-ring/K-theory equality, l-adic sheaf/weight theorem, CM-motive comparison or motivic L-function statement follows. DP\_MULT Q3e and THETA3 G1/A7 gain the exact F-sector restrictions above. CX\_GRH CX-ALG-024 and GCT\_PLETH T-052 retain their separate arithmetic comparison gates.

The signed sufficient inequality, prescribed Bessel/Mram period, fully admissible mixed-CM source, selected-image content, actual Read/STiesR and032 value, and DP44 cost comparison remain open. The companion checks exact directed bounds, all old/new cells and complete error sums. The whole proof check separately verifies that K069's parameter-dependent original count has been specialized without changing its geometric or arithmetic hypotheses. No Lean build or archive audit is performed.

\subsection{190. LIFT074: positive source scalar and conditional prescribed detector}\label{lift074-positive-source-scalar-and-conditional-prescribed-detector}

Conditional on a nonzero global Bessel functional on the two specified embedded automorphic occurrences and the genuine source/local-vector premises, the prescribed initial Mram column and admitted bottom row have \textless F,Pu(T)\textgreater=c\_Per\emph{conjugate(b)!=0, where c\_Per=c1}mu\_f(K2)\textgreater0 is the actual source Haar-disintegration scalar. The finite tests survive and the complete real Wu contraction equals conjugate(b), even if W(1)=0. The normalized diagonal Gram is at least c\_Per\^{}2\emph{\textbar b\textbar\^{}2/(c0}\textbar\textbar F\textbar\textbar\^{}2)\textgreater0. Global distinction, actual arithmetic source/Read/STiesR/032 transport, the unevaluated coupling b and all DP44 costs remain open.

This theorem is an analytic contraction of the same original Hermitian/Picard/Hom data and its literal source correspondence; it does not identify that correspondence with an integral theta image or CM-motive morphism. The coefficient realization and selected metric are unchanged. Exact contextual anchors remain SGA3 I4.7.3/4.7.3.1; SGA6 VI Theorem3.3; FGA232 Theorem3.1 pp148--149; FGA236 Corollaries1.7/2.3 pp229/232; SGA7 I VII2.7/2.9 pp169--174 and VIII7.4(b) pp310--311; cote4 scan003 §§8.1.1--8.1.2, scans005--006 §§8.4.1--8.4.3 and cote54 scans005--007. No new scheme comparison, lambda-ring/K-theory equality, l-adic sheaf/weight theorem, motive or motivic L-function assertion follows. DP\_MULT Q3e and THETA3 G1/A7 gain this conditional detector transfer, while the literal arithmetic pin and value remain. CX\_GRH CX-ALG-024 and GCT\_PLETH T-052 retain their distinct arithmetic comparison gates.

\subsection{191. K093: rank-one torsion translations and residual ordinary derivatives}\label{k093-rank-one-torsion-translations-and-residual-ordinary-derivatives}

The actual scalar-decoration matrix has the exact finite Fourier/augmentation expansion RT13, represented by rational rank-one first-edge Weyl translations with the original central phase, full Y, class-dependent covectors and contracted full-middle mass. Every nonzero finite Fourier coefficient is an11-unit. The actual selected symbol along a rank-one increment has degree\textless=nu\_tau,1, so its ordinary translated evaluation retains at least D\_tau-nu\_tau,1\textgreater0 first-edge derivatives at every label. The pure-exponential term is zero; the complete differentiated germ and class sum remain unevaluated. No selected-value/content, global nonzero, arithmetic inverse, source family, Read/032 or DP44 conclusion is asserted.

The rational rank-one lift uses the canonical tensor-Hom isomorphism for the original finite projective modules and module CRT at paired CM primes. Its translated section lives in the actual rigidified torsion difference line in the Picard/Poincare framework; the complete Mumford lift fixes the central scalar. The integral selected Weyl symbol restricts to the determinantal rank-one locus with the stated degree bound, but evaluation still uses the original ordinary differential functional. Exact carried anchors are FGA232 Theorem3.1 pp148--149; FGA236 Corollaries1.7/2.3 pp229/232; SGA7 I VII2.7/2.9 pp169--174 and VIII7.4(b) pp310--311; SGA3 I4.7.3/4.7.3.1; SGA6 VI Theorem3.3; cote4 scan003 §§8.1.1--8.1.2, scans005--006 §§8.4.1--8.4.3 and cote54 scans005--007. No new lambda-ring/K-theory identity, l-adic purity/weight theorem, CM-motive comparison or motivic L-function claim follows. DP\_MULT Q3e and THETA3 G1/A7 retain the exact selected derivative/class gate; CX\_GRH CX-ALG-024 and GCT\_PLETH T-052 retain their separate arithmetic gates.

\subsection{192. DO091: integral transverse pairings and a real-ideal norm count}\label{do091-integral-transverse-pairings-and-a-real-ideal-norm-count}

The entire actual all-top H133 globally rank-one-unipotent fixed-W0 class at support23 has absolute contribution\textless2\^{}-1544. The original transverse pairings give a complete witness-or-norm-equation split. The norm sector retains the actual affine difference ideal, epsilon1/12 E/F divisor counts, real pure-imaginary ideal covolume, all24 marked sizes, original type/pair weights and selected near-point metric. Global unipotence is a hypothesis. The complete all-top floor remains23 and the uniform floor11; the other structural classes and all signed/global/arithmetic/DP44 gates remain open.

The new counts use the original Hermitian modules, actual integral pairing sections and relative norm ideals; the pure-imaginary line is the real ideal i Lambda\_F with its real different and covolume. No moving integral basis, changed source normalization or arithmetic theta-image inverse is introduced. Carried exact contextual anchors are SGA3 I4.7.3/4.7.3.1; SGA6 VI Theorem3.3; FGA232 Theorem3.1 pp148--149; FGA236 Corollaries1.7/2.3 pp229/232; SGA7 I VII2.7/2.9 pp169--174 and VIII7.4(b) pp310--311; cote4 scan003 §§8.1.1--8.1.2, scans005--006 §§8.4.1--8.4.3 and cote54 scans005--007. No new lambda-ring/K-theory identity, l-adic sheaf or weight theorem, CM motive or motivic L-function assertion follows. DP\_MULT Q3e and THETA3 G1/A7 gain the stated containing-set removal; CX\_GRH CX-ALG-024 and GCT\_PLETH T-052 retain their arithmetic gates.

\subsection{193. R069: symmetric compression and paired cross-column ideals}\label{r069-symmetric-compression-and-paired-cross-column-ideals}

The entire original IC8 positive-compression rank1/2 class at supports23..28 of the unchanged actual all-top H133 kernel has absolute contribution\textless2\^{}-23546. The symmetric partial-monomial compression preserves rank at every G factor. Finite rank2 then forces every untouched original cross column to vanish at both CM partners, supplying a paired ideal in the unchanged joint160-coordinate ball count. All original columns, fibres, target and pair weights are retained. No global unipotence is assumed. The complete all-top floor remains23 and uniform floor11; identity-compression and higher-support classes, signed positivity, actual source/arithmetic matching and DP44 remain open.

The new input is the symmetric compression on the same original Hermitian module and the resulting paired residue-field kernel statement. Paths/cycles prove a determinantal rank identity, and the count stays in the original product of fractional ideals. No integral moving basis, new CM field, source normalization, labelled correspondence, or selected projector is introduced. Exact carried contextual anchors are SGA3 I4.7.3/4.7.3.1; SGA6 VI Theorem3.3; FGA232 Theorem3.1 pp148--149; FGA236 Corollaries1.7/2.3 pp229/232; SGA7 I VII2.7/2.9 pp169--174 and VIII7.4(b) pp310--311; cote4 scan003 §§8.1.1--8.1.2, scans005--006 §§8.4.1--8.4.3 and cote54 scans005--007. No new lambda-ring/K-theory identity, l-adic sheaf/weight theorem, CM motive comparison or motivic L-function assertion follows.

DP\_MULT Q3e and THETA3 G1/A7 gain precisely this smaller containing remainder for the unchanged kernel. The signed sufficient inequality, actual fully admissible mixed-CM source, prescribed global Bessel/Mram value, selected derivative/class content, actual Read/STiesR/specific032 transport and DP44 cost difference remain open. CX\_GRH CX-ALG-024 and GCT\_PLETH T-052 retain their separate arithmetic requirements. The source Green and lattice transference inputs are carried through the already checked R063/R066/DO064/DO085 interfaces. New finite diagnostics do not enumerate the actual rational kernel or substitute for this whole proof.

\subsection{194. DO092: conjugate content and the selected radial integral}\label{do092-conjugate-content-and-the-selected-radial-integral}

The entire actual all-top H133 globally rank-one-unipotent fixed-W0 class at support24 has absolute contribution\textless2\^{}-516. The exact conjugate-content constraint and split-only relative norm count give constants\textless2\^{}361 and\textless2\^{}244. Integrating the retained actual selected radial density gives eta\_a=4\^{}(a-1)/{[}2(a-1)binom(2a-2,a-1){]}, with a full630-bit improvement at a=f168. The weighted G patterns, real eta metric, original affine ideal and all25 marked/type/pair cells are retained. Global unipotence and rank one are premises. Complete floors remain23/11; signed/global/arithmetic/source/DP44 gates remain open.

This refines the count of original content ideals C under C bar(C) subset A and of actual relative norm ideals, without assuming principality. The complex realization retains its Hermitian bundle metric and selected trace. The original modules, finite correspondences, Picard classes, polarizations and biextensions are unchanged. Exact carried contextual anchors are SGA3 I4.7.3/4.7.3.1; SGA6 VI Theorem3.3; FGA232 Theorem3.1 pp148--149; FGA236 Corollaries1.7/2.3 pp229/232; SGA7 I VII2.7/2.9 pp169--174 and VIII7.4(b) pp310--311; cote4 scan003 §§8.1.1--8.1.2, scans005--006 §§8.4.1--8.4.3 and cote54 scans005--007. No new lambda-ring/K-theory identity, l-adic sheaf/weight theorem, CM motive or motivic L-function assertion follows. DP\_MULT Q3e and THETA3 G1/A7 gain precisely this containing-set removal; CX\_GRH CX-ALG-024 and GCT\_PLETH T-052 retain their separate arithmetic gates.

\subsection{195. AM-LIFT-075: Every positive-gap direct theta output is excluded by the actual good projector}\label{am-lift-075-every-positive-gap-direct-theta-output-is-excluded-by-the-actual-good-projector}

For a split type-II theta quotient from GL\_a to GL\_r with d=r−a\textgreater0, every nonzero principal-level finite constituent has one linear Green factor of multiplicity at least d.~For d\textgreater=2 squarefreeness kills it at every selected G place; for d=1 the actual goodTuple degree-one witness kills each tensor block. Thus the genuine full-rank e0 kills the actual compact-source direct theta cuspidal projection for every permitted even m\textgreater=2, where d=m−1. This excludes that direct construction only. It proves no vanishing of the prescribed Bessel period, coupled theta map or transported-leading-r intermediate projector. Fixed-real-type, fixed-finite-level global distinction remains open.

This is a local representation-theoretic obstruction in the existing genuine finite projector and a compact-source cuspidal theta pairing, not a new morphism of the CM Jacobian or its Picard scheme. Finite parabolic induction respects the separate Green orbit grading; this grading is not identified with a motivic weight or a new lambda-ring identity. Exact contextual anchors are SGA3 I4.7.3/4.7.3.1; SGA6 VI Theorem3.3; FGA232 Theorem3.1 pp148--149; FGA236 Corollaries1.7/2.3 pp229/232; SGA7 I VII2.7/2.9 pp169--174 and VIII7.4(b) pp310--311; cote4 scan003 §§8.1.1--8.1.2, scans005--006 §§8.4.1--8.4.3 and cote54 scans005--007. No new scheme comparison, K-theory equality, l-adic sheaf/weight theorem, CM-motive or motivic L-function assertion follows. DP\_MULT Q3e and THETA3 G1/A7 lose the specified direct compact-to-bottom candidate, while their coupled arithmetic family, normalization and value remain open. CX\_GRH CX-ALG-024 and GCT\_PLETH T-052 retain their separate gates.

\subsection{196. AM-K094: The actual full covector fibre gives an exact integral corner-column matrix at11}\label{am-k094-the-actual-full-covector-fibre-gives-an-exact-integral-corner-column-matrix-at11}

For each actual coarse initial class retaining all non-T data, G-neatness makes the full strict covector fibre the product of all nonzero covectors. Its varying global coefficient lattices have one common localization at11. Punctured Fourier transform is an integral isomorphism of the line-augmentation modules there, with inverse T\^{}(−a)F\_minus and no division by T−1. The actual transformed Mram basis has exact class norm T\^{}(af/2)2\^{}(f/2)\textbar\textbar P\textbar\textbar. After the common full-middle contracted cut, its image is exactly the signed2\^{}f corner column FC13, with original masses, central phases, coordinate pullbacks, one shared coefficient P and all residual ordinary derivatives retained. FC14 identifies the actual localized image/content ideal with the corner-column span, even when that image is zero. Global clearing sizes, corner/middle-class noncancellation, source/Read/032 matching and every DP44 cost remain open.

The argument is an integral change of basis at11 on the actual finite locally free coefficient lattices over the fixed covector fibre. The finite Fourier transform uses the genuine torsion character pairing, and its rational rank-one representatives retain the rigidified torsion difference line and Poincare biextension central scalar. It leaves the actual ordinary differential functional unchanged. Exact contextual anchors are SGA3 I4.7.3/4.7.3.1; SGA6 VI Theorem3.3; FGA232 Theorem3.1 pp148--149; FGA236 Corollaries1.7/2.3 pp229/232; SGA7 I VII2.7/2.9 pp169--174 and VIII7.4(b) pp310--311; cote4 scan003 §§8.1.1--8.1.2, scans005--006 §§8.4.1--8.4.3 and cote54 scans005--007. No new lambda-ring/K-theory identity, l-adic purity or weight theorem, CM-motive comparison or motivic L-function statement is inferred. DP\_MULT Q3e and THETA3 G1/A7 gain the exact corner-image/content criterion, while their evaluated derivative/class gate remains. CX\_GRH CX-ALG-024 and GCT\_PLETH T-052 retain their separate arithmetic gates.

\subsection{197. AM-LIFT-076: Every existing genuine full-e0 bottom constituent has simple tempered base change}\label{am-lift-076-every-existing-genuine-full-e0-bottom-constituent-has-simple-tempered-base-change}

Every existing cuspidal bottom constituent selected by the genuine full e0 in the stated actual setting has simple generic global Arthur parameter and conjugate-selfdual cuspidal, cohomological base change, tempered at every finite place. Squarefree finite Green types pin the complete distinct inertial multiset; actual unitarity forces all real exponents to zero. The actual ramified auxiliary-theta comparison excludes Arthur SL2 lengths greater than one, and the literal all-degree divisor witnesses exclude multiple global blocks. Original odd-rank parity permits unit Harish--Chandra gaps. Generic packet membership gives tempered local packets at every place and spectral multiplicity one. No one-dimensional finite-level or coefficient-space assertion, source-family existence, globally distinguished pair, arithmetic image/value, Read/032 or DP44 conclusion follows.

This classifies actual automorphic constituents of the existing Hermitian unitary group through its finite Green types and global Arthur/base-change parameters. It does not identify those automorphic representations with a motive of the particular CM Jacobian or construct the requested theta-cycle map. The coefficient parity is a statement about the original algebraic highest weight; the temperedness theorem is not a new l-adic sheaf-purity assertion. Exact contextual anchors are SGA3 I4.7.3/4.7.3.1; SGA6 VI Theorem3.3; FGA232 Theorem3.1 pp148--149; FGA236 Corollaries1.7/2.3 pp229/232; SGA7 I VII2.7/2.9 pp169--174 and VIII7.4(b) pp310--311; cote4 scan003 §§8.1.1--8.1.2, scans005--006 §§8.4.1--8.4.3 and cote54 scans005--007. No new scheme comparison, lambda-ring or K-theory identity, CM-motive comparison or motivic L-function theorem is inferred. DP\_MULT Q3e and THETA3 G1/A7 gain the larger-representation local temperedness premise for every existing actual bottom constituent, while their global distinction, integral image and value gates remain. CX\_GRH CX-ALG-024 and GCT\_PLETH T-052 retain their separate arithmetic comparisons.

\subsection{198. R070: integral short compression and a rank-three finite baseline}\label{r070-integral-short-compression-and-a-rank-three-finite-baseline}

At actual all-top H133 supports23/24, the entire C-transverse scalar-cross CaseA and the plane-C subcase with at least one nonzero original2x2 short compression have combined absolute contribution\textless2\^{}-31703. The four integral pairings have conjugates\textless2\^{}47. The complete paired finite lemma restores a rank-three Green baseline through one whole-C ideal and its uniformly norm-bounded compression divisor, with no extra witness or compression-pattern factor. All original type/pair, selected-trace, target-column and source-fibre factors remain. No global unipotence is assumed. The all-four-zero plane compression subcase remains; complete floors23/11 and signed/global/arithmetic/DP44 gates remain open.

The new witness is a matrix of integral pairings between the original source and target Hermitian modules over Spec(O\_E). Its simultaneous vanishing at assigned CM primes is proved by paired nilpotent linear algebra, not by choosing a new integral frame or fixing the individual short vectors. The norm-divisor estimate and original fractional-ideal packing use the same complex Hermitian realization and actual selected trace. Exact contextual anchors remain SGA3 I4.7.3/4.7.3.1; SGA6 VI Theorem3.3; FGA232 Theorem3.1 pp148--149; FGA236 Corollaries1.7/2.3 pp229/232; SGA7 I VII2.7/2.9 pp169--174 and VIII7.4(b) pp310--311; cote4 scan003 §§8.1.1--8.1.2, scans005--006 §§8.4.1--8.4.3 and cote54 scans005--007. No new scheme comparison, lambda-ring/K-theory equality, l-adic sheaf/weight theorem, CM-motive identification or motivic L-function assertion is inferred. DP\_MULT Q3e and THETA3 G1/A7 gain this strict containing-set removal; CX\_GRH CX-ALG-024 and GCT\_PLETH T-052 keep their separate arithmetic comparisons.

\subsection{199. DO093: exact inverse correspondence and positive165-dimensional kernel}\label{do093-exact-inverse-correspondence-and-positive165-dimensional-kernel}

In the actual K068 finite-unipotent remainder at all-top H133 supports23/24, the plane-C but inverse-C-transverse scalar-cross terms have total absolute contribution\textless2\^{}-885948. Literal inversion conjugates the actual normalized summand, preserves the signed pencil coefficient, all source/target weights, central/class/Haar normalizations and the cutoff, and injects this set into the complete original transverse CaseA count. Together with R070 the remaining scalar-cross terms have both cross blocks in the packet plane and all four short compressions zero. They fix an E-defined positive165-dimensional subspace and have global rank(h−I)\textless=4; individual test vectors need not be fixed and no global unipotence is assumed. Complete floors23/11, signed/global/arithmetic and DP44 gates remain open.

The inversion is the original Hermitian adjoint inverse on the same source/target modules over Spec(O\_E), and the exact T correspondence is transposed together with its symmetric projector coefficients. The residual fixed space is the kernel of an E-linear rank-one map on the original positive166-space; no integral moving basis or controlled-denominator isometry is constructed. Exact contextual anchors remain SGA3 I4.7.3/4.7.3.1; SGA6 VI Theorem3.3; FGA232 Theorem3.1 pp148--149; FGA236 Corollaries1.7/2.3 pp229/232; SGA7 I VII2.7/2.9 pp169--174 and VIII7.4(b) pp310--311; cote4 scan003 §§8.1.1--8.1.2, scans005--006 §§8.4.1--8.4.3 and cote54 scans005--007. No new scheme comparison, lambda-ring/K-theory identity, l-adic sheaf/weight theorem, CM-motive comparison or motivic L-function theorem follows. DP\_MULT Q3e and THETA3 G1/A7 gain this containing-set removal and explicit positive fixed-space reduction; CX\_GRH CX-ALG-024 and GCT\_PLETH T-052 retain their separate arithmetic gates.

\subsection{200. AM-DO-094: Paired finite ranks and the last transverse scalar}\label{am-do-094-paired-finite-ranks-and-the-last-transverse-scalar-1}

Under the full original/inverse doubly-plane block conditions, finite nilpotent rank at most3 and both nonzero paired cross blocks force the transverse scalar alpha=1. In the actual global residual, retained norm-one and polynomial wedge identities also give alpha=1 at rank\textless=2 with at least one nonzero cross block, and rank\textless=1 without a cross condition. The explicit original G-unit basis makes delta\_z=(alpha−1)U/M integral, with all conjugates\textless2\^{}71. Exact counterexamples exclude the weaker block hypothesis and the two unclaimed rank/cross cells. This supplies no complete restoration/count, charge, global nilpotence or support-floor change.

The statement concerns paired determinantal conditions on the original Hermitian modules over Spec(O\_E). Its rational orthogonal basis is proved integral and invertible only at the actual G primes; the scalar pairing uses separate source and target selfdual lattices. Exact carried contextual anchors are SGA3 I4.7.3/4.7.3.1; SGA6 VI Theorem3.3; FGA232 Theorem3.1 pp148--149; FGA236 Corollaries1.7/2.3 pp229/232; SGA7 I VII2.7/2.9 pp169--174 and VIII7.4(b) pp310--311; cote4 scan003 §§8.1.1--8.1.2, scans005--006 §§8.4.1--8.4.3 and cote54 scans005--007. No new scheme comparison, lambda-ring/K-theory identity, l-adic sheaf/weight theorem, CM-motive or motivic L-function assertion follows. DP\_MULT Q3e and THETA3 G1/A7 gain the displayed local conditions, while CX\_GRH CX-ALG-024 and GCT\_PLETH T-052 retain their arithmetic gates.

\subsection{201. AM-K095: The actual complete middle fibre becomes one composite incidence}\label{am-k095-the-actual-complete-middle-fibre-becomes-one-composite-incidence}

The full middle T-covector fibre, with its genuine fine stabilizer masses, sums exactly to the composite incidence u2 Z X=alpha. The weighted coarse factor1/\textbar Gamma\_q\textbar{} is represented by a full coarse-stabilizer G cut only after invariant middle coevaluation. The finite norm is divided by \textbar Gamma\_q\textbar, and the auxiliary family is split at11 before the prescribed ordinary derivative. The exact K094 Fourier columns become the same-input coarse-class quadratic-shear formula MF12, with polar rank2b and explicit central phase. All original coefficient clearings, norms, phases, derivatives, cover pullbacks and terminal operators remain; selected content and coarse-class noncancellation are open.

This is a finite-fibre orbit--stabilizer identity and integral splitting of the auxiliary theta-section module at11 on the existing Picard/Hom and theta-group data. The finite quadratic shear retains its Heisenberg central phase; it is not identified with an algebraic additive translation family or an independent jet source. Exact carried contextual anchors are SGA3 I4.7.3/4.7.3.1; SGA6 VI Theorem3.3; FGA232 Theorem3.1 pp148--149; FGA236 Corollaries1.7/2.3 pp229/232; SGA7 I VII2.7/2.9 pp169--174 and VIII7.4(b) pp310--311; cote4 scan003 §§8.1.1--8.1.2, scans005--006 §§8.4.1--8.4.3 and cote54 scans005--007. No new lambda-ring/K-theory equality, l-adic weight theorem, CM-motive comparison or motivic L-function theorem follows. DP\_MULT Q3e and THETA3 G1/A7 retain the prescribed derivative/coarse-class and arithmetic comparison gates; CX\_GRH CX-ALG-024 and GCT\_PLETH T-052 keep their separate requirements.

\subsection{202. AM-R-071: Integral norm rigidity eliminates the last compact transverse scalar}\label{am-r-071-integral-norm-rigidity-eliminates-the-last-compact-transverse-scalar}

In the actual DO093 support23/24 doubly-plane/four-zero scalar-cross remainder, the original integral delta\_z=(U/M)(alpha−1) and alpha\emph{bar(alpha)=1 imply delta\_z=0 and alpha=1 globally. Here gcd(U,M)=9 and36U/M²\textless1 exactly. Integer Bezout makes delta\_z}bar(delta\_z)/(U/9) an algebraic integer in F; every real conjugate lies in{[}0,36U/M²{]}, so its positive integral norm forces zero. The actual h fixes a positive E-space of dimension166 and has rank(h−I)\textless=3, with remaining signature(2,1). No finite-rank hypothesis, integral moving frame, global nilpotence or new error charge is used. Complete support floors and all signed/global/arithmetic/DP44 gates remain open.

The new step is an integral trace-and-norm identity over Spec(O\_F) for the original scalar pairing section over Spec(O\_E). Its field norm is the determinant of multiplication, and all real realizations give a strict bound below1. The fixed space is the kernel of the actual rank-one map on the original positive167-space. Exact carried contextual anchors are SGA3 I4.7.3/4.7.3.1; SGA6 VI Theorem3.3; FGA232 Theorem3.1 pp148--149; FGA236 Corollaries1.7/2.3 pp229/232; SGA7 I VII2.7/2.9 pp169--174 and VIII7.4(b) pp310--311; cote4 scan003 §§8.1.1--8.1.2, scans005--006 §§8.4.1--8.4.3 and cote54 scans005--007. These are contextual frameworks; no new lambda-ring/K-theory equality, l-adic weight theorem, CM-motive comparison or motivic L-function assertion is invoked. DP\_MULT Q3e and THETA3 G1/A7 gain the strict smaller structural remainder. CX\_GRH CX-ALG-024 and GCT\_PLETH T-052 retain their arithmetic gates.

\subsection{203. AM-LIFT-077: The prescribed global coupling has an exact central-L-value calibration}\label{am-lift-077-the-prescribed-global-coupling-has-an-exact-central-l-value-calibration}

For an existing genuine initial/bottom embedded cuspidal pair in the LIFT070/LIFT076 scopes and the prescribed finite vectors, GC17 gives \textbar b\textbar²/(nT\emph{nF)={[}aS²/(aH\emph{aG}a0\emph{d\_sigma}2\^{}(t\_pi+1)){]}}Lratio\^{}(G union infinity)(1/2)\emph{product\_G d\_varsigma\^{}(-1)}product\_T q/{[}(q−1)(q\^{}k−1){]}. The formula includes b=0, retains the conjugated bottom slot and physical character, and uses exactly one formal degree, whole-reference Haar ratio and component factor. GC19 gives the corresponding source Gram lower bound after nF cancels. Central nonvanishing, its size, source existence, actual arithmetic Read/032 matching and DP44 costs remain open; no finite-kernel charge follows.

This is an analytic period identity for the existing unitary groups and physical Bessel subgroup, with all global/reference Haar and representation embedding norms distinguished. It identifies a central automorphic L-ratio; it does not identify that ratio with the motive of the specified CM Jacobian or construct the required theta-cycle comparison. Exact carried contextual anchors are SGA3 I4.7.3/4.7.3.1; SGA6 VI Theorem3.3; FGA232 Theorem3.1 pp148--149; FGA236 Corollaries1.7/2.3 pp229/232; SGA7 I VII2.7/2.9 pp169--174 and VIII7.4(b) pp310--311; cote4 scan003 §§8.1.1--8.1.2, scans005--006 §§8.4.1--8.4.3 and cote54 scans005--007. No new scheme isomorphism, lambda-ring/K-theory identity, l-adic weight or CM-motive theorem is claimed. DP\_MULT Q3e and THETA3 G1/A7 gain the exact conditional coupling normalization, while their central nonzero/arithmetic value gates remain. CX\_GRH CX-ALG-024 and GCT\_PLETH T-052 remain separate.

\subsection{204. R072: projected packet-plane lattice and joint cross-row count}\label{r072-projected-packet-plane-lattice-and-joint-cross-row-count}

At the unchanged actual all-top H133 endpoint, R071 scalar-cross support23 terms with at least six nonzero original S scalar coordinates have total absolute contribution\textless2\^{}-3600. The full positive167 compression supplies three original integral transverse coefficients, while the exact negative-row bound is H/(2sqrt U). Spherical-cap/Gaussian packing counts the first row jointly. The fixed packet-plane projection retains its U denominator, exact paired target indexT\^{}(4s), and positive-rank-two source fibre. All159 k-cells6..164, all source/target types, actual selected trace, class weights and finite projectors are retained. The remaining R071 support23 cross row has k\textless=5. No support24 removal, global-unipotence inference or complete support-floor increase follows; signed/global/source/arithmetic/DP44 gates remain open.

The fixed rational projection of the original Hermitian module has image (U\^{}-1 O\_E)u plus O\_E e\_n.~Its denominator and determinant index are retained, and the G/T frame is proved invertible only at the required primes. The cross ideal is a paired determinantal condition over Spec(O\_E), and the joint count uses the unchanged complex half-trace metric. Exact contextual anchors remain SGA3 I4.7.3/4.7.3.1; SGA6 VI Theorem3.3; FGA232 Theorem3.1 pp148--149; FGA236 Corollaries1.7/2.3 pp229/232; SGA7 I VII2.7/2.9 pp169--174 and VIII7.4(b) pp310--311; cote4 scan003 §§8.1.1--8.1.2, scans005--006 §§8.4.1--8.4.3 and cote54 scans005--007. No new lambda-ring/K-theory equality, l-adic weight theorem or motivic L-function assertion is added. DP\_MULT Q3e and THETA3 G1/A7 gain the smaller finite remainder; CX\_GRH CX-ALG-024 and GCT\_PLETH T-052 retain their arithmetic gates.

\subsection{205. AM-DO-095: the actual three-dimensional residual}\label{am-do-095-the-actual-three-dimensional-residual}

In the actual R071 signature(2,1) residual, BC is skew-adjoint rank one exactly at global rank3; BC=0 is exactly global rank2 with real reciprocal eigenvalue lambda. Every original S entry of BC is an exact G-unit three-minor, has conjugates\textless2\^{}9 and only simple paired typeIII poles. The pointwise Green envelope is2 Q\^{}(-501\emph{167)2\^{}(9}167f)T\^{}(167t), with no C-pattern/witness factor. In rank2, doubled G divisibility of (lambda-1)\^{}2/lambda forces global unipotence for t\textless=22. The existing DO082 removal therefore leaves rank-two support23 only nonunipotent with all23 active places typeIII; at24 its nonunipotent branch needs t\textgreater=23. No new charge or complete-floor change. Actual central value, arithmetic/source and DP44 gates remain open.

Original Hermitian Hom/determinant modules and paired finite valuations over Spec(O\_E), with their fixed complex metric. Exact carried SGA/FGA/cote locators and the limits on motive/weight claims are in the complete canonical proof; no new comparison theorem is asserted.

\subsection{206. AM-R-073: the actual three-dimensional residual}\label{am-r-073-the-actual-three-dimensional-residual}

The actual global rank3 R071 terms at s23,k\textless=63 and s24,k\textless=51 have combined absolute contribution\textless2\^{}-21500. The bare support-ideal count retains all original sparse cross coordinates, three transverse values, pivot, fixed packet-plane U denominator, paired T\^{}4 column index, positive rank-two source fibre and complete pair weights. The actual DO095 witness gives the uniform no-pattern rank-three Green bound. All114 cells are checked. Combined with R072 this removes every rank3 support23 term; DO095 then leaves the support23 scalar-cross remainder rank2 nonunipotent, k\textless=5, t23. At support24 remaining rank3 has k\textgreater=52. Complete floors23/11 and signed/global/source/arithmetic/DP44 gates remain open.

Original Hermitian Hom/determinant modules and paired finite valuations over Spec(O\_E), with their fixed complex metric. Exact carried SGA/FGA/cote locators and the limits on motive/weight claims are in the complete canonical proof; no new comparison theorem is asserted.

\subsection{207. K096: reciprocal poles and original scalar ideals}\label{k096-reciprocal-poles-and-original-scalar-ideals}

The actual R071 support23 residual has no globally nonunipotent global-rank-two term. The nonzero reciprocal trace has exactly23 actual simple F-poles; principal-norm parity and inversion give twelve positive and eleven negative eigenvalue valuations. Both actual arrows then make every original cross scalar integral at the twelve positive places and at most simple-polar at the eleven negative places, with full paired ideal normT\^{}-22 and no U denominator. A nonzero scalar contradicts H\^{}336 T\^{}22\textless(4U)\^{}168, with exact margin374045/512. No sparsity premise or numerical charge is used. DO082 retains the complementary unipotent exclusion. All structural classes and the complete support-floor update require a separate assembly; support24 and signed/global/source/arithmetic/DP44 gates remain open.

This is a norm and valuation argument for the original scalar fractional ideals over Spec(O\_E) and Spec(O\_F). Both source/target module inclusions and both CM partners are retained; a primitive integer extraction avoids a new determinant denominator. The fixed CM Jacobian, Picard classes, polarization, theta line and biextensions are unchanged. Exact carried contextual locators are SGA3 I4.7.3/4.7.3.1; SGA6 VI Theorem3.3; FGA232 Theorem3.1 pp148--149; FGA236 Corollaries1.7/2.3 pp229/232; SGA7 I VII2.7/2.9 pp169--174 and VIII7.4(b) pp310--311; cote4 scan003 §§8.1.1--8.1.2, scans005--006 §§8.4.1--8.4.3 and cote54 scans005--007. These are contextual placements, not newly invoked comparison theorems. No new lambda-ring/K-theory equality, l-adic weight theorem or motivic L-function assertion follows. DP\_MULT Q3e and THETA3 G1/A7 gain the exact empty scalar-cross subcase; CX\_GRH CX-ALG-024 and GCT\_PLETH T-052 keep their distinct arithmetic gates.

\subsection{208. AM-R-074: complete original support23 exhaustion}\label{am-r-074-complete-original-support23-exhaustion}

Every actual all-top support23 term is exhausted by the complete original F/IC8 partition at all compressed ranks1..33 and every fixed-S/scalar-cross alternative. R072/R073 cover every global-rank3 scalar cross, DO095 verifies all precise DO082 unipotent premises, and K096 removes the entire nonunipotent rank-two residual without sparsity. The complete all-top support floor is24, with uniform floor11. All77/31 chronological errors remain unchanged and both full sums are below2\^{}-230. No new charge is added. The signed inequality J24.2, actual source, selected arithmetic content, global Bessel/Mram value, specific032 matching/value and DP44 remain open.

This is an exhaustion of determinantal and fixed/isotropic-submodule strata of the same original integral Hermitian correspondences over Spec(O\_E), together with their specified complex realizations. The fractional ideals, both CM partners, coefficient modules, polarization, theta line, CM Jacobian, Picard data and biextensions are unchanged. Rational reconstruction in the cited counts supplies no inverse on the arithmetic theta image.

The exact carried module and determinantal contextual anchors are SGA3, Expose I,4.7.3/4.7.3.1, and SGA6, Expose VI, Theorem3.3. The Picard anchors are FGA, Bourbaki232, Theorem3.1, pp148--149, and Bourbaki236, Corollaries1.7/2.3, pp229/232. Biextensions retain SGA7 I, Expose VII,2.7/2.9, pp169--174, and Expose VIII,7.4(b), pp310--311. Formal groups, Dieudonne crystals and Cartier duality retain cote4 scan003, §§8.1.1--8.1.2, scans005--006, §§8.4.1--8.4.3, and cote54 scans005--007. These are contextual placements, not new comparison theorems.

No new scheme identification, lambda-ring/K-theory equality, l-adic sheaf/weight statement, CM-motive comparison or motivic L-function theorem follows. DP\_MULT Q3e and THETA3 G1/A7 gain exactly the smaller finite/real gate J24.1--2. CX\_GRH CX-ALG-024 and GCT\_PLETH T-052 retain their distinct arithmetic comparisons. Every cited quantitative theorem keeps its own exact source locators and whole-proof checks. This new proof checks complete partition coverage and chronological accounting; no Lean build or archive audit is performed.

\subsection{209. AM-DO-096: paired reciprocal-pole cap at support24}\label{am-do-096-paired-reciprocal-pole-cap-at-support24}

The entire globally nonunipotent rank-two R071 support24 sector has absolute contribution\textless2\^{}-2962487. The nonzero reciprocal trace has23 or24 actual poles; principal-norm parity and one exactly normalized inversion leave at most twelve paired cross-pole places, including the possible active eigenvalue-unit place. Primitive original extraction gives normT\^{}-24 without a U denominator. The full original R059 count retains both five-coordinate columns, source fibre, all G patterns, padded subset choices and pair weights; the complete exponent is -60671740183/20480. Exactly one2962487 charge is added, leaving78/31 entries with both full sums below2\^{}-230. Complete floors remain24/11. Globally unipotent rank2, rank3, signed/global/source/arithmetic and DP44 gates remain open.

prescribed global Bessel/Mram period, selected-image arithmetic content, literal admissible source constructor, concrete cycle/value and DP44 sign remain separate. The algebraic-geometric setting is the original Hermitian modules and scalar fractional ideals over Spec(O\_E), with their exact CM pairing and unchanged half-trace metric. Carried contextual anchors remain SGA3 I4.7.3/4.7.3.1; SGA6 VI Theorem3.3; FGA232 Theorem3.1 pp148--149; FGA236 Corollaries1.7/2.3 pp229/232; SGA7 I VII2.7/2.9 pp169--174 and VIII7.4(b) pp310--311; cote4 scan003 §§8.1.1--8.1.2, scans005--006 §§8.4.1--8.4.3 and cote54 scans005--007. No new lambda-ring, K-theory, l-adic or motivic conclusion is imported.

Direct inputs at their established scope are R071(GR6--12), DO095(RA9--10,RA21--23), K096(EP2--6,EP8--13), DO093(IV4--15), LIFT44(PR4--7), K070(PT5--8), and the complete original R059(GB7--21). The new result is the even/odd support24 pole classification, its twelve-place cross ideal, and the exact full-count substitution. The companion diagnostics supplement that general proof and do not claim to enumerate actual rational unitary arrows.

\subsection{210. R075: symmetric compression and paired minimal-rank columns}\label{r075-symmetric-compression-and-paired-minimal-rank-columns}

In the original IC8 class at compressed ranks3..31, the full symmetric compression has rank r at every actual G factor. Finite unipotence therefore forces full rank\textgreater=r+1; at equality the unique nontrivial Jordan block fixes every untouched original column at both CM partners. The resulting paired ideal and baseline rank r+2 retain every original DO085 coordinate, volume, target, source fibre and pair factor. All145 stated fresh cells contribute\textless2\^{}-8866, with exact maximum -36351867/4096. The complete rank3 ceiling becomes28, and with the earlier F/IC8 cuts every remaining term at supports24..28 has A=I. One8866 charge gives79/31 entries with both full sums\textless2\^{}-230. Complete floors24/11 remain; signed/global/source/selected-image/032/DP44 gates are open.

the determinantal identity NR3 and the single-Jordan-block argument NR5--7 in reductions of the original Hermitian modules over Spec(O\_E). The original scalar ideals, Picard data, CM polarization, theta line, biextensions and coefficient metrics are unchanged. Exact contextual anchors are SGA3 I4.7.3/4.7.3.1; SGA6 VI Theorem3.3; FGA232 Theorem3.1 pp148--149; FGA236 Corollaries1.7/2.3 pp229/232; SGA7 I VII2.7/2.9 pp169--174 and VIII7.4(b) pp310--311; cote4 scan003 §§8.1.1--8.1.2, scans005--006 §§8.4.1--8.4.3 and cote54 scans005--007. No new lambda-ring/K-theory equality, l-adic sheaf/weight theorem, CM-motive comparison or motivic L-function conclusion is imported. DP\_MULT Q3e and THETA3 G1/A7 gain precisely the stated smaller remainder; CX\_GRH CX-ALG-024 and GCT\_PLETH T-052 keep their separate arithmetic comparisons.

The external Green and transference inputs are used only through the already whole-checked DO085 interfaces: Gorodetsky and Hazan (2017), Theorem2.1/equation(2.1), through K067; Miller and Stephens-Davidowitz (2019), Theorem3.2/equation(3.3), through DO070. The companion is an exact endpoint/accounting diagnostic. The whole proof separately checks the new rank and paired-coordinate arguments. No rational-matrix enumeration, Lean build or archive audit is claimed.

\subsection{211. DO097: original positive-input integrality and unipotent norm}\label{do097-original-positive-input-integrality-and-unipotent-norm}

The entire globally unipotent rank-two R071 support24 sector with at least one type-I or type-III place has absolute contribution\textless2\^{}-71310. The exact block normal form gives N²Z=0 and alpha=Uq\textless=H+H\^{}-1-2, so all original S and fixed transverse inputs have only type-II poles and the full covector radius is sqrt(R/U). At most24 padded23-place ideals cover every non-all-II profile, with all164 coordinate sectors, common G patterns, genuine selected trace, true projected-column U denominator, paired T\^{}4 index and source fibre retained. The complete maximum is -292137655/4096; one71310 charge gives80/31 entries, both full sums\textless2\^{}-230. Complete floors24/11 remain. The all-II unipotent case, rank3 and fixed-block residuals, signed/global/source/arithmetic/032/DP44 gates remain open.

The module and fractional-ideal setting remains over the unchanged Spec(O\_E), with its fixed CM pairing and complex half-trace realization. Carried contextual anchors remain SGA3 I4.7.3/4.7.3.1; SGA6 VI Theorem3.3; FGA232 Theorem3.1 pp148--149; FGA236 Corollaries1.7/2.3 pp229/232; SGA7 I VII2.7/2.9 pp169--174 and VIII7.4(b) pp310--311; cote4 scan003 §§8.1.1--8.1.2, scans005--006 §§8.4.1--8.4.3 and cote54 scans005--007. No new K-theory, lambda-ring, l-adic or motivic assertion is made. The signed sufficient inequality, prescribed global period, actual selected arithmetic image, source constructor and DP44 comparison remain open.

\subsection{212. R076: the full positive dual module and exact relative norm fibre}\label{r076-the-full-positive-dual-module-and-exact-relative-norm-fibre}

The entire global-rank-three R071 support24 cross sector has absolute contribution\textless2\^{}-5710, with no sparsity premise. The full positive167 covector is counted in its actual Gram metric, with exact determinant4U and arbitrary-ideal covering radius; its last scalar is counted by the exact relative norm fibre, including nonprincipal and non-conjugation-stable ideals. Two pointwise Green bounds on the same original G patterns meet at z=2\^{}-100; full low/high exponents are -23390187/4096 and -361072601/4096. All original trace, paired targets, projected U denominator, source fibre, support, pivots and incidence weights are retained. One5710 append gives81/31 complete errors below2\^{}-230 and unchanged floors24/11. Rank-two all-II and fixed-block residuals, signed/global/source/selected-image/032/DP44 gates remain open.

The new geometric count is a fixed archimedean realization of the dual of the original positive167-module. Its determinant is kept exactly in CV9, and the rational Gram transformation is never called integral. CV11--13 count the exact scalar norm fibre over the original fractional ideal, including its ideal classes and the finite CM norm-one unit group. Both CM components, the polarization, Picard data, theta line, biextensions and selected coefficient metric remain unchanged.

Exact carried contextual anchors are SGA3 I4.7.3/4.7.3.1; SGA6 VI Theorem3.3; FGA232 Theorem3.1 pp148--149; FGA236 Corollaries1.7/2.3 pp229/232; SGA7 I VII2.7/2.9 pp169--174 and VIII7.4(b) pp310--311; cote4 scan003 §§8.1.1--8.1.2, scans005--006 §§8.4.1--8.4.3 and cote54 scans005--007. No new lambda-ring/K-theory equality, l-adic sheaf/weight theorem, CM-motive comparison or motivic L-function result follows. DP\_MULT Q3e and THETA3 G1/A7 gain exactly this smaller finite/real remainder; CX\_GRH CX-ALG-024 and GCT\_PLETH T-052 keep their distinct arithmetic comparisons.

The signed sufficient inequality, actual fully admissible mixed-CM source, selected arithmetic-image content, prescribed global Bessel/Mram value, Read/STiesR/specific032 pairing and every changed-family DP44 term remain open. The external arbitrary-ideal transference theorem and exact relative-ideal inputs are used only through the already checked DO085/K086 and K075/DO091 interfaces. The companion verifies the actual Gram minors, determinant, volume, divisor constant, cutoff inequalities, complete exponents and chronological error sums. It does not enumerate the arithmetic kernel or prove a global period by finite calculation. No Lean build or archive audit is performed.

\subsection{213. DO098: the affine central lattice and exact scalar-cross norm fibre}\label{do098-the-affine-central-lattice-and-exact-scalar-cross-norm-fibre}

The entire all-II globally rank-two-unipotent R071 scalar-cross support24 sector has absolute contribution\textless2\^{}-2879. The actual cross map and original central entry w reconstruct the full matrix. Its affine difference lattice iJ0 I\_F² retains the rank-zero G congruence, real codifferent and full discriminant, giving the complete central count2\^{}479 Q\^{}167 z\^{}-2f. The true transverse Gram volume, exact split-relative norm fibre, every original all-II paired weight and the actual selected trace are retained. All164 coordinate sectors are exhausted by64/99/1 complete monotone ranges with strict maxima -2010899/512,-41042253/4096,-11800269/4096. One2879 append gives82/31 errors below2\^{}-230 and unchanged floors24/11. Together with DO097 and R076/DO096 this removes every R071 scalar-cross support24 residual. The fixed-S rank-two sector, signed/global/source/selected-image/032/DP44 gates remain open.

The modules and ideals remain over the original Spec(O\_E), with its fixed CM pairing and complex half-trace/real trace realizations. Carried contextual anchors remain SGA3 I4.7.3/4.7.3.1; SGA6 VI Theorem3.3; FGA232 Theorem3.1 pp148--149; FGA236 Corollaries1.7/2.3 pp229/232; SGA7 I VII2.7/2.9 pp169--174 and VIII7.4(b) pp310--311; cote4 scan003 §§8.1.1--8.1.2, scans005--006 §§8.4.1--8.4.3 and cote54 scans005--007. No new K-theoretic, l-adic, motivic or source-constructor assertion follows. The signed sufficient inequality, selected arithmetic image, prescribed periods and DP44 comparison remain open.

\subsection{214. R077: fixed positive dual module and affine central congruence}\label{r077-fixed-positive-dual-module-and-affine-central-congruence}

The entire original fixed-S globally rank-two-unipotent support24 sector has absolute contribution\textless2\^{}-6845. The complete R058 high finite-rank envelope extends with exponent -119372763/4096. At low finite rank the original square pairings force the null image into the packet plane; a type-I place is impossible by the unchanged pole-budget46 norm contradiction. All remaining II/III profiles are counted through the true transverse3 dual Gram metric, exact pivot relative norm fibre and original affine central-entry lattice, retaining both CM partners, full discriminant and actual incidence weights. The complete low exponent is -168257413/24576; the high/low union contributes one6845 charge. Every old charge, including DO098, is retained, giving83/31 errors below2\^{}-230 and unchanged complete floors24/11 until a separate assembly. Signed/global/source/selected-image/032/DP44 gates remain open.

The new result counts the dual of the original positive three-module with its exact Gram determinant, an exact scalar relative norm fibre over its fractional ideal, and an affine real central-entry lattice with its full discriminant and codifferent. These live on the original Hermitian modules over Spec(O\_E), with both CM partners and the polarization retained. Contextual module/exterior anchors are SGA3 I4.7.3/4.7.3.1 and SGA6 VI Theorem3.3. Picard and biextension anchors remain FGA232 Theorem3.1 pp148--149; FGA236 Corollaries1.7/2.3 pp229/232; SGA7 I VII2.7/2.9 pp169--174 and VIII7.4(b) pp310--311. Crystals and Cartier duality retain cote4 scan003 §§8.1.1--8.1.2, scans005--006 §§8.4.1--8.4.3, and cote54 scans005--007.

No new lambda-ring equality, K-theory comparison, l-adic sheaf or weight theorem, motivic projector, CM-motive identification or L-function conclusion follows. DP\_MULT Q3e and THETA3 G1/A7 gain only the stated smaller finite/real remainder. Actual mixed-CM source construction, selected arithmetic-image content, prescribed Bessel/Mram and specific032 pairings, every actual DP44 term, CX\_GRH CX-ALG-024, and the global GCT\_PLETH T-052 comparison remain separately open. The companion is an exact arithmetic diagnostic, not an enumeration of the arithmetic kernel. No Lean build or archive audit is performed.

\subsection{215. R078: complete original support24 exhaustion}\label{r078-complete-original-support24-exhaustion}

Every actual all-top support24 term is exhausted by the complete original F/IC8 partition at compressed ranks1..33, every fixed-S nonunipotent/rank-one/rank-two class, and every scalar-cross global-rank/type/coordinate alternative. R075 covers all nonidentity positive compressions; DO090/DO092/R077 exhaust fixed S; R070/DO093/R071 reduce cross terms to the precise DO095 rank2/3 alternatives, exhausted by R076,DO096,DO097,DO098. The complete all-top support floor is25, uniform floor11. All83/31 chronological errors are unchanged and both full sums remain below2\^{}-230. No new charge. The actual remaining signed gate is J25.2; prescribed global/source/selected-image/specific032/DP44 and RH/global oscillator gates remain open.

This is an exhaustion of the original determinantal, fixed-submodule and isotropic-cross strata of integral Hermitian correspondences over Spec(O\_E), together with their actual selected complex realizations. Both CM partners, fractional ideals, coefficient modules, polarization, theta line, Picard classes and biextensions remain unchanged. Rational reconstruction in the cited counts supplies no inverse on the arithmetic theta image.

The exact contextual module/determinantal anchors are SGA3, ExposeI,4.7.3/4.7.3.1, and SGA6, ExposeVI, Theorem3.3. The Picard anchors are FGA, Bourbaki232, Theorem3.1, pp148--149, and Bourbaki236, Corollaries1.7/2.3, pp229/232. Biextensions retain SGA7 I, ExposeVII,2.7/2.9, pp169--174, and ExposeVIII,7.4(b), pp310--311. Formal groups, Dieudonne crystals and Cartier duality retain cote4 scan003, §§8.1.1--8.1.2, scans005--006, §§8.4.1--8.4.3, and cote54 scans005--007. These are contextual placements, not new comparison theorems.

No new scheme identification, lambda-ring/K-theory equality, l-adic sheaf or weight result, CM-motive comparison or motivic L-function theorem follows. DP\_MULT Q3e and THETA3 G1/A7 gain exactly J25.1--2; CX\_GRH CX-ALG-024 and GCT\_PLETH T-052 keep their distinct arithmetic gates. Settled dependencies retain their exact source locators and whole-proof checks. The present proof checks their complete partition coverage and chronological sum. No Lean build or archive audit is performed.

\subsection{216. DO099: original joint rows and the support25 structural remainder}\label{do099-original-joint-rows-and-the-support25-structural-remainder}

Both complete C-transverse directions, with every incidence profile and no global-unipotence restriction, are removed at supports25..70. At support25 the plane-C/nonzero-original-short-compression sector is also removed. The full joint row, both original columns and complete source fibre give transverse maximum -49537435/4096; the complete plane row, exact relative norm fibre and actual selected trace give -435793141/4096. Their entire union costs one12079 charge. The remaining scalar-cross support25 matrices have both cross blocks in the original packet plane, all four short compressions zero, positive166-dimensional fixed space and global rank at most3 by the support-free norm-one lemma. At26..70 only the two plane conditions are concluded. Actual84/31 ordered errors remain below2\^{}-230 with complete floors25/11. No signed/global/source/selected-image/specific032/DP44 or RH conclusion.

The original modules over Spec(O\_E), fixed CM pairing, coefficient bundle, Picard data, polarizations and biextensions are unchanged. Carried anchors remain SGA3 I4.7.3/4.7.3.1; SGA6 VI Theorem3.3; FGA232 Theorem3.1 pp148--149; FGA236 Corollaries1.7/2.3 pp229/232; SGA7 I VII2.7/2.9 pp169--174 and VIII7.4(b) pp310--311; cote4 scan003 §§8.1.1--8.1.2, scans005--006 §§8.4.1--8.4.3 and cote54 scans005--007. No new K-theory, l-adic, motivic or source-constructor assertion follows. The signed inequality, prescribed global periods, selected arithmetic content and DP44 comparison remain open.

\subsection{217. R079: actual profile cancellation and the true last-coordinate metric}\label{r079-actual-profile-cancellation-and-the-true-last-coordinate-metric}

The entire original fixed-S globally rank-two-unipotent sector at support25 is removed for every actual incidence profile. Also removed at26 are D\_G\textgreater=33 and the D\_G\textless=32 covector sector r(y1)=r(y2)=0. The original central affine lattice has poles I1/II2/III1; its T\^{}(u-s) factor cancels actual pair weights to T\^{}(-t)\textless=1. The true three-covector short pivots and one-coordinate Schur-complement last pivot retain all source/norm/trace/metric factors. Complete low exponent at25 is -27553337/24576; high endpoints are -77307919/4096 and -35243075/4096; the last-only26 endpoint is -78939175/8192. One1120 charge covers the complete union, giving85/31 errors below2\^{}-230 with unchanged complete floors25/11. Short-pivot low-rank26, fixed-S rank-one25, remaining scalar-cross25, signed/global/source/selected-image/specific032/DP44 and RH gates remain open.

The covector and its ideals live on the original Hermitian modules over Spec(O\_E); their actual dual metric and the one-coordinate Schur complement retain the polarization. The relative norm fibre and affine imaginary central line keep both CM partners and the real discriminant. The contextual module/exterior anchors remain SGA3 I4.7.3/4.7.3.1 and SGA6 VI Theorem3.3. Picard and biextension anchors are FGA232 Theorem3.1 pp148--149; FGA236 Corollaries1.7/2.3 pp229/232; SGA7 I VII2.7/2.9 pp169--174 and VIII7.4(b) pp310--311. Crystals and Cartier duality retain cote4 scan003 §§8.1.1--8.1.2, scans005--006 §§8.4.1--8.4.3 and cote54 scans005--007.

These are unchanged geometric frameworks, with no new lambda-ring equality, K-theory comparison, l-adic sheaf or weight statement, motivic projector or CM-motive identification. DP\_MULT Q3e and THETA3 G1/A7 gain only the stated smaller actual finite/real remainder. Literal admissible mixed-CM source construction, prescribed global periods and selected arithmetic-image content, the specific032 matching/value, all actual DP44 terms, CX\_GRH CX-ALG-024 and the global arithmetic comparison for GCT\_PLETH T-052 remain open. The companion is exact arithmetic verification of constants and bookkeeping, not rational-arrow enumeration. No Lean build or archive audit is performed.

\subsection{218. DO100: the canonical cross ideal and its complete large-ideal endpoint}\label{do100-the-canonical-cross-ideal-and-its-complete-large-ideal-endpoint}

The complete global-rank-three scalar-cross support25 sector with the deterministic original common-C ideal satisfying2\^{}32928 N\_EJ\textgreater=T\^{}50 is removed with one29550 charge. The whole original167-coordinate dual-metric ball, relative norm fibre, selected trace, paired projected column, source fibre, all patterns and complete actual incidence weights remain. The full monotone high-ideal endpoint is -121040253/4096. Every remaining rank-three support25 term satisfies2\^{}32928 N\_EJ\textless T\^{}50. The weaker low-ideal extension is a failed sufficient bound4090765885/12288\textgreater0; no actual contribution is inferred. A complete local typeII shear model shows why C-integrality cannot be substituted. Actual86/31 errors remain below2\^{}-230 with complete floors25/11. No whole support25, signed/global/source/selected-image/specific032/DP44 or RH conclusion.

The objects remain the same Hermitian source/target modules, original Hom and determinant sections, their paired prime ideals and their unchanged complex half-trace metric. The fixed Gram transformation and rational packet projection retain their exact determinants; neither supplies a new integral frame. The carried algebraic-geometric contextual anchors remain those explicitly recorded in R076 section6: SGA3 I4.7.3/4.7.3.1, SGA6 VI Theorem3.3, FGA232 Theorem3.1, FGA236 Corollaries1.7/2.3, SGA7 I VII2.7/2.9 and VIII7.4(b), and the already pinned cote4/cote54 locators. No new comparison theorem, lambda-ring/K-theory identity, l-adic weight assertion or CM-motive result is inferred.

DP\_MULT Q3e and THETA3 G1/A7 gain exactly HI4's smaller finite/real remainder. The signed sufficient inequality, actual fully admissible source, prescribed global Bessel/Mram value, selected arithmetic-image content, Read/STiesR/specific032 pairing and changed-family DP44 comparison remain open. The certificate verifies the actual constants, original Gram, all powers, endpoint, local scope example and chronological sums; it does not enumerate the arithmetic kernel or establish a global period by calculation.

\subsection{219. DO101: reciprocal norms and all paired pole positions}\label{do101-reciprocal-norms-and-all-paired-pole-positions}

The entire remaining globally nonunipotent rank-two scalar-cross interval25..32 is removed after the retained DO099 transverse cut. The support-independent DO076 reciprocal-line theorem gives lambda in F and inverse-C=-lambda\^{}-1 C without requiring full-Z at26+. At least23 actual lambda poles and one global inversion give at least12 positive valuations; all active lambda-unit places are retained in the complete paired C-pole cap s-12. Every original column, source fibre, pattern, pivot, support and labelled-pair weight remains. The complete worst exponent is -2610113463/20480 at32, and eight supports cost one127443 charge. Actual87/31 ordered errors remain below2\^{}-230 with complete floors25/11. Support33 is a failed sufficient bound only. Global-unipotent rank2, rank3, signed/global/source/selected-image/specific032/DP44 and RH gates remain separately open.

The original Hermitian modules over Spec(O\_E), paired ideals over Spec(O\_F), polarizations, Picard data, theta line and biextensions are unchanged. Primitive integer extraction and the actual source/target inclusions retain all local denominator information. The contextual anchors remain SGA3 I4.7.3/4.7.3.1, SGA6 VI Theorem3.3, FGA232 Theorem3.1 pp148--149, FGA236 Corollaries1.7/2.3 pp229/232, SGA7 I VII2.7/2.9 pp169--174 and VIII7.4(b) pp310--311, and the already pinned cote4/cote54 locators. No new K-theory, l-adic, motivic or source-constructor statement follows.

The signed sufficient inequality, prescribed global Bessel/Mram period, selected arithmetic-image content, literal source construction, Read/STiesR/specific032 pairing and changed-family DP44 integer comparison remain open. The exact certificate supplements the complete proof by checking every sign/pole case, all actual prime powers, all eight complete endpoints and the historical error arrays; it does not enumerate rational unitary arrows or establish a global period from finite diagnostics.

\subsection{220. LIFT078: averaging actual divisor ideals over every real sublattice}\label{lift078-averaging-actual-divisor-ideals-over-every-real-sublattice}

For every fractional ideal of the actual F and bounded nonzero real element, the entire original conjugate-content ideal sum is\textless2\^{}1065 y\^{}(21021/125). Ordered divisor ideals and full real subideal packing yield3\^{}f y\^{}(f(1+delta)) zeta\_F(1+delta)\^{}2. Atdelta1/1000 the actual imprimitive-character Euler product and all ramified factors give zeta\_F(1.001)\textless3003·5\^{}167 without GRH or a class-number estimate. Primitive v congruences also retain the exact original q affine ideal and both true eta-metric q/r packing bounds. The lemma has no sector charge, floor change or R080 dependency; the canonical87/31 error arrays remain unchanged pending a separately checked full application.

The count lives on fractional invertible modules and divisor ideals over Spec(O\_F), extended to the original Spec(O\_E) Hermitian modules. Actual ideal classes, the CM involution and the polarization are retained. It supplies an unconditional counting constant in the absolutely convergent range, not a motivic or critical-line L-function comparison. Exact carried contextual anchors are SGA3 I4.7.3/4.7.3.1; SGA6 VI Theorem3.3; FGA232 Theorem3.1 pp148--149; FGA236 Corollaries1.7/2.3 pp229/232; SGA7 I VII2.7/2.9 pp169--174 and VIII7.4(b) pp310--311; cote4 scan003 §§8.1.1--8.1.2, scans005--006 §§8.4.1--8.4.3 and cote54 scans005--007. No new lambda-ring/K-theory identity, l-adic sheaf or weight theorem, CM-motive identification or arithmetic source constructor follows. DP\_MULT Q3e and THETA3 G1/A7 await its separately counted application; CX\_GRH CX-ALG-024 and GCT\_PLETH T-052 retain their existing gates.

\subsection{221. R080: averaged original content and the complete rank-one support25 sector}\label{r080-averaged-original-content-and-the-complete-rank-one-support25-sector}

Every original fixed-S globally rank-one-unipotent support25 term is removed with one537 charge. The primitive q coordinate retains the original affine difference ideal, so both q/r true eta-metric element counts replace the older relative-norm fibre. Independently checked LIFT078 averages every conjugate-content ideal over eta via actual subideal packing and the unconditional actual-field bound zeta\_F(1.001)\textless3003·5\^{}167. Its full constant2\^{}1065 and exponent21021/125 are retained before all Green patterns, selected traces, supports, marked subsets and original paired type/label sums. All26 double-pole cells pass; worst margin538, next1530, complete witness exponent -25183329/64. Actual88/31 errors remain below2\^{}-230 with complete floors25/11. This is a full rank-one sector result, not whole support25 or signed/global/source/selected-image/specific032/DP44/RH nonvanishing.

The affine condition lives in the original Hom module over Spec(O\_E) and its fractional ideal, with both CM partners and exact content kept. The averaging sums actual ideal divisibility lattices rather than changing the coefficient field. The true eta-dependent positive metric is used only for counting actual differences. The polarization, Picard class, biextension and selected complex realization are unchanged. Exact carried contextual anchors are SGA3 I4.7.3/4.7.3.1; SGA6 VI Theorem3.3; FGA232 Theorem3.1 pp148--149; FGA236 Corollaries1.7/2.3 pp229/232; SGA7 I VII2.7/2.9 pp169--174 and VIII7.4(b) pp310--311; cote4 scan003 §§8.1.1--8.1.2, scans005--006 §§8.4.1--8.4.3 and cote54 scans005--007.

No new lambda-ring identity, K-theory map, l-adic weight statement, motivic projector or CM-motive comparison follows. DP\_MULT Q3e and THETA3 G1/A7 gain precisely the stated smaller actual finite/real remainder. Literal mixed-CM admissible source construction, selected arithmetic-image content, prescribed global periods and the specific032 pairing/value, every actual DP44 term, CX\_GRH CX-ALG-024 and the global GCT\_PLETH T-052 comparison remain separate open gates. No Lean build or archive audit is performed.

\subsection{222. DO102: padded cross ideals and the original affine central fibre}\label{do102-padded-cross-ideals-and-the-original-affine-central-fibre}

Every globally unipotent rank-two scalar-cross support25 term with at least one type-I or type-III place is removed with one881 charge. The actual full-positive normal form makes all original cross coordinates integral outside typeII. Its original central coordinate has affine difference lattice iJ0 I\_Sigma,F I\_D,F, retaining the shift, full real discriminant and750 transference coefficient. The complete u24 sector costs less than2\^{}-882, including50 type-position choices and both extra T factors. The padded u\textless=23 sector includes300 sets, all3\^{}25 type assignments and T\^{}25 pair cost, with explicitly charged cap30 removal and full sum below2\^{}-174292. Actual89/31 ordered errors remain below2\^{}-230 with complete floors25/11. The all-II support25 sector remains open; no signed/global/source/selected-image/specific032/DP44 or RH conclusion.

The original Hom modules and determinantal conditions remain over Spec(O\_E), and the affine central fibre is a torsor under a fractional ideal in the fixed real trace realization. Its contraction and norm retain both CM partners; no integral moving frame is introduced. The polarization, theta line, Picard classes and biextensions remain unchanged. Exact carried contextual anchors: SGA3 I4.7.3/4.7.3.1; SGA6 VI Theorem3.3; FGA232 Theorem3.1 pp148--149; FGA236 Corollaries1.7/2.3 pp229/232; SGA7 I VII2.7/2.9 pp169--174 and VIII7.4(b) pp310--311; cote4 scan003 §§8.1.1--8.1.2, scans005--006 §§8.4.1--8.4.3; cote54 scans005--007. These retain the formal-group, Dieudonne and Cartier-duality context without a new comparison theorem. No new lambda-ring/K-theory equality, l-adic sheaf/weight theorem, CM-motive or motivic L-function assertion is inferred. DP\_MULT Q3e and THETA3 G1/A7 gain exactly the stated finite/real gate reduction; CX\_GRH CX-ALG-024 and GCT\_PLETH T-052 retain their separate arithmetic gates.

\subsection{223. R081: original coordinate submodules and the paired common-ideal cutoff}\label{r081-original-coordinate-submodules-and-the-paired-common-ideal-cutoff}

Every remaining rank-three scalar-cross support25 term now has at least123 nonzero original S coordinates and2\^{}33096 N\_EJ\textless T\^{}50 for the same deterministic common-C ideal. The original(k+3)-coordinate covector ball keeps the true Gramdet4U, arbitrary-ideal cover, exact beta norm fibre, actual selected trace, full paired plane-column U/T index and source fibre. At z*=2\^{}(-197/2), all122 low-ideal cells have maximum -7111295/4096; the complete high-ideal endpoint is -5721469/4096. Both complete unions cost one1395 charge, retaining the old overlapping DO10029550. Actual90/31 errors remain below2\^{}-230, complete floors25/11. Dense smaller-ideal rank3, all-II unipotent rank2 and signed/global/source/selected-image/specific032/DP44/RH gates remain open.

The fixed support condition defines coordinate submodules inside the original positive Hom module over Spec(O\_E). Its dual metric and determinant4U remain explicit; the common-C ideal is a paired determinantal condition, and the norm fibre retains actual ideal classes and CM norm-one units. Rational metric transformations never supply new integral frames. The original polarization, theta line, Picard classes and biextensions remain unchanged. Contextual anchors are SGA3 I4.7.3/4.7.3.1; SGA6 VI Theorem3.3; FGA232 Theorem3.1 pp148--149; FGA236 Corollaries1.7/2.3 pp229/232; SGA7 I VII2.7/2.9 pp169--174 and VIII7.4(b) pp310--311; cote4 scan003 §§8.1.1--8.1.2,scans005--006 §§8.4.1--8.4.3; cote54 scans005--007. These carry the formal-group, Dieudonne-crystal and Cartier-duality context without a new comparison theorem. No new lambda-ring/K-theory equality, l-adic sheaf/weight assertion, CM-motive or motivic L-function conclusion follows.

DP\_MULT Q3e and THETA3 G1/A7 gain exactly DC5's smaller actual finite/real gate. The signed sufficient inequality R078 J25.2, literal fully admissible mixed-CM source, selected arithmetic-image content, prescribed global Bessel/Mram value, Read/STiesR/specific032 pairing and changed-family DP44 comparison remain open. CX\_GRH CX-ALG-024 has no new RH theorem, and GCT\_PLETH T-052 retains its separate global arithmetic charge comparison. No Lean build or archive audit is performed.

\subsection{224. R082: fixed integral incidences and actual type-conditioned counts}\label{r082-fixed-integral-incidences-and-actual-type-conditioned-counts}

R082 TP1--13 restrict every remaining rank-three support25 term, with its actual t=11..25 type-III places, to k\textgreater K\_t and2\^{}(42p\_t)N\_EJ\textless T\^{}50 for the same common-C ideal. In increasing t order, (p\_t,K\_t) are(1120,149),(1096,148),(1073,147),(1049,147),(1026,146),(1002,145),(979,144),(955,142),(932,141),(908,139),(885,137),(861,135),(838,132),(814,129),(791,125). The complete general type/pair weights and actual-t minor poles give maxima-9691203/4096 over2106 low cells and-11857677/4096 over15 high cells, one2353 charge. Actual91/31 errors stay below2\^{}-230; complete floors25/11. The all-II unipotent rank-two branch and surviving dense smaller-ideal rank-three profiles remain; signed/global/source/selected-image/specific032/DP44/RH gates remain open.

The type decomposition consists of incidences between the same integral pencil lattices over Spec(O\_E). The paired determinantal ideal and its original-coordinate submodule retain the same norms and fixed complex realization. Conditioning on t changes no source/target object, polarization, Picard class, theta line, biextension, metric or Haar measure. The two inequalities are counts of the actual integral correspondences with their prescribed normalized coefficients.

Exact contextual anchors remain SGA3 I4.7.3/4.7.3.1; SGA6 VI Theorem3.3; FGA232 Theorem3.1 pp148--149; FGA236 Corollaries1.7/2.3 pp229/232; SGA7 I VII2.7/2.9 pp169--174 and VIII7.4(b) pp310--311; cote4 scan003 §§8.1.1--8.1.2,scans005--006 §§8.4.1--8.4.3; cote54 scans005--007. No new lambda-ring/K-theory identity, l-adic sheaf or weight theorem, CM-motive comparison or motivic L-function conclusion is inferred. DP\_MULT Q3e and THETA3 G1/A7 gain exactly TP2's smaller finite/real gate; CX\_GRH CX-ALG-024 and GCT\_PLETH T-052 retain their distinct arithmetic gates. No Lean build or archive audit is performed.

\subsection{225. AM-DO-103}\label{am-do-103}

DO103 removes the complete all-II globally rank-two-unipotent scalar-cross support25 subset with N\_E denominator(eta)\textless=2\^{}120. Both actual CM pole lines force the same selected delta=eta-Q(-7+2sqrt(-7)) divisibility. The globally integral denominator-cleared delta ideal has norm at least7\^{}84 T\^{}p and below{[}4(U+14){]}\^{}168 N\_E D\_eta, so the exact host inequality gives p\textless=24. The shared cross/double-pole containing set, original affine central counts and all164 CF endpoints give25T*2\^{}-2879\textless2\^{}-1878. Actual92/31 ordered errors remain below2\^{}-230, complete floors25/11. Every remaining all-II support25 term has N\_E D\_eta\textgreater2\^{}120. The large-denominator branch, rank-three profiles, signed/global/source/selected-image/specific032/DP44/RH gates remain open.

The denominator and common pole conditions are fractional ideal conditions on the original integral correspondence over Spec(O\_E). The imaginary central affine lattice retains its integral model and discriminant; the inequality changes a counted subset, not the polarization, Poincare biextension, theta line or Haar measure. The fixed complex norm is the same realization of those original modules. Exact inherited contextual anchors: SGA3 I4.7.3/4.7.3.1; SGA6 VI Theorem3.3; FGA232 Theorem3.1 pp148--149; FGA236 Corollaries1.7/2.3 pp229/232; SGA7 I VII2.7/2.9 pp169--174 and VIII7.4(b) pp310--311; cote4 scan003 §§8.1.1--8.1.2 and scans005--006 §§8.4.1--8.4.3; cote54 scans005--007. These are contextual locators, not newly established comparison theorems. DP\_MULT Q3e and THETA3 G1/A7 retain their actual-source/period/accounting obligations; CX\_GRH CX-ALG-024 and GCT\_PLETH T-052 retain their separate gates. No new lambda-ring or K-theory operation, l-adic weight theorem or motivic L-function identity is inferred.

\subsection{226. AM-LIFT-079}\label{am-lift-079}

LIFT079 CT1--10 proves that literal CMObj and, independently, both actual LamObjGood copies force the complex target to be isogenous to a power of one CM elliptic curve, using pure complete quadratic eigenbases and the actual surjective/quasi-inverse sum. The actual LamPtClauses fibre identity transfers this constraint to the prescribed point; a point with a nonelliptic simple factor has no such witness. Thus dropping IsTen or bypassing CMObj while retaining Good, or substituting an unrelated elliptic Good object, does not supply the mixed-J point hypothesis. This is conditional on the stated point/frame premises and is not a contradiction to the accepted original source theorem. R003/DP25 simplicity is carried, not counted again. No mixed replacement constructor, finite charge, period/image/value, DP44 or RH conclusion follows.

The conclusion concerns characteristic-zero fibres of the given abelian schemes with their endomorphism actions. Up to rational isogeny their degree-one CM motives have only the quadratic elliptic simple factor; this statement forgets the specified integral model and polarization. The Picard-dual and biextension structures remain those of the actual varieties and cannot be transferred to the primitive-Jacobian fibre merely by changing a label. No claim is made that higher tensor motives or reconstructed mixed moduli are impossible. Exact inherited contextual anchors: SGA3 I4.7.3/4.7.3.1; SGA6 VI Theorem3.3; FGA232 Theorem3.1 pp148--149; FGA236 Corollaries1.7/2.3 pp229/232; SGA7 I VII2.7/2.9 pp169--174 and VIII7.4(b) pp310--311; cote4 scan003 §§8.1.1--8.1.2 and scans005--006 §§8.4.1--8.4.3; cote54 scans005--007. These are contextual locators, not newly established comparison theorems. DP\_MULT Q3e and THETA3 G1/A7 retain their actual-source/period/accounting obligations; CX\_GRH CX-ALG-024 and GCT\_PLETH T-052 retain their separate gates. No new lambda-ring or K-theory operation, l-adic weight theorem or motivic L-function identity is inferred.

\subsection{227. AM-K097}\label{am-k097}

AM-K097 removes the complete actual global-rank-three scalar-cross support25 subset with N\_E D\_eta\textless=2\^{}120, using the cross-only pole lemma, at most24 eligible places, all25 padding choices and all15 actual type profiles. The complete normalized bound is below2\^{}-11658. No rank-three full-double-pole or unipotent central-fibre assertion is imported. One11658 charge appends to the actual DO10392/31 prefix, giving93/31; both full exact sums remain below2\^{}-230 and complete floors remain25/11. Remaining rank-three terms have N\_E D\_eta\textgreater2\^{}120 and retain R082 cuts. The signed kernel, actual selected arithmetic image, prescribed global period, mixed source construction, specific032/Read, changed-family DP44 and RH remain open.

All modules remain inside the existing original Hom space over Spec(O\_E); only a containing fractional ideal changes in the count. The theta line, Picard/biextension data, polarization and complex realization are unaltered. Carried contextual anchors remain SGA3 I4.7.3/4.7.3.1; SGA6 VI Theorem3.3; FGA232 Theorem3.1 pp148--149; FGA236 Corollaries1.7/2.3 pp229/232; SGA7 I VII2.7/2.9 pp169--174 and VIII7.4(b) pp310--311; cote4 scan003 §§8.1.1--8.1.2, scans005--006 §§8.4.1--8.4.3 and cote54 scans005--007. No new comparison, motivic, K-theoretic, l-adic or source-constructor theorem is asserted.

\subsection{228. AM-R-083}\label{am-r-083}

AM-R-083 removes the entire actual global-rank-three scalar-cross support25 tail N\_E D\_eta\textgreater=2\^{}4864 with k\textgreater=126, the latter already forced on the current remainder by R082. Its full167-dimensional arbitrary-ideal covector count, exact per-denominator eta enumeration and complete ideal-zeta tail retain all original type weights, G patterns, selected traces, projected U/T column and source fibre. The complete exponent is -22330263/4096\textless-5451, charged conservatively once as5400. The actual K09793/31 prefix becomes94/31; full sums remain below2\^{}-230 and floors25/11 are unchanged. Together with K097, remaining rank-three support25 terms satisfy2\^{}120\textless N\_E D\_eta\textless2\^{}4864 and all R082 profile cuts. The all-II unipotent rank-two, signed kernel, selected arithmetic image, mixed source, prescribed global period, specific032/Read, DP44 and RH gates remain separate and open.

The exact denominator is an invertible ideal over Spec(O\_E), and L intersection eta\^{}-1 L=LD is a primewise identity in the unchanged Hom module. The Hermitian metric and original CM involution, theta/Picard and biextension data are unchanged. The finite Green exponent and convergent whole-ideal sum do not construct a new motive, l-adic weight comparison or source object. Contextual locators: SGA3 I4.7.3/4.7.3.1; SGA6 VI Theorem3.3; FGA232 Theorem3.1 pp148--149; FGA236 Corollaries1.7/2.3 pp229/232; SGA7 I VII2.7/2.9 pp169--174 and VIII7.4(b) pp310--311; cote4 scan003 sections8.1.1--8.1.2, scans005--006 sections8.4.1--8.4.3; cote54 scans005--007. These are carried context, not fresh primary-source readings.

\subsection{229. AM-DO-104}\label{am-do-104}

AM-DO-104 removes the complete actual all-II globally rank-two-unipotent scalar-cross support25 tail N\_E D\_eta\textgreater=2\^{}1820. Every original coordinate lies in LD\_eta; the full original4-by4 block and both mixed blocks give the exact affine central difference ideal iJ0 B\_eta I\_Sigma,F\^{}2 with D\_eta bar(D\_eta)=B\_eta O\_E. The full real750 cover, all167 covector values, actual selected trace, all164 row cells, original weights and the entire ideal-zeta tail are retained. The complete maximum -10346711/4096\textless-2526 yields one2518 charge. Actual R08394/31 becomes95/31, complete sums remain below2\^{}-230 and floors25/11 remain. Together with DO103 the remaining all-II rank-two branch has2\^{}120\textless N\_E D\_eta\textless2\^{}1820. No whole support25, signed-kernel, prescribed selected arithmetic image/period, mixed source, specific032/Read, DP44 or RH conclusion.

The ideals are invertible modules on the unchanged Spec(O\_E)/Spec(O\_F), with the original CM involution and polarization. Carried contextual anchors remain SGA3 I4.7.3/4.7.3.1; SGA6 VI Theorem3.3; FGA232 Theorem3.1 pp148--149; FGA236 Corollaries1.7/2.3 pp229/232; SGA7 I VII2.7/2.9 pp169--174 and VIII7.4(b) pp310--311; cote4 scan003 sections8.1.1--8.1.2, scans005--006 sections8.4.1--8.4.3, and cote54 scans005--007. No new source constructor, motivic comparison, signed-kernel conclusion, selected derivative-image primitivity, prescribed Bessel period or DP44/RH conclusion follows.

\subsection{230. AM-LIFT-080}\label{am-lift-080}

LIFT080 MT1--17 proves a generic Step5 transfer for actual mixed-type comparison lines with individually bounded normalized degrees d\_j\textgreater=d\_-. Rescaling only their archimedean norms by exp(d\_j-e) makes all auxiliary degrees e, keeps every finite ideal and local polynomial quantifier unchanged, and replaces La=L by L exp((e-d\_-)\emph{+). The exact cost is C(e)=(1/2+\textbar e\textbar+log(L)/2)(D+gn)+D log(rL)+(3D+gn)(e-d}-)\emph{+/2. Its optimum within this comparison family is C0+(-d}-)\_+ min(D+gn,(3D+gn)/2). Arbitrary-type isogeny derivatives and the literal arbitrary-set form\_letter\_le supply independence and norm mechanisms; explicit mixed products J\^{}(14h p) with complementary-action factors supply geometric inclusions/projections with N=L=1 and permutation U. The actual prescribed point/model/section, all-extension polynomial tests, effective eigenline degrees, coefficient-field/frame ties and fully revised source/DP44 room inequality remain required. The old LamPtClauses/LamObjGood pure-type boundary remains; no main-family constructor, finite charge, period/image/value, specific032/Read or RH conclusion follows.

The comparison lines are invertible fractional-ideal modules on the unchanged spectrum of the coefficient integers, equipped with positive metrics at every counted embedding. The auxiliary metric change preserves their algebraic modules, determinant functor, integral Hom maps and all finite local tests. It is an archimedean arithmetic-degree comparison, not an isomorphism of the original metrized Hodge lines and not a new CM height equality. The explicit product abelian schemes, their duals, Picard data and polarization biextensions retain their actual geometry; the isogeny derivative is a characteristic-zero tangent construction. Formal groups/Dieudonne crystals at finite places and the existing lambda-ring/K-theory and motivic realizations are unchanged by the auxiliary norms. No new l-adic weight, motivic L-function or mixed source comparison is inferred. Carried contextual locators: SGA3 I4.7.3/4.7.3.1; SGA6 VI Theorem3.3; FGA232 Theorem3.1 pp148--149; FGA236 Corollaries1.7/2.3 pp229/232; SGA7 I VII2.7/2.9 pp169--174 and VIII7.4(b) pp310--311; cote4 scan003 sections8.1.1--8.1.2, scans005--006 sections8.4.1--8.4.3; cote54 scans005--007. Programme rows: DP\_MULT Q3e/DP44/DP52/DP54--65; THETA3 G1/A7; CX\_GRH CX-ALG-024; GCT\_PLETH T-052. These locators supply context, not a claimed newly checked comparison theorem.

\subsection{231. AM-R-084}\label{am-r-084}

AM-R-084 exhausts every remaining global-rank-three scalar-cross support25 term. The complete K097 count applies to \textless=24 eligible places without a denominator premise; the exact DO093 inversion carries this count and R083 to the inverse null line with the same original S-supportk\textgreater=126, actual weights, masses, phases and cutoff. Outside the four charged sets both slopes are fully eligible and each denominator norm is\textless2\^{}4864. Their denominator-cleared difference has the full paired T\^{}50 divisor but norm\textless(4U)\^{}168*2\^{}9728\textless T\^{}50, so the two lines coincide. DO095 then forces global rank2, a contradiction. The entire four-sector absolute union is\textless2\^{}-5398, with both inverse costs included. Actual DO10495/31 becomes96/31; every old overlapping charge remains, full sums are\textless2\^{}-230, and complete floors25/11 remain unchanged. At support25 only the separately retained all-II globally rank-two-unipotent scalar-cross branch remains, with2\^{}120\textless N\_E D\_eta\textless2\^{}1820. Higher supports and signed-kernel, selected arithmetic-image, prescribed global period, actual mixed source, specific032/Read, DP44 and RH gates remain open.

The integral ideal IL15 is an invertible module on the unchanged Spec(O\_E). Its norm is the determinant of multiplication combined with the original ideal norm, and the nonzero section is tested at every complex realization and all selected paired finite points. Inversion is the actual transpose of the same Hermitian correspondence and its symmetric pencil coefficients. The null-line equality concerns the rational Hom module and the resulting rank stratum; it supplies no globally integral null basis. The CM abelian variety, theta line, Picard scheme, polarization and biextension remain unchanged. Formal groups/Dieudonne crystals and Cartier duality remain those already constructed; no new lambda-ring/K-theory equality, l-adic weight theorem, motivic comparison or motivic L-function result is inferred.

Carried exact contextual locators are SGA3 I4.7.3/4.7.3.1; SGA6 VI Theorem3.3; FGA232 Theorem3.1 pp148--149; FGA236 Corollaries1.7/2.3 pp229/232; SGA7 I VII2.7/2.9 pp169--174 and VIII7.4(b) pp310--311; cote4 scan003 sections8.1.1--8.1.2, scans005--006 sections8.4.1--8.4.3; cote54 scans005--007. DP\_MULT Q3e and THETA3 G1/A7 gain the exhaustive rank-three support25 reduction. CX\_GRH CX-ALG-024 and GCT\_PLETH T-052 retain their separate arithmetic gates. No accepted Alpöge theorem or archive is audited, and no Lean build is performed.

\subsection{232. AM-LIFT-081}\label{am-lift-081}

LIFT081 ED1--23 supplies effective intrinsic line-degree lower bounds for every actual comparison component eigenline on J14/J56, with the full integral CM action and the literal source metric n\^{}2=pi H. The good model is obtained over K0(alpha,beta), alpha\^{}7=-4 and beta\^{}2=(zeta7-1)\^{}7, of degree at most14 over K0. The actual integral vector 2q\_h p\_tau e\_i has norm below sqrt(2)q\_h at every counted embedding, with q1=406 and q4=26390. Therefore ndeg\textgreater-log(sqrt(2)q\_h), and the weaker common bounds are -10log(2) for J14 and -16log(2) for J56. These normalized degrees are invariant under further finite base change of the actual good model. Substitution into LIFT080 makes its sufficient conditional Step5 cost fully explicit: Copt=1/2(1+log L)(D+V)+D log(rL)+b\_h min(D+V,(3D+V)/2), V=gn and b\_h=10log(2) or16log(2). This resolves the effective comparison-line-degree premise only. The prescribed mixed point, model/level/coefficient-field ties, nonzero frame polynomial, all residue-extension Step5Loc conditions, source-room data, selected arithmetic image/content, global period and DP44 evaluation remain required. No claim is made about an arbitrary line in a larger multiplicity space, equal degrees, a primitive coordinate-ideal equality, an old LamPtClauses witness, a finite charge, or RH.

The saturated eigenline is the intersection of a characteristic-zero CM eigenspace with the actual tangent bundle of the good abelian scheme over Spec(O\_K). The trace-dual separability projector is made integral by the displayed inverse-different clearing integer; no integral splitting at ramified primes is asserted. Its Rosati-compatible complex realization is an orthogonal projection. One global integral vector controls the same metrized invertible module at all embeddings; the fractional-ideal norm and product formula give its lower arithmetic degree. This preserves the original Jacobian, dual/Picard data, principal polarization, Poincare biextension and theta line. Auxiliary metric equalization in LIFT080 changes an archimedean comparison cost only. Formal groups and their Dieudonne crystals, existing lambda-ring/K-theory operations and the motivic realizations remain the same. No new l-adic weight, motivic L-function identity or global arithmetic source comparison is inferred. Carried contextual locators: SGA3 I4.7.3/4.7.3.1; SGA6 VI Theorem3.3; FGA232 Theorem3.1 pp148--149; FGA236 Corollaries1.7/2.3 pp229/232; SGA7 I VII2.7/2.9 pp169--174 and VIII7.4(b) pp310--311; cote4 scan003 sections8.1.1--8.1.2, scans005--006 sections8.4.1--8.4.3; cote54 scans005--007. Programme rows: DP\_MULT Q3e/DP44/DP52/DP54--65; THETA3 G1/A7; CX\_GRH CX-ALG-024; GCT\_PLETH T-052. These are contextual locators, not newly proved comparison theorems.

\subsection{233. AM-DO-105}\label{am-do-105}

AM-DO-105 removes the entire actual all-II globally rank-two-unipotent scalar-cross support25 subset with at most24 eligible places, without any denominator premise, together with the fully eligible middle window2\^{}120\textless N\_E D\_eta\textless2\^{}1820 having1\textless=k\textless=21 nonzero original S coordinates. For each fixed support there is at most one fully eligible eta in this window and at most one actual diagonal pencil label per active place, including central-only poles. The original weighted label sum is at most(T-1)\^{}-25. The joint S/transverse beta integral, moving central cover and simultaneous x/z endpoint retain all full metric/phase/class/Haar factors. The worst claimed cell -1852891/4096\textless-452 and full fresh union give one451 charge. Actual R08496/31 becomes97/31, both complete sums remain below2\^{}-230 and complete floors25/11 remain. After R084, DO103/104 and this theorem the only support25 branch left is fully eligible all-II rank-two unipotent, with2\^{}120\textless N\_E D\_eta\textless2\^{}1820 and22\textless=k\textless=164, unique eta per support and at most one label per active place. The dense range, all higher-support cuts, signed inequality, actual mixed source, selected arithmetic image/content, prescribed global period, specific032/Read, DP44 and RH gates remain.

The original fractional ideals and affine fibres are invertible-module conditions on Spec(O\_E)/Spec(O\_F), with the same CM involution, Hermitian polarization, theta line, biextension and complex metric. Exact inherited contextual anchors remain SGA3 I4.7.3/4.7.3.1; SGA6 VI Theorem3.3; FGA232 Theorem3.1 pp148--149; FGA236 Corollaries1.7/2.3 pp229/232; SGA7 I VII2.7/2.9 pp169--174 and VIII7.4(b) pp310--311; cote4 scan003 sections8.1.1--8.1.2 and scans005--006 sections8.4.1--8.4.3; cote54 scans005--007. No new lambda-ring/K-theory equality, l-adic or motivic comparison is inferred. DP\_MULT Q3e/DP44/DP52/DP54--65, THETA3 G1/A7, CX\_GRH CX-ALG-024 and GCT\_PLETH T-052 retain their actual arithmetic and global comparison gates.

\subsection{234. AM-LIFT-082}\label{am-lift-082}

LIFT082 FT1--29 transfers one actual nonzero frame polynomial to every Step5Loc prime and complete-DVR extension, and every Step5Arch embedding, from explicit arbitrary-type geometric inputs. Generic FrPt polynomial existence and same-point nonzero evaluation produce one G before any completion. Integral eigenmodule bases commute with arbitrary compatible DVR extension by flatness; no integral eigenspace splitting is assumed. The exact relation Ub\^{}t Y=X, intrinsic theta generator unit, determinant inclusion and physical row coefficient-clearing indices give all local clauses. The original pi-normalized Hermitian metric, actual fibre/section maps and goodOf conjugation give all archimedean clauses, retaining lieFac and the dual theta norm. There is no new transfer cost beyond LIFT080/081. The actual prescribed point with Dw.mult, good model/level/Key/section/coefficient-field ties, valSet arithmetic-lattice identification, same-point nonzero F value, Host/Sizes, Read/STiesR, selected arithmetic image/content, global period, specific032 and DP44 remain required. No old CMPt.Good/LamObjGood/LamPtClauses instance, finite charge, support-floor change or RH assertion follows.

The actual abelian scheme, its tangent bundle and intrinsic theta/determinant invertible modules live on Spec(O\_K), with the given full CM order action. Integral eigenmodules are common kernels over each DVR and commute with flat base change, including ramified extensions; they need not split the integral tangent lattice. The section is the actual generator of the polarized theta line and is transported with its rigidification, Picard data and Poincare biextension. The same algebraic frame polynomial is evaluated under every local and complex base change. The pi-Hermitian form and determinant volume retain the original complex realization and arithmetic degree. Formal groups/Dieudonne crystals, lambda-ring/K-theory operations and motivic realizations are not altered by this conditional transfer. No new l-adic weight theorem, motivic L-function identity or global arithmetic comparison is inferred. Exact carried contextual locators: SGA3 I4.7.3/4.7.3.1; SGA6 VI Theorem3.3; FGA232 Theorem3.1 pp148--149; FGA236 Corollaries1.7/2.3 pp229/232; SGA7 I VII2.7/2.9 pp169--174 and VIII7.4(b) pp310--311; cote4 scan003 sections8.1.1--8.1.2, scans005--006 sections8.4.1--8.4.3; cote54 scans005--007. Programme rows DP\_MULT Q3e/DP44/DP52/DP54--65, THETA3 G1/A7, CX\_GRH CX-ALG-024 and GCT\_PLETH T-052 retain their stated remaining gates.

\subsection{235. AM-LIFT-083}\label{am-lift-083}

LIFT083 JS1--25 proves the actual pinned signature obstruction for every J-isotypic complex fibre, independently of CMObj/IsTen/LamObjGood. For any quadratic Host action into End0(J\^{}N)=M\_N(Q(zeta7)), d=7 gives Lie multiplicities(2a+b,a+2b), a+b=N, and every other Host quadratic field gives balanced(3N/2,3N/2). Actual DetCond requires(f\_E p,f\_E q). Hence d=7 needs p\textless=2q and q\textless=2p, while d!=7 needs p=q. The literal pinned profiles p\_j=k+j−1,q\_j=m−j are unbalanced at every stage and terminal q\_m=0; under full SizesAt, k\textgreater m\^{}5 f\_R\textgreater=2m forces p\_j\textgreater2q\_j at every stage. Thus no solely J-isotypic point can satisfy any pinned stage under those full premises, even with enlarged field/changed primes and the old good-point wrappers bypassed. An actual enlargement A\textasciitilde J\^{}N×B must have N\textless=f\_E min(p,q) and dim B\textgreater=f\_E max(p−2q,q−2p,0); the terminal fibre contains no J factor. These are necessary geometric dimensions, not height or DP44 costs. Other CM isogeny classes, their actual full point/action/model/level/section and same-point nonzero value remain open, together with selected arithmetic content, prescribed global period, Read/STiesR/specific032 and DP44. No finite charge, support-floor change or RH conclusion follows.

The obstruction is a statement about the characteristic-zero geometric fibre of the actual abelian scheme and its full CM action: the rational isotypic projector of h\^{}1(J) is central, and the Hodge/Lie multiplicities of every quadratic action constrain the determinant condition. It persists under geometric isogeny and coefficient-field extension but does not assert any integral eigenspace splitting or identify a new polarized lattice. The complementary weight-one factor and its Picard/polarization/Poincare-biextension data still require construction. Formal groups and Dieudonne crystals of the J factors retain the accepted DP slopes; the dimension inequality supplies no stronger divided-power gain, lambda-ring or K-theory identity, l-adic weight bound or motivic L-function identity. Exact inherited contextual locators: SGA3 I4.7.3/4.7.3.1; SGA6 VI Theorem3.3; FGA232 Theorem3.1 pp148--149; FGA236 Corollaries1.7/2.3 pp229/232; SGA7 I VII2.7/2.9 pp169--174 and VIII7.4(b) pp310--311; cote4 scan003 sections8.1.1--8.1.2 and scans005--006 sections8.4.1--8.4.3; cote54 scans005--007. DP\_MULT Q3e/DP44/DP52/DP54--65 and THETA3 G1/A7 now need a point in a sufficiently enlarged isogeny class or a changed signature interface; CX\_GRH CX-ALG-024 and GCT\_PLETH T-052 retain their separate open gates.

\subsection{236. AM-K099}\label{am-k099}

AM-K099 GE1--13 restricts the actual remaining fully eligible support25 branch to total eligible count q=25 or26. The norm bounds give x\textgreater2\^{}120 for q25, x\textgreater2\^{}1116 for q26, and x\textgreater2\^{}2111 for q27, excluding q27 under x\textless2\^{}1820. Over K0=Q(sqrt(-7)), the relative Galois group of order168 acts regularly on the selected T primes. At q25 the scalar has trivial relative stabilizer and generates E over K0; at q26 the relative stabilizer has order at most2. The absolute degrees are at least168 and84 respectively. Distinct actual norm-U scalars in the bounded denominator window have eligible-set overlap at most14; set and scalar stabilizers agree. A relative stabilizer of order2 forces the denominator ideal to extend from the fixed field and x to be a square, even at ramified primes. This gives necessary structural restrictions only, no numerical sector removal, new error charge, source construction, signed positivity, period, DP44 or RH claim. The existing97/31 error arrays and complete floors25/11 remain unchanged.

The scalar and its denominator define an actual fractional ideal on Spec(O\_E); denominator descent is valuationwise and includes ramification. The Galois argument concerns this arithmetic datum only and asserts no symmetry of the selected kernel. The original fractional ideals and affine fibres are invertible-module conditions on Spec(O\_E)/Spec(O\_F), with the same CM involution, Hermitian polarization, theta line, biextension and complex metric. Exact inherited contextual anchors remain SGA3 I4.7.3/4.7.3.1; SGA6 VI Theorem3.3; FGA232 Theorem3.1 pp148--149; FGA236 Corollaries1.7/2.3 pp229/232; SGA7 I VII2.7/2.9 pp169--174 and VIII7.4(b) pp310--311; cote4 scan003 sections8.1.1--8.1.2 and scans005--006 sections8.4.1--8.4.3; cote54 scans005--007. No new lambda-ring/K-theory equality, l-adic or motivic comparison is inferred. DP\_MULT Q3e/DP44/DP52/DP54--65, THETA3 G1/A7, CX\_GRH CX-ALG-024 and GCT\_PLETH T-052 retain their actual arithmetic and global comparison gates.

\subsection{237. AM-K100}\label{am-k100}

AM-K100 QE1--16 removes the complete actual fully eligible all-II globally rank-two-unipotent scalar-cross support25 subset with2\^{}1116\textless=x=N\_E D\_eta\textless2\^{}1820 and22\textless=k\textless=48, including either total eligible count25 or26. The exact1116 endpoint, cap1117, complete joint row/transverse beta count, central volume\textless2\^{}679 and unique-eta/label weight give worst k48 ceiling-1730011/4096\textless-422; the26 smaller cells and this cell form one complete2\^{}-421 union. All k\textless=21 contributions remain in DO105451 and are not added again here. Append421 to the actual DO10597/31 ordered totals, giving98/31; both entire sums remain below2\^{}-230 with complete floors25/11 unchanged. In particular all q26 cells up through k48 are removed. The surviving support25 possibilities have2\^{}120\textless x\textless2\^{}1116 with22\textless=k\textless=164, or2\^{}1116\textless=x\textless2\^{}1820 with49\textless=k\textless=164, in addition to K099 structural restrictions. No claim for these surviving cells, no actual source, selected arithmetic image, prescribed period, signed positivity, specific032, DP44 or RH conclusion follows.

The original fractional ideals and affine fibres are invertible-module conditions on Spec(O\_E)/Spec(O\_F), with the same CM involution, Hermitian polarization, theta line, biextension and complex metric. Exact inherited contextual anchors remain SGA3 I4.7.3/4.7.3.1; SGA6 VI Theorem3.3; FGA232 Theorem3.1 pp148--149; FGA236 Corollaries1.7/2.3 pp229/232; SGA7 I VII2.7/2.9 pp169--174 and VIII7.4(b) pp310--311; cote4 scan003 sections8.1.1--8.1.2 and scans005--006 sections8.4.1--8.4.3; cote54 scans005--007. No new lambda-ring/K-theory equality, l-adic or motivic comparison is inferred. DP\_MULT Q3e/DP44/DP52/DP54--65, THETA3 G1/A7, CX\_GRH CX-ALG-024 and GCT\_PLETH T-052 retain their actual arithmetic and global comparison gates.

\subsection{238. AM-DO-106}\label{am-do-106}

AM-DO-106 CFI1--24 proves the complete actual central fibre Gamma=B\_eta I\_Sigma,F squared, covolume sqrt(D\_F) x T\^{}-50, sufficient period Q T squared B\_eta of index Q\^{}168 T\^{}386, and the full-real codifferent dual. The original active type-II factor has Fourier l1 norm2/T for integral cross rows and1/(T-1) for a fixed cross pole. The actual mixed G projectors and every inactive-T factor are retained by the finite positive-projection trace lemma. The center-independent whole-C Green budget gives total finite Fourier l1 at most2 Q\^{}(-335*167)(N\_E J\_C)\^{}(335/2)(2/T)\^{}j(T-1)\^{}(-(25-j)), with no finite-index loss. The complete orbit lies in the original finite support; regularized Poisson retains its exact signs and covolume. Extending a disjoint collection past H133 uses at most the single complete2\^{}-2000 tail once, subject to reconciliation of earlier removals. No numerical alias estimate, sector removal, new charge or floor change follows; current98/31 errors remain. Actual mixed source, selected arithmetic content, prescribed period, signed inequality, specific032, DP44 and RH remain open.

The exact fractional ideals and codifferent define the original arithmetic central correspondence over Spec(O\_F), including ramified and dyadic rational primes; no rational moving basis is declared globally integral. The original hermitian vector bundle and finite representation are unchanged. The original fractional ideals and affine fibres are invertible-module conditions on Spec(O\_E)/Spec(O\_F), with the same CM involution, Hermitian polarization, theta line, biextension and complex metric. Exact inherited contextual anchors remain SGA3 I4.7.3/4.7.3.1; SGA6 VI Theorem3.3; FGA232 Theorem3.1 pp148--149; FGA236 Corollaries1.7/2.3 pp229/232; SGA7 I VII2.7/2.9 pp169--174 and VIII7.4(b) pp310--311; cote4 scan003 sections8.1.1--8.1.2 and scans005--006 sections8.4.1--8.4.3; cote54 scans005--007. No new lambda-ring/K-theory equality, l-adic or motivic comparison is inferred. DP\_MULT Q3e/DP44/DP52/DP54--65, THETA3 G1/A7, CX\_GRH CX-ALG-024 and GCT\_PLETH T-052 retain their actual arithmetic and global comparison gates.

\subsection{239. AM-R-085}\label{am-r-085}

AM-R-085 RF1--31 gives the exact genuine selected vector-valued central coefficient c(t)=P(t)/(z0-epsilon i gamma t)\^{}s, with s=kappa+K+2N and degree P\textless=s-kappa, using the original Pieri component and its exact Hilbert/PBW norm. The finite coefficient rule, partial fractions, signed gamma Fourier transform and affine Poisson phases are explicit. The positive central density for the same selected component has mass one and pointwise dominates the transformed commuting fibre, without assuming its projection commutes. Products over all168 labels lie in the actual signed orthant and vanish on its boundary faces. Ordinary absolutely convergent Poisson holds for the actual lattice and finite characters supplied separately by DO106, with no quotient-cardinality loss. This removes the unspecified analytic-model step but proves no numerical dual-lattice alias bound, finite sector removal or new error charge. Current98/31 totals and floors25/11 remain; the signed inequality, actual source, selected arithmetic image, prescribed period, specific032, DP44 and RH remain open.

The finite polynomial extraction takes place in the prescribed homogeneous Hermitian bundle and its selected Schur module; the central fibre and Fourier trace pairing are the arithmetic ones. The result constructs an analytic model only, not a new arithmetic or motivic comparison. The original fractional ideals and affine fibres are invertible-module conditions on Spec(O\_E)/Spec(O\_F), with the same CM involution, Hermitian polarization, theta line, biextension and complex metric. Exact inherited contextual anchors remain SGA3 I4.7.3/4.7.3.1; SGA6 VI Theorem3.3; FGA232 Theorem3.1 pp148--149; FGA236 Corollaries1.7/2.3 pp229/232; SGA7 I VII2.7/2.9 pp169--174 and VIII7.4(b) pp310--311; cote4 scan003 sections8.1.1--8.1.2 and scans005--006 sections8.4.1--8.4.3; cote54 scans005--007. No new lambda-ring/K-theory equality, l-adic or motivic comparison is inferred. DP\_MULT Q3e/DP44/DP52/DP54--65, THETA3 G1/A7, CX\_GRH CX-ALG-024 and GCT\_PLETH T-052 retain their actual arithmetic and global comparison gates.

\subsection{240. AM-K101}\label{am-k101}

AM-K101 UC1--17 removes the complete actual fully eligible all-II globally rank-two-unipotent scalar-cross support25 subset with2\^{}1585\textless=x=N\_E D\_eta\textless2\^{}1820 for every22\textless=k\textless=164. It uses the original DO105 all-factor bound with a single common denominator endpoint, exact capmax(161,k), full selected trace and moving affine central cover. The143-cell union is below2\^{}-254, with nondense maximum-1412571/4096\textless-344 and dense maximum-1048027/4096\textless-255. Append254 once to actual K10098/31, giving99/31 and both whole sums below2\^{}-230; all earlier overlaps,451 and421 are retained and complete floors25/11 remain. The only still-unremoved support25 cells therefore have2\^{}120\textless x\textless2\^{}1585,22\textless=k\textless=164, with K100 additionally excluding22\textless=k\textless=48 when x\textgreater=2\^{}1116, and all K099 structural restrictions retained. No prospective inactive-T gain, signed positivity, actual source, selected arithmetic content, prescribed period, specific032, DP44 or RH result is inferred.

The original fractional ideals and affine fibres are invertible-module conditions on Spec(O\_E)/Spec(O\_F), with the same CM involution, Hermitian polarization, theta line, biextension and complex metric. Exact inherited contextual anchors remain SGA3 I4.7.3/4.7.3.1; SGA6 VI Theorem3.3; FGA232 Theorem3.1 pp148--149; FGA236 Corollaries1.7/2.3 pp229/232; SGA7 I VII2.7/2.9 pp169--174 and VIII7.4(b) pp310--311; cote4 scan003 sections8.1.1--8.1.2 and scans005--006 sections8.4.1--8.4.3; cote54 scans005--007. No new lambda-ring/K-theory equality, l-adic or motivic comparison is inferred. DP\_MULT Q3e/DP44/DP52/DP54--65, THETA3 G1/A7, CX\_GRH CX-ALG-024 and GCT\_PLETH T-052 retain their actual arithmetic and global comparison gates.

\subsection{241. AM-DO-107}\label{am-do-107}

AM-DO-107 IP1--23 removes three complete subsets of the actual fully eligible all-II globally rank-two-unipotent scalar-cross support25 middle sector: any inactive paired-C-zero place for all k22..164; no paired zero with x\textgreater=2\^{}744 for all k22..164; and no paired zero with x\textless2\^{}744 and22\textless=k\textless=99. The literal included-Lambda projector gives rho=1/(T-1) at every noneligible inactive place with a nonzero paired C block. A paired zero instead supplies an actual extra normT² row ideal, giving T\^{}-335 while retaining the old central affine cover. The complete143/143/78-cell union is below2\^{}-392. Append392 to the actual K10199/31 arrays, giving100/31; all historical overlaps including451,421,254 remain, both full sums are below2\^{}-230, and complete floors25/11 are unchanged. The only support25 branch left has no inactive paired C zero, total eligibility exactly25,2\^{}120\textless N\_E D\_eta\textless2\^{}744 and100\textless=k\textless=164, with K099 field/stabilizer restrictions retained. This is a strict smaller finite gate, not support25 exhaustion, signed positivity, an admissible mixed source, selected arithmetic image, prescribed global period, specific032, DP44 or RH.

The original fractional ideals and affine fibres are invertible-module conditions on Spec(O\_E)/Spec(O\_F), with the same CM involution, Hermitian polarization, theta line, biextension and complex metric. Exact inherited contextual anchors remain SGA3 I4.7.3/4.7.3.1; SGA6 VI Theorem3.3; FGA232 Theorem3.1 pp148--149; FGA236 Corollaries1.7/2.3 pp229/232; SGA7 I VII2.7/2.9 pp169--174 and VIII7.4(b) pp310--311; cote4 scan003 sections8.1.1--8.1.2 and scans005--006 sections8.4.1--8.4.3; cote54 scans005--007. No new lambda-ring/K-theory equality, l-adic or motivic comparison is inferred. DP\_MULT Q3e/DP44/DP52/DP54--65, THETA3 G1/A7, CX\_GRH CX-ALG-024 and GCT\_PLETH T-052 retain their actual arithmetic and global comparison gates.

\subsection{242. AM-DO-108}\label{am-do-108}

AM-DO-108 AC1--15 removes the entire actual support25 fully eligible all-II globally rank-two-unipotent scalar-cross subset with any active paired-C-integral place, at cost less than2\^{}-55281. The exact extra normT² row ideal gives T\^{}-335; the central cover uses the old scale. Its143 complete cells append55281 to actual DO107100/31, giving101/31, all historical charges retained and both sums below2\^{}-230, floors25/11 unchanged. Remaining middle data have j=b=0, exactly25 total eligible places,2\^{}120\textless x\textless2\^{}744 and100\textless=k\textless=164. Every complete actual central orbit then keeps all168 original T factors equal to1/(T-1), with finite weight(T-1)\^{}-168 a(h(t)), sufficient periodQGamma, indexQ\^{}168, exact trace dual and ordinary Poisson formula AC15. This strictly removes active-integral profiles and removes their indicator from the exact Fourier interface; it gives no alias estimate or automatic short-sum/fibre replacement. Admissible source, selected arithmetic content, signed nonvanishing, global period, specific032, DP44 and RH remain open.

The original fractional ideals and affine fibres are invertible-module conditions on Spec(O\_E)/Spec(O\_F), with the same CM involution, Hermitian polarization, theta line, biextension and complex metric. Exact inherited contextual anchors remain SGA3 I4.7.3/4.7.3.1; SGA6 VI Theorem3.3; FGA232 Theorem3.1 pp148--149; FGA236 Corollaries1.7/2.3 pp229/232; SGA7 I VII2.7/2.9 pp169--174 and VIII7.4(b) pp310--311; cote4 scan003 sections8.1.1--8.1.2 and scans005--006 sections8.4.1--8.4.3; cote54 scans005--007. No new lambda-ring/K-theory equality, l-adic or motivic comparison is inferred. DP\_MULT Q3e/DP44/DP52/DP54--65, THETA3 G1/A7, CX\_GRH CX-ALG-024 and GCT\_PLETH T-052 retain their actual arithmetic and global comparison gates.

\subsection{243. AM-LIFT-084}\label{am-lift-084}

AM-LIFT-084 MP1--25 constructs a literal field-valued Obj and FrPt for any given actual source Host, compatible Par.PinHost/PinLevel/PinDat, specified stage Datum, existing Moduli representation and same Coeff field, after a finite extension. The actual principal ideal identity is derived from Host fields. A pure principally polarized CM block and its conjugate action supply every uniform(p,q); the explicit trace-one semicharacter supplies actual SAB prim/adjoint/cubical lines. Relative-unramified selfdual Hermitian lattice uniqueness, including dyadic places, gives the required profinite Tate isometry; the literal commutator sign and CRT preserve the exact finite level M=g0g1,A=1. All-a normalized line identities descend over one finite field; only geometric full Tate systems are asserted. Separately, on the existing E\_h fields the explicit pure block and three J-action twists give J\^{}(fq) times C\_h\^{}(p-2q) for p\textgreater=2q, attaining LIFT083 complement dimension f(p-2q). Those fixed fields still fail full Sizes. No effective field/frame/line costs, good-model/old LamPt identification, prescribed same-point F nonzero, actual full Host/Sizes constructor, selected arithmetic content, prescribed period, specific032, DP44 or RH is supplied. No numerical error charge;101/31 arrays and floors25/11 remain.

This is a principally polarized abelian scheme over a number field, with full CM integer-ring action, rigidified symmetric cubical theta line, literal adjoint biextension and finite level; a spread over an open part of Spec(O\_K) is qualitative only and no needed bad-prime control is inferred. Exact contextual anchors are FGA232 Theorem3.1 pp148--149 and FGA236 Corollaries1.7/2.3 pp229/232 for Picard representability; SGA7 I VII2.7/2.9 pp169--174 and VIII7.4(b) pp310--311 for biextensions/duality; cote4 scan003 sections8.1.1--8.1.2 and scans005--006 sections8.4.1--8.4.3, and cote54 scans005--007 for the carried formal-group/Cartier-dual language; SGA3 I4.7.3/4.7.3.1 and SGA6 VI Theorem3.3 remain contextual. The isogeny decomposition describes CM h1 motives rationally, but no new integral crystal comparison, lambda-ring/K-theory identity, l-adic weight estimate or motivic L-function consequence is asserted. DP\_MULT Q3e/DP44/DP52/DP54--65, THETA3 G1/A7, CX\_GRH CX-ALG-024 and GCT\_PLETH T-052 retain their effective arithmetic/global gates.

\subsection{244. AM-R-086}\label{am-r-086}

AM-R-086 ZJ1--21 constructs principal source points with actual J content over any fixed actual Host E under the retained LIFT084 input hypotheses. The auxiliary compositum L=E Q(zeta7), c={[}L:E{]} in\{1,2,3,6\}, has an explicit integral CM trace polarization of degree\textbar D\_L\textbar7\^{}(cf/3). Zarhin Theorem1.1 gives a full integer-ring compatible principal block. Its exact conjugate-action balancing, with the even-c improvement, yields block sizes a\_c=(8,8,24,24). For p\textgreater=q the resulting exact J exponent is N=(2f/3)(q-r), r=q mod a\_c, and pure complementary dimension f(p-q+2r). If d=7 and p\textgreater=2q, a principal(16,8) block gives N=f(q-r7), r7=q mod8, complement dimension f(p-2q+3r7). Zero remainders attain the corresponding LIFT083 bounds. The unchanged Datum, finite level and Coeff give actual Obj/FrPt over a finite extension. A separate full O\_K7 lattice-approximation argument gives an underlying isogeny J\^{}N to the actual J factor prime to any prescribed finite set, and an isomorphism of p-divisible group schemes after good reduction. No old theta-line/frame compatibility is inferred. Effective field/isogeny/line/frame costs, numerical Host/Sizes existence, same-point F nonzero, actual lift identities/arithmetic image, global period, specific032 and DP44 remain open; no RH or finite numerical charge. Current101/31 arrays and complete floors25/11 are unchanged.

The principal polarized abelian scheme over a number field retains the source O\_E action, actual Picard class, trace-one rigidified symmetric cubical line and Poincare biextension. The auxiliary field does not change the Host. The rational J-isotypic h1 projector gives the displayed CM-motive decomposition; after selected good reduction the prime-avoiding isogeny yields an actual p-divisible group-scheme isomorphism, hence the corresponding formal-group/Dieudonne isomorphism. It does not identify the selected theta line or source-normalized differential operation. Exact inherited contextual anchors: FGA232 Theorem3.1 pp148--149;FGA236 Corollaries1.7/2.3 pp229/232;SGA7 I VII2.7/2.9 pp169--174 and VIII7.4(b) pp310--311;SGA3 I4.7.3/4.7.3.1;SGA6 VI Theorem3.3;cote4 scan003 sections8.1.1--8.1.2 and scans005--006 sections8.4.1--8.4.3;cote54 scans005--007. No new lambda-ring/K-theory identity, l-adic weight bound or motivic L-function assertion. DP\_MULT Q3e/DP44/DP52/DP54--65 and THETA3 G1/A7 retain their effective costs/nonvanishing/arithmetic gates;CX\_GRH CX-ALG-024 and GCT\_PLETH T-052 remain open in their global scopes.

\subsection{245. AM-LIFT-085}\label{am-lift-085}

AM-LIFT-085 HF1--26 proves that, for an existing literal Host/Datum/Moduli/Coeff, actual CM seed Obj, original finite level A=1 and a given nonzero natural frame rule, a point in that same OE-isogeny class has a nonzero literal FrPt evaluation at the unchanged finite level. The isogeny u has degree N\^{}(2g), N prime to any prescribed finite set S containing the level primes, N=1 mod M, and u*L1=L0\^{}(N\^{}2) rigidified and cubical. Every output object and a nonzero frame descend over K1=K0(A0{[}N\^{}2{]}), with {[}K1:K0{]}\textless=N\^{}(8g\^{}2). The argument meets every analytic moduli component through an algebraic family based on its own target lattice. A weight-rule frame can be chosen by integer matrices with entries0..D+1 relative to any initial rational eigenframe. N itself, seed degree, initial-frame/theta-generator heights, models and actual line costs remain unevaluated. The given rule must itself be nonzero; no prescribed arithmetic-image nonvanishing is proved. No old pure LamObjGood or selected derivative identification, global period, specific032 pairing, DP44 or RH conclusion follows. This removes the qualitative same-rule same-point existence gate within the fixed CM isogeny class under the precise assumptions. Current101/31 error arrays and complete floors25/11 remain unchanged.

The map is an actual isogeny of abelian varieties over a number field, with OE action, finite level, principal polarization and the uniquely reconstructed trace-one cubical line. Its construction uses the rigidified relative Picard functor and normalized Poincare biextension: inherited FGA232 Theorem3.1 pp148--149, FGA236 Corollaries1.7/2.3 pp229/232; SGA7 I VII2.7/2.9 pp169--174 and VIII7.4(b) pp310--311. Actual eigenbundles and frame torsors retain SGA3 I4.7.3/4.7.3.1 and SGA6 VI Theorem3.3 as context. At protected primes the p-divisible consequence assumes extension to actual abelian models; cote4 scan003 sections8.1.1--8.1.2, scans005--006 sections8.4.1--8.4.3 and cote54 scans005--007 remain contextual. No new integral crystal, lambda-ring/K-theory, l-adic weight or CM-motivic L-function comparison is claimed. DP\_MULT Q3e/DP44/DP52/DP54--65 and THETA3 G1/A7 retain the quantitative and arithmetic-image gates; CX\_GRH CX-ALG-024 and GCT\_PLETH T-052 retain their distinct global comparison gates.

\subsection{246. AM-K-102}\label{am-k-102}

AM-K-102 PC1--27 bounds the entire actual DO107 IP3 support25 remainder by2\^{}-3550: all-II globally rank-two-unipotent scalar-cross, exactly25 total eligible places, no inactive paired cross zero,2\^{}120\textless N\_E D\_eta\textless2\^{}744 and100\textless=k\textless=164. The relative-norm phase count, positive-cone Gaussian/Poisson majorant and true transverse metric use one joint energy constraint; the complete original complex-volume tail covers every larger ideal-norm parameter. The original25 active weighted-label factors and143 inactive pointwise factors are retained without assuming complete-orbit constancy. All195 endpoint bounds are below-3557; the65-cell union costs3550 once. Actual DO108101/31 becomes102/31, every old overlapping charge retained, both complete sums below2\^{}-230 and floors25/11 until a separate exhaustive assembly. No signed positivity, prescribed period, selected arithmetic-image witness, effective admitted source cost, specific032, DP44 or RH claim follows.

The positive real norms are sections of the original relative-norm ideal C with C O\_E=L bar(L), and its full trace dual is Different\_F\^{}-1 C\^{}-1. The original Hermitian metric, paired source lattices, polarization, theta line, Picard/biextension data, selected coefficient and finite projectors remain fixed. Exact inherited contextual anchors are SGA3 I4.7.3/4.7.3.1; SGA6 VI Theorem3.3; FGA232 Theorem3.1 pp148--149; FGA236 Corollaries1.7/2.3 pp229/232; SGA7 I VII2.7/2.9 pp169--174 and VIII7.4(b) pp310--311; cote4 scan003 §§8.1.1--8.1.2 and scans005--006 §§8.4.1--8.4.3; cote54 scans005--007. These are retained framework locators, not new lambda-ring/K-theory, l-adic weight, CM-motive or L-function conclusions. DP\_MULT Q3e and THETA3 G1/A7 gain this finite subset removal; CX\_GRH CX-ALG-024 and GCT\_PLETH T-052 retain their distinct gates.

\subsection{247. R087: complete support25 exhaustion}\label{r087-complete-support25-exhaustion}

AM-R-087 J26.1--5 exhausts every original all-top support25 class at the unchanged H133 cutoff. The complete F/IC8 rank1..33 ranges force A=I; DO090/R080/R079 exhaust fixed S; DO099/DO095 give the precise cross rank2/3 alternatives; R082/R084 exhaust rank3; DO101 and DO102 remove nonunipotent rank2 and all non-all-II unipotent profiles; the exact DO103/104/105/107/K102 denominator, eligibility and coordinate partition exhausts the final all-II branch, including every equality boundary. The complete all-top floor is26 and the independent uniform floor remains11. Every102/31 chronological error entry is unchanged, both full sums stay below2\^{}-230 and no assembly charge is added. The actual sufficient signed gate is J26.2; global period, fully sized admissible source, selected arithmetic image, specific032/Read, actual DP44, Kummer-Hecke RH and the global oscillator comparison remain open.

This is an exhaustion of the original determinantal, fixed-submodule and isotropic-cross strata of Hermitian correspondences over Spec(O\_E), with the original paired incidence schemes and chosen complex realizations retained. The fractional denominator and relative-norm ideals remain actual invertible modules; no principality, globally integral moving frame, or inverse on the arithmetic theta image is inferred. The polarization, theta line, Picard data and biextensions are unchanged.

The exact contextual module and determinantal anchors are SGA3, ExposeI,4.7.3/4.7.3.1, and SGA6, ExposeVI, Theorem3.3. The Picard anchors are FGA, Bourbaki232, Theorem3.1, pp148--149, and Bourbaki236, Corollaries1.7/2.3, pp229/232. Biextensions retain SGA7 I, ExposeVII,2.7/2.9, pp169--174, and ExposeVIII,7.4(b), pp310--311. Formal groups, Dieudonne crystals and Cartier duality retain cote4 scan003, §§8.1.1--8.1.2, scans005--006, §§8.4.1--8.4.3, and cote54 scans005--007. These are inherited contextual placements, not newly proved comparison theorems.

No new lambda-ring/K-theory equality, l-adic sheaf or weight statement, CM-motive comparison, or motivic L-function theorem follows. DP\_MULT Q3e, DP44/DP52/DP54--65 and THETA3 G1/A7 gain precisely J26.1--2 as a smaller finite sufficient gate. CX\_GRH CX-ALG-024 and GCT\_PLETH T-052 retain their separate arithmetic gates. The known source theorems are used at their exact displayed locators, without rederiving them. No Lean build or archive audit is performed.

\subsection{248. AM-R-088}\label{am-r-088}

AM-R-088 GZ1--11 removes the actual fixed-S rank-one-unipotent support26 integral-witness branch, every reduced double-pole cell d\textless=22, and every reduced pattern with at least one nonidentity unmixed G reduction. The exact averaged-content and original affine-coordinate count retains every paired type/label factor and the actual selected trace. The full334-place nonempty weighted G-pattern sum is below2\^{}-5533, replacing only the former factor2. All27 complete endpoints and the witness give one523 charge. Only the reduced all334Gidentity pattern with d23..26 remains in this branch. Actual R087102/31 becomes103/31, every old overlap retained, both whole sums below2\^{}-230, and complete floors26/11 unchanged. No full support26 exhaustion, signed positivity, actual source, selected arithmetic content, prescribed period, specific032/Read, DP44 or RH conclusion follows.

The reductions are the actual integral Hermitian modules at unmixed finite points; content and coordinate ideals remain fractional invertible modules over Spec(O\_E) and Spec(O\_F), with the same paired CM action, polarization, theta line, Picard scheme and biextension. Carried exact anchors: SGA3 I4.7.3/4.7.3.1; SGA6 VI Theorem3.3; FGA232 Theorem3.1 pp148--149; FGA236 Corollaries1.7/2.3 pp229/232; SGA7 I VII2.7/2.9 pp169--174 and VIII7.4(b) pp310--311; cote4 scan003 sections8.1.1--8.1.2 and scans005--006 sections8.4.1--8.4.3; cote54 scans005--007. No new lambda-ring/K-theory, l-adic weight or CM-motive comparison is asserted. DP\_MULT Q3e, THETA3 G1/A7, CX\_GRH CX-ALG-024 and GCT\_PLETH T-052 retain their stated global gates.

\subsection{249. AM-LIFT-086}\label{am-lift-086}

AM-LIFT-086 IM1--27 constructs actual everywhere-good polarized models for the LIFT085 geometric CM seed and target over K=K0(A0{[}15N\^{}2{]}), with {[}K:K0{]}\textless=(15N\^{}2)\^{}(4g\^{}2). The full integral OE action, principal polarization, canonical trace-one cubical line, isogeny and adjoint extend over the same ring. Actual intrinsic tangent-line degrees change by an amount in{[}-log N,log N{]}, the full determinant degree by at most g log N, and the theta degree by at most g log N/2 under KeyCore. Target intrinsic lines have length constant1. If the actual seed lines have degree at least-b and all genuine rule/integral/cusp/source premises hold, the sufficient Step5 cost is C\_N=(D+gn)/2+D log r+(b+log N)min(D+gn,(3D+gn)/2). No effective bound on N, initial field or arbitrary-Host b is supplied. Host/Sizes, actual rule nonzero/content, prescribed period, specific032/Read, evaluated DP44 and RH remain open. No new finite error charge; actual103/31 arrays and complete floors26/11 remain unchanged.

The actual abelian schemes, integer-ring action, principal Picard/Poincare construction, normalized cubical line, finite-flat isogeny and its kernel live over Spec(O\_K). Away from N the displayed adjoint inverse gives an actual p-divisible-group and formal-group/Dieudonne comparison, without an operator or slope gain. Carried exact anchors: FGA232 Theorem3.1 pp148--149; FGA236 Corollaries1.7/2.3 pp229/232; SGA7 I VII2.7/2.9 pp169--174 and VIII7.4(b) pp310--311; SGA3 I4.7.3/4.7.3.1; SGA6 VI Theorem3.3; cote4 scan003 sections8.1.1--8.1.2 and scans005--006 sections8.4.1--8.4.3; cote54 scans005--007. No new lambda-ring/K-theory, l-adic weight or CM-motive identification. DP\_MULT Q3e/DP44/DP52/DP54--65, THETA3 G1/A7, CX\_GRH CX-ALG-024 and GCT\_PLETH T-052 retain their separate arithmetic comparisons.

\subsection{250. AM-K-103}\label{am-k-103}

AM-K-103 SC26.1--28 removes nine complete subsets of the actual all-II globally rank-two-unipotent scalar-cross support26 class. Its support-free full-Z proof, new eta window and uniqueness, exact paired ideals and inactive weights, full infinite ideal tail, joint energy and positive-norm-cone counts retain every original factor. All1,229 cells give one293 charge. The residual necessarily has v=q26, 2\^{}1116\textless N\_E D\_eta\textless2\^{}2108, b=j0 and81\textless=k\textless=164. Actual R088103/31 becomes104/31, all earlier overlapping charges retained, both entire sums below2\^{}-230 and complete floors26/11 unchanged. Other types/classes, signed positivity, actual source, selected arithmetic content, prescribed period, specific032/Read, DP44 and RH remain open.

The reductions are the actual integral Hermitian modules at unmixed finite points; content and coordinate ideals remain fractional invertible modules over Spec(O\_E) and Spec(O\_F), with the same paired CM action, polarization, theta line, Picard scheme and biextension. Carried exact anchors: SGA3 I4.7.3/4.7.3.1; SGA6 VI Theorem3.3; FGA232 Theorem3.1 pp148--149; FGA236 Corollaries1.7/2.3 pp229/232; SGA7 I VII2.7/2.9 pp169--174 and VIII7.4(b) pp310--311; cote4 scan003 sections8.1.1--8.1.2 and scans005--006 sections8.4.1--8.4.3; cote54 scans005--007. No new lambda-ring/K-theory, l-adic weight or CM-motive comparison is asserted. DP\_MULT Q3e, THETA3 G1/A7, CX\_GRH CX-ALG-024 and GCT\_PLETH T-052 retain their stated global gates.

\subsection{251. AM-DO-109}\label{am-do-109}

AM-DO-109 SP26.1--25 removes the entire actual fixed-S globally rank-two-unipotent support26 low-D\_G\textless=32 short-pivot class. The original plane reduction and paired alpha budget give2i+t\textless=5. The actual V-input eligibility and central-only label proof, true Gram metric, affine central difference lattice and full weighted four-state G sum are retained. The three complete p\textless=24,p25,p26 endpoints give one820 charge; the p26 branch proves eta uniqueness before summing actual labels. Together with already charged R079 high-D\_G and last-only pivots, this completes the whole fixed-S rank-two-unipotent support26 structural class. Actual K103104/31 becomes105/31, every overlap retained, both entire sums below2\^{}-230 and complete floors26/11 unchanged. Other support26 classes, signed positivity, actual source, selected arithmetic content, prescribed period, specific032/Read, DP44 and RH remain open.

The reductions are the actual integral Hermitian modules at unmixed finite points; content and coordinate ideals remain fractional invertible modules over Spec(O\_E) and Spec(O\_F), with the same paired CM action, polarization, theta line, Picard scheme and biextension. Carried exact anchors: SGA3 I4.7.3/4.7.3.1; SGA6 VI Theorem3.3; FGA232 Theorem3.1 pp148--149; FGA236 Corollaries1.7/2.3 pp229/232; SGA7 I VII2.7/2.9 pp169--174 and VIII7.4(b) pp310--311; cote4 scan003 sections8.1.1--8.1.2 and scans005--006 sections8.4.1--8.4.3; cote54 scans005--007. No new lambda-ring/K-theory, l-adic weight or CM-motive comparison is asserted. DP\_MULT Q3e, THETA3 G1/A7, CX\_GRH CX-ALG-024 and GCT\_PLETH T-052 retain their stated global gates.

\subsection{252. AM-K-104}\label{am-k-104}

AM-K-104 FV1--14 removes the complete new coordinate interval81\textless=k\textless=95 from the actual K103 all-II globally rank-two-unipotent scalar-cross support26 residual. The original full joint-volume lattice, actual central affine fibre and selected trace are retained. The full334-place four-state weighted G product is below2; the exact support binomial is below2\^{}101. Complete coupled-x monotonicity gives the actual lower denominator endpoint without a missing range. All15 cells, with worst-1128627/4096, give one274 charge. The all-II residual now has v=q26, 2\^{}1116\textless x\textless2\^{}2108, b=j0 and96\textless=k\textless=164. Actual DO109105/31 becomes106/31, all prior overlaps retained, both entire error sums below2\^{}-230 and complete floors26/11 unchanged. Other classes, signed positivity, actual source, selected arithmetic content, prescribed period, specific032/Read, DP44 and RH remain open.

The reductions are the actual integral Hermitian modules at unmixed finite points; content and coordinate ideals remain fractional invertible modules over Spec(O\_E) and Spec(O\_F), with the same paired CM action, polarization, theta line, Picard scheme and biextension. Carried exact anchors: SGA3 I4.7.3/4.7.3.1; SGA6 VI Theorem3.3; FGA232 Theorem3.1 pp148--149; FGA236 Corollaries1.7/2.3 pp229/232; SGA7 I VII2.7/2.9 pp169--174 and VIII7.4(b) pp310--311; cote4 scan003 sections8.1.1--8.1.2 and scans005--006 sections8.4.1--8.4.3; cote54 scans005--007. No new lambda-ring/K-theory, l-adic weight or CM-motive comparison is asserted. DP\_MULT Q3e, THETA3 G1/A7, CX\_GRH CX-ALG-024 and GCT\_PLETH T-052 retain their stated global gates.

\subsection{253. AM-R-089}\label{am-r-089}

AM-R-089 WN1--18 removes the complete actual fixed-S rank-one-unipotent support26 remainder: all334Gidentity and intrinsic double-pole sizes23..26. Its universal reciprocal-norm weighted eta/content lemma uses determinant-one rectangular rescaling and full dyadic shell packing, including every unit and ideal class. The original selected trace retains its actual eta norm before summation; both genuine affine coordinates and the whole actual support/marked/type/pair sum remain. The complete rational upper exponent is-146697349/4096\textless-35000, giving one35000 charge. Together with already charged R088 this completes the whole fixed-S rank-one-unipotent support26 structural class. Actual K104106/31 becomes107/31, every old overlapping error retained, both entire sums below2\^{}-230 and complete floors26/11 unchanged. Other support26 classes, signed positivity, actual source, selected arithmetic content, prescribed period, specific032/Read, DP44 and RH remain open.

The reductions are the actual integral Hermitian modules at unmixed finite points; content and coordinate ideals remain fractional invertible modules over Spec(O\_E) and Spec(O\_F), with the same paired CM action, polarization, theta line, Picard scheme and biextension. Carried exact anchors: SGA3 I4.7.3/4.7.3.1; SGA6 VI Theorem3.3; FGA232 Theorem3.1 pp148--149; FGA236 Corollaries1.7/2.3 pp229/232; SGA7 I VII2.7/2.9 pp169--174 and VIII7.4(b) pp310--311; cote4 scan003 sections8.1.1--8.1.2 and scans005--006 sections8.4.1--8.4.3; cote54 scans005--007. No new lambda-ring/K-theory, l-adic weight or CM-motive comparison is asserted. DP\_MULT Q3e, THETA3 G1/A7, CX\_GRH CX-ALG-024 and GCT\_PLETH T-052 retain their stated global gates.

\subsection{254. AM-DO-110}\label{am-do-110}

AM-DO-110 NT26.1--31 removes the complete actual non-all-II globally rank-two-unipotent scalar-cross support26 sector. The support-independent full-positive normal form, original paired cross/full-entry poles, 24-place eta uniqueness, actual one-label cancellation including central-only poles, full affine central lattice, common-G ideal and exact joint count are retained. Seven exhaustive branches and all1148 cells give one561 charge. Actual R089107/31 becomes108/31, every overlap retained, both complete sums below2\^{}-230 and full floors26/11 unchanged. The remaining all-II unipotent26 class is the K103/K104 dense window; other scalar-cross classes, signed positivity, actual source, arithmetic content, prescribed period, specific032/Read, DP44 and RH remain open.

The actual fractional coordinate and content ideals remain invertible modules over Spec(O\_E) and Spec(O\_F), with both CM components, original polarizations, Picard classes, theta lines, biextensions, metric and Haar normalization unchanged. Exact carried anchors: SGA3 I4.7.3/4.7.3.1; SGA6 VI Theorem3.3; FGA232 Theorem3.1 pp148--149; FGA236 Corollaries1.7/2.3 pp229/232; SGA7 I VII2.7/2.9 pp169--174 and VIII7.4(b) pp310--311; cote4 scan003 sections8.1.1--8.1.2 and scans005--006 sections8.4.1--8.4.3; cote54 scans005--007. No new lambda-ring/K-theory, l-adic weight or CM-motive comparison follows. DP\_MULT Q3e, THETA3 G1/A7, CX\_GRH CX-ALG-024 and GCT\_PLETH T-052 retain their stated global gates.

\subsection{255. AM-LIFT-087}\label{am-lift-087}

AM-LIFT-087 GS1--20 chooses an actual new general-Host Serre seed X=O\_L tensor\_\{O\_K7\} J and the actual pure block b\^{}-1 tensor\_\{O\_k\} C0, preserving the Host ideal, Datum, level, coefficient premises and R086 signature/J-content formulas. Every constructed intrinsic component eigenline has degree at least-b, with b=max(log D\_L+.5 log(2t),(sqrt(d)+1)/3+log D\_E+.5 log f). Literal integral b gives1 in b\^{}-1; full-order discriminant-cleared projectors and the exact Zarhin pullback isometry retain all integral models and source pi H norms. The new trace-polarization degree is D\_L/7\^{}(5t), distinct from R086. Under all retained actual rule, KeyCore, integrality/cusp/source premises, LIFT086 now has explicit seed b in C\_N=(D+gn)/2+D log r+(b+log N)min(D+gn,(3D+gn)/2). Effective N, initial field, actual Host/Sizes constructor, selected image, prescribed period, specific032/Read, DP44 and RH remain open. No finite charge; current108/31 arrays and complete floors26/11 unchanged.

The actual abelian schemes, integer-ring action, principal Picard/Poincare construction, normalized cubical line, finite-flat isogeny and its kernel live over Spec(O\_K). Away from N the displayed adjoint inverse gives an actual p-divisible-group and formal-group/Dieudonne comparison, without an operator or slope gain. Carried exact anchors: FGA232 Theorem3.1 pp148--149; FGA236 Corollaries1.7/2.3 pp229/232; SGA7 I VII2.7/2.9 pp169--174 and VIII7.4(b) pp310--311; SGA3 I4.7.3/4.7.3.1; SGA6 VI Theorem3.3; cote4 scan003 sections8.1.1--8.1.2 and scans005--006 sections8.4.1--8.4.3; cote54 scans005--007. No new lambda-ring/K-theory, l-adic weight or CM-motive identification. DP\_MULT Q3e/DP44/DP52/DP54--65, THETA3 G1/A7, CX\_GRH CX-ALG-024 and GCT\_PLETH T-052 retain their separate arithmetic comparisons.

\subsection{256. AM-R-090}\label{am-r-090}

AM-R-090 IS1--25 removes the entire actual full-positive globally rank-three scalar-cross support26 class with nonzero original S-cross, without a sparsity hypothesis. The complete eligible-set count at most24 costs6670; the complete all-k denominator tail at least2\^{}8192 costs11965. Literal inversion applies each containing-set bound once more. Outside them, the two eligible sets intersect in at least24 paired places, forcing equal null lines by the exact6395-bit denominator-cleared norm gap and contradicting global rank three. The full four-sector union is below2\^{}-6668. Actual DO110108/31 becomes109/31, every previous overlapping charge retained, both complete sums below2\^{}-230 and floors26/11 unchanged. The nonzero short-compression and dense all-II rank-two classes, signed positivity, actual source, selected arithmetic content, prescribed period, specific032/Read, DP44 and RH retain their gates.

The reductions are the actual integral Hermitian modules at unmixed finite points; content and coordinate ideals remain fractional invertible modules over Spec(O\_E) and Spec(O\_F), with the same paired CM action, polarization, theta line, Picard scheme and biextension. Carried exact anchors: SGA3 I4.7.3/4.7.3.1; SGA6 VI Theorem3.3; FGA232 Theorem3.1 pp148--149; FGA236 Corollaries1.7/2.3 pp229/232; SGA7 I VII2.7/2.9 pp169--174 and VIII7.4(b) pp310--311; cote4 scan003 sections8.1.1--8.1.2 and scans005--006 sections8.4.1--8.4.3; cote54 scans005--007. No new lambda-ring/K-theory, l-adic weight or CM-motive comparison is asserted. DP\_MULT Q3e, THETA3 G1/A7, CX\_GRH CX-ALG-024 and GCT\_PLETH T-052 retain their stated global gates.

\subsection{257. AM-DO-111}\label{am-do-111}

AM-DO-111 VC1--25 proves the actual positive transverse compression is unitary when both original scalar-cross directions lie in the retained plane. Its four original integral short witnesses have the sharper all-embedding bound2\^{}41. The complete simultaneous finite restoration, arbitrary-ideal row volume, relative norm fibre, original selected trace, full W0 column, source fibre and all paired type weights give exact exponent-87517717/4096\textless-21366 for the entire nonzero-short-compression support26 sector. The complementary integral norm argument gives full positive167 compressionI and a positive166 fixed space, hence global rank at most3; global unipotence is not inferred. Actual R090109/31 becomes110/31 with one21366 append, every old overlapping charge retained, both entire error sums below2\^{}-230 and complete floors26/11 unchanged. Whole support26 assembly remains separate; dense all-II rank-two, signed positivity, actual admissible source, selected arithmetic content, prescribed period, specific032/Read, DP44 and RH remain open.

The reductions are the actual integral Hermitian modules at unmixed finite points; content and coordinate ideals remain fractional invertible modules over Spec(O\_E) and Spec(O\_F), with the same paired CM action, polarization, theta line, Picard scheme and biextension. Carried exact anchors: SGA3 I4.7.3/4.7.3.1; SGA6 VI Theorem3.3; FGA232 Theorem3.1 pp148--149; FGA236 Corollaries1.7/2.3 pp229/232; SGA7 I VII2.7/2.9 pp169--174 and VIII7.4(b) pp310--311; cote4 scan003 sections8.1.1--8.1.2 and scans005--006 sections8.4.1--8.4.3; cote54 scans005--007. No new lambda-ring/K-theory, l-adic weight or CM-motive comparison is asserted. DP\_MULT Q3e, THETA3 G1/A7, CX\_GRH CX-ALG-024 and GCT\_PLETH T-052 retain their stated global gates.

\subsection{258. AM-LIFT-088}\label{am-lift-088}

AM-LIFT-088 SF1--16 supplies an explicit field for the entire newly chosen LIFT087 general-Host seed over the same given Coeff field B: {[}K\_seed:B{]}\textless=168\emph{23040}h(-d)\emph{M\^{}(4g\^{}2)\textless=3870720}d*M\^{}(4g\^{}2), g=f(p+q). The actual prescribed elliptic CM action descends over B(j) by scalar differential injectivity. The explicit J good-model radicals, elliptic15-torsion, actual projective Serre modules, equivariant principalization, canonical trace-one cubical line, original finite M-level and literal Obj/FrPt all lie over the stated field. Under every retained LIFT085 nonzero-rule/source hypothesis, {[}K\_out:B{]} is at most that bound timesN\^{}(8g\^{}2), with the existing intrinsic Step5 cost unchanged. EffectiveN, absolute degree of the givenB, field discriminants, coefficient/section heights, numerical Host/Sizes existence, selected arithmetic content, prescribed period, specific032/Read, evaluated DP44 and RH remain open. No finite error charge; current110/31 arrays and complete floors26/11 are unchanged.

The actual abelian schemes, integer-ring action, principal Picard/Poincare construction, normalized cubical line, finite-flat isogeny and its kernel live over Spec(O\_K). Away from N the displayed adjoint inverse gives an actual p-divisible-group and formal-group/Dieudonne comparison, without an operator or slope gain. Carried exact anchors: FGA232 Theorem3.1 pp148--149; FGA236 Corollaries1.7/2.3 pp229/232; SGA7 I VII2.7/2.9 pp169--174 and VIII7.4(b) pp310--311; SGA3 I4.7.3/4.7.3.1; SGA6 VI Theorem3.3; cote4 scan003 sections8.1.1--8.1.2 and scans005--006 sections8.4.1--8.4.3; cote54 scans005--007. No new lambda-ring/K-theory, l-adic weight or CM-motive identification. DP\_MULT Q3e/DP44/DP52/DP54--65, THETA3 G1/A7, CX\_GRH CX-ALG-024 and GCT\_PLETH T-052 retain their separate arithmetic comparisons. Precise primary locator clarification: Zarhin authorv5 Theorem1.1 is printedp2; section4 construction is pp4--6 and equations16--19 with equivariance are pp5--6. The compressed page range in the frozen paragraph is read with this exact split. The same-field Serre construction is Milne CM2020v0.10 Proposition7.33/Remark7.34/Proposition7.36 pp62--63. The initial elliptic degree uses Sutherland2023 Lecture20 Theorem20.12 pp5--6; its elementary class-number majorant specializes Milne ANT2020v3.08 Theorem4.3 printedp70 and its proof pp80--81.

\subsection{259. AM-R-091}\label{am-r-091}

AM-R-091 DR1--7 consolidates every checked support26 cut through DO111 into one actual-remainder theorem. The only possible support26 terms are all-type-II globally rank-two-unipotent scalar-cross matrices with full positive167 compression, nonzero original S cross, v=q26, b=j0, strict denominator window2\^{}1116\textless N\_E D\_eta\textless2\^{}2108 and96\textless=k\textless=164. The actual retained support26 set and actual higher-support remainder split the unchanged energy exactly, with the complete110/31 errors below2\^{}-230 and floors26/11. Every earlier cut and overlap charge is retained. No charge or independent new movement is assigned to this consolidation. DR5, actual fully sized source, selected arithmetic content, prescribed period, specific032/Read, evaluated DP44, Kummer-Hecke RH and global oscillator comparison remain open.

This is the original fixed-submodule/determinantal/isotropic-cross stratification over Spec(O\_E), with actual fractional ideals and paired incidence modules. Carried anchors: SGA3 I4.7.3/4.7.3.1; SGA6 VI Theorem3.3; FGA232 Theorem3.1 pp148--149; FGA236 Corollaries1.7/2.3 pp229/232; SGA7 I VII2.7/2.9 pp169--174 and VIII7.4(b) pp310--311; cote4 scan003 sections8.1.1--8.1.2, scans005--006 sections8.4.1--8.4.3, cote54 scans005--007. No new lambda-ring/K-theory, l-adic weight, formal-group or CM-motive identification. DP\_MULT Q3e/DP44/DP52/DP54--65 and THETA3 G1/A7 receive the consolidated finite gate; CX\_GRH CX-ALG-024 and GCT\_PLETH T-052 retain their arithmetic comparisons.

\subsection{260. AM-K-105}\label{am-k-105}

AM-K-105 CG1--25 removes all39 complete coordinate sectors96\textless=k\textless=134 from the actual all-II globally rank-two-unipotent support26 scalar-cross remainder. The scalar real-dimension336 Poisson bound, full coupled Gram determinant, original selected trace, full affine central covolume and every finite/source factor remain. Restoring(N\_FJ0)\^{}(-1) and its actual common-G correlation gives the full weighted pattern factor\textless2Q\^{}(-167); complete coupled-x monotonicity gives the exact lower denominator endpoint. The worst cell is-9009839/6144\textless-1466 and the preceding38 are below-2831, so one1465 charge covers the union. Actual DO111110/31 becomes111/31, every overlap retained and both complete sums below2\^{}-230; floors26/11 stay fixed. With R091, the only possible support26 branch now has135\textless=k\textless=164, the same v=q26,b=j0 and strict denominator window2\^{}1116\textless x\textless2\^{}2108. Signed positivity, actual source, selected arithmetic content, period, specific032/Read, DP44 and RH remain open.

The original fractional invertible modules, Hermitian form, involution, polarization, theta line, paired incidences and selected complex realization are unchanged. Carried exact anchors: SGA3 I4.7.3/4.7.3.1; SGA6 VI Theorem3.3; FGA232 Theorem3.1 pp148--149; FGA236 Corollaries1.7/2.3 pp229/232; SGA7 I VII2.7/2.9 pp169--174 and VIII7.4(b) pp310--311; cote4 scan003 sections8.1.1--8.1.2, scans005--006 sections8.4.1--8.4.3 and cote54 scans005--007. No new lambda-ring/K-theory, l-adic weight, formal-group or CM-motive comparison. DP\_MULT Q3e/DP44/DP52/DP54--65 and THETA3 G1/A7 gain the stated smaller finite gate; CX\_GRH CX-ALG-024 and GCT\_PLETH T-052 retain their arithmetic gates. Primary bibliographic precision: Stephen D. Miller and Noah Stephens-Davidowitz (2019), arXiv1802.05708v2, equations1.1/2.1, printedpp1--2/4. The frozen source manifest misspells the first given name as Steven; its S. D. proof citation and mathematics are unchanged.

\subsection{261. AM-LIFT-089}\label{am-lift-089}

AM-LIFT-089 HF89.1--18 sharpens the degree of exactly the same LIFT088 seed and output fields. For r=p+q,g=fr, the actual projective OE-lattice is free modulo each finite ideal, and the full defined action bounds its torsion Galois image. The literal m/bar(m) level halves are perfect Weil duals; Coeff.zeta\_mem already gives mu\_M inB. Thus Kgood(A{[}M{]})=Kgood(A{[}m{]}) and {[}Kseed:Kgood{]}\textless=M\^{}(f r\^{}2). The same target field Kseed(A{[}N\^{}2{]}) has relative degree\textless=N\^{}(4f r\^{}2), including ramifiedN andN1. Consequently {[}Kout:B{]}\textless=3870720\emph{d}M\^{}(f r\^{}2)*N\^{}(4f r\^{}2), under every actual LIFT085 nonzero-rule/source hypothesis. All actions, lines, level, Obj/FrPt and intrinsic costs are unchanged. EffectiveN, absoluteBdegree, discriminants, coefficient/frame heights, actual Host/Sizes existence, selected arithmetic content, period, specific032/Read, DP44 and RH remain open. No error charge; current111/31 arrays and complete floors26/11 are unchanged.

The finite-etale torsion modules carry the actual OE-action and perfect polarization pairing. The two literal idempotent halves are Cartier dual through the existing roots of unity; all descended abelian schemes and cubical Picard/Poincare data are those already constructed. Carried exact anchors: FGA232 Theorem3.1 pp148--149; FGA236 Corollaries1.7/2.3 pp229/232; SGA7 I VII2.7/2.9 pp169--174 and VIII7.4(b) pp310--311; SGA3 I4.7.3/4.7.3.1; SGA6 VI Theorem3.3; cote4 scan003 sections8.1.1--8.1.2 and scans005--006 sections8.4.1--8.4.3; cote54 scans005--007. No new formal-group slope, lambda-ring/K-theory identity, l-adic weight theorem or CM-motive comparison follows. DP\_MULT Q3e/DP44/DP52/DP54--65, THETA3 G1/A7, CX\_GRH CX-ALG-024 and GCT\_PLETH T-052 retain their separate arithmetic comparisons.

\subsection{262. AM-DO-112}\label{am-do-112}

AM-DO-112 EC1--22 removes both complete coordinate sectors k135 and k136 from the actual all-II globally rank-two-unipotent support26 scalar-cross remainder. The larger Gaussian cap59/(50*65536), with the full real dimension336 tail and scalar theta bound26, preserves the original row metric, affine central divisor, selected trace, every finite factor and complete coupled pattern/denominator monotonicities. Exact whole-sector exponents are-691859209/307200\textless-2252 and-283283209/307200\textless-922; their union is below2\^{}-921. Actual K105111/31 becomes112/31, every historical overlap retained and both complete sums below2\^{}-230; floors26/11 stay fixed. The only possible support26 branch now has137\textless=k\textless=164, v=q26,b=j0 and strict denominator window2\^{}1116\textless x\textless2\^{}2108. The positive k137 diagnostic is unclaimed. Signed positivity, actual source, selected arithmetic content, period, specific032/Read, actual DP44 and RH remain open.

The original fractional invertible modules, Hermitian form, involution, polarization, theta line, paired incidences and selected complex realization are unchanged. Carried exact anchors: SGA3 I4.7.3/4.7.3.1; SGA6 VI Theorem3.3; FGA232 Theorem3.1 pp148--149; FGA236 Corollaries1.7/2.3 pp229/232; SGA7 I VII2.7/2.9 pp169--174 and VIII7.4(b) pp310--311; cote4 scan003 sections8.1.1--8.1.2, scans005--006 sections8.4.1--8.4.3 and cote54 scans005--007. No new lambda-ring/K-theory, l-adic weight, formal-group or CM-motive comparison. DP\_MULT Q3e/DP44/DP52/DP54--65 and THETA3 G1/A7 gain the stated smaller finite gate; CX\_GRH CX-ALG-024 and GCT\_PLETH T-052 retain their arithmetic gates. Primary Gaussian/Poisson context remains Miller, S. D., and Stephens-Davidowitz, N. (2019), arXiv1802.05708v2, equations1.1/2.1, printedpp1--2/4; the new-cap specialization is proved in full here.

\subsection{263. AM-LIFT-090}\label{am-lift-090}

AM-LIFT-090 SG1--21 sharply bounds the same actual LIFT087 seeds and LIFT088 fields through the two original generating CM varieties J and prescribed C0. Integral projective split maps and their duals put every seed n-torsion in Kgood(J{[}n{]},C0{[}n{]}) at every n, with no prime-to-index restriction; full maximal-order rank-one actions bound that field by n\^{}8. Thus the same seed and output fields cost M\^{}8 and N\^{}16, and a finite list of stages sharing J,C0,Kgood costs one M\^{}8 and one lcm(N\_j)\^{}16. The same C0{[}15{]} field has degree\textless=225 before good reduction, sharpening the same ground factor to37800\emph{h(-d)\textless=37800}d.~The complete relative bound is37800\emph{d}M\^{}8*lcm(N\_j)\^{}16. Retain the smaller LIFT089 bounds where applicable. The actual models, actions, polarizations, canonical line, finite level, Obj/FrPt and intrinsic cost remain unchanged. EffectiveN, absoluteBdegree, discriminants, coefficient heights, all nonzero-rule/source premises, selected arithmetic content, prescribed period, specific032/Read, actual DP44 and RH remain open. No error charge; current112/31 arrays and complete floors26/11 remain.

The finite-etale torsion modules carry the actual OE-action and perfect polarization pairing. The two literal idempotent halves are Cartier dual through the existing roots of unity; all descended abelian schemes and cubical Picard/Poincare data are those already constructed. Carried exact anchors: FGA232 Theorem3.1 pp148--149; FGA236 Corollaries1.7/2.3 pp229/232; SGA7 I VII2.7/2.9 pp169--174 and VIII7.4(b) pp310--311; SGA3 I4.7.3/4.7.3.1; SGA6 VI Theorem3.3; cote4 scan003 sections8.1.1--8.1.2 and scans005--006 sections8.4.1--8.4.3; cote54 scans005--007. No new formal-group slope, lambda-ring/K-theory identity, l-adic weight theorem or CM-motive comparison follows. DP\_MULT Q3e/DP44/DP52/DP54--65, THETA3 G1/A7, CX\_GRH CX-ALG-024 and GCT\_PLETH T-052 retain their separate arithmetic comparisons. The new split inclusions and their duals are integral homomorphisms of the existing abelian schemes; finite-etale geometric torsion fields are bounded through the two original rank-one maximal-order Galois representations. No new CM reciprocity or motive theorem is required.

\subsection{264. AM-K-106}\label{am-k-106}

AM-K-106 NT1--14 removes the complete coordinate sector k137 from the actual all-II globally rank-two-unipotent support26 scalar-cross remainder. The exact universal radial-profile factor carries through the original rank-two near-point trace and yields630 bits. Every nonzero original S coordinate lies in the same fractional ideal; the proved product-of-sums inequality gives k\^{}-84 and596 further bits at137. Applied before the full DO112 Gaussian cardinality enlargement, these constant factors preserve every actual metric, finite weight, central affine covolume and complete pattern/denominator count. The exact whole-sector exponent is-251641609/307200\textless-819. Actual DO112112/31 becomes113/31, all previous overlapping charges retained, both complete sums below2\^{}-230 and floors26/11 unchanged. The only possible support26 branch now has138\textless=k\textless=164 with the same v=q26,b=j0 and2\^{}1116\textless x\textless2\^{}2108. Signed positivity, selected arithmetic content, prescribed period, full effective source, specific032/Read, actual DP44 and RH remain open.

The original fractional invertible modules, Hermitian form, involution, polarization, theta line, paired incidences and selected complex realization are unchanged. Carried exact anchors: SGA3 I4.7.3/4.7.3.1; SGA6 VI Theorem3.3; FGA232 Theorem3.1 pp148--149; FGA236 Corollaries1.7/2.3 pp229/232; SGA7 I VII2.7/2.9 pp169--174 and VIII7.4(b) pp310--311; cote4 scan003 sections8.1.1--8.1.2, scans005--006 sections8.4.1--8.4.3 and cote54 scans005--007. No new lambda-ring/K-theory, l-adic weight, formal-group or CM-motive comparison. DP\_MULT Q3e/DP44/DP52/DP54--65 and THETA3 G1/A7 gain the stated smaller finite gate; CX\_GRH CX-ALG-024 and GCT\_PLETH T-052 retain their arithmetic gates. Primary Gaussian/Poisson context remains Miller, S. D., and Stephens-Davidowitz, N. (2019), arXiv1802.05708v2, equations1.1/2.1, printedpp1--2/4; the new-cap specialization is proved in full here. The sharper trace is the same genuine vector-valued selected normalized trace and original near-point argument; no operator-norm or rank-one transport substitute is used. Nakahama(2016), Journal of Lie Theory26(4),927--990, equation5.7 and Theorem5.1 printed960--961, supplies the carried representation convention.

\subsection{265. AM-LIFT-091}\label{am-lift-091}

AM-LIFT-091 RD1--18 bounds root discriminants of the same actual LIFT088--090 finite torsion fields. For a finite Galois extension of actual degree D, unramified outside an integral ideal S, rd(L)/rd(K)\textless=D\emph{Norm(rad S)\^{}(1/{[}K:Q{]}); the complete base-normalized local different and actual-degree wild contribution are proved. Thus the simultaneous seed satisfies rd(KSigma)\textless=rd(Kgood)}rad(M)\emph{M\^{}8, and the same actual output satisfies rd(KSigma,out)\textless=rd(Kgood)}rad(M L\_N)\emph{lcm(M,L\_N²)\^{}8\textless=rd(Kgood)}M\^{}9\emph{L\_N\^{}17. Its relative increase over KSigma is at most rad(L\_N)}L\_N\^{}16. The literal N\_j=1 mod M congruence is retained, so no false level/output-overlap saving is claimed. Full all-object descent and the37800\emph{d}M\^{}8*L\_N\^{}16 relative degree bound are unchanged. Initial rd(Kgood), absolute B degree/discriminant, effective N\_j, heights, all nonzero-rule/source conditions, selected content, prescribed period, specific032/Read and evaluated DP44 remain open. No numerical error charge;113/31 full arrays and floors26/11 remain.

The actual everywhere-good abelian schemes and finite-etale prime-to-order torsion give finite Galois covers outside the stated torsion primes. The different and trace pairing quantify the same finite base changes. Picard/Poincare, normalized cubical line, full actions and finite-flat quotients are unchanged. Carried anchors: FGA232 Theorem3.1 pp148--149; FGA236 Corollaries1.7/2.3 pp229/232; SGA7 I VII2.7/2.9 pp169--174 and VIII7.4(b) pp310--311; SGA3 I4.7.3/4.7.3.1; SGA6 VI Theorem3.3; cote4 scan003 sections8.1.1--8.1.2 and scans005--006 sections8.4.1--8.4.3; cote54 scans005--007. No new formal-group slope, lambda-ring/K-theory, l-adic weight or CM-motive comparison. DP\_MULT Q3e/DP44/DP52/DP54--65 receives a relative field-discriminant reduction; THETA3 G1/A7, CX\_GRH CX-ALG-024 and GCT\_PLETH T-052 retain their arithmetic comparison gates.

\subsection{266. AM-R-092}\label{am-r-092}

AM-R-092 ID1--18 bounds the initial root discriminant for the exact SF5 seed C0 and radical/torsion tower of LIFT088, over the actual source coefficient field B. The maximal-order class-field step B(j)/B is unramified; the actual radical tower has relative root-discriminant factor at most76832. Every CM conjugate satisfies \textbar sigma(j)\textbar\textless32 exp(pi sqrt(d)), and the actual C0{[}15{]} extension is treated at its true degree before applying the bound225. Consequently rd(Kgood)\textless2\^{}50 exp(2 pi sqrt(d)) rd(B). Composed with LIFT091, rd(KSigma,out)\textless2\^{}50 exp(2 pi sqrt(d)) rd(B) rad(M L\_N) lcm(M,L\_N\^{}2)\^{}8\textless=2\^{}50 exp(2 pi sqrt(d)) rd(B) M\^{}9 L\_N\^{}17. The stated log absolute discriminant bound retains the exact relative degree37800\emph{d}M\^{}8*L\_N\^{}16 and actual {[}B:Q{]},rd(B). This removes the separate initial-good-field discriminant premise for this fixed seed, not for arbitrary twists. Absolute B degree/discriminant, effective N\_j, heights, Host/Sizes existence, selected nonzero arithmetic rule/content, period, specific032/Read and evaluated DP44 remain open. No new finite-error charge;113/31 full arrays and complete floors26/11 unchanged.

The exact CM elliptic curve, its maximal-order action and the fixed Jacobian radical model give the same actual abelian schemes after the established good-reduction base change. Unramified class-field base change, integral Weierstrass smoothness, finite-etale prime-to-order torsion and the different/trace pairing control the finite covers and root discriminants. Picard/Poincare data, normalized cubical line, quotient/action/level/FrPt descent and integral maps remain those already checked. Carried anchors: FGA232 Theorem3.1 pp148--149; FGA236 Corollaries1.7/2.3 pp229/232; SGA7 I VII2.7/2.9 pp169--174 and VIII7.4(b) pp310--311; SGA3 I4.7.3/4.7.3.1; SGA6 VI Theorem3.3; cote4 scan003 sections8.1.1--8.1.2 and scans005--006 sections8.4.1--8.4.3; cote54 scans005--007. No new lambda-ring/K-theory identity, crystalline slope, l-adic weight theorem, motivic comparison or Hecke RH statement. DP\_MULT Q3e/DP44/DP52/DP54--65 receives this effective initial-field cost reduction; THETA3 G1/A7, CX\_GRH CX-ALG-024 and GCT\_PLETH T-052 retain their actual arithmetic comparison gates.

\subsection{267. AM-LIFT-092}\label{am-lift-092}

AM-LIFT-092 CD1--19 strengthens the root-discriminant cost of the same actual LIFT088--090 fields. For the two full-maximal-order CM progenitors J,C0 over the actual everywhere-good Kgood, every actual F inside Kgood(J{[}n{]},C0{[}n{]}) is abelian and has conductor dividing n O\_Kgood and discriminant dividing (n O\_Kgood)\^{}(D-1), D={[}F:Kgood{]}. Thus rd(F)/rd(Kgood)\textless=n\^{}(1-1/D)\textless=n, including all dyadic/CM-ramified primes and D=1. The same simultaneous output satisfies rd(KSigma,out)\textless=rd(Kgood)\emph{lcm(M,L\_N\^{}2)=rd(Kgood)}M\emph{L\_N\^{}2 under the actual N\_j=1 mod M; its relative output increase over KSigma is at most L\_N\^{}2. Combining the separate canonical R092 bound gives rd(KSigma,out)\textless2\^{}50 exp(2 pi sqrt(d))}rd(B)\emph{M}L\_N\^{}2. LIFT090 relative degree37800\emph{d}M\^{}8*L\_N\^{}16 is unchanged; the sufficient absolute-discriminant logarithm is this degree times {[}B:Q{]} times (log rd(B)+50 log2+2 pi sqrt(d)+log M+2 log L\_N). All original fields, varieties, polarizations, line, level, quotient, action, source points and frames are retained. Actual B degree/discriminant, effective N\_j, metric/height costs, Host/Sizes, selected nonzero arithmetic content, period, specific032/Read and evaluated DP44 remain open. No finite-error charge;113/31 full arrays and complete floors26/11 unchanged.

The actual full CM Tate lattices carry integral maximal-order actions at every prime. CM reciprocity and good reduction make their same torsion covers factor through the finite ray-class modulus of the torsion order; the character-conductor product measures the different of each actual abelian subcover. The Picard/Poincare and origin-rigidified cubical line, finite-flat quotients, full actions, finite level and framed points are unchanged. Carried anchors: FGA232 Theorem3.1 pp148--149; FGA236 Corollaries1.7/2.3 pp229/232; SGA7 I VII2.7/2.9 pp169--174 and VIII7.4(b) pp310--311; SGA3 I4.7.3/4.7.3.1; SGA6 VI Theorem3.3; cote4 scan003 sections8.1.1--8.1.2 and scans005--006 sections8.4.1--8.4.3; cote54 scans005--007. No new crystalline slope, lambda-ring/K-theory identity, l-adic weight theorem, CM-motive comparison or Hecke RH result. DP\_MULT Q3e/DP44/DP52/DP54--65 receives a sharper finite-field ramification cost; THETA3 G1/A7, CX\_GRH CX-ALG-024 and GCT\_PLETH T-052 retain their separate arithmetic comparisons.

\subsection{268. AM-K-107}\label{am-k-107}

AM-K-107 SV1--36 removes the entire k138 cell of the actual all-II globally rank-two-unipotent scalar-cross support26 remainder. The genuine selected density yields a full tail-inclusive squared coefficient below2\^{}-933 and a single actual near-point product below2. Together with the exact radial factor, all-coordinate norm and original saturated transverse lattice, the complete count is bounded by2\^{}(-50428781/204800)\textless2\^{}-246. The original scalar theta is in real dimension336, with full21\^{}(k+3) and all original central, support, row, finite and pattern factors. The remaining branch has139\textless=k\textless=164 and2\^{}(71443/64)\textless x\textless2\^{}2108, v=q26,b=j0. Exactly one246 append gives114/31 ordered errors with all old overlaps retained, both full sums below2\^{}-230, unique minimum231 and next minimum246; complete floors26/11 remain. Mandatory corrected boundary conventions: SV6 is strict for h\textgreater=1 and equality at h0; SV11 quotient is for r\textgreater0 and the bound at r0 is1. The frozen proof is read with the explicit boundary addendum reproduced below. Signed positivity, selected arithmetic content, period, effective source, specific032/Read, actual DP44 and RH remain open.

The original fractional invertible modules, Hermitian form, involution, polarization, theta line, paired incidences and selected complex realization are unchanged. Carried exact anchors: SGA3 I4.7.3/4.7.3.1; SGA6 VI Theorem3.3; FGA232 Theorem3.1 pp148--149; FGA236 Corollaries1.7/2.3 pp229/232; SGA7 I VII2.7/2.9 pp169--174 and VIII7.4(b) pp310--311; cote4 scan003 sections8.1.1--8.1.2, scans005--006 sections8.4.1--8.4.3 and cote54 scans005--007. No new lambda-ring/K-theory, l-adic weight, formal-group or CM-motive comparison. DP\_MULT Q3e/DP44/DP52/DP54--65 and THETA3 G1/A7 gain the stated smaller finite gate; CX\_GRH CX-ALG-024 and GCT\_PLETH T-052 retain their arithmetic gates. Primary Gaussian/Poisson context remains Miller, S. D., and Stephens-Davidowitz, N. (2019), arXiv1802.05708v2, equations1.1/2.1, printedpp1--2/4; the new-cap specialization is proved in full here. The sharper trace is the same genuine vector-valued selected normalized trace and original near-point argument; no operator-norm or rank-one transport substitute is used. Nakahama(2016), Journal of Lie Theory26(4),927--990, equation5.7 and Theorem5.1 printed960--961, supplies the carried representation convention. The new bound uses the same selected Hilbert metric and integral saturated kernel of the original primitive packet row; the projection norm U and central affine difference module remain. The exact tail is absorbed before counting.

\subsubsection{Mandatory checked boundary qualification}\label{mandatory-checked-boundary-qualification-1}

\section{AM-K-107 --- explicit boundary precision for the preserved frozen proof}\label{am-k-107-explicit-boundary-precision-for-the-preserved-frozen-proof-1}

Author: /root/kummer\_extension. Independent observations: /root/lift\_accounting.

The preserved subject is sharpened\_radial\_variance\_support\_twenty\_six.md, SHA256 3ad12d483d85b2f0e20917d5f78c0b5a31f4e9a726ea1fab18c746c4788dda2a. The parent requested that these frozen bytes remain unchanged. Read that subject with the following two precise qualifications.

\begin{enumerate}
\def\labelenumi{\arabic{enumi}.}
\tightlist
\item
  In SV6 the final strict inequality holds for integer h\textgreater=1. At h=0 both sides are exactly1. Equivalently the final displayed sign is \textless= for h\textgreater=0, and is strict for h\textgreater=1. Every use in SV7 has h=2\^{}19, so its strict probability bound and all subsequent constants are unchanged.
\item
  The quotient in SV11 is stated for r(g)\textgreater0. At r(g)=0 retain the actual normalized bound \textbar c\_W(g)\textbar\textless=1. The nonidentity rank-two application SV14--16 has positive alpha and hence positive Cartan radius; no division by zero occurs in that application.
\end{enumerate}

These are boundary qualifications, not a new estimate or sector charge. The full lower-tail absorption, saturated original lattice, theta21, complete exponent -50428781/204800 and charge246 are unchanged. The mathematical subject for the whole author/peer/root verdict is the frozen proof together with this explicit qualification. The script remains unchanged, and its final6,118 exact checks already include the h=0 factor identity; its tail calculation uses h=2\^{}19. This note does not amend a canonical file or error array.

\subsection{269. AM-DO-113}\label{am-do-113}

AM-DO-113 LS1--4 proves local sufficiency of the retained simple cross and double central pole allowances in the original paired unipotent normal form, conditional on the actually admitted cancellation label and the unitary relation. Both cancellation factors are divisible by the uniformizer, all three conjugated terms are integral separately, and nilpotence gives an integral inverse. Thus reimposing this same active-pencil preservation alone supplies no additional local linear or quadratic congruence beyond those allowances and the unitary relation. The existing pending row is closed in this local scope. No arbitrary local parameter is asserted to globalize, satisfy shortness/G/total support, or attain the enlarged count. No numerical charge or main finite-sector removal; complete114/31 arrays and floors26/11 are unchanged.

The proof computes the integral stabilizer of the original paired lattices over the selected split local ring, retaining the Hermitian form, involution, admitted pencil label, unipotent group coordinates and polarization. No positive residue-field metric is assumed. It gives no new lambda-ring/K-theory identity, l-adic weight theorem, CM-motive or oscillator comparison. Carried exact anchors: SGA3 I4.7.3/4.7.3.1; SGA6 VI Theorem3.3; FGA232 Theorem3.1 pp148--149; FGA236 Corollaries1.7/2.3 pp229/232; SGA7 I VII2.7/2.9 pp169--174 and VIII7.4(b) pp310--311; cote4 scan003 sections8.1.1--8.1.2 and scans005--006 sections8.4.1--8.4.3; cote54 scans005--007. DP\_MULT Q3e/DP44 and THETA3 G1/A7 gain a checked boundary for this specific congruence-count route, not a completed source/period construction; CX\_GRH CX-ALG-024 and GCT\_PLETH T-052 retain their separate arithmetic gates.

\subsection{270. AM-DO-114}\label{am-do-114}

AM-DO-114 SP1-27 proves a uniform scalar theta below3/2 in real dimension336 for A\textless=1/43500 by a complete finite spherical-cap and nonzero-shell argument. Applied to the original saturated transverse Gram, actual selected trace, full affine central fibre and every finite pattern, the whole139..141 union is below2\^{}-354. The necessary actual support26 all-II rank-two-unipotent range becomes142\textless=k\textless=164 with2\^{}(71443/64)\textless x\textless2\^{}2108,v=q26,b=j0. Exactly one354 append gives115/31 complete ordered errors, both full sums below2\^{}-230, with every historical overlap retained and complete floors26/11 unchanged. No support27, signed positivity, selected arithmetic content, period, effective source, specific032/Read, evaluated DP44 or RH conclusion.

The original invertible fractional modules, Hermitian form, involution, paired incidences, polarization, theta line and genuine selected complex realization remain fixed. The exact dimension336 half-trace dual is2 D\_E\^{}-1 conjugate(L)\^{}-1; the saturated original transverse Gram and central affine difference module are retained. The finite cap proof is given in full, with classical context Rankin(1955), Theorem1(iv), printed139 and proof printed142, doi10.1017/S2040618500033219. Gaussian convention: Miller and Stephens-Davidowitz(2019), arXiv1802.05708v2,equations1.1/2.1, carried directly checked scope. Contextual exact anchors: SGA3 I4.7.3/4.7.3.1; SGA6 VI Theorem3.3; FGA232 Theorem3.1pp148-149; FGA236 Corollaries1.7/2.3pp229/232; SGA7 I VII2.7/2.9pp169-174 and VIII7.4(b)pp310-311; cote4 scan003 sections8.1.1-8.1.2/scans005-006 sections8.4.1-8.4.3, cote54 scans005-007. No new crystal, lambda-ring/K-theory, l-adic weight or CM-motive comparison. DP\_MULT Q3e/DP44/DP52/DP54-65 and THETA3 G1/A7 gain only this finite reduction; CX\_GRH CX-ALG-024 and GCT\_PLETH T-052 retain their global comparison gates.

\subsection{271. AM-K-108}\label{am-k-108}

AM-K-108 AE1-30 retains the full actual Cartesian conditional angular density, cancels its D ratio in detailed balance and inserts exp(-2U\textbar\textbar r\textbar\textbar\^{}2) in the original lattice sum. With the separately checked DO114 scalar theta, original transverse Gram, central affine fibre, common ideal correlation and all finite factors, the complete142..159 union is below2\^{}-585. The necessary support26 all-II rank-two-unipotent remainder becomes160\textless=k\textless=164, with2\^{}(71443/64)\textless x\textless2\^{}2108,v=q26,b=j0. Exactly one585 charge gives116/31 ordered errors, both full sums below2\^{}-230; floors26/11 remain. No160+, signed positivity, arithmetic-content nonvanishing, period, effective source, specific032/Read, evaluated DP44 or RH conclusion.

The original paired invertible modules, Hermitian metric, saturated transverse basis, theta line and actual selected complex realization remain fixed. The new result is an analytic density and energy estimate on that realization, not a new algebraic correspondence. Nakahama(2016), Journal of Lie Theory26(4),927-990, equation5.7 and Theorem5.1, printed960-961, supplies the carried actual invariant selected norm; the named passage was directly reread. The finite cap proof of DO114 is retained with its Rankin(1955), Theorem1(iv), pp139,142 context. Exact contextual anchors remain SGA3 I4.7.3/4.7.3.1; SGA6 VI Theorem3.3; FGA232 Theorem3.1pp148-149; FGA236 Corollaries1.7/2.3pp229/232; SGA7 I VII2.7/2.9pp169-174 and VIII7.4(b)pp310-311; cote4 scan003 sections8.1.1-8.1.2/scans005-006 sections8.4.1-8.4.3 and cote54 scans005-007. No new crystal, lambda-ring/K-theory, l-adic weight, CM-motive or global period comparison follows. DP\_MULT Q3e/DP44/DP52/DP54-65 and THETA3 G1/A7 gain only this finite gate reduction; CX\_GRH CX-ALG-024 and GCT\_PLETH T-052 retain their global gates.

\subsection{272. AM-R-093}\label{am-r-093}

AM-R-093 DC1-19, on the actual d=7 branch, proves one finite-prime-unramified quadratic character epsilon with N\_E7/k(rho\_J)=epsilon\emph{chi\_cyc}rho\_C\_plus at every Tate prime. The integral rank-three norm at all primes, including2/7, gives K(J{[}n{]},C0{[}n{]}) contained in K(J{[}n{]})K\_epsilon, with relative degree at most2, simultaneously for every n.~Same actual projective/dual seeds and source congruences then give {[}Kgood:B{]}\textless=16128 and {[}K\_Sigma,out:B{]}\textless=32256\emph{M\^{}3}L\_N\^{}(129/16). The norm/type computation and arithmetic Artin convention are explicit; epsilon is neither evaluated nor assumed trivial. Root-discriminant bounds remain R092/LIFT092 on the same fields. Absolute coefficient-field degree/discriminant, effectiveN, Host/Sizes, selected arithmetic nonzero content, period, Read/specific032, evaluatedDP44 and RH remain open. No finite-error charge;116/31 and floors26/11 unchanged.

The actual full-order l-adic realizations of J and C0 and their principal Weil pairings yield an integral determinant-character identity in the k-component of the CM torus, up to one unramified quadratic Artin character. Coefficient conjugation only orients the existing C0 character; no curve, model, action, ideal class, seed or source coefficient field is replaced. The identity is not asserted to arise from an algebraic correspondence or isomorphism of CM motives. Conrad, Main theorem of complex multiplication, Theorem3.2, Remarks3.3-3.4 and Theorem3.5(1), printedpp8-10, with Remark3.6pp10-11; Lombardo(2017), Definitions2.7-2.9,p476; Milne CFTv4.03,V Theorem1.7pp150-151 and Propositions4.6-4.7pp172-174 are directly checked in their indicated scopes. Carried geometric anchors: FGA232 Theorem3.1pp148-149; FGA236 Corollaries1.7/2.3pp229/232; SGA7 I VII2.7/2.9pp169-174 and VIII7.4(b)pp310-311; SGA3 I4.7.3/4.7.3.1; SGA6 VI Theorem3.3; cote4 scan003 sections8.1.1-8.1.2/scans005-006 sections8.4.1-8.4.3, cote54 scans005-007. No new crystalline slope, lambda-ring/K-theory or weight theorem follows. DP\_MULT Q3e and DP44/DP52/DP54-65 gain the d7 same-field cost reduction; THETA3 G1/A7, CX\_GRH CX-ALG-024 and GCT\_PLETH T-052 retain their global comparison gates.

\subsection{273. AM-LIFT-093}\label{am-lift-093}

AM-LIFT-093 PA1-27 refines the general-d degree for the exact same Kgood, J,C0 and original source seeds. Principal Weil pairings impose one shared scalar multiplier. Exact all-prime norm counts, including2/7/d and the overlappingd7 case, give {[}F(J{[}n{]},C0{[}n{]}):F{]}\textless=U(n)=n\^{}5\emph{product\_(p\textbar n,p=6 mod7)(1+1/p)\^{}3 for every finiteF/Kgood, and U(N\^{}2)\textless=N\^{}(81/8). The actual source level has multiplier1 and boundV(M)=M\^{}4}product\_(p\textbar M)F\_J(p)(1-1/p)\textless2M\^{}4; no universal g mod7 condition is inserted. The existing C0{[}15{]} extension costs at most192, so {[}Kgood:B{]}\textless=32256h(-d)\textless=32256d and the actual common output satisfies {[}Kout\_Sigma:B{]}\textless=32256d V(M)U(L\_N\^{}2)\textless64512d M\^{}4 L\_N\^{}(81/8). R093 retains its independent strongerd7 bound. Same fields, all objects and intrinsic Step5 costs remain; effectiveN, absoluteB, Host/Sizes, selected nonzero arithmetic content, period, Read/specific032, evaluatedDP44 and RH are open. No finite-error charge;116/31 and floors26/11 unchanged.

The actual principal Picard/Poincare pairings and full maximal-order scalar actions constrain the l-adic and finite-etale torsion realizations of the same two CM progenitors. The new count is an image-container estimate, not an image-surjectivity or motivic-comparison theorem. Literal source precision: in the frozen PushInvertible.lean, Coeff.zeta\_mem is line44, not the isolated45 printed in the frozen author body; the fullCoeff40-46 locator was already correct. Read that occurrence with the immutable locator-precision addendum. No source or mathematical premise changes. Fresh primary anchors: Milne AVv2.00(2008),I13 Proposition13.2pp57-58,I14p61,III6.6p105 and Summary6.11pp107-108; Milne ANT v3.08(2020),Proposition6.2pp96-98,Theorem6.4p99,Example8.18p143; Milne CMv0.10(2020),Corollary1.40p21,Example2.9pp25-26. Carried geometric anchors: FGA232 Theorem3.1pp148-149; FGA236 Corollaries1.7/2.3pp229/232; SGA7 I VII2.7/2.9pp169-174 and VIII7.4(b)pp310-311; SGA3 I4.7.3/4.7.3.1; SGA6 VI Theorem3.3; cote4 scan003 sections8.1.1-8.1.2/scans005-006 sections8.4.1-8.4.3 and cote54 scans005-007. There is no new crystalline slope, lambda-ring/K-theory identity, weight or CM-motive map. DP\_MULT Q3e/DP44/DP52/DP54-65 gain a smaller same-field cost; THETA3 G1/A7, CX\_GRH CX-ALG-024 and GCT\_PLETH T-052 retain their global gates.

\subsection{274. AM-K-109}\label{am-k-109}

AM-K-109 TG1-19 retains the full all-A scalar theta polynomial and the original transverse Gram G\textgreater2\^{}36 I, with true determinantU. Using the entire actual K108 energy, it removes the complete k160 support26 cell with exact exponent-163509/512\textless-319. Original module membership, central affine fibre, common ideal correlation, all finite factors and both q/x monotonicities are retained. The necessary actual all-II rank-two-unipotent remainder becomes161\textless=k\textless=164 with2\^{}(71443/64)\textless x\textless2\^{}2108,v=q26,b=j0. Exactly one319 append gives117/31 complete ordered errors, both entire sums below2\^{}-230; complete floors26/11 remain. No161+, signed positivity, selected arithmetic content, effective source, period, specific032/Read, evaluatedDP44 or RH conclusion.

The original paired fractional modules and exact integral saturated transverse lattice retain their true Hermitian Gram and projection denominator. The stronger real dual-metric inequality and Gaussian count act on the same actual selected complex realization and theta line. No rational orthogonal frame is declared integral. The previous exact Nakahama(2016)equation5.7/Theorem5.1pp960-961, Gaussian Fourier convention of Miller and Stephens-Davidowitz(2019)equations1.1/2.1, and Rankin(1955)Theorem1(iv) context remain in the directly checked scopes of K108/DO114. Carried geometric anchors: SGA3 I4.7.3/4.7.3.1; SGA6 VI Theorem3.3; FGA232 Theorem3.1pp148-149; FGA236 Corollaries1.7/2.3pp229/232; SGA7 I VII2.7/2.9pp169-174 and VIII7.4(b)pp310-311; cote4 scan003 sections8.1.1-8.1.2/scans005-006 sections8.4.1-8.4.3 and cote54 scans005-007. No new formal-group slope, lambda-ring/K-theory identity, l-adic weight or CM-motive comparison. DP\_MULT Q3e/DP44/DP52/DP54-65 and THETA3 G1/A7 gain this finite reduction; CX\_GRH CX-ALG-024 and GCT\_PLETH T-052 retain their global gates.

\subsection{275. AM-DO-115}\label{am-do-115}

AM-DO-115 DC1-24 proves an all-parameter scalar theta estimate using the exact positive cubic kernel, finite cap occupancy and nonzero-shell integral in real dimension336. The original selected angular profile gives beta777/352, and the unchanged saturated transverse Gram gives a complete257/256 theta factor. With every original paired ideal, affine central fibre, finite factor and both continuum monotonicities retained, the complete k161 support26 cell is below2\^{}-233, with exact exponent-119477/512. The necessary actual all-II rank-two-unipotent remainder becomes162\textless=k\textless=164 with2\^{}(71443/64)\textless x\textless2\^{}2108,v=q26,b=j0. Exactly one233 append gives118/31 complete ordered errors, both entire sums below2\^{}-230; unique minimum231,next233 and complete floors26/11 remain. No162+, signed positivity, selected arithmetic content, effective source, period, specific032/Read, evaluatedDP44 or RH conclusion.

The universal positive-kernel and Gaussian estimates act on the unchanged actual selected complex realization. Original projective integral modules, paired CM lattices, saturated transverse basis, Hermitian metric, projection denominator and theta line are retained. Delsarte, Goethals, and Seidel(1977),Theorems3.3/3.6/4.3,Corollaries3.7-3.8 and Examples4.6-4.7,printed367-370,provide the exact primary context; the needed degree-three kernels and full cap/count form are proved explicitly. The previous exact Nakahama(2016)equation5.7/Theorem5.1pp960-961, Gaussian Fourier convention of Miller and Stephens-Davidowitz(2019)equations1.1/2.1, and Rankin(1955)Theorem1(iv) context remain in the directly checked scopes of K108/DO114. Carried geometric anchors: SGA3 I4.7.3/4.7.3.1; SGA6 VI Theorem3.3; FGA232 Theorem3.1pp148-149; FGA236 Corollaries1.7/2.3pp229/232; SGA7 I VII2.7/2.9pp169-174 and VIII7.4(b)pp310-311; cote4 scan003 sections8.1.1-8.1.2/scans005-006 sections8.4.1-8.4.3 and cote54 scans005-007. No new formal-group slope, lambda-ring/K-theory identity, l-adic weight or CM-motive comparison. DP\_MULT Q3e/DP44/DP52/DP54-65 and THETA3 G1/A7 gain this finite reduction; CX\_GRH CX-ALG-024 and GCT\_PLETH T-052 retain their global gates.

\subsection{276. AM-K-110}\label{am-k-110}

AM-K-110 PM1-22 retains the full actual angular profile as a positive Laplace mixture, includes the entire low-parameter tail in an explicit atom, and computes the exact inverse third moment. Weighted AM-GM and the all-minimum/all-parameter DO115 scalar theorem control arbitrary embedding weights in the unchanged original ideal and transverse Gram. Every transverse correction is retained, giving a uniform2180-bit row improvement. With all original affine central, paired ideal, finite and source factors and both continuum monotonicities, the complete k162 support26 cell is below2\^{}-1071, with exact exponent-548533/512. The necessary actual all-II rank-two-unipotent remainder becomes163\textless=k\textless=164 with2\^{}(71443/64)\textless x\textless2\^{}2108,v=q26,b=j0. Exactly one1071 append gives119/31 complete ordered errors, both entire sums below2\^{}-230; minimum231,next233 and complete floors26/11 remain. No k163/164, signed positivity, selected arithmetic content, effective source, period, specific032/Read, evaluatedDP44 or RH conclusion.

The positive real mixture is an estimate in the existing selected complex realization; it introduces no arithmetic scalar, field extension, new projector or freely chosen source vector. Original integral projective modules and paired CM ideals retain their exact half-trace dual and true transverse covolume. NIST DLMF equations10.32.10/10.39.2 supply the exact positive-real integral context, and the full normalized elementary proof is included in PM3. Delsarte, Goethals, and Seidel(1977),Theorems3.3/3.6/4.3,Corollaries3.7-3.8 and Examples4.6-4.7,printed367-370,provide the exact primary context; the needed degree-three kernels and full cap/count form are proved explicitly. The previous exact Nakahama(2016)equation5.7/Theorem5.1pp960-961, Gaussian Fourier convention of Miller and Stephens-Davidowitz(2019)equations1.1/2.1, and Rankin(1955)Theorem1(iv) context remain in the directly checked scopes of K108/DO114. Carried geometric anchors: SGA3 I4.7.3/4.7.3.1; SGA6 VI Theorem3.3; FGA232 Theorem3.1pp148-149; FGA236 Corollaries1.7/2.3pp229/232; SGA7 I VII2.7/2.9pp169-174 and VIII7.4(b)pp310-311; cote4 scan003 sections8.1.1-8.1.2/scans005-006 sections8.4.1-8.4.3 and cote54 scans005-007. No new formal-group slope, lambda-ring/K-theory identity, l-adic weight or CM-motive comparison. DP\_MULT Q3e/DP44/DP52/DP54-65 and THETA3 G1/A7 gain this finite reduction; CX\_GRH CX-ALG-024 and GCT\_PLETH T-052 retain their global gates.

\subsection{277. AM-DO-116}\label{am-do-116}

AM-DO-116 QC1--18 proves a universal real336-dimensional lattice theta bound for every positive Gaussian parameter from its actual nonzero minimum, using an exact positive degree-five spherical polynomial and the full cap/nonzero-shell argument. The same theorem applies to arbitrary positive embedding weights with the original half-trace minimum. Its exact coefficient K5 satisfies128K5²\textless Kcub², a scalar improvement of more than3.5 bits. No complete arithmetic sector, numerical error charge, support-floor change or independent main-gate movement is asserted. Current119/31 errors, floors26/11 and remainingk163..164 are unchanged.

The positive real metric estimate leaves original integral Hermitian modules, paired fractional ideals, polarization, theta line, selected complex/arithmetic realizations and source data unchanged. Delsarte, Goethals, and Seidel(1977),Definition2.1,Theorems2.4/3.3/4.3 and Corollary3.8,printed364--369,are the exact primary context; QC6--15 supplies the normalized proof and new explicit polynomial. Carried anchors: SGA3 I4.7.3/4.7.3.1; SGA6 VI Theorem3.3; FGA232 Theorem3.1pp148--149; FGA236 Corollaries1.7/2.3pp229/232; SGA7 I VII2.7/2.9pp169--174 and VIII7.4(b)pp310--311; cote4 scan003 sections8.1.1--8.1.2/scans005--006 sections8.4.1--8.4.3 and cote54 scans005--007. No new formal-group slope, lambda/K-theory identity, l-adic weight or CM-motive comparison. DP\_MULT Q3e,THETA3 G1/A7,CX\_GRH CX-ALG-024 andGCT\_PLETH T-052 retain their global gates.

\subsection{278. AM-K-111}\label{am-k-111}

AM-K-111 HP1--18 proves the universal real336-dimensional lattice-minimum/all-positive-parameter scalar theorem using a fixed exact degree-fifteen polynomial with16 strictly positive normalized Gegenbauer coefficients. Its full harmonic positivity, cap projection, nonzero-shell integral and original half-trace/anisotropic normalization give K15²\textless2\^{}-21 Kcub². This is only a scalar coefficient improvement; the positive zero-dual term, all arithmetic patterns and entire row endpoint remain required. No error charge or standalone main-gate movement. Current119/31 arrays,floors26/11 and actual residualk163..164 remain unchanged.

The polynomial and cap proof estimate the existing selected real metric and leave every original integral module, paired ideal, polarization, theta line and arithmetic source fixed. Delsarte,Goethals,andSeidel(1977),Definition2.1,Theorems2.4/3.3/4.3 andCorollary3.8,printed364--369,are the exact primary context; HP8--16 proves the normalized degree0--15 and all-minimum/all-parameter form explicitly. Carried anchors: SGA3 I4.7.3/4.7.3.1; SGA6 VI Theorem3.3; FGA232 Theorem3.1pp148--149; FGA236 Corollaries1.7/2.3pp229/232; SGA7 I VII2.7/2.9pp169--174 and VIII7.4(b)pp310--311; cote4 scan003 sections8.1.1--8.1.2/scans005--006 sections8.4.1--8.4.3 and cote54 scans005--007. No new formal-group slope, lambda/K-theory identity, l-adic weight or CM-motive comparison. DP\_MULT Q3e,THETA3 G1/A7,CX\_GRH CX-ALG-024 andGCT\_PLETH T-052 retain their global gates.

\subsection{279. AM-LIFT-094}\label{am-lift-094}

AM-LIFT-094 HE0-22 proves conditional effective algebraic Hecke escape on a fixed indefinite arithmetic family. For a reduced algebraic bad locus of total degreeD and dimension\textless=s, finite equal-ample-pullback correspondences of degree\textless=delta with no nonempty proper closed forward-invariant locus admit escape in at mostB\_delta,s(D) steps, where A\_delta(D)=2 delta\^{}D D\^{}2, B\_delta,0(D)=D and B\_delta,s(D)=D+B\_delta,s-1(A\_delta(D)). For actual projective O\_E lattice rankr\textgreater=3 indefinite at every real place, f={[}F:Q{]}, an explicitly cleared rational Witt frame gives t=1+3McQ\_S, two physical original-M-level Hecke correspondences with delta=t\^{}(4fr\^{}2) and multiplier t\^{}2. Every forward orbit is dense; the conditional accumulatedN\textless=t\^{}B is obtained. The actual regular LIFT085 period-family map and nonempty inverse image of the prescribed nonvanishing open remain hypotheses; H,D,c,N0 and admitted Host/Sizes costs are not evaluated. Same original fields, quotient/action/line/level/frame conditions and independent entry multiplierN0 remain. No selected arithmetic nonzero content, global period, specific032/Read/STiesR, DP44 or RH conclusion. No finite-error charge;119/31 and floors26/11 unchanged.

The new argument is algebraic degree descent on the actual regular period family, with finite Hecke correspondences, equal ample Chern classes and proper pushforward. It preserves the physical CM Picard/Poincare cubical line and its multiplier under original-level isogenies. Mandatory exact locator: Datum in TwoWeightEqOneAddRoot.lean is lines30-49, replacing the isolated30-47 in the frozen proof; the final condition at49 was directly read. The complete immutable locator addendum accompanies this theorem. Fresh primary anchors: Deligne(1982) Proposition4.1/Corollary4.2pp28-29; Klingler-Yafaev(2014) Proposition5.3.2(1),Remark5.3.4 and extension after(5.1)pp24-25; Daw-Javanpeykar-Kuhne(2020) sections3.1-3.2,Proposition10.1 with its distinct positive-dimensional scope; Stacks0B1N Lemma43.23.1 and02SU Lemma42.26.4. Carried geometric anchors from LIFT085: FGA232 Theorem3.1pp148-149, FGA236 Corollaries1.7/2.3pp229/232; SGA7 I VII2.7/2.9pp169-174,VIII7.4(b)pp310-311; cote4 scan003 sections8.1.1-8.1.2/scans005-006 sections8.4.1-8.4.3 and cote54 scans005-007. No new crystalline slope, lambda-ring/K-theory identity, weight or CM-motive comparison is established. DP\_MULT Q3e/DP44/DP52/DP54-65 gain a conditional finite algebraic escape input; THETA3 G1/A7,CX\_GRH CX-ALG-024 and GCT\_PLETH T-052 retain their global gates.

\subsection{280. AM-R-094}\label{am-r-094}

AM-R-094 PP1-24 removes the entire remaining conditional J56 all-top support26 class. R091 plus K109/DO115/K110 leaves only actual all-II global-rank-two-unipotent scalar-cross cells k163/164; the full correlated finite-pattern containing polynomial and low-parameter angular mixture give complete exponents -199349/512\textless-389 and -159413/512\textless-311, using DO116 and K111 respectively. Their union is below2\^{}-310. Append310 once to the exact K110119/31 arrays:120/31, both entire sums\textless2\^{}-230, complete floors27/11. The actual remainder now has support at least27 with every earlier higher cut and original finite/real/arithmetic convention retained. PP3, the sufficient signed higher-support inequality, remains open, as do actual Host/Sizes, selected nonzero arithmetic content, global period, specific032/Read/STiesR and evaluatedDP44/RH. No actual shorter mixed-CM chain or universal strict-negative cost theorem is asserted.

This is a bound on the same selected real theta coefficient and finite correspondences for the original integral Hermitian modules. The local paired fractional ideals and affine central difference ideal are retained before nonnegative enlargement; the containing polynomial is not an equality for the original signed coefficients. Picard/Poincare line, polarization, finite-etale torsion realizations, projectors and arithmetic source are unchanged. Primary scalar context carried through fully checked DO116/K111: Delsarte,Goethals,andSeidel(1977),Definition2.1,Theorems2.4/3.3/4.3 and Corollary3.8,pp364-369. Geometric anchors carried through R091 and LIFT085: FGA232 Theorem3.1pp148-149,FGA236 Corollaries1.7/2.3pp229/232; SGA7 I VII2.7/2.9pp169-174,VIII7.4(b)pp310-311; SGA3 I4.7.3/4.7.3.1; SGA6 VI Theorem3.3; cote4 scan003 sections8.1.1-8.1.2/scans005-006 sections8.4.1-8.4.3 and cote54 scans005-007. No new Dieudonne slope, lambda/K-theory relation, l-adic weight or CM-motive comparison is asserted. DP\_MULT Q3e/DP44/DP52/DP54-65 and THETA3 G1/A7 have a strictly smaller conditional finite remainder; their actual arithmetic chain/pairing gates, CX\_GRH CX-ALG-024 and GCT\_PLETH T-052 remain open globally.

\subsection{281. AM-DO-117}\label{am-do-117}

AM-DO-117 FC1-25 removes the entire original conditional J56 all-top H133 fixed-S nonunipotent support27 sector, at every global rank and incidence profile. The five complete containing-set exponents are -5531348963/4096, -3078643005/4096, -2767418421/1024, -9563431915/2048 and -120294901/4096; their union is below2\^{}-29365. The new small scalar remains in the full paired I\_Sigma, while its exact original near-point transport uses only the universal RT2 selected Hilbert trace bound. Append29365 once to the complete R094120/31 predecessor:121/31, both full sums\textless2\^{}-230, complete floors27/11 unchanged. This is a strictly smaller conditional support27 remainder, with no full support27 exhaustion, actual admitted Host/Sizes, selected arithmetic content, prescribed global period, specific032 pairing, evaluatedDP44 or RH conclusion.

This is an estimate for the same selected real theta coefficient on the original integral Hermitian modules and finite correspondences. Exact null-image transport preserves the original fractional ideal, packet, compact factors, vector-valued metric and Hilbert/PBW normalization. The carried primary analytic source is Nakahama(2016), equation(5.7) and Theorem5.1, through the checked K043/K046 dictionary and DO080 RT1-24. The unchanged geometric anchors are SGA3 I4.7.3/4.7.3.1; SGA6 VI Theorem3.3; FGA232 Theorem3.1 pp148-149; FGA236 Corollaries1.7/2.3 pp229/232; SGA7 I VII2.7/2.9 pp169-174 and VIII7.4(b) pp310-311; cote4 scan003 sections8.1.1-8.1.2 and scans005-006 sections8.4.1-8.4.3; cote54 scans005-007. These are inherited locators, not new whole-source readings or constructions. No new Dieudonne slope, lambda/K-theory relation, l-adic weight, CM-motive map or global arithmetic seesaw is asserted. DP\_MULT Q3e/DP44 and THETA3 G1 retain their actual arithmetic gates; CX\_GRH CX-ALG-024 and GCT\_PLETH T-052 retain their established scope.

\subsection{282. AM-K-112}\label{am-k-112}

AM-K-112 UE1-14 proves that every active type-II place in the exact full-positive globally rank-two-unipotent form is eligible, without an eta-denominator bound. The noneligible active subset is empty, with zero error charge. In the actual all-II scalar-cross support27 remainder, v=s=27, x\textgreater2\^{}2111; x\textless2\^{}3107 forces total eligibility27; x\textless2\^{}12000 gives total eligibility at most36 and at most one eta per fixed support. Fixed-eta active labels contribute at most(T-1)\^{}-27. The full168-factor bound additionally requires total eligibility27 and no inactive paired-C-zero place. Original paired row and central difference ideals retain their support27 denominator power54. Current DO117121/31 exact errors and complete floors27/11 remain unchanged; no full support27 estimate, actual source/content/period, specific032, DP44 or RH conclusion is asserted.

The theorem computes the integral stabilizer of the unchanged paired Hermitian lattices over the original selected split DVR. The original finite correspondences, omitted labels, inverse-adjoint companion, Picard/Poincare polarization data and selected real coefficient are retained. It constructs no new arithmetic realization. Carried contextual anchors: SGA3 I4.7.3/4.7.3.1; SGA6 VI Theorem3.3; FGA232 Theorem3.1 pp148-149; FGA236 Corollaries1.7/2.3 pp229/232; SGA7 I VII2.7/2.9 pp169-174 and VIII7.4(b) pp310-311; cote4 scan003 sections8.1.1-8.1.2 and scans005-006 sections8.4.1-8.4.3; cote54 scans005-007. These are inherited exact locators, not fresh primary readings or new lambda/K-theory, l-adic weight, Dieudonne or CM-motive theorems. DP\_MULT Q3e/DP44 and THETA3 G1 have a smaller conditional finite remainder; their actual arithmetic constructions, CX\_GRH CX-ALG-024 and GCT\_PLETH T-052 retain their separate global gates.

\subsection{283. AM-LIFT-095}\label{am-lift-095}

AM-LIFT-095 IH1-6 proves an integral hyperbolic pair inside every stated selfdual relative-unramified indefinite Hermitian lattice of rank at least3, by local integral pairs and strong approximation on SU. The original nonterminal source W\_j.L0 satisfies these hypotheses conditionally on the admitted full record, and the same E-linear operators preserve its actual AA multiple. Choosing auxiliary doubled horizontal and integral imaginary root bases gives c=1 in LIFT094: t=1+3M Q\_S, delta=t\^{}(4fr\^{}2), L=t; any separately required larger congruence modulus is retained. This removes the unevaluated auxiliary clearing integer without a replacement lattice, seed translation or extra isogeny. The actual very ample line H, same-rule bad-locus degree D, nonempty family intersection, component-entry multiplier N0, selected nonzero rule and full Host/Sizes remain unresolved inputs. No evaluated total N, period, specific032, DP44 or RH conclusion; no finite charge. Current DO117121/31 arrays and floors27/11 remain unchanged.

The construction is an integral orthogonal direct-summand decomposition of the same finite projective Hermitian module over Spec(O\_E), followed by integral E-linear root endomorphisms and the same fractional-ideal multiple. It retains the source alternating form signed Tr(xi0 h), the original polarization and exact M-level. The resulting isogenies use the unchanged rigidified Poincare/cubical-line transport of LIFT085/LIFT094. Fresh primary group input: Rapinchuk(2013),Theorem2.3,p281; Milne(2017),Definition3.18,p40 and Theorem4.16,p48. Carried rational classification: Deligne(1982),Proposition4.1,p28. Carried projective-module/exterior context: SGA3 I4.7.3/4.7.3.1 and SGA6 VI Theorem3.3; Picard/biextension context: FGA232 Theorem3.1 pp148-149,FGA236 Corollaries1.7/2.3 pp229/232,SGA7 I VII2.7/2.9 pp169-174 and VIII7.4(b) pp310-311. Formal-crystal/Cartier-duality anchors remain cote4 scan003 sections8.1.1-8.1.2,scans005-006 sections8.4.1-8.4.3 and cote54 scans005-007. No new formal slope, motivic comparison, l-adic weight or global oscillator map is inferred. DP\_MULT Q3e/DP44/DP52/DP54-65 and THETA3 G1/A7 receive a sharper conditional actual-source cost interface; CX\_GRH CX-ALG-024 and GCT\_PLETH T-052 keep their separate global gates.

\subsection{284. AM-DO-118}\label{am-do-118}

AM-DO-118 UB1-29 removes the entire actual conditional J56 fixed-S globally rank-two-unipotent support27 sector, all original profiles and labels. The rank35 high envelope retains the full five-column and same-pair isometry fibre. In the rank34 low branch, the true row metric, full beta norm fibre, affine central cosets and actual selected trace are retained; fixing row and beta before eta-dependent labels permits the one-label bound without global eta uniqueness. The six complete strict dyadic bounds have exponents8792,4491,3928,23649,17510,11371; their union is below2\^{}-3927. Append3927 once to DO117121/31, giving122/31 and both full sums below2\^{}-230; complete floors27/11 remain unchanged. This does not remove the other support27 classes, establish a support28 floor, or solve the signed inequality, actual admitted source, selected arithmetic content, prescribed period, specific032 pairing, evaluatedDP44 or RH.

The count uses the same integral Hermitian modules and original finite correspondences. The row and central difference ideals remain fractional projective modules on Spec(O\_E) and Spec(O\_F); the exact relative norm fibre includes all ideal classes and norm-one phases. The metric is the true fixed Gram/Schur metric, and the actual selected theta trace retains its Hilbert normalization and packet transport. Carried exact anchors are SGA3 I4.7.3/4.7.3.1; SGA6 VI Theorem3.3; FGA232 Theorem3.1 pp148-149; FGA236 Corollaries1.7/2.3 pp229/232; SGA7 I VII2.7/2.9 pp169-174 and VIII7.4(b) pp310-311; cote4 scan003 sections8.1.1-8.1.2 and scans005-006 sections8.4.1-8.4.3; cote54 scans005-007. No new primary reading, formal slope, lambda/K-theory relation, l-adic weight or CM-motive comparison is inferred. DP\_MULT Q3e/DP44 and THETA3 G1/A7 receive only this smaller conditional finite remainder; CX\_GRH CX-ALG-024 and GCT\_PLETH T-052 retain their separate global gates.

\subsection{285. AM-K-113}\label{am-k-113}

AM-K-113 SC27.1-27 removes seven complete containing groups in the actual conditional J56 all-II global-rank-two-unipotent scalar-cross support27 sector. The33 complete endpoints retain all four transverse mixture terms, every G pattern, all inactive-zero subsets, chosen active-integral positions, coordinate sizes and a separately enumerated infinite eta-ideal tail. Their full union is below2\^{}-304. The necessary residual is b=j0,2\^{}2111\textless x\textless2\^{}3328, either total eligibility27 and100\textless=k\textless=164, or total eligibility28,k164 with x\textgreater2\^{}3107. Append304 once to actual DO118122/31, giving123/31 with both full sums\textless2\^{}-230; complete floors27/11 remain unchanged. No whole support27 exhaustion, actual admitted Host/Sizes, signed inequality, selected arithmetic content, prescribed period, specific032 pairing, evaluatedDP44 or RH conclusion.

All local row and central difference ideals are the original fractional projective modules over Spec(O\_E)/Spec(O\_F), with the same CM involution, polarization, Picard/Poincare line, finite correspondences and selected real metric. Positive mixture, scalar packing and finite-pattern polynomials estimate these original objects without replacing their arithmetic source. Primary analytic context is carried through checked K107/K108/K110: Nakahama(2016),equation5.7 and Theorem5.1 pp960-961; NIST DLMF10.32.10/10.39.2. Scalar context is Delsarte,Goethals,andSeidel(1977),Theorems3.3/3.6/4.3 through fully checked K111 HP1-18. Carried geometric anchors: SGA3 I4.7.3/4.7.3.1; SGA6 VI Theorem3.3; FGA232 Theorem3.1 pp148-149; FGA236 Corollaries1.7/2.3 pp229/232; SGA7 I VII2.7/2.9 pp169-174 and VIII7.4(b) pp310-311; cote4 scan003 sections8.1.1-8.1.2/scans005-006 sections8.4.1-8.4.3 and cote54 scans005-007. No new formal slope, lambda/K-theory relation, l-adic weight or CM-motive comparison is inferred. DP\_MULT Q3e/DP44 and THETA3 G1/A7 receive only the stated smaller conditional finite gate; CX\_GRH CX-ALG-024 and GCT\_PLETH T-052 remain globally open at their established gates.

\subsection{286. AM-R-095}\label{am-r-095}

AM-R-095 FW1-21 removes the entire actual conditional J56 fixed-S globally rank-one-unipotent support27 sector, all original profiles, labelled pairs, double-pole sizes and unmixed G patterns. A uniform weighted eta/content cutoff precedes the complete334-place binary pattern product; the remaining z exponent is165f, both genuine affine coordinates and the whole support/marked/type sum remain. The reduced exponent is-54242909/4096 and the complete witness exponent is-95126187/256; their union is below2\^{}-13000. Append13000 once to actual K113123/31, giving124/31 with both full sums\textless2\^{}-230; complete floors27/11 remain. No whole support27 exhaustion, signed inequality, actual admitted Host/Sizes, selected arithmetic content, prescribed period, specific032 pairing, evaluatedDP44 or RH conclusion.

The entry, content and coordinate ideals remain the same fractional invertible modules over Spec(O\_F)/Spec(O\_E), with the original paired CM action. R089 real-embedding packing bounds the reciprocal norm on these ideals without changing the polarization, Picard class, theta line, Poincare biextension or selected realization. Carried exact anchors: SGA3 I4.7.3/4.7.3.1; SGA6 VI Theorem3.3; FGA232 Theorem3.1 pp148-149; FGA236 Corollaries1.7/2.3 pp229/232; SGA7 I VII2.7/2.9 pp169-174 and VIII7.4(b) pp310-311; cote4 scan003 sections8.1.1-8.1.2/scans005-006 sections8.4.1-8.4.3 and cote54 scans005-007. No new external theorem, Dieudonne slope, lambda/K-theory relation, l-adic weight or CM-motive comparison. DP\_MULT Q3e/DP44/DP52/DP54-65 and THETA3 G1/A7 receive only this smaller conditional finite remainder; CX\_GRH CX-ALG-024 and GCT\_PLETH T-052 retain their global gates.

\subsection{287. AM-DO-119}\label{am-do-119}

AM-DO-119 NA1-29 proves every active place of the actual full-positive globally rank-two-unipotent form is eligible, including distinct source/target labels and every type. Its complete non-all-II scalar-cross support27 bound retains all original paired rows, full-entry poles, affine central cosets, finite weights, relative norm fibres and the entire denominator-ideal tail. Six exhaustive branches give one412 charge. Append412 once to actual R095124/31, yielding125/31 with both complete sums below2\^{}-230 and floors27/11 unchanged. The dense all-II27 and separate nonzero-short-compression sector remain open; no full support28, actual admitted source/Host/Sizes, selected arithmetic content, prescribed period, specific032 pairing, evaluatedDP44 or RH conclusion.

The new lemma concerns the original labelled lattice correspondence over each selected split DVR. Its fractional row and central ideals remain the same invertible modules over Spec(O\_E)/Spec(O\_F); the real central fibre is an affine torsor with the original shift, metric and codifferent normalization. Actual paired CM involution, polarizations, Picard and theta lines, Poincare biextensions, selected realization, classes and Haar measures are unchanged. Exact carried anchors: SGA3 I4.7.3/4.7.3.1; SGA6 VI Theorem3.3; FGA232 Theorem3.1 pp148-149; FGA236 Corollaries1.7/2.3 pp229/232; SGA7 I VII2.7/2.9 pp169-174 and VIII7.4(b) pp310-311; cote4 scan003 sections8.1.1-8.1.2 and scans005-006 sections8.4.1-8.4.3; cote54 scans005-007. No new lambda/K-theory, l-adic weight, Dieudonne or CM-motive theorem is inferred. DP\_MULT Q3e/DP44 and THETA3 G1/A7 gain only this smaller conditional finite remainder; CX\_GRH CX-ALG-024 and GCT\_PLETH T-052 retain their global gates.

\subsection{288. AM-K-114}\label{am-k-114}

AM-K-114 FP27.1-29 removes the entire full-positive-compression global-rank-three scalar-cross support27 sector with nonzero original S-cross. All34 complete ideal/type cells,2925 padded sets, original full167-covector metric, relative norm fibre, projected U denominator/T108 column and G2 source fibre are retained. The complete eta-first all-ideal tail requires no coordinate sparsity. Literal inversion and23 common paired eligible positions force equal null lines outside four containing sectors, contradicting rank3. Their full union costs1650 once. Append1650 to actual DO119125/31, yielding126/31, both entire sums below2\^{}-230 and complete floors27/11 unchanged. Full positive compression remains a premise; the separate nonzero-short-compression27 and dense all-II27 sectors remain open. No full support28, actual admitted Host/Sizes, selected content, prescribed period, specific032, evaluatedDP44 or RH conclusion.

The actual fractional coordinate and denominator ideals remain invertible modules over Spec(O\_E)/Spec(O\_F). Relative ideal norm and the paired divisor of the line difference use those same modules; the projected U denominator, original Hermitian Gram, source/target lattice pair, selected metric, finite projectors and class/Haar factors are retained. Inversion is the original adjoint correspondence with its theta/Picard/Poincare-biextension data, not a new integral model. Exact inherited anchors: SGA3 I4.7.3/4.7.3.1; SGA6 VI Theorem3.3; FGA232 Theorem3.1 pp148-149; FGA236 Corollaries1.7/2.3 pp229/232; SGA7 I VII2.7/2.9 pp169-174 and VIII7.4(b) pp310-311; cote4 scan003 sections8.1.1-8.1.2 and scans005-006 sections8.4.1-8.4.3; cote54 scans005-007. No new lambda/K-theory, l-adic, Dieudonne or CM-motive comparison is inferred. DP\_MULT Q3e/DP44 and THETA3 G1/A7 receive only this conditional finite reduction; CX\_GRH CX-ALG-024 and GCT\_PLETH T-052 retain their arithmetic gates.

\subsection{289. AM-LIFT-096}\label{am-lift-096}

AM-LIFT-096 TJ1-18 proves the same-rule bad-locus degree bound on the actual principal-3M special-unitary family of LIFT085/LIFT094, under its stated nonempty same-rule open-locus hypothesis. For g=fr,d=fpq, total frame weight w and theta weight n, take k=gw+n/2 for the literal even-n Dom; squaring the same rule gives k=2gw+n in general without changing its reduced zero set. A finite proof cover of degree at most e\_g=2\^{}(8g\^{}2+2g+4), first jets of the fourth power of the original symmetric principal line, and exact theta multiplier identities produce scalar cuts. For H=(psi\_bar\^{}\emph{L\_S)\^{}h very ample and V=deg\_H(Ybar), proper pushforward gives D\textless=D\_}=e\_g V sum\_\{j=0\}\^{}d k\^{}j, including mixed-dimensional components. When k=0 and the bad locus is proper it is empty. With LIFT095 c=1, LIFT094 then gives the conditional within-family multiplier N\textless=t\^{}\{B\_\{t\^{}\{4fr\^{}2\},fpq-1\}(D\_\emph{)\}, t=1+3M Q\_S; any independent component-entry multiplier N0 is retained. The proof cover contributes only to D\_}, not to the physical isogeny or output field. The numerical h,V, nonempty actual-family intersection, controlled N0, selected arithmetic nonzero rule and admitted full Host/Sizes remain open. No evaluated total N, six-identity chain, actual period, specific032 pairing, DP44 sign or RH conclusion; no finite error charge. Current K114126/31 exact arrays and complete floors27/11 are unchanged.

The construction uses the relative first-principal-parts sequence of the original rigidified cubical line, finite etale theta-structure and multiplier torsors, the descended invariant theta line, the representable frame torsor on the smooth literal moduli scheme, and proper first-Chern-class pushforward on the normal Baily-Borel compactification. Fresh primary locators: Candelori(2016) Lemma2.5.3/Corollary2.5.4 pp9-10, Definition3.2.2 p20, Theorem3.3.1 proof p21, Theorems4.1.3/4.2.1 pp28-29; Lubicz-Robert(2022) Introduction1.1 p1; Klingler-Yafaev(2014) Proposition5.3.2(1), Remark5.3.4, Proposition5.3.6 pp24-26; Milne(2017) Theorem3.12/Remark3.13 pp38-39; Kieffer-Page-Robert(2025) Section4.3/Proposition4.8 pp24-25 of the45-page author version; Stacks Tags0EBJ/02SU. The standard ample-plus-nef step has the additional exact root/peer locator Lazarsfeld(2004), Positivity in Algebraic Geometry I, Corollary1.4.10 pp53-54; this is a locator addition, not a change to the frozen proof. Carried Picard/biextension anchors: FGA232 Theorem3.1 pp148-149; FGA236 Corollaries1.7/2.3 pp229/232; SGA7 I VII2.7/2.9 pp169-174 and VIII7.4(b) pp310-311. Projective-module/exterior context: SGA3 I4.7.3/4.7.3.1 and SGA6 VI Theorem3.3. Formal-crystal and Cartier-duality context remains cote4 scan003 sections8.1.1-8.1.2,scans005-006 sections8.4.1-8.4.3, and cote54 scans005-007. These carried anchors do not turn the construction into a new slope, lambda-ring, K-theory, l-adic weight, CM-motive or global oscillator comparison. DP\_MULT Q3e/DP44/DP52/DP54-65 and THETA3 G1/A7 receive a more explicit conditional geometric-cost interface; CX\_GRH CX-ALG-024 and GCT\_PLETH T-052 retain their separate global gates.

\subsection{290. AM-DO-120}\label{am-do-120}

AM-DO-120 QW1-29 removes the entire original doubly-plane scalar-cross sector with nonzero four-entry short compression at every support27..31, including all global characteristic types, actual coordinate supports and paired labels. The original integral positive F witness q\_j=\textbar\textbar(A-I)y\_j\textbar\textbar\^{}2 has every conjugate below2\^{}42. The complete paired finite low-rank proof restores the rank-four Green baseline666 with coefficients331/2 and335/2; both original full source/target costs, selected trace, arbitrary-ideal row, relative norm fibre and all incidence weights remain. All five full endpoint bounds give one270871 charge. Appending it once to K114126/31 yields127/31, with both entire sums below2\^{}-230 and complete floors27/11 unchanged. The complementary scalar-cross remainder at supports27..31 has full positive compressionI, a positive166 fixed space and global rank at most3 by the support-free original integral scalar proof; global unipotence is not inferred. The dense all-II27 window remains open; no complete support28, actual admitted Host/Sizes, selected arithmetic content, prescribed period, specific032 value, evaluatedDP44 or RH conclusion.

The new q\_j are integral trace/norm pairing sections over Spec(O\_F) from the original Hermitian modules over Spec(O\_E). Their assigned F-divisor is simultaneous for both original short vectors; its norm is the same determinant of multiplication. The finite low-rank flags and Schur minors live in the actual G-unit reductions and assert no integral global quotient or moving integral frame. The original fractional row ideals, CM involution, metric, finite projectors, classes, Haar measures, polarization, theta line, Picard data and Poincare biextensions remain unchanged. Carried exact anchors: SGA3 I4.7.3/4.7.3.1; SGA6 VI Theorem3.3; FGA232 Theorem3.1 pp148-149; FGA236 Corollaries1.7/2.3 pp229/232; SGA7 I VII2.7/2.9 pp169-174 and VIII7.4(b) pp310-311; cote4 scan003 sections8.1.1-8.1.2 and scans005-006 sections8.4.1-8.4.3; cote54 scans005-007. No new lambda-ring/K-theory, l-adic weight, Dieudonne or CM-motive comparison is inferred. DP\_MULT Q3e/DP44 and THETA3 G1/A7 receive this strictly smaller conditional finite remainder; CX\_GRH CX-ALG-024 and GCT\_PLETH T-052 retain their global gates.

\subsection{291. AM-LIFT-097}\label{am-lift-097}

AM-LIFT-097 VA1-15 evaluates the previously unspecified very-ampleness exponent on the same actual principal-3M arithmetic compactification X. For d=fpq,g=fr and A=psi\_bar\^{}\emph{L\_S, the exact integer h\_d=2\^{}(d+1)(d+1)(d+2)! makes H=A\^{}h\_d very ample. The no-PGL2 boundary condition and canonical-line extension identify K\_X as Cartier ample; the neat Baily-Borel pair is log canonical, A-K\_X is nef, and both applications of effective freeness/very ampleness retain log bigness on every lc/slc stratum. The original Hodge restriction and Hecke equal-pullback identities are unchanged. Writing v\_A=int\_X c1(A)\^{}d, the full ambient degree is h\_d\^{}d v\_A, so LIFT096 gives D\_}=2\^{}(8g\^{}2+2g+4)h\_d\^{}d v\_A sum\_\{j=0\}\^{}d k\^{}j, with k=gw+n/2 for even n or2gw+n after the same-zero square. The independent quantity v\_A is a positive unevaluated integer. LIFT094/095 then retain N\_within\textless=t\^{}\{B\_\{t\^{}(4fr\^{}2),d-1\}(D\_*)\},t=1+3M Q\_S, with every protected congruence condition and any separate component-entry N0. The same-family nonempty open, selected arithmetic nonzero rule and admitted Host/Sizes remain open. No evaluated actual N, new physical field/level/isogeny, six-identity chain, period, specific032 pairing, DP44 or RH conclusion; no finite error charge. Current DO120127/31 arrays and complete floors27/11 are unchanged.

This is a coherent-sheaf and algebraic compactification result for the same original family: rank-one reflexive canonical extension over a normal scheme, actual Cartier positivity, log-canonical vanishing, algebraic section-ring generation and top first-Chern-class intersection. Exact fresh primary inputs are Alexeev(1996), Section3 Theorem3.2/Corollary3.3 pp9-10, independent of the conditional general-moduli discussion; Fujino(2009), Theorem1.1/Remark1.2 p475 and log-bigness definition p476; Fujino(2017), independent all-dimensional Section7 Lemma7.1 and full proof pp367-368; Klingler-Yafaev(2014), Proposition5.3.2(1), Remark5.3.4, Proposition5.3.6 authorpp24-26; Milne(2017), Theorem3.12/Remark3.13 pp38-39. The standard ample-plus-nef implication is carried with Lazarsfeld(2004), Corollary1.4.10 pp53-54. The prescribed theta line, OE action, Picard data and Poincare/cubical-line transports are unchanged. Carried contextual anchors: FGA232 Theorem3.1 pp148-149; FGA236 Corollaries1.7/2.3 pp229/232; SGA7 I VII2.7/2.9 pp169-174 and VIII7.4(b) pp310-311; SGA3 I4.7.3/4.7.3.1; SGA6 VI Theorem3.3; cote4 scan003 sections8.1.1-8.1.2 and scans005-006 sections8.4.1-8.4.3, cote54 scans005-007. No new lambda-ring/K-theory equality, l-adic sheaf/weight, Dieudonne-crystal or CM-motive comparison follows. DP\_MULT Q3e/DP44/DP52/DP54-65 and THETA3 G1/A7 receive a conditional cost bound with one fewer geometric input; CX\_GRH CX-ALG-024 and GCT\_PLETH T-052 retain their arithmetic/global gates.

\subsection{292. AM-R-097}\label{am-r-097}

AM-R-097 NW1-9 exhaustively composes the checked structural cuts at support27. Every remaining actual support27 arrow is full-positive, scalar-cross, globally rank-two-unipotent with all27 active places typeII, b=j0 and one of two necessary alternatives: total eligibility27,100\textless=k\textless=164,2\^{}2111\textless x\textless2\^{}3328; or total eligibility28,k164,2\^{}3107\textless x\textless2\^{}3328. DO120 supplies full positive compression before K114 is applied. The signed energy retains the actual R27 subset of this containing window and R\_ge28, with \textbar epsilon\textbar\textless2\^{}-230. No set equality, existence, emptiness, sign or floor28 conclusion. No new charge; actual DO120127/31 ordered arrays and complete floors27/11 are unchanged. Actual mixed-CM source, selected arithmetic content, six identities, prescribed period, specific032 pairing, evaluatedDP44, Kummer-Hecke RH and global oscillator comparison remain open.

This assembly concerns the same integral projective OE modules, paired CM ideals, polarization and finite correspondences over Spec(OE)/Spec(OF), with the same selected Hilbert realization. Positive compression is proved over the original global field; it is not transferred to residue-field positivity. Original Picard/theta/Poincare data and actual signed coefficients are retained. Carried exact anchors: SGA3 I4.7.3/4.7.3.1; SGA6 VI Theorem3.3; FGA232 Theorem3.1 pp148-149; FGA236 Corollaries1.7/2.3 pp229/232; SGA7 I VII2.7/2.9 pp169-174 and VIII7.4(b) pp310-311; cote4 scan003 sections8.1.1-8.1.2 and scans005-006 sections8.4.1-8.4.3; cote54 scans005-007. These are unchanged contextual interfaces, not fresh primary-source readings or new formal-group/Dieudonne, lambda-ring/K-theory, l-adic-weight or CM-motive comparison claims. DP\_MULT Q3e/DP44 and THETA3 G1/A7 retain the actual-family gates; CX\_GRH CX-ALG-024 and GCT\_PLETH T-052 remain open at their recorded global gates.

\subsection{293. AM-K-115}\label{am-k-115}

AM-K-115 GM1-23 proves the cap-dependent shape83/rate7/2 positive Gamma majorant, retains its entire low-weight interval by a paid positive shift, and controls every inverse moment0..k+3. All scalar and transverse terms, original saturated metric, central affine coset, source reconstruction,168 T factors and four-state finite patterns remain. The complete ten cells with total eligibility27,b=j0,2\^{}2111\textless x\textless2\^{}3328,k100..109 have total absolute contribution below2\^{}-945. The worst full exponent is-2907383/3072\textless-946, every other is below-3166. One945 append changes DO120127/31 to128/31, both entire sums remain below2\^{}-230 and complete floors27/11 stay. Combined with R097, the remaining actual support27 set is contained in the full-positive scalar-cross rank-two-unipotent all-II class with either total eligibility27,110\textless=k\textless=164,2\^{}2111\textless x\textless2\^{}3328, or eligibility28,k164,2\^{}3107\textless x\textless2\^{}3328. The next positive sufficient envelope is not an actual lower bound. No floor28, signed nonvanishing, admitted source, selected arithmetic content, prescribed period, specific032 value, evaluatedDP44 or RH conclusion.

The analytic probability measure is a bound on the same selected real realization; it introduces no arithmetic parameter, field extension, physical source vector, lattice rescaling or projector. All original projective OE modules, fractional paired ideals, CM involution, polarization, theta line, Picard and Poincare data remain. The new Gamma/beta/Holder and positive-integral arguments are proved directly; original representation and harmonic-code primary scopes are inherited through K107/K108/K110/K111/K113. Carried exact anchors: SGA3 I4.7.3/4.7.3.1; SGA6 VI Theorem3.3; FGA232 Theorem3.1 pp148-149; FGA236 Corollaries1.7/2.3 pp229/232; SGA7 I VII2.7/2.9 pp169-174 and VIII7.4(b) pp310-311; cote4 scan003 sections8.1.1-8.1.2 and scans005-006 sections8.4.1-8.4.3; cote54 scans005-007. No new formal-group/Dieudonne, lambda-ring/K-theory, l-adic-weight or CM-motive comparison. DP\_MULT Q3e/DP44 and THETA3 G1/A7 receive a strictly smaller conditional finite gate; CX\_GRH CX-ALG-024 and GCT\_PLETH T-052 retain their global gates.

\subsection{294. AM-LIFT-098}\label{am-lift-098}

AM-LIFT-098 AV1-20 evaluates the previously unevaluated automorphic volume for the same actual principal-3M special-unitary arithmetic family. Put n=3M, r=p+q\textgreater=3, f={[}F:Q{]}, g=fr, d=fpq, a=r\^{}2-1, a\_r=product\_\{j=1\}\^{}\{r-1\} j!/(2pi)\^{}(j+1), b\_pq=product\_\{i=0\}\^{}\{p-1\}i!/(q+i)!, and e\_v=product\_\{j=2\}\^{}r(1-epsilon\_v\^{}j q\_v\^{}(-j))\^{}(-1), epsilon\_v=1 at split and -1 at inert finite places. The literal relative-unramified CM and ordinary selfdual-lattice assumptions give hyperspecial models even at dyadic places. The original invertible ideal multiple changes no stabilizer or congruence index. Exact strong approximation and finite orders give mu\_Gamma=D\_F\^{}(a/2)a\_r\^{}f n\^{}(fa) product\_\{v not dividing n\}e\_v with compact SU mass one. The physical Hodge representation S+Q\^{}dual gives lambda\^{}r=K\_Y\^{}2 and A\^{}r=K\_X\^{}(2(g+1)); Mumford proportionality yields v\_A={[}2(g+1){]}\^{}d d! b\_pq\^{}f mu\_Gamma. An explicit integer upper bound is Vcal=n\^{}(fa)D\_F\^{}ceil(a/2)2\^{}(f(r-1)){[}2(g+1){]}\^{}d d!(product\_\{j=1\}\^{}\{r-1\}j!)\^{}f.~Thus D\_arith=2\^{}(8g\^{}2+2g+4)h\_d\^{}d Vcal sum\_\{j=0\}\^{}d k\^{}j, h\_d=2\^{}(d+1)(d+1)(d+2)!, retains the full same-rule cost, and N\_within\textless=t\^{}\{B\_\{t\^{}(4fr\^{}2),d-1\}(D\_arith)\}, t=1+3M Q\_S. The independent volume input has been replaced by actual field, rank, signature and level data. Any component-entry N0, same-family nonempty open, selected arithmetic nonzero rule and admitted Host/Sizes remain open. No evaluated actual N, six-identity mixed-CM chain, period, specific032 pairing, DP44 difference or RH conclusion. No finite error charge; current K115128/31 arrays and complete floors27/11 unchanged.

The calculation concerns the same smooth reductive unitary group scheme over the integer ring, obtained by finite-etale descent, its smooth congruence reductions, the actual Picard/Hodge bundle and the top intersection of its canonical extension on the normal arithmetic compactification. Exact fresh primary locators: Prasad(1989), Sections1.5 and3.3-3.7 pp96,109-111 and Remark3.11 pp112-113; Borel-Prasad(1989), Section4.2 p137; Mumford(1977), Theorem3.2 and proof pp262-263 and Proposition3.4(a,b) and proof pp264-265; Tits(1979), Section3.8.1 p55; Cameron, Section5.1 Proposition5.1 and proof p58; Smirnov(2015), Section2.9 Theorems2.31-2.32 pp18-19. Rapinchuk(2013), Theorem2.3 p281 and Milne(2017), Theorem4.16 p48 retain the accepted LIFT095 checks. The original OE action, determinant line, Picard data and Poincare/cubical-line transports are unchanged. Carried contextual anchors: FGA232 Theorem3.1 pp148-149; FGA236 Corollaries1.7/2.3 pp229/232; SGA7 I VII2.7/2.9 pp169-174 and VIII7.4(b) pp310-311; SGA3 I4.7.3/4.7.3.1; SGA6 VI Theorem3.3; cote4 scan003 sections8.1.1-8.1.2 and scans005-006 sections8.4.1-8.4.3, cote54 scans005-007. These are contextual uses, not newly reread Grothendieck primary proofs. No new lambda-ring/K-theory equality, l-adic weight, Dieudonne-crystal or CM-motive comparison follows. DP\_MULT Q3e/DP44/DP52/DP54-65 and THETA3 G1/A7 receive a conditional explicit arithmetic bound with the volume parameter eliminated; CX\_GRH CX-ALG-024 and GCT\_PLETH T-052 retain their arithmetic/global gates.

\subsection{295. AM-K-116}\label{am-k-116}

AM-K-116 RE1-28 retains the actual reciprocal-energy factor in a shifted positive half-Gamma convolution. The infinite total mass is explicit: every positive inverse moment is bounded by a convergent Gamma/Beta formula, and the unique polynomial degree0 term is made integrable by reserving theta=2\^{}-20 of the actual product-energy exponential before lattice enlargement. Every scalar/transverse term, original saturated metric and covolume, full167-entry reconstruction, central affine750 cover, all168 T factors and all paired four-state G patterns remain. The complete total-eligibility27,b=j\_int=0,2\^{}2111\textless x\textless2\^{}3328,k110 cell has full exponent-1539841/3072\textless-501. One501 append changes actual K115128/31 to129/31, retaining every old overlapping charge; both full sums remain below2\^{}-230 and complete floors27/11 stay. With R097/K115 the remaining actual support27 set is contained in the full-positive scalar-cross globally rank-two-unipotent all-II class with either total eligibility27,111\textless=k\textless=164,2\^{}2111\textless x\textless2\^{}3328 or total eligibility28,k164,2\^{}3107\textless x\textless2\^{}3328. The k111 positive sufficient envelope5269759/3072 is not a lower bound or impossibility theorem. No complete support28, signed nonvanishing, admitted source, selected arithmetic content, prescribed period, specific032 value, evaluated DP44 or RH conclusion.

The positive infinite measure is an analytic majorant of the same selected real realization. Its reserved damping is imposed on the actual original rows before enlargement; it introduces no arithmetic parameter, field extension, physical source vector, lattice rescaling or projector. All original projective OE modules, paired fractional ideals, CM involution, polarization, theta line, Picard and Poincare data remain. Fresh primary normalization: NIST DLMF equations5.9.1,5.12.1/5.12.3 and5.4.6, with their positive-real domains; exact URLs and separate root/peer reading scopes are preserved. The universal moment and zero-degree arguments are proved fully. Original representation and harmonic-code inputs retain their checked K107/K108/K110/K111/K113/K115 scopes. Carried exact anchors: SGA3 I4.7.3/4.7.3.1; SGA6 VI Theorem3.3; FGA232 Theorem3.1 pp148-149; FGA236 Corollaries1.7/2.3 pp229/232; SGA7 I VII2.7/2.9 pp169-174 and VIII7.4(b) pp310-311; cote4 scan003 sections8.1.1-8.1.2 and scans005-006 sections8.4.1-8.4.3; cote54 scans005-007. No new formal-group/Dieudonne, lambda-ring/K-theory, l-adic-weight or CM-motive comparison. DP\_MULT Q3e/DP44 and THETA3 G1/A7 receive the smaller conditional finite gate; CX\_GRH CX-ALG-024 and GCT\_PLETH T-052 retain their arithmetic/global gates.

\subsection{296. AM-LIFT-099}\label{am-lift-099}

AM-LIFT-099 RE1-26 supplies bounded component entry and effective same-rule escape for a supplied actual literal admitted Host/Datum/Level/Moduli/Coeff with A=1, fixed CM seed, relative-unramified E/F, ordinary selfdual IsDatW ideal multiple, and uniform indefinite rank r=p+q\textgreater=3. Let f={[}F:Q{]}, g=fr, d=fpq, n=3M, m=3M Q\_S, D\_E=D\_F\^{}2. Exact determinant congruence images and SU strong approximation identify the unitary finite double quotient with the norm-one determinant quotient; Cl\_E(m) surjects onto it. Half-trace duality and all-shape transference give a prime-to-m representative B of every ray class with N\_E B\textless=(2m)\^{}(2f)D\_E. The actual lattice B e+bar(B)\^{}(-1)f+Lprime, multiplied by the original AA and retaining J0, admits an original-level entry isogeny with N0=1 mod m and N0\textless=C\_ent=m(2m)\^{}(2f)D\_F\^{}2. All such bounded families cover the literal moduli complex points, so one meets the same supplied nonzero-rule open. With t=1+m, delta=t\^{}(4fr\^{}2), k=gw+n\_F/2 for even n\_F or k=2gw+n\_F after the same-zero square, Vcal=(3M)\^{}(f(r\^{}2-1))D\_F\^{}ceil((r\^{}2-1)/2)2\^{}(f(r-1)){[}2(g+1){]}\^{}d d!(product\_(j=1)\^{}(r-1)j!)\^{}f, h\_d=2\^{}(d+1)(d+1)(d+2)!, e\_g=2\^{}(8g\^{}2+2g+4), and D\_\emph{=e\_g h\_d\^{}d Vcal sum\_(j=0)\^{}d k\^{}j, the entire scale satisfies N=N0 t\^{}ell\textless=C\_ent t\^{}\{B\_delta,d-1(D\_})\}, B\_delta,0(D)=D and B\_delta,s(D)=D+B\_delta,s-1(2 delta\^{}D D\^{}2). The exact action, principal trace-one cubical line, finite eta and nonzero literal frame descend over the original K0(A0{[}N\^{}2{]}); generic relative degree\textless=N\^{}(8g\^{}2). Under precisely the original J,C0,Kgood same-seed hypotheses, {[}Kout:Bcoef{]}\textless64512 d\_CM M\^{}4 N\^{}(81/8) and rd(Kout)\textless2\^{}50 exp(2pi sqrt(d\_CM))rd(Bcoef) M N\^{}2, with the full composed N and actual coefficient-field degree/discriminant retained. No ray-class or proof-cover field is adjoined. Selected arithmetic-rule nonzero, actual Host/Sizes admission, remaining analytic/model/section/fibre Step5 inputs, Read/specific032, evaluated DP44 and RH remain open. Definite terminal signatures and rank below3 are excluded. No finite error charge; complete129/31 arrays and floors27/11 unchanged.

The new reduction is a finite ray-class quotient controlling integral projective Hermitian lattices in the same rational genus over Spec(O\_E), with exact congruence level and a bounded ordinary ideal norm. The chosen lattice becomes an actual principally polarized CM abelian variety over the original seed torsion field; its line-power identity is the normalized rigidified Poincare/cubical identity, and the nonzero rule is evaluated on an actual eigenframe torsor. New primary locators: Miller and Stephens-Davidowitz(2019),Theorem3.2/equation3.3 and proof,p8; Milne(2020),CFT II1.22-1.23 p67,III1.1 pp97-98,V4.6 pp172-173; Milne(2017),SVI Lemma8.20 p89. Strong approximation is carried Rapinchuk(2013),Theorem2.3 p281 and Milne(2017),Theorem4.16 p48 through LIFT095. Carried Grothendieck anchors: projective modules/exterior operations SGA3 I4.7.3/4.7.3.1 and SGA6 VI Theorem3.3; Picard/Poincare FGA232 Theorem3.1 pp148-149,FGA236 Corollaries1.7/2.3 pp229/232,SGA7 I VII2.7/2.9 pp169-174 and VIII7.4(b) pp310-311; formal crystals and Cartier duality cote4 scan003 sections8.1.1-8.1.2,scans005-006 sections8.4.1-8.4.3 and cote54 scans005-007. These are the existing contextual anchors, not new proofs of crystalline slopes, motivic realization comparisons, lambda-ring positivity or l-adic weights. DP\_MULT Q3e/DP44/DP52/DP54-65 and THETA3 G1/A7 receive an explicit conditional source-point cost, with the actual selected rule and six-identity/period matching still separate. CX\_GRH CX-ALG-024 and GCT\_PLETH T-052 retain their own global gates.

\subsection{297. AM-DO-121}\label{am-do-121}

AM-DO-121 EC1-24 retains the exact original unipotent identity cosh r\textgreater=1+alpha/2, hence sinh squared r\textgreater=alpha(1+alpha/4), in the genuine selected coefficient. Its positive infinite-mass Gamma167/rate4 convolution has every required convergent inverse moment; original product-energy reservation repairs the sole degree0 term before enlargement. The complete actual total-eligibility27/support27 all-II globally rank-two-unipotent scalar-cross union k111..139 has absolute normalized contribution below2\^{}-993. Every scalar/transverse polynomial term, original saturated metric and paired ideal, full167-entry reconstruction, unique eta, central affine750 cover, all168 finite factors and correlated four-state patterns is paid. One993 append changes actual K116129/31 to130/31, with both entire sums\textless2\^{}-230, all old overlapping errors retained and floors27/11 unchanged. The remaining necessary support27 alternatives are total eligibility27,140\textless=k\textless=164,2\^{}2111\textless x\textless2\^{}3328, or total eligibility28,k164,2\^{}3107\textless x\textless2\^{}3328; b=j\_int=0 throughout. The next k140 sufficient envelope767061/1024 is positive and is not a lower bound or impossibility theorem. No complete support28, signed nonvanishing, admitted source, selected arithmetic content, prescribed period, specific032 pairing value, evaluated DP44 saving or RH consequence.

The strengthened majorant applies to the same actual selected real realization and original unipotent endomorphism. The positive infinite measure is only an analytic integration device; it adds no arithmetic parameter, field, physical vector, projector or lattice rescaling. All original projective OE modules, CM involution, polarization, cubical theta line, Picard and Poincare data remain. The new identities and ordinary convergent inverse moments are proved in full, using NIST DLMF5.9.1,5.12.1/5.12.3 and5.4.6 with explicit positive arguments. Genuine vector-valued action/metric remains the established Nakahama2016 section5.1, equations5.1/5.3/5.5/5.7 and Theorem5.1 pp960-961; independent fresh versus carried source scopes are recorded. Inherited context only: SGA3 I4.7.3/4.7.3.1; SGA6 VI Theorem3.3; FGA232 Theorem3.1 pp148-149; FGA236 Corollaries1.7/2.3 pp229/232; SGA7 I VII2.7/2.9 pp169-174 and VIII7.4(b) pp310-311; cote4 scan003 sections8.1.1-8.1.2 and scans005-006 sections8.4.1-8.4.3; cote54 scans005-007. No new formal-group/Dieudonne, lambda-ring/K-theory, l-adic-weight or CM-motive comparison. DP\_MULT Q3e/DP44 and THETA3 G1/A7 receive the strictly smaller conditional finite gate; CX\_GRH CX-ALG-024 and GCT\_PLETH T-052 retain their arithmetic/global gates.

\subsection{298. AM-R-098}\label{am-r-098}

AM-R-098 WJ1-28 proves a new bound for the same actual selected trace with factor(cosh r/sinh r){[}2/(cosh r+1){]}\^{}(333/2)/sqrt30, retaining both compact factors and the full probability tail. The original unipotent energy yields shape166 with whole-field coefficient(2/15)\^{}84 and the single old near-point2. A positive clamp at delta=1/(1000R), imposed on actual rows before enlargement, pays less than125/104 and makes every required inverse moment through167 finite. The original un-clamped167th moment and total mass diverge; the highest moment includes both clamp atom and full upper integral, and the sole polynomial degree0 is repaired by original product-energy reservation. Every term of the all-positive-weight scalar/transverse polynomial is retained. This is an analytic interface only: it introduces no charge, source admission, central-fibre gain, high-degree finite-pattern shortcut, signed nonvanishing, selected arithmetic-image detection, specific032 value, evaluated DP44 or RH result. Actual130/31 error arrays and floors27/11 remain unchanged.

This is an estimate at the archimedean realization of the same projective OE modules and polarized CM abelian data. The positive measure, analytic weights and clamp do not alter the scheme over its original arithmetic base, the involution, Hermitian lattice, theta line, Picard scheme or Poincare biextension. Carried contextual anchors are SGA3 I4.7.3/4.7.3.1; SGA6 VI Theorem3.3; FGA232 Theorem3.1 pp148--149; FGA236 Corollaries1.7/2.3 pp229/232; SGA7 I VII2.7/2.9 pp169--174 and VIII7.4(b) pp310--311; cote4 scan003 sections8.1.1--8.1.2 and scans005--006 sections8.4.1--8.4.3; cote54 scans005--007. These locators are inherited context, not new primary readings or newly proved comparisons. No new formal-group/Dieudonne, lambda-ring/K-theory, l-adic weight, CM-motive or motivic L-function result follows from this real analytic inequality alone. DP\_MULT Q3e/DP44, THETA3 G1/A7, CX\_GRH CX-ALG-024 and GCT\_PLETH T-052 retain their current arithmetic and global gates until a separately checked composition changes them. No Lean build or archive theorem-validity audit is performed.

\subsection{299. AM-K-117}\label{am-k-117}

AM-K-117 D29.1-27 proves the universal degree29 scalar Gaussian lattice bound from the fixed exact normalized Gegenbauer polynomial and, with AM-R-098 WJ1-28, removes the complete actual support27 total-eligibility28 branch k164,b=j\_int0,2\^{}3107\textless x\textless2\^{}3328 with absolute normalized contribution below2\^{}-278. The entire660-term polynomial includes the repaired zero degree and the unique degree167 term with its full Q\^{}167 pattern cost. Every original167-entry reconstruction, paired ideal, selected metric, central affine750 cover, unique eta, support and all finite factors is paid; the one eligible inactive place keeps contraction1. One278 append preserves all130/31 predecessor entries, giving131/31 with both entire sums below2\^{}-230 and floors27/11 unchanged. Only the eligibility27 alternative140\textless=k\textless=164,2\^{}2111\textless x\textless2\^{}3328,b=j\_int0 remains at actual support27. No complete support28, signed nonvanishing, actual admitted source, selected arithmetic nonzero rule, prescribed period, evaluated DP44 cost or RH conclusion.

The fixed harmonic positive-kernel polynomial controls the same real scalar dual lattice after arbitrary positive diagonal rescaling; no transformed lattice is asserted to be an arithmetic ideal. The actual projective OE modules, half-trace pairing, CM involution, polarized theta line, Picard and Poincare data stay unchanged. Delsarte, Goethals, and Seidel (1977), Definition2.1, Theorem2.4, Theorems3.2-3.3, Corollary3.8 and Theorem4.3, printed364-368, provide the exact primary normalization context. NIST DLMF5.9.1,5.12.1/3,5.4.6 are used only with positive real integral arguments, including the full clamped degree167 atom and tail. Inherited context only: SGA3 I4.7.3/4.7.3.1; SGA6 VI Theorem3.3; FGA232 Theorem3.1 pp148-149; FGA236 Corollaries1.7/2.3 pp229/232; SGA7 I VII2.7/2.9 pp169-174 and VIII7.4(b) pp310-311; cote4 scan003 sections8.1.1-8.1.2 and scans005-006 sections8.4.1-8.4.3; cote54 scans005-007. No new formal-group/Dieudonne, lambda-ring/K-theory, l-adic-weight or CM-motive comparison. DP\_MULT Q3e/DP44 and THETA3 G1/A7 receive the strictly smaller conditional finite gate; CX\_GRH CX-ALG-024 and GCT\_PLETH T-052 retain their arithmetic/global gates.

\subsection{300. AM-DO-122}\label{am-do-122}

AM-DO-122 CE1-29 removes the complete actual support27 total-eligibility27 all-II globally rank-two-unipotent scalar-cross k140 sector with absolute contribution below2\^{}-359. MANDATORY CUTOFF QUALIFICATION: the prose after CE8 is read as norm(h\_tau)\textless=H=133, as proved by R048-RC RC48.4/RC48.6 and the paragraph after RC48.7; CE8 already has the correct non-strict sign. The frozen proof is read together with recovery\_pending/definite/DO122\_CUTOFF\_SOURCE\_ADDENDUM.md, whose full text is included in the checked paper. The joint physical central variable w=-U Im d, original affine ideal, full750 cover, weighted Voronoi integral and U\^{}f physical Jacobian are retained. The central integral consumes one profile power, giving shape166; all moments used and every positive polynomial term are paid, with original-row zero repair, all finite factors and correlated patterns. One359 append to actual K117131/31 gives132/31; both entire sums remain below2\^{}-230 and floors27/11 stay unchanged. The sole support27 necessary alternative is now total eligibility27,141\textless=k\textless=164,2\^{}2111\textless x\textless2\^{}3328,b=j\_int0. The positive k141 envelope is only a failed sufficient bound. No full support28, signed nonvanishing, actual admitted source, selected arithmetic content, prescribed period, evaluated DP44 or RH claim.

The physical coordinate change is a real analytic treatment of the same original central affine module, preserving its shift, CM partner conditions and polarization. It adds no arithmetic parameter or field and changes no selected metric. NIST DLMF5.9.1,5.12.1/3 and5.4.6 are applied with explicit positive integral arguments. Inherited context only: SGA3 I4.7.3/4.7.3.1; SGA6 VI Theorem3.3; FGA232 Theorem3.1 pp148-149; FGA236 Corollaries1.7/2.3 pp229/232; SGA7 I VII2.7/2.9 pp169-174 and VIII7.4(b) pp310-311; cote4 scan003 sections8.1.1-8.1.2 and scans005-006 sections8.4.1-8.4.3; cote54 scans005-007. No new formal-group/Dieudonne, lambda-ring/K-theory, l-adic-weight or CM-motive comparison. DP\_MULT Q3e/DP44 and THETA3 G1/A7 receive the strictly smaller conditional finite gate; CX\_GRH CX-ALG-024 and GCT\_PLETH T-052 retain their arithmetic/global gates.

\subsection{301. AM-LIFT-100}\label{am-lift-100}

AM-LIFT-100 HN1-19 proves the literal ThNe0 law given an actual algebraic FrPt over a finite number-field extensionK of the coefficient field, the original symmetric principal rigidified line and odd ideal half-level, literal IsNab with its prescribed conjugation, and literal ThValTot. For every unchanged nonzero P in harmA of total degreeD=DF with D+2\textless=M, some individual permitted n-tuple produces a nonzero actual thGen rule. The original ideal b satisfies b bar(b)=(M); C=A{[}b{]} has orderM\^{}g, H=A{[}b\^{}2{]} has orderM\^{}(2g), and the actual eta identifies La/bLa withC. Unique normalized odd lifts of H/C form exactlyM\^{}g permitted sections spanning the pullback of H0(A/C,L\_B\^{}M), with pi\^{}*L\_B=L\^{}M and L\_B principal. Bauer-Szemberg interpolation, the original common IsNab chart and ordinary highest symbols prove detection, retaining the literal M\^{}(-D) and factorials. The construction occurs on a reduced test-ring presentation ofK(A{[}M\^{}2{]}); no additional scalar-root field is required. At mostM\^{}(gn) tuples suffice existentially; no particular tuple or its norm/content is evaluated. Dom, n\textgreater0 and r\textless=n remain as in ThNe0; cusp-valuedness uses a separately supplied ThInCusp. This is a conditional named source-law constructor, not nonvanishing of the restricted original coupled selected image. Actual Moduli/FrPt existence, IsNab, ThValTot, other LawsP/LawsSh, full Host/Sizes, prescribed period, specific032, actual g\_j,e\_j,H and evaluatedDP44 remain open. No charge; current132/31 arrays and complete floors27/11 unchanged.

This is a characteristic-zero generic-fibre statement for the original polarized CM abelian scheme and its projective OE lattice; no extension over all of Spec(Z), integral content or new reduction-at11 theorem is asserted. The finite etale ideal subgroup C gives an actual quotient isogeny and descent of the cubical theta line, with exact Picard identity pi\^{}*L\_B=L\^{}M. The deck-character decomposition is in the finite-dimensional section representation; it does not replace the physical selected vector. Primary exact locators: Mumford(1966),section1(II)-(V) p289,Proposition1 and degree computation p291,Theorem1 pp293-294,Proposition3 pp295-296,Theorem2/invariant-section proof p297; Bauer and Szemberg(1997),Theorem1 PDFp2 and Corollary2.2 PDFp6,with full Theorem2.1 proof PDFpp3-5 and complex-ground-field convention. Actual IsVtOver32-41,ThNe030-32,DF28,IsNab208-217,IsNormMul103-107 and SymLift44-51 are pinned in the exact source record. Carried contextual anchors: SGA3 I4.7.3/4.7.3.1; SGA6 VI Theorem3.3; FGA232 Theorem3.1 pp148-149; FGA236 Corollaries1.7/2.3 pp229/232; SGA7 I VII2.7/2.9 pp169-174 and VIII7.4(b) pp310-311; cote4 scan003 sections8.1.1-8.1.2/scans005-006 sections8.4.1-8.4.3 and cote54 scans005-007. Formal derivatives retain the actual formal group and gauge, but no new Dieudonne slope or Cartier-dual divisibility follows. Section dimensions and exact line identities are compatible with the existing K-theory/exterior-power framework; no lambda-positivity, l-adic weight or new CM-motive L-function comparison is proved. The auxiliary quotient is isogenous to the original CM variety, without proving the required arithmetic oscillator or period matching. DP\_MULT Q3e/DP44/DP52/DP54-65 and THETA3 G1/A7 receive this named conditional source-law construction; CX\_GRH CX-ALG-024 and GCT\_PLETH T-052 retain their global gates.

\subsection{302. AM-K-118}\label{am-k-118}

AM-K-118 BT1-22 proves that the entire actual all-top selected probability outsideI={[}1/35,1/31{]} is below2\^{}-40000000, uniformly in every original word and twist. Both high moments are in their true convergence domain. Fidelity restriction costs at most2\^{}9 times its square root over all168 embeddings, including radius0 and both compact factors. The full bounded original all-II globally rank-two-unipotent scalar-cross support27 family with nonzeroS cross and x\textless2\^{}3328 has complete nonarchimedean normalized mass below2\^{}11394877; original167 inputs, full affine central coset, supports, labels, G patterns and source/Haar factors are counted. The entire discarded arithmetic contribution is therefore below2\^{}-8605114. The retained-band main term remains to be estimated and this input appends no standalone charge or cell removal. Current132/31 arrays and complete floors27/11 remain unchanged; no unbounded-denominator, support28, signed nonvanishing, source admission, selected-image content, period, specific032, evaluatedDP44 or RH conclusion.

This is an estimate at the archimedean realization of the same projective OE modules and polarized CM abelian data. The positive measure and simultaneous radial restriction do not alter the scheme over its original arithmetic base, the involution, Hermitian lattice, theta line, Picard scheme or Poincare biextension. Carried contextual anchors are SGA3 I4.7.3/4.7.3.1; SGA6 VI Theorem3.3; FGA232 Theorem3.1 pp148--149; FGA236 Corollaries1.7/2.3 pp229/232; SGA7 I VII2.7/2.9 pp169--174 and VIII7.4(b) pp310--311; cote4 scan003 sections8.1.1--8.1.2 and scans005--006 sections8.4.1--8.4.3; cote54 scans005--007. These locators are inherited context, not new primary readings or newly proved comparisons. No new formal-group/Dieudonne, lambda-ring/K-theory, l-adic weight, CM-motive or motivic L-function result follows from this real analytic inequality alone. DP\_MULT Q3e/DP44, THETA3 G1/A7, CX\_GRH CX-ALG-024 and GCT\_PLETH T-052 retain their current arithmetic and global gates until a separately checked composition changes them. No Lean build or archive theorem-validity audit is performed.

\subsection{303. AM-R-099}\label{am-r-099}

AM-R-099 AH1-25 derives an extra exp{[}-(111/50)(cosh r-6/5){]} from the same actual selected angular kernel restricted toI={[}1/35,1/31{]} at both radial arguments. The old and weighted Jacobians give separate alternative bounds, retaining the exact detailed-balance cancellation and whole radial spread. The original unipotent Cartan identity yields joint J\_lambda(t,w)=t\^{}(-1/2)(1+t/4)\^{}(-167)(1+w\^{}2/{[}4(1+t/4)\^{}2{]})\^{}(-333/4)exp(-111t/100), with the original physical w=-U Im d.~The non-strict original norm\textless=133 cutoff gives a single full near-point2, and the entire intercept costs less than2\^{}108. Original central integration may consume exactly one radial power, leaving shape166; its arithmetic application remains separate. The weighted energy-only alternative retains its own(2/15)\^{}84 coefficient and is not multiplied into the joint bound. The exact shifted measure has infinite mass but finite every positive integer inverse moment; the zero-degree term needs original-row damping before enlargement. The actual trace is bounded only after adding the complete K118 tail2\^{}-8605114 in its precise bounded support27 scope. No standalone error charge or coordinate removal; current132/31 arrays and complete floors27/11 unchanged.

This is an estimate at the archimedean realization of the same projective OE modules and polarized CM abelian data. The positive shifted measure and simultaneous radial restriction do not alter the scheme over its original arithmetic base, the involution, Hermitian lattice, theta line, Picard scheme or Poincare biextension. Carried contextual anchors are SGA3 I4.7.3/4.7.3.1; SGA6 VI Theorem3.3; FGA232 Theorem3.1 pp148--149; FGA236 Corollaries1.7/2.3 pp229/232; SGA7 I VII2.7/2.9 pp169--174 and VIII7.4(b) pp310--311; cote4 scan003 sections8.1.1--8.1.2 and scans005--006 sections8.4.1--8.4.3; cote54 scans005--007. These locators are inherited context, not new primary readings or newly proved comparisons. No new formal-group/Dieudonne, lambda-ring/K-theory, l-adic weight, CM-motive or motivic L-function result follows from this real analytic inequality alone. DP\_MULT Q3e/DP44, THETA3 G1/A7, CX\_GRH CX-ALG-024 and GCT\_PLETH T-052 retain their current arithmetic and global gates until a separately checked composition changes them. No Lean build or archive theorem-validity audit is performed.

\subsection{304. AM-K-119}\label{am-k-119}

AM-K-119 RM1-19 proves the global positive majorant t\^{}(-1/2)(1+t/4)\^{}(-166)exp(-111t/100)\textless=t\^{}(-1/2)(1+t/4)\^{}(-166)(1+t/400)\^{}(-444). Its Gamma166/rate4 plus Gamma444/rate400 plus half-Gamma convolution has infinite mass and inverse moments exactly for1/2\textless j\textless610+1/2. Every required integer degree1..167 is evaluated as pi times an explicit positive rational, with all full coefficients and strict product168 bounds recorded. Finite principal parts plus integration-by-parts recurrence agree exactly with the independent73704-cell rectangular-integral/binomial calculation at every degree. The complete scalar/transverse row polynomial retains all terms and repairs its sole zero degree by original-row energy reservation before enlargement. Actual selected-kernel, intercept, near-point, centralU scaling, all finite/source/class/Haar/support/pattern and complete band-tail costs remain separate. This analytic input appends no charge and removes no sector; current132/31 arrays and complete floors27/11 unchanged.

This is an estimate at the archimedean realization of the same projective OE modules and polarized CM abelian data. The positive Gamma convolution and exact inverse-moment evaluation do not alter the scheme over its original arithmetic base, the involution, Hermitian lattice, theta line, Picard scheme or Poincare biextension. Carried contextual anchors are SGA3 I4.7.3/4.7.3.1; SGA6 VI Theorem3.3; FGA232 Theorem3.1 pp148--149; FGA236 Corollaries1.7/2.3 pp229/232; SGA7 I VII2.7/2.9 pp169--174 and VIII7.4(b) pp310--311; cote4 scan003 sections8.1.1--8.1.2 and scans005--006 sections8.4.1--8.4.3; cote54 scans005--007. These locators are inherited context, not new primary readings or newly proved comparisons. No new formal-group/Dieudonne, lambda-ring/K-theory, l-adic weight, CM-motive or motivic L-function result follows from this real analytic inequality alone. DP\_MULT Q3e/DP44, THETA3 G1/A7, CX\_GRH CX-ALG-024 and GCT\_PLETH T-052 retain their current arithmetic and global gates until a separately checked composition changes them. No Lean build or archive theorem-validity audit is performed.

\subsection{305. AM-R-100}\label{am-r-100}

AM-R-100 NB1-16 proves the actual selected fidelity restricted to I*={[}16/529,16/527{]} has the same original joint physical-central kernel with lambda32/25. The universal angular loss coefficient527/33858 exceeds1/65 and both conditional-density branches, zero height, full radial spread and detailed balance are retained. Original non-strict norm\textless=133 yields precisely one full near-point2; the entire intercept exp(128f/13175) costs less than2\^{}3. Central integration may remove one radial power to shape166 while retaining its full physical U\^{}f covariance and original finite costs. Passage from restricted fidelity to actual selected traces explicitly requires the new narrow probability NB15; under it K118 band-independent fidelity/mass reasoning gives the full additive tail2\^{}-8605114. No wide-band probability implication, K120 moment premise, standalone charge or sector removal is inferred. Existing132/31 arrays and complete floors27/11 remain unchanged.

This is an estimate at the archimedean realization of the same projective OE modules and polarized CM abelian data. The positive shifted measure and simultaneous radial restriction do not alter the scheme over its original arithmetic base, the involution, Hermitian lattice, theta line, Picard scheme or Poincare biextension. Carried contextual anchors are SGA3 I4.7.3/4.7.3.1; SGA6 VI Theorem3.3; FGA232 Theorem3.1 pp148--149; FGA236 Corollaries1.7/2.3 pp229/232; SGA7 I VII2.7/2.9 pp169--174 and VIII7.4(b) pp310--311; cote4 scan003 sections8.1.1--8.1.2 and scans005--006 sections8.4.1--8.4.3; cote54 scans005--007. These locators are inherited context, not new primary readings or newly proved comparisons. No new formal-group/Dieudonne, lambda-ring/K-theory, l-adic weight, CM-motive or motivic L-function result follows from this real analytic inequality alone. DP\_MULT Q3e/DP44, THETA3 G1/A7, CX\_GRH CX-ALG-024 and GCT\_PLETH T-052 retain their current arithmetic and global gates until a separately checked composition changes them. No Lean build or archive theorem-validity audit is performed.

\subsection{306. AM-DO-123}\label{am-do-123}

AM-DO-123 BC1-27 removes the entire original bounded support27/total-eligibility27 all-II full-positive scalar-cross globally rank-two-unipotent k141..142 union, with b=j\_int0 and2\^{}2111\textless x\textless2\^{}3328. The same R099 selected angular kernel, original physical central coordinate w=-U Im d, full affine covolume/U scaling/750 cover, shape166/shift111/100 moments and actual-row zero-degree repair retain every polynomial, finite, label, support, pivot, class/Haar and pattern factor. Full exponents-3344299/1024 and-1688491/1024 plus the entire K118 tail2\^{}-8605114 yield one1647 charge. Appending it once to actual DO122132/31 gives133/31 with both entire sums below2\^{}-230 and complete floors27/11 unchanged. The sole remaining support27 necessary alternative is eligibility27,143\textless=k\textless=164,2\^{}2111\textless x\textless2\^{}3328,b=j\_int0. The k143 envelope-28587/1024 is insufficient and unclaimed. No full support28, signed nonvanishing, source admission, selected arithmetic-image content, prescribed period, specific032 value, evaluatedDP44 or RH conclusion.

The physical coordinate change is a real analytic treatment of the same original central affine module, preserving its shift, CM partner conditions and polarization. It adds no arithmetic parameter or field and changes no selected metric. NIST DLMF5.9.1,5.12.1/3 and5.4.6 are applied with explicit positive integral arguments. Inherited context only: SGA3 I4.7.3/4.7.3.1; SGA6 VI Theorem3.3; FGA232 Theorem3.1 pp148-149; FGA236 Corollaries1.7/2.3 pp229/232; SGA7 I VII2.7/2.9 pp169-174 and VIII7.4(b) pp310-311; cote4 scan003 sections8.1.1-8.1.2 and scans005-006 sections8.4.1-8.4.3; cote54 scans005-007. No new formal-group/Dieudonne, lambda-ring/K-theory, l-adic-weight or CM-motive comparison. DP\_MULT Q3e/DP44 and THETA3 G1/A7 receive the strictly smaller conditional finite gate; CX\_GRH CX-ALG-024 and GCT\_PLETH T-052 retain their arithmetic/global gates.

\subsection{307. AM-K-120}\label{am-k-120}

AM-K-120 NB1-24 proves the actual all-word/twist mean2111/64\textless M1\textless2113/64 and full probability outside I\emph{={[}16/529,16/527{]} below2\^{}-40000000, using h=2\^{}35 in the true negative-moment chamber. Both event complements and all168 embeddings, multiplied by the complete original bounded mass2\^{}11394877, give the full additive arithmetic tail2\^{}-8605114. Separately the shifted shape166 lambda32/25 profile is dominated by the positive Gamma166/rate4, Gamma512/rate400 and half-Gamma convolution. All167 exact rational inverse moments pi}r\_j are evaluated by two independent rational integral methods; the real chamber is1/2\textless j\textless678.5. Infinite mass and the original-row first-moment repair for the sole zero degree are retained, with every scalar/transverse polynomial term and true primal covolume. This theorem supplies neither the same-kernel angular implication nor the complete sector count; all intercept, central, finite and high-degree pattern costs remain separately payable. No standalone charge or main-gate change; current133/31 arrays and complete floors27/11 unchanged.

This is an estimate at the archimedean realization of the same projective OE modules and polarized CM abelian data. The positive Gamma convolution and exact inverse-moment evaluation do not alter the scheme over its original arithmetic base, the involution, Hermitian lattice, theta line, Picard scheme or Poincare biextension. Carried contextual anchors are SGA3 I4.7.3/4.7.3.1; SGA6 VI Theorem3.3; FGA232 Theorem3.1 pp148--149; FGA236 Corollaries1.7/2.3 pp229/232; SGA7 I VII2.7/2.9 pp169--174 and VIII7.4(b) pp310--311; cote4 scan003 sections8.1.1--8.1.2 and scans005--006 sections8.4.1--8.4.3; cote54 scans005--007. These locators are inherited context, not new primary readings or newly proved comparisons. No new formal-group/Dieudonne, lambda-ring/K-theory, l-adic weight, CM-motive or motivic L-function result follows from this real analytic inequality alone. DP\_MULT Q3e/DP44, THETA3 G1/A7, CX\_GRH CX-ALG-024 and GCT\_PLETH T-052 retain their current arithmetic and global gates until a separately checked composition changes them. No Lean build or archive theorem-validity audit is performed.

\subsection{308. AM-LIFT-101}\label{am-lift-101}

AM-LIFT-101 AT1-25 proves two precise conditional constructors. Given actual Host, Datum, odd Level, Coeff and Moduli records, every literal IsNab yields all-ring ThValTot, actual FrMap naturality and harmonic independence. With the accepted CanonGaugeSt theorem, those same records permit a choice of NabV satisfying IsNab, including NabZero at all integer powers. The proof supplies coherent formal charts and frames over arbitrary coefficient algebras, global primitive/adjunction normalization, finite-order corrected additive frames, ordinary derivatives, exact harmonic cancellation and faithfully-flat descent of the original odd ideal-half-level theta sections. Where the operator is free to be chosen, independent nabPin, nabUniq and thTot premises are removed. Combined with AM-LIFT-100 and an actual FrPt over a finite number-field extension, unchanged nonzero harmonic P and DF+2\textless=M, the literal ThNe0 law follows without separate chosen-IsNab/ThValTot assumptions. This does not identify a separately prescribed operator, construct the actual arithmetic source or Host/Sizes, prove the other source laws, or establish nonvanishing of the original coupled selected image. No error charge; complete133/31 arrays and floors27/11 unchanged.

The construction lives over arbitrary commutative coefficient algebras of characteristic zero, including nonreduced rings. Proper SAB records give abelian schemes; universal degree-zero cohomology normalizes global line automorphisms, formal smoothness supplies compatible jets, and nilpotent lifting supplies line generators. The cubical line and its biextension induce an additive corrected frame; the OE eigenspace decomposition and trace-one CM weight give the exact quadratic harmonic cancellation. Odd normalized theta lifts and finite etale ideal-half root torsors yield actual sections and flat descent, without principalizing the ideal or identifying global units away from the chosen point. Primary locators: Stacks Lemmas36.32.5-6(tag0BDM),Definition37.11.1/Lemma37.11.7(tag02GZ),Lemma30.5.2(tag02KH); Mumford(1966),section1II-V p289,Proposition1 pp290-291,Proposition2 pp291-292,Theorem1 pp293-294,Proposition3 pp295-296,Theorem2 p297. Exact Alpoge literal interfaces and scoped reads are recorded in the root source record and mandatory locator addendum. Carried contextual locators: SGA7 I VII2.7/2.9 pp169-174 and VIII7.4(b) pp310-311; FGA232 Theorem3.1 pp148-149 and FGA236 Corollaries1.7/2.3 pp229/232; cote4 scan003 sections8.1.1-8.1.2/scans005-006 sections8.4.1-8.4.3 and cote54 scans005-007; SGA3 I4.7.3/4.7.3.1 and SGA6 VI Theorem3.3. These are geometric/source-law results, with no new integral Dieudonne slope, lambda-positivity, l-adic weight theorem or CM-motive L-function comparison. DP\_MULT Q3e/DP44/DP52/DP54-65 and THETA3 G1/A7 gain the stated source-law reduction; CX\_GRH CX-ALG-024 and GCT\_PLETH T-052 keep their arithmetic gates.

\subsection{309. AM-DO-124}\label{am-do-124}

AM-DO-124 CB1-27 removes the entire original bounded support27/total-eligibility27 all-II full-positive scalar-cross globally rank-two-unipotent k143 sector, with b=j\_int0 and2\^{}2111\textless x\textless2\^{}3328. The same accepted R100 narrow-band angular kernel and K120 actual narrow probability/moments retain the original physical central coordinate w=-U Im d, complete affine covolume/U scaling/750 cover, shape166/lambda32/25 profile, all167 inverse moments and actual-row zero-degree repair. Every original polynomial, finite, label, support, pivot, class/Haar and pattern factor is paid. The full exponent-725931/1024\textless-708 plus the entire new narrow-band tail2\^{}-8605114 gives one707 charge. Appending exactly once to actual DO123133/31 gives134/31 with both entire sums below2\^{}-230 and complete floors27/11 unchanged. The sole remaining support27 necessary alternative is total eligibility27,144\textless=k\textless=164,2\^{}2111\textless x\textless2\^{}3328,b=j\_int0. The k144 sufficient envelope924757/1024\textgreater0 is explicitly unclaimed. No full support28, signed nonvanishing, source admission, selected arithmetic-image content, prescribed period, specific032 value, evaluatedDP44 or RH conclusion.

The physical coordinate change is a real analytic treatment of the same original central affine module, preserving its shift, CM partner conditions and polarization. It adds no arithmetic parameter or field and changes no selected metric. NIST DLMF5.9.1,5.12.1/3 and5.4.6 are applied with explicit positive integral arguments. Inherited context only: SGA3 I4.7.3/4.7.3.1; SGA6 VI Theorem3.3; FGA232 Theorem3.1 pp148-149; FGA236 Corollaries1.7/2.3 pp229/232; SGA7 I VII2.7/2.9 pp169-174 and VIII7.4(b) pp310-311; cote4 scan003 sections8.1.1-8.1.2 and scans005-006 sections8.4.1-8.4.3; cote54 scans005-007. No new formal-group/Dieudonne, lambda-ring/K-theory, l-adic-weight or CM-motive comparison. DP\_MULT Q3e/DP44 and THETA3 G1/A7 receive the strictly smaller conditional finite gate; CX\_GRH CX-ALG-024 and GCT\_PLETH T-052 retain their arithmetic/global gates.

\subsection{310. AM-LIFT-102}\label{am-lift-102}

AM-LIFT-102 TheoremsA-E/EC1-16 proves the prescribed-frame converse and the actual-package harmonic comparison. Given the actual geometric records, every coherent IsExpChart/IsExpGauge is realized by the original Poincare biframe with h0, so its nonnegative ordinary Taylor operator satisfies literal IsNab with the exact product factorial. The normalized Appell-Humbert reduced frame at every actual complex point gives such data, for mixed signatures as well as pure types, and admits a pinned all-ring extension. The entire ShortValAt and class relations are unchanged between two pinned operators on every harmK input. In an actually supplied OMPins/OneMapTieP package, CovPin already pins the operative DeltaDat.nab2 and Pins.io supplies range(io)=span(HHa) contained in harmK. Thus no additional equality of individual ordinary coefficients is needed to compare its original thLaValAt, IsThClsD and OneMapRow relation in the reduced gauge. All original witnesses, roots, theta-group sections, beta/index/cardinality/N-degree normalizations, coefficient eF map, vol/gam/ep/cdelta weights, phases and selected maps remain fixed. The package itself, absent short witnesses and sec/cdelta laws, mixed integral coefficient descent and original coupled selected-image nonvanishing remain unconstructed. Unknown unpinned operators still require their own witness. No error charge; complete134/31 arrays and floors27/11 unchanged.

The construction is over arbitrary characteristic-zero coefficient algebras and retains the given integral CM action, polarization, cubical maps, global adjunctions and primitive. A balanced quadratic change is realized by twisting the same formal Poincare biframe; it changes no abelian scheme, Picard point, source field, level or chosen arithmetic model. The ordinary factorial and the harmonic contraction are exact. Bannai and Kobayashi(2010), Algebraic theta functions and the p-adic interpolation of Eisenstein-Kronecker numbers, Duke Mathematical Journal153(2),229-295, doi10.1215/00127094-2010-024; author arXiv math/0610163v4 section1.2, Theorem1.7/Definition1.8, printedpp9-10, supplies the freshly read Appell-Humbert classification/multiplier statement. Universal global functions and formal smoothness are the checked LIFT101 inputs, Stacks Lemmas36.32.5-6(tag0BDM), Definition37.11.1/Lemma37.11.7(tag02GZ), Lemma30.5.2(tag02KH). Carried exact geometric context: SGA7 I, VII2.7/2.9 pp169-174 and VIII7.4(b) pp310-311; FGA232 Theorem3.1 pp148-149 and FGA236 Corollaries1.7/2.3 pp229/232; SGA3 I4.7.3/4.7.3.1; SGA6 VI Theorem3.3; cote4 scan003 sections8.1.1-8.1.2 and scans005-006 sections8.4.1-8.4.3; cote54 scans005-007. These contextual locators are carried from the documented earlier readings, not claimed as new whole-text checks. No new integral extension over Z, Dieudonne slope or Cartier-duality formula, lambda-ring positivity, K-theory filtration, l-adic weight theorem or CM-motive L-function comparison is obtained. DP\_MULT Q3e/DP44/DP52/DP54-65 and THETA3 G1/A7 gain only the stated source comparison; CX\_GRH CX-ALG-024 and GCT\_PLETH T-052 retain their global arithmetic gates.

\subsection{311. AM-K-121}\label{am-k-121}

AM-K-121 KP1-34 proves the exact complete coarea kernel and all-c paired-height bound for the original J56 selected density restricted to16/529\textless=Y\textless=16/527. Both real branches and both directions retain one common Beta expectation bounded by23479/23450; the paired factor is(yyprime)\^{}1/4*c\^{}-1/2, and the same kernel has height{[}1-(8/513)(c-1){]}\_+\^{}(333/2), supported at c\textless521/8. The original near-point and physical central affine lattice are transported by explicit additive cutoff broadenings2\^{}-340 and2\^{}-74, retaining the complete U\^{}f covolume, affine shift and750 cover. The full central integral has a genuine positive containing u-mixture after a globally paid exponential tangent; the compact profile itself is not asserted to be a positive Laplace transform. Every positive inverse moment is finite; the infinite zero degree still requires the actual-row reserve before enlargement. The same selected-spread coefficient is refined using every actual word/twist and S\_epsilon\textless97/(100B), with the entire old lower tail absorbed and the separate complete band-complement2\^{}-8605114 tail retained. The conservative coarea and spread field ceilings are-423 and-15, respectively; they are analytic factors and supply no sector charge alone. A complete original polynomial/moment/pattern/endpoint count remains required. No error charge; complete134/31 arrays, floors27/11 and k144..164 finite gate remain unchanged.

The theorem is a real-place estimate on the unchanged selected theta coefficient and its original arithmetic affine fibre. It retains the same polarized projective order modules, CM involution, saturated lattices, cubical/Poincare data, finite projectors, phases, Haar measures and arithmetic periods. NIST DLMF5.9.1 and5.12.1/3 are used with positive real Gamma/Beta arguments; Nakahama(2016), Norm computation and analytic continuation of vector valued holomorphic discrete series representations, Journal of Lie Theory26(4),927-990, equation5.7 and Theorem5.1, printed960-961, remains the actual Hilbert-metric input through K043/K046/DO080. Carried exact geometric context: SGA7 I, VII2.7/2.9 pp169-174 and VIII7.4(b) pp310-311; FGA232 Theorem3.1 pp148-149 and FGA236 Corollaries1.7/2.3 pp229/232; SGA3 I4.7.3/4.7.3.1; SGA6 VI Theorem3.3; cote4 scan003 sections8.1.1-8.1.2 and scans005-006 sections8.4.1-8.4.3; cote54 scans005-007. These contextual locators are carried from the documented earlier readings, not claimed as new whole-text checks. No scheme over Z or Picard/motive is reconstructed, and no new Dieudonne slope, Cartier-duality formula, lambda-ring/K-theory positivity, l-adic weight or CM-motive L-function theorem follows. DP\_MULT Q3e/DP44, THETA3 G1/A7, CX\_GRH CX-ALG-024 and GCT\_PLETH T-052 retain the same actual source, arithmetic image, period and global comparison gates.

\subsection{312. AM-DO-125}\label{am-do-125}

AM-DO-125 CM1--28 removes the complete original bounded support27/total-eligibility27 all-II full-positive scalar-cross globally rank-two-unipotent k144 cell, b=j\_int0 and2\^{}2111\textless x\textless2\^{}3328. Accepted K121 same-kernel coarea/height/variance is globally bounded by a paid exponential tangent before transport, then the exact physical central integral is bounded by a positive Gamma166/rate4, half-Gamma/rate2, half-Gamma/rate3514/29, infinite half-Gamma and shift29/20 convolution. Both complete quadrature tails and8064 intervals for degree147 yield a directed rational moment enclosure; finite-endpoint Hölder covers every lower positive degree and the actual-row reserve repairs degree0. All579 corrections, original supports/pivots/Green/T factors, full750 cover/U scaling and four finite small inflations remain. The full exponent is below−7622655/32768\textless−232.5, and the entire2\^{}-8605114 tail is paid once, giving one232 charge. The exact135/31 full error sums stay below2\^{}-230 with floors27/11. The sole remaining support27 necessary alternative is145\textless=k\textless=164 with the same denominator interval; the positive k145 envelope removes nothing. Signed PP3, higher-support remainder, actual source and selected arithmetic image, prescribed period, specific032 value, actual lift data/evaluatedDP44 and Hecke RH remain open.

The exact positive profile treats the same original central affine module, real shift, CM partner incidence, polarization and selected Hilbert metric. The full central Jacobian and original projective row covolume are retained. Ordinary positive Gamma/Beta formulas are NIST DLMF5.9.1 and5.12.1/3; the Beta--Gamma/Jensen reduction and continuous exact quadrature are proved in CM11--20. Inherited context only: SGA3 I4.7.3/4.7.3.1; SGA6 VI Theorem3.3; FGA232 Theorem3.1 pp148-149; FGA236 Corollaries1.7/2.3 pp229/232; SGA7 I VII2.7/2.9 pp169-174 and VIII7.4(b) pp310-311; cote4 scan003 sections8.1.1-8.1.2 and scans005-006 sections8.4.1-8.4.3; cote54 scans005-007. No new formal-group/Dieudonne, lambda-ring/K-theory, l-adic-weight or CM-motive comparison. DP\_MULT Q3e/DP44 and THETA3 G1/A7 receive the strictly smaller conditional finite gate; CX\_GRH CX-ALG-024 and GCT\_PLETH T-052 retain their arithmetic/global gates.

\subsection{313. AM-LIFT-103}\label{am-lift-103}

AM-LIFT-103 ST1--19 constructs the entire first all-ring existence/uniqueness conjunct of literal ShortValTot from the actual geometric records, Lambda.Unimod/TrivAt(M)/LevOK(N,M), an operative IsNab pin and P in identity harmK. The canonical finite-etale dual-row kernel has exactly the source zCard, and every GeoInj witness is a full enumeration. Allowed half-primary root differences lie in that subgroup by TrivAt(M); unique normalized odd theta lifts give the exact complete section-trace identity with the original U\_t U\_z order. Harmonic contraction, a finite-etale construction cover, flat H0/eigenbasis base change and faithfully flat element descent produce the unique value over every coefficient algebra, retaining nilpotents and all original beta/mIdx/cardinality/N-degree/weight factors. Every original faithfully flat witness agrees. Arbitrary-map naturality compares supplied actual FrMap points and does not assume FrPt.gen under nonflat pullback. Applying the theorem to the actual cover-stage data removes the short-value scalar constructor from the conditional source gate. The second original-base sec/cdelta/IsShortRule conjunct remains open, as do source/package existence, selected arithmetic content/nonzero, cusp/spanning laws, period, specific032 value and evaluatedDP44/RH. No numerical error charge;135/31 complete arrays and floors27/11 unchanged.

The kernel is a canonical finite-etale subgroup of the original abelian scheme, defined by integral dual-row endomorphisms. It uses the actual projective CM module, alternating adjoint form, principal polarization and cubical line; no global Tate basis or changed field/source is introduced. The original theta-group trace is an equality of sections and its scalar descends faithfully flatly over arbitrary characteristic-zero coefficient rings. Exact primary applications: Milne(2008), Abelian varieties Version2.00, ChapterI Proposition7.1/Theorem7.2, printedpp32--33; Stacks Lemma59.64.4(tag03RV), Lemma41.18.3(tag04HN), Lemma35.3.6(tag023M), Lemma59.16.3(tag03O9), Lemma30.5.2(tag02KH). Relative and all-witness applications are proved in ST3--18. Carried exact geometric context: SGA7 I, VII2.7/2.9 pp169-174 and VIII7.4(b) pp310-311; FGA232 Theorem3.1 pp148-149 and FGA236 Corollaries1.7/2.3 pp229/232; SGA3 I4.7.3/4.7.3.1; SGA6 VI Theorem3.3; cote4 scan003 sections8.1.1-8.1.2 and scans005-006 sections8.4.1-8.4.3; cote54 scans005-007. These contextual locators are carried from the documented earlier readings, not claimed as new whole-text checks. No new integral extension over Z, Dieudonne slope or Cartier-duality formula, lambda-ring positivity, K-theory filtration, l-adic weight theorem or CM-motive L-function comparison is obtained. DP\_MULT Q3e/DP44/DP52/DP54-65 and THETA3 G1/A7 gain only the stated source comparison; CX\_GRH CX-ALG-024 and GCT\_PLETH T-052 retain their global arithmetic gates.

\subsection{314. AM-LIFT-104}\label{am-lift-104}

AM-LIFT-104 CV1--21 constructs the exact cover-stage Rules/sec(a,n) object from LIFT103 for the original HHa polynomial and operative pin. Infinite-field RowInv identities extend to every coefficient algebra; commuting invariant derivations and the literal physical frame change give the positive row character. The original theta generator contributes t\_none\^{}n, with every full-trace normalization retained. On FromTest the actual evaluator equals the unique scalar for every presentation. Actual IsLev nonzero evR(cdelta) supplies FromTest at every required algebraically closed cover frame, so H2=cdelta*G2 belongs to cover sec(a,n) and has precisely the required field evaluations. IsShortRule is equivalent to original-base descent of this particular H2 under all original FrOver quantifiers. The actual CovPin/LawsFull suppliers force both auxiliary levels A=A2=1, removing that potential mismatch in those scopes. Original-base descent itself, source/package existence, integral content, cusp/spanning, selected arithmetic nonzero, period, specific032 value and evaluatedDP44/RH remain open. No numerical charge;135/31 complete arrays and floors27/11 unchanged.

The construction uses the original CM abelian scheme, projective OE lattice, rigidified cubical line and actual physical tangent eigenframes. Formal invariant differentiation and identity-matrix harmonic contraction commute with the stated coefficient maps, retaining nilpotents. The cover rule is a natural function on the literal reduced finite-type frame tests with exactly the source Borel character; its evaluation boundary is derived from the source FromTest branch. The original-base section descent remains explicit. Exact Alpoge interfaces: RowDeg.lean33--48; IsCombAt.lean67--127; EmbKCoe.lean13--15; \nolinkurl{ThGerm.lean56–74/189–217;} \nolinkurl{CocycleRightZeroThd.lean25–36/166–190;} AlgebraMapCAt.lean9--98; \nolinkurl{ThShort.lean155–181/263–293;} ThClassPinsRec.lean62--63; FullExists.lean40; \nolinkurl{TwoWeightEqOneAddRoot.lean53–66.} The non-elementary geometric inputs are accepted LIFT101--103 at their documented primary scopes. Carried exact context: SGA7 I VII2.7/2.9 pp169--174 and VIII7.4(b) pp310--311; FGA232 Theorem3.1 pp148--149 and FGA236 Corollaries1.7/2.3 pp229/232; SGA3 I4.7.3/4.7.3.1; SGA6 VI Theorem3.3; cote4 scan003 sections8.1.1--8.1.2 and scans005--006 sections8.4.1--8.4.3; cote54 scans005--007. These are carried contextual readings, not new whole-text checks. No new integral extension over Z, Dieudonne slope, lambda-ring positivity, K-theory filtration, l-adic weight theorem or CM-motive L-function comparison follows. DP\_MULT Q3e/DP44/DP52/DP54--65 and THETA3 G1/A7 gain this conditional source-rule reduction; CX\_GRH CX-ALG-024 and GCT\_PLETH T-052 retain their global gates.

\subsection{315. AM-K-122}\label{am-k-122}

AM-K-122 HC1--18, with the mandatory explicit nonzero-original-S-cross scope addendum, bounds the actual selected trace by its simultaneous narrow-band high-coherence integral plus the complete band error and2\^{}(-B/4) locally. Algebraic evaluation has plus-N circle weight; its genuine metric adjoint range lies in the selected high-weight spectral subspace. Two-projector exponential factorization, exact two-singular-value calculation and actual all-word/twist weights give overlap below2\^{}(-B/4) whenever projective angular coherence is at most9/10. Insertion into the original finite trace before absolute values retains both actual fibre metrics, isometric transport and every compact factor. The full original bounded support27 all-II globally rank-two-unipotent scalar-cross family with nonzero original S-cross and x\textless2\^{}3328 has paid global errors2\^{}-8605114 and2\^{}(11394885-B/4), the latter below2\^{}-40000000. The actual high-coherence main integral remains unbounded. This analytic input adds no sector charge, removes no k145..164 cell and changes no complete support floor;135/31 arrays and floors27/11 remain. Source admission, selected arithmetic image/content/nonzero, signedPP3, period, specific032 and evaluatedDP44/RH remain open.

This is a real-place spectral/evaluation estimate on the unchanged selected theta coefficient, retaining the original polarized projective order modules, CM involution, cubical/Poincare data, finite projectors, phases, Haar measures and arithmetic periods. The range of the true evaluation adjoint and the actual invariant Hermitian fibre metric determine the angular cut. Nakahama(2016), Norm computation and analytic continuation of vector valued holomorphic discrete series representations, Journal of Lie Theory26(4),927--990, equations5.5/5.7 and Theorem5.1, printed959--961, is freshly read in the retained published primary with the exact K043 PB2--8 convention; DO080 RT1--8/RT22 supplies the source-normalized trace and corrected word formula. K118 BT5/BT15--22 and K120 NB1--16 supply the precise bounded nonzero-S-cross family and full arithmetic mass/tail. Carried exact geometric context: SGA7 I, VII2.7/2.9 pp169-174 and VIII7.4(b) pp310-311; FGA232 Theorem3.1 pp148-149 and FGA236 Corollaries1.7/2.3 pp229/232; SGA3 I4.7.3/4.7.3.1; SGA6 VI Theorem3.3; cote4 scan003 sections8.1.1-8.1.2 and scans005-006 sections8.4.1-8.4.3; cote54 scans005-007. These contextual locators are carried from the documented earlier readings, not claimed as new whole-text checks. No scheme over Z or Picard/motive is reconstructed, and no new Dieudonne slope, Cartier-duality formula, lambda-ring/K-theory positivity, l-adic weight or CM-motive L-function theorem follows. DP\_MULT Q3e/DP44, THETA3 G1/A7, CX\_GRH CX-ALG-024 and GCT\_PLETH T-052 retain the same actual source, arithmetic image, period and global comparison gates.

\subsection{316. AM-DO-126}\label{am-do-126}

AM-DO-126 IQ1--36 removes the complete original bounded support27/total-eligibility27 all-II full-positive scalar-cross globally rank-two-unipotent union145\textless=k\textless=159, b=j\_int0 and2\^{}2111\textless x\textless2\^{}3328. The saturated original isotropic quotient is a projective selfdual O\_E-module, including at2/U; every actual row lies in J I\_Sigma D\_eta times its ordinary row dual. Its standard S164 summand and selfdual rank3 complement give exact transverse covolume D0\^{}3 z\^{}(6f), while the independent physical central U\^{}-f remains. A genuine real-dimension1008 Gaussian bound retains both transverse terms. The unchanged positive central mixture, exact continuous moments148/162, finite-endpoint Holder, repaired zero term, all4590 positive terms, all original finite/source factors and the entire band tail yield one582 charge. Worst directed term is−19510525/32768\textless−595; the entire union is\textless2\^{}-582. Exact136/31 full error sums remain\textless2\^{}-230 with floors27/11. The sole remaining support27 necessary alternative is160\textless=k\textless=164 with the same denominator interval. The positive k160 sufficient envelope removes nothing. Signed PP3, higher-support remainder, actual source and selected arithmetic image, prescribed period, specific032 value, actual lift data/evaluatedDP44 and Hecke RH remain open.

The new object is the saturated isotropic quotient (ell-perpendicular intersect Lambda)/(ell intersect Lambda) over Spec(O\_E), with its descended perfect Hermitian pairing and nonprincipal projective module. Integral local plane splitting uses determinant−1 and no division by2. Its positive metric and exact codifferent covolume retain the original polarization, CM involution, cubical line, Picard/Poincare data and all finite phases. The actual arithmetic arrow is only placed in a necessary containing module. Inherited context only: SGA3 I4.7.3/4.7.3.1; SGA6 VI Theorem3.3; FGA232 Theorem3.1 pp148--149; FGA236 Corollaries1.7/2.3 pp229/232; SGA7 I VII2.7/2.9 pp169--174 and VIII7.4(b) pp310--311; cote4 scan003 sections8.1.1--8.1.2 and scans005--006 sections8.4.1--8.4.3; cote54 scans005--007. No new Dieudonne, lambda-ring/K-theory, l-adic-weight or CM-motive comparison. DP\_MULT Q3e/DP44 and THETA3 G1/A7 gain the strictly smaller conditional finite gate; CX\_GRH CX-ALG-024 and GCT\_PLETH T-052 retain their arithmetic/global gates.

\subsection{317. AM-LIFT-105}\label{am-lift-105}

AM-LIFT-105 SU1-22 proves the literal ShortThetaUnit for the standard auxiliary lattice O\^{}n and any actual operative IsNab on the original geometric data. The same unit is the inverse of the actual selected beta0, independent of N, x0, polynomial and every original ring/witness. Global normalized multiplication/primitive identities give same-point ordinary derivative scaling without added truncation coherence. The complete standard A{[}N{]}\^{}n deck trace cancels exactly mIdx in the source action order, every allowed N-root gives its allowed M-root, and RowDeg cancels the full N/M degree factor. Harmonic all-ring comparison and faithful-flat injectivity retain nilpotents and every original witness. Nonstandard lattice beta laws, original-base descent, source/package, cusp/spanning, selected arithmetic nonzero/content, period, specific032 value and evaluatedDP44/RH remain separate. No numerical charge;136/31 complete arrays and floors27/11 unchanged.

The new identity is formed on the actual rigidified Picard/Poincare line and cubical biextension over the original characteristic-zero moduli base. The original projective CM lattice is retained; only the separate auxiliary standard lattice is used. Normalized multiplication and odd lifts are universal line/group-scheme identities before the full harmonic contraction. Exact Alpoge interfaces: \nolinkurl{ThShort.lean20-32,37-106,244-261,313-323;} \nolinkurl{IsThRule.lean32-43,57-80,85-88;} RowDeg.lean40-48; ThGerm.lean108-221; recovered \nolinkurl{ThetaPt.lean44-70,73-95,152-161} and ThetaVV.lean103-108. Primary normalization input: Stacks36.32.5-6, tag0BDM. Accepted LIFT101 AT3,AT8-21, LIFT102 TheoremD and LIFT103 ST5-17 retain their complete prior scope. Carried exact context: SGA7 I VII2.7/2.9 pp169--174 and VIII7.4(b) pp310--311; FGA232 Theorem3.1 pp148--149 and FGA236 Corollaries1.7/2.3 pp229/232; SGA3 I4.7.3/4.7.3.1; SGA6 VI Theorem3.3; cote4 scan003 sections8.1.1--8.1.2 and scans005--006 sections8.4.1--8.4.3; cote54 scans005--007. These are carried contextual readings, not new whole-text checks. No new integral extension over Z, Dieudonne slope, lambda-ring positivity, K-theory filtration, l-adic weight theorem or CM-motive L-function comparison follows. DP\_MULT Q3e/DP44/DP52/DP54--65 and THETA3 G1/A7 gain this conditional source-rule reduction; CX\_GRH CX-ALG-024 and GCT\_PLETH T-052 retain their global gates.

\subsection{318. AM-R-101}\label{am-r-101}

AM-R-101 FD1-23 constructs both literal ShortValTot conjuncts for every original admissible lattice, polynomial and frame witness when actual IsLev, new A=1, operative IsNab and PhiPin/PsiTie are supplied. The principal theta section module is an invertible line with arbitrary base change; the exact physical frame scheme represents the literal Rules and field evaluator. Half-level rigidity and reduced-base Moduli.rep construct a canonical finite-etale half-level torsor independently of stored d.r. Its exact Sim/LevAut action gives the positive psi Fourier character, inverse psi full trace and cancelling PsiTie scalar. The full harmonic cover rule descends effectively to an original-base sec(a,n), and every original FrOver witness is identified with its canonical framed pullback. No global d.r equality or circular full LawsSh premise is assumed. Actual CovPin supplies the first three hypotheses; dPhi supplies the character package only for n\textgreater=r. The separate small-n supplier, source/package existence, cusp/spanning, selected arithmetic nonzero/content, period, specific032 value and evaluatedDP44/RH remain open here. No numerical charge;136/31 arrays and floors27/11 unchanged.

This is effective descent of a Borel-equivariant regular function on the genuine physical frame scheme over the original characteristic-zero moduli base. The Picard theta section line commutes with arbitrary base change, while its canonical finite-etale half-level torsor is built from the original projective OE-module and rigidified cubical/Poincare line. The Tate and inverse-adjoint data preserve the original polarization. Exact primary locators: Milne Abelian varieties Theorem13.3(a), Proposition17.5(a,b) and Lemma17.6, printed25/37; Stacks30.17.1 (01XO),36.30.4 (0B91),36.31.3 (0BDK),36.32.5-6 (0BDM),59.64.4 (03RV),41.18.3 (04HN),35.3.6 (023M),59.16.3 (03O9). Exact Alpoge interfaces and root fresh/carried extents are the eighteen-file canonical\_half\_level\_short\_rule\_descent\_sources.json; in particular ThShort DeltaDat.hF/EtaOK/FrOver/IsLev/PhiPin/PsiTie/ShortValTot, reduced-base PushInvertible.Moduli.rep, HeckeKer.Sim/InGam/LevAut, and actual CubeMap.IsPullCubical. Carried exact context: SGA7 I VII2.7/2.9 pp169--174 and VIII7.4(b) pp310--311; FGA232 Theorem3.1 pp148--149 and FGA236 Corollaries1.7/2.3 pp229/232; SGA3 I4.7.3/4.7.3.1; SGA6 VI Theorem3.3; cote4 scan003 sections8.1.1--8.1.2 and scans005--006 sections8.4.1--8.4.3; cote54 scans005--007. These are carried contextual readings, not new whole-text checks. No new integral extension over Z, Dieudonne slope, lambda-ring positivity, K-theory filtration, l-adic weight theorem or CM-motive L-function comparison follows. DP\_MULT Q3e/DP44/DP52/DP54--65 and THETA3 G1/A7 gain this conditional source-rule reduction; CX\_GRH CX-ALG-024 and GCT\_PLETH T-052 retain their global gates.

\subsection{319. AM-DO-127}\label{am-do-127}

AM-DO-127 GC1-16 removes five complete containing subsets of the original support27/eligibility27 full-positive scalar-cross globally rank-two-unipotent all-II branch, with b=j\_int0: k160/161/162/163/164 and log2x\textgreater=2120/2130/2140/2150/2165 respectively, inside the original x\textless2\^{}3328 window. The exact normalized finite Green formula gives2 Q\^{}(-335*167)(NEJ)\^{}167 NFJ0, whose NFJ0 cancels the independent physical central divisor. The complete pattern is product(1+q\_v\^{}(j-167))\^{}2, including exactly2\^{}668 at degree167. The actual projective selfdual quotient, full1008 transverse theta, physical U\^{}-f, same positive profile, all1630 terms and one entire band tail remain. Worst directed term -13968857/32768\textless-426 gives the full union\textless2\^{}-415. Exactly one415 append gives137/31 full sums\textless2\^{}-230 and unchanged floors27/11. The sole support27 necessary branch has2\^{}2111\textless x\textless2\^{}(X\_k) with the five respective upper exponents. Low-x cells, signed PP3, higher supports, actual source/selected arithmetic image, period, specific032 value, evaluatedDP44 and Hecke RH remain open.

The refined finite formula is the character of the original cuspidal finite representation at its unipotent reduction, correlated with the actual row ideal and central congruence ideal. It is applied to the same saturated isotropic quotient over Spec(O\_E), projective Hermitian dual and physical real central coset. No lattice, polarization, cubical/Poincare line, source arrow, compact factor or selected CM phase is replaced. Primary formula: Gorodetsky and Hazan(2017), arXiv1707.07308v2, Theorem2.1/equation2.1 and Corollary2.2, printed5, retaining Green attribution. Inherited context only: SGA3 I4.7.3/4.7.3.1; SGA6 VI Theorem3.3; FGA232 Theorem3.1 pp148--149; FGA236 Corollaries1.7/2.3 pp229/232; SGA7 I VII2.7/2.9 pp169--174 and VIII7.4(b) pp310--311; cote4 scan003 sections8.1.1--8.1.2 and scans005--006 sections8.4.1--8.4.3; cote54 scans005--007. No new Dieudonne, lambda-ring/K-theory, l-adic-weight or CM-motive comparison. DP\_MULT Q3e/DP44 and THETA3 G1/A7 gain the strictly smaller conditional finite gate; CX\_GRH CX-ALG-024 and GCT\_PLETH T-052 retain their arithmetic/global gates.

\subsection{320. AM-LIFT-106}\label{am-lift-106}

AM-LIFT-106 AR1-18 proves that actual AdmX forces n\textgreater=p+q for both Boolean branches on the original finite projective OE lattice. Below that rank shortGen is empty and both original ShortValTot conjuncts are vacuous for every already selected DeltaDat. Actual PinDat gives p+q=R.r, including the terminal definite datum. Above that rank the actual dPhi range supplies exactly the character hypotheses of accepted R101; unbundled dLev/dNab/newA1/PinDat therefore fill all-n shTot with the original data unchanged. Carrying the other eight supplied fields completes LawsSh. Cusp and uniqueness below r are vacuous; upper-range cusp/spanning, selected arithmetic nonzero/content, source/package existence, period,032 and evaluatedDP44/RH remain separate. No new charge;137/31 arrays and27/11 floors unchanged.

Over Spec OE the physical CM lattice is a vector bundle of constant rank p+q. The proper original half-level ideal has a residue point, and generation of its restriction requires at least that many sections. The false-Boolean branch is an actual basis of the restricted vector bundle. This finite-flat rank statement closes the missing lower-column range of the characteristic-zero framed theta descent without changing any Picard/Poincare or class-character data. Primary locator: Stacks Lemma10.78.2 (00NX). Exact source interfaces: TwoWeightEqOneAddRoot14-62; ThShort93,107-109,271-298,322-333; PartitionExtendByZeroAntitone12,24-49; Params32,34,46; DepRun87-92; HasDataCondition57-73; FullExists40-47. Carried exact context: SGA7 I VII2.7/2.9 pp169--174 and VIII7.4(b) pp310--311; FGA232 Theorem3.1 pp148--149 and FGA236 Corollaries1.7/2.3 pp229/232; SGA3 I4.7.3/4.7.3.1; SGA6 VI Theorem3.3; cote4 scan003 sections8.1.1--8.1.2 and scans005--006 sections8.4.1--8.4.3; cote54 scans005--007. These are carried contextual readings, not new whole-text checks. No new integral extension over Z, Dieudonne slope, lambda-ring positivity, K-theory filtration, l-adic weight theorem or CM-motive L-function comparison follows. DP\_MULT Q3e/DP44/DP52/DP54--65 and THETA3 G1/A7 gain this conditional source-rule reduction; CX\_GRH CX-ALG-024 and GCT\_PLETH T-052 retain their global gates.

\subsection{321. AM-R-102}\label{am-r-102}

AM-R-102 CB1-8 applies accepted GlobalTheory.BetaEx.betaExists directly to every original T.Q.Dw(j), T.Q.c, n and badN(D), filling LawsSh.beta. KPsiIn remains the internal antecedent; the original selected betaH has its successful-choice specification only when that antecedent is supplied. Together with accepted LIFT106, the seven other supplied clauses and separate PinDat/newA1 assemble LawsSh without any DeltaDat change. Accepted kPsiIn\_iff/reflexRoots\_iff/reflexLe\_of\_basis identify the exact same-field character supplier; the actual class/root data are not constructed. The accepted coefficient constructor supplies an eventual extension-field law, not the prescribed eO/PinField identity. No duplicate determinant theorem or archive audit is credited. No charge;137/31 arrays and27/11 floors unchanged; actual source, selected nonzero/content, cusp/spanning, period,032 and evaluatedDP44/RH remain separate.

This is applicability of the accepted determinant-index/type-norm trivialization over the original number-ring schemes, followed by exact assembly of the existing characteristic-zero theta source record. The prescribed integer-ring isomorphism and complex embedding are retained. Accepted source commit7cbf7a0c32fca1badf05931879a37c3f8a2b03b4: ExistsIndexIdeal.betaExists322-324; MultAddMultConjugate.kPsiIn\_iff269-272,reflexLe\_of\_basis447-458,reflexRoots\_iff460-474,typings\_agree557-589; ReflexNormIndexIdeal64-72; ExistsCoeffTypeNormsPrincipal23-25. Literal consumers: ThShort74-85,277-285; HasDataCondition55,57-73; PartitionExtendByZeroAntitone20,45; FullExists40. Carried exact context: SGA7 I VII2.7/2.9 pp169--174 and VIII7.4(b) pp310--311; FGA232 Theorem3.1 pp148--149 and FGA236 Corollaries1.7/2.3 pp229/232; SGA3 I4.7.3/4.7.3.1; SGA6 VI Theorem3.3; cote4 scan003 sections8.1.1--8.1.2 and scans005--006 sections8.4.1--8.4.3; cote54 scans005--007. These are carried contextual readings, not new whole-text checks. No new integral extension over Z, Dieudonne slope, lambda-ring positivity, K-theory filtration, l-adic weight theorem or CM-motive L-function comparison follows. DP\_MULT Q3e/DP44/DP52/DP54--65 and THETA3 G1/A7 gain this conditional source-rule reduction; CX\_GRH CX-ALG-024 and GCT\_PLETH T-052 retain their global gates.

\subsection{322. AM-K-123}\label{am-k-123}

AM-K-123 HCX1-26 removes the complete original bounded support27/eligibility27 full-positive scalar-cross globally rank-two-unipotent all-II k160 cell, b=j\_int0 and2\^{}2111\textless x\textless2\^{}3328. The actual all-word selected evaluation gives a chi\textgreater937/1000 restriction with both full arithmetic errors paid. A525772-test outward interval cover proves the complete physical-coherence cap; exact three-real-dimensional Beta disintegration gives a fraction below2\^{}-10 with a paid margin. The same-kernel positive four-component central mixture retains the original affine U scaling, full750 cover, Young regularization and extra transport2. Nine full quadratures, finite Holder anchors, actual-row zero repair, the true1008-dimensional transverse theta, exact Green/J0 correlation, all322 terms and both full tails give one600 charge. Exact138/31 full sums remain below2\^{}-230 and floors27/11 remain. The sole support27 necessary alternative is161\textless=k\textless=164,2\^{}2111\textless x, with upper exponents2130/2140/2150/2165 respectively. Signed PP3, higher supports, actual source/selected arithmetic image/content/nonzero, prescribed period, specific032 value, evaluatedDP44 and RH remain open.

The new estimate is on the original real-place selected theta coefficient and its actual homogeneous Hermitian bundle. Simultaneous compact conjugation and conditional sphere disintegration preserve the original polarized projective OE modules, selected CM phase, finite projector, cubical/Poincare line, Haar/class normalization and prescribed arithmetic arrow. The same projective selfdual isotropic quotient over Spec(OE) and original affine central ideal remain in the count. The positive real Laplace mixture is an analytic majorant, not a new algebraic theta correspondence. Primary normalization: Nakahama(2016), equations5.5/5.7 and Theorem5.1, printed959-961; exact finite Green specialization: Gorodetsky and Hazan(2017), Theorem2.1/equation2.1 and Corollary2.2, printed5, through the accepted inputs. Inherited context only: SGA3 I4.7.3/4.7.3.1; SGA6 VI Theorem3.3; FGA232 Theorem3.1 pp148--149; FGA236 Corollaries1.7/2.3 pp229/232; SGA7 I VII2.7/2.9 pp169--174 and VIII7.4(b) pp310--311; cote4 scan003 sections8.1.1--8.1.2 and scans005--006 sections8.4.1--8.4.3; cote54 scans005--007. No new Dieudonne, lambda-ring/K-theory, l-adic-weight or CM-motive comparison. DP\_MULT Q3e/DP44 and THETA3 G1/A7 gain the strictly smaller conditional finite gate; CX\_GRH CX-ALG-024 and GCT\_PLETH T-052 retain their arithmetic/global gates.

\subsection{323. AM-LIFT-107}\label{am-lift-107}

AM-LIFT-107 AG1-19 proves that, for the original Par/PinDat and coefficient field F, every signature KPsiIn is equivalent to the two actual stages j1/j2 and to one normalized rank-one CM-type generator law. Explicit adjacent Bezout extraction preserves F and every embedded absolute norm, with no relation imposed between independently chosen witness phases. A finite normalized class-generator table is necessary and sufficient; a fixed principal-power root table is sufficient only, retaining its unit-phase distinction. Accepted R102 then supplies the unchanged beta implication. Actual table existence, original source/Host/Sizes and selected arithmetic nonzero/content, period,032 and evaluatedDP44/RH remain open. No new charge; current138/31 arrays and27/11 floors unchanged. Frozen137/31 author context is historical.

Over Spec(O\_F), the two signature-weighted invertible ideals determine the CM half-type ideal by an explicit integral Bezout combination and the rational norm line. The finite table is a normalized trivialization of these ideal classes with all prescribed archimedean norms, not ordinary principalization alone. The original polarized CM projective lattices, Poincare/cubical lines, finite phases and coefficient-field pin remain unchanged. Exact literal locators: Params6-22/32/46; PartitionExtendByZeroAntitone12/24-31/47; Datum.mult in FlagBasis76; Coeff in PushInvertible40-46; EmbKCoe25-33; ThShort277-285; HasDataCondition55; accepted MultAddMultConjugate262-272/460-474; MapRingEquivTransApply52-69. Primary ideal/norm/class/unit locators are Sutherland2021 Lectures3/4/6/14/15 as listed in the proof. Carried exact context: SGA7 I VII2.7/2.9 pp169--174 and VIII7.4(b) pp310--311; FGA232 Theorem3.1 pp148--149 and FGA236 Corollaries1.7/2.3 pp229/232; SGA3 I4.7.3/4.7.3.1; SGA6 VI Theorem3.3; cote4 scan003 sections8.1.1--8.1.2 and scans005--006 sections8.4.1--8.4.3; cote54 scans005--007. These are carried contextual readings, not new whole-text checks. No new integral extension over Z, Dieudonne slope, lambda-ring positivity, K-theory filtration, l-adic weight theorem or CM-motive L-function comparison follows. DP\_MULT Q3e/DP44/DP52/DP54--65 and THETA3 G1/A7 gain this conditional source-rule reduction; CX\_GRH CX-ALG-024 and GCT\_PLETH T-052 retain their global gates.

\subsection{324. AM-DO-128}\label{am-do-128}

AM-DO-128 HZ1-24 gives the exact signed full-Heisenberg zero mode and additive Fourier coefficients, qualitative paired CRT completion preserving every original T incidence, full original-kernel positive definiteness and a convergent full-group nonnegative sum, and exact physical-frequency covering bounds. Identity repetition, cutoff, coset and actual selected-transform limitations are explicit. This structural theorem alone makes no numerical main-gate reduction or new charge. Current K123138/31 arrays, absolute floors27/11 and support27 k161-164 remain unchanged; frozen137/31 and k160 statements are historical. Original source/Host/Sizes, selected arithmetic image/nonzero, period,032 and DP44/RH remain open.

The arithmetic group completion retains the original paired Hermitian projective lattice, finite flags, induced CM field, positive definite quotient and finite projector lines. The real central trace dual is the inverse different times the inverse fractional ideal; its covolume and covering radius have distinct roles. Full positive-definite matrix coefficients and their arithmetic sums do not supply positivity of a selected support mask or a global arithmetic theta cycle. Exact internal locators are K062 AT1-25; DO106 CFI1-24; DO108 AC8-15; DO126 IQ1-36; DO127 GC1-16; DO107 TD20-22, together with Gorodetsky and Hazan(2017), Theorem2.1/equation2.1/Corollary2.2, printed5. Carried exact context: SGA7 I VII2.7/2.9 pp169--174 and VIII7.4(b) pp310--311; FGA232 Theorem3.1 pp148--149 and FGA236 Corollaries1.7/2.3 pp229/232; SGA3 I4.7.3/4.7.3.1; SGA6 VI Theorem3.3; cote4 scan003 sections8.1.1--8.1.2 and scans005--006 sections8.4.1--8.4.3; cote54 scans005--007. These are carried contextual readings, not new whole-text checks. No new integral extension over Z, Dieudonne slope, lambda-ring positivity, K-theory filtration, l-adic weight theorem or CM-motive L-function comparison follows. DP\_MULT Q3e/DP44/DP52/DP54--65 and THETA3 G1/A7 gain this conditional source-rule reduction; CX\_GRH CX-ALG-024 and GCT\_PLETH T-052 retain their global gates.

\subsection{325. AM-LIFT-108}\label{am-lift-108}

AM-LIFT-108 NI1-21 identifies the actual induced quadratic type ideal with the extension of N\_E/K(I), for K=Q(H.rt) embedded as specified in the original F. The complete original kPsi law is equivalent to a finite normalized table on G=image Cl(E)-\textgreater Cl(K), with arbitrary integral representatives of generator classes. Every nonnegative supplied signature is covered in the forward direction. Small quadratic representatives have norm≤(2/pi)sqrt(d), scalar height≤half log of that bound, and explicit class-power roots supply a new coefficient record of relative degree≤\textbar G\textbar≤h\_K, persistent over further upper fields. The original-field table remains unevaluated; the new field does not transport eO/PinField or changed DP44 costs. The geometric d7 class-number-one specialization is excluded by the full source d\textgreater365 and cannot relabel J14/J56 Fields.Kq. No actual proper-image claim or new charge; current138/31 and27/11 floors remain.

On Spec(O\_K), relative ideal norm pushes invertible ideal classes from E to the finite Picard group of the specified imaginary quadratic field; extension along its original embedding gives the exact CM-type line over Spec(O\_F). The table is a trivialization with every archimedean modulus, not principalization alone. The new coefficient field is an explicit root tower with bounded relative degree and normalized scalar height, while its arithmetic discriminant and prescribed field identity remain separate. Exact source locators are TwoWeightEqOneAddRoot14-26/30-49; FlagBasis76; PushInvertible40-46; EmbKCoe25-33; ThShort277-285; HasDataCondition55; Partition24-49; MapRing52-69; DataHost39-54 and Params6-22/32. Primary norm/class/Minkowski locators are Sutherland2021 Lectures3/4/6/14 and Milne2020 ANT Example2.41, as stated in the proof. Carried exact context: SGA7 I VII2.7/2.9 pp169--174 and VIII7.4(b) pp310--311; FGA232 Theorem3.1 pp148--149 and FGA236 Corollaries1.7/2.3 pp229/232; SGA3 I4.7.3/4.7.3.1; SGA6 VI Theorem3.3; cote4 scan003 sections8.1.1--8.1.2 and scans005--006 sections8.4.1--8.4.3; cote54 scans005--007. These are carried contextual readings, not new whole-text checks. No new integral extension over Z, Dieudonne slope, lambda-ring positivity, K-theory filtration, l-adic weight theorem or CM-motive L-function comparison follows. DP\_MULT Q3e/DP44/DP52/DP54--65 and THETA3 G1/A7 gain this conditional source-rule reduction; CX\_GRH CX-ALG-024 and GCT\_PLETH T-052 retain their global gates.

\end{document}
